% generated by GAPDoc2LaTeX from XML source (Frank Luebeck)
\documentclass[a4paper,11pt]{report}

\usepackage[top=37mm,bottom=37mm,left=27mm,right=27mm]{geometry}
\sloppy
\pagestyle{myheadings}
\usepackage{amssymb}
\usepackage[latin1]{inputenc}
\usepackage{makeidx}
\makeindex
\usepackage{color}
\definecolor{FireBrick}{rgb}{0.5812,0.0074,0.0083}
\definecolor{RoyalBlue}{rgb}{0.0236,0.0894,0.6179}
\definecolor{RoyalGreen}{rgb}{0.0236,0.6179,0.0894}
\definecolor{RoyalRed}{rgb}{0.6179,0.0236,0.0894}
\definecolor{LightBlue}{rgb}{0.8544,0.9511,1.0000}
\definecolor{Black}{rgb}{0.0,0.0,0.0}

\definecolor{linkColor}{rgb}{0.0,0.0,0.554}
\definecolor{citeColor}{rgb}{0.0,0.0,0.554}
\definecolor{fileColor}{rgb}{0.0,0.0,0.554}
\definecolor{urlColor}{rgb}{0.0,0.0,0.554}
\definecolor{promptColor}{rgb}{0.0,0.0,0.589}
\definecolor{brkpromptColor}{rgb}{0.589,0.0,0.0}
\definecolor{gapinputColor}{rgb}{0.589,0.0,0.0}
\definecolor{gapoutputColor}{rgb}{0.0,0.0,0.0}

%%  for a long time these were red and blue by default,
%%  now black, but keep variables to overwrite
\definecolor{FuncColor}{rgb}{0.0,0.0,0.0}
%% strange name because of pdflatex bug:
\definecolor{Chapter }{rgb}{0.0,0.0,0.0}
\definecolor{DarkOlive}{rgb}{0.1047,0.2412,0.0064}


\usepackage{fancyvrb}

\usepackage{mathptmx,helvet}
\usepackage[T1]{fontenc}
\usepackage{textcomp}


\usepackage[
            pdftex=true,
            bookmarks=true,        
            a4paper=true,
            pdftitle={Written with GAPDoc},
            pdfcreator={LaTeX with hyperref package / GAPDoc},
            colorlinks=true,
            backref=page,
            breaklinks=true,
            linkcolor=linkColor,
            citecolor=citeColor,
            filecolor=fileColor,
            urlcolor=urlColor,
            pdfpagemode={UseNone}, 
           ]{hyperref}

\newcommand{\maintitlesize}{\fontsize{36}{38}\selectfont}

% write page numbers to a .pnr log file for online help
\newwrite\pagenrlog
\immediate\openout\pagenrlog =\jobname.pnr
\immediate\write\pagenrlog{PAGENRS := [}
\newcommand{\logpage}[1]{\protect\write\pagenrlog{#1, \thepage,}}
%% were never documented, give conflicts with some additional packages

\newcommand{\GAP}{\textsf{GAP}}

%% nicer description environments, allows long labels
\usepackage{enumitem}
\setdescription{style=nextline}

%% depth of toc
\setcounter{tocdepth}{1}




\makeatletter
\renewcommand*\l@section{\@dottedtocline{1}{1.5em}{2.8em}}
\renewcommand*\l@subsection{\@dottedtocline{2}{3.8em}{4.0em}}
\def\@pnumwidth{2.1em}
\makeatother


%% command for ColorPrompt style examples
\newcommand{\gapprompt}[1]{\color{promptColor}{\bfseries #1}}
\newcommand{\gapbrkprompt}[1]{\color{brkpromptColor}{\bfseries #1}}
\newcommand{\gapinput}[1]{\color{gapinputColor}{#1}}


\begin{document}

\logpage{[ 0, 0, 0 ]}
\begin{titlepage}
\mbox{}\vfill

\begin{center}{\maintitlesize \textbf{GAP \texttt{\symbol{45}} Reference Manual\\
\mbox{}}}\\
\vfill

\hypersetup{pdftitle={GAP \texttt{\symbol{45}} Reference Manual}}
\markright{\scriptsize \mbox{}\hfill GAP \texttt{\symbol{45}} Reference Manual \hfill\mbox{}}
{\Huge Release 4.17dev, today\mbox{}}\\[1cm]
\mbox{}\\[2cm]
{\Large \textbf{\strut The GAP Group    \strut\mbox{}}}\\
\hypersetup{pdfauthor={The GAP Group   }}
\end{center}\vfill

\mbox{}\\
{\mbox{}\\
\small \noindent \textbf{The GAP Group   }  Email: \href{mailto://support@gap-system.org} {\texttt{support@gap\texttt{\symbol{45}}system.org}}\\
  Homepage: \href{https://www.gap-system.org} {\texttt{https://www.gap\texttt{\symbol{45}}system.org}}}\\
\end{titlepage}

\newpage\setcounter{page}{2}
{\small 
\section*{Copyright}
\logpage{[ 0, 0, 1 ]}
 Copyright {\copyright} (1987\texttt{\symbol{45}}this year) for the core part
of the \textsf{GAP} system by the \textsf{GAP} Group. 

 Most parts of this distribution, including the core part of the \textsf{GAP} system are distributed under the terms of the GNU General Public License
Version 2, see \href{https://www.gnu.org/licenses/old-licenses/gpl-2.0.html} {\texttt{https://www.gnu.org/licenses/old\texttt{\symbol{45}}licenses/gpl\texttt{\symbol{45}}2.0.html}} or the \texttt{LICENSE} file in the root directory of the \textsf{GAP} installation. 

 More detailed information about copyright and licenses of parts of this
distribution can be found in Section \ref{Copyright and License} of this manual. 

 \textsf{GAP} has been developed over a long time and has many authors and contributors.
More detailed information can be found in Section \ref{Authors and Maintainers} of this manual. \mbox{}}\\[1cm]
\newpage

\def\contentsname{Contents\logpage{[ 0, 0, 2 ]}}

\tableofcontents
\newpage

           
\chapter{\textcolor{Chapter }{Preface}}\label{Preface}
\logpage{[ 1, 0, 0 ]}
\hyperdef{L}{X874E1D45845007FE}{}
{
  \index{About GAP manual@About \textsf{GAP} manual} Welcome to \textsf{GAP}. This is one of the manuals documenting the core part of \textsf{GAP}, the others being the \emph{\textsf{GAP} Tutorial}  and \emph{HPC\texttt{\symbol{45}}GAP Reference Manual}.  

 There is also a document \texttt{CHANGES.md} on the changes from earlier versions in the root directory. This preface
serves not only to introduce ``The \textsf{GAP} Reference Manual'', but also as an introduction to the whole system. 

 \textsf{GAP} stands for \emph{Groups, Algorithms and Programming}. The name was chosen to reflect the original aim of the system, which is
introduced in this reference manual. Since that choice, the system has become
somewhat broader, and you will also find information about algorithms and
programming for other algebraic structures, such as semigroups and algebras. 

 This manual, the \emph{\textsf{GAP} Reference Manual} contains the official definitions of \textsf{GAP} functions. It should contain all the information needed to use \textsf{GAP}, and is not intended to be read
cover\texttt{\symbol{45}}to\texttt{\symbol{45}}cover. 

 To get started a new user may first look at parts of the \emph{\textsf{GAP} Tutorial}. 

 A lot of the functionality of the system and a number of contributed
extensions are provided as ``\textsf{GAP} packages'' which are developed independently of the core part of \textsf{GAP} and can be loaded into a \textsf{GAP} session. Each package comes with a its own manual which is also available
through the \textsf{GAP} help system. 

 This manual is divided into chapters, sections and subsections.
Chapter{\nobreakspace}\ref{The Help System} describes the \emph{help system}, which provides access to all the manuals from a running \textsf{GAP} session. Chapter{\nobreakspace}\ref{Running GAP} gives technical advice for \emph{running} \textsf{GAP}. Chapter{\nobreakspace}\ref{The Programming Language} introduces the \textsf{GAP} language, and the next chapters deal with the \emph{environment} provided by \textsf{GAP} for the user. These are followed by the main bulk of chapters which are
devoted to the various mathematical structures that \textsf{GAP} can handle. 

 Subsequent sections of this preface explain the structure of the system and
provide copyright and licensing information. 

  
\section{\textcolor{Chapter }{The \textsf{GAP} System}}\label{The GAP System}
\logpage{[ 1, 1, 0 ]}
\hyperdef{L}{X863F306C7D32F4B0}{}
{
            \textsf{GAP} is a \emph{free}, \emph{open} and \emph{extensible} software package for computation in discrete abstract algebra. The terms ``free'' and ``open'' describe the conditions under which the system is distributed
\texttt{\symbol{45}}\texttt{\symbol{45}} in brief, it is \emph{free of charge} (except possibly for the immediate costs of delivering it to you), you are \emph{free to pass it on} within certain limits, and all of the workings of the system are \emph{open for you to examine and change}. Details of these conditions can be found in Section  (\textbf{Reference: Copyright and License}). 

 The system is ``extensible'' in that you can write your own programs in the \textsf{GAP} language, and use them in just the same way as the programs which form part of
the system (the ``library''). Indeed, we actively support the contribution, refereeing and distribution
of extensions to the system, in the form of ``\textsf{GAP} packages''. Further details of this can be found in Chapter  (\textbf{Reference: Using and Developing GAP Packages}), and on our website. 

 Development of \textsf{GAP} began at Lehrstuhl D f{\"u}r Mathematik, RWTH\texttt{\symbol{45}}Aachen, under
the leadership of Joachim Neub{\"u}ser in 1985. Version 2.4 was released in
1988 and version 3.1 in 1992. In 1997 coordination of \textsf{GAP} development, now very much an international effort, was transferred to St
Andrews. A complete internal redesign and almost complete rewrite of the
system was completed over the following years and version 4.1 was released in
July 1999. A sign of the further internationalization of the project was the \textsf{GAP}{\nobreakspace}4.4 release in 2004, which has been coordinated from Colorado
State University, Fort Collins. 

 More information on the motivation and development of \textsf{GAP} to date, can be found on our website in a section entitled ``Some History of \textsf{GAP}'': \href{https://www.gap-system.org/about/history/} {\texttt{https://www.gap\texttt{\symbol{45}}system.org/about/history/}}. 

 For those readers who have used an earlier version of \textsf{GAP}, an overview of the changes from \textsf{GAP}{\nobreakspace}4.4 and a brief summary of changes from earlier versions is
given in \texttt{CHANGES.md} file in the main directory. 

 The system that you are getting now consists of a ``core system'' and a number of packages. The core system consists of four main parts. 
\begin{enumerate}
\item  A \emph{kernel}, written in C, which provides the user with 
\begin{itemize}
\item  automatic dynamic storage management, which the user needn't bother about when
programming; 
\item  a set of time\texttt{\symbol{45}}critical basic functions, e.g. ``arithmetic'', operations for integers, finite fields, permutations and words, as well as
natural operations for lists and records; 
\item  an interpreter for the \textsf{GAP} language, an untyped imperative programming language with functions as first
class objects and some extra built\texttt{\symbol{45}}in data types such as
permutations and finite field elements. The language supports a form of
object\texttt{\symbol{45}}oriented programming, similar to that supported by
languages like C++ and Java but with some important differences. 
\item  a small set of system functions allowing the \textsf{GAP} programmer to handle files and execute external programs in a uniform way,
regardless of the particular operating system in use. 
\item  a set of programming tools for testing, debugging, and timing algorithms. 
\item  a ``read\texttt{\symbol{45}}eval\texttt{\symbol{45}}print loop'' (REPL) user interface. 
\end{itemize}
 
\item  A much larger \emph{library of \textsf{GAP} functions} that implement algebraic and other algorithms. Since this is written entirely
in the \textsf{GAP} language, the \textsf{GAP} language is both the main implementation language and the user language of the
system. Therefore a user can, as easily as the original programmers,
investigate and vary algorithms of the library and add new ones to it, first
for their own use and eventually for the benefit of all \textsf{GAP} users. 
\item  A \emph{library of group theoretical data} which contains various libraries of groups, including the library of small
groups (containing all groups of order at most 2000, except those of order
1024) and others. Large libraries of ordinary and Brauer character tables and
Tables of Marks are included as packages. 
\item  The \emph{documentation}. This is available as on\texttt{\symbol{45}}line help, as printable files in
PDF format and as HTML for viewing with a web browser. 
\end{enumerate}
 

 Also included with the core system are some test files and a few small
utilities which we hope you will find useful. 

 \textsf{GAP} \emph{packages} are self\texttt{\symbol{45}}contained extensions to the core system. A package
contains \textsf{GAP} code and its own documentation and may also contain data files or external
programs to which the \textsf{GAP} code provides an interface. These packages may be loaded into \textsf{GAP} using the \texttt{LoadPackage} (\textbf{Reference: LoadPackage}) command, and both the package and its documentation are then available just as
if they were parts of the core system. Some packages may be loaded
automatically, when \textsf{GAP} is started, if they are present. Some packages, because they depend on
external programs, may only be available on the operating systems where those
programs are available (usually UNIX). You should note that, while the
packages included with this release are the most recent versions ready for
release at this time, new packages and new versions may be released at any
time and can be easily installed in your copy of \textsf{GAP}. 

  With \textsf{GAP} there are two packages (the library of ordinary and Brauer character tables,
and the library of tables of marks) which contain functionality developed from
parts of the \textsf{GAP} core system. These have been moved into packages for ease of maintenance and
to allow new versions to be released independently of new releases of the core
system. The library of small groups should also be regarded as a package,
although it does not currently use the standard package mechanism. Other
packages contain functionality which has never been part of the core system,
and may extend it substantially, implementing specific algorithms to enhance
its capabilities, providing data libraries, interfaces to other computer
algebra systems and data sources such as the electronic version of the Atlas
of Finite Group Representations; therefore, installation and usage of packages
is recommended. 

 Further details about \textsf{GAP} packages can be found in chapter  (\textbf{Reference: Using and Developing GAP Packages}), and on the \textsf{GAP} website at \href{https://www.gap-system.org/packages/authors/} {\texttt{https://www.gap\texttt{\symbol{45}}system.org/packages/authors/}}.    }

  
\section{\textcolor{Chapter }{Authors and Maintainers}}\label{Authors and Maintainers}
\logpage{[ 1, 2, 0 ]}
\hyperdef{L}{X877A62A1781C2147}{}
{
  \textsf{GAP} is the work of very many people, some of whom still maintain parts of the
system. A complete list of authors can be found at \href{https://github.com/gap-system/gap/blob/master/COPYRIGHT} {\texttt{https://github.com/gap\texttt{\symbol{45}}system/gap/blob/master/COPYRIGHT}}. All \textsf{GAP} packages have their own authors and maintainers. It should however be noted
that some packages provide interfaces between \textsf{GAP} and an external program, a copy of which is included for convenience, and
that, in these cases, we do not claim that the package authors or maintainers
wrote, or maintain, this external program. Similarly, the system and some
packages include large data libraries that may have been computed by many
people. We try to make clear in each case what credit is attributable to whom. 

 We have, for some time, operated a refereeing system for contributed packages,
both to ensure the quality of the software we distribute, and to provide
recognition for the authors. We now consider this to be a refereeing system
for modules, and we would note, in particular that, although it does not use
the standard package interface, the library of small groups has been refereed
and accepted on exactly the same basis as the accepted packages.  

 We also include with this distribution a number of packages which have not
(yet) gone through our refereeing process. Some may be accepted in the future,
in other cases the authors have chosen not to submit them. More information
can be found on our website \href{https://www.gap-system.org} {\texttt{https://www.gap\texttt{\symbol{45}}system.org}}. 

 }

  
\section{\textcolor{Chapter }{Acknowledgements}}\label{Acknowledgements}
\logpage{[ 1, 3, 0 ]}
\hyperdef{L}{X82A988D47DFAFCFA}{}
{
  Very many people have worked on, and contributed to, \textsf{GAP} over the years since its inception. On our website you will find the prefaces
to the previous releases, each of which acknowledges people who have made
special contributions to that release. Even so, it is appropriate to mention
here Joachim Neub{\"u}ser whose vision of a free, open and extensible system
for computational algebra inspired \textsf{GAP} in the first place, and Martin Sch{\"o}nert, who was the technical architect
of \textsf{GAP} 3 and \textsf{GAP} 4. 

 }

  
\section{\textcolor{Chapter }{Copyright and License}}\label{Copyright and License}
\logpage{[ 1, 4, 0 ]}
\hyperdef{L}{X7950EFA183E3F666}{}
{
          Copyright {\copyright} 1987\texttt{\symbol{45}}this year by its authors, which
are listed at \href{https://github.com/gap-system/gap/blob/master/COPYRIGHT} {\texttt{https://github.com/gap\texttt{\symbol{45}}system/gap/blob/master/COPYRIGHT}}, except for files in the distribution, which have an explicit different
copyright statement. In particular, the copyright of packages distributed with \textsf{GAP} is usually with the package authors or their institutions. 

 \textsf{GAP} is free software; you can redistribute it and/or modify it under the terms of
the GNU General Public License as published by the Free Software Foundation;
either version 2 of the License, or (at your option) any later version. 

 \textsf{GAP} is distributed in the hope that it will be useful, but WITHOUT ANY WARRANTY;
without even the implied warranty of MERCHANTABILITY or FITNESS FOR A
PARTICULAR PURPOSE. See the GNU General Public License for more details. 

 You should have received a copy of the GNU General Public License along with \textsf{GAP}, see the file \texttt{LICENSE} in the root directory of the \textsf{GAP} distribution or see \href{https://www.gnu.org/licenses/gpl.html} {\texttt{https://www.gnu.org/licenses/gpl.html}}. 

 If you publish a mathematical result that was partly obtained using \textsf{GAP}, please cite \textsf{GAP}, just as you would cite another paper that you used (see below for a sample
citation). 

 Specifically, please refer to 

 
\begin{Verbatim}[commandchars=!@|,fontsize=\small,frame=single,label=]
  [GAP] The GAP Group, GAP - Groups, Algorithms, and Programming,
        Version 4.17dev; this year (https://www.gap-system.org)
\end{Verbatim}
 

 \textsf{GAP} for Windows is packaged with several packages from \textsf{Cygwin} (\href{https://cygwin.com} {\texttt{https://cygwin.com}}). Files from \textsf{Cygwin} are provided under their respective licenses, which are included in the
installation. The \textsf{GAP} for Windows executables that we distribute are themselves released under the
terms of the GNU General Public License (GNU GPL); either version 2 of the
License, or (at your option) any later version. 

 Please contact \href{mailto://support@gap-system.org} {\texttt{support@gap\texttt{\symbol{45}}system.org}} if you need further information. }

 }

          
\chapter{\textcolor{Chapter }{The Help System}}\label{The Help System}
\logpage{[ 2, 0, 0 ]}
\hyperdef{L}{X8755A2C67B197C63}{}
{
  This chapter describes the \textsf{GAP} help system. The help system lets you read the documentation interactively.  
\section{\textcolor{Chapter }{Invoking the Help}}\label{Invoking the Help}
\logpage{[ 2, 1, 0 ]}
\hyperdef{L}{X7E2C53D2844DD8C3}{}
{
  The basic command to read \textsf{GAP}'s documentation from within a \textsf{GAP} session is as follows. 

 \index{getting help} \texttt{?[\mbox{\texttt{\mdseries\slshape book}}:][?]\mbox{\texttt{\mdseries\slshape topic}}} 

 For an explanation and some examples see{\nobreakspace} (\textbf{Tutorial: Help}). 

 Note that the first question mark must appear in the \emph{first position} after the \texttt{gap{\textgreater} } prompt. The search strings \mbox{\texttt{\mdseries\slshape book}} and \mbox{\texttt{\mdseries\slshape topic}} are normalized in a certain way (see the end of this section for details)
before the search starts. This makes the search case insensitive and there can
be arbitrary white space after the first question mark. 

 When there are several manual sections that match the query a numbered list of
topics is displayed. These matches can be accessed with \texttt{?\mbox{\texttt{\mdseries\slshape number}}}. 

 There are some further specially handled commands which start with a question
mark. They are explained in Section{\nobreakspace}\ref{Browsing through the Sections}. 

 By default \textsf{GAP} shows the help sections as text in the terminal (window), page by page if the
shown text does not fit on the screen. But there are several other choices to
read (other formats of) the documents: via a viewer for \texttt{pdf} files or via a web browser. This is explained below in Section{\nobreakspace}\ref{Changing the Help Viewer}. 

 \emph{Details of the string normalization process} 

 Here is a precise description how the search strings \mbox{\texttt{\mdseries\slshape book}} and \mbox{\texttt{\mdseries\slshape topic}} are normalized before a search starts: backslashes and double or single quotes
are removed, parentheses and braces are substituted by blanks,
non\texttt{\symbol{45}}ASCII characters are considered as
ISO\texttt{\symbol{45}}latin1 characters and the accented letters are
substituted by their non\texttt{\symbol{45}}accented counterpart. Finally
white space is normalized. }

  
\section{\textcolor{Chapter }{Browsing through the Sections}}\label{Browsing through the Sections}
\logpage{[ 2, 2, 0 ]}
\hyperdef{L}{X7BE8068878B7D7D1}{}
{
  Help books for \textsf{GAP} are organized in chapters, sections, and subsections. There are a few special
commands starting with a question mark (in the first position after the \texttt{gap{\textgreater} } prompt) which allow browsing a book section or chapter wise. 

 \index{browsing forward} \texttt{?{\textgreater}} 

 \index{browsing backwards} \texttt{?{\textless}} 

 The two help commands \texttt{?{\textless}} and \texttt{?{\textgreater}} allow one to browse through a whole help book. \texttt{?{\textless}} displays the section or subsection preceding the previously shown
(sub)section, and \texttt{?{\textgreater}} takes you to the section or subsection following the previously shown one. 

 \index{browsing forward one chapter} \texttt{?{\textgreater}{\textgreater}} 

 \index{browsing backwards one chapter} \texttt{?{\textless}{\textless}} 

 \texttt{?{\textless}{\textless}} takes you back to the beginning of the current chapter. If you are already at
the start of a chapter \texttt{?{\textless}{\textless}} takes you to the beginning of the previous chapter. \texttt{?{\textgreater}{\textgreater}} takes you to the beginning of the next chapter. 

 \index{browsing the previous section browsed} \texttt{?\texttt{\symbol{45}}} 

 \index{browsing the next section browsed} \texttt{?+} 

 \textsf{GAP} remembers the last few sections that you have read. \texttt{?\texttt{\symbol{45}}} takes you to the one that you have read before the current one, and displays
it again. Further applications of \texttt{?\texttt{\symbol{45}}} take you further back in this history. \texttt{?+} reverses this process, i.e., it takes you back to the section that you have
read after the current one. It is important to note that \texttt{?\texttt{\symbol{45}}} and \texttt{?+} do not alter the history like the other help commands. 

 \index{list of available books} \texttt{?books} 

 This command shows a list of the books which are currently known to the help
system. For each book there is a short name which is used with the \mbox{\texttt{\mdseries\slshape book}} part of the basic help query and there is a long name which hopefully tells
you what this book is about. 

 A short name which ends in \texttt{(not loaded)} refers to a \textsf{GAP} package whose documentation is loaded but which needs a call of \texttt{LoadPackage} (\ref{LoadPackage}) before you can use the described functions. 

 \index{table of sections for help books} \texttt{?[\mbox{\texttt{\mdseries\slshape book}}:]sections} 

 \index{table of chapters for help books} \texttt{?[\mbox{\texttt{\mdseries\slshape book}}:][chapters]} 

 These commands show tables of contents for all available, respectively the
matching books. For some books these commands show the same, namely the whole
table of contents. 

 \index{redisplay a help section} \texttt{?} 

 \index{redisplay with next help viewer} \texttt{?\&} 

 These commands redisplay the last shown help section. In the form \texttt{?\&} the next preferred help viewer is used for the display (provided one has
chosen several viewers), see{\nobreakspace}\texttt{SetHelpViewer} (\ref{SetHelpViewer}) below. }

  
\section{\textcolor{Chapter }{Changing the Help Viewer}}\label{Changing the Help Viewer}
\logpage{[ 2, 3, 0 ]}
\hyperdef{L}{X863FF9087EDA8DF9}{}
{
  \index{document formats (text, dvi, ps, pdf, HTML)} Books of the \textsf{GAP} help system or package manuals can be available in several formats. Currently
the following formats occur (not all of them may be available for all books): 

 
\begin{description}
\item[{text}]  This is used for display in the terminal window in which \textsf{GAP} is running. Complicated mathematical expressions may not be easy to read in
this format. 
\item[{pdf}]  Adobe's \texttt{pdf} format. Can be used for printing and onscreen reading on most current systems
(with freely available software). Some manual books contain hyperlinks in this
format. 
\item[{HTML}]  The format of web pages. Can be used with any web browser. There may be
hyperlink information available which allows a convenient browsing through the
book via cross\texttt{\symbol{45}}references. This format has the problem that
complicated formulae may be not be easy to read since there is no syntax for
formulae in HTML. (Some older manual books use special symbol fonts for
formulae and need a particular configuration of the web browser for correct
display. Some manuals may use technology for quite sophisticated formula
display.) 
\end{description}
 

 Depending on your operating system and available additional software you can
use several of these formats with \textsf{GAP}'s help system. This is configured with the following command. 

\subsection{\textcolor{Chapter }{SetHelpViewer}}
\logpage{[ 2, 3, 1 ]}\nobreak
\hyperdef{L}{X87C1BFB2826488B0}{}
{\noindent\textcolor{FuncColor}{$\triangleright$\enspace\texttt{SetHelpViewer({\mdseries\slshape viewer1, viewer2, ...})\index{SetHelpViewer@\texttt{SetHelpViewer}}
\label{SetHelpViewer}
}\hfill{\scriptsize (function)}}\\


 This command takes an arbitrary number of arguments which must be strings
describing a viewer. The recognized viewers are explained below. A call with
no arguments shows the current setting. 

 The first given arguments are those with higher priority. So, if a help
section is available in the format needed by \mbox{\texttt{\mdseries\slshape viewer1}}, this viewer is used. If not, availability of the format for \mbox{\texttt{\mdseries\slshape viewer2}} is checked and so on. Recall that the command \texttt{?\&} displays the last seen section again but with the next possible viewer in your
list, see{\nobreakspace}\ref{Browsing through the Sections}. 

 The viewer \texttt{"screen"} (see below) is always silently appended since we assume that each help book is
available in text format. 

 If you want to change the default setting you can use a call of \texttt{SetUserPreference( "HelpViewers", [ ... ] );} (the list in the second argument containing the viewers you want) in your \texttt{gap.ini} file (see{\nobreakspace}\ref{sect:gap.ini}). 

 
\begin{description}
\item[{\texttt{"screen"}}]  This is the default setting. The help is shown in text format using the \texttt{Pager} (\ref{Pager}) command. Hint: Text versions of manuals are formatted assuming that your
terminal displays at least 80 characters per line, if this is not the case
some sections may look very bad. We suggest to use a terminal in \texttt{UTF\texttt{\symbol{45}}8} encoding with a fixed width font (this is the default on most modern
Linux/Windows/Mac systems anyway). Terminals in \texttt{ISO\texttt{\symbol{45}}8859\texttt{\symbol{45}}X} encoding will also work reasonably well (so far, since we do not yet use many
special characters which such terminals could not display). 
\item[{\texttt{"firefox"}, \texttt{"chrome"}, \texttt{"mozilla"}, \texttt{"netscape"}, \texttt{"konqueror"}}]  If a book is available in HTML format this is shown using the corresponding
web browser. How well this works, for example by using a running instance of
this browser, depends on your particular start script of this browser. (Note,
that for some old books the browser must be configured to use symbol fonts.) 
\item[{\texttt{"browser"}}]  (for MS Windows) If a book is available in HTML format, it will be opened
using the Windows default application (typically, a web browser). 
\item[{\texttt{"links2"}, \texttt{"w3m"}, \texttt{"lynx"}}]  If a book is available in HTML format this is shown using the text based \texttt{"links2"} (in graphics mode), \texttt{w3m} or \texttt{lynx} web browser, respectively, inside the terminal running \textsf{GAP}. (Formulae in some older books which use symbol fonts may be unreadable.) 
\item[{\texttt{"mac default browser"}, \texttt{"browser"}, \texttt{"safari"}, \texttt{"firefox"}}]  (for macOS) If a book is available in HTML format this is shown in a web
browser. The options \texttt{"safari"} and \texttt{"firefox"} use the corresponding browsers. The other two options use the program default
browser (which can be set in Safari's preferences, in the "General" tab). 
\item[{\texttt{"xpdf"}}]  (on X window systems) If a book is available in pdf format it is shown with
the onscreen viewer program \texttt{xpdf} (which must be installed on your system). This is a nice program, once it is
running it is reused by \textsf{GAP} for the next displays of help sections. 
\item[{\texttt{"acroread"}}]  If a book is available in pdf format it is shown with the onscreen viewer
program \texttt{acroread} (which must be available on your system). This program does not allow remote
commands or startup with a given page. Therefore the page numbers you have to
visit are just printed on the screen. When you are looking at several sections
of the same book, this viewer assumes that the \texttt{acroread} window still exists. When you go to another book a new acroread window is
launched. 
\item[{\texttt{"pdf viewer"}, \texttt{"skim"}, \texttt{"preview"}, \texttt{"adobe reader"}}]  (for macOS) If a book is available in pdf format this is shown in a pdf
viewer. The options \texttt{"skim"}, \texttt{"preview"} and \texttt{"adobe reader"} use the corresponding viewers. The other two options use the pdf viewer which
you have chosen to open pdf files from the Finder. Note that only \texttt{"Skim"} seems to be capable to open a pdf file on a given page. For the other help
viewers, the page numbers where the information can be found will just be
printed on the screen. None of the help viewers seems to be capable of opening
a pdf at a given named destination (i. e., jump to precisely the place where
the information can be found). The pdf viewer \texttt{"Skim"} is open source software, it can be downloaded from \texttt{https://skim\texttt{\symbol{45}}app.sourceforge.io/}. 
\item[{\texttt{"less"} or \texttt{"more"}}]  This is the same as \texttt{"screen"} but additionally the user preferences \texttt{"Pager"} and \texttt{"PagerOptions"} are set, see the section{\nobreakspace}\ref{The Pager Command} for more details. 
\end{description}
 

 Please, send ideas for further viewer commands to \href{mailto://support@gap-system.org} {\texttt{support@gap\texttt{\symbol{45}}system.org}}. }

 }

  
\section{\textcolor{Chapter }{The Pager Command}}\label{The Pager Command}
\logpage{[ 2, 4, 0 ]}
\hyperdef{L}{X84AFFC817B282359}{}
{
  \textsf{GAP} contains a builtin pager which shows a text string which does not fit on the
screen page by page. Its functionality is very rudimentary and
self\texttt{\symbol{45}}explaining. This is because (at least under UNIX)
there are powerful external standard programs which do this job. 

\subsection{\textcolor{Chapter }{Pager}}
\logpage{[ 2, 4, 1 ]}\nobreak
\hyperdef{L}{X7ED03E41792C3840}{}
{\noindent\textcolor{FuncColor}{$\triangleright$\enspace\texttt{Pager({\mdseries\slshape lines})\index{Pager@\texttt{Pager}}
\label{Pager}
}\hfill{\scriptsize (function)}}\\


 This function can be used to display a text on screen using a pager, i.e., the
text is shown page by page. 

 There is a default builtin pager in \textsf{GAP} which has very limited capabilities but should work on any system. 

 At least on a UNIX system one should use an external pager program like \texttt{less} or \texttt{more}. \textsf{GAP} assumes that this program has command line options \texttt{+nr} and \texttt{+/str} which start the display of the text with line number \texttt{nr} or at the line with the first occurrence of the string \texttt{str}. 

 Which pager is used can be controlled by setting the user preference \texttt{"Pager"}. The default value is \texttt{"builtin"} which means that the internal pager is used. 

 On UNIX systems you probably want to set the user preference \texttt{"Pager"} to the value \texttt{"less"} or \texttt{"more"}, you can do this for example in your \texttt{gap.ini} file (see \ref{sect:gap.ini}). In that case you can also tell \textsf{GAP} a list of standard options for the external pager, via the user preference \texttt{"PagerOptions"}. 

 
\begin{Verbatim}[commandchars=!@|,fontsize=\small,frame=single,label=Example]
    SetUserPreference( "Pager", "less" );
    SetUserPreference( "PagerOptions", ["-f","-r","-a","-i","-M","-j2"] );
\end{Verbatim}
 

 The argument \mbox{\texttt{\mdseries\slshape lines}} can have one of the following forms: 

 
\begin{enumerate}
\item  a string (i.e., lines are separated by newline characters) 
\item  a list of strings (without newline characters) which are interpreted as lines
of the text to be shown 
\item  a record with component \texttt{lines} as in 1. or 2. and optional further components 
\end{enumerate}
 

 In case{\nobreakspace}3. currently the following additional components are
used: 

 
\begin{description}
\item[{\texttt{formatted}}]  can be \texttt{false} or \texttt{true}. If set to \texttt{true} the builtin pager tries to show the text exactly as it is given (avoiding \textsf{GAP}'s automatic line breaking), 
\item[{\texttt{start}}]  must be a positive integer or a string. An integer is interpreted as the
number of the first line shown by the pager, a string is interpreted as a
search string such that the first line containing this string is the first
line shown by the pager (in both cases, one may see the beginning of the text
via back scrolling), 
\item[{\texttt{exitAtEnd}}]  can be \texttt{false} or \texttt{true}. If set to \texttt{true} (the default), the builtin pager is terminated as soon as the end of the list
is shown; otherwise entering the \textsc{q} key is necessary in order to return from the pager. 
\end{description}
 

 The \texttt{Pager} command is used by \textsf{GAP}'s help system for displaying help sections in text format. But, of course, it
may be used for other purposes as well. 

 
\begin{Verbatim}[commandchars=!@|,fontsize=\small,frame=single,label=Example]
  !gapprompt@gap>| !gapinput@s6 := SymmetricGroup(6);;|
  !gapprompt@gap>| !gapinput@words := ["This", "is", "a", "very", "stupid", "example"];;|
  !gapprompt@gap>| !gapinput@l := List(s6, p-> Permuted(words, p));;|
  !gapprompt@gap>| !gapinput@Pager(List(l, a-> JoinStringsWithSeparator(a," ")));;|
\end{Verbatim}
 }

 }

 }

                  
\chapter{\textcolor{Chapter }{Running GAP}}\label{Running GAP}
\logpage{[ 3, 0, 0 ]}
\hyperdef{L}{X79CCD3A6821E5A37}{}
{
  \index{options} This chapter contains information about the command line options for \textsf{GAP} (see{\nobreakspace}\ref{Command Line Options}), about some files in user specific \textsf{GAP} root directory (see{\nobreakspace}\ref{sect:gap.ini}) and about saving and loading a \textsf{GAP} workspace (see{\nobreakspace}\ref{Saving and Loading a Workspace}).  
\section{\textcolor{Chapter }{Command Line Options}}\label{Command Line Options}
\logpage{[ 3, 1, 0 ]}
\hyperdef{L}{X782751D5858A6EAF}{}
{
  \index{features!under UNIX} \index{UNIX!features} \index{options!under UNIX} \index{UNIX!options} \index{GAPInfo.CommandLineOptions@\texttt{GAPInfo.CommandLineOptions}} When you start \textsf{GAP} from a command line or from a script you may specify a number of options on
the command\texttt{\symbol{45}}line to change the default behaviour of \textsf{GAP}. All these options start with a hyphen \texttt{\texttt{\symbol{45}}}, followed by a single letter. Options must not be grouped, e.g., \texttt{gap \texttt{\symbol{45}}gq} is invalid, use \texttt{gap \texttt{\symbol{45}}g \texttt{\symbol{45}}q} instead. Some options require an argument, this must follow the option and
must be separated by whitespace, e.g., \texttt{gap \texttt{\symbol{45}}m 256m}, it is not correct to say \texttt{gap \texttt{\symbol{45}}m256m} instead. Certain boolean options (\texttt{\texttt{\symbol{45}}b}, \texttt{\texttt{\symbol{45}}q}, \texttt{\texttt{\symbol{45}}e}, \texttt{\texttt{\symbol{45}}r}, \texttt{\texttt{\symbol{45}}A}, \texttt{\texttt{\symbol{45}}D}, \texttt{\texttt{\symbol{45}}M}, \texttt{\texttt{\symbol{45}}T}, \texttt{\texttt{\symbol{45}}X}, \texttt{\texttt{\symbol{45}}Y}) toggle the current value so that \texttt{gap \texttt{\symbol{45}}b \texttt{\symbol{45}}b} is equivalent to \texttt{gap} and to \texttt{gap \texttt{\symbol{45}}b \texttt{\symbol{45}}q \texttt{\symbol{45}}b
\texttt{\symbol{45}}q} etc. 

 \textsf{GAP} for UNIX will distinguish between upper and lower case options. 

 As described in the \textsf{GAP} installation instructions (see the \texttt{INSTALL.md} file in the \textsf{GAP} root directory), usually you will not execute \textsf{GAP} directly. Instead you will call a (shell) script, with the name \texttt{gap}, which in turn executes \textsf{GAP}. This script sets some options which are necessary to make \textsf{GAP} work on your system. This means that the default settings mentioned below may
not be what you experience when you execute \textsf{GAP} on your system. 

 During a \textsf{GAP} session, one can find the current values of command line options in the record \texttt{GAPInfo.CommandLineOptions} (see \texttt{GAPInfo} (\ref{GAPInfo})), whose component names are the command line options (without the leading \texttt{\texttt{\symbol{45}}}). 

 
\begin{description}
\item[{\index{-A@\texttt{\texttt{\symbol{45}}A}} \texttt{\texttt{\symbol{45}}A}}]  By default, some needed and suggested \textsf{GAP} packages (see \ref{Using and Developing GAP Packages}) are loaded, if present, into the \textsf{GAP} session when it starts. This option disables (actually toggles) the loading of
suggested packages, which can be useful for debugging or testing. The needed
packages (and their needed packages, and so on) are loaded in any case. 
\item[{\index{-b@\texttt{\texttt{\symbol{45}}b}} \texttt{\texttt{\symbol{45}}b}}]  tells \textsf{GAP} to suppress the banner. That means that \textsf{GAP} immediately prints the prompt. This is useful when, after a while, you get
tired of the banner. This option can be repeated to enable the banner; each \texttt{\texttt{\symbol{45}}b} toggles the state of banner display. 
\item[{\index{-c@\texttt{\texttt{\symbol{45}}c}} \texttt{\texttt{\symbol{45}}c }\mbox{\texttt{\mdseries\slshape gapcode}}}]  tells \textsf{GAP} to execute the given \textsf{GAP} code as if it was entered into a temporary file which then is processed
together with the other files given to GAP (see the explanation at the end of
this list for further details on how filenames are processed). 
\item[{\index{-D@\texttt{\texttt{\symbol{45}}D}} \texttt{\texttt{\symbol{45}}D}}]  The \texttt{\texttt{\symbol{45}}D} option tells \textsf{GAP} to print short messages when it is reading files or loading modules. This
option may be repeated to toggle this behavior on and off. The message, 

 
\begin{Verbatim}[commandchars=!@|,fontsize=\small,frame=single,label=Example]
  #I  READ_GAP_ROOT: loading 'lib/kernel.g' as GAP file
\end{Verbatim}
 

 tells you that \textsf{GAP} has started to read the library file \texttt{lib/kernel.g}. 

 
\begin{Verbatim}[commandchars=!@|,fontsize=\small,frame=single,label=Example]
  #I  READ_GAP_ROOT: loading 'lib/kernel.g' statically
\end{Verbatim}
 

 tells you that \textsf{GAP} has used the compiled version of the library file \texttt{lib/kernel.g}. This compiled module was statically linked to the \textsf{GAP} kernel at the time the kernel was created. 

 
\begin{Verbatim}[commandchars=!@|,fontsize=\small,frame=single,label=Example]
  #I  READ_GAP_ROOT: loading 'lib/kernel.g' dynamically
\end{Verbatim}
 

 tells you that \textsf{GAP} has loaded the compiled version of the library file \texttt{lib/kernel.g}. This compiled module was dynamically loaded to the \textsf{GAP} kernel at runtime from a corresponding \texttt{.so} file. 

 Obviously, this is a debugging option and most users will not need it. 
\item[{\index{-E@\texttt{\texttt{\symbol{45}}E}} \texttt{\texttt{\symbol{45}}E}}]  If your \textsf{GAP} installation uses the \texttt{readline} library for command line editing (see \ref{sec:readline}), this may be disabled by using \texttt{\texttt{\symbol{45}}E} option. This option may be repeated to toggle this behavior on and off. If
your \textsf{GAP} installation does not use the \texttt{readline} library (you can check by \texttt{IsBound(GAPInfo.UseReadline);} if this is the case), this option will have no effect at all. 
\item[{\index{-e@\texttt{\texttt{\symbol{45}}e}} \texttt{\texttt{\symbol{45}}e}}]  tells \textsf{GAP} not to quit when receiving a \textsc{Ctrl\texttt{\symbol{45}}D} on an empty input line (see{\nobreakspace}\ref{quit}). This option should not be used when the input is a file or pipe. This
option may be repeated to toggle this behavior on and off. 
\item[{\index{-f@\texttt{\texttt{\symbol{45}}f}} \texttt{\texttt{\symbol{45}}f}}]  tells \textsf{GAP} to enable the line editing and history (see \ref{Line Editing}). 

 In general line editing will be enabled if the input is connected to a
terminal. There are rare circumstances, for example when using a remote
session with a corrupted telnet implementation, when this detection fails. Try
using \texttt{\texttt{\symbol{45}}f} in this case to enable line editing. This option does not toggle; you must use \texttt{\texttt{\symbol{45}}n} to disable line editing. 
\item[{\index{-g@\texttt{\texttt{\symbol{45}}g}} \texttt{\texttt{\symbol{45}}g}}]  tells \textsf{GAP} to print a message every time a full garbage collection is performed. (This is
available only if the \texttt{GASMAN} garbage collector is used, see \ref{Garbage Collection}.) 

 
\begin{Verbatim}[commandchars=!@|,fontsize=\small,frame=single,label=Example]
  #G  FULL 44580/2479kb live   57304/4392kb dead   734/4096kb free
\end{Verbatim}
 

 For example, this tells you that there are 44580 live objects that survived a
full garbage collection, that 57304 unused objects were reclaimed by it, and
that 734 kilobytes from a total allocated memory of 4096 kilobytes are
available afterwards. 
\item[{\index{-g -g@\texttt{\texttt{\symbol{45}}g \texttt{\symbol{45}}g}} \texttt{\texttt{\symbol{45}}g \texttt{\symbol{45}}g}}]  If you give the option \texttt{\texttt{\symbol{45}}g} twice, \textsf{GAP} prints a information message every time a partial or full garbage collection
is performed. (This is available only if the \texttt{GASMAN} garbage collector is used, see \ref{Garbage Collection}.) The message, 

 
\begin{Verbatim}[commandchars=!@|,fontsize=\small,frame=single,label=Example]
  #G  PART 9405/961kb+live   7525/1324kb+dead   2541/4096kb free
\end{Verbatim}
 

 for example, tells you that 9405 objects survived the partial garbage
collection and 7525 objects were reclaimed, and that 2541 kilobytes from a
total allocated memory of 4096 kilobytes are available afterwards. 
\item[{\index{-h@\texttt{\texttt{\symbol{45}}h}} \texttt{\texttt{\symbol{45}}h}}] 
\item[{\index{--help@\texttt{\texttt{\symbol{45}}\texttt{\symbol{45}}help}} \texttt{\texttt{\symbol{45}}\texttt{\symbol{45}}help}}]  tells \textsf{GAP} to print a summary of all available options. \textsf{GAP} exits after printing the summary, all options coming after are ignored. 
\item[{\index{-K@\texttt{\texttt{\symbol{45}}K}} \texttt{\texttt{\symbol{45}}K }\mbox{\texttt{\mdseries\slshape memory}}}] 
\item[{\index{--limitworkspace@\texttt{\texttt{\symbol{45}}\texttt{\symbol{45}}limitworkspace}} \texttt{\texttt{\symbol{45}}\texttt{\symbol{45}}limitworkspace }\mbox{\texttt{\mdseries\slshape memory}}}]  is like the \texttt{\texttt{\symbol{45}}o} option. But while the latter actually allocates more memory if the system
allows it and then prints a warning inside a break loop the \texttt{\texttt{\symbol{45}}K} options tells \textsf{GAP} not even to try to allocate more memory. Instead \textsf{GAP} just exits with an appropriate message. The default is that this feature is
switched off. You have to set it explicitly when you want to enable it. 
\item[{\index{-L@\texttt{\texttt{\symbol{45}}L}} \texttt{\texttt{\symbol{45}}L }\mbox{\texttt{\mdseries\slshape filename}}}]  The option \texttt{\texttt{\symbol{45}}L} tells \textsf{GAP} to load a saved workspace. See section{\nobreakspace}\ref{Saving and Loading a Workspace}. (This is available only if the \texttt{GASMAN} garbage collector is used, see \ref{Garbage Collection}.) 
\item[{\index{-l@\texttt{\texttt{\symbol{45}}l}} \texttt{\texttt{\symbol{45}}l }\mbox{\texttt{\mdseries\slshape path{\textunderscore}list}}}] 
\item[{\index{--roots@\texttt{\texttt{\symbol{45}}\texttt{\symbol{45}}roots}} \texttt{\texttt{\symbol{45}}\texttt{\symbol{45}}roots }\mbox{\texttt{\mdseries\slshape path{\textunderscore}list}}}]  can be used to set or modify \textsf{GAP}'s list of root directories (see \ref{GAP Root Directories}). The default if no \texttt{\texttt{\symbol{45}}l } option is given is the current directory \texttt{./}. This option can be used several times. Depending on the \texttt{\texttt{\symbol{45}}r} option a further user specific path is prepended to the list of root
directories (the path in \texttt{GAPInfo.UserGapRoot}). 

 \mbox{\texttt{\mdseries\slshape path{\textunderscore}list}} should be a list of directories separated by semicolons. No whitespace is
permitted before or after a semicolon. If \mbox{\texttt{\mdseries\slshape path{\textunderscore}list}} does not start or end with a semicolon, then \mbox{\texttt{\mdseries\slshape path{\textunderscore}list}} replaces the existing list of root directories. If \mbox{\texttt{\mdseries\slshape path{\textunderscore}list}} starts with a semicolon, then \mbox{\texttt{\mdseries\slshape path{\textunderscore}list}} is appended to the existing list of root directories. If \mbox{\texttt{\mdseries\slshape path{\textunderscore}list}} ends with a semicolon and does not start with one, then the new list of root
directories is the concatenation of \mbox{\texttt{\mdseries\slshape path{\textunderscore}list}} and the existing list of root directories. After \textsf{GAP} has completed its startup procedure and displays the prompt, the list of root
directories can be seen in the variable \texttt{GAPInfo.RootPaths}, see \texttt{GAPInfo} (\ref{GAPInfo}). 

 Usually this option is used inside a startup script to specify where \textsf{GAP} is installed on the system. The \texttt{\texttt{\symbol{45}}l} option can also be used by individual users to tell \textsf{GAP} about privately installed modifications of the library, additional \textsf{GAP} packages and so on. Section \ref{GAP Root Directories} explains how several root paths can be used to do this. 

 \textsf{GAP} will attempt to read the file \mbox{\texttt{\mdseries\slshape root{\textunderscore}dir}}\texttt{/lib/init.g} during startup where \mbox{\texttt{\mdseries\slshape root{\textunderscore}dir}} is one of the directories in its list of root directories. If \textsf{GAP} cannot find its \texttt{init.g} file it will print the following warning. 

 
\begin{Verbatim}[commandchars=!@|,fontsize=\small,frame=single,label=Example]
  gap: hmm, I cannot find 'lib/init.g' maybe use option '-l <gaproot>'?
\end{Verbatim}
 

 It is not possible to use \textsf{GAP} without the library files, so you must not ignore this warning. You should
leave \textsf{GAP} and start it again, specifying the correct root path using the \texttt{\texttt{\symbol{45}}l} option. 
\item[{\index{--packagedirs@\texttt{\texttt{\symbol{45}}\texttt{\symbol{45}}packagedirs}} \texttt{\texttt{\symbol{45}}\texttt{\symbol{45}}packagedirs }\mbox{\texttt{\mdseries\slshape path{\textunderscore}list}}}]  can be used to add paths to \textsf{GAP}'s list of package directories (see \ref{GAP Package Directories}). The list always contains all subdirectories \texttt{pkg} in a \textsf{GAP} root directory. 

 \mbox{\texttt{\mdseries\slshape path{\textunderscore}list}} should be a list of directories separated by semicolons. No whitespace is
permitted before or after a semicolon, and the first and last character of \mbox{\texttt{\mdseries\slshape path{\textunderscore}list}} may not be a semicolon. After \textsf{GAP} has completed its startup procedure and displays the prompt, the list of
package directories can be seen in the variable \texttt{GAPInfo.PackageDirectories}, see \texttt{GAPInfo} (\ref{GAPInfo}). 

 Usually this option is used inside a startup script to specify where
additional \textsf{GAP} packages are located on the system. The \texttt{\texttt{\symbol{45}}\texttt{\symbol{45}}packagedirs} option can also be used by individual users to tell \textsf{GAP} about additional \textsf{GAP} packages, without the need to set up a complete root directory structure. 
\item[{\index{-M@\texttt{\texttt{\symbol{45}}M}} \texttt{\texttt{\symbol{45}}M}}]  tells \textsf{GAP} not to check for, nor to use, compiled versions of library files. This option
may be repeated to toggle this behavior on and off. 
\item[{\index{-m@\texttt{\texttt{\symbol{45}}m}} \texttt{\texttt{\symbol{45}}m }\mbox{\texttt{\mdseries\slshape memory}}}] 
\item[{\index{--minworkspace@\texttt{\texttt{\symbol{45}}\texttt{\symbol{45}}minworkspace}} \texttt{\texttt{\symbol{45}}\texttt{\symbol{45}}minworkspace }\mbox{\texttt{\mdseries\slshape memory}}}]  tells \textsf{GAP} to allocate \mbox{\texttt{\mdseries\slshape memory}} bytes at startup time. If the last character of \mbox{\texttt{\mdseries\slshape memory}} is \texttt{k} or \texttt{K} it is taken as kilobytes, if the last character is \texttt{m} or \texttt{M} \mbox{\texttt{\mdseries\slshape memory}} is taken as megabytes and if it is \texttt{g} or \texttt{G} it is taken as gigabytes. 

 This amount of memory should be large enough so that computations do not
require too many garbage collections. On the other hand, if \textsf{GAP} allocates more memory than is physically available, it will spend most of the
time paging. 
\item[{\index{-n@\texttt{\texttt{\symbol{45}}n}} \texttt{\texttt{\symbol{45}}n}}]  tells \textsf{GAP} to disable the line editing and history (see \ref{Line Editing}). 

 You may want to do this if the command line editing is incompatible with
another program that is used to run \textsf{GAP}. For example if \textsf{GAP} is run from inside a GNU Emacs shell window, \texttt{\texttt{\symbol{45}}n} should be used since otherwise every input line will be echoed twice, once by
Emacs and once by \textsf{GAP}. This option does not toggle; you must use \texttt{\texttt{\symbol{45}}f} to enable line editing. 
\item[{\index{-O@\texttt{\texttt{\symbol{45}}O}} \texttt{\texttt{\symbol{45}}O}}]  disables loading obsolete variables (see Chapter \ref{Replaced and Removed Command Names}). This option is used mainly for testing purposes, for example in order to
make sure that a \textsf{GAP} package or one's own \textsf{GAP} code does not rely on the obsolete variables. 
\item[{\index{-o@\texttt{\texttt{\symbol{45}}o}} \texttt{\texttt{\symbol{45}}o }\mbox{\texttt{\mdseries\slshape memory}}}] 
\item[{\index{--maxworkspace@\texttt{\texttt{\symbol{45}}\texttt{\symbol{45}}maxworkspace}} \texttt{\texttt{\symbol{45}}\texttt{\symbol{45}}maxworkspace }\mbox{\texttt{\mdseries\slshape memory}}}]  tells \textsf{GAP} to allocate at most \mbox{\texttt{\mdseries\slshape memory}} bytes without asking. The option argument \mbox{\texttt{\mdseries\slshape memory}} is specified as with the \texttt{\texttt{\symbol{45}}m} option. 

 If more than this amount is required during the \textsf{GAP} session, \textsf{GAP} prints an error message and enters a break loop. In that case you can enter \texttt{return;} which implicitly doubles the amount given with this option. 
\item[{\index{-q@\texttt{\texttt{\symbol{45}}q}} \texttt{\texttt{\symbol{45}}q}}]  tells \textsf{GAP} to be quiet. This means that \textsf{GAP} displays neither the banner nor the prompt \texttt{gap{\textgreater}}. This is useful if you want to run \textsf{GAP} as a filter with input and output redirection and want to avoid the banner and
the prompts appearing in the output file. This option may be repeated to
disable quiet mode; each \texttt{\texttt{\symbol{45}}q} toggles quiet mode. 
\item[{\index{-R@\texttt{\texttt{\symbol{45}}R}} \texttt{\texttt{\symbol{45}}R}}]  The option \texttt{\texttt{\symbol{45}}R} tells \textsf{GAP} not to load a saved workspace previously specified via the \texttt{\texttt{\symbol{45}}L} option. This option does not toggle. 
\item[{\index{-r@\texttt{\texttt{\symbol{45}}r}} \texttt{\texttt{\symbol{45}}r}}]  The option \texttt{\texttt{\symbol{45}}r} tells \textsf{GAP} to ignore any user specific configuration files. In particular, the user
specific root directory \texttt{GAPInfo.UserGapRoot} is not added to the \textsf{GAP} root directories and so \texttt{gap.ini} and \texttt{gaprc} files that may be contained in that directory are not read, see \ref{sect:gap.ini}. Multiple \texttt{\texttt{\symbol{45}}r} options toggle this behaviour. 
\item[{\index{-s@\texttt{\texttt{\symbol{45}}s}} \texttt{\texttt{\symbol{45}}s }\mbox{\texttt{\mdseries\slshape memory}}}]  With this option \textsf{GAP} does not use \texttt{sbrk} to get memory from the operating system. Instead it uses \texttt{mmap}, \texttt{malloc} or some other command for the amount given with this option to allocate space
for the \texttt{GASMAN} memory manager. Usually \textsf{GAP} does not really use all of this memory, the options \texttt{\texttt{\symbol{45}}m}, \texttt{\texttt{\symbol{45}}o}, \texttt{\texttt{\symbol{45}}K} still work as documented. This feature assumes that the operating system only
assigns physical memory to the \textsf{GAP} process when it is accessed, so that specifying a large amount of memory with \texttt{\texttt{\symbol{45}}s} should not cause any performance problem. The advantage of using this option
is that \textsf{GAP} can work together with kernel modules which allocate a lot of memory with \texttt{malloc}. 

 The option argument \mbox{\texttt{\mdseries\slshape memory}} is specified as with the \texttt{\texttt{\symbol{45}}m} option. 
\item[{\index{-T@\texttt{\texttt{\symbol{45}}T}} \texttt{\texttt{\symbol{45}}T}}]  suppresses the usual break loop behaviour of \textsf{GAP}. With this option \textsf{GAP} behaves as if the user \texttt{quit} immediately from every break loop, and also suppresses displaying any error
backtrace. This is intended for automated testing of \textsf{GAP}. This option may be repeated to toggle this behavior on and off. 
\item[{\index{--version@\texttt{\texttt{\symbol{45}}\texttt{\symbol{45}}version}} \texttt{\texttt{\symbol{45}}\texttt{\symbol{45}}version}}]  tells \textsf{GAP} to print its version and then exit immediately, all options coming after are
ignored. 
\item[{\index{-x@\texttt{\texttt{\symbol{45}}x}} \texttt{\texttt{\symbol{45}}x }\mbox{\texttt{\mdseries\slshape length}}}] 
\item[{\index{--width@\texttt{\texttt{\symbol{45}}\texttt{\symbol{45}}width}} \texttt{\texttt{\symbol{45}}\texttt{\symbol{45}}width }\mbox{\texttt{\mdseries\slshape length}}}]  With this option you can tell \textsf{GAP} how long lines are. \textsf{GAP} uses this value to decide when to split long lines. After starting \textsf{GAP} you may use \texttt{SizeScreen} (\ref{SizeScreen}) to alter the line length. 

 The default value is 80, unless another value can be obtained from the
Operating System, which is the right value if you have a standard terminal
application. If you have a larger monitor, or use a smaller font, or redirect
the output to a printer, you may want to increase this value. 
\item[{\index{-y@\texttt{\texttt{\symbol{45}}y}} \texttt{\texttt{\symbol{45}}y }\mbox{\texttt{\mdseries\slshape length}}}] 
\item[{\index{--line@\texttt{\texttt{\symbol{45}}\texttt{\symbol{45}}line}} \texttt{\texttt{\symbol{45}}\texttt{\symbol{45}}line }\mbox{\texttt{\mdseries\slshape length}}}]  With this option you can tell \textsf{GAP} how many lines your screen has. \textsf{GAP} uses this value to decide after how many lines of on\texttt{\symbol{45}}line
help it should wait. After starting \textsf{GAP} you may use \texttt{SizeScreen} (\ref{SizeScreen}) to alter the number of lines. 

 The default value is 24, unless another value can be obtained from the
Operating System, which is the right value if you have a standard terminal
application. If you have a larger monitor, or use a smaller font, or redirect
the output to a printer, you may want to increase this value. 
\item[{\index{options!command line, filenames} \mbox{\texttt{\mdseries\slshape filename}} ...}]  Further arguments are taken as filenames of files that are read by \textsf{GAP} during startup, after the system and private init files are read, but before
the first prompt is printed. Filenames ending with \texttt{.tst} are processed by \texttt{Test} (\ref{Test}), all other files by \texttt{Read} (\ref{Read}). These files and also commands specified via the \texttt{\texttt{\symbol{45}}c} option are processed in the order in which they appear on the command line. If
a file cannot be opened or if executing the code in it raises an error, then
the usual error handling for \texttt{Read} (\ref{Read}) respectively \texttt{Test} (\ref{Test}) kicks in. If this enters a break loop, then exiting that break loop also exits \textsf{GAP}. 
\end{description}
 

 \index{options!command line, internal} \index{-C@\texttt{\texttt{\symbol{45}}C}} \index{-P@\texttt{\texttt{\symbol{45}}P}}\index{-p@\texttt{\texttt{\symbol{45}}p}} Additional options, such as \texttt{\texttt{\symbol{45}}C}, \texttt{\texttt{\symbol{45}}P} and\texttt{\texttt{\symbol{45}}p} are used internally by the \textsf{gac} script (see \ref{Kernel modules}) and/or on specific operating systems. }

  
\section{\textcolor{Chapter }{The gap.ini and gaprc files}}\label{sect:gap.ini}
\logpage{[ 3, 2, 0 ]}
\hyperdef{L}{X7FD66F977A3B02DF}{}
{
  When you start \textsf{GAP}, it looks for files with the names \texttt{gap.ini} and \texttt{gaprc} in its root directories (see \ref{GAP Root Directories}), and reads the first \texttt{gap.ini} and the first \texttt{gaprc} file it finds. These files are used for certain initializations, as follows. 

 The file \texttt{gap.ini} is read early in the startup process. Therefore, the parameters set in this
file can influence the startup process, such as which packages are
automatically loaded (see \texttt{LoadPackage} (\ref{LoadPackage})) and whether library files containing obsolete variables are read (see
Chapter \ref{Replaced and Removed Command Names}). On the other hand, only calls to a restricted set of \textsf{GAP} functions are allowed in a \texttt{gap.ini} file. Usually, it should only contain calls of \texttt{SetUserPreference} (\ref{SetUserPreference}). This file can be generated (or updated when new releases introduce further
user preferences) with the command \texttt{WriteGapIniFile} (\ref{WriteGapIniFile}). This file is read whenever \textsf{GAP} is started, with or without a workspace. 

 The file \texttt{gaprc} is read after the startup process, before the first input file given on the
command line (see \ref{Command Line Options}). So the contents of this file cannot influence the startup process, but all \textsf{GAP} library functions can be called in this file. When \textsf{GAP} is started with a workspace then the file is read only if no \texttt{gaprc} file had been read before the workspace was created. (With this setup, it is
on the one hand possible that administrators provide a \textsf{GAP} workspace for several users such that the user's \texttt{gaprc} file is read when \textsf{GAP} is started with the workspace, and on the other hand one can start \textsf{GAP}, read one's \texttt{gaprc} file, save a workspace, and then start from this workspace \emph{without} reading one's \texttt{gaprc} file again.) 

 Note that by default, the user specific \textsf{GAP} root directory \texttt{GAPInfo.UserGapRoot} is the first \textsf{GAP} root directory. So you can put your \texttt{gap.ini} and \texttt{gaprc} files into this directory. 

 This mechanism substitutes the much less flexible reading of a users \texttt{.gaprc} file in versions of \textsf{GAP} up to 4.4. For compatibility this \texttt{.gaprc} file is still read if the directory \texttt{GAPInfo.UserGapRoot} does not exist, see \ref{obsolete-gaprc} how to migrate your old setup.  
\subsection{\textcolor{Chapter }{The gap.ini file}}\label{subsect:gap.ini file}
\logpage{[ 3, 2, 1 ]}
\hyperdef{L}{X87DF11C885E73583}{}
{
  \index{gap.ini@\texttt{gap.ini}} The file \texttt{gap.ini} is read after the declaration part of the \textsf{GAP} library is read, before the declaration parts of the packages needed and
suggested by \textsf{GAP} are read, and before the implementation parts of \textsf{GAP} and of the packages are read. 

 The file \texttt{gap.ini} is expected to consist of calls to the function \texttt{SetUserPreference} (\ref{SetUserPreference}), see Section{\nobreakspace}\texttt{SetUserPreference} (\ref{SetUserPreference}). 

 Since the file \texttt{gap.ini} is read before the implementation part of \textsf{GAP} is read, not all \textsf{GAP} functions may be called in the file. Assignments of numbers, lists, and
records are admissible as well as calls to basic functions such as \texttt{Concatenation} (\ref{Concatenation:for several lists}) and \texttt{JoinStringsWithSeparator} (\ref{JoinStringsWithSeparator}). 

 Note that the file \texttt{gap.ini} is read also when \textsf{GAP} is started with a workspace. }

  
\subsection{\textcolor{Chapter }{The gaprc file}}\label{subsect:gaprc file}
\logpage{[ 3, 2, 2 ]}
\hyperdef{L}{X84D4CF587D437C00}{}
{
  If a file \texttt{gaprc} is found it is read after \textsf{GAP}'s \texttt{init.g}, but before any of the files mentioned on the command line are read. You can
use this file for your private customizations. (Many users may be happy with
using just user preferences in the \texttt{gap.ini} file (see above) for private customization.) For example, if you have a file
containing functions or data that you always need, you could read this from \texttt{gaprc}. Or if you find some of the names in the library too long, you could define
abbreviations for those names in \texttt{gaprc}. The following sample \texttt{gaprc} file does both. 

 
\begin{Verbatim}[commandchars=!@|,fontsize=\small,frame=single,label=Example]
  Read( "/usr/you/dat/mygroups.grp" );
  Ac := Action;
  AcHom := ActionHomomorphism;
  RepAc := RepresentativeAction;
\end{Verbatim}
 

 Note that only one \texttt{gaprc} file is read when \textsf{GAP} is started. When a workspace is created in a \textsf{GAP} session after a \texttt{gaprc} file has been read then no more \texttt{gaprc} file will be read when \textsf{GAP} is started with this workspace. 

 Also note that the file must be called \texttt{gaprc}. If you use a Windows text editor, in particular if your default is not to
show file suffixes, you might accidentally create a file \texttt{gaprc.txt} or \texttt{gaprc.doc} which GAP will not recognize. }

 
\subsection{\textcolor{Chapter }{Configuring User preferences}}\logpage{[ 3, 2, 3 ]}
\hyperdef{L}{X7B0AD104839B6C3C}{}
{
\noindent\textcolor{FuncColor}{$\triangleright$\enspace\texttt{SetUserPreference({\mdseries\slshape [package, ]name, value})\index{SetUserPreference@\texttt{SetUserPreference}}
\label{SetUserPreference}
}\hfill{\scriptsize (function)}}\\
\noindent\textcolor{FuncColor}{$\triangleright$\enspace\texttt{UserPreference({\mdseries\slshape [package, ]name})\index{UserPreference@\texttt{UserPreference}}
\label{UserPreference}
}\hfill{\scriptsize (function)}}\\
\noindent\textcolor{FuncColor}{$\triangleright$\enspace\texttt{ShowUserPreferences({\mdseries\slshape package1, package2, ...})\index{ShowUserPreferences@\texttt{ShowUserPreferences}}
\label{ShowUserPreferences}
}\hfill{\scriptsize (function)}}\\
\noindent\textcolor{FuncColor}{$\triangleright$\enspace\texttt{WriteGapIniFile({\mdseries\slshape [dir][,] [ignorecurrent]})\index{WriteGapIniFile@\texttt{WriteGapIniFile}}
\label{WriteGapIniFile}
}\hfill{\scriptsize (function)}}\\


 Some aspects of the behaviour of \textsf{GAP} can be customized by the user via \emph{user preferences}. Examples include the way help sections are displayed or the use of colors in
the terminal. 

 User preferences are specified via a pair of strings, the first is the (case
insensitive) name of a package (or \texttt{"GAP"} for the core \textsf{GAP} library) and the second is some arbitrary case sensitive string. 

 User preferences can be set to some \mbox{\texttt{\mdseries\slshape value}} with \texttt{SetUserPreference}. The current value of a user preference can be found with \texttt{UserPreference}. In both cases, if no package name is given the default \texttt{"GAP"} is used. If a user preference is not known or not set then \texttt{UserPreference} returns \texttt{fail}. 

 The stored values of user preferences are always immutable, see Section \ref{Mutability and Copyability}. 

 The function \texttt{ShowUserPreferences} with no argument shows in a pager an overview of all known user preferences
together with some explanation and the current value. If one or more strings \mbox{\texttt{\mdseries\slshape package1}}, ... are given then only the user preferences for these packages are shown.
The \textsf{Browse} package provides the function \texttt{BrowseUserPreferences} (\textbf{Browse: BrowseUserPreferences}) which gives an overview of the known user preferenes and also admits editing
the values of the preferences. 

 The easiest way to make use of user preferences is probably to use the
function \texttt{WriteGapIniFile}, usually without argument. This function creates a file \texttt{gap.ini} in your user specific \textsf{GAP} root directory (\texttt{GAPInfo.UserGapRoot}). If such a file already exists the function will make a backup of it first.
This newly created file contains descriptions of all known user preferences
and also calls of \texttt{SetUserPreference} for those user preferences which currently do not have their default value.
You can then edit that file to customize (further) the user preferences for
future \textsf{GAP} sessions. 

 Should a later version of \textsf{GAP} or some packages introduce new user preferences then you can call \texttt{WriteGapIniFile} again since it will set the previously known user preferences to their current
values. 

 Optionally, a different directory for the resulting \texttt{gap.ini} file can be specified as argument \mbox{\texttt{\mdseries\slshape dir}} to \texttt{WriteGapIniFile}. Another optional argument is the boolean value \texttt{true}, if this is given, the settings of all user preferences in the current
session are ignored. 

 Note that your \texttt{gap.ini} file is read by \textsf{GAP} very early during its startup process. A consequence is that the \mbox{\texttt{\mdseries\slshape value}} argument in a call of \texttt{SetUserPreference} must be some very basic \textsf{GAP} object, usually a boolean, a number, a string or a list of those. A few user
preferences support more complicated settings. For example, the user
preference \texttt{"UseColorPrompt"} admits a record as its value whose components are available only after the \textsf{GAPDoc} package has been loaded, see{\nobreakspace}\texttt{ColorPrompt} (\ref{ColorPrompt}). If you want to specify such a complicated value, then move the corresponding
call of \texttt{SetUserPreference} from your \texttt{gap.ini} file into your \texttt{gaprc} file (also in the directory \texttt{GAPInfo.UserGapRoot}). This file is read much later. 

 
\begin{Verbatim}[commandchars=!@|,fontsize=\small,frame=single,label=Example]
  !gapprompt@gap>| !gapinput@SetUserPreference( "Pager", "less" );|
  !gapprompt@gap>| !gapinput@SetUserPreference("PagerOptions",|
  !gapprompt@>| !gapinput@                     [ "-f", "-r", "-a", "-i", "-M", "-j2" ] );|
  !gapprompt@gap>| !gapinput@UserPreference("Pager");|
  "less"
\end{Verbatim}
 The first two lines of this example will cause \textsf{GAP} to use the programm \texttt{less} as a pager. This is highly recommended if \texttt{less} is available on your system. The last line displays the current setting. 

 }

 

\subsection{\textcolor{Chapter }{DeclareUserPreference}}
\logpage{[ 3, 2, 4 ]}\nobreak
\hyperdef{L}{X7F1DF6757B248014}{}
{\noindent\textcolor{FuncColor}{$\triangleright$\enspace\texttt{DeclareUserPreference({\mdseries\slshape record})\index{DeclareUserPreference@\texttt{DeclareUserPreference}}
\label{DeclareUserPreference}
}\hfill{\scriptsize (function)}}\\


 This function can be used (also in packages) to introduce new user
preferences. It declares a user preference, determines a default value and
contains documentation of the user preference. After declaration a user
preference will be shown with \texttt{ShowUserPreferences} (\ref{ShowUserPreferences}) and \texttt{WriteGapIniFile} (\ref{WriteGapIniFile}). 

 When this declaration is evaluated it is checked, if this user preference is
already set in the current session. If not the value of the user preference is
set to its default. (Do not use \texttt{fail} as default value since this indicated that a user preference is not set.)

 The argument \mbox{\texttt{\mdseries\slshape record}} of \texttt{DeclareUserPreference} must be a record with the following components. 

 
\begin{description}
\item[{\texttt{name}}]  a string or a list of strings, the latter meaning several preferences which
belong together, 
\item[{\texttt{description}}]  a list of strings describing the preference(s), one string for each paragraph;
if several preferences are declared together then the description refers to
all of them, 
\item[{\texttt{default}}]  the default value that is used, or a function without arguments that computes
this default value; if several preferences are declared together then the
value of this component must be the list of default values for the individual
preferences. 
\end{description}
 

 The following components of \mbox{\texttt{\mdseries\slshape record}} are optional. 

 
\begin{description}
\item[{\texttt{check}}]  a function that takes a value as its argument and returns either \texttt{true} or \texttt{false}, depending on whether the given value is admissible for this preference; if
several preferences are declared together then the number of arguments of the
function must equal the length of the \texttt{name} list, 
\item[{\texttt{values}}]  the list of admissible values, or a function without arguments that returns
this list, 
\item[{\texttt{multi}}]  \texttt{true} or \texttt{false}, depending on whether one may choose several values from the given list or
just one; needed (and useful only) if the \texttt{values} component is present, 
\item[{\texttt{package}}]  the name of the \textsf{GAP} package to which the preference is assigned; if the declaration happens inside
a file that belongs to this package then the value of this component is
computed, using \texttt{GAPInfo.PackageCurrent}; otherwise, the default value for \texttt{package} is \texttt{"GAP"}, 
\item[{\texttt{omitFromGapIniFile}}]  if the value is \texttt{true} then this user preference is ignored by \texttt{WriteGapIniFile} (\ref{WriteGapIniFile}). 
\end{description}
 

 
\begin{Verbatim}[commandchars=!@|,fontsize=\small,frame=single,label=Example]
  !gapprompt@gap>| !gapinput@UserPreference( "MyFavouritePrime" );|
  fail
  !gapprompt@gap>| !gapinput@DeclareUserPreference( rec(|
  !gapprompt@>| !gapinput@       name:= "MyFavouritePrime",|
  !gapprompt@>| !gapinput@       description:= [ "is not used, serves as an example" ],|
  !gapprompt@>| !gapinput@       default:= 2,|
  !gapprompt@>| !gapinput@       omitFromGapIniFile:= true ) );|
  !gapprompt@gap>| !gapinput@UserPreference( "MyFavouritePrime" );|
  2
  !gapprompt@gap>| !gapinput@SetUserPreference( "MyFavouritePrime", 17 );|
  !gapprompt@gap>| !gapinput@UserPreference( "MyFavouritePrime" );|
  17
\end{Verbatim}
 }

  
\subsection{\textcolor{Chapter }{User Preferences Defined by \textsf{GAP}}}\label{subsect:User Preferences Defined by GAP}
\logpage{[ 3, 2, 5 ]}
\hyperdef{L}{X870A11E7864F9CA7}{}
{
  Here is the list of those user preferences that are currently declared via \texttt{DeclareUserPreference} (\ref{DeclareUserPreference}) for \textsf{GAP} itself. The preferences that are declared for \textsf{GAP} packages belong to the documentation of these packages. 
\begin{description}
\item[{ \index{Autocompleter@\texttt{Autocompleter}} \texttt{Autocompleter} }]  Set how names are filtered during tab\texttt{\symbol{45}}autocomplete, this
can be: \texttt{"default"}: case\texttt{\symbol{45}}sensitive matching. \texttt{"case\texttt{\symbol{45}}insensitive"}: case\texttt{\symbol{45}}insensitive matching, or a record with two
components named \texttt{filter} and \texttt{completer}, which are both functions which take two arguments. \texttt{filter} takes a list of names and a partial identifier and returns all the members of \texttt{names} which are a valid extension of the partial identifier. \texttt{completer} takes a list of names and a partial identifier and returns the partial
identifier as extended as possible (it may also change the identifier, for
example to correct the case, or spelling mistakes), or returns \texttt{fail} to leave the existing partial identifier. 

 This preference is ignored if \textsf{GAP} was not compiled with readline support. 

 Default: \texttt{"default"}. 
\item[{ \index{Editor@\texttt{Editor}} \index{EditorOptions@\texttt{EditorOptions}} \texttt{Editor}, \texttt{EditorOptions} }]  Determines the editor and options (used by \textsf{GAP}'s \texttt{Edit} (\ref{Edit}) command). Under macOS, the value \texttt{"open"} for \texttt{Editor} will work. For further options, see the \textsf{GAP} help for \texttt{Edit} (\ref{Edit}). If you want to use the editor defined in your (shell) environment then leave
the \texttt{Editor} and \texttt{EditorOptions} preferences empty. 

 The defaults are computed at runtime. 
\item[{ \index{ExcludeFromAutoload@\texttt{ExcludeFromAutoload}} \texttt{ExcludeFromAutoload} }]  These packages are not loaded at \textsf{GAP} startup. This doesn't work for packages which are needed by the \textsf{GAP} library, or which are already loaded in a workspace. 

 Default: \texttt{""}. 
\item[{ \index{HelpViewers@\texttt{HelpViewers}} \index{XpdfOptions@\texttt{XpdfOptions}} \index{XdviOptions@\texttt{XdviOptions}} \texttt{HelpViewers}, \texttt{XpdfOptions}, \texttt{XdviOptions} }]  Here you can choose your preferred help viewers. See the help for \texttt{SetHelpViewer} (\ref{SetHelpViewer}) for further options. 

 Try \texttt{HelpViewers:= [ "screen", "firefox", "xpdf" ];}. 

 (For \texttt{"screen"} we also suggest to set the \texttt{Pager} entry to \texttt{"less"}.) 

 Defaults: \texttt{[ [ "screen" ], "", "" ]}. 
\item[{ \index{HistoryBackwardSearchSkipIdenticalEntries@\texttt{HistoryBackwardSearchSkipIdenticalEntries}} \texttt{HistoryBackwardSearchSkipIdenticalEntries} }]  When a command is executed multiple times, it is also stored in history
multiple times. Setting this option to \texttt{true} skips identical entries when searching backwards in history. 

 Admissible values: \texttt{true}, \texttt{false}. 

 Default: \texttt{false}. 
\item[{ \index{HistoryMaxLines@\texttt{HistoryMaxLines}} \index{SaveAndRestoreHistory@\texttt{SaveAndRestoreHistory}} \texttt{HistoryMaxLines}, \texttt{SaveAndRestoreHistory} }]  \texttt{HistoryMaxLines} is the maximal amount of input lines held in \textsf{GAP}'s command line history. 

 If \texttt{SaveAndRestoreHistory} is \texttt{true} then \textsf{GAP} saves its command line history before terminating a \textsf{GAP} session, and prepends the stored history when \textsf{GAP} is started. If this is enabled it is suggested to set \texttt{HistoryMaxLines} to some finite value. It is also possible to set \texttt{HistoryMaxLines} to \texttt{infinity} (\ref{infinity}) to keep arbitrarily many lines. 

 These preferences are ignored if \textsf{GAP} was not compiled with readline support. 

 Defaults: \texttt{[ 10000, true ]}. 
\item[{ \index{InfoPackageLoadingLevel@\texttt{InfoPackageLoadingLevel}} \texttt{InfoPackageLoadingLevel} }]  Info messages concerning package loading up to this level are printed. The
level can be changed in a running session using \texttt{SetInfoLevel} (\ref{SetInfoLevel}). 

 Admissible values: \texttt{1}, \texttt{2}, \texttt{3}, \texttt{4}. 

 Default: \texttt{1}. 
\item[{ \index{MaxBitsIntView@\texttt{MaxBitsIntView}} \texttt{MaxBitsIntView} }]  Maximal bit length of integers to \texttt{View} unabbreviated. Default is about $30$ lines of a $80$ character wide terminal. Set this to \texttt{0} to avoid abbreviated ints. 

 Default: \texttt{8000}. 
\item[{ \index{PartialPermDisplayLimit@\texttt{PartialPermDisplayLimit}} \index{NotationForPartialPerms@\texttt{NotationForPartialPerms}} \texttt{PartialPermDisplayLimit}, \texttt{NotationForPartialPerms} }]  options for the display of partial perms 

 Defaults: \texttt{[ 100, "component" ]}. 
\item[{ \index{TransformationDisplayLimit@\texttt{TransformationDisplayLimit}} \index{NotationForTransformations@\texttt{NotationForTransformations}} \texttt{TransformationDisplayLimit}, \texttt{NotationForTransformations} }]  options for the display of transformations 

 Defaults: \texttt{[ 100, "input" ]}. 
\item[{ \index{PackagesToIgnore@\texttt{PackagesToIgnore}} \texttt{PackagesToIgnore} }]  These packages are not regarded as available. This doesn't work for packages
which are needed by the \textsf{GAP} library, or which are already loaded in a workspace. 

 Default: \texttt{""}. 
\item[{ \index{PackagesToLoad@\texttt{PackagesToLoad}} \texttt{PackagesToLoad} }]  A list of names of packages which should be loaded during startup. For
backwards compatibility, the default lists most of packages that were
autoloaded in \textsf{GAP} 4.4 (add or remove packages as you like). 

 Default: \texttt{[ "autpgrp", "alnuth", "crisp", "ctbllib", "factint", "fga", "irredsol",
"laguna", "PackageManager", "polenta", "polycyclic", "resclasses", "sophus",
"tomlib" ]}. 
\item[{ \index{Pager@\texttt{Pager}} \index{PagerOptions@\texttt{PagerOptions}} \texttt{Pager}, \texttt{PagerOptions} }]  For displaying help pages on screen and other things \textsf{GAP} has a rudimentary builtin pager. We recommend using a more sophisticated
external program. For example, when you have the program \texttt{less} on your computer we recommend: 

 \texttt{Pager := "less";} 

 \texttt{PagerOptions := ["\texttt{\symbol{45}}f", "\texttt{\symbol{45}}r",
"\texttt{\symbol{45}}a", "\texttt{\symbol{45}}i", "\texttt{\symbol{45}}M",
"\texttt{\symbol{45}}j2"];} 

 If you want to use \texttt{more}, we suggest to use the \texttt{\texttt{\symbol{45}}f} option. If you want to use the pager defined in your environment then leave
the \texttt{Pager} and \texttt{PagerOptions} preferences empty. 

 The defaults are computed at runtime. 
\item[{ \index{ReadObsolete@\texttt{ReadObsolete}} \texttt{ReadObsolete} }]  May be useful to say \texttt{false} here to check if you are using commands which may vanish in a future version
of \textsf{GAP} 

 Admissible values: \texttt{true}, \texttt{false}. 

 Default: \texttt{true}. 
\item[{ \index{ReproducibleBehaviour@\texttt{ReproducibleBehaviour}} \texttt{ReproducibleBehaviour} }]  This preference disables code in \textsf{GAP} which changes behaviour based on time spent, and therefore can produce
different results depending on how much time is taken by other programs
running on the same computer. This option may lead to slower or
lower\texttt{\symbol{45}}quality results. Note that many algorithms in \textsf{GAP} use the global random number generator, which is NOT affected by this option.
This only tries to ensure the same version of \textsf{GAP}, with the same package versions loaded, on the same machine, running the same
code, in a fresh \textsf{GAP} session, will produce the same results. 

 Admissible values: \texttt{true}, \texttt{false}. 

 Default: \texttt{false}. 
\item[{ \index{ShortBanners@\texttt{ShortBanners}} \texttt{ShortBanners} }]  If this option is set to \texttt{true}, package banners printed during loading will only show the name, version and
description of a package. 

 Admissible values: \texttt{true}, \texttt{false}. 

 Default: \texttt{false}. 
\item[{ \index{UseColorPrompt@\texttt{UseColorPrompt}} \texttt{UseColorPrompt} }]  In a color capable terminal (almost any terminal application) you can run \textsf{GAP} such that the prompts, the input and output are distinguished by colors.
Options are \texttt{true}, \texttt{false} or some record as explained in the help section for \texttt{ColorPrompt} (\ref{ColorPrompt}). 

 Default: \texttt{true}. 
\item[{ \index{UseColorsInTerminal@\texttt{UseColorsInTerminal}} \texttt{UseColorsInTerminal} }]  Almost all current terminal emulations support color display, setting this to \texttt{true} implies a default display of most manuals with color markup. It may influence
the display of other things in the future. 

 Admissible values: \texttt{true}, \texttt{false}. 

 Default: \texttt{true}. 
\item[{ \index{ViewLength@\texttt{ViewLength}} \texttt{ViewLength} }]  A bound for the number of lines printed when \texttt{View}ing some large objects. 

 Default: \texttt{3}. 
\item[{ \index{WhereDepth@\texttt{WhereDepth}} \texttt{WhereDepth} }]  The number of stack frames shown by \texttt{Where} and \texttt{WhereWithVars} when called without an explicit depth argument, e.g. in the default \texttt{OnBreak} handler. 

 Default: \texttt{5}. 
\end{description}
 }

 }

  
\section{\textcolor{Chapter }{Saving and Loading a Workspace}}\label{Saving and Loading a Workspace}
\logpage{[ 3, 3, 0 ]}
\hyperdef{L}{X7CB282757ACB1C09}{}
{
  \textsf{GAP} workspace files are binary files that contain the data of a \textsf{GAP} session. Currently saving and loading workspace files are supported only when
the \texttt{GASMAN} garbage collector is used, see Section \ref{Garbage Collection}. 

 One can produce a workspace file with \texttt{SaveWorkspace} (\ref{SaveWorkspace}), and load it into a new \textsf{GAP} session using the \texttt{\texttt{\symbol{45}}L} command line option, see Section \ref{Command Line Options}. 

 One purpose of workspace files is of course the possibility to save a ``snapshot'' image of the current \textsf{GAP} workspace in a file. 

 The recommended way to start \textsf{GAP} is to load an existing workspace file, because this reduces the startup time
of \textsf{GAP} drastically. So if you have installed \textsf{GAP} yourself then you should think about creating a workspace file immediately
after you have started \textsf{GAP}, and then using this workspace file later on, whenever you start \textsf{GAP}. If your \textsf{GAP} installation is shared between several users, the system administrator should
think about providing such a workspace file. 

\subsection{\textcolor{Chapter }{SaveWorkspace}}
\logpage{[ 3, 3, 1 ]}\nobreak
\hyperdef{L}{X876544A57C73C488}{}
{\noindent\textcolor{FuncColor}{$\triangleright$\enspace\texttt{SaveWorkspace({\mdseries\slshape filename})\index{SaveWorkspace@\texttt{SaveWorkspace}}
\label{SaveWorkspace}
}\hfill{\scriptsize (function)}}\\


 \index{save} will save a ``snapshot'' image of the current \textsf{GAP} workspace in the file \mbox{\texttt{\mdseries\slshape filename}}. This image then can be loaded by another copy of \textsf{GAP} which then will behave as at the point when \texttt{SaveWorkspace} was called. 

 
\begin{Verbatim}[commandchars=!@|,fontsize=\small,frame=single,label=Example]
  !gapprompt@gap>| !gapinput@a:=1;|
  !gapprompt@gap>| !gapinput@SaveWorkspace("savefile");|
  true
  !gapprompt@gap>| !gapinput@quit;|
\end{Verbatim}
 

 \texttt{SaveWorkspace} can only be used at the main \texttt{gap{\textgreater}} prompt. It cannot be included in the body of a loop or function, or called
from a break loop. }

 }

  
\section{\textcolor{Chapter }{Testing for the System Architecture}}\label{Testing for the System Architecture}
\logpage{[ 3, 4, 0 ]}
\hyperdef{L}{X83BF07587F2CC6CD}{}
{
  

\subsection{\textcolor{Chapter }{ARCH{\textunderscore}IS{\textunderscore}UNIX}}
\logpage{[ 3, 4, 1 ]}\nobreak
\hyperdef{L}{X7C825AF087A27884}{}
{\noindent\textcolor{FuncColor}{$\triangleright$\enspace\texttt{ARCH{\textunderscore}IS{\textunderscore}UNIX({\mdseries\slshape })\index{ARCH{\textunderscore}IS{\textunderscore}UNIX@\texttt{ARC}\-\texttt{H{\textunderscore}}\-\texttt{I}\-\texttt{S{\textunderscore}}\-\texttt{UNIX}}
\label{ARCHuScoreISuScoreUNIX}
}\hfill{\scriptsize (function)}}\\


 tests whether \textsf{GAP} is running on a UNIX system (including macOS). }

 

\subsection{\textcolor{Chapter }{ARCH{\textunderscore}IS{\textunderscore}MAC{\textunderscore}OS{\textunderscore}X}}
\logpage{[ 3, 4, 2 ]}\nobreak
\hyperdef{L}{X82A6893A7EC8FA72}{}
{\noindent\textcolor{FuncColor}{$\triangleright$\enspace\texttt{ARCH{\textunderscore}IS{\textunderscore}MAC{\textunderscore}OS{\textunderscore}X({\mdseries\slshape })\index{ARCH{\textunderscore}IS{\textunderscore}MAC{\textunderscore}OS{\textunderscore}X@\texttt{ARC}\-\texttt{H{\textunderscore}}\-\texttt{I}\-\texttt{S{\textunderscore}}\-\texttt{M}\-\texttt{A}\-\texttt{C{\textunderscore}}\-\texttt{O}\-\texttt{S{\textunderscore}X}}
\label{ARCHuScoreISuScoreMACuScoreOSuScoreX}
}\hfill{\scriptsize (function)}}\\


 tests whether \textsf{GAP} is running on macOS. Note that on macOS, also \texttt{ARCH{\textunderscore}IS{\textunderscore}UNIX} (\ref{ARCHuScoreISuScoreUNIX}) will be \texttt{true}. }

 

\subsection{\textcolor{Chapter }{ARCH{\textunderscore}IS{\textunderscore}WINDOWS}}
\logpage{[ 3, 4, 3 ]}\nobreak
\hyperdef{L}{X7A14B659847B8627}{}
{\noindent\textcolor{FuncColor}{$\triangleright$\enspace\texttt{ARCH{\textunderscore}IS{\textunderscore}WINDOWS({\mdseries\slshape })\index{ARCH{\textunderscore}IS{\textunderscore}WINDOWS@\texttt{ARC}\-\texttt{H{\textunderscore}}\-\texttt{I}\-\texttt{S{\textunderscore}}\-\texttt{W}\-\texttt{I}\-\texttt{N}\-\texttt{DOWS}}
\label{ARCHuScoreISuScoreWINDOWS}
}\hfill{\scriptsize (function)}}\\


 tests whether \textsf{GAP} is running on a Windows system without standard POSIX tools available (such as
a shell). }

 

\subsection{\textcolor{Chapter }{ARCH{\textunderscore}IS{\textunderscore}WSL}}
\logpage{[ 3, 4, 4 ]}\nobreak
\hyperdef{L}{X87E7CC3B8395BBB3}{}
{\noindent\textcolor{FuncColor}{$\triangleright$\enspace\texttt{ARCH{\textunderscore}IS{\textunderscore}WSL({\mdseries\slshape })\index{ARCH{\textunderscore}IS{\textunderscore}WSL@\texttt{ARC}\-\texttt{H{\textunderscore}}\-\texttt{I}\-\texttt{S{\textunderscore}WSL}}
\label{ARCHuScoreISuScoreWSL}
}\hfill{\scriptsize (function)}}\\


 tests whether \textsf{GAP} is running on a Windows system inside the 'Windows Subsystem for Linux'. Note
that in this case \texttt{ARCH{\textunderscore}IS{\textunderscore}UNIX} (\ref{ARCHuScoreISuScoreUNIX}) will be \texttt{true}, and in most situations WSL can be treated identically to Linux. }

 }

  
\section{\textcolor{Chapter }{Global Values that Control the \textsf{GAP} Session}}\label{GAPInfo_sect}
\logpage{[ 3, 5, 0 ]}
\hyperdef{L}{X8719B2118511645F}{}
{
  

\subsection{\textcolor{Chapter }{GAPInfo}}
\logpage{[ 3, 5, 1 ]}\nobreak
\hyperdef{L}{X8354754E7935F935}{}
{\noindent\textcolor{FuncColor}{$\triangleright$\enspace\texttt{GAPInfo\index{GAPInfo@\texttt{GAPInfo}}
\label{GAPInfo}
}\hfill{\scriptsize (global variable)}}\\


 Several global values control the \textsf{GAP} session, such as the command line, the architecture, or the information about
available and loaded packages. Many of these values are accessible as
components of the global record \texttt{GAPInfo}. Typically, these components are set and read in low level \textsf{GAP} functions, so changing the values of existing components of \texttt{GAPInfo} ``by hand'' is not recommended. 

 Important components are documented via index entries, try the input \texttt{??GAPInfo} for getting an overview of these components. }

  }

  
\section{\textcolor{Chapter }{Coloring the Prompt and Input}}\label{Coloring the Prompt and Input}
\logpage{[ 3, 6, 0 ]}
\hyperdef{L}{X818F2DDC863C381E}{}
{
  \textsf{GAP} provides hooks for functions which are called when the prompt is to be printed
and when an input line is finished. 

 An example of using this feature is the following function. 

\subsection{\textcolor{Chapter }{ColorPrompt}}
\logpage{[ 3, 6, 1 ]}\nobreak
\hyperdef{L}{X84F3481C8466C7FC}{}
{\noindent\textcolor{FuncColor}{$\triangleright$\enspace\texttt{ColorPrompt({\mdseries\slshape bool[, optrec]})\index{ColorPrompt@\texttt{ColorPrompt}}
\label{ColorPrompt}
}\hfill{\scriptsize (function)}}\\


 \texttt{ColorPrompt} changes \textsf{GAP}'s user interface: After calling \texttt{ColorPrompt(true);}, the prompts and the user input are displayed in colors different from the
color that is used for the output. This is also the default for a \textsf{GAP} session. Switch off these colorings with \texttt{ColorPrompt(false);}. 

 Note that colors will only work if your terminal emulation in which you run \textsf{GAP} understands the so called ANSI color escape sequences {\textendash}almost all
terminal emulations on current UNIX/Linux (\texttt{xterm}, \texttt{rxvt}, \texttt{konsole}, ...) systems do so. 

 The colors shown depend on the terminal configuration and cannot be forced
from an application. If your terminal follows the ANSI conventions you see the
standard prompt in bold blue and the break loop prompt in bold red, as well as
your input in red. 

 If you prefer to switch off colors for prompts and input at the start of your \textsf{GAP} sessions, put a call of \texttt{SetUserPreference("UseColorPrompt", false);} in your \texttt{gap.ini} file. If you want a more complicated setting as explained below then put your \texttt{SetUserPreference("UseColorPrompt", rec( ... ) );} call into your \texttt{gaprc} file. 

 The optional second argument \mbox{\texttt{\mdseries\slshape optrec}} allows one to further customize the behaviour. It must be a record from which
the following components are recognized: 

 
\begin{description}
\item[{\texttt{MarkupStdPrompt}}]  a string or no argument function returning a string containing the escape
sequence used for the main prompt \texttt{gap{\textgreater} }. 
\item[{\texttt{MarkupContPrompt}}]  a string or no argument function returning a string containing the escape
sequence used for the continuation prompt \texttt{{\textgreater} }. 
\item[{\texttt{MarkupBrkPrompt}}]  a string or no argument function returning a string containing the escape
sequence used for the break prompt \texttt{brk...{\textgreater} }. 
\item[{\texttt{MarkupInput}}]  a string or no argument function returning a string containing the escape
sequence used for user input. 
\item[{\texttt{TextPrompt}}]  a no argument function returning the string with the text of the prompt, but
without any escape sequences. The current standard prompt is returned by \texttt{CPROMPT()}. But note that changing the standard prompts makes the automatic removal of
prompts from input lines impossible (see{\nobreakspace}\ref{Special Rules for Input Lines}). 
\item[{\texttt{PrePrompt}}]  a function called before printing a prompt. 
\end{description}
 

 Here is an example. 

 
\begin{Verbatim}[commandchars=!@|,fontsize=\small,frame=single,label=]
  LoadPackage("GAPDoc");
  timeSHOWMIN := 100;
  ColorPrompt(true, rec(
     # usually cyan bold, see ?TextAttr
     MarkupStdPrompt := Concatenation(TextAttr.bold, TextAttr.6),
     MarkupContPrompt := Concatenation(TextAttr.bold, TextAttr.6),
     PrePrompt := function()
       # show the 'time' automatically if at least timeSHOWMIN
       if CPROMPT() = "gap> " and time >= timeSHOWMIN then
         Print("Time of last command: ", time, " ms\n");
       fi;
     end)    );
\end{Verbatim}
 }

 }

 }

          
\chapter{\textcolor{Chapter }{The Programming Language}}\label{The Programming Language}
\logpage{[ 4, 0, 0 ]}
\hyperdef{L}{X7FE7C0C17E1ED118}{}
{
  This chapter describes the \textsf{GAP} programming language. It should allow you, in principle, to predict the result
of each and every input. In order to know what we are talking about, we first
have to look more closely at the process of interpretation and the various
representations of data involved.    
\section{\textcolor{Chapter }{Language Overview}}\label{Language Overview}
\logpage{[ 4, 1, 0 ]}
\hyperdef{L}{X7B5FF6827DFBDF20}{}
{
  First we have the input to \textsf{GAP}, given as a string of characters. How those characters enter \textsf{GAP} is operating system dependent, e.g., they might be entered at a terminal,
pasted with a mouse into a window, or read from a file. The mechanism does not
matter. This representation of expressions by characters is called the \emph{external representation} of the expression. Every expression has at least one external representation
that can be entered to get exactly this expression. 

 The input, i.e., the external representation, is transformed in a process
called \emph{reading} to an internal representation. At this point the input is analyzed and inputs
that are not legal external representations, according to the rules given
below, are rejected as errors. Those rules are usually called the \emph{syntax} of a programming language. 

 The internal representation created by reading is called either an \emph{expression} or a \emph{statement}. Later we will distinguish between those two terms. However for now we will
use them interchangeably. The exact form of the internal representation does
not matter. It could be a string of characters equal to the external
representation, in which case the reading would only need to check for errors.
It could be a series of machine instructions for the processor on which \textsf{GAP} is running, in which case the reading would more appropriately be called
compilation. It is in fact a tree\texttt{\symbol{45}}like structure. 

 After the input has been read it is again transformed in a process called \emph{evaluation} or \emph{execution}. Later we will distinguish between those two terms too, but for the moment we
will use them interchangeably. The name hints at the nature of this process,
it replaces an expression with the value of the expression. This works
recursively, i.e., to evaluate an expression first the subexpressions are
evaluated and then the value of the expression is computed from those values
according to rules given below. Those rules are usually called the \emph{semantics} of a programming language. 

 The result of the evaluation is, not surprisingly, called a \emph{value}.     Again the form in which such a value is represented internally does not
matter. It is in fact a tree\texttt{\symbol{45}}like structure again. 

 The last process is called \emph{printing}. It takes the value produced by the evaluation and creates an external
representation, i.e., a string of characters again. What you do with this
external representation is up to you. You can look at it, paste it with the
mouse into another window, or write it to a file. 

 Lets look at an example to make this more clear. Suppose you type in the
following string of 8 characters 

 
\begin{Verbatim}[commandchars=!@|,fontsize=\small,frame=single,label=]
  1 + 2 * 3;
\end{Verbatim}
 

 \textsf{GAP} takes this external representation and creates a tree\texttt{\symbol{45}}like
internal representation, which we can picture as follows 

 
\begin{Verbatim}[commandchars=!@|,fontsize=\small,frame=single,label=]
    +
   / \
  1   *
     / \
    2   3
\end{Verbatim}
 

 This expression is then evaluated. To do this \textsf{GAP} first evaluates the right subexpression \texttt{2*3}. Again, to do this \textsf{GAP} first evaluates its subexpressions 2 and 3. However they are so simple that
they are their own value, we say that they are
self\texttt{\symbol{45}}evaluating. After this has been done, the rule for \texttt{*} tells us that the value is the product of the values of the two
subexpressions, which in this case is clearly 6. Combining this with the value
of the left operand of the \texttt{+}, which is self\texttt{\symbol{45}}evaluating, too, gives us the value of the
whole expression 7. This is then printed, i.e., converted into the external
representation consisting of the single character \texttt{7}. 

 In this fashion we can predict the result of every input when we know the
syntactic rules that govern the process of reading and the semantic rules that
tell us for every expression how its value is computed in terms of the values
of the subexpressions. The syntactic rules are given in sections \ref{Lexical Structure}, \ref{Symbols}, \ref{Whitespaces}, \ref{Keywords}, and \ref{Identifiers}, the semantic rules are given in sections \ref{Expressions}, \ref{Variables}, \ref{Function Calls}, \ref{Comparisons}, \ref{Arithmetic Operators}, \ref{Statements}, \ref{Assignments}, \ref{Procedure Calls}, \ref{If}, \ref{While}, \ref{Repeat}, \ref{For}, \ref{Function}, and the chapters describing the individual data types. }

  
\section{\textcolor{Chapter }{Lexical Structure}}\label{Lexical Structure}
\logpage{[ 4, 2, 0 ]}
\hyperdef{L}{X80A85A707B6F4BE7}{}
{
  Most input of \textsf{GAP} consists of sequences of the following characters. 

 Digits, uppercase and lowercase letters, \textsc{Space}, \textsc{Tab}, \textsc{Newline}, \textsc{Return} and the special characters 

 
\begin{Verbatim}[commandchars=@|A,fontsize=\small,frame=single,label=]
  "    '    (    )    *    +    ,    -    #
  .    /    :    ;    <    =    >    ~
  [    \    ]    ^    _    {    }    !
\end{Verbatim}
 

 It is possible to use other characters in identifiers by escaping them with
backslashes, but we do not recommend the use of this feature. Inside strings
(see section{\nobreakspace}\ref{Symbols} and chapter{\nobreakspace}\ref{Strings and Characters}) and comments (see{\nobreakspace}\ref{Whitespaces}) the full character set supported by the computer is allowed. }

  
\section{\textcolor{Chapter }{Symbols}}\label{Symbols}
\logpage{[ 4, 3, 0 ]}
\hyperdef{L}{X7E90E6607F4E4943}{}
{
  The process of reading, i.e., of assembling the input into expressions, has a
subprocess, called \emph{scanning}, that assembles the characters into symbols. A \emph{symbol} is a sequence of characters that form a lexical unit. The set of symbols
consists of keywords, identifiers, strings, integers, and operator and
delimiter symbols. 

 A \emph{keyword} is a reserved word (see \ref{Keywords}). An \emph{identifier} is a sequence of letters, digits and underscores (or other characters escaped
by backslashes) that contains at least one non\texttt{\symbol{45}}digit and is
not a keyword (see \ref{Identifiers}). An integer is a sequence of digits (see \ref{Integers}), possibly prepended by \texttt{\texttt{\symbol{45}}} and \texttt{+} sign characters. A \emph{string} is a sequence of arbitrary characters enclosed in double quotes (see \ref{Strings and Characters}). 

 Operator and delimiter symbols are 

 
\begin{Verbatim}[commandchars=@|A,fontsize=\small,frame=single,label=]
  +    -    *    /    ^    ~   !.
  =    <>   <    <=   >    >=  ![
  :=   .    ..   ->   ,    ;   [
  ]    {    }    (    )    :
\end{Verbatim}
 

 Note also that during the process of scanning all whitespace is removed (see \ref{Whitespaces}). }

  
\section{\textcolor{Chapter }{Whitespaces}}\label{Whitespaces}
\logpage{[ 4, 4, 0 ]}
\hyperdef{L}{X7C53CEFC8641B919}{}
{
  \index{space} \index{blank}\index{tabulator}\index{newline} \index{comments}\index{\#} The characters \textsc{Space}, \textsc{Tab}, \textsc{Newline}, and \textsc{Return} are called \emph{whitespace characters}. Whitespace is used as necessary to separate lexical symbols, such as
integers, identifiers, or keywords. For example \texttt{Thorondor} is a single identifier, while \texttt{Th or ondor} is the keyword \texttt{or} between the two identifiers \texttt{Th} and \texttt{ondor}. Whitespace may occur between any two symbols, but not within a symbol. Two
or more adjacent whitespace characters are equivalent to a single whitespace.
Apart from the role as separator of symbols, whitespace characters are
otherwise insignificant. Whitespace characters may also occur inside a string,
where they are significant. Whitespace characters should also be used freely
for improved readability. 

 A \emph{comment} starts with the character \texttt{\#}, which is sometimes called sharp or hatch, and continues to the end of the
line on which the comment character appears. The whole comment, including \texttt{\#} and the \textsc{Newline} character is treated as a single whitespace. Inside a string, the comment
character \texttt{\#} loses its role and is just an ordinary character. 

 For example, the following statement 

 
\begin{Verbatim}[commandchars=!@|,fontsize=\small,frame=single,label=]
  if i<0 then a:=-i;else a:=i;fi;
\end{Verbatim}
 

 is equivalent to 

 
\begin{Verbatim}[commandchars=!@|,fontsize=\small,frame=single,label=]
  if i < 0 then   # if i is negative
    a := -i;      #   take its additive inverse
  else            # otherwise
    a := i;       #   take itself
  fi;
\end{Verbatim}
 

 (which by the way shows that it is possible to write superfluous comments).
However the first statement is \emph{not} equivalent to 

 
\begin{Verbatim}[commandchars=!@|,fontsize=\small,frame=single,label=]
  ifi<0thena:=-i;elsea:=i;fi;
\end{Verbatim}
 

 since the keyword \texttt{if} must be separated from the identifier \texttt{i} by a whitespace, and similarly \texttt{then} and \texttt{a}, and \texttt{else} and \texttt{a} must be separated. }

  
\section{\textcolor{Chapter }{Keywords}}\label{Keywords}
\logpage{[ 4, 5, 0 ]}
\hyperdef{L}{X87506BDC7D5F789E}{}
{
  \index{GAPInfo.Keywords@\texttt{GAPInfo.Keywords}} \emph{Keywords} are reserved words that are used to denote special operations or are part of
statements. They must not be used as identifiers. The list of keywords is
contained in the \texttt{GAPInfo.Keywords} component of the \texttt{GAPInfo} record (see \ref{GAPInfo}). We will show how to print it in a nice table, demonstrating at the same
time some list manipulation techniques: 

 
\begin{Verbatim}[commandchars=!@|,fontsize=\small,frame=single,label=Example]
  !gapprompt@gap>| !gapinput@keys:=SortedList( GAPInfo.Keywords );; l:=Length( keys );;|
  !gapprompt@gap>| !gapinput@arr:= List( [ 0 .. Int( l/4 )-1 ], i-> keys{ 4*i + [ 1 .. 4 ] } );;|
  !gapprompt@gap>| !gapinput@if l mod 4 <> 0 then Add( arr, keys{[ 4*Int(l/4) + 1 .. l ]} ); fi;|
  !gapprompt@gap>| !gapinput@Length( keys ); PrintArray( arr );|
  35
  [ [         Assert,           Info,        IsBound,           QUIT ],
    [  TryNextMethod,         Unbind,            and,         atomic ],
    [          break,       continue,             do,           elif ],
    [           else,            end,          false,             fi ],
    [            for,       function,             if,             in ],
    [          local,            mod,            not,             od ],
    [             or,           quit,       readonly,      readwrite ],
    [            rec,         repeat,         return,           then ],
    [           true,          until,          while ] ]
\end{Verbatim}
 

 Note that (almost) all keywords are written in lowercase and that they are
case sensitive. For example \texttt{else} is a keyword; \texttt{Else}, \texttt{eLsE}, \texttt{ELSE} and so forth are ordinary identifiers. Keywords must not contain whitespace,
for example \texttt{el if} is not the same as \texttt{elif}. 

 \emph{Note}: Several tokens from the list of keywords above may appear to be normal
identifiers representing functions or literals of various kinds but are
actually implemented as keywords for technical reasons. The only consequence
of this is that those identifiers cannot be re\texttt{\symbol{45}}assigned,
and do not actually have function objects bound to them, which could be
assigned to other variables or passed to functions. These keywords are \texttt{true}, \texttt{false}, \texttt{Assert} (\ref{Assert}), \texttt{IsBound} (\ref{IsBound:for a global variable}), \texttt{Unbind} (\ref{Unbind:unbind a variable}), \texttt{Info} (\ref{Info}) and \texttt{TryNextMethod} (\ref{TryNextMethod}). 

  Keywords \texttt{atomic}, \texttt{readonly}, \texttt{readwrite} are not used at the moment. They are reserved for the future version of GAP to
prevent their accidental use as identifiers. }

  
\section{\textcolor{Chapter }{Identifiers}}\label{Identifiers}
\logpage{[ 4, 6, 0 ]}
\hyperdef{L}{X860313A179A5163F}{}
{
  An \emph{identifier} is used to refer to a variable (see \ref{Variables}). An identifier usually consists of letters, digits, underscores \texttt{{\textunderscore}}, and ``at''\texttt{\symbol{45}}characters \texttt{@}, and must contain at least one non\texttt{\symbol{45}}digit. An identifier is
terminated by the first character not in this class. Note that the ``at''\texttt{\symbol{45}}character \texttt{@} is used to implement namespaces, see Section \ref{Namespaces} for details. 

 Examples of valid identifiers are 

 
\begin{Verbatim}[commandchars=!|A,fontsize=\small,frame=single,label=]
  a           foo         aLongIdentifier
  hello       Hello       HELLO
  x100        100x       _100
  some_people_prefer_underscores_to_separate_words
  WePreferMixedCaseToSeparateWords
  abc@def
\end{Verbatim}
 

 Note that case is significant, so the three identifiers in the second line are
distinguished. 

 The backslash \texttt{\texttt{\symbol{92}}} can be used to include other characters in identifiers; a backslash followed
by a character is equivalent to the character, except that this escape
sequence is considered to be an ordinary letter. For example 
\begin{Verbatim}[commandchars=!@|,fontsize=\small,frame=single,label=]
  G\(2\,5\)
\end{Verbatim}
 is an identifier, not a call to a function \texttt{G}. 

 An identifier that starts with a backslash is never a keyword, so for example \texttt{\texttt{\symbol{92}}*} and \texttt{\texttt{\symbol{92}}mod} are identifiers. 

 The length of identifiers is not limited, however only the first $1023$ characters are significant. The escape sequence \texttt{\texttt{\symbol{92}}}\textsc{newline} is ignored, making it possible to split long identifiers over multiple lines. 

\subsection{\textcolor{Chapter }{IsValidIdentifier}}
\logpage{[ 4, 6, 1 ]}\nobreak
\hyperdef{L}{X85CF993B7D19F2C4}{}
{\noindent\textcolor{FuncColor}{$\triangleright$\enspace\texttt{IsValidIdentifier({\mdseries\slshape str})\index{IsValidIdentifier@\texttt{IsValidIdentifier}}
\label{IsValidIdentifier}
}\hfill{\scriptsize (function)}}\\


 returns \texttt{true} if the string \mbox{\texttt{\mdseries\slshape str}} would form a valid identifier consisting of letters, digits and underscores;
otherwise it returns \texttt{false}. It does not check whether \mbox{\texttt{\mdseries\slshape str}} contains characters escaped by a backslash \texttt{\texttt{\symbol{92}}}. }

 Note that the ``at''\texttt{\symbol{45}}character is used to implement namespaces for global
variables in packages. See \ref{Namespaces} for details. \index{namespace} 
\subsection{\textcolor{Chapter }{Conventions about Identifiers}}\label{Conventions about Identifiers}
\logpage{[ 4, 6, 2 ]}
\hyperdef{L}{X839A7F8E84BBCA57}{}
{
  \index{variable names} (The following rule is stated also in Section  (\textbf{Tutorial: Variables versus Objects}).) 

 The name of almost every global variable in the \textsf{GAP} library and in \textsf{GAP} packages starts with a \emph{capital letter}. (See Section \ref{Main Loop} for the few exceptions.) For user variables, we recommend only choosing names
that start with a \emph{lower case letter}, in order to avoid name clashes. 

 For example, valid \textsf{GAP} input which assigns some user variables whose names start with capital letters
may run into errors with a newer version of \textsf{GAP} or in a \textsf{GAP} session with more or newer packages, because it may happen that these
variables are predefined global variables in this situation. }

 }

  
\section{\textcolor{Chapter }{Expressions}}\label{Expressions}
\logpage{[ 4, 7, 0 ]}
\hyperdef{L}{X7BAFE9C1817253C6}{}
{
  \index{evaluation} An \emph{expression} is a construct that evaluates to a value. Syntactic constructs that are
executed to produce a side effect and return no value are called \emph{statements} (see \ref{Statements}). Expressions appear as right hand sides of assignments (see \ref{Assignments}), as actual arguments in function calls (see \ref{Function Calls}), and in statements. 

 Note that an expression is not the same as a value. For example \texttt{1 + 11} is an expression, whose value is the integer 12. The external representation
of this integer is the character sequence \texttt{12}, i.e., this sequence is output if the integer is printed. This sequence is
another expression whose value is the integer $12$. The process of finding the value of an expression is done by the interpreter
and is called the \emph{evaluation} of the expression. 

 The simplest cases of expressions are the following: 
\begin{itemize}
\item variables (see Section{\nobreakspace}\ref{Variables}),
\item function literals (see Section{\nobreakspace}\ref{Function}),
\item function calls (see Section{\nobreakspace}\ref{Function Calls}),
\item integer literals (see Chapter{\nobreakspace}\ref{Integers}),
\item floating point literals (see Chapter{\nobreakspace}\ref{Floats}),
\item permutation literals (see Chapter{\nobreakspace}\ref{Permutations}),
\item string literals (see Chapter{\nobreakspace}\ref{Strings and Characters}),
\item character literals (see Chapter{\nobreakspace}\ref{Strings and Characters}),
\item list literals (see Chapter{\nobreakspace}\ref{Lists}), and
\item record literals (see Chapter{\nobreakspace}\ref{Records}).
\end{itemize}
 Expressions, for example the simple expressions mentioned above, can be
combined with the operators to form more complex expressions. Of course those
expressions can then be combined further with the operators to form even more
complex expressions. The \emph{operators} fall into three classes. \index{operators} The \emph{comparisons} are \texttt{=}, \texttt{{\textless}{\textgreater}}, \texttt{{\textless}}, \texttt{{\textless}=}, \texttt{{\textgreater}}, \texttt{{\textgreater}=}, and \texttt{in} (see \ref{Comparisons} and \ref{Membership Test for Collections}). The \emph{arithmetic operators} are \texttt{+}, \texttt{\texttt{\symbol{45}}}, \texttt{*}, \texttt{/}, \texttt{mod}, and \texttt{\texttt{\symbol{94}}} (see{\nobreakspace}\ref{Arithmetic Operators}). The \emph{logical operators} are \texttt{not}, \texttt{and}, and \texttt{or} (see{\nobreakspace}\ref{Operations for Booleans}). 

 The following example shows a very simple expression with value 4 and a more
complex expression. 

 
\begin{Verbatim}[commandchars=!@|,fontsize=\small,frame=single,label=Example]
  !gapprompt@gap>| !gapinput@2 * 2;|
  4
  !gapprompt@gap>| !gapinput@2 * 2 + 9 = Fibonacci(7) and Fibonacci(13) in Primes;|
  true
\end{Verbatim}
 

 The following table lists all operators by precedence, from highest to lowest,
and also indicates whether the operator is left associative (aka
left\texttt{\symbol{45}}to\texttt{\symbol{45}}right) or right associative (aka
right\texttt{\symbol{45}}to\texttt{\symbol{45}}left) or neither. \index{precedence} \index{associativity} \begin{center}
\begin{tabular}{ll}\emph{operator}&
\emph{associativity}\\
\hline
\emph{arithmetic (see{\nobreakspace}\ref{Arithmetic Operators})}&
\\
{\nobreakspace}{\nobreakspace}{\nobreakspace}{\nobreakspace}\texttt{\texttt{\symbol{94}}}&
none\\
{\nobreakspace}{\nobreakspace}{\nobreakspace}{\nobreakspace}unary \texttt{+}, unary \texttt{\texttt{\symbol{45}}}&
right\texttt{\symbol{45}}to\texttt{\symbol{45}}left\\
{\nobreakspace}{\nobreakspace}{\nobreakspace}{\nobreakspace}\texttt{*}, \texttt{/}, \texttt{mod}&
left\texttt{\symbol{45}}to\texttt{\symbol{45}}right\\
{\nobreakspace}{\nobreakspace}{\nobreakspace}{\nobreakspace}binary \texttt{+}, binary \texttt{\texttt{\symbol{45}}}&
left\texttt{\symbol{45}}to\texttt{\symbol{45}}right\\
\hline
\emph{comparison (see{\nobreakspace}\ref{Comparisons})}&
\\
{\nobreakspace}{\nobreakspace}{\nobreakspace}{\nobreakspace}\texttt{=}, \texttt{{\textless}{\textgreater}}, \texttt{{\textless}}, \texttt{{\textless}=}, \texttt{{\textgreater}}, \texttt{{\textgreater}=}, and \texttt{in}&
none\\
\hline
\emph{logical (see{\nobreakspace}\ref{Operations for Booleans})}&
\\
{\nobreakspace}{\nobreakspace}{\nobreakspace}{\nobreakspace}\texttt{not}&
right\texttt{\symbol{45}}to\texttt{\symbol{45}}left\\
{\nobreakspace}{\nobreakspace}{\nobreakspace}{\nobreakspace}\texttt{and}&
left\texttt{\symbol{45}}to\texttt{\symbol{45}}right\\
{\nobreakspace}{\nobreakspace}{\nobreakspace}{\nobreakspace}\texttt{or}&
left\texttt{\symbol{45}}to\texttt{\symbol{45}}right\\
\end{tabular}\\[2mm]
\end{center}

 }

  
\section{\textcolor{Chapter }{Variables}}\label{Variables}
\logpage{[ 4, 8, 0 ]}
\hyperdef{L}{X7A4C2D0E7E286B4F}{}
{
  \index{scope}\index{bound} A \emph{variable} is a location in a \textsf{GAP} program that points to a value. We say the variable is \emph{bound} to this value. If a variable is evaluated it evaluates to this value. 

 Initially an ordinary variable is not bound to any value. The variable can be
bound to a value by \emph{assigning} this value to the variable (see \ref{Assignments}). Because of this we sometimes say that a variable that is not bound to any
value has no assigned value. Assignment is in fact the only way by which a
variable, which is not an argument of a function, can be bound to a value.
After a variable has been bound to a value an assignment can also be used to
bind the variable to another value. 

 A special class of variables is the class of \emph{arguments} of functions. They behave similarly to other variables, except they are bound
to the value of the actual arguments upon a function call (see \ref{Function Calls}). 

 Each variable has a name that is also called its \emph{identifier}. This is because in a given scope an identifier identifies a unique variable
(see \ref{Identifiers}). A \emph{scope} is a lexical part of a program text. There is the \emph{global scope} that encloses the entire program text, and there are local scopes that range
from the \texttt{function} keyword, denoting the beginning of a function definition, to the corresponding \texttt{end} keyword. A \emph{local scope} introduces new variables, whose identifiers are given in the formal argument
list and the \texttt{local} declaration of the function (see \ref{Function}). Usage of an identifier in a program text refers to the variable in the
innermost scope that has this identifier as its name. Because this mapping
from identifiers to variables is done when the program is read, not when it is
executed, \textsf{GAP} is said to have \emph{lexical scoping}. The following example shows how one identifier refers to different variables
at different points in the program text. 

 
\begin{Verbatim}[commandchars=!@|,fontsize=\small,frame=single,label=]
  g := 0;      # global variable g
  x := function ( a, b, c )
    local  y;
    g := c;     # c refers to argument c of function x
    y := function ( y )
      local d, e, f;
      d := y;   # y refers to argument y of function y
      e := b;   # b refers to argument b of function x
      f := g;   # g refers to global variable g
      return d + e + f;
    end;
    return y( a ); # y refers to local y of function x
  end;
\end{Verbatim}
 

 It is important to note that the concept of a variable in \textsf{GAP} is quite different from the concept of a variable in most compiled programming
languages. 

 In those languages a variable denotes a block of memory. The value of the
variable is stored in this block. So in those languages two variables can have
the same value, but they can never have identical values, because they denote
different blocks of memory. Note that some languages have the concept of a
reference argument. It seems as if such an argument and the variable used in
the actual function call have the same value, since changing the argument's
value also changes the value of the variable used in the actual function call.
But this is not so; the reference argument is actually a pointer to the
variable used in the actual function call, and it is the compiler that inserts
enough magic to make the pointer invisible. In order for this to work the
compiler needs enough information to compute the amount of memory needed for
each variable in a program, which is readily available in the declarations. 

 In \textsf{GAP} on the other hand each variable just points to a value, and different
variables can share the same value. 

\subsection{\textcolor{Chapter }{IsBound (for a global variable)}}
\logpage{[ 4, 8, 1 ]}\nobreak
\hyperdef{L}{X842B89D4860FD5DB}{}
{\noindent\textcolor{FuncColor}{$\triangleright$\enspace\texttt{IsBound({\mdseries\slshape ident})\index{IsBound@\texttt{IsBound}!for a global variable}
\label{IsBound:for a global variable}
}\hfill{\scriptsize (function)}}\\


 \texttt{IsBound} returns \texttt{true} if the variable \mbox{\texttt{\mdseries\slshape ident}} points to a value, and \texttt{false} otherwise. 

 For records and lists \texttt{IsBound} can be used to check whether components or entries, respectively, are bound
(see Chapters{\nobreakspace}\ref{Records} and \ref{Lists}). }

 

\subsection{\textcolor{Chapter }{Unbind (unbind a variable)}}
\logpage{[ 4, 8, 2 ]}\nobreak
\hyperdef{L}{X7BABB3E77F52626C}{}
{\noindent\textcolor{FuncColor}{$\triangleright$\enspace\texttt{Unbind({\mdseries\slshape ident})\index{Unbind@\texttt{Unbind}!unbind a variable}
\label{Unbind:unbind a variable}
}\hfill{\scriptsize (function)}}\\


 deletes the identifier \mbox{\texttt{\mdseries\slshape ident}}. If there is no other variable pointing to the same value as \mbox{\texttt{\mdseries\slshape ident}} was, this value will be removed by the next garbage collection. Therefore \texttt{Unbind} can be used to get rid of unwanted large objects. 

 For records and lists \texttt{Unbind} can be used to delete components or entries, respectively (see
Chapters{\nobreakspace}\ref{Records} and \ref{Lists}). }

 }

  
\section{\textcolor{Chapter }{More About Global Variables}}\label{More About Global Variables}
\logpage{[ 4, 9, 0 ]}
\hyperdef{L}{X816FBEEA85782EC2}{}
{
  The vast majority of variables in \textsf{GAP} are defined at the outer level (the global scope). They are used to access
functions and other objects created either in the \textsf{GAP} library or packages or in the user's code. 

 Note that for packages there is a mechanism to implement package local
namespaces on top of this global namespace. See Section \ref{Namespaces} for details. \index{namespace} 

 Certain special facilities are provided for manipulating global variables
which are not available for other types of variable (such as local variables
or function arguments). 

 First, such variables may be marked \emph{read\texttt{\symbol{45}}only} using \texttt{MakeReadOnlyGlobal} (\ref{MakeReadOnlyGlobal}). In which case attempts to change them will fail. Most of the global
variables defined in the \textsf{GAP} library are so marked. \emph{read\texttt{\symbol{45}}only} variables can be made read\texttt{\symbol{45}}write again by calling \texttt{MakeReadWriteGlobal} (\ref{MakeReadWriteGlobal}). GAP also features \emph{constant} variables, which are created by calling \texttt{MakeConstantGlobal} (\ref{MakeConstantGlobal}). Constant variables can never be changed. In some cases, GAP can optimise
code which uses \emph{constant} variables, as their value never changes. In this version GAP these
optimisations can be observed by printing the function back out, but this
behaviour may change in future. 
\begin{Verbatim}[commandchars=!@|,fontsize=\small,frame=single,label=Example]
  !gapprompt@gap>| !gapinput@globali := 1 + 2;;|
  !gapprompt@gap>| !gapinput@globalb := true;;|
  !gapprompt@gap>| !gapinput@MakeConstantGlobal("globali");|
  !gapprompt@gap>| !gapinput@MakeConstantGlobal("globalb");|
  !gapprompt@gap>| !gapinput@f := function()|
  !gapprompt@>| !gapinput@    if globalb then|
  !gapprompt@>| !gapinput@        return globali + 1;|
  !gapprompt@>| !gapinput@    else|
  !gapprompt@>| !gapinput@        return globali + 2;|
  !gapprompt@>| !gapinput@    fi;|
  !gapprompt@>| !gapinput@end;;|
  !gapprompt@gap>| !gapinput@Print(f);|
  function (  )
      return 3 + 1;
  end
\end{Verbatim}
 

 Second, a group of functions are supplied for accessing and altering the
values assigned to global variables. Use of these functions differs from the
use of assignment, \texttt{Unbind} (\ref{Unbind:unbind a variable}) and \texttt{IsBound} (\ref{IsBound:for a global variable}) statements, in two ways. First, these functions always affect global
variables, even if local variables of the same names exist. Second, the
variable names are passed as strings, rather than being written directly into
the statements. 

 Note that the functions \texttt{NamesGVars} (\ref{NamesGVars}), \texttt{NamesSystemGVars} (\ref{NamesSystemGVars}), and \texttt{NamesUserGVars} (\ref{NamesUserGVars}), deal with the \emph{global namespace}. 

\subsection{\textcolor{Chapter }{IsReadOnlyGlobal}}
\logpage{[ 4, 9, 1 ]}\nobreak
\hyperdef{L}{X7CD3523B84744EB2}{}
{\noindent\textcolor{FuncColor}{$\triangleright$\enspace\texttt{IsReadOnlyGlobal({\mdseries\slshape name})\index{IsReadOnlyGlobal@\texttt{IsReadOnlyGlobal}}
\label{IsReadOnlyGlobal}
}\hfill{\scriptsize (function)}}\\


 returns \texttt{true} if the global variable named by the string \mbox{\texttt{\mdseries\slshape name}} is read\texttt{\symbol{45}}only and \texttt{false} otherwise (the default). }

 

\subsection{\textcolor{Chapter }{MakeReadOnlyGlobal}}
\logpage{[ 4, 9, 2 ]}\nobreak
\hyperdef{L}{X850CE44478254F27}{}
{\noindent\textcolor{FuncColor}{$\triangleright$\enspace\texttt{MakeReadOnlyGlobal({\mdseries\slshape name})\index{MakeReadOnlyGlobal@\texttt{MakeReadOnlyGlobal}}
\label{MakeReadOnlyGlobal}
}\hfill{\scriptsize (function)}}\\


 marks the global variable named by the string \mbox{\texttt{\mdseries\slshape name}} as read\texttt{\symbol{45}}only. 

 A warning is given if \mbox{\texttt{\mdseries\slshape name}} has no value bound to it or if it is already read\texttt{\symbol{45}}only. }

 

\subsection{\textcolor{Chapter }{MakeReadWriteGlobal}}
\logpage{[ 4, 9, 3 ]}\nobreak
\hyperdef{L}{X832AAF13861968BE}{}
{\noindent\textcolor{FuncColor}{$\triangleright$\enspace\texttt{MakeReadWriteGlobal({\mdseries\slshape name})\index{MakeReadWriteGlobal@\texttt{MakeReadWriteGlobal}}
\label{MakeReadWriteGlobal}
}\hfill{\scriptsize (function)}}\\


 marks the global variable named by the string \mbox{\texttt{\mdseries\slshape name}} as read\texttt{\symbol{45}}write. 

 A warning is given if \mbox{\texttt{\mdseries\slshape name}} is already read\texttt{\symbol{45}}write. 

 
\begin{Verbatim}[commandchars=!@|,fontsize=\small,frame=single,label=Example]
  !gapprompt@gap>| !gapinput@xx := 17;|
  17
  !gapprompt@gap>| !gapinput@IsReadOnlyGlobal("xx");|
  false
  !gapprompt@gap>| !gapinput@xx := 15;|
  15
  !gapprompt@gap>| !gapinput@MakeReadOnlyGlobal("xx");|
  !gapprompt@gap>| !gapinput@xx := 16;|
  Error, Variable: 'xx' is read only
  not in any function at *stdin*:3
  you can enter 'quit;' to quit to outer loop
  !gapbrkprompt@brk>| !gapinput@quit;|
  !gapprompt@gap>| !gapinput@IsReadOnlyGlobal("xx");|
  true
  !gapprompt@gap>| !gapinput@MakeReadWriteGlobal("xx");|
  !gapprompt@gap>| !gapinput@xx := 16;|
  16
  !gapprompt@gap>| !gapinput@IsReadOnlyGlobal("xx");|
  false
\end{Verbatim}
 }

 

\subsection{\textcolor{Chapter }{MakeConstantGlobal}}
\logpage{[ 4, 9, 4 ]}\nobreak
\hyperdef{L}{X847706237E72418F}{}
{\noindent\textcolor{FuncColor}{$\triangleright$\enspace\texttt{MakeConstantGlobal({\mdseries\slshape name})\index{MakeConstantGlobal@\texttt{MakeConstantGlobal}}
\label{MakeConstantGlobal}
}\hfill{\scriptsize (function)}}\\


 MakeConstantGlobal ( \mbox{\texttt{\mdseries\slshape name}} ) marks the global variable named by the string \mbox{\texttt{\mdseries\slshape name}} as constant. A constant variable can never be reassigned or made
read\texttt{\symbol{45}}write again. 

 A warning is given if \mbox{\texttt{\mdseries\slshape name}} is already constant. }

 

\subsection{\textcolor{Chapter }{ValueGlobal}}
\logpage{[ 4, 9, 5 ]}\nobreak
\hyperdef{L}{X84BB4B1E872849FF}{}
{\noindent\textcolor{FuncColor}{$\triangleright$\enspace\texttt{ValueGlobal({\mdseries\slshape name})\index{ValueGlobal@\texttt{ValueGlobal}}
\label{ValueGlobal}
}\hfill{\scriptsize (function)}}\\


 returns the value currently bound to the global variable named by the string \mbox{\texttt{\mdseries\slshape name}}. An error is raised if no value is currently bound. }

 

\subsection{\textcolor{Chapter }{IsBoundGlobal}}
\logpage{[ 4, 9, 6 ]}\nobreak
\hyperdef{L}{X823D4BC378395B32}{}
{\noindent\textcolor{FuncColor}{$\triangleright$\enspace\texttt{IsBoundGlobal({\mdseries\slshape name})\index{IsBoundGlobal@\texttt{IsBoundGlobal}}
\label{IsBoundGlobal}
}\hfill{\scriptsize (function)}}\\


 returns \texttt{true} if a value currently bound to the global variable named by the string \mbox{\texttt{\mdseries\slshape name}} and \texttt{false} otherwise. }

 

\subsection{\textcolor{Chapter }{UnbindGlobal}}
\logpage{[ 4, 9, 7 ]}\nobreak
\hyperdef{L}{X829A5F0E811F77D3}{}
{\noindent\textcolor{FuncColor}{$\triangleright$\enspace\texttt{UnbindGlobal({\mdseries\slshape name})\index{UnbindGlobal@\texttt{UnbindGlobal}}
\label{UnbindGlobal}
}\hfill{\scriptsize (function)}}\\


 removes any value currently bound to the global variable named by the string \mbox{\texttt{\mdseries\slshape name}}. Nothing is returned. 

 A warning is given if \mbox{\texttt{\mdseries\slshape name}} was not bound. The global variable named by \mbox{\texttt{\mdseries\slshape name}} must be writable, otherwise an error is raised. }

 

\subsection{\textcolor{Chapter }{BindGlobal}}
\logpage{[ 4, 9, 8 ]}\nobreak
\hyperdef{L}{X7D39D3E17CF49F5B}{}
{\noindent\textcolor{FuncColor}{$\triangleright$\enspace\texttt{BindGlobal({\mdseries\slshape name, val})\index{BindGlobal@\texttt{BindGlobal}}
\label{BindGlobal}
}\hfill{\scriptsize (function)}}\\
\noindent\textcolor{FuncColor}{$\triangleright$\enspace\texttt{BindConstant({\mdseries\slshape name, val})\index{BindConstant@\texttt{BindConstant}}
\label{BindConstant}
}\hfill{\scriptsize (function)}}\\


 \texttt{BindGlobal} and \texttt{BindConstant} set the global variable named by the string \mbox{\texttt{\mdseries\slshape name}} to the value \mbox{\texttt{\mdseries\slshape val}}, provided that variable is writable. \texttt{BindGlobal} makes the resulting variable read\texttt{\symbol{45}}only, while \texttt{BindConstant} makes it constant. If \mbox{\texttt{\mdseries\slshape name}} already had a value, a warning message is printed. 

 This is intended to be the normal way to create and set ``official'' global variables (such as operations, filters and constants). 

 Caution should be exercised in using these functions, especially \texttt{UnbindGlobal} (\ref{UnbindGlobal}) as unexpected changes in global variables can be very confusing for the user. 

 
\begin{Verbatim}[commandchars=!@|,fontsize=\small,frame=single,label=Example]
  !gapprompt@gap>| !gapinput@xx := 16;|
  16
  !gapprompt@gap>| !gapinput@IsReadOnlyGlobal("xx");|
  false
  !gapprompt@gap>| !gapinput@ValueGlobal("xx");|
  16
  !gapprompt@gap>| !gapinput@IsBoundGlobal("xx");|
  true
  !gapprompt@gap>| !gapinput@BindGlobal("xx",17);|
  #W BIND_GLOBAL: variable `xx' already has a value
  !gapprompt@gap>| !gapinput@xx;|
  17
  !gapprompt@gap>| !gapinput@IsReadOnlyGlobal("xx");|
  true
  !gapprompt@gap>| !gapinput@MakeReadWriteGlobal("xx");|
  !gapprompt@gap>| !gapinput@Unbind(xx);|
\end{Verbatim}
 }

 

\subsection{\textcolor{Chapter }{NamesGVars}}
\logpage{[ 4, 9, 9 ]}\nobreak
\hyperdef{L}{X876A6EB68745A510}{}
{\noindent\textcolor{FuncColor}{$\triangleright$\enspace\texttt{NamesGVars({\mdseries\slshape })\index{NamesGVars@\texttt{NamesGVars}}
\label{NamesGVars}
}\hfill{\scriptsize (function)}}\\


 This function returns an immutable (see{\nobreakspace}\ref{Mutability and Copyability}) sorted (see{\nobreakspace}\ref{Sorted Lists and Sets}) list of all the global variable names known to the system. This includes
names of variables which were bound but have now been unbound and some other
names which have never been bound but have become known to the system by
various routes. }

 

\subsection{\textcolor{Chapter }{NamesSystemGVars}}
\logpage{[ 4, 9, 10 ]}\nobreak
\hyperdef{L}{X7E604AF579A7BC92}{}
{\noindent\textcolor{FuncColor}{$\triangleright$\enspace\texttt{NamesSystemGVars({\mdseries\slshape })\index{NamesSystemGVars@\texttt{NamesSystemGVars}}
\label{NamesSystemGVars}
}\hfill{\scriptsize (function)}}\\


 This function returns an immutable sorted list of all the global variable
names created by the \textsf{GAP} library when \textsf{GAP} was started. }

 

\subsection{\textcolor{Chapter }{NamesUserGVars}}
\logpage{[ 4, 9, 11 ]}\nobreak
\hyperdef{L}{X870169447AF490D8}{}
{\noindent\textcolor{FuncColor}{$\triangleright$\enspace\texttt{NamesUserGVars({\mdseries\slshape })\index{NamesUserGVars@\texttt{NamesUserGVars}}
\label{NamesUserGVars}
}\hfill{\scriptsize (function)}}\\


 This function returns an immutable sorted list of the global variable names
created since the library was read, to which a value is currently bound. }

  }

  
\section{\textcolor{Chapter }{Namespaces for \textsf{GAP} packages}}\label{Namespaces}
\logpage{[ 4, 10, 0 ]}
\hyperdef{L}{X7DF8774F7D542298}{}
{
  As mentioned in Section \ref{More About Global Variables} above all global variables share a common namespace. This can relatively
easily lead to name clashes, in particular when many \textsf{GAP} packages are loaded at the same time. To give package code a way to have a
package local namespace without breaking backward compatibility of the \textsf{GAP} language, the following simple rule has been devised: 

 If in package code a global variable that ends with an ``at''\texttt{\symbol{45}}character \texttt{@} is accessed in any way, the name of the package is appended before accessing
it. Here, ``package code'' refers to everything which is read with \texttt{ReadPackage} (\ref{ReadPackage}). As the name of the package the entry \texttt{PackageName} in its \texttt{PackageInfo.g} file is taken. As for all identifiers, this name is case sensitive. 

 For example, if the following is done in the code of a package with name \texttt{xYz}: 
\begin{Verbatim}[commandchars=!|A,fontsize=\small,frame=single,label=Example]
  !gapprompt|gap>A !gapinput|a@ := 12;A
  
\end{Verbatim}
 Then actually the global variable \texttt{a@xYz} is assigned. Further accesses to \texttt{a@} within the package code will all be redirected to \texttt{a@xYz}. This includes all the functions described in Section \ref{More About Global Variables} and indeed all the functions described Section \ref{Global Variables in the Library} like for example \texttt{DeclareCategory} (\ref{DeclareCategory}). Note that from code in the same package it is still possible to access the
same global variable via \texttt{a@xYz} explicitly. 

 All other code outside the package as well as interactive user input that
wants to refer to that variable \texttt{a@xYz} must do so explicitly by using \texttt{a@xYz}. 

 Since in earlier releases of \textsf{GAP} the ``at''\texttt{\symbol{45}}character \texttt{@} was not a legal character (without using backslashes), this small extension of
the language does not break any old code. }

  
\section{\textcolor{Chapter }{Function}}\label{Function}
\logpage{[ 4, 11, 0 ]}
\hyperdef{L}{X815F71EA7BC0EB6F}{}
{
  \index{functions!definition of} \index{end@\texttt{end}} \index{local@\texttt{local}} \index{recursion} \index{functions!recursive} \index{environment}\index{body} \texttt{function( [ \mbox{\texttt{\mdseries\slshape arg\texttt{\symbol{45}}ident}} \texttt{\symbol{123}}, \mbox{\texttt{\mdseries\slshape arg\texttt{\symbol{45}}ident}}\texttt{\symbol{125}} ] )} 

 \texttt{{\nobreakspace}{\nobreakspace}[local \mbox{\texttt{\mdseries\slshape loc\texttt{\symbol{45}}ident}} \texttt{\symbol{123}}, \mbox{\texttt{\mdseries\slshape loc\texttt{\symbol{45}}ident}}\texttt{\symbol{125}} ; ]} 

 \texttt{{\nobreakspace}{\nobreakspace}\mbox{\texttt{\mdseries\slshape statements}}} 

 \texttt{end} 

 A function literal can be assigned to a variable or to a list element or a
record component. Later this function can be called as described in \ref{Function Calls}. 

 The following is an example of a function definition. It is a function to
compute values of the Fibonacci sequence (see \texttt{Fibonacci} (\ref{Fibonacci})). 

 
\begin{Verbatim}[commandchars=!@|,fontsize=\small,frame=single,label=Example]
  !gapprompt@gap>| !gapinput@fib := function ( n )|
  !gapprompt@>| !gapinput@    local f1, f2, f3, i;|
  !gapprompt@>| !gapinput@    f1 := 1; f2 := 1;|
  !gapprompt@>| !gapinput@    for i in [3..n] do|
  !gapprompt@>| !gapinput@      f3 := f1 + f2;|
  !gapprompt@>| !gapinput@      f1 := f2;|
  !gapprompt@>| !gapinput@      f2 := f3;|
  !gapprompt@>| !gapinput@    od;|
  !gapprompt@>| !gapinput@    return f2;|
  !gapprompt@>| !gapinput@  end;;|
  !gapprompt@gap>| !gapinput@List( [1..10], fib );|
  [ 1, 1, 2, 3, 5, 8, 13, 21, 34, 55 ]
\end{Verbatim}
 

 Because for each of the formal arguments \mbox{\texttt{\mdseries\slshape arg\texttt{\symbol{45}}ident}} and for each of the formal locals \mbox{\texttt{\mdseries\slshape loc\texttt{\symbol{45}}ident}} a new variable is allocated when the function is called (see \ref{Function Calls}), it is possible that a function calls itself. This is usually called \emph{recursion}. The following is a recursive function that computes values of the Fibonacci
sequence. 

 
\begin{Verbatim}[commandchars=!@|,fontsize=\small,frame=single,label=Example]
  !gapprompt@gap>| !gapinput@fib := function ( n )|
  !gapprompt@>| !gapinput@    if n < 3 then|
  !gapprompt@>| !gapinput@      return 1;|
  !gapprompt@>| !gapinput@    else|
  !gapprompt@>| !gapinput@      return fib(n-1) + fib(n-2);|
  !gapprompt@>| !gapinput@    fi;|
  !gapprompt@>| !gapinput@  end;;|
  !gapprompt@gap>| !gapinput@List( [1..10], fib );|
  [ 1, 1, 2, 3, 5, 8, 13, 21, 34, 55 ]
\end{Verbatim}
 

 Note that the recursive version needs \texttt{2 * fib(\mbox{\texttt{\mdseries\slshape n}})\texttt{\symbol{45}}1} steps to compute \texttt{fib(\mbox{\texttt{\mdseries\slshape n}})}, while the iterative version of \texttt{fib} needs only \texttt{\mbox{\texttt{\mdseries\slshape n}}\texttt{\symbol{45}}2} steps. Both are not optimal however, the library function \texttt{Fibonacci} (\ref{Fibonacci}) only needs about \texttt{Log(\mbox{\texttt{\mdseries\slshape n}})} steps. 

 \index{functions!with a variable number of arguments} \index{arg@\texttt{arg}!special function argument} As noted in Section{\nobreakspace}\ref{Function Calls}, the case where a function's last argument is followed by \texttt{...} is special. It provides a way of defining a function with a variable number of
arguments. The values of the actual arguments are computed and the first ones
are assigned to the new variables corresponding to the formal arguments before
the last argument, if any. The values of all the remaining actual arguments
are stored in a list and this list is assigned to the new variable
corresponding to the final formal argument. There are two typical scenarios
for wanting such a possibility: having optional arguments and having any
number of arguments. 

 The following example shows one way that the function \texttt{Position} (\ref{Position}) might be encoded and demonstrates the ``optional argument'' scenario. 

 
\begin{Verbatim}[commandchars=!@|,fontsize=\small,frame=single,label=Example]
  !gapprompt@gap>| !gapinput@position := function ( list, obj, arg... )|
  !gapprompt@>| !gapinput@    local pos;|
  !gapprompt@>| !gapinput@    if 0 = Length(arg) then|
  !gapprompt@>| !gapinput@      pos := 0;|
  !gapprompt@>| !gapinput@    else|
  !gapprompt@>| !gapinput@      pos := arg[1];|
  !gapprompt@>| !gapinput@    fi;|
  !gapprompt@>| !gapinput@    repeat|
  !gapprompt@>| !gapinput@      pos := pos + 1;|
  !gapprompt@>| !gapinput@      if pos > Length(list) then|
  !gapprompt@>| !gapinput@        return fail;|
  !gapprompt@>| !gapinput@      fi;|
  !gapprompt@>| !gapinput@    until list[pos] = obj;|
  !gapprompt@>| !gapinput@    return pos;|
  !gapprompt@>| !gapinput@   end;|
  function( list, obj, arg... ) ... end
  !gapprompt@gap>| !gapinput@position([1, 4, 2], 4);|
  2
  !gapprompt@gap>| !gapinput@position([1, 4, 2], 3);|
  fail
  !gapprompt@gap>| !gapinput@position([1, 4, 2], 4, 2);|
  fail
\end{Verbatim}
 

 The following example demonstrates the ``any number of arguments'' scenario. 

 
\begin{Verbatim}[commandchars=!@|,fontsize=\small,frame=single,label=Example]
  !gapprompt@gap>| !gapinput@sum := function ( l... )|
  !gapprompt@>| !gapinput@    local total, x;|
  !gapprompt@>| !gapinput@    total := 0;|
  !gapprompt@>| !gapinput@    for x in l do|
  !gapprompt@>| !gapinput@      total := total + x;|
  !gapprompt@>| !gapinput@    od;|
  !gapprompt@>| !gapinput@    return total;|
  !gapprompt@>| !gapinput@   end;|
  function( l... ) ... end
  !gapprompt@gap>| !gapinput@sum(1, 2, 3);|
  6
  !gapprompt@gap>| !gapinput@sum(1, 2, 3, 4);|
  10
  !gapprompt@gap>| !gapinput@sum();|
  0
\end{Verbatim}
 

 The user should compare the above with the \textsf{GAP} function \texttt{Sum} (\ref{Sum}) which, for example, may take a list argument and optionally an initial element
(which zero should the sum of an empty list return?). 

 GAP will also special case a function with a single argument with the name \texttt{arg} as function with a variable length list of arguments, as if the user had
written \texttt{arg...}. 

 Note that if a function \mbox{\texttt{\mdseries\slshape f}} is defined as above then \texttt{NumberArgumentsFunction(\mbox{\texttt{\mdseries\slshape f}})} returns minus the number of formal arguments (including the final argument)
(see{\nobreakspace}\texttt{NumberArgumentsFunction} (\ref{NumberArgumentsFunction})). 

 Using the \texttt{...} notation on a function \mbox{\texttt{\mdseries\slshape f}} with only a single named argument tells \textsf{GAP} that when it encounters \mbox{\texttt{\mdseries\slshape f}} that it should form a list out of the arguments of \mbox{\texttt{\mdseries\slshape f}}. What if one wishes to do the ``opposite'': tell \textsf{GAP} that a list should be ``unwrapped'' and passed as several arguments to a function. The function \texttt{CallFuncList} (\ref{CallFuncList}) is provided for this purpose. 

 Also see Chapter{\nobreakspace}\ref{Functions}. 

 \index{functions!definition by arrow notation} \index{arrow notation for functions} \texttt{\texttt{\symbol{123}} \mbox{\texttt{\mdseries\slshape arg\texttt{\symbol{45}}list}} \texttt{\symbol{125}} \texttt{\symbol{45}}{\textgreater} \mbox{\texttt{\mdseries\slshape expr}}} 

 This is a shorthand for 

 \texttt{function ( \mbox{\texttt{\mdseries\slshape arg\texttt{\symbol{45}}list}} ) return \mbox{\texttt{\mdseries\slshape expr}}; end.} 

 \mbox{\texttt{\mdseries\slshape arg\texttt{\symbol{45}}list}} is a (possibly empty) argument list. Any arguments list which would be valid
for a normal GAP function is also valid here (including variadic arguments). 

 The following gives a couple of examples of a typical use of such a function 

 
\begin{Verbatim}[commandchars=!@|,fontsize=\small,frame=single,label=Example]
  !gapprompt@gap>| !gapinput@Sum( List( [1..100], {x} -> x^2 ) );|
  338350
  !gapprompt@gap>| !gapinput@list := [3, 5, 2, 1, 3];;|
  !gapprompt@gap>| !gapinput@Sort(list, {x,y} -> x > y);|
  !gapprompt@gap>| !gapinput@list;|
  [ 5, 3, 3, 2, 1 ]
  !gapprompt@gap>| !gapinput@f := {x,y...} -> y;;|
  !gapprompt@gap>| !gapinput@f(1,2,3,4);|
  [ 2, 3, 4 ]
  !gapprompt@gap>| !gapinput@f := {} -> 2;|
  function(  ) ... end
  !gapprompt@gap>| !gapinput@Print(f);|
  function (  )
      return 2;
  end
  !gapprompt@gap>| !gapinput@f();|
  2
\end{Verbatim}
 

 The \texttt{\texttt{\symbol{123}}} and \texttt{\texttt{\symbol{125}}} may be omitted for functions with one argument: 
\begin{Verbatim}[commandchars=!@|,fontsize=\small,frame=single,label=Example]
  !gapprompt@gap>| !gapinput@Sum( List( [1..100], {x} -> x^2 ) );|
  338350
  !gapprompt@gap>| !gapinput@Sum( List( [1..100], x -> x^2 ) );|
  338350
\end{Verbatim}
 

 When the definition of a function \mbox{\texttt{\mdseries\slshape fun1}} is evaluated inside another function \mbox{\texttt{\mdseries\slshape fun2}}, \textsf{GAP} binds all the identifiers inside the function \mbox{\texttt{\mdseries\slshape fun1}} that are identifiers of an argument or a local of \mbox{\texttt{\mdseries\slshape fun2}} to the corresponding variable. This set of bindings is called the environment
of the function \mbox{\texttt{\mdseries\slshape fun1}}. When \mbox{\texttt{\mdseries\slshape fun1}} is called, its body is executed in this environment. The following
implementation of a simple stack uses this. Values can be pushed onto the
stack and then later be popped off again. The interesting thing here is that
the functions \texttt{push} and \texttt{pop} in the record returned by \texttt{Stack} access the local variable \texttt{stack} of \texttt{Stack}. When \texttt{Stack} is called, a new variable for the identifier \texttt{stack} is created. When the function definitions of \texttt{push} and \texttt{pop} are then evaluated (as part of the \texttt{return} statement) each reference to \texttt{stack} is bound to this new variable. Note also that the two stacks \texttt{A} and \texttt{B} do not interfere, because each call of \texttt{Stack} creates a new variable for \texttt{stack}. 

 
\begin{Verbatim}[commandchars=!@|,fontsize=\small,frame=single,label=Example]
  !gapprompt@gap>| !gapinput@Stack := function()|
  !gapprompt@>| !gapinput@    local  stack;|
  !gapprompt@>| !gapinput@    stack := [];|
  !gapprompt@>| !gapinput@    return rec(|
  !gapprompt@>| !gapinput@      push := function( value )|
  !gapprompt@>| !gapinput@        Add( stack, value );|
  !gapprompt@>| !gapinput@      end,|
  !gapprompt@>| !gapinput@      pop := function()|
  !gapprompt@>| !gapinput@        return Remove( stack) ;|
  !gapprompt@>| !gapinput@      end|
  !gapprompt@>| !gapinput@    );|
  !gapprompt@>| !gapinput@ end;;|
  !gapprompt@gap>| !gapinput@A := Stack();;|
  !gapprompt@gap>| !gapinput@B := Stack();;|
  !gapprompt@gap>| !gapinput@A.push( 1 ); A.push( 2 ); A.push( 3 );|
  !gapprompt@gap>| !gapinput@B.push( 4 ); B.push( 5 ); B.push( 6 );|
  !gapprompt@gap>| !gapinput@A.pop(); A.pop(); A.pop();|
  3
  2
  1
  !gapprompt@gap>| !gapinput@B.pop(); B.pop(); B.pop();|
  6
  5
  4
\end{Verbatim}
 

 This feature should be used rarely, since its implementation in \textsf{GAP} is not very efficient. }

  
\section{\textcolor{Chapter }{Function Calls}}\label{Function Calls}
\logpage{[ 4, 12, 0 ]}
\hyperdef{L}{X78C70489791FDF43}{}
{
  
\subsection{\textcolor{Chapter }{Function Call With Arguments}}\label{Function Call With Arguments}
\logpage{[ 4, 12, 1 ]}
\hyperdef{L}{X80B93A9C7E0A57F4}{}
{
  \texttt{\mbox{\texttt{\mdseries\slshape function\texttt{\symbol{45}}var}}( [\mbox{\texttt{\mdseries\slshape arg\texttt{\symbol{45}}expr}}[, \mbox{\texttt{\mdseries\slshape arg\texttt{\symbol{45}}expr}}, ...]] )} 

 The function call has the effect of calling the function \mbox{\texttt{\mdseries\slshape function\texttt{\symbol{45}}var}}. The precise semantics are as follows. 

 First \textsf{GAP} evaluates the \mbox{\texttt{\mdseries\slshape function\texttt{\symbol{45}}var}}. Usually \mbox{\texttt{\mdseries\slshape function\texttt{\symbol{45}}var}} is a variable, and \textsf{GAP} does nothing more than taking the value of this variable. It is allowed though
that \mbox{\texttt{\mdseries\slshape function\texttt{\symbol{45}}var}} is a more complex expression, such as a reference to an element of a list (see
Chapter{\nobreakspace}\ref{Lists}) \texttt{\mbox{\texttt{\mdseries\slshape list\texttt{\symbol{45}}var}}[\mbox{\texttt{\mdseries\slshape int\texttt{\symbol{45}}expr}}]}, or to a component of a record (see Chapter{\nobreakspace}\ref{Records}) \texttt{\mbox{\texttt{\mdseries\slshape record\texttt{\symbol{45}}var}}.\mbox{\texttt{\mdseries\slshape ident}}}. In any case \textsf{GAP} tests whether the value is a function. If it is not, \textsf{GAP} signals an error. 

 \index{functions!with a variable number of arguments, calling} \index{arg@\texttt{arg}!special function argument, calling with} Next \textsf{GAP} checks that the number of actual arguments \mbox{\texttt{\mdseries\slshape arg\texttt{\symbol{45}}expr}}s agrees with the number of \emph{formal arguments} as given in the function definition. If they do not agree \textsf{GAP} signals an error. An exception is the case when the function has a variable
length argument list, which is denoted by adding \texttt{...} after the final argument. In this case there must be at least as many actual
arguments as there are formal arguments \emph{before the final argument} and can be any larger number (see{\nobreakspace}\ref{Function} for examples). 

 Now \textsf{GAP} allocates for each formal argument and for each \emph{formal local} (that is, the identifiers in the \texttt{local} declaration) a new variable. Remember that a variable is a location in a \textsf{GAP} program that points to a value. Thus for each formal argument and for each
formal local such a location is allocated. 

 Next the arguments \mbox{\texttt{\mdseries\slshape arg\texttt{\symbol{45}}expr}}s are evaluated from left to right, and the values are assigned to the newly
created variables corresponding to the formal arguments. Of course the first
value is assigned to the new variable corresponding to the first formal
argument, the second value is assigned to the new variable corresponding to
the second formal argument, and so on. An exception again occurs if the last
formal argument has the name \texttt{arg}. In this case the values of all the actual arguments not assigned to the
other formal parameters are stored in a list and this list is assigned to the
new variable corresponding to the formal argument \texttt{arg}. 

 The new variables corresponding to the formal locals are initially not bound
to any value. So trying to evaluate those variables before something has been
assigned to them will signal an error. 

 Now the body of the function, which is a statement, is executed. If the
identifier of one of the formal arguments or formal locals appears in the body
of the function it refers to the new variable that was allocated for this
formal argument or formal local, and evaluates to the value of this variable. 

 If during the execution of the body of the function a \texttt{return} statement with an expression (see \ref{Return}) is executed, execution of the body is terminated and the value of the
function call is the value of the expression of the \texttt{return}. If during the execution of the body a \texttt{return} statement without an expression is executed, execution of the body is
terminated and the function call does not produce a value, in which case we
call this call a procedure call (see \ref{Procedure Calls}). If the execution of the body completes without execution of a \texttt{return} statement, the function call again produces no value, and again we talk about
a procedure call. 

 
\begin{Verbatim}[commandchars=!@|,fontsize=\small,frame=single,label=Example]
  !gapprompt@gap>| !gapinput@Fibonacci( 11 );|
  89
\end{Verbatim}
 

 The above example shows a call to the function \texttt{Fibonacci} (\ref{Fibonacci}) with actual argument \texttt{11}, the following one shows a call to the operation \texttt{RightCosets} (\ref{RightCosets}) where the second actual argument is another function call. 

 
\begin{Verbatim}[commandchars=!@|,fontsize=\small,frame=single,label=Example]
  !gapprompt@gap>| !gapinput@RightCosets( G, Intersection( U, V ) );;|
\end{Verbatim}
 }

 
\subsection{\textcolor{Chapter }{Function Call With Options}}\label{Function Call With Options}
\logpage{[ 4, 12, 2 ]}
\hyperdef{L}{X867D54987EF86D1D}{}
{
  \texttt{\mbox{\texttt{\mdseries\slshape function\texttt{\symbol{45}}var}}( \mbox{\texttt{\mdseries\slshape arg\texttt{\symbol{45}}expr}}[, \mbox{\texttt{\mdseries\slshape arg\texttt{\symbol{45}}expr}}, ...][ : [ \mbox{\texttt{\mdseries\slshape option\texttt{\symbol{45}}expr}} [,\mbox{\texttt{\mdseries\slshape option\texttt{\symbol{45}}expr}}, ....]]])} 

 As well as passing arguments to a function, providing the mathematical input
to its calculation, it is sometimes useful to supply ``hints'' suggesting to \textsf{GAP} how the desired result may be computed more quickly, or specifying a level of
tolerance for random errors in a Monte Carlo algorithm. 

 Such hints may be supplied to a function\texttt{\symbol{45}}call \emph{and to all subsidiary functions called from that call} using the options mechanism. Options are separated from the actual arguments
by a colon \texttt{:} and have much the same syntax as the components of a record expression. The
one exception to this is that a component name may appear without a value, in
which case the value \texttt{true} is silently inserted. 

 Options are evaluated from left to right, but only after all arguments have
been evaluated. 

 The following example shows a call to \texttt{Size} (\ref{Size}) passing the options \texttt{hard} (with the value \texttt{true}) and \texttt{tcselection} (with the string \texttt{"external"} as value). 

 
\begin{Verbatim}[commandchars=!@|,fontsize=\small,frame=single,label=Example]
  !gapprompt@gap>| !gapinput@Size( fpgrp : hard, tcselection := "external" );|
\end{Verbatim}
 

 Options supplied with function calls in this way are passed down using the
global options stack described in chapter \ref{Options Stack}, so that the call above is exactly equivalent to 

 
\begin{Verbatim}[commandchars=!@|,fontsize=\small,frame=single,label=Example]
  !gapprompt@gap>| !gapinput@PushOptions( rec( hard := true, tcselection := "external") );|
  !gapprompt@gap>| !gapinput@Size( fpgrp );|
  !gapprompt@gap>| !gapinput@PopOptions( );|
\end{Verbatim}
 

 \emph{Note} that any option may be passed with any function, whether or not it has any
actual meaning for that function, or any function called by it. The system
provides no safeguard against misspelled option names. }

 }

  
\section{\textcolor{Chapter }{Comparisons}}\label{Comparisons}
\logpage{[ 4, 13, 0 ]}
\hyperdef{L}{X7A274A1F8553B7E6}{}
{
  \index{equality test} \texttt{\mbox{\texttt{\mdseries\slshape left\texttt{\symbol{45}}expr}} = \mbox{\texttt{\mdseries\slshape right\texttt{\symbol{45}}expr}}} 

 \index{inequality test} \texttt{\mbox{\texttt{\mdseries\slshape left\texttt{\symbol{45}}expr}} {\textless}{\textgreater} \mbox{\texttt{\mdseries\slshape right\texttt{\symbol{45}}expr}}} 

 The operator \texttt{=} tests for equality of its two operands and evaluates to \texttt{true} if they are equal and to \texttt{false} otherwise. Likewise \texttt{{\textless}{\textgreater}} tests for inequality of its two operands. For each type of objects the
definition of equality is given in the respective chapter. Objects in
different families (see{\nobreakspace}\ref{Families}) are never equal, i.e., \texttt{=} evaluates in this case to \texttt{false}, and \texttt{{\textless}{\textgreater}} evaluates to \texttt{true}. 

 \index{smaller test} \texttt{\mbox{\texttt{\mdseries\slshape left\texttt{\symbol{45}}expr}} {\textless} \mbox{\texttt{\mdseries\slshape right\texttt{\symbol{45}}expr}}} 

 \index{larger test} \texttt{\mbox{\texttt{\mdseries\slshape left\texttt{\symbol{45}}expr}} {\textgreater} \mbox{\texttt{\mdseries\slshape right\texttt{\symbol{45}}expr}}} 

 \index{smaller or equal} \texttt{\mbox{\texttt{\mdseries\slshape left\texttt{\symbol{45}}expr}} {\textless}= \mbox{\texttt{\mdseries\slshape right\texttt{\symbol{45}}expr}}} 

 \index{larger or equal} \texttt{\mbox{\texttt{\mdseries\slshape left\texttt{\symbol{45}}expr}} {\textgreater}= \mbox{\texttt{\mdseries\slshape right\texttt{\symbol{45}}expr}}} 

 \texttt{{\textless}} denotes less than, \texttt{{\textless}=} less than or equal, \texttt{{\textgreater}} greater than, and \texttt{{\textgreater}=} greater than or equal of its two operands. For each kind of objects the
definition of the ordering is given in the respective chapter. 

 Note that \texttt{{\textless}} implements a \emph{total ordering} of objects (which can be used for example to sort a list of elements).
Therefore in general \texttt{{\textless}} will not be compatible with any inclusion relation (which can be tested using \texttt{IsSubset} (\ref{IsSubset})). (For example, it is possible to compare permutation groups with \texttt{{\textless}} in a total ordering of all permutation groups, but this ordering is not
compatible with the relation of being a subgroup.) 

 Only for the following kinds of objects, an ordering via \texttt{{\textless}} of objects in \emph{different} families (see{\nobreakspace}\ref{Families}) is supported. Rationals (see{\nobreakspace}\texttt{IsRat} (\ref{IsRat})) are smallest, next are cyclotomics (see{\nobreakspace}\texttt{IsCyclotomic} (\ref{IsCyclotomic})), followed by finite field elements (see{\nobreakspace}\texttt{IsFFE} (\ref{IsFFE})); finite field elements in different characteristics are compared via their
characteristics, next are permutations (see{\nobreakspace}\texttt{IsPerm} (\ref{IsPerm})), followed by the boolean values \texttt{true}, \texttt{false}, and \texttt{fail} (see{\nobreakspace}\texttt{IsBool} (\ref{IsBool})), characters (such as \texttt{\texttt{\symbol{123}}}\texttt{\symbol{125}}a\texttt{\symbol{123}}'\texttt{\symbol{125}}',
see{\nobreakspace}\texttt{IsChar} (\ref{IsChar})), and lists (see{\nobreakspace}\texttt{IsList} (\ref{IsList})) are largest; note that two lists can be compared with \texttt{{\textless}} if and only if their elements are again objects that can be compared with \texttt{{\textless}}. 

 For other objects, \textsf{GAP} does \emph{not} provide an ordering via \texttt{{\textless}}. The reason for this is that a total ordering of all \textsf{GAP} objects would be hard to maintain when new kinds of objects are introduced,
and such a total ordering is hardly used in its full generality. 

 However, for objects in the filters listed above, the ordering via \texttt{{\textless}} has turned out to be useful. For example, one can form \emph{sorted lists} containing integers and nested lists of integers, and then search in them
using \texttt{PositionSorted} (see{\nobreakspace}\ref{Finding Positions in Lists}). 

 Of course it would in principle be possible to define an ordering via \texttt{{\textless}} also for certain other objects, by installing appropriate methods for the
operation \texttt{\texttt{\symbol{92}}{\textless}}. But this may lead to problems at least as soon as one loads \textsf{GAP} code in which the same is done, under the assumption that one is completely
free to define an ordering via \texttt{{\textless}} for other objects than the ones for which the ``official'' \textsf{GAP} provides already an ordering via \texttt{{\textless}}. 

 Comparison operators, including the operator \texttt{in} (see{\nobreakspace}\ref{Membership Test for Lists}), are not associative, Hence it is not allowed to write \texttt{\mbox{\texttt{\mdseries\slshape a}} = \mbox{\texttt{\mdseries\slshape b}} {\textless}{\textgreater} \mbox{\texttt{\mdseries\slshape c}} = \mbox{\texttt{\mdseries\slshape d}}}, you must use \texttt{(\mbox{\texttt{\mdseries\slshape a}} = \mbox{\texttt{\mdseries\slshape b}}) {\textless}{\textgreater} (\mbox{\texttt{\mdseries\slshape c}} = \mbox{\texttt{\mdseries\slshape d}})} instead. The comparison operators have higher precedence than the logical
operators \index{operators!precedence} (see{\nobreakspace}\ref{Operations for Booleans}), but lower precedence than the arithmetic operators (see{\nobreakspace}\ref{Arithmetic Operators}). Thus, for instance, \texttt{\mbox{\texttt{\mdseries\slshape a}} * \mbox{\texttt{\mdseries\slshape b}} = \mbox{\texttt{\mdseries\slshape c}} and \mbox{\texttt{\mdseries\slshape d}}} is interpreted as \texttt{((\mbox{\texttt{\mdseries\slshape a}} * \mbox{\texttt{\mdseries\slshape b}}) = \mbox{\texttt{\mdseries\slshape c}}) and \mbox{\texttt{\mdseries\slshape d}})}. 

 The following example shows a comparison where the left operand is an
expression. 

 
\begin{Verbatim}[commandchars=!@|,fontsize=\small,frame=single,label=Example]
  !gapprompt@gap>| !gapinput@2 * 2 + 9 = Fibonacci(7);|
  true
\end{Verbatim}
 

 For the underlying operations of the operators introduced above,
see{\nobreakspace}\ref{Comparison Operations for Elements}. }

  
\section{\textcolor{Chapter }{Arithmetic Operators}}\label{Arithmetic Operators}
\logpage{[ 4, 14, 0 ]}
\hyperdef{L}{X7B66C8707B5DE10A}{}
{
  \index{precedence} \index{associativity} \index{operators!arithmetic} \index{+@\texttt{+}} \index{-@\texttt{\texttt{\symbol{45}}}} \index{*@\texttt{*}} \index{/@\texttt{/}} \index{^@\texttt{\texttt{\symbol{94}}}} \index{mod@\texttt{mod}!arithmetic operators} \index{modulo} \index{modulo!arithmetic operators} \index{positive number} \texttt{+ \mbox{\texttt{\mdseries\slshape right\texttt{\symbol{45}}expr}}} 

 \index{negative number} \texttt{\texttt{\symbol{45}} \mbox{\texttt{\mdseries\slshape right\texttt{\symbol{45}}expr}}} 

 \index{addition} \texttt{\mbox{\texttt{\mdseries\slshape left\texttt{\symbol{45}}expr}} + \mbox{\texttt{\mdseries\slshape right\texttt{\symbol{45}}expr}}} 

 \index{subtraction} \texttt{\mbox{\texttt{\mdseries\slshape left\texttt{\symbol{45}}expr}} \texttt{\symbol{45}} \mbox{\texttt{\mdseries\slshape right\texttt{\symbol{45}}expr}}} 

 \index{multiplication} \texttt{\mbox{\texttt{\mdseries\slshape left\texttt{\symbol{45}}expr}} * \mbox{\texttt{\mdseries\slshape right\texttt{\symbol{45}}expr}}} 

 \index{division} \texttt{\mbox{\texttt{\mdseries\slshape left\texttt{\symbol{45}}expr}} / \mbox{\texttt{\mdseries\slshape right\texttt{\symbol{45}}expr}}} 

 \index{mod@\texttt{mod}} \texttt{\mbox{\texttt{\mdseries\slshape left\texttt{\symbol{45}}expr}} mod \mbox{\texttt{\mdseries\slshape right\texttt{\symbol{45}}expr}}} 

 \index{power} \texttt{\mbox{\texttt{\mdseries\slshape left\texttt{\symbol{45}}expr}} \texttt{\symbol{94}} \mbox{\texttt{\mdseries\slshape right\texttt{\symbol{45}}expr}}} 

 The arithmetic operators are \texttt{+}, \texttt{\texttt{\symbol{45}}}, \texttt{*}, \texttt{/}, \texttt{mod}, and \texttt{\texttt{\symbol{94}}}. The meanings (semantics) of those operators generally depend on the types of
the operands involved, and they are defined in the various chapters describing
the types. However basically the meanings are as follows. 

 \texttt{\mbox{\texttt{\mdseries\slshape a}} + \mbox{\texttt{\mdseries\slshape b}}} denotes the addition of additive elements \mbox{\texttt{\mdseries\slshape a}} and \mbox{\texttt{\mdseries\slshape b}}. 

 \texttt{\mbox{\texttt{\mdseries\slshape a}} \texttt{\symbol{45}} \mbox{\texttt{\mdseries\slshape b}}} denotes the addition of \mbox{\texttt{\mdseries\slshape a}} and the additive inverse of \mbox{\texttt{\mdseries\slshape b}}. 

 \texttt{\mbox{\texttt{\mdseries\slshape a}} * \mbox{\texttt{\mdseries\slshape b}}} denotes the multiplication of multiplicative elements \mbox{\texttt{\mdseries\slshape a}} and \mbox{\texttt{\mdseries\slshape b}}. 

 \texttt{\mbox{\texttt{\mdseries\slshape a}} / \mbox{\texttt{\mdseries\slshape b}}} denotes the multiplication of \mbox{\texttt{\mdseries\slshape a}} with the multiplicative inverse of \mbox{\texttt{\mdseries\slshape b}}. 

 \index{mod!rationals} \texttt{\mbox{\texttt{\mdseries\slshape a}} mod \mbox{\texttt{\mdseries\slshape b}}}, for integer or rational left operand \mbox{\texttt{\mdseries\slshape a}} and for non\texttt{\symbol{45}}zero integer right operand \mbox{\texttt{\mdseries\slshape b}}, is defined as follows. If \mbox{\texttt{\mdseries\slshape a}} and \mbox{\texttt{\mdseries\slshape b}} are both integers, \texttt{\mbox{\texttt{\mdseries\slshape a}} mod \mbox{\texttt{\mdseries\slshape b}}} is the integer \mbox{\texttt{\mdseries\slshape r}} in the integer range \texttt{0 .. |\mbox{\texttt{\mdseries\slshape b}}| \texttt{\symbol{45}} 1} satisfying \texttt{\mbox{\texttt{\mdseries\slshape a}} = \mbox{\texttt{\mdseries\slshape r}} + \mbox{\texttt{\mdseries\slshape b}}\mbox{\texttt{\mdseries\slshape q}}}, for some integer \mbox{\texttt{\mdseries\slshape q}} (where the operations occurring have their usual meaning over the integers, of
course). 

 \index{modular remainder}\index{modular inverse} \index{coprime}\index{relatively prime} If \mbox{\texttt{\mdseries\slshape a}} is a rational number and \mbox{\texttt{\mdseries\slshape b}} is a non\texttt{\symbol{45}}zero integer, and \texttt{\mbox{\texttt{\mdseries\slshape a}} = \mbox{\texttt{\mdseries\slshape m}} / \mbox{\texttt{\mdseries\slshape n}}} where \mbox{\texttt{\mdseries\slshape m}} and \mbox{\texttt{\mdseries\slshape n}} are coprime integers with \mbox{\texttt{\mdseries\slshape n}} positive, then \texttt{\mbox{\texttt{\mdseries\slshape a}} mod \mbox{\texttt{\mdseries\slshape b}}} is the integer \mbox{\texttt{\mdseries\slshape r}} in the integer range \texttt{0 .. |\mbox{\texttt{\mdseries\slshape b}}| \texttt{\symbol{45}} 1} such that \mbox{\texttt{\mdseries\slshape m}} is congruent to \texttt{\mbox{\texttt{\mdseries\slshape r}}\mbox{\texttt{\mdseries\slshape n}}} modulo \mbox{\texttt{\mdseries\slshape b}}, and \mbox{\texttt{\mdseries\slshape r}} is called the ``modular remainder'' of \mbox{\texttt{\mdseries\slshape a}} modulo \mbox{\texttt{\mdseries\slshape b}}. Also, \texttt{1 / \mbox{\texttt{\mdseries\slshape n}} mod \mbox{\texttt{\mdseries\slshape b}}} is called the ``modular inverse'' of \mbox{\texttt{\mdseries\slshape n}} modulo \mbox{\texttt{\mdseries\slshape b}}. (A pair of integers is said to be \emph{coprime} (or \emph{relatively prime}) if their greatest common divisor is 1.) 

 With the above definition, \texttt{4 / 6 mod 32} equals \texttt{2 / 3 mod 32} and hence exists (and is equal to 22), despite the fact that 6 has no inverse
modulo 32. 

 \emph{Note:} For rational \mbox{\texttt{\mdseries\slshape a}}, \texttt{\mbox{\texttt{\mdseries\slshape a}} mod \mbox{\texttt{\mdseries\slshape b}}} could have been defined to be the non\texttt{\symbol{45}}negative rational \mbox{\texttt{\mdseries\slshape c}} less than \texttt{|\mbox{\texttt{\mdseries\slshape b}}|} such that \texttt{\mbox{\texttt{\mdseries\slshape a}} \texttt{\symbol{45}} \mbox{\texttt{\mdseries\slshape c}}} is a multiple of \mbox{\texttt{\mdseries\slshape b}}. However this definition is seldom useful and \emph{not} the one chosen for \textsf{GAP}. 

 \texttt{+} and \texttt{\texttt{\symbol{45}}} can also be used as unary operations. The unary \texttt{+} is ignored. The unary \texttt{\texttt{\symbol{45}}} returns the additive inverse of its operand; over the integers it is
equivalent to multiplication by \texttt{\texttt{\symbol{45}}1}. 

 \index{^@\texttt{\texttt{\symbol{94}}}!for two group elements} \index{conjugation!with a group element} \texttt{\texttt{\symbol{94}}} denotes powering of a multiplicative element if the right operand is an
integer, and is also used to denote the action of a group element on a point
of a set if the right operand is a group element. In the special case that
both operands are group elements, \texttt{\texttt{\symbol{94}}} denotes conjugation, that is, $g$\texttt{\texttt{\symbol{94}}}$h = h^{{-1}} g h$. 

 \index{arithmetic operators!precedence} The \emph{precedence} of those operators is as follows. The powering operator \texttt{\texttt{\symbol{94}}} has the highest precedence, followed by the unary operators \texttt{+} and \texttt{\texttt{\symbol{45}}}, which are followed by the multiplicative operators \texttt{*}, \texttt{/}, and \texttt{mod}, and the additive binary operators \texttt{+} and \texttt{\texttt{\symbol{45}}} have the lowest precedence. That means that the expression \texttt{\texttt{\symbol{45}}2 \texttt{\symbol{94}} \texttt{\symbol{45}}2 * 3 + 1} is interpreted as \texttt{(\texttt{\symbol{45}}(2 \texttt{\symbol{94}} (\texttt{\symbol{45}}2)) * 3) + 1}. If in doubt use parentheses to clarify your intention. 

 \index{operators!associativity} The \emph{associativity} of the arithmetic operators is as follows. \texttt{\texttt{\symbol{94}}} is not associative, i.e., it is invalid to write \texttt{2\texttt{\symbol{94}}3\texttt{\symbol{94}}4}, use parentheses to clarify whether you mean \texttt{(2\texttt{\symbol{94}}3)\texttt{\symbol{94}}4} or \texttt{2\texttt{\symbol{94}}(3\texttt{\symbol{94}}4)}. The unary operators \texttt{+} and \texttt{\texttt{\symbol{45}}} are right associative, because they are written to the left of their operands. \texttt{*}, \texttt{/}, \texttt{mod}, \texttt{+}, and \texttt{\texttt{\symbol{45}}} are all left associative, i.e., \texttt{1\texttt{\symbol{45}}2\texttt{\symbol{45}}3} is interpreted as \texttt{(1\texttt{\symbol{45}}2)\texttt{\symbol{45}}3} not as \texttt{1\texttt{\symbol{45}}(2\texttt{\symbol{45}}3)}. Again, if in doubt use parentheses to clarify your intentions. 

 The arithmetic operators have higher precedence than the comparison operators
(see{\nobreakspace}\ref{Comparisons} and{\nobreakspace}\ref{Membership Test for Collections}) and the logical operators (see \ref{Operations for Booleans}). Thus, for example, \texttt{\mbox{\texttt{\mdseries\slshape a}} * \mbox{\texttt{\mdseries\slshape b}} = \mbox{\texttt{\mdseries\slshape c}} and \mbox{\texttt{\mdseries\slshape d}}} is interpreted, \texttt{((\mbox{\texttt{\mdseries\slshape a}} * \mbox{\texttt{\mdseries\slshape b}}) = \mbox{\texttt{\mdseries\slshape c}}) and \mbox{\texttt{\mdseries\slshape d}}}. 

 
\begin{Verbatim}[commandchars=!@|,fontsize=\small,frame=single,label=Example]
  !gapprompt@gap>| !gapinput@2 * 2 + 9;  # a very simple arithmetic expression|
  13
\end{Verbatim}
 

 For other arithmetic operations, and for the underlying operations of the
operators introduced above, see{\nobreakspace}\ref{Arithmetic Operations for Elements}. }

  
\section{\textcolor{Chapter }{Statements}}\label{Statements}
\logpage{[ 4, 15, 0 ]}
\hyperdef{L}{X8543285D87361BE6}{}
{
  \index{execution} \textsf{GAP} programs consist of a sequence of so\texttt{\symbol{45}}called \emph{statements}. The following types of statements exist: 
\begin{itemize}
\item Assignments (see Section{\nobreakspace}\ref{Assignments}),
\item Procedure calls (see Section{\nobreakspace}\ref{Procedure Calls}),
\item \texttt{if} statements (see Section{\nobreakspace}\ref{If}),
\item \texttt{while} loops (see Section{\nobreakspace}\ref{While}),
\item \texttt{repeat} loops (see Section{\nobreakspace}\ref{Repeat}),
\item \texttt{for} loops (see Section{\nobreakspace}\ref{For}),
\item \texttt{break} statements (see Section{\nobreakspace}\ref{Break}),
\item \texttt{continue} statements (see Section{\nobreakspace}\ref{Continue}), and
\item \texttt{return} statements (see Section{\nobreakspace}\ref{Return}).
\end{itemize}
 They can be entered interactively or be part of a function definition. Every
statement must be terminated by a semicolon. 

 Statements, unlike expressions, have no value. They are executed only to
produce an effect. For example an assignment has the effect of assigning a
value to a variable, a \texttt{for} loop has the effect of executing a statement sequence for all elements in a
list and so on. We will talk about \emph{evaluation} of expressions but about \emph{execution} of statements to emphasize this difference. 

 Using expressions as statements is treated as syntax error. 

 
\begin{Verbatim}[commandchars=!@|,fontsize=\small,frame=single,label=Example]
  !gapprompt@gap>| !gapinput@i := 7;;|
  !gapprompt@gap>| !gapinput@if i <> 0 then k = 16/i; fi;|
  Syntax error: := expected
  if i <> 0 then k = 16/i; fi;
                   ^
  gap>
\end{Verbatim}
 

 As you can see from the example this warning does in particular address those
users who are used to languages where \texttt{=} instead of \texttt{:=} denotes assignment. 

 Empty statements are permitted and have no effect. 

 A sequence of one or more statements is a \emph{statement sequence}, and may occur everywhere instead of a single statement. Each construct is
terminated by a keyword. The simplest statement sequence is a single
semicolon, which can be used as an empty statement sequence. In fact an empty
statement sequence as in \texttt{for i in [ 1 .. 2 ] do od} is also permitted and is silently translated into the sequence containing just
a semicolon.  
\subsection{\textcolor{Chapter }{Assignments}}\label{Assignments}
\logpage{[ 4, 15, 1 ]}
\hyperdef{L}{X7E6A50307F4D3FAE}{}
{
  \index{assignment!variable} \texttt{\mbox{\texttt{\mdseries\slshape var}} := \mbox{\texttt{\mdseries\slshape expr}};} 

 The \emph{assignment} has the effect of assigning the value of the expressions \mbox{\texttt{\mdseries\slshape expr}} to the variable \mbox{\texttt{\mdseries\slshape var}}. 

 The variable \mbox{\texttt{\mdseries\slshape var}} may be an ordinary variable (see \ref{Variables}), a list element selection \texttt{\mbox{\texttt{\mdseries\slshape list\texttt{\symbol{45}}var}}[\mbox{\texttt{\mdseries\slshape int\texttt{\symbol{45}}expr}}]} (see \ref{List Assignment}) or a record component selection \texttt{\mbox{\texttt{\mdseries\slshape record\texttt{\symbol{45}}var}}.\mbox{\texttt{\mdseries\slshape ident}}} (see \ref{Record Assignment}). Since a list element or a record component may itself be a list or a record
the left hand side of an assignment may be arbitrarily complex. 

 Note that variables do not have a type. Thus any value may be assigned to any
variable. For example a variable with an integer value may be assigned a
permutation or a list or anything else. 

 
\begin{Verbatim}[commandchars=!@|,fontsize=\small,frame=single,label=Example]
  !gapprompt@gap>| !gapinput@data:= rec( numbers:= [ 1, 2, 3 ] );|
  rec( numbers := [ 1, 2, 3 ] )
  !gapprompt@gap>| !gapinput@data.string:= "string";; data;|
  rec( numbers := [ 1, 2, 3 ], string := "string" )
  !gapprompt@gap>| !gapinput@data.numbers[2]:= 4;; data;|
  rec( numbers := [ 1, 4, 3 ], string := "string" )
\end{Verbatim}
 

 If the expression \mbox{\texttt{\mdseries\slshape expr}} is a function call then this function must return a value. If the function
does not return a value an error is signalled and you enter a break loop (see \ref{Break Loops}). As usual you can leave the break loop with \texttt{quit;}. 

 
\begin{Verbatim}[commandchars=!|A,fontsize=\small,frame=single,label=Example]
  !gapprompt|gap>A !gapinput|f1:= function( x ) Print( "value: ", x, "\n" ); end;;A
  !gapprompt|gap>A !gapinput|f2:= function( x ) return f1( x ); end;;A
  !gapprompt|gap>A !gapinput|f2( 4 );A
  value: 4
  Error, Function Calls: <func> must return a value
  Stack trace:
  *[1] f1( x )
     @ *stdin*:2
  <function "f2">( <arguments> )
   called from read-eval loop at *stdin*:3
  you can enter 'quit;' to quit to outer loop
  !gapbrkprompt|brk>A !gapinput|A
\end{Verbatim}
 

 In the above example, the function \texttt{f2} calls \texttt{f1} with argument \texttt{4}, and since \texttt{f1} does not return a value (but only prints a line ``\texttt{value: ...}''), the \texttt{return} statement of \texttt{f2} cannot be executed. The error message says that it is possible to return an
appropriate value, and the returned string \texttt{"hello"} is used by \texttt{f2} instead of the missing return value of \texttt{f1}. }

  
\subsection{\textcolor{Chapter }{Procedure Calls}}\label{Procedure Calls}
\logpage{[ 4, 15, 2 ]}
\hyperdef{L}{X825803DE78251DA6}{}
{
  \index{procedure call} \index{procedure call with arguments} \texttt{\mbox{\texttt{\mdseries\slshape procedure\texttt{\symbol{45}}var}}( [\mbox{\texttt{\mdseries\slshape arg\texttt{\symbol{45}}expr}} [,\mbox{\texttt{\mdseries\slshape arg\texttt{\symbol{45}}expr}}, ...]] );} 

 The \emph{procedure call} has the effect of calling the procedure \mbox{\texttt{\mdseries\slshape procedure\texttt{\symbol{45}}var}}. A procedure call is done exactly like a function call (see \ref{Function Calls}). The distinction between functions and procedures is only for the sake of
the discussion, \textsf{GAP} does not distinguish between them. So we state the following conventions. 

 A \emph{function} does return a value but does not produce a side effect. As a convention the
name of a function is a noun, denoting what the function returns, e.g., \texttt{"Length"}, \texttt{"Concatenation"} and \texttt{"Order"}. 

 A \emph{procedure} is a function that does not return a value but produces some effect.
Procedures are called only for this effect. As a convention the name of a
procedure is a verb, denoting what the procedure does, e.g., \texttt{"Print"}, \texttt{"Append"} and \texttt{"Sort"}. 

 
\begin{Verbatim}[commandchars=!@|,fontsize=\small,frame=single,label=Example]
  !gapprompt@gap>| !gapinput@Read( "myfile.g" );   # a call to the procedure Read|
  !gapprompt@gap>| !gapinput@l := [ 1, 2 ];;|
  !gapprompt@gap>| !gapinput@Append( l, [3,4,5] );  # a call to the procedure Append|
\end{Verbatim}
 

 There are a few exceptions of \textsf{GAP} functions that do both return a value and produce some effect. An example is \texttt{Sortex} (\ref{Sortex}) which sorts a list and returns the corresponding permutation of the entries. }

  
\subsection{\textcolor{Chapter }{If}}\label{If}
\logpage{[ 4, 15, 3 ]}
\hyperdef{L}{X875000188622700D}{}
{
  \index{fi@\texttt{fi}} \index{then@\texttt{then}} \index{else@\texttt{else}} \index{elif@\texttt{elif}} \texttt{if \mbox{\texttt{\mdseries\slshape bool\texttt{\symbol{45}}expr1}} then \mbox{\texttt{\mdseries\slshape statements1}} \texttt{\symbol{123}} elif \mbox{\texttt{\mdseries\slshape bool\texttt{\symbol{45}}expr2}} then \mbox{\texttt{\mdseries\slshape statements2}} \texttt{\symbol{125}}[ else \mbox{\texttt{\mdseries\slshape statements3}} ] fi;} \index{if statement@\texttt{if} statement} 

 The \texttt{if} statement allows one to execute statements depending on the value of some
boolean expression. The execution is done as follows. 

 First the expression \mbox{\texttt{\mdseries\slshape bool\texttt{\symbol{45}}expr1}} following the \texttt{if} is evaluated. If it evaluates to \texttt{true} the statement sequence \mbox{\texttt{\mdseries\slshape statements1}} after the first \texttt{then} is executed, and the execution of the \texttt{if} statement is complete. 

 Otherwise the expressions \mbox{\texttt{\mdseries\slshape bool\texttt{\symbol{45}}expr2}} following the \texttt{elif} are evaluated in turn. There may be any number of \texttt{elif} parts, possibly none at all. As soon as an expression evaluates to \texttt{true} the corresponding statement sequence \mbox{\texttt{\mdseries\slshape statements2}} is executed and execution of the \texttt{if} statement is complete. 

 If the \texttt{if} expression and all, if any, \texttt{elif} expressions evaluate to \texttt{false} and there is an \texttt{else} part, which is optional, its statement sequence \mbox{\texttt{\mdseries\slshape statements3}} is executed and the execution of the \texttt{if} statement is complete. If there is no \texttt{else} part the \texttt{if} statement is complete without executing any statement sequence. 

 Since the \texttt{if} statement is terminated by the \texttt{fi} keyword there is no question where an \texttt{else} part belongs, i.e., \textsf{GAP} has no ``dangling else''. In 

 
\begin{Verbatim}[commandchars=!@|,fontsize=\small,frame=single,label=]
  if expr1 then if expr2 then stats1 else stats2 fi; fi;
\end{Verbatim}
 

 the \texttt{else} part belongs to the second \texttt{if} statement, whereas in 

 
\begin{Verbatim}[commandchars=!@|,fontsize=\small,frame=single,label=]
  if expr1 then if expr2 then stats1 fi; else stats2 fi;
\end{Verbatim}
 

 the \texttt{else} part belongs to the first \texttt{if} statement. 

 Since an \texttt{if} statement is not an expression it is not possible to write 

 
\begin{Verbatim}[commandchars=!@|,fontsize=\small,frame=single,label=]
  abs := if x > 0 then x; else -x; fi;
\end{Verbatim}
 

 which would, even if legal syntax, be meaningless, since the \texttt{if} statement does not produce a value that could be assigned to \texttt{abs}. 

 If one of the expressions \mbox{\texttt{\mdseries\slshape bool\texttt{\symbol{45}}expr1}}, \mbox{\texttt{\mdseries\slshape bool\texttt{\symbol{45}}expr2}} is evaluated and its value is neither \texttt{true} nor \texttt{false} an error is signalled and a break loop (see \ref{Break Loops}) is entered. As usual you can leave the break loop with \texttt{quit;}. If you enter \texttt{return true;}, execution of the \texttt{if} statement continues as if the expression whose evaluation failed had evaluated
to \texttt{true}. Likewise, if you enter \texttt{return false;}, execution of the \texttt{if} statement continues as if the expression whose evaluation failed had evaluated
to \texttt{false}. 

 
\begin{Verbatim}[commandchars=!@|,fontsize=\small,frame=single,label=Example]
  !gapprompt@gap>| !gapinput@i := 10;;|
  !gapprompt@gap>| !gapinput@if 0 < i then|
  !gapprompt@>| !gapinput@   s := 1;|
  !gapprompt@>| !gapinput@ elif i < 0 then|
  !gapprompt@>| !gapinput@   s := -1;|
  !gapprompt@>| !gapinput@ else|
  !gapprompt@>| !gapinput@   s := 0;|
  !gapprompt@>| !gapinput@ fi;|
  !gapprompt@gap>| !gapinput@s;  # the sign of i|
  1
\end{Verbatim}
 }

  
\subsection{\textcolor{Chapter }{While}}\label{While}
\logpage{[ 4, 15, 4 ]}
\hyperdef{L}{X87AA46408783383F}{}
{
  \index{loop!while} \index{while loop@\texttt{while} loop} \texttt{while \mbox{\texttt{\mdseries\slshape bool\texttt{\symbol{45}}expr}} do \mbox{\texttt{\mdseries\slshape statements}} od;} 

 The \texttt{while} loop executes the statement sequence \mbox{\texttt{\mdseries\slshape statements}} while the condition \mbox{\texttt{\mdseries\slshape bool\texttt{\symbol{45}}expr}} evaluates to \texttt{true}. 

 First \mbox{\texttt{\mdseries\slshape bool\texttt{\symbol{45}}expr}} is evaluated. If it evaluates to \texttt{false} execution of the \texttt{while} loop terminates and the statement immediately following the \texttt{while} loop is executed next. Otherwise if it evaluates to \texttt{true} the \mbox{\texttt{\mdseries\slshape statements}} are executed and the whole process begins again. 

 The difference between the \texttt{while} loop and the \texttt{repeat} \texttt{until} loop (see \ref{Repeat}) is that the \mbox{\texttt{\mdseries\slshape statements}} in the \texttt{repeat} \texttt{until} loop are executed at least once, while the \mbox{\texttt{\mdseries\slshape statements}} in the \texttt{while} loop are not executed at all if \mbox{\texttt{\mdseries\slshape bool\texttt{\symbol{45}}expr}} is \texttt{false} at the first iteration. 

 If \mbox{\texttt{\mdseries\slshape bool\texttt{\symbol{45}}expr}} does not evaluate to \texttt{true} or \texttt{false} an error is signalled and a break loop (see \ref{Break Loops}) is entered. As usual you can leave the break loop with \texttt{quit;}. If you enter \texttt{return false;}, execution continues with the next statement immediately following the \texttt{while} loop. If you enter \texttt{return true;}, execution continues at \mbox{\texttt{\mdseries\slshape statements}}, after which the next evaluation of \mbox{\texttt{\mdseries\slshape bool\texttt{\symbol{45}}expr}} may cause another error. 

 The following example shows a \texttt{while} loop that sums up the squares $1^2, 2^2, \ldots$ until the sum exceeds $200$. 

 
\begin{Verbatim}[commandchars=!@|,fontsize=\small,frame=single,label=Example]
  !gapprompt@gap>| !gapinput@i := 0;; s := 0;;|
  !gapprompt@gap>| !gapinput@while s <= 200 do|
  !gapprompt@>| !gapinput@   i := i + 1; s := s + i^2;|
  !gapprompt@>| !gapinput@ od;|
  !gapprompt@gap>| !gapinput@s;|
  204
\end{Verbatim}
 

 A \texttt{while} loop may be left prematurely using \texttt{break}, see \ref{Break}. }

  
\subsection{\textcolor{Chapter }{Repeat}}\label{Repeat}
\logpage{[ 4, 15, 5 ]}
\hyperdef{L}{X8295CBF47FAA05C9}{}
{
  \index{loop!repeat} \index{until@\texttt{until}} \index{repeat loop@\texttt{repeat} loop} \texttt{repeat \mbox{\texttt{\mdseries\slshape statements}} until \mbox{\texttt{\mdseries\slshape bool\texttt{\symbol{45}}expr}};} 

 The \texttt{repeat} loop executes the statement sequence \mbox{\texttt{\mdseries\slshape statements}} until the condition \mbox{\texttt{\mdseries\slshape bool\texttt{\symbol{45}}expr}} evaluates to \texttt{true}. 

 First \mbox{\texttt{\mdseries\slshape statements}} are executed. Then \mbox{\texttt{\mdseries\slshape bool\texttt{\symbol{45}}expr}} is evaluated. If it evaluates to \texttt{true} the \texttt{repeat} loop terminates and the statement immediately following the \texttt{repeat} loop is executed next. Otherwise if it evaluates to \texttt{false} the whole process begins again with the execution of the \mbox{\texttt{\mdseries\slshape statements}}. 

 The difference between the \texttt{while} loop (see \ref{While}) and the \texttt{repeat} \texttt{until} loop is that the \mbox{\texttt{\mdseries\slshape statements}} in the \texttt{repeat} \texttt{until} loop are executed at least once, while the \mbox{\texttt{\mdseries\slshape statements}} in the \texttt{while} loop are not executed at all if \mbox{\texttt{\mdseries\slshape bool\texttt{\symbol{45}}expr}} is \texttt{false} at the first iteration. 

 If \mbox{\texttt{\mdseries\slshape bool\texttt{\symbol{45}}expr}} does not evaluate to \texttt{true} or \texttt{false} an error is signalled and a break loop (see \ref{Break Loops}) is entered. As usual you can leave the break loop with \texttt{quit;}. If you enter \texttt{return true;}, execution continues with the next statement immediately following the \texttt{repeat} loop. If you enter \texttt{return false;}, execution continues at \mbox{\texttt{\mdseries\slshape statements}}, after which the next evaluation of \mbox{\texttt{\mdseries\slshape bool\texttt{\symbol{45}}expr}} may cause another error. 

 The \texttt{repeat} loop in the following example has the same purpose as the \texttt{while} loop in the preceding example, namely to sum up the squares $1^2, 2^2, \ldots$ until the sum exceeds $200$. 

 
\begin{Verbatim}[commandchars=!@|,fontsize=\small,frame=single,label=Example]
  !gapprompt@gap>| !gapinput@i := 0;; s := 0;;|
  !gapprompt@gap>| !gapinput@repeat|
  !gapprompt@>| !gapinput@   i := i + 1; s := s + i^2;|
  !gapprompt@>| !gapinput@ until s > 200;|
  !gapprompt@gap>| !gapinput@s;|
  204
\end{Verbatim}
 

 A \texttt{repeat} loop may be left prematurely using \texttt{break}, see \ref{Break}. }

  
\subsection{\textcolor{Chapter }{For}}\label{For}
\logpage{[ 4, 15, 6 ]}
\hyperdef{L}{X78783E777867638A}{}
{
  \index{loop!for} \index{do@\texttt{do}} \index{od@\texttt{od}} \texttt{for \mbox{\texttt{\mdseries\slshape simple\texttt{\symbol{45}}var}} in \mbox{\texttt{\mdseries\slshape list\texttt{\symbol{45}}expr}} do \mbox{\texttt{\mdseries\slshape statements}} od;} \index{for loop@\texttt{for} loop} 

 The \texttt{for} loop executes the statement sequence \mbox{\texttt{\mdseries\slshape statements}} for every element of the list \mbox{\texttt{\mdseries\slshape list\texttt{\symbol{45}}expr}}. 

 The statement sequence \mbox{\texttt{\mdseries\slshape statements}} is first executed with \mbox{\texttt{\mdseries\slshape simple\texttt{\symbol{45}}var}} bound to the first element of the list \mbox{\texttt{\mdseries\slshape list\texttt{\symbol{45}}expr}}, then with \mbox{\texttt{\mdseries\slshape simple\texttt{\symbol{45}}var}} bound to the second element of \mbox{\texttt{\mdseries\slshape list\texttt{\symbol{45}}expr}} and so on. \mbox{\texttt{\mdseries\slshape simple\texttt{\symbol{45}}var}} must be a simple variable, it must not be a list element selection \texttt{\mbox{\texttt{\mdseries\slshape list\texttt{\symbol{45}}var}}[\mbox{\texttt{\mdseries\slshape int\texttt{\symbol{45}}expr}}]} or a record component selection \texttt{\mbox{\texttt{\mdseries\slshape record\texttt{\symbol{45}}var}}.\mbox{\texttt{\mdseries\slshape ident}}}. 

 The execution of the \texttt{for} loop over a list is exactly equivalent to the following \texttt{while} loop. 

 
\begin{Verbatim}[commandchars=!@|,fontsize=\small,frame=single,label=]
  loop_list := list;
  loop_index := 1;
  while loop_index <= Length(loop_list) do
    variable := loop_list[loop_index];
    statements
    loop_index := loop_index + 1;
  od;
\end{Verbatim}
 

 with the exception that ``loop{\textunderscore}list'' and ``loop{\textunderscore}index'' are different variables for each \texttt{for} loop, i.e., these variables of different \texttt{for} loops do not interfere with each other. 

 The list \mbox{\texttt{\mdseries\slshape list\texttt{\symbol{45}}expr}} is very often a range (see{\nobreakspace}\ref{Ranges}). 

 

 \index{loop over range} \texttt{for \mbox{\texttt{\mdseries\slshape variable}} in [\mbox{\texttt{\mdseries\slshape from}}..\mbox{\texttt{\mdseries\slshape to}}] do \mbox{\texttt{\mdseries\slshape statements}} od;} 

 corresponds to the more common 

 \texttt{for \mbox{\texttt{\mdseries\slshape variable}} from \mbox{\texttt{\mdseries\slshape from}} to \mbox{\texttt{\mdseries\slshape to}} do \mbox{\texttt{\mdseries\slshape statements}} od;} 

 in other programming languages. 

 
\begin{Verbatim}[commandchars=!@|,fontsize=\small,frame=single,label=Example]
  !gapprompt@gap>| !gapinput@s := 0;;|
  !gapprompt@gap>| !gapinput@for i in [1..100] do|
  !gapprompt@>| !gapinput@   s := s + i;|
  !gapprompt@>| !gapinput@od;|
  !gapprompt@gap>| !gapinput@s;|
  5050
\end{Verbatim}
 

 Note in the following example how the modification of the \emph{list} in the loop body causes the loop body also to be executed for the new values. 

 
\begin{Verbatim}[commandchars=!@|,fontsize=\small,frame=single,label=Example]
  !gapprompt@gap>| !gapinput@l := [ 1, 2, 3, 4, 5, 6 ];;|
  !gapprompt@gap>| !gapinput@for i in l do|
  !gapprompt@>| !gapinput@   Print( i, " " );|
  !gapprompt@>| !gapinput@   if i mod 2 = 0 then Add( l, 3 * i / 2 ); fi;|
  !gapprompt@>| !gapinput@od; Print( "\n" );|
  1 2 3 4 5 6 3 6 9 9
  !gapprompt@gap>| !gapinput@l;|
  [ 1, 2, 3, 4, 5, 6, 3, 6, 9, 9 ]
\end{Verbatim}
 

 Note in the following example that the modification of the \emph{variable} that holds the list has no influence on the loop. 

 
\begin{Verbatim}[commandchars=!@|,fontsize=\small,frame=single,label=Example]
  !gapprompt@gap>| !gapinput@l := [ 1, 2, 3, 4, 5, 6 ];;|
  !gapprompt@gap>| !gapinput@for i in l do|
  !gapprompt@>| !gapinput@   Print( i, " " );|
  !gapprompt@>| !gapinput@   l := [];|
  !gapprompt@>| !gapinput@od; Print( "\n" );|
  1 2 3 4 5 6
  !gapprompt@gap>| !gapinput@l;|
  [  ]
\end{Verbatim}
 

 \index{loop over iterator} \texttt{for \mbox{\texttt{\mdseries\slshape variable}} in \mbox{\texttt{\mdseries\slshape iterator}} do \mbox{\texttt{\mdseries\slshape statements}} od;} 

 It is also possible to have a \texttt{for}\texttt{\symbol{45}}loop run over an iterator (see{\nobreakspace}\ref{Iterators}). In this case the \texttt{for}\texttt{\symbol{45}}loop is equivalent to 

 
\begin{Verbatim}[commandchars=!@|,fontsize=\small,frame=single,label=]
  while not IsDoneIterator(iterator) do
    variable := NextIterator(iterator)
    statements
  od;
\end{Verbatim}
 

 \index{loop over object} \texttt{for \mbox{\texttt{\mdseries\slshape variable}} in \mbox{\texttt{\mdseries\slshape object}} do \mbox{\texttt{\mdseries\slshape statements}} od;} 

 Finally, if an object \mbox{\texttt{\mdseries\slshape object}} which is not a list or an iterator appears in a \texttt{for}\texttt{\symbol{45}}loop, then \textsf{GAP} will attempt to evaluate the function call \texttt{Iterator(\mbox{\texttt{\mdseries\slshape object}})}. If this is successful then the loop is taken to run over the iterator
returned. 

 
\begin{Verbatim}[commandchars=!@|,fontsize=\small,frame=single,label=Example]
  !gapprompt@gap>| !gapinput@g := Group((1,2,3,4,5),(1,2)(3,4)(5,6));|
  Group([ (1,2,3,4,5), (1,2)(3,4)(5,6) ])
  !gapprompt@gap>| !gapinput@count := 0;; sumord := 0;;|
  !gapprompt@gap>| !gapinput@for x in g do|
  !gapprompt@>| !gapinput@count := count + 1; sumord := sumord + Order(x); od;|
  !gapprompt@gap>| !gapinput@count;|
  120
  !gapprompt@gap>| !gapinput@sumord;|
  471
\end{Verbatim}
 

 The effect of 

 \texttt{for \mbox{\texttt{\mdseries\slshape variable}} in \mbox{\texttt{\mdseries\slshape domain}} do} 

 should thus normally be the same as 

 \texttt{for \mbox{\texttt{\mdseries\slshape variable}} in AsList(\mbox{\texttt{\mdseries\slshape domain}}) do} 

 but may use much less storage, as the iterator may be more compact than a list
of all the elements. 

 See \ref{Iterators} for details about iterators. 

 A \texttt{for} loop may be left prematurely using \texttt{break}, see \ref{Break}. This combines especially well with a loop over an iterator, as a way of
searching through a domain for an element with some useful property. }

  
\subsection{\textcolor{Chapter }{Break}}\label{Break}
\logpage{[ 4, 15, 7 ]}
\hyperdef{L}{X7B60C6127E183021}{}
{
  \index{loops!leaving} \texttt{break;} \index{break statement@\texttt{break} statement} 

 The statement \texttt{break;} causes an immediate exit from the innermost loop enclosing it. 

 
\begin{Verbatim}[commandchars=!@|,fontsize=\small,frame=single,label=Example]
  !gapprompt@gap>| !gapinput@g := Group((1,2,3,4,5),(1,2)(3,4)(5,6));|
  Group([ (1,2,3,4,5), (1,2)(3,4)(5,6) ])
  !gapprompt@gap>| !gapinput@for x in g do|
  !gapprompt@>| !gapinput@if Order(x) = 3 then|
  !gapprompt@>| !gapinput@break;|
  !gapprompt@>| !gapinput@fi; od;|
  !gapprompt@gap>| !gapinput@x;|
  (1,5,2)(3,4,6)
\end{Verbatim}
 

 It is an error to use this statement other than inside a loop. 

 
\begin{Verbatim}[commandchars=!@|,fontsize=\small,frame=single,label=Example]
  !gapprompt@gap>| !gapinput@break;|
  Syntax error: 'break' statement not enclosed in a loop
\end{Verbatim}
 }

  
\subsection{\textcolor{Chapter }{Continue}}\label{Continue}
\logpage{[ 4, 15, 8 ]}
\hyperdef{L}{X7CCBA2247AA366BD}{}
{
  \index{loops!restarting} \texttt{continue;} \index{continue statement@\texttt{continue} statement} 

 The statement \texttt{continue;} causes the rest of the current iteration of the innermost loop enclosing it to
be skipped. 

 
\begin{Verbatim}[commandchars=!@|,fontsize=\small,frame=single,label=Example]
  !gapprompt@gap>| !gapinput@g := Group((1,2,3),(1,2));|
  Group([ (1,2,3), (1,2) ])
  !gapprompt@gap>| !gapinput@for x in g do|
  !gapprompt@>| !gapinput@if Order(x) = 3 then|
  !gapprompt@>| !gapinput@continue;|
  !gapprompt@>| !gapinput@fi; Print(x,"\n"); od;|
  ()
  (2,3)
  (1,3)
  (1,2)
\end{Verbatim}
 

 It is an error to use this statement other than inside a loop. 

 
\begin{Verbatim}[commandchars=!@|,fontsize=\small,frame=single,label=Example]
  !gapprompt@gap>| !gapinput@continue;|
  Syntax error: 'continue' statement not enclosed in a loop
\end{Verbatim}
 }

  
\subsection{\textcolor{Chapter }{Return (With or without Value)}}\label{Return}
\logpage{[ 4, 15, 9 ]}
\hyperdef{L}{X812C6ABC7A182E9E}{}
{
  \index{return@\texttt{return}!no value} \texttt{return;} 

 In this form \texttt{return} terminates the call of the innermost function that is currently executing, and
control returns to the calling function. An error is signalled if no function
is currently executing. No value is returned by the function. 

 \texttt{return \mbox{\texttt{\mdseries\slshape expr}};} \index{return@\texttt{return}!with value} 

 In this form \texttt{return} terminates the call of the innermost function that is currently executing, and
returns the value of the expression \mbox{\texttt{\mdseries\slshape expr}}. Control returns to the calling function. An error is signalled if no
function is currently executing. 

 Both statements can also be used in break loops (see \ref{Break Loops}). \texttt{return;} has the effect that the computation continues where it was interrupted by an
error or the user hitting \textsc{Ctrl\texttt{\symbol{45}}C}. \texttt{return \mbox{\texttt{\mdseries\slshape expr}};} can be used to continue execution after an error. What happens with the value \mbox{\texttt{\mdseries\slshape expr}} depends on the particular error. 

 For examples of \texttt{return} statements, see the functions \texttt{fib} and \texttt{Stack} in Section{\nobreakspace}\ref{Function}. }

 }

 
\section{\textcolor{Chapter }{Syntax Trees}}\label{Syntax Trees}
\logpage{[ 4, 16, 0 ]}
\hyperdef{L}{X8732D9257FFCEA1B}{}
{
  This section describes the tools available to handle \textsf{GAP} syntax trees. 

\subsection{\textcolor{Chapter }{SyntaxTree}}
\logpage{[ 4, 16, 1 ]}\nobreak
\hyperdef{L}{X81558D66810BEA67}{}
{\noindent\textcolor{FuncColor}{$\triangleright$\enspace\texttt{SyntaxTree({\mdseries\slshape f})\index{SyntaxTree@\texttt{SyntaxTree}}
\label{SyntaxTree}
}\hfill{\scriptsize (function)}}\\


 Takes a GAP function \mbox{\texttt{\mdseries\slshape f}} and returns its syntax tree. }

 }

   }

             
\chapter{\textcolor{Chapter }{Functions}}\label{Functions}
\logpage{[ 5, 0, 0 ]}
\hyperdef{L}{X86FA580F8055B274}{}
{
  \index{functions!as in programming language} The section{\nobreakspace}\ref{Function} describes how to define a function. In this chapter we describe functions that
give information about functions, and various utility functions used either
when defining functions or calling functions.  
\section{\textcolor{Chapter }{Information about a function}}\label{Information about a function}
\logpage{[ 5, 1, 0 ]}
\hyperdef{L}{X870553EF8605792F}{}
{
  

\subsection{\textcolor{Chapter }{NameFunction}}
\logpage{[ 5, 1, 1 ]}\nobreak
\hyperdef{L}{X79C3BDC4781FA0FD}{}
{\noindent\textcolor{FuncColor}{$\triangleright$\enspace\texttt{NameFunction({\mdseries\slshape func})\index{NameFunction@\texttt{NameFunction}}
\label{NameFunction}
}\hfill{\scriptsize (attribute)}}\\


 returns the name of a function. For operations, this is the name used in their
declaration. For functions, this is the variable name they were first assigned
to. (For some internal functions, this might be a name \emph{different} from the name that is documented.) If no such name exists, the string \texttt{"unknown"} is returned. 

 
\begin{Verbatim}[commandchars=!@|,fontsize=\small,frame=single,label=Example]
  !gapprompt@gap>| !gapinput@NameFunction(SylowSubgroup);|
  "SylowSubgroup"
  !gapprompt@gap>| !gapinput@Blubberflutsch:=x->x;;|
  !gapprompt@gap>| !gapinput@HasNameFunction(Blubberflutsch);|
  true
  !gapprompt@gap>| !gapinput@NameFunction(Blubberflutsch);|
  "Blubberflutsch"
  !gapprompt@gap>| !gapinput@a:=Blubberflutsch;;|
  !gapprompt@gap>| !gapinput@NameFunction(a);|
  "Blubberflutsch"
  !gapprompt@gap>| !gapinput@SetNameFunction(a, "f");|
  !gapprompt@gap>| !gapinput@NameFunction(a);|
  "f"
  !gapprompt@gap>| !gapinput@HasNameFunction(x->x);|
  false
  !gapprompt@gap>| !gapinput@NameFunction(x->x);|
  "unknown"
\end{Verbatim}
 }

 

\subsection{\textcolor{Chapter }{NumberArgumentsFunction}}
\logpage{[ 5, 1, 2 ]}\nobreak
\hyperdef{L}{X877F03F77FD74C98}{}
{\noindent\textcolor{FuncColor}{$\triangleright$\enspace\texttt{NumberArgumentsFunction({\mdseries\slshape func})\index{NumberArgumentsFunction@\texttt{NumberArgumentsFunction}}
\label{NumberArgumentsFunction}
}\hfill{\scriptsize (operation)}}\\


 returns the number of arguments the function \mbox{\texttt{\mdseries\slshape func}} accepts. \texttt{\symbol{45}}1 is returned for all operations. For functions
that use \texttt{...} or \texttt{arg} to take a variable number of arguments, the number returned is
\texttt{\symbol{45}}1 times the total number of parameters. For attributes, 1
is returned. 

 
\begin{Verbatim}[commandchars=!@|,fontsize=\small,frame=single,label=Example]
  !gapprompt@gap>| !gapinput@NumberArgumentsFunction(function(a,b,c,d,e,f,g,h,i,j,k)return 1;end);|
  11
  !gapprompt@gap>| !gapinput@NumberArgumentsFunction(Size);|
  1
  !gapprompt@gap>| !gapinput@NumberArgumentsFunction(IsCollsCollsElms);|
  3
  !gapprompt@gap>| !gapinput@NumberArgumentsFunction(Sum);|
  -1
  !gapprompt@gap>| !gapinput@NumberArgumentsFunction(function(a, x...) return 1; end);|
  -2
\end{Verbatim}
 }

 

\subsection{\textcolor{Chapter }{NamesLocalVariablesFunction}}
\logpage{[ 5, 1, 3 ]}\nobreak
\hyperdef{L}{X818BAB817A4FB346}{}
{\noindent\textcolor{FuncColor}{$\triangleright$\enspace\texttt{NamesLocalVariablesFunction({\mdseries\slshape func})\index{NamesLocalVariablesFunction@\texttt{NamesLocalVariablesFunction}}
\label{NamesLocalVariablesFunction}
}\hfill{\scriptsize (operation)}}\\


 returns a mutable list of strings; the first entries are the names of the
arguments of the function \mbox{\texttt{\mdseries\slshape func}}, in the same order as they were entered in the definition of \mbox{\texttt{\mdseries\slshape func}}, and the remaining ones are the local variables as given in the \texttt{local} statement in \mbox{\texttt{\mdseries\slshape func}}. (The number of arguments can be computed with \texttt{NumberArgumentsFunction} (\ref{NumberArgumentsFunction}).) 

 
\begin{Verbatim}[commandchars=!@|,fontsize=\small,frame=single,label=Example]
  !gapprompt@gap>| !gapinput@NamesLocalVariablesFunction(function( a, b ) local c; return 1; end);|
  [ "a", "b", "c" ]
  !gapprompt@gap>| !gapinput@NamesLocalVariablesFunction(function( arg ) local a; return 1; end);|
  [ "arg", "a" ]
  !gapprompt@gap>| !gapinput@NamesLocalVariablesFunction( Size );|
  fail
\end{Verbatim}
 }

 

\subsection{\textcolor{Chapter }{FilenameFunc}}
\logpage{[ 5, 1, 4 ]}\nobreak
\hyperdef{L}{X80E108C57F90FAA3}{}
{\noindent\textcolor{FuncColor}{$\triangleright$\enspace\texttt{FilenameFunc({\mdseries\slshape func})\index{FilenameFunc@\texttt{FilenameFunc}}
\label{FilenameFunc}
}\hfill{\scriptsize (function)}}\\


 For a function \mbox{\texttt{\mdseries\slshape func}}, \texttt{FilenameFunc} returns either \texttt{fail} or the absolute path of the file from which \mbox{\texttt{\mdseries\slshape func}} has been read. The return value \texttt{fail} occurs if \mbox{\texttt{\mdseries\slshape func}} is a compiled function or an operation. For functions that have been entered
interactively, the string \texttt{"*stdin*"} is returned, see Section \ref{Special Filenames}. 

 
\begin{Verbatim}[commandchars=!@|,fontsize=\small,frame=single,label=Example]
  !gapprompt@gap>| !gapinput@FilenameFunc( LEN_LIST );  # a kernel function|
  fail
  !gapprompt@gap>| !gapinput@FilenameFunc( Size );      # an operation|
  fail
  !gapprompt@gap>| !gapinput@FilenameFunc( x -> x^2 );  # an interactively entered function|
  "*stdin*"
  !gapprompt@gap>| !gapinput@meth:= ApplicableMethod( Size, [ Group( () ) ] );;|
  !gapprompt@gap>| !gapinput@FilenameFunc( meth );|
  "... some path .../grpperm.gi"
\end{Verbatim}
 }

 

\subsection{\textcolor{Chapter }{StartlineFunc}}
\logpage{[ 5, 1, 5 ]}\nobreak
\hyperdef{L}{X7FF7643781D2C194}{}
{\noindent\textcolor{FuncColor}{$\triangleright$\enspace\texttt{StartlineFunc({\mdseries\slshape func})\index{StartlineFunc@\texttt{StartlineFunc}}
\label{StartlineFunc}
}\hfill{\scriptsize (function)}}\\
\noindent\textcolor{FuncColor}{$\triangleright$\enspace\texttt{EndlineFunc({\mdseries\slshape func})\index{EndlineFunc@\texttt{EndlineFunc}}
\label{EndlineFunc}
}\hfill{\scriptsize (function)}}\\


 Let \mbox{\texttt{\mdseries\slshape func}} be a function. If \texttt{FilenameFunc} (\ref{FilenameFunc}) returns \texttt{fail} for \mbox{\texttt{\mdseries\slshape func}} then also \texttt{StartlineFunc} returns \texttt{fail}. If \texttt{FilenameFunc} (\ref{FilenameFunc}) returns a filename for \mbox{\texttt{\mdseries\slshape func}} then \texttt{StartlineFunc} returns the line number in this file where the definition of \mbox{\texttt{\mdseries\slshape func}} starts. 

 \texttt{EndlineFunc} behaves similarly and returns the line number in this file where the
definition of \mbox{\texttt{\mdseries\slshape func}} ends. 

 
\begin{Verbatim}[commandchars=!@|,fontsize=\small,frame=single,label=Example]
  !gapprompt@gap>| !gapinput@meth:= ApplicableMethod( Size, [ Group( () ) ] );;|
  !gapprompt@gap>| !gapinput@FilenameFunc( meth );|
  "... some path ... /lib/grpperm.gi"
  !gapprompt@gap>| !gapinput@StartlineFunc( meth );|
  487
  !gapprompt@gap>| !gapinput@EndlineFunc( meth );|
  487
\end{Verbatim}
 }

 

\subsection{\textcolor{Chapter }{LocationFunc}}
\logpage{[ 5, 1, 6 ]}\nobreak
\hyperdef{L}{X844F95767C74834F}{}
{\noindent\textcolor{FuncColor}{$\triangleright$\enspace\texttt{LocationFunc({\mdseries\slshape func})\index{LocationFunc@\texttt{LocationFunc}}
\label{LocationFunc}
}\hfill{\scriptsize (function)}}\\


 Let \mbox{\texttt{\mdseries\slshape func}} be a function. Returns a string describing the location of \mbox{\texttt{\mdseries\slshape func}}, or \texttt{fail} if the information cannot be found. This uses the information provided by \texttt{FilenameFunc} (\ref{FilenameFunc}) and \texttt{StartlineFunc} (\ref{StartlineFunc}) 

 
\begin{Verbatim}[commandchars=!@|,fontsize=\small,frame=single,label=Example]
  !gapprompt@gap>| !gapinput@LocationFunc( Intersection );|
  "... some path ... gap/lib/coll.gi:2467"
  # String is an attribute, so no information is stored
  !gapprompt@gap>| !gapinput@LocationFunc( String );|
  fail
\end{Verbatim}
 }

 

\subsection{\textcolor{Chapter }{PageSource}}
\logpage{[ 5, 1, 7 ]}\nobreak
\hyperdef{L}{X845A929B83D46E01}{}
{\noindent\textcolor{FuncColor}{$\triangleright$\enspace\texttt{PageSource({\mdseries\slshape func[, nr]})\index{PageSource@\texttt{PageSource}}
\label{PageSource}
}\hfill{\scriptsize (function)}}\\


 This shows the file containing the source code of the function or method \mbox{\texttt{\mdseries\slshape func}} in a pager (see \texttt{Pager} (\ref{Pager})). The display starts at a line shortly before the code of \mbox{\texttt{\mdseries\slshape func}}.

 For operations \mbox{\texttt{\mdseries\slshape func}} the function shows the source code of the declaration of \mbox{\texttt{\mdseries\slshape func}}. Operations can have several declarations, use the optional second argument
to specify which one should be shown (in the order the declarations were
read); the default is to show the first.

 For kernel functions the function tries to show the C source code.

 If GAP cannot find a file containing the source code this will be indicated. 

 Usage examples:\\
 \texttt{met := ApplicableMethod(\texttt{\symbol{92}}\texttt{\symbol{94}},
[(1,2),2743527]); PageSource(met);}\\
 \texttt{PageSource(Combinations);}\\
 \texttt{PageSource(SORT{\textunderscore}LIST); }\\
 \texttt{PageSource(Size, 2);}\\
 \texttt{ct := CharacterTable(Group((1,2,3))); }\\
 \texttt{met := ApplicableMethod(Size,[ct]); PageSource(met); } 

 }

 }

  
\section{\textcolor{Chapter }{Calling a function with a list argument that is interpreted as several
arguments}}\label{Calling a function with a list argument that is interpreted as
several arguments}
\logpage{[ 5, 2, 0 ]}
\hyperdef{L}{X851B58408520700D}{}
{
  

\subsection{\textcolor{Chapter }{CallFuncList}}
\logpage{[ 5, 2, 1 ]}\nobreak
\hyperdef{L}{X7CF4DDB97D65AE52}{}
{\noindent\textcolor{FuncColor}{$\triangleright$\enspace\texttt{CallFuncList({\mdseries\slshape func, args})\index{CallFuncList@\texttt{CallFuncList}}
\label{CallFuncList}
}\hfill{\scriptsize (operation)}}\\
\noindent\textcolor{FuncColor}{$\triangleright$\enspace\texttt{CallFuncListWrap({\mdseries\slshape func, args})\index{CallFuncListWrap@\texttt{CallFuncListWrap}}
\label{CallFuncListWrap}
}\hfill{\scriptsize (operation)}}\\


 returns the result, when calling function \mbox{\texttt{\mdseries\slshape func}} with the arguments given in the list \mbox{\texttt{\mdseries\slshape args}}, i.e.{\nobreakspace}\mbox{\texttt{\mdseries\slshape args}} is ``unwrapped'' so that \mbox{\texttt{\mdseries\slshape args}} appears as several arguments to \mbox{\texttt{\mdseries\slshape func}}. 

 
\begin{Verbatim}[commandchars=!@|,fontsize=\small,frame=single,label=Example]
  !gapprompt@gap>| !gapinput@CallFuncList(\+, [6, 7]);|
  13
  !gapprompt@gap>| !gapinput@#is equivalent to:|
  !gapprompt@gap>| !gapinput@\+(6, 7);|
  13
\end{Verbatim}
 

 A more useful application of \texttt{CallFuncList} is for a function \texttt{g} that is called in the body of a function \texttt{f} with (a sublist of) the arguments of \texttt{f}, where \texttt{f} has been defined with a single formal argument \texttt{arg} (see{\nobreakspace}\ref{Function}), as in the following code fragment. 

 
\begin{Verbatim}[commandchars=!@|,fontsize=\small,frame=single,label=Example]
  f := function ( arg )
         CallFuncList(g, arg);
         ...
       end;
\end{Verbatim}
 

 In the body of \texttt{f} the several arguments passed to \texttt{f} become a list \texttt{arg}. If \texttt{g} were called instead via \texttt{g( arg )} then \texttt{g} would see a single list argument, so that \texttt{g} would, in general, have to ``unwrap'' the passed list. The following (not particularly useful) example demonstrates
both described possibilities for the call to \texttt{g}. 

 
\begin{Verbatim}[commandchars=!@|,fontsize=\small,frame=single,label=Example]
  !gapprompt@gap>| !gapinput@PrintNumberFromDigits := function ( arg )|
  !gapprompt@>| !gapinput@    CallFuncList( Print, arg );|
  !gapprompt@>| !gapinput@    Print( "\n" );|
  !gapprompt@>| !gapinput@   end;|
  function( arg... ) ... end
  !gapprompt@gap>| !gapinput@PrintNumberFromDigits( 1, 9, 7, 3, 2 );|
  19732
  !gapprompt@gap>| !gapinput@PrintDigits := function ( arg )|
  !gapprompt@>| !gapinput@    Print( arg );|
  !gapprompt@>| !gapinput@    Print( "\n" );|
  !gapprompt@>| !gapinput@   end;|
  function( arg... ) ... end
  !gapprompt@gap>| !gapinput@PrintDigits( 1, 9, 7, 3, 2 );|
  [ 1, 9, 7, 3, 2 ]
\end{Verbatim}
 \texttt{CallFuncListWrap} differs only in that the result is a list. This returned list is empty if the
called function returned no value, else it contains the returned value as its
single member. This allows wrapping functions which may, or may not return a
value. 
\begin{Verbatim}[commandchars=!@|,fontsize=\small,frame=single,label=Example]
  !gapprompt@gap>| !gapinput@CallFuncListWrap( x -> x, [1] );|
  [ 1 ]
  !gapprompt@gap>| !gapinput@CallFuncListWrap( function(x) end, [1] );|
  [ ]
\end{Verbatim}
 }

 }

 
\section{\textcolor{Chapter }{Wrapping a function, so the values produced are cached}}\label{Memoizing a function}
\logpage{[ 5, 3, 0 ]}
\hyperdef{L}{X83066E5A80B5FB71}{}
{
  

\subsection{\textcolor{Chapter }{MemoizePosIntFunction}}
\logpage{[ 5, 3, 1 ]}\nobreak
\hyperdef{L}{X817ED3B280A64601}{}
{\noindent\textcolor{FuncColor}{$\triangleright$\enspace\texttt{MemoizePosIntFunction({\mdseries\slshape function[, options]})\index{MemoizePosIntFunction@\texttt{MemoizePosIntFunction}}
\label{MemoizePosIntFunction}
}\hfill{\scriptsize (function)}}\\


 \texttt{MemoizePosIntFunction} returns a function which behaves the same as \mbox{\texttt{\mdseries\slshape function}}, except it caches the results for any inputs that are positive integers. Thus
if the new function is called multiple times with the same input, then any
call after the first will return the cached value, instead of recomputing it.
By default, the cache can be flushed by calling \texttt{FlushCaches} (\ref{FlushCaches}). 

 The returned function will by default only accept positive integers. 

 This function does not promise to never call \mbox{\texttt{\mdseries\slshape function}} more than once for any input \texttt{\symbol{45}}\texttt{\symbol{45}} values
may be removed if the cache gets too large, or if \texttt{FlushCaches} (\ref{FlushCaches}) is called, or if multiple threads try to calculate the same value
simultaneously. 

 The optional second argument is a record which provides a number of
configuration options. The following options are supported. 
\begin{description}
\item[{\texttt{defaults} (default an empty list)}]  Used to initialise the cache, both initially and after each flush. If \texttt{defaults[i]} is bound, then this is used as default value for the input \texttt{i}. 
\item[{\texttt{flush} (default \texttt{true})}]  If this is \texttt{true}, the cache is emptied whenever \texttt{FlushCaches} (\ref{FlushCaches}) is called; if false, then the cache cannot be flushed. 
\item[{\texttt{errorHandler} (defaults to \texttt{Error} (\ref{Error}))}]  A function to be called when an input which is not a positive integer is
passed to the cache. The function can either raise an error, or else return a
value which is then returned by the cache. Note that such a value does not get
cached itself. 
\end{description}
 

 
\begin{Verbatim}[commandchars=!@|,fontsize=\small,frame=single,label=Example]
  !gapprompt@gap>| !gapinput@f := MemoizePosIntFunction(|
  !gapprompt@>| !gapinput@          function(i) Print("Check: ",i,"\n"); return i*i; end,|
  !gapprompt@>| !gapinput@          rec(defaults := [,,50], errorHandler := x -> "Bad") );;|
  !gapprompt@gap>| !gapinput@f(2);|
  Check: 2
  4
  !gapprompt@gap>| !gapinput@f(2);|
  4
  !gapprompt@gap>| !gapinput@f(3);|
  50
  !gapprompt@gap>| !gapinput@f(-3);|
  "Bad"
  !gapprompt@gap>| !gapinput@FlushCaches();|
  !gapprompt@gap>| !gapinput@f(2);|
  Check: 2
  4
  !gapprompt@gap>| !gapinput@f(3);|
  50
\end{Verbatim}
 }

 }

  
\section{\textcolor{Chapter }{Functions that do nothing}}\label{Functions that do nothing}
\logpage{[ 5, 4, 0 ]}
\hyperdef{L}{X7EB0A85F7D128BE0}{}
{
  The following functions return fixed results (or just their own argument).
They can be useful in places when the syntax requires a function, but actually
no functionality is required. So \texttt{ReturnTrue} (\ref{ReturnTrue}) is often used as family predicate in \texttt{InstallMethod} (\ref{InstallMethod}). 

\subsection{\textcolor{Chapter }{ReturnTrue}}
\logpage{[ 5, 4, 1 ]}\nobreak
\hyperdef{L}{X7DB422A2876CCC4D}{}
{\noindent\textcolor{FuncColor}{$\triangleright$\enspace\texttt{ReturnTrue({\mdseries\slshape ...})\index{ReturnTrue@\texttt{ReturnTrue}}
\label{ReturnTrue}
}\hfill{\scriptsize (function)}}\\


 This function takes any number of arguments, and always returns \texttt{true}. 

 
\begin{Verbatim}[commandchars=!@|,fontsize=\small,frame=single,label=Example]
  !gapprompt@gap>| !gapinput@f:=ReturnTrue;|
  function( arg... ) ... end
  !gapprompt@gap>| !gapinput@f();|
  true
  !gapprompt@gap>| !gapinput@f(42);|
  true
\end{Verbatim}
 }

 

\subsection{\textcolor{Chapter }{ReturnFalse}}
\logpage{[ 5, 4, 2 ]}\nobreak
\hyperdef{L}{X7C131FB17D7518FC}{}
{\noindent\textcolor{FuncColor}{$\triangleright$\enspace\texttt{ReturnFalse({\mdseries\slshape ...})\index{ReturnFalse@\texttt{ReturnFalse}}
\label{ReturnFalse}
}\hfill{\scriptsize (function)}}\\


 This function takes any number of arguments, and always returns \texttt{false}. 

 
\begin{Verbatim}[commandchars=!@|,fontsize=\small,frame=single,label=Example]
  !gapprompt@gap>| !gapinput@f:=ReturnFalse;|
  function( arg... ) ... end
  !gapprompt@gap>| !gapinput@f();|
  false
  !gapprompt@gap>| !gapinput@f("any_string");|
  false
\end{Verbatim}
 }

 

\subsection{\textcolor{Chapter }{ReturnFail}}
\logpage{[ 5, 4, 3 ]}\nobreak
\hyperdef{L}{X7A0994DE7C258E55}{}
{\noindent\textcolor{FuncColor}{$\triangleright$\enspace\texttt{ReturnFail({\mdseries\slshape ...})\index{ReturnFail@\texttt{ReturnFail}}
\label{ReturnFail}
}\hfill{\scriptsize (function)}}\\


 This function takes any number of arguments, and always returns \texttt{fail}. 

 
\begin{Verbatim}[commandchars=!@|,fontsize=\small,frame=single,label=Example]
  !gapprompt@gap>| !gapinput@oops:=ReturnFail;|
  function( arg... ) ... end
  !gapprompt@gap>| !gapinput@oops();|
  fail
  !gapprompt@gap>| !gapinput@oops(-42);|
  fail
\end{Verbatim}
 }

 

\subsection{\textcolor{Chapter }{ReturnNothing}}
\logpage{[ 5, 4, 4 ]}\nobreak
\hyperdef{L}{X818EA8C47B46A634}{}
{\noindent\textcolor{FuncColor}{$\triangleright$\enspace\texttt{ReturnNothing({\mdseries\slshape ...})\index{ReturnNothing@\texttt{ReturnNothing}}
\label{ReturnNothing}
}\hfill{\scriptsize (function)}}\\


 This function takes any number of arguments, and always returns nothing. 

 
\begin{Verbatim}[commandchars=!@|,fontsize=\small,frame=single,label=Example]
  !gapprompt@gap>| !gapinput@n:=ReturnNothing;|
  function( object... ) ... end
  !gapprompt@gap>| !gapinput@n();|
  !gapprompt@gap>| !gapinput@n(-42);|
\end{Verbatim}
 }

 

\subsection{\textcolor{Chapter }{ReturnFirst}}
\logpage{[ 5, 4, 5 ]}\nobreak
\hyperdef{L}{X8788D7D780FCE169}{}
{\noindent\textcolor{FuncColor}{$\triangleright$\enspace\texttt{ReturnFirst({\mdseries\slshape ...})\index{ReturnFirst@\texttt{ReturnFirst}}
\label{ReturnFirst}
}\hfill{\scriptsize (function)}}\\


 This function takes one or more arguments, and always returns the first
argument. \texttt{IdFunc} (\ref{IdFunc}) behaves similarly, but only accepts a single argument. 

 
\begin{Verbatim}[commandchars=!@|,fontsize=\small,frame=single,label=Example]
  !gapprompt@gap>| !gapinput@f:=ReturnFirst;|
  function( first, rest... ) ... end
  !gapprompt@gap>| !gapinput@f(1);|
  1
  !gapprompt@gap>| !gapinput@f(2,3,4);|
  2
  !gapprompt@gap>| !gapinput@f();|
  Error, Function: number of arguments must be at least 1 (not 0)
\end{Verbatim}
 }

 

\subsection{\textcolor{Chapter }{IdFunc}}
\logpage{[ 5, 4, 6 ]}\nobreak
\hyperdef{L}{X810325697BDEF899}{}
{\noindent\textcolor{FuncColor}{$\triangleright$\enspace\texttt{IdFunc({\mdseries\slshape obj})\index{IdFunc@\texttt{IdFunc}}
\label{IdFunc}
}\hfill{\scriptsize (function)}}\\


 returns \mbox{\texttt{\mdseries\slshape obj}}. \texttt{ReturnFirst} (\ref{ReturnFirst}) is similar, but accepts one or more arguments, returning only the first. 

 
\begin{Verbatim}[commandchars=!@|,fontsize=\small,frame=single,label=Example]
  !gapprompt@gap>| !gapinput@id:=IdFunc;|
  function( object ) ... end
  !gapprompt@gap>| !gapinput@id(42);|
  42
  !gapprompt@gap>| !gapinput@f:=id(SymmetricGroup(3));|
  Sym( [ 1 .. 3 ] )
  !gapprompt@gap>| !gapinput@s:=One(AutomorphismGroup(SymmetricGroup(3)));|
  IdentityMapping( Sym( [ 1 .. 3 ] ) )
  !gapprompt@gap>| !gapinput@f=s;|
  false
\end{Verbatim}
 }

 }

  
\section{\textcolor{Chapter }{Function Types}}\label{Function Types}
\logpage{[ 5, 5, 0 ]}
\hyperdef{L}{X80FE39D27CE3DE1B}{}
{
  Functions are \textsf{GAP} objects and thus have categories and a family. 

\subsection{\textcolor{Chapter }{IsFunction}}
\logpage{[ 5, 5, 1 ]}\nobreak
\hyperdef{L}{X85E40340806C2B8C}{}
{\noindent\textcolor{FuncColor}{$\triangleright$\enspace\texttt{IsFunction({\mdseries\slshape obj})\index{IsFunction@\texttt{IsFunction}}
\label{IsFunction}
}\hfill{\scriptsize (Category)}}\\


 is the category of functions. 

 
\begin{Verbatim}[commandchars=!@|,fontsize=\small,frame=single,label=Example]
  !gapprompt@gap>| !gapinput@IsFunction(x->x^2);|
  true
  !gapprompt@gap>| !gapinput@IsFunction(Factorial);|
  true
  !gapprompt@gap>| !gapinput@f:=One(AutomorphismGroup(SymmetricGroup(3)));|
  IdentityMapping( Sym( [ 1 .. 3 ] ) )
  !gapprompt@gap>| !gapinput@IsFunction(f);|
  false
\end{Verbatim}
 }

 

\subsection{\textcolor{Chapter }{FunctionsFamily}}
\logpage{[ 5, 5, 2 ]}\nobreak
\hyperdef{L}{X87838FE885A9AAF9}{}
{\noindent\textcolor{FuncColor}{$\triangleright$\enspace\texttt{FunctionsFamily\index{FunctionsFamily@\texttt{FunctionsFamily}}
\label{FunctionsFamily}
}\hfill{\scriptsize (family)}}\\


 is the family of all functions. }

 }

  
\section{\textcolor{Chapter }{Naming Conventions}}\label{Naming Conventions}
\logpage{[ 5, 6, 0 ]}
\hyperdef{L}{X81F732457F7BC851}{}
{
  The way functions are named in \textsf{GAP} might help to memorize or even guess names of library functions. 

 If a variable name consists of several words then the first letter of each
word is capitalized. 

 If the first part of the name of a function is a verb then the function may
modify its argument(s) but does not return anything, for example \texttt{Append} (\ref{Append}) appends the list given as second argument to the list given as first argument.
Otherwise the function returns an object without changing the arguments, for
example \texttt{Concatenation} (\ref{Concatenation:for several lists}) returns the concatenation of the lists given as arguments. 

 If the name of a function contains the word ``\texttt{Of}'' then the return value is thought of as information deduced from the arguments.
Usually such functions are attributes (see{\nobreakspace}\ref{Attributes}). Examples are \texttt{GeneratorsOfGroup} (\ref{GeneratorsOfGroup}), which returns a list of generators for the group entered as argument, or \texttt{DiagonalOfMat} (\ref{DiagonalOfMat}). 

 For the setter and tester functions of an attribute \texttt{Attr} the names \texttt{SetAttr} resp.{\nobreakspace}\texttt{HasAttr} are available (see{\nobreakspace}\ref{Attributes}). 

 If the name of a function contains the word ``\texttt{By}'' then the return value is thought of as built in a certain way from the parts
given as arguments. For example, creating a group as a factor group of a given
group by a normal subgroup can be done by taking the image of \texttt{NaturalHomomorphismByNormalSubgroup} (\ref{NaturalHomomorphismByNormalSubgroup}). Other examples of ``\texttt{By}'' functions are \texttt{GroupHomomorphismByImages} (\ref{GroupHomomorphismByImages}) and \texttt{LaurentPolynomialByCoefficients} (\ref{LaurentPolynomialByCoefficients}). 

 Often such functions construct an algebraic structure given by its generators
(for example, \texttt{RingByGenerators} (\ref{RingByGenerators})). In some cases, ``\texttt{By}'' may be replaced by ``\texttt{With}'' (like e.g. \texttt{GroupWithGenerators} (\ref{GroupWithGenerators})) or even both versions of the name may be used. The difference between \texttt{StructByGenerators} and \texttt{StructWithGenerators} is that the latter guarantees that the \texttt{GeneratorsOfStruct} value of the result is equal to the given set of generators (see \ref{Constructing Domains}). 

 If the name of a function has the form ``\texttt{AsSomething}'' then the return value is an object (usually a collection which has the same
family of elements), which may, for example: 
\begin{itemize}
\item  know more about its own structure (and so support more operations) than its
input (e.g. if the elements of the collection form a group, then this group
can be constructed using \texttt{AsGroup} (\ref{AsGroup})); 
\item  discard its additional structure (e.g. \texttt{AsList} (\ref{AsList}) applied to a group will return a list of its elements); 
\item  contain all elements of the original object without duplicates (like e.g. \texttt{AsSet} (\ref{AsSet}) does if its argument is a list of elements from the same family); 
\item  remain unchanged (like e.g. \texttt{AsSemigroup} (\ref{AsSemigroup}) does if its argument is a group). 
\end{itemize}
 If \texttt{Something} and the argument of \texttt{AsSomething} are domains, some further rules apply as explained in  \textbf{Tutorial: Changing the Structure}. 

 If the name of a function \texttt{fun1} ends with ``\texttt{NC}'' then there is another function \texttt{fun2} with the same name except that the \texttt{NC} is missing. \texttt{NC} stands for ``no check''. When \texttt{fun2} is called then it checks whether its arguments are valid, and if so then it
calls \texttt{fun1}. The functions \texttt{SubgroupNC} (\ref{SubgroupNC}) and \texttt{Subgroup} (\ref{Subgroup}) are a typical example. 

 The idea is that the possibly time consuming check of the arguments can be
omitted if one is sure that they are unnecessary. For example, if an algorithm
produces generators of the derived subgroup of a group then it is guaranteed
that they lie in the original group; \texttt{Subgroup} (\ref{Subgroup}) would check this, and \texttt{SubgroupNC} (\ref{SubgroupNC}) omits the check. 

 Needless to say, all these rules are not followed slavishly, for example there
is one operation \texttt{Zero} (\ref{Zero}) instead of two operations \texttt{ZeroOfElement} and \texttt{ZeroOfAdditiveGroup}. }

  
\section{\textcolor{Chapter }{Code annotations (pragmas)}}\label{Pragmas}
\logpage{[ 5, 7, 0 ]}
\hyperdef{L}{X7A1721CD79F08E71}{}
{
  \index{Pragmas} \index{Code annotations} \index{\#\%} GAP supports the use of code annotations (pragmas) in functions, i.e., adding
comments to functions that are stored in the function object itself, unlike
regular comments. Pragmas are single\texttt{\symbol{45}}line comments,
starting with \texttt{\#\%}: 
\begin{Verbatim}[commandchars=!@|,fontsize=\small,frame=single,label=Example]
  !gapprompt@gap>| !gapinput@function()|
  !gapprompt@>| !gapinput@      #% This is a pragma|
  !gapprompt@>| !gapinput@      #  This is not a pragma|
  !gapprompt@>| !gapinput@      return;|
  !gapprompt@>| !gapinput@end;;|
  !gapprompt@gap>| !gapinput@Display( last );|
  function (  )
      #% This is a pragma
      return;
  end
\end{Verbatim}
 

 Pragmas can be used to mark parts of functions that should later be
manipulated using \ref{Syntax Trees}. 

 Please note that heavy use of pragmas in functions slows down the execution of
your function in the same way as adding empty \texttt{;} statements to your code. 
\begin{Verbatim}[commandchars=!@|,fontsize=\small,frame=single,label=Example]
  !gapprompt@gap>| !gapinput@a := function( )|
  !gapprompt@>| !gapinput@      local i;|
  !gapprompt@>| !gapinput@      for i in [ 1 .. 1000000 ] do|
  !gapprompt@>| !gapinput@              i := i + 1;|
  !gapprompt@>| !gapinput@      od;|
  !gapprompt@>| !gapinput@end;|
  function(  ) ... end
  !gapprompt@gap>| !gapinput@a();|
  !gapprompt@gap>| !gapinput@time;|
  14
  !gapprompt@gap>| !gapinput@b := function( )|
  !gapprompt@>| !gapinput@      local i;|
  !gapprompt@>| !gapinput@      for i in [ 1 .. 1000000 ] do|
  !gapprompt@>| !gapinput@              i := i + 1;|
  !gapprompt@>| !gapinput@              #% pragma|
  !gapprompt@>| !gapinput@              #% pragma|
  !gapprompt@>| !gapinput@              #% pragma|
  !gapprompt@>| !gapinput@              #% pragma|
  !gapprompt@>| !gapinput@              #% pragma|
  !gapprompt@>| !gapinput@      od;|
  !gapprompt@>| !gapinput@end;|
  function(  ) ... end
  !gapprompt@gap>| !gapinput@b();|
  !gapprompt@gap>| !gapinput@time;|
  25
\end{Verbatim}
 }

 }

        
\chapter{\textcolor{Chapter }{Main Loop and Break Loop}}\label{Main Loop and Break Loop}
\logpage{[ 6, 0, 0 ]}
\hyperdef{L}{X7DB71A2A841CADA5}{}
{
  This chapter is a first of a series of chapters that describe the interactive
environment in which you use \textsf{GAP}.  
\section{\textcolor{Chapter }{Main Loop}}\label{Main Loop}
\logpage{[ 6, 1, 0 ]}
\hyperdef{L}{X81667F568237B232}{}
{
  \index{read eval print loop} \index{loop!read eval print} \index{prompt} \index{prompt!partial} \index{syntax errors} \index{errors!syntax} \index{output!suppressing} \index{last@\texttt{last}} \index{last2@\texttt{last2}} \index{last3@\texttt{last3}} \index{time@\texttt{time}} \index{memory_allocated@\texttt{memory{\textunderscore}allocated}} \index{previous result} The normal interaction with \textsf{GAP} happens in the so\texttt{\symbol{45}}called \emph{read\texttt{\symbol{45}}eval\texttt{\symbol{45}}print} loop. This means that you type an input, \textsf{GAP} first reads it, evaluates it, and then shows the result. Note that the term \emph{print} may be confusing since there is a \textsf{GAP} function called \texttt{Print} (\ref{Print}) (see{\nobreakspace}\ref{View and Print}) which is in fact \emph{not} used in the read\texttt{\symbol{45}}eval\texttt{\symbol{45}}print loop, but
traditions are hard to break. In the following, whenever we want to express
that \textsf{GAP} places some characters on the standard output, we will say that \textsf{GAP} \emph{shows} something. 

 The exact sequence in the
read\texttt{\symbol{45}}eval\texttt{\symbol{45}}print loop is as follows. 

 To signal that it is ready to accept your input, \textsf{GAP} shows the \emph{prompt} \texttt{gap{\textgreater}}. When you see this, you know that \textsf{GAP} is waiting for your input. 

 Note that every statement must be terminated by a semicolon. You must also
enter \textsc{Return} (i.e., strike the \textsc{Return} key) before \textsf{GAP} starts to read and evaluate your input. (The \textsc{Return} key may actually be marked with the word \textsc{Enter} and a returning arrow on your terminal.) Because \textsf{GAP} does not do anything until you enter \textsc{Return}, you can edit your input to fix typos and only when everything is correct
enter \textsc{Return} and have \textsf{GAP} take a look at it (see{\nobreakspace}\ref{Line Editing}). It is also possible to enter several statements as input on a single line.
Of course each statement must be terminated by a semicolon. 

 It is absolutely acceptable to enter a single statement on several lines. When
you have entered the beginning of a statement, but the statement is not yet
complete, and you enter \textsc{Return}, \textsf{GAP} will show the \emph{partial prompt} \texttt{{\textgreater}}. When you see this, you know that \textsf{GAP} is waiting for the rest of the statement. This happens also when you forget
the semicolon \texttt{;} that terminates every \textsf{GAP} statement. Note that when \textsc{Return} has been entered and the current statement is not yet complete, \textsf{GAP} will already evaluate those parts of the input that are complete, for example
function calls that appear as arguments in another function call which needs
several input lines. So it may happen that one has to wait some time for the
partial prompt.   

 When you enter \textsc{Return}, \textsf{GAP} first checks your input to see if it is syntactically correct (see
Chapter{\nobreakspace}\ref{The Programming Language} for the definition of syntactically correct). If it is not, \textsf{GAP} prints an error message of the following form 

 
\begin{Verbatim}[commandchars=!@|,fontsize=\small,frame=single,label=Example]
  !gapprompt@gap>| !gapinput@1 * ;|
  Syntax error: Expression expected
  1 * ;
      ^
\end{Verbatim}
 

 The first line tells you what is wrong about the input, in this case the \texttt{*} operator takes two expressions as operands, so obviously the right one is
missing. If the input came from a file (see{\nobreakspace}\texttt{Read} (\ref{Read})), this line will also contain the filename and the line number. The second
line is a copy of the input. And the third line contains a caret pointing to
the place in the previous line where \textsf{GAP} realized that something is wrong. This need not be the exact place where the
error is, but it is usually quite close. 

 Sometimes, you will also see a partial prompt after you have entered an input
that is syntactically incorrect. This is because \textsf{GAP} is so confused by your input, that it thinks that there is still something to
follow. In this case you should enter \texttt{;}\textsc{Return} repeatedly, ignoring further error messages, until you see the full prompt
again. When you see the full prompt, you know that \textsf{GAP} forgave you and is now ready to accept your next {\textendash}hopefully
correct{\textendash} input. 

 If your input is syntactically correct, \textsf{GAP} evaluates or executes it, i.e., performs the required computations (see
Chapter{\nobreakspace}\ref{The Programming Language} for the definition of the evaluation). 

 If you do not see a prompt, you know that \textsf{GAP} is still working on your last input. Of course, you can \emph{type ahead}, i.e., already start entering new input, but it will not be accepted by \textsf{GAP} until \textsf{GAP} has completed the ongoing computation. 

 When \textsf{GAP} is ready it will usually show the result of the computation, i.e., the value
computed. Note that not all statements produce a value, for example, if you
enter a \texttt{for} loop, nothing will be printed, because the \texttt{for} loop does not produce a value that could be shown. 

 Also sometimes you do not want to see the result. For example if you have
computed a value and now want to assign the result to a variable, you probably
do not want to see the value again. You can terminate statements by \emph{two semicolons} to suppress showing the result. 

 If you have entered several statements on a single line \textsf{GAP} will first read, evaluate, and show the first one, then read, evaluate, and
show the second one, and so on. This means that the second statement will not
even be checked for syntactical correctness until \textsf{GAP} has completed the first computation. 

 After the result has been shown \textsf{GAP} will display another prompt, and wait for your next input. And the whole
process starts all over again. Note that if you have entered several
statements on a single line, a new prompt will only be printed after \textsf{GAP} has read, evaluated, and shown the last statement. 

 In each statement that you enter, the result of the previous statement that
produced a value is available in the variable \texttt{last}. The next to previous result is available in \texttt{last2} and the result produced before that is available in \texttt{last3}. 

 
\begin{Verbatim}[commandchars=!@|,fontsize=\small,frame=single,label=Example]
  !gapprompt@gap>| !gapinput@1;2;3;|
  1
  2
  3
  !gapprompt@gap>| !gapinput@last3 + last2 * last;|
  7
\end{Verbatim}
 

 Also in each statement the time spent by the last statement, whether it
produced a value or not, is available in the variable \texttt{time} (\ref{time}). This is an integer that holds the number of milliseconds. Similarly the
amount of memory allocated during that statement (in bytes) is stored in the
variable \texttt{memory{\textunderscore}allocated} (\ref{memoryuScoreallocated}). The variables \texttt{last}, \texttt{last2}, \texttt{last3}, \texttt{time} (\ref{time}) and \texttt{memory{\textunderscore}allocated} (\ref{memoryuScoreallocated}) are all write\texttt{\symbol{45}}protected. }

  
\section{\textcolor{Chapter }{Special Rules for Input Lines}}\label{Special Rules for Input Lines}
\logpage{[ 6, 2, 0 ]}
\hyperdef{L}{X866092F281910B74}{}
{
  The input for some \textsf{GAP} objects may not fit on one line, in particular big integers, long strings or
long identifiers. In these cases you can still type or paste them in long
single lines.  For nicer display you can also specify the input on several lines. This is
achieved by ending a line by a backslash or by a backslash and a carriage
return character, then continue the input on the beginning of the next line.
When reading this \textsf{GAP} will ignore such continuation backslashes, carriage return characters and
newline characters. \textsf{GAP} also prints long strings and integers this way. 

 
\begin{Verbatim}[commandchars=!@|,fontsize=\small,frame=single,label=Example]
  !gapprompt@gap>| !gapinput@n := 1234\|
  !gapprompt@>| !gapinput@567890;|
  1234567890
  !gapprompt@gap>| !gapinput@"This is a very long string that does not fit on a line \|
  !gapprompt@>| !gapinput@and is therefore continued on the next line.";|
  "This is a very long string that does not fit on a line and is therefo\
  re continued on the next line."
  !gapprompt@gap>| !gapinput@bla\|
  !gapprompt@>| !gapinput@bla := 5;; blabla;|
  5
\end{Verbatim}
 

 There is a special rule about \textsf{GAP} prompts in input lines: In line editing mode (usual user input and \textsf{GAP} started without \texttt{\texttt{\symbol{45}}n}) in lines starting with whitespace following \texttt{gap{\textgreater} }, \texttt{{\textgreater} } or \texttt{brk{\textgreater} } this beginning part is removed. This rule is very convenient because it allows
to cut and paste input from other \textsf{GAP} sessions or manual examples easily into your current session. }

  
\section{\textcolor{Chapter }{View and Print}}\label{View and Print}
\logpage{[ 6, 3, 0 ]}
\hyperdef{L}{X8074A8387C9DB9A8}{}
{
  \textsf{GAP} has three different operations to display or print objects: \texttt{Display} (\ref{Display}), \texttt{ViewObj} (\ref{ViewObj}) and \texttt{PrintObj} (\ref{PrintObj}), and these three have different purposes as follows. The first, \texttt{Display} (\ref{Display}), should print the object to the standard output in a
human\texttt{\symbol{45}}readable relatively complete and verbose form. The
second, \texttt{ViewObj} (\ref{ViewObj}), should print the object to the standard output in a short and concise form,
it is used in the main read\texttt{\symbol{45}}eval\texttt{\symbol{45}}print
loop to display the resulting object of a computation. The third, \texttt{PrintObj} (\ref{PrintObj}), should print the object to the standard output in a complete form which is \textsf{GAP}\texttt{\symbol{45}}readable if at all possible, such that reading the output
into \textsf{GAP} produces an object which is equal to the original one. 

 All three operations have corresponding operations which do not print anything
to standard output but return the output as a string. These are \texttt{DisplayString} (\ref{DisplayString}), \texttt{ViewString} (\ref{ViewString}) and \texttt{PrintString} (\ref{PrintString}) (corresponding to \texttt{PrintObj} (\ref{PrintObj})). Additionally, there is \texttt{String} (\ref{String}) which is very similar to \texttt{PrintString} (\ref{PrintString}) but does not insert control characters for line breaks. 

 For implementation convenience it is allowed that some of these operations
have methods which delegate to some other of these operations. However, the
rules for this are that a method may only delegate to another operation which
appears further down in the following table: 

 \begin{center}
\begin{tabular}{|c|}\hline
\texttt{Display} (\ref{Display})\\
\hline
\texttt{ViewObj} (\ref{ViewObj})\\
\hline
\texttt{PrintObj} (\ref{PrintObj})\\
\hline
\texttt{DisplayString} (\ref{DisplayString})\\
\hline
\texttt{ViewString} (\ref{ViewString})\\
\hline
\texttt{PrintString} (\ref{PrintString})\\
\hline
\texttt{String} (\ref{String})\\
\hline
\end{tabular}\\[2mm]
\end{center}

 

 This is to avoid circular delegations. 

 Note in particular that none of the methods of the string producing operations
may delegate to the corresponding printing operations. Note also that the
above mentioned purposes of the different operations suggest that delegations
between different operations will be sub\texttt{\symbol{45}}optimal in most
scenarios. 
\subsection{\textcolor{Chapter }{Default delegations in the library}}\logpage{[ 6, 3, 1 ]}
\hyperdef{L}{X8082880F824292E9}{}
{
  The library contains the following low ranked default methods: 
\begin{itemize}
\item  A method for \texttt{DisplayString} (\ref{DisplayString}) which returns the constant value of the global variable \texttt{DEFAULTDISPLAYSTRING} (\ref{DEFAULTDISPLAYSTRING}). 
\item  A method for \texttt{ViewString} (\ref{ViewString}) which returns the constant value of the global variable \texttt{DEFAULTVIEWSTRING} (\ref{DEFAULTVIEWSTRING}). 
\item  A method for \texttt{Display} (\ref{Display}) which first calls \texttt{DisplayString} (\ref{DisplayString}) and prints the result, if it is a different object than \texttt{DEFAULTDISPLAYSTRING} (\ref{DEFAULTDISPLAYSTRING}). Otherwise the method delegates to \texttt{PrintObj} (\ref{PrintObj}). 
\item  A method for \texttt{ViewObj} (\ref{ViewObj}) which first calls \texttt{ViewString} (\ref{ViewString}) and prints the result, if it is a different object than \texttt{DEFAULTVIEWSTRING} (\ref{DEFAULTVIEWSTRING}). Otherwise the method delegates to \texttt{PrintObj} (\ref{PrintObj}). 
\item  A method for \texttt{PrintObj} (\ref{PrintObj}) which prints the result of \texttt{PrintString} (\ref{PrintString}). 
\item  A method for \texttt{PrintString} (\ref{PrintString}) which returns the result of \texttt{String} (\ref{String}) 
\end{itemize}
 }

 
\subsection{\textcolor{Chapter }{Recommendations for the implementation}}\logpage{[ 6, 3, 2 ]}
\hyperdef{L}{X87D445D37B31DADB}{}
{
  This subsection describes what methods for printing and viewing one should
implement for new \textsf{GAP} objects. 

 One should at the very least install a \texttt{String} (\ref{String}) method to allow printing. Using the standard delegations this enables a
limited form of viewing, displaying and printing. 

 If, for larger objects, nicer line breaks are needed, one should install a
separate \texttt{PrintString} (\ref{PrintString}) method which puts in positions for good line breaks using the control
characters \texttt{\texttt{\symbol{92}}{\textless}} (ASCII 1) and \texttt{\texttt{\symbol{92}}{\textgreater}} (ASCII 2). 

 If, for even larger objects, output performance and memory usage matters, one
should install a separate \texttt{PrintObj} (\ref{PrintObj}) method. 

 One should usually install a \texttt{ViewString} (\ref{ViewString}) method, unless the above \texttt{String} (\ref{String}) method is good enough for \texttt{ViewObj} (\ref{ViewObj}) purposes. Performance and memory should never matter here, so it is usually
unnecessary to install a separate \texttt{ViewObj} (\ref{ViewObj}) method. 

 If the type of object calls for it one should install a \texttt{DisplayString} (\ref{DisplayString}) method. This is the case if a human readable verbose form is required. 

 If the performance and memory usage for \texttt{Display} (\ref{Display}) matters, one should install a separate \texttt{Display} (\ref{Display}) method. 

 Note that if only a \texttt{String} (\ref{String}) method is installed, then \texttt{ViewObj} (\ref{ViewObj}) works and \texttt{ViewString} (\ref{ViewString}) returns \texttt{DEFAULTVIEWSTRING} (\ref{DEFAULTVIEWSTRING}). Likewise, \texttt{Display} (\ref{Display}) works and \texttt{DisplayString} (\ref{DisplayString}) returns \texttt{DEFAULTDISPLAYSTRING} (\ref{DEFAULTDISPLAYSTRING}). If you want to avoid this then install methods for these operations as well. }

 

\subsection{\textcolor{Chapter }{View}}
\logpage{[ 6, 3, 3 ]}\nobreak
\hyperdef{L}{X851902C583B84CDC}{}
{\noindent\textcolor{FuncColor}{$\triangleright$\enspace\texttt{View({\mdseries\slshape obj1, obj2...})\index{View@\texttt{View}}
\label{View}
}\hfill{\scriptsize (function)}}\\


 \texttt{View} shows the objects \mbox{\texttt{\mdseries\slshape obj1}}, \mbox{\texttt{\mdseries\slshape obj2}}... etc. \emph{in a short form} on the standard output by calling the \texttt{ViewObj} (\ref{ViewObj}) operation on each of them. \texttt{View} is called in the read\texttt{\symbol{45}}eval\texttt{\symbol{45}}print loop,
thus the output looks exactly like the representation of the objects shown by
the main loop. Note that no space or newline is printed between the objects. }

 

\subsection{\textcolor{Chapter }{Print}}
\logpage{[ 6, 3, 4 ]}\nobreak
\hyperdef{L}{X7AFA64D97A1F39A3}{}
{\noindent\textcolor{FuncColor}{$\triangleright$\enspace\texttt{Print({\mdseries\slshape obj1, obj2, ...})\index{Print@\texttt{Print}}
\label{Print}
}\hfill{\scriptsize (function)}}\\


 Also \texttt{Print} shows the objects \mbox{\texttt{\mdseries\slshape obj1}}, \mbox{\texttt{\mdseries\slshape obj2}}... etc. on the standard output. The difference compared to \texttt{View} (\ref{View}) is in general that the shown form is not required to be short, and that in
many cases the form shown by \texttt{Print} is \textsf{GAP} readable. 

    
\begin{Verbatim}[commandchars=!@|,fontsize=\small,frame=single,label=Example]
  !gapprompt@gap>| !gapinput@z:= Z(2);|
  Z(2)^0
  !gapprompt@gap>| !gapinput@v:= [ z, z, z, z, z, z, z ];|
  [ Z(2)^0, Z(2)^0, Z(2)^0, Z(2)^0, Z(2)^0, Z(2)^0, Z(2)^0 ]
  !gapprompt@gap>| !gapinput@ConvertToVectorRep(v);; v;|
  <a GF2 vector of length 7>
  !gapprompt@gap>| !gapinput@Print( v, "\n" );|
  [ Z(2)^0, Z(2)^0, Z(2)^0, Z(2)^0, Z(2)^0, Z(2)^0, Z(2)^0 ]
\end{Verbatim}
 

 Another difference is that \texttt{Print} shows strings without the enclosing quotes, so \texttt{Print} can be used to produce formatted text on the standard output (see also
chapter{\nobreakspace}\ref{Strings and Characters}). Some characters preceded by a backslash, such as \texttt{\texttt{\symbol{92}}n}, are processed specially (see chapter{\nobreakspace}\ref{Special Characters}). \texttt{PrintTo} (\ref{PrintTo}) can be used to print to a file. 

 
\begin{Verbatim}[commandchars=!@|,fontsize=\small,frame=single,label=Example]
  !gapprompt@gap>| !gapinput@for i in [1..5] do|
  !gapprompt@>| !gapinput@     Print( i, " ", i^2, " ", i^3, "\n" );|
  !gapprompt@>| !gapinput@   od;|
  1 1 1
  2 4 8
  3 9 27
  4 16 64
  5 25 125
  !gapprompt@gap>| !gapinput@g:= SmallGroup(12,5);|
  <pc group of size 12 with 3 generators>
  !gapprompt@gap>| !gapinput@Print( g, "\n" );|
  Group( [ f1, f2, f3 ] )
  !gapprompt@gap>| !gapinput@View( g );  Print( "\n" );|
  <pc group of size 12 with 3 generators>
\end{Verbatim}
 }

 

 

\subsection{\textcolor{Chapter }{ViewObj}}
\logpage{[ 6, 3, 5 ]}\nobreak
\hyperdef{L}{X815BF22186FD43C9}{}
{\noindent\textcolor{FuncColor}{$\triangleright$\enspace\texttt{ViewObj({\mdseries\slshape obj})\index{ViewObj@\texttt{ViewObj}}
\label{ViewObj}
}\hfill{\scriptsize (operation)}}\\
\noindent\textcolor{FuncColor}{$\triangleright$\enspace\texttt{PrintObj({\mdseries\slshape obj})\index{PrintObj@\texttt{PrintObj}}
\label{PrintObj}
}\hfill{\scriptsize (operation)}}\\


 The functions \texttt{View} (\ref{View}) and \texttt{Print} (\ref{Print}) actually call the operations \texttt{ViewObj} and{\nobreakspace}\texttt{PrintObj}, respectively, for each argument. By installing special methods for these
operations, it is possible to achieve special printing behavior for certain
objects (see chapter{\nobreakspace}\ref{Method Selection}). The only exceptions are strings (see Chapter{\nobreakspace}\ref{Strings and Characters}), for which the default \texttt{PrintObj} and \texttt{ViewObj} methods as well as the function \texttt{View} (\ref{View}) print also the enclosing doublequotes, whereas \texttt{Print} (\ref{Print}) strips the doublequotes. 

 The default method for \texttt{ViewObj} is to call \texttt{PrintObj}. So it is sufficient to have a \texttt{PrintObj} method for an object in order to \texttt{View} (\ref{View}) it. If one wants to supply a ``short form'' for \texttt{View} (\ref{View}), one can install additionally a method for \texttt{ViewObj}. }

 

\subsection{\textcolor{Chapter }{Display}}
\logpage{[ 6, 3, 6 ]}\nobreak
\hyperdef{L}{X83A5C59278E13248}{}
{\noindent\textcolor{FuncColor}{$\triangleright$\enspace\texttt{Display({\mdseries\slshape obj})\index{Display@\texttt{Display}}
\label{Display}
}\hfill{\scriptsize (operation)}}\\


 Displays the object \mbox{\texttt{\mdseries\slshape obj}} in a nice, formatted way which is easy to read (but might be difficult for
machines to understand). The actual format used for this depends on the type
of \mbox{\texttt{\mdseries\slshape obj}}. Each method should print a newline character as last character. 
\begin{Verbatim}[commandchars=!@|,fontsize=\small,frame=single,label=Example]
  !gapprompt@gap>| !gapinput@Display( [ [ 1, 2, 3 ], [ 4, 5, 6 ] ] * Z(5) );|
   2 4 1
   3 . 2
\end{Verbatim}
 

 One can assign a string to an object that \texttt{Print} (\ref{Print}) will use instead of the default used by \texttt{Print} (\ref{Print}), via \texttt{SetName} (\ref{SetName}). Also, \texttt{Name} (\ref{Name}) returns the string previously assigned to the object for printing, via \texttt{SetName} (\ref{SetName}). The following is an example in the context of domains. 

 
\begin{Verbatim}[commandchars=!@|,fontsize=\small,frame=single,label=Example]
  !gapprompt@gap>| !gapinput@g:= Group( (1,2,3,4) );|
  Group([ (1,2,3,4) ])
  !gapprompt@gap>| !gapinput@SetName( g, "C4" ); g;|
  C4
  !gapprompt@gap>| !gapinput@Name( g );|
  "C4"
\end{Verbatim}
 }

 When setting up examples, in particular if for beginning users, it sometimes
can be convenient to hide the structure behind a printing name. For many
objects, such as groups, this can be done using \texttt{SetName} (\ref{SetName}). If the objects however is represented internally, for example permutations
representing group elements, this function is not applicable. Instead the
function \texttt{SetNameObject} (\ref{SetNameObject}) can be used to interface with the display routines on a lower level. 

\subsection{\textcolor{Chapter }{SetNameObject}}
\logpage{[ 6, 3, 7 ]}\nobreak
\hyperdef{L}{X87E546E27A1F1FAB}{}
{\noindent\textcolor{FuncColor}{$\triangleright$\enspace\texttt{SetNameObject({\mdseries\slshape o, s})\index{SetNameObject@\texttt{SetNameObject}}
\label{SetNameObject}
}\hfill{\scriptsize (function)}}\\


 \texttt{SetNameObject} sets the string \mbox{\texttt{\mdseries\slshape s}} as display name for object \mbox{\texttt{\mdseries\slshape o}} in an interactive session. When applying \texttt{View} (\ref{View}) to object \mbox{\texttt{\mdseries\slshape o}}, for example in the system's main loop, \textsf{GAP} will print the string \mbox{\texttt{\mdseries\slshape s}}. Calling \texttt{SetNameObject} for the same object \mbox{\texttt{\mdseries\slshape o}} with \mbox{\texttt{\mdseries\slshape s}} set to \texttt{fail} (\ref{fail}) deletes the special viewing setup. Since use of this features potentially
slows down the whole print process, this function should be used sparingly. 
\begin{Verbatim}[commandchars=!@|,fontsize=\small,frame=single,label=Example]
  !gapprompt@gap>| !gapinput@SetNameObject(3,"three");|
  !gapprompt@gap>| !gapinput@Filtered([1..10],IsPrimeInt);|
  [ 2, three, 5, 7 ]
  !gapprompt@gap>| !gapinput@SetNameObject(3,fail);|
  !gapprompt@gap>| !gapinput@Filtered([1..10],IsPrimeInt);|
  [ 2, 3, 5, 7 ]
\end{Verbatim}
 }

 }

  
\section{\textcolor{Chapter }{Break Loops}}\label{Break Loops}
\logpage{[ 6, 4, 0 ]}
\hyperdef{L}{X8593B49F8705B486}{}
{
  When an error has occurred or when you interrupt \textsf{GAP} (usually by hitting \textsc{Ctrl\texttt{\symbol{45}}C}) \textsf{GAP} enters a break loop, that is in most respects like the main read eval print
loop (see{\nobreakspace}\ref{Main Loop}). That is, you can enter statements, \textsf{GAP} reads them, evaluates them, and shows the result if any. However those
evaluations happen within the context in which the error occurred. So you can
look at the arguments and local variables of the functions that were active
when the error happened and even change them. The prompt is changed from \texttt{gap{\textgreater}} to \texttt{brk{\textgreater}} to indicate that you are in a break loop. 

 
\begin{Verbatim}[commandchars=!@|,fontsize=\small,frame=single,label=Example]
  !gapprompt@gap>| !gapinput@1/0;|
  Error, Rational operations: <divisor> must not be zero
  not in any function at *stdin*:2
  you can enter 'quit;' to quit to outer loop
\end{Verbatim}
 

 If errors occur within a break loop \textsf{GAP} enters another break loop at a \emph{deeper level}. This is indicated by a number appended to \texttt{brk}: 

 
\begin{Verbatim}[commandchars=!@|,fontsize=\small,frame=single,label=Example]
  !gapbrkprompt@brk>| !gapinput@1/0;|
  Error, Rational operations: <divisor> must not be zero
  not in any function at *errin*:1
  you can enter 'quit;' to quit to outer loop
  !gapbrkprompt@brk_2>| !gapinput@|
\end{Verbatim}
 

 There are two ways to leave a break loop, see \ref{quit} and \ref{subsect:return}. 
\subsection{\textcolor{Chapter }{quit from a break loop}}\label{quit}
\logpage{[ 6, 4, 1 ]}
\hyperdef{L}{X83033EEB81CF4F49}{}
{
  The first way to leave a break loop is to \emph{quit} the break loop. To do this you enter \texttt{quit;} or type the \emph{eof} (\emph{e}nd \emph{o}f \emph{f}ile) character, which is usually \textsc{Ctrl\texttt{\symbol{45}}D} except when using the \texttt{\texttt{\symbol{45}}e} option (see Section{\nobreakspace}\ref{Command Line Options}). Note that \textsf{GAP} code between \texttt{quit;} and the end of the input line is ignored. 

 
\begin{Verbatim}[commandchars=!@|,fontsize=\small,frame=single,label=Example]
  !gapbrkprompt@brk_2>| !gapinput@quit;|
  !gapbrkprompt@brk>| !gapinput@|
\end{Verbatim}
 

 In this case control returns to the break loop one level above or to the main
loop, respectively. So iterated break loops must be left iteratively. Note
also that if you type \texttt{quit;} from a \texttt{gap{\textgreater}} prompt, \textsf{GAP} will exit (see{\nobreakspace}\ref{Leaving GAP}). 

 \emph{Note:} If you leave a break loop with \texttt{quit} without completing a command it is possible (though not very likely) that data
structures will be corrupted or incomplete data have been stored in objects.
Therefore no guarantee can be given that calculations afterwards will return
correct results! If you have been using options \texttt{quit}ting a break loop generally leaves the options stack with options you no
longer want. The function \texttt{ResetOptionsStack} (\ref{ResetOptionsStack}) removes all options on the options stack, and this is the sole intended
purpose of this function. }

 
\subsection{\textcolor{Chapter }{return from a break loop}}\label{subsect:return}
\logpage{[ 6, 4, 2 ]}
\hyperdef{L}{X7A388B808167FE09}{}
{
  \index{return@\texttt{return}} \index{return from break loop} The other way to leave a break loop is to \emph{return} from a break loop. If the break loop was entered because you interrupted \textsf{GAP}, then you can continue by typing \texttt{return;}. If the break loop was entered due to an error, this option may or may not be
available, depending on the kind of error. The message printed on entering the
break loop will tell you whether this possible. }

 

\subsection{\textcolor{Chapter }{OnBreak}}
\logpage{[ 6, 4, 3 ]}\nobreak
\hyperdef{L}{X82EBF01181C3C859}{}
{\noindent\textcolor{FuncColor}{$\triangleright$\enspace\texttt{OnBreak({\mdseries\slshape })\index{OnBreak@\texttt{OnBreak}}
\label{OnBreak}
}\hfill{\scriptsize (function)}}\\


 By default, when a break loop is entered, \textsf{GAP} prints a trace of the innermost 5 commands currently being executed. This
behaviour can be configured by changing the value of the global variable \texttt{OnBreak}. When a break loop is entered, the value of \texttt{OnBreak} is checked. If it is a function, then it is called with no arguments. By
default, the value of \texttt{OnBreak} is \texttt{Where} (\ref{Where}). 

 
\begin{Verbatim}[commandchars=!@|,fontsize=\small,frame=single,label=Example]
  !gapprompt@gap>| !gapinput@OnBreak := function() Print("Hello\n"); end;|
  function(  ) ... end
\end{Verbatim}
 

 
\begin{Verbatim}[commandchars=@|A,fontsize=\small,frame=single,label=Example]
  @gapprompt|gap>A @gapinput|Error("!\n");A
  Error, !
  Hello
  you can enter 'quit;' to quit to outer loop, or
  you can enter 'return;' to continue
  @gapbrkprompt|brk>A @gapinput|quit;A
\end{Verbatim}
 

 In cases where a break loop is entered during a function that was called with
options (see Chapter{\nobreakspace}\ref{Options Stack}), a \texttt{quit;} will also cause the options stack to be reset and an \texttt{Info}\texttt{\symbol{45}}ed warning stating this is emitted at \texttt{InfoWarning} (\ref{InfoWarning}) level 1 (see Chapter{\nobreakspace}\ref{Info Functions}).            

 Note that for break loops entered by a call to \texttt{Error} (\ref{Error}), the lines just before the \texttt{brk{\textgreater}} prompt can also be customised, namely by redefining \texttt{OnBreakMessage} (\ref{OnBreakMessage}). 

 \index{ErrorNoTraceBack@\texttt{ErrorNoTraceBack}} Also, note that one can achieve the effect of changing \texttt{OnBreak} \emph{locally}. As mentioned above, the default value of \texttt{OnBreak} is \texttt{Where} (\ref{Where}). Thus, a call to \texttt{Error} (\ref{Error}) generally gives a trace back up to five levels of calling functions.
Conceivably, we might like to have a function like \texttt{Error} (\ref{Error}) that does not trace back without globally changing \texttt{OnBreak}. Such a function we might call \texttt{ErrorNoTraceBack} and here is how we might define it. (Note \texttt{ErrorNoTraceBack} is \emph{not} a \textsf{GAP} function.) 

 
\begin{Verbatim}[commandchars=!@|,fontsize=\small,frame=single,label=Example]
  !gapprompt@gap>| !gapinput@ErrorNoTraceBack := function(arg) # arg is special variable that GAP|
  !gapprompt@>| !gapinput@                                     # knows to treat as list of arg's|
  !gapprompt@>| !gapinput@     local SavedOnBreak, ENTBOnBreak;|
  !gapprompt@>| !gapinput@     SavedOnBreak := OnBreak;        # save current value of OnBreak|
  !gapprompt@>| !gapinput@|
  !gapprompt@>| !gapinput@     ENTBOnBreak := function()       # our `local' OnBreak|
  !gapprompt@>| !gapinput@     local s;|
  !gapprompt@>| !gapinput@       for s in arg do|
  !gapprompt@>| !gapinput@         Print(s);|
  !gapprompt@>| !gapinput@       od;|
  !gapprompt@>| !gapinput@       OnBreak := SavedOnBreak;      # restore OnBreak afterwards|
  !gapprompt@>| !gapinput@     end;|
  !gapprompt@>| !gapinput@|
  !gapprompt@>| !gapinput@     OnBreak := ENTBOnBreak;|
  !gapprompt@>| !gapinput@     Error();|
  !gapprompt@>| !gapinput@   end;|
  function( arg... ) ... end
\end{Verbatim}
 

 Here is a somewhat trivial demonstration of the use of \texttt{ErrorNoTraceBack}. 

 
\begin{Verbatim}[commandchars=@|A,fontsize=\small,frame=single,label=Example]
  @gapprompt|gap>A @gapinput|ErrorNoTraceBack("Gidday!", " How's", " it", " going?\n");A
  Error,
  Gidday! How's it going?
  you can enter 'quit;' to quit to outer loop, or
  you can enter 'return;' to continue
  @gapbrkprompt|brk>A @gapinput|quit;A
\end{Verbatim}
 

 Now we call \texttt{Error} (\ref{Error}) with the same arguments to show the difference. 

 
\begin{Verbatim}[commandchars=@|A,fontsize=\small,frame=single,label=Example]
  @gapprompt|gap>A @gapinput|Error("Gidday!", " How's", " it", " going?\n");A
  Error, Gidday! How's it going?
  Hello
  you can enter 'quit;' to quit to outer loop, or
  you can enter 'return;' to continue
  @gapbrkprompt|brk>A @gapinput|quit;A
\end{Verbatim}
 

 Observe that the value of \texttt{OnBreak} before the \texttt{ErrorNoTraceBack} call was restored. However, we had changed \texttt{OnBreak} from its default value; to restore \texttt{OnBreak} to its default value, we should do the following. 

 
\begin{Verbatim}[commandchars=!@|,fontsize=\small,frame=single,label=Example]
  !gapprompt@gap>| !gapinput@OnBreak := Where;;|
\end{Verbatim}
 }

 

\subsection{\textcolor{Chapter }{OnBreakMessage}}
\logpage{[ 6, 4, 4 ]}\nobreak
\hyperdef{L}{X80711C807C99C220}{}
{\noindent\textcolor{FuncColor}{$\triangleright$\enspace\texttt{OnBreakMessage({\mdseries\slshape })\index{OnBreakMessage@\texttt{OnBreakMessage}}
\label{OnBreakMessage}
}\hfill{\scriptsize (function)}}\\


 \index{Break loop message} When a break loop is entered by a call to \texttt{Error} (\ref{Error}) the message at the end starting with ``\texttt{you can enter 'quit;' to ...}'' is produced by the function \texttt{OnBreakMessage}, which just like \texttt{OnBreak} (\ref{OnBreak}) is a user\texttt{\symbol{45}}configurable global variable that is a \emph{function} with \emph{no arguments}. 

 
\begin{Verbatim}[commandchars=!@|,fontsize=\small,frame=single,label=Example]
  !gapprompt@gap>| !gapinput@OnBreakMessage(); # By default, OnBreakMessage prints the following|
  you can enter 'quit;' to quit to outer loop, or
  you can enter 'return;' to continue
\end{Verbatim}
 

 Perhaps you are familiar with what's possible in a break loop, and so don't
need to be reminded. In this case, you might wish to do the following (the
first line just makes it easy to restore the default value later). 

 
\begin{Verbatim}[commandchars=!@|,fontsize=\small,frame=single,label=Example]
  !gapprompt@gap>| !gapinput@NormalOnBreakMessage := OnBreakMessage;; # save the default value|
  !gapprompt@gap>| !gapinput@OnBreakMessage := function() end;        # do-nothing function|
  function(  ) ... end
  !gapprompt@gap>| !gapinput@OnBreakMessage();|
  !gapprompt@gap>| !gapinput@OnBreakMessage := NormalOnBreakMessage;; # reset|
\end{Verbatim}
 

 With \texttt{OnBreak} (\ref{OnBreak}) still set away from its default value, calling \texttt{Error} (\ref{Error}) as we did above, now produces: 

 
\begin{Verbatim}[commandchars=@|A,fontsize=\small,frame=single,label=Example]
  @gapprompt|gap>A @gapinput|Error("!\n");A
  Error, !
  Hello
  @gapbrkprompt|brk>A @gapinput|quit; # to get back to outer loopA
\end{Verbatim}
 

 However, suppose you are writing a function which detects an error condition
and \texttt{OnBreakMessage} needs to be changed only \emph{locally}, i.e., the instructions on how to recover from the break loop need to be
specific to that function. The same idea used to define \texttt{ErrorNoTraceBack} (see \texttt{OnBreak} (\ref{OnBreak})) can be adapted to achieve this. The function \texttt{CosetTableFromGensAndRels} (\ref{CosetTableFromGensAndRels}) is an example in the \textsf{GAP} code where the idea is actually used. }

 

 

\subsection{\textcolor{Chapter }{Where}}
\logpage{[ 6, 4, 5 ]}\nobreak
\hyperdef{L}{X7A7FFA2B7C1EF5A3}{}
{\noindent\textcolor{FuncColor}{$\triangleright$\enspace\texttt{Where({\mdseries\slshape nr})\index{Where@\texttt{Where}}
\label{Where}
}\hfill{\scriptsize (function)}}\\
\noindent\textcolor{FuncColor}{$\triangleright$\enspace\texttt{WhereWithVars({\mdseries\slshape nr})\index{WhereWithVars@\texttt{WhereWithVars}}
\label{WhereWithVars}
}\hfill{\scriptsize (function)}}\\


 \index{Backtrace!GAP3 name for Where} \index{Stack trace} shows the last \mbox{\texttt{\mdseries\slshape nr}} commands on the execution stack during whose execution the error occurred. If
not given, \mbox{\texttt{\mdseries\slshape nr}} defaults to the value of the \texttt{WhereDepth} user preference (initially 5). If \mbox{\texttt{\mdseries\slshape nr}} or the \texttt{WhereDepth} preference is not a non\texttt{\symbol{45}}negative integer, it is treated as
5. (Assume, for the following example, that after the last example \texttt{OnBreak} (\ref{OnBreak}) has been set back to its default value.). \texttt{WhereWithVars} acts the same as \texttt{Where} while also showing the arguments and local variables of each function. 

 \texttt{WhereWithVars} acts the same as \texttt{Where} while also showing the arguments and local variables of each function. 

 The following example assumes that \texttt{OnBreak} (\ref{OnBreak}) is set to its default value. 
\begin{Verbatim}[commandchars=!|D,fontsize=\small,frame=single,label=Example]
  !gapprompt|gap>D !gapinput|StabChain(SymmetricGroup(1000)); # After this we typed ^CD
  Error, user interrupt
  *[1] genlabels := Filtered( genlabels, function ( x )
          return pnt ^ sgs[x] = pnt;
      end );
     @ GAPROOT/lib/stbc.gi:426
   [2] StabChainBaseStrongGenerators( dom{[ 1 .. Length( dom ) - 1 ]}, sgs, () )
     @ GAPROOT/lib/gpprmsya.gi:2006
   [3] StabChainOp( G, rec(
       ) )
     @ GAPROOT/lib/stbc.gi:28
   [4] StabChainImmutable( arg[1] )
     @ GAPROOT/lib/stbc.gi:18
  <function "StabChain">( <arguments> )
   called from read-eval loop at *stdin*:2
  you can enter 'quit;' to quit to outer loop, or
  you can enter 'return;' to continue
  !gapbrkprompt|brk>D !gapinput|Where(2);D
  *[1] genlabels := Filtered( genlabels, function ( x )
          return pnt ^ sgs[x] = pnt;
      end );
     @ GAPROOT/lib/stbc.gi:426
   [2] StabChainBaseStrongGenerators( dom{[ 1 .. Length( dom ) - 1 ]}, sgs, () )
     @ GAPROOT/lib/gpprmsya.gi:2006
  ...  at *errin*:1
\end{Verbatim}
 

 Note that the variables displayed even in the first line of the \texttt{Where} list (after the \texttt{called from} line) may be already one environment level higher and \texttt{DownEnv} (\ref{DownEnv}) may be necessary to access them. 

 At the moment this backtrace does not work from within compiled code (this
includes the method selection which by default is compiled into the kernel).
If this creates problems for debugging, call \textsf{GAP} with the \texttt{\texttt{\symbol{45}}M} option (see{\nobreakspace}\ref{Command Line Options}) to avoid loading compiled code. 

 (Function calls to \texttt{Info} (\ref{Info}) and methods installed for binary operations are handled in a special way. In
rare circumstances it is possible therefore that they do not show up in a \texttt{Where} log but the log refers to the \emph{last} proper function call that happened before.) 

 The command line option \texttt{\texttt{\symbol{45}}T} to \textsf{GAP} disables the break loop. This is mainly intended for testing purposes and for
special applications. If this option is given then errors simply cause \textsf{GAP} to return to the main loop. }

 }

  
\section{\textcolor{Chapter }{Variable Access in a Break Loop}}\label{Variable Access in a Break Loop}
\logpage{[ 6, 5, 0 ]}
\hyperdef{L}{X7EE5CF2C8419F061}{}
{
  In a break loop access to variables of the current break level and higher
levels is possible, but if the same variable name is used for different
objects or if a function calls itself recursively, of course only the variable
at the lowest level can be accessed. 
\subsection{\textcolor{Chapter }{DownEnv and UpEnv}}\logpage{[ 6, 5, 1 ]}
\hyperdef{L}{X79E66DA2875303B0}{}
{
\noindent\textcolor{FuncColor}{$\triangleright$\enspace\texttt{DownEnv({\mdseries\slshape nr})\index{DownEnv@\texttt{DownEnv}}
\label{DownEnv}
}\hfill{\scriptsize (function)}}\\
\noindent\textcolor{FuncColor}{$\triangleright$\enspace\texttt{UpEnv({\mdseries\slshape nr})\index{UpEnv@\texttt{UpEnv}}
\label{UpEnv}
}\hfill{\scriptsize (function)}}\\
\noindent\textcolor{FuncColor}{$\triangleright$\enspace\texttt{CurrentEnv({\mdseries\slshape })\index{CurrentEnv@\texttt{CurrentEnv}}
\label{CurrentEnv}
}\hfill{\scriptsize (function)}}\\


 \texttt{DownEnv} moves down \mbox{\texttt{\mdseries\slshape nr}} steps in the environment and allows one to inspect variables on this level; if \mbox{\texttt{\mdseries\slshape nr}} is negative it steps up in the environment again; \mbox{\texttt{\mdseries\slshape nr}} defaults to 1 if not given. \texttt{UpEnv} acts similarly to \texttt{DownEnv} but in the reverse direction (the mnemonic rule to remember the difference
between \texttt{DownEnv} and \texttt{UpEnv} is the order in which commands on the execution stack are displayed by \texttt{Where} (\ref{Where})). \texttt{CurrentEnv} returns the currently selected environment. At present this is the local
variables bag used by the break loop for local variable lookup; outside such a
context it returns \texttt{fail}. 

 
\begin{Verbatim}[commandchars=|AC,fontsize=\small,frame=single,label=Example]
  |gappromptAgap>C |gapinputAOnBreak := function() end;; # eliminate back-tracing on entry to break loopC
  |gappromptAgap>C |gapinputAtest:= function( n )C
  |gappromptA>C |gapinputA   if n > 3 then Error( "!\n" ); fi; test( n+1 ); end;;C
  |gappromptAgap>C |gapinputAtest( 1 );C
  Error, !
  you can enter 'quit;' to quit to outer loop, or
  you can enter 'return;' to continue
  |gapbrkpromptAbrk>C |gapinputAWhere();C
  Stack trace:
  *[1] Error( "!\n" );
     @ *stdin*:3
   [2] test( n + 1 );
     @ *stdin*:3
   [3] test( n + 1 );
     @ *stdin*:3
   [4] test( n + 1 );
     @ *stdin*:3
  <function "test">( <arguments> )
   called from read-eval loop at *errin*:1
  |gapbrkpromptAbrk>C |gapinputAn;C
  4
  |gapbrkpromptAbrk>C |gapinputADownEnv();C
  |gapbrkpromptAbrk>C |gapinputAn;C
  3
  |gapbrkpromptAbrk>C |gapinputAWhere();C
  Stack trace:
   [1] Error( "!\n" );
     @ *stdin*:3
  *[2] test( n + 1 );
     @ *stdin*:3
   [3] test( n + 1 );
     @ *stdin*:3
   [4] test( n + 1 );
     @ *stdin*:3
  <function "test">( <arguments> )
   called from read-eval loop at *errin*:5
  |gapbrkpromptAbrk>C |gapinputADownEnv( 2 );C
  |gapbrkpromptAbrk>C |gapinputAn;C
  1
  |gapbrkpromptAbrk>C |gapinputAWhere();C
  Stack trace:
   [1] Error( "!\n" );
     @ *stdin*:3
   [2] test( n + 1 );
     @ *stdin*:3
   [3] test( n + 1 );
     @ *stdin*:3
  *[4] test( n + 1 );
     @ *stdin*:3
  <function "test">( <arguments> )
   called from read-eval loop at *errin*:8
  |gapbrkpromptAbrk>C |gapinputADownEnv( -2 );C
  |gapbrkpromptAbrk>C |gapinputAn;C
  3
  |gapbrkpromptAbrk>C |gapinputAquit;C
  |gappromptAgap>C |gapinputAOnBreak := Where;; # restore OnBreak to its default valueC
\end{Verbatim}
 

 Note that the change of the environment caused by \texttt{DownEnv} only affects variable access in the break loop. If you use \texttt{return} to continue a calculation \textsf{GAP} automatically jumps to the right environment level again. 

 Note also that search for variables looks first in the chain of outer
functions which enclosed the definition of a currently executing function,
before it looks at the chain of calling functions which led to the current
invocation of the function. 

 
\begin{Verbatim}[commandchars=|AB,fontsize=\small,frame=single,label=Example]
  |gappromptAgap>B |gapinputAfoo := function()B
  |gappromptA>B |gapinputAlocal x; x := 1;B
  |gappromptA>B |gapinputAreturn function() local y; y := x*x; Error("!!\n"); end;B
  |gappromptA>B |gapinputAend;B
  function(  ) ... end
  |gappromptAgap>B |gapinputAbar := foo();B
  function(  ) ... end
  |gappromptAgap>B |gapinputAfun := function() local x; x := 3; bar(); end;B
  function(  ) ... end
  |gappromptAgap>B |gapinputAfun();B
  Error, !!
  Stack trace:
  *[1] Error( "!!\n" );
     @ *stdin*:4
   [2] bar(  );
     @ *stdin*:7
  <function "fun">( <arguments> )
   called from read-eval loop at *stdin*:8
  you can enter 'quit;' to quit to outer loop, or
  you can enter 'return;' to continue
  |gapbrkpromptAbrk>B |gapinputAx;B
  1
  |gapbrkpromptAbrk>B |gapinputADownEnv(1);B
  |gapbrkpromptAbrk>B |gapinputAx;B
  3
\end{Verbatim}
 

 Here the \texttt{x} of \texttt{foo} which contained the definition of \texttt{bar} is found before that of \texttt{fun} which caused its execution. Using \texttt{DownEnv} we can access the \texttt{x} from \texttt{fun}. }

 }

  
\section{\textcolor{Chapter }{Error and ErrorCount}}\label{sect:Error}
\logpage{[ 6, 6, 0 ]}
\hyperdef{L}{X7BC8D2E37ADE9062}{}
{
  

\subsection{\textcolor{Chapter }{Error}}
\logpage{[ 6, 6, 1 ]}\nobreak
\hyperdef{L}{X7E7AD8D87EBA1A08}{}
{\noindent\textcolor{FuncColor}{$\triangleright$\enspace\texttt{Error({\mdseries\slshape messages, ...})\index{Error@\texttt{Error}}
\label{Error}
}\hfill{\scriptsize (function)}}\\


 \texttt{Error} signals an error from within a function. First the messages \mbox{\texttt{\mdseries\slshape messages}} are printed, this is done exactly as if \texttt{Print} (\ref{Print}) (see{\nobreakspace}\ref{View and Print}) were called with these arguments. Then a break loop (see{\nobreakspace}\ref{Break Loops}) is entered, unless the standard error output is not connected to a terminal.
You can leave this break loop with \texttt{return;} to continue execution with the statement following the call to \texttt{Error}. \texttt{ErrorNoReturn} (\ref{ErrorNoReturn}) operates identically to \texttt{Error}, except it does not allow using \texttt{return;} to continue execution. }

 

\subsection{\textcolor{Chapter }{ErrorNoReturn}}
\logpage{[ 6, 6, 2 ]}\nobreak
\hyperdef{L}{X7A5C000D7E4984DD}{}
{\noindent\textcolor{FuncColor}{$\triangleright$\enspace\texttt{ErrorNoReturn({\mdseries\slshape messages, ...})\index{ErrorNoReturn@\texttt{ErrorNoReturn}}
\label{ErrorNoReturn}
}\hfill{\scriptsize (function)}}\\


 \texttt{ErrorNoReturn} signals an error from within a function. First the messages \mbox{\texttt{\mdseries\slshape messages}} are printed, this is done exactly as if \texttt{Print} (\ref{Print}) (see{\nobreakspace}\ref{View and Print}) were called with these arguments. Then a break loop (see{\nobreakspace}\ref{Break Loops}) is entered, unless the standard error output is not connected to a terminal.
This break loop can only be exited with \texttt{quit;}. The function differs from \texttt{Error} (\ref{Error}) by not allowing execution to continue. }

 

\subsection{\textcolor{Chapter }{ErrorCount}}
\logpage{[ 6, 6, 3 ]}\nobreak
\hyperdef{L}{X86A11BCC7FECEEA4}{}
{\noindent\textcolor{FuncColor}{$\triangleright$\enspace\texttt{ErrorCount({\mdseries\slshape })\index{ErrorCount@\texttt{ErrorCount}}
\label{ErrorCount}
}\hfill{\scriptsize (function)}}\\


 \texttt{ErrorCount} returns a count of the number of errors (including user interruptions) which
have occurred in the \textsf{GAP} session so far. The count is incremented by each error, even if \textsf{GAP} was started with the \texttt{\texttt{\symbol{45}}T} option to disable the break loop. }

 }

  
\section{\textcolor{Chapter }{Leaving GAP}}\label{Leaving GAP}
\logpage{[ 6, 7, 0 ]}
\hyperdef{L}{X83704B1080FD9B40}{}
{
  \index{quit!in emergency} \index{exit} \index{at exit functions} \index{saving on exit} The normal way to terminate a \textsf{GAP} session is to enter either \texttt{quit;} (note the semicolon) or an end\texttt{\symbol{45}}of\texttt{\symbol{45}}file
character (usually \textsc{Ctrl\texttt{\symbol{45}}D}) at the \texttt{gap{\textgreater} } prompt in the main read eval print loop. 

\subsection{\textcolor{Chapter }{QUIT}}
\logpage{[ 6, 7, 1 ]}\nobreak
\hyperdef{L}{X7ECC75048583853B}{}
{\noindent\textcolor{FuncColor}{$\triangleright$\enspace\texttt{QUIT\index{QUIT@\texttt{QUIT}}
\label{QUIT}
}\hfill{\scriptsize (global variable)}}\\


 An emergency way to leave \textsf{GAP} is to enter \index{QUIT@\texttt{QUIT}!emergency quit} \texttt{QUIT} at any \texttt{gap{\textgreater}} or \texttt{brk{\textgreater}} or \texttt{brk{\textunderscore}\mbox{\texttt{\mdseries\slshape nn}}{\textgreater}} prompt. }

 

\subsection{\textcolor{Chapter }{GapExitCode}}
\logpage{[ 6, 7, 2 ]}\nobreak
\hyperdef{L}{X838B50A9790DE55B}{}
{\noindent\textcolor{FuncColor}{$\triangleright$\enspace\texttt{GapExitCode({\mdseries\slshape [ret]})\index{GapExitCode@\texttt{GapExitCode}}
\label{GapExitCode}
}\hfill{\scriptsize (function)}}\\


 \texttt{GapExitCode} sets the exit value which is returned to the operating system (or parent
process) when \textsf{GAP} exits. This may be an integer in the range [\texttt{\symbol{45}}128..127]
(other values are reduced modulo 256), or a boolean. \texttt{true} corresponds to the return value 0, which by convention is treated as
"success". \texttt{false} corresponds to the return value 1, which by convention is treated as
"failure". The exit value is not changed if no argument is given. 

 The \emph{previous} exit code is returned. }

 

\subsection{\textcolor{Chapter }{QuitGap}}
\logpage{[ 6, 7, 3 ]}\nobreak
\hyperdef{L}{X7AB1567987922580}{}
{\noindent\textcolor{FuncColor}{$\triangleright$\enspace\texttt{QuitGap({\mdseries\slshape [ret]})\index{QuitGap@\texttt{QuitGap}}
\label{QuitGap}
}\hfill{\scriptsize (function)}}\\


 \texttt{QuitGap} acts similarly to the keyword \texttt{QUIT}, except \texttt{QUIT} cannot be called from a function. It exits \textsf{GAP} cleanly, calling any function installed using \texttt{InstallAtExit} (\ref{InstallAtExit}). The optional argument \mbox{\texttt{\mdseries\slshape ret}} will be passed to \texttt{GapExitCode} (\ref{GapExitCode}). 

 }

 

\subsection{\textcolor{Chapter }{ForceQuitGap}}
\logpage{[ 6, 7, 4 ]}\nobreak
\hyperdef{L}{X85A8DD6B7A20DD89}{}
{\noindent\textcolor{FuncColor}{$\triangleright$\enspace\texttt{ForceQuitGap({\mdseries\slshape [ret]})\index{ForceQuitGap@\texttt{ForceQuitGap}}
\label{ForceQuitGap}
}\hfill{\scriptsize (function)}}\\


 \texttt{ForceQuitGap} is similar to \texttt{QuitGap} (\ref{QuitGap}), except it ignores any functions installed with \texttt{InstallAtExit} (\ref{InstallAtExit}), or any other functions normally run at GAP exit, such as flushing any
partially outputted lines to both the screen and files, and exits GAP
immediately. The optional argument \mbox{\texttt{\mdseries\slshape ret}} will be passed to \texttt{GapExitCode} (\ref{GapExitCode}). 

 }

 

\subsection{\textcolor{Chapter }{InstallAtExit}}
\logpage{[ 6, 7, 5 ]}\nobreak
\hyperdef{L}{X7A2C380986F46FEE}{}
{\noindent\textcolor{FuncColor}{$\triangleright$\enspace\texttt{InstallAtExit({\mdseries\slshape func})\index{InstallAtExit@\texttt{InstallAtExit}}
\label{InstallAtExit}
}\hfill{\scriptsize (function)}}\\
\noindent\textcolor{FuncColor}{$\triangleright$\enspace\texttt{QUITTING\index{QUITTING@\texttt{QUITTING}}
\label{QUITTING}
}\hfill{\scriptsize (global variable)}}\\


 Before actually terminating, \textsf{GAP} will call (with no arguments) all of the functions that have been installed
using \texttt{InstallAtExit}. These typically perform tasks such as cleaning up temporary files created
during the session, and closing open files. If an error occurs during the
execution of one of these functions, that function is simply abandoned, no
break loop is entered. 

  
\begin{Verbatim}[commandchars=!@|,fontsize=\small,frame=single,label=Example]
  !gapprompt@gap>| !gapinput@InstallAtExit(function() Print("bye\n"); end);|
  !gapprompt@gap>| !gapinput@quit;|
  bye
\end{Verbatim}
 

 During execution of these functions, the global variable \texttt{QUITTING} will be set to \texttt{true} if \textsf{GAP} is exiting because the user typed \texttt{QUIT} and \texttt{false} otherwise. Since \texttt{QUIT} is considered as an emergency measure, different action may be appropriate. }

 

\subsection{\textcolor{Chapter }{SaveOnExitFile}}
\logpage{[ 6, 7, 6 ]}\nobreak
\hyperdef{L}{X843C07A4869EAA1D}{}
{\noindent\textcolor{FuncColor}{$\triangleright$\enspace\texttt{SaveOnExitFile\index{SaveOnExitFile@\texttt{SaveOnExitFile}}
\label{SaveOnExitFile}
}\hfill{\scriptsize (global variable)}}\\


 If, when \textsf{GAP} is exiting due to a \texttt{quit} or end\texttt{\symbol{45}}of\texttt{\symbol{45}}file (i.e. not due to a \texttt{QUIT}) the variable \texttt{SaveOnExitFile} is bound to a string value, then the system will try to save the \textsf{GAP} workspace to that file, see \texttt{SaveWorkspace} (\ref{SaveWorkspace}). }

 }

  
\section{\textcolor{Chapter }{Line Editing}}\label{Line Editing}
\logpage{[ 6, 8, 0 ]}
\hyperdef{L}{X82234FD181899530}{}
{
  In most installations \textsf{GAP} will be compiled to use the Gnu readline library (see the line \texttt{Libs used:} on \textsf{GAP} startup). In that case skip to the next section \ref{sec:readline}. (The line editing commands described in the rest of this section were
available in previous versions of \textsf{GAP}, they will work almost the same in the standard configuration of the Gnu
readline library.) 

 \textsf{GAP} allows one you to edit the current input line with a number of editing
commands. Those commands are accessible either as \emph{control keys} or as \emph{escape keys}. You enter a control key by pressing the \textsc{Ctrl} key, and, while still holding the \textsc{Ctrl} key down, hitting another key \texttt{key}. You enter an escape key by hitting \textsc{Esc} and then hitting another key \texttt{key}. Below we denote control keys by \textsc{Ctrl\texttt{\symbol{45}}}\texttt{key} and escape keys by \textsc{Esc\texttt{\symbol{45}}}\texttt{key}. The case of \texttt{key} does not matter, i.e., \textsc{Ctrl\texttt{\symbol{45}}A} and \textsc{Ctrl\texttt{\symbol{45}}a} are equivalent. 

 Normally, line editing will be enabled if the input is connected to a
terminal. Line editing can be enabled or disabled using the command line
options \texttt{\texttt{\symbol{45}}f} and \texttt{\texttt{\symbol{45}}n} respectively (see{\nobreakspace}\ref{Command Line Options}), however this is a machine dependent feature of \textsf{GAP}. 

 Typing \textsc{Ctrl\texttt{\symbol{45}}key} or \textsc{Esc\texttt{\symbol{45}}key} for characters not mentioned below always inserts \textsc{Ctrl\texttt{\symbol{45}}}\texttt{key} resp.{\nobreakspace}\textsc{Esc\texttt{\symbol{45}}}\texttt{key} at the current cursor position. 

 The first few commands allow you to move the cursor on the current line. 

 
\begin{description}
\item[{\textsc{Ctrl\texttt{\symbol{45}}A}}]  move the cursor to the beginning of the line. 
\item[{\textsc{Esc\texttt{\symbol{45}}B}}]  move the cursor to the beginning of the previous word. 
\item[{\textsc{Ctrl\texttt{\symbol{45}}B}}]  move the cursor backward one character. 
\item[{\textsc{Ctrl\texttt{\symbol{45}}F}}]  move the cursor forward one character. 
\item[{\textsc{Esc\texttt{\symbol{45}}F}}]  move the cursor to the end of the next word. 
\item[{\textsc{Ctrl\texttt{\symbol{45}}E}}]  move the cursor to the end of the line. 
\end{description}
 

 The next commands delete or kill text. The last killed text can be reinserted,
possibly at a different position, with the ``yank'' command \textsc{Ctrl\texttt{\symbol{45}}Y}. 
\begin{description}
\item[{\textsc{Ctrl\texttt{\symbol{45}}H} or \mbox{\texttt{\mdseries\slshape del}}}]  delete the character left of the cursor. 
\item[{\textsc{Ctrl\texttt{\symbol{45}}D}}]  delete the character under the cursor. 
\item[{\textsc{Ctrl\texttt{\symbol{45}}K}}]  kill up to the end of the line. 
\item[{\textsc{Esc\texttt{\symbol{45}}D}}]  kill forward to the end of the next word. 
\item[{\textsc{Esc\texttt{\symbol{45}}del}}]  kill backward to the beginning of the last word. 
\item[{\textsc{Ctrl\texttt{\symbol{45}}X}}]  kill entire input line, and discard all pending input. 
\item[{\textsc{Ctrl\texttt{\symbol{45}}Y}}]  insert (yank) a just killed text. 
\end{description}
 

 The next commands allow you to change the input. 

 
\begin{description}
\item[{\textsc{Ctrl\texttt{\symbol{45}}T}}]  exchange (twiddle) current and previous character. 
\item[{\textsc{Esc\texttt{\symbol{45}}U}}]  uppercase next word. 
\item[{\textsc{Esc\texttt{\symbol{45}}L}}]  lowercase next word. 
\item[{\textsc{Esc\texttt{\symbol{45}}C}}]  capitalize next word. 
\end{description}
 

 The \textsc{Tab} character, which is in fact the control key \textsc{Ctrl\texttt{\symbol{45}}I}, looks at the characters before the cursor, interprets them as the beginning
of an identifier and tries to complete this identifier. If there is more than
one possible completion, it completes to the longest common prefix of all
those completions. If the characters to the left of the cursor are already the
longest common prefix of all completions hitting \textsc{Tab} a second time will display all possible completions. 

 
\begin{description}
\item[{\textsc{tab}}]  complete the identifier before the cursor. 
\end{description}
 

 The next commands allow you to fetch previous lines, e.g., to correct typos,
etc. 

 
\begin{description}
\item[{\textsc{Ctrl\texttt{\symbol{45}}L}}]  insert last input line before current character. 
\item[{\textsc{Ctrl\texttt{\symbol{45}}P}}]  redisplay the last input line, another \textsc{Ctrl\texttt{\symbol{45}}P} will redisplay the line before that, etc. If the cursor is not in the first
column only the lines starting with the string to the left of the cursor are
taken. 
\item[{\textsc{Ctrl\texttt{\symbol{45}}N}}]  Like \textsc{Ctrl\texttt{\symbol{45}}P} but goes the other way round through the history. 
\item[{\textsc{Esc\texttt{\symbol{45}}{\textless}}}]  goes to the beginning of the history. 
\item[{\textsc{Esc\texttt{\symbol{45}}{\textgreater}}}]  goes to the end of the history. 
\item[{\textsc{Ctrl\texttt{\symbol{45}}O}}]  accepts this line and perform a \textsc{Ctrl\texttt{\symbol{45}}N}. 
\end{description}
 

 Finally there are a few miscellaneous commands. 

 
\begin{description}
\item[{\textsc{Ctrl\texttt{\symbol{45}}V}}]  enter next character literally, i.e., enter it even if it is one of the
control keys. 
\item[{\textsc{Ctrl\texttt{\symbol{45}}U}}]  execute the next line editing command 4 times. 
\item[{\textsc{Esc\texttt{\symbol{45}}}\texttt{num}}]  execute the next line editing command \texttt{num} times. 
\item[{\textsc{Esc\texttt{\symbol{45}}Ctrl\texttt{\symbol{45}}L}}]  redisplay input line. 
\end{description}
 

 The four arrow keys (cursor keys) can be used instead of \textsc{Ctrl\texttt{\symbol{45}}B}, \textsc{Ctrl\texttt{\symbol{45}}F}, \textsc{Ctrl\texttt{\symbol{45}}P}, and \textsc{Ctrl\texttt{\symbol{45}}N}, respectively. }

 
\section{\textcolor{Chapter }{Editing using the \texttt{readline} library}}\label{sec:readline}
\logpage{[ 6, 9, 0 ]}
\hyperdef{L}{X7AD8D65F7BA1C3E0}{}
{
  The descriptions in this section are valid only if your \textsf{GAP} installation uses the \texttt{readline} library for command line editing. You can check by \texttt{IsBound(GAPInfo.UseReadline);} if this is the case. 

 You can use all the features of \texttt{readline}, as for example explained in \href{https://tiswww.case.edu/php/chet/readline/rluserman.html} {\texttt{https://tiswww.case.edu/php/chet/readline/rluserman.html}}. Therefore the command line editing in \textsf{GAP} is similar to the \texttt{bash} shell and many other programs. On a Unix/Linux system you may also have a
manpage, try \texttt{man readline}. 

 Compared to the command line editing which was used in \textsf{GAP} up to version{\nobreakspace}4.4 (or compared to not using the \texttt{readline} library) using \texttt{readline} has several advantages: 
\begin{itemize}
\item Most keys still do the same as explained in \ref{Line Editing} (in the default configuration). 
\item There are many additional commands, e.g. undoing (\textsc{Ctrl\texttt{\symbol{45}}{\textunderscore}}, keyboard macros (\textsc{Ctrl\texttt{\symbol{45}}x(}, \textsc{Ctrl\texttt{\symbol{45}}x)} and \textsc{Ctrl\texttt{\symbol{45}}xe}), file name completion (hit \textsc{Esc} two or four times), showing matching parentheses, \texttt{vi}\texttt{\symbol{45}}style key bindings, deleting and yanking text, ...
\item Lines which are longer than a physical terminal row can be edited more
conveniently.
\item Arbitrary unicode characters can be typed into string literals. 
\item The key bindings can be configured, either via your \texttt{\texttt{\symbol{126}}/.inputrc} file or by \textsf{GAP} commands, see \ref{ssec:readlineCustom}.
\item The command line history can be saved to and read from a file, see \ref{ssec:cmdlinehistory}.
\item Adventurous users can even implement completely new command line editing
functions on \textsf{GAP} level, see \ref{ssec:readlineUserFuncs}.
\end{itemize}
 

 
\subsection{\textcolor{Chapter }{Readline customization}}\label{ssec:readlineCustom}
\logpage{[ 6, 9, 1 ]}
\hyperdef{L}{X7C38F9E0783D9442}{}
{
 \index{ReadlineInitLine@\texttt{ReadlineInitLine}}  You can use your readline init file (by default \texttt{\texttt{\symbol{126}}/.inputrc} on Unix/Linux) to customize key bindings. If you want settings be used only
within \textsf{GAP} you can write them between lines containing \texttt{\$if GAP} and \texttt{\$endif}. For a detailed documentation of the available settings and functions see \href{ https://tiswww.case.edu/php/chet/readline/rluserman.html} {here}. 
\begin{Verbatim}[commandchars=!@|,fontsize=\small,frame=single,label=From readline init file]
  $if GAP
    set blink-matching-paren on
    "\C-x\C-o": dump-functions
    "\ep": kill-region
  $endif
\end{Verbatim}
 Alternatively, from within \textsf{GAP} the command \texttt{ReadlineInitLine(\mbox{\texttt{\mdseries\slshape line}});} can be used, where \mbox{\texttt{\mdseries\slshape line}} is a string containing a line as in the init file. 

 Caveat: \textsf{GAP} overwrites the following keys (after reading the \texttt{\texttt{\symbol{126}}/.inputrc} file): \texttt{\texttt{\symbol{92}}C\texttt{\symbol{45}}g}, \texttt{\texttt{\symbol{92}}C\texttt{\symbol{45}}i}, \texttt{\texttt{\symbol{92}}C\texttt{\symbol{45}}n}, \texttt{\texttt{\symbol{92}}C\texttt{\symbol{45}}o}, \texttt{\texttt{\symbol{92}}C\texttt{\symbol{45}}p}, \texttt{\texttt{\symbol{92}}C\texttt{\symbol{45}}r}, \texttt{\texttt{\symbol{92}}C\texttt{\symbol{45}}\texttt{\symbol{92}}}, \texttt{\texttt{\symbol{92}}e{\textless}}, \texttt{\texttt{\symbol{92}}e{\textgreater}}, \texttt{Up}, \texttt{Down}, \texttt{TAB}, \texttt{Space}, \texttt{PageUp}, \texttt{PageDown}. So, do not redefine these in your \texttt{\texttt{\symbol{126}}/.inputrc}. 

 Note that after pressing \textsc{Ctrl\texttt{\symbol{45}}v} the next special character is input verbatim. This is very useful to bind keys
or key sequences. For example, binding the function key \textsc{F3} to the command \texttt{kill\texttt{\symbol{45}}whole\texttt{\symbol{45}}line} by using the sequence \textsc{Ctrl\texttt{\symbol{45}}v} \textsc{F3} looks on many terminals like this: \texttt{ReadlineInitLine("\texttt{\symbol{92}}"\texttt{\symbol{94}}[OR\texttt{\symbol{92}}":kill\texttt{\symbol{45}}whole\texttt{\symbol{45}}line");}. (You can get the line back later with \textsc{Ctrl\texttt{\symbol{45}}y}.) 

 The \textsc{Ctrl\texttt{\symbol{45}}g} key can be used to type any unicode character by its code point. The number of
the character can either be given as a count, or if the count is one the input
characters before the cursor are taken (as decimal number or as hex number
which starts with \texttt{0x}. For example, the double stroke character {\ensuremath{\mathbb Z}} can be
input by any of the three key sequences \textsc{Esc 8484 Ctrl\texttt{\symbol{45}}g}, \textsc{8484 Ctrl\texttt{\symbol{45}}g} or \textsc{0x2124 Ctrl\texttt{\symbol{45}}g}. 

 Some terminals bind the \textsc{Ctrl\texttt{\symbol{45}}s} and \textsc{Ctrl\texttt{\symbol{45}}q} keys to stop and restart terminal output. Furthermore, sometimes \textsc{Ctrl\texttt{\symbol{45}}\texttt{\symbol{92}}} quits a program. To disable this behaviour (and maybe use these keys for
command line editing) you can use \texttt{Exec("stty stop undef; stty start undef; stty quit undef");} in your \textsf{GAP} session or your \texttt{gaprc} file (see \ref{sect:gap.ini}). 

 }

 
\subsection{\textcolor{Chapter }{The command line history}}\label{ssec:cmdlinehistory}
\logpage{[ 6, 9, 2 ]}
\hyperdef{L}{X846C3DED84AD7593}{}
{
  \textsf{GAP} can save your input lines for later reuse. The keys \textsc{Ctrl\texttt{\symbol{45}}p} (or \textsc{Up}), \textsc{Ctrl\texttt{\symbol{45}}n} (or \textsc{Down}), \textsc{ESC{\textless}} and \textsc{ESC{\textgreater}} work as documented in \ref{Line Editing}, that is they scroll backward and forward in the history or go to its
beginning or end. Also, \textsc{Ctrl\texttt{\symbol{45}}o} works as documented, it is useful for repeating a sequence of previous lines.
(But \textsc{Ctrl\texttt{\symbol{45}}l} clears the screen as in other programs.) 

 The command line history can be used across several instances of \textsf{GAP} via the following two commands. }

 

\subsection{\textcolor{Chapter }{SaveCommandLineHistory}}
\logpage{[ 6, 9, 3 ]}\nobreak
\hyperdef{L}{X7C1F4D04861C1197}{}
{\noindent\textcolor{FuncColor}{$\triangleright$\enspace\texttt{SaveCommandLineHistory({\mdseries\slshape [fname][,] [app]})\index{SaveCommandLineHistory@\texttt{SaveCommandLineHistory}}
\label{SaveCommandLineHistory}
}\hfill{\scriptsize (function)}}\\
\textbf{\indent Returns:\ }
\texttt{fail} or number of saved lines

\noindent\textcolor{FuncColor}{$\triangleright$\enspace\texttt{ReadCommandLineHistory({\mdseries\slshape [fname][,] [app]})\index{ReadCommandLineHistory@\texttt{ReadCommandLineHistory}}
\label{ReadCommandLineHistory}
}\hfill{\scriptsize (function)}}\\
\textbf{\indent Returns:\ }
\texttt{fail} or number of added lines



 The first command saves the lines in the command line history to the file
given by the string \mbox{\texttt{\mdseries\slshape fname}}. The default for \mbox{\texttt{\mdseries\slshape fname}} is \texttt{history} in the user's \textsf{GAP} root path \texttt{GAPInfo.UserGapRoot} or \texttt{"\texttt{\symbol{126}}/.gap{\textunderscore}hist"} if this directory does not exist. If the optional argument \mbox{\texttt{\mdseries\slshape app}} is \texttt{true} then the lines are appended to that file otherwise the file is overwritten. 

 The second command is the converse, it reads the lines from file \mbox{\texttt{\mdseries\slshape fname}}. If the optional argument \mbox{\texttt{\mdseries\slshape app}} is true the lines are appended to the history, else it \emph{prepends} them. 

 By default, the command line history stores up to 1000 input lines. command
line history. This number may be restricted or enlarged via via \texttt{SetUserPreference("HistoryMaxLines", num);} which may be set to a non negative number \texttt{num} to store up to \texttt{num} input lines or to \texttt{infinity} to store arbitrarily many lines. An automatic storing and restoring of the
command line history can be configured via \texttt{SetUserPreference("SaveAndRestoreHistory", true);}. 

 Note that these functions are only available if your \textsf{GAP} is configured to use the \texttt{readline} library. }

 
\subsection{\textcolor{Chapter }{Writing your own command line editing functions}}\label{ssec:readlineUserFuncs}
\logpage{[ 6, 9, 4 ]}
\hyperdef{L}{X87D4EA197A263FB7}{}
{
 \index{InstallReadlineMacro@\texttt{InstallReadlineMacro}} \index{InvocationReadlineMacro@\texttt{InvocationReadlineMacro}}  It is possible to write new command line editing functions in \textsf{GAP} as follows. 

 The functions have one argument \mbox{\texttt{\mdseries\slshape l}} which is a list with five entries of the form \texttt{[count, key, line, cursorpos, markpos]} where \texttt{count} and \texttt{key} are the last pressed key and its count (these are not so useful here because
users probably do not want to overwrite the binding of a single key), then \texttt{line} is a string containing the line typed so far, \texttt{cursorpos} is the current position of the cursor (point), and \texttt{markpos} the current position of the mark. 

 The result of such a function must be a list which can have various forms: 
\begin{description}
\item[{\texttt{[str]}}] with a string \texttt{str}. In this case the text \texttt{str} is inserted at the cursor position.
\item[{\texttt{[kill, begin, end]}}]  where \texttt{kill} is \texttt{true} or \texttt{false} and \texttt{begin} and \texttt{end} are positions on the input line. This removes the text from the lower position
to before the higher position. If \texttt{kill} is \texttt{true} the text is killed, i.e. put in the kill ring for later yanking. 
\item[{\texttt{[begin, end, str]}}] where \texttt{begin} and \texttt{end} are positions on the input line and \texttt{str} is a string. Then the text from position \texttt{begin} to before \texttt{end} is substituted by \texttt{str}. 
\item[{\texttt{[1, lstr]}}]  where \texttt{lstr} is a list of strings. Then these strings are displayed like a list of possible
completions. The input line is not changed. 
\item[{\texttt{[2, chars]}}] where \texttt{chars} is a string. The characters from \texttt{chars} are used as the next characters from the input. (At most 512 characters are
possible.)
\item[{\texttt{[100]}}] This rings the bell as configured in the terminal.
\end{description}
 In the first three cases the result list can contain a position as a further
entry, this becomes the new cursor position. Or it can contain two positions
as further entries, these become the new cursor position and the new position
of the mark. 

 Such a function can be installed as a macro for \texttt{readline} via \texttt{InstallReadlineMacro(name, fun);} where \texttt{name} is a string used as name of the macro and \texttt{fun} is a function as above. This macro can be called by a key sequence which is
returned by \texttt{InvocationReadlineMacro(name);}. 

 As an example we define a function which puts double quotes around the word
under or before the cursor position. The space character, the characters in \texttt{"(,)"}, and the beginning and end of the line are considered as word boundaries. The
function is then installed as a macro and bound to the key sequence \textsc{Esc} \textsc{Q}. 

 
\begin{Verbatim}[commandchars=!@|,fontsize=\small,frame=single,label=Example]
  !gapprompt@gap>| !gapinput@EditAddQuotes := function(l)|
  !gapprompt@>| !gapinput@  local str, pos, i, j, new;|
  !gapprompt@>| !gapinput@  str := l[3];|
  !gapprompt@>| !gapinput@  pos := l[4];|
  !gapprompt@>| !gapinput@  i := pos;|
  !gapprompt@>| !gapinput@  while i > 1 and (not str[i-1] in ",( ") do|
  !gapprompt@>| !gapinput@    i := i-1;|
  !gapprompt@>| !gapinput@  od;|
  !gapprompt@>| !gapinput@  j := pos;|
  !gapprompt@>| !gapinput@  while IsBound(str[j]) and not str[j] in ",) " do|
  !gapprompt@>| !gapinput@    j := j+1;|
  !gapprompt@>| !gapinput@  od;|
  !gapprompt@>| !gapinput@  new := "\"";|
  !gapprompt@>| !gapinput@  Append(new, str{[i..j-1]});|
  !gapprompt@>| !gapinput@  Append(new, "\"");|
  !gapprompt@>| !gapinput@  return [i, j, new];|
  !gapprompt@>| !gapinput@end;;|
  !gapprompt@gap>| !gapinput@InstallReadlineMacro("addquotes", EditAddQuotes);|
  !gapprompt@gap>| !gapinput@invl := InvocationReadlineMacro("addquotes");;|
  !gapprompt@gap>| !gapinput@ReadlineInitLine(Concatenation("\"\\eQ\":\"",invl,"\""));;|
\end{Verbatim}
 }

 }

  
\section{\textcolor{Chapter }{Editing Files}}\label{Editing Files}
\logpage{[ 6, 10, 0 ]}
\hyperdef{L}{X7D8E1CF47E97A764}{}
{
  In most cases, it is preferable to create longer input (in particular \textsf{GAP} programs) separately in an editor, and to read in the result via \texttt{Read} (\ref{Read}). Note that \texttt{Read} (\ref{Read}) by default reads from the directory in which \textsf{GAP} was started (respectively under Windows the directory containing the \textsf{GAP} binary), so you might have to give an absolute path to the file. 

 If you cannot create several windows, the \texttt{Edit} (\ref{Edit}) command may be used to leave \textsf{GAP}, start an editor, and read in the edited file automatically. 

\subsection{\textcolor{Chapter }{Edit}}
\logpage{[ 6, 10, 1 ]}\nobreak
\hyperdef{L}{X82E5859C8113BA4D}{}
{\noindent\textcolor{FuncColor}{$\triangleright$\enspace\texttt{Edit({\mdseries\slshape filename})\index{Edit@\texttt{Edit}}
\label{Edit}
}\hfill{\scriptsize (function)}}\\


 \texttt{Edit} starts an editor with the file whose filename is given by the string \mbox{\texttt{\mdseries\slshape filename}}, and reads the file back into \textsf{GAP} when you exit the editor again. 

 \textsf{GAP} will call your preferred editor if you call \texttt{SetUserPreference("Editor", \mbox{\texttt{\mdseries\slshape path}});} where \mbox{\texttt{\mdseries\slshape path}} is the path to your editor, e.g., \texttt{/usr/bin/vim}. On Windows you can use \texttt{edit.com}. 

 Under macOS, you should use \texttt{SetUserPreference("Editor", "open");}, this will open the file in the default editor. If you call \texttt{SetUserPreference("EditorOptions", ["\texttt{\symbol{45}}t"]);}, the file will open in \texttt{TextEdit}, and \texttt{SetUserPreference("EditorOptions", ["\texttt{\symbol{45}}a",
"{\textless}appl{\textgreater}"]);} will open the file using the application \texttt{{\textless}appl{\textgreater}}. 

 This can for example be done in your \texttt{gap.ini} file, see Section \ref{subsect:gap.ini file}. }

 }

  
\section{\textcolor{Chapter }{Editor Support}}\label{Editor Support}
\logpage{[ 6, 11, 0 ]}
\hyperdef{L}{X7B67FF1E87FE67D1}{}
{
  \index{utilities for editing GAP files} \index{vi@\texttt{vi}} \index{vim@\texttt{vim}} \index{emacs@\texttt{emacs}} In the \texttt{etc} subdirectory of the \textsf{GAP} installation we provide some setup files for the editor \texttt{vim}. 

 \texttt{vim} is a powerful editor that understands the basic \texttt{vi} commands but provides much more functionality. You can find more information
about it (and download it) from \href{https://www.vim.org} {\texttt{https://www.vim.org}}. 

 To get support for \textsf{GAP} syntax in vim, create in your home directory a directory \texttt{.vim} with subdirectories \texttt{.vim/syntax} and \texttt{.vim/indent} (If you are not using Unix, refer to the \texttt{vim} documentation on where to place syntax files). Then copy the file \texttt{etc/vim/gap.vim} to \texttt{.vim/syntax/gap.vim} and the file \texttt{etc/vim/gap{\textunderscore}indent.vim} to \texttt{.vim/indent/gap.vim}. 

 Then edit the \texttt{.vimrc} file in your home directory. Add lines as in the following example: 

 
\begin{Verbatim}[commandchars=@|C,fontsize=\small,frame=single,label=Example]
  if has("syntax")
    syntax on             " Default to no syntax highlighting
  endif
  
  " For GAP files
  augroup gap
    " Remove all gap autocommands
    au!
  autocmd BufRead,BufNewFile *.g,*.gi,*.gd set filetype=gap comments=s:##\ \ ,m:##\ \ ,e:##\ \ b:#
  
  " I'm using the external program `par' for formatting comment lines starting
  " with `##  '. Include these lines only when you have par installed.
    autocmd BufRead,BufNewFile *.g,*.gi,*.gd set formatprg="par w76p4s0j"
    autocmd BufWritePost,FileWritePost *.g,*.gi,*.gd set formatprg="par w76p0s0j"
  augroup END
\end{Verbatim}
 

 See the headers of the two mentioned files for additional comments and adjust
details according to your personal taste. Send comments and suggestions to \href{mailto://support@gap-system.org} {\texttt{support@gap\texttt{\symbol{45}}system.org}}. 

 Users of \texttt{emacs}/\texttt{xemacs} may wish to take a look at the \href{https://melpa.org/#/gap-mode} {major-mode for editing GAP files} by Ivan Andrus. }

  
\section{\textcolor{Chapter }{Changing the Screen Size}}\label{sect:SizeScreen}
\logpage{[ 6, 12, 0 ]}
\hyperdef{L}{X83279E897ACCFFFA}{}
{
  

\subsection{\textcolor{Chapter }{SizeScreen}}
\logpage{[ 6, 12, 1 ]}\nobreak
\hyperdef{L}{X8723E0A1837894F3}{}
{\noindent\textcolor{FuncColor}{$\triangleright$\enspace\texttt{SizeScreen({\mdseries\slshape [sz]})\index{SizeScreen@\texttt{SizeScreen}}
\label{SizeScreen}
}\hfill{\scriptsize (function)}}\\


 Called with no arguments, \texttt{SizeScreen} returns the size of the screen as a list with two entries. The first is the
length of each line, the second is the number of lines. 

 Called with one argument that is a list \mbox{\texttt{\mdseries\slshape sz}}, \texttt{SizeScreen} sets the size of the screen; The first entry of \mbox{\texttt{\mdseries\slshape sz}}, if bound, is the length of each line, and the second entry of \mbox{\texttt{\mdseries\slshape sz}}, if bound, is the number of lines. The values for unbound entries of \mbox{\texttt{\mdseries\slshape sz}} are left unaffected. The function returns the new values. 

 Note that those parameters can also be set with the command line options \texttt{\texttt{\symbol{45}}x} for the line length and \texttt{\texttt{\symbol{45}}y} for the number of lines (see Section{\nobreakspace}\ref{Command Line Options}). 

 To check/change whether line breaking occurs for files and streams
see{\nobreakspace}\texttt{PrintFormattingStatus} (\ref{PrintFormattingStatus}) and{\nobreakspace}\texttt{SetPrintFormattingStatus} (\ref{SetPrintFormattingStatus}). 

 The line length must be between $20$ and $4096$ characters (inclusive) and the number of lines must be at least $10$. Values outside this range will be adjusted to the nearest endpoint of the
range. }

 }

  
\section{\textcolor{Chapter }{Teaching Mode}}\label{Teaching Mode}
\logpage{[ 6, 13, 0 ]}
\hyperdef{L}{X87847E5087D6F47D}{}
{
  When using \textsf{GAP} in the context of (undergraduate) teaching it is often desirable to simplify
some of the system output and functionality defaults (potentially at the cost
of making the printing of objects more expensive). This can be achieved by
turning on a teaching mode: 

\subsection{\textcolor{Chapter }{TeachingMode}}
\logpage{[ 6, 13, 1 ]}\nobreak
\hyperdef{L}{X7BE2515F82425404}{}
{\noindent\textcolor{FuncColor}{$\triangleright$\enspace\texttt{TeachingMode({\mdseries\slshape [switch]})\index{TeachingMode@\texttt{TeachingMode}}
\label{TeachingMode}
}\hfill{\scriptsize (function)}}\\


 When called with a boolean argument \mbox{\texttt{\mdseries\slshape switch}}, this function will turn teaching mode respectively on or off. 
\begin{Verbatim}[commandchars=!@|,fontsize=\small,frame=single,label=Example]
  !gapprompt@gap>| !gapinput@a:=Z(11)^3;|
  Z(11)^3
  !gapprompt@gap>| !gapinput@TeachingMode(true);|
  #I  Teaching mode is turned ON
  !gapprompt@gap>| !gapinput@a;|
  ZmodnZObj(8,11)
  !gapprompt@gap>| !gapinput@TeachingMode(false);|
  #I  Teaching mode is turned OFF
  !gapprompt@gap>| !gapinput@a;|
  Z(11)^3
\end{Verbatim}
 At the moment, teaching mode changes the following things 

 
\begin{description}
\item[{Prime Field Elements}]  Elements of fields of prime order are printed as \texttt{ZmodnZObj} (\ref{ZmodnZObj:for two integers}) instead as power of a primitive root. 
\item[{Quadratic Irrationalities}]  Elements of a quadratic extension of the rationals are printed using the
square root \texttt{ER} (\ref{ER}) instead of using roots of unity. 
\item[{Creation of some small groups}]  The group creator functions \texttt{CyclicGroup} (\ref{CyclicGroup}), \texttt{AbelianGroup} (\ref{AbelianGroup}), \texttt{ElementaryAbelianGroup} (\ref{ElementaryAbelianGroup}), and \texttt{DihedralGroup} (\ref{DihedralGroup}) create by default (if no other representation is specified) not a pc group,
but a finitely presented group, which makes the generators easier to
interpret. 
\end{description}
 }

 }

 }

            
\chapter{\textcolor{Chapter }{Debugging and Profiling Facilities}}\label{Debugging and Profiling Facilities}
\logpage{[ 7, 0, 0 ]}
\hyperdef{L}{X8345F6817DFD6394}{}
{
  This chapter describes some functions that are useful mainly for debugging and
profiling purposes. 

 Probably the most important debugging tool in \textsf{GAP} is the break loop (see Section \ref{Break Loops}) which can be entered by putting an \texttt{Error} (\ref{Error}) statement into your code or by hitting Control\texttt{\symbol{45}}C. In the
break loop one can inspect variables, stack traces and issue commands as usual
in an interactive \textsf{GAP} session. See also the \texttt{DownEnv} (\ref{DownEnv}), \texttt{UpEnv} (\ref{UpEnv}), \texttt{Where} (\ref{Where}) and \texttt{WhereWithVars} (\ref{WhereWithVars}) functions. 

 Sections{\nobreakspace}\ref{sect:ApplicableMethod} and{\nobreakspace}\ref{Tracing Methods} show how to get information about the methods chosen by the method selection
mechanism (see chapter{\nobreakspace}\ref{Method Selection}). 

 The final sections describe functions for collecting statistics about
computations (see \texttt{Runtime} (\ref{Runtime}), \ref{Profiling}).  
\section{\textcolor{Chapter }{Recovery from NoMethodFound\texttt{\symbol{45}}Errors}}\label{Recovery from NoMethodFound-Errors}
\logpage{[ 7, 1, 0 ]}
\hyperdef{L}{X83C45B0A797AAF96}{}
{
  When the method selection fails because there is no applicable method, an
error as in the following example occurs and a break loop is entered: 

 
\begin{Verbatim}[commandchars=@|A,fontsize=\small,frame=single,label=Example]
  @gapprompt|gap>A @gapinput|IsNormal(2,2);A
  Error, no method found! For debugging hints type ?Recovery from NoMethodFound
  Error, no 1st choice method found for `IsNormal' on 2 arguments
  called from read-eval loop at *stdin*:1
  you can enter 'quit;' to quit to outer loop
  @gapbrkprompt|brk>A @gapinput|A
\end{Verbatim}
 

 This only says, that the method selection tried to find a method for \texttt{IsNormal} on two arguments and failed. In this situation it is crucial to find out, why
this happened. Therefore there are a few functions which can display further
information. Note that you can leave the break loop by the \texttt{quit} command (see{\nobreakspace}\ref{quit}) and that the information about the incident is no longer accessible
afterwards. 

   

\subsection{\textcolor{Chapter }{ShowArguments}}
\logpage{[ 7, 1, 1 ]}\nobreak
\hyperdef{L}{X86B5FEC67A9394DC}{}
{\noindent\textcolor{FuncColor}{$\triangleright$\enspace\texttt{ShowArguments({\mdseries\slshape })\index{ShowArguments@\texttt{ShowArguments}}
\label{ShowArguments}
}\hfill{\scriptsize (function)}}\\


 This function is only available within a break loop caused by a ``No Method Found''\texttt{\symbol{45}}error. It prints as a list the arguments of the operation
call for which no method was found.   }

 

\subsection{\textcolor{Chapter }{ShowArgument}}
\logpage{[ 7, 1, 2 ]}\nobreak
\hyperdef{L}{X834BD9928773DCC1}{}
{\noindent\textcolor{FuncColor}{$\triangleright$\enspace\texttt{ShowArgument({\mdseries\slshape nr})\index{ShowArgument@\texttt{ShowArgument}}
\label{ShowArgument}
}\hfill{\scriptsize (function)}}\\


 This function is only available within a break loop caused by a ``No Method Found''\texttt{\symbol{45}}error. It prints the \mbox{\texttt{\mdseries\slshape nr}}\texttt{\symbol{45}}th arguments of the operation call for which no method was
found. \texttt{ShowArgument} needs exactly one argument which is an integer between 0 and the number of
arguments the operation was called with. }

 

\subsection{\textcolor{Chapter }{ShowDetails}}
\logpage{[ 7, 1, 3 ]}\nobreak
\hyperdef{L}{X7D25D904800D5CBA}{}
{\noindent\textcolor{FuncColor}{$\triangleright$\enspace\texttt{ShowDetails({\mdseries\slshape })\index{ShowDetails@\texttt{ShowDetails}}
\label{ShowDetails}
}\hfill{\scriptsize (function)}}\\


 This function is only available within a break loop caused by a ``No Method Found''\texttt{\symbol{45}}error. It prints the details of this error: The operation,
the number of arguments, a flag which indicates whether the operation is being
traced, a flag which indicates whether the operation is a constructor method,
and the number of methods that refused to apply by calling \texttt{TryNextMethod} (\ref{TryNextMethod}). The last number is called \texttt{Choice} and is printed as an ordinal. So if exactly $k$ methods were found but called \texttt{TryNextMethod} (\ref{TryNextMethod}) and there were no more methods it says \texttt{Choice: }$k$\texttt{th}. }

 

\subsection{\textcolor{Chapter }{ShowMethods}}
\logpage{[ 7, 1, 4 ]}\nobreak
\hyperdef{L}{X7F6996CA872478B8}{}
{\noindent\textcolor{FuncColor}{$\triangleright$\enspace\texttt{ShowMethods({\mdseries\slshape [verbosity]})\index{ShowMethods@\texttt{ShowMethods}}
\label{ShowMethods}
}\hfill{\scriptsize (function)}}\\


 This function is only available within a break loop caused by a ``No Method Found''\texttt{\symbol{45}}error. It prints an overview about the installed methods
for those arguments the operation was called with (using \ref{sect:ApplicableMethod}. The verbosity can be controlled by the optional integer parameter \mbox{\texttt{\mdseries\slshape verbosity}}. The default is 2, which lists all applicable methods. With verbosity 1 \texttt{ShowMethods} only shows the number of installed methods and the methods matching, which can
only be those that were already called but refused to work by calling \texttt{TryNextMethod} (\ref{TryNextMethod}). With verbosity 3 not only all installed methods but also the reasons why
they do not match are displayed. }

 

\subsection{\textcolor{Chapter }{ShowOtherMethods}}
\logpage{[ 7, 1, 5 ]}\nobreak
\hyperdef{L}{X7E5E2E7B85029E34}{}
{\noindent\textcolor{FuncColor}{$\triangleright$\enspace\texttt{ShowOtherMethods({\mdseries\slshape [verbosity]})\index{ShowOtherMethods@\texttt{ShowOtherMethods}}
\label{ShowOtherMethods}
}\hfill{\scriptsize (function)}}\\


 This function is only available within a break loop caused by a ``No Method Found''\texttt{\symbol{45}}error. It prints an overview about the installed methods
for a different number of arguments than the number of arguments the operation
was called with (using \ref{sect:ApplicableMethod}. The verbosity can be controlled by the optional integer parameter \mbox{\texttt{\mdseries\slshape verbosity}}. The default is 1 which lists only the number of applicable methods. With
verbosity 2 \texttt{ShowOtherMethods} lists all installed methods and with verbosity 3 also the reasons, why they
are not applicable. Calling \texttt{ShowOtherMethods} with verbosity 3 in this function will normally not make any sense, because
the different numbers of arguments are simulated by supplying the
corresponding number of ones, for which normally no reasonable methods will be
installed. }

 }

  
\section{\textcolor{Chapter }{Inspecting Applicable Methods}}\label{sect:ApplicableMethod}
\logpage{[ 7, 2, 0 ]}
\hyperdef{L}{X7FDA1D4B87BD25A8}{}
{
  

\subsection{\textcolor{Chapter }{ApplicableMethod}}
\logpage{[ 7, 2, 1 ]}\nobreak
\hyperdef{L}{X80848FF486BD6F9F}{}
{\noindent\textcolor{FuncColor}{$\triangleright$\enspace\texttt{ApplicableMethod({\mdseries\slshape opr, args[, printlevel[, nr]]})\index{ApplicableMethod@\texttt{ApplicableMethod}}
\label{ApplicableMethod}
}\hfill{\scriptsize (function)}}\\
\noindent\textcolor{FuncColor}{$\triangleright$\enspace\texttt{ApplicableMethodTypes({\mdseries\slshape opr, args[, printlevel[, nr]]})\index{ApplicableMethodTypes@\texttt{ApplicableMethodTypes}}
\label{ApplicableMethodTypes}
}\hfill{\scriptsize (function)}}\\


 Called with two arguments, \texttt{ApplicableMethod} returns the method of highest rank that is applicable for the operation \mbox{\texttt{\mdseries\slshape opr}} with the arguments in the list \mbox{\texttt{\mdseries\slshape args}}. The default \mbox{\texttt{\mdseries\slshape printlevel}} is \texttt{0}. If no method is applicable then \texttt{fail} is returned. 

 If a positive integer is given as the fourth argument \mbox{\texttt{\mdseries\slshape nr}} then \texttt{ApplicableMethod} returns the \mbox{\texttt{\mdseries\slshape nr}}\texttt{\symbol{45}}th applicable method for the operation \mbox{\texttt{\mdseries\slshape opr}} with the arguments in the list \mbox{\texttt{\mdseries\slshape args}}, where the methods are ordered according to descending rank. If less than \mbox{\texttt{\mdseries\slshape nr}} methods are applicable then \texttt{fail} is returned. 

 If the fourth argument \mbox{\texttt{\mdseries\slshape nr}} is the string \texttt{"all"} then \texttt{ApplicableMethod} returns a list of all applicable methods for \mbox{\texttt{\mdseries\slshape opr}} with arguments \mbox{\texttt{\mdseries\slshape args}}, ordered according to descending rank. 

 Depending on the integer value \mbox{\texttt{\mdseries\slshape printlevel}}, additional information is printed. Admissible values and their meaning are
as follows. 

 
\begin{description}
\item[{0}]  no information, 
\item[{1}]  information about the applicable method, 
\item[{2}]  also information about the not applicable methods of higher rank, 
\item[{3}]  also for each not applicable method the first reason why it is not applicable, 
\item[{4}]  also for each not applicable method all reasons why it is not applicable. 
\item[{6}]  also the function body of the selected method(s) 
\end{description}
 

 When a method returned by \texttt{ApplicableMethod} is called then it returns either the desired result or the string \texttt{"TRY{\textunderscore}NEXT{\textunderscore}METHOD"}, which corresponds to a call to \texttt{TryNextMethod} (\ref{TryNextMethod}) in the method and means that the method selection would call the next
applicable method. 

 \emph{Note:} The \textsf{GAP} kernel provides special treatment for the infix operations \texttt{\texttt{\symbol{92}}+}, \texttt{\texttt{\symbol{92}}\texttt{\symbol{45}}}, \texttt{\texttt{\symbol{92}}*}, \texttt{\texttt{\symbol{92}}/}, \texttt{\texttt{\symbol{92}}\texttt{\symbol{94}}}, \texttt{\texttt{\symbol{92}}mod} and \texttt{\texttt{\symbol{92}}in}. For some kernel objects (notably cyclotomic numbers, finite field elements
and row vectors thereof) it calls kernel methods circumventing the method
selection mechanism. Therefore for these operations \texttt{ApplicableMethod} may return a method which is not the kernel method actually used. 

 The function \texttt{ApplicableMethodTypes} takes the \emph{types} or \emph{filters} of the arguments as argument (if only filters are given of course family
predicates cannot be tested). }

 

\subsection{\textcolor{Chapter }{MethodsOperation}}
\logpage{[ 7, 2, 2 ]}\nobreak
\hyperdef{L}{X7D88B5FE839FFD02}{}
{\noindent\textcolor{FuncColor}{$\triangleright$\enspace\texttt{MethodsOperation({\mdseries\slshape oper, n})\index{MethodsOperation@\texttt{MethodsOperation}}
\label{MethodsOperation}
}\hfill{\scriptsize (function)}}\\


 Return a list of records that describe the installed methods for the operation \mbox{\texttt{\mdseries\slshape oper}}, with \mbox{\texttt{\mdseries\slshape n}} arguments. 

 Each record has at least the following components. 

 
\begin{description}
\item[{func}] the method itself, a \textsf{GAP} function,
\item[{rank}] the rank of the method, an integer or infinity,
\item[{info}] a string that describes when the method is applicable,
\item[{location}] filename and start line where the method gets installed.
\end{description}
 }

 }

  
\section{\textcolor{Chapter }{Tracing Methods}}\label{Tracing Methods}
\logpage{[ 7, 3, 0 ]}
\hyperdef{L}{X7D43A2D885B37739}{}
{
  

\subsection{\textcolor{Chapter }{TraceMethods (for operations)}}
\logpage{[ 7, 3, 1 ]}\nobreak
\hyperdef{L}{X80B044017C9E4137}{}
{\noindent\textcolor{FuncColor}{$\triangleright$\enspace\texttt{TraceMethods({\mdseries\slshape opr1, opr2, ...})\index{TraceMethods@\texttt{TraceMethods}!for operations}
\label{TraceMethods:for operations}
}\hfill{\scriptsize (function)}}\\
\noindent\textcolor{FuncColor}{$\triangleright$\enspace\texttt{TraceMethods({\mdseries\slshape oprs})\index{TraceMethods@\texttt{TraceMethods}!for a list of operations}
\label{TraceMethods:for a list of operations}
}\hfill{\scriptsize (function)}}\\


 After the call of \texttt{TraceMethods}, whenever a method of one of the operations \mbox{\texttt{\mdseries\slshape opr1}}, \mbox{\texttt{\mdseries\slshape opr2}}, ... is called, the information string used in the installation of the method
is printed. The second form has the same effect for each operation from the
list \mbox{\texttt{\mdseries\slshape oprs}} of operations. }

 

\subsection{\textcolor{Chapter }{TraceAllMethods}}
\logpage{[ 7, 3, 2 ]}\nobreak
\hyperdef{L}{X7D34CADB813A4AF1}{}
{\noindent\textcolor{FuncColor}{$\triangleright$\enspace\texttt{TraceAllMethods({\mdseries\slshape })\index{TraceAllMethods@\texttt{TraceAllMethods}}
\label{TraceAllMethods}
}\hfill{\scriptsize (function)}}\\


 Invokes \texttt{TraceMethods} for all operations. }

 

\subsection{\textcolor{Chapter }{UntraceMethods (for operations)}}
\logpage{[ 7, 3, 3 ]}\nobreak
\hyperdef{L}{X7EB04D387C53E4C1}{}
{\noindent\textcolor{FuncColor}{$\triangleright$\enspace\texttt{UntraceMethods({\mdseries\slshape opr1, opr2, ...})\index{UntraceMethods@\texttt{UntraceMethods}!for operations}
\label{UntraceMethods:for operations}
}\hfill{\scriptsize (function)}}\\
\noindent\textcolor{FuncColor}{$\triangleright$\enspace\texttt{UntraceMethods({\mdseries\slshape oprs})\index{UntraceMethods@\texttt{UntraceMethods}!for a list of operations}
\label{UntraceMethods:for a list of operations}
}\hfill{\scriptsize (function)}}\\


 turns the tracing off for all operations \mbox{\texttt{\mdseries\slshape opr1}}, \mbox{\texttt{\mdseries\slshape opr2}}, ... or in the second form, for all operations in the list \mbox{\texttt{\mdseries\slshape oprs}}. 
\begin{Verbatim}[commandchars=!@|,fontsize=\small,frame=single,label=Example]
  !gapprompt@gap>| !gapinput@TraceMethods( [ Size ] );|
  !gapprompt@gap>| !gapinput@g:= Group( (1,2,3), (1,2) );;|
  !gapprompt@gap>| !gapinput@Size( g );|
  #I  Size: for a permutation group at /gap5/lib/grpperm.gi:487
  #I  Setter(Size): system setter
  #I  Size: system getter
  #I  Size: system getter
  6
  !gapprompt@gap>| !gapinput@UntraceMethods( [ Size ] );|
\end{Verbatim}
 }

 

\subsection{\textcolor{Chapter }{UntraceAllMethods}}
\logpage{[ 7, 3, 4 ]}\nobreak
\hyperdef{L}{X7B3018AA82D55949}{}
{\noindent\textcolor{FuncColor}{$\triangleright$\enspace\texttt{UntraceAllMethods({\mdseries\slshape })\index{UntraceAllMethods@\texttt{UntraceAllMethods}}
\label{UntraceAllMethods}
}\hfill{\scriptsize (function)}}\\


 Equivalent to calling \texttt{UntraceMethods} for all operations. }

 

\subsection{\textcolor{Chapter }{TraceImmediateMethods}}
\logpage{[ 7, 3, 5 ]}\nobreak
\hyperdef{L}{X81078D3387A38E31}{}
{\noindent\textcolor{FuncColor}{$\triangleright$\enspace\texttt{TraceImmediateMethods({\mdseries\slshape [flag]})\index{TraceImmediateMethods@\texttt{TraceImmediateMethods}}
\label{TraceImmediateMethods}
}\hfill{\scriptsize (function)}}\\
\noindent\textcolor{FuncColor}{$\triangleright$\enspace\texttt{UntraceImmediateMethods({\mdseries\slshape })\index{UntraceImmediateMethods@\texttt{UntraceImmediateMethods}}
\label{UntraceImmediateMethods}
}\hfill{\scriptsize (function)}}\\


 \texttt{TraceImmediateMethods} enables tracing for all immediate methods if \mbox{\texttt{\mdseries\slshape flag}} is either \texttt{true}, or not present. \texttt{UntraceImmediateMethods}, or \texttt{TraceImmediateMethods} with \mbox{\texttt{\mdseries\slshape flag}} equal \texttt{false} turns tracing off. (There is no facility to trace \emph{specific} immediate methods.) 
\begin{Verbatim}[commandchars=!@|,fontsize=\small,frame=single,label=Example]
  !gapprompt@gap>| !gapinput@TraceImmediateMethods( );|
  !gapprompt@gap>| !gapinput@g:= Group( (1,2,3), (1,2) );;|
  #I RunImmediateMethods
  #I  immediate: Size
  #I  immediate: IsCyclic
  #I  immediate: IsCommutative
  #I  immediate: IsTrivial
  !gapprompt@gap>| !gapinput@Size( g );|
  #I  immediate: IsPerfectGroup
  #I  immediate: IsNonTrivial
  #I  immediate: Size
  #I  immediate: IsFreeAbelian
  #I  immediate: IsTorsionFree
  #I  immediate: IsNonTrivial
  #I  immediate: IsPerfectGroup
  #I  immediate: GeneralizedPcgs
  #I  immediate: IsEmpty
  6
  !gapprompt@gap>| !gapinput@UntraceImmediateMethods( );|
  !gapprompt@gap>| !gapinput@UntraceMethods( [ Size ] );|
\end{Verbatim}
 

 This example gives an explanation for the two calls of the ``system getter'' for \texttt{Size} (\ref{Size}). Namely, there are immediate methods that access the known size of the group.
Note that the group \texttt{g} was known to be finitely generated already before the size was computed, the
calls of the immediate method for \texttt{IsFinitelyGeneratedGroup} (\ref{IsFinitelyGeneratedGroup}) after the call of \texttt{Size} (\ref{Size}) have other arguments than \texttt{g}. }

 

\subsection{\textcolor{Chapter }{TraceInternalMethods}}
\logpage{[ 7, 3, 6 ]}\nobreak
\hyperdef{L}{X81B000CF86BA1534}{}
{\noindent\textcolor{FuncColor}{$\triangleright$\enspace\texttt{TraceInternalMethods({\mdseries\slshape })\index{TraceInternalMethods@\texttt{TraceInternalMethods}}
\label{TraceInternalMethods}
}\hfill{\scriptsize (function)}}\\
\noindent\textcolor{FuncColor}{$\triangleright$\enspace\texttt{UntraceInternalMethods({\mdseries\slshape })\index{UntraceInternalMethods@\texttt{UntraceInternalMethods}}
\label{UntraceInternalMethods}
}\hfill{\scriptsize (function)}}\\
\noindent\textcolor{FuncColor}{$\triangleright$\enspace\texttt{GetTraceInternalMethodsCounts({\mdseries\slshape })\index{GetTraceInternalMethodsCounts@\texttt{GetTraceInternalMethodsCounts}}
\label{GetTraceInternalMethodsCounts}
}\hfill{\scriptsize (function)}}\\
\noindent\textcolor{FuncColor}{$\triangleright$\enspace\texttt{ClearTraceInternalMethodsCounts({\mdseries\slshape })\index{ClearTraceInternalMethodsCounts@\texttt{ClearTraceInternalMethodsCounts}}
\label{ClearTraceInternalMethodsCounts}
}\hfill{\scriptsize (function)}}\\


 \texttt{TraceInternalMethods} enables tracing for all internal methods. Internal methods are methods which
implement many fundamental operations in GAP. In this version of GAP, the
internal methods which can be traced are: 
\begin{description}
\item[{Zero, ZeroMut}] Mutable and Immutable \texttt{Zero} (\ref{Zero})
\item[{AInv, AInvMut}] Mutable and Immutable \texttt{AdditiveInverse} (\ref{AdditiveInverse})
\item[{One, OneMut}] Mutable and Immutable \texttt{One} (\ref{One})
\item[{Inv, InvMut}] Mutable and Immutable \texttt{Inverse} (\ref{Inverse})
\item[{Sum}] The operator \texttt{\texttt{\symbol{92}}+} (\ref{bSlash+})
\item[{Diff}] The operator \texttt{\texttt{\symbol{45}}} operator
\item[{Prod}] The operator \texttt{\texttt{\symbol{92}}*} (\ref{bSlash*})
\item[{Quo}] The operator \texttt{\texttt{\symbol{92}}/} (\ref{bSlash/})
\item[{LQuo}] The left\texttt{\symbol{45}}quotient operator
\item[{Pow}] The operator \texttt{\texttt{\symbol{92}}\texttt{\symbol{94}}} (\ref{bSlash^})
\item[{Comm}] The operator \texttt{Comm} (\ref{Comm})
\item[{Mod}] The operator \texttt{\texttt{\symbol{92}}mod} (\ref{bSlashmod})
\end{description}
 

 \texttt{UntraceInternalMethods} turns tracing off. As these methods can be called hundreds of thousands of
times in simple GAP code, there isn't a statement printed each time one is
called. Instead, the method \texttt{GetTraceInternalMethodsCounts} returns how many times each operation has been applied to each type of
variable (the type of a variable can be found with the \texttt{TNAM{\textunderscore}OBJ} method). The return value for two argument operators is a record of records \texttt{r}, where \texttt{r.op} stores information about operator \texttt{op}. For one argument operators \texttt{r.op.i} stores how many times \texttt{op} was called with an argument of type \texttt{i}, while for two argument operators \texttt{r.op.i.j} stores how many times \texttt{op} was called with arguments of type \texttt{i} and \texttt{j}. 
\begin{Verbatim}[commandchars=!@|,fontsize=\small,frame=single,label=Example]
  !gapprompt@gap>| !gapinput@TraceInternalMethods();|
  true
  !gapprompt@gap>| !gapinput@2+3+4+5+6;;|
  !gapprompt@gap>| !gapinput@2.0+2.0;;|
  !gapprompt@gap>| !gapinput@3^(1,2,3);;|
  !gapprompt@gap>| !gapinput@GetTraceInternalMethodsCounts();|
  rec( Pow := rec( integer := rec( ("permutation (small)") := 1 ) ),
  Sum := rec( integer := rec( integer := 4 ),
      macfloat := rec( macfloat := 1 ) ) )
  # 'macfloat' is a floating point number
  !gapprompt@gap>| !gapinput@UntraceInternalMethods();|
\end{Verbatim}
 }

 }

  
\section{\textcolor{Chapter }{Info Functions}}\label{Info Functions}
\logpage{[ 7, 4, 0 ]}
\hyperdef{L}{X7A9C902479CB6F7C}{}
{
  The \index{verbosity of GAP output}\texttt{Info} (\ref{Info}) mechanism permits operations to display intermediate results or information
about the progress of the algorithms. Information is always given according to
one or more \emph{info classes}. Each of the info classes defined in the \textsf{GAP} library usually covers a certain range of algorithms, so for example \texttt{InfoLattice} covers all the cyclic extension algorithms for the computation of a subgroup
lattice. 

 Note that not all info classes defined in the \textsf{GAP} library are currently documented. Many \textsf{GAP} packages define additional info classes, which are typically documented in the
corresponding package documentation. The function \texttt{ShowUsedInfoClasses} (\ref{ShowUsedInfoClasses}) will show all info classes which \textsf{GAP} considers while executing code. 

 The amount of information to be displayed by each info class can be separately
specified by the user. This is done by selecting a
non\texttt{\symbol{45}}negative integer \emph{level} for the info class: no information will be displayed at level 0, and the
higher the level, the more information that will be displayed. At creation, an
info class has level 0. By default, all built\texttt{\symbol{45}}in GAP info
classes have level 0, except for the following info classes, which have level
1: 
\begin{itemize}
\item \texttt{InfoWarning} (\ref{InfoWarning}),
\item \texttt{InfoPackageLoading} (\ref{InfoPackageLoading}),
\item \texttt{InfoDebug},
\item \texttt{InfoPerformance},
\item \texttt{InfoTempDirectories},
\item \texttt{InfoPrimeInt}, and
\item \texttt{InfoSLP}.
\end{itemize}
 

 

\subsection{\textcolor{Chapter }{NewInfoClass}}
\logpage{[ 7, 4, 1 ]}\nobreak
\hyperdef{L}{X7AA1A1CF79F20790}{}
{\noindent\textcolor{FuncColor}{$\triangleright$\enspace\texttt{NewInfoClass({\mdseries\slshape name})\index{NewInfoClass@\texttt{NewInfoClass}}
\label{NewInfoClass}
}\hfill{\scriptsize (operation)}}\\


 creates a new info class with name \mbox{\texttt{\mdseries\slshape name}}. }

 

 

\subsection{\textcolor{Chapter }{DeclareInfoClass}}
\logpage{[ 7, 4, 2 ]}\nobreak
\hyperdef{L}{X7B3709C584B3DA1E}{}
{\noindent\textcolor{FuncColor}{$\triangleright$\enspace\texttt{DeclareInfoClass({\mdseries\slshape name})\index{DeclareInfoClass@\texttt{DeclareInfoClass}}
\label{DeclareInfoClass}
}\hfill{\scriptsize (function)}}\\


 creates a new info class with name \mbox{\texttt{\mdseries\slshape name}} and binds it to the global variable \mbox{\texttt{\mdseries\slshape name}}. The variable must previously be writable, and is made
read\texttt{\symbol{45}}only by this function. }

 

 

\subsection{\textcolor{Chapter }{SetInfoLevel}}
\logpage{[ 7, 4, 3 ]}\nobreak
\hyperdef{L}{X7A43B9E68765EE9E}{}
{\noindent\textcolor{FuncColor}{$\triangleright$\enspace\texttt{SetInfoLevel({\mdseries\slshape infoclass, level})\index{SetInfoLevel@\texttt{SetInfoLevel}}
\label{SetInfoLevel}
}\hfill{\scriptsize (operation)}}\\


 Sets the info level for \mbox{\texttt{\mdseries\slshape infoclass}} to the non\texttt{\symbol{45}}negative integer \mbox{\texttt{\mdseries\slshape level}}. }

 

 

\subsection{\textcolor{Chapter }{InfoLevel}}
\logpage{[ 7, 4, 4 ]}\nobreak
\hyperdef{L}{X7B2ADC37783104B9}{}
{\noindent\textcolor{FuncColor}{$\triangleright$\enspace\texttt{InfoLevel({\mdseries\slshape infoclass})\index{InfoLevel@\texttt{InfoLevel}}
\label{InfoLevel}
}\hfill{\scriptsize (operation)}}\\


 returns the info level of \mbox{\texttt{\mdseries\slshape infoclass}}. }

 

 

\subsection{\textcolor{Chapter }{ShowUsedInfoClasses}}
\logpage{[ 7, 4, 5 ]}\nobreak
\hyperdef{L}{X7BA636EF80A1435A}{}
{\noindent\textcolor{FuncColor}{$\triangleright$\enspace\texttt{ShowUsedInfoClasses({\mdseries\slshape infoclass})\index{ShowUsedInfoClasses@\texttt{ShowUsedInfoClasses}}
\label{ShowUsedInfoClasses}
}\hfill{\scriptsize (function)}}\\


 Called with argument \texttt{true}, this makes \textsf{GAP} print the info class and level of any executed \texttt{Info} (\ref{Info}) statement. Calling with the argument \texttt{false} stops this printing. Each level of each info class is only printed once. The
history of printed info classes and levels is reset whenever \texttt{true} is passed. 

 
\begin{Verbatim}[commandchars=!@|,fontsize=\small,frame=single,label=Example]
  !gapprompt@gap>| !gapinput@ShowUsedInfoClasses(true);|
  !gapprompt@gap>| !gapinput@Intersection(Group((1,3,2,4,5,6)), Group((1,2,3,4,5,6)));|
  #I Would print info with SetInfoLevel(InfoBckt,1)
  #I Would print info with SetInfoLevel(InfoBckt,3)
  #I Would print info with SetInfoLevel(InfoBckt,5)
  Group(())
  !gapprompt@gap>| !gapinput@Intersection(Group((1,3,2,4,5,6)), Group((1,2,3,4,5,6)));|
  Group(())
  !gapprompt@gap>| !gapinput@ShowUsedInfoClasses(false);|
\end{Verbatim}
 }

 

 

\subsection{\textcolor{Chapter }{Info}}
\logpage{[ 7, 4, 6 ]}\nobreak
\hyperdef{L}{X864E4B6886E2697D}{}
{\noindent\textcolor{FuncColor}{$\triangleright$\enspace\texttt{Info({\mdseries\slshape infoclass, level, info[, moreinfo, ...]})\index{Info@\texttt{Info}}
\label{Info}
}\hfill{\scriptsize (function)}}\\


 If the info level of \mbox{\texttt{\mdseries\slshape infoclass}} is at least \mbox{\texttt{\mdseries\slshape level}}, then the remaining arguments, \mbox{\texttt{\mdseries\slshape info}}, and possibly \mbox{\texttt{\mdseries\slshape moreinfo}} and so on, are evaluated. (Technically, \texttt{Info} is a keyword and not a function.) 

 By default, the results of these evaluations are viewed, preceded by the
string \texttt{"\#I "} and followed by a newline. 

 If the info level of \mbox{\texttt{\mdseries\slshape infoclass}} is strictly less than \mbox{\texttt{\mdseries\slshape level}}, then the third and subsequent arguments are not evaluated. (The latter can
save substantial time when displaying difficult results.) 

 The behaviour can be customized with \texttt{SetInfoHandler} (\ref{SetInfoHandler}). 

 
\begin{Verbatim}[commandchars=!@|,fontsize=\small,frame=single,label=Example]
  !gapprompt@gap>| !gapinput@InfoExample:=NewInfoClass("InfoExample");;|
  !gapprompt@gap>| !gapinput@Info(InfoExample,1,"one");Info(InfoExample,2,"two");|
  !gapprompt@gap>| !gapinput@SetInfoLevel(InfoExample,1);|
  !gapprompt@gap>| !gapinput@Info(InfoExample,1,"one");Info(InfoExample,2,"two");|
  #I  one
  !gapprompt@gap>| !gapinput@SetInfoLevel(InfoExample,2);|
  !gapprompt@gap>| !gapinput@Info(InfoExample,1,"one");Info(InfoExample,2,"two");|
  #I  one
  #I  two
  !gapprompt@gap>| !gapinput@InfoLevel(InfoExample);|
  2
  !gapprompt@gap>| !gapinput@Info(InfoExample,3,Length(Combinations([1..9999])));|
\end{Verbatim}
 

 Note that the last \texttt{Info} call is executed without problems, since the actual level \texttt{2} of \texttt{InfoExample} causes \texttt{Info} to ignore the last argument, which prevents \texttt{Length(Combinations([1..9999]))} from being evaluated; note that an evaluation would be impossible due to
memory restrictions. 

 A set of info classes (called an \emph{info selector}) may be passed to a single \texttt{Info} statement. As a shorthand, info classes and selectors may be combined with \texttt{+} rather than \texttt{Union} (\ref{Union:for a list}). In this case, the message is triggered if the level of \emph{any} of the classes is high enough. 

 
\begin{Verbatim}[commandchars=!@|,fontsize=\small,frame=single,label=Example]
  !gapprompt@gap>| !gapinput@InfoExample:=NewInfoClass("InfoExample");;|
  !gapprompt@gap>| !gapinput@SetInfoLevel(InfoExample,0);|
  !gapprompt@gap>| !gapinput@Info(InfoExample + InfoWarning, 1, "hello");|
  #I  hello
  !gapprompt@gap>| !gapinput@Info(InfoExample + InfoWarning, 2, "hello");|
  !gapprompt@gap>| !gapinput@SetInfoLevel(InfoExample,2);|
  !gapprompt@gap>| !gapinput@Info(InfoExample + InfoWarning, 2, "hello");|
  #I  hello
  !gapprompt@gap>| !gapinput@InfoLevel(InfoWarning);|
  1
\end{Verbatim}
 }

 
\subsection{\textcolor{Chapter }{Customizing \texttt{Info} (\ref{Info}) statements}}\logpage{[ 7, 4, 7 ]}
\hyperdef{L}{X800234B5815CAC97}{}
{
\noindent\textcolor{FuncColor}{$\triangleright$\enspace\texttt{SetInfoHandler({\mdseries\slshape infoclass, handler})\index{SetInfoHandler@\texttt{SetInfoHandler}}
\label{SetInfoHandler}
}\hfill{\scriptsize (function)}}\\
\noindent\textcolor{FuncColor}{$\triangleright$\enspace\texttt{SetInfoOutput({\mdseries\slshape infoclass, out})\index{SetInfoOutput@\texttt{SetInfoOutput}}
\label{SetInfoOutput}
}\hfill{\scriptsize (function)}}\\
\noindent\textcolor{FuncColor}{$\triangleright$\enspace\texttt{UnbindInfoOutput({\mdseries\slshape infoclass})\index{UnbindInfoOutput@\texttt{UnbindInfoOutput}}
\label{UnbindInfoOutput}
}\hfill{\scriptsize (function)}}\\
\noindent\textcolor{FuncColor}{$\triangleright$\enspace\texttt{InfoOutput({\mdseries\slshape infoclass})\index{InfoOutput@\texttt{InfoOutput}}
\label{InfoOutput}
}\hfill{\scriptsize (function)}}\\
\noindent\textcolor{FuncColor}{$\triangleright$\enspace\texttt{SetDefaultInfoOutput({\mdseries\slshape out})\index{SetDefaultInfoOutput@\texttt{SetDefaultInfoOutput}}
\label{SetDefaultInfoOutput}
}\hfill{\scriptsize (function)}}\\
\textbf{\indent Returns:\ }
nothing



 This allows one to customize what happens in an \texttt{Info(\mbox{\texttt{\mdseries\slshape infoclass}}, \mbox{\texttt{\mdseries\slshape level}}, ...)} statement.

 In the first function, \mbox{\texttt{\mdseries\slshape handler}} must be a function with three arguments \mbox{\texttt{\mdseries\slshape infoclass}}, \mbox{\texttt{\mdseries\slshape level}}, \mbox{\texttt{\mdseries\slshape list}}. Here \mbox{\texttt{\mdseries\slshape list}} is the list containing the third argument and any subsequent optional
arguments of the \texttt{Info} (\ref{Info}) call. 

 The default handler is the function \texttt{DefaultInfoHandler}. \index{DefaultInfoHandler@\texttt{DefaultInfoHandler}} It prints \texttt{"\#I "}, then the third and further arguments of the info statement, and finally a \texttt{"\texttt{\symbol{92}}n"}. 

 If the first argument of an \texttt{Info} (\ref{Info}) statement is a sum of Info classes, the handler of the first summand is used. 

 The file or stream to which \texttt{Info} (\ref{Info}) statements for individual \texttt{Info} (\ref{Info}) classes print can be overridden with \texttt{SetInfoOutput}, retrieved with \texttt{InfoOutput} and reset to the default with \texttt{UnbindInfoOutput}. The initial default for all \texttt{Info} (\ref{Info}) classes is the string \texttt{"*Print*"} which means the current output file. The default can be changed with \texttt{SetDefaultInfoOutput}. The argument \mbox{\texttt{\mdseries\slshape out}} can be a filename or an open stream, the special names \texttt{"*Print*"}, \texttt{"*errout*} and \texttt{"*stdout*} are also recognized.

 For example, \texttt{SetDefaultInfoOutput("*errout*");} would send \texttt{Info} (\ref{Info}) output to standard error, which can be interesting if \textsf{GAP}s output is redirected. }

 

\subsection{\textcolor{Chapter }{InfoWarning}}
\logpage{[ 7, 4, 8 ]}\nobreak
\hyperdef{L}{X7A28F77C82D6A3E0}{}
{\noindent\textcolor{FuncColor}{$\triangleright$\enspace\texttt{InfoWarning\index{InfoWarning@\texttt{InfoWarning}}
\label{InfoWarning}
}\hfill{\scriptsize (info class)}}\\


 is an info class to which general warnings are sent at level 1, which is its
default level. More specialised warnings are shown via calls of \texttt{Info} (\ref{Info}) at \texttt{InfoWarning} level 2, e.g.{\nobreakspace}information about the autoloading of \textsf{GAP} packages and the initial line matched when displaying an
on\texttt{\symbol{45}}line help topic. }

 }

  
\section{\textcolor{Chapter }{Assertions}}\label{Assertions}
\logpage{[ 7, 5, 0 ]}
\hyperdef{L}{X86425F067FC63A4C}{}
{
  Assertions are used to find errors in algorithms. They test whether
intermediate results conform to required conditions and issue an error if not. 

 

\subsection{\textcolor{Chapter }{SetAssertionLevel}}
\logpage{[ 7, 5, 1 ]}\nobreak
\hyperdef{L}{X7C7596418423660B}{}
{\noindent\textcolor{FuncColor}{$\triangleright$\enspace\texttt{SetAssertionLevel({\mdseries\slshape lev})\index{SetAssertionLevel@\texttt{SetAssertionLevel}}
\label{SetAssertionLevel}
}\hfill{\scriptsize (function)}}\\


 assigns the global assertion level to \mbox{\texttt{\mdseries\slshape lev}}. By default it is zero. }

 

 

\subsection{\textcolor{Chapter }{AssertionLevel}}
\logpage{[ 7, 5, 2 ]}\nobreak
\hyperdef{L}{X876C83707F13A0FD}{}
{\noindent\textcolor{FuncColor}{$\triangleright$\enspace\texttt{AssertionLevel({\mdseries\slshape })\index{AssertionLevel@\texttt{AssertionLevel}}
\label{AssertionLevel}
}\hfill{\scriptsize (function)}}\\


 returns the current assertion level. }

 

 

\subsection{\textcolor{Chapter }{Assert}}
\logpage{[ 7, 5, 3 ]}\nobreak
\hyperdef{L}{X830E443284780FB9}{}
{\noindent\textcolor{FuncColor}{$\triangleright$\enspace\texttt{Assert({\mdseries\slshape lev, cond[, message]})\index{Assert@\texttt{Assert}}
\label{Assert}
}\hfill{\scriptsize (function)}}\\


 With two arguments, if the global assertion level is at least \mbox{\texttt{\mdseries\slshape lev}}, condition \mbox{\texttt{\mdseries\slshape cond}} is tested and if it does not return \texttt{true} an error is raised. Thus \texttt{Assert(lev, \mbox{\texttt{\mdseries\slshape cond}})} is equivalent to the code 
\begin{Verbatim}[commandchars=!@|,fontsize=\small,frame=single,label=Example]
  if AssertionLevel() >= lev and not <cond> then
    Error("Assertion failure");
  fi;
\end{Verbatim}
 

 If the \mbox{\texttt{\mdseries\slshape message}} argument form of the \texttt{Assert} statement is provided, and if an error is raised, then this message is printed
as part of the error. 

 Assertions are used at various places in the library. Thus turning assertions
on can slow code execution significantly. }

 }

  
\section{\textcolor{Chapter }{Timing}}\label{Timing}
\logpage{[ 7, 6, 0 ]}
\hyperdef{L}{X792BA9A67E64CDED}{}
{
  

\subsection{\textcolor{Chapter }{Runtimes}}
\logpage{[ 7, 6, 1 ]}\nobreak
\hyperdef{L}{X80355C9282B35673}{}
{\noindent\textcolor{FuncColor}{$\triangleright$\enspace\texttt{Runtimes({\mdseries\slshape })\index{Runtimes@\texttt{Runtimes}}
\label{Runtimes}
}\hfill{\scriptsize (function)}}\\


 \texttt{Runtimes} returns a record with components bound to integers or \texttt{fail}. Each integer is the cpu time (processor time) in milliseconds spent by \textsf{GAP} in a certain status: 

 
\begin{description}
\item[{\texttt{user{\textunderscore}time}}]  cpu time spent with \textsf{GAP} functions (without child processes). 
\item[{\texttt{system{\textunderscore}time}}]  cpu time spent in system calls, e.g., file access (\texttt{fail} if not available). 
\item[{\texttt{user{\textunderscore}time{\textunderscore}children}}]  cpu time spent in child processes (\texttt{fail} if not available). 
\item[{\texttt{system{\textunderscore}time{\textunderscore}children}}]  cpu time spent in system calls by child processes (\texttt{fail} if not available). 
\end{description}
 

 Note that this function is not fully supported on all systems. Only the \texttt{user{\textunderscore}time} component is (and may on some systems include the system time). 

 The following example demonstrates tasks which contribute to the different
time components: 

 
\begin{Verbatim}[commandchars=!@B,fontsize=\small,frame=single,label=Example]
  !gapprompt@gap>B !gapinput@Runtimes(); # after startupB
  rec( user_time := 3980, system_time := 60, user_time_children := 0,
    system_time_children := 0 )
  !gapprompt@gap>B !gapinput@Exec("cat /usr/bin/*||wc"); # child process with a lot of file accessB
   893799 7551659 200928302
  !gapprompt@gap>B !gapinput@Runtimes();B
  rec( user_time := 3990, system_time := 60, user_time_children := 1590,
    system_time_children := 600 )
  !gapprompt@gap>B !gapinput@a:=0;;for i in [1..100000000] do a:=a+1; od; # GAP user timeB
  !gapprompt@gap>B !gapinput@Runtimes();B
  rec( user_time := 12980, system_time := 70, user_time_children := 1590,
    system_time_children := 600 )
  !gapprompt@gap>B !gapinput@?blabla  # first call of help, a lot of file accessB
  Help: no matching entry found
  !gapprompt@gap>B !gapinput@Runtimes();B
  rec( user_time := 13500, system_time := 440, user_time_children := 1590,
    system_time_children := 600 )
\end{Verbatim}
 }

 

\subsection{\textcolor{Chapter }{Runtime}}
\logpage{[ 7, 6, 2 ]}\nobreak
\hyperdef{L}{X7E32B27F81870D24}{}
{\noindent\textcolor{FuncColor}{$\triangleright$\enspace\texttt{Runtime({\mdseries\slshape })\index{Runtime@\texttt{Runtime}}
\label{Runtime}
}\hfill{\scriptsize (function)}}\\


 \texttt{Runtime} returns the time spent by \textsf{GAP} in milliseconds as an integer. It is the same as the value of the \texttt{user{\textunderscore}time} component given by \texttt{Runtimes} (\ref{Runtimes}), as explained above. 

 See \texttt{StringTime} (\ref{StringTime}) for a translation from milliseconds into hour/minute format. }

 

\subsection{\textcolor{Chapter }{NanosecondsSinceEpoch}}
\logpage{[ 7, 6, 3 ]}\nobreak
\hyperdef{L}{X844E1CFE80F41760}{}
{\noindent\textcolor{FuncColor}{$\triangleright$\enspace\texttt{NanosecondsSinceEpoch({\mdseries\slshape })\index{NanosecondsSinceEpoch@\texttt{NanosecondsSinceEpoch}}
\label{NanosecondsSinceEpoch}
}\hfill{\scriptsize (function)}}\\
\noindent\textcolor{FuncColor}{$\triangleright$\enspace\texttt{NanosecondsSinceEpochInfo({\mdseries\slshape })\index{NanosecondsSinceEpochInfo@\texttt{NanosecondsSinceEpochInfo}}
\label{NanosecondsSinceEpochInfo}
}\hfill{\scriptsize (function)}}\\


 \texttt{NanosecondsSinceEpoch} returns the time in nanoseconds that has passed since some fixed, but
unspecified time in the past. This function is appropriate for doing wallclock
time measurements. The actual resolution depends on the system that \textsf{GAP} is run on. Information about the used timers can be obtained by calling \texttt{NanosecondsSinceEpochInfo}, which returns a record containing members \texttt{Method}, \texttt{Monotonic}, \texttt{Reliable} and \texttt{Resolution}. 

 \texttt{Method} is a string describing the method used to obtain timer values. This will
usually contain the name of the syscall used. 

 \texttt{Monotonic} is a boolean. If it is \texttt{true}, then the values returned by \texttt{NanosecondsSinceEpoch} are guaranteed to be strictly monotonically increasing between two calls, if
it is \texttt{false} then there is no such guarantee. 

 \texttt{Resolution} is an integer reflecting the resolution of the timer used in nanoseconds. 

 \texttt{Reliable} is a boolean. If it is \texttt{true} then the value \texttt{Resolution} is deemed reliable in the sense that it was obtained by querying the operating
system, otherwise \texttt{Resolution} should be treated as an estimate. }

 

\subsection{\textcolor{Chapter }{time}}
\logpage{[ 7, 6, 4 ]}\nobreak
\hyperdef{L}{X7C0F91F982189624}{}
{\noindent\textcolor{FuncColor}{$\triangleright$\enspace\texttt{time\index{time@\texttt{time}}
\label{time}
}\hfill{\scriptsize (global variable)}}\\


 In the read\texttt{\symbol{45}}eval\texttt{\symbol{45}}print loop, \texttt{time} stores the number of milliseconds the last command took (see also \texttt{memory{\textunderscore}allocated} (\ref{memoryuScoreallocated}) for the number of bytes of memory it allocated). }

 

\subsection{\textcolor{Chapter }{Sleep}}
\logpage{[ 7, 6, 5 ]}\nobreak
\hyperdef{L}{X7B543F357C7202CF}{}
{\noindent\textcolor{FuncColor}{$\triangleright$\enspace\texttt{Sleep({\mdseries\slshape time})\index{Sleep@\texttt{Sleep}}
\label{Sleep}
}\hfill{\scriptsize (function)}}\\
\noindent\textcolor{FuncColor}{$\triangleright$\enspace\texttt{MicroSleep({\mdseries\slshape time})\index{MicroSleep@\texttt{MicroSleep}}
\label{MicroSleep}
}\hfill{\scriptsize (function)}}\\


 These functions make GAP stop execution for a given period of time. The time
to stop is given to \texttt{Sleep} in seconds and \texttt{MicroSleep} in microseconds. }

 }

  
\section{\textcolor{Chapter }{Tracking Memory Usage}}\label{Tracking Memory Usage}
\logpage{[ 7, 7, 0 ]}
\hyperdef{L}{X844CB04081A771FB}{}
{
  

\subsection{\textcolor{Chapter }{TotalMemoryAllocated}}
\logpage{[ 7, 7, 1 ]}\nobreak
\hyperdef{L}{X8077B50B844C4EFC}{}
{\noindent\textcolor{FuncColor}{$\triangleright$\enspace\texttt{TotalMemoryAllocated({\mdseries\slshape })\index{TotalMemoryAllocated@\texttt{TotalMemoryAllocated}}
\label{TotalMemoryAllocated}
}\hfill{\scriptsize (function)}}\\


 \texttt{TotalMemoryAllocated} returns the total amount of memory in bytes allocated by the \textsf{GAP} memory manager since \textsf{GAP} started. }

 

\subsection{\textcolor{Chapter }{memory{\textunderscore}allocated}}
\logpage{[ 7, 7, 2 ]}\nobreak
\hyperdef{L}{X8156D7208591460F}{}
{\noindent\textcolor{FuncColor}{$\triangleright$\enspace\texttt{memory{\textunderscore}allocated\index{memory{\textunderscore}allocated@\texttt{memory{\textunderscore}allocated}}
\label{memoryuScoreallocated}
}\hfill{\scriptsize (global variable)}}\\


 In the read\texttt{\symbol{45}}eval\texttt{\symbol{45}}print loop, \texttt{memory{\textunderscore}allocated} stores the number of bytes of memory allocated by the last completed statement
(see also \texttt{time} (\ref{time}) for the number of milliseconds it took). }

 }

  
\section{\textcolor{Chapter }{Profiling}}\label{Profiling}
\logpage{[ 7, 8, 0 ]}
\hyperdef{L}{X7FDF923D7D2937A1}{}
{
  Profiling of code can be used to determine in which parts of a program how
much time has been spent and how much memory has been allocated during
runtime. GAP has two different methods of profiling. GAP can either profile by
function, or line\texttt{\symbol{45}}by\texttt{\symbol{45}}line. Line by line
profiling is currently only used for code coverage, while function profiling
tracks memory and time usage. 
\subsection{\textcolor{Chapter }{Function Profiling}}\label{FunctionProfiling}
\logpage{[ 7, 8, 1 ]}
\hyperdef{L}{X7939F6F182FDA5F1}{}
{
  This section describes how to profiling at the function level. The idea is
that 
\begin{itemize}
\item  first one switches on profiling for those \textsf{GAP} functions the performance of which one wants to check, 
\item  then one runs some \textsf{GAP} computations, 
\item  then one looks at the profile information collected during these computations, 
\item  then one runs more computations (perhaps clearing all profile information
before, see \texttt{ClearProfile} (\ref{ClearProfile})), 
\item  and finally one switches off profiling. 
\end{itemize}
 

 For switching on and off profiling, \textsf{GAP} supports entering a list of functions (see \texttt{ProfileFunctions} (\ref{ProfileFunctions}), \texttt{UnprofileFunctions} (\ref{UnprofileFunctions})) or a list of operations whose methods shall be (un)profiled (\texttt{ProfileMethods} (\ref{ProfileMethods}), \texttt{UnprofileMethods} (\ref{UnprofileMethods})), and \texttt{DisplayProfile} (\ref{DisplayProfile}) can be used to show profile information about functions in a given list. 

 Besides these functions, \texttt{ProfileGlobalFunctions} (\ref{ProfileGlobalFunctions}), \texttt{ProfileOperations} (\ref{ProfileOperations}), and \texttt{ProfileOperationsAndMethods} (\ref{ProfileOperationsAndMethods}) can be used for switching on or off profiling for \emph{all} global functions, operations, and operations together with all their methods,
respectively, and for showing profile information about these functions. 

 Note that \textsf{GAP} will perform more slowly when profiling than when not. }

 

\subsection{\textcolor{Chapter }{ProfileGlobalFunctions}}
\logpage{[ 7, 8, 2 ]}\nobreak
\hyperdef{L}{X79D6CB927BBEB940}{}
{\noindent\textcolor{FuncColor}{$\triangleright$\enspace\texttt{ProfileGlobalFunctions({\mdseries\slshape [bool]})\index{ProfileGlobalFunctions@\texttt{ProfileGlobalFunctions}}
\label{ProfileGlobalFunctions}
}\hfill{\scriptsize (function)}}\\


 Called with argument \texttt{true}, \texttt{ProfileGlobalFunctions} starts profiling of all functions that have been declared via \texttt{DeclareGlobalFunction} (\ref{DeclareGlobalFunction}). Old profile information for all these functions is cleared. A function call
with the argument \texttt{false} stops profiling of all these functions. Recorded information is still kept, so
you can display it even after turning the profiling off. 

 When \texttt{ProfileGlobalFunctions} is called without argument, profile information for all global functions is
displayed, see{\nobreakspace}\texttt{DisplayProfile} (\ref{DisplayProfile}). }

 

\subsection{\textcolor{Chapter }{ProfileOperations}}
\logpage{[ 7, 8, 3 ]}\nobreak
\hyperdef{L}{X7C893F68841B990B}{}
{\noindent\textcolor{FuncColor}{$\triangleright$\enspace\texttt{ProfileOperations({\mdseries\slshape [bool]})\index{ProfileOperations@\texttt{ProfileOperations}}
\label{ProfileOperations}
}\hfill{\scriptsize (function)}}\\


 Called with argument \texttt{true}, \texttt{ProfileOperations} starts profiling of all operations. Old profile information for all operations
is cleared. A function call with the argument \texttt{false} stops profiling of all operations. Recorded information is still kept, so you
can display it even after turning the profiling off. 

 When \texttt{ProfileOperations} is called without argument, profile information for all operations is
displayed (see{\nobreakspace}\texttt{DisplayProfile} (\ref{DisplayProfile})). }

 

\subsection{\textcolor{Chapter }{ProfileOperationsAndMethods}}
\logpage{[ 7, 8, 4 ]}\nobreak
\hyperdef{L}{X79D41E977DCA2BEE}{}
{\noindent\textcolor{FuncColor}{$\triangleright$\enspace\texttt{ProfileOperationsAndMethods({\mdseries\slshape [bool]})\index{ProfileOperationsAndMethods@\texttt{ProfileOperationsAndMethods}}
\label{ProfileOperationsAndMethods}
}\hfill{\scriptsize (function)}}\\


 Called with argument \texttt{true}, \texttt{ProfileOperationsAndMethods} starts profiling of all operations and their methods. Old profile information
for these functions is cleared. A function call with the argument \texttt{false} stops profiling of all operations and their methods. Recorded information is
still kept, so you can display it even after turning the profiling off. 

 When \texttt{ProfileOperationsAndMethods} is called without argument, profile information for all operations and their
methods is displayed, see{\nobreakspace}\texttt{DisplayProfile} (\ref{DisplayProfile}). }

 

\subsection{\textcolor{Chapter }{ProfileFunctions}}
\logpage{[ 7, 8, 5 ]}\nobreak
\hyperdef{L}{X81E8A8627C34FD3B}{}
{\noindent\textcolor{FuncColor}{$\triangleright$\enspace\texttt{ProfileFunctions({\mdseries\slshape funcs})\index{ProfileFunctions@\texttt{ProfileFunctions}}
\label{ProfileFunctions}
}\hfill{\scriptsize (function)}}\\


 starts profiling for all function in the list \mbox{\texttt{\mdseries\slshape funcs}}. You can use \texttt{ProfileGlobalFunctions} (\ref{ProfileGlobalFunctions}) to turn profiling on for all globally declared functions simultaneously. }

 

\subsection{\textcolor{Chapter }{UnprofileFunctions}}
\logpage{[ 7, 8, 6 ]}\nobreak
\hyperdef{L}{X79D394EC7BE8D008}{}
{\noindent\textcolor{FuncColor}{$\triangleright$\enspace\texttt{UnprofileFunctions({\mdseries\slshape funcs})\index{UnprofileFunctions@\texttt{UnprofileFunctions}}
\label{UnprofileFunctions}
}\hfill{\scriptsize (function)}}\\


 stops profiling for all function in the list \mbox{\texttt{\mdseries\slshape funcs}}. Recorded information is still kept, so you can display it even after turning
the profiling off. }

 

\subsection{\textcolor{Chapter }{ProfileMethods}}
\logpage{[ 7, 8, 7 ]}\nobreak
\hyperdef{L}{X787AC3BE7F991344}{}
{\noindent\textcolor{FuncColor}{$\triangleright$\enspace\texttt{ProfileMethods({\mdseries\slshape ops})\index{ProfileMethods@\texttt{ProfileMethods}}
\label{ProfileMethods}
}\hfill{\scriptsize (function)}}\\


 starts profiling of the methods for all operations in the list \mbox{\texttt{\mdseries\slshape ops}}. }

 

\subsection{\textcolor{Chapter }{UnprofileMethods}}
\logpage{[ 7, 8, 8 ]}\nobreak
\hyperdef{L}{X87A05F977F033693}{}
{\noindent\textcolor{FuncColor}{$\triangleright$\enspace\texttt{UnprofileMethods({\mdseries\slshape ops})\index{UnprofileMethods@\texttt{UnprofileMethods}}
\label{UnprofileMethods}
}\hfill{\scriptsize (function)}}\\


 stops profiling of the methods for all operations in the list \mbox{\texttt{\mdseries\slshape ops}}. Recorded information is still kept, so you can display it even after turning
the profiling off. }

 

\subsection{\textcolor{Chapter }{DisplayProfile}}
\logpage{[ 7, 8, 9 ]}\nobreak
\hyperdef{L}{X80FEA6A08775A48E}{}
{\noindent\textcolor{FuncColor}{$\triangleright$\enspace\texttt{DisplayProfile({\mdseries\slshape [functions][,] [mincount, mintime]})\index{DisplayProfile@\texttt{DisplayProfile}}
\label{DisplayProfile}
}\hfill{\scriptsize (function)}}\\
\noindent\textcolor{FuncColor}{$\triangleright$\enspace\texttt{GAPInfo.ProfileThreshold\index{GAPInfo.ProfileThreshold@\texttt{GAPInfo.ProfileThreshold}}
\label{GAPInfo.ProfileThreshold}
}\hfill{\scriptsize (global variable)}}\\


 Called without arguments, \texttt{DisplayProfile} displays the profile information for profiled operations, methods and
functions. If an argument \mbox{\texttt{\mdseries\slshape functions}} is given, only profile information for the functions in the list \mbox{\texttt{\mdseries\slshape functions}} is shown. If two integer values \mbox{\texttt{\mdseries\slshape mincount}}, \mbox{\texttt{\mdseries\slshape mintime}} are given as arguments then the output is restricted to those functions that
were called at least \mbox{\texttt{\mdseries\slshape mincount}} times or for which the total time spent (see below) was at least \mbox{\texttt{\mdseries\slshape mintime}} milliseconds. The defaults for \mbox{\texttt{\mdseries\slshape mincount}} and \mbox{\texttt{\mdseries\slshape mintime}} are the entries of the list stored in the global variable \texttt{GAPInfo.ProfileThreshold}. 

 The default value of \texttt{GAPInfo.ProfileThreshold} is \texttt{[ 10000, 30 ]}. 

 Profile information is displayed in a list of lines for all functions
(including operations and methods) which are profiled. For each function, ``count'' gives the number of times the function has been called. ``self/ms'' gives the time (in milliseconds) spent in the function itself, ``chld/ms'' the time (in milliseconds) spent in profiled functions called from within this
function, ``stor/kb'' the amount of storage (in kilobytes) allocated by the function itself, ``chld/kb'' the amount of storage (in kilobytes) allocated by profiled functions called
from within this function, and ``package'' the name of the \textsf{GAP} package to which the function belongs; the entry ``GAP'' in this column means that the function belongs to the \textsf{GAP} library, the entry ``(oprt.)'' means that the function is an operation (which may belong to several
packages), and an empty entry means that \texttt{FilenameFunc} (\ref{FilenameFunc}) cannot determine in which file the function is defined. 

 The list is sorted according to the total time spent in the functions, that is
the sum of the values in the columns ``self/ms'' and ``chld/ms''. 

 At the end of the list, two lines are printed that show the total time used
and the total memory allocated by the profiled functions not shown in the list
(label \texttt{OTHER}) and by all profiled functions (label \texttt{TOTAL}), respectively. 

 An interactive variant of \texttt{DisplayProfile} is the function \texttt{BrowseProfile} (\textbf{Browse: BrowseProfile}) that is provided by the \textsf{GAP} package \textsf{Browse}. }

 

\subsection{\textcolor{Chapter }{ClearProfile}}
\logpage{[ 7, 8, 10 ]}\nobreak
\hyperdef{L}{X7DAF9AB9793AE203}{}
{\noindent\textcolor{FuncColor}{$\triangleright$\enspace\texttt{ClearProfile({\mdseries\slshape })\index{ClearProfile@\texttt{ClearProfile}}
\label{ClearProfile}
}\hfill{\scriptsize (function)}}\\


 clears all stored profile information. }

  
\subsection{\textcolor{Chapter }{An Example of Function Profiling}}\label{subsect:profiling_example}
\logpage{[ 7, 8, 11 ]}
\hyperdef{L}{X7C5CE32579891120}{}
{
  Let us suppose we want to get information about the computation of the
conjugacy classes of a certain permutation group. For that, first we create
the group, then we start profiling for all global functions and for all
operations and their methods, then we compute the conjugacy classes, and then
we stop profiling. 

 
\begin{Verbatim}[commandchars=!@|,fontsize=\small,frame=single,label=Example]
  !gapprompt@gap>| !gapinput@g:= PrimitiveGroup( 24, 1 );;|
  !gapprompt@gap>| !gapinput@ProfileGlobalFunctions( true );|
  !gapprompt@gap>| !gapinput@ProfileOperationsAndMethods( true );|
  !gapprompt@gap>| !gapinput@ConjugacyClasses( g );;|
  !gapprompt@gap>| !gapinput@ProfileGlobalFunctions( false );|
  !gapprompt@gap>| !gapinput@ProfileOperationsAndMethods( false );|
\end{Verbatim}
 

 Now the profile information is available. We can list the information for all
profiled functions with \texttt{DisplayProfile} (\ref{DisplayProfile}). 

 
\begin{Verbatim}[commandchars=!@|,fontsize=\small,frame=single,label=Example]
  !gapprompt@gap>| !gapinput@DisplayProfile();|
    count  self/ms  chld/ms  stor/kb  chld/kb  package  function
    17647        0        0      275        0  GAP      BasePoint
    10230        0        0      226        0  (oprt.)  ShallowCopy
    10139        0        0        0        0           PositionSortedOp: for*
    10001        0        0      688        0           UniteSet: for two int*
    10001        8        0       28      688  (oprt.)  UniteSet
    14751       12        0        0        0           =: for two families: *
    10830        8        4      182      276  GAP      Concatenation
     2700       20       12      313       55  GAP      AddRefinement
     2444       28        4     3924      317  GAP      ConjugateStabChain
     4368        0       32        7      714  (oprt.)  Size
     2174       32        4     1030      116  GAP      List
      585        4       32       45      742  GAP      RRefine
     1532       32        8      194       56  GAP      AddGeneratorsExtendSc*
     1221        8       32      349      420  GAP      Partition
   185309       28       12        0        0  (oprt.)  Length
      336        4       40       95      817  GAP      ExtendSeriesPermGroup
        4       28       20      488      454  (oprt.)  Sortex
     2798        0       52       54      944  GAP      StabChainForcePoint
      560        4       48       83      628  GAP      StabChainSwap
      432       16       40      259      461  GAP      SubmagmaWithInversesNC
   185553       48        8      915       94  (oprt.)  Add
       26        0       64        0     2023  (oprt.)  CentralizerOp
       26        0       64        0     2023  GAP      CentralizerOp: perm g*
       26        0       64        0     2023  GAP      Centralizer: try to e*
      152        4       64        0     2024  (oprt.)  Centralizer
     1605        0       68        0     2032  (oprt.)  StabilizerOfExternalS*
       26        0       68        0     2024  GAP      Meth(StabilizerOfExte*
      382        0       96       69     1922  GAP      TryPcgsPermGroup
     5130        4       96      309     3165  GAP      ForAll
     7980       24      116      330     6434  GAP      ChangeStabChain
    12076       12      136      351     6478  GAP      ProcessFixpoint
      192        0      148        4     3029  GAP      StabChainMutable: cal*
     2208        4      148        3     3083  (oprt.)  StabChainMutable
      217        0      160        0     3177  (oprt.)  StabChainOp
      217       12      148       60     3117  GAP      StabChainOp: group an*
      216       36      464      334    12546  GAP      PartitionBacktrack
     1479       12      668      566    18474  GAP      RepOpElmTuplesPermGro*
     1453       12      684       56    18460  GAP      in: perm class rep
      126        0      728       13    19233  GAP      ConjugacyClassesTry
        1        0      736        0    19671  GAP      ConjugacyClassesByRan*
        2        0      736        2    19678  (oprt.)  ConjugacyClasses
        1        0      736        0    19675  GAP      ConjugacyClasses: per*
    13400     1164        0        0        0  (oprt.)  Position
               484             12052                    OTHER
              2048             23319                    TOTAL
\end{Verbatim}
 

 We can restrict the list to global functions with \texttt{ProfileGlobalFunctions} (\ref{ProfileGlobalFunctions}). 

 
\begin{Verbatim}[commandchars=!@|,fontsize=\small,frame=single,label=Example]
  !gapprompt@gap>| !gapinput@ProfileGlobalFunctions();|
    count  self/ms  chld/ms  stor/kb  chld/kb  package  function
    17647        0        0      275        0  GAP      BasePoint
    10830        8        4      182      276  GAP      Concatenation
     2700       20       12      313       55  GAP      AddRefinement
     2444       28        4     3924      317  GAP      ConjugateStabChain
     2174       32        4     1030      116  GAP      List
      585        4       32       45      742  GAP      RRefine
     1532       32        8      194       56  GAP      AddGeneratorsExtendSc*
     1221        8       32      349      420  GAP      Partition
      336        4       40       95      817  GAP      ExtendSeriesPermGroup
     2798        0       52       54      944  GAP      StabChainForcePoint
      560        4       48       83      628  GAP      StabChainSwap
      432       16       40      259      461  GAP      SubmagmaWithInversesNC
      382        0       96       69     1922  GAP      TryPcgsPermGroup
     5130        4       96      309     3165  GAP      ForAll
     7980       24      116      330     6434  GAP      ChangeStabChain
    12076       12      136      351     6478  GAP      ProcessFixpoint
      216       36      464      334    12546  GAP      PartitionBacktrack
     1479       12      668      566    18474  GAP      RepOpElmTuplesPermGro*
      126        0      728       13    19233  GAP      ConjugacyClassesTry
        1        0      736        0    19671  GAP      ConjugacyClassesByRan*
              1804             14536                    OTHER
              2048             23319                    TOTAL
\end{Verbatim}
 

 We can restrict the list to operations with \texttt{ProfileOperations} (\ref{ProfileOperations}). 

 
\begin{Verbatim}[commandchars=!@|,fontsize=\small,frame=single,label=Example]
  !gapprompt@gap>| !gapinput@ProfileOperations();|
    count  self/ms  chld/ms  stor/kb  chld/kb  package  function
    10230        0        0      226        0  (oprt.)  ShallowCopy
    10001        8        0       28      688  (oprt.)  UniteSet
     4368        0       32        7      714  (oprt.)  Size
   185309       28       12        0        0  (oprt.)  Length
        4       28       20      488      454  (oprt.)  Sortex
   185553       48        8      915       94  (oprt.)  Add
       26        0       64        0     2023  (oprt.)  CentralizerOp
      152        4       64        0     2024  (oprt.)  Centralizer
     1605        0       68        0     2032  (oprt.)  StabilizerOfExternalS*
     2208        4      148        3     3083  (oprt.)  StabChainMutable
      217        0      160        0     3177  (oprt.)  StabChainOp
        2        0      736        2    19678  (oprt.)  ConjugacyClasses
    13400     1164        0        0        0  (oprt.)  Position
               764             21646                    OTHER
              2048             23319                    TOTAL
\end{Verbatim}
 

 We can restrict the list to operations and their methods with \texttt{ProfileOperationsAndMethods} (\ref{ProfileOperationsAndMethods}). 

 
\begin{Verbatim}[commandchars=!@|,fontsize=\small,frame=single,label=Example]
  !gapprompt@gap>| !gapinput@ProfileOperationsAndMethods();|
    count  self/ms  chld/ms  stor/kb  chld/kb  package  function
    10230        0        0      226        0  (oprt.)  ShallowCopy
    10139        0        0        0        0           PositionSortedOp: for*
    10001        0        0      688        0           UniteSet: for two int*
    10001        8        0       28      688  (oprt.)  UniteSet
    14751       12        0        0        0           =: for two families: *
     4368        0       32        7      714  (oprt.)  Size
   185309       28       12        0        0  (oprt.)  Length
        4       28       20      488      454  (oprt.)  Sortex
   185553       48        8      915       94  (oprt.)  Add
       26        0       64        0     2023  (oprt.)  CentralizerOp
       26        0       64        0     2023  GAP      CentralizerOp: perm g*
       26        0       64        0     2023  GAP      Centralizer: try to e*
      152        4       64        0     2024  (oprt.)  Centralizer
     1605        0       68        0     2032  (oprt.)  StabilizerOfExternalS*
       26        0       68        0     2024  GAP      Meth(StabilizerOfExte*
      192        0      148        4     3029  GAP      StabChainMutable: cal*
     2208        4      148        3     3083  (oprt.)  StabChainMutable
      217        0      160        0     3177  (oprt.)  StabChainOp
      217       12      148       60     3117  GAP      StabChainOp: group an*
     1453       12      684       56    18460  GAP      in: perm class rep
        2        0      736        2    19678  (oprt.)  ConjugacyClasses
        1        0      736        0    19675  GAP      ConjugacyClasses: per*
    13400     1164        0        0        0  (oprt.)  Position
               728             20834                    OTHER
              2048             23319                    TOTAL
\end{Verbatim}
 

 Finally, we can restrict the list to explicitly given functions with \texttt{DisplayProfile} (\ref{DisplayProfile}), by entering the list of functions as an argument. 

 
\begin{Verbatim}[commandchars=!@|,fontsize=\small,frame=single,label=Example]
  !gapprompt@gap>| !gapinput@DisplayProfile( [ StabChainOp, Centralizer ] );|
    count  self/ms  chld/ms  stor/kb  chld/kb  package  function
      152        4       64        0     2024  (oprt.)  Centralizer
      217        0      160        0     3177  (oprt.)  StabChainOp
              2044             23319                    OTHER
              2048             23319                    TOTAL
\end{Verbatim}
  }

 
\subsection{\textcolor{Chapter }{Line By Line Profiling}}\label{linebylineprofiling}
\logpage{[ 7, 8, 12 ]}
\hyperdef{L}{X812F9CE0817110EA}{}
{
  Line By Line profiling tracks which lines have been executed in a piece of GAP
code. Built into GAP are the methods necessary to generate profiles, the
resulting profiles can be displayed with the 'profiling' package. }

 
\subsection{\textcolor{Chapter }{Line by Line profiling example}}\label{subsect:linebylineprofexample}
\logpage{[ 7, 8, 13 ]}
\hyperdef{L}{X7E9C65B17B8EF993}{}
{
  There are two kinds of profiles GAP can build: 
\begin{itemize}
\item  Coverage : This records which lines of code are executed
\item  Timing : This records how much time is spend executing each line of code 
\end{itemize}
 A timing profile provides more information, but will take longer to generate
and parse. A timing profile is generated using the functions \texttt{ProfileLineByLine} (\ref{ProfileLineByLine}) and \texttt{UnprofileLineByLine} (\ref{UnprofileLineByLine}), as follows: 
\begin{Verbatim}[commandchars=!@|,fontsize=\small,frame=single,label=Example]
  !gapprompt@gap>| !gapinput@ProfileLineByLine("output.gz");|
  !gapprompt@gap>| !gapinput@Size(AlternatingGroup(10)); ; # Execute some GAP code you want to profile|
  !gapprompt@gap>| !gapinput@UnprofileLineByLine();|
\end{Verbatim}
 For code coverage, use instead the functions \texttt{CoverageLineByLine} (\ref{CoverageLineByLine}) and \texttt{UncoverageLineByLine} (\ref{UncoverageLineByLine}). The profiler will only record lines which are read and executed while the
profiler is running. If you want to perform code coverage or profile GAP's
library, then you can use the GAP command line option
'\texttt{\symbol{45}}\texttt{\symbol{45}}cover filename.gz', which executes \texttt{CoverageLineByLine} (\ref{CoverageLineByLine}) before GAP starts. Similarly the option
'\texttt{\symbol{45}}\texttt{\symbol{45}}prof filename.gz' executes \texttt{ProfileLineByLine} (\ref{ProfileLineByLine}) before GAP starts. The profiler is designed for high performance, because of
this, there are some limitations which users should be aware of: 
\begin{itemize}
\item  By default the profiler records the wall\texttt{\symbol{45}}clock time which
has passed, rather than the CPU time taken (because it is lower overhead), so
any time taken writing commands will be charged to the last GAP statement
which was executed. Therefore it is better to write a function which starts
profiling, executes your code, and then stops profiling. 
\item  If you end the filename with ".gz", the resulting file will automatically be
compressed. This is highly recommended! 
\item  The profiler can only track GAP code which occurs in a function
\texttt{\symbol{45}}\texttt{\symbol{45}} this is most obvious when looking at
code coverage examples, which will appear to miss lines of code in files not
in a function. 
\item  If the current GAP is forked, using the \texttt{IO{\textunderscore}fork} function in the \textsf{IO} package, a new profile output file will be created for the new child process,
with the process ID of the child attached to the end of the filename. 
\end{itemize}
 Profiles are transformed into a human\texttt{\symbol{45}}readable form with
'profiling' package, for example with the 'OutputAnnotatedCodeCoverageFiles'
function. }

 

\subsection{\textcolor{Chapter }{ProfileLineByLine}}
\logpage{[ 7, 8, 14 ]}\nobreak
\hyperdef{L}{X86557887796F66FA}{}
{\noindent\textcolor{FuncColor}{$\triangleright$\enspace\texttt{ProfileLineByLine({\mdseries\slshape filename[, options]})\index{ProfileLineByLine@\texttt{ProfileLineByLine}}
\label{ProfileLineByLine}
}\hfill{\scriptsize (function)}}\\


 \texttt{ProfileLineByLine} begins GAP recording profiling data to the file \mbox{\texttt{\mdseries\slshape filename}}. This file will get *very* large very quickly. This file is compressed using
gzip to reduce its size. \mbox{\texttt{\mdseries\slshape options}} is an optional dictionary, which sets various configuration options. These are 
\begin{description}
\item[{coverage}]  Boolean (defaults to false). If this is enabled, only information about which
lines are read and executed is stored. Enabling this is the same as calling \texttt{CoverageLineByLine} (\ref{CoverageLineByLine}). Using this ignores all other options.
\item[{wallTime}]  Boolean (defaults to true). Sets if time should be measured using
wall\texttt{\symbol{45}}clock time (true) or CPU time (false). (measuring
CPU\texttt{\symbol{45}}time has a higher overhead). 
\item[{recordMem}]  Boolean (defaults to false). Instead of recording the CPU time taken by
statements, record the total size of all new objects created by each line. 
\item[{resolution}]  Integer (defaults to 0). By default profiling will record a trace of all
executed code. When \mbox{\texttt{\mdseries\slshape resolution}} non\texttt{\symbol{45}}zero, GAP instead samples which piece of code is being
executed every \mbox{\texttt{\mdseries\slshape resolution}} nanoseconds. Increasing this improves performance and produces smaller traces,
at the cost of accuracy. GAP will still accurately record which statements are
executed at least once.
\end{description}
 }

 

\subsection{\textcolor{Chapter }{CoverageLineByLine}}
\logpage{[ 7, 8, 15 ]}\nobreak
\hyperdef{L}{X87CC48807DB4C008}{}
{\noindent\textcolor{FuncColor}{$\triangleright$\enspace\texttt{CoverageLineByLine({\mdseries\slshape filename})\index{CoverageLineByLine@\texttt{CoverageLineByLine}}
\label{CoverageLineByLine}
}\hfill{\scriptsize (function)}}\\


 \texttt{CoverageLineByLine} begins GAP recording code coverage to the file \mbox{\texttt{\mdseries\slshape filename}}. This is equivalent to calling \texttt{ProfileLineByLine} (\ref{ProfileLineByLine}) with coverage=true. }

 

\subsection{\textcolor{Chapter }{UnprofileLineByLine}}
\logpage{[ 7, 8, 16 ]}\nobreak
\hyperdef{L}{X7C5DED9C7CC77504}{}
{\noindent\textcolor{FuncColor}{$\triangleright$\enspace\texttt{UnprofileLineByLine({\mdseries\slshape })\index{UnprofileLineByLine@\texttt{UnprofileLineByLine}}
\label{UnprofileLineByLine}
}\hfill{\scriptsize (function)}}\\


 Stops profiling which was previously started with \texttt{ProfileLineByLine} (\ref{ProfileLineByLine}) or \texttt{CoverageLineByLine} (\ref{CoverageLineByLine}). }

 

\subsection{\textcolor{Chapter }{UncoverageLineByLine}}
\logpage{[ 7, 8, 17 ]}\nobreak
\hyperdef{L}{X7B705B2D8670A9C5}{}
{\noindent\textcolor{FuncColor}{$\triangleright$\enspace\texttt{UncoverageLineByLine({\mdseries\slshape })\index{UncoverageLineByLine@\texttt{UncoverageLineByLine}}
\label{UncoverageLineByLine}
}\hfill{\scriptsize (function)}}\\


 Stops profiling which was previously started with \texttt{ProfileLineByLine} (\ref{ProfileLineByLine}) or \texttt{CoverageLineByLine} (\ref{CoverageLineByLine}). }

 

\subsection{\textcolor{Chapter }{IsLineByLineProfileActive}}
\logpage{[ 7, 8, 18 ]}\nobreak
\hyperdef{L}{X7823C83D79B36D3B}{}
{\noindent\textcolor{FuncColor}{$\triangleright$\enspace\texttt{IsLineByLineProfileActive({\mdseries\slshape })\index{IsLineByLineProfileActive@\texttt{IsLineByLineProfileActive}}
\label{IsLineByLineProfileActive}
}\hfill{\scriptsize (function)}}\\


 \texttt{IsLineByLineProfileActive} returns if line\texttt{\symbol{45}}by\texttt{\symbol{45}}line profiling is
currently activated. }

 

\subsection{\textcolor{Chapter }{DisplayCacheStats}}
\logpage{[ 7, 8, 19 ]}\nobreak
\hyperdef{L}{X83D8A42B7BB92F5B}{}
{\noindent\textcolor{FuncColor}{$\triangleright$\enspace\texttt{DisplayCacheStats({\mdseries\slshape })\index{DisplayCacheStats@\texttt{DisplayCacheStats}}
\label{DisplayCacheStats}
}\hfill{\scriptsize (function)}}\\


 displays statistics about the different caches used by the method selection. }

 

\subsection{\textcolor{Chapter }{ClearCacheStats}}
\logpage{[ 7, 8, 20 ]}\nobreak
\hyperdef{L}{X79C58704838232CC}{}
{\noindent\textcolor{FuncColor}{$\triangleright$\enspace\texttt{ClearCacheStats({\mdseries\slshape })\index{ClearCacheStats@\texttt{ClearCacheStats}}
\label{ClearCacheStats}
}\hfill{\scriptsize (function)}}\\


 clears all statistics about the different caches used by the method selection. }

 }

  
\section{\textcolor{Chapter }{Information about the version used}}\label{Information about the version used}
\logpage{[ 7, 9, 0 ]}
\hyperdef{L}{X7EE874867C0BEEDD}{}
{
  \index{GAPInfo.Version@\texttt{GAPInfo.Version}} The global variable \texttt{GAPInfo.Version} (see \texttt{GAPInfo} (\ref{GAPInfo})) contains the version number of the version of \textsf{GAP}. Its value can be checked other version number using \texttt{CompareVersionNumbers} (\ref{CompareVersionNumbers}). 

 To produce sample citations for the used version of \textsf{GAP} or for a package available in this \textsf{GAP} installation, use \texttt{Cite} (\ref{Cite}). 

 If you wish to report a problem to \textsf{GAP} Support or \textsf{GAP} Forum, it may be useful to not only report the version used, but also to
include the \textsf{GAP} banner displays the information about the architecture for which the \textsf{GAP} binary is built, used libraries and loaded packages. }

  
\section{\textcolor{Chapter }{Test Files}}\label{Test Files}
\logpage{[ 7, 10, 0 ]}
\hyperdef{L}{X801051CC86594630}{}
{
  Test files are used to check that \textsf{GAP} produces correct results in certain computations. A selection of test files
for the library can be found in the \texttt{tst} directory of the \textsf{GAP} distribution. 
\subsection{\textcolor{Chapter }{Starting and stopping test}}\logpage{[ 7, 10, 1 ]}
\hyperdef{L}{X8213757B7ACC76E6}{}
{
\noindent\textcolor{FuncColor}{$\triangleright$\enspace\texttt{START{\textunderscore}TEST({\mdseries\slshape name})\index{START{\textunderscore}TEST@\texttt{START{\textunderscore}TEST}}
\label{STARTuScoreTEST}
}\hfill{\scriptsize (function)}}\\
\noindent\textcolor{FuncColor}{$\triangleright$\enspace\texttt{STOP{\textunderscore}TEST({\mdseries\slshape name})\index{STOP{\textunderscore}TEST@\texttt{STOP{\textunderscore}TEST}}
\label{STOPuScoreTEST}
}\hfill{\scriptsize (function)}}\\


 \texttt{START{\textunderscore}TEST} and \texttt{STOP{\textunderscore}TEST} may be optionally used in files that are read via \texttt{Test} (\ref{Test}). If used, \texttt{START{\textunderscore}TEST} reinitializes 
\begin{itemize}
\item  the caches (see \texttt{FlushCaches} (\ref{FlushCaches})), 
\item  the global random number generators (see \texttt{GlobalMersenneTwister} (\ref{GlobalMersenneTwister})), and 
\item  the global options stack (see Chapter \ref{Options Stack}), 
\end{itemize}
 in order to be independent of the reading order of several test files. 

 Furthermore, the assertion level (see{\nobreakspace}\texttt{Assert} (\ref{Assert})) is set to $2$ (if it was lower before) by \texttt{START{\textunderscore}TEST} and set back to the previous value in the subsequent \texttt{STOP{\textunderscore}TEST} call. 

 Also \texttt{STOP{\textunderscore}TEST} reinitializes the global options stack, otherwise it could happen that options
which got set during the tests are not reset after the tests. 

 To use this feature, a test file should be started with a line 

 
\begin{Verbatim}[commandchars=!@|,fontsize=\small,frame=single,label=Example]
  !gapprompt@gap>| !gapinput@START_TEST( "arbitrary identifier string" );|
\end{Verbatim}
 

 (Note that the \texttt{gap{\textgreater} } prompt is part of the line!) 

 and should be finished with a line 

 
\begin{Verbatim}[commandchars=!@|,fontsize=\small,frame=single,label=Example]
  !gapprompt@gap>| !gapinput@STOP_TEST( "same identifier string as for START_TEST" );|
\end{Verbatim}
 }

 If you want to run a quick test of your \textsf{GAP} installation (though this is not required), you can read in a test script that
exercises some \textsf{GAP}'s capabilities. 

 
\begin{Verbatim}[commandchars=!@|,fontsize=\small,frame=single,label=Example]
  !gapprompt@gap>| !gapinput@Read( Filename( DirectoriesLibrary( "tst" ), "testinstall.g" ) );|
\end{Verbatim}
 

 
\begin{Verbatim}[commandchars=!@|,fontsize=\small,frame=single,label=Example]
  test file         time(msec)
  -------------------------------------------
  testing: ................/gap4r5/tst/zlattice.tst
  zlattice.tst               0
  testing: ................/gap4r5/tst/gaussian.tst
  gaussian.tst              10
  [ further lines deleted ]
\end{Verbatim}
 If you want to run a more advanced check (this is not required and make take
up to an hour), you can read \texttt{teststandard.g} which is an extended test script performing all tests from the \texttt{tst} directory. 

 
\begin{Verbatim}[commandchars=!@|,fontsize=\small,frame=single,label=Example]
  !gapprompt@gap>| !gapinput@Read( Filename( DirectoriesLibrary( "tst" ), "teststandard.g" ) );|
\end{Verbatim}
 

 

 

\subsection{\textcolor{Chapter }{Test}}
\logpage{[ 7, 10, 2 ]}\nobreak
\hyperdef{L}{X87712F9D8732193C}{}
{\noindent\textcolor{FuncColor}{$\triangleright$\enspace\texttt{Test({\mdseries\slshape fname[, optrec]})\index{Test@\texttt{Test}}
\label{Test}
}\hfill{\scriptsize (function)}}\\
\textbf{\indent Returns:\ }
\texttt{true} or \texttt{false}.



 The argument \mbox{\texttt{\mdseries\slshape fname}} must be the name of a file or an open input stream. The content of this file
or stream should contain \textsf{GAP} input and output. The function \texttt{Test} runs the input lines, compares the actual output with the output stored in \mbox{\texttt{\mdseries\slshape fname}} and reports differences. With an optional record as argument \mbox{\texttt{\mdseries\slshape optrec}} details of this process can be adjusted. Note that the \texttt{rewriteToFile} option is especially useful for generating test files. 

 More precisely, the content of \mbox{\texttt{\mdseries\slshape fname}} must have the following format. \\
 Lines starting with \texttt{"gap{\textgreater}{\nobreakspace}"} are considered as \textsf{GAP} input, they can be followed by lines starting with \texttt{"{\textgreater}{\nobreakspace}"} if the input is continued over several lines. \\
 To allow for comments in \mbox{\texttt{\mdseries\slshape fname}} the following lines are ignored by default: lines at the beginning of \mbox{\texttt{\mdseries\slshape fname}} that start with \texttt{"\#"} or are empty, and one empty line together with one or more lines starting with \texttt{"\#"}.\\
 All other lines are considered as \textsf{GAP} output from the preceding \textsf{GAP} input. 

 Lines which begin "\#@" define special configuration options for tests. The \texttt{\#@local} and \texttt{\#@exec} options can only be used before any \textsf{GAP} input, and the other commands can only be used between individual tests (just
before a line starting \texttt{gap{\textgreater}}, or at end of the file). Currently defined options are: 
\begin{description}
\item[{\#@local identifierlist}] Run all the tests in the input as if it is in a function with local variable
list \texttt{identifierlist}, which is a comma\texttt{\symbol{45}}separated list of identifiers. Multiple
\#@local lines may be used. These lines should \emph{not} end with a comma or semicolon. If this option is used then an error will occur
unless \emph{all} the variables used are included in the local list. 

 As an example, the \textsf{Utils} package has a test file \texttt{tst/iterator.tst} which starts with the lines: 
\begin{Verbatim}[commandchars=!|A,fontsize=\small,frame=single,label=Example]
  #@local  c3c3, cart, G, h, it1, it2, iter, iter0, iter4, iterL
  #@local  L, n, pairs0, pairs4, pairsL, s3, s4
\end{Verbatim}
 
\item[{\#@exec gapcode}] Execute the code \texttt{gapcode} before any test in the input is run. This allows defining global variables
when using \texttt{\#@local}. 
\item[{\#@if EXPR ... \texttt{\symbol{123}}\#@elif EXPR ...\texttt{\symbol{125}}
[\#@else ...] \#@fi}] A \texttt{\#@if} allows to conditionally skip parts of the test input depending on the value of
a boolean expression. The exact behavior is done as follows: 

 If the first \textsf{GAP} expression \texttt{EXPR} evaluates to \texttt{true}, then the lines after the \texttt{\#@if} are used until either a \texttt{\#@elif}, \texttt{\#@else} or \texttt{\#@fi} is reached. If an \texttt{\#@elif} is present, then the lines after \texttt{\#@elif} are used if and only if its \texttt{EXPR} evaluates to \texttt{true} and all previous \texttt{\#@if} and \texttt{\#@elif} clauses had expressions evaluating to \texttt{false}. In this case the lines after \texttt{\#@elif} are used until either a \texttt{\#@elif}, \texttt{\#@else} or \texttt{\#@fi} is reached. If an \texttt{\#@else} is present then the lines after the \texttt{\#@else} are used if and only if \texttt{EXPR} evaluated to \texttt{false} in all \texttt{\#@if} and \texttt{\#@elif} clauses. Finally, once \texttt{\#fi} is reached, evaluation continues normally. 

 Note that each \texttt{EXPR} is evaluated after all \texttt{\#@exec} lines have been executed but before any tests are run. Thus, it cannot depend
on test results or packages loaded in tests, but it can depend on packages
loaded via \texttt{\#@exec}. 

 In addition \texttt{\#@if} clauses cannot be nested within each other. 

 As an example, the \textsf{GAP} test suite contains the test file \texttt{tst/testinstall/pperm.tst} which contains the lines: 
\begin{Verbatim}[commandchars=!|D,fontsize=\small,frame=single,label=Example]
  #@if GAPInfo.BytesPerVariable = 8
  !gapprompt|gap>D !gapinput|HASH_FUNC_FOR_PPERM(f, 10 ^ 6) in [260581, 402746];D
  true
  #@else
  !gapprompt|gap>D !gapinput|HASH_FUNC_FOR_PPERM(f, 10 ^ 6);D
  953600
  #@fi
\end{Verbatim}
 
\end{description}
 By default the actual \textsf{GAP} output is compared exactly with the stored output, and if these are different
some information about the differences is printed. 

 If any differences are found then \texttt{Test} returns \texttt{false}, otherwise \texttt{true}. 

 If the optional argument \mbox{\texttt{\mdseries\slshape optrec}} is given it must be a record. The following components of \mbox{\texttt{\mdseries\slshape optrec}} are recognized and can change the default behaviour of \texttt{Test}: 
\begin{description}
\item[{\texttt{ignoreComments}}] If set to \texttt{false} then no lines in \mbox{\texttt{\mdseries\slshape fname}} are ignored as explained above (default is \texttt{true}).
\item[{\texttt{width}}] The screen width used for the new output (default is \texttt{80}). 
\item[{\texttt{compareFunction}}] This must be a function that gets two strings as input, the newly generated
and the stored output of some \textsf{GAP} input. The function must return \texttt{true} or \texttt{false}, indicating if the strings should be considered equivalent or not. By default \texttt{\texttt{\symbol{92}}=} (\ref{bSlash=}) is used. \\
 Two strings are recognized as abbreviations in this component: \texttt{"uptowhitespace"} checks if the two strings become equal after removing all white space. And \texttt{"uptonl"} compares the string up to trailing newline characters. 
\item[{\texttt{transformFunction}}] This must be a function that gets one string as input, either the newly
generated or the stored output of some \textsf{GAP} input. The function must return a new string which will be used to compare the
actual and the expected output. By default \texttt{IdFunc} (\ref{IdFunc}) is used. \\
 Two strings are recognized as abbreviations in this component: \texttt{"removewhitespace"} removes all white space. And \texttt{"removenl"} removes all trailing newline characters. 
\item[{\texttt{reportDiff}}] A function that gets six arguments and reports a difference in the output: the \textsf{GAP} input, the expected \textsf{GAP} output, the newly generated output, the name of tested file, the line number
of the input, the time to run the input. (The default is demonstrated in the
example below.)
\item[{\texttt{rewriteToFile}}] If this is bound to a string it is considered as a file name and that file is
written with the same input and comment lines as \mbox{\texttt{\mdseries\slshape fname}} but the output substituted by the newly generated version; if it is bound to \texttt{true}, then this is treated as if it was bound to \mbox{\texttt{\mdseries\slshape fname}} (default is \texttt{false}). This is especially useful for generating test files because it ensures that
the test files are formatted exactly as \texttt{Test} expects them to be. 
\item[{\texttt{writeTimings}}] If this is bound to a string it is considered as a file name, that file is
written and contains timing information for each input in \mbox{\texttt{\mdseries\slshape fname}}. 
\item[{\texttt{compareTimings}}] If this is bound to a string it is considered as name of a file to which
timing information was stored via \texttt{writeTimings} in a previous call. The new timings are compared to the stored ones. By
default only commands which take more than a threshold of 100 milliseconds are
considered, and only differences of more than 20\% are considered significant.
These defaults can be overwritten by assigning a list \texttt{[timingfile, threshold, percentage]} to this component. (The default of \texttt{compareTimings} is \texttt{false}.)
\item[{\texttt{reportTimeDiff}}] This component can be used to overwrite the default function to display timing
differences. It must be a function with 5 arguments: \textsf{GAP} input, name of test file, line number, stored time, new time. 
\item[{\texttt{ignoreSTOP{\textunderscore}TEST}}] By default set to \texttt{true}, in that case the output of \textsf{GAP} input starting with \texttt{"STOP{\textunderscore}TEST"} is not checked.
\item[{\texttt{showProgress}}] If this is \texttt{true} then \textsf{GAP} prints position information and the input line before it is processed; if set
to \texttt{"some"}, then GAP shows the current line number of the test being processed; if set
to \texttt{false}, no progress updates are displayed (default is \texttt{"some"} if GAP's output goes to a terminal, otherwise \texttt{false}). 
\item[{\texttt{subsWindowsLineBreaks}}] If this is \texttt{true} then \textsf{GAP} substitutes DOS/Windows style line breaks
"\texttt{\symbol{92}}r\texttt{\symbol{92}}n" by UNIX style line breaks
"\texttt{\symbol{92}}n" after reading the test file. (default is \texttt{true}).
\item[{\texttt{returnNumFailures}}] If this is \texttt{true} then \textsf{GAP} returns the number of input lines of the test file which had differences in
their output, instead of returning \texttt{true} or \texttt{false}.
\end{description}
 
\begin{Verbatim}[commandchars=!@A,fontsize=\small,frame=single,label=Example]
  !gapprompt@gap>A !gapinput@tnam := Filename(DirectoriesLibrary(), "../doc/ref/demo.tst");;A
  !gapprompt@gap>A !gapinput@mask := function(str) return Concatenation("| ",A
  !gapprompt@>A !gapinput@         JoinStringsWithSeparator(SplitString(str, "\n", ""), "\n| "),A
  !gapprompt@>A !gapinput@         "\n"); end;;A
  !gapprompt@gap>A !gapinput@Print(mask(StringFile(tnam)));A
  | # this is a demo file for the 'Test' function
  | #
  | gap> g := Group((1,2), (1,2,3));
  | Group([ (1,2), (1,2,3) ])
  |
  | # another comment following an empty line
  | # the following fails:
  | gap> a := 13+29;
  | 41
  !gapprompt@gap>A !gapinput@ss := InputTextString(StringFile(tnam));;A
  !gapprompt@gap>A !gapinput@Test(ss);A
  ########> Diff in test stream, line 8:
  # Input is:
  a := 13+29;
  # Expected output:
  41
  # But found:
  42
  ########
  false
  !gapprompt@gap>A !gapinput@RewindStream(ss);A
  true
  !gapprompt@gap>A !gapinput@dtmp := DirectoryTemporary();;A
  !gapprompt@gap>A !gapinput@ftmp := Filename(dtmp,"demo.tst");;A
  !gapprompt@gap>A !gapinput@Test(ss, rec(reportDiff := Ignore, rewriteToFile := ftmp));A
  false
  !gapprompt@gap>A !gapinput@Test(ftmp);A
  true
  !gapprompt@gap>A !gapinput@Print(mask(StringFile(ftmp)));A
  | # this is a demo file for the 'Test' function
  | #
  | gap> g := Group((1,2), (1,2,3));
  | Group([ (1,2), (1,2,3) ])
  |
  | # another comment following an empty line
  | # the following fails:
  | gap> a := 13+29;
  | 42
\end{Verbatim}
 }

 

\subsection{\textcolor{Chapter }{TestDirectory}}
\logpage{[ 7, 10, 3 ]}\nobreak
\hyperdef{L}{X87AF67528799481F}{}
{\noindent\textcolor{FuncColor}{$\triangleright$\enspace\texttt{TestDirectory({\mdseries\slshape inlist[, optrec]})\index{TestDirectory@\texttt{TestDirectory}}
\label{TestDirectory}
}\hfill{\scriptsize (function)}}\\
\textbf{\indent Returns:\ }
\texttt{true} or \texttt{false}.



 The argument \mbox{\texttt{\mdseries\slshape inlist}} must be either a single filename or directory name, or a list of filenames and
directories. The function \texttt{TestDirectory} will create a list of files to be tested by taking any files in \mbox{\texttt{\mdseries\slshape inlist}}, and recursively searching any directories in \mbox{\texttt{\mdseries\slshape inlist}} for files ending in \texttt{.tst}. Each of these files is then run through \texttt{Test} (\ref{Test}), and the results printed, and \texttt{true} returned if all tests passed. 

 If the optional argument \mbox{\texttt{\mdseries\slshape optrec}} is given it must be a record. Note that the \texttt{rewriteToFile} option is especially useful for generating test files. The following
components of \mbox{\texttt{\mdseries\slshape optrec}} are recognized and can change the default behaviour of \texttt{TestDirectory}: 
\begin{description}
\item[{\texttt{testOptions}}] A record which will be passed on as the second argument of \texttt{Test} (\ref{Test}) if present.
\item[{\texttt{earlyStop}}] If \texttt{true}, stop as soon as any \texttt{Test} (\ref{Test}) fails (defaults to \texttt{false}). 
\item[{\texttt{showProgress}}] Print information about how tests are progressing (defaults to \texttt{"some"} if GAP's output goes to a terminal, otherwise \texttt{false}). 
\item[{\texttt{suppressStatusMessage}}] suppress displaying status messages \texttt{\#I Errors detected while testing} and \texttt{\#I No errors detected while testing} after the test (defaults to \texttt{false}). 
\item[{\texttt{rewriteToFile}}] If \texttt{true}, then rewrite each test file to disc, with the output substituted by the
results of running the test (defaults to \texttt{false}). This is especially useful for generating test files because it ensures that
the test files are formatted exactly as \texttt{Test} (\ref{Test}) expects them to be. 
\item[{\texttt{exclude}}] A list of file and directory names which will be excluded from testing
(defaults to \texttt{[]}). 
\item[{\texttt{exitGAP}}] Rather than returning \texttt{true} or \texttt{false}, exit GAP with the return value of GAP set to success or fail, depending on
if all tests passed (defaults to \texttt{false}). 
\end{description}
 }

 See also \texttt{TestPackage} (\ref{TestPackage}) for the information on running standard tests for \textsf{GAP} packages. }

  
\section{\textcolor{Chapter }{Debugging Recursion}}\label{Debugging Recursion}
\logpage{[ 7, 11, 0 ]}
\hyperdef{L}{X85FF55448787CCA0}{}
{
  The \textsf{GAP} interpreter monitors the level of nesting of \textsf{GAP} functions during execution. By default, whenever this nesting reaches a
multiple of $5000$, \textsf{GAP} enters a break loop (\ref{Break Loops}) allowing you to terminate the calculation, or enter \textsc{Return}\texttt{;} to continue it. 

 
\begin{Verbatim}[commandchars=!|A,fontsize=\small,frame=single,label=Example]
  !gapprompt|gap>A !gapinput|dive:= function(depth) if depth>1 then dive(depth-1); fi; return; end;A
  function( depth ) ... end
  !gapprompt|gap>A !gapinput|dive(100);A
  !gapprompt|gap>A !gapinput|OnBreak:= function() Where(2); end;; # shorter tracebackA
  !gapprompt|gap>A !gapinput|dive(6000);A
  Error, recursion depth trap (5000)
  Stack trace:
  *[1] dive( depth - 1 );
     @ *stdin*:1
   [2] dive( depth - 1 );
     @ *stdin*:1
  ...  at *stdin*:4
  you can enter 'quit;' to quit to outer loop, or
  you can enter 'return;' to continue
  !gapbrkprompt|brk>A !gapinput|return;A
  !gapprompt|gap>A !gapinput|dive(11000);A
  Error, recursion depth trap (5000)
  Stack trace:
  *[1] dive( depth - 1 );
     @ *stdin*:1
   [2] dive( depth - 1 );
     @ *stdin*:1
  ...  at *stdin*:5
  you can enter 'quit;' to quit to outer loop, or
  you can enter 'return;' to continue
  !gapbrkprompt|brk>A !gapinput|return;A
  Error, recursion depth trap (10000)
  Stack trace:
  *[1] dive( depth - 1 );
     @ *stdin*:1
   [2] dive( depth - 1 );
     @ *stdin*:1
  ...  at *stdin*:5
  you can enter 'quit;' to quit to outer loop, or
  you can enter 'return;' to continue
  !gapbrkprompt|brk>A !gapinput|return;A
  gap>
\end{Verbatim}
 

 This behaviour can be controlled using the following procedures. 

\subsection{\textcolor{Chapter }{SetRecursionTrapInterval}}
\logpage{[ 7, 11, 1 ]}\nobreak
\hyperdef{L}{X7D8968FC7E24A4E5}{}
{\noindent\textcolor{FuncColor}{$\triangleright$\enspace\texttt{SetRecursionTrapInterval({\mdseries\slshape interval})\index{SetRecursionTrapInterval@\texttt{SetRecursionTrapInterval}}
\label{SetRecursionTrapInterval}
}\hfill{\scriptsize (function)}}\\
\noindent\textcolor{FuncColor}{$\triangleright$\enspace\texttt{GetRecursionDepth({\mdseries\slshape })\index{GetRecursionDepth@\texttt{GetRecursionDepth}}
\label{GetRecursionDepth}
}\hfill{\scriptsize (function)}}\\


 \texttt{GetRecursionDepth} returns the nesting level of the GAP interpreter. This is reset to 0 every
time the break loop is entered. \texttt{SetRecursionTrapInterval} sets the depth of the stack at which GAP will enter the Break loop. \mbox{\texttt{\mdseries\slshape interval}} must be a non\texttt{\symbol{45}}negative small integer (between 0 and $2^{28}$). An \mbox{\texttt{\mdseries\slshape interval}} of 0 suppresses the monitoring of recursion altogether. In this case excessive
recursion may cause \textsf{GAP} to crash. 

 
\begin{Verbatim}[commandchars=!|A,fontsize=\small,frame=single,label=Example]
  !gapprompt|gap>A !gapinput|GetRecursionDepth();A
  0
  !gapprompt|gap>A !gapinput|dive := function(depth)A
  !gapprompt|>A !gapinput| if depth>1 thenA
  !gapprompt|>A !gapinput|   dive(depth-1);A
  !gapprompt|>A !gapinput| elseA
  !gapprompt|>A !gapinput|   Print("Depth ", GetRecursionDepth());A
  !gapprompt|>A !gapinput| fi;A
  !gapprompt|>A !gapinput|end;;A
  !gapprompt|gap>A !gapinput|SetRecursionTrapInterval(1000);A
  !gapprompt|gap>A !gapinput|OnBreak:= function() Where(2); end;; # shorter tracebackA
  !gapprompt|gap>A !gapinput|dive(100);A
  Depth 100
  !gapprompt|gap>A !gapinput|dive(2500);A
  Error, recursion depth trap (1000)
  Stack trace:
  *[1] dive( depth - 1 );
     @ *stdin*:4
   [2] dive( depth - 1 );
     @ *stdin*:4
  ...  at *stdin*:12
  you can enter 'quit;' to quit to outer loop, or
  you can enter 'return;' to continue
  !gapbrkprompt|brk>A !gapinput|return;A
  Error, recursion depth trap (2000)
  Stack trace:
  *[1] dive( depth - 1 );
     @ *stdin*:4
   [2] dive( depth - 1 );
     @ *stdin*:4
  ...  at *stdin*:12
  you can enter 'quit;' to quit to outer loop, or
  you can enter 'return;' to continue
  !gapbrkprompt|brk>A !gapinput|GetRecursionDepth();A
  0
  !gapbrkprompt|brk>A !gapinput|return;A
  !gapprompt|gap>A !gapinput|SetRecursionTrapInterval(2^50);A
  !gapprompt|gap>A !gapinput|dive(20000);A
  Depth 20000
  !gapprompt|gap>A !gapinput|dive(2000000);A
  Segmentation fault
\end{Verbatim}
 }

 }

  
\section{\textcolor{Chapter }{Global Memory Information}}\label{Global Memory Information}
\logpage{[ 7, 12, 0 ]}
\hyperdef{L}{X85679F17791D9B63}{}
{
  
\subsection{\textcolor{Chapter }{Garbage Collection}}\label{Garbage Collection}
\logpage{[ 7, 12, 1 ]}
\hyperdef{L}{X7F1F741D7F0899D1}{}
{
  \index{GASMAN@\texttt{GASMAN}} The \textsf{GAP} environment provides automatic memory management, so that the programmer does
not need to concern themselves with allocating space for objects, or
recovering space when objects are no longer needed. The memory manager that
shall be used by \textsf{GAP} is specified at compile time. One of the choices is called \texttt{GASMAN} (\textsf{GAP} Storage MANager). (The name of the currently used garbage collector is stored
in the variable \texttt{GAPInfo.KernelInfo.GC}.) 

 If \textsf{GAP} uses \texttt{GASMAN} then messages reporting garbage collections performed by \texttt{GASMAN} can be switched on by the \texttt{\texttt{\symbol{45}}g} command line option (see section \ref{Command Line Options}). There are also some facilities to access information from \texttt{GASMAN} from \textsf{GAP} programs, see below. }

 

\subsection{\textcolor{Chapter }{CollectGarbage}}
\logpage{[ 7, 12, 2 ]}\nobreak
\hyperdef{L}{X7848AB367F3A1221}{}
{\noindent\textcolor{FuncColor}{$\triangleright$\enspace\texttt{CollectGarbage({\mdseries\slshape full})\index{CollectGarbage@\texttt{CollectGarbage}}
\label{CollectGarbage}
}\hfill{\scriptsize (function)}}\\
\textbf{\indent Returns:\ }
 nothing. 



 This function forces a garbage collection. If \mbox{\texttt{\mdseries\slshape full}} is \texttt{true} then it triggers a full garbage collection, otherwise a partial one. 

 \textsf{GAP} invokes its garbage collector automatically, thus there is normally no need to
call \texttt{CollectGarbage}. 

 The function \texttt{CollectGarbage} was introduced in \textsf{GAP}{\nobreakspace}4.12. In older \textsf{GAP} versions, one can use \texttt{GASMAN( "collect" )} (if \mbox{\texttt{\mdseries\slshape full}} is \texttt{true}) or \texttt{GASMAN( "partial" )} (if \mbox{\texttt{\mdseries\slshape full}} is not \texttt{true}) instead. 

 
\begin{Verbatim}[commandchars=!@|,fontsize=\small,frame=single,label=Example]
  !gapprompt@gap>| !gapinput@CollectGarbage( false );|
  !gapprompt@gap>| !gapinput@CollectGarbage( true );|
\end{Verbatim}
 }

 

\subsection{\textcolor{Chapter }{GasmanStatistics}}
\logpage{[ 7, 12, 3 ]}\nobreak
\hyperdef{L}{X836977DE80416F3D}{}
{\noindent\textcolor{FuncColor}{$\triangleright$\enspace\texttt{GasmanStatistics({\mdseries\slshape })\index{GasmanStatistics@\texttt{GasmanStatistics}}
\label{GasmanStatistics}
}\hfill{\scriptsize (function)}}\\


 This function is meaningful only if \texttt{GASMAN} is the garbage collector used by \textsf{GAP}, see Section \ref{Garbage Collection}. 

 \texttt{GasmanStatistics} returns a record containing some information from the garbage collection
mechanism. The record may contain up to four components: \texttt{full}, \texttt{partial}, \texttt{npartial}, and \texttt{nfull}. 

 The \texttt{full} component will be present if a full garbage collection has taken place since \textsf{GAP} started. It contains information about the most recent full garbage
collection. It is a record, with eight components: \texttt{livebags} contains the number of bags which survived the garbage collection; \texttt{livekb} contains the total number of kilobytes occupied by those bags; \texttt{deadbags} contains the total number of bags which were reclaimed by that garbage
collection and all the partial garbage collections preceding it, since the
previous full garbage collection; \texttt{deadkb} contains the total number of kilobytes occupied by those bags; \texttt{freekb} reports the total number of kilobytes available in the \textsf{GAP} workspace for new objects; \texttt{totalkb} reports the actual size of the workspace; \texttt{time} reports the CPU time in milliseconds spent on the last garbage collection and \texttt{cumulative} the total CPU time in milliseconds spent on that type of garbage collection
since \textsf{GAP} started. 

 These figures should be viewed with some caution. They are stored internally
in fixed length integer formats, and \texttt{deadkb} and \texttt{deadbags} are liable to overflow if there are many partial collections before a full
collection. Also, note that \texttt{livekb} and \texttt{freekb} will not usually add up to \texttt{totalkb}. The difference is essentially the space overhead of the memory management
system. 

 The \texttt{partial} component will be present if there has been a partial garbage collection since
the last full one. It is also a record with the same six components as \texttt{full}. In this case \texttt{deadbags} and \texttt{deadkb} refer only to the number and total size of the garbage bags reclaimed in this
partial garbage collection and \texttt{livebags}and \texttt{livekb} only to the numbers and total size of the young bags that were considered for
garbage collection, and survived. 

 The \texttt{npartial} and \texttt{nfull} components will contain the number of full and partial garbage collections
performed since \textsf{GAP} started. }

 

\subsection{\textcolor{Chapter }{GasmanMessageStatus}}
\logpage{[ 7, 12, 4 ]}\nobreak
\hyperdef{L}{X85327FA5872E0356}{}
{\noindent\textcolor{FuncColor}{$\triangleright$\enspace\texttt{GasmanMessageStatus({\mdseries\slshape })\index{GasmanMessageStatus@\texttt{GasmanMessageStatus}}
\label{GasmanMessageStatus}
}\hfill{\scriptsize (function)}}\\
\noindent\textcolor{FuncColor}{$\triangleright$\enspace\texttt{SetGasmanMessageStatus({\mdseries\slshape stat})\index{SetGasmanMessageStatus@\texttt{SetGasmanMessageStatus}}
\label{SetGasmanMessageStatus}
}\hfill{\scriptsize (function)}}\\


 This function is meaningful only if \texttt{GASMAN} is the garbage collector used by \textsf{GAP}, see Section \ref{Garbage Collection}. 

 \texttt{GasmanMessageStatus} returns one of the strings \texttt{"none"}, \texttt{"full"}, or \texttt{"all"}, depending on whether the garbage collector is currently set to print
messages on no collections, full collections only, or all collections,
respectively. 

 Calling \texttt{SetGasmanMessageStatus} with the argument \mbox{\texttt{\mdseries\slshape stat}}, which should be one of the three strings mentioned above, sets the garbage
collector messaging level. }

 

\subsection{\textcolor{Chapter }{GasmanLimits}}
\logpage{[ 7, 12, 5 ]}\nobreak
\hyperdef{L}{X80C683247E94769F}{}
{\noindent\textcolor{FuncColor}{$\triangleright$\enspace\texttt{GasmanLimits({\mdseries\slshape })\index{GasmanLimits@\texttt{GasmanLimits}}
\label{GasmanLimits}
}\hfill{\scriptsize (function)}}\\


 This function is meaningful only if \texttt{GASMAN} is the garbage collector used by \textsf{GAP}, see Section \ref{Garbage Collection}. 

 \texttt{GasmanLimits} returns a record with three components: \texttt{min} is the minimum workspace size as set by the \texttt{\texttt{\symbol{45}}m} command line option in kilobytes. The workspace size will never be reduced
below this by the garbage collector. \texttt{max} is the maximum workspace size, as set by the \texttt{\texttt{\symbol{45}}o} command line option, also in kilobytes. If the workspace would need to grow
past this point, \textsf{GAP} will enter a break loop to warn the user. A value of 0 indicates no limit. \texttt{kill} is the absolute maximum, set by the \texttt{\texttt{\symbol{45}}K} command line option. The workspace will never be allowed to grow past this
limit. }

 }

 }

            
\chapter{\textcolor{Chapter }{Options Stack}}\label{Options Stack}
\logpage{[ 8, 0, 0 ]}
\hyperdef{L}{X7FD84061873F72A2}{}
{
  \textsf{GAP} supports a \emph{global options system}. This is intended as a way for the user to provide guidance to various
algorithms that might be used in a computation. Such guidance should not
change mathematically the specification of the computation to be performed,
although it may change the algorithm used. A typical example is the selection
of a strategy for the Todd\texttt{\symbol{45}}Coxeter coset enumeration
procedure. An example of something not suited to the options mechanism is the
imposition of exponent laws in the $p$\texttt{\symbol{45}}Quotient algorithm. 

 The basis of this system is a global stack of records. All the entries of each
record are thought of as options settings, and the effective setting of an
option is given by the topmost record in which the relevant field is bound. 

 The reason for the choice of a stack is the intended pattern of use: 

 \texttt{PushOptions( rec( \mbox{\texttt{\mdseries\slshape stuff}} ) );} 

 \texttt{DoSomething( \mbox{\texttt{\mdseries\slshape args}} );} 

 \texttt{PopOptions();} 

 This can be abbreviated, to \texttt{DoSomething( \mbox{\texttt{\mdseries\slshape args}} : \mbox{\texttt{\mdseries\slshape stuff}} );} with a small additional abbreviation of \mbox{\texttt{\mdseries\slshape stuff}} permitted. See \ref{Function Call With Options} for details. The full form can be used where the same options are to run
across several calls, or where the \texttt{DoSomething} procedure is actually an infix operator, or other function with special
syntax. 

 An alternative to this system is the use of additional optional arguments in
procedure calls. This is not felt to be sufficient because many procedure
calls might cause, for example, a coset enumeration and each would need to
make provision for the possibility of extra arguments. In this system the
options are pushed when the user\texttt{\symbol{45}}level procedure is called,
and remain in effect (unless altered) for all procedures called by it. 

 Note that in some places in the system optional records containing options
which are valid only for the immediate function or method call are in fact
used.  
\section{\textcolor{Chapter }{Functions Dealing with the Options Stack}}\label{sect:Options Stack}
\logpage{[ 8, 1, 0 ]}
\hyperdef{L}{X794C5B5A80203CF9}{}
{
  

\subsection{\textcolor{Chapter }{PushOptions}}
\logpage{[ 8, 1, 1 ]}\nobreak
\hyperdef{L}{X7D4939FF7FB37FBE}{}
{\noindent\textcolor{FuncColor}{$\triangleright$\enspace\texttt{PushOptions({\mdseries\slshape options{\textunderscore}record})\index{PushOptions@\texttt{PushOptions}}
\label{PushOptions}
}\hfill{\scriptsize (function)}}\\


 This function pushes a record of options onto the global option stack. Note
that \texttt{PushOptions( rec( \mbox{\texttt{\mdseries\slshape opt}}:= fail ) )} has the effect of resetting the option \mbox{\texttt{\mdseries\slshape opt}}, since an option that has never been set has the value \texttt{fail} returned by \texttt{ValueOption} (\ref{ValueOption}). 

 Note that there is no check for misspelt or undefined options. }

 

\subsection{\textcolor{Chapter }{PopOptions}}
\logpage{[ 8, 1, 2 ]}\nobreak
\hyperdef{L}{X7818A5278679FD43}{}
{\noindent\textcolor{FuncColor}{$\triangleright$\enspace\texttt{PopOptions({\mdseries\slshape })\index{PopOptions@\texttt{PopOptions}}
\label{PopOptions}
}\hfill{\scriptsize (function)}}\\


 This function removes the top\texttt{\symbol{45}}most options record from the
options stack if there is one. }

 

\subsection{\textcolor{Chapter }{ResetOptionsStack}}
\logpage{[ 8, 1, 3 ]}\nobreak
\hyperdef{L}{X83D1190984DA3B85}{}
{\noindent\textcolor{FuncColor}{$\triangleright$\enspace\texttt{ResetOptionsStack({\mdseries\slshape })\index{ResetOptionsStack@\texttt{ResetOptionsStack}}
\label{ResetOptionsStack}
}\hfill{\scriptsize (function)}}\\


 unbinds (i.e. removes) all the options records from the options stack. 

 \emph{Note:} \texttt{ResetOptionsStack} should \emph{not} be used within a function. Its intended use is to clean up the options stack
in the event that the user has \texttt{quit} from a \texttt{break} loop, so leaving a stack of
no\texttt{\symbol{45}}longer\texttt{\symbol{45}}needed options
(see{\nobreakspace}\ref{quit}). }

 

\subsection{\textcolor{Chapter }{OnQuit}}
\logpage{[ 8, 1, 4 ]}\nobreak
\hyperdef{L}{X78D87D1081BF99FE}{}
{\noindent\textcolor{FuncColor}{$\triangleright$\enspace\texttt{OnQuit({\mdseries\slshape })\index{OnQuit@\texttt{OnQuit}}
\label{OnQuit}
}\hfill{\scriptsize (function)}}\\


 called when a user selects to \texttt{quit;} a break loop entered via execution of \texttt{Error} (\ref{Error}). As \textsf{GAP} starts up, \texttt{OnQuit} is defined to do nothing, in case an error is encountered during \textsf{GAP} start\texttt{\symbol{45}}up. Later in the loading process we redefine \texttt{OnQuit} to do a variant of \texttt{ResetOptionsStack} (\ref{ResetOptionsStack}) to ensure the options stack is empty after a user quits an \texttt{Error} (\ref{Error})\texttt{\symbol{45}}induced break loop. (\texttt{OnQuit} differs from \texttt{ResetOptionsStack} (\ref{ResetOptionsStack}) in that it warns when it does something rather than the other way round.)
Currently, \texttt{OnQuit} is not advertised, since exception handling may make it obsolete. }

 

\subsection{\textcolor{Chapter }{ValueOption}}
\logpage{[ 8, 1, 5 ]}\nobreak
\hyperdef{L}{X7F9373AD7DB88D1F}{}
{\noindent\textcolor{FuncColor}{$\triangleright$\enspace\texttt{ValueOption({\mdseries\slshape opt})\index{ValueOption@\texttt{ValueOption}}
\label{ValueOption}
}\hfill{\scriptsize (function)}}\\


 This function is a method for accessing the options stack without changing it; \mbox{\texttt{\mdseries\slshape opt}} should be the name of an option, i.e.{\nobreakspace}a string. A function which
makes decisions that might be affected by options should examine the result of \texttt{ValueOption}. If \mbox{\texttt{\mdseries\slshape opt}} is currently not set then \texttt{fail} is returned. }

 

\subsection{\textcolor{Chapter }{DisplayOptionsStack}}
\logpage{[ 8, 1, 6 ]}\nobreak
\hyperdef{L}{X7EDA4EB67D43FE33}{}
{\noindent\textcolor{FuncColor}{$\triangleright$\enspace\texttt{DisplayOptionsStack({\mdseries\slshape })\index{DisplayOptionsStack@\texttt{DisplayOptionsStack}}
\label{DisplayOptionsStack}
}\hfill{\scriptsize (function)}}\\


 This function prints a human\texttt{\symbol{45}}readable display of the
complete options stack. 

 }

 

\subsection{\textcolor{Chapter }{InfoOptions}}
\logpage{[ 8, 1, 7 ]}\nobreak
\hyperdef{L}{X832F41187B150C19}{}
{\noindent\textcolor{FuncColor}{$\triangleright$\enspace\texttt{InfoOptions\index{InfoOptions@\texttt{InfoOptions}}
\label{InfoOptions}
}\hfill{\scriptsize (info class)}}\\


 This info class can be used to enable messages about options being changed
(level 1) or accessed (level 2). }

 }

  
\section{\textcolor{Chapter }{Options Stack {\textendash} an Example}}\label{sect:Options Stack Example}
\logpage{[ 8, 2, 0 ]}
\hyperdef{L}{X7BB781647CAAE9B4}{}
{
  The example below shows simple manipulation of the Options Stack, first using \texttt{PushOptions} (\ref{PushOptions}) and \texttt{PopOptions} (\ref{PopOptions}) and then using the special function calling syntax. 

 
\begin{Verbatim}[commandchars=!@|,fontsize=\small,frame=single,label=Example]
  !gapprompt@gap>| !gapinput@foo := function()|
  !gapprompt@>| !gapinput@Print("myopt1 = ", ValueOption("myopt1"),|
  !gapprompt@>| !gapinput@      " myopt2 = ",ValueOption("myopt2"),"\n");|
  !gapprompt@>| !gapinput@end;|
  function(  ) ... end
  !gapprompt@gap>| !gapinput@foo();|
  myopt1 = fail myopt2 = fail
  !gapprompt@gap>| !gapinput@PushOptions(rec(myopt1 := 17));|
  !gapprompt@gap>| !gapinput@foo();|
  myopt1 = 17 myopt2 = fail
  !gapprompt@gap>| !gapinput@DisplayOptionsStack();|
  [ rec(
        myopt1 := 17 ) ]
  !gapprompt@gap>| !gapinput@PopOptions();|
  !gapprompt@gap>| !gapinput@foo();|
  myopt1 = fail myopt2 = fail
  !gapprompt@gap>| !gapinput@foo( : myopt1, myopt2 := [Z(3),"aardvark"]);|
  myopt1 = true myopt2 = [ Z(3), "aardvark" ]
  !gapprompt@gap>| !gapinput@DisplayOptionsStack();|
  [  ]
  !gapprompt@gap>| !gapinput@|
\end{Verbatim}
 }

 }

                
\chapter{\textcolor{Chapter }{Files and Filenames}}\label{Files and Filenames}
\logpage{[ 9, 0, 0 ]}
\hyperdef{L}{X82BCD4297920C903}{}
{
  Files are identified by filenames, which are represented in \textsf{GAP} as strings. Filenames can be created directly by the user or a program, but of
course this is operating system dependent. 

 Filenames for some files can be constructed in a system independent way using
the following functions. This is done by first getting a directory object for
the directory the file shall reside in, and then constructing the filename.
However, it is sometimes necessary to construct filenames of files in
subdirectories relative to a given directory object. In this case the
directory separator is \emph{always} \texttt{/} even under DOS or MacOS. 

 Section \ref{Directories} describes how to construct directory objects for the common \textsf{GAP} and system directories. Using the command \texttt{Filename} (\ref{Filename:for a directory and a string}) it is possible to construct a filename pointing to a file in these
directories. There are also functions to test for accessibility of files,
see{\nobreakspace}\ref{File Access}.  
\section{\textcolor{Chapter }{Portability}}\label{Portability}
\logpage{[ 9, 1, 0 ]}
\hyperdef{L}{X83D8AAA484EE95D9}{}
{
  For portability filenames and directory names should be restricted to at most
8 alphanumerical characters optionally followed by a dot \texttt{.} and between 1 and 3 alphanumerical characters. Upper case letters should be
avoided because some operating systems do not make any distinction between
case, so that \texttt{NaMe}, \texttt{Name} and \texttt{name} all refer to the same file whereas some operating systems are case sensitive.
To avoid problems only lower case characters should be used. 

 Another function which is system\texttt{\symbol{45}}dependent is \texttt{LastSystemError} (\ref{LastSystemError}). 

\subsection{\textcolor{Chapter }{LastSystemError}}
\logpage{[ 9, 1, 1 ]}\nobreak
\hyperdef{L}{X87D278437A916905}{}
{\noindent\textcolor{FuncColor}{$\triangleright$\enspace\texttt{LastSystemError({\mdseries\slshape })\index{LastSystemError@\texttt{LastSystemError}}
\label{LastSystemError}
}\hfill{\scriptsize (function)}}\\


 \texttt{LastSystemError} returns a record describing the last system error that has occurred. This
record contains at least the component \texttt{message} which is a string. This message is, however, highly operating system dependent
and should only be used as an informational message for the user. }

 }

  
\section{\textcolor{Chapter }{GAP Root Directories}}\label{GAP Root Directories}
\logpage{[ 9, 2, 0 ]}
\hyperdef{L}{X7A4973627A5DB27D}{}
{
  \index{GAPInfo.RootPaths@\texttt{GAPInfo.RootPaths}} \index{GAPInfo.UserGapRoot@\texttt{GAPInfo.UserGapRoot}} When \textsf{GAP} is started it determines a list of directories which we call the \emph{\textsf{GAP} root directories}. In a running \textsf{GAP} session this list can be found in \texttt{GAPInfo.RootPaths}. 

 The core part of \textsf{GAP} knows which files to read relative to its root directories. For example when \textsf{GAP} wants to read its library file \texttt{lib/group.gd}, it appends this path to each path in \texttt{GAPInfo.RootPaths} until it finds the path of an existing file. The first file found this way is
read. 

 Any subdirectories named \texttt{pkg/} in one of the directories in \texttt{GAPInfo.RootPaths} are added to \texttt{GAPInfo.PackageDirectories} (see \ref{GAP Package Directories}), which controls where \textsf{GAP} looks for available packages. 

 The root directories are specified via one or several of the \texttt{\texttt{\symbol{45}}l paths} command line options, see \ref{Command Line Options}. Furthermore, by default \textsf{GAP} automatically prepends a user specific \textsf{GAP} root directory to the list; this can be avoided by calling \textsf{GAP} with the \texttt{\texttt{\symbol{45}}r} option. The name of this user specific directory depends on your operating
system, it can be found in \texttt{GAPInfo.UserGapRoot}. This directory can be used to tell \textsf{GAP} about personal preferences, to always load some additional code, to install
additional packages, or to overwrite some \textsf{GAP} files. See \ref{sect:gap.ini} for more information how to do this. After \textsf{GAP} has been started, one can add additional root directories via the function \texttt{ExtendRootDirectories} (\ref{ExtendRootDirectories}). 

 }

  
\section{\textcolor{Chapter }{GAP Package Directories}}\label{GAP Package Directories}
\logpage{[ 9, 3, 0 ]}
\hyperdef{L}{X8223D52E78AF4420}{}
{
  \index{GAPInfo.PackageDirectories@\texttt{GAPInfo.PackageDirectories}} When \textsf{GAP} is started it determines a list of directories potentially containing
packages. We refer to these as the \emph{\textsf{GAP} package directories}. In a running \textsf{GAP} session this list can be found in \texttt{GAPInfo.PackageDirectories}. 

 Every subdirectory \texttt{pkg} in a \textsf{GAP} root directory is automatically added to this list. Further package
directories can be specified via one or several \texttt{\texttt{\symbol{45}}\texttt{\symbol{45}}packagedirs paths} command line options, see \ref{Command Line Options}, or after \textsf{GAP} has been started via the function \texttt{ExtendPackageDirectories} (\ref{ExtendPackageDirectories}). The order of the directories in \texttt{GAPInfo.PackageDirectories} is as follows: first the package directories specified via the command line
option \texttt{\texttt{\symbol{45}}\texttt{\symbol{45}}packagedirs}, then the subdirectories \texttt{pkg} of the \textsf{GAP} root directories that were known at startup in the same order, and finally the
directories added after \textsf{GAP} has been started. 

 \textsf{GAP} looks for available packages by examining each of the directories in \texttt{GAPInfo.PackageDirectories}. }

  
\section{\textcolor{Chapter }{Directories}}\label{Directories}
\logpage{[ 9, 4, 0 ]}
\hyperdef{L}{X85030B35865A1080}{}
{
  

\subsection{\textcolor{Chapter }{IsDirectory}}
\logpage{[ 9, 4, 1 ]}\nobreak
\hyperdef{L}{X82B3E24683942597}{}
{\noindent\textcolor{FuncColor}{$\triangleright$\enspace\texttt{IsDirectory({\mdseries\slshape obj})\index{IsDirectory@\texttt{IsDirectory}}
\label{IsDirectory}
}\hfill{\scriptsize (Category)}}\\


 \texttt{IsDirectory} is a category of directories. }

 

\subsection{\textcolor{Chapter }{Directory}}
\logpage{[ 9, 4, 2 ]}\nobreak
\hyperdef{L}{X86A71E927EEC7EAD}{}
{\noindent\textcolor{FuncColor}{$\triangleright$\enspace\texttt{Directory({\mdseries\slshape string})\index{Directory@\texttt{Directory}}
\label{Directory}
}\hfill{\scriptsize (operation)}}\\


 returns a directory object for the string \mbox{\texttt{\mdseries\slshape string}}. \texttt{Directory} understands \texttt{"."} for ``current directory'', that is, the directory in which \textsf{GAP} was started. It also understands absolute paths. 

 If the variable \texttt{GAPInfo.UserHome} is defined (this may depend on the operating system) then \texttt{Directory} understands a string with a leading \texttt{\texttt{\symbol{126}}} (tilde) character for a path relative to the user's home directory (but a
string beginning with \texttt{"\texttt{\symbol{126}}other{\textunderscore}user"} is \emph{not} interpreted as a path relative to \texttt{other{\textunderscore}user}'s home directory, as in a UNIX shell). 

 Paths are otherwise taken relative to the current directory. }

 

\subsection{\textcolor{Chapter }{DirectoryTemporary}}
\logpage{[ 9, 4, 3 ]}\nobreak
\hyperdef{L}{X8222B1A886E6195E}{}
{\noindent\textcolor{FuncColor}{$\triangleright$\enspace\texttt{DirectoryTemporary({\mdseries\slshape })\index{DirectoryTemporary@\texttt{DirectoryTemporary}}
\label{DirectoryTemporary}
}\hfill{\scriptsize (function)}}\\


 returns a directory object in the category \texttt{IsDirectory} (\ref{IsDirectory}) for a \emph{new} temporary directory. This is guaranteed to be newly created and empty
immediately after the call to \texttt{DirectoryTemporary}. \textsf{GAP} will make a reasonable effort to remove this directory upon termination of the \textsf{GAP} job that created the directory. 

 If \texttt{DirectoryTemporary} is unable to create a new directory, \texttt{fail} is returned. In this case \texttt{LastSystemError} (\ref{LastSystemError}) can be used to get information about the error. 

 A warning message is given if more than 1000 temporary directories are created
in any \textsf{GAP} session. }

 

\subsection{\textcolor{Chapter }{DirectoryCurrent}}
\logpage{[ 9, 4, 4 ]}\nobreak
\hyperdef{L}{X7BAD8036849E8430}{}
{\noindent\textcolor{FuncColor}{$\triangleright$\enspace\texttt{DirectoryCurrent({\mdseries\slshape })\index{DirectoryCurrent@\texttt{DirectoryCurrent}}
\label{DirectoryCurrent}
}\hfill{\scriptsize (function)}}\\


 returns the directory object for the current directory. }

 

\subsection{\textcolor{Chapter }{ChangeDirectoryCurrent}}
\logpage{[ 9, 4, 5 ]}\nobreak
\hyperdef{L}{X81DDD2E87F68E086}{}
{\noindent\textcolor{FuncColor}{$\triangleright$\enspace\texttt{ChangeDirectoryCurrent({\mdseries\slshape path})\index{ChangeDirectoryCurrent@\texttt{ChangeDirectoryCurrent}}
\label{ChangeDirectoryCurrent}
}\hfill{\scriptsize (function)}}\\


 Changes the current directory. Returns \texttt{true} on success and \texttt{fail} on failure. }

 

\subsection{\textcolor{Chapter }{DirectoriesLibrary}}
\logpage{[ 9, 4, 6 ]}\nobreak
\hyperdef{L}{X87ED469A85343A3C}{}
{\noindent\textcolor{FuncColor}{$\triangleright$\enspace\texttt{DirectoriesLibrary({\mdseries\slshape [name]})\index{DirectoriesLibrary@\texttt{DirectoriesLibrary}}
\label{DirectoriesLibrary}
}\hfill{\scriptsize (function)}}\\


 \texttt{DirectoriesLibrary} returns the directory objects for the \textsf{GAP} library \mbox{\texttt{\mdseries\slshape name}} as a list. \mbox{\texttt{\mdseries\slshape name}} must be one of \texttt{"lib"} (the default), \texttt{"doc"}, \texttt{"tst"}, and so on. 

 The string \texttt{""} is also legal and with this argument \texttt{DirectoriesLibrary} returns the list of \textsf{GAP} root directories. The return value of this call differs from \texttt{GAPInfo.RootPaths} in that the former is a list of directory objects and the latter a list of
strings. 

 The directory \mbox{\texttt{\mdseries\slshape name}} must exist in at least one of the root directories, otherwise \texttt{fail} is returned.   

 As the files in the \textsf{GAP} root directories (see{\nobreakspace}\ref{GAP Root Directories}) can be distributed into different directories in the filespace a list of
directories is returned. In order to find an existing file in a \textsf{GAP} root directory you should pass that list to \texttt{Filename} (\ref{Filename:for a directory and a string}) as the first argument. In order to create a filename for a new file inside a \textsf{GAP} root directory you should pass the first entry of that list. However, creating
files inside the \textsf{GAP} root directory is not recommended, you should use \texttt{DirectoryTemporary} (\ref{DirectoryTemporary}) instead. }

 

\subsection{\textcolor{Chapter }{DirectoriesSystemPrograms}}
\logpage{[ 9, 4, 7 ]}\nobreak
\hyperdef{L}{X808E2C187DD984B4}{}
{\noindent\textcolor{FuncColor}{$\triangleright$\enspace\texttt{DirectoriesSystemPrograms({\mdseries\slshape })\index{DirectoriesSystemPrograms@\texttt{DirectoriesSystemPrograms}}
\label{DirectoriesSystemPrograms}
}\hfill{\scriptsize (function)}}\\


 \texttt{DirectoriesSystemPrograms} returns the directory objects for the list of directories where the system
programs reside, as a list. Under UNIX this would usually represent \texttt{\$PATH}. }

 

\subsection{\textcolor{Chapter }{DirectoryContents}}
\logpage{[ 9, 4, 8 ]}\nobreak
\hyperdef{L}{X7B225E5282534EDA}{}
{\noindent\textcolor{FuncColor}{$\triangleright$\enspace\texttt{DirectoryContents({\mdseries\slshape dir})\index{DirectoryContents@\texttt{DirectoryContents}}
\label{DirectoryContents}
}\hfill{\scriptsize (function)}}\\


 This function returns a list of filenames/directory names that reside in the
directory \mbox{\texttt{\mdseries\slshape dir}}. The argument \mbox{\texttt{\mdseries\slshape dir}} can either be given as a string indicating the name of the directory or as a
directory object (see \texttt{IsDirectory} (\ref{IsDirectory})). If an error occurs (the specified directory does not exist or has no read
permissions), \texttt{fail} is returned. In this case \texttt{LastSystemError} (\ref{LastSystemError}) can be used to get information about the error. 

 The ordering of the list entries can depend on the operating system. 

 An interactive way to show the contents of a directory is provided by the
function \texttt{BrowseDirectory} (\textbf{Browse: BrowseDirectory}) from the \textsf{GAP} package \textsf{Browse}. }

 

\subsection{\textcolor{Chapter }{DirectoryDesktop}}
\logpage{[ 9, 4, 9 ]}\nobreak
\hyperdef{L}{X86F4A32C83B82369}{}
{\noindent\textcolor{FuncColor}{$\triangleright$\enspace\texttt{DirectoryDesktop({\mdseries\slshape })\index{DirectoryDesktop@\texttt{DirectoryDesktop}}
\label{DirectoryDesktop}
}\hfill{\scriptsize (function)}}\\


 returns a directory object for the user's desktop directory as defined on many
modern operating systems. The function is intended to provide a
cross\texttt{\symbol{45}}platform interface to a directory that is easily
accessible by the user. 

 Under Unix systems (including macOS) this will be the \texttt{Desktop} directory in the user's home directory if it exists, and the user's home
directory otherwise. Under Windows it will be the user's \texttt{Desktop} folder (or the appropriate name under different languages). }

 

\subsection{\textcolor{Chapter }{DirectoryHome}}
\logpage{[ 9, 4, 10 ]}\nobreak
\hyperdef{L}{X7B0D818A808A3481}{}
{\noindent\textcolor{FuncColor}{$\triangleright$\enspace\texttt{DirectoryHome({\mdseries\slshape })\index{DirectoryHome@\texttt{DirectoryHome}}
\label{DirectoryHome}
}\hfill{\scriptsize (function)}}\\


 returns a directory object for the user's home directory, defined as a
directory in which the user will typically have full read and write access.
The function is intended to provide a cross\texttt{\symbol{45}}platform
interface to a directory that is easily accessible by the user. 

 Under Unix systems (including macOS) this will be the usual user home
directory. Under Windows it will be the user's \texttt{My Documents} folder (or the appropriate name under different languages). }

 }

  
\section{\textcolor{Chapter }{File Names}}\label{File Names}
\logpage{[ 9, 5, 0 ]}
\hyperdef{L}{X8545E03E7D651456}{}
{
  
\subsection{\textcolor{Chapter }{Filename}}\logpage{[ 9, 5, 1 ]}
\hyperdef{L}{X7E352E1F87060602}{}
{
\noindent\textcolor{FuncColor}{$\triangleright$\enspace\texttt{Filename({\mdseries\slshape dir, name})\index{Filename@\texttt{Filename}!for a directory and a string}
\label{Filename:for a directory and a string}
}\hfill{\scriptsize (operation)}}\\
\noindent\textcolor{FuncColor}{$\triangleright$\enspace\texttt{Filename({\mdseries\slshape list\texttt{\symbol{45}}of\texttt{\symbol{45}}dirs, name})\index{Filename@\texttt{Filename}!for a list of directories and a string}
\label{Filename:for a list of directories and a string}
}\hfill{\scriptsize (operation)}}\\


 If the first argument is a directory object \mbox{\texttt{\mdseries\slshape dir}}, \texttt{Filename} returns the (system dependent) filename as a string for the file with name \mbox{\texttt{\mdseries\slshape name}} in the directory \mbox{\texttt{\mdseries\slshape dir}}. \texttt{Filename} returns the filename regardless of whether the directory contains a file with
name \mbox{\texttt{\mdseries\slshape name}} or not. 

 If the first argument is a list \mbox{\texttt{\mdseries\slshape list\texttt{\symbol{45}}of\texttt{\symbol{45}}dirs}} (possibly of length 1) of directory objects, then \texttt{Filename} searches the directories in order, and returns the filename for the file \mbox{\texttt{\mdseries\slshape name}} in the first directory which contains a file \mbox{\texttt{\mdseries\slshape name}} or \texttt{fail} if no directory contains a file \mbox{\texttt{\mdseries\slshape name}}. 

 \emph{For example}, in order to locate the system program \texttt{date} use \texttt{PathSystemProgram} (\ref{PathSystemProgram}). 

 
\begin{Verbatim}[commandchars=!@|,fontsize=\small,frame=single,label=Example]
  !gapprompt@gap>| !gapinput@date := PathSystemProgram( "date" );|
  "/bin/date"
\end{Verbatim}
 

 In order to locate the library file \texttt{files.gd} use \texttt{DirectoriesLibrary} (\ref{DirectoriesLibrary}) together with the second form of \texttt{Filename}. 

 
\begin{Verbatim}[commandchars=!@|,fontsize=\small,frame=single,label=Example]
  !gapprompt@gap>| !gapinput@path := DirectoriesLibrary();;|
  !gapprompt@gap>| !gapinput@Filename( path, "files.gd" );|
  "./lib/files.gd"
\end{Verbatim}
 

 In order to construct filenames for new files in a temporary directory use \texttt{DirectoryTemporary} (\ref{DirectoryTemporary}) together with the first form of \texttt{Filename}. 

 
\begin{Verbatim}[commandchars=!@|,fontsize=\small,frame=single,label=Example]
  !gapprompt@gap>| !gapinput@tmpdir := DirectoryTemporary();;|
  !gapprompt@gap>| !gapinput@Filename( [ tmpdir ], "file.new" );|
  fail
  !gapprompt@gap>| !gapinput@Filename( tmpdir, "file.new" );|
  "/var/tmp/tmp.0.021738.0001/file.new"
\end{Verbatim}
 }

 

\subsection{\textcolor{Chapter }{PathSystemProgram}}
\logpage{[ 9, 5, 2 ]}\nobreak
\hyperdef{L}{X86C7683E7A2A2146}{}
{\noindent\textcolor{FuncColor}{$\triangleright$\enspace\texttt{PathSystemProgram({\mdseries\slshape name})\index{PathSystemProgram@\texttt{PathSystemProgram}}
\label{PathSystemProgram}
}\hfill{\scriptsize (function)}}\\


 \texttt{PathSystemProgram} returns either the path of the first executable file \mbox{\texttt{\mdseries\slshape name}} in one of the directories returned by \texttt{DirectoriesSystemPrograms} (\ref{DirectoriesSystemPrograms}), or \texttt{fail} if no such file exists. }

 }

  
\section{\textcolor{Chapter }{Special Filenames}}\label{Special Filenames}
\logpage{[ 9, 6, 0 ]}
\hyperdef{L}{X85EC7D9087C481B0}{}
{
  The special filename \texttt{"*stdin*"} denotes the standard input, i.e., the stream through which the user enters
commands to \textsf{GAP}. The exact behaviour of reading from \texttt{"*stdin*"} is operating system dependent, but usually the following happens. If \textsf{GAP} was started with no input redirection, statements are read from the terminal
stream until the user enters the end of file character, which is usually \textsc{Ctrl\texttt{\symbol{45}}D}. Note that terminal streams are special, in that they may yield ordinary
input \emph{after} an end of file. Thus when control returns to the main
read\texttt{\symbol{45}}eval\texttt{\symbol{45}}print loop the user can
continue with \textsf{GAP}. If \textsf{GAP} was started with an input redirection, statements are read from the current
position in the input file up to the end of the file. When control returns to
the main read eval view loop the input stream will still return end of file,
and \textsf{GAP} will terminate. 

 The special filename \texttt{"*errin*"} denotes the stream connected to the UNIX \texttt{stderr} output. This stream is usually connected to the terminal, even if the standard
input was redirected, unless the standard error stream was also redirected, in
which case opening of \texttt{"*errin*"} fails. 

 The special filename \texttt{"*stdout*"} can be used to print to the standard output. 

 The special filename \texttt{"*errout*"} can be used to print to the standard error output file, which is usually
connected to the terminal, even if the standard output was redirected. }

  
\section{\textcolor{Chapter }{File Access}}\label{File Access}
\logpage{[ 9, 7, 0 ]}
\hyperdef{L}{X87271FEF86A6A0F9}{}
{
  When the following functions return \texttt{false} one can use \texttt{LastSystemError} (\ref{LastSystemError}) to find out the reason (as provided by the operating system), see the
examples. 

\subsection{\textcolor{Chapter }{IsExistingFile}}
\logpage{[ 9, 7, 1 ]}\nobreak
\hyperdef{L}{X8269697A7B927AF1}{}
{\noindent\textcolor{FuncColor}{$\triangleright$\enspace\texttt{IsExistingFile({\mdseries\slshape filename})\index{IsExistingFile@\texttt{IsExistingFile}}
\label{IsExistingFile}
}\hfill{\scriptsize (function)}}\\


 \texttt{IsExistingFile} returns \texttt{true} if a file with the filename \mbox{\texttt{\mdseries\slshape filename}} exists and can be seen by the \textsf{GAP} process. Otherwise \texttt{false} is returned. 

 
\begin{Verbatim}[commandchars=!@|,fontsize=\small,frame=single,label=Example]
  !gapprompt@gap>| !gapinput@IsExistingFile( "/bin/date" );     # file `/bin/date' exists|
  true
  !gapprompt@gap>| !gapinput@IsExistingFile( "/bin/date.new" ); # non existing `/bin/date.new'|
  false
  !gapprompt@gap>| !gapinput@IsExistingFile( "/bin/date/new" ); # `/bin/date' is not a directory|
  false
  !gapprompt@gap>| !gapinput@LastSystemError().message;|
  "Not a directory"
\end{Verbatim}
 }

 

\subsection{\textcolor{Chapter }{IsReadableFile}}
\logpage{[ 9, 7, 2 ]}\nobreak
\hyperdef{L}{X7E156EC886E11BBC}{}
{\noindent\textcolor{FuncColor}{$\triangleright$\enspace\texttt{IsReadableFile({\mdseries\slshape filename})\index{IsReadableFile@\texttt{IsReadableFile}}
\label{IsReadableFile}
}\hfill{\scriptsize (function)}}\\


 \texttt{IsReadableFile} returns \texttt{true} if a file with the filename \mbox{\texttt{\mdseries\slshape filename}} exists \emph{and} the \textsf{GAP} process has read permissions for the file, or \texttt{false} if this is not the case. 

 
\begin{Verbatim}[commandchars=!@|,fontsize=\small,frame=single,label=Example]
  !gapprompt@gap>| !gapinput@IsReadableFile( "/bin/date" );     # file `/bin/date' is readable|
  true
  !gapprompt@gap>| !gapinput@IsReadableFile( "/bin/date.new" ); # non-existing `/bin/date.new'|
  false
  !gapprompt@gap>| !gapinput@LastSystemError().message;|
  "No such file or directory"
\end{Verbatim}
 }

 

\subsection{\textcolor{Chapter }{IsWritableFile}}
\logpage{[ 9, 7, 3 ]}\nobreak
\hyperdef{L}{X8412F485796B25F5}{}
{\noindent\textcolor{FuncColor}{$\triangleright$\enspace\texttt{IsWritableFile({\mdseries\slshape filename})\index{IsWritableFile@\texttt{IsWritableFile}}
\label{IsWritableFile}
}\hfill{\scriptsize (function)}}\\


 \texttt{IsWritableFile} returns \texttt{true} if a file with the filename \mbox{\texttt{\mdseries\slshape filename}} exists \emph{and} the \textsf{GAP} process has write permissions for the file, or \texttt{false} if this is not the case. 

 
\begin{Verbatim}[commandchars=!@|,fontsize=\small,frame=single,label=Example]
  !gapprompt@gap>| !gapinput@IsWritableFile( "/bin/date" );  # file `/bin/date' is not writable|
  false
\end{Verbatim}
 

 }

 

\subsection{\textcolor{Chapter }{IsExecutableFile}}
\logpage{[ 9, 7, 4 ]}\nobreak
\hyperdef{L}{X83A1AAD58435FC4C}{}
{\noindent\textcolor{FuncColor}{$\triangleright$\enspace\texttt{IsExecutableFile({\mdseries\slshape filename})\index{IsExecutableFile@\texttt{IsExecutableFile}}
\label{IsExecutableFile}
}\hfill{\scriptsize (function)}}\\


 \texttt{IsExecutableFile} returns \texttt{true} if a file with the filename \mbox{\texttt{\mdseries\slshape filename}} exists \emph{and} the \textsf{GAP} process has execute permissions for the file, or \texttt{false} if this is not the case. Note that execute permissions do not imply that it is
possible to execute the file, e.g., it may only be executable on a different
machine. 

 
\begin{Verbatim}[commandchars=!@|,fontsize=\small,frame=single,label=Example]
  !gapprompt@gap>| !gapinput@IsExecutableFile( "/bin/date" );   # ... but executable|
  true
\end{Verbatim}
 }

 

\subsection{\textcolor{Chapter }{IsDirectoryPath}}
\logpage{[ 9, 7, 5 ]}\nobreak
\hyperdef{L}{X7D1BE00F83C4EEE8}{}
{\noindent\textcolor{FuncColor}{$\triangleright$\enspace\texttt{IsDirectoryPath({\mdseries\slshape filename})\index{IsDirectoryPath@\texttt{IsDirectoryPath}}
\label{IsDirectoryPath}
}\hfill{\scriptsize (function)}}\\


 \texttt{IsDirectoryPath} returns \texttt{true} if the file with the filename \mbox{\texttt{\mdseries\slshape filename}} exists \emph{and} is a directory, and \texttt{false} otherwise. Note that this function does not check if the \textsf{GAP} process actually has write or execute permissions for the directory. You can
use \texttt{IsWritableFile} (\ref{IsWritableFile}), resp. \texttt{IsExecutableFile} (\ref{IsExecutableFile}) to check such permissions. }

 }

  
\section{\textcolor{Chapter }{File Operations}}\label{File Operations}
\logpage{[ 9, 8, 0 ]}
\hyperdef{L}{X81A0A4FF842B039B}{}
{
  

\subsection{\textcolor{Chapter }{Read}}
\logpage{[ 9, 8, 1 ]}\nobreak
\hyperdef{L}{X8373AC6B7D5F9167}{}
{\noindent\textcolor{FuncColor}{$\triangleright$\enspace\texttt{Read({\mdseries\slshape filename})\index{Read@\texttt{Read}}
\label{Read}
}\hfill{\scriptsize (operation)}}\\


 reads the input from the file with the filename \mbox{\texttt{\mdseries\slshape filename}}, which must be given as a string. 

 \texttt{Read} first opens the file \mbox{\texttt{\mdseries\slshape filename}}. If the file does not exist, or if \textsf{GAP} cannot open it, e.g., because of access restrictions, an error is signalled. 

 Then the contents of the file are read and evaluated, but the results are not
printed. The reading and evaluations happens exactly as described for the main
loop (see \ref{Main Loop}). 

 If a statement in the file causes an error a break loop is entered
(see{\nobreakspace}\ref{Break Loops}). The input for this break loop is not taken from the file, but from the
input connected to the \texttt{stderr} output of \textsf{GAP}. If \texttt{stderr} is not connected to a terminal, no break loop is entered. If this break loop
is left with \texttt{quit} (or \textsc{Ctrl\texttt{\symbol{45}}D}), \textsf{GAP} exits from the \texttt{Read} command, and from all enclosing \texttt{Read} commands, so that control is normally returned to an interactive prompt. The \texttt{QUIT} statement (see{\nobreakspace}\ref{Leaving GAP}) can also be used in the break loop to exit \textsf{GAP} immediately. 

 Note that a statement must not begin in one file and end in another. I.e., \emph{eof} (\emph{e}nd\texttt{\symbol{45}}\emph{o}f\texttt{\symbol{45}}\emph{f}ile) is not treated as whitespace, but as a special symbol that must not
appear inside any statement. 

 Note that one file may very well contain a read statement causing another file
to be read, before input is again taken from the first file. There is an upper
limit of 15 on the number of files that may be open simultaneously. }

 

\subsection{\textcolor{Chapter }{ReadAsFunction}}
\logpage{[ 9, 8, 2 ]}\nobreak
\hyperdef{L}{X7824CB7D7D4BAFBC}{}
{\noindent\textcolor{FuncColor}{$\triangleright$\enspace\texttt{ReadAsFunction({\mdseries\slshape filename})\index{ReadAsFunction@\texttt{ReadAsFunction}}
\label{ReadAsFunction}
}\hfill{\scriptsize (operation)}}\\


 reads the file with filename \mbox{\texttt{\mdseries\slshape filename}} as a function and returns this function. 

 \emph{Example} 

 Suppose that the file \texttt{/tmp/example.g} contains the following 

 
\begin{Verbatim}[commandchars=!@|,fontsize=\small,frame=single,label=Example]
  local a;
  
  a := 10;
  return a*10;
\end{Verbatim}
 

 Reading the file as a function will not affect a global variable \texttt{a}. 

 
\begin{Verbatim}[commandchars=!@|,fontsize=\small,frame=single,label=Example]
  !gapprompt@gap>| !gapinput@a := 1;|
  1
  !gapprompt@gap>| !gapinput@ReadAsFunction("/tmp/example.g")();|
  100
  !gapprompt@gap>| !gapinput@a;|
  1
\end{Verbatim}
 }

 
\subsection{\textcolor{Chapter }{PrintTo and AppendTo}}\logpage{[ 9, 8, 3 ]}
\hyperdef{L}{X86956C577FFEE1F9}{}
{
\noindent\textcolor{FuncColor}{$\triangleright$\enspace\texttt{PrintTo({\mdseries\slshape filename[, obj1, ...]})\index{PrintTo@\texttt{PrintTo}}
\label{PrintTo}
}\hfill{\scriptsize (function)}}\\
\noindent\textcolor{FuncColor}{$\triangleright$\enspace\texttt{AppendTo({\mdseries\slshape filename[, obj1, ...]})\index{AppendTo@\texttt{AppendTo}}
\label{AppendTo}
}\hfill{\scriptsize (function)}}\\


 \texttt{PrintTo} works like \texttt{Print} (\ref{Print}), except that the arguments \mbox{\texttt{\mdseries\slshape obj1}}, $\ldots$ (if present) are printed to the file with the name \mbox{\texttt{\mdseries\slshape filename}} instead of the standard output. This file must of course be writable by \textsf{GAP}. Otherwise an error is signalled. Note that \texttt{PrintTo} will \emph{overwrite} the previous contents of this file if it already existed; in particular, \texttt{PrintTo} with just the \mbox{\texttt{\mdseries\slshape filename}} argument empties that file. 

 \texttt{AppendTo} works like \texttt{PrintTo}, except that the output does not overwrite the previous contents of the file,
but is appended to the file. 

 There is an upper limit of 15 on the number of output files that may be open
simultaneously. 

 \emph{Note} that one should be careful not to write to a logfile (see \texttt{LogTo} (\ref{LogTo:for a filename})) with \texttt{PrintTo} or \texttt{AppendTo}.  }

 
\subsection{\textcolor{Chapter }{LogTo}}\logpage{[ 9, 8, 4 ]}
\hyperdef{L}{X79813A6686894960}{}
{
\noindent\textcolor{FuncColor}{$\triangleright$\enspace\texttt{LogTo({\mdseries\slshape filename})\index{LogTo@\texttt{LogTo}!for a filename}
\label{LogTo:for a filename}
}\hfill{\scriptsize (operation)}}\\
\noindent\textcolor{FuncColor}{$\triangleright$\enspace\texttt{LogTo({\mdseries\slshape })\index{LogTo@\texttt{LogTo}!stop logging}
\label{LogTo:stop logging}
}\hfill{\scriptsize (operation)}}\\


 Calling \texttt{LogTo} with a string \mbox{\texttt{\mdseries\slshape filename}} causes the subsequent interaction to be logged to the file with the name \mbox{\texttt{\mdseries\slshape filename}}, i.e., everything you see on your terminal will also appear in this file. (\texttt{LogTo} (\ref{LogTo:for streams}) may also be used to log to a stream.) This file must of course be writable by \textsf{GAP}, otherwise an error is signalled. Note that \texttt{LogTo} will overwrite the previous contents of this file if it already existed. 

 Called without arguments, \texttt{LogTo} stops logging to a file or stream. }

 
\subsection{\textcolor{Chapter }{InputLogTo}}\logpage{[ 9, 8, 5 ]}
\hyperdef{L}{X7CAB119378B075B7}{}
{
\noindent\textcolor{FuncColor}{$\triangleright$\enspace\texttt{InputLogTo({\mdseries\slshape filename})\index{InputLogTo@\texttt{InputLogTo}!for a filename}
\label{InputLogTo:for a filename}
}\hfill{\scriptsize (operation)}}\\
\noindent\textcolor{FuncColor}{$\triangleright$\enspace\texttt{InputLogTo({\mdseries\slshape })\index{InputLogTo@\texttt{InputLogTo}!stop logging input}
\label{InputLogTo:stop logging input}
}\hfill{\scriptsize (operation)}}\\


 Calling \texttt{InputLogTo} with a string \mbox{\texttt{\mdseries\slshape filename}} causes the subsequent input to be logged to the file with the name \mbox{\texttt{\mdseries\slshape filename}}, i.e., everything you type on your terminal will also appear in this file.
Note that \texttt{InputLogTo} and \texttt{LogTo} (\ref{LogTo:for a filename}) cannot be used at the same time while \texttt{InputLogTo} and \texttt{OutputLogTo} (\ref{OutputLogTo:for a filename}) can. Note that \texttt{InputLogTo} will overwrite the previous contents of this file if it already existed. 

 Called without arguments, \texttt{InputLogTo} stops logging to a file or stream. }

 
\subsection{\textcolor{Chapter }{OutputLogTo}}\logpage{[ 9, 8, 6 ]}
\hyperdef{L}{X7A5591D87EAFA6CC}{}
{
\noindent\textcolor{FuncColor}{$\triangleright$\enspace\texttt{OutputLogTo({\mdseries\slshape filename})\index{OutputLogTo@\texttt{OutputLogTo}!for a filename}
\label{OutputLogTo:for a filename}
}\hfill{\scriptsize (operation)}}\\
\noindent\textcolor{FuncColor}{$\triangleright$\enspace\texttt{OutputLogTo({\mdseries\slshape })\index{OutputLogTo@\texttt{OutputLogTo}!stop logging output}
\label{OutputLogTo:stop logging output}
}\hfill{\scriptsize (operation)}}\\


 Calling \texttt{OutputLogTo} with a string \mbox{\texttt{\mdseries\slshape filename}} causes the subsequent output to be logged to the file with the name \mbox{\texttt{\mdseries\slshape filename}}, i.e., everything \textsf{GAP} prints on your terminal will also appear in this file. Note that \texttt{OutputLogTo} and \texttt{LogTo} (\ref{LogTo:for a filename}) cannot be used at the same time while \texttt{InputLogTo} (\ref{InputLogTo:for a filename}) and \texttt{OutputLogTo} can. Note that \texttt{OutputLogTo} will overwrite the previous contents of this file if it already existed. 

 Called without arguments, \texttt{OutputLogTo} stops logging to a file or stream. }

 

\subsection{\textcolor{Chapter }{CrcFile}}
\logpage{[ 9, 8, 7 ]}\nobreak
\hyperdef{L}{X8241CEAD80415BB9}{}
{\noindent\textcolor{FuncColor}{$\triangleright$\enspace\texttt{CrcFile({\mdseries\slshape filename})\index{CrcFile@\texttt{CrcFile}}
\label{CrcFile}
}\hfill{\scriptsize (function)}}\\


 \index{hash function} \index{checksum} This function computes a CRC (cyclic redundancy check) number for the content
of the file \mbox{\texttt{\mdseries\slshape filename}}. 

 \texttt{CrcFile} computes a CRC (cyclic redundancy check) checksum value for the file with
filename \mbox{\texttt{\mdseries\slshape filename}} and returns this value as an integer. The function returns \texttt{fail} if an error occurred, for example, if \mbox{\texttt{\mdseries\slshape filename}} does not exist. In this case the function \texttt{LastSystemError} (\ref{LastSystemError}) can be used to get information about the error. See also \texttt{CrcFile} and \texttt{HexSHA256} (\ref{HexSHA256}). 

 
\begin{Verbatim}[commandchars=!@|,fontsize=\small,frame=single,label=Example]
  !gapprompt@gap>| !gapinput@CrcFile( "lib/morpheus.gi" );|
  2705743645
\end{Verbatim}
 }

 

\subsection{\textcolor{Chapter }{RemoveFile}}
\logpage{[ 9, 8, 8 ]}\nobreak
\hyperdef{L}{X7E63ACA38142BE96}{}
{\noindent\textcolor{FuncColor}{$\triangleright$\enspace\texttt{RemoveFile({\mdseries\slshape filename})\index{RemoveFile@\texttt{RemoveFile}}
\label{RemoveFile}
}\hfill{\scriptsize (function)}}\\


 will remove the file with filename \mbox{\texttt{\mdseries\slshape filename}} and returns \texttt{true} in case of success. The function returns \texttt{fail} if a system error occurred, for example, if your permissions do not allow the
removal of \mbox{\texttt{\mdseries\slshape filename}}. In this case the function \texttt{LastSystemError} (\ref{LastSystemError}) can be used to get information about the error. }

 

\subsection{\textcolor{Chapter }{UserHomeExpand}}
\logpage{[ 9, 8, 9 ]}\nobreak
\hyperdef{L}{X83F3B0337C7EA5CC}{}
{\noindent\textcolor{FuncColor}{$\triangleright$\enspace\texttt{UserHomeExpand({\mdseries\slshape str})\index{UserHomeExpand@\texttt{UserHomeExpand}}
\label{UserHomeExpand}
}\hfill{\scriptsize (function)}}\\


 If the string \mbox{\texttt{\mdseries\slshape str}} starts with a \texttt{'\texttt{\symbol{126}}'} character this function returns a new string with the leading \texttt{'\texttt{\symbol{126}}'} substituted by the user's home directory as stored in \texttt{GAPInfo.UserHome}. Otherwise \mbox{\texttt{\mdseries\slshape str}} is returned unchanged. 

 This function is the counterpart of \texttt{UserHomeShorten} (\ref{UserHomeShorten}). }

 

\subsection{\textcolor{Chapter }{UserHomeShorten}}
\logpage{[ 9, 8, 10 ]}\nobreak
\hyperdef{L}{X7B450DAD8651FC2B}{}
{\noindent\textcolor{FuncColor}{$\triangleright$\enspace\texttt{UserHomeShorten({\mdseries\slshape str})\index{UserHomeShorten@\texttt{UserHomeShorten}}
\label{UserHomeShorten}
}\hfill{\scriptsize (function)}}\\


 If the string \mbox{\texttt{\mdseries\slshape str}} starts with the user's home directory as stored in \texttt{GAPInfo.UserHome} then this function returns a new string with that prefix replaced by a leading \texttt{'\texttt{\symbol{126}}'} character. Otherwise \mbox{\texttt{\mdseries\slshape str}} is returned unchanged. 

 This function is the counterpart of \texttt{UserHomeExpand} (\ref{UserHomeExpand}). }

 

\subsection{\textcolor{Chapter }{Reread}}
\logpage{[ 9, 8, 11 ]}\nobreak
\hyperdef{L}{X79EE267A7FAF28A6}{}
{\noindent\textcolor{FuncColor}{$\triangleright$\enspace\texttt{Reread({\mdseries\slshape filename})\index{Reread@\texttt{Reread}}
\label{Reread}
}\hfill{\scriptsize (function)}}\\
\noindent\textcolor{FuncColor}{$\triangleright$\enspace\texttt{REREADING\index{REREADING@\texttt{REREADING}}
\label{REREADING}
}\hfill{\scriptsize (global variable)}}\\


 In general, it is not possible to read the same \textsf{GAP} library file twice, or to read a compiled version after reading a \textsf{GAP} version, because crucial global variables are made
read\texttt{\symbol{45}}only (see \ref{More About Global Variables}) and filters and methods are added to global tables. 

 A partial solution to this problem is provided by the function \texttt{Reread} (and related functions \texttt{RereadLib} etc.). \texttt{Reread( \mbox{\texttt{\mdseries\slshape filename}} )} sets the global variable \texttt{REREADING} to \texttt{true}, reads the file named by \mbox{\texttt{\mdseries\slshape filename}} and then resets \texttt{REREADING}. Various system functions behave differently when \texttt{REREADING} is set to \texttt{true}. In particular, assignment to read\texttt{\symbol{45}}only global variables
is permitted, calls to \texttt{NewRepresentation} (\ref{NewRepresentation}) and \texttt{NewInfoClass} (\ref{NewInfoClass}) with parameters identical to those of an existing representation or info class
will return the existing object, and methods installed with \texttt{InstallMethod} (\ref{InstallMethod}) may sometimes displace existing methods. 

 This function may not entirely produce the intended results, especially if
what has changed is the super\texttt{\symbol{45}}representation of a
representation or the requirements of a method. In these cases, it is
necessary to restart \textsf{GAP} to read the modified file. 

 An additional use of \texttt{Reread} is to load the compiled version of a file for which the \textsf{GAP} language version had previously been read (or perhaps was included in a saved
workspace). See \ref{Kernel modules} and \ref{Saving and Loading a Workspace} for more information. 

 It is not advisable to use \texttt{Reread} programmatically. For example, if a file that contains calls to \texttt{Reread} is read with \texttt{Reread} then \texttt{REREADING} may be reset too early. }

 }

 }

                 
\chapter{\textcolor{Chapter }{Streams}}\label{Streams}
\logpage{[ 10, 0, 0 ]}
\hyperdef{L}{X839725177BF8B5B4}{}
{
  \emph{Streams} provide flexible access to \textsf{GAP}'s input and output processing. An \emph{input stream} takes characters from some source and delivers them to \textsf{GAP} which \emph{reads} them from the stream. When an input stream has delivered all characters it is
at \texttt{end\texttt{\symbol{45}}of\texttt{\symbol{45}}stream}. An \emph{output stream} receives characters from \textsf{GAP} which \emph{writes} them to the stream, and delivers them to some destination. 

 A major use of streams is to provide efficient and flexible access to files.
Files can be read and written using \texttt{Read} (\ref{Read}) and \texttt{AppendTo} (\ref{AppendTo}), however the former only allows a complete file to be read as \textsf{GAP} input and the latter imposes a high time penalty if many small pieces of
output are written to a large file. Streams allow input files in other formats
to be read and processed, and files to be built up efficiently from small
pieces of output. Streams may also be used for other purposes, for example to
read from and print to \textsf{GAP} strings, or to read input directly from the user. 

 Any stream is either a \emph{text stream}, which translates the \texttt{end\texttt{\symbol{45}}of\texttt{\symbol{45}}line} character (\texttt{\texttt{\symbol{92}}n}) to or from the system's representation of \texttt{end\texttt{\symbol{45}}of\texttt{\symbol{45}}line} (e.g., \emph{new\texttt{\symbol{45}}line} under UNIX and \emph{carriage\texttt{\symbol{45}}return}\texttt{\symbol{45}}\emph{new\texttt{\symbol{45}}line} under DOS), or a \emph{binary stream}, which does not translate the \texttt{end\texttt{\symbol{45}}of\texttt{\symbol{45}}line} character. The processing of other unprintable characters by text streams is
undefined. Binary streams pass them unchanged. 

 Whereas it is cheap to append to a stream, streams do consume system
resources, and only a limited number can be open at any time, therefore it is
necessary to close a stream as soon as possible using \texttt{CloseStream} (\ref{CloseStream}). If creating a stream failed then \texttt{LastSystemError} (\ref{LastSystemError}) can be used to get information about the failure.  
\section{\textcolor{Chapter }{Categories for Streams and the StreamsFamily}}\label{Categories for Streams and the StreamsFamily}
\logpage{[ 10, 1, 0 ]}
\hyperdef{L}{X7F89070B7CF52DE0}{}
{
  

\subsection{\textcolor{Chapter }{IsStream}}
\logpage{[ 10, 1, 1 ]}\nobreak
\hyperdef{L}{X7E974B96785E91A8}{}
{\noindent\textcolor{FuncColor}{$\triangleright$\enspace\texttt{IsStream({\mdseries\slshape obj})\index{IsStream@\texttt{IsStream}}
\label{IsStream}
}\hfill{\scriptsize (Category)}}\\


 Streams are \textsf{GAP} objects and all open streams, input, output, text and binary, lie in this
category. }

 

\subsection{\textcolor{Chapter }{IsClosedStream}}
\logpage{[ 10, 1, 2 ]}\nobreak
\hyperdef{L}{X7FE4096F8497B7F2}{}
{\noindent\textcolor{FuncColor}{$\triangleright$\enspace\texttt{IsClosedStream({\mdseries\slshape obj})\index{IsClosedStream@\texttt{IsClosedStream}}
\label{IsClosedStream}
}\hfill{\scriptsize (Category)}}\\


 When a stream is closed, its type changes to lie in \texttt{IsClosedStream}. This category is used to install methods that trap accesses to closed
streams. }

 

\subsection{\textcolor{Chapter }{IsInputStream}}
\logpage{[ 10, 1, 3 ]}\nobreak
\hyperdef{L}{X7FB4391283847C3A}{}
{\noindent\textcolor{FuncColor}{$\triangleright$\enspace\texttt{IsInputStream({\mdseries\slshape obj})\index{IsInputStream@\texttt{IsInputStream}}
\label{IsInputStream}
}\hfill{\scriptsize (Category)}}\\


 All input streams lie in this category, and support input operations such as \texttt{ReadByte} (\ref{ReadByte}) (see \ref{Operations for Input Streams}) }

 

\subsection{\textcolor{Chapter }{IsInputTextStream}}
\logpage{[ 10, 1, 4 ]}\nobreak
\hyperdef{L}{X7C8956BB7FE2A89C}{}
{\noindent\textcolor{FuncColor}{$\triangleright$\enspace\texttt{IsInputTextStream({\mdseries\slshape obj})\index{IsInputTextStream@\texttt{IsInputTextStream}}
\label{IsInputTextStream}
}\hfill{\scriptsize (Category)}}\\


 All \emph{text} input streams lie in this category. They translate new\texttt{\symbol{45}}line
characters read. }

 

\subsection{\textcolor{Chapter }{IsInputTextNone}}
\logpage{[ 10, 1, 5 ]}\nobreak
\hyperdef{L}{X7DCD6ADC86CF2472}{}
{\noindent\textcolor{FuncColor}{$\triangleright$\enspace\texttt{IsInputTextNone({\mdseries\slshape obj})\index{IsInputTextNone@\texttt{IsInputTextNone}}
\label{IsInputTextNone}
}\hfill{\scriptsize (Category)}}\\


 It is convenient to use a category to distinguish dummy streams (see \ref{Dummy Streams}) from others. Other distinctions are usually made using representations }

 

\subsection{\textcolor{Chapter }{IsOutputStream}}
\logpage{[ 10, 1, 6 ]}\nobreak
\hyperdef{L}{X7D357CA07E7B1E78}{}
{\noindent\textcolor{FuncColor}{$\triangleright$\enspace\texttt{IsOutputStream({\mdseries\slshape obj})\index{IsOutputStream@\texttt{IsOutputStream}}
\label{IsOutputStream}
}\hfill{\scriptsize (Category)}}\\


 All output streams lie in this category and support basic operations such as \texttt{WriteByte} (\ref{WriteByte}) (see Section \ref{Operations for Output Streams}). }

 

\subsection{\textcolor{Chapter }{IsOutputTextStream}}
\logpage{[ 10, 1, 7 ]}\nobreak
\hyperdef{L}{X8248B8A4844CB8AB}{}
{\noindent\textcolor{FuncColor}{$\triangleright$\enspace\texttt{IsOutputTextStream({\mdseries\slshape obj})\index{IsOutputTextStream@\texttt{IsOutputTextStream}}
\label{IsOutputTextStream}
}\hfill{\scriptsize (Category)}}\\


 All \emph{text} output streams lie in this category and translate new\texttt{\symbol{45}}line
characters on output. }

 

\subsection{\textcolor{Chapter }{IsOutputTextNone}}
\logpage{[ 10, 1, 8 ]}\nobreak
\hyperdef{L}{X7C89CDD47E33E741}{}
{\noindent\textcolor{FuncColor}{$\triangleright$\enspace\texttt{IsOutputTextNone({\mdseries\slshape obj})\index{IsOutputTextNone@\texttt{IsOutputTextNone}}
\label{IsOutputTextNone}
}\hfill{\scriptsize (Category)}}\\


 It is convenient to use a category to distinguish dummy streams (see \ref{Dummy Streams}) from others. Other distinctions are usually made using representations }

 

\subsection{\textcolor{Chapter }{StreamsFamily}}
\logpage{[ 10, 1, 9 ]}\nobreak
\hyperdef{L}{X7F0F9DD47DE16DAB}{}
{\noindent\textcolor{FuncColor}{$\triangleright$\enspace\texttt{StreamsFamily\index{StreamsFamily@\texttt{StreamsFamily}}
\label{StreamsFamily}
}\hfill{\scriptsize (family)}}\\


 All streams lie in the \texttt{StreamsFamily}. }

 }

  
\section{\textcolor{Chapter }{Operations applicable to All Streams}}\label{Operations applicable to All Streams}
\logpage{[ 10, 2, 0 ]}
\hyperdef{L}{X8461F4DF7FC20C4B}{}
{
  

\subsection{\textcolor{Chapter }{CloseStream}}
\logpage{[ 10, 2, 1 ]}\nobreak
\hyperdef{L}{X786E5520803FDE00}{}
{\noindent\textcolor{FuncColor}{$\triangleright$\enspace\texttt{CloseStream({\mdseries\slshape stream})\index{CloseStream@\texttt{CloseStream}}
\label{CloseStream}
}\hfill{\scriptsize (operation)}}\\


 In order to preserve system resources and to flush output streams every stream
should be closed as soon as it is no longer used using \texttt{CloseStream}. 

 It is an error to try to read characters from or write characters to a closed
stream. Closing a stream tells the \textsf{GAP} kernel and/or the operating system kernel that the file is no longer needed.
This may be necessary because the \textsf{GAP} kernel and/or the operating system may impose a limit on how many streams may
be open simultaneously. }

 

\subsection{\textcolor{Chapter }{FileDescriptorOfStream}}
\logpage{[ 10, 2, 2 ]}\nobreak
\hyperdef{L}{X7F0459287E717456}{}
{\noindent\textcolor{FuncColor}{$\triangleright$\enspace\texttt{FileDescriptorOfStream({\mdseries\slshape stream})\index{FileDescriptorOfStream@\texttt{FileDescriptorOfStream}}
\label{FileDescriptorOfStream}
}\hfill{\scriptsize (operation)}}\\


 returns the UNIX file descriptor of the underlying file. This is mainly useful
for the \texttt{UNIXSelect} (\ref{UNIXSelect}) function call. This is as of now only available on
UNIX\texttt{\symbol{45}}like operating systems and only for streams to local
processes and local files. }

 

\subsection{\textcolor{Chapter }{UNIXSelect}}
\logpage{[ 10, 2, 3 ]}\nobreak
\hyperdef{L}{X87BC257A78F96828}{}
{\noindent\textcolor{FuncColor}{$\triangleright$\enspace\texttt{UNIXSelect({\mdseries\slshape inlist, outlist, exclist, timeoutsec, timeoutusec})\index{UNIXSelect@\texttt{UNIXSelect}}
\label{UNIXSelect}
}\hfill{\scriptsize (function)}}\\


 makes the UNIX C\texttt{\symbol{45}}library function \texttt{select} accessible from \textsf{GAP} for streams. The functionality is as described in the man page (see UNIX file
descriptors (integers) for streams. They can be obtained via \texttt{FileDescriptorOfStream} (\ref{FileDescriptorOfStream}) for streams to local processes and to local files. The argument \mbox{\texttt{\mdseries\slshape timeoutsec}} is a timeout in seconds as in the \texttt{struct timeval} on the C level. The argument \mbox{\texttt{\mdseries\slshape timeoutusec}} is analogously in microseconds. The total timeout is the sum of both. If one
of those timeout arguments is not a small integer then no timeout is
applicable (\texttt{fail} is allowed for the timeout arguments). 

 The return value is the number of streams that are ready, this may be 0 if a
timeout was specified. All file descriptors in the three lists that are not
yet ready are replaced by \texttt{fail} in this function. So the lists are changed! 

 This function is only available if your operating system has \texttt{select}, which is detected during compilation of \textsf{GAP}. }

 }

  
\section{\textcolor{Chapter }{Operations for Input Streams}}\label{Operations for Input Streams}
\logpage{[ 10, 3, 0 ]}
\hyperdef{L}{X7D1D33A587BFD93D}{}
{
  Two operations normally used to read files: \texttt{Read} (\ref{Read}) and \texttt{ReadAsFunction} (\ref{ReadAsFunction}) can also be used to read \textsf{GAP} input from a stream. The input is immediately parsed and executed. When
reading from a stream \mbox{\texttt{\mdseries\slshape str}}, the \textsf{GAP} kernel generates calls to \texttt{ReadLine(\mbox{\texttt{\mdseries\slshape str}})} to supply text to the parser. 

 Three further operations: \texttt{ReadByte} (\ref{ReadByte}), \texttt{ReadLine} (\ref{ReadLine}) and \texttt{ReadAll} (\ref{ReadAll}), support reading characters from an input stream without parsing them. This
can be used to read data in any format and process it in \textsf{GAP}. 

 Additional operations for input streams support detection of end of stream,
and (for those streams for which it is appropriate) random access to the data. 

 

\subsection{\textcolor{Chapter }{Read (for streams)}}
\logpage{[ 10, 3, 1 ]}\nobreak
\hyperdef{L}{X7A5DC83D7E295568}{}
{\noindent\textcolor{FuncColor}{$\triangleright$\enspace\texttt{Read({\mdseries\slshape input\texttt{\symbol{45}}text\texttt{\symbol{45}}stream})\index{Read@\texttt{Read}!for streams}
\label{Read:for streams}
}\hfill{\scriptsize (operation)}}\\


 reads the input\texttt{\symbol{45}}text\texttt{\symbol{45}}stream as input
until \texttt{end\texttt{\symbol{45}}of\texttt{\symbol{45}}stream} occurs. See \ref{File Operations} for details. }

 

\subsection{\textcolor{Chapter }{ReadAsFunction (for streams)}}
\logpage{[ 10, 3, 2 ]}\nobreak
\hyperdef{L}{X7D62F2877F0E45A7}{}
{\noindent\textcolor{FuncColor}{$\triangleright$\enspace\texttt{ReadAsFunction({\mdseries\slshape input\texttt{\symbol{45}}text\texttt{\symbol{45}}stream})\index{ReadAsFunction@\texttt{ReadAsFunction}!for streams}
\label{ReadAsFunction:for streams}
}\hfill{\scriptsize (operation)}}\\


 reads the input\texttt{\symbol{45}}text\texttt{\symbol{45}}stream as function
and returns this function. See \ref{File Operations} for details. 
\begin{Verbatim}[commandchars=!@|,fontsize=\small,frame=single,label=Example]
  !gapprompt@gap>| !gapinput@# a function with local `a' does not change the global one|
  !gapprompt@gap>| !gapinput@a := 1;;|
  !gapprompt@gap>| !gapinput@i := InputTextString( "local a; a := 10; return a*10;" );;|
  !gapprompt@gap>| !gapinput@ReadAsFunction(i)();|
  100
  !gapprompt@gap>| !gapinput@a;|
  1
  !gapprompt@gap>| !gapinput@# reading it via `Read' does|
  !gapprompt@gap>| !gapinput@i := InputTextString( "a := 10;" );;|
  !gapprompt@gap>| !gapinput@Read(i);|
  !gapprompt@gap>| !gapinput@a;|
  10
\end{Verbatim}
 }

 

\subsection{\textcolor{Chapter }{ReadByte}}
\logpage{[ 10, 3, 3 ]}\nobreak
\hyperdef{L}{X79E1E6A57AE58BB8}{}
{\noindent\textcolor{FuncColor}{$\triangleright$\enspace\texttt{ReadByte({\mdseries\slshape input\texttt{\symbol{45}}stream})\index{ReadByte@\texttt{ReadByte}}
\label{ReadByte}
}\hfill{\scriptsize (operation)}}\\


 \texttt{ReadByte} returns one character (returned as integer) from the input stream \mbox{\texttt{\mdseries\slshape input\texttt{\symbol{45}}stream}}. \texttt{ReadByte} returns \texttt{fail} if there is no character available, in particular if it is at the end of a
file. 

 If \mbox{\texttt{\mdseries\slshape input\texttt{\symbol{45}}stream}} is the input stream of a input/output process, \texttt{ReadByte} may also return \texttt{fail} if no byte is currently available. 

 \texttt{ReadByte} is the basic operation for input streams. If a \texttt{ReadByte} method is installed for a user\texttt{\symbol{45}}defined type of stream which
does not block, then all the other input stream operations will work (although
possibly not at peak efficiency). 

 \texttt{ReadByte} will wait (block) until a byte is available. For instance if the stream is a
connection to another process, it will wait for the process to output a byte. }

 

\subsection{\textcolor{Chapter }{ReadLine}}
\logpage{[ 10, 3, 4 ]}\nobreak
\hyperdef{L}{X7D2CA44C7D110C4F}{}
{\noindent\textcolor{FuncColor}{$\triangleright$\enspace\texttt{ReadLine({\mdseries\slshape input\texttt{\symbol{45}}stream})\index{ReadLine@\texttt{ReadLine}}
\label{ReadLine}
}\hfill{\scriptsize (operation)}}\\


 \texttt{ReadLine} returns one line (returned as string \emph{with} the newline) from the input stream \mbox{\texttt{\mdseries\slshape input\texttt{\symbol{45}}stream}}. \texttt{ReadLine} reads in the input until a newline is read or the
end\texttt{\symbol{45}}of\texttt{\symbol{45}}stream is encountered. 

 If \mbox{\texttt{\mdseries\slshape input\texttt{\symbol{45}}stream}} is the input stream of a input/output process, \texttt{ReadLine} may also return \texttt{fail} or return an incomplete line if the other process has not yet written any
more. It will always wait (block) for at least one byte to be available, but
will then return as much input as is available, up to a limit of one line. 

 A default method is supplied for \texttt{ReadLine} which simply calls \texttt{ReadByte} (\ref{ReadByte}) repeatedly. This is only safe for streams that cannot block. The kernel uses
calls to \texttt{ReadLine} to supply input to the parser when reading from a stream. }

 

\subsection{\textcolor{Chapter }{ReadAll}}
\logpage{[ 10, 3, 5 ]}\nobreak
\hyperdef{L}{X85C603D7867430D0}{}
{\noindent\textcolor{FuncColor}{$\triangleright$\enspace\texttt{ReadAll({\mdseries\slshape input\texttt{\symbol{45}}stream[, limit]})\index{ReadAll@\texttt{ReadAll}}
\label{ReadAll}
}\hfill{\scriptsize (operation)}}\\


 \texttt{ReadAll} returns all characters as string from the input stream \mbox{\texttt{\mdseries\slshape stream\texttt{\symbol{45}}in}}. It waits (blocks) until at least one character is available from the stream,
or until there is evidence that no characters will ever be available again.
This last indicates that the stream is at
end\texttt{\symbol{45}}of\texttt{\symbol{45}}stream. Otherwise, it reads as
much input as it can from the stream without blocking further and returns it
to the user. If the stream is already at end of file, so that no bytes are
available, \texttt{fail} is returned. In the case of a file stream connected to a normal file (not a
pseudo\texttt{\symbol{45}}tty or named pipe or similar), all the bytes should
be immediately available and this function will read the remainder of the
file. 

 With a second argument, at most \mbox{\texttt{\mdseries\slshape limit}} bytes will be returned. Depending on the stream a bounded number of additional
bytes may have been read into an internal buffer. 

 A default method is supplied for \texttt{ReadAll} which simply calls \texttt{ReadLine} (\ref{ReadLine}) repeatedly. This is only really safe for streams which cannot block. Other
streams should install a method for \texttt{ReadAll}. 

 
\begin{Verbatim}[commandchars=!@|,fontsize=\small,frame=single,label=Example]
  !gapprompt@gap>| !gapinput@i := InputTextString( "1Hallo\nYou\n1" );;|
  !gapprompt@gap>| !gapinput@ReadByte(i);|
  49
  !gapprompt@gap>| !gapinput@CHAR_INT(last);|
  '1'
  !gapprompt@gap>| !gapinput@ReadLine(i);|
  "Hallo\n"
  !gapprompt@gap>| !gapinput@ReadLine(i);|
  "You\n"
  !gapprompt@gap>| !gapinput@ReadLine(i);|
  "1"
  !gapprompt@gap>| !gapinput@ReadLine(i);|
  fail
  !gapprompt@gap>| !gapinput@ReadAll(i);|
  ""
  !gapprompt@gap>| !gapinput@RewindStream(i);;|
  !gapprompt@gap>| !gapinput@ReadAll(i);|
  "1Hallo\nYou\n1"
\end{Verbatim}
 }

 

\subsection{\textcolor{Chapter }{IsEndOfStream}}
\logpage{[ 10, 3, 6 ]}\nobreak
\hyperdef{L}{X81D4FB097F631A79}{}
{\noindent\textcolor{FuncColor}{$\triangleright$\enspace\texttt{IsEndOfStream({\mdseries\slshape input\texttt{\symbol{45}}stream})\index{IsEndOfStream@\texttt{IsEndOfStream}}
\label{IsEndOfStream}
}\hfill{\scriptsize (operation)}}\\


 \texttt{IsEndOfStream} returns \texttt{true} if the input stream is at \emph{end\texttt{\symbol{45}}of\texttt{\symbol{45}}stream}, and \texttt{false} otherwise. Note that \texttt{IsEndOfStream} might return \texttt{false} even if the next \texttt{ReadByte} (\ref{ReadByte}) fails. }

 

\subsection{\textcolor{Chapter }{PositionStream}}
\logpage{[ 10, 3, 7 ]}\nobreak
\hyperdef{L}{X7B646FA3860521D1}{}
{\noindent\textcolor{FuncColor}{$\triangleright$\enspace\texttt{PositionStream({\mdseries\slshape input\texttt{\symbol{45}}stream})\index{PositionStream@\texttt{PositionStream}}
\label{PositionStream}
}\hfill{\scriptsize (operation)}}\\


 Some input streams, such as string streams and file streams attached to disk
files, support a form of random access by way of the operations \texttt{PositionStream}, \texttt{SeekPositionStream} (\ref{SeekPositionStream}) and \texttt{RewindStream} (\ref{RewindStream}). \texttt{PositionStream} returns a non\texttt{\symbol{45}}negative integer denoting the current
position in the stream (usually the number of characters \emph{before} the next one to be read. 

 If this is not possible, for example for an input stream attached to standard
input (normally the keyboard), then \texttt{fail} is returned }

 

\subsection{\textcolor{Chapter }{RewindStream}}
\logpage{[ 10, 3, 8 ]}\nobreak
\hyperdef{L}{X7A777E1186EB330B}{}
{\noindent\textcolor{FuncColor}{$\triangleright$\enspace\texttt{RewindStream({\mdseries\slshape input\texttt{\symbol{45}}stream})\index{RewindStream@\texttt{RewindStream}}
\label{RewindStream}
}\hfill{\scriptsize (operation)}}\\


 \texttt{RewindStream} attempts to return an input stream to its starting condition, so that all the
same characters can be read again. It returns \texttt{true} if the rewind succeeds and \texttt{fail} otherwise 

 A default method implements RewindStream using \texttt{SeekPositionStream} (\ref{SeekPositionStream}). }

 

\subsection{\textcolor{Chapter }{SeekPositionStream}}
\logpage{[ 10, 3, 9 ]}\nobreak
\hyperdef{L}{X7A60AD8C7E0D0507}{}
{\noindent\textcolor{FuncColor}{$\triangleright$\enspace\texttt{SeekPositionStream({\mdseries\slshape input\texttt{\symbol{45}}stream, pos})\index{SeekPositionStream@\texttt{SeekPositionStream}}
\label{SeekPositionStream}
}\hfill{\scriptsize (operation)}}\\


 \texttt{SeekPositionStream} attempts to rewind or wind forward an input stream to the specified position.
This is not possible for all streams. It returns \texttt{true} if the seek is successful and \texttt{fail} otherwise. }

 }

  
\section{\textcolor{Chapter }{Operations for Output Streams}}\label{Operations for Output Streams}
\logpage{[ 10, 4, 0 ]}
\hyperdef{L}{X7F454EB286947C85}{}
{
  

\subsection{\textcolor{Chapter }{WriteByte}}
\logpage{[ 10, 4, 1 ]}\nobreak
\hyperdef{L}{X7D37C7A07E9C319C}{}
{\noindent\textcolor{FuncColor}{$\triangleright$\enspace\texttt{WriteByte({\mdseries\slshape output\texttt{\symbol{45}}stream, byte})\index{WriteByte@\texttt{WriteByte}}
\label{WriteByte}
}\hfill{\scriptsize (operation)}}\\


 writes the next character (given as \emph{integer}) to the output stream \mbox{\texttt{\mdseries\slshape output\texttt{\symbol{45}}stream}}. The function returns \texttt{true} if the write succeeds and \texttt{fail} otherwise. 

 \texttt{WriteByte} is the basic operation for output streams. If a \texttt{WriteByte} method is installed for a user\texttt{\symbol{45}}defined type of stream, then
all the other output stream operations will work (although possibly not at
peak efficiency). }

 

\subsection{\textcolor{Chapter }{WriteLine}}
\logpage{[ 10, 4, 2 ]}\nobreak
\hyperdef{L}{X79FA85498596CC99}{}
{\noindent\textcolor{FuncColor}{$\triangleright$\enspace\texttt{WriteLine({\mdseries\slshape output\texttt{\symbol{45}}stream, string})\index{WriteLine@\texttt{WriteLine}}
\label{WriteLine}
}\hfill{\scriptsize (operation)}}\\


 appends \mbox{\texttt{\mdseries\slshape string}} to \mbox{\texttt{\mdseries\slshape output\texttt{\symbol{45}}stream}}. A final newline is written. The function returns \texttt{true} if the write succeeds and \texttt{fail} otherwise. 

 A default method is installed which implements \texttt{WriteLine} by repeated calls to \texttt{WriteByte} (\ref{WriteByte}). }

 

\subsection{\textcolor{Chapter }{WriteAll}}
\logpage{[ 10, 4, 3 ]}\nobreak
\hyperdef{L}{X78C113917936058D}{}
{\noindent\textcolor{FuncColor}{$\triangleright$\enspace\texttt{WriteAll({\mdseries\slshape output\texttt{\symbol{45}}stream, string})\index{WriteAll@\texttt{WriteAll}}
\label{WriteAll}
}\hfill{\scriptsize (operation)}}\\


 appends \mbox{\texttt{\mdseries\slshape string}} to \mbox{\texttt{\mdseries\slshape output\texttt{\symbol{45}}stream}}. No final newline is written. The function returns \texttt{true} if the write succeeds and \texttt{fail} otherwise. It will block as long as necessary for the write operation to
complete (for example for a child process to clear its input buffer ) 

 A default method is installed which implements \texttt{WriteAll} by repeated calls to \texttt{WriteByte} (\ref{WriteByte}). 

 When printing or appending to a stream (using \texttt{PrintTo} (\ref{PrintTo}), or \texttt{AppendTo} (\ref{AppendTo}) or when logging to a stream), the kernel generates a call to \texttt{WriteAll} for each line output. 

 
\begin{Verbatim}[commandchars=!@|,fontsize=\small,frame=single,label=Example]
  !gapprompt@gap>| !gapinput@str := "";; a := OutputTextString(str,true);;|
  !gapprompt@gap>| !gapinput@WriteByte(a,INT_CHAR('H'));|
  true
  !gapprompt@gap>| !gapinput@WriteLine(a,"allo");|
  true
  !gapprompt@gap>| !gapinput@WriteAll(a,"You\n");|
  true
  !gapprompt@gap>| !gapinput@CloseStream(a);|
  !gapprompt@gap>| !gapinput@Print(str);|
  Hallo
  You
\end{Verbatim}
 }

 
\subsection{\textcolor{Chapter }{PrintTo and AppendTo (for streams)}}\logpage{[ 10, 4, 4 ]}
\hyperdef{L}{X7F4E090C86AACCF7}{}
{
\noindent\textcolor{FuncColor}{$\triangleright$\enspace\texttt{PrintTo({\mdseries\slshape output\texttt{\symbol{45}}stream, arg1, ...})\index{PrintTo@\texttt{PrintTo}!for streams}
\label{PrintTo:for streams}
}\hfill{\scriptsize (function)}}\\
\noindent\textcolor{FuncColor}{$\triangleright$\enspace\texttt{AppendTo({\mdseries\slshape output\texttt{\symbol{45}}stream, arg1, ...})\index{AppendTo@\texttt{AppendTo}!for streams}
\label{AppendTo:for streams}
}\hfill{\scriptsize (function)}}\\


 These functions work like \texttt{Print} (\ref{Print}), except that the output is appended to the output stream \mbox{\texttt{\mdseries\slshape output\texttt{\symbol{45}}stream}}. 

 
\begin{Verbatim}[commandchars=!@|,fontsize=\small,frame=single,label=Example]
  !gapprompt@gap>| !gapinput@str := "";; a := OutputTextString(str,true);;|
  !gapprompt@gap>| !gapinput@AppendTo( a, (1,2,3), ":", Z(3) );|
  !gapprompt@gap>| !gapinput@CloseStream(a);|
  !gapprompt@gap>| !gapinput@Print( str, "\n" );|
  (1,2,3):Z(3)
\end{Verbatim}
 }

 

\subsection{\textcolor{Chapter }{LogTo (for streams)}}
\logpage{[ 10, 4, 5 ]}\nobreak
\hyperdef{L}{X7BF4E44C7D51E085}{}
{\noindent\textcolor{FuncColor}{$\triangleright$\enspace\texttt{LogTo({\mdseries\slshape stream})\index{LogTo@\texttt{LogTo}!for streams}
\label{LogTo:for streams}
}\hfill{\scriptsize (operation)}}\\


 causes the subsequent interaction to be logged to the output stream \mbox{\texttt{\mdseries\slshape stream}}. It works in precisely the same way as it does for files (see{\nobreakspace}\texttt{LogTo} (\ref{LogTo:for a filename})). }

 

\subsection{\textcolor{Chapter }{InputLogTo (for streams)}}
\logpage{[ 10, 4, 6 ]}\nobreak
\hyperdef{L}{X7B843516796B2A18}{}
{\noindent\textcolor{FuncColor}{$\triangleright$\enspace\texttt{InputLogTo({\mdseries\slshape stream})\index{InputLogTo@\texttt{InputLogTo}!for streams}
\label{InputLogTo:for streams}
}\hfill{\scriptsize (operation)}}\\


 causes the subsequent input to be logged to the output stream \mbox{\texttt{\mdseries\slshape stream}}. It works just like it does for files (see{\nobreakspace}\texttt{InputLogTo} (\ref{InputLogTo:for a filename})). }

 

\subsection{\textcolor{Chapter }{OutputLogTo (for streams)}}
\logpage{[ 10, 4, 7 ]}\nobreak
\hyperdef{L}{X834A6DD17B0E2062}{}
{\noindent\textcolor{FuncColor}{$\triangleright$\enspace\texttt{OutputLogTo({\mdseries\slshape stream})\index{OutputLogTo@\texttt{OutputLogTo}!for streams}
\label{OutputLogTo:for streams}
}\hfill{\scriptsize (operation)}}\\


 causes the subsequent output to be logged to the output stream \mbox{\texttt{\mdseries\slshape stream}}. It works just like it does for files (see{\nobreakspace}\texttt{OutputLogTo} (\ref{OutputLogTo:for a filename})). 

 }

 

\subsection{\textcolor{Chapter }{SetPrintFormattingStatus}}
\logpage{[ 10, 4, 8 ]}\nobreak
\hyperdef{L}{X8663FCD57E8BC390}{}
{\noindent\textcolor{FuncColor}{$\triangleright$\enspace\texttt{SetPrintFormattingStatus({\mdseries\slshape stream, newstatus})\index{SetPrintFormattingStatus@\texttt{SetPrintFormattingStatus}}
\label{SetPrintFormattingStatus}
}\hfill{\scriptsize (operation)}}\\
\noindent\textcolor{FuncColor}{$\triangleright$\enspace\texttt{PrintFormattingStatus({\mdseries\slshape stream})\index{PrintFormattingStatus@\texttt{PrintFormattingStatus}}
\label{PrintFormattingStatus}
}\hfill{\scriptsize (operation)}}\\


 When text is being sent to an output text stream via \texttt{PrintTo} (\ref{PrintTo}), \texttt{AppendTo} (\ref{AppendTo}), \texttt{LogTo} (\ref{LogTo:for streams}), etc., it is by default formatted just as it would be were it being printed
to the screen. Thus, it is broken into lines of reasonable length at (where
possible) sensible places, lines containing elements of lists or records are
indented, and so forth. This is appropriate if the output is eventually to be
viewed by a human, and harmless if it to passed as input to \textsf{GAP}, but may be unhelpful if the output is to be passed as input to another
program. It is possible to turn off this behaviour for a stream using the \texttt{SetPrintFormattingStatus} operation, and to test whether it is on or off using \texttt{PrintFormattingStatus}. 

 \texttt{SetPrintFormattingStatus} sets whether output sent to the output stream \mbox{\texttt{\mdseries\slshape stream}} via \texttt{PrintTo} (\ref{PrintTo}), \texttt{AppendTo} (\ref{AppendTo}), etc. will be formatted with line breaks and indentation. If the second
argument \mbox{\texttt{\mdseries\slshape newstatus}} is \texttt{true} then output will be so formatted, and if \texttt{false} then it will not. If the stream is not a text stream, only \texttt{false} is allowed. 

 \texttt{PrintFormattingStatus} returns \texttt{true} if output sent to the output text stream \mbox{\texttt{\mdseries\slshape stream}} via \texttt{PrintTo} (\ref{PrintTo}), \texttt{AppendTo} (\ref{AppendTo}), etc. will be formatted with line breaks and indentation, and \texttt{false} otherwise. For non\texttt{\symbol{45}}text streams, it returns \texttt{false}. If as argument \mbox{\texttt{\mdseries\slshape stream}} the string \texttt{"*stdout*"} is given, these functions refer to the formatting status of the standard
output (so usually the user's terminal screen).

 Similarly, the string \texttt{"*errout*"} refers to the formatting status of the standard error output, which influences
how error messages are printed.

 These functions do not influence the behaviour of the low level functions \texttt{WriteByte} (\ref{WriteByte}), \texttt{WriteLine} (\ref{WriteLine}) or \texttt{WriteAll} (\ref{WriteAll}) which always write without formatting. 

 
\begin{Verbatim}[commandchars=!@|,fontsize=\small,frame=single,label=Example]
  !gapprompt@gap>| !gapinput@s := "";; str := OutputTextString(s,false);;|
  !gapprompt@gap>| !gapinput@PrintTo(str,Primes{[1..30]});|
  !gapprompt@gap>| !gapinput@s;|
  "[ 2, 3, 5, 7, 11, 13, 17, 19, 23, 29, 31, 37, 41, 43, 47, 53, 59, 61,\
   \n  67, 71, 73, 79, 83, 89, 97, 101, 103, 107, 109, 113 ]"
  !gapprompt@gap>| !gapinput@Print(s,"\n");|
  [ 2, 3, 5, 7, 11, 13, 17, 19, 23, 29, 31, 37, 41, 43, 47, 53, 59, 61,
    67, 71, 73, 79, 83, 89, 97, 101, 103, 107, 109, 113 ]
  !gapprompt@gap>| !gapinput@SetPrintFormattingStatus(str, false);|
  !gapprompt@gap>| !gapinput@PrintTo(str,Primes{[1..30]});|
  !gapprompt@gap>| !gapinput@s;|
  "[ 2, 3, 5, 7, 11, 13, 17, 19, 23, 29, 31, 37, 41, 43, 47, 53, 59, 61,\
   \n  67, 71, 73, 79, 83, 89, 97, 101, 103, 107, 109, 113 ][ 2, 3, 5, 7\
  , 11, 13, 17, 19, 23, 29, 31, 37, 41, 43, 47, 53, 59, 61, 67, 71, 73, \
  79, 83, 89, 97, 101, 103, 107, 109, 113 ]"
  !gapprompt@gap>| !gapinput@Print(s,"\n");|
  [ 2, 3, 5, 7, 11, 13, 17, 19, 23, 29, 31, 37, 41, 43, 47, 53, 59, 61,
    67, 71, 73, 79, 83, 89, 97, 101, 103, 107, 109, 113 ][ 2, 3, 5, 7, 1\
  1, 13, 17, 19, 23, 29, 31, 37, 41, 43, 47, 53, 59, 61, 67, 71, 73, 79,\
   83, 89, 97, 101, 103, 107, 109, 113 ]
\end{Verbatim}
 

 }

 }

  
\section{\textcolor{Chapter }{File Streams}}\label{File Streams}
\logpage{[ 10, 5, 0 ]}
\hyperdef{L}{X80B5F2E4856D8980}{}
{
  File streams are streams associated with files. An input file stream reads the
characters it delivers from a file, an output file stream prints the
characters it receives to a file. The following functions can be used to
create such streams. They return \texttt{fail} if an error occurred, in this case \texttt{LastSystemError} (\ref{LastSystemError}) can be used to get information about the error. 

\subsection{\textcolor{Chapter }{InputTextFile}}
\logpage{[ 10, 5, 1 ]}\nobreak
\hyperdef{L}{X8343D04981128784}{}
{\noindent\textcolor{FuncColor}{$\triangleright$\enspace\texttt{InputTextFile({\mdseries\slshape filename})\index{InputTextFile@\texttt{InputTextFile}}
\label{InputTextFile}
}\hfill{\scriptsize (operation)}}\\


 \texttt{InputTextFile( \mbox{\texttt{\mdseries\slshape filename}} )} returns an input stream in the category \texttt{IsInputTextStream} (\ref{IsInputTextStream}) that delivers the characters from the file \mbox{\texttt{\mdseries\slshape filename}}. If \mbox{\texttt{\mdseries\slshape filename}} ends in \texttt{.gz} and the file is a valid gzipped file, then the file will be transparently
uncompressed. 

 \texttt{InputTextFile} is designed for use with text files and automatically handles
windows\texttt{\symbol{45}}style line endings. This means it should \emph{not} be used for binary data. The \texttt{IO{\textunderscore}File} (\textbf{IO: IO{\textunderscore}File mode}) function from the \textsf{IO} package should be used to access binary data. 

 Note: At most 256 files may be open for reading or writing at the same time.
Use \texttt{CloseStream} (\ref{CloseStream}) to close the input stream once you have finished reading from it. }

 

\subsection{\textcolor{Chapter }{OutputTextFile}}
\logpage{[ 10, 5, 2 ]}\nobreak
\hyperdef{L}{X83F53291822B7126}{}
{\noindent\textcolor{FuncColor}{$\triangleright$\enspace\texttt{OutputTextFile({\mdseries\slshape filename, append})\index{OutputTextFile@\texttt{OutputTextFile}}
\label{OutputTextFile}
}\hfill{\scriptsize (operation)}}\\
\noindent\textcolor{FuncColor}{$\triangleright$\enspace\texttt{OutputGzipFile({\mdseries\slshape filename, append})\index{OutputGzipFile@\texttt{OutputGzipFile}}
\label{OutputGzipFile}
}\hfill{\scriptsize (operation)}}\\


 \texttt{OutputTextFile( \mbox{\texttt{\mdseries\slshape filename}}, \mbox{\texttt{\mdseries\slshape append}} )} returns an output stream in the category \texttt{IsOutputTextFile} that writes received characters to the file \mbox{\texttt{\mdseries\slshape filename}}. If \mbox{\texttt{\mdseries\slshape append}} is \texttt{false}, then the file is emptied first, otherwise received characters are added at
the end of the file. \texttt{OutputGzipFile} acts identically to \texttt{OutputTextFile}, except it compresses the output with gzip. 

 Note: At most 256 files may be open for reading or writing at the same time.
Use \texttt{CloseStream} (\ref{CloseStream}) to close the output stream once you have finished writing to it. 

 
\begin{Verbatim}[commandchars=!@|,fontsize=\small,frame=single,label=Example]
  !gapprompt@gap>| !gapinput@# use a temporary directory|
  !gapprompt@gap>| !gapinput@name := Filename( DirectoryTemporary(), "test" );;|
  !gapprompt@gap>| !gapinput@# create an output stream, append output, and close again|
  !gapprompt@gap>| !gapinput@output := OutputTextFile( name, true );;|
  !gapprompt@gap>| !gapinput@AppendTo( output, "Hallo\n", "You\n" );|
  !gapprompt@gap>| !gapinput@CloseStream(output);|
  !gapprompt@gap>| !gapinput@# create an input, print complete contents of file, and close|
  !gapprompt@gap>| !gapinput@input := InputTextFile(name);;|
  !gapprompt@gap>| !gapinput@Print( ReadAll(input) );|
  Hallo
  You
  !gapprompt@gap>| !gapinput@CloseStream(input);|
  !gapprompt@gap>| !gapinput@# append a single line|
  !gapprompt@gap>| !gapinput@output := OutputTextFile( name, true );;|
  !gapprompt@gap>| !gapinput@AppendTo( output, "AppendLine\n" );|
  !gapprompt@gap>| !gapinput@# close output stream to flush the output|
  !gapprompt@gap>| !gapinput@CloseStream(output);|
  !gapprompt@gap>| !gapinput@# create an input, print complete contents of file, and close|
  !gapprompt@gap>| !gapinput@input := InputTextFile(name);;|
  !gapprompt@gap>| !gapinput@Print( ReadAll(input) );|
  Hallo
  You
  AppendLine
  !gapprompt@gap>| !gapinput@CloseStream(input);|
\end{Verbatim}
 }

 }

  
\section{\textcolor{Chapter }{User Streams}}\label{User Streams}
\logpage{[ 10, 6, 0 ]}
\hyperdef{L}{X808348977A05477A}{}
{
  The commands described in this section create streams which accept characters
from, or deliver characters to, the user, via the keyboard or the \textsf{GAP} session display. 

\subsection{\textcolor{Chapter }{InputTextUser}}
\logpage{[ 10, 6, 1 ]}\nobreak
\hyperdef{L}{X83531E4C7C53544F}{}
{\noindent\textcolor{FuncColor}{$\triangleright$\enspace\texttt{InputTextUser({\mdseries\slshape })\index{InputTextUser@\texttt{InputTextUser}}
\label{InputTextUser}
}\hfill{\scriptsize (function)}}\\


 returns an input text stream which delivers characters typed by the user (or
from the standard input device if it has been redirected). In normal
circumstances, characters are delivered one by one as they are typed, without
waiting until the end of a line. No prompts are printed. }

 

\subsection{\textcolor{Chapter }{OutputTextUser}}
\logpage{[ 10, 6, 2 ]}\nobreak
\hyperdef{L}{X83E5FC9487766297}{}
{\noindent\textcolor{FuncColor}{$\triangleright$\enspace\texttt{OutputTextUser({\mdseries\slshape })\index{OutputTextUser@\texttt{OutputTextUser}}
\label{OutputTextUser}
}\hfill{\scriptsize (function)}}\\


 returns an output stream which delivers characters to the user's display (or
the standard output device if it has been redirected). Each character is
delivered immediately it is written, without waiting for a full line of
output. Text written in this way is \emph{not} written to the session log (see \texttt{LogTo} (\ref{LogTo:for a filename})). }

 

\subsection{\textcolor{Chapter }{InputFromUser}}
\logpage{[ 10, 6, 3 ]}\nobreak
\hyperdef{L}{X7DAF5B7085F4F893}{}
{\noindent\textcolor{FuncColor}{$\triangleright$\enspace\texttt{InputFromUser({\mdseries\slshape arg})\index{InputFromUser@\texttt{InputFromUser}}
\label{InputFromUser}
}\hfill{\scriptsize (function)}}\\


 prints the \mbox{\texttt{\mdseries\slshape arg}} as a prompt, then waits until a text is typed by the user (or from the
standard input device if it has been redirected). This text must be a \emph{single} expression, followed by one \emph{enter}. This is evaluated (see \texttt{EvalString} (\ref{EvalString})) and the result is returned. }

 }

  
\section{\textcolor{Chapter }{String Streams}}\label{String Streams}
\logpage{[ 10, 7, 0 ]}
\hyperdef{L}{X8028E1D87CE2F059}{}
{
  String streams are streams associated with strings. An input string stream
reads the characters it delivers from a string, an output string stream
appends the characters it receives to a string. The following functions can be
used to create such streams. 

\subsection{\textcolor{Chapter }{InputTextString}}
\logpage{[ 10, 7, 1 ]}\nobreak
\hyperdef{L}{X7ABABCDF7ED81F7F}{}
{\noindent\textcolor{FuncColor}{$\triangleright$\enspace\texttt{InputTextString({\mdseries\slshape string})\index{InputTextString@\texttt{InputTextString}}
\label{InputTextString}
}\hfill{\scriptsize (operation)}}\\


 \texttt{InputTextString( \mbox{\texttt{\mdseries\slshape string}} )} returns an input stream that delivers the characters from the string \mbox{\texttt{\mdseries\slshape string}}. The \mbox{\texttt{\mdseries\slshape string}} is not changed when reading characters from it and changing the \mbox{\texttt{\mdseries\slshape string}} after the call to \texttt{InputTextString} has no influence on the input stream. }

 

\subsection{\textcolor{Chapter }{OutputTextString}}
\logpage{[ 10, 7, 2 ]}\nobreak
\hyperdef{L}{X7FEDA5167979B74D}{}
{\noindent\textcolor{FuncColor}{$\triangleright$\enspace\texttt{OutputTextString({\mdseries\slshape list, append})\index{OutputTextString@\texttt{OutputTextString}}
\label{OutputTextString}
}\hfill{\scriptsize (operation)}}\\


 returns an output stream that puts all received characters into the list \mbox{\texttt{\mdseries\slshape list}}. If \mbox{\texttt{\mdseries\slshape append}} is \texttt{false}, then the list is emptied first, otherwise received characters are added at
the end of the list. 

 
\begin{Verbatim}[commandchars=!@|,fontsize=\small,frame=single,label=Example]
  !gapprompt@gap>| !gapinput@# read input from a string|
  !gapprompt@gap>| !gapinput@input := InputTextString( "Hallo\nYou\n" );;|
  !gapprompt@gap>| !gapinput@ReadLine(input);|
  "Hallo\n"
  !gapprompt@gap>| !gapinput@ReadLine(input);|
  "You\n"
  !gapprompt@gap>| !gapinput@# print to a string|
  !gapprompt@gap>| !gapinput@str := "";;|
  !gapprompt@gap>| !gapinput@out := OutputTextString( str, true );;|
  !gapprompt@gap>| !gapinput@PrintTo( out, 1, "\n", (1,2,3,4)(5,6), "\n" );|
  !gapprompt@gap>| !gapinput@CloseStream(out);|
  !gapprompt@gap>| !gapinput@Print( str );|
  1
  (1,2,3,4)(5,6)
\end{Verbatim}
 }

 }

  
\section{\textcolor{Chapter }{Input\texttt{\symbol{45}}Output Streams}}\label{Input-Output Streams}
\logpage{[ 10, 8, 0 ]}
\hyperdef{L}{X8563EF8387236417}{}
{
  Input\texttt{\symbol{45}}output streams capture bidirectional communications
between \textsf{GAP} and another process, either locally or (@as yet unimplemented@) remotely. 

 Such streams support the basic operations of both input and output streams.
They should provide some buffering, allowing output data to be written to the
stream, even when input data is waiting to be read, but the amount of this
buffering is operating system dependent, and the user should take care not to
get too far ahead in writing, or behind in reading, or deadlock may occur. 

 At present the only type of Input\texttt{\symbol{45}}Output streams that are
implemented provide communication with a local child process, using a
pseudo\texttt{\symbol{45}}tty. 

 Like other streams, write operations are blocking, read operations will block
to get the first character, but not thereafter. 

 As far as possible, no translation is done on characters written to, or read
from the stream, and no control characters have special effects, but the
details of particular pseudo\texttt{\symbol{45}}tty implementations may effect
this. 

\subsection{\textcolor{Chapter }{IsInputOutputStream}}
\logpage{[ 10, 8, 1 ]}\nobreak
\hyperdef{L}{X82822D3D8339F635}{}
{\noindent\textcolor{FuncColor}{$\triangleright$\enspace\texttt{IsInputOutputStream({\mdseries\slshape obj})\index{IsInputOutputStream@\texttt{IsInputOutputStream}}
\label{IsInputOutputStream}
}\hfill{\scriptsize (Category)}}\\


 \texttt{IsInputOutputStream} is the Category of Input\texttt{\symbol{45}}Output Streams; it returns \texttt{true} if the \mbox{\texttt{\mdseries\slshape obj}} is an input\texttt{\symbol{45}}output stream and \texttt{false} otherwise. }

 

\subsection{\textcolor{Chapter }{InputOutputLocalProcess}}
\logpage{[ 10, 8, 2 ]}\nobreak
\hyperdef{L}{X820799A3824684AC}{}
{\noindent\textcolor{FuncColor}{$\triangleright$\enspace\texttt{InputOutputLocalProcess({\mdseries\slshape dir, executable, args})\index{InputOutputLocalProcess@\texttt{InputOutputLocalProcess}}
\label{InputOutputLocalProcess}
}\hfill{\scriptsize (function)}}\\


 starts up a child process, whose executable file is \mbox{\texttt{\mdseries\slshape executable}}, with ``command line'' arguments \mbox{\texttt{\mdseries\slshape args}} in the directory \mbox{\texttt{\mdseries\slshape dir}}. (Suitable choices for \mbox{\texttt{\mdseries\slshape dir}} are \texttt{DirectoryCurrent()} or \texttt{DirectoryTemporary()} (see Section{\nobreakspace}\ref{Directories}); \texttt{DirectoryTemporary()} may be a good choice when \mbox{\texttt{\mdseries\slshape executable}} generates output files that it doesn't itself remove afterwards.) \texttt{InputOutputLocalProcess} returns an InputOutputStream object. Bytes written to this stream are received
by the child process as if typed at a terminal on standard input. Bytes
written to standard output by the child process can be read from the stream. 

 When the stream is closed, the signal SIGTERM is delivered to the child
process, which is expected to exit. 
\begin{Verbatim}[commandchars=!@|,fontsize=\small,frame=single,label=Example]
  !gapprompt@gap>| !gapinput@d := DirectoryCurrent();|
  dir("./")
  !gapprompt@gap>| !gapinput@f := PathSystemProgram("rev");|
  "/usr/bin/rev"
  !gapprompt@gap>| !gapinput@s := InputOutputLocalProcess(d,f,[]);|
  < input/output stream to rev >
  !gapprompt@gap>| !gapinput@WriteLine(s,"The cat sat on the mat");|
  true
  !gapprompt@gap>| !gapinput@Print(ReadLine(s));|
  tam eht no tas tac ehT
  !gapprompt@gap>| !gapinput@x := ListWithIdenticalEntries(10000,'x');;|
  !gapprompt@gap>| !gapinput@ConvertToStringRep(x);|
  !gapprompt@gap>| !gapinput@WriteLine(s,x);|
  true
  !gapprompt@gap>| !gapinput@WriteByte(s,INT_CHAR('\n'));|
  true
  !gapprompt@gap>| !gapinput@y := ReadAll(s);;|
  !gapprompt@gap>| !gapinput@Length(y);|
  10002
  !gapprompt@gap>| !gapinput@CloseStream(s);|
  !gapprompt@gap>| !gapinput@s;|
  < closed input/output stream to rev >
\end{Verbatim}
 }

 

\subsection{\textcolor{Chapter }{ReadAllLine}}
\logpage{[ 10, 8, 3 ]}\nobreak
\hyperdef{L}{X7CDF48447E823977}{}
{\noindent\textcolor{FuncColor}{$\triangleright$\enspace\texttt{ReadAllLine({\mdseries\slshape iostream[, nofail][, IsAllLine]})\index{ReadAllLine@\texttt{ReadAllLine}}
\label{ReadAllLine}
}\hfill{\scriptsize (operation)}}\\


 For an input/output stream \mbox{\texttt{\mdseries\slshape iostream}} \texttt{ReadAllLine} reads until a newline character if any input is found or returns \texttt{fail} if no input is found, i.e.{\nobreakspace}if any input is found \texttt{ReadAllLine} is non\texttt{\symbol{45}}blocking. 

 If the argument \mbox{\texttt{\mdseries\slshape nofail}} (which must be \texttt{false} or \texttt{true}) is provided and it is set to \texttt{true} then \texttt{ReadAllLine} will wait, if necessary, for input and never return \texttt{fail}. 

 If the argument \mbox{\texttt{\mdseries\slshape IsAllLine}} (which must be a function that takes a string argument and returns either \texttt{true} or \texttt{false}) then it is used to determine what constitutes a whole line. The default
behaviour is equivalent to passing the function 

 
\begin{Verbatim}[commandchars=!@|,fontsize=\small,frame=single,label=Example]
  line -> 0 < Length(line) and Last(line) = '\n'
\end{Verbatim}
 

 for the \mbox{\texttt{\mdseries\slshape IsAllLine}} argument. The purpose of the \mbox{\texttt{\mdseries\slshape IsAllLine}} argument is to cater for the case where the input being read is from an
external process that writes a ``prompt'' for data that does not terminate with a newline. 

 If the first argument is an input stream but not an input/output stream then \texttt{ReadAllLine} behaves as if \texttt{ReadLine} (\ref{ReadLine}) was called with just the first argument and any additional arguments are
ignored. }

 }

  
\section{\textcolor{Chapter }{Dummy Streams}}\label{Dummy Streams}
\logpage{[ 10, 9, 0 ]}
\hyperdef{L}{X8724699C7D67BA47}{}
{
  The following two commands create dummy streams which will consume all
characters and never deliver one. 

\subsection{\textcolor{Chapter }{InputTextNone}}
\logpage{[ 10, 9, 1 ]}\nobreak
\hyperdef{L}{X7C732324806716C6}{}
{\noindent\textcolor{FuncColor}{$\triangleright$\enspace\texttt{InputTextNone({\mdseries\slshape })\index{InputTextNone@\texttt{InputTextNone}}
\label{InputTextNone}
}\hfill{\scriptsize (function)}}\\


 returns a dummy input text stream, which delivers no characters, i.e., it is
always at end of stream. Its main use is for calls to \texttt{Process} (\ref{Process}) when the started program does not read anything. }

 

\subsection{\textcolor{Chapter }{OutputTextNone}}
\logpage{[ 10, 9, 2 ]}\nobreak
\hyperdef{L}{X7CC5C1FC81715E38}{}
{\noindent\textcolor{FuncColor}{$\triangleright$\enspace\texttt{OutputTextNone({\mdseries\slshape })\index{OutputTextNone@\texttt{OutputTextNone}}
\label{OutputTextNone}
}\hfill{\scriptsize (function)}}\\


 returns a dummy output stream, which discards all received characters. Its
main use is for calls to \texttt{Process} (\ref{Process}) when the started program does not write anything. }

 }

  
\section{\textcolor{Chapter }{Handling of Streams in the Background}}\label{Handling of Streams in the Background}
\logpage{[ 10, 10, 0 ]}
\hyperdef{L}{X7CB5832F8721ADF3}{}
{
  This section describes a feature of the \textsf{GAP} kernel that can be used to handle pending streams somehow ``in the background''. This is only available on operating systems that have \texttt{select}. 

 Right before \textsf{GAP} reads a keypress from the keyboard it calls a little subroutine that can
handle streams that are ready to be read or ready to be written. This means
that \textsf{GAP} can handle these streams during user input on the command line. Note that this
does not work when \textsf{GAP} is in the middle of some calculation. 

 This feature is used in the following way. One can install handler functions
for reading or writing streams via \texttt{InstallCharReadHookFunc} (\ref{InstallCharReadHookFunc}). Handlers can be removed via \texttt{UnInstallCharReadHookFunc} (\ref{UnInstallCharReadHookFunc}) 

 Note that handler functions must not return anything and get one integer
argument, which refers to an index in one of the following arrays (according
to whether the function was installed for input, output or exceptions on the
stream). Handler functions usually should not output anything on the standard
output because this ruins the command line during command line editing. 

\subsection{\textcolor{Chapter }{InstallCharReadHookFunc}}
\logpage{[ 10, 10, 1 ]}\nobreak
\hyperdef{L}{X81FB42517E3EA96D}{}
{\noindent\textcolor{FuncColor}{$\triangleright$\enspace\texttt{InstallCharReadHookFunc({\mdseries\slshape stream, mode, func})\index{InstallCharReadHookFunc@\texttt{InstallCharReadHookFunc}}
\label{InstallCharReadHookFunc}
}\hfill{\scriptsize (function)}}\\


 installs the function \mbox{\texttt{\mdseries\slshape func}} as a handler function for the stream \mbox{\texttt{\mdseries\slshape stream}}. The argument \mbox{\texttt{\mdseries\slshape mode}} decides, for what operations on the stream this function is installed. \mbox{\texttt{\mdseries\slshape mode}} must be a string, in which a letter \texttt{r} means ``read'', \texttt{w} means ``write'' and \texttt{x} means ``exception'', according to the \texttt{select} function call in the UNIX C\texttt{\symbol{45}}library (see \texttt{man select} and \texttt{UNIXSelect} (\ref{UNIXSelect})). More than one letter is allowed in \mbox{\texttt{\mdseries\slshape mode}}. As described above the function is called in a situation when \textsf{GAP} is reading a character from the keyboard. Handler functions should not use
much time to complete. 

 This functionality only works if the operating system has a \texttt{select} function. }

 

\subsection{\textcolor{Chapter }{UnInstallCharReadHookFunc}}
\logpage{[ 10, 10, 2 ]}\nobreak
\hyperdef{L}{X8492474C7A0B10AD}{}
{\noindent\textcolor{FuncColor}{$\triangleright$\enspace\texttt{UnInstallCharReadHookFunc({\mdseries\slshape stream, func})\index{UnInstallCharReadHookFunc@\texttt{UnInstallCharReadHookFunc}}
\label{UnInstallCharReadHookFunc}
}\hfill{\scriptsize (function)}}\\


 uninstalls the function \mbox{\texttt{\mdseries\slshape func}} as a handler function for the stream \mbox{\texttt{\mdseries\slshape stream}}. All instances are deinstalled, regardless of the mode of operation (read,
write, exception). 

 This functionality only works if the operating system has a \texttt{select} function. }

 }

  
\section{\textcolor{Chapter }{Comma separated files}}\label{Comma separated files}
\logpage{[ 10, 11, 0 ]}
\hyperdef{L}{X848DD7DC79363341}{}
{
  \index{Spreadsheet} \index{Excel} In some situations it can be desirable to process data given in the form of a
spreadsheet (such as Excel). \textsf{GAP} can do this using the CSV (comma separated values) format, which spreadsheet
programs can usually read in or write out. 

 The first line of the spreadsheet is used as labels of record components, each
subsequent line then corresponds to a record. Entries enclosed in double
quotes are considered as strings and are permitted to contain the separation
character (usually a comma). 

\subsection{\textcolor{Chapter }{ReadCSV}}
\logpage{[ 10, 11, 1 ]}\nobreak
\hyperdef{L}{X86FDC1EF82CAD2DA}{}
{\noindent\textcolor{FuncColor}{$\triangleright$\enspace\texttt{ReadCSV({\mdseries\slshape filename[, nohead][, separator]})\index{ReadCSV@\texttt{ReadCSV}}
\label{ReadCSV}
}\hfill{\scriptsize (function)}}\\


 This function reads in a spreadsheet, saved in CSV format (\emph{c}omma \emph{s}eparated \emph{v}alues) and returns its entries as a list of records. The entries of the first
line of the spreadsheet are used to denote the names of the record components.
Blanks will be translated into underscore characters. If the parameter \mbox{\texttt{\mdseries\slshape nohead}} is given as \texttt{true}, instead the record components will be called \texttt{fieldn}. Each subsequent line will create one record. If given, \mbox{\texttt{\mdseries\slshape separator}} is the character used to separate fields. Otherwise it defaults to a comma. }

 

\subsection{\textcolor{Chapter }{PrintCSV}}
\logpage{[ 10, 11, 2 ]}\nobreak
\hyperdef{L}{X8779DAC585E05A47}{}
{\noindent\textcolor{FuncColor}{$\triangleright$\enspace\texttt{PrintCSV({\mdseries\slshape filename, list[, fields]})\index{PrintCSV@\texttt{PrintCSV}}
\label{PrintCSV}
}\hfill{\scriptsize (function)}}\\


 This function prints a list of records as a spreadsheet in CSV format (which
can be read in for example into Excel). The names of the record components
will be printed as entries in the first line. If the argument \mbox{\texttt{\mdseries\slshape fields}} is given only the record fields listed in this list will be printed and they
will be printed in the same arrangement as given in this list. If the option
noheader is set to true the line with the record field names will not be
printed. }

 }

  
\section{\textcolor{Chapter }{Opening files in the Operating System}}\label{Opening files in the Operating System}
\logpage{[ 10, 12, 0 ]}
\hyperdef{L}{X87396F857ADA3F97}{}
{
  In some situations it can be desirable to open a file outside of \textsf{GAP}, for example HTML files, PDFs, or pictures. 

\subsection{\textcolor{Chapter }{OpenExternal}}
\logpage{[ 10, 12, 1 ]}\nobreak
\hyperdef{L}{X86B98E287AD42BE8}{}
{\noindent\textcolor{FuncColor}{$\triangleright$\enspace\texttt{OpenExternal({\mdseries\slshape filename})\index{OpenExternal@\texttt{OpenExternal}}
\label{OpenExternal}
}\hfill{\scriptsize (function)}}\\


 Open the file \mbox{\texttt{\mdseries\slshape filename}} using the default application for this file in the operating system. This can
be used to open files like HTML and PDF files in the GUI. }

 }

 }

               
\chapter{\textcolor{Chapter }{Processes}}\label{Processes}
\logpage{[ 11, 0, 0 ]}
\hyperdef{L}{X7882133B7BDD51BC}{}
{
  \textsf{GAP} can call other programs, such programs are called \emph{processes}. There are two kinds of processes: first there are processes that are
started, run and return a result, while \textsf{GAP} is suspended until the process terminates. Then there are processes that will
run in parallel to \textsf{GAP} as subprocesses and \textsf{GAP} can communicate and control the processes using streams (see{\nobreakspace}\texttt{InputOutputLocalProcess} (\ref{InputOutputLocalProcess})).  
\section{\textcolor{Chapter }{Process and Exec}}\label{Process and Exec}
\logpage{[ 11, 1, 0 ]}
\hyperdef{L}{X8390266186E61CCE}{}
{
  

\subsection{\textcolor{Chapter }{Process}}
\logpage{[ 11, 1, 1 ]}\nobreak
\hyperdef{L}{X7B09033178D1107A}{}
{\noindent\textcolor{FuncColor}{$\triangleright$\enspace\texttt{Process({\mdseries\slshape dir, prg, stream\texttt{\symbol{45}}in, stream\texttt{\symbol{45}}out, options})\index{Process@\texttt{Process}}
\label{Process}
}\hfill{\scriptsize (operation)}}\\


 \texttt{Process} runs a new process and returns when the process terminates. It returns the
return value of the process if the operating system supports such a concept. 

 The first argument \mbox{\texttt{\mdseries\slshape dir}} is a directory object (see{\nobreakspace}\ref{Directories}) which will be the current directory (in the usual UNIX or
MS\texttt{\symbol{45}}DOS sense) when the program is run. This will only
matter if the program accesses files (including running other programs) via
relative path names. In particular, it has nothing to do with finding the
binary to run. 

 In general the directory will either be the current directory, which is
returned by \texttt{DirectoryCurrent} (\ref{DirectoryCurrent}) {\textendash}this was the behaviour of \textsf{GAP}{\nobreakspace}3{\textendash} or a temporary directory returned by \texttt{DirectoryTemporary} (\ref{DirectoryTemporary}). If one expects that the process creates temporary or log files the latter
should be used because \textsf{GAP} will attempt to remove these directories together with all the files in them
when quitting. 

 If a program of a \textsf{GAP} package which does not only consist of \textsf{GAP} code needs to be launched in a directory relative to certain data libraries,
then the first entry of \texttt{DirectoriesPackageLibrary} (\ref{DirectoriesPackageLibrary}) should be used. The argument of \texttt{DirectoriesPackageLibrary} (\ref{DirectoriesPackageLibrary}) should be the path to the data library relative to the package directory. 

 If a program calls other programs and needs to be launched in a directory
containing the executables for such a \textsf{GAP} package then the first entry of \texttt{DirectoriesPackagePrograms} (\ref{DirectoriesPackagePrograms}) should be used. 

 The latter two alternatives should only be used if absolutely necessary
because otherwise one risks accumulating log or core files in the package
directory. 

 
\begin{Verbatim}[commandchars=!@|,fontsize=\small,frame=single,label=Example]
  !gapprompt@gap>| !gapinput@ls := PathSystemProgram( "ls" );;|
  !gapprompt@gap>| !gapinput@stdin := InputTextUser();;|
  !gapprompt@gap>| !gapinput@stdout := OutputTextUser();;|
  !gapprompt@gap>| !gapinput@path := DirectoriesSystemPrograms();;|
  !gapprompt@gap>| !gapinput@Process( path[1], ls, stdin, stdout, ["-c"] );;|
  awk    ls     mkdir
  !gapprompt@gap>| !gapinput@# current directory, here the root directory|
  !gapprompt@gap>| !gapinput@Process( DirectoryCurrent(), ls, stdin, stdout, ["-c"] );;|
  bin    lib    trans  tst    CVS    grp    prim   thr    two
  src    dev    etc    tbl    doc    pkg    small  tom
  !gapprompt@gap>| !gapinput@# create a temporary directory|
  !gapprompt@gap>| !gapinput@tmpdir := DirectoryTemporary();;|
  !gapprompt@gap>| !gapinput@Process( tmpdir, ls, stdin, stdout, ["-c"] );;|
  !gapprompt@gap>| !gapinput@PrintTo( Filename( tmpdir, "emil" ) );|
  !gapprompt@gap>| !gapinput@Process( tmpdir, ls, stdin, stdout, ["-c"] );;|
  emil
\end{Verbatim}
 

 \mbox{\texttt{\mdseries\slshape prg}} is the filename of the program to launch, for portability it should be the
result of \texttt{Filename} (\ref{Filename:for a directory and a string}) and should pass \texttt{IsExecutableFile} (\ref{IsExecutableFile}). Note that \texttt{Process} does \emph{no} searching through a list of directories, this is done by \texttt{Filename} (\ref{Filename:for a directory and a string}). 

 \mbox{\texttt{\mdseries\slshape stream\texttt{\symbol{45}}in}} is the input stream that delivers the characters to the process. For
portability it should either be \texttt{InputTextNone} (\ref{InputTextNone}) (if the process reads no characters), \texttt{InputTextUser} (\ref{InputTextUser}), the result of a call to \texttt{InputTextFile} (\ref{InputTextFile}) from which no characters have been read, or the result of a call to \texttt{InputTextString} (\ref{InputTextString}). 

 \texttt{Process} is free to consume \emph{all} the input even if the program itself does not require any input at all. 

 \mbox{\texttt{\mdseries\slshape stream\texttt{\symbol{45}}out}} is the output stream which receives the characters from the process. For
portability it should either be \texttt{OutputTextNone} (\ref{OutputTextNone}) (if the process writes no characters), \texttt{OutputTextUser} (\ref{OutputTextUser}), the result of a call to \texttt{OutputTextFile} (\ref{OutputTextFile}) to which no characters have been written, or the result of a call to \texttt{OutputTextString} (\ref{OutputTextString}). 

 \mbox{\texttt{\mdseries\slshape options}} is a list of strings which are passed to the process as command line argument.
Note that no substitutions are performed on the strings, i.e., they are passed
immediately to the process and are not processed by a command interpreter
(shell). Further note that each string is passed as one argument, even if it
contains \emph{space} characters. Note that input/output redirection commands are \emph{not} allowed as \mbox{\texttt{\mdseries\slshape options}}. 

 In order to find a system program use \texttt{PathSystemProgram} (\ref{PathSystemProgram}). 

 
\begin{Verbatim}[commandchars=!@|,fontsize=\small,frame=single,label=Example]
  !gapprompt@gap>| !gapinput@date := PathSystemProgram( "date" );|
  "/bin/date"
\end{Verbatim}
 

 The next example shows how to execute \texttt{date} with no argument and no input, and collect the output into a string stream. 

 
\begin{Verbatim}[commandchars=!@|,fontsize=\small,frame=single,label=Example]
  !gapprompt@gap>| !gapinput@str := "";; a := OutputTextString(str,true);;|
  !gapprompt@gap>| !gapinput@Process( DirectoryCurrent(), date, InputTextNone(), a, [] );|
  0
  !gapprompt@gap>| !gapinput@CloseStream(a);|
  !gapprompt@gap>| !gapinput@Print(str);|
  Fri Jul 11 09:04:23 MET DST 1997
\end{Verbatim}
 }

 

\subsection{\textcolor{Chapter }{Exec}}
\logpage{[ 11, 1, 2 ]}\nobreak
\hyperdef{L}{X81402C91833986FC}{}
{\noindent\textcolor{FuncColor}{$\triangleright$\enspace\texttt{Exec({\mdseries\slshape cmd, option1, ..., optionN})\index{Exec@\texttt{Exec}}
\label{Exec}
}\hfill{\scriptsize (function)}}\\


 \texttt{Exec} runs a shell in the current directory to execute the command given by the
string \mbox{\texttt{\mdseries\slshape cmd}} with options \mbox{\texttt{\mdseries\slshape option1}}, ..., \mbox{\texttt{\mdseries\slshape optionN}}. 

 
\begin{Verbatim}[commandchars=!@|,fontsize=\small,frame=single,label=Example]
  !gapprompt@gap>| !gapinput@Exec( "date" );|
  Thu Jul 24 10:04:13 BST 1997
\end{Verbatim}
 

 \mbox{\texttt{\mdseries\slshape cmd}} is interpreted by the shell and therefore we can make use of the various
features that a shell offers as in following example. 

 
\begin{Verbatim}[commandchars=@|B,fontsize=\small,frame=single,label=Example]
  @gapprompt|gap>B @gapinput|Exec( "echo \"GAP is great!\" > foo" );B
  @gapprompt|gap>B @gapinput|Exec( "cat foo" );B
  GAP is great!
  @gapprompt|gap>B @gapinput|Exec( "rm foo" );B
\end{Verbatim}
 

 \texttt{Exec} calls the more general operation \texttt{Process} (\ref{Process}). The function \texttt{Edit} (\ref{Edit}) should be used to call an editor from within \textsf{GAP}. }

 }

 }

              
\chapter{\textcolor{Chapter }{Objects and Elements}}\label{Objects and Elements}
\logpage{[ 12, 0, 0 ]}
\hyperdef{L}{X86710F997832ABA4}{}
{
  An \emph{object} is anything in \textsf{GAP} that can be assigned to a variable, so nearly everything in \textsf{GAP} is an object. 

 Different objects can be regarded as equal with respect to the equivalence
relation ``\texttt{=}'', in this case we say that the objects describe the same \emph{element}.  
\section{\textcolor{Chapter }{Objects}}\label{Objects}
\logpage{[ 12, 1, 0 ]}
\hyperdef{L}{X78497E777FB3E402}{}
{
  Nearly all things one deals with in \textsf{GAP} are \emph{objects}. For example, an integer is an object, as is a list of integers, a matrix, a
permutation, a function, a list of functions, a record, a group, a coset or a
conjugacy class in a group. 

 Examples of things that are not objects are comments which are only lexical
constructs, \texttt{while} loops which are only syntactical constructs, and expressions, such as \texttt{1 + 1}; but note that the value of an expression, in this case the integer \texttt{2}, is an object. 

 Objects can be assigned to variables, and everything that can be assigned to a
variable is an object. Analogously, objects can be used as arguments of
functions, and can be returned by functions. 

\subsection{\textcolor{Chapter }{IsObject}}
\logpage{[ 12, 1, 1 ]}\nobreak
\hyperdef{L}{X7B130AC98415CAFB}{}
{\noindent\textcolor{FuncColor}{$\triangleright$\enspace\texttt{IsObject({\mdseries\slshape obj})\index{IsObject@\texttt{IsObject}}
\label{IsObject}
}\hfill{\scriptsize (Category)}}\\


 \texttt{IsObject} returns \texttt{true} if the object \mbox{\texttt{\mdseries\slshape obj}} is an object. Obviously it can never return \texttt{false}. 

 It can be used as a filter in \texttt{InstallMethod} (\ref{InstallMethod}) when one of the arguments can be anything. }

 }

  
\section{\textcolor{Chapter }{Elements as equivalence classes}}\label{Elements as equivalence classes}
\logpage{[ 12, 2, 0 ]}
\hyperdef{L}{X780C66027A49D110}{}
{
  \index{elements!definition} The equality operation ``\texttt{=}'' defines an equivalence relation on all \textsf{GAP} objects. The equivalence classes are called \emph{elements}. 

 There are basically three reasons to regard different objects as equal.
Firstly the same information may be stored in different places. Secondly the
same information may be stored in different ways; for example, a polynomial
can be stored sparsely or densely. Thirdly different information may be equal
modulo a mathematical equivalence relation. For example, in a finitely
presented group with the relation $a^2 = 1$ the different objects $a$ and $a^3$ describe the same element. 

 As an example of all three reasons, consider the possibility of storing an
integer in several places of the memory, of representing it as a fraction with
denominator 1, or of representing it as a fraction with any denominator, and
numerator a suitable multiple of the denominator. }

  
\section{\textcolor{Chapter }{Sets}}\label{Sets}
\logpage{[ 12, 3, 0 ]}
\hyperdef{L}{X83BE0C20875DD285}{}
{
  In \textsf{GAP} there is no category whose definition corresponds to the mathematical property
of being a set, however in the manual we will often refer to an object as a \emph{set} in order to convey the fact that mathematically, we are thinking of it as a
set. In particular, two sets $A$ and $B$ are equal if and only if, $x \in A \iff x \in B$. 

 There are two types of object in \textsf{GAP} which exhibit this kind of behaviour with respect to equality, namely domains
(see Section{\nobreakspace}\ref{Domains}) and lists whose elements are strictly sorted see \texttt{IsSSortedList} (\ref{IsSSortedList}). In general, \emph{set} in this manual will mean an object of one of these types. 

 More precisely: two domains can be compared with ``\texttt{\symbol{123}}\texttt{=}\texttt{\symbol{125}}'', the answer being \texttt{true} if and only if the sets of elements are equal (regardless of any additional
structure) and; a domain and a list can be compared with ``\texttt{=}'', the answer being \texttt{true} if and only if the list is equal to the strictly sorted list of elements of
the domain. 

 A discussion about sorted lists and sets can be found in Section \ref{Sorted Lists and Sets}. }

  
\section{\textcolor{Chapter }{Domains}}\label{Domains}
\logpage{[ 12, 4, 0 ]}
\hyperdef{L}{X7BAF69417BB925F6}{}
{
  An especially important class of objects in \textsf{GAP} are those whose underlying mathematical abstraction is that of a structured
set, for example a group, a conjugacy class, or a vector space. Such objects
are called \emph{domains}. The equality relation between domains is always equality \emph{as sets}, so that two domains are equal if and only if they contain the same elements. 

 Domains play a central role in \textsf{GAP}. In a sense, the only reason that \textsf{GAP} supports objects such as integers and permutations is the wish to form domains
of them and compute the properties of those domains. 

 Domains are described in Chapter{\nobreakspace}\ref{Domains and their Elements}. }

  
\section{\textcolor{Chapter }{Identical Objects}}\label{Identical Objects}
\logpage{[ 12, 5, 0 ]}
\hyperdef{L}{X84545F3985C60F5B}{}
{
  Two objects that are equal \emph{as objects} (that is they actually refer to the same area of computer memory) and not only
w.r.t. the equality relation ``\texttt{=}'' are called \emph{identical}. Identical objects do of course describe the same element. 

\subsection{\textcolor{Chapter }{IsIdenticalObj}}
\logpage{[ 12, 5, 1 ]}\nobreak
\hyperdef{L}{X7961183378DFB902}{}
{\noindent\textcolor{FuncColor}{$\triangleright$\enspace\texttt{IsIdenticalObj({\mdseries\slshape obj1, obj2})\index{IsIdenticalObj@\texttt{IsIdenticalObj}}
\label{IsIdenticalObj}
}\hfill{\scriptsize (function)}}\\


 \texttt{IsIdenticalObj} tests whether the objects \mbox{\texttt{\mdseries\slshape obj1}} and \mbox{\texttt{\mdseries\slshape obj2}} are identical (that is they are either equal immediate objects or are both
stored at the same location in memory. 

 If two copies of a simple constant object (see section \ref{Mutability and Copyability}) are created, it is not defined whether \textsf{GAP} will actually store two equal but non\texttt{\symbol{45}}identical objects, or
just a single object. For mutable objects, however, it is important to know
whether two values refer to identical or non\texttt{\symbol{45}}identical
objects, and the documentation of operations that return mutable values should
make clear whether the values returned are new, or may be identical to values
stored elsewhere. 

 
\begin{Verbatim}[commandchars=!@|,fontsize=\small,frame=single,label=Example]
  !gapprompt@gap>| !gapinput@IsIdenticalObj( 10^6, 10^6);|
  true
  !gapprompt@gap>| !gapinput@IsIdenticalObj( 10^30, 10^30);|
  false
  !gapprompt@gap>| !gapinput@IsIdenticalObj( true, true);|
  true
\end{Verbatim}
 

 Generally, one may compute with objects but think of the results in terms of
the underlying elements because one is not interested in locations in memory,
data formats or information beyond underlying equivalence relations. But there
are cases where it is important to distinguish the relations identity and
equality. This is best illustrated with an example. (The reader who is not
familiar with lists in \textsf{GAP}, in particular element access and assignment, is referred to
Chapter{\nobreakspace}\ref{Lists}.) 
\begin{Verbatim}[commandchars=!@|,fontsize=\small,frame=single,label=Example]
  !gapprompt@gap>| !gapinput@l1:= [ 1, 2, 3 ];; l2:= [ 1, 2, 3 ];;|
  !gapprompt@gap>| !gapinput@l1 = l2;|
  true
  !gapprompt@gap>| !gapinput@IsIdenticalObj( l1, l2 );|
  false
  !gapprompt@gap>| !gapinput@l1[3]:= 4;; l1; l2;|
  [ 1, 2, 4 ]
  [ 1, 2, 3 ]
  !gapprompt@gap>| !gapinput@l1 = l2;|
  false
\end{Verbatim}
 The two lists \texttt{l1} and \texttt{l2} are equal but not identical. Thus a change in \texttt{l1} does not affect \texttt{l2}. 
\begin{Verbatim}[commandchars=!@|,fontsize=\small,frame=single,label=Example]
  !gapprompt@gap>| !gapinput@l1:= [ 1, 2, 3 ];; l2:= l1;;|
  !gapprompt@gap>| !gapinput@l1 = l2;|
  true
  !gapprompt@gap>| !gapinput@IsIdenticalObj( l1, l2 );|
  true
  !gapprompt@gap>| !gapinput@l1[3]:= 4;; l1; l2;|
  [ 1, 2, 4 ]
  [ 1, 2, 4 ]
  !gapprompt@gap>| !gapinput@l1 = l2;|
  true
\end{Verbatim}
 Here, \texttt{l1} and \texttt{l2} are identical objects, so changing \texttt{l1} means a change to \texttt{l2} as well. }

 

\subsection{\textcolor{Chapter }{IsNotIdenticalObj}}
\logpage{[ 12, 5, 2 ]}\nobreak
\hyperdef{L}{X811976EC78EC5E29}{}
{\noindent\textcolor{FuncColor}{$\triangleright$\enspace\texttt{IsNotIdenticalObj({\mdseries\slshape obj1, obj2})\index{IsNotIdenticalObj@\texttt{IsNotIdenticalObj}}
\label{IsNotIdenticalObj}
}\hfill{\scriptsize (function)}}\\


 tests whether the objects \mbox{\texttt{\mdseries\slshape obj1}} and \mbox{\texttt{\mdseries\slshape obj2}} are not identical. }

 }

  
\section{\textcolor{Chapter }{Mutability and Copyability}}\label{Mutability and Copyability}
\logpage{[ 12, 6, 0 ]}
\hyperdef{L}{X7F0C119682196D65}{}
{
  An object in \textsf{GAP} is said to be \emph{immutable} if its mathematical value (as defined by $=$) does not change under any operation. More explicitly, suppose $a$ is immutable and $O$ is some operation on $a$, then if $a = b$ evaluates to \texttt{true} before executing $O(a)$, $a = b$ also evaluates to \texttt{true} afterwards. (Examples for operations $O$ that change mutable objects are \texttt{Add} (\ref{Add}) and \texttt{Unbind} (\ref{Unbind:unbind a list entry}) which are used to change list objects, see Chapter{\nobreakspace}\ref{Lists}.) An immutable object \emph{may} change, for example to store new information, or to adopt a more efficient
representation, but this does not affect its behaviour under $=$. 

 There are two points here to note. Firstly, ``operation'' above refers to the functions and methods which can legitimately be applied to
the object, and not the \texttt{!.} operation whereby virtually any aspect of any \textsf{GAP} level object may be changed. The second point which follows from this, is that
when implementing new types of objects, it is the programmer's responsibility
to ensure that the functions and methods they write never change immutable
objects mathematically. 

 In fact, most objects with which one deals in \textsf{GAP} are immutable. For instance, the permutation \texttt{(1,2)} will never become a different permutation or a
non\texttt{\symbol{45}}permutation (although a variable which previously had \texttt{(1,2)} stored in it may subsequently have some other value). 

 For many purposes, however, \emph{mutable} objects are useful. These objects may be changed to represent different
mathematical objects during their life. For example, mutable lists can be
changed by assigning values to positions or by unbinding values at certain
positions. Similarly, one can assign values to components of a mutable record,
or unbind them. 

\subsection{\textcolor{Chapter }{IsCopyable}}
\logpage{[ 12, 6, 1 ]}\nobreak
\hyperdef{L}{X811EFD727EBD1ADC}{}
{\noindent\textcolor{FuncColor}{$\triangleright$\enspace\texttt{IsCopyable({\mdseries\slshape obj})\index{IsCopyable@\texttt{IsCopyable}}
\label{IsCopyable}
}\hfill{\scriptsize (Category)}}\\


 If a mutable form of an object \mbox{\texttt{\mdseries\slshape obj}} can be made in \textsf{GAP}, the object is called \emph{copyable}. Examples of copyable objects are of course lists and records. A new mutable
version of the object can always be obtained by the operation \texttt{ShallowCopy} (\ref{ShallowCopy}). 

 Objects for which only an immutable form exists in \textsf{GAP} are called \emph{constants}. Examples of constants are integers, permutations, and domains. Called with a
constant as argument, \texttt{Immutable} (\ref{Immutable}) and \texttt{ShallowCopy} (\ref{ShallowCopy}) return this argument. }

 

\subsection{\textcolor{Chapter }{IsMutable}}
\logpage{[ 12, 6, 2 ]}\nobreak
\hyperdef{L}{X7999AD1D7A4F1F46}{}
{\noindent\textcolor{FuncColor}{$\triangleright$\enspace\texttt{IsMutable({\mdseries\slshape obj})\index{IsMutable@\texttt{IsMutable}}
\label{IsMutable}
}\hfill{\scriptsize (Category)}}\\


 tests whether \mbox{\texttt{\mdseries\slshape obj}} is mutable. 

 If an object is mutable then it is also copyable (see{\nobreakspace}\texttt{IsCopyable} (\ref{IsCopyable})), and a \texttt{ShallowCopy} (\ref{ShallowCopy}) method should be supplied for it. Note that \texttt{IsMutable} must not be implied by another filter, since otherwise \texttt{Immutable} (\ref{Immutable}) would be able to create paradoxical objects in the sense that \texttt{IsMutable} for such an object is \texttt{false} but the filter that implies \texttt{IsMutable} is \texttt{true}. 

 In many situations, however, one wants to ensure that objects are \emph{immutable}. For example, take the identity of a matrix group. Since this matrix may be
referred to as the identity of the group in several places, it would be fatal
to modify its entries, or add or unbind rows. We can obtain an immutable copy
of an object with \texttt{Immutable} (\ref{Immutable}). }

 

\subsection{\textcolor{Chapter }{Immutable}}
\logpage{[ 12, 6, 3 ]}\nobreak
\hyperdef{L}{X7F0ABF2C870B0CBB}{}
{\noindent\textcolor{FuncColor}{$\triangleright$\enspace\texttt{Immutable({\mdseries\slshape obj})\index{Immutable@\texttt{Immutable}}
\label{Immutable}
}\hfill{\scriptsize (function)}}\\


 returns an immutable structural copy (see{\nobreakspace}\texttt{StructuralCopy} (\ref{StructuralCopy})) of \mbox{\texttt{\mdseries\slshape obj}} in which the subobjects are immutable \emph{copies} of the subobjects of \mbox{\texttt{\mdseries\slshape obj}}. If \mbox{\texttt{\mdseries\slshape obj}} is immutable then \texttt{Immutable} returns \mbox{\texttt{\mdseries\slshape obj}} itself. 

 \textsf{GAP} will complain with an error if one tries to change an immutable object. }

 

\subsection{\textcolor{Chapter }{MakeImmutable}}
\logpage{[ 12, 6, 4 ]}\nobreak
\hyperdef{L}{X80CE136D804097C7}{}
{\noindent\textcolor{FuncColor}{$\triangleright$\enspace\texttt{MakeImmutable({\mdseries\slshape obj})\index{MakeImmutable@\texttt{MakeImmutable}}
\label{MakeImmutable}
}\hfill{\scriptsize (function)}}\\


 One can turn the (mutable or immutable) object \mbox{\texttt{\mdseries\slshape obj}} into an immutable one with \texttt{MakeImmutable}; note that this also makes all subobjects of \mbox{\texttt{\mdseries\slshape obj}} immutable, so one should call \texttt{MakeImmutable} only if \mbox{\texttt{\mdseries\slshape obj}} and its mutable subobjects are newly created. If one is not sure about this, \texttt{Immutable} (\ref{Immutable}) should be used. 

 Note that it is \emph{not} possible to turn an immutable object into a mutable one; only mutable copies
can be made (see{\nobreakspace}\ref{Duplication of Objects}). 

 Using \texttt{Immutable} (\ref{Immutable}), it is possible to store an immutable identity matrix or an immutable list of
generators, and to pass around references to this immutable object safely.
Only when a mutable copy is really needed does the actual object have to be
duplicated. Compared to the situation without immutable objects, much
unnecessary copying is avoided this way. Another advantage of immutability is
that lists of immutable objects may remember whether they are sorted
(see{\nobreakspace}\ref{Sorted Lists and Sets}), which is not possible for lists of mutable objects. 

 Since the operation \texttt{Immutable} (\ref{Immutable}) must work for any object in \textsf{GAP}, it follows that an immutable form of every object must be possible, even if
it is not sensible, and user\texttt{\symbol{45}}defined objects must allow for
the possibility of becoming immutable without notice. }

 
\subsection{\textcolor{Chapter }{Mutability of Iterators}}\label{subsect:mutability iterators}
\logpage{[ 12, 6, 5 ]}
\hyperdef{L}{X7FBA5F4D7C6872BD}{}
{
  An interesting example of mutable (and thus copyable) objects is provided by \emph{iterators}, see{\nobreakspace}\ref{Iterators}. (Of course an immutable form of an iterator is not very useful, but clearly \texttt{Immutable} (\ref{Immutable}) will yield such an object.) Every call of \texttt{NextIterator} (\ref{NextIterator}) changes a mutable iterator until it is exhausted, and this is the only way to
change an iterator. \texttt{ShallowCopy} (\ref{ShallowCopy}) for an iterator \mbox{\texttt{\mdseries\slshape iter}} is defined so as to return a mutable iterator that has no mutable data in
common with \mbox{\texttt{\mdseries\slshape iter}}, and that behaves equally to \mbox{\texttt{\mdseries\slshape iter}} w.r.t.{\nobreakspace}\texttt{IsDoneIterator} (\ref{IsDoneIterator}) and (if \mbox{\texttt{\mdseries\slshape iter}} is mutable) \texttt{NextIterator} (\ref{NextIterator}). Note that this meaning of the ``shallow copy'' of an iterator that is returned by \texttt{ShallowCopy} (\ref{ShallowCopy}) is not as obvious as for lists and records, and must be explicitly defined. }

 
\subsection{\textcolor{Chapter }{Mutability of Results of Arithmetic Operations}}\label{subsect:mutability arithmetic operations}
\logpage{[ 12, 6, 6 ]}
\hyperdef{L}{X7ADB82997A16E853}{}
{
  Many operations return immutable results, among those in particular attributes
(see{\nobreakspace}\ref{Attributes}). Examples of attributes are \texttt{Size} (\ref{Size}), \texttt{Zero} (\ref{Zero}), \texttt{AdditiveInverse} (\ref{AdditiveInverse}), \texttt{One} (\ref{One}), and \texttt{Inverse} (\ref{Inverse}). Arithmetic operations, such as the binary infix operations \texttt{+}, \texttt{\texttt{\symbol{45}}}, \texttt{*}, \texttt{/}, \texttt{\texttt{\symbol{94}}}, \texttt{mod}, the unary \texttt{\texttt{\symbol{45}}}, and operations such as \texttt{Comm} (\ref{Comm}) and \texttt{LeftQuotient} (\ref{LeftQuotient}), return \emph{mutable} results, \emph{except} if all arguments are immutable. So the product of two matrices or of a vector
and a matrix is immutable if and only if the two matrices or both the vector
and the matrix are immutable (see also{\nobreakspace}\ref{Arithmetic for Lists}). There is one exception to this rule, which arises where the result is less
deeply nested than at least one of the argument, where mutable arguments may
sometimes lead to an immutable result. For instance, a mutable matrix with
immutable rows, multiplied by an immutable vector gives an immutable vector
result. The exact rules are given in{\nobreakspace}\ref{Arithmetic for Lists}. 

 It should be noted that \texttt{0 * \mbox{\texttt{\mdseries\slshape obj}}} is equivalent to \texttt{ZeroSameMutability( \mbox{\texttt{\mdseries\slshape obj}} )}, \texttt{\texttt{\symbol{45}}\mbox{\texttt{\mdseries\slshape obj}}} is equivalent to \texttt{AdditiveInverseSameMutability( \mbox{\texttt{\mdseries\slshape obj}} )}, \texttt{\mbox{\texttt{\mdseries\slshape obj}}\texttt{\symbol{94}}0} is equivalent to \texttt{OneSameMutability( \mbox{\texttt{\mdseries\slshape obj}})}, and \texttt{\mbox{\texttt{\mdseries\slshape obj}}\texttt{\symbol{94}}\texttt{\symbol{45}}1} is equivalent to \texttt{InverseSameMutability( \mbox{\texttt{\mdseries\slshape obj}} )}. The suffix ``SameMutability'' indicates that the result is mutable if and only if the argument is mutable. 

 The operations \texttt{ZeroOp} (\ref{ZeroOp}), \texttt{AdditiveInverseOp} (\ref{AdditiveInverseOp}), \texttt{OneOp} (\ref{OneOp}), and \texttt{InverseOp} (\ref{InverseOp}) return \emph{mutable} results whenever a mutable version of the result exists, contrary to the
attributes \texttt{Zero} (\ref{Zero}), \texttt{AdditiveInverse} (\ref{AdditiveInverse}), \texttt{One} (\ref{One}), and \texttt{Inverse} (\ref{Inverse}). 

 If one introduces new arithmetic objects then one need not install methods for
the attributes \texttt{One} (\ref{One}), \texttt{Zero} (\ref{Zero}), etc. The methods for the associated operations \texttt{OneOp} (\ref{OneOp}) and \texttt{ZeroOp} (\ref{ZeroOp}) will be called, and then the results made immutable. 

 All methods installed for the arithmetic operations must obey the rule about
the mutability of the result. This means that one may try to avoid the perhaps
expensive creation of a new object if both operands are immutable, and of
course no problems of this kind arise at all in the (usual) case that the
objects in question do not admit a mutable form, i.e., that these objects are
not copyable. 

 In a few, relatively low\texttt{\symbol{45}}level algorithms, one wishes to
treat a matrix partly as a data structure, and manipulate and change its
entries. For this, the matrix needs to be mutable, and the rule that attribute
values are immutable is an obstacle. For these situations, a number of
additional operations are provided, for example \texttt{TransposedMatMutable} (\ref{TransposedMatMutable}) constructs a mutable matrix (contrary to the attribute \texttt{TransposedMat} (\ref{TransposedMat})), while \texttt{TriangulizeMat} (\ref{TriangulizeMat}) modifies a mutable matrix (in place) into upper triangular form. 

 Note that being immutable does not forbid an object to store knowledge. For
example, if it is found out that an immutable list is strictly sorted then the
list may store this information. More precisely, an immutable object may
change in any way, provided that it continues to represent the same
mathematical object. }

 }

  
\section{\textcolor{Chapter }{Duplication of Objects}}\label{Duplication of Objects}
\logpage{[ 12, 7, 0 ]}
\hyperdef{L}{X786B942B82D684BD}{}
{
  \index{Copy} \index{copy!an object} \index{clone!an object} 

\subsection{\textcolor{Chapter }{ShallowCopy}}
\logpage{[ 12, 7, 1 ]}\nobreak
\hyperdef{L}{X846BC7107C352031}{}
{\noindent\textcolor{FuncColor}{$\triangleright$\enspace\texttt{ShallowCopy({\mdseries\slshape obj})\index{ShallowCopy@\texttt{ShallowCopy}}
\label{ShallowCopy}
}\hfill{\scriptsize (operation)}}\\


 \texttt{ShallowCopy} returns a \emph{new mutable} object \emph{equal} to its argument, if this is possible. The subobjects of \texttt{ShallowCopy( \mbox{\texttt{\mdseries\slshape obj}} )} are \emph{identical} to the subobjects of \mbox{\texttt{\mdseries\slshape obj}}. 

 If \textsf{GAP} does not support a mutable form of the immutable object \mbox{\texttt{\mdseries\slshape obj}} (see{\nobreakspace}\ref{Mutability and Copyability}) then \texttt{ShallowCopy} returns \mbox{\texttt{\mdseries\slshape obj}} itself. 

 Since \texttt{ShallowCopy} is an operation, the concrete meaning of ``subobject'' depends on the type of \mbox{\texttt{\mdseries\slshape obj}}. But for any copyable object \mbox{\texttt{\mdseries\slshape obj}}, the definition should reflect the idea of ``first level copying''. 

 The definition of \texttt{ShallowCopy} for lists (in particular for matrices) can be found in{\nobreakspace}\ref{Duplication of Lists}. }

 

\subsection{\textcolor{Chapter }{StructuralCopy}}
\logpage{[ 12, 7, 2 ]}\nobreak
\hyperdef{L}{X7C1E70587EBDD2CB}{}
{\noindent\textcolor{FuncColor}{$\triangleright$\enspace\texttt{StructuralCopy({\mdseries\slshape obj})\index{StructuralCopy@\texttt{StructuralCopy}}
\label{StructuralCopy}
}\hfill{\scriptsize (function)}}\\


 In a few situations, one wants to make a \emph{structural copy} \texttt{scp} of an object \mbox{\texttt{\mdseries\slshape obj}}. This is defined as follows. \texttt{scp} and \mbox{\texttt{\mdseries\slshape obj}} are identical if \mbox{\texttt{\mdseries\slshape obj}} is immutable. Otherwise, \texttt{scp} is a mutable copy of \mbox{\texttt{\mdseries\slshape obj}} such that each subobject of \texttt{scp} is a structural copy of the corresponding subobject of \mbox{\texttt{\mdseries\slshape obj}}. Furthermore, if two subobjects of \mbox{\texttt{\mdseries\slshape obj}} are identical then also the corresponding subobjects of \texttt{scp} are identical. 
\begin{Verbatim}[commandchars=!@|,fontsize=\small,frame=single,label=Example]
  !gapprompt@gap>| !gapinput@obj:= [ [ 0, 1 ] ];;|
  !gapprompt@gap>| !gapinput@obj[2]:= obj[1];;|
  !gapprompt@gap>| !gapinput@obj[3]:= Immutable( obj[1] );;|
  !gapprompt@gap>| !gapinput@scp:= StructuralCopy( obj );;|
  !gapprompt@gap>| !gapinput@scp = obj; IsIdenticalObj( scp, obj );|
  true
  false
  !gapprompt@gap>| !gapinput@IsIdenticalObj( scp[1], obj[1] );|
  false
  !gapprompt@gap>| !gapinput@IsIdenticalObj( scp[3], obj[3] );|
  true
  !gapprompt@gap>| !gapinput@IsIdenticalObj( scp[1], scp[2] );|
  true
\end{Verbatim}
 

 That both \texttt{ShallowCopy} (\ref{ShallowCopy}) and \texttt{StructuralCopy} return the argument \mbox{\texttt{\mdseries\slshape obj}} itself if it is not copyable is consistent with this definition, since there
is no way to change \mbox{\texttt{\mdseries\slshape obj}} by modifying the result of any of the two functions, because in fact there is
no way to change this result at all. }

 }

  
\section{\textcolor{Chapter }{Other Operations Applicable to any Object}}\label{Other Operations Applicable to any Object}
\logpage{[ 12, 8, 0 ]}
\hyperdef{L}{X86E7193D848C53FC}{}
{
  There are a number of general operations which can be applied, in principle,
to any object in \textsf{GAP}. Some of these are documented elsewhere {\textendash}see \texttt{String} (\ref{String}), \texttt{PrintObj} (\ref{PrintObj}) and \texttt{Display} (\ref{Display}). Others are mainly somewhat technical. 

\subsection{\textcolor{Chapter }{SetName}}
\logpage{[ 12, 8, 1 ]}\nobreak
\hyperdef{L}{X85D6D47B83BD02A1}{}
{\noindent\textcolor{FuncColor}{$\triangleright$\enspace\texttt{SetName({\mdseries\slshape obj, name})\index{SetName@\texttt{SetName}}
\label{SetName}
}\hfill{\scriptsize (operation)}}\\


 for a suitable object \mbox{\texttt{\mdseries\slshape obj}} sets that object to have name \mbox{\texttt{\mdseries\slshape name}} (a string). }

 

\subsection{\textcolor{Chapter }{Name}}
\logpage{[ 12, 8, 2 ]}\nobreak
\hyperdef{L}{X7F14EF9D81432113}{}
{\noindent\textcolor{FuncColor}{$\triangleright$\enspace\texttt{Name({\mdseries\slshape obj})\index{Name@\texttt{Name}}
\label{Name}
}\hfill{\scriptsize (attribute)}}\\


 returns the name, a string, previously assigned to \mbox{\texttt{\mdseries\slshape obj}} via a call to \texttt{SetName} (\ref{SetName}). The name of an object is used \emph{only} for viewing the object via this name. 

 There are no methods installed for computing names of objects, but the name
may be set for suitable objects, using \texttt{SetName} (\ref{SetName}). 
\begin{Verbatim}[commandchars=!@|,fontsize=\small,frame=single,label=Example]
  !gapprompt@gap>| !gapinput@R := PolynomialRing(Integers,2);|
  Integers[x_1,x_2]
  !gapprompt@gap>| !gapinput@SetName(R,"Z[x,y]");|
  !gapprompt@gap>| !gapinput@R;|
  Z[x,y]
  !gapprompt@gap>| !gapinput@Name(R);|
  "Z[x,y]"
\end{Verbatim}
 }

 

\subsection{\textcolor{Chapter }{InfoText}}
\logpage{[ 12, 8, 3 ]}\nobreak
\hyperdef{L}{X871562FD7F982C12}{}
{\noindent\textcolor{FuncColor}{$\triangleright$\enspace\texttt{InfoText({\mdseries\slshape obj})\index{InfoText@\texttt{InfoText}}
\label{InfoText}
}\hfill{\scriptsize (attribute)}}\\


 is a mutable string with information about the object \mbox{\texttt{\mdseries\slshape obj}}. There is no default method to create an info text. }

 

\subsection{\textcolor{Chapter }{IsInternallyConsistent}}
\logpage{[ 12, 8, 4 ]}\nobreak
\hyperdef{L}{X7F6C5C3287E8B816}{}
{\noindent\textcolor{FuncColor}{$\triangleright$\enspace\texttt{IsInternallyConsistent({\mdseries\slshape obj})\index{IsInternallyConsistent@\texttt{IsInternallyConsistent}}
\label{IsInternallyConsistent}
}\hfill{\scriptsize (operation)}}\\


 For debugging purposes, it may be useful to check the consistency of an object \mbox{\texttt{\mdseries\slshape obj}} that is composed from other (composed) objects. 

 There is a default method of \texttt{IsInternallyConsistent}, with rank zero, that returns \texttt{true}. So it is possible (and recommended) to check the consistency of subobjects
of \mbox{\texttt{\mdseries\slshape obj}} recursively by \texttt{IsInternallyConsistent}. 

 (Note that \texttt{IsInternallyConsistent} is not an attribute.) }

 

\subsection{\textcolor{Chapter }{MemoryUsage}}
\logpage{[ 12, 8, 5 ]}\nobreak
\hyperdef{L}{X7F4D216B7DF7BE9D}{}
{\noindent\textcolor{FuncColor}{$\triangleright$\enspace\texttt{MemoryUsage({\mdseries\slshape obj})\index{MemoryUsage@\texttt{MemoryUsage}}
\label{MemoryUsage}
}\hfill{\scriptsize (function)}}\\


 Returns the amount of memory in bytes used by the object \mbox{\texttt{\mdseries\slshape obj}} and its subobjects. Note that in general, objects can reference each other in
very difficult ways such that determining the memory usage is a recursive
procedure. In particular, computing the memory usage of a complicated
structure itself uses some additional memory, which is however no longer used
after completion of this operation. This procedure descends into lists and
records, positional and component objects; however it deliberately does not
take into account the type and family objects. For functions, it only takes
the memory usage of the function body, not of the local context the function
was created in, although the function keeps a reference to that as well. }

 }

 }

           
\chapter{\textcolor{Chapter }{Types of Objects}}\label{Types of Objects}
\logpage{[ 13, 0, 0 ]}
\hyperdef{L}{X7E8202627B421DB1}{}
{
  Every \textsf{GAP} object has a \emph{type}. The type of an object is the information which is used to decide whether an
operation is admissible or possible with that object as an argument, and if
so, how it is to be performed (see Chapter{\nobreakspace}\ref{Method Selection}). 

 For example, the types determine whether two objects can be multiplied and
what function is called to compute the product. Analogously, the type of an
object determines whether and how the size of the object can be computed. It
is sometimes useful in discussing the type system, to identify types with the
set of objects that have this type. Partial types can then also be regarded as
sets, such that any type is the intersection of its parts. 

 The type of an object consists of two main parts, which describe different
aspects of the object. 

 The \emph{family} determines the relation of the object to other objects. For example, all
permutations form a family. Another family consists of all collections of
permutations, this family contains the set of permutation groups as a subset.
A third family consists of all rational functions with coefficients in a
certain family. 

 The other part of a type is a collection of \emph{filters} (actually stored as a bit\texttt{\symbol{45}}list indicating, from the
complete set of possible filters, which are included in this particular type).
These filters are all treated equally by the method selection, but, from the
viewpoint of their creation and use, they can be divided (with a small number
of unimportant exceptions) into categories, representations, attribute testers
and properties. Each of these is described in more detail below. 

 This chapter does not describe how types and their constituent parts can be
created. Information about this topic can be found in Chapter{\nobreakspace}\ref{Creating New Objects}. 

 \emph{Note:} Detailed understanding of the type system is not required to use \textsf{GAP}. It can be helpful, however, to understand how things work and why \textsf{GAP} behaves the way it does. 

   A discussion of the type system can be found in \cite{BreuerLinton98}.  
\section{\textcolor{Chapter }{Families}}\label{Families}
\logpage{[ 13, 1, 0 ]}
\hyperdef{L}{X846063757EC05986}{}
{
  The family of an object determines its relationship to other objects. 

 More precisely, the families form a partition of all \textsf{GAP} objects such that the following two conditions hold: objects that are equal
w.r.t. \texttt{=} lie in the same family; and the family of the result of an operation depends
only on the families of its operands. 

 The first condition means that a family can be regarded as a set of elements
instead of a set of objects. Note that this does not hold for categories and
representations (see below), two objects that are equal w.r.t. \texttt{=} need not lie in the same categories and representations. For example, a
sparsely represented matrix can be equal to a densely represented matrix.
Similarly, each domain is equal w.r.t. \texttt{=} to the sorted list of its elements, but a domain is not a list, and a list is
not a domain. 

 Families are probably the least obvious part of the \textsf{GAP} type system, so some remarks about the role of families are necessary. When
one uses \textsf{GAP} as it is, one will (better: should) not meet families at all. The two
situations where families come into play are the following. 

 First, since families are used to describe relations between arguments of
operations in the method selection mechanism (see Chapter{\nobreakspace}\ref{Method Selection}, and also Chapter{\nobreakspace}\ref{Types of Objects}), one has to prescribe such a relation in each method installation
(see{\nobreakspace}\ref{Method Installation}); usual relations are \texttt{ReturnTrue} (\ref{ReturnTrue}) (which means that any relation of the actual arguments is admissible), \texttt{IsIdenticalObj} (\ref{IsIdenticalObj}) (which means that there are two arguments that lie in the same family), and \texttt{IsCollsElms} (which means that there are two arguments, the first being a collection of
elements that lie in the same family as the second argument). 

 Second {\textendash}and this is the more complicated situation{\textendash}
whenever one creates a new kind of objects, one has to decide what its family
shall be. If the new object shall be equal to existing objects, for example if
it is just represented in a different way, there is no choice: The new object
must lie in the same family as all objects that shall be equal to it. So only
if the new object is different (w.r.t.{\nobreakspace}the equality ``\texttt{=}'') from all other \textsf{GAP} objects, we are likely to create a new family for it. Note that enlarging an
existing family by such new objects may be problematic because of implications
that have been installed for all objects of the family in question. The choice
of families depends on the applications one has in mind. For example, if the
new objects in question are not likely to be arguments of operations for which
family relations are relevant (for example binary arithmetic operations), one
could create one family for all such objects, and regard it as ``the family of all those \textsf{GAP} objects that would in fact not need a family''. On the other extreme, if one wants to create domains of the new objects then
one has to choose the family in such a way that all intended elements of a
domain do in fact lie in the same family. (Remember that a domain is a
collection, see Chapter{\nobreakspace}\ref{Domains}, and that a collection consists of elements in the same family, see
Chapter{\nobreakspace}\ref{Collections} and Section{\nobreakspace}\ref{Families}.) 

 Let us look at an example. Suppose that no permutations are available in \textsf{GAP}, and that we want to implement permutations. Clearly we want to support
permutation groups, but it is not a priori clear how to distribute the new
permutations into families. We can put all permutations into one family; this
is how in fact permutations are implemented in \textsf{GAP}. But it would also be possible to put all permutations of a given degree into
a family of their own; this would for example mean that for each degree, there
would be distinguished trivial permutations, and that the stabilizer of the
point \texttt{5} in the symmetric group on the points \texttt{1}, \texttt{2}, $\ldots$, \texttt{5} is not regarded as equal to the symmetric group on \texttt{1}, \texttt{2}, \texttt{3}, \texttt{4}. Note that the latter approach would have the advantage that it is no problem
to construct permutations and permutation groups acting on arbitrary (finite)
sets, for example by constructing first the symmetric group on the set and
then generating any desired permutation group as a subgroup of this symmetric
group. 

 So one aspect concerning a reasonable choice of families is to make the
families large enough for being able to form interesting domains of elements
in the family. But on the other hand, it is useful to choose the families
small enough for admitting meaningful relations between objects. For example,
the elements of different free groups in \textsf{GAP} lie in different families; the multiplication of free group elements is
installed only for the case that the two operands lie in the same family, with
the effect that one cannot erroneously form the product of elements from
different free groups. In this case, families appear as a tool for providing
useful restrictions. 

 As another example, note that an element and a collection containing this
element never lie in the same family, by the general implementation of
collections; namely, the family of a collection of elements in the family \mbox{\texttt{\mdseries\slshape Fam}} is the collections family of \mbox{\texttt{\mdseries\slshape Fam}} (see{\nobreakspace}\texttt{CollectionsFamily} (\ref{CollectionsFamily})). This means that for a collection, we need not (because we cannot) decide
about its family. 

 A few functions in \textsf{GAP} return families, see \texttt{CollectionsFamily} (\ref{CollectionsFamily}) and \texttt{ElementsFamily} (\ref{ElementsFamily}). 

\subsection{\textcolor{Chapter }{FamilyObj}}
\logpage{[ 13, 1, 1 ]}\nobreak
\hyperdef{L}{X7CF70EAC84284919}{}
{\noindent\textcolor{FuncColor}{$\triangleright$\enspace\texttt{FamilyObj({\mdseries\slshape obj})\index{FamilyObj@\texttt{FamilyObj}}
\label{FamilyObj}
}\hfill{\scriptsize (function)}}\\


 returns the family of the object \mbox{\texttt{\mdseries\slshape obj}}. 

 The family of the object \mbox{\texttt{\mdseries\slshape obj}} is itself an object, its family is \texttt{FamilyOfFamilies}. 

 It should be emphasized that families may be created when they are needed. For
example, the family of elements of a finitely presented group is created only
after the presentation has been constructed. Thus families are the dynamic
part of the type system, that is, the part that is not fixed after the
initialisation of \textsf{GAP}. 

 Families can be parametrized. For example, the elements of each finitely
presented group form a family of their own. Here the family of elements and
the finitely presented group coincide when viewed as sets. Note that elements
in different finitely presented groups lie in different families. This
distinction allows \textsf{GAP} to forbid multiplications of elements in different finitely presented groups. 

 As a special case, families can be parametrized by other families. An
important example is the family of \emph{collections} that can be formed for each family. A collection consists of objects that lie
in the same family, it is either a nonempty dense list of objects from the
same family or a domain. 

 Note that every domain is a collection, that is, it is not possible to
construct domains whose elements lie in different families. For example, a
polynomial ring over the rationals cannot contain the integer \texttt{0} because the family that contains the integers does not contain polynomials. So
one has to distinguish the integer zero from each zero polynomial. 

 Let us look at this example from a different viewpoint. A polynomial ring and
its coefficients ring lie in different families, hence the coefficients ring
cannot be embedded ``naturally'' into the polynomial ring in the sense that it is a subset. But it is possible
to allow, e.g., the multiplication of an integer and a polynomial over the
integers. The relation between the arguments, namely that one is a coefficient
and the other a polynomial, can be detected from the relation of their
families. Moreover, this analysis is easier than in a situation where the
rationals would lie in one family together with all polynomials over the
rationals, because then the relation of families would not distinguish the
multiplication of two polynomials, the multiplication of two coefficients, and
the multiplication of a coefficient with a polynomial. So the wish to describe
relations between elements can be taken as a motivation for the introduction
of families. }

 

\subsection{\textcolor{Chapter }{NewFamily}}
\logpage{[ 13, 1, 2 ]}\nobreak
\hyperdef{L}{X7FB4123E7E22137D}{}
{\noindent\textcolor{FuncColor}{$\triangleright$\enspace\texttt{NewFamily({\mdseries\slshape name[, req[, imp[, famfilter]]]})\index{NewFamily@\texttt{NewFamily}}
\label{NewFamily}
}\hfill{\scriptsize (function)}}\\


 \texttt{NewFamily} returns a new family \texttt{fam} with name \mbox{\texttt{\mdseries\slshape name}}. The arguments \mbox{\texttt{\mdseries\slshape req}} and \mbox{\texttt{\mdseries\slshape imp}}, if present, are filters of which \texttt{fam} shall be a subset. Any type that is created for \texttt{fam} (see \texttt{NewType} (\ref{NewType})) will describe objects in these filters, and thus any object that is created
in \texttt{fam} (see \texttt{Objectify} (\ref{Objectify})) will lie automatically in \mbox{\texttt{\mdseries\slshape req}}\texttt{ and }\mbox{\texttt{\mdseries\slshape imp}}. (Thus it does not matter whether one uses \mbox{\texttt{\mdseries\slshape req}} or \mbox{\texttt{\mdseries\slshape imp}} or both in the call, the distinction has only historical reasons.) 

 The filter \mbox{\texttt{\mdseries\slshape famfilter}}, if given, specifies a filter that will hold for the object \texttt{fam} itself (not for objects in \texttt{fam}). 

 Families are always represented as component objects (see{\nobreakspace}\ref{Component Objects}). This means that individual components can be set in \texttt{fam} in order to store and access useful information about it. }

 }

  
\section{\textcolor{Chapter }{Filters}}\label{Filters}
\logpage{[ 13, 2, 0 ]}
\hyperdef{L}{X84EFA4C07D4277BB}{}
{
  A \emph{filter} is a special unary \textsf{GAP} function that returns either \texttt{true} or \texttt{false}, depending on whether or not the argument lies in the set defined by the
filter. Filters are used to express different aspects of information about a \textsf{GAP} object, which are described below (see{\nobreakspace}\ref{Categories}, \ref{Representation}, \ref{Attributes}, \ref{Setter and Tester for Attributes}, \ref{Properties}, \ref{Other Filters}). 

 Presently any filter in \textsf{GAP} is implemented as a function which corresponds to a set of positions in the
bitlist which forms part of the type of each \textsf{GAP} object, and returns \texttt{true} if and only if the bitlist of the type of the argument has the value \texttt{true} at all of these positions. 

 The intersection (or meet) of two filters \mbox{\texttt{\mdseries\slshape filt1}}, \mbox{\texttt{\mdseries\slshape filt2}} is again a filter, it can be formed as 

 \index{and@\texttt{and}!for filters} \mbox{\texttt{\mdseries\slshape filt1}} \texttt{and} \mbox{\texttt{\mdseries\slshape filt2}} 

 See \ref{Operations for Booleans} for more details. 

 For example, \texttt{IsList and IsEmpty} is a filter that returns \texttt{true} if its argument is an empty list, and \texttt{false} otherwise. The filter \texttt{IsGroup} (\ref{IsGroup}) is defined as the intersection of the category \texttt{IsMagmaWithInverses} (\ref{IsMagmaWithInverses}) and the property \texttt{IsAssociative} (\ref{IsAssociative}). 

 A filter that is not the meet of other filters is called a \emph{simple filter}. For example, each attribute tester (see{\nobreakspace}\ref{Setter and Tester for Attributes}) is a simple filter. Each simple filter corresponds to a position in the
bitlist currently used as part of the data structure representing a type. 

 Every filter has a \emph{rank}, which is used to define a ranking of the methods installed for an operation,
see Section{\nobreakspace}\ref{Method Installation}. The rank of a filter can be accessed with \texttt{RankFilter} (\ref{RankFilter}). 

\subsection{\textcolor{Chapter }{RankFilter}}
\logpage{[ 13, 2, 1 ]}\nobreak
\hyperdef{L}{X82E62B997C05E05E}{}
{\noindent\textcolor{FuncColor}{$\triangleright$\enspace\texttt{RankFilter({\mdseries\slshape filt})\index{RankFilter@\texttt{RankFilter}}
\label{RankFilter}
}\hfill{\scriptsize (function)}}\\


 For simple filters, an \emph{incremental rank} is defined when the filter is created, see the sections about the creation of
filters: \texttt{NewCategory} (\ref{NewCategory}), \texttt{NewRepresentation} (\ref{NewRepresentation}), \texttt{NewAttribute} (\ref{NewAttribute}), \texttt{NewProperty} (\ref{NewProperty}), \texttt{NewFilter} (\ref{NewFilter}). For an arbitrary filter, its rank is given by the sum of the incremental
ranks of the \emph{involved} simple filters; in addition to the implied filters, these are also the
required filters of attributes (again see the sections about the creation of
filters). In other words, for the purpose of computing the rank and \emph{only} for this purpose, attribute testers are treated as if they would imply the
requirements of their attributes.  }

 

\subsection{\textcolor{Chapter }{NamesFilter}}
\logpage{[ 13, 2, 2 ]}\nobreak
\hyperdef{L}{X7A78ECC67E2C9D78}{}
{\noindent\textcolor{FuncColor}{$\triangleright$\enspace\texttt{NamesFilter({\mdseries\slshape filt})\index{NamesFilter@\texttt{NamesFilter}}
\label{NamesFilter}
}\hfill{\scriptsize (function)}}\\


 \texttt{NamesFilter} returns a list of names of the \emph{implied} simple filters of the filter \mbox{\texttt{\mdseries\slshape filt}}, these are all those simple filters \texttt{imp} such that every object in \mbox{\texttt{\mdseries\slshape filt}} also lies in \texttt{imp}. For implications between filters, see \texttt{ShowImpliedFilters} (\ref{ShowImpliedFilters}) as well as sections{\nobreakspace}\ref{Logical Implications}, \texttt{NewCategory} (\ref{NewCategory}), \texttt{NewRepresentation} (\ref{NewRepresentation}), \texttt{NewAttribute} (\ref{NewAttribute}), \texttt{NewProperty} (\ref{NewProperty}). }

 

\subsection{\textcolor{Chapter }{FilterByName}}
\logpage{[ 13, 2, 3 ]}\nobreak
\hyperdef{L}{X7F6645D87DD26CF0}{}
{\noindent\textcolor{FuncColor}{$\triangleright$\enspace\texttt{FilterByName({\mdseries\slshape name})\index{FilterByName@\texttt{FilterByName}}
\label{FilterByName}
}\hfill{\scriptsize (function)}}\\


 finds the filter with name \mbox{\texttt{\mdseries\slshape name}} in the global FILTERS list. This is useful to find filters that were created
but not bound to a global variable. }

 

\subsection{\textcolor{Chapter }{ShowImpliedFilters}}
\logpage{[ 13, 2, 4 ]}\nobreak
\hyperdef{L}{X7F9568A67F3840DE}{}
{\noindent\textcolor{FuncColor}{$\triangleright$\enspace\texttt{ShowImpliedFilters({\mdseries\slshape filter})\index{ShowImpliedFilters@\texttt{ShowImpliedFilters}}
\label{ShowImpliedFilters}
}\hfill{\scriptsize (function)}}\\


 Displays information about the filters that may be implied by \mbox{\texttt{\mdseries\slshape filter}}. They are given by their names. \texttt{ShowImpliedFilters} first displays the names of all filters that are unconditionally implied by \mbox{\texttt{\mdseries\slshape filter}}. It then displays implications that require further filters to be present
(indicating by \texttt{+} the required further filters). 
\begin{Verbatim}[commandchars=!@|,fontsize=\small,frame=single,label=Example]
  !gapprompt@gap>| !gapinput@ShowImpliedFilters(IsNilpotentGroup);|
  Implies:
     IsListOrCollection
     IsCollection
     IsDuplicateFree
     IsExtLElement
     CategoryCollections(IsExtLElement)
     IsExtRElement
     CategoryCollections(IsExtRElement)
     CategoryCollections(IsMultiplicativeElement)
     CategoryCollections(IsMultiplicativeElementWithOne)
     CategoryCollections(IsMultiplicativeElementWithInverse)
     IsGeneralizedDomain
     IsMagma
     IsMagmaWithOne
     IsMagmaWithInversesIfNonzero
     IsMagmaWithInverses
     IsAssociative
     HasMultiplicativeNeutralElement
     IsGeneratorsOfSemigroup
     IsSimpleSemigroup
     IsRegularSemigroup
     IsInverseSemigroup
     IsCompletelyRegularSemigroup
     IsGroupAsSemigroup
     IsMonoidAsSemigroup
     IsOrthodoxSemigroup
     IsSupersolvableGroup
     IsSolvableGroup
     IsNilpotentByFinite
  
  
  May imply with:
  +IsFinitelyGeneratedGroup
     IsPolycyclicGroup
  
\end{Verbatim}
 }

 

\subsection{\textcolor{Chapter }{FiltersType}}
\logpage{[ 13, 2, 5 ]}\nobreak
\hyperdef{L}{X836FAA18861BE387}{}
{\noindent\textcolor{FuncColor}{$\triangleright$\enspace\texttt{FiltersType({\mdseries\slshape type})\index{FiltersType@\texttt{FiltersType}}
\label{FiltersType}
}\hfill{\scriptsize (operation)}}\\
\noindent\textcolor{FuncColor}{$\triangleright$\enspace\texttt{FiltersObj({\mdseries\slshape object})\index{FiltersObj@\texttt{FiltersObj}}
\label{FiltersObj}
}\hfill{\scriptsize (operation)}}\\


 returns a list of the filters in the type \mbox{\texttt{\mdseries\slshape type}}, or in the type of the object \mbox{\texttt{\mdseries\slshape object}} respectively. 
\begin{Verbatim}[commandchars=!@|,fontsize=\small,frame=single,label=Example]
  !gapprompt@gap>| !gapinput@FiltersObj(fail);|
  [ <Category "IsBool">, <Representation "IsInternalRep"> ]
  !gapprompt@gap>| !gapinput@FiltersType(TypeOfTypes);|
  [ <Representation "IsPositionalObjectRep">, <Category "IsType">, <Representation "IsTypeDefaultRep"> ]
\end{Verbatim}
 }

 }

  
\section{\textcolor{Chapter }{Categories}}\label{Categories}
\logpage{[ 13, 3, 0 ]}
\hyperdef{L}{X7CC6903E78F24167}{}
{
  The \emph{categories} of an object are filters (see{\nobreakspace}\ref{Filters}) that determine what operations an object admits. For example, all integers
form a category, all rationals form a category, and all rational functions
form a category. An object which claims to lie in a certain category is
accepting the requirement that it should have methods for certain operations
(and perhaps that their behaviour should satisfy certain axioms). For example,
an object lying in the category \texttt{IsList} (\ref{IsList}) must have methods for \texttt{Length} (\ref{Length}), \texttt{IsBound\texttt{\symbol{92}}[\texttt{\symbol{92}}]} (\ref{IsBoundbSlash[bSlash]}) and the list element access operation \texttt{\texttt{\symbol{92}}[\texttt{\symbol{92}}]} (\ref{bSlash[bSlash]}). 

 An object can lie in several categories. For example, a row vector lies in the
categories \texttt{IsList} (\ref{IsList}) and \texttt{IsVector} (\ref{IsVector}); each list lies in the category \texttt{IsCopyable} (\ref{IsCopyable}), and depending on whether or not it is mutable, it may lie in the category \texttt{IsMutable} (\ref{IsMutable}). Every domain lies in the category \texttt{IsDomain} (\ref{IsDomain}). 

 Of course some categories of a mutable object may change when the object is
changed. For example, after assigning values to positions of a mutable
non\texttt{\symbol{45}}dense list, this list may become part of the category \texttt{IsDenseList} (\ref{IsDenseList}). 

 However, if an object is immutable then the set of categories it lies in is
fixed. 

 All categories in the library are created during initialization, in particular
they are not created dynamically at runtime. 

 The following list gives an overview of important categories of arithmetic
objects. Indented categories are to be understood as subcategories of the non
indented category listed above it. 
\begin{Verbatim}[commandchars=!@|,fontsize=\small,frame=single,label=Example]
      IsObject
          IsExtLElement
          IsExtRElement
              IsMultiplicativeElement
                  IsMultiplicativeElementWithOne
                      IsMultiplicativeElementWithInverse
          IsExtAElement
              IsAdditiveElement
                  IsAdditiveElementWithZero
                      IsAdditiveElementWithInverse
\end{Verbatim}
 Every object lies in the category \texttt{IsObject} (\ref{IsObject}). 

 The categories \texttt{IsExtLElement} (\ref{IsExtLElement}) and \texttt{IsExtRElement} (\ref{IsExtRElement}) contain objects that can be multiplied with other objects via \texttt{*} from the left and from the right, respectively. These categories are required
for the operands of the operation \texttt{*}. 

 The category \texttt{IsMultiplicativeElement} (\ref{IsMultiplicativeElement}) contains objects that can be multiplied from the left and from the right with
objects from the same family. \texttt{IsMultiplicativeElementWithOne} (\ref{IsMultiplicativeElementWithOne}) contains objects \texttt{obj} for which a multiplicatively neutral element can be obtained by taking the $0$\texttt{\symbol{45}}th power \texttt{obj\texttt{\symbol{94}}0}. \texttt{IsMultiplicativeElementWithInverse} (\ref{IsMultiplicativeElementWithInverse}) contains objects \texttt{obj} for which a multiplicative inverse can be obtained by forming \texttt{obj\texttt{\symbol{94}}\texttt{\symbol{45}}1}. 

 Likewise, the categories \texttt{IsExtAElement} (\ref{IsExtAElement}), \texttt{IsAdditiveElement} (\ref{IsAdditiveElement}), \texttt{IsAdditiveElementWithZero} (\ref{IsAdditiveElementWithZero}) and \texttt{IsAdditiveElementWithInverse} (\ref{IsAdditiveElementWithInverse}) contain objects that can be added via \texttt{+} to other objects, objects that can be added to objects of the same family,
objects for which an additively neutral element can be obtained by
multiplication with zero, and objects for which an additive inverse can be
obtained by multiplication with \texttt{\texttt{\symbol{45}}1}. 

 So a vector lies in \texttt{IsExtLElement} (\ref{IsExtLElement}), \texttt{IsExtRElement} (\ref{IsExtRElement}) and \texttt{IsAdditiveElementWithInverse} (\ref{IsAdditiveElementWithInverse}). A ring element must additionally lie in \texttt{IsMultiplicativeElement} (\ref{IsMultiplicativeElement}). 

 As stated above it is not guaranteed by the categories of objects whether the
result of an operation with these objects as arguments is defined. For
example, the category \texttt{IsMatrix} (\ref{IsMatrix}) is a subcategory of \texttt{IsMultiplicativeElementWithInverse} (\ref{IsMultiplicativeElementWithInverse}). Clearly not every matrix has a multiplicative inverse. But the category \texttt{IsMatrix} (\ref{IsMatrix}) makes each matrix an admissible argument of the operation \texttt{Inverse} (\ref{Inverse}), which may sometimes return \texttt{fail}. Likewise, two matrices can be multiplied only if they are of appropriate
shapes. 

 Analogous to the categories of arithmetic elements, there are categories of
domains of these elements. 
\begin{Verbatim}[commandchars=!@|,fontsize=\small,frame=single,label=Example]
      IsObject
          IsDomain
              IsMagma
                  IsMagmaWithOne
                      IsMagmaWithInversesIfNonzero
                          IsMagmaWithInverses
              IsAdditiveMagma
                  IsAdditiveMagmaWithZero
                      IsAdditiveMagmaWithInverses
              IsExtLSet
              IsExtRSet
\end{Verbatim}
 Of course \texttt{IsDomain} (\ref{IsDomain}) is a subcategory of \texttt{IsObject} (\ref{IsObject}). A domain that is closed under multiplication \texttt{*} is called a magma and it lies in the category \texttt{IsMagma} (\ref{IsMagma}). If a magma is closed under taking the identity, it lies in \texttt{IsMagmaWithOne} (\ref{IsMagmaWithOne}), and if it is also closed under taking inverses, it lies in \texttt{IsMagmaWithInverses} (\ref{IsMagmaWithInverses}). The category \texttt{IsMagmaWithInversesIfNonzero} (\ref{IsMagmaWithInversesIfNonzero}) denotes closure under taking inverses only for nonzero elements, every
division ring lies in this category. 

 Note that every set of categories constitutes its own notion of generation,
for example a group may be generated as a magma with inverses by some
elements, but to generate it as a magma with one it may be necessary to take
the union of these generators and their inverses. 

\subsection{\textcolor{Chapter }{IsCategory}}
\logpage{[ 13, 3, 1 ]}\nobreak
\hyperdef{L}{X792A23BF82BDF66B}{}
{\noindent\textcolor{FuncColor}{$\triangleright$\enspace\texttt{IsCategory({\mdseries\slshape object})\index{IsCategory@\texttt{IsCategory}}
\label{IsCategory}
}\hfill{\scriptsize (function)}}\\


 returns \texttt{true} if \mbox{\texttt{\mdseries\slshape object}} is a category (see{\nobreakspace}\ref{Categories}), and \texttt{false} otherwise. 

 Note that \textsf{GAP} categories are \emph{not} categories in the usual mathematical sense. }

 

\subsection{\textcolor{Chapter }{CategoriesOfObject}}
\logpage{[ 13, 3, 2 ]}\nobreak
\hyperdef{L}{X85C6EB707A406A5A}{}
{\noindent\textcolor{FuncColor}{$\triangleright$\enspace\texttt{CategoriesOfObject({\mdseries\slshape object})\index{CategoriesOfObject@\texttt{CategoriesOfObject}}
\label{CategoriesOfObject}
}\hfill{\scriptsize (operation)}}\\


 returns a list of the names of the categories in which \mbox{\texttt{\mdseries\slshape object}} lies. 
\begin{Verbatim}[commandchars=!@|,fontsize=\small,frame=single,label=Example]
  !gapprompt@gap>| !gapinput@g:=Group((1,2),(1,2,3));;|
  !gapprompt@gap>| !gapinput@CategoriesOfObject(g);|
  [ "IsListOrCollection", "IsCollection", "IsExtLElement",
    "CategoryCollections(IsExtLElement)", "IsExtRElement",
    "CategoryCollections(IsExtRElement)",
    "CategoryCollections(IsMultiplicativeElement)",
    "CategoryCollections(IsMultiplicativeElementWithOne)",
    "CategoryCollections(IsMultiplicativeElementWithInverse)",
    "CategoryCollections(IsAssociativeElement)",
    "CategoryCollections(IsFiniteOrderElement)", "IsGeneralizedDomain",
    "CategoryCollections(IsPerm)", "IsMagma", "IsMagmaWithOne",
    "IsMagmaWithInversesIfNonzero", "IsMagmaWithInverses" ]
\end{Verbatim}
 }

 

\subsection{\textcolor{Chapter }{CategoryByName}}
\logpage{[ 13, 3, 3 ]}\nobreak
\hyperdef{L}{X85D07C3E7F4D4043}{}
{\noindent\textcolor{FuncColor}{$\triangleright$\enspace\texttt{CategoryByName({\mdseries\slshape name})\index{CategoryByName@\texttt{CategoryByName}}
\label{CategoryByName}
}\hfill{\scriptsize (function)}}\\


 returns the category with name \mbox{\texttt{\mdseries\slshape name}} if it is found, or fail otherwise. }

 

\subsection{\textcolor{Chapter }{NewCategory}}
\logpage{[ 13, 3, 4 ]}\nobreak
\hyperdef{L}{X87F68F887B44DBBD}{}
{\noindent\textcolor{FuncColor}{$\triangleright$\enspace\texttt{NewCategory({\mdseries\slshape name, super[, rank]})\index{NewCategory@\texttt{NewCategory}}
\label{NewCategory}
}\hfill{\scriptsize (function)}}\\


 \texttt{NewCategory} returns a new category \mbox{\texttt{\mdseries\slshape cat}} that has the name \mbox{\texttt{\mdseries\slshape name}} and is contained in the filter \mbox{\texttt{\mdseries\slshape super}}, see{\nobreakspace}\ref{Filters}. This means that every object in \mbox{\texttt{\mdseries\slshape cat}} lies automatically also in \mbox{\texttt{\mdseries\slshape super}}. We say also that \mbox{\texttt{\mdseries\slshape super}} is an implied filter of \mbox{\texttt{\mdseries\slshape cat}}. 

 For example, if one wants to create a category of group elements then \mbox{\texttt{\mdseries\slshape super}} should be \texttt{IsMultiplicativeElementWithInverse} (\ref{IsMultiplicativeElementWithInverse}) or a subcategory of it. If no specific supercategory of \mbox{\texttt{\mdseries\slshape cat}} is known, \mbox{\texttt{\mdseries\slshape super}} may be \texttt{IsObject} (\ref{IsObject}). 

 The optional third argument \mbox{\texttt{\mdseries\slshape rank}} denotes the incremental rank (see{\nobreakspace}\ref{Filters}) of \mbox{\texttt{\mdseries\slshape cat}}, the default value is 1. }

 

\subsection{\textcolor{Chapter }{DeclareCategory}}
\logpage{[ 13, 3, 5 ]}\nobreak
\hyperdef{L}{X879DE2A17A6C6E92}{}
{\noindent\textcolor{FuncColor}{$\triangleright$\enspace\texttt{DeclareCategory({\mdseries\slshape name, super[, rank]})\index{DeclareCategory@\texttt{DeclareCategory}}
\label{DeclareCategory}
}\hfill{\scriptsize (function)}}\\


 does the same as \texttt{NewCategory} (\ref{NewCategory}) and then binds the result to the global variable \mbox{\texttt{\mdseries\slshape name}}. The variable must previously be writable, and is made
read\texttt{\symbol{45}}only by this function. }

 

\subsection{\textcolor{Chapter }{CategoryFamily}}
\logpage{[ 13, 3, 6 ]}\nobreak
\hyperdef{L}{X787BACEE7937EF01}{}
{\noindent\textcolor{FuncColor}{$\triangleright$\enspace\texttt{CategoryFamily({\mdseries\slshape cat})\index{CategoryFamily@\texttt{CategoryFamily}}
\label{CategoryFamily}
}\hfill{\scriptsize (function)}}\\


 For a category \mbox{\texttt{\mdseries\slshape cat}}, \texttt{CategoryFamily} returns the \emph{family category} of \mbox{\texttt{\mdseries\slshape cat}}. This is a category in which all families lie that know from their creation
that all their elements are in the category \mbox{\texttt{\mdseries\slshape cat}}, see{\nobreakspace}\ref{Families}. 

 For example, a family of associative words is in the category \texttt{CategoryFamily( IsAssocWord )}, and one can distinguish such a family from others by this category. So it is
possible to install methods for operations that require one argument to be a
family of associative words. 

 \texttt{CategoryFamily} is quite technical, and in fact of minor importance. }

 

 See also \texttt{CategoryCollections} (\ref{CategoryCollections}). }

  
\section{\textcolor{Chapter }{Representation}}\label{Representation}
\logpage{[ 13, 4, 0 ]}
\hyperdef{L}{X8698205F8648EB33}{}
{
  The \emph{representation} of an object is a set of filters (see{\nobreakspace}\ref{Filters}) that determines how an object is actually represented. For example, a matrix
or a polynomial can be stored sparsely or densely; all dense polynomials form
a representation. An object which claims to lie in a certain representation is
accepting the requirement that certain fields in the data structure be present
and have specified meanings. 
\subsection{\textcolor{Chapter }{Basic Representations of Objects}}\logpage{[ 13, 4, 1 ]}
\hyperdef{L}{X805F1C3B7C730062}{}
{
\noindent\textcolor{FuncColor}{$\triangleright$\enspace\texttt{IsInternalRep({\mdseries\slshape obj})\index{IsInternalRep@\texttt{IsInternalRep}}
\label{IsInternalRep}
}\hfill{\scriptsize (Representation)}}\\
\noindent\textcolor{FuncColor}{$\triangleright$\enspace\texttt{IsDataObjectRep({\mdseries\slshape obj})\index{IsDataObjectRep@\texttt{IsDataObjectRep}}
\label{IsDataObjectRep}
}\hfill{\scriptsize (Representation)}}\\
\noindent\textcolor{FuncColor}{$\triangleright$\enspace\texttt{IsPositionalObjectRep({\mdseries\slshape obj})\index{IsPositionalObjectRep@\texttt{IsPositionalObjectRep}}
\label{IsPositionalObjectRep}
}\hfill{\scriptsize (Representation)}}\\
\noindent\textcolor{FuncColor}{$\triangleright$\enspace\texttt{IsComponentObjectRep({\mdseries\slshape obj})\index{IsComponentObjectRep@\texttt{IsComponentObjectRep}}
\label{IsComponentObjectRep}
}\hfill{\scriptsize (Representation)}}\\


 \textsf{GAP} distinguishes four essentially different ways to represent objects. First
there are the representations \texttt{IsInternalRep} for internal objects such as integers and permutations, and \texttt{IsDataObjectRep} for other objects that are created and whose data are accessible only by
kernel functions. The data structures underlying such objects cannot be
manipulated at the \textsf{GAP} level. 

 All other objects are either in the representation \texttt{IsComponentObjectRep} or in the representation \texttt{IsPositionalObjectRep}, see{\nobreakspace}\ref{Component Objects} and{\nobreakspace}\ref{Positional Objects}. 

 An object can belong to several representations in the sense that it lies in
several subrepresentations of \texttt{IsComponentObjectRep} or of \texttt{IsPositionalObjectRep}. The representations to which an object belongs should form a chain and
either two representations are disjoint or one is contained in the other. So
the subrepresentations of \texttt{IsComponentObjectRep} and \texttt{IsPositionalObjectRep} each form trees. In the language of Object Oriented Programming, we support
only single inheritance for representations. 

 These trees are typically rather shallow, since for one representation to be
contained in another implies that all the components of the data structure
implied by the containing representation, are present in, and have the same
meaning in, the smaller representation (whose data structure presumably
contains some additional components). 

 Objects may change their representation, for example a mutable list of
characters can be converted into a string. 

 All representations in the library are created during initialization, in
particular they are not created dynamically at runtime. 

 Examples of subrepresentations of \texttt{IsPositionalObjectRep} are \texttt{IsModulusRep}, which is used for residue classes in the ring of integers, and \texttt{IsDenseCoeffVectorRep}, which is used for elements of algebras that are defined by structure
constants. 

 An important subrepresentation of \texttt{IsComponentObjectRep} is \texttt{IsAttributeStoringRep} (\ref{IsAttributeStoringRep}), which is used for many domains and some other objects. It provides automatic
storing of all attribute values (see Section \ref{Attributes}). }

 

\subsection{\textcolor{Chapter }{IsRepresentation}}
\logpage{[ 13, 4, 2 ]}\nobreak
\hyperdef{L}{X86D42C7783ACA5F4}{}
{\noindent\textcolor{FuncColor}{$\triangleright$\enspace\texttt{IsRepresentation({\mdseries\slshape object})\index{IsRepresentation@\texttt{IsRepresentation}}
\label{IsRepresentation}
}\hfill{\scriptsize (function)}}\\


 returns \texttt{true} if \mbox{\texttt{\mdseries\slshape object}} is a representation (see{\nobreakspace}\ref{Representation}), and \texttt{false} otherwise. }

 

\subsection{\textcolor{Chapter }{RepresentationsOfObject}}
\logpage{[ 13, 4, 3 ]}\nobreak
\hyperdef{L}{X7BBE93BE7977750F}{}
{\noindent\textcolor{FuncColor}{$\triangleright$\enspace\texttt{RepresentationsOfObject({\mdseries\slshape object})\index{RepresentationsOfObject@\texttt{RepresentationsOfObject}}
\label{RepresentationsOfObject}
}\hfill{\scriptsize (operation)}}\\


 returns a list of the names of the representations \mbox{\texttt{\mdseries\slshape object}} has. 
\begin{Verbatim}[commandchars=!@|,fontsize=\small,frame=single,label=Example]
  !gapprompt@gap>| !gapinput@g:=Group((1,2),(1,2,3));;|
  !gapprompt@gap>| !gapinput@RepresentationsOfObject(g);|
  [ "IsComponentObjectRep", "IsAttributeStoringRep" ]
\end{Verbatim}
 }

 

\subsection{\textcolor{Chapter }{NewRepresentation}}
\logpage{[ 13, 4, 4 ]}\nobreak
\hyperdef{L}{X7CC8106F809E15CF}{}
{\noindent\textcolor{FuncColor}{$\triangleright$\enspace\texttt{NewRepresentation({\mdseries\slshape name, super[, slots[, req]]})\index{NewRepresentation@\texttt{NewRepresentation}}
\label{NewRepresentation}
}\hfill{\scriptsize (function)}}\\


 \texttt{NewRepresentation} returns a new representation \mbox{\texttt{\mdseries\slshape rep}} that has the name \mbox{\texttt{\mdseries\slshape name}} and is a subrepresentation of the representation \mbox{\texttt{\mdseries\slshape super}}. This means that every object in \mbox{\texttt{\mdseries\slshape rep}} lies automatically also in \mbox{\texttt{\mdseries\slshape super}}. We say also that \mbox{\texttt{\mdseries\slshape super}} is an implied filter of \mbox{\texttt{\mdseries\slshape rep}}. 

 Each representation in \textsf{GAP} is a subrepresentation of exactly one of the four representations \texttt{IsInternalRep} (\ref{IsInternalRep}), \texttt{IsDataObjectRep} (\ref{IsDataObjectRep}), \texttt{IsComponentObjectRep} (\ref{IsComponentObjectRep}), \texttt{IsPositionalObjectRep} (\ref{IsPositionalObjectRep}). The data describing objects in the former two can be accessed only via \textsf{GAP} kernel functions, the data describing objects in the latter two is accessible
also in library functions, see{\nobreakspace}\ref{Component Objects} and{\nobreakspace}\ref{Positional Objects} for the details. 

 The optional third and fourth arguments \mbox{\texttt{\mdseries\slshape slots}} and \mbox{\texttt{\mdseries\slshape req}} are (and always were) unused and are only provided for backwards
compatibility. Note that \mbox{\texttt{\mdseries\slshape slots}} was required (but still unused) before GAP 4.12. 

 The incremental rank (see{\nobreakspace}\ref{Filters}) of \mbox{\texttt{\mdseries\slshape rep}} is 1. 

 Examples for the use of \texttt{NewRepresentation} can be found in{\nobreakspace}\ref{Component Objects}, \ref{Positional Objects}, and also in \ref{A Second Attempt to Implement Elements of Residue Class Rings}. }

 

\subsection{\textcolor{Chapter }{DeclareRepresentation}}
\logpage{[ 13, 4, 5 ]}\nobreak
\hyperdef{L}{X7C81FB2682AE54CD}{}
{\noindent\textcolor{FuncColor}{$\triangleright$\enspace\texttt{DeclareRepresentation({\mdseries\slshape name, super[, slots[, req]]})\index{DeclareRepresentation@\texttt{DeclareRepresentation}}
\label{DeclareRepresentation}
}\hfill{\scriptsize (function)}}\\


 does the same as \texttt{NewRepresentation} (\ref{NewRepresentation}) and then binds the result to the global variable \mbox{\texttt{\mdseries\slshape name}}. The variable must previously be writable, and is made
read\texttt{\symbol{45}}only by this function. }

 }

  
\section{\textcolor{Chapter }{Attributes}}\label{Attributes}
\logpage{[ 13, 5, 0 ]}
\hyperdef{L}{X7C701DBF7BAE649A}{}
{
  The attributes of an object describe knowledge about it. 

 An attribute is a unary operation without side\texttt{\symbol{45}}effects. 

 An object may store values of its attributes once they have been computed, and
claim that it knows these values. Note that ``store'' and ``know'' have to be understood in the sense that it is very cheap to get such a value
when the attribute is called again. 

 The value returned by an attribute is in general immutable (see{\nobreakspace}\ref{Mutability and Copyability}), except if the attribute had been specially constructed as ``mutable attribute'', see \texttt{NewAttribute} (\ref{NewAttribute}). 

 Note that functions installed as methods for attributes may return mutable
objects, the attribute returns the value of \texttt{Immutable} (\ref{Immutable}) for the result. Thus it is recommended that these methods ensure their return
values are immutable, e.g., by calling \texttt{MakeImmutable} (\ref{MakeImmutable}) if this is possible, in order to avoid that \texttt{Immutable} (\ref{Immutable}) creates an unnecessary copy of the object. 

 It depends on the representation of an object (see{\nobreakspace}\ref{Representation}) which attribute values it stores. An immutable object in the representation \texttt{IsAttributeStoringRep} (\ref{IsAttributeStoringRep}) stores \emph{all} attribute values once they are computed. 

 Note that it is impossible to get rid of a stored attribute value because the
system may have drawn conclusions from the old attribute value, and just
removing the value might leave the data structures in an inconsistent state.
If necessary, a new object can be constructed. 

 Each method that is installed for an attribute via \texttt{InstallMethod} (\ref{InstallMethod}) must require exactly one argument, and this must lie in the filter \mbox{\texttt{\mdseries\slshape filter}} that was entered as second argument of \texttt{NewAttribute} (\ref{NewAttribute}) resp. \texttt{NewProperty} (\ref{NewProperty}). 

 As for any operation, for attributes one can install a method taking an
argument that does not lie in \mbox{\texttt{\mdseries\slshape filt}} via \texttt{InstallOtherMethod} (\ref{InstallOtherMethod}), or a method for more than one argument. For example, \texttt{IsTransitive} (\ref{IsTransitive:for a group, an action domain, etc.}) is an attribute for a $G$\texttt{\symbol{45}}set that can also be called for the two arguments, being a
group $G$ and its action domain. If attributes are called with more than one argument
then the return value is not stored in any of the arguments. 

 Properties are a special form of attributes that have the value \texttt{true} or \texttt{false}, see section{\nobreakspace}\ref{Properties}. 

 Examples of attributes for multiplicative elements are \texttt{Inverse} (\ref{Inverse}), \texttt{One} (\ref{One}), and \texttt{Order} (\ref{Order}). \texttt{Size} (\ref{Size}) is an attribute for domains, \texttt{Centre} (\ref{Centre}) is an attribute for magmas, and \texttt{DerivedSubgroup} (\ref{DerivedSubgroup}) is an attribute for groups. 

\subsection{\textcolor{Chapter }{IsAttribute}}
\logpage{[ 13, 5, 1 ]}\nobreak
\hyperdef{L}{X7973C8F4782D15A1}{}
{\noindent\textcolor{FuncColor}{$\triangleright$\enspace\texttt{IsAttribute({\mdseries\slshape object})\index{IsAttribute@\texttt{IsAttribute}}
\label{IsAttribute}
}\hfill{\scriptsize (function)}}\\


 returns \texttt{true} if \mbox{\texttt{\mdseries\slshape object}} is an attribute (see{\nobreakspace}\ref{Attributes}), and \texttt{false} otherwise. }

 

\subsection{\textcolor{Chapter }{KnownAttributesOfObject}}
\logpage{[ 13, 5, 2 ]}\nobreak
\hyperdef{L}{X7F7960338163AA88}{}
{\noindent\textcolor{FuncColor}{$\triangleright$\enspace\texttt{KnownAttributesOfObject({\mdseries\slshape object})\index{KnownAttributesOfObject@\texttt{KnownAttributesOfObject}}
\label{KnownAttributesOfObject}
}\hfill{\scriptsize (operation)}}\\


 returns a list of the names of the attributes whose values are known for \mbox{\texttt{\mdseries\slshape object}}. 
\begin{Verbatim}[commandchars=!@|,fontsize=\small,frame=single,label=Example]
  !gapprompt@gap>| !gapinput@g:=Group((1,2),(1,2,3));;Size(g);;|
  !gapprompt@gap>| !gapinput@KnownAttributesOfObject(g);|
  [ "Size", "OneImmutable", "NrMovedPoints", "MovedPoints",
    "GeneratorsOfMagmaWithInverses", "MultiplicativeNeutralElement",
    "HomePcgs", "Pcgs", "StabChainMutable", "StabChainOptions" ]
\end{Verbatim}
 }

 

\subsection{\textcolor{Chapter }{NewAttribute}}
\logpage{[ 13, 5, 3 ]}\nobreak
\hyperdef{L}{X7B9654807858A3B0}{}
{\noindent\textcolor{FuncColor}{$\triangleright$\enspace\texttt{NewAttribute({\mdseries\slshape name, filter[, "mutable"][, rank]})\index{NewAttribute@\texttt{NewAttribute}}
\label{NewAttribute}
}\hfill{\scriptsize (function)}}\\


 \texttt{NewAttribute} returns a new attribute getter with name \mbox{\texttt{\mdseries\slshape name}} that is applicable to objects with the property \mbox{\texttt{\mdseries\slshape filter}}. 

 Contrary to the situation with categories and representations, the tester of
the new attribute does \emph{not} imply \mbox{\texttt{\mdseries\slshape filter}}. This is exactly because of the possibility to install methods that do not
require \mbox{\texttt{\mdseries\slshape filter}}. 

 For example, the attribute \texttt{Size} (\ref{Size}) was created with second argument a list or a collection, but there is also a
method for \texttt{Size} (\ref{Size}) that is applicable to a character table, which is neither a list nor a
collection. 

 For the optional third and fourth arguments, there are the following
possibilities. 
\begin{itemize}
\item  The integer argument \mbox{\texttt{\mdseries\slshape rank}} causes the attribute tester to have this incremental rank (see{\nobreakspace}\ref{Filters}), 
\item  If the argument \mbox{\texttt{\mdseries\slshape mutable}} is the string \texttt{"mutable"} or the boolean \texttt{true}, then the values of the attribute are mutable. 
\item  If the argument \mbox{\texttt{\mdseries\slshape mutable}} is the boolean \texttt{false}, then the values of the attribute are immutable. 
\end{itemize}
 

 When a value of such mutable attribute is set then this value itself is
stored, not an immutable copy of it, and it is the user's responsibility to
set an object that is mutable. This is useful for an attribute whose value is
some partial information that may be completed later. For example, there is an
attribute \texttt{ComputedSylowSubgroups} for the list holding those Sylow subgroups of a group that have been computed
already by the function \texttt{SylowSubgroup} (\ref{SylowSubgroup}), and this list is mutable because one may want to enter groups into it as
they are computed. 

 Mutable and immutable attributes behave differently w. r. t. repeated calls of
the function that sets the attribute value, see \texttt{Setter} (\ref{Setter}). 

 If no argument for \mbox{\texttt{\mdseries\slshape rank}} is given, then the rank of the tester is 1. 

 Each method for the new attribute that does \emph{not} require its argument to lie in \mbox{\texttt{\mdseries\slshape filter}} must be installed using \texttt{InstallOtherMethod} (\ref{InstallOtherMethod}). }

 

\subsection{\textcolor{Chapter }{DeclareAttribute}}
\logpage{[ 13, 5, 4 ]}\nobreak
\hyperdef{L}{X7A00FC8A7A677A56}{}
{\noindent\textcolor{FuncColor}{$\triangleright$\enspace\texttt{DeclareAttribute({\mdseries\slshape name, filter[, "mutable"][, rank]})\index{DeclareAttribute@\texttt{DeclareAttribute}}
\label{DeclareAttribute}
}\hfill{\scriptsize (function)}}\\


 does the same as \texttt{NewAttribute} (\ref{NewAttribute}) and then binds the result to the global variable \mbox{\texttt{\mdseries\slshape name}}. The variable must previously be writable, and is made
read\texttt{\symbol{45}}only by this function. It also binds
read\texttt{\symbol{45}}only global variables with names \texttt{Has\mbox{\texttt{\mdseries\slshape name}}} and \texttt{Set\mbox{\texttt{\mdseries\slshape name}}} for the tester and setter of the attribute (see Section \ref{Setter and Tester for Attributes}). }

 

\subsection{\textcolor{Chapter }{IsAttributeStoringRep}}
\logpage{[ 13, 5, 5 ]}\nobreak
\hyperdef{L}{X7A951C33839AF2C1}{}
{\noindent\textcolor{FuncColor}{$\triangleright$\enspace\texttt{IsAttributeStoringRep({\mdseries\slshape obj})\index{IsAttributeStoringRep@\texttt{IsAttributeStoringRep}}
\label{IsAttributeStoringRep}
}\hfill{\scriptsize (Representation)}}\\


 Objects in this representation have default methods to get stored values of
attributes and {\textendash}if they are immutable{\textendash} to store
attribute values automatically once they have been computed. \index{system getter} \index{system setter} (These methods are called the ``system getter'' and the ``system setter'' of the attribute, respectively.) 

 As a consequence, for immutable objects in \texttt{IsAttributeStoringRep}, subsequent calls to an attribute will return the \emph{same} object. 

 \emph{Mutable} objects in \texttt{IsAttributeStoringRep} are allowed, but attribute values are not stored automatically in them. Such
objects are useful because they may later be made immutable using \texttt{MakeImmutable} (\ref{MakeImmutable}), at which point they will start storing all attribute values. 

 Note that one can force an attribute value to be stored in a mutable object in \texttt{IsAttributeStoringRep}, by explicitly calling the attribute setter. This feature should be used with
care. For example, think of a mutable matrix whose rank or trace gets stored,
and the values later become wrong when somebody changes the matrix entries. 

 
\begin{Verbatim}[commandchars=!@|,fontsize=\small,frame=single,label=Example]
  !gapprompt@gap>| !gapinput@g:= Group( (1,2)(3,4), (1,3)(2,4) );;|
  !gapprompt@gap>| !gapinput@IsAttributeStoringRep( g );|
  true
  !gapprompt@gap>| !gapinput@HasSize( g );  Size( g );  HasSize( g );|
  false
  4
  true
  !gapprompt@gap>| !gapinput@r:= 7/4;;|
  !gapprompt@gap>| !gapinput@IsAttributeStoringRep( r );|
  false
  !gapprompt@gap>| !gapinput@Int( r );  HasInt( r );|
  1
  false
\end{Verbatim}
 }

 }

  
\section{\textcolor{Chapter }{Setter and Tester for Attributes}}\label{Setter and Tester for Attributes}
\logpage{[ 13, 6, 0 ]}
\hyperdef{L}{X79DE5208877AE42A}{}
{
  \index{setter} \index{tester} For every attribute, the \emph{attribute setter} and the \emph{attribute tester} are defined. 

 To check whether an object belongs to an attribute \mbox{\texttt{\mdseries\slshape attr}}, the tester of the attribute is used, see \texttt{Tester} (\ref{Tester}). To store a value for the attribute \mbox{\texttt{\mdseries\slshape attr}} in an object, the setter of the attribute is used, see \texttt{Setter} (\ref{Setter}). 

\subsection{\textcolor{Chapter }{Tester}}
\logpage{[ 13, 6, 1 ]}\nobreak
\hyperdef{L}{X87D5B5AC7DAF932D}{}
{\noindent\textcolor{FuncColor}{$\triangleright$\enspace\texttt{Tester({\mdseries\slshape attr})\index{Tester@\texttt{Tester}}
\label{Tester}
}\hfill{\scriptsize (function)}}\\


 For an attribute \mbox{\texttt{\mdseries\slshape attr}}, \texttt{Tester(\mbox{\texttt{\mdseries\slshape attr}})} is a filter (see{\nobreakspace}\ref{Filters}) that returns \texttt{true} or \texttt{false}, depending on whether or not the value of \mbox{\texttt{\mdseries\slshape attr}} for the object is known. For example, \texttt{Tester( Size )( \mbox{\texttt{\mdseries\slshape obj}} )} is \texttt{true} if the size of the object \mbox{\texttt{\mdseries\slshape obj}} is known. }

 

\subsection{\textcolor{Chapter }{Setter}}
\logpage{[ 13, 6, 2 ]}\nobreak
\hyperdef{L}{X7FD8952C841D2B1F}{}
{\noindent\textcolor{FuncColor}{$\triangleright$\enspace\texttt{Setter({\mdseries\slshape attr})\index{Setter@\texttt{Setter}}
\label{Setter}
}\hfill{\scriptsize (function)}}\\


 For an attribute \mbox{\texttt{\mdseries\slshape attr}}, \texttt{Setter(\mbox{\texttt{\mdseries\slshape attr}})} is called automatically when the attribute value has been computed for an
immutable object which does not already have a value stored for \mbox{\texttt{\mdseries\slshape attr}}. One can also call the setter explicitly, for example, \texttt{Setter( Size )( \mbox{\texttt{\mdseries\slshape obj}}, \mbox{\texttt{\mdseries\slshape val}} )} sets \mbox{\texttt{\mdseries\slshape val}} as size of the object \mbox{\texttt{\mdseries\slshape obj}} if the size was not yet known. 

 Calling the setter an of attribute for an object which stores already a value
for this attribute \emph{has no effect}, except if the attribute has been created as a \emph{mutable} attribute (see \texttt{NewAttribute} (\ref{NewAttribute})). For mutable attributes, each call of the setter \emph{replaces the stored value}. (Usually mutable attributes are used to store a mutable list or record once,
and later add data to this value, and then one should take care not to call
the setter of the attribute again.) 

 
\begin{Verbatim}[commandchars=!@|,fontsize=\small,frame=single,label=Example]
  !gapprompt@gap>| !gapinput@imm_attr:= NewAttribute( "imm_attr", IsGroup );;|
  !gapprompt@gap>| !gapinput@mut_attr:= NewAttribute( "mut_attr", IsGroup, "mutable" );;|
  !gapprompt@gap>| !gapinput@G:= SymmetricGroup( 4 );;|
  !gapprompt@gap>| !gapinput@Setter( imm_attr )( G, 0 );|
  !gapprompt@gap>| !gapinput@imm_attr( G );|
  0
  !gapprompt@gap>| !gapinput@Setter( imm_attr )( G, 1 );  # This call has no effect.|
  !gapprompt@gap>| !gapinput@imm_attr( G );|
  0
  !gapprompt@gap>| !gapinput@Setter( mut_attr )( G, 0 );|
  !gapprompt@gap>| !gapinput@mut_attr( G );|
  0
  !gapprompt@gap>| !gapinput@Setter( mut_attr )( G, 1 );  # This call replaces the stored value.|
  !gapprompt@gap>| !gapinput@mut_attr( G );|
  1
\end{Verbatim}
 

 For each attribute \mbox{\texttt{\mdseries\slshape attr}} that is declared with \texttt{DeclareAttribute} (\ref{DeclareAttribute}) resp.{\nobreakspace}\texttt{DeclareProperty} (\ref{DeclareProperty}), tester and setter are automatically made accessible by the names \texttt{Has\mbox{\texttt{\mdseries\slshape attr}}} and \texttt{Set\mbox{\texttt{\mdseries\slshape attr}}}, respectively. For example, the tester for \texttt{Size} (\ref{Size}) is called \texttt{HasSize}, and the setter is called \texttt{SetSize}. 

 
\begin{Verbatim}[commandchars=!@|,fontsize=\small,frame=single,label=Example]
  !gapprompt@gap>| !gapinput@g:=Group((1,2,3,4),(1,2));;Size(g);|
  24
  !gapprompt@gap>| !gapinput@HasSize(g);|
  true
  !gapprompt@gap>| !gapinput@SetSize(g,99);|
  !gapprompt@gap>| !gapinput@Size(g);|
  24
\end{Verbatim}
 

 For two properties \mbox{\texttt{\mdseries\slshape prop1}} and \mbox{\texttt{\mdseries\slshape prop2}}, the intersection \texttt{\mbox{\texttt{\mdseries\slshape prop1}} and \mbox{\texttt{\mdseries\slshape prop2}}} (see{\nobreakspace}\ref{Filters}) is again a property for which a setter and a tester exist. Setting the value
of this intersection to \texttt{true} for a \textsf{GAP} object means to set the values of \mbox{\texttt{\mdseries\slshape prop1}} and \mbox{\texttt{\mdseries\slshape prop2}} to \texttt{true} for this object. 

 
\begin{Verbatim}[commandchars=!@|,fontsize=\small,frame=single,label=Example]
  !gapprompt@gap>| !gapinput@prop:= IsFinite and IsCommutative;|
  <Property "(IsFinite and IsCommutative)">
  !gapprompt@gap>| !gapinput@g:= Group( (1,2,3), (4,5) );;|
  !gapprompt@gap>| !gapinput@Tester( prop )( g );|
  false
  !gapprompt@gap>| !gapinput@Setter( prop )( g, true );|
  !gapprompt@gap>| !gapinput@Tester( prop )( g );  prop( g );|
  true
  true
\end{Verbatim}
 

 It is \emph{not allowed} to set the value of such an intersection to \texttt{false} for an object. 

 
\begin{Verbatim}[commandchars=!@|,fontsize=\small,frame=single,label=Example]
  !gapprompt@gap>| !gapinput@Setter( prop )( Rationals, false );|
  Error, You cannot set an "and-filter" except to true
  not in any function at *stdin*:8
  you can enter 'quit;' to quit to outer loop
  !gapbrkprompt@brk>| !gapinput@|
\end{Verbatim}
 }

 

\subsection{\textcolor{Chapter }{AttributeValueNotSet}}
\logpage{[ 13, 6, 3 ]}\nobreak
\hyperdef{L}{X8529F8A17884A32C}{}
{\noindent\textcolor{FuncColor}{$\triangleright$\enspace\texttt{AttributeValueNotSet({\mdseries\slshape attr, obj})\index{AttributeValueNotSet@\texttt{AttributeValueNotSet}}
\label{AttributeValueNotSet}
}\hfill{\scriptsize (function)}}\\


 If the value of the attribute \mbox{\texttt{\mdseries\slshape attr}} is already stored for \mbox{\texttt{\mdseries\slshape obj}}, \texttt{AttributeValueNotSet} simply returns this value. Otherwise the value of \texttt{\mbox{\texttt{\mdseries\slshape attr}}( \mbox{\texttt{\mdseries\slshape obj}} )} is computed and returned \emph{without storing it} in \mbox{\texttt{\mdseries\slshape obj}}. This can be useful when ``large'' attribute values (such as element lists) are needed only once and shall not be
stored in the object. 

 
\begin{Verbatim}[commandchars=!@|,fontsize=\small,frame=single,label=Example]
  !gapprompt@gap>| !gapinput@HasAsSSortedList(g);|
  false
  !gapprompt@gap>| !gapinput@AttributeValueNotSet(AsSSortedList,g);|
  [ (), (4,5), (1,2,3), (1,2,3)(4,5), (1,3,2), (1,3,2)(4,5) ]
  !gapprompt@gap>| !gapinput@HasAsSSortedList(g);|
  false
\end{Verbatim}
 

 The normal behaviour of attributes (when called with just one argument) is
that once a method has been selected and executed, and has returned a value
the setter of the attribute is called, to (possibly) store the computed value.
In special circumstances, this behaviour can be altered dynamically on an
attribute\texttt{\symbol{45}}by\texttt{\symbol{45}}attribute basis, using the
functions \texttt{DisableAttributeValueStoring} (\ref{DisableAttributeValueStoring}) and \texttt{EnableAttributeValueStoring} (\ref{EnableAttributeValueStoring}). 

 In general, the code in the library assumes, for efficiency, but not for
correctness, that attribute values \emph{will} be stored (in suitable objects), so disabling storing may cause substantial
computations to be repeated. }

 

\subsection{\textcolor{Chapter }{InfoAttributes}}
\logpage{[ 13, 6, 4 ]}\nobreak
\hyperdef{L}{X79120CE37BB69D11}{}
{\noindent\textcolor{FuncColor}{$\triangleright$\enspace\texttt{InfoAttributes\index{InfoAttributes@\texttt{InfoAttributes}}
\label{InfoAttributes}
}\hfill{\scriptsize (info class)}}\\


 This info class (together with \texttt{InfoWarning} (\ref{InfoWarning})) is used for messages about attributes. Messages are shown under the
following circumstances: 
\begin{itemize}
\item  \texttt{EnableAttributeValueStoring} (\ref{EnableAttributeValueStoring}) is used (level 2).
\item  \texttt{DisableAttributeValueStoring} (\ref{DisableAttributeValueStoring}) is used (level 3).
\item  When trying to assign to non\texttt{\symbol{45}}mutable attribute which
already is set to a different value (level 3).
\item  When the test filter for an attribute (i.e., \texttt{HasFOO}) is set, but no value is assigned (level 3).
\end{itemize}
 }

 

\subsection{\textcolor{Chapter }{DisableAttributeValueStoring}}
\logpage{[ 13, 6, 5 ]}\nobreak
\hyperdef{L}{X7851E2DB79656DB0}{}
{\noindent\textcolor{FuncColor}{$\triangleright$\enspace\texttt{DisableAttributeValueStoring({\mdseries\slshape attr})\index{DisableAttributeValueStoring@\texttt{DisableAttributeValueStoring}}
\label{DisableAttributeValueStoring}
}\hfill{\scriptsize (function)}}\\


 disables the usual call of \texttt{Setter( \mbox{\texttt{\mdseries\slshape attr}} )} when a method for \mbox{\texttt{\mdseries\slshape attr}} returns a value. In consequence the values will never be stored. Note that \mbox{\texttt{\mdseries\slshape attr}} must be an attribute and \emph{not} a property. }

 

\subsection{\textcolor{Chapter }{EnableAttributeValueStoring}}
\logpage{[ 13, 6, 6 ]}\nobreak
\hyperdef{L}{X7E5DACBE7A9A9AD1}{}
{\noindent\textcolor{FuncColor}{$\triangleright$\enspace\texttt{EnableAttributeValueStoring({\mdseries\slshape attr})\index{EnableAttributeValueStoring@\texttt{EnableAttributeValueStoring}}
\label{EnableAttributeValueStoring}
}\hfill{\scriptsize (function)}}\\


 enables the usual call of \texttt{Setter( \mbox{\texttt{\mdseries\slshape attr}} )} when a method for \mbox{\texttt{\mdseries\slshape attr}} returns a value. In consequence the values may be stored. This will usually
have no effect unless \texttt{DisableAttributeValueStoring} (\ref{DisableAttributeValueStoring}) has previously been used for \mbox{\texttt{\mdseries\slshape attr}}. Note that \mbox{\texttt{\mdseries\slshape attr}} must be an attribute and \emph{not} a property. 
\begin{Verbatim}[commandchars=!@|,fontsize=\small,frame=single,label=Example]
  !gapprompt@gap>| !gapinput@g := Group((1,2,3,4,5),(1,2,3));|
  Group([ (1,2,3,4,5), (1,2,3) ])
  !gapprompt@gap>| !gapinput@KnownAttributesOfObject(g);|
  [ "LargestMovedPoint", "GeneratorsOfMagmaWithInverses",
    "MultiplicativeNeutralElement" ]
  !gapprompt@gap>| !gapinput@SetInfoLevel(InfoAttributes,3);|
  !gapprompt@gap>| !gapinput@DisableAttributeValueStoring(Size);|
  #I  Disabling value storing for Size
  !gapprompt@gap>| !gapinput@Size(g);|
  60
  !gapprompt@gap>| !gapinput@KnownAttributesOfObject(g);|
  [ "OneImmutable", "LargestMovedPoint", "NrMovedPoints",
    "MovedPoints", "GeneratorsOfMagmaWithInverses",
    "MultiplicativeNeutralElement", "StabChainMutable",
    "StabChainOptions" ]
  !gapprompt@gap>| !gapinput@Size(g);|
  60
  !gapprompt@gap>| !gapinput@EnableAttributeValueStoring(Size);|
  #I  Enabling value storing for Size
  !gapprompt@gap>| !gapinput@Size(g);|
  60
  !gapprompt@gap>| !gapinput@KnownAttributesOfObject(g);|
  [ "Size", "OneImmutable", "LargestMovedPoint", "NrMovedPoints",
    "MovedPoints", "GeneratorsOfMagmaWithInverses",
    "MultiplicativeNeutralElement", "StabChainMutable",
    "StabChainOptions" ]
\end{Verbatim}
 }

 }

  
\section{\textcolor{Chapter }{Properties}}\label{Properties}
\logpage{[ 13, 7, 0 ]}
\hyperdef{L}{X871597447BB998A1}{}
{
  The properties of an object are those of its attributes (see{\nobreakspace}\ref{Attributes}) whose values can only be \texttt{true} or \texttt{false}. 

 The main difference between attributes and properties is that a property
defines two sets of objects, namely the usual set of all objects for which the
value is known, and the set of all objects for which the value is known to be \texttt{true}. 

 (Note that it makes no sense to consider a third set, namely the set of
objects for which the value of a property is \texttt{true} whether or not it is known, since there may be objects for which the
containment in this set cannot be decided.) 

 For a property \mbox{\texttt{\mdseries\slshape prop}}, the containment of an object \mbox{\texttt{\mdseries\slshape obj}} in the first set is checked again by applying \texttt{Tester( \mbox{\texttt{\mdseries\slshape prop}} )} to \mbox{\texttt{\mdseries\slshape obj}}, and \mbox{\texttt{\mdseries\slshape obj}} lies in the second set if and only if \texttt{Tester( \mbox{\texttt{\mdseries\slshape prop}} )( \mbox{\texttt{\mdseries\slshape obj}} ) and \mbox{\texttt{\mdseries\slshape prop}}( \mbox{\texttt{\mdseries\slshape obj}} )} is \texttt{true}. 

 If a property value is known for an immutable object then this value is also
stored, as part of the type of the object. To some extent, property values of
mutable objects also can be stored, for example a mutable list all of whose
entries are immutable can store whether it is strictly sorted. When the object
is mutated (for example by list assignment) the type may need to be adjusted. 

 Important properties for domains are \texttt{IsAssociative} (\ref{IsAssociative}), \texttt{IsCommutative} (\ref{IsCommutative}), \texttt{IsAnticommutative} (\ref{IsAnticommutative}), \texttt{IsLDistributive} (\ref{IsLDistributive}) and \texttt{IsRDistributive} (\ref{IsRDistributive}), which mean that the multiplication of elements in the domain satisfies $( a * b ) * c = a * ( b * c )$, $a * b = b * a$, $a * b = - ( b * a )$, $a * ( b + c ) = a * b + a * c$, and $( a + b ) * c = a * c + b * c$, respectively, for all $a$, $b$, $c$ in the domain. 

 

\subsection{\textcolor{Chapter }{IsProperty}}
\logpage{[ 13, 7, 1 ]}\nobreak
\hyperdef{L}{X81F1C3EE83003FA0}{}
{\noindent\textcolor{FuncColor}{$\triangleright$\enspace\texttt{IsProperty({\mdseries\slshape object})\index{IsProperty@\texttt{IsProperty}}
\label{IsProperty}
}\hfill{\scriptsize (function)}}\\


 returns \texttt{true} if \mbox{\texttt{\mdseries\slshape object}} is a property (see{\nobreakspace}\ref{Properties}), and \texttt{false} otherwise. }

 

\subsection{\textcolor{Chapter }{KnownPropertiesOfObject}}
\logpage{[ 13, 7, 2 ]}\nobreak
\hyperdef{L}{X7E51C08286E03E7F}{}
{\noindent\textcolor{FuncColor}{$\triangleright$\enspace\texttt{KnownPropertiesOfObject({\mdseries\slshape object})\index{KnownPropertiesOfObject@\texttt{KnownPropertiesOfObject}}
\label{KnownPropertiesOfObject}
}\hfill{\scriptsize (operation)}}\\


 returns a list of the names of the properties whose values are known for \mbox{\texttt{\mdseries\slshape object}}. }

 

\subsection{\textcolor{Chapter }{KnownTruePropertiesOfObject}}
\logpage{[ 13, 7, 3 ]}\nobreak
\hyperdef{L}{X86711BC77B62EB02}{}
{\noindent\textcolor{FuncColor}{$\triangleright$\enspace\texttt{KnownTruePropertiesOfObject({\mdseries\slshape object})\index{KnownTruePropertiesOfObject@\texttt{KnownTruePropertiesOfObject}}
\label{KnownTruePropertiesOfObject}
}\hfill{\scriptsize (operation)}}\\


 returns a list of the names of the properties known to be \texttt{true} for \mbox{\texttt{\mdseries\slshape object}}. 
\begin{Verbatim}[commandchars=!@|,fontsize=\small,frame=single,label=Example]
  !gapprompt@gap>| !gapinput@g:=Group((1,2),(1,2,3));;|
  !gapprompt@gap>| !gapinput@KnownPropertiesOfObject(g);|
  [ "IsEmpty", "IsTrivial", "IsNonTrivial", "IsFinite",
    "CanEasilyCompareElements", "CanEasilySortElements",
    "IsDuplicateFree", "IsGeneratorsOfMagmaWithInverses",
    "IsAssociative", "IsFinitelyGeneratedMagma",
    "IsGeneratorsOfSemigroup", "IsSimpleSemigroup",
    "IsRegularSemigroup", "IsInverseSemigroup",
    "IsCompletelyRegularSemigroup", "IsCompletelySimpleSemigroup",
    "IsGroupAsSemigroup", "IsMonoidAsSemigroup", "IsOrthodoxSemigroup",
    "IsFinitelyGeneratedMonoid", "IsFinitelyGeneratedGroup",
    "IsSubsetLocallyFiniteGroup", "KnowsHowToDecompose",
    "IsInfiniteAbelianizationGroup", "IsNilpotentByFinite",
    "IsTorsionFree", "IsFreeAbelian" ]
  !gapprompt@gap>| !gapinput@Size(g);|
  6
  !gapprompt@gap>| !gapinput@KnownPropertiesOfObject(g);|
  [ "IsEmpty", "IsTrivial", "IsNonTrivial", "IsFinite",
    "CanEasilyCompareElements", "CanEasilySortElements",
    "IsDuplicateFree", "IsGeneratorsOfMagmaWithInverses",
    "IsAssociative", "IsFinitelyGeneratedMagma",
    "IsGeneratorsOfSemigroup", "IsSimpleSemigroup",
    "IsRegularSemigroup", "IsInverseSemigroup",
    "IsCompletelyRegularSemigroup", "IsCompletelySimpleSemigroup",
    "IsGroupAsSemigroup", "IsMonoidAsSemigroup", "IsOrthodoxSemigroup",
    "IsFinitelyGeneratedMonoid", "IsFinitelyGeneratedGroup",
    "IsSubsetLocallyFiniteGroup", "KnowsHowToDecompose",
    "IsPerfectGroup", "IsSolvableGroup", "IsPolycyclicGroup",
    "IsInfiniteAbelianizationGroup", "IsNilpotentByFinite",
    "IsTorsionFree", "IsFreeAbelian" ]
  !gapprompt@gap>| !gapinput@KnownTruePropertiesOfObject(g);|
  [ "IsNonTrivial", "IsFinite", "CanEasilyCompareElements",
    "CanEasilySortElements", "IsDuplicateFree",
    "IsGeneratorsOfMagmaWithInverses", "IsAssociative",
    "IsFinitelyGeneratedMagma", "IsGeneratorsOfSemigroup",
    "IsSimpleSemigroup", "IsRegularSemigroup", "IsInverseSemigroup",
    "IsCompletelyRegularSemigroup", "IsCompletelySimpleSemigroup",
    "IsGroupAsSemigroup", "IsMonoidAsSemigroup", "IsOrthodoxSemigroup",
    "IsFinitelyGeneratedMonoid", "IsFinitelyGeneratedGroup",
    "IsSubsetLocallyFiniteGroup", "KnowsHowToDecompose",
    "IsSolvableGroup", "IsPolycyclicGroup", "IsNilpotentByFinite" ]
\end{Verbatim}
 }

 

\subsection{\textcolor{Chapter }{NewProperty}}
\logpage{[ 13, 7, 4 ]}\nobreak
\hyperdef{L}{X7F2D6FD979FE23DD}{}
{\noindent\textcolor{FuncColor}{$\triangleright$\enspace\texttt{NewProperty({\mdseries\slshape name, filter[, rank]})\index{NewProperty@\texttt{NewProperty}}
\label{NewProperty}
}\hfill{\scriptsize (function)}}\\


 \texttt{NewProperty} returns a new property \mbox{\texttt{\mdseries\slshape prop}} with name \mbox{\texttt{\mdseries\slshape name}} (see also{\nobreakspace}\ref{Properties}). The filter \mbox{\texttt{\mdseries\slshape filter}} describes the involved filters of \mbox{\texttt{\mdseries\slshape prop}}. As in the case of attributes, \mbox{\texttt{\mdseries\slshape filter}} is not implied by \mbox{\texttt{\mdseries\slshape prop}}. 

 The optional third argument \mbox{\texttt{\mdseries\slshape rank}} denotes the incremental rank (see{\nobreakspace}\ref{Filters}) of the property \mbox{\texttt{\mdseries\slshape prop}} itself, i.e. \emph{not} of its tester; the default value is 1. }

 

\subsection{\textcolor{Chapter }{DeclareProperty}}
\logpage{[ 13, 7, 5 ]}\nobreak
\hyperdef{L}{X7F4602F082682A04}{}
{\noindent\textcolor{FuncColor}{$\triangleright$\enspace\texttt{DeclareProperty({\mdseries\slshape name, filter[, rank]})\index{DeclareProperty@\texttt{DeclareProperty}}
\label{DeclareProperty}
}\hfill{\scriptsize (function)}}\\


 does the same as \texttt{NewProperty} (\ref{NewProperty}) and then binds the result to the global variable \mbox{\texttt{\mdseries\slshape name}}. The variable must previously be writable, and is made
read\texttt{\symbol{45}}only by this function. It also binds
read\texttt{\symbol{45}}only global variables with names \texttt{Has\mbox{\texttt{\mdseries\slshape name}}} and \texttt{Set\mbox{\texttt{\mdseries\slshape name}}} for the tester and setter of the property (see Section \ref{Setter and Tester for Attributes}). }

 }

  
\section{\textcolor{Chapter }{Other Filters}}\label{Other Filters}
\logpage{[ 13, 8, 0 ]}
\hyperdef{L}{X7997705185C7E720}{}
{
  There are situations where one wants to express a kind of knowledge that is
based on some heuristic. 

 For example, the filters (see{\nobreakspace}\ref{Filters}) \texttt{CanEasilyTestMembership} (\ref{CanEasilyTestMembership}) and \texttt{CanEasilyComputePcgs} (\ref{CanEasilyComputePcgs}) are defined in the \textsf{GAP} library. Note that such filters do not correspond to a mathematical concept,
contrary to properties (see{\nobreakspace}\ref{Properties}). Also it need not be defined what ``easily'' means for an arbitrary \textsf{GAP} object, and in this case one cannot compute the value for an arbitrary \textsf{GAP} object. In order to access this kind of knowledge as a part of the type of an
object, \textsf{GAP} provides filters for which the value is \texttt{false} by default, and it is changed to \texttt{true} in certain situations, either explicitly (for the given object) or via a
logical implication (see{\nobreakspace}\ref{Logical Implications}) from other filters. 

 For example, a \texttt{true} value of \texttt{CanEasilyComputePcgs} (\ref{CanEasilyComputePcgs}) for a group means that certain methods are applicable that use a pcgs
(see{\nobreakspace}\ref{Polycyclic Generating Systems}) for the group. There are logical implications to set the filter value to \texttt{true} for permutation groups that are known to be solvable, and for groups that have
already a (sufficiently nice) pcgs stored. In the case one has a solvable
matrix group and wants to enable methods that use a pcgs, one can set the \texttt{CanEasilyComputePcgs} (\ref{CanEasilyComputePcgs}) value to \texttt{true} for this particular group. 

 A filter \mbox{\texttt{\mdseries\slshape filt}} of the kind described here is different from the other filters introduced in
the previous sections. In particular, \mbox{\texttt{\mdseries\slshape filt}} is not a category (see{\nobreakspace}\ref{Categories}) or a property (see{\nobreakspace}\ref{Properties}) because its value may change for a given object, and \mbox{\texttt{\mdseries\slshape filt}} is not a representation (see{\nobreakspace}\ref{Representation}) because it has nothing to do with the way an object is made up from some
data. \mbox{\texttt{\mdseries\slshape filt}} is similar to an attribute tester (see{\nobreakspace}\ref{Setter and Tester for Attributes}), the only difference is that \mbox{\texttt{\mdseries\slshape filt}} does not refer to an attribute value; note that \mbox{\texttt{\mdseries\slshape filt}} is also used in the same way as an attribute tester; namely, the \texttt{true} value may be required for certain methods to be applicable. 

 In order to change the value of \mbox{\texttt{\mdseries\slshape filt}} for an object \mbox{\texttt{\mdseries\slshape obj}}, one can use logical implications (see{\nobreakspace}\ref{Logical Implications}) or \texttt{SetFilterObj} (\ref{SetFilterObj}), \texttt{ResetFilterObj} (\ref{ResetFilterObj}). 

\subsection{\textcolor{Chapter }{NewFilter}}
\logpage{[ 13, 8, 1 ]}\nobreak
\hyperdef{L}{X821635DA7821ED74}{}
{\noindent\textcolor{FuncColor}{$\triangleright$\enspace\texttt{NewFilter({\mdseries\slshape name[, implied][, rank]})\index{NewFilter@\texttt{NewFilter}}
\label{NewFilter}
}\hfill{\scriptsize (function)}}\\


 \texttt{NewFilter} returns a simple filter with name \mbox{\texttt{\mdseries\slshape name}} (see{\nobreakspace}\ref{Other Filters}). 

 The optional argument \mbox{\texttt{\mdseries\slshape implied}}, if given, must be a filter, meaning that for each object in the new filter,
also \mbox{\texttt{\mdseries\slshape implied}} will be set. Note that resetting the new filter with \texttt{ResetFilterObj} (\ref{ResetFilterObj}) does \emph{not} reset \mbox{\texttt{\mdseries\slshape implied}}. If the new filter is intended to be set or reset manually for existing
objects then the argument \mbox{\texttt{\mdseries\slshape implied}} will cause trouble; if the filter is not intended to be set or reset manually
then perhaps calling \texttt{NewCategory} (\ref{NewCategory}) is more appropriate than calling \texttt{NewFilter}. 

 The optional argument \mbox{\texttt{\mdseries\slshape rank}} denotes the incremental rank (see{\nobreakspace}\ref{Filters}) of the filter, the default value is 1. 

 The default value of the new simple filter for each object is \texttt{false}. }

 

\subsection{\textcolor{Chapter }{DeclareFilter}}
\logpage{[ 13, 8, 2 ]}\nobreak
\hyperdef{L}{X846EA18A7D36626C}{}
{\noindent\textcolor{FuncColor}{$\triangleright$\enspace\texttt{DeclareFilter({\mdseries\slshape name[, implied][, rank]})\index{DeclareFilter@\texttt{DeclareFilter}}
\label{DeclareFilter}
}\hfill{\scriptsize (function)}}\\


 does the same as \texttt{NewFilter} (\ref{NewFilter}) and then binds the result to the global variable \mbox{\texttt{\mdseries\slshape name}}. The variable must previously be writable, and is made
read\texttt{\symbol{45}}only by this function. }

 

\subsection{\textcolor{Chapter }{SetFilterObj}}
\logpage{[ 13, 8, 3 ]}\nobreak
\hyperdef{L}{X7C92D53E7920CE02}{}
{\noindent\textcolor{FuncColor}{$\triangleright$\enspace\texttt{SetFilterObj({\mdseries\slshape obj, filter})\index{SetFilterObj@\texttt{SetFilterObj}}
\label{SetFilterObj}
}\hfill{\scriptsize (function)}}\\


 \texttt{SetFilterObj} sets the value of \mbox{\texttt{\mdseries\slshape filter}} (and of all filters implied by \mbox{\texttt{\mdseries\slshape filter}}) for \mbox{\texttt{\mdseries\slshape obj}} to \texttt{true}. 

 This may trigger immediate methods. }

 

\subsection{\textcolor{Chapter }{ResetFilterObj}}
\logpage{[ 13, 8, 4 ]}\nobreak
\hyperdef{L}{X8117FD03870FB02E}{}
{\noindent\textcolor{FuncColor}{$\triangleright$\enspace\texttt{ResetFilterObj({\mdseries\slshape obj, filter})\index{ResetFilterObj@\texttt{ResetFilterObj}}
\label{ResetFilterObj}
}\hfill{\scriptsize (function)}}\\


 \texttt{ResetFilterObj} sets the value of \mbox{\texttt{\mdseries\slshape filter}} for \mbox{\texttt{\mdseries\slshape obj}} to \texttt{false}. (Implied filters of \mbox{\texttt{\mdseries\slshape filt}} are not touched. This might create inconsistent situations if applied
carelessly). }

   }

  
\section{\textcolor{Chapter }{Types}}\label{Types}
\logpage{[ 13, 9, 0 ]}
\hyperdef{L}{X7E340B8C833BC440}{}
{
  We stated above (see \ref{Types of Objects}) that, for an object \mbox{\texttt{\mdseries\slshape obj}}, its \emph{type} is formed from its family and its filters. There is also a third component,
used in a few situations, namely defining data of the type. 

\subsection{\textcolor{Chapter }{TypeObj}}
\logpage{[ 13, 9, 1 ]}\nobreak
\hyperdef{L}{X7D3E6B6482BE5B16}{}
{\noindent\textcolor{FuncColor}{$\triangleright$\enspace\texttt{TypeObj({\mdseries\slshape obj})\index{TypeObj@\texttt{TypeObj}}
\label{TypeObj}
}\hfill{\scriptsize (function)}}\\


 returns the type of the object \mbox{\texttt{\mdseries\slshape obj}}. 

 The type of an object is itself an object. 

 Two types are equal if and only if the two families are identical, the filters
are equal, and, if present, also the defining data of the types are equal. }

 

\subsection{\textcolor{Chapter }{DataType}}
\logpage{[ 13, 9, 2 ]}\nobreak
\hyperdef{L}{X85A60A7F8083C1C4}{}
{\noindent\textcolor{FuncColor}{$\triangleright$\enspace\texttt{DataType({\mdseries\slshape type})\index{DataType@\texttt{DataType}}
\label{DataType}
}\hfill{\scriptsize (function)}}\\


 The last part of the type, defining data, has not been mentioned before and
seems to be of minor importance. It can be used, e.g., for cosets $U g$ of a group $U$, where the type of each coset may contain the group $U$ as defining data. As a consequence, two such cosets mod $U$ and $V$ can have the same type only if $U = V$. The defining data of the type \mbox{\texttt{\mdseries\slshape type}} can be accessed via \texttt{DataType}. }

 

\subsection{\textcolor{Chapter }{NewType}}
\logpage{[ 13, 9, 3 ]}\nobreak
\hyperdef{L}{X7CE39E9478AEC826}{}
{\noindent\textcolor{FuncColor}{$\triangleright$\enspace\texttt{NewType({\mdseries\slshape family, filter[, data]})\index{NewType@\texttt{NewType}}
\label{NewType}
}\hfill{\scriptsize (function)}}\\


 \texttt{NewType} returns the type given by the family \mbox{\texttt{\mdseries\slshape family}} and the filter \mbox{\texttt{\mdseries\slshape filter}}. The optional third argument \mbox{\texttt{\mdseries\slshape data}} is any object that denotes defining data of the desired type. 

 For examples where \texttt{NewType} is used, see{\nobreakspace}\ref{Component Objects}, \ref{Positional Objects}, and the example in Chapter \ref{An Example -- Residue Class Rings}. }

 }

 }

             
\chapter{\textcolor{Chapter }{Integers}}\label{Integers}
\logpage{[ 14, 0, 0 ]}
\hyperdef{L}{X853DF11B80068ED5}{}
{
  One of the most fundamental datatypes in every programming language is the
integer type. \textsf{GAP} is no exception. 

 \textsf{GAP} integers are entered as a sequence of decimal digits optionally preceded by a ``\texttt{+}'' sign for positive integers or a ``\texttt{\texttt{\symbol{45}}}'' sign for negative integers. The size of integers in \textsf{GAP} is only limited by the amount of available memory, so you can compute with
integers having thousands of digits. 

 
\begin{Verbatim}[commandchars=!@|,fontsize=\small,frame=single,label=Example]
  !gapprompt@gap>| !gapinput@-1234;|
  -1234
  !gapprompt@gap>| !gapinput@123456789012345678901234567890123456789012345678901234567890;|
  123456789012345678901234567890123456789012345678901234567890
\end{Verbatim}
 Note that in a few places, only certain small integer values can be used. A \emph{small integer} (also referred to as \emph{immediate integer}) is an integer $n$ satisfying \texttt{INTOBJ{\textunderscore}MIN} $ \leq n \leq $ \texttt{INTOBJ{\textunderscore}MAX}, where \texttt{INTOBJ{\textunderscore}MIN} and \texttt{INTOBJ{\textunderscore}MAX} equal either $-2^{28}$ and $2^{28}-1$ (on 32\texttt{\symbol{45}}bit systems) or $-2^{60}$ and $2^{60}-1$ (on 64\texttt{\symbol{45}}bit systems). \index{small integer} \index{immediate integer} \index{INTOBJ{\textunderscore}MIN} \index{INTOBJ{\textunderscore}MAX} For example, the elements of a range are restricted to small integers
(see{\nobreakspace}\ref{Ranges}). 

 Many more functions that are mainly related to the prime residue group of
integers modulo an integer are described in chapter{\nobreakspace}\ref{Number Theory}, and functions dealing with combinatorics can be found in
chapter{\nobreakspace}\ref{Combinatorics}.  
\section{\textcolor{Chapter }{Integers: Global Variables}}\label{sect:Integers: Global Variables}
\logpage{[ 14, 1, 0 ]}
\hyperdef{L}{X838230CE810107A3}{}
{
  

\subsection{\textcolor{Chapter }{Integers (global variable)}}
\logpage{[ 14, 1, 1 ]}\nobreak
\hyperdef{L}{X7E20D82B79DE5129}{}
{\noindent\textcolor{FuncColor}{$\triangleright$\enspace\texttt{Integers\index{Integers@\texttt{Integers}!global variable}
\label{Integers:global variable}
}\hfill{\scriptsize (global variable)}}\\
\noindent\textcolor{FuncColor}{$\triangleright$\enspace\texttt{PositiveIntegers\index{PositiveIntegers@\texttt{PositiveIntegers}}
\label{PositiveIntegers}
}\hfill{\scriptsize (global variable)}}\\
\noindent\textcolor{FuncColor}{$\triangleright$\enspace\texttt{NonnegativeIntegers\index{NonnegativeIntegers@\texttt{NonnegativeIntegers}}
\label{NonnegativeIntegers}
}\hfill{\scriptsize (global variable)}}\\


 These global variables represent the ring of integers and the semirings of
positive and nonnegative integers, respectively. 
\begin{Verbatim}[commandchars=!@|,fontsize=\small,frame=single,label=Example]
  !gapprompt@gap>| !gapinput@Size( Integers ); 2 in Integers;|
  infinity
  true
\end{Verbatim}
 

 \texttt{Integers} is a subset of \texttt{Rationals} (\ref{Rationals}), which is a subset of \texttt{Cyclotomics} (\ref{Cyclotomics}). See Chapter{\nobreakspace}\ref{Cyclotomic Numbers} for arithmetic operations and comparison of integers. }

 

\subsection{\textcolor{Chapter }{IsIntegers}}
\logpage{[ 14, 1, 2 ]}\nobreak
\hyperdef{L}{X818683B17F8C97F3}{}
{\noindent\textcolor{FuncColor}{$\triangleright$\enspace\texttt{IsIntegers({\mdseries\slshape obj})\index{IsIntegers@\texttt{IsIntegers}}
\label{IsIntegers}
}\hfill{\scriptsize (Category)}}\\
\noindent\textcolor{FuncColor}{$\triangleright$\enspace\texttt{IsPositiveIntegers({\mdseries\slshape obj})\index{IsPositiveIntegers@\texttt{IsPositiveIntegers}}
\label{IsPositiveIntegers}
}\hfill{\scriptsize (Category)}}\\
\noindent\textcolor{FuncColor}{$\triangleright$\enspace\texttt{IsNonnegativeIntegers({\mdseries\slshape obj})\index{IsNonnegativeIntegers@\texttt{IsNonnegativeIntegers}}
\label{IsNonnegativeIntegers}
}\hfill{\scriptsize (Category)}}\\


 are the defining categories for the domains \texttt{Integers} (\ref{Integers:global variable}), \texttt{PositiveIntegers} (\ref{PositiveIntegers}), and \texttt{NonnegativeIntegers} (\ref{NonnegativeIntegers}). 
\begin{Verbatim}[commandchars=!@|,fontsize=\small,frame=single,label=Example]
  !gapprompt@gap>| !gapinput@IsIntegers( Integers );  IsIntegers( Rationals );  IsIntegers( 7 );|
  true
  false
  false
\end{Verbatim}
 }

 }

  
\section{\textcolor{Chapter }{Elementary Operations for Integers}}\label{Elementary Operations for Integers}
\logpage{[ 14, 2, 0 ]}
\hyperdef{L}{X80CF510B8080C7CA}{}
{
  

\subsection{\textcolor{Chapter }{IsInt}}
\logpage{[ 14, 2, 1 ]}\nobreak
\hyperdef{L}{X87AEADF07DC8303B}{}
{\noindent\textcolor{FuncColor}{$\triangleright$\enspace\texttt{IsInt({\mdseries\slshape obj})\index{IsInt@\texttt{IsInt}}
\label{IsInt}
}\hfill{\scriptsize (Category)}}\\


 Every rational integer lies in the category \texttt{IsInt}, which is a subcategory of \texttt{IsRat} (\ref{IsRat}). }

 

\subsection{\textcolor{Chapter }{IsPosInt}}
\logpage{[ 14, 2, 2 ]}\nobreak
\hyperdef{L}{X82A854757DFA9C76}{}
{\noindent\textcolor{FuncColor}{$\triangleright$\enspace\texttt{IsPosInt({\mdseries\slshape obj})\index{IsPosInt@\texttt{IsPosInt}}
\label{IsPosInt}
}\hfill{\scriptsize (Category)}}\\


 Every positive integer lies in the category \texttt{IsPosInt}. }

 

\subsection{\textcolor{Chapter }{Int}}
\logpage{[ 14, 2, 3 ]}\nobreak
\hyperdef{L}{X87CA734380B5F68C}{}
{\noindent\textcolor{FuncColor}{$\triangleright$\enspace\texttt{Int({\mdseries\slshape elm})\index{Int@\texttt{Int}}
\label{Int}
}\hfill{\scriptsize (attribute)}}\\


 \texttt{Int} returns an integer \texttt{int} whose meaning depends on the type of \mbox{\texttt{\mdseries\slshape elm}}. For example: 

 If \mbox{\texttt{\mdseries\slshape elm}} is a rational number (see Chapter{\nobreakspace}\ref{Rational Numbers}) then \texttt{int} is the integer part of the quotient of numerator and denominator of \mbox{\texttt{\mdseries\slshape elm}} (see{\nobreakspace}\texttt{QuoInt} (\ref{QuoInt})). 

 If \mbox{\texttt{\mdseries\slshape elm}} is an element of a finite prime field (see Chapter{\nobreakspace}\ref{Finite Fields}) then \texttt{int} is the smallest nonnegative integer such that \texttt{\mbox{\texttt{\mdseries\slshape elm}} = int * One( \mbox{\texttt{\mdseries\slshape elm}} )}. 

 If \mbox{\texttt{\mdseries\slshape elm}} is a string (see Chapter{\nobreakspace}\ref{Strings and Characters}) consisting entirely of decimal digits \texttt{'0'}, \texttt{'1'}, $\ldots$, \texttt{'9'}, and optionally a sign \texttt{'\texttt{\symbol{45}}'} (at the first position), then \texttt{int} is the integer described by this string. For all other strings, \texttt{fail} is returned. See \texttt{Int} (\ref{Int:for strings}). 

 The operation \texttt{String} (\ref{String}) can be used to compute a string for rational integers, in fact for all
cyclotomics. 

 
\begin{Verbatim}[commandchars=!@|,fontsize=\small,frame=single,label=Example]
  !gapprompt@gap>| !gapinput@Int( 4/3 );  Int( -2/3 );|
  1
  0
  !gapprompt@gap>| !gapinput@int:= Int( Z(5) );  int * One( Z(5) );|
  2
  Z(5)
  !gapprompt@gap>| !gapinput@Int( "12345" );  Int( "-27" );  Int( "-27/3" );|
  12345
  -27
  fail
\end{Verbatim}
 }

 

\subsection{\textcolor{Chapter }{IsEvenInt}}
\logpage{[ 14, 2, 4 ]}\nobreak
\hyperdef{L}{X87DD1EEE7EF18036}{}
{\noindent\textcolor{FuncColor}{$\triangleright$\enspace\texttt{IsEvenInt({\mdseries\slshape n})\index{IsEvenInt@\texttt{IsEvenInt}}
\label{IsEvenInt}
}\hfill{\scriptsize (function)}}\\


 tests if the integer \mbox{\texttt{\mdseries\slshape n}} is divisible by 2. }

 

\subsection{\textcolor{Chapter }{IsOddInt}}
\logpage{[ 14, 2, 5 ]}\nobreak
\hyperdef{L}{X8621BA927CD12EFB}{}
{\noindent\textcolor{FuncColor}{$\triangleright$\enspace\texttt{IsOddInt({\mdseries\slshape n})\index{IsOddInt@\texttt{IsOddInt}}
\label{IsOddInt}
}\hfill{\scriptsize (function)}}\\


 tests if the integer \mbox{\texttt{\mdseries\slshape n}} is not divisible by 2. }

 

\subsection{\textcolor{Chapter }{AbsInt}}
\logpage{[ 14, 2, 6 ]}\nobreak
\hyperdef{L}{X782095927FB9F1DB}{}
{\noindent\textcolor{FuncColor}{$\triangleright$\enspace\texttt{AbsInt({\mdseries\slshape n})\index{AbsInt@\texttt{AbsInt}}
\label{AbsInt}
}\hfill{\scriptsize (function)}}\\


 \index{absolute value of an integer} \texttt{AbsInt} returns the absolute value of the integer \mbox{\texttt{\mdseries\slshape n}}, i.e., \mbox{\texttt{\mdseries\slshape n}} if \mbox{\texttt{\mdseries\slshape n}} is positive, \texttt{\symbol{45}}\mbox{\texttt{\mdseries\slshape n}} if \mbox{\texttt{\mdseries\slshape n}} is negative and 0 if \mbox{\texttt{\mdseries\slshape n}} is 0. 

 \texttt{AbsInt} is a special case of the general operation \texttt{EuclideanDegree} (\ref{EuclideanDegree}). 

 See also \texttt{AbsoluteValue} (\ref{AbsoluteValue}). 
\begin{Verbatim}[commandchars=!@|,fontsize=\small,frame=single,label=Example]
  !gapprompt@gap>| !gapinput@AbsInt( 33 );|
  33
  !gapprompt@gap>| !gapinput@AbsInt( -214378 );|
  214378
  !gapprompt@gap>| !gapinput@AbsInt( 0 );|
  0
\end{Verbatim}
 }

 

\subsection{\textcolor{Chapter }{SignInt}}
\logpage{[ 14, 2, 7 ]}\nobreak
\hyperdef{L}{X842614817FE48D62}{}
{\noindent\textcolor{FuncColor}{$\triangleright$\enspace\texttt{SignInt({\mdseries\slshape n})\index{SignInt@\texttt{SignInt}}
\label{SignInt}
}\hfill{\scriptsize (function)}}\\


 \index{sign!of an integer} \texttt{SignInt} returns the sign of the integer \mbox{\texttt{\mdseries\slshape n}}, i.e., 1 if \mbox{\texttt{\mdseries\slshape n}} is positive, \texttt{\symbol{45}}1 if \mbox{\texttt{\mdseries\slshape n}} is negative and 0 if \mbox{\texttt{\mdseries\slshape n}} is 0. 
\begin{Verbatim}[commandchars=!@|,fontsize=\small,frame=single,label=Example]
  !gapprompt@gap>| !gapinput@SignInt( 33 );|
  1
  !gapprompt@gap>| !gapinput@SignInt( -214378 );|
  -1
  !gapprompt@gap>| !gapinput@SignInt( 0 );|
  0
\end{Verbatim}
 }

 

\subsection{\textcolor{Chapter }{LogInt}}
\logpage{[ 14, 2, 8 ]}\nobreak
\hyperdef{L}{X8197C4E882BAF14E}{}
{\noindent\textcolor{FuncColor}{$\triangleright$\enspace\texttt{LogInt({\mdseries\slshape n, base})\index{LogInt@\texttt{LogInt}}
\label{LogInt}
}\hfill{\scriptsize (function)}}\\


 \texttt{LogInt} returns the integer part of the logarithm of the positive integer \mbox{\texttt{\mdseries\slshape n}} with respect to the positive integer \mbox{\texttt{\mdseries\slshape base}}, i.e., the largest positive integer $e$ such that $\mbox{\texttt{\mdseries\slshape base}}^e \leq \mbox{\texttt{\mdseries\slshape n}}$. The function \texttt{LogInt} will signal an error if either \mbox{\texttt{\mdseries\slshape n}} or \mbox{\texttt{\mdseries\slshape base}} is not positive. 

 For \mbox{\texttt{\mdseries\slshape base}} $= 2$ this is very efficient because the internal binary representation of the
integer is used. 

 
\begin{Verbatim}[commandchars=!@|,fontsize=\small,frame=single,label=Example]
  !gapprompt@gap>| !gapinput@LogInt( 1030, 2 );|
  10
  !gapprompt@gap>| !gapinput@2^10;|
  1024
  !gapprompt@gap>| !gapinput@LogInt( 1, 10 );|
  0
\end{Verbatim}
 }

 

\subsection{\textcolor{Chapter }{RootInt}}
\logpage{[ 14, 2, 9 ]}\nobreak
\hyperdef{L}{X83D9B5C87EEA2A77}{}
{\noindent\textcolor{FuncColor}{$\triangleright$\enspace\texttt{RootInt({\mdseries\slshape n[, k]})\index{RootInt@\texttt{RootInt}}
\label{RootInt}
}\hfill{\scriptsize (function)}}\\


 \index{root!of an integer} \index{square root!of an integer} \texttt{RootInt} returns the integer part of the \mbox{\texttt{\mdseries\slshape k}}th root of the integer \mbox{\texttt{\mdseries\slshape n}}. If the optional integer argument \mbox{\texttt{\mdseries\slshape k}} is not given it defaults to 2, i.e., \texttt{RootInt} returns the integer part of the square root in this case. 

 If \mbox{\texttt{\mdseries\slshape n}} is positive, \texttt{RootInt} returns the largest positive integer $r$ such that $r^{\mbox{\texttt{\mdseries\slshape k}}} \leq \mbox{\texttt{\mdseries\slshape n}}$. If \mbox{\texttt{\mdseries\slshape n}} is negative and \mbox{\texttt{\mdseries\slshape k}} is odd \texttt{RootInt} returns \texttt{\texttt{\symbol{45}}RootInt( \texttt{\symbol{45}}\mbox{\texttt{\mdseries\slshape n}}, \mbox{\texttt{\mdseries\slshape k}} )}. If \mbox{\texttt{\mdseries\slshape n}} is negative and \mbox{\texttt{\mdseries\slshape k}} is even \texttt{RootInt} will cause an error. \texttt{RootInt} will also cause an error if \mbox{\texttt{\mdseries\slshape k}} is 0 or negative. 
\begin{Verbatim}[commandchars=!@|,fontsize=\small,frame=single,label=Example]
  !gapprompt@gap>| !gapinput@RootInt( 361 );|
  19
  !gapprompt@gap>| !gapinput@RootInt( 2 * 10^12 );|
  1414213
  !gapprompt@gap>| !gapinput@RootInt( 17000, 5 );|
  7
  !gapprompt@gap>| !gapinput@7^5;|
  16807
\end{Verbatim}
 }

 

\subsection{\textcolor{Chapter }{SmallestRootInt}}
\logpage{[ 14, 2, 10 ]}\nobreak
\hyperdef{L}{X7F98A0CE7B9FD366}{}
{\noindent\textcolor{FuncColor}{$\triangleright$\enspace\texttt{SmallestRootInt({\mdseries\slshape n})\index{SmallestRootInt@\texttt{SmallestRootInt}}
\label{SmallestRootInt}
}\hfill{\scriptsize (function)}}\\


 \index{root!of an integer, smallest} \texttt{SmallestRootInt} returns the smallest root of the integer \mbox{\texttt{\mdseries\slshape n}}. 

 The smallest root of an integer \mbox{\texttt{\mdseries\slshape n}} is the integer $r$ of smallest absolute value for which a positive integer $k$ exists such that $\mbox{\texttt{\mdseries\slshape n}} = r^k$. 
\begin{Verbatim}[commandchars=!@|,fontsize=\small,frame=single,label=Example]
  !gapprompt@gap>| !gapinput@SmallestRootInt( 2^30 );|
  2
  !gapprompt@gap>| !gapinput@SmallestRootInt( -(2^30) );|
  -4
\end{Verbatim}
 

 Note that $(-2)^{30} = +(2^{30})$. 

 
\begin{Verbatim}[commandchars=!@|,fontsize=\small,frame=single,label=Example]
  !gapprompt@gap>| !gapinput@SmallestRootInt( 279936 );|
  6
  !gapprompt@gap>| !gapinput@LogInt( 279936, 6 );|
  7
  !gapprompt@gap>| !gapinput@SmallestRootInt( 1001 );|
  1001
\end{Verbatim}
 }

 

\subsection{\textcolor{Chapter }{IsSquareInt}}
\logpage{[ 14, 2, 11 ]}\nobreak
\hyperdef{L}{X83B998E486893FED}{}
{\noindent\textcolor{FuncColor}{$\triangleright$\enspace\texttt{IsSquareInt({\mdseries\slshape n})\index{IsSquareInt@\texttt{IsSquareInt}}
\label{IsSquareInt}
}\hfill{\scriptsize (function)}}\\


 \texttt{IsSquareInt} tests whether the integer \mbox{\texttt{\mdseries\slshape n}} is the square of an integer or not. This test is much faster than the simpler \texttt{RootInt}$(n)^2=n$ because it first tests whether \mbox{\texttt{\mdseries\slshape n}} is a square residue modulo some small integers. }

 

\subsection{\textcolor{Chapter }{ListOfDigits}}
\logpage{[ 14, 2, 12 ]}\nobreak
\hyperdef{L}{X862D1BD786EFFDA9}{}
{\noindent\textcolor{FuncColor}{$\triangleright$\enspace\texttt{ListOfDigits({\mdseries\slshape n})\index{ListOfDigits@\texttt{ListOfDigits}}
\label{ListOfDigits}
}\hfill{\scriptsize (function)}}\\


 For a positive integer \mbox{\texttt{\mdseries\slshape n}} this function returns a list \mbox{\texttt{\mdseries\slshape l}}, consisting of the digits of \mbox{\texttt{\mdseries\slshape n}} in decimal representation. 
\begin{Verbatim}[commandchars=!@|,fontsize=\small,frame=single,label=Example]
  !gapprompt@gap>| !gapinput@ListOfDigits(3142);|
  [ 3, 1, 4, 2 ]
\end{Verbatim}
 }

 

\subsection{\textcolor{Chapter }{Random (for integers)}}
\logpage{[ 14, 2, 13 ]}\nobreak
\hyperdef{L}{X8185784B7E228DEA}{}
{\noindent\textcolor{FuncColor}{$\triangleright$\enspace\texttt{Random({\mdseries\slshape Integers})\index{Random@\texttt{Random}!for integers}
\label{Random:for integers}
}\hfill{\scriptsize (method)}}\\


 \texttt{Random} for integers returns pseudo random integers between $-10$ and $10$ distributed according to a binomial distribution. To generate uniformly
distributed integers from a range, use the construction \texttt{Random( [ \mbox{\texttt{\mdseries\slshape low}} .. \mbox{\texttt{\mdseries\slshape high}} ] )} {\nobreakspace}(see{\nobreakspace}\texttt{Random} (\ref{Random:for lower and upper bound})). }

 }

  
\section{\textcolor{Chapter }{Quotients and Remainders}}\label{Quotients and Remainders}
\logpage{[ 14, 3, 0 ]}
\hyperdef{L}{X7A9FD25D81D88D1B}{}
{
  

\subsection{\textcolor{Chapter }{QuoInt}}
\logpage{[ 14, 3, 1 ]}\nobreak
\hyperdef{L}{X849D0F807F697D35}{}
{\noindent\textcolor{FuncColor}{$\triangleright$\enspace\texttt{QuoInt({\mdseries\slshape n, m})\index{QuoInt@\texttt{QuoInt}}
\label{QuoInt}
}\hfill{\scriptsize (function)}}\\


 \index{integer part of a quotient} \texttt{QuoInt} returns the integer part of the quotient of its integer operands. 

 If \mbox{\texttt{\mdseries\slshape n}} and \mbox{\texttt{\mdseries\slshape m}} are positive, \texttt{QuoInt} returns the largest positive integer $q$ such that $q * \mbox{\texttt{\mdseries\slshape m}} \leq \mbox{\texttt{\mdseries\slshape n}}$. If \mbox{\texttt{\mdseries\slshape n}} or \mbox{\texttt{\mdseries\slshape m}} or both are negative the absolute value of the integer part of the quotient is
the quotient of the absolute values of \mbox{\texttt{\mdseries\slshape n}} and \mbox{\texttt{\mdseries\slshape m}}, and the sign of it is the product of the signs of \mbox{\texttt{\mdseries\slshape n}} and \mbox{\texttt{\mdseries\slshape m}}. 

 \texttt{QuoInt} is used in a method for the general operation \texttt{EuclideanQuotient} (\ref{EuclideanQuotient}). 
\begin{Verbatim}[commandchars=!@|,fontsize=\small,frame=single,label=Example]
  !gapprompt@gap>| !gapinput@QuoInt(5,3);  QuoInt(-5,3);  QuoInt(5,-3);  QuoInt(-5,-3);|
  1
  -1
  -1
  1
\end{Verbatim}
 }

 

\subsection{\textcolor{Chapter }{BestQuoInt}}
\logpage{[ 14, 3, 2 ]}\nobreak
\hyperdef{L}{X795170A385AC8FEE}{}
{\noindent\textcolor{FuncColor}{$\triangleright$\enspace\texttt{BestQuoInt({\mdseries\slshape n, m})\index{BestQuoInt@\texttt{BestQuoInt}}
\label{BestQuoInt}
}\hfill{\scriptsize (function)}}\\


 \texttt{BestQuoInt} returns the best quotient $q$ of the integers \mbox{\texttt{\mdseries\slshape n}} and \mbox{\texttt{\mdseries\slshape m}}. This is the quotient such that $\mbox{\texttt{\mdseries\slshape n}}-q*\mbox{\texttt{\mdseries\slshape m}}$ has minimal absolute value. If there are two quotients whose remainders have
the same absolute value, then the quotient with the smaller absolute value is
chosen. 
\begin{Verbatim}[commandchars=!@|,fontsize=\small,frame=single,label=Example]
  !gapprompt@gap>| !gapinput@BestQuoInt( 5, 3 );  BestQuoInt( -5, 3 );|
  2
  -2
\end{Verbatim}
 }

 

\subsection{\textcolor{Chapter }{RemInt}}
\logpage{[ 14, 3, 3 ]}\nobreak
\hyperdef{L}{X805ADD5A826D844D}{}
{\noindent\textcolor{FuncColor}{$\triangleright$\enspace\texttt{RemInt({\mdseries\slshape n, m})\index{RemInt@\texttt{RemInt}}
\label{RemInt}
}\hfill{\scriptsize (function)}}\\


 \index{remainder of a quotient} \texttt{RemInt} returns the remainder of its two integer operands. 

 If \mbox{\texttt{\mdseries\slshape m}} is not equal to zero, \texttt{RemInt} returns \texttt{\mbox{\texttt{\mdseries\slshape n}} \texttt{\symbol{45}} \mbox{\texttt{\mdseries\slshape m}} * QuoInt( \mbox{\texttt{\mdseries\slshape n}}, \mbox{\texttt{\mdseries\slshape m}} )}. Note that the rules given for \texttt{QuoInt} (\ref{QuoInt}) imply that the return value of \texttt{RemInt} has the same sign as \mbox{\texttt{\mdseries\slshape n}} and its absolute value is strictly less than the absolute value of \mbox{\texttt{\mdseries\slshape m}}. Note also that the return value equals \texttt{\mbox{\texttt{\mdseries\slshape n}} mod \mbox{\texttt{\mdseries\slshape m}}} when both \mbox{\texttt{\mdseries\slshape n}} and \mbox{\texttt{\mdseries\slshape m}} are nonnegative. Dividing by \texttt{0} signals an error. 

 \texttt{RemInt} is used in a method for the general operation \texttt{EuclideanRemainder} (\ref{EuclideanRemainder}). 
\begin{Verbatim}[commandchars=!@|,fontsize=\small,frame=single,label=Example]
  !gapprompt@gap>| !gapinput@RemInt(5,3);  RemInt(-5,3);  RemInt(5,-3);  RemInt(-5,-3);|
  2
  -2
  2
  -2
\end{Verbatim}
 }

 

\subsection{\textcolor{Chapter }{GcdInt}}
\logpage{[ 14, 3, 4 ]}\nobreak
\hyperdef{L}{X7A4FEFCA8128E3C3}{}
{\noindent\textcolor{FuncColor}{$\triangleright$\enspace\texttt{GcdInt({\mdseries\slshape m, n})\index{GcdInt@\texttt{GcdInt}}
\label{GcdInt}
}\hfill{\scriptsize (function)}}\\


 \texttt{GcdInt} returns the greatest common divisor of its two integer operands \mbox{\texttt{\mdseries\slshape m}} and \mbox{\texttt{\mdseries\slshape n}}, i.e., the greatest integer that divides both \mbox{\texttt{\mdseries\slshape m}} and \mbox{\texttt{\mdseries\slshape n}}. The greatest common divisor is never negative, even if the arguments are. We
define \texttt{GcdInt( \mbox{\texttt{\mdseries\slshape m}}, 0 ) = GcdInt( 0, \mbox{\texttt{\mdseries\slshape m}} ) = AbsInt( \mbox{\texttt{\mdseries\slshape m}} )} and \texttt{GcdInt( 0, 0 ) = 0}. 

 \texttt{GcdInt} is a method used by the general function \texttt{Gcd} (\ref{Gcd:for (a ring and) several elements}). 

 
\begin{Verbatim}[commandchars=!@|,fontsize=\small,frame=single,label=Example]
  !gapprompt@gap>| !gapinput@GcdInt( 123, 66 );|
  3
\end{Verbatim}
 }

 

\subsection{\textcolor{Chapter }{Gcdex}}
\logpage{[ 14, 3, 5 ]}\nobreak
\hyperdef{L}{X8775930486BD0C5B}{}
{\noindent\textcolor{FuncColor}{$\triangleright$\enspace\texttt{Gcdex({\mdseries\slshape m, n})\index{Gcdex@\texttt{Gcdex}}
\label{Gcdex}
}\hfill{\scriptsize (function)}}\\


 returns a record \texttt{g} describing the extended greatest common divisor of \mbox{\texttt{\mdseries\slshape m}} and \mbox{\texttt{\mdseries\slshape n}}. The component \texttt{gcd} is this gcd, the components \texttt{coeff1} and \texttt{coeff2} are integer cofactors such that \texttt{g.gcd = g.coeff1 * \mbox{\texttt{\mdseries\slshape m}} + g.coeff2 * \mbox{\texttt{\mdseries\slshape n}}}, and the components \texttt{coeff3} and \texttt{coeff4} are integer cofactors such that \texttt{0 = g.coeff3 * \mbox{\texttt{\mdseries\slshape m}} + g.coeff4 * \mbox{\texttt{\mdseries\slshape n}}}. 

 If \mbox{\texttt{\mdseries\slshape m}} and \mbox{\texttt{\mdseries\slshape n}} both are nonzero, \texttt{AbsInt( g.coeff1 )} is less than or equal to \texttt{AbsInt(\mbox{\texttt{\mdseries\slshape n}}) / (2 * g.gcd)}, and \texttt{AbsInt( g.coeff2 )} is less than or equal to \texttt{AbsInt(\mbox{\texttt{\mdseries\slshape m}}) / (2 * g.gcd)}. 

 If \mbox{\texttt{\mdseries\slshape m}} or \mbox{\texttt{\mdseries\slshape n}} or both are zero \texttt{coeff3} is \texttt{\texttt{\symbol{45}}\mbox{\texttt{\mdseries\slshape n}} / g.gcd} and \texttt{coeff4} is \texttt{\mbox{\texttt{\mdseries\slshape m}} / g.gcd}. 

 The coefficients always form a unimodular matrix, i.e., the determinant \texttt{g.coeff1 * g.coeff4 \texttt{\symbol{45}} g.coeff3 * g.coeff2} is $1$ or $-1$. 
\begin{Verbatim}[commandchars=!@|,fontsize=\small,frame=single,label=Example]
  !gapprompt@gap>| !gapinput@Gcdex( 123, 66 );|
  rec( coeff1 := 7, coeff2 := -13, coeff3 := -22, coeff4 := 41,
    gcd := 3 )
\end{Verbatim}
 This means $3 = 7 * 123 - 13 * 66$, $0 = -22 * 123 + 41 * 66$. 
\begin{Verbatim}[commandchars=!@|,fontsize=\small,frame=single,label=Example]
  !gapprompt@gap>| !gapinput@Gcdex( 0, -3 );|
  rec( coeff1 := 0, coeff2 := -1, coeff3 := 1, coeff4 := 0, gcd := 3 )
  !gapprompt@gap>| !gapinput@Gcdex( 0, 0 );|
  rec( coeff1 := 1, coeff2 := 0, coeff3 := 0, coeff4 := 1, gcd := 0 )
\end{Verbatim}
 

 \texttt{GcdRepresentation} (\ref{GcdRepresentation:for (a ring and) several elements}) provides similar functionality over arbitrary Euclidean rings. }

 

\subsection{\textcolor{Chapter }{LcmInt}}
\logpage{[ 14, 3, 6 ]}\nobreak
\hyperdef{L}{X7B33143E78A8DDE3}{}
{\noindent\textcolor{FuncColor}{$\triangleright$\enspace\texttt{LcmInt({\mdseries\slshape m, n})\index{LcmInt@\texttt{LcmInt}}
\label{LcmInt}
}\hfill{\scriptsize (function)}}\\


 returns the least common multiple of the integers \mbox{\texttt{\mdseries\slshape m}} and \mbox{\texttt{\mdseries\slshape n}}. 

 \texttt{LcmInt} is a method used by the general operation \texttt{Lcm} (\ref{Lcm:for (a ring and) several elements}). 
\begin{Verbatim}[commandchars=!@|,fontsize=\small,frame=single,label=Example]
  !gapprompt@gap>| !gapinput@LcmInt( 123, 66 );|
  2706
\end{Verbatim}
 }

 

\subsection{\textcolor{Chapter }{CoefficientsQadic}}
\logpage{[ 14, 3, 7 ]}\nobreak
\hyperdef{L}{X79B466E984CD52D4}{}
{\noindent\textcolor{FuncColor}{$\triangleright$\enspace\texttt{CoefficientsQadic({\mdseries\slshape i, q})\index{CoefficientsQadic@\texttt{CoefficientsQadic}}
\label{CoefficientsQadic}
}\hfill{\scriptsize (operation)}}\\


 returns the \mbox{\texttt{\mdseries\slshape q}}\texttt{\symbol{45}}adic representation of the integer \mbox{\texttt{\mdseries\slshape i}} as a list $l$ of coefficients satisfying the equality $\mbox{\texttt{\mdseries\slshape i}} = \sum_{{j = 0}} \mbox{\texttt{\mdseries\slshape q}}^j \cdot l[j+1]$ for an integer $\mbox{\texttt{\mdseries\slshape q}} > 1$. 
\begin{Verbatim}[commandchars=!@|,fontsize=\small,frame=single,label=Example]
  !gapprompt@gap>| !gapinput@l:=CoefficientsQadic(462,3);|
  [ 0, 1, 0, 2, 2, 1 ]
\end{Verbatim}
 }

 

\subsection{\textcolor{Chapter }{CoefficientsMultiadic}}
\logpage{[ 14, 3, 8 ]}\nobreak
\hyperdef{L}{X83124F86839DC7E6}{}
{\noindent\textcolor{FuncColor}{$\triangleright$\enspace\texttt{CoefficientsMultiadic({\mdseries\slshape ints, int})\index{CoefficientsMultiadic@\texttt{CoefficientsMultiadic}}
\label{CoefficientsMultiadic}
}\hfill{\scriptsize (function)}}\\


 returns the multiadic expansion of the integer \mbox{\texttt{\mdseries\slshape int}} modulo the integers given in \mbox{\texttt{\mdseries\slshape ints}} (in ascending order). It returns a list of coefficients in the \emph{reverse} order to that in \mbox{\texttt{\mdseries\slshape ints}}.  }

 

\subsection{\textcolor{Chapter }{ChineseRem}}
\logpage{[ 14, 3, 9 ]}\nobreak
\hyperdef{L}{X84A1900E82902B5F}{}
{\noindent\textcolor{FuncColor}{$\triangleright$\enspace\texttt{ChineseRem({\mdseries\slshape moduli, residues})\index{ChineseRem@\texttt{ChineseRem}}
\label{ChineseRem}
}\hfill{\scriptsize (function)}}\\


 \index{Chinese remainder} \texttt{ChineseRem} returns the combination of the \mbox{\texttt{\mdseries\slshape residues}} modulo the \mbox{\texttt{\mdseries\slshape moduli}}, i.e., the unique integer \texttt{c} from \texttt{[0..Lcm(\mbox{\texttt{\mdseries\slshape moduli}})\texttt{\symbol{45}}1]} such that \texttt{c = \mbox{\texttt{\mdseries\slshape residues}}[i]} modulo \texttt{\mbox{\texttt{\mdseries\slshape moduli}}[i]} for all \texttt{i}, if it exists. If no such combination exists \texttt{ChineseRem} signals an error. 

 Such a combination does exist if and only if \texttt{\mbox{\texttt{\mdseries\slshape residues}}[i] = \mbox{\texttt{\mdseries\slshape residues}}[k] mod Gcd( \mbox{\texttt{\mdseries\slshape moduli}}[i], \mbox{\texttt{\mdseries\slshape moduli}}[k] )} for every pair \texttt{i}, \texttt{k}. Note that this implies that such a combination exists if the moduli are
pairwise relatively prime. This is called the Chinese remainder theorem. 
\begin{Verbatim}[commandchars=!@|,fontsize=\small,frame=single,label=Example]
  !gapprompt@gap>| !gapinput@ChineseRem( [ 2, 3, 5, 7 ], [ 1, 2, 3, 4 ] );|
  53
  !gapprompt@gap>| !gapinput@ChineseRem( [ 6, 10, 14 ], [ 1, 3, 5 ] );|
  103
\end{Verbatim}
 
\begin{Verbatim}[commandchars=!|B,fontsize=\small,frame=single,label=Example]
  !gapprompt|gap>B !gapinput|ChineseRem( [ 6, 10, 14 ], [ 1, 2, 3 ] );B
  Error, the residues must be equal modulo 2
  Stack trace:
  *[1] Error( "the residues must be equal modulo ", g.gcd );
     @ GAPROOT/lib/integer.gi:391
  <function "ChineseRem">( <arguments> )
   called from read-eval loop at *stdin*:2
  you can enter 'quit;' to quit to outer loop, or
  you can enter 'return;' to continue
  !gapbrkprompt|brk>B !gapinput|B
\end{Verbatim}
 }

 

\subsection{\textcolor{Chapter }{PowerModInt}}
\logpage{[ 14, 3, 10 ]}\nobreak
\hyperdef{L}{X7E404B1183DBC82A}{}
{\noindent\textcolor{FuncColor}{$\triangleright$\enspace\texttt{PowerModInt({\mdseries\slshape r, e, m})\index{PowerModInt@\texttt{PowerModInt}}
\label{PowerModInt}
}\hfill{\scriptsize (function)}}\\


 returns $\mbox{\texttt{\mdseries\slshape r}}^{\mbox{\texttt{\mdseries\slshape e}}} \pmod{\mbox{\texttt{\mdseries\slshape m}}}$ for integers \mbox{\texttt{\mdseries\slshape r}}, \mbox{\texttt{\mdseries\slshape e}} and \mbox{\texttt{\mdseries\slshape m}}. 

 Note that \texttt{PowerModInt} can reduce intermediate results and thus will generally be faster than using \mbox{\texttt{\mdseries\slshape r}}\texttt{\texttt{\symbol{94}}}\mbox{\texttt{\mdseries\slshape e}}\texttt{ mod }\mbox{\texttt{\mdseries\slshape m}}, which would compute $\mbox{\texttt{\mdseries\slshape r}}^{\mbox{\texttt{\mdseries\slshape e}}}$ first and reduces the result afterwards. 

 \texttt{PowerModInt} is a method for the general operation \texttt{PowerMod} (\ref{PowerMod}). }

 }

  
\section{\textcolor{Chapter }{Prime Integers and Factorization}}\label{Prime Integers and Factorization}
\logpage{[ 14, 4, 0 ]}
\hyperdef{L}{X82005E587F0CB02A}{}
{
  

\subsection{\textcolor{Chapter }{Primes}}
\logpage{[ 14, 4, 1 ]}\nobreak
\hyperdef{L}{X86F5E4CD82FEB9F4}{}
{\noindent\textcolor{FuncColor}{$\triangleright$\enspace\texttt{Primes\index{Primes@\texttt{Primes}}
\label{Primes}
}\hfill{\scriptsize (global variable)}}\\


 \texttt{Primes} is a strictly sorted list of the 168 primes less than 1000. 

 This is used in \texttt{IsPrimeInt} (\ref{IsPrimeInt}) and \texttt{FactorsInt} (\ref{FactorsInt}) to cast out small primes quickly. 
\begin{Verbatim}[commandchars=!@|,fontsize=\small,frame=single,label=Example]
  !gapprompt@gap>| !gapinput@Primes[1];|
  2
  !gapprompt@gap>| !gapinput@Primes[100];|
  541
\end{Verbatim}
 }

 

\subsection{\textcolor{Chapter }{IsPrimeInt}}
\logpage{[ 14, 4, 2 ]}\nobreak
\hyperdef{L}{X78FDA4437EDCA70C}{}
{\noindent\textcolor{FuncColor}{$\triangleright$\enspace\texttt{IsPrimeInt({\mdseries\slshape n})\index{IsPrimeInt@\texttt{IsPrimeInt}}
\label{IsPrimeInt}
}\hfill{\scriptsize (function)}}\\
\noindent\textcolor{FuncColor}{$\triangleright$\enspace\texttt{IsProbablyPrimeInt({\mdseries\slshape n})\index{IsProbablyPrimeInt@\texttt{IsProbablyPrimeInt}}
\label{IsProbablyPrimeInt}
}\hfill{\scriptsize (function)}}\\


 \texttt{IsPrimeInt} returns \texttt{false} if it can prove that the integer \mbox{\texttt{\mdseries\slshape n}} is composite and \texttt{true} otherwise. By convention \texttt{IsPrimeInt(0) = IsPrimeInt(1) = false} and we define \texttt{IsPrimeInt(\texttt{\symbol{45}}}\mbox{\texttt{\mdseries\slshape n}}\texttt{) = IsPrimeInt(}\mbox{\texttt{\mdseries\slshape n}}\texttt{)}. 

 \texttt{IsPrimeInt} will return \texttt{true} for every prime \mbox{\texttt{\mdseries\slshape n}}. \texttt{IsPrimeInt} will return \texttt{false} for all composite \mbox{\texttt{\mdseries\slshape n}} $< 10^{18}$ and for all composite \mbox{\texttt{\mdseries\slshape n}} that have a factor $p < 1000$. So for integers \mbox{\texttt{\mdseries\slshape n}} $< 10^{18}$, \texttt{IsPrimeInt} is a proper primality test. It is conceivable that \texttt{IsPrimeInt} may return \texttt{true} for some composite \mbox{\texttt{\mdseries\slshape n}} $> 10^{18}$, but no such \mbox{\texttt{\mdseries\slshape n}} is currently known. So for integers \mbox{\texttt{\mdseries\slshape n}} $> 10^{18}$, \texttt{IsPrimeInt} is a probable\texttt{\symbol{45}}primality test. \texttt{IsPrimeInt} will issue a warning when its argument is probably prime but not a proven
prime. (The function \texttt{IsProbablyPrimeInt} will do a similar calculation but not issue a warning.) The warning can be
switched off by \texttt{SetInfoLevel( InfoPrimeInt, 0 );}, the default level is $1$ (also see \texttt{SetInfoLevel} (\ref{SetInfoLevel}) ). 

 If composites that fool \texttt{IsPrimeInt} do exist, they would be extremely rare, and finding one by pure chance might
be less likely than finding a bug in \textsf{GAP}. We would appreciate being informed about any example of a composite number \mbox{\texttt{\mdseries\slshape n}} for which \texttt{IsPrimeInt} returns \texttt{true}. 

 \texttt{IsPrimeInt} is a deterministic algorithm, i.e., the computations involve no random
numbers, and repeated calls will always return the same result. \texttt{IsPrimeInt} first does trial divisions by the primes less than 1000. Then it tests that \mbox{\texttt{\mdseries\slshape n}} is a strong pseudoprime w.r.t. the base 2. Finally it tests whether \mbox{\texttt{\mdseries\slshape n}} is a Lucas pseudoprime w.r.t. the smallest quadratic nonresidue of \mbox{\texttt{\mdseries\slshape n}}. A better description can be found in the comment in the library file \texttt{primality.gi}. 

 The time taken by \texttt{IsPrimeInt} is approximately proportional to the third power of the number of digits of \mbox{\texttt{\mdseries\slshape n}}. Testing numbers with several hundreds digits is quite feasible. 

 \texttt{IsPrimeInt} is a method for the general operation \texttt{IsPrime} (\ref{IsPrime}). 

 Remark: In future versions of \textsf{GAP} we hope to change the definition of \texttt{IsPrimeInt} to return \texttt{true} only for proven primes (currently, we lack a sufficiently good primality
proving function). In applications, use explicitly \texttt{IsPrimeInt} or \texttt{IsProbablyPrimeInt} with this change in mind. 
\begin{Verbatim}[commandchars=!@|,fontsize=\small,frame=single,label=Example]
  !gapprompt@gap>| !gapinput@IsPrimeInt( 2^31 - 1 );|
  true
  !gapprompt@gap>| !gapinput@IsPrimeInt( 10^42 + 1 );|
  false
\end{Verbatim}
 }

 

\subsection{\textcolor{Chapter }{PrimalityProof}}
\logpage{[ 14, 4, 3 ]}\nobreak
\hyperdef{L}{X7CD977B17B4A7A4B}{}
{\noindent\textcolor{FuncColor}{$\triangleright$\enspace\texttt{PrimalityProof({\mdseries\slshape n})\index{PrimalityProof@\texttt{PrimalityProof}}
\label{PrimalityProof}
}\hfill{\scriptsize (function)}}\\


 Construct a machine verifiable proof of the primality of (the probable prime) \mbox{\texttt{\mdseries\slshape n}}, following the ideas of \cite{BLS1975}. The proof consists of various Fermat and Lucas pseudoprimality tests, which
taken as a whole prove the primality. The proof is represented as a list of
witnesses of two kinds. The first kind, \texttt{[ "F", divisor, base ]}, indicates a successful Fermat pseudoprimality test, where \mbox{\texttt{\mdseries\slshape n}} is a strong pseudoprime at \texttt{base} with order not divisible by $(\mbox{\texttt{\mdseries\slshape n}}-1)/divisor$. The second kind, \texttt{[ "L", divisor, discriminant, P ]} indicates a successful Lucas pseudoprimality test, for a quadratic form of
given \texttt{discriminant} and middle term \texttt{P} with an extra check at $(\mbox{\texttt{\mdseries\slshape n}}+1)/divisor$. }

 

\subsection{\textcolor{Chapter }{IsPrimePowerInt}}
\logpage{[ 14, 4, 4 ]}\nobreak
\hyperdef{L}{X8443125D7FD6F2A6}{}
{\noindent\textcolor{FuncColor}{$\triangleright$\enspace\texttt{IsPrimePowerInt({\mdseries\slshape n})\index{IsPrimePowerInt@\texttt{IsPrimePowerInt}}
\label{IsPrimePowerInt}
}\hfill{\scriptsize (function)}}\\


 \texttt{IsPrimePowerInt} returns \texttt{true} if the integer \mbox{\texttt{\mdseries\slshape n}} is a prime power and \texttt{false} otherwise. 

 An integer $n$ is a \emph{prime power} if there exists a prime $p$ and a positive integer $i$ such that $p^i = n$. If $n$ is negative the condition is that there must exist a negative prime $p$ and an odd positive integer $i$ such that $p^i = n$. The integers 1 and \texttt{\symbol{45}}1 are not prime powers. 

 Note that \texttt{IsPrimePowerInt} uses \texttt{SmallestRootInt} (\ref{SmallestRootInt}) and a probable\texttt{\symbol{45}}primality test (see \texttt{IsPrimeInt} (\ref{IsPrimeInt})). 
\begin{Verbatim}[commandchars=!@|,fontsize=\small,frame=single,label=Example]
  !gapprompt@gap>| !gapinput@IsPrimePowerInt( 31^5 );|
  true
  !gapprompt@gap>| !gapinput@IsPrimePowerInt( 2^31-1 );  # 2^31-1 is actually a prime|
  true
  !gapprompt@gap>| !gapinput@IsPrimePowerInt( 2^63-1 );|
  false
  !gapprompt@gap>| !gapinput@Filtered( [-10..10], IsPrimePowerInt );|
  [ -8, -7, -5, -3, -2, 2, 3, 4, 5, 7, 8, 9 ]
\end{Verbatim}
 }

 

\subsection{\textcolor{Chapter }{NextPrimeInt}}
\logpage{[ 14, 4, 5 ]}\nobreak
\hyperdef{L}{X78744C367A94C69F}{}
{\noindent\textcolor{FuncColor}{$\triangleright$\enspace\texttt{NextPrimeInt({\mdseries\slshape n})\index{NextPrimeInt@\texttt{NextPrimeInt}}
\label{NextPrimeInt}
}\hfill{\scriptsize (function)}}\\


 \texttt{NextPrimeInt} returns the smallest prime which is strictly larger than the integer \mbox{\texttt{\mdseries\slshape n}}. 

 Note that \texttt{NextPrimeInt} uses a probable\texttt{\symbol{45}}primality test (see \texttt{IsPrimeInt} (\ref{IsPrimeInt})). 
\begin{Verbatim}[commandchars=!@|,fontsize=\small,frame=single,label=Example]
  !gapprompt@gap>| !gapinput@NextPrimeInt( 541 ); NextPrimeInt( -1 );|
  547
  2
\end{Verbatim}
 }

 

\subsection{\textcolor{Chapter }{PrevPrimeInt}}
\logpage{[ 14, 4, 6 ]}\nobreak
\hyperdef{L}{X819060E17E83728A}{}
{\noindent\textcolor{FuncColor}{$\triangleright$\enspace\texttt{PrevPrimeInt({\mdseries\slshape n})\index{PrevPrimeInt@\texttt{PrevPrimeInt}}
\label{PrevPrimeInt}
}\hfill{\scriptsize (function)}}\\


 \texttt{PrevPrimeInt} returns the largest prime which is strictly smaller than the integer \mbox{\texttt{\mdseries\slshape n}}. 

 Note that \texttt{PrevPrimeInt} uses a probable\texttt{\symbol{45}}primality test (see \texttt{IsPrimeInt} (\ref{IsPrimeInt})). 
\begin{Verbatim}[commandchars=!@|,fontsize=\small,frame=single,label=Example]
  !gapprompt@gap>| !gapinput@PrevPrimeInt( 541 ); PrevPrimeInt( 1 );|
  523
  -2
\end{Verbatim}
 }

 

\subsection{\textcolor{Chapter }{FactorsInt}}
\logpage{[ 14, 4, 7 ]}\nobreak
\hyperdef{L}{X82C989DB84744B36}{}
{\noindent\textcolor{FuncColor}{$\triangleright$\enspace\texttt{FactorsInt({\mdseries\slshape n})\index{FactorsInt@\texttt{FactorsInt}}
\label{FactorsInt}
}\hfill{\scriptsize (function)}}\\
\noindent\textcolor{FuncColor}{$\triangleright$\enspace\texttt{FactorsInt({\mdseries\slshape n: RhoTrials := trials})\index{FactorsInt@\texttt{FactorsInt}!using Pollard's Rho}
\label{FactorsInt:using Pollard's Rho}
}\hfill{\scriptsize (function)}}\\


 \texttt{FactorsInt} returns a list of factors of a given integer \mbox{\texttt{\mdseries\slshape n}} such that \texttt{Product( FactorsInt( \mbox{\texttt{\mdseries\slshape n}} ) ) = \mbox{\texttt{\mdseries\slshape n}}}. If $|n| \leq 1$ the list \texttt{[\mbox{\texttt{\mdseries\slshape n}}]} is returned. Otherwise the result contains probable primes, sorted by absolute
value. The entries will all be positive except for the first one in case of a
negative \mbox{\texttt{\mdseries\slshape n}}. 

 See \texttt{PrimeDivisors} (\ref{PrimeDivisors}) for a function that returns a set of (probable) primes dividing \mbox{\texttt{\mdseries\slshape n}}. 

 Note that \texttt{FactorsInt} uses a probable\texttt{\symbol{45}}primality test (see{\nobreakspace}\texttt{IsPrimeInt} (\ref{IsPrimeInt})). Thus \texttt{FactorsInt} might return a list which contains composite integers. In such a case you will
get a warning about the use of a probable prime. You can switch off these
warnings by \texttt{SetInfoLevel( InfoPrimeInt, 0 );} (also see \texttt{SetInfoLevel} (\ref{SetInfoLevel})). 

 The time taken by \texttt{FactorsInt} is approximately proportional to the square root of the second largest prime
factor of \mbox{\texttt{\mdseries\slshape n}}, which is the last one that \texttt{FactorsInt} has to find, since the largest factor is simply what remains when all others
have been removed. Thus the time is roughly bounded by the fourth root of \mbox{\texttt{\mdseries\slshape n}}. \texttt{FactorsInt} is guaranteed to find all factors less than $10^6$ and will find most factors less than $10^{10}$. If \mbox{\texttt{\mdseries\slshape n}} contains multiple factors larger than that \texttt{FactorsInt} may not be able to factor \mbox{\texttt{\mdseries\slshape n}} and will then signal an error. 

 \texttt{FactorsInt} is used in a method for the general operation \texttt{Factors} (\ref{Factors}). 

 In the second form, \texttt{FactorsInt} calls \texttt{FactorsRho} with a limit of \mbox{\texttt{\mdseries\slshape trials}} on the number of trials it performs. The default is 8192. Note that Pollard's
Rho is the fastest method for finding prime factors with roughly
5\texttt{\symbol{45}}10 decimal digits, but becomes more and more inferior to
other factorization techniques like e.g. the Elliptic Curves Method (ECM) the
bigger the prime factors are. Therefore instead of performing a huge number of
Rho \mbox{\texttt{\mdseries\slshape trials}}, it is usually more advisable to install the \textsf{FactInt} package and then simply to use the operation \texttt{Factors} (\ref{Factors}). The factorization of the 8\texttt{\symbol{45}}th Fermat number by Pollard's
Rho below takes already a while. 
\begin{Verbatim}[commandchars=!@|,fontsize=\small,frame=single,label=Example]
  !gapprompt@gap>| !gapinput@FactorsInt( -Factorial(6) );|
  [ -2, 2, 2, 2, 3, 3, 5 ]
  !gapprompt@gap>| !gapinput@Set( FactorsInt( Factorial(13)/11 ) );|
  [ 2, 3, 5, 7, 13 ]
  !gapprompt@gap>| !gapinput@FactorsInt( 2^63 - 1 );|
  [ 7, 7, 73, 127, 337, 92737, 649657 ]
  !gapprompt@gap>| !gapinput@FactorsInt( 10^42 + 1 );|
  [ 29, 101, 281, 9901, 226549, 121499449, 4458192223320340849 ]
  !gapprompt@gap>| !gapinput@FactorsInt(2^256+1:RhoTrials:=100000000);|
  [ 1238926361552897,
    93461639715357977769163558199606896584051237541638188580280321 ]
\end{Verbatim}
 }

 

\subsection{\textcolor{Chapter }{PrimeDivisors}}
\logpage{[ 14, 4, 8 ]}\nobreak
\hyperdef{L}{X80E7A5D381C64CC9}{}
{\noindent\textcolor{FuncColor}{$\triangleright$\enspace\texttt{PrimeDivisors({\mdseries\slshape n})\index{PrimeDivisors@\texttt{PrimeDivisors}}
\label{PrimeDivisors}
}\hfill{\scriptsize (attribute)}}\\


 \texttt{PrimeDivisors} returns for a non\texttt{\symbol{45}}zero integer \mbox{\texttt{\mdseries\slshape n}} a set of its positive (probable) primes divisors. In rare cases the result
could contain a composite number which passed certain primality tests, see \texttt{IsProbablyPrimeInt} (\ref{IsProbablyPrimeInt}) and \texttt{FactorsInt} (\ref{FactorsInt}) for more details. 
\begin{Verbatim}[commandchars=!@|,fontsize=\small,frame=single,label=Example]
  !gapprompt@gap>| !gapinput@PrimeDivisors(-12);|
  [ 2, 3 ]
  !gapprompt@gap>| !gapinput@PrimeDivisors(1);|
  [  ]
\end{Verbatim}
 }

 

\subsection{\textcolor{Chapter }{PartialFactorization}}
\logpage{[ 14, 4, 9 ]}\nobreak
\hyperdef{L}{X786FF92C7C54BF97}{}
{\noindent\textcolor{FuncColor}{$\triangleright$\enspace\texttt{PartialFactorization({\mdseries\slshape n[, effort]})\index{PartialFactorization@\texttt{PartialFactorization}}
\label{PartialFactorization}
}\hfill{\scriptsize (operation)}}\\


 \texttt{PartialFactorization} returns a partial factorization of the integer \mbox{\texttt{\mdseries\slshape n}}. No assertions are made about the primality of the factors, except of those
mentioned below. 

 The argument \mbox{\texttt{\mdseries\slshape effort}}, if given, specifies how intensively the function should try to determine
factors of \mbox{\texttt{\mdseries\slshape n}}. The default is \mbox{\texttt{\mdseries\slshape effort}}{\nobreakspace}={\nobreakspace}5. 

 
\begin{itemize}
\item  If \mbox{\texttt{\mdseries\slshape effort}}{\nobreakspace}={\nobreakspace}0, trial division by the primes below 100 is
done. Returned factors below $10^4$ are guaranteed to be prime. 
\item  If \mbox{\texttt{\mdseries\slshape effort}}{\nobreakspace}={\nobreakspace}1, trial division by the primes below 1000 is
done. Returned factors below $10^6$ are guaranteed to be prime. 
\item  If \mbox{\texttt{\mdseries\slshape effort}}{\nobreakspace}={\nobreakspace}2, additionally trial division by the numbers
in the lists \texttt{Primes2} and \texttt{ProbablePrimes2} is done, and perfect powers are detected. Returned factors below $10^6$ are guaranteed to be prime. 
\item  If \mbox{\texttt{\mdseries\slshape effort}}{\nobreakspace}={\nobreakspace}3, additionally \texttt{FactorsRho} (Pollard's Rho) with \texttt{RhoTrials} = 256 is used. 
\item  If \mbox{\texttt{\mdseries\slshape effort}}{\nobreakspace}={\nobreakspace}4, as above, but \texttt{RhoTrials} = 2048. 
\item  If \mbox{\texttt{\mdseries\slshape effort}}{\nobreakspace}={\nobreakspace}5, as above, but \texttt{RhoTrials} = 8192. Returned factors below $10^{12}$ are guaranteed to be prime, and all prime factors below $10^6$ are guaranteed to be found. 
\item  If \mbox{\texttt{\mdseries\slshape effort}}{\nobreakspace}={\nobreakspace}6 and the package \textsf{FactInt} is loaded, in addition to the above quite a number of special cases are
handled. 
\item  If \mbox{\texttt{\mdseries\slshape effort}}{\nobreakspace}={\nobreakspace}7 and the package \textsf{FactInt} is loaded, the only thing which is not attempted to obtain a full
factorization into
Baillie\texttt{\symbol{45}}Pomerance\texttt{\symbol{45}}Selfridge\texttt{\symbol{45}}Wagstaff
pseudoprimes is the application of the MPQS to a remaining composite with more
than 50 decimal digits. 
\end{itemize}
 

 Increasing the value of the argument \mbox{\texttt{\mdseries\slshape effort}} by one usually results in an increase of the runtime requirements by a factor
of (very roughly!) 3 to{\nobreakspace}10. (Also see \texttt{CheapFactorsInt} (\textbf{EDIM: CheapFactorsInt})). 
\begin{Verbatim}[commandchars=!@|,fontsize=\small,frame=single,label=Example]
  !gapprompt@gap>| !gapinput@List([0..5],i->PartialFactorization(97^35-1,i));|
  [ [ 2, 2, 2, 2, 2, 3, 11, 31, 43,
        2446338959059521520901826365168917110105972824229555319002965029 ],
    [ 2, 2, 2, 2, 2, 3, 11, 31, 43, 967,
        2529823122088440042297648774735177983563570655873376751812787 ],
    [ 2, 2, 2, 2, 2, 3, 11, 31, 43, 967,
        2529823122088440042297648774735177983563570655873376751812787 ],
    [ 2, 2, 2, 2, 2, 3, 11, 31, 43, 967, 39761, 262321,
        242549173950325921859769421435653153445616962914227 ],
    [ 2, 2, 2, 2, 2, 3, 11, 31, 43, 967, 39761, 262321, 687121,
        352993394104278463123335513593170858474150787 ],
    [ 2, 2, 2, 2, 2, 3, 11, 31, 43, 967, 39761, 262321, 687121,
        20241187, 504769301, 34549173843451574629911361501 ] ]
\end{Verbatim}
 }

 

\subsection{\textcolor{Chapter }{PrintFactorsInt}}
\logpage{[ 14, 4, 10 ]}\nobreak
\hyperdef{L}{X803D431087B6FF28}{}
{\noindent\textcolor{FuncColor}{$\triangleright$\enspace\texttt{PrintFactorsInt({\mdseries\slshape n})\index{PrintFactorsInt@\texttt{PrintFactorsInt}}
\label{PrintFactorsInt}
}\hfill{\scriptsize (function)}}\\


 prints the prime factorization of the integer \mbox{\texttt{\mdseries\slshape n}} in human\texttt{\symbol{45}}readable form. See also \texttt{StringPP} (\ref{StringPP}). 
\begin{Verbatim}[commandchars=!@|,fontsize=\small,frame=single,label=Example]
  !gapprompt@gap>| !gapinput@PrintFactorsInt( Factorial( 7 ) ); Print( "\n" );|
  2^4*3^2*5*7
\end{Verbatim}
 }

 

\subsection{\textcolor{Chapter }{PrimePowersInt}}
\logpage{[ 14, 4, 11 ]}\nobreak
\hyperdef{L}{X82148B347E294C87}{}
{\noindent\textcolor{FuncColor}{$\triangleright$\enspace\texttt{PrimePowersInt({\mdseries\slshape n})\index{PrimePowersInt@\texttt{PrimePowersInt}}
\label{PrimePowersInt}
}\hfill{\scriptsize (function)}}\\


 returns the prime factorization of the integer \mbox{\texttt{\mdseries\slshape n}} as a list $[ p_1, e_1, \ldots, p_k, e_k ]$ with \mbox{\texttt{\mdseries\slshape n}} = $p_1^{{e_1}} \cdot p_2^{{e_2}} \cdot ... \cdot p_k^{{e_k}}$. 

 For negative integers, the absolute value is taken. Zero is not allowed as
input. 
\begin{Verbatim}[commandchars=!@|,fontsize=\small,frame=single,label=Example]
  !gapprompt@gap>| !gapinput@PrimePowersInt( Factorial( 7 ) );|
  [ 2, 4, 3, 2, 5, 1, 7, 1 ]
  !gapprompt@gap>| !gapinput@PrimePowersInt( 1 );|
  [  ]
\end{Verbatim}
 }

 

\subsection{\textcolor{Chapter }{DivisorsInt}}
\logpage{[ 14, 4, 12 ]}\nobreak
\hyperdef{L}{X809E0E1B83AF7695}{}
{\noindent\textcolor{FuncColor}{$\triangleright$\enspace\texttt{DivisorsInt({\mdseries\slshape n})\index{DivisorsInt@\texttt{DivisorsInt}}
\label{DivisorsInt}
}\hfill{\scriptsize (function)}}\\


 \index{divisors!of an integer} \texttt{DivisorsInt} returns a list of all divisors of the integer \mbox{\texttt{\mdseries\slshape n}}. The list is sorted, so that it starts with 1 and ends with \mbox{\texttt{\mdseries\slshape n}}. We define that \texttt{DivisorsInt( \texttt{\symbol{45}}\mbox{\texttt{\mdseries\slshape n}} ) = DivisorsInt( \mbox{\texttt{\mdseries\slshape n}} )}. 

 Since the set of divisors of 0 is infinite calling \texttt{DivisorsInt( 0 )} causes an error. 

 \texttt{DivisorsInt} may call \texttt{FactorsInt} (\ref{FactorsInt}) to obtain the prime factors. \texttt{Sigma} (\ref{Sigma}) and \texttt{Tau} (\ref{Tau}) compute the sum and the number of positive divisors, respectively. 
\begin{Verbatim}[commandchars=!@|,fontsize=\small,frame=single,label=Example]
  !gapprompt@gap>| !gapinput@DivisorsInt( 1 ); DivisorsInt( 20 ); DivisorsInt( 541 );|
  [ 1 ]
  [ 1, 2, 4, 5, 10, 20 ]
  [ 1, 541 ]
\end{Verbatim}
 }

 }

  
\section{\textcolor{Chapter }{Residue Class Rings}}\label{Residue Class Rings}
\logpage{[ 14, 5, 0 ]}
\hyperdef{L}{X864BF040862409FC}{}
{
  \index{mod@\texttt{mod}!residue class rings} \index{modulo!residue class rings} \texttt{ZmodnZ} (\ref{ZmodnZ}) returns a residue class ring of \texttt{Integers} (\ref{Integers}) modulo an ideal. These residue class rings are rings, thus all operations for
rings (see Chapter{\nobreakspace}\ref{Rings}) apply. See also Chapters{\nobreakspace}\ref{Finite Fields} and \ref{Number Theory}. 

\subsection{\textcolor{Chapter }{\texttt{\symbol{92}}mod (for residue class rings)}}
\logpage{[ 14, 5, 1 ]}\nobreak
\hyperdef{L}{X87B1210B8581D5B2}{}
{\noindent\textcolor{FuncColor}{$\triangleright$\enspace\texttt{\texttt{\symbol{92}}mod({\mdseries\slshape r/s, n})\index{\texttt{\symbol{92}}mod@\texttt{\texttt{\symbol{92}}mod}!for residue class rings}
\label{bSlashmod:for residue class rings}
}\hfill{\scriptsize (operation)}}\\


 If \mbox{\texttt{\mdseries\slshape r}}, \mbox{\texttt{\mdseries\slshape s}} and \mbox{\texttt{\mdseries\slshape n}} are integers, \texttt{\mbox{\texttt{\mdseries\slshape r}} / \mbox{\texttt{\mdseries\slshape s}}} as a reduced fraction is \texttt{p/q}, where \texttt{q} and \mbox{\texttt{\mdseries\slshape n}} are coprime, then \texttt{\mbox{\texttt{\mdseries\slshape r}} / \mbox{\texttt{\mdseries\slshape s}} mod \mbox{\texttt{\mdseries\slshape n}}} is defined to be the product of \texttt{p} and the inverse of \texttt{q} modulo \mbox{\texttt{\mdseries\slshape n}}. See Section{\nobreakspace}\ref{Arithmetic Operators} for more details and definitions. 

 With the above definition, \texttt{4 / 6 mod 32} equals \texttt{2 / 3 mod 32} and hence exists (and is equal to 22), despite the fact that 6 has no inverse
modulo 32. }

 

\subsection{\textcolor{Chapter }{ZmodnZ}}
\logpage{[ 14, 5, 2 ]}\nobreak
\hyperdef{L}{X79CE76AD82B3E2B2}{}
{\noindent\textcolor{FuncColor}{$\triangleright$\enspace\texttt{ZmodnZ({\mdseries\slshape n})\index{ZmodnZ@\texttt{ZmodnZ}}
\label{ZmodnZ}
}\hfill{\scriptsize (function)}}\\
\noindent\textcolor{FuncColor}{$\triangleright$\enspace\texttt{ZmodpZ({\mdseries\slshape p})\index{ZmodpZ@\texttt{ZmodpZ}}
\label{ZmodpZ}
}\hfill{\scriptsize (function)}}\\
\noindent\textcolor{FuncColor}{$\triangleright$\enspace\texttt{ZmodpZNC({\mdseries\slshape p})\index{ZmodpZNC@\texttt{ZmodpZNC}}
\label{ZmodpZNC}
}\hfill{\scriptsize (function)}}\\


 \texttt{ZmodnZ} returns a ring $R$ isomorphic to the residue class ring of the integers modulo the ideal
generated by \mbox{\texttt{\mdseries\slshape n}}. The element corresponding to the residue class of the integer $i$ in this ring can be obtained by \texttt{i * One( R )}, and a representative of the residue class corresponding to the element $x \in R$ can be computed by \texttt{Int}$( x )$. 

 \index{mod!Integers} \texttt{ZmodnZ( \mbox{\texttt{\mdseries\slshape n}} )} is equal to \texttt{Integers mod \mbox{\texttt{\mdseries\slshape n}}}. 

 \texttt{ZmodpZ} does the same if the argument \mbox{\texttt{\mdseries\slshape p}} is a prime integer, additionally the result is a field. \texttt{ZmodpZNC} omits the check whether \mbox{\texttt{\mdseries\slshape p}} is a prime. 

 Each ring returned by these functions contains the whole family of its
elements if \mbox{\texttt{\mdseries\slshape n}} is not a prime, and is embedded into the family of finite field elements of
characteristic \mbox{\texttt{\mdseries\slshape n}} if \mbox{\texttt{\mdseries\slshape n}} is a prime. }

 

\subsection{\textcolor{Chapter }{ZmodnZObj (for a residue class family and integer)}}
\logpage{[ 14, 5, 3 ]}\nobreak
\hyperdef{L}{X838F36507D985EDA}{}
{\noindent\textcolor{FuncColor}{$\triangleright$\enspace\texttt{ZmodnZObj({\mdseries\slshape Fam, r})\index{ZmodnZObj@\texttt{ZmodnZObj}!for a residue class family and integer}
\label{ZmodnZObj:for a residue class family and integer}
}\hfill{\scriptsize (operation)}}\\
\noindent\textcolor{FuncColor}{$\triangleright$\enspace\texttt{ZmodnZObj({\mdseries\slshape r, n})\index{ZmodnZObj@\texttt{ZmodnZObj}!for two integers}
\label{ZmodnZObj:for two integers}
}\hfill{\scriptsize (operation)}}\\


 If the first argument is a residue class family \mbox{\texttt{\mdseries\slshape Fam}} then \texttt{ZmodnZObj} returns the element in \mbox{\texttt{\mdseries\slshape Fam}} whose coset is represented by the integer \mbox{\texttt{\mdseries\slshape r}}. 

 If the two arguments are an integer \mbox{\texttt{\mdseries\slshape r}} and a positive integer \mbox{\texttt{\mdseries\slshape n}} then \texttt{ZmodnZObj} returns the element in \texttt{ZmodnZ( \mbox{\texttt{\mdseries\slshape n}} )} (see{\nobreakspace}\texttt{ZmodnZ} (\ref{ZmodnZ})) whose coset is represented by the integer \mbox{\texttt{\mdseries\slshape r}}. 

 
\begin{Verbatim}[commandchars=!@|,fontsize=\small,frame=single,label=Example]
  !gapprompt@gap>| !gapinput@r:= ZmodnZ(15);|
  (Integers mod 15)
  !gapprompt@gap>| !gapinput@fam:=ElementsFamily(FamilyObj(r));;|
  !gapprompt@gap>| !gapinput@a:= ZmodnZObj(fam,9);|
  ZmodnZObj( 9, 15 )
  !gapprompt@gap>| !gapinput@a+a;|
  ZmodnZObj( 3, 15 )
  !gapprompt@gap>| !gapinput@Int(a+a);|
  3
\end{Verbatim}
 }

 

\subsection{\textcolor{Chapter }{IsZmodnZObj}}
\logpage{[ 14, 5, 4 ]}\nobreak
\hyperdef{L}{X7D0107DD79753901}{}
{\noindent\textcolor{FuncColor}{$\triangleright$\enspace\texttt{IsZmodnZObj({\mdseries\slshape obj})\index{IsZmodnZObj@\texttt{IsZmodnZObj}}
\label{IsZmodnZObj}
}\hfill{\scriptsize (Category)}}\\
\noindent\textcolor{FuncColor}{$\triangleright$\enspace\texttt{IsZmodnZObjNonprime({\mdseries\slshape obj})\index{IsZmodnZObjNonprime@\texttt{IsZmodnZObjNonprime}}
\label{IsZmodnZObjNonprime}
}\hfill{\scriptsize (Category)}}\\
\noindent\textcolor{FuncColor}{$\triangleright$\enspace\texttt{IsZmodpZObj({\mdseries\slshape obj})\index{IsZmodpZObj@\texttt{IsZmodpZObj}}
\label{IsZmodpZObj}
}\hfill{\scriptsize (Category)}}\\
\noindent\textcolor{FuncColor}{$\triangleright$\enspace\texttt{IsZmodpZObjSmall({\mdseries\slshape obj})\index{IsZmodpZObjSmall@\texttt{IsZmodpZObjSmall}}
\label{IsZmodpZObjSmall}
}\hfill{\scriptsize (Category)}}\\
\noindent\textcolor{FuncColor}{$\triangleright$\enspace\texttt{IsZmodpZObjLarge({\mdseries\slshape obj})\index{IsZmodpZObjLarge@\texttt{IsZmodpZObjLarge}}
\label{IsZmodpZObjLarge}
}\hfill{\scriptsize (Category)}}\\


 The elements in the rings $Z / n Z$ are in the category \texttt{IsZmodnZObj}. If $n$ is a prime then the elements are of course also in the category \texttt{IsFFE} (\ref{IsFFE}), otherwise they are in \texttt{IsZmodnZObjNonprime}. \texttt{IsZmodpZObj} is an abbreviation of \texttt{IsZmodnZObj and IsFFE}. This category is the disjoint union of \texttt{IsZmodpZObjSmall} and \texttt{IsZmodpZObjLarge}, the former containing all elements with $n$ at most \texttt{MAXSIZE{\textunderscore}GF{\textunderscore}INTERNAL}. 

 The reasons to distinguish the prime case from the nonprime case are 
\begin{itemize}
\item  that objects in \texttt{IsZmodnZObjNonprime} have an external representation (namely the residue in the range $[ 0, 1, \ldots, n-1 ]$), 
\item  that the comparison of elements can be defined as comparison of the residues,
and 
\item  that the elements lie in a family of type \texttt{IsZmodnZObjNonprimeFamily} (note that for prime $n$, the family must be an \texttt{IsFFEFamily}). 
\end{itemize}
 

 The reasons to distinguish the small and the large case are that for small $n$ the elements must be compatible with the internal representation of finite
field elements, whereas we are free to define comparison as comparison of
residues for large $n$. 

 Note that we \emph{cannot} claim that every finite field element of degree 1 is in \texttt{IsZmodnZObj}, since finite field elements in internal representation may not know that
they lie in the prime field. }

 }

  
\section{\textcolor{Chapter }{Check Digits}}\label{Check Digits}
\logpage{[ 14, 6, 0 ]}
\hyperdef{L}{X7904B6D681EBF091}{}
{
  

\subsection{\textcolor{Chapter }{CheckDigitISBN}}
\logpage{[ 14, 6, 1 ]}\nobreak
\hyperdef{L}{X82BABA8F868BD425}{}
{\noindent\textcolor{FuncColor}{$\triangleright$\enspace\texttt{CheckDigitISBN({\mdseries\slshape n})\index{CheckDigitISBN@\texttt{CheckDigitISBN}}
\label{CheckDigitISBN}
}\hfill{\scriptsize (function)}}\\
\noindent\textcolor{FuncColor}{$\triangleright$\enspace\texttt{CheckDigitISBN13({\mdseries\slshape n})\index{CheckDigitISBN13@\texttt{CheckDigitISBN13}}
\label{CheckDigitISBN13}
}\hfill{\scriptsize (function)}}\\
\noindent\textcolor{FuncColor}{$\triangleright$\enspace\texttt{CheckDigitPostalMoneyOrder({\mdseries\slshape n})\index{CheckDigitPostalMoneyOrder@\texttt{CheckDigitPostalMoneyOrder}}
\label{CheckDigitPostalMoneyOrder}
}\hfill{\scriptsize (function)}}\\
\noindent\textcolor{FuncColor}{$\triangleright$\enspace\texttt{CheckDigitUPC({\mdseries\slshape n})\index{CheckDigitUPC@\texttt{CheckDigitUPC}}
\label{CheckDigitUPC}
}\hfill{\scriptsize (function)}}\\


 These functions can be used to compute, or check, check digits for some
everyday items. In each case what is submitted as input is either the number
with check digit (in which case the function returns \texttt{true} or \texttt{false}), or the number without check digit (in which case the function returns the
missing check digit). The number can be specified as integer, as string (for
example in case of leading zeros) or as a sequence of arguments, each
representing a single digit. The check digits tested are the
10\texttt{\symbol{45}}digit ISBN (International Standard Book Number) using \texttt{CheckDigitISBN} (since arithmetic is module 11, a digit 11 is represented by an X); the newer
13\texttt{\symbol{45}}digit ISBN\texttt{\symbol{45}}13 using \texttt{CheckDigitISBN13}; the numbers of 11\texttt{\symbol{45}}digit US postal money orders using \texttt{CheckDigitPostalMoneyOrder}; and the 12\texttt{\symbol{45}}digit UPC bar code found on groceries using \texttt{CheckDigitUPC}. 
\begin{Verbatim}[commandchars=@|A,fontsize=\small,frame=single,label=Example]
  @gapprompt|gap>A @gapinput|CheckDigitISBN("052166103");A
  Check Digit is 'X'
  'X'
  @gapprompt|gap>A @gapinput|CheckDigitISBN("052166103X");A
  Checksum test satisfied
  true
  @gapprompt|gap>A @gapinput|CheckDigitISBN(0,5,2,1,6,6,1,0,3,1);A
  Checksum test failed
  false
  @gapprompt|gap>A @gapinput|CheckDigitISBN(0,5,2,1,6,6,1,0,3,'X'); # note single quotes!A
  Checksum test satisfied
  true
  @gapprompt|gap>A @gapinput|CheckDigitISBN13("9781420094527");A
  Checksum test satisfied
  true
  @gapprompt|gap>A @gapinput|CheckDigitUPC("07164183001");A
  Check Digit is 1
  1
  @gapprompt|gap>A @gapinput|CheckDigitPostalMoneyOrder(16786457155);A
  Checksum test satisfied
  true
\end{Verbatim}
 }

 

\subsection{\textcolor{Chapter }{CheckDigitTestFunction}}
\logpage{[ 14, 6, 2 ]}\nobreak
\hyperdef{L}{X85F1A6A5870485B9}{}
{\noindent\textcolor{FuncColor}{$\triangleright$\enspace\texttt{CheckDigitTestFunction({\mdseries\slshape l, m, f})\index{CheckDigitTestFunction@\texttt{CheckDigitTestFunction}}
\label{CheckDigitTestFunction}
}\hfill{\scriptsize (function)}}\\


 This function creates check digit test functions such as \texttt{CheckDigitISBN} (\ref{CheckDigitISBN}) for check digit schemes that use the inner products with a fixed vector modulo
a number. The scheme creates will use strings of \mbox{\texttt{\mdseries\slshape l}} digits (including the check digits), the check consists of taking the standard
product of the vector of digits with the fixed vector \mbox{\texttt{\mdseries\slshape f}} modulo \mbox{\texttt{\mdseries\slshape m}}; the result needs to be 0. The function returns a function that then can be
used for testing or determining check digits. 
\begin{Verbatim}[commandchars=!@|,fontsize=\small,frame=single,label=Example]
  !gapprompt@gap>| !gapinput@isbntest:=CheckDigitTestFunction(10,11,[1,2,3,4,5,6,7,8,9,-1]);|
  function( arg... ) ... end
  !gapprompt@gap>| !gapinput@isbntest("038794680");|
  Check Digit is 2
  2
\end{Verbatim}
 }

 }

  
\section{\textcolor{Chapter }{Random Sources}}\label{Random Sources}
\logpage{[ 14, 7, 0 ]}
\hyperdef{L}{X85361FAE8088C006}{}
{
  \textsf{GAP} provides \texttt{Random} (\ref{Random:for a list or collection}) methods for many collections of objects. On a lower level these methods use \emph{random sources} which provide random integers and random choices from lists. 

 See \texttt{IsRandomSource} (\ref{IsRandomSource}) for the user interface for random sources, and Section \ref{Implementing new kinds of random sources} for information about developing new kinds of random sources. 

\subsection{\textcolor{Chapter }{IsRandomSource}}
\logpage{[ 14, 7, 1 ]}\nobreak
\hyperdef{L}{X82E31A697E389F1D}{}
{\noindent\textcolor{FuncColor}{$\triangleright$\enspace\texttt{IsRandomSource({\mdseries\slshape obj})\index{IsRandomSource@\texttt{IsRandomSource}}
\label{IsRandomSource}
}\hfill{\scriptsize (Category)}}\\


 This is the category of random source objects. The \emph{user interface} for these objects consists of the following functions. 

 \texttt{RandomSource} (\ref{RandomSource}) creates a new random source \mbox{\texttt{\mdseries\slshape rs}}, say. 

 \texttt{Random( \mbox{\texttt{\mdseries\slshape rs}}, \mbox{\texttt{\mdseries\slshape list}} )} yields a random element of the list \mbox{\texttt{\mdseries\slshape list}}, and \texttt{Random( \mbox{\texttt{\mdseries\slshape rs}}, \mbox{\texttt{\mdseries\slshape low}}, \mbox{\texttt{\mdseries\slshape high}} )} yields a random integer between \mbox{\texttt{\mdseries\slshape low}} and \mbox{\texttt{\mdseries\slshape high}} (inclusive), see \texttt{Random} (\ref{Random:for random source and list}). 

 If \mbox{\texttt{\mdseries\slshape rs}} supports resetting (see \texttt{State} (\ref{State})) then \texttt{State( \mbox{\texttt{\mdseries\slshape rs}} )} yields a copy \mbox{\texttt{\mdseries\slshape state}}, say, of the current state of \mbox{\texttt{\mdseries\slshape rs}} such that \texttt{Reset( \mbox{\texttt{\mdseries\slshape rs}}, \mbox{\texttt{\mdseries\slshape state}} )} resets \mbox{\texttt{\mdseries\slshape rs}} to the given state. 

 One idea behind providing several independent (pseudo) random sources is to
make algorithms which use some sort of random choices deterministic. They can
use their own new random source created with a fixed seed and so do exactly
the same in different calls. 

 Random source objects lie in the family \texttt{RandomSourcesFamily}. }

 

\subsection{\textcolor{Chapter }{Random (for random source and list)}}
\logpage{[ 14, 7, 2 ]}\nobreak
\hyperdef{L}{X821004F286282D49}{}
{\noindent\textcolor{FuncColor}{$\triangleright$\enspace\texttt{Random({\mdseries\slshape rs, list})\index{Random@\texttt{Random}!for random source and list}
\label{Random:for random source and list}
}\hfill{\scriptsize (operation)}}\\
\noindent\textcolor{FuncColor}{$\triangleright$\enspace\texttt{Random({\mdseries\slshape rs, coll})\index{Random@\texttt{Random}!for random source and collection}
\label{Random:for random source and collection}
}\hfill{\scriptsize (operation)}}\\
\noindent\textcolor{FuncColor}{$\triangleright$\enspace\texttt{Random({\mdseries\slshape rs, low, high})\index{Random@\texttt{Random}!for random source and two integers}
\label{Random:for random source and two integers}
}\hfill{\scriptsize (operation)}}\\


 This operation returns a random element from the dense, nonempty list \mbox{\texttt{\mdseries\slshape list}} or the nonempty collection \mbox{\texttt{\mdseries\slshape coll}}, or an integer in the range from the given (possibly large) integers \mbox{\texttt{\mdseries\slshape low}} to \mbox{\texttt{\mdseries\slshape high}}, respectively. 

 The choice should only depend on the random source \mbox{\texttt{\mdseries\slshape rs}} and have no effect on other random sources. 

 It is not defined what happens if \mbox{\texttt{\mdseries\slshape list}} or \mbox{\texttt{\mdseries\slshape coll}} is empty, \mbox{\texttt{\mdseries\slshape list}} is not dense, or \mbox{\texttt{\mdseries\slshape low}} is larger than \mbox{\texttt{\mdseries\slshape high}}. 
\begin{Verbatim}[commandchars=!@|,fontsize=\small,frame=single,label=Example]
  !gapprompt@gap>| !gapinput@mysource := RandomSource(IsMersenneTwister, 42);;|
  !gapprompt@gap>| !gapinput@Random(mysource, 1, 10^60);|
  999331861769949319194941485000557997842686717712198687315183
\end{Verbatim}
 }

 
\subsection{\textcolor{Chapter }{State and Reset for Random Sources}}\label{State and Reset for Random Sources}
\logpage{[ 14, 7, 3 ]}
\hyperdef{L}{X86FFFBC9790F9742}{}
{
\noindent\textcolor{FuncColor}{$\triangleright$\enspace\texttt{State({\mdseries\slshape rs})\index{State@\texttt{State}}
\label{State}
}\hfill{\scriptsize (operation)}}\\
\noindent\textcolor{FuncColor}{$\triangleright$\enspace\texttt{Reset({\mdseries\slshape rs[, seed]})\index{Reset@\texttt{Reset}}
\label{Reset}
}\hfill{\scriptsize (operation)}}\\


 These are the basic operations for random sources (see \texttt{IsRandomSource} (\ref{IsRandomSource})). 

 \texttt{State} returns a data structure which admits recovering the state of the random
source such that a sequence of random calls using this random source can be
reproduced. If a random source cannot be reset (say, it uses truly random
physical data) then \texttt{State} returns \texttt{fail}. 

 \texttt{Reset( \mbox{\texttt{\mdseries\slshape rs}}, \mbox{\texttt{\mdseries\slshape seed}} )} resets the random source \mbox{\texttt{\mdseries\slshape rs}} to a state described by \mbox{\texttt{\mdseries\slshape seed}}, if the random source can be reset; otherwise it does nothing. Here \mbox{\texttt{\mdseries\slshape seed}} can be an output of \texttt{State} and then \mbox{\texttt{\mdseries\slshape rs}} gets reset to that state. For historical reasons, random sources accept
integer values as \mbox{\texttt{\mdseries\slshape seed}}. We recommend that new code should not rely on this; always use the output of
a prior call to \texttt{State} as \mbox{\texttt{\mdseries\slshape seed}}, or omit it. Without the \mbox{\texttt{\mdseries\slshape seed}} argument a fixed default seed is used. \texttt{Reset} returns the state of \mbox{\texttt{\mdseries\slshape rs}} before the call. 

 Most methods for \texttt{Random} (\ref{Random:for a list or collection}) in the \textsf{GAP} library that do not take a random source as argument use the \texttt{GlobalMersenneTwister} (\ref{GlobalMersenneTwister}) as random source. It can be reset into a known state as in the following
example. 
\begin{Verbatim}[commandchars=!@|,fontsize=\small,frame=single,label=Example]
  !gapprompt@gap>| !gapinput@seed := Reset(GlobalMersenneTwister);;|
  !gapprompt@gap>| !gapinput@seed = State(GlobalMersenneTwister);|
  true
  !gapprompt@gap>| !gapinput@List([1..10],i->Random(Integers));|
  [ -3, 2, -1, -2, -1, -1, 1, -4, 1, 0 ]
  !gapprompt@gap>| !gapinput@List([1..10],i->Random(Integers));|
  [ -1, -1, -1, 1, -1, 1, -2, -1, -2, 0 ]
  !gapprompt@gap>| !gapinput@Reset(GlobalMersenneTwister, seed);;|
  !gapprompt@gap>| !gapinput@List([1..10],i->Random(Integers));|
  [ -3, 2, -1, -2, -1, -1, 1, -4, 1, 0 ]
\end{Verbatim}
 }

 
\subsection{\textcolor{Chapter }{Kinds of Random Sources}}\label{Kinds of Random Sources}
\logpage{[ 14, 7, 4 ]}
\hyperdef{L}{X7AC96008820FAF1F}{}
{
\noindent\textcolor{FuncColor}{$\triangleright$\enspace\texttt{IsMersenneTwister({\mdseries\slshape rs})\index{IsMersenneTwister@\texttt{IsMersenneTwister}}
\label{IsMersenneTwister}
}\hfill{\scriptsize (Category)}}\\
\noindent\textcolor{FuncColor}{$\triangleright$\enspace\texttt{IsGAPRandomSource({\mdseries\slshape rs})\index{IsGAPRandomSource@\texttt{IsGAPRandomSource}}
\label{IsGAPRandomSource}
}\hfill{\scriptsize (Category)}}\\
\noindent\textcolor{FuncColor}{$\triangleright$\enspace\texttt{IsGlobalRandomSource({\mdseries\slshape rs})\index{IsGlobalRandomSource@\texttt{IsGlobalRandomSource}}
\label{IsGlobalRandomSource}
}\hfill{\scriptsize (Category)}}\\
\noindent\textcolor{FuncColor}{$\triangleright$\enspace\texttt{GlobalMersenneTwister\index{GlobalMersenneTwister@\texttt{GlobalMersenneTwister}}
\label{GlobalMersenneTwister}
}\hfill{\scriptsize (global variable)}}\\
\noindent\textcolor{FuncColor}{$\triangleright$\enspace\texttt{GlobalRandomSource\index{GlobalRandomSource@\texttt{GlobalRandomSource}}
\label{GlobalRandomSource}
}\hfill{\scriptsize (global variable)}}\\


 Currently, the \textsf{GAP} library provides three types of random sources, distinguished by the three
listed categories. 

 \texttt{IsMersenneTwister} are random sources which use a fast random generator of 32 bit numbers, called
the Mersenne twister. The pseudo random sequence has a period of $2^{19937}-1$ and the numbers have a $623$\texttt{\symbol{45}}dimensional equidistribution. For more details and the
origin of the code used in the \textsf{GAP} kernel, see: \href{http://www.math.sci.hiroshima-u.ac.jp/~m-mat/MT/emt.html} {\texttt{http://www.math.sci.hiroshima\texttt{\symbol{45}}u.ac.jp/\texttt{\symbol{126}}m\texttt{\symbol{45}}mat/MT/emt.html}}. 

 Use the Mersenne twister if possible, in particular for generating many large
random integers. 

 There is also a predefined global random source \texttt{GlobalMersenneTwister} which is used as the default random source by those library methods for \texttt{Random} (\ref{Random:for a list or collection}) that do not take a random source as an argument. 

 \texttt{IsGAPRandomSource} uses the same number generator as \texttt{IsGlobalRandomSource}, but you can create several of these random sources which generate their
random numbers independently of all other random sources. 

 \texttt{IsGlobalRandomSource} gives access to the \emph{classical} global random generator which was used by \textsf{GAP} in former releases. You do not need to construct new random sources of this
kind which would all use the same global data structure. Just use the existing
random source \texttt{GlobalRandomSource}. This uses the additive random number generator described in \cite{TACP2} (Algorithm A in{\nobreakspace}3.2.2 with lag $30$). 

 Other kinds of random sources are implemented by \textsf{GAP} packages. }

 

\subsection{\textcolor{Chapter }{RandomSource}}
\logpage{[ 14, 7, 5 ]}\nobreak
\hyperdef{L}{X7CB0B5BC82F8FD8F}{}
{\noindent\textcolor{FuncColor}{$\triangleright$\enspace\texttt{RandomSource({\mdseries\slshape cat[, seed]})\index{RandomSource@\texttt{RandomSource}}
\label{RandomSource}
}\hfill{\scriptsize (operation)}}\\


 This operation is used to create new random sources. The first argument \mbox{\texttt{\mdseries\slshape cat}} is the category describing the type of the random generator, for example one
of the categories listed in Section \ref{Kinds of Random Sources}. 

 An optional \mbox{\texttt{\mdseries\slshape seed}} can be given to specify the initial state. For details, see Section \ref{State and Reset for Random Sources}. 

 
\begin{Verbatim}[commandchars=!@|,fontsize=\small,frame=single,label=Example]
  !gapprompt@gap>| !gapinput@rs1 := RandomSource(IsMersenneTwister);|
  <RandomSource in IsMersenneTwister>
  !gapprompt@gap>| !gapinput@l1 := List([1..10000], i-> Random(rs1, [1..6]));;|
  !gapprompt@gap>| !gapinput@state1 := State(rs1);;|
  !gapprompt@gap>| !gapinput@rs2 := RandomSource(IsMersenneTwister);;|
  !gapprompt@gap>| !gapinput@l2 := List([1..10000], i-> Random(rs2, [1..6]));;|
  !gapprompt@gap>| !gapinput@l1 = l2;|
  true
  !gapprompt@gap>| !gapinput@l3 := List([1..10000], i-> Random(rs1, [1..6]));;|
  !gapprompt@gap>| !gapinput@l1 = l3;|
  false
  !gapprompt@gap>| !gapinput@rs3 := RandomSource(IsMersenneTwister, state1);;|
  !gapprompt@gap>| !gapinput@l4 := List([1..10000], i-> Random(rs3, [1..6]));;|
  !gapprompt@gap>| !gapinput@l3 = l4;|
  true
  !gapprompt@gap>| !gapinput@n := Random(rs1, 1, 2^220);|
  1077726777923092117987668044202944212469136000816111066409337432400
\end{Verbatim}
 }

 
\subsection{\textcolor{Chapter }{Implementing new kinds of random sources}}\label{Implementing new kinds of random sources}
\logpage{[ 14, 7, 6 ]}
\hyperdef{L}{X8653AE447D94C1DC}{}
{
  If one wants to implement a new kind of random sources then the first step is
the declaration of a new category \texttt{C}, say, that implies \texttt{IsRandomSource} (\ref{IsRandomSource}), analogous to the categories listed in Section \ref{Kinds of Random Sources}, as follows. 

 \texttt{DeclareCategory( "C", IsRandomSource );}. 

 Then the following method installations are needed. 

 \index{Init@\texttt{Init} (initialize a random source object)} \texttt{InstallMethod( Init, [ C, IsObject ], function( prers, seed ) ... end );} 

 Here \texttt{prers} is an empty component object (which has already the filter \texttt{C}), and \texttt{seed} is an integer or a state value as returned by \texttt{State} (\ref{State}) that describes the initial state of the random source. The function should
fill in the actual data and then return the (now initialized) object \texttt{prers}. The default used for \texttt{seed} is the integer \texttt{1}. A given state value need not be copied by the function. 

 \texttt{InstallMethod( Random, [ C, IsInt, IsInt ], function( rs, low, high ) ... end
);} 

 Here \texttt{rs} is an already initialized random source object in the category \texttt{C}, and the function returns an integer between \texttt{low} and \texttt{high} (inclusive). It is not defined what happens when \texttt{low} is larger than \texttt{high}. 

 \texttt{InstallMethod( State, [ C ], function( rs ) ... end );} 

 If \texttt{rs} supports resetting then the function must return an object that describes the
current state of \texttt{rs}. This object must be an independent copy, that is, calling \texttt{Random} (\ref{Random:for random source and list}) for \texttt{rs} must not change the object that was returned by \texttt{State} (\ref{State}); otherwise \texttt{ReturnFail} (\ref{ReturnFail}) should be installed. 

 \texttt{InstallMethod( Reset, [ C, IsObject ], function( rs, seed ) ... end );} 

 If \texttt{rs} supports resetting then the function must reinitialize \texttt{rs} to the integer or \texttt{State} (\ref{State}) value \texttt{seed} and must return the \texttt{State} (\ref{State}) value of \texttt{rs} before these changes; if resetting is not supported then \texttt{ReturnNothing} (\ref{ReturnNothing}) should be installed. Reset need not copy a given state. Note that the generic
unary \texttt{Reset} (\ref{Reset}) method uses the default seed \texttt{1}. 

 Examples of implementations as described here are given by the random sources
with defining filter \texttt{IsMersenneTwister} (\ref{IsMersenneTwister}) or \texttt{IsRealRandomSource}. (For the latter, see \texttt{RandomSource} (\textbf{IO: RandomSource}) in the \textsf{GAP} package \textsf{IO}.) }

 }

  
\section{\textcolor{Chapter }{Bitfields}}\label{Bitfields}
\logpage{[ 14, 8, 0 ]}
\hyperdef{L}{X7A0311DF78DB4FD8}{}
{
  Bitfields are a low\texttt{\symbol{45}}level feature intended to support
efficient subdivision of immediate integers into bitfields of various widths.
This is typically useful in implementing space\texttt{\symbol{45}}efficient
and/or cache\texttt{\symbol{45}}efficient data structures. This feature should
be used with care because (\emph{inter alia}) it has different limitations on 32\texttt{\symbol{45}}bit and
64\texttt{\symbol{45}}bit architectures. 

\subsection{\textcolor{Chapter }{MakeBitfields}}
\logpage{[ 14, 8, 1 ]}\nobreak
\hyperdef{L}{X85C7BD9E7FCC6C10}{}
{\noindent\textcolor{FuncColor}{$\triangleright$\enspace\texttt{MakeBitfields({\mdseries\slshape width....})\index{MakeBitfields@\texttt{MakeBitfields}}
\label{MakeBitfields}
}\hfill{\scriptsize (function)}}\\


 This function sets up the machinery for a set of bitfields of the given
widths. All bitfield values are treated as unsigned. The total of the widths
must not exceed 60 bits on 64\texttt{\symbol{45}}bit architecture or 28 bits
on a 32\texttt{\symbol{45}}bit architecture. For performance reasons some
checks that one might wish to do are omitted. In particular, the builder and
setter functions do not check if the value[s] passed to them are negative or
too large (unless \textsf{GAP} is specially compiled for debugging). Behaviour when such arguments are passed
is undefined. You can tell which type of architecture you are running on by
accessing \texttt{GAPInfo.BytesPerVariable} which is 8 on 64\texttt{\symbol{45}}bits and 4 on 32. The return value when $n$ widths are given is a record whose fields are 
\begin{description}
\item[{\texttt{widths}}] a copy of the arguments, for convenience,
\item[{\texttt{getters}}]  a list of $n$ functions of one argument each of which extracts one of the fields from an
immediate integer
\item[{\texttt{setters}}]  a list of $n$ functions each taking two arguments: a packed value and a new value for one of
its fields and returning a new packed value. The $i$th function returned the new packed value in which the $i$th field has been replaced by the new value. Note that this does NOT modify
the original packed value.
\end{description}
 Two additional fields may be present if any of the field widths is one. Each
is a list and only has entried bound in the positions corresponding to the
width 1 fields. 
\begin{description}
\item[{\texttt{booleanGetters}}] if the $i$th position of this list is set, it contains a function which extracts the $i$th field (which will have width one) and returns \texttt{true} if it contains 1 and \texttt{false} if it contains 0
\item[{\texttt{booleanSetters}}] if the $i$th position of this list is set, it contains a function of two arguments. The
first argument is a packed value, the second is \texttt{true} or \texttt{false}. It returns a new packed value in which the $i$th field is set to 1 if the second argument was \texttt{true} and 0 if it was \texttt{false}. Behaviour for any other value is undefined.
\end{description}
 }

 

\subsection{\textcolor{Chapter }{BuildBitfields}}
\logpage{[ 14, 8, 2 ]}\nobreak
\hyperdef{L}{X8068CE3781F4003C}{}
{\noindent\textcolor{FuncColor}{$\triangleright$\enspace\texttt{BuildBitfields({\mdseries\slshape widths, vals...})\index{BuildBitfields@\texttt{BuildBitfields}}
\label{BuildBitfields}
}\hfill{\scriptsize (function)}}\\


 This function takes one or more argument. Its first argument is a list of
field widths, as found in the \texttt{widths} entry of a record returned by \texttt{MakeBitfields}. The remaining arguments are unsigned integer values, equal in number to the
entries of the list of field widths. It returns a small integer in which those
entries are packed into bitfields of the given widths. The first entry
occupies the least significant bits. DeclareGlobalFunction("BuildBitfields"); 
\begin{Verbatim}[commandchars=!@|,fontsize=\small,frame=single,label=Example]
  !gapprompt@gap>| !gapinput@bf := MakeBitfields(1,2,3);|
  rec( booleanGetters := [ function( data ) ... end ],
    booleanSetters := [ function( data, val ) ... end ],
    getters := [ function( data ) ... end, function( data ) ... end,
        function( data ) ... end ],
    setters := [ function( data, val ) ... end, function( data, val ) ... end,
        function( data, val ) ... end ], widths := [ 1, 2, 3 ] )
  !gapprompt@gap>| !gapinput@x := BuildBitfields(bf.widths,0,3,5);|
  46
  !gapprompt@gap>| !gapinput@bf.getters[3](x);|
  5
  !gapprompt@gap>| !gapinput@y := bf.setters[1](x,1);|
  47
  !gapprompt@gap>| !gapinput@x;|
  46
  !gapprompt@gap>| !gapinput@bf.booleanGetters[1](x);|
  false
  !gapprompt@gap>| !gapinput@bf.booleanGetters[1](y);|
  true
\end{Verbatim}
 }

 }

 }

          
\chapter{\textcolor{Chapter }{Number Theory}}\label{Number Theory}
\logpage{[ 15, 0, 0 ]}
\hyperdef{L}{X7FB995737B7ED8A2}{}
{
  \index{prime residue group} \textsf{GAP} provides a couple of elementary number theoretic functions. Most of these deal
with the group of integers coprime to $m$, called the \emph{prime residue group}. The order of this group is $\phi(m)$ (see{\nobreakspace}\texttt{Phi} (\ref{Phi})), and $\lambda(m)$ (see{\nobreakspace}\texttt{Lambda} (\ref{Lambda})) is its exponent. This group is cyclic if and only if $m$ is 2, 4, an odd prime power $p^n$, or twice an odd prime power $2 p^n$. In this case the generators of the group, i.e., elements of order $\phi(m)$, are called \emph{primitive roots} (see{\nobreakspace}\texttt{PrimitiveRootMod} (\ref{PrimitiveRootMod})). 

 Note that neither the arguments nor the return values of the functions listed
below are groups or group elements in the sense of \textsf{GAP}. The arguments are simply integers.  
\section{\textcolor{Chapter }{InfoNumtheor (Info Class)}}\label{sect:InfoNumtheor}
\logpage{[ 15, 1, 0 ]}
\hyperdef{L}{X7845C1F97A1742C7}{}
{
  

\subsection{\textcolor{Chapter }{InfoNumtheor}}
\logpage{[ 15, 1, 1 ]}\nobreak
\hyperdef{L}{X796F0DFE7D5D211C}{}
{\noindent\textcolor{FuncColor}{$\triangleright$\enspace\texttt{InfoNumtheor\index{InfoNumtheor@\texttt{InfoNumtheor}}
\label{InfoNumtheor}
}\hfill{\scriptsize (info class)}}\\


 \texttt{InfoNumtheor} is the info class (see{\nobreakspace}\ref{Info Functions}) for the functions in the number theory chapter. }

 }

  
\section{\textcolor{Chapter }{Prime Residues}}\label{Prime Residues}
\logpage{[ 15, 2, 0 ]}
\hyperdef{L}{X823386567DAC22E6}{}
{
  \index{prime residue group} 

\subsection{\textcolor{Chapter }{PrimeResidues}}
\logpage{[ 15, 2, 1 ]}\nobreak
\hyperdef{L}{X7FA3F5347B7004BA}{}
{\noindent\textcolor{FuncColor}{$\triangleright$\enspace\texttt{PrimeResidues({\mdseries\slshape m})\index{PrimeResidues@\texttt{PrimeResidues}}
\label{PrimeResidues}
}\hfill{\scriptsize (function)}}\\


 \texttt{PrimeResidues} returns the set of integers from the range \texttt{[ 0 .. Abs( \mbox{\texttt{\mdseries\slshape m}} )\texttt{\symbol{45}}1 ]} that are coprime to the integer \mbox{\texttt{\mdseries\slshape m}}. 

 \texttt{Abs(\mbox{\texttt{\mdseries\slshape m}})} must be less than $2^{28}$, otherwise the set would probably be too large anyhow. 

 
\begin{Verbatim}[commandchars=!@|,fontsize=\small,frame=single,label=Example]
  !gapprompt@gap>| !gapinput@PrimeResidues( 0 );  PrimeResidues( 1 );  PrimeResidues( 20 );|
  [  ]
  [ 0 ]
  [ 1, 3, 7, 9, 11, 13, 17, 19 ]
\end{Verbatim}
 }

 

\subsection{\textcolor{Chapter }{Phi}}
\logpage{[ 15, 2, 2 ]}\nobreak
\hyperdef{L}{X85A0C67982D9057A}{}
{\noindent\textcolor{FuncColor}{$\triangleright$\enspace\texttt{Phi({\mdseries\slshape m})\index{Phi@\texttt{Phi}}
\label{Phi}
}\hfill{\scriptsize (operation)}}\\


 \index{order!of the prime residue group} \index{prime residue group!order} \index{Euler's totient function} \texttt{Phi} returns the number $\phi(\mbox{\texttt{\mdseries\slshape m}})$ of positive integers less than the positive integer \mbox{\texttt{\mdseries\slshape m}} that are coprime to \mbox{\texttt{\mdseries\slshape m}}. 

 Suppose that $m = p_1^{{e_1}} p_2^{{e_2}} \cdots p_k^{{e_k}}$. Then $\phi(m)$ is $p_1^{{e_1-1}} (p_1-1) p_2^{{e_2-1}} (p_2-1) \cdots p_k^{{e_k-1}} (p_k-1)$. 
\begin{Verbatim}[commandchars=!@|,fontsize=\small,frame=single,label=Example]
  !gapprompt@gap>| !gapinput@Phi( 12 );|
  4
  !gapprompt@gap>| !gapinput@Phi( 2^13-1 );  # this proves that 2^(13)-1 is a prime|
  8190
  !gapprompt@gap>| !gapinput@Phi( 2^15-1 );|
  27000
\end{Verbatim}
 }

 

\subsection{\textcolor{Chapter }{Lambda}}
\logpage{[ 15, 2, 3 ]}\nobreak
\hyperdef{L}{X85296F3087611B03}{}
{\noindent\textcolor{FuncColor}{$\triangleright$\enspace\texttt{Lambda({\mdseries\slshape m})\index{Lambda@\texttt{Lambda}}
\label{Lambda}
}\hfill{\scriptsize (operation)}}\\


 \index{Carmichael's lambda function} \index{prime residue group!exponent} \index{exponent!of the prime residue group} \texttt{Lambda} returns the exponent $\lambda(\mbox{\texttt{\mdseries\slshape m}})$ of the group of prime residues modulo the integer \mbox{\texttt{\mdseries\slshape m}}. 

 $\lambda(\mbox{\texttt{\mdseries\slshape m}})$ is the smallest positive integer $l$ such that for every $a$ relatively prime to \mbox{\texttt{\mdseries\slshape m}} we have $a^l \equiv 1 \pmod{\mbox{\texttt{\mdseries\slshape m}}}$. Fermat's theorem asserts $a^{{\phi(\mbox{\texttt{\mdseries\slshape m}})}} \equiv 1 \pmod{\mbox{\texttt{\mdseries\slshape m}}}$; thus $\lambda(\mbox{\texttt{\mdseries\slshape m}})$ divides $\phi(\mbox{\texttt{\mdseries\slshape m}})$ (see{\nobreakspace}\texttt{Phi} (\ref{Phi})). 

 Carmichael's theorem states that $\lambda$ can be computed as follows: $\lambda(2) = 1$, $\lambda(4) = 2$ and $\lambda(2^e) = 2^{{e-2}}$ if $3 \leq e$, $\lambda(p^e) = (p-1) p^{{e-1}}$ (i.e. $\phi(m)$) if $p$ is an odd prime and $\lambda(m*n) = $\texttt{Lcm}$( \lambda(m), \lambda(n) )$ if $m, n$ are coprime. 

 Composites for which $\lambda(m)$ divides $m - 1$ are called Carmichaels. If $6k+1$, $12k+1$ and $18k+1$ are primes their product is such a number. There are only 1547 Carmichaels
below $10^{10}$ but 455052511 primes. 
\begin{Verbatim}[commandchars=!@|,fontsize=\small,frame=single,label=Example]
  !gapprompt@gap>| !gapinput@Lambda( 10 );|
  4
  !gapprompt@gap>| !gapinput@Lambda( 30 );|
  4
  !gapprompt@gap>| !gapinput@Lambda( 561 );  # 561 is the smallest Carmichael number|
  80
\end{Verbatim}
 }

 

\subsection{\textcolor{Chapter }{GeneratorsPrimeResidues}}
\logpage{[ 15, 2, 4 ]}\nobreak
\hyperdef{L}{X7D191CF67E5018BE}{}
{\noindent\textcolor{FuncColor}{$\triangleright$\enspace\texttt{GeneratorsPrimeResidues({\mdseries\slshape n})\index{GeneratorsPrimeResidues@\texttt{GeneratorsPrimeResidues}}
\label{GeneratorsPrimeResidues}
}\hfill{\scriptsize (function)}}\\


 Let \mbox{\texttt{\mdseries\slshape n}} be a positive integer. \texttt{GeneratorsPrimeResidues} returns a description of generators of the group of prime residues modulo \mbox{\texttt{\mdseries\slshape n}}. The return value is a record with components 
\begin{description}
\item[{\texttt{primes}: }]  a list of the prime factors of \mbox{\texttt{\mdseries\slshape n}}, 
\item[{\texttt{exponents}: }]  a list of the exponents of these primes in the factorization of \mbox{\texttt{\mdseries\slshape n}}, and 
\item[{\texttt{generators}: }]  a list describing generators of the group of prime residues; for the prime
factor $2$, either a primitive root or a list of two generators is stored, for each
other prime factor of \mbox{\texttt{\mdseries\slshape n}}, a primitive root is stored. 
\end{description}
 
\begin{Verbatim}[commandchars=!@|,fontsize=\small,frame=single,label=Example]
  !gapprompt@gap>| !gapinput@GeneratorsPrimeResidues( 1 );|
  rec( exponents := [  ], generators := [  ], primes := [  ] )
  !gapprompt@gap>| !gapinput@GeneratorsPrimeResidues( 4*3 );|
  rec( exponents := [ 2, 1 ], generators := [ 7, 5 ],
    primes := [ 2, 3 ] )
  !gapprompt@gap>| !gapinput@GeneratorsPrimeResidues( 8*9*5 );|
  rec( exponents := [ 3, 2, 1 ],
    generators := [ [ 271, 181 ], 281, 217 ], primes := [ 2, 3, 5 ] )
\end{Verbatim}
 }

 }

  
\section{\textcolor{Chapter }{Primitive Roots and Discrete Logarithms}}\label{Primitive Roots and Discrete Logarithms}
\logpage{[ 15, 3, 0 ]}
\hyperdef{L}{X83103A5385821BAE}{}
{
  

\subsection{\textcolor{Chapter }{OrderMod}}
\logpage{[ 15, 3, 1 ]}\nobreak
\hyperdef{L}{X82373F3D8277EE9E}{}
{\noindent\textcolor{FuncColor}{$\triangleright$\enspace\texttt{OrderMod({\mdseries\slshape n, m[, bound]})\index{OrderMod@\texttt{OrderMod}}
\label{OrderMod}
}\hfill{\scriptsize (function)}}\\


 \index{multiplicative order of an integer} \texttt{OrderMod} returns the multiplicative order of the integer \mbox{\texttt{\mdseries\slshape n}} modulo the positive integer \mbox{\texttt{\mdseries\slshape m}}. If \mbox{\texttt{\mdseries\slshape n}} and \mbox{\texttt{\mdseries\slshape m}} are not coprime the order of \mbox{\texttt{\mdseries\slshape n}} is not defined and \texttt{OrderMod} will return \texttt{0}. 

 If \mbox{\texttt{\mdseries\slshape n}} and \mbox{\texttt{\mdseries\slshape m}} are relatively prime the multiplicative order of \mbox{\texttt{\mdseries\slshape n}} modulo \mbox{\texttt{\mdseries\slshape m}} is the smallest positive integer $i$ such that $\mbox{\texttt{\mdseries\slshape n}}^i \equiv 1 \pmod{\mbox{\texttt{\mdseries\slshape m}}}$. If the group of prime residues modulo \mbox{\texttt{\mdseries\slshape m}} is cyclic then each element of maximal order is called a primitive root modulo \mbox{\texttt{\mdseries\slshape m}} (see{\nobreakspace}\texttt{IsPrimitiveRootMod} (\ref{IsPrimitiveRootMod})). 

 If no a priori known multiple \mbox{\texttt{\mdseries\slshape bound}} of the desired order is given, \texttt{OrderMod} usually spends most of its time factoring \mbox{\texttt{\mdseries\slshape m}} for computing $\lambda(\mbox{\texttt{\mdseries\slshape m}})$ (see \texttt{Lambda} (\ref{Lambda})) as the default for \mbox{\texttt{\mdseries\slshape bound}}, and then factoring \mbox{\texttt{\mdseries\slshape bound}} (see{\nobreakspace}\texttt{FactorsInt} (\ref{FactorsInt})). 

 If an incorrect \mbox{\texttt{\mdseries\slshape bound}} is given then the result will be wrong. 
\begin{Verbatim}[commandchars=!@|,fontsize=\small,frame=single,label=Example]
  !gapprompt@gap>| !gapinput@OrderMod( 2, 7 );|
  3
  !gapprompt@gap>| !gapinput@OrderMod( 3, 7 );  # 3 is a primitive root modulo 7|
  6
  !gapprompt@gap>| !gapinput@m:= (5^166-1) / 167;;   # about 10^113|
  !gapprompt@gap>| !gapinput@OrderMod( 5, m, 166 );  # needs minutes without third argument|
  166
\end{Verbatim}
 }

 

\subsection{\textcolor{Chapter }{LogMod}}
\logpage{[ 15, 3, 2 ]}\nobreak
\hyperdef{L}{X81AD9C7779A7BA89}{}
{\noindent\textcolor{FuncColor}{$\triangleright$\enspace\texttt{LogMod({\mdseries\slshape n, r, m})\index{LogMod@\texttt{LogMod}}
\label{LogMod}
}\hfill{\scriptsize (function)}}\\
\noindent\textcolor{FuncColor}{$\triangleright$\enspace\texttt{LogModShanks({\mdseries\slshape n, r, m})\index{LogModShanks@\texttt{LogModShanks}}
\label{LogModShanks}
}\hfill{\scriptsize (function)}}\\


 \index{logarithm!discrete} computes the discrete \mbox{\texttt{\mdseries\slshape r}}\texttt{\symbol{45}}logarithm of the integer \mbox{\texttt{\mdseries\slshape n}} modulo the integer \mbox{\texttt{\mdseries\slshape m}}. It returns a number \mbox{\texttt{\mdseries\slshape l}} such that $\mbox{\texttt{\mdseries\slshape r}}^{\mbox{\texttt{\mdseries\slshape l}}} \equiv \mbox{\texttt{\mdseries\slshape n}} \pmod{\mbox{\texttt{\mdseries\slshape m}}}$ if such a number exists. Otherwise \texttt{fail} is returned. 

 \texttt{LogModShanks} uses the Baby Step \texttt{\symbol{45}} Giant Step Method of Shanks (see for
example \cite[section 5.4.1]{Coh93}) and in general requires more memory than a call to \texttt{LogMod}. 
\begin{Verbatim}[commandchars=!@|,fontsize=\small,frame=single,label=Example]
  !gapprompt@gap>| !gapinput@l:= LogMod( 2, 5, 7 );  5^l mod 7 = 2;|
  4
  true
  !gapprompt@gap>| !gapinput@LogMod( 1, 3, 3 );  LogMod( 2, 3, 3 );|
  0
  fail
\end{Verbatim}
 }

 

\subsection{\textcolor{Chapter }{DLog}}
\logpage{[ 15, 3, 3 ]}\nobreak
\hyperdef{L}{X84A138947E8C49A8}{}
{\noindent\textcolor{FuncColor}{$\triangleright$\enspace\texttt{DLog({\mdseries\slshape base, x[, m]})\index{DLog@\texttt{DLog}}
\label{DLog}
}\hfill{\scriptsize (function)}}\\
\textbf{\indent Returns:\ }
an integer



 \index{logarithm!discrete} The argument \mbox{\texttt{\mdseries\slshape base}} must be a multiplicative element and \mbox{\texttt{\mdseries\slshape x}} must lie in the cyclic group generated by \mbox{\texttt{\mdseries\slshape base}}. The third argument \mbox{\texttt{\mdseries\slshape m}} must be the order of \mbox{\texttt{\mdseries\slshape base}} or its factorization. If \mbox{\texttt{\mdseries\slshape m}} is not given, it is computed first. This function returns the discrete
logarithm, that is an integer $e$ such that \mbox{\texttt{\mdseries\slshape base}}$^e = $ \mbox{\texttt{\mdseries\slshape x}}. 

 If \mbox{\texttt{\mdseries\slshape m}} is prime then Shanks' algorithm is used (which needs $O(\sqrt{\mbox{\texttt{\mdseries\slshape m}}})$ space and time). Otherwise let \mbox{\texttt{\mdseries\slshape m}} $ = r l$ and $e = a + b r$ with $0 \leq a < r$. Then $a =$ \texttt{DLog}$(\mbox{\texttt{\mdseries\slshape base}}^l, \mbox{\texttt{\mdseries\slshape x}}^l, r)$ and $b = $ \texttt{DLog}$(\mbox{\texttt{\mdseries\slshape base}}^r, \mbox{\texttt{\mdseries\slshape x}}/\mbox{\texttt{\mdseries\slshape base}}^a, l)$. 

 This function is used for a method of \texttt{LogFFE} (\ref{LogFFE}). 
\begin{Verbatim}[commandchars=!@|,fontsize=\small,frame=single,label=Example]
  !gapprompt@gap>| !gapinput@q:= 67^12;|
  8182718904632857144561
  !gapprompt@gap>| !gapinput@z:= Z(q);;|
  !gapprompt@gap>| !gapinput@DLog(z, z+1);|
  2874413785388345993274
  !gapprompt@gap>| !gapinput@DLog(z, z^2+1);|
  1667375214152688471247
  !gapprompt@gap>| !gapinput@DLog(z, Z(67));|
  123980589464134199160
\end{Verbatim}
 }

 

\subsection{\textcolor{Chapter }{PrimitiveRootMod}}
\logpage{[ 15, 3, 4 ]}\nobreak
\hyperdef{L}{X82440BB9812FF148}{}
{\noindent\textcolor{FuncColor}{$\triangleright$\enspace\texttt{PrimitiveRootMod({\mdseries\slshape m[, start]})\index{PrimitiveRootMod@\texttt{PrimitiveRootMod}}
\label{PrimitiveRootMod}
}\hfill{\scriptsize (function)}}\\


 \index{primitive root modulo an integer} \index{prime residue group!generator} \index{generator!of the prime residue group} \texttt{PrimitiveRootMod} returns the smallest primitive root modulo the positive integer \mbox{\texttt{\mdseries\slshape m}} and \texttt{fail} if no such primitive root exists. If the optional second integer argument \mbox{\texttt{\mdseries\slshape start}} is given \texttt{PrimitiveRootMod} returns the smallest primitive root that is strictly larger than \mbox{\texttt{\mdseries\slshape start}}. 
\begin{Verbatim}[commandchars=!@|,fontsize=\small,frame=single,label=Example]
  !gapprompt@gap>| !gapinput@# largest primitive root for a prime less than 2000:|
  !gapprompt@gap>| !gapinput@PrimitiveRootMod( 409 );|
  21
  !gapprompt@gap>| !gapinput@PrimitiveRootMod( 541, 2 );|
  10
  !gapprompt@gap>| !gapinput@# 327 is the largest primitive root mod 337:|
  !gapprompt@gap>| !gapinput@PrimitiveRootMod( 337, 327 );|
  fail
  !gapprompt@gap>| !gapinput@# there exists no primitive root modulo 30:|
  !gapprompt@gap>| !gapinput@PrimitiveRootMod( 30 );|
  fail
\end{Verbatim}
 }

 

\subsection{\textcolor{Chapter }{IsPrimitiveRootMod}}
\logpage{[ 15, 3, 5 ]}\nobreak
\hyperdef{L}{X790466C07BD90E20}{}
{\noindent\textcolor{FuncColor}{$\triangleright$\enspace\texttt{IsPrimitiveRootMod({\mdseries\slshape r, m})\index{IsPrimitiveRootMod@\texttt{IsPrimitiveRootMod}}
\label{IsPrimitiveRootMod}
}\hfill{\scriptsize (function)}}\\


 \index{test!for a primitive root} \index{prime residue group!generator} \index{generator!of the prime residue group} \texttt{IsPrimitiveRootMod} returns \texttt{true} if the integer \mbox{\texttt{\mdseries\slshape r}} is a primitive root modulo the positive integer \mbox{\texttt{\mdseries\slshape m}}, and \texttt{false} otherwise. If \mbox{\texttt{\mdseries\slshape r}} is less than 0 or larger than \mbox{\texttt{\mdseries\slshape m}} it is replaced by its remainder. 
\begin{Verbatim}[commandchars=!@|,fontsize=\small,frame=single,label=Example]
  !gapprompt@gap>| !gapinput@IsPrimitiveRootMod( 2, 541 );|
  true
  !gapprompt@gap>| !gapinput@IsPrimitiveRootMod( -539, 541 );  # same computation as above;|
  true
  !gapprompt@gap>| !gapinput@IsPrimitiveRootMod( 4, 541 );|
  false
  !gapprompt@gap>| !gapinput@ForAny( [1..29], r -> IsPrimitiveRootMod( r, 30 ) );|
  false
  !gapprompt@gap>| !gapinput@# there is no a primitive root modulo 30|
\end{Verbatim}
 }

 }

  
\section{\textcolor{Chapter }{Roots Modulo Integers}}\label{Roots Modulo Integers}
\logpage{[ 15, 4, 0 ]}
\hyperdef{L}{X7F9069D77AC48054}{}
{
  

\subsection{\textcolor{Chapter }{Jacobi}}
\logpage{[ 15, 4, 1 ]}\nobreak
\hyperdef{L}{X83449DBC80495971}{}
{\noindent\textcolor{FuncColor}{$\triangleright$\enspace\texttt{Jacobi({\mdseries\slshape n, m})\index{Jacobi@\texttt{Jacobi}}
\label{Jacobi}
}\hfill{\scriptsize (function)}}\\


 \index{quadratic residue} \index{residue!quadratic} \texttt{Jacobi} returns the value of the \emph{Kronecker\texttt{\symbol{45}}Jacobi symbol} $J(\mbox{\texttt{\mdseries\slshape n}},\mbox{\texttt{\mdseries\slshape m}})$ of the integer \mbox{\texttt{\mdseries\slshape n}} modulo the integer \mbox{\texttt{\mdseries\slshape m}}. It is defined as follows: 

 If $n$ and $m$ are not coprime then $J(n,m) = 0$. Furthermore, $J(n,1) = 1$ and $J(n,-1) = -1$ if $m < 0$ and $+1$ otherwise. And for odd $n$ it is $J(n,2) = (-1)^k$ with $k = (n^2-1)/8$. For odd primes $m$ which are coprime to $n$ the Kronecker\texttt{\symbol{45}}Jacobi symbol has the same value as the
Legendre symbol (see{\nobreakspace}\texttt{Legendre} (\ref{Legendre})). 

 For the general case suppose that $m = p_1 \cdot p_2 \cdots p_k$ is a product of $-1$ and of primes, not necessarily distinct, and that $n$ is coprime to $m$. Then $J(n,m) = J(n,p_1) \cdot J(n,p_2) \cdots J(n,p_k)$. 

 Note that the Kronecker\texttt{\symbol{45}}Jacobi symbol coincides with the
Jacobi symbol that is defined for odd $m$ in many number theory books. For odd primes $m$ and $n$ coprime to $m$ it coincides with the Legendre symbol. 

 \texttt{Jacobi} is very efficient, even for large values of \mbox{\texttt{\mdseries\slshape n}} and \mbox{\texttt{\mdseries\slshape m}}, it is about as fast as the Euclidean algorithm (see{\nobreakspace}\texttt{Gcd} (\ref{Gcd:for (a ring and) several elements})). 
\begin{Verbatim}[commandchars=!@|,fontsize=\small,frame=single,label=Example]
  !gapprompt@gap>| !gapinput@Jacobi( 11, 35 );  # 9^2 = 11 mod 35|
  1
  !gapprompt@gap>| !gapinput@# this is -1, thus there is no r such that r^2 = 6 mod 35|
  !gapprompt@gap>| !gapinput@Jacobi( 6, 35 );|
  -1
  !gapprompt@gap>| !gapinput@# this is 1 even though there is no r with r^2 = 3 mod 35|
  !gapprompt@gap>| !gapinput@Jacobi( 3, 35 );|
  1
\end{Verbatim}
 }

 

\subsection{\textcolor{Chapter }{Legendre}}
\logpage{[ 15, 4, 2 ]}\nobreak
\hyperdef{L}{X81464ABF7F10E544}{}
{\noindent\textcolor{FuncColor}{$\triangleright$\enspace\texttt{Legendre({\mdseries\slshape n, m})\index{Legendre@\texttt{Legendre}}
\label{Legendre}
}\hfill{\scriptsize (function)}}\\


 \index{quadratic residue} \index{residue!quadratic} \texttt{Legendre} returns the value of the \emph{Legendre symbol} of the integer \mbox{\texttt{\mdseries\slshape n}} modulo the positive integer \mbox{\texttt{\mdseries\slshape m}}. 

 The value of the Legendre symbol $L(n/m)$ is 1 if $n$ is a \emph{quadratic residue} modulo $m$, i.e., if there exists an integer $r$ such that $r^2 \equiv n \pmod{m}$ and $-1$ otherwise. 

 If a root of \mbox{\texttt{\mdseries\slshape n}} exists it can be found by \texttt{RootMod} (\ref{RootMod}). 

 While the value of the Legendre symbol usually is only defined for \mbox{\texttt{\mdseries\slshape m}} a prime, we have extended the definition to include composite moduli too. The
Jacobi symbol (see \texttt{Jacobi} (\ref{Jacobi})) is another generalization of the Legendre symbol for composite moduli that
is much cheaper to compute, because it does not need the factorization of \mbox{\texttt{\mdseries\slshape m}} (see \texttt{FactorsInt} (\ref{FactorsInt})). 

 A description of the Jacobi symbol, the Legendre symbol, and related topics
can be found in \cite{Baker84}. 
\begin{Verbatim}[commandchars=!@|,fontsize=\small,frame=single,label=Example]
  !gapprompt@gap>| !gapinput@Legendre( 5, 11 );  # 4^2 = 5 mod 11|
  1
  !gapprompt@gap>| !gapinput@# this is -1, thus there is no r such that r^2 = 6 mod 11|
  !gapprompt@gap>| !gapinput@Legendre( 6, 11 );|
  -1
  !gapprompt@gap>| !gapinput@# this is -1, thus there is no r such that r^2 = 3 mod 35|
  !gapprompt@gap>| !gapinput@Legendre( 3, 35 );|
  -1
\end{Verbatim}
 }

 

\subsection{\textcolor{Chapter }{RootMod}}
\logpage{[ 15, 4, 3 ]}\nobreak
\hyperdef{L}{X83E3ED577B7A04ED}{}
{\noindent\textcolor{FuncColor}{$\triangleright$\enspace\texttt{RootMod({\mdseries\slshape n[, k], m})\index{RootMod@\texttt{RootMod}}
\label{RootMod}
}\hfill{\scriptsize (function)}}\\


 \index{quadratic residue} \index{residue!quadratic} \index{root!of an integer modulo another} \texttt{RootMod} computes a \mbox{\texttt{\mdseries\slshape k}}th root of the integer \mbox{\texttt{\mdseries\slshape n}} modulo the positive integer \mbox{\texttt{\mdseries\slshape m}}, i.e., a $r$ such that $r^{\mbox{\texttt{\mdseries\slshape k}}} \equiv \mbox{\texttt{\mdseries\slshape n}} \pmod{\mbox{\texttt{\mdseries\slshape m}}}$. If no such root exists \texttt{RootMod} returns \texttt{fail}. If only the arguments \mbox{\texttt{\mdseries\slshape n}} and \mbox{\texttt{\mdseries\slshape m}} are given, the default value for \mbox{\texttt{\mdseries\slshape k}} is $2$. 

 A square root of \mbox{\texttt{\mdseries\slshape n}} exists only if \texttt{Legendre(\mbox{\texttt{\mdseries\slshape n}},\mbox{\texttt{\mdseries\slshape m}}) = 1} (see{\nobreakspace}\texttt{Legendre} (\ref{Legendre})). If \mbox{\texttt{\mdseries\slshape m}} has $r$ different prime factors then there are $2^r$ different roots of \mbox{\texttt{\mdseries\slshape n}} mod \mbox{\texttt{\mdseries\slshape m}}. It is unspecified which one \texttt{RootMod} returns. You can, however, use \texttt{RootsMod} (\ref{RootsMod}) to compute the full set of roots. 

 \texttt{RootMod} is efficient even for large values of \mbox{\texttt{\mdseries\slshape m}}, in fact the most time is usually spent factoring \mbox{\texttt{\mdseries\slshape m}} (see \texttt{FactorsInt} (\ref{FactorsInt})). 
\begin{Verbatim}[commandchars=!@|,fontsize=\small,frame=single,label=Example]
  !gapprompt@gap>| !gapinput@# note 'RootMod' does not return 8 in this case but -8:|
  !gapprompt@gap>| !gapinput@RootMod( 64, 1009 );|
  1001
  !gapprompt@gap>| !gapinput@RootMod( 64, 3, 1009 );|
  518
  !gapprompt@gap>| !gapinput@RootMod( 64, 5, 1009 );|
  656
  !gapprompt@gap>| !gapinput@List( RootMod( 64, 1009 ) * RootsUnityMod( 1009 ),|
  !gapprompt@>| !gapinput@      x -> x mod 1009 );  # set of all square roots of 64 mod 1009|
  [ 1001, 8 ]
\end{Verbatim}
 }

 

\subsection{\textcolor{Chapter }{RootsMod}}
\logpage{[ 15, 4, 4 ]}\nobreak
\hyperdef{L}{X84D3F03B862841F8}{}
{\noindent\textcolor{FuncColor}{$\triangleright$\enspace\texttt{RootsMod({\mdseries\slshape n[, k], m})\index{RootsMod@\texttt{RootsMod}}
\label{RootsMod}
}\hfill{\scriptsize (function)}}\\


 \texttt{RootsMod} computes the set of \mbox{\texttt{\mdseries\slshape k}}th roots of the integer \mbox{\texttt{\mdseries\slshape n}} modulo the positive integer \mbox{\texttt{\mdseries\slshape m}}, i.e., the list of all $r$ such that $r^{\mbox{\texttt{\mdseries\slshape k}}} \equiv \mbox{\texttt{\mdseries\slshape n}} \pmod{\mbox{\texttt{\mdseries\slshape m}}}$. If only the arguments \mbox{\texttt{\mdseries\slshape n}} and \mbox{\texttt{\mdseries\slshape m}} are given, the default value for \mbox{\texttt{\mdseries\slshape k}} is $2$. 
\begin{Verbatim}[commandchars=!@|,fontsize=\small,frame=single,label=Example]
  !gapprompt@gap>| !gapinput@RootsMod( 1, 7*31 );  # the same as `RootsUnityMod( 7*31 )'|
  [ 1, 92, 125, 216 ]
  !gapprompt@gap>| !gapinput@RootsMod( 7, 7*31 );|
  [ 21, 196 ]
  !gapprompt@gap>| !gapinput@RootsMod( 5, 7*31 );|
  [  ]
  !gapprompt@gap>| !gapinput@RootsMod( 1, 5, 7*31 );|
  [ 1, 8, 64, 78, 190 ]
\end{Verbatim}
 }

 

\subsection{\textcolor{Chapter }{RootsUnityMod}}
\logpage{[ 15, 4, 5 ]}\nobreak
\hyperdef{L}{X81F856E682A8ECBA}{}
{\noindent\textcolor{FuncColor}{$\triangleright$\enspace\texttt{RootsUnityMod({\mdseries\slshape [k, ]m})\index{RootsUnityMod@\texttt{RootsUnityMod}}
\label{RootsUnityMod}
}\hfill{\scriptsize (function)}}\\


 \index{modular roots} \index{root!of 1 modulo an integer} \texttt{RootsUnityMod} returns the set of \mbox{\texttt{\mdseries\slshape k}}\texttt{\symbol{45}}th roots of unity modulo the positive integer \mbox{\texttt{\mdseries\slshape m}}, i.e., the list of all solutions $r$ of $r^{\mbox{\texttt{\mdseries\slshape k}}} \equiv \mbox{\texttt{\mdseries\slshape n}} \pmod{\mbox{\texttt{\mdseries\slshape m}}}$. If only the argument \mbox{\texttt{\mdseries\slshape m}} is given, the default value for \mbox{\texttt{\mdseries\slshape k}} is $2$. 

 In general there are $\mbox{\texttt{\mdseries\slshape k}}^n$ such roots if the modulus \mbox{\texttt{\mdseries\slshape m}} has $n$ different prime factors $p$ such that $p \equiv 1 \pmod{\mbox{\texttt{\mdseries\slshape k}}}$. If $\mbox{\texttt{\mdseries\slshape k}}^2$ divides \mbox{\texttt{\mdseries\slshape m}} then there are $\mbox{\texttt{\mdseries\slshape k}}^{{n+1}}$ such roots; and especially if $\mbox{\texttt{\mdseries\slshape k}} = 2$ and 8 divides \mbox{\texttt{\mdseries\slshape m}} there are $2^{{n+2}}$ such roots. 

 In the current implementation \mbox{\texttt{\mdseries\slshape k}} must be a prime. 
\begin{Verbatim}[commandchars=!@|,fontsize=\small,frame=single,label=Example]
  !gapprompt@gap>| !gapinput@RootsUnityMod( 7*31 );  RootsUnityMod( 3, 7*31 );|
  [ 1, 92, 125, 216 ]
  [ 1, 25, 32, 36, 67, 149, 156, 191, 211 ]
  !gapprompt@gap>| !gapinput@RootsUnityMod( 5, 7*31 );|
  [ 1, 8, 64, 78, 190 ]
  !gapprompt@gap>| !gapinput@List( RootMod( 64, 1009 ) * RootsUnityMod( 1009 ),|
  !gapprompt@>| !gapinput@         x -> x mod 1009 );  # set of all square roots of 64 mod 1009|
  [ 1001, 8 ]
\end{Verbatim}
 }

 }

  
\section{\textcolor{Chapter }{Multiplicative Arithmetic Functions}}\label{Multiplicative Arithmetic Functions}
\logpage{[ 15, 5, 0 ]}
\hyperdef{L}{X7B3A5A0378A32F83}{}
{
  

\subsection{\textcolor{Chapter }{Sigma}}
\logpage{[ 15, 5, 1 ]}\nobreak
\hyperdef{L}{X823707DF821E79A0}{}
{\noindent\textcolor{FuncColor}{$\triangleright$\enspace\texttt{Sigma({\mdseries\slshape n})\index{Sigma@\texttt{Sigma}}
\label{Sigma}
}\hfill{\scriptsize (operation)}}\\


 \texttt{Sigma} returns the sum of the positive divisors of the nonzero integer \mbox{\texttt{\mdseries\slshape n}}. 

 \texttt{Sigma} is a multiplicative arithmetic function, i.e., if $n$ and $m$ are relatively prime we have that $\sigma(n \cdot m) = \sigma(n) \sigma(m)$. 

 Together with the formula $\sigma(p^k) = (p^{{k+1}}-1) / (p-1)$ this allows us to compute $\sigma(\mbox{\texttt{\mdseries\slshape n}})$. 

 Integers \mbox{\texttt{\mdseries\slshape n}} for which $\sigma(\mbox{\texttt{\mdseries\slshape n}}) = 2 \mbox{\texttt{\mdseries\slshape n}}$ are called perfect. Even perfect integers are exactly of the form $2^{{\mbox{\texttt{\mdseries\slshape n}}-1}}(2^{\mbox{\texttt{\mdseries\slshape n}}}-1)$ where $2^{\mbox{\texttt{\mdseries\slshape n}}}-1$ is prime. Primes of the form $2^{\mbox{\texttt{\mdseries\slshape n}}}-1$ are called \emph{Mersenne primes}, and 42 among the known Mersenne primes are obtained for \mbox{\texttt{\mdseries\slshape n}} $=$ 2, 3, 5, 7, 13, 17, 19, 31, 61, 89, 107, 127, 521, 607, 1279, 2203, 2281,
3217, 4253, 4423, 9689, 9941, 11213, 19937, 21701, 23209, 44497, 86243,
110503, 132049, 216091, 756839, 859433, 1257787, 1398269, 2976221, 3021377,
6972593, 13466917, 20996011, 24036583 and 25964951. Please find more up to
date information about Mersenne primes at \href{https://www.mersenne.org} {\texttt{https://www.mersenne.org}}. It is not known whether odd perfect integers exist, however{\nobreakspace}\cite{BC89} show that any such integer must have at least 300 decimal digits. 

 \texttt{Sigma} usually spends most of its time factoring \mbox{\texttt{\mdseries\slshape n}} (see{\nobreakspace}\texttt{FactorsInt} (\ref{FactorsInt})). 

 
\begin{Verbatim}[commandchars=!@|,fontsize=\small,frame=single,label=Example]
  !gapprompt@gap>| !gapinput@Sigma( 1 );|
  1
  !gapprompt@gap>| !gapinput@Sigma( 1009 );  # 1009 is a prime|
  1010
  !gapprompt@gap>| !gapinput@Sigma( 8128 ) = 2*8128;  # 8128 is a perfect number|
  true
\end{Verbatim}
 }

 

\subsection{\textcolor{Chapter }{Tau}}
\logpage{[ 15, 5, 2 ]}\nobreak
\hyperdef{L}{X798C62847EE0372E}{}
{\noindent\textcolor{FuncColor}{$\triangleright$\enspace\texttt{Tau({\mdseries\slshape n})\index{Tau@\texttt{Tau}}
\label{Tau}
}\hfill{\scriptsize (operation)}}\\


 \texttt{Tau} returns the number of the positive divisors of the nonzero integer \mbox{\texttt{\mdseries\slshape n}}. 

 \texttt{Tau} is a multiplicative arithmetic function, i.e., if $n$ and $m$ are relative prime we have $\tau(n \cdot m) = \tau(n) \tau(m)$. Together with the formula $\tau(p^k) = k+1$ this allows us to compute $\tau(\mbox{\texttt{\mdseries\slshape n}})$. 

 \texttt{Tau} usually spends most of its time factoring \mbox{\texttt{\mdseries\slshape n}} (see{\nobreakspace}\texttt{FactorsInt} (\ref{FactorsInt})). 
\begin{Verbatim}[commandchars=!@|,fontsize=\small,frame=single,label=Example]
  !gapprompt@gap>| !gapinput@Tau( 1 );|
  1
  !gapprompt@gap>| !gapinput@Tau( 1013 );  # thus 1013 is a prime|
  2
  !gapprompt@gap>| !gapinput@Tau( 8128 );|
  14
  !gapprompt@gap>| !gapinput@# result is odd if and only if argument is a perfect square:|
  !gapprompt@gap>| !gapinput@Tau( 36 );|
  9
\end{Verbatim}
 }

 

\subsection{\textcolor{Chapter }{MoebiusMu}}
\logpage{[ 15, 5, 3 ]}\nobreak
\hyperdef{L}{X79C1DA36827C2959}{}
{\noindent\textcolor{FuncColor}{$\triangleright$\enspace\texttt{MoebiusMu({\mdseries\slshape n})\index{MoebiusMu@\texttt{MoebiusMu}}
\label{MoebiusMu}
}\hfill{\scriptsize (function)}}\\


 \texttt{MoebiusMu} computes the value of Moebius inversion function for the nonzero integer \mbox{\texttt{\mdseries\slshape n}}. This is 0 for integers which are not squarefree, i.e., which are divided by
a square $r^2$. Otherwise it is 1 if \mbox{\texttt{\mdseries\slshape n}} has a even number and $-1$ if \mbox{\texttt{\mdseries\slshape n}} has an odd number of prime factors. 

 The importance of $\mu$ stems from the so called inversion formula. Suppose $f$ is a multiplicative arithmetic function defined on the positive integers and
let $g(n) = \sum_{{d \mid n}} f(d)$. Then $f(n) = \sum_{{d \mid n}} \mu(d) g(n/d)$. As a special case we have $\phi(n) = \sum_{{d \mid n}} \mu(d) n/d$ since $n = \sum_{{d \mid n}} \phi(d)$ (see{\nobreakspace}\texttt{Phi} (\ref{Phi})). 

 \texttt{MoebiusMu} usually spends all of its time factoring \mbox{\texttt{\mdseries\slshape n}} (see \texttt{FactorsInt} (\ref{FactorsInt})). 
\begin{Verbatim}[commandchars=!@|,fontsize=\small,frame=single,label=Example]
  !gapprompt@gap>| !gapinput@MoebiusMu( 60 );  MoebiusMu( 61 );  MoebiusMu( 62 );|
  0
  -1
  1
\end{Verbatim}
 }

 }

  
\section{\textcolor{Chapter }{Continued Fractions}}\label{Continued Fractions}
\logpage{[ 15, 6, 0 ]}
\hyperdef{L}{X7B2E061C835159B9}{}
{
  

\subsection{\textcolor{Chapter }{ContinuedFractionExpansionOfRoot}}
\logpage{[ 15, 6, 1 ]}\nobreak
\hyperdef{L}{X874C161B83416092}{}
{\noindent\textcolor{FuncColor}{$\triangleright$\enspace\texttt{ContinuedFractionExpansionOfRoot({\mdseries\slshape f, n})\index{ContinuedFractionExpansionOfRoot@\texttt{ContinuedFractionExpansionOfRoot}}
\label{ContinuedFractionExpansionOfRoot}
}\hfill{\scriptsize (function)}}\\


 The first \mbox{\texttt{\mdseries\slshape n}} terms of the continued fraction expansion of the only positive real root of
the polynomial \mbox{\texttt{\mdseries\slshape f}} with integer coefficients. The leading coefficient of \mbox{\texttt{\mdseries\slshape f}} must be positive and the value of \mbox{\texttt{\mdseries\slshape f}} at 0 must be negative. If the degree of \mbox{\texttt{\mdseries\slshape f}} is 2 and \mbox{\texttt{\mdseries\slshape n}} = 0, the function computes one period of the continued fraction expansion of
the root in question. Anything may happen if \mbox{\texttt{\mdseries\slshape f}} has three or more positive real roots. 
\begin{Verbatim}[commandchars=!@|,fontsize=\small,frame=single,label=Example]
  !gapprompt@gap>| !gapinput@x := Indeterminate(Integers);;|
  !gapprompt@gap>| !gapinput@ContinuedFractionExpansionOfRoot(x^2-7,20);|
  [ 2, 1, 1, 1, 4, 1, 1, 1, 4, 1, 1, 1, 4, 1, 1, 1, 4, 1, 1, 1 ]
  !gapprompt@gap>| !gapinput@ContinuedFractionExpansionOfRoot(x^2-7,0);|
  [ 2, 1, 1, 1, 4 ]
  !gapprompt@gap>| !gapinput@ContinuedFractionExpansionOfRoot(x^3-2,20);|
  [ 1, 3, 1, 5, 1, 1, 4, 1, 1, 8, 1, 14, 1, 10, 2, 1, 4, 12, 2, 3 ]
  !gapprompt@gap>| !gapinput@ContinuedFractionExpansionOfRoot(x^5-x-1,50);|
  [ 1, 5, 1, 42, 1, 3, 24, 2, 2, 1, 16, 1, 11, 1, 1, 2, 31, 1, 12, 5,
    1, 7, 11, 1, 4, 1, 4, 2, 2, 3, 4, 2, 1, 1, 11, 1, 41, 12, 1, 8, 1,
    1, 1, 1, 1, 9, 2, 1, 5, 4 ]
\end{Verbatim}
 }

 

\subsection{\textcolor{Chapter }{ContinuedFractionApproximationOfRoot}}
\logpage{[ 15, 6, 2 ]}\nobreak
\hyperdef{L}{X8059667580A039A6}{}
{\noindent\textcolor{FuncColor}{$\triangleright$\enspace\texttt{ContinuedFractionApproximationOfRoot({\mdseries\slshape f, n})\index{ContinuedFractionApproximationOfRoot@\texttt{Continued}\-\texttt{Fraction}\-\texttt{Approximation}\-\texttt{Of}\-\texttt{Root}}
\label{ContinuedFractionApproximationOfRoot}
}\hfill{\scriptsize (function)}}\\


 The \mbox{\texttt{\mdseries\slshape n}}th continued fraction approximation of the only positive real root of the
polynomial \mbox{\texttt{\mdseries\slshape f}} with integer coefficients. The leading coefficient of \mbox{\texttt{\mdseries\slshape f}} must be positive and the value of \mbox{\texttt{\mdseries\slshape f}} at 0 must be negative. Anything may happen if \mbox{\texttt{\mdseries\slshape f}} has three or more positive real roots. 
\begin{Verbatim}[commandchars=!@|,fontsize=\small,frame=single,label=Example]
  !gapprompt@gap>| !gapinput@ContinuedFractionApproximationOfRoot(x^2-2,10);|
  3363/2378
  !gapprompt@gap>| !gapinput@3363^2-2*2378^2;|
  1
  !gapprompt@gap>| !gapinput@z := ContinuedFractionApproximationOfRoot(x^5-x-1,20);|
  499898783527/428250732317
  !gapprompt@gap>| !gapinput@z^5-z-1;|
  486192462527432755459620441970617283/
  14404247382319842421697357558805709031116987826242631261357
\end{Verbatim}
 }

 }

  
\section{\textcolor{Chapter }{Miscellaneous}}\label{Miscellaneous}
\logpage{[ 15, 7, 0 ]}
\hyperdef{L}{X7C5563A37D566DA5}{}
{
  

\subsection{\textcolor{Chapter }{PValuation}}
\logpage{[ 15, 7, 1 ]}\nobreak
\hyperdef{L}{X8243EAA586D78ED4}{}
{\noindent\textcolor{FuncColor}{$\triangleright$\enspace\texttt{PValuation({\mdseries\slshape n, p})\index{PValuation@\texttt{PValuation}}
\label{PValuation}
}\hfill{\scriptsize (function)}}\\


 For an integer \mbox{\texttt{\mdseries\slshape n}} and a prime \mbox{\texttt{\mdseries\slshape p}} this function returns the \mbox{\texttt{\mdseries\slshape p}}\texttt{\symbol{45}}valuation of \mbox{\texttt{\mdseries\slshape n}}, that is the exponent $e$ such that $p^e$ is the largest power of \mbox{\texttt{\mdseries\slshape p}} that divides \mbox{\texttt{\mdseries\slshape n}}. The valuation of zero is infinity. 
\begin{Verbatim}[commandchars=!@|,fontsize=\small,frame=single,label=Example]
  !gapprompt@gap>| !gapinput@PValuation(100,2);|
  2
  !gapprompt@gap>| !gapinput@PValuation(100,3);|
  0
\end{Verbatim}
 }

 

\subsection{\textcolor{Chapter }{TwoSquares}}
\logpage{[ 15, 7, 2 ]}\nobreak
\hyperdef{L}{X85E1EFC484F648A4}{}
{\noindent\textcolor{FuncColor}{$\triangleright$\enspace\texttt{TwoSquares({\mdseries\slshape n})\index{TwoSquares@\texttt{TwoSquares}}
\label{TwoSquares}
}\hfill{\scriptsize (function)}}\\


 \index{representation!as a sum of two squares} \texttt{TwoSquares} returns a list of two integers $x \leq y$ such that the sum of the squares of $x$ and $y$ is equal to the nonnegative integer \mbox{\texttt{\mdseries\slshape n}}, i.e., $n = x^2 + y^2$. If no such representation exists \texttt{TwoSquares} will return \texttt{fail}. \texttt{TwoSquares} will return a representation for which the gcd of $x$ and $y$ is as small as possible. It is not specified which representation \texttt{TwoSquares} returns if there is more than one. 

 Let $a$ be the product of all maximal powers of primes of the form $4k+3$ dividing \mbox{\texttt{\mdseries\slshape n}}. A representation of \mbox{\texttt{\mdseries\slshape n}} as a sum of two squares exists if and only if $a$ is a perfect square. Let $b$ be the maximal power of $2$ dividing \mbox{\texttt{\mdseries\slshape n}} or its half, whichever is a perfect square. Then the minimal possible gcd of $x$ and $y$ is the square root $c$ of $a \cdot b$. The number of different minimal representation with $x \leq y$ is $2^{{l-1}}$, where $l$ is the number of different prime factors of the form $4k+1$ of \mbox{\texttt{\mdseries\slshape n}}. 

 The algorithm first finds a square root $r$ of $-1$ modulo $\mbox{\texttt{\mdseries\slshape n}} / (a \cdot b)$, which must exist, and applies the Euclidean algorithm to $r$ and \mbox{\texttt{\mdseries\slshape n}}. The first residues in the sequence that are smaller than $\sqrt{{\mbox{\texttt{\mdseries\slshape n}}/(a \cdot b)}}$ times $c$ are a possible pair $x$ and $y$. 

 Better descriptions of the algorithm and related topics can be found in \cite{Wagon90} and \cite{Zagier90}. 
\begin{Verbatim}[commandchars=!@|,fontsize=\small,frame=single,label=Example]
  !gapprompt@gap>| !gapinput@TwoSquares( 5 );|
  [ 1, 2 ]
  !gapprompt@gap>| !gapinput@TwoSquares( 11 );  # there is no representation|
  fail
  !gapprompt@gap>| !gapinput@TwoSquares( 16 );|
  [ 0, 4 ]
  !gapprompt@gap>| !gapinput@# 3 is the minimal possible gcd because 9 divides 45:|
  !gapprompt@gap>| !gapinput@TwoSquares( 45 );|
  [ 3, 6 ]
  !gapprompt@gap>| !gapinput@# it is not [5,10] because their gcd is not minimal:|
  !gapprompt@gap>| !gapinput@TwoSquares( 125 );|
  [ 2, 11 ]
  !gapprompt@gap>| !gapinput@# [10,11] would be the other possible representation:|
  !gapprompt@gap>| !gapinput@TwoSquares( 13*17 );|
  [ 5, 14 ]
  !gapprompt@gap>| !gapinput@TwoSquares( 848654483879497562821 );  # argument is prime|
  [ 6305894639, 28440994650 ]
\end{Verbatim}
 }

 }

 }

             
\chapter{\textcolor{Chapter }{Combinatorics}}\label{Combinatorics}
\logpage{[ 16, 0, 0 ]}
\hyperdef{L}{X7BDA99EE7CEADA7C}{}
{
  This chapter describes functions that deal with combinatorics. We mainly
concentrate on two areas. One is about \emph{selections}, that is the ways one can select elements from a set. The other is about \emph{partitions}, that is the ways one can partition a set into the union of pairwise disjoint
subsets.  
\section{\textcolor{Chapter }{Combinatorial Numbers}}\label{Combinatorial Numbers}
\logpage{[ 16, 1, 0 ]}
\hyperdef{L}{X800E48927D5C83F5}{}
{
  

\subsection{\textcolor{Chapter }{Factorial}}
\logpage{[ 16, 1, 1 ]}\nobreak
\hyperdef{L}{X87665F748594BF29}{}
{\noindent\textcolor{FuncColor}{$\triangleright$\enspace\texttt{Factorial({\mdseries\slshape n})\index{Factorial@\texttt{Factorial}}
\label{Factorial}
}\hfill{\scriptsize (function)}}\\


 returns the \emph{factorial} $n!$ of the positive integer \mbox{\texttt{\mdseries\slshape n}}, which is defined as the product $1 \cdot 2 \cdot 3 \cdots n$. 

 $n!$ is the number of permutations of a set of $n$ elements. $1 / n!$ is the coefficient of $x^n$ in the formal series $\exp(x)$, which is the generating function for factorial. 

 
\begin{Verbatim}[commandchars=!@|,fontsize=\small,frame=single,label=Example]
  !gapprompt@gap>| !gapinput@List( [0..10], Factorial );|
  [ 1, 1, 2, 6, 24, 120, 720, 5040, 40320, 362880, 3628800 ]
  !gapprompt@gap>| !gapinput@Factorial( 30 );|
  265252859812191058636308480000000
\end{Verbatim}
 

 \texttt{PermutationsList} (\ref{PermutationsList}) computes the set of all permutations of a list. }

 

\subsection{\textcolor{Chapter }{Binomial}}
\logpage{[ 16, 1, 2 ]}\nobreak
\hyperdef{L}{X7A9AF5F58682819D}{}
{\noindent\textcolor{FuncColor}{$\triangleright$\enspace\texttt{Binomial({\mdseries\slshape n, k})\index{Binomial@\texttt{Binomial}}
\label{Binomial}
}\hfill{\scriptsize (function)}}\\


 returns the \emph{binomial coefficient} \index{coefficient!binomial} \index{number!binomial} ${{n \choose k}}$ of integers \mbox{\texttt{\mdseries\slshape n}} and \mbox{\texttt{\mdseries\slshape k}}. This is defined by the conditions ${{n \choose k}} = 0$ for $k < 0$, ${{0 \choose k}} = 0$ for $k \neq 0$, ${{0 \choose 0}} = 1$ and the relation ${{n \choose k}} = {{n-1 \choose k}} + {{n-1 \choose k-1}}$ for all $n$ and $k$. 

 There are many ways of describing this function. For example, if $n \geq 0$ and $0 \leq k \leq n$, then ${{n \choose k}} = n! / (k! (n-k)!)$ and for $n < 0$ and $k \geq 0$ we have ${{n \choose k}} = (-1)^k {{-n+k-1 \choose k}}$. 

 If $n \geq 0$ then ${{n \choose k}}$ is the number of subsets with $k$ elements of a set with $n$ elements. Also, ${{n \choose k}}$ is the coefficient of $x^k$ in the polynomial $(x + 1)^n$, which is the generating function for ${{n \choose .}}$, hence the name. 

 
\begin{Verbatim}[commandchars=!@|,fontsize=\small,frame=single,label=Example]
  !gapprompt@gap>| !gapinput@# Knuth calls this the trademark of Binomial:|
  !gapprompt@gap>| !gapinput@List( [0..4], k->Binomial( 4, k ) );|
  [ 1, 4, 6, 4, 1 ]
  !gapprompt@gap>| !gapinput@List( [0..6], n->List( [0..6], k->Binomial( n, k ) ) );;|
  !gapprompt@gap>| !gapinput@# the lower triangle is called Pascal's triangle:|
  !gapprompt@gap>| !gapinput@PrintArray( last );|
  [ [   1,   0,   0,   0,   0,   0,   0 ],
    [   1,   1,   0,   0,   0,   0,   0 ],
    [   1,   2,   1,   0,   0,   0,   0 ],
    [   1,   3,   3,   1,   0,   0,   0 ],
    [   1,   4,   6,   4,   1,   0,   0 ],
    [   1,   5,  10,  10,   5,   1,   0 ],
    [   1,   6,  15,  20,  15,   6,   1 ] ]
  !gapprompt@gap>| !gapinput@Binomial( 50, 10 );|
  10272278170
\end{Verbatim}
 

 \texttt{NrCombinations} (\ref{NrCombinations}) is the generalization of \texttt{Binomial} for multisets. \texttt{Combinations} (\ref{Combinations}) computes the set of all combinations of a multiset. }

 

\subsection{\textcolor{Chapter }{Bell}}
\logpage{[ 16, 1, 3 ]}\nobreak
\hyperdef{L}{X7DC5667580522BDA}{}
{\noindent\textcolor{FuncColor}{$\triangleright$\enspace\texttt{Bell({\mdseries\slshape n})\index{Bell@\texttt{Bell}}
\label{Bell}
}\hfill{\scriptsize (function)}}\\


 returns the \emph{Bell number} \index{number!Bell} $B(n)$. The Bell numbers are defined by $B(0) = 1$ and the recurrence $B(n+1) = \sum_{{k = 0}}^n {{n \choose k}} B(k)$. 

 $B(n)$ is the number of ways to partition a set of \mbox{\texttt{\mdseries\slshape n}} elements into pairwise disjoint nonempty subsets (see \texttt{PartitionsSet} (\ref{PartitionsSet})). This implies of course that $B(n) = \sum_{{k = 0}}^n S_2(n,k)$ (see \texttt{Stirling2} (\ref{Stirling2})). $B(n)/n!$ is the coefficient of $x^n$ in the formal series $\exp( \exp(x)-1 )$, which is the generating function for $B(n)$. 

 
\begin{Verbatim}[commandchars=!@|,fontsize=\small,frame=single,label=Example]
  !gapprompt@gap>| !gapinput@List( [0..6], n -> Bell( n ) );|
  [ 1, 1, 2, 5, 15, 52, 203 ]
  !gapprompt@gap>| !gapinput@Bell( 14 );|
  190899322
\end{Verbatim}
 }

 

\subsection{\textcolor{Chapter }{Bernoulli}}
\logpage{[ 16, 1, 4 ]}\nobreak
\hyperdef{L}{X792FF6EA786A5C2B}{}
{\noindent\textcolor{FuncColor}{$\triangleright$\enspace\texttt{Bernoulli({\mdseries\slshape n})\index{Bernoulli@\texttt{Bernoulli}}
\label{Bernoulli}
}\hfill{\scriptsize (function)}}\\


 returns the \mbox{\texttt{\mdseries\slshape n}}\texttt{\symbol{45}}th \emph{Bernoulli number} \index{sequence!Bernoulli} $B_n$, which is defined by $B_0 = 1$ and $B_n = -\sum_{{k = 0}}^{{n-1}} {{n+1 \choose k}} B_k/(n+1)$. 

 $B_n / n!$ is the coefficient of $x^n$ in the power series of $x / (\exp(x)-1)$. Except for $B_1 = -1/2$ the Bernoulli numbers for odd indices are zero. 

 
\begin{Verbatim}[commandchars=!@|,fontsize=\small,frame=single,label=Example]
  !gapprompt@gap>| !gapinput@Bernoulli( 4 );|
  -1/30
  !gapprompt@gap>| !gapinput@Bernoulli( 10 );|
  5/66
  !gapprompt@gap>| !gapinput@Bernoulli( 12 );  # there is no simple pattern in Bernoulli numbers|
  -691/2730
  !gapprompt@gap>| !gapinput@Bernoulli( 50 );  # and they grow fairly fast|
  495057205241079648212477525/66
\end{Verbatim}
 }

 

\subsection{\textcolor{Chapter }{Stirling1}}
\logpage{[ 16, 1, 5 ]}\nobreak
\hyperdef{L}{X85037456785BB33C}{}
{\noindent\textcolor{FuncColor}{$\triangleright$\enspace\texttt{Stirling1({\mdseries\slshape n, k})\index{Stirling1@\texttt{Stirling1}}
\label{Stirling1}
}\hfill{\scriptsize (function)}}\\


 returns the \emph{Stirling number of the first kind} \index{Stirling number of the first kind} \index{number!Stirling, of the first kind} $S_1(n,k)$ of the integers \mbox{\texttt{\mdseries\slshape n}} and \mbox{\texttt{\mdseries\slshape k}}. Stirling numbers of the first kind are defined by $S_1(0,0) = 1$, $S_1(n,0) = S_1(0,k) = 0$ if $n, k \ne 0$ and the recurrence $S_1(n,k) = (n-1) S_1(n-1,k) + S_1(n-1,k-1)$. 

 $S_1(n,k)$ is the number of permutations of \mbox{\texttt{\mdseries\slshape n}} points with \mbox{\texttt{\mdseries\slshape k}} cycles. Stirling numbers of the first kind appear as coefficients in the
series $n! {{x \choose n}} = \sum_{{k = 0}}^n S_1(n,k) x^k$ which is the generating function for Stirling numbers of the first kind. Note
the similarity to $x^n = \sum_{{k = 0}}^n S_2(n,k) k! {{x \choose k}}$ (see \texttt{Stirling2} (\ref{Stirling2})). Also the definition of $S_1$ implies $S_1(n,k) = S_2(-k,-n)$ if $n, k < 0$. There are many formulae relating Stirling numbers of the first kind to
Stirling numbers of the second kind, Bell numbers, and Binomial coefficients. 

 
\begin{Verbatim}[commandchars=!@|,fontsize=\small,frame=single,label=Example]
  !gapprompt@gap>| !gapinput@# Knuth calls this the trademark of S_1:|
  !gapprompt@gap>| !gapinput@List( [0..4], k -> Stirling1( 4, k ) );|
  [ 0, 6, 11, 6, 1 ]
  !gapprompt@gap>| !gapinput@List( [0..6], n->List( [0..6], k->Stirling1( n, k ) ) );;|
  !gapprompt@gap>| !gapinput@# note the similarity with Pascal's triangle for Binomial numbers|
  !gapprompt@gap>| !gapinput@PrintArray( last );|
  [ [    1,    0,    0,    0,    0,    0,    0 ],
    [    0,    1,    0,    0,    0,    0,    0 ],
    [    0,    1,    1,    0,    0,    0,    0 ],
    [    0,    2,    3,    1,    0,    0,    0 ],
    [    0,    6,   11,    6,    1,    0,    0 ],
    [    0,   24,   50,   35,   10,    1,    0 ],
    [    0,  120,  274,  225,   85,   15,    1 ] ]
  !gapprompt@gap>| !gapinput@Stirling1(50,10);|
  101623020926367490059043797119309944043405505380503665627365376
\end{Verbatim}
 }

 

\subsection{\textcolor{Chapter }{Stirling2}}
\logpage{[ 16, 1, 6 ]}\nobreak
\hyperdef{L}{X7C93E14D7BC360F0}{}
{\noindent\textcolor{FuncColor}{$\triangleright$\enspace\texttt{Stirling2({\mdseries\slshape n, k})\index{Stirling2@\texttt{Stirling2}}
\label{Stirling2}
}\hfill{\scriptsize (function)}}\\


 returns the \emph{Stirling number of the second kind} \index{Stirling number of the second kind} \index{number!Stirling, of the second kind} $S_2(n,k)$ of the integers \mbox{\texttt{\mdseries\slshape n}} and \mbox{\texttt{\mdseries\slshape k}}. Stirling numbers of the second kind are defined by $S_2(0,0) = 1$, $S_2(n,0) = S_2(0,k) = 0$ if $n, k \ne 0$ and the recurrence $S_2(n,k) = k S_2(n-1,k) + S_2(n-1,k-1)$. 

 $S_2(n,k)$ is the number of ways to partition a set of \mbox{\texttt{\mdseries\slshape n}} elements into \mbox{\texttt{\mdseries\slshape k}} pairwise disjoint nonempty subsets (see \texttt{PartitionsSet} (\ref{PartitionsSet})). Stirling numbers of the second kind appear as coefficients in the expansion
of $x^n = \sum_{{k = 0}}^n S_2(n,k) k! {{x \choose k}}$. Note the similarity to $n! {{x \choose n}} = \sum_{{k = 0}}^n S_1(n,k) x^k$ (see \texttt{Stirling1} (\ref{Stirling1})). Also the definition of $S_2$ implies $S_2(n,k) = S_1(-k,-n)$ if $n, k < 0$. There are many formulae relating Stirling numbers of the second kind to
Stirling numbers of the first kind, Bell numbers, and Binomial coefficients. 

 
\begin{Verbatim}[commandchars=!@|,fontsize=\small,frame=single,label=Example]
  !gapprompt@gap>| !gapinput@# Knuth calls this the trademark of S_2:|
  !gapprompt@gap>| !gapinput@List( [0..4], k->Stirling2( 4, k ) );|
  [ 0, 1, 7, 6, 1 ]
  !gapprompt@gap>| !gapinput@List( [0..6], n->List( [0..6], k->Stirling2( n, k ) ) );;|
  !gapprompt@gap>| !gapinput@# note the similarity with Pascal's triangle for Binomial numbers|
  !gapprompt@gap>| !gapinput@PrintArray( last );|
  [ [   1,   0,   0,   0,   0,   0,   0 ],
    [   0,   1,   0,   0,   0,   0,   0 ],
    [   0,   1,   1,   0,   0,   0,   0 ],
    [   0,   1,   3,   1,   0,   0,   0 ],
    [   0,   1,   7,   6,   1,   0,   0 ],
    [   0,   1,  15,  25,  10,   1,   0 ],
    [   0,   1,  31,  90,  65,  15,   1 ] ]
  !gapprompt@gap>| !gapinput@Stirling2( 50, 10 );|
  26154716515862881292012777396577993781727011
\end{Verbatim}
 }

 }

  
\section{\textcolor{Chapter }{Combinations, Arrangements and Tuples}}\label{Combinations, Arrangements and Tuples}
\logpage{[ 16, 2, 0 ]}
\hyperdef{L}{X81B4696585C38147}{}
{
  

\subsection{\textcolor{Chapter }{Combinations}}
\logpage{[ 16, 2, 1 ]}\nobreak
\hyperdef{L}{X8770F16D794C0ADB}{}
{\noindent\textcolor{FuncColor}{$\triangleright$\enspace\texttt{Combinations({\mdseries\slshape mset[, k]})\index{Combinations@\texttt{Combinations}}
\label{Combinations}
}\hfill{\scriptsize (function)}}\\


 returns the set of all combinations of the multiset \mbox{\texttt{\mdseries\slshape mset}} (a list of objects which may contain the same object several times) with \mbox{\texttt{\mdseries\slshape k}} elements; if \mbox{\texttt{\mdseries\slshape k}} is not given it returns all combinations of \mbox{\texttt{\mdseries\slshape mset}}. 

 A \emph{combination} of \mbox{\texttt{\mdseries\slshape mset}} is an unordered selection without repetitions and is represented by a sorted
sublist of \mbox{\texttt{\mdseries\slshape mset}}. If \mbox{\texttt{\mdseries\slshape mset}} is a proper set, there are ${{|\mbox{\texttt{\mdseries\slshape mset}}| \choose \mbox{\texttt{\mdseries\slshape k}}}}$ (see \texttt{Binomial} (\ref{Binomial})) combinations with \mbox{\texttt{\mdseries\slshape k}} elements, and the set of all combinations is just the \emph{power set} \index{power set} \index{subsets} of \mbox{\texttt{\mdseries\slshape mset}}, which contains all \emph{subsets} of \mbox{\texttt{\mdseries\slshape mset}} and has cardinality $2^{{|\mbox{\texttt{\mdseries\slshape mset}}|}}$. 

 To loop over combinations of a larger multiset use \texttt{IteratorOfCombinations} (\ref{IteratorOfCombinations}) which produces combinations one by one and may save a lot of memory. Another
memory efficient representation of the list of all combinations is provided by \texttt{EnumeratorOfCombinations} (\ref{EnumeratorOfCombinations}). }

 
\subsection{\textcolor{Chapter }{Iterator and enumerator of combinations}}\logpage{[ 16, 2, 2 ]}
\hyperdef{L}{X78DD5C0D81057540}{}
{
\noindent\textcolor{FuncColor}{$\triangleright$\enspace\texttt{IteratorOfCombinations({\mdseries\slshape mset[, k]})\index{IteratorOfCombinations@\texttt{IteratorOfCombinations}}
\label{IteratorOfCombinations}
}\hfill{\scriptsize (function)}}\\
\noindent\textcolor{FuncColor}{$\triangleright$\enspace\texttt{EnumeratorOfCombinations({\mdseries\slshape mset})\index{EnumeratorOfCombinations@\texttt{EnumeratorOfCombinations}}
\label{EnumeratorOfCombinations}
}\hfill{\scriptsize (function)}}\\


 \texttt{IteratorOfCombinations} returns an \texttt{Iterator} (\ref{Iterator}) for combinations (see \texttt{Combinations} (\ref{Combinations})) of the given multiset \mbox{\texttt{\mdseries\slshape mset}}. If a non\texttt{\symbol{45}}negative integer \mbox{\texttt{\mdseries\slshape k}} is given as second argument then only the combinations with \mbox{\texttt{\mdseries\slshape k}} entries are produced, otherwise all combinations. 

 \texttt{EnumeratorOfCombinations} returns an \texttt{Enumerator} (\ref{Enumerator}) of the given multiset \mbox{\texttt{\mdseries\slshape mset}}. Currently only a variant without second argument \mbox{\texttt{\mdseries\slshape k}} is implemented. 

 The ordering of combinations from these functions can be different and also
different from the list returned by \texttt{Combinations} (\ref{Combinations}). 

 
\begin{Verbatim}[commandchars=!@|,fontsize=\small,frame=single,label=Example]
  !gapprompt@gap>| !gapinput@m:=[1..15];; Add(m, 15);|
  !gapprompt@gap>| !gapinput@NrCombinations(m);|
  49152
  !gapprompt@gap>| !gapinput@i := 0;; for c in Combinations(m) do i := i+1; od;|
  !gapprompt@gap>| !gapinput@i;|
  49152
  !gapprompt@gap>| !gapinput@cm := EnumeratorOfCombinations(m);;|
  !gapprompt@gap>| !gapinput@cm[1000];|
  [ 1, 2, 3, 6, 7, 8, 9, 10 ]
  !gapprompt@gap>| !gapinput@Position(cm, [1,13,15,15]);|
  36866
\end{Verbatim}
 }

 

\subsection{\textcolor{Chapter }{NrCombinations}}
\logpage{[ 16, 2, 3 ]}\nobreak
\hyperdef{L}{X82A6E98C85714FD0}{}
{\noindent\textcolor{FuncColor}{$\triangleright$\enspace\texttt{NrCombinations({\mdseries\slshape mset[, k]})\index{NrCombinations@\texttt{NrCombinations}}
\label{NrCombinations}
}\hfill{\scriptsize (function)}}\\


 returns the number of \texttt{Combinations(\mbox{\texttt{\mdseries\slshape mset}},\mbox{\texttt{\mdseries\slshape k}})}. 

 
\begin{Verbatim}[commandchars=!@|,fontsize=\small,frame=single,label=Example]
  !gapprompt@gap>| !gapinput@Combinations( [1,2,2,3] );|
  [ [  ], [ 1 ], [ 1, 2 ], [ 1, 2, 2 ], [ 1, 2, 2, 3 ], [ 1, 2, 3 ],
    [ 1, 3 ], [ 2 ], [ 2, 2 ], [ 2, 2, 3 ], [ 2, 3 ], [ 3 ] ]
  !gapprompt@gap>| !gapinput@# number of different hands in a game of poker:|
  !gapprompt@gap>| !gapinput@NrCombinations( [1..52], 5 );|
  2598960
\end{Verbatim}
 

 The function \texttt{Arrangements} (\ref{Arrangements}) computes ordered selections without repetitions, \texttt{UnorderedTuples} (\ref{UnorderedTuples}) computes unordered selections with repetitions, and \texttt{Tuples} (\ref{Tuples}) computes ordered selections with repetitions. }

 

\subsection{\textcolor{Chapter }{Arrangements}}
\logpage{[ 16, 2, 4 ]}\nobreak
\hyperdef{L}{X7837B3357C7566C8}{}
{\noindent\textcolor{FuncColor}{$\triangleright$\enspace\texttt{Arrangements({\mdseries\slshape mset[, k]})\index{Arrangements@\texttt{Arrangements}}
\label{Arrangements}
}\hfill{\scriptsize (function)}}\\


 returns the set of arrangements of the multiset \mbox{\texttt{\mdseries\slshape mset}} that contain \mbox{\texttt{\mdseries\slshape k}} elements. If \mbox{\texttt{\mdseries\slshape k}} is not given it returns all arrangements of \mbox{\texttt{\mdseries\slshape mset}}. 

 An \emph{arrangement} of \mbox{\texttt{\mdseries\slshape mset}} is an ordered selection without repetitions and is represented by a list that
contains only elements from \mbox{\texttt{\mdseries\slshape mset}}, but maybe in a different order. If \mbox{\texttt{\mdseries\slshape mset}} is a proper set there are $|mset|! / (|mset|-k)!$ (see \texttt{Factorial} (\ref{Factorial})) arrangements with \mbox{\texttt{\mdseries\slshape k}} elements. }

 

\subsection{\textcolor{Chapter }{NrArrangements}}
\logpage{[ 16, 2, 5 ]}\nobreak
\hyperdef{L}{X7DE1ABD47D19F140}{}
{\noindent\textcolor{FuncColor}{$\triangleright$\enspace\texttt{NrArrangements({\mdseries\slshape mset[, k]})\index{NrArrangements@\texttt{NrArrangements}}
\label{NrArrangements}
}\hfill{\scriptsize (function)}}\\


 returns the number of \texttt{Arrangements(\mbox{\texttt{\mdseries\slshape mset}},\mbox{\texttt{\mdseries\slshape k}})}. 

 As an example of arrangements of a multiset, think of the game Scrabble.
Suppose you have the six characters of the word \texttt{"settle"} and you have to make a four letter word. Then the possibilities are given by 

 
\begin{Verbatim}[commandchars=!@|,fontsize=\small,frame=single,label=Example]
  !gapprompt@gap>| !gapinput@Arrangements( ["s","e","t","t","l","e"], 4 );|
  [ [ "e", "e", "l", "s" ], [ "e", "e", "l", "t" ], [ "e", "e", "s", "l" ],
    [ "e", "e", "s", "t" ], [ "e", "e", "t", "l" ], [ "e", "e", "t", "s" ],
    ... 93 more possibilities ...
    [ "t", "t", "l", "s" ], [ "t", "t", "s", "e" ], [ "t", "t", "s", "l" ] ]
\end{Verbatim}
 

 Can you find the five proper English words, where \texttt{"lets"} does not count? Note that the fact that the list returned by \texttt{Arrangements} (\ref{Arrangements}) is a proper set means in this example that the possibilities are listed in the
same order as they appear in the dictionary. 

 
\begin{Verbatim}[commandchars=!@|,fontsize=\small,frame=single,label=Example]
  !gapprompt@gap>| !gapinput@NrArrangements( ["s","e","t","t","l","e"] );|
  523
\end{Verbatim}
 

 The function \texttt{Combinations} (\ref{Combinations}) computes unordered selections without repetitions, \texttt{UnorderedTuples} (\ref{UnorderedTuples}) computes unordered selections with repetitions, and \texttt{Tuples} (\ref{Tuples}) computes ordered selections with repetitions. }

 

\subsection{\textcolor{Chapter }{UnorderedTuples}}
\logpage{[ 16, 2, 6 ]}\nobreak
\hyperdef{L}{X81601C6786120DDC}{}
{\noindent\textcolor{FuncColor}{$\triangleright$\enspace\texttt{UnorderedTuples({\mdseries\slshape set, k})\index{UnorderedTuples@\texttt{UnorderedTuples}}
\label{UnorderedTuples}
}\hfill{\scriptsize (function)}}\\


 returns the set of all unordered tuples of length \mbox{\texttt{\mdseries\slshape k}} of the set \mbox{\texttt{\mdseries\slshape set}}. 

 An \emph{unordered tuple} of length \mbox{\texttt{\mdseries\slshape k}} of \mbox{\texttt{\mdseries\slshape set}} is an unordered selection with repetitions of \mbox{\texttt{\mdseries\slshape set}} and is represented by a sorted list of length \mbox{\texttt{\mdseries\slshape k}} containing elements from \mbox{\texttt{\mdseries\slshape set}}. There are ${{|set| + k - 1 \choose k}}$ (see \texttt{Binomial} (\ref{Binomial})) such unordered tuples. 

 Note that the fact that \texttt{UnorderedTuples} returns a set implies that the last index runs fastest. That means the first
tuple contains the smallest element from \mbox{\texttt{\mdseries\slshape set}} \mbox{\texttt{\mdseries\slshape k}} times, the second tuple contains the smallest element of \mbox{\texttt{\mdseries\slshape set}} at all positions except at the last positions, where it contains the second
smallest element from \mbox{\texttt{\mdseries\slshape set}} and so on. }

 

\subsection{\textcolor{Chapter }{NrUnorderedTuples}}
\logpage{[ 16, 2, 7 ]}\nobreak
\hyperdef{L}{X7959281584C42C52}{}
{\noindent\textcolor{FuncColor}{$\triangleright$\enspace\texttt{NrUnorderedTuples({\mdseries\slshape set, k})\index{NrUnorderedTuples@\texttt{NrUnorderedTuples}}
\label{NrUnorderedTuples}
}\hfill{\scriptsize (function)}}\\


 returns the number of \texttt{UnorderedTuples(\mbox{\texttt{\mdseries\slshape set}},\mbox{\texttt{\mdseries\slshape k}})}. 

 As an example for unordered tuples think of a poker\texttt{\symbol{45}}like
game played with 5 dice. Then each possible hand corresponds to an unordered
five\texttt{\symbol{45}}tuple from the set $\{ 1, 2, \ldots, 6 \}$. 

 
\begin{Verbatim}[commandchars=!@|,fontsize=\small,frame=single,label=Example]
  !gapprompt@gap>| !gapinput@NrUnorderedTuples( [1..6], 5 );|
  252
  !gapprompt@gap>| !gapinput@UnorderedTuples( [1..6], 5 );|
  [ [ 1, 1, 1, 1, 1 ], [ 1, 1, 1, 1, 2 ], [ 1, 1, 1, 1, 3 ], [ 1, 1, 1, 1, 4 ],
    [ 1, 1, 1, 1, 5 ], [ 1, 1, 1, 1, 6 ], [ 1, 1, 1, 2, 2 ], [ 1, 1, 1, 2, 3 ],
    ... 100 more tuples ...
    [ 1, 3, 5, 5, 6 ], [ 1, 3, 5, 6, 6 ], [ 1, 3, 6, 6, 6 ], [ 1, 4, 4, 4, 4 ],
    ... 100 more tuples ...
    [ 3, 3, 5, 5, 5 ], [ 3, 3, 5, 5, 6 ], [ 3, 3, 5, 6, 6 ], [ 3, 3, 6, 6, 6 ],
    ... 32 more tuples ...
    [ 5, 5, 5, 6, 6 ], [ 5, 5, 6, 6, 6 ], [ 5, 6, 6, 6, 6 ], [ 6, 6, 6, 6, 6 ] ]
\end{Verbatim}
 

 The function \texttt{Combinations} (\ref{Combinations}) computes unordered selections without repetitions, \texttt{Arrangements} (\ref{Arrangements}) computes ordered selections without repetitions, and \texttt{Tuples} (\ref{Tuples}) computes ordered selections with repetitions. }

 

\subsection{\textcolor{Chapter }{Tuples}}
\logpage{[ 16, 2, 8 ]}\nobreak
\hyperdef{L}{X86A3CA0F7CC8C320}{}
{\noindent\textcolor{FuncColor}{$\triangleright$\enspace\texttt{Tuples({\mdseries\slshape set, k})\index{Tuples@\texttt{Tuples}}
\label{Tuples}
}\hfill{\scriptsize (function)}}\\


 returns the set of all ordered tuples of length \mbox{\texttt{\mdseries\slshape k}} of the set \mbox{\texttt{\mdseries\slshape set}}. 

 An \emph{ordered tuple} of length \mbox{\texttt{\mdseries\slshape k}} of \mbox{\texttt{\mdseries\slshape set}} is an ordered selection with repetition and is represented by a list of length \mbox{\texttt{\mdseries\slshape k}} containing elements of \mbox{\texttt{\mdseries\slshape set}}. There are $|\mbox{\texttt{\mdseries\slshape set}}|^{\mbox{\texttt{\mdseries\slshape k}}}$ such ordered tuples. 

 Note that the fact that \texttt{Tuples} returns a set implies that the last index runs fastest. That means the first
tuple contains the smallest element from \mbox{\texttt{\mdseries\slshape set}} \mbox{\texttt{\mdseries\slshape k}} times, the second tuple contains the smallest element of \mbox{\texttt{\mdseries\slshape set}} at all positions except at the last positions, where it contains the second
smallest element from \mbox{\texttt{\mdseries\slshape set}} and so on. }

 

\subsection{\textcolor{Chapter }{EnumeratorOfTuples}}
\logpage{[ 16, 2, 9 ]}\nobreak
\hyperdef{L}{X7BA135297E8DA819}{}
{\noindent\textcolor{FuncColor}{$\triangleright$\enspace\texttt{EnumeratorOfTuples({\mdseries\slshape set, k})\index{EnumeratorOfTuples@\texttt{EnumeratorOfTuples}}
\label{EnumeratorOfTuples}
}\hfill{\scriptsize (function)}}\\


 This function is referred to as an example of enumerators that are defined by
functions but are not constructed from a domain. The result is equal to that
of \texttt{Tuples( \mbox{\texttt{\mdseries\slshape set}}, \mbox{\texttt{\mdseries\slshape k}} )}. However, the entries are not stored physically in the list but are
created/identified on demand. 

 }

 

\subsection{\textcolor{Chapter }{IteratorOfTuples}}
\logpage{[ 16, 2, 10 ]}\nobreak
\hyperdef{L}{X86416A31807B0086}{}
{\noindent\textcolor{FuncColor}{$\triangleright$\enspace\texttt{IteratorOfTuples({\mdseries\slshape set, k})\index{IteratorOfTuples@\texttt{IteratorOfTuples}}
\label{IteratorOfTuples}
}\hfill{\scriptsize (function)}}\\


 For a set \mbox{\texttt{\mdseries\slshape set}} and a positive integer \mbox{\texttt{\mdseries\slshape k}}, \texttt{IteratorOfTuples} returns an iterator (see{\nobreakspace}\ref{Iterators}) of the set of all ordered tuples (see{\nobreakspace}\texttt{Tuples} (\ref{Tuples})) of length \mbox{\texttt{\mdseries\slshape k}} of the set \mbox{\texttt{\mdseries\slshape set}}. The tuples are returned in lexicographic order. }

 

\subsection{\textcolor{Chapter }{NrTuples}}
\logpage{[ 16, 2, 11 ]}\nobreak
\hyperdef{L}{X85E18A9A87FD4CA2}{}
{\noindent\textcolor{FuncColor}{$\triangleright$\enspace\texttt{NrTuples({\mdseries\slshape set, k})\index{NrTuples@\texttt{NrTuples}}
\label{NrTuples}
}\hfill{\scriptsize (function)}}\\


 returns the number of \texttt{Tuples(\mbox{\texttt{\mdseries\slshape set}},\mbox{\texttt{\mdseries\slshape k}})}. 
\begin{Verbatim}[commandchars=!@|,fontsize=\small,frame=single,label=Example]
  !gapprompt@gap>| !gapinput@Tuples( [1,2,3], 2 );|
  [ [ 1, 1 ], [ 1, 2 ], [ 1, 3 ], [ 2, 1 ], [ 2, 2 ], [ 2, 3 ],
    [ 3, 1 ], [ 3, 2 ], [ 3, 3 ] ]
  !gapprompt@gap>| !gapinput@NrTuples( [1..10], 5 );|
  100000
\end{Verbatim}
 

 \texttt{Tuples(\mbox{\texttt{\mdseries\slshape set}},\mbox{\texttt{\mdseries\slshape k}})} can also be viewed as the \mbox{\texttt{\mdseries\slshape k}}\texttt{\symbol{45}}fold cartesian product of \mbox{\texttt{\mdseries\slshape set}} (see \texttt{Cartesian} (\ref{Cartesian:for a list})). 

 The function \texttt{Combinations} (\ref{Combinations}) computes unordered selections without repetitions, \texttt{Arrangements} (\ref{Arrangements}) computes ordered selections without repetitions, and finally the function \texttt{UnorderedTuples} (\ref{UnorderedTuples}) computes unordered selections with repetitions. }

 

\subsection{\textcolor{Chapter }{PermutationsList}}
\logpage{[ 16, 2, 12 ]}\nobreak
\hyperdef{L}{X7B0143FB83F359B7}{}
{\noindent\textcolor{FuncColor}{$\triangleright$\enspace\texttt{PermutationsList({\mdseries\slshape mset})\index{PermutationsList@\texttt{PermutationsList}}
\label{PermutationsList}
}\hfill{\scriptsize (function)}}\\


 \texttt{PermutationsList} returns the set of permutations of the multiset \mbox{\texttt{\mdseries\slshape mset}}. 

 A \emph{permutation} is represented by a list that contains exactly the same elements as \mbox{\texttt{\mdseries\slshape mset}}, but possibly in different order. If \mbox{\texttt{\mdseries\slshape mset}} is a proper set there are $|\mbox{\texttt{\mdseries\slshape mset}}| !$ (see \texttt{Factorial} (\ref{Factorial})) such permutations. Otherwise if the first elements appears $k_1$ times, the second element appears $k_2$ times and so on, the number of permutations is $|\mbox{\texttt{\mdseries\slshape mset}}| ! / (k_1! k_2! \ldots)$, which is sometimes called multinomial coefficient. }

 

\subsection{\textcolor{Chapter }{NrPermutationsList}}
\logpage{[ 16, 2, 13 ]}\nobreak
\hyperdef{L}{X8629A2908050EB3A}{}
{\noindent\textcolor{FuncColor}{$\triangleright$\enspace\texttt{NrPermutationsList({\mdseries\slshape mset})\index{NrPermutationsList@\texttt{NrPermutationsList}}
\label{NrPermutationsList}
}\hfill{\scriptsize (function)}}\\


 returns the number of \texttt{PermutationsList(\mbox{\texttt{\mdseries\slshape mset}})}. 
\begin{Verbatim}[commandchars=!@|,fontsize=\small,frame=single,label=Example]
  !gapprompt@gap>| !gapinput@PermutationsList( [1,2,3] );|
  [ [ 1, 2, 3 ], [ 1, 3, 2 ], [ 2, 1, 3 ], [ 2, 3, 1 ], [ 3, 1, 2 ],
    [ 3, 2, 1 ] ]
  !gapprompt@gap>| !gapinput@PermutationsList( [1,1,2,2] );|
  [ [ 1, 1, 2, 2 ], [ 1, 2, 1, 2 ], [ 1, 2, 2, 1 ], [ 2, 1, 1, 2 ],
    [ 2, 1, 2, 1 ], [ 2, 2, 1, 1 ] ]
  !gapprompt@gap>| !gapinput@NrPermutationsList( [1,2,2,3,3,3,4,4,4,4] );|
  12600
\end{Verbatim}
 

 The function \texttt{Arrangements} (\ref{Arrangements}) is the generalization of \texttt{PermutationsList} (\ref{PermutationsList}) that allows you to specify the size of the permutations. \texttt{Derangements} (\ref{Derangements}) computes permutations that have no fixed points. }

 

\subsection{\textcolor{Chapter }{Derangements}}
\logpage{[ 16, 2, 14 ]}\nobreak
\hyperdef{L}{X79C159507B2BF1C9}{}
{\noindent\textcolor{FuncColor}{$\triangleright$\enspace\texttt{Derangements({\mdseries\slshape list})\index{Derangements@\texttt{Derangements}}
\label{Derangements}
}\hfill{\scriptsize (function)}}\\


 returns the set of all derangements of the list \mbox{\texttt{\mdseries\slshape list}}. 

 A \emph{derangement} is a fixpointfree permutation of \mbox{\texttt{\mdseries\slshape list}} and is represented by a list that contains exactly the same elements as \mbox{\texttt{\mdseries\slshape list}}, but in such an order that the derangement has at no position the same
element as \mbox{\texttt{\mdseries\slshape list}}. If the list \mbox{\texttt{\mdseries\slshape list}} contains no element twice there are exactly $|\mbox{\texttt{\mdseries\slshape list}}|! (1/2! - 1/3! + 1/4! - \cdots + (-1)^n / n!)$ derangements. 

 Note that the ratio \texttt{NrPermutationsList( [ 1 .. n ] ) / NrDerangements( [ 1 .. n ] )}, which is $n! / (n! (1/2! - 1/3! + 1/4! - \cdots + (-1)^n / n!))$ is an approximation for the base of the natural logarithm $e = 2.7182818285\ldots$, which is correct to about $n$ digits. }

 

\subsection{\textcolor{Chapter }{NrDerangements}}
\logpage{[ 16, 2, 15 ]}\nobreak
\hyperdef{L}{X7C1741B181A9AB9C}{}
{\noindent\textcolor{FuncColor}{$\triangleright$\enspace\texttt{NrDerangements({\mdseries\slshape list})\index{NrDerangements@\texttt{NrDerangements}}
\label{NrDerangements}
}\hfill{\scriptsize (function)}}\\


 returns the number of \texttt{Derangements(\mbox{\texttt{\mdseries\slshape list}})}. 

 As an example of derangements suppose that you have to send four different
letters to four different people. Then a derangement corresponds to a way to
send those letters such that no letter reaches the intended person. 

 
\begin{Verbatim}[commandchars=!@|,fontsize=\small,frame=single,label=Example]
  !gapprompt@gap>| !gapinput@Derangements( [1,2,3,4] );|
  [ [ 2, 1, 4, 3 ], [ 2, 3, 4, 1 ], [ 2, 4, 1, 3 ], [ 3, 1, 4, 2 ],
    [ 3, 4, 1, 2 ], [ 3, 4, 2, 1 ], [ 4, 1, 2, 3 ], [ 4, 3, 1, 2 ],
    [ 4, 3, 2, 1 ] ]
  !gapprompt@gap>| !gapinput@NrDerangements( [1..10] );|
  1334961
  !gapprompt@gap>| !gapinput@Int( 10^7*NrPermutationsList([1..10])/last );|
  27182816
  !gapprompt@gap>| !gapinput@Derangements( [1,1,2,2,3,3] );|
  [ [ 2, 2, 3, 3, 1, 1 ], [ 2, 3, 1, 3, 1, 2 ], [ 2, 3, 1, 3, 2, 1 ],
    [ 2, 3, 3, 1, 1, 2 ], [ 2, 3, 3, 1, 2, 1 ], [ 3, 2, 1, 3, 1, 2 ],
    [ 3, 2, 1, 3, 2, 1 ], [ 3, 2, 3, 1, 1, 2 ], [ 3, 2, 3, 1, 2, 1 ],
    [ 3, 3, 1, 1, 2, 2 ] ]
  !gapprompt@gap>| !gapinput@NrDerangements( [1,2,2,3,3,3,4,4,4,4] );|
  338
\end{Verbatim}
 

 The function \texttt{PermutationsList} (\ref{PermutationsList}) computes all permutations of a list. }

 

\subsection{\textcolor{Chapter }{PartitionsSet}}
\logpage{[ 16, 2, 16 ]}\nobreak
\hyperdef{L}{X7A13D8DC8204525F}{}
{\noindent\textcolor{FuncColor}{$\triangleright$\enspace\texttt{PartitionsSet({\mdseries\slshape set[, k]})\index{PartitionsSet@\texttt{PartitionsSet}}
\label{PartitionsSet}
}\hfill{\scriptsize (function)}}\\


 returns the set of all unordered partitions of the set \mbox{\texttt{\mdseries\slshape set}} into \mbox{\texttt{\mdseries\slshape k}} pairwise disjoint nonempty sets. If \mbox{\texttt{\mdseries\slshape k}} is not given it returns all unordered partitions of \mbox{\texttt{\mdseries\slshape set}} for all \mbox{\texttt{\mdseries\slshape k}}. 

 An \emph{unordered partition} of \mbox{\texttt{\mdseries\slshape set}} is a set of pairwise disjoint nonempty sets with union \mbox{\texttt{\mdseries\slshape set}} and is represented by a sorted list of such sets. There are $B( |set| )$ (see \texttt{Bell} (\ref{Bell})) partitions of the set \mbox{\texttt{\mdseries\slshape set}} and $S_2( |set|, k )$ (see \texttt{Stirling2} (\ref{Stirling2})) partitions with \mbox{\texttt{\mdseries\slshape k}} elements. }

 

\subsection{\textcolor{Chapter }{NrPartitionsSet}}
\logpage{[ 16, 2, 17 ]}\nobreak
\hyperdef{L}{X7BCD7FC2876386F1}{}
{\noindent\textcolor{FuncColor}{$\triangleright$\enspace\texttt{NrPartitionsSet({\mdseries\slshape set[, k]})\index{NrPartitionsSet@\texttt{NrPartitionsSet}}
\label{NrPartitionsSet}
}\hfill{\scriptsize (function)}}\\


 returns the number of \texttt{PartitionsSet(\mbox{\texttt{\mdseries\slshape set}},\mbox{\texttt{\mdseries\slshape k}})}. 
\begin{Verbatim}[commandchars=!@|,fontsize=\small,frame=single,label=Example]
  !gapprompt@gap>| !gapinput@PartitionsSet( [1,2,3] );|
  [ [ [ 1 ], [ 2 ], [ 3 ] ], [ [ 1 ], [ 2, 3 ] ], [ [ 1, 2 ], [ 3 ] ],
    [ [ 1, 2, 3 ] ], [ [ 1, 3 ], [ 2 ] ] ]
  !gapprompt@gap>| !gapinput@PartitionsSet( [1,2,3,4], 2 );|
  [ [ [ 1 ], [ 2, 3, 4 ] ], [ [ 1, 2 ], [ 3, 4 ] ],
    [ [ 1, 2, 3 ], [ 4 ] ], [ [ 1, 2, 4 ], [ 3 ] ],
    [ [ 1, 3 ], [ 2, 4 ] ], [ [ 1, 3, 4 ], [ 2 ] ],
    [ [ 1, 4 ], [ 2, 3 ] ] ]
  !gapprompt@gap>| !gapinput@NrPartitionsSet( [1..6] );|
  203
  !gapprompt@gap>| !gapinput@NrPartitionsSet( [1..10], 3 );|
  9330
\end{Verbatim}
 

 Note that \texttt{PartitionsSet} (\ref{PartitionsSet}) does currently not support multisets and that there is currently no ordered
counterpart. }

 

\subsection{\textcolor{Chapter }{Partitions}}
\logpage{[ 16, 2, 18 ]}\nobreak
\hyperdef{L}{X84A6D15F8107008B}{}
{\noindent\textcolor{FuncColor}{$\triangleright$\enspace\texttt{Partitions({\mdseries\slshape n[, k]})\index{Partitions@\texttt{Partitions}}
\label{Partitions}
}\hfill{\scriptsize (function)}}\\


 returns the set of all (unordered) partitions of the positive integer \mbox{\texttt{\mdseries\slshape n}} into sums with \mbox{\texttt{\mdseries\slshape k}} summands. If \mbox{\texttt{\mdseries\slshape k}} is not given it returns all unordered partitions of \mbox{\texttt{\mdseries\slshape n}} for all \mbox{\texttt{\mdseries\slshape k}}. 

 An \emph{unordered partition} is an unordered sum $n = p_1 + p_2 + \cdots + p_k$ of positive integers and is represented by the list $p = [ p_1, p_2, \ldots, p_k ]$, in nonincreasing order, i.e., $p_1 \geq p_2 \geq \ldots \geq p_k$. We write $p \vdash n$. There are approximately $\exp(\pi \sqrt{{2/3 n}}) / (4 \sqrt{{3}} n)$ such partitions, use \texttt{NrPartitions} (\ref{NrPartitions}) to compute the precise number. 

 If you want to loop over all partitions of some larger \mbox{\texttt{\mdseries\slshape n}} use the more memory efficient \texttt{IteratorOfPartitions} (\ref{IteratorOfPartitions}). 

 It is possible to associate with every partition of the integer \mbox{\texttt{\mdseries\slshape n}} a conjugacy class of permutations in the symmetric group on \mbox{\texttt{\mdseries\slshape n}} points and vice versa. Therefore $p(n) := $\texttt{NrPartitions}$(n)$ is the number of conjugacy classes of the symmetric group on \mbox{\texttt{\mdseries\slshape n}} points. 

 Ramanujan found the identities $p(5i+4) = 0$ mod 5, $p(7i+5) = 0$ mod 7 and $p(11i+6) = 0$ mod 11 and many other fascinating things about the number of partitions. }

 

\subsection{\textcolor{Chapter }{IteratorOfPartitions}}
\logpage{[ 16, 2, 19 ]}\nobreak
\hyperdef{L}{X8793AEBD7E529E1D}{}
{\noindent\textcolor{FuncColor}{$\triangleright$\enspace\texttt{IteratorOfPartitions({\mdseries\slshape n})\index{IteratorOfPartitions@\texttt{IteratorOfPartitions}}
\label{IteratorOfPartitions}
}\hfill{\scriptsize (function)}}\\


 For a positive integer \mbox{\texttt{\mdseries\slshape n}}, \texttt{IteratorOfPartitions} returns an iterator (see{\nobreakspace}\ref{Iterators}) of the set of partitions of \mbox{\texttt{\mdseries\slshape n}} (see{\nobreakspace}\texttt{Partitions} (\ref{Partitions})). The partitions of \mbox{\texttt{\mdseries\slshape n}} are returned in lexicographic order. }

 
\subsection{\textcolor{Chapter }{Iterator and enumerator of unordered set partitions}}\logpage{[ 16, 2, 20 ]}
\hyperdef{L}{X7ECDBC227BFB2EC6}{}
{
\noindent\textcolor{FuncColor}{$\triangleright$\enspace\texttt{IteratorOfPartitionsSet({\mdseries\slshape set[, k[, flag]]})\index{IteratorOfPartitionsSet@\texttt{IteratorOfPartitionsSet}}
\label{IteratorOfPartitionsSet}
}\hfill{\scriptsize (function)}}\\
\noindent\textcolor{FuncColor}{$\triangleright$\enspace\texttt{EnumeratorOfPartitionsSet({\mdseries\slshape set[, k[, flag]]})\index{EnumeratorOfPartitionsSet@\texttt{EnumeratorOfPartitionsSet}}
\label{EnumeratorOfPartitionsSet}
}\hfill{\scriptsize (function)}}\\


 \texttt{IteratorOfPartitionsSet} returns an iterator (see{\nobreakspace}\ref{Iterators}) for all unordered partitions of the set \mbox{\texttt{\mdseries\slshape set}} into pairwise disjoint nonempty sets (see \texttt{PartitionsSet} (\ref{PartitionsSet})). If \mbox{\texttt{\mdseries\slshape k}} given and \mbox{\texttt{\mdseries\slshape flag}} is omitted or equal to \texttt{false}, then only partitions of size \mbox{\texttt{\mdseries\slshape k}} are computed. If \mbox{\texttt{\mdseries\slshape k}} is given and \mbox{\texttt{\mdseries\slshape flag}} is equal to \texttt{true}, then only partitions of size at most \mbox{\texttt{\mdseries\slshape k}} are computed. 

 \texttt{EnumeratorOfPartitionsSet} returns an enumerator (see{\nobreakspace}\texttt{Enumerator} (\ref{Enumerator})) for all unordered partitions of the set \mbox{\texttt{\mdseries\slshape set}}. The arguments \mbox{\texttt{\mdseries\slshape k}} and \mbox{\texttt{\mdseries\slshape flag}} function in the same manner as for the iterator. 

 The ordering of the partitions from these functions can be different and also
different from the list returned by \texttt{PartitionsSet} (\ref{PartitionsSet}). 

 
\begin{Verbatim}[commandchars=!@|,fontsize=\small,frame=single,label=Example]
  !gapprompt@gap>| !gapinput@m:=[1..9];;|
  !gapprompt@gap>| !gapinput@Bell(9);|
  21147
  !gapprompt@gap>| !gapinput@Length(PartitionsSet(m));|
  21147
  !gapprompt@gap>| !gapinput@cm := EnumeratorOfPartitionsSet(m);;|
  !gapprompt@gap>| !gapinput@cm[1000];|
  [ [ 1, 2, 4, 9 ], [ 3, 6, 8 ], [ 5, 7 ] ]
  !gapprompt@gap>| !gapinput@Position(cm, last);|
  1000
\end{Verbatim}
 }

 

\subsection{\textcolor{Chapter }{NrPartitions}}
\logpage{[ 16, 2, 21 ]}\nobreak
\hyperdef{L}{X86933C4F795C4EBD}{}
{\noindent\textcolor{FuncColor}{$\triangleright$\enspace\texttt{NrPartitions({\mdseries\slshape n[, k]})\index{NrPartitions@\texttt{NrPartitions}}
\label{NrPartitions}
}\hfill{\scriptsize (function)}}\\


 returns the number of \texttt{Partitions(\mbox{\texttt{\mdseries\slshape set}},\mbox{\texttt{\mdseries\slshape k}})}. 
\begin{Verbatim}[commandchars=!@|,fontsize=\small,frame=single,label=Example]
  !gapprompt@gap>| !gapinput@Partitions( 7 );|
  [ [ 1, 1, 1, 1, 1, 1, 1 ], [ 2, 1, 1, 1, 1, 1 ], [ 2, 2, 1, 1, 1 ],
    [ 2, 2, 2, 1 ], [ 3, 1, 1, 1, 1 ], [ 3, 2, 1, 1 ], [ 3, 2, 2 ],
    [ 3, 3, 1 ], [ 4, 1, 1, 1 ], [ 4, 2, 1 ], [ 4, 3 ], [ 5, 1, 1 ],
    [ 5, 2 ], [ 6, 1 ], [ 7 ] ]
  !gapprompt@gap>| !gapinput@Partitions( 8, 3 );|
  [ [ 3, 3, 2 ], [ 4, 2, 2 ], [ 4, 3, 1 ], [ 5, 2, 1 ], [ 6, 1, 1 ] ]
  !gapprompt@gap>| !gapinput@NrPartitions( 7 );|
  15
  !gapprompt@gap>| !gapinput@NrPartitions( 100 );|
  190569292
\end{Verbatim}
 

 The function \texttt{OrderedPartitions} (\ref{OrderedPartitions}) is the ordered counterpart of \texttt{Partitions} (\ref{Partitions}). }

 

\subsection{\textcolor{Chapter }{OrderedPartitions}}
\logpage{[ 16, 2, 22 ]}\nobreak
\hyperdef{L}{X820DF201871F2723}{}
{\noindent\textcolor{FuncColor}{$\triangleright$\enspace\texttt{OrderedPartitions({\mdseries\slshape n[, k]})\index{OrderedPartitions@\texttt{OrderedPartitions}}
\label{OrderedPartitions}
}\hfill{\scriptsize (function)}}\\


 returns the set of all ordered partitions \index{partitions!ordered, of an integer} \index{partitions!improper, of an integer} of the positive integer \mbox{\texttt{\mdseries\slshape n}} into sums with \mbox{\texttt{\mdseries\slshape k}} summands. If \mbox{\texttt{\mdseries\slshape k}} is not given it returns all ordered partitions of \mbox{\texttt{\mdseries\slshape set}} for all \mbox{\texttt{\mdseries\slshape k}}. 

 An \emph{ordered partition} is an ordered sum $n = p_1 + p_2 + \ldots + p_k$ of positive integers and is represented by the list $[ p_1, p_2, \ldots, p_k ]$. There are totally $2^{{n-1}}$ ordered partitions and ${{n-1 \choose k-1}}$ (see \texttt{Binomial} (\ref{Binomial})) ordered partitions with \mbox{\texttt{\mdseries\slshape k}} summands. 

 Do not call \texttt{OrderedPartitions} with an \mbox{\texttt{\mdseries\slshape n}} much larger than $15$, the list will simply become too large. }

 

\subsection{\textcolor{Chapter }{NrOrderedPartitions}}
\logpage{[ 16, 2, 23 ]}\nobreak
\hyperdef{L}{X80BB9F4982CA1E8B}{}
{\noindent\textcolor{FuncColor}{$\triangleright$\enspace\texttt{NrOrderedPartitions({\mdseries\slshape n[, k]})\index{NrOrderedPartitions@\texttt{NrOrderedPartitions}}
\label{NrOrderedPartitions}
}\hfill{\scriptsize (function)}}\\


 returns the number of \texttt{OrderedPartitions(\mbox{\texttt{\mdseries\slshape set}},\mbox{\texttt{\mdseries\slshape k}})}. 

 
\begin{Verbatim}[commandchars=!@|,fontsize=\small,frame=single,label=Example]
  !gapprompt@gap>| !gapinput@OrderedPartitions( 5 );|
  [ [ 1, 1, 1, 1, 1 ], [ 1, 1, 1, 2 ], [ 1, 1, 2, 1 ], [ 1, 1, 3 ],
    [ 1, 2, 1, 1 ], [ 1, 2, 2 ], [ 1, 3, 1 ], [ 1, 4 ], [ 2, 1, 1, 1 ],
    [ 2, 1, 2 ], [ 2, 2, 1 ], [ 2, 3 ], [ 3, 1, 1 ], [ 3, 2 ],
    [ 4, 1 ], [ 5 ] ]
  !gapprompt@gap>| !gapinput@OrderedPartitions( 6, 3 );|
  [ [ 1, 1, 4 ], [ 1, 2, 3 ], [ 1, 3, 2 ], [ 1, 4, 1 ], [ 2, 1, 3 ],
    [ 2, 2, 2 ], [ 2, 3, 1 ], [ 3, 1, 2 ], [ 3, 2, 1 ], [ 4, 1, 1 ] ]
  !gapprompt@gap>| !gapinput@NrOrderedPartitions(20);|
  524288
\end{Verbatim}
 

 The function \texttt{Partitions} (\ref{Partitions}) is the unordered counterpart of \texttt{OrderedPartitions} (\ref{OrderedPartitions}). }

 

\subsection{\textcolor{Chapter }{PartitionsGreatestLE}}
\logpage{[ 16, 2, 24 ]}\nobreak
\hyperdef{L}{X8009520C82942461}{}
{\noindent\textcolor{FuncColor}{$\triangleright$\enspace\texttt{PartitionsGreatestLE({\mdseries\slshape n, m})\index{PartitionsGreatestLE@\texttt{PartitionsGreatestLE}}
\label{PartitionsGreatestLE}
}\hfill{\scriptsize (function)}}\\


 returns the set of all (unordered) partitions of the integer \mbox{\texttt{\mdseries\slshape n}} having parts less or equal to the integer \mbox{\texttt{\mdseries\slshape m}}. For $\mbox{\texttt{\mdseries\slshape n}} = 0$ and $\mbox{\texttt{\mdseries\slshape m}} \geq 0$, this returns the empty partition \texttt{[[]]}. Negative arguments yield the empty list. 
\begin{Verbatim}[commandchars=!@|,fontsize=\small,frame=single,label=Example]
  !gapprompt@gap>| !gapinput@PartitionsGreatestLE( 6, 3 );|
  [ [ 1, 1, 1, 1, 1, 1 ], [ 2, 1, 1, 1, 1 ], [ 2, 2, 1, 1 ], [ 2, 2, 2 ],
    [ 3, 1, 1, 1 ], [ 3, 2, 1 ], [ 3, 3 ] ]
  !gapprompt@gap>| !gapinput@List( [ 0 .. 3 ], m -> PartitionsGreatestLE( 0, m ) );|
  [ [ [  ] ], [ [  ] ], [ [  ] ], [ [  ] ] ]
\end{Verbatim}
 }

 

\subsection{\textcolor{Chapter }{PartitionsGreatestEQ}}
\logpage{[ 16, 2, 25 ]}\nobreak
\hyperdef{L}{X7CB8D4FF8592A9BB}{}
{\noindent\textcolor{FuncColor}{$\triangleright$\enspace\texttt{PartitionsGreatestEQ({\mdseries\slshape n, m})\index{PartitionsGreatestEQ@\texttt{PartitionsGreatestEQ}}
\label{PartitionsGreatestEQ}
}\hfill{\scriptsize (function)}}\\


 returns the set of all (unordered) partitions of the integer \mbox{\texttt{\mdseries\slshape n}} having greatest part equal to the integer \mbox{\texttt{\mdseries\slshape m}}. Since partitions use positive parts, this returns the empty list if $\mbox{\texttt{\mdseries\slshape n}} \leq 0$ or $\mbox{\texttt{\mdseries\slshape m}} \leq 0$. 
\begin{Verbatim}[commandchars=!@|,fontsize=\small,frame=single,label=Example]
  !gapprompt@gap>| !gapinput@PartitionsGreatestEQ( 6, 3 );|
  [ [ 3, 1, 1, 1 ], [ 3, 2, 1 ], [ 3, 3 ] ]
  !gapprompt@gap>| !gapinput@PartitionsGreatestEQ( 0, 1 );|
  [  ]
\end{Verbatim}
 }

 

\subsection{\textcolor{Chapter }{RestrictedPartitions}}
\logpage{[ 16, 2, 26 ]}\nobreak
\hyperdef{L}{X7A70D4F3809494E7}{}
{\noindent\textcolor{FuncColor}{$\triangleright$\enspace\texttt{RestrictedPartitions({\mdseries\slshape n, set[, k]})\index{RestrictedPartitions@\texttt{RestrictedPartitions}}
\label{RestrictedPartitions}
}\hfill{\scriptsize (function)}}\\


 In the first form \texttt{RestrictedPartitions} returns the set of all restricted partitions \index{partitions!restricted, of an integer} of the positive integer \mbox{\texttt{\mdseries\slshape n}} into sums with \mbox{\texttt{\mdseries\slshape k}} summands with the summands of the partition coming from the set \mbox{\texttt{\mdseries\slshape set}}. If \mbox{\texttt{\mdseries\slshape k}} is not given all restricted partitions for all \mbox{\texttt{\mdseries\slshape k}} are returned. 

 A \emph{restricted partition} is like an ordinary partition (see \texttt{Partitions} (\ref{Partitions})) an unordered sum $n = p_1 + p_2 + \ldots + p_k$ of positive integers and is represented by the list $p = [ p_1, p_2, \ldots, p_k ]$, in nonincreasing order. The difference is that here the $p_i$ must be elements from the set \mbox{\texttt{\mdseries\slshape set}}, while for ordinary partitions they may be elements from \texttt{[ 1 .. n ]}. }

 

\subsection{\textcolor{Chapter }{NrRestrictedPartitions}}
\logpage{[ 16, 2, 27 ]}\nobreak
\hyperdef{L}{X800B43838742FBF4}{}
{\noindent\textcolor{FuncColor}{$\triangleright$\enspace\texttt{NrRestrictedPartitions({\mdseries\slshape n, set[, k]})\index{NrRestrictedPartitions@\texttt{NrRestrictedPartitions}}
\label{NrRestrictedPartitions}
}\hfill{\scriptsize (function)}}\\


 returns the number of \texttt{RestrictedPartitions(\mbox{\texttt{\mdseries\slshape n}},\mbox{\texttt{\mdseries\slshape set}},\mbox{\texttt{\mdseries\slshape k}})}. 

 
\begin{Verbatim}[commandchars=!@|,fontsize=\small,frame=single,label=Example]
  !gapprompt@gap>| !gapinput@RestrictedPartitions( 8, [1,3,5,7] );|
  [ [ 1, 1, 1, 1, 1, 1, 1, 1 ], [ 3, 1, 1, 1, 1, 1 ], [ 3, 3, 1, 1 ],
    [ 5, 1, 1, 1 ], [ 5, 3 ], [ 7, 1 ] ]
  !gapprompt@gap>| !gapinput@NrRestrictedPartitions(50,[1,2,5,10,20,50]);|
  451
\end{Verbatim}
 

 The last example tells us that there are 451 ways to return 50 pence change
using 1, 2, 5, 10, 20 and 50 pence coins. }

 

\subsection{\textcolor{Chapter }{SignPartition}}
\logpage{[ 16, 2, 28 ]}\nobreak
\hyperdef{L}{X7F4EDCCA780B469D}{}
{\noindent\textcolor{FuncColor}{$\triangleright$\enspace\texttt{SignPartition({\mdseries\slshape pi})\index{SignPartition@\texttt{SignPartition}}
\label{SignPartition}
}\hfill{\scriptsize (function)}}\\


 returns the sign of a permutation with cycle structure \mbox{\texttt{\mdseries\slshape pi}}. 

 This function actually describes a homomorphism from the symmetric group $S_n$ into the cyclic group of order 2, whose kernel is exactly the alternating
group $A_n$ (see \texttt{SignPerm} (\ref{SignPerm})). Partitions of sign 1 are called \emph{even} partitions while partitions of sign $-1$ are called \emph{odd}. 
\begin{Verbatim}[commandchars=!@|,fontsize=\small,frame=single,label=Example]
  !gapprompt@gap>| !gapinput@SignPartition([6,5,4,3,2,1]);|
  -1
\end{Verbatim}
 }

 

\subsection{\textcolor{Chapter }{AssociatedPartition}}
\logpage{[ 16, 2, 29 ]}\nobreak
\hyperdef{L}{X7DB9BEB6856EC03D}{}
{\noindent\textcolor{FuncColor}{$\triangleright$\enspace\texttt{AssociatedPartition({\mdseries\slshape pi})\index{AssociatedPartition@\texttt{AssociatedPartition}}
\label{AssociatedPartition}
}\hfill{\scriptsize (function)}}\\


 \texttt{AssociatedPartition} returns the associated partition of the partition \mbox{\texttt{\mdseries\slshape pi}} which is obtained by transposing the corresponding Young diagram. 

 
\begin{Verbatim}[commandchars=!@|,fontsize=\small,frame=single,label=Example]
  !gapprompt@gap>| !gapinput@AssociatedPartition([4,2,1]);|
  [ 3, 2, 1, 1 ]
  !gapprompt@gap>| !gapinput@AssociatedPartition([6]);|
  [ 1, 1, 1, 1, 1, 1 ]
\end{Verbatim}
 }

 

\subsection{\textcolor{Chapter }{PowerPartition}}
\logpage{[ 16, 2, 30 ]}\nobreak
\hyperdef{L}{X7A95D8A6820363A8}{}
{\noindent\textcolor{FuncColor}{$\triangleright$\enspace\texttt{PowerPartition({\mdseries\slshape pi, k})\index{PowerPartition@\texttt{PowerPartition}}
\label{PowerPartition}
}\hfill{\scriptsize (function)}}\\


 \texttt{PowerPartition} returns the partition corresponding to the \mbox{\texttt{\mdseries\slshape k}}\texttt{\symbol{45}}th power of a permutation with cycle structure \mbox{\texttt{\mdseries\slshape pi}}. 

 Each part $l$ of \mbox{\texttt{\mdseries\slshape pi}} is replaced by $d = \gcd(l, k)$ parts $l/d$. So if \mbox{\texttt{\mdseries\slshape pi}} is a partition of $n$ then $\mbox{\texttt{\mdseries\slshape pi}}^{\mbox{\texttt{\mdseries\slshape k}}}$ also is a partition of $n$. \texttt{PowerPartition} describes the power map \index{symmetric group!power map} of symmetric groups. 

 
\begin{Verbatim}[commandchars=!@|,fontsize=\small,frame=single,label=Example]
  !gapprompt@gap>| !gapinput@PowerPartition([6,5,4,3,2,1], 3);|
  [ 5, 4, 2, 2, 2, 2, 1, 1, 1, 1 ]
\end{Verbatim}
 }

 

\subsection{\textcolor{Chapter }{PartitionTuples}}
\logpage{[ 16, 2, 31 ]}\nobreak
\hyperdef{L}{X877D997B7F66A119}{}
{\noindent\textcolor{FuncColor}{$\triangleright$\enspace\texttt{PartitionTuples({\mdseries\slshape n, r})\index{PartitionTuples@\texttt{PartitionTuples}}
\label{PartitionTuples}
}\hfill{\scriptsize (function)}}\\


 \texttt{PartitionTuples} returns the list of all \mbox{\texttt{\mdseries\slshape r}}\texttt{\symbol{45}}tuples of partitions which together form a partition of \mbox{\texttt{\mdseries\slshape n}}. 

 \mbox{\texttt{\mdseries\slshape r}}\texttt{\symbol{45}}tuples of partitions describe the classes and the
characters of wreath products of groups with \mbox{\texttt{\mdseries\slshape r}} conjugacy classes with the symmetric group on \mbox{\texttt{\mdseries\slshape n}} points, see \texttt{CharacterTableWreathSymmetric} (\ref{CharacterTableWreathSymmetric}) and \texttt{CharacterValueWreathSymmetric} (\ref{CharacterValueWreathSymmetric}). }

 

\subsection{\textcolor{Chapter }{NrPartitionTuples}}
\logpage{[ 16, 2, 32 ]}\nobreak
\hyperdef{L}{X7F44AD098561DE32}{}
{\noindent\textcolor{FuncColor}{$\triangleright$\enspace\texttt{NrPartitionTuples({\mdseries\slshape n, r})\index{NrPartitionTuples@\texttt{NrPartitionTuples}}
\label{NrPartitionTuples}
}\hfill{\scriptsize (function)}}\\


 returns the number of \texttt{PartitionTuples( \mbox{\texttt{\mdseries\slshape n}}, \mbox{\texttt{\mdseries\slshape r}} )}. 
\begin{Verbatim}[commandchars=!@|,fontsize=\small,frame=single,label=Example]
  !gapprompt@gap>| !gapinput@PartitionTuples(3, 2);|
  [ [ [ 1, 1, 1 ], [  ] ], [ [ 1, 1 ], [ 1 ] ], [ [ 1 ], [ 1, 1 ] ],
    [ [  ], [ 1, 1, 1 ] ], [ [ 2, 1 ], [  ] ], [ [ 1 ], [ 2 ] ],
    [ [ 2 ], [ 1 ] ], [ [  ], [ 2, 1 ] ], [ [ 3 ], [  ] ],
    [ [  ], [ 3 ] ] ]
\end{Verbatim}
 }

 

\subsection{\textcolor{Chapter }{BetaSet}}
\logpage{[ 16, 2, 33 ]}\nobreak
\hyperdef{L}{X8796C1D783ED9CB4}{}
{\noindent\textcolor{FuncColor}{$\triangleright$\enspace\texttt{BetaSet({\mdseries\slshape alpha})\index{BetaSet@\texttt{BetaSet}}
\label{BetaSet}
}\hfill{\scriptsize (function)}}\\


 For a list \mbox{\texttt{\mdseries\slshape alpha}} that describes a partition of a nonnegative integer (see \texttt{Partitions} (\ref{Partitions})), \texttt{BetaSet} returns the list of integers obtained by reversing the order of \mbox{\texttt{\mdseries\slshape alpha}} and then adding the sequence \texttt{[ 0, 1, 2, ... ]} of the same length, cf. \cite[Section 2.7]{JK81}. 

 
\begin{Verbatim}[commandchars=!@|,fontsize=\small,frame=single,label=Example]
  !gapprompt@gap>| !gapinput@BetaSet( [ 4, 2, 1 ] );|
  [ 1, 3, 6 ]
  !gapprompt@gap>| !gapinput@BetaSet( [] );|
  [  ]
\end{Verbatim}
 }

 }

  
\section{\textcolor{Chapter }{Fibonacci and Lucas Sequences}}\label{Fibonacci and Lucas Sequences}
\logpage{[ 16, 3, 0 ]}
\hyperdef{L}{X83DC50B67D74E674}{}
{
  

\subsection{\textcolor{Chapter }{Fibonacci}}
\logpage{[ 16, 3, 1 ]}\nobreak
\hyperdef{L}{X85AE1D70803A886C}{}
{\noindent\textcolor{FuncColor}{$\triangleright$\enspace\texttt{Fibonacci({\mdseries\slshape n})\index{Fibonacci@\texttt{Fibonacci}}
\label{Fibonacci}
}\hfill{\scriptsize (function)}}\\


 returns the \mbox{\texttt{\mdseries\slshape n}}th number of the \emph{Fibonacci sequence}. The Fibonacci sequence $F_n$ \index{sequence!Fibonacci} is defined by the initial conditions $F_1 = F_2 = 1$ and the recurrence relation $F_{{n+2}} = F_{{n+1}} + F_n$. For negative $n$ we define $F_n = (-1)^{{n+1}} F_{{-n}}$, which is consistent with the recurrence relation. 

 Using generating functions one can prove that $F_n = \phi^n - 1/\phi^n$, where $\phi$ is $(\sqrt{{5}} + 1)/2$, i.e., one root of $x^2 - x - 1 = 0$. Fibonacci numbers have the property $\gcd( F_m, F_n ) = F_{{\gcd(m,n)}}$. But a pair of Fibonacci numbers requires more division steps in Euclid's
algorithm (see{\nobreakspace}\texttt{Gcd} (\ref{Gcd:for (a ring and) several elements})) than any other pair of integers of the same size. \texttt{Fibonacci(\mbox{\texttt{\mdseries\slshape k}})} is the special case \texttt{Lucas(1,\texttt{\symbol{45}}1,\mbox{\texttt{\mdseries\slshape k}})[1]} (see \texttt{Lucas} (\ref{Lucas})). 

 
\begin{Verbatim}[commandchars=!@|,fontsize=\small,frame=single,label=Example]
  !gapprompt@gap>| !gapinput@Fibonacci( 10 );|
  55
  !gapprompt@gap>| !gapinput@Fibonacci( 35 );|
  9227465
  !gapprompt@gap>| !gapinput@Fibonacci( -10 );|
  -55
\end{Verbatim}
 }

 

\subsection{\textcolor{Chapter }{Lucas}}
\logpage{[ 16, 3, 2 ]}\nobreak
\hyperdef{L}{X7830A03181D67192}{}
{\noindent\textcolor{FuncColor}{$\triangleright$\enspace\texttt{Lucas({\mdseries\slshape P, Q, k})\index{Lucas@\texttt{Lucas}}
\label{Lucas}
}\hfill{\scriptsize (function)}}\\


 returns the \mbox{\texttt{\mdseries\slshape k}}\texttt{\symbol{45}}th values of the \emph{Lucas sequence} \index{sequence!Lucas} with parameters \mbox{\texttt{\mdseries\slshape P}} and \mbox{\texttt{\mdseries\slshape Q}}, which must be integers, as a list of three integers. If \mbox{\texttt{\mdseries\slshape k}} is a negative integer, then the values of the Lucas sequence may be
nonintegral rational numbers, with denominator roughly \mbox{\texttt{\mdseries\slshape Q}}\texttt{\symbol{94}}\mbox{\texttt{\mdseries\slshape k}}. 

 Let $\alpha, \beta$ be the two roots of $x^2 - P x + Q$ then we define \texttt{Lucas( \mbox{\texttt{\mdseries\slshape P}}, \mbox{\texttt{\mdseries\slshape Q}}, \mbox{\texttt{\mdseries\slshape k}} )[1]} $= U_k = (\alpha^k - \beta^k) / (\alpha - \beta)$ and \texttt{Lucas( \mbox{\texttt{\mdseries\slshape P}}, \mbox{\texttt{\mdseries\slshape Q}}, \mbox{\texttt{\mdseries\slshape k}} )[2]} $= V_k = (\alpha^k + \beta^k)$ and as a convenience \texttt{Lucas( \mbox{\texttt{\mdseries\slshape P}}, \mbox{\texttt{\mdseries\slshape Q}}, \mbox{\texttt{\mdseries\slshape k}} )[3]} $= Q^k$. 

 The following recurrence relations are easily derived from the definition $U_0 = 0, U_1 = 1, U_k = P U_{{k-1}} - Q U_{{k-2}}$ and $V_0 = 2, V_1 = P, V_k = P V_{{k-1}} - Q V_{{k-2}}$. Those relations are actually used to define \texttt{Lucas} if $\alpha = \beta$. 

 Also the more complex relations used in \texttt{Lucas} can be easily derived $U_{2k} = U_k V_k$, $U_{{2k+1}} = (P U_{2k} + V_{2k}) / 2$ and $V_{2k} = V_k^2 - 2 Q^k$, $V_{{2k+1}} = ((P^2-4Q) U_{2k} + P V_{2k}) / 2$. 

 \texttt{Fibonacci(\mbox{\texttt{\mdseries\slshape k}})} (see \texttt{Fibonacci} (\ref{Fibonacci})) is simply \texttt{Lucas(1,\texttt{\symbol{45}}1,\mbox{\texttt{\mdseries\slshape k}})[1]}. In an abuse of notation, the sequence \texttt{Lucas(1,\texttt{\symbol{45}}1,\mbox{\texttt{\mdseries\slshape k}})[2]} is sometimes called the Lucas sequence. 

 
\begin{Verbatim}[commandchars=!@|,fontsize=\small,frame=single,label=Example]
  !gapprompt@gap>| !gapinput@List( [0..10], i -> Lucas(1,-2,i)[1] );     # 2^k - (-1)^k)/3|
  [ 0, 1, 1, 3, 5, 11, 21, 43, 85, 171, 341 ]
  !gapprompt@gap>| !gapinput@List( [0..10], i -> Lucas(1,-2,i)[2] );     # 2^k + (-1)^k|
  [ 2, 1, 5, 7, 17, 31, 65, 127, 257, 511, 1025 ]
  !gapprompt@gap>| !gapinput@List( [0..10], i -> Lucas(1,-1,i)[1] );     # Fibonacci sequence|
  [ 0, 1, 1, 2, 3, 5, 8, 13, 21, 34, 55 ]
  !gapprompt@gap>| !gapinput@List( [0..10], i -> Lucas(2,1,i)[1] );      # the roots are equal|
  [ 0, 1, 2, 3, 4, 5, 6, 7, 8, 9, 10 ]
\end{Verbatim}
 }

 }

  
\section{\textcolor{Chapter }{Permanent of a Matrix}}\label{Permanent of a Matrix}
\logpage{[ 16, 4, 0 ]}
\hyperdef{L}{X821888E77EB43F67}{}
{
  

\subsection{\textcolor{Chapter }{Permanent}}
\logpage{[ 16, 4, 1 ]}\nobreak
\hyperdef{L}{X7F0942DD83BBAB7A}{}
{\noindent\textcolor{FuncColor}{$\triangleright$\enspace\texttt{Permanent({\mdseries\slshape mat})\index{Permanent@\texttt{Permanent}}
\label{Permanent}
}\hfill{\scriptsize (attribute)}}\\


 returns the \emph{permanent} of the matrix \mbox{\texttt{\mdseries\slshape mat}}. The permanent is defined by $\sum_{{p \in Sym(n)}} \prod_{{i = 1}}^n mat[i][i^p]$. 

 Note the similarity of the definition of the permanent to the definition of
the determinant (see{\nobreakspace}\texttt{DeterminantMat} (\ref{DeterminantMat})). In fact the only difference is the missing sign of the permutation. However
the permanent is quite unlike the determinant, for example it is not
multilinear or alternating. It has however important combinatorial properties. 

 
\begin{Verbatim}[commandchars=!@|,fontsize=\small,frame=single,label=Example]
  !gapprompt@gap>| !gapinput@Permanent( [[0,1,1,1],|
  !gapprompt@>| !gapinput@     [1,0,1,1],|
  !gapprompt@>| !gapinput@     [1,1,0,1],|
  !gapprompt@>| !gapinput@     [1,1,1,0]] );  # inefficient way to compute NrDerangements([1..4])|
  9
  !gapprompt@gap>| !gapinput@# 24 permutations fit the projective plane of order 2:|
  !gapprompt@gap>| !gapinput@Permanent( [[1,1,0,1,0,0,0],|
  !gapprompt@>| !gapinput@     [0,1,1,0,1,0,0],|
  !gapprompt@>| !gapinput@     [0,0,1,1,0,1,0],|
  !gapprompt@>| !gapinput@     [0,0,0,1,1,0,1],|
  !gapprompt@>| !gapinput@     [1,0,0,0,1,1,0],|
  !gapprompt@>| !gapinput@     [0,1,0,0,0,1,1],|
  !gapprompt@>| !gapinput@     [1,0,1,0,0,0,1]] );|
  24
\end{Verbatim}
 }

 }

 }

         
\chapter{\textcolor{Chapter }{Rational Numbers}}\label{Rational Numbers}
\logpage{[ 17, 0, 0 ]}
\hyperdef{L}{X87003045878E74DF}{}
{
  The \emph{rationals} form a very important field. On the one hand it is the quotient field of the
integers (see chapter{\nobreakspace}\ref{Integers}). On the other hand it is the prime field of the fields of characteristic
zero (see chapter{\nobreakspace}\ref{Abelian Number Fields}). 

 The former comment suggests the representation actually used. A rational is
represented as a pair of integers, called \emph{numerator} and \emph{denominator}. Numerator and denominator are \emph{reduced}, i.e., their greatest common divisor is 1. If the denominator is 1, the
rational is in fact an integer and is represented as such. The numerator holds
the sign of the rational, thus the denominator is always positive. 

 Because the underlying integer arithmetic can compute with arbitrary size
integers, the rational arithmetic is always exact, even for rationals whose
numerators and denominators have thousands of digits. 

 
\begin{Verbatim}[commandchars=!@|,fontsize=\small,frame=single,label=Example]
  !gapprompt@gap>| !gapinput@2/3;|
  2/3
  !gapprompt@gap>| !gapinput@66/123;  # numerator and denominator are made relatively prime|
  22/41
  !gapprompt@gap>| !gapinput@17/-13;  # the numerator carries the sign;|
  -17/13
  !gapprompt@gap>| !gapinput@121/11;  # rationals with denominator 1 (when canceled) are integers|
  11
\end{Verbatim}
  
\section{\textcolor{Chapter }{Rationals: Global Variables}}\label{sect:Rationals: Global Variables}
\logpage{[ 17, 1, 0 ]}
\hyperdef{L}{X7A76497986DA921F}{}
{
  

\subsection{\textcolor{Chapter }{Rationals}}
\logpage{[ 17, 1, 1 ]}\nobreak
\hyperdef{L}{X7B6029D18570C08A}{}
{\noindent\textcolor{FuncColor}{$\triangleright$\enspace\texttt{Rationals\index{Rationals@\texttt{Rationals}}
\label{Rationals}
}\hfill{\scriptsize (global variable)}}\\
\noindent\textcolor{FuncColor}{$\triangleright$\enspace\texttt{IsRationals({\mdseries\slshape obj})\index{IsRationals@\texttt{IsRationals}}
\label{IsRationals}
}\hfill{\scriptsize (filter)}}\\


 \texttt{Rationals} is the field ${\ensuremath{\mathbb Q}}$ of rational integers, as a set of cyclotomic numbers, see
Chapter{\nobreakspace}\ref{Cyclotomic Numbers} for basic operations, Functions for the field \texttt{Rationals} can be found in the chapters{\nobreakspace}\ref{Fields and Division Rings} and{\nobreakspace}\ref{Abelian Number Fields}. 

 \texttt{IsRationals} returns \texttt{true} for a prime field that consists of cyclotomic numbers {\textendash}for example
the \textsf{GAP} object \texttt{Rationals}{\textendash} and \texttt{false} for all other \textsf{GAP} objects. 

 
\begin{Verbatim}[commandchars=!@|,fontsize=\small,frame=single,label=Example]
  !gapprompt@gap>| !gapinput@Size( Rationals ); 2/3 in Rationals;|
  infinity
  true
\end{Verbatim}
 }

 }

  
\section{\textcolor{Chapter }{Elementary Operations for Rationals}}\label{Elementary Operations for Rationals}
\logpage{[ 17, 2, 0 ]}
\hyperdef{L}{X826E2AA88679B3DF}{}
{
  

\subsection{\textcolor{Chapter }{IsRat}}
\logpage{[ 17, 2, 1 ]}\nobreak
\hyperdef{L}{X7ED018F5794935F7}{}
{\noindent\textcolor{FuncColor}{$\triangleright$\enspace\texttt{IsRat({\mdseries\slshape obj})\index{IsRat@\texttt{IsRat}}
\label{IsRat}
}\hfill{\scriptsize (Category)}}\\


 \index{test!for a rational} Every rational number lies in the category \texttt{IsRat}, which is a subcategory of \texttt{IsCyc} (\ref{IsCyc}). 

 
\begin{Verbatim}[commandchars=!@|,fontsize=\small,frame=single,label=Example]
  !gapprompt@gap>| !gapinput@IsRat( 2/3 );|
  true
  !gapprompt@gap>| !gapinput@IsRat( 17/-13 );|
  true
  !gapprompt@gap>| !gapinput@IsRat( 11 );|
  true
  !gapprompt@gap>| !gapinput@IsRat( IsRat );  # `IsRat' is a function, not a rational|
  false
\end{Verbatim}
 }

 

\subsection{\textcolor{Chapter }{IsPosRat}}
\logpage{[ 17, 2, 2 ]}\nobreak
\hyperdef{L}{X7BD6E170840F045D}{}
{\noindent\textcolor{FuncColor}{$\triangleright$\enspace\texttt{IsPosRat({\mdseries\slshape obj})\index{IsPosRat@\texttt{IsPosRat}}
\label{IsPosRat}
}\hfill{\scriptsize (Category)}}\\


 Every positive rational number lies in the category \texttt{IsPosRat}. }

 

\subsection{\textcolor{Chapter }{IsNegRat}}
\logpage{[ 17, 2, 3 ]}\nobreak
\hyperdef{L}{X81179AC87AC951A8}{}
{\noindent\textcolor{FuncColor}{$\triangleright$\enspace\texttt{IsNegRat({\mdseries\slshape obj})\index{IsNegRat@\texttt{IsNegRat}}
\label{IsNegRat}
}\hfill{\scriptsize (Category)}}\\


 Every negative rational number lies in the category \texttt{IsNegRat}. }

 

\subsection{\textcolor{Chapter }{NumeratorRat}}
\logpage{[ 17, 2, 4 ]}\nobreak
\hyperdef{L}{X7D830E7482E7F528}{}
{\noindent\textcolor{FuncColor}{$\triangleright$\enspace\texttt{NumeratorRat({\mdseries\slshape rat})\index{NumeratorRat@\texttt{NumeratorRat}}
\label{NumeratorRat}
}\hfill{\scriptsize (function)}}\\


 \index{numerator!of a rational} \texttt{NumeratorRat} returns the numerator of the rational \mbox{\texttt{\mdseries\slshape rat}}. Because the numerator holds the sign of the rational it may be any integer.
Integers are rationals with denominator $1$, thus \texttt{NumeratorRat} is the identity function for integers. 

 
\begin{Verbatim}[commandchars=!@|,fontsize=\small,frame=single,label=Example]
  !gapprompt@gap>| !gapinput@NumeratorRat( 2/3 );|
  2
  !gapprompt@gap>| !gapinput@# numerator and denominator are made relatively prime:|
  !gapprompt@gap>| !gapinput@NumeratorRat( 66/123 );|
  22
  !gapprompt@gap>| !gapinput@NumeratorRat( 17/-13 );  # numerator holds the sign of the rational|
  -17
  !gapprompt@gap>| !gapinput@NumeratorRat( 11 );      # integers are rationals with denominator 1|
  11
\end{Verbatim}
 }

 

\subsection{\textcolor{Chapter }{DenominatorRat}}
\logpage{[ 17, 2, 5 ]}\nobreak
\hyperdef{L}{X81F6B5877A81E727}{}
{\noindent\textcolor{FuncColor}{$\triangleright$\enspace\texttt{DenominatorRat({\mdseries\slshape rat})\index{DenominatorRat@\texttt{DenominatorRat}}
\label{DenominatorRat}
}\hfill{\scriptsize (function)}}\\


 \index{denominator!of a rational} \texttt{DenominatorRat} returns the denominator of the rational \mbox{\texttt{\mdseries\slshape rat}}. Because the numerator holds the sign of the rational the denominator is
always a positive integer. Integers are rationals with the denominator 1, thus \texttt{DenominatorRat} returns 1 for integers. 

 
\begin{Verbatim}[commandchars=!@|,fontsize=\small,frame=single,label=Example]
  !gapprompt@gap>| !gapinput@DenominatorRat( 2/3 );|
  3
  !gapprompt@gap>| !gapinput@# numerator and denominator are made relatively prime:|
  !gapprompt@gap>| !gapinput@DenominatorRat( 66/123 );|
  41
  !gapprompt@gap>| !gapinput@# the denominator holds the sign of the rational:|
  !gapprompt@gap>| !gapinput@DenominatorRat( 17/-13 );|
  13
  !gapprompt@gap>| !gapinput@DenominatorRat( 11 ); # integers are rationals with denominator 1|
  1
\end{Verbatim}
 }

 

\subsection{\textcolor{Chapter }{Rat}}
\logpage{[ 17, 2, 6 ]}\nobreak
\hyperdef{L}{X7EB4C646806A2BDE}{}
{\noindent\textcolor{FuncColor}{$\triangleright$\enspace\texttt{Rat({\mdseries\slshape elm})\index{Rat@\texttt{Rat}}
\label{Rat}
}\hfill{\scriptsize (attribute)}}\\


 \texttt{Rat} returns a rational number \mbox{\texttt{\mdseries\slshape rat}} whose meaning depends on the type of \mbox{\texttt{\mdseries\slshape elm}}. 

 If \mbox{\texttt{\mdseries\slshape elm}} is a string consisting of digits \texttt{'0'}, \texttt{'1'}, $\ldots$, \texttt{'9'} and \texttt{'\texttt{\symbol{45}}'} (at the first position), \texttt{'/'} and the decimal dot \texttt{'.'} then \mbox{\texttt{\mdseries\slshape rat}} is the rational described by this string. If \mbox{\texttt{\mdseries\slshape elm}} is a rational number, then \texttt{Rat} returns \mbox{\texttt{\mdseries\slshape elm}}. The operation \texttt{String} (\ref{String}) can be used to compute a string for rational numbers, in fact for all
cyclotomics. 

 
\begin{Verbatim}[commandchars=!@|,fontsize=\small,frame=single,label=Example]
  !gapprompt@gap>| !gapinput@Rat( "1/2" );  Rat( "35/14" );  Rat( "35/-27" );  Rat( "3.14159" );|
  1/2
  5/2
  -35/27
  314159/100000
\end{Verbatim}
 }

 

\subsection{\textcolor{Chapter }{Random (for rationals)}}
\logpage{[ 17, 2, 7 ]}\nobreak
\hyperdef{L}{X7C8F8693825C28A4}{}
{\noindent\textcolor{FuncColor}{$\triangleright$\enspace\texttt{Random({\mdseries\slshape Rationals})\index{Random@\texttt{Random}!for rationals}
\label{Random:for rationals}
}\hfill{\scriptsize (operation)}}\\


 \texttt{Random} for rationals returns pseudo random rationals which are the quotient of two
random integers. See the description of \texttt{Random} (\ref{Random:for integers}) for details. (Also see{\nobreakspace}\texttt{Random} (\ref{Random:for a list or collection}).) }

 }

 }

               
\chapter{\textcolor{Chapter }{Cyclotomic Numbers}}\label{Cyclotomic Numbers}
\logpage{[ 18, 0, 0 ]}
\hyperdef{L}{X7DFC03C187DE4841}{}
{
  \index{type!cyclotomic} \index{irrationalities} \index{cyclotomic field elements} \textsf{GAP} admits computations in abelian extension fields of the rational number field ${\ensuremath{\mathbb Q}}$, that is fields with abelian Galois group over ${\ensuremath{\mathbb Q}}$. These fields are subfields of \emph{cyclotomic fields} ${\ensuremath{\mathbb Q}}(e_n)$ where $e_n = \exp(2 \pi i/n)$ is a primitive complex $n$\texttt{\symbol{45}}th root of unity. The elements of these fields are called \emph{cyclotomics}. 

 Information concerning operations for domains of cyclotomics, for example
certain integral bases of fields of cyclotomics, can be found in
Chapter{\nobreakspace}\ref{Abelian Number Fields}. For more general operations that take a field extension as a
{\textendash}possibly optional{\textendash} argument, e.g., \texttt{Trace} (\ref{Trace:for a field element}) or \texttt{Coefficients} (\ref{Coefficients}), see Chapter{\nobreakspace}\ref{Fields and Division Rings}.  
\section{\textcolor{Chapter }{Operations for Cyclotomics}}\label{Operations for Cyclotomics}
\logpage{[ 18, 1, 0 ]}
\hyperdef{L}{X79E25C3085AA568F}{}
{
  

\subsection{\textcolor{Chapter }{E}}
\logpage{[ 18, 1, 1 ]}\nobreak
\hyperdef{L}{X8631458886314588}{}
{\noindent\textcolor{FuncColor}{$\triangleright$\enspace\texttt{E({\mdseries\slshape n})\index{E@\texttt{E}}
\label{E}
}\hfill{\scriptsize (operation)}}\\


 \index{roots of unity} \texttt{E} returns the primitive \mbox{\texttt{\mdseries\slshape n}}\texttt{\symbol{45}}th root of unity $e_n = \exp(2\pi i/n)$. Cyclotomics are usually entered as sums of roots of unity, with rational
coefficients, and irrational cyclotomics are displayed in such a way. (For
special cyclotomics, see{\nobreakspace}\ref{ATLAS Irrationalities}.) 

 
\begin{Verbatim}[commandchars=!@|,fontsize=\small,frame=single,label=Example]
  !gapprompt@gap>| !gapinput@E(9); E(9)^3; E(6); E(12) / 3;|
  -E(9)^4-E(9)^7
  E(3)
  -E(3)^2
  -1/3*E(12)^7
\end{Verbatim}
 

 A particular basis is used to express cyclotomics, see{\nobreakspace}\ref{Integral Bases of Abelian Number Fields}; note that \texttt{E(9)} is \emph{not} a basis element, as the above example shows. }

 

\subsection{\textcolor{Chapter }{Cyclotomics}}
\logpage{[ 18, 1, 2 ]}\nobreak
\hyperdef{L}{X863D1E017BC9EB7F}{}
{\noindent\textcolor{FuncColor}{$\triangleright$\enspace\texttt{Cyclotomics\index{Cyclotomics@\texttt{Cyclotomics}}
\label{Cyclotomics}
}\hfill{\scriptsize (global variable)}}\\


 is the domain of all cyclotomics. 

 
\begin{Verbatim}[commandchars=!@|,fontsize=\small,frame=single,label=Example]
  !gapprompt@gap>| !gapinput@E(9) in Cyclotomics; 37 in Cyclotomics; true in Cyclotomics;|
  true
  true
  false
\end{Verbatim}
 

 As the cyclotomics are field elements, the usual arithmetic operators \texttt{+}, \texttt{\texttt{\symbol{45}}}, \texttt{*} and \texttt{/} (and \texttt{\texttt{\symbol{94}}} to take powers by integers) are applicable. Note that \texttt{\texttt{\symbol{94}}} does \emph{not} denote the conjugation of group elements, so it is \emph{not} possible to explicitly construct groups of cyclotomics. (However, it is
possible to compute the inverse and the multiplicative order of a nonzero
cyclotomic.) Also, taking the $k$\texttt{\symbol{45}}th power of a root of unity $z$ defines a Galois automorphism if and only if $k$ is coprime to the conductor (see \texttt{Conductor} (\ref{Conductor:for a cyclotomic})) of $z$. 

 
\begin{Verbatim}[commandchars=!@|,fontsize=\small,frame=single,label=Example]
  !gapprompt@gap>| !gapinput@E(5) + E(3); (E(5) + E(5)^4) ^ 2; E(5) / E(3); E(5) * E(3);|
  -E(15)^2-2*E(15)^8-E(15)^11-E(15)^13-E(15)^14
  -2*E(5)-E(5)^2-E(5)^3-2*E(5)^4
  E(15)^13
  E(15)^8
  !gapprompt@gap>| !gapinput@Order( E(5) ); Order( 1+E(5) );|
  5
  infinity
\end{Verbatim}
 }

 

\subsection{\textcolor{Chapter }{IsCyclotomic}}
\logpage{[ 18, 1, 3 ]}\nobreak
\hyperdef{L}{X841C425281A6F775}{}
{\noindent\textcolor{FuncColor}{$\triangleright$\enspace\texttt{IsCyclotomic({\mdseries\slshape obj})\index{IsCyclotomic@\texttt{IsCyclotomic}}
\label{IsCyclotomic}
}\hfill{\scriptsize (Category)}}\\
\noindent\textcolor{FuncColor}{$\triangleright$\enspace\texttt{IsCyc({\mdseries\slshape obj})\index{IsCyc@\texttt{IsCyc}}
\label{IsCyc}
}\hfill{\scriptsize (Category)}}\\


 \index{CyclotomicsFamily@\texttt{CyclotomicsFamily}} Every object in the family \texttt{CyclotomicsFamily} lies in the category \texttt{IsCyclotomic}. This covers integers, rationals, proper cyclotomics, the object \texttt{infinity} (\ref{infinity}), and unknowns (see Chapter{\nobreakspace}\ref{Unknowns}). All these objects except \texttt{infinity} (\ref{infinity}) and unknowns lie also in the category \texttt{IsCyc}, \texttt{infinity} (\ref{infinity}) lies in (and can be detected from) the category \texttt{IsInfinity} (\ref{IsInfinity}), and unknowns lie in \texttt{IsUnknown} (\ref{IsUnknown}). 

 
\begin{Verbatim}[commandchars=!@|,fontsize=\small,frame=single,label=Example]
  !gapprompt@gap>| !gapinput@IsCyclotomic(0); IsCyclotomic(1/2*E(3)); IsCyclotomic( infinity );|
  true
  true
  true
  !gapprompt@gap>| !gapinput@IsCyc(0); IsCyc(1/2*E(3)); IsCyc( infinity );|
  true
  true
  false
\end{Verbatim}
 }

 

\subsection{\textcolor{Chapter }{IsIntegralCyclotomic}}
\logpage{[ 18, 1, 4 ]}\nobreak
\hyperdef{L}{X869750DA81EA0E67}{}
{\noindent\textcolor{FuncColor}{$\triangleright$\enspace\texttt{IsIntegralCyclotomic({\mdseries\slshape obj})\index{IsIntegralCyclotomic@\texttt{IsIntegralCyclotomic}}
\label{IsIntegralCyclotomic}
}\hfill{\scriptsize (property)}}\\


 A cyclotomic is called \emph{integral} or a \emph{cyclotomic integer} if all coefficients of its minimal polynomial over the rationals are integers.
Since the underlying basis of the external representation of cyclotomics is an
integral basis (see{\nobreakspace}\ref{Integral Bases of Abelian Number Fields}), the subring of cyclotomic integers in a cyclotomic field is formed by those
cyclotomics for which the external representation is a list of integers. For
example, square roots of integers are cyclotomic integers (see{\nobreakspace}\ref{ATLAS Irrationalities}), any root of unity is a cyclotomic integer, character values are always
cyclotomic integers, but all rationals which are not integers are not
cyclotomic integers. 

 
\begin{Verbatim}[commandchars=!@|,fontsize=\small,frame=single,label=Example]
  !gapprompt@gap>| !gapinput@r:= ER( 5 );               # The square root of 5 ...|
  E(5)-E(5)^2-E(5)^3+E(5)^4
  !gapprompt@gap>| !gapinput@IsIntegralCyclotomic( r ); # ... is a cyclotomic integer.|
  true
  !gapprompt@gap>| !gapinput@r2:= 1/2 * r;              # This is not a cyclotomic integer, ...|
  1/2*E(5)-1/2*E(5)^2-1/2*E(5)^3+1/2*E(5)^4
  !gapprompt@gap>| !gapinput@IsIntegralCyclotomic( r2 );|
  false
  !gapprompt@gap>| !gapinput@r3:= 1/2 * r - 1/2;        # ... but this is one.|
  E(5)+E(5)^4
  !gapprompt@gap>| !gapinput@IsIntegralCyclotomic( r3 );|
  true
\end{Verbatim}
 }

 

\subsection{\textcolor{Chapter }{Int (for a cyclotomic)}}
\logpage{[ 18, 1, 5 ]}\nobreak
\hyperdef{L}{X7DD6B95F79321D23}{}
{\noindent\textcolor{FuncColor}{$\triangleright$\enspace\texttt{Int({\mdseries\slshape cyc})\index{Int@\texttt{Int}!for a cyclotomic}
\label{Int:for a cyclotomic}
}\hfill{\scriptsize (method)}}\\


 The operation \texttt{Int} can be used to find a cyclotomic integer near to an arbitrary cyclotomic, by
applying \texttt{Int} (\ref{Int}) to the coefficients. 

 
\begin{Verbatim}[commandchars=!@|,fontsize=\small,frame=single,label=Example]
  !gapprompt@gap>| !gapinput@Int( E(5)+1/2*E(5)^2 ); Int( 2/3*E(7)-3/2*E(4) );|
  E(5)
  -E(4)
\end{Verbatim}
 }

 

\subsection{\textcolor{Chapter }{String (for a cyclotomic)}}
\logpage{[ 18, 1, 6 ]}\nobreak
\hyperdef{L}{X7CBA6CB678E2B143}{}
{\noindent\textcolor{FuncColor}{$\triangleright$\enspace\texttt{String({\mdseries\slshape cyc})\index{String@\texttt{String}!for a cyclotomic}
\label{String:for a cyclotomic}
}\hfill{\scriptsize (method)}}\\


 The operation \texttt{String} returns for a cyclotomic \mbox{\texttt{\mdseries\slshape cyc}} a string corresponding to the way the cyclotomic is printed by \texttt{ViewObj} (\ref{ViewObj}) and \texttt{PrintObj} (\ref{PrintObj}). 

 
\begin{Verbatim}[commandchars=!@|,fontsize=\small,frame=single,label=Example]
  !gapprompt@gap>| !gapinput@String( E(5)+1/2*E(5)^2 ); String( 17/3 );|
  "E(5)+1/2*E(5)^2"
  "17/3"
\end{Verbatim}
 }

 

\subsection{\textcolor{Chapter }{Conductor (for a cyclotomic)}}
\logpage{[ 18, 1, 7 ]}\nobreak
\hyperdef{L}{X815D6EC57CBA9827}{}
{\noindent\textcolor{FuncColor}{$\triangleright$\enspace\texttt{Conductor({\mdseries\slshape cyc})\index{Conductor@\texttt{Conductor}!for a cyclotomic}
\label{Conductor:for a cyclotomic}
}\hfill{\scriptsize (attribute)}}\\
\noindent\textcolor{FuncColor}{$\triangleright$\enspace\texttt{Conductor({\mdseries\slshape C})\index{Conductor@\texttt{Conductor}!for a collection of cyclotomics}
\label{Conductor:for a collection of cyclotomics}
}\hfill{\scriptsize (attribute)}}\\


 For an element \mbox{\texttt{\mdseries\slshape cyc}} of a cyclotomic field, \texttt{Conductor} returns the smallest integer $n$ such that \mbox{\texttt{\mdseries\slshape cyc}} is contained in the $n$\texttt{\symbol{45}}th cyclotomic field. For a collection \mbox{\texttt{\mdseries\slshape C}} of cyclotomics (for example a dense list of cyclotomics or a field of
cyclotomics), \texttt{Conductor} returns the smallest integer $n$ such that all elements of \mbox{\texttt{\mdseries\slshape C}} are contained in the $n$\texttt{\symbol{45}}th cyclotomic field. 

 
\begin{Verbatim}[commandchars=!@|,fontsize=\small,frame=single,label=Example]
  !gapprompt@gap>| !gapinput@Conductor( 0 ); Conductor( E(10) ); Conductor( E(12) );|
  1
  5
  12
\end{Verbatim}
 }

 

\subsection{\textcolor{Chapter }{AbsoluteValue}}
\logpage{[ 18, 1, 8 ]}\nobreak
\hyperdef{L}{X81DD58BB81FB3426}{}
{\noindent\textcolor{FuncColor}{$\triangleright$\enspace\texttt{AbsoluteValue({\mdseries\slshape cyc})\index{AbsoluteValue@\texttt{AbsoluteValue}}
\label{AbsoluteValue}
}\hfill{\scriptsize (attribute)}}\\


 returns the absolute value of a cyclotomic number \mbox{\texttt{\mdseries\slshape cyc}}. At the moment only methods for rational numbers exist. 
\begin{Verbatim}[commandchars=!@|,fontsize=\small,frame=single,label=Example]
  !gapprompt@gap>| !gapinput@AbsoluteValue(-3);|
  3
\end{Verbatim}
 }

 

\subsection{\textcolor{Chapter }{RoundCyc}}
\logpage{[ 18, 1, 9 ]}\nobreak
\hyperdef{L}{X7808ECF37AA9004D}{}
{\noindent\textcolor{FuncColor}{$\triangleright$\enspace\texttt{RoundCyc({\mdseries\slshape cyc})\index{RoundCyc@\texttt{RoundCyc}}
\label{RoundCyc}
}\hfill{\scriptsize (operation)}}\\


 is a cyclotomic integer $z$ (see \texttt{IsIntegralCyclotomic} (\ref{IsIntegralCyclotomic})) near to the cyclotomic \mbox{\texttt{\mdseries\slshape cyc}} in the following sense: Let \texttt{c} be the $i$\texttt{\symbol{45}}th coefficient in the external representation
(see{\nobreakspace}\texttt{CoeffsCyc} (\ref{CoeffsCyc})) of \mbox{\texttt{\mdseries\slshape cyc}}. Then the $i$\texttt{\symbol{45}}th coefficient in the external representation of $z$ is \texttt{Int( c + 1/2 )} or \texttt{Int( c \texttt{\symbol{45}} 1/2 )}, depending on whether \texttt{c} is nonnegative or negative, respectively. 

 Expressed in terms of the Zumbroich basis (see{\nobreakspace}\ref{Integral Bases of Abelian Number Fields}), rounding the coefficients of \mbox{\texttt{\mdseries\slshape cyc}} w.r.t.{\nobreakspace}this basis to the nearest integer yields the coefficients
of $z$. 

 
\begin{Verbatim}[commandchars=!@|,fontsize=\small,frame=single,label=Example]
  !gapprompt@gap>| !gapinput@RoundCyc( E(5)+1/2*E(5)^2 ); RoundCyc( 2/3*E(7)+3/2*E(4) );|
  E(5)+E(5)^2
  -2*E(28)^3+E(28)^4-2*E(28)^11-2*E(28)^15-2*E(28)^19-2*E(28)^23
   -2*E(28)^27
\end{Verbatim}
 }

 

\subsection{\textcolor{Chapter }{CoeffsCyc}}
\logpage{[ 18, 1, 10 ]}\nobreak
\hyperdef{L}{X7AE2933985BE4C3E}{}
{\noindent\textcolor{FuncColor}{$\triangleright$\enspace\texttt{CoeffsCyc({\mdseries\slshape cyc, N})\index{CoeffsCyc@\texttt{CoeffsCyc}}
\label{CoeffsCyc}
}\hfill{\scriptsize (function)}}\\


 \index{coefficients!for cyclotomics} Let \mbox{\texttt{\mdseries\slshape cyc}} be a cyclotomic with conductor $n$ (see \texttt{Conductor} (\ref{Conductor:for a cyclotomic})). If \mbox{\texttt{\mdseries\slshape N}} is not a multiple of $n$ then \texttt{CoeffsCyc} returns \texttt{fail} because \mbox{\texttt{\mdseries\slshape cyc}} cannot be expressed in terms of \mbox{\texttt{\mdseries\slshape N}}\texttt{\symbol{45}}th roots of unity. Otherwise \texttt{CoeffsCyc} returns a list of length \mbox{\texttt{\mdseries\slshape N}} with entry at position $j$ equal to the coefficient of $\exp(2 \pi i (j-1)/\mbox{\texttt{\mdseries\slshape N}})$ if this root belongs to the \mbox{\texttt{\mdseries\slshape N}}\texttt{\symbol{45}}th Zumbroich basis (see{\nobreakspace}\ref{Integral Bases of Abelian Number Fields}), and equal to zero otherwise. So we have \mbox{\texttt{\mdseries\slshape cyc}} = \texttt{CoeffsCyc(} \mbox{\texttt{\mdseries\slshape cyc}}, \mbox{\texttt{\mdseries\slshape N}} \texttt{) * List( [1..}\mbox{\texttt{\mdseries\slshape N}}\texttt{], j \texttt{\symbol{45}}{\textgreater} E(}\mbox{\texttt{\mdseries\slshape N}}\texttt{)\texttt{\symbol{94}}(j\texttt{\symbol{45}}1) )}. 

 
\begin{Verbatim}[commandchars=!@|,fontsize=\small,frame=single,label=Example]
  !gapprompt@gap>| !gapinput@cyc:= E(5)+E(5)^2;|
  E(5)+E(5)^2
  !gapprompt@gap>| !gapinput@CoeffsCyc( cyc, 5 );  CoeffsCyc( cyc, 15 );  CoeffsCyc( cyc, 7 );|
  [ 0, 1, 1, 0, 0 ]
  [ 0, -1, 0, 0, 0, 0, 0, 0, -1, 0, 0, -1, 0, -1, 0 ]
  fail
\end{Verbatim}
 }

 

\subsection{\textcolor{Chapter }{DenominatorCyc}}
\logpage{[ 18, 1, 11 ]}\nobreak
\hyperdef{L}{X803478CA7D2D830F}{}
{\noindent\textcolor{FuncColor}{$\triangleright$\enspace\texttt{DenominatorCyc({\mdseries\slshape cyc})\index{DenominatorCyc@\texttt{DenominatorCyc}}
\label{DenominatorCyc}
}\hfill{\scriptsize (function)}}\\


 For a cyclotomic number \mbox{\texttt{\mdseries\slshape cyc}} (see{\nobreakspace}\texttt{IsCyclotomic} (\ref{IsCyclotomic})), this function returns the smallest positive integer $n$ such that $n$\texttt{ * }\mbox{\texttt{\mdseries\slshape cyc}} is a cyclotomic integer (see{\nobreakspace}\texttt{IsIntegralCyclotomic} (\ref{IsIntegralCyclotomic})). For rational numbers \mbox{\texttt{\mdseries\slshape cyc}}, the result is the same as that of \texttt{DenominatorRat} (\ref{DenominatorRat}). }

 

\subsection{\textcolor{Chapter }{ExtRepOfObj (for a cyclotomic)}}
\logpage{[ 18, 1, 12 ]}\nobreak
\hyperdef{L}{X785F2CAB805DE1BE}{}
{\noindent\textcolor{FuncColor}{$\triangleright$\enspace\texttt{ExtRepOfObj({\mdseries\slshape cyc})\index{ExtRepOfObj@\texttt{ExtRepOfObj}!for a cyclotomic}
\label{ExtRepOfObj:for a cyclotomic}
}\hfill{\scriptsize (method)}}\\


 The external representation of a cyclotomic \mbox{\texttt{\mdseries\slshape cyc}} with conductor $n$ (see \texttt{Conductor} (\ref{Conductor:for a cyclotomic}) is the list returned by \texttt{CoeffsCyc} (\ref{CoeffsCyc}), called with \mbox{\texttt{\mdseries\slshape cyc}} and $n$. 

 
\begin{Verbatim}[commandchars=!@|,fontsize=\small,frame=single,label=Example]
  !gapprompt@gap>| !gapinput@ExtRepOfObj( E(5) ); CoeffsCyc( E(5), 5 );|
  [ 0, 1, 0, 0, 0 ]
  [ 0, 1, 0, 0, 0 ]
  !gapprompt@gap>| !gapinput@CoeffsCyc( E(5), 15 );|
  [ 0, 0, 0, 0, 0, 0, 0, 0, -1, 0, 0, 0, 0, -1, 0 ]
\end{Verbatim}
 }

 

\subsection{\textcolor{Chapter }{DescriptionOfRootOfUnity}}
\logpage{[ 18, 1, 13 ]}\nobreak
\hyperdef{L}{X7DDD51B983D5BC44}{}
{\noindent\textcolor{FuncColor}{$\triangleright$\enspace\texttt{DescriptionOfRootOfUnity({\mdseries\slshape root})\index{DescriptionOfRootOfUnity@\texttt{DescriptionOfRootOfUnity}}
\label{DescriptionOfRootOfUnity}
}\hfill{\scriptsize (function)}}\\


 \index{logarithm!of a root of unity} 

 Given a cyclotomic \mbox{\texttt{\mdseries\slshape root}} that is known to be a root of unity (this is \emph{not} checked), \texttt{DescriptionOfRootOfUnity} returns a list $[ n, e ]$ of coprime positive integers such that \mbox{\texttt{\mdseries\slshape root}} $=$ \texttt{E}$(n)^e$ holds. 

 
\begin{Verbatim}[commandchars=!@|,fontsize=\small,frame=single,label=Example]
  !gapprompt@gap>| !gapinput@E(9);  DescriptionOfRootOfUnity( E(9) );|
  -E(9)^4-E(9)^7
  [ 9, 1 ]
  !gapprompt@gap>| !gapinput@DescriptionOfRootOfUnity( -E(3) );|
  [ 6, 5 ]
\end{Verbatim}
 }

 

\subsection{\textcolor{Chapter }{IsGaussInt}}
\logpage{[ 18, 1, 14 ]}\nobreak
\hyperdef{L}{X8712419182ECD8DD}{}
{\noindent\textcolor{FuncColor}{$\triangleright$\enspace\texttt{IsGaussInt({\mdseries\slshape x})\index{IsGaussInt@\texttt{IsGaussInt}}
\label{IsGaussInt}
}\hfill{\scriptsize (function)}}\\


 \texttt{IsGaussInt} returns \texttt{true} if the object \mbox{\texttt{\mdseries\slshape x}} is a Gaussian integer (see{\nobreakspace}\texttt{GaussianIntegers} (\ref{GaussianIntegers})), and \texttt{false} otherwise. Gaussian integers are of the form $a + b$\texttt{*E(4)}, where $a$ and $b$ are integers. }

 

\subsection{\textcolor{Chapter }{IsGaussRat}}
\logpage{[ 18, 1, 15 ]}\nobreak
\hyperdef{L}{X7E6CF4947D0A56F7}{}
{\noindent\textcolor{FuncColor}{$\triangleright$\enspace\texttt{IsGaussRat({\mdseries\slshape x})\index{IsGaussRat@\texttt{IsGaussRat}}
\label{IsGaussRat}
}\hfill{\scriptsize (function)}}\\


 \texttt{IsGaussRat} returns \texttt{true} if the object \mbox{\texttt{\mdseries\slshape x}} is a Gaussian rational (see{\nobreakspace}\texttt{GaussianRationals} (\ref{GaussianRationals})), and \texttt{false} otherwise. Gaussian rationals are of the form $a + b$\texttt{*E(4)}, where $a$ and $b$ are rationals. }

 

\subsection{\textcolor{Chapter }{DefaultField (for cyclotomics)}}
\logpage{[ 18, 1, 16 ]}\nobreak
\hyperdef{L}{X7FE3D5637B5485D0}{}
{\noindent\textcolor{FuncColor}{$\triangleright$\enspace\texttt{DefaultField({\mdseries\slshape list})\index{DefaultField@\texttt{DefaultField}!for cyclotomics}
\label{DefaultField:for cyclotomics}
}\hfill{\scriptsize (function)}}\\


 \texttt{DefaultField} for cyclotomics is defined to return the smallest \emph{cyclotomic} field containing the given elements. 

 Note that \texttt{Field} (\ref{Field:for several generators}) returns the smallest field containing all given elements, which need not be a
cyclotomic field. In both cases, the fields represent vector spaces over the
rationals (see{\nobreakspace}\ref{Integral Bases of Abelian Number Fields}). 

 
\begin{Verbatim}[commandchars=!@|,fontsize=\small,frame=single,label=Example]
  !gapprompt@gap>| !gapinput@Field( E(5)+E(5)^4 );  DefaultField( E(5)+E(5)^4 );|
  NF(5,[ 1, 4 ])
  CF(5)
\end{Verbatim}
 }

  }

  
\section{\textcolor{Chapter }{Infinity and negative Infinity}}\label{Infinity}
\logpage{[ 18, 2, 0 ]}
\hyperdef{L}{X7EE5FB7181125E02}{}
{
  

\subsection{\textcolor{Chapter }{IsInfinity}}
\logpage{[ 18, 2, 1 ]}\nobreak
\hyperdef{L}{X8511B8DF83324C27}{}
{\noindent\textcolor{FuncColor}{$\triangleright$\enspace\texttt{IsInfinity({\mdseries\slshape obj})\index{IsInfinity@\texttt{IsInfinity}}
\label{IsInfinity}
}\hfill{\scriptsize (Category)}}\\
\noindent\textcolor{FuncColor}{$\triangleright$\enspace\texttt{IsNegInfinity({\mdseries\slshape obj})\index{IsNegInfinity@\texttt{IsNegInfinity}}
\label{IsNegInfinity}
}\hfill{\scriptsize (Category)}}\\
\noindent\textcolor{FuncColor}{$\triangleright$\enspace\texttt{infinity\index{infinity@\texttt{infinity}}
\label{infinity}
}\hfill{\scriptsize (global variable)}}\\
\noindent\textcolor{FuncColor}{$\triangleright$\enspace\texttt{\texttt{\symbol{45}}infinity\index{\texttt{\symbol{45}}infinity@\texttt{\texttt{\symbol{45}}infinity}}
\label{-infinity}
}\hfill{\scriptsize (global variable)}}\\


 \texttt{infinity} and \texttt{\texttt{\symbol{45}}infinity} are special \textsf{GAP} objects that lie in \texttt{CyclotomicsFamily}. They are larger or smaller than all other objects in this family
respectively. \texttt{infinity} is mainly used as return value of operations such as \texttt{Size} (\ref{Size}) and \texttt{Dimension} (\ref{Dimension}) for infinite and infinite dimensional domains, respectively. 

 Some arithmetic operations are provided for convenience when using \texttt{infinity} and \texttt{\texttt{\symbol{45}}infinity} as top and bottom element respectively. 
\begin{Verbatim}[commandchars=!@|,fontsize=\small,frame=single,label=Example]
  !gapprompt@gap>| !gapinput@-infinity + 1;|
  -infinity
  !gapprompt@gap>| !gapinput@infinity + infinity;|
  infinity
\end{Verbatim}
 Often it is useful to distinguish \texttt{infinity} from ``proper'' cyclotomics. For that, \texttt{infinity} lies in the category \texttt{IsInfinity} but not in \texttt{IsCyc} (\ref{IsCyc}), and the other cyclotomics lie in the category \texttt{IsCyc} (\ref{IsCyc}) but not in \texttt{IsInfinity}. 

 
\begin{Verbatim}[commandchars=!@|,fontsize=\small,frame=single,label=Example]
  !gapprompt@gap>| !gapinput@s:= Size( Rationals );|
  infinity
  !gapprompt@gap>| !gapinput@s = infinity; IsCyclotomic( s ); IsCyc( s ); IsInfinity( s );|
  true
  true
  false
  true
  !gapprompt@gap>| !gapinput@s in Rationals; s > 17;|
  false
  true
  !gapprompt@gap>| !gapinput@Set( [ s, 2, s, E(17), s, 19 ] );|
  [ 2, 19, E(17), infinity ]
\end{Verbatim}
 }

 }

  
\section{\textcolor{Chapter }{Comparisons of Cyclotomics}}\label{Comparisons of Cyclotomics}
\logpage{[ 18, 3, 0 ]}
\hyperdef{L}{X7F66A62384329705}{}
{
  \index{operators!for cyclotomics} To compare cyclotomics, the operators \texttt{{\textless}}, \texttt{{\textless}=}, \texttt{=}, \texttt{{\textgreater}=}, \texttt{{\textgreater}}, and \texttt{{\textless}{\textgreater}} can be used, the result will be \texttt{true} if the first operand is smaller, smaller or equal, equal, larger or equal,
larger, or unequal, respectively, and \texttt{false} otherwise. 

 Cyclotomics are ordered as follows: The relation between rationals is the
natural one, rationals are smaller than irrational cyclotomics, and \texttt{infinity} (\ref{infinity}) is the largest cyclotomic. For two irrational cyclotomics with different
conductors (see \texttt{Conductor} (\ref{Conductor:for a cyclotomic})), the one with smaller conductor is regarded as smaller. Two irrational
cyclotomics with same conductor are compared via their external representation
(see \texttt{ExtRepOfObj} (\ref{ExtRepOfObj:for a cyclotomic})). 

 For comparisons of cyclotomics and other \textsf{GAP} objects, see Section{\nobreakspace}\ref{Comparisons}. 

 
\begin{Verbatim}[commandchars=!@|,fontsize=\small,frame=single,label=Example]
  !gapprompt@gap>| !gapinput@E(5) < E(6);      # the latter value has conductor 3|
  false
  !gapprompt@gap>| !gapinput@E(3) < E(3)^2;    # both have conductor 3, compare the ext. repr.|
  false
  !gapprompt@gap>| !gapinput@3 < E(3); E(5) < E(7);|
  true
  true
\end{Verbatim}
 }

  
\section{\textcolor{Chapter }{ATLAS Irrationalities}}\label{ATLAS Irrationalities}
\logpage{[ 18, 4, 0 ]}
\hyperdef{L}{X7B242083873DD74F}{}
{
  \index{atomic irrationalities} 
\subsection{\textcolor{Chapter }{EB, EC, $\ldots$, EH}}\logpage{[ 18, 4, 1 ]}
\hyperdef{L}{X8414ED887AF36359}{}
{
\noindent\textcolor{FuncColor}{$\triangleright$\enspace\texttt{EB({\mdseries\slshape N})\index{EB@\texttt{EB}}
\label{EB}
}\hfill{\scriptsize (function)}}\\
\noindent\textcolor{FuncColor}{$\triangleright$\enspace\texttt{EC({\mdseries\slshape N})\index{EC@\texttt{EC}}
\label{EC}
}\hfill{\scriptsize (function)}}\\
\noindent\textcolor{FuncColor}{$\triangleright$\enspace\texttt{ED({\mdseries\slshape N})\index{ED@\texttt{ED}}
\label{ED}
}\hfill{\scriptsize (function)}}\\
\noindent\textcolor{FuncColor}{$\triangleright$\enspace\texttt{EE({\mdseries\slshape N})\index{EE@\texttt{EE}}
\label{EE}
}\hfill{\scriptsize (function)}}\\
\noindent\textcolor{FuncColor}{$\triangleright$\enspace\texttt{EF({\mdseries\slshape N})\index{EF@\texttt{EF}}
\label{EF}
}\hfill{\scriptsize (function)}}\\
\noindent\textcolor{FuncColor}{$\triangleright$\enspace\texttt{EG({\mdseries\slshape N})\index{EG@\texttt{EG}}
\label{EG}
}\hfill{\scriptsize (function)}}\\
\noindent\textcolor{FuncColor}{$\triangleright$\enspace\texttt{EH({\mdseries\slshape N})\index{EH@\texttt{EH}}
\label{EH}
}\hfill{\scriptsize (function)}}\\


 \index{b_N@$b_N$ (irrational value)} \index{c_N@$c_N$ (irrational value)} \index{d_N@$d_N$ (irrational value)} \index{e_N@$e_N$ (irrational value)} \index{f_N@$f_N$ (irrational value)} \index{g_N@$g_N$ (irrational value)} \index{h_N@$h_N$ (irrational value)} For a positive integer \mbox{\texttt{\mdseries\slshape N}}, let $z =$ \texttt{E(}\mbox{\texttt{\mdseries\slshape N}}\texttt{)} $= \exp(2 \pi i/\mbox{\texttt{\mdseries\slshape N}})$. The following so\texttt{\symbol{45}}called \emph{atomic irrationalities} (see \cite[Chapter 7, Section 10]{CCN85}) can be entered using functions. (Note that the values are not necessary
irrational.) 

 \begin{center}
\begin{tabular}{lclclcl}\texttt{EB(}\mbox{\texttt{\mdseries\slshape N}}\texttt{)}&
=&
$b_{\mbox{\texttt{\mdseries\slshape N}}}$&
=&
$\left( \sum_{{j = 1}}^{{\mbox{\texttt{\mdseries\slshape N}}-1}} z^{{j^2}} \right) / 2$ &
,&
$\mbox{\texttt{\mdseries\slshape N}} \equiv 1 \pmod{2}$\\
\texttt{EC(}\mbox{\texttt{\mdseries\slshape N}}\texttt{)}&
=&
$c_{\mbox{\texttt{\mdseries\slshape N}}}$&
=&
$\left( \sum_{{j = 1}}^{{\mbox{\texttt{\mdseries\slshape N}}-1}} z^{{j^3}} \right) / 3$ &
,&
$\mbox{\texttt{\mdseries\slshape N}} \equiv 1 \pmod{3}$\\
\texttt{ED(}\mbox{\texttt{\mdseries\slshape N}}\texttt{)}&
=&
$d_{\mbox{\texttt{\mdseries\slshape N}}}$&
=&
$\left( \sum_{{j = 1}}^{{\mbox{\texttt{\mdseries\slshape N}}-1}} z^{{j^4}} \right) / 4$ &
,&
$\mbox{\texttt{\mdseries\slshape N}} \equiv 1 \pmod{4}$\\
\texttt{EE(}\mbox{\texttt{\mdseries\slshape N}}\texttt{)}&
=&
$e_{\mbox{\texttt{\mdseries\slshape N}}}$&
=&
$\left( \sum_{{j = 1}}^{{\mbox{\texttt{\mdseries\slshape N}}-1}} z^{{j^5}} \right) / 5$ &
,&
$\mbox{\texttt{\mdseries\slshape N}} \equiv 1 \pmod{5}$\\
\texttt{EF(}\mbox{\texttt{\mdseries\slshape N}}\texttt{)}&
=&
$f_{\mbox{\texttt{\mdseries\slshape N}}}$&
=&
$\left( \sum_{{j = 1}}^{{\mbox{\texttt{\mdseries\slshape N}}-1}} z^{{j^6}} \right) / 6$ &
,&
$\mbox{\texttt{\mdseries\slshape N}} \equiv 1 \pmod{6}$\\
\texttt{EG(}\mbox{\texttt{\mdseries\slshape N}}\texttt{)}&
=&
$g_{\mbox{\texttt{\mdseries\slshape N}}}$&
=&
$\left( \sum_{{j = 1}}^{{\mbox{\texttt{\mdseries\slshape N}}-1}} z^{{j^7}} \right) / 7$ &
,&
$\mbox{\texttt{\mdseries\slshape N}} \equiv 1 \pmod{7}$\\
\texttt{EH(}\mbox{\texttt{\mdseries\slshape N}}\texttt{)}&
=&
$h_{\mbox{\texttt{\mdseries\slshape N}}}$&
=&
$\left( \sum_{{j = 1}}^{{\mbox{\texttt{\mdseries\slshape N}}-1}} z^{{j^8}} \right) / 8$ &
,&
$\mbox{\texttt{\mdseries\slshape N}} \equiv 1 \pmod{8}$\\
\end{tabular}\\[2mm]
\end{center}

 (Note that in \texttt{EC(}\mbox{\texttt{\mdseries\slshape N}}\texttt{)}, $\ldots$, \texttt{EH(}\mbox{\texttt{\mdseries\slshape N}}\texttt{)}, \mbox{\texttt{\mdseries\slshape N}} must be a prime.) 

 
\begin{Verbatim}[commandchars=!@|,fontsize=\small,frame=single,label=Example]
  !gapprompt@gap>| !gapinput@EB(5);  EB(9);|
  E(5)+E(5)^4
  1
\end{Verbatim}
 }

 
\subsection{\textcolor{Chapter }{EI and ER}}\logpage{[ 18, 4, 2 ]}
\hyperdef{L}{X813CF4327C4B4D29}{}
{
\noindent\textcolor{FuncColor}{$\triangleright$\enspace\texttt{EI({\mdseries\slshape N})\index{EI@\texttt{EI}}
\label{EI}
}\hfill{\scriptsize (function)}}\\
\noindent\textcolor{FuncColor}{$\triangleright$\enspace\texttt{ER({\mdseries\slshape N})\index{ER@\texttt{ER}}
\label{ER}
}\hfill{\scriptsize (function)}}\\


 \index{i_N@$i_N$ (irrational value)} \index{r_N@$r_N$ (irrational value)} For a rational number \mbox{\texttt{\mdseries\slshape N}}, \texttt{ER} returns the square root $\sqrt{{\mbox{\texttt{\mdseries\slshape N}}}}$ of \mbox{\texttt{\mdseries\slshape N}}, and \texttt{EI} returns $\sqrt{{-\mbox{\texttt{\mdseries\slshape N}}}}$. By the chosen embedding of cyclotomic fields into the complex numbers, \texttt{ER} returns the positive square root if \mbox{\texttt{\mdseries\slshape N}} is positive, and if \mbox{\texttt{\mdseries\slshape N}} is negative then \texttt{ER(}\mbox{\texttt{\mdseries\slshape N}}\texttt{) = EI(\texttt{\symbol{45}}}\mbox{\texttt{\mdseries\slshape N}}\texttt{)} holds. In any case, \texttt{EI(}\mbox{\texttt{\mdseries\slshape N}}\texttt{) = E(4) * ER(}\mbox{\texttt{\mdseries\slshape N}}\texttt{)}. 

 \texttt{ER} is installed as method for the operation \texttt{Sqrt} (\ref{Sqrt}), for rational argument. 

 From a theorem of Gauss (see \cite[QS4 in Ch. IV, Par. 3]{Lan70}) we know that $b_{\mbox{\texttt{\mdseries\slshape N}}} =$ \begin{center}
\begin{tabular}{lcl}$(-1 + \sqrt{{\mbox{\texttt{\mdseries\slshape N}}}}) / 2$&
if&
$\mbox{\texttt{\mdseries\slshape N}} \equiv 1 \pmod 4$\\
$(-1 + i \sqrt{{\mbox{\texttt{\mdseries\slshape N}}}}) / 2$&
if&
$\mbox{\texttt{\mdseries\slshape N}} \equiv -1 \pmod 4$\\
\end{tabular}\\[2mm]
\end{center}

 So $\sqrt{{\mbox{\texttt{\mdseries\slshape N}}}}$ can be computed from $b_{\mbox{\texttt{\mdseries\slshape N}}}$, see{\nobreakspace}\texttt{EB} (\ref{EB}). 

 
\begin{Verbatim}[commandchars=!@|,fontsize=\small,frame=single,label=Example]
  !gapprompt@gap>| !gapinput@ER(3); EI(3);|
  -E(12)^7+E(12)^11
  E(3)-E(3)^2
\end{Verbatim}
 }

 
\subsection{\textcolor{Chapter }{EY, EX, $\ldots$, ES}}\logpage{[ 18, 4, 3 ]}
\hyperdef{L}{X8672D7F986CBA116}{}
{
\noindent\textcolor{FuncColor}{$\triangleright$\enspace\texttt{EY({\mdseries\slshape N[, d]})\index{EY@\texttt{EY}}
\label{EY}
}\hfill{\scriptsize (function)}}\\
\noindent\textcolor{FuncColor}{$\triangleright$\enspace\texttt{EX({\mdseries\slshape N[, d]})\index{EX@\texttt{EX}}
\label{EX}
}\hfill{\scriptsize (function)}}\\
\noindent\textcolor{FuncColor}{$\triangleright$\enspace\texttt{EW({\mdseries\slshape N[, d]})\index{EW@\texttt{EW}}
\label{EW}
}\hfill{\scriptsize (function)}}\\
\noindent\textcolor{FuncColor}{$\triangleright$\enspace\texttt{EV({\mdseries\slshape N[, d]})\index{EV@\texttt{EV}}
\label{EV}
}\hfill{\scriptsize (function)}}\\
\noindent\textcolor{FuncColor}{$\triangleright$\enspace\texttt{EU({\mdseries\slshape N[, d]})\index{EU@\texttt{EU}}
\label{EU}
}\hfill{\scriptsize (function)}}\\
\noindent\textcolor{FuncColor}{$\triangleright$\enspace\texttt{ET({\mdseries\slshape N[, d]})\index{ET@\texttt{ET}}
\label{ET}
}\hfill{\scriptsize (function)}}\\
\noindent\textcolor{FuncColor}{$\triangleright$\enspace\texttt{ES({\mdseries\slshape N[, d]})\index{ES@\texttt{ES}}
\label{ES}
}\hfill{\scriptsize (function)}}\\


 \index{s_N@$s_N$ (irrational value)} \index{t_N@$t_N$ (irrational value)} \index{u_N@$u_N$ (irrational value)} \index{v_N@$v_N$ (irrational value)} \index{w_N@$w_N$ (irrational value)} \index{x_N@$x_N$ (irrational value)} \index{y_N@$y_N$ (irrational value)} For the given integer \mbox{\texttt{\mdseries\slshape N}} $> 2$, let $\mbox{\texttt{\mdseries\slshape N}}_k$ denote the first integer with multiplicative order exactly $k$ modulo \mbox{\texttt{\mdseries\slshape N}}, chosen in the order of preference 
\[ 1, -1, 2, -2, 3, -3, 4, -4, \ldots . \]
 

 We define (with $z = \exp(2 \pi i/\mbox{\texttt{\mdseries\slshape N}})$) \begin{center}
\begin{tabular}{lclcll}\texttt{EY(}\mbox{\texttt{\mdseries\slshape N}}\texttt{)}&
=&
$y_{\mbox{\texttt{\mdseries\slshape N}}}$&
=&
$z + z^n$&
$(n = \mbox{\texttt{\mdseries\slshape N}}_2)$\\
\texttt{EX(}\mbox{\texttt{\mdseries\slshape N}}\texttt{)}&
=&
$x_{\mbox{\texttt{\mdseries\slshape N}}}$&
=&
$z + z^n + z^{{n^2}}$&
$(n = \mbox{\texttt{\mdseries\slshape N}}_3)$\\
\texttt{EW}(\mbox{\texttt{\mdseries\slshape N}}\texttt{)}&
=&
$w_{\mbox{\texttt{\mdseries\slshape N}}}$&
=&
$z + z^n + z^{{n^2}} + z^{{n^3}}$&
$(n = \mbox{\texttt{\mdseries\slshape N}}_4)$\\
\texttt{EV(}\mbox{\texttt{\mdseries\slshape N}}\texttt{)}&
=&
$v_{\mbox{\texttt{\mdseries\slshape N}}}$&
=&
$z + z^n + z^{{n^2}} + z^{{n^3}} + z^{{n^4}}$&
$(n = \mbox{\texttt{\mdseries\slshape N}}_5)$\\
\texttt{EU(}\mbox{\texttt{\mdseries\slshape N}}\texttt{)}&
=&
$u_{\mbox{\texttt{\mdseries\slshape N}}}$&
=&
$z + z^n + z^{{n^2}} + \ldots + z^{{n^5}}$&
$(n = \mbox{\texttt{\mdseries\slshape N}}_6)$\\
\texttt{ET(}\mbox{\texttt{\mdseries\slshape N}}\texttt{)}&
=&
$t_{\mbox{\texttt{\mdseries\slshape N}}}$&
=&
$z + z^n + z^{{n^2}} + \ldots + z^{{n^6}}$&
$(n = \mbox{\texttt{\mdseries\slshape N}}_7)$\\
\texttt{ES(}\mbox{\texttt{\mdseries\slshape N}}\texttt{)}&
=&
$s_{\mbox{\texttt{\mdseries\slshape N}}}$&
=&
$z + z^n + z^{{n^2}} + \ldots + z^{{n^7}}$&
$(n = \mbox{\texttt{\mdseries\slshape N}}_8)$\\
\end{tabular}\\[2mm]
\end{center}

 

 For the two\texttt{\symbol{45}}argument versions of the functions, see Section \texttt{NK} (\ref{NK}). 

 
\begin{Verbatim}[commandchars=!@|,fontsize=\small,frame=single,label=Example]
  !gapprompt@gap>| !gapinput@EY(5);|
  E(5)+E(5)^4
  !gapprompt@gap>| !gapinput@EW(16,3); EW(17,2);|
  0
  E(17)+E(17)^4+E(17)^13+E(17)^16
\end{Verbatim}
 }

 
\subsection{\textcolor{Chapter }{EM, EL, $\ldots$, EJ}}\logpage{[ 18, 4, 4 ]}
\hyperdef{L}{X7E5985FC846C5201}{}
{
\noindent\textcolor{FuncColor}{$\triangleright$\enspace\texttt{EM({\mdseries\slshape N[, d]})\index{EM@\texttt{EM}}
\label{EM}
}\hfill{\scriptsize (function)}}\\
\noindent\textcolor{FuncColor}{$\triangleright$\enspace\texttt{EL({\mdseries\slshape N[, d]})\index{EL@\texttt{EL}}
\label{EL}
}\hfill{\scriptsize (function)}}\\
\noindent\textcolor{FuncColor}{$\triangleright$\enspace\texttt{EK({\mdseries\slshape N[, d]})\index{EK@\texttt{EK}}
\label{EK}
}\hfill{\scriptsize (function)}}\\
\noindent\textcolor{FuncColor}{$\triangleright$\enspace\texttt{EJ({\mdseries\slshape N[, d]})\index{EJ@\texttt{EJ}}
\label{EJ}
}\hfill{\scriptsize (function)}}\\


 Let \mbox{\texttt{\mdseries\slshape N}} be an integer, \mbox{\texttt{\mdseries\slshape N}} $> 2$. We define (with $z = \exp(2 \pi i/\mbox{\texttt{\mdseries\slshape N}})$) \index{j_N@$j_N$ (irrational value)} \index{k_N@$k_N$ (irrational value)} \index{l_N@$l_N$ (irrational value)} \index{m_N@$m_N$ (irrational value)} \begin{center}
\begin{tabular}{lclcll}\texttt{EM(}\mbox{\texttt{\mdseries\slshape N}}\texttt{)}&
=&
$m_{\mbox{\texttt{\mdseries\slshape N}}}$&
=&
$z - z^n$&
$(n = \mbox{\texttt{\mdseries\slshape N}}_2)$\\
\texttt{EL(}\mbox{\texttt{\mdseries\slshape N}}\texttt{)}&
=&
$l_{\mbox{\texttt{\mdseries\slshape N}}}$&
=&
$z - z^n + z^{{n^2}} - z^{{n^3}}$&
$(n = \mbox{\texttt{\mdseries\slshape N}}_4)$\\
\texttt{EK(}\mbox{\texttt{\mdseries\slshape N}}\texttt{)}&
=&
$k_{\mbox{\texttt{\mdseries\slshape N}}}$&
=&
$z - z^n + \ldots - z^{{n^5}}$&
$(n = \mbox{\texttt{\mdseries\slshape N}}_6)$\\
\texttt{EJ(}\mbox{\texttt{\mdseries\slshape N}}\texttt{)}&
=&
$j_{\mbox{\texttt{\mdseries\slshape N}}}$&
=&
$z - z^n + \ldots - z^{{n^7}}$&
$(n = \mbox{\texttt{\mdseries\slshape N}}_8)$\\
\end{tabular}\\[2mm]
\end{center}

 

 For the two\texttt{\symbol{45}}argument versions of the functions, see Section \texttt{NK} (\ref{NK}). }

 

\subsection{\textcolor{Chapter }{NK}}
\logpage{[ 18, 4, 5 ]}\nobreak
\hyperdef{L}{X844F0EBF849EDEB3}{}
{\noindent\textcolor{FuncColor}{$\triangleright$\enspace\texttt{NK({\mdseries\slshape N, k, d})\index{NK@\texttt{NK}}
\label{NK}
}\hfill{\scriptsize (function)}}\\


 Let $\mbox{\texttt{\mdseries\slshape N}}_{\mbox{\texttt{\mdseries\slshape k}}}^{(\mbox{\texttt{\mdseries\slshape d}})}$ be the $(\mbox{\texttt{\mdseries\slshape d}}+1)$\texttt{\symbol{45}}th integer with multiplicative order exactly \mbox{\texttt{\mdseries\slshape k}} modulo \mbox{\texttt{\mdseries\slshape N}}, chosen in the order of preference defined in Section \ref{EY}; \texttt{NK} returns $\mbox{\texttt{\mdseries\slshape N}}_{\mbox{\texttt{\mdseries\slshape k}}}^{(\mbox{\texttt{\mdseries\slshape d}})}$; if there is no integer with the required multiplicative order, \texttt{NK} returns \texttt{fail}. 

 We write $\mbox{\texttt{\mdseries\slshape N}}_{\mbox{\texttt{\mdseries\slshape k}}} = \mbox{\texttt{\mdseries\slshape N}}_{\mbox{\texttt{\mdseries\slshape k}}}^{(0)}, \mbox{\texttt{\mdseries\slshape N}}_{\mbox{\texttt{\mdseries\slshape k}}}^{\prime} = \mbox{\texttt{\mdseries\slshape N}}_{\mbox{\texttt{\mdseries\slshape k}}}^{(1)}, \mbox{\texttt{\mdseries\slshape N}}_{\mbox{\texttt{\mdseries\slshape k}}}^{\prime\prime} = \mbox{\texttt{\mdseries\slshape N}}_{\mbox{\texttt{\mdseries\slshape k}}}^{(2)}$ and so on. 

 The algebraic numbers 
\[ y_{\mbox{\texttt{\mdseries\slshape N}}}^{\prime} = y_{\mbox{\texttt{\mdseries\slshape N}}}^{(1)}, y_{\mbox{\texttt{\mdseries\slshape N}}}^{\prime\prime} = y_{\mbox{\texttt{\mdseries\slshape N}}}^{(2)}, \ldots, x_{\mbox{\texttt{\mdseries\slshape N}}}^{\prime}, x_{\mbox{\texttt{\mdseries\slshape N}}}^{\prime\prime}, \ldots, j_{\mbox{\texttt{\mdseries\slshape N}}}^{\prime}, j_{\mbox{\texttt{\mdseries\slshape N}}}^{\prime\prime}, \ldots \]
 are obtained on replacing $\mbox{\texttt{\mdseries\slshape N}}_{\mbox{\texttt{\mdseries\slshape k}}}$ in the definitions in the sections \ref{EY} and \ref{EM} by $\mbox{\texttt{\mdseries\slshape N}}_{\mbox{\texttt{\mdseries\slshape k}}}^{\prime}, \mbox{\texttt{\mdseries\slshape N}}_{\mbox{\texttt{\mdseries\slshape k}}}^{\prime\prime}, \ldots$; they can be entered as 

 \begin{center}
\begin{tabular}{lcl}\texttt{EY(}\mbox{\texttt{\mdseries\slshape N}},\mbox{\texttt{\mdseries\slshape d}}\texttt{)}&
=&
$y_{\mbox{\texttt{\mdseries\slshape N}}}^{(\mbox{\texttt{\mdseries\slshape d}})}$\\
\texttt{EX(}\mbox{\texttt{\mdseries\slshape N}},\mbox{\texttt{\mdseries\slshape d}}\texttt{)}&
=&
$x_{\mbox{\texttt{\mdseries\slshape N}}}^{(\mbox{\texttt{\mdseries\slshape d}})}$\\
&
$\ldots$&
\\
\texttt{EJ(}\mbox{\texttt{\mdseries\slshape N}},\mbox{\texttt{\mdseries\slshape d}}\texttt{)}&
=&
$j_{\mbox{\texttt{\mdseries\slshape N}}}^{(\mbox{\texttt{\mdseries\slshape d}})}$\\
\end{tabular}\\[2mm]
\end{center}

 }

 

\subsection{\textcolor{Chapter }{AtlasIrrationality}}
\logpage{[ 18, 4, 6 ]}\nobreak
\hyperdef{L}{X812E334E7A869D33}{}
{\noindent\textcolor{FuncColor}{$\triangleright$\enspace\texttt{AtlasIrrationality({\mdseries\slshape irratname})\index{AtlasIrrationality@\texttt{AtlasIrrationality}}
\label{AtlasIrrationality}
}\hfill{\scriptsize (function)}}\\


 Let \mbox{\texttt{\mdseries\slshape irratname}} be a string that describes an irrational value as a linear combination in
terms of the atomic irrationalities introduced in the sections \ref{EB}, \ref{EI}, \ref{EY}, \ref{EM}. These irrational values are defined in \cite[Chapter 6, Section 10]{CCN85}, and the following description is mainly copied from there. If $q_N$ is such a value (e.g. $y_{24}^{\prime\prime}$) then linear combinations of algebraic conjugates of $q_N$ are abbreviated as in the following examples: 

 \begin{center}
\begin{tabular}{lcl}\texttt{2qN+3\&5\texttt{\symbol{45}}4\&7+\&9}&
means&
$2 q_N + 3 q_N^{{*5}} - 4 q_N^{{*7}} + q_N^{{*9}}$ \\
\texttt{4qN\&3\&5\&7\texttt{\symbol{45}}3\&4}&
means&
$4 (q_N + q_N^{{*3}} + q_N^{{*5}} + q_N^{{*7}}) - 3 q_N^{{*11}}$\\
\texttt{4qN*3\&5+\&7}&
means&
$4 (q_N^{{*3}} + q_N^{{*5}}) + q_N^{{*7}}$\\
\end{tabular}\\[2mm]
\end{center}

 

 To explain the ``ampersand'' syntax in general we remark that ``\&k'' is interpreted as $q_N^{{*k}}$, where $q_N$ is the most recently named atomic irrationality, and that the scope of any
premultiplying coefficient is broken by a $+$ or $-$ sign, but not by $\&$ or $*k$. The algebraic conjugations indicated by the ampersands apply directly to the \emph{atomic} irrationality $q_N$, even when, as in the last example, $q_N$ first appears with another conjugacy $*k$. 

 
\begin{Verbatim}[commandchars=!@|,fontsize=\small,frame=single,label=Example]
  !gapprompt@gap>| !gapinput@AtlasIrrationality( "b7*3" );|
  E(7)^3+E(7)^5+E(7)^6
  !gapprompt@gap>| !gapinput@AtlasIrrationality( "y'''24" );|
  E(24)-E(24)^19
  !gapprompt@gap>| !gapinput@AtlasIrrationality( "-3y'''24*13&5" );|
  3*E(8)-3*E(8)^3
  !gapprompt@gap>| !gapinput@AtlasIrrationality( "3y'''24*13-2&5" );|
  -3*E(24)-2*E(24)^11+2*E(24)^17+3*E(24)^19
  !gapprompt@gap>| !gapinput@AtlasIrrationality( "3y'''24*13-&5" );|
  -3*E(24)-E(24)^11+E(24)^17+3*E(24)^19
  !gapprompt@gap>| !gapinput@AtlasIrrationality( "3y'''24*13-4&5&7" );|
  -7*E(24)-4*E(24)^11+4*E(24)^17+7*E(24)^19
  !gapprompt@gap>| !gapinput@AtlasIrrationality( "3y'''24&7" );|
  6*E(24)-6*E(24)^19
\end{Verbatim}
 }

 }

  
\section{\textcolor{Chapter }{Galois Conjugacy of Cyclotomics}}\label{Galois Conjugacy of Cyclotomics}
\logpage{[ 18, 5, 0 ]}
\hyperdef{L}{X79FE34337DF2CD10}{}
{
  

\subsection{\textcolor{Chapter }{GaloisCyc (for a cyclotomic)}}
\logpage{[ 18, 5, 1 ]}\nobreak
\hyperdef{L}{X79EE9097783128C4}{}
{\noindent\textcolor{FuncColor}{$\triangleright$\enspace\texttt{GaloisCyc({\mdseries\slshape cyc, k})\index{GaloisCyc@\texttt{GaloisCyc}!for a cyclotomic}
\label{GaloisCyc:for a cyclotomic}
}\hfill{\scriptsize (operation)}}\\
\noindent\textcolor{FuncColor}{$\triangleright$\enspace\texttt{GaloisCyc({\mdseries\slshape list, k})\index{GaloisCyc@\texttt{GaloisCyc}!for a list of cyclotomics}
\label{GaloisCyc:for a list of cyclotomics}
}\hfill{\scriptsize (operation)}}\\


 For a cyclotomic \mbox{\texttt{\mdseries\slshape cyc}} and an integer \mbox{\texttt{\mdseries\slshape k}}, \texttt{GaloisCyc} returns the cyclotomic obtained by raising the roots of unity in the Zumbroich
basis representation of \mbox{\texttt{\mdseries\slshape cyc}} to the \mbox{\texttt{\mdseries\slshape k}}\texttt{\symbol{45}}th power. If \mbox{\texttt{\mdseries\slshape k}} is coprime to the integer $n$, \texttt{GaloisCyc( ., \mbox{\texttt{\mdseries\slshape k}} )} acts as a Galois automorphism of the $n$\texttt{\symbol{45}}th cyclotomic field (see{\nobreakspace}\ref{Galois Groups of Abelian Number Fields}); to get the Galois automorphisms themselves, use \texttt{GaloisGroup} (\ref{GaloisGroup:of field}). 

 The \emph{complex conjugate} of \mbox{\texttt{\mdseries\slshape cyc}} is \texttt{GaloisCyc( \mbox{\texttt{\mdseries\slshape cyc}}, \texttt{\symbol{45}}1 )}, which can also be computed using \texttt{ComplexConjugate} (\ref{ComplexConjugate}). 

 For a list or matrix \mbox{\texttt{\mdseries\slshape list}} of cyclotomics, \texttt{GaloisCyc} returns the list obtained by applying \texttt{GaloisCyc} to the entries of \mbox{\texttt{\mdseries\slshape list}}. }

 

\subsection{\textcolor{Chapter }{ComplexConjugate}}
\logpage{[ 18, 5, 2 ]}\nobreak
\hyperdef{L}{X7BE001A0811CD599}{}
{\noindent\textcolor{FuncColor}{$\triangleright$\enspace\texttt{ComplexConjugate({\mdseries\slshape z})\index{ComplexConjugate@\texttt{ComplexConjugate}}
\label{ComplexConjugate}
}\hfill{\scriptsize (attribute)}}\\
\noindent\textcolor{FuncColor}{$\triangleright$\enspace\texttt{RealPart({\mdseries\slshape z})\index{RealPart@\texttt{RealPart}}
\label{RealPart}
}\hfill{\scriptsize (attribute)}}\\
\noindent\textcolor{FuncColor}{$\triangleright$\enspace\texttt{ImaginaryPart({\mdseries\slshape z})\index{ImaginaryPart@\texttt{ImaginaryPart}}
\label{ImaginaryPart}
}\hfill{\scriptsize (attribute)}}\\


 For a cyclotomic number \mbox{\texttt{\mdseries\slshape z}}, \texttt{ComplexConjugate} returns \texttt{GaloisCyc( \mbox{\texttt{\mdseries\slshape z}}, \texttt{\symbol{45}}1 )}, see{\nobreakspace}\texttt{GaloisCyc} (\ref{GaloisCyc:for a cyclotomic}). For a quaternion $\mbox{\texttt{\mdseries\slshape z}} = c_1 e + c_2 i + c_3 j + c_4 k$, \texttt{ComplexConjugate} returns $c_1 e - c_2 i - c_3 j - c_4 k$, see{\nobreakspace}\texttt{IsQuaternion} (\ref{IsQuaternion}). 

 When \texttt{ComplexConjugate} is called with a list then the result is the list of return values of \texttt{ComplexConjugate} for the list entries in the corresponding positions. 

 When \texttt{ComplexConjugate} is defined for an object \mbox{\texttt{\mdseries\slshape z}} then \texttt{RealPart} and \texttt{ImaginaryPart} return \texttt{(\mbox{\texttt{\mdseries\slshape z}} + ComplexConjugate( \mbox{\texttt{\mdseries\slshape z}} )) / 2} and \texttt{(\mbox{\texttt{\mdseries\slshape z}} \texttt{\symbol{45}} ComplexConjugate( \mbox{\texttt{\mdseries\slshape z}} )) / 2 i}, respectively, where \texttt{i} denotes the corresponding imaginary unit. 

 
\begin{Verbatim}[commandchars=!@|,fontsize=\small,frame=single,label=Example]
  !gapprompt@gap>| !gapinput@GaloisCyc( E(5) + E(5)^4, 2 );|
  E(5)^2+E(5)^3
  !gapprompt@gap>| !gapinput@GaloisCyc( E(5), -1 );           # the complex conjugate|
  E(5)^4
  !gapprompt@gap>| !gapinput@GaloisCyc( E(5) + E(5)^4, -1 );  # this value is real|
  E(5)+E(5)^4
  !gapprompt@gap>| !gapinput@GaloisCyc( E(15) + E(15)^4, 3 );|
  E(5)+E(5)^4
  !gapprompt@gap>| !gapinput@ComplexConjugate( E(7) );|
  E(7)^6
\end{Verbatim}
 }

 

\subsection{\textcolor{Chapter }{StarCyc}}
\logpage{[ 18, 5, 3 ]}\nobreak
\hyperdef{L}{X7E361C057E97CA66}{}
{\noindent\textcolor{FuncColor}{$\triangleright$\enspace\texttt{StarCyc({\mdseries\slshape cyc})\index{StarCyc@\texttt{StarCyc}}
\label{StarCyc}
}\hfill{\scriptsize (function)}}\\


 If the cyclotomic \mbox{\texttt{\mdseries\slshape cyc}} is an irrational element of a quadratic extension of the rationals then \texttt{StarCyc} returns the unique Galois conjugate of \mbox{\texttt{\mdseries\slshape cyc}} that is different from \mbox{\texttt{\mdseries\slshape cyc}}, otherwise \texttt{fail} is returned. In the first case, the return value is often called \mbox{\texttt{\mdseries\slshape cyc}}$*$ (see{\nobreakspace}\ref{Printing Character Tables}). 

 
\begin{Verbatim}[commandchars=!@|,fontsize=\small,frame=single,label=Example]
  !gapprompt@gap>| !gapinput@StarCyc( EB(5) ); StarCyc( E(5) );|
  E(5)^2+E(5)^3
  fail
\end{Verbatim}
 }

 

\subsection{\textcolor{Chapter }{Quadratic}}
\logpage{[ 18, 5, 4 ]}\nobreak
\hyperdef{L}{X84438F867B0CC299}{}
{\noindent\textcolor{FuncColor}{$\triangleright$\enspace\texttt{Quadratic({\mdseries\slshape cyc})\index{Quadratic@\texttt{Quadratic}}
\label{Quadratic}
}\hfill{\scriptsize (function)}}\\


 Let \mbox{\texttt{\mdseries\slshape cyc}} be a cyclotomic integer that lies in a quadratic extension field of the
rationals. Then we have \mbox{\texttt{\mdseries\slshape cyc}}$ = (a + b \sqrt{{n}}) / d$, for integers $a$, $b$, $n$, $d$, such that $d$ is either $1$ or $2$. In this case, \texttt{Quadratic} returns a record with the components \texttt{a}, \texttt{b}, \texttt{root}, \texttt{d}, \texttt{ATLAS}, and \texttt{display}; the values of the first four are $a$, $b$, $n$, and $d$, the \texttt{ATLAS} value is a (not necessarily shortest) representation of \mbox{\texttt{\mdseries\slshape cyc}} in terms of the \textsf{Atlas} irrationalities $b_{{|n|}}$, $i_{{|n|}}$, $r_{{|n|}}$, and the \texttt{display} value is a string that expresses \mbox{\texttt{\mdseries\slshape cyc}} in \textsf{GAP} notation, corresponding to the value of the \texttt{ATLAS} component. 

 If \mbox{\texttt{\mdseries\slshape cyc}} is not a cyclotomic integer or does not lie in a quadratic extension field of
the rationals then \texttt{fail} is returned. 

 If the denominator $d$ is $2$ then necessarily $n$ is congruent to $1$ modulo $4$, and $r_n$, $i_n$ are not possible; we have \texttt{\mbox{\texttt{\mdseries\slshape cyc}} = x + y * EB( root )} with \texttt{y = b}, \texttt{x = ( a + b ) / 2}. 

 If $d = 1$, we have the possibilities $i_{{|n|}}$ for $n < -1$, $a + b * i$ for $n = -1$, $a + b * r_n$ for $n > 0$. Furthermore if $n$ is congruent to $1$ modulo $4$, also \mbox{\texttt{\mdseries\slshape cyc}} $= (a+b) + 2 * b * b_{{|n|}}$ is possible; the shortest string of these is taken as the value for the
component \texttt{ATLAS}. 

 
\begin{Verbatim}[commandchars=!@|,fontsize=\small,frame=single,label=Example]
  !gapprompt@gap>| !gapinput@Quadratic( EB(5) ); Quadratic( EB(27) );|
  rec( ATLAS := "b5", a := -1, b := 1, d := 2,
    display := "(-1+Sqrt(5))/2", root := 5 )
  rec( ATLAS := "1+3b3", a := -1, b := 3, d := 2,
    display := "(-1+3*Sqrt(-3))/2", root := -3 )
  !gapprompt@gap>| !gapinput@Quadratic(0); Quadratic( E(5) );|
  rec( ATLAS := "0", a := 0, b := 0, d := 1, display := "0", root := 1 )
  fail
\end{Verbatim}
 }

 

\subsection{\textcolor{Chapter }{GaloisMat}}
\logpage{[ 18, 5, 5 ]}\nobreak
\hyperdef{L}{X7DDDEC3F80543B7D}{}
{\noindent\textcolor{FuncColor}{$\triangleright$\enspace\texttt{GaloisMat({\mdseries\slshape mat})\index{GaloisMat@\texttt{GaloisMat}}
\label{GaloisMat}
}\hfill{\scriptsize (attribute)}}\\


 Let \mbox{\texttt{\mdseries\slshape mat}} be a matrix of cyclotomics. \texttt{GaloisMat} calculates the complete orbits under the operation of the Galois group of the
(irrational) entries of \mbox{\texttt{\mdseries\slshape mat}}, and the permutations of rows corresponding to the generators of the Galois
group. 

 If some rows of \mbox{\texttt{\mdseries\slshape mat}} are identical, only the first one is considered for the permutations, and a
warning will be printed. 

 \texttt{GaloisMat} returns a record with the components \texttt{mat}, \texttt{galoisfams}, and \texttt{generators}. 

 
\begin{description}
\item[{\texttt{mat}}]  a list with initial segment being the rows of \mbox{\texttt{\mdseries\slshape mat}} (\emph{not} shallow copies of these rows); the list consists of full orbits under the
action of the Galois group of the entries of \mbox{\texttt{\mdseries\slshape mat}} defined above. The last rows in the list are those not contained in \mbox{\texttt{\mdseries\slshape mat}} but must be added in order to complete the orbits; so if the orbits were
already complete, \mbox{\texttt{\mdseries\slshape mat}} and \texttt{mat} have identical rows. 
\item[{\texttt{galoisfams}}]  a list that has the same length as the \texttt{mat} component, its entries are either 1, 0, \texttt{\symbol{45}}1, or lists. 
\begin{description}
\item[{\texttt{galoisfams[i] = 1}}]  means that \texttt{mat[i]} consists of rationals, i.e., \texttt{[ mat[i] ]} forms an orbit; 
\item[{\texttt{galoisfams[i] = \texttt{\symbol{45}}1}}]  means that \texttt{mat[i]} contains unknowns (see Chapter{\nobreakspace}\ref{Unknowns}); in this case \texttt{[ mat[i] ]} is regarded as an orbit, too, even if \texttt{mat[i]} contains irrational entries; 
\item[{\texttt{galoisfams[i] = }$[ l_1, l_2 ]$}]  (a list) means that \texttt{mat[i]} is the first element of its orbit in \texttt{mat}, $l_1$ is the list of positions of rows that form the orbit, and $l_2$ is the list of corresponding Galois automorphisms (as exponents, not as
functions); so we have \texttt{mat}$[ l_1[j] ][k] = $ \texttt{GaloisCyc( mat}$[i][k], l_2[j]$\texttt{ )}; 
\item[{\texttt{galoisfams[i] = 0}}]  means that \texttt{mat[i]} is an element of a nontrivial orbit but not the first element of it. 
\end{description}
 
\item[{\texttt{generators}}]  a list of permutations generating the permutation group corresponding to the
action of the Galois group on the rows of \texttt{mat}. 
\end{description}
 

 
\begin{Verbatim}[commandchars=!@|,fontsize=\small,frame=single,label=Example]
  !gapprompt@gap>| !gapinput@GaloisMat( [ [ E(3), E(4) ] ] );|
  rec( galoisfams := [ [ [ 1, 2, 3, 4 ], [ 1, 7, 5, 11 ] ], 0, 0, 0 ],
    generators := [ (1,2)(3,4), (1,3)(2,4) ],
    mat := [ [ E(3), E(4) ], [ E(3), -E(4) ], [ E(3)^2, E(4) ],
        [ E(3)^2, -E(4) ] ] )
  !gapprompt@gap>| !gapinput@GaloisMat( [ [ 1, 1, 1 ], [ 1, E(3), E(3)^2 ] ] );|
  rec( galoisfams := [ 1, [ [ 2, 3 ], [ 1, 2 ] ], 0 ],
    generators := [ (2,3) ],
    mat := [ [ 1, 1, 1 ], [ 1, E(3), E(3)^2 ], [ 1, E(3)^2, E(3) ] ] )
\end{Verbatim}
 }

 

\subsection{\textcolor{Chapter }{RationalizedMat}}
\logpage{[ 18, 5, 6 ]}\nobreak
\hyperdef{L}{X7BB9F5957AA8C082}{}
{\noindent\textcolor{FuncColor}{$\triangleright$\enspace\texttt{RationalizedMat({\mdseries\slshape mat})\index{RationalizedMat@\texttt{RationalizedMat}}
\label{RationalizedMat}
}\hfill{\scriptsize (attribute)}}\\


 returns the list of rationalized rows of \mbox{\texttt{\mdseries\slshape mat}}, which must be a matrix of cyclotomics. This is the set of sums over orbits
under the action of the Galois group of the entries of \mbox{\texttt{\mdseries\slshape mat}} (see \texttt{GaloisMat} (\ref{GaloisMat})), so the operation may be viewed as a kind of trace on the rows. 

 Note that no two rows of \mbox{\texttt{\mdseries\slshape mat}} should be equal. 

 
\begin{Verbatim}[commandchars=!@|,fontsize=\small,frame=single,label=Example]
  !gapprompt@gap>| !gapinput@mat:= [ [ 1, 1, 1 ], [ 1, E(3), E(3)^2 ], [ 1, E(3)^2, E(3) ] ];;|
  !gapprompt@gap>| !gapinput@RationalizedMat( mat );|
  [ [ 1, 1, 1 ], [ 2, -1, -1 ] ]
\end{Verbatim}
 }

 }

  
\section{\textcolor{Chapter }{Internally Represented Cyclotomics}}\label{Internally Represented Cyclotomics}
\logpage{[ 18, 6, 0 ]}
\hyperdef{L}{X8557FC2D7ACD6105}{}
{
  The implementation of an \emph{internally represented cyclotomic} is based on a list of length equal to its conductor. This means that the
internal representation of a cyclotomic does \emph{not} refer to the smallest number field but the smallest \emph{cyclotomic} field containing it. The reason for this is the wish to reflect the natural
embedding of two cyclotomic fields into a larger one that contains both. With
such embeddings, it is easy to construct the sum or the product of two
arbitrary cyclotomics (in possibly different fields) as an element of a
cyclotomic field. 

 The disadvantage of this approach is that the arithmetical operations are
quite expensive, so the use of internally represented cyclotomics is not
recommended for doing arithmetics over number fields, such as calculations
with matrices of cyclotomics. But internally represented cyclotomics are good
enough for dealing with irrationalities in character tables (see
chapter{\nobreakspace}\ref{Character Tables}). 

 For the representation of cyclotomics one has to recall that the $n$\texttt{\symbol{45}}th cyclotomic field ${\ensuremath{\mathbb Q}}(e_n)$ is a vector space of dimension $\varphi(n)$ over the rationals where $\varphi$ denotes Euler's phi\texttt{\symbol{45}}function (see{\nobreakspace}\texttt{Phi} (\ref{Phi})). 

 A special integral basis of cyclotomic fields is chosen that allows one to
easily convert arbitrary sums of roots of unity into the basis, as well as to
convert a cyclotomic represented w.r.t.{\nobreakspace}the basis into the
smallest possible cyclotomic field. This basis is accessible in \textsf{GAP}, see{\nobreakspace}\ref{Integral Bases of Abelian Number Fields} for more information and references. 

 Note that the set of all $n$\texttt{\symbol{45}}th roots of unity is linearly dependent for $n > 1$, so multiplication is \emph{not} the multiplication of the group ring ${\ensuremath{\mathbb Q}}\langle e_n \rangle$; given a ${\ensuremath{\mathbb Q}}$\texttt{\symbol{45}}basis of ${\ensuremath{\mathbb Q}}(e_n)$ the result of the multiplication (computed as multiplication of polynomials in $e_n$, using $(e_n)^n = 1$) will be converted to the basis. 

 
\begin{Verbatim}[commandchars=!@|,fontsize=\small,frame=single,label=Example]
  !gapprompt@gap>| !gapinput@E(5) * E(5)^2; ( E(5) + E(5)^4 ) * E(5)^2;|
  E(5)^3
  E(5)+E(5)^3
  !gapprompt@gap>| !gapinput@( E(5) + E(5)^4 ) * E(5);|
  -E(5)-E(5)^3-E(5)^4
\end{Verbatim}
 

 An internally represented cyclotomic is always represented in the smallest
cyclotomic field it is contained in. The internal coefficients list coincides
with the external representation returned by \texttt{ExtRepOfObj} (\ref{ExtRepOfObj:for a cyclotomic}). 

 To avoid calculations becoming unintentionally very long, or consuming very
large amounts of memory, there is a limit on the conductor of internally
represented cyclotomics, by default set to one million. This can be raised
(although not lowered) using \texttt{SetCyclotomicsLimit} (\ref{SetCyclotomicsLimit}) and accessed using \texttt{GetCyclotomicsLimit} (\ref{GetCyclotomicsLimit}). The maximum value of the limit is $2^{28}-1$ on $32$ bit systems, and $2^{32}-1$ on $64$ bit systems. So the maximal cyclotomic field implemented in \textsf{GAP} is not really the field ${\ensuremath{\mathbb Q}}^{ab}$. 

 It should be emphasized that one disadvantage of representing a cyclotomic in
the smallest \emph{cyclotomic} field (and not in the smallest field) is that arithmetic operations in a fixed
small extension field of the rational number field are comparatively
expensive. For example, take a prime integer $p$ and suppose that we want to work with a matrix group over the field ${\ensuremath{\mathbb Q}}(\sqrt{{p}})$. Then each matrix entry could be described by two rational coefficients,
whereas the representation in the smallest cyclotomic field requires $p-1$ rational coefficients for each entry. So it is worth thinking about using
elements in a field constructed with \texttt{AlgebraicExtension} (\ref{AlgebraicExtension}) when natural embeddings of cyclotomic fields are not needed.  

\subsection{\textcolor{Chapter }{SetCyclotomicsLimit}}
\logpage{[ 18, 6, 1 ]}\nobreak
\hyperdef{L}{X7D3028777DE39709}{}
{\noindent\textcolor{FuncColor}{$\triangleright$\enspace\texttt{SetCyclotomicsLimit({\mdseries\slshape newlimit})\index{SetCyclotomicsLimit@\texttt{SetCyclotomicsLimit}}
\label{SetCyclotomicsLimit}
}\hfill{\scriptsize (function)}}\\
\noindent\textcolor{FuncColor}{$\triangleright$\enspace\texttt{GetCyclotomicsLimit({\mdseries\slshape })\index{GetCyclotomicsLimit@\texttt{GetCyclotomicsLimit}}
\label{GetCyclotomicsLimit}
}\hfill{\scriptsize (function)}}\\


 \texttt{GetCyclotomicsLimit} returns the current limit on conductors of internally represented cyclotomic
numbers 

 \texttt{SetCyclotomicsLimit} can be called to increase the limit on conductors of internally represented
cyclotomic numbers. Note that computing in large cyclotomic fields using this
representation can be both slow and memory\texttt{\symbol{45}}consuming, and
that other approaches may be better for some problems. See \ref{Internally Represented Cyclotomics}. }

 }

 }

    
\chapter{\textcolor{Chapter }{Floats}}\label{Floats}
\logpage{[ 19, 0, 0 ]}
\hyperdef{L}{X81AA901181CA568F}{}
{
  Starting with version 4.5, \textsf{GAP} has built\texttt{\symbol{45}}in support for floating\texttt{\symbol{45}}point
numbers in machine format, and allows package to implement
arbitrary\texttt{\symbol{45}}precision floating\texttt{\symbol{45}}point
arithmetic in a uniform manner. For now, one such package, \textsf{Float} exists, and is based on the arbitrary\texttt{\symbol{45}}precision routines in \textsf{mpfr}. 

 A word of caution: \textsf{GAP} deals primarily with algebraic objects, which can be represented exactly in a
computer. Numerical imprecision means that floating\texttt{\symbol{45}}point
numbers do not form a ring in the strict \textsf{GAP} sense, because addition is in general not associative (\texttt{(1.0e\texttt{\symbol{45}}100+1.0)\texttt{\symbol{45}}1.0} is not the same as \texttt{1.0e\texttt{\symbol{45}}100+(1.0\texttt{\symbol{45}}1.0)}, in the default precision setting). 

 Most algorithms in \textsf{GAP} which require ring elements will therefore not be applicable to
floating\texttt{\symbol{45}}point elements. In some cases, such a notion would
not even make any sense (what is the greatest common divisor of two
floating\texttt{\symbol{45}}point numbers?) 
\section{\textcolor{Chapter }{A sample run}}\logpage{[ 19, 1, 0 ]}
\hyperdef{L}{X7B4092CA7ABB93B0}{}
{
 Floating\texttt{\symbol{45}}point numbers can be input into \textsf{GAP} in the standard floating\texttt{\symbol{45}}point notation: 

 
\begin{Verbatim}[commandchars=!@|,fontsize=\small,frame=single,label=Example]
  !gapprompt@gap>| !gapinput@3.14;|
  3.14
  !gapprompt@gap>| !gapinput@last^2/6;|
  1.64327
  !gapprompt@gap>| !gapinput@h := 6.62606896e-34;|
  6.62607e-34
  !gapprompt@gap>| !gapinput@pi := 4*Atan(1.0);|
  3.14159
  !gapprompt@gap>| !gapinput@hbar := h/(2*pi);|
  1.05457e-34
\end{Verbatim}
 

 Floating\texttt{\symbol{45}}point numbers can also be created using \texttt{Float}, from strings or rational numbers; and can be converted back using \texttt{String,Rat,Int}. 

 \textsf{GAP} allows rational and floating\texttt{\symbol{45}}point numbers to be mixed in
the elementary operations \texttt{+,\texttt{\symbol{45}},*,/}. However, floating\texttt{\symbol{45}}point numbers and rational numbers may
not be compared. Conversions are performed using the creator \texttt{Float}: 

 
\begin{Verbatim}[commandchars=!@|,fontsize=\small,frame=single,label=Example]
  !gapprompt@gap>| !gapinput@Float("3.1416");|
  3.1416
  !gapprompt@gap>| !gapinput@Float(355/113);|
  3.14159
  !gapprompt@gap>| !gapinput@Rat(last);|
  355/113
  !gapprompt@gap>| !gapinput@Rat(0.33333);|
  1/3
  !gapprompt@gap>| !gapinput@Int(1.e10);|
  10000000000
  !gapprompt@gap>| !gapinput@Int(1.e20);|
  100000000000000000000
  !gapprompt@gap>| !gapinput@Int(1.e30);|
  1000000000000000019884624838656
\end{Verbatim}
 }

 
\section{\textcolor{Chapter }{Methods}}\logpage{[ 19, 2, 0 ]}
\hyperdef{L}{X8606FDCE878850EF}{}
{
 Floating\texttt{\symbol{45}}point numbers may be directly input, as in any
usual mathematical software or language; with the exception that every
floating\texttt{\symbol{45}}point number must contain a decimal digit.
Therefore \texttt{.1}, \texttt{.1e1}, \texttt{\texttt{\symbol{45}}.999} etc. are all valid \textsf{GAP} inputs. 

 Floating\texttt{\symbol{45}}point numbers so entered in \textsf{GAP} are stored as strings. They are converted to floating\texttt{\symbol{45}}point
when they are first used. This means that, if the
floating\texttt{\symbol{45}}point precision is increased, the constants are
reevaluated to fit the new format. 

 Floating\texttt{\symbol{45}}point numbers may be followed by an underscore, as
in \texttt{1.{\textunderscore}}. This means that they are to be immediately converted to the current
floating\texttt{\symbol{45}}point format. The underscore may be followed by a
single letter, which specifies which format/precision to use. By default, \textsf{GAP} has a single floating\texttt{\symbol{45}}point handler, with fixed (53 bits)
precision, and its format specifier is \texttt{'l'} as in \texttt{1.{\textunderscore}l}. Higher\texttt{\symbol{45}}precision floating\texttt{\symbol{45}}point
computations is available via external packages; \textsf{float} for example. 

 A record, \texttt{FLOAT} (\ref{FLOAT:constants}), contains all relevant constants for the current
floating\texttt{\symbol{45}}point format; see its documentation for details.
Typical fields are \texttt{FLOAT.MANT{\textunderscore}DIG=53}, the constant \texttt{FLOAT.VIEW{\textunderscore}DIG=6} specifying the number of digits to view, and \texttt{FLOAT.PI} for the constant $\pi$. The constants have the same name as their C counterparts, except for the
missing initial \texttt{DBL{\textunderscore}} or \texttt{M{\textunderscore}}. 

 Floating\texttt{\symbol{45}}point numbers may be created using the single
function \texttt{Float} (\ref{Float}), which accepts as arguments rational, string, or
floating\texttt{\symbol{45}}point numbers. Floating\texttt{\symbol{45}}point
numbers may also be created, in any floating\texttt{\symbol{45}}point
representation, using \texttt{NewFloat} (\ref{NewFloat}) as in \texttt{NewFloat(IsIEEE754FloatRep,355/113)}, by supplying the category filter of the desired new
floating\texttt{\symbol{45}}point number; or using \texttt{MakeFloat} (\ref{MakeFloat}) as in \texttt{MakeFloat(1.0,355/113)}, by supplying a sample floating\texttt{\symbol{45}}point number. 

 Floating\texttt{\symbol{45}}point numbers may also be converted to other \textsf{GAP} formats using the usual commands \texttt{Int} (\ref{Int}), \texttt{Rat} (\ref{Rat}), \texttt{String} (\ref{String}). 

 Exact conversion to and from floating\texttt{\symbol{45}}point format may be
done using external representations. The "external representation" of a
floating\texttt{\symbol{45}}point number \texttt{x} is a pair \texttt{[m,e]} of integers, such that \texttt{x=m*2\texttt{\symbol{94}}(\texttt{\symbol{45}}1+e\texttt{\symbol{45}}LogInt(AbsInt(m),2))}. Conversion to and from external representation is performed as usual using \texttt{ExtRepOfObj} (\ref{ExtRepOfObj}) and \texttt{ObjByExtRep} (\ref{ObjByExtRep}): 
\begin{Verbatim}[commandchars=!@|,fontsize=\small,frame=single,label=Example]
  !gapprompt@gap>| !gapinput@ExtRepOfObj(3.14);|
  [ 7070651414971679, 2 ]
  !gapprompt@gap>| !gapinput@ObjByExtRep(IEEE754FloatsFamily,last);|
  3.14
\end{Verbatim}
 

 Computations with floating\texttt{\symbol{45}}point numbers never raise any
error. Division by zero is allowed, and produces a signed infinity. Illegal
operations, such as \texttt{0./0.}, produce \texttt{NaN}'s (not\texttt{\symbol{45}}a\texttt{\symbol{45}}number); this is the only
floating\texttt{\symbol{45}}point number \texttt{x} such that \texttt{not EqFloat(x+0.0,x)}. 

 The IEEE754 standard requires \texttt{NaN} to be non\texttt{\symbol{45}}equal to itself. On the other hand, \textsf{GAP} requires every object to be equal to itself. To respect the IEEE754 standard,
the function \texttt{EqFloat} (\ref{EqFloat}) should be used instead of \texttt{=}. 

 The category a floating\texttt{\symbol{45}}point belongs to can be checked
using the filters \texttt{IsFinite} (\ref{IsFinite}), \texttt{IsPInfinity} (\ref{IsPInfinity}), \texttt{IsNInfinity} (\ref{IsNInfinity}), \texttt{IsXInfinity} (\ref{IsXInfinity}), \texttt{IsNaN} (\ref{IsNaN}). 

 Comparisons between floating\texttt{\symbol{45}}point numbers and rationals
are explicitly forbidden. The rationale is that objects belonging to different
families should in general not be comparable in \textsf{GAP}. Floating\texttt{\symbol{45}}point numbers are also approximations of real
numbers, and don't follow the same rules; consider for example, using the
default \textsf{GAP} implementation of floating\texttt{\symbol{45}}point numbers, 
\begin{Verbatim}[commandchars=!@|,fontsize=\small,frame=single,label=Example]
  !gapprompt@gap>| !gapinput@1.0/3.0 = Float(1/3);|
  true
  !gapprompt@gap>| !gapinput@(1.0/3.0)^5 = Float((1/3)^5);|
  false
\end{Verbatim}
 

 
\subsection{\textcolor{Chapter }{Float creators}}\logpage{[ 19, 2, 1 ]}
\hyperdef{L}{X86D5EA93813FB6C4}{}
{
\noindent\textcolor{FuncColor}{$\triangleright$\enspace\texttt{Float({\mdseries\slshape obj})\index{Float@\texttt{Float}}
\label{Float}
}\hfill{\scriptsize (function)}}\\
\noindent\textcolor{FuncColor}{$\triangleright$\enspace\texttt{NewFloat({\mdseries\slshape filter, obj})\index{NewFloat@\texttt{NewFloat}}
\label{NewFloat}
}\hfill{\scriptsize (constructor)}}\\
\noindent\textcolor{FuncColor}{$\triangleright$\enspace\texttt{MakeFloat({\mdseries\slshape sample, obj, obj})\index{MakeFloat@\texttt{MakeFloat}}
\label{MakeFloat}
}\hfill{\scriptsize (operation)}}\\
\textbf{\indent Returns:\ }
A new floating\texttt{\symbol{45}}point number, based on \mbox{\texttt{\mdseries\slshape obj}}



 This function creates a new floating\texttt{\symbol{45}}point number. 

 If \mbox{\texttt{\mdseries\slshape obj}} is a rational number, the created number is created with sufficient precision
so that the number can (usually) be converted back to the original number (see \texttt{Rat} (\textbf{Reference: Rat}) and \texttt{Rat} (\ref{Rat})). For an integer, the precision, if unspecified, is chosen sufficient so that \texttt{Int(Float(obj))=obj} always holds, but at least 64 bits. 

 \mbox{\texttt{\mdseries\slshape obj}} may also be a string, which may be of the form \texttt{"3.14e0"} or \texttt{".314e1"} or \texttt{".314@1"} etc. 

 An option may be passed to specify, it bits, a desired precision. The format
is \texttt{Float("3.14":PrecisionFloat:=1000)} to create a 1000\texttt{\symbol{45}}bit approximation of $3.14$. 

 In particular, if \mbox{\texttt{\mdseries\slshape obj}} is already a floating\texttt{\symbol{45}}point number, then \texttt{Float(obj:PrecisionFloat:=prec)} creates a copy of \mbox{\texttt{\mdseries\slshape obj}} with a new precision. prec }

 

\subsection{\textcolor{Chapter }{Rat (for floats)}}
\logpage{[ 19, 2, 2 ]}\nobreak
\hyperdef{L}{X7BCD34DC7B5A0521}{}
{\noindent\textcolor{FuncColor}{$\triangleright$\enspace\texttt{Rat({\mdseries\slshape f})\index{Rat@\texttt{Rat}!for floats}
\label{Rat:for floats}
}\hfill{\scriptsize (attribute)}}\\
\textbf{\indent Returns:\ }
A rational approximation to \mbox{\texttt{\mdseries\slshape f}}



 This command constructs a rational approximation to the
floating\texttt{\symbol{45}}point number \mbox{\texttt{\mdseries\slshape f}}. Of course, it is not guaranteed to return the original rational number \mbox{\texttt{\mdseries\slshape f}} was created from, though it returns the most `reasonable' one given the
precision of \mbox{\texttt{\mdseries\slshape f}}. 

 Two options control the precision of the rational approximation: In the form \texttt{Rat(f:maxdenom:=md,maxpartial:=mp)}, the rational returned is such that the denominator is at most \mbox{\texttt{\mdseries\slshape md}} and the partials in its continued fraction expansion are at most \mbox{\texttt{\mdseries\slshape mp}}. The default values are \texttt{maxpartial:=10000} and \texttt{maxdenom:=2\texttt{\symbol{94}}(precision/2)}. }

 

\subsection{\textcolor{Chapter }{Cyc (for floats)}}
\logpage{[ 19, 2, 3 ]}\nobreak
\hyperdef{L}{X7D1EAE11844625F4}{}
{\noindent\textcolor{FuncColor}{$\triangleright$\enspace\texttt{Cyc({\mdseries\slshape f[, degree]})\index{Cyc@\texttt{Cyc}!for floats}
\label{Cyc:for floats}
}\hfill{\scriptsize (operation)}}\\
\textbf{\indent Returns:\ }
A cyclotomic approximation to \mbox{\texttt{\mdseries\slshape f}}



 This command constructs a cyclotomic approximation to the
floating\texttt{\symbol{45}}point number \mbox{\texttt{\mdseries\slshape f}}. Of course, it is not guaranteed to return the original rational number \mbox{\texttt{\mdseries\slshape f}} was created from, though it returns the most `reasonable' one given the
precision of \mbox{\texttt{\mdseries\slshape f}}. An optional argument \mbox{\texttt{\mdseries\slshape degree}} specifies the maximal degree of the cyclotomic to be constructed. 

 The method used is LLL lattice reduction. }

 

\subsection{\textcolor{Chapter }{SetFloats}}
\logpage{[ 19, 2, 4 ]}\nobreak
\hyperdef{L}{X7A962B0983FA66E8}{}
{\noindent\textcolor{FuncColor}{$\triangleright$\enspace\texttt{SetFloats({\mdseries\slshape rec[, bits][, install]})\index{SetFloats@\texttt{SetFloats}}
\label{SetFloats}
}\hfill{\scriptsize (function)}}\\


 Installs a new interface to floating\texttt{\symbol{45}}point numbers in \textsf{GAP}, optionally with a desired precision \mbox{\texttt{\mdseries\slshape bits}} in binary digits. The last optional argument \mbox{\texttt{\mdseries\slshape install}} is a boolean value; if false, it only installs the eager handler and the
precision for the floateans, without making them the default. }

 

\subsection{\textcolor{Chapter }{FLOAT (constants)}}
\logpage{[ 19, 2, 5 ]}\nobreak
\hyperdef{L}{X819050BF8403806E}{}
{\noindent\textcolor{FuncColor}{$\triangleright$\enspace\texttt{FLOAT\index{FLOAT@\texttt{FLOAT}!constants}
\label{FLOAT:constants}
}\hfill{\scriptsize (global variable)}}\\


 This record contains useful floating\texttt{\symbol{45}}point constants: 
\begin{description}
\item[{DECIMAL{\textunderscore}DIG}] Maximal number of useful digits;
\item[{DIG}] Number of significant digits;
\item[{VIEW{\textunderscore}DIG}] Number of digits to print in short view;
\item[{EPSILON}] Smallest number such that $1\neq1+\epsilon$;
\item[{MANT{\textunderscore}DIG}] Number of bits in the mantissa;
\item[{MAX}] Maximal representable number;
\item[{MAX{\textunderscore}10{\textunderscore}EXP}] Maximal decimal exponent;
\item[{MAX{\textunderscore}EXP}] Maximal binary exponent;
\item[{MIN}] Minimal positive representable number;
\item[{MIN{\textunderscore}10{\textunderscore}EXP}] Minimal decimal exponent;
\item[{MIN{\textunderscore}EXP}] Minimal exponent;
\item[{INFINITY}] Positive infinity;
\item[{NINFINITY}] Negative infinity;
\item[{NAN}] Not\texttt{\symbol{45}}a\texttt{\symbol{45}}number,
\end{description}
 as well as mathematical constants \texttt{E}, \texttt{LOG2E}, \texttt{LOG10E}, \texttt{LN2}, \texttt{LN10}, \texttt{PI}, \texttt{PI{\textunderscore}2}, \texttt{PI{\textunderscore}4}, \texttt{1{\textunderscore}PI}, \texttt{2{\textunderscore}PI}, \texttt{2{\textunderscore}SQRTPI}, \texttt{SQRT2}, \texttt{SQRT1{\textunderscore}2}. }

 

\subsection{\textcolor{Chapter }{EqFloat}}
\logpage{[ 19, 2, 6 ]}\nobreak
\hyperdef{L}{X7BD96E0585D5A1EE}{}
{\noindent\textcolor{FuncColor}{$\triangleright$\enspace\texttt{EqFloat({\mdseries\slshape x, y})\index{EqFloat@\texttt{EqFloat}}
\label{EqFloat}
}\hfill{\scriptsize (operation)}}\\
\textbf{\indent Returns:\ }
Whether the floateans \mbox{\texttt{\mdseries\slshape x}} and \mbox{\texttt{\mdseries\slshape y}} are equal



 This function compares two floating\texttt{\symbol{45}}point numbers, and
returns \texttt{true} if they are equal, and \texttt{false} otherwise; with the exception that \texttt{NaN} is always considered to be different from itself. }

 

\subsection{\textcolor{Chapter }{PrecisionFloat}}
\logpage{[ 19, 2, 7 ]}\nobreak
\hyperdef{L}{X7B3133497DDE839B}{}
{\noindent\textcolor{FuncColor}{$\triangleright$\enspace\texttt{PrecisionFloat({\mdseries\slshape x})\index{PrecisionFloat@\texttt{PrecisionFloat}}
\label{PrecisionFloat}
}\hfill{\scriptsize (attribute)}}\\
\textbf{\indent Returns:\ }
The precision of \mbox{\texttt{\mdseries\slshape x}}



 This function returns the precision, counted in number of binary digits, of
the floating\texttt{\symbol{45}}point number \mbox{\texttt{\mdseries\slshape x}}. }

 

\subsection{\textcolor{Chapter }{SignBit}}
\logpage{[ 19, 2, 8 ]}\nobreak
\hyperdef{L}{X801753137949DD78}{}
{\noindent\textcolor{FuncColor}{$\triangleright$\enspace\texttt{SignBit({\mdseries\slshape x})\index{SignBit@\texttt{SignBit}}
\label{SignBit}
}\hfill{\scriptsize (attribute)}}\\
\noindent\textcolor{FuncColor}{$\triangleright$\enspace\texttt{SignFloat({\mdseries\slshape x})\index{SignFloat@\texttt{SignFloat}}
\label{SignFloat}
}\hfill{\scriptsize (attribute)}}\\
\textbf{\indent Returns:\ }
The sign of \mbox{\texttt{\mdseries\slshape x}}.



 The first function \texttt{SignBit} returns the sign bit of the floating\texttt{\symbol{45}}point number \mbox{\texttt{\mdseries\slshape x}}: \texttt{true} if \mbox{\texttt{\mdseries\slshape x}} is negative (including \texttt{\texttt{\symbol{45}}0.}) and \texttt{false} otherwise. 

 The second function \texttt{SignFloat} returns the integer \texttt{\texttt{\symbol{45}}1} if \mbox{\texttt{\mdseries\slshape x{\textless}0}}, \texttt{0} if \mbox{\texttt{\mdseries\slshape x=0}} and \texttt{1} if \mbox{\texttt{\mdseries\slshape x{\textgreater}0}}. }

 

\subsection{\textcolor{Chapter }{SinCos}}
\logpage{[ 19, 2, 9 ]}\nobreak
\hyperdef{L}{X7935C65D7B0F47C7}{}
{\noindent\textcolor{FuncColor}{$\triangleright$\enspace\texttt{SinCos({\mdseries\slshape x})\index{SinCos@\texttt{SinCos}}
\label{SinCos}
}\hfill{\scriptsize (attribute)}}\\
\textbf{\indent Returns:\ }
The list \texttt{[sin(x), cos(x)]}.



 The function returns a list with \texttt{sin} and \texttt{cos} of \mbox{\texttt{\mdseries\slshape x}}. }

 

\subsection{\textcolor{Chapter }{Atan2}}
\logpage{[ 19, 2, 10 ]}\nobreak
\hyperdef{L}{X846E1196844B9E11}{}
{\noindent\textcolor{FuncColor}{$\triangleright$\enspace\texttt{Atan2({\mdseries\slshape y, x})\index{Atan2@\texttt{Atan2}}
\label{Atan2}
}\hfill{\scriptsize (operation)}}\\
\textbf{\indent Returns:\ }
The polar angle of \mbox{\texttt{\mdseries\slshape (x, y)}} in the plane as float.



 Returns the principal value of the argument (polar angle) of $(\mbox{\texttt{\mdseries\slshape x}}, \mbox{\texttt{\mdseries\slshape y}})$ in the plane. The returned value will always be in $(-\pi , \pi]$ and is not defined on $(0,0)$. This function is defined in accordance with IEEE
1788\texttt{\symbol{45}}2015 and imported from IEEE 754. }

 

\subsection{\textcolor{Chapter }{Log1p}}
\logpage{[ 19, 2, 11 ]}\nobreak
\hyperdef{L}{X7981510D826EE3E5}{}
{\noindent\textcolor{FuncColor}{$\triangleright$\enspace\texttt{Log1p({\mdseries\slshape x})\index{Log1p@\texttt{Log1p}}
\label{Log1p}
}\hfill{\scriptsize (attribute)}}\\
\noindent\textcolor{FuncColor}{$\triangleright$\enspace\texttt{Expm1({\mdseries\slshape x})\index{Expm1@\texttt{Expm1}}
\label{Expm1}
}\hfill{\scriptsize (attribute)}}\\
\textbf{\indent Returns:\ }
The natural logarithm of $\mbox{\texttt{\mdseries\slshape x}}+1$ or exponential $-1$ of \mbox{\texttt{\mdseries\slshape x}} respectively.



 The first function \texttt{Log1p} returns the natural logarithm $log(\mbox{\texttt{\mdseries\slshape x}}+1)$. 

 The second function \texttt{Expm1} returns the exponential function $exp(\mbox{\texttt{\mdseries\slshape x}})-1$ 

 These two functions are inverse to each other. }

 

\subsection{\textcolor{Chapter }{Erf}}
\logpage{[ 19, 2, 12 ]}\nobreak
\hyperdef{L}{X86073E147FB3C0EA}{}
{\noindent\textcolor{FuncColor}{$\triangleright$\enspace\texttt{Erf({\mdseries\slshape x})\index{Erf@\texttt{Erf}}
\label{Erf}
}\hfill{\scriptsize (operation)}}\\
\textbf{\indent Returns:\ }
The error function given by the Gaussian integral



 Returns the error function imported from IEEE 754 given by the formula: 
\[ Erf(x) := \frac{2}{\sqrt{\pi}} \int_{0}^{x} exp(- t^2 ) dt \]
 }

 
\subsection{\textcolor{Chapter }{Infinity testers}}\logpage{[ 19, 2, 13 ]}
\hyperdef{L}{X7E03FDEE824D1E8E}{}
{
\noindent\textcolor{FuncColor}{$\triangleright$\enspace\texttt{IsPInfinity({\mdseries\slshape x})\index{IsPInfinity@\texttt{IsPInfinity}}
\label{IsPInfinity}
}\hfill{\scriptsize (property)}}\\
\noindent\textcolor{FuncColor}{$\triangleright$\enspace\texttt{IsNInfinity({\mdseries\slshape x})\index{IsNInfinity@\texttt{IsNInfinity}}
\label{IsNInfinity}
}\hfill{\scriptsize (property)}}\\
\noindent\textcolor{FuncColor}{$\triangleright$\enspace\texttt{IsXInfinity({\mdseries\slshape x})\index{IsXInfinity@\texttt{IsXInfinity}}
\label{IsXInfinity}
}\hfill{\scriptsize (property)}}\\
\noindent\textcolor{FuncColor}{$\triangleright$\enspace\texttt{IsFinite({\mdseries\slshape x})\index{IsFinite@\texttt{IsFinite}!for floats}
\label{IsFinite:for floats}
}\hfill{\scriptsize (property)}}\\
\noindent\textcolor{FuncColor}{$\triangleright$\enspace\texttt{IsNaN({\mdseries\slshape x})\index{IsNaN@\texttt{IsNaN}}
\label{IsNaN}
}\hfill{\scriptsize (property)}}\\


 Returns \texttt{true} if the floating\texttt{\symbol{45}}point number \mbox{\texttt{\mdseries\slshape x}} is respectively $+\infty$, $-\infty$, $\pm\infty$, finite, or `not a number', such as the result of \texttt{0.0/0.0}. }

 
\subsection{\textcolor{Chapter }{Standard mathematical operations}}\logpage{[ 19, 2, 14 ]}
\hyperdef{L}{X8151581186F75BA3}{}
{
\noindent\textcolor{FuncColor}{$\triangleright$\enspace\texttt{Sin({\mdseries\slshape f})\index{Sin@\texttt{Sin}}
\label{Sin}
}\hfill{\scriptsize (attribute)}}\\
\noindent\textcolor{FuncColor}{$\triangleright$\enspace\texttt{Cos({\mdseries\slshape f})\index{Cos@\texttt{Cos}}
\label{Cos}
}\hfill{\scriptsize (attribute)}}\\
\noindent\textcolor{FuncColor}{$\triangleright$\enspace\texttt{Tan({\mdseries\slshape f})\index{Tan@\texttt{Tan}}
\label{Tan}
}\hfill{\scriptsize (attribute)}}\\
\noindent\textcolor{FuncColor}{$\triangleright$\enspace\texttt{Sec({\mdseries\slshape f})\index{Sec@\texttt{Sec}}
\label{Sec}
}\hfill{\scriptsize (attribute)}}\\
\noindent\textcolor{FuncColor}{$\triangleright$\enspace\texttt{Csc({\mdseries\slshape f})\index{Csc@\texttt{Csc}}
\label{Csc}
}\hfill{\scriptsize (attribute)}}\\
\noindent\textcolor{FuncColor}{$\triangleright$\enspace\texttt{Cot({\mdseries\slshape f})\index{Cot@\texttt{Cot}}
\label{Cot}
}\hfill{\scriptsize (attribute)}}\\
\noindent\textcolor{FuncColor}{$\triangleright$\enspace\texttt{Asin({\mdseries\slshape f})\index{Asin@\texttt{Asin}}
\label{Asin}
}\hfill{\scriptsize (attribute)}}\\
\noindent\textcolor{FuncColor}{$\triangleright$\enspace\texttt{Acos({\mdseries\slshape f})\index{Acos@\texttt{Acos}}
\label{Acos}
}\hfill{\scriptsize (attribute)}}\\
\noindent\textcolor{FuncColor}{$\triangleright$\enspace\texttt{Atan({\mdseries\slshape f})\index{Atan@\texttt{Atan}}
\label{Atan}
}\hfill{\scriptsize (attribute)}}\\
\noindent\textcolor{FuncColor}{$\triangleright$\enspace\texttt{Sinh({\mdseries\slshape f})\index{Sinh@\texttt{Sinh}}
\label{Sinh}
}\hfill{\scriptsize (attribute)}}\\
\noindent\textcolor{FuncColor}{$\triangleright$\enspace\texttt{Cosh({\mdseries\slshape f})\index{Cosh@\texttt{Cosh}}
\label{Cosh}
}\hfill{\scriptsize (attribute)}}\\
\noindent\textcolor{FuncColor}{$\triangleright$\enspace\texttt{Tanh({\mdseries\slshape f})\index{Tanh@\texttt{Tanh}}
\label{Tanh}
}\hfill{\scriptsize (attribute)}}\\
\noindent\textcolor{FuncColor}{$\triangleright$\enspace\texttt{Sech({\mdseries\slshape f})\index{Sech@\texttt{Sech}}
\label{Sech}
}\hfill{\scriptsize (attribute)}}\\
\noindent\textcolor{FuncColor}{$\triangleright$\enspace\texttt{Csch({\mdseries\slshape f})\index{Csch@\texttt{Csch}}
\label{Csch}
}\hfill{\scriptsize (attribute)}}\\
\noindent\textcolor{FuncColor}{$\triangleright$\enspace\texttt{Coth({\mdseries\slshape f})\index{Coth@\texttt{Coth}}
\label{Coth}
}\hfill{\scriptsize (attribute)}}\\
\noindent\textcolor{FuncColor}{$\triangleright$\enspace\texttt{Asinh({\mdseries\slshape f})\index{Asinh@\texttt{Asinh}}
\label{Asinh}
}\hfill{\scriptsize (attribute)}}\\
\noindent\textcolor{FuncColor}{$\triangleright$\enspace\texttt{Acosh({\mdseries\slshape f})\index{Acosh@\texttt{Acosh}}
\label{Acosh}
}\hfill{\scriptsize (attribute)}}\\
\noindent\textcolor{FuncColor}{$\triangleright$\enspace\texttt{Atanh({\mdseries\slshape f})\index{Atanh@\texttt{Atanh}}
\label{Atanh}
}\hfill{\scriptsize (attribute)}}\\
\noindent\textcolor{FuncColor}{$\triangleright$\enspace\texttt{Log({\mdseries\slshape f})\index{Log@\texttt{Log}}
\label{Log}
}\hfill{\scriptsize (operation)}}\\
\noindent\textcolor{FuncColor}{$\triangleright$\enspace\texttt{Log2({\mdseries\slshape f})\index{Log2@\texttt{Log2}}
\label{Log2}
}\hfill{\scriptsize (attribute)}}\\
\noindent\textcolor{FuncColor}{$\triangleright$\enspace\texttt{Log10({\mdseries\slshape f})\index{Log10@\texttt{Log10}}
\label{Log10}
}\hfill{\scriptsize (attribute)}}\\
\noindent\textcolor{FuncColor}{$\triangleright$\enspace\texttt{Exp({\mdseries\slshape f})\index{Exp@\texttt{Exp}}
\label{Exp}
}\hfill{\scriptsize (attribute)}}\\
\noindent\textcolor{FuncColor}{$\triangleright$\enspace\texttt{Exp2({\mdseries\slshape f})\index{Exp2@\texttt{Exp2}}
\label{Exp2}
}\hfill{\scriptsize (attribute)}}\\
\noindent\textcolor{FuncColor}{$\triangleright$\enspace\texttt{Exp10({\mdseries\slshape f})\index{Exp10@\texttt{Exp10}}
\label{Exp10}
}\hfill{\scriptsize (attribute)}}\\
\noindent\textcolor{FuncColor}{$\triangleright$\enspace\texttt{CubeRoot({\mdseries\slshape f})\index{CubeRoot@\texttt{CubeRoot}}
\label{CubeRoot}
}\hfill{\scriptsize (attribute)}}\\
\noindent\textcolor{FuncColor}{$\triangleright$\enspace\texttt{Square({\mdseries\slshape f})\index{Square@\texttt{Square}}
\label{Square}
}\hfill{\scriptsize (attribute)}}\\
\noindent\textcolor{FuncColor}{$\triangleright$\enspace\texttt{Hypothenuse({\mdseries\slshape x, y})\index{Hypothenuse@\texttt{Hypothenuse}}
\label{Hypothenuse}
}\hfill{\scriptsize (operation)}}\\
\noindent\textcolor{FuncColor}{$\triangleright$\enspace\texttt{Ceil({\mdseries\slshape f})\index{Ceil@\texttt{Ceil}}
\label{Ceil}
}\hfill{\scriptsize (attribute)}}\\
\noindent\textcolor{FuncColor}{$\triangleright$\enspace\texttt{Floor({\mdseries\slshape f})\index{Floor@\texttt{Floor}}
\label{Floor}
}\hfill{\scriptsize (attribute)}}\\
\noindent\textcolor{FuncColor}{$\triangleright$\enspace\texttt{Round({\mdseries\slshape f})\index{Round@\texttt{Round}}
\label{Round}
}\hfill{\scriptsize (attribute)}}\\
\noindent\textcolor{FuncColor}{$\triangleright$\enspace\texttt{Trunc({\mdseries\slshape f})\index{Trunc@\texttt{Trunc}}
\label{Trunc}
}\hfill{\scriptsize (attribute)}}\\
\noindent\textcolor{FuncColor}{$\triangleright$\enspace\texttt{FrExp({\mdseries\slshape f})\index{FrExp@\texttt{FrExp}}
\label{FrExp}
}\hfill{\scriptsize (attribute)}}\\
\noindent\textcolor{FuncColor}{$\triangleright$\enspace\texttt{LdExp({\mdseries\slshape f, exp})\index{LdExp@\texttt{LdExp}}
\label{LdExp}
}\hfill{\scriptsize (operation)}}\\
\noindent\textcolor{FuncColor}{$\triangleright$\enspace\texttt{AbsoluteValue({\mdseries\slshape f})\index{AbsoluteValue@\texttt{AbsoluteValue}!for floats}
\label{AbsoluteValue:for floats}
}\hfill{\scriptsize (attribute)}}\\
\noindent\textcolor{FuncColor}{$\triangleright$\enspace\texttt{Norm({\mdseries\slshape f})\index{Norm@\texttt{Norm}!for floats}
\label{Norm:for floats}
}\hfill{\scriptsize (attribute)}}\\
\noindent\textcolor{FuncColor}{$\triangleright$\enspace\texttt{Frac({\mdseries\slshape f})\index{Frac@\texttt{Frac}}
\label{Frac}
}\hfill{\scriptsize (attribute)}}\\
\noindent\textcolor{FuncColor}{$\triangleright$\enspace\texttt{Zeta({\mdseries\slshape f})\index{Zeta@\texttt{Zeta}}
\label{Zeta}
}\hfill{\scriptsize (attribute)}}\\
\noindent\textcolor{FuncColor}{$\triangleright$\enspace\texttt{Gamma({\mdseries\slshape f})\index{Gamma@\texttt{Gamma}}
\label{Gamma}
}\hfill{\scriptsize (attribute)}}\\


 Standard math functions. Functions ending in an integer like \texttt{Log2}, \texttt{Log10}, \texttt{Exp2} and \texttt{Exp10} indicate the base used, in \texttt{log} and \texttt{exp} the natural base is used, i.e. $e$. }

 }

 
\section{\textcolor{Chapter }{High\texttt{\symbol{45}}precision\texttt{\symbol{45}}specific methods}}\logpage{[ 19, 3, 0 ]}
\hyperdef{L}{X845ACF3A78BD2771}{}
{
 \textsf{GAP} provides a mechanism for packages to implement new
floating\texttt{\symbol{45}}point numerical interfaces. The following
describes that mechanism, actual examples of packages are documented
separately. 

 A package must create a record with fields (all optional) 
\begin{description}
\item[{creator}] a function converting strings to floating\texttt{\symbol{45}}point;
\item[{eager}] a character allowing immediate conversion to
floating\texttt{\symbol{45}}point;
\item[{objbyextrep}] a function creating a floating\texttt{\symbol{45}}point number out of a list \texttt{[mantissa,exponent]};
\item[{filter}] a filter for the new floating\texttt{\symbol{45}}point objects;
\item[{constants}] a record containing numerical constants, such as \texttt{MANT{\textunderscore}DIG}, \texttt{MAX}, \texttt{MIN}, \texttt{NAN}.
\end{description}
 

 The package must install methods \texttt{Int}, \texttt{Rat}, \texttt{String} for its objects, and creators \texttt{NewFloat(filter,IsRat)}, \texttt{NewFloat(IsString)}. 

 It must then install methods for all arithmetic and numerical operations: \texttt{SUM}, \texttt{Exp}, ... 

 The user chooses that implementation by calling \texttt{SetFloats} (\ref{SetFloats}) with the record as argument, and with an optional second argument requesting a
precision in binary digits. }

 
\section{\textcolor{Chapter }{Complex arithmetic}}\logpage{[ 19, 4, 0 ]}
\hyperdef{L}{X7E8F6EFB87A65F78}{}
{
 Complex arithmetic may be implemented in packages, and is present in \textsf{float}. Complex numbers are treated as usual numbers; they may be input with an
extra "i" as in \texttt{\texttt{\symbol{45}}0.5+0.866i}. They may also be created using \texttt{NewFloat} (\ref{NewFloat}) with three arguments: the float filter, the real part, and the imaginary part. 

 Methods should then be implemented for \texttt{Norm}, \texttt{RealPart}, \texttt{ImaginaryPart}, \texttt{ComplexConjugate}, ... 

\subsection{\textcolor{Chapter }{Argument (for complex floats)}}
\logpage{[ 19, 4, 1 ]}\nobreak
\hyperdef{L}{X7B0269D983F96677}{}
{\noindent\textcolor{FuncColor}{$\triangleright$\enspace\texttt{Argument({\mdseries\slshape z})\index{Argument@\texttt{Argument}!for complex floats}
\label{Argument:for complex floats}
}\hfill{\scriptsize (attribute)}}\\


 Returns the argument of the complex number \mbox{\texttt{\mdseries\slshape z}}, namely the value \texttt{Atan2(ImaginaryPart(z),RealPart(z))}. }

 }

 
\section{\textcolor{Chapter }{Interval\texttt{\symbol{45}}specific methods}}\logpage{[ 19, 5, 0 ]}
\hyperdef{L}{X7E57B09C80136484}{}
{
 Interval arithmetic may also be implemented in packages. Intervals are in fact
efficient implementations of sets of real numbers. The only
non\texttt{\symbol{45}}trivial issue is how they should be compared. The
standard \texttt{EQ} tests if the intervals are equal; however, it is usually more useful to know
if intervals overlap, or are disjoint, or are contained in each other. 

 Note the usual convention that intervals are compared as in $[a,b]\leq[c,d]$ if and only if $a\leq c$ and $b\leq d$. 

\subsection{\textcolor{Chapter }{Sup}}
\logpage{[ 19, 5, 1 ]}\nobreak
\hyperdef{L}{X7C34D1D185802F2F}{}
{\noindent\textcolor{FuncColor}{$\triangleright$\enspace\texttt{Sup({\mdseries\slshape x})\index{Sup@\texttt{Sup}}
\label{Sup}
}\hfill{\scriptsize (attribute)}}\\


 Returns the supremum of the interval \mbox{\texttt{\mdseries\slshape x}}. }

 

\subsection{\textcolor{Chapter }{Inf}}
\logpage{[ 19, 5, 2 ]}\nobreak
\hyperdef{L}{X78F1E457814FD1FD}{}
{\noindent\textcolor{FuncColor}{$\triangleright$\enspace\texttt{Inf({\mdseries\slshape x})\index{Inf@\texttt{Inf}}
\label{Inf}
}\hfill{\scriptsize (attribute)}}\\


 Returns the infimum of the interval \mbox{\texttt{\mdseries\slshape x}}. }

 

\subsection{\textcolor{Chapter }{Mid}}
\logpage{[ 19, 5, 3 ]}\nobreak
\hyperdef{L}{X829581A485F55996}{}
{\noindent\textcolor{FuncColor}{$\triangleright$\enspace\texttt{Mid({\mdseries\slshape x})\index{Mid@\texttt{Mid}}
\label{Mid}
}\hfill{\scriptsize (attribute)}}\\


 Returns the midpoint of the interval \mbox{\texttt{\mdseries\slshape x}}. }

 

\subsection{\textcolor{Chapter }{AbsoluteDiameter}}
\logpage{[ 19, 5, 4 ]}\nobreak
\hyperdef{L}{X7FE540B387B0012C}{}
{\noindent\textcolor{FuncColor}{$\triangleright$\enspace\texttt{AbsoluteDiameter({\mdseries\slshape x})\index{AbsoluteDiameter@\texttt{AbsoluteDiameter}}
\label{AbsoluteDiameter}
}\hfill{\scriptsize (attribute)}}\\
\noindent\textcolor{FuncColor}{$\triangleright$\enspace\texttt{Diameter({\mdseries\slshape x})\index{Diameter@\texttt{Diameter}}
\label{Diameter}
}\hfill{\scriptsize (operation)}}\\


 Returns the absolute diameter of the interval \mbox{\texttt{\mdseries\slshape x}}, namely the difference \texttt{Sup(x)\texttt{\symbol{45}}Inf(x)}. }

 

\subsection{\textcolor{Chapter }{RelativeDiameter}}
\logpage{[ 19, 5, 5 ]}\nobreak
\hyperdef{L}{X7CA771757F441592}{}
{\noindent\textcolor{FuncColor}{$\triangleright$\enspace\texttt{RelativeDiameter({\mdseries\slshape x})\index{RelativeDiameter@\texttt{RelativeDiameter}}
\label{RelativeDiameter}
}\hfill{\scriptsize (attribute)}}\\


 Returns the relative diameter of the interval \mbox{\texttt{\mdseries\slshape x}}, namely \texttt{(Sup(x)\texttt{\symbol{45}}Inf(x))/AbsoluteValue(Min(x))}. }

 

\subsection{\textcolor{Chapter }{IsDisjoint}}
\logpage{[ 19, 5, 6 ]}\nobreak
\hyperdef{L}{X86D22AE57E2D84B2}{}
{\noindent\textcolor{FuncColor}{$\triangleright$\enspace\texttt{IsDisjoint({\mdseries\slshape x1, x2})\index{IsDisjoint@\texttt{IsDisjoint}}
\label{IsDisjoint}
}\hfill{\scriptsize (operation)}}\\


 Returns \texttt{true} if the two intervals \mbox{\texttt{\mdseries\slshape x1}}, \mbox{\texttt{\mdseries\slshape x2}} are disjoint. }

 

\subsection{\textcolor{Chapter }{IsSubset (for interval floats)}}
\logpage{[ 19, 5, 7 ]}\nobreak
\hyperdef{L}{X7A5E0C3E79837EB8}{}
{\noindent\textcolor{FuncColor}{$\triangleright$\enspace\texttt{IsSubset({\mdseries\slshape x1, x2})\index{IsSubset@\texttt{IsSubset}!for interval floats}
\label{IsSubset:for interval floats}
}\hfill{\scriptsize (operation)}}\\


 Returns \texttt{true} if the interval \mbox{\texttt{\mdseries\slshape x1}} contains \mbox{\texttt{\mdseries\slshape x2}}. }

 

\subsection{\textcolor{Chapter }{IncreaseInterval}}
\logpage{[ 19, 5, 8 ]}\nobreak
\hyperdef{L}{X85191E1679936CE9}{}
{\noindent\textcolor{FuncColor}{$\triangleright$\enspace\texttt{IncreaseInterval({\mdseries\slshape x, delta})\index{IncreaseInterval@\texttt{IncreaseInterval}}
\label{IncreaseInterval}
}\hfill{\scriptsize (operation)}}\\


 Returns an interval with same midpoint as \mbox{\texttt{\mdseries\slshape x}} but absolute diameter increased by \mbox{\texttt{\mdseries\slshape delta}}. }

 

\subsection{\textcolor{Chapter }{BlowupInterval}}
\logpage{[ 19, 5, 9 ]}\nobreak
\hyperdef{L}{X879EE14282DD1539}{}
{\noindent\textcolor{FuncColor}{$\triangleright$\enspace\texttt{BlowupInterval({\mdseries\slshape x, ratio})\index{BlowupInterval@\texttt{BlowupInterval}}
\label{BlowupInterval}
}\hfill{\scriptsize (operation)}}\\


 Returns an interval with same midpoint as \mbox{\texttt{\mdseries\slshape x}} but relative diameter increased by \mbox{\texttt{\mdseries\slshape ratio}}. }

 

\subsection{\textcolor{Chapter }{BisectInterval}}
\logpage{[ 19, 5, 10 ]}\nobreak
\hyperdef{L}{X7EC15DAE7CBBB42E}{}
{\noindent\textcolor{FuncColor}{$\triangleright$\enspace\texttt{BisectInterval({\mdseries\slshape x})\index{BisectInterval@\texttt{BisectInterval}}
\label{BisectInterval}
}\hfill{\scriptsize (operation)}}\\


 Returns a list of two intervals whose union equals the interval \mbox{\texttt{\mdseries\slshape x}}. }

 }

 }

          
\chapter{\textcolor{Chapter }{Booleans}}\label{Booleans}
\logpage{[ 20, 0, 0 ]}
\hyperdef{L}{X787B4AB77A2F5E14}{}
{
  \index{type!boolean} \index{logical} The two main \emph{boolean} values are \texttt{true} and \texttt{false}. They stand for the \emph{logical} values of the same name. They appear as values of the conditions in \texttt{if}\texttt{\symbol{45}}statements and \texttt{while}\texttt{\symbol{45}}loops. Booleans are also important as return values of \emph{filters} (see{\nobreakspace}\ref{Filters}) such as \texttt{IsFinite} (\ref{IsFinite}) and \texttt{IsBool} (\ref{IsBool}). Note that it is a convention that the name of a function that returns \texttt{true} or \texttt{false} according to the outcome, starts with \texttt{Is}. 

 For technical reasons, also the value \texttt{fail} (see{\nobreakspace}\ref{Fail}) is regarded as a boolean.  
\section{\textcolor{Chapter }{IsBool (Filter)}}\label{sect:IsBool}
\logpage{[ 20, 1, 0 ]}
\hyperdef{L}{X87F9AF65832E7AD2}{}
{
  

\subsection{\textcolor{Chapter }{IsBool}}
\logpage{[ 20, 1, 1 ]}\nobreak
\hyperdef{L}{X7D58580284CF7894}{}
{\noindent\textcolor{FuncColor}{$\triangleright$\enspace\texttt{IsBool({\mdseries\slshape obj})\index{IsBool@\texttt{IsBool}}
\label{IsBool}
}\hfill{\scriptsize (Category)}}\\


 tests whether \mbox{\texttt{\mdseries\slshape obj}} is \texttt{true}, \texttt{false} or \texttt{fail}. 
\begin{Verbatim}[commandchars=!@|,fontsize=\small,frame=single,label=Example]
  !gapprompt@gap>| !gapinput@IsBool( true );  IsBool( false );  IsBool( 17 );|
  true
  true
  false
\end{Verbatim}
 }

 }

  
\section{\textcolor{Chapter }{Fail (Variable)}}\label{Fail}
\logpage{[ 20, 2, 0 ]}
\hyperdef{L}{X85E648AA8414F303}{}
{
  

\subsection{\textcolor{Chapter }{fail}}
\logpage{[ 20, 2, 1 ]}\nobreak
\hyperdef{L}{X8294AAC9860E87E5}{}
{\noindent\textcolor{FuncColor}{$\triangleright$\enspace\texttt{fail\index{fail@\texttt{fail}}
\label{fail}
}\hfill{\scriptsize (global variable)}}\\


 The value \texttt{fail} is used to indicate situations when an operation could not be performed for
the given arguments, either because of shortcomings of the arguments or
because of restrictions in the implementation or computability. So for example \texttt{Position} (\ref{Position}) will return \texttt{fail} if the point searched for is not in the list. 

 \texttt{fail} is simply an object that is different from every other object than itself. 

 For technical reasons, \texttt{fail} is a boolean value. But note that \texttt{fail} cannot be used to form boolean expressions with \texttt{and}, \texttt{or}, and \texttt{not} (see{\nobreakspace}\ref{Operations for Booleans} below), and \texttt{fail} cannot appear in boolean lists (see Chapter{\nobreakspace}\ref{Boolean Lists}). }

 }

  
\section{\textcolor{Chapter }{Comparisons of Booleans}}\label{Comparisons of Booleans}
\logpage{[ 20, 3, 0 ]}
\hyperdef{L}{X862F17B68465B399}{}
{
  \index{comparisons!of booleans} 
\subsection{\textcolor{Chapter }{Equality and inequality of Booleans}}\label{Equality and inequality of Booleans}
\logpage{[ 20, 3, 1 ]}
\hyperdef{L}{X79305F9780394190}{}
{
  \index{equality!of booleans} \index{inequality!of booleans} \texttt{\mbox{\texttt{\mdseries\slshape bool1}} = \mbox{\texttt{\mdseries\slshape bool2}}} 

 \noindent \texttt{\mbox{\texttt{\mdseries\slshape bool1}} {\textless}{\textgreater} \mbox{\texttt{\mdseries\slshape bool2}}} 

 The equality operator \texttt{=} evaluates to \texttt{true} if the two boolean values \mbox{\texttt{\mdseries\slshape bool1}} and \mbox{\texttt{\mdseries\slshape bool2}} are equal, i.e., both are \texttt{true} or both are \texttt{false} or both \texttt{fail}, and \texttt{false} otherwise. The inequality operator \texttt{{\textless}{\textgreater}} evaluates to \texttt{true} if the two boolean values \mbox{\texttt{\mdseries\slshape bool1}}, \mbox{\texttt{\mdseries\slshape bool2}} are different, and \texttt{false} otherwise. This operation is also called the \emph{exclusive or}, because its value is \texttt{true} if exactly one of \mbox{\texttt{\mdseries\slshape bool1}} or \mbox{\texttt{\mdseries\slshape bool2}} is \texttt{true}. 

 You can compare boolean values with objects of other types. Of course they are
never equal. 

 
\begin{Verbatim}[commandchars=!@|,fontsize=\small,frame=single,label=Example]
  !gapprompt@gap>| !gapinput@true = false;|
  false
  !gapprompt@gap>| !gapinput@false = (true = fail);|
  true
  !gapprompt@gap>| !gapinput@true <> 17;|
  true
\end{Verbatim}
 }

 
\subsection{\textcolor{Chapter }{Ordering of Booleans}}\label{Ordering of Booleans}
\logpage{[ 20, 3, 2 ]}
\hyperdef{L}{X7FEF019482AF5923}{}
{
  \index{ordering!booleans} \mbox{\texttt{\mdseries\slshape bool1}} \texttt{{\textless}} \mbox{\texttt{\mdseries\slshape bool2}} 

 The ordering of boolean values is defined by \texttt{true} \texttt{{\textless}} \texttt{false} \texttt{{\textless}} \texttt{fail}. For the comparison of booleans with other \textsf{GAP} objects, see Section{\nobreakspace}\ref{Comparisons}. 

 
\begin{Verbatim}[commandchars=!@|,fontsize=\small,frame=single,label=Example]
  !gapprompt@gap>| !gapinput@true < false;  fail >= false;|
  true
  true
\end{Verbatim}
 }

 }

  
\section{\textcolor{Chapter }{Operations for Booleans}}\label{Operations for Booleans}
\logpage{[ 20, 4, 0 ]}
\hyperdef{L}{X79AD41A185FD7213}{}
{
  \index{operations!for booleans} \index{logical operations} The following boolean operations are only applicable to \texttt{true} and \texttt{false}. 

 
\subsection{\textcolor{Chapter }{Logical disjunction}}\label{Logical disjunction}
\logpage{[ 20, 4, 1 ]}
\hyperdef{L}{X7DFE7E518088AA89}{}
{
  \index{Logical disjunction} \index{or@\texttt{or}} \mbox{\texttt{\mdseries\slshape bool1}} \texttt{or} \mbox{\texttt{\mdseries\slshape bool2}} 

 The logical operator \texttt{or} evaluates to \texttt{true} if at least one of the two boolean operands \mbox{\texttt{\mdseries\slshape bool1}} and \mbox{\texttt{\mdseries\slshape bool2}} is \texttt{true}, and to \texttt{false} otherwise. 

 \texttt{or} first evaluates \mbox{\texttt{\mdseries\slshape bool1}}. If the value is neither \texttt{true} nor \texttt{false} an error is signalled. If the value is \texttt{true}, then \texttt{or} returns \texttt{true} \emph{without} evaluating \mbox{\texttt{\mdseries\slshape bool2}}. If the value is \texttt{false}, then \texttt{or} evaluates \mbox{\texttt{\mdseries\slshape bool2}}. Again, if the value is neither \texttt{true} nor \texttt{false} an error is signalled. Otherwise \texttt{or} returns the value of \mbox{\texttt{\mdseries\slshape bool2}}. This \emph{short\texttt{\symbol{45}}circuited} evaluation is important if the value of \mbox{\texttt{\mdseries\slshape bool1}} is \texttt{true} and evaluation of \mbox{\texttt{\mdseries\slshape bool2}} would take much time or cause an error. 

 \texttt{or} is associative, i.e., it is allowed to write \mbox{\texttt{\mdseries\slshape b1}} \texttt{or} \mbox{\texttt{\mdseries\slshape b2}} \texttt{or} \mbox{\texttt{\mdseries\slshape b3}}, which is interpreted as (\mbox{\texttt{\mdseries\slshape b1}} \texttt{or} \mbox{\texttt{\mdseries\slshape b2}}) \texttt{or} \mbox{\texttt{\mdseries\slshape b3}}. \texttt{or} has the lowest precedence of the logical operators. All logical operators have
lower precedence than the comparison operators \texttt{=}, \texttt{{\textless}}, \texttt{in}, etc. 

 
\begin{Verbatim}[commandchars=!@|,fontsize=\small,frame=single,label=Example]
  !gapprompt@gap>| !gapinput@true or false;|
  true
  !gapprompt@gap>| !gapinput@false or false;|
  false
  !gapprompt@gap>| !gapinput@i := -1;; l := [1,2,3];;|
  !gapprompt@gap>| !gapinput@if i <= 0 or l[i] = false then     # this does not cause an error,|
  !gapprompt@>| !gapinput@   Print("aha\n"); fi;                # because `l[i]' is not evaluated|
  aha
\end{Verbatim}
 }

 
\subsection{\textcolor{Chapter }{Logical conjunction}}\label{Logical conjunction}
\logpage{[ 20, 4, 2 ]}
\hyperdef{L}{X7A64D25F804973CD}{}
{
  \index{Logical conjunction} \index{and@\texttt{and}} \mbox{\texttt{\mdseries\slshape bool1}} \texttt{and} \mbox{\texttt{\mdseries\slshape bool2}} 

 \index{and@\texttt{and}!for filters} \noindent \mbox{\texttt{\mdseries\slshape fil1}} \texttt{and} \mbox{\texttt{\mdseries\slshape fil2}} 

 The logical operator \texttt{and} evaluates to \texttt{true} if both boolean operands \mbox{\texttt{\mdseries\slshape bool1}}, \mbox{\texttt{\mdseries\slshape bool2}} are \texttt{true}, and to \texttt{false} otherwise. 

 \texttt{and} first evaluates \mbox{\texttt{\mdseries\slshape bool1}}. If the value is neither \texttt{true} nor \texttt{false} an error is signalled. If the value is \texttt{false}, then \texttt{and} returns \texttt{false} \emph{without} evaluating \mbox{\texttt{\mdseries\slshape bool2}}. If the value is \texttt{true}, then \texttt{and} evaluates \mbox{\texttt{\mdseries\slshape bool2}}. Again, if the value is neither \texttt{true} nor \texttt{false} an error is signalled. Otherwise \texttt{and} returns the value of \mbox{\texttt{\mdseries\slshape bool2}}. This \emph{short\texttt{\symbol{45}}circuited} evaluation is important if the value of \mbox{\texttt{\mdseries\slshape bool1}} is \texttt{false} and evaluation of \mbox{\texttt{\mdseries\slshape bool2}} would take much time or cause an error. 

 \texttt{and} is associative, i.e., it is allowed to write \mbox{\texttt{\mdseries\slshape b1}} \texttt{and} \mbox{\texttt{\mdseries\slshape b2}} \texttt{and} \mbox{\texttt{\mdseries\slshape b3}}, which is interpreted as (\mbox{\texttt{\mdseries\slshape b1}} \texttt{and} \mbox{\texttt{\mdseries\slshape b2}}) \texttt{and} \mbox{\texttt{\mdseries\slshape b3}}. \texttt{and} has higher precedence than the logical \texttt{or} operator, but lower than the unary logical \texttt{not} operator. All logical operators have lower precedence than the comparison
operators \texttt{=}, \texttt{{\textless}}, \texttt{in}, etc. 

 
\begin{Verbatim}[commandchars=!@|,fontsize=\small,frame=single,label=Example]
  !gapprompt@gap>| !gapinput@true and false;|
  false
  !gapprompt@gap>| !gapinput@true and true;|
  true
  !gapprompt@gap>| !gapinput@false and 17;  # does not cause error, because 17 is never looked at|
  false
\end{Verbatim}
 

 \texttt{and} can also be applied to filters. It returns a filter that when applied to some
argument \mbox{\texttt{\mdseries\slshape x}}, tests \mbox{\texttt{\mdseries\slshape fil1}}$(x)$ \texttt{and} \mbox{\texttt{\mdseries\slshape fil2}}$(x)$. 

 
\begin{Verbatim}[commandchars=!@|,fontsize=\small,frame=single,label=Example]
  !gapprompt@gap>| !gapinput@andfilt:= IsPosRat and IsInt;;|
  !gapprompt@gap>| !gapinput@andfilt( 17 );  andfilt( 1/2 );|
  true
  false
\end{Verbatim}
 }

 
\subsection{\textcolor{Chapter }{Logical negation}}\label{Logical negation}
\logpage{[ 20, 4, 3 ]}
\hyperdef{L}{X84F5034185D7EC3C}{}
{
  \index{Logical negation} \index{not@\texttt{not}} \texttt{not} \mbox{\texttt{\mdseries\slshape bool}} 

 The logical operator \texttt{not} returns \texttt{true} if the boolean value \mbox{\texttt{\mdseries\slshape bool}} is \texttt{false}, and \texttt{true} otherwise. An error is signalled if \mbox{\texttt{\mdseries\slshape bool}} does not evaluate to \texttt{true} or \texttt{false}. 

 \texttt{not} has higher precedence than the other logical operators, \texttt{or} and \texttt{and}. All logical operators have lower precedence than the comparison operators \texttt{=}, \texttt{{\textless}}, \texttt{in}, etc. 

 
\begin{Verbatim}[commandchars=!@|,fontsize=\small,frame=single,label=Example]
  !gapprompt@gap>| !gapinput@true and false;|
  false
  !gapprompt@gap>| !gapinput@not true;|
  false
  !gapprompt@gap>| !gapinput@not false;|
  true
\end{Verbatim}
 }

 }

 }

             
\chapter{\textcolor{Chapter }{Lists}}\label{Lists}
\logpage{[ 21, 0, 0 ]}
\hyperdef{L}{X7B256AE5780F140A}{}
{
  Lists are the most important way to treat objects together. A \emph{list} arranges objects in a definite order. So each list implies a partial mapping
from the integers to the elements of the list. I.e., there is a first element
of a list, a second, a third, and so on. Lists can occur in mutable or
immutable form, see{\nobreakspace}\ref{Mutability and Copyability} for the concept of mutability, and{\nobreakspace}\ref{Duplication of Lists} for the case of lists. 

 This chapter deals mainly with the aspect of lists in \textsf{GAP} as \emph{data structures}. Chapter{\nobreakspace}\ref{Collections} tells more about the \emph{collection} aspect of certain lists, and more about lists as \emph{arithmetic objects} can be found in the chapters \ref{Row Vectors} and \ref{Matrices}. 

 Lists are used to implement ranges (see{\nobreakspace}\ref{Ranges}), sets (see{\nobreakspace}\ref{Sorted Lists and Sets}),\index{Sets} strings (see{\nobreakspace}\ref{Strings and Characters}), row vectors and matrices (see{\nobreakspace}\ref{Row Vectors} and \ref{Matrices}, but note that \textsf{GAP} supports also linear algebra for objects which are \emph{not} lists, see \ref{Vector and Matrix Objects}); boolean lists (see{\nobreakspace}\ref{Boolean Lists}) are a further special kind of lists. 

 Several operations for lists, such as \texttt{Intersection} (\ref{Intersection:for a list}) and \texttt{Random} (\ref{Random:for a list or collection}), will be described in Chapter{\nobreakspace}\ref{Collections}, in particular see{\nobreakspace}\ref{Lists and Collections}.  
\section{\textcolor{Chapter }{List Categories}}\label{List Categories}
\logpage{[ 21, 1, 0 ]}
\hyperdef{L}{X86B28F5B781FFD31}{}
{
  A list can be written by writing down the elements in order between square
brackets \texttt{[}, \texttt{]}, and separating them with commas \texttt{,}. An \emph{empty list}, i.e., a list with no elements, is written as \texttt{[]}. 

 
\begin{Verbatim}[commandchars=!@|,fontsize=\small,frame=single,label=Example]
  !gapprompt@gap>| !gapinput@[ 1, 2, 3 ];              # a list with three elements|
  [ 1, 2, 3 ]
  !gapprompt@gap>| !gapinput@[ [], [ 1 ], [ 1, 2 ] ];  # a list may contain other lists|
  [ [  ], [ 1 ], [ 1, 2 ] ]
\end{Verbatim}
 

 Each list constructed this way is mutable (see{\nobreakspace}\ref{Mutability and Copyability}). 

\subsection{\textcolor{Chapter }{IsList}}
\logpage{[ 21, 1, 1 ]}\nobreak
\hyperdef{L}{X7C4CC4EA8299701E}{}
{\noindent\textcolor{FuncColor}{$\triangleright$\enspace\texttt{IsList({\mdseries\slshape obj})\index{IsList@\texttt{IsList}}
\label{IsList}
}\hfill{\scriptsize (Category)}}\\


 tests whether \mbox{\texttt{\mdseries\slshape obj}} is a list. 

 
\begin{Verbatim}[commandchars=!@|,fontsize=\small,frame=single,label=Example]
  !gapprompt@gap>| !gapinput@IsList( [ 1, 3, 5, 7 ] );  IsList( 1 );|
  true
  false
\end{Verbatim}
 }

 

\subsection{\textcolor{Chapter }{IsDenseList}}
\logpage{[ 21, 1, 2 ]}\nobreak
\hyperdef{L}{X870AA9D8798C93DD}{}
{\noindent\textcolor{FuncColor}{$\triangleright$\enspace\texttt{IsDenseList({\mdseries\slshape obj})\index{IsDenseList@\texttt{IsDenseList}}
\label{IsDenseList}
}\hfill{\scriptsize (Category)}}\\


 A list is \emph{dense} if it has no holes, i.e., contains an element at every position up to the
length. It is absolutely legal to have lists with holes. They are created by
leaving the entry between the commas empty. Holes at the end of a list are
ignored. Lists with holes are sometimes convenient when the list represents a
mapping from a finite, but not consecutive, subset of the positive integers. 
\begin{Verbatim}[commandchars=!@|,fontsize=\small,frame=single,label=Example]
  !gapprompt@gap>| !gapinput@IsDenseList( [ 1, 2, 3 ] );|
  true
  !gapprompt@gap>| !gapinput@l := [ , 4, 9,, 25,, 49,,,, 121 ];;  IsDenseList( l );|
  false
  !gapprompt@gap>| !gapinput@l[3];|
  9
  !gapprompt@gap>| !gapinput@l[4];|
  Error, List Element: <list>[4] must have an assigned value
\end{Verbatim}
 

 Observe that requesting the value of \texttt{l[4]}, which was not assigned, caused the entry of a \texttt{break}\texttt{\symbol{45}}loop (see Section{\nobreakspace}\ref{Break Loops}). }

 

\subsection{\textcolor{Chapter }{IsHomogeneousList}}
\logpage{[ 21, 1, 3 ]}\nobreak
\hyperdef{L}{X7C71596C82B6EF35}{}
{\noindent\textcolor{FuncColor}{$\triangleright$\enspace\texttt{IsHomogeneousList({\mdseries\slshape obj})\index{IsHomogeneousList@\texttt{IsHomogeneousList}}
\label{IsHomogeneousList}
}\hfill{\scriptsize (Category)}}\\


 returns \texttt{true} if \mbox{\texttt{\mdseries\slshape obj}} is a list and it is homogeneous, and \texttt{false} otherwise. 

 A \emph{homogeneous} list is a dense list whose elements lie in the same family (see{\nobreakspace}\ref{Families}). The empty list is homogeneous but not a collection (see{\nobreakspace}\ref{Collections}), a nonempty homogeneous list is also a collection.  
\begin{Verbatim}[commandchars=!@|,fontsize=\small,frame=single,label=Example]
  !gapprompt@gap>| !gapinput@IsHomogeneousList( [ 1, 2, 3 ] );  IsHomogeneousList( [] );|
  true
  true
  !gapprompt@gap>| !gapinput@IsHomogeneousList( [ 1, false, () ] );|
  false
\end{Verbatim}
 }

 

\subsection{\textcolor{Chapter }{IsTable}}
\logpage{[ 21, 1, 4 ]}\nobreak
\hyperdef{L}{X80872FAF80EB5DF9}{}
{\noindent\textcolor{FuncColor}{$\triangleright$\enspace\texttt{IsTable({\mdseries\slshape obj})\index{IsTable@\texttt{IsTable}}
\label{IsTable}
}\hfill{\scriptsize (Category)}}\\


 A \emph{table} is a nonempty list of homogeneous lists which lie in the same family. Typical
examples of tables are matrices (see{\nobreakspace}\ref{Matrices}). 
\begin{Verbatim}[commandchars=!@|,fontsize=\small,frame=single,label=Example]
  !gapprompt@gap>| !gapinput@IsTable( [ [ 1, 2 ], [ 3, 4 ] ] );    # in fact a matrix|
  true
  !gapprompt@gap>| !gapinput@IsTable( [ [ 1 ], [ 2, 3 ] ] );       # not rectangular but a table|
  true
  !gapprompt@gap>| !gapinput@IsTable( [ [ 1, 2 ], [ () , (1,2) ] ] );  # not homogeneous|
  false
\end{Verbatim}
 }

 

\subsection{\textcolor{Chapter }{IsRectangularTable}}
\logpage{[ 21, 1, 5 ]}\nobreak
\hyperdef{L}{X79581E0387F7F7A9}{}
{\noindent\textcolor{FuncColor}{$\triangleright$\enspace\texttt{IsRectangularTable({\mdseries\slshape list})\index{IsRectangularTable@\texttt{IsRectangularTable}}
\label{IsRectangularTable}
}\hfill{\scriptsize (property)}}\\


 A list lies in \texttt{IsRectangularTable} when it is nonempty and its elements are all homogeneous lists of the same
family and the same length. }

 

\subsection{\textcolor{Chapter }{IsConstantTimeAccessList}}
\logpage{[ 21, 1, 6 ]}\nobreak
\hyperdef{L}{X7C84E16A85C99C8C}{}
{\noindent\textcolor{FuncColor}{$\triangleright$\enspace\texttt{IsConstantTimeAccessList({\mdseries\slshape list})\index{IsConstantTimeAccessList@\texttt{IsConstantTimeAccessList}}
\label{IsConstantTimeAccessList}
}\hfill{\scriptsize (Category)}}\\


 This category indicates whether the access to each element of the list \mbox{\texttt{\mdseries\slshape list}} will take roughly the same time. This is implied for example by \texttt{IsList and IsInternalRep}, so all strings, Boolean lists, ranges, and internally represented plain
lists are in this category. 

 But also other enumerators (see{\nobreakspace}\ref{Enumerators}) can lie in this category if they guarantee constant time access to their
elements. }

 }

  
\section{\textcolor{Chapter }{Basic Operations for Lists}}\label{Basic Operations for Lists}
\logpage{[ 21, 2, 0 ]}
\hyperdef{L}{X7B202D147A5C2884}{}
{
  The basic operations for lists are element access (see{\nobreakspace}\ref{List Elements}), assignment of elements to a list (see{\nobreakspace}\ref{List Assignment}), fetching the length of a list (see{\nobreakspace}\texttt{Length} (\ref{Length})), the test for a hole at a given position, and unbinding an element at a
given position (see{\nobreakspace}\ref{IsBound and Unbind for Lists}). 

 The term basic operation means that each other list operation can be
formulated in terms of the basic operations. (But note that often a more
efficient method than this one is implemented.) 

 Any \textsf{GAP} object \mbox{\texttt{\mdseries\slshape list}} in the category \texttt{IsList} (\ref{IsList}) is regarded as a list, and if methods for the basic list operations are
installed for \mbox{\texttt{\mdseries\slshape list}} then \mbox{\texttt{\mdseries\slshape list}} can be used also for the other list operations. 

 For internally represented lists, kernel methods are provided for the basic
list operations with positive integer indices. For other lists or other
indices, it is possible to install appropriate methods for these operations.
This permits the implementation of lists that do not need to store all list
elements (see also{\nobreakspace}\ref{Enumerators}); for example, the elements might be described by an algorithm, such as the
elements list of a group. For this reduction of space requirements, however, a
price in access time may have to be paid (see{\nobreakspace}\texttt{ConstantTimeAccessList} (\ref{ConstantTimeAccessList})). 

 \index{list element!operation} \index{list boundedness test!operation} \index{list assignment!operation} \index{list unbind!operation} 

\subsection{\textcolor{Chapter }{\texttt{\symbol{92}}[\texttt{\symbol{92}}]}}
\logpage{[ 21, 2, 1 ]}\nobreak
\hyperdef{L}{X8297BBCD79642BE6}{}
{\noindent\textcolor{FuncColor}{$\triangleright$\enspace\texttt{\texttt{\symbol{92}}[\texttt{\symbol{92}}]({\mdseries\slshape list, ix})\index{\texttt{\symbol{92}}[\texttt{\symbol{92}}]@\texttt{\texttt{\symbol{92}}[\texttt{\symbol{92}}]}}
\label{bSlash[bSlash]}
}\hfill{\scriptsize (operation)}}\\
\noindent\textcolor{FuncColor}{$\triangleright$\enspace\texttt{IsBound\texttt{\symbol{92}}[\texttt{\symbol{92}}]({\mdseries\slshape list, ix})\index{IsBound\texttt{\symbol{92}}[\texttt{\symbol{92}}]@\texttt{IsBound\texttt{\symbol{92}}[\texttt{\symbol{92}}]}}
\label{IsBoundbSlash[bSlash]}
}\hfill{\scriptsize (operation)}}\\
\noindent\textcolor{FuncColor}{$\triangleright$\enspace\texttt{\texttt{\symbol{92}}[\texttt{\symbol{92}}]\texttt{\symbol{92}}:\texttt{\symbol{92}}=({\mdseries\slshape list, pos, ix})\index{\texttt{\symbol{92}}[\texttt{\symbol{92}}]\texttt{\symbol{92}}:\texttt{\symbol{92}}=@\texttt{\texttt{\symbol{92}}[\texttt{\symbol{92}}]\texttt{\symbol{92}}:\texttt{\symbol{92}}=}}
\label{bSlash[bSlash]bSlash:bSlash=}
}\hfill{\scriptsize (operation)}}\\
\noindent\textcolor{FuncColor}{$\triangleright$\enspace\texttt{Unbind\texttt{\symbol{92}}[\texttt{\symbol{92}}]({\mdseries\slshape list, ix})\index{Unbind\texttt{\symbol{92}}[\texttt{\symbol{92}}]@\texttt{Unbind\texttt{\symbol{92}}[\texttt{\symbol{92}}]}}
\label{UnbindbSlash[bSlash]}
}\hfill{\scriptsize (operation)}}\\


 These operations implement element access, test for element boundedness, list
element assignment, and removal of the element with index \mbox{\texttt{\mdseries\slshape ix}}. 

 Note that the special characters \texttt{[}, \texttt{]}, \texttt{:}, and \texttt{=} must be escaped with a backslash \texttt{\texttt{\symbol{92}}} (see{\nobreakspace}\ref{Symbols}); so \texttt{\texttt{\symbol{92}}[\texttt{\symbol{92}}]} denotes the operation for element access in a list, whereas \texttt{[]} denotes an empty list. (Maybe the variable names involving special characters
look strange, but nevertheless they are quite suggestive.) 

 \texttt{\texttt{\symbol{92}}[\texttt{\symbol{92}}]( \mbox{\texttt{\mdseries\slshape list}}, \mbox{\texttt{\mdseries\slshape ix}} )} is equivalent to \texttt{\mbox{\texttt{\mdseries\slshape list}}[ \mbox{\texttt{\mdseries\slshape ix}} ]}, which clearly will usually be preferred; the former is useful mainly if one
wants to access the operation itself, for example if one wants to install a
method for element access in a special kind of lists. 

 Similarly, \texttt{IsBound\texttt{\symbol{92}}[\texttt{\symbol{92}}]} is used explicitly mainly in method installations. In other situations, one
can simply call \texttt{IsBound} (\ref{IsBound:for a list index}), which then delegates to \texttt{IsBound\texttt{\symbol{92}}[\texttt{\symbol{92}}]} if the first argument is a list, and to \texttt{IsBound\texttt{\symbol{92}}.} (\ref{IsBoundbSlash.}) if the first argument is a record. 

 Analogous statements hold for \texttt{\texttt{\symbol{92}}[\texttt{\symbol{92}}]\texttt{\symbol{92}}:\texttt{\symbol{92}}=} and \texttt{Unbind\texttt{\symbol{92}}[\texttt{\symbol{92}}]}. }

 }

  
\section{\textcolor{Chapter }{List Elements}}\label{List Elements}
\logpage{[ 21, 3, 0 ]}
\hyperdef{L}{X7921047F83F5FA28}{}
{
  \index{accessing!list elements} \index{list element!access} \texttt{\mbox{\texttt{\mdseries\slshape list}}[ \mbox{\texttt{\mdseries\slshape ix}} ]} 

 The above construct evaluates to the element of the list \mbox{\texttt{\mdseries\slshape list}} with index \mbox{\texttt{\mdseries\slshape ix}}. For built\texttt{\symbol{45}}in list types and collections, indexing is done
with origin 1, i.e., the first element of the list is the element with index
1. 
\begin{Verbatim}[commandchars=!@|,fontsize=\small,frame=single,label=Example]
  !gapprompt@gap>| !gapinput@l := [ 2, 3, 5, 7, 11, 13 ];;  l[1];  l[2];  l[6];|
  2
  3
  13
\end{Verbatim}
 If \mbox{\texttt{\mdseries\slshape list}} is not a built\texttt{\symbol{45}}in list, or \mbox{\texttt{\mdseries\slshape ix}} does not evaluate to a positive integer, method selection is invoked to try
and find a way of indexing \mbox{\texttt{\mdseries\slshape list}} with index \mbox{\texttt{\mdseries\slshape ix}}. If this fails, or the selected method finds that \texttt{\mbox{\texttt{\mdseries\slshape list}}[\mbox{\texttt{\mdseries\slshape ix}}]} is unbound, an error is signalled. 

 \index{sublist} \index{sublist!access} \texttt{\mbox{\texttt{\mdseries\slshape list}}\texttt{\symbol{123}} \mbox{\texttt{\mdseries\slshape poss}} \texttt{\symbol{125}}} 

 The above construct evaluates to a new list \mbox{\texttt{\mdseries\slshape new}} whose first element is \texttt{\mbox{\texttt{\mdseries\slshape list}}[\mbox{\texttt{\mdseries\slshape poss}}[1]]}, whose second element is \texttt{\mbox{\texttt{\mdseries\slshape list}}[\mbox{\texttt{\mdseries\slshape poss}}[2]]}, and so on. However, it does not need to be sorted and may contain duplicate
elements. If for any $i$, \texttt{\mbox{\texttt{\mdseries\slshape list}}[ \mbox{\texttt{\mdseries\slshape poss}}[}$i$\texttt{] ]} is unbound, an error is signalled. 

 
\begin{Verbatim}[commandchars=!@|,fontsize=\small,frame=single,label=Example]
  !gapprompt@gap>| !gapinput@l := [ 2, 3, 5, 7, 11, 13, 17, 19 ];;|
  !gapprompt@gap>| !gapinput@l{[4..6]};  l{[1,7,1,8]};|
  [ 7, 11, 13 ]
  [ 2, 17, 2, 19 ]
\end{Verbatim}
 

 The result is a \emph{new} list, that is not identical to any other list. The elements of that list,
however, are identical to the corresponding elements of the left operand
(see{\nobreakspace}\ref{Identical Lists}). 

 It is possible to nest such \emph{sublist extractions}, as can be seen in the example below. 

 
\begin{Verbatim}[commandchars=!@|,fontsize=\small,frame=single,label=Example]
  !gapprompt@gap>| !gapinput@m := [ [1,2,3], [4,5,6], [7,8,9], [10,11,12] ];;  m{[1,2,3]}{[3,2]};|
  [ [ 3, 2 ], [ 6, 5 ], [ 9, 8 ] ]
  !gapprompt@gap>| !gapinput@l := m{[1,2,3]};; l{[3,2]};|
  [ [ 7, 8, 9 ], [ 4, 5, 6 ] ]
\end{Verbatim}
 

 Note the difference between the two examples. The latter extracts elements 1,
2, and 3 from \mbox{\texttt{\mdseries\slshape m}} and then extracts the elements 3 and 2 from \emph{this list}. The former extracts elements 1, 2, and 3 from \mbox{\texttt{\mdseries\slshape m}} and then extracts the elements 3 and 2 from \emph{each of those element lists}. 

 To be precise: With each selector \texttt{[\mbox{\texttt{\mdseries\slshape pos}}]} or \texttt{\texttt{\symbol{123}}\mbox{\texttt{\mdseries\slshape poss}}\texttt{\symbol{125}}} we associate a \emph{level} that is defined as the number of selectors of the form \texttt{\texttt{\symbol{123}}\mbox{\texttt{\mdseries\slshape poss}}\texttt{\symbol{125}}} to its left in the same expression. For example 

 
\begin{Verbatim}[commandchars=!@|,fontsize=\small,frame=single,label=]
      l[pos1]{poss2}{poss3}[pos4]{poss5}[pos6]
  level   0      0      1     2      2     3
\end{Verbatim}
 

 Then a selector \texttt{\mbox{\texttt{\mdseries\slshape list}}[\mbox{\texttt{\mdseries\slshape pos}}]} of level \mbox{\texttt{\mdseries\slshape level}} is computed as \texttt{ListElement(\mbox{\texttt{\mdseries\slshape list}},\mbox{\texttt{\mdseries\slshape pos}},\mbox{\texttt{\mdseries\slshape level}})}, where \texttt{ListElement} is defined as follows. (Note that \texttt{ListElement} is \emph{not} a \textsf{GAP} function.) 

 
\begin{Verbatim}[commandchars=!@|,fontsize=\small,frame=single,label=Example]
  ListElement := function ( list, pos, level )
   if level = 0 then
    return list[pos];
   else
    return List( list, elm -> ListElement(elm,pos,level-1) );
   fi;
  end;
\end{Verbatim}
 

 and a selector \texttt{\mbox{\texttt{\mdseries\slshape list}}\texttt{\symbol{123}}\mbox{\texttt{\mdseries\slshape poss}}\texttt{\symbol{125}}} of level \mbox{\texttt{\mdseries\slshape level}} is computed as \texttt{ListElements(\mbox{\texttt{\mdseries\slshape list}},\mbox{\texttt{\mdseries\slshape poss}},\mbox{\texttt{\mdseries\slshape level}})}, where \texttt{ListElements} is defined as follows. (Note that \texttt{ListElements} is \emph{not} a \textsf{GAP} function.) 

 
\begin{Verbatim}[commandchars=!@|,fontsize=\small,frame=single,label=Example]
  ListElements := function ( list, poss, level )
   if level = 0 then
    return list{poss};
    else
     return List( list, elm -> ListElements(elm,poss,level-1) );
    fi;
  end;
\end{Verbatim}
 

 \index{sublist!operation} 

\subsection{\textcolor{Chapter }{\texttt{\symbol{92}}\texttt{\symbol{123}}\texttt{\symbol{92}}\texttt{\symbol{125}}}}
\logpage{[ 21, 3, 1 ]}\nobreak
\hyperdef{L}{X78791B8B838A8BA0}{}
{\noindent\textcolor{FuncColor}{$\triangleright$\enspace\texttt{\texttt{\symbol{92}}\texttt{\symbol{123}}\texttt{\symbol{92}}\texttt{\symbol{125}}({\mdseries\slshape list, poss})\index{\texttt{\symbol{92}}\texttt{\symbol{123}}\texttt{\symbol{92}}\texttt{\symbol{125}}@\texttt{\texttt{\symbol{92}}\texttt{\symbol{123}}\texttt{\symbol{92}}\texttt{\symbol{125}}}}
\label{bSlash{bSlash}}
}\hfill{\scriptsize (operation)}}\\


 This operation implements \emph{sublist access}. For any list, the default method is to loop over the entries in the list \mbox{\texttt{\mdseries\slshape poss}}, and to delegate to the element access operation. (For nested sublist
extractions, cf.{\nobreakspace}\ref{List Elements}. For the somewhat strange variable name, cf.{\nobreakspace}\ref{Basic Operations for Lists}.) }

 }

  
\section{\textcolor{Chapter }{List Assignment}}\label{List Assignment}
\logpage{[ 21, 4, 0 ]}
\hyperdef{L}{X8611EF768210625B}{}
{
  \index{assignment!to a list} \index{list element!assignment} \texttt{\mbox{\texttt{\mdseries\slshape list}}[ \mbox{\texttt{\mdseries\slshape ix}} ] := \mbox{\texttt{\mdseries\slshape object}};} 

 The list element assignment assigns the object \mbox{\texttt{\mdseries\slshape object}}, which can be of any type, to the list with index \mbox{\texttt{\mdseries\slshape ix}}, in the mutable (see{\nobreakspace}\ref{Mutability and Copyability}) list \mbox{\texttt{\mdseries\slshape list}}. That means that accessing the \mbox{\texttt{\mdseries\slshape ix}}\texttt{\symbol{45}}th element of the list \mbox{\texttt{\mdseries\slshape list}} will return \mbox{\texttt{\mdseries\slshape object}} after this assignment. 

 
\begin{Verbatim}[commandchars=!@|,fontsize=\small,frame=single,label=Example]
  !gapprompt@gap>| !gapinput@l := [ 1, 2, 3 ];;|
  !gapprompt@gap>| !gapinput@l[1] := 3;; l;             # assign a new object|
  [ 3, 2, 3 ]
  !gapprompt@gap>| !gapinput@l[2] := [ 4, 5, 6 ];; l;   # <object> may be of any type|
  [ 3, [ 4, 5, 6 ], 3 ]
  !gapprompt@gap>| !gapinput@l[ l[1] ] := 10;; l;       # <index> may be an expression|
  [ 3, [ 4, 5, 6 ], 10 ]
\end{Verbatim}
 

 If the index \mbox{\texttt{\mdseries\slshape ix}} is an integer larger than the length of the list \mbox{\texttt{\mdseries\slshape list}} (see \texttt{Length} (\ref{Length})), the list is automatically enlarged to make room for the new element. Note
that it is possible to generate lists with holes that way. 

 
\begin{Verbatim}[commandchars=!@|,fontsize=\small,frame=single,label=Example]
  !gapprompt@gap>| !gapinput@l[4] := "another entry";; l;  # <list> is enlarged|
  [ 3, [ 4, 5, 6 ], 10, "another entry" ]
  !gapprompt@gap>| !gapinput@l[ 10 ] := 1;; l;             # now <list> has a hole|
  [ 3, [ 4, 5, 6 ], 10, "another entry",,,,,, 1 ]
\end{Verbatim}
 

 The function \texttt{Add} (\ref{Add}) should be used if you want to add an element to the end of the list. 

 Note that assigning to a list changes the list, thus this list must be mutable
(see{\nobreakspace}\ref{Mutability and Copyability}). See{\nobreakspace}\ref{Identical Lists} for subtleties of changing lists. 

 If \mbox{\texttt{\mdseries\slshape list}} does not evaluate to a list, \mbox{\texttt{\mdseries\slshape pos}} does not evaluate to a positive integer, method selection is invoked to try
and find a way of indexing \mbox{\texttt{\mdseries\slshape list}} with index \mbox{\texttt{\mdseries\slshape pos}}. If this fails, or the selected method finds that \texttt{\mbox{\texttt{\mdseries\slshape list}}[\mbox{\texttt{\mdseries\slshape pos}}]} is unbound, or if \mbox{\texttt{\mdseries\slshape object}} is a call to a function which does not return a value (for example \texttt{Print}) an error is signalled. 

 \index{sublist!assignment} \texttt{\mbox{\texttt{\mdseries\slshape list}}\texttt{\symbol{123}} \mbox{\texttt{\mdseries\slshape poss}} \texttt{\symbol{125}} := \mbox{\texttt{\mdseries\slshape objects}};} 

 The sublist assignment assigns the object \texttt{\mbox{\texttt{\mdseries\slshape objects}}[1]}, which can be of any type, to the list \mbox{\texttt{\mdseries\slshape list}} at the position \texttt{\mbox{\texttt{\mdseries\slshape poss}}[1]}, the object \texttt{\mbox{\texttt{\mdseries\slshape objects}}[2]} to \texttt{\mbox{\texttt{\mdseries\slshape list}}[\mbox{\texttt{\mdseries\slshape poss}}[2]]}, and so on. \mbox{\texttt{\mdseries\slshape poss}} must be a dense list of positive integers, it need, however, not be sorted and
may contain duplicate elements. \mbox{\texttt{\mdseries\slshape objects}} must be a dense list and must have the same length as \mbox{\texttt{\mdseries\slshape poss}}. 

 
\begin{Verbatim}[commandchars=!@|,fontsize=\small,frame=single,label=Example]
  !gapprompt@gap>| !gapinput@l := [ 2, 3, 5, 7, 11, 13, 17, 19 ];;|
  !gapprompt@gap>| !gapinput@l{[1..4]} := [10..13];; l;|
  [ 10, 11, 12, 13, 11, 13, 17, 19 ]
  !gapprompt@gap>| !gapinput@l{[1,7,1,10]} := [ 1, 2, 3, 4 ];; l;|
  [ 3, 11, 12, 13, 11, 13, 2, 19,, 4 ]
\end{Verbatim}
 

 The next example shows that it is possible to nest such sublist assignments. 

 
\begin{Verbatim}[commandchars=!@|,fontsize=\small,frame=single,label=Example]
  !gapprompt@gap>| !gapinput@m := [ [1,2,3], [4,5,6], [7,8,9], [10,11,12] ];;|
  !gapprompt@gap>| !gapinput@m{[1,2,3]}{[3,2]} := [ [11,12], [13,14], [15,16] ];; m;|
  [ [ 1, 12, 11 ], [ 4, 14, 13 ], [ 7, 16, 15 ], [ 10, 11, 12 ] ]
\end{Verbatim}
 

 The exact behaviour is defined in the same way as for list extractions (see \ref{List Elements}). Namely, with each selector \texttt{[\mbox{\texttt{\mdseries\slshape pos}}]} or \texttt{\texttt{\symbol{123}}\mbox{\texttt{\mdseries\slshape poss}}\texttt{\symbol{125}}} we associate a \emph{level} that is defined as the number of selectors of the form \texttt{\texttt{\symbol{123}}\mbox{\texttt{\mdseries\slshape poss}}\texttt{\symbol{125}}} to its left in the same expression. For example 

 
\begin{Verbatim}[commandchars=!@|,fontsize=\small,frame=single,label=Example]
      l[pos1]{poss2}{poss3}[pos4]{poss5}[pos6]
  level   0      0      1     1      1     2
\end{Verbatim}
 

 Then a list assignment \texttt{\mbox{\texttt{\mdseries\slshape list}}[\mbox{\texttt{\mdseries\slshape pos}}] := \mbox{\texttt{\mdseries\slshape vals}};} of level \mbox{\texttt{\mdseries\slshape level}} is computed as \texttt{ListAssignment( \mbox{\texttt{\mdseries\slshape list}}, \mbox{\texttt{\mdseries\slshape pos}}, \mbox{\texttt{\mdseries\slshape vals}}, \mbox{\texttt{\mdseries\slshape level}} )}, where \texttt{ListAssignment} is defined as follows. (Note that \texttt{ListAssignment} is \emph{not} a \textsf{GAP} function.) 

 
\begin{Verbatim}[commandchars=!@|,fontsize=\small,frame=single,label=Example]
  ListAssignment := function ( list, pos, vals, level )
   local i;
   if level = 0 then
    list[pos] := vals;
   else
    for i in [1..Length(list)] do
     ListAssignment( list[i], pos, vals[i], level-1 );
    od;
   fi;
  end;
\end{Verbatim}
 

 and a list assignment \texttt{\mbox{\texttt{\mdseries\slshape list}}\texttt{\symbol{123}}\mbox{\texttt{\mdseries\slshape poss}}\texttt{\symbol{125}} := \mbox{\texttt{\mdseries\slshape vals}}} of level \mbox{\texttt{\mdseries\slshape level}} is computed as \texttt{ListAssignments( \mbox{\texttt{\mdseries\slshape list}}, \mbox{\texttt{\mdseries\slshape poss}}, \mbox{\texttt{\mdseries\slshape vals}}, \mbox{\texttt{\mdseries\slshape level}} )}, where \texttt{ListAssignments} is defined as follows. (Note that \texttt{ListAssignments} is \emph{not} a \textsf{GAP} function.) 

 
\begin{Verbatim}[commandchars=!@|,fontsize=\small,frame=single,label=Example]
  ListAssignments := function ( list, poss, vals, level )
   local i;
   if level = 0 then
    list{poss} := vals;
   else
    for i in [1..Length(list)] do
     ListAssignments( list[i], poss, vals[i], level-1 );
    od;
   fi;
  end;
\end{Verbatim}
 

 \index{sublist assignment!operation} 

\subsection{\textcolor{Chapter }{\texttt{\symbol{92}}\texttt{\symbol{123}}\texttt{\symbol{92}}\texttt{\symbol{125}}\texttt{\symbol{92}}:\texttt{\symbol{92}}=}}
\logpage{[ 21, 4, 1 ]}\nobreak
\hyperdef{L}{X813FF1637F8D2B7F}{}
{\noindent\textcolor{FuncColor}{$\triangleright$\enspace\texttt{\texttt{\symbol{92}}\texttt{\symbol{123}}\texttt{\symbol{92}}\texttt{\symbol{125}}\texttt{\symbol{92}}:\texttt{\symbol{92}}=({\mdseries\slshape list, poss, val})\index{\texttt{\symbol{92}}\texttt{\symbol{123}}\texttt{\symbol{92}}\texttt{\symbol{125}}\texttt{\symbol{92}}:\texttt{\symbol{92}}=@\texttt{\texttt{\symbol{92}}\texttt{\symbol{123}}\texttt{\symbol{92}}\texttt{\symbol{125}}\texttt{\symbol{92}}:\texttt{\symbol{92}}=}}
\label{bSlash{bSlash}bSlash:bSlash=}
}\hfill{\scriptsize (operation)}}\\


 This operation implements sublist assignment. For any list, the default method
is to loop over the entries in the list \mbox{\texttt{\mdseries\slshape poss}}, and to delegate to the element assignment operation. (For the somewhat
strange variable name, cf.{\nobreakspace}\ref{Basic Operations for Lists}.) }

 

\subsection{\textcolor{Chapter }{Add}}
\logpage{[ 21, 4, 2 ]}\nobreak
\hyperdef{L}{X795EC9D67E34DAB0}{}
{\noindent\textcolor{FuncColor}{$\triangleright$\enspace\texttt{Add({\mdseries\slshape list, obj[, pos]})\index{Add@\texttt{Add}}
\label{Add}
}\hfill{\scriptsize (operation)}}\\


 adds the element \mbox{\texttt{\mdseries\slshape obj}} to the mutable list \mbox{\texttt{\mdseries\slshape list}}. The two argument version adds \mbox{\texttt{\mdseries\slshape obj}} at the end of \mbox{\texttt{\mdseries\slshape list}}, i.e., it is equivalent to the assignment \texttt{\mbox{\texttt{\mdseries\slshape list}}[ Length(\mbox{\texttt{\mdseries\slshape list}}) + 1 ] := \mbox{\texttt{\mdseries\slshape obj}}}, see{\nobreakspace}\ref{List Assignment}. 

 The three argument version adds \mbox{\texttt{\mdseries\slshape obj}} in position \mbox{\texttt{\mdseries\slshape pos}}, moving all later elements of the list (if any) up by one position. Any holes
at or after position \mbox{\texttt{\mdseries\slshape pos}} are also moved up by one position, and new holes are created before \mbox{\texttt{\mdseries\slshape pos}} if they are needed. 

 Nothing is returned by \texttt{Add}, the function is only called for its side effect. }

 

\subsection{\textcolor{Chapter }{Remove}}
\logpage{[ 21, 4, 3 ]}\nobreak
\hyperdef{L}{X7E98B11B79BA9167}{}
{\noindent\textcolor{FuncColor}{$\triangleright$\enspace\texttt{Remove({\mdseries\slshape list[, pos]})\index{Remove@\texttt{Remove}}
\label{Remove}
}\hfill{\scriptsize (operation)}}\\


 removes an element from \mbox{\texttt{\mdseries\slshape list}}. The one argument form removes the last element. The two argument form
removes the element in position \mbox{\texttt{\mdseries\slshape pos}}, moving all subsequent elements down one position. Any holes after position \mbox{\texttt{\mdseries\slshape pos}} are also moved down by one position. 

 The one argument form always returns the removed element. In this case \mbox{\texttt{\mdseries\slshape list}} must be non\texttt{\symbol{45}}empty. 

 The two argument form returns the old value of \mbox{\texttt{\mdseries\slshape list}}[\mbox{\texttt{\mdseries\slshape pos}}] if it was bound, and nothing if it was not. Note that accessing or assigning
the return value of this form of the \texttt{Remove} operation is only safe when you \emph{know} that there will be a value, otherwise it will cause an error. 

 
\begin{Verbatim}[commandchars=!@|,fontsize=\small,frame=single,label=Example]
  !gapprompt@gap>| !gapinput@l := [ 2, 3, 5 ];; Add( l, 7 ); l;|
  [ 2, 3, 5, 7 ]
  !gapprompt@gap>| !gapinput@Add(l,4,2); l;|
  [ 2, 4, 3, 5, 7 ]
  !gapprompt@gap>| !gapinput@Remove(l,2); l;|
  4
  [ 2, 3, 5, 7 ]
  !gapprompt@gap>| !gapinput@Remove(l); l;|
  7
  [ 2, 3, 5 ]
  !gapprompt@gap>| !gapinput@Remove(l,5); l;|
  [ 2, 3, 5 ]
\end{Verbatim}
 }

 

\subsection{\textcolor{Chapter }{CopyListEntries}}
\logpage{[ 21, 4, 4 ]}\nobreak
\hyperdef{L}{X79D7E96F80A2D7C0}{}
{\noindent\textcolor{FuncColor}{$\triangleright$\enspace\texttt{CopyListEntries({\mdseries\slshape fromlst, fromind, fromstep, tolst, toind, tostep, n})\index{CopyListEntries@\texttt{CopyListEntries}}
\label{CopyListEntries}
}\hfill{\scriptsize (function)}}\\


 This function copies \mbox{\texttt{\mdseries\slshape n}} elements from \mbox{\texttt{\mdseries\slshape fromlst}}, starting at position \mbox{\texttt{\mdseries\slshape fromind}} and incrementing the position by \mbox{\texttt{\mdseries\slshape fromstep}} each time, into \mbox{\texttt{\mdseries\slshape tolst}} starting at position \mbox{\texttt{\mdseries\slshape toind}} and incrementing the position by \mbox{\texttt{\mdseries\slshape tostep}} each time. \mbox{\texttt{\mdseries\slshape fromlst}} and \mbox{\texttt{\mdseries\slshape tolst}} must be plain lists. \mbox{\texttt{\mdseries\slshape fromstep}} and/or \mbox{\texttt{\mdseries\slshape tostep}} can be negative. Unbound positions of \mbox{\texttt{\mdseries\slshape fromlst}} are simply copied to \mbox{\texttt{\mdseries\slshape tolst}}. 

 \texttt{CopyListEntries} is used in methods for the operations \texttt{Add} (\ref{Add}) and \texttt{Remove} (\ref{Remove}). }

 

\subsection{\textcolor{Chapter }{Append}}
\logpage{[ 21, 4, 5 ]}\nobreak
\hyperdef{L}{X79E31DB27C82D6E1}{}
{\noindent\textcolor{FuncColor}{$\triangleright$\enspace\texttt{Append({\mdseries\slshape list1, list2})\index{Append@\texttt{Append}}
\label{Append}
}\hfill{\scriptsize (operation)}}\\


 adds the elements of the list \mbox{\texttt{\mdseries\slshape list2}} to the end of the mutable list \mbox{\texttt{\mdseries\slshape list1}}, see{\nobreakspace}\ref{List Assignment}. \mbox{\texttt{\mdseries\slshape list2}} may contain holes, in which case the corresponding entries in \mbox{\texttt{\mdseries\slshape list1}} will be left unbound. \texttt{Append} returns nothing, it is only called for its side effect. 

 Note that \texttt{Append} changes its first argument, while \texttt{Concatenation} (\ref{Concatenation:for a list of lists}) creates a new list and leaves its arguments unchanged. 

 
\begin{Verbatim}[commandchars=!@|,fontsize=\small,frame=single,label=Example]
  !gapprompt@gap>| !gapinput@l := [ 2, 3, 5 ];; Append( l, [ 7, 11, 13 ] ); l;|
  [ 2, 3, 5, 7, 11, 13 ]
  !gapprompt@gap>| !gapinput@Append( l, [ 17,, 23 ] ); l;|
  [ 2, 3, 5, 7, 11, 13, 17,, 23 ]
\end{Verbatim}
 }

 }

  
\section{\textcolor{Chapter }{IsBound and Unbind for Lists}}\label{IsBound and Unbind for Lists}
\logpage{[ 21, 5, 0 ]}
\hyperdef{L}{X7963C8E17EFF86DB}{}
{
  

\subsection{\textcolor{Chapter }{IsBound (for a list index)}}
\logpage{[ 21, 5, 1 ]}\nobreak
\hyperdef{L}{X79EC565A7DCEC938}{}
{\noindent\textcolor{FuncColor}{$\triangleright$\enspace\texttt{IsBound({\mdseries\slshape list[, n]})\index{IsBound@\texttt{IsBound}!for a list index}
\label{IsBound:for a list index}
}\hfill{\scriptsize (operation)}}\\


 \texttt{IsBound} returns \texttt{true} if the list \mbox{\texttt{\mdseries\slshape list}} has an element at index \mbox{\texttt{\mdseries\slshape n}}, and \texttt{false} otherwise. \mbox{\texttt{\mdseries\slshape list}} must evaluate to a list, or to an object for which a suitable method for \texttt{IsBound\texttt{\symbol{92}}[\texttt{\symbol{92}}]} has been installed, otherwise an error is signalled. 

 
\begin{Verbatim}[commandchars=!@|,fontsize=\small,frame=single,label=Example]
  !gapprompt@gap>| !gapinput@l := [ , 2, 3, , 5, , 7, , , , 11 ];;|
  !gapprompt@gap>| !gapinput@IsBound( l[7] );|
  true
  !gapprompt@gap>| !gapinput@IsBound( l[4] );|
  false
  !gapprompt@gap>| !gapinput@IsBound( l[101] );|
  false
\end{Verbatim}
 }

 

\subsection{\textcolor{Chapter }{GetWithDefault}}
\logpage{[ 21, 5, 2 ]}\nobreak
\hyperdef{L}{X866F45D3797FDA00}{}
{\noindent\textcolor{FuncColor}{$\triangleright$\enspace\texttt{GetWithDefault({\mdseries\slshape list, n, default})\index{GetWithDefault@\texttt{GetWithDefault}}
\label{GetWithDefault}
}\hfill{\scriptsize (operation)}}\\


 \texttt{GetWithDefault} returns the \mbox{\texttt{\mdseries\slshape n}}th element of the list \mbox{\texttt{\mdseries\slshape list}}, if \mbox{\texttt{\mdseries\slshape list}} has a value at index \mbox{\texttt{\mdseries\slshape n}}, and \mbox{\texttt{\mdseries\slshape default}} otherwise. 

 While this method can be used on any list, it is particularly useful for Weak
Pointer lists \ref{Weak Pointer Objects} where the value of the list can change. 

 To distinguish between the \mbox{\texttt{\mdseries\slshape n}}th element being unbound, or \mbox{\texttt{\mdseries\slshape default}} being in \mbox{\texttt{\mdseries\slshape list}}, users can create a new mutable object, such as a string. \texttt{IsIdenticalObj} (\ref{IsIdenticalObj}) returns \texttt{false} for different mutable strings, even if their contents are the same. 
\begin{Verbatim}[commandchars=!@|,fontsize=\small,frame=single,label=Example]
  !gapprompt@gap>| !gapinput@l := [1,2,,"a"];|
  [ 1, 2,, "a" ]
  !gapprompt@gap>| !gapinput@newobj := "a";|
  "a"
  !gapprompt@gap>| !gapinput@GetWithDefault(l, 2, newobj);|
  2
  !gapprompt@gap>| !gapinput@GetWithDefault(l, 3, newobj);|
  "a"
  !gapprompt@gap>| !gapinput@GetWithDefault(l, 4, newobj);|
  "a"
  !gapprompt@gap>| !gapinput@IsIdenticalObj(GetWithDefault(l, 3, newobj), newobj);|
  true
  !gapprompt@gap>| !gapinput@IsIdenticalObj(GetWithDefault(l, 4, newobj), newobj);|
  false
\end{Verbatim}
 }

 

\subsection{\textcolor{Chapter }{Unbind (unbind a list entry)}}
\logpage{[ 21, 5, 3 ]}\nobreak
\hyperdef{L}{X78B72FDF7BD63C0B}{}
{\noindent\textcolor{FuncColor}{$\triangleright$\enspace\texttt{Unbind({\mdseries\slshape list[, n]})\index{Unbind@\texttt{Unbind}!unbind a list entry}
\label{Unbind:unbind a list entry}
}\hfill{\scriptsize (operation)}}\\


 \texttt{Unbind} deletes the element with index \mbox{\texttt{\mdseries\slshape n}} in the mutable list \mbox{\texttt{\mdseries\slshape list}}. That is, after execution of \texttt{Unbind}, \mbox{\texttt{\mdseries\slshape list}} no longer has an assigned value with index \mbox{\texttt{\mdseries\slshape n}}. Thus \texttt{Unbind} can be used to produce holes in a list. Note that it is not an error to unbind
a nonexistent list element. \mbox{\texttt{\mdseries\slshape list}} must evaluate to a list, or to an object for which a suitable method for \texttt{Unbind\texttt{\symbol{92}}[\texttt{\symbol{92}}]} has been installed, otherwise an error is signalled. 

 
\begin{Verbatim}[commandchars=!@|,fontsize=\small,frame=single,label=Example]
  !gapprompt@gap>| !gapinput@l := [ , 2, 3, 5, , 7, , , , 11 ];;|
  !gapprompt@gap>| !gapinput@Unbind( l[3] ); l;|
  [ , 2,, 5,, 7,,,, 11 ]
  !gapprompt@gap>| !gapinput@Unbind( l[4] ); l;|
  [ , 2,,,, 7,,,, 11 ]
\end{Verbatim}
 

 Note that \texttt{IsBound} (\ref{IsBound:for a list index}) and \texttt{Unbind} are special in that they do not evaluate their argument, otherwise \texttt{IsBound} (\ref{IsBound:for a list index}) would always signal an error when it is supposed to return \texttt{false} and there would be no way to tell \texttt{Unbind} which component to remove. }

 }

  
\section{\textcolor{Chapter }{Identical Lists}}\label{Identical Lists}
\logpage{[ 21, 6, 0 ]}
\hyperdef{L}{X7DD65BEA7EDB0CD7}{}
{
  With the list assignment (see{\nobreakspace}\ref{List Assignment}) it is possible to change a mutable list. This section describes the semantic
consequences of this fact. (See also{\nobreakspace}\ref{Identical Objects}.) 

 First we define what it means when we say that ``an object is changed''. You may think that in the following example the second assignment changes
the integer. 

 
\begin{Verbatim}[commandchars=!@|,fontsize=\small,frame=single,label=Example]
  i := 3;
  i := i + 1;
\end{Verbatim}
 

 But in this example it is not the \emph{integer} \texttt{3} which is changed, by adding one to it. Instead the \emph{variable} \texttt{i} is changed by assigning the value of \texttt{i+1}, which happens to be \texttt{4}, to \texttt{i}. The same thing happens in the example below. 

 
\begin{Verbatim}[commandchars=!@|,fontsize=\small,frame=single,label=Example]
  l := [ 1, 2 ];
  l := [ 1, 2, 3 ];
\end{Verbatim}
 

 The second assignment does not change the first list, instead it assigns a new
list to the variable \texttt{l}. On the other hand, in the following example the list \emph{is} changed by the second assignment. 

 
\begin{Verbatim}[commandchars=!@|,fontsize=\small,frame=single,label=Example]
  l := [ 1, 2 ];
  l[3] := 3;
\end{Verbatim}
 

 To understand the difference, think of a variable as a name for an object. The
important point is that a list can have several names at the same time. An
assignment \texttt{\mbox{\texttt{\mdseries\slshape var}}:= \mbox{\texttt{\mdseries\slshape list}};} means in this interpretation that \mbox{\texttt{\mdseries\slshape var}} is a name for the object \mbox{\texttt{\mdseries\slshape list}}. At the end of the following example \texttt{l2} still has the value \texttt{[ 1, 2 ]} as this list has not been changed and nothing else has been assigned to it. 

 
\begin{Verbatim}[commandchars=!@|,fontsize=\small,frame=single,label=Example]
  l1 := [ 1, 2 ];
  l2 := l1;
  l1 := [ 1, 2, 3 ];
\end{Verbatim}
 

 But after the following example the list for which \texttt{l2} is a name has been changed and thus the value of \texttt{l2} is now \texttt{[ 1, 2, 3 ]}. 

 
\begin{Verbatim}[commandchars=!@|,fontsize=\small,frame=single,label=Example]
  l1 := [ 1, 2 ];
  l2 := l1;
  l1[3] := 3;
\end{Verbatim}
 

 We say that two lists are \emph{identical} if changing one of them by a list assignment also changes the other one. This
is slightly incorrect, because if \emph{two} lists are identical, there are actually only two names for \emph{one} list. However, the correct usage would be very awkward and would only add to
the confusion. Note that two identical lists must be equal, because there is
only one list with two different names. Thus identity is an equivalence
relation that is a refinement of equality. Identity of objects can be detected
using \texttt{IsIdenticalObj} (\ref{IsIdenticalObj}). 

 Let us now consider under which circumstances two lists are identical. 

 If you enter a list literal then the list denoted by this literal is a new
list that is not identical to any other list. Thus in the following example \texttt{l1} and \texttt{l2} are not identical, though they are equal of course. 

 
\begin{Verbatim}[commandchars=!@|,fontsize=\small,frame=single,label=Example]
  l1 := [ 1, 2 ];
  l2 := [ 1, 2 ];
\end{Verbatim}
 

 Also in the following example, no lists in the list \texttt{l} are identical. 

 
\begin{Verbatim}[commandchars=!@|,fontsize=\small,frame=single,label=Example]
  l := [];
  for i in [1..10] do l[i] := [ 1, 2 ]; od;
\end{Verbatim}
 

 If you assign a list to a variable no new list is created. Thus the list value
of the variable on the left hand side and the list on the right hand side of
the assignment are identical. So in the following example \texttt{l1} and \texttt{l2} are identical lists. 

 
\begin{Verbatim}[commandchars=!@|,fontsize=\small,frame=single,label=Example]
  l1 := [ 1, 2 ];
  l2 := l1;
\end{Verbatim}
 

 If you pass a list as an argument, the old list and the argument of the
function are identical. Also if you return a list from a function, the old
list and the value of the function call are identical. So in the following
example \texttt{l1} and \texttt{l2} are identical lists: 

 
\begin{Verbatim}[commandchars=!@|,fontsize=\small,frame=single,label=Example]
  l1 := [ 1, 2 ];
  f := function ( l ) return l; end;
  l2 := f( l1 );
\end{Verbatim}
 

 If you change a list it keeps its identity. Thus if two lists are identical
and you change one of them, you also change the other, and they are still
identical afterwards. On the other hand, two lists that are not identical will
never become identical if you change one of them. So in the following example
both \texttt{l1} and \texttt{l2} are changed, and are still identical. 

 
\begin{Verbatim}[commandchars=!@|,fontsize=\small,frame=single,label=Example]
  l1 := [ 1, 2 ];
  l2 := l1;
  l1[1] := 2;
\end{Verbatim}
 }

  
\section{\textcolor{Chapter }{Duplication of Lists}}\label{Duplication of Lists}
\logpage{[ 21, 7, 0 ]}
\hyperdef{L}{X7ED7C0738495556F}{}
{
  Here we describe the meaning of \texttt{ShallowCopy} (\ref{ShallowCopy}) and \texttt{StructuralCopy} (\ref{StructuralCopy}) for lists. For the general definition of these functions, see{\nobreakspace}\ref{Duplication of Objects}. 

 \index{ShallowCopy@\texttt{ShallowCopy}!for lists} The subobjects (see{\nobreakspace}\texttt{ShallowCopy} (\ref{ShallowCopy})) of a list are exactly its elements. 

 This means that for any list \mbox{\texttt{\mdseries\slshape list}}, \texttt{ShallowCopy} (\ref{ShallowCopy}) returns a mutable \emph{new} list \mbox{\texttt{\mdseries\slshape new}} that is \emph{not identical} to any other list (see{\nobreakspace}\ref{Identical Lists}), and whose elements are identical to the elements of \mbox{\texttt{\mdseries\slshape list}}. 

 \index{StructuralCopy@\texttt{StructuralCopy}!for lists} Analogously, for a \emph{mutable} list \mbox{\texttt{\mdseries\slshape list}}, \texttt{StructuralCopy} (\ref{StructuralCopy}) returns a mutable \emph{new} list \mbox{\texttt{\mdseries\slshape scp}} that is \emph{not identical} to any other list, and whose elements are structural copies (defined
recursively) of the elements of \mbox{\texttt{\mdseries\slshape list}}; an element of \mbox{\texttt{\mdseries\slshape scp}} is mutable (and then a \emph{new} list) if and only if the corresponding element of \mbox{\texttt{\mdseries\slshape list}} is mutable. 

 In both cases, modifying the copy \mbox{\texttt{\mdseries\slshape new}} resp.{\nobreakspace}\mbox{\texttt{\mdseries\slshape scp}} by assignments (see{\nobreakspace}\ref{List Assignment}) does not modify the original object \mbox{\texttt{\mdseries\slshape list}}. 

 \texttt{ShallowCopy} (\ref{ShallowCopy}) basically executes the following code for lists. 
\begin{Verbatim}[commandchars=!@|,fontsize=\small,frame=single,label=Example]
  new := [];
  for i in [ 1 .. Length( list ) ] do
    if IsBound( list[i] ) then
      new[i] := list[i];
    fi;
  od;
\end{Verbatim}
 

 
\begin{Verbatim}[commandchars=!@|,fontsize=\small,frame=single,label=Example]
  !gapprompt@gap>| !gapinput@list1 := [ [ 1, 2 ], [ 3, 4 ] ];;  list2 := ShallowCopy( list1 );;|
  !gapprompt@gap>| !gapinput@IsIdenticalObj( list1, list2 );|
  false
  !gapprompt@gap>| !gapinput@IsIdenticalObj( list1[1], list2[1] );|
  true
  !gapprompt@gap>| !gapinput@list2[1] := 0;;  list1;  list2;|
  [ [ 1, 2 ], [ 3, 4 ] ]
  [ 0, [ 3, 4 ] ]
\end{Verbatim}
 

 \texttt{StructuralCopy} (\ref{StructuralCopy}) basically executes the following code for lists. 
\begin{Verbatim}[commandchars=!@|,fontsize=\small,frame=single,label=Example]
  new := [];
  for i in [ 1 .. Length( list ) ] do
    if IsBound( list[i] ) then
      new[i] := StructuralCopy( list[i] );
    fi;
  od;
\end{Verbatim}
 

 
\begin{Verbatim}[commandchars=!@|,fontsize=\small,frame=single,label=Example]
  !gapprompt@gap>| !gapinput@list1 := [ [ 1, 2 ], [ 3, 4 ] ];;  list2 := StructuralCopy( list1 );;|
  !gapprompt@gap>| !gapinput@IsIdenticalObj( list1, list2 );|
  false
  !gapprompt@gap>| !gapinput@IsIdenticalObj( list1[1], list2[1] );|
  false
  !gapprompt@gap>| !gapinput@list2[1][1] := 0;;  list1;  list2;|
  [ [ 1, 2 ], [ 3, 4 ] ]
  [ [ 0, 2 ], [ 3, 4 ] ]
\end{Verbatim}
 

 The above code is not entirely correct. If the object \mbox{\texttt{\mdseries\slshape list}} contains a mutable object twice this object is not copied twice, as would
happen with the above definition, but only once. This means that the copy \mbox{\texttt{\mdseries\slshape new}} and the object \mbox{\texttt{\mdseries\slshape list}} have exactly the same structure when viewed as a general graph. 

 
\begin{Verbatim}[commandchars=!@|,fontsize=\small,frame=single,label=Example]
  !gapprompt@gap>| !gapinput@sub := [ 1, 2 ];; list1 := [ sub, sub ];;|
  !gapprompt@gap>| !gapinput@list2 := StructuralCopy( list1 );|
  [ [ 1, 2 ], [ 1, 2 ] ]
  !gapprompt@gap>| !gapinput@list2[1][1] := 0;; list2;|
  [ [ 0, 2 ], [ 0, 2 ] ]
  !gapprompt@gap>| !gapinput@list1;|
  [ [ 1, 2 ], [ 1, 2 ] ]
\end{Verbatim}
 }

  
\section{\textcolor{Chapter }{Membership Test for Lists}}\label{Membership Test for Lists}
\logpage{[ 21, 8, 0 ]}
\hyperdef{L}{X808A207182B2F84F}{}
{
  

\subsection{\textcolor{Chapter }{\texttt{\symbol{92}}in (element test for lists)}}
\logpage{[ 21, 8, 1 ]}\nobreak
\hyperdef{L}{X7B914A287F88ED0A}{}
{\noindent\textcolor{FuncColor}{$\triangleright$\enspace\texttt{\texttt{\symbol{92}}in({\mdseries\slshape obj, list})\index{\texttt{\symbol{92}}in@\texttt{\texttt{\symbol{92}}in}!element test for lists}
\label{bSlashin:element test for lists}
}\hfill{\scriptsize (operation)}}\\


 \index{in@\texttt{in}!for lists} \index{element test!for lists} This function call or the infix variant \mbox{\texttt{\mdseries\slshape obj}}{\nobreakspace}\texttt{in}{\nobreakspace}\mbox{\texttt{\mdseries\slshape list}} tests whether there is a positive integer $i$ such that \mbox{\texttt{\mdseries\slshape list}}$[i] =$ \mbox{\texttt{\mdseries\slshape obj}} holds. 

 If the list \mbox{\texttt{\mdseries\slshape list}} knows that it is strictly sorted (see{\nobreakspace}\texttt{IsSSortedList} (\ref{IsSSortedList})), the membership test is much quicker, because a binary search can be used
instead of the linear search used for arbitrary lists, see \texttt{\texttt{\symbol{92}}in} (\ref{bSlashin:for strictly sorted lists}). 

 
\begin{Verbatim}[commandchars=!@|,fontsize=\small,frame=single,label=Example]
  !gapprompt@gap>| !gapinput@1 in [ 2, 2, 1, 3 ];  1 in [ 4, -1, 0, 3 ];|
  true
  false
  !gapprompt@gap>| !gapinput@s := SSortedList( [2,4,6,8,10,12,14,16,18,20,22,24,26,28,30,32] );;|
  !gapprompt@gap>| !gapinput@17 in s;  # uses binary search and only 4 comparisons|
  false
\end{Verbatim}
 

 For finding the position of an element in a list, see{\nobreakspace}\ref{Finding Positions in Lists}. }

 }

  
\section{\textcolor{Chapter }{Enlarging Internally Represented Lists}}\label{Enlarging Internally Represented Lists}
\logpage{[ 21, 9, 0 ]}
\hyperdef{L}{X84D6FC7E7E39ED33}{}
{
  Section{\nobreakspace}\ref{List Assignment} told you (among other things) that it is possible to assign beyond the logical
end of a mutable list, automatically enlarging the list. This section tells
you how this is done for internally represented lists. 

 It would be extremely wasteful to make all lists large enough so that there is
room for all assignments, because some lists may have more than 100000
elements, while most lists have less than 10 elements. 

 On the other hand suppose every assignment beyond the end of a list would be
done by allocating new space for the list and copying all entries to the new
space. Then creating a list of 1000 elements by assigning them in order, would
take half a million copy operations and also create a lot of garbage that the
garbage collector would have to reclaim. 

 So the following strategy is used. If a list is created it is created with
exactly the correct size. If a list is enlarged, because of an assignment
beyond the end of the list, it is enlarged by at least \texttt{\mbox{\texttt{\mdseries\slshape length}}/8 + 4} entries. Therefore the next assignments beyond the end of the list do not need
to enlarge the list. For example creating a list of 1000 elements by assigning
them in order, would now take only 32 enlargements. 

 The result of this is of course that the \emph{physical length} of a list may be larger than the \emph{logical length}, which is usually called simply the length of the list. Aside from the
implications for the performance you need not be aware of the physical length.
In fact all you can ever observe, for example by calling \texttt{Length} (\ref{Length}), is the logical length. 

 Suppose that \texttt{Length} (\ref{Length}) would have to take the physical length and then test how many entries at the
end of a list are unassigned, to compute the logical length of the list. That
would take too much time. In order to make \texttt{Length} (\ref{Length}), and other functions that need to know the logical length, more efficient,
the length of a list is stored along with the list. 

 For fine tuning code dealing with plain lists we provide the following two
functions. 

 

\subsection{\textcolor{Chapter }{EmptyPlist}}
\logpage{[ 21, 9, 1 ]}\nobreak
\hyperdef{L}{X78BF67A5802E93AD}{}
{\noindent\textcolor{FuncColor}{$\triangleright$\enspace\texttt{EmptyPlist({\mdseries\slshape len})\index{EmptyPlist@\texttt{EmptyPlist}}
\label{EmptyPlist}
}\hfill{\scriptsize (function)}}\\
\textbf{\indent Returns:\ }
a plain list

\noindent\textcolor{FuncColor}{$\triangleright$\enspace\texttt{ShrinkAllocationPlist({\mdseries\slshape l})\index{ShrinkAllocationPlist@\texttt{ShrinkAllocationPlist}}
\label{ShrinkAllocationPlist}
}\hfill{\scriptsize (function)}}\\
\textbf{\indent Returns:\ }
nothing



 The function \texttt{EmptyPlist} returns an empty plain list which has enough memory allocated for \mbox{\texttt{\mdseries\slshape len}} entries. This can be useful for creating and filling a plain list with a known
number of entries. 

 The function \texttt{ShrinkAllocationPlist} gives back to \textsf{GAP}'s memory manager the physical memory which is allocated for the plain list \mbox{\texttt{\mdseries\slshape l}} but not needed by the current number of entries. 

 Note that there are similar functions \texttt{EmptyString} (\ref{EmptyString}) and \texttt{ShrinkAllocationString} (\ref{ShrinkAllocationString}) for strings instead of plain lists. 
\begin{Verbatim}[commandchars=!@|,fontsize=\small,frame=single,label=Example]
  !gapprompt@gap>| !gapinput@l:=[]; for i in [1..160] do Add(l, i^2); od;|
  [  ]
  !gapprompt@gap>| !gapinput@m:=EmptyPlist(160); for i in [1..160] do Add(m, i^2); od;|
  [  ]
  !gapprompt@gap>| !gapinput@# now l uses about 25% more memory than the equal list m|
  !gapprompt@gap>| !gapinput@ShrinkAllocationPlist(l);|
  !gapprompt@gap>| !gapinput@# now l and m use the same amount of memory|
\end{Verbatim}
 }

 }

  
\section{\textcolor{Chapter }{Comparisons of Lists}}\label{Comparisons of Lists}
\logpage{[ 21, 10, 0 ]}
\hyperdef{L}{X8016D50F85147A77}{}
{
  \index{comparisons!of lists} \index{list equal!comparison} \texttt{\mbox{\texttt{\mdseries\slshape list1}} = \mbox{\texttt{\mdseries\slshape list2}}} 

 \texttt{\mbox{\texttt{\mdseries\slshape list1}} {\textless}{\textgreater} \mbox{\texttt{\mdseries\slshape list2}}} 

 Two lists \mbox{\texttt{\mdseries\slshape list1}} and \mbox{\texttt{\mdseries\slshape list2}} are equal if and only if for every index $i$, either both entries \mbox{\texttt{\mdseries\slshape list1}}$[i]$ and \mbox{\texttt{\mdseries\slshape list2}}$[i]$ are unbound, or both are bound and are equal, i.e., \mbox{\texttt{\mdseries\slshape list1}}$[i] =$ \mbox{\texttt{\mdseries\slshape list2}}$[i]$ is \texttt{true}. 

 
\begin{Verbatim}[commandchars=!@|,fontsize=\small,frame=single,label=Example]
  !gapprompt@gap>| !gapinput@[ 1, 2, 3 ] = [ 1, 2, 3 ];|
  true
  !gapprompt@gap>| !gapinput@[ , 2, 3 ] = [ 1, 2, ];|
  false
  !gapprompt@gap>| !gapinput@[ 1, 2, 3 ] = [ 3, 2, 1 ];|
  false
\end{Verbatim}
 

 This definition will cause problems with lists which are their own entries.
Comparing two such lists for equality may lead to an infinite recursion in the
kernel if the list comparison has to compare the list entries which are in
fact the lists themselves, and then \textsf{GAP} crashes. 

 \index{list smaller!comparison} \texttt{\mbox{\texttt{\mdseries\slshape list1}} {\textless} \mbox{\texttt{\mdseries\slshape list2}}} 

 \texttt{\mbox{\texttt{\mdseries\slshape list1}} {\textless}= \mbox{\texttt{\mdseries\slshape list2}}} 

 Lists are ordered \emph{lexicographically}. Unbound entries are smaller than any bound entry. That implies the following
behaviour. Let $i$ be the smallest positive integer $i$ such that \mbox{\texttt{\mdseries\slshape list1}} and \mbox{\texttt{\mdseries\slshape list2}} at position $i$ differ, i.e., either exactly one of \mbox{\texttt{\mdseries\slshape list1}}$[i]$, \mbox{\texttt{\mdseries\slshape list2}}$[i]$ is bound or both entries are bound and differ. Then \mbox{\texttt{\mdseries\slshape list1}} is less than \mbox{\texttt{\mdseries\slshape list2}} if either \mbox{\texttt{\mdseries\slshape list1}}$[i]$ is unbound (and \mbox{\texttt{\mdseries\slshape list2}}$[i]$ is not) or both are bound and \mbox{\texttt{\mdseries\slshape list1}}$[i]$ {\textless} \mbox{\texttt{\mdseries\slshape list2}}$[i]$ is \texttt{true}. 

 
\begin{Verbatim}[commandchars=!@|,fontsize=\small,frame=single,label=Example]
  !gapprompt@gap>| !gapinput@[ 1, 2, 3, 4 ] < [ 1, 2, 4, 8 ]; # <list1>[3] < <list2>[3]|
  true
  !gapprompt@gap>| !gapinput@[ 1, 2, 3 ] < [ 1, 2, 3, 5 ];  # <list1>[4] is unbound and thus < 5|
  true
  !gapprompt@gap>| !gapinput@[ 1, , 3, 4 ] < [ 1, -1, 3 ];  # <list1>[2] is unbound and thus < -1|
  true
\end{Verbatim}
 

 Note that for comparing two lists with \texttt{{\textless}} or \texttt{{\textless}=}, the (relevant) list elements must be comparable with \texttt{{\textless}}, which is usually \emph{not} the case for objects in different families, see{\nobreakspace}\ref{Families}. Also for the possibility to compare lists with other objects,
see{\nobreakspace}\ref{Families}. }

  
\section{\textcolor{Chapter }{Arithmetic for Lists}}\label{Arithmetic for Lists}
\logpage{[ 21, 11, 0 ]}
\hyperdef{L}{X845EEAF083D43CCE}{}
{
  \index{operators!for lists} It is convenient to have arithmetic operations for lists, in particular
because in \textsf{GAP} row vectors and matrices are special kinds of lists. However, it is the wide
variety of list objects because of which we prescribe arithmetic operations \emph{not for all} of them. (Keep in mind that ``list'' means just an object in the category \texttt{IsList} (\ref{IsList}).) 

 (Due to the intended generality and flexibility, the definitions given in the
following sections are quite technical. But for not too complicated cases such
as matrices (see{\nobreakspace}\ref{Operators for Matrices}) and row vectors (see{\nobreakspace}\ref{Operators for Row Vectors}) whose entries aren't lists, the resulting behaviour should be intuitive.) 

 For example, we want to deal with matrices which can be added and multiplied
in the usual way, via the infix operators \texttt{+} and \texttt{*}; and we want also Lie matrices, with the same additive behaviour but with the
multiplication defined by the Lie bracket. Both kinds of matrices shall be
lists, with the usual access to their rows, with \texttt{Length} (\ref{Length}) returning the number of rows etc. 

 For the categories and attributes that control the arithmetic behaviour of
lists, see{\nobreakspace}\ref{Filters Controlling the Arithmetic Behaviour of Lists}. 

 For the definition of return values of additive and multiplicative operations
whose arguments are lists in these filters, see{\nobreakspace}\ref{Additive Arithmetic for Lists} and \ref{Multiplicative Arithmetic for Lists}, respectively. It should be emphasized that these sections describe only what
the return values are, and not how they are computed. 

 For the mutability status of the return values, see{\nobreakspace}\ref{Mutability Status and List Arithmetic}. (Note that this is not dealt with in the sections about the result values.) 

 Further details about the special cases of row vectors and matrices can be
found in{\nobreakspace}\ref{Operators for Row Vectors} and in{\nobreakspace}\ref{Operators for Matrices}, the compression status is dealt with in{\nobreakspace}\ref{Row Vectors over Finite Fields} and{\nobreakspace}\ref{Matrices over Finite Fields}. }

  
\section{\textcolor{Chapter }{Filters Controlling the Arithmetic Behaviour of Lists}}\label{Filters Controlling the Arithmetic Behaviour of Lists}
\logpage{[ 21, 12, 0 ]}
\hyperdef{L}{X84D642967B8546B7}{}
{
  The arithmetic behaviour of lists is controlled by their types. The following
categories and attributes are used for that. 

 Note that we distinguish additive and multiplicative behaviour. For example,
Lie matrices have the usual additive behaviour but not the usual
multiplicative behaviour. 

\subsection{\textcolor{Chapter }{IsGeneralizedRowVector}}
\logpage{[ 21, 12, 1 ]}\nobreak
\hyperdef{L}{X87ABCEE9809585A0}{}
{\noindent\textcolor{FuncColor}{$\triangleright$\enspace\texttt{IsGeneralizedRowVector({\mdseries\slshape list})\index{IsGeneralizedRowVector@\texttt{IsGeneralizedRowVector}}
\label{IsGeneralizedRowVector}
}\hfill{\scriptsize (Category)}}\\


 For a list \mbox{\texttt{\mdseries\slshape list}}, the value \texttt{true} for \texttt{IsGeneralizedRowVector} indicates that the additive arithmetic behaviour of \mbox{\texttt{\mdseries\slshape list}} is as defined in \ref{Additive Arithmetic for Lists}, and that the attribute \texttt{NestingDepthA} (\ref{NestingDepthA}) will return a nonzero value when called with \mbox{\texttt{\mdseries\slshape list}}. 

 
\begin{Verbatim}[commandchars=!@|,fontsize=\small,frame=single,label=Example]
  !gapprompt@gap>| !gapinput@IsList( "abc" ); IsGeneralizedRowVector( "abc" );|
  true
  false
  !gapprompt@gap>| !gapinput@liemat:= LieObject( [ [ 1, 2 ], [ 3, 4 ] ] );|
  LieObject( [ [ 1, 2 ], [ 3, 4 ] ] )
  !gapprompt@gap>| !gapinput@IsGeneralizedRowVector( liemat );|
  true
\end{Verbatim}
 }

 

\subsection{\textcolor{Chapter }{IsMultiplicativeGeneralizedRowVector}}
\logpage{[ 21, 12, 2 ]}\nobreak
\hyperdef{L}{X7FBCA5B58308C158}{}
{\noindent\textcolor{FuncColor}{$\triangleright$\enspace\texttt{IsMultiplicativeGeneralizedRowVector({\mdseries\slshape list})\index{IsMultiplicativeGeneralizedRowVector@\texttt{IsMultiplicative}\-\texttt{Generalized}\-\texttt{Row}\-\texttt{Vector}}
\label{IsMultiplicativeGeneralizedRowVector}
}\hfill{\scriptsize (Category)}}\\


 For a list \mbox{\texttt{\mdseries\slshape list}}, the value \texttt{true} for \texttt{IsMultiplicativeGeneralizedRowVector} indicates that the multiplicative arithmetic behaviour of \mbox{\texttt{\mdseries\slshape list}} is as defined in \ref{Multiplicative Arithmetic for Lists}, and that the attribute \texttt{NestingDepthM} (\ref{NestingDepthM}) will return a nonzero value when called with \mbox{\texttt{\mdseries\slshape list}}. 

 
\begin{Verbatim}[commandchars=!@|,fontsize=\small,frame=single,label=Example]
  !gapprompt@gap>| !gapinput@IsMultiplicativeGeneralizedRowVector( liemat );|
  false
  !gapprompt@gap>| !gapinput@bas:= CanonicalBasis( FullRowSpace( Rationals, 3 ) );|
  CanonicalBasis( ( Rationals^3 ) )
  !gapprompt@gap>| !gapinput@IsMultiplicativeGeneralizedRowVector( bas );|
  true
\end{Verbatim}
 

 Note that the filters \texttt{IsGeneralizedRowVector} (\ref{IsGeneralizedRowVector}), \texttt{IsMultiplicativeGeneralizedRowVector} do \emph{not} enable default methods for addition or multiplication (cf.{\nobreakspace}\texttt{IsListDefault} (\ref{IsListDefault})). }

 

\subsection{\textcolor{Chapter }{IsListDefault}}
\logpage{[ 21, 12, 3 ]}\nobreak
\hyperdef{L}{X7BAD12E67BFC90DE}{}
{\noindent\textcolor{FuncColor}{$\triangleright$\enspace\texttt{IsListDefault({\mdseries\slshape list})\index{IsListDefault@\texttt{IsListDefault}}
\label{IsListDefault}
}\hfill{\scriptsize (Category)}}\\


 For a list \mbox{\texttt{\mdseries\slshape list}}, \texttt{IsListDefault} indicates that the default methods for arithmetic operations of lists, such as
pointwise addition and multiplication as inner product or matrix product,
shall be applicable to \mbox{\texttt{\mdseries\slshape list}}. 

 \texttt{IsListDefault} implies \texttt{IsGeneralizedRowVector} (\ref{IsGeneralizedRowVector}) and \texttt{IsMultiplicativeGeneralizedRowVector} (\ref{IsMultiplicativeGeneralizedRowVector}). 

 All internally represented lists are in this category, and also all lists in
the representations \texttt{IsGF2VectorRep}, \texttt{Is8BitVectorRep}, \texttt{IsGF2MatrixRep}, and \texttt{Is8BitMatrixRep} (see{\nobreakspace}\ref{Row Vectors over Finite Fields} and \ref{Matrices over Finite Fields}).   Note that the result of an arithmetic operation with lists in \texttt{IsListDefault} will in general be an internally represented list, so most ``wrapped list objects'' will not lie in \texttt{IsListDefault}. 

 
\begin{Verbatim}[commandchars=!@|,fontsize=\small,frame=single,label=Example]
  !gapprompt@gap>| !gapinput@v:= [ 1, 2 ];;  m:= [ v, 2*v ];;|
  !gapprompt@gap>| !gapinput@IsListDefault( v );  IsListDefault( m );|
  true
  true
  !gapprompt@gap>| !gapinput@IsListDefault( bas );  IsListDefault( liemat );|
  true
  false
\end{Verbatim}
 }

 

\subsection{\textcolor{Chapter }{NestingDepthA}}
\logpage{[ 21, 12, 4 ]}\nobreak
\hyperdef{L}{X8428E77B86722D52}{}
{\noindent\textcolor{FuncColor}{$\triangleright$\enspace\texttt{NestingDepthA({\mdseries\slshape obj})\index{NestingDepthA@\texttt{NestingDepthA}}
\label{NestingDepthA}
}\hfill{\scriptsize (attribute)}}\\


 For a \textsf{GAP} object \mbox{\texttt{\mdseries\slshape obj}}, \texttt{NestingDepthA} returns the \emph{additive nesting depth} of \mbox{\texttt{\mdseries\slshape obj}}. This is defined recursively as the integer $0$ if \mbox{\texttt{\mdseries\slshape obj}} is not in \texttt{IsGeneralizedRowVector} (\ref{IsGeneralizedRowVector}), as the integer $1$ if \mbox{\texttt{\mdseries\slshape obj}} is an empty list in \texttt{IsGeneralizedRowVector} (\ref{IsGeneralizedRowVector}), and as $1$ plus the additive nesting depth of the first bound entry in \mbox{\texttt{\mdseries\slshape obj}} otherwise. }

 

\subsection{\textcolor{Chapter }{NestingDepthM}}
\logpage{[ 21, 12, 5 ]}\nobreak
\hyperdef{L}{X84B383B97FD986CD}{}
{\noindent\textcolor{FuncColor}{$\triangleright$\enspace\texttt{NestingDepthM({\mdseries\slshape obj})\index{NestingDepthM@\texttt{NestingDepthM}}
\label{NestingDepthM}
}\hfill{\scriptsize (attribute)}}\\


 For a \textsf{GAP} object \mbox{\texttt{\mdseries\slshape obj}}, \texttt{NestingDepthM} returns the \emph{multiplicative nesting depth} of \mbox{\texttt{\mdseries\slshape obj}}. This is defined recursively as the integer $0$ if \mbox{\texttt{\mdseries\slshape obj}} is not in \texttt{IsMultiplicativeGeneralizedRowVector} (\ref{IsMultiplicativeGeneralizedRowVector}), as the integer $1$ if \mbox{\texttt{\mdseries\slshape obj}} is an empty list in \texttt{IsMultiplicativeGeneralizedRowVector} (\ref{IsMultiplicativeGeneralizedRowVector}), and as $1$ plus the multiplicative nesting depth of the first bound entry in \mbox{\texttt{\mdseries\slshape obj}} otherwise. 
\begin{Verbatim}[commandchars=!@|,fontsize=\small,frame=single,label=Example]
  !gapprompt@gap>| !gapinput@NestingDepthA( v );  NestingDepthM( v );|
  1
  1
  !gapprompt@gap>| !gapinput@NestingDepthA( m );  NestingDepthM( m );|
  2
  2
  !gapprompt@gap>| !gapinput@NestingDepthA( liemat );  NestingDepthM( liemat );|
  2
  0
  !gapprompt@gap>| !gapinput@l1:= [ [ 1, 2 ], 3 ];;  l2:= [ 1, [ 2, 3 ] ];;|
  !gapprompt@gap>| !gapinput@NestingDepthA( l1 );  NestingDepthM( l1 );|
  2
  2
  !gapprompt@gap>| !gapinput@NestingDepthA( l2 );  NestingDepthM( l2 );|
  1
  1
\end{Verbatim}
 }

 }

  
\section{\textcolor{Chapter }{Additive Arithmetic for Lists}}\label{Additive Arithmetic for Lists}
\logpage{[ 21, 13, 0 ]}
\hyperdef{L}{X7E6A1F66781BE923}{}
{
  In this general context, we define the results of additive operations only in
the following situations. For unary operations (zero and additive inverse),
the unique argument must be in \texttt{IsGeneralizedRowVector} (\ref{IsGeneralizedRowVector}); for binary operations (addition and subtraction), at least one argument must
be in \texttt{IsGeneralizedRowVector} (\ref{IsGeneralizedRowVector}), and the other either is not a list or also in \texttt{IsGeneralizedRowVector} (\ref{IsGeneralizedRowVector}). 

 (For non\texttt{\symbol{45}}list \textsf{GAP} objects, defining the results of unary operations is not an issue here, and if
at least one argument is a list not in \texttt{IsGeneralizedRowVector} (\ref{IsGeneralizedRowVector}), it shall be left to this argument whether the result in question is defined
and what it is.) 
\subsection{\textcolor{Chapter }{Zero for lists}}\label{Zero for lists}
\logpage{[ 21, 13, 1 ]}
\hyperdef{L}{X86A85ADC85C451DC}{}
{
  The zero (see{\nobreakspace}\texttt{Zero} (\ref{Zero})) of a list $x$ in \texttt{IsGeneralizedRowVector} (\ref{IsGeneralizedRowVector}) is defined as the list whose entry at position $i$ is the zero of $x[i]$ if this entry is bound, and is unbound otherwise. 

 
\begin{Verbatim}[commandchars=!@|,fontsize=\small,frame=single,label=Example]
  !gapprompt@gap>| !gapinput@Zero( [ 1, 2, 3 ] );  Zero( [ [ 1, 2 ], 3 ] );  Zero( liemat );|
  [ 0, 0, 0 ]
  [ [ 0, 0 ], 0 ]
  LieObject( [ [ 0, 0 ], [ 0, 0 ] ] )
\end{Verbatim}
 }

 
\subsection{\textcolor{Chapter }{AdditiveInverse for lists}}\label{AdditiveInverse for lists}
\logpage{[ 21, 13, 2 ]}
\hyperdef{L}{X7B91CE4D814C2D08}{}
{
  The additive inverse (see{\nobreakspace}\texttt{AdditiveInverse} (\ref{AdditiveInverse})) of a list $x$ in \texttt{IsGeneralizedRowVector} (\ref{IsGeneralizedRowVector}) is defined as the list whose entry at position $i$ is the additive inverse of $x[i]$ if this entry is bound, and is unbound otherwise. 

 
\begin{Verbatim}[commandchars=!@|,fontsize=\small,frame=single,label=Example]
  !gapprompt@gap>| !gapinput@AdditiveInverse( [ 1, 2, 3 ] );  AdditiveInverse( [ [ 1, 2 ], 3 ] );|
  [ -1, -2, -3 ]
  [ [ -1, -2 ], -3 ]
\end{Verbatim}
 }

 
\subsection{\textcolor{Chapter }{Addition of lists}}\label{Addition of lists}
\logpage{[ 21, 13, 3 ]}
\hyperdef{L}{X842D123E7EE5E3DB}{}
{
  \index{addition!list and non-list} If $x$ and $y$ are in \texttt{IsGeneralizedRowVector} (\ref{IsGeneralizedRowVector}) and have the same additive nesting depth (see{\nobreakspace}\texttt{NestingDepthA} (\ref{NestingDepthA})),  the sum $x + y$ is defined \emph{pointwise}, in the sense that the result is a list whose entry at position $i$ is $x[i] + y[i]$ if these entries are bound, is a shallow copy (see{\nobreakspace}\texttt{ShallowCopy} (\ref{ShallowCopy})) of $x[i]$ or $y[i]$ if the other argument is not bound at position $i$, and is unbound if both $x$ and $y$ are unbound at position $i$. 

 If $x$ is in \texttt{IsGeneralizedRowVector} (\ref{IsGeneralizedRowVector}) and $y$ is in \texttt{IsGeneralizedRowVector} (\ref{IsGeneralizedRowVector}) and has lower additive nesting depth, or is neither a list nor a domain, the
sum $x + y$ is defined as a list whose entry at position $i$ is $x[i] + y$ if $x$ is bound at position $i$, and is unbound if not. The equivalent holds in the reversed case, where the
order of the summands is kept, as addition is not always commutative. 

 
\begin{Verbatim}[commandchars=!@|,fontsize=\small,frame=single,label=Example]
  !gapprompt@gap>| !gapinput@1 + [ 1, 2, 3 ];  [ 1, 2, 3 ] + [ 0, 2, 4 ];  [ 1, 2 ] + [ Z(2) ];|
  [ 2, 3, 4 ]
  [ 1, 4, 7 ]
  [ 0*Z(2), 2 ]
  !gapprompt@gap>| !gapinput@l1:= [ 1, , 3, 4 ];;             l2:= [ , 2, 3, 4, 5 ];;|
  !gapprompt@gap>| !gapinput@l3:= [ [ 1, 2 ], , [ 5, 6 ] ];;  l4:= [ , [ 3, 4 ], [ 5, 6 ] ];;|
  !gapprompt@gap>| !gapinput@NestingDepthA( l1 );  NestingDepthA( l2 );|
  1
  1
  !gapprompt@gap>| !gapinput@NestingDepthA( l3 );  NestingDepthA( l4 );|
  2
  2
  !gapprompt@gap>| !gapinput@l1 + l2;|
  [ 1, 2, 6, 8, 5 ]
  !gapprompt@gap>| !gapinput@l1 + l3;|
  [ [ 2, 2, 3, 4 ],, [ 6, 6, 3, 4 ] ]
  !gapprompt@gap>| !gapinput@l2 + l4;|
  [ , [ 3, 6, 3, 4, 5 ], [ 5, 8, 3, 4, 5 ] ]
  !gapprompt@gap>| !gapinput@l3 + l4;|
  [ [ 1, 2 ], [ 3, 4 ], [ 10, 12 ] ]
  !gapprompt@gap>| !gapinput@l1 + [];|
  [ 1,, 3, 4 ]
\end{Verbatim}
 }

 
\subsection{\textcolor{Chapter }{Subtraction of lists}}\label{Subtraction of lists}
\logpage{[ 21, 13, 4 ]}
\hyperdef{L}{X7C3DC8BE78DEECDE}{}
{
  \index{list and non\texttt{\symbol{45}}list!difference} For two \textsf{GAP} objects $x$ and $y$ of which one is in \texttt{IsGeneralizedRowVector} (\ref{IsGeneralizedRowVector}) and the other is also in \texttt{IsGeneralizedRowVector} (\ref{IsGeneralizedRowVector}) or is neither a list nor a domain, $x - y$ is defined as $x + (-y)$. 

 
\begin{Verbatim}[commandchars=!@|,fontsize=\small,frame=single,label=Example]
  !gapprompt@gap>| !gapinput@l1 - l2;|
  [ 1, -2, 0, 0, -5 ]
  !gapprompt@gap>| !gapinput@l1 - l3;|
  [ [ 0, -2, 3, 4 ],, [ -4, -6, 3, 4 ] ]
  !gapprompt@gap>| !gapinput@l2 - l4;|
  [ , [ -3, -2, 3, 4, 5 ], [ -5, -4, 3, 4, 5 ] ]
  !gapprompt@gap>| !gapinput@l3 - l4;|
  [ [ 1, 2 ], [ -3, -4 ], [ 0, 0 ] ]
  !gapprompt@gap>| !gapinput@l1 - [];|
  [ 1,, 3, 4 ]
\end{Verbatim}
 }

 }

  
\section{\textcolor{Chapter }{Multiplicative Arithmetic for Lists}}\label{Multiplicative Arithmetic for Lists}
\logpage{[ 21, 14, 0 ]}
\hyperdef{L}{X782ED7F27D8C7FC1}{}
{
  In this general context, we define the results of multiplicative operations
only in the following situations. For unary operations (one and inverse), the
unique argument must be in \texttt{IsMultiplicativeGeneralizedRowVector} (\ref{IsMultiplicativeGeneralizedRowVector}); for binary operations (multiplication and division), at least one argument
must be in \texttt{IsMultiplicativeGeneralizedRowVector} (\ref{IsMultiplicativeGeneralizedRowVector}), and the other either not a list or also in \texttt{IsMultiplicativeGeneralizedRowVector} (\ref{IsMultiplicativeGeneralizedRowVector}). 

 (For non\texttt{\symbol{45}}list \textsf{GAP} objects, defining the results of unary operations is not an issue here, and if
at least one argument is a list not in \texttt{IsMultiplicativeGeneralizedRowVector} (\ref{IsMultiplicativeGeneralizedRowVector}), it shall be left to this argument whether the result in question is defined
and what it is.) 
\subsection{\textcolor{Chapter }{One for lists}}\label{One for lists}
\logpage{[ 21, 14, 1 ]}
\hyperdef{L}{X79A8A5627FD42FA5}{}
{
  The one (see{\nobreakspace}\texttt{One} (\ref{One})) of a dense list \mbox{\texttt{\mdseries\slshape x}} in \texttt{IsMultiplicativeGeneralizedRowVector} (\ref{IsMultiplicativeGeneralizedRowVector}) such that \mbox{\texttt{\mdseries\slshape x}} has even multiplicative nesting depth and has the same length as each of its
rows is defined as the usual identity matrix on the outer two levels, that is,
an identity matrix of the same dimensions, with diagonal entries \texttt{One( \mbox{\texttt{\mdseries\slshape x}}[1][1] )} and off\texttt{\symbol{45}}diagonal entries \texttt{Zero( \mbox{\texttt{\mdseries\slshape x}}[1][1] )}. 

 
\begin{Verbatim}[commandchars=!@|,fontsize=\small,frame=single,label=Example]
  !gapprompt@gap>| !gapinput@One( [ [ 1, 2 ], [ 3, 4 ] ] );|
  [ [ 1, 0 ], [ 0, 1 ] ]
  !gapprompt@gap>| !gapinput@One( [ [ [ [ 1 ] ], [ [ 2 ] ] ], [ [ [ 3 ] ], [ [ 4 ] ] ] ] );|
  [ [ [ [ 1 ] ], [ [ 0 ] ] ], [ [ [ 0 ] ], [ [ 1 ] ] ] ]
\end{Verbatim}
 }

 
\subsection{\textcolor{Chapter }{Inverse for lists}}\label{Inverse for lists}
\logpage{[ 21, 14, 2 ]}
\hyperdef{L}{X78C6C1E2849D303A}{}
{
  The inverse (see{\nobreakspace}\texttt{Inverse} (\ref{Inverse})) of an invertible square table \mbox{\texttt{\mdseries\slshape x}} in \texttt{IsMultiplicativeGeneralizedRowVector} (\ref{IsMultiplicativeGeneralizedRowVector}) whose entries lie in a common field is defined as the usual inverse $y$, i.e., a square matrix over the same field such that $\mbox{\texttt{\mdseries\slshape x}} y$ and $y \mbox{\texttt{\mdseries\slshape x}}$ is equal to \texttt{One( \mbox{\texttt{\mdseries\slshape x}} )}. 

 
\begin{Verbatim}[commandchars=!@|,fontsize=\small,frame=single,label=Example]
  !gapprompt@gap>| !gapinput@Inverse( [ [ 1, 2 ], [ 3, 4 ] ] );|
  [ [ -2, 1 ], [ 3/2, -1/2 ] ]
\end{Verbatim}
 }

 
\subsection{\textcolor{Chapter }{Multiplication of lists}}\label{Multiplication of lists}
\logpage{[ 21, 14, 3 ]}
\hyperdef{L}{X84FDB95179BFE4CD}{}
{
  \index{list and non\texttt{\symbol{45}}list!product} There are three possible computations that might be triggered by a
multiplication involving a list in \texttt{IsMultiplicativeGeneralizedRowVector} (\ref{IsMultiplicativeGeneralizedRowVector}). Namely, $x * y$ might be 
\begin{description}
\item[{(I)}]  the inner product $x[1] * y[1] + x[2] * y[2] + \cdots + x[n] * y[n]$, where summands are omitted for which the entry in $x$ or $y$ is unbound (if this leaves no summand then the multiplication is an error), or 
\item[{(L)}]  the left scalar multiple, i.e., a list whose entry at position $i$ is $x * y[i]$ if $y$ is bound at position $i$, and is unbound if not, or 
\item[{(R)}]  the right scalar multiple, i.e., a list whose entry at position $i$ is $x[i] * y$ if $x$ is bound at position $i$, and is unbound if not. 
\end{description}
 

 Our aim is to generalize the basic arithmetic of simple row vectors and
matrices, so we first summarize the situations that shall be covered. 

 \begin{center}
\begin{tabular}{c|ccc} &
scl&
vec&
mat\\
\hline
scl&
 &
(L)&
(L)\\
vec&
(R)&
(I)&
(I)\\
mat&
(R)&
(R)&
(R)\\
\end{tabular}\\[2mm]
\end{center}

 

 This means for example that the product of a scalar (scl) with a vector (vec)
or a matrix (mat) is computed according to (L). Note that this is asymmetric. 

 Now we can state the general multiplication rules. 

 If exactly one argument is in \texttt{IsMultiplicativeGeneralizedRowVector} (\ref{IsMultiplicativeGeneralizedRowVector}) then we regard the other argument (which is then neither a list nor a domain)
as a scalar, and specify result (L) or (R), depending on ordering. 

 In the remaining cases, both $x$ and $y$ are in \texttt{IsMultiplicativeGeneralizedRowVector} (\ref{IsMultiplicativeGeneralizedRowVector}), and we distinguish the possibilities by their multiplicative nesting depths.
An argument with \emph{odd} multiplicative nesting depth is regarded as a vector, and an argument with \emph{even} multiplicative nesting depth is regarded as a scalar or a matrix. 

 So if both arguments have odd multiplicative nesting depth, we specify result
(I). 

 If exactly one argument has odd nesting depth, the other is treated as a
scalar if it has lower multiplicative nesting depth, and as a matrix
otherwise. In the former case, we specify result (L) or (R), depending on
ordering; in the latter case, we specify result (L) or (I), depending on
ordering. 

 We are left with the case that each argument has even multiplicative nesting
depth.  If the two depths are equal, we treat the computation as a matrix product, and
specify result (R). Otherwise, we treat the less deeply nested argument as a
scalar and the other as a matrix, and specify result (L) or (R), depending on
ordering. 

 
\begin{Verbatim}[commandchars=!@|,fontsize=\small,frame=single,label=Example]
  !gapprompt@gap>| !gapinput@[ (), (2,3), (1,2), (1,2,3), (1,3,2), (1,3) ] * (1,4);|
  [ (1,4), (1,4)(2,3), (1,2,4), (1,2,3,4), (1,3,2,4), (1,3,4) ]
  !gapprompt@gap>| !gapinput@[ 1, 2, , 4 ] * 2;|
  [ 2, 4,, 8 ]
  !gapprompt@gap>| !gapinput@[ 1, 2, 3 ] * [ 1, 3, 5, 7 ];|
  22
  !gapprompt@gap>| !gapinput@m:= [ [ 1, 2 ], 3 ];;  m * m;|
  [ [ 7, 8 ], [ [ 3, 6 ], 9 ] ]
  !gapprompt@gap>| !gapinput@m * m = [ m[1] * m, m[2] * m ];|
  true
  !gapprompt@gap>| !gapinput@n:= [ 1, [ 2, 3 ] ];;  n * n;|
  14
  !gapprompt@gap>| !gapinput@n * n = n[1] * n[1] + n[2] * n[2];|
  true
\end{Verbatim}
 }

 
\subsection{\textcolor{Chapter }{Division of lists}}\label{Division of lists}
\logpage{[ 21, 14, 4 ]}
\hyperdef{L}{X82EA2A5B786181C7}{}
{
  \index{list and non\texttt{\symbol{45}}list!quotient} For two \textsf{GAP} objects $x$ and $y$ of which one is in \texttt{IsMultiplicativeGeneralizedRowVector} (\ref{IsMultiplicativeGeneralizedRowVector}) and the other is also in \texttt{IsMultiplicativeGeneralizedRowVector} (\ref{IsMultiplicativeGeneralizedRowVector}) or is neither a list nor a domain, $x / y$ is defined as $x * y^{{-1}}$. 

 
\begin{Verbatim}[commandchars=!@|,fontsize=\small,frame=single,label=Example]
  !gapprompt@gap>| !gapinput@[ 1, 2, 3 ] / 2;  [ 1, 2 ] / [ [ 1, 2 ], [ 3, 4 ] ];|
  [ 1/2, 1, 3/2 ]
  [ 1, 0 ]
\end{Verbatim}
 }

 
\subsection{\textcolor{Chapter }{mod for lists}}\label{mod for lists}
\logpage{[ 21, 14, 5 ]}
\hyperdef{L}{X7A0FD70C80B95C00}{}
{
  \index{list and non\texttt{\symbol{45}}list!mod} \index{mod!lists} If \mbox{\texttt{\mdseries\slshape x}} and \mbox{\texttt{\mdseries\slshape y}} are in \texttt{IsMultiplicativeGeneralizedRowVector} (\ref{IsMultiplicativeGeneralizedRowVector}) and have the same multiplicative nesting depth (see{\nobreakspace}\texttt{NestingDepthM} (\ref{NestingDepthM})),  \texttt{\mbox{\texttt{\mdseries\slshape x}} mod \mbox{\texttt{\mdseries\slshape y}}} is defined \emph{pointwise}, in the sense that the result is a list whose entry at position $i$ is \texttt{\mbox{\texttt{\mdseries\slshape x}}[i] mod \mbox{\texttt{\mdseries\slshape y}}[i]} if these entries are bound, is a shallow copy (see{\nobreakspace}\texttt{ShallowCopy} (\ref{ShallowCopy})) of $x[i]$ or $y[i]$ if the other argument is not bound at position $i$, and is unbound if both $x$ and $y$ are unbound at position $i$. 

 If $x$ is in \texttt{IsMultiplicativeGeneralizedRowVector} (\ref{IsMultiplicativeGeneralizedRowVector}) and $y$ is in \texttt{IsMultiplicativeGeneralizedRowVector} (\ref{IsMultiplicativeGeneralizedRowVector}) and has lower multiplicative nesting depth or is neither a list nor a domain, \texttt{\mbox{\texttt{\mdseries\slshape x}} mod \mbox{\texttt{\mdseries\slshape y}}} is defined as a list whose entry at position $i$ is \texttt{\mbox{\texttt{\mdseries\slshape x}}[i] mod \mbox{\texttt{\mdseries\slshape y}}} if \mbox{\texttt{\mdseries\slshape x}} is bound at position $i$, and is unbound if not. The equivalent holds in the reversed case, where the
order of the arguments is kept. 

 
\begin{Verbatim}[commandchars=!@|,fontsize=\small,frame=single,label=Example]
  !gapprompt@gap>| !gapinput@4711 mod [ 2, 3,, 5, 7 ];|
  [ 1, 1,, 1, 0 ]
  !gapprompt@gap>| !gapinput@[ 2, 3, 4, 5, 6 ] mod 3;|
  [ 2, 0, 1, 2, 0 ]
  !gapprompt@gap>| !gapinput@[ 10, 12, 14, 16 ] mod [ 3, 5, 7 ];|
  [ 1, 2, 0, 16 ]
\end{Verbatim}
 }

 
\subsection{\textcolor{Chapter }{Left quotients of lists}}\label{Left quotients of lists}
\logpage{[ 21, 14, 6 ]}
\hyperdef{L}{X84BB2DFB8432A1A4}{}
{
  \index{list and non\texttt{\symbol{45}}list!left quotient} For two \textsf{GAP} objects $x$ and $y$ of which one is in \texttt{IsMultiplicativeGeneralizedRowVector} (\ref{IsMultiplicativeGeneralizedRowVector}) and the other is also in \texttt{IsMultiplicativeGeneralizedRowVector} (\ref{IsMultiplicativeGeneralizedRowVector}) or is neither a list nor a domain, \texttt{LeftQuotient( \mbox{\texttt{\mdseries\slshape x}}, \mbox{\texttt{\mdseries\slshape y}} )} is defined as $x^{{-1}} * y$. 

 
\begin{Verbatim}[commandchars=!@|,fontsize=\small,frame=single,label=Example]
  !gapprompt@gap>| !gapinput@LeftQuotient( [ [ 1, 2 ], [ 3, 4 ] ], [ 1, 2 ] );|
  [ 0, 1/2 ]
\end{Verbatim}
 }

 }

  
\section{\textcolor{Chapter }{Mutability Status and List Arithmetic}}\label{Mutability Status and List Arithmetic}
\logpage{[ 21, 15, 0 ]}
\hyperdef{L}{X8676EFE67972FD06}{}
{
  Many results of arithmetic operations, when applied to lists, are again lists,
and it is of interest whether their entries are mutable or not (if
applicable). Note that the mutability status of the result itself is already
defined by the general rule for any result of an arithmetic operation, not
only for lists (see{\nobreakspace}\ref{Mutability and Copyability}). 

 However, we do \emph{not} define exactly the mutability status for each element on each level of a
nested list returned by an arithmetic operation. (Of course it would be
possible to define this recursively, but since the methods used are in general
not recursive, in particular for efficient multiplication of compressed
matrices, such a general definition would be a burden in these cases.) Instead
we consider, for a list $x$ in \texttt{IsGeneralizedRowVector} (\ref{IsGeneralizedRowVector}), the sequence $x = x_1, x_2, \ldots x_n$ where $x_{{i+1}}$ is the first bound entry in $x_i$ if exists (that is, if $x_i$ is a nonempty list), and $n$ is the largest $i$ such that $x_i$ lies in \texttt{IsGeneralizedRowVector} (\ref{IsGeneralizedRowVector}). The \emph{immutability level} of $x$ is defined as infinity if $x$ is immutable, and otherwise the number of $x_i$ which are immutable. (So the immutability level of a mutable empty list is $0$.) 

 Thus a fully mutable matrix has immutability level $0$, and a mutable matrix with immutable first row has immutability level $1$ (independent of the mutability of other rows). 

 The immutability level of the result of any of the binary operations discussed
here is the minimum of the immutability levels of the arguments, provided that
objects of the required mutability status exist in \textsf{GAP}.  

 Moreover, the results have a ``homogeneous'' mutability status, that is, if the first bound entry at nesting depth $i$ is immutable (mutable) then all entries at nesting depth $i$ are immutable (mutable, provided that a mutable version of this entry exists
in \textsf{GAP}). 

 Thus the sum of two mutable matrices whose first rows are mutable is a matrix
all of whose rows are mutable, and the product of two matrices whose first
rows are immutable is a matrix all of whose rows are immutable, independent of
the mutability status of the other rows of the arguments.  

 For example, the sum of a matrix (mutable or immutable, i.e., of immutability
level one of $0$, $1$, or $2$) and a mutable row vector (i.e., immutability level $0$) is a fully mutable matrix. The product of two mutable row vectors of
integers is an integer, and since \textsf{GAP} does not support mutable integers, the result is immutable. 

 For unary arithmetic operations, there are three operations available, an
attribute that returns an immutable result (\texttt{Zero} (\ref{Zero}), \texttt{AdditiveInverse} (\ref{AdditiveInverse}), \texttt{One} (\ref{One}), \texttt{Inverse} (\ref{Inverse})), an operation that returns a result that is mutable  (\texttt{ZeroOp} (\ref{ZeroOp}), \texttt{AdditiveInverseOp} (\ref{AdditiveInverseOp}), \texttt{OneOp} (\ref{OneOp}), \texttt{InverseOp} (\ref{InverseOp})), and an operation whose result has the same immutability level as the
argument (\texttt{ZeroSameMutability} (\ref{ZeroSameMutability}), \texttt{AdditiveInverseSameMutability} (\ref{AdditiveInverseSameMutability}), \texttt{OneSameMutability} (\ref{OneSameMutability}), \texttt{InverseSameMutability} (\ref{InverseSameMutability})). The last kind of operations is equivalent to the corresponding infix
operations \texttt{0 * \mbox{\texttt{\mdseries\slshape list}}}, \texttt{\texttt{\symbol{45}} \mbox{\texttt{\mdseries\slshape list}}}, \texttt{\mbox{\texttt{\mdseries\slshape list}}\texttt{\symbol{94}}0}, and \texttt{\mbox{\texttt{\mdseries\slshape list}}\texttt{\symbol{94}}\texttt{\symbol{45}}1}. (This holds not only for lists, see{\nobreakspace}\ref{Mutability and Copyability}.) 

 
\begin{Verbatim}[commandchars=!@|,fontsize=\small,frame=single,label=Example]
  !gapprompt@gap>| !gapinput@IsMutable( l1 );  IsMutable( 2 * Immutable( [ 1, 2, 3 ] ) );|
  true
  false
  !gapprompt@gap>| !gapinput@IsMutable( l2 );  IsMutable( l3 );|
  true
  true
\end{Verbatim}
 

 An example motivating the mutability rule is the use of syntactic constructs
such as \texttt{\mbox{\texttt{\mdseries\slshape obj}} * \mbox{\texttt{\mdseries\slshape list}}} and \texttt{\texttt{\symbol{45}} \mbox{\texttt{\mdseries\slshape list}}} as an elegant and efficient way to create mutable lists needed for further
manipulations from mutable lists. In particular one can construct a mutable
zero vector of length $n$ by \texttt{0 * [ 1 .. }$n$\texttt{ ]}. The latter can be done also using \texttt{ListWithIdenticalEntries} (\ref{ListWithIdenticalEntries}). 

\subsection{\textcolor{Chapter }{ListWithIdenticalEntries}}
\logpage{[ 21, 15, 1 ]}\nobreak
\hyperdef{L}{X80FDB1457FF582E7}{}
{\noindent\textcolor{FuncColor}{$\triangleright$\enspace\texttt{ListWithIdenticalEntries({\mdseries\slshape n, obj})\index{ListWithIdenticalEntries@\texttt{ListWithIdenticalEntries}}
\label{ListWithIdenticalEntries}
}\hfill{\scriptsize (function)}}\\


 is a list \mbox{\texttt{\mdseries\slshape list}} of length \mbox{\texttt{\mdseries\slshape n}} that has the object \mbox{\texttt{\mdseries\slshape obj}} stored at each of the positions from \texttt{1} to \mbox{\texttt{\mdseries\slshape n}}. Note that all elements of \mbox{\texttt{\mdseries\slshape lists}} are identical, see{\nobreakspace}\ref{Identical Lists}. 

 
\begin{Verbatim}[commandchars=!@|,fontsize=\small,frame=single,label=Example]
  !gapprompt@gap>| !gapinput@ListWithIdenticalEntries( 10, 0 );|
  [ 0, 0, 0, 0, 0, 0, 0, 0, 0, 0 ]
\end{Verbatim}
 }

 }

  
\section{\textcolor{Chapter }{Finding Positions in Lists}}\label{Finding Positions in Lists}
\logpage{[ 21, 16, 0 ]}
\hyperdef{L}{X8196FD4779BCCA0C}{}
{
  

\subsection{\textcolor{Chapter }{Position}}
\logpage{[ 21, 16, 1 ]}\nobreak
\hyperdef{L}{X79975EC6783B4293}{}
{\noindent\textcolor{FuncColor}{$\triangleright$\enspace\texttt{Position({\mdseries\slshape list, obj[, from]})\index{Position@\texttt{Position}}
\label{Position}
}\hfill{\scriptsize (operation)}}\\


 returns the position of the first occurrence \mbox{\texttt{\mdseries\slshape obj}} in \mbox{\texttt{\mdseries\slshape list}}, or \texttt{fail} if \mbox{\texttt{\mdseries\slshape obj}} is not contained in \mbox{\texttt{\mdseries\slshape list}}. If a starting index \mbox{\texttt{\mdseries\slshape from}} is given, it returns the position of the first occurrence starting the search \emph{after} position \mbox{\texttt{\mdseries\slshape from}}. 

 Each call to the two argument version is translated into a call of the three
argument version, with third argument the integer zero \texttt{0}. (Methods for the two argument version must be installed as methods for the
version with three arguments, the third being described by \texttt{IsZeroCyc}.) 

 
\begin{Verbatim}[commandchars=!@|,fontsize=\small,frame=single,label=Example]
  !gapprompt@gap>| !gapinput@Position( [ 2, 2, 1, 3 ], 1 );|
  3
  !gapprompt@gap>| !gapinput@Position( [ 2, 1, 1, 3 ], 1 );|
  2
  !gapprompt@gap>| !gapinput@Position( [ 2, 1, 1, 3 ], 1, 2 );|
  3
  !gapprompt@gap>| !gapinput@Position( [ 2, 1, 1, 3 ], 1, 3 );|
  fail
\end{Verbatim}
 }

 

\subsection{\textcolor{Chapter }{Positions}}
\logpage{[ 21, 16, 2 ]}\nobreak
\hyperdef{L}{X7FA9648883AE1B88}{}
{\noindent\textcolor{FuncColor}{$\triangleright$\enspace\texttt{Positions({\mdseries\slshape list, obj})\index{Positions@\texttt{Positions}}
\label{Positions}
}\hfill{\scriptsize (function)}}\\


 returns the set of positions of \emph{all} occurrences of \mbox{\texttt{\mdseries\slshape obj}} in \mbox{\texttt{\mdseries\slshape list}}. 

 Developers who wish to adapt this for custom list types need to install
suitable methods for the operation \texttt{PositionsOp}. \index{PositionsOp@\texttt{PositionsOp}} 

 
\begin{Verbatim}[commandchars=!@|,fontsize=\small,frame=single,label=Example]
  !gapprompt@gap>| !gapinput@Positions([1,2,1,2,3,2,2],2);|
  [ 2, 4, 6, 7 ]
  !gapprompt@gap>| !gapinput@Positions([1,2,1,2,3,2,2],4);|
  [  ]
\end{Verbatim}
 }

 

\subsection{\textcolor{Chapter }{PositionCanonical}}
\logpage{[ 21, 16, 3 ]}\nobreak
\hyperdef{L}{X7B4B10AE81602D4E}{}
{\noindent\textcolor{FuncColor}{$\triangleright$\enspace\texttt{PositionCanonical({\mdseries\slshape list, obj})\index{PositionCanonical@\texttt{PositionCanonical}}
\label{PositionCanonical}
}\hfill{\scriptsize (operation)}}\\


 returns the position of the canonical associate of \mbox{\texttt{\mdseries\slshape obj}} in \mbox{\texttt{\mdseries\slshape list}}. The definition of this associate depends on \mbox{\texttt{\mdseries\slshape list}}. For internally represented lists it is defined as the element itself (and \texttt{PositionCanonical} thus defaults to \texttt{Position} (\ref{Position}), but for example for certain enumerators (see{\nobreakspace}\ref{Enumerators}) other canonical associates can be defined. 

 For example \texttt{RightTransversal} (\ref{RightTransversal}) defines the canonical associate to be the element in the transversal defining
the same coset of a subgroup in a group. 

 
\begin{Verbatim}[commandchars=!@|,fontsize=\small,frame=single,label=Example]
  !gapprompt@gap>| !gapinput@g:=Group((1,2,3,4),(1,2));;u:=Subgroup(g,[(1,2)(3,4),(1,3)(2,4)]);;|
  !gapprompt@gap>| !gapinput@rt:=RightTransversal(g,u);;AsList(rt);|
  [ (), (3,4), (2,3), (2,3,4), (2,4,3), (2,4) ]
  !gapprompt@gap>| !gapinput@Position(rt,(1,2));|
  fail
  !gapprompt@gap>| !gapinput@PositionCanonical(rt,(1,2));|
  2
\end{Verbatim}
 }

 

\subsection{\textcolor{Chapter }{PositionNthOccurrence}}
\logpage{[ 21, 16, 4 ]}\nobreak
\hyperdef{L}{X7D2B25B484591506}{}
{\noindent\textcolor{FuncColor}{$\triangleright$\enspace\texttt{PositionNthOccurrence({\mdseries\slshape list, obj, n})\index{PositionNthOccurrence@\texttt{PositionNthOccurrence}}
\label{PositionNthOccurrence}
}\hfill{\scriptsize (operation)}}\\


 returns the position of the \mbox{\texttt{\mdseries\slshape n}}\texttt{\symbol{45}}th occurrence of \mbox{\texttt{\mdseries\slshape obj}} in \mbox{\texttt{\mdseries\slshape list}} and returns \texttt{fail} if \mbox{\texttt{\mdseries\slshape obj}} does not occur \mbox{\texttt{\mdseries\slshape n}} times. 

 
\begin{Verbatim}[commandchars=!@|,fontsize=\small,frame=single,label=Example]
  !gapprompt@gap>| !gapinput@PositionNthOccurrence([1,2,3,2,4,2,1],1,1);|
  1
  !gapprompt@gap>| !gapinput@PositionNthOccurrence([1,2,3,2,4,2,1],1,2);|
  7
  !gapprompt@gap>| !gapinput@PositionNthOccurrence([1,2,3,2,4,2,1],2,3);|
  6
  !gapprompt@gap>| !gapinput@PositionNthOccurrence([1,2,3,2,4,2,1],2,4);|
  fail
\end{Verbatim}
 }

 

\subsection{\textcolor{Chapter }{PositionSorted}}
\logpage{[ 21, 16, 5 ]}\nobreak
\hyperdef{L}{X7A122E848464E534}{}
{\noindent\textcolor{FuncColor}{$\triangleright$\enspace\texttt{PositionSorted({\mdseries\slshape list, elm[, func]})\index{PositionSorted@\texttt{PositionSorted}}
\label{PositionSorted}
}\hfill{\scriptsize (function)}}\\


 Called with two arguments, \texttt{PositionSorted} returns the position of the element \mbox{\texttt{\mdseries\slshape elm}} in the sorted list \mbox{\texttt{\mdseries\slshape list}}. 

 Called with three arguments, \texttt{PositionSorted} returns the position of the element \mbox{\texttt{\mdseries\slshape elm}} in the list \mbox{\texttt{\mdseries\slshape list}}, which must be sorted with respect to \mbox{\texttt{\mdseries\slshape func}}. \mbox{\texttt{\mdseries\slshape func}} must be a function of two arguments that returns \texttt{true} if the first argument is less than the second argument, and \texttt{false} otherwise. 

 \texttt{PositionSorted} returns \mbox{\texttt{\mdseries\slshape pos}} such that $\mbox{\texttt{\mdseries\slshape list}}[\mbox{\texttt{\mdseries\slshape pos}}-1] < \mbox{\texttt{\mdseries\slshape elm}} \leq \mbox{\texttt{\mdseries\slshape list}}[\mbox{\texttt{\mdseries\slshape pos}}]$ holds. That means, if \mbox{\texttt{\mdseries\slshape elm}} appears once in \mbox{\texttt{\mdseries\slshape list}}, its position is returned. If \mbox{\texttt{\mdseries\slshape elm}} appears several times in \mbox{\texttt{\mdseries\slshape list}}, the position of the first occurrence is returned. If \mbox{\texttt{\mdseries\slshape elm}} is not an element of \mbox{\texttt{\mdseries\slshape list}}, the index where \mbox{\texttt{\mdseries\slshape elm}} must be inserted to keep the list sorted is returned. 

 \texttt{PositionSorted} uses binary search, whereas \texttt{Position} (\ref{Position}) can in general use only linear search, see the remark at the beginning
of{\nobreakspace}\ref{Sorted Lists and Sets}. For sorting lists, see{\nobreakspace}\ref{Sorting Lists}, for testing whether a list is sorted, see{\nobreakspace}\texttt{IsSortedList} (\ref{IsSortedList}) and \texttt{IsSSortedList} (\ref{IsSSortedList}). 

 Developers who wish to adapt this for custom list types need to install
suitable methods for the operation \texttt{PositionSortedOp}. \index{PositionSortedOp@\texttt{PositionSortedOp}} 

 
\begin{Verbatim}[commandchars=!@|,fontsize=\small,frame=single,label=Example]
  !gapprompt@gap>| !gapinput@PositionSorted( [1,4,5,5,6,7], 0 );|
  1
  !gapprompt@gap>| !gapinput@PositionSorted( [1,4,5,5,6,7], 2 );|
  2
  !gapprompt@gap>| !gapinput@PositionSorted( [1,4,5,5,6,7], 4 );|
  2
  !gapprompt@gap>| !gapinput@PositionSorted( [1,4,5,5,6,7], 5 );|
  3
  !gapprompt@gap>| !gapinput@PositionSorted( [1,4,5,5,6,7], 8 );|
  7
\end{Verbatim}
 }

 

\subsection{\textcolor{Chapter }{PositionSortedBy}}
\logpage{[ 21, 16, 6 ]}\nobreak
\hyperdef{L}{X820BA44D85930EBF}{}
{\noindent\textcolor{FuncColor}{$\triangleright$\enspace\texttt{PositionSortedBy({\mdseries\slshape list, val, func})\index{PositionSortedBy@\texttt{PositionSortedBy}}
\label{PositionSortedBy}
}\hfill{\scriptsize (function)}}\\


 This function returns the same value that would be returned by \texttt{PositionSorted(List(list, func), val)}, but computes it in a more efficient way. 

 To be more precise, \mbox{\texttt{\mdseries\slshape func}} must be a function on one argument which returns values that can be compared
to \mbox{\texttt{\mdseries\slshape val}}, and \mbox{\texttt{\mdseries\slshape list}} must be a list for which \texttt{func(list[i]) {\textless}= func(list[i+1])} holds for all relevant \mbox{\texttt{\mdseries\slshape i}}. This property is not verified, and if the input violates it, then the result
is undefined. 

 \texttt{PositionSortedBy} returns \mbox{\texttt{\mdseries\slshape pos}} such that $\mbox{\texttt{\mdseries\slshape func}}(\mbox{\texttt{\mdseries\slshape list}}[\mbox{\texttt{\mdseries\slshape pos}}-1]) < \mbox{\texttt{\mdseries\slshape val}} \leq \mbox{\texttt{\mdseries\slshape func}}(\mbox{\texttt{\mdseries\slshape list}}[\mbox{\texttt{\mdseries\slshape pos}}])$ holds. That means, if there are elements \texttt{elm} in \mbox{\texttt{\mdseries\slshape list}} for which $\mbox{\texttt{\mdseries\slshape func}}(elm) = \mbox{\texttt{\mdseries\slshape val}}$ holds, then the position of the first such element is returned. If no element
of \mbox{\texttt{\mdseries\slshape list}} satisfies this condition, then the lowest index where an element \mbox{\texttt{\mdseries\slshape elm}} satisfying $\mbox{\texttt{\mdseries\slshape func}}(elm) = \mbox{\texttt{\mdseries\slshape val}}$ must be inserted to preserve the property \texttt{func(list[i]) {\textless}= func(list[i+1])} is returned. 

 \texttt{PositionSortedBy} uses binary search. Each \texttt{func(list[i])} is computed at most once. 

 Developers who wish to adapt this for custom list types need to install
suitable methods for the operation \texttt{PositionSortedByOp}. \index{PositionSortedByOp@\texttt{PositionSortedByOp}} 

 
\begin{Verbatim}[commandchars=!@|,fontsize=\small,frame=single,label=Example]
  !gapprompt@gap>| !gapinput@PositionSortedBy( [ "", "ab", ], -1, Length );|
  1
  !gapprompt@gap>| !gapinput@PositionSortedBy( [ "", "ab", ], 0, Length );|
  1
  !gapprompt@gap>| !gapinput@PositionSortedBy( [ "", "ab", ], 1, Length );|
  2
  !gapprompt@gap>| !gapinput@PositionSortedBy( [ "", "ab", ], 2, Length );|
  2
  !gapprompt@gap>| !gapinput@PositionSortedBy( [ "", "ab", ], 3, Length );|
  3
\end{Verbatim}
 }

 

\subsection{\textcolor{Chapter }{PositionSet}}
\logpage{[ 21, 16, 7 ]}\nobreak
\hyperdef{L}{X78BFE9D78347C0DA}{}
{\noindent\textcolor{FuncColor}{$\triangleright$\enspace\texttt{PositionSet({\mdseries\slshape list, obj[, func]})\index{PositionSet@\texttt{PositionSet}}
\label{PositionSet}
}\hfill{\scriptsize (function)}}\\


 \texttt{PositionSet} is a slight variation of \texttt{PositionSorted} (\ref{PositionSorted}). The only difference to \texttt{PositionSorted} (\ref{PositionSorted}) is that \texttt{PositionSet} returns \texttt{fail} if \mbox{\texttt{\mdseries\slshape obj}} is not in \mbox{\texttt{\mdseries\slshape list}}. 

 
\begin{Verbatim}[commandchars=!@|,fontsize=\small,frame=single,label=Example]
  !gapprompt@gap>| !gapinput@PositionSet( [1,4,5,5,6,7], 0 );|
  fail
  !gapprompt@gap>| !gapinput@PositionSet( [1,4,5,5,6,7], 2 );|
  fail
  !gapprompt@gap>| !gapinput@PositionSet( [1,4,5,5,6,7], 4 );|
  2
  !gapprompt@gap>| !gapinput@PositionSet( [1,4,5,5,6,7], 5 );|
  3
  !gapprompt@gap>| !gapinput@PositionSet( [1,4,5,5,6,7], 8 );|
  fail
\end{Verbatim}
 }

 

\subsection{\textcolor{Chapter }{PositionMaximum}}
\logpage{[ 21, 16, 8 ]}\nobreak
\hyperdef{L}{X7FD9C1D37F300206}{}
{\noindent\textcolor{FuncColor}{$\triangleright$\enspace\texttt{PositionMaximum({\mdseries\slshape list[, func]})\index{PositionMaximum@\texttt{PositionMaximum}}
\label{PositionMaximum}
}\hfill{\scriptsize (function)}}\\
\noindent\textcolor{FuncColor}{$\triangleright$\enspace\texttt{PositionMinimum({\mdseries\slshape list[, func]})\index{PositionMinimum@\texttt{PositionMinimum}}
\label{PositionMinimum}
}\hfill{\scriptsize (function)}}\\


 returns the position of maximum (with \texttt{PositionMaximum}) or minimum (with \texttt{PositionMinimum}) entry in the list \mbox{\texttt{\mdseries\slshape list}}. If a second argument \mbox{\texttt{\mdseries\slshape func}} is passed, then return instead the position of the largest/smallest entry in \texttt{List( \mbox{\texttt{\mdseries\slshape list}} , \mbox{\texttt{\mdseries\slshape func}} )}. If several entries of the list are equal to the maximum/minimum, the first
such position is returned. 

 
\begin{Verbatim}[commandchars=!@|,fontsize=\small,frame=single,label=Example]
  !gapprompt@gap>| !gapinput@PositionMaximum( [2,4,-6,2,4] );|
  2
  !gapprompt@gap>| !gapinput@PositionMaximum( [2,4,-6,2,4], x -> -x);|
  3
  !gapprompt@gap>| !gapinput@PositionMinimum( [2,4,-6,2,4] );|
  3
  !gapprompt@gap>| !gapinput@PositionMinimum( [2,4,-6,2,4], x -> -x);|
  2
\end{Verbatim}
 

 \texttt{Maximum} (\ref{Maximum:for various objects}) and \texttt{Minimum} (\ref{Minimum:for various objects}) allow you to find the maximum or minimum element of a list directly. }

 

\subsection{\textcolor{Chapter }{PositionProperty}}
\logpage{[ 21, 16, 9 ]}\nobreak
\hyperdef{L}{X7E6C763A82C6153B}{}
{\noindent\textcolor{FuncColor}{$\triangleright$\enspace\texttt{PositionProperty({\mdseries\slshape list, func[, from]})\index{PositionProperty@\texttt{PositionProperty}}
\label{PositionProperty}
}\hfill{\scriptsize (operation)}}\\


 returns the position of the first entry in the list \mbox{\texttt{\mdseries\slshape list}} for which the property tester function \mbox{\texttt{\mdseries\slshape func}} returns \texttt{true}, or \texttt{fail} if no such entry exists. If a starting index \mbox{\texttt{\mdseries\slshape from}} is given, it returns the position of the first entry satisfying \mbox{\texttt{\mdseries\slshape func}}, starting the search \emph{after} position \mbox{\texttt{\mdseries\slshape from}}. 

 
\begin{Verbatim}[commandchars=!@|,fontsize=\small,frame=single,label=Example]
  !gapprompt@gap>| !gapinput@PositionProperty( [10^7..10^8], IsPrime );|
  20
  !gapprompt@gap>| !gapinput@PositionProperty( [10^5..10^6],|
  !gapprompt@>| !gapinput@       n -> not IsPrime(n) and IsPrimePowerInt(n) );|
  490
\end{Verbatim}
 

 \texttt{First} (\ref{First}) allows you to extract the first element of a list that satisfies a certain
property. }

 

\subsection{\textcolor{Chapter }{PositionsProperty}}
\logpage{[ 21, 16, 10 ]}\nobreak
\hyperdef{L}{X7DA94D278304EC3D}{}
{\noindent\textcolor{FuncColor}{$\triangleright$\enspace\texttt{PositionsProperty({\mdseries\slshape list, func})\index{PositionsProperty@\texttt{PositionsProperty}}
\label{PositionsProperty}
}\hfill{\scriptsize (operation)}}\\


 returns the set of all those positions in the list \mbox{\texttt{\mdseries\slshape list}} which are bound and for which the property tester function \mbox{\texttt{\mdseries\slshape func}} returns \texttt{true}. 

 
\begin{Verbatim}[commandchars=!@|,fontsize=\small,frame=single,label=Example]
  !gapprompt@gap>| !gapinput@l:= [ -5 .. 5 ];;|
  !gapprompt@gap>| !gapinput@PositionsProperty( l, IsPosInt );|
  [ 7, 8, 9, 10, 11 ]
  !gapprompt@gap>| !gapinput@PositionsProperty( l, IsPrimeInt );|
  [ 1, 3, 4, 8, 9, 11 ]
\end{Verbatim}
 

 \texttt{PositionProperty} (\ref{PositionProperty}) allows you to extract the position of the first element in a list that
satisfies a certain property. }

 

\subsection{\textcolor{Chapter }{PositionBound}}
\logpage{[ 21, 16, 11 ]}\nobreak
\hyperdef{L}{X86C9E5C3863B3C03}{}
{\noindent\textcolor{FuncColor}{$\triangleright$\enspace\texttt{PositionBound({\mdseries\slshape list})\index{PositionBound@\texttt{PositionBound}}
\label{PositionBound}
}\hfill{\scriptsize (operation)}}\\


 returns the first bound position of the list \mbox{\texttt{\mdseries\slshape list}}. For the empty list it returns \texttt{fail}. 

 
\begin{Verbatim}[commandchars=!@|,fontsize=\small,frame=single,label=Example]
  !gapprompt@gap>| !gapinput@PositionBound([1,2,3]);|
  1
  !gapprompt@gap>| !gapinput@PositionBound([,1,2,3]);|
  2
\end{Verbatim}
 }

 

\subsection{\textcolor{Chapter }{PositionsBound}}
\logpage{[ 21, 16, 12 ]}\nobreak
\hyperdef{L}{X819F71047AABEA2F}{}
{\noindent\textcolor{FuncColor}{$\triangleright$\enspace\texttt{PositionsBound({\mdseries\slshape list})\index{PositionsBound@\texttt{PositionsBound}}
\label{PositionsBound}
}\hfill{\scriptsize (function)}}\\


 returns the set of all bound positions in the list \mbox{\texttt{\mdseries\slshape list}}. 

 
\begin{Verbatim}[commandchars=!@|,fontsize=\small,frame=single,label=Example]
  !gapprompt@gap>| !gapinput@PositionsBound([1,2,3]);|
  [ 1 .. 3 ]
  !gapprompt@gap>| !gapinput@PositionsBound([,1,,3]);|
  [ 2, 4 ]
  !gapprompt@gap>| !gapinput@PositionsBound([]);|
  []
\end{Verbatim}
 }

 

\subsection{\textcolor{Chapter }{PositionNot}}
\logpage{[ 21, 16, 13 ]}\nobreak
\hyperdef{L}{X865EF45D87ED1384}{}
{\noindent\textcolor{FuncColor}{$\triangleright$\enspace\texttt{PositionNot({\mdseries\slshape list, val[, from]})\index{PositionNot@\texttt{PositionNot}}
\label{PositionNot}
}\hfill{\scriptsize (operation)}}\\


 For a list \mbox{\texttt{\mdseries\slshape list}} and an object \mbox{\texttt{\mdseries\slshape val}}, \texttt{PositionNot} returns the smallest nonnegative integer $n$ such that $\mbox{\texttt{\mdseries\slshape list}}[n]$ is either unbound or not equal to \mbox{\texttt{\mdseries\slshape val}}. If a starting index \mbox{\texttt{\mdseries\slshape from}} is given, it returns the first position with this property starting the search \emph{after} position \mbox{\texttt{\mdseries\slshape from}}. 

 
\begin{Verbatim}[commandchars=!@|,fontsize=\small,frame=single,label=Example]
  !gapprompt@gap>| !gapinput@l:= [ 1, 1, 2, 3, 2 ];;  PositionNot( l, 1 );|
  3
  !gapprompt@gap>| !gapinput@PositionNot( l, 1, 4 );  PositionNot( l, 2, 4 );|
  5
  6
\end{Verbatim}
 }

 

\subsection{\textcolor{Chapter }{PositionNonZero}}
\logpage{[ 21, 16, 14 ]}\nobreak
\hyperdef{L}{X7F42E5AD87EC9D5A}{}
{\noindent\textcolor{FuncColor}{$\triangleright$\enspace\texttt{PositionNonZero({\mdseries\slshape vec[, from]})\index{PositionNonZero@\texttt{PositionNonZero}}
\label{PositionNonZero}
}\hfill{\scriptsize (operation)}}\\


 For a row vector \mbox{\texttt{\mdseries\slshape vec}}, \texttt{PositionNonZero} returns the position of the first non\texttt{\symbol{45}}zero element of \mbox{\texttt{\mdseries\slshape vec}}, or \texttt{Length(} \mbox{\texttt{\mdseries\slshape vec}} \texttt{)+1} if all entries of \mbox{\texttt{\mdseries\slshape vec}} are zero. 

 If a starting index \mbox{\texttt{\mdseries\slshape from}} is given, it returns the position of the first occurrence starting the search \emph{after} position \mbox{\texttt{\mdseries\slshape from}}. 

 \texttt{PositionNonZero} implements a special case of \texttt{PositionNot} (\ref{PositionNot}). Namely, the element to be avoided is the zero element, and the list must be
(at least) homogeneous because otherwise the zero element cannot be specified
implicitly. 

 
\begin{Verbatim}[commandchars=!@|,fontsize=\small,frame=single,label=Example]
  !gapprompt@gap>| !gapinput@PositionNonZero( [ 1, 1, 2, 3, 2 ] );|
  1
  !gapprompt@gap>| !gapinput@PositionNonZero( [ 2, 3, 4, 5 ] * Z(2) );|
  2
\end{Verbatim}
 }

 

\subsection{\textcolor{Chapter }{PositionSublist}}
\logpage{[ 21, 16, 15 ]}\nobreak
\hyperdef{L}{X87A8C62A867D6DA4}{}
{\noindent\textcolor{FuncColor}{$\triangleright$\enspace\texttt{PositionSublist({\mdseries\slshape list, sub[, from]})\index{PositionSublist@\texttt{PositionSublist}}
\label{PositionSublist}
}\hfill{\scriptsize (operation)}}\\


 returns the smallest index in the list \mbox{\texttt{\mdseries\slshape list}} at which a sublist equal to \mbox{\texttt{\mdseries\slshape sub}} starts. If \mbox{\texttt{\mdseries\slshape sub}} does not occur the operation returns \texttt{fail}. The version with given \mbox{\texttt{\mdseries\slshape from}} starts searching \emph{after} position \mbox{\texttt{\mdseries\slshape from}}. 

 To determine whether \mbox{\texttt{\mdseries\slshape sub}} matches \mbox{\texttt{\mdseries\slshape list}} at a particular position, use \texttt{IsMatchingSublist} (\ref{IsMatchingSublist}) instead. }

 }

  
\section{\textcolor{Chapter }{Properties and Attributes for Lists}}\label{Properties and Attributes for Lists}
\logpage{[ 21, 17, 0 ]}
\hyperdef{L}{X7865747A7CCF5812}{}
{
  A list that contains mutable objects (like lists or records) \emph{cannot} store attribute values that depend on the values of its entries, such as
whether it is homogeneous, sorted, or strictly sorted, as changes in any of
its entries could change such property values, like the following example
shows. 

 
\begin{Verbatim}[commandchars=!@|,fontsize=\small,frame=single,label=Example]
  !gapprompt@gap>| !gapinput@l:=[[1],[2]];|
  [ [ 1 ], [ 2 ] ]
  !gapprompt@gap>| !gapinput@IsSSortedList(l);|
  true
  !gapprompt@gap>| !gapinput@l[1][1]:=3;|
  3
  !gapprompt@gap>| !gapinput@IsSSortedList(l);|
  false
\end{Verbatim}
 

 For such lists these property values must be computed anew each time the
property is asked for. For example, if \mbox{\texttt{\mdseries\slshape list}} is a list of mutable row vectors then the call of \texttt{Position} (\ref{Position}) with \mbox{\texttt{\mdseries\slshape list}} as first argument cannot take advantage of the fact that \mbox{\texttt{\mdseries\slshape list}} is in fact sorted. One solution is to call explicitly \texttt{PositionSorted} (\ref{PositionSorted}) in such a situation, another solution is to replace \mbox{\texttt{\mdseries\slshape list}} by an immutable copy using \texttt{Immutable} (\ref{Immutable}). 

\subsection{\textcolor{Chapter }{IsMatchingSublist}}
\logpage{[ 21, 17, 1 ]}\nobreak
\hyperdef{L}{X83F8EC7C7BF27EFC}{}
{\noindent\textcolor{FuncColor}{$\triangleright$\enspace\texttt{IsMatchingSublist({\mdseries\slshape list, sub[, at]})\index{IsMatchingSublist@\texttt{IsMatchingSublist}}
\label{IsMatchingSublist}
}\hfill{\scriptsize (operation)}}\\


 returns \texttt{true} if \mbox{\texttt{\mdseries\slshape sub}} matches a sublist of \mbox{\texttt{\mdseries\slshape list}} from position \texttt{1} (or position \mbox{\texttt{\mdseries\slshape at}}, in the case of three arguments), or \texttt{false}, otherwise. If \mbox{\texttt{\mdseries\slshape sub}} is empty \texttt{true} is returned. If \mbox{\texttt{\mdseries\slshape list}} is empty but \mbox{\texttt{\mdseries\slshape sub}} is non\texttt{\symbol{45}}empty \texttt{false} is returned. 

 If you actually want to know whether there is an \mbox{\texttt{\mdseries\slshape at}} for which \texttt{IsMatchingSublist( \mbox{\texttt{\mdseries\slshape list}}, \mbox{\texttt{\mdseries\slshape sub}}, \mbox{\texttt{\mdseries\slshape at}} )} is true, use a construction like \texttt{PositionSublist( \mbox{\texttt{\mdseries\slshape list}}, \mbox{\texttt{\mdseries\slshape sub}} ) {\textless}{\textgreater} fail} instead (see \texttt{PositionSublist} (\ref{PositionSublist})); it's more efficient. }

 

\subsection{\textcolor{Chapter }{IsDuplicateFree}}
\logpage{[ 21, 17, 2 ]}\nobreak
\hyperdef{L}{X7FA892828252BB3B}{}
{\noindent\textcolor{FuncColor}{$\triangleright$\enspace\texttt{IsDuplicateFree({\mdseries\slshape obj})\index{IsDuplicateFree@\texttt{IsDuplicateFree}}
\label{IsDuplicateFree}
}\hfill{\scriptsize (property)}}\\
\noindent\textcolor{FuncColor}{$\triangleright$\enspace\texttt{IsDuplicateFreeList({\mdseries\slshape obj})\index{IsDuplicateFreeList@\texttt{IsDuplicateFreeList}}
\label{IsDuplicateFreeList}
}\hfill{\scriptsize (filter)}}\\


 \texttt{IsDuplicateFree} returns \texttt{true} if \mbox{\texttt{\mdseries\slshape obj}} is both a list or collection, and it is duplicate free; \index{duplicate free} otherwise it returns \texttt{false}. \texttt{IsDuplicateFreeList} is a synonym for \texttt{IsDuplicateFree and IsList}. 

 A list is \emph{duplicate free} if it is dense and does not contain equal entries in different positions.
Every domain (see{\nobreakspace}\ref{Domains}) is duplicate free. 

 Note that \textsf{GAP} cannot compare arbitrary objects (by equality). This can cause that \texttt{IsDuplicateFree} runs into an error, if \mbox{\texttt{\mdseries\slshape obj}} is a list with some non\texttt{\symbol{45}}comparable entries. }

 

\subsection{\textcolor{Chapter }{IsSortedList}}
\logpage{[ 21, 17, 3 ]}\nobreak
\hyperdef{L}{X7BAA9B0E81D4A884}{}
{\noindent\textcolor{FuncColor}{$\triangleright$\enspace\texttt{IsSortedList({\mdseries\slshape obj})\index{IsSortedList@\texttt{IsSortedList}}
\label{IsSortedList}
}\hfill{\scriptsize (property)}}\\


 returns \texttt{true} if \mbox{\texttt{\mdseries\slshape obj}} is a list and it is sorted, \index{list!sorted} and \texttt{false} otherwise. 

 A list \mbox{\texttt{\mdseries\slshape list}} is \emph{sorted} if it is dense (see{\nobreakspace}\texttt{IsDenseList} (\ref{IsDenseList})) and satisfies the relation $\mbox{\texttt{\mdseries\slshape list}}[i] \leq \mbox{\texttt{\mdseries\slshape list}}[j]$ whenever $i < j$. Note that a sorted list is not necessarily duplicate free
(see{\nobreakspace}\texttt{IsDuplicateFree} (\ref{IsDuplicateFree}) and \texttt{IsSSortedList} (\ref{IsSSortedList})). 

 Many sorted lists are in fact homogeneous (see{\nobreakspace}\texttt{IsHomogeneousList} (\ref{IsHomogeneousList})), but also non\texttt{\symbol{45}}homogeneous lists may be sorted
(see{\nobreakspace}\ref{Comparison Operations for Elements}). 

 In sorted lists, membership test and computing of positions can be done by
binary search, see{\nobreakspace}\ref{Sorted Lists and Sets}. 

 Note that \textsf{GAP} cannot compare (by less than) arbitrary objects. This can cause that \texttt{IsSortedList} runs into an error, if \mbox{\texttt{\mdseries\slshape obj}} is a list with some non\texttt{\symbol{45}}comparable entries. }

 

\subsection{\textcolor{Chapter }{IsSSortedList}}
\logpage{[ 21, 17, 4 ]}\nobreak
\hyperdef{L}{X80CDAF45782E8DCB}{}
{\noindent\textcolor{FuncColor}{$\triangleright$\enspace\texttt{IsSSortedList({\mdseries\slshape obj})\index{IsSSortedList@\texttt{IsSSortedList}}
\label{IsSSortedList}
}\hfill{\scriptsize (property)}}\\
\noindent\textcolor{FuncColor}{$\triangleright$\enspace\texttt{IsSet({\mdseries\slshape obj})\index{IsSet@\texttt{IsSet}}
\label{IsSet}
}\hfill{\scriptsize (property)}}\\


 returns \texttt{true} if \mbox{\texttt{\mdseries\slshape obj}} is a list and it is strictly sorted, \index{strictly sorted list} and \texttt{false} otherwise. \texttt{IsSSortedList} is short for ``is strictly sorted list''; \texttt{IsSet} is just a synonym for \texttt{IsSSortedList}. 

 A list \mbox{\texttt{\mdseries\slshape list}} is \emph{strictly sorted} if it is sorted (see{\nobreakspace}\texttt{IsSortedList} (\ref{IsSortedList})) and satisfies the relation $\mbox{\texttt{\mdseries\slshape list}}[i] < \mbox{\texttt{\mdseries\slshape list}}[j]$ whenever $i < j$. In particular, such lists are duplicate free (see{\nobreakspace}\texttt{IsDuplicateFree} (\ref{IsDuplicateFree})). 

 (Currently there is little special treatment of lists that are sorted but not
strictly sorted. In particular, internally represented lists will \emph{not} store that they are sorted but not strictly sorted.) 

 Note that \textsf{GAP} cannot compare (by less than) arbitrary objects. This can cause that \texttt{IsSSortedList} runs into an error, if \mbox{\texttt{\mdseries\slshape obj}} is a list with some non\texttt{\symbol{45}}comparable entries. }

 

\subsection{\textcolor{Chapter }{Length}}
\logpage{[ 21, 17, 5 ]}\nobreak
\hyperdef{L}{X780769238600AFD1}{}
{\noindent\textcolor{FuncColor}{$\triangleright$\enspace\texttt{Length({\mdseries\slshape list})\index{Length@\texttt{Length}}
\label{Length}
}\hfill{\scriptsize (attribute)}}\\


 returns the \emph{length} of the list \mbox{\texttt{\mdseries\slshape list}}, which is defined to be the index of the last bound entry in \mbox{\texttt{\mdseries\slshape list}}. }

 

\subsection{\textcolor{Chapter }{ConstantTimeAccessList}}
\logpage{[ 21, 17, 6 ]}\nobreak
\hyperdef{L}{X7B55FB967CDEF468}{}
{\noindent\textcolor{FuncColor}{$\triangleright$\enspace\texttt{ConstantTimeAccessList({\mdseries\slshape list})\index{ConstantTimeAccessList@\texttt{ConstantTimeAccessList}}
\label{ConstantTimeAccessList}
}\hfill{\scriptsize (attribute)}}\\


 \texttt{ConstantTimeAccessList} returns an immutable list containing the same elements as the list \mbox{\texttt{\mdseries\slshape list}} (which may have holes) in the same order. If \mbox{\texttt{\mdseries\slshape list}} is already a constant time access list, \texttt{ConstantTimeAccessList} returns an immutable copy of \mbox{\texttt{\mdseries\slshape list}} directly. Otherwise it puts all elements and holes of \mbox{\texttt{\mdseries\slshape list}} into a new list and makes that list immutable. }

 }

  
\section{\textcolor{Chapter }{Sorting Lists}}\label{Sorting Lists}
\logpage{[ 21, 18, 0 ]}
\hyperdef{L}{X83E558E37D1B44D4}{}
{
  GAP implements three different families of sorting algorithms. The default
algorithm is pattern\texttt{\symbol{45}}defeating quicksort, a variant of
quicksort which performs better on partially sorted lists and has good
worst\texttt{\symbol{45}}case behaviour. The functions which begin \texttt{Stable} are stable (equal elements keep the same relative order in the sorted list)
and use merge sort. Finally, the functions which begin \texttt{Shell} use the shell sort which was GAP's default search algorithm before 4.9. \texttt{Sortex} (\ref{Sortex}) and \texttt{SortingPerm} (\ref{SortingPerm}) are also stable. 

\subsection{\textcolor{Chapter }{Sort}}
\logpage{[ 21, 18, 1 ]}\nobreak
\hyperdef{L}{X7FE4975F8166884D}{}
{\noindent\textcolor{FuncColor}{$\triangleright$\enspace\texttt{Sort({\mdseries\slshape list[, func]})\index{Sort@\texttt{Sort}}
\label{Sort}
}\hfill{\scriptsize (operation)}}\\
\noindent\textcolor{FuncColor}{$\triangleright$\enspace\texttt{SortBy({\mdseries\slshape list, func})\index{SortBy@\texttt{SortBy}}
\label{SortBy}
}\hfill{\scriptsize (operation)}}\\
\noindent\textcolor{FuncColor}{$\triangleright$\enspace\texttt{StableSort({\mdseries\slshape list[, func]})\index{StableSort@\texttt{StableSort}}
\label{StableSort}
}\hfill{\scriptsize (operation)}}\\
\noindent\textcolor{FuncColor}{$\triangleright$\enspace\texttt{StableSortBy({\mdseries\slshape list[, func]})\index{StableSortBy@\texttt{StableSortBy}}
\label{StableSortBy}
}\hfill{\scriptsize (operation)}}\\


 \texttt{Sort} sorts the list \mbox{\texttt{\mdseries\slshape list}} in increasing order. In the one argument form \texttt{Sort} uses the operator \texttt{{\textless}} to compare the elements. (If the list is not homogeneous it is the user's
responsibility to ensure that \texttt{{\textless}} is defined for all element pairs, see{\nobreakspace}\ref{Comparison Operations for Elements}) In the two argument form \texttt{Sort} uses the function \mbox{\texttt{\mdseries\slshape func}} to compare elements. \mbox{\texttt{\mdseries\slshape func}} must be a function taking two arguments that returns \texttt{true} if the first is regarded as strictly smaller than the second, and \texttt{false} otherwise. 

 \texttt{StableSort} behaves identically to \texttt{Sort}, except that \texttt{StableSort} will keep elements which compare equal in the same relative order, while \texttt{Sort} may change their relative order. 

 \texttt{Sort} does not return anything, it just changes the argument \mbox{\texttt{\mdseries\slshape list}}. Use \texttt{ShallowCopy} (\ref{ShallowCopy}) if you want to keep \mbox{\texttt{\mdseries\slshape list}}. Use \texttt{Reversed} (\ref{Reversed}) if you want to get a new list that is sorted in decreasing order. 

 \texttt{SortBy} sorts the list \mbox{\texttt{\mdseries\slshape list}} into an order such that \texttt{func(list[i]) {\textless}= func(list[i+1])} for all relevant \mbox{\texttt{\mdseries\slshape i}}. \mbox{\texttt{\mdseries\slshape func}} must thus be a function on one argument which returns values that can be
compared. Each \texttt{func(list[i])} is computed just once and stored, making this more efficient than using the
two\texttt{\symbol{45}}argument version of \texttt{Sort} in many cases. 

 \texttt{StableSortBy} behaves the same as \texttt{SortBy} except that, like \texttt{StableSort}, it keeps pairs of values which compare equal when \texttt{func} is applied to them in the same relative order. 

 
\begin{Verbatim}[commandchars=!@|,fontsize=\small,frame=single,label=Example]
  !gapprompt@gap>| !gapinput@list := [ 5, 4, 6, 1, 7, 5 ];; Sort( list ); list;|
  [ 1, 4, 5, 5, 6, 7 ]
  !gapprompt@gap>| !gapinput@SortBy(list, x -> x mod 3);|
  !gapprompt@gap>| !gapinput@list; # Sorted by mod 3|
  [ 6, 1, 4, 7, 5, 5]
  !gapprompt@gap>| !gapinput@list := [ [0,6], [1,2], [1,3], [1,5], [0,4], [3,4] ];;|
  !gapprompt@gap>| !gapinput@Sort( list, function(v,w) return v*v < w*w; end );|
  !gapprompt@gap>| !gapinput@list;  # sorted according to the Euclidean distance from [0,0]|
  [ [ 1, 2 ], [ 1, 3 ], [ 0, 4 ], [ 3, 4 ], [ 1, 5 ], [ 0, 6 ] ]
  !gapprompt@gap>| !gapinput@SortBy( list, function(v) return v[1] + v[2]; end );|
  !gapprompt@gap>| !gapinput@list;  # sorted according to Manhattan distance from [0,0]|
  [ [ 1, 2 ], [ 1, 3 ], [ 0, 4 ], [ 1, 5 ], [ 0, 6 ], [ 3, 4 ] ]
  !gapprompt@gap>| !gapinput@list := [ [0,6], [1,3], [3,4], [1,5], [1,2], [0,4], ];;|
  !gapprompt@gap>| !gapinput@Sort( list, function(v,w) return v[1] < w[1]; end );|
  !gapprompt@gap>| !gapinput@# note the random order of the elements with equal first component:|
  !gapprompt@gap>| !gapinput@list;|
  [ [ 0, 6 ], [ 0, 4 ], [ 1, 3 ], [ 1, 5 ], [ 1, 2 ], [ 3, 4 ] ]
\end{Verbatim}
 }

 

\subsection{\textcolor{Chapter }{SortParallel}}
\logpage{[ 21, 18, 2 ]}\nobreak
\hyperdef{L}{X791F2B2C7E9B9A46}{}
{\noindent\textcolor{FuncColor}{$\triangleright$\enspace\texttt{SortParallel({\mdseries\slshape list1, list2[, func]})\index{SortParallel@\texttt{SortParallel}}
\label{SortParallel}
}\hfill{\scriptsize (operation)}}\\
\noindent\textcolor{FuncColor}{$\triangleright$\enspace\texttt{StableSortParallel({\mdseries\slshape list1, list2[, func]})\index{StableSortParallel@\texttt{StableSortParallel}}
\label{StableSortParallel}
}\hfill{\scriptsize (operation)}}\\


 \texttt{SortParallel} sorts the list \mbox{\texttt{\mdseries\slshape list1}} in increasing order just as \texttt{Sort} (\ref{Sort}) does. In parallel it applies the same exchanges that are necessary to sort \mbox{\texttt{\mdseries\slshape list1}} to the list \mbox{\texttt{\mdseries\slshape list2}}, which must of course have at least as many elements as \mbox{\texttt{\mdseries\slshape list1}} does. 

 \texttt{StableSortParallel} behaves identically to \texttt{SortParallel}, except it keeps elements in \mbox{\texttt{\mdseries\slshape list1}} which compare equal in the same relative order. 

 
\begin{Verbatim}[commandchars=!@|,fontsize=\small,frame=single,label=Example]
  !gapprompt@gap>| !gapinput@list1 := [ 5, 4, 6, 1, 7, 5 ];;|
  !gapprompt@gap>| !gapinput@list2 := [ 2, 3, 5, 7, 8, 9 ];;|
  !gapprompt@gap>| !gapinput@SortParallel( list1, list2 );|
  !gapprompt@gap>| !gapinput@list1;|
  [ 1, 4, 5, 5, 6, 7 ]
  !gapprompt@gap>| !gapinput@list2;|
  [ 7, 3, 2, 9, 5, 8 ]
\end{Verbatim}
 

 Note that \texttt{[ 7, 3, 2, 9, 5, 8 ]} or \texttt{[ 7, 3, 9, 2, 5, 8 ]} are possible results. \texttt{StableSortParallel} will always return \texttt{[ 7, 3, 2, 9, 5, 8]}. }

 

\subsection{\textcolor{Chapter }{Sortex}}
\logpage{[ 21, 18, 3 ]}\nobreak
\hyperdef{L}{X87287FCA81E2B06A}{}
{\noindent\textcolor{FuncColor}{$\triangleright$\enspace\texttt{Sortex({\mdseries\slshape list[, func]})\index{Sortex@\texttt{Sortex}}
\label{Sortex}
}\hfill{\scriptsize (operation)}}\\


 sorts the list \mbox{\texttt{\mdseries\slshape list}} and returns a permutation that can be applied to \mbox{\texttt{\mdseries\slshape list}} to obtain the sorted list. The one argument form sorts via the operator \texttt{{\textless}}, the two argument form sorts w.r.t. the function \mbox{\texttt{\mdseries\slshape func}}. The permutation returned by \texttt{Sortex} will keep elements which compare equal in the same relative order. (If the
list is not homogeneous it is the user's responsibility to ensure that \texttt{{\textless}} is defined for all element pairs, see{\nobreakspace}\ref{Comparison Operations for Elements}) 

 \texttt{Permuted} (\ref{Permuted}) allows you to rearrange a list according to a given permutation. 

 
\begin{Verbatim}[commandchars=!@|,fontsize=\small,frame=single,label=Example]
  !gapprompt@gap>| !gapinput@list1 := [ 5, 4, 6, 1, 7, 5 ];;|
  !gapprompt@gap>| !gapinput@list2 := ShallowCopy( list1 );;|
  !gapprompt@gap>| !gapinput@perm := Sortex( list1 );|
  (1,3,5,6,4)
  !gapprompt@gap>| !gapinput@list1;|
  [ 1, 4, 5, 5, 6, 7 ]
  !gapprompt@gap>| !gapinput@Permuted( list2, perm );|
  [ 1, 4, 5, 5, 6, 7 ]
\end{Verbatim}
 }

 

\subsection{\textcolor{Chapter }{SortingPerm}}
\logpage{[ 21, 18, 4 ]}\nobreak
\hyperdef{L}{X800209E881E7CECB}{}
{\noindent\textcolor{FuncColor}{$\triangleright$\enspace\texttt{SortingPerm({\mdseries\slshape list})\index{SortingPerm@\texttt{SortingPerm}}
\label{SortingPerm}
}\hfill{\scriptsize (attribute)}}\\


 \texttt{SortingPerm} returns the same as \texttt{Sortex} (\ref{Sortex}) but does \emph{not} change the argument. 

 
\begin{Verbatim}[commandchars=!@|,fontsize=\small,frame=single,label=Example]
  !gapprompt@gap>| !gapinput@list1 := [ 5, 4, 6, 1, 7, 5 ];;|
  !gapprompt@gap>| !gapinput@list2 := ShallowCopy( list1 );;|
  !gapprompt@gap>| !gapinput@perm := SortingPerm( list1 );|
  (1,3,5,6,4)
  !gapprompt@gap>| !gapinput@list1;|
  [ 5, 4, 6, 1, 7, 5 ]
  !gapprompt@gap>| !gapinput@Permuted( list2, perm );|
  [ 1, 4, 5, 5, 6, 7 ]
\end{Verbatim}
 }

 }

  
\section{\textcolor{Chapter }{Sorted Lists and Sets}}\label{Sorted Lists and Sets}
\logpage{[ 21, 19, 0 ]}
\hyperdef{L}{X80ABC25582343910}{}
{
  \index{sets} \index{multisets} Searching objects in a list works much quicker if the list is known to be
sorted. Currently \textsf{GAP} exploits the sortedness of a list automatically only if the list is \emph{strictly sorted}, which is indicated by the property \texttt{IsSSortedList} (\ref{IsSSortedList}). 

 Remember that a list of \emph{mutable} objects cannot store that it is strictly sorted but has to test it anew
whenever it is asked whether it is sorted, see the remark in{\nobreakspace}\ref{Properties and Attributes for Lists}. Therefore \textsf{GAP} cannot take advantage of the sortedness of a list if this list has mutable
entries. Moreover, if a sorted list \mbox{\texttt{\mdseries\slshape list}} with mutable elements is used as an argument of a function that \emph{expects} this argument to be sorted, for example \texttt{UniteSet} (\ref{UniteSet}) or \texttt{RemoveSet} (\ref{RemoveSet}), then it is checked whether \mbox{\texttt{\mdseries\slshape list}} is in fact sorted; this check can have the effect actually to slow down the
computations, compared to computations with sorted lists of immutable elements
or computations that do not involve functions that do automatically check
sortedness. 

 Strictly sorted lists are used to represent \emph{sets} in \textsf{GAP}. More precisely, a strictly sorted list is called a \emph{proper set} in the following, in order to avoid confusion with domains (see{\nobreakspace}\ref{Domains}) which also represent sets. 

 In short proper sets are represented by sorted lists without holes and
duplicates in \textsf{GAP}. Note that we guarantee this representation, so you may make use of the fact
that a set is represented by a sorted list in your functions. 

 In some contexts (for example see{\nobreakspace}\ref{Combinatorics}), we also want to talk about multisets. A \emph{multiset} is like a set, except that an element may appear several times in a multiset.
Such multisets are represented by sorted lists without holes that may have
duplicates. 

 This section lists only those functions that are defined exclusively for
proper sets. Set theoretic functions for general collections, such as \texttt{Intersection} (\ref{Intersection:for a list}) and \texttt{Union} (\ref{Union:for a list}), are described in Chapter{\nobreakspace}\ref{Collections}. In particular, for the construction of proper sets, see{\nobreakspace}\texttt{SSortedList} (\ref{SSortedList}) and \texttt{AsSSortedList} (\ref{AsSSortedList}). For finding positions in sorted lists, see{\nobreakspace}\texttt{PositionSorted} (\ref{PositionSorted}). 

 There are nondestructive counterparts of the functions \texttt{UniteSet} (\ref{UniteSet}), \texttt{IntersectSet} (\ref{IntersectSet}), and \texttt{SubtractSet} (\ref{SubtractSet}) available for proper sets. These are \texttt{UnionSet}, \texttt{IntersectionSet}, and \texttt{Difference} (\ref{Difference}). The former two are methods for the more general operations \texttt{Union} (\ref{Union:for a list}) and \texttt{Intersection} (\ref{Intersection:for a list}), the latter is itself an operation (see{\nobreakspace}\texttt{Difference} (\ref{Difference})). 

 The result of \texttt{IntersectionSet} and \texttt{UnionSet} is always a new list, that is not identical to any other list. The elements of
that list however are identical to the corresponding elements of the first
argument \mbox{\texttt{\mdseries\slshape set}}. If \mbox{\texttt{\mdseries\slshape set}} is not a proper set it is not specified to which of a number of equal elements
in \mbox{\texttt{\mdseries\slshape set}} the element in the result is identical (see{\nobreakspace}\ref{Identical Lists}). The following functions, if not explicitly stated differently, take two
arguments, \mbox{\texttt{\mdseries\slshape set}} and \mbox{\texttt{\mdseries\slshape obj}}, where \mbox{\texttt{\mdseries\slshape set}} must be a proper set, otherwise an error is signalled; If the second argument \mbox{\texttt{\mdseries\slshape obj}} is a list that is not a proper set then \texttt{Set} (\ref{Set}) is silently applied to it first. 

\subsection{\textcolor{Chapter }{\texttt{\symbol{92}}in (for strictly sorted lists)}}
\logpage{[ 21, 19, 1 ]}\nobreak
\hyperdef{L}{X7B16AD597CB12305}{}
{\noindent\textcolor{FuncColor}{$\triangleright$\enspace\texttt{\texttt{\symbol{92}}in({\mdseries\slshape obj, list})\index{\texttt{\symbol{92}}in@\texttt{\texttt{\symbol{92}}in}!for strictly sorted lists}
\label{bSlashin:for strictly sorted lists}
}\hfill{\scriptsize (method)}}\\


 For a list \mbox{\texttt{\mdseries\slshape list}} that stores that it is strictly sorted, the test with \texttt{\texttt{\symbol{92}}in} whether the object \mbox{\texttt{\mdseries\slshape obj}} is an entry of \mbox{\texttt{\mdseries\slshape list}} uses binary search. This test can be entered also with the infix notation \mbox{\texttt{\mdseries\slshape obj}} \texttt{in} \mbox{\texttt{\mdseries\slshape list}}. }

 

\subsection{\textcolor{Chapter }{IsEqualSet}}
\logpage{[ 21, 19, 2 ]}\nobreak
\hyperdef{L}{X7B4C0FEE7CDF6F2A}{}
{\noindent\textcolor{FuncColor}{$\triangleright$\enspace\texttt{IsEqualSet({\mdseries\slshape list1, list2})\index{IsEqualSet@\texttt{IsEqualSet}}
\label{IsEqualSet}
}\hfill{\scriptsize (operation)}}\\


 \index{test!for set equality} tests whether \mbox{\texttt{\mdseries\slshape list1}} and \mbox{\texttt{\mdseries\slshape list2}} are equal \emph{when viewed as sets}, that is if every element of \mbox{\texttt{\mdseries\slshape list1}} is an element of \mbox{\texttt{\mdseries\slshape list2}} and vice versa. Either argument of \texttt{IsEqualSet} may also be a list that is not a proper set, in which case \texttt{Set} (\ref{Set}) is applied to it first. 

 If both lists are proper sets then they are of course equal if and only if
they are also equal as lists. Thus \texttt{IsEqualSet( \mbox{\texttt{\mdseries\slshape list1}}, \mbox{\texttt{\mdseries\slshape list2}} )} is equivalent to \texttt{Set( \mbox{\texttt{\mdseries\slshape list1}} ) = Set( \mbox{\texttt{\mdseries\slshape list2}} )} (see{\nobreakspace}\texttt{Set} (\ref{Set})), but the former is more efficient. 

 
\begin{Verbatim}[commandchars=!@|,fontsize=\small,frame=single,label=Example]
  !gapprompt@gap>| !gapinput@IsEqualSet( [2,3,5,7,11], [11,7,5,3,2] );|
  true
  !gapprompt@gap>| !gapinput@IsEqualSet( [2,3,5,7,11], [2,3,5,7,11,13] );|
  false
\end{Verbatim}
 }

 

\subsection{\textcolor{Chapter }{IsSubsetSet}}
\logpage{[ 21, 19, 3 ]}\nobreak
\hyperdef{L}{X79B940567A849216}{}
{\noindent\textcolor{FuncColor}{$\triangleright$\enspace\texttt{IsSubsetSet({\mdseries\slshape list1, list2})\index{IsSubsetSet@\texttt{IsSubsetSet}}
\label{IsSubsetSet}
}\hfill{\scriptsize (operation)}}\\


 tests whether every element of \mbox{\texttt{\mdseries\slshape list2}} is contained in \mbox{\texttt{\mdseries\slshape list1}}. Either argument of \texttt{IsSubsetSet} may also be a list that is not a proper set, in which case \texttt{Set} (\ref{Set}) is applied to it first. }

 

\subsection{\textcolor{Chapter }{AddSet}}
\logpage{[ 21, 19, 4 ]}\nobreak
\hyperdef{L}{X832C23CC7FCD8892}{}
{\noindent\textcolor{FuncColor}{$\triangleright$\enspace\texttt{AddSet({\mdseries\slshape set, obj})\index{AddSet@\texttt{AddSet}}
\label{AddSet}
}\hfill{\scriptsize (operation)}}\\


 \index{add!an element to a set} adds the element \mbox{\texttt{\mdseries\slshape obj}} to the proper set \mbox{\texttt{\mdseries\slshape set}}. If \mbox{\texttt{\mdseries\slshape obj}} is already contained in \mbox{\texttt{\mdseries\slshape set}} then \mbox{\texttt{\mdseries\slshape set}} is not changed. Otherwise \mbox{\texttt{\mdseries\slshape obj}} is inserted at the correct position such that \mbox{\texttt{\mdseries\slshape set}} is again a proper set afterwards. 

 Note that \mbox{\texttt{\mdseries\slshape obj}} must be in the same family as each element of \mbox{\texttt{\mdseries\slshape set}}. 
\begin{Verbatim}[commandchars=!@|,fontsize=\small,frame=single,label=Example]
  !gapprompt@gap>| !gapinput@s := [2,3,7,11];;|
  !gapprompt@gap>| !gapinput@AddSet( s, 5 );  s;|
  [ 2, 3, 5, 7, 11 ]
  !gapprompt@gap>| !gapinput@AddSet( s, 13 );  s;|
  [ 2, 3, 5, 7, 11, 13 ]
  !gapprompt@gap>| !gapinput@AddSet( s, 3 );  s;|
  [ 2, 3, 5, 7, 11, 13 ]
\end{Verbatim}
 }

 

\subsection{\textcolor{Chapter }{RemoveSet}}
\logpage{[ 21, 19, 5 ]}\nobreak
\hyperdef{L}{X7FCA282E789A4F4B}{}
{\noindent\textcolor{FuncColor}{$\triangleright$\enspace\texttt{RemoveSet({\mdseries\slshape set, obj})\index{RemoveSet@\texttt{RemoveSet}}
\label{RemoveSet}
}\hfill{\scriptsize (operation)}}\\


 \index{remove!an element from a set} removes the element \mbox{\texttt{\mdseries\slshape obj}} from the proper set \mbox{\texttt{\mdseries\slshape set}}. If \mbox{\texttt{\mdseries\slshape obj}} is not contained in \mbox{\texttt{\mdseries\slshape set}} then \mbox{\texttt{\mdseries\slshape set}} is not changed. If \mbox{\texttt{\mdseries\slshape obj}} is an element of \mbox{\texttt{\mdseries\slshape set}} it is removed and all the following elements in the list are moved one
position forward. 

 
\begin{Verbatim}[commandchars=!@|,fontsize=\small,frame=single,label=Example]
  !gapprompt@gap>| !gapinput@s := [ 2, 3, 4, 5, 6, 7 ];;|
  !gapprompt@gap>| !gapinput@RemoveSet( s, 6 ); s;|
  [ 2, 3, 4, 5, 7 ]
  !gapprompt@gap>| !gapinput@RemoveSet( s, 10 ); s;|
  [ 2, 3, 4, 5, 7 ]
\end{Verbatim}
 }

 

\subsection{\textcolor{Chapter }{UniteSet}}
\logpage{[ 21, 19, 6 ]}\nobreak
\hyperdef{L}{X7B3469CD7EFC1A87}{}
{\noindent\textcolor{FuncColor}{$\triangleright$\enspace\texttt{UniteSet({\mdseries\slshape set, list})\index{UniteSet@\texttt{UniteSet}}
\label{UniteSet}
}\hfill{\scriptsize (operation)}}\\


 \index{union!of sets} unites the proper set \mbox{\texttt{\mdseries\slshape set}} with \mbox{\texttt{\mdseries\slshape list}}. This is equivalent to adding all elements of \mbox{\texttt{\mdseries\slshape list}} to \mbox{\texttt{\mdseries\slshape set}} (see{\nobreakspace}\texttt{AddSet} (\ref{AddSet})). 

 
\begin{Verbatim}[commandchars=!@|,fontsize=\small,frame=single,label=Example]
  !gapprompt@gap>| !gapinput@set := [ 2, 3, 5, 7, 11 ];;|
  !gapprompt@gap>| !gapinput@UniteSet( set, [ 4, 8, 9 ] );  set;|
  [ 2, 3, 4, 5, 7, 8, 9, 11 ]
  !gapprompt@gap>| !gapinput@UniteSet( set, [ 16, 9, 25, 13, 16 ] );  set;|
  [ 2, 3, 4, 5, 7, 8, 9, 11, 13, 16, 25 ]
\end{Verbatim}
 }

 

\subsection{\textcolor{Chapter }{IntersectSet}}
\logpage{[ 21, 19, 7 ]}\nobreak
\hyperdef{L}{X8473AA657FEC3D4D}{}
{\noindent\textcolor{FuncColor}{$\triangleright$\enspace\texttt{IntersectSet({\mdseries\slshape set, list})\index{IntersectSet@\texttt{IntersectSet}}
\label{IntersectSet}
}\hfill{\scriptsize (operation)}}\\


 \index{intersection!of sets} intersects the proper set \mbox{\texttt{\mdseries\slshape set}} with \mbox{\texttt{\mdseries\slshape list}}. This is equivalent to removing from \mbox{\texttt{\mdseries\slshape set}} all elements of \mbox{\texttt{\mdseries\slshape set}} that are not contained in \mbox{\texttt{\mdseries\slshape list}}. 

 
\begin{Verbatim}[commandchars=!@|,fontsize=\small,frame=single,label=Example]
  !gapprompt@gap>| !gapinput@set := [ 2, 3, 4, 5, 7, 8, 9, 11, 13, 16 ];;|
  !gapprompt@gap>| !gapinput@IntersectSet( set, [ 3, 5, 7, 9, 11, 13, 15, 17 ] );  set;|
  [ 3, 5, 7, 9, 11, 13 ]
  !gapprompt@gap>| !gapinput@IntersectSet( set, [ 9, 4, 6, 8 ] );  set;|
  [ 9 ]
\end{Verbatim}
 }

 

\subsection{\textcolor{Chapter }{SubtractSet}}
\logpage{[ 21, 19, 8 ]}\nobreak
\hyperdef{L}{X80B427537EB07D09}{}
{\noindent\textcolor{FuncColor}{$\triangleright$\enspace\texttt{SubtractSet({\mdseries\slshape set, list})\index{SubtractSet@\texttt{SubtractSet}}
\label{SubtractSet}
}\hfill{\scriptsize (operation)}}\\


 \index{subtract!a set from another} subtracts \mbox{\texttt{\mdseries\slshape list}} from the proper set \mbox{\texttt{\mdseries\slshape set}}. This is equivalent to removing from \mbox{\texttt{\mdseries\slshape set}} all elements of \mbox{\texttt{\mdseries\slshape list}}. 

 
\begin{Verbatim}[commandchars=!@|,fontsize=\small,frame=single,label=Example]
  !gapprompt@gap>| !gapinput@set := [ 2, 3, 4, 5, 6, 7, 8, 9, 10, 11 ];;|
  !gapprompt@gap>| !gapinput@SubtractSet( set, [ 6, 10 ] );  set;|
  [ 2, 3, 4, 5, 7, 8, 9, 11 ]
  !gapprompt@gap>| !gapinput@SubtractSet( set, [ 9, 4, 6, 8 ] );  set;|
  [ 2, 3, 5, 7, 11 ]
\end{Verbatim}
 }

                                                          }

  
\section{\textcolor{Chapter }{Operations for Lists}}\label{Operations for Lists}
\logpage{[ 21, 20, 0 ]}
\hyperdef{L}{X7DF510F7848CBBFD}{}
{
  Several of the following functions expect the first argument to be either a
list or a collection (see{\nobreakspace}\ref{Collections}), with possibly slightly different meaning for lists and
non\texttt{\symbol{45}}list collections. \index{concatenation!of lists} 

\subsection{\textcolor{Chapter }{Concatenation (for several lists)}}
\logpage{[ 21, 20, 1 ]}\nobreak
\hyperdef{L}{X840C55A77D1BB2E1}{}
{\noindent\textcolor{FuncColor}{$\triangleright$\enspace\texttt{Concatenation({\mdseries\slshape list1, list2, ...})\index{Concatenation@\texttt{Concatenation}!for several lists}
\label{Concatenation:for several lists}
}\hfill{\scriptsize (function)}}\\
\noindent\textcolor{FuncColor}{$\triangleright$\enspace\texttt{Concatenation({\mdseries\slshape list})\index{Concatenation@\texttt{Concatenation}!for a list of lists}
\label{Concatenation:for a list of lists}
}\hfill{\scriptsize (function)}}\\


 In the first form \texttt{Concatenation} returns the concatenation of the lists \mbox{\texttt{\mdseries\slshape list1}}, \mbox{\texttt{\mdseries\slshape list2}}, etc. The \emph{concatenation} is the list that begins with the elements of \mbox{\texttt{\mdseries\slshape list1}}, followed by the elements of \mbox{\texttt{\mdseries\slshape list2}}, and so on. Each list may also contain holes, in which case the concatenation
also contains holes at the corresponding positions. 

 In the second form \mbox{\texttt{\mdseries\slshape list}} must be a dense list of lists \mbox{\texttt{\mdseries\slshape list1}}, \mbox{\texttt{\mdseries\slshape list2}}, etc., and \texttt{Concatenation} returns the concatenation of those lists. 

 The result is a new mutable list, that is not identical to any other list. The
elements of that list however are identical to the corresponding elements of \mbox{\texttt{\mdseries\slshape list1}}, \mbox{\texttt{\mdseries\slshape list2}}, etc. (see{\nobreakspace}\ref{Identical Lists}). 

 Note that \texttt{Concatenation} creates a new list and leaves its arguments unchanged, while \texttt{Append} (\ref{Append}) changes its first argument. For computing the union of proper sets, \texttt{Union} (\ref{Union:for a list}) can be used, see also \ref{Sorted Lists and Sets}. 

 
\begin{Verbatim}[commandchars=!@|,fontsize=\small,frame=single,label=Example]
  !gapprompt@gap>| !gapinput@Concatenation( [ 1, 2, 3 ], [ 4, 5 ] );|
  [ 1, 2, 3, 4, 5 ]
  !gapprompt@gap>| !gapinput@Concatenation( [2,3,,5,,7], [11,,13,,,,17,,19] );|
  [ 2, 3,, 5,, 7, 11,, 13,,,, 17,, 19 ]
  !gapprompt@gap>| !gapinput@Concatenation( [ [1,2,3], [2,3,4], [3,4,5] ] );|
  [ 1, 2, 3, 2, 3, 4, 3, 4, 5 ]
\end{Verbatim}
 }

 

\subsection{\textcolor{Chapter }{Compacted}}
\logpage{[ 21, 20, 2 ]}\nobreak
\hyperdef{L}{X7CB0A6AF87C7FAF7}{}
{\noindent\textcolor{FuncColor}{$\triangleright$\enspace\texttt{Compacted({\mdseries\slshape list})\index{Compacted@\texttt{Compacted}}
\label{Compacted}
}\hfill{\scriptsize (operation)}}\\


 returns a new mutable list that contains the elements of \mbox{\texttt{\mdseries\slshape list}} in the same order but omitting the holes. 

 
\begin{Verbatim}[commandchars=!@|,fontsize=\small,frame=single,label=Example]
  !gapprompt@gap>| !gapinput@l:=[,1,,,3,,,4,[5,,,6],7];;  Compacted( l );|
  [ 1, 3, 4, [ 5,,, 6 ], 7 ]
\end{Verbatim}
 }

 

\subsection{\textcolor{Chapter }{Collected}}
\logpage{[ 21, 20, 3 ]}\nobreak
\hyperdef{L}{X7ECE9056792F28BA}{}
{\noindent\textcolor{FuncColor}{$\triangleright$\enspace\texttt{Collected({\mdseries\slshape list})\index{Collected@\texttt{Collected}}
\label{Collected}
}\hfill{\scriptsize (operation)}}\\


 returns a new list \mbox{\texttt{\mdseries\slshape new}} that contains for each element \mbox{\texttt{\mdseries\slshape elm}} of the list \mbox{\texttt{\mdseries\slshape list}} a list of length two, the first element of this is \mbox{\texttt{\mdseries\slshape elm}} itself and the second element is the number of times \mbox{\texttt{\mdseries\slshape elm}} appears in \mbox{\texttt{\mdseries\slshape list}}. The order of those pairs in \mbox{\texttt{\mdseries\slshape new}} corresponds to the ordering of the elements elm, so that the result is sorted. 

 For all pairs of elements in \mbox{\texttt{\mdseries\slshape list}} the comparison via \texttt{{\textless}} must be defined. 

 
\begin{Verbatim}[commandchars=!@|,fontsize=\small,frame=single,label=Example]
  !gapprompt@gap>| !gapinput@Factors( Factorial( 10 ) );|
  [ 2, 2, 2, 2, 2, 2, 2, 2, 3, 3, 3, 3, 5, 5, 7 ]
  !gapprompt@gap>| !gapinput@Collected( last );|
  [ [ 2, 8 ], [ 3, 4 ], [ 5, 2 ], [ 7, 1 ] ]
  !gapprompt@gap>| !gapinput@Collected( last );|
  [ [ [ 2, 8 ], 1 ], [ [ 3, 4 ], 1 ], [ [ 5, 2 ], 1 ], [ [ 7, 1 ], 1 ] ]
\end{Verbatim}
 }

 

\subsection{\textcolor{Chapter }{DuplicateFreeList}}
\logpage{[ 21, 20, 4 ]}\nobreak
\hyperdef{L}{X8727F2928467C2F9}{}
{\noindent\textcolor{FuncColor}{$\triangleright$\enspace\texttt{DuplicateFreeList({\mdseries\slshape list})\index{DuplicateFreeList@\texttt{DuplicateFreeList}}
\label{DuplicateFreeList}
}\hfill{\scriptsize (operation)}}\\
\noindent\textcolor{FuncColor}{$\triangleright$\enspace\texttt{Unique({\mdseries\slshape list})\index{Unique@\texttt{Unique}}
\label{Unique}
}\hfill{\scriptsize (operation)}}\\


 returns a new mutable list whose entries are the elements of the list \mbox{\texttt{\mdseries\slshape list}} with duplicates removed. \texttt{DuplicateFreeList} only uses the \texttt{=} comparison and will not sort the result. Therefore \texttt{DuplicateFreeList} can be used even if the elements of \mbox{\texttt{\mdseries\slshape list}} do not lie in the same family. Otherwise, if \mbox{\texttt{\mdseries\slshape list}} contains objects that can be compared with \texttt{\texttt{\symbol{92}}{\textless}} (\ref{bSlash<}) then it is much more efficient to use \texttt{Set} (\ref{Set}) instead of \texttt{DuplicateFreeList}. 

 \texttt{Unique} is a synonym for \texttt{DuplicateFreeList}. 

 
\begin{Verbatim}[commandchars=!@|,fontsize=\small,frame=single,label=Example]
  !gapprompt@gap>| !gapinput@l:=[1,Z(3),1,"abc",Group((1,2,3),(1,2)),Z(3),Group((1,2),(2,3))];;|
  !gapprompt@gap>| !gapinput@DuplicateFreeList( l );|
  [ 1, Z(3), "abc", Group([ (1,2,3), (1,2) ]) ]
\end{Verbatim}
 }

 

\subsection{\textcolor{Chapter }{AsDuplicateFreeList}}
\logpage{[ 21, 20, 5 ]}\nobreak
\hyperdef{L}{X7F5D4DD87E4378AC}{}
{\noindent\textcolor{FuncColor}{$\triangleright$\enspace\texttt{AsDuplicateFreeList({\mdseries\slshape list})\index{AsDuplicateFreeList@\texttt{AsDuplicateFreeList}}
\label{AsDuplicateFreeList}
}\hfill{\scriptsize (attribute)}}\\


 returns the same result as \texttt{DuplicateFreeList} (\ref{DuplicateFreeList}), except that the result is immutable. }

 

\subsection{\textcolor{Chapter }{Flat}}
\logpage{[ 21, 20, 6 ]}\nobreak
\hyperdef{L}{X7FA272D984EF82ED}{}
{\noindent\textcolor{FuncColor}{$\triangleright$\enspace\texttt{Flat({\mdseries\slshape list})\index{Flat@\texttt{Flat}}
\label{Flat}
}\hfill{\scriptsize (operation)}}\\


 returns the list of all elements that are contained in the list \mbox{\texttt{\mdseries\slshape list}} or its sublists. That is, \texttt{Flat} first makes a new empty list \mbox{\texttt{\mdseries\slshape new}}. Then it loops over the elements \mbox{\texttt{\mdseries\slshape elm}} of \mbox{\texttt{\mdseries\slshape list}}. If \mbox{\texttt{\mdseries\slshape elm}} is not a list it is added to \mbox{\texttt{\mdseries\slshape new}}, otherwise \texttt{Flat} appends \texttt{Flat( \mbox{\texttt{\mdseries\slshape elm}} )} to \mbox{\texttt{\mdseries\slshape new}}. 

 
\begin{Verbatim}[commandchars=!@|,fontsize=\small,frame=single,label=Example]
  !gapprompt@gap>| !gapinput@Flat( [ 1, [ 2, 3 ], [ [ 1, 2 ], 3 ] ] );|
  [ 1, 2, 3, 1, 2, 3 ]
  !gapprompt@gap>| !gapinput@Flat( [ ] );|
  [  ]
\end{Verbatim}
 

 To reconstruct a matrix from the list obtained by applying \texttt{Flat} to the matrix, the sublist operator can be used, as follows. 

 
\begin{Verbatim}[commandchars=!@|,fontsize=\small,frame=single,label=Example]
  !gapprompt@gap>| !gapinput@l:=[9..14];;w:=2;; # w is the length of each row|
  !gapprompt@gap>| !gapinput@sub:=[1..w];;List([1..Length(l)/w],i->l{(i-1)*w+sub});|
  [ [ 9, 10 ], [ 11, 12 ], [ 13, 14 ] ]
\end{Verbatim}
 }

 

\subsection{\textcolor{Chapter }{Reversed}}
\logpage{[ 21, 20, 7 ]}\nobreak
\hyperdef{L}{X7C4FDB007C3F54A1}{}
{\noindent\textcolor{FuncColor}{$\triangleright$\enspace\texttt{Reversed({\mdseries\slshape list})\index{Reversed@\texttt{Reversed}}
\label{Reversed}
}\hfill{\scriptsize (function)}}\\


 returns a new mutable list, containing the elements of the dense list \mbox{\texttt{\mdseries\slshape list}} in reversed order. 

 The argument list is unchanged. The result list is a new list, that is not
identical to any other list. The elements of that list however are identical
to the corresponding elements of the argument list (see{\nobreakspace}\ref{Identical Lists}). 

 \texttt{Reversed} implements a special case of list assignment, which can also be formulated in
terms of the \texttt{\texttt{\symbol{123}}\texttt{\symbol{125}}} operator (see{\nobreakspace}\ref{List Assignment}). 

 Developers who wish to adapt this for custom list types need to install
suitable methods for the operation \texttt{ReversedOp}. \index{ReversedOp@\texttt{ReversedOp}} 

 
\begin{Verbatim}[commandchars=!@|,fontsize=\small,frame=single,label=Example]
  !gapprompt@gap>| !gapinput@Reversed( [ 1, 4, 9, 5, 6, 7 ] );|
  [ 7, 6, 5, 9, 4, 1 ]
\end{Verbatim}
 }

 

\subsection{\textcolor{Chapter }{Shuffle}}
\logpage{[ 21, 20, 8 ]}\nobreak
\hyperdef{L}{X8057372F83374193}{}
{\noindent\textcolor{FuncColor}{$\triangleright$\enspace\texttt{Shuffle({\mdseries\slshape list})\index{Shuffle@\texttt{Shuffle}}
\label{Shuffle}
}\hfill{\scriptsize (operation)}}\\


 The argument \mbox{\texttt{\mdseries\slshape list}} must be a dense mutable list. This operation permutes the entries of \mbox{\texttt{\mdseries\slshape list}} randomly (in place), and returns \mbox{\texttt{\mdseries\slshape list}}. 
\begin{Verbatim}[commandchars=!@|,fontsize=\small,frame=single,label=Example]
  !gapprompt@gap>| !gapinput@Reset(GlobalMersenneTwister, 12345);; # make manual tester happy|
  !gapprompt@gap>| !gapinput@l := [1..20];|
  [ 1 .. 20 ]
  !gapprompt@gap>| !gapinput@m := Shuffle(ShallowCopy(l));|
  [ 8, 13, 1, 3, 20, 15, 4, 7, 5, 18, 6, 12, 16, 11, 2, 10, 19, 17, 9,
    14 ]
  !gapprompt@gap>| !gapinput@l;|
  [ 1 .. 20 ]
  !gapprompt@gap>| !gapinput@Shuffle(l);;|
  !gapprompt@gap>| !gapinput@l;|
  [ 19, 5, 7, 20, 16, 1, 10, 15, 12, 11, 13, 2, 14, 3, 4, 17, 6, 8, 9,
    18 ]
\end{Verbatim}
 }

 

\subsection{\textcolor{Chapter }{Apply}}
\logpage{[ 21, 20, 9 ]}\nobreak
\hyperdef{L}{X8075FBDE7B81B4C8}{}
{\noindent\textcolor{FuncColor}{$\triangleright$\enspace\texttt{Apply({\mdseries\slshape list, func})\index{Apply@\texttt{Apply}}
\label{Apply}
}\hfill{\scriptsize (function)}}\\


 \texttt{Apply} applies the function \mbox{\texttt{\mdseries\slshape func}} to every element of the dense and mutable list \mbox{\texttt{\mdseries\slshape list}}, and replaces each element entry by the corresponding return value. 

 \texttt{Apply} changes its argument. The nondestructive counterpart of \texttt{Apply} is \texttt{List} (\ref{List:for a collection}). 

 
\begin{Verbatim}[commandchars=!@|,fontsize=\small,frame=single,label=Example]
  !gapprompt@gap>| !gapinput@l:= [ 1, 2, 3 ];;  Apply( l, i -> i^2 );  l;|
  [ 1, 4, 9 ]
\end{Verbatim}
 }

 

\subsection{\textcolor{Chapter }{Perform}}
\logpage{[ 21, 20, 10 ]}\nobreak
\hyperdef{L}{X7EF6E2BC81DBF6FB}{}
{\noindent\textcolor{FuncColor}{$\triangleright$\enspace\texttt{Perform({\mdseries\slshape list, func})\index{Perform@\texttt{Perform}}
\label{Perform}
}\hfill{\scriptsize (function)}}\\


 \texttt{Perform} applies the function \mbox{\texttt{\mdseries\slshape func}} to every element of the list \mbox{\texttt{\mdseries\slshape list}}, discarding any return values. It does not return a value. 

 
\begin{Verbatim}[commandchars=!@|,fontsize=\small,frame=single,label=Example]
  !gapprompt@gap>| !gapinput@l := [1, 2, 3];; Perform(l,|
  !gapprompt@>| !gapinput@function(x) if IsPrimeInt(x) then Print(x,"\n"); fi; end);|
  2
  3
\end{Verbatim}
 }

 

\subsection{\textcolor{Chapter }{PermListList}}
\logpage{[ 21, 20, 11 ]}\nobreak
\hyperdef{L}{X8763882A7D65F979}{}
{\noindent\textcolor{FuncColor}{$\triangleright$\enspace\texttt{PermListList({\mdseries\slshape list1, list2})\index{PermListList@\texttt{PermListList}}
\label{PermListList}
}\hfill{\scriptsize (function)}}\\


 returns a permutation $p$ of \texttt{[ 1 .. Length( \mbox{\texttt{\mdseries\slshape list1}} ) ]} such that \mbox{\texttt{\mdseries\slshape list1}}$[i$\texttt{\texttt{\symbol{94}}}$p] =$ \mbox{\texttt{\mdseries\slshape list2}}$[i]$. It returns \texttt{fail} if there is no such permutation. 

 
\begin{Verbatim}[commandchars=!@|,fontsize=\small,frame=single,label=Example]
  !gapprompt@gap>| !gapinput@list1 := [ 5, 4, 6, 1, 7, 5 ];;|
  !gapprompt@gap>| !gapinput@list2 := [ 4, 1, 7, 5, 5, 6 ];;|
  !gapprompt@gap>| !gapinput@perm := PermListList(list1, list2);|
  (1,2,4)(3,5,6)
  !gapprompt@gap>| !gapinput@Permuted( list2, perm );|
  [ 5, 4, 6, 1, 7, 5 ]
\end{Verbatim}
 }

 
\subsection{\textcolor{Chapter }{Maximum}}\logpage{[ 21, 20, 12 ]}
\hyperdef{L}{X82CE0DE8828E4303}{}
{
\noindent\textcolor{FuncColor}{$\triangleright$\enspace\texttt{Maximum({\mdseries\slshape obj1, obj2, ...})\index{Maximum@\texttt{Maximum}!for various objects}
\label{Maximum:for various objects}
}\hfill{\scriptsize (function)}}\\
\noindent\textcolor{FuncColor}{$\triangleright$\enspace\texttt{Maximum({\mdseries\slshape list})\index{Maximum@\texttt{Maximum}!for a list}
\label{Maximum:for a list}
}\hfill{\scriptsize (function)}}\\


 In the first form \texttt{Maximum} returns the \emph{maximum} of its arguments, i.e., one argument \mbox{\texttt{\mdseries\slshape obj}} for which $\mbox{\texttt{\mdseries\slshape obj}} \geq \mbox{\texttt{\mdseries\slshape obj1}}$, $\mbox{\texttt{\mdseries\slshape obj}} \geq \mbox{\texttt{\mdseries\slshape obj2}}$ etc. 

 In the second form \texttt{Maximum} takes a homogeneous list \mbox{\texttt{\mdseries\slshape list}} and returns the maximum of the elements in this list. 

 
\begin{Verbatim}[commandchars=!@|,fontsize=\small,frame=single,label=Example]
  !gapprompt@gap>| !gapinput@Maximum( -123, 700, 123, 0, -1000 );|
  700
  !gapprompt@gap>| !gapinput@Maximum( [ -123, 700, 123, 0, -1000 ] );|
  700
  !gapprompt@gap>| !gapinput@# lists are compared elementwise:|
  !gapprompt@gap>| !gapinput@Maximum( [1,2], [0,15], [1,5], [2,-11] );|
  [ 2, -11 ]
\end{Verbatim}
 To get the index of the maximum element use \texttt{PositionMaximum} (\ref{PositionMaximum}) }

 
\subsection{\textcolor{Chapter }{Minimum}}\logpage{[ 21, 20, 13 ]}
\hyperdef{L}{X82F133EC7F89665F}{}
{
\noindent\textcolor{FuncColor}{$\triangleright$\enspace\texttt{Minimum({\mdseries\slshape obj1, obj2, ...})\index{Minimum@\texttt{Minimum}!for various objects}
\label{Minimum:for various objects}
}\hfill{\scriptsize (function)}}\\
\noindent\textcolor{FuncColor}{$\triangleright$\enspace\texttt{Minimum({\mdseries\slshape list})\index{Minimum@\texttt{Minimum}!for a list}
\label{Minimum:for a list}
}\hfill{\scriptsize (function)}}\\


 In the first form \texttt{Minimum} returns the \emph{minimum} of its arguments, i.e., one argument \mbox{\texttt{\mdseries\slshape obj}} for which $\mbox{\texttt{\mdseries\slshape obj}} \leq \mbox{\texttt{\mdseries\slshape obj1}}$, $\mbox{\texttt{\mdseries\slshape obj}} \leq \mbox{\texttt{\mdseries\slshape obj2}}$ etc. 

 In the second form \texttt{Minimum} takes a homogeneous list \mbox{\texttt{\mdseries\slshape list}} and returns the minimum of the elements in this list. 

 Note that for both \texttt{Maximum} (\ref{Maximum:for various objects}) and \texttt{Minimum} the comparison of the objects \mbox{\texttt{\mdseries\slshape obj1}}, \mbox{\texttt{\mdseries\slshape obj2}} etc.{\nobreakspace}must be defined; for that, usually they must lie in the
same family (see{\nobreakspace}\ref{Families}). 

 
\begin{Verbatim}[commandchars=!@|,fontsize=\small,frame=single,label=Example]
  !gapprompt@gap>| !gapinput@Minimum( -123, 700, 123, 0, -1000 );|
  -1000
  !gapprompt@gap>| !gapinput@Minimum( [ -123, 700, 123, 0, -1000 ] );|
  -1000
  !gapprompt@gap>| !gapinput@Minimum( [ 1, 2 ], [ 0, 15 ], [ 1, 5 ], [ 2, -11 ] );|
  [ 0, 15 ]
\end{Verbatim}
 To get the index of the minimum element use \texttt{PositionMinimum} (\ref{PositionMinimum}) }

 
\subsection{\textcolor{Chapter }{MaximumList and MinimumList}}\logpage{[ 21, 20, 14 ]}
\hyperdef{L}{X842851EB7E0969F7}{}
{
\noindent\textcolor{FuncColor}{$\triangleright$\enspace\texttt{MaximumList({\mdseries\slshape list[, seed]})\index{MaximumList@\texttt{MaximumList}}
\label{MaximumList}
}\hfill{\scriptsize (operation)}}\\
\noindent\textcolor{FuncColor}{$\triangleright$\enspace\texttt{MinimumList({\mdseries\slshape list[, seed]})\index{MinimumList@\texttt{MinimumList}}
\label{MinimumList}
}\hfill{\scriptsize (operation)}}\\


 return the maximum resp.{\nobreakspace}the minimum of the elements in the list \mbox{\texttt{\mdseries\slshape list}}. They are the operations called by \texttt{Maximum} (\ref{Maximum:for various objects}) resp.{\nobreakspace}\texttt{Minimum} (\ref{Minimum:for various objects}). Methods can be installed for special kinds of lists. For example, there are
special methods to compute the maximum resp.{\nobreakspace}the minimum of a
range (see{\nobreakspace}\ref{Ranges}). 

 If a second argument \mbox{\texttt{\mdseries\slshape seed}} is supplied, then the result is the maximum resp.{\nobreakspace}minimum of the
union of \mbox{\texttt{\mdseries\slshape list}} and \mbox{\texttt{\mdseries\slshape seed}}. In this manner, the operations may be applied to empty lists. }

 
\subsection{\textcolor{Chapter }{Cartesian}}\logpage{[ 21, 20, 15 ]}
\hyperdef{L}{X7E1593B979BDF2CD}{}
{
\noindent\textcolor{FuncColor}{$\triangleright$\enspace\texttt{Cartesian({\mdseries\slshape list1, list2, ...})\index{Cartesian@\texttt{Cartesian}!for various objects}
\label{Cartesian:for various objects}
}\hfill{\scriptsize (function)}}\\
\noindent\textcolor{FuncColor}{$\triangleright$\enspace\texttt{Cartesian({\mdseries\slshape list})\index{Cartesian@\texttt{Cartesian}!for a list}
\label{Cartesian:for a list}
}\hfill{\scriptsize (function)}}\\


 In the first form \texttt{Cartesian} returns the cartesian product of the lists \mbox{\texttt{\mdseries\slshape list1}}, \mbox{\texttt{\mdseries\slshape list2}}, etc. 

 In the second form \mbox{\texttt{\mdseries\slshape list}} must be a list of lists \mbox{\texttt{\mdseries\slshape list1}}, \mbox{\texttt{\mdseries\slshape list2}}, etc., and \texttt{Cartesian} returns the cartesian product of those lists. 

 The \emph{cartesian product} is a list \mbox{\texttt{\mdseries\slshape cart}} of lists \mbox{\texttt{\mdseries\slshape tup}}, such that the first element of \mbox{\texttt{\mdseries\slshape tup}} is an element of \mbox{\texttt{\mdseries\slshape list1}}, the second element of \mbox{\texttt{\mdseries\slshape tup}} is an element of \mbox{\texttt{\mdseries\slshape list2}}, and so on. The total number of elements in \mbox{\texttt{\mdseries\slshape cart}} is the product of the lengths of the argument lists. In particular \mbox{\texttt{\mdseries\slshape cart}} is empty if and only if at least one of the argument lists is empty. Also \mbox{\texttt{\mdseries\slshape cart}} contains duplicates if and only if no argument list is empty and at least one
contains duplicates. 

 The last index runs fastest. That means that the first element \mbox{\texttt{\mdseries\slshape tup1}} of \mbox{\texttt{\mdseries\slshape cart}} contains the first element from \mbox{\texttt{\mdseries\slshape list1}}, from \mbox{\texttt{\mdseries\slshape list2}} and so on. The second element \mbox{\texttt{\mdseries\slshape tup2}} of \mbox{\texttt{\mdseries\slshape cart}} contains the first element from \mbox{\texttt{\mdseries\slshape list1}}, the first from \mbox{\texttt{\mdseries\slshape list2}}, and so on, but the last element of \mbox{\texttt{\mdseries\slshape tup2}} is the second element of the last argument list. This implies that \mbox{\texttt{\mdseries\slshape cart}} is a proper set if and only if all argument lists are proper sets
(see{\nobreakspace}\ref{Sorted Lists and Sets}). 

 The function \texttt{Tuples} (\ref{Tuples}) computes the \mbox{\texttt{\mdseries\slshape k}}\texttt{\symbol{45}}fold cartesian product of a list. 

 
\begin{Verbatim}[commandchars=!@|,fontsize=\small,frame=single,label=Example]
  !gapprompt@gap>| !gapinput@Cartesian( [1,2], [3,4], [5,6] );|
  [ [ 1, 3, 5 ], [ 1, 3, 6 ], [ 1, 4, 5 ], [ 1, 4, 6 ], [ 2, 3, 5 ],
    [ 2, 3, 6 ], [ 2, 4, 5 ], [ 2, 4, 6 ] ]
  !gapprompt@gap>| !gapinput@Cartesian( [1,2,2], [1,1,2] );|
  [ [ 1, 1 ], [ 1, 1 ], [ 1, 2 ], [ 2, 1 ], [ 2, 1 ], [ 2, 2 ],
    [ 2, 1 ], [ 2, 1 ], [ 2, 2 ] ]
\end{Verbatim}
 }

 
\subsection{\textcolor{Chapter }{IteratorOfCartesianProduct}}\logpage{[ 21, 20, 16 ]}
\hyperdef{L}{X7E76F5A782184823}{}
{
\noindent\textcolor{FuncColor}{$\triangleright$\enspace\texttt{IteratorOfCartesianProduct({\mdseries\slshape list1, list2, ...})\index{IteratorOfCartesianProduct@\texttt{IteratorOfCartesianProduct}!for several lists}
\label{IteratorOfCartesianProduct:for several lists}
}\hfill{\scriptsize (function)}}\\
\noindent\textcolor{FuncColor}{$\triangleright$\enspace\texttt{IteratorOfCartesianProduct({\mdseries\slshape list})\index{IteratorOfCartesianProduct@\texttt{IteratorOfCartesianProduct}!for a list of lists}
\label{IteratorOfCartesianProduct:for a list of lists}
}\hfill{\scriptsize (function)}}\\


 In the first form \texttt{IteratorOfCartesianProduct} returns an iterator (see{\nobreakspace}\ref{Iterators}) of all elements of the cartesian product (see{\nobreakspace}\texttt{Cartesian} (\ref{Cartesian:for a list})) of the lists \mbox{\texttt{\mdseries\slshape list1}}, \mbox{\texttt{\mdseries\slshape list2}}, etc. 

 In the second form \mbox{\texttt{\mdseries\slshape list}} must be a list of lists \mbox{\texttt{\mdseries\slshape list1}}, \mbox{\texttt{\mdseries\slshape list2}}, etc., and \texttt{IteratorOfCartesianProduct} returns an iterator of the cartesian product of those lists. 

 Resulting tuples will be returned in the lexicographic order. Usage of
iterators of cartesian products is recommended in the case when the resulting
cartesian product is big enough, so its generating and storage will require
essential amount of runtime and memory. For smaller cartesian products it is
faster to generate the full set of tuples using \texttt{Cartesian} (\ref{Cartesian:for a list}) and then loop over its elements (with some minor overhead of needing more
memory). }

 

\subsection{\textcolor{Chapter }{Permuted}}
\logpage{[ 21, 20, 17 ]}\nobreak
\hyperdef{L}{X7B5A19098406347A}{}
{\noindent\textcolor{FuncColor}{$\triangleright$\enspace\texttt{Permuted({\mdseries\slshape list, perm})\index{Permuted@\texttt{Permuted}}
\label{Permuted}
}\hfill{\scriptsize (operation)}}\\


 returns a new list \mbox{\texttt{\mdseries\slshape new}} that contains the elements of the list \mbox{\texttt{\mdseries\slshape list}} permuted according to the permutation \mbox{\texttt{\mdseries\slshape perm}}. That is \texttt{\mbox{\texttt{\mdseries\slshape new}}[\mbox{\texttt{\mdseries\slshape i}}\texttt{\symbol{94}}\mbox{\texttt{\mdseries\slshape perm}}] = \mbox{\texttt{\mdseries\slshape list}}[\mbox{\texttt{\mdseries\slshape i}}]} whenever \texttt{\mbox{\texttt{\mdseries\slshape list}}[\mbox{\texttt{\mdseries\slshape i}}]} is bound. 

 \texttt{Sortex} (\ref{Sortex}) allows you to compute a permutation that must be applied to a list in order to
get the sorted list. 

 
\begin{Verbatim}[commandchars=!@|,fontsize=\small,frame=single,label=Example]
  !gapprompt@gap>| !gapinput@Permuted( [ 5, 4, 6, 1, 7, 5 ], (1,3,5,6,4) );|
  [ 1, 4, 5, 5, 6, 7 ]
  !gapprompt@gap>| !gapinput@Permuted( [ 5, 4, 6,, 7, 5 ], (1,3,5,6,4) );|
  [ , 4, 5, 5, 6, 7 ]
\end{Verbatim}
 }

 

\subsection{\textcolor{Chapter }{List (for a list (and a function))}}
\logpage{[ 21, 20, 18 ]}\nobreak
\hyperdef{L}{X86CB7DCE8510F977}{}
{\noindent\textcolor{FuncColor}{$\triangleright$\enspace\texttt{List({\mdseries\slshape list[, func]})\index{List@\texttt{List}!for a list (and a function)}
\label{List:for a list (and a function)}
}\hfill{\scriptsize (function)}}\\


 This function returns a new mutable list \texttt{new} of the same length as the list \mbox{\texttt{\mdseries\slshape list}} (which may have holes). The entry \texttt{new[i]} is unbound if \texttt{\mbox{\texttt{\mdseries\slshape list}}[i]} is unbound. Otherwise \texttt{new[i] = \mbox{\texttt{\mdseries\slshape func}}(\mbox{\texttt{\mdseries\slshape list}}[i])}. If the argument \mbox{\texttt{\mdseries\slshape func}} is omitted, its default is \texttt{IdFunc} (\ref{IdFunc}), so this function does the same as \texttt{ShallowCopy} (\ref{ShallowCopy}) (see also \ref{Duplication of Lists}). 

 Developers who wish to adapt this for custom list or collection types need to
install suitable methods for the operation \texttt{ListOp}. \index{ListOp@\texttt{ListOp}} 

 
\begin{Verbatim}[commandchars=!@|,fontsize=\small,frame=single,label=Example]
  !gapprompt@gap>| !gapinput@List( [1,2,3], i -> i^2 );|
  [ 1, 4, 9 ]
  !gapprompt@gap>| !gapinput@List( [1..10], IsPrime );|
  [ false, true, true, false, true, false, true, false, false, false ]
  !gapprompt@gap>| !gapinput@List([,1,,3,4], x-> x > 2);|
  [ , false,, true, true ]
\end{Verbatim}
 

 (See also \texttt{List} (\ref{List:for a collection}).) }

 

\subsection{\textcolor{Chapter }{Filtered}}
\logpage{[ 21, 20, 19 ]}\nobreak
\hyperdef{L}{X7C86D7F7795125F0}{}
{\noindent\textcolor{FuncColor}{$\triangleright$\enspace\texttt{Filtered({\mdseries\slshape listorcoll, func})\index{Filtered@\texttt{Filtered}}
\label{Filtered}
}\hfill{\scriptsize (function)}}\\


 returns a new list that contains those elements of the list or collection \mbox{\texttt{\mdseries\slshape listorcoll}} (see{\nobreakspace}\ref{Collections}), respectively, for which the unary function \mbox{\texttt{\mdseries\slshape func}} returns \texttt{true}. 

 If the first argument is a list, the order of the elements in the result is
the same as the order of the corresponding elements of this list. If an
element for which \mbox{\texttt{\mdseries\slshape func}} returns \texttt{true} appears several times in the list it will also appear the same number of times
in the result. The argument list may contain holes, they are ignored by \texttt{Filtered}. 

 For each element of \mbox{\texttt{\mdseries\slshape listorcoll}}, \mbox{\texttt{\mdseries\slshape func}} must return either \texttt{true} or \texttt{false}, otherwise an error is signalled. 

 The result is a new list that is not identical to any other list. The elements
of that list however are identical to the corresponding elements of the
argument list (see{\nobreakspace}\ref{Identical Lists}). 

 List assignment using the operator \texttt{\texttt{\symbol{92}}\texttt{\symbol{123}}\texttt{\symbol{92}}\texttt{\symbol{125}}} (\ref{bSlash{bSlash}}) (see{\nobreakspace}\ref{List Assignment}) can be used to extract elements of a list according to indices given in
another list. 

 Developers who wish to adapt this for custom list or collection types need to
install suitable methods for the operation \texttt{FilteredOp}. \index{FilteredOp@\texttt{FilteredOp}} 

 
\begin{Verbatim}[commandchars=!@|,fontsize=\small,frame=single,label=Example]
  !gapprompt@gap>| !gapinput@Filtered( [1..20], IsPrime );|
  [ 2, 3, 5, 7, 11, 13, 17, 19 ]
  !gapprompt@gap>| !gapinput@Filtered( [ 1, 3, 4, -4, 4, 7, 10, 6 ], IsPrimePowerInt );|
  [ 3, 4, 4, 7 ]
  !gapprompt@gap>| !gapinput@Filtered( [ 1, 3, 4, -4, 4, 7, 10, 6 ],|
  !gapprompt@>| !gapinput@             n -> IsPrimePowerInt(n) and n mod 2 <> 0 );|
  [ 3, 7 ]
  !gapprompt@gap>| !gapinput@Filtered( Group( (1,2), (1,2,3) ), x -> Order( x ) = 2 );|
  [ (2,3), (1,2), (1,3) ]
\end{Verbatim}
 }

 

\subsection{\textcolor{Chapter }{Number}}
\logpage{[ 21, 20, 20 ]}\nobreak
\hyperdef{L}{X8179B13D80E935FC}{}
{\noindent\textcolor{FuncColor}{$\triangleright$\enspace\texttt{Number({\mdseries\slshape listorcoll[, func]})\index{Number@\texttt{Number}}
\label{Number}
}\hfill{\scriptsize (function)}}\\


 Called with a list \mbox{\texttt{\mdseries\slshape listorcoll}}, \texttt{Number} returns the number of bound entries in this list. For dense lists \texttt{Number}, \texttt{Length} (\ref{Length}), and \texttt{Size} (\ref{Size}) return the same value; for lists with holes \texttt{Number} returns the number of bound entries, \texttt{Length} (\ref{Length}) returns the largest index of a bound entry, and \texttt{Size} (\ref{Size}) signals an error. 

 Called with two arguments, a list or collection \mbox{\texttt{\mdseries\slshape listorcoll}} and a unary function \mbox{\texttt{\mdseries\slshape func}}, \texttt{Number} returns the number of elements of \mbox{\texttt{\mdseries\slshape listorcoll}} for which \mbox{\texttt{\mdseries\slshape func}} returns \texttt{true}. If an element for which \mbox{\texttt{\mdseries\slshape func}} returns \texttt{true} appears several times in \mbox{\texttt{\mdseries\slshape listorcoll}} it will also be counted the same number of times. 

 For each element of \mbox{\texttt{\mdseries\slshape listorcoll}}, \mbox{\texttt{\mdseries\slshape func}} must return either \texttt{true} or \texttt{false}, otherwise an error is signalled. 

 \texttt{Filtered} (\ref{Filtered}) allows you to extract the elements of a list that have a certain property. 

 Developers who wish to adapt this for custom list or collection types need to
install suitable methods for the operation \texttt{NumberOp}. \index{NumberOp@\texttt{NumberOp}} 

 
\begin{Verbatim}[commandchars=!@|,fontsize=\small,frame=single,label=Example]
  !gapprompt@gap>| !gapinput@Number( [ 2, 3, 5, 7 ] );|
  4
  !gapprompt@gap>| !gapinput@Number( [, 2, 3,, 5,, 7,,,, 11 ] );|
  5
  !gapprompt@gap>| !gapinput@Number( [1..20], IsPrime );|
  8
  !gapprompt@gap>| !gapinput@Number( [ 1, 3, 4, -4, 4, 7, 10, 6 ], IsPrimePowerInt );|
  4
  !gapprompt@gap>| !gapinput@Number( [ 1, 3, 4, -4, 4, 7, 10, 6 ],|
  !gapprompt@>| !gapinput@           n -> IsPrimePowerInt(n) and n mod 2 <> 0 );|
  2
  !gapprompt@gap>| !gapinput@Number( Group( (1,2), (1,2,3) ), x -> Order( x ) = 2 );|
  3
\end{Verbatim}
 }

 

\subsection{\textcolor{Chapter }{First}}
\logpage{[ 21, 20, 21 ]}\nobreak
\hyperdef{L}{X82801DFA84E11272}{}
{\noindent\textcolor{FuncColor}{$\triangleright$\enspace\texttt{First({\mdseries\slshape list[, func]})\index{First@\texttt{First}}
\label{First}
}\hfill{\scriptsize (operation)}}\\


 \texttt{First} returns the first element of the list \mbox{\texttt{\mdseries\slshape list}} for which the unary function \mbox{\texttt{\mdseries\slshape func}} returns \texttt{true}; if \mbox{\texttt{\mdseries\slshape func}} is not given, the first element is returned. \mbox{\texttt{\mdseries\slshape list}} may contain holes. \mbox{\texttt{\mdseries\slshape func}} must return either \texttt{true} or \texttt{false} for each element of \mbox{\texttt{\mdseries\slshape list}}, otherwise an error is signalled. If \mbox{\texttt{\mdseries\slshape func}} returns \texttt{false} for all elements of \mbox{\texttt{\mdseries\slshape list}} then \texttt{First} returns \texttt{fail}. 

 \texttt{PositionProperty} (\ref{PositionProperty}) allows you to find the position of the first element in a list that satisfies
a certain property. 

 Before \textsf{GAP} 4.12, developers who wished to adapt this for custom list types needed to
install suitable methods for the operation \texttt{FirstOp}. This is still possible for backwards compatibility, but \texttt{FirstOp} now is just a synonym for \texttt{First}. 

 
\begin{Verbatim}[commandchars=!@|,fontsize=\small,frame=single,label=Example]
  !gapprompt@gap>| !gapinput@First( [10^7..10^8], IsPrime );|
  10000019
  !gapprompt@gap>| !gapinput@First( [10^5..10^6],|
  !gapprompt@>| !gapinput@     n -> not IsPrime(n) and IsPrimePowerInt(n) );|
  100489
  !gapprompt@gap>| !gapinput@First( [ 1 .. 20 ], x -> x < 0 );|
  fail
  !gapprompt@gap>| !gapinput@First( [ fail ], x -> x = fail );|
  fail
\end{Verbatim}
 }

 

\subsection{\textcolor{Chapter }{Last}}
\logpage{[ 21, 20, 22 ]}\nobreak
\hyperdef{L}{X7E5B62E780421CE9}{}
{\noindent\textcolor{FuncColor}{$\triangleright$\enspace\texttt{Last({\mdseries\slshape list[, func]})\index{Last@\texttt{Last}}
\label{Last}
}\hfill{\scriptsize (function)}}\\


 \texttt{Last} returns the last element of the list \mbox{\texttt{\mdseries\slshape list}} for which the unary function \mbox{\texttt{\mdseries\slshape func}} returns \texttt{true}; if \mbox{\texttt{\mdseries\slshape func}} is not given, the last element is returned. \mbox{\texttt{\mdseries\slshape list}} may contain holes. \mbox{\texttt{\mdseries\slshape func}} must return either \texttt{true} or \texttt{false} for each element of \mbox{\texttt{\mdseries\slshape list}}, otherwise an error is signalled. If \mbox{\texttt{\mdseries\slshape func}} returns \texttt{false} for all elements of \mbox{\texttt{\mdseries\slshape list}} then \texttt{Last} returns \texttt{fail}. 

 Developers who wish to adapt this for custom list types need to install
suitable methods for the operation \texttt{LastOp}. \index{LastOp@\texttt{LastOp}} 

 
\begin{Verbatim}[commandchars=!@|,fontsize=\small,frame=single,label=Example]
  !gapprompt@gap>| !gapinput@Last( [10^7..10^8], IsPrime );|
  99999989
  !gapprompt@gap>| !gapinput@Last( [10^5..10^6],|
  !gapprompt@>| !gapinput@     n -> not IsPrime(n) and IsPrimePowerInt(n) );|
  994009
  !gapprompt@gap>| !gapinput@Last( [ 1 .. 20 ], x -> x < 0 );|
  fail
  !gapprompt@gap>| !gapinput@Last( [ fail ], x -> x = fail );|
  fail
\end{Verbatim}
 }

 

\subsection{\textcolor{Chapter }{ForAll}}
\logpage{[ 21, 20, 23 ]}\nobreak
\hyperdef{L}{X7F06961278166671}{}
{\noindent\textcolor{FuncColor}{$\triangleright$\enspace\texttt{ForAll({\mdseries\slshape listorcoll, func})\index{ForAll@\texttt{ForAll}}
\label{ForAll}
}\hfill{\scriptsize (function)}}\\


 tests whether the unary function \mbox{\texttt{\mdseries\slshape func}} returns \texttt{true} for all elements in the list or collection \mbox{\texttt{\mdseries\slshape listorcoll}}. 

 Developers who wish to adapt this for custom list or collection types need to
install suitable methods for the operation \texttt{ForAllOp}. \index{ForAllOp@\texttt{ForAllOp}} 

 
\begin{Verbatim}[commandchars=!@|,fontsize=\small,frame=single,label=Example]
  !gapprompt@gap>| !gapinput@ForAll( [1..20], IsPrime );|
  false
  !gapprompt@gap>| !gapinput@ForAll( [2,3,4,5,8,9], IsPrimePowerInt );|
  true
  !gapprompt@gap>| !gapinput@ForAll( [2..14], n -> IsPrimePowerInt(n) or n mod 2 = 0 );|
  true
  !gapprompt@gap>| !gapinput@ForAll( Group( (1,2), (1,2,3) ), i -> SignPerm(i) = 1 );|
  false
\end{Verbatim}
 }

 

\subsection{\textcolor{Chapter }{ForAny}}
\logpage{[ 21, 20, 24 ]}\nobreak
\hyperdef{L}{X7AF82E747A8BDA75}{}
{\noindent\textcolor{FuncColor}{$\triangleright$\enspace\texttt{ForAny({\mdseries\slshape listorcoll, func})\index{ForAny@\texttt{ForAny}}
\label{ForAny}
}\hfill{\scriptsize (function)}}\\


 tests whether the unary function \mbox{\texttt{\mdseries\slshape func}} returns \texttt{true} for at least one element in the list or collection \mbox{\texttt{\mdseries\slshape listorcoll}}. 

 Developers who wish to adapt this for custom list or collection types need to
install suitable methods for the operation \texttt{ForAnyOp}. \index{ForAnyOp@\texttt{ForAnyOp}} 

 
\begin{Verbatim}[commandchars=!@|,fontsize=\small,frame=single,label=Example]
  !gapprompt@gap>| !gapinput@ForAny( [1..20], IsPrime );|
  true
  !gapprompt@gap>| !gapinput@ForAny( [2,3,4,5,8,9], IsPrimePowerInt );|
  true
  !gapprompt@gap>| !gapinput@ForAny( [2..14],|
  !gapprompt@>| !gapinput@   n -> IsPrimePowerInt(n) and n mod 5 = 0 and not IsPrime(n) );|
  false
  !gapprompt@gap>| !gapinput@ForAny( Integers, i ->     i > 0|
  !gapprompt@>| !gapinput@                          and ForAll( [0,2..4], j -> IsPrime(i+j) ) );|
  true
\end{Verbatim}
 }

 

\subsection{\textcolor{Chapter }{Product}}
\logpage{[ 21, 20, 25 ]}\nobreak
\hyperdef{L}{X7E5C72F27B657948}{}
{\noindent\textcolor{FuncColor}{$\triangleright$\enspace\texttt{Product({\mdseries\slshape listorcoll[, func][, init]})\index{Product@\texttt{Product}}
\label{Product}
}\hfill{\scriptsize (function)}}\\


 Called with one argument, a dense list or collection \mbox{\texttt{\mdseries\slshape listorcoll}}, \texttt{Product} returns the product of the elements of \mbox{\texttt{\mdseries\slshape listorcoll}} (see{\nobreakspace}\ref{Collections}). 

 Called with a dense list or collection \mbox{\texttt{\mdseries\slshape listorcoll}} and a function \mbox{\texttt{\mdseries\slshape func}}, which must be a function taking one argument, \texttt{Product} applies the function \mbox{\texttt{\mdseries\slshape func}} to the elements of \mbox{\texttt{\mdseries\slshape listorcoll}}, and returns the product of the results. In either case \texttt{Product} returns \texttt{1} if the first argument is empty. 

 The general rules for arithmetic operations apply (see{\nobreakspace}\ref{Mutability Status and List Arithmetic}), so the result is immutable if and only if all summands are immutable. 

 If \mbox{\texttt{\mdseries\slshape listorcoll}} contains exactly one element then this element (or its image under \mbox{\texttt{\mdseries\slshape func}} if applicable) itself is returned, not a shallow copy of this element. 

 If an additional initial value \mbox{\texttt{\mdseries\slshape init}} is given, \texttt{Product} returns the product of \mbox{\texttt{\mdseries\slshape init}} and the elements of the first argument resp.{\nobreakspace}of their images
under the function \mbox{\texttt{\mdseries\slshape func}}. This is useful for example if the first argument is empty and a different
identity than \texttt{1} is desired, in which case \mbox{\texttt{\mdseries\slshape init}} is returned. 

 Developers who wish to adapt this for custom list or collection types need to
install suitable methods for the operation \texttt{ProductOp}. \index{ProductOp@\texttt{ProductOp}} 

 
\begin{Verbatim}[commandchars=!@|,fontsize=\small,frame=single,label=Example]
  !gapprompt@gap>| !gapinput@Product( [ 2, 3, 5, 7, 11, 13, 17, 19 ] );|
  9699690
  !gapprompt@gap>| !gapinput@Product( [1..10], x->x^2 );|
  13168189440000
  !gapprompt@gap>| !gapinput@Product( [ (1,2), (1,3), (1,4), (2,3), (2,4), (3,4) ] );|
  (1,4)(2,3)
  !gapprompt@gap>| !gapinput@Product( GF(8) );|
  0*Z(2)
\end{Verbatim}
 }

 

\subsection{\textcolor{Chapter }{Sum}}
\logpage{[ 21, 20, 26 ]}\nobreak
\hyperdef{L}{X7A04B71C84CFCC2D}{}
{\noindent\textcolor{FuncColor}{$\triangleright$\enspace\texttt{Sum({\mdseries\slshape listorcoll[, func][, init]})\index{Sum@\texttt{Sum}}
\label{Sum}
}\hfill{\scriptsize (function)}}\\


 Called with one argument, a dense list or collection \mbox{\texttt{\mdseries\slshape listorcoll}}, \texttt{Sum} returns the sum of the elements of \mbox{\texttt{\mdseries\slshape listorcoll}} (see{\nobreakspace}\ref{Collections}). 

 Called with a dense list or collection \mbox{\texttt{\mdseries\slshape listorcoll}} and a function \mbox{\texttt{\mdseries\slshape func}}, which must be a function taking one argument, \texttt{Sum} applies the function \mbox{\texttt{\mdseries\slshape func}} to the elements of \mbox{\texttt{\mdseries\slshape listorcoll}}, and returns the sum of the results. In either case \texttt{Sum} returns \texttt{0} if the first argument is empty. 

 The general rules for arithmetic operations apply (see{\nobreakspace}\ref{Mutability Status and List Arithmetic}), so the result is immutable if and only if all summands are immutable. 

 If \mbox{\texttt{\mdseries\slshape listorcoll}} contains exactly one element then this element (or its image under \mbox{\texttt{\mdseries\slshape func}} if applicable) itself is returned, not a shallow copy of this element. 

 If an additional initial value \mbox{\texttt{\mdseries\slshape init}} is given, \texttt{Sum} returns the sum of \mbox{\texttt{\mdseries\slshape init}} and the elements of the first argument resp.{\nobreakspace}of their images
under the function \mbox{\texttt{\mdseries\slshape func}}. This is useful for example if the first argument is empty and a different
zero than \texttt{0} is desired, in which case \mbox{\texttt{\mdseries\slshape init}} is returned. 

 Developers who wish to adapt this for custom list or collection types need to
install suitable methods for the operation \texttt{SumOp}. \index{SumOp@\texttt{SumOp}} 

 
\begin{Verbatim}[commandchars=!@|,fontsize=\small,frame=single,label=Example]
  !gapprompt@gap>| !gapinput@Sum( [ 2, 3, 5, 7, 11, 13, 17, 19 ] );|
  77
  !gapprompt@gap>| !gapinput@Sum( [1..10], x->x^2 );|
  385
  !gapprompt@gap>| !gapinput@Sum( [ [1,2], [3,4], [5,6] ] );|
  [ 9, 12 ]
  !gapprompt@gap>| !gapinput@Sum( GF(8) );|
  0*Z(2)
\end{Verbatim}
 }

 

\subsection{\textcolor{Chapter }{Iterated}}
\logpage{[ 21, 20, 27 ]}\nobreak
\hyperdef{L}{X834E4DF57F3A20F0}{}
{\noindent\textcolor{FuncColor}{$\triangleright$\enspace\texttt{Iterated({\mdseries\slshape list, f})\index{Iterated@\texttt{Iterated}}
\label{Iterated}
}\hfill{\scriptsize (operation)}}\\


 returns the result of the iterated application of the function \mbox{\texttt{\mdseries\slshape f}}, which must take two arguments, to the elements of the list \mbox{\texttt{\mdseries\slshape list}}. More precisely, if \mbox{\texttt{\mdseries\slshape list}} has length $n$ then \texttt{Iterated} returns the result of the following application, $\mbox{\texttt{\mdseries\slshape f}}( \ldots \mbox{\texttt{\mdseries\slshape f}}( \mbox{\texttt{\mdseries\slshape f}}( \mbox{\texttt{\mdseries\slshape list}}[1], \mbox{\texttt{\mdseries\slshape list}}[2] ), \mbox{\texttt{\mdseries\slshape list}}[3] ), \ldots, \mbox{\texttt{\mdseries\slshape list}}[n] )$. 

 
\begin{Verbatim}[commandchars=!@|,fontsize=\small,frame=single,label=Example]
  !gapprompt@gap>| !gapinput@Iterated( [ 126, 66, 105 ], Gcd );|
  3
\end{Verbatim}
 }

 

\subsection{\textcolor{Chapter }{ListN}}
\logpage{[ 21, 20, 28 ]}\nobreak
\hyperdef{L}{X7D150C2881881139}{}
{\noindent\textcolor{FuncColor}{$\triangleright$\enspace\texttt{ListN({\mdseries\slshape list1, list2, ..., listn, f})\index{ListN@\texttt{ListN}}
\label{ListN}
}\hfill{\scriptsize (function)}}\\


 applies the $n$\texttt{\symbol{45}}argument function \mbox{\texttt{\mdseries\slshape f}} to the lists. That is, \texttt{ListN} returns the list whose $i$\texttt{\symbol{45}}th entry is $\mbox{\texttt{\mdseries\slshape f}}(\mbox{\texttt{\mdseries\slshape list1}}[i], \mbox{\texttt{\mdseries\slshape list2}}[i], \ldots, \mbox{\texttt{\mdseries\slshape listn}}[i])$. 

 
\begin{Verbatim}[commandchars=!@|,fontsize=\small,frame=single,label=Example]
  !gapprompt@gap>| !gapinput@ListN( [1,2], [3,4], \+ );|
  [ 4, 6 ]
\end{Verbatim}
 }

 }

  
\section{\textcolor{Chapter }{Advanced List Manipulations}}\label{Advanced List Manipulations}
\logpage{[ 21, 21, 0 ]}
\hyperdef{L}{X805CA0B68029B47A}{}
{
  The following functions are generalizations of \texttt{List} (\ref{List:for a collection}), \texttt{Set} (\ref{Set}), \texttt{Sum} (\ref{Sum}), \texttt{Product} (\ref{Product}), \texttt{ForAll} (\ref{ForAll}), \texttt{ForAny} (\ref{ForAny}), \texttt{Filtered} (\ref{Filtered}), \texttt{Number} (\ref{Number}), and \texttt{Perform} (\ref{Perform}). 

\subsection{\textcolor{Chapter }{FoldLeft}}
\logpage{[ 21, 21, 1 ]}\nobreak
\hyperdef{L}{X799A66417CDC1D4F}{}
{\noindent\textcolor{FuncColor}{$\triangleright$\enspace\texttt{FoldLeft({\mdseries\slshape C, func[, init]})\index{FoldLeft@\texttt{FoldLeft}}
\label{FoldLeft}
}\hfill{\scriptsize (operation)}}\\


 \texttt{FoldLeft} applies the binary function \mbox{\texttt{\mdseries\slshape func}} left\texttt{\symbol{45}}associatively to the elements of the list or
collection \mbox{\texttt{\mdseries\slshape C}}. 

 If no initial accumulator \mbox{\texttt{\mdseries\slshape init}} is given, the first element of \mbox{\texttt{\mdseries\slshape C}} is used as the initial accumulator and the remaining elements are combined
with it from left to right. In that case \mbox{\texttt{\mdseries\slshape C}} must not be empty. 

 If \mbox{\texttt{\mdseries\slshape init}} is given, the accumulator is initialized with \mbox{\texttt{\mdseries\slshape init}}. Then an empty list or collection is allowed, and \mbox{\texttt{\mdseries\slshape init}} is returned unchanged in that case. 

 If \mbox{\texttt{\mdseries\slshape C}} is a list with holes, these holes are ignored, just as when iterating over the
list with a \texttt{for}\texttt{\symbol{45}}loop. 

 This operation can be used to express reductions similar to \texttt{Sum} (\ref{Sum}) or \texttt{Product} (\ref{Product}), but it is more general because the accumulator is updated by an arbitrary
binary function. 

 
\begin{Verbatim}[commandchars=!@|,fontsize=\small,frame=single,label=Example]
  !gapprompt@gap>| !gapinput@FoldLeft( [ 1 .. 10 ], \+ );|
  55
  !gapprompt@gap>| !gapinput@FoldLeft( [ 1 .. 4 ], \*, 2 );|
  48
  !gapprompt@gap>| !gapinput@FoldLeft( [ 1,, 3 ], { a, b } -> a + b );|
  4
\end{Verbatim}
 }

 

\subsection{\textcolor{Chapter }{FoldLeftX}}
\logpage{[ 21, 21, 2 ]}\nobreak
\hyperdef{L}{X83A2CF7A7F6C1DB9}{}
{\noindent\textcolor{FuncColor}{$\triangleright$\enspace\texttt{FoldLeftX({\mdseries\slshape gens, func, init[, abortValue]})\index{FoldLeftX@\texttt{FoldLeftX}}
\label{FoldLeftX}
}\hfill{\scriptsize (function)}}\\


 \texttt{FoldLeftX} is a generalization of \texttt{FoldLeft} (\ref{FoldLeft}) which applies an accumulation function \mbox{\texttt{\mdseries\slshape func}} to a bunch of inputs which are derived from \mbox{\texttt{\mdseries\slshape gens}}. 

 The argument \mbox{\texttt{\mdseries\slshape gens}} must be a plain list whose entries describe a sequence of nested loops and
filters. Specifically, let \texttt{n} denote the length of \mbox{\texttt{\mdseries\slshape gens}}. Then each of the entries \mbox{\texttt{\mdseries\slshape gens}}\texttt{[1]}, $\ldots$ \mbox{\texttt{\mdseries\slshape gens}}\texttt{[n]} must be one of the following: 
\begin{description}
\item[{a list or collection}]  this introduces a new for\texttt{\symbol{45}}loop in the sequence of nested
for\texttt{\symbol{45}}loops and if\texttt{\symbol{45}}statements; 
\item[{a function returning a list or collection}]  this introduces a new for\texttt{\symbol{45}}loop in the sequence of nested
for\texttt{\symbol{45}}loops and if\texttt{\symbol{45}}statements, where the
loop\texttt{\symbol{45}}range depends on the values of the outer
loop\texttt{\symbol{45}}variables; or 
\item[{a function returning \texttt{true} or \texttt{false}}]  this introduces a new if\texttt{\symbol{45}}statement in the sequence of
nested for\texttt{\symbol{45}}loops and if\texttt{\symbol{45}}statements. 
\end{description}
 

 Each function occurring in \mbox{\texttt{\mdseries\slshape gens}} is called with the values selected by the preceding entries of \mbox{\texttt{\mdseries\slshape gens}}, in order. Thus the function in \mbox{\texttt{\mdseries\slshape gens}}\texttt{[k]} receives one argument for each of \mbox{\texttt{\mdseries\slshape gens}}\texttt{[1]}, $\ldots$, \mbox{\texttt{\mdseries\slshape gens}}\texttt{[k\texttt{\symbol{45}}1]} that introduced a loop variable. 

 The argument \mbox{\texttt{\mdseries\slshape func}} must be a binary function, whose first argument is an accumulator variable,
and the second argument is the tuple of values of the
loop\texttt{\symbol{45}}variables. 

 The accumulator is initialized with \mbox{\texttt{\mdseries\slshape init}}. For every tuple that is produced by the loops and filters described by \mbox{\texttt{\mdseries\slshape gens}}, \mbox{\texttt{\mdseries\slshape func}} is called with the current accumulator and that tuple. Its return value
becomes the new accumulator. 

 If the optional argument \mbox{\texttt{\mdseries\slshape abortValue}} is given, iteration stops as soon as the accumulator becomes identical to \mbox{\texttt{\mdseries\slshape abortValue}}. This is useful for short\texttt{\symbol{45}}circuiting computations such as \texttt{ForAllX} (\ref{ForAllX}) and \texttt{ForAnyX} (\ref{ForAnyX}). 

 
\begin{Verbatim}[commandchars=!@|,fontsize=\small,frame=single,label=Example]
  !gapprompt@gap>| !gapinput@FoldLeftX( [ [ 1 .. 3 ], i -> [ 1 .. i ],|
  !gapprompt@>| !gapinput@               { i, j } -> i <> j ],|
  !gapprompt@>| !gapinput@             { acc, x } -> acc + 1, 0 );|
  3
  !gapprompt@gap>| !gapinput@FoldLeftX( [ [ 1 .. 3 ], [ 1 .. 3 ], \< ],|
  !gapprompt@>| !gapinput@             function( acc, x )|
  !gapprompt@>| !gapinput@               Add( acc, ShallowCopy( x ) );|
  !gapprompt@>| !gapinput@               return acc;|
  !gapprompt@>| !gapinput@             end, [] );|
  [ [ 1, 2 ], [ 1, 3 ], [ 2, 3 ] ]
\end{Verbatim}
 }

 

\subsection{\textcolor{Chapter }{ListX}}
\logpage{[ 21, 21, 3 ]}\nobreak
\hyperdef{L}{X8258477D7F72171B}{}
{\noindent\textcolor{FuncColor}{$\triangleright$\enspace\texttt{ListX({\mdseries\slshape arg1, arg2, ..., argn, func})\index{ListX@\texttt{ListX}}
\label{ListX}
}\hfill{\scriptsize (function)}}\\


 \texttt{ListX} returns a new list constructed from the arguments. The general mechanism
underlying this and the related ``X''\texttt{\symbol{45}}functions is described in \texttt{FoldLeftX} (\ref{FoldLeftX}). 

 Each of the arguments \mbox{\texttt{\mdseries\slshape arg1}}, \mbox{\texttt{\mdseries\slshape arg2}}, $\ldots$ \mbox{\texttt{\mdseries\slshape argn}} must be one of the following: 
\begin{description}
\item[{a list or collection}]  this introduces a new for\texttt{\symbol{45}}loop in the sequence of nested
for\texttt{\symbol{45}}loops and if\texttt{\symbol{45}}statements; 
\item[{a function returning a list or collection}]  this introduces a new for\texttt{\symbol{45}}loop in the sequence of nested
for\texttt{\symbol{45}}loops and if\texttt{\symbol{45}}statements, where the
loop\texttt{\symbol{45}}range depends on the values of the outer
loop\texttt{\symbol{45}}variables; or 
\item[{a function returning \texttt{true} or \texttt{false}}]  this introduces a new if\texttt{\symbol{45}}statement in the sequence of
nested for\texttt{\symbol{45}}loops and if\texttt{\symbol{45}}statements. 
\end{description}
 

 Each function occurring among the arguments is called with the values selected
by the preceding arguments, in order. Thus the function in the \mbox{\texttt{\mdseries\slshape k}}\texttt{\symbol{45}}th argument position receives one argument for each
preceding argument that introduced a loop variable. 

 The last argument \mbox{\texttt{\mdseries\slshape func}} must be a function, it is applied to the values of the
loop\texttt{\symbol{45}}variables and the results are collected. 

 Thus \texttt{ListX( \mbox{\texttt{\mdseries\slshape list}}, \mbox{\texttt{\mdseries\slshape func}} )} is the same as \texttt{List( \mbox{\texttt{\mdseries\slshape list}}, \mbox{\texttt{\mdseries\slshape func}} )}, and \texttt{ListX( \mbox{\texttt{\mdseries\slshape list}}, \mbox{\texttt{\mdseries\slshape func}}, x \texttt{\symbol{45}}{\textgreater} x )} is the same as \texttt{Filtered( \mbox{\texttt{\mdseries\slshape list}}, \mbox{\texttt{\mdseries\slshape func}} )}. 

 As a more elaborate example, assume \mbox{\texttt{\mdseries\slshape arg1}} is a list or collection, \mbox{\texttt{\mdseries\slshape arg2}} is a function returning \texttt{true} or \texttt{false}, \mbox{\texttt{\mdseries\slshape arg3}} is a function returning a list or collection, and \mbox{\texttt{\mdseries\slshape arg4}} is another function returning \texttt{true} or \texttt{false}, then 

 \texttt{\mbox{\texttt{\mdseries\slshape result}} := ListX( \mbox{\texttt{\mdseries\slshape arg1}}, \mbox{\texttt{\mdseries\slshape arg2}}, \mbox{\texttt{\mdseries\slshape arg3}}, \mbox{\texttt{\mdseries\slshape arg4}}, \mbox{\texttt{\mdseries\slshape func}} );} 

 is equivalent to 

 
\begin{Verbatim}[commandchars=!@|,fontsize=\small,frame=single,label=]
  result := [];
  for v1 in arg1 do
    if arg2( v1 ) then
      for v2 in arg3( v1 ) do
        if arg4( v1, v2 ) then
          Add( result, func( v1, v2 ) );
        fi;
      od;
    fi;
  od;
\end{Verbatim}
 

 The following example shows how \texttt{ListX} can be used to compute all pairs and all strictly sorted pairs of elements in
a list. 

 
\begin{Verbatim}[commandchars=!@|,fontsize=\small,frame=single,label=Example]
  !gapprompt@gap>| !gapinput@l:= [ 1, 2, 3, 4 ];;|
  !gapprompt@gap>| !gapinput@pair:= function( x, y ) return [ x, y ]; end;;|
  !gapprompt@gap>| !gapinput@ListX( l, l, {x,y} -> [x,y] );|
  [ [ 1, 1 ], [ 1, 2 ], [ 1, 3 ], [ 1, 4 ], [ 2, 1 ], [ 2, 2 ],
    [ 2, 3 ], [ 2, 4 ], [ 3, 1 ], [ 3, 2 ], [ 3, 3 ], [ 3, 4 ],
    [ 4, 1 ], [ 4, 2 ], [ 4, 3 ], [ 4, 4 ] ]
\end{Verbatim}
 

 In the following example, \texttt{\texttt{\symbol{92}}{\textless}} (\ref{bSlash<}) is the comparison operation: 

 
\begin{Verbatim}[commandchars=!@|,fontsize=\small,frame=single,label=Example]
  !gapprompt@gap>| !gapinput@ListX( l, l, \<, pair );|
  [ [ 1, 2 ], [ 1, 3 ], [ 1, 4 ], [ 2, 3 ], [ 2, 4 ], [ 3, 4 ] ]
\end{Verbatim}
 }

 

\subsection{\textcolor{Chapter }{SetX}}
\logpage{[ 21, 21, 4 ]}\nobreak
\hyperdef{L}{X7AC321B87A2DCAF5}{}
{\noindent\textcolor{FuncColor}{$\triangleright$\enspace\texttt{SetX({\mdseries\slshape arg1, arg2, ..., func})\index{SetX@\texttt{SetX}}
\label{SetX}
}\hfill{\scriptsize (function)}}\\


 The only difference between \texttt{SetX} and \texttt{ListX} (\ref{ListX}) is that the result list of \texttt{SetX} is strictly sorted. 

 
\begin{Verbatim}[commandchars=!@|,fontsize=\small,frame=single,label=Example]
  !gapprompt@gap>| !gapinput@SetX( [ 1 .. 3 ], [ 1 .. 3 ], \+ );|
  [ 2, 3, 4, 5, 6 ]
\end{Verbatim}
 }

 

\subsection{\textcolor{Chapter }{SumX}}
\logpage{[ 21, 21, 5 ]}\nobreak
\hyperdef{L}{X82B1411E7FBE925F}{}
{\noindent\textcolor{FuncColor}{$\triangleright$\enspace\texttt{SumX({\mdseries\slshape arg1, arg2, ..., func})\index{SumX@\texttt{SumX}}
\label{SumX}
}\hfill{\scriptsize (function)}}\\


 \texttt{SumX} returns the sum of the elements in the list obtained by \texttt{ListX} (\ref{ListX}) when this is called with the same arguments. 

 
\begin{Verbatim}[commandchars=!@|,fontsize=\small,frame=single,label=Example]
  !gapprompt@gap>| !gapinput@SumX( [ 1 .. 3 ], [ 1 .. 3 ], \+ );|
  36
\end{Verbatim}
 }

 

\subsection{\textcolor{Chapter }{ProductX}}
\logpage{[ 21, 21, 6 ]}\nobreak
\hyperdef{L}{X7FB318B47D8783DA}{}
{\noindent\textcolor{FuncColor}{$\triangleright$\enspace\texttt{ProductX({\mdseries\slshape arg1, arg2, ..., func})\index{ProductX@\texttt{ProductX}}
\label{ProductX}
}\hfill{\scriptsize (function)}}\\


 \texttt{ProductX} returns the product of the elements in the list obtained by \texttt{ListX} (\ref{ListX}) when this is called with the same arguments. 

 
\begin{Verbatim}[commandchars=!@|,fontsize=\small,frame=single,label=Example]
  !gapprompt@gap>| !gapinput@ProductX( [ 1 .. 3 ], [ 1 .. 3 ], \+ );|
  172800
\end{Verbatim}
 }

 

\subsection{\textcolor{Chapter }{ForAllX}}
\logpage{[ 21, 21, 7 ]}\nobreak
\hyperdef{L}{X86CF6FC87C5C9A08}{}
{\noindent\textcolor{FuncColor}{$\triangleright$\enspace\texttt{ForAllX({\mdseries\slshape arg1, arg2, ..., func})\index{ForAllX@\texttt{ForAllX}}
\label{ForAllX}
}\hfill{\scriptsize (function)}}\\


 \texttt{ForAllX} returns \texttt{true} if all elements are \texttt{true} in the list obtained by calling \texttt{ListX} (\ref{ListX}) with the same arguments. Otherwise \texttt{false} is returned. As with \texttt{ForAll} (\ref{ForAll}), the last argument must return either \texttt{true} or \texttt{false}, and evaluation stops as soon as the result is known. 

 
\begin{Verbatim}[commandchars=!@|,fontsize=\small,frame=single,label=Example]
  !gapprompt@gap>| !gapinput@ForAllX( [ 1 .. 3 ], [ 1 .. 3 ], \< );|
  false
\end{Verbatim}
 }

 

\subsection{\textcolor{Chapter }{ForAnyX}}
\logpage{[ 21, 21, 8 ]}\nobreak
\hyperdef{L}{X858C7CEF7EC1260C}{}
{\noindent\textcolor{FuncColor}{$\triangleright$\enspace\texttt{ForAnyX({\mdseries\slshape arg1, arg2, ..., func})\index{ForAnyX@\texttt{ForAnyX}}
\label{ForAnyX}
}\hfill{\scriptsize (function)}}\\


 \texttt{ForAnyX} returns \texttt{true} if any element is \texttt{true} in the list obtained by calling \texttt{ListX} (\ref{ListX}) with the same arguments. Otherwise \texttt{false} is returned. As with \texttt{ForAny} (\ref{ForAny}), the last argument must return either \texttt{true} or \texttt{false}, and evaluation stops as soon as the result is known. 

 
\begin{Verbatim}[commandchars=!@|,fontsize=\small,frame=single,label=Example]
  !gapprompt@gap>| !gapinput@ForAnyX( [ 1 .. 3 ], [ 1 .. 3 ], \< );|
  true
\end{Verbatim}
 }

 

\subsection{\textcolor{Chapter }{FilteredX}}
\logpage{[ 21, 21, 9 ]}\nobreak
\hyperdef{L}{X877CF5DE7AE12506}{}
{\noindent\textcolor{FuncColor}{$\triangleright$\enspace\texttt{FilteredX({\mdseries\slshape arg1, arg2, ..., func})\index{FilteredX@\texttt{FilteredX}}
\label{FilteredX}
}\hfill{\scriptsize (function)}}\\


 \texttt{FilteredX} returns the tuples of loop\texttt{\symbol{45}}variable values for which the
last argument returns \texttt{true}. In other words, it keeps precisely those tuples that would pass the filters
when calling \texttt{ListX} (\ref{ListX}) with the same arguments. 

 Even with only one generator, the result consists of singleton tuples. Thus \texttt{FilteredX( [ 1 .. 4 ], IsEvenInt )} returns \texttt{[ [ 2 ], [ 4 ] ]}, not \texttt{[ 2, 4 ]}. Use \texttt{Filtered} (\ref{Filtered}) if you want to filter the elements of a single list or collection directly. 

 
\begin{Verbatim}[commandchars=!@|,fontsize=\small,frame=single,label=Example]
  !gapprompt@gap>| !gapinput@FilteredX( [ 1 .. 4 ], IsEvenInt );|
  [ [ 2 ], [ 4 ] ]
  !gapprompt@gap>| !gapinput@FilteredX( [ 1 .. 4 ], [ 1 .. 4 ], \< );|
  [ [ 1, 2 ], [ 1, 3 ], [ 1, 4 ], [ 2, 3 ], [ 2, 4 ], [ 3, 4 ] ]
\end{Verbatim}
 }

 

\subsection{\textcolor{Chapter }{NumberX}}
\logpage{[ 21, 21, 10 ]}\nobreak
\hyperdef{L}{X7D6C138F84A3C985}{}
{\noindent\textcolor{FuncColor}{$\triangleright$\enspace\texttt{NumberX({\mdseries\slshape arg1, arg2, ..., func})\index{NumberX@\texttt{NumberX}}
\label{NumberX}
}\hfill{\scriptsize (function)}}\\


 \texttt{NumberX} returns the number of tuples selected by the last argument, using the same
arguments as \texttt{ListX} (\ref{ListX}). Equivalently, it counts the entries of the result that \texttt{FilteredX} (\ref{FilteredX}) would return with the same arguments. 

 
\begin{Verbatim}[commandchars=!@|,fontsize=\small,frame=single,label=Example]
  !gapprompt@gap>| !gapinput@NumberX( [ 1 .. 4 ], IsEvenInt );|
  2
  !gapprompt@gap>| !gapinput@NumberX( [ 1 .. 4 ], [ 1 .. 4 ], \< );|
  6
\end{Verbatim}
 }

 

\subsection{\textcolor{Chapter }{PerformX}}
\logpage{[ 21, 21, 11 ]}\nobreak
\hyperdef{L}{X82580BC887390C69}{}
{\noindent\textcolor{FuncColor}{$\triangleright$\enspace\texttt{PerformX({\mdseries\slshape arg1, arg2, ..., func})\index{PerformX@\texttt{PerformX}}
\label{PerformX}
}\hfill{\scriptsize (function)}}\\


 \texttt{PerformX} works like \texttt{ListX} (\ref{ListX}) except that it returns nothing and ignores the return values of \mbox{\texttt{\mdseries\slshape func}}. It is useful for iterating through the tuples described by the arguments
purely for their side effects. 

 
\begin{Verbatim}[commandchars=!@|,fontsize=\small,frame=single,label=Example]
  !gapprompt@gap>| !gapinput@PerformX( [ 1 .. 2 ], [ 3 .. 4 ], Print );  Print( "\n" );|
  13142324
\end{Verbatim}
 }

 }

  
\section{\textcolor{Chapter }{Ranges}}\label{Ranges}
\logpage{[ 21, 22, 0 ]}
\hyperdef{L}{X79596BDE7CAF8491}{}
{
  \index{range} A \emph{range} is a dense list of integers in arithmetic progression (or degression). This is
a list of integers such that the difference between consecutive elements is a
nonzero constant. Ranges can be abbreviated with the syntactic construct 

 \texttt{[ \mbox{\texttt{\mdseries\slshape first}}, \mbox{\texttt{\mdseries\slshape second}} .. \mbox{\texttt{\mdseries\slshape last}} ]} 

 or, if the difference between consecutive elements is 1, as 

 \texttt{[ \mbox{\texttt{\mdseries\slshape first}} .. \mbox{\texttt{\mdseries\slshape last}} ]}. 

 If \texttt{\mbox{\texttt{\mdseries\slshape first}} {\textgreater} \mbox{\texttt{\mdseries\slshape last}}}, \texttt{[ \mbox{\texttt{\mdseries\slshape first}} .. \mbox{\texttt{\mdseries\slshape last}} ]} is the empty list, which by definition is also a range; also, if \texttt{\mbox{\texttt{\mdseries\slshape second}} {\textgreater} \mbox{\texttt{\mdseries\slshape first}} {\textgreater} \mbox{\texttt{\mdseries\slshape last}}} or \texttt{\mbox{\texttt{\mdseries\slshape second}} {\textless} \mbox{\texttt{\mdseries\slshape first}} {\textless} \mbox{\texttt{\mdseries\slshape last}}}, then \texttt{[ \mbox{\texttt{\mdseries\slshape first}}, \mbox{\texttt{\mdseries\slshape second}} .. \mbox{\texttt{\mdseries\slshape last}} ]} is the empty list. If \texttt{\mbox{\texttt{\mdseries\slshape first}} = \mbox{\texttt{\mdseries\slshape last}}}, \texttt{[ \mbox{\texttt{\mdseries\slshape first}}, \mbox{\texttt{\mdseries\slshape second}} .. \mbox{\texttt{\mdseries\slshape last}} ]} is a singleton list, which is a range, too. Note that \texttt{\mbox{\texttt{\mdseries\slshape last}} \texttt{\symbol{45}} \mbox{\texttt{\mdseries\slshape first}}} must be divisible by the increment \texttt{\mbox{\texttt{\mdseries\slshape second}} \texttt{\symbol{45}} \mbox{\texttt{\mdseries\slshape first}}}, otherwise an error is signalled. 

 Currently, the integers \mbox{\texttt{\mdseries\slshape first}}, \mbox{\texttt{\mdseries\slshape second}} and \mbox{\texttt{\mdseries\slshape last}} and the length of a range must be \emph{small integers} as defined in chapter{\nobreakspace}\ref{Integers}. 

 Note also that a range is just a special case of a list. Thus you can access
elements in a range (see \ref{List Elements}), test for membership etc. You can even assign to such a range if it is
mutable (see{\nobreakspace}\ref{List Assignment}). Of course, unless you assign \texttt{\mbox{\texttt{\mdseries\slshape last}} + \mbox{\texttt{\mdseries\slshape second}} \texttt{\symbol{45}} \mbox{\texttt{\mdseries\slshape first}}} to the entry \texttt{\mbox{\texttt{\mdseries\slshape range}}[ Length( \mbox{\texttt{\mdseries\slshape range}} ) + 1 ]}, the resulting list will no longer be a range. Note that assigning to an
entry of \mbox{\texttt{\mdseries\slshape range}} will convert it back into a plain list. 

 
\begin{Verbatim}[commandchars=!@|,fontsize=\small,frame=single,label=Example]
  !gapprompt@gap>| !gapinput@r := [10..20];|
  [ 10 .. 20 ]
  !gapprompt@gap>| !gapinput@Length( r );|
  11
  !gapprompt@gap>| !gapinput@r[3];|
  12
  !gapprompt@gap>| !gapinput@17 in r;|
  true
  !gapprompt@gap>| !gapinput@# r still is a range but is now represented as a plain list|
  !gapprompt@gap>| !gapinput@r[1] := 10;; r;|
  [ 10, 11, 12, 13, 14, 15, 16, 17, 18, 19, 20 ]
  !gapprompt@gap>| !gapinput@IsRange(r);|
  true
  !gapprompt@gap>| !gapinput@# r is no longer a range|
  !gapprompt@gap>| !gapinput@r[12] := 25;;|
  !gapprompt@gap>| !gapinput@IsRange(r);|
  false
  !gapprompt@gap>| !gapinput@r := [1,3..17];|
  [ 1, 3 .. 17 ]
  !gapprompt@gap>| !gapinput@Length( r );|
  9
  !gapprompt@gap>| !gapinput@r[4];|
  7
  !gapprompt@gap>| !gapinput@r := [0,-1..-9];|
  [ 0, -1 .. -9 ]
  !gapprompt@gap>| !gapinput@r[5];|
  -4
  !gapprompt@gap>| !gapinput@r := [ 1, 4 .. 32 ];|
  Error, Range: <last>-<first> (31) must be divisible by <inc> (3)
\end{Verbatim}
 

 Most often ranges are used in connection with the \texttt{for}\texttt{\symbol{45}}loop see{\nobreakspace}\ref{For}). Here the construct 

 \texttt{for \mbox{\texttt{\mdseries\slshape var}} in [ \mbox{\texttt{\mdseries\slshape first}} .. \mbox{\texttt{\mdseries\slshape last}} ] do \mbox{\texttt{\mdseries\slshape statements}} od} 

 replaces the 

 \texttt{for \mbox{\texttt{\mdseries\slshape var}} from \mbox{\texttt{\mdseries\slshape first}} to \mbox{\texttt{\mdseries\slshape last}} do \mbox{\texttt{\mdseries\slshape statements}} od} 

 which is more usual in other programming languages. 

 
\begin{Verbatim}[commandchars=!@|,fontsize=\small,frame=single,label=Example]
  !gapprompt@gap>| !gapinput@s := [];; for i in [10..20] do Add( s, i^2 ); od; s;|
  [ 100, 121, 144, 169, 196, 225, 256, 289, 324, 361, 400 ]
\end{Verbatim}
 

 Note that a range with \texttt{\mbox{\texttt{\mdseries\slshape last}} {\textgreater}= \mbox{\texttt{\mdseries\slshape first}}} is at the same time also a proper set (see{\nobreakspace}\ref{Sorted Lists and Sets}), because it contains no holes or duplicates and is sorted, and also a row
vector (see{\nobreakspace}\ref{Row Vectors}), because it contains no holes and all elements are integers. 

 

\subsection{\textcolor{Chapter }{IsRange}}
\logpage{[ 21, 22, 1 ]}\nobreak
\hyperdef{L}{X86DDC2FF7A50FBEE}{}
{\noindent\textcolor{FuncColor}{$\triangleright$\enspace\texttt{IsRange({\mdseries\slshape obj})\index{IsRange@\texttt{IsRange}}
\label{IsRange}
}\hfill{\scriptsize (Category)}}\\


 tests if the object \mbox{\texttt{\mdseries\slshape obj}} is a range, i.e. is a dense list of integers that is also a range
(see{\nobreakspace}\ref{Ranges} for a definition of ``range''). 
\begin{Verbatim}[commandchars=!@|,fontsize=\small,frame=single,label=Example]
  !gapprompt@gap>| !gapinput@IsRange( [1,2,3] );  IsRange( [7,5,3,1] );|
  true
  true
  !gapprompt@gap>| !gapinput@IsRange( [1 .. 3] );|
  true
  !gapprompt@gap>| !gapinput@IsRange( [1,2,4,5] );  IsRange( [1,,3,,5,,7] );|
  false
  false
  !gapprompt@gap>| !gapinput@IsRange( [] );  IsRange( [1] );|
  true
  true
\end{Verbatim}
 }

 

\subsection{\textcolor{Chapter }{IsRangeRep}}
\logpage{[ 21, 22, 2 ]}\nobreak
\hyperdef{L}{X83896BC481536B07}{}
{\noindent\textcolor{FuncColor}{$\triangleright$\enspace\texttt{IsRangeRep({\mdseries\slshape obj})\index{IsRangeRep@\texttt{IsRangeRep}}
\label{IsRangeRep}
}\hfill{\scriptsize (Representation)}}\\


 Tests whether \mbox{\texttt{\mdseries\slshape obj}} is represented as a range, that is by internally storing only the first value,
the in\texttt{\symbol{45}} or decrement, and the last value of the range. 

 To test whether a list is a range in the mathematical sense see \texttt{IsRange} (\ref{IsRange}). 

 Lists created by the syntactic construct \texttt{[ \mbox{\texttt{\mdseries\slshape first}}, \mbox{\texttt{\mdseries\slshape second}} .. \mbox{\texttt{\mdseries\slshape last}} ]}, see \ref{Ranges}, are in \texttt{IsRangeRep}. 

 Note that if you modify an \texttt{IsRangeRep} object by assigning to one of its entries, or by using \texttt{Add} (\ref{Add}) or \texttt{Append} (\ref{Append}), then the range may be converted into a plain list, even though the resulting
list may still be a range, mathematically. 

 
\begin{Verbatim}[commandchars=!@|,fontsize=\small,frame=single,label=Example]
  !gapprompt@gap>| !gapinput@IsRangeRep( [1 .. 3] );|
  true
  !gapprompt@gap>| !gapinput@IsRangeRep( [1, 2, 3] );|
  false
  !gapprompt@gap>| !gapinput@l := [1..3];;|
  !gapprompt@gap>| !gapinput@l[1] := 1;;|
  !gapprompt@gap>| !gapinput@l;|
  [ 1, 2, 3 ]
\end{Verbatim}
 }

 

\subsection{\textcolor{Chapter }{ConvertToRangeRep}}
\logpage{[ 21, 22, 3 ]}\nobreak
\hyperdef{L}{X7D22B2298167A58F}{}
{\noindent\textcolor{FuncColor}{$\triangleright$\enspace\texttt{ConvertToRangeRep({\mdseries\slshape list})\index{ConvertToRangeRep@\texttt{ConvertToRangeRep}}
\label{ConvertToRangeRep}
}\hfill{\scriptsize (function)}}\\


 For some lists the \textsf{GAP} kernel knows that they are in fact ranges. Those lists are represented
internally in a compact way, namely in \texttt{IsRangeRep} (\ref{IsRangeRep}), instead of as plain lists. A list that is represented as a plain list might
still be a range but \textsf{GAP} may not know this. 

 If \mbox{\texttt{\mdseries\slshape list}} is a range then \texttt{ConvertToRangeRep} changes the representation of \mbox{\texttt{\mdseries\slshape list}} to \texttt{IsRangeRep} (\ref{IsRangeRep}). A call of \texttt{ConvertToRangeRep} for a list that is not a range is ignored. 

 
\begin{Verbatim}[commandchars=!@|,fontsize=\small,frame=single,label=Example]
  !gapprompt@gap>| !gapinput@r:= [ 1, 2, 3, 4, 5, 6, 7, 8, 9, 10 ];|
  [ 1, 2, 3, 4, 5, 6, 7, 8, 9, 10 ]
  !gapprompt@gap>| !gapinput@ConvertToRangeRep( r );  r;|
  [ 1 .. 10 ]
  !gapprompt@gap>| !gapinput@l:= [ 1, 2, 4, 5 ];;  ConvertToRangeRep( l );  l;|
  [ 1, 2, 4, 5 ]
\end{Verbatim}
 }

 }

  
\section{\textcolor{Chapter }{Enumerators}}\label{Enumerators}
\logpage{[ 21, 23, 0 ]}
\hyperdef{L}{X7EA3ACE27E43D174}{}
{
  An \emph{enumerator} is an immutable list that need not store its elements explicitly but knows,
from a set of basic data, how to determine the $i$\texttt{\symbol{45}}th element and the position of a given object. A typical
example of this is a vector space over a finite field with $q$ elements for which it is very easy to enumerate all elements using $q$\texttt{\symbol{45}}adic expansions of integers. 

 Using this enumeration can be even quicker than a binary search in a sorted
list of vectors, see \texttt{IsQuickPositionList} (\ref{IsQuickPositionList}). 

 On the one hand, element access to an enumerator may take more time than
element access to an internally represented list containing the same elements.
On the other hand, an enumerator may save a vast amount of memory. Take for
example a permutation group of size a few millions. Even for moderate degree
it is unlikely that a list of all its elements will fit into memory whereas it
is no problem to construct an enumerator from a stabilizer chain
(see{\nobreakspace}\ref{Stabilizer Chains}). 

 There are situations where one only wants to loop over the elements of a
domain, without using the special facilities of an enumerator, namely the
particular order of elements and the possibility to find the position of
elements. For such cases, \textsf{GAP} provides iterators (see{\nobreakspace}\ref{Iterators}). 

 The functions \texttt{Enumerator} (\ref{Enumerator}) and \texttt{EnumeratorSorted} (\ref{EnumeratorSorted}) return enumerators of domains. Most of the special implementations of
enumerators in the \textsf{GAP} library are based on the general interface that is provided by \texttt{EnumeratorByFunctions} (\ref{EnumeratorByFunctions:for a domain and a record}); one generic example is \texttt{EnumeratorByBasis} (\ref{EnumeratorByBasis}), which can be used to get an enumerator of a finite dimensional free module. 

 Also enumerators for non\texttt{\symbol{45}}domains can be implemented via \texttt{EnumeratorByFunctions} (\ref{EnumeratorByFunctions:for a domain and a record}); for a discussion, see{\nobreakspace}\ref{Example -- Constructing Enumerators}. 

\subsection{\textcolor{Chapter }{IsQuickPositionList}}
\logpage{[ 21, 23, 1 ]}\nobreak
\hyperdef{L}{X7BB462C17962647F}{}
{\noindent\textcolor{FuncColor}{$\triangleright$\enspace\texttt{IsQuickPositionList({\mdseries\slshape list})\index{IsQuickPositionList@\texttt{IsQuickPositionList}}
\label{IsQuickPositionList}
}\hfill{\scriptsize (filter)}}\\


 This filter indicates that a position test in \mbox{\texttt{\mdseries\slshape list}} is quicker than about 5 or 6 element comparisons for ``smaller''. If this is the case it can be beneficial to use \texttt{Position} (\ref{Position}) in \mbox{\texttt{\mdseries\slshape list}} and a bit list than ordered lists to represent subsets of \mbox{\texttt{\mdseries\slshape list}}. }

 }

  
\section{\textcolor{Chapter }{Plain Lists}}\label{Plain Lists}
\logpage{[ 21, 24, 0 ]}
\hyperdef{L}{X81ECC2077D88E112}{}
{
  Plain lists are the default kind of lists in \textsf{GAP}, in the sense that \textsf{GAP} stores the list entries and does not know how to do better (as opposed to
ranges or strings, which are also lists). Often it is not necessary to know
how a given list is represented internally, the operations defined for lists
apply to all lists. 

 Typical situations where the representation matters are when one wants to make
sure that the given list is \emph{not} a plain list and thus will be handled more efficiently, for example when one
installs a method for a particular operation, where an argument is required to
be a list in a particular representation. 

\subsection{\textcolor{Chapter }{PlainListCopy}}
\logpage{[ 21, 24, 1 ]}\nobreak
\hyperdef{L}{X8438CB908367254C}{}
{\noindent\textcolor{FuncColor}{$\triangleright$\enspace\texttt{PlainListCopy({\mdseries\slshape list})\index{PlainListCopy@\texttt{PlainListCopy}}
\label{PlainListCopy}
}\hfill{\scriptsize (function)}}\\


 This function returns a list equal to its argument, in a plain list
representation (see \texttt{IsPlistRep} (\ref{IsPlistRep})). This is intended for use in certain rare situations, such as before
objectifying, or calling some kernel functions. }

 

\subsection{\textcolor{Chapter }{IsPlistRep}}
\logpage{[ 21, 24, 2 ]}\nobreak
\hyperdef{L}{X87BA4EBF80F16B72}{}
{\noindent\textcolor{FuncColor}{$\triangleright$\enspace\texttt{IsPlistRep({\mdseries\slshape obj})\index{IsPlistRep@\texttt{IsPlistRep}}
\label{IsPlistRep}
}\hfill{\scriptsize (Representation)}}\\


 \textsf{GAP} lists created by entering comma separated values in square brackets are
usually represented internally as so\texttt{\symbol{45}}called \emph{plain lists}. Other representations of lists are \texttt{IsBlistRep} (\ref{IsBlistRep}), \texttt{IsRangeRep} (\ref{IsRangeRep}), \texttt{IsStringRep} (\ref{IsStringRep}), or the ones that are chosen for implementing enumerators, see Section \ref{Enumerators}. 

 
\begin{Verbatim}[commandchars=!@|,fontsize=\small,frame=single,label=Example]
  !gapprompt@gap>| !gapinput@IsPlistRep( [ 1, 2, 3 ] );|
  true
  !gapprompt@gap>| !gapinput@IsPlistRep( "abc" );|
  false
  !gapprompt@gap>| !gapinput@IsPlistRep( [ 1 .. 5 ] );|
  false
  !gapprompt@gap>| !gapinput@IsPlistRep( BlistList( [ 1 .. 5 ], [ 1 ] ) );|
  false
  !gapprompt@gap>| !gapinput@IsPlistRep( 0 );|
  false
\end{Verbatim}
 }

 }

 }

            
\chapter{\textcolor{Chapter }{Boolean Lists}}\label{Boolean Lists}
\logpage{[ 22, 0, 0 ]}
\hyperdef{L}{X7AC531DD79B6938E}{}
{
  This chapter describes boolean lists. A \emph{boolean list} is a list that has no holes and contains only the boolean values \texttt{true} and \texttt{false} (see Chapter{\nobreakspace}\ref{Booleans}). In function names we call boolean lists \emph{blists} for brevity.  
\section{\textcolor{Chapter }{IsBlist (Filter)}}\label{sect:IsBlist}
\logpage{[ 22, 1, 0 ]}
\hyperdef{L}{X7E7832B0804221AE}{}
{
  

\subsection{\textcolor{Chapter }{IsBlist}}
\logpage{[ 22, 1, 1 ]}\nobreak
\hyperdef{L}{X7BE078187A08DCEA}{}
{\noindent\textcolor{FuncColor}{$\triangleright$\enspace\texttt{IsBlist({\mdseries\slshape obj})\index{IsBlist@\texttt{IsBlist}}
\label{IsBlist}
}\hfill{\scriptsize (Category)}}\\


 A boolean list (``blist'') is a list that has no holes and contains only \texttt{true} and \texttt{false}. Boolean lists can be represented in an efficient compact form, see \ref{More about Boolean Lists} for details. 

 
\begin{Verbatim}[commandchars=!@|,fontsize=\small,frame=single,label=Example]
  !gapprompt@gap>| !gapinput@IsBlist( [ true, true, false, false ] );|
  true
  !gapprompt@gap>| !gapinput@IsBlist( [] );|
  true
  !gapprompt@gap>| !gapinput@IsBlist( [false,,true] );  # has holes|
  false
  !gapprompt@gap>| !gapinput@IsBlist( [1,1,0,0] );      # contains not only boolean values|
  false
  !gapprompt@gap>| !gapinput@IsBlist( 17 );             # is not even a list|
  false
\end{Verbatim}
 

 Boolean lists are lists and all operations for lists are therefore applicable
to boolean lists. 

 Boolean lists can be used in various ways, but maybe the most important
application is their use for the description of \emph{subsets} of finite sets. Suppose $set$ is a finite set, represented as a list. Then a subset $sub$ of $set$ is represented by a boolean list $blist$ of the same length as $set$ such that $blist[i]$ is \texttt{true} if $set[i]$ is in $sub$, and \texttt{false} otherwise. }

 }

  
\section{\textcolor{Chapter }{Boolean Lists Representing Subsets}}\label{Boolean Lists Representing Subsets}
\logpage{[ 22, 2, 0 ]}
\hyperdef{L}{X7CC745317FE54C14}{}
{
  

\subsection{\textcolor{Chapter }{BlistList}}
\logpage{[ 22, 2, 1 ]}\nobreak
\hyperdef{L}{X7C597B2D87CA2E6E}{}
{\noindent\textcolor{FuncColor}{$\triangleright$\enspace\texttt{BlistList({\mdseries\slshape list, sub})\index{BlistList@\texttt{BlistList}}
\label{BlistList}
}\hfill{\scriptsize (function)}}\\


 returns a new boolean list that describes the list \mbox{\texttt{\mdseries\slshape sub}} as a sublist of the dense list \mbox{\texttt{\mdseries\slshape list}}. That is \texttt{BlistList} returns a boolean list $blist$ of the same length as \mbox{\texttt{\mdseries\slshape list}} such that $blist[i]$ is \texttt{true} if \mbox{\texttt{\mdseries\slshape list}}$[i]$ is in \mbox{\texttt{\mdseries\slshape sub}} and \texttt{false} otherwise. 

 \mbox{\texttt{\mdseries\slshape list}} need not be a proper set (see{\nobreakspace}\ref{Sorted Lists and Sets}), even though in this case \texttt{BlistList} is most efficient. In particular \mbox{\texttt{\mdseries\slshape list}} may contain duplicates. \mbox{\texttt{\mdseries\slshape sub}} need not be a proper sublist of \mbox{\texttt{\mdseries\slshape list}}, i.e., \mbox{\texttt{\mdseries\slshape sub}} may contain elements that are not in \mbox{\texttt{\mdseries\slshape list}}. Those elements of course have no influence on the result of \texttt{BlistList}. 

 
\begin{Verbatim}[commandchars=!@|,fontsize=\small,frame=single,label=Example]
  !gapprompt@gap>| !gapinput@BlistList( [1..10], [2,3,5,7] );|
  [ false, true, true, false, true, false, true, false, false, false ]
  !gapprompt@gap>| !gapinput@BlistList( [1,2,3,4,5,2,8,6,4,10], [4,8,9,16] );|
  [ false, false, false, true, false, false, true, false, true, false ]
\end{Verbatim}
 

 See also{\nobreakspace}\texttt{UniteBlistList} (\ref{UniteBlistList}). }

 

\subsection{\textcolor{Chapter }{ListBlist}}
\logpage{[ 22, 2, 2 ]}\nobreak
\hyperdef{L}{X874BEF63785AB439}{}
{\noindent\textcolor{FuncColor}{$\triangleright$\enspace\texttt{ListBlist({\mdseries\slshape list, blist})\index{ListBlist@\texttt{ListBlist}}
\label{ListBlist}
}\hfill{\scriptsize (function)}}\\


 returns the sublist $sub$ of the list \mbox{\texttt{\mdseries\slshape list}}, which must have no holes, represented by the boolean list \mbox{\texttt{\mdseries\slshape blist}}, which must have the same length as \mbox{\texttt{\mdseries\slshape list}}. 

 $sub$ contains the element \mbox{\texttt{\mdseries\slshape list}}$[i]$ if \mbox{\texttt{\mdseries\slshape blist}}$[i]$ is \texttt{true} and does not contain the element if \mbox{\texttt{\mdseries\slshape blist}}$[i]$ is \texttt{false}. The order of the elements in $sub$ is the same as the order of the corresponding elements in \mbox{\texttt{\mdseries\slshape list}}. 

 
\begin{Verbatim}[commandchars=!@|,fontsize=\small,frame=single,label=Example]
  !gapprompt@gap>| !gapinput@ListBlist([1..8],[false,true,true,true,true,false,true,true]);|
  [ 2, 3, 4, 5, 7, 8 ]
  !gapprompt@gap>| !gapinput@ListBlist( [1,2,3,4,5,2,8,6,4,10],|
  !gapprompt@>| !gapinput@[false,false,false,true,false,false,true,false,true,false] );|
  [ 4, 8, 4 ]
\end{Verbatim}
 }

 

\subsection{\textcolor{Chapter }{SizeBlist}}
\logpage{[ 22, 2, 3 ]}\nobreak
\hyperdef{L}{X85AD5EF77EFD7451}{}
{\noindent\textcolor{FuncColor}{$\triangleright$\enspace\texttt{SizeBlist({\mdseries\slshape blist})\index{SizeBlist@\texttt{SizeBlist}}
\label{SizeBlist}
}\hfill{\scriptsize (function)}}\\


 returns the number of entries of the boolean list \mbox{\texttt{\mdseries\slshape blist}} that are \texttt{true}. This is the size of the subset represented by the boolean list \mbox{\texttt{\mdseries\slshape blist}}. 

 
\begin{Verbatim}[commandchars=!@|,fontsize=\small,frame=single,label=Example]
  !gapprompt@gap>| !gapinput@SizeBlist( [ false, true, false, true, false ] );|
  2
\end{Verbatim}
 }

 

\subsection{\textcolor{Chapter }{IsSubsetBlist}}
\logpage{[ 22, 2, 4 ]}\nobreak
\hyperdef{L}{X7BA42D03796ED4B3}{}
{\noindent\textcolor{FuncColor}{$\triangleright$\enspace\texttt{IsSubsetBlist({\mdseries\slshape blist1, blist2})\index{IsSubsetBlist@\texttt{IsSubsetBlist}}
\label{IsSubsetBlist}
}\hfill{\scriptsize (function)}}\\


 returns \texttt{true} if the boolean list \mbox{\texttt{\mdseries\slshape blist2}} is a subset of the boolean list \mbox{\texttt{\mdseries\slshape blist1}}, which must have equal length, and \texttt{false} otherwise. \mbox{\texttt{\mdseries\slshape blist2}} is a subset of \mbox{\texttt{\mdseries\slshape blist1}} if \mbox{\texttt{\mdseries\slshape blist1}}$[i] =$ \mbox{\texttt{\mdseries\slshape blist1}}$[i]$ \texttt{or} \mbox{\texttt{\mdseries\slshape blist2}}$[i]$ for all $i$. 

 
\begin{Verbatim}[commandchars=!@|,fontsize=\small,frame=single,label=Example]
  !gapprompt@gap>| !gapinput@blist1 := [ true, true, false, false ];;|
  !gapprompt@gap>| !gapinput@blist2 := [ true, false, true, false ];;|
  !gapprompt@gap>| !gapinput@IsSubsetBlist( blist1, blist2 );|
  false
  !gapprompt@gap>| !gapinput@blist2 := [ true, false, false, false ];;|
  !gapprompt@gap>| !gapinput@IsSubsetBlist( blist1, blist2 );|
  true
\end{Verbatim}
 }

 }

  
\section{\textcolor{Chapter }{Set Operations via Boolean Lists}}\label{Set Operations via Boolean Lists}
\logpage{[ 22, 3, 0 ]}
\hyperdef{L}{X8100080382AECFF9}{}
{
  
\subsection{\textcolor{Chapter }{UnionBlist}}\logpage{[ 22, 3, 1 ]}
\hyperdef{L}{X7970BD3883C42D91}{}
{
\noindent\textcolor{FuncColor}{$\triangleright$\enspace\texttt{UnionBlist({\mdseries\slshape blist1, blist2[, ...]})\index{UnionBlist@\texttt{UnionBlist}!for various boolean lists}
\label{UnionBlist:for various boolean lists}
}\hfill{\scriptsize (function)}}\\
\noindent\textcolor{FuncColor}{$\triangleright$\enspace\texttt{UnionBlist({\mdseries\slshape list})\index{UnionBlist@\texttt{UnionBlist}!for a list}
\label{UnionBlist:for a list}
}\hfill{\scriptsize (function)}}\\


 In the first form \texttt{UnionBlist} returns the union of the boolean lists \mbox{\texttt{\mdseries\slshape blist1}}, \mbox{\texttt{\mdseries\slshape blist2}}, etc., which must have equal length. The \emph{union} is a new boolean list that contains at position $i$ the value \mbox{\texttt{\mdseries\slshape blist1}}$[i]$ \texttt{or} \mbox{\texttt{\mdseries\slshape blist2}}$[i]$ \texttt{or} $\ldots$. 

 The second form takes the union of all blists (which as for the first form
must have equal length) in the list \mbox{\texttt{\mdseries\slshape list}}. }

 
\subsection{\textcolor{Chapter }{IntersectionBlist}}\logpage{[ 22, 3, 2 ]}
\hyperdef{L}{X86E1F8DE85E1EE1E}{}
{
\noindent\textcolor{FuncColor}{$\triangleright$\enspace\texttt{IntersectionBlist({\mdseries\slshape blist1, blist2[, ...]})\index{IntersectionBlist@\texttt{IntersectionBlist}!for various boolean lists}
\label{IntersectionBlist:for various boolean lists}
}\hfill{\scriptsize (function)}}\\
\noindent\textcolor{FuncColor}{$\triangleright$\enspace\texttt{IntersectionBlist({\mdseries\slshape list})\index{IntersectionBlist@\texttt{IntersectionBlist}!for a list}
\label{IntersectionBlist:for a list}
}\hfill{\scriptsize (function)}}\\


 In the first form \texttt{IntersectionBlist} returns the intersection of the boolean lists \mbox{\texttt{\mdseries\slshape blist1}}, \mbox{\texttt{\mdseries\slshape blist2}}, etc., which must have equal length. The \emph{intersection} is a new blist that contains at position $i$ the value \mbox{\texttt{\mdseries\slshape blist1}}$[i]$ \texttt{and} \mbox{\texttt{\mdseries\slshape blist2}}$[i]$ \texttt{and} $\ldots$. 

 In the second form \mbox{\texttt{\mdseries\slshape list}} must be a list of boolean lists \mbox{\texttt{\mdseries\slshape blist1}}, \mbox{\texttt{\mdseries\slshape blist2}}, etc., which must have equal length, and \texttt{IntersectionBlist} returns the intersection of those boolean lists. }

 

\subsection{\textcolor{Chapter }{DifferenceBlist}}
\logpage{[ 22, 3, 3 ]}\nobreak
\hyperdef{L}{X7D6FC2C58725708C}{}
{\noindent\textcolor{FuncColor}{$\triangleright$\enspace\texttt{DifferenceBlist({\mdseries\slshape blist1, blist2})\index{DifferenceBlist@\texttt{DifferenceBlist}}
\label{DifferenceBlist}
}\hfill{\scriptsize (function)}}\\


 returns the asymmetric set difference of the two boolean lists \mbox{\texttt{\mdseries\slshape blist1}} and \mbox{\texttt{\mdseries\slshape blist2}}, which must have equal length. The \emph{asymmetric set difference} is a new boolean list that contains at position $i$ the value \mbox{\texttt{\mdseries\slshape blist1}}$[i]$ \texttt{and} \texttt{not} \mbox{\texttt{\mdseries\slshape blist2}}$[i]$. 

 
\begin{Verbatim}[commandchars=!@|,fontsize=\small,frame=single,label=Example]
  !gapprompt@gap>| !gapinput@blist1 := [ true, true, false, false ];;|
  !gapprompt@gap>| !gapinput@blist2 := [ true, false, true, false ];;|
  !gapprompt@gap>| !gapinput@UnionBlist( blist1, blist2 );|
  [ true, true, true, false ]
  !gapprompt@gap>| !gapinput@IntersectionBlist( blist1, blist2 );|
  [ true, false, false, false ]
  !gapprompt@gap>| !gapinput@DifferenceBlist( blist1, blist2 );|
  [ false, true, false, false ]
\end{Verbatim}
 }

 }

  
\section{\textcolor{Chapter }{Function that Modify Boolean Lists}}\label{Function that Modify Boolean Lists}
\logpage{[ 22, 4, 0 ]}
\hyperdef{L}{X8634D25D7B4C6151}{}
{
  

\subsection{\textcolor{Chapter }{UniteBlist}}
\logpage{[ 22, 4, 1 ]}\nobreak
\hyperdef{L}{X79815EB77CC8A389}{}
{\noindent\textcolor{FuncColor}{$\triangleright$\enspace\texttt{UniteBlist({\mdseries\slshape blist1, blist2})\index{UniteBlist@\texttt{UniteBlist}}
\label{UniteBlist}
}\hfill{\scriptsize (function)}}\\


 \texttt{UniteBlist} unites the boolean list \mbox{\texttt{\mdseries\slshape blist1}} with the boolean list \mbox{\texttt{\mdseries\slshape blist2}}, which must have the same length. This is equivalent to assigning \mbox{\texttt{\mdseries\slshape blist1}}$[i] :=$ \mbox{\texttt{\mdseries\slshape blist1}}$[i]$ \texttt{or} \mbox{\texttt{\mdseries\slshape blist2}}$[i]$ for all $i$. 

 \texttt{UniteBlist} returns nothing, it is only called to change \mbox{\texttt{\mdseries\slshape blist1}}. 

 
\begin{Verbatim}[commandchars=!@|,fontsize=\small,frame=single,label=Example]
  !gapprompt@gap>| !gapinput@blist1 := [ true, true, false, false ];;|
  !gapprompt@gap>| !gapinput@blist2 := [ true, false, true, false ];;|
  !gapprompt@gap>| !gapinput@UniteBlist( blist1, blist2 );|
  !gapprompt@gap>| !gapinput@blist1;|
  [ true, true, true, false ]
\end{Verbatim}
 

 The function \texttt{UnionBlist} (\ref{UnionBlist:for a list}) is the nondestructive counterpart to \texttt{UniteBlist}. }

 

\subsection{\textcolor{Chapter }{UniteBlistList}}
\logpage{[ 22, 4, 2 ]}\nobreak
\hyperdef{L}{X7C86C8D3853BE5EB}{}
{\noindent\textcolor{FuncColor}{$\triangleright$\enspace\texttt{UniteBlistList({\mdseries\slshape list, blist, sub})\index{UniteBlistList@\texttt{UniteBlistList}}
\label{UniteBlistList}
}\hfill{\scriptsize (function)}}\\


 works like \texttt{UniteBlist(\mbox{\texttt{\mdseries\slshape blist}},BlistList(\mbox{\texttt{\mdseries\slshape list}},\mbox{\texttt{\mdseries\slshape sub}}))}. As no intermediate blist is created, the performance is better than the
separate function calls. }

 

\subsection{\textcolor{Chapter }{IntersectBlist}}
\logpage{[ 22, 4, 3 ]}\nobreak
\hyperdef{L}{X84EB70D37EB275DF}{}
{\noindent\textcolor{FuncColor}{$\triangleright$\enspace\texttt{IntersectBlist({\mdseries\slshape blist1, blist2})\index{IntersectBlist@\texttt{IntersectBlist}}
\label{IntersectBlist}
}\hfill{\scriptsize (function)}}\\


 intersects the boolean list \mbox{\texttt{\mdseries\slshape blist1}} with the boolean list \mbox{\texttt{\mdseries\slshape blist2}}, which must have the same length. This is equivalent to assigning \mbox{\texttt{\mdseries\slshape blist1}}$[i]:=$ \mbox{\texttt{\mdseries\slshape blist1}}$[i]$ \texttt{and} \mbox{\texttt{\mdseries\slshape blist2}}$[i]$ for all $i$. 

 \texttt{IntersectBlist} returns nothing, it is only called to change \mbox{\texttt{\mdseries\slshape blist1}}. 

 
\begin{Verbatim}[commandchars=!@|,fontsize=\small,frame=single,label=Example]
  !gapprompt@gap>| !gapinput@blist1 := [ true, true, false, false ];;|
  !gapprompt@gap>| !gapinput@blist2 := [ true, false, true, false ];;|
  !gapprompt@gap>| !gapinput@IntersectBlist( blist1, blist2 );|
  !gapprompt@gap>| !gapinput@blist1;|
  [ true, false, false, false ]
\end{Verbatim}
 

 The function \texttt{IntersectionBlist} (\ref{IntersectionBlist:for a list}) is the nondestructive counterpart to \texttt{IntersectBlist}. }

 

\subsection{\textcolor{Chapter }{SubtractBlist}}
\logpage{[ 22, 4, 4 ]}\nobreak
\hyperdef{L}{X7AA138407D5A3BAC}{}
{\noindent\textcolor{FuncColor}{$\triangleright$\enspace\texttt{SubtractBlist({\mdseries\slshape blist1, blist2})\index{SubtractBlist@\texttt{SubtractBlist}}
\label{SubtractBlist}
}\hfill{\scriptsize (function)}}\\


 subtracts the boolean list \mbox{\texttt{\mdseries\slshape blist2}} from the boolean list \mbox{\texttt{\mdseries\slshape blist1}}, which must have equal length. This is equivalent to assigning \mbox{\texttt{\mdseries\slshape blist1}}$[i]:=$ \mbox{\texttt{\mdseries\slshape blist1}}$[i]$ \texttt{and} \texttt{not} \mbox{\texttt{\mdseries\slshape blist2}}$[i]$ for all $i$. 

 \texttt{SubtractBlist} returns nothing, it is only called to change \mbox{\texttt{\mdseries\slshape blist1}}. 

 
\begin{Verbatim}[commandchars=!@|,fontsize=\small,frame=single,label=Example]
  !gapprompt@gap>| !gapinput@blist1 := [ true, true, false, false ];;|
  !gapprompt@gap>| !gapinput@blist2 := [ true, false, true, false ];;|
  !gapprompt@gap>| !gapinput@SubtractBlist( blist1, blist2 );|
  !gapprompt@gap>| !gapinput@blist1;|
  [ false, true, false, false ]
\end{Verbatim}
 

 The function \texttt{DifferenceBlist} (\ref{DifferenceBlist}) is the nondestructive counterpart to \texttt{SubtractBlist}. }

 

\subsection{\textcolor{Chapter }{MeetBlist}}
\logpage{[ 22, 4, 5 ]}\nobreak
\hyperdef{L}{X830645EC846B2E3C}{}
{\noindent\textcolor{FuncColor}{$\triangleright$\enspace\texttt{MeetBlist({\mdseries\slshape blist1, blist2})\index{MeetBlist@\texttt{MeetBlist}}
\label{MeetBlist}
}\hfill{\scriptsize (function)}}\\


 Returns \texttt{true} if there is a position at which both \mbox{\texttt{\mdseries\slshape blist1}} and \mbox{\texttt{\mdseries\slshape blist2}} contain \texttt{true}, and \texttt{false} otherwise. It is equivalent to, but faster than \texttt{SizeBlist(IntersectionBlist(blist1, blist2)) {\textless}{\textgreater} 0}. An error is thrown if the lists do not have the same length. 
\begin{Verbatim}[commandchars=!@|,fontsize=\small,frame=single,label=Example]
  !gapprompt@gap>| !gapinput@blist1 := [ true, true, true, true ];;|
  !gapprompt@gap>| !gapinput@blist2 := [ true, false, true, false ];;|
  !gapprompt@gap>| !gapinput@MeetBlist( blist1, blist2 );|
  true
  !gapprompt@gap>| !gapinput@FlipBlist( blist1 );|
  !gapprompt@gap>| !gapinput@MeetBlist( blist1, blist2 );|
  false
\end{Verbatim}
 }

 

\subsection{\textcolor{Chapter }{FlipBlist}}
\logpage{[ 22, 4, 6 ]}\nobreak
\hyperdef{L}{X7F14FF35786DAEF3}{}
{\noindent\textcolor{FuncColor}{$\triangleright$\enspace\texttt{FlipBlist({\mdseries\slshape blist})\index{FlipBlist@\texttt{FlipBlist}}
\label{FlipBlist}
}\hfill{\scriptsize (function)}}\\


 Changes every entry in \mbox{\texttt{\mdseries\slshape blist}} that equals \texttt{true} to \texttt{false} and vice versa. If \texttt{blist1} and \texttt{blist2} are boolean lists with equal length and every value in \texttt{blist2} is \texttt{true}, then \texttt{FlipBlist( blist1 )} is equivalent to \texttt{SubtractBlist( blist2, blist1 ); blist1 := blist2;} but \texttt{FlipBlist} is faster, and simpler to type. 

 \texttt{FlipBlist} returns nothing, it is only called to change \mbox{\texttt{\mdseries\slshape blist}} in\texttt{\symbol{45}}place. 

 
\begin{Verbatim}[commandchars=!@|,fontsize=\small,frame=single,label=Example]
  !gapprompt@gap>| !gapinput@blist1 := [ true, true, true, true ];;|
  !gapprompt@gap>| !gapinput@blist2 := [ true, false, true, false ];;|
  !gapprompt@gap>| !gapinput@SubtractBlist( blist1, blist2 );|
  !gapprompt@gap>| !gapinput@blist1;|
  [ false, true, false, true ]
  !gapprompt@gap>| !gapinput@FlipBlist( blist2 );|
  !gapprompt@gap>| !gapinput@blist2;|
  [ false, true, false, true ]
\end{Verbatim}
 }

 

\subsection{\textcolor{Chapter }{SetAllBlist}}
\logpage{[ 22, 4, 7 ]}\nobreak
\hyperdef{L}{X7E9F6C197A79098F}{}
{\noindent\textcolor{FuncColor}{$\triangleright$\enspace\texttt{SetAllBlist({\mdseries\slshape blist})\index{SetAllBlist@\texttt{SetAllBlist}}
\label{SetAllBlist}
}\hfill{\scriptsize (function)}}\\


 Changes every entry in \mbox{\texttt{\mdseries\slshape blist}} to \texttt{true}. If \texttt{blist1} and \texttt{blist2} are boolean lists with equal length and every value in \texttt{blist2} is \texttt{true}, then \texttt{SetAllBlist( blist1 )} is equivalent to \texttt{UniteBlist( blist1, blist2 );} but is faster, and simpler to type. 

 \texttt{SetAllBlist} returns nothing, it is only called to change \mbox{\texttt{\mdseries\slshape blist}} in\texttt{\symbol{45}}place. 

 
\begin{Verbatim}[commandchars=!@|,fontsize=\small,frame=single,label=Example]
  !gapprompt@gap>| !gapinput@blist1 := [ true, true, true, true ];;|
  !gapprompt@gap>| !gapinput@blist2 := [ true, false, true, false ];;|
  !gapprompt@gap>| !gapinput@SetAllBlist( blist1 );|
  !gapprompt@gap>| !gapinput@blist1;|
  [ true, true, true, true ]
  !gapprompt@gap>| !gapinput@SetAllBlist(blist2);|
  !gapprompt@gap>| !gapinput@blist2;|
  [ true, true, true, true ]
\end{Verbatim}
 }

 

\subsection{\textcolor{Chapter }{ClearAllBlist}}
\logpage{[ 22, 4, 8 ]}\nobreak
\hyperdef{L}{X87ED45A88688AE8E}{}
{\noindent\textcolor{FuncColor}{$\triangleright$\enspace\texttt{ClearAllBlist({\mdseries\slshape blist})\index{ClearAllBlist@\texttt{ClearAllBlist}}
\label{ClearAllBlist}
}\hfill{\scriptsize (function)}}\\


 Changes every entry in \mbox{\texttt{\mdseries\slshape blist}} to \texttt{false}. If \texttt{blist1} and \texttt{blist2} are boolean lists with equal length and every value in \texttt{blist2} is \texttt{false}, then \texttt{ClearAllBlist( blist1 )} is equivalent to \texttt{IntersectBlist( blist1, blist2 );} but is faster, and simpler to type. 

 \texttt{ClearAllBlist} returns nothing, it is only called to change \mbox{\texttt{\mdseries\slshape blist}} in\texttt{\symbol{45}}place. 

 
\begin{Verbatim}[commandchars=!@|,fontsize=\small,frame=single,label=Example]
  !gapprompt@gap>| !gapinput@blist1 := [ true, true, true, true ];;|
  !gapprompt@gap>| !gapinput@blist2 := [ true, false, true, false ];;|
  !gapprompt@gap>| !gapinput@ClearAllBlist( blist1 );|
  !gapprompt@gap>| !gapinput@blist1;|
  [ false, false, false, false ]
  !gapprompt@gap>| !gapinput@ClearAllBlist(blist2);|
  !gapprompt@gap>| !gapinput@blist2;|
  [ false, false, false, false ]
\end{Verbatim}
 }

 }

  
\section{\textcolor{Chapter }{More about Boolean Lists}}\label{More about Boolean Lists}
\logpage{[ 22, 5, 0 ]}
\hyperdef{L}{X7C71B225841DFC0F}{}
{
  We defined a boolean list as a list that has no holes and contains only \texttt{true} and \texttt{false}. There is a special internal representation for boolean lists that needs only
1 bit for each entry. This bit is set if the entry is \texttt{true} and reset if the entry is \texttt{false}. This representation is of course much more compact than the ordinary
representation of lists, which needs 32 or 64 bits per entry. 

       

 Not every boolean list is represented in this compact representation. It would
be too much work to test every time a list is changed, whether this list has
become a boolean list. This section tells you under which circumstances a
boolean list is represented in the compact representation, so you can write
your functions in such a way that you make best use of the compact
representation. 

 If a dense list containing only \texttt{true} and \texttt{false} is read, it is stored in the compact representation. Furthermore, the results
of \texttt{BlistList} (\ref{BlistList}), \texttt{UnionBlist} (\ref{UnionBlist:for a list}), \texttt{IntersectionBlist} (\ref{IntersectionBlist:for a list}) and \texttt{DifferenceBlist} (\ref{DifferenceBlist}) are known to be boolean lists by construction, and thus are represented in the
compact representation upon creation. 

 If an argument of \texttt{IsSubsetBlist} (\ref{IsSubsetBlist}), \texttt{ListBlist} (\ref{ListBlist}), \texttt{UnionBlist} (\ref{UnionBlist:for a list}), \texttt{IntersectionBlist} (\ref{IntersectionBlist:for a list}), \texttt{DifferenceBlist} (\ref{DifferenceBlist}), \texttt{UniteBlist} (\ref{UniteBlist}), \texttt{IntersectBlist} (\ref{IntersectBlist}) and \texttt{SubtractBlist} (\ref{SubtractBlist}) is a list represented in the ordinary representation, it is tested to see if
it is in fact a boolean list. If it is not, an error is signalled. If it is,
the representation of the list is changed to the compact representation. 

 If you change a boolean list that is represented in the compact representation
by assignment (see \ref{List Assignment}) or \texttt{Add} (\ref{Add}) in such a way that the list remains a boolean list it will remain represented
in the compact representation. Note that changing a list that is not
represented in the compact representation, whether it is a boolean list or
not, in such a way that the resulting list becomes a boolean list, will never
change the representation of the list. 

\subsection{\textcolor{Chapter }{IsBlistRep}}
\logpage{[ 22, 5, 1 ]}\nobreak
\hyperdef{L}{X8453ADDA810B4C03}{}
{\noindent\textcolor{FuncColor}{$\triangleright$\enspace\texttt{IsBlistRep({\mdseries\slshape obj})\index{IsBlistRep@\texttt{IsBlistRep}}
\label{IsBlistRep}
}\hfill{\scriptsize (Representation)}}\\
\noindent\textcolor{FuncColor}{$\triangleright$\enspace\texttt{ConvertToBlistRep({\mdseries\slshape blist})\index{ConvertToBlistRep@\texttt{ConvertToBlistRep}}
\label{ConvertToBlistRep}
}\hfill{\scriptsize (function)}}\\
\textbf{\indent Returns:\ }
\texttt{true} or \texttt{false}



 The first function is a filter that returns \texttt{true} if the object \mbox{\texttt{\mdseries\slshape obj}} is a boolean list in compact representation and \texttt{false} otherwise, see \ref{More about Boolean Lists}.

 The second function converts the object \mbox{\texttt{\mdseries\slshape blist}} to a boolean list in compact representation and returns \texttt{true} if this is possible. Otherwise \mbox{\texttt{\mdseries\slshape blist}} is unchanged and \texttt{false} is returned. 
\begin{Verbatim}[commandchars=!@|,fontsize=\small,frame=single,label=Example]
  !gapprompt@gap>| !gapinput@l := [true, false, true];|
  [ true, false, true ]
  !gapprompt@gap>| !gapinput@IsBlistRep(l);|
  true
  !gapprompt@gap>| !gapinput@l := [true, false, 1];|
  [ true, false, 1 ]
  !gapprompt@gap>| !gapinput@l[3] := false;|
  false
  !gapprompt@gap>| !gapinput@IsBlistRep(l);|
  false
  !gapprompt@gap>| !gapinput@ConvertToBlistRep(l);|
  true
\end{Verbatim}
 }

 }

 }

          
\chapter{\textcolor{Chapter }{Row Vectors}}\label{Row Vectors}
\logpage{[ 23, 0, 0 ]}
\hyperdef{L}{X82C7E6CF7BA03391}{}
{
  Just as in mathematics, a vector in \textsf{GAP} is any object which supports appropriate addition and scalar multiplication
operations (see Chapter{\nobreakspace}\ref{Vector Spaces}). As in mathematics, an especially important class of vectors are those
represented by a list of coefficients with respect to some basis. These
correspond roughly to the \textsf{GAP} concept of \emph{row vectors}.       
\section{\textcolor{Chapter }{IsRowVector (Filter)}}\label{sect:IsRowVector}
\logpage{[ 23, 1, 0 ]}
\hyperdef{L}{X7E383689817D2371}{}
{
  

\subsection{\textcolor{Chapter }{IsRowVector}}
\logpage{[ 23, 1, 1 ]}\nobreak
\hyperdef{L}{X7DFB22A07836A7A9}{}
{\noindent\textcolor{FuncColor}{$\triangleright$\enspace\texttt{IsRowVector({\mdseries\slshape obj})\index{IsRowVector@\texttt{IsRowVector}}
\label{IsRowVector}
}\hfill{\scriptsize (Category)}}\\


 A \emph{row vector} is a vector (see{\nobreakspace}\texttt{IsVector} (\ref{IsVector})) that is also a homogeneous list of odd additive nesting depth
(see{\nobreakspace}\ref{Filters Controlling the Arithmetic Behaviour of Lists}). Typical examples are lists of integers and rationals, lists of finite field
elements of the same characteristic, and lists of polynomials from a common
polynomial ring. Note that matrices are \emph{not} regarded as row vectors, because they have even additive nesting depth. 

 The additive operations of the vector must thus be compatible with that for
lists, implying that the list entries are the coefficients of the vector with
respect to some basis. 

 Note that not all row vectors admit a multiplication via \texttt{*} (which is to be understood as a scalar product); for example, class functions
are row vectors but the product of two class functions is defined in a
different way. For the installation of a scalar product of row vectors, the
entries of the vector must be ring elements; note that the default method
expects the row vectors to lie in \texttt{IsRingElementList}, and this category may not be implied by \texttt{IsRingElement} (\ref{IsRingElement}) for all entries of the row vector (see the comment in \texttt{IsVector} (\ref{IsVector})). 

 Note that methods for special types of row vectors really must be installed
with the requirement \texttt{IsRowVector}, since \texttt{IsVector} (\ref{IsVector}) may lead to a rank of the method below that of the default method for row
vectors (see file \texttt{lib/vecmat.gi}). 

 
\begin{Verbatim}[commandchars=!@|,fontsize=\small,frame=single,label=Example]
  !gapprompt@gap>| !gapinput@IsRowVector([1,2,3]);|
  true
\end{Verbatim}
 

 Because row vectors are just a special case of lists, all operations and
functions for lists are applicable to row vectors as well (see
Chapter{\nobreakspace}\ref{Lists}). This especially includes accessing elements of a row vector (see \ref{List Elements}), changing elements of a mutable row vector (see \ref{List Assignment}), and comparing row vectors (see \ref{Comparisons of Lists}). 

 Note that, unless your algorithms specifically require you to be able to
change entries of your vectors, it is generally better and faster to work with
immutable row vectors. See Section{\nobreakspace}\ref{Mutability and Copyability} for more details. }

 }

  
\section{\textcolor{Chapter }{Operators for Row Vectors}}\label{Operators for Row Vectors}
\logpage{[ 23, 2, 0 ]}
\hyperdef{L}{X85516C3179C229DB}{}
{
  The rules for arithmetic operations involving row vectors are in fact special
cases of those for the arithmetic of lists, as given in Section{\nobreakspace}\ref{Arithmetic for Lists} and the following sections, here we reiterate that definition, in the language
of vectors. 

 Note that the additive behaviour sketched below is defined only for lists in
the category \texttt{IsGeneralizedRowVector} (\ref{IsGeneralizedRowVector}), and the multiplicative behaviour is defined only for lists in the category \texttt{IsMultiplicativeGeneralizedRowVector} (\ref{IsMultiplicativeGeneralizedRowVector}). 

 \index{addition!vectors} \texttt{\mbox{\texttt{\mdseries\slshape vec1}} + \mbox{\texttt{\mdseries\slshape vec2}}} 

 returns the sum of the two row vectors \mbox{\texttt{\mdseries\slshape vec1}} and \mbox{\texttt{\mdseries\slshape vec2}}. Probably the most usual situation is that \mbox{\texttt{\mdseries\slshape vec1}} and \mbox{\texttt{\mdseries\slshape vec2}} have the same length and are defined over a common field; in this case the sum
is a new row vector over the same field where each entry is the sum of the
corresponding entries of the vectors. 

 In more general situations, the sum of two row vectors need not be a row
vector, for example adding an integer vector \mbox{\texttt{\mdseries\slshape vec1}} and a vector \mbox{\texttt{\mdseries\slshape vec2}} over a finite field yields the list of pointwise sums, which will be a mixture
of finite field elements and integers if \mbox{\texttt{\mdseries\slshape vec1}} is longer than \mbox{\texttt{\mdseries\slshape vec2}}. 

 \index{addition!vector and scalar} \texttt{\mbox{\texttt{\mdseries\slshape scalar}} + \mbox{\texttt{\mdseries\slshape vec}}} 

 \texttt{\mbox{\texttt{\mdseries\slshape vec}} + \mbox{\texttt{\mdseries\slshape scalar}}} 

 returns the sum of the scalar \mbox{\texttt{\mdseries\slshape scalar}} and the row vector \mbox{\texttt{\mdseries\slshape vec}}. Probably the most usual situation is that the elements of \mbox{\texttt{\mdseries\slshape vec}} lie in a common field with \mbox{\texttt{\mdseries\slshape scalar}}; in this case the sum is a new row vector over the same field where each
entry is the sum of the scalar and the corresponding entry of the vector. 

 More general situations are for example the sum of an integer scalar and a
vector over a finite field, or the sum of a finite field element and an
integer vector. 

 
\begin{Verbatim}[commandchars=!@|,fontsize=\small,frame=single,label=Example]
  !gapprompt@gap>| !gapinput@[ 1, 2, 3 ] + [ 1/2, 1/3, 1/4 ];|
  [ 3/2, 7/3, 13/4 ]
  !gapprompt@gap>| !gapinput@ [ 1/2, 3/2, 1/2 ] + 1/2;|
  [ 1, 2, 1 ]
\end{Verbatim}
 

 \index{subtraction!vectors} \index{subtraction!scalar and vector} \index{subtraction!vector and scalar} \texttt{\mbox{\texttt{\mdseries\slshape vec1}} \texttt{\symbol{45}} \mbox{\texttt{\mdseries\slshape vec2}}} 

 \texttt{\mbox{\texttt{\mdseries\slshape scalar}} \texttt{\symbol{45}} \mbox{\texttt{\mdseries\slshape vec}}} 

 \texttt{\mbox{\texttt{\mdseries\slshape vec}} \texttt{\symbol{45}} \mbox{\texttt{\mdseries\slshape scalar}}} 

 Subtracting a vector or scalar is defined as adding its additive inverse, so
the statements for the addition hold likewise. 

 
\begin{Verbatim}[commandchars=!@|,fontsize=\small,frame=single,label=Example]
  !gapprompt@gap>| !gapinput@[ 1, 2, 3 ] - [ 1/2, 1/3, 1/4 ];|
  [ 1/2, 5/3, 11/4 ]
  !gapprompt@gap>| !gapinput@[ 1/2, 3/2, 1/2 ] - 1/2;|
  [ 0, 1, 0 ]
\end{Verbatim}
 

 \index{multiplication!scalar and vector} \index{multiplication!vector and scalar} \texttt{\mbox{\texttt{\mdseries\slshape scalar}} * \mbox{\texttt{\mdseries\slshape vec}}} 

 \texttt{\mbox{\texttt{\mdseries\slshape vec}} * \mbox{\texttt{\mdseries\slshape scalar}}} 

 returns the product of the scalar \mbox{\texttt{\mdseries\slshape scalar}} and the row vector \mbox{\texttt{\mdseries\slshape vec}}. Probably the most usual situation is that the elements of \mbox{\texttt{\mdseries\slshape vec}} lie in a common field with \mbox{\texttt{\mdseries\slshape scalar}}; in this case the product is a new row vector over the same field where each
entry is the product of the scalar and the corresponding entry of the vector. 

 More general situations are for example the product of an integer scalar and a
vector over a finite field, or the product of a finite field element and an
integer vector. 

 
\begin{Verbatim}[commandchars=!@|,fontsize=\small,frame=single,label=Example]
  !gapprompt@gap>| !gapinput@[ 1/2, 3/2, 1/2 ] * 2;|
  [ 1, 3, 1 ]
\end{Verbatim}
 

 \index{multiplication!vectors} \texttt{\mbox{\texttt{\mdseries\slshape vec1}} * \mbox{\texttt{\mdseries\slshape vec2}}} 

 returns the standard scalar product of \mbox{\texttt{\mdseries\slshape vec1}} and \mbox{\texttt{\mdseries\slshape vec2}}, i.e., the sum of the products of the corresponding entries of the vectors.
Probably the most usual situation is that \mbox{\texttt{\mdseries\slshape vec1}} and \mbox{\texttt{\mdseries\slshape vec2}} have the same length and are defined over a common field; in this case the sum
is an element of this field. 

 More general situations are for example the inner product of an integer vector
and a vector over a finite field, or the inner product of two row vectors of
different lengths. 

 
\begin{Verbatim}[commandchars=!@|,fontsize=\small,frame=single,label=Example]
  !gapprompt@gap>| !gapinput@[ 1, 2, 3 ] * [ 1/2, 1/3, 1/4 ];|
  23/12
\end{Verbatim}
 

 For the mutability of results of arithmetic operations, see{\nobreakspace}\ref{Mutability and Copyability}. 

 Further operations with vectors as operands are defined by the matrix
operations, see{\nobreakspace}\ref{Operators for Matrices}. 

\subsection{\textcolor{Chapter }{NormedRowVector}}
\logpage{[ 23, 2, 1 ]}\nobreak
\hyperdef{L}{X785DC60D8482695D}{}
{\noindent\textcolor{FuncColor}{$\triangleright$\enspace\texttt{NormedRowVector({\mdseries\slshape v})\index{NormedRowVector@\texttt{NormedRowVector}}
\label{NormedRowVector}
}\hfill{\scriptsize (attribute)}}\\


 returns a scalar multiple \texttt{\mbox{\texttt{\mdseries\slshape w}} = \mbox{\texttt{\mdseries\slshape c}} * \mbox{\texttt{\mdseries\slshape v}}} of the row vector \mbox{\texttt{\mdseries\slshape v}} with the property that the first nonzero entry of \mbox{\texttt{\mdseries\slshape w}} is an identity element in the sense of \texttt{IsOne} (\ref{IsOne}). 

 
\begin{Verbatim}[commandchars=!@|,fontsize=\small,frame=single,label=Example]
  !gapprompt@gap>| !gapinput@NormedRowVector( [ 5, 2, 3 ] );|
  [ 1, 2/5, 3/5 ]
\end{Verbatim}
 }

 }

  
\section{\textcolor{Chapter }{Row Vectors over Finite Fields}}\label{Row Vectors over Finite Fields}
\logpage{[ 23, 3, 0 ]}
\hyperdef{L}{X8679F7DD7DFCBD9C}{}
{
  \textsf{GAP} can use compact formats to store row vectors over fields of order at most 256,
based on those used by the Meat\texttt{\symbol{45}}Axe \cite{Rin93}. This format also permits extremely efficient vector arithmetic. On the other
hand element access and assignment is significantly slower than for plain
lists. 

 The function \texttt{ConvertToVectorRep} (\ref{ConvertToVectorRep:for a list (and a field)}) is used to convert a list into a compressed vector, or to rewrite a compressed
vector over another field. Note that this function is \emph{much} faster when it is given a field (or field size) as an argument, rather than
having to scan the vector and try to decide the field. Supplying the field can
also avoid errors and/or loss of performance, when one vector from some
collection happens to have all of its entries over a smaller field than the ``natural'' field of the problem. 
\subsection{\textcolor{Chapter }{ConvertToVectorRep}}\logpage{[ 23, 3, 1 ]}
\hyperdef{L}{X810E46927F9E8F75}{}
{
\noindent\textcolor{FuncColor}{$\triangleright$\enspace\texttt{ConvertToVectorRep({\mdseries\slshape list[, field]})\index{ConvertToVectorRep@\texttt{ConvertToVectorRep}!for a list (and a field)}
\label{ConvertToVectorRep:for a list (and a field)}
}\hfill{\scriptsize (function)}}\\
\noindent\textcolor{FuncColor}{$\triangleright$\enspace\texttt{ConvertToVectorRep({\mdseries\slshape list[, fieldsize]})\index{ConvertToVectorRep@\texttt{ConvertToVectorRep}!for a list (and a prime power)}
\label{ConvertToVectorRep:for a list (and a prime power)}
}\hfill{\scriptsize (function)}}\\
\noindent\textcolor{FuncColor}{$\triangleright$\enspace\texttt{ConvertToVectorRepNC({\mdseries\slshape list[, field]})\index{ConvertToVectorRepNC@\texttt{ConvertToVectorRepNC}!for a list (and a field)}
\label{ConvertToVectorRepNC:for a list (and a field)}
}\hfill{\scriptsize (function)}}\\
\noindent\textcolor{FuncColor}{$\triangleright$\enspace\texttt{ConvertToVectorRepNC({\mdseries\slshape list[, fieldsize]})\index{ConvertToVectorRepNC@\texttt{ConvertToVectorRepNC}!for a list (and a prime power)}
\label{ConvertToVectorRepNC:for a list (and a prime power)}
}\hfill{\scriptsize (function)}}\\


 Called with one argument \mbox{\texttt{\mdseries\slshape list}}, \texttt{ConvertToVectorRep} converts \mbox{\texttt{\mdseries\slshape list}} to an internal row vector representation if possible. 

 Called with a list \mbox{\texttt{\mdseries\slshape list}} and a finite field \mbox{\texttt{\mdseries\slshape field}}, \texttt{ConvertToVectorRep} converts \mbox{\texttt{\mdseries\slshape list}} to an internal row vector representation appropriate for a row vector over \mbox{\texttt{\mdseries\slshape field}}. 

 Instead of a \mbox{\texttt{\mdseries\slshape field}} also its size \mbox{\texttt{\mdseries\slshape fieldsize}} may be given. 

 It is forbidden to call this function unless \mbox{\texttt{\mdseries\slshape list}} is a plain list or a row vector, \mbox{\texttt{\mdseries\slshape field}} is a field, and all elements of \mbox{\texttt{\mdseries\slshape list}} lie in \mbox{\texttt{\mdseries\slshape field}}. Violation of this condition can lead to unpredictable behaviour or a system
crash. (Setting the assertion level to at least 2 might catch some violations
before a crash, see{\nobreakspace}\texttt{SetAssertionLevel} (\ref{SetAssertionLevel}).) 

 \mbox{\texttt{\mdseries\slshape list}} may already be a compressed vector. In this case, if no \mbox{\texttt{\mdseries\slshape field}} or \mbox{\texttt{\mdseries\slshape fieldsize}} is given, then nothing happens. If one is given then the vector is rewritten
as a compressed vector over the given \mbox{\texttt{\mdseries\slshape field}} unless it has the filter \texttt{IsLockedRepresentationVector}, in which case it is not changed. 

 The return value is the size of the field over which the vector ends up
written, if it is written in a compressed representation. 

 In this example, we first create a row vector and then ask \textsf{GAP} to rewrite it, first over \texttt{GF(2)} and then over \texttt{GF(4)}. 

 
\begin{Verbatim}[commandchars=!@|,fontsize=\small,frame=single,label=Example]
  !gapprompt@gap>| !gapinput@v := [Z(2)^0,Z(2),Z(2),0*Z(2)];|
  [ Z(2)^0, Z(2)^0, Z(2)^0, 0*Z(2) ]
  !gapprompt@gap>| !gapinput@RepresentationsOfObject(v);|
  [ "IsPlistRep", "IsInternalRep" ]
  !gapprompt@gap>| !gapinput@ConvertToVectorRep(v);|
  2
  !gapprompt@gap>| !gapinput@v;|
  <a GF2 vector of length 4>
  !gapprompt@gap>| !gapinput@ConvertToVectorRep(v,4);|
  4
  !gapprompt@gap>| !gapinput@v;|
  [ Z(2)^0, Z(2)^0, Z(2)^0, 0*Z(2) ]
  !gapprompt@gap>| !gapinput@RepresentationsOfObject(v);|
  [ "IsDataObjectRep", "Is8BitVectorRep" ]
\end{Verbatim}
 

 A vector in the special representation over \texttt{GF(2)} is always viewed as \texttt{{\textless}a GF2 vector of length ...{\textgreater}}. Over fields of orders 3 to 256, a vector of length 10 or less is viewed as
the list of its coefficients, but a longer one is abbreviated. 

 Arithmetic operations (see{\nobreakspace}\ref{Arithmetic for Lists} and the following sections) preserve the compression status of row vectors in
the sense that if all arguments are compressed row vectors written over the
same field and the result is a row vector then also the result is a compressed
row vector written over this field. }

 

\subsection{\textcolor{Chapter }{ImmutableVector}}
\logpage{[ 23, 3, 2 ]}\nobreak
\hyperdef{L}{X83D8F5BB80089279}{}
{\noindent\textcolor{FuncColor}{$\triangleright$\enspace\texttt{ImmutableVector({\mdseries\slshape field, vector[, change]})\index{ImmutableVector@\texttt{ImmutableVector}}
\label{ImmutableVector}
}\hfill{\scriptsize (operation)}}\\


 Let \mbox{\texttt{\mdseries\slshape vector}} be an object for which \texttt{IsRowVector} (\ref{IsRowVector}) or \texttt{IsVectorObj} (\ref{IsVectorObj}) returns \texttt{true}. In the former case, \mbox{\texttt{\mdseries\slshape vector}} is a list, and \texttt{ImmutableVector} returns an immutable object for which \texttt{IsRowVector} (\ref{IsRowVector}) returns \texttt{true} (in particular again a list), which is equal to \mbox{\texttt{\mdseries\slshape vector}}, and which is in the optimal (concerning space and runtime) representation
for vectors defined over \mbox{\texttt{\mdseries\slshape field}}, provided that the entries of \mbox{\texttt{\mdseries\slshape vector}} lie in \mbox{\texttt{\mdseries\slshape field}}. In the latter case, if \mbox{\texttt{\mdseries\slshape vector}} is not in \texttt{IsRowVector} (\ref{IsRowVector}), \texttt{ImmutableVector} returns an immutable object that is equal to the result of \texttt{ChangedBaseDomain} (\ref{ChangedBaseDomain:for a vector object}) when this is called with \mbox{\texttt{\mdseries\slshape vector}} and \mbox{\texttt{\mdseries\slshape field}}. 

 This means that vectors obtained by several calls of \texttt{ImmutableVector} for the same \mbox{\texttt{\mdseries\slshape field}} are compatible for fast arithmetic without need for field conversion. 

 If the input vector \mbox{\texttt{\mdseries\slshape vector}} is in \texttt{IsRowVector} (\ref{IsRowVector}) then it might change its representation as a side effect of this function.
However, one cannot rely on this side effect. Also, if \mbox{\texttt{\mdseries\slshape vector}} is already immutable and the result of \texttt{ImmutableVector} has the same internal representation as \mbox{\texttt{\mdseries\slshape vector}}, the result is not necessarily \emph{identical} to \mbox{\texttt{\mdseries\slshape vector}}. 

 If \mbox{\texttt{\mdseries\slshape change}} is \texttt{true}, then \mbox{\texttt{\mdseries\slshape vector}} may be changed to become immutable; otherwise it is copied first. }

 

\subsection{\textcolor{Chapter }{NumberFFVector}}
\logpage{[ 23, 3, 3 ]}\nobreak
\hyperdef{L}{X872E17FF829DB50F}{}
{\noindent\textcolor{FuncColor}{$\triangleright$\enspace\texttt{NumberFFVector({\mdseries\slshape vec, sz})\index{NumberFFVector@\texttt{NumberFFVector}}
\label{NumberFFVector}
}\hfill{\scriptsize (operation)}}\\


 returns an integer that gives the position minus one of the finite field row
vector \mbox{\texttt{\mdseries\slshape vec}} in the sorted list of all row vectors over the field with \mbox{\texttt{\mdseries\slshape sz}} elements in the same dimension as \mbox{\texttt{\mdseries\slshape vec}}. \texttt{NumberFFVector} returns \texttt{fail} if the vector cannot be represented over the field with \mbox{\texttt{\mdseries\slshape sz}} elements. 

 
\begin{Verbatim}[commandchars=!@|,fontsize=\small,frame=single,label=Example]
  !gapprompt@gap>| !gapinput@v:=[0,1,2,0]*Z(3);;|
  !gapprompt@gap>| !gapinput@NumberFFVector(v, 3);|
  21
  !gapprompt@gap>| !gapinput@NumberFFVector(Zero(v),3);|
  0
  !gapprompt@gap>| !gapinput@V:=EnumeratorSorted(GF(3)^4);|
  <enumerator of ( GF(3)^4 )>
  !gapprompt@gap>| !gapinput@V[21+1] = v;|
  true
\end{Verbatim}
 }

 }

  
\section{\textcolor{Chapter }{Coefficient List Arithmetic}}\label{Coefficient List Arithmetic}
\logpage{[ 23, 4, 0 ]}
\hyperdef{L}{X85C68AED805E4B9C}{}
{
  The following operations all perform arithmetic on row vectors, given as
homogeneous lists of the same length, containing elements of a commutative
ring. 

 There are two reasons for using \texttt{AddRowVector} (\ref{AddRowVector}) in preference to arithmetic operators. Firstly, the three argument form has no
single\texttt{\symbol{45}}step equivalent. Secondly \texttt{AddRowVector} (\ref{AddRowVector}) changes its first argument in\texttt{\symbol{45}}place, rather than allocating
a new vector to hold the result, and may thus produce less garbage. 

\subsection{\textcolor{Chapter }{AddVector}}
\logpage{[ 23, 4, 1 ]}\nobreak
\hyperdef{L}{X800EC03F7E0A5F23}{}
{\noindent\textcolor{FuncColor}{$\triangleright$\enspace\texttt{AddVector({\mdseries\slshape dst, src[, mul[, from, to]]})\index{AddVector@\texttt{AddVector}}
\label{AddVector}
}\hfill{\scriptsize (operation)}}\\
\noindent\textcolor{FuncColor}{$\triangleright$\enspace\texttt{AddRowVector({\mdseries\slshape dst, src[, mul[, from, to]]})\index{AddRowVector@\texttt{AddRowVector}}
\label{AddRowVector}
}\hfill{\scriptsize (operation)}}\\


 Adds the product of \mbox{\texttt{\mdseries\slshape src}} and \mbox{\texttt{\mdseries\slshape mul}} to \mbox{\texttt{\mdseries\slshape dst}}, changing \mbox{\texttt{\mdseries\slshape dst}}. If \mbox{\texttt{\mdseries\slshape from}} and \mbox{\texttt{\mdseries\slshape to}} are given then only the index range \texttt{[ \mbox{\texttt{\mdseries\slshape from}} .. \mbox{\texttt{\mdseries\slshape to}} ]} is guaranteed to be affected. Other indices \emph{may} be affected, if it is more convenient to do so. Even when \mbox{\texttt{\mdseries\slshape from}} and \mbox{\texttt{\mdseries\slshape to}} are given, \mbox{\texttt{\mdseries\slshape dst}} and \mbox{\texttt{\mdseries\slshape src}} must be row vectors of the \emph{same} length. 

 If \mbox{\texttt{\mdseries\slshape mul}} is not given either then this operation simply adds \mbox{\texttt{\mdseries\slshape src}} to \mbox{\texttt{\mdseries\slshape dst}}. }

 

\subsection{\textcolor{Chapter }{AddCoeffs}}
\logpage{[ 23, 4, 2 ]}\nobreak
\hyperdef{L}{X7854B2B67E3FE2CA}{}
{\noindent\textcolor{FuncColor}{$\triangleright$\enspace\texttt{AddCoeffs({\mdseries\slshape list1[, poss1], list2[, poss2[, mul]]})\index{AddCoeffs@\texttt{AddCoeffs}}
\label{AddCoeffs}
}\hfill{\scriptsize (operation)}}\\


 \texttt{AddCoeffs} adds the entries of \mbox{\texttt{\mdseries\slshape list2}}\texttt{\texttt{\symbol{123}}}\mbox{\texttt{\mdseries\slshape poss2}}\texttt{\texttt{\symbol{125}}}, multiplied by the scalar \mbox{\texttt{\mdseries\slshape mul}}, to \mbox{\texttt{\mdseries\slshape list1}}\texttt{\texttt{\symbol{123}}}\mbox{\texttt{\mdseries\slshape poss1}}\texttt{\texttt{\symbol{125}}}. Unbound entries in \mbox{\texttt{\mdseries\slshape list1}} are assumed to be zero. The position of the right\texttt{\symbol{45}}most
non\texttt{\symbol{45}}zero element is returned. 

 If the ranges \mbox{\texttt{\mdseries\slshape poss1}} and \mbox{\texttt{\mdseries\slshape poss2}} are not given, they are assumed to span the whole vectors. If the scalar \mbox{\texttt{\mdseries\slshape mul}} is omitted, one is used as a default. 

 Note that it is the responsibility of the caller to ensure that \mbox{\texttt{\mdseries\slshape list2}} has elements at position \mbox{\texttt{\mdseries\slshape poss2}} and that the result (in \mbox{\texttt{\mdseries\slshape list1}}) will be a dense list. 

 The function is free to remove trailing (right\texttt{\symbol{45}}most) zeros. 
\begin{Verbatim}[commandchars=!@|,fontsize=\small,frame=single,label=Example]
  !gapprompt@gap>| !gapinput@l:=[1,2,3,4];;m:=[5,6,7];;AddCoeffs(l,m);|
  4
  !gapprompt@gap>| !gapinput@l;|
  [ 6, 8, 10, 4 ]
\end{Verbatim}
 }

 

\subsection{\textcolor{Chapter }{MultVector}}
\logpage{[ 23, 4, 3 ]}\nobreak
\hyperdef{L}{X7BEF28C981C42E16}{}
{\noindent\textcolor{FuncColor}{$\triangleright$\enspace\texttt{MultVector({\mdseries\slshape list, mul})\index{MultVector@\texttt{MultVector}}
\label{MultVector}
}\hfill{\scriptsize (operation)}}\\
\noindent\textcolor{FuncColor}{$\triangleright$\enspace\texttt{MultVectorLeft({\mdseries\slshape list, mul})\index{MultVectorLeft@\texttt{MultVectorLeft}}
\label{MultVectorLeft}
}\hfill{\scriptsize (operation)}}\\
\noindent\textcolor{FuncColor}{$\triangleright$\enspace\texttt{MultVectorRight({\mdseries\slshape list, mul})\index{MultVectorRight@\texttt{MultVectorRight}}
\label{MultVectorRight}
}\hfill{\scriptsize (operation)}}\\
\textbf{\indent Returns:\ }
nothing



 The operation \texttt{MultVectorLeft} multiplies the entries of \mbox{\texttt{\mdseries\slshape list}} by the scalar \mbox{\texttt{\mdseries\slshape mul}} from the left, in\texttt{\symbol{45}}place. 

 The operation \texttt{MultVectorRight} multiplies the entries of \mbox{\texttt{\mdseries\slshape list}} by the scalar \mbox{\texttt{\mdseries\slshape mul}} from the right, in\texttt{\symbol{45}}place. 

 Note that \texttt{MultVector} is just a synonym for \texttt{MultVectorLeft}. 
\begin{Verbatim}[commandchars=!@|,fontsize=\small,frame=single,label=Example]
  !gapprompt@gap>| !gapinput@A:= FreeAssociativeAlgebraWithOne( Rationals, "a", "b", "c" );;|
  !gapprompt@gap>| !gapinput@l:=[A.1,A.2];; MultVectorLeft(l, A.3); l;|
  [ (1)*c*a, (1)*c*b ]
  !gapprompt@gap>| !gapinput@l:=[A.1,A.2];; MultVectorRight(l, A.3); l;|
  [ (1)*a*c, (1)*b*c ]
\end{Verbatim}
 }

 

\subsection{\textcolor{Chapter }{CoeffsMod}}
\logpage{[ 23, 4, 4 ]}\nobreak
\hyperdef{L}{X8264B3EE7D56EEDD}{}
{\noindent\textcolor{FuncColor}{$\triangleright$\enspace\texttt{CoeffsMod({\mdseries\slshape list1[, len1], modulus})\index{CoeffsMod@\texttt{CoeffsMod}}
\label{CoeffsMod}
}\hfill{\scriptsize (operation)}}\\


 returns the coefficient list obtained by reducing the entries in \mbox{\texttt{\mdseries\slshape list1}} modulo \mbox{\texttt{\mdseries\slshape modulus}}. After reducing it shrinks the list to remove trailing zeroes. If the
optional argument \mbox{\texttt{\mdseries\slshape len1}} is used, it reduces only first \mbox{\texttt{\mdseries\slshape len1}} elements of the list. 
\begin{Verbatim}[commandchars=!@|,fontsize=\small,frame=single,label=Example]
  !gapprompt@gap>| !gapinput@l:=[1,2,3,4];;CoeffsMod(l,2);|
  [ 1, 0, 1 ]
\end{Verbatim}
 }

 }

  
\section{\textcolor{Chapter }{Shifting and Trimming Coefficient Lists}}\label{Shifting and Trimming Coefficient Lists}
\logpage{[ 23, 5, 0 ]}
\hyperdef{L}{X7D287281781E16A2}{}
{
  The following functions change coefficient lists by shifting or trimming. 

\subsection{\textcolor{Chapter }{LeftShiftRowVector}}
\logpage{[ 23, 5, 1 ]}\nobreak
\hyperdef{L}{X80465E9B7A38C176}{}
{\noindent\textcolor{FuncColor}{$\triangleright$\enspace\texttt{LeftShiftRowVector({\mdseries\slshape list, shift})\index{LeftShiftRowVector@\texttt{LeftShiftRowVector}}
\label{LeftShiftRowVector}
}\hfill{\scriptsize (operation)}}\\


 changes \mbox{\texttt{\mdseries\slshape list}} by assigning \mbox{\texttt{\mdseries\slshape list}}$[i]$\texttt{:= }\mbox{\texttt{\mdseries\slshape list}}$[i+\mbox{\texttt{\mdseries\slshape shift}}]$ and removing the last \mbox{\texttt{\mdseries\slshape shift}} entries of the result. }

 

\subsection{\textcolor{Chapter }{RightShiftRowVector}}
\logpage{[ 23, 5, 2 ]}\nobreak
\hyperdef{L}{X822CCA4781D5C5EC}{}
{\noindent\textcolor{FuncColor}{$\triangleright$\enspace\texttt{RightShiftRowVector({\mdseries\slshape list, shift, fill})\index{RightShiftRowVector@\texttt{RightShiftRowVector}}
\label{RightShiftRowVector}
}\hfill{\scriptsize (operation)}}\\


 changes \mbox{\texttt{\mdseries\slshape list}} by assigning \mbox{\texttt{\mdseries\slshape list}}$[i+\mbox{\texttt{\mdseries\slshape shift}}]$\texttt{:= }\mbox{\texttt{\mdseries\slshape list}}$[i]$ and filling each of the \mbox{\texttt{\mdseries\slshape shift}} first entries with \mbox{\texttt{\mdseries\slshape fill}}. }

 

\subsection{\textcolor{Chapter }{ShrinkRowVector}}
\logpage{[ 23, 5, 3 ]}\nobreak
\hyperdef{L}{X78951C0E86D857B5}{}
{\noindent\textcolor{FuncColor}{$\triangleright$\enspace\texttt{ShrinkRowVector({\mdseries\slshape list})\index{ShrinkRowVector@\texttt{ShrinkRowVector}}
\label{ShrinkRowVector}
}\hfill{\scriptsize (operation)}}\\


 removes trailing zeroes from the list \mbox{\texttt{\mdseries\slshape list}}. 
\begin{Verbatim}[commandchars=!@|,fontsize=\small,frame=single,label=Example]
  !gapprompt@gap>| !gapinput@l:=[1,0,0];;ShrinkRowVector(l);l;|
  [ 1 ]
\end{Verbatim}
 }

 

\subsection{\textcolor{Chapter }{RemoveOuterCoeffs}}
\logpage{[ 23, 5, 4 ]}\nobreak
\hyperdef{L}{X85796B6079581023}{}
{\noindent\textcolor{FuncColor}{$\triangleright$\enspace\texttt{RemoveOuterCoeffs({\mdseries\slshape list, coef})\index{RemoveOuterCoeffs@\texttt{RemoveOuterCoeffs}}
\label{RemoveOuterCoeffs}
}\hfill{\scriptsize (operation)}}\\


 removes \mbox{\texttt{\mdseries\slshape coef}} at the beginning and at the end of \mbox{\texttt{\mdseries\slshape list}} and returns the number of elements removed at the beginning. 
\begin{Verbatim}[commandchars=!@|,fontsize=\small,frame=single,label=Example]
  !gapprompt@gap>| !gapinput@l:=[1,1,2,1,2,1,1,2,1];; RemoveOuterCoeffs(l,1);|
  2
  !gapprompt@gap>| !gapinput@l;|
  [ 2, 1, 2, 1, 1, 2 ]
\end{Verbatim}
 }

 }

  
\section{\textcolor{Chapter }{Functions for Coding Theory}}\label{Functions for Coding Theory}
\logpage{[ 23, 6, 0 ]}
\hyperdef{L}{X7B63F1EB83FA0CF6}{}
{
  The following functions perform operations on finite fields vectors considered
as code words in a linear code. 

\subsection{\textcolor{Chapter }{WeightVecFFE}}
\logpage{[ 23, 6, 1 ]}\nobreak
\hyperdef{L}{X7C9F4D657F9BA5A1}{}
{\noindent\textcolor{FuncColor}{$\triangleright$\enspace\texttt{WeightVecFFE({\mdseries\slshape vec})\index{WeightVecFFE@\texttt{WeightVecFFE}}
\label{WeightVecFFE}
}\hfill{\scriptsize (operation)}}\\


 returns the weight of the finite field vector \mbox{\texttt{\mdseries\slshape vec}}, i.e. the number of nonzero entries. }

 

\subsection{\textcolor{Chapter }{DistanceVecFFE}}
\logpage{[ 23, 6, 2 ]}\nobreak
\hyperdef{L}{X85AA5C6587559C1C}{}
{\noindent\textcolor{FuncColor}{$\triangleright$\enspace\texttt{DistanceVecFFE({\mdseries\slshape vec1, vec2})\index{DistanceVecFFE@\texttt{DistanceVecFFE}}
\label{DistanceVecFFE}
}\hfill{\scriptsize (operation)}}\\


 returns the distance between the two vectors \mbox{\texttt{\mdseries\slshape vec1}} and \mbox{\texttt{\mdseries\slshape vec2}}, which must have the same length and whose elements must lie in a common
field. The distance is the number of places where \mbox{\texttt{\mdseries\slshape vec1}} and \mbox{\texttt{\mdseries\slshape vec2}} differ. }

 

\subsection{\textcolor{Chapter }{DistancesDistributionVecFFEsVecFFE}}
\logpage{[ 23, 6, 3 ]}\nobreak
\hyperdef{L}{X7F2F630984A9D3D6}{}
{\noindent\textcolor{FuncColor}{$\triangleright$\enspace\texttt{DistancesDistributionVecFFEsVecFFE({\mdseries\slshape vecs, vec})\index{DistancesDistributionVecFFEsVecFFE@\texttt{DistancesDistributionVecFFEsVecFFE}}
\label{DistancesDistributionVecFFEsVecFFE}
}\hfill{\scriptsize (operation)}}\\


 returns the distances distribution of the vector \mbox{\texttt{\mdseries\slshape vec}} to the vectors in the list \mbox{\texttt{\mdseries\slshape vecs}}. All vectors must have the same length, and all elements must lie in a common
field. The distances distribution is a list $d$ of length \texttt{Length(\mbox{\texttt{\mdseries\slshape vec}})+1}, such that the value $d[i]$ is the number of vectors in \mbox{\texttt{\mdseries\slshape vecs}} that have distance $i+1$ to \mbox{\texttt{\mdseries\slshape vec}}. }

 

\subsection{\textcolor{Chapter }{DistancesDistributionMatFFEVecFFE}}
\logpage{[ 23, 6, 4 ]}\nobreak
\hyperdef{L}{X85135CEB86E61D49}{}
{\noindent\textcolor{FuncColor}{$\triangleright$\enspace\texttt{DistancesDistributionMatFFEVecFFE({\mdseries\slshape mat, F, vec})\index{DistancesDistributionMatFFEVecFFE@\texttt{DistancesDistributionMatFFEVecFFE}}
\label{DistancesDistributionMatFFEVecFFE}
}\hfill{\scriptsize (operation)}}\\


 returns the distances distribution of the vector \mbox{\texttt{\mdseries\slshape vec}} to the vectors in the vector space generated by the rows of the matrix \mbox{\texttt{\mdseries\slshape mat}} over the finite field \mbox{\texttt{\mdseries\slshape F}}. The length of the rows of \mbox{\texttt{\mdseries\slshape mat}} and the length of \mbox{\texttt{\mdseries\slshape vec}} must be equal, and all entries must lie in \mbox{\texttt{\mdseries\slshape F}}. The rows of \mbox{\texttt{\mdseries\slshape mat}} must be linearly independent. The distances distribution is a list $d$ of length \texttt{Length(\mbox{\texttt{\mdseries\slshape vec}})+1}, such that the value $d[i]$ is the number of vectors in the vector space generated by the rows of \mbox{\texttt{\mdseries\slshape mat}} that have distance $i+1$ to \mbox{\texttt{\mdseries\slshape vec}}. }

 

\subsection{\textcolor{Chapter }{AClosestVectorCombinationsMatFFEVecFFE}}
\logpage{[ 23, 6, 5 ]}\nobreak
\hyperdef{L}{X82E5987E81487D18}{}
{\noindent\textcolor{FuncColor}{$\triangleright$\enspace\texttt{AClosestVectorCombinationsMatFFEVecFFE({\mdseries\slshape mat, f, vec, cnt, stop})\index{AClosestVectorCombinationsMatFFEVecFFE@\texttt{AClosest}\-\texttt{Vector}\-\texttt{Combinations}\-\texttt{Mat}\-\texttt{F}\-\texttt{F}\-\texttt{E}\-\texttt{VecFFE}}
\label{AClosestVectorCombinationsMatFFEVecFFE}
}\hfill{\scriptsize (operation)}}\\
\noindent\textcolor{FuncColor}{$\triangleright$\enspace\texttt{AClosestVectorCombinationsMatFFEVecFFECoords({\mdseries\slshape mat, f, vec, cnt, stop})\index{AClosestVectorCombinationsMatFFEVecFFECoords@\texttt{AClosest}\-\texttt{Vector}\-\texttt{Combinations}\-\texttt{Mat}\-\texttt{F}\-\texttt{F}\-\texttt{E}\-\texttt{Vec}\-\texttt{F}\-\texttt{F}\-\texttt{E}\-\texttt{Coords}}
\label{AClosestVectorCombinationsMatFFEVecFFECoords}
}\hfill{\scriptsize (operation)}}\\


 These functions run through the \mbox{\texttt{\mdseries\slshape f}}\texttt{\symbol{45}}linear combinations of the vectors in the rows of the
matrix \mbox{\texttt{\mdseries\slshape mat}} that can be written as linear combinations of exactly \mbox{\texttt{\mdseries\slshape cnt}} rows (that is without using zero as a coefficient). The length of the rows of \mbox{\texttt{\mdseries\slshape mat}} and the length of \mbox{\texttt{\mdseries\slshape vec}} must be equal, and all elements must lie in the field \mbox{\texttt{\mdseries\slshape f}}. The rows of \mbox{\texttt{\mdseries\slshape mat}} must be linearly independent. \texttt{AClosestVectorCombinationsMatFFEVecFFE} returns a vector from these that is closest to the vector \mbox{\texttt{\mdseries\slshape vec}}. If it finds a vector of distance at most \mbox{\texttt{\mdseries\slshape stop}}, which must be a nonnegative integer, then it stops immediately and returns
this vector. 

 \texttt{AClosestVectorCombinationsMatFFEVecFFECoords} returns a length 2 list containing the same closest vector and also a vector \mbox{\texttt{\mdseries\slshape v}} with exactly \mbox{\texttt{\mdseries\slshape cnt}} non\texttt{\symbol{45}}zero entries, such that \mbox{\texttt{\mdseries\slshape v}} times \mbox{\texttt{\mdseries\slshape mat}} is the closest vector. }

 

\subsection{\textcolor{Chapter }{CosetLeadersMatFFE}}
\logpage{[ 23, 6, 6 ]}\nobreak
\hyperdef{L}{X7C88671678A2BEB4}{}
{\noindent\textcolor{FuncColor}{$\triangleright$\enspace\texttt{CosetLeadersMatFFE({\mdseries\slshape mat, f})\index{CosetLeadersMatFFE@\texttt{CosetLeadersMatFFE}}
\label{CosetLeadersMatFFE}
}\hfill{\scriptsize (operation)}}\\


 Let $C$ be the linear code over \mbox{\texttt{\mdseries\slshape f}} with check matrix \mbox{\texttt{\mdseries\slshape mat}}, that is, $C$ is the right null space of \mbox{\texttt{\mdseries\slshape mat}}: 
\[ C = \{ v \mid mat \cdot v = 0 \}. \]
 For each syndrome $s$, the solutions of the non\texttt{\symbol{45}}homogeneous system $\mbox{\texttt{\mdseries\slshape mat}} * v = s$ form one coset of $C$. Equivalently, two vectors lie in the same coset if their difference lies in $C$, just as in group theory. 

 \texttt{CosetLeadersMatFFE} returns, for each coset of $C$, a vector of smallest Hamming weight in that coset. Such a vector is called a \emph{coset leader}. 

 If \texttt{L := CosetLeadersMatFFE( \mbox{\texttt{\mdseries\slshape mat}}, \mbox{\texttt{\mdseries\slshape f}} )} and \texttt{q := Size( \mbox{\texttt{\mdseries\slshape f}} )}, then \texttt{L[i]} is the coset leader whose syndrome is the vector with number \texttt{i\texttt{\symbol{45}}1}; in other words, \texttt{NumberFFVector( mat * L[i], q ) = i \texttt{\symbol{45}} 1} for all valid indices \texttt{i}. 

 Thus the list is indexed by syndromes, not by codewords. 

 All rows of \mbox{\texttt{\mdseries\slshape mat}} must have the same length, all entries must lie in \mbox{\texttt{\mdseries\slshape f}}, and the rows of \mbox{\texttt{\mdseries\slshape mat}} must be linearly independent. }

 }

  
\section{\textcolor{Chapter }{Vectors as coefficients of polynomials}}\label{Vectors as coefficients of polynomials}
\logpage{[ 23, 7, 0 ]}
\hyperdef{L}{X87FEC1927B3A63C8}{}
{
  A list of ring elements can be interpreted as a row vector or the list of
coefficients of a polynomial. There are a couple of functions that implement
arithmetic operations based on these interpretations. \textsf{GAP} contains proper support for polynomials (see{\nobreakspace}\ref{Polynomials and Rational Functions}), the operations described in this section are on a lower level. 

 The following operations all perform arithmetic on univariate polynomials
given by their coefficient lists. These lists can have different lengths but
must be dense homogeneous lists containing elements of a commutative ring. Not
all input lists may be empty. 

 In the following descriptions we will always assume that \mbox{\texttt{\mdseries\slshape list1}} is the coefficient list of the polynomial \mbox{\texttt{\mdseries\slshape pol1}} and so forth. If length parameter \mbox{\texttt{\mdseries\slshape leni}} is not given, it is set to the length of \mbox{\texttt{\mdseries\slshape listi}} by default. 

\subsection{\textcolor{Chapter }{ValuePol}}
\logpage{[ 23, 7, 1 ]}\nobreak
\hyperdef{L}{X84DE99D57C29D47F}{}
{\noindent\textcolor{FuncColor}{$\triangleright$\enspace\texttt{ValuePol({\mdseries\slshape coeff, x})\index{ValuePol@\texttt{ValuePol}}
\label{ValuePol}
}\hfill{\scriptsize (operation)}}\\


 Let \mbox{\texttt{\mdseries\slshape coeff}} be the coefficients list of a univariate polynomial $f$, and \mbox{\texttt{\mdseries\slshape x}} a ring element. Then \texttt{ValuePol} returns the value $f(\mbox{\texttt{\mdseries\slshape x}})$. 

 The coefficient of $\mbox{\texttt{\mdseries\slshape x}}^i$ is assumed to be stored at position $i+1$ in the coefficients list. 
\begin{Verbatim}[commandchars=!@|,fontsize=\small,frame=single,label=Example]
  !gapprompt@gap>| !gapinput@ValuePol([1,2,3],4);|
  57
\end{Verbatim}
 }

         

\subsection{\textcolor{Chapter }{ProductCoeffs}}
\logpage{[ 23, 7, 2 ]}\nobreak
\hyperdef{L}{X8328088C807AFFAF}{}
{\noindent\textcolor{FuncColor}{$\triangleright$\enspace\texttt{ProductCoeffs({\mdseries\slshape list1[, len1], list2[, len2]})\index{ProductCoeffs@\texttt{ProductCoeffs}}
\label{ProductCoeffs}
}\hfill{\scriptsize (operation)}}\\


 Let $p1$ (and $p2$) be polynomials given by the first \mbox{\texttt{\mdseries\slshape len1}} (\mbox{\texttt{\mdseries\slshape len2}}) entries of the coefficient list \mbox{\texttt{\mdseries\slshape list2}} (\mbox{\texttt{\mdseries\slshape list2}}). If \mbox{\texttt{\mdseries\slshape len1}} and \mbox{\texttt{\mdseries\slshape len2}} are omitted, they default to the lengths of \mbox{\texttt{\mdseries\slshape list1}} and \mbox{\texttt{\mdseries\slshape list2}}. This operation returns the coefficient list of the product of $p1$ and $p2$. 
\begin{Verbatim}[commandchars=!@|,fontsize=\small,frame=single,label=Example]
  !gapprompt@gap>| !gapinput@l:=[1,2,3,4];;m:=[5,6,7];;ProductCoeffs(l,m);|
  [ 5, 16, 34, 52, 45, 28 ]
\end{Verbatim}
 }

 

\subsection{\textcolor{Chapter }{ReduceCoeffs}}
\logpage{[ 23, 7, 3 ]}\nobreak
\hyperdef{L}{X87248AA27F05BDCC}{}
{\noindent\textcolor{FuncColor}{$\triangleright$\enspace\texttt{ReduceCoeffs({\mdseries\slshape list1[, len1], list2[, len2]})\index{ReduceCoeffs@\texttt{ReduceCoeffs}}
\label{ReduceCoeffs}
}\hfill{\scriptsize (operation)}}\\


 Let $p1$ (and $p2$) be polynomials given by the first \mbox{\texttt{\mdseries\slshape len1}} (\mbox{\texttt{\mdseries\slshape len2}}) entries of the coefficient list \mbox{\texttt{\mdseries\slshape list1}} (\mbox{\texttt{\mdseries\slshape list2}}). If \mbox{\texttt{\mdseries\slshape len1}} and \mbox{\texttt{\mdseries\slshape len2}} are omitted, they default to the lengths of \mbox{\texttt{\mdseries\slshape list1}} and \mbox{\texttt{\mdseries\slshape list2}}. \texttt{ReduceCoeffs} changes \mbox{\texttt{\mdseries\slshape list1}} to the coefficient list of the remainder when dividing \mbox{\texttt{\mdseries\slshape p1}} by \mbox{\texttt{\mdseries\slshape p2}}. This operation changes \mbox{\texttt{\mdseries\slshape list1}} which therefore must be a mutable list. The operation returns the position of
the last non\texttt{\symbol{45}}zero entry of the result but is not guaranteed
to remove trailing zeroes. 
\begin{Verbatim}[commandchars=!@|,fontsize=\small,frame=single,label=Example]
  !gapprompt@gap>| !gapinput@l:=[1,2,3,4];;m:=[5,6,7];;ReduceCoeffs(l,m);|
  2
  !gapprompt@gap>| !gapinput@l;|
  [ 64/49, -24/49, 0, 0 ]
\end{Verbatim}
 }

 

\subsection{\textcolor{Chapter }{ReduceCoeffsMod}}
\logpage{[ 23, 7, 4 ]}\nobreak
\hyperdef{L}{X7F74B1637CB13B7B}{}
{\noindent\textcolor{FuncColor}{$\triangleright$\enspace\texttt{ReduceCoeffsMod({\mdseries\slshape list1[, len1], list2[, len2], modulus})\index{ReduceCoeffsMod@\texttt{ReduceCoeffsMod}}
\label{ReduceCoeffsMod}
}\hfill{\scriptsize (operation)}}\\


 Let $p1$ (and $p2$) be polynomials given by the first \mbox{\texttt{\mdseries\slshape len1}} (\mbox{\texttt{\mdseries\slshape len2}}) entries of the coefficient list \mbox{\texttt{\mdseries\slshape list1}} (\mbox{\texttt{\mdseries\slshape list2}}). If \mbox{\texttt{\mdseries\slshape len1}} and \mbox{\texttt{\mdseries\slshape len2}} are omitted, they default to the lengths of \mbox{\texttt{\mdseries\slshape list1}} and \mbox{\texttt{\mdseries\slshape list2}}. \texttt{ReduceCoeffsMod} changes \mbox{\texttt{\mdseries\slshape list1}} to the coefficient list of the remainder when dividing \mbox{\texttt{\mdseries\slshape p1}} by \mbox{\texttt{\mdseries\slshape p2}} modulo \mbox{\texttt{\mdseries\slshape modulus}}, which must be a positive integer. This operation changes \mbox{\texttt{\mdseries\slshape list1}} which therefore must be a mutable list. The operation returns the position of
the last non\texttt{\symbol{45}}zero entry of the result but is not guaranteed
to remove trailing zeroes. 
\begin{Verbatim}[commandchars=!@|,fontsize=\small,frame=single,label=Example]
  !gapprompt@gap>| !gapinput@l:=[1,2,3,4];;m:=[5,6,7];;ReduceCoeffsMod(l,m,3);|
  1
  !gapprompt@gap>| !gapinput@l;|
  [ 1, 0, 0, 0 ]
\end{Verbatim}
 }

 

\subsection{\textcolor{Chapter }{PowerModCoeffs}}
\logpage{[ 23, 7, 5 ]}\nobreak
\hyperdef{L}{X825F8F357FB1BF56}{}
{\noindent\textcolor{FuncColor}{$\triangleright$\enspace\texttt{PowerModCoeffs({\mdseries\slshape list1[, len1], exp, list2[, len2]})\index{PowerModCoeffs@\texttt{PowerModCoeffs}}
\label{PowerModCoeffs}
}\hfill{\scriptsize (operation)}}\\


 Let $p1$ and $p2$ be polynomials whose coefficients are given by the first \mbox{\texttt{\mdseries\slshape len1}} resp. \mbox{\texttt{\mdseries\slshape len2}} entries of the lists \mbox{\texttt{\mdseries\slshape list1}} and \mbox{\texttt{\mdseries\slshape list2}}, respectively. If \mbox{\texttt{\mdseries\slshape len1}} and \mbox{\texttt{\mdseries\slshape len2}} are omitted, they default to the lengths of \mbox{\texttt{\mdseries\slshape list1}} and \mbox{\texttt{\mdseries\slshape list2}}. Let \mbox{\texttt{\mdseries\slshape exp}} be a positive integer. \texttt{PowerModCoeffs} returns the coefficient list of the remainder when dividing the \mbox{\texttt{\mdseries\slshape exp}}\texttt{\symbol{45}}th power of $p1$ by $p2$. The coefficients are reduced already while powers are computed, therefore
avoiding an explosion in list length. 
\begin{Verbatim}[commandchars=!@|,fontsize=\small,frame=single,label=Example]
  !gapprompt@gap>| !gapinput@l:=[1,2,3,4];;m:=[5,6,7];;PowerModCoeffs(l,5,m);|
  [ -839462813696/678223072849, -7807439437824/678223072849 ]
\end{Verbatim}
 }

 

\subsection{\textcolor{Chapter }{ShiftedCoeffs}}
\logpage{[ 23, 7, 6 ]}\nobreak
\hyperdef{L}{X833EF7AE80CE8B3C}{}
{\noindent\textcolor{FuncColor}{$\triangleright$\enspace\texttt{ShiftedCoeffs({\mdseries\slshape list, shift})\index{ShiftedCoeffs@\texttt{ShiftedCoeffs}}
\label{ShiftedCoeffs}
}\hfill{\scriptsize (operation)}}\\


 produces a new coefficient list \texttt{new} obtained by the rule \texttt{new[i+\mbox{\texttt{\mdseries\slshape shift}}]:= \mbox{\texttt{\mdseries\slshape list}}[i]} and filling initial holes by the appropriate zero. 
\begin{Verbatim}[commandchars=!@|,fontsize=\small,frame=single,label=Example]
  !gapprompt@gap>| !gapinput@l:=[1,2,3];;ShiftedCoeffs(l,2);ShiftedCoeffs(l,-2);|
  [ 0, 0, 1, 2, 3 ]
  [ 3 ]
\end{Verbatim}
 }

 }

 }

             
\chapter{\textcolor{Chapter }{Matrices}}\label{Matrices}
\logpage{[ 24, 0, 0 ]}
\hyperdef{L}{X812CCAB278643A59}{}
{
  In \textsf{GAP}, Matrices can be represented by lists of row vectors, see \ref{Row Vectors}. (For a more general way to represent vectors and matrices, see Chapter \ref{Vector and Matrix Objects}). The row vectors must all have the same length, and their elements must lie
in a common ring. However, since checking rectangularness can be expensive
functions and methods of operations for matrices often will not give an error
message for non\texttt{\symbol{45}}rectangular lists of lists {\textendash}in
such cases the result is undefined. 

 Because matrices are just a special case of lists, all operations and
functions for lists are applicable to matrices also (see chapter \ref{Lists}). This especially includes accessing elements of a matrix (see \ref{List Elements}), changing elements of a matrix (see \ref{List Assignment}), and comparing matrices (see \ref{Comparisons of Lists}). 

 Note that, since a matrix is a list of lists, the behaviour of \texttt{ShallowCopy} (\ref{ShallowCopy}) for matrices is just a special case of \texttt{ShallowCopy} (\ref{ShallowCopy}) for lists (see{\nobreakspace}\ref{Duplication of Lists}); called with an immutable matrix \mbox{\texttt{\mdseries\slshape mat}}, \texttt{ShallowCopy} (\ref{ShallowCopy}) returns a mutable matrix whose rows are identical to the rows of \mbox{\texttt{\mdseries\slshape mat}}. In particular the rows are still immutable. To get a matrix whose rows are
mutable, one can use \texttt{List( \mbox{\texttt{\mdseries\slshape mat}}, ShallowCopy )}.  
\section{\textcolor{Chapter }{InfoMatrix (Info Class)}}\label{sect:InfoMatrix}
\logpage{[ 24, 1, 0 ]}
\hyperdef{L}{X801E1B5D7EC8DDD3}{}
{
  

\subsection{\textcolor{Chapter }{InfoMatrix}}
\logpage{[ 24, 1, 1 ]}\nobreak
\hyperdef{L}{X78EC82D27B4191DA}{}
{\noindent\textcolor{FuncColor}{$\triangleright$\enspace\texttt{InfoMatrix\index{InfoMatrix@\texttt{InfoMatrix}}
\label{InfoMatrix}
}\hfill{\scriptsize (info class)}}\\


 The info class for matrix operations is \texttt{InfoMatrix}. }

 }

  
\section{\textcolor{Chapter }{Categories of Matrices}}\label{Categories of Matrices}
\logpage{[ 24, 2, 0 ]}
\hyperdef{L}{X866E55A58164FAED}{}
{
  

\subsection{\textcolor{Chapter }{IsMatrix}}
\logpage{[ 24, 2, 1 ]}\nobreak
\hyperdef{L}{X7E1AE46B862B185F}{}
{\noindent\textcolor{FuncColor}{$\triangleright$\enspace\texttt{IsMatrix({\mdseries\slshape obj})\index{IsMatrix@\texttt{IsMatrix}}
\label{IsMatrix}
}\hfill{\scriptsize (Category)}}\\


 By convention \emph{matrix} is a list of lists of equal length whose entries lie in a common ring. 

 For technical reasons laid out at the top of Chapter \ref{Matrices}, the filter \texttt{IsMatrix} is a synonym for a table of ring elements, (see \texttt{IsTable} (\ref{IsTable}) and \texttt{IsRingElement} (\ref{IsRingElement})). This means that \texttt{IsMatrix} returns \texttt{true} for tables such as \texttt{[[1,2],[3]]}. If necessary, \texttt{IsRectangularTable} (\ref{IsRectangularTable}) can be used to test whether an object is a list of homogeneous lists of equal
lengths manually. 

 Note that matrices may have different multiplications, besides the usual
matrix product there is for example the Lie product. So there are categories
such as \texttt{IsOrdinaryMatrix} (\ref{IsOrdinaryMatrix}) and \texttt{IsLieMatrix} (\ref{IsLieMatrix}) that describe the matrix multiplication. One can form the product of two
matrices only if they support the same multiplication. 

 
\begin{Verbatim}[commandchars=!@|,fontsize=\small,frame=single,label=Example]
  !gapprompt@gap>| !gapinput@mat:=[[1,2,3],[4,5,6],[7,8,9]];|
  [ [ 1, 2, 3 ], [ 4, 5, 6 ], [ 7, 8, 9 ] ]
  !gapprompt@gap>| !gapinput@IsMatrix(mat);|
  true
  !gapprompt@gap>| !gapinput@mat:=[[1,2],[3]];|
  [ [ 1, 2 ], [ 3 ] ]
  !gapprompt@gap>| !gapinput@IsMatrix(mat);|
  true
  !gapprompt@gap>| !gapinput@IsRectangularTable(mat);|
  false
\end{Verbatim}
 

 Note that the empty list \texttt{[]} and more complex ``empty'' structures such as \texttt{[[]]} are \emph{not} matrices, although special methods allow them be used in place of matrices in
some situations. See \texttt{EmptyMatrix} (\ref{EmptyMatrix}) below. 

 
\begin{Verbatim}[commandchars=!@|,fontsize=\small,frame=single,label=Example]
  !gapprompt@gap>| !gapinput@[[0]]*[[]];|
  [ [  ] ]
  !gapprompt@gap>| !gapinput@IsMatrix([[]]);|
  false
\end{Verbatim}
 }

 

\subsection{\textcolor{Chapter }{IsOrdinaryMatrix}}
\logpage{[ 24, 2, 2 ]}\nobreak
\hyperdef{L}{X7CF42B8A845BC6A9}{}
{\noindent\textcolor{FuncColor}{$\triangleright$\enspace\texttt{IsOrdinaryMatrix({\mdseries\slshape obj})\index{IsOrdinaryMatrix@\texttt{IsOrdinaryMatrix}}
\label{IsOrdinaryMatrix}
}\hfill{\scriptsize (Category)}}\\


 An \emph{ordinary matrix} is a matrix whose multiplication is the ordinary matrix multiplication. 

 Each matrix in internal representation is in the category \texttt{IsOrdinaryMatrix}, and arithmetic operations with objects in \texttt{IsOrdinaryMatrix} produce again matrices in \texttt{IsOrdinaryMatrix}. 

 Note that we want that Lie matrices shall be matrices that behave in the same
way as ordinary matrices, except that they have a different multiplication. So
we must distinguish the different matrix multiplications, in order to be able
to describe the applicability of multiplication, and also in order to form a
matrix of the appropriate type as the sum, difference etc.{\nobreakspace}of
two matrices which have the same multiplication. }

 

\subsection{\textcolor{Chapter }{IsLieMatrix}}
\logpage{[ 24, 2, 3 ]}\nobreak
\hyperdef{L}{X86EC33E17DD12D0E}{}
{\noindent\textcolor{FuncColor}{$\triangleright$\enspace\texttt{IsLieMatrix({\mdseries\slshape mat})\index{IsLieMatrix@\texttt{IsLieMatrix}}
\label{IsLieMatrix}
}\hfill{\scriptsize (Category)}}\\


 A \emph{Lie matrix} is a matrix whose multiplication is given by the Lie bracket. (Note that a
matrix with ordinary matrix multiplication is in the category \texttt{IsOrdinaryMatrix} (\ref{IsOrdinaryMatrix}).) 

 Each matrix created by \texttt{LieObject} (\ref{LieObject}) is in the category \texttt{IsLieMatrix}, and arithmetic operations with objects in \texttt{IsLieMatrix} produce again matrices in \texttt{IsLieMatrix}.  }

 }

  
\section{\textcolor{Chapter }{Operators for Matrices}}\label{Operators for Matrices}
\logpage{[ 24, 3, 0 ]}
\hyperdef{L}{X7899335779A39A95}{}
{
  The rules for arithmetic operations involving matrices are in fact special
cases of those for the arithmetic of lists, given in Section{\nobreakspace}\ref{Arithmetic for Lists} and the following sections, here we reiterate that definition, in the language
of vectors and matrices. 

 Note that the additive behaviour sketched below is defined only for lists in
the category \texttt{IsGeneralizedRowVector} (\ref{IsGeneralizedRowVector}), and the multiplicative behaviour is defined only for lists in the category \texttt{IsMultiplicativeGeneralizedRowVector} (\ref{IsMultiplicativeGeneralizedRowVector}) (see{\nobreakspace}\ref{Filters Controlling the Arithmetic Behaviour of Lists}). 

 \index{addition!matrices} \texttt{\mbox{\texttt{\mdseries\slshape mat1}} + \mbox{\texttt{\mdseries\slshape mat2}}} 

 returns the sum of the two matrices \mbox{\texttt{\mdseries\slshape mat1}} and \mbox{\texttt{\mdseries\slshape mat2}}, Probably the most usual situation is that \mbox{\texttt{\mdseries\slshape mat1}} and \mbox{\texttt{\mdseries\slshape mat2}} have the same dimensions and are defined over a common field; in this case the
sum is a new matrix over the same field where each entry is the sum of the
corresponding entries of the matrices. 

 In more general situations, the sum of two matrices need not be a matrix, for
example adding an integer matrix \mbox{\texttt{\mdseries\slshape mat1}} and a matrix \mbox{\texttt{\mdseries\slshape mat2}} over a finite field yields the table of pointwise sums, which will be a
mixture of finite field elements and integers if \mbox{\texttt{\mdseries\slshape mat1}} has bigger dimensions than \mbox{\texttt{\mdseries\slshape mat2}}. 

 \index{addition!scalar and matrix} \texttt{\mbox{\texttt{\mdseries\slshape scalar}} + \mbox{\texttt{\mdseries\slshape mat}}} 

 \index{addition!matrix and scalar} \texttt{\mbox{\texttt{\mdseries\slshape mat}} + \mbox{\texttt{\mdseries\slshape scalar}}} 

 returns the sum of the scalar \mbox{\texttt{\mdseries\slshape scalar}} and the matrix \mbox{\texttt{\mdseries\slshape mat}}. Probably the most usual situation is that the entries of \mbox{\texttt{\mdseries\slshape mat}} lie in a common field with \mbox{\texttt{\mdseries\slshape scalar}}; in this case the sum is a new matrix over the same field where each entry is
the sum of the scalar and the corresponding entry of the matrix. 

 More general situations are for example the sum of an integer scalar and a
matrix over a finite field, or the sum of a finite field element and an
integer matrix. 

 \index{subtraction!matrices} \texttt{\mbox{\texttt{\mdseries\slshape mat1}} \texttt{\symbol{45}} \mbox{\texttt{\mdseries\slshape mat2}}} 

 \index{subtraction!scalar and matrix} \texttt{\mbox{\texttt{\mdseries\slshape scalar}} \texttt{\symbol{45}} \mbox{\texttt{\mdseries\slshape mat}}} 

 \index{subtraction!matrix and scalar} \texttt{\mbox{\texttt{\mdseries\slshape mat}} \texttt{\symbol{45}} \mbox{\texttt{\mdseries\slshape scalar}}} 

 Subtracting a matrix or scalar is defined as adding its additive inverse, so
the statements for the addition hold likewise. 

 \index{multiplication!scalar and matrix} \texttt{\mbox{\texttt{\mdseries\slshape scalar}} * \mbox{\texttt{\mdseries\slshape mat}}} 

 \index{multiplication!matrix and scalar} \texttt{\mbox{\texttt{\mdseries\slshape mat}} * \mbox{\texttt{\mdseries\slshape scalar}}} 

 returns the product of the scalar \mbox{\texttt{\mdseries\slshape scalar}} and the matrix \mbox{\texttt{\mdseries\slshape mat}}. Probably the most usual situation is that the elements of \mbox{\texttt{\mdseries\slshape mat}} lie in a common field with \mbox{\texttt{\mdseries\slshape scalar}}; in this case the product is a new matrix over the same field where each
entry is the product of the scalar and the corresponding entry of the matrix. 

 More general situations are for example the product of an integer scalar and a
matrix over a finite field, or the product of a finite field element and an
integer matrix. 

 \index{multiplication!vector and matrix} \texttt{\mbox{\texttt{\mdseries\slshape vec}} * \mbox{\texttt{\mdseries\slshape mat}}} 

 returns the product of the row vector \mbox{\texttt{\mdseries\slshape vec}} and the matrix \mbox{\texttt{\mdseries\slshape mat}}. Probably the most usual situation is that \mbox{\texttt{\mdseries\slshape vec}} and \mbox{\texttt{\mdseries\slshape mat}} have the same lengths and are defined over a common field, and that all rows
of \mbox{\texttt{\mdseries\slshape mat}} have some common length $m$; in this case the product is a new row vector of length $m$ over the same field which is the sum of the scalar multiples of the rows of \mbox{\texttt{\mdseries\slshape mat}} with the corresponding entries of \mbox{\texttt{\mdseries\slshape vec}}. 

 More general situations are for example the product of an integer vector and a
matrix over a finite field, or the product of a vector over a finite field and
an integer matrix. 

 \index{multiplication!matrix and vector} \texttt{\mbox{\texttt{\mdseries\slshape mat}} * \mbox{\texttt{\mdseries\slshape vec}}} 

 returns the product of the matrix \mbox{\texttt{\mdseries\slshape mat}} and the row vector \mbox{\texttt{\mdseries\slshape vec}}. (This is the standard product of a matrix with a \emph{column} vector.) Probably the most usual situation is that the length of \mbox{\texttt{\mdseries\slshape vec}} and of all rows of \mbox{\texttt{\mdseries\slshape mat}} are equal, and that the elements of \mbox{\texttt{\mdseries\slshape mat}} and \mbox{\texttt{\mdseries\slshape vec}} lie in a common field; in this case the product is a new row vector of the
same length as \mbox{\texttt{\mdseries\slshape mat}} and over the same field which is the sum of the scalar multiples of the
columns of \mbox{\texttt{\mdseries\slshape mat}} with the corresponding entries of \mbox{\texttt{\mdseries\slshape vec}}. 

 More general situations are for example the product of an integer matrix and a
vector over a finite field, or the product of a matrix over a finite field and
an integer vector. 

 \index{multiplication!matrices} \texttt{\mbox{\texttt{\mdseries\slshape mat1}} * \mbox{\texttt{\mdseries\slshape mat2}}} 

 This form evaluates to the (Cauchy) product of the two matrices \mbox{\texttt{\mdseries\slshape mat1}} and \mbox{\texttt{\mdseries\slshape mat2}}. Probably the most usual situation is that the number of columns of \mbox{\texttt{\mdseries\slshape mat1}} equals the number of rows of \mbox{\texttt{\mdseries\slshape mat2}}, and that the elements of \mbox{\texttt{\mdseries\slshape mat}} and \mbox{\texttt{\mdseries\slshape vec}} lie in a common field; if \mbox{\texttt{\mdseries\slshape mat1}} is a matrix with $m$ rows and $n$ columns and \mbox{\texttt{\mdseries\slshape mat2}} is a matrix with $n$ rows and $o$ columns, the result is a new matrix with $m$ rows and $o$ columns. The element in row $i$ at position $j$ of the product is the sum of $\mbox{\texttt{\mdseries\slshape mat1}}[i][l] * \mbox{\texttt{\mdseries\slshape mat2}}[l][j]$, with $l$ running from $1$ to $n$. 

 \index{inverse!matrix} \texttt{Inverse( \mbox{\texttt{\mdseries\slshape mat}} )} 

 returns the inverse of the matrix \mbox{\texttt{\mdseries\slshape mat}}, which must be an invertible square matrix. If \mbox{\texttt{\mdseries\slshape mat}} is not invertible then \texttt{fail} is returned. 

 \index{quotient!matrices} \texttt{\mbox{\texttt{\mdseries\slshape mat1}} / \mbox{\texttt{\mdseries\slshape mat2}}} 

 \index{quotient!scalar and matrix} \texttt{\mbox{\texttt{\mdseries\slshape scalar}} / \mbox{\texttt{\mdseries\slshape mat}}} 

 \index{quotient!matrix and scalar} \texttt{\mbox{\texttt{\mdseries\slshape mat}} / \mbox{\texttt{\mdseries\slshape scalar}}} 

 \index{quotient!vector and matrix} \texttt{\mbox{\texttt{\mdseries\slshape vec}} / \mbox{\texttt{\mdseries\slshape mat}}} 

 In general, \texttt{\mbox{\texttt{\mdseries\slshape left}} / \mbox{\texttt{\mdseries\slshape right}}} is defined as \texttt{\mbox{\texttt{\mdseries\slshape left}} * \mbox{\texttt{\mdseries\slshape right}}\texttt{\symbol{94}}\texttt{\symbol{45}}1}. Thus in the above forms the right operand must always be invertible. 

 \index{power!matrix} \texttt{\mbox{\texttt{\mdseries\slshape mat}} \texttt{\symbol{94}} \mbox{\texttt{\mdseries\slshape int}}} 

 \index{conjugate!matrix} \texttt{\mbox{\texttt{\mdseries\slshape mat1}} \texttt{\symbol{94}} \mbox{\texttt{\mdseries\slshape mat2}}} 

 \index{image!vector under matrix} \texttt{\mbox{\texttt{\mdseries\slshape vec}} \texttt{\symbol{94}} \mbox{\texttt{\mdseries\slshape mat}}} 

 Powering a square matrix \mbox{\texttt{\mdseries\slshape mat}} by an integer \mbox{\texttt{\mdseries\slshape int}} yields the \mbox{\texttt{\mdseries\slshape int}}\texttt{\symbol{45}}th power of \mbox{\texttt{\mdseries\slshape mat}}; if \mbox{\texttt{\mdseries\slshape int}} is negative then \mbox{\texttt{\mdseries\slshape mat}} must be invertible, if \mbox{\texttt{\mdseries\slshape int}} is \texttt{0} then the result is the identity matrix \texttt{One( \mbox{\texttt{\mdseries\slshape mat}} )}, even if \mbox{\texttt{\mdseries\slshape mat}} is not invertible. 

 Powering a square matrix \mbox{\texttt{\mdseries\slshape mat1}} by an invertible square matrix \mbox{\texttt{\mdseries\slshape mat2}} of the same dimensions yields the conjugate of \mbox{\texttt{\mdseries\slshape mat1}} by \mbox{\texttt{\mdseries\slshape mat2}}, i.e., the matrix \texttt{\mbox{\texttt{\mdseries\slshape mat2}}\texttt{\symbol{94}}\texttt{\symbol{45}}1 * \mbox{\texttt{\mdseries\slshape mat1}} * \mbox{\texttt{\mdseries\slshape mat2}}}. 

 Powering a row vector \mbox{\texttt{\mdseries\slshape vec}} by a matrix \mbox{\texttt{\mdseries\slshape mat}} is in every respect equivalent to \texttt{\mbox{\texttt{\mdseries\slshape vec}} * \mbox{\texttt{\mdseries\slshape mat}}}. This operations reflects the fact that matrices act naturally on row vectors
by multiplication from the right, and that the powering operator is \textsf{GAP}'s standard for group actions. 

 \index{matrices!commutator} \texttt{Comm( \mbox{\texttt{\mdseries\slshape mat1}}, \mbox{\texttt{\mdseries\slshape mat2}} )} 

 returns the commutator of the square invertible matrices \mbox{\texttt{\mdseries\slshape mat1}} and \mbox{\texttt{\mdseries\slshape mat2}} of the same dimensions and over a common field, which is the matrix \texttt{\mbox{\texttt{\mdseries\slshape mat1}}\texttt{\symbol{94}}\texttt{\symbol{45}}1 * \mbox{\texttt{\mdseries\slshape mat2}}\texttt{\symbol{94}}\texttt{\symbol{45}}1 * \mbox{\texttt{\mdseries\slshape mat1}} * \mbox{\texttt{\mdseries\slshape mat2}}}. 

 The following cases are still special cases of the general list arithmetic
defined in{\nobreakspace}\ref{Arithmetic for Lists}. 

 \index{addition!scalar and matrix list} \texttt{\mbox{\texttt{\mdseries\slshape scalar}} + \mbox{\texttt{\mdseries\slshape matlist}}} 

 \index{addition!scalar and matrix list} \texttt{\mbox{\texttt{\mdseries\slshape matlist}} + \mbox{\texttt{\mdseries\slshape scalar}}} 

 \index{subtraction!scalar and matrix list} \texttt{\mbox{\texttt{\mdseries\slshape scalar}} \texttt{\symbol{45}} \mbox{\texttt{\mdseries\slshape matlist}}} 

 \index{subtraction!scalar and matrix list} \texttt{\mbox{\texttt{\mdseries\slshape matlist}} \texttt{\symbol{45}} \mbox{\texttt{\mdseries\slshape scalar}}} 

 \index{multiplication!scalar and matrix list} \texttt{\mbox{\texttt{\mdseries\slshape scalar}} * \mbox{\texttt{\mdseries\slshape matlist}}} 

 \index{multiplication!scalar and matrix list} \texttt{\mbox{\texttt{\mdseries\slshape matlist}} * \mbox{\texttt{\mdseries\slshape scalar}}} 

 \index{quotient!scalar and matrix list} \texttt{\mbox{\texttt{\mdseries\slshape matlist}} / \mbox{\texttt{\mdseries\slshape scalar}}} 

 A scalar \mbox{\texttt{\mdseries\slshape scalar}} may also be added, subtracted, multiplied with, or divided into a list \mbox{\texttt{\mdseries\slshape matlist}} of matrices. The result is a new list of matrices where each matrix is the
result of performing the operation with the corresponding matrix in \mbox{\texttt{\mdseries\slshape matlist}}. 

 \index{multiplication!matrix and matrix list} \texttt{\mbox{\texttt{\mdseries\slshape mat}} * \mbox{\texttt{\mdseries\slshape matlist}}} 

 \index{multiplication!matrix and matrix list} \texttt{\mbox{\texttt{\mdseries\slshape matlist}} * \mbox{\texttt{\mdseries\slshape mat}}} 

 A matrix \mbox{\texttt{\mdseries\slshape mat}} may also be multiplied with a list \mbox{\texttt{\mdseries\slshape matlist}} of matrices. The result is a new list of matrices, where each entry is the
product of \mbox{\texttt{\mdseries\slshape mat}} and the corresponding entry in \mbox{\texttt{\mdseries\slshape matlist}}. 

 \index{quotient!matrix and matrix list} \texttt{\mbox{\texttt{\mdseries\slshape matlist}} / \mbox{\texttt{\mdseries\slshape mat}}} 

 Dividing a list \mbox{\texttt{\mdseries\slshape matlist}} of matrices by an invertible matrix \mbox{\texttt{\mdseries\slshape mat}} evaluates to \texttt{\mbox{\texttt{\mdseries\slshape matlist}} * \mbox{\texttt{\mdseries\slshape mat}}\texttt{\symbol{94}}\texttt{\symbol{45}}1}. 

 \index{multiplication!vector and matrix list} \texttt{\mbox{\texttt{\mdseries\slshape vec}} * \mbox{\texttt{\mdseries\slshape matlist}}} 

 returns the product of the vector \mbox{\texttt{\mdseries\slshape vec}} and the list of matrices \mbox{\texttt{\mdseries\slshape mat}}. The lengths \mbox{\texttt{\mdseries\slshape l}} of \mbox{\texttt{\mdseries\slshape vec}} and \mbox{\texttt{\mdseries\slshape matlist}} must be equal. All matrices in \mbox{\texttt{\mdseries\slshape matlist}} must have the same dimensions. The elements of \mbox{\texttt{\mdseries\slshape vec}} and the elements of the matrices in \mbox{\texttt{\mdseries\slshape matlist}} must lie in a common ring. The product is the sum over \texttt{\mbox{\texttt{\mdseries\slshape vec}}[\mbox{\texttt{\mdseries\slshape i}}] * \mbox{\texttt{\mdseries\slshape matlist}}[\mbox{\texttt{\mdseries\slshape i}}]} with \mbox{\texttt{\mdseries\slshape i}} running from 1 to \mbox{\texttt{\mdseries\slshape l}}. 

 For the mutability of results of arithmetic operations, see{\nobreakspace}\ref{Mutability and Copyability}. }

  
\section{\textcolor{Chapter }{Properties and Attributes of Matrices}}\label{Properties and Attributes of Matrices}
\logpage{[ 24, 4, 0 ]}
\hyperdef{L}{X7F5AD28E869B66CB}{}
{
  

\subsection{\textcolor{Chapter }{DimensionsMat}}
\logpage{[ 24, 4, 1 ]}\nobreak
\hyperdef{L}{X83A9DC2085D3A972}{}
{\noindent\textcolor{FuncColor}{$\triangleright$\enspace\texttt{DimensionsMat({\mdseries\slshape mat})\index{DimensionsMat@\texttt{DimensionsMat}}
\label{DimensionsMat}
}\hfill{\scriptsize (attribute)}}\\


 is a list of length 2, the first being the number of rows, the second being
the number of columns of the matrix \mbox{\texttt{\mdseries\slshape mat}}. If \mbox{\texttt{\mdseries\slshape mat}} is malformed, that is, it is not a \texttt{IsRectangularTable} (\ref{IsRectangularTable}), returns \texttt{fail}. 
\begin{Verbatim}[commandchars=!@|,fontsize=\small,frame=single,label=Example]
  !gapprompt@gap>| !gapinput@DimensionsMat([[1,2,3],[4,5,6]]);|
  [ 2, 3 ]
  !gapprompt@gap>| !gapinput@DimensionsMat([[1,2,3],[4,5]]);|
  fail
\end{Verbatim}
 }

 

\subsection{\textcolor{Chapter }{DefaultFieldOfMatrix}}
\logpage{[ 24, 4, 2 ]}\nobreak
\hyperdef{L}{X80AE547B8095A5CB}{}
{\noindent\textcolor{FuncColor}{$\triangleright$\enspace\texttt{DefaultFieldOfMatrix({\mdseries\slshape mat})\index{DefaultFieldOfMatrix@\texttt{DefaultFieldOfMatrix}}
\label{DefaultFieldOfMatrix}
}\hfill{\scriptsize (attribute)}}\\


 For a matrix \mbox{\texttt{\mdseries\slshape mat}}, \texttt{DefaultFieldOfMatrix} returns either a field (not necessarily the smallest one) containing all
entries of \mbox{\texttt{\mdseries\slshape mat}}, or \texttt{fail}. 

 If \mbox{\texttt{\mdseries\slshape mat}} is a matrix of finite field elements or a matrix of cyclotomics, \texttt{DefaultFieldOfMatrix} returns the default field generated by the matrix entries (see{\nobreakspace}\ref{Creating Finite Fields} and \ref{Operations for Cyclotomics}). 
\begin{Verbatim}[commandchars=!@|,fontsize=\small,frame=single,label=Example]
  !gapprompt@gap>| !gapinput@DefaultFieldOfMatrix([[Z(4),Z(8)]]);|
  GF(2^6)
\end{Verbatim}
 }

 

\subsection{\textcolor{Chapter }{TraceMatrix}}
\logpage{[ 24, 4, 3 ]}\nobreak
\hyperdef{L}{X784EC2777C06AFE4}{}
{\noindent\textcolor{FuncColor}{$\triangleright$\enspace\texttt{TraceMatrix({\mdseries\slshape mat})\index{TraceMatrix@\texttt{TraceMatrix}}
\label{TraceMatrix}
}\hfill{\scriptsize (attribute)}}\\
\noindent\textcolor{FuncColor}{$\triangleright$\enspace\texttt{TraceMat({\mdseries\slshape mat})\index{TraceMat@\texttt{TraceMat}}
\label{TraceMat}
}\hfill{\scriptsize (attribute)}}\\
\noindent\textcolor{FuncColor}{$\triangleright$\enspace\texttt{Trace({\mdseries\slshape mat})\index{Trace@\texttt{Trace}!of a matrix}
\label{Trace:of a matrix}
}\hfill{\scriptsize (attribute)}}\\


 The trace of a square matrix is the sum of its diagonal entries. 
\begin{Verbatim}[commandchars=!@|,fontsize=\small,frame=single,label=Example]
  !gapprompt@gap>| !gapinput@TraceMatrix([[1,2,3],[4,5,6],[7,8,9]]);|
  15
\end{Verbatim}
 }

 

\subsection{\textcolor{Chapter }{DeterminantMatrix}}
\logpage{[ 24, 4, 4 ]}\nobreak
\hyperdef{L}{X8488D69A7ADDB4E2}{}
{\noindent\textcolor{FuncColor}{$\triangleright$\enspace\texttt{DeterminantMatrix({\mdseries\slshape mat})\index{DeterminantMatrix@\texttt{DeterminantMatrix}}
\label{DeterminantMatrix}
}\hfill{\scriptsize (attribute)}}\\
\noindent\textcolor{FuncColor}{$\triangleright$\enspace\texttt{DeterminantMat({\mdseries\slshape mat})\index{DeterminantMat@\texttt{DeterminantMat}}
\label{DeterminantMat}
}\hfill{\scriptsize (attribute)}}\\
\noindent\textcolor{FuncColor}{$\triangleright$\enspace\texttt{Determinant({\mdseries\slshape mat})\index{Determinant@\texttt{Determinant}}
\label{Determinant}
}\hfill{\scriptsize (operation)}}\\


 returns the determinant of the square matrix \mbox{\texttt{\mdseries\slshape mat}}. 

 These methods assume implicitly that \mbox{\texttt{\mdseries\slshape mat}} is defined over an integral domain whose quotient field is implemented in \textsf{GAP}. For matrices defined over an arbitrary commutative ring with one
see{\nobreakspace}\texttt{DeterminantMatDivFree} (\ref{DeterminantMatDivFree}). }

 

\subsection{\textcolor{Chapter }{DeterminantMatrixDestructive}}
\logpage{[ 24, 4, 5 ]}\nobreak
\hyperdef{L}{X824B5DC2875118B3}{}
{\noindent\textcolor{FuncColor}{$\triangleright$\enspace\texttt{DeterminantMatrixDestructive({\mdseries\slshape mat})\index{DeterminantMatrixDestructive@\texttt{DeterminantMatrixDestructive}}
\label{DeterminantMatrixDestructive}
}\hfill{\scriptsize (operation)}}\\
\noindent\textcolor{FuncColor}{$\triangleright$\enspace\texttt{DeterminantMatDestructive({\mdseries\slshape mat})\index{DeterminantMatDestructive@\texttt{DeterminantMatDestructive}}
\label{DeterminantMatDestructive}
}\hfill{\scriptsize (operation)}}\\


 Does the same as \texttt{DeterminantMatrix} (\ref{DeterminantMatrix}), with the difference that it may destroy its argument. The matrix \mbox{\texttt{\mdseries\slshape mat}} must be mutable. 
\begin{Verbatim}[commandchars=!@|,fontsize=\small,frame=single,label=Example]
  !gapprompt@gap>| !gapinput@DeterminantMatrix([[1,2],[2,1]]);|
  -3
  !gapprompt@gap>| !gapinput@mm:= [[1,2],[2,1]];;|
  !gapprompt@gap>| !gapinput@DeterminantMatrixDestructive( mm );|
  -3
  !gapprompt@gap>| !gapinput@mm;|
  [ [ 1, 2 ], [ 0, -3 ] ]
\end{Verbatim}
 }

 

\subsection{\textcolor{Chapter }{DeterminantMatrixDivFree}}
\logpage{[ 24, 4, 6 ]}\nobreak
\hyperdef{L}{X80693FAB7D541804}{}
{\noindent\textcolor{FuncColor}{$\triangleright$\enspace\texttt{DeterminantMatrixDivFree({\mdseries\slshape mat})\index{DeterminantMatrixDivFree@\texttt{DeterminantMatrixDivFree}}
\label{DeterminantMatrixDivFree}
}\hfill{\scriptsize (operation)}}\\
\noindent\textcolor{FuncColor}{$\triangleright$\enspace\texttt{DeterminantMatDivFree({\mdseries\slshape mat})\index{DeterminantMatDivFree@\texttt{DeterminantMatDivFree}}
\label{DeterminantMatDivFree}
}\hfill{\scriptsize (operation)}}\\


 return the determinant of a square matrix \mbox{\texttt{\mdseries\slshape mat}} over an arbitrary commutative ring with one using the division free method of
Mahajan and Vinay \cite{MV97}. }

 

\subsection{\textcolor{Chapter }{IsEmptyMatrix (for a matrix object)}}
\logpage{[ 24, 4, 7 ]}\nobreak
\hyperdef{L}{X7F8D25897EC1630B}{}
{\noindent\textcolor{FuncColor}{$\triangleright$\enspace\texttt{IsEmptyMatrix({\mdseries\slshape M})\index{IsEmptyMatrix@\texttt{IsEmptyMatrix}!for a matrix object}
\label{IsEmptyMatrix:for a matrix object}
}\hfill{\scriptsize (property)}}\\
\textbf{\indent Returns:\ }
A boolean



 Is \texttt{true} if the matrix object \mbox{\texttt{\mdseries\slshape M}} either has zero columns or zero rows, and \texttt{false} otherwise. In other words, a matrix object is empty if it has no entries. }

 

\subsection{\textcolor{Chapter }{IsMonomialMatrix}}
\logpage{[ 24, 4, 8 ]}\nobreak
\hyperdef{L}{X848B80437CE65FF3}{}
{\noindent\textcolor{FuncColor}{$\triangleright$\enspace\texttt{IsMonomialMatrix({\mdseries\slshape mat})\index{IsMonomialMatrix@\texttt{IsMonomialMatrix}}
\label{IsMonomialMatrix}
}\hfill{\scriptsize (property)}}\\


 A matrix is monomial if and only if it has exactly one nonzero entry in every
row and every column. 
\begin{Verbatim}[commandchars=!@|,fontsize=\small,frame=single,label=Example]
  !gapprompt@gap>| !gapinput@IsMonomialMatrix([[0,1],[1,0]]);|
  true
\end{Verbatim}
 }

 

\subsection{\textcolor{Chapter }{IsDiagonalMatrix}}
\logpage{[ 24, 4, 9 ]}\nobreak
\hyperdef{L}{X7EEC8E768178696E}{}
{\noindent\textcolor{FuncColor}{$\triangleright$\enspace\texttt{IsDiagonalMatrix({\mdseries\slshape mat})\index{IsDiagonalMatrix@\texttt{IsDiagonalMatrix}}
\label{IsDiagonalMatrix}
}\hfill{\scriptsize (property)}}\\
\noindent\textcolor{FuncColor}{$\triangleright$\enspace\texttt{IsDiagonalMat({\mdseries\slshape mat})\index{IsDiagonalMat@\texttt{IsDiagonalMat}}
\label{IsDiagonalMat}
}\hfill{\scriptsize (property)}}\\


 return \texttt{true} if the matrix \mbox{\texttt{\mdseries\slshape mat}} has only zero entries off the main diagonal, and \texttt{false} otherwise. 
\begin{Verbatim}[commandchars=!@|,fontsize=\small,frame=single,label=Example]
  !gapprompt@gap>| !gapinput@IsDiagonalMatrix( [ [ 1 ] ] );|
  true
  !gapprompt@gap>| !gapinput@IsDiagonalMatrix( [ [ 1, 0, 0 ], [ 0, 1, 0 ] ] );|
  true
  !gapprompt@gap>| !gapinput@IsDiagonalMatrix( [ [ 0, 1 ], [ 1, 0 ] ] );|
  false
\end{Verbatim}
 }

 

\subsection{\textcolor{Chapter }{IsUpperTriangularMatrix}}
\logpage{[ 24, 4, 10 ]}\nobreak
\hyperdef{L}{X8740E71C799C0BCC}{}
{\noindent\textcolor{FuncColor}{$\triangleright$\enspace\texttt{IsUpperTriangularMatrix({\mdseries\slshape mat})\index{IsUpperTriangularMatrix@\texttt{IsUpperTriangularMatrix}}
\label{IsUpperTriangularMatrix}
}\hfill{\scriptsize (property)}}\\
\noindent\textcolor{FuncColor}{$\triangleright$\enspace\texttt{IsUpperTriangularMat({\mdseries\slshape mat})\index{IsUpperTriangularMat@\texttt{IsUpperTriangularMat}}
\label{IsUpperTriangularMat}
}\hfill{\scriptsize (property)}}\\


 return \texttt{true} if the matrix \mbox{\texttt{\mdseries\slshape mat}} has only zero entries below the main diagonal, and \texttt{false} otherwise. 
\begin{Verbatim}[commandchars=!@|,fontsize=\small,frame=single,label=Example]
  !gapprompt@gap>| !gapinput@IsUpperTriangularMatrix( [ [ 1 ] ] );|
  true
  !gapprompt@gap>| !gapinput@IsUpperTriangularMatrix( [ [ 1, 2, 3 ], [ 0, 5, 6 ] ] );|
  true
  !gapprompt@gap>| !gapinput@IsUpperTriangularMatrix( [ [ 0, 1 ], [ 1, 0 ] ] );|
  false
\end{Verbatim}
 }

 

\subsection{\textcolor{Chapter }{IsLowerTriangularMatrix}}
\logpage{[ 24, 4, 11 ]}\nobreak
\hyperdef{L}{X853A5B988306DBFE}{}
{\noindent\textcolor{FuncColor}{$\triangleright$\enspace\texttt{IsLowerTriangularMatrix({\mdseries\slshape mat})\index{IsLowerTriangularMatrix@\texttt{IsLowerTriangularMatrix}}
\label{IsLowerTriangularMatrix}
}\hfill{\scriptsize (property)}}\\
\noindent\textcolor{FuncColor}{$\triangleright$\enspace\texttt{IsLowerTriangularMat({\mdseries\slshape mat})\index{IsLowerTriangularMat@\texttt{IsLowerTriangularMat}}
\label{IsLowerTriangularMat}
}\hfill{\scriptsize (property)}}\\


 return \texttt{true} if the matrix \mbox{\texttt{\mdseries\slshape mat}} has only zero entries above the main diagonal, and \texttt{false} otherwise. 
\begin{Verbatim}[commandchars=!@|,fontsize=\small,frame=single,label=Example]
  !gapprompt@gap>| !gapinput@IsLowerTriangularMatrix( [ [ 1 ] ] );|
  true
  !gapprompt@gap>| !gapinput@IsLowerTriangularMatrix( [ [ 1, 0, 0 ], [ 2, 3, 0 ] ] );|
  true
  !gapprompt@gap>| !gapinput@IsLowerTriangularMatrix( [ [ 0, 1 ], [ 1, 0 ] ] );|
  false
\end{Verbatim}
 }

 

\subsection{\textcolor{Chapter }{IsSquareMatrix}}
\logpage{[ 24, 4, 12 ]}\nobreak
\hyperdef{L}{X7FD5345487011416}{}
{\noindent\textcolor{FuncColor}{$\triangleright$\enspace\texttt{IsSquareMatrix({\mdseries\slshape mat})\index{IsSquareMatrix@\texttt{IsSquareMatrix}}
\label{IsSquareMatrix}
}\hfill{\scriptsize (property)}}\\
\noindent\textcolor{FuncColor}{$\triangleright$\enspace\texttt{IsSquareMat({\mdseries\slshape mat})\index{IsSquareMat@\texttt{IsSquareMat}}
\label{IsSquareMat}
}\hfill{\scriptsize (property)}}\\


 return \texttt{true} if the matrix \mbox{\texttt{\mdseries\slshape mat}} has the same number of rows and columns, and \texttt{false} otherwise. 
\begin{Verbatim}[commandchars=!@|,fontsize=\small,frame=single,label=Example]
  !gapprompt@gap>| !gapinput@IsSquareMatrix( [ [ 1 ] ] );|
  true
  !gapprompt@gap>| !gapinput@IsSquareMatrix( [ [ 1, 2 ], [ 3, 4 ] ] );|
  true
  !gapprompt@gap>| !gapinput@IsSquareMatrix( [ [ 1, 2, 3 ], [ 4, 5, 6 ] ] );|
  false
\end{Verbatim}
 }

 

\subsection{\textcolor{Chapter }{IsSymmetricMatrix}}
\logpage{[ 24, 4, 13 ]}\nobreak
\hyperdef{L}{X7ACD31A77E8458AA}{}
{\noindent\textcolor{FuncColor}{$\triangleright$\enspace\texttt{IsSymmetricMatrix({\mdseries\slshape mat})\index{IsSymmetricMatrix@\texttt{IsSymmetricMatrix}}
\label{IsSymmetricMatrix}
}\hfill{\scriptsize (property)}}\\
\noindent\textcolor{FuncColor}{$\triangleright$\enspace\texttt{IsSymmetricMat({\mdseries\slshape mat})\index{IsSymmetricMat@\texttt{IsSymmetricMat}}
\label{IsSymmetricMat}
}\hfill{\scriptsize (property)}}\\


 return \texttt{true} if the matrix \mbox{\texttt{\mdseries\slshape mat}} is a square matrix and satisfies \texttt{\mbox{\texttt{\mdseries\slshape mat}}[i,j] = \mbox{\texttt{\mdseries\slshape mat}}[j,i]} for all $i, j$, and \texttt{false} otherwise. 
\begin{Verbatim}[commandchars=!@|,fontsize=\small,frame=single,label=Example]
  !gapprompt@gap>| !gapinput@IsSymmetricMatrix( [ [ 1 ] ] );|
  true
  !gapprompt@gap>| !gapinput@IsSymmetricMatrix( [ [ 1, 2 ], [ 2, 1 ] ] );|
  true
  !gapprompt@gap>| !gapinput@IsSymmetricMatrix( [ [ 0, 1 ], [ 2, 0 ] ] );|
  false
  !gapprompt@gap>| !gapinput@IsSymmetricMatrix( [ [ 1, 2, 3 ], [ 2, 4, 5 ] ] );|
  false
\end{Verbatim}
 }

 

\subsection{\textcolor{Chapter }{IsAntisymmetricMatrix}}
\logpage{[ 24, 4, 14 ]}\nobreak
\hyperdef{L}{X833DA1E07B0B23EF}{}
{\noindent\textcolor{FuncColor}{$\triangleright$\enspace\texttt{IsAntisymmetricMatrix({\mdseries\slshape mat})\index{IsAntisymmetricMatrix@\texttt{IsAntisymmetricMatrix}}
\label{IsAntisymmetricMatrix}
}\hfill{\scriptsize (property)}}\\
\noindent\textcolor{FuncColor}{$\triangleright$\enspace\texttt{IsAntisymmetricMat({\mdseries\slshape mat})\index{IsAntisymmetricMat@\texttt{IsAntisymmetricMat}}
\label{IsAntisymmetricMat}
}\hfill{\scriptsize (property)}}\\


 return \texttt{true} if the matrix \mbox{\texttt{\mdseries\slshape mat}} is a square matrix and satisfies \texttt{\mbox{\texttt{\mdseries\slshape mat}}[i,j] = \texttt{\symbol{45}}\mbox{\texttt{\mdseries\slshape mat}}[j,i]} for all $i, j$, and \texttt{false} otherwise. (In particular, the diagonal entries must be zero.) 
\begin{Verbatim}[commandchars=!@|,fontsize=\small,frame=single,label=Example]
  !gapprompt@gap>| !gapinput@IsAntisymmetricMatrix( [ [ 0 ] ] );|
  true
  !gapprompt@gap>| !gapinput@IsAntisymmetricMatrix( [ [ 0, 1 ], [ -1, 0 ] ] );|
  true
  !gapprompt@gap>| !gapinput@IsAntisymmetricMatrix( [ [ 0, 1 ], [ 1, 0 ] ] );|
  false
  !gapprompt@gap>| !gapinput@IsAntisymmetricMatrix( [ [ 1, 2, 3 ], [ 4, 5, 6 ] ] );|
  false
\end{Verbatim}
 }

 }

  
\section{\textcolor{Chapter }{Matrix Constructions}}\label{Matrix Constructions}
\logpage{[ 24, 5, 0 ]}
\hyperdef{L}{X823FB2398697B957}{}
{
  

\subsection{\textcolor{Chapter }{IdentityMat}}
\logpage{[ 24, 5, 1 ]}\nobreak
\hyperdef{L}{X7DB902CE848D1524}{}
{\noindent\textcolor{FuncColor}{$\triangleright$\enspace\texttt{IdentityMat({\mdseries\slshape m[, R]})\index{IdentityMat@\texttt{IdentityMat}}
\label{IdentityMat}
}\hfill{\scriptsize (function)}}\\


 returns a (mutable) \mbox{\texttt{\mdseries\slshape m}}$\times$\mbox{\texttt{\mdseries\slshape m}} identity matrix over the ring given by \mbox{\texttt{\mdseries\slshape R}}. Here, \mbox{\texttt{\mdseries\slshape R}} can be either a ring, or an element of a ring. By default, an integer matrix
is created. 
\begin{Verbatim}[commandchars=!@|,fontsize=\small,frame=single,label=Example]
  !gapprompt@gap>| !gapinput@IdentityMat(3);|
  [ [ 1, 0, 0 ], [ 0, 1, 0 ], [ 0, 0, 1 ] ]
  !gapprompt@gap>| !gapinput@IdentityMat(2,Integers mod 15);|
  [ [ ZmodnZObj( 1, 15 ), ZmodnZObj( 0, 15 ) ],
    [ ZmodnZObj( 0, 15 ), ZmodnZObj( 1, 15 ) ] ]
  !gapprompt@gap>| !gapinput@IdentityMat(2,Z(3));|
  [ [ Z(3)^0, 0*Z(3) ], [ 0*Z(3), Z(3)^0 ] ]
\end{Verbatim}
 }

 

\subsection{\textcolor{Chapter }{NullMat}}
\logpage{[ 24, 5, 2 ]}\nobreak
\hyperdef{L}{X86D343A77D9B3D4D}{}
{\noindent\textcolor{FuncColor}{$\triangleright$\enspace\texttt{NullMat({\mdseries\slshape m, n[, R]})\index{NullMat@\texttt{NullMat}}
\label{NullMat}
}\hfill{\scriptsize (function)}}\\


 returns a (mutable) \mbox{\texttt{\mdseries\slshape m}}$\times$\mbox{\texttt{\mdseries\slshape n}} null matrix over the ring given by by \mbox{\texttt{\mdseries\slshape R}}. Here, \mbox{\texttt{\mdseries\slshape R}} can be either a ring, or an element of a ring. By default, an integer matrix
is created. 
\begin{Verbatim}[commandchars=!@|,fontsize=\small,frame=single,label=Example]
  !gapprompt@gap>| !gapinput@NullMat(3,2);|
  [ [ 0, 0 ], [ 0, 0 ], [ 0, 0 ] ]
  !gapprompt@gap>| !gapinput@NullMat(2,2,Integers mod 15);|
  [ [ ZmodnZObj( 0, 15 ), ZmodnZObj( 0, 15 ) ],
    [ ZmodnZObj( 0, 15 ), ZmodnZObj( 0, 15 ) ] ]
  !gapprompt@gap>| !gapinput@NullMat(3,2,Z(3));|
  [ [ 0*Z(3), 0*Z(3) ], [ 0*Z(3), 0*Z(3) ], [ 0*Z(3), 0*Z(3) ] ]
\end{Verbatim}
 }

 

\subsection{\textcolor{Chapter }{EmptyMatrix}}
\logpage{[ 24, 5, 3 ]}\nobreak
\hyperdef{L}{X8508A7EA812BA0CC}{}
{\noindent\textcolor{FuncColor}{$\triangleright$\enspace\texttt{EmptyMatrix({\mdseries\slshape char})\index{EmptyMatrix@\texttt{EmptyMatrix}}
\label{EmptyMatrix}
}\hfill{\scriptsize (function)}}\\


 is an empty (ordinary) matrix in characteristic \mbox{\texttt{\mdseries\slshape char}} that can be added to or multiplied with empty lists (representing
zero\texttt{\symbol{45}}dimensional row vectors). It also acts (via the
operation \texttt{\texttt{\symbol{92}}\texttt{\symbol{94}}} (\ref{bSlash^})) on empty lists. 

  
\begin{Verbatim}[commandchars=!@|,fontsize=\small,frame=single,label=Example]
  !gapprompt@gap>| !gapinput@EmptyMatrix(5);|
  EmptyMatrix( 5 )
  !gapprompt@gap>| !gapinput@AsList(last);|
  [  ]
\end{Verbatim}
 }

 

\subsection{\textcolor{Chapter }{DiagonalMat}}
\logpage{[ 24, 5, 4 ]}\nobreak
\hyperdef{L}{X81042E7A7F247ADE}{}
{\noindent\textcolor{FuncColor}{$\triangleright$\enspace\texttt{DiagonalMat({\mdseries\slshape vector})\index{DiagonalMat@\texttt{DiagonalMat}}
\label{DiagonalMat}
}\hfill{\scriptsize (function)}}\\


 returns a diagonal matrix \mbox{\texttt{\mdseries\slshape mat}} with the diagonal entries given by \mbox{\texttt{\mdseries\slshape vector}}. 
\begin{Verbatim}[commandchars=!@|,fontsize=\small,frame=single,label=Example]
  !gapprompt@gap>| !gapinput@DiagonalMat([1,2,3]);|
  [ [ 1, 0, 0 ], [ 0, 2, 0 ], [ 0, 0, 3 ] ]
\end{Verbatim}
 }

 

\subsection{\textcolor{Chapter }{DiagonalMatrix (with base domain)}}
\logpage{[ 24, 5, 5 ]}\nobreak
\hyperdef{L}{X87BADF217C19CBE1}{}
{\noindent\textcolor{FuncColor}{$\triangleright$\enspace\texttt{DiagonalMatrix({\mdseries\slshape [filt, ]R, vector})\index{DiagonalMatrix@\texttt{DiagonalMatrix}!with base domain}
\label{DiagonalMatrix:with base domain}
}\hfill{\scriptsize (operation)}}\\
\noindent\textcolor{FuncColor}{$\triangleright$\enspace\texttt{DiagonalMatrix({\mdseries\slshape vector[, M]})\index{DiagonalMatrix@\texttt{DiagonalMatrix}!with example matrix}
\label{DiagonalMatrix:with example matrix}
}\hfill{\scriptsize (operation)}}\\
\textbf{\indent Returns:\ }
 a square matrix or matrix object with column number equal to the length of the
dense list \mbox{\texttt{\mdseries\slshape vector}}, whose diagonal entries are given by the entries of \mbox{\texttt{\mdseries\slshape vector}}, and whose off\texttt{\symbol{45}}diagonal entries are zero. 



 If a semiring \mbox{\texttt{\mdseries\slshape R}} is given then it will be the base domain (see \texttt{BaseDomain} (\ref{BaseDomain:for a matrix object})) of the returned matrix. In this case, a filter \mbox{\texttt{\mdseries\slshape filt}} can be specified that defines the internal representation of the result (see \texttt{ConstructingFilter} (\ref{ConstructingFilter:for a matrix object})). The default value for \mbox{\texttt{\mdseries\slshape filt}} is determined from \mbox{\texttt{\mdseries\slshape R}}. 

 If a matrix object \mbox{\texttt{\mdseries\slshape M}} is given then the returned matrix will have the same internal representation
and the same base domain as \mbox{\texttt{\mdseries\slshape M}}. 

 If only \mbox{\texttt{\mdseries\slshape vector}} is given then it is used to compute a default for \mbox{\texttt{\mdseries\slshape R}}. 

 If the \texttt{ConstructingFilter} (\ref{ConstructingFilter:for a matrix object}) value of the result implies \texttt{IsCopyable} (\ref{IsCopyable}) then the result is fully mutable. 

 
\begin{Verbatim}[commandchars=!@|,fontsize=\small,frame=single,label=Example]
  !gapprompt@gap>| !gapinput@d1:= DiagonalMatrix( GF(9), [ 1, 2 ] * Z(3)^0 );|
  [ [ Z(3)^0, 0*Z(3) ], [ 0*Z(3), Z(3) ] ]
  !gapprompt@gap>| !gapinput@Is8BitMatrixRep( d1 );|
  true
  !gapprompt@gap>| !gapinput@d2:= DiagonalMatrix( IsPlistMatrixRep, GF(9), [ 1, 2 ] * Z(3)^0 );|
  <2x2-matrix over GF(3^2)>
  !gapprompt@gap>| !gapinput@IsPlistMatrixRep( d2 );|
  true
  !gapprompt@gap>| !gapinput@DiagonalMatrix( [ 1, 2 ] );|
  <2x2-matrix over Rationals>
  !gapprompt@gap>| !gapinput@DiagonalMatrix( [ 1, 2 ], Matrix( Integers, [ [ 1 ] ], 1 ) );|
  <2x2-matrix over Integers>
  !gapprompt@gap>| !gapinput@DiagonalMatrix( [ 1, 2 ], [ [ 1 ] ] );|
  [ [ 1, 0 ], [ 0, 2 ] ]
\end{Verbatim}
 }

 

\subsection{\textcolor{Chapter }{PermutationMat}}
\logpage{[ 24, 5, 6 ]}\nobreak
\hyperdef{L}{X806C62A67A7D5379}{}
{\noindent\textcolor{FuncColor}{$\triangleright$\enspace\texttt{PermutationMat({\mdseries\slshape perm, dim[, F]})\index{PermutationMat@\texttt{PermutationMat}}
\label{PermutationMat}
}\hfill{\scriptsize (function)}}\\


 returns a matrix in dimension \mbox{\texttt{\mdseries\slshape dim}} over the field given by \mbox{\texttt{\mdseries\slshape F}} (i.e. the smallest field containing the element \mbox{\texttt{\mdseries\slshape F}} or \mbox{\texttt{\mdseries\slshape F}} itself if it is a field) that represents the permutation \mbox{\texttt{\mdseries\slshape perm}} acting by permuting the basis vectors as it permutes points. 
\begin{Verbatim}[commandchars=!@|,fontsize=\small,frame=single,label=Example]
  !gapprompt@gap>| !gapinput@PermutationMat((1,2,3),4,1);|
  [ [ 0, 1, 0, 0 ], [ 0, 0, 1, 0 ], [ 1, 0, 0, 0 ], [ 0, 0, 0, 1 ] ]
\end{Verbatim}
 }

 

\subsection{\textcolor{Chapter }{TransposedMatImmutable}}
\logpage{[ 24, 5, 7 ]}\nobreak
\hyperdef{L}{X7C52A38C79C36C35}{}
{\noindent\textcolor{FuncColor}{$\triangleright$\enspace\texttt{TransposedMatImmutable({\mdseries\slshape mat})\index{TransposedMatImmutable@\texttt{TransposedMatImmutable}}
\label{TransposedMatImmutable}
}\hfill{\scriptsize (attribute)}}\\
\noindent\textcolor{FuncColor}{$\triangleright$\enspace\texttt{TransposedMat({\mdseries\slshape mat})\index{TransposedMat@\texttt{TransposedMat}}
\label{TransposedMat}
}\hfill{\scriptsize (attribute)}}\\
\noindent\textcolor{FuncColor}{$\triangleright$\enspace\texttt{TransposedMatMutable({\mdseries\slshape mat})\index{TransposedMatMutable@\texttt{TransposedMatMutable}}
\label{TransposedMatMutable}
}\hfill{\scriptsize (operation)}}\\
\noindent\textcolor{FuncColor}{$\triangleright$\enspace\texttt{TransposedMatOp({\mdseries\slshape mat})\index{TransposedMatOp@\texttt{TransposedMatOp}}
\label{TransposedMatOp}
}\hfill{\scriptsize (operation)}}\\


 These functions all return the transposed of the matrix object \mbox{\texttt{\mdseries\slshape mat}}, i.e., a matrix object $trans$ such that $trans[i,k] = \mbox{\texttt{\mdseries\slshape mat}}[k,i]$ holds. 

 They differ only w.r.t. the mutability of the result. 

 \texttt{TransposedMat} is an attribute and hence returns an immutable result. \texttt{TransposedMatMutable} is guaranteed to return a new \emph{mutable} matrix. 

 \texttt{TransposedMatImmutable} is a synonym of \texttt{TransposedMat}, and \texttt{TransposedMatOp} is a synonym of \texttt{TransposedMatMutable}, in analogy to operations such as \texttt{Zero} (\ref{Zero}). }

 

\subsection{\textcolor{Chapter }{TransposedMatDestructive}}
\logpage{[ 24, 5, 8 ]}\nobreak
\hyperdef{L}{X7DBB40847E2B6252}{}
{\noindent\textcolor{FuncColor}{$\triangleright$\enspace\texttt{TransposedMatDestructive({\mdseries\slshape mat})\index{TransposedMatDestructive@\texttt{TransposedMatDestructive}}
\label{TransposedMatDestructive}
}\hfill{\scriptsize (operation)}}\\


 If \mbox{\texttt{\mdseries\slshape mat}} is a mutable matrix, then the transposed is computed by swapping the entries
in \mbox{\texttt{\mdseries\slshape mat}}. In this way \mbox{\texttt{\mdseries\slshape mat}} gets changed. In all other cases the transposed is computed by \texttt{TransposedMat} (\ref{TransposedMat}). 
\begin{Verbatim}[commandchars=!@|,fontsize=\small,frame=single,label=Example]
  !gapprompt@gap>| !gapinput@TransposedMat([[1,2,3],[4,5,6],[7,8,9]]);|
  [ [ 1, 4, 7 ], [ 2, 5, 8 ], [ 3, 6, 9 ] ]
  !gapprompt@gap>| !gapinput@mm:= [[1,2,3],[4,5,6],[7,8,9]];;|
  !gapprompt@gap>| !gapinput@TransposedMatDestructive( mm );|
  [ [ 1, 4, 7 ], [ 2, 5, 8 ], [ 3, 6, 9 ] ]
  !gapprompt@gap>| !gapinput@mm;|
  [ [ 1, 4, 7 ], [ 2, 5, 8 ], [ 3, 6, 9 ] ]
\end{Verbatim}
 }

 

\subsection{\textcolor{Chapter }{KroneckerProduct}}
\logpage{[ 24, 5, 9 ]}\nobreak
\hyperdef{L}{X8634C79E7DB22934}{}
{\noindent\textcolor{FuncColor}{$\triangleright$\enspace\texttt{KroneckerProduct({\mdseries\slshape mat1, mat2})\index{KroneckerProduct@\texttt{KroneckerProduct}}
\label{KroneckerProduct}
}\hfill{\scriptsize (operation)}}\\


 The Kronecker product of two matrices is the matrix obtained when replacing
each entry \mbox{\texttt{\mdseries\slshape a}} of \mbox{\texttt{\mdseries\slshape mat1}} by the product \texttt{\mbox{\texttt{\mdseries\slshape a}}*\mbox{\texttt{\mdseries\slshape mat2}}} in one matrix. 
\begin{Verbatim}[commandchars=!@|,fontsize=\small,frame=single,label=Example]
  !gapprompt@gap>| !gapinput@KroneckerProduct([[1,2]],[[5,7],[9,2]]);|
  [ [ 5, 7, 10, 14 ], [ 9, 2, 18, 4 ] ]
\end{Verbatim}
 }

 

\subsection{\textcolor{Chapter }{ReflectionMat}}
\logpage{[ 24, 5, 10 ]}\nobreak
\hyperdef{L}{X845EC4D18054D140}{}
{\noindent\textcolor{FuncColor}{$\triangleright$\enspace\texttt{ReflectionMat({\mdseries\slshape coeffs[, conj][, root]})\index{ReflectionMat@\texttt{ReflectionMat}}
\label{ReflectionMat}
}\hfill{\scriptsize (function)}}\\


 Let \mbox{\texttt{\mdseries\slshape coeffs}} be a row vector. \texttt{ReflectionMat} returns the matrix of the reflection in this vector. 

 More precisely, if \mbox{\texttt{\mdseries\slshape coeffs}} is the coefficients list of a vector $v$ w.r.t. a basis $B$ (see{\nobreakspace}\texttt{Basis} (\ref{Basis})) then the returned matrix describes the reflection in $v$ w.r.t. $B$ as a map on a row space, with action from the right. 

 The optional argument \mbox{\texttt{\mdseries\slshape root}} is a root of unity that determines the order of the reflection. The default is
a reflection of order 2. For triflections one should choose a third root of
unity etc. (see{\nobreakspace}\texttt{E} (\ref{E})). 

 \mbox{\texttt{\mdseries\slshape conj}} is a function of one argument that conjugates a ring element. The default is \texttt{ComplexConjugate} (\ref{ComplexConjugate}). 

 The matrix of the reflection in $v$ is defined as 
\[ M = I_n + \mbox{\texttt{\mdseries\slshape conj}}(v^{tr}) \cdot (\mbox{\texttt{\mdseries\slshape root}}-1) / (v \cdot \mbox{\texttt{\mdseries\slshape conj}}(v^{tr})) \cdot v \]
 where $n$ is the length of the coefficient list. 

 So $v$ is mapped to \mbox{\texttt{\mdseries\slshape root}}$ \cdot v$, with default $-v$, and any vector $x$ with the property $x \cdot $\mbox{\texttt{\mdseries\slshape conj}}$(v^{tr}) = 0$ is fixed by the reflection. }

 

\subsection{\textcolor{Chapter }{PrintArray}}
\logpage{[ 24, 5, 11 ]}\nobreak
\hyperdef{L}{X7DEBC9967DFDFC18}{}
{\noindent\textcolor{FuncColor}{$\triangleright$\enspace\texttt{PrintArray({\mdseries\slshape array})\index{PrintArray@\texttt{PrintArray}}
\label{PrintArray}
}\hfill{\scriptsize (function)}}\\


 pretty\texttt{\symbol{45}}prints the array \mbox{\texttt{\mdseries\slshape array}}. }

 }

  
\section{\textcolor{Chapter }{Random Matrices}}\label{Random Matrices}
\logpage{[ 24, 6, 0 ]}
\hyperdef{L}{X79CC5F568252D341}{}
{
  

\subsection{\textcolor{Chapter }{RandomMat}}
\logpage{[ 24, 6, 1 ]}\nobreak
\hyperdef{L}{X7F957F0280A87961}{}
{\noindent\textcolor{FuncColor}{$\triangleright$\enspace\texttt{RandomMat({\mdseries\slshape [rs, ]m, n[, R]})\index{RandomMat@\texttt{RandomMat}}
\label{RandomMat}
}\hfill{\scriptsize (function)}}\\


 \texttt{RandomMat} returns a new mutable random matrix with \mbox{\texttt{\mdseries\slshape m}} rows and \mbox{\texttt{\mdseries\slshape n}} columns with elements taken from the ring \mbox{\texttt{\mdseries\slshape R}}, which defaults to \texttt{Integers} (\ref{Integers}). Optionally, a random source \mbox{\texttt{\mdseries\slshape rs}} can be supplied. 
\begin{Verbatim}[commandchars=!@|,fontsize=\small,frame=single,label=Example]
  !gapprompt@gap>| !gapinput@RandomMat(2,3,GF(3));|
  [ [ Z(3), Z(3), 0*Z(3) ], [ Z(3), Z(3)^0, Z(3) ] ]
\end{Verbatim}
 }

 

\subsection{\textcolor{Chapter }{RandomMatrix (with base domain)}}
\logpage{[ 24, 6, 2 ]}\nobreak
\hyperdef{L}{X7D31ACA4809B321B}{}
{\noindent\textcolor{FuncColor}{$\triangleright$\enspace\texttt{RandomMatrix({\mdseries\slshape [filt, ][rs, ]R, m, n})\index{RandomMatrix@\texttt{RandomMatrix}!with base domain}
\label{RandomMatrix:with base domain}
}\hfill{\scriptsize (operation)}}\\
\noindent\textcolor{FuncColor}{$\triangleright$\enspace\texttt{RandomMatrix({\mdseries\slshape [rs, ]m, n, M})\index{RandomMatrix@\texttt{RandomMatrix}!with example matrix}
\label{RandomMatrix:with example matrix}
}\hfill{\scriptsize (operation)}}\\
\textbf{\indent Returns:\ }
 a matrix or matrix object with \mbox{\texttt{\mdseries\slshape m}} rows and \mbox{\texttt{\mdseries\slshape n}} columns, with base domain \mbox{\texttt{\mdseries\slshape R}} or equal to \texttt{BaseDomain( \mbox{\texttt{\mdseries\slshape M}} )}. 



 If a semiring \mbox{\texttt{\mdseries\slshape R}} is given then it will be the base domain (see \texttt{BaseDomain} (\ref{BaseDomain:for a matrix object})) of the returned matrix. In this case, a filter \mbox{\texttt{\mdseries\slshape filt}} can be specified that defines the internal representation of the result (see \texttt{ConstructingFilter} (\ref{ConstructingFilter:for a matrix object})). The default value for \mbox{\texttt{\mdseries\slshape filt}} is determined from \mbox{\texttt{\mdseries\slshape R}}. 

 If a matrix object \mbox{\texttt{\mdseries\slshape M}} is given then the returned matrix will have the same internal representation
and the same base domain as \mbox{\texttt{\mdseries\slshape M}}. 

 If a random source \mbox{\texttt{\mdseries\slshape rs}} is given (see \ref{Random Sources}) then the entries of the result are computed using \mbox{\texttt{\mdseries\slshape rs}}, the default random source is \texttt{GlobalMersenneTwister} (\ref{GlobalMersenneTwister}). 

 If the \texttt{ConstructingFilter} (\ref{ConstructingFilter:for a matrix object}) value of the result implies \texttt{IsCopyable} (\ref{IsCopyable}) then the result is fully mutable. 

 
\begin{Verbatim}[commandchars=!@|,fontsize=\small,frame=single,label=Example]
  !gapprompt@gap>| !gapinput@m:= RandomMatrix( GF(9), 2, 3 );;|
  !gapprompt@gap>| !gapinput@DimensionsMat( m );|
  [ 2, 3 ]
  !gapprompt@gap>| !gapinput@BaseDomain( m );|
  GF(3^2)
  !gapprompt@gap>| !gapinput@Is8BitMatrixRep( m );|
  true
  !gapprompt@gap>| !gapinput@m1:= RandomMatrix( IsPlistMatrixRep, GF(9), 2, 3 );;|
  !gapprompt@gap>| !gapinput@IsPlistMatrixRep( m1 );|
  true
  !gapprompt@gap>| !gapinput@m2:= RandomMatrix( 3, 4, m1 );;|
  !gapprompt@gap>| !gapinput@DimensionsMat( m2 );|
  [ 3, 4 ]
  !gapprompt@gap>| !gapinput@IsPlistMatrixRep( m2 );|
  true
\end{Verbatim}
 }

 

\subsection{\textcolor{Chapter }{RandomInvertibleMat}}
\logpage{[ 24, 6, 3 ]}\nobreak
\hyperdef{L}{X7C939B4A7EDF015D}{}
{\noindent\textcolor{FuncColor}{$\triangleright$\enspace\texttt{RandomInvertibleMat({\mdseries\slshape [rs, ]m[, R]})\index{RandomInvertibleMat@\texttt{RandomInvertibleMat}}
\label{RandomInvertibleMat}
}\hfill{\scriptsize (function)}}\\


 \texttt{RandomInvertibleMat} returns a new mutable invertible random matrix with \mbox{\texttt{\mdseries\slshape m}} rows and columns with elements taken from the ring \mbox{\texttt{\mdseries\slshape R}}, which defaults to \texttt{Integers} (\ref{Integers}). Optionally, a random source \mbox{\texttt{\mdseries\slshape rs}} can be supplied. 

 Note that the result need \emph{not} be invertible over the prescribed ring \mbox{\texttt{\mdseries\slshape R}}. Use \texttt{RandomInvertibleMatrix} (\ref{RandomInvertibleMatrix:with base domain}) for creating random matrices with this stronger property. 

 
\begin{Verbatim}[commandchars=!@|,fontsize=\small,frame=single,label=Example]
  !gapprompt@gap>| !gapinput@m := RandomInvertibleMat( 4, Integers );|
  [ [ 1, -4, 0, 1 ], [ 0, -2, 3, 1 ], [ 0, -2, -1, -1 ],
    [ 0, -2, 2, 4 ] ]
  !gapprompt@gap>| !gapinput@m^-1;|
  [ [ 1, -1/14, -19/14, -4/7 ], [ 0, -1/14, -5/14, -1/14 ],
    [ 0, 5/14, -3/14, -1/7 ], [ 0, -3/14, -1/14, 2/7 ] ]
\end{Verbatim}
 }

 

\subsection{\textcolor{Chapter }{RandomInvertibleMatrix (with base domain)}}
\logpage{[ 24, 6, 4 ]}\nobreak
\hyperdef{L}{X8743657384B15B5B}{}
{\noindent\textcolor{FuncColor}{$\triangleright$\enspace\texttt{RandomInvertibleMatrix({\mdseries\slshape [filt, ][rs, ]R, m})\index{RandomInvertibleMatrix@\texttt{RandomInvertibleMatrix}!with base domain}
\label{RandomInvertibleMatrix:with base domain}
}\hfill{\scriptsize (operation)}}\\
\noindent\textcolor{FuncColor}{$\triangleright$\enspace\texttt{RandomInvertibleMatrix({\mdseries\slshape [rs, ]m, M})\index{RandomInvertibleMatrix@\texttt{RandomInvertibleMatrix}!with example matrix}
\label{RandomInvertibleMatrix:with example matrix}
}\hfill{\scriptsize (operation)}}\\
\textbf{\indent Returns:\ }
 a matrix or matrix object $I$ with \mbox{\texttt{\mdseries\slshape m}} columns, with base domain \mbox{\texttt{\mdseries\slshape R}} or equal to \texttt{BaseDomain( \mbox{\texttt{\mdseries\slshape M}} )}, such that $I$ is invertible over its base domain. 



 If a semiring \mbox{\texttt{\mdseries\slshape R}} is given then it will be the base domain (see \texttt{BaseDomain} (\ref{BaseDomain:for a matrix object})) of the returned matrix. In this case, a filter \mbox{\texttt{\mdseries\slshape filt}} can be specified that defines the internal representation of the result (see \texttt{ConstructingFilter} (\ref{ConstructingFilter:for a matrix object})). The default value for \mbox{\texttt{\mdseries\slshape filt}} is determined from \mbox{\texttt{\mdseries\slshape R}}. 

 If a matrix object \mbox{\texttt{\mdseries\slshape M}} is given then the returned matrix will have the same internal representation
and the same base domain as \mbox{\texttt{\mdseries\slshape M}}. 

 If a random source \mbox{\texttt{\mdseries\slshape rs}} is given (see \ref{Random Sources}) then the entries of the result are computed using \mbox{\texttt{\mdseries\slshape rs}}, the default random source is \texttt{GlobalMersenneTwister} (\ref{GlobalMersenneTwister}). 

 If the \texttt{ConstructingFilter} (\ref{ConstructingFilter:for a matrix object}) value of the result implies \texttt{IsCopyable} (\ref{IsCopyable}) then the result is fully mutable. 

 
\begin{Verbatim}[commandchars=!@|,fontsize=\small,frame=single,label=Example]
  !gapprompt@gap>| !gapinput@m:= RandomInvertibleMatrix( GF(9), 2 );;|
  !gapprompt@gap>| !gapinput@RankMatrix( m );|
  2
  !gapprompt@gap>| !gapinput@DimensionsMat( m );|
  [ 2, 2 ]
  !gapprompt@gap>| !gapinput@BaseDomain( m );|
  GF(3^2)
  !gapprompt@gap>| !gapinput@Is8BitMatrixRep( m );|
  true
  !gapprompt@gap>| !gapinput@m1:= RandomInvertibleMatrix( IsPlistMatrixRep, GF(9), 2 );;|
  !gapprompt@gap>| !gapinput@IsPlistMatrixRep( m1 );|
  true
  !gapprompt@gap>| !gapinput@m2:= RandomInvertibleMatrix( 2, m1 );;|
  !gapprompt@gap>| !gapinput@IsPlistMatrixRep( m2 );|
  true
  !gapprompt@gap>| !gapinput@m:= RandomInvertibleMatrix( Integers, 2 );;|
  !gapprompt@gap>| !gapinput@IsUnit( Integers, DeterminantMatrix( m ) );|
  true
\end{Verbatim}
 }

 

\subsection{\textcolor{Chapter }{RandomUnimodularMat}}
\logpage{[ 24, 6, 5 ]}\nobreak
\hyperdef{L}{X84743732846ACB44}{}
{\noindent\textcolor{FuncColor}{$\triangleright$\enspace\texttt{RandomUnimodularMat({\mdseries\slshape [rs, ]m})\index{RandomUnimodularMat@\texttt{RandomUnimodularMat}}
\label{RandomUnimodularMat}
}\hfill{\scriptsize (function)}}\\


 returns a new random mutable \mbox{\texttt{\mdseries\slshape m}}$\times$\mbox{\texttt{\mdseries\slshape m}} matrix with integer entries that is invertible over the integers. Optionally,
a random source \mbox{\texttt{\mdseries\slshape rs}} can be supplied. If the option \mbox{\texttt{\mdseries\slshape domain}} is given, random selection is made from \mbox{\texttt{\mdseries\slshape domain}}, otherwise from \mbox{\texttt{\mdseries\slshape Integers}} 
\begin{Verbatim}[commandchars=!@|,fontsize=\small,frame=single,label=Example]
  !gapprompt@gap>| !gapinput@m := RandomUnimodularMat(3);|
  [ [ 1, -5, 26 ], [ 0, 1, -6 ], [ 0, 0, 1 ] ]
  !gapprompt@gap>| !gapinput@m^-1;|
  [ [ 1, 5, 4 ], [ 0, 1, 6 ], [ 0, 0, 1 ] ]
  !gapprompt@gap>| !gapinput@RandomUnimodularMat(3:domain:=[-1000..1000]);|
  [ [ 6001314, 239514853714734, -44705745791413 ],
    [ -1448123744, -57795200635226587, 10787546189680040 ],
    [ -671, -26780118200, 4998542420 ] ]
\end{Verbatim}
 }

 }

  
\section{\textcolor{Chapter }{Matrices Representing Linear Equations and the Gaussian Algorithm}}\label{Matrices Representing Linear Equations and the Gaussian Algorithm}
\logpage{[ 24, 7, 0 ]}
\hyperdef{L}{X85485DCE809E323A}{}
{
  \index{Gaussian algorithm} 

\subsection{\textcolor{Chapter }{RankMatrix}}
\logpage{[ 24, 7, 1 ]}\nobreak
\hyperdef{L}{X7A995A74838950E6}{}
{\noindent\textcolor{FuncColor}{$\triangleright$\enspace\texttt{RankMatrix({\mdseries\slshape mat})\index{RankMatrix@\texttt{RankMatrix}}
\label{RankMatrix}
}\hfill{\scriptsize (attribute)}}\\
\noindent\textcolor{FuncColor}{$\triangleright$\enspace\texttt{RankMat({\mdseries\slshape mat})\index{RankMat@\texttt{RankMat}}
\label{RankMat}
}\hfill{\scriptsize (attribute)}}\\


 If \mbox{\texttt{\mdseries\slshape mat}} is a matrix object representing a matrix whose rows span a free module over
the ring generated by the matrix entries and their inverses then \texttt{RankMatrix} returns the dimension of this free module. Otherwise \texttt{fail} is returned. 

 Note that \texttt{RankMatrix} may perform a Gaussian elimination. For large rational matrices this may take
very long, because the entries may become very large. 
\begin{Verbatim}[commandchars=!@|,fontsize=\small,frame=single,label=Example]
  !gapprompt@gap>| !gapinput@mat:=[[1,2,3],[4,5,6],[7,8,9]];;|
  !gapprompt@gap>| !gapinput@RankMatrix( mat );|
  2
\end{Verbatim}
 }

 

\subsection{\textcolor{Chapter }{TriangulizedMat}}
\logpage{[ 24, 7, 2 ]}\nobreak
\hyperdef{L}{X7BA26C3387AB434E}{}
{\noindent\textcolor{FuncColor}{$\triangleright$\enspace\texttt{TriangulizedMat({\mdseries\slshape mat})\index{TriangulizedMat@\texttt{TriangulizedMat}}
\label{TriangulizedMat}
}\hfill{\scriptsize (operation)}}\\
\noindent\textcolor{FuncColor}{$\triangleright$\enspace\texttt{RREF({\mdseries\slshape mat})\index{RREF@\texttt{RREF}}
\label{RREF}
}\hfill{\scriptsize (operation)}}\\


 Computes an upper triangular form of the matrix \mbox{\texttt{\mdseries\slshape mat}} via the Gaussian Algorithm. It returns a mutable matrix in upper triangular
form. This is sometimes also called ``Hermite normal form'' or ``Reduced Row Echelon Form''. \texttt{RREF} is a synonym for \texttt{TriangulizedMat}. }

 

\subsection{\textcolor{Chapter }{TriangulizeMat}}
\logpage{[ 24, 7, 3 ]}\nobreak
\hyperdef{L}{X8384CA8E7B3850D3}{}
{\noindent\textcolor{FuncColor}{$\triangleright$\enspace\texttt{TriangulizeMat({\mdseries\slshape mat})\index{TriangulizeMat@\texttt{TriangulizeMat}}
\label{TriangulizeMat}
}\hfill{\scriptsize (operation)}}\\


 Applies the Gaussian Algorithm to the mutable matrix \mbox{\texttt{\mdseries\slshape mat}} and changes \mbox{\texttt{\mdseries\slshape mat}} such that it is in upper triangular normal form (sometimes called ``Hermite normal form'' or ``Reduced Row Echelon Form''). 
\begin{Verbatim}[commandchars=!@|,fontsize=\small,frame=single,label=Example]
  !gapprompt@gap>| !gapinput@m:=TransposedMatMutable(mat);|
  [ [ 1, 4, 7 ], [ 2, 5, 8 ], [ 3, 6, 9 ] ]
  !gapprompt@gap>| !gapinput@TriangulizeMat(m);m;|
  [ [ 1, 0, -1 ], [ 0, 1, 2 ], [ 0, 0, 0 ] ]
  !gapprompt@gap>| !gapinput@m:=TransposedMatMutable(mat);|
  [ [ 1, 4, 7 ], [ 2, 5, 8 ], [ 3, 6, 9 ] ]
  !gapprompt@gap>| !gapinput@TriangulizedMat(m);m;|
  [ [ 1, 0, -1 ], [ 0, 1, 2 ], [ 0, 0, 0 ] ]
  [ [ 1, 4, 7 ], [ 2, 5, 8 ], [ 3, 6, 9 ] ]
\end{Verbatim}
 }

 

\subsection{\textcolor{Chapter }{NullspaceMat}}
\logpage{[ 24, 7, 4 ]}\nobreak
\hyperdef{L}{X7DA0D5887DB12DC4}{}
{\noindent\textcolor{FuncColor}{$\triangleright$\enspace\texttt{NullspaceMat({\mdseries\slshape mat})\index{NullspaceMat@\texttt{NullspaceMat}}
\label{NullspaceMat}
}\hfill{\scriptsize (attribute)}}\\
\noindent\textcolor{FuncColor}{$\triangleright$\enspace\texttt{TriangulizedNullspaceMat({\mdseries\slshape mat})\index{TriangulizedNullspaceMat@\texttt{TriangulizedNullspaceMat}}
\label{TriangulizedNullspaceMat}
}\hfill{\scriptsize (attribute)}}\\


 \index{kernel!of a matrix} returns a list of row vectors that form a basis of the vector space of
solutions to the equation \texttt{\mbox{\texttt{\mdseries\slshape vec}}*\mbox{\texttt{\mdseries\slshape mat}}=0}. The result is an immutable matrix. This basis is not guaranteed to be in any
specific form. 

 The variant \texttt{TriangulizedNullspaceMat} returns a basis of the nullspace in triangulized form as is often needed for
algorithms. }

 

\subsection{\textcolor{Chapter }{NullspaceMatDestructive}}
\logpage{[ 24, 7, 5 ]}\nobreak
\hyperdef{L}{X87684B0F7AB7B7DB}{}
{\noindent\textcolor{FuncColor}{$\triangleright$\enspace\texttt{NullspaceMatDestructive({\mdseries\slshape mat})\index{NullspaceMatDestructive@\texttt{NullspaceMatDestructive}}
\label{NullspaceMatDestructive}
}\hfill{\scriptsize (operation)}}\\
\noindent\textcolor{FuncColor}{$\triangleright$\enspace\texttt{TriangulizedNullspaceMatDestructive({\mdseries\slshape mat})\index{TriangulizedNullspaceMatDestructive@\texttt{TriangulizedNullspaceMatDestructive}}
\label{TriangulizedNullspaceMatDestructive}
}\hfill{\scriptsize (operation)}}\\


 This function does the same as \texttt{NullspaceMat} (\ref{NullspaceMat}). However, the latter function makes a copy of \mbox{\texttt{\mdseries\slshape mat}} to avoid having to change it. This function does not do that; it returns the
nullspace and may destroy \mbox{\texttt{\mdseries\slshape mat}}; this saves a lot of memory in case \mbox{\texttt{\mdseries\slshape mat}} is big. The matrix \mbox{\texttt{\mdseries\slshape mat}} must be mutable. 

 The variant \texttt{TriangulizedNullspaceMatDestructive} returns a basis of the nullspace in triangulized form. It may destroy the
matrix \mbox{\texttt{\mdseries\slshape mat}}. 
\begin{Verbatim}[commandchars=!@|,fontsize=\small,frame=single,label=Example]
  !gapprompt@gap>| !gapinput@mat:=[[1,2,3],[4,5,6],[7,8,9]];;|
  !gapprompt@gap>| !gapinput@NullspaceMat(mat);|
  [ [ 1, -2, 1 ] ]
  !gapprompt@gap>| !gapinput@mm:=[[1,2,3],[4,5,6],[7,8,9]];;|
  !gapprompt@gap>| !gapinput@NullspaceMatDestructive( mm );|
  [ [ 1, -2, 1 ] ]
  !gapprompt@gap>| !gapinput@mm;|
  [ [ 1, 2, 3 ], [ 0, -3, -6 ], [ 0, 0, 0 ] ]
\end{Verbatim}
 }

 

\subsection{\textcolor{Chapter }{SolutionMat}}
\logpage{[ 24, 7, 6 ]}\nobreak
\hyperdef{L}{X838A519C7CD2969E}{}
{\noindent\textcolor{FuncColor}{$\triangleright$\enspace\texttt{SolutionMat({\mdseries\slshape mat, vec})\index{SolutionMat@\texttt{SolutionMat}}
\label{SolutionMat}
}\hfill{\scriptsize (operation)}}\\


 returns a row vector \mbox{\texttt{\mdseries\slshape x}} that is a solution of the equation \texttt{\mbox{\texttt{\mdseries\slshape x}} * \mbox{\texttt{\mdseries\slshape mat}} = \mbox{\texttt{\mdseries\slshape vec}}}. It returns \texttt{fail} if no such vector exists. }

 

\subsection{\textcolor{Chapter }{SolutionMatDestructive}}
\logpage{[ 24, 7, 7 ]}\nobreak
\hyperdef{L}{X7A7880D27CE7C1FE}{}
{\noindent\textcolor{FuncColor}{$\triangleright$\enspace\texttt{SolutionMatDestructive({\mdseries\slshape mat, vec})\index{SolutionMatDestructive@\texttt{SolutionMatDestructive}}
\label{SolutionMatDestructive}
}\hfill{\scriptsize (operation)}}\\


 Does the same as \texttt{SolutionMat( \mbox{\texttt{\mdseries\slshape mat}}, \mbox{\texttt{\mdseries\slshape vec}} )} except that it may destroy the matrix \mbox{\texttt{\mdseries\slshape mat}} and the vector \mbox{\texttt{\mdseries\slshape vec}}. The matrix \mbox{\texttt{\mdseries\slshape mat}} must be mutable. 
\begin{Verbatim}[commandchars=!@|,fontsize=\small,frame=single,label=Example]
  !gapprompt@gap>| !gapinput@mat:=[[1,2,3],[4,5,6],[7,8,9]];;|
  !gapprompt@gap>| !gapinput@SolutionMat(mat,[3,5,7]);|
  [ 5/3, 1/3, 0 ]
  !gapprompt@gap>| !gapinput@mm:= [[1,2,3],[4,5,6],[7,8,9]];;|
  !gapprompt@gap>| !gapinput@v:= [3,5,7];;|
  !gapprompt@gap>| !gapinput@SolutionMatDestructive( mm, v );|
  [ 5/3, 1/3, 0 ]
  !gapprompt@gap>| !gapinput@mm;|
  [ [ 1, 2, 3 ], [ 0, -3, -6 ], [ 0, 0, 0 ] ]
  !gapprompt@gap>| !gapinput@v;|
  [ 0, 0, 0 ]
\end{Verbatim}
 }

 

\subsection{\textcolor{Chapter }{BaseFixedSpace}}
\logpage{[ 24, 7, 8 ]}\nobreak
\hyperdef{L}{X7AB5AC547809F999}{}
{\noindent\textcolor{FuncColor}{$\triangleright$\enspace\texttt{BaseFixedSpace({\mdseries\slshape mats})\index{BaseFixedSpace@\texttt{BaseFixedSpace}}
\label{BaseFixedSpace}
}\hfill{\scriptsize (function)}}\\


 \texttt{BaseFixedSpace} returns a list of row vectors that form a base of the vector space $V$ such that $v M = v$ for all $v$ in $V$ and all matrices $M$ in the list \mbox{\texttt{\mdseries\slshape mats}}. (This is the common eigenspace of all matrices in \mbox{\texttt{\mdseries\slshape mats}} for the eigenvalue 1.) 
\begin{Verbatim}[commandchars=!@|,fontsize=\small,frame=single,label=Example]
  !gapprompt@gap>| !gapinput@BaseFixedSpace([[[1,2],[0,1]]]);|
  [ [ 0, 1 ] ]
\end{Verbatim}
 }

 }

  
\section{\textcolor{Chapter }{Eigenvectors and eigenvalues}}\label{Eigenvectors and eigenvalues}
\logpage{[ 24, 8, 0 ]}
\hyperdef{L}{X871FCAA97C60B2BA}{}
{
  

\subsection{\textcolor{Chapter }{GeneralisedEigenvalues}}
\logpage{[ 24, 8, 1 ]}\nobreak
\hyperdef{L}{X7A2462CC7B0C9D66}{}
{\noindent\textcolor{FuncColor}{$\triangleright$\enspace\texttt{GeneralisedEigenvalues({\mdseries\slshape F, A})\index{GeneralisedEigenvalues@\texttt{GeneralisedEigenvalues}}
\label{GeneralisedEigenvalues}
}\hfill{\scriptsize (operation)}}\\
\noindent\textcolor{FuncColor}{$\triangleright$\enspace\texttt{GeneralizedEigenvalues({\mdseries\slshape F, A})\index{GeneralizedEigenvalues@\texttt{GeneralizedEigenvalues}}
\label{GeneralizedEigenvalues}
}\hfill{\scriptsize (operation)}}\\


 The generalised eigenvalues of the matrix \mbox{\texttt{\mdseries\slshape A}} over the field \mbox{\texttt{\mdseries\slshape F}}. }

 

\subsection{\textcolor{Chapter }{GeneralisedEigenspaces}}
\logpage{[ 24, 8, 2 ]}\nobreak
\hyperdef{L}{X845CA0457D65876D}{}
{\noindent\textcolor{FuncColor}{$\triangleright$\enspace\texttt{GeneralisedEigenspaces({\mdseries\slshape F, A})\index{GeneralisedEigenspaces@\texttt{GeneralisedEigenspaces}}
\label{GeneralisedEigenspaces}
}\hfill{\scriptsize (operation)}}\\
\noindent\textcolor{FuncColor}{$\triangleright$\enspace\texttt{GeneralizedEigenspaces({\mdseries\slshape F, A})\index{GeneralizedEigenspaces@\texttt{GeneralizedEigenspaces}}
\label{GeneralizedEigenspaces}
}\hfill{\scriptsize (operation)}}\\


 The generalised eigenspaces of the matrix \mbox{\texttt{\mdseries\slshape A}} over the field \mbox{\texttt{\mdseries\slshape F}}. }

 

\subsection{\textcolor{Chapter }{Eigenvalues}}
\logpage{[ 24, 8, 3 ]}\nobreak
\hyperdef{L}{X8413C6FB7CEE9D59}{}
{\noindent\textcolor{FuncColor}{$\triangleright$\enspace\texttt{Eigenvalues({\mdseries\slshape F, A})\index{Eigenvalues@\texttt{Eigenvalues}}
\label{Eigenvalues}
}\hfill{\scriptsize (operation)}}\\


 The eigenvalues of the matrix \mbox{\texttt{\mdseries\slshape A}} over the field \mbox{\texttt{\mdseries\slshape F}}. }

 

\subsection{\textcolor{Chapter }{Eigenspaces}}
\logpage{[ 24, 8, 4 ]}\nobreak
\hyperdef{L}{X7A6B047281B52FD7}{}
{\noindent\textcolor{FuncColor}{$\triangleright$\enspace\texttt{Eigenspaces({\mdseries\slshape F, A})\index{Eigenspaces@\texttt{Eigenspaces}}
\label{Eigenspaces}
}\hfill{\scriptsize (operation)}}\\


 The eigenspaces of the matrix \mbox{\texttt{\mdseries\slshape A}} over the field \mbox{\texttt{\mdseries\slshape F}}. }

 

\subsection{\textcolor{Chapter }{Eigenvectors}}
\logpage{[ 24, 8, 5 ]}\nobreak
\hyperdef{L}{X8506584579D4EA18}{}
{\noindent\textcolor{FuncColor}{$\triangleright$\enspace\texttt{Eigenvectors({\mdseries\slshape F, A})\index{Eigenvectors@\texttt{Eigenvectors}}
\label{Eigenvectors}
}\hfill{\scriptsize (operation)}}\\


 The eigenvectors of the matrix \mbox{\texttt{\mdseries\slshape A}} over the field \mbox{\texttt{\mdseries\slshape F}}. }

 }

  
\section{\textcolor{Chapter }{Elementary Divisors}}\label{Elementary Divisors}
\logpage{[ 24, 9, 0 ]}
\hyperdef{L}{X7E5405D085661B29}{}
{
  See also chapter \ref{Integral matrices and lattices}. 

\subsection{\textcolor{Chapter }{ElementaryDivisorsMat}}
\logpage{[ 24, 9, 1 ]}\nobreak
\hyperdef{L}{X7AC4D74F81908109}{}
{\noindent\textcolor{FuncColor}{$\triangleright$\enspace\texttt{ElementaryDivisorsMat({\mdseries\slshape [ring, ]mat})\index{ElementaryDivisorsMat@\texttt{ElementaryDivisorsMat}}
\label{ElementaryDivisorsMat}
}\hfill{\scriptsize (operation)}}\\
\noindent\textcolor{FuncColor}{$\triangleright$\enspace\texttt{ElementaryDivisorsMatDestructive({\mdseries\slshape ring, mat})\index{ElementaryDivisorsMatDestructive@\texttt{ElementaryDivisorsMatDestructive}}
\label{ElementaryDivisorsMatDestructive}
}\hfill{\scriptsize (function)}}\\


 returns a list of the elementary divisors, i.e., the unique $d$ with $d[i]$ divides $d[i+1]$ and \mbox{\texttt{\mdseries\slshape mat}} is equivalent to a diagonal matrix with the elements $d[i]$ on the diagonal. The operations are performed over the euclidean ring \mbox{\texttt{\mdseries\slshape ring}}, which must contain all matrix entries. For compatibility reasons it can be
omitted and defaults to the \texttt{DefaultRing} (\ref{DefaultRing:for ring elements}) of the matrix entries. 

 The function \texttt{ElementaryDivisorsMatDestructive} produces the same result but in the process may destroy the contents of \mbox{\texttt{\mdseries\slshape mat}}. 
\begin{Verbatim}[commandchars=!@|,fontsize=\small,frame=single,label=Example]
  !gapprompt@gap>| !gapinput@mat:=[[1,2,3],[4,5,6],[7,8,9]];;|
  !gapprompt@gap>| !gapinput@ElementaryDivisorsMat(mat);|
  [ 1, 3, 0 ]
  !gapprompt@gap>| !gapinput@x:=Indeterminate(Rationals,"x");;|
  !gapprompt@gap>| !gapinput@mat:=mat*One(x)-x*mat^0;|
  [ [ -x+1, 2, 3 ], [ 4, -x+5, 6 ], [ 7, 8, -x+9 ] ]
  !gapprompt@gap>| !gapinput@ElementaryDivisorsMat(PolynomialRing(Rationals,1),mat);|
  [ 1, 1, x^3-15*x^2-18*x ]
  !gapprompt@gap>| !gapinput@mat:=KroneckerProduct(CompanionMat((x-1)^2),|
  !gapprompt@>| !gapinput@                         CompanionMat((x^3-1)*(x-1)));;|
  !gapprompt@gap>| !gapinput@mat:=mat*One(x)-x*mat^0;|
  [ [ -x, 0, 0, 0, 0, 0, 0, 1 ], [ 0, -x, 0, 0, -1, 0, 0, -1 ],
    [ 0, 0, -x, 0, 0, -1, 0, 0 ], [ 0, 0, 0, -x, 0, 0, -1, -1 ],
    [ 0, 0, 0, -1, -x, 0, 0, -2 ], [ 1, 0, 0, 1, 2, -x, 0, 2 ],
    [ 0, 1, 0, 0, 0, 2, -x, 0 ], [ 0, 0, 1, 1, 0, 0, 2, -x+2 ] ]
  !gapprompt@gap>| !gapinput@ElementaryDivisorsMat(PolynomialRing(Rationals,1),mat);|
  [ 1, 1, 1, 1, 1, 1, x-1, x^7-x^6-2*x^4+2*x^3+x-1 ]
\end{Verbatim}
 }

 

\subsection{\textcolor{Chapter }{ElementaryDivisorsTransformationsMat}}
\logpage{[ 24, 9, 2 ]}\nobreak
\hyperdef{L}{X7AA1C9047B102204}{}
{\noindent\textcolor{FuncColor}{$\triangleright$\enspace\texttt{ElementaryDivisorsTransformationsMat({\mdseries\slshape [ring, ]mat})\index{ElementaryDivisorsTransformationsMat@\texttt{Elementary}\-\texttt{Divisors}\-\texttt{TransformationsMat}}
\label{ElementaryDivisorsTransformationsMat}
}\hfill{\scriptsize (operation)}}\\
\noindent\textcolor{FuncColor}{$\triangleright$\enspace\texttt{ElementaryDivisorsTransformationsMatDestructive({\mdseries\slshape ring, mat})\index{ElementaryDivisorsTransformationsMatDestructive@\texttt{Elementary}\-\texttt{Divisors}\-\texttt{Transformations}\-\texttt{Mat}\-\texttt{Destructive}}
\label{ElementaryDivisorsTransformationsMatDestructive}
}\hfill{\scriptsize (function)}}\\


 \texttt{ElementaryDivisorsTransformations}, in addition to the tasks done by \texttt{ElementaryDivisorsMat}, also calculates transforming matrices. It returns a record with components \texttt{normal} (a matrix $S$), \texttt{rowtrans} (a matrix $P$), and \texttt{coltrans} (a matrix $Q$) such that $P A Q = S$. The operations are performed over the euclidean ring \mbox{\texttt{\mdseries\slshape ring}}, which must contain all matrix entries. For compatibility reasons it can be
omitted and defaults to the \texttt{DefaultRing} (\ref{DefaultRing:for ring elements}) of the matrix entries. 

 The function \texttt{ElementaryDivisorsTransformationsMatDestructive} produces the same result but in the process destroys the contents of \mbox{\texttt{\mdseries\slshape mat}}. 
\begin{Verbatim}[commandchars=!@|,fontsize=\small,frame=single,label=Example]
  !gapprompt@gap>| !gapinput@mat:=KroneckerProduct(CompanionMat((x-1)^2),CompanionMat((x^3-1)*(x-1)));;|
  !gapprompt@gap>| !gapinput@mat:=mat*One(x)-x*mat^0;|
  [ [ -x, 0, 0, 0, 0, 0, 0, 1 ], [ 0, -x, 0, 0, -1, 0, 0, -1 ],
    [ 0, 0, -x, 0, 0, -1, 0, 0 ], [ 0, 0, 0, -x, 0, 0, -1, -1 ],
    [ 0, 0, 0, -1, -x, 0, 0, -2 ], [ 1, 0, 0, 1, 2, -x, 0, 2 ],
    [ 0, 1, 0, 0, 0, 2, -x, 0 ], [ 0, 0, 1, 1, 0, 0, 2, -x+2 ] ]
  !gapprompt@gap>| !gapinput@t:=ElementaryDivisorsTransformationsMat(PolynomialRing(Rationals,1),mat);|
  rec( coltrans := [ [ 0, 0, 0, 0, 0, 0, 1/6*x^2-7/9*x-1/18, -3*x^3-x^2-x-1 ],
        [ 0, 0, 0, 0, 0, 0, -1/6*x^2+x-1, 3*x^3-3*x^2 ],
        [ 0, 0, 0, 0, 0, 1, -1/18*x^4+1/3*x^3-1/3*x^2-1/9*x, x^5-x^4+2*x^2-2*x
           ], [ 0, 0, 0, 0, -1, 0, -1/9*x^3+1/2*x^2+1/9*x, 2*x^4+x^3+x^2+2*x ],
        [ 0, -1, 0, 0, 0, 0, -2/9*x^2+19/18*x, 4*x^3+x^2+x ],
        [ 0, 0, -1, 0, 0, -x, 1/18*x^5-1/3*x^4+1/3*x^3+1/9*x^2,
            -x^6+x^5-2*x^3+2*x^2 ],
        [ 0, 0, 0, -1, x, 0, 1/9*x^4-2/3*x^3+2/3*x^2+1/18*x,
            -2*x^5+2*x^4-x^2+x ],
        [ 1, 0, 0, 0, 0, 0, 1/6*x^3-7/9*x^2-1/18*x, -3*x^4-x^3-x^2-x ] ],
    normal := [ [ 1, 0, 0, 0, 0, 0, 0, 0 ], [ 0, 1, 0, 0, 0, 0, 0, 0 ],
        [ 0, 0, 1, 0, 0, 0, 0, 0 ], [ 0, 0, 0, 1, 0, 0, 0, 0 ],
        [ 0, 0, 0, 0, 1, 0, 0, 0 ], [ 0, 0, 0, 0, 0, 1, 0, 0 ],
        [ 0, 0, 0, 0, 0, 0, x-1, 0 ],
        [ 0, 0, 0, 0, 0, 0, 0, x^7-x^6-2*x^4+2*x^3+x-1 ] ],
    rowtrans := [ [ 1, 0, 0, 0, 0, 0, 0, 0 ], [ 1, 1, 0, 0, 0, 0, 0, 0 ],
        [ 0, 0, 1, 0, 0, 0, 0, 0 ], [ 1, 0, 0, 1, 0, 0, 0, 0 ],
        [ -x+2, -x, 0, 0, 1, 0, 0, 0 ],
        [ 2*x^2-4*x+2, 2*x^2-x, 0, 2, -2*x+1, 0, 0, 1 ],
        [ 3*x^3-6*x^2+3*x, 3*x^3-2*x^2, 2, 3*x, -3*x^2+2*x, 0, 1, 2*x ],
        [ 1/6*x^8-7/6*x^7+2*x^6-4/3*x^5+7/3*x^4-4*x^3+13/6*x^2-7/6*x+2,
            1/6*x^8-17/18*x^7+13/18*x^6-5/18*x^5+35/18*x^4-31/18*x^3+1/9*x^2-x+\
  2, 1/9*x^5-5/9*x^4+1/9*x^3-1/9*x^2+14/9*x-1/9,
            1/6*x^6-5/6*x^5+1/6*x^4-1/6*x^3+11/6*x^2-1/6*x,
            -1/6*x^7+17/18*x^6-13/18*x^5+5/18*x^4-35/18*x^3+31/18*x^2-1/9*x+1,
            1, 1/18*x^5-5/18*x^4+1/18*x^3-1/18*x^2+23/18*x-1/18,
            1/9*x^6-5/9*x^5+1/9*x^4-1/9*x^3+14/9*x^2-1/9*x ] ] )
  !gapprompt@gap>| !gapinput@t.rowtrans*mat*t.coltrans;|
  [ [ 1, 0, 0, 0, 0, 0, 0, 0 ], [ 0, 1, 0, 0, 0, 0, 0, 0 ],
    [ 0, 0, 1, 0, 0, 0, 0, 0 ], [ 0, 0, 0, 1, 0, 0, 0, 0 ],
    [ 0, 0, 0, 0, 1, 0, 0, 0 ], [ 0, 0, 0, 0, 0, 1, 0, 0 ],
    [ 0, 0, 0, 0, 0, 0, x-1, 0 ],
    [ 0, 0, 0, 0, 0, 0, 0, x^7-x^6-2*x^4+2*x^3+x-1 ] ]
\end{Verbatim}
 }

 

\subsection{\textcolor{Chapter }{DiagonalizeMat}}
\logpage{[ 24, 9, 3 ]}\nobreak
\hyperdef{L}{X85819D3F7A582180}{}
{\noindent\textcolor{FuncColor}{$\triangleright$\enspace\texttt{DiagonalizeMat({\mdseries\slshape ring, mat})\index{DiagonalizeMat@\texttt{DiagonalizeMat}}
\label{DiagonalizeMat}
}\hfill{\scriptsize (operation)}}\\


 brings the mutable matrix \mbox{\texttt{\mdseries\slshape mat}}, considered as a matrix over \mbox{\texttt{\mdseries\slshape ring}}, into diagonal form by elementary row and column operations. 
\begin{Verbatim}[commandchars=!@|,fontsize=\small,frame=single,label=Example]
  !gapprompt@gap>| !gapinput@m:=[[1,2],[2,1]];;|
  !gapprompt@gap>| !gapinput@DiagonalizeMat(Integers,m);m;|
  [ [ 1, 0 ], [ 0, 3 ] ]
\end{Verbatim}
 }

 }

  
\section{\textcolor{Chapter }{Echelonized Matrices}}\label{Echelonized Matrices}
\logpage{[ 24, 10, 0 ]}
\hyperdef{L}{X7CA6B51D7AE3172B}{}
{
  

\subsection{\textcolor{Chapter }{SemiEchelonMat}}
\logpage{[ 24, 10, 1 ]}\nobreak
\hyperdef{L}{X7D5D6BD07B7E981B}{}
{\noindent\textcolor{FuncColor}{$\triangleright$\enspace\texttt{SemiEchelonMat({\mdseries\slshape mat})\index{SemiEchelonMat@\texttt{SemiEchelonMat}}
\label{SemiEchelonMat}
}\hfill{\scriptsize (attribute)}}\\


 A matrix over a field $F$ is in semi\texttt{\symbol{45}}echelon form if the first nonzero element in
each row is the identity of $F$, and all values exactly below these pivots are the zero of $F$. 

 \texttt{SemiEchelonMat} returns a record that contains information about a
semi\texttt{\symbol{45}}echelonized form of the matrix \mbox{\texttt{\mdseries\slshape mat}}. 

 The components of this record are 

 
\begin{description}
\item[{\texttt{vectors}}]  list of row vectors, each with pivot element the identity of $F$, 
\item[{\texttt{heads}}]  list that contains at position \mbox{\texttt{\mdseries\slshape i}}, if nonzero, the number of the row for that the pivot element is in column \mbox{\texttt{\mdseries\slshape i}}. 
\end{description}
 }

 

\subsection{\textcolor{Chapter }{SemiEchelonMatDestructive}}
\logpage{[ 24, 10, 2 ]}\nobreak
\hyperdef{L}{X8251F6F57D346385}{}
{\noindent\textcolor{FuncColor}{$\triangleright$\enspace\texttt{SemiEchelonMatDestructive({\mdseries\slshape mat})\index{SemiEchelonMatDestructive@\texttt{SemiEchelonMatDestructive}}
\label{SemiEchelonMatDestructive}
}\hfill{\scriptsize (operation)}}\\


 This does the same as \texttt{SemiEchelonMat( \mbox{\texttt{\mdseries\slshape mat}} )}, except that it may (and probably will) destroy the matrix \mbox{\texttt{\mdseries\slshape mat}}. 
\begin{Verbatim}[commandchars=!@|,fontsize=\small,frame=single,label=Example]
  !gapprompt@gap>| !gapinput@mm:=[[1,2,3],[4,5,6],[7,8,9]];;|
  !gapprompt@gap>| !gapinput@SemiEchelonMatDestructive( mm );|
  rec( heads := [ 1, 2, 0 ], vectors := [ [ 1, 2, 3 ], [ 0, 1, 2 ] ] )
  !gapprompt@gap>| !gapinput@mm;|
  [ [ 1, 2, 3 ], [ 0, 1, 2 ], [ 0, 0, 0 ] ]
\end{Verbatim}
 }

 

\subsection{\textcolor{Chapter }{SemiEchelonMatTransformation}}
\logpage{[ 24, 10, 3 ]}\nobreak
\hyperdef{L}{X7EFD1DB5861A54F0}{}
{\noindent\textcolor{FuncColor}{$\triangleright$\enspace\texttt{SemiEchelonMatTransformation({\mdseries\slshape mat})\index{SemiEchelonMatTransformation@\texttt{SemiEchelonMatTransformation}}
\label{SemiEchelonMatTransformation}
}\hfill{\scriptsize (attribute)}}\\


 does the same as \texttt{SemiEchelonMat} (\ref{SemiEchelonMat}) but additionally stores the linear transformation $T$ performed on the matrix. The additional components of the result are 

 
\begin{description}
\item[{\texttt{coeffs}}]  a list of coefficients vectors of the \texttt{vectors} component, with respect to the rows of \mbox{\texttt{\mdseries\slshape mat}}, that is, \texttt{coeffs * mat} is the \texttt{vectors} component. 
\item[{\texttt{relations}}]  a list of basis vectors for the (left) null space of \mbox{\texttt{\mdseries\slshape mat}}. 
\end{description}
 
\begin{Verbatim}[commandchars=!@|,fontsize=\small,frame=single,label=Example]
  !gapprompt@gap>| !gapinput@SemiEchelonMatTransformation([[1,2,3],[0,0,1]]);|
  rec( coeffs := [ [ 1, 0 ], [ 0, 1 ] ], heads := [ 1, 0, 2 ],
    relations := [  ], vectors := [ [ 1, 2, 3 ], [ 0, 0, 1 ] ] )
\end{Verbatim}
 }

 

\subsection{\textcolor{Chapter }{SemiEchelonMats}}
\logpage{[ 24, 10, 4 ]}\nobreak
\hyperdef{L}{X827D7971800DB661}{}
{\noindent\textcolor{FuncColor}{$\triangleright$\enspace\texttt{SemiEchelonMats({\mdseries\slshape mats})\index{SemiEchelonMats@\texttt{SemiEchelonMats}}
\label{SemiEchelonMats}
}\hfill{\scriptsize (operation)}}\\


 A list of matrices over a field $F$ is in semi\texttt{\symbol{45}}echelon form if the list of row vectors obtained
on concatenating the rows of each matrix is a
semi\texttt{\symbol{45}}echelonized matrix (see \texttt{SemiEchelonMat} (\ref{SemiEchelonMat})). 

 \texttt{SemiEchelonMats} returns a record that contains information about a
semi\texttt{\symbol{45}}echelonized form of the list \mbox{\texttt{\mdseries\slshape mats}} of matrices. 

 The components of this record are 

 
\begin{description}
\item[{\texttt{vectors}}]  list of matrices, each with pivot element the identity of $F$, 
\item[{\texttt{heads}}]  matrix that contains at position [\mbox{\texttt{\mdseries\slshape i}},\mbox{\texttt{\mdseries\slshape j}}], if nonzero, the number of the matrix that has the pivot element in this
position 
\end{description}
 }

 

\subsection{\textcolor{Chapter }{SemiEchelonMatsDestructive}}
\logpage{[ 24, 10, 5 ]}\nobreak
\hyperdef{L}{X808F493B839BC7A6}{}
{\noindent\textcolor{FuncColor}{$\triangleright$\enspace\texttt{SemiEchelonMatsDestructive({\mdseries\slshape mats})\index{SemiEchelonMatsDestructive@\texttt{SemiEchelonMatsDestructive}}
\label{SemiEchelonMatsDestructive}
}\hfill{\scriptsize (operation)}}\\


 Does the same as \texttt{SemiEchelonMats} (\ref{SemiEchelonMats}), except that it may destroy its argument. Therefore the argument must be a
list of matrices that are mutable. }

 }

  
\section{\textcolor{Chapter }{Matrices as Basis of a Row Space}}\label{Matrices as Basis of a Row Space}
\logpage{[ 24, 11, 0 ]}
\hyperdef{L}{X86B0D4A886BC0C6E}{}
{
  See also chapter \ref{Integral matrices and lattices} 

\subsection{\textcolor{Chapter }{BaseMat}}
\logpage{[ 24, 11, 1 ]}\nobreak
\hyperdef{L}{X7AD6B5F5794D9E46}{}
{\noindent\textcolor{FuncColor}{$\triangleright$\enspace\texttt{BaseMat({\mdseries\slshape mat})\index{BaseMat@\texttt{BaseMat}}
\label{BaseMat}
}\hfill{\scriptsize (attribute)}}\\


 returns a basis for the row space generated by the rows of \mbox{\texttt{\mdseries\slshape mat}} in the form of an immutable matrix. }

 

\subsection{\textcolor{Chapter }{BaseMatDestructive}}
\logpage{[ 24, 11, 2 ]}\nobreak
\hyperdef{L}{X78B094597E382A5F}{}
{\noindent\textcolor{FuncColor}{$\triangleright$\enspace\texttt{BaseMatDestructive({\mdseries\slshape mat})\index{BaseMatDestructive@\texttt{BaseMatDestructive}}
\label{BaseMatDestructive}
}\hfill{\scriptsize (operation)}}\\


 Does the same as \texttt{BaseMat} (\ref{BaseMat}), with the difference that it may destroy the matrix \mbox{\texttt{\mdseries\slshape mat}}. The matrix \mbox{\texttt{\mdseries\slshape mat}} must be mutable. 
\begin{Verbatim}[commandchars=!@|,fontsize=\small,frame=single,label=Example]
  !gapprompt@gap>| !gapinput@mat:=[[1,2,3],[4,5,6],[7,8,9]];;|
  !gapprompt@gap>| !gapinput@BaseMat(mat);|
  [ [ 1, 2, 3 ], [ 0, 1, 2 ] ]
  !gapprompt@gap>| !gapinput@mm:= [[1,2,3],[4,5,6],[5,7,9]];;|
  !gapprompt@gap>| !gapinput@BaseMatDestructive( mm );|
  [ [ 1, 2, 3 ], [ 0, 1, 2 ] ]
  !gapprompt@gap>| !gapinput@mm;|
  [ [ 1, 2, 3 ], [ 0, 1, 2 ], [ 0, 0, 0 ] ]
\end{Verbatim}
 }

 

\subsection{\textcolor{Chapter }{BaseOrthogonalSpaceMat}}
\logpage{[ 24, 11, 3 ]}\nobreak
\hyperdef{L}{X78B94EFF87A455BE}{}
{\noindent\textcolor{FuncColor}{$\triangleright$\enspace\texttt{BaseOrthogonalSpaceMat({\mdseries\slshape mat})\index{BaseOrthogonalSpaceMat@\texttt{BaseOrthogonalSpaceMat}}
\label{BaseOrthogonalSpaceMat}
}\hfill{\scriptsize (attribute)}}\\


 Let $V$ be the row space generated by the rows of \mbox{\texttt{\mdseries\slshape mat}} (over any field that contains all entries of \mbox{\texttt{\mdseries\slshape mat}}). \texttt{BaseOrthogonalSpaceMat( \mbox{\texttt{\mdseries\slshape mat}} )} computes a base of the orthogonal space of $V$. 

 The rows of \mbox{\texttt{\mdseries\slshape mat}} need not be linearly independent. 

  }

 

\subsection{\textcolor{Chapter }{SumIntersectionMat}}
\logpage{[ 24, 11, 4 ]}\nobreak
\hyperdef{L}{X7AFF8BCF80C88B45}{}
{\noindent\textcolor{FuncColor}{$\triangleright$\enspace\texttt{SumIntersectionMat({\mdseries\slshape M1, M2})\index{SumIntersectionMat@\texttt{SumIntersectionMat}}
\label{SumIntersectionMat}
}\hfill{\scriptsize (operation)}}\\


 performs Zassenhaus' algorithm to compute bases for the sum and the
intersection of spaces generated by the rows of the matrices \mbox{\texttt{\mdseries\slshape M1}}, \mbox{\texttt{\mdseries\slshape M2}}. 

 returns a list of length 2, at first position a base of the sum, at second
position a base of the intersection. Both bases are in
semi\texttt{\symbol{45}}echelon form (see{\nobreakspace}\ref{Echelonized Matrices}). 
\begin{Verbatim}[commandchars=!@|,fontsize=\small,frame=single,label=Example]
  !gapprompt@gap>| !gapinput@SumIntersectionMat(mat,[[2,7,6],[5,9,4]]);|
  [ [ [ 1, 2, 3 ], [ 0, 1, 2 ], [ 0, 0, 1 ] ], [ [ 1, -3/4, -5/2 ] ] ]
\end{Verbatim}
 }

 

\subsection{\textcolor{Chapter }{BaseSteinitzVectors}}
\logpage{[ 24, 11, 5 ]}\nobreak
\hyperdef{L}{X8245D54F7AC532EB}{}
{\noindent\textcolor{FuncColor}{$\triangleright$\enspace\texttt{BaseSteinitzVectors({\mdseries\slshape bas, mat})\index{BaseSteinitzVectors@\texttt{BaseSteinitzVectors}}
\label{BaseSteinitzVectors}
}\hfill{\scriptsize (function)}}\\


 find vectors extending mat to a basis spanning the span of \mbox{\texttt{\mdseries\slshape bas}}. Both \mbox{\texttt{\mdseries\slshape bas}} and \mbox{\texttt{\mdseries\slshape mat}} must be matrices of full (row) rank. It returns a record with the following
components: 
\begin{description}
\item[{\texttt{subspace}}]  is a basis of the space spanned by \mbox{\texttt{\mdseries\slshape mat}} in upper triangular form with leading ones at all echelon steps and zeroes
above these ones. 
\item[{\texttt{factorspace}}]  is a list of extending vectors in upper triangular form. 
\item[{\texttt{factorzero}}]  is a zero vector. 
\item[{\texttt{heads}}]  is a list of integers which can be used to decompose vectors in the basis
vectors. The \mbox{\texttt{\mdseries\slshape i}}th entry indicating the vector that gives an echelon step at position \mbox{\texttt{\mdseries\slshape i}}. A negative number indicates an echelon step in the subspace, a positive
number an echelon step in the complement, the absolute value gives the
position of the vector in the lists \texttt{subspace} and \texttt{factorspace}. 
\end{description}
 
\begin{Verbatim}[commandchars=!@|,fontsize=\small,frame=single,label=Example]
  !gapprompt@gap>| !gapinput@BaseSteinitzVectors(IdentityMat(3,1),[[11,13,15]]);|
  rec( factorspace := [ [ 0, 1, 15/13 ], [ 0, 0, 1 ] ],
    factorzero := [ 0, 0, 0 ], heads := [ -1, 1, 2 ],
    subspace := [ [ 1, 13/11, 15/11 ] ] )
\end{Verbatim}
 }

 }

  
\section{\textcolor{Chapter }{Triangular Matrices}}\label{Triangular Matrices}
\logpage{[ 24, 12, 0 ]}
\hyperdef{L}{X79D5E53685F0FBEE}{}
{
  

\subsection{\textcolor{Chapter }{DiagonalOfMatrix}}
\logpage{[ 24, 12, 1 ]}\nobreak
\hyperdef{L}{X7A9139D686ACB7D8}{}
{\noindent\textcolor{FuncColor}{$\triangleright$\enspace\texttt{DiagonalOfMatrix({\mdseries\slshape mat})\index{DiagonalOfMatrix@\texttt{DiagonalOfMatrix}}
\label{DiagonalOfMatrix}
}\hfill{\scriptsize (function)}}\\
\noindent\textcolor{FuncColor}{$\triangleright$\enspace\texttt{DiagonalOfMat({\mdseries\slshape mat})\index{DiagonalOfMat@\texttt{DiagonalOfMat}}
\label{DiagonalOfMat}
}\hfill{\scriptsize (function)}}\\


 return the diagonal of the matrix \mbox{\texttt{\mdseries\slshape mat}}. If \mbox{\texttt{\mdseries\slshape mat}} is not a square matrix, then the result has the same length as the rows of \mbox{\texttt{\mdseries\slshape mat}}, and is padded with zeros if \mbox{\texttt{\mdseries\slshape mat}} has fewer rows than columns. 
\begin{Verbatim}[commandchars=!@|,fontsize=\small,frame=single,label=Example]
  !gapprompt@gap>| !gapinput@DiagonalOfMatrix( [ [ 1, 2, 3 ], [ 4, 5, 6 ] ] );|
  [ 1, 5, 0 ]
\end{Verbatim}
 }

 

\subsection{\textcolor{Chapter }{UpperSubdiagonal}}
\logpage{[ 24, 12, 2 ]}\nobreak
\hyperdef{L}{X84A78C057F9DAE5E}{}
{\noindent\textcolor{FuncColor}{$\triangleright$\enspace\texttt{UpperSubdiagonal({\mdseries\slshape mat, pos})\index{UpperSubdiagonal@\texttt{UpperSubdiagonal}}
\label{UpperSubdiagonal}
}\hfill{\scriptsize (operation)}}\\


 returns a mutable list containing the entries of the \mbox{\texttt{\mdseries\slshape pos}}th upper subdiagonal of the matrix \mbox{\texttt{\mdseries\slshape mat}}. 
\begin{Verbatim}[commandchars=!@|,fontsize=\small,frame=single,label=Example]
  !gapprompt@gap>| !gapinput@UpperSubdiagonal( [ [ 1, 2, 3 ], [ 4, 5, 6 ], [ 7, 8, 9 ] ], 1 );|
  [ 2, 6 ]
  !gapprompt@gap>| !gapinput@UpperSubdiagonal( [ [ 1, 2 ], [ 3, 4 ], [ 5, 6 ] ], 1 );|
  [ 2 ]
  !gapprompt@gap>| !gapinput@UpperSubdiagonal( [ [ 1, 2, 3, 4 ], [ 5, 6, 7, 8 ] ], 1 );|
  [ 2, 7 ]
\end{Verbatim}
 }

 

\subsection{\textcolor{Chapter }{DepthOfUpperTriangularMatrix}}
\logpage{[ 24, 12, 3 ]}\nobreak
\hyperdef{L}{X84D74DEA798A9094}{}
{\noindent\textcolor{FuncColor}{$\triangleright$\enspace\texttt{DepthOfUpperTriangularMatrix({\mdseries\slshape mat})\index{DepthOfUpperTriangularMatrix@\texttt{DepthOfUpperTriangularMatrix}}
\label{DepthOfUpperTriangularMatrix}
}\hfill{\scriptsize (attribute)}}\\


 If \mbox{\texttt{\mdseries\slshape mat}} is an upper triangular matrix this attribute returns the index of the first
nonzero diagonal. 
\begin{Verbatim}[commandchars=!@|,fontsize=\small,frame=single,label=Example]
  !gapprompt@gap>| !gapinput@DepthOfUpperTriangularMatrix([[0,1,2],[0,0,1],[0,0,0]]);|
  1
  !gapprompt@gap>| !gapinput@DepthOfUpperTriangularMatrix([[0,0,2],[0,0,0],[0,0,0]]);|
  2
\end{Verbatim}
 }

 }

  
\section{\textcolor{Chapter }{Matrices as Linear Mappings}}\label{Matrices as Linear Mappings}
\logpage{[ 24, 13, 0 ]}
\hyperdef{L}{X85B403857F2855F7}{}
{
  

\subsection{\textcolor{Chapter }{CharacteristicPolynomial}}
\logpage{[ 24, 13, 1 ]}\nobreak
\hyperdef{L}{X87FA0A727CDB060B}{}
{\noindent\textcolor{FuncColor}{$\triangleright$\enspace\texttt{CharacteristicPolynomial({\mdseries\slshape [F, E, ]mat[, ind]})\index{CharacteristicPolynomial@\texttt{CharacteristicPolynomial}}
\label{CharacteristicPolynomial}
}\hfill{\scriptsize (attribute)}}\\


 For a square matrix \mbox{\texttt{\mdseries\slshape mat}}, \texttt{CharacteristicPolynomial} returns the \emph{characteristic polynomial} of \mbox{\texttt{\mdseries\slshape mat}}, that is, the \texttt{StandardAssociate} (\ref{StandardAssociate}) of the determinant of the matrix $\mbox{\texttt{\mdseries\slshape mat}} - X \cdot I$, where $X$ is an indeterminate and $I$ is the appropriate identity matrix. 

 If fields \mbox{\texttt{\mdseries\slshape F}} and \mbox{\texttt{\mdseries\slshape E}} are given, then \mbox{\texttt{\mdseries\slshape F}} must be a subfield of \mbox{\texttt{\mdseries\slshape E}}, and \mbox{\texttt{\mdseries\slshape mat}} must have entries in \mbox{\texttt{\mdseries\slshape E}}. Then \texttt{CharacteristicPolynomial} returns the characteristic polynomial of the \mbox{\texttt{\mdseries\slshape F}}\texttt{\symbol{45}}linear mapping induced by \mbox{\texttt{\mdseries\slshape mat}} on the underlying \mbox{\texttt{\mdseries\slshape E}}\texttt{\symbol{45}}vector space of \mbox{\texttt{\mdseries\slshape mat}}. In this case, the characteristic polynomial is computed using \texttt{BlownUpMat} (\ref{BlownUpMat}) for the field extension of $E/F$ generated by the default field. Thus, if $F = E$, the result is the same as for the one argument version. 

 The returned polynomials are expressed in the indeterminate number \mbox{\texttt{\mdseries\slshape ind}}. If \mbox{\texttt{\mdseries\slshape ind}} is not given, it defaults to $1$. 

 \texttt{CharacteristicPolynomial(\mbox{\texttt{\mdseries\slshape F}}, \mbox{\texttt{\mdseries\slshape E}}, \mbox{\texttt{\mdseries\slshape mat}})} is a multiple of the minimal polynomial \texttt{MinimalPolynomial(\mbox{\texttt{\mdseries\slshape F}}, \mbox{\texttt{\mdseries\slshape mat}})} (see{\nobreakspace}\texttt{MinimalPolynomial} (\ref{MinimalPolynomial})). 

 Note that, up to \textsf{GAP} version 4.4.6, \texttt{CharacteristicPolynomial} only allowed to specify one field (corresponding to \mbox{\texttt{\mdseries\slshape F}}) as an argument. That usage has been disabled because its definition turned
out to be ambiguous and may have lead to unexpected results. (To ensure
backward compatibility, it is still possible to use the old form if \mbox{\texttt{\mdseries\slshape F}} contains the default field of the matrix, see{\nobreakspace}\texttt{DefaultFieldOfMatrix} (\ref{DefaultFieldOfMatrix}), but this feature will disappear in future versions of \textsf{GAP}.) 
\begin{Verbatim}[commandchars=!@|,fontsize=\small,frame=single,label=Example]
  !gapprompt@gap>| !gapinput@CharacteristicPolynomial( [ [ 1, 1 ], [ 0, 1 ] ] );|
  x^2-2*x+1
  !gapprompt@gap>| !gapinput@mat := [[0,1],[E(4)-1,E(4)]];;|
  !gapprompt@gap>| !gapinput@CharacteristicPolynomial( mat );|
  x^2+(-E(4))*x+(1-E(4))
  !gapprompt@gap>| !gapinput@CharacteristicPolynomial( Rationals, CF(4), mat );|
  x^4+3*x^2+2*x+2
  !gapprompt@gap>| !gapinput@mat:= [ [ E(4), 1 ], [ 0, -E(4) ] ];;|
  !gapprompt@gap>| !gapinput@CharacteristicPolynomial( mat );|
  x^2+1
  !gapprompt@gap>| !gapinput@CharacteristicPolynomial( Rationals, CF(4), mat );|
  x^4+2*x^2+1
\end{Verbatim}
 }

 

\subsection{\textcolor{Chapter }{RationalCanonicalFormTransform}}
\logpage{[ 24, 13, 2 ]}\nobreak
\hyperdef{L}{X7B52560C792C1A0F}{}
{\noindent\textcolor{FuncColor}{$\triangleright$\enspace\texttt{RationalCanonicalFormTransform({\mdseries\slshape mat})\index{RationalCanonicalFormTransform@\texttt{RationalCanonicalFormTransform}}
\label{RationalCanonicalFormTransform}
}\hfill{\scriptsize (function)}}\\


 \index{Frobenius Normal Form} For a matrix \texttt{A}, return a matrix \texttt{P} such that $A^{P}$ is in rational canonical form (also called Frobenius normal form). The
algorithm used is the basic textbook version and thus not of optimal
complexity. 
\begin{Verbatim}[commandchars=!@|,fontsize=\small,frame=single,label=Example]
  !gapprompt@gap>| !gapinput@aa:=[[0,-8,12,40,-36,4,0,59,15,-9],[-2,-2,-2,6,-11,1,-1,10,1,0],|
  !gapprompt@>| !gapinput@[1,5,0,-6,12,-2,0,-12,-4,2],[0,0,0,2,0,0,0,7,0,0],|
  !gapprompt@>| !gapinput@[0,2,-3,-7,8,-1,0,-7,-3,2],[-5,-4,-6,18,-30,2,-2,35,5,-1],|
  !gapprompt@>| !gapinput@[-1,-6,6,20,-28,3,0,24,10,-6],[0,0,0,-1,0,0,0,-3,0,0],|
  !gapprompt@>| !gapinput@[0,0,-1,-2,-2,0,-1,-7,0,0],[0,-8,9,21,-36,4,-2,12,12,-8]];;|
  !gapprompt@gap>| !gapinput@t:=RationalCanonicalFormTransform(aa);;|
  !gapprompt@gap>| !gapinput@aa^t;|
  [ [ 0, 0, 0, 1, 0, 0, 0, 0, 0, 0 ], [ 1, 0, 0, 0, 0, 0, 0, 0, 0, 0 ],
    [ 0, 1, 0, 0, 0, 0, 0, 0, 0, 0 ], [ 0, 0, 1, 0, 0, 0, 0, 0, 0, 0 ],
    [ 0, 0, 0, 0, 0, 0, 0, 0, 0, 1 ], [ 0, 0, 0, 0, 1, 0, 0, 0, 0, 1 ],
    [ 0, 0, 0, 0, 0, 1, 0, 0, 0, 1 ], [ 0, 0, 0, 0, 0, 0, 1, 0, 0, 0 ],
    [ 0, 0, 0, 0, 0, 0, 0, 1, 0, -1 ], [ 0, 0, 0, 0, 0, 0, 0, 0, 1, -1 ] ]
\end{Verbatim}
 }

 

\subsection{\textcolor{Chapter }{JordanDecomposition}}
\logpage{[ 24, 13, 3 ]}\nobreak
\hyperdef{L}{X83F55D4E79BA5D1B}{}
{\noindent\textcolor{FuncColor}{$\triangleright$\enspace\texttt{JordanDecomposition({\mdseries\slshape mat})\index{JordanDecomposition@\texttt{JordanDecomposition}}
\label{JordanDecomposition}
}\hfill{\scriptsize (attribute)}}\\


 \texttt{JordanDecomposition( \mbox{\texttt{\mdseries\slshape mat }} )} returns a list \texttt{[S,N]} such that \texttt{S} is a semisimple matrix and \texttt{N} is nilpotent. Furthermore, \texttt{S} and \texttt{N} commute and \texttt{\mbox{\texttt{\mdseries\slshape mat}}=S+N}. 
\begin{Verbatim}[commandchars=!@|,fontsize=\small,frame=single,label=Example]
  !gapprompt@gap>| !gapinput@mat:=[[1,2,3],[4,5,6],[7,8,9]];;|
  !gapprompt@gap>| !gapinput@JordanDecomposition(mat);|
  [ [ [ 1, 2, 3 ], [ 4, 5, 6 ], [ 7, 8, 9 ] ],
    [ [ 0, 0, 0 ], [ 0, 0, 0 ], [ 0, 0, 0 ] ] ]
\end{Verbatim}
 }

 

\subsection{\textcolor{Chapter }{BlownUpMat}}
\logpage{[ 24, 13, 4 ]}\nobreak
\hyperdef{L}{X85923C107A4569D0}{}
{\noindent\textcolor{FuncColor}{$\triangleright$\enspace\texttt{BlownUpMat({\mdseries\slshape B, mat})\index{BlownUpMat@\texttt{BlownUpMat}}
\label{BlownUpMat}
}\hfill{\scriptsize (function)}}\\


 Let \mbox{\texttt{\mdseries\slshape B}} be a basis of a field extension $F / K$, and \mbox{\texttt{\mdseries\slshape mat}} a matrix whose entries are all in $F$. (This is not checked.) \texttt{BlownUpMat} returns a matrix over $K$ that is obtained by replacing each entry of \mbox{\texttt{\mdseries\slshape mat}} by its regular representation w.r.t.{\nobreakspace}\mbox{\texttt{\mdseries\slshape B}}. 

 More precisely, regard \mbox{\texttt{\mdseries\slshape mat}} as the matrix of a linear transformation on the row space $F^n$ w.r.t.{\nobreakspace}the $F$\texttt{\symbol{45}}basis with vectors $(v_1, \ldots, v_n)$ and suppose that the basis \mbox{\texttt{\mdseries\slshape B}} consists of the vectors $(b_1, \ldots, b_m)$; then the returned matrix is the matrix of the linear transformation on the
row space $K^{mn}$ w.r.t.{\nobreakspace}the $K$\texttt{\symbol{45}}basis whose vectors are $(b_1 v_1, \ldots b_m v_1, \ldots, b_m v_n)$. 

 Note that the linear transformations act on \emph{row} vectors, i.e., each row of the matrix is a concatenation of vectors of \mbox{\texttt{\mdseries\slshape B}}\texttt{\symbol{45}}coefficients. }

 

\subsection{\textcolor{Chapter }{BlownUpVector}}
\logpage{[ 24, 13, 5 ]}\nobreak
\hyperdef{L}{X82AC277D84EC5749}{}
{\noindent\textcolor{FuncColor}{$\triangleright$\enspace\texttt{BlownUpVector({\mdseries\slshape B, vector})\index{BlownUpVector@\texttt{BlownUpVector}}
\label{BlownUpVector}
}\hfill{\scriptsize (function)}}\\


 Let \mbox{\texttt{\mdseries\slshape B}} be a basis of a field extension $F / K$, and \mbox{\texttt{\mdseries\slshape vector}} a row vector whose entries are all in $F$. \texttt{BlownUpVector} returns a row vector over $K$ that is obtained by replacing each entry of \mbox{\texttt{\mdseries\slshape vector}} by its coefficients w.r.t.{\nobreakspace}\mbox{\texttt{\mdseries\slshape B}}. 

 So \texttt{BlownUpVector} and \texttt{BlownUpMat} (\ref{BlownUpMat}) are compatible in the sense that for a matrix \mbox{\texttt{\mdseries\slshape mat}} over $F$, \texttt{BlownUpVector( \mbox{\texttt{\mdseries\slshape B}}, \mbox{\texttt{\mdseries\slshape mat}} * \mbox{\texttt{\mdseries\slshape vector}} )} is equal to \texttt{BlownUpMat( \mbox{\texttt{\mdseries\slshape B}}, \mbox{\texttt{\mdseries\slshape mat}} ) * BlownUpVector( \mbox{\texttt{\mdseries\slshape B}}, \mbox{\texttt{\mdseries\slshape vector}} )}. 
\begin{Verbatim}[commandchars=!@|,fontsize=\small,frame=single,label=Example]
  !gapprompt@gap>| !gapinput@B:= Basis( CF(4), [ 1, E(4) ] );;|
  !gapprompt@gap>| !gapinput@mat:= [ [ 1, E(4) ], [ 0, 1 ] ];;  vec:= [ 1, E(4) ];;|
  !gapprompt@gap>| !gapinput@bmat:= BlownUpMat( B, mat );;  bvec:= BlownUpVector( B, vec );;|
  !gapprompt@gap>| !gapinput@Display( bmat );  bvec;|
  [ [   1,   0,   0,   1 ],
    [   0,   1,  -1,   0 ],
    [   0,   0,   1,   0 ],
    [   0,   0,   0,   1 ] ]
  [ 1, 0, 0, 1 ]
  !gapprompt@gap>| !gapinput@bvec * bmat = BlownUpVector( B, vec * mat );|
  true
\end{Verbatim}
 }

 

\subsection{\textcolor{Chapter }{CompanionMatrix}}
\logpage{[ 24, 13, 6 ]}\nobreak
\hyperdef{L}{X7E06762479A00DF4}{}
{\noindent\textcolor{FuncColor}{$\triangleright$\enspace\texttt{CompanionMatrix({\mdseries\slshape poly})\index{CompanionMatrix@\texttt{CompanionMatrix}}
\label{CompanionMatrix}
}\hfill{\scriptsize (operation)}}\\
\noindent\textcolor{FuncColor}{$\triangleright$\enspace\texttt{CompanionMat({\mdseries\slshape poly})\index{CompanionMat@\texttt{CompanionMat}}
\label{CompanionMat}
}\hfill{\scriptsize (operation)}}\\


 Return a fully mutable matrix that is the companion matrix of the polynomial \mbox{\texttt{\mdseries\slshape poly}}. The negatives of the coefficients of \mbox{\texttt{\mdseries\slshape poly}} appear in the last column of the result. 

 The companion matrix of \mbox{\texttt{\mdseries\slshape poly}} has \mbox{\texttt{\mdseries\slshape poly}} as its minimal polynomial (see \texttt{MinimalPolynomial} (\ref{MinimalPolynomial})) and as its characteristic polynomial (see \texttt{CharacteristicPolynomial} (\ref{CharacteristicPolynomial})). 

 
\begin{Verbatim}[commandchars=!@|,fontsize=\small,frame=single,label=Example]
  !gapprompt@gap>| !gapinput@x:= X( Rationals );;  pol:= x^3 + x^2 + 2*x + 3;;|
  !gapprompt@gap>| !gapinput@M:= CompanionMatrix( pol );;|
  !gapprompt@gap>| !gapinput@Display( M );|
  [ [   0,   0,  -3 ],
    [   1,   0,  -2 ],
    [   0,   1,  -1 ] ]
  !gapprompt@gap>| !gapinput@MinimalPolynomial( M ) = pol;|
  true
\end{Verbatim}
 }

 }

  
\section{\textcolor{Chapter }{Matrices over Finite Fields}}\label{Matrices over Finite Fields}
\logpage{[ 24, 14, 0 ]}
\hyperdef{L}{X873822B6830CE367}{}
{
  Just as for row vectors, (see section \ref{Row Vectors over Finite Fields}), \textsf{GAP} has a special representation for matrices over small finite fields. 

 To be eligible to be represented in this way, each row of a matrix must be
able to be represented as a compact row vector of the same length over \emph{the same} finite field. 

 
\begin{Verbatim}[commandchars=!@|,fontsize=\small,frame=single,label=Example]
  !gapprompt@gap>| !gapinput@v := Z(2)*[1,0,0,1,1];|
  [ Z(2)^0, 0*Z(2), 0*Z(2), Z(2)^0, Z(2)^0 ]
  !gapprompt@gap>| !gapinput@ConvertToVectorRep(v,2);|
  2
  !gapprompt@gap>| !gapinput@v;|
  <a GF2 vector of length 5>
  !gapprompt@gap>| !gapinput@m := [v];; ConvertToMatrixRep(m,GF(2));; m;|
  <a 1x5 matrix over GF2>
  !gapprompt@gap>| !gapinput@m := [v,v];; ConvertToMatrixRep(m,GF(2));; m;|
  <a 2x5 matrix over GF2>
  !gapprompt@gap>| !gapinput@m := [v,v,v];; ConvertToMatrixRep(m,GF(2));; m;|
  <a 3x5 matrix over GF2>
  !gapprompt@gap>| !gapinput@v := Z(3)*[1..8];|
  [ Z(3), Z(3)^0, 0*Z(3), Z(3), Z(3)^0, 0*Z(3), Z(3), Z(3)^0 ]
  !gapprompt@gap>| !gapinput@ConvertToVectorRep(v);|
  3
  !gapprompt@gap>| !gapinput@m := [v];; ConvertToMatrixRep(m,GF(3));; m;|
  [ [ Z(3), Z(3)^0, 0*Z(3), Z(3), Z(3)^0, 0*Z(3), Z(3), Z(3)^0 ] ]
  !gapprompt@gap>| !gapinput@RepresentationsOfObject(m);|
  [ "IsPositionalObjectRep", "Is8BitMatrixRep" ]
  !gapprompt@gap>| !gapinput@m := [v,v,v,v];; ConvertToMatrixRep(m,GF(3));; m;|
  < mutable compressed matrix 4x8 over GF(3) >
\end{Verbatim}
 

 All compressed matrices over GF(2) are viewed as \texttt{{\textless}a \mbox{\texttt{\mdseries\slshape n}}x\mbox{\texttt{\mdseries\slshape m}} matrix over GF2{\textgreater}}, while over fields GF(q) for q between 3 and 256, matrices with 25 or more
entries are viewed in this way, and smaller ones as lists of lists. 

 Matrices can be converted to this special representation via the following
functions. 

 Note that the main advantage of this special representation of matrices is in
low dimensions, where various overheads can be reduced. In higher dimensions,
a list of compressed vectors will be almost as fast. Note also that list
access and assignment will be somewhat slower for compressed matrices than for
plain lists. 

 In order to form a row of a compressed matrix a vector must accept certain
restrictions. Specifically, it cannot change its length or change the field
over which it is compressed. The main consequences of this are: that only
elements of the appropriate field can be assigned to entries of the vector,
and only to positions between 1 and the original length; that the vector
cannot be shared between two matrices compressed over different fields. 

 This is enforced by the filter \texttt{IsLockedRepresentationVector}. When a vector becomes part of a compressed matrix, this filter is set for
it. Assignment, \texttt{Unbind} (\ref{Unbind:unbind a list entry}), \texttt{ConvertToVectorRep} (\ref{ConvertToVectorRep:for a list (and a field)}) and \texttt{ConvertToMatrixRep} (\ref{ConvertToMatrixRep:for a list (and a field)}) are all prevented from altering a vector with this filter. 

 
\begin{Verbatim}[commandchars=!@|,fontsize=\small,frame=single,label=Example]
  !gapprompt@gap>| !gapinput@v := [Z(2),Z(2)];; ConvertToVectorRep(v,GF(2));; v;|
  <a GF2 vector of length 2>
  !gapprompt@gap>| !gapinput@m := [v,v];|
  [ <a GF2 vector of length 2>, <a GF2 vector of length 2> ]
  !gapprompt@gap>| !gapinput@ConvertToMatrixRep(m,GF(2));|
  2
  !gapprompt@gap>| !gapinput@m2 := [m[1], [Z(4),Z(4)]]; # now try and mix in some GF(4)|
  [ <a GF2 vector of length 2>, [ Z(2^2), Z(2^2) ] ]
  !gapprompt@gap>| !gapinput@ConvertToMatrixRep(m2); # but m2[1] is locked|
  #I  ConvertToVectorRep: locked vector not converted to different field
  fail
  !gapprompt@gap>| !gapinput@m2 := [ShallowCopy(m[1]), [Z(4),Z(4)]]; # a fresh copy of row 1|
  [ <a GF2 vector of length 2>, [ Z(2^2), Z(2^2) ] ]
  !gapprompt@gap>| !gapinput@ConvertToMatrixRep(m2); # now it works|
  4
  !gapprompt@gap>| !gapinput@m2;|
  [ [ Z(2)^0, Z(2)^0 ], [ Z(2^2), Z(2^2) ] ]
  !gapprompt@gap>| !gapinput@RepresentationsOfObject(m2);|
  [ "IsPositionalObjectRep", "Is8BitMatrixRep" ]
\end{Verbatim}
 

 Arithmetic operations (see{\nobreakspace}\ref{Arithmetic for Lists} and the following sections) preserve the compression status of matrices in the
sense that if all arguments are compressed matrices written over the same
field and the result is a matrix then also the result is a compressed matrix
written over this field. 

 There are also two operations that are only available for matrices written
over finite fields. 

\subsection{\textcolor{Chapter }{ImmutableMatrix}}
\logpage{[ 24, 14, 1 ]}\nobreak
\hyperdef{L}{X7DED2522828B6C30}{}
{\noindent\textcolor{FuncColor}{$\triangleright$\enspace\texttt{ImmutableMatrix({\mdseries\slshape field, matrix[, change]})\index{ImmutableMatrix@\texttt{ImmutableMatrix}}
\label{ImmutableMatrix}
}\hfill{\scriptsize (operation)}}\\


 Let \mbox{\texttt{\mdseries\slshape matrix}} be an object for which either \texttt{IsMatrix} (\ref{IsMatrix}) or \texttt{IsMatrixObj} (\ref{IsMatrixObj}) returns \texttt{true}. In the former case, \mbox{\texttt{\mdseries\slshape matrix}} is a list of lists, and \texttt{ImmutableMatrix} returns an immutable object for which \texttt{IsMatrix} (\ref{IsMatrix}) returns \texttt{true} (in particular again a list of lists), which is equal to \mbox{\texttt{\mdseries\slshape matrix}}, and which is in the optimal (concerning space and runtime) representation
for matrices defined over \mbox{\texttt{\mdseries\slshape field}}, provided that the entries of \mbox{\texttt{\mdseries\slshape matrix}} lie in \mbox{\texttt{\mdseries\slshape field}}. In the latter case, \texttt{ImmutableMatrix} returns an immutable object that is equal to the result of \texttt{ChangedBaseDomain} (\ref{ChangedBaseDomain:for a matrix object}) when this is called with \mbox{\texttt{\mdseries\slshape matrix}} and \mbox{\texttt{\mdseries\slshape field}}. 

 This means that matrices obtained by several calls of \texttt{ImmutableMatrix} for the same \mbox{\texttt{\mdseries\slshape field}} are compatible for fast arithmetic without need for field conversion. 

 If the input matrix \mbox{\texttt{\mdseries\slshape matrix}} is in \texttt{IsMatrix} (\ref{IsMatrix}) then it or its rows might change their representation as a side effect of this
function. However, one cannot rely on this side effect. Also, if \mbox{\texttt{\mdseries\slshape matrix}} is already immutable and the result of \texttt{ImmutableMatrix} has the same internal representation as \mbox{\texttt{\mdseries\slshape matrix}}, the result is not necessarily \emph{identical} to \mbox{\texttt{\mdseries\slshape matrix}}. 

 If \mbox{\texttt{\mdseries\slshape change}} is \texttt{true}, \mbox{\texttt{\mdseries\slshape matrix}} or its rows (if there are subobjects that represent rows) may be changed to
become immutable; otherwise the rows of \mbox{\texttt{\mdseries\slshape matrix}} are copied first. }

 

\subsection{\textcolor{Chapter }{ConvertToMatrixRep (for a list (and a field))}}
\logpage{[ 24, 14, 2 ]}\nobreak
\hyperdef{L}{X8587A62F818AA0D6}{}
{\noindent\textcolor{FuncColor}{$\triangleright$\enspace\texttt{ConvertToMatrixRep({\mdseries\slshape list[, field]})\index{ConvertToMatrixRep@\texttt{ConvertToMatrixRep}!for a list (and a field)}
\label{ConvertToMatrixRep:for a list (and a field)}
}\hfill{\scriptsize (function)}}\\
\noindent\textcolor{FuncColor}{$\triangleright$\enspace\texttt{ConvertToMatrixRep({\mdseries\slshape list[, fieldsize]})\index{ConvertToMatrixRep@\texttt{ConvertToMatrixRep}!for a list (and a prime power)}
\label{ConvertToMatrixRep:for a list (and a prime power)}
}\hfill{\scriptsize (function)}}\\
\noindent\textcolor{FuncColor}{$\triangleright$\enspace\texttt{ConvertToMatrixRepNC({\mdseries\slshape list[, field]})\index{ConvertToMatrixRepNC@\texttt{ConvertToMatrixRepNC}!for a list (and a field)}
\label{ConvertToMatrixRepNC:for a list (and a field)}
}\hfill{\scriptsize (function)}}\\
\noindent\textcolor{FuncColor}{$\triangleright$\enspace\texttt{ConvertToMatrixRepNC({\mdseries\slshape list[, fieldsize]})\index{ConvertToMatrixRepNC@\texttt{ConvertToMatrixRepNC}!for a list (and a prime power)}
\label{ConvertToMatrixRepNC:for a list (and a prime power)}
}\hfill{\scriptsize (function)}}\\


 This function is more technical version of \texttt{ImmutableMatrix} (\ref{ImmutableMatrix}), which will never copy a matrix (or any rows of it) but may fail if it
encounters rows locked in the wrong representation, or various other more
technical problems. Most users should use \texttt{ImmutableMatrix} (\ref{ImmutableMatrix}) instead. The NC versions of the function do less checking of the argument and
may cause unpredictable results or crashes if given unsuitable arguments.
Called with one argument \mbox{\texttt{\mdseries\slshape list}}, \texttt{ConvertToMatrixRep} converts \mbox{\texttt{\mdseries\slshape list}} to an internal matrix representation if possible. 

 Called with a list \mbox{\texttt{\mdseries\slshape list}} and a finite field \mbox{\texttt{\mdseries\slshape field}}, \texttt{ConvertToMatrixRep} converts \mbox{\texttt{\mdseries\slshape list}} to an internal matrix representation appropriate for a matrix over \mbox{\texttt{\mdseries\slshape field}}. 

 Instead of a \mbox{\texttt{\mdseries\slshape field}} also its size \mbox{\texttt{\mdseries\slshape fieldsize}} may be given. 

 It is forbidden to call this function unless all elements of \mbox{\texttt{\mdseries\slshape list}} are row vectors with entries in the field \mbox{\texttt{\mdseries\slshape field}}. Violation of this condition can lead to unpredictable behaviour or a system
crash. (Setting the assertion level to at least 2 might catch some violations
before a crash, see{\nobreakspace}\texttt{SetAssertionLevel} (\ref{SetAssertionLevel}).) 

 \mbox{\texttt{\mdseries\slshape list}} may already be a compressed matrix. In this case, if no \mbox{\texttt{\mdseries\slshape field}} or \mbox{\texttt{\mdseries\slshape fieldsize}} is given, then nothing happens. 

 The return value is the size of the field over which the matrix ends up
written, if it is written in a compressed representation. 

 }

 

\subsection{\textcolor{Chapter }{ProjectiveOrder}}
\logpage{[ 24, 14, 3 ]}\nobreak
\hyperdef{L}{X84A76F7A7B4166BC}{}
{\noindent\textcolor{FuncColor}{$\triangleright$\enspace\texttt{ProjectiveOrder({\mdseries\slshape mat})\index{ProjectiveOrder@\texttt{ProjectiveOrder}}
\label{ProjectiveOrder}
}\hfill{\scriptsize (attribute)}}\\


 Returns an integer n and a finite field element e such that \mbox{\texttt{\mdseries\slshape A}}\texttt{\symbol{94}}n = eI. \mbox{\texttt{\mdseries\slshape mat}} must be a matrix defined over a finite field. 
\begin{Verbatim}[commandchars=!@|,fontsize=\small,frame=single,label=Example]
  !gapprompt@gap>| !gapinput@ProjectiveOrder([[1,4],[5,2]]*Z(11)^0);|
  [ 5, Z(11)^5 ]
\end{Verbatim}
 }

 

\subsection{\textcolor{Chapter }{SimultaneousEigenvalues}}
\logpage{[ 24, 14, 4 ]}\nobreak
\hyperdef{L}{X847ADC6779E33A1C}{}
{\noindent\textcolor{FuncColor}{$\triangleright$\enspace\texttt{SimultaneousEigenvalues({\mdseries\slshape matlist, expo})\index{SimultaneousEigenvalues@\texttt{SimultaneousEigenvalues}}
\label{SimultaneousEigenvalues}
}\hfill{\scriptsize (function)}}\\


 The matrices in \mbox{\texttt{\mdseries\slshape matlist}} must be matrices over GF(\mbox{\texttt{\mdseries\slshape q}}) for some prime \mbox{\texttt{\mdseries\slshape q}}. Together, they must generate an abelian p\texttt{\symbol{45}}group of
exponent \mbox{\texttt{\mdseries\slshape expo}}. Then the eigenvalues of \mbox{\texttt{\mdseries\slshape mat}} in the splitting field \texttt{GF(\mbox{\texttt{\mdseries\slshape q}}\texttt{\symbol{94}}\mbox{\texttt{\mdseries\slshape r}})} for some \mbox{\texttt{\mdseries\slshape r}} are powers of an element $\xi$ in the splitting field, which is of order \mbox{\texttt{\mdseries\slshape expo}}. \texttt{SimultaneousEigenvalues} returns a matrix of integers mod \mbox{\texttt{\mdseries\slshape expo}} $(a_{{i,j}})$, such that the power $\xi^{{a_{{i,j}}}}$ is an eigenvalue of the \mbox{\texttt{\mdseries\slshape i}}\texttt{\symbol{45}}th matrix in \mbox{\texttt{\mdseries\slshape matlist}} and the eigenspaces of the different matrices to the eigenvalues $\xi^{{a_{{i,j}}}}$ for fixed \mbox{\texttt{\mdseries\slshape j}} are equal. }

 }

  
\section{\textcolor{Chapter }{Inverse and Nullspace of an Integer Matrix Modulo an Ideal}}\label{Inverse and Nullspace of an Integer Matrix Modulo an Ideal}
\logpage{[ 24, 15, 0 ]}
\hyperdef{L}{X8593A5337D3B2C70}{}
{
  The following operations deal with matrices over a ring, but only care about
the residues of their entries modulo some ring element. In the case of the
integers and a prime number $p$, this is effectively computation in a matrix over the prime field in
characteristic $p$. 

\subsection{\textcolor{Chapter }{InverseMatMod}}
\logpage{[ 24, 15, 1 ]}\nobreak
\hyperdef{L}{X7D8D1E0E83C7F872}{}
{\noindent\textcolor{FuncColor}{$\triangleright$\enspace\texttt{InverseMatMod({\mdseries\slshape mat, obj})\index{InverseMatMod@\texttt{InverseMatMod}}
\label{InverseMatMod}
}\hfill{\scriptsize (operation)}}\\


 For a square matrix \mbox{\texttt{\mdseries\slshape mat}}, \texttt{InverseMatMod} returns a matrix \mbox{\texttt{\mdseries\slshape inv}} such that \texttt{\mbox{\texttt{\mdseries\slshape inv}} * \mbox{\texttt{\mdseries\slshape mat}}} is congruent to the identity matrix modulo \mbox{\texttt{\mdseries\slshape obj}}, if such a matrix exists, and \texttt{fail} otherwise. 
\begin{Verbatim}[commandchars=!@|,fontsize=\small,frame=single,label=Example]
  !gapprompt@gap>| !gapinput@mat:= [ [ 1, 2 ], [ 3, 4 ] ];;  inv:= InverseMatMod( mat, 5 );|
  [ [ 3, 1 ], [ 4, 2 ] ]
  !gapprompt@gap>| !gapinput@mat * inv;|
  [ [ 11, 5 ], [ 25, 11 ] ]
\end{Verbatim}
 }

 

\subsection{\textcolor{Chapter }{BasisNullspaceModN}}
\logpage{[ 24, 15, 2 ]}\nobreak
\hyperdef{L}{X7D7DF873826A7C20}{}
{\noindent\textcolor{FuncColor}{$\triangleright$\enspace\texttt{BasisNullspaceModN({\mdseries\slshape M, n})\index{BasisNullspaceModN@\texttt{BasisNullspaceModN}}
\label{BasisNullspaceModN}
}\hfill{\scriptsize (function)}}\\


 \mbox{\texttt{\mdseries\slshape M}} must be a matrix of integers and \mbox{\texttt{\mdseries\slshape n}} a positive integer. Then \texttt{BasisNullspaceModN} returns a set \mbox{\texttt{\mdseries\slshape B}} of vectors such that every vector \mbox{\texttt{\mdseries\slshape v}} of integer modulo \mbox{\texttt{\mdseries\slshape n}} satisfying \mbox{\texttt{\mdseries\slshape v}} \mbox{\texttt{\mdseries\slshape M}} = 0 modulo \mbox{\texttt{\mdseries\slshape n}} can be expressed by a Z\texttt{\symbol{45}}linear combination of elements of \mbox{\texttt{\mdseries\slshape B}}. }

 

\subsection{\textcolor{Chapter }{NullspaceModQ}}
\logpage{[ 24, 15, 3 ]}\nobreak
\hyperdef{L}{X86AE919983B242E2}{}
{\noindent\textcolor{FuncColor}{$\triangleright$\enspace\texttt{NullspaceModQ({\mdseries\slshape M, q})\index{NullspaceModQ@\texttt{NullspaceModQ}}
\label{NullspaceModQ}
}\hfill{\scriptsize (function)}}\\
\noindent\textcolor{FuncColor}{$\triangleright$\enspace\texttt{NullspaceModN({\mdseries\slshape M, n})\index{NullspaceModN@\texttt{NullspaceModN}}
\label{NullspaceModN}
}\hfill{\scriptsize (function)}}\\


 \mbox{\texttt{\mdseries\slshape M}} must be a matrix of integers and \mbox{\texttt{\mdseries\slshape n}} a positive integer. Then \texttt{NullspaceModN} returns the set of all vectors of integers modulo \mbox{\texttt{\mdseries\slshape n}}, which solve the homogeneous equation system \mbox{\texttt{\mdseries\slshape v}} \mbox{\texttt{\mdseries\slshape M}} = 0 modulo \mbox{\texttt{\mdseries\slshape n}}. 

 \texttt{NullspaceModQ} is a synonym for \texttt{NullspaceModN}. 
\begin{Verbatim}[commandchars=!@|,fontsize=\small,frame=single,label=Example]
  !gapprompt@gap>| !gapinput@NullspaceModN( [ [ 2 ] ], 8 );|
  [ [ 0 ], [ 4 ] ]
  !gapprompt@gap>| !gapinput@NullspaceModN( [ [ 2, 1 ], [ 0, 2 ] ], 6 );|
  [ [ 0, 0 ], [ 0, 3 ] ]
  !gapprompt@gap>| !gapinput@mat:= [ [ 1, 3 ], [ 1, 2 ], [ 1, 1 ] ];;|
  !gapprompt@gap>| !gapinput@NullspaceModN( mat, 5 );|
  [ [ 0, 0, 0 ], [ 1, 3, 1 ], [ 2, 1, 2 ], [ 3, 4, 3 ], [ 4, 2, 4 ] ]
\end{Verbatim}
 }

 }

  
\section{\textcolor{Chapter }{Special Multiplication Algorithms for Matrices over GF(2)}}\label{Special Multiplication Algorithms for Matrices over GF(2)}
\logpage{[ 24, 16, 0 ]}
\hyperdef{L}{X787DF5F07DC7D86E}{}
{
  When multiplying two compressed matrices $M$ and $N$ over GF(2) of dimensions $a \times b$ and $b \times c$, where $a$, $b$ and $c$ are all greater than or equal to 128, \textsf{GAP} by default uses a more sophisticated matrix multiplication algorithm, in which
linear combinations of groups of 8 rows of $M$ are remembered and re\texttt{\symbol{45}}used in constructing various rows of
the product. This is called level 8 grease. To optimise memory access
patterns, these combinations are stored for $(b+255)/256$ sets of 8 rows at once. This number is called the blocking level. 

 These levels of grease and blocking are found experimentally to give good
performance across a range of processors and matrix sizes, but other levels
may do even better in some cases. You can control the levels exactly using the
functions below. 

 We plan to include greased blocked matrix multiplication for other finite
fields, and greased blocked algorithms for inversion and other matrix
operations in a future release. 

 

\subsection{\textcolor{Chapter }{PROD{\textunderscore}GF2MAT{\textunderscore}GF2MAT{\textunderscore}SIMPLE}}
\logpage{[ 24, 16, 1 ]}\nobreak
\hyperdef{L}{X7C0C26027FAE0C83}{}
{\noindent\textcolor{FuncColor}{$\triangleright$\enspace\texttt{PROD{\textunderscore}GF2MAT{\textunderscore}GF2MAT{\textunderscore}SIMPLE({\mdseries\slshape m1, m2})\index{PROD{\textunderscore}GF2MAT{\textunderscore}GF2MAT{\textunderscore}SIMPLE@\texttt{PRO}\-\texttt{D{\textunderscore}}\-\texttt{G}\-\texttt{F2}\-\texttt{M}\-\texttt{A}\-\texttt{T{\textunderscore}}\-\texttt{G}\-\texttt{F2}\-\texttt{M}\-\texttt{A}\-\texttt{T{\textunderscore}}\-\texttt{S}\-\texttt{I}\-\texttt{MPLE}}
\label{PRODuScoreGF2MATuScoreGF2MATuScoreSIMPLE}
}\hfill{\scriptsize (function)}}\\


 This function performs the standard unblocked and ungreased matrix
multiplication for matrices of any size. }

 

 

\subsection{\textcolor{Chapter }{PROD{\textunderscore}GF2MAT{\textunderscore}GF2MAT{\textunderscore}ADVANCED}}
\logpage{[ 24, 16, 2 ]}\nobreak
\hyperdef{L}{X81965B7D7F45E088}{}
{\noindent\textcolor{FuncColor}{$\triangleright$\enspace\texttt{PROD{\textunderscore}GF2MAT{\textunderscore}GF2MAT{\textunderscore}ADVANCED({\mdseries\slshape m1, m2, g, b})\index{PROD{\textunderscore}GF2MAT{\textunderscore}GF2MAT{\textunderscore}ADVANCED@\texttt{PRO}\-\texttt{D{\textunderscore}}\-\texttt{G}\-\texttt{F2}\-\texttt{M}\-\texttt{A}\-\texttt{T{\textunderscore}}\-\texttt{G}\-\texttt{F2}\-\texttt{M}\-\texttt{A}\-\texttt{T{\textunderscore}}\-\texttt{A}\-\texttt{D}\-\texttt{V}\-\texttt{A}\-\texttt{NCED}}
\label{PRODuScoreGF2MATuScoreGF2MATuScoreADVANCED}
}\hfill{\scriptsize (function)}}\\


 This function computes the product of \mbox{\texttt{\mdseries\slshape m1}} and \mbox{\texttt{\mdseries\slshape m2}}, which must be compressed matrices over GF(2) of compatible dimensions, using
level \mbox{\texttt{\mdseries\slshape g}} grease and level \mbox{\texttt{\mdseries\slshape b}} blocking. }

 }

  
\section{\textcolor{Chapter }{Block Matrices}}\label{Block Matrices}
\logpage{[ 24, 17, 0 ]}
\hyperdef{L}{X7F8A71F38201A250}{}
{
  Block matrices are a special representation of matrices which can save a lot
of memory if large matrices have a block structure with lots of zero blocks. \textsf{GAP} uses the representation \texttt{IsBlockMatrixRep} to store block matrices. \index{IsBlockMatrixRep@IsBlockMatrixRep} 

\subsection{\textcolor{Chapter }{AsBlockMatrix}}
\logpage{[ 24, 17, 1 ]}\nobreak
\hyperdef{L}{X7D675B3C79CF8871}{}
{\noindent\textcolor{FuncColor}{$\triangleright$\enspace\texttt{AsBlockMatrix({\mdseries\slshape m, nrb, ncb})\index{AsBlockMatrix@\texttt{AsBlockMatrix}}
\label{AsBlockMatrix}
}\hfill{\scriptsize (function)}}\\


 returns a block matrix with \mbox{\texttt{\mdseries\slshape nrb}} row blocks and \mbox{\texttt{\mdseries\slshape ncb}} column blocks which is equal to the ordinary matrix \mbox{\texttt{\mdseries\slshape m}}. }

 

\subsection{\textcolor{Chapter }{BlockMatrix}}
\logpage{[ 24, 17, 2 ]}\nobreak
\hyperdef{L}{X8633538685551E7A}{}
{\noindent\textcolor{FuncColor}{$\triangleright$\enspace\texttt{BlockMatrix({\mdseries\slshape blocks, nrb, ncb[, rpb, cpb, zero]})\index{BlockMatrix@\texttt{BlockMatrix}}
\label{BlockMatrix}
}\hfill{\scriptsize (function)}}\\


 \texttt{BlockMatrix} returns an immutable matrix in the sparse representation \texttt{IsBlockMatrixRep}. The nonzero blocks are described by the list \mbox{\texttt{\mdseries\slshape blocks}} of triples $[ \mbox{\texttt{\mdseries\slshape i}}, \mbox{\texttt{\mdseries\slshape j}}, M(i,j) ]$ each consisting of a matrix $M(i,j)$ and its block coordinates in the block matrix to be constructed. All matrices $M(i,j)$ must have the same dimensions. As usual the first coordinate specifies the row
and the second one the column. The resulting matrix has \mbox{\texttt{\mdseries\slshape nrb}} row blocks and \mbox{\texttt{\mdseries\slshape ncb}} column blocks. 

 If \mbox{\texttt{\mdseries\slshape blocks}} is empty (i.e., if the matrix is a zero matrix) then the dimensions of the
blocks must be entered as \mbox{\texttt{\mdseries\slshape rpb}} and \mbox{\texttt{\mdseries\slshape cpb}}, and the zero element as \mbox{\texttt{\mdseries\slshape zero}}. 

 Note that all blocks must be ordinary matrices (see{\nobreakspace}\texttt{IsOrdinaryMatrix} (\ref{IsOrdinaryMatrix})), and also the block matrix is an ordinary matrix. 
\begin{Verbatim}[commandchars=!@|,fontsize=\small,frame=single,label=Example]
  !gapprompt@gap>| !gapinput@M := BlockMatrix([[1,1,[[1, 2],[ 3, 4]]],|
  !gapprompt@>| !gapinput@                     [1,2,[[9,10],[11,12]]],|
  !gapprompt@>| !gapinput@                     [2,2,[[5, 6],[ 7, 8]]]],2,2);|
  <block matrix of dimensions (2*2)x(2*2)>
  !gapprompt@gap>| !gapinput@Display(M);|
  [ [   1,   2,   9,  10 ],
    [   3,   4,  11,  12 ],
    [   0,   0,   5,   6 ],
    [   0,   0,   7,   8 ] ]
\end{Verbatim}
 }

 

\subsection{\textcolor{Chapter }{MatrixByBlockMatrix}}
\logpage{[ 24, 17, 3 ]}\nobreak
\hyperdef{L}{X83FAF4158180041F}{}
{\noindent\textcolor{FuncColor}{$\triangleright$\enspace\texttt{MatrixByBlockMatrix({\mdseries\slshape blockmat})\index{MatrixByBlockMatrix@\texttt{MatrixByBlockMatrix}}
\label{MatrixByBlockMatrix}
}\hfill{\scriptsize (attribute)}}\\


 returns a plain ordinary matrix that is equal to the block matrix \mbox{\texttt{\mdseries\slshape blockmat}}. }

 }

  
\section{\textcolor{Chapter }{Linear Programming}}\label{Linear Programming}
\logpage{[ 24, 18, 0 ]}
\hyperdef{L}{X782F2EBF80C431D0}{}
{
  

\subsection{\textcolor{Chapter }{SimplexMethod}}
\logpage{[ 24, 18, 1 ]}\nobreak
\hyperdef{L}{X845D5F8D7D905CB8}{}
{\noindent\textcolor{FuncColor}{$\triangleright$\enspace\texttt{SimplexMethod({\mdseries\slshape A, b, c})\index{SimplexMethod@\texttt{SimplexMethod}}
\label{SimplexMethod}
}\hfill{\scriptsize (function)}}\\


 Find a rational vector \mbox{\texttt{\mdseries\slshape x}} that maximizes $\mbox{\texttt{\mdseries\slshape x}}\cdot\mbox{\texttt{\mdseries\slshape c}}$, subject to the constraint $\mbox{\texttt{\mdseries\slshape A}}\cdot\mbox{\texttt{\mdseries\slshape x}}\le\mbox{\texttt{\mdseries\slshape b}}$. 
\begin{Verbatim}[commandchars=!@|,fontsize=\small,frame=single,label=Example]
  !gapprompt@gap>| !gapinput@A:=[[3,1,1,4],[1,-3,2,3],[2,1,3,-1]];;|
  !gapprompt@gap>| !gapinput@b:=[12,7,10];;c:=[2,4,3,1];;|
  !gapprompt@gap>| !gapinput@SimplexMethod(A,b,c);|
  [ [ 0, 52/5, 0, 2/5 ], 42 ]
\end{Verbatim}
 }

 }

 }

              
\chapter{\textcolor{Chapter }{Integral matrices and lattices}}\label{Integral matrices and lattices}
\logpage{[ 25, 0, 0 ]}
\hyperdef{L}{X8414F20D8412DDA4}{}
{
   
\section{\textcolor{Chapter }{Linear equations over the integers and Integral Matrices}}\label{Linear equations over the integers and Integral Matrices}
\logpage{[ 25, 1, 0 ]}
\hyperdef{L}{X786A64B983339767}{}
{
  

\subsection{\textcolor{Chapter }{NullspaceIntMat}}
\logpage{[ 25, 1, 1 ]}\nobreak
\hyperdef{L}{X792315717F5B0294}{}
{\noindent\textcolor{FuncColor}{$\triangleright$\enspace\texttt{NullspaceIntMat({\mdseries\slshape mat})\index{NullspaceIntMat@\texttt{NullspaceIntMat}}
\label{NullspaceIntMat}
}\hfill{\scriptsize (attribute)}}\\


 If \mbox{\texttt{\mdseries\slshape mat}} is a matrix with integral entries, this function returns a list of vectors
that forms a basis of the integral nullspace of \mbox{\texttt{\mdseries\slshape mat}}, i.e., of those vectors in the nullspace of \mbox{\texttt{\mdseries\slshape mat}} that have integral entries. 
\begin{Verbatim}[commandchars=!@|,fontsize=\small,frame=single,label=Example]
  !gapprompt@gap>| !gapinput@mat:=[[1,2,7],[4,5,6],[7,8,9],[10,11,19],[5,7,12]];;|
  !gapprompt@gap>| !gapinput@NullspaceMat(mat);|
  [ [ -7/4, 9/2, -15/4, 1, 0 ], [ -3/4, -3/2, 1/4, 0, 1 ] ]
  !gapprompt@gap>| !gapinput@NullspaceIntMat(mat);|
  [ [ 1, 18, -9, 2, -6 ], [ 0, 24, -13, 3, -7 ] ]
\end{Verbatim}
 }

 

\subsection{\textcolor{Chapter }{SolutionIntMat}}
\logpage{[ 25, 1, 2 ]}\nobreak
\hyperdef{L}{X7D749F317DBD1E69}{}
{\noindent\textcolor{FuncColor}{$\triangleright$\enspace\texttt{SolutionIntMat({\mdseries\slshape mat, vec})\index{SolutionIntMat@\texttt{SolutionIntMat}}
\label{SolutionIntMat}
}\hfill{\scriptsize (operation)}}\\


 If \mbox{\texttt{\mdseries\slshape mat}} is a matrix with integral entries and \mbox{\texttt{\mdseries\slshape vec}} a vector with integral entries, this function returns a vector $x$ with integer entries that is a solution of the equation $x$ \texttt{* \mbox{\texttt{\mdseries\slshape mat}} = \mbox{\texttt{\mdseries\slshape vec}}}. It returns \texttt{fail} if no such vector exists. 
\begin{Verbatim}[commandchars=!@|,fontsize=\small,frame=single,label=Example]
  !gapprompt@gap>| !gapinput@mat:=[[1,2,7],[4,5,6],[7,8,9],[10,11,19],[5,7,12]];;|
  !gapprompt@gap>| !gapinput@SolutionMat(mat,[95,115,182]);|
  [ 47/4, -17/2, 67/4, 0, 0 ]
  !gapprompt@gap>| !gapinput@SolutionIntMat(mat,[95,115,182]);|
  [ 2285, -5854, 4888, -1299, 0 ]
\end{Verbatim}
 }

 

\subsection{\textcolor{Chapter }{SolutionNullspaceIntMat}}
\logpage{[ 25, 1, 3 ]}\nobreak
\hyperdef{L}{X82CECB6E7D515CD2}{}
{\noindent\textcolor{FuncColor}{$\triangleright$\enspace\texttt{SolutionNullspaceIntMat({\mdseries\slshape mat, vec})\index{SolutionNullspaceIntMat@\texttt{SolutionNullspaceIntMat}}
\label{SolutionNullspaceIntMat}
}\hfill{\scriptsize (operation)}}\\


 This function returns a list of length two, its first entry being the result
of a call to \texttt{SolutionIntMat} (\ref{SolutionIntMat}) with same arguments, the second the result of \texttt{NullspaceIntMat} (\ref{NullspaceIntMat}) applied to the matrix \mbox{\texttt{\mdseries\slshape mat}}. The calculation is performed faster than if two separate calls would be
used. 
\begin{Verbatim}[commandchars=!@|,fontsize=\small,frame=single,label=Example]
  !gapprompt@gap>| !gapinput@mat:=[[1,2,7],[4,5,6],[7,8,9],[10,11,19],[5,7,12]];;|
  !gapprompt@gap>| !gapinput@SolutionNullspaceIntMat(mat,[95,115,182]);|
  [ [ 2285, -5854, 4888, -1299, 0 ],
    [ [ 1, 18, -9, 2, -6 ], [ 0, 24, -13, 3, -7 ] ] ]
\end{Verbatim}
 }

 

\subsection{\textcolor{Chapter }{BaseIntMat}}
\logpage{[ 25, 1, 4 ]}\nobreak
\hyperdef{L}{X7F66E8EA7D1AA2C1}{}
{\noindent\textcolor{FuncColor}{$\triangleright$\enspace\texttt{BaseIntMat({\mdseries\slshape mat})\index{BaseIntMat@\texttt{BaseIntMat}}
\label{BaseIntMat}
}\hfill{\scriptsize (attribute)}}\\


 If \mbox{\texttt{\mdseries\slshape mat}} is a matrix with integral entries, this function returns a list of vectors
that forms a basis of the integral row space of \mbox{\texttt{\mdseries\slshape mat}}, i.e. of the set of integral linear combinations of the rows of \mbox{\texttt{\mdseries\slshape mat}}. 
\begin{Verbatim}[commandchars=!@|,fontsize=\small,frame=single,label=Example]
  !gapprompt@gap>| !gapinput@mat:=[[1,2,7],[4,5,6],[10,11,19]];;|
  !gapprompt@gap>| !gapinput@BaseIntMat(mat);|
  [ [ 1, 2, 7 ], [ 0, 3, 7 ], [ 0, 0, 15 ] ]
\end{Verbatim}
 }

 

\subsection{\textcolor{Chapter }{BaseIntersectionIntMats}}
\logpage{[ 25, 1, 5 ]}\nobreak
\hyperdef{L}{X8771349D865C9179}{}
{\noindent\textcolor{FuncColor}{$\triangleright$\enspace\texttt{BaseIntersectionIntMats({\mdseries\slshape m, n})\index{BaseIntersectionIntMats@\texttt{BaseIntersectionIntMats}}
\label{BaseIntersectionIntMats}
}\hfill{\scriptsize (operation)}}\\


 If \mbox{\texttt{\mdseries\slshape m}} and \mbox{\texttt{\mdseries\slshape n}} are matrices with integral entries, this function returns a list of vectors
that forms a basis of the intersection of the integral row spaces of \mbox{\texttt{\mdseries\slshape m}} and \mbox{\texttt{\mdseries\slshape n}}. 
\begin{Verbatim}[commandchars=!@|,fontsize=\small,frame=single,label=Example]
  !gapprompt@gap>| !gapinput@nat:=[[5,7,2],[4,2,5],[7,1,4]];;|
  !gapprompt@gap>| !gapinput@BaseIntMat(nat);|
  [ [ 1, 1, 15 ], [ 0, 2, 55 ], [ 0, 0, 64 ] ]
  !gapprompt@gap>| !gapinput@BaseIntersectionIntMats(mat,nat);|
  [ [ 1, 5, 509 ], [ 0, 6, 869 ], [ 0, 0, 960 ] ]
\end{Verbatim}
 }

 

\subsection{\textcolor{Chapter }{ComplementIntMat}}
\logpage{[ 25, 1, 6 ]}\nobreak
\hyperdef{L}{X7848EF9F83D491C1}{}
{\noindent\textcolor{FuncColor}{$\triangleright$\enspace\texttt{ComplementIntMat({\mdseries\slshape full, sub})\index{ComplementIntMat@\texttt{ComplementIntMat}}
\label{ComplementIntMat}
}\hfill{\scriptsize (operation)}}\\


 Let \mbox{\texttt{\mdseries\slshape full}} be a list of integer vectors generating an integral row module $M$ and \mbox{\texttt{\mdseries\slshape sub}} a list of vectors defining a submodule $S$ of $M$. This function computes a free basis for $M$ that extends $S$. I.e., if the dimension of $S$ is $n$ it determines a basis $B = \{ b_1, \ldots, b_m \}$ for $M$, as well as $n$ integers $x_i$ such that the $n$ vectors $s_i:= x_i \cdot b_i$ form a basis for $S$. 

 It returns a record with the following components: 
\begin{description}
\item[{\texttt{complement}}]  the vectors $b_{{n+1}}$ up to $b_m$ (they generate a complement to $S$). 
\item[{\texttt{sub}}]  the vectors $s_i$ (a basis for $S$). 
\item[{\texttt{moduli}}]  the factors $x_i$. 
\end{description}
 
\begin{Verbatim}[commandchars=!@|,fontsize=\small,frame=single,label=Example]
  !gapprompt@gap>| !gapinput@m:=IdentityMat(3);;|
  !gapprompt@gap>| !gapinput@n:=[[1,2,3],[4,5,6]];;|
  !gapprompt@gap>| !gapinput@ComplementIntMat(m,n);|
  rec( complement := [ [ 0, 0, 1 ] ], moduli := [ 1, 3 ],
    sub := [ [ 1, 2, 3 ], [ 0, 3, 6 ] ] )
\end{Verbatim}
 }

 }

  
\section{\textcolor{Chapter }{Normal Forms over the Integers}}\label{Normal Forms over the Integers}
\logpage{[ 25, 2, 0 ]}
\hyperdef{L}{X8143C1448069D846}{}
{
  This section describes the computation of the Hermite and Smith normal form of
integer matrices. 

 The Hermite Normal Form (HNF) $H$ of an integer matrix $A$ is a row equivalent upper triangular form such that all
off\texttt{\symbol{45}}diagonal entries are reduced modulo the diagonal entry
of the column they are in. There exists a unique unimodular matrix $Q$ such that $Q A = H$. 

 The Smith Normal Form $S$ of an integer matrix $A$ is the unique equivalent diagonal form with $S_i$ dividing $S_j$ for $i < j$. There exist unimodular integer matrices $P, Q$ such that $P A Q = S$. 

 All routines described in this section build on the ``workhorse'' routine \texttt{NormalFormIntMat} (\ref{NormalFormIntMat}). 

\subsection{\textcolor{Chapter }{TriangulizedIntegerMat}}
\logpage{[ 25, 2, 1 ]}\nobreak
\hyperdef{L}{X783CEC847D81F22A}{}
{\noindent\textcolor{FuncColor}{$\triangleright$\enspace\texttt{TriangulizedIntegerMat({\mdseries\slshape mat})\index{TriangulizedIntegerMat@\texttt{TriangulizedIntegerMat}}
\label{TriangulizedIntegerMat}
}\hfill{\scriptsize (operation)}}\\


 Computes an upper triangular form of a matrix with integer entries. It returns
a mutable matrix in upper triangular form. }

 

\subsection{\textcolor{Chapter }{TriangulizedIntegerMatTransform}}
\logpage{[ 25, 2, 2 ]}\nobreak
\hyperdef{L}{X7DBE174E8625AFA5}{}
{\noindent\textcolor{FuncColor}{$\triangleright$\enspace\texttt{TriangulizedIntegerMatTransform({\mdseries\slshape mat})\index{TriangulizedIntegerMatTransform@\texttt{TriangulizedIntegerMatTransform}}
\label{TriangulizedIntegerMatTransform}
}\hfill{\scriptsize (operation)}}\\


 Computes an upper triangular form of a matrix with integer entries. It returns
a record with a component \texttt{normal} (an immutable matrix in upper triangular form) and a component \texttt{rowtrans} that gives the transformations done to the original matrix to bring it into
upper triangular form. }

 

\subsection{\textcolor{Chapter }{TriangulizeIntegerMat}}
\logpage{[ 25, 2, 3 ]}\nobreak
\hyperdef{L}{X78CD40A687FE2311}{}
{\noindent\textcolor{FuncColor}{$\triangleright$\enspace\texttt{TriangulizeIntegerMat({\mdseries\slshape mat})\index{TriangulizeIntegerMat@\texttt{TriangulizeIntegerMat}}
\label{TriangulizeIntegerMat}
}\hfill{\scriptsize (operation)}}\\


 Changes \mbox{\texttt{\mdseries\slshape mat}} to be in upper triangular form. (The result is the same as that of \texttt{TriangulizedIntegerMat} (\ref{TriangulizedIntegerMat}), but \mbox{\texttt{\mdseries\slshape mat}} will be modified, thus using less memory.) If \mbox{\texttt{\mdseries\slshape mat}} is immutable an error will be triggered. 
\begin{Verbatim}[commandchars=!@|,fontsize=\small,frame=single,label=Example]
  !gapprompt@gap>| !gapinput@m:=[[1,15,28],[4,5,6],[7,8,9]];;|
  !gapprompt@gap>| !gapinput@TriangulizedIntegerMat(m);|
  [ [ 1, 15, 28 ], [ 0, 1, 1 ], [ 0, 0, 3 ] ]
  !gapprompt@gap>| !gapinput@n:=TriangulizedIntegerMatTransform(m);|
  rec( normal := [ [ 1, 15, 28 ], [ 0, 1, 1 ], [ 0, 0, 3 ] ],
    rank := 3, rowC := [ [ 1, 0, 0 ], [ 0, 1, 0 ], [ 0, 0, 1 ] ],
    rowQ := [ [ 1, 0, 0 ], [ 1, -30, 17 ], [ -3, 97, -55 ] ],
    rowtrans := [ [ 1, 0, 0 ], [ 1, -30, 17 ], [ -3, 97, -55 ] ],
    signdet := 1 )
  !gapprompt@gap>| !gapinput@n.rowtrans*m=n.normal;|
  true
  !gapprompt@gap>| !gapinput@TriangulizeIntegerMat(m); m;|
  [ [ 1, 15, 28 ], [ 0, 1, 1 ], [ 0, 0, 3 ] ]
\end{Verbatim}
 }

 

\subsection{\textcolor{Chapter }{HermiteNormalFormIntegerMat}}
\logpage{[ 25, 2, 4 ]}\nobreak
\hyperdef{L}{X8535AC327932B89F}{}
{\noindent\textcolor{FuncColor}{$\triangleright$\enspace\texttt{HermiteNormalFormIntegerMat({\mdseries\slshape mat})\index{HermiteNormalFormIntegerMat@\texttt{HermiteNormalFormIntegerMat}}
\label{HermiteNormalFormIntegerMat}
}\hfill{\scriptsize (operation)}}\\


 This operation computes the Hermite normal form of a matrix \mbox{\texttt{\mdseries\slshape mat}} with integer entries. It returns a immutable matrix in HNF. }

 

\subsection{\textcolor{Chapter }{HermiteNormalFormIntegerMatTransform}}
\logpage{[ 25, 2, 5 ]}\nobreak
\hyperdef{L}{X7FDA78F979574ACC}{}
{\noindent\textcolor{FuncColor}{$\triangleright$\enspace\texttt{HermiteNormalFormIntegerMatTransform({\mdseries\slshape mat})\index{HermiteNormalFormIntegerMatTransform@\texttt{Hermite}\-\texttt{Normal}\-\texttt{Form}\-\texttt{Integer}\-\texttt{Mat}\-\texttt{Transform}}
\label{HermiteNormalFormIntegerMatTransform}
}\hfill{\scriptsize (operation)}}\\


 This operation computes the Hermite normal form of a matrix \mbox{\texttt{\mdseries\slshape mat}} with integer entries. It returns a record with components \texttt{normal} (a matrix $H$) and \texttt{rowtrans} (a matrix $Q$) such that $Q A = H$. 
\begin{Verbatim}[commandchars=!@|,fontsize=\small,frame=single,label=Example]
  !gapprompt@gap>| !gapinput@m:=[[1,15,28],[4,5,6],[7,8,9]];;|
  !gapprompt@gap>| !gapinput@HermiteNormalFormIntegerMat(m);|
  [ [ 1, 0, 1 ], [ 0, 1, 1 ], [ 0, 0, 3 ] ]
  !gapprompt@gap>| !gapinput@n:=HermiteNormalFormIntegerMatTransform(m);|
  rec( normal := [ [ 1, 0, 1 ], [ 0, 1, 1 ], [ 0, 0, 3 ] ], rank := 3,
    rowC := [ [ 1, 0, 0 ], [ 0, 1, 0 ], [ 0, 0, 1 ] ],
    rowQ := [ [ -2, 62, -35 ], [ 1, -30, 17 ], [ -3, 97, -55 ] ],
    rowtrans := [ [ -2, 62, -35 ], [ 1, -30, 17 ], [ -3, 97, -55 ] ],
    signdet := 1 )
  !gapprompt@gap>| !gapinput@n.rowtrans*m=n.normal;|
  true
\end{Verbatim}
 }

 

\subsection{\textcolor{Chapter }{SmithNormalFormIntegerMat}}
\logpage{[ 25, 2, 6 ]}\nobreak
\hyperdef{L}{X87089FEC7FBEEA8F}{}
{\noindent\textcolor{FuncColor}{$\triangleright$\enspace\texttt{SmithNormalFormIntegerMat({\mdseries\slshape mat})\index{SmithNormalFormIntegerMat@\texttt{SmithNormalFormIntegerMat}}
\label{SmithNormalFormIntegerMat}
}\hfill{\scriptsize (operation)}}\\


 This operation computes the Smith normal form of a matrix \mbox{\texttt{\mdseries\slshape mat}} with integer entries. It returns a new immutable matrix in the Smith normal
form. }

 

\subsection{\textcolor{Chapter }{SmithNormalFormIntegerMatTransforms}}
\logpage{[ 25, 2, 7 ]}\nobreak
\hyperdef{L}{X839C1F9E87273A93}{}
{\noindent\textcolor{FuncColor}{$\triangleright$\enspace\texttt{SmithNormalFormIntegerMatTransforms({\mdseries\slshape mat})\index{SmithNormalFormIntegerMatTransforms@\texttt{SmithNormalFormIntegerMatTransforms}}
\label{SmithNormalFormIntegerMatTransforms}
}\hfill{\scriptsize (operation)}}\\


 This operation computes the Smith normal form of a matrix \mbox{\texttt{\mdseries\slshape mat}} with integer entries. It returns a record with components \texttt{normal} (a matrix $S$), \texttt{rowtrans} (a matrix $P$), and \texttt{coltrans} (a matrix $Q$) such that $P A Q = S$. }

 

\subsection{\textcolor{Chapter }{DiagonalizeIntMat}}
\logpage{[ 25, 2, 8 ]}\nobreak
\hyperdef{L}{X80EF38737F6D61DB}{}
{\noindent\textcolor{FuncColor}{$\triangleright$\enspace\texttt{DiagonalizeIntMat({\mdseries\slshape mat})\index{DiagonalizeIntMat@\texttt{DiagonalizeIntMat}}
\label{DiagonalizeIntMat}
}\hfill{\scriptsize (function)}}\\


 This function changes \mbox{\texttt{\mdseries\slshape mat}} to its SNF. (The result is the same as that of \texttt{SmithNormalFormIntegerMat} (\ref{SmithNormalFormIntegerMat}), but \mbox{\texttt{\mdseries\slshape mat}} will be modified, thus using less memory.) If \mbox{\texttt{\mdseries\slshape mat}} is immutable an error will be triggered. 
\begin{Verbatim}[commandchars=!@|,fontsize=\small,frame=single,label=Example]
  !gapprompt@gap>| !gapinput@m:=[[1,15,28],[4,5,6],[7,8,9]];;|
  !gapprompt@gap>| !gapinput@SmithNormalFormIntegerMat(m);|
  [ [ 1, 0, 0 ], [ 0, 1, 0 ], [ 0, 0, 3 ] ]
  !gapprompt@gap>| !gapinput@n:=SmithNormalFormIntegerMatTransforms(m);|
  rec( colC := [ [ 1, 0, 0 ], [ 0, 1, 0 ], [ 0, 0, 1 ] ],
    colQ := [ [ 1, 0, -1 ], [ 0, 1, -1 ], [ 0, 0, 1 ] ],
    coltrans := [ [ 1, 0, -1 ], [ 0, 1, -1 ], [ 0, 0, 1 ] ],
    normal := [ [ 1, 0, 0 ], [ 0, 1, 0 ], [ 0, 0, 3 ] ], rank := 3,
    rowC := [ [ 1, 0, 0 ], [ 0, 1, 0 ], [ 0, 0, 1 ] ],
    rowQ := [ [ -2, 62, -35 ], [ 1, -30, 17 ], [ -3, 97, -55 ] ],
    rowtrans := [ [ -2, 62, -35 ], [ 1, -30, 17 ], [ -3, 97, -55 ] ],
    signdet := 1 )
  !gapprompt@gap>| !gapinput@n.rowtrans*m*n.coltrans=n.normal;|
  true
  !gapprompt@gap>| !gapinput@DiagonalizeIntMat(m);m;|
  [ [ 1, 0, 0 ], [ 0, 1, 0 ], [ 0, 0, 3 ] ]
\end{Verbatim}
 }

 

\subsection{\textcolor{Chapter }{NormalFormIntMat}}
\logpage{[ 25, 2, 9 ]}\nobreak
\hyperdef{L}{X81FB746E82BE6CDA}{}
{\noindent\textcolor{FuncColor}{$\triangleright$\enspace\texttt{NormalFormIntMat({\mdseries\slshape mat, options})\index{NormalFormIntMat@\texttt{NormalFormIntMat}}
\label{NormalFormIntMat}
}\hfill{\scriptsize (function)}}\\


 This general operation for computation of various Normal Forms is probably the
most efficient. 

 Options bit values: 
\begin{description}
\item[{0/1}]  Triangular Form / Smith Normal Form. 
\item[{2}]  Reduce off diagonal entries. 
\item[{4}]  Row Transformations. 
\item[{8}]  Col Transformations. 
\item[{16}]  Destructive (the original matrix may be destroyed) 
\end{description}
 

 Compute a Triangular, Hermite or Smith form of the $n \times m$ integer input matrix $A$. Optionally, compute $n \times n$ and $m \times m$ unimodular transforming matrices $Q, P$ which satisfy $Q A = H$ or $Q A P = S$.   

 Note option is a value ranging from 0 to 15 but not all options make sense
(e.g., reducing off diagonal entries with SNF option selected already). If an
option makes no sense it is ignored. 

 Returns a record with component \texttt{normal} containing the computed normal form and optional components \texttt{rowtrans} and/or \texttt{coltrans} which hold the respective transformation matrix. Also in the record are
components holding the sign of the determinant, \texttt{signdet}, and the rank of the matrix, \texttt{rank}. 
\begin{Verbatim}[commandchars=!@|,fontsize=\small,frame=single,label=Example]
  !gapprompt@gap>| !gapinput@m:=[[1,15,28],[4,5,6],[7,8,9]];;|
  !gapprompt@gap>| !gapinput@NormalFormIntMat(m,0);  # Triangular, no transforms|
  rec( normal := [ [ 1, 15, 28 ], [ 0, 1, 1 ], [ 0, 0, 3 ] ],
    rank := 3, signdet := 1 )
  !gapprompt@gap>| !gapinput@NormalFormIntMat(m,6);  # Hermite Normal Form with row transforms|
  rec( normal := [ [ 1, 0, 1 ], [ 0, 1, 1 ], [ 0, 0, 3 ] ], rank := 3,
    rowC := [ [ 1, 0, 0 ], [ 0, 1, 0 ], [ 0, 0, 1 ] ],
    rowQ := [ [ -2, 62, -35 ], [ 1, -30, 17 ], [ -3, 97, -55 ] ],
    rowtrans := [ [ -2, 62, -35 ], [ 1, -30, 17 ], [ -3, 97, -55 ] ],
    signdet := 1 )
  !gapprompt@gap>| !gapinput@NormalFormIntMat(m,13); # Smith Normal Form with both transforms|
  rec( colC := [ [ 1, 0, 0 ], [ 0, 1, 0 ], [ 0, 0, 1 ] ],
    colQ := [ [ 1, 0, -1 ], [ 0, 1, -1 ], [ 0, 0, 1 ] ],
    coltrans := [ [ 1, 0, -1 ], [ 0, 1, -1 ], [ 0, 0, 1 ] ],
    normal := [ [ 1, 0, 0 ], [ 0, 1, 0 ], [ 0, 0, 3 ] ], rank := 3,
    rowC := [ [ 1, 0, 0 ], [ 0, 1, 0 ], [ 0, 0, 1 ] ],
    rowQ := [ [ -2, 62, -35 ], [ 1, -30, 17 ], [ -3, 97, -55 ] ],
    rowtrans := [ [ -2, 62, -35 ], [ 1, -30, 17 ], [ -3, 97, -55 ] ],
    signdet := 1 )
  !gapprompt@gap>| !gapinput@last.rowtrans*m*last.coltrans;|
  [ [ 1, 0, 0 ], [ 0, 1, 0 ], [ 0, 0, 3 ] ]
\end{Verbatim}
 }

 

\subsection{\textcolor{Chapter }{AbelianInvariantsOfList}}
\logpage{[ 25, 2, 10 ]}\nobreak
\hyperdef{L}{X8221694D7C99197A}{}
{\noindent\textcolor{FuncColor}{$\triangleright$\enspace\texttt{AbelianInvariantsOfList({\mdseries\slshape list})\index{AbelianInvariantsOfList@\texttt{AbelianInvariantsOfList}}
\label{AbelianInvariantsOfList}
}\hfill{\scriptsize (attribute)}}\\


 Given a list of nonnegative integers, this routine returns a sorted list
containing the prime power factors of the positive entries in the original
list, as well as all zeroes of the original list. 
\begin{Verbatim}[commandchars=!@|,fontsize=\small,frame=single,label=Example]
  !gapprompt@gap>| !gapinput@AbelianInvariantsOfList([4,6,2,0,12]);|
  [ 0, 2, 2, 3, 3, 4, 4 ]
\end{Verbatim}
 }

     }

  
\section{\textcolor{Chapter }{Determinant of an integer matrix}}\label{Determinant of an integer matrix}
\logpage{[ 25, 3, 0 ]}
\hyperdef{L}{X80F6990983C979FB}{}
{
  

\subsection{\textcolor{Chapter }{DeterminantIntMat}}
\logpage{[ 25, 3, 1 ]}\nobreak
\hyperdef{L}{X787599E087F4C0BA}{}
{\noindent\textcolor{FuncColor}{$\triangleright$\enspace\texttt{DeterminantIntMat({\mdseries\slshape mat})\index{DeterminantIntMat@\texttt{DeterminantIntMat}}
\label{DeterminantIntMat}
}\hfill{\scriptsize (function)}}\\


 \index{determinant!integer matrix} Computes the determinant of an integer matrix using the same strategy as \texttt{NormalFormIntMat} (\ref{NormalFormIntMat}). This method is faster in general for matrices greater than $20 \times 20$ but quite a lot slower for smaller matrices. It therefore passes the work to
the more general \texttt{DeterminantMat} (\ref{DeterminantMat}) for these smaller matrices. }

 }

  
\section{\textcolor{Chapter }{Decompositions}}\label{Decompositions}
\logpage{[ 25, 4, 0 ]}
\hyperdef{L}{X79F2EFEC7C3EA80C}{}
{
  \index{decomposition matrix} \index{DEC} For computing the decomposition of a vector of integers into the rows of a
matrix of integers, with integral coefficients, one can use $p$\texttt{\symbol{45}}adic approximations, as follows. 

 Let $A$ be a square integral matrix, and $p$ an odd prime. The reduction of $A$ modulo $p$ is $\overline{A}$, its entries are chosen in the interval $[ -(p-1)/2, (p-1)/2 ]$. If $\overline{A}$ is regular over the field with $p$ elements, we can form $A' = \overline{A}^{{-1}}$. Now we consider the integral linear equation system $x A = b$, i.e., we look for an integral solution $x$. Define $b_0 = b$, and then iteratively compute 
\[ x_i = (b_i A') \bmod p, b_{{i+1}} = (b_i - x_i A) / p, i = 0, 1, 2, \ldots . \]
 By induction, we get 
\[ p^{{i+1}} b_{{i+1}} + \left( \sum_{{j = 0}}^i p^j x_j \right) A = b. \]
 If there is an integral solution $x$ then it is unique, and there is an index $l$ such that $b_{{l+1}}$ is zero and $x = \sum_{{j = 0}}^l p^j x_j$. 

 There are two useful generalizations of this idea. First, $A$ need not be square; it is only necessary that there is a square regular matrix
formed by a subset of columns of $A$. Second, $A$ does not need to be integral; the entries may be cyclotomic integers as well,
in this case one can replace each column of $A$ by the columns formed by the coefficients w.r.t.{\nobreakspace}an integral
basis (which are integers). Note that this preprocessing must be performed
compatibly for $A$ and $b$. 

 \textsf{GAP} provides the following functions for this purpose (see also{\nobreakspace}\texttt{InverseMatMod} (\ref{InverseMatMod})). 

\subsection{\textcolor{Chapter }{Decomposition}}
\logpage{[ 25, 4, 1 ]}\nobreak
\hyperdef{L}{X7911A60384C511AB}{}
{\noindent\textcolor{FuncColor}{$\triangleright$\enspace\texttt{Decomposition({\mdseries\slshape A, B, depth})\index{Decomposition@\texttt{Decomposition}}
\label{Decomposition}
}\hfill{\scriptsize (operation)}}\\


 For a $m \times n$ matrix \mbox{\texttt{\mdseries\slshape A}} of cyclotomics that has rank $m \leq n$, and a list \mbox{\texttt{\mdseries\slshape B}} of cyclotomic vectors, each of length $n$, \texttt{Decomposition} tries to find integral solutions of the linear equation systems \texttt{\mbox{\texttt{\mdseries\slshape x}} * \mbox{\texttt{\mdseries\slshape A}} = \mbox{\texttt{\mdseries\slshape B}}[i]}, by computing the $p$\texttt{\symbol{45}}adic series of hypothetical solutions. 

 \texttt{Decomposition( \mbox{\texttt{\mdseries\slshape A}}, \mbox{\texttt{\mdseries\slshape B}}, \mbox{\texttt{\mdseries\slshape depth}} )}, where \mbox{\texttt{\mdseries\slshape depth}} is a nonnegative integer, computes for each vector \texttt{\mbox{\texttt{\mdseries\slshape B}}[i]} the initial part $\sum_{{k = 0}}^{\mbox{\texttt{\mdseries\slshape depth}}} x_k p^k$, with all $x_k$ vectors of integers with entries bounded by $\pm (p-1)/2$. The prime $p$ is set to 83 first; if the reduction of \mbox{\texttt{\mdseries\slshape A}} modulo $p$ is singular, the next prime is chosen automatically. 

 A list \mbox{\texttt{\mdseries\slshape X}} is returned. If the computed initial part for \texttt{\mbox{\texttt{\mdseries\slshape x}} * \mbox{\texttt{\mdseries\slshape A}} = \mbox{\texttt{\mdseries\slshape B}}[i]} \emph{is} a solution, we have \texttt{\mbox{\texttt{\mdseries\slshape X}}[i] = \mbox{\texttt{\mdseries\slshape x}}}, otherwise \texttt{\mbox{\texttt{\mdseries\slshape X}}[i] = fail}. 

 If \mbox{\texttt{\mdseries\slshape depth}} is not an integer then it must be the string \texttt{"nonnegative"}. \texttt{Decomposition( \mbox{\texttt{\mdseries\slshape A}}, \mbox{\texttt{\mdseries\slshape B}}, "nonnegative" )} assumes that the solutions have only nonnegative entries, and that the first
column of \mbox{\texttt{\mdseries\slshape A}} consists of positive integers. This is satisfied, e.g., for the decomposition
of ordinary characters into Brauer characters. In this case the necessary
number \mbox{\texttt{\mdseries\slshape depth}} of iterations can be computed; the \texttt{i}\texttt{\symbol{45}}th entry of the returned list is \texttt{fail} if there \emph{exists} no nonnegative integral solution of the system \texttt{\mbox{\texttt{\mdseries\slshape x}} * \mbox{\texttt{\mdseries\slshape A}} = \mbox{\texttt{\mdseries\slshape B}}[i]}, and it is the solution otherwise. 

 \emph{Note} that the result is a list of \texttt{fail} if \mbox{\texttt{\mdseries\slshape A}} has not full rank, even if there might be a unique integral solution for some
equation system. }

 

\subsection{\textcolor{Chapter }{LinearIndependentColumns}}
\logpage{[ 25, 4, 2 ]}\nobreak
\hyperdef{L}{X843A976787600F13}{}
{\noindent\textcolor{FuncColor}{$\triangleright$\enspace\texttt{LinearIndependentColumns({\mdseries\slshape mat})\index{LinearIndependentColumns@\texttt{LinearIndependentColumns}}
\label{LinearIndependentColumns}
}\hfill{\scriptsize (function)}}\\


 Called with a matrix \mbox{\texttt{\mdseries\slshape mat}}, \texttt{LinearIndependentColumns} returns a maximal list of column positions such that the restriction of \mbox{\texttt{\mdseries\slshape mat}} to these columns has the same rank as \mbox{\texttt{\mdseries\slshape mat}}. }

 

\subsection{\textcolor{Chapter }{PadicCoefficients}}
\logpage{[ 25, 4, 3 ]}\nobreak
\hyperdef{L}{X8285776B7DD86925}{}
{\noindent\textcolor{FuncColor}{$\triangleright$\enspace\texttt{PadicCoefficients({\mdseries\slshape A, Amodpinv, b, prime, depth})\index{PadicCoefficients@\texttt{PadicCoefficients}}
\label{PadicCoefficients}
}\hfill{\scriptsize (function)}}\\


 Let \mbox{\texttt{\mdseries\slshape A}} be an integral matrix, \mbox{\texttt{\mdseries\slshape prime}} a prime integer, \mbox{\texttt{\mdseries\slshape Amodpinv}} an inverse of \mbox{\texttt{\mdseries\slshape A}} modulo \mbox{\texttt{\mdseries\slshape prime}}, \mbox{\texttt{\mdseries\slshape b}} an integral vector, and \mbox{\texttt{\mdseries\slshape depth}} a nonnegative integer. \texttt{PadicCoefficients} returns the list $[ x_0, x_1, \ldots, x_l, b_{{l+1}} ]$ describing the \mbox{\texttt{\mdseries\slshape prime}}\texttt{\symbol{45}}adic approximation of \mbox{\texttt{\mdseries\slshape b}} (see above), where $l = \mbox{\texttt{\mdseries\slshape depth}}$ or $l$ is minimal with the property that $b_{{l+1}} = 0$. }

 

\subsection{\textcolor{Chapter }{IntegralizedMat}}
\logpage{[ 25, 4, 4 ]}\nobreak
\hyperdef{L}{X7F5C619B7A9C3EB9}{}
{\noindent\textcolor{FuncColor}{$\triangleright$\enspace\texttt{IntegralizedMat({\mdseries\slshape A[, inforec]})\index{IntegralizedMat@\texttt{IntegralizedMat}}
\label{IntegralizedMat}
}\hfill{\scriptsize (function)}}\\


 \texttt{IntegralizedMat} returns, for a matrix \mbox{\texttt{\mdseries\slshape A}} of cyclotomics, a record \texttt{intmat} with components \texttt{mat} and \texttt{inforec}. Each family of algebraic conjugate columns of \mbox{\texttt{\mdseries\slshape A}} is encoded in a set of columns of the rational matrix \texttt{intmat.mat} by replacing cyclotomics in \mbox{\texttt{\mdseries\slshape A}} by their coefficients w.r.t.{\nobreakspace}an integral basis. \texttt{intmat.inforec} is a record containing the information how to encode the columns. 

 If the only argument is \mbox{\texttt{\mdseries\slshape A}}, the value of the component \texttt{inforec} is computed that can be entered as second argument \mbox{\texttt{\mdseries\slshape inforec}} in a later call of \texttt{IntegralizedMat} with a matrix \mbox{\texttt{\mdseries\slshape B}} that shall be encoded compatibly with \mbox{\texttt{\mdseries\slshape A}}. }

 

\subsection{\textcolor{Chapter }{DecompositionInt}}
\logpage{[ 25, 4, 5 ]}\nobreak
\hyperdef{L}{X8512FB69824AE353}{}
{\noindent\textcolor{FuncColor}{$\triangleright$\enspace\texttt{DecompositionInt({\mdseries\slshape A, B, depth})\index{DecompositionInt@\texttt{DecompositionInt}}
\label{DecompositionInt}
}\hfill{\scriptsize (function)}}\\


 \texttt{DecompositionInt} does the same as \texttt{Decomposition} (\ref{Decomposition}), except that \mbox{\texttt{\mdseries\slshape A}} and \mbox{\texttt{\mdseries\slshape B}} must be integral matrices, and \mbox{\texttt{\mdseries\slshape depth}} must be a nonnegative integer. }

 }

  
\section{\textcolor{Chapter }{Lattice Reduction}}\label{Lattice Reduction}
\logpage{[ 25, 5, 0 ]}
\hyperdef{L}{X839C6ABE829355F4}{}
{
  

\subsection{\textcolor{Chapter }{LLLReducedBasis}}
\logpage{[ 25, 5, 1 ]}\nobreak
\hyperdef{L}{X7D0FCEF8859E8637}{}
{\noindent\textcolor{FuncColor}{$\triangleright$\enspace\texttt{LLLReducedBasis({\mdseries\slshape [L, ]vectors[, y][, "linearcomb"][, lllout]})\index{LLLReducedBasis@\texttt{LLLReducedBasis}}
\label{LLLReducedBasis}
}\hfill{\scriptsize (function)}}\\


 \index{LLL algorithm!for vectors} \index{short vectors spanning a lattice} \index{lattice base reduction} provides an implementation of the \emph{LLL algorithm} by Lenstra, Lenstra and Lov{\a'a}sz (see{\nobreakspace}\cite{LLL82}, \cite{Poh87}). The implementation follows the description in{\nobreakspace}\cite[p. 94f.]{Coh93}. 

 \texttt{LLLReducedBasis} returns a record whose component \texttt{basis} is a list of LLL reduced linearly independent vectors spanning the same
lattice as the list \mbox{\texttt{\mdseries\slshape vectors}}. \mbox{\texttt{\mdseries\slshape L}} must be a lattice, with scalar product of the vectors \mbox{\texttt{\mdseries\slshape v}} and \mbox{\texttt{\mdseries\slshape w}} given by \texttt{ScalarProduct( \mbox{\texttt{\mdseries\slshape L}}, \mbox{\texttt{\mdseries\slshape v}}, \mbox{\texttt{\mdseries\slshape w}} )}. If no lattice is specified then the scalar product of vectors given by \texttt{ScalarProduct( \mbox{\texttt{\mdseries\slshape v}}, \mbox{\texttt{\mdseries\slshape w}} )} is used. 

 In the case of the option \texttt{"linearcomb"}, the result record contains also the components \texttt{relations} and \texttt{transformation}, with the following meaning. \texttt{relations} is a basis of the relation space of \mbox{\texttt{\mdseries\slshape vectors}}, i.e., of vectors \mbox{\texttt{\mdseries\slshape x}} such that \texttt{\mbox{\texttt{\mdseries\slshape x}} * \mbox{\texttt{\mdseries\slshape vectors}}} is zero. \texttt{transformation} gives the expression of the new lattice basis in terms of the old, i.e., \texttt{transformation * \mbox{\texttt{\mdseries\slshape vectors}}} equals the \texttt{basis} component of the result. 

 Another optional argument is \mbox{\texttt{\mdseries\slshape y}}, the ``sensitivity'' of the algorithm, a rational number between $1/4$ and $1$ (the default value is $3/4$). 

 The optional argument \mbox{\texttt{\mdseries\slshape lllout}} is a record with the components \texttt{mue} and \texttt{B}, both lists of length $k$, with the meaning that if \mbox{\texttt{\mdseries\slshape lllout}} is present then the first $k$ vectors in \mbox{\texttt{\mdseries\slshape vectors}} form an LLL reduced basis of the lattice they generate, and \texttt{\mbox{\texttt{\mdseries\slshape lllout}}.mue} and \texttt{\mbox{\texttt{\mdseries\slshape lllout}}.B} contain their scalar products and norms used internally in the algorithm,
which are also present in the output of \texttt{LLLReducedBasis}. So \mbox{\texttt{\mdseries\slshape lllout}} can be used for ``incremental'' calls of \texttt{LLLReducedBasis}. 

 The function \texttt{LLLReducedGramMat} (\ref{LLLReducedGramMat}) computes an LLL reduced Gram matrix. 

 
\begin{Verbatim}[commandchars=!@|,fontsize=\small,frame=single,label=Example]
  !gapprompt@gap>| !gapinput@vectors:= [ [ 9, 1, 0, -1, -1 ], [ 15, -1, 0, 0, 0 ],|
  !gapprompt@>| !gapinput@               [ 16, 0, 1, 1, 1 ], [ 20, 0, -1, 0, 0 ],|
  !gapprompt@>| !gapinput@               [ 25, 1, 1, 0, 0 ] ];;|
  !gapprompt@gap>| !gapinput@LLLReducedBasis( vectors, "linearcomb" );|
  rec( B := [ 5, 36/5, 12, 50/3 ],
    basis := [ [ 1, 1, 1, 1, 1 ], [ 1, 1, -2, 1, 1 ],
        [ -1, 3, -1, -1, -1 ], [ -3, 1, 0, 2, 2 ] ],
    mue := [ [  ], [ 2/5 ], [ -1/5, 1/3 ], [ 2/5, 1/6, 1/6 ] ],
    relations := [ [ -1, 0, -1, 0, 1 ] ],
    transformation := [ [ 0, -1, 1, 0, 0 ], [ -1, -2, 0, 2, 0 ],
        [ 1, -2, 0, 1, 0 ], [ -1, -2, 1, 1, 0 ] ] )
\end{Verbatim}
 }

 

\subsection{\textcolor{Chapter }{LLLReducedGramMat}}
\logpage{[ 25, 5, 2 ]}\nobreak
\hyperdef{L}{X86D23EB885EDE60E}{}
{\noindent\textcolor{FuncColor}{$\triangleright$\enspace\texttt{LLLReducedGramMat({\mdseries\slshape G[, y]})\index{LLLReducedGramMat@\texttt{LLLReducedGramMat}}
\label{LLLReducedGramMat}
}\hfill{\scriptsize (function)}}\\


 \index{LLL algorithm!for Gram matrices} \index{lattice base reduction} \texttt{LLLReducedGramMat} provides an implementation of the \emph{LLL algorithm} by Lenstra, Lenstra and Lov{\a'a}sz (see{\nobreakspace}\cite{LLL82},{\nobreakspace}\cite{Poh87}). The implementation follows the description in \cite[p. 94f.]{Coh93}. 

 Let \mbox{\texttt{\mdseries\slshape G}} the Gram matrix of the vectors $(b_1, b_2, \ldots, b_n)$; this means \mbox{\texttt{\mdseries\slshape G}} is either a square symmetric matrix or lower triangular matrix (only the
entries in the lower triangular half are used by the program). 

 \texttt{LLLReducedGramMat} returns a record whose component \texttt{remainder} is the Gram matrix of the LLL reduced basis corresponding to $(b_1, b_2, \ldots, b_n)$. If \mbox{\texttt{\mdseries\slshape G}} is a lower triangular matrix then also the \texttt{remainder} component of the result record is a lower triangular matrix. 

 The result record contains also the components \texttt{relations} and \texttt{transformation}, which have the following meaning. 

 \texttt{relations} is a basis of the space of vectors $(x_1, x_2, \ldots, x_n)$ such that $\sum_{{i = 1}}^n x_i b_i$ is zero, and \texttt{transformation} gives the expression of the new lattice basis in terms of the old, i.e., \texttt{transformation} is the matrix $T$ such that $T \cdot \mbox{\texttt{\mdseries\slshape G}} \cdot T^{tr}$ is the \texttt{remainder} component of the result. 

 The optional argument \mbox{\texttt{\mdseries\slshape y}} denotes the ``sensitivity'' of the algorithm, it must be a rational number between $1/4$ and $1$; the default value is $\mbox{\texttt{\mdseries\slshape y}} = 3/4$. 

 The function \texttt{LLLReducedBasis} (\ref{LLLReducedBasis}) computes an LLL reduced basis. 

 
\begin{Verbatim}[commandchars=!@|,fontsize=\small,frame=single,label=Example]
  !gapprompt@gap>| !gapinput@g:= [ [ 4, 6, 5, 2, 2 ], [ 6, 13, 7, 4, 4 ],|
  !gapprompt@>| !gapinput@   [ 5, 7, 11, 2, 0 ], [ 2, 4, 2, 8, 4 ], [ 2, 4, 0, 4, 8 ] ];;|
  !gapprompt@gap>| !gapinput@LLLReducedGramMat( g );|
  rec( B := [ 4, 4, 75/16, 168/25, 32/7 ],
    mue := [ [  ], [ 1/2 ], [ 1/4, -1/8 ], [ 1/2, 1/4, -2/25 ],
        [ -1/4, 1/8, 37/75, 8/21 ] ], relations := [  ],
    remainder := [ [ 4, 2, 1, 2, -1 ], [ 2, 5, 0, 2, 0 ],
        [ 1, 0, 5, 0, 2 ], [ 2, 2, 0, 8, 2 ], [ -1, 0, 2, 2, 7 ] ],
    transformation := [ [ 1, 0, 0, 0, 0 ], [ -1, 1, 0, 0, 0 ],
        [ -1, 0, 1, 0, 0 ], [ 0, 0, 0, 1, 0 ], [ -2, 0, 1, 0, 1 ] ] )
\end{Verbatim}
 }

 }

  
\section{\textcolor{Chapter }{Orthogonal Embeddings}}\label{Orthogonal Embeddings}
\logpage{[ 25, 6, 0 ]}
\hyperdef{L}{X871DB00B803D5177}{}
{
  

\subsection{\textcolor{Chapter }{OrthogonalEmbeddings}}
\logpage{[ 25, 6, 1 ]}\nobreak
\hyperdef{L}{X842280C2808FF05D}{}
{\noindent\textcolor{FuncColor}{$\triangleright$\enspace\texttt{OrthogonalEmbeddings({\mdseries\slshape gram[, "positive"][, maxdim]})\index{OrthogonalEmbeddings@\texttt{OrthogonalEmbeddings}}
\label{OrthogonalEmbeddings}
}\hfill{\scriptsize (function)}}\\


 computes all possible orthogonal embeddings of a lattice given by its Gram
matrix \mbox{\texttt{\mdseries\slshape gram}}, which must be a regular symmetric matrix of integers. In other words, all
integral solutions $X$ of the equation $X^{tr} \cdot X = $\mbox{\texttt{\mdseries\slshape gram}} are calculated. The implementation follows the description in{\nobreakspace}\cite{Ple90}. 

 Usually there are many solutions $X$ but all their rows belong to a small set of vectors, so \texttt{OrthogonalEmbeddings} returns the solutions encoded by a record with the following components. 

 
\begin{description}
\item[{\texttt{vectors}}]  the list $L = [ x_1, x_2, \ldots, x_n ]$ of vectors that may be rows of a solution, up to sign; these are exactly the
vectors with the property $x_i \cdot $\mbox{\texttt{\mdseries\slshape gram}}$^{{-1}} \cdot x_i^{tr} \leq 1$, see{\nobreakspace}\texttt{ShortestVectors} (\ref{ShortestVectors}), 
\item[{\texttt{norms}}]  the list of values $x_i \cdot $\mbox{\texttt{\mdseries\slshape gram}}$^{{-1}} \cdot x_i^{tr}$, and 
\item[{\texttt{solutions}}]  a list $S$ of index lists; the $i$\texttt{\symbol{45}}th solution matrix is $L$\texttt{\texttt{\symbol{123}} }$S[i]$\texttt{ \texttt{\symbol{125}}}, so the dimension of the \mbox{\texttt{\mdseries\slshape i}}\texttt{\symbol{45}}th solution is the length of $S[i]$, and we have \mbox{\texttt{\mdseries\slshape gram}}$ = \sum_{{j \in S[i]}} x_j^{tr} \cdot x_j$, 
\end{description}
 

 The optional argument \texttt{"positive"} will cause \texttt{OrthogonalEmbeddings} to compute only vectors $x_i$ with nonnegative entries. In the context of characters this is allowed (and
useful) if \mbox{\texttt{\mdseries\slshape gram}} is the matrix of scalar products of ordinary characters. 

 When \texttt{OrthogonalEmbeddings} is called with the optional argument \mbox{\texttt{\mdseries\slshape maxdim}} (a positive integer), only solutions up to dimension \mbox{\texttt{\mdseries\slshape maxdim}} are computed; this may accelerate the algorithm. 

 
\begin{Verbatim}[commandchars=!@|,fontsize=\small,frame=single,label=Example]
  !gapprompt@gap>| !gapinput@b:= [ [ 3, -1, -1 ], [ -1, 3, -1 ], [ -1, -1, 3 ] ];;|
  !gapprompt@gap>| !gapinput@c:=OrthogonalEmbeddings( b );|
  rec( norms := [ 1, 1, 1, 1/2, 1/2, 1/2, 1/2, 1/2, 1/2 ],
    solutions := [ [ 1, 2, 3 ], [ 1, 6, 6, 7, 7 ], [ 2, 5, 5, 8, 8 ],
        [ 3, 4, 4, 9, 9 ], [ 4, 5, 6, 7, 8, 9 ] ],
    vectors := [ [ -1, 1, 1 ], [ 1, -1, 1 ], [ -1, -1, 1 ],
        [ -1, 1, 0 ], [ -1, 0, 1 ], [ 1, 0, 0 ], [ 0, -1, 1 ],
        [ 0, 1, 0 ], [ 0, 0, 1 ] ] )
  !gapprompt@gap>| !gapinput@c.vectors{ c.solutions[1] };|
  [ [ -1, 1, 1 ], [ 1, -1, 1 ], [ -1, -1, 1 ] ]
\end{Verbatim}
 

 \mbox{\texttt{\mdseries\slshape gram}} may be the matrix of scalar products of some virtual characters. From the
characters and the embedding given by the matrix $X$, \texttt{Decreased} (\ref{Decreased}) may be able to compute irreducibles. }

 

\subsection{\textcolor{Chapter }{ShortestVectors}}
\logpage{[ 25, 6, 2 ]}\nobreak
\hyperdef{L}{X79A692B6819353D4}{}
{\noindent\textcolor{FuncColor}{$\triangleright$\enspace\texttt{ShortestVectors({\mdseries\slshape G, m[, "positive"]})\index{ShortestVectors@\texttt{ShortestVectors}}
\label{ShortestVectors}
}\hfill{\scriptsize (function)}}\\


 Let \mbox{\texttt{\mdseries\slshape G}} be a regular matrix of a symmetric bilinear form, and \mbox{\texttt{\mdseries\slshape m}} a nonnegative integer. \texttt{ShortestVectors} computes the vectors $x$ that satisfy $x \cdot \mbox{\texttt{\mdseries\slshape G}} \cdot x^{tr} \leq \mbox{\texttt{\mdseries\slshape m}}$, and returns a record describing these vectors. The result record has the
components 
\begin{description}
\item[{\texttt{vectors}}]  list of the nonzero vectors $x$, but only one of each pair $(x,-x)$, 
\item[{\texttt{norms}}]  list of norms of the vectors according to the Gram matrix \mbox{\texttt{\mdseries\slshape G}}. 
\end{description}
 If the optional argument \texttt{"positive"} is entered, only those vectors $x$ with nonnegative entries are computed. 
\begin{Verbatim}[commandchars=!@|,fontsize=\small,frame=single,label=Example]
  !gapprompt@gap>| !gapinput@g:= [ [ 2, 1, 1 ], [ 1, 2, 1 ], [ 1, 1, 2 ] ];;|
  !gapprompt@gap>| !gapinput@ShortestVectors(g,4);|
  rec( norms := [ 4, 2, 2, 4, 2, 4, 2, 2, 2 ],
    vectors := [ [ -1, 1, 1 ], [ 0, 0, 1 ], [ -1, 0, 1 ], [ 1, -1, 1 ],
        [ 0, -1, 1 ], [ -1, -1, 1 ], [ 0, 1, 0 ], [ -1, 1, 0 ],
        [ 1, 0, 0 ] ] )
\end{Verbatim}
 }

 }

 }

         
\chapter{\textcolor{Chapter }{Vector and Matrix Objects}}\label{Vector and Matrix Objects}
\logpage{[ 26, 0, 0 ]}
\hyperdef{L}{X856C23B87E50F118}{}
{
  This chapter describes an interface to vector and matrix objects which are not
represented by plain lists (of plain lists), cf. Chapters \ref{Row Vectors} and \ref{Matrices}.  
\section{\textcolor{Chapter }{Concepts and Rules for Vector and Matrix Objects}}\label{Concepts and Rules for Vector and Matrix Objects}
\logpage{[ 26, 1, 0 ]}
\hyperdef{L}{X7A7275C27EC61ACE}{}
{
  Traditionally, vectors and matrices in \textsf{GAP} have been represented by (lists of) lists, see the chapters \ref{Row Vectors} and \ref{Matrices}. More precisely, the term ``vector'' (corresponding to the filter \texttt{IsVector} (\ref{IsVector})) is used in the abstract sense of an ``element of a vector space'', the term ``row vector'' (corresponding to \texttt{IsRowVector} (\ref{IsRowVector})) is used to denote a ``coordinate vector'' which is represented by a \textsf{GAP} list (see \texttt{IsList} (\ref{IsList})), and the term ``matrix'' is used to denote a list of lists, with additional properties (see \texttt{IsMatrix} (\ref{IsMatrix})). 

 Unfortunately, such lists (objects in \texttt{IsPlistRep} (\ref{IsPlistRep})) cannot store their type, and so it is impossible to use the advantages of \textsf{GAP}'s method selection on them. This situation is unsustainable in the long run
since more special representations (compressed, sparse, etc.) have already
been and even more will be implemented. Here we describe a programming
interface to vectors and matrices, which solves this problem, 

 The idea of this interface is that \textsf{GAP} should be able to represent vectors and matrices by objects that store their
type, in order to benefit from method selection. These objects are created by \texttt{Objectify} (\ref{Objectify}), we therefore refer to the them as ``vector objects'' and ``matrix objects'' respectively. 

 (Of course the terminology is somewhat confusing: An ``abstract matrix'' in \textsf{GAP} can be represented either by a list of lists or by a matrix object. It can be
detected from the filter \texttt{IsMatrixOrMatrixObj} (\ref{IsMatrixOrMatrixObj}); this is the union of the filters \texttt{IsMatrix} (\ref{IsMatrix}) {\textendash}which denotes those matrices that are represented by lists of
lists{\textendash} and the filter \texttt{IsMatrixObj} (\ref{IsMatrixObj}) {\textendash}which defines ``proper'' matrix objects in the above sense. In particular, we do \emph{not} regard the objects in \texttt{IsMatrix} (\ref{IsMatrix}) as special cases of objects in \texttt{IsMatrixObj} (\ref{IsMatrixObj}), or vice versa. Thus one can install specific methods for all three
situations: just for ``proper'' matrix objects, just for matrices represented by lists of lists, or for both
kinds of matrices. For example, a \textsf{GAP} package may decide to accept only ``proper'' matrix objects as arguments of its functions, or it may try to support also
objects in \texttt{IsMatrix} (\ref{IsMatrix}) as far as this is possible.) 

 We want to be able to write (efficient) code that is independent of the actual
representation (in the sense of \textsf{GAP}'s representation filters, see Section \ref{Representation}) and preserves it. 

 This latter requirement makes it necessary to distinguish between different
representations of matrices: ``Row list'' matrices (see \texttt{IsRowListMatrix} (\ref{IsRowListMatrix}) behave basically like lists of rows, in particular the rows are individual \textsf{GAP} objects that can be shared between different matrix objects. One can think of
other representations of matrices, such as matrices whose subobjects represent
columns, or ``flat'' matrices which do not have subobjects like rows or columns at all. The
different kinds of matrices have to be distinguished already with respect to
the definition of the operations for them. 

 In particular vector and matrix objects know their base domain (see \texttt{BaseDomain} (\ref{BaseDomain:for a vector object})) and their dimensions. The basic condition is that the entries of vector and
matrix objects must either lie in the base domain or naturally embed in the
sense that addition and multiplication automatically work with elements of the
base domain; for example, a matrix object over a polynomial ring may also
contain entries from the coefficient ring. 

 Vector and matrix objects may be mutable or immutable. Of course all
operations changing an object are only allowed/implemented for mutable
variants. 

 Vector objects are equal with respect to \texttt{\texttt{\symbol{92}}=} (\ref{bSlash=}) if they have the same length and the same entries. It is not necessary that
they have the same base domain. Matrices are equal with respect to \texttt{\texttt{\symbol{92}}=} (\ref{bSlash=}) if they have the same dimensions and the same entries. 

 For a row list matrix object, it is not guaranteed that all its rows have the
same vector type. It is for example thinkable that a matrix object stores some
of its rows in a sparse representation and some in a dense one. However, it is
guaranteed that the rows of two matrices in the same representation are
compatible in the sense that all vector operations defined in this interface
can be applied to them and that new matrices in the same representation as the
original matrix can be formed out of them. 

 Note that there is neither a default mapping from the set of matrix object
representations to the set of vector representations nor one in the reverse
direction. There is in general no ``associated'' vector object representation to a matrix object representation or vice versa.
(However, \texttt{CompatibleVectorFilter} (\ref{CompatibleVectorFilter:for a matrix object}) may describe a vector object representation that is compatible with a given
matrix object.) 

 The recommended way to write code that preserves the representation basically
works by using constructing operations that take template objects to decide
about the intended representation for the new object. 

 Vector and matrix objects do not have to be \textsf{GAP} lists in the sense of \texttt{IsList} (\ref{IsList}). Note that objects not in the filter \texttt{IsList} (\ref{IsList}) need not support all list operations, and their behaviour is not prescribed by
the rules for lists, e.g., behaviour w.r.t. arithmetic operations. However,
row list matrices behave nearly like lists of row vectors that insist on being
dense and containing only vectors of the same length and with the same base
domain. 

 Vector and matrix objects are not likely to benefit from \textsf{GAP}'s immediate methods (see section \ref{Immediate Methods}). Therefore it may be useful to set the filter \texttt{IsNoImmediateMethodsObject} (\ref{IsNoImmediateMethodsObject}) in the definition of new kinds of vector and matrix objects. 

 For information on how to implement new \texttt{IsMatrixObj} (\ref{IsMatrixObj}) and \texttt{IsVectorObj} (\ref{IsVectorObj}) representations see Section \ref{Implementing New Vector and Matrix Objects Types}. }

  
\section{\textcolor{Chapter }{Categories of Vector and Matrix Objects}}\label{Categories of Vector and Matrix Objects}
\logpage{[ 26, 2, 0 ]}
\hyperdef{L}{X7C6CDBFE7EB083A5}{}
{
  Currently the following categories of vector and matrix objects are supported
in \textsf{GAP}. More can be added as soon as there is need for them.  

\subsection{\textcolor{Chapter }{IsVectorObj}}
\logpage{[ 26, 2, 1 ]}\nobreak
\hyperdef{L}{X7D963FCC7E849BE0}{}
{\noindent\textcolor{FuncColor}{$\triangleright$\enspace\texttt{IsVectorObj({\mdseries\slshape obj})\index{IsVectorObj@\texttt{IsVectorObj}}
\label{IsVectorObj}
}\hfill{\scriptsize (Category)}}\\


 The idea behind \emph{vector objects} is that one wants to deal with objects like coefficient lists of fixed length
over a given domain $R$, say, which can be added and can be multiplied from the left with elements
from $R$. A vector object $v$, say, is always a copyable object (see \texttt{IsCopyable} (\ref{IsCopyable})) in \texttt{IsVector} (\ref{IsVector}), which knows the values of \texttt{BaseDomain} (\ref{BaseDomain:for a vector object}) (with value $R$) and \texttt{Length} (\ref{Length}), where $R$ is a domain (see Chapter \ref{Domains}) that has methods for \texttt{Zero} (\ref{Zero}), \texttt{One} (\ref{One}), \texttt{\texttt{\symbol{92}}in} (\ref{bSlashin:for a collection}), \texttt{Characteristic} (\ref{Characteristic}), \texttt{IsFinite} (\ref{IsFinite}). We say that $v$ is defined over $R$. Typically, $R$ will be at least a semiring. 

 For creating new vector objects compatible with $v$, \texttt{NewVector} (\ref{NewVector}) requires that also the value of \texttt{ConstructingFilter} (\ref{ConstructingFilter:for a vector object}) is known for $v$. 

 Further, entry access $v[i]$ is expected to return a \textsf{GAP} object, for $1 \leq i \leq$\texttt{ Length}$( v )$, and that these entries of $v$ belong to the base domain $R$. 

 Note that we do \emph{not} require that $v$ is a list in the sense of \texttt{IsList} (\ref{IsList}), in particular the rules of list arithmetic (see the sections \ref{Additive Arithmetic for Lists} and \ref{Multiplicative Arithmetic for Lists}) need \emph{not} hold. For example, the sum of two vector objects of different lengths or
defined over different base domains is not defined, and a plain list of vector
objects is not a matrix. Also unbinding entries of vector objects is not
defined. 

 Scalar multiplication from the left is defined only with elements from $R$. 

 The family of $v$ (see \texttt{FamilyObj} (\ref{FamilyObj})) is the same as the family of its base domain $R$. However, it is \emph{not} required that the entries lie in $R$ in the sense of \texttt{\texttt{\symbol{92}}in} (\ref{bSlashin:for a collection}), also values may occur that can be naturally embedded into $R$. For example, if $R$ is a polynomial ring then some entries in $v$ may be elements of the coefficient ring of $R$. }

 

\subsection{\textcolor{Chapter }{IsMatrixObj}}
\logpage{[ 26, 2, 2 ]}\nobreak
\hyperdef{L}{X7E7617A0781D1E4B}{}
{\noindent\textcolor{FuncColor}{$\triangleright$\enspace\texttt{IsMatrixObj({\mdseries\slshape obj})\index{IsMatrixObj@\texttt{IsMatrixObj}}
\label{IsMatrixObj}
}\hfill{\scriptsize (Category)}}\\


 The idea behind \emph{matrix objects} is that one wants to deal with objects like $m$ by $n$ arrays over a given domain $R$, say, which can be added and multiplied and can be multiplied from the left
with elements from $R$. A matrix object $M$, say, is always a copyable object (see \texttt{IsCopyable} (\ref{IsCopyable})) in \texttt{IsVector} (\ref{IsVector}) and \texttt{IsScalar} (\ref{IsScalar}), which knows the values of \texttt{BaseDomain} (\ref{BaseDomain:for a matrix object}) (with value $R$), \texttt{NumberRows} (\ref{NumberRows:for a matrix object}) (with value $m$), \texttt{NumberColumns} (\ref{NumberColumns:for a matrix object}) (with value $n$), where $R$ is a domain (see Chapter \ref{Domains}) that has methods for \texttt{Zero} (\ref{Zero}), \texttt{One} (\ref{One}), \texttt{\texttt{\symbol{92}}in} (\ref{bSlashin:for a collection}), \texttt{Characteristic} (\ref{Characteristic}), \texttt{IsFinite} (\ref{IsFinite}). We say that $v$ is defined over $R$. Typically, $R$ will be at least a semiring. 

 For creating new matrix objects compatible with $M$, \texttt{NewMatrix} (\ref{NewMatrix}) requires that also the value of \texttt{ConstructingFilter} (\ref{ConstructingFilter:for a matrix object}) is known for $M$. 

 Further, entry access $M[i,j]$ is expected to return a \textsf{GAP} object, for $1 \leq i \leq m$ and $1 \leq j \leq n$, and that these entries of $M$ belong to the base domain $R$. 

 Note that we do \emph{not} require that $M$ is a list in the sense of \texttt{IsList} (\ref{IsList}), in particular the rules of list arithmetic (see the sections \ref{Additive Arithmetic for Lists} and \ref{Multiplicative Arithmetic for Lists}) need \emph{not} hold. For example, accessing ``rows'' of $M$ via \texttt{\texttt{\symbol{92}}[\texttt{\symbol{92}}]} (\ref{bSlash[bSlash]}) is in general not possible, and the sum of two matrix objects with different
numbers of rows or columns is not defined. Also unbinding entries of matrix
objects is not defined. 

 Scalar multiplication from the left is defined only with elements from $R$. 

 It is not assumed that the multiplication in $R$ is associative, and we do not define what the $k$\texttt{\symbol{45}}th power of a matrix object is in this case, for positive
integers $k$. (However, a default powering method is available.) 

 The filter \texttt{IsMatrixObj} alone does \emph{not} imply that the multiplication is the usual matrix multiplication. This
multiplication can be defined via the filter \texttt{IsOrdinaryMatrix} (\ref{IsOrdinaryMatrix}); this filter together with the associativity of the base domain also implies
the associativity of matrix multiplication. For example, elements of matrix
Lie algebras (see \texttt{LieObject} (\ref{LieObject})) lie in \texttt{IsMatrixObj} but not in \texttt{IsOrdinaryMatrix} (\ref{IsOrdinaryMatrix}). 

 The family of $M$ (see \texttt{FamilyObj} (\ref{FamilyObj})) is the collections family (see \texttt{CollectionsFamily} (\ref{CollectionsFamily})) of its base domain $R$. However, it is \emph{not} required that the entries lie in $R$ in the sense of \texttt{\texttt{\symbol{92}}in} (\ref{bSlashin:for a collection}), also values may occur that can be naturally embedded into $R$. For example, if $R$ is a polynomial ring then some entries in $M$ may be elements of the coefficient ring of $R$. }

 

\subsection{\textcolor{Chapter }{IsMatrixOrMatrixObj}}
\logpage{[ 26, 2, 3 ]}\nobreak
\hyperdef{L}{X7D1ACCBE7E9CF501}{}
{\noindent\textcolor{FuncColor}{$\triangleright$\enspace\texttt{IsMatrixOrMatrixObj({\mdseries\slshape obj})\index{IsMatrixOrMatrixObj@\texttt{IsMatrixOrMatrixObj}}
\label{IsMatrixOrMatrixObj}
}\hfill{\scriptsize (Category)}}\\


 Several functions are defined for objects in \texttt{IsMatrix} (\ref{IsMatrix}) and objects in \texttt{IsMatrixObj} (\ref{IsMatrixObj}). All these objects lie in the filter \texttt{IsMatrixOrMatrixObj}. It should be used in situations where an object can be either a list of
lists in \texttt{IsMatrix} (\ref{IsMatrix}) or a ``proper'' matrix object in \texttt{IsMatrixObj} (\ref{IsMatrixObj}), for example as a requirement in the installation of a method for such an
argument. 

 
\begin{Verbatim}[commandchars=!@|,fontsize=\small,frame=single,label=Example]
  !gapprompt@gap>| !gapinput@m:= IdentityMat( 2, GF(2) );;|
  !gapprompt@gap>| !gapinput@IsMatrix( m );  IsMatrixObj( m ); IsMatrixOrMatrixObj( m );|
  true
  false
  true
  !gapprompt@gap>| !gapinput@m:= NewIdentityMatrix( IsPlistMatrixRep, GF(2), 2 );;|
  !gapprompt@gap>| !gapinput@IsMatrix( m );  IsMatrixObj( m ); IsMatrixOrMatrixObj( m );|
  false
  true
  true
\end{Verbatim}
 }

 

\subsection{\textcolor{Chapter }{IsRowListMatrix}}
\logpage{[ 26, 2, 4 ]}\nobreak
\hyperdef{L}{X78CD88A283330E72}{}
{\noindent\textcolor{FuncColor}{$\triangleright$\enspace\texttt{IsRowListMatrix({\mdseries\slshape obj})\index{IsRowListMatrix@\texttt{IsRowListMatrix}}
\label{IsRowListMatrix}
}\hfill{\scriptsize (Category)}}\\


 A \emph{row list matrix object} is a matrix object (see \texttt{IsMatrixObj} (\ref{IsMatrixObj})) $M$ which admits access to its rows, that is, list access $M[i]$ (see \texttt{\texttt{\symbol{92}}[\texttt{\symbol{92}}]} (\ref{bSlash[bSlash]})) yields the $i$\texttt{\symbol{45}}th row of $M$, for $1 \leq i \leq$ \texttt{NumberRows( }$M$\texttt{ )}. 

 All rows are \texttt{IsVectorObj} (\ref{IsVectorObj}) objects in the same representation. Several rows of a row list matrix object
can be identical objects, and different row list matrices may share rows. Row
access just gives a reference to the row object, without copying the row. 

 Matrix objects in \texttt{IsRowListMatrix} are \emph{not} necessarily in \texttt{IsList} (\ref{IsList}), and then they need not obey the general rules for lists. }

 }

  
\section{\textcolor{Chapter }{Defining Attributes of Vector and Matrix Objects}}\label{Defining Attributes of Vector and Matrix Objects}
\logpage{[ 26, 3, 0 ]}
\hyperdef{L}{X877A706186C89ADB}{}
{
  
\subsection{\textcolor{Chapter }{BaseDomain}}\logpage{[ 26, 3, 1 ]}
\hyperdef{L}{X8662026C7CCDB446}{}
{
\noindent\textcolor{FuncColor}{$\triangleright$\enspace\texttt{BaseDomain({\mdseries\slshape vector})\index{BaseDomain@\texttt{BaseDomain}!for a vector object}
\label{BaseDomain:for a vector object}
}\hfill{\scriptsize (attribute)}}\\
\noindent\textcolor{FuncColor}{$\triangleright$\enspace\texttt{BaseDomain({\mdseries\slshape matrix})\index{BaseDomain@\texttt{BaseDomain}!for a matrix object}
\label{BaseDomain:for a matrix object}
}\hfill{\scriptsize (attribute)}}\\


 The vector object \mbox{\texttt{\mdseries\slshape vector}} or matrix object \mbox{\texttt{\mdseries\slshape matrix}}, respectively, is defined over the domain given by its \texttt{BaseDomain} value. 

 Note that not all entries of the object necessarily lie in its base domain
with respect to \texttt{\texttt{\symbol{92}}in} (\ref{bSlashin:for a collection}), see Section \ref{Concepts and Rules for Vector and Matrix Objects}. }

 
\subsection{\textcolor{Chapter }{ConstructingFilter}}\logpage{[ 26, 3, 2 ]}
\hyperdef{L}{X85ABF33684865ED5}{}
{
\noindent\textcolor{FuncColor}{$\triangleright$\enspace\texttt{ConstructingFilter({\mdseries\slshape v})\index{ConstructingFilter@\texttt{ConstructingFilter}!for a vector object}
\label{ConstructingFilter:for a vector object}
}\hfill{\scriptsize (attribute)}}\\
\noindent\textcolor{FuncColor}{$\triangleright$\enspace\texttt{ConstructingFilter({\mdseries\slshape M})\index{ConstructingFilter@\texttt{ConstructingFilter}!for a matrix object}
\label{ConstructingFilter:for a matrix object}
}\hfill{\scriptsize (attribute)}}\\
\textbf{\indent Returns:\ }
a filter



 Called with a vector object \mbox{\texttt{\mdseries\slshape v}} or a matrix object \mbox{\texttt{\mdseries\slshape M}}, respectively, \texttt{ConstructingFilter} returns a filter \texttt{f} such that when \texttt{NewVector} (\ref{NewVector}) or \texttt{NewMatrix} (\ref{NewMatrix}), respectively, is called with \texttt{f} then a vector object or a matrix object, respectively, in the same
representation as the argument is produced. 

 If the \texttt{ConstructingFilter} value of \mbox{\texttt{\mdseries\slshape v}} or \mbox{\texttt{\mdseries\slshape M}} implies \texttt{IsCopyable} (\ref{IsCopyable}) then mutable versions of \mbox{\texttt{\mdseries\slshape v}} or \mbox{\texttt{\mdseries\slshape M}} can be created, otherwise all vector or matrix objects with this filter are
immutable. }

 
\subsection{\textcolor{Chapter }{CompatibleVectorFilter}}\logpage{[ 26, 3, 3 ]}
\hyperdef{L}{X818702FD7A2E9D90}{}
{
\noindent\textcolor{FuncColor}{$\triangleright$\enspace\texttt{CompatibleVectorFilter({\mdseries\slshape M})\index{CompatibleVectorFilter@\texttt{CompatibleVectorFilter}!for a matrix object}
\label{CompatibleVectorFilter:for a matrix object}
}\hfill{\scriptsize (attribute)}}\\
\textbf{\indent Returns:\ }
a filter



 Called with a matrix object \mbox{\texttt{\mdseries\slshape M}}, \texttt{CompatibleVectorFilter} returns either a filter \texttt{f} such that vector objects with \texttt{ConstructingFilter} (\ref{ConstructingFilter:for a vector object}) value \texttt{f} are compatible in the sense that \mbox{\texttt{\mdseries\slshape M}} can be multiplied with these vector objects, of \texttt{fail} if no such filter is known. }

 

\subsection{\textcolor{Chapter }{Length (for a vector object)}}
\logpage{[ 26, 3, 4 ]}\nobreak
\hyperdef{L}{X828BA5E1849E3D06}{}
{\noindent\textcolor{FuncColor}{$\triangleright$\enspace\texttt{Length({\mdseries\slshape v})\index{Length@\texttt{Length}!for a vector object}
\label{Length:for a vector object}
}\hfill{\scriptsize (attribute)}}\\


 returns the length of the vector object \mbox{\texttt{\mdseries\slshape v}}, which is defined to be the number of entries of \mbox{\texttt{\mdseries\slshape v}}. }

 
\subsection{\textcolor{Chapter }{NumberRows and NumberColumns}}\logpage{[ 26, 3, 5 ]}
\hyperdef{L}{X820ED34380C10E19}{}
{
\noindent\textcolor{FuncColor}{$\triangleright$\enspace\texttt{NumberRows({\mdseries\slshape M})\index{NumberRows@\texttt{NumberRows}!for a matrix object}
\label{NumberRows:for a matrix object}
}\hfill{\scriptsize (attribute)}}\\
\noindent\textcolor{FuncColor}{$\triangleright$\enspace\texttt{NrRows({\mdseries\slshape M})\index{NrRows@\texttt{NrRows}!for a matrix object}
\label{NrRows:for a matrix object}
}\hfill{\scriptsize (attribute)}}\\
\noindent\textcolor{FuncColor}{$\triangleright$\enspace\texttt{NumberColumns({\mdseries\slshape M})\index{NumberColumns@\texttt{NumberColumns}!for a matrix object}
\label{NumberColumns:for a matrix object}
}\hfill{\scriptsize (attribute)}}\\
\noindent\textcolor{FuncColor}{$\triangleright$\enspace\texttt{NrCols({\mdseries\slshape M})\index{NrCols@\texttt{NrCols}!for a matrix object}
\label{NrCols:for a matrix object}
}\hfill{\scriptsize (attribute)}}\\


 For a matrix object \mbox{\texttt{\mdseries\slshape M}}, \texttt{NumberRows} and \texttt{NumberColumns} store the number of rows and columns of \mbox{\texttt{\mdseries\slshape M}}, respectively. 

 \texttt{NrRows} and \texttt{NrCols} are synonyms of \texttt{NumberRows} and \texttt{NumberColumns}, respectively. }

 }

  
\section{\textcolor{Chapter }{Constructing Vector and Matrix Objects}}\label{Constructing Vector and Matrix Objects}
\logpage{[ 26, 4, 0 ]}
\hyperdef{L}{X7BD7D2837BFDE649}{}
{
  
\subsection{\textcolor{Chapter }{NewVector and NewZeroVector}}\logpage{[ 26, 4, 1 ]}
\hyperdef{L}{X860E84397BD148E9}{}
{
\noindent\textcolor{FuncColor}{$\triangleright$\enspace\texttt{NewVector({\mdseries\slshape filt, R, list})\index{NewVector@\texttt{NewVector}}
\label{NewVector}
}\hfill{\scriptsize (operation)}}\\
\noindent\textcolor{FuncColor}{$\triangleright$\enspace\texttt{NewZeroVector({\mdseries\slshape filt, R, n})\index{NewZeroVector@\texttt{NewZeroVector}}
\label{NewZeroVector}
}\hfill{\scriptsize (operation)}}\\


 For a filter \mbox{\texttt{\mdseries\slshape filt}}, a semiring \mbox{\texttt{\mdseries\slshape R}}, and a list \mbox{\texttt{\mdseries\slshape list}} of elements that belong to \mbox{\texttt{\mdseries\slshape R}}, \texttt{NewVector} returns a vector object which has the \texttt{ConstructingFilter} (\ref{ConstructingFilter:for a vector object}) \mbox{\texttt{\mdseries\slshape filt}}, the \texttt{BaseDomain} (\ref{BaseDomain:for a vector object}) \mbox{\texttt{\mdseries\slshape R}}, and the entries in \mbox{\texttt{\mdseries\slshape list}}. The list \mbox{\texttt{\mdseries\slshape list}} is guaranteed not to be changed by this operation. 

 If the global option \texttt{check} is set to \texttt{false} then \texttt{NewVector} need not perform consistency checks. 

 Similarly, \texttt{NewZeroVector} returns a vector object of length \mbox{\texttt{\mdseries\slshape n}} which has \mbox{\texttt{\mdseries\slshape filt}} and \mbox{\texttt{\mdseries\slshape R}} as \texttt{ConstructingFilter} (\ref{ConstructingFilter:for a vector object}) and \texttt{BaseDomain} (\ref{BaseDomain:for a vector object}) values, and contains the zero of \mbox{\texttt{\mdseries\slshape R}} in each position. 

 The returned object is mutable if and only if \mbox{\texttt{\mdseries\slshape filt}} implies \texttt{IsCopyable} (\ref{IsCopyable}). }

 
\subsection{\textcolor{Chapter }{Vector}}\logpage{[ 26, 4, 2 ]}
\hyperdef{L}{X79A6544D86261E82}{}
{
\noindent\textcolor{FuncColor}{$\triangleright$\enspace\texttt{Vector({\mdseries\slshape filt, R, list})\index{Vector@\texttt{Vector}!for filter, base domain, and list}
\label{Vector:for filter, base domain, and list}
}\hfill{\scriptsize (operation)}}\\
\noindent\textcolor{FuncColor}{$\triangleright$\enspace\texttt{Vector({\mdseries\slshape filt, R, v})\index{Vector@\texttt{Vector}!for filter, base domain, and vector object}
\label{Vector:for filter, base domain, and vector object}
}\hfill{\scriptsize (operation)}}\\
\noindent\textcolor{FuncColor}{$\triangleright$\enspace\texttt{Vector({\mdseries\slshape R, list})\index{Vector@\texttt{Vector}!for base domain and list}
\label{Vector:for base domain and list}
}\hfill{\scriptsize (operation)}}\\
\noindent\textcolor{FuncColor}{$\triangleright$\enspace\texttt{Vector({\mdseries\slshape R, v})\index{Vector@\texttt{Vector}!for base domain and vector object}
\label{Vector:for base domain and vector object}
}\hfill{\scriptsize (operation)}}\\
\noindent\textcolor{FuncColor}{$\triangleright$\enspace\texttt{Vector({\mdseries\slshape list, v})\index{Vector@\texttt{Vector}!for a list and a vector object}
\label{Vector:for a list and a vector object}
}\hfill{\scriptsize (operation)}}\\
\noindent\textcolor{FuncColor}{$\triangleright$\enspace\texttt{Vector({\mdseries\slshape v1, v2})\index{Vector@\texttt{Vector}!for two vector objects}
\label{Vector:for two vector objects}
}\hfill{\scriptsize (operation)}}\\
\noindent\textcolor{FuncColor}{$\triangleright$\enspace\texttt{Vector({\mdseries\slshape list})\index{Vector@\texttt{Vector}!for a list}
\label{Vector:for a list}
}\hfill{\scriptsize (operation)}}\\
\textbf{\indent Returns:\ }
a vector object



 If a filter \mbox{\texttt{\mdseries\slshape filt}} is given as the first argument then a vector object is returned that has \texttt{ConstructingFilter} (\ref{ConstructingFilter:for a vector object}) value \mbox{\texttt{\mdseries\slshape filt}}, is defined over the base domain \mbox{\texttt{\mdseries\slshape R}}, and has the entries given by the list \mbox{\texttt{\mdseries\slshape list}} or the vector object \mbox{\texttt{\mdseries\slshape v}}, respectively. 

 If a semiring \mbox{\texttt{\mdseries\slshape R}} is given as the first argument then a vector object is returned whose \texttt{ConstructingFilter} (\ref{ConstructingFilter:for a vector object}) value is guessed from \mbox{\texttt{\mdseries\slshape R}}, again with base domain \mbox{\texttt{\mdseries\slshape R}} and entries given by the last argument. 

 In the remaining cases with two arguments, the first argument is a list or a
vector object that defines the entries of the result, and the second argument
is a vector object whose \texttt{ConstructingFilter} (\ref{ConstructingFilter:for a vector object}) and \texttt{BaseDomain} (\ref{BaseDomain:for a vector object}) are taken for the result. 

 If only a list \mbox{\texttt{\mdseries\slshape list}} is given then both the \texttt{ConstructingFilter} (\ref{ConstructingFilter:for a vector object}) and the \texttt{BaseDomain} (\ref{BaseDomain:for a vector object}) are guessed from this list. 

 The variant \texttt{Vector( }\mbox{\texttt{\mdseries\slshape v1}}\texttt{, }\mbox{\texttt{\mdseries\slshape v2}}\texttt{ )} is supported also for the case that \mbox{\texttt{\mdseries\slshape v2}} is a row vector but not a vector object. In this situation, the result is a
row vector that is equal to \mbox{\texttt{\mdseries\slshape v1}} and whose internal representation fits to that of \mbox{\texttt{\mdseries\slshape v2}}. 

 If the global option \texttt{check} is set to \texttt{false} then \texttt{Vector} need not perform consistency checks. 

 If the \texttt{ConstructingFilter} (\ref{ConstructingFilter:for a vector object}) value of the result implies \texttt{IsCopyable} (\ref{IsCopyable}) then the result is mutable if and only if the argument that determines the
entries of the result (\mbox{\texttt{\mdseries\slshape list}}, \mbox{\texttt{\mdseries\slshape v}}, \mbox{\texttt{\mdseries\slshape v1}}) is mutable. 

 In the case of a mutable result, it is \emph{not} guaranteed that the given list of entries is copied. 

 Default methods for \texttt{Vector} delegate to \texttt{NewVector} (\ref{NewVector}). }

 
\subsection{\textcolor{Chapter }{ZeroVector}}\logpage{[ 26, 4, 3 ]}
\hyperdef{L}{X7DBA8BF5844F3281}{}
{
\noindent\textcolor{FuncColor}{$\triangleright$\enspace\texttt{ZeroVector({\mdseries\slshape filt, R, len})\index{ZeroVector@\texttt{ZeroVector}!for filter, base domain and length}
\label{ZeroVector:for filter, base domain and length}
}\hfill{\scriptsize (operation)}}\\
\noindent\textcolor{FuncColor}{$\triangleright$\enspace\texttt{ZeroVector({\mdseries\slshape R, len})\index{ZeroVector@\texttt{ZeroVector}!for base domain and length}
\label{ZeroVector:for base domain and length}
}\hfill{\scriptsize (operation)}}\\
\noindent\textcolor{FuncColor}{$\triangleright$\enspace\texttt{ZeroVector({\mdseries\slshape len, v})\index{ZeroVector@\texttt{ZeroVector}!for length and vector object}
\label{ZeroVector:for length and vector object}
}\hfill{\scriptsize (operation)}}\\
\noindent\textcolor{FuncColor}{$\triangleright$\enspace\texttt{ZeroVector({\mdseries\slshape len, M})\index{ZeroVector@\texttt{ZeroVector}!for length and matrix object}
\label{ZeroVector:for length and matrix object}
}\hfill{\scriptsize (operation)}}\\
\textbf{\indent Returns:\ }
a vector object



 For a filter \mbox{\texttt{\mdseries\slshape filt}}, a semiring \mbox{\texttt{\mdseries\slshape R}} and a nonnegative integer \mbox{\texttt{\mdseries\slshape len}}, this operation returns a new vector object of length \mbox{\texttt{\mdseries\slshape len}} over \mbox{\texttt{\mdseries\slshape R}} in the representation \mbox{\texttt{\mdseries\slshape filt}} containing only zeros. 

 If only \mbox{\texttt{\mdseries\slshape R}} and \mbox{\texttt{\mdseries\slshape len}} are given, then \textsf{GAP} guesses a suitable representation. 

 For a vector object \mbox{\texttt{\mdseries\slshape v}} and a nonnegative integer \mbox{\texttt{\mdseries\slshape len}}, this operation returns a new vector object of length \mbox{\texttt{\mdseries\slshape len}} in the same representation as \mbox{\texttt{\mdseries\slshape v}} containing only zeros. 

 For a matrix object \mbox{\texttt{\mdseries\slshape M}} and a nonnegative integer \mbox{\texttt{\mdseries\slshape len}}, this operation returns a new zero vector object of length \mbox{\texttt{\mdseries\slshape len}} in the representation given by the \texttt{CompatibleVectorFilter} (\ref{CompatibleVectorFilter:for a matrix object}) value of \mbox{\texttt{\mdseries\slshape M}}, provided that such a representation exists. 

 If the \texttt{ConstructingFilter} (\ref{ConstructingFilter:for a vector object}) value of the result implies \texttt{IsCopyable} (\ref{IsCopyable}) then the result is mutable. 

 Default methods for \texttt{ZeroVector} delegate to \texttt{NewZeroVector} (\ref{NewZeroVector}). }

 
\subsection{\textcolor{Chapter }{NewMatrix, NewZeroMatrix, NewIdentityMatrix}}\logpage{[ 26, 4, 4 ]}
\hyperdef{L}{X7AD2210B8047FB01}{}
{
\noindent\textcolor{FuncColor}{$\triangleright$\enspace\texttt{NewMatrix({\mdseries\slshape filt, R, ncols, list})\index{NewMatrix@\texttt{NewMatrix}}
\label{NewMatrix}
}\hfill{\scriptsize (operation)}}\\
\noindent\textcolor{FuncColor}{$\triangleright$\enspace\texttt{NewZeroMatrix({\mdseries\slshape filt, R, m, n})\index{NewZeroMatrix@\texttt{NewZeroMatrix}}
\label{NewZeroMatrix}
}\hfill{\scriptsize (operation)}}\\
\noindent\textcolor{FuncColor}{$\triangleright$\enspace\texttt{NewIdentityMatrix({\mdseries\slshape filt, R, n})\index{NewIdentityMatrix@\texttt{NewIdentityMatrix}}
\label{NewIdentityMatrix}
}\hfill{\scriptsize (operation)}}\\


 For a filter \mbox{\texttt{\mdseries\slshape filt}}, a semiring \mbox{\texttt{\mdseries\slshape R}}, a positive integer \mbox{\texttt{\mdseries\slshape ncols}}, and a list \mbox{\texttt{\mdseries\slshape list}}, \texttt{NewMatrix} returns a matrix object which has the \texttt{ConstructingFilter} (\ref{ConstructingFilter:for a vector object}) \mbox{\texttt{\mdseries\slshape filt}}, the \texttt{BaseDomain} (\ref{BaseDomain:for a matrix object}) \mbox{\texttt{\mdseries\slshape R}}, \mbox{\texttt{\mdseries\slshape n}} columns (see \texttt{NumberColumns} (\ref{NumberColumns:for a matrix object})), and the entries described by \mbox{\texttt{\mdseries\slshape list}}, which can be either a plain list of vector objects of length \mbox{\texttt{\mdseries\slshape ncols}} or a plain list of plain lists of length \mbox{\texttt{\mdseries\slshape ncols}} or a plain list of length a multiple of \mbox{\texttt{\mdseries\slshape ncols}} containing the entries in row major order. The list \mbox{\texttt{\mdseries\slshape list}} is guaranteed not to be changed by this operation. 

 The corresponding entries must be in or compatible with \mbox{\texttt{\mdseries\slshape R}}. If \mbox{\texttt{\mdseries\slshape list}} already contains vector objects, they are copied. 

 If the global option \texttt{check} is set to \texttt{false} then \texttt{NewMatrix} need not perform consistency checks. 

 Similarly, \texttt{NewZeroMatrix} returns a zero matrix object with \mbox{\texttt{\mdseries\slshape m}} rows and \mbox{\texttt{\mdseries\slshape n}} columns which has \mbox{\texttt{\mdseries\slshape filt}} and \mbox{\texttt{\mdseries\slshape R}} as \texttt{ConstructingFilter} (\ref{ConstructingFilter:for a vector object}) and \texttt{BaseDomain} (\ref{BaseDomain:for a vector object}) values. 

 Similarly, \texttt{NewIdentityMatrix} returns an identity matrix object with \mbox{\texttt{\mdseries\slshape n}} rows and columns which has \mbox{\texttt{\mdseries\slshape filt}} and \mbox{\texttt{\mdseries\slshape R}} as \texttt{ConstructingFilter} (\ref{ConstructingFilter:for a vector object}) and \texttt{BaseDomain} (\ref{BaseDomain:for a vector object}) values, and contains the identity element of \mbox{\texttt{\mdseries\slshape R}} in the diagonal and the zero of \mbox{\texttt{\mdseries\slshape R}} in each off\texttt{\symbol{45}}diagonal position. 

 The returned object is mutable if and only if \mbox{\texttt{\mdseries\slshape filt}} implies \texttt{IsCopyable} (\ref{IsCopyable}). }

 
\subsection{\textcolor{Chapter }{Matrix}}\logpage{[ 26, 4, 5 ]}
\hyperdef{L}{X879384D479EB1D82}{}
{
\noindent\textcolor{FuncColor}{$\triangleright$\enspace\texttt{Matrix({\mdseries\slshape [filt, ]R, list[, ncols]})\index{Matrix@\texttt{Matrix}!for filter, base domain, list, ncols}
\label{Matrix:for filter, base domain, list, ncols}
}\hfill{\scriptsize (operation)}}\\
\noindent\textcolor{FuncColor}{$\triangleright$\enspace\texttt{Matrix({\mdseries\slshape [filt, ]R, matobj})\index{Matrix@\texttt{Matrix}!for filter, base domain, and matrix object}
\label{Matrix:for filter, base domain, and matrix object}
}\hfill{\scriptsize (operation)}}\\
\noindent\textcolor{FuncColor}{$\triangleright$\enspace\texttt{Matrix({\mdseries\slshape list[, ncols], example{\textunderscore}matobj})\index{Matrix@\texttt{Matrix}!for a list, ncols, and a matrix object}
\label{Matrix:for a list, ncols, and a matrix object}
}\hfill{\scriptsize (operation)}}\\
\noindent\textcolor{FuncColor}{$\triangleright$\enspace\texttt{Matrix({\mdseries\slshape matobj, example{\textunderscore}matobj})\index{Matrix@\texttt{Matrix}!for two matrix objects}
\label{Matrix:for two matrix objects}
}\hfill{\scriptsize (operation)}}\\
\noindent\textcolor{FuncColor}{$\triangleright$\enspace\texttt{Matrix({\mdseries\slshape list[, ncols]})\index{Matrix@\texttt{Matrix}!for a list and ncols}
\label{Matrix:for a list and ncols}
}\hfill{\scriptsize (operation)}}\\
\textbf{\indent Returns:\ }
a matrix object



 If a filter \mbox{\texttt{\mdseries\slshape filt}} is given as the first argument then a matrix object is returned that has \texttt{ConstructingFilter} (\ref{ConstructingFilter:for a matrix object}) value \mbox{\texttt{\mdseries\slshape filt}}, is defined over the base domain \mbox{\texttt{\mdseries\slshape R}}, and has the entries given by the list \mbox{\texttt{\mdseries\slshape list}} or the matrix object \mbox{\texttt{\mdseries\slshape matobj}}, respectively. Here \mbox{\texttt{\mdseries\slshape list}} can be either a list of plain lists that describe the entries of the rows, or
a flat list of the entries in row major order, where \mbox{\texttt{\mdseries\slshape ncols}} defines the number of columns. 

 If a semiring \mbox{\texttt{\mdseries\slshape R}} is given as the first argument then a matrix object is returned whose \texttt{ConstructingFilter} (\ref{ConstructingFilter:for a matrix object}) value is guessed from \mbox{\texttt{\mdseries\slshape R}}, again with base domain \mbox{\texttt{\mdseries\slshape R}} and entries given by the last argument. 

 In those remaining cases where the last argument is a matrix object, the first
argument is a list or a matrix object that defines (together with \mbox{\texttt{\mdseries\slshape ncols}} if applicable) the entries of the result, and the \texttt{ConstructingFilter} (\ref{ConstructingFilter:for a matrix object}) and \texttt{BaseDomain} (\ref{BaseDomain:for a matrix object}) of the last argument are taken for the result. 

 Finally, if only a list \mbox{\texttt{\mdseries\slshape list}} and perhaps \mbox{\texttt{\mdseries\slshape ncols}} is given then both the \texttt{ConstructingFilter} (\ref{ConstructingFilter:for a matrix object}) and the \texttt{BaseDomain} (\ref{BaseDomain:for a vector object}) are guessed from the list. 

 If the global option \texttt{check} is set to \texttt{false} then \texttt{Matrix} need not perform consistency checks. 

 If the \texttt{ConstructingFilter} (\ref{ConstructingFilter:for a matrix object}) value of the result implies \texttt{IsCopyable} (\ref{IsCopyable}) then the result is mutable if and only if the argument that determines the
entries of the result (\mbox{\texttt{\mdseries\slshape list}} or \mbox{\texttt{\mdseries\slshape matobj}}) is mutable. 

 In the case of a mutable result, it is guaranteed that the given list \mbox{\texttt{\mdseries\slshape list}} is copied in the sense of \texttt{ShallowCopy} (\ref{ShallowCopy}), and if \mbox{\texttt{\mdseries\slshape list}} is a nested list then it is \emph{not} guaranteed that also the entries of \mbox{\texttt{\mdseries\slshape list}} are copied. 

 Default methods for \texttt{Matrix} delegate to \texttt{NewMatrix} (\ref{NewMatrix}). }

 
\subsection{\textcolor{Chapter }{ZeroMatrix}}\logpage{[ 26, 4, 6 ]}
\hyperdef{L}{X838F5B6C7C87C8E1}{}
{
\noindent\textcolor{FuncColor}{$\triangleright$\enspace\texttt{ZeroMatrix({\mdseries\slshape m, n, M})\index{ZeroMatrix@\texttt{ZeroMatrix}!for dimensions and matrix object}
\label{ZeroMatrix:for dimensions and matrix object}
}\hfill{\scriptsize (operation)}}\\
\noindent\textcolor{FuncColor}{$\triangleright$\enspace\texttt{ZeroMatrix({\mdseries\slshape R, m, n})\index{ZeroMatrix@\texttt{ZeroMatrix}!for base domain and dimensions}
\label{ZeroMatrix:for base domain and dimensions}
}\hfill{\scriptsize (operation)}}\\
\noindent\textcolor{FuncColor}{$\triangleright$\enspace\texttt{ZeroMatrix({\mdseries\slshape filt, R, m, n})\index{ZeroMatrix@\texttt{ZeroMatrix}!for filter, base domain, and dimensions}
\label{ZeroMatrix:for filter, base domain, and dimensions}
}\hfill{\scriptsize (operation)}}\\
\textbf{\indent Returns:\ }
a matrix object



 For a matrix object \mbox{\texttt{\mdseries\slshape M}} and two nonnegative integers \mbox{\texttt{\mdseries\slshape m}} and \mbox{\texttt{\mdseries\slshape n}}, this operation returns a new matrix object with \mbox{\texttt{\mdseries\slshape m}} rows and \mbox{\texttt{\mdseries\slshape n}} columns in the same representation and over the same base domain as \mbox{\texttt{\mdseries\slshape M}} containing only zeros. 

 If a semiring \mbox{\texttt{\mdseries\slshape R}} and two nonnegative integers \mbox{\texttt{\mdseries\slshape m}} and \mbox{\texttt{\mdseries\slshape n}} are given, the representation of the result is guessed from \mbox{\texttt{\mdseries\slshape R}}. 

 If a filter \mbox{\texttt{\mdseries\slshape filt}} and a semiring \mbox{\texttt{\mdseries\slshape R}} are given as the first and second argument, they are taken as the values of \texttt{ConstructingFilter} (\ref{ConstructingFilter:for a matrix object}) and \texttt{BaseDomain} (\ref{BaseDomain:for a matrix object}) of the result. 

 If the \texttt{ConstructingFilter} (\ref{ConstructingFilter:for a matrix object}) value of the result implies \texttt{IsCopyable} (\ref{IsCopyable}) then the result is fully mutable. 

 Default methods for \texttt{ZeroMatrix} delegate to \texttt{NewZeroMatrix} (\ref{NewZeroMatrix}). }

 
\subsection{\textcolor{Chapter }{IdentityMatrix}}\logpage{[ 26, 4, 7 ]}
\hyperdef{L}{X7D807ABC7FCB4E77}{}
{
\noindent\textcolor{FuncColor}{$\triangleright$\enspace\texttt{IdentityMatrix({\mdseries\slshape n, M})\index{IdentityMatrix@\texttt{IdentityMatrix}!for dimension and matrix object}
\label{IdentityMatrix:for dimension and matrix object}
}\hfill{\scriptsize (operation)}}\\
\noindent\textcolor{FuncColor}{$\triangleright$\enspace\texttt{IdentityMatrix({\mdseries\slshape R, n})\index{IdentityMatrix@\texttt{IdentityMatrix}!for base domain and dimension}
\label{IdentityMatrix:for base domain and dimension}
}\hfill{\scriptsize (operation)}}\\
\noindent\textcolor{FuncColor}{$\triangleright$\enspace\texttt{IdentityMatrix({\mdseries\slshape filt, R, n})\index{IdentityMatrix@\texttt{IdentityMatrix}!for filter, base domain, and dimension}
\label{IdentityMatrix:for filter, base domain, and dimension}
}\hfill{\scriptsize (operation)}}\\
\textbf{\indent Returns:\ }
a matrix object



 For a matrix object \mbox{\texttt{\mdseries\slshape M}} and a nonnegative integer \mbox{\texttt{\mdseries\slshape n}}, this operation returns a new identity matrix object with \mbox{\texttt{\mdseries\slshape n}} rows and columns in the same representation and over the same base domain as \mbox{\texttt{\mdseries\slshape M}}. 

 If a semiring \mbox{\texttt{\mdseries\slshape R}} and a nonnegative integer \mbox{\texttt{\mdseries\slshape n}} is given, the representation of the result is guessed from \mbox{\texttt{\mdseries\slshape R}}. 

 If a filter \mbox{\texttt{\mdseries\slshape filt}} and a semiring \mbox{\texttt{\mdseries\slshape R}} are given as the first and second argument, they are taken as the values of \texttt{ConstructingFilter} (\ref{ConstructingFilter:for a matrix object}) and \texttt{BaseDomain} (\ref{BaseDomain:for a matrix object}) of the result. 

 If the \texttt{ConstructingFilter} (\ref{ConstructingFilter:for a matrix object}) value of the result implies \texttt{IsCopyable} (\ref{IsCopyable}) then the result is fully mutable. 

 Default methods for \texttt{IdentityMatrix} delegate to \texttt{NewIdentityMatrix} (\ref{NewIdentityMatrix}). }

 }

  
\section{\textcolor{Chapter }{Operations for Base Domains of Vector and Matrix Objects}}\label{Operations for Base Domains of Vector and Matrix Objects}
\logpage{[ 26, 5, 0 ]}
\hyperdef{L}{X7C7F5250855C4371}{}
{
  
\subsection{\textcolor{Chapter }{OneOfBaseDomain and ZeroOfBaseDomain}}\logpage{[ 26, 5, 1 ]}
\hyperdef{L}{X85D7A6A782B21E5C}{}
{
\noindent\textcolor{FuncColor}{$\triangleright$\enspace\texttt{OneOfBaseDomain({\mdseries\slshape v})\index{OneOfBaseDomain@\texttt{OneOfBaseDomain}!for a vector object}
\label{OneOfBaseDomain:for a vector object}
}\hfill{\scriptsize (attribute)}}\\
\noindent\textcolor{FuncColor}{$\triangleright$\enspace\texttt{OneOfBaseDomain({\mdseries\slshape M})\index{OneOfBaseDomain@\texttt{OneOfBaseDomain}!for a matrix object}
\label{OneOfBaseDomain:for a matrix object}
}\hfill{\scriptsize (attribute)}}\\
\noindent\textcolor{FuncColor}{$\triangleright$\enspace\texttt{ZeroOfBaseDomain({\mdseries\slshape v})\index{ZeroOfBaseDomain@\texttt{ZeroOfBaseDomain}!for a vector object}
\label{ZeroOfBaseDomain:for a vector object}
}\hfill{\scriptsize (attribute)}}\\
\noindent\textcolor{FuncColor}{$\triangleright$\enspace\texttt{ZeroOfBaseDomain({\mdseries\slshape M})\index{ZeroOfBaseDomain@\texttt{ZeroOfBaseDomain}!for a matrix object}
\label{ZeroOfBaseDomain:for a matrix object}
}\hfill{\scriptsize (attribute)}}\\


 These attributes return the identity element and the zero element of the \texttt{BaseDomain} (\ref{BaseDomain:for a vector object}) value of the vector object \mbox{\texttt{\mdseries\slshape v}} or the matrix object \mbox{\texttt{\mdseries\slshape M}}, respectively. 

 If \mbox{\texttt{\mdseries\slshape v}} or \mbox{\texttt{\mdseries\slshape M}}, respectively, is a plain list (see \texttt{IsPlistRep} (\ref{IsPlistRep})) then computing its \texttt{BaseDomain} (\ref{BaseDomain:for a vector object}) value can be regarded as expensive, whereas calling \texttt{OneOfBaseDomain} or \texttt{ZeroOfBaseDomain} can be regarded as cheap. If \mbox{\texttt{\mdseries\slshape v}} or \mbox{\texttt{\mdseries\slshape M}}, respectively, is not a plain list then one can also call \texttt{BaseDomain} (\ref{BaseDomain:for a vector object}) first, without loss of performance. }

 }

  
\section{\textcolor{Chapter }{Operations for Vector and Matrix Objects}}\label{Operations for Vector and Matrix Objects}
\logpage{[ 26, 6, 0 ]}
\hyperdef{L}{X7954E20987E0B260}{}
{
  
\subsection{\textcolor{Chapter }{Comparison of Vector and Matrix Objects}}\logpage{[ 26, 6, 1 ]}
\hyperdef{L}{X7FFC60A27FE6FA97}{}
{
\noindent\textcolor{FuncColor}{$\triangleright$\enspace\texttt{\texttt{\symbol{92}}=({\mdseries\slshape v1, v2})\index{\texttt{\symbol{92}}=@\texttt{\texttt{\symbol{92}}=}!for two vector objects}
\label{bSlash=:for two vector objects}
}\hfill{\scriptsize (operation)}}\\
\noindent\textcolor{FuncColor}{$\triangleright$\enspace\texttt{\texttt{\symbol{92}}=({\mdseries\slshape M1, M2})\index{\texttt{\symbol{92}}=@\texttt{\texttt{\symbol{92}}=}!for two matrix objects}
\label{bSlash=:for two matrix objects}
}\hfill{\scriptsize (operation)}}\\
\noindent\textcolor{FuncColor}{$\triangleright$\enspace\texttt{\texttt{\symbol{92}}{\textless}({\mdseries\slshape v1, v2})\index{\texttt{\symbol{92}}{\textless}@\texttt{\texttt{\symbol{92}}{\textless}}!for two vector objects}
\label{bSlash<:for two vector objects}
}\hfill{\scriptsize (operation)}}\\
\noindent\textcolor{FuncColor}{$\triangleright$\enspace\texttt{\texttt{\symbol{92}}{\textless}({\mdseries\slshape M1, M2})\index{\texttt{\symbol{92}}{\textless}@\texttt{\texttt{\symbol{92}}{\textless}}!for two matrix objects}
\label{bSlash<:for two matrix objects}
}\hfill{\scriptsize (operation)}}\\


 Two vector objects in \texttt{IsList} (\ref{IsList}) are equal if they are equal as lists. Two matrix objects in \texttt{IsList} (\ref{IsList}) are equal if they are equal as lists. 

 Two vector objects of which at least one is not in \texttt{IsList} (\ref{IsList}) are equal with respect to \texttt{\texttt{\symbol{92}}=} (\ref{bSlash=}) if they have the same \texttt{ConstructingFilter} (\ref{ConstructingFilter:for a vector object}) value, the same \texttt{BaseDomain} (\ref{BaseDomain:for a vector object}) value, the same length, and the same entries. 

 Two matrix objects of which at least one is not in \texttt{IsList} (\ref{IsList}) are equal with respect to \texttt{\texttt{\symbol{92}}=} (\ref{bSlash=}) if they have the same \texttt{ConstructingFilter} (\ref{ConstructingFilter:for a matrix object}) value, the same \texttt{BaseDomain} (\ref{BaseDomain:for a matrix object}) value, the same dimensions, and the same entries. 

 We do \emph{not} state a general rule how vector and matrix objects shall behave w.r.t. the
comparison by \texttt{\texttt{\symbol{92}}{\textless}} (\ref{bSlash<}). Note that a ``row lexicographic order'' would be quite unnatural for matrices that are internally represented via a
list of columns. 

 Note that the operations \texttt{\texttt{\symbol{92}}=} (\ref{bSlash=}) and \texttt{\texttt{\symbol{92}}{\textless}} (\ref{bSlash<}) are used to form sorted lists and sets of objects, see for example \texttt{Sort} (\ref{Sort}) and \texttt{Set} (\ref{Set}). }

 
\subsection{\textcolor{Chapter }{Unpack}}\logpage{[ 26, 6, 2 ]}
\hyperdef{L}{X7FBBE79478012648}{}
{
\noindent\textcolor{FuncColor}{$\triangleright$\enspace\texttt{Unpack({\mdseries\slshape v})\index{Unpack@\texttt{Unpack}!for a vector object}
\label{Unpack:for a vector object}
}\hfill{\scriptsize (operation)}}\\
\noindent\textcolor{FuncColor}{$\triangleright$\enspace\texttt{Unpack({\mdseries\slshape M})\index{Unpack@\texttt{Unpack}!for a matrix object}
\label{Unpack:for a matrix object}
}\hfill{\scriptsize (operation)}}\\
\textbf{\indent Returns:\ }
A plain list



 Returns a new mutable plain list (see \texttt{IsPlistRep} (\ref{IsPlistRep})) containing the entries of the vector object \mbox{\texttt{\mdseries\slshape v}} or the matrix object \mbox{\texttt{\mdseries\slshape M}}, respectively. In the case of a matrix object, the result is a plain list of
plain lists. 

 Changing the result does not change \mbox{\texttt{\mdseries\slshape v}} or \mbox{\texttt{\mdseries\slshape M}}, respectively. The entries themselves are not copied. }

 
\subsection{\textcolor{Chapter }{ChangedBaseDomain}}\logpage{[ 26, 6, 3 ]}
\hyperdef{L}{X85E896F67CE2F925}{}
{
\noindent\textcolor{FuncColor}{$\triangleright$\enspace\texttt{ChangedBaseDomain({\mdseries\slshape v, R})\index{ChangedBaseDomain@\texttt{ChangedBaseDomain}!for a vector object}
\label{ChangedBaseDomain:for a vector object}
}\hfill{\scriptsize (operation)}}\\
\noindent\textcolor{FuncColor}{$\triangleright$\enspace\texttt{ChangedBaseDomain({\mdseries\slshape M, R})\index{ChangedBaseDomain@\texttt{ChangedBaseDomain}!for a matrix object}
\label{ChangedBaseDomain:for a matrix object}
}\hfill{\scriptsize (operation)}}\\


 For a vector object \mbox{\texttt{\mdseries\slshape v}} (a matrix object \mbox{\texttt{\mdseries\slshape M}}) and a semiring \mbox{\texttt{\mdseries\slshape R}}, \texttt{ChangedBaseDomain} returns a new vector object (matrix object) with \texttt{BaseDomain} (\ref{BaseDomain:for a vector object}) value \mbox{\texttt{\mdseries\slshape R}}, \texttt{ConstructingFilter} (\ref{ConstructingFilter:for a vector object}) value equal to that of \mbox{\texttt{\mdseries\slshape v}} (\mbox{\texttt{\mdseries\slshape M}}), and the same entries as \mbox{\texttt{\mdseries\slshape v}} (\mbox{\texttt{\mdseries\slshape M}}). 

 The result is mutable if and only if \mbox{\texttt{\mdseries\slshape v}} (\mbox{\texttt{\mdseries\slshape M}}) is mutable. 

 For example, one can create a vector defined over \texttt{GF(4)} from a vector defined over \texttt{GF(2)} with this operation. }

 
\subsection{\textcolor{Chapter }{Randomize}}\logpage{[ 26, 6, 4 ]}
\hyperdef{L}{X83DD8B39864A2C94}{}
{
\noindent\textcolor{FuncColor}{$\triangleright$\enspace\texttt{Randomize({\mdseries\slshape [Rs, ]v})\index{Randomize@\texttt{Randomize}!for a vector object}
\label{Randomize:for a vector object}
}\hfill{\scriptsize (operation)}}\\
\noindent\textcolor{FuncColor}{$\triangleright$\enspace\texttt{Randomize({\mdseries\slshape [Rs, ]M})\index{Randomize@\texttt{Randomize}!for a matrix object}
\label{Randomize:for a matrix object}
}\hfill{\scriptsize (operation)}}\\


 Replaces every entry in the mutable vector object \mbox{\texttt{\mdseries\slshape v}} or matrix object \mbox{\texttt{\mdseries\slshape M}}, respectively, with a random one from the base domain of \mbox{\texttt{\mdseries\slshape v}} or \mbox{\texttt{\mdseries\slshape M}}, respectively, and returns the argument. 

 If given, the random source \mbox{\texttt{\mdseries\slshape Rs}} is used to compute the random elements. Note that in this case, a \texttt{Random} (\ref{Random:for random source and collection}) method must be available that takes a random source as its first argument and
the base domain as its second argument. }

 }

  
\section{\textcolor{Chapter }{List Like Operations for Vector Objects}}\label{List Like Operations for Vector Objects}
\logpage{[ 26, 7, 0 ]}
\hyperdef{L}{X7FE662477F36A21F}{}
{
  The following operations that are defined for lists are useful also for vector
objects. (More such operations can be added if this is appropriate.) 
\subsection{\textcolor{Chapter }{Element Access and Assignment for Vector Objects}}\logpage{[ 26, 7, 1 ]}
\hyperdef{L}{X7D5DF49C7ADB6986}{}
{
\noindent\textcolor{FuncColor}{$\triangleright$\enspace\texttt{\texttt{\symbol{92}}[\texttt{\symbol{92}}]({\mdseries\slshape v, i})\index{\texttt{\symbol{92}}[\texttt{\symbol{92}}]@\texttt{\texttt{\symbol{92}}[\texttt{\symbol{92}}]}!for a vector object and an integer}
\label{bSlash[bSlash]:for a vector object and an integer}
}\hfill{\scriptsize (operation)}}\\
\noindent\textcolor{FuncColor}{$\triangleright$\enspace\texttt{\texttt{\symbol{92}}[\texttt{\symbol{92}}]\texttt{\symbol{92}}:\texttt{\symbol{92}}=({\mdseries\slshape v, i, obj})\index{\texttt{\symbol{92}}[\texttt{\symbol{92}}]\texttt{\symbol{92}}:\texttt{\symbol{92}}=@\texttt{\texttt{\symbol{92}}[\texttt{\symbol{92}}]\texttt{\symbol{92}}:\texttt{\symbol{92}}=}!for a vector object and an integer}
\label{bSlash[bSlash]bSlash:bSlash=:for a vector object and an integer}
}\hfill{\scriptsize (operation)}}\\
\noindent\textcolor{FuncColor}{$\triangleright$\enspace\texttt{\texttt{\symbol{92}}\texttt{\symbol{123}}\texttt{\symbol{92}}\texttt{\symbol{125}}({\mdseries\slshape v, list})\index{\texttt{\symbol{92}}\texttt{\symbol{123}}\texttt{\symbol{92}}\texttt{\symbol{125}}@\texttt{\texttt{\symbol{92}}\texttt{\symbol{123}}\texttt{\symbol{92}}\texttt{\symbol{125}}}!for a vector object and a list}
\label{bSlash{bSlash}:for a vector object and a list}
}\hfill{\scriptsize (operation)}}\\


 For a vector object \mbox{\texttt{\mdseries\slshape v}} and a positive integer \mbox{\texttt{\mdseries\slshape i}} that is not larger than the length of \mbox{\texttt{\mdseries\slshape v}} (see \texttt{Length} (\ref{Length:for a vector object})), \mbox{\texttt{\mdseries\slshape v}}\texttt{[}\mbox{\texttt{\mdseries\slshape i}}\texttt{]} is the entry at position \mbox{\texttt{\mdseries\slshape i}}. 

 If \mbox{\texttt{\mdseries\slshape v}} is mutable, \mbox{\texttt{\mdseries\slshape i}} is as above, and \mbox{\texttt{\mdseries\slshape obj}} is an object from the base domain of \mbox{\texttt{\mdseries\slshape v}} then \mbox{\texttt{\mdseries\slshape v}}\texttt{[}\mbox{\texttt{\mdseries\slshape i}}\texttt{]:= }\mbox{\texttt{\mdseries\slshape obj}} assigns \mbox{\texttt{\mdseries\slshape obj}} to the \mbox{\texttt{\mdseries\slshape i}}\texttt{\symbol{45}}th position of \mbox{\texttt{\mdseries\slshape v}}. 

 If \mbox{\texttt{\mdseries\slshape list}} is a list of positive integers that are not larger than the length of \mbox{\texttt{\mdseries\slshape v}} then \mbox{\texttt{\mdseries\slshape v}}\texttt{\texttt{\symbol{123}}}\mbox{\texttt{\mdseries\slshape list}}\texttt{\texttt{\symbol{125}}} returns a new mutable vector object in the same representation as \mbox{\texttt{\mdseries\slshape v}} (see \texttt{ConstructingFilter} (\ref{ConstructingFilter:for a vector object})) that contains the \mbox{\texttt{\mdseries\slshape list}}$[ k ]$\texttt{\symbol{45}}th entry of \mbox{\texttt{\mdseries\slshape v}} at position $k$. 

 If the global option \texttt{check} is set to \texttt{false} then \texttt{\texttt{\symbol{92}}[\texttt{\symbol{92}}]\texttt{\symbol{92}}:\texttt{\symbol{92}}=} need not perform consistency checks. 

 Note that the sublist assignment operation \texttt{\texttt{\symbol{92}}\texttt{\symbol{123}}\texttt{\symbol{92}}\texttt{\symbol{125}}\texttt{\symbol{92}}:\texttt{\symbol{92}}=} (\ref{bSlash{bSlash}bSlash:bSlash=}) is left out here since it tempts the programmer to use constructions like \texttt{v\texttt{\symbol{123}} [ 1 .. 3 ] \texttt{\symbol{125}}:=
w\texttt{\symbol{123}} [ 4 .. 6 ] \texttt{\symbol{125}}} which produces an unnecessary intermediate object; one should use \texttt{CopySubVector} (\ref{CopySubVector}) instead. }

 

\subsection{\textcolor{Chapter }{PositionNonZero (for a vector object)}}
\logpage{[ 26, 7, 2 ]}\nobreak
\hyperdef{L}{X7A21731C83EE3BB0}{}
{\noindent\textcolor{FuncColor}{$\triangleright$\enspace\texttt{PositionNonZero({\mdseries\slshape v})\index{PositionNonZero@\texttt{PositionNonZero}!for a vector object}
\label{PositionNonZero:for a vector object}
}\hfill{\scriptsize (operation)}}\\
\textbf{\indent Returns:\ }
An integer



 Returns the index of the first entry in the vector object \mbox{\texttt{\mdseries\slshape v}} that is not zero. If all entries are zero, the function returns \texttt{Length(\mbox{\texttt{\mdseries\slshape v}}) + 1}. }

 

\subsection{\textcolor{Chapter }{PositionLastNonZero (for a vector object)}}
\logpage{[ 26, 7, 3 ]}\nobreak
\hyperdef{L}{X7ABDE1B685A78326}{}
{\noindent\textcolor{FuncColor}{$\triangleright$\enspace\texttt{PositionLastNonZero({\mdseries\slshape v})\index{PositionLastNonZero@\texttt{PositionLastNonZero}!for a vector object}
\label{PositionLastNonZero:for a vector object}
}\hfill{\scriptsize (operation)}}\\
\textbf{\indent Returns:\ }
An integer



 Returns the index of the last entry in the vector object \mbox{\texttt{\mdseries\slshape v}} that is not zero. If all entries are zero, the function returns $0$. }

 

\subsection{\textcolor{Chapter }{ListOp (for vector object and function)}}
\logpage{[ 26, 7, 4 ]}\nobreak
\hyperdef{L}{X790013817E314B2D}{}
{\noindent\textcolor{FuncColor}{$\triangleright$\enspace\texttt{ListOp({\mdseries\slshape v[, func]})\index{ListOp@\texttt{ListOp}!for vector object and function}
\label{ListOp:for vector object and function}
}\hfill{\scriptsize (operation)}}\\
\textbf{\indent Returns:\ }
A plain list



 Applies the function \mbox{\texttt{\mdseries\slshape func}} to each entry of the vector object \mbox{\texttt{\mdseries\slshape v}} and returns the results as a mutable plain list. This allows for calling \texttt{List} (\ref{List:for a collection}) on vector objects. 

 If the argument \mbox{\texttt{\mdseries\slshape func}} is not given, applies \texttt{IdFunc} (\ref{IdFunc}) to all entries. }

 }

  
\section{\textcolor{Chapter }{Arithmetical Operations for Vector Objects}}\label{Arithmetical Operations for Vector Objects}
\logpage{[ 26, 8, 0 ]}
\hyperdef{L}{X7FDF7655852AEAAE}{}
{
  
\subsection{\textcolor{Chapter }{Unary Arithmetical Operations for Vector Objects}}\logpage{[ 26, 8, 1 ]}
\hyperdef{L}{X7F8CE23F7A250072}{}
{
\noindent\textcolor{FuncColor}{$\triangleright$\enspace\texttt{AdditiveInverseMutable({\mdseries\slshape v})\index{AdditiveInverseMutable@\texttt{AdditiveInverseMutable}!for vector object}
\label{AdditiveInverseMutable:for vector object}
}\hfill{\scriptsize (operation)}}\\
\noindent\textcolor{FuncColor}{$\triangleright$\enspace\texttt{AdditiveInverseSameMutability({\mdseries\slshape v})\index{AdditiveInverseSameMutability@\texttt{AdditiveInverseSameMutability}!for vector object}
\label{AdditiveInverseSameMutability:for vector object}
}\hfill{\scriptsize (operation)}}\\
\noindent\textcolor{FuncColor}{$\triangleright$\enspace\texttt{ZeroMutable({\mdseries\slshape v})\index{ZeroMutable@\texttt{ZeroMutable}!for vector object}
\label{ZeroMutable:for vector object}
}\hfill{\scriptsize (operation)}}\\
\noindent\textcolor{FuncColor}{$\triangleright$\enspace\texttt{ZeroSameMutability({\mdseries\slshape v})\index{ZeroSameMutability@\texttt{ZeroSameMutability}!for vector object}
\label{ZeroSameMutability:for vector object}
}\hfill{\scriptsize (operation)}}\\
\noindent\textcolor{FuncColor}{$\triangleright$\enspace\texttt{IsZero({\mdseries\slshape v})\index{IsZero@\texttt{IsZero}!for vector object}
\label{IsZero:for vector object}
}\hfill{\scriptsize (property)}}\\
\noindent\textcolor{FuncColor}{$\triangleright$\enspace\texttt{Characteristic({\mdseries\slshape v})\index{Characteristic@\texttt{Characteristic}!for vector object}
\label{Characteristic:for vector object}
}\hfill{\scriptsize (attribute)}}\\
\textbf{\indent Returns:\ }
a vector object



 For a vector object \mbox{\texttt{\mdseries\slshape v}}, the operations for computing the additive inverse with prescribed mutability
return a vector object with the same \texttt{ConstructingFilter} (\ref{ConstructingFilter:for a vector object}) and \texttt{BaseDomain} (\ref{BaseDomain:for a vector object}) values, such that the sum with \mbox{\texttt{\mdseries\slshape v}} is a zero vector. It is not specified what happens if the base domain does not
admit the additive inverses of the entries. 

 Analogously, the operations for computing a zero vector with prescribed
mutability return a vector object compatible with \mbox{\texttt{\mdseries\slshape v}}. 

 \texttt{IsZero} returns \texttt{true} if all entries in \mbox{\texttt{\mdseries\slshape v}} are zero, and \texttt{false} otherwise. 

 \texttt{Characteristic} returns the corresponding value of the \texttt{BaseDomain} (\ref{BaseDomain:for a vector object}) value of \mbox{\texttt{\mdseries\slshape v}}. }

 
\subsection{\textcolor{Chapter }{Binary Arithmetical Operations for Vector Objects}}\logpage{[ 26, 8, 2 ]}
\hyperdef{L}{X85A815CA790094CC}{}
{
\noindent\textcolor{FuncColor}{$\triangleright$\enspace\texttt{\texttt{\symbol{92}}+({\mdseries\slshape v1, v2})\index{\texttt{\symbol{92}}+@\texttt{\texttt{\symbol{92}}+}!for two vector objects}
\label{bSlash+:for two vector objects}
}\hfill{\scriptsize (method)}}\\
\noindent\textcolor{FuncColor}{$\triangleright$\enspace\texttt{\texttt{\symbol{92}}\texttt{\symbol{45}}({\mdseries\slshape v1, v2})\index{\texttt{\symbol{92}}\texttt{\symbol{45}}@\texttt{\texttt{\symbol{92}}\texttt{\symbol{45}}}!for two vector objects}
\label{bSlash-:for two vector objects}
}\hfill{\scriptsize (method)}}\\
\noindent\textcolor{FuncColor}{$\triangleright$\enspace\texttt{\texttt{\symbol{92}}*({\mdseries\slshape s, v})\index{\texttt{\symbol{92}}*@\texttt{\texttt{\symbol{92}}*}!for scalar and vector object}
\label{bSlash*:for scalar and vector object}
}\hfill{\scriptsize (method)}}\\
\noindent\textcolor{FuncColor}{$\triangleright$\enspace\texttt{\texttt{\symbol{92}}*({\mdseries\slshape v, s})\index{\texttt{\symbol{92}}*@\texttt{\texttt{\symbol{92}}*}!for vector object and scalar}
\label{bSlash*:for vector object and scalar}
}\hfill{\scriptsize (method)}}\\
\noindent\textcolor{FuncColor}{$\triangleright$\enspace\texttt{\texttt{\symbol{92}}*({\mdseries\slshape v1, v2})\index{\texttt{\symbol{92}}*@\texttt{\texttt{\symbol{92}}*}!for two vector objects}
\label{bSlash*:for two vector objects}
}\hfill{\scriptsize (method)}}\\
\noindent\textcolor{FuncColor}{$\triangleright$\enspace\texttt{ScalarProduct({\mdseries\slshape v1, v2})\index{ScalarProduct@\texttt{ScalarProduct}!for two vector objects}
\label{ScalarProduct:for two vector objects}
}\hfill{\scriptsize (method)}}\\
\noindent\textcolor{FuncColor}{$\triangleright$\enspace\texttt{\texttt{\symbol{92}}/({\mdseries\slshape v, s})\index{\texttt{\symbol{92}}/@\texttt{\texttt{\symbol{92}}/}!for vector object and scalar}
\label{bSlash/:for vector object and scalar}
}\hfill{\scriptsize (method)}}\\


 The sum and the difference, respectively, of two vector objects \mbox{\texttt{\mdseries\slshape v1}} and \mbox{\texttt{\mdseries\slshape v2}} is a new mutable vector object whose entries are the sums and the differences
of the entries of the arguments. 

 The product of a scalar \mbox{\texttt{\mdseries\slshape s}} and a vector object \mbox{\texttt{\mdseries\slshape v}} (from the left or from the right) is a new mutable vector object whose entries
are the corresponding products. 

 The quotient of a vector object \mbox{\texttt{\mdseries\slshape v}} and a scalar \mbox{\texttt{\mdseries\slshape s}} is a new mutable vector object whose entries are the corresponding quotients. 

 The product of two vector objects \mbox{\texttt{\mdseries\slshape v1}} and \mbox{\texttt{\mdseries\slshape v2}} as well as the result of \texttt{ScalarProduct} is the standard scalar product of the two arguments (an element of the base
domain of the vector objects). 

 All this is defined only if the vector objects have the same length and are
defined over the same base domain and have the same representation, and if the
products with the given scalar belong to the base domain; otherwise it is not
specified what happens. If the result is a vector object then it has the same
representation and the same base domain as the given vector object(s). }

 

\subsection{\textcolor{Chapter }{AddVector (for two vector objects)}}
\logpage{[ 26, 8, 3 ]}\nobreak
\hyperdef{L}{X876090A684E71C93}{}
{\noindent\textcolor{FuncColor}{$\triangleright$\enspace\texttt{AddVector({\mdseries\slshape dst, src[, mul[, from, to]]})\index{AddVector@\texttt{AddVector}!for two vector objects}
\label{AddVector:for two vector objects}
}\hfill{\scriptsize (operation)}}\\
\noindent\textcolor{FuncColor}{$\triangleright$\enspace\texttt{AddVector({\mdseries\slshape dst, mul, src[, from, to]})\index{AddVector@\texttt{AddVector}!for two vector objects and a scalar}
\label{AddVector:for two vector objects and a scalar}
}\hfill{\scriptsize (operation)}}\\
\textbf{\indent Returns:\ }
nothing



 Called with two vector objects \mbox{\texttt{\mdseries\slshape dst}} and \mbox{\texttt{\mdseries\slshape src}}, this function replaces the entries of \mbox{\texttt{\mdseries\slshape dst}} in\texttt{\symbol{45}}place by the entries of the sum \mbox{\texttt{\mdseries\slshape dst}}\texttt{ + }\mbox{\texttt{\mdseries\slshape src}}. 

 If a scalar \mbox{\texttt{\mdseries\slshape mul}} is given as the third or second argument, respectively, then the entries of \mbox{\texttt{\mdseries\slshape dst}} get replaced by those of \mbox{\texttt{\mdseries\slshape dst}}\texttt{ + }\mbox{\texttt{\mdseries\slshape src}}\texttt{ * }\mbox{\texttt{\mdseries\slshape mul}} or \mbox{\texttt{\mdseries\slshape dst}}\texttt{ + }\mbox{\texttt{\mdseries\slshape mul}}\texttt{ * }\mbox{\texttt{\mdseries\slshape src}}, respectively. 

 If the optional parameters \mbox{\texttt{\mdseries\slshape from}} and \mbox{\texttt{\mdseries\slshape to}} are given then only the index range \texttt{[\mbox{\texttt{\mdseries\slshape from}}..\mbox{\texttt{\mdseries\slshape to}}]} is guaranteed to be affected. Other indices \emph{may} be affected, if it is more convenient to do so. This can be helpful if entries
of \mbox{\texttt{\mdseries\slshape src}} are known to be zero. 

 If \mbox{\texttt{\mdseries\slshape from}} is bigger than \mbox{\texttt{\mdseries\slshape to}}, the operation does nothing. }

 

\subsection{\textcolor{Chapter }{MultVector (for a vector object)}}
\logpage{[ 26, 8, 4 ]}\nobreak
\hyperdef{L}{X8039D013817317C3}{}
{\noindent\textcolor{FuncColor}{$\triangleright$\enspace\texttt{MultVector({\mdseries\slshape v, mul[, from, to]})\index{MultVector@\texttt{MultVector}!for a vector object}
\label{MultVector:for a vector object}
}\hfill{\scriptsize (operation)}}\\
\noindent\textcolor{FuncColor}{$\triangleright$\enspace\texttt{MultVectorLeft({\mdseries\slshape v, mul[, from, to]})\index{MultVectorLeft@\texttt{MultVectorLeft}!for a vector object}
\label{MultVectorLeft:for a vector object}
}\hfill{\scriptsize (operation)}}\\
\noindent\textcolor{FuncColor}{$\triangleright$\enspace\texttt{MultVectorRight({\mdseries\slshape v, mul[, from, to]})\index{MultVectorRight@\texttt{MultVectorRight}!for a vector object}
\label{MultVectorRight:for a vector object}
}\hfill{\scriptsize (operation)}}\\
\textbf{\indent Returns:\ }
nothing



 These operations multiply \mbox{\texttt{\mdseries\slshape v}} by \mbox{\texttt{\mdseries\slshape mul}} in\texttt{\symbol{45}}place where \texttt{MultVectorLeft} multiplies with \mbox{\texttt{\mdseries\slshape mul}} from the left and \texttt{MultVectorRight} does so from the right. 

 Note that \texttt{MultVector} is just a synonym for \texttt{MultVectorLeft}. This was chosen because vectors in \textsf{GAP} are by default row vectors and scalar multiplication is usually written as $a \cdot v = a \cdot [v_1, ..., v_n] = [a \cdot v_1, ..., a \cdot v_n]$ with scalars being applied from the left. 

 If the optional parameters \mbox{\texttt{\mdseries\slshape from}} and \mbox{\texttt{\mdseries\slshape to}} are given then only the index range \texttt{[\mbox{\texttt{\mdseries\slshape from}}..\mbox{\texttt{\mdseries\slshape to}}]} is guaranteed to be affected. Other indices \emph{may} be affected, if it is more convenient to do so. This can be helpful if entries
of \mbox{\texttt{\mdseries\slshape v}} are known to be zero. If \mbox{\texttt{\mdseries\slshape from}} is bigger than \mbox{\texttt{\mdseries\slshape to}}, the operation does nothing. 

 }

 }

  
\section{\textcolor{Chapter }{Operations for Vector Objects}}\label{Operations for Vector Objects}
\logpage{[ 26, 9, 0 ]}
\hyperdef{L}{X7BE9D278852C13BC}{}
{
  
\subsection{\textcolor{Chapter }{ConcatenationOfVectors}}\logpage{[ 26, 9, 1 ]}
\hyperdef{L}{X7AC470557EC90714}{}
{
\noindent\textcolor{FuncColor}{$\triangleright$\enspace\texttt{ConcatenationOfVectors({\mdseries\slshape v1, v2, ...})\index{ConcatenationOfVectors@\texttt{ConcatenationOfVectors}!for arbitrary many vector objects}
\label{ConcatenationOfVectors:for arbitrary many vector objects}
}\hfill{\scriptsize (function)}}\\
\noindent\textcolor{FuncColor}{$\triangleright$\enspace\texttt{ConcatenationOfVectors({\mdseries\slshape vlist})\index{ConcatenationOfVectors@\texttt{ConcatenationOfVectors}!for a list of vector objects}
\label{ConcatenationOfVectors:for a list of vector objects}
}\hfill{\scriptsize (function)}}\\
\textbf{\indent Returns:\ }
a vector object



 Returns a new mutable vector object in the representation of \mbox{\texttt{\mdseries\slshape v1}} or the first entry of the nonempty list \mbox{\texttt{\mdseries\slshape vlist}} of vector objects, respectively, such that the entries are the concatenation
of the given vector objects. 

 (Note that \texttt{Concatenation} (\ref{Concatenation:for several lists}) is a function for which no methods can be installed.) }

 

\subsection{\textcolor{Chapter }{ExtractSubVector}}
\logpage{[ 26, 9, 2 ]}\nobreak
\hyperdef{L}{X7DBE956E7F9C700E}{}
{\noindent\textcolor{FuncColor}{$\triangleright$\enspace\texttt{ExtractSubVector({\mdseries\slshape v, l})\index{ExtractSubVector@\texttt{ExtractSubVector}}
\label{ExtractSubVector}
}\hfill{\scriptsize (operation)}}\\
\textbf{\indent Returns:\ }
a vector object



 Returns a new mutable vector object of the same vector representation as \mbox{\texttt{\mdseries\slshape v}}, containing the entries of \mbox{\texttt{\mdseries\slshape v}} at the positions in the list \mbox{\texttt{\mdseries\slshape l}}. 

 This is the same as \mbox{\texttt{\mdseries\slshape v}}\texttt{\texttt{\symbol{123}}}\mbox{\texttt{\mdseries\slshape l}}\texttt{\texttt{\symbol{125}}}, the name \texttt{ExtractSubVector} was introduced in analogy to \texttt{ExtractSubMatrix} (\ref{ExtractSubMatrix}), for which no equivalent syntax using curly brackets is available. }

 

\subsection{\textcolor{Chapter }{CopySubVector}}
\logpage{[ 26, 9, 3 ]}\nobreak
\hyperdef{L}{X80EC354D78D7B5A6}{}
{\noindent\textcolor{FuncColor}{$\triangleright$\enspace\texttt{CopySubVector({\mdseries\slshape src, dst, scols, dcols})\index{CopySubVector@\texttt{CopySubVector}}
\label{CopySubVector}
}\hfill{\scriptsize (operation)}}\\
\textbf{\indent Returns:\ }
nothing



 For two row vectors or vector objects \mbox{\texttt{\mdseries\slshape src}} and \mbox{\texttt{\mdseries\slshape dst}}, such that \mbox{\texttt{\mdseries\slshape dst}} is mutable, and two lists \mbox{\texttt{\mdseries\slshape scols}} and \mbox{\texttt{\mdseries\slshape dcols}} of positions, \texttt{CopySubVector} assigns the entries \mbox{\texttt{\mdseries\slshape src}}\texttt{\texttt{\symbol{123}} }\mbox{\texttt{\mdseries\slshape scols}}\texttt{ \texttt{\symbol{125}}} (see \texttt{ExtractSubVector} (\ref{ExtractSubVector})) to the positions \mbox{\texttt{\mdseries\slshape dcols}} in \mbox{\texttt{\mdseries\slshape dst}}, but without creating an intermediate object and thus {\textendash}at least
in special cases{\textendash} much more efficiently. 

 For certain objects like compressed vectors this might be significantly more
efficient if \mbox{\texttt{\mdseries\slshape scols}} and \mbox{\texttt{\mdseries\slshape dcols}} are ranges with increment 1. 

 If the global option \texttt{check} is set to \texttt{false} then \texttt{CopySubVector} need not perform consistency checks. }

 

\subsection{\textcolor{Chapter }{WeightOfVector (for a vector object)}}
\logpage{[ 26, 9, 4 ]}\nobreak
\hyperdef{L}{X866366E587991171}{}
{\noindent\textcolor{FuncColor}{$\triangleright$\enspace\texttt{WeightOfVector({\mdseries\slshape v})\index{WeightOfVector@\texttt{WeightOfVector}!for a vector object}
\label{WeightOfVector:for a vector object}
}\hfill{\scriptsize (operation)}}\\
\textbf{\indent Returns:\ }
an integer



 returns the Hamming weight of the vector object \mbox{\texttt{\mdseries\slshape v}}, i.e., the number of nonzero entries in \mbox{\texttt{\mdseries\slshape v}}. }

 

\subsection{\textcolor{Chapter }{DistanceOfVectors (for two vector objects)}}
\logpage{[ 26, 9, 5 ]}\nobreak
\hyperdef{L}{X81ACAE017C00F782}{}
{\noindent\textcolor{FuncColor}{$\triangleright$\enspace\texttt{DistanceOfVectors({\mdseries\slshape v1, v2})\index{DistanceOfVectors@\texttt{DistanceOfVectors}!for two vector objects}
\label{DistanceOfVectors:for two vector objects}
}\hfill{\scriptsize (operation)}}\\
\textbf{\indent Returns:\ }
an integer



 returns the Hamming distance of the vector objects \mbox{\texttt{\mdseries\slshape v1}} and \mbox{\texttt{\mdseries\slshape v2}}, i.e., the number of entries in which the vectors differ. The vectors must
have equal length. }

 }

  
\section{\textcolor{Chapter }{Arithmetical Operations for Matrix Objects}}\label{Arithmetical Operations for Matrix Objects}
\logpage{[ 26, 10, 0 ]}
\hyperdef{L}{X81CC13CA7A1FF4AA}{}
{
  
\subsection{\textcolor{Chapter }{Unary Arithmetical Operations for Matrix Objects}}\logpage{[ 26, 10, 1 ]}
\hyperdef{L}{X819F87A07DA7E2DC}{}
{
\noindent\textcolor{FuncColor}{$\triangleright$\enspace\texttt{AdditiveInverseMutable({\mdseries\slshape M})\index{AdditiveInverseMutable@\texttt{AdditiveInverseMutable}!for matrix object}
\label{AdditiveInverseMutable:for matrix object}
}\hfill{\scriptsize (operation)}}\\
\noindent\textcolor{FuncColor}{$\triangleright$\enspace\texttt{AdditiveInverseSameMutability({\mdseries\slshape M})\index{AdditiveInverseSameMutability@\texttt{AdditiveInverseSameMutability}!for matrix object}
\label{AdditiveInverseSameMutability:for matrix object}
}\hfill{\scriptsize (operation)}}\\
\noindent\textcolor{FuncColor}{$\triangleright$\enspace\texttt{ZeroMutable({\mdseries\slshape M})\index{ZeroMutable@\texttt{ZeroMutable}!for matrix object}
\label{ZeroMutable:for matrix object}
}\hfill{\scriptsize (operation)}}\\
\noindent\textcolor{FuncColor}{$\triangleright$\enspace\texttt{ZeroSameMutability({\mdseries\slshape M})\index{ZeroSameMutability@\texttt{ZeroSameMutability}!for matrix object}
\label{ZeroSameMutability:for matrix object}
}\hfill{\scriptsize (operation)}}\\
\noindent\textcolor{FuncColor}{$\triangleright$\enspace\texttt{OneMutable({\mdseries\slshape M})\index{OneMutable@\texttt{OneMutable}!for matrix object}
\label{OneMutable:for matrix object}
}\hfill{\scriptsize (operation)}}\\
\noindent\textcolor{FuncColor}{$\triangleright$\enspace\texttt{OneSameMutability({\mdseries\slshape M})\index{OneSameMutability@\texttt{OneSameMutability}!for matrix object}
\label{OneSameMutability:for matrix object}
}\hfill{\scriptsize (operation)}}\\
\noindent\textcolor{FuncColor}{$\triangleright$\enspace\texttt{InverseMutable({\mdseries\slshape M})\index{InverseMutable@\texttt{InverseMutable}!for matrix object}
\label{InverseMutable:for matrix object}
}\hfill{\scriptsize (operation)}}\\
\noindent\textcolor{FuncColor}{$\triangleright$\enspace\texttt{InverseSameMutability({\mdseries\slshape M})\index{InverseSameMutability@\texttt{InverseSameMutability}!for matrix object}
\label{InverseSameMutability:for matrix object}
}\hfill{\scriptsize (operation)}}\\
\noindent\textcolor{FuncColor}{$\triangleright$\enspace\texttt{IsZero({\mdseries\slshape M})\index{IsZero@\texttt{IsZero}!for matrix object}
\label{IsZero:for matrix object}
}\hfill{\scriptsize (property)}}\\
\noindent\textcolor{FuncColor}{$\triangleright$\enspace\texttt{IsOne({\mdseries\slshape M})\index{IsOne@\texttt{IsOne}!for matrix object}
\label{IsOne:for matrix object}
}\hfill{\scriptsize (property)}}\\
\noindent\textcolor{FuncColor}{$\triangleright$\enspace\texttt{Characteristic({\mdseries\slshape M})\index{Characteristic@\texttt{Characteristic}!for matrix object}
\label{Characteristic:for matrix object}
}\hfill{\scriptsize (attribute)}}\\
\textbf{\indent Returns:\ }
a matrix object



 For a vector object \mbox{\texttt{\mdseries\slshape M}}, the operations for computing the additive inverse with prescribed mutability
return a matrix object with the same \texttt{ConstructingFilter} (\ref{ConstructingFilter:for a matrix object}) and \texttt{BaseDomain} (\ref{BaseDomain:for a matrix object}) values, such that the sum with \mbox{\texttt{\mdseries\slshape M}} is a zero matrix. It is not specified what happens if the base domain does not
admit the additive inverses of the entries. 

 Analogously, the operations for computing a zero matrix with prescribed
mutability return a matrix object compatible with \mbox{\texttt{\mdseries\slshape M}}. 

 The operations for computing an identity matrix with prescribed mutability
return a matrix object compatible with \mbox{\texttt{\mdseries\slshape M}}, provided that the base domain admits this and \mbox{\texttt{\mdseries\slshape M}} is square and nonempty. 

 Analogously, the operations for computing an inverse matrix with prescribed
mutability return a matrix object compatible with \mbox{\texttt{\mdseries\slshape M}}, provided that \mbox{\texttt{\mdseries\slshape M}} is invertible.  (If \mbox{\texttt{\mdseries\slshape M}} is not invertible then the operations return \texttt{fail}.) 

 \texttt{IsZero} returns \texttt{true} if all entries in \mbox{\texttt{\mdseries\slshape M}} are zero, and \texttt{false} otherwise. \texttt{IsOne} returns \texttt{true} if \mbox{\texttt{\mdseries\slshape M}} is nonempty and square and contains the identity of the base domain in the
diagonal, and zero in all other places. 

 \texttt{Characteristic} returns the corresponding value of the \texttt{BaseDomain} (\ref{BaseDomain:for a matrix object}) value of \mbox{\texttt{\mdseries\slshape M}}. }

 
\subsection{\textcolor{Chapter }{Binary Arithmetical Operations for Matrix Objects}}\logpage{[ 26, 10, 2 ]}
\hyperdef{L}{X7BBB70557A7A9591}{}
{
\noindent\textcolor{FuncColor}{$\triangleright$\enspace\texttt{\texttt{\symbol{92}}+({\mdseries\slshape M1, M2})\index{\texttt{\symbol{92}}+@\texttt{\texttt{\symbol{92}}+}!for two matrix objects}
\label{bSlash+:for two matrix objects}
}\hfill{\scriptsize (method)}}\\
\noindent\textcolor{FuncColor}{$\triangleright$\enspace\texttt{\texttt{\symbol{92}}\texttt{\symbol{45}}({\mdseries\slshape M1, M2})\index{\texttt{\symbol{92}}\texttt{\symbol{45}}@\texttt{\texttt{\symbol{92}}\texttt{\symbol{45}}}!for two matrix objects}
\label{bSlash-:for two matrix objects}
}\hfill{\scriptsize (method)}}\\
\noindent\textcolor{FuncColor}{$\triangleright$\enspace\texttt{\texttt{\symbol{92}}*({\mdseries\slshape s, M})\index{\texttt{\symbol{92}}*@\texttt{\texttt{\symbol{92}}*}!for scalar and matrix object}
\label{bSlash*:for scalar and matrix object}
}\hfill{\scriptsize (method)}}\\
\noindent\textcolor{FuncColor}{$\triangleright$\enspace\texttt{\texttt{\symbol{92}}*({\mdseries\slshape M, s})\index{\texttt{\symbol{92}}*@\texttt{\texttt{\symbol{92}}*}!for Matrix object and scalar}
\label{bSlash*:for Matrix object and scalar}
}\hfill{\scriptsize (method)}}\\
\noindent\textcolor{FuncColor}{$\triangleright$\enspace\texttt{\texttt{\symbol{92}}*({\mdseries\slshape M1, M2})\index{\texttt{\symbol{92}}*@\texttt{\texttt{\symbol{92}}*}!for two matrix objects}
\label{bSlash*:for two matrix objects}
}\hfill{\scriptsize (method)}}\\
\noindent\textcolor{FuncColor}{$\triangleright$\enspace\texttt{\texttt{\symbol{92}}/({\mdseries\slshape M, s})\index{\texttt{\symbol{92}}/@\texttt{\texttt{\symbol{92}}/}!for matrix object and scalar}
\label{bSlash/:for matrix object and scalar}
}\hfill{\scriptsize (method)}}\\
\noindent\textcolor{FuncColor}{$\triangleright$\enspace\texttt{\texttt{\symbol{92}}\texttt{\symbol{94}}({\mdseries\slshape M, n})\index{\texttt{\symbol{92}}\texttt{\symbol{94}}@\texttt{\texttt{\symbol{92}}\texttt{\symbol{94}}}!for matrix object and integer}
\label{bSlash^:for matrix object and integer}
}\hfill{\scriptsize (method)}}\\


 The sum and the difference, respectively, of two matrix objects \mbox{\texttt{\mdseries\slshape M1}} and \mbox{\texttt{\mdseries\slshape M2}} is a new fully mutable matrix object whose entries are the sums and the
differences of the entries of the arguments. 

 The product of a scalar \mbox{\texttt{\mdseries\slshape s}} and a matrix object \mbox{\texttt{\mdseries\slshape M}} (from the left or from the right) is a new fully mutable matrix object whose
entries are the corresponding products. 

 The product of two matrix objects \mbox{\texttt{\mdseries\slshape M1}} and \mbox{\texttt{\mdseries\slshape M2}} is a new fully mutable matrix object; if both \mbox{\texttt{\mdseries\slshape M1}} and \mbox{\texttt{\mdseries\slshape M2}} are in the filter \texttt{IsOrdinaryMatrix} (\ref{IsOrdinaryMatrix}) then the entries of the result are those of the ordinary matrix product. 

 The quotient of a matrix object \mbox{\texttt{\mdseries\slshape M}} and a scalar \mbox{\texttt{\mdseries\slshape s}} is a new fully mutable matrix object whose entries are the corresponding
quotients. 

 For a nonempty square matrix object \mbox{\texttt{\mdseries\slshape M}} over an associative base domain, and a positive integer \mbox{\texttt{\mdseries\slshape n}}, \mbox{\texttt{\mdseries\slshape M}}\texttt{\texttt{\symbol{94}}}\mbox{\texttt{\mdseries\slshape n}} is a fully mutable matrix object whose entries are those of the \mbox{\texttt{\mdseries\slshape n}}\texttt{\symbol{45}}th power of \mbox{\texttt{\mdseries\slshape M}}. If \mbox{\texttt{\mdseries\slshape n}} is zero then \mbox{\texttt{\mdseries\slshape M}}\texttt{\texttt{\symbol{94}}}\mbox{\texttt{\mdseries\slshape n}} is an identity matrix, and if \mbox{\texttt{\mdseries\slshape n}} is a negative integer and \mbox{\texttt{\mdseries\slshape M}} is invertible then \mbox{\texttt{\mdseries\slshape M}}\texttt{\texttt{\symbol{94}}}\mbox{\texttt{\mdseries\slshape n}} is the (\texttt{\texttt{\symbol{45}}}\mbox{\texttt{\mdseries\slshape n}})\texttt{\symbol{45}}th power of the inverse of \mbox{\texttt{\mdseries\slshape M}}. 

 All this is defined only if the matrix objects have the same dimensions and
are defined over the same base domain and have the same representation, and if
the products with the given scalar belong to the base domain; otherwise it is
not specified what happens. If the result is a matrix object then it has the
same representation and the same base domain as the given matrix object(s). }

 }

  
\section{\textcolor{Chapter }{Operations for Matrix Objects}}\label{Operations for Matrix Objects}
\logpage{[ 26, 11, 0 ]}
\hyperdef{L}{X85FAB7E778A71C19}{}
{
  

\subsection{\textcolor{Chapter }{MatElm}}
\logpage{[ 26, 11, 1 ]}\nobreak
\hyperdef{L}{X870FBE817C884AB5}{}
{\noindent\textcolor{FuncColor}{$\triangleright$\enspace\texttt{MatElm({\mdseries\slshape M, row, col})\index{MatElm@\texttt{MatElm}}
\label{MatElm}
}\hfill{\scriptsize (operation)}}\\
\textbf{\indent Returns:\ }
an entry of the matrix object



 For a matrix object \mbox{\texttt{\mdseries\slshape M}}, this operation returns the entry in row \mbox{\texttt{\mdseries\slshape row}} and column \mbox{\texttt{\mdseries\slshape col}}. 

 Also the syntax \mbox{\texttt{\mdseries\slshape M}}\texttt{[ }\mbox{\texttt{\mdseries\slshape row}}\texttt{, }\mbox{\texttt{\mdseries\slshape col}}\texttt{ ]} is supported. 

 Note that this is \emph{not} equivalent to \mbox{\texttt{\mdseries\slshape M}}\texttt{[ }\mbox{\texttt{\mdseries\slshape row}}\texttt{ ][ }\mbox{\texttt{\mdseries\slshape col}}\texttt{ ]}, which would first try to access \mbox{\texttt{\mdseries\slshape M}}\texttt{[ }\mbox{\texttt{\mdseries\slshape row}}\texttt{ ]}, and this is in general not possible. }

 

\subsection{\textcolor{Chapter }{SetMatElm}}
\logpage{[ 26, 11, 2 ]}\nobreak
\hyperdef{L}{X7C33059984635480}{}
{\noindent\textcolor{FuncColor}{$\triangleright$\enspace\texttt{SetMatElm({\mdseries\slshape M, row, col, obj})\index{SetMatElm@\texttt{SetMatElm}}
\label{SetMatElm}
}\hfill{\scriptsize (operation)}}\\
\textbf{\indent Returns:\ }
nothing



 For a mutable matrix object \mbox{\texttt{\mdseries\slshape M}}, this operation assigns the object \mbox{\texttt{\mdseries\slshape obj}} to the position in row \mbox{\texttt{\mdseries\slshape row}} and column \mbox{\texttt{\mdseries\slshape col}}, provided that \mbox{\texttt{\mdseries\slshape obj}} is compatible with the \texttt{BaseDomain} (\ref{BaseDomain:for a matrix object}) value of \mbox{\texttt{\mdseries\slshape M}}. 

 Also the syntax \mbox{\texttt{\mdseries\slshape M}}\texttt{[ }\mbox{\texttt{\mdseries\slshape row}}\texttt{, }\mbox{\texttt{\mdseries\slshape col}}\texttt{ ]:= }\mbox{\texttt{\mdseries\slshape obj}} is supported. 

 Note that this is \emph{not} equivalent to \mbox{\texttt{\mdseries\slshape M}}\texttt{[ }\mbox{\texttt{\mdseries\slshape row}}\texttt{ ][ }\mbox{\texttt{\mdseries\slshape col}}\texttt{ ]:= }\mbox{\texttt{\mdseries\slshape obj}}, which would first try to access \mbox{\texttt{\mdseries\slshape M}}\texttt{[ }\mbox{\texttt{\mdseries\slshape row}}\texttt{ ]}, and this is in general not possible. 

 If the global option \texttt{check} is set to \texttt{false} then \texttt{SetMatElm} need not perform consistency checks. }

 

\subsection{\textcolor{Chapter }{ExtractSubMatrix}}
\logpage{[ 26, 11, 3 ]}\nobreak
\hyperdef{L}{X838B45F7790E9FDF}{}
{\noindent\textcolor{FuncColor}{$\triangleright$\enspace\texttt{ExtractSubMatrix({\mdseries\slshape M, rows, cols})\index{ExtractSubMatrix@\texttt{ExtractSubMatrix}}
\label{ExtractSubMatrix}
}\hfill{\scriptsize (operation)}}\\


 Creates a copy of the submatrix described by the two lists, which mean subsets
of row and column positions, respectively. This does \mbox{\texttt{\mdseries\slshape M}}\texttt{\symbol{123}}\mbox{\texttt{\mdseries\slshape rows}}\texttt{\symbol{125}}\texttt{\symbol{123}}\mbox{\texttt{\mdseries\slshape cols}}\texttt{\symbol{125}} and returns the result. It preserves the representation
of the matrix. 

 If the \texttt{ConstructingFilter} (\ref{ConstructingFilter:for a matrix object}) value of the result implies \texttt{IsCopyable} (\ref{IsCopyable}) then the result is fully mutable. }

 

\subsection{\textcolor{Chapter }{MutableCopyMatrix (for a matrix object)}}
\logpage{[ 26, 11, 4 ]}\nobreak
\hyperdef{L}{X793CD4637F237915}{}
{\noindent\textcolor{FuncColor}{$\triangleright$\enspace\texttt{MutableCopyMatrix({\mdseries\slshape M})\index{MutableCopyMatrix@\texttt{MutableCopyMatrix}!for a matrix object}
\label{MutableCopyMatrix:for a matrix object}
}\hfill{\scriptsize (operation)}}\\


 For a matrix object \mbox{\texttt{\mdseries\slshape M}}, this operation returns a fully mutable copy of \mbox{\texttt{\mdseries\slshape M}}, with the same \texttt{ConstructingFilter} (\ref{ConstructingFilter:for a matrix object}) and \texttt{BaseDomain} (\ref{BaseDomain:for a matrix object}) values, }

 

\subsection{\textcolor{Chapter }{CopySubMatrix}}
\logpage{[ 26, 11, 5 ]}\nobreak
\hyperdef{L}{X7ED9E5D4809E3B50}{}
{\noindent\textcolor{FuncColor}{$\triangleright$\enspace\texttt{CopySubMatrix({\mdseries\slshape src, dst, srows, drows, scols, dcols})\index{CopySubMatrix@\texttt{CopySubMatrix}}
\label{CopySubMatrix}
}\hfill{\scriptsize (operation)}}\\
\textbf{\indent Returns:\ }
nothing



 Does \texttt{\mbox{\texttt{\mdseries\slshape dst}}\texttt{\symbol{123}}\mbox{\texttt{\mdseries\slshape drows}}\texttt{\symbol{125}}\texttt{\symbol{123}}\mbox{\texttt{\mdseries\slshape dcols}}\texttt{\symbol{125}} := \mbox{\texttt{\mdseries\slshape src}}\texttt{\symbol{123}}\mbox{\texttt{\mdseries\slshape srows}}\texttt{\symbol{125}}\texttt{\symbol{123}}\mbox{\texttt{\mdseries\slshape scols}}\texttt{\symbol{125}}} without creating an intermediate object and thus {\textendash}at least in
special cases{\textendash} much more efficiently. For certain objects like
compressed vectors this might be significantly more efficient if \mbox{\texttt{\mdseries\slshape scols}} and \mbox{\texttt{\mdseries\slshape dcols}} are ranges with increment 1. 

 If the global option \texttt{check} is set to \texttt{false} then \texttt{CopySubMatrix} need not perform consistency checks. }

 

\subsection{\textcolor{Chapter }{CompatibleVector (for a matrix object)}}
\logpage{[ 26, 11, 6 ]}\nobreak
\hyperdef{L}{X809A6B3F7EA5E7D8}{}
{\noindent\textcolor{FuncColor}{$\triangleright$\enspace\texttt{CompatibleVector({\mdseries\slshape M})\index{CompatibleVector@\texttt{CompatibleVector}!for a matrix object}
\label{CompatibleVector:for a matrix object}
}\hfill{\scriptsize (operation)}}\\
\textbf{\indent Returns:\ }
a vector object



 Called with a matrix object \mbox{\texttt{\mdseries\slshape M}} with $m$ rows, this operation returns a mutable zero vector object $v$ of length $m$ and in the representation given by the \texttt{CompatibleVectorFilter} (\ref{CompatibleVectorFilter:for a matrix object}) value of \mbox{\texttt{\mdseries\slshape M}} (provided that such a representation exists). 

 The idea is that there should be an efficient way to form the product $v$\mbox{\texttt{\mdseries\slshape M}}. }

 

\subsection{\textcolor{Chapter }{RowsOfMatrix (for a matrix object)}}
\logpage{[ 26, 11, 7 ]}\nobreak
\hyperdef{L}{X7EE70D5A81E9ED72}{}
{\noindent\textcolor{FuncColor}{$\triangleright$\enspace\texttt{RowsOfMatrix({\mdseries\slshape M})\index{RowsOfMatrix@\texttt{RowsOfMatrix}!for a matrix object}
\label{RowsOfMatrix:for a matrix object}
}\hfill{\scriptsize (attribute)}}\\
\textbf{\indent Returns:\ }
a plain list



 Called with a matrix object \mbox{\texttt{\mdseries\slshape M}}, this operation returns a plain list of objects in the representation given
by the \texttt{CompatibleVectorFilter} (\ref{CompatibleVectorFilter:for a matrix object}) value of \mbox{\texttt{\mdseries\slshape M}} (provided that such a representation exists), where the $i$\texttt{\symbol{45}}th entry describes the $i$\texttt{\symbol{45}}th row of the input. }

 
\subsection{\textcolor{Chapter }{CompanionMatrix}}\logpage{[ 26, 11, 8 ]}
\hyperdef{L}{X7E06762479A00DF4}{}
{
\noindent\textcolor{FuncColor}{$\triangleright$\enspace\texttt{CompanionMatrix({\mdseries\slshape pol, M})\index{CompanionMatrix@\texttt{CompanionMatrix}!for polynomial and matrix object}
\label{CompanionMatrix:for polynomial and matrix object}
}\hfill{\scriptsize (operation)}}\\
\noindent\textcolor{FuncColor}{$\triangleright$\enspace\texttt{CompanionMatrix({\mdseries\slshape filt, pol, R})\index{CompanionMatrix@\texttt{CompanionMatrix}!for filter, polynomial, and semiring}
\label{CompanionMatrix:for filter, polynomial, and semiring}
}\hfill{\scriptsize (operation)}}\\
\noindent\textcolor{FuncColor}{$\triangleright$\enspace\texttt{CompanionMatrix({\mdseries\slshape pol, R})\index{CompanionMatrix@\texttt{CompanionMatrix}!for polynomial and semiring}
\label{CompanionMatrix:for polynomial and semiring}
}\hfill{\scriptsize (operation)}}\\
\textbf{\indent Returns:\ }
a matrix object



 For a monic, univariate polynomial \mbox{\texttt{\mdseries\slshape pol}} whose coefficients lie in the base domain of the matrix object \mbox{\texttt{\mdseries\slshape M}}, \texttt{CompanionMatrix} returns the companion matrix of \mbox{\texttt{\mdseries\slshape pol}}, as a matrix object with the same \texttt{ConstructingFilter} (\ref{ConstructingFilter:for a matrix object}) and \texttt{BaseDomain} (\ref{BaseDomain:for a matrix object}) values as \mbox{\texttt{\mdseries\slshape M}}. 

 We use column convention, that is, the negatives of the coefficients of \mbox{\texttt{\mdseries\slshape pol}} appear in the last column of the result. 

 If a filter \mbox{\texttt{\mdseries\slshape filt}} and a semiring \mbox{\texttt{\mdseries\slshape R}} are given then the companion matrix is returned as a matrix object with \texttt{ConstructingFilter} (\ref{ConstructingFilter:for a matrix object}) value \mbox{\texttt{\mdseries\slshape filt}} and \texttt{BaseDomain} (\ref{BaseDomain:for a matrix object}) value \mbox{\texttt{\mdseries\slshape R}}. 

 If only \mbox{\texttt{\mdseries\slshape pol}} and a semiring \mbox{\texttt{\mdseries\slshape R}} are given, the representation of the result is guessed from \mbox{\texttt{\mdseries\slshape R}}. 

 If the \texttt{ConstructingFilter} (\ref{ConstructingFilter:for a matrix object}) value of the result implies \texttt{IsCopyable} (\ref{IsCopyable}) then the result is fully mutable. 

 
\begin{Verbatim}[commandchars=!@|,fontsize=\small,frame=single,label=Example]
  !gapprompt@gap>| !gapinput@x:= X( GF(5) );;  pol:= x^3 + x^2 + 2*x + 3;;|
  !gapprompt@gap>| !gapinput@M:= CompanionMatrix( IsPlistMatrixRep, pol, GF(25) );;|
  !gapprompt@gap>| !gapinput@Display( M );|
  <3x3-matrix over GF(5^2):
  [[ 0*Z(5), 0*Z(5), Z(5) ]
   [ Z(5)^0, 0*Z(5), Z(5)^3 ]
   [ 0*Z(5), Z(5)^0, Z(5)^2 ]
  ]>
\end{Verbatim}
 }

 }

  
\section{\textcolor{Chapter }{Operations for Row List Matrix Objects}}\label{Operations for Row List Matrix Objects}
\logpage{[ 26, 12, 0 ]}
\hyperdef{L}{X7D40EE2084A6C976}{}
{
  In general, matrix objects are not lists in the sense of \texttt{IsList} (\ref{IsList}), and they need not behave like lists, that is, they need not obey all the
rules for lists that are stated in Chapter \ref{Lists}. There are situations where one wants to have matrix objects that can on the
one hand benefit from \textsf{GAP}'s method selection, as is explained in Section \ref{Concepts and Rules for Vector and Matrix Objects}, and do on the other hands support access to \textsf{GAP} objects that represent their rows (which are suitable vector objects). Matrix
objects whose \texttt{ConstructingFilter} (\ref{ConstructingFilter:for a matrix object}) value implies \texttt{IsRowListMatrix} (\ref{IsRowListMatrix}) support the operations described in this section. 

 One implementation of such matrices is given by the \texttt{ConstructingFilter} (\ref{ConstructingFilter:for a matrix object}) value \texttt{IsPlistMatrixRep} (\ref{IsPlistMatrixRep}), and any row of these matrices is a vector object in \texttt{IsPlistVectorRep} (\ref{IsPlistVectorRep}). Note that these objects do \emph{not} lie in \texttt{IsList} (\ref{IsList}) (and in particular not in \texttt{IsPlistRep} (\ref{IsPlistRep})), thus we are allowed to define the above operations only restrictively, as
follows. 

 Unbinding an entry in a row or unbinding a row in a matrix is allowed only in
the last position, that is, the vector and matrix objects insist on being
dense. All rows of a matrix must have the same length and the same base
domain. 
\subsection{\textcolor{Chapter }{List Access for a Row List Matrix}}\logpage{[ 26, 12, 1 ]}
\hyperdef{L}{X82C4FCFA808010F8}{}
{
\noindent\textcolor{FuncColor}{$\triangleright$\enspace\texttt{\texttt{\symbol{92}}[\texttt{\symbol{92}}]({\mdseries\slshape M, pos})\index{\texttt{\symbol{92}}[\texttt{\symbol{92}}]@\texttt{\texttt{\symbol{92}}[\texttt{\symbol{92}}]}!for a row list matrix}
\label{bSlash[bSlash]:for a row list matrix}
}\hfill{\scriptsize (operation)}}\\
\textbf{\indent Returns:\ }
a vector object



 If \mbox{\texttt{\mdseries\slshape M}} is a row list matrix and if \mbox{\texttt{\mdseries\slshape pos}} is a positive integer not larger than the number of rows of \mbox{\texttt{\mdseries\slshape M}}, this operation returns the \mbox{\texttt{\mdseries\slshape pos}}\texttt{\symbol{45}}th row of \mbox{\texttt{\mdseries\slshape M}}. 

 It is not specified what happens if \mbox{\texttt{\mdseries\slshape pos}} is larger. }

 
\subsection{\textcolor{Chapter }{List Assignment for a Row List Matrix}}\logpage{[ 26, 12, 2 ]}
\hyperdef{L}{X7F89BB2482D28AAE}{}
{
\noindent\textcolor{FuncColor}{$\triangleright$\enspace\texttt{\texttt{\symbol{92}}[\texttt{\symbol{92}}]\texttt{\symbol{92}}:\texttt{\symbol{92}}=({\mdseries\slshape M, pos, v})\index{\texttt{\symbol{92}}[\texttt{\symbol{92}}]\texttt{\symbol{92}}:\texttt{\symbol{92}}=@\texttt{\texttt{\symbol{92}}[\texttt{\symbol{92}}]\texttt{\symbol{92}}:\texttt{\symbol{92}}=}!for a row list matrix and a vector object}
\label{bSlash[bSlash]bSlash:bSlash=:for a row list matrix and a vector object}
}\hfill{\scriptsize (operation)}}\\
\textbf{\indent Returns:\ }
nothing



 If \mbox{\texttt{\mdseries\slshape M}} is a row list matrix, \mbox{\texttt{\mdseries\slshape v}} is a vector object that can occur as a row in \mbox{\texttt{\mdseries\slshape M}} (that is, \mbox{\texttt{\mdseries\slshape v}} has the same base domain, the right length, and the right vector
representation), and if \mbox{\texttt{\mdseries\slshape pos}} is a positive integer not larger than the number of rows of \mbox{\texttt{\mdseries\slshape M}} plus 1, this operation sets \mbox{\texttt{\mdseries\slshape v}} as the \mbox{\texttt{\mdseries\slshape pos}}\texttt{\symbol{45}}th row of \mbox{\texttt{\mdseries\slshape M}}. 

 In all other situations, it is not specified what happens. }

 
\subsection{\textcolor{Chapter }{Sublist Access for a Row List Matrix}}\logpage{[ 26, 12, 3 ]}
\hyperdef{L}{X807518367C96516F}{}
{
\noindent\textcolor{FuncColor}{$\triangleright$\enspace\texttt{\texttt{\symbol{92}}\texttt{\symbol{123}}\texttt{\symbol{92}}\texttt{\symbol{125}}({\mdseries\slshape M, poss})\index{\texttt{\symbol{92}}\texttt{\symbol{123}}\texttt{\symbol{92}}\texttt{\symbol{125}}@\texttt{\texttt{\symbol{92}}\texttt{\symbol{123}}\texttt{\symbol{92}}\texttt{\symbol{125}}}!for a row list matrix}
\label{bSlash{bSlash}:for a row list matrix}
}\hfill{\scriptsize (operation)}}\\
\textbf{\indent Returns:\ }
a row list matrix



 For a row list matrix \mbox{\texttt{\mdseries\slshape M}} and a list \mbox{\texttt{\mdseries\slshape poss}} of positions, \mbox{\texttt{\mdseries\slshape M}}\texttt{\texttt{\symbol{123}} }\mbox{\texttt{\mdseries\slshape poss}}\texttt{ \texttt{\symbol{125}}} returns a new mutable row list matrix with the same representation as \mbox{\texttt{\mdseries\slshape M}}, whose rows are identical to the rows at the positions in the list \mbox{\texttt{\mdseries\slshape poss}} in \mbox{\texttt{\mdseries\slshape M}}. }

 
\subsection{\textcolor{Chapter }{Sublist Assignment for a Row List Matrix}}\logpage{[ 26, 12, 4 ]}
\hyperdef{L}{X8371789181FA136B}{}
{
\noindent\textcolor{FuncColor}{$\triangleright$\enspace\texttt{\texttt{\symbol{92}}\texttt{\symbol{123}}\texttt{\symbol{92}}\texttt{\symbol{125}}\texttt{\symbol{92}}:\texttt{\symbol{92}}=({\mdseries\slshape M, poss, M2})\index{\texttt{\symbol{92}}\texttt{\symbol{123}}\texttt{\symbol{92}}\texttt{\symbol{125}}\texttt{\symbol{92}}:\texttt{\symbol{92}}=@\texttt{\texttt{\symbol{92}}\texttt{\symbol{123}}\texttt{\symbol{92}}\texttt{\symbol{125}}\texttt{\symbol{92}}:\texttt{\symbol{92}}=}!for row list matrices}
\label{bSlash{bSlash}bSlash:bSlash=:for row list matrices}
}\hfill{\scriptsize (operation)}}\\
\textbf{\indent Returns:\ }
nothing



 For a mutable row list matrix \mbox{\texttt{\mdseries\slshape M}}, a list \mbox{\texttt{\mdseries\slshape poss}} of positions, and a row list matrix \mbox{\texttt{\mdseries\slshape M2}} of the same vector type and with the same base domain, \mbox{\texttt{\mdseries\slshape M}}\texttt{\texttt{\symbol{123}} }\mbox{\texttt{\mdseries\slshape poss}}\texttt{ \texttt{\symbol{125}}:= }\mbox{\texttt{\mdseries\slshape M2}} assigns the rows of \mbox{\texttt{\mdseries\slshape M2}} to the positions \mbox{\texttt{\mdseries\slshape poss}} in the list of rows of \mbox{\texttt{\mdseries\slshape M}}. 

 It is not specified what happens if the resulting range of row positions is
not dense. }

 

\subsection{\textcolor{Chapter }{IsBound\texttt{\symbol{92}}[\texttt{\symbol{92}}] (for a row list matrix)}}
\logpage{[ 26, 12, 5 ]}\nobreak
\hyperdef{L}{X872E63867803ED78}{}
{\noindent\textcolor{FuncColor}{$\triangleright$\enspace\texttt{IsBound\texttt{\symbol{92}}[\texttt{\symbol{92}}]({\mdseries\slshape M, pos})\index{IsBound\texttt{\symbol{92}}[\texttt{\symbol{92}}]@\texttt{IsBound\texttt{\symbol{92}}[\texttt{\symbol{92}}]}!for a row list matrix}
\label{IsBoundbSlash[bSlash]:for a row list matrix}
}\hfill{\scriptsize (operation)}}\\
\textbf{\indent Returns:\ }
\texttt{true} or \texttt{false}



 For a row list matrix \mbox{\texttt{\mdseries\slshape M}} and a positive integer \mbox{\texttt{\mdseries\slshape pos}}, \texttt{IsBound( }\mbox{\texttt{\mdseries\slshape M}}\texttt{[ }\mbox{\texttt{\mdseries\slshape pos}}\texttt{ ] )} returns \texttt{true} if \mbox{\texttt{\mdseries\slshape pos}} is at most the number of rows of \mbox{\texttt{\mdseries\slshape M}}, and \texttt{false} otherwise. }

 

\subsection{\textcolor{Chapter }{Unbind\texttt{\symbol{92}}[\texttt{\symbol{92}}] (for a row list matrix)}}
\logpage{[ 26, 12, 6 ]}\nobreak
\hyperdef{L}{X79328CB280C71DDB}{}
{\noindent\textcolor{FuncColor}{$\triangleright$\enspace\texttt{Unbind\texttt{\symbol{92}}[\texttt{\symbol{92}}]({\mdseries\slshape M, pos})\index{Unbind\texttt{\symbol{92}}[\texttt{\symbol{92}}]@\texttt{Unbind\texttt{\symbol{92}}[\texttt{\symbol{92}}]}!for a row list matrix}
\label{UnbindbSlash[bSlash]:for a row list matrix}
}\hfill{\scriptsize (operation)}}\\
\textbf{\indent Returns:\ }
nothing



 For a mutable row list matrix \mbox{\texttt{\mdseries\slshape M}} with \mbox{\texttt{\mdseries\slshape pos}} rows, \texttt{Unbind( }\mbox{\texttt{\mdseries\slshape M}}\texttt{[ }\mbox{\texttt{\mdseries\slshape pos}}\texttt{ ] )} removes the last row. It is not specified what happens if \mbox{\texttt{\mdseries\slshape pos}} has another value. }

 

\subsection{\textcolor{Chapter }{Add (for a row list matrix and a vector object)}}
\logpage{[ 26, 12, 7 ]}\nobreak
\hyperdef{L}{X7BDD838579E4D2D6}{}
{\noindent\textcolor{FuncColor}{$\triangleright$\enspace\texttt{Add({\mdseries\slshape M, v[, pos]})\index{Add@\texttt{Add}!for a row list matrix and a vector object}
\label{Add:for a row list matrix and a vector object}
}\hfill{\scriptsize (operation)}}\\
\textbf{\indent Returns:\ }
nothing



 For a mutable row list matrix \mbox{\texttt{\mdseries\slshape M}} and a vector object \mbox{\texttt{\mdseries\slshape v}} that is compatible with the rows of \mbox{\texttt{\mdseries\slshape M}}, the two argument version adds \mbox{\texttt{\mdseries\slshape v}} at the end of the list of rows of \mbox{\texttt{\mdseries\slshape M}}. 

 If a positive integer \mbox{\texttt{\mdseries\slshape pos}} is given then \mbox{\texttt{\mdseries\slshape v}} is added in position \mbox{\texttt{\mdseries\slshape pos}}, and all later rows are shifted up by one position. }

 

\subsection{\textcolor{Chapter }{Remove (for a row list matrix)}}
\logpage{[ 26, 12, 8 ]}\nobreak
\hyperdef{L}{X86E355D07A41C025}{}
{\noindent\textcolor{FuncColor}{$\triangleright$\enspace\texttt{Remove({\mdseries\slshape M[, pos]})\index{Remove@\texttt{Remove}!for a row list matrix}
\label{Remove:for a row list matrix}
}\hfill{\scriptsize (operation)}}\\
\textbf{\indent Returns:\ }
a vector object if the removed row exists, otherwise nothing



 For a mutable row list matrix \mbox{\texttt{\mdseries\slshape M}}, this operation removes the \mbox{\texttt{\mdseries\slshape pos}}\texttt{\symbol{45}}th row and shifts the later rows down by one position. The
default for \mbox{\texttt{\mdseries\slshape pos}} is the number of rows of \mbox{\texttt{\mdseries\slshape M}}. 

 If the \mbox{\texttt{\mdseries\slshape pos}}\texttt{\symbol{45}}th row existed in \mbox{\texttt{\mdseries\slshape M}} then it is returned, otherwise nothing is returned. }

 

\subsection{\textcolor{Chapter }{Append (for two row list matrices)}}
\logpage{[ 26, 12, 9 ]}\nobreak
\hyperdef{L}{X82D0359B81F8D442}{}
{\noindent\textcolor{FuncColor}{$\triangleright$\enspace\texttt{Append({\mdseries\slshape M1, M2})\index{Append@\texttt{Append}!for two row list matrices}
\label{Append:for two row list matrices}
}\hfill{\scriptsize (operation)}}\\
\textbf{\indent Returns:\ }
nothing



 For two row list matrices \mbox{\texttt{\mdseries\slshape M1}}, \mbox{\texttt{\mdseries\slshape M2}} such that \mbox{\texttt{\mdseries\slshape M1}} is mutable and such that the \texttt{ConstructingFilter} (\ref{ConstructingFilter:for a matrix object}) and \texttt{BaseDomain} (\ref{BaseDomain:for a matrix object}) values are equal, this operation appends the rows of \mbox{\texttt{\mdseries\slshape M2}} to the rows of \mbox{\texttt{\mdseries\slshape M1}}. }

 

\subsection{\textcolor{Chapter }{ShallowCopy (for a row list matrix)}}
\logpage{[ 26, 12, 10 ]}\nobreak
\hyperdef{L}{X7E234F717BE333EA}{}
{\noindent\textcolor{FuncColor}{$\triangleright$\enspace\texttt{ShallowCopy({\mdseries\slshape M})\index{ShallowCopy@\texttt{ShallowCopy}!for a row list matrix}
\label{ShallowCopy:for a row list matrix}
}\hfill{\scriptsize (operation)}}\\
\textbf{\indent Returns:\ }
a matrix object



 For a row list matrix \mbox{\texttt{\mdseries\slshape M}}, this operation returns a new mutable matrix with the same \texttt{ConstructingFilter} (\ref{ConstructingFilter:for a matrix object}) and \texttt{BaseDomain} (\ref{BaseDomain:for a matrix object}) values as \mbox{\texttt{\mdseries\slshape M}}, which shares its rows with \mbox{\texttt{\mdseries\slshape M}}. }

 

\subsection{\textcolor{Chapter }{ListOp (for a row list matrix)}}
\logpage{[ 26, 12, 11 ]}\nobreak
\hyperdef{L}{X7E9F095E85DED480}{}
{\noindent\textcolor{FuncColor}{$\triangleright$\enspace\texttt{ListOp({\mdseries\slshape M[, func]})\index{ListOp@\texttt{ListOp}!for a row list matrix}
\label{ListOp:for a row list matrix}
}\hfill{\scriptsize (operation)}}\\
\textbf{\indent Returns:\ }
a plain list



 For a row list matrix \mbox{\texttt{\mdseries\slshape M}}, the variant with one argument returns the plain list (see \texttt{IsPlistRep} (\ref{IsPlistRep})) of its rows, and the variant with two arguments returns the plain list of
values of these rows under the function \mbox{\texttt{\mdseries\slshape func}}. }

 }

  
\section{\textcolor{Chapter }{Basic operations for row/column reductions}}\label{Basic operations for row/column reductions}
\logpage{[ 26, 13, 0 ]}
\hyperdef{L}{X7B86A8487B12F9BD}{}
{
  

\subsection{\textcolor{Chapter }{MultMatrixRowLeft}}
\logpage{[ 26, 13, 1 ]}\nobreak
\hyperdef{L}{X7B3997D37CC44FCA}{}
{\noindent\textcolor{FuncColor}{$\triangleright$\enspace\texttt{MultMatrixRowLeft({\mdseries\slshape mat, i, elm})\index{MultMatrixRowLeft@\texttt{MultMatrixRowLeft}}
\label{MultMatrixRowLeft}
}\hfill{\scriptsize (operation)}}\\
\noindent\textcolor{FuncColor}{$\triangleright$\enspace\texttt{MultMatrixRow({\mdseries\slshape mat, i, elm})\index{MultMatrixRow@\texttt{MultMatrixRow}}
\label{MultMatrixRow}
}\hfill{\scriptsize (operation)}}\\
\textbf{\indent Returns:\ }
nothing



 

 Multiplies the \mbox{\texttt{\mdseries\slshape i}}\texttt{\symbol{45}}th row of the mutable matrix \mbox{\texttt{\mdseries\slshape mat}} with the scalar \mbox{\texttt{\mdseries\slshape elm}} from the left in\texttt{\symbol{45}}place. 

 \texttt{MultMatrixRow} is a synonym of \texttt{MultMatrixRowLeft}. This was chosen because linear combinations of rows of matrices are usually
written as $ v \cdot A = [v_1, ... ,v_n] \cdot A$ which multiplies scalars from the left. }

 

\subsection{\textcolor{Chapter }{MultMatrixRowRight}}
\logpage{[ 26, 13, 2 ]}\nobreak
\hyperdef{L}{X794636447E8C5553}{}
{\noindent\textcolor{FuncColor}{$\triangleright$\enspace\texttt{MultMatrixRowRight({\mdseries\slshape M, i, elm})\index{MultMatrixRowRight@\texttt{MultMatrixRowRight}}
\label{MultMatrixRowRight}
}\hfill{\scriptsize (operation)}}\\
\textbf{\indent Returns:\ }
nothing



 

 Multiplies the \mbox{\texttt{\mdseries\slshape i}}\texttt{\symbol{45}}th row of the mutable matrix \mbox{\texttt{\mdseries\slshape M}} with the scalar \mbox{\texttt{\mdseries\slshape elm}} from the right in\texttt{\symbol{45}}place. }

 

\subsection{\textcolor{Chapter }{MultMatrixColumnRight}}
\logpage{[ 26, 13, 3 ]}\nobreak
\hyperdef{L}{X80AF7B267E6B9CE0}{}
{\noindent\textcolor{FuncColor}{$\triangleright$\enspace\texttt{MultMatrixColumnRight({\mdseries\slshape M, i, elm})\index{MultMatrixColumnRight@\texttt{MultMatrixColumnRight}}
\label{MultMatrixColumnRight}
}\hfill{\scriptsize (operation)}}\\
\noindent\textcolor{FuncColor}{$\triangleright$\enspace\texttt{MultMatrixColumn({\mdseries\slshape M, i, elm})\index{MultMatrixColumn@\texttt{MultMatrixColumn}}
\label{MultMatrixColumn}
}\hfill{\scriptsize (operation)}}\\
\textbf{\indent Returns:\ }
nothing



 

 Multiplies the \mbox{\texttt{\mdseries\slshape i}}\texttt{\symbol{45}}th column of the mutable matrix \mbox{\texttt{\mdseries\slshape M}} with the scalar \mbox{\texttt{\mdseries\slshape elm}} from the right in\texttt{\symbol{45}}place. 

 \texttt{MultMatrixColumn} is a synonym of \texttt{MultMatrixColumnRight}. This was chosen because linear combinations of columns of matrices are
usually written as $A \cdot v^T = A \cdot [v_1, ... ,v_n]^T$ which multiplies scalars from the right. }

 

\subsection{\textcolor{Chapter }{MultMatrixColumnLeft}}
\logpage{[ 26, 13, 4 ]}\nobreak
\hyperdef{L}{X843DAFE37F347471}{}
{\noindent\textcolor{FuncColor}{$\triangleright$\enspace\texttt{MultMatrixColumnLeft({\mdseries\slshape M, i, elm})\index{MultMatrixColumnLeft@\texttt{MultMatrixColumnLeft}}
\label{MultMatrixColumnLeft}
}\hfill{\scriptsize (operation)}}\\
\textbf{\indent Returns:\ }
nothing



 

 Multiplies the \mbox{\texttt{\mdseries\slshape i}}\texttt{\symbol{45}}th column of the mutable matrix \mbox{\texttt{\mdseries\slshape M}} with the scalar \mbox{\texttt{\mdseries\slshape elm}} from the left in\texttt{\symbol{45}}place. }

 

\subsection{\textcolor{Chapter }{AddMatrixRowsLeft}}
\logpage{[ 26, 13, 5 ]}\nobreak
\hyperdef{L}{X8662EB748629502F}{}
{\noindent\textcolor{FuncColor}{$\triangleright$\enspace\texttt{AddMatrixRowsLeft({\mdseries\slshape M, i, j, elm})\index{AddMatrixRowsLeft@\texttt{AddMatrixRowsLeft}}
\label{AddMatrixRowsLeft}
}\hfill{\scriptsize (operation)}}\\
\noindent\textcolor{FuncColor}{$\triangleright$\enspace\texttt{AddMatrixRows({\mdseries\slshape M, i, j, elm})\index{AddMatrixRows@\texttt{AddMatrixRows}}
\label{AddMatrixRows}
}\hfill{\scriptsize (operation)}}\\
\textbf{\indent Returns:\ }
nothing



 

 Adds the product of \mbox{\texttt{\mdseries\slshape elm}} with the \mbox{\texttt{\mdseries\slshape j}}\texttt{\symbol{45}}th row of the mutable matrix \mbox{\texttt{\mdseries\slshape M}} to its \mbox{\texttt{\mdseries\slshape i}}\texttt{\symbol{45}}th row in\texttt{\symbol{45}}place. The \mbox{\texttt{\mdseries\slshape j}}\texttt{\symbol{45}}th row is multiplied with \mbox{\texttt{\mdseries\slshape elm}} from the left. 

 \texttt{AddMatrixRows} is a synonym of \texttt{AddMatrixRowsLeft}. This was chosen because linear combinations of rows of matrices are usually
written as $ v \cdot A = [v_1, ... ,v_n] \cdot A$ which multiplies scalars from the left. }

 

\subsection{\textcolor{Chapter }{AddMatrixRowsRight}}
\logpage{[ 26, 13, 6 ]}\nobreak
\hyperdef{L}{X7CD05EE984614AB6}{}
{\noindent\textcolor{FuncColor}{$\triangleright$\enspace\texttt{AddMatrixRowsRight({\mdseries\slshape M, i, j, elm})\index{AddMatrixRowsRight@\texttt{AddMatrixRowsRight}}
\label{AddMatrixRowsRight}
}\hfill{\scriptsize (operation)}}\\
\textbf{\indent Returns:\ }
nothing



 

 Adds the product of \mbox{\texttt{\mdseries\slshape elm}} with the \mbox{\texttt{\mdseries\slshape j}}\texttt{\symbol{45}}th row of the mutable matrix \mbox{\texttt{\mdseries\slshape M}} to its \mbox{\texttt{\mdseries\slshape i}}\texttt{\symbol{45}}th row in\texttt{\symbol{45}}place. The \mbox{\texttt{\mdseries\slshape j}}\texttt{\symbol{45}}th row is multiplied with \mbox{\texttt{\mdseries\slshape elm}} from the right. }

 

\subsection{\textcolor{Chapter }{AddMatrixColumnsRight}}
\logpage{[ 26, 13, 7 ]}\nobreak
\hyperdef{L}{X7B1E1E417CA267A3}{}
{\noindent\textcolor{FuncColor}{$\triangleright$\enspace\texttt{AddMatrixColumnsRight({\mdseries\slshape M, i, j, elm})\index{AddMatrixColumnsRight@\texttt{AddMatrixColumnsRight}}
\label{AddMatrixColumnsRight}
}\hfill{\scriptsize (operation)}}\\
\noindent\textcolor{FuncColor}{$\triangleright$\enspace\texttt{AddMatrixColumns({\mdseries\slshape M, i, j, elm})\index{AddMatrixColumns@\texttt{AddMatrixColumns}}
\label{AddMatrixColumns}
}\hfill{\scriptsize (operation)}}\\
\textbf{\indent Returns:\ }
nothing



 

 Adds the product of \mbox{\texttt{\mdseries\slshape elm}} with the \mbox{\texttt{\mdseries\slshape j}}\texttt{\symbol{45}}th column of the mutable matrix \mbox{\texttt{\mdseries\slshape M}} to its \mbox{\texttt{\mdseries\slshape i}}\texttt{\symbol{45}}th column in\texttt{\symbol{45}}place. The \mbox{\texttt{\mdseries\slshape j}}\texttt{\symbol{45}}th column is multiplied with \mbox{\texttt{\mdseries\slshape elm}} from the right. 

 \texttt{AddMatrixColumns} is a synonym of \texttt{AddMatrixColumnsRight}. This was chosen because linear combinations of columns of matrices are
usually written as $A \cdot v^T = A \cdot [v_1, ... ,v_n]^T$ which multiplies scalars from the right. }

 

\subsection{\textcolor{Chapter }{AddMatrixColumnsLeft}}
\logpage{[ 26, 13, 8 ]}\nobreak
\hyperdef{L}{X85ECB8C87DFD8F32}{}
{\noindent\textcolor{FuncColor}{$\triangleright$\enspace\texttt{AddMatrixColumnsLeft({\mdseries\slshape M, i, j, elm})\index{AddMatrixColumnsLeft@\texttt{AddMatrixColumnsLeft}}
\label{AddMatrixColumnsLeft}
}\hfill{\scriptsize (operation)}}\\
\textbf{\indent Returns:\ }
nothing



 

 Adds the product of \mbox{\texttt{\mdseries\slshape elm}} with the \mbox{\texttt{\mdseries\slshape j}}\texttt{\symbol{45}}th column of the mutable matrix \mbox{\texttt{\mdseries\slshape M}} to its \mbox{\texttt{\mdseries\slshape i}}\texttt{\symbol{45}}th column in\texttt{\symbol{45}}place. The \mbox{\texttt{\mdseries\slshape j}}\texttt{\symbol{45}}th column is multiplied with \mbox{\texttt{\mdseries\slshape elm}} from the left. }

 

\subsection{\textcolor{Chapter }{PositionNonZeroInRow}}
\logpage{[ 26, 13, 9 ]}\nobreak
\hyperdef{L}{X86B9C2EC8794BFE0}{}
{\noindent\textcolor{FuncColor}{$\triangleright$\enspace\texttt{PositionNonZeroInRow({\mdseries\slshape M, i[, from]})\index{PositionNonZeroInRow@\texttt{PositionNonZeroInRow}}
\label{PositionNonZeroInRow}
}\hfill{\scriptsize (operation)}}\\
\textbf{\indent Returns:\ }
a positive integer



 

 Returns the position of the first nonzero entry in the \mbox{\texttt{\mdseries\slshape i}}\texttt{\symbol{45}}th row of the matrix \mbox{\texttt{\mdseries\slshape M}}, or \texttt{NrCols( M ) + 1} if the row is zero. If the optional argument \mbox{\texttt{\mdseries\slshape from}} is given, the search starts after position \mbox{\texttt{\mdseries\slshape from}}. }

 

\subsection{\textcolor{Chapter }{SwapMatrixRows}}
\logpage{[ 26, 13, 10 ]}\nobreak
\hyperdef{L}{X87CCA3117F6B3F0D}{}
{\noindent\textcolor{FuncColor}{$\triangleright$\enspace\texttt{SwapMatrixRows({\mdseries\slshape M, i, j})\index{SwapMatrixRows@\texttt{SwapMatrixRows}}
\label{SwapMatrixRows}
}\hfill{\scriptsize (operation)}}\\
\textbf{\indent Returns:\ }
nothing



 

 Swaps the \mbox{\texttt{\mdseries\slshape i}}\texttt{\symbol{45}}th row and \mbox{\texttt{\mdseries\slshape j}}\texttt{\symbol{45}}th row of a mutable matrix \mbox{\texttt{\mdseries\slshape M}}. }

 

\subsection{\textcolor{Chapter }{SwapMatrixColumns}}
\logpage{[ 26, 13, 11 ]}\nobreak
\hyperdef{L}{X824C8A347EB9D499}{}
{\noindent\textcolor{FuncColor}{$\triangleright$\enspace\texttt{SwapMatrixColumns({\mdseries\slshape M, i, j})\index{SwapMatrixColumns@\texttt{SwapMatrixColumns}}
\label{SwapMatrixColumns}
}\hfill{\scriptsize (operation)}}\\
\textbf{\indent Returns:\ }
nothing



 

 Swaps the \mbox{\texttt{\mdseries\slshape i}}\texttt{\symbol{45}}th column and \mbox{\texttt{\mdseries\slshape j}}\texttt{\symbol{45}}th column of a mutable matrix \mbox{\texttt{\mdseries\slshape M}}. }

 

\subsection{\textcolor{Chapter }{AddMatrix}}
\logpage{[ 26, 13, 12 ]}\nobreak
\hyperdef{L}{X7E3B10A67BC7F006}{}
{\noindent\textcolor{FuncColor}{$\triangleright$\enspace\texttt{AddMatrix({\mdseries\slshape M, N[, c]})\index{AddMatrix@\texttt{AddMatrix}}
\label{AddMatrix}
}\hfill{\scriptsize (operation)}}\\
\textbf{\indent Returns:\ }
nothing



 Computes the calculation $M + N \cdot c$ in\texttt{\symbol{45}}place, storing the result in \mbox{\texttt{\mdseries\slshape M}}. If the optional argument \mbox{\texttt{\mdseries\slshape c}} is omitted, then \mbox{\texttt{\mdseries\slshape N}} is added directly. The matrices must have the same dimensions, otherwise the
result is undefined. Specialized methods may be defined only when \mbox{\texttt{\mdseries\slshape M}} and \mbox{\texttt{\mdseries\slshape N}} have the same representation. If both of the matrices are
lists\texttt{\symbol{45}}of\texttt{\symbol{45}}lists, then the operation is
delegated row by row to \texttt{AddRowVector} (\ref{AddRowVector}). 
\begin{Verbatim}[commandchars=!@|,fontsize=\small,frame=single,label=Example]
  !gapprompt@gap>| !gapinput@mat1 := [ [ 1, 2 ], [ 3, 4 ] ];|
  [ [ 1, 2 ], [ 3, 4 ] ]
  !gapprompt@gap>| !gapinput@mat2 := [ [ 1, 0 ], [ 3, -1 ] ];|
  [ [ 1, 0 ], [ 3, -1 ] ]
  !gapprompt@gap>| !gapinput@AddMatrix( mat1, mat2, 2 );|
  !gapprompt@gap>| !gapinput@mat1;|
  [ [ 3, 2 ], [ 9, 2 ] ]
  !gapprompt@gap>| !gapinput@mat2;|
  [ [ 1, 0 ], [ 3, -1 ] ]
  !gapprompt@gap>| !gapinput@AddMatrix( mat1, [ [ 1, 0], [ 3, -1] ] );|
  !gapprompt@gap>| !gapinput@mat1;|
  [ [ 4, 2 ], [ 12, 1 ] ]
\end{Verbatim}
 }

 

\subsection{\textcolor{Chapter }{MultMatrix}}
\logpage{[ 26, 13, 13 ]}\nobreak
\hyperdef{L}{X85DAF850790CD476}{}
{\noindent\textcolor{FuncColor}{$\triangleright$\enspace\texttt{MultMatrix({\mdseries\slshape mat, c})\index{MultMatrix@\texttt{MultMatrix}}
\label{MultMatrix}
}\hfill{\scriptsize (operation)}}\\
\noindent\textcolor{FuncColor}{$\triangleright$\enspace\texttt{MultMatrixLeft({\mdseries\slshape mat, c})\index{MultMatrixLeft@\texttt{MultMatrixLeft}}
\label{MultMatrixLeft}
}\hfill{\scriptsize (operation)}}\\
\noindent\textcolor{FuncColor}{$\triangleright$\enspace\texttt{MultMatrixRight({\mdseries\slshape mat, c})\index{MultMatrixRight@\texttt{MultMatrixRight}}
\label{MultMatrixRight}
}\hfill{\scriptsize (operation)}}\\
\textbf{\indent Returns:\ }
nothing



 These functions multiply the entries of \mbox{\texttt{\mdseries\slshape mat}} by \mbox{\texttt{\mdseries\slshape c}} in\texttt{\symbol{45}}place. \texttt{MultMatrixRight} performs the operation \mbox{\texttt{\mdseries\slshape mat}}\texttt{*}\mbox{\texttt{\mdseries\slshape c}}, whereas \texttt{MultMatrixLeft} performs the operation \mbox{\texttt{\mdseries\slshape c}}\texttt{*}\mbox{\texttt{\mdseries\slshape mat}} and \texttt{MultMatrix} is an alias for \texttt{MultMatrixLeft}. In all of these, if the matrix \mbox{\texttt{\mdseries\slshape mat}} is a lists\texttt{\symbol{45}}of\texttt{\symbol{45}}lists, then the operation
is delegated row by row to \texttt{MultVectorRight} (\ref{MultVectorRight}) and \texttt{MultVectorLeft} (\ref{MultVectorLeft}). 
\begin{Verbatim}[commandchars=!@|,fontsize=\small,frame=single,label=Example]
  !gapprompt@gap>| !gapinput@mat1 := [ [ 1, 2 ], [ 3, 4 ] ];|
  [ [ 1, 2 ], [ 3, 4 ] ]
  !gapprompt@gap>| !gapinput@MultMatrixRight(mat1, -2);|
  !gapprompt@gap>| !gapinput@mat1;|
  [ [ -2, -4 ], [ -6, -8 ] ]
  !gapprompt@gap>| !gapinput@MultMatrix(mat1, -2); # Note that this is the same as calling MultMatrixLeft(mat1, -2)|
  !gapprompt@gap>| !gapinput@mat1;|
  [ [ 4, 8 ], [ 12, 16 ] ]
  !gapprompt@gap>| !gapinput@A := FreeAssociativeAlgebra(Rationals, 2);|
  <algebra over Rationals, with 2 generators>
  !gapprompt@gap>| !gapinput@mat2 := [ [ A.1, A.2 ], [ A.1 * 2, A.2 * 3 ] ];|
  [ [ (1)*x.1, (1)*x.2 ], [ (2)*x.1, (3)*x.2 ] ]
  !gapprompt@gap>| !gapinput@MultMatrixLeft(mat2, A.1);|
  !gapprompt@gap>| !gapinput@mat2;|
  [ [ (1)*x.1^2, (1)*x.1*x.2 ], [ (2)*x.1^2, (3)*x.1*x.2 ] ]
  !gapprompt@gap>| !gapinput@mat2 := [ [ A.1, A.2 ], [ A.1 * 2, A.2 * 3 ] ];|
  [ [ (1)*x.1, (1)*x.2 ], [ (2)*x.1, (3)*x.2 ] ]
  !gapprompt@gap>| !gapinput@MultMatrixRight(mat2, A.1);|
  !gapprompt@gap>| !gapinput@mat2;|
  [ [ (1)*x.1^2, (1)*x.2*x.1 ], [ (2)*x.1^2, (3)*x.2*x.1 ] ]
\end{Verbatim}
 }

 }

  
\section{\textcolor{Chapter }{Implementing New Vector and Matrix Objects Types}}\label{Implementing New Vector and Matrix Objects Types}
\logpage{[ 26, 14, 0 ]}
\hyperdef{L}{X7BEE647484978886}{}
{
  The first step in the design of a new type of vector or matrix objects is to
create a new filter that serves as the \texttt{ConstructingFilter} (\ref{ConstructingFilter:for a vector object}) of the new objects, see the sections \ref{Available Representations of Vector Objects} and \ref{Available Representations of Matrix Objects} for an overview of such filters that are already available. 

 Here we list those operations for vector and matrix objects for which no
default methods can be installed. When one implements a new type of vector or
matrix objects then one has to install specific methods at least for these
operations, in order to make the objects behave as described in this chapter.
It is advisable to install specific methods also for other operations, for
performance reasons. The installations of default methods can be found in the
file \texttt{lib/matobj.gi} of the \textsf{GAP} distribution. There one can check for which operations it makes sense to
overload them for the new type of vector or matrix objects. Note that the
specific methods must be installed with \texttt{InstallTagBasedMethod} (\ref{InstallTagBasedMethod}) whenever the default method is installed with this function. 

 \emph{Vector objects} 
\begin{itemize}
\item  \texttt{BaseDomain} (\ref{BaseDomain:for a vector object}), 
\item  \texttt{Length} (\ref{Length:for a vector object}), 
\item  \texttt{\texttt{\symbol{92}}[\texttt{\symbol{92}}]} (\ref{bSlash[bSlash]:for a vector object and an integer}), 
\item  \texttt{\texttt{\symbol{92}}[\texttt{\symbol{92}}]\texttt{\symbol{92}}:\texttt{\symbol{92}}=} (\ref{bSlash[bSlash]bSlash:bSlash=:for a vector object and an integer}) (with consistency checks if the global option \texttt{check} is not set to \texttt{false}), 
\item  \texttt{\texttt{\symbol{92}}{\textless}} (\ref{bSlash<}) (see \texttt{\texttt{\symbol{92}}{\textless}} (\ref{bSlash<:for two vector objects})), 
\item  \texttt{ConstructingFilter} (\ref{ConstructingFilter:for a vector object}), 
\item  \texttt{NewVector} (\ref{NewVector}) (with consistency checks if the global option \texttt{check} is not set to \texttt{false}, install the method with \texttt{InstallTagBasedMethod} (\ref{InstallTagBasedMethod})). 
\end{itemize}
 

 \emph{Matrix objects} 
\begin{itemize}
\item  \texttt{BaseDomain} (\ref{BaseDomain:for a matrix object}), 
\item  \texttt{NumberRows} (\ref{NumberRows:for a matrix object}), 
\item  \texttt{NumberColumns} (\ref{NumberColumns:for a matrix object}), 
\item  \texttt{MatElm} (\ref{MatElm}), 
\item  \texttt{SetMatElm} (\ref{SetMatElm}) (with consistency checks if the global option \texttt{check} is not set to \texttt{false}), 
\item  \texttt{\texttt{\symbol{92}}{\textless}} (\ref{bSlash<}) (see \texttt{\texttt{\symbol{92}}{\textless}} (\ref{bSlash<:for two matrix objects})), 
\item  \texttt{ConstructingFilter} (\ref{ConstructingFilter:for a matrix object}), 
\item  \texttt{CompatibleVectorFilter} (\ref{CompatibleVectorFilter:for a matrix object}), 
\item  \texttt{NewMatrix} (\ref{NewMatrix}) (with consistency checks if the global option \texttt{check} is not set to \texttt{false}, install the method with \texttt{InstallTagBasedMethod} (\ref{InstallTagBasedMethod})). 
\end{itemize}
 

 Methods for \texttt{NewVector} (\ref{NewVector}) and \texttt{NewMatrix} (\ref{NewMatrix}) must check their arguments for consistency (do the given filter and base
domain fit together, are the entries compatible with the given base domain, is
the number of matrix entries a multiple of the given number of columns) except
if the global option \texttt{check} is set to \texttt{false}. (See Chapter \ref{Options Stack} for information about global options.) The same holds for methods for
operations that modify mutable vector or matrix objects, such as \texttt{\texttt{\symbol{92}}[\texttt{\symbol{92}}]\texttt{\symbol{92}}:\texttt{\symbol{92}}=} (\ref{bSlash[bSlash]bSlash:bSlash=:for a vector object and an integer}), \texttt{SetMatElm} (\ref{SetMatElm}), \texttt{CopySubVector} (\ref{CopySubVector}), \texttt{CopySubMatrix} (\ref{CopySubMatrix}), and for those methods of \texttt{Vector} (\ref{Vector:for filter, base domain, and list}) and \texttt{Matrix} (\ref{Matrix:for filter, base domain, list, ncols}) that do not delegate to \texttt{NewVector} (\ref{NewVector}) and \texttt{NewMatrix} (\ref{NewMatrix}), respectively. }

  
\section{\textcolor{Chapter }{Available Representations of Vector Objects}}\label{Available Representations of Vector Objects}
\logpage{[ 26, 15, 0 ]}
\hyperdef{L}{X82EEE1D37A94F807}{}
{
  The following filters define vector objects for which the the functionality
described in this chapter is supported. 

\subsection{\textcolor{Chapter }{IsGF2VectorRep}}
\logpage{[ 26, 15, 1 ]}\nobreak
\hyperdef{L}{X7C8050938691A914}{}
{\noindent\textcolor{FuncColor}{$\triangleright$\enspace\texttt{IsGF2VectorRep({\mdseries\slshape obj})\index{IsGF2VectorRep@\texttt{IsGF2VectorRep}}
\label{IsGF2VectorRep}
}\hfill{\scriptsize (Representation)}}\\


 An object \mbox{\texttt{\mdseries\slshape obj}} in \texttt{IsGF2VectorRep} describes a vector object (see \texttt{IsVectorObj} (\ref{IsVectorObj})) with entries in the finite field with $2$ elements. 

 \texttt{IsGF2VectorRep} implies \texttt{IsCopyable} (\ref{IsCopyable}), thus vector objects in this representation can be mutable. }

 

\subsection{\textcolor{Chapter }{Is8BitVectorRep}}
\logpage{[ 26, 15, 2 ]}\nobreak
\hyperdef{L}{X82A643007EC6D1CA}{}
{\noindent\textcolor{FuncColor}{$\triangleright$\enspace\texttt{Is8BitVectorRep({\mdseries\slshape obj})\index{Is8BitVectorRep@\texttt{Is8BitVectorRep}}
\label{Is8BitVectorRep}
}\hfill{\scriptsize (Representation)}}\\


 An object \mbox{\texttt{\mdseries\slshape obj}} in \texttt{Is8BitVectorRep} describes a vector object (see \texttt{IsVectorObj} (\ref{IsVectorObj})) with entries in a finite field with $q$ elements, for $3 \leq q \leq 256$. The base domain of \mbox{\texttt{\mdseries\slshape obj}} is not necessarily the smallest field that contains all matrix entries. 

 \texttt{Is8BitVectorRep} implies \texttt{IsCopyable} (\ref{IsCopyable}), thus vector objects in this representation can be mutable. }

 

\subsection{\textcolor{Chapter }{IsPlistVectorRep}}
\logpage{[ 26, 15, 3 ]}\nobreak
\hyperdef{L}{X83262B7085FA94E3}{}
{\noindent\textcolor{FuncColor}{$\triangleright$\enspace\texttt{IsPlistVectorRep({\mdseries\slshape obj})\index{IsPlistVectorRep@\texttt{IsPlistVectorRep}}
\label{IsPlistVectorRep}
}\hfill{\scriptsize (Representation)}}\\


 An object \mbox{\texttt{\mdseries\slshape obj}} in \texttt{IsPlistVectorRep} describes a vector object (see \texttt{IsVectorObj} (\ref{IsVectorObj})) that can occur as a row in a row list matrix (see Section \ref{Operations for Row List Matrix Objects}). 

 \texttt{IsPlistVectorRep} implies \texttt{IsCopyable} (\ref{IsCopyable}), thus vector objects in this representation can be mutable. }

 

\subsection{\textcolor{Chapter }{IsZmodnZVectorRep}}
\logpage{[ 26, 15, 4 ]}\nobreak
\hyperdef{L}{X8730DB7D7E7DA883}{}
{\noindent\textcolor{FuncColor}{$\triangleright$\enspace\texttt{IsZmodnZVectorRep({\mdseries\slshape obj})\index{IsZmodnZVectorRep@\texttt{IsZmodnZVectorRep}}
\label{IsZmodnZVectorRep}
}\hfill{\scriptsize (Representation)}}\\


 An object \mbox{\texttt{\mdseries\slshape obj}} in \texttt{IsZmodnZVectorRep} describes a vector object (see \texttt{IsVectorObj} (\ref{IsVectorObj})) with entries in a residue class ring of the ring of integers (see \texttt{ZmodnZ} (\ref{ZmodnZ})). This ring is the base domain (see \texttt{BaseDomain} (\ref{BaseDomain:for a vector object})) of \mbox{\texttt{\mdseries\slshape obj}}. 

 \texttt{IsZmodnZVectorRep} implies \texttt{IsCopyable} (\ref{IsCopyable}), thus matrix objects in this representation can be mutable. }

 }

  
\section{\textcolor{Chapter }{Available Representations of Matrix Objects}}\label{Available Representations of Matrix Objects}
\logpage{[ 26, 16, 0 ]}
\hyperdef{L}{X7CFD844C7D80D541}{}
{
  The following filters define matrix objects for which the the functionality
described in this chapter is supported. 

\subsection{\textcolor{Chapter }{IsGF2MatrixRep}}
\logpage{[ 26, 16, 1 ]}\nobreak
\hyperdef{L}{X7F6078FF81E912E7}{}
{\noindent\textcolor{FuncColor}{$\triangleright$\enspace\texttt{IsGF2MatrixRep({\mdseries\slshape obj})\index{IsGF2MatrixRep@\texttt{IsGF2MatrixRep}}
\label{IsGF2MatrixRep}
}\hfill{\scriptsize (Representation)}}\\


 An object \mbox{\texttt{\mdseries\slshape obj}} in \texttt{IsGF2MatrixRep} describes a matrix object (see \texttt{IsMatrixObj} (\ref{IsMatrixObj})) with entries in the finite field with $2$ elements, which behaves like the list of its rows (see \texttt{IsRowListMatrix} (\ref{IsRowListMatrix})). The base domain of \mbox{\texttt{\mdseries\slshape obj}} is the field with $2$ elements. 

 \texttt{IsGF2MatrixRep} implies \texttt{IsCopyable} (\ref{IsCopyable}), thus vector objects in this representation can be mutable. }

 

\subsection{\textcolor{Chapter }{Is8BitMatrixRep}}
\logpage{[ 26, 16, 2 ]}\nobreak
\hyperdef{L}{X81466B6C7CAC3A7B}{}
{\noindent\textcolor{FuncColor}{$\triangleright$\enspace\texttt{Is8BitMatrixRep({\mdseries\slshape obj})\index{Is8BitMatrixRep@\texttt{Is8BitMatrixRep}}
\label{Is8BitMatrixRep}
}\hfill{\scriptsize (Representation)}}\\


 An object \mbox{\texttt{\mdseries\slshape obj}} in \texttt{Is8BitMatrixRep} describes a matrix object (see \texttt{IsMatrixObj} (\ref{IsMatrixObj})) that behaves like the list of its rows (see \texttt{IsRowListMatrix} (\ref{IsRowListMatrix})). The base domain of \mbox{\texttt{\mdseries\slshape obj}} is a field that contains all matrix entries (but not necessarily the smallest
such field), it must be a finite field with $q$ elements, for $3 \leq q \leq 256$. 

 \texttt{Is8BitMatrixRep} implies \texttt{IsCopyable} (\ref{IsCopyable}), thus matrix objects in this representation can be mutable. }

 

\subsection{\textcolor{Chapter }{IsPlistMatrixRep}}
\logpage{[ 26, 16, 3 ]}\nobreak
\hyperdef{L}{X80C6031C7DB31A15}{}
{\noindent\textcolor{FuncColor}{$\triangleright$\enspace\texttt{IsPlistMatrixRep({\mdseries\slshape obj})\index{IsPlistMatrixRep@\texttt{IsPlistMatrixRep}}
\label{IsPlistMatrixRep}
}\hfill{\scriptsize (Representation)}}\\


 An object \mbox{\texttt{\mdseries\slshape obj}} in \texttt{IsPlistMatrixRep} describes a matrix object (see \texttt{IsMatrixObj} (\ref{IsMatrixObj})) that behaves similar to a list of its rows, in the sense of \texttt{IsRowListMatrix} (\ref{IsRowListMatrix}). 

 In particular, the rows can be accessed via $M[i]$, and they are vector objects in \texttt{IsPlistVectorRep} (\ref{IsPlistVectorRep}). This representation is useful if one wants to work explicitly with such row
objects. 

 For most purposes, \texttt{IsGenericMatrixRep} (\ref{IsGenericMatrixRep}) is the better choice. Unlike \texttt{IsPlistMatrixRep}, it does not support direct row access, but it is intended as the
general\texttt{\symbol{45}}purpose representation for matrix objects. 

 \texttt{IsPlistMatrixRep} implies \texttt{IsCopyable} (\ref{IsCopyable}), thus matrix objects in this representation can be mutable. }

 

\subsection{\textcolor{Chapter }{IsGenericMatrixRep}}
\logpage{[ 26, 16, 4 ]}\nobreak
\hyperdef{L}{X8206177679151859}{}
{\noindent\textcolor{FuncColor}{$\triangleright$\enspace\texttt{IsGenericMatrixRep({\mdseries\slshape obj})\index{IsGenericMatrixRep@\texttt{IsGenericMatrixRep}}
\label{IsGenericMatrixRep}
}\hfill{\scriptsize (Representation)}}\\


 An object \mbox{\texttt{\mdseries\slshape obj}} in \texttt{IsGenericMatrixRep} describes a matrix object (see \texttt{IsMatrixObj} (\ref{IsMatrixObj})) whose entries are stored as a dense plain list of dense plain row lists. 

 Unlike \texttt{IsPlistMatrixRep} (\ref{IsPlistMatrixRep}), this representation is not a row list matrix, so direct row access via $M[i]$ is not supported. Instead, it is intended as the
general\texttt{\symbol{45}}purpose representation for matrix objects. 

 This representation supports the usual matrix operations and is often more
efficient than \texttt{IsPlistMatrixRep} (\ref{IsPlistMatrixRep}). If existing code uses row access, it is often possible to replace this by
the operations from Section \ref{Basic operations for row/column reductions}. 

 Use \texttt{IsPlistMatrixRep} (\ref{IsPlistMatrixRep}) instead if you need the rows to be available as vector objects in \texttt{IsPlistVectorRep} (\ref{IsPlistVectorRep}). }

 

\subsection{\textcolor{Chapter }{IsZmodnZMatrixRep}}
\logpage{[ 26, 16, 5 ]}\nobreak
\hyperdef{L}{X84D0F3117DA86850}{}
{\noindent\textcolor{FuncColor}{$\triangleright$\enspace\texttt{IsZmodnZMatrixRep({\mdseries\slshape obj})\index{IsZmodnZMatrixRep@\texttt{IsZmodnZMatrixRep}}
\label{IsZmodnZMatrixRep}
}\hfill{\scriptsize (Representation)}}\\


 An object \mbox{\texttt{\mdseries\slshape obj}} in \texttt{IsZmodnZMatrixRep} describes a matrix object (see \texttt{IsMatrixObj} (\ref{IsMatrixObj})) that behaves like the list of its rows (see \texttt{IsRowListMatrix} (\ref{IsRowListMatrix})). The matrix entries lie in a residue class ring of the ring of integers (see \texttt{ZmodnZ} (\ref{ZmodnZ})). This ring is the base domain (see \texttt{BaseDomain} (\ref{BaseDomain:for a vector object})) of \mbox{\texttt{\mdseries\slshape obj}}. 

 \texttt{IsZmodnZMatrixRep} implies \texttt{IsCopyable} (\ref{IsCopyable}), thus matrix objects in this representation can be mutable. }

 }

  
\section{\textcolor{Chapter }{Groups Consisting of Matrix Objects}}\label{Groups Consisting of Matrix Objects}
\logpage{[ 26, 17, 0 ]}
\hyperdef{L}{X82BD536483B78DB9}{}
{
  One application of matrix objects is to construct groups whose elements are
matrix objects. The following conventions hold for such a group \texttt{G}. 
\begin{itemize}
\item  All elements of \texttt{G} have the same \texttt{ConstructingFilter} (\ref{ConstructingFilter:for a matrix object}) and the same \texttt{BaseDomain} (\ref{BaseDomain:for a matrix object}). For a given matrix or matrix object \texttt{M}, a matrix object that can be multiplied with elements of \texttt{G} can be constructed as \texttt{Matrix( M, Representative( G ) )}, and the result can be tested for membership in \texttt{G}. 
\end{itemize}
 

\subsection{\textcolor{Chapter }{ConstructingFiltersForMatrixGroupElements}}
\logpage{[ 26, 17, 1 ]}\nobreak
\hyperdef{L}{X7F7E8B6F7D008CF2}{}
{\noindent\textcolor{FuncColor}{$\triangleright$\enspace\texttt{ConstructingFiltersForMatrixGroupElements\index{ConstructingFiltersForMatrixGroupElements@\texttt{Constructing}\-\texttt{Filters}\-\texttt{For}\-\texttt{Matrix}\-\texttt{Group}\-\texttt{Elements}}
\label{ConstructingFiltersForMatrixGroupElements}
}\hfill{\scriptsize (global variable)}}\\


 This is a list of filters \texttt{f} that are supported as the first argument of group constructors such as \texttt{SymplecticGroup} (\ref{SymplecticGroup:for dimension and a ring}), meaning that one asks for a matrix group whose elements are in \texttt{f}. 

 If one creates a new representation of matrix objects such that it makes sense
to create matrix groups from them then the underlying filter should be added
to \texttt{ConstructingFiltersForMatrixGroupElements}. }

 }

  
\section{\textcolor{Chapter }{How to Write Code for Vector and Matrix Objects}}\label{How to Write Code for Vector and Matrix Objects}
\logpage{[ 26, 18, 0 ]}
\hyperdef{L}{X786024B686AFAD91}{}
{
  Vector and matrix objects have a number of different representations in \textsf{GAP} which are optimised for different things; see the list \texttt{ConstructingFiltersForMatrixGroupElements} (\ref{ConstructingFiltersForMatrixGroupElements}) for examples of available \texttt{ConstructingFilter} (\ref{ConstructingFilter:for a matrix object}) values of matrix objects. Most \textsf{GAP} functions accepting vector or matrix objects as arguments do not depend on any
particular representation for these objects. If the output of such a function
involves vector and matrix objects, then these are expected to have the same
representations as the inputs. For example, a function that computes the
Kronecker product of two matrix objects \texttt{mat1}, \texttt{mat2} may take the two inputs, which must have the same \texttt{ConstructingFilter} (\ref{ConstructingFilter:for a matrix object}) value, create a new matrix object which also lies in this filter, for example
by calling \texttt{ZeroMatrix( m, n, mat1 )}, and then set the entries in this matrix. 

 Functions like \texttt{ZeroMatrix} (\ref{ZeroMatrix:for dimensions and matrix object}), \texttt{IdentityMatrix} (\ref{IdentityMatrix:for dimension and matrix object}), and \texttt{Matrix} (\ref{Matrix:for filter, base domain, list, ncols}) admit different kinds of inputs. \texttt{ZeroMatrix( R, m, n )} can be used, for example, if we want to leave the decision about the
representation of the result to \textsf{GAP}. This might be useful if the result of \texttt{ZeroMatrix( R, m, n )} is the first matrix we create, and then later matrices are created relative to
this matrix. If we already have a matrix object \texttt{M} and want \texttt{ZeroMatrix} to return a new matrix object that has the same representation as \texttt{M}, then \texttt{ZeroMatrix( R, m, n )} may not return a matrix with the expected representation. For example, even if
we use the same \texttt{R} as before, \textsf{GAP} might decide to use a sparse representation for large enough \texttt{m} and \texttt{n}. It is possible to create a \texttt{ZeroMatrix} in the same representation as \texttt{M} by fully specifying this representation as follows \texttt{ZeroMatrix( ConstructingFilter( M ), BaseDomain( M ), m, n )}. It is also possible to produce the same result using the more convenient \texttt{ZeroMatrix( m, n, M )}. 

 This approach works also in many situation where the input involves the ``list of list'' matrices in the filter \texttt{IsMatrix} (\ref{IsMatrix}). In practice, the question is often the other way round: one has old \textsf{GAP} code that was written for objects in \texttt{IsMatrix} (\ref{IsMatrix}), and wants to rewrite it such that it works for general matrix objects also.
In such cases, the following guidelines may be useful. 
\begin{itemize}
\item  Use \texttt{M[i, j]} not \texttt{M[i][j]} for accessing/assigning matrix entries. 

 \emph{Reason:} \texttt{M[i][j]} means to fetch or even create \texttt{M[i]} and then take the \texttt{j}\texttt{\symbol{45}}th entry of it. 
\item  Use \texttt{ExtractSubMatrix( M, rows, cols )} not \texttt{M\texttt{\symbol{123}} rows \texttt{\symbol{125}}\texttt{\symbol{123}} cols
\texttt{\symbol{125}}} for accessing submatrices, similarly use \texttt{CopySubMatrix( src, dst, srows, drows, scols, dcols )}. 

 \emph{Reason:} The \texttt{M\texttt{\symbol{123}} rows \texttt{\symbol{125}}} syntax is supported only for row\texttt{\symbol{45}}list matrices, see \texttt{IsRowListMatrix} (\ref{IsRowListMatrix}). 
\item  Use \texttt{ZeroVector( R, n )} not \texttt{[ 1 .. n ] * Zero( R )} for creating a zero vector over a given domain \texttt{R}. 

 \emph{Reason:} The latter syntax creates an unnecessary new object \texttt{[ 1 .. n ]}, multiplies each of its entries with \texttt{Zero( R )} (not knowing that the result will be \texttt{Zero( R )} in each case), and creates not a vector object but a plain list of the
results. 
\item  Use \texttt{ZeroVector( n, M )} not \texttt{0 * M[1]} for a given matrix object \texttt{M} with \texttt{n} columns. 

 \emph{Reason:} Calling \texttt{M[1]} may have to create an unnecessary new object \texttt{M[1]}. 
\item  Use \texttt{BaseDomain} (\ref{BaseDomain:for a matrix object}) not \texttt{DefaultFieldOfMatrix} (\ref{DefaultFieldOfMatrix}). 

 \emph{Reason:} This is just a conceptual issue, the results should in fact coincide. A matrix
object \texttt{M} which is not a list of lists stores its \texttt{BaseDomain} (\ref{BaseDomain:for a matrix object}) value. For convenience, \texttt{DefaultFieldOfMatrix( M )} is defined to return this value (although this need not be a field). On the
other hand, a list of lists \texttt{M} in \texttt{IsMatrix} (\ref{IsMatrix}) does not store the two attribute values, \texttt{DefaultFieldOfMatrix} (\ref{DefaultFieldOfMatrix}) value is computed from the entries of \texttt{M}, and for convenience, \texttt{BaseDomain( M )} is defined to return the \texttt{DefaultFieldOfMatrix} (\ref{DefaultFieldOfMatrix}) value. 
\item  Do not use the unary versions of \texttt{ConvertToVectorRep} (\ref{ConvertToVectorRep:for a list (and a field)}) and \texttt{ConvertToMatrixRep} (\ref{ConvertToMatrixRep:for a list (and a field)}). 

 \emph{Reason:} For example, consider the two matrices \texttt{M1 = [ [ Z(4) ] ]} and \texttt{M2 = [ [ Z(4)\texttt{\symbol{94}}3 ] ]} over the field with four elements. The latter is in fact a matrix over the
field with two elements, and \texttt{ConvertToMatrixRep( M2 )} turns it into a matrix in \texttt{IsGF2MatrixRep} (\ref{IsGF2MatrixRep}), whereas \texttt{ConvertToMatrixRep( M1 )} yields a matrix in \texttt{Is8BitMatrixRep} (\ref{Is8BitMatrixRep}). Thus computing products or sums of these matrices is more expensive than
computations with two matrices in \texttt{Is8BitMatrixRep} (\ref{Is8BitMatrixRep}). 
\item  Use \texttt{ImmutableMatrix} (\ref{ImmutableMatrix}) and the binary versions of \texttt{ConvertToVectorRep} (\ref{ConvertToVectorRep:for a list (and a field)}) and \texttt{ConvertToMatrixRep} (\ref{ConvertToMatrixRep:for a list (and a field)}) only when creating initial objects which need not be compatible with given
vectors or matrices. 

 \emph{Reason:} The idea behind these functions is to choose a good representation for initial
data (for example for some expensive MeatAxe computations, see Chapter \ref{The MeatAxe}). The results of computations with these data should then automatically be in
the same representation. If not then the code in question has problems, and ``adjusting'' the representation of intermediate results by calling the abovementioned
conversion functions just hides these problems. 
\end{itemize}
 Conversely, do \emph{not} use functions for vector and matrix objects when you want to create an object
that shall be used just as a list. For example, use \texttt{ListWithIdenticalEntries( n, Zero( R ) )} not \texttt{Vector( R, n )} in this case. For creating a list with \texttt{n} entries, you can also first call \texttt{EmptyPlist} (\ref{EmptyPlist}) for creating a big enough list, and then enter the values. }

 }

          
\chapter{\textcolor{Chapter }{Strings and Characters}}\label{Strings and Characters}
\logpage{[ 27, 0, 0 ]}
\hyperdef{L}{X7D28329B7EDB8F47}{}
{
   
\section{\textcolor{Chapter }{IsChar and IsString}}\label{sect:IsChar}
\logpage{[ 27, 1, 0 ]}
\hyperdef{L}{X7A90690B78260194}{}
{
  \index{type!strings} \index{doublequotes} \index{singlequotes} 

\subsection{\textcolor{Chapter }{IsChar}}
\logpage{[ 27, 1, 1 ]}\nobreak
\hyperdef{L}{X80CFAE128560E064}{}
{\noindent\textcolor{FuncColor}{$\triangleright$\enspace\texttt{IsChar({\mdseries\slshape obj})\index{IsChar@\texttt{IsChar}}
\label{IsChar}
}\hfill{\scriptsize (Category)}}\\
\noindent\textcolor{FuncColor}{$\triangleright$\enspace\texttt{IsCharCollection({\mdseries\slshape obj})\index{IsCharCollection@\texttt{IsCharCollection}}
\label{IsCharCollection}
}\hfill{\scriptsize (Category)}}\\


 A \emph{character} is simply an object in \textsf{GAP} that represents an arbitrary character from the character set of the operating
system. Character literals can be entered in \textsf{GAP} by enclosing the character in \emph{singlequotes} \texttt{'}. 
\begin{Verbatim}[commandchars=!@|,fontsize=\small,frame=single,label=Example]
  !gapprompt@gap>| !gapinput@x:= 'a';  IsChar( x );|
  'a'
  true
  !gapprompt@gap>| !gapinput@'*';|
  '*'
\end{Verbatim}
 }

 

\subsection{\textcolor{Chapter }{IsString}}
\logpage{[ 27, 1, 2 ]}\nobreak
\hyperdef{L}{X78723B5D795A3B6D}{}
{\noindent\textcolor{FuncColor}{$\triangleright$\enspace\texttt{IsString({\mdseries\slshape obj})\index{IsString@\texttt{IsString}}
\label{IsString}
}\hfill{\scriptsize (filter)}}\\


 A \emph{string} is a dense list (see{\nobreakspace}\texttt{IsList} (\ref{IsList}), \texttt{IsDenseList} (\ref{IsDenseList})) of characters (see{\nobreakspace}\texttt{IsChar} (\ref{IsChar})); thus strings are always homogeneous (see{\nobreakspace}\texttt{IsHomogeneousList} (\ref{IsHomogeneousList})). 

 A string literal can either be entered as the list of characters or by writing
the characters between \emph{doublequotes} \texttt{"}. \textsf{GAP} will always output strings in the latter format. However, the input via the
double quote syntax enables \textsf{GAP} to store the string in an efficient compact internal representation. See \texttt{IsStringRep} (\ref{IsStringRep}) below for more details. 

 Each character, in particular those which cannot be typed directly from the
keyboard, can also be typed in three digit octal notation, or two digit
hexadecimal notation. And for some special characters (like the newline
character) there is a further possibility to type them, see section \ref{Special Characters}. 

 
\begin{Verbatim}[commandchars=!@|,fontsize=\small,frame=single,label=Example]
  !gapprompt@gap>| !gapinput@s1 := ['H','e','l','l','o',' ','w','o','r','l','d','.'];|
  "Hello world."
  !gapprompt@gap>| !gapinput@IsString( s1 );|
  true
  !gapprompt@gap>| !gapinput@s2 := "Hello world.";|
  "Hello world."
  !gapprompt@gap>| !gapinput@s1 = s2;|
  true
  !gapprompt@gap>| !gapinput@s3 := "";  # the empty string|
  ""
  !gapprompt@gap>| !gapinput@s3 = [];|
  true
  !gapprompt@gap>| !gapinput@IsString( [] );|
  true
  !gapprompt@gap>| !gapinput@IsString( "123" );  IsString( 123 );|
  true
  false
  !gapprompt@gap>| !gapinput@IsString( [ '1', '2', '3' ] );|
  true
  !gapprompt@gap>| !gapinput@IsString( [ '1', '2', , '4' ] );  # strings must be dense|
  false
  !gapprompt@gap>| !gapinput@IsString( [ '1', '2', 3 ] );  # strings must only contain characters|
  false
\end{Verbatim}
 }

 
\subsection{\textcolor{Chapter }{Strings As Lists}}\label{subsect:strings as lists}
\logpage{[ 27, 1, 3 ]}
\hyperdef{L}{X7B1B45C587A72F96}{}
{
  Note that a string is just a special case of a list. So everything that is
possible for lists (see{\nobreakspace}\ref{Lists}) is also possible for strings. Thus you can access the characters in such a
string (see{\nobreakspace}\ref{List Elements}), test for membership (see{\nobreakspace}\ref{Membership Test for Collections}), ask for the length, concatenate strings (see{\nobreakspace}\texttt{Concatenation} (\ref{Concatenation:for several lists})), form substrings etc. You can even assign to a mutable string
(see{\nobreakspace}\ref{List Assignment}). Of course unless you assign a character in such a way that the list stays
dense, the resulting list will no longer be a string. 

 
\begin{Verbatim}[commandchars=!@|,fontsize=\small,frame=single,label=Example]
  !gapprompt@gap>| !gapinput@Length( s2 );|
  12
  !gapprompt@gap>| !gapinput@s2[2];|
  'e'
  !gapprompt@gap>| !gapinput@'a' in s2;|
  false
  !gapprompt@gap>| !gapinput@s2[2] := 'a';;  s2;|
  "Hallo world."
  !gapprompt@gap>| !gapinput@s1{ [1..4] };|
  "Hell"
  !gapprompt@gap>| !gapinput@Concatenation( s1{ [ 1 .. 6 ] }, s1{ [ 1 .. 4 ] } );|
  "Hello Hell"
\end{Verbatim}
 }

 
\subsection{\textcolor{Chapter }{Printing Strings}}\logpage{[ 27, 1, 4 ]}
\hyperdef{L}{X7EA6CA7486D7E9DD}{}
{
\noindent\textcolor{FuncColor}{$\triangleright$\enspace\texttt{ViewObj({\mdseries\slshape str})\index{ViewObj@\texttt{ViewObj}!for a string}
\label{ViewObj:for a string}
}\hfill{\scriptsize (method)}}\\
\noindent\textcolor{FuncColor}{$\triangleright$\enspace\texttt{PrintObj({\mdseries\slshape str})\index{PrintObj@\texttt{PrintObj}!for a string}
\label{PrintObj:for a string}
}\hfill{\scriptsize (method)}}\\


 If a string is displayed by \texttt{View} (\ref{View}), for example as result of an evaluation (see{\nobreakspace}\ref{Main Loop}), or by \texttt{ViewObj} (\ref{ViewObj}) and \texttt{PrintObj} (\ref{PrintObj}), it is displayed with enclosing doublequotes. (But note that there is an
ambiguity for the empty string which is also an empty list of arbitrary \textsf{GAP} objects; it is only printed like a string if it was input as empty string or
converted to a string with \texttt{ConvertToStringRep} (\ref{ConvertToStringRep}).)  The output of \texttt{PrintObj} can be read back into \textsf{GAP}. 

 Strings behave differently from other \textsf{GAP} objects with respect to \texttt{Print} (\ref{Print}), \texttt{PrintTo} (\ref{PrintTo}), or \texttt{AppendTo} (\ref{AppendTo}). These commands \emph{interpret} a string in the sense that they essentially send the characters of the string
directly to the output stream/file. (But depending on the type of the stream
and the presence of some special characters used as hints for line breaks
there may be sent some additional newline (or backslash and newline)
characters.  

 
\begin{Verbatim}[commandchars=!@|,fontsize=\small,frame=single,label=Example]
  !gapprompt@gap>| !gapinput@s4:= "abc\"def\nghi";;|
  !gapprompt@gap>| !gapinput@View( s4 );  Print( "\n" );|
  "abc\"def\nghi"
  !gapprompt@gap>| !gapinput@ViewObj( s4 );  Print( "\n" );|
  "abc\"def\nghi"
  !gapprompt@gap>| !gapinput@PrintObj( s4 );  Print( "\n" );|
  "abc\"def\nghi"
  !gapprompt@gap>| !gapinput@Print( s4 );  Print( "\n" );|
  abc"def
  ghi
  !gapprompt@gap>| !gapinput@s := "German uses strange characters: ����\n";|
  "German uses strange characters: ����\n"
  !gapprompt@gap>| !gapinput@Print(s);|
  German uses strange characters: ����
  !gapprompt@gap>| !gapinput@PrintObj(s);  Print( "\n" );|
  "German uses strange characters: \303\244\303\266\303\274\303\237\n"
\end{Verbatim}
 

 
\begin{Verbatim}[commandchars=!@|,fontsize=\small,frame=single,label=Example]
  !gapprompt@gap>| !gapinput@s := "\007";|
  "\007"
  !gapprompt@gap>| !gapinput@Print(s); # rings bell in many terminals|
\end{Verbatim}
 

 Note that only those line breaks are printed by \texttt{Print} (\ref{Print}) that are contained in the string (\texttt{\texttt{\symbol{92}}n} characters, see{\nobreakspace}\ref{Special Characters}), as is shown in the example below. 

 
\begin{Verbatim}[commandchars=!@|,fontsize=\small,frame=single,label=Example]
  !gapprompt@gap>| !gapinput@s1;|
  "Hello world."
  !gapprompt@gap>| !gapinput@Print( s1 );|
  Hello world.gap> Print( s1, "\n" );
  Hello world.
  !gapprompt@gap>| !gapinput@Print( s1, "\nnext line\n" );|
  Hello world.
  next line
\end{Verbatim}
 }

 }

  
\section{\textcolor{Chapter }{Special Characters}}\label{Special Characters}
\logpage{[ 27, 2, 0 ]}
\hyperdef{L}{X82E5F5AB818F32DB}{}
{
  \index{escaped characters} \index{special character sequences} There are a number of \emph{special character sequences} that can be used between the singlequotes of a character literal or between
the doublequotes of a string literal to specify characters. They consist of a
backslash \texttt{\texttt{\symbol{92}}} followed by a second character indicating the type of special character
sequence, and possibly more characters. The following special character
sequences are currently defined. For any other sequence starting with a
backslash, the backslash is ignored. 

 
\begin{description}
\item[{ \index{\\n@\texttt{\texttt{\symbol{92}}n}} \index{newline character} \texttt{\texttt{\symbol{92}}n}}]  \emph{newline character}. This is the character that, at least on UNIX systems, separates lines in a
text file. Printing of this character in a string has the effect of moving the
cursor down one line and back to the beginning of the line. 
\item[{ \index{\""@\texttt{\texttt{\symbol{92}}""}} \index{doublequote character} \texttt{\texttt{\symbol{92}}"}}]  \emph{doublequote character}. Inside a string a doublequote must be escaped by the backslash, because it
is otherwise interpreted as end of the string. 
\item[{ \index{\'@\texttt{\texttt{\symbol{92}}'}} \index{singlequote character} \texttt{\texttt{\symbol{92}}'}}]  \emph{singlequote character}. Inside a character a singlequote must escaped by the backslash, because it
is otherwise interpreted as end of the character. 
\item[{ \index{\\@\texttt{\texttt{\symbol{92}}\texttt{\symbol{92}}}} \index{backslash character} \texttt{\texttt{\symbol{92}}\texttt{\symbol{92}}}}]  \emph{backslash character}. Inside a string a backslash must be escaped by another backslash, because it
is otherwise interpreted as first character of an escape sequence. 
\item[{ \index{\b@\texttt{\texttt{\symbol{92}}b}} \index{backspace character} \texttt{\texttt{\symbol{92}}b}}]  \emph{backspace character}. Printing this character should have the effect of moving the cursor back one
character. Whether it works or not is system dependent and should not be
relied upon. 
\item[{ \index{\r@\texttt{\texttt{\symbol{92}}r}} \index{carriage return character} \texttt{\texttt{\symbol{92}}r}}]  \emph{carriage return character}. Printing this character should have the effect of moving the cursor back to
the beginning of the same line. Whether this works or not is again system
dependent. 
\item[{ \index{\c@\texttt{\texttt{\symbol{92}}c}} \index{flush character} \texttt{\texttt{\symbol{92}}c}}]  \emph{flush character}. This character is not printed. Its purpose is to flush the output queue.
Usually \textsf{GAP} waits until it sees a \texttt{newline} before it prints a string. If you want to display a string that does not
include this character use \texttt{\texttt{\symbol{92}}c}. 
\item[{ \index{\\XYZ@\texttt{\texttt{\symbol{92}}XYZ}} \index{octal character codes} \texttt{\texttt{\symbol{92}}XYZ}}]  with \texttt{X}, \texttt{Y}, \texttt{Z} three octal digits, that is one of \texttt{"01234567"}. This is translated to the character corresponding to the number \texttt{X * 64 + Y * 8 + Z modulo 256}. This can be used to specify and store arbitrary binary data as a string in \textsf{GAP}. 
\item[{ \index{\\0xYZ@\texttt{\texttt{\symbol{92}}0xYZ}} \index{hexadecimal character codes} \texttt{\texttt{\symbol{92}}0xYZ}}]  with \texttt{Y}, and \texttt{Z} hexadecimal digits, that is one of \texttt{"0123456789ABCDEFabcdef"}, where \texttt{a} to \texttt{f} and \texttt{A} to \texttt{F} are interpreted as the numbers \texttt{10} to \texttt{15}. This is translated to the character corresponding to the number \texttt{Y*16 + Z}. 
\item[{ \index{escaping non\texttt{\symbol{45}}special characters} other}]  For any other character the backslash is ignored.  
\end{description}
 

 Again, if the line is displayed as result of an evaluation, those escape
sequences are displayed in the same way that they are input. 

 Only \texttt{Print} (\ref{Print}), \texttt{PrintTo} (\ref{PrintTo}), or \texttt{AppendTo} (\ref{AppendTo}) send the characters directly to the output stream. 

 
\begin{Verbatim}[commandchars=!@|,fontsize=\small,frame=single,label=Example]
  !gapprompt@gap>| !gapinput@"This is one line.\nThis is another line.\n";|
  "This is one line.\nThis is another line.\n"
  !gapprompt@gap>| !gapinput@Print( last );|
  This is one line.
  This is another line.
\end{Verbatim}
 

 Note in particular that it is not allowed to enclose a \mbox{\texttt{\mdseries\slshape newline}} inside the string. You can use the special character sequence \texttt{\texttt{\symbol{92}}n} to write strings that include \mbox{\texttt{\mdseries\slshape newline}} characters. If, however, an input string is too long to fit on a single line
it is possible to \emph{continue} it over several lines. In this case the last character of each input line,
except the last line must be a backslash. Both backslash and \mbox{\texttt{\mdseries\slshape newline}} are thrown away by \textsf{GAP} while reading the string. Note that the same continuation mechanism is
available for identifiers and integers, see{\nobreakspace}\ref{Special Rules for Input Lines}. The rules on escaping are ignored in a triple quoted string, see \ref{Triple Quoted String} }

 
\section{\textcolor{Chapter }{Triple Quoted Strings}}\label{Triple Quoted String}
\logpage{[ 27, 3, 0 ]}
\hyperdef{L}{X7E70384E7D0B7083}{}
{
  Another method of entering strings in GAP is triple quoted strings. Triple
quoted strings ignore the rules on escaping given in \ref{Special Characters}. Triple quoted strings begin an end with three doublequotes. Inside the
triple quotes no escaping is done, and the string continues, including
newlines, until three doublequotes are found. 

 
\begin{Verbatim}[commandchars=!@|,fontsize=\small,frame=single,label=Example]
  !gapprompt@gap>| !gapinput@"""Print("\n")""";|
  "Print(\"\\n\")"
\end{Verbatim}
 Triple quoted strings are represented internally identically to all other
strings, they only provide an alternative method of giving strings to GAP.
Triple quoted strings still follow GAP's line editing rules (\ref{Special Rules for Input Lines}), which state that in normal line editing mode, lines starting \texttt{gap{\textgreater} }, \texttt{{\textgreater} } or \texttt{brk{\textgreater} } will have this beginning part removed. }

  
\section{\textcolor{Chapter }{Internally Represented Strings}}\label{Internally Represented Strings}
\logpage{[ 27, 4, 0 ]}
\hyperdef{L}{X82AEC07487C45ECD}{}
{
  \index{convert!to a string} 

\subsection{\textcolor{Chapter }{IsStringRep}}
\logpage{[ 27, 4, 1 ]}\nobreak
\hyperdef{L}{X7A17EDF8785C9F58}{}
{\noindent\textcolor{FuncColor}{$\triangleright$\enspace\texttt{IsStringRep({\mdseries\slshape obj})\index{IsStringRep@\texttt{IsStringRep}}
\label{IsStringRep}
}\hfill{\scriptsize (Representation)}}\\


 \texttt{IsStringRep} is a special (internal) representation of dense lists of characters. Dense
lists of characters can be converted into this representation using \texttt{ConvertToStringRep} (\ref{ConvertToStringRep}). Note that calling \texttt{IsString} (\ref{IsString}) does \emph{not} change the representation. }

 

\subsection{\textcolor{Chapter }{ConvertToStringRep}}
\logpage{[ 27, 4, 2 ]}\nobreak
\hyperdef{L}{X7CE2415F7FEC5809}{}
{\noindent\textcolor{FuncColor}{$\triangleright$\enspace\texttt{ConvertToStringRep({\mdseries\slshape obj})\index{ConvertToStringRep@\texttt{ConvertToStringRep}}
\label{ConvertToStringRep}
}\hfill{\scriptsize (function)}}\\


 If \mbox{\texttt{\mdseries\slshape obj}} is a dense internally represented list of characters then \texttt{ConvertToStringRep} changes the representation to \texttt{IsStringRep} (\ref{IsStringRep}). This is useful in particular for converting the empty list \texttt{[]}, which usually is in \texttt{IsPlistRep} (\ref{IsPlistRep}), to \texttt{IsStringRep} (\ref{IsStringRep}). If \mbox{\texttt{\mdseries\slshape obj}} is not a string then \texttt{ConvertToStringRep} signals an error. }

 

\subsection{\textcolor{Chapter }{CopyToStringRep}}
\logpage{[ 27, 4, 3 ]}\nobreak
\hyperdef{L}{X7FFC464683CC8023}{}
{\noindent\textcolor{FuncColor}{$\triangleright$\enspace\texttt{CopyToStringRep({\mdseries\slshape obj})\index{CopyToStringRep@\texttt{CopyToStringRep}}
\label{CopyToStringRep}
}\hfill{\scriptsize (function)}}\\


 If \mbox{\texttt{\mdseries\slshape obj}} is a dense internally represented list of characters then \texttt{CopyToStringRep} copies \mbox{\texttt{\mdseries\slshape obj}} to a new object with representation \texttt{IsStringRep} (\ref{IsStringRep}). If \mbox{\texttt{\mdseries\slshape obj}} is not a string then \texttt{CopyToStringRep} signals an error. }

 

\subsection{\textcolor{Chapter }{IsEmptyString}}
\logpage{[ 27, 4, 4 ]}\nobreak
\hyperdef{L}{X7D944D507CBB24CD}{}
{\noindent\textcolor{FuncColor}{$\triangleright$\enspace\texttt{IsEmptyString({\mdseries\slshape str})\index{IsEmptyString@\texttt{IsEmptyString}}
\label{IsEmptyString}
}\hfill{\scriptsize (function)}}\\


 \texttt{IsEmptyString} returns \texttt{true} if \mbox{\texttt{\mdseries\slshape str}} is the empty string in the representation \texttt{IsStringRep} (\ref{IsStringRep}), and \texttt{false} otherwise. Note that the empty list \texttt{[]} and the empty string \texttt{""} have the same type, the recommended way to distinguish them is via \texttt{IsEmptyString}. For formatted printing, this distinction is sometimes necessary.     

 
\begin{Verbatim}[commandchars=!@|,fontsize=\small,frame=single,label=Example]
  !gapprompt@gap>| !gapinput@l:= [];;  IsString( l );  IsEmptyString( l );  IsEmpty( l );|
  true
  false
  true
  !gapprompt@gap>| !gapinput@l;  ConvertToStringRep( l );  l;|
  [  ]
  ""
  !gapprompt@gap>| !gapinput@IsEmptyString( l );  IsEmptyString( "" );  IsEmptyString( "abc" );|
  true
  true
  false
  !gapprompt@gap>| !gapinput@ll:= [ 'a', 'b' ];  IsStringRep( ll );  ConvertToStringRep( ll );|
  "ab"
  false
  !gapprompt@gap>| !gapinput@ll;  IsStringRep( ll );|
  "ab"
  true
\end{Verbatim}
 }

 

\subsection{\textcolor{Chapter }{EmptyString}}
\logpage{[ 27, 4, 5 ]}\nobreak
\hyperdef{L}{X836078DC829A8221}{}
{\noindent\textcolor{FuncColor}{$\triangleright$\enspace\texttt{EmptyString({\mdseries\slshape len})\index{EmptyString@\texttt{EmptyString}}
\label{EmptyString}
}\hfill{\scriptsize (function)}}\\
\textbf{\indent Returns:\ }
a string

\noindent\textcolor{FuncColor}{$\triangleright$\enspace\texttt{ShrinkAllocationString({\mdseries\slshape str})\index{ShrinkAllocationString@\texttt{ShrinkAllocationString}}
\label{ShrinkAllocationString}
}\hfill{\scriptsize (function)}}\\
\textbf{\indent Returns:\ }
nothing



 The function \texttt{EmptyString} returns an empty string in internal representation which has enough memory
allocated for \mbox{\texttt{\mdseries\slshape len}} characters. This can be useful for creating and filling a string with a known
number of entries. 

 The function \texttt{ShrinkAllocationString} gives back to \textsf{GAP}s memory manager the physical memory which is allocated for the string \mbox{\texttt{\mdseries\slshape str}} in internal representation but not needed by its current number of characters. 

 These functions are intended for saving some of \textsf{GAP}s memory in certain situations, see the explanations and the example for the
analogous functions \texttt{EmptyPlist} (\ref{EmptyPlist}) and \texttt{ShrinkAllocationPlist} (\ref{ShrinkAllocationPlist}) for plain lists. }

 

\subsection{\textcolor{Chapter }{CharsFamily}}
\logpage{[ 27, 4, 6 ]}\nobreak
\hyperdef{L}{X7DA671FC7F490C16}{}
{\noindent\textcolor{FuncColor}{$\triangleright$\enspace\texttt{CharsFamily\index{CharsFamily@\texttt{CharsFamily}}
\label{CharsFamily}
}\hfill{\scriptsize (family)}}\\


 Each character lies in the family \texttt{CharsFamily}, each nonempty string lies in the collections family of this family. Note the
subtle differences between the empty list \texttt{[]} and the empty string \texttt{""} when both are printed. }

 }

  
\section{\textcolor{Chapter }{Recognizing Characters}}\label{Recognizing Characters}
\logpage{[ 27, 5, 0 ]}
\hyperdef{L}{X82F980A17FE84AA4}{}
{
  

\subsection{\textcolor{Chapter }{IsDigitChar}}
\logpage{[ 27, 5, 1 ]}\nobreak
\hyperdef{L}{X78566FD57B95ECBE}{}
{\noindent\textcolor{FuncColor}{$\triangleright$\enspace\texttt{IsDigitChar({\mdseries\slshape c})\index{IsDigitChar@\texttt{IsDigitChar}}
\label{IsDigitChar}
}\hfill{\scriptsize (function)}}\\


 checks whether the character \mbox{\texttt{\mdseries\slshape c}} is a digit, i.e., occurs in the string \texttt{"0123456789"}. }

 

\subsection{\textcolor{Chapter }{IsLowerAlphaChar}}
\logpage{[ 27, 5, 2 ]}\nobreak
\hyperdef{L}{X854114A97BAFEAEA}{}
{\noindent\textcolor{FuncColor}{$\triangleright$\enspace\texttt{IsLowerAlphaChar({\mdseries\slshape c})\index{IsLowerAlphaChar@\texttt{IsLowerAlphaChar}}
\label{IsLowerAlphaChar}
}\hfill{\scriptsize (function)}}\\


 checks whether the character \mbox{\texttt{\mdseries\slshape c}} is a lowercase alphabet letter, i.e., occurs in the string \texttt{"abcdefghijklmnopqrstuvwxyz"}. }

 

\subsection{\textcolor{Chapter }{IsUpperAlphaChar}}
\logpage{[ 27, 5, 3 ]}\nobreak
\hyperdef{L}{X87B1A13D81353AD8}{}
{\noindent\textcolor{FuncColor}{$\triangleright$\enspace\texttt{IsUpperAlphaChar({\mdseries\slshape c})\index{IsUpperAlphaChar@\texttt{IsUpperAlphaChar}}
\label{IsUpperAlphaChar}
}\hfill{\scriptsize (function)}}\\


 checks whether the character \mbox{\texttt{\mdseries\slshape c}} is an uppercase alphabet letter, i.e., occurs in the string \texttt{"ABCDEFGHIJKLMNOPQRSTUVWXYZ"}. }

 

\subsection{\textcolor{Chapter }{IsAlphaChar}}
\logpage{[ 27, 5, 4 ]}\nobreak
\hyperdef{L}{X84634DF67A431D26}{}
{\noindent\textcolor{FuncColor}{$\triangleright$\enspace\texttt{IsAlphaChar({\mdseries\slshape c})\index{IsAlphaChar@\texttt{IsAlphaChar}}
\label{IsAlphaChar}
}\hfill{\scriptsize (function)}}\\


 checks whether the character \mbox{\texttt{\mdseries\slshape c}} is either a lowercase or an uppercase alphabet letter. }

 }

  
\section{\textcolor{Chapter }{Comparisons of Strings}}\label{Comparisons of Strings}
\logpage{[ 27, 6, 0 ]}
\hyperdef{L}{X8127954B79B8A0DA}{}
{
  

\subsection{\textcolor{Chapter }{\texttt{\symbol{92}}= (for two strings)}}
\logpage{[ 27, 6, 1 ]}\nobreak
\hyperdef{L}{X79538F138286739A}{}
{\noindent\textcolor{FuncColor}{$\triangleright$\enspace\texttt{\texttt{\symbol{92}}=({\mdseries\slshape string1, string2})\index{\texttt{\symbol{92}}=@\texttt{\texttt{\symbol{92}}=}!for two strings}
\label{bSlash=:for two strings}
}\hfill{\scriptsize (method)}}\\


 \index{strings!equality of} \index{strings!inequality of} The equality operator \texttt{=} returns \texttt{true} if the two strings \mbox{\texttt{\mdseries\slshape string1}} and \mbox{\texttt{\mdseries\slshape string2}} are equal and \texttt{false} otherwise. The inequality operator \texttt{{\textless}{\textgreater}} returns \texttt{true} if the two strings \mbox{\texttt{\mdseries\slshape string1}} and \mbox{\texttt{\mdseries\slshape string2}} are not equal and \texttt{false} otherwise. 

 
\begin{Verbatim}[commandchars=!@|,fontsize=\small,frame=single,label=Example]
  !gapprompt@gap>| !gapinput@"Hello world.\n" = "Hello world.\n";|
  true
  !gapprompt@gap>| !gapinput@"Hello World.\n" = "Hello world.\n"; # comparison is case sensitive|
  false
  !gapprompt@gap>| !gapinput@"Hello world." = "Hello world.\n";  # first string has no <newline>|
  false
  !gapprompt@gap>| !gapinput@"Goodbye world.\n" = "Hello world.\n";|
  false
  !gapprompt@gap>| !gapinput@[ 'a', 'b' ] = "ab";|
  true
\end{Verbatim}
 }

 

\subsection{\textcolor{Chapter }{\texttt{\symbol{92}}{\textless} (for two strings)}}
\logpage{[ 27, 6, 2 ]}\nobreak
\hyperdef{L}{X8129E3A785F60093}{}
{\noindent\textcolor{FuncColor}{$\triangleright$\enspace\texttt{\texttt{\symbol{92}}{\textless}({\mdseries\slshape string1, string2})\index{\texttt{\symbol{92}}{\textless}@\texttt{\texttt{\symbol{92}}{\textless}}!for two strings}
\label{bSlash<:for two strings}
}\hfill{\scriptsize (method)}}\\


 \index{strings!lexicographic ordering of} The ordering of strings is lexicographically according to the order implied by
the underlying, system dependent, character set. 

 
\begin{Verbatim}[commandchars=!@|,fontsize=\small,frame=single,label=Example]
  !gapprompt@gap>| !gapinput@"Hello world.\n" < "Hello world.\n";  # the strings are equal|
  false
  !gapprompt@gap>| !gapinput@# in ASCII capitals range before small letters:|
  !gapprompt@gap>| !gapinput@"Hello World." < "Hello world.";|
  true
  !gapprompt@gap>| !gapinput@"Hello world." < "Hello world.\n";  # prefixes are always smaller|
  true
  !gapprompt@gap>| !gapinput@# G comes before H, in ASCII at least:|
  !gapprompt@gap>| !gapinput@"Goodbye world.\n" < "Hello world.\n";|
  true
\end{Verbatim}
 

 Strings can be compared via \texttt{{\textless}} with certain \textsf{GAP} objects that are not strings, see{\nobreakspace}\ref{Comparisons} for the details. }

 }

  
\section{\textcolor{Chapter }{Operations to Produce or Manipulate Strings}}\label{Operations to Produce or Manipulate Strings}
\logpage{[ 27, 7, 0 ]}
\hyperdef{L}{X7E72717A82A309F5}{}
{
  For the possibility to print \textsf{GAP} objects to strings, see{\nobreakspace}\ref{String Streams}. 

\subsection{\textcolor{Chapter }{DisplayString}}
\logpage{[ 27, 7, 1 ]}\nobreak
\hyperdef{L}{X792FB3A1849FD739}{}
{\noindent\textcolor{FuncColor}{$\triangleright$\enspace\texttt{DisplayString({\mdseries\slshape obj})\index{DisplayString@\texttt{DisplayString}}
\label{DisplayString}
}\hfill{\scriptsize (operation)}}\\


 Returns a string which could be used to display the object \mbox{\texttt{\mdseries\slshape obj}} in a nice, formatted way which is easy to read (but might be difficult for
machines to understand). The actual format used for this depends on the type
of \mbox{\texttt{\mdseries\slshape obj}}. Each method should include a newline character as last character. Note that
no method for \texttt{DisplayString} may delegate to any of the operations \texttt{Display} (\ref{Display}), \texttt{ViewObj} (\ref{ViewObj}) or \texttt{PrintObj} (\ref{PrintObj}) to avoid circular delegations. }

 

\subsection{\textcolor{Chapter }{DEFAULTDISPLAYSTRING}}
\logpage{[ 27, 7, 2 ]}\nobreak
\hyperdef{L}{X8482132779EA7A23}{}
{\noindent\textcolor{FuncColor}{$\triangleright$\enspace\texttt{DEFAULTDISPLAYSTRING\index{DEFAULTDISPLAYSTRING@\texttt{DEFAULTDISPLAYSTRING}}
\label{DEFAULTDISPLAYSTRING}
}\hfill{\scriptsize (global variable)}}\\


 This is the default value for \texttt{DisplayString} (\ref{DisplayString}). }

 

\subsection{\textcolor{Chapter }{ViewString}}
\logpage{[ 27, 7, 3 ]}\nobreak
\hyperdef{L}{X7803FBCA79DB5529}{}
{\noindent\textcolor{FuncColor}{$\triangleright$\enspace\texttt{ViewString({\mdseries\slshape obj})\index{ViewString@\texttt{ViewString}}
\label{ViewString}
}\hfill{\scriptsize (operation)}}\\


 \texttt{ViewString} returns a string which would be displayed by \texttt{ViewObj} (\ref{ViewObj}) for an object. Note that no method for \texttt{ViewString} may delegate to any of the operations \texttt{Display} (\ref{Display}), \texttt{ViewObj} (\ref{ViewObj}), \texttt{DisplayString} (\ref{DisplayString}) or \texttt{PrintObj} (\ref{PrintObj}) to avoid circular delegations. }

 

\subsection{\textcolor{Chapter }{DEFAULTVIEWSTRING}}
\logpage{[ 27, 7, 4 ]}\nobreak
\hyperdef{L}{X7BBDF9D383595425}{}
{\noindent\textcolor{FuncColor}{$\triangleright$\enspace\texttt{DEFAULTVIEWSTRING\index{DEFAULTVIEWSTRING@\texttt{DEFAULTVIEWSTRING}}
\label{DEFAULTVIEWSTRING}
}\hfill{\scriptsize (global variable)}}\\


 This is the default value for \texttt{ViewString} (\ref{ViewString}). }

 

\subsection{\textcolor{Chapter }{PrintString}}
\logpage{[ 27, 7, 5 ]}\nobreak
\hyperdef{L}{X7B3CC87285DEC23D}{}
{\noindent\textcolor{FuncColor}{$\triangleright$\enspace\texttt{PrintString({\mdseries\slshape obj[, length]})\index{PrintString@\texttt{PrintString}}
\label{PrintString}
}\hfill{\scriptsize (operation)}}\\


 \texttt{PrintString} returns a representation of \mbox{\texttt{\mdseries\slshape obj}}, which may be an object of arbitrary type, as a string. This string should
approximate as closely as possible the character sequence you see if you print \mbox{\texttt{\mdseries\slshape obj}} using \texttt{PrintObj} (\ref{PrintObj}). 

 If \mbox{\texttt{\mdseries\slshape length}} is given it must be an integer. The absolute value gives the minimal length of
the result. If the string representation of \mbox{\texttt{\mdseries\slshape obj}} takes less than that many characters it is filled with blanks. If \mbox{\texttt{\mdseries\slshape length}} is positive it is filled on the left, if \mbox{\texttt{\mdseries\slshape length}} is negative it is filled on the right. 

 In the two argument case, the string returned is a new mutable string (in
particular not a part of any other object); it can be modified safely, and \texttt{MakeImmutable} (\ref{MakeImmutable}) may be safely applied to it. 
\begin{Verbatim}[commandchars=!@|,fontsize=\small,frame=single,label=Example]
  !gapprompt@gap>| !gapinput@PrintString(123);PrintString([1,2,3]);|
  "123"
  "[ 1, 2, 3 ]"
\end{Verbatim}
 \texttt{PrintString} is entitled to put in additional control characters \texttt{\texttt{\symbol{92}}{\textless}} (ASCII 1) and \texttt{\texttt{\symbol{92}}{\textgreater}} (ASCII 2) that allow proper line breaks. See \texttt{StripLineBreakCharacters} (\ref{StripLineBreakCharacters}) for a function to get rid of these control characters. }

 

\subsection{\textcolor{Chapter }{String}}
\logpage{[ 27, 7, 6 ]}\nobreak
\hyperdef{L}{X81FB5BE27903EC32}{}
{\noindent\textcolor{FuncColor}{$\triangleright$\enspace\texttt{String({\mdseries\slshape obj[, length]})\index{String@\texttt{String}}
\label{String}
}\hfill{\scriptsize (attribute)}}\\


 \texttt{String} returns a representation of \mbox{\texttt{\mdseries\slshape obj}}, which may be an object of arbitrary type, as a string. This string should
approximate as closely as possible the character sequence you see if you print \mbox{\texttt{\mdseries\slshape obj}}. 

 If \mbox{\texttt{\mdseries\slshape length}} is given it must be an integer. The absolute value gives the minimal length of
the result. If the string representation of \mbox{\texttt{\mdseries\slshape obj}} takes less than that many characters it is filled with blanks. If \mbox{\texttt{\mdseries\slshape length}} is positive it is filled on the left, if \mbox{\texttt{\mdseries\slshape length}} is negative it is filled on the right. 

 In the two argument case, the string returned is a new mutable string (in
particular not a part of any other object); it can be modified safely, and \texttt{MakeImmutable} (\ref{MakeImmutable}) may be safely applied to it. 
\begin{Verbatim}[commandchars=!@|,fontsize=\small,frame=single,label=Example]
  !gapprompt@gap>| !gapinput@String(123);String([1,2,3]);|
  "123"
  "[ 1, 2, 3 ]"
\end{Verbatim}
 \texttt{String} must not put in additional control characters \texttt{\texttt{\symbol{92}}{\textless}} (ASCII 1) and \texttt{\texttt{\symbol{92}}{\textgreater}} (ASCII 2) that allow proper line breaks. }

 

\subsection{\textcolor{Chapter }{StripLineBreakCharacters}}
\logpage{[ 27, 7, 7 ]}\nobreak
\hyperdef{L}{X86AACCE987F74FA5}{}
{\noindent\textcolor{FuncColor}{$\triangleright$\enspace\texttt{StripLineBreakCharacters({\mdseries\slshape st})\index{StripLineBreakCharacters@\texttt{StripLineBreakCharacters}}
\label{StripLineBreakCharacters}
}\hfill{\scriptsize (function)}}\\


 This function takes a string \mbox{\texttt{\mdseries\slshape st}} as an argument and removes all control characters \texttt{\texttt{\symbol{92}}{\textless}} (ASCII 1) and \texttt{\texttt{\symbol{92}}{\textgreater}} (ASCII 2) which are used by \texttt{PrintString} (\ref{PrintString}) and \texttt{PrintObj} (\ref{PrintObj}) to ensure proper line breaking. A new string with these characters removed is
returned. }

 

\subsection{\textcolor{Chapter }{HexStringInt}}
\logpage{[ 27, 7, 8 ]}\nobreak
\hyperdef{L}{X865FBB7E788017DD}{}
{\noindent\textcolor{FuncColor}{$\triangleright$\enspace\texttt{HexStringInt({\mdseries\slshape int})\index{HexStringInt@\texttt{HexStringInt}}
\label{HexStringInt}
}\hfill{\scriptsize (function)}}\\


 returns a string which represents the integer \mbox{\texttt{\mdseries\slshape int}} with hexadecimal digits (using \texttt{A} to \texttt{F} as digits \texttt{10} to \texttt{15}). The inverse translation can be achieved with \texttt{IntHexString} (\ref{IntHexString}). }

 

\subsection{\textcolor{Chapter }{StringPP}}
\logpage{[ 27, 7, 9 ]}\nobreak
\hyperdef{L}{X7BB1059185AB4F84}{}
{\noindent\textcolor{FuncColor}{$\triangleright$\enspace\texttt{StringPP({\mdseries\slshape int})\index{StringPP@\texttt{StringPP}}
\label{StringPP}
}\hfill{\scriptsize (function)}}\\


 returns a string representing the prime factor decomposition of the integer \mbox{\texttt{\mdseries\slshape int}}. See also \texttt{PrintFactorsInt} (\ref{PrintFactorsInt}). 
\begin{Verbatim}[commandchars=!@|,fontsize=\small,frame=single,label=Example]
  !gapprompt@gap>| !gapinput@StringPP(40320);|
  "2^7*3^2*5*7"
\end{Verbatim}
 }

 

\subsection{\textcolor{Chapter }{WordAlp}}
\logpage{[ 27, 7, 10 ]}\nobreak
\hyperdef{L}{X79C8280A853D8FA9}{}
{\noindent\textcolor{FuncColor}{$\triangleright$\enspace\texttt{WordAlp({\mdseries\slshape alpha, nr})\index{WordAlp@\texttt{WordAlp}}
\label{WordAlp}
}\hfill{\scriptsize (function)}}\\


 returns a string that is the \mbox{\texttt{\mdseries\slshape nr}}\texttt{\symbol{45}}th word over the alphabet list \mbox{\texttt{\mdseries\slshape alpha}}, w.r.t. word length and lexicographical order. The empty word is \texttt{WordAlp( \mbox{\texttt{\mdseries\slshape alpha}}, 0 )}. 
\begin{Verbatim}[commandchars=!@|,fontsize=\small,frame=single,label=Example]
  !gapprompt@gap>| !gapinput@List([0..5],i->WordAlp("abc",i));|
  [ "", "a", "b", "c", "aa", "ab" ]
\end{Verbatim}
 }

 

\subsection{\textcolor{Chapter }{LowercaseString}}
\logpage{[ 27, 7, 11 ]}\nobreak
\hyperdef{L}{X798A0F35852ABDAD}{}
{\noindent\textcolor{FuncColor}{$\triangleright$\enspace\texttt{LowercaseString({\mdseries\slshape string})\index{LowercaseString@\texttt{LowercaseString}}
\label{LowercaseString}
}\hfill{\scriptsize (function)}}\\


 Returns a lowercase version of the string \mbox{\texttt{\mdseries\slshape string}}, that is, a string in which each uppercase alphabet character is replaced by
the corresponding lowercase character. 
\begin{Verbatim}[commandchars=!@|,fontsize=\small,frame=single,label=Example]
  !gapprompt@gap>| !gapinput@LowercaseString("This Is UpperCase");|
  "this is uppercase"
\end{Verbatim}
 }

 

\subsection{\textcolor{Chapter }{LowercaseChar}}
\logpage{[ 27, 7, 12 ]}\nobreak
\hyperdef{L}{X87A2F2557DE7EE08}{}
{\noindent\textcolor{FuncColor}{$\triangleright$\enspace\texttt{LowercaseChar({\mdseries\slshape character})\index{LowercaseChar@\texttt{LowercaseChar}}
\label{LowercaseChar}
}\hfill{\scriptsize (function)}}\\


 Returns the lowercase version of the character \mbox{\texttt{\mdseries\slshape character}}. }

 

\subsection{\textcolor{Chapter }{UppercaseString}}
\logpage{[ 27, 7, 13 ]}\nobreak
\hyperdef{L}{X7E7E5F5B7FED56A0}{}
{\noindent\textcolor{FuncColor}{$\triangleright$\enspace\texttt{UppercaseString({\mdseries\slshape string})\index{UppercaseString@\texttt{UppercaseString}}
\label{UppercaseString}
}\hfill{\scriptsize (function)}}\\


 Returns a uppercase version of the string \mbox{\texttt{\mdseries\slshape string}}, that is, a string in which each lowercase alphabet character is replaced by
the corresponding uppercase character. 
\begin{Verbatim}[commandchars=!@|,fontsize=\small,frame=single,label=Example]
  !gapprompt@gap>| !gapinput@UppercaseString("This Is UpperCase");|
  "THIS IS UPPERCASE"
\end{Verbatim}
 }

 

\subsection{\textcolor{Chapter }{UppercaseChar}}
\logpage{[ 27, 7, 14 ]}\nobreak
\hyperdef{L}{X81E0AEE687200505}{}
{\noindent\textcolor{FuncColor}{$\triangleright$\enspace\texttt{UppercaseChar({\mdseries\slshape character})\index{UppercaseChar@\texttt{UppercaseChar}}
\label{UppercaseChar}
}\hfill{\scriptsize (function)}}\\


 Returns the uppercase version of the character \mbox{\texttt{\mdseries\slshape character}}. }

 

\subsection{\textcolor{Chapter }{SplitString}}
\logpage{[ 27, 7, 15 ]}\nobreak
\hyperdef{L}{X86E897D486DCFEAB}{}
{\noindent\textcolor{FuncColor}{$\triangleright$\enspace\texttt{SplitString({\mdseries\slshape string, seps[, wspace]})\index{SplitString@\texttt{SplitString}}
\label{SplitString}
}\hfill{\scriptsize (operation)}}\\


 This function accepts a string \mbox{\texttt{\mdseries\slshape string}} and lists \mbox{\texttt{\mdseries\slshape seps}} and, optionally, \mbox{\texttt{\mdseries\slshape wspace}} of characters. Now \mbox{\texttt{\mdseries\slshape string}} is split into substrings at each occurrence of a character in \mbox{\texttt{\mdseries\slshape seps}} or \mbox{\texttt{\mdseries\slshape wspace}}. The characters in \mbox{\texttt{\mdseries\slshape wspace}} are interpreted as white space characters. Substrings of characters in \mbox{\texttt{\mdseries\slshape wspace}} are treated as one white space character and they are ignored at the beginning
and end of a string. 

 Both arguments \mbox{\texttt{\mdseries\slshape seps}} and \mbox{\texttt{\mdseries\slshape wspace}} can be single characters. 

 Each string in the resulting list of substring does not contain any characters
in \mbox{\texttt{\mdseries\slshape seps}} or \mbox{\texttt{\mdseries\slshape wspace}}. 

 A character that occurs both in \mbox{\texttt{\mdseries\slshape seps}} and \mbox{\texttt{\mdseries\slshape wspace}} is treated as a white space character. 

 A separator at the end of a string is interpreted as a terminator; in this
case, the separator does not produce a trailing empty string. Also
see{\nobreakspace}\texttt{Chomp} (\ref{Chomp}). 
\begin{Verbatim}[commandchars=!@|,fontsize=\small,frame=single,label=Example]
  !gapprompt@gap>| !gapinput@SplitString( "substr1:substr2::substr4", ":" );|
  [ "substr1", "substr2", "", "substr4" ]
  !gapprompt@gap>| !gapinput@SplitString( "a;b;c;d;", ";" );|
  [ "a", "b", "c", "d" ]
  !gapprompt@gap>| !gapinput@SplitString( "/home//user//dir/", "", "/" );|
  [ "home", "user", "dir" ]
\end{Verbatim}
 }

 

\subsection{\textcolor{Chapter }{ReplacedString}}
\logpage{[ 27, 7, 16 ]}\nobreak
\hyperdef{L}{X864F0A9078D4DE0E}{}
{\noindent\textcolor{FuncColor}{$\triangleright$\enspace\texttt{ReplacedString({\mdseries\slshape string, old, new})\index{ReplacedString@\texttt{ReplacedString}}
\label{ReplacedString}
}\hfill{\scriptsize (function)}}\\


 replaces occurrences of the string \mbox{\texttt{\mdseries\slshape old}} in \mbox{\texttt{\mdseries\slshape string}} by \mbox{\texttt{\mdseries\slshape new}}, starting from the left and always replacing the first occurrence. To avoid
infinite recursion, characters which have been replaced already, are not
subject to renewed replacement. 
\begin{Verbatim}[commandchars=!@|,fontsize=\small,frame=single,label=Example]
  !gapprompt@gap>| !gapinput@ReplacedString("abacab","a","zl");|
  "zlbzlczlb"
  !gapprompt@gap>| !gapinput@ReplacedString("ababa", "aba","c");|
  "cba"
  !gapprompt@gap>| !gapinput@ReplacedString("abacab","a","ba");|
  "babbacbab"
\end{Verbatim}
 }

 

\subsection{\textcolor{Chapter }{NormalizeWhitespace}}
\logpage{[ 27, 7, 17 ]}\nobreak
\hyperdef{L}{X806379367A53D171}{}
{\noindent\textcolor{FuncColor}{$\triangleright$\enspace\texttt{NormalizeWhitespace({\mdseries\slshape string})\index{NormalizeWhitespace@\texttt{NormalizeWhitespace}}
\label{NormalizeWhitespace}
}\hfill{\scriptsize (function)}}\\


 This function changes the string \mbox{\texttt{\mdseries\slshape string}} in place. The characters \texttt{ } (space), \texttt{\texttt{\symbol{92}}n}, \texttt{\texttt{\symbol{92}}r} and \texttt{\texttt{\symbol{92}}t} are considered as \emph{white space}. Leading and trailing white space characters in \mbox{\texttt{\mdseries\slshape string}} are removed. Sequences of white space characters between other characters are
replaced by a single space character. 

 See \texttt{NormalizedWhitespace} (\ref{NormalizedWhitespace}) for a non\texttt{\symbol{45}}destructive version. 
\begin{Verbatim}[commandchars=!@|,fontsize=\small,frame=single,label=Example]
  !gapprompt@gap>| !gapinput@s := "   x y \n\n\t\r  z\n   \n";|
  "   x y \n\n\t\r  z\n   \n"
  !gapprompt@gap>| !gapinput@NormalizeWhitespace(s);|
  !gapprompt@gap>| !gapinput@s;|
  "x y z"
\end{Verbatim}
 }

 

\subsection{\textcolor{Chapter }{NormalizedWhitespace}}
\logpage{[ 27, 7, 18 ]}\nobreak
\hyperdef{L}{X8685DE9386E57771}{}
{\noindent\textcolor{FuncColor}{$\triangleright$\enspace\texttt{NormalizedWhitespace({\mdseries\slshape str})\index{NormalizedWhitespace@\texttt{NormalizedWhitespace}}
\label{NormalizedWhitespace}
}\hfill{\scriptsize (function)}}\\


 This function returns a copy of string \mbox{\texttt{\mdseries\slshape str}} to which \texttt{NormalizeWhitespace} (\ref{NormalizeWhitespace}) was applied. }

 

\subsection{\textcolor{Chapter }{RemoveCharacters}}
\logpage{[ 27, 7, 19 ]}\nobreak
\hyperdef{L}{X86EBB6EB829723E4}{}
{\noindent\textcolor{FuncColor}{$\triangleright$\enspace\texttt{RemoveCharacters({\mdseries\slshape string, chars})\index{RemoveCharacters@\texttt{RemoveCharacters}}
\label{RemoveCharacters}
}\hfill{\scriptsize (function)}}\\


 Both arguments must be strings. This function efficiently removes all
characters given in \mbox{\texttt{\mdseries\slshape chars}} from \mbox{\texttt{\mdseries\slshape string}}. 
\begin{Verbatim}[commandchars=!@|,fontsize=\small,frame=single,label=Example]
  !gapprompt@gap>| !gapinput@s := "ab c\ndef\n\ng    h i .\n";|
  "ab c\ndef\n\ng    h i .\n"
  !gapprompt@gap>| !gapinput@RemoveCharacters(s, " \n\t\r"); # remove all whitespace characters|
  !gapprompt@gap>| !gapinput@s;|
  "abcdefghi."
\end{Verbatim}
 }

 

\subsection{\textcolor{Chapter }{JoinStringsWithSeparator}}
\logpage{[ 27, 7, 20 ]}\nobreak
\hyperdef{L}{X84624FEB825EC4B5}{}
{\noindent\textcolor{FuncColor}{$\triangleright$\enspace\texttt{JoinStringsWithSeparator({\mdseries\slshape list[, sep]})\index{JoinStringsWithSeparator@\texttt{JoinStringsWithSeparator}}
\label{JoinStringsWithSeparator}
}\hfill{\scriptsize (function)}}\\


 joins \mbox{\texttt{\mdseries\slshape list}} (a list of strings) after interpolating \mbox{\texttt{\mdseries\slshape sep}} (or \texttt{","} if the second argument is omitted) between each adjacent pair of strings; \mbox{\texttt{\mdseries\slshape sep}} should be a string. 

 
\begin{Verbatim}[commandchars=!@|,fontsize=\small,frame=single,label=Example]
  !gapprompt@gap>| !gapinput@list := List([1..10], String);|
  [ "1", "2", "3", "4", "5", "6", "7", "8", "9", "10" ]
  !gapprompt@gap>| !gapinput@JoinStringsWithSeparator(list);|
  "1,2,3,4,5,6,7,8,9,10"
  !gapprompt@gap>| !gapinput@JoinStringsWithSeparator(["The", "quick", "brown", "fox"], " ");|
  "The quick brown fox"
  !gapprompt@gap>| !gapinput@new:= JoinStringsWithSeparator(["a", "b", "c", "d"], ",\n    ");|
  "a,\n    b,\n    c,\n    d"
  !gapprompt@gap>| !gapinput@Print("    ", new, "\n");|
      a,
      b,
      c,
      d
\end{Verbatim}
 }

 

\subsection{\textcolor{Chapter }{Chomp}}
\logpage{[ 27, 7, 21 ]}\nobreak
\hyperdef{L}{X79F8FFC5876D854A}{}
{\noindent\textcolor{FuncColor}{$\triangleright$\enspace\texttt{Chomp({\mdseries\slshape str})\index{Chomp@\texttt{Chomp}}
\label{Chomp}
}\hfill{\scriptsize (function)}}\\


 Like the similarly named Perl function, \texttt{Chomp} removes a trailing newline character (or carriage\texttt{\symbol{45}}return
line\texttt{\symbol{45}}feed couplet) from a string argument \mbox{\texttt{\mdseries\slshape str}} if present and returns the result. If \mbox{\texttt{\mdseries\slshape str}} is not a string or does not have such trailing character(s) it is returned
unchanged. This latter property means that \texttt{Chomp} is safe to use in cases where one is manipulating the result of another
function which might sometimes return \texttt{fail}. 

 
\begin{Verbatim}[commandchars=!@|,fontsize=\small,frame=single,label=Example]
  !gapprompt@gap>| !gapinput@Chomp("The quick brown fox jumps over the lazy dog.\n");|
  "The quick brown fox jumps over the lazy dog."
  !gapprompt@gap>| !gapinput@Chomp("The quick brown fox jumps over the lazy dog.\r\n");|
  "The quick brown fox jumps over the lazy dog."
  !gapprompt@gap>| !gapinput@Chomp("The quick brown fox jumps over the lazy dog.");|
  "The quick brown fox jumps over the lazy dog."
  !gapprompt@gap>| !gapinput@Chomp(fail);|
  fail
  !gapprompt@gap>| !gapinput@Chomp(32);|
  32
\end{Verbatim}
 

 \emph{Note:} \texttt{Chomp} only removes a trailing newline character from \mbox{\texttt{\mdseries\slshape str}}. If your string contains several newline characters and you really want to
split \mbox{\texttt{\mdseries\slshape str}} into lines at the newline characters (and remove those newline characters)
then you should use \texttt{SplitString} (\ref{SplitString}), e.g. 

 
\begin{Verbatim}[commandchars=!@|,fontsize=\small,frame=single,label=Example]
  !gapprompt@gap>| !gapinput@str := "The quick brown fox\njumps over the lazy dog.\n";|
  "The quick brown fox\njumps over the lazy dog.\n"
  !gapprompt@gap>| !gapinput@SplitString(str, "", "\n");|
  [ "The quick brown fox", "jumps over the lazy dog." ]
  !gapprompt@gap>| !gapinput@Chomp(str);|
  "The quick brown fox\njumps over the lazy dog."
\end{Verbatim}
 }

 

\subsection{\textcolor{Chapter }{StartsWith}}
\logpage{[ 27, 7, 22 ]}\nobreak
\hyperdef{L}{X855820848179CC28}{}
{\noindent\textcolor{FuncColor}{$\triangleright$\enspace\texttt{StartsWith({\mdseries\slshape string, prefix})\index{StartsWith@\texttt{StartsWith}}
\label{StartsWith}
}\hfill{\scriptsize (function)}}\\
\noindent\textcolor{FuncColor}{$\triangleright$\enspace\texttt{EndsWith({\mdseries\slshape string, suffix})\index{EndsWith@\texttt{EndsWith}}
\label{EndsWith}
}\hfill{\scriptsize (function)}}\\


 \index{Prefix} \index{Suffix} Determines whether a string starts or ends with another string. }

 

\subsection{\textcolor{Chapter }{StringFormatted}}
\logpage{[ 27, 7, 23 ]}\nobreak
\hyperdef{L}{X8235AD797868E872}{}
{\noindent\textcolor{FuncColor}{$\triangleright$\enspace\texttt{StringFormatted({\mdseries\slshape string, data...})\index{StringFormatted@\texttt{StringFormatted}}
\label{StringFormatted}
}\hfill{\scriptsize (function)}}\\
\noindent\textcolor{FuncColor}{$\triangleright$\enspace\texttt{PrintFormatted({\mdseries\slshape string, data...})\index{PrintFormatted@\texttt{PrintFormatted}}
\label{PrintFormatted}
}\hfill{\scriptsize (function)}}\\
\noindent\textcolor{FuncColor}{$\triangleright$\enspace\texttt{PrintToFormatted({\mdseries\slshape stream, string, data...})\index{PrintToFormatted@\texttt{PrintToFormatted}}
\label{PrintToFormatted}
}\hfill{\scriptsize (function)}}\\


 These functions perform a string formatting operation. They accept a format
string, which can contain replacement fields which are delimited by braces
\texttt{\symbol{123}}\texttt{\symbol{125}}. Each replacement field contains a
numeric or positional argument, describing the element of \mbox{\texttt{\mdseries\slshape data}} to replace the braces with. 

 There are three formatting functions, which differ only in how they output the
formatted string. \texttt{StringFormatted} returns the formatted string, \texttt{PrintFormatted} prints the formatted string and \texttt{PrintToFormatted} appends the formatted string to \mbox{\texttt{\mdseries\slshape stream}}, which can be either an output stream or a filename. 

 The arguments after \mbox{\texttt{\mdseries\slshape string}} form a list \mbox{\texttt{\mdseries\slshape data}} of values used to substitute the replacement fields in \mbox{\texttt{\mdseries\slshape string}}, using the following formatting rules: 

 \mbox{\texttt{\mdseries\slshape string}} is treated as a normal string, except for occurrences of \texttt{\texttt{\symbol{123}}} and \texttt{\texttt{\symbol{125}}}, which follow special rules, as follows: 

 The contents of \texttt{\texttt{\symbol{123}} \texttt{\symbol{125}}} is split by a \texttt{!} into \texttt{\texttt{\symbol{123}}id!format\texttt{\symbol{125}}}, where both \texttt{id} and \texttt{format} are optional. If the \texttt{!} is omitted, the bracket is treated as \texttt{\texttt{\symbol{123}}id\texttt{\symbol{125}}} with no \texttt{format}. 

 \texttt{id} is interpreted as follows: 
\begin{description}
\item[{An integer \texttt{i}}]  Take the \texttt{i}th element of \mbox{\texttt{\mdseries\slshape data}}. 
\item[{A string \texttt{str}}]  If this is used, the first element of \mbox{\texttt{\mdseries\slshape data}} must be a record \texttt{r}. In this case, the value \texttt{r.(str)} is taken. 
\item[{No id given}]  Take the \texttt{j}th element of \mbox{\texttt{\mdseries\slshape data}}, where \texttt{j} is the number of replacement fields with no id in the format string so far. If
any replacement field has no id, then all replacement fields must have no id. 
\end{description}
 A single brace can be outputted by doubling, so \texttt{\texttt{\symbol{123}}\texttt{\symbol{123}}} in the format string produces \texttt{\texttt{\symbol{123}}} and \texttt{\texttt{\symbol{125}}\texttt{\symbol{125}}} produces \texttt{\texttt{\symbol{125}}}. 

 The \texttt{format} decides how the variable is printed. \texttt{format} must be one of \texttt{s} (which uses \texttt{String} (\ref{String})), \texttt{v} (which uses \texttt{ViewString} (\ref{ViewString})) or \texttt{d} (which calls \texttt{DisplayString} (\ref{DisplayString})). The default value for \texttt{format} is \texttt{s}. }

 
\begin{Verbatim}[commandchars=@|A,fontsize=\small,frame=single,label=Example]
  @gapprompt|gap>A @gapinput|StringFormatted("I have {} cats and {} dogs", 4, 5);A
  "I have 4 cats and 5 dogs"
  @gapprompt|gap>A @gapinput|StringFormatted("I have {2} cats and {1} dogs", 4, 5);A
  "I have 5 cats and 4 dogs"
  @gapprompt|gap>A @gapinput|StringFormatted("I have {cats} cats and {dogs} dogs", rec(cats:=3, dogs:=2));A
  "I have 3 cats and 2 dogs"
  @gapprompt|gap>A @gapinput|StringFormatted("We use {{ and }} to mark {dogs} dogs", rec(cats:=3, dogs:=2));A
  "We use { and } to mark 2 dogs"
  @gapprompt|gap>A @gapinput|sym3 := SymmetricGroup(3);;A
  @gapprompt|gap>A @gapinput|StringFormatted("String: {1!s}, ViewString: {1!v}", sym3);A
  "String: SymmetricGroup( [ 1 .. 3 ] ), ViewString: Sym( [ 1 .. 3 ] )"
\end{Verbatim}
 The following two functions convert basic strings to lists of numbers and vice
versa. They are useful for examples of text encryption. 

\subsection{\textcolor{Chapter }{NumbersString}}
\logpage{[ 27, 7, 24 ]}\nobreak
\hyperdef{L}{X7848A9D878FD59BB}{}
{\noindent\textcolor{FuncColor}{$\triangleright$\enspace\texttt{NumbersString({\mdseries\slshape s, m[, table]})\index{NumbersString@\texttt{NumbersString}}
\label{NumbersString}
}\hfill{\scriptsize (function)}}\\


 \texttt{NumbersString} takes a string message \mbox{\texttt{\mdseries\slshape s}} and returns a list of integers, each not exceeding the integer \mbox{\texttt{\mdseries\slshape m}} that encode the message using the scheme $A=11$, $B=12$ and so on (and converting lower case to upper case). If a list of characters
is given in \mbox{\texttt{\mdseries\slshape table}}, it is used instead for encoding). 
\begin{Verbatim}[commandchars=!@|,fontsize=\small,frame=single,label=Example]
  !gapprompt@gap>| !gapinput@l:=NumbersString("Twas brillig and the slithy toves",1000000);|
  [ 303311, 291012, 281922, 221917, 101124, 141030, 181510, 292219,
    301835, 103025, 321529 ]
\end{Verbatim}
 }

 

\subsection{\textcolor{Chapter }{StringNumbers}}
\logpage{[ 27, 7, 25 ]}\nobreak
\hyperdef{L}{X787EAB117816578E}{}
{\noindent\textcolor{FuncColor}{$\triangleright$\enspace\texttt{StringNumbers({\mdseries\slshape l, m[, table]})\index{StringNumbers@\texttt{StringNumbers}}
\label{StringNumbers}
}\hfill{\scriptsize (function)}}\\


 \texttt{StringNumbers} takes a list \mbox{\texttt{\mdseries\slshape l}} of integers that was encoded using \texttt{NumbersString} (\ref{NumbersString}) and the size integer \mbox{\texttt{\mdseries\slshape m}}, and returns a message string, using the scheme $A=11$, $B=12$ and so on. If a list of characters is given in \mbox{\texttt{\mdseries\slshape table}}, it is used instead for decoding). 
\begin{Verbatim}[commandchars=!@|,fontsize=\small,frame=single,label=Example]
  !gapprompt@gap>| !gapinput@StringNumbers(l,1000000);|
  "TWAS BRILLIG AND THE SLITHY TOVES"
\end{Verbatim}
 }

 

\subsection{\textcolor{Chapter }{StringOfMemoryAmount}}
\logpage{[ 27, 7, 26 ]}\nobreak
\hyperdef{L}{X82975B6480932683}{}
{\noindent\textcolor{FuncColor}{$\triangleright$\enspace\texttt{StringOfMemoryAmount({\mdseries\slshape numbytes})\index{StringOfMemoryAmount@\texttt{StringOfMemoryAmount}}
\label{StringOfMemoryAmount}
}\hfill{\scriptsize (function)}}\\


 This function returns a human\texttt{\symbol{45}}readable string representing \mbox{\texttt{\mdseries\slshape numbytes}} of memory. It is used in printing amounts of memory allocated by tests and
benchmarks. Binary prefixes (representing powers of 1024) are used. }

 
\begin{Verbatim}[commandchars=!@|,fontsize=\small,frame=single,label=Example]
  !gapprompt@gap>| !gapinput@StringOfMemoryAmount(123456789);|
  "117MB"
\end{Verbatim}
 }

  
\section{\textcolor{Chapter }{Character Conversion}}\label{Character Conversion}
\logpage{[ 27, 8, 0 ]}
\hyperdef{L}{X844BDC8578A3B508}{}
{
  The following functions convert characters in their internal integer values
and vice versa. Note that the number corresponding to a particular character
might depend on the system used. While most systems use an extension of ASCII,
in particular character values outside the range \texttt{[ 32 .. 126 ]} might differ between architectures. 

 

\subsection{\textcolor{Chapter }{IntChar}}
\logpage{[ 27, 8, 1 ]}\nobreak
\hyperdef{L}{X826D95D680F87D23}{}
{\noindent\textcolor{FuncColor}{$\triangleright$\enspace\texttt{IntChar({\mdseries\slshape char})\index{IntChar@\texttt{IntChar}}
\label{IntChar}
}\hfill{\scriptsize (function)}}\\


 returns an integer value in the range \texttt{[ 0 .. 255 ]} that corresponds to \mbox{\texttt{\mdseries\slshape char}}. }

 

\subsection{\textcolor{Chapter }{CharInt}}
\logpage{[ 27, 8, 2 ]}\nobreak
\hyperdef{L}{X87B6C1AF7E4A6639}{}
{\noindent\textcolor{FuncColor}{$\triangleright$\enspace\texttt{CharInt({\mdseries\slshape int})\index{CharInt@\texttt{CharInt}}
\label{CharInt}
}\hfill{\scriptsize (function)}}\\


 returns a character that corresponds to the integer value \mbox{\texttt{\mdseries\slshape int}}, which must be in the range \texttt{[ 0 .. 255 ]}. 

 
\begin{Verbatim}[commandchars=!@|,fontsize=\small,frame=single,label=Example]
  !gapprompt@gap>| !gapinput@c:=CharInt(65);|
  'A'
  !gapprompt@gap>| !gapinput@IntChar(c);|
  65
\end{Verbatim}
 }

 

\subsection{\textcolor{Chapter }{SIntChar}}
\logpage{[ 27, 8, 3 ]}\nobreak
\hyperdef{L}{X8159CE81798DDA76}{}
{\noindent\textcolor{FuncColor}{$\triangleright$\enspace\texttt{SIntChar({\mdseries\slshape char})\index{SIntChar@\texttt{SIntChar}}
\label{SIntChar}
}\hfill{\scriptsize (function)}}\\


 returns a signed integer value in the range \texttt{[ \texttt{\symbol{45}}128 .. 127 ]} that corresponds to \mbox{\texttt{\mdseries\slshape char}}. }

 

\subsection{\textcolor{Chapter }{CharSInt}}
\logpage{[ 27, 8, 4 ]}\nobreak
\hyperdef{L}{X78E6611A829DDA3E}{}
{\noindent\textcolor{FuncColor}{$\triangleright$\enspace\texttt{CharSInt({\mdseries\slshape int})\index{CharSInt@\texttt{CharSInt}}
\label{CharSInt}
}\hfill{\scriptsize (function)}}\\


 returns a character which corresponds to the signed integer value \mbox{\texttt{\mdseries\slshape int}}, which must be in the range \texttt{[ \texttt{\symbol{45}}128 .. 127 ]}. 

 The signed and unsigned integer functions behave the same for values in the
range \texttt{[ 0 .. 127 ]}. 

 
\begin{Verbatim}[commandchars=!@|,fontsize=\small,frame=single,label=Example]
  !gapprompt@gap>| !gapinput@SIntChar(c);|
  65
  !gapprompt@gap>| !gapinput@c:=CharSInt(-20);;|
  !gapprompt@gap>| !gapinput@SIntChar(c);|
  -20
  !gapprompt@gap>| !gapinput@IntChar(c);|
  236
  !gapprompt@gap>| !gapinput@SIntChar(CharInt(255));|
  -1
\end{Verbatim}
 }

 }

  
\section{\textcolor{Chapter }{Operations to Evaluate Strings}}\label{Operations to Evaluate Strings}
\logpage{[ 27, 9, 0 ]}
\hyperdef{L}{X78D9BD857F890C0A}{}
{
  

\subsection{\textcolor{Chapter }{Int (for strings)}}
\logpage{[ 27, 9, 1 ]}\nobreak
\hyperdef{L}{X7B6D118184F692A0}{}
{\noindent\textcolor{FuncColor}{$\triangleright$\enspace\texttt{Int({\mdseries\slshape str})\index{Int@\texttt{Int}!for strings}
\label{Int:for strings}
}\hfill{\scriptsize (attribute)}}\\


 \index{evaluation!strings} returns an integer as represented by the string \mbox{\texttt{\mdseries\slshape str}}. The argument string may optionally start with the sign character \texttt{\texttt{\symbol{45}}}, followed by a sequence of decimal digits. For any other input \texttt{fail} is returned. 

 For backwards compatibility, the empty string is accepted, in which case $0$ is returned as result. 

 
\begin{Verbatim}[commandchars=!@|,fontsize=\small,frame=single,label=Example]
  !gapprompt@gap>| !gapinput@Int("12345");|
  12345
  !gapprompt@gap>| !gapinput@Int("123/45");|
  fail
  !gapprompt@gap>| !gapinput@Int("1+2");|
  fail
  !gapprompt@gap>| !gapinput@Int("-12");|
  -12
  !gapprompt@gap>| !gapinput@Int("");|
  0
\end{Verbatim}
 }

 

\subsection{\textcolor{Chapter }{Rat (for strings)}}
\logpage{[ 27, 9, 2 ]}\nobreak
\hyperdef{L}{X87AD395584294FF2}{}
{\noindent\textcolor{FuncColor}{$\triangleright$\enspace\texttt{Rat({\mdseries\slshape str})\index{Rat@\texttt{Rat}!for strings}
\label{Rat:for strings}
}\hfill{\scriptsize (attribute)}}\\


 \index{evaluation!strings} returns a rational as represented by the string \mbox{\texttt{\mdseries\slshape str}}. The argument string may optionally start with the sign character \texttt{\texttt{\symbol{45}}}, followed by either a sequence of decimal digits or by two sequences of
decimal digits that are separated by one of the characters \texttt{/} or \texttt{.}, where the latter stands for a decimal dot. For any other input \texttt{fail} is returned. 

 
\begin{Verbatim}[commandchars=!@|,fontsize=\small,frame=single,label=Example]
  !gapprompt@gap>| !gapinput@Rat("123/45");|
  41/15
  !gapprompt@gap>| !gapinput@Rat("-123.45");|
  -2469/20
\end{Verbatim}
 }

 

\subsection{\textcolor{Chapter }{IntHexString}}
\logpage{[ 27, 9, 3 ]}\nobreak
\hyperdef{L}{X796D366B7DDEFF67}{}
{\noindent\textcolor{FuncColor}{$\triangleright$\enspace\texttt{IntHexString({\mdseries\slshape str})\index{IntHexString@\texttt{IntHexString}}
\label{IntHexString}
}\hfill{\scriptsize (function)}}\\


 \index{evaluation!strings} returns an integer as represented by the string \mbox{\texttt{\mdseries\slshape str}}. The argument string may optionally start with the sign character \texttt{\texttt{\symbol{45}}}, followed by a sequence of hexadecimal digits. Here the letters \texttt{a}\texttt{\symbol{45}}\texttt{f} or \texttt{A}\texttt{\symbol{45}}\texttt{F} are used as \emph{digits} $10$ to $15$. Any other input results in an error. 

 This function can be used (together with \texttt{HexStringInt} (\ref{HexStringInt})) for efficiently storing and reading large integers from respectively into \textsf{GAP}. Note that the translation between integers and their hexadecimal
representation costs linear computation time in terms of the number of digits,
while translation from and into decimal representation needs substantial
computations. 

 
\begin{Verbatim}[commandchars=!@|,fontsize=\small,frame=single,label=Example]
  !gapprompt@gap>| !gapinput@IntHexString("-abcdef0123456789");|
  -12379813738877118345
  !gapprompt@gap>| !gapinput@HexStringInt(last);|
  "-ABCDEF0123456789"
\end{Verbatim}
 }

 

\subsection{\textcolor{Chapter }{Ordinal}}
\logpage{[ 27, 9, 4 ]}\nobreak
\hyperdef{L}{X7C0C29C87CBA97B7}{}
{\noindent\textcolor{FuncColor}{$\triangleright$\enspace\texttt{Ordinal({\mdseries\slshape n})\index{Ordinal@\texttt{Ordinal}}
\label{Ordinal}
}\hfill{\scriptsize (function)}}\\


 returns the ordinal of the integer \mbox{\texttt{\mdseries\slshape n}} as a string. 
\begin{Verbatim}[commandchars=!@|,fontsize=\small,frame=single,label=Example]
  !gapprompt@gap>| !gapinput@Ordinal(2);  Ordinal(21);  Ordinal(33);  Ordinal(-33);|
  "2nd"
  "21st"
  "33rd"
  "-33rd"
\end{Verbatim}
 }

 

\subsection{\textcolor{Chapter }{EvalString}}
\logpage{[ 27, 9, 5 ]}\nobreak
\hyperdef{L}{X7DE4CCD285440659}{}
{\noindent\textcolor{FuncColor}{$\triangleright$\enspace\texttt{EvalString({\mdseries\slshape expr})\index{EvalString@\texttt{EvalString}}
\label{EvalString}
}\hfill{\scriptsize (function)}}\\


 passes the string \mbox{\texttt{\mdseries\slshape expr}} through an input text stream so that \textsf{GAP} interprets it, and returns the result. 

 
\begin{Verbatim}[commandchars=!@|,fontsize=\small,frame=single,label=Example]
  !gapprompt@gap>| !gapinput@a:=10;|
  10
  !gapprompt@gap>| !gapinput@EvalString("a^2");|
  100
\end{Verbatim}
 

 \texttt{EvalString} is intended for \emph{single} expressions. A sequence of commands may be interpreted by using the functions \texttt{InputTextString} (\ref{InputTextString}) and \texttt{ReadAsFunction} (\ref{ReadAsFunction:for streams}) together; see \ref{Operations for Input Streams} for an example. 

 If \texttt{EvalString} is used inside a function, then it doesn't know about the local variables and
the arguments of the function. A possible workaround is to define global
variables in advance, and then to assign the values of the local variables to
the global ones, like in the example below. 

 
\begin{Verbatim}[commandchars=!@|,fontsize=\small,frame=single,label=Example]
  !gapprompt@gap>| !gapinput@global_a := 0;;|
  !gapprompt@gap>| !gapinput@global_b := 0;;|
  !gapprompt@gap>| !gapinput@example := function ( local_a )|
  !gapprompt@>| !gapinput@    local  local_b;|
  !gapprompt@>| !gapinput@    local_b := 5;|
  !gapprompt@>| !gapinput@    global_a := local_a;|
  !gapprompt@>| !gapinput@    global_b := local_b;|
  !gapprompt@>| !gapinput@    return EvalString( "global_a * global_b" );|
  !gapprompt@>| !gapinput@end;;|
  !gapprompt@gap>| !gapinput@example( 2 );|
  10
\end{Verbatim}
 }

 

\subsection{\textcolor{Chapter }{CrcString}}
\logpage{[ 27, 9, 6 ]}\nobreak
\hyperdef{L}{X7D4B9D7A7995C55D}{}
{\noindent\textcolor{FuncColor}{$\triangleright$\enspace\texttt{CrcString({\mdseries\slshape str})\index{CrcString@\texttt{CrcString}}
\label{CrcString}
}\hfill{\scriptsize (function)}}\\
\textbf{\indent Returns:\ }
an integer



 \index{hash function} \index{checksum} This function computes a CRC (cyclic redundancy check) number from a string \mbox{\texttt{\mdseries\slshape str}}. See also \texttt{CrcFile} (\ref{CrcFile}) and \texttt{HexSHA256} (\ref{HexSHA256}). 
\begin{Verbatim}[commandchars=!@|,fontsize=\small,frame=single,label=Example]
  !gapprompt@gap>| !gapinput@CrcString("GAP example string");|
  -50451670
\end{Verbatim}
 }

 

\subsection{\textcolor{Chapter }{HexSHA256}}
\logpage{[ 27, 9, 7 ]}\nobreak
\hyperdef{L}{X7873D1F28779B490}{}
{\noindent\textcolor{FuncColor}{$\triangleright$\enspace\texttt{HexSHA256({\mdseries\slshape string})\index{HexSHA256@\texttt{HexSHA256}}
\label{HexSHA256}
}\hfill{\scriptsize (function)}}\\
\noindent\textcolor{FuncColor}{$\triangleright$\enspace\texttt{HexSHA256({\mdseries\slshape stream})\index{HexSHA256@\texttt{HexSHA256}!for a stream}
\label{HexSHA256:for a stream}
}\hfill{\scriptsize (function)}}\\


 \index{hash function} \index{checksum} Return the SHA\texttt{\symbol{45}}256 cryptographic checksum of the bytes in \mbox{\texttt{\mdseries\slshape string}}, resp. of the data in the input stream object \mbox{\texttt{\mdseries\slshape stream}} (see Chapter{\nobreakspace}\ref{Streams} to learn about streams) when read from the current position until EOF
(end\texttt{\symbol{45}}of\texttt{\symbol{45}}file). 

 The checksum is returned as string with 64 lowercase hexadecimal digits. 
\begin{Verbatim}[commandchars=!@|,fontsize=\small,frame=single,label=Example]
  !gapprompt@gap>| !gapinput@HexSHA256("abcd");|
  "88d4266fd4e6338d13b845fcf289579d209c897823b9217da3e161936f031589"
  !gapprompt@gap>| !gapinput@HexSHA256(InputTextString("abcd"));|
  "88d4266fd4e6338d13b845fcf289579d209c897823b9217da3e161936f031589"
\end{Verbatim}
 }

 

\subsection{\textcolor{Chapter }{Pluralize}}
\logpage{[ 27, 9, 8 ]}\nobreak
\hyperdef{L}{X83FF17E782E6FFF3}{}
{\noindent\textcolor{FuncColor}{$\triangleright$\enspace\texttt{Pluralize({\mdseries\slshape [count, ]string[, plural]})\index{Pluralize@\texttt{Pluralize}}
\label{Pluralize}
}\hfill{\scriptsize (function)}}\\
\textbf{\indent Returns:\ }
A string



 This function returns an attempt at the appropriate pluralization of a string
(considered as a singular English noun), using several rules and heuristics of
English grammar. 

 The arguments to this function are an optional non\texttt{\symbol{45}}negative
integer \mbox{\texttt{\mdseries\slshape count}} (the number of objects in question), a non\texttt{\symbol{45}}empty string \mbox{\texttt{\mdseries\slshape string}} (the singular form of the object in question), and an optional additional
string \mbox{\texttt{\mdseries\slshape plural}} (the plural form of \mbox{\texttt{\mdseries\slshape string}}). 

 If \mbox{\texttt{\mdseries\slshape plural}} is given, then \texttt{Pluralize} uses it as the plural form of \mbox{\texttt{\mdseries\slshape string}}, otherwise \texttt{Pluralize} makes an informed guess at the plural. 

 If \mbox{\texttt{\mdseries\slshape count}} is not given, then \texttt{Pluralize} returns this plural form of \mbox{\texttt{\mdseries\slshape string}}. If \mbox{\texttt{\mdseries\slshape count}} is given and has value $n \neq 1$, then this string is prepended by
"\texttt{\symbol{92}}{\textgreater}n\texttt{\symbol{92}}{\textless} "; else if \mbox{\texttt{\mdseries\slshape count}} has value $1$, then \texttt{Pluralize} returns \mbox{\texttt{\mdseries\slshape string}}, prepended by
"\texttt{\symbol{92}}{\textgreater}1\texttt{\symbol{92}}{\textless} ". 

 Note that \texttt{StripLineBreakCharacters} (\ref{StripLineBreakCharacters}) can be used to remove the control characters \texttt{\texttt{\symbol{92}}{\textless}} and \texttt{\texttt{\symbol{92}}{\textgreater}} from the return value. }

 
\begin{Verbatim}[commandchars=!@|,fontsize=\small,frame=single,label=Example]
  !gapprompt@gap>| !gapinput@Pluralize( "generator" );|
  "generators"
  !gapprompt@gap>| !gapinput@Pluralize( 1, "generator" );|
  "\>1\< generator"
  !gapprompt@gap>| !gapinput@Pluralize( 0, "generator" );|
  "\>0\< generators"
  !gapprompt@gap>| !gapinput@Pluralize( "man", "men" );|
  "men"
  !gapprompt@gap>| !gapinput@Pluralize( 1, "man", "men" );|
  "\>1\< man"
  !gapprompt@gap>| !gapinput@Print( Pluralize( 2, "man", "men" ) );|
  2 men
  !gapprompt@gap>| !gapinput@Print( Pluralize( 2, "vertex" ) );|
  2 vertices
  !gapprompt@gap>| !gapinput@Print( Pluralize( 3, "matrix" ) );|
  3 matrices
  !gapprompt@gap>| !gapinput@Print( Pluralize( 4, "battery" ) );|
  4 batteries
\end{Verbatim}
 }

  
\section{\textcolor{Chapter }{Calendar Arithmetic}}\label{Calendar Arithmetic}
\logpage{[ 27, 10, 0 ]}
\hyperdef{L}{X78F20AA1804D524F}{}
{
  All calendar functions use the Gregorian calendar. 

\subsection{\textcolor{Chapter }{DaysInYear}}
\logpage{[ 27, 10, 1 ]}\nobreak
\hyperdef{L}{X87BA46787FF000E8}{}
{\noindent\textcolor{FuncColor}{$\triangleright$\enspace\texttt{DaysInYear({\mdseries\slshape year})\index{DaysInYear@\texttt{DaysInYear}}
\label{DaysInYear}
}\hfill{\scriptsize (function)}}\\


 returns the number of days in the year \mbox{\texttt{\mdseries\slshape year}}. }

 

\subsection{\textcolor{Chapter }{DaysInMonth}}
\logpage{[ 27, 10, 2 ]}\nobreak
\hyperdef{L}{X8791B0B386D59ADB}{}
{\noindent\textcolor{FuncColor}{$\triangleright$\enspace\texttt{DaysInMonth({\mdseries\slshape month, year})\index{DaysInMonth@\texttt{DaysInMonth}}
\label{DaysInMonth}
}\hfill{\scriptsize (function)}}\\


 returns the number of days in month number \mbox{\texttt{\mdseries\slshape month}} of \mbox{\texttt{\mdseries\slshape year}}, and \texttt{fail} if \texttt{month} is not in the valid range. 
\begin{Verbatim}[commandchars=!@|,fontsize=\small,frame=single,label=Example]
  !gapprompt@gap>| !gapinput@DaysInYear(1998);|
  365
  !gapprompt@gap>| !gapinput@DaysInMonth(3,1998);|
  31
\end{Verbatim}
 }

 

\subsection{\textcolor{Chapter }{DMYDay}}
\logpage{[ 27, 10, 3 ]}\nobreak
\hyperdef{L}{X7CED84C07CD5E2CF}{}
{\noindent\textcolor{FuncColor}{$\triangleright$\enspace\texttt{DMYDay({\mdseries\slshape day})\index{DMYDay@\texttt{DMYDay}}
\label{DMYDay}
}\hfill{\scriptsize (function)}}\\


 converts a number of days, starting
1\texttt{\symbol{45}}Jan\texttt{\symbol{45}}1970, to a list \texttt{[ day, month, year ]} in Gregorian calendar counting. }

 

\subsection{\textcolor{Chapter }{DayDMY}}
\logpage{[ 27, 10, 4 ]}\nobreak
\hyperdef{L}{X7A79DEE07A41B8EF}{}
{\noindent\textcolor{FuncColor}{$\triangleright$\enspace\texttt{DayDMY({\mdseries\slshape dmy})\index{DayDMY@\texttt{DayDMY}}
\label{DayDMY}
}\hfill{\scriptsize (function)}}\\


 returns the number of days from
01\texttt{\symbol{45}}Jan\texttt{\symbol{45}}1970 to the day given by \mbox{\texttt{\mdseries\slshape dmy}}, which must be a list of the form \texttt{[ day, month, year ]} in Gregorian calendar counting. The result is \texttt{fail} on input outside valid ranges. 

 Note that this makes not much sense for early dates like: before 1582 (no
Gregorian calendar at all), or before 1753 in many English speaking countries
or before 1917 in Russia. }

 

\subsection{\textcolor{Chapter }{WeekDay}}
\logpage{[ 27, 10, 5 ]}\nobreak
\hyperdef{L}{X87D03FC0809DB6EC}{}
{\noindent\textcolor{FuncColor}{$\triangleright$\enspace\texttt{WeekDay({\mdseries\slshape date})\index{WeekDay@\texttt{WeekDay}}
\label{WeekDay}
}\hfill{\scriptsize (function)}}\\


 returns the weekday of a day given by \mbox{\texttt{\mdseries\slshape date}}, which can be a number of days since
1\texttt{\symbol{45}}Jan\texttt{\symbol{45}}1970 or a list \texttt{[ day, month, year ]}. }

 

\subsection{\textcolor{Chapter }{StringDate}}
\logpage{[ 27, 10, 6 ]}\nobreak
\hyperdef{L}{X7C74C33784CDED6C}{}
{\noindent\textcolor{FuncColor}{$\triangleright$\enspace\texttt{StringDate({\mdseries\slshape date})\index{StringDate@\texttt{StringDate}}
\label{StringDate}
}\hfill{\scriptsize (function)}}\\


 converts \mbox{\texttt{\mdseries\slshape date}} to a readable string. \mbox{\texttt{\mdseries\slshape date}} can be a number of days since 1\texttt{\symbol{45}}Jan\texttt{\symbol{45}}1970
or a list \texttt{[ day, month, year ]}. 
\begin{Verbatim}[commandchars=!@|,fontsize=\small,frame=single,label=Example]
  !gapprompt@gap>| !gapinput@DayDMY([1,1,1970]);DayDMY([2,1,1970]);|
  0
  1
  !gapprompt@gap>| !gapinput@DMYDay(12345);|
  [ 20, 10, 2003 ]
  !gapprompt@gap>| !gapinput@WeekDay([11,3,1998]);|
  "Wed"
  !gapprompt@gap>| !gapinput@StringDate([11,3,1998]);|
  "11-Mar-1998"
\end{Verbatim}
 }

 

\subsection{\textcolor{Chapter }{HMSMSec}}
\logpage{[ 27, 10, 7 ]}\nobreak
\hyperdef{L}{X84A6A2637FB35A32}{}
{\noindent\textcolor{FuncColor}{$\triangleright$\enspace\texttt{HMSMSec({\mdseries\slshape msec})\index{HMSMSec@\texttt{HMSMSec}}
\label{HMSMSec}
}\hfill{\scriptsize (function)}}\\


 converts a number \mbox{\texttt{\mdseries\slshape msec}} of milliseconds into a list \texttt{[ hour, min, sec, milli ]}. }

 

\subsection{\textcolor{Chapter }{SecHMSM}}
\logpage{[ 27, 10, 8 ]}\nobreak
\hyperdef{L}{X879461D77C81100B}{}
{\noindent\textcolor{FuncColor}{$\triangleright$\enspace\texttt{SecHMSM({\mdseries\slshape hmsm})\index{SecHMSM@\texttt{SecHMSM}}
\label{SecHMSM}
}\hfill{\scriptsize (function)}}\\


 is the reverse of \texttt{HMSMSec} (\ref{HMSMSec}). }

 

\subsection{\textcolor{Chapter }{StringTime}}
\logpage{[ 27, 10, 9 ]}\nobreak
\hyperdef{L}{X802469C47F886A59}{}
{\noindent\textcolor{FuncColor}{$\triangleright$\enspace\texttt{StringTime({\mdseries\slshape time})\index{StringTime@\texttt{StringTime}}
\label{StringTime}
}\hfill{\scriptsize (function)}}\\


 converts \mbox{\texttt{\mdseries\slshape time}} (given as a number of milliseconds or a list \texttt{[ hour, min, sec, milli ]}) to a readable string. 
\begin{Verbatim}[commandchars=!@|,fontsize=\small,frame=single,label=Example]
  !gapprompt@gap>| !gapinput@HMSMSec(Factorial(10));|
  [ 1, 0, 28, 800 ]
  !gapprompt@gap>| !gapinput@SecHMSM([1,10,5,13]);|
  4205013
  !gapprompt@gap>| !gapinput@StringTime([1,10,5,13]);|
  " 1:10:05.013"
\end{Verbatim}
 }

 

\subsection{\textcolor{Chapter }{SecondsDMYhms}}
\logpage{[ 27, 10, 10 ]}\nobreak
\hyperdef{L}{X870A71D47B0E936E}{}
{\noindent\textcolor{FuncColor}{$\triangleright$\enspace\texttt{SecondsDMYhms({\mdseries\slshape DMYhms})\index{SecondsDMYhms@\texttt{SecondsDMYhms}}
\label{SecondsDMYhms}
}\hfill{\scriptsize (function)}}\\


 returns the number of seconds from
01\texttt{\symbol{45}}Jan\texttt{\symbol{45}}1970, 00:00:00, to the time given
by \mbox{\texttt{\mdseries\slshape DMYhms}}, which must be a list of the form \texttt{[ day, month, year, hour, minute, second ]}. The remarks on the Gregorian calendar in the section on \texttt{DayDMY} (\ref{DayDMY}) apply here as well. The last three arguments must lie in the appropriate
ranges. }

 

\subsection{\textcolor{Chapter }{DMYhmsSeconds}}
\logpage{[ 27, 10, 11 ]}\nobreak
\hyperdef{L}{X78AF8EA887532B5B}{}
{\noindent\textcolor{FuncColor}{$\triangleright$\enspace\texttt{DMYhmsSeconds({\mdseries\slshape secs})\index{DMYhmsSeconds@\texttt{DMYhmsSeconds}}
\label{DMYhmsSeconds}
}\hfill{\scriptsize (function)}}\\


 This is the inverse function to \texttt{SecondsDMYhms} (\ref{SecondsDMYhms}). 
\begin{Verbatim}[commandchars=!@|,fontsize=\small,frame=single,label=Example]
  !gapprompt@gap>| !gapinput@SecondsDMYhms([ 9, 9, 2001, 1, 46, 40 ]);|
  1000000000
  !gapprompt@gap>| !gapinput@DMYhmsSeconds(-1000000000);|
  [ 24, 4, 1938, 22, 13, 20 ]
\end{Verbatim}
 }

 }

  
\section{\textcolor{Chapter }{Obtaining LaTeX Representations of Objects}}\label{Obtaining LaTeX Representations of Objects}
\logpage{[ 27, 11, 0 ]}
\hyperdef{L}{X78024C8087F3E07F}{}
{
  \index{LaTeX!for GAP objects} For the purpose of generating {\LaTeX} source code with \textsf{GAP} it is recommended to add new functions which will print the {\LaTeX} source or return {\LaTeX} strings for further processing. 

 An alternative approach could be based on methods for the default {\LaTeX} representation for each appropriate type of objects. However, there is no
clear notion of a default {\LaTeX} code for any non\texttt{\symbol{45}}trivial mathematical object; moreover,
different output may be required in different contexts. 

 While customisation of such an operation may require changes in a variety of
methods that may be distributed all over the library, the user will have a
clean overview of the whole process of {\LaTeX} code generation if it is contained in a single function. Furthermore, there
may be kinds of objects which are not detected by the method selection, or
there may be a need in additional parameters specifying requirements for the
output. 

 This is why having a special purpose function for each particular case is more
suitable. \textsf{GAP} provides several functions that produce {\LaTeX} strings for those situations where this is nontrivial and reasonable. A useful
example is \texttt{LaTeXStringDecompositionMatrix} (\ref{LaTeXStringDecompositionMatrix}) from the \textsf{GAP} library, others can be found entering \texttt{?LaTeX} at the \textsf{GAP} prompt. Package authors are encouraged to add an index entry \texttt{LaTeX} to the documentation of all {\LaTeX} string producing functions. This way, entering \texttt{?LaTeX} will give an overview of all documented functionality in this direction. }

 }

              
\chapter{\textcolor{Chapter }{Dictionaries and General Hash Tables}}\label{Dictionaries and General Hash Tables}
\logpage{[ 28, 0, 0 ]}
\hyperdef{L}{X867203C5877489A2}{}
{
  People and computers spend a large amount of time with searching. Dictionaries
are an abstract data structure which facilitates searching for objects.
Depending on the kind of objects the implementation will use a variety of
possible internal storage methods that will aim to provide the fastest
possible access to objects. These internal methods include 

 
\begin{description}
\item[{Hash Tables}]  for objects for which a hash function has been defined. 
\item[{Direct Indexing}]  if the domain is small and cheaply enumerable 
\item[{Sorted Lists}]  if a total order can be computed easily 
\item[{Plain lists}]  for objects for which nothing but an equality test is available. 
\end{description}
  
\section{\textcolor{Chapter }{Using Dictionaries}}\label{Using Dictionaries}
\logpage{[ 28, 1, 0 ]}
\hyperdef{L}{X81560C4083E27955}{}
{
  The standard way to use dictionaries is to first create a dictionary (using \texttt{NewDictionary} (\ref{NewDictionary}), and then to store objects (and associated information) in it and look them
up. 

 For the creation of objects the user has to make a few choices: Is the
dictionary only to be used to check whether objects are known already, or
whether associated information is to be stored with the objects. This second
case is called a \emph{lookup dictionary} and is selected by the second parameter of \texttt{NewDictionary} (\ref{NewDictionary}). 

 The second choice is to indicate which kind of objects are to be stored. This
choice will decide the internal storage used. This kind of objects is
specified by the first parameter to \texttt{NewDictionary} (\ref{NewDictionary}), which is a ``sample'' object. 

 In some cases however such a sample object is not specific enough. For example
when storing vectors over a finite field, it would not be clear whether all
vectors will be over a prime field or over a field extension. Such an issue
can be resolved by indicating in an (optional) third parameter to \texttt{NewDictionary} (\ref{NewDictionary}) a \emph{domain} which has to be a collection that will contain all objects to be used with
this dictionary. (Such a domain may also be used internally to decide that
direct indexing can be used). 

 The reason for this choice of giving two parameters is that in some cases no
suitable collection of objects has been defined in \textsf{GAP} \texttt{\symbol{45}} for example for permutations there is no object
representing the symmetric group on infinitely many points. 

 Once a dictionary has been created, it is possible to use \texttt{RepresentationsOfObject} (\ref{RepresentationsOfObject}) to check which representation is used by \textsf{GAP}. 

 In the following example, we create a dictionary to store permutations with
associated information. 
\begin{Verbatim}[commandchars=!@|,fontsize=\small,frame=single,label=Example]
  !gapprompt@gap>| !gapinput@d:=NewDictionary((1,2,3),true);;|
  !gapprompt@gap>| !gapinput@AddDictionary(d,(1,2),1);|
  !gapprompt@gap>| !gapinput@AddDictionary(d,(5,6),9);|
  !gapprompt@gap>| !gapinput@AddDictionary(d,(4,7),2);|
  !gapprompt@gap>| !gapinput@LookupDictionary(d,(5,6));|
  9
  !gapprompt@gap>| !gapinput@LookupDictionary(d,(5,7));|
  fail
\end{Verbatim}
 A typical example of this use would be in an orbit algorithm. The dictionary
would be used to store the elements known in the orbit together with their
respective orbit positions. 

 We observe that this dictionary is stored internally by a sorted list. On the
other hand, if we have an explicit, sorted element list, direct indexing is to
be used. 
\begin{Verbatim}[commandchars=!@|,fontsize=\small,frame=single,label=Example]
  !gapprompt@gap>| !gapinput@RepresentationsOfObject(d);|
  [ "IsComponentObjectRep", "IsNonAtomicComponentObjectRep",
    "IsDictionaryDefaultRep", "IsListDictionary",
    "IsListLookupDictionary", "IsSortDictionary",
    "IsSortLookupDictionary" ]
  !gapprompt@gap>| !gapinput@d:=NewDictionary((1,2,3),true,Elements(SymmetricGroup(5)));;|
  !gapprompt@gap>| !gapinput@RepresentationsOfObject(d);|
  [ "IsComponentObjectRep", "IsNonAtomicComponentObjectRep",
    "IsDictionaryDefaultRep", "IsPositionDictionary",
    "IsPositionLookupDictionary" ]
\end{Verbatim}
 (Just indicating \texttt{SymmetricGroup(5)} as a third parameter would still keep the first storage method, as indexing
would be too expensive if no explicit element list is known.) 

 The same effect happens in the following example, in which we work with
vectors: Indicating only a vector only enables sorted index, as it cannot be
known whether all vectors will be defined over the prime field. On the other
hand, providing the vector space (and thus limiting the domain) enables the
use of hashing (which will be faster). 
\begin{Verbatim}[commandchars=!@|,fontsize=\small,frame=single,label=Example]
  !gapprompt@gap>| !gapinput@v:=GF(2)^7;;|
  !gapprompt@gap>| !gapinput@d:=NewDictionary(Zero(v),true);;|
  !gapprompt@gap>| !gapinput@RepresentationsOfObject(d);|
  [ "IsComponentObjectRep", "IsNonAtomicComponentObjectRep",
    "IsDictionaryDefaultRep", "IsListDictionary",
    "IsListLookupDictionary", "IsSortDictionary",
    "IsSortLookupDictionary" ]
  !gapprompt@gap>| !gapinput@d:=NewDictionary(Zero(v),true,v);;|
  !gapprompt@gap>| !gapinput@RepresentationsOfObject(d);|
  [ "IsComponentObjectRep", "IsNonAtomicComponentObjectRep",
    "IsDictionaryDefaultRep", "IsPositionDictionary",
    "IsPositionLookupDictionary" ]
\end{Verbatim}
 }

  
\section{\textcolor{Chapter }{Dictionaries}}\label{Dictionaries}
\logpage{[ 28, 2, 0 ]}
\hyperdef{L}{X7B571EA282AF70D7}{}
{
  This section contains the formal declarations for dictionaries. For
information on how to use them, please refer to the previous
section{\nobreakspace}\ref{Using Dictionaries}. There are several ways how dictionaries are implemented: As lists, as sorted
lists, as hash tables or via binary lists. A user however will just have to
call \texttt{NewDictionary} (\ref{NewDictionary}) and obtain a ``suitable'' dictionary for the kind of objects she wants to create. It is possible however
to create hash tables (see{\nobreakspace}\ref{General Hash Tables}) and dictionaries using binary lists (see{\nobreakspace}\texttt{DictionaryByPosition} (\ref{DictionaryByPosition})). 

 The use of two objects, \mbox{\texttt{\mdseries\slshape obj}} and \mbox{\texttt{\mdseries\slshape objcoll}} to parametrize the objects a dictionary is able to store might look confusing.
However there are situations where either of them might be needed: 

 The first situation is that of objects, for which no formal ``collection object'' has been defined. A typical example here might be subspaces of a vector space. \textsf{GAP} does not formally define a ``Grassmannian'' or anything else to represent the multitude of all subspaces. So it is only
possible to give the dictionary a ``sample object''. 

 The other situation is that of an object which might represent quite varied
domains. The permutation $(1,10^6)$ might be the nontrivial element of a cyclic group of order 2, it might be a
representative of $S_{{10^6}}$. In the first situation the best approach might be just to have two entries
for the two possible objects, in the second situation a much more elaborate
approach might be needed. 

 An algorithm that creates a dictionary will usually know a priori, from what
domain all the objects will be, giving this domain permits to use a more
efficient dictionary. 

 This is particularly true for vectors. From a single vector one cannot decide
whether a calculation will take place over the smallest field containing all
its entries or over a larger field. 

\subsection{\textcolor{Chapter }{NewDictionary}}
\logpage{[ 28, 2, 1 ]}\nobreak
\hyperdef{L}{X7E78E3E983A5C895}{}
{\noindent\textcolor{FuncColor}{$\triangleright$\enspace\texttt{NewDictionary({\mdseries\slshape obj, look[, objcoll]})\index{NewDictionary@\texttt{NewDictionary}}
\label{NewDictionary}
}\hfill{\scriptsize (function)}}\\


 creates a new dictionary for objects such as \mbox{\texttt{\mdseries\slshape obj}}. If \mbox{\texttt{\mdseries\slshape objcoll}} is given the dictionary will be for objects only from this collection, knowing
this can improve the performance. If \mbox{\texttt{\mdseries\slshape objcoll}} is given, \mbox{\texttt{\mdseries\slshape obj}} may be replaced by \texttt{false}, i.e. no sample object is needed. 

 The function tries to find the right kind of dictionary for the basic
dictionary functions to be quick. If \mbox{\texttt{\mdseries\slshape look}} is \texttt{true}, the dictionary will be a lookup dictionary, otherwise it is an ordinary
dictionary. }

 }

  
\section{\textcolor{Chapter }{Dictionaries via Binary Lists}}\label{Dictionaries via Binary Lists}
\logpage{[ 28, 3, 0 ]}
\hyperdef{L}{X86BD015B7B889329}{}
{
  As there are situations where the approach via binary lists is explicitly
desired, such dictionaries can be created deliberately. 

\subsection{\textcolor{Chapter }{DictionaryByPosition}}
\logpage{[ 28, 3, 1 ]}\nobreak
\hyperdef{L}{X865D5BE1830A448D}{}
{\noindent\textcolor{FuncColor}{$\triangleright$\enspace\texttt{DictionaryByPosition({\mdseries\slshape list, lookup})\index{DictionaryByPosition@\texttt{DictionaryByPosition}}
\label{DictionaryByPosition}
}\hfill{\scriptsize (function)}}\\


 creates a new (lookup) dictionary which uses \texttt{PositionCanonical} (\ref{PositionCanonical}) in \mbox{\texttt{\mdseries\slshape list}} for indexing. The dictionary will have an entry \mbox{\texttt{\mdseries\slshape dict}}\texttt{!.blist} which is a bit list corresponding to \mbox{\texttt{\mdseries\slshape list}} indicating the known values. If \mbox{\texttt{\mdseries\slshape look}} is \texttt{true}, the dictionary will be a lookup dictionary, otherwise it is an ordinary
dictionary. }

 

\subsection{\textcolor{Chapter }{IsDictionary}}
\logpage{[ 28, 3, 2 ]}\nobreak
\hyperdef{L}{X87F7247E784021C2}{}
{\noindent\textcolor{FuncColor}{$\triangleright$\enspace\texttt{IsDictionary({\mdseries\slshape obj})\index{IsDictionary@\texttt{IsDictionary}}
\label{IsDictionary}
}\hfill{\scriptsize (Category)}}\\


 A dictionary is a growable collection of objects that permits to add objects
(with associated values) and to check whether an object is already known. }

 

\subsection{\textcolor{Chapter }{IsLookupDictionary}}
\logpage{[ 28, 3, 3 ]}\nobreak
\hyperdef{L}{X7D776BC67ABDDCCE}{}
{\noindent\textcolor{FuncColor}{$\triangleright$\enspace\texttt{IsLookupDictionary({\mdseries\slshape obj})\index{IsLookupDictionary@\texttt{IsLookupDictionary}}
\label{IsLookupDictionary}
}\hfill{\scriptsize (Category)}}\\


 A \emph{lookup dictionary} is a dictionary, which permits not only to check whether an object is
contained, but also to retrieve associated values, using the operation \texttt{LookupDictionary} (\ref{LookupDictionary}). }

 

\subsection{\textcolor{Chapter }{AddDictionary}}
\logpage{[ 28, 3, 4 ]}\nobreak
\hyperdef{L}{X86C4F0507AD98B8A}{}
{\noindent\textcolor{FuncColor}{$\triangleright$\enspace\texttt{AddDictionary({\mdseries\slshape dict, key[, val]})\index{AddDictionary@\texttt{AddDictionary}}
\label{AddDictionary}
}\hfill{\scriptsize (operation)}}\\


 adds \mbox{\texttt{\mdseries\slshape key}} to the dictionary \mbox{\texttt{\mdseries\slshape dict}}, storing the associated value \mbox{\texttt{\mdseries\slshape val}} in case \mbox{\texttt{\mdseries\slshape dict}} is a lookup dictionary. If \mbox{\texttt{\mdseries\slshape key}} is not an object of the kind for which the dictionary was specified, or if \mbox{\texttt{\mdseries\slshape key}} is known already to \mbox{\texttt{\mdseries\slshape dict}}, the results are unpredictable. }

 

\subsection{\textcolor{Chapter }{KnowsDictionary}}
\logpage{[ 28, 3, 5 ]}\nobreak
\hyperdef{L}{X808C885D7E267285}{}
{\noindent\textcolor{FuncColor}{$\triangleright$\enspace\texttt{KnowsDictionary({\mdseries\slshape dict, key})\index{KnowsDictionary@\texttt{KnowsDictionary}}
\label{KnowsDictionary}
}\hfill{\scriptsize (operation)}}\\


 checks, whether \mbox{\texttt{\mdseries\slshape key}} is known to the dictionary \mbox{\texttt{\mdseries\slshape dict}}, and returns \texttt{true} or \texttt{false} accordingly. \mbox{\texttt{\mdseries\slshape key}} \emph{must} be an object of the kind for which the dictionary was specified, otherwise the
results are unpredictable. }

 

\subsection{\textcolor{Chapter }{LookupDictionary}}
\logpage{[ 28, 3, 6 ]}\nobreak
\hyperdef{L}{X863706BF847A47EB}{}
{\noindent\textcolor{FuncColor}{$\triangleright$\enspace\texttt{LookupDictionary({\mdseries\slshape dict, key})\index{LookupDictionary@\texttt{LookupDictionary}}
\label{LookupDictionary}
}\hfill{\scriptsize (operation)}}\\


 looks up \mbox{\texttt{\mdseries\slshape key}} in the lookup dictionary \mbox{\texttt{\mdseries\slshape dict}} and returns the associated value. If \mbox{\texttt{\mdseries\slshape key}} is not known to the dictionary, \texttt{fail} is returned. }

 }

  
\section{\textcolor{Chapter }{General Hash Tables}}\label{General Hash Tables}
\logpage{[ 28, 4, 0 ]}
\hyperdef{L}{X8444087381BBA88A}{}
{
  These sections describe some particularities for hash tables. These are
intended mainly for extending the implementation \texttt{\symbol{45}} programs
requiring hash functionality ought to use the dictionary interface described
above. 

 We hash by keys and also store a value. Keys cannot be removed from the table,
but the corresponding value can be changed. Fast access to last hash index
allows you to efficiently store more than one array of values
{\textendash}this facility should be used with care. 

 This code works for any kind of object, provided you have a \texttt{DenseIntKey} (\ref{DenseIntKey}) method to convert the key into a positive integer. This method should ideally
be implemented efficiently in the core. 

Note that, for efficiency, it is currently impossible to create a hash table
with non\texttt{\symbol{45}}positive integers. }

  
\section{\textcolor{Chapter }{Hash keys}}\label{Hash keys}
\logpage{[ 28, 5, 0 ]}
\hyperdef{L}{X85CD6C9B85DE7C54}{}
{
  The crucial step of hashing is to transform key objects into integers such
that equal objects produce the same integer. 

 The actual function used will vary very much on the type of objects. However \textsf{GAP} provides already key functions for some commonly encountered objects. 

\subsection{\textcolor{Chapter }{DenseIntKey}}
\logpage{[ 28, 5, 1 ]}\nobreak
\hyperdef{L}{X79DEEB5783513838}{}
{\noindent\textcolor{FuncColor}{$\triangleright$\enspace\texttt{DenseIntKey({\mdseries\slshape objcoll, obj})\index{DenseIntKey@\texttt{DenseIntKey}}
\label{DenseIntKey}
}\hfill{\scriptsize (operation)}}\\


 returns a function that can be used as hash key function for objects such as \mbox{\texttt{\mdseries\slshape obj}} in the collection \mbox{\texttt{\mdseries\slshape objcoll}}. Typically, \mbox{\texttt{\mdseries\slshape objcoll}} will be a large domain. If the domain is not available, it can be given as \texttt{false} in which case the hash key function will be determined only based on \mbox{\texttt{\mdseries\slshape obj}}. (For a further discussion of these two arguments see{\nobreakspace}\texttt{NewDictionary} (\ref{NewDictionary})). 

 The function returned by \texttt{DenseIntKey} is guaranteed to give different values for different objects. If no suitable
hash key function has been predefined, \texttt{fail} is returned. }

 

\subsection{\textcolor{Chapter }{SparseIntKey}}
\logpage{[ 28, 5, 2 ]}\nobreak
\hyperdef{L}{X87FC10AC81E5F6BA}{}
{\noindent\textcolor{FuncColor}{$\triangleright$\enspace\texttt{SparseIntKey({\mdseries\slshape objcoll, obj})\index{SparseIntKey@\texttt{SparseIntKey}}
\label{SparseIntKey}
}\hfill{\scriptsize (operation)}}\\


 returns a function that can be used as hash key function for objects such as \mbox{\texttt{\mdseries\slshape obj}} in the collection \mbox{\texttt{\mdseries\slshape objcoll}}. In contrast to \texttt{DenseIntKey} (\ref{DenseIntKey}), the function returned may return the same key value for different objects.
If no suitable hash key function has been predefined, \texttt{fail} is returned. }

 }

  
\section{\textcolor{Chapter }{Dense hash tables}}\label{Dense hash tables}
\logpage{[ 28, 6, 0 ]}
\hyperdef{L}{X84D1A83C8247E7FB}{}
{
  Dense hash tables are used for hashing dense sets without collisions, in
particular integers. Keys are stored as an unordered list and values as an
array with holes. The position of a value is given by the function returned by \texttt{DenseIntKey} (\ref{DenseIntKey}), and so \texttt{KeyIntDense} must be one\texttt{\symbol{45}}to\texttt{\symbol{45}}one. 

\subsection{\textcolor{Chapter }{DenseHashTable}}
\logpage{[ 28, 6, 1 ]}\nobreak
\hyperdef{L}{X874FAA447930C7DA}{}
{\noindent\textcolor{FuncColor}{$\triangleright$\enspace\texttt{DenseHashTable({\mdseries\slshape })\index{DenseHashTable@\texttt{DenseHashTable}}
\label{DenseHashTable}
}\hfill{\scriptsize (function)}}\\


 Construct an empty dense hash table. This is the only correct way to construct
such a table. }

 }

  
\section{\textcolor{Chapter }{Sparse hash tables}}\label{Sparse hash tables}
\logpage{[ 28, 7, 0 ]}
\hyperdef{L}{X7FDB74417A19E674}{}
{
  Sparse hash tables are used for hashing sparse sets. Stores keys as an array
with fail denoting an empty position, stores values as an array with holes.
Uses the result of calling \texttt{SparseIntKey} (\ref{SparseIntKey})) of the key. \texttt{DefaultHashLength} is the default starting hash table length; the table is doubled when it
becomes half full. 

 In sparse hash tables, the integer obtained from the hash key is then
transformed to an index position by taking it modulo the length of the hash
array. 

\subsection{\textcolor{Chapter }{SparseHashTable}}
\logpage{[ 28, 7, 1 ]}\nobreak
\hyperdef{L}{X8757D3A785290640}{}
{\noindent\textcolor{FuncColor}{$\triangleright$\enspace\texttt{SparseHashTable({\mdseries\slshape [intkeyfun[, startlength]]})\index{SparseHashTable@\texttt{SparseHashTable}}
\label{SparseHashTable}
}\hfill{\scriptsize (function)}}\\


 Construct an empty sparse hash table. This is the only correct way to
construct such a table. If the argument \mbox{\texttt{\mdseries\slshape intkeyfun}} is given, this function will be used to obtain numbers for the keys passed to
it. If also \mbox{\texttt{\mdseries\slshape startlength}} is given, the hash table will be initialized at that size. }

 

\subsection{\textcolor{Chapter }{DoubleHashArraySize}}
\logpage{[ 28, 7, 2 ]}\nobreak
\hyperdef{L}{X80FDDF957887B4FC}{}
{\noindent\textcolor{FuncColor}{$\triangleright$\enspace\texttt{DoubleHashArraySize({\mdseries\slshape hash})\index{DoubleHashArraySize@\texttt{DoubleHashArraySize}}
\label{DoubleHashArraySize}
}\hfill{\scriptsize (function)}}\\


 Double the size of the hash array and rehash all the entries. This will also
happen automatically when the hash array is half full. }

 }

 }

             
\chapter{\textcolor{Chapter }{Records}}\label{Records}
\logpage{[ 29, 0, 0 ]}
\hyperdef{L}{X7AA1073C7E943DD7}{}
{
  \index{type!records} \emph{Records} are next to lists the most important way to collect objects together. A record
is a collection of \emph{components}. Each component has a unique \emph{name}, which is an identifier that distinguishes this component, and a \emph{value}, which is an object of arbitrary type. We often abbreviate \emph{value of a component} to \emph{element}. We also say that a record \emph{contains} its elements. You can access and change the elements of a record using its
name. 

 Record literals are written by writing down the components in order between ``\texttt{rec(}'' and ``\texttt{)}'', and separating them by commas ``\texttt{,}''. Each component consists of the name, the assignment operator ``\texttt{:=}'', and the value. The \emph{empty record}, i.e., the record with no components, is written as \texttt{rec()}. 

 
\begin{Verbatim}[commandchars=!@|,fontsize=\small,frame=single,label=Example]
  !gapprompt@gap>| !gapinput@rec( a := 1, b := "2" ); # a record with two components|
  rec( a := 1, b := "2" )
  !gapprompt@gap>| !gapinput@rec( a := 1, b := rec( c := 2 ) ); # record may contain records|
  rec( a := 1, b := rec( c := 2 ) )
\end{Verbatim}
 

 We may use the \texttt{Display} (\ref{Display}) function to illustrate the hierarchy of the record components. 

 
\begin{Verbatim}[commandchars=!@|,fontsize=\small,frame=single,label=Example]
  !gapprompt@gap>| !gapinput@Display( last );|
  rec(
    a := 1,
    b := rec(
        c := 2 ) )
\end{Verbatim}
 

 Records usually contain elements of various types, i.e., they are usually not
homogeneous like lists.  
\section{\textcolor{Chapter }{IsRecord and RecNames}}\label{sect:IsRecord}
\logpage{[ 29, 1, 0 ]}
\hyperdef{L}{X864F92347B5A3FF0}{}
{
  

\subsection{\textcolor{Chapter }{IsRecord}}
\logpage{[ 29, 1, 1 ]}\nobreak
\hyperdef{L}{X782A998E7D9EC406}{}
{\noindent\textcolor{FuncColor}{$\triangleright$\enspace\texttt{IsRecord({\mdseries\slshape obj})\index{IsRecord@\texttt{IsRecord}}
\label{IsRecord}
}\hfill{\scriptsize (Category)}}\\
\noindent\textcolor{FuncColor}{$\triangleright$\enspace\texttt{IsRecordCollection({\mdseries\slshape obj})\index{IsRecordCollection@\texttt{IsRecordCollection}}
\label{IsRecordCollection}
}\hfill{\scriptsize (Category)}}\\
\noindent\textcolor{FuncColor}{$\triangleright$\enspace\texttt{IsRecordCollColl({\mdseries\slshape obj})\index{IsRecordCollColl@\texttt{IsRecordCollColl}}
\label{IsRecordCollColl}
}\hfill{\scriptsize (Category)}}\\


 \index{test!for records} 
\begin{Verbatim}[commandchars=!@|,fontsize=\small,frame=single,label=Example]
  !gapprompt@gap>| !gapinput@IsRecord( rec( a := 1, b := 2 ) );|
  true
  !gapprompt@gap>| !gapinput@IsRecord( IsRecord );|
  false
\end{Verbatim}
 }

 

\subsection{\textcolor{Chapter }{RecNames}}
\logpage{[ 29, 1, 2 ]}\nobreak
\hyperdef{L}{X837F1E1F866FB1A0}{}
{\noindent\textcolor{FuncColor}{$\triangleright$\enspace\texttt{RecNames({\mdseries\slshape record})\index{RecNames@\texttt{RecNames}}
\label{RecNames}
}\hfill{\scriptsize (attribute)}}\\


 returns a duplicate\texttt{\symbol{45}}free list of strings corresponding to
the names of the record components of the record \mbox{\texttt{\mdseries\slshape record}}. 

 The list is sorted by \texttt{RNamObj} (\ref{RNamObj:for a string}) for reasons of efficiency; see \texttt{SortBy} (\ref{SortBy}). Therefore the ordering is consistent within one \textsf{GAP} session, but it is not necessarily consistent across different sessions. To
obtain a result that is consistent across \textsf{GAP} sessions, apply \texttt{Set} (\ref{Set}) to the returned list. 

 Note that given a string \texttt{name} containing the name of a record component, you can access the record component
via \texttt{\mbox{\texttt{\mdseries\slshape record}}.(name)}, see Sections{\nobreakspace}\ref{Accessing Record Elements} and{\nobreakspace}\ref{Record Assignment}. 

 
\begin{Verbatim}[commandchars=!@|,fontsize=\small,frame=single,label=Example]
  !gapprompt@gap>| !gapinput@r := rec( a := 1, b := 2 );;|
  !gapprompt@gap>| !gapinput@Set(RecNames( r )); # 'Set' because ordering depends on GAP session|
  [ "a", "b" ]
  !gapprompt@gap>| !gapinput@Set(RecNames( r ), x -> r.(x));|
  [ 1, 2 ]
\end{Verbatim}
 

 Note that it is most efficient to assign components to a (new or old) record
in the order given by \texttt{RecNames}. }

 }

  
\section{\textcolor{Chapter }{Accessing Record Elements}}\label{Accessing Record Elements}
\logpage{[ 29, 2, 0 ]}
\hyperdef{L}{X7EAAE25D7A17F778}{}
{
  \index{accessing!record elements} \index{record!component access} \texttt{\mbox{\texttt{\mdseries\slshape r}}.\mbox{\texttt{\mdseries\slshape name}}} 

 The above construct evaluates to the value of the record component with the
name \mbox{\texttt{\mdseries\slshape name}} in the record \mbox{\texttt{\mdseries\slshape r}}. Note that the \mbox{\texttt{\mdseries\slshape name}} is not evaluated, i.e. it is taken literal. 

 
\begin{Verbatim}[commandchars=!@|,fontsize=\small,frame=single,label=Example]
  !gapprompt@gap>| !gapinput@r := rec( a := 1, b := 2 );;|
  !gapprompt@gap>| !gapinput@r.a;|
  1
  !gapprompt@gap>| !gapinput@r.b;|
  2
\end{Verbatim}
 

 \index{record!component variable} \texttt{\mbox{\texttt{\mdseries\slshape r}}.(\mbox{\texttt{\mdseries\slshape name}})} 

 This construct is similar to the above construct. The difference is that the
second operand \mbox{\texttt{\mdseries\slshape name}} is evaluated. It must evaluate to a string or an integer otherwise an error is
signalled. The construct then evaluates to the element of the record \mbox{\texttt{\mdseries\slshape r}} whose name is, as a string, equal to \mbox{\texttt{\mdseries\slshape name}}. 

 
\begin{Verbatim}[commandchars=!@|,fontsize=\small,frame=single,label=Example]
  !gapprompt@gap>| !gapinput@old := rec( a := 1, b := 2 );;|
  !gapprompt@gap>| !gapinput@new := rec();|
  rec(  )
  !gapprompt@gap>| !gapinput@for i in RecNames( old ) do|
  !gapprompt@>| !gapinput@ new.(i) := old.(i);|
  !gapprompt@>| !gapinput@od;|
  !gapprompt@gap>| !gapinput@Display( new );|
  rec(
    a := 1,
    b := 2 )
\end{Verbatim}
 }

  
\section{\textcolor{Chapter }{Record Assignment}}\label{Record Assignment}
\logpage{[ 29, 3, 0 ]}
\hyperdef{L}{X806DE3BD78742CA4}{}
{
  \index{assignment!to a record} \index{record!component assignment} \texttt{\mbox{\texttt{\mdseries\slshape r}}.\mbox{\texttt{\mdseries\slshape name}} := \mbox{\texttt{\mdseries\slshape obj}}} 

 The record assignment assigns the object \mbox{\texttt{\mdseries\slshape obj}}, which may be an object of arbitrary type, to the record component with the
name \mbox{\texttt{\mdseries\slshape name}}, which must be an identifier, of the record \mbox{\texttt{\mdseries\slshape r}}. That means that accessing the element with name \mbox{\texttt{\mdseries\slshape name}} of the record \mbox{\texttt{\mdseries\slshape r}} will return \mbox{\texttt{\mdseries\slshape obj}} after this assignment. If the record \mbox{\texttt{\mdseries\slshape r}} has no component with the name \mbox{\texttt{\mdseries\slshape name}}, the record is automatically extended to make room for the new component. 

 
\begin{Verbatim}[commandchars=!@|,fontsize=\small,frame=single,label=Example]
  !gapprompt@gap>| !gapinput@r := rec( a := 1, b := 2 );;|
  !gapprompt@gap>| !gapinput@r.a := 10;;|
  !gapprompt@gap>| !gapinput@Display( r );|
  rec(
    a := 10,
    b := 2 )
  !gapprompt@gap>| !gapinput@r.c := 3;;|
  !gapprompt@gap>| !gapinput@Display( r );|
  rec(
    a := 10,
    b := 2,
    c := 3 )
\end{Verbatim}
 

 Note that assigning to a record changes the record. 

 The function \texttt{IsBound} (\ref{IsBound:for a record component}) can be used to test if a record has a component with a certain name, the
function \texttt{Unbind} (\ref{Unbind:unbind a record component}) can be used to remove a component with a certain name again. 

 
\begin{Verbatim}[commandchars=!@|,fontsize=\small,frame=single,label=Example]
  !gapprompt@gap>| !gapinput@IsBound(r.a);|
  true
  !gapprompt@gap>| !gapinput@IsBound(r.d);|
  false
  !gapprompt@gap>| !gapinput@Unbind(r.b);|
  !gapprompt@gap>| !gapinput@Display( r );|
  rec(
    a := 10,
    c := 3 )
\end{Verbatim}
 

 \index{record!component variable assignment} \texttt{\mbox{\texttt{\mdseries\slshape r}}.(\mbox{\texttt{\mdseries\slshape name}}) := \mbox{\texttt{\mdseries\slshape obj}}} 

 This construct is similar to the above construct. The difference is that the
second operand \mbox{\texttt{\mdseries\slshape name}} is evaluated. It must evaluate to a string or an integer otherwise an error is
signalled. The construct then assigns \mbox{\texttt{\mdseries\slshape obj}} to the record component of the record \mbox{\texttt{\mdseries\slshape r}} whose name is, as a string, equal to \mbox{\texttt{\mdseries\slshape name}}. }

  
\section{\textcolor{Chapter }{Identical Records}}\label{Identical Records}
\logpage{[ 29, 4, 0 ]}
\hyperdef{L}{X86BC2672803863FB}{}
{
  With the record assignment (see \ref{Record Assignment}) it is possible to change a record. This section describes the semantic
consequences of this fact which are essentially the same as for lists
(see{\nobreakspace}\ref{Identical Lists}). 

 
\begin{Verbatim}[commandchars=!@|,fontsize=\small,frame=single,label=Example]
  r := rec( a := 1 );
  r := rec( a := 1, b := 2 );
\end{Verbatim}
 

 The second assignment does not change the first record, instead it assigns a
new record to the variable \texttt{r}. On the other hand, in the following example the record is changed by the
second assignment. 

 
\begin{Verbatim}[commandchars=!@|,fontsize=\small,frame=single,label=Example]
  r := rec( a := 1 );
  r.b := 2;
\end{Verbatim}
 

 To understand the difference first think of a variable as a name for an
object. The important point is that a record can have several names at the
same time. An assignment \texttt{\mbox{\texttt{\mdseries\slshape var}} := \mbox{\texttt{\mdseries\slshape r}}} means in this interpretation that \mbox{\texttt{\mdseries\slshape var}} is a name for the object \mbox{\texttt{\mdseries\slshape r}}. At the end of the following example \texttt{r2} still has the value \texttt{rec( a := 1 )} as this record has not been changed and nothing else has been assigned to \texttt{r2}. 

 
\begin{Verbatim}[commandchars=!@|,fontsize=\small,frame=single,label=Example]
  r1 := rec( a := 1 );
  r2 := r1;
  r1 := rec( a := 1, b := 2 );
\end{Verbatim}
 

 But after the following example the record for which \texttt{r2} is a name has been changed and thus the value of \texttt{r2} is now \texttt{rec( a := 1, b := 2 )}. 

 
\begin{Verbatim}[commandchars=!@|,fontsize=\small,frame=single,label=Example]
  r1 := rec( a := 1 );
  r2 := r1;
  r1.b := 2;
\end{Verbatim}
 

 We shall say that two records are \emph{identical} if changing one of them by a record assignment also changes the other one.
This is slightly incorrect, because if \emph{two} records are identical, there are actually only two names for \emph{one} record. However, the correct usage would be very awkward and would only add to
the confusion. Note that two identical records must be equal, because there is
only one records with two different names. Thus identity is an equivalence
relation that is a refinement of equality. 

 Let us now consider under which circumstances two records are identical. 

 If you enter a record literal then the record denoted by this literal is a new
record that is not identical to any other record. Thus in the following
example \texttt{r1} and \texttt{r2} are not identical, though they are equal of course. 

 
\begin{Verbatim}[commandchars=!@|,fontsize=\small,frame=single,label=Example]
  r1 := rec( a := 1 );
  r2 := rec( a := 1 );
\end{Verbatim}
 

 Also in the following example, no records in the list \texttt{l} are identical. 

 
\begin{Verbatim}[commandchars=!@|,fontsize=\small,frame=single,label=Example]
  l := [];
  for i in [1..10] do
    l[i] := rec( a := 1 );
  od;
\end{Verbatim}
 

 If you assign a record to a variable no new record is created. Thus the record
value of the variable on the left hand side and the record on the right hand
side of the assignment are identical. So in the following example \texttt{r1} and \texttt{r2} are identical records. 

 
\begin{Verbatim}[commandchars=!@|,fontsize=\small,frame=single,label=Example]
  r1 := rec( a := 1 );
  r2 := r1;
\end{Verbatim}
 

 If you pass a record as argument, the old record and the argument of the
function are identical. Also if you return a record from a function, the old
record and the value of the function call are identical. So in the following
example \texttt{r1} and \texttt{r2} are identical records. 

 
\begin{Verbatim}[commandchars=!@|,fontsize=\small,frame=single,label=Example]
  r1 := rec( a := 1 );
  f := function ( r ) return r; end;
  r2 := f( r1 );
\end{Verbatim}
 

 The functions \texttt{StructuralCopy} (\ref{StructuralCopy}) and \texttt{ShallowCopy} (\ref{ShallowCopy}) accept a record and return a new record that is equal to the old record but
that is \emph{not} identical to the old record. The difference between \texttt{StructuralCopy} (\ref{StructuralCopy}) and \texttt{ShallowCopy} (\ref{ShallowCopy}) is that in the case of \texttt{ShallowCopy} (\ref{ShallowCopy}) the corresponding components of the new and the old records will be identical,
whereas in the case of \texttt{StructuralCopy} (\ref{StructuralCopy}) they will only be equal. So in the following example \texttt{r1} and \texttt{r2} are not identical records. 

 
\begin{Verbatim}[commandchars=!@|,fontsize=\small,frame=single,label=Example]
  r1 := rec( a := 1 );
  r2 := ShallowCopy( r1 );
\end{Verbatim}
 

 If you change a record it keeps its identity. Thus if two records are
identical and you change one of them, you also change the other, and they are
still identical afterwards. On the other hand, two records that are not
identical will never become identical if you change one of them. So in the
following example both \texttt{r1} and \texttt{r2} are changed, and are still identical. 

 
\begin{Verbatim}[commandchars=!@|,fontsize=\small,frame=single,label=Example]
  r1 := rec( a := 1 );
  r2 := r1;
  r1.b := 2;
\end{Verbatim}
 

 }

  
\section{\textcolor{Chapter }{Comparisons of Records}}\label{Comparisons of Records}
\logpage{[ 29, 5, 0 ]}
\hyperdef{L}{X83A7E6607B1D63BC}{}
{
  \index{equality!of records} \index{inequality!of records} \texttt{\mbox{\texttt{\mdseries\slshape rec1}} = \mbox{\texttt{\mdseries\slshape rec2}}} 

 \texttt{\mbox{\texttt{\mdseries\slshape rec1}} {\textless}{\textgreater} \mbox{\texttt{\mdseries\slshape rec2}}} 

 Two records are considered equal, if for each component of one record the
other record has a component of the same name with an equal value and vice
versa. 

 
\begin{Verbatim}[commandchars=!@|,fontsize=\small,frame=single,label=Example]
  !gapprompt@gap>| !gapinput@rec( a := 1, b := 2 ) = rec( b := 2, a := 1 );|
  true
  !gapprompt@gap>| !gapinput@rec( a := 1, b := 2 ) = rec( a := 2, b := 1 );|
  false
  !gapprompt@gap>| !gapinput@rec( a := 1 ) = rec( a := 1, b := 2 );|
  false
  !gapprompt@gap>| !gapinput@rec( a := 1 ) = 1;|
  false
\end{Verbatim}
 

 \index{ordering!of records} \texttt{\mbox{\texttt{\mdseries\slshape rec1}} {\textless} \mbox{\texttt{\mdseries\slshape rec2}}} 

 \texttt{\mbox{\texttt{\mdseries\slshape rec1}} {\textless}= \mbox{\texttt{\mdseries\slshape rec2}}} 

 To compare records we imagine that the components of both records are sorted
according to their names (the sorting depends on the \textsf{GAP} session, more precisely the order in which component names were first used).
Then the records are compared lexicographically with unbound elements
considered smaller than anything else. Precisely one record \mbox{\texttt{\mdseries\slshape rec1}} is considered less than another record \mbox{\texttt{\mdseries\slshape rec2}} if \mbox{\texttt{\mdseries\slshape rec2}} has a component with name \mbox{\texttt{\mdseries\slshape name2}} and either \mbox{\texttt{\mdseries\slshape rec1}} has no component with this name or \texttt{\mbox{\texttt{\mdseries\slshape rec1}}.\mbox{\texttt{\mdseries\slshape name2}} {\textless} \mbox{\texttt{\mdseries\slshape rec2}}.\mbox{\texttt{\mdseries\slshape name2}}} and for each component of \mbox{\texttt{\mdseries\slshape rec1}} with name \texttt{\mbox{\texttt{\mdseries\slshape name1}} {\textless} \mbox{\texttt{\mdseries\slshape name2}}} \mbox{\texttt{\mdseries\slshape rec2}} has a component with this name and \texttt{\mbox{\texttt{\mdseries\slshape rec1}}.\mbox{\texttt{\mdseries\slshape name1}} = \mbox{\texttt{\mdseries\slshape rec2}}.\mbox{\texttt{\mdseries\slshape name1}}}. 

 
\begin{Verbatim}[commandchars=!@|,fontsize=\small,frame=single,label=Example]
  !gapprompt@gap>| !gapinput@rec( axy := 1, bxy := 2 ) < rec( bxy := 2, axy := 1 ); # are equal|
  false
  !gapprompt@gap>| !gapinput@rec( axy := 1 ) < rec( axy := 1, bxy := 2 ); # unbound is < 2|
  true
  !gapprompt@gap>| !gapinput@# in new session the .axy components are compared first|
  !gapprompt@gap>| !gapinput@rec( axy := 1, bxy := 2 ) < rec( axy := 2, bxy := 0 ); # 1 < 2|
  true
  !gapprompt@gap>| !gapinput@rec( axy := 1 ) < rec( axy := 0, bxy := 2 ); # 0 < 1|
  false
  !gapprompt@gap>| !gapinput@rec( bxy := 1 ) < rec( bxy := 0, axy := 2 ); # unbound is < 2|
  true
\end{Verbatim}
 }

  
\section{\textcolor{Chapter }{IsBound and Unbind for Records}}\label{IsBound and Unbind for Records}
\logpage{[ 29, 6, 0 ]}
\hyperdef{L}{X79BE8D0E829E7ACE}{}
{
  

\subsection{\textcolor{Chapter }{IsBound (for a record component)}}
\logpage{[ 29, 6, 1 ]}\nobreak
\hyperdef{L}{X7A13E8F87CAAA0AF}{}
{\noindent\textcolor{FuncColor}{$\triangleright$\enspace\texttt{IsBound({\mdseries\slshape r.name})\index{IsBound@\texttt{IsBound}!for a record component}
\label{IsBound:for a record component}
}\hfill{\scriptsize (operation)}}\\


 \texttt{IsBound} returns \texttt{true} if the record \mbox{\texttt{\mdseries\slshape r}} has a component with the name \mbox{\texttt{\mdseries\slshape name}} (which must be an identifier) and \texttt{false} otherwise. \mbox{\texttt{\mdseries\slshape r}} must evaluate to a record, otherwise an error is signalled. 

 
\begin{Verbatim}[commandchars=!@|,fontsize=\small,frame=single,label=Example]
  !gapprompt@gap>| !gapinput@r := rec( a := 1, b := 2 );;|
  !gapprompt@gap>| !gapinput@IsBound( r.a );|
  true
  !gapprompt@gap>| !gapinput@IsBound( r.c );|
  false
\end{Verbatim}
 }

 

\subsection{\textcolor{Chapter }{Unbind (unbind a record component)}}
\logpage{[ 29, 6, 2 ]}\nobreak
\hyperdef{L}{X7CA9AEFE7DB71604}{}
{\noindent\textcolor{FuncColor}{$\triangleright$\enspace\texttt{Unbind({\mdseries\slshape r.name})\index{Unbind@\texttt{Unbind}!unbind a record component}
\label{Unbind:unbind a record component}
}\hfill{\scriptsize (operation)}}\\


 \texttt{Unbind} deletes the component with the name \mbox{\texttt{\mdseries\slshape name}} in the record \mbox{\texttt{\mdseries\slshape r}}. That is, after execution of \texttt{Unbind}, \mbox{\texttt{\mdseries\slshape r}} no longer has a record component with this name. Note that it is not an error
to unbind a nonexisting record component. \mbox{\texttt{\mdseries\slshape r}} must evaluate to a record, otherwise an error is signalled. 

 
\begin{Verbatim}[commandchars=!@|,fontsize=\small,frame=single,label=Example]
  !gapprompt@gap>| !gapinput@r := rec( a := 1, b := 2 );;|
  !gapprompt@gap>| !gapinput@Unbind( r.a ); r;|
  rec( b := 2 )
  !gapprompt@gap>| !gapinput@Unbind( r.c ); r;|
  rec( b := 2 )
\end{Verbatim}
 

 Note that \texttt{IsBound} (\ref{IsBound:for a record component}) and \texttt{Unbind} are special in that they do not evaluate their argument, otherwise \texttt{IsBound} (\ref{IsBound:for a record component}) would always signal an error when it is supposed to return \texttt{false} and there would be no way to tell \texttt{Unbind} which component to remove. }

 }

  
\section{\textcolor{Chapter }{Record Access Operations}}\label{Record Access Operations}
\logpage{[ 29, 7, 0 ]}
\hyperdef{L}{X784897E180815EDA}{}
{
  Internally, record accesses are done using the operations listed in this
section. For the records implemented in the kernel, kernel methods are
provided for all these operations but otherwise it is possible to install
methods for these operations for any object. This permits objects to simulate
record behavior. 

 To save memory, records do not store a list of all component names, but only
numbers identifying the components. There numbers are called \emph{RNams}. \textsf{GAP} keeps a global list of all RNams that are used and provides the functions \texttt{NameRNam} (\ref{NameRNam}) and \texttt{RNamObj} (\ref{RNamObj:for a string}) to translate RNams to strings that give the component names and vice versa. 

\subsection{\textcolor{Chapter }{NameRNam}}
\logpage{[ 29, 7, 1 ]}\nobreak
\hyperdef{L}{X87BF90FA7F7A3B1B}{}
{\noindent\textcolor{FuncColor}{$\triangleright$\enspace\texttt{NameRNam({\mdseries\slshape nr})\index{NameRNam@\texttt{NameRNam}}
\label{NameRNam}
}\hfill{\scriptsize (function)}}\\


 returns a string representing the component name corresponding to the RNam \mbox{\texttt{\mdseries\slshape nr}}. }

 

\subsection{\textcolor{Chapter }{RNamObj (for a string)}}
\logpage{[ 29, 7, 2 ]}\nobreak
\hyperdef{L}{X78199B6B84A017B9}{}
{\noindent\textcolor{FuncColor}{$\triangleright$\enspace\texttt{RNamObj({\mdseries\slshape str})\index{RNamObj@\texttt{RNamObj}!for a string}
\label{RNamObj:for a string}
}\hfill{\scriptsize (function)}}\\
\noindent\textcolor{FuncColor}{$\triangleright$\enspace\texttt{RNamObj({\mdseries\slshape int})\index{RNamObj@\texttt{RNamObj}!for a positive integer}
\label{RNamObj:for a positive integer}
}\hfill{\scriptsize (function)}}\\


 returns a number (the RNam) corresponding to the string \mbox{\texttt{\mdseries\slshape str}}. It is also possible to pass a positive integer \mbox{\texttt{\mdseries\slshape int}} in which case the decimal expansion of \mbox{\texttt{\mdseries\slshape int}} is used as a string. 

 
\begin{Verbatim}[commandchars=!@|,fontsize=\small,frame=single,label=Example]
  !gapprompt@gap>| !gapinput@NameRNam(798);|
  "BravaisSupergroups"
  !gapprompt@gap>| !gapinput@RNamObj("blubberflutsch");|
  2075
  !gapprompt@gap>| !gapinput@NameRNam(last);|
  "blubberflutsch"
\end{Verbatim}
 

 The correspondence between strings and RNams is not predetermined ab initio,
but RNams are assigned to component names dynamically on a ``first come, first serve'' basis. Therefore, depending on the version of the library you are using and on
the assignments done so far, the \emph{same} component name may be represented by \emph{different} RNams in different \textsf{GAP} sessions. }

 

\subsection{\textcolor{Chapter }{\texttt{\symbol{92}}.}}
\logpage{[ 29, 7, 3 ]}\nobreak
\hyperdef{L}{X7821AC097821AC09}{}
{\noindent\textcolor{FuncColor}{$\triangleright$\enspace\texttt{\texttt{\symbol{92}}.({\mdseries\slshape obj, rnam})\index{\texttt{\symbol{92}}.@\texttt{\texttt{\symbol{92}}.}}
\label{bSlash.}
}\hfill{\scriptsize (operation)}}\\
\noindent\textcolor{FuncColor}{$\triangleright$\enspace\texttt{IsBound\texttt{\symbol{92}}.({\mdseries\slshape obj, rnam})\index{IsBound\texttt{\symbol{92}}.@\texttt{IsBound\texttt{\symbol{92}}.}}
\label{IsBoundbSlash.}
}\hfill{\scriptsize (operation)}}\\
\noindent\textcolor{FuncColor}{$\triangleright$\enspace\texttt{\texttt{\symbol{92}}.\texttt{\symbol{92}}:\texttt{\symbol{92}}=({\mdseries\slshape obj, rnam, val})\index{\texttt{\symbol{92}}.\texttt{\symbol{92}}:\texttt{\symbol{92}}=@\texttt{\texttt{\symbol{92}}.\texttt{\symbol{92}}:\texttt{\symbol{92}}=}}
\label{bSlash.bSlash:bSlash=}
}\hfill{\scriptsize (operation)}}\\
\noindent\textcolor{FuncColor}{$\triangleright$\enspace\texttt{Unbind\texttt{\symbol{92}}.({\mdseries\slshape obj, rnam})\index{Unbind\texttt{\symbol{92}}.@\texttt{Unbind\texttt{\symbol{92}}.}}
\label{UnbindbSlash.}
}\hfill{\scriptsize (operation)}}\\


 \index{record component!operation} \index{record boundness test!operation} \index{record assignment!operation} \index{record unbind!operation} These operations are called for record accesses to arbitrary objects. If
applicable methods are installed, they are called when the object is accessed
as a record. 

 For records, the operations implement component access, test for element
boundness, component assignment and removal of the component represented by
the RNam \mbox{\texttt{\mdseries\slshape rnam}}. 

 The component identifier \mbox{\texttt{\mdseries\slshape rnam}} is always required to be in \texttt{IsPosInt} (\ref{IsPosInt}). 

 
\begin{Verbatim}[commandchars=!@|,fontsize=\small,frame=single,label=Example]
  !gapprompt@gap>| !gapinput@r:= rec( a:= 1 );;|
  !gapprompt@gap>| !gapinput@IsBound\.( r, RNamObj( "a" ) );|
  true
  !gapprompt@gap>| !gapinput@\.( r, RNamObj( "a" ) );|
  1
  !gapprompt@gap>| !gapinput@IsBound\.( r, RNamObj( "b" ) );|
  false
  !gapprompt@gap>| !gapinput@\.\:\=( r, RNamObj( "b" ), 2 );|
  !gapprompt@gap>| !gapinput@r;|
  rec( a := 1, b := 2 )
  !gapprompt@gap>| !gapinput@Unbind\.( r, RNamObj( "b" ) );|
  !gapprompt@gap>| !gapinput@r;|
  rec( a := 1 )
  !gapprompt@gap>| !gapinput@G:= SymmetricGroup( 4 );;|
  !gapprompt@gap>| !gapinput@G.1;|
  (1,2,3,4)
  !gapprompt@gap>| !gapinput@\.( G, RNamObj( 1 ) );|
  (1,2,3,4)
  !gapprompt@gap>| !gapinput@meth:= ApplicableMethod( \., [ G, 4711 ] );;|
  !gapprompt@gap>| !gapinput@meth( G, RNamObj( 1 ) );|
  (1,2,3,4)
\end{Verbatim}
 }

 }

 }

         
\chapter{\textcolor{Chapter }{Collections}}\label{Collections}
\logpage{[ 30, 0, 0 ]}
\hyperdef{L}{X8050A8037984E5B6}{}
{
  A \emph{collection} in \textsf{GAP} consists of elements in the same family (see{\nobreakspace}\ref{Families}). The most important kinds of collections are \emph{homogeneous lists} (see{\nobreakspace}\ref{Lists}) and \emph{domains} (see{\nobreakspace}\ref{Domains}). Note that a list is never a domain, and a domain is never a list. A list is
a collection if and only if it is nonempty and homogeneous. 

 Basic operations for collections are \texttt{Size} (\ref{Size}) and \texttt{Enumerator} (\ref{Enumerator}); for \emph{finite} collections, \texttt{Enumerator} (\ref{Enumerator}) admits to delegate the other operations for collections (see{\nobreakspace}\ref{Attributes and Properties for Collections} and{\nobreakspace}\ref{Operations for Collections}) to functions for lists (see{\nobreakspace}\ref{Lists}). Obviously, special methods depending on the arguments are needed for the
computation of e.g.{\nobreakspace}the intersection of two \emph{infinite} domains.  
\section{\textcolor{Chapter }{IsCollection (Filter)}}\label{sect:IsCollection}
\logpage{[ 30, 1, 0 ]}
\hyperdef{L}{X8084F03A78ABD4F8}{}
{
  

\subsection{\textcolor{Chapter }{IsCollection}}
\logpage{[ 30, 1, 1 ]}\nobreak
\hyperdef{L}{X79C9FC7F86E2738C}{}
{\noindent\textcolor{FuncColor}{$\triangleright$\enspace\texttt{IsCollection({\mdseries\slshape obj})\index{IsCollection@\texttt{IsCollection}}
\label{IsCollection}
}\hfill{\scriptsize (Category)}}\\


 tests whether an object is a collection. 

 Some of the functions for lists and collections are described in the chapter
about lists, mainly in Section{\nobreakspace}\ref{Operations for Lists}. In the current chapter, we describe those functions for which the ``collection aspect'' seems to be more important than the ``list aspect''. }

 }

  
\section{\textcolor{Chapter }{Collection Families}}\label{Collection Families}
\logpage{[ 30, 2, 0 ]}
\hyperdef{L}{X85D8D8F684B02DDF}{}
{
  

\subsection{\textcolor{Chapter }{CollectionsFamily}}
\logpage{[ 30, 2, 1 ]}\nobreak
\hyperdef{L}{X84E5A67E87D8DD66}{}
{\noindent\textcolor{FuncColor}{$\triangleright$\enspace\texttt{CollectionsFamily({\mdseries\slshape Fam})\index{CollectionsFamily@\texttt{CollectionsFamily}}
\label{CollectionsFamily}
}\hfill{\scriptsize (attribute)}}\\


 For a family \mbox{\texttt{\mdseries\slshape Fam}}, \texttt{CollectionsFamily} returns the family of all collections over \mbox{\texttt{\mdseries\slshape Fam}}, that is, of all dense lists and domains that consist of objects in \mbox{\texttt{\mdseries\slshape Fam}}. 

 The \texttt{NewFamily} (\ref{NewFamily}) call in the standard method of \texttt{CollectionsFamily} is executed with second argument \texttt{IsCollection} (\ref{IsCollection}), since every object in the collections family must be a collection, and with
third argument the collections categories of the involved categories in the
implied filter of \mbox{\texttt{\mdseries\slshape Fam}}. 

 Note that families (see{\nobreakspace}\ref{Families}) are used to describe relations between objects. Important such relations are
that between an element $e$ and each collection of elements that lie in the same family as $e$, and that between two collections whose elements lie in the same family.
Therefore, all collections of elements in the family \mbox{\texttt{\mdseries\slshape Fam}} form the new family \texttt{CollectionsFamily( \mbox{\texttt{\mdseries\slshape Fam}} )}. }

 

\subsection{\textcolor{Chapter }{IsCollectionFamily}}
\logpage{[ 30, 2, 2 ]}\nobreak
\hyperdef{L}{X856AC2DF7F7CBAAF}{}
{\noindent\textcolor{FuncColor}{$\triangleright$\enspace\texttt{IsCollectionFamily({\mdseries\slshape obj})\index{IsCollectionFamily@\texttt{IsCollectionFamily}}
\label{IsCollectionFamily}
}\hfill{\scriptsize (Category)}}\\


 is \texttt{true} if \mbox{\texttt{\mdseries\slshape Fam}} is a family of collections, and \texttt{false} otherwise. }

 

\subsection{\textcolor{Chapter }{ElementsFamily}}
\logpage{[ 30, 2, 3 ]}\nobreak
\hyperdef{L}{X864BB3748546F63F}{}
{\noindent\textcolor{FuncColor}{$\triangleright$\enspace\texttt{ElementsFamily({\mdseries\slshape Fam})\index{ElementsFamily@\texttt{ElementsFamily}}
\label{ElementsFamily}
}\hfill{\scriptsize (attribute)}}\\


 If \mbox{\texttt{\mdseries\slshape Fam}} is a collections family (see{\nobreakspace}\texttt{IsCollectionFamily} (\ref{IsCollectionFamily})) then \texttt{ElementsFamily} returns the family from which \mbox{\texttt{\mdseries\slshape Fam}} was created by \texttt{CollectionsFamily} (\ref{CollectionsFamily}). The way a collections family is created, it always has its elements family
stored. If \mbox{\texttt{\mdseries\slshape Fam}} is not a collections family then an error is signalled. 

 
\begin{Verbatim}[commandchars=!@|,fontsize=\small,frame=single,label=Example]
  !gapprompt@gap>| !gapinput@fam:= FamilyObj( (1,2) );;|
  !gapprompt@gap>| !gapinput@collfam:= CollectionsFamily( fam );;|
  !gapprompt@gap>| !gapinput@fam = collfam;  fam = ElementsFamily( collfam );|
  false
  true
  !gapprompt@gap>| !gapinput@collfam = FamilyObj( [ (1,2,3) ] );|
  true
  !gapprompt@gap>| !gapinput@collfam = FamilyObj( Group( () ) );|
  true
  !gapprompt@gap>| !gapinput@collfam = CollectionsFamily( collfam );|
  false
\end{Verbatim}
 }

 

\subsection{\textcolor{Chapter }{CategoryCollections}}
\logpage{[ 30, 2, 4 ]}\nobreak
\hyperdef{L}{X78C38017804B2EA7}{}
{\noindent\textcolor{FuncColor}{$\triangleright$\enspace\texttt{CategoryCollections({\mdseries\slshape filter})\index{CategoryCollections@\texttt{CategoryCollections}}
\label{CategoryCollections}
}\hfill{\scriptsize (function)}}\\


 Let \mbox{\texttt{\mdseries\slshape filter}} be a filter that is \texttt{true} for all elements of a family \mbox{\texttt{\mdseries\slshape Fam}}, by the construction of \mbox{\texttt{\mdseries\slshape Fam}}. Then \texttt{CategoryCollections} returns the \emph{collections category} of \mbox{\texttt{\mdseries\slshape filter}}. This is a category that is \texttt{true} for all elements in \texttt{CollectionsFamily( \mbox{\texttt{\mdseries\slshape Fam}} )}. 

 For example, the construction of \texttt{PermutationsFamily} (\ref{PermutationsFamily}) guarantees that each of its elements lies in the filter \texttt{IsPerm} (\ref{IsPerm}), and each collection of permutations (permutation group or dense list of
permutations) lies in the category \texttt{CategoryCollections( IsPerm )}. \texttt{CategoryCollections( IsPerm )}. Note that this works only if the collections category is created \emph{before} the collections family. So it is necessary to construct interesting
collections categories immediately after the underlying category has been
created. }

 

\subsection{\textcolor{Chapter }{DeclareCategoryCollections}}
\logpage{[ 30, 2, 5 ]}\nobreak
\hyperdef{L}{X85ABC1B4829778C7}{}
{\noindent\textcolor{FuncColor}{$\triangleright$\enspace\texttt{DeclareCategoryCollections({\mdseries\slshape name})\index{DeclareCategoryCollections@\texttt{DeclareCategoryCollections}}
\label{DeclareCategoryCollections}
}\hfill{\scriptsize (function)}}\\


 Calls \texttt{CategoryCollections} (\ref{CategoryCollections}) on the category that is bound to the global variable with name \mbox{\texttt{\mdseries\slshape name}} to obtain its collections category, and binds the latter to the global
variable with name \texttt{nname}. This name is defined as follows: If \mbox{\texttt{\mdseries\slshape name}} is of the form \texttt{\mbox{\texttt{\mdseries\slshape Something}}Collection} then \texttt{nname} is set to \texttt{\mbox{\texttt{\mdseries\slshape Something}}CollColl}, if \mbox{\texttt{\mdseries\slshape name}} is of the form \texttt{\mbox{\texttt{\mdseries\slshape Something}}Coll} then \texttt{nname} is set to \texttt{\mbox{\texttt{\mdseries\slshape Something}}CollColl}, otherwise we set \texttt{nname} to \texttt{\mbox{\texttt{\mdseries\slshape name}}Collection}. }

 }

  
\section{\textcolor{Chapter }{Lists and Collections}}\label{Lists and Collections}
\logpage{[ 30, 3, 0 ]}
\hyperdef{L}{X7C3722DF8736FFDB}{}
{
  \index{Sorted Lists as Collections} The following functions take a \emph{list or collection} as argument, and return a corresponding \emph{list}. They differ in whether or not the result is mutable or immutable
(see{\nobreakspace}\ref{Mutability and Copyability}), guaranteed to be sorted, or guaranteed to admit list access in constant
time (see{\nobreakspace}\texttt{IsConstantTimeAccessList} (\ref{IsConstantTimeAccessList})). 

\subsection{\textcolor{Chapter }{IsListOrCollection}}
\logpage{[ 30, 3, 1 ]}\nobreak
\hyperdef{L}{X877128A77826DD69}{}
{\noindent\textcolor{FuncColor}{$\triangleright$\enspace\texttt{IsListOrCollection({\mdseries\slshape obj})\index{IsListOrCollection@\texttt{IsListOrCollection}}
\label{IsListOrCollection}
}\hfill{\scriptsize (Category)}}\\


 Several functions are defined for both lists and collections, for example \texttt{Intersection} (\ref{Intersection:for a list}), \texttt{Iterator} (\ref{Iterator}), and \texttt{Random} (\ref{Random:for a list or collection}). \texttt{IsListOrCollection} is a supercategory of \texttt{IsList} (\ref{IsList}) and \texttt{IsCollection} (\ref{IsCollection}) (that is, all lists and collections lie in this category), which is used to
describe the arguments of functions such as the ones listed above. }

 

\subsection{\textcolor{Chapter }{Enumerator}}
\logpage{[ 30, 3, 2 ]}\nobreak
\hyperdef{L}{X7EF8910F82B45EC7}{}
{\noindent\textcolor{FuncColor}{$\triangleright$\enspace\texttt{Enumerator({\mdseries\slshape listorcoll})\index{Enumerator@\texttt{Enumerator}}
\label{Enumerator}
}\hfill{\scriptsize (attribute)}}\\


 \texttt{Enumerator} returns an immutable list $enum$. If the argument is a list (which may contain holes), then \texttt{Length( }$enum$\texttt{ )} is the length of this list, and $enum$ contains the elements (and holes) of this list in the same order. If the
argument is a collection that is not a list, then \texttt{Length( }$enum$\texttt{ )} is the number of different elements of \mbox{\texttt{\mdseries\slshape C}}, and $enum$ contains the different elements of the collection in an unspecified order,
which may change for repeated calls of \texttt{Enumerator}. $enum[pos]$ may not execute in constant time (see{\nobreakspace}\texttt{IsConstantTimeAccessList} (\ref{IsConstantTimeAccessList})), and the size of $enum$ in memory is as small as is feasible. 

 For lists, the default method is \texttt{Immutable} (\ref{Immutable}). For collections that are not lists, there is no default method. }

 

\subsection{\textcolor{Chapter }{EnumeratorSorted}}
\logpage{[ 30, 3, 3 ]}\nobreak
\hyperdef{L}{X80CD7DDC7D0C60D5}{}
{\noindent\textcolor{FuncColor}{$\triangleright$\enspace\texttt{EnumeratorSorted({\mdseries\slshape listorcoll})\index{EnumeratorSorted@\texttt{EnumeratorSorted}}
\label{EnumeratorSorted}
}\hfill{\scriptsize (attribute)}}\\


 \texttt{EnumeratorSorted} returns an immutable list $enum$. The argument must be a collection or a list \mbox{\texttt{\mdseries\slshape listorcoll}} which may contain holes but whose elements lie in the same family
(see{\nobreakspace}\ref{Families}). \texttt{Length( }$enum$\texttt{ )} is the number of different elements of the argument, and $enum$ contains the different elements in sorted order, w.r.t.{\nobreakspace}\texttt{{\textless}}. $enum[pos]$ may not execute in constant time (see{\nobreakspace}\texttt{IsConstantTimeAccessList} (\ref{IsConstantTimeAccessList})), and the size of $enum$ in memory is as small as is feasible. 

 
\begin{Verbatim}[commandchars=!@|,fontsize=\small,frame=single,label=Example]
  !gapprompt@gap>| !gapinput@Enumerator( [ 1, 3,, 2 ] );|
  [ 1, 3,, 2 ]
  !gapprompt@gap>| !gapinput@enum:= Enumerator( Rationals );;  elm:= enum[ 10^6 ];|
  -69/907
  !gapprompt@gap>| !gapinput@Position( enum, elm );|
  1000000
  !gapprompt@gap>| !gapinput@IsMutable( enum );  IsSortedList( enum );|
  false
  false
  !gapprompt@gap>| !gapinput@IsConstantTimeAccessList( enum );|
  false
  !gapprompt@gap>| !gapinput@EnumeratorSorted( [ 1, 3,, 2 ] );|
  [ 1, 2, 3 ]
\end{Verbatim}
 }

  

\subsection{\textcolor{Chapter }{EnumeratorByFunctions (for a domain and a record)}}
\logpage{[ 30, 3, 4 ]}\nobreak
\hyperdef{L}{X85E149177AC547C3}{}
{\noindent\textcolor{FuncColor}{$\triangleright$\enspace\texttt{EnumeratorByFunctions({\mdseries\slshape D, record})\index{EnumeratorByFunctions@\texttt{EnumeratorByFunctions}!for a domain and a record}
\label{EnumeratorByFunctions:for a domain and a record}
}\hfill{\scriptsize (function)}}\\
\noindent\textcolor{FuncColor}{$\triangleright$\enspace\texttt{EnumeratorByFunctions({\mdseries\slshape Fam, record})\index{EnumeratorByFunctions@\texttt{EnumeratorByFunctions}!for a family and a record}
\label{EnumeratorByFunctions:for a family and a record}
}\hfill{\scriptsize (function)}}\\


 \texttt{EnumeratorByFunctions} returns an immutable, dense, and duplicate\texttt{\symbol{45}}free list $enum$ for which \texttt{IsBound} (\ref{IsBound:for a list index}), element access via \texttt{\texttt{\symbol{92}}[\texttt{\symbol{92}}]} (\ref{bSlash[bSlash]}), \texttt{Length} (\ref{Length}), and \texttt{Position} (\ref{Position}) are computed via prescribed functions. 

 Let \mbox{\texttt{\mdseries\slshape record}} be a record with at least the following components. 
\begin{description}
\item[{\texttt{ElementNumber}}]  a function taking two arguments \mbox{\texttt{\mdseries\slshape enum}} and \mbox{\texttt{\mdseries\slshape pos}}, which returns \texttt{\mbox{\texttt{\mdseries\slshape enum}}[ \mbox{\texttt{\mdseries\slshape pos}} ]} (see{\nobreakspace}\ref{Basic Operations for Lists}); it can be assumed that the argument \mbox{\texttt{\mdseries\slshape pos}} is a positive integer, but \mbox{\texttt{\mdseries\slshape pos}} may be larger than the length of \mbox{\texttt{\mdseries\slshape enum}} (in which case an error must be signalled); note that the result must be
immutable since \mbox{\texttt{\mdseries\slshape enum}} itself is immutable, 
\item[{\texttt{NumberElement}}]  a function taking two arguments \mbox{\texttt{\mdseries\slshape enum}} and \mbox{\texttt{\mdseries\slshape elm}}, which returns \texttt{Position( \mbox{\texttt{\mdseries\slshape enum}}, \mbox{\texttt{\mdseries\slshape elm}} )} (see{\nobreakspace}\texttt{Position} (\ref{Position})); it cannot be assumed that \mbox{\texttt{\mdseries\slshape elm}} is really contained in \mbox{\texttt{\mdseries\slshape enum}} (and \texttt{fail} must be returned if not); note that for the three argument version of \texttt{Position} (\ref{Position}), the method that is available for duplicate\texttt{\symbol{45}}free lists
suffices. 
\end{description}
 

 Further (data) components may be contained in \mbox{\texttt{\mdseries\slshape record}} which can be used by these function. 

 If the first argument is a domain \mbox{\texttt{\mdseries\slshape D}} then \mbox{\texttt{\mdseries\slshape enum}} lists the elements of \mbox{\texttt{\mdseries\slshape D}} (in general \mbox{\texttt{\mdseries\slshape enum}} is \emph{not} sorted), and methods for \texttt{Length} (\ref{Length}), \texttt{IsBound} (\ref{IsBound:for a list index}), and \texttt{PrintObj} (\ref{PrintObj}) may use \mbox{\texttt{\mdseries\slshape D}}.  

 If one wants to describe the result without creating a domain then the
elements are given implicitly by the functions in \mbox{\texttt{\mdseries\slshape record}}, and the first argument must be a family \mbox{\texttt{\mdseries\slshape Fam}} which will become the family of \mbox{\texttt{\mdseries\slshape enum}}; if \mbox{\texttt{\mdseries\slshape enum}} is not homogeneous then \mbox{\texttt{\mdseries\slshape Fam}} must be \texttt{ListsFamily}, otherwise it must be the collections family of any element in \mbox{\texttt{\mdseries\slshape enum}}. In this case, additionally the following component in \mbox{\texttt{\mdseries\slshape record}} is needed. 

 
\begin{description}
\item[{\texttt{Length}}]  a function taking the argument \mbox{\texttt{\mdseries\slshape enum}}, which returns the length of \mbox{\texttt{\mdseries\slshape enum}} (see{\nobreakspace}\texttt{Length} (\ref{Length})). 
\end{description}
 

 The following components are optional; they are used if they are present but
default methods are installed for the case that they are missing. 
\begin{description}
\item[{\texttt{IsBound\texttt{\symbol{92}}[\texttt{\symbol{92}}]}}]  a function taking two arguments \mbox{\texttt{\mdseries\slshape enum}} and \mbox{\texttt{\mdseries\slshape k}}, which returns \texttt{IsBound( \mbox{\texttt{\mdseries\slshape enum}}[ \mbox{\texttt{\mdseries\slshape k}} ] )} (see{\nobreakspace}\ref{Basic Operations for Lists}); if this component is missing then \texttt{Length} (\ref{Length}) is used for computing the result, 
\item[{\texttt{Membership}}]  a function taking two arguments \mbox{\texttt{\mdseries\slshape elm}} and \mbox{\texttt{\mdseries\slshape enum}}, which returns \texttt{true} is \mbox{\texttt{\mdseries\slshape elm}} is an element of \mbox{\texttt{\mdseries\slshape enum}}, and \texttt{false} otherwise (see{\nobreakspace}\ref{Basic Operations for Lists}); if this component is missing then \texttt{NumberElement} is used for computing the result, 
\item[{\texttt{AsList}}]  a function taking one argument \mbox{\texttt{\mdseries\slshape enum}}, which returns a list with the property that the access to each of its
elements will take roughly the same time (see{\nobreakspace}\texttt{IsConstantTimeAccessList} (\ref{IsConstantTimeAccessList})); if this component is missing then \texttt{ConstantTimeAccessList} (\ref{ConstantTimeAccessList}) is used for computing the result, 
\item[{\texttt{ViewObj} and \texttt{PrintObj}}]  two functions that print what one wants to be printed when \texttt{View( \mbox{\texttt{\mdseries\slshape enum}} )} or \texttt{Print( \mbox{\texttt{\mdseries\slshape enum}} )} is called (see{\nobreakspace}\ref{View and Print}), if the \texttt{ViewObj} component is missing then the \texttt{PrintObj} method is used as a default. 
\end{description}
 

 If the result is known to have additional properties such as being strictly
sorted (see{\nobreakspace}\texttt{IsSSortedList} (\ref{IsSSortedList})) then it can be useful to set these properties after the construction of the
enumerator, before it is used for the first time. And in the case that a new
sorted enumerator of a domain is implemented via \texttt{EnumeratorByFunctions}, and this construction is installed as a method for the operation \texttt{Enumerator} (\ref{Enumerator}), then it should be installed also as a method for \texttt{EnumeratorSorted} (\ref{EnumeratorSorted}). 

 Note that it is \emph{not} checked that \texttt{EnumeratorByFunctions} really returns a dense and duplicate\texttt{\symbol{45}}free list. \texttt{EnumeratorByFunctions} does \emph{not} make a shallow copy of \mbox{\texttt{\mdseries\slshape record}}, this record is changed in place, see{\nobreakspace}\ref{Creating Objects}. 

 It would be easy to implement a slightly generalized setup for enumerators
that need not be duplicate\texttt{\symbol{45}}free (where the three argument
version of \texttt{Position} (\ref{Position}) is supported), but the resulting overhead for the methods seems not to be
justified. }

 

\subsection{\textcolor{Chapter }{List (for a collection)}}
\logpage{[ 30, 3, 5 ]}\nobreak
\hyperdef{L}{X7F12F40E87F3C3A7}{}
{\noindent\textcolor{FuncColor}{$\triangleright$\enspace\texttt{List({\mdseries\slshape C})\index{List@\texttt{List}!for a collection}
\label{List:for a collection}
}\hfill{\scriptsize (function)}}\\


 For a collection \mbox{\texttt{\mdseries\slshape C}} (see{\nobreakspace}\ref{Collections}) that is not a list, \texttt{List} returns a new mutable list \mbox{\texttt{\mdseries\slshape new}} such that \texttt{Length( \mbox{\texttt{\mdseries\slshape new}} )} is the number of different elements of \mbox{\texttt{\mdseries\slshape C}}, and \mbox{\texttt{\mdseries\slshape new}} contains the different elements of \mbox{\texttt{\mdseries\slshape C}} in an unspecified order which may change for repeated calls. \texttt{\mbox{\texttt{\mdseries\slshape new}}[\mbox{\texttt{\mdseries\slshape pos}}]} executes in constant time (see{\nobreakspace}\texttt{IsConstantTimeAccessList} (\ref{IsConstantTimeAccessList})), and the size of \mbox{\texttt{\mdseries\slshape new}} is proportional to its length. The generic method for this case is \texttt{ShallowCopy( Enumerator( \mbox{\texttt{\mdseries\slshape C}} ) )}.   

 Developers who wish to adapt this for custom list or collection types need to
install suitable methods for the operation \texttt{ListOp}. \index{ListOp@\texttt{ListOp}} 

 
\begin{Verbatim}[commandchars=!@|,fontsize=\small,frame=single,label=Example]
  !gapprompt@gap>| !gapinput@l:= List( Group( (1,2,3) ) );|
  [ (), (1,3,2), (1,2,3) ]
  !gapprompt@gap>| !gapinput@IsMutable( l );  IsSortedList( l );  IsConstantTimeAccessList( l );|
  true
  false
  true
\end{Verbatim}
 

 (See also \texttt{List} (\ref{List:for a list (and a function)}).) }

 

\subsection{\textcolor{Chapter }{SortedList}}
\logpage{[ 30, 3, 6 ]}\nobreak
\hyperdef{L}{X82CE157A7FAD8036}{}
{\noindent\textcolor{FuncColor}{$\triangleright$\enspace\texttt{SortedList({\mdseries\slshape listorcoll[, func]})\index{SortedList@\texttt{SortedList}}
\label{SortedList}
}\hfill{\scriptsize (operation)}}\\
\noindent\textcolor{FuncColor}{$\triangleright$\enspace\texttt{SortedListBy({\mdseries\slshape listorcoll, func})\index{SortedListBy@\texttt{SortedListBy}}
\label{SortedListBy}
}\hfill{\scriptsize (operation)}}\\


 \texttt{SortedList} returns a new mutable and dense list \mbox{\texttt{\mdseries\slshape new}}. The argument must be a collection or list \mbox{\texttt{\mdseries\slshape listorcoll}} which may contain holes but whose elements lie in the same family
(see{\nobreakspace}\ref{Families}). \texttt{Length( \mbox{\texttt{\mdseries\slshape new}} )} is the number of elements of \mbox{\texttt{\mdseries\slshape listorcoll}}, and \mbox{\texttt{\mdseries\slshape new}} contains the elements in sorted order, w.r.t.{\nobreakspace}\texttt{{\textless}} or \mbox{\texttt{\mdseries\slshape func}} if it is specified. For details, please refer to \texttt{Sort} (\ref{Sort}). 

 \texttt{SortedListBy} also returns a new mutable and dense list \mbox{\texttt{\mdseries\slshape new}}. The first argument must be a collection or list \mbox{\texttt{\mdseries\slshape listorcoll}} which may contain holes but whose elements lie in the same family
(see{\nobreakspace}\ref{Families}). \texttt{Length( \mbox{\texttt{\mdseries\slshape new}} )} is the number of elements of \mbox{\texttt{\mdseries\slshape listorcoll}}, and \mbox{\texttt{\mdseries\slshape new}} has elements in an order such that \texttt{func(new[i]) {\textless}= func(new[i+1])} for all relevant \mbox{\texttt{\mdseries\slshape i}}. \mbox{\texttt{\mdseries\slshape func}} must thus be a function on one argument which returns values that can be
compared. Each \texttt{func(listorcoll[i])} is computed just once and stored, making this more efficient than using the
two\texttt{\symbol{45}}argument version of \texttt{SortedList} in many cases. 

 In both cases, \texttt{\mbox{\texttt{\mdseries\slshape new}}[\mbox{\texttt{\mdseries\slshape pos}}]} executes in constant time (see{\nobreakspace}\texttt{IsConstantTimeAccessList} (\ref{IsConstantTimeAccessList})), and the size of \mbox{\texttt{\mdseries\slshape new}} in memory is proportional to its length. 

 
\begin{Verbatim}[commandchars=!@|,fontsize=\small,frame=single,label=Example]
  !gapprompt@gap>| !gapinput@l:= SortedList( Group( (1,2,3) ) );|
  [ (), (1,2,3), (1,3,2) ]
  !gapprompt@gap>| !gapinput@IsMutable( l );  IsSortedList( l );  IsConstantTimeAccessList( l );|
  true
  true
  true
  !gapprompt@gap>| !gapinput@SortedList( [ 1, 2, 1,, 3, 2 ] );|
  [ 1, 1, 2, 2, 3 ]
  !gapprompt@gap>| !gapinput@SortedListBy(SymmetricGroup(3), Order);|
  [ (), (1,3), (2,3), (1,2), (1,2,3), (1,3,2) ]
  !gapprompt@gap>| !gapinput@SortedListBy( [ 3, 1,,, -6, 5, -2,, -5, 3, 3,,, -1 ], x -> x^2 );|
  [ 1, -1, -2, 3, 3, 3, 5, -5, -6 ]
  !gapprompt@gap>| !gapinput@l:=SortedListBy( [ 2, 1,, 1, -5,, 3, -5 ], x -> x^3);|
  [ -5, -5, 1, 1, 2, 3 ]
  !gapprompt@gap>| !gapinput@IsMutable( l );  IsSortedList( l );  IsConstantTimeAccessList( l );|
  true
  true
  true
\end{Verbatim}
 }

 

\subsection{\textcolor{Chapter }{SSortedList}}
\logpage{[ 30, 3, 7 ]}\nobreak
\hyperdef{L}{X7E399AC97FD98217}{}
{\noindent\textcolor{FuncColor}{$\triangleright$\enspace\texttt{SSortedList({\mdseries\slshape listorcoll[, fun]})\index{SSortedList@\texttt{SSortedList}}
\label{SSortedList}
}\hfill{\scriptsize (operation)}}\\
\noindent\textcolor{FuncColor}{$\triangleright$\enspace\texttt{Set({\mdseries\slshape C[, fun]})\index{Set@\texttt{Set}}
\label{Set}
}\hfill{\scriptsize (operation)}}\\


 \texttt{SSortedList} (``strictly sorted list'') returns a new dense, mutable, and duplicate free list \mbox{\texttt{\mdseries\slshape new}}. The argument must be a collection or list \mbox{\texttt{\mdseries\slshape listorcoll}} which may contain holes. 

 If the optional argument \mbox{\texttt{\mdseries\slshape fun}} is not given then \texttt{Length( \mbox{\texttt{\mdseries\slshape new}} )} is the number of different elements of \mbox{\texttt{\mdseries\slshape listorcoll}}, and \mbox{\texttt{\mdseries\slshape new}} contains the different elements in strictly sorted order,
w.r.t.{\nobreakspace}\texttt{\texttt{\symbol{92}}{\textless}} (\ref{bSlash<}). For that, any two entries of \mbox{\texttt{\mdseries\slshape listorcoll}} must be comparable via \texttt{\texttt{\symbol{92}}{\textless}} (\ref{bSlash<}). (Typically, the entries lie in the same family, see{\nobreakspace}\ref{Families}.) 

 If \mbox{\texttt{\mdseries\slshape fun}} is given then it must be a unary function. In this case, \mbox{\texttt{\mdseries\slshape fun}} is applied to all elements of \mbox{\texttt{\mdseries\slshape listorcoll}}, \mbox{\texttt{\mdseries\slshape new}} contains the different return values in strictly sorted order, and \texttt{Length( \mbox{\texttt{\mdseries\slshape new}} )} is the number of different such values. For that, any two return values must
be comparable via \texttt{\texttt{\symbol{92}}{\textless}} (\ref{bSlash<}). 

 \texttt{\mbox{\texttt{\mdseries\slshape new}}[\mbox{\texttt{\mdseries\slshape pos}}]} executes in constant time (see{\nobreakspace}\texttt{IsConstantTimeAccessList} (\ref{IsConstantTimeAccessList})), and the size of \mbox{\texttt{\mdseries\slshape new}} in memory is proportional to its length. 

 \texttt{Set} is simply a synonym for \texttt{SSortedList}.    

 
\begin{Verbatim}[commandchars=!@|,fontsize=\small,frame=single,label=Example]
  !gapprompt@gap>| !gapinput@l:= SSortedList( Group( (1,2,3) ) );|
  [ (), (1,2,3), (1,3,2) ]
  !gapprompt@gap>| !gapinput@IsMutable( l );  IsSSortedList( l );  IsConstantTimeAccessList( l );|
  true
  true
  true
  !gapprompt@gap>| !gapinput@SSortedList( Group( (1,2,3) ), Order );|
  [ 1, 3 ]
  !gapprompt@gap>| !gapinput@SSortedList( [ 1, 2, 1,, 3, 2 ] );|
  [ 1, 2, 3 ]
  !gapprompt@gap>| !gapinput@SSortedList( [ 1, 2, 1,, 3, 2 ], x -> x^2 );|
  [ 1, 4, 9 ]
\end{Verbatim}
 }

 

\subsection{\textcolor{Chapter }{AsList}}
\logpage{[ 30, 3, 8 ]}\nobreak
\hyperdef{L}{X8289FCCC8274C89D}{}
{\noindent\textcolor{FuncColor}{$\triangleright$\enspace\texttt{AsList({\mdseries\slshape listorcoll})\index{AsList@\texttt{AsList}}
\label{AsList}
}\hfill{\scriptsize (attribute)}}\\


 \texttt{AsList} returns a immutable list \mbox{\texttt{\mdseries\slshape imm}}. If the argument is a list (which may contain holes), then \texttt{Length( \mbox{\texttt{\mdseries\slshape imm}} )} is the \texttt{Length} (\ref{Length}) value of this list, and \mbox{\texttt{\mdseries\slshape imm}} contains the elements (and holes) of the list in the same order. If the
argument is a collection that is not a list, then \texttt{Length( \mbox{\texttt{\mdseries\slshape imm}} )} is the number of different elements of this collection, and \mbox{\texttt{\mdseries\slshape imm}} contains the different elements of the collection in an unspecified order,
which may change for repeated calls of \texttt{AsList}. \texttt{\mbox{\texttt{\mdseries\slshape imm}}[\mbox{\texttt{\mdseries\slshape pos}}]} executes in constant time (see{\nobreakspace}\texttt{IsConstantTimeAccessList} (\ref{IsConstantTimeAccessList})), and the size of \mbox{\texttt{\mdseries\slshape imm}} in memory is proportional to its length. 

 If you expect to do many element tests in the resulting list, it might be
worth to use a sorted list instead, using \texttt{AsSSortedList} (\ref{AsSSortedList}).    

 
\begin{Verbatim}[commandchars=!@|,fontsize=\small,frame=single,label=Example]
  !gapprompt@gap>| !gapinput@l:= AsList( [ 1, 3, 3,, 2 ] );|
  [ 1, 3, 3,, 2 ]
  !gapprompt@gap>| !gapinput@IsMutable( l );  IsSortedList( l );  IsConstantTimeAccessList( l );|
  false
  false
  true
  !gapprompt@gap>| !gapinput@AsList( Group( (1,2,3), (1,2) ) );|
  [ (), (2,3), (1,3,2), (1,3), (1,2,3), (1,2) ]
\end{Verbatim}
 }

 

\subsection{\textcolor{Chapter }{AsSortedList}}
\logpage{[ 30, 3, 9 ]}\nobreak
\hyperdef{L}{X7BCA5C6181391007}{}
{\noindent\textcolor{FuncColor}{$\triangleright$\enspace\texttt{AsSortedList({\mdseries\slshape listorcoll})\index{AsSortedList@\texttt{AsSortedList}}
\label{AsSortedList}
}\hfill{\scriptsize (attribute)}}\\


 \texttt{AsSortedList} returns a dense and immutable list \mbox{\texttt{\mdseries\slshape imm}}. The argument must be a collection or list \mbox{\texttt{\mdseries\slshape listorcoll}} which may contain holes but whose elements lie in the same family
(see{\nobreakspace}\ref{Families}). \texttt{Length( \mbox{\texttt{\mdseries\slshape imm}} )} is the number of elements of the argument, and \mbox{\texttt{\mdseries\slshape imm}} contains the elements in sorted order, w.r.t.{\nobreakspace}\texttt{{\textless}=}. \texttt{\mbox{\texttt{\mdseries\slshape new}}[\mbox{\texttt{\mdseries\slshape pos}}]} executes in constant time (see{\nobreakspace}\texttt{IsConstantTimeAccessList} (\ref{IsConstantTimeAccessList})), and the size of \mbox{\texttt{\mdseries\slshape imm}} in memory is proportional to its length. 

 The only difference to the operation \texttt{SortedList} (\ref{SortedList}) is that \texttt{AsSortedList} returns an \emph{immutable} list. 

 
\begin{Verbatim}[commandchars=!@|,fontsize=\small,frame=single,label=Example]
  !gapprompt@gap>| !gapinput@l:= AsSortedList( [ 1, 3, 3,, 2 ] );|
  [ 1, 2, 3, 3 ]
  !gapprompt@gap>| !gapinput@IsMutable( l );  IsSortedList( l );  IsConstantTimeAccessList( l );|
  false
  true
  true
  !gapprompt@gap>| !gapinput@IsSSortedList( l );|
  false
\end{Verbatim}
 }

 

\subsection{\textcolor{Chapter }{AsSSortedList}}
\logpage{[ 30, 3, 10 ]}\nobreak
\hyperdef{L}{X856D927378C33548}{}
{\noindent\textcolor{FuncColor}{$\triangleright$\enspace\texttt{AsSSortedList({\mdseries\slshape listorcoll})\index{AsSSortedList@\texttt{AsSSortedList}}
\label{AsSSortedList}
}\hfill{\scriptsize (attribute)}}\\
\noindent\textcolor{FuncColor}{$\triangleright$\enspace\texttt{AsSet({\mdseries\slshape listorcoll})\index{AsSet@\texttt{AsSet}}
\label{AsSet}
}\hfill{\scriptsize (attribute)}}\\


 \index{elements!of a list or collection} \texttt{AsSSortedList} (``as strictly sorted list'') returns a dense, immutable, and duplicate free list \mbox{\texttt{\mdseries\slshape imm}}. The argument must be a collection or list \mbox{\texttt{\mdseries\slshape listorcoll}} which may contain holes but whose elements lie in the same family
(see{\nobreakspace}\ref{Families}). \texttt{Length( \mbox{\texttt{\mdseries\slshape imm}} )} is the number of different elements of \mbox{\texttt{\mdseries\slshape listorcoll}}, and \mbox{\texttt{\mdseries\slshape imm}} contains the different elements in strictly sorted order,
w.r.t.{\nobreakspace}\texttt{\texttt{\symbol{92}}{\textless}} (\ref{bSlash<}). \texttt{\mbox{\texttt{\mdseries\slshape imm}}[\mbox{\texttt{\mdseries\slshape pos}}]} executes in constant time (see{\nobreakspace}\texttt{IsConstantTimeAccessList} (\ref{IsConstantTimeAccessList})), and the size of \mbox{\texttt{\mdseries\slshape imm}} in memory is proportional to its length. 

 Because the comparisons required for sorting can be very expensive for some
kinds of objects, you should use \texttt{AsList} (\ref{AsList}) instead if you do not require the result to be sorted. 

 The only difference to the operation \texttt{SSortedList} (\ref{SSortedList}) is that \texttt{AsSSortedList} returns an \emph{immutable} list. 

 \texttt{AsSet} is simply a synonym for \texttt{AsSSortedList}. 

 In general a function that returns a set of elements is free, in fact
encouraged, to return a domain instead of the proper set of its elements. This
allows one to keep a given structure, and moreover the representation by a
domain object is usually more space efficient. \texttt{AsSSortedList} must of course \emph{not} do this, its only purpose is to create the proper set of elements.    

 
\begin{Verbatim}[commandchars=!@|,fontsize=\small,frame=single,label=Example]
  !gapprompt@gap>| !gapinput@l:= AsSSortedList( l );|
  [ 1, 2, 3 ]
  !gapprompt@gap>| !gapinput@IsMutable( l );  IsSSortedList( l );  IsConstantTimeAccessList( l );|
  false
  true
  true
  !gapprompt@gap>| !gapinput@AsSSortedList( Group( (1,2,3), (1,2) ) );|
  [ (), (2,3), (1,2), (1,2,3), (1,3,2), (1,3) ]
\end{Verbatim}
 }

 

\subsection{\textcolor{Chapter }{Elements}}
\logpage{[ 30, 3, 11 ]}\nobreak
\hyperdef{L}{X79B130FC7906FB4C}{}
{\noindent\textcolor{FuncColor}{$\triangleright$\enspace\texttt{Elements({\mdseries\slshape C})\index{Elements@\texttt{Elements}}
\label{Elements}
}\hfill{\scriptsize (function)}}\\


 \texttt{Elements} does the same as \texttt{AsSSortedList} (\ref{AsSSortedList}), that is, the return value is a strictly sorted list of the elements in the
list or collection \mbox{\texttt{\mdseries\slshape C}}. 

 \texttt{Elements} is only supported for backwards compatibility. In many situations, the
sortedness of the ``element list'' for a collection is in fact not needed, and one can save a lot of time by
asking for a list that is \emph{not} necessarily sorted, using \texttt{AsList} (\ref{AsList}). If one is really interested in the strictly sorted list of elements in \mbox{\texttt{\mdseries\slshape C}} then one should use \texttt{AsSet} (\ref{AsSet}) or \texttt{AsSSortedList} (\ref{AsSSortedList}) instead. }

 }

  
\section{\textcolor{Chapter }{Attributes and Properties for Collections}}\label{Attributes and Properties for Collections}
\logpage{[ 30, 4, 0 ]}
\hyperdef{L}{X79AD18737E70B414}{}
{
  

\subsection{\textcolor{Chapter }{IsEmpty}}
\logpage{[ 30, 4, 1 ]}\nobreak
\hyperdef{L}{X7969C48780C5C1BC}{}
{\noindent\textcolor{FuncColor}{$\triangleright$\enspace\texttt{IsEmpty({\mdseries\slshape listorcoll})\index{IsEmpty@\texttt{IsEmpty}}
\label{IsEmpty}
}\hfill{\scriptsize (property)}}\\


 \texttt{IsEmpty} returns \texttt{true} if the collection or list \mbox{\texttt{\mdseries\slshape listorcoll}} is \emph{empty} (that is it contains no elements), and \texttt{false} otherwise. }

 

\subsection{\textcolor{Chapter }{IsFinite}}
\logpage{[ 30, 4, 2 ]}\nobreak
\hyperdef{L}{X808A4061809A6E67}{}
{\noindent\textcolor{FuncColor}{$\triangleright$\enspace\texttt{IsFinite({\mdseries\slshape C})\index{IsFinite@\texttt{IsFinite}}
\label{IsFinite}
}\hfill{\scriptsize (property)}}\\


 \index{finiteness test!for a list or collection} \texttt{IsFinite} returns \texttt{true} if the collection \mbox{\texttt{\mdseries\slshape C}} is finite, and \texttt{false} otherwise. 

 The default method for \texttt{IsFinite} checks the size (see{\nobreakspace}\texttt{Size} (\ref{Size})) of \mbox{\texttt{\mdseries\slshape C}}. 

 Methods for \texttt{IsFinite} may call \texttt{Size} (\ref{Size}), but methods for \texttt{Size} (\ref{Size}) must \emph{not} call \texttt{IsFinite}. }

 

\subsection{\textcolor{Chapter }{IsTrivial}}
\logpage{[ 30, 4, 3 ]}\nobreak
\hyperdef{L}{X7E3402D6799D3C24}{}
{\noindent\textcolor{FuncColor}{$\triangleright$\enspace\texttt{IsTrivial({\mdseries\slshape C})\index{IsTrivial@\texttt{IsTrivial}}
\label{IsTrivial}
}\hfill{\scriptsize (property)}}\\


 \texttt{IsTrivial} returns \texttt{true} if the collection \mbox{\texttt{\mdseries\slshape C}} consists of exactly one element.  }

 

\subsection{\textcolor{Chapter }{IsNonTrivial}}
\logpage{[ 30, 4, 4 ]}\nobreak
\hyperdef{L}{X7F192373850B85B9}{}
{\noindent\textcolor{FuncColor}{$\triangleright$\enspace\texttt{IsNonTrivial({\mdseries\slshape C})\index{IsNonTrivial@\texttt{IsNonTrivial}}
\label{IsNonTrivial}
}\hfill{\scriptsize (property)}}\\


 \texttt{IsNonTrivial} returns \texttt{true} if the collection \mbox{\texttt{\mdseries\slshape C}} is empty or consists of at least two elements (see{\nobreakspace}\texttt{IsTrivial} (\ref{IsTrivial})). 

    
\begin{Verbatim}[commandchars=!@|,fontsize=\small,frame=single,label=Example]
  !gapprompt@gap>| !gapinput@IsEmpty( [] ); IsEmpty( [ 1 .. 100 ] ); IsEmpty( Group( (1,2,3) ) );|
  true
  false
  false
  !gapprompt@gap>| !gapinput@IsFinite( [ 1 .. 100 ] );  IsFinite( Integers );|
  true
  false
  !gapprompt@gap>| !gapinput@IsTrivial( Integers );  IsTrivial( Group( () ) );|
  false
  true
  !gapprompt@gap>| !gapinput@IsNonTrivial( Integers );  IsNonTrivial( Group( () ) );|
  true
  false
\end{Verbatim}
 }

 

\subsection{\textcolor{Chapter }{IsWholeFamily}}
\logpage{[ 30, 4, 5 ]}\nobreak
\hyperdef{L}{X78EF6A137E8F66B0}{}
{\noindent\textcolor{FuncColor}{$\triangleright$\enspace\texttt{IsWholeFamily({\mdseries\slshape C})\index{IsWholeFamily@\texttt{IsWholeFamily}}
\label{IsWholeFamily}
}\hfill{\scriptsize (property)}}\\


 \texttt{IsWholeFamily} returns \texttt{true} if the collection \mbox{\texttt{\mdseries\slshape C}} contains the whole family (see{\nobreakspace}\ref{Families}) of its elements. 

 
\begin{Verbatim}[commandchars=!@|,fontsize=\small,frame=single,label=Example]
  !gapprompt@gap>| !gapinput@IsWholeFamily( Integers )|
  !gapprompt@>| !gapinput@   ;  # all rationals and cyclotomics lie in the family|
  false
  !gapprompt@gap>| !gapinput@IsWholeFamily( Integers mod 3 )|
  !gapprompt@>| !gapinput@   ;  # all finite field elements in char. 3 lie in this family|
  false
  !gapprompt@gap>| !gapinput@IsWholeFamily( Integers mod 4 );|
  true
  !gapprompt@gap>| !gapinput@IsWholeFamily( FreeGroup( 2 ) );|
  true
\end{Verbatim}
 }

 

\subsection{\textcolor{Chapter }{Size}}
\logpage{[ 30, 4, 6 ]}\nobreak
\hyperdef{L}{X858ADA3B7A684421}{}
{\noindent\textcolor{FuncColor}{$\triangleright$\enspace\texttt{Size({\mdseries\slshape listorcoll})\index{Size@\texttt{Size}}
\label{Size}
}\hfill{\scriptsize (attribute)}}\\


 \index{size!of a list or collection} \index{order!of a list, collection or domain} \texttt{Size} returns the size of the list or collection \mbox{\texttt{\mdseries\slshape listorcoll}}, which is either an integer or \texttt{infinity} (\ref{infinity}). If the argument is a list then the result is its length (see{\nobreakspace}\texttt{Length} (\ref{Length})). 

 The default method for \texttt{Size} checks the length of an enumerator of \mbox{\texttt{\mdseries\slshape listorcoll}}. 

 Methods for \texttt{IsFinite} (\ref{IsFinite}) may call \texttt{Size}, but methods for \texttt{Size} must not call \texttt{IsFinite} (\ref{IsFinite}). 

 
\begin{Verbatim}[commandchars=!@|,fontsize=\small,frame=single,label=Example]
  !gapprompt@gap>| !gapinput@Size( [1,2,3] );  Size( Group( () ) );  Size( Integers );|
  3
  1
  infinity
\end{Verbatim}
 }

 

\subsection{\textcolor{Chapter }{Representative}}
\logpage{[ 30, 4, 7 ]}\nobreak
\hyperdef{L}{X865507568182424E}{}
{\noindent\textcolor{FuncColor}{$\triangleright$\enspace\texttt{Representative({\mdseries\slshape C})\index{Representative@\texttt{Representative}}
\label{Representative}
}\hfill{\scriptsize (attribute)}}\\


 \texttt{Representative} returns a \emph{representative} of the collection \mbox{\texttt{\mdseries\slshape C}}. 

 Note that \texttt{Representative} is free in choosing a representative if there are several elements in \mbox{\texttt{\mdseries\slshape C}}. It is not even guaranteed that \texttt{Representative} returns the same representative if it is called several times for one
collection. The main difference between \texttt{Representative} and \texttt{Random} (\ref{Random:for a list or collection}) is that \texttt{Representative} is free to choose a value that is cheap to compute, while \texttt{Random} (\ref{Random:for a list or collection}) must make an effort to randomly distribute its answers. 

 If \mbox{\texttt{\mdseries\slshape C}} is a domain then there are methods for \texttt{Representative} that try to fetch an element from any known generator list of \mbox{\texttt{\mdseries\slshape C}}, see{\nobreakspace}\ref{Domains and their Elements}. Note that \texttt{Representative} does not try to \emph{compute} generators of \mbox{\texttt{\mdseries\slshape C}}, thus \texttt{Representative} may give up and signal an error if \mbox{\texttt{\mdseries\slshape C}} has no generators stored at all. }

 

\subsection{\textcolor{Chapter }{RepresentativeSmallest}}
\logpage{[ 30, 4, 8 ]}\nobreak
\hyperdef{L}{X8026085680270D37}{}
{\noindent\textcolor{FuncColor}{$\triangleright$\enspace\texttt{RepresentativeSmallest({\mdseries\slshape C})\index{RepresentativeSmallest@\texttt{RepresentativeSmallest}}
\label{RepresentativeSmallest}
}\hfill{\scriptsize (attribute)}}\\


 \index{representative!of a list or collection} returns the smallest element in the collection \mbox{\texttt{\mdseries\slshape C}}, w.r.t.{\nobreakspace}the ordering \texttt{\texttt{\symbol{92}}{\textless}} (\ref{bSlash<}). While the operation defaults to comparing all elements, better methods are
installed for some collections. 

 
\begin{Verbatim}[commandchars=!@|,fontsize=\small,frame=single,label=Example]
  !gapprompt@gap>| !gapinput@Representative( Rationals );|
  0
  !gapprompt@gap>| !gapinput@Representative( [ -1, -2 .. -100 ] );|
  -1
  !gapprompt@gap>| !gapinput@RepresentativeSmallest( [ -1, -2 .. -100 ] );|
  -100
\end{Verbatim}
 }

 }

  
\section{\textcolor{Chapter }{Operations for Collections}}\label{Operations for Collections}
\logpage{[ 30, 5, 0 ]}
\hyperdef{L}{X7F8FEA3278239ADE}{}
{
  

\subsection{\textcolor{Chapter }{IsSubset}}
\logpage{[ 30, 5, 1 ]}\nobreak
\hyperdef{L}{X79CA175481F8105F}{}
{\noindent\textcolor{FuncColor}{$\triangleright$\enspace\texttt{IsSubset({\mdseries\slshape C1, C2})\index{IsSubset@\texttt{IsSubset}}
\label{IsSubset}
}\hfill{\scriptsize (operation)}}\\


 \index{subset test!for collections} \texttt{IsSubset} returns \texttt{true} if \mbox{\texttt{\mdseries\slshape C2}}, which must be a collection, is a \emph{subset} of \mbox{\texttt{\mdseries\slshape C1}}, which also must be a collection, and \texttt{false} otherwise. 

 \mbox{\texttt{\mdseries\slshape C2}} is considered a subset of \mbox{\texttt{\mdseries\slshape C1}} if and only if each element of \mbox{\texttt{\mdseries\slshape C2}} is also an element of \mbox{\texttt{\mdseries\slshape C1}}. That is \texttt{IsSubset} behaves as if implemented as \texttt{IsSubsetSet( AsSSortedList( \mbox{\texttt{\mdseries\slshape C1}} ), AsSSortedList( \mbox{\texttt{\mdseries\slshape C2}} ) )}, except that it will also sometimes, but not always, work for infinite
collections, and that it will usually work much faster than the above
definition. Either argument may also be a proper set (see{\nobreakspace}\ref{Sorted Lists and Sets}). 

 
\begin{Verbatim}[commandchars=!@|,fontsize=\small,frame=single,label=Example]
  !gapprompt@gap>| !gapinput@IsSubset( Rationals, Integers );|
  true
  !gapprompt@gap>| !gapinput@IsSubset( Integers, [ 1, 2, 3 ] );|
  true
  !gapprompt@gap>| !gapinput@IsSubset( Group( (1,2,3,4) ), [ (1,2,3) ] );|
  false
\end{Verbatim}
 }

 
\subsection{\textcolor{Chapter }{Intersection}}\logpage{[ 30, 5, 2 ]}
\hyperdef{L}{X851069107CACF98E}{}
{
\noindent\textcolor{FuncColor}{$\triangleright$\enspace\texttt{Intersection({\mdseries\slshape C1, C2, ...})\index{Intersection@\texttt{Intersection}!for various collections}
\label{Intersection:for various collections}
}\hfill{\scriptsize (function)}}\\
\noindent\textcolor{FuncColor}{$\triangleright$\enspace\texttt{Intersection({\mdseries\slshape list})\index{Intersection@\texttt{Intersection}!for a list}
\label{Intersection:for a list}
}\hfill{\scriptsize (function)}}\\
\noindent\textcolor{FuncColor}{$\triangleright$\enspace\texttt{Intersection2({\mdseries\slshape C1, C2})\index{Intersection2@\texttt{Intersection2}}
\label{Intersection2}
}\hfill{\scriptsize (operation)}}\\


 \index{intersection!of collections} In the first form \texttt{Intersection} returns the intersection of the collections \mbox{\texttt{\mdseries\slshape C1}}, \mbox{\texttt{\mdseries\slshape C2}}, etc. In the second form \mbox{\texttt{\mdseries\slshape list}} must be a \emph{nonempty} list of collections and \texttt{Intersection} returns the intersection of those collections. Each argument or element of \mbox{\texttt{\mdseries\slshape list}} respectively may also be a homogeneous list that is not a proper set, in which
case \texttt{Intersection} silently applies \texttt{Set} (\ref{Set}) to it first. 

 The result of \texttt{Intersection} is the set of elements that lie in every of the collections \mbox{\texttt{\mdseries\slshape C1}}, \mbox{\texttt{\mdseries\slshape C2}}, etc. If the result is a list then it is mutable and new, i.e., not identical
to any of \mbox{\texttt{\mdseries\slshape C1}}, \mbox{\texttt{\mdseries\slshape C2}}, etc. 

 Methods can be installed for the operation \texttt{Intersection2} that takes only two arguments. \texttt{Intersection} calls \texttt{Intersection2}. 

 Methods for \texttt{Intersection2} should try to maintain as much structure as possible, for example the
intersection of two permutation groups is again a permutation group. 

 
\begin{Verbatim}[commandchars=!@|,fontsize=\small,frame=single,label=Example]
  !gapprompt@gap>| !gapinput@# this is one of the rare cases where the intersection of two|
  !gapprompt@gap>| !gapinput@# infinite domains works ('CF' is a shorthand for 'CyclotomicField'):|
  !gapprompt@gap>| !gapinput@Intersection( CyclotomicField(9), CyclotomicField(12) );|
  CF(3)
  !gapprompt@gap>| !gapinput@D12 := Group( (2,6)(3,5), (1,2)(3,6)(4,5) );;|
  !gapprompt@gap>| !gapinput@Intersection( D12, Group( (1,2), (1,2,3,4,5) ) );|
  Group([ (1,5)(2,4) ])
  !gapprompt@gap>| !gapinput@Intersection( D12, [ (1,3)(4,6), (1,2)(3,4) ] )|
  !gapprompt@>| !gapinput@   ;  # note that the second argument is not a proper set|
  [ (1,3)(4,6) ]
  !gapprompt@gap>| !gapinput@# although the result is mathematically a group it is returned as a|
  !gapprompt@gap>| !gapinput@# proper set because the second argument is not regarded as a group:|
  !gapprompt@gap>| !gapinput@Intersection( D12, [ (), (1,2)(3,4), (1,3)(4,6), (1,4)(5,6) ] );|
  [ (), (1,3)(4,6) ]
  !gapprompt@gap>| !gapinput@Intersection( Group( () ), [1,2,3] );|
  [  ]
  !gapprompt@gap>| !gapinput@Intersection( [2,4,6,8,10], [3,6,9,12,15], [5,10,15,20,25] )|
  !gapprompt@>| !gapinput@   ;  # two or more lists or collections as arguments are legal|
  [  ]
  !gapprompt@gap>| !gapinput@Intersection( [ [1,2,4], [2,3,4], [1,3,4] ] )|
  !gapprompt@>| !gapinput@   ;  # or one list of lists or collections|
  [ 4 ]
\end{Verbatim}
 }

 
\subsection{\textcolor{Chapter }{Union}}\logpage{[ 30, 5, 3 ]}
\hyperdef{L}{X799F0E2F7A502DBA}{}
{
\noindent\textcolor{FuncColor}{$\triangleright$\enspace\texttt{Union({\mdseries\slshape C1, C2, ...})\index{Union@\texttt{Union}!for various collections}
\label{Union:for various collections}
}\hfill{\scriptsize (function)}}\\
\noindent\textcolor{FuncColor}{$\triangleright$\enspace\texttt{Union({\mdseries\slshape list})\index{Union@\texttt{Union}!for a list}
\label{Union:for a list}
}\hfill{\scriptsize (function)}}\\
\noindent\textcolor{FuncColor}{$\triangleright$\enspace\texttt{Union2({\mdseries\slshape C1, C2})\index{Union2@\texttt{Union2}}
\label{Union2}
}\hfill{\scriptsize (operation)}}\\


 \index{union!of collections} In the first form \texttt{Union} returns the union of the collections \mbox{\texttt{\mdseries\slshape C1}}, \mbox{\texttt{\mdseries\slshape C2}}, etc. In the second form \mbox{\texttt{\mdseries\slshape list}} must be a list of collections and \texttt{Union} returns the union of those collections. Each argument or element of \mbox{\texttt{\mdseries\slshape list}} respectively may also be a homogeneous list that is not a proper set, in which
case \texttt{Union} silently applies \texttt{Set} (\ref{Set}) to it first. 

 The result of \texttt{Union} is the set of elements that lie in any of the collections \mbox{\texttt{\mdseries\slshape C1}}, \mbox{\texttt{\mdseries\slshape C2}}, etc. If the result is a list then it is mutable and new, i.e., not identical
to any of \mbox{\texttt{\mdseries\slshape C1}}, \mbox{\texttt{\mdseries\slshape C2}}, etc. 

 Methods can be installed for the operation \texttt{Union2} that takes only two arguments. \texttt{Union} calls \texttt{Union2}. 

 
\begin{Verbatim}[commandchars=!@|,fontsize=\small,frame=single,label=Example]
  !gapprompt@gap>| !gapinput@Union( [ (1,2,3), (1,2,3,4) ], Group( (1,2,3), (1,2) ) );|
  [ (), (2,3), (1,2), (1,2,3), (1,2,3,4), (1,3,2), (1,3) ]
  !gapprompt@gap>| !gapinput@Union( [2,4,6,8,10], [3,6,9,12,15], [5,10,15,20,25] )|
  !gapprompt@>| !gapinput@   ;  # two or more lists or collections as arguments are legal|
  [ 2, 3, 4, 5, 6, 8, 9, 10, 12, 15, 20, 25 ]
  !gapprompt@gap>| !gapinput@Union( [ [1,2,4], [2,3,4], [1,3,4] ] )|
  !gapprompt@>| !gapinput@   ;  # or one list of lists or collections|
  [ 1 .. 4 ]
  !gapprompt@gap>| !gapinput@Union( [ ] );|
  [  ]
\end{Verbatim}


 When computing the Union of lists or sets of small integers and ranges, every
attempt is made to return the result as a range and to avoid expanding ranges
provided as input. }

 

\subsection{\textcolor{Chapter }{Difference}}
\logpage{[ 30, 5, 4 ]}\nobreak
\hyperdef{L}{X825AC0F07E010B07}{}
{\noindent\textcolor{FuncColor}{$\triangleright$\enspace\texttt{Difference({\mdseries\slshape C1, C2})\index{Difference@\texttt{Difference}}
\label{Difference}
}\hfill{\scriptsize (operation)}}\\


 \index{set difference!of collections} \texttt{Difference} returns the set difference of the collections \mbox{\texttt{\mdseries\slshape C1}} and \mbox{\texttt{\mdseries\slshape C2}}. Either argument may also be a homogeneous list that is not a proper set, in
which case \texttt{Difference} silently applies \texttt{Set} (\ref{Set}) to it first. 

 The result of \texttt{Difference} is the set of elements that lie in \mbox{\texttt{\mdseries\slshape C1}} but not in \mbox{\texttt{\mdseries\slshape C2}}. Note that \mbox{\texttt{\mdseries\slshape C2}} need not be a subset of \mbox{\texttt{\mdseries\slshape C1}}. The elements of \mbox{\texttt{\mdseries\slshape C2}}, however, that are not elements of \mbox{\texttt{\mdseries\slshape C1}} play no role for the result. If the result is a list then it is mutable and
new, i.e., not identical to \mbox{\texttt{\mdseries\slshape C1}} or \mbox{\texttt{\mdseries\slshape C2}}. 

 
\begin{Verbatim}[commandchars=!@|,fontsize=\small,frame=single,label=Example]
  !gapprompt@gap>| !gapinput@Difference( [ (1,2,3), (1,2,3,4) ], Group( (1,2,3), (1,2) ) );|
  [ (1,2,3,4) ]
\end{Verbatim}
 }

 }

  
\section{\textcolor{Chapter }{Membership Test for Collections}}\label{Membership Test for Collections}
\logpage{[ 30, 6, 0 ]}
\hyperdef{L}{X82D39CF980FDBFFA}{}
{
  \index{in@\texttt{\texttt{\symbol{92}}in}!operation for testing membership} \index{in@\texttt{in}!operation for} 

\subsection{\textcolor{Chapter }{\texttt{\symbol{92}}in (for a collection)}}
\logpage{[ 30, 6, 1 ]}\nobreak
\hyperdef{L}{X84B7FA8C7C94400F}{}
{\noindent\textcolor{FuncColor}{$\triangleright$\enspace\texttt{\texttt{\symbol{92}}in({\mdseries\slshape obj, C})\index{\texttt{\symbol{92}}in@\texttt{\texttt{\symbol{92}}in}!for a collection}
\label{bSlashin:for a collection}
}\hfill{\scriptsize (operation)}}\\


 returns \texttt{true} if the object \mbox{\texttt{\mdseries\slshape obj}} lies in the collection \mbox{\texttt{\mdseries\slshape C}}, and \texttt{false} otherwise. 

 The infix version of the command 

 \mbox{\texttt{\mdseries\slshape obj}} \texttt{in} \mbox{\texttt{\mdseries\slshape C}} 

 calls the operation \texttt{\texttt{\symbol{92}}in}, for which methods can be installed. 

 
\begin{Verbatim}[commandchars=!@|,fontsize=\small,frame=single,label=Example]
  !gapprompt@gap>| !gapinput@13 in Integers;  [ 1, 2 ] in Integers;|
  true
  false
  !gapprompt@gap>| !gapinput@g:= Group( (1,2) );;  (1,2) in g;  (1,2,3) in g;|
  true
  false
\end{Verbatim}
 }

 }

  
\section{\textcolor{Chapter }{Random Elements}}\label{Random Elements}
\logpage{[ 30, 7, 0 ]}
\hyperdef{L}{X8151A51884B7EE2C}{}
{
  The method used by \textsf{GAP} to obtain random elements may depend on the type object. 

 Most methods which produce random elements in \textsf{GAP} use a global random number generator (see \texttt{GlobalMersenneTwister} (\ref{GlobalMersenneTwister})). This random number generator is (deliberately) initialized to the same
values when \textsf{GAP} is started, so different runs of \textsf{GAP} with the same input will always produce the same result, even if random
calculations are involved. 

 See \texttt{Reset} (\ref{Reset}) for a description of how to reset the random number generator to a previous
state. 

 

\subsection{\textcolor{Chapter }{Random (for a list or collection)}}
\logpage{[ 30, 7, 1 ]}\nobreak
\hyperdef{L}{X7FF906E57D6936F8}{}
{\noindent\textcolor{FuncColor}{$\triangleright$\enspace\texttt{Random({\mdseries\slshape listorcoll})\index{Random@\texttt{Random}!for a list or collection}
\label{Random:for a list or collection}
}\hfill{\scriptsize (operation)}}\\
\noindent\textcolor{FuncColor}{$\triangleright$\enspace\texttt{Random({\mdseries\slshape from, to})\index{Random@\texttt{Random}!for lower and upper bound}
\label{Random:for lower and upper bound}
}\hfill{\scriptsize (operation)}}\\


  \index{Random@\texttt{Random}} \texttt{Random} returns a (pseudo\texttt{\symbol{45}})random element of the dense, nonempty
list or nonempty collection \mbox{\texttt{\mdseries\slshape listorcoll}}. The behaviour for non\texttt{\symbol{45}}dense or empty lists, and for empty
collections (see \texttt{IsDenseList} (\ref{IsDenseList}), \texttt{IsEmpty} (\ref{IsEmpty})) is undefined. 

 As lists or ranges are restricted in length ($2^{28}-1$ or $2^{60}-1$ depending on your system), the second form returns a random integer in the
range \mbox{\texttt{\mdseries\slshape from}} to \mbox{\texttt{\mdseries\slshape to}} (inclusive) for arbitrary integers \mbox{\texttt{\mdseries\slshape from}} and \mbox{\texttt{\mdseries\slshape to}}. The behaviour in the case that \mbox{\texttt{\mdseries\slshape from}} is larger than \mbox{\texttt{\mdseries\slshape to}} is undefined. 

 See Section \ref{Random Sources} for more about computing random elements, in particular for \texttt{Random} (\ref{Random:for random source and list}) methods that take a random source as the first argument. 

 The distribution of elements returned by \texttt{Random} depends on the argument. For a dense, nonempty list the distribution is
uniform (all elements are equally likely). The same holds usually for finite
collections that are not lists. For infinite collections some reasonable
distribution is used. 

 See the chapters of the various collections to find out which distribution is
being used. 

 For some collections ensuring a reasonable distribution can be difficult and
require substantial runtime (for example for large finite groups). If speed is
more important than a guaranteed distribution, the operation \texttt{PseudoRandom} (\ref{PseudoRandom}) should be used instead. 

 Note that \texttt{Random} is of course \emph{not} an attribute. 

 
\begin{Verbatim}[commandchars=!@|,fontsize=\small,frame=single,label=Example]
  !gapprompt@gap>| !gapinput@Random([1..6]);|
  6
  !gapprompt@gap>| !gapinput@Random(1, 2^100);|
  866227015645295902682304086250
  !gapprompt@gap>| !gapinput@g:= Group( (1,2,3) );;  Random( g );  Random( g );|
  (1,3,2)
  ()
  !gapprompt@gap>| !gapinput@Random(Rationals);|
  -4
\end{Verbatim}
 }

 

\subsection{\textcolor{Chapter }{PseudoRandom}}
\logpage{[ 30, 7, 2 ]}\nobreak
\hyperdef{L}{X811B5BD47DC5356B}{}
{\noindent\textcolor{FuncColor}{$\triangleright$\enspace\texttt{PseudoRandom({\mdseries\slshape listorcoll})\index{PseudoRandom@\texttt{PseudoRandom}}
\label{PseudoRandom}
}\hfill{\scriptsize (operation)}}\\


 \texttt{PseudoRandom} returns a pseudo random element of the list or collection \mbox{\texttt{\mdseries\slshape listorcoll}}, which can be roughly described as follows. For a list, \texttt{PseudoRandom} returns the same as \texttt{Random} (\ref{Random:for a list or collection}). For collections that are not lists, the elements returned by \texttt{PseudoRandom} are \emph{not} necessarily equally distributed, even for finite collections; the idea is that \texttt{Random} (\ref{Random:for a list or collection}) returns elements according to a reasonable distribution, \texttt{PseudoRandom} returns elements that are cheap to compute but need not satisfy this strong
condition, and \texttt{Representative} (\ref{Representative}) returns arbitrary elements, probably the same element for each call. }

 

\subsection{\textcolor{Chapter }{RandomList}}
\logpage{[ 30, 7, 3 ]}\nobreak
\hyperdef{L}{X7EBA01EB83BC65A9}{}
{\noindent\textcolor{FuncColor}{$\triangleright$\enspace\texttt{RandomList({\mdseries\slshape [rs, ]list})\index{RandomList@\texttt{RandomList}}
\label{RandomList}
}\hfill{\scriptsize (function)}}\\


 \index{random seed} For a dense list \mbox{\texttt{\mdseries\slshape list}}, \texttt{RandomList} returns a (pseudo\texttt{\symbol{45}})random element with equal distribution. 

 The random source \mbox{\texttt{\mdseries\slshape rs}} (see \ref{Random Sources}) is used to choose a random number. If \mbox{\texttt{\mdseries\slshape rs}} is absent, this function uses the \texttt{GlobalMersenneTwister} (\ref{GlobalMersenneTwister}) to produce the random elements (a source of high quality random numbers). 

 
\begin{Verbatim}[commandchars=!@|,fontsize=\small,frame=single,label=Example]
  !gapprompt@gap>| !gapinput@RandomList( [ 1 .. 6 ] );|
  3
  !gapprompt@gap>| !gapinput@elms:= AsList( Group( (1,2,3) ) );;|
  !gapprompt@gap>| !gapinput@RandomList( elms );  RandomList( elms );|
  (1,3,2)
  (1,2,3)
  !gapprompt@gap>| !gapinput@rs:= RandomSource( IsMersenneTwister, 1 );|
  <RandomSource in IsMersenneTwister>
  !gapprompt@gap>| !gapinput@RandomList( rs, elms );|
  (1,3,2)
\end{Verbatim}
 }

 }

  
\section{\textcolor{Chapter }{Iterators}}\label{Iterators}
\logpage{[ 30, 8, 0 ]}
\hyperdef{L}{X85A3F00985453F95}{}
{
  

\subsection{\textcolor{Chapter }{Iterator}}
\logpage{[ 30, 8, 1 ]}\nobreak
\hyperdef{L}{X83ADF8287ED0668E}{}
{\noindent\textcolor{FuncColor}{$\triangleright$\enspace\texttt{Iterator({\mdseries\slshape listorcoll})\index{Iterator@\texttt{Iterator}}
\label{Iterator}
}\hfill{\scriptsize (operation)}}\\
\noindent\textcolor{FuncColor}{$\triangleright$\enspace\texttt{IsStandardIterator({\mdseries\slshape listorcoll})\index{IsStandardIterator@\texttt{IsStandardIterator}}
\label{IsStandardIterator}
}\hfill{\scriptsize (filter)}}\\


 Iterators provide a possibility to loop over the elements of a (countable)
collection or list \mbox{\texttt{\mdseries\slshape listorcoll}}, without repetition. For many collections $C$, an iterator of $C$ need not store all elements of $C$, for example it is possible to construct an iterator of some infinite
domains, such as the field of rational numbers. 

 \texttt{Iterator} returns a mutable \emph{iterator} $iter$ for its argument. If this argument is a list (which may contain holes), then $iter$ iterates over the elements (but not the holes) of this list in the same order
(see{\nobreakspace}\texttt{IteratorList} (\ref{IteratorList}) for details). If this argument is a collection but not a list then $iter$ iterates over the elements of this collection in an unspecified order, which
may change for repeated calls of \texttt{Iterator}. Because iterators returned by \texttt{Iterator} are mutable (see{\nobreakspace}\ref{Mutability and Copyability}), each call of \texttt{Iterator} for the same argument returns a \emph{new} iterator. Therefore \texttt{Iterator} is not an attribute (see{\nobreakspace}\ref{Attributes}). 

 The only operations for iterators are \texttt{IsDoneIterator} (\ref{IsDoneIterator}), \texttt{NextIterator} (\ref{NextIterator}), and \texttt{ShallowCopy} (\ref{ShallowCopy}). In particular, it is only possible to access the next element of the
iterator with \texttt{NextIterator} (\ref{NextIterator}) if there is one, and this can be checked with \texttt{IsDoneIterator} (\ref{IsDoneIterator}) For an iterator $iter$, \texttt{ShallowCopy} (\ref{ShallowCopy}) returns a mutable iterator $new$ that iterates over the remaining elements independent of $iter$; the results of \texttt{IsDoneIterator} (\ref{IsDoneIterator}) for $iter$ and $new$ are equal, and if $iter$ is mutable then also the results of \texttt{NextIterator} (\ref{NextIterator}) for $iter$ and $new$ are equal; note that \texttt{=} is not defined for iterators, so the equality of two iterators cannot be
checked with \texttt{=}. 

 When \texttt{Iterator} is called for a \emph{mutable} collection $C$ then it is not defined whether $iter$ respects changes to $C$ occurring after the construction of $iter$, except if the documentation explicitly promises a certain behaviour. The
latter is the case if the argument is a mutable list (see{\nobreakspace}\texttt{IteratorList} (\ref{IteratorList}) for subtleties in this case). 

 It is possible to have \texttt{for}\texttt{\symbol{45}}loops run over mutable iterators instead of lists. 

 In some situations, one can construct iterators with a special succession of
elements, see{\nobreakspace}\texttt{IteratorByBasis} (\ref{IteratorByBasis}) for the possibility to loop over the elements of a vector space
w.r.t.{\nobreakspace}a given basis.  

 For lists, \texttt{Iterator} is implemented by \texttt{IteratorList} (\ref{IteratorList}). For collections $C$ that are not lists, the default method is \texttt{IteratorList( Enumerator( }$C$\texttt{ ) )}. Better methods depending on $C$ should be provided if possible. 

 For random access to the elements of a (possibly infinite) collection, \emph{enumerators} are used. See{\nobreakspace}\ref{Enumerators} for the facility to compute a list from $C$, which provides a (partial) mapping from $C$ to the positive integers. 

 The filter \texttt{IsStandardIterator} means that the iterator is implemented as a component object and has
components \texttt{IsDoneIterator} and \texttt{NextIterator} which are bound to the methods of the operations of the same name for this
iterator.          
\begin{Verbatim}[commandchars=!@|,fontsize=\small,frame=single,label=Example]
  !gapprompt@gap>| !gapinput@iter:= Iterator( GF(5) );|
  <iterator>
  !gapprompt@gap>| !gapinput@l:= [];;|
  !gapprompt@gap>| !gapinput@for i in iter do Add( l, i ); od; l;|
  [ 0*Z(5), Z(5)^0, Z(5), Z(5)^2, Z(5)^3 ]
  !gapprompt@gap>| !gapinput@iter:= Iterator( [ 1, 2, 3, 4 ] );;  l:= [];;|
  !gapprompt@gap>| !gapinput@for i in iter do|
  !gapprompt@>| !gapinput@     new:= ShallowCopy( iter );|
  !gapprompt@>| !gapinput@     for j in new do Add( l, j ); od;|
  !gapprompt@>| !gapinput@   od; l;|
  [ 2, 3, 4, 3, 4, 4 ]
\end{Verbatim}
 }

 

\subsection{\textcolor{Chapter }{IteratorSorted}}
\logpage{[ 30, 8, 2 ]}\nobreak
\hyperdef{L}{X8688C20B828FC129}{}
{\noindent\textcolor{FuncColor}{$\triangleright$\enspace\texttt{IteratorSorted({\mdseries\slshape listorcoll})\index{IteratorSorted@\texttt{IteratorSorted}}
\label{IteratorSorted}
}\hfill{\scriptsize (operation)}}\\


 \texttt{IteratorSorted} returns a mutable iterator. The argument must be a collection or a list that
is not necessarily dense but whose elements lie in the same family
(see{\nobreakspace}\ref{Families}). It loops over the different elements in sorted order. 

 For a collection $C$ that is not a list, the generic method is \texttt{IteratorList( EnumeratorSorted( }\mbox{\texttt{\mdseries\slshape C}}\texttt{ ) )}. }

 

\subsection{\textcolor{Chapter }{IsIterator}}
\logpage{[ 30, 8, 3 ]}\nobreak
\hyperdef{L}{X87168A827E5B28E4}{}
{\noindent\textcolor{FuncColor}{$\triangleright$\enspace\texttt{IsIterator({\mdseries\slshape obj})\index{IsIterator@\texttt{IsIterator}}
\label{IsIterator}
}\hfill{\scriptsize (Category)}}\\


 Every iterator lies in the category \texttt{IsIterator}. }

 

\subsection{\textcolor{Chapter }{IsDoneIterator}}
\logpage{[ 30, 8, 4 ]}\nobreak
\hyperdef{L}{X8055FC557B5D899E}{}
{\noindent\textcolor{FuncColor}{$\triangleright$\enspace\texttt{IsDoneIterator({\mdseries\slshape iter})\index{IsDoneIterator@\texttt{IsDoneIterator}}
\label{IsDoneIterator}
}\hfill{\scriptsize (operation)}}\\


 If \mbox{\texttt{\mdseries\slshape iter}} is an iterator for the list or collection $C$ then \texttt{IsDoneIterator( \mbox{\texttt{\mdseries\slshape iter}} )} is \texttt{true} if all elements of $C$ have been returned already by \texttt{NextIterator( \mbox{\texttt{\mdseries\slshape iter}} )}, and \texttt{false} otherwise. }

 

\subsection{\textcolor{Chapter }{NextIterator}}
\logpage{[ 30, 8, 5 ]}\nobreak
\hyperdef{L}{X879F62F77D1D1179}{}
{\noindent\textcolor{FuncColor}{$\triangleright$\enspace\texttt{NextIterator({\mdseries\slshape iter})\index{NextIterator@\texttt{NextIterator}}
\label{NextIterator}
}\hfill{\scriptsize (operation)}}\\


 Let \mbox{\texttt{\mdseries\slshape iter}} be a mutable iterator for the list or collection $C$. If \texttt{IsDoneIterator( \mbox{\texttt{\mdseries\slshape iter}} )} is \texttt{false} then \texttt{NextIterator} is applicable to \mbox{\texttt{\mdseries\slshape iter}}, and the result is the next element of $C$, according to the succession defined by \mbox{\texttt{\mdseries\slshape iter}}. 

 If \texttt{IsDoneIterator( \mbox{\texttt{\mdseries\slshape iter}} )} is \texttt{true} then it is not defined what happens when \texttt{NextIterator} is called for \mbox{\texttt{\mdseries\slshape iter}}; that is, it may happen that an error is signalled or that something
meaningless is returned, or even that \textsf{GAP} crashes. 

 
\begin{Verbatim}[commandchars=!@|,fontsize=\small,frame=single,label=Example]
  !gapprompt@gap>| !gapinput@iter:= Iterator( [ 1, 2, 3, 4 ] );|
  <iterator>
  !gapprompt@gap>| !gapinput@sum:= 0;;|
  !gapprompt@gap>| !gapinput@while not IsDoneIterator( iter ) do|
  !gapprompt@>| !gapinput@     sum:= sum + NextIterator( iter );|
  !gapprompt@>| !gapinput@   od;|
  !gapprompt@gap>| !gapinput@IsDoneIterator( iter ); sum;|
  true
  10
  !gapprompt@gap>| !gapinput@ir:= Iterator( Rationals );;|
  !gapprompt@gap>| !gapinput@l:= [];; for i in [1..20] do Add( l, NextIterator( ir ) ); od; l;|
  [ 0, 1, -1, 1/2, 2, -1/2, -2, 1/3, 2/3, 3/2, 3, -1/3, -2/3, -3/2, -3,
    1/4, 3/4, 4/3, 4, -1/4 ]
  !gapprompt@gap>| !gapinput@for i in ir do|
  !gapprompt@>| !gapinput@     if DenominatorRat( i ) > 10 then break; fi;|
  !gapprompt@>| !gapinput@   od;|
  !gapprompt@gap>| !gapinput@i;|
  1/11
\end{Verbatim}
 }

 

\subsection{\textcolor{Chapter }{IteratorList}}
\logpage{[ 30, 8, 6 ]}\nobreak
\hyperdef{L}{X858A28667D137C4B}{}
{\noindent\textcolor{FuncColor}{$\triangleright$\enspace\texttt{IteratorList({\mdseries\slshape list})\index{IteratorList@\texttt{IteratorList}}
\label{IteratorList}
}\hfill{\scriptsize (function)}}\\


 \texttt{IteratorList} returns a new iterator that allows iteration over the elements of the list \mbox{\texttt{\mdseries\slshape list}} (which may have holes) in the same order. 

 If \mbox{\texttt{\mdseries\slshape list}} is mutable then it is in principle possible to change \mbox{\texttt{\mdseries\slshape list}} after the call of \texttt{IteratorList}. In this case all changes concerning positions that have not yet been reached
in the iteration will also affect the iterator. For example, if \mbox{\texttt{\mdseries\slshape list}} is enlarged then the iterator will iterate also over the new elements at the
end of the changed list. 

 \emph{Note} that changes of \mbox{\texttt{\mdseries\slshape list}} will also affect all shallow copies of \mbox{\texttt{\mdseries\slshape list}}. }

 

\subsection{\textcolor{Chapter }{TrivialIterator}}
\logpage{[ 30, 8, 7 ]}\nobreak
\hyperdef{L}{X7DB80BE68271247E}{}
{\noindent\textcolor{FuncColor}{$\triangleright$\enspace\texttt{TrivialIterator({\mdseries\slshape elm})\index{TrivialIterator@\texttt{TrivialIterator}}
\label{TrivialIterator}
}\hfill{\scriptsize (function)}}\\


 is a mutable iterator for the collection \texttt{[ \mbox{\texttt{\mdseries\slshape elm}} ]} that consists of exactly one element \mbox{\texttt{\mdseries\slshape elm}} (see{\nobreakspace}\texttt{IsTrivial} (\ref{IsTrivial})). }

 

\subsection{\textcolor{Chapter }{IteratorByFunctions}}
\logpage{[ 30, 8, 8 ]}\nobreak
\hyperdef{L}{X82677D8F817D6701}{}
{\noindent\textcolor{FuncColor}{$\triangleright$\enspace\texttt{IteratorByFunctions({\mdseries\slshape record})\index{IteratorByFunctions@\texttt{IteratorByFunctions}}
\label{IteratorByFunctions}
}\hfill{\scriptsize (function)}}\\


 \texttt{IteratorByFunctions} returns a (mutable) iterator \mbox{\texttt{\mdseries\slshape iter}} for which \texttt{NextIterator} (\ref{NextIterator}), \texttt{IsDoneIterator} (\ref{IsDoneIterator}), and \texttt{ShallowCopy} (\ref{ShallowCopy}) are computed via prescribed functions. 

 Let \mbox{\texttt{\mdseries\slshape record}} be a record with at least the following components. 
\begin{description}
\item[{\texttt{NextIterator}}]  a function taking one argument \mbox{\texttt{\mdseries\slshape iter}}, which returns the next element of \mbox{\texttt{\mdseries\slshape iter}} (see{\nobreakspace}\texttt{NextIterator} (\ref{NextIterator})); for that, the components of \mbox{\texttt{\mdseries\slshape iter}} are changed, 
\item[{\texttt{IsDoneIterator}}]  a function taking one argument \mbox{\texttt{\mdseries\slshape iter}}, which returns the \texttt{IsDoneIterator} (\ref{IsDoneIterator}) value of \mbox{\texttt{\mdseries\slshape iter}}, 
\item[{\texttt{ShallowCopy}}]  a function taking one argument \mbox{\texttt{\mdseries\slshape iter}}, which returns a record for which \texttt{IteratorByFunctions} can be called in order to create a new iterator that is independent of \mbox{\texttt{\mdseries\slshape iter}} but behaves like \mbox{\texttt{\mdseries\slshape iter}} w.r.t. the operations \texttt{NextIterator} (\ref{NextIterator}) and \texttt{IsDoneIterator} (\ref{IsDoneIterator}). 
\item[{\texttt{ViewObj} and \texttt{PrintObj}}]  two functions that print what one wants to be printed when \texttt{View( \mbox{\texttt{\mdseries\slshape iter}} )} or \texttt{Print( \mbox{\texttt{\mdseries\slshape item}} )} is called (see{\nobreakspace}\ref{View and Print}), if the \texttt{ViewObj} component is missing then the \texttt{PrintObj} method is used as a default. 
\end{description}
 Further (data) components may be contained in \mbox{\texttt{\mdseries\slshape record}} which can be used by these function. 

 \texttt{IteratorByFunctions} does \emph{not} make a shallow copy of \mbox{\texttt{\mdseries\slshape record}}, this record is changed in place (see Section {\nobreakspace}\ref{Creating Objects}). 

 Iterators constructed with \texttt{IteratorByFunctions} are in the filter \texttt{IsStandardIterator} (\ref{IsStandardIterator}). }

 }

 }

            
\chapter{\textcolor{Chapter }{Domains and their Elements}}\label{Domains and their Elements}
\logpage{[ 31, 0, 0 ]}
\hyperdef{L}{X7E651AC287AFDCC1}{}
{
  \emph{Domain} is \textsf{GAP}'s name for structured sets. The ring of Gaussian integers ${\ensuremath{\mathbb Z}}[\sqrt{{-1}}]$ is an example of a domain, the group $D_{12}$ of symmetries of a regular hexahedron is another. 

 The \textsf{GAP} library predefines some domains. For example the ring of Gaussian integers is
predefined as \texttt{GaussianIntegers} (\ref{GaussianIntegers}) (see{\nobreakspace}\ref{Gaussians}) and the field of rationals is predefined as \texttt{Rationals} (\ref{Rationals}) (see{\nobreakspace}\ref{Rational Numbers}). Most domains are constructed by functions, which are called \emph{domain constructors} (see{\nobreakspace}\ref{Constructing Domains}). For example the group $D_{12}$ is constructed by the construction \texttt{Group( (1,2,3,4,5,6), (2,6)(3,5) )} (see{\nobreakspace}\texttt{Group} (\ref{Group:for several generators})) and the finite field with 16 elements is constructed by \texttt{GaloisField( 16 )} (see{\nobreakspace}\texttt{GaloisField} (\ref{GaloisField:for field size})). 

 The first place where you need domains in \textsf{GAP} is the obvious one. Sometimes you simply want to deal with a domain. For
example if you want to compute the size of the group $D_{12}$, you had better be able to represent this group in a way that the \texttt{Size} (\ref{Size}) function can understand. 

 The second place where you need domains in \textsf{GAP} is when you want to be able to specify that an operation or computation takes
place in a certain domain. For example suppose you want to factor 10 in the
ring of Gaussian integers. Saying \texttt{Factors( 10 )} will not do, because this will return the factorization \texttt{[ 2, 5 ]} in the ring of integers. To allow operations and computations to happen in a
specific domain, \texttt{Factors} (\ref{Factors}), and many other functions as well, accept this domain as optional first
argument. Thus \texttt{Factors( GaussianIntegers, 10 )} yields the desired result \texttt{[ 1+E(4), 1\texttt{\symbol{45}}E(4), 2+E(4), 2\texttt{\symbol{45}}E(4) ]}. (The imaginary unit $\sqrt{{-1}}$ is written as \texttt{E(4)} in \textsf{GAP}, see \texttt{E} (\ref{E}).) 

 An introduction to the most important facts about domains is given in
Chapter{\nobreakspace} (\textbf{Tutorial: Domains}). 

 There are only few \emph{operations} especially for domains (see{\nobreakspace}\ref{Operations for Domains}), operations such as \texttt{Intersection} (\ref{Intersection:for a list}) and \texttt{Random} (\ref{Random:for a list or collection}) are defined for the more general situation of collections (see
Chapter{\nobreakspace}\ref{Collections}).    
\section{\textcolor{Chapter }{Operational Structure of Domains}}\label{Operational Structure of Domains}
\logpage{[ 31, 1, 0 ]}
\hyperdef{L}{X859C7AB97B34F55F}{}
{
  Domains have an \emph{operational structure}, that is, a collection of operations under which the domain is closed. For
example, a group is closed under multiplication, taking the zeroth power of
elements, and taking inverses of elements. The operational structure may be
empty, examples of domains without additional structure are the underlying
relations of general mappings (see{\nobreakspace}\ref{Properties and Attributes of (General) Mappings}). 

 The operations under which a domain is closed are a subset of the operations
that the elements of a domain admit. It is possible that the elements admit
more operations. For example, matrices can be multiplied and added. But
addition plays no role in a group of matrices, and multiplication plays no
role in a vector space of matrices. In particular, a matrix group is not
closed under addition. 

 Note that the elements of a domain exist independently of this domain, usually
they existed already before the domain was created. So it makes sense to say
that a domain is \emph{generated} by some elements with respect to certain operations. 

 Of course, different sets of operations yield different notions of generation.
For example, the group generated by some matrices is different from the ring
generated by these matrices, and these two will in general be different from
the vector space generated by the same matrices, over a suitable field. 

 The other way round, the same set of elements may be obtained by generation
w.r.t.{\nobreakspace}different notions of generation. For example, one can get
the group generated by two elements $g$ and $h$ also as the monoid generated by the elements $g$, $g^{{-1}}$, $h$, $h^{{-1}}$; if both $g$ and $h$ have finite order then of course the group generated by $g$ and $h$ coincides with the monoid generated by $g$ and $h$. 

 Additionally to the operational structure, a domain can have properties. For
example, the multiplication of a group is associative, and the multiplication
in a field is commutative. 

 Note that associativity and commutativity depend on the set of elements for
which one considers the multiplication, i.e., it depends on the domain. For
example, the multiplication in a full matrix ring over a field is not
commutative, whereas its restriction to the set of diagonal matrices is
commutative. 

 One important difference between the operational structure and the properties
of a domain is that the operational structure is fixed when the domain is
constructed, whereas properties can be discovered later. For example, take a
domain whose operational structure is given by closure under multiplication.
If it is discovered that the inverses of all its elements also do (by chance)
lie in this domain, being closed under taking inverses is \emph{not} added to the operational structure. But a domain with operational structure of
multiplication, taking the identity, and taking inverses will be treated as a
group as soon as the multiplication is found out to be associative for this
domain. 

 The operational structures available in \textsf{GAP} form a hierarchy, which is explicitly formulated in terms of domain
categories, see{\nobreakspace}\ref{Domain Categories}. }

  
\section{\textcolor{Chapter }{Equality and Comparison of Domains}}\label{Equality and Comparison of Domains}
\logpage{[ 31, 2, 0 ]}
\hyperdef{L}{X84FA03F87A17B059}{}
{
  \emph{Equality} and \emph{comparison} of domains are defined as follows. 

 Two domains are considered \emph{equal} if and only if the sets of their elements as computed by \texttt{AsSSortedList} (\ref{AsSSortedList})) are equal. Thus, in general \texttt{=} behaves as if each domain operand were replaced by its set of elements. Except
that \texttt{=} will also sometimes, but not always, work for infinite domains, for which of
course \textsf{GAP} cannot compute the set of elements. Note that this implies that domains with
different algebraic structure may well be equal. As a special case of this,
either operand of \texttt{=} may also be a proper set (see{\nobreakspace}\ref{Sorted Lists and Sets}), i.e., a sorted list without holes or duplicates (see \texttt{AsSSortedList} (\ref{AsSSortedList})), and \texttt{=} will return \texttt{true} if and only if this proper set is equal to the set of elements of the argument
that is a domain. 

    \emph{No} general \emph{ordering} of arbitrary domains via \texttt{{\textless}} is defined in \textsf{GAP}{\nobreakspace}4. This is because a well\texttt{\symbol{45}}defined \texttt{{\textless}} for domains or, more general, for collections, would have to be compatible
with \texttt{=} and would need to be transitive and antisymmetric in order to be used to form
ordered sets. In particular, \texttt{{\textless}} would have to be independent of the algebraic structure of its arguments
because this holds for \texttt{=}, and thus there would be hardly a situation where one could implement an
efficient comparison method. (Note that in the case that two domains are
comparable with \texttt{{\textless}}, the result is in general \emph{not} compatible with the set theoretical subset relation, which can be decided with \texttt{IsSubset} (\ref{IsSubset}).) }

  
\section{\textcolor{Chapter }{Constructing Domains}}\label{Constructing Domains}
\logpage{[ 31, 3, 0 ]}
\hyperdef{L}{X82039A218274826F}{}
{
  For several operational structures (see{\nobreakspace}\ref{Operational Structure of Domains}), \textsf{GAP} provides functions to construct domains with this structure (note that such
functions do not exist for all operational structures). For example, \texttt{Group} (\ref{Group:for several generators}) returns groups, \texttt{VectorSpace} (\ref{VectorSpace}) returns vector spaces etc.: 

 \index{Struct@\mbox{\texttt{\mdseries\slshape Struct}}} \texttt{\mbox{\texttt{\mdseries\slshape Struct}}( \mbox{\texttt{\mdseries\slshape arg1}}, \mbox{\texttt{\mdseries\slshape arg2}}, ... )} 

 The syntax of these functions may vary, dependent on the structure in
question. Usually a domain is constructed as the closure of some elements
under the given operations, that is, the domain is given by its \emph{generators}. For example, a group can be constructed from a list of generating
permutations or matrices or whatever is admissible as group elements, and a
vector space over a given field $F$ can be constructed from $F$ and a list of appropriate vectors. 

 The idea of generation and generators in \textsf{GAP} is that the domain returned by a function such as \texttt{Group}, \texttt{Algebra}, or \texttt{FreeLeftModule} \emph{contains} the given generators. This implies that the generators of a group must know
how they are multiplied and inverted, the generators of a module must know how
they are added and how scalar multiplication works, and so on. Thus one cannot
use for example permutations as generators of a vector space. 

 The function \mbox{\texttt{\mdseries\slshape Struct}} first checks whether the arguments admit the construction of a domain with the
desired structure. This is done by calling the operation 

 \texttt{IsGeneratorsOf\mbox{\texttt{\mdseries\slshape Struct}}( [\mbox{\texttt{\mdseries\slshape info}}, ]\mbox{\texttt{\mdseries\slshape gens}} )} 

 \index{IsGeneratorsOfStruct@\texttt{IsGeneratorsOf\mbox{\texttt{\mdseries\slshape Struct}}}} 

 where \mbox{\texttt{\mdseries\slshape arglist}} is the list of given generators and \mbox{\texttt{\mdseries\slshape info}} an argument of \mbox{\texttt{\mdseries\slshape Struct}}, for example the field of scalars in the case that a vector space shall be
constructed. If the check failed then \mbox{\texttt{\mdseries\slshape Struct}} returns \texttt{fail}, otherwise it returns the result of \texttt{\mbox{\texttt{\mdseries\slshape Struct}}ByGenerators} (see below). (So if one wants to omit the check then one should call \texttt{\mbox{\texttt{\mdseries\slshape Struct}}ByGenerators} directly.) 

 \index{GeneratorsOfStruct@\texttt{GeneratorsOf\mbox{\texttt{\mdseries\slshape Struct}}}} \texttt{GeneratorsOf\mbox{\texttt{\mdseries\slshape Struct}}( \mbox{\texttt{\mdseries\slshape D}})} 

 For a domain \mbox{\texttt{\mdseries\slshape D}} with operational structure corresponding to \mbox{\texttt{\mdseries\slshape Struct}}, the attribute \texttt{GeneratorsOf\mbox{\texttt{\mdseries\slshape Struct}}} returns a list of corresponding generators of \mbox{\texttt{\mdseries\slshape D}}. If these generators were not yet stored in \mbox{\texttt{\mdseries\slshape D}} then \mbox{\texttt{\mdseries\slshape D}} must know \emph{some} generators if \texttt{GeneratorsOf\mbox{\texttt{\mdseries\slshape Struct}}} shall have a chance to compute the desired result; for example, monoid
generators of a group can be computed from known group generators (and vice
versa). Note that several notions of generation may be meaningful for a given
domain, so it makes no sense to ask for ``the generators of a domain''. Further note that the generators may depend on other information about \mbox{\texttt{\mdseries\slshape D}}. For example the generators of a vector space depend on the underlying field
of scalars; the vector space generators of a vector space over the field with
four elements need not generate the same vector space when this is viewed as a
space over the field with two elements. 

 \index{StructByGenerators@\texttt{\mbox{\texttt{\mdseries\slshape Struct}}ByGenerators}} \texttt{\mbox{\texttt{\mdseries\slshape Struct}}ByGenerators( [\mbox{\texttt{\mdseries\slshape info}}, ]\mbox{\texttt{\mdseries\slshape gens}} )} 

 Domain construction from generators \mbox{\texttt{\mdseries\slshape gens}} is implemented by operations \texttt{\mbox{\texttt{\mdseries\slshape Struct}}ByGenerators}, which are called by the simple functions \mbox{\texttt{\mdseries\slshape Struct}}; methods can be installed only for the operations. Note that additional
information \mbox{\texttt{\mdseries\slshape info}} may be necessary to construct the domain; for example, a vector space needs
the underlying field of scalars in addition to the list of vector space
generators. The \texttt{GeneratorsOf\mbox{\texttt{\mdseries\slshape Struct}}} value of the returned domain need \emph{not} be equal to \mbox{\texttt{\mdseries\slshape gens}}. But if a domain \mbox{\texttt{\mdseries\slshape D}} is printed as \texttt{\mbox{\texttt{\mdseries\slshape Struct}}([\mbox{\texttt{\mdseries\slshape a}}, \mbox{\texttt{\mdseries\slshape b}}, ...])} and if there is an attribute \texttt{GeneratorsOf\mbox{\texttt{\mdseries\slshape Struct}}} then the list \texttt{GeneratorsOf\mbox{\texttt{\mdseries\slshape Struct}}( \mbox{\texttt{\mdseries\slshape D}} )} is guaranteed to be equal to \texttt{[ \mbox{\texttt{\mdseries\slshape a}}, \mbox{\texttt{\mdseries\slshape b}}, ... ]}. 

 \index{StructWithGenerators@\texttt{\mbox{\texttt{\mdseries\slshape Struct}}WithGenerators}} \texttt{\mbox{\texttt{\mdseries\slshape Struct}}WithGenerators( [\mbox{\texttt{\mdseries\slshape info}}, ]\mbox{\texttt{\mdseries\slshape gens}} )} 

 The only difference between \texttt{\mbox{\texttt{\mdseries\slshape Struct}}ByGenerators} and \texttt{\mbox{\texttt{\mdseries\slshape Struct}}WithGenerators} is that the latter guarantees that the \texttt{GeneratorsOf\mbox{\texttt{\mdseries\slshape Struct}}} value of the result is equal to the given generators \mbox{\texttt{\mdseries\slshape gens}}. 

 \index{ClosureStruct@\texttt{Closure\mbox{\texttt{\mdseries\slshape Struct}}}} \texttt{Closure\mbox{\texttt{\mdseries\slshape Struct}}( \mbox{\texttt{\mdseries\slshape D}}, \mbox{\texttt{\mdseries\slshape obj}} )} 

 For constructing a domain as the closure of a given domain with an element or
another domain, one can use the operation \texttt{Closure\mbox{\texttt{\mdseries\slshape Struct}}}. It returns the smallest domain with operational structure corresponding to \mbox{\texttt{\mdseries\slshape Struct}} that contains \mbox{\texttt{\mdseries\slshape D}} as a subset and \mbox{\texttt{\mdseries\slshape obj}} as an element. }

  
\section{\textcolor{Chapter }{Changing the Structure}}\label{Changing the Structure}
\logpage{[ 31, 4, 0 ]}
\hyperdef{L}{X7EA77DE17DD8A231}{}
{
  The same set of elements can have different operational structures. For
example, it may happen that a monoid $M$ does in fact contain the inverses of all of its elements; if $M$ has not been constructed as a group (see{\nobreakspace}\ref{Domain Categories}) then it is reasonable to ask for the group that is equal to $M$. 

 \index{AsStruct@\texttt{As\mbox{\texttt{\mdseries\slshape Struct}}}} \texttt{As\mbox{\texttt{\mdseries\slshape Struct}}( [\mbox{\texttt{\mdseries\slshape info}}, ]\mbox{\texttt{\mdseries\slshape D}} )} 

 If \mbox{\texttt{\mdseries\slshape D}} is a domain that is closed under the operational structure given by \mbox{\texttt{\mdseries\slshape Struct}} then \texttt{As\mbox{\texttt{\mdseries\slshape Struct}}} returns a domain \mbox{\texttt{\mdseries\slshape E}} that consists of the same elements (that is, \texttt{\mbox{\texttt{\mdseries\slshape D}} = \mbox{\texttt{\mdseries\slshape E}}}) and that has this operational structure (that is, \texttt{Is\mbox{\texttt{\mdseries\slshape Struct}}( \mbox{\texttt{\mdseries\slshape E}} )} is \texttt{true}); if \mbox{\texttt{\mdseries\slshape D}} is not closed under the structure given by \mbox{\texttt{\mdseries\slshape Struct}} then \texttt{As\mbox{\texttt{\mdseries\slshape Struct}}} returns \texttt{fail}. 

 If additional information besides generators are necessary to define \mbox{\texttt{\mdseries\slshape D}} then the argument \mbox{\texttt{\mdseries\slshape info}} describes the value of this information for the desired domain. For example,
if we want to view \mbox{\texttt{\mdseries\slshape D}} as a vector space over the field with two elements then we may call \texttt{AsVectorSpace( GF(2), \mbox{\texttt{\mdseries\slshape D}} )}; this allows us to change the underlying field of scalars, for example if \mbox{\texttt{\mdseries\slshape D}} is a vector space over the field with four elements. Again, if \mbox{\texttt{\mdseries\slshape D}} is not equal to a domain with the desired structure and additional information
then \texttt{fail} is returned. 

 In the case that no additional information \mbox{\texttt{\mdseries\slshape info}} is related to the structure given by \mbox{\texttt{\mdseries\slshape Struct}}, the operation \texttt{As\mbox{\texttt{\mdseries\slshape Struct}}} is in fact an attribute (see{\nobreakspace}\ref{Attributes}). 

 See the index of the \textsf{GAP} Reference Manual for an overview of the available \texttt{As\mbox{\texttt{\mdseries\slshape Struct}}} functions. }

  
\section{\textcolor{Chapter }{Changing the Representation}}\label{Changing the Representation}
\logpage{[ 31, 5, 0 ]}
\hyperdef{L}{X860FCCBE7A41412F}{}
{
  Often it is useful to answer questions about a domain via computations in a
different but isomorphic domain. In the sense that this approach keeps the
structure and changes the underlying set of elements, it can be viewed as a
counterpart of keeping the set of elements and changing its structure
(see{\nobreakspace}\ref{Changing the Structure}). 

 One reason for doing so can be that computations with the elements in the
given domain are not very efficient. For example, if one is given a solvable
matrix group (see Chapter{\nobreakspace}\ref{Matrix Groups}) then one can compute an isomorphism to a polycyclicly presented group $G$, say (see Chapter{\nobreakspace}\ref{Polycyclic Groups}); the multiplication of two matrices {\textendash}which is essentially
determined by the dimension of the matrices{\textendash} is much more
expensive than the multiplication of two elements in $G$ {\textendash}which is essentially determined by the composition length of $G$. 

 \index{IsomorphismRepStruct@\texttt{Isomorphism\mbox{\texttt{\mdseries\slshape Rep}}\mbox{\texttt{\mdseries\slshape Struct}}}} \texttt{Isomorphism\mbox{\texttt{\mdseries\slshape Rep}}\mbox{\texttt{\mdseries\slshape Struct}}( \mbox{\texttt{\mdseries\slshape D}} )} 

 If \mbox{\texttt{\mdseries\slshape D}} is a domain that is closed under the operational structure given by \mbox{\texttt{\mdseries\slshape Struct}} then \texttt{Isomorphism\mbox{\texttt{\mdseries\slshape Rep}}\mbox{\texttt{\mdseries\slshape Struct}}} returns a mapping \mbox{\texttt{\mdseries\slshape hom}} from \mbox{\texttt{\mdseries\slshape D}} to a domain $E$ having structure given by \mbox{\texttt{\mdseries\slshape Struct}}, such that \mbox{\texttt{\mdseries\slshape hom}} respects the structure \mbox{\texttt{\mdseries\slshape Struct}} and \mbox{\texttt{\mdseries\slshape Rep}} describes the representation of the elements in $E$. If no domain $E$ with the required properties exists then \texttt{fail} is returned. 

 For example, \texttt{IsomorphismPermGroup} (\ref{IsomorphismPermGroup}) takes a group as its argument and returns a group homomorphism
(see{\nobreakspace}\ref{Group Homomorphisms}) onto an isomorphic permutation group (see Chapter{\nobreakspace}\ref{Permutation Groups}) provided the original group is finite; for infinite groups, \texttt{IsomorphismPermGroup} (\ref{IsomorphismPermGroup}) returns \texttt{fail}. Similarly, \texttt{IsomorphismPcGroup} (\ref{IsomorphismPcGroup}) returns a group homomorphism from its argument to a polycyclicly presented
group (see{\nobreakspace}\ref{Pc Groups}) if the argument is polycyclic, and \texttt{fail} otherwise. 

 See the index of the \textsf{GAP} Reference Manual for an overview of the available \texttt{Isomorphism\mbox{\texttt{\mdseries\slshape Rep}}\mbox{\texttt{\mdseries\slshape Struct}}} functions. }

  
\section{\textcolor{Chapter }{Domain Categories}}\label{Domain Categories}
\logpage{[ 31, 6, 0 ]}
\hyperdef{L}{X7D72F11B82F4A036}{}
{
  As mentioned in{\nobreakspace}\ref{Operational Structure of Domains}, the operational structure of a domain is fixed when the domain is
constructed. For example, if \mbox{\texttt{\mdseries\slshape D}} was constructed by \texttt{Monoid} (\ref{Monoid:for various generators}) then \mbox{\texttt{\mdseries\slshape D}} is in general not regarded as a group in \textsf{GAP}, even if \mbox{\texttt{\mdseries\slshape D}} is in fact closed under taking inverses. In this case, \texttt{IsGroup} (\ref{IsGroup}) returns \texttt{false} for \mbox{\texttt{\mdseries\slshape D}}. The operational structure determines which operations are applicable for a
domain, so for example \texttt{SylowSubgroup} (\ref{SylowSubgroup}) is not defined for \mbox{\texttt{\mdseries\slshape D}} and therefore will signal an error. 

 \index{IsStruct@\texttt{Is\mbox{\texttt{\mdseries\slshape Struct}}}} \texttt{Is\mbox{\texttt{\mdseries\slshape Struct}}( \mbox{\texttt{\mdseries\slshape D}} )} 

 The functions \texttt{Is\mbox{\texttt{\mdseries\slshape Struct}}} implement the tests whether a domain \mbox{\texttt{\mdseries\slshape D}} has the respective operational structure (upon construction). \texttt{Is\mbox{\texttt{\mdseries\slshape Struct}}} is a filter (see{\nobreakspace}\ref{Types of Objects}) that involves certain categories (see{\nobreakspace}\ref{Categories}) and usually also certain properties (see{\nobreakspace}\ref{Properties}). For example, \texttt{IsGroup} (\ref{IsGroup}) is equivalent to \texttt{IsMagmaWithInverses and IsAssociative}, the first being a category and the second being a property. 

 Implications between domain categories describe the hierarchy of operational
structures available in \textsf{GAP}. Here are some typical examples. 

 
\begin{itemize}
\item  \texttt{IsDomain} (\ref{IsDomain}) is implied by each domain category, 
\item  \texttt{IsMagma} (\ref{IsMagma}) is implied by each category that describes the closure under multiplication \texttt{*}, 
\item  \texttt{IsAdditiveMagma} (\ref{IsAdditiveMagma}) is implied by each category that describes the closure under addition \texttt{+}, 
\item  \texttt{IsMagmaWithOne} (\ref{IsMagmaWithOne}) implies \texttt{IsMagma} (\ref{IsMagma}); a \emph{magma\texttt{\symbol{45}}with\texttt{\symbol{45}}one} is a magma such that each element (and thus also the magma itself) can be
asked for its zeroth power, 
\item  \texttt{IsMagmaWithInverses} (\ref{IsMagmaWithInverses}) implies \texttt{IsMagmaWithOne} (\ref{IsMagmaWithOne}); a \emph{magma\texttt{\symbol{45}}with\texttt{\symbol{45}}inverses} is a magma such that each element can be asked for its inverse; important
special cases are \emph{groups}, which in addition are associative, 
\item  a \emph{ring} is a magma that is also an additive group, 
\item  a \emph{ring\texttt{\symbol{45}}with\texttt{\symbol{45}}one} is a ring that is also a magma\texttt{\symbol{45}}with\texttt{\symbol{45}}one, 
\item  a \emph{division ring} is a ring\texttt{\symbol{45}}with\texttt{\symbol{45}}one that is also closed
under taking inverses of nonzero elements, 
\item  a \emph{field} is a commutative division ring. 
\end{itemize}
 

 Each operational structure \mbox{\texttt{\mdseries\slshape Struct}} has associated with it a domain category \texttt{Is\mbox{\texttt{\mdseries\slshape Struct}}}, and operations \texttt{\mbox{\texttt{\mdseries\slshape Struct}}ByGenerators} for constructing a domain from generators, \texttt{GeneratorsOf\mbox{\texttt{\mdseries\slshape Struct}}} for storing and accessing generators w.r.t.{\nobreakspace}this structure, \texttt{Closure\mbox{\texttt{\mdseries\slshape Struct}}} for forming the closure, and \texttt{As\mbox{\texttt{\mdseries\slshape Struct}}} for getting a domain with the desired structure from one with weaker
operational structure and for testing whether a given domain can be regarded
as a domain with \mbox{\texttt{\mdseries\slshape Struct}}. 

 The functions applicable to domains with the various structures are described
in the corresponding chapters of the Reference Manual. For example, functions
for rings, fields, groups, and vector spaces are described in
Chapters{\nobreakspace}\ref{Rings}, \ref{Fields and Division Rings}, \ref{Groups}, and \ref{Vector Spaces}, respectively. More general functions for arbitrary collections can be found
in Chapter{\nobreakspace}\ref{Collections}. }

  
\section{\textcolor{Chapter }{Parents}}\label{Parents}
\logpage{[ 31, 7, 0 ]}
\hyperdef{L}{X7CBDD36E7B7BE286}{}
{
  

\subsection{\textcolor{Chapter }{Parent}}
\logpage{[ 31, 7, 1 ]}\nobreak
\hyperdef{L}{X7BC856CC7F116BB0}{}
{\noindent\textcolor{FuncColor}{$\triangleright$\enspace\texttt{Parent({\mdseries\slshape D})\index{Parent@\texttt{Parent}}
\label{Parent}
}\hfill{\scriptsize (function)}}\\
\noindent\textcolor{FuncColor}{$\triangleright$\enspace\texttt{SetParent({\mdseries\slshape D, P})\index{SetParent@\texttt{SetParent}}
\label{SetParent}
}\hfill{\scriptsize (operation)}}\\
\noindent\textcolor{FuncColor}{$\triangleright$\enspace\texttt{HasParent({\mdseries\slshape D})\index{HasParent@\texttt{HasParent}}
\label{HasParent}
}\hfill{\scriptsize (filter)}}\\


 It is possible to assign to a domain \mbox{\texttt{\mdseries\slshape D}} one other domain \mbox{\texttt{\mdseries\slshape P}} containing \mbox{\texttt{\mdseries\slshape D}} as a subset, in order to exploit this subset relation between \mbox{\texttt{\mdseries\slshape D}} and \mbox{\texttt{\mdseries\slshape P}}. Note that \mbox{\texttt{\mdseries\slshape P}} need not have the same operational structure as \mbox{\texttt{\mdseries\slshape D}}, for example \mbox{\texttt{\mdseries\slshape P}} may be a magma and \mbox{\texttt{\mdseries\slshape D}} a field. 

 The assignment is done by calling \texttt{SetParent}, and \mbox{\texttt{\mdseries\slshape P}} is called the \emph{parent} of \mbox{\texttt{\mdseries\slshape D}}. If \mbox{\texttt{\mdseries\slshape D}} has already a parent, calls to \texttt{SetParent} will be ignored. 

 If \mbox{\texttt{\mdseries\slshape D}} has a parent \mbox{\texttt{\mdseries\slshape P}} {\textendash}this can be checked with \texttt{HasParent}{\textendash} then \mbox{\texttt{\mdseries\slshape P}} can be used to gain information about \mbox{\texttt{\mdseries\slshape D}}. First, the call of \texttt{SetParent} causes \texttt{UseSubsetRelation} (\ref{UseSubsetRelation}) to be called. Second, for a domain \mbox{\texttt{\mdseries\slshape D}} with parent, information relative to the parent can be stored in \mbox{\texttt{\mdseries\slshape D}}; for example, there is an attribute \texttt{NormalizerInParent} for storing \texttt{Normalizer( \mbox{\texttt{\mdseries\slshape P}}, \mbox{\texttt{\mdseries\slshape D}} )} in the case that \mbox{\texttt{\mdseries\slshape D}} is a group. (More about such parent dependent attributes can be found in \ref{In Parent Attributes}.)  Note that because of this relative information, one cannot change the parent;
that is, one can set the parent only once, subsequent calls to \texttt{SetParent} for the same domain \mbox{\texttt{\mdseries\slshape D}} are ignored.  Further note that contrary to \texttt{UseSubsetRelation} (\ref{UseSubsetRelation}), also knowledge about the parent \mbox{\texttt{\mdseries\slshape P}} might be used that is discovered after the \texttt{SetParent} call. 

 A stored parent can be accessed using \texttt{Parent}. If \mbox{\texttt{\mdseries\slshape D}} has no parent then \texttt{Parent} returns \mbox{\texttt{\mdseries\slshape D}} itself, and \texttt{HasParent} will return \texttt{false} also after a call to \texttt{Parent}. So \texttt{Parent} is \emph{not} an attribute, the underlying attribute to store the parent is \texttt{ParentAttr}.  

 Certain functions that return domains with parent already set, for example \texttt{Subgroup} (\ref{Subgroup}), are described in Section{\nobreakspace}\ref{Constructing Subdomains}. Whenever a function has this property, the \textsf{GAP} Reference Manual states this explicitly. Note that these functions \emph{do not guarantee} a certain parent, for example \texttt{DerivedSubgroup} (\ref{DerivedSubgroup}) for a perfect group $G$ may return $G$ itself, and if $G$ had already a parent then this is not replaced by $G$. As a rule of thumb, \textsf{GAP} avoids to set a domain as its own parent, which is consistent with the
behaviour of \texttt{Parent}, at least until a parent is set explicitly with \texttt{SetParent}. 

 
\begin{Verbatim}[commandchars=!@|,fontsize=\small,frame=single,label=Example]
  !gapprompt@gap>| !gapinput@g:= Group( (1,2,3), (1,2) );; h:= Group( (1,2) );;|
  !gapprompt@gap>| !gapinput@HasParent( g );  HasParent( h );|
  false
  false
  !gapprompt@gap>| !gapinput@SetParent( h, g );|
  !gapprompt@gap>| !gapinput@Parent( g );  Parent( h );|
  Group([ (1,2,3), (1,2) ])
  Group([ (1,2,3), (1,2) ])
  !gapprompt@gap>| !gapinput@HasParent( g );  HasParent( h );|
  false
  true
\end{Verbatim}
 }

 }

  
\section{\textcolor{Chapter }{Constructing Subdomains}}\label{Constructing Subdomains}
\logpage{[ 31, 8, 0 ]}
\hyperdef{L}{X7B58FDEF80338DD6}{}
{
  \index{Subdomains} For many domains \mbox{\texttt{\mdseries\slshape D}}, there are functions that construct certain subsets \mbox{\texttt{\mdseries\slshape S}} of \mbox{\texttt{\mdseries\slshape D}} as domains with parent (see{\nobreakspace}\ref{Parents}) already set to \mbox{\texttt{\mdseries\slshape D}}. For example, if \mbox{\texttt{\mdseries\slshape G}} is a group that contains the elements in the list \mbox{\texttt{\mdseries\slshape gens}} then \texttt{Subgroup( \mbox{\texttt{\mdseries\slshape G}}, \mbox{\texttt{\mdseries\slshape gens}} )} returns a group \mbox{\texttt{\mdseries\slshape S}} that is generated by the elements in \mbox{\texttt{\mdseries\slshape gens}} and with \texttt{Parent( \mbox{\texttt{\mdseries\slshape S}} ) = \mbox{\texttt{\mdseries\slshape G}}}. 

 \index{Substruct@\texttt{Sub\mbox{\texttt{\mdseries\slshape struct}}}} \texttt{Sub\mbox{\texttt{\mdseries\slshape struct}}( \mbox{\texttt{\mdseries\slshape D}}, \mbox{\texttt{\mdseries\slshape gens}} )} 

 More general, if \mbox{\texttt{\mdseries\slshape D}} is a domain whose algebraic structure is given by the function \mbox{\texttt{\mdseries\slshape Struct}} (for example \texttt{Group}, \texttt{Algebra}, \texttt{Field}) then the function \texttt{Sub\mbox{\texttt{\mdseries\slshape struct}}} (for example \texttt{Subgroup}, \texttt{Subalgebra}, \texttt{Subfield}) returns domains with structure \mbox{\texttt{\mdseries\slshape Struct}} and parent set to the first argument. 

 \index{SubstructNC@\texttt{Sub\mbox{\texttt{\mdseries\slshape struct}}NC}} \texttt{Sub\mbox{\texttt{\mdseries\slshape struct}}NC( \mbox{\texttt{\mdseries\slshape D}}, \mbox{\texttt{\mdseries\slshape gens}} )} 

 Each function \texttt{Sub\mbox{\texttt{\mdseries\slshape struct}}} checks that the \mbox{\texttt{\mdseries\slshape Struct}} generated by \mbox{\texttt{\mdseries\slshape gens}} is in fact a subset of \mbox{\texttt{\mdseries\slshape D}}. If one wants to omit this check then one can call \texttt{Sub\mbox{\texttt{\mdseries\slshape struct}}NC} instead; the suffix \texttt{NC} stands for ``no check''. 

 \index{AsSubstruct@\texttt{AsSub\mbox{\texttt{\mdseries\slshape struct}}}} \texttt{AsSub\mbox{\texttt{\mdseries\slshape struct}}( \mbox{\texttt{\mdseries\slshape D}}, \mbox{\texttt{\mdseries\slshape S}} )} 

 first constructs \texttt{As\mbox{\texttt{\mdseries\slshape Struct}}( [\mbox{\texttt{\mdseries\slshape info}}, ]\mbox{\texttt{\mdseries\slshape S}} )}, where \mbox{\texttt{\mdseries\slshape info}} depends on \mbox{\texttt{\mdseries\slshape D}} and \mbox{\texttt{\mdseries\slshape S}}, and then sets the parent (see{\nobreakspace}\ref{Parents}) of this new domain to \mbox{\texttt{\mdseries\slshape D}}. 

 \index{IsSubstruct@\texttt{IsSub\mbox{\texttt{\mdseries\slshape struct}}}} \texttt{IsSub\mbox{\texttt{\mdseries\slshape struct}}( \mbox{\texttt{\mdseries\slshape D}}, \mbox{\texttt{\mdseries\slshape S}} )} 

 There is no real need for functions that check whether a domain \mbox{\texttt{\mdseries\slshape S}} is a \texttt{Sub\mbox{\texttt{\mdseries\slshape struct}}} of a domain \mbox{\texttt{\mdseries\slshape D}}, since this is equivalent to the checks whether \mbox{\texttt{\mdseries\slshape S}} is a \mbox{\texttt{\mdseries\slshape Struct}} and \mbox{\texttt{\mdseries\slshape S}} is a subset of \mbox{\texttt{\mdseries\slshape D}}. Note that in many cases, only the subset relation is what one really wants
to check, and that appropriate methods for the operation \texttt{IsSubset} (\ref{IsSubset}) are available for many special situations, such as the test whether a group is
contained in another group, where only generators need to be checked. 

 If a function \texttt{IsSub\mbox{\texttt{\mdseries\slshape struct}}} is available in \textsf{GAP} then it is implemented as first a call to \texttt{Is\mbox{\texttt{\mdseries\slshape Struct}}} for the second argument and then a call to \texttt{IsSubset} (\ref{IsSubset}) for the two arguments. }

  
\section{\textcolor{Chapter }{Operations for Domains}}\label{Operations for Domains}
\logpage{[ 31, 9, 0 ]}
\hyperdef{L}{X86D579707B112970}{}
{
  For the meaning of the attributes \texttt{Characteristic} (\ref{Characteristic}), \texttt{One} (\ref{One}), \texttt{Zero} (\ref{Zero}) in the case of a domain argument, see{\nobreakspace}\ref{Attributes and Properties of Elements}. 

\subsection{\textcolor{Chapter }{IsGeneralizedDomain}}
\logpage{[ 31, 9, 1 ]}\nobreak
\hyperdef{L}{X86B4AC017FAF4D12}{}
{\noindent\textcolor{FuncColor}{$\triangleright$\enspace\texttt{IsGeneralizedDomain({\mdseries\slshape obj})\index{IsGeneralizedDomain@\texttt{IsGeneralizedDomain}}
\label{IsGeneralizedDomain}
}\hfill{\scriptsize (Category)}}\\
\noindent\textcolor{FuncColor}{$\triangleright$\enspace\texttt{IsDomain({\mdseries\slshape obj})\index{IsDomain@\texttt{IsDomain}}
\label{IsDomain}
}\hfill{\scriptsize (Category)}}\\


 For some purposes, it is useful to deal with objects that are similar to
domains but that are not collections in the sense of \textsf{GAP} because their elements may lie in different families; such objects are called \emph{generalized domains}. An instance of generalized domains are ``operation domains'', for example any $G$\texttt{\symbol{45}}set for a permutation group $G$ consisting of some union of points, sets of points, sets of sets of points
etc., under a suitable action. 

 \texttt{IsDomain} is a synonym for \texttt{IsGeneralizedDomain and IsCollection}. }

 

\subsection{\textcolor{Chapter }{GeneratorsOfDomain}}
\logpage{[ 31, 9, 2 ]}\nobreak
\hyperdef{L}{X7E353DD1838AB223}{}
{\noindent\textcolor{FuncColor}{$\triangleright$\enspace\texttt{GeneratorsOfDomain({\mdseries\slshape D})\index{GeneratorsOfDomain@\texttt{GeneratorsOfDomain}}
\label{GeneratorsOfDomain}
}\hfill{\scriptsize (attribute)}}\\


 For a domain \mbox{\texttt{\mdseries\slshape D}}, \texttt{GeneratorsOfDomain} returns a list containing generators of \mbox{\texttt{\mdseries\slshape D}} with respect to the trivial operational structure, that is interpreting \mbox{\texttt{\mdseries\slshape D}} as a set. The returned list may contain repetitions. 

 See{\nobreakspace}\ref{Constructing Domains} and for \texttt{GeneratorsOf\mbox{\texttt{\mdseries\slshape Struct}}} methods with respect to other available operational structures. 

 For many domains that have \emph{natural generators by construction} (for example, the natural generators of a free group of rank two are the two
generators stored as value of the attribute \texttt{GeneratorsOfGroup} (\ref{GeneratorsOfGroup}), and the natural generators of a free associative algebra are those
generators stored as value of the attribute \texttt{GeneratorsOfAlgebra} (\ref{GeneratorsOfAlgebra})), each \emph{natural} generator can be accessed using the \texttt{.} operator. For a domain \mbox{\texttt{\mdseries\slshape D}}, \texttt{\mbox{\texttt{\mdseries\slshape D}}.i} returns the $i$\texttt{\symbol{45}}th generator if $i$ is a positive integer, and if \texttt{name} is the name of a generator of \mbox{\texttt{\mdseries\slshape D}} then \texttt{\mbox{\texttt{\mdseries\slshape D}}.name} returns this generator. 

 
\begin{Verbatim}[commandchars=!@|,fontsize=\small,frame=single,label=Example]
  !gapprompt@gap>| !gapinput@G := DihedralGroup(IsPermGroup, 4);;|
  !gapprompt@gap>| !gapinput@GeneratorsOfGroup(G);|
  [ (1,2), (3,4) ]
  !gapprompt@gap>| !gapinput@GeneratorsOfDomain(G);|
  [ (), (3,4), (1,2), (1,2)(3,4) ]
  !gapprompt@gap>| !gapinput@F := FreeGroup("x");|
  <free group on the generators [ x ]>
  !gapprompt@gap>| !gapinput@GeneratorsOfGroup(F);|
  [ x ]
  !gapprompt@gap>| !gapinput@GeneratorsOfDomain(F);|
  Error, resulting list would be too large (length infinity)
\end{Verbatim}
 }

 

\subsection{\textcolor{Chapter }{Domain}}
\logpage{[ 31, 9, 3 ]}\nobreak
\hyperdef{L}{X826A21287FD3ACC0}{}
{\noindent\textcolor{FuncColor}{$\triangleright$\enspace\texttt{Domain({\mdseries\slshape [Fam, ]generators})\index{Domain@\texttt{Domain}}
\label{Domain}
}\hfill{\scriptsize (function)}}\\
\noindent\textcolor{FuncColor}{$\triangleright$\enspace\texttt{DomainByGenerators({\mdseries\slshape Fam, generators})\index{DomainByGenerators@\texttt{DomainByGenerators}}
\label{DomainByGenerators}
}\hfill{\scriptsize (operation)}}\\


 \texttt{Domain} returns the domain consisting of the elements in the homogeneous list \mbox{\texttt{\mdseries\slshape generators}}. If \mbox{\texttt{\mdseries\slshape generators}} is empty then a family \mbox{\texttt{\mdseries\slshape Fam}} must be entered as the first argument, and the returned (empty) domain lies in
the collections family of \mbox{\texttt{\mdseries\slshape Fam}}. 

 \texttt{DomainByGenerators} is the operation called by \texttt{Domain}. }

 }

  
\section{\textcolor{Chapter }{Attributes and Properties of Elements}}\label{Attributes and Properties of Elements}
\logpage{[ 31, 10, 0 ]}
\hyperdef{L}{X7C2B0C1280237CB0}{}
{
  The following attributes and properties for elements and domains correspond to
the operational structure. 

\subsection{\textcolor{Chapter }{Characteristic}}
\logpage{[ 31, 10, 1 ]}\nobreak
\hyperdef{L}{X81278E53800BF64D}{}
{\noindent\textcolor{FuncColor}{$\triangleright$\enspace\texttt{Characteristic({\mdseries\slshape obj})\index{Characteristic@\texttt{Characteristic}}
\label{Characteristic}
}\hfill{\scriptsize (attribute)}}\\


 \texttt{Characteristic} returns the \emph{characteristic} of \mbox{\texttt{\mdseries\slshape obj}}. 

 If \mbox{\texttt{\mdseries\slshape obj}} is a family, all of whose elements lie in \texttt{IsAdditiveElementWithZero} (\ref{IsAdditiveElementWithZero}) then its characteristic is the least positive integer $n$, if any, such that \texttt{IsZero(n*x)} is \texttt{true} for all \texttt{x} in the family \mbox{\texttt{\mdseries\slshape obj}}, otherwise it is $0$. 

 If \mbox{\texttt{\mdseries\slshape obj}} is a collections family of a family $g$ which has a characteristic, then the characteristic of \mbox{\texttt{\mdseries\slshape obj}} is the same as the characteristic of $g$. 

 For other families \mbox{\texttt{\mdseries\slshape obj}} the characteristic is not defined and \texttt{fail} will be returned. 

 For any object \mbox{\texttt{\mdseries\slshape obj}} which is in the filter \texttt{IsAdditiveElementWithZero} (\ref{IsAdditiveElementWithZero}) or in the filter \texttt{IsAdditiveMagmaWithZero} (\ref{IsAdditiveMagmaWithZero}) the characteristic of \mbox{\texttt{\mdseries\slshape obj}} is the same as the characteristic of its family if that is defined and
undefined otherwise. 

 For all other objects \mbox{\texttt{\mdseries\slshape obj}} the characteristic is undefined and may return \texttt{fail} or a ``no method found'' error. }

 

\subsection{\textcolor{Chapter }{OneImmutable}}
\logpage{[ 31, 10, 2 ]}\nobreak
\hyperdef{L}{X8046262384895B2A}{}
{\noindent\textcolor{FuncColor}{$\triangleright$\enspace\texttt{OneImmutable({\mdseries\slshape obj})\index{OneImmutable@\texttt{OneImmutable}}
\label{OneImmutable}
}\hfill{\scriptsize (attribute)}}\\
\noindent\textcolor{FuncColor}{$\triangleright$\enspace\texttt{One({\mdseries\slshape obj})\index{One@\texttt{One}}
\label{One}
}\hfill{\scriptsize (attribute)}}\\
\noindent\textcolor{FuncColor}{$\triangleright$\enspace\texttt{Identity({\mdseries\slshape obj})\index{Identity@\texttt{Identity}}
\label{Identity}
}\hfill{\scriptsize (attribute)}}\\
\noindent\textcolor{FuncColor}{$\triangleright$\enspace\texttt{OneMutable({\mdseries\slshape obj})\index{OneMutable@\texttt{OneMutable}}
\label{OneMutable}
}\hfill{\scriptsize (operation)}}\\
\noindent\textcolor{FuncColor}{$\triangleright$\enspace\texttt{OneOp({\mdseries\slshape obj})\index{OneOp@\texttt{OneOp}}
\label{OneOp}
}\hfill{\scriptsize (operation)}}\\
\noindent\textcolor{FuncColor}{$\triangleright$\enspace\texttt{OneSameMutability({\mdseries\slshape obj})\index{OneSameMutability@\texttt{OneSameMutability}}
\label{OneSameMutability}
}\hfill{\scriptsize (operation)}}\\


 \texttt{OneImmutable}, \texttt{OneMutable}, and \texttt{OneSameMutability} return the multiplicative neutral element of the multiplicative element \mbox{\texttt{\mdseries\slshape obj}}. 

 They differ only w.r.t. the mutability of the result. \texttt{OneImmutable} is an attribute and hence returns an immutable result. \texttt{OneMutable} is guaranteed to return a new \emph{mutable} object whenever a mutable version of the required element exists in \textsf{GAP} (see{\nobreakspace}\texttt{IsCopyable} (\ref{IsCopyable})). \texttt{OneSameMutability} returns a result that is mutable if \mbox{\texttt{\mdseries\slshape obj}} is mutable and if a mutable version of the required element exists in \textsf{GAP}; for lists, it returns a result of the same immutability level as the
argument. For instance, if the argument is a mutable matrix with immutable
rows, it returns a similar object. 

 If \mbox{\texttt{\mdseries\slshape obj}} is a multiplicative element then \texttt{OneSameMutability( \mbox{\texttt{\mdseries\slshape obj}} )} is equivalent to \texttt{\mbox{\texttt{\mdseries\slshape obj}}\texttt{\symbol{94}}0}. 

 \texttt{One} and \texttt{Identity} are synonyms of \texttt{OneImmutable}. \texttt{OneOp} is a synonym of \texttt{OneMutable}. 

 If \mbox{\texttt{\mdseries\slshape obj}} is a domain or a family then \texttt{One} is defined as the identity element of all elements in \mbox{\texttt{\mdseries\slshape obj}}, provided that all these elements have the same identity. For example, the
family of all cyclotomics has the identity element \texttt{1}, but a collections family (see{\nobreakspace}\texttt{CollectionsFamily} (\ref{CollectionsFamily})) may contain matrices of all dimensions and then it cannot have a unique
identity element. Note that \texttt{One} is applicable to a domain only if it is a
magma\texttt{\symbol{45}}with\texttt{\symbol{45}}one (see{\nobreakspace}\texttt{IsMagmaWithOne} (\ref{IsMagmaWithOne})); use \texttt{MultiplicativeNeutralElement} (\ref{MultiplicativeNeutralElement}) otherwise. 

 The identity of an object need not be distinct from its zero, so for example a
ring consisting of a single element can be regarded as a
ring\texttt{\symbol{45}}with\texttt{\symbol{45}}one (see{\nobreakspace}\ref{Rings}). This is particularly useful in the case of finitely presented algebras,
where any factor of a free
algebra\texttt{\symbol{45}}with\texttt{\symbol{45}}one is again an
algebra\texttt{\symbol{45}}with\texttt{\symbol{45}}one, no matter whether or
not it is a zero algebra. 

 The default method of \texttt{One} for multiplicative elements calls \texttt{OneMutable} (note that methods for \texttt{OneMutable} must \emph{not} delegate to \texttt{One}); so other methods to compute identity elements need to be installed only for \texttt{OneOp} and (in the case of copyable objects) \texttt{OneSameMutability}. 

 For domains, \texttt{One} may call \texttt{Representative} (\ref{Representative}), but \texttt{Representative} (\ref{Representative}) is allowed to fetch the identity of a domain \mbox{\texttt{\mdseries\slshape D}} only if \texttt{HasOne( \mbox{\texttt{\mdseries\slshape D}} )} is \texttt{true}. }

 

\subsection{\textcolor{Chapter }{ZeroImmutable}}
\logpage{[ 31, 10, 3 ]}\nobreak
\hyperdef{L}{X8040AC7A79FFC442}{}
{\noindent\textcolor{FuncColor}{$\triangleright$\enspace\texttt{ZeroImmutable({\mdseries\slshape obj})\index{ZeroImmutable@\texttt{ZeroImmutable}}
\label{ZeroImmutable}
}\hfill{\scriptsize (attribute)}}\\
\noindent\textcolor{FuncColor}{$\triangleright$\enspace\texttt{Zero({\mdseries\slshape obj})\index{Zero@\texttt{Zero}}
\label{Zero}
}\hfill{\scriptsize (attribute)}}\\
\noindent\textcolor{FuncColor}{$\triangleright$\enspace\texttt{ZeroMutable({\mdseries\slshape obj})\index{ZeroMutable@\texttt{ZeroMutable}}
\label{ZeroMutable}
}\hfill{\scriptsize (operation)}}\\
\noindent\textcolor{FuncColor}{$\triangleright$\enspace\texttt{ZeroOp({\mdseries\slshape obj})\index{ZeroOp@\texttt{ZeroOp}}
\label{ZeroOp}
}\hfill{\scriptsize (operation)}}\\
\noindent\textcolor{FuncColor}{$\triangleright$\enspace\texttt{ZeroSameMutability({\mdseries\slshape obj})\index{ZeroSameMutability@\texttt{ZeroSameMutability}}
\label{ZeroSameMutability}
}\hfill{\scriptsize (operation)}}\\


 \texttt{ZeroImmutable}, \texttt{ZeroMutable}, and \texttt{ZeroSameMutability} all return the additive neutral element of the additive element \mbox{\texttt{\mdseries\slshape obj}}. 

 They differ only w.r.t. the mutability of the result. \texttt{ZeroImmutable} is an attribute and hence returns an immutable result. \texttt{ZeroMutable} is guaranteed to return a new \emph{mutable} object whenever a mutable version of the required element exists in \textsf{GAP} (see{\nobreakspace}\texttt{IsCopyable} (\ref{IsCopyable})). \texttt{ZeroSameMutability} returns a result that is mutable if \mbox{\texttt{\mdseries\slshape obj}} is mutable and if a mutable version of the required element exists in \textsf{GAP}; for lists, it returns a result of the same immutability level as the
argument. For instance, if the argument is a mutable matrix with immutable
rows, it returns a similar object. 

 \texttt{ZeroSameMutability( \mbox{\texttt{\mdseries\slshape obj}} )} is equivalent to \texttt{0 * \mbox{\texttt{\mdseries\slshape obj}}}. 

 \texttt{Zero} is a synonym of \texttt{ZeroImmutable}. \texttt{ZeroOp} is a synonym of \texttt{ZeroMutable}. 

 If \mbox{\texttt{\mdseries\slshape obj}} is a domain or a family then \texttt{Zero} is defined as the zero element of all elements in \mbox{\texttt{\mdseries\slshape obj}}, provided that all these elements have the same zero. For example, the family
of all cyclotomics has the zero element \texttt{0}, but a collections family (see{\nobreakspace}\texttt{CollectionsFamily} (\ref{CollectionsFamily})) may contain matrices of all dimensions and then it cannot have a unique zero
element. Note that \texttt{Zero} is applicable to a domain only if it is an additive
magma\texttt{\symbol{45}}with\texttt{\symbol{45}}zero (see{\nobreakspace}\texttt{IsAdditiveMagmaWithZero} (\ref{IsAdditiveMagmaWithZero})); use \texttt{AdditiveNeutralElement} (\ref{AdditiveNeutralElement}) otherwise. 

 The default method of \texttt{Zero} for additive elements calls \texttt{ZeroMutable} (note that methods for \texttt{ZeroMutable} must \emph{not} delegate to \texttt{Zero}); so other methods to compute zero elements need to be installed only for \texttt{ZeroMutable} and (in the case of copyable objects) \texttt{ZeroSameMutability}. 

 For domains, \texttt{Zero} may call \texttt{Representative} (\ref{Representative}), but \texttt{Representative} (\ref{Representative}) is allowed to fetch the zero of a domain \mbox{\texttt{\mdseries\slshape D}} only if \texttt{HasZero( \mbox{\texttt{\mdseries\slshape D}} )} is \texttt{true}. }

 

\subsection{\textcolor{Chapter }{MultiplicativeZeroOp}}
\logpage{[ 31, 10, 4 ]}\nobreak
\hyperdef{L}{X86DEB543824C40EB}{}
{\noindent\textcolor{FuncColor}{$\triangleright$\enspace\texttt{MultiplicativeZeroOp({\mdseries\slshape elt})\index{MultiplicativeZeroOp@\texttt{MultiplicativeZeroOp}}
\label{MultiplicativeZeroOp}
}\hfill{\scriptsize (operation)}}\\
\textbf{\indent Returns:\ }
A multiplicative zero element.



 for an element \mbox{\texttt{\mdseries\slshape elt}} in the category \texttt{IsMultiplicativeElementWithZero} (\ref{IsMultiplicativeElementWithZero}), \texttt{MultiplicativeZeroOp} returns the element $z$ in the family $F$ of \mbox{\texttt{\mdseries\slshape elt}} with the property that $z * m = z = m * z$ holds for all $m \in F$, if such an element can be determined. 

 Families of elements in the category \texttt{IsMultiplicativeElementWithZero} (\ref{IsMultiplicativeElementWithZero}) often arise from adjoining a new zero to an existing magma. See{\nobreakspace}\texttt{InjectionZeroMagma} (\ref{InjectionZeroMagma}) or \texttt{MagmaWithZeroAdjoined} (\ref{MagmaWithZeroAdjoined}) for details. 
\begin{Verbatim}[commandchars=!@|,fontsize=\small,frame=single,label=Example]
  !gapprompt@gap>| !gapinput@G:=AlternatingGroup(5);;|
  !gapprompt@gap>| !gapinput@x:=Representative(MagmaWithZeroAdjoined(G));|
  <group with 0 adjoined elt: ()>
  !gapprompt@gap>| !gapinput@MultiplicativeZeroOp(x);|
  <group with 0 adjoined elt: 0>
\end{Verbatim}
 }

 

\subsection{\textcolor{Chapter }{IsOne}}
\logpage{[ 31, 10, 5 ]}\nobreak
\hyperdef{L}{X814D78347858EC13}{}
{\noindent\textcolor{FuncColor}{$\triangleright$\enspace\texttt{IsOne({\mdseries\slshape elm})\index{IsOne@\texttt{IsOne}}
\label{IsOne}
}\hfill{\scriptsize (property)}}\\


 is \texttt{true} if \texttt{\mbox{\texttt{\mdseries\slshape elm}} = One( \mbox{\texttt{\mdseries\slshape elm}} )}, and \texttt{false} otherwise. }

 

\subsection{\textcolor{Chapter }{IsZero}}
\logpage{[ 31, 10, 6 ]}\nobreak
\hyperdef{L}{X82BDA47282F9BBA7}{}
{\noindent\textcolor{FuncColor}{$\triangleright$\enspace\texttt{IsZero({\mdseries\slshape elm})\index{IsZero@\texttt{IsZero}}
\label{IsZero}
}\hfill{\scriptsize (property)}}\\


 is \texttt{true} if \texttt{\mbox{\texttt{\mdseries\slshape elm}} = Zero( \mbox{\texttt{\mdseries\slshape elm}} )}, and \texttt{false} otherwise. }

 

\subsection{\textcolor{Chapter }{IsIdempotent}}
\logpage{[ 31, 10, 7 ]}\nobreak
\hyperdef{L}{X7CB5896082D29173}{}
{\noindent\textcolor{FuncColor}{$\triangleright$\enspace\texttt{IsIdempotent({\mdseries\slshape elt})\index{IsIdempotent@\texttt{IsIdempotent}}
\label{IsIdempotent}
}\hfill{\scriptsize (property)}}\\


 returns \texttt{true} iff \mbox{\texttt{\mdseries\slshape elt}} is its own square. (Even if \texttt{IsZero} (\ref{IsZero}) returns \texttt{true} for \mbox{\texttt{\mdseries\slshape elt}}.) }

 

\subsection{\textcolor{Chapter }{InverseImmutable}}
\logpage{[ 31, 10, 8 ]}\nobreak
\hyperdef{L}{X78EE524E83624057}{}
{\noindent\textcolor{FuncColor}{$\triangleright$\enspace\texttt{InverseImmutable({\mdseries\slshape elm})\index{InverseImmutable@\texttt{InverseImmutable}}
\label{InverseImmutable}
}\hfill{\scriptsize (attribute)}}\\
\noindent\textcolor{FuncColor}{$\triangleright$\enspace\texttt{Inverse({\mdseries\slshape elm})\index{Inverse@\texttt{Inverse}}
\label{Inverse}
}\hfill{\scriptsize (attribute)}}\\
\noindent\textcolor{FuncColor}{$\triangleright$\enspace\texttt{InverseMutable({\mdseries\slshape elm})\index{InverseMutable@\texttt{InverseMutable}}
\label{InverseMutable}
}\hfill{\scriptsize (operation)}}\\
\noindent\textcolor{FuncColor}{$\triangleright$\enspace\texttt{InverseOp({\mdseries\slshape elm})\index{InverseOp@\texttt{InverseOp}}
\label{InverseOp}
}\hfill{\scriptsize (operation)}}\\
\noindent\textcolor{FuncColor}{$\triangleright$\enspace\texttt{InverseSameMutability({\mdseries\slshape elm})\index{InverseSameMutability@\texttt{InverseSameMutability}}
\label{InverseSameMutability}
}\hfill{\scriptsize (operation)}}\\


 \texttt{InverseImmutable}, \texttt{InverseMutable}, and \texttt{InverseSameMutability} all return the multiplicative inverse of an element \mbox{\texttt{\mdseries\slshape elm}}, that is, an element \mbox{\texttt{\mdseries\slshape inv}} such that \texttt{\mbox{\texttt{\mdseries\slshape elm}} * \mbox{\texttt{\mdseries\slshape inv}} = \mbox{\texttt{\mdseries\slshape inv}} * \mbox{\texttt{\mdseries\slshape elm}} = One( \mbox{\texttt{\mdseries\slshape elm}} )} holds; if \mbox{\texttt{\mdseries\slshape elm}} is not invertible then \texttt{fail} (see{\nobreakspace}\ref{Fail}) is returned. 

 Note that the above definition implies that a (general) mapping is invertible
in the sense of \texttt{Inverse} only if its source equals its range (see{\nobreakspace}\ref{Technical Matters Concerning General Mappings}). For a bijective mapping $f$ whose source and range differ, \texttt{InverseGeneralMapping} (\ref{InverseGeneralMapping}) can be used to construct a mapping $g$ with the property that $f$ \texttt{*} $g$ is the identity mapping on the source of $f$ and $g$ \texttt{*} $f$ is the identity mapping on the range of $f$. 

 The operations differ only w.r.t. the mutability of the result. \texttt{InverseImmutable} is an attribute and hence returns an immutable result. \texttt{InverseMutable} is guaranteed to return a new \emph{mutable} object whenever a mutable version of the required element exists in \textsf{GAP}. \texttt{InverseSameMutability} returns a result that is mutable if \mbox{\texttt{\mdseries\slshape elm}} is mutable and if a mutable version of the required element exists in \textsf{GAP}; for lists, it returns a result of the same immutability level as the
argument. For instance, if the argument is a mutable matrix with immutable
rows, it returns a similar object. 

 \texttt{InverseSameMutability( \mbox{\texttt{\mdseries\slshape elm}} )} is equivalent to \texttt{\mbox{\texttt{\mdseries\slshape elm}}\texttt{\symbol{94}}\texttt{\symbol{45}}1}. 

 \texttt{Inverse} is a synonym of \texttt{InverseImmutable}. \texttt{InverseOp} is a synonym of \texttt{InverseMutable}. 

 The default method of \texttt{InverseImmutable} calls \texttt{InverseMutable} (note that methods for \texttt{InverseMutable} must \emph{not} delegate to \texttt{InverseImmutable}); other methods to compute inverses need to be installed only for \texttt{InverseMutable} and (in the case of copyable objects) \texttt{InverseSameMutability}. }

 

\subsection{\textcolor{Chapter }{AdditiveInverseImmutable}}
\logpage{[ 31, 10, 9 ]}\nobreak
\hyperdef{L}{X84BB723C81D55D63}{}
{\noindent\textcolor{FuncColor}{$\triangleright$\enspace\texttt{AdditiveInverseImmutable({\mdseries\slshape elm})\index{AdditiveInverseImmutable@\texttt{AdditiveInverseImmutable}}
\label{AdditiveInverseImmutable}
}\hfill{\scriptsize (attribute)}}\\
\noindent\textcolor{FuncColor}{$\triangleright$\enspace\texttt{AdditiveInverse({\mdseries\slshape elm})\index{AdditiveInverse@\texttt{AdditiveInverse}}
\label{AdditiveInverse}
}\hfill{\scriptsize (attribute)}}\\
\noindent\textcolor{FuncColor}{$\triangleright$\enspace\texttt{AdditiveInverseMutable({\mdseries\slshape elm})\index{AdditiveInverseMutable@\texttt{AdditiveInverseMutable}}
\label{AdditiveInverseMutable}
}\hfill{\scriptsize (operation)}}\\
\noindent\textcolor{FuncColor}{$\triangleright$\enspace\texttt{AdditiveInverseOp({\mdseries\slshape elm})\index{AdditiveInverseOp@\texttt{AdditiveInverseOp}}
\label{AdditiveInverseOp}
}\hfill{\scriptsize (operation)}}\\
\noindent\textcolor{FuncColor}{$\triangleright$\enspace\texttt{AdditiveInverseSameMutability({\mdseries\slshape elm})\index{AdditiveInverseSameMutability@\texttt{AdditiveInverseSameMutability}}
\label{AdditiveInverseSameMutability}
}\hfill{\scriptsize (operation)}}\\


 \texttt{AdditiveInverseImmutable}, \texttt{AdditiveInverseMutable}, and \texttt{AdditiveInverseSameMutability} all return the additive inverse of \mbox{\texttt{\mdseries\slshape elm}}. 

 They differ only w.r.t. the mutability of the result. \texttt{AdditiveInverseImmutable} is an attribute and hence returns an immutable result. \texttt{AdditiveInverseMutable} is guaranteed to return a new \emph{mutable} object whenever a mutable version of the required element exists in \textsf{GAP} (see{\nobreakspace}\texttt{IsCopyable} (\ref{IsCopyable})). \texttt{AdditiveInverseSameMutability} returns a result that is mutable if \mbox{\texttt{\mdseries\slshape elm}} is mutable and if a mutable version of the required element exists in \textsf{GAP}; for lists, it returns a result of the same immutability level as the
argument. For instance, if the argument is a mutable matrix with immutable
rows, it returns a similar object. 

 \texttt{AdditiveInverseSameMutability( \mbox{\texttt{\mdseries\slshape elm}} )} is equivalent to \texttt{\texttt{\symbol{45}}\mbox{\texttt{\mdseries\slshape elm}}}. 

 \texttt{AdditiveInverse} is a synonym of \texttt{AdditiveInverseImmutable}. \texttt{AdditiveInverseOp} is a synonym of \texttt{AdditiveInverseMutable}. 

 The default method of \texttt{AdditiveInverse} calls \texttt{AdditiveInverseMutable} (note that methods for \texttt{AdditiveInverseMutable} must \emph{not} delegate to \texttt{AdditiveInverse}); so other methods to compute additive inverses need to be installed only for \texttt{AdditiveInverseMutable} and (in the case of copyable objects) \texttt{AdditiveInverseSameMutability}. }

 

\subsection{\textcolor{Chapter }{Order}}
\logpage{[ 31, 10, 10 ]}\nobreak
\hyperdef{L}{X84F59A2687C62763}{}
{\noindent\textcolor{FuncColor}{$\triangleright$\enspace\texttt{Order({\mdseries\slshape elm})\index{Order@\texttt{Order}}
\label{Order}
}\hfill{\scriptsize (attribute)}}\\


 is the multiplicative order of \mbox{\texttt{\mdseries\slshape elm}}. This is the smallest positive integer $n$ such that \mbox{\texttt{\mdseries\slshape elm}} \texttt{\texttt{\symbol{94}}} $n$ \texttt{= One( \mbox{\texttt{\mdseries\slshape elm}} )} if such an integer exists. If the order is infinite, \texttt{Order} may return the value \texttt{infinity} (\ref{infinity}), but it also might run into an infinite loop trying to test the order. }

 }

  
\section{\textcolor{Chapter }{Comparison Operations for Elements}}\label{Comparison Operations for Elements}
\logpage{[ 31, 11, 0 ]}
\hyperdef{L}{X7B3BC7BA7BB2646D}{}
{
  Binary comparison operations have been introduced already in \ref{Comparisons}. The underlying operations for which methods can be installed are the
following. 
\subsection{\textcolor{Chapter }{\texttt{\symbol{92}}= and \texttt{\symbol{92}}{\textless}}}\logpage{[ 31, 11, 1 ]}
\hyperdef{L}{X7EF67D047F03CA6F}{}
{
\noindent\textcolor{FuncColor}{$\triangleright$\enspace\texttt{\texttt{\symbol{92}}=({\mdseries\slshape left\texttt{\symbol{45}}expr, right\texttt{\symbol{45}}expr})\index{\texttt{\symbol{92}}=@\texttt{\texttt{\symbol{92}}=}}
\label{bSlash=}
}\hfill{\scriptsize (operation)}}\\
\noindent\textcolor{FuncColor}{$\triangleright$\enspace\texttt{\texttt{\symbol{92}}{\textless}({\mdseries\slshape left\texttt{\symbol{45}}expr, right\texttt{\symbol{45}}expr})\index{\texttt{\symbol{92}}{\textless}@\texttt{\texttt{\symbol{92}}{\textless}}}
\label{bSlash<}
}\hfill{\scriptsize (operation)}}\\


 \index{equality!operation} \index{comparison!operation} Note that the comparisons via \texttt{{\textless}{\textgreater}}, \texttt{{\textless}=}, \texttt{{\textgreater}}, and \texttt{{\textgreater}=} are delegated to the operations \texttt{\texttt{\symbol{92}}=} and \texttt{\texttt{\symbol{92}}{\textless}}. 

 In general, objects in \emph{different} families cannot be compared with \texttt{\texttt{\symbol{92}}{\textless}}. For the reason and for exceptions from this rule, see{\nobreakspace}\ref{Comparisons}. }

 

\subsection{\textcolor{Chapter }{CanEasilyCompareElements}}
\logpage{[ 31, 11, 2 ]}\nobreak
\hyperdef{L}{X7EFE013B8634D214}{}
{\noindent\textcolor{FuncColor}{$\triangleright$\enspace\texttt{CanEasilyCompareElements({\mdseries\slshape obj})\index{CanEasilyCompareElements@\texttt{CanEasilyCompareElements}}
\label{CanEasilyCompareElements}
}\hfill{\scriptsize (property)}}\\
\noindent\textcolor{FuncColor}{$\triangleright$\enspace\texttt{CanEasilyCompareElementsFamily({\mdseries\slshape fam})\index{CanEasilyCompareElementsFamily@\texttt{CanEasilyCompareElementsFamily}}
\label{CanEasilyCompareElementsFamily}
}\hfill{\scriptsize (function)}}\\
\noindent\textcolor{FuncColor}{$\triangleright$\enspace\texttt{CanEasilySortElements({\mdseries\slshape obj})\index{CanEasilySortElements@\texttt{CanEasilySortElements}}
\label{CanEasilySortElements}
}\hfill{\scriptsize (property)}}\\
\noindent\textcolor{FuncColor}{$\triangleright$\enspace\texttt{CanEasilySortElementsFamily({\mdseries\slshape fam})\index{CanEasilySortElementsFamily@\texttt{CanEasilySortElementsFamily}}
\label{CanEasilySortElementsFamily}
}\hfill{\scriptsize (function)}}\\


 For some objects a ``normal form'' is hard to compute and thus equality of elements of a domain might be
expensive to test. Therefore \textsf{GAP} provides a (slightly technical) property with which an algorithm can test
whether an efficient equality test is available for elements of a certain
kind. 

 \texttt{CanEasilyCompareElements} indicates whether the elements in the family \mbox{\texttt{\mdseries\slshape fam}} of \mbox{\texttt{\mdseries\slshape obj}} can be easily compared with \texttt{\texttt{\symbol{92}}=} (\ref{bSlash=}). 

 The default method for this property is to ask the family of \mbox{\texttt{\mdseries\slshape obj}}, the default method for the family is to return \texttt{false}. 

 The ability to compare elements may depend on the successful computation of
certain information. (For example for finitely presented groups it might
depend on the knowledge of a faithful permutation representation.) This
information might change over time and thus it might not be a good idea to
store a value \texttt{false} too early in a family. Instead the function \texttt{CanEasilyCompareElementsFamily} should be called for the family of \mbox{\texttt{\mdseries\slshape obj}} which returns \texttt{false} if the value of \texttt{CanEasilyCompareElements} is not known for the family without computing it. (This is in fact what the
above mentioned family dispatch does.) 

 If a family knows ab initio that it can compare elements this property should
be set as implied filter \emph{and} filter for the family (the 3rd and 4th argument of \texttt{NewFamily} (\ref{NewFamily}) respectively). This guarantees that code which directly asks the family gets a
right answer. 

 The property \texttt{CanEasilySortElements} and the function \texttt{CanEasilySortElementsFamily} behave exactly in the same way, except that they indicate that objects can be
compared via \texttt{\texttt{\symbol{92}}{\textless}} (\ref{bSlash<}). This property implies \texttt{CanEasilyCompareElements}, as the ordering must be total. }

 }

  
\section{\textcolor{Chapter }{Arithmetic Operations for Elements}}\label{Arithmetic Operations for Elements}
\logpage{[ 31, 12, 0 ]}
\hyperdef{L}{X7A2914307963E370}{}
{
  \emph{Binary} arithmetic operations have been introduced already in \ref{Arithmetic Operators}. The underlying operations for which methods can be installed are the
following. 
\subsection{\textcolor{Chapter }{\texttt{\symbol{92}}+, \texttt{\symbol{92}}*, \texttt{\symbol{92}}/,
\texttt{\symbol{92}}\texttt{\symbol{94}}, \texttt{\symbol{92}}mod}}\logpage{[ 31, 12, 1 ]}
\hyperdef{L}{X8481C9B97B214C23}{}
{
\noindent\textcolor{FuncColor}{$\triangleright$\enspace\texttt{\texttt{\symbol{92}}+({\mdseries\slshape left\texttt{\symbol{45}}expr, right\texttt{\symbol{45}}expr})\index{\texttt{\symbol{92}}+@\texttt{\texttt{\symbol{92}}+}}
\label{bSlash+}
}\hfill{\scriptsize (operation)}}\\
\noindent\textcolor{FuncColor}{$\triangleright$\enspace\texttt{\texttt{\symbol{92}}*({\mdseries\slshape left\texttt{\symbol{45}}expr, right\texttt{\symbol{45}}expr})\index{\texttt{\symbol{92}}*@\texttt{\texttt{\symbol{92}}*}}
\label{bSlash*}
}\hfill{\scriptsize (operation)}}\\
\noindent\textcolor{FuncColor}{$\triangleright$\enspace\texttt{\texttt{\symbol{92}}/({\mdseries\slshape left\texttt{\symbol{45}}expr, right\texttt{\symbol{45}}expr})\index{\texttt{\symbol{92}}/@\texttt{\texttt{\symbol{92}}/}}
\label{bSlash/}
}\hfill{\scriptsize (operation)}}\\
\noindent\textcolor{FuncColor}{$\triangleright$\enspace\texttt{\texttt{\symbol{92}}\texttt{\symbol{94}}({\mdseries\slshape left\texttt{\symbol{45}}expr, right\texttt{\symbol{45}}expr})\index{\texttt{\symbol{92}}\texttt{\symbol{94}}@\texttt{\texttt{\symbol{92}}\texttt{\symbol{94}}}}
\label{bSlash^}
}\hfill{\scriptsize (operation)}}\\
\noindent\textcolor{FuncColor}{$\triangleright$\enspace\texttt{\texttt{\symbol{92}}mod({\mdseries\slshape left\texttt{\symbol{45}}expr, right\texttt{\symbol{45}}expr})\index{\texttt{\symbol{92}}mod@\texttt{\texttt{\symbol{92}}mod}}
\label{bSlashmod}
}\hfill{\scriptsize (operation)}}\\


 \index{addition!operation} \index{multiplication!operation} \index{division!operation} \index{exponentiation!operation} \index{remainder!operation} For details about special methods for \texttt{\texttt{\symbol{92}}*}, \texttt{\texttt{\symbol{92}}/}, \texttt{\texttt{\symbol{92}}\texttt{\symbol{94}}} and \texttt{\texttt{\symbol{92}}mod}, consult the appropriate index entries for them.  }

 

\subsection{\textcolor{Chapter }{LeftQuotient}}
\logpage{[ 31, 12, 2 ]}\nobreak
\hyperdef{L}{X7A37082878DB3930}{}
{\noindent\textcolor{FuncColor}{$\triangleright$\enspace\texttt{LeftQuotient({\mdseries\slshape elm1, elm2})\index{LeftQuotient@\texttt{LeftQuotient}}
\label{LeftQuotient}
}\hfill{\scriptsize (operation)}}\\


 returns the product \texttt{\mbox{\texttt{\mdseries\slshape elm1}}\texttt{\symbol{94}}(\texttt{\symbol{45}}1) * \mbox{\texttt{\mdseries\slshape elm2}}}. For some types of objects (for example permutations) this product can be
evaluated more efficiently than by first inverting \mbox{\texttt{\mdseries\slshape elm1}} and then forming the product with \mbox{\texttt{\mdseries\slshape elm2}}. }

 

\subsection{\textcolor{Chapter }{Comm}}
\logpage{[ 31, 12, 3 ]}\nobreak
\hyperdef{L}{X80761843831B468E}{}
{\noindent\textcolor{FuncColor}{$\triangleright$\enspace\texttt{Comm({\mdseries\slshape elm1, elm2})\index{Comm@\texttt{Comm}}
\label{Comm}
}\hfill{\scriptsize (operation)}}\\


 returns the \emph{commutator} of \mbox{\texttt{\mdseries\slshape elm1}} and \mbox{\texttt{\mdseries\slshape elm2}}. The commutator is defined as the product $\mbox{\texttt{\mdseries\slshape elm1}}^{{-1}} * \mbox{\texttt{\mdseries\slshape elm2}}^{{-1}} * \mbox{\texttt{\mdseries\slshape elm1}} * \mbox{\texttt{\mdseries\slshape elm2}}$. 

 
\begin{Verbatim}[commandchars=!@|,fontsize=\small,frame=single,label=Example]
  !gapprompt@gap>| !gapinput@a:= (1,3)(4,6);; b:= (1,6,5,4,3,2);;|
  !gapprompt@gap>| !gapinput@Comm( a, b );|
  (1,5,3)(2,6,4)
  !gapprompt@gap>| !gapinput@LeftQuotient( a, b );|
  (1,2)(3,6)(4,5)
\end{Verbatim}
 }

 

\subsection{\textcolor{Chapter }{LieBracket}}
\logpage{[ 31, 12, 4 ]}\nobreak
\hyperdef{L}{X86A62A937A42B82E}{}
{\noindent\textcolor{FuncColor}{$\triangleright$\enspace\texttt{LieBracket({\mdseries\slshape elm1, elm2})\index{LieBracket@\texttt{LieBracket}}
\label{LieBracket}
}\hfill{\scriptsize (operation)}}\\


 returns the element \texttt{\mbox{\texttt{\mdseries\slshape elm1}} * \mbox{\texttt{\mdseries\slshape elm2}} \texttt{\symbol{45}} \mbox{\texttt{\mdseries\slshape elm2}} * \mbox{\texttt{\mdseries\slshape elm1}}}. 

 The addition \texttt{\texttt{\symbol{92}}+} (\ref{bSlash+}) is assumed to be associative but \emph{not} assumed to be commutative (see{\nobreakspace}\texttt{IsAdditivelyCommutative} (\ref{IsAdditivelyCommutative})). The multiplication \texttt{\texttt{\symbol{92}}*} (\ref{bSlash*}) is \emph{not} assumed to be commutative or associative (see{\nobreakspace}\texttt{IsCommutative} (\ref{IsCommutative}), \texttt{IsAssociative} (\ref{IsAssociative})). }

 

\subsection{\textcolor{Chapter }{Sqrt}}
\logpage{[ 31, 12, 5 ]}\nobreak
\hyperdef{L}{X7E8F1FB87C229BB0}{}
{\noindent\textcolor{FuncColor}{$\triangleright$\enspace\texttt{Sqrt({\mdseries\slshape obj})\index{Sqrt@\texttt{Sqrt}}
\label{Sqrt}
}\hfill{\scriptsize (operation)}}\\


 \texttt{Sqrt} returns a square root of \mbox{\texttt{\mdseries\slshape obj}}, that is, an object $x$ with the property that $x \cdot x = \mbox{\texttt{\mdseries\slshape obj}}$ holds. If such an $x$ is not unique then the choice of $x$ depends on the type of \mbox{\texttt{\mdseries\slshape obj}}. For example, \texttt{ER} (\ref{ER}) is the \texttt{Sqrt} method for rationals (see{\nobreakspace}\texttt{IsRat} (\ref{IsRat})). }

  }

  
\section{\textcolor{Chapter }{Relations Between Domains}}\label{Relations Between Domains}
\logpage{[ 31, 13, 0 ]}
\hyperdef{L}{X80A2D8A7874B268B}{}
{
  Domains are often constructed relative to other domains. The probably most
usual case is to form a \emph{subset} of a domain, for example the intersection (see{\nobreakspace}\texttt{Intersection} (\ref{Intersection:for a list})) of two domains, or a Sylow subgroup of a given group (see{\nobreakspace}\texttt{SylowSubgroup} (\ref{SylowSubgroup})). 

 In such a situation, the new domain can gain knowledge by exploiting that
several attributes are maintained under taking subsets. For example, the
intersection of an arbitrary domain with a finite domain is clearly finite, a
Sylow subgroup of an abelian group is abelian, too, and so on. 

 Since usually the new domain has access to the knowledge of the old domain(s)
only when it is created (see{\nobreakspace}\ref{Constructing Subdomains} for the exception), this is the right moment to take advantage of the subset
relation, using \texttt{UseSubsetRelation} (\ref{UseSubsetRelation}). 

 Analogous relations occur when a \emph{factor structure} is created from a domain and a subset (see \texttt{UseFactorRelation} (\ref{UseFactorRelation})), and when a domain \emph{isomorphic} to a given one is created (see \texttt{UseIsomorphismRelation} (\ref{UseIsomorphismRelation})). 

 The functions \texttt{InstallSubsetMaintenance} (\ref{InstallSubsetMaintenance}), \texttt{InstallIsomorphismMaintenance} (\ref{InstallIsomorphismMaintenance}), and \texttt{InstallFactorMaintenance} (\ref{InstallFactorMaintenance}) are used to tell \textsf{GAP} under what conditions an attribute is maintained under taking subsets, or
forming factor structures or isomorphic domains. This is used only when a new
attribute or property is created, see{\nobreakspace}\texttt{NewAttribute} (\ref{NewAttribute}) and \texttt{NewProperty} (\ref{NewProperty}). For the attributes already available, such as \texttt{IsFinite} (\ref{IsFinite}) and \texttt{IsCommutative} (\ref{IsCommutative}), the maintenances are already notified. 

\subsection{\textcolor{Chapter }{UseSubsetRelation}}
\logpage{[ 31, 13, 1 ]}\nobreak
\hyperdef{L}{X7C03098C838ADE40}{}
{\noindent\textcolor{FuncColor}{$\triangleright$\enspace\texttt{UseSubsetRelation({\mdseries\slshape super, sub})\index{UseSubsetRelation@\texttt{UseSubsetRelation}}
\label{UseSubsetRelation}
}\hfill{\scriptsize (operation)}}\\


 Methods for this operation transfer possibly useful information from the
domain \mbox{\texttt{\mdseries\slshape super}} to its subset \mbox{\texttt{\mdseries\slshape sub}}, and vice versa. 

 \texttt{UseSubsetRelation} is designed to be called automatically whenever substructures of domains are
constructed. So the methods must be \emph{cheap}, and the requirements should be as sharp as possible! 

 To achieve that \emph{all} applicable methods are executed, all methods for this operation except the
default method must end with \texttt{TryNextMethod()}. This default method deals with the information that is available by the
calls of \texttt{InstallSubsetMaintenance} (\ref{InstallSubsetMaintenance}) in the \textsf{GAP} library. 

 
\begin{Verbatim}[commandchars=!@|,fontsize=\small,frame=single,label=Example]
  !gapprompt@gap>| !gapinput@g:= Group( (1,2), (3,4), (5,6) );; h:= Group( (1,2), (3,4) );;|
  !gapprompt@gap>| !gapinput@IsAbelian( g );  HasIsAbelian( h );|
  true
  false
  !gapprompt@gap>| !gapinput@UseSubsetRelation( g, h );;  HasIsAbelian( h );  IsAbelian( h );|
  true
  true
\end{Verbatim}
 }

 

\subsection{\textcolor{Chapter }{UseFactorRelation}}
\logpage{[ 31, 13, 2 ]}\nobreak
\hyperdef{L}{X78039B628262BFA8}{}
{\noindent\textcolor{FuncColor}{$\triangleright$\enspace\texttt{UseFactorRelation({\mdseries\slshape numer, denom, factor})\index{UseFactorRelation@\texttt{UseFactorRelation}}
\label{UseFactorRelation}
}\hfill{\scriptsize (operation)}}\\


 Methods for this operation transfer possibly useful information from the
domain \mbox{\texttt{\mdseries\slshape numer}} or its subset \mbox{\texttt{\mdseries\slshape denom}} to the domain \mbox{\texttt{\mdseries\slshape factor}} that is isomorphic to the factor of \mbox{\texttt{\mdseries\slshape numer}} by \mbox{\texttt{\mdseries\slshape denom}}, and vice versa. \mbox{\texttt{\mdseries\slshape denom}} may be \texttt{fail}, for example if \mbox{\texttt{\mdseries\slshape factor}} is just known to be a factor of \mbox{\texttt{\mdseries\slshape numer}} but \mbox{\texttt{\mdseries\slshape denom}} is not available as a \textsf{GAP} object; in this case those factor relations are used that are installed
without special requirements for \mbox{\texttt{\mdseries\slshape denom}}. 

 \texttt{UseFactorRelation} is designed to be called automatically whenever factor structures of domains
are constructed. So the methods must be \emph{cheap}, and the requirements should be as sharp as possible! 

 To achieve that \emph{all} applicable methods are executed, all methods for this operation except the
default method must end with a call to \texttt{TryNextMethod} (\ref{TryNextMethod}). This default method deals with the information that is available by the
calls of \texttt{InstallFactorMaintenance} (\ref{InstallFactorMaintenance}) in the \textsf{GAP} library. 

 
\begin{Verbatim}[commandchars=!@|,fontsize=\small,frame=single,label=Example]
  !gapprompt@gap>| !gapinput@g:= Group( (1,2,3,4), (1,2) );; h:= Group( (1,2,3), (1,2) );;|
  !gapprompt@gap>| !gapinput@IsSolvableGroup( g );  HasIsSolvableGroup( h );|
  true
  false
  !gapprompt@gap>| !gapinput@UseFactorRelation(g, Subgroup( g, [ (1,2)(3,4), (1,3)(2,4) ] ), h);;|
  !gapprompt@gap>| !gapinput@HasIsSolvableGroup( h );  IsSolvableGroup( h );|
  true
  true
\end{Verbatim}
 }

 

\subsection{\textcolor{Chapter }{UseIsomorphismRelation}}
\logpage{[ 31, 13, 3 ]}\nobreak
\hyperdef{L}{X839BE6467E8474D9}{}
{\noindent\textcolor{FuncColor}{$\triangleright$\enspace\texttt{UseIsomorphismRelation({\mdseries\slshape old, new})\index{UseIsomorphismRelation@\texttt{UseIsomorphismRelation}}
\label{UseIsomorphismRelation}
}\hfill{\scriptsize (operation)}}\\


 Methods for this operation transfer possibly useful information from the
domain \mbox{\texttt{\mdseries\slshape old}} to the isomorphic domain \mbox{\texttt{\mdseries\slshape new}}. 

 \texttt{UseIsomorphismRelation} is designed to be called automatically whenever isomorphic structures of
domains are constructed. So the methods must be \emph{cheap}, and the requirements should be as sharp as possible! 

 To achieve that \emph{all} applicable methods are executed, all methods for this operation except the
default method must end with a call to \texttt{TryNextMethod} (\ref{TryNextMethod}). This default method deals with the information that is available by the
calls of \texttt{InstallIsomorphismMaintenance} (\ref{InstallIsomorphismMaintenance}) in the \textsf{GAP} library. 

 
\begin{Verbatim}[commandchars=!@|,fontsize=\small,frame=single,label=Example]
  !gapprompt@gap>| !gapinput@g:= Group( (1,2) );;  h:= Group( [ [ -1 ] ] );;|
  !gapprompt@gap>| !gapinput@Size( g );  HasSize( h );|
  2
  false
  !gapprompt@gap>| !gapinput@UseIsomorphismRelation( g, h );;  HasSize( h );  Size( h );|
  true
  2
\end{Verbatim}
 }

 

\subsection{\textcolor{Chapter }{InstallSubsetMaintenance}}
\logpage{[ 31, 13, 4 ]}\nobreak
\hyperdef{L}{X863C35007C7AA914}{}
{\noindent\textcolor{FuncColor}{$\triangleright$\enspace\texttt{InstallSubsetMaintenance({\mdseries\slshape opr, super{\textunderscore}req, sub{\textunderscore}req})\index{InstallSubsetMaintenance@\texttt{InstallSubsetMaintenance}}
\label{InstallSubsetMaintenance}
}\hfill{\scriptsize (function)}}\\


 \mbox{\texttt{\mdseries\slshape opr}} must be a property or an attribute. The call of \texttt{InstallSubsetMaintenance} has the effect that for a domain $D$ in the filter \mbox{\texttt{\mdseries\slshape super{\textunderscore}req}}, and a domain $S$ in the filter \mbox{\texttt{\mdseries\slshape sub{\textunderscore}req}}, the call \texttt{UseSubsetRelation}$( D, S )$ (see{\nobreakspace}\texttt{UseSubsetRelation} (\ref{UseSubsetRelation})) sets a known value of \mbox{\texttt{\mdseries\slshape opr}} for $D$ as value of \mbox{\texttt{\mdseries\slshape opr}} also for $S$. A typical example for which \texttt{InstallSubsetMaintenance} is applied is given by \mbox{\texttt{\mdseries\slshape opr}} \texttt{= IsFinite}, \mbox{\texttt{\mdseries\slshape super{\textunderscore}req}} \texttt{= IsCollection and IsFinite}, and \mbox{\texttt{\mdseries\slshape sub{\textunderscore}req}} \texttt{= IsCollection}. 

 If \mbox{\texttt{\mdseries\slshape opr}} is a property and the filter \mbox{\texttt{\mdseries\slshape super{\textunderscore}req}} lies in the filter \mbox{\texttt{\mdseries\slshape opr}} then we can use also the following inverse implication. If $D$ is in the filter whose intersection with \mbox{\texttt{\mdseries\slshape opr}} is \mbox{\texttt{\mdseries\slshape super{\textunderscore}req}} and if $S$ is in the filter \mbox{\texttt{\mdseries\slshape sub{\textunderscore}req}}, $S$ is a subset of $D$, and the value of \mbox{\texttt{\mdseries\slshape opr}} for $S$ is \texttt{false} then the value of \mbox{\texttt{\mdseries\slshape opr}} for $D$ is also \texttt{false}.  }

 

\subsection{\textcolor{Chapter }{InstallFactorMaintenance}}
\logpage{[ 31, 13, 5 ]}\nobreak
\hyperdef{L}{X7BB7EE5078EF6F47}{}
{\noindent\textcolor{FuncColor}{$\triangleright$\enspace\texttt{InstallFactorMaintenance({\mdseries\slshape opr, numer{\textunderscore}req, denom{\textunderscore}req, factor{\textunderscore}req})\index{InstallFactorMaintenance@\texttt{InstallFactorMaintenance}}
\label{InstallFactorMaintenance}
}\hfill{\scriptsize (function)}}\\


 \mbox{\texttt{\mdseries\slshape opr}} must be a property or an attribute. The call of \texttt{InstallFactorMaintenance} has the effect that for collections $N$, $D$, $F$ in the filters \mbox{\texttt{\mdseries\slshape numer{\textunderscore}req}}, \mbox{\texttt{\mdseries\slshape denom{\textunderscore}req}}, and \mbox{\texttt{\mdseries\slshape factor{\textunderscore}req}}, respectively, the call \texttt{UseFactorRelation}$( N, D, F )$ (see{\nobreakspace}\texttt{UseFactorRelation} (\ref{UseFactorRelation})) sets a known value of \mbox{\texttt{\mdseries\slshape opr}} for $N$ as value of \mbox{\texttt{\mdseries\slshape opr}} also for $F$. A typical example for which \texttt{InstallFactorMaintenance} is applied is given by \mbox{\texttt{\mdseries\slshape opr}} \texttt{= IsFinite}, \mbox{\texttt{\mdseries\slshape numer{\textunderscore}req}} \texttt{= IsCollection and IsFinite}, \mbox{\texttt{\mdseries\slshape denom{\textunderscore}req}} \texttt{= IsCollection}, and \mbox{\texttt{\mdseries\slshape factor{\textunderscore}req}} \texttt{= IsCollection}. 

 For the other direction, if \mbox{\texttt{\mdseries\slshape numer{\textunderscore}req}} involves the filter \mbox{\texttt{\mdseries\slshape opr}} then a known \texttt{false} value of \mbox{\texttt{\mdseries\slshape opr}} for $F$ implies a \texttt{false} value for $D$ provided that $D$ lies in the filter obtained from \mbox{\texttt{\mdseries\slshape numer{\textunderscore}req}} by removing \mbox{\texttt{\mdseries\slshape opr}}. 

 Note that an implication of a factor relation holds in particular for the case
of isomorphisms. So one need \emph{not} install an isomorphism maintained method when a factor maintained method is
already installed. For example, \texttt{UseIsomorphismRelation} (\ref{UseIsomorphismRelation}) will transfer a known \texttt{IsFinite} (\ref{IsFinite}) value because of the installed factor maintained method. }

 

\subsection{\textcolor{Chapter }{InstallIsomorphismMaintenance}}
\logpage{[ 31, 13, 6 ]}\nobreak
\hyperdef{L}{X79F97F0F78D89186}{}
{\noindent\textcolor{FuncColor}{$\triangleright$\enspace\texttt{InstallIsomorphismMaintenance({\mdseries\slshape opr, old{\textunderscore}req, new{\textunderscore}req})\index{InstallIsomorphismMaintenance@\texttt{InstallIsomorphismMaintenance}}
\label{InstallIsomorphismMaintenance}
}\hfill{\scriptsize (function)}}\\


 \mbox{\texttt{\mdseries\slshape opr}} must be a property or an attribute. The call of \texttt{InstallIsomorphismMaintenance} has the effect that for a domain $D$ in the filter \mbox{\texttt{\mdseries\slshape old{\textunderscore}req}}, and a domain $E$ in the filter \mbox{\texttt{\mdseries\slshape new{\textunderscore}req}}, the call \texttt{UseIsomorphismRelation}$( D, E )$ (see{\nobreakspace}\texttt{UseIsomorphismRelation} (\ref{UseIsomorphismRelation})) sets a known value of \mbox{\texttt{\mdseries\slshape opr}} for $D$ as value of \mbox{\texttt{\mdseries\slshape opr}} also for $E$. A typical example for which \texttt{InstallIsomorphismMaintenance} is applied is given by \mbox{\texttt{\mdseries\slshape opr}} \texttt{= Size}, \mbox{\texttt{\mdseries\slshape old{\textunderscore}req}} \texttt{= IsCollection}, and \mbox{\texttt{\mdseries\slshape new{\textunderscore}req}} \texttt{= IsCollection}.    }

 }

  
\section{\textcolor{Chapter }{Useful Categories of Elements}}\label{Useful Categories of Elements}
\logpage{[ 31, 14, 0 ]}
\hyperdef{L}{X7B97A0307EA161E5}{}
{
  This section and the following one are rather technical, and may be
interesting only for those \textsf{GAP} users who want to implement new kinds of elements. 

 It deals with certain categories of elements that are useful mainly for the
design of elements, from the viewpoint that one wants to form certain domains
of these elements. For example, a domain closed under multiplication \texttt{*} (a so\texttt{\symbol{45}}called magma, see Chapter{\nobreakspace}\ref{Magmas}) makes sense only if its elements can be multiplied, and the latter is
indicated by the category \texttt{IsMultiplicativeElement} (\ref{IsMultiplicativeElement}) for each element. Again note that the underlying idea is that a domain is
regarded as \emph{generated} by given elements, and that these elements carry information about the desired
domain. For general information on categories and their hierarchies,
see{\nobreakspace}\ref{Categories}. 

 More special categories of this kind are described in the contexts where they
arise, they are \texttt{IsRowVector} (\ref{IsRowVector}), \texttt{IsMatrix} (\ref{IsMatrix}), \texttt{IsOrdinaryMatrix} (\ref{IsOrdinaryMatrix}), and \texttt{IsLieMatrix} (\ref{IsLieMatrix}). 

 

\subsection{\textcolor{Chapter }{IsExtAElement}}
\logpage{[ 31, 14, 1 ]}\nobreak
\hyperdef{L}{X7FBD4F65861C2DF2}{}
{\noindent\textcolor{FuncColor}{$\triangleright$\enspace\texttt{IsExtAElement({\mdseries\slshape obj})\index{IsExtAElement@\texttt{IsExtAElement}}
\label{IsExtAElement}
}\hfill{\scriptsize (Category)}}\\


 An \emph{external additive element} is an object that can be added via \texttt{+} with other elements (not necessarily in the same family, see{\nobreakspace}\ref{Families}). }

 

\subsection{\textcolor{Chapter }{IsNearAdditiveElement}}
\logpage{[ 31, 14, 2 ]}\nobreak
\hyperdef{L}{X7F346AA47AEC39AB}{}
{\noindent\textcolor{FuncColor}{$\triangleright$\enspace\texttt{IsNearAdditiveElement({\mdseries\slshape obj})\index{IsNearAdditiveElement@\texttt{IsNearAdditiveElement}}
\label{IsNearAdditiveElement}
}\hfill{\scriptsize (Category)}}\\


 A \emph{near\texttt{\symbol{45}}additive element} is an object that can be added via \texttt{+} with elements in its family (see{\nobreakspace}\ref{Families}); this addition is not necessarily commutative. }

 

\subsection{\textcolor{Chapter }{IsAdditiveElement}}
\logpage{[ 31, 14, 3 ]}\nobreak
\hyperdef{L}{X78D042B486E1D7F7}{}
{\noindent\textcolor{FuncColor}{$\triangleright$\enspace\texttt{IsAdditiveElement({\mdseries\slshape obj})\index{IsAdditiveElement@\texttt{IsAdditiveElement}}
\label{IsAdditiveElement}
}\hfill{\scriptsize (Category)}}\\


 An \emph{additive element} is an object that can be added via \texttt{+} with elements in its family (see{\nobreakspace}\ref{Families}); this addition is commutative. }

 

\subsection{\textcolor{Chapter }{IsNearAdditiveElementWithZero}}
\logpage{[ 31, 14, 4 ]}\nobreak
\hyperdef{L}{X7CE2353F836F6E0A}{}
{\noindent\textcolor{FuncColor}{$\triangleright$\enspace\texttt{IsNearAdditiveElementWithZero({\mdseries\slshape obj})\index{IsNearAdditiveElementWithZero@\texttt{IsNearAdditiveElementWithZero}}
\label{IsNearAdditiveElementWithZero}
}\hfill{\scriptsize (Category)}}\\


 A \emph{near\texttt{\symbol{45}}additive
element\texttt{\symbol{45}}with\texttt{\symbol{45}}zero} is an object that can be added via \texttt{+} with elements in its family (see{\nobreakspace}\ref{Families}), and that is an admissible argument for the operation \texttt{Zero} (\ref{Zero}); this addition is not necessarily commutative. }

 

\subsection{\textcolor{Chapter }{IsAdditiveElementWithZero}}
\logpage{[ 31, 14, 5 ]}\nobreak
\hyperdef{L}{X87F3552A789C572D}{}
{\noindent\textcolor{FuncColor}{$\triangleright$\enspace\texttt{IsAdditiveElementWithZero({\mdseries\slshape obj})\index{IsAdditiveElementWithZero@\texttt{IsAdditiveElementWithZero}}
\label{IsAdditiveElementWithZero}
}\hfill{\scriptsize (Category)}}\\


 An \emph{additive element\texttt{\symbol{45}}with\texttt{\symbol{45}}zero} is an object that can be added via \texttt{+} with elements in its family (see{\nobreakspace}\ref{Families}), and that is an admissible argument for the operation \texttt{Zero} (\ref{Zero}); this addition is commutative. }

 

\subsection{\textcolor{Chapter }{IsNearAdditiveElementWithInverse}}
\logpage{[ 31, 14, 6 ]}\nobreak
\hyperdef{L}{X84B0929982B51CB4}{}
{\noindent\textcolor{FuncColor}{$\triangleright$\enspace\texttt{IsNearAdditiveElementWithInverse({\mdseries\slshape obj})\index{IsNearAdditiveElementWithInverse@\texttt{IsNearAdditiveElementWithInverse}}
\label{IsNearAdditiveElementWithInverse}
}\hfill{\scriptsize (Category)}}\\


 A \emph{near\texttt{\symbol{45}}additive
element\texttt{\symbol{45}}with\texttt{\symbol{45}}inverse} is an object that can be added via \texttt{+} with elements in its family (see{\nobreakspace}\ref{Families}), and that is an admissible argument for the operations \texttt{Zero} (\ref{Zero}) and \texttt{AdditiveInverse} (\ref{AdditiveInverse}); this addition is not necessarily commutative. }

 

\subsection{\textcolor{Chapter }{IsAdditiveElementWithInverse}}
\logpage{[ 31, 14, 7 ]}\nobreak
\hyperdef{L}{X7C0E4AE883947778}{}
{\noindent\textcolor{FuncColor}{$\triangleright$\enspace\texttt{IsAdditiveElementWithInverse({\mdseries\slshape obj})\index{IsAdditiveElementWithInverse@\texttt{IsAdditiveElementWithInverse}}
\label{IsAdditiveElementWithInverse}
}\hfill{\scriptsize (Category)}}\\


 An \emph{additive element\texttt{\symbol{45}}with\texttt{\symbol{45}}inverse} is an object that can be added via \texttt{+} with elements in its family (see{\nobreakspace}\ref{Families}), and that is an admissible argument for the operations \texttt{Zero} (\ref{Zero}) and \texttt{AdditiveInverse} (\ref{AdditiveInverse}); this addition is commutative. }

 

\subsection{\textcolor{Chapter }{IsExtLElement}}
\logpage{[ 31, 14, 8 ]}\nobreak
\hyperdef{L}{X860D1E387DD5CCCF}{}
{\noindent\textcolor{FuncColor}{$\triangleright$\enspace\texttt{IsExtLElement({\mdseries\slshape obj})\index{IsExtLElement@\texttt{IsExtLElement}}
\label{IsExtLElement}
}\hfill{\scriptsize (Category)}}\\


 An \emph{external left element} is an object that can be multiplied from the left, via \texttt{*}, with other elements (not necessarily in the same family, see{\nobreakspace}\ref{Families}). }

 

\subsection{\textcolor{Chapter }{IsExtRElement}}
\logpage{[ 31, 14, 9 ]}\nobreak
\hyperdef{L}{X809E0C097E480AF1}{}
{\noindent\textcolor{FuncColor}{$\triangleright$\enspace\texttt{IsExtRElement({\mdseries\slshape obj})\index{IsExtRElement@\texttt{IsExtRElement}}
\label{IsExtRElement}
}\hfill{\scriptsize (Category)}}\\


 An \emph{external right element} is an object that can be multiplied from the right, via \texttt{*}, with other elements (not necessarily in the same family, see{\nobreakspace}\ref{Families}). }

 

\subsection{\textcolor{Chapter }{IsMultiplicativeElement}}
\logpage{[ 31, 14, 10 ]}\nobreak
\hyperdef{L}{X797D3B2A7A2B2F53}{}
{\noindent\textcolor{FuncColor}{$\triangleright$\enspace\texttt{IsMultiplicativeElement({\mdseries\slshape obj})\index{IsMultiplicativeElement@\texttt{IsMultiplicativeElement}}
\label{IsMultiplicativeElement}
}\hfill{\scriptsize (Category)}}\\


 A \emph{multiplicative element} is an object that can be multiplied via \texttt{*} with elements in its family (see{\nobreakspace}\ref{Families}). }

 

\subsection{\textcolor{Chapter }{IsMultiplicativeElementWithOne}}
\logpage{[ 31, 14, 11 ]}\nobreak
\hyperdef{L}{X82BC294F7D388AE8}{}
{\noindent\textcolor{FuncColor}{$\triangleright$\enspace\texttt{IsMultiplicativeElementWithOne({\mdseries\slshape obj})\index{IsMultiplicativeElementWithOne@\texttt{IsMultiplicativeElementWithOne}}
\label{IsMultiplicativeElementWithOne}
}\hfill{\scriptsize (Category)}}\\


 A \emph{multiplicative element\texttt{\symbol{45}}with\texttt{\symbol{45}}one} is an object that can be multiplied via \texttt{*} with elements in its family (see{\nobreakspace}\ref{Families}), and that is an admissible argument for the operation \texttt{One} (\ref{One}). }

 

\subsection{\textcolor{Chapter }{IsMultiplicativeElementWithZero}}
\logpage{[ 31, 14, 12 ]}\nobreak
\hyperdef{L}{X8703BFC2841BBD63}{}
{\noindent\textcolor{FuncColor}{$\triangleright$\enspace\texttt{IsMultiplicativeElementWithZero({\mdseries\slshape elt})\index{IsMultiplicativeElementWithZero@\texttt{IsMultiplicativeElementWithZero}}
\label{IsMultiplicativeElementWithZero}
}\hfill{\scriptsize (Category)}}\\
\textbf{\indent Returns:\ }
\texttt{true} or \texttt{false}.



 This is the category of elements in a family which can be the operands of \texttt{*} (multiplication) and the operation \texttt{MultiplicativeZeroOp} (\ref{MultiplicativeZeroOp}). 
\begin{Verbatim}[commandchars=!@|,fontsize=\small,frame=single,label=Example]
  !gapprompt@gap>| !gapinput@S:=Semigroup(Transformation( [ 1, 1, 1 ] ));;|
  !gapprompt@gap>| !gapinput@M:=MagmaWithZeroAdjoined(S);|
  <<commutative transformation semigroup of degree 3 with 1 generator>
    with 0 adjoined>
  !gapprompt@gap>| !gapinput@x:=Representative(M);|
  <semigroup with 0 adjoined elt: Transformation( [ 1, 1, 1 ] )>
  !gapprompt@gap>| !gapinput@IsMultiplicativeElementWithZero(x);|
  true
  !gapprompt@gap>| !gapinput@MultiplicativeZeroOp(x);|
  <semigroup with 0 adjoined elt: 0>
\end{Verbatim}
 }

 

\subsection{\textcolor{Chapter }{IsMultiplicativeElementWithInverse}}
\logpage{[ 31, 14, 13 ]}\nobreak
\hyperdef{L}{X7FDB14E57814FA3B}{}
{\noindent\textcolor{FuncColor}{$\triangleright$\enspace\texttt{IsMultiplicativeElementWithInverse({\mdseries\slshape obj})\index{IsMultiplicativeElementWithInverse@\texttt{IsMultiplicativeElementWithInverse}}
\label{IsMultiplicativeElementWithInverse}
}\hfill{\scriptsize (Category)}}\\


 A \emph{multiplicative element\texttt{\symbol{45}}with\texttt{\symbol{45}}inverse} is an object that can be multiplied via \texttt{*} with elements in its family (see{\nobreakspace}\ref{Families}), and that is an admissible argument for the operations \texttt{One} (\ref{One}) and \texttt{Inverse} (\ref{Inverse}). (Note the word ``admissible'': an object in this category does not necessarily have an inverse, \texttt{Inverse} (\ref{Inverse}) may return \texttt{fail}.) }

 

\subsection{\textcolor{Chapter }{IsVector}}
\logpage{[ 31, 14, 14 ]}\nobreak
\hyperdef{L}{X802F34F280B29DF4}{}
{\noindent\textcolor{FuncColor}{$\triangleright$\enspace\texttt{IsVector({\mdseries\slshape obj})\index{IsVector@\texttt{IsVector}}
\label{IsVector}
}\hfill{\scriptsize (Category)}}\\


 A \emph{vector} is an
additive\texttt{\symbol{45}}element\texttt{\symbol{45}}with\texttt{\symbol{45}}inverse
that can be multiplied from the left and right with other objects (not
necessarily of the same type). Examples are cyclotomics, finite field
elements, and of course row vectors (see below). 

 Note that not all lists of ring elements are regarded as vectors, for example
lists of matrices are not vectors. This is because although the category \texttt{IsAdditiveElementWithInverse} (\ref{IsAdditiveElementWithInverse}) is implied by the meet of its collections category and \texttt{IsList} (\ref{IsList}), the family of a list entry may not imply \texttt{IsAdditiveElementWithInverse} (\ref{IsAdditiveElementWithInverse}) for all its elements. }

 

\subsection{\textcolor{Chapter }{IsNearRingElement}}
\logpage{[ 31, 14, 15 ]}\nobreak
\hyperdef{L}{X799AEDE180C31276}{}
{\noindent\textcolor{FuncColor}{$\triangleright$\enspace\texttt{IsNearRingElement({\mdseries\slshape obj})\index{IsNearRingElement@\texttt{IsNearRingElement}}
\label{IsNearRingElement}
}\hfill{\scriptsize (Category)}}\\


 \texttt{IsNearRingElement} is just a synonym for the meet of \texttt{IsNearAdditiveElementWithInverse} (\ref{IsNearAdditiveElementWithInverse}) and \texttt{IsMultiplicativeElement} (\ref{IsMultiplicativeElement}). }

 

\subsection{\textcolor{Chapter }{IsRingElement}}
\logpage{[ 31, 14, 16 ]}\nobreak
\hyperdef{L}{X84BF40CA86C07361}{}
{\noindent\textcolor{FuncColor}{$\triangleright$\enspace\texttt{IsRingElement({\mdseries\slshape obj})\index{IsRingElement@\texttt{IsRingElement}}
\label{IsRingElement}
}\hfill{\scriptsize (Category)}}\\


 \texttt{IsRingElement} is just a synonym for the meet of \texttt{IsAdditiveElementWithInverse} (\ref{IsAdditiveElementWithInverse}) and \texttt{IsMultiplicativeElement} (\ref{IsMultiplicativeElement}). }

 

\subsection{\textcolor{Chapter }{IsNearRingElementWithOne}}
\logpage{[ 31, 14, 17 ]}\nobreak
\hyperdef{L}{X7C724689784EEF3D}{}
{\noindent\textcolor{FuncColor}{$\triangleright$\enspace\texttt{IsNearRingElementWithOne({\mdseries\slshape obj})\index{IsNearRingElementWithOne@\texttt{IsNearRingElementWithOne}}
\label{IsNearRingElementWithOne}
}\hfill{\scriptsize (Category)}}\\


 \texttt{IsNearRingElementWithOne} is just a synonym for the meet of \texttt{IsNearAdditiveElementWithInverse} (\ref{IsNearAdditiveElementWithInverse}) and \texttt{IsMultiplicativeElementWithOne} (\ref{IsMultiplicativeElementWithOne}). }

 

\subsection{\textcolor{Chapter }{IsRingElementWithOne}}
\logpage{[ 31, 14, 18 ]}\nobreak
\hyperdef{L}{X875B67208017608E}{}
{\noindent\textcolor{FuncColor}{$\triangleright$\enspace\texttt{IsRingElementWithOne({\mdseries\slshape obj})\index{IsRingElementWithOne@\texttt{IsRingElementWithOne}}
\label{IsRingElementWithOne}
}\hfill{\scriptsize (Category)}}\\


 \texttt{IsRingElementWithOne} is just a synonym for the meet of \texttt{IsAdditiveElementWithInverse} (\ref{IsAdditiveElementWithInverse}) and \texttt{IsMultiplicativeElementWithOne} (\ref{IsMultiplicativeElementWithOne}). }

 

\subsection{\textcolor{Chapter }{IsNearRingElementWithInverse}}
\logpage{[ 31, 14, 19 ]}\nobreak
\hyperdef{L}{X80CD04ED85B6B2F9}{}
{\noindent\textcolor{FuncColor}{$\triangleright$\enspace\texttt{IsNearRingElementWithInverse({\mdseries\slshape obj})\index{IsNearRingElementWithInverse@\texttt{IsNearRingElementWithInverse}}
\label{IsNearRingElementWithInverse}
}\hfill{\scriptsize (Category)}}\\


 \texttt{IsNearRingElementWithInverse} is just a synonym for the meet of \texttt{IsNearAdditiveElementWithInverse} (\ref{IsNearAdditiveElementWithInverse}) and \texttt{IsMultiplicativeElementWithInverse} (\ref{IsMultiplicativeElementWithInverse}). }

 

\subsection{\textcolor{Chapter }{IsRingElementWithInverse}}
\logpage{[ 31, 14, 20 ]}\nobreak
\hyperdef{L}{X8113834E84FD0435}{}
{\noindent\textcolor{FuncColor}{$\triangleright$\enspace\texttt{IsRingElementWithInverse({\mdseries\slshape obj})\index{IsRingElementWithInverse@\texttt{IsRingElementWithInverse}}
\label{IsRingElementWithInverse}
}\hfill{\scriptsize (Category)}}\\
\noindent\textcolor{FuncColor}{$\triangleright$\enspace\texttt{IsScalar({\mdseries\slshape obj})\index{IsScalar@\texttt{IsScalar}}
\label{IsScalar}
}\hfill{\scriptsize (Category)}}\\


 \texttt{IsRingElementWithInverse} and \texttt{IsScalar} are just synonyms for the meet of \texttt{IsAdditiveElementWithInverse} (\ref{IsAdditiveElementWithInverse}) and \texttt{IsMultiplicativeElementWithInverse} (\ref{IsMultiplicativeElementWithInverse}). }

 }

  
\section{\textcolor{Chapter }{Useful Categories for all Elements of a Family}}\label{Useful Categories for all Elements of a Family}
\logpage{[ 31, 15, 0 ]}
\hyperdef{L}{X7ABEF00C870789D2}{}
{
  The following categories of elements are to be understood mainly as categories
for all objects in a family, they are usually used as third argument of \texttt{NewFamily} (\ref{NewFamily}). The purpose of each of the following categories is then to guarantee that
each collection of its elements automatically lies in its collections category
(see{\nobreakspace}\texttt{CategoryCollections} (\ref{CategoryCollections})). 

 For example, the multiplication of permutations is associative, and it is
stored in the family of permutations that each permutation lies in \texttt{IsAssociativeElement} (\ref{IsAssociativeElement}). As a consequence, each magma consisting of permutations (more precisely:
each collection that lies in the family \texttt{CollectionsFamily( PermutationsFamily )}, see{\nobreakspace}\texttt{CollectionsFamily} (\ref{CollectionsFamily})) automatically lies in \texttt{CategoryCollections( IsAssociativeElement )}. A magma in this category is always known to be associative, via a logical
implication (see{\nobreakspace}\ref{Logical Implications}). 

 Similarly, if a family knows that all its elements are in the categories \texttt{IsJacobianElement} (\ref{IsJacobianElement}) and \texttt{IsZeroSquaredElement} (\ref{IsZeroSquaredElement}), then each algebra of these elements is automatically known to be a Lie
algebra (see Chapter{\nobreakspace}\ref{Algebras}). 

\subsection{\textcolor{Chapter }{IsAssociativeElement}}
\logpage{[ 31, 15, 1 ]}\nobreak
\hyperdef{L}{X7979AFAA80FF795A}{}
{\noindent\textcolor{FuncColor}{$\triangleright$\enspace\texttt{IsAssociativeElement({\mdseries\slshape obj})\index{IsAssociativeElement@\texttt{IsAssociativeElement}}
\label{IsAssociativeElement}
}\hfill{\scriptsize (Category)}}\\
\noindent\textcolor{FuncColor}{$\triangleright$\enspace\texttt{IsAssociativeElementCollection({\mdseries\slshape obj})\index{IsAssociativeElementCollection@\texttt{IsAssociativeElementCollection}}
\label{IsAssociativeElementCollection}
}\hfill{\scriptsize (Category)}}\\
\noindent\textcolor{FuncColor}{$\triangleright$\enspace\texttt{IsAssociativeElementCollColl({\mdseries\slshape obj})\index{IsAssociativeElementCollColl@\texttt{IsAssociativeElementCollColl}}
\label{IsAssociativeElementCollColl}
}\hfill{\scriptsize (Category)}}\\


 An element \mbox{\texttt{\mdseries\slshape obj}} in the category \texttt{IsAssociativeElement} knows that the multiplication of any elements in the family of \mbox{\texttt{\mdseries\slshape obj}} is associative. For example, all permutations lie in this category, as well as
those ordinary matrices (see{\nobreakspace}\texttt{IsOrdinaryMatrix} (\ref{IsOrdinaryMatrix})) whose entries are also in \texttt{IsAssociativeElement}. }

 

\subsection{\textcolor{Chapter }{IsAdditivelyCommutativeElement}}
\logpage{[ 31, 15, 2 ]}\nobreak
\hyperdef{L}{X78A286418205CE44}{}
{\noindent\textcolor{FuncColor}{$\triangleright$\enspace\texttt{IsAdditivelyCommutativeElement({\mdseries\slshape obj})\index{IsAdditivelyCommutativeElement@\texttt{IsAdditivelyCommutativeElement}}
\label{IsAdditivelyCommutativeElement}
}\hfill{\scriptsize (Category)}}\\
\noindent\textcolor{FuncColor}{$\triangleright$\enspace\texttt{IsAdditivelyCommutativeElementCollection({\mdseries\slshape obj})\index{IsAdditivelyCommutativeElementCollection@\texttt{IsAdditively}\-\texttt{Commutative}\-\texttt{Element}\-\texttt{Collection}}
\label{IsAdditivelyCommutativeElementCollection}
}\hfill{\scriptsize (Category)}}\\
\noindent\textcolor{FuncColor}{$\triangleright$\enspace\texttt{IsAdditivelyCommutativeElementCollColl({\mdseries\slshape obj})\index{IsAdditivelyCommutativeElementCollColl@\texttt{IsAdditively}\-\texttt{Commutative}\-\texttt{Element}\-\texttt{Coll}\-\texttt{Coll}}
\label{IsAdditivelyCommutativeElementCollColl}
}\hfill{\scriptsize (Category)}}\\
\noindent\textcolor{FuncColor}{$\triangleright$\enspace\texttt{IsAdditivelyCommutativeElementFamily({\mdseries\slshape obj})\index{IsAdditivelyCommutativeElementFamily@\texttt{IsAdditively}\-\texttt{Commutative}\-\texttt{Element}\-\texttt{Family}}
\label{IsAdditivelyCommutativeElementFamily}
}\hfill{\scriptsize (Category)}}\\


 An element \mbox{\texttt{\mdseries\slshape obj}} in the category \texttt{IsAdditivelyCommutativeElement} knows that the addition of any elements in the family of \mbox{\texttt{\mdseries\slshape obj}} is commutative. For example, each finite field element and each rational
number lies in this category. }

 

\subsection{\textcolor{Chapter }{IsCommutativeElement}}
\logpage{[ 31, 15, 3 ]}\nobreak
\hyperdef{L}{X8137FA8D86714AC0}{}
{\noindent\textcolor{FuncColor}{$\triangleright$\enspace\texttt{IsCommutativeElement({\mdseries\slshape obj})\index{IsCommutativeElement@\texttt{IsCommutativeElement}}
\label{IsCommutativeElement}
}\hfill{\scriptsize (Category)}}\\
\noindent\textcolor{FuncColor}{$\triangleright$\enspace\texttt{IsCommutativeElementCollection({\mdseries\slshape obj})\index{IsCommutativeElementCollection@\texttt{IsCommutativeElementCollection}}
\label{IsCommutativeElementCollection}
}\hfill{\scriptsize (Category)}}\\
\noindent\textcolor{FuncColor}{$\triangleright$\enspace\texttt{IsCommutativeElementCollColl({\mdseries\slshape obj})\index{IsCommutativeElementCollColl@\texttt{IsCommutativeElementCollColl}}
\label{IsCommutativeElementCollColl}
}\hfill{\scriptsize (Category)}}\\


 An element \mbox{\texttt{\mdseries\slshape obj}} in the category \texttt{IsCommutativeElement} knows that the multiplication of any elements in the family of \mbox{\texttt{\mdseries\slshape obj}} is commutative. For example, each finite field element and each rational
number lies in this category. }

 

\subsection{\textcolor{Chapter }{IsFiniteOrderElement}}
\logpage{[ 31, 15, 4 ]}\nobreak
\hyperdef{L}{X810D2E5E832594AA}{}
{\noindent\textcolor{FuncColor}{$\triangleright$\enspace\texttt{IsFiniteOrderElement({\mdseries\slshape obj})\index{IsFiniteOrderElement@\texttt{IsFiniteOrderElement}}
\label{IsFiniteOrderElement}
}\hfill{\scriptsize (Category)}}\\
\noindent\textcolor{FuncColor}{$\triangleright$\enspace\texttt{IsFiniteOrderElementCollection({\mdseries\slshape obj})\index{IsFiniteOrderElementCollection@\texttt{IsFiniteOrderElementCollection}}
\label{IsFiniteOrderElementCollection}
}\hfill{\scriptsize (Category)}}\\
\noindent\textcolor{FuncColor}{$\triangleright$\enspace\texttt{IsFiniteOrderElementCollColl({\mdseries\slshape obj})\index{IsFiniteOrderElementCollColl@\texttt{IsFiniteOrderElementCollColl}}
\label{IsFiniteOrderElementCollColl}
}\hfill{\scriptsize (Category)}}\\


 An element \mbox{\texttt{\mdseries\slshape obj}} in the category \texttt{IsFiniteOrderElement} knows that it has finite multiplicative order. For example, each finite field
element and each permutation lies in this category. However the value may be \texttt{false} even if \mbox{\texttt{\mdseries\slshape obj}} has finite order, but if this was not known when \mbox{\texttt{\mdseries\slshape obj}} was constructed. 

 Although it is legal to set this filter for any object with finite order, this
is really useful only in the case that all elements of a family are known to
have finite order. }

 

\subsection{\textcolor{Chapter }{IsJacobianElement}}
\logpage{[ 31, 15, 5 ]}\nobreak
\hyperdef{L}{X796957D0805A0221}{}
{\noindent\textcolor{FuncColor}{$\triangleright$\enspace\texttt{IsJacobianElement({\mdseries\slshape obj})\index{IsJacobianElement@\texttt{IsJacobianElement}}
\label{IsJacobianElement}
}\hfill{\scriptsize (Category)}}\\
\noindent\textcolor{FuncColor}{$\triangleright$\enspace\texttt{IsJacobianElementCollection({\mdseries\slshape obj})\index{IsJacobianElementCollection@\texttt{IsJacobianElementCollection}}
\label{IsJacobianElementCollection}
}\hfill{\scriptsize (Category)}}\\
\noindent\textcolor{FuncColor}{$\triangleright$\enspace\texttt{IsJacobianElementCollColl({\mdseries\slshape obj})\index{IsJacobianElementCollColl@\texttt{IsJacobianElementCollColl}}
\label{IsJacobianElementCollColl}
}\hfill{\scriptsize (Category)}}\\
\noindent\textcolor{FuncColor}{$\triangleright$\enspace\texttt{IsRestrictedJacobianElement({\mdseries\slshape obj})\index{IsRestrictedJacobianElement@\texttt{IsRestrictedJacobianElement}}
\label{IsRestrictedJacobianElement}
}\hfill{\scriptsize (Category)}}\\
\noindent\textcolor{FuncColor}{$\triangleright$\enspace\texttt{IsRestrictedJacobianElementCollection({\mdseries\slshape obj})\index{IsRestrictedJacobianElementCollection@\texttt{IsRestricted}\-\texttt{Jacobian}\-\texttt{Element}\-\texttt{Collection}}
\label{IsRestrictedJacobianElementCollection}
}\hfill{\scriptsize (Category)}}\\
\noindent\textcolor{FuncColor}{$\triangleright$\enspace\texttt{IsRestrictedJacobianElementCollColl({\mdseries\slshape obj})\index{IsRestrictedJacobianElementCollColl@\texttt{IsRestrictedJacobianElementCollColl}}
\label{IsRestrictedJacobianElementCollColl}
}\hfill{\scriptsize (Category)}}\\


 An element \mbox{\texttt{\mdseries\slshape obj}} in the category \texttt{IsJacobianElement} knows that the multiplication of any elements in the family $F$ of \mbox{\texttt{\mdseries\slshape obj}} satisfies the Jacobi identity, that is, $x * y * z + z * x * y + y * z * x$ is zero for all $x$, $y$, $z$ in $F$. 

 For example, each Lie matrix (see{\nobreakspace}\texttt{IsLieMatrix} (\ref{IsLieMatrix})) lies in this category. }

 

\subsection{\textcolor{Chapter }{IsZeroSquaredElement}}
\logpage{[ 31, 15, 6 ]}\nobreak
\hyperdef{L}{X7844399D7847AB24}{}
{\noindent\textcolor{FuncColor}{$\triangleright$\enspace\texttt{IsZeroSquaredElement({\mdseries\slshape obj})\index{IsZeroSquaredElement@\texttt{IsZeroSquaredElement}}
\label{IsZeroSquaredElement}
}\hfill{\scriptsize (Category)}}\\
\noindent\textcolor{FuncColor}{$\triangleright$\enspace\texttt{IsZeroSquaredElementCollection({\mdseries\slshape obj})\index{IsZeroSquaredElementCollection@\texttt{IsZeroSquaredElementCollection}}
\label{IsZeroSquaredElementCollection}
}\hfill{\scriptsize (Category)}}\\
\noindent\textcolor{FuncColor}{$\triangleright$\enspace\texttt{IsZeroSquaredElementCollColl({\mdseries\slshape obj})\index{IsZeroSquaredElementCollColl@\texttt{IsZeroSquaredElementCollColl}}
\label{IsZeroSquaredElementCollColl}
}\hfill{\scriptsize (Category)}}\\


 An element \mbox{\texttt{\mdseries\slshape obj}} in the category \texttt{IsZeroSquaredElement} knows that \texttt{\mbox{\texttt{\mdseries\slshape obj}}\texttt{\symbol{94}}2 = Zero( \mbox{\texttt{\mdseries\slshape obj}} )}. For example, each Lie matrix (see{\nobreakspace}\texttt{IsLieMatrix} (\ref{IsLieMatrix})) lies in this category. 

 Although it is legal to set this filter for any zero squared object, this is
really useful only in the case that all elements of a family are known to have
square zero. }

 }

 }

            
\chapter{\textcolor{Chapter }{Mappings}}\label{Mappings}
\logpage{[ 32, 0, 0 ]}
\hyperdef{L}{X7C9734B880042C73}{}
{
  \index{functions!as in mathematics} \index{relations} A \emph{mapping} in \textsf{GAP} is what is called a ``function'' in mathematics. \textsf{GAP} also implements \emph{generalized mappings} in which one element might have several images, these can be imagined as
subsets of the cartesian product and are often called ``relations''. 

 Most operations are declared for general mappings and therefore this manual
often refers to ``(general) mappings'', unless you deliberately need the generalization you can ignore the ``general'' bit and just read it as ``mappings''. 

 A \emph{general mapping} $F$ in \textsf{GAP} is described by its source $S$, its range $R$, and a subset $Rel$ of the direct product $S \times R$, which is called the underlying relation of $F$. $S$, $R$, and $Rel$ are generalized domains (see \ref{Domains}). The corresponding attributes for general mappings are \texttt{Source} (\ref{Source}), \texttt{Range} (\ref{Range:of a general mapping}), and \texttt{UnderlyingRelation} (\ref{UnderlyingRelation}).   

 Note that general mappings themselves are \emph{not} domains. One reason for this is that two general mappings with same underlying
relation are regarded as equal only if also the sources are equal and the
ranges are equal. Other, more technical, reasons are that general mappings and
domains have different basic operations, and that general mappings are
arithmetic objects (see{\nobreakspace}\ref{Arithmetic Operations for General Mappings}); both should not apply to domains. 

 Each element of an underlying relation of a general mapping lies in the
category of direct product elements (see{\nobreakspace}\texttt{IsDirectProductElement} (\ref{IsDirectProductElement})). 

 For each $s \in S$, the set $\{ r \in R | (s,r) \in Rel \}$ is called the set of \emph{images} of $s$. Analogously, for $r \in R$, the set $\{ s \in S | (s,r) \in Rel \}$ is called the set of \emph{preimages} of $r$. 

 The \emph{ordering} of general mappings via \texttt{{\textless}} is defined by the ordering of source, range, and underlying relation.
Specifically, if the source and range domains of \mbox{\texttt{\mdseries\slshape map1}} and \mbox{\texttt{\mdseries\slshape map2}} are the same, then one considers the union of the preimages of \mbox{\texttt{\mdseries\slshape map1}} and \mbox{\texttt{\mdseries\slshape map2}} as a strictly ordered set. The smaller of \mbox{\texttt{\mdseries\slshape map1}} and \mbox{\texttt{\mdseries\slshape map2}} is the one whose image is smaller on the first point of this sequence where
they differ. 

 For mappings which preserve an algebraic structure a \emph{kernel} is defined. Depending on the structure preserved the operation to compute this
kernel is called differently, see Section{\nobreakspace}\ref{Mappings which are Compatible with Algebraic Structures}. 

 Some technical details of general mappings are described in
section{\nobreakspace}\ref{Technical Matters Concerning General Mappings}.   
\section{\textcolor{Chapter }{Direct Products and their Elements}}\label{sect:Direct Products and their Elements}
\logpage{[ 32, 1, 0 ]}
\hyperdef{L}{X859A13548515A5D7}{}
{
  

\subsection{\textcolor{Chapter }{IsDirectProductElement}}
\logpage{[ 32, 1, 1 ]}\nobreak
\hyperdef{L}{X87FD9FE787023FF0}{}
{\noindent\textcolor{FuncColor}{$\triangleright$\enspace\texttt{IsDirectProductElement({\mdseries\slshape obj})\index{IsDirectProductElement@\texttt{IsDirectProductElement}}
\label{IsDirectProductElement}
}\hfill{\scriptsize (Category)}}\\


 \texttt{IsDirectProductElement} is a subcategory of the meet of \texttt{IsDenseList} (\ref{IsDenseList}), \texttt{IsMultiplicativeElementWithInverse} (\ref{IsMultiplicativeElementWithInverse}), \texttt{IsAdditiveElementWithInverse} (\ref{IsAdditiveElementWithInverse}), and \texttt{IsCopyable} (\ref{IsCopyable}), where the arithmetic operations (addition, zero, additive inverse,
multiplication, powering, one, inverse) are defined componentwise. 

 Note that each of these operations will cause an error message if its result
for at least one component cannot be formed. 

 For an object in the filter \texttt{IsDirectProductElement}, \texttt{ShallowCopy} (\ref{ShallowCopy}) returns a mutable plain list with the same entries. The sum and the product of
a direct product element and a list in \texttt{IsListDefault} (\ref{IsListDefault}) is the list of sums and products, respectively. The sum and the product of a
direct product element and an object that is neither a list nor a collection
is the direct product element of componentwise sums and products,
respectively. }

 

\subsection{\textcolor{Chapter }{DirectProductElement}}
\logpage{[ 32, 1, 2 ]}\nobreak
\hyperdef{L}{X8226334A862C5219}{}
{\noindent\textcolor{FuncColor}{$\triangleright$\enspace\texttt{DirectProductElement({\mdseries\slshape [fam, ]objlist})\index{DirectProductElement@\texttt{DirectProductElement}}
\label{DirectProductElement}
}\hfill{\scriptsize (operation)}}\\
\noindent\textcolor{FuncColor}{$\triangleright$\enspace\texttt{DirectProductElementNC({\mdseries\slshape fam, objlist})\index{DirectProductElementNC@\texttt{DirectProductElementNC}}
\label{DirectProductElementNC}
}\hfill{\scriptsize (operation)}}\\


 \texttt{DirectProductElement} returns a direct product element whose components are the entries of the list \mbox{\texttt{\mdseries\slshape objlist}}. If the optional argument \mbox{\texttt{\mdseries\slshape fam}} is given, it must be the family of direct product elements with the
appropriate component families; otherwise the family is determined from \mbox{\texttt{\mdseries\slshape objlist}}. 

 \texttt{DirectProductElementNC} does the same but omits the checks that \mbox{\texttt{\mdseries\slshape fam}} is consistent with \mbox{\texttt{\mdseries\slshape objlist}}. 

 Direct product elements behave like lists: they can be indexed with \texttt{\mbox{\texttt{\mdseries\slshape elm}}[i]} to access the $i$\texttt{\symbol{45}}th component, and \texttt{Length} (\ref{Length}) returns the number of components. Arithmetic operations (multiplication,
powering, one, inverse, addition, zero, additive inverse) are performed
componentwise. 
\begin{Verbatim}[commandchars=!@|,fontsize=\small,frame=single,label=Example]
  !gapprompt@gap>| !gapinput@a := (1,2,3);; b := (1,2);;|
  !gapprompt@gap>| !gapinput@elm := DirectProductElement([a, b]);|
  DirectProductElement( [ (1,2,3), (1,2) ] )
  !gapprompt@gap>| !gapinput@elm[1];|
  (1,2,3)
  !gapprompt@gap>| !gapinput@elm[2];|
  (1,2)
  !gapprompt@gap>| !gapinput@Length(elm);|
  2
  !gapprompt@gap>| !gapinput@elm2 := DirectProductElement([b, a]);;|
  !gapprompt@gap>| !gapinput@elm * elm2;|
  DirectProductElement( [ (2,3), (1,3) ] )
  !gapprompt@gap>| !gapinput@One(elm);|
  DirectProductElement( [ (), () ] )
  !gapprompt@gap>| !gapinput@Inverse(elm);|
  DirectProductElement( [ (1,3,2), (1,2) ] )
\end{Verbatim}
 

 The following example shows \texttt{DirectProductElement} in the context of a direct product $K = G \times H$ built with \texttt{DirectProduct} (\ref{DirectProduct}), where $G$ is a permutation group and $H$ is a pc group. The projections of $K$ are used to move between the direct product and its component groups. 
\begin{Verbatim}[commandchars=!@|,fontsize=\small,frame=single,label=Example]
  !gapprompt@gap>| !gapinput@G := DihedralGroup(IsPermGroup, 8);;|
  !gapprompt@gap>| !gapinput@H := CyclicGroup(IsPcGroup, 4);;|
  !gapprompt@gap>| !gapinput@K := DirectProduct(G, H);;|
  !gapprompt@gap>| !gapinput@proj1 := Projection(K, 1);; proj2 := Projection(K, 2);;|
  !gapprompt@gap>| !gapinput@g := (1,2,3,4);; h := GeneratorsOfGroup(H)[1];;|
  !gapprompt@gap>| !gapinput@elm := DirectProductElement([g, h]);|
  DirectProductElement( [ (1,2,3,4), f1 ] )
  !gapprompt@gap>| !gapinput@Image(proj1, elm) = g;|
  true
  !gapprompt@gap>| !gapinput@Image(proj2, elm) = h;|
  true
\end{Verbatim}
 }

 

\subsection{\textcolor{Chapter }{DirectProductFamily}}
\logpage{[ 32, 1, 3 ]}\nobreak
\hyperdef{L}{X78F8A1168280E06D}{}
{\noindent\textcolor{FuncColor}{$\triangleright$\enspace\texttt{DirectProductFamily({\mdseries\slshape args})\index{DirectProductFamily@\texttt{DirectProductFamily}}
\label{DirectProductFamily}
}\hfill{\scriptsize (function)}}\\


 \mbox{\texttt{\mdseries\slshape args}} must be a dense list of \texttt{CollectionsFamily} (\ref{CollectionsFamily}) families, otherwise the function raises an error. 

 \texttt{DirectProductFamily} returns \texttt{fam}, a collections family of \texttt{IsDirectProductElement} (\ref{IsDirectProductElement}) objects. 

 \texttt{fam} is the \texttt{CollectionsFamily} (\ref{CollectionsFamily}) of \texttt{IsDirectProductElement} (\ref{IsDirectProductElement}) objects whose \texttt{i}\texttt{\symbol{45}}th component is in \texttt{ElementsFamily(args[i])}. 

 Note that a collection in \texttt{fam} may not itself be a direct product; it just is a subcollection of a direct
product. 

 
\begin{Verbatim}[commandchars=!@|,fontsize=\small,frame=single,label=Example]
  !gapprompt@gap>| !gapinput@D8 := DihedralGroup(IsPermGroup, 8);;|
  !gapprompt@gap>| !gapinput@FamilyObj(D8) = CollectionsFamily(PermutationsFamily);|
  true
  !gapprompt@gap>| !gapinput@fam := DirectProductFamily([FamilyObj(D8), FamilyObj(D8)]);;|
  !gapprompt@gap>| !gapinput@ComponentsOfDirectProductElementsFamily(ElementsFamily(fam));|
  [ <Family: "PermutationsFamily">, <Family: "PermutationsFamily"> ]
\end{Verbatim}
 Also note that not all direct products in \textsf{GAP} are created via these families. For example if the arguments to \texttt{DirectProduct} (\ref{DirectProduct}) are permutation groups, then it returns a permutation group as well, whose
elements are not \texttt{IsDirectProductElement} (\ref{IsDirectProductElement}) objects. 

 
\begin{Verbatim}[commandchars=!@|,fontsize=\small,frame=single,label=Example]
  !gapprompt@gap>| !gapinput@fam = FamilyObj(DirectProduct(D8, D8));|
  false
  !gapprompt@gap>| !gapinput@D4 := DihedralGroup(IsPcGroup, 4);;|
  !gapprompt@gap>| !gapinput@fam2 := DirectProductFamily([FamilyObj(D8), FamilyObj(D4)]);;|
  !gapprompt@gap>| !gapinput@fam2 = FamilyObj(DirectProduct(D8, D4));|
  true
\end{Verbatim}
 }

 }

  
\section{\textcolor{Chapter }{Creating Mappings}}\label{Creating Mappings}
\logpage{[ 32, 2, 0 ]}
\hyperdef{L}{X7CF6FEFB8290D5CB}{}
{
  

\subsection{\textcolor{Chapter }{GeneralMappingByElements}}
\logpage{[ 32, 2, 1 ]}\nobreak
\hyperdef{L}{X79D0D2F07A14D039}{}
{\noindent\textcolor{FuncColor}{$\triangleright$\enspace\texttt{GeneralMappingByElements({\mdseries\slshape S, R, elms})\index{GeneralMappingByElements@\texttt{GeneralMappingByElements}}
\label{GeneralMappingByElements}
}\hfill{\scriptsize (function)}}\\


 is the general mapping with source \mbox{\texttt{\mdseries\slshape S}} and range \mbox{\texttt{\mdseries\slshape R}}, and with underlying relation consisting of the collection \mbox{\texttt{\mdseries\slshape elms}} of direct product elements. }

 
\subsection{\textcolor{Chapter }{MappingByFunction}}\logpage{[ 32, 2, 2 ]}
\hyperdef{L}{X7D55E1977ED70E01}{}
{
\noindent\textcolor{FuncColor}{$\triangleright$\enspace\texttt{MappingByFunction({\mdseries\slshape S, R, fun[, invfun]})\index{MappingByFunction@\texttt{MappingByFunction}!by function (and inverse function) between two domains}
\label{MappingByFunction:by function (and inverse function) between two domains}
}\hfill{\scriptsize (function)}}\\
\noindent\textcolor{FuncColor}{$\triangleright$\enspace\texttt{MappingByFunction({\mdseries\slshape S, R, fun, false, prefun})\index{MappingByFunction@\texttt{MappingByFunction}!by function and function that computes one preimage}
\label{MappingByFunction:by function and function that computes one preimage}
}\hfill{\scriptsize (function)}}\\


 \texttt{MappingByFunction} returns a mapping \texttt{map} with source \mbox{\texttt{\mdseries\slshape S}} and range \mbox{\texttt{\mdseries\slshape R}}, such that each element $s$ of \mbox{\texttt{\mdseries\slshape S}} is mapped to the element \mbox{\texttt{\mdseries\slshape fun}}$( s )$, where \mbox{\texttt{\mdseries\slshape fun}} is a \textsf{GAP} function. 

 If the argument \mbox{\texttt{\mdseries\slshape invfun}} is bound then \texttt{map} is a bijection between \mbox{\texttt{\mdseries\slshape S}} and \mbox{\texttt{\mdseries\slshape R}}, and the preimage of each element $r$ of \mbox{\texttt{\mdseries\slshape R}} is given by \mbox{\texttt{\mdseries\slshape invfun}}$( r )$, where \mbox{\texttt{\mdseries\slshape invfun}} is a \textsf{GAP} function. 

 If five arguments are given and the fourth argument is \texttt{false} then the \textsf{GAP} function \mbox{\texttt{\mdseries\slshape prefun}} can be used to compute a single preimage also if \texttt{map} is not bijective.    

 The mapping returned by \texttt{MappingByFunction} lies in the filter \texttt{IsNonSPGeneralMapping} (\ref{IsNonSPGeneralMapping}), see{\nobreakspace}\ref{Technical Matters Concerning General Mappings}. }

 

\subsection{\textcolor{Chapter }{InverseGeneralMapping}}
\logpage{[ 32, 2, 3 ]}\nobreak
\hyperdef{L}{X865FC25A87D36F3D}{}
{\noindent\textcolor{FuncColor}{$\triangleright$\enspace\texttt{InverseGeneralMapping({\mdseries\slshape map})\index{InverseGeneralMapping@\texttt{InverseGeneralMapping}}
\label{InverseGeneralMapping}
}\hfill{\scriptsize (attribute)}}\\


 The \emph{inverse general mapping} of a general mapping \mbox{\texttt{\mdseries\slshape map}} is the general mapping whose underlying relation (see{\nobreakspace}\texttt{UnderlyingRelation} (\ref{UnderlyingRelation})) contains a pair $(r,s)$ if and only if the underlying relation of \mbox{\texttt{\mdseries\slshape map}} contains the pair $(s,r)$. 

 Note that the inverse general mapping of a mapping \mbox{\texttt{\mdseries\slshape map}} is in general only a general mapping. If \mbox{\texttt{\mdseries\slshape map}} knows to be bijective, then its inverse general mapping will know to be a
mapping. However, \texttt{Inverse} (\ref{Inverse}) works for a bijective general mapping \mbox{\texttt{\mdseries\slshape map}} only if the source and range of \mbox{\texttt{\mdseries\slshape map}} coincide. This behaviour agrees with the definition of \mbox{\texttt{\mdseries\slshape Inverse}}. }

 

\subsection{\textcolor{Chapter }{RestrictedInverseGeneralMapping}}
\logpage{[ 32, 2, 4 ]}\nobreak
\hyperdef{L}{X7BD2D5A87CD6B213}{}
{\noindent\textcolor{FuncColor}{$\triangleright$\enspace\texttt{RestrictedInverseGeneralMapping({\mdseries\slshape map})\index{RestrictedInverseGeneralMapping@\texttt{RestrictedInverseGeneralMapping}}
\label{RestrictedInverseGeneralMapping}
}\hfill{\scriptsize (attribute)}}\\


 The \emph{restricted inverse general mapping} of a general mapping \mbox{\texttt{\mdseries\slshape map}} is the general mapping whose underlying relation (see{\nobreakspace}\texttt{UnderlyingRelation} (\ref{UnderlyingRelation})) contains a pair $(r,s)$ if and only if the underlying relation of \mbox{\texttt{\mdseries\slshape map}} contains the pair $(s,r)$, and whose domain is restricted to the image of \mbox{\texttt{\mdseries\slshape map}} and whose range is the domain of \mbox{\texttt{\mdseries\slshape map}}. 

 }

 

\subsection{\textcolor{Chapter }{CompositionMapping}}
\logpage{[ 32, 2, 5 ]}\nobreak
\hyperdef{L}{X7ED1E4E27CCE2DCA}{}
{\noindent\textcolor{FuncColor}{$\triangleright$\enspace\texttt{CompositionMapping({\mdseries\slshape map1, map2, ...})\index{CompositionMapping@\texttt{CompositionMapping}}
\label{CompositionMapping}
}\hfill{\scriptsize (function)}}\\


 \texttt{CompositionMapping} allows one to compose arbitrarily many general mappings, and delegates each
step to \texttt{CompositionMapping2} (\ref{CompositionMapping2}). The result is a map that maps an element first under the last argument, then
under the penultimate argument and so forth. 

 Additionally, the properties \texttt{IsInjective} (\ref{IsInjective}) and \texttt{IsSingleValued} (\ref{IsSingleValued}) are maintained. If the range of the $i+1$\texttt{\symbol{45}}th argument is identical to the range of the $i$\texttt{\symbol{45}}th argument, also \texttt{IsTotal} (\ref{IsTotal}) and \texttt{IsSurjective} (\ref{IsSurjective}) are maintained. (So one should not call \texttt{CompositionMapping2} (\ref{CompositionMapping2}) directly if one wants to maintain these properties.) 

 Depending on the types of \mbox{\texttt{\mdseries\slshape map1}} and \mbox{\texttt{\mdseries\slshape map2}}, the returned mapping might be constructed completely new (for example by
giving domain generators and their images, this is for example the case if
both mappings preserve the same algebraic structures and \textsf{GAP} can decompose elements of the source of \mbox{\texttt{\mdseries\slshape map2}} into generators) or as an (iterated) composition (see{\nobreakspace}\texttt{IsCompositionMappingRep} (\ref{IsCompositionMappingRep})). 

 
\begin{Verbatim}[commandchars=!@|,fontsize=\small,frame=single,label=Example]
  !gapprompt@gap>| !gapinput@f := GroupHomomorphismByImages(CyclicGroup(IsPermGroup, 2),|
  !gapprompt@>| !gapinput@                                  CyclicGroup(IsPermGroup, 1));|
  [ (1,2) ] -> [ () ]
  !gapprompt@gap>| !gapinput@g := GroupHomomorphismByImages(CyclicGroup(IsPermGroup, 6),|
  !gapprompt@>| !gapinput@                                  CyclicGroup(IsPermGroup, 2));|
  [ (1,2,3,4,5,6) ] -> [ (1,2) ]
  !gapprompt@gap>| !gapinput@CompositionMapping(f, g);|
  [ (1,2,3,4,5,6) ] -> [ () ]
  !gapprompt@gap>| !gapinput@CompositionMapping(g, f);|
  [ (1,2) ] -> [ () ]
\end{Verbatim}
 

 }

 

\subsection{\textcolor{Chapter }{CompositionMapping2}}
\logpage{[ 32, 2, 6 ]}\nobreak
\hyperdef{L}{X86486B687B7077AC}{}
{\noindent\textcolor{FuncColor}{$\triangleright$\enspace\texttt{CompositionMapping2({\mdseries\slshape map2, map1})\index{CompositionMapping2@\texttt{CompositionMapping2}}
\label{CompositionMapping2}
}\hfill{\scriptsize (operation)}}\\
\noindent\textcolor{FuncColor}{$\triangleright$\enspace\texttt{CompositionMapping2General({\mdseries\slshape map2, map1})\index{CompositionMapping2General@\texttt{CompositionMapping2General}}
\label{CompositionMapping2General}
}\hfill{\scriptsize (function)}}\\


 \texttt{CompositionMapping2} returns the composition of \mbox{\texttt{\mdseries\slshape map2}} and \mbox{\texttt{\mdseries\slshape map1}}, this is the general mapping that maps an element first under \mbox{\texttt{\mdseries\slshape map1}}, and then maps the images under \mbox{\texttt{\mdseries\slshape map2}}. 

 (Note the reverse ordering of arguments in the composition via the
multiplication \texttt{\texttt{\symbol{92}}*} (\ref{bSlash*}). 

 \texttt{CompositionMapping2General} is the method that forms a composite mapping with two constituent mappings.
(This is used in some algorithms.) }

 

\subsection{\textcolor{Chapter }{IsCompositionMappingRep}}
\logpage{[ 32, 2, 7 ]}\nobreak
\hyperdef{L}{X7A926D167C3155F6}{}
{\noindent\textcolor{FuncColor}{$\triangleright$\enspace\texttt{IsCompositionMappingRep({\mdseries\slshape map})\index{IsCompositionMappingRep@\texttt{IsCompositionMappingRep}}
\label{IsCompositionMappingRep}
}\hfill{\scriptsize (Representation)}}\\


 Mappings in this representation are stored as composition of two mappings,
(pre)images of elements are computed in a two\texttt{\symbol{45}}step process.
The constituent mappings of the composition can be obtained via \texttt{ConstituentsCompositionMapping} (\ref{ConstituentsCompositionMapping}). }

 

\subsection{\textcolor{Chapter }{ConstituentsCompositionMapping}}
\logpage{[ 32, 2, 8 ]}\nobreak
\hyperdef{L}{X87775B438008DCA5}{}
{\noindent\textcolor{FuncColor}{$\triangleright$\enspace\texttt{ConstituentsCompositionMapping({\mdseries\slshape map})\index{ConstituentsCompositionMapping@\texttt{ConstituentsCompositionMapping}}
\label{ConstituentsCompositionMapping}
}\hfill{\scriptsize (function)}}\\


 If \mbox{\texttt{\mdseries\slshape map}} is stored in the representation \texttt{IsCompositionMappingRep} (\ref{IsCompositionMappingRep}) as composition of two mappings \mbox{\texttt{\mdseries\slshape map1}} and \mbox{\texttt{\mdseries\slshape map2}}, this function returns the two constituent mappings in a list \texttt{[ \mbox{\texttt{\mdseries\slshape map1}}, \mbox{\texttt{\mdseries\slshape map2}} ]}. }

 

\subsection{\textcolor{Chapter }{ZeroMapping}}
\logpage{[ 32, 2, 9 ]}\nobreak
\hyperdef{L}{X795FF8DC785F110A}{}
{\noindent\textcolor{FuncColor}{$\triangleright$\enspace\texttt{ZeroMapping({\mdseries\slshape S, R})\index{ZeroMapping@\texttt{ZeroMapping}}
\label{ZeroMapping}
}\hfill{\scriptsize (operation)}}\\


 A zero mapping is a total general mapping that maps each element of its source
to the zero element of its range. 

 (Each mapping with empty source is a zero mapping.) }

 

\subsection{\textcolor{Chapter }{IdentityMapping}}
\logpage{[ 32, 2, 10 ]}\nobreak
\hyperdef{L}{X7EBAE0368470A603}{}
{\noindent\textcolor{FuncColor}{$\triangleright$\enspace\texttt{IdentityMapping({\mdseries\slshape D})\index{IdentityMapping@\texttt{IdentityMapping}}
\label{IdentityMapping}
}\hfill{\scriptsize (attribute)}}\\


 is the bijective mapping with source and range equal to the collection \mbox{\texttt{\mdseries\slshape D}}, which maps each element of \mbox{\texttt{\mdseries\slshape D}} to itself. }

 
\subsection{\textcolor{Chapter }{Embedding}}\logpage{[ 32, 2, 11 ]}
\hyperdef{L}{X86452F8587CBAEA0}{}
{
\noindent\textcolor{FuncColor}{$\triangleright$\enspace\texttt{Embedding({\mdseries\slshape S, T})\index{Embedding@\texttt{Embedding}!for two domains}
\label{Embedding:for two domains}
}\hfill{\scriptsize (operation)}}\\
\noindent\textcolor{FuncColor}{$\triangleright$\enspace\texttt{Embedding({\mdseries\slshape S, i})\index{Embedding@\texttt{Embedding}!for a domain and a positive integer}
\label{Embedding:for a domain and a positive integer}
}\hfill{\scriptsize (operation)}}\\


 returns the embedding of the domain \mbox{\texttt{\mdseries\slshape S}} in the domain \mbox{\texttt{\mdseries\slshape T}}, or in the second form, some domain indexed by the positive integer \mbox{\texttt{\mdseries\slshape i}}. The precise natures of the various methods are described elsewhere: for Lie
algebras, see \texttt{LieFamily} (\ref{LieFamily}); for group products, see{\nobreakspace}\ref{Embeddings and Projections for Group Products} for a general description, or for examples see{\nobreakspace}\ref{Direct Products} for direct products, \ref{Semidirect Products} for semidirect products, or{\nobreakspace}\ref{Wreath Products} for wreath products; or for magma rings see{\nobreakspace}\ref{Natural Embeddings related to Magma Rings}. }

 
\subsection{\textcolor{Chapter }{Projection}}\logpage{[ 32, 2, 12 ]}
\hyperdef{L}{X8769E8DA80BC96C1}{}
{
\noindent\textcolor{FuncColor}{$\triangleright$\enspace\texttt{Projection({\mdseries\slshape S, T})\index{Projection@\texttt{Projection}!for two domains}
\label{Projection:for two domains}
}\hfill{\scriptsize (operation)}}\\
\noindent\textcolor{FuncColor}{$\triangleright$\enspace\texttt{Projection({\mdseries\slshape S, i})\index{Projection@\texttt{Projection}!for a domain and a positive integer}
\label{Projection:for a domain and a positive integer}
}\hfill{\scriptsize (operation)}}\\
\noindent\textcolor{FuncColor}{$\triangleright$\enspace\texttt{Projection({\mdseries\slshape S})\index{Projection@\texttt{Projection}!for a domain}
\label{Projection:for a domain}
}\hfill{\scriptsize (operation)}}\\


 returns the projection of the domain \mbox{\texttt{\mdseries\slshape S}} onto the domain \mbox{\texttt{\mdseries\slshape T}}, or in the second form, some domain indexed by the positive integer \mbox{\texttt{\mdseries\slshape i}}, or in the third form some natural quotient domain of \mbox{\texttt{\mdseries\slshape S}}. Various methods are defined for group products; see{\nobreakspace}\ref{Embeddings and Projections for Group Products} for a general description, or for examples see{\nobreakspace}\ref{Direct Products} for direct products, \ref{Semidirect Products} for semidirect products, \ref{Subdirect Products} for subdirect products, or{\nobreakspace}\ref{Wreath Products} for wreath products. }

 

\subsection{\textcolor{Chapter }{RestrictedMapping}}
\logpage{[ 32, 2, 13 ]}\nobreak
\hyperdef{L}{X800014D683A81009}{}
{\noindent\textcolor{FuncColor}{$\triangleright$\enspace\texttt{RestrictedMapping({\mdseries\slshape map, subdom})\index{RestrictedMapping@\texttt{RestrictedMapping}}
\label{RestrictedMapping}
}\hfill{\scriptsize (operation)}}\\


 If \mbox{\texttt{\mdseries\slshape subdom}} is a subdomain of the source of the general mapping \mbox{\texttt{\mdseries\slshape map}}, this operation returns the restriction of \mbox{\texttt{\mdseries\slshape map}} to \mbox{\texttt{\mdseries\slshape subdom}}.  }

 }

  
\section{\textcolor{Chapter }{Properties and Attributes of (General) Mappings}}\label{Properties and Attributes of (General) Mappings}
\logpage{[ 32, 3, 0 ]}
\hyperdef{L}{X7E5A430D7F838F1C}{}
{
  

\subsection{\textcolor{Chapter }{IsTotal}}
\logpage{[ 32, 3, 1 ]}\nobreak
\hyperdef{L}{X83C7494E828CC9C8}{}
{\noindent\textcolor{FuncColor}{$\triangleright$\enspace\texttt{IsTotal({\mdseries\slshape map})\index{IsTotal@\texttt{IsTotal}}
\label{IsTotal}
}\hfill{\scriptsize (property)}}\\


 is \texttt{true} if each element in the source $S$ of the general mapping \mbox{\texttt{\mdseries\slshape map}} has images, i.e., $s^{\mbox{\texttt{\mdseries\slshape map}}} \neq \emptyset$ for all $s \in S$, and \texttt{false} otherwise. }

 

\subsection{\textcolor{Chapter }{IsSingleValued}}
\logpage{[ 32, 3, 2 ]}\nobreak
\hyperdef{L}{X86D44C8A78BF1981}{}
{\noindent\textcolor{FuncColor}{$\triangleright$\enspace\texttt{IsSingleValued({\mdseries\slshape map})\index{IsSingleValued@\texttt{IsSingleValued}}
\label{IsSingleValued}
}\hfill{\scriptsize (property)}}\\


 is \texttt{true} if each element in the source $S$ of the general mapping \mbox{\texttt{\mdseries\slshape map}} has at most one image, i.e., $|s^{\mbox{\texttt{\mdseries\slshape map}}}| \leq 1$ for all $s \in S$, and \texttt{false} otherwise. 

 Equivalently, \texttt{IsSingleValued( \mbox{\texttt{\mdseries\slshape map}} )} is \texttt{true} if and only if the preimages of different elements in $R$ are disjoint. }

 

\subsection{\textcolor{Chapter }{IsMapping}}
\logpage{[ 32, 3, 3 ]}\nobreak
\hyperdef{L}{X7CC95EB282854385}{}
{\noindent\textcolor{FuncColor}{$\triangleright$\enspace\texttt{IsMapping({\mdseries\slshape map})\index{IsMapping@\texttt{IsMapping}}
\label{IsMapping}
}\hfill{\scriptsize (filter)}}\\


 A \emph{mapping} \mbox{\texttt{\mdseries\slshape map}} is a general mapping that assigns to each element \texttt{elm} of its source a unique element \texttt{Image( \mbox{\texttt{\mdseries\slshape map}}, elm )} of its range. 

 Equivalently, the general mapping \mbox{\texttt{\mdseries\slshape map}} is a mapping if and only if it is total and single\texttt{\symbol{45}}valued
(see{\nobreakspace}\texttt{IsTotal} (\ref{IsTotal}), \texttt{IsSingleValued} (\ref{IsSingleValued})). }

 

\subsection{\textcolor{Chapter }{IsInjective}}
\logpage{[ 32, 3, 4 ]}\nobreak
\hyperdef{L}{X7F065FD7822C0A12}{}
{\noindent\textcolor{FuncColor}{$\triangleright$\enspace\texttt{IsInjective({\mdseries\slshape map})\index{IsInjective@\texttt{IsInjective}}
\label{IsInjective}
}\hfill{\scriptsize (property)}}\\


 is \texttt{true} if the images of different elements in the source $S$ of the general mapping \mbox{\texttt{\mdseries\slshape map}} are disjoint, i.e., $x^{\mbox{\texttt{\mdseries\slshape map}}} \cap y^{\mbox{\texttt{\mdseries\slshape map}}} = \emptyset$ for $x \neq y \in S$, and \texttt{false} otherwise. 

 Equivalently, \texttt{IsInjective( \mbox{\texttt{\mdseries\slshape map}} )} is \texttt{true} if and only if each element in the range of \mbox{\texttt{\mdseries\slshape map}} has at most one preimage in $S$. }

 

\subsection{\textcolor{Chapter }{IsSurjective}}
\logpage{[ 32, 3, 5 ]}\nobreak
\hyperdef{L}{X784ECE847E005B8F}{}
{\noindent\textcolor{FuncColor}{$\triangleright$\enspace\texttt{IsSurjective({\mdseries\slshape map})\index{IsSurjective@\texttt{IsSurjective}}
\label{IsSurjective}
}\hfill{\scriptsize (property)}}\\


 is \texttt{true} if each element in the range $R$ of the general mapping \mbox{\texttt{\mdseries\slshape map}} has preimages in the source $S$ of \mbox{\texttt{\mdseries\slshape map}}, i.e., $\{ s \in S \mid x \in s^{\mbox{\texttt{\mdseries\slshape map}}} \} \neq \emptyset$ for all $x \in R$, and \texttt{false} otherwise. }

 

\subsection{\textcolor{Chapter }{IsBijective}}
\logpage{[ 32, 3, 6 ]}\nobreak
\hyperdef{L}{X878F56AB7B342767}{}
{\noindent\textcolor{FuncColor}{$\triangleright$\enspace\texttt{IsBijective({\mdseries\slshape map})\index{IsBijective@\texttt{IsBijective}}
\label{IsBijective}
}\hfill{\scriptsize (property)}}\\


 A general mapping \mbox{\texttt{\mdseries\slshape map}} is \emph{bijective} if and only if it is an injective and surjective mapping (see{\nobreakspace}\texttt{IsMapping} (\ref{IsMapping}), \texttt{IsInjective} (\ref{IsInjective}), \texttt{IsSurjective} (\ref{IsSurjective})). }

 

\subsection{\textcolor{Chapter }{Range (of a general mapping)}}
\logpage{[ 32, 3, 7 ]}\nobreak
\hyperdef{L}{X7B6FD7277CDE9FCB}{}
{\noindent\textcolor{FuncColor}{$\triangleright$\enspace\texttt{Range({\mdseries\slshape map})\index{Range@\texttt{Range}!of a general mapping}
\label{Range:of a general mapping}
}\hfill{\scriptsize (attribute)}}\\


 The range of a general mapping. }

 

\subsection{\textcolor{Chapter }{Source}}
\logpage{[ 32, 3, 8 ]}\nobreak
\hyperdef{L}{X7DE8173F80E07AB1}{}
{\noindent\textcolor{FuncColor}{$\triangleright$\enspace\texttt{Source({\mdseries\slshape map})\index{Source@\texttt{Source}}
\label{Source}
}\hfill{\scriptsize (attribute)}}\\


 The source of a general mapping. }

 

\subsection{\textcolor{Chapter }{UnderlyingRelation}}
\logpage{[ 32, 3, 9 ]}\nobreak
\hyperdef{L}{X784F871383FB599B}{}
{\noindent\textcolor{FuncColor}{$\triangleright$\enspace\texttt{UnderlyingRelation({\mdseries\slshape map})\index{UnderlyingRelation@\texttt{UnderlyingRelation}}
\label{UnderlyingRelation}
}\hfill{\scriptsize (attribute)}}\\


 The \emph{underlying relation} of a general mapping \mbox{\texttt{\mdseries\slshape map}} is the domain of pairs $(s,r)$, with $s$ in the source and $r$ in the range of \mbox{\texttt{\mdseries\slshape map}} (see{\nobreakspace}\texttt{Source} (\ref{Source}), \texttt{Range} (\ref{Range:of a general mapping})), and $r \in$ \texttt{ImagesElm( \mbox{\texttt{\mdseries\slshape map}}, }$s$\texttt{ )}. 

 Each element of the underlying relation is represented by a direct product
element (see{\nobreakspace}\texttt{IsDirectProductElement} (\ref{IsDirectProductElement})). }

 

\subsection{\textcolor{Chapter }{UnderlyingGeneralMapping}}
\logpage{[ 32, 3, 10 ]}\nobreak
\hyperdef{L}{X786581DE871A47D0}{}
{\noindent\textcolor{FuncColor}{$\triangleright$\enspace\texttt{UnderlyingGeneralMapping({\mdseries\slshape map})\index{UnderlyingGeneralMapping@\texttt{UnderlyingGeneralMapping}}
\label{UnderlyingGeneralMapping}
}\hfill{\scriptsize (attribute)}}\\


 attribute for underlying relations of general mappings }

 }

  
\section{\textcolor{Chapter }{Images under Mappings}}\label{Images under Mappings}
\logpage{[ 32, 4, 0 ]}
\hyperdef{L}{X83B4FF15847F06FC}{}
{
  

\subsection{\textcolor{Chapter }{ImagesSource}}
\logpage{[ 32, 4, 1 ]}\nobreak
\hyperdef{L}{X7D23C1CE863DACD8}{}
{\noindent\textcolor{FuncColor}{$\triangleright$\enspace\texttt{ImagesSource({\mdseries\slshape map})\index{ImagesSource@\texttt{ImagesSource}}
\label{ImagesSource}
}\hfill{\scriptsize (attribute)}}\\


 is the set of images of the source of the general mapping \mbox{\texttt{\mdseries\slshape map}}. 

 \texttt{ImagesSource} delegates to \texttt{ImagesSet} (\ref{ImagesSet}), it is introduced only to store the image of \mbox{\texttt{\mdseries\slshape map}} as attribute value. }

 

\subsection{\textcolor{Chapter }{ImagesRepresentative}}
\logpage{[ 32, 4, 2 ]}\nobreak
\hyperdef{L}{X85ADB89B7C8DD7D0}{}
{\noindent\textcolor{FuncColor}{$\triangleright$\enspace\texttt{ImagesRepresentative({\mdseries\slshape map, elm})\index{ImagesRepresentative@\texttt{ImagesRepresentative}}
\label{ImagesRepresentative}
}\hfill{\scriptsize (operation)}}\\


 If \mbox{\texttt{\mdseries\slshape elm}} is an element of the source of the general mapping \mbox{\texttt{\mdseries\slshape map}} then \texttt{ImagesRepresentative} returns either a representative of the set of images of \mbox{\texttt{\mdseries\slshape elm}} under \mbox{\texttt{\mdseries\slshape map}} or \texttt{fail}, the latter if and only if \mbox{\texttt{\mdseries\slshape elm}} has no images under \mbox{\texttt{\mdseries\slshape map}}. 

 Anything may happen if \mbox{\texttt{\mdseries\slshape elm}} is not an element of the source of \mbox{\texttt{\mdseries\slshape map}}. }

 

\subsection{\textcolor{Chapter }{ImagesElm}}
\logpage{[ 32, 4, 3 ]}\nobreak
\hyperdef{L}{X7D51184B7EE5B2CF}{}
{\noindent\textcolor{FuncColor}{$\triangleright$\enspace\texttt{ImagesElm({\mdseries\slshape map, elm})\index{ImagesElm@\texttt{ImagesElm}}
\label{ImagesElm}
}\hfill{\scriptsize (operation)}}\\


 If \mbox{\texttt{\mdseries\slshape elm}} is an element of the source of the general mapping \mbox{\texttt{\mdseries\slshape map}} then \texttt{ImagesElm} returns the set of all images of \mbox{\texttt{\mdseries\slshape elm}} under \mbox{\texttt{\mdseries\slshape map}}. 

 Anything may happen if \mbox{\texttt{\mdseries\slshape elm}} is not an element of the source of \mbox{\texttt{\mdseries\slshape map}}. }

 

\subsection{\textcolor{Chapter }{ImagesSet}}
\logpage{[ 32, 4, 4 ]}\nobreak
\hyperdef{L}{X8781348F7F5796A0}{}
{\noindent\textcolor{FuncColor}{$\triangleright$\enspace\texttt{ImagesSet({\mdseries\slshape map, elms})\index{ImagesSet@\texttt{ImagesSet}}
\label{ImagesSet}
}\hfill{\scriptsize (operation)}}\\


 If \mbox{\texttt{\mdseries\slshape elms}} is a subset of the source of the general mapping \mbox{\texttt{\mdseries\slshape map}} then \texttt{ImagesSet} returns the set of all images of \mbox{\texttt{\mdseries\slshape elms}} under \mbox{\texttt{\mdseries\slshape map}}. 

 The result will be either a proper set or a domain. Anything may happen if \mbox{\texttt{\mdseries\slshape elms}} is not a subset of the source of \mbox{\texttt{\mdseries\slshape map}}. }

 

\subsection{\textcolor{Chapter }{ImageElm}}
\logpage{[ 32, 4, 5 ]}\nobreak
\hyperdef{L}{X7CFAB0157BFB1806}{}
{\noindent\textcolor{FuncColor}{$\triangleright$\enspace\texttt{ImageElm({\mdseries\slshape map, elm})\index{ImageElm@\texttt{ImageElm}}
\label{ImageElm}
}\hfill{\scriptsize (operation)}}\\


 If \mbox{\texttt{\mdseries\slshape elm}} is an element of the source of the total and single\texttt{\symbol{45}}valued
mapping \mbox{\texttt{\mdseries\slshape map}} then \texttt{ImageElm} returns the unique image of \mbox{\texttt{\mdseries\slshape elm}} under \mbox{\texttt{\mdseries\slshape map}}. 

 Anything may happen if \mbox{\texttt{\mdseries\slshape elm}} is not an element of the source of \mbox{\texttt{\mdseries\slshape map}}. }

 
\subsection{\textcolor{Chapter }{Image}}\logpage{[ 32, 4, 6 ]}
\hyperdef{L}{X87F4D35A826599C6}{}
{
\noindent\textcolor{FuncColor}{$\triangleright$\enspace\texttt{Image({\mdseries\slshape map})\index{Image@\texttt{Image}!set of images of the source of a general mapping}
\label{Image:set of images of the source of a general mapping}
}\hfill{\scriptsize (function)}}\\
\noindent\textcolor{FuncColor}{$\triangleright$\enspace\texttt{Image({\mdseries\slshape map, elm})\index{Image@\texttt{Image}!unique image of an element under a mapping}
\label{Image:unique image of an element under a mapping}
}\hfill{\scriptsize (function)}}\\
\noindent\textcolor{FuncColor}{$\triangleright$\enspace\texttt{Image({\mdseries\slshape map, coll})\index{Image@\texttt{Image}!set of images of a collection under a mapping}
\label{Image:set of images of a collection under a mapping}
}\hfill{\scriptsize (function)}}\\


 \texttt{Image( \mbox{\texttt{\mdseries\slshape map}} )} is the \emph{image} of the general mapping \mbox{\texttt{\mdseries\slshape map}}, i.e., the subset of elements of the range of \mbox{\texttt{\mdseries\slshape map}} that are actually values of \mbox{\texttt{\mdseries\slshape map}}. \emph{Note} that in this case the argument may also be multi\texttt{\symbol{45}}valued. 

 \texttt{Image( \mbox{\texttt{\mdseries\slshape map}}, \mbox{\texttt{\mdseries\slshape elm}} )} is the image of the element \mbox{\texttt{\mdseries\slshape elm}} of the source of the mapping \mbox{\texttt{\mdseries\slshape map}} under \mbox{\texttt{\mdseries\slshape map}}, i.e., the unique element of the range to which \mbox{\texttt{\mdseries\slshape map}} maps \mbox{\texttt{\mdseries\slshape elm}}. This can also be expressed as \mbox{\texttt{\mdseries\slshape elm}}\texttt{\texttt{\symbol{94}}}\mbox{\texttt{\mdseries\slshape map}} or as \mbox{\texttt{\mdseries\slshape map}}\texttt{( }\mbox{\texttt{\mdseries\slshape elm}}\texttt{ )}. 

 Note that \mbox{\texttt{\mdseries\slshape map}} must be total and single valued, a multi valued general mapping is not allowed
(see{\nobreakspace}\texttt{Images} (\ref{Images:set of images of the source of a general mapping})). 

 \texttt{Image( \mbox{\texttt{\mdseries\slshape map}}, \mbox{\texttt{\mdseries\slshape coll}} )} is the image of the subset \mbox{\texttt{\mdseries\slshape coll}} of the source of the mapping \mbox{\texttt{\mdseries\slshape map}} under \mbox{\texttt{\mdseries\slshape map}}, i.e., the subset of the range to which \mbox{\texttt{\mdseries\slshape map}} maps elements of \mbox{\texttt{\mdseries\slshape coll}}. 

 \mbox{\texttt{\mdseries\slshape coll}} may be a proper set or a domain. The result will be either a proper set or a
domain. Note that in this case \mbox{\texttt{\mdseries\slshape map}} may also be multi\texttt{\symbol{45}}valued. (If \mbox{\texttt{\mdseries\slshape coll}} and the result are lists then the positions of entries do in general \emph{not} correspond.) 

 \texttt{Image( \mbox{\texttt{\mdseries\slshape map}}, \mbox{\texttt{\mdseries\slshape coll}} )} can also be expressed as \texttt{\mbox{\texttt{\mdseries\slshape map}}(\mbox{\texttt{\mdseries\slshape coll}})} and \texttt{Image( \mbox{\texttt{\mdseries\slshape map}}, \mbox{\texttt{\mdseries\slshape elm}} )} as \texttt{\mbox{\texttt{\mdseries\slshape map}}(\mbox{\texttt{\mdseries\slshape elm}})}. Those using this notation should remember that composition of mappings in \textsf{GAP} still follows the conventions appropriate for mapping acting from the right,
so that \texttt{(\mbox{\texttt{\mdseries\slshape map1}}*\mbox{\texttt{\mdseries\slshape map2}})(\mbox{\texttt{\mdseries\slshape x}})} is equivalent to \texttt{\mbox{\texttt{\mdseries\slshape map2}}(\mbox{\texttt{\mdseries\slshape map1}}(\mbox{\texttt{\mdseries\slshape x}}))} 

 \texttt{Image} delegates to \texttt{ImagesSource} (\ref{ImagesSource}) when called with one argument, and to \texttt{ImageElm} (\ref{ImageElm}) resp. \texttt{ImagesSet} (\ref{ImagesSet}) when called with two arguments. 

 If the second argument is not an element or a subset of the source of the
first argument, an error is signalled. }

 
\subsection{\textcolor{Chapter }{Images}}\logpage{[ 32, 4, 7 ]}
\hyperdef{L}{X86114B2E7E77488C}{}
{
\noindent\textcolor{FuncColor}{$\triangleright$\enspace\texttt{Images({\mdseries\slshape map})\index{Images@\texttt{Images}!set of images of the source of a general mapping}
\label{Images:set of images of the source of a general mapping}
}\hfill{\scriptsize (function)}}\\
\noindent\textcolor{FuncColor}{$\triangleright$\enspace\texttt{Images({\mdseries\slshape map, elm})\index{Images@\texttt{Images}!set of images of an element under a mapping}
\label{Images:set of images of an element under a mapping}
}\hfill{\scriptsize (function)}}\\
\noindent\textcolor{FuncColor}{$\triangleright$\enspace\texttt{Images({\mdseries\slshape map, coll})\index{Images@\texttt{Images}!set of images of a collection under a mapping}
\label{Images:set of images of a collection under a mapping}
}\hfill{\scriptsize (function)}}\\


 \texttt{Images( \mbox{\texttt{\mdseries\slshape map}} )} is the \emph{image} of the general mapping \mbox{\texttt{\mdseries\slshape map}}, i.e., the subset of elements of the range of \mbox{\texttt{\mdseries\slshape map}} that are actually values of \mbox{\texttt{\mdseries\slshape map}}. 

 \texttt{Images( \mbox{\texttt{\mdseries\slshape map}}, \mbox{\texttt{\mdseries\slshape elm}} )} is the set of images of the element \mbox{\texttt{\mdseries\slshape elm}} of the source of the general mapping \mbox{\texttt{\mdseries\slshape map}} under \mbox{\texttt{\mdseries\slshape map}}, i.e., the set of elements of the range to which \mbox{\texttt{\mdseries\slshape map}} maps \mbox{\texttt{\mdseries\slshape elm}}. 

 \texttt{Images( \mbox{\texttt{\mdseries\slshape map}}, \mbox{\texttt{\mdseries\slshape coll}} )} is the set of images of the subset \mbox{\texttt{\mdseries\slshape coll}} of the source of the general mapping \mbox{\texttt{\mdseries\slshape map}} under \mbox{\texttt{\mdseries\slshape map}}, i.e., the subset of the range to which \mbox{\texttt{\mdseries\slshape map}} maps elements of \mbox{\texttt{\mdseries\slshape coll}}. \mbox{\texttt{\mdseries\slshape coll}} may be a proper set or a domain. The result will be either a proper set or a
domain. (If \mbox{\texttt{\mdseries\slshape coll}} and the result are lists then the positions of entries do in general \emph{not} correspond.) 

 \texttt{Images} delegates to \texttt{ImagesSource} (\ref{ImagesSource}) when called with one argument, and to \texttt{ImagesElm} (\ref{ImagesElm}) resp. \texttt{ImagesSet} (\ref{ImagesSet}) when called with two arguments. 

 If the second argument is not an element or a subset of the source of the
first argument, an error is signalled. }

 }

  
\section{\textcolor{Chapter }{Preimages under Mappings}}\label{Preimages under Mappings}
\logpage{[ 32, 5, 0 ]}
\hyperdef{L}{X79BB1EC07C828667}{}
{
  

\subsection{\textcolor{Chapter }{PreImagesRange}}
\logpage{[ 32, 5, 1 ]}\nobreak
\hyperdef{L}{X78EF1FE77B0973C0}{}
{\noindent\textcolor{FuncColor}{$\triangleright$\enspace\texttt{PreImagesRange({\mdseries\slshape map})\index{PreImagesRange@\texttt{PreImagesRange}}
\label{PreImagesRange}
}\hfill{\scriptsize (attribute)}}\\


 is the set of preimages of the range of the general mapping \mbox{\texttt{\mdseries\slshape map}}. 

 \texttt{PreImagesRange} delegates to \texttt{PreImagesSet} (\ref{PreImagesSet}), it is introduced only to store the preimage of \mbox{\texttt{\mdseries\slshape map}} as attribute value. }

 

\subsection{\textcolor{Chapter }{PreImagesElm}}
\logpage{[ 32, 5, 2 ]}\nobreak
\hyperdef{L}{X7FBB830C8729E995}{}
{\noindent\textcolor{FuncColor}{$\triangleright$\enspace\texttt{PreImagesElm({\mdseries\slshape map, elm})\index{PreImagesElm@\texttt{PreImagesElm}}
\label{PreImagesElm}
}\hfill{\scriptsize (operation)}}\\


 If \mbox{\texttt{\mdseries\slshape elm}} is an element of the range of the general mapping \mbox{\texttt{\mdseries\slshape map}} then \texttt{PreImagesElm} returns the set of all preimages of \mbox{\texttt{\mdseries\slshape elm}} under \mbox{\texttt{\mdseries\slshape map}}. 

 Anything may happen if \mbox{\texttt{\mdseries\slshape elm}} is not an element of the range of \mbox{\texttt{\mdseries\slshape map}}. }

 

\subsection{\textcolor{Chapter }{PreImageElm}}
\logpage{[ 32, 5, 3 ]}\nobreak
\hyperdef{L}{X7D212F727CAE971A}{}
{\noindent\textcolor{FuncColor}{$\triangleright$\enspace\texttt{PreImageElm({\mdseries\slshape map, elm})\index{PreImageElm@\texttt{PreImageElm}}
\label{PreImageElm}
}\hfill{\scriptsize (operation)}}\\


 If \mbox{\texttt{\mdseries\slshape elm}} is an element of the range of the injective and surjective general mapping \mbox{\texttt{\mdseries\slshape map}} then \texttt{PreImageElm} returns the unique preimage of \mbox{\texttt{\mdseries\slshape elm}} under \mbox{\texttt{\mdseries\slshape map}}. 

 Anything may happen if \mbox{\texttt{\mdseries\slshape elm}} is not an element of the range of \mbox{\texttt{\mdseries\slshape map}}. }

 

\subsection{\textcolor{Chapter }{PreImagesRepresentative}}
\logpage{[ 32, 5, 4 ]}\nobreak
\hyperdef{L}{X7AE24A1586B7DE79}{}
{\noindent\textcolor{FuncColor}{$\triangleright$\enspace\texttt{PreImagesRepresentative({\mdseries\slshape map, elm})\index{PreImagesRepresentative@\texttt{PreImagesRepresentative}}
\label{PreImagesRepresentative}
}\hfill{\scriptsize (operation)}}\\


 If \mbox{\texttt{\mdseries\slshape elm}} is an element of the range of the general mapping \mbox{\texttt{\mdseries\slshape map}} then \texttt{PreImagesRepresentative} returns either a representative of the set of preimages of \mbox{\texttt{\mdseries\slshape elm}} under \mbox{\texttt{\mdseries\slshape map}} or \texttt{fail}, the latter if and only if \mbox{\texttt{\mdseries\slshape elm}} has no preimages under \mbox{\texttt{\mdseries\slshape map}}. 

 Anything may happen if \mbox{\texttt{\mdseries\slshape elm}} is not an element of the range of \mbox{\texttt{\mdseries\slshape map}}. }

 

\subsection{\textcolor{Chapter }{PreImagesSet}}
\logpage{[ 32, 5, 5 ]}\nobreak
\hyperdef{L}{X856BAFC87B2D2811}{}
{\noindent\textcolor{FuncColor}{$\triangleright$\enspace\texttt{PreImagesSet({\mdseries\slshape map, elms})\index{PreImagesSet@\texttt{PreImagesSet}}
\label{PreImagesSet}
}\hfill{\scriptsize (operation)}}\\


 If \mbox{\texttt{\mdseries\slshape elms}} is a subset of the range of the general mapping \mbox{\texttt{\mdseries\slshape map}} then \texttt{PreImagesSet} returns the set of all preimages of \mbox{\texttt{\mdseries\slshape elms}} under \mbox{\texttt{\mdseries\slshape map}}. 

 Anything may happen if \mbox{\texttt{\mdseries\slshape elms}} is not a subset of the range of \mbox{\texttt{\mdseries\slshape map}}. }

 
\subsection{\textcolor{Chapter }{PreImage}}\logpage{[ 32, 5, 6 ]}
\hyperdef{L}{X836FAEAC78B55BF4}{}
{
\noindent\textcolor{FuncColor}{$\triangleright$\enspace\texttt{PreImage({\mdseries\slshape map})\index{PreImage@\texttt{PreImage}!set of preimages of the range of a general mapping}
\label{PreImage:set of preimages of the range of a general mapping}
}\hfill{\scriptsize (function)}}\\
\noindent\textcolor{FuncColor}{$\triangleright$\enspace\texttt{PreImage({\mdseries\slshape map, elm})\index{PreImage@\texttt{PreImage}!unique preimage of an element under a general mapping}
\label{PreImage:unique preimage of an element under a general mapping}
}\hfill{\scriptsize (function)}}\\
\noindent\textcolor{FuncColor}{$\triangleright$\enspace\texttt{PreImage({\mdseries\slshape map, coll})\index{PreImage@\texttt{PreImage}!set of preimages of a collection under a general mapping}
\label{PreImage:set of preimages of a collection under a general mapping}
}\hfill{\scriptsize (function)}}\\


 \texttt{PreImage( \mbox{\texttt{\mdseries\slshape map}} )} is the preimage of the general mapping \mbox{\texttt{\mdseries\slshape map}}, i.e., the subset of elements of the source of \mbox{\texttt{\mdseries\slshape map}} that actually have values under \mbox{\texttt{\mdseries\slshape map}}. Note that in this case the argument may also be
non\texttt{\symbol{45}}injective or non\texttt{\symbol{45}}surjective. 

 \texttt{PreImage( \mbox{\texttt{\mdseries\slshape map}}, \mbox{\texttt{\mdseries\slshape elm}} )} is the preimage of the element \mbox{\texttt{\mdseries\slshape elm}} of the range of the injective and surjective mapping \mbox{\texttt{\mdseries\slshape map}} under \mbox{\texttt{\mdseries\slshape map}}, i.e., the unique element of the source which is mapped under \mbox{\texttt{\mdseries\slshape map}} to \mbox{\texttt{\mdseries\slshape elm}}. Note that \mbox{\texttt{\mdseries\slshape map}} must be injective and surjective (see{\nobreakspace}\texttt{PreImages} (\ref{PreImages:set of preimages of the range of a general mapping})). 

 \texttt{PreImage( \mbox{\texttt{\mdseries\slshape map}}, \mbox{\texttt{\mdseries\slshape coll}} )} is the preimage of the subset \mbox{\texttt{\mdseries\slshape coll}} of the range of the general mapping \mbox{\texttt{\mdseries\slshape map}} under \mbox{\texttt{\mdseries\slshape map}}, i.e., the subset of the source which is mapped under \mbox{\texttt{\mdseries\slshape map}} to elements of \mbox{\texttt{\mdseries\slshape coll}}. \mbox{\texttt{\mdseries\slshape coll}} may be a proper set or a domain. The result will be either a proper set or a
domain. Note that in this case \mbox{\texttt{\mdseries\slshape map}} may also be non\texttt{\symbol{45}}injective or
non\texttt{\symbol{45}}surjective. (If \mbox{\texttt{\mdseries\slshape coll}} and the result are lists then the positions of entries do in general \emph{not} correspond.) 

 \texttt{PreImage} delegates to \texttt{PreImagesRange} (\ref{PreImagesRange}) when called with one argument, and to \texttt{PreImageElm} (\ref{PreImageElm}) resp. \texttt{PreImagesSet} (\ref{PreImagesSet}) when called with two arguments. 

 If the second argument is not an element or a subset of the range of the first
argument, an error is signalled. }

 
\subsection{\textcolor{Chapter }{PreImages}}\logpage{[ 32, 5, 7 ]}
\hyperdef{L}{X85C8590E832002EF}{}
{
\noindent\textcolor{FuncColor}{$\triangleright$\enspace\texttt{PreImages({\mdseries\slshape map})\index{PreImages@\texttt{PreImages}!set of preimages of the range of a general mapping}
\label{PreImages:set of preimages of the range of a general mapping}
}\hfill{\scriptsize (function)}}\\
\noindent\textcolor{FuncColor}{$\triangleright$\enspace\texttt{PreImages({\mdseries\slshape map, elm})\index{PreImages@\texttt{PreImages}!set of preimages of an elm under a general mapping}
\label{PreImages:set of preimages of an elm under a general mapping}
}\hfill{\scriptsize (function)}}\\
\noindent\textcolor{FuncColor}{$\triangleright$\enspace\texttt{PreImages({\mdseries\slshape map, coll})\index{PreImages@\texttt{PreImages}!set of preimages of a collection under a general mapping}
\label{PreImages:set of preimages of a collection under a general mapping}
}\hfill{\scriptsize (function)}}\\


 \texttt{PreImages( \mbox{\texttt{\mdseries\slshape map}} )} is the preimage of the general mapping \mbox{\texttt{\mdseries\slshape map}}, i.e., the subset of elements of the source of \mbox{\texttt{\mdseries\slshape map}} that have actually values under \mbox{\texttt{\mdseries\slshape map}}. 

 \texttt{PreImages( \mbox{\texttt{\mdseries\slshape map}}, \mbox{\texttt{\mdseries\slshape elm}} )} is the set of preimages of the element \mbox{\texttt{\mdseries\slshape elm}} of the range of the general mapping \mbox{\texttt{\mdseries\slshape map}} under \mbox{\texttt{\mdseries\slshape map}}, i.e., the set of elements of the source which \mbox{\texttt{\mdseries\slshape map}} maps to \mbox{\texttt{\mdseries\slshape elm}}. 

 \texttt{PreImages( \mbox{\texttt{\mdseries\slshape map}}, \mbox{\texttt{\mdseries\slshape coll}} )} is the set of images of the subset \mbox{\texttt{\mdseries\slshape coll}} of the range of the general mapping \mbox{\texttt{\mdseries\slshape map}} under \mbox{\texttt{\mdseries\slshape map}}, i.e., the subset of the source which \mbox{\texttt{\mdseries\slshape map}} maps to elements of \mbox{\texttt{\mdseries\slshape coll}}. \mbox{\texttt{\mdseries\slshape coll}} may be a proper set or a domain. The result will be either a proper set or a
domain. (If \mbox{\texttt{\mdseries\slshape coll}} and the result are lists then the positions of entries do in general \emph{not} correspond.) 

 \texttt{PreImages} delegates to \texttt{PreImagesRange} (\ref{PreImagesRange}) when called with one argument, and to \texttt{PreImagesElm} (\ref{PreImagesElm}) resp. \texttt{PreImagesSet} (\ref{PreImagesSet}) when called with two arguments. 

 If the second argument is not an element or a subset of the range of the first
argument, an error is signalled. }

 }

  
\section{\textcolor{Chapter }{Arithmetic Operations for General Mappings}}\label{Arithmetic Operations for General Mappings}
\logpage{[ 32, 6, 0 ]}
\hyperdef{L}{X7E2E16277940FA0B}{}
{
  General mappings are arithmetic objects. One can form groups and vector spaces
of general mappings provided that they are invertible or can be added and
admit scalar multiplication, respectively. 

 For two general mappings with same source, range, preimage, and image, the \emph{sum} is defined pointwise, i.e., the images of a point under the sum is the set of
all sums with first summand in the images of the first general mapping and
second summand in the images of the second general mapping. 

 \emph{Scalar multiplication} of general mappings is defined likewise. 

 The \emph{product} of two general mappings is defined as the composition. This multiplication is
always associative. In addition to the composition via \texttt{*}, general mappings can be composed {\textendash}in reversed order{\textendash}
via \texttt{CompositionMapping} (\ref{CompositionMapping}). 

 General mappings are in the category of multiplicative elements with inverses.
Similar to matrices, not every general mapping has an inverse or an identity,
and we define the behaviour of \texttt{One} (\ref{One}) and \texttt{Inverse} (\ref{Inverse}) for general mappings as follows. \texttt{One} (\ref{One}) returns \texttt{fail} when called for a general mapping whose source and range differ, otherwise \texttt{One} (\ref{One}) returns the identity mapping of the source. (Note that the source may differ
from the preimage). \texttt{Inverse} (\ref{Inverse}) returns \texttt{fail} when called for a non\texttt{\symbol{45}}bijective general mapping or for a
general mapping whose source and range differ; otherwise \texttt{Inverse} (\ref{Inverse}) returns the inverse mapping. 

 Besides the usual inverse of multiplicative elements, which means that \texttt{Inverse( \mbox{\texttt{\mdseries\slshape g}} ) * \mbox{\texttt{\mdseries\slshape g}} = \mbox{\texttt{\mdseries\slshape g}} * Inverse( \mbox{\texttt{\mdseries\slshape g}} ) = One( \mbox{\texttt{\mdseries\slshape g}} )}, for general mappings we have the attribute \texttt{InverseGeneralMapping} (\ref{InverseGeneralMapping}). If \mbox{\texttt{\mdseries\slshape F}} is a general mapping with source $S$, range $R$, and underlying relation $Rel$ then \texttt{InverseGeneralMapping( \mbox{\texttt{\mdseries\slshape F}} )} has source $R$, range $S$, and underlying relation $\{ (r,s) \mid (s,r) \in Rel \}$. For a general mapping that has an inverse in the usual sense, i.e., for a
bijection of the source, of course both concepts coincide. 

 \texttt{Inverse} (\ref{Inverse}) may delegate to \texttt{InverseGeneralMapping} (\ref{InverseGeneralMapping}). \texttt{InverseGeneralMapping} (\ref{InverseGeneralMapping}) must not delegate to \texttt{Inverse} (\ref{Inverse}), but a known value of \texttt{Inverse} (\ref{Inverse}) may be fetched. So methods to compute the inverse of a general mapping should
be installed for \texttt{InverseGeneralMapping} (\ref{InverseGeneralMapping}). 

 (Note that in many respects, general mappings behave similar to matrices, for
example one can define left and right identities and inverses, which do not
fit into the current concepts of \textsf{GAP}.) }

  
\section{\textcolor{Chapter }{Mappings which are Compatible with Algebraic Structures}}\label{Mappings which are Compatible with Algebraic Structures}
\logpage{[ 32, 7, 0 ]}
\hyperdef{L}{X834E02BB7D4B4AE5}{}
{
  From an algebraical point of view, the most important mappings are those which
are compatible with a structure. For Magmas, Groups and Rings, \textsf{GAP} supports the following four types of such mappings: 

 
\begin{enumerate}
\item  General mappings that respect multiplication 
\item  General mappings that respect addition 
\item  General mappings that respect scalar mult. 
\item  General mappings that respect multiplicative and additive structure 
\end{enumerate}
 

 (Very technical note: \textsf{GAP} defines categories \texttt{IsSPGeneralMapping} and \texttt{IsNonSPGeneralMapping}. The distinction between these is orthogonal to the structure compatibility
described here and should not be confused.) }

  
\section{\textcolor{Chapter }{Magma Homomorphisms}}\label{Magma Homomorphisms}
\logpage{[ 32, 8, 0 ]}
\hyperdef{L}{X8008FCCC7F4C731F}{}
{
  

\subsection{\textcolor{Chapter }{IsMagmaHomomorphism}}
\logpage{[ 32, 8, 1 ]}\nobreak
\hyperdef{L}{X7DC72CF28539A251}{}
{\noindent\textcolor{FuncColor}{$\triangleright$\enspace\texttt{IsMagmaHomomorphism({\mdseries\slshape mapp})\index{IsMagmaHomomorphism@\texttt{IsMagmaHomomorphism}}
\label{IsMagmaHomomorphism}
}\hfill{\scriptsize (filter)}}\\


 A \emph{magma homomorphism} is a total single valued mapping which respects multiplication. }

 

\subsection{\textcolor{Chapter }{MagmaHomomorphismByFunctionNC}}
\logpage{[ 32, 8, 2 ]}\nobreak
\hyperdef{L}{X8181676787E760A2}{}
{\noindent\textcolor{FuncColor}{$\triangleright$\enspace\texttt{MagmaHomomorphismByFunctionNC({\mdseries\slshape G, H, fn})\index{MagmaHomomorphismByFunctionNC@\texttt{MagmaHomomorphismByFunctionNC}}
\label{MagmaHomomorphismByFunctionNC}
}\hfill{\scriptsize (function)}}\\


 Creates the homomorphism from \mbox{\texttt{\mdseries\slshape G}} to \mbox{\texttt{\mdseries\slshape H}} without checking that \mbox{\texttt{\mdseries\slshape fn}} is a homomorphism. }

 

\subsection{\textcolor{Chapter }{NaturalHomomorphismByGenerators}}
\logpage{[ 32, 8, 3 ]}\nobreak
\hyperdef{L}{X79D0216E871B7051}{}
{\noindent\textcolor{FuncColor}{$\triangleright$\enspace\texttt{NaturalHomomorphismByGenerators({\mdseries\slshape f, s})\index{NaturalHomomorphismByGenerators@\texttt{NaturalHomomorphismByGenerators}}
\label{NaturalHomomorphismByGenerators}
}\hfill{\scriptsize (operation)}}\\


 returns a mapping from the magma \mbox{\texttt{\mdseries\slshape f}} with $n$ generators to the magma \mbox{\texttt{\mdseries\slshape s}} with $n$ generators, which maps the $i$\texttt{\symbol{45}}th generator of \mbox{\texttt{\mdseries\slshape f}} to the $i$\texttt{\symbol{45}}th generator of \mbox{\texttt{\mdseries\slshape s}}. }

 }

  
\section{\textcolor{Chapter }{Mappings that Respect Multiplication}}\label{Mappings that Respect Multiplication}
\logpage{[ 32, 9, 0 ]}
\hyperdef{L}{X806F892C862F29F9}{}
{
  

\subsection{\textcolor{Chapter }{RespectsMultiplication}}
\logpage{[ 32, 9, 1 ]}\nobreak
\hyperdef{L}{X7BEFF95883EAEC78}{}
{\noindent\textcolor{FuncColor}{$\triangleright$\enspace\texttt{RespectsMultiplication({\mdseries\slshape mapp})\index{RespectsMultiplication@\texttt{RespectsMultiplication}}
\label{RespectsMultiplication}
}\hfill{\scriptsize (property)}}\\


 Let \mbox{\texttt{\mdseries\slshape mapp}} be a general mapping with underlying relation $F \subseteq S \times R$, where $S$ and $R$ are the source and the range of \mbox{\texttt{\mdseries\slshape mapp}}, respectively. Then \texttt{RespectsMultiplication} returns \texttt{true} if $S$ and $R$ are magmas such that $(s_1,r_1), (s_2,r_2) \in F$ implies $(s_1 * s_2,r_1 * r_2) \in F$, and \texttt{false} otherwise. 

 If \mbox{\texttt{\mdseries\slshape mapp}} is single\texttt{\symbol{45}}valued then \texttt{RespectsMultiplication} returns \texttt{true} if and only if the equation \texttt{\mbox{\texttt{\mdseries\slshape s1}}\texttt{\symbol{94}}\mbox{\texttt{\mdseries\slshape mapp}} * \mbox{\texttt{\mdseries\slshape s2}}\texttt{\symbol{94}}\mbox{\texttt{\mdseries\slshape mapp}} = (\mbox{\texttt{\mdseries\slshape s1}} * \mbox{\texttt{\mdseries\slshape s2}})\texttt{\symbol{94}}\mbox{\texttt{\mdseries\slshape mapp}}} holds for all \mbox{\texttt{\mdseries\slshape s1}}, \mbox{\texttt{\mdseries\slshape s2}} in $S$. }

 

\subsection{\textcolor{Chapter }{RespectsOne}}
\logpage{[ 32, 9, 2 ]}\nobreak
\hyperdef{L}{X7EE4DA097AE9CBC1}{}
{\noindent\textcolor{FuncColor}{$\triangleright$\enspace\texttt{RespectsOne({\mdseries\slshape mapp})\index{RespectsOne@\texttt{RespectsOne}}
\label{RespectsOne}
}\hfill{\scriptsize (property)}}\\


 Let \mbox{\texttt{\mdseries\slshape mapp}} be a general mapping with underlying relation $F \subseteq \mbox{\texttt{\mdseries\slshape S}} \times \mbox{\texttt{\mdseries\slshape R}}$, where \mbox{\texttt{\mdseries\slshape S}} and \mbox{\texttt{\mdseries\slshape R}} are the source and the range of \mbox{\texttt{\mdseries\slshape mapp}}, respectively. Then \texttt{RespectsOne} returns \texttt{true} if \mbox{\texttt{\mdseries\slshape S}} and \mbox{\texttt{\mdseries\slshape R}} are magmas\texttt{\symbol{45}}with\texttt{\symbol{45}}one such that $( $\texttt{One(\mbox{\texttt{\mdseries\slshape S}})}$, $\texttt{One(\mbox{\texttt{\mdseries\slshape R}})}$ ) \in F$, and \texttt{false} otherwise. 

 If \mbox{\texttt{\mdseries\slshape mapp}} is single\texttt{\symbol{45}}valued then \texttt{RespectsOne} returns \texttt{true} if and only if the equation \texttt{One( \mbox{\texttt{\mdseries\slshape S}} )\texttt{\symbol{94}}\mbox{\texttt{\mdseries\slshape mapp}} = One( \mbox{\texttt{\mdseries\slshape R}} )} holds. }

 

\subsection{\textcolor{Chapter }{RespectsInverses}}
\logpage{[ 32, 9, 3 ]}\nobreak
\hyperdef{L}{X7F27AE9C84A4DF90}{}
{\noindent\textcolor{FuncColor}{$\triangleright$\enspace\texttt{RespectsInverses({\mdseries\slshape mapp})\index{RespectsInverses@\texttt{RespectsInverses}}
\label{RespectsInverses}
}\hfill{\scriptsize (property)}}\\


 Let \mbox{\texttt{\mdseries\slshape mapp}} be a general mapping with underlying relation $F \subseteq \mbox{\texttt{\mdseries\slshape S}} \times \mbox{\texttt{\mdseries\slshape R}}$, where \mbox{\texttt{\mdseries\slshape S}} and \mbox{\texttt{\mdseries\slshape R}} are the source and the range of \mbox{\texttt{\mdseries\slshape mapp}}, respectively. Then \texttt{RespectsInverses} returns \texttt{true} if \mbox{\texttt{\mdseries\slshape S}} and \mbox{\texttt{\mdseries\slshape R}} are magmas\texttt{\symbol{45}}with\texttt{\symbol{45}}inverses such that, for $s \in \mbox{\texttt{\mdseries\slshape S}}$ and $r \in \mbox{\texttt{\mdseries\slshape R}}$, $(s,r) \in F$ implies $(s^{{-1}},r^{{-1}}) \in F$, and \texttt{false} otherwise. 

 If \mbox{\texttt{\mdseries\slshape mapp}} is single\texttt{\symbol{45}}valued then \texttt{RespectsInverses} returns \texttt{true} if and only if the equation \texttt{Inverse( \mbox{\texttt{\mdseries\slshape s}} )\texttt{\symbol{94}}\mbox{\texttt{\mdseries\slshape mapp}} = Inverse( \mbox{\texttt{\mdseries\slshape s}}\texttt{\symbol{94}}\mbox{\texttt{\mdseries\slshape mapp}} )} holds for all \mbox{\texttt{\mdseries\slshape s}} in $S$. }

 

\subsection{\textcolor{Chapter }{IsGroupGeneralMapping}}
\logpage{[ 32, 9, 4 ]}\nobreak
\hyperdef{L}{X819DD174829BF3AE}{}
{\noindent\textcolor{FuncColor}{$\triangleright$\enspace\texttt{IsGroupGeneralMapping({\mdseries\slshape mapp})\index{IsGroupGeneralMapping@\texttt{IsGroupGeneralMapping}}
\label{IsGroupGeneralMapping}
}\hfill{\scriptsize (filter)}}\\
\noindent\textcolor{FuncColor}{$\triangleright$\enspace\texttt{IsGroupHomomorphism({\mdseries\slshape mapp})\index{IsGroupHomomorphism@\texttt{IsGroupHomomorphism}}
\label{IsGroupHomomorphism}
}\hfill{\scriptsize (filter)}}\\


 A \emph{group general mapping} is a mapping which respects multiplication and inverses. If it is total and
single valued it is called a \emph{group homomorphism}. 

 Chapter{\nobreakspace}\ref{Group Homomorphisms} explains group homomorphisms in more detail. }

 

\subsection{\textcolor{Chapter }{KernelOfMultiplicativeGeneralMapping}}
\logpage{[ 32, 9, 5 ]}\nobreak
\hyperdef{L}{X81A5A5CF846E5FBF}{}
{\noindent\textcolor{FuncColor}{$\triangleright$\enspace\texttt{KernelOfMultiplicativeGeneralMapping({\mdseries\slshape mapp})\index{KernelOfMultiplicativeGeneralMapping@\texttt{Kernel}\-\texttt{Of}\-\texttt{Multiplicative}\-\texttt{General}\-\texttt{Mapping}}
\label{KernelOfMultiplicativeGeneralMapping}
}\hfill{\scriptsize (attribute)}}\\


 Let \mbox{\texttt{\mdseries\slshape mapp}} be a general mapping. Then \texttt{KernelOfMultiplicativeGeneralMapping} returns the set of all elements in the source of \mbox{\texttt{\mdseries\slshape mapp}} that have the identity of the range in their set of images. 

 (This is a monoid if \mbox{\texttt{\mdseries\slshape mapp}} respects multiplication and one, and if the source of \mbox{\texttt{\mdseries\slshape mapp}} is associative.) }

 

\subsection{\textcolor{Chapter }{CoKernelOfMultiplicativeGeneralMapping}}
\logpage{[ 32, 9, 6 ]}\nobreak
\hyperdef{L}{X7F09B6E28080DCB4}{}
{\noindent\textcolor{FuncColor}{$\triangleright$\enspace\texttt{CoKernelOfMultiplicativeGeneralMapping({\mdseries\slshape mapp})\index{CoKernelOfMultiplicativeGeneralMapping@\texttt{CoKernel}\-\texttt{Of}\-\texttt{Multiplicative}\-\texttt{General}\-\texttt{Mapping}}
\label{CoKernelOfMultiplicativeGeneralMapping}
}\hfill{\scriptsize (attribute)}}\\


 Let \mbox{\texttt{\mdseries\slshape mapp}} be a general mapping. Then \texttt{CoKernelOfMultiplicativeGeneralMapping} returns the set of all elements in the range of \mbox{\texttt{\mdseries\slshape mapp}} that have the identity of the source in their set of preimages. 

 (This is a monoid if \mbox{\texttt{\mdseries\slshape mapp}} respects multiplication and one, and if the range of \mbox{\texttt{\mdseries\slshape mapp}} is associative.) }

 }

  
\section{\textcolor{Chapter }{Mappings that Respect Addition}}\label{Mappings that Respect Addition}
\logpage{[ 32, 10, 0 ]}
\hyperdef{L}{X8455A5A67C35178B}{}
{
  

\subsection{\textcolor{Chapter }{RespectsAddition}}
\logpage{[ 32, 10, 1 ]}\nobreak
\hyperdef{L}{X7A3321E878925C3A}{}
{\noindent\textcolor{FuncColor}{$\triangleright$\enspace\texttt{RespectsAddition({\mdseries\slshape mapp})\index{RespectsAddition@\texttt{RespectsAddition}}
\label{RespectsAddition}
}\hfill{\scriptsize (property)}}\\


 Let \mbox{\texttt{\mdseries\slshape mapp}} be a general mapping with underlying relation $F \subseteq S \times R$, where $S$ and $R$ are the source and the range of \mbox{\texttt{\mdseries\slshape mapp}}, respectively. Then \texttt{RespectsAddition} returns \texttt{true} if $S$ and $R$ are additive magmas such that $(s_1,r_1), (s_2,r_2) \in F$ implies $(s_1 + s_2,r_1 + r_2) \in F$, and \texttt{false} otherwise. 

 If \mbox{\texttt{\mdseries\slshape mapp}} is single\texttt{\symbol{45}}valued then \texttt{RespectsAddition} returns \texttt{true} if and only if the equation \texttt{\mbox{\texttt{\mdseries\slshape s1}}\texttt{\symbol{94}}\mbox{\texttt{\mdseries\slshape mapp}} + \mbox{\texttt{\mdseries\slshape s2}}\texttt{\symbol{94}}\mbox{\texttt{\mdseries\slshape mapp}} = (\mbox{\texttt{\mdseries\slshape s1}}+\mbox{\texttt{\mdseries\slshape s2}})\texttt{\symbol{94}}\mbox{\texttt{\mdseries\slshape mapp}}} holds for all \mbox{\texttt{\mdseries\slshape s1}}, \mbox{\texttt{\mdseries\slshape s2}} in $S$. }

 

\subsection{\textcolor{Chapter }{RespectsAdditiveInverses}}
\logpage{[ 32, 10, 2 ]}\nobreak
\hyperdef{L}{X8130D8907B92F746}{}
{\noindent\textcolor{FuncColor}{$\triangleright$\enspace\texttt{RespectsAdditiveInverses({\mdseries\slshape mapp})\index{RespectsAdditiveInverses@\texttt{RespectsAdditiveInverses}}
\label{RespectsAdditiveInverses}
}\hfill{\scriptsize (property)}}\\


 Let \mbox{\texttt{\mdseries\slshape mapp}} be a general mapping with underlying relation $F \subseteq S \times R$, where $S$ and $R$ are the source and the range of \mbox{\texttt{\mdseries\slshape mapp}}, respectively. Then \texttt{RespectsAdditiveInverses} returns \texttt{true} if $S$ and $R$ are
additive\texttt{\symbol{45}}magmas\texttt{\symbol{45}}with\texttt{\symbol{45}}inverses
such that $(s,r) \in F$ implies $(-s,-r) \in F$, and \texttt{false} otherwise. 

 If \mbox{\texttt{\mdseries\slshape mapp}} is single\texttt{\symbol{45}}valued then \texttt{RespectsAdditiveInverses} returns \texttt{true} if and only if the equation \texttt{AdditiveInverse( \mbox{\texttt{\mdseries\slshape s}} )\texttt{\symbol{94}}\mbox{\texttt{\mdseries\slshape mapp}} = AdditiveInverse( \mbox{\texttt{\mdseries\slshape s}}\texttt{\symbol{94}}\mbox{\texttt{\mdseries\slshape mapp}} )} holds for all \mbox{\texttt{\mdseries\slshape s}} in $S$. }

 

\subsection{\textcolor{Chapter }{RespectsZero}}
\logpage{[ 32, 10, 3 ]}\nobreak
\hyperdef{L}{X7D342736781EB280}{}
{\noindent\textcolor{FuncColor}{$\triangleright$\enspace\texttt{RespectsZero({\mdseries\slshape mapp})\index{RespectsZero@\texttt{RespectsZero}}
\label{RespectsZero}
}\hfill{\scriptsize (property)}}\\


 Let \mbox{\texttt{\mdseries\slshape mapp}} be a general mapping with underlying relation $F \subseteq \mbox{\texttt{\mdseries\slshape S}} \times \mbox{\texttt{\mdseries\slshape R}}$, where \mbox{\texttt{\mdseries\slshape S}} and \mbox{\texttt{\mdseries\slshape R}} are the source and the range of \mbox{\texttt{\mdseries\slshape mapp}}, respectively. Then \texttt{RespectsZero} returns \texttt{true} if \mbox{\texttt{\mdseries\slshape S}} and \mbox{\texttt{\mdseries\slshape R}} are
additive\texttt{\symbol{45}}magmas\texttt{\symbol{45}}with\texttt{\symbol{45}}zero
such that $( $\texttt{Zero(\mbox{\texttt{\mdseries\slshape S}})}$, $\texttt{Zero(\mbox{\texttt{\mdseries\slshape R}})}$ ) \in F$, and \texttt{false} otherwise. 

 If \mbox{\texttt{\mdseries\slshape mapp}} is single\texttt{\symbol{45}}valued then \texttt{RespectsZero} returns \texttt{true} if and only if the equation \texttt{Zero( \mbox{\texttt{\mdseries\slshape S}} )\texttt{\symbol{94}}\mbox{\texttt{\mdseries\slshape mapp}} = Zero( \mbox{\texttt{\mdseries\slshape R}} )} holds. }

 

\subsection{\textcolor{Chapter }{IsAdditiveGroupGeneralMapping}}
\logpage{[ 32, 10, 4 ]}\nobreak
\hyperdef{L}{X7B99EF287A8A0BD9}{}
{\noindent\textcolor{FuncColor}{$\triangleright$\enspace\texttt{IsAdditiveGroupGeneralMapping({\mdseries\slshape mapp})\index{IsAdditiveGroupGeneralMapping@\texttt{IsAdditiveGroupGeneralMapping}}
\label{IsAdditiveGroupGeneralMapping}
}\hfill{\scriptsize (filter)}}\\
\noindent\textcolor{FuncColor}{$\triangleright$\enspace\texttt{IsAdditiveGroupHomomorphism({\mdseries\slshape mapp})\index{IsAdditiveGroupHomomorphism@\texttt{IsAdditiveGroupHomomorphism}}
\label{IsAdditiveGroupHomomorphism}
}\hfill{\scriptsize (filter)}}\\


 \texttt{IsAdditiveGroupGeneralMapping} specifies whether a general mapping \mbox{\texttt{\mdseries\slshape mapp}} respects addition (see \texttt{RespectsAddition} (\ref{RespectsAddition})) and respects additive inverses (see \texttt{RespectsAdditiveInverses} (\ref{RespectsAdditiveInverses})). 

 \texttt{IsAdditiveGroupHomomorphism} is a synonym for the meet of \texttt{IsAdditiveGroupGeneralMapping} and \texttt{IsMapping} (\ref{IsMapping}). }

 

\subsection{\textcolor{Chapter }{KernelOfAdditiveGeneralMapping}}
\logpage{[ 32, 10, 5 ]}\nobreak
\hyperdef{L}{X7EC0E9907D6631D6}{}
{\noindent\textcolor{FuncColor}{$\triangleright$\enspace\texttt{KernelOfAdditiveGeneralMapping({\mdseries\slshape mapp})\index{KernelOfAdditiveGeneralMapping@\texttt{KernelOfAdditiveGeneralMapping}}
\label{KernelOfAdditiveGeneralMapping}
}\hfill{\scriptsize (attribute)}}\\


 Let \mbox{\texttt{\mdseries\slshape mapp}} be a general mapping. Then \texttt{KernelOfAdditiveGeneralMapping} returns the set of all elements in the source of \mbox{\texttt{\mdseries\slshape mapp}} that have the zero of the range in their set of images. }

 

\subsection{\textcolor{Chapter }{CoKernelOfAdditiveGeneralMapping}}
\logpage{[ 32, 10, 6 ]}\nobreak
\hyperdef{L}{X813C6D7980213F41}{}
{\noindent\textcolor{FuncColor}{$\triangleright$\enspace\texttt{CoKernelOfAdditiveGeneralMapping({\mdseries\slshape mapp})\index{CoKernelOfAdditiveGeneralMapping@\texttt{CoKernelOfAdditiveGeneralMapping}}
\label{CoKernelOfAdditiveGeneralMapping}
}\hfill{\scriptsize (attribute)}}\\


 Let \mbox{\texttt{\mdseries\slshape mapp}} be a general mapping. Then \texttt{CoKernelOfAdditiveGeneralMapping} returns the set of all elements in the range of \mbox{\texttt{\mdseries\slshape mapp}} that have the zero of the source in their set of preimages. }

 }

  
\section{\textcolor{Chapter }{Linear Mappings}}\label{Linear Mappings}
\logpage{[ 32, 11, 0 ]}
\hyperdef{L}{X7C24431C81532575}{}
{
  Also see Sections{\nobreakspace}\ref{Mappings that Respect Multiplication}, \ref{Mappings that Respect Addition}, and \texttt{KernelOfMultiplicativeGeneralMapping} (\ref{KernelOfMultiplicativeGeneralMapping}), \texttt{CoKernelOfMultiplicativeGeneralMapping} (\ref{CoKernelOfMultiplicativeGeneralMapping}). 

\subsection{\textcolor{Chapter }{RespectsScalarMultiplication}}
\logpage{[ 32, 11, 1 ]}\nobreak
\hyperdef{L}{X87842ED97FA19973}{}
{\noindent\textcolor{FuncColor}{$\triangleright$\enspace\texttt{RespectsScalarMultiplication({\mdseries\slshape mapp})\index{RespectsScalarMultiplication@\texttt{RespectsScalarMultiplication}}
\label{RespectsScalarMultiplication}
}\hfill{\scriptsize (property)}}\\


 Let \mbox{\texttt{\mdseries\slshape mapp}} be a general mapping, with underlying relation $F \subseteq S \times R$, where $S$ and $R$ are the source and the range of \mbox{\texttt{\mdseries\slshape mapp}}, respectively. Then \texttt{RespectsScalarMultiplication} returns \texttt{true} if $S$ and $R$ are left modules with the left acting domain $D$ of $S$ contained in the left acting domain of $R$ and such that $(s,r) \in F$ implies $(c * s,c * r) \in F$ for all $c \in D$, and \texttt{false} otherwise. 

 If \mbox{\texttt{\mdseries\slshape mapp}} is single\texttt{\symbol{45}}valued then \texttt{RespectsScalarMultiplication} returns \texttt{true} if and only if the equation \texttt{\mbox{\texttt{\mdseries\slshape c}} * \mbox{\texttt{\mdseries\slshape s}}\texttt{\symbol{94}}\mbox{\texttt{\mdseries\slshape mapp}} = (\mbox{\texttt{\mdseries\slshape c}} * \mbox{\texttt{\mdseries\slshape s}})\texttt{\symbol{94}}\mbox{\texttt{\mdseries\slshape mapp}}} holds for all \mbox{\texttt{\mdseries\slshape c}} in $D$ and \mbox{\texttt{\mdseries\slshape s}} in $S$. }

 

\subsection{\textcolor{Chapter }{IsLeftModuleGeneralMapping}}
\logpage{[ 32, 11, 2 ]}\nobreak
\hyperdef{L}{X780BE6307A3271A9}{}
{\noindent\textcolor{FuncColor}{$\triangleright$\enspace\texttt{IsLeftModuleGeneralMapping({\mdseries\slshape mapp})\index{IsLeftModuleGeneralMapping@\texttt{IsLeftModuleGeneralMapping}}
\label{IsLeftModuleGeneralMapping}
}\hfill{\scriptsize (filter)}}\\
\noindent\textcolor{FuncColor}{$\triangleright$\enspace\texttt{IsLeftModuleHomomorphism({\mdseries\slshape mapp})\index{IsLeftModuleHomomorphism@\texttt{IsLeftModuleHomomorphism}}
\label{IsLeftModuleHomomorphism}
}\hfill{\scriptsize (filter)}}\\


 \texttt{IsLeftModuleGeneralMapping} specifies whether a general mapping \mbox{\texttt{\mdseries\slshape mapp}} satisfies the property \texttt{IsAdditiveGroupGeneralMapping} (\ref{IsAdditiveGroupGeneralMapping}) and respects scalar multiplication (see \texttt{RespectsScalarMultiplication} (\ref{RespectsScalarMultiplication})). 

 \texttt{IsLeftModuleHomomorphism} is a synonym for the meet of \texttt{IsLeftModuleGeneralMapping} and \texttt{IsMapping} (\ref{IsMapping}). }

 

\subsection{\textcolor{Chapter }{IsLinearMapping}}
\logpage{[ 32, 11, 3 ]}\nobreak
\hyperdef{L}{X7F6841107E59107F}{}
{\noindent\textcolor{FuncColor}{$\triangleright$\enspace\texttt{IsLinearMapping({\mdseries\slshape F, mapp})\index{IsLinearMapping@\texttt{IsLinearMapping}}
\label{IsLinearMapping}
}\hfill{\scriptsize (operation)}}\\


 For a field \mbox{\texttt{\mdseries\slshape F}} and a general mapping \mbox{\texttt{\mdseries\slshape mapp}}, \texttt{IsLinearMapping} returns \texttt{true} if \mbox{\texttt{\mdseries\slshape mapp}} is an \mbox{\texttt{\mdseries\slshape F}}\texttt{\symbol{45}}linear mapping, and \texttt{false} otherwise. 

 A mapping $f$ is a linear mapping (or vector space homomorphism) if the source and range are
vector spaces over the same division ring $D$, and if $f( a + b ) = f(a) + f(b)$ and $f( s * a ) = s * f(a)$ hold for all elements $a$, $b$ in the source of $f$ and $s \in D$. }

 }

  
\section{\textcolor{Chapter }{Ring Homomorphisms}}\label{Ring Homomorphisms}
\logpage{[ 32, 12, 0 ]}
\hyperdef{L}{X7E88C32A82E942DA}{}
{
  

\subsection{\textcolor{Chapter }{IsRingGeneralMapping}}
\logpage{[ 32, 12, 1 ]}\nobreak
\hyperdef{L}{X7C8DA031799B79D5}{}
{\noindent\textcolor{FuncColor}{$\triangleright$\enspace\texttt{IsRingGeneralMapping({\mdseries\slshape mapp})\index{IsRingGeneralMapping@\texttt{IsRingGeneralMapping}}
\label{IsRingGeneralMapping}
}\hfill{\scriptsize (filter)}}\\
\noindent\textcolor{FuncColor}{$\triangleright$\enspace\texttt{IsRingHomomorphism({\mdseries\slshape mapp})\index{IsRingHomomorphism@\texttt{IsRingHomomorphism}}
\label{IsRingHomomorphism}
}\hfill{\scriptsize (filter)}}\\


 \texttt{IsRingGeneralMapping} specifies whether a general mapping \mbox{\texttt{\mdseries\slshape mapp}} satisfies the property \texttt{IsAdditiveGroupGeneralMapping} (\ref{IsAdditiveGroupGeneralMapping}) and respects multiplication (see \texttt{RespectsMultiplication} (\ref{RespectsMultiplication})). 

 \texttt{IsRingHomomorphism} is a synonym for the meet of \texttt{IsRingGeneralMapping} and \texttt{IsMapping} (\ref{IsMapping}). }

 

\subsection{\textcolor{Chapter }{IsRingWithOneGeneralMapping}}
\logpage{[ 32, 12, 2 ]}\nobreak
\hyperdef{L}{X7988102883675606}{}
{\noindent\textcolor{FuncColor}{$\triangleright$\enspace\texttt{IsRingWithOneGeneralMapping({\mdseries\slshape mapp})\index{IsRingWithOneGeneralMapping@\texttt{IsRingWithOneGeneralMapping}}
\label{IsRingWithOneGeneralMapping}
}\hfill{\scriptsize (filter)}}\\
\noindent\textcolor{FuncColor}{$\triangleright$\enspace\texttt{IsRingWithOneHomomorphism({\mdseries\slshape mapp})\index{IsRingWithOneHomomorphism@\texttt{IsRingWithOneHomomorphism}}
\label{IsRingWithOneHomomorphism}
}\hfill{\scriptsize (filter)}}\\


 }

 

\subsection{\textcolor{Chapter }{IsAlgebraGeneralMapping}}
\logpage{[ 32, 12, 3 ]}\nobreak
\hyperdef{L}{X86B14F908601DEA9}{}
{\noindent\textcolor{FuncColor}{$\triangleright$\enspace\texttt{IsAlgebraGeneralMapping({\mdseries\slshape mapp})\index{IsAlgebraGeneralMapping@\texttt{IsAlgebraGeneralMapping}}
\label{IsAlgebraGeneralMapping}
}\hfill{\scriptsize (filter)}}\\
\noindent\textcolor{FuncColor}{$\triangleright$\enspace\texttt{IsAlgebraHomomorphism({\mdseries\slshape mapp})\index{IsAlgebraHomomorphism@\texttt{IsAlgebraHomomorphism}}
\label{IsAlgebraHomomorphism}
}\hfill{\scriptsize (filter)}}\\


 \texttt{IsAlgebraGeneralMapping} specifies whether a general mapping \mbox{\texttt{\mdseries\slshape mapp}} satisfies both properties \texttt{IsRingGeneralMapping} (\ref{IsRingGeneralMapping}) and (see \texttt{IsLeftModuleGeneralMapping} (\ref{IsLeftModuleGeneralMapping})). 

 \texttt{IsAlgebraHomomorphism} is a synonym for the meet of \texttt{IsAlgebraGeneralMapping} and \texttt{IsMapping} (\ref{IsMapping}). }

 

\subsection{\textcolor{Chapter }{IsAlgebraWithOneGeneralMapping}}
\logpage{[ 32, 12, 4 ]}\nobreak
\hyperdef{L}{X842AD44679C5BDC2}{}
{\noindent\textcolor{FuncColor}{$\triangleright$\enspace\texttt{IsAlgebraWithOneGeneralMapping({\mdseries\slshape mapp})\index{IsAlgebraWithOneGeneralMapping@\texttt{IsAlgebraWithOneGeneralMapping}}
\label{IsAlgebraWithOneGeneralMapping}
}\hfill{\scriptsize (filter)}}\\
\noindent\textcolor{FuncColor}{$\triangleright$\enspace\texttt{IsAlgebraWithOneHomomorphism({\mdseries\slshape mapp})\index{IsAlgebraWithOneHomomorphism@\texttt{IsAlgebraWithOneHomomorphism}}
\label{IsAlgebraWithOneHomomorphism}
}\hfill{\scriptsize (filter)}}\\


 \texttt{IsAlgebraWithOneGeneralMapping} specifies whether a general mapping \mbox{\texttt{\mdseries\slshape mapp}} satisfies both properties \texttt{IsAlgebraGeneralMapping} (\ref{IsAlgebraGeneralMapping}) and \texttt{RespectsOne} (\ref{RespectsOne}). 

 \texttt{IsAlgebraWithOneHomomorphism} is a synonym for the meet of \texttt{IsAlgebraWithOneGeneralMapping} and \texttt{IsMapping} (\ref{IsMapping}). }

 

\subsection{\textcolor{Chapter }{IsFieldHomomorphism}}
\logpage{[ 32, 12, 5 ]}\nobreak
\hyperdef{L}{X8324DA78879DF4D7}{}
{\noindent\textcolor{FuncColor}{$\triangleright$\enspace\texttt{IsFieldHomomorphism({\mdseries\slshape mapp})\index{IsFieldHomomorphism@\texttt{IsFieldHomomorphism}}
\label{IsFieldHomomorphism}
}\hfill{\scriptsize (property)}}\\


 A general mapping is a field homomorphism if and only if it is a ring
homomorphism with source a field. }

 }

  
\section{\textcolor{Chapter }{General Mappings}}\label{General Mappings}
\logpage{[ 32, 13, 0 ]}
\hyperdef{L}{X7E4A55567BED0F88}{}
{
  

\subsection{\textcolor{Chapter }{IsGeneralMapping}}
\logpage{[ 32, 13, 1 ]}\nobreak
\hyperdef{L}{X8656AB8A7D672CAE}{}
{\noindent\textcolor{FuncColor}{$\triangleright$\enspace\texttt{IsGeneralMapping({\mdseries\slshape map})\index{IsGeneralMapping@\texttt{IsGeneralMapping}}
\label{IsGeneralMapping}
}\hfill{\scriptsize (Category)}}\\


 Each general mapping lies in the category \texttt{IsGeneralMapping}. It implies the categories \texttt{IsMultiplicativeElementWithInverse} (\ref{IsMultiplicativeElementWithInverse}) and \texttt{IsAssociativeElement} (\ref{IsAssociativeElement}); for a discussion of these implications, see{\nobreakspace}\ref{Arithmetic Operations for General Mappings}. }

 

\subsection{\textcolor{Chapter }{IsConstantTimeAccessGeneralMapping}}
\logpage{[ 32, 13, 2 ]}\nobreak
\hyperdef{L}{X791690817E23D90C}{}
{\noindent\textcolor{FuncColor}{$\triangleright$\enspace\texttt{IsConstantTimeAccessGeneralMapping({\mdseries\slshape map})\index{IsConstantTimeAccessGeneralMapping@\texttt{IsConstantTimeAccessGeneralMapping}}
\label{IsConstantTimeAccessGeneralMapping}
}\hfill{\scriptsize (property)}}\\


 is \texttt{true} if the underlying relation of the general mapping \mbox{\texttt{\mdseries\slshape map}} knows its \texttt{AsList} (\ref{AsList}) value, and \texttt{false} otherwise. 

 In the former case, \mbox{\texttt{\mdseries\slshape map}} is allowed to use this list for calls to \texttt{ImagesElm} (\ref{ImagesElm}) etc. }

 

\subsection{\textcolor{Chapter }{IsEndoGeneralMapping}}
\logpage{[ 32, 13, 3 ]}\nobreak
\hyperdef{L}{X81CFF5F87BBEA8AD}{}
{\noindent\textcolor{FuncColor}{$\triangleright$\enspace\texttt{IsEndoGeneralMapping({\mdseries\slshape obj})\index{IsEndoGeneralMapping@\texttt{IsEndoGeneralMapping}}
\label{IsEndoGeneralMapping}
}\hfill{\scriptsize (property)}}\\


 If a general mapping has this property then its source and range are equal. }

 }

  
\section{\textcolor{Chapter }{Technical Matters Concerning General Mappings}}\label{Technical Matters Concerning General Mappings}
\logpage{[ 32, 14, 0 ]}
\hyperdef{L}{X7D6F78587C00CDD0}{}
{
  \texttt{Source} (\ref{Source}) and \texttt{Range} (\ref{Range:of a general mapping}) are basic operations for general mappings. \texttt{UnderlyingRelation} (\ref{UnderlyingRelation}) is secondary, its default method sets up a domain that delegates tasks to the
general mapping. (Note that this allows one to handle also infinite relations
by generic methods if source or range of the general mapping is finite.) 

 The distinction between basic operations and secondary operations for general
mappings may be a little bit complicated. Namely, each general mapping must be
in one of the two categories \texttt{IsNonSPGeneralMapping} (\ref{IsNonSPGeneralMapping}), \texttt{IsSPGeneralMapping} (\ref{IsSPGeneralMapping}). (The category \texttt{IsGeneralMapping} (\ref{IsGeneralMapping}) is defined as the disjoint union of these two.) 

 For general mappings of the first category, \texttt{ImagesElm} (\ref{ImagesElm}) and \texttt{PreImagesElm} (\ref{PreImagesElm}) are basic operations. (Note that in principle it is possible to delegate from \texttt{PreImagesElm} (\ref{PreImagesElm}) to \texttt{ImagesElm} (\ref{ImagesElm}).) Methods for the secondary operations \texttt{ImageElm} (\ref{ImageElm}), \texttt{PreImageElm} (\ref{PreImageElm}), \texttt{ImagesSet} (\ref{ImagesSet}), \texttt{PreImagesSet} (\ref{PreImagesSet}), \texttt{ImagesRepresentative} (\ref{ImagesRepresentative}), and \texttt{PreImagesRepresentative} (\ref{PreImagesRepresentative}) may use \texttt{ImagesElm} (\ref{ImagesElm}) and \texttt{PreImagesElm} (\ref{PreImagesElm}), respectively, and methods for \texttt{ImagesElm} (\ref{ImagesElm}), \texttt{PreImagesElm} (\ref{PreImagesElm}) must \emph{not} call the secondary operations. In particular, there are no generic methods for \texttt{ImagesElm} (\ref{ImagesElm}) and \texttt{PreImagesElm} (\ref{PreImagesElm}). 

 Methods for \texttt{ImagesSet} (\ref{ImagesSet}) and \texttt{PreImagesSet} (\ref{PreImagesSet}) must \emph{not} use \texttt{PreImagesRange} (\ref{PreImagesRange}) and \texttt{ImagesSource} (\ref{ImagesSource}), e.g., compute the intersection of the set in question with the preimage of
the range resp. the image of the source. 

 For general mappings of the second category (which means structure preserving
general mappings), the situation is different. The set of preimages under a
group homomorphism, for example, is either empty or can be described as a
coset of the (multiplicative) kernel. So it is reasonable to have \texttt{ImagesRepresentative} (\ref{ImagesRepresentative}), \texttt{PreImagesRepresentative} (\ref{PreImagesRepresentative}), \texttt{KernelOfMultiplicativeGeneralMapping} (\ref{KernelOfMultiplicativeGeneralMapping}), and \texttt{CoKernelOfMultiplicativeGeneralMapping} (\ref{CoKernelOfMultiplicativeGeneralMapping}) as basic operations here, and to make \texttt{ImagesElm} (\ref{ImagesElm}) and \texttt{PreImagesElm} (\ref{PreImagesElm}) secondary operations that may delegate to these. 

 In order to avoid infinite recursions, we must distinguish between the two
different types of mappings. 

 (Note that the basic domain operations such as \texttt{AsList} (\ref{AsList}) for the underlying relation of a general mapping may use either \texttt{ImagesElm} (\ref{ImagesElm}) or \texttt{ImagesRepresentative} (\ref{ImagesRepresentative}) and the appropriate cokernel. Conversely, if \texttt{AsList} (\ref{AsList}) for the underlying relation is known then \texttt{ImagesElm} (\ref{ImagesElm}) resp. \texttt{ImagesRepresentative} (\ref{ImagesRepresentative}) may delegate to it, the general mapping gets the property \texttt{IsConstantTimeAccessGeneralMapping} (\ref{IsConstantTimeAccessGeneralMapping}) for this; note that this is not allowed if only an enumerator of the
underlying relation is known.) 

 Secondary operations are \texttt{IsInjective} (\ref{IsInjective}), \texttt{IsSingleValued} (\ref{IsSingleValued}), \texttt{IsSurjective} (\ref{IsSurjective}), \texttt{IsTotal} (\ref{IsTotal}); they may use the basic operations, and must not be used by them. 

 Methods for the operations \texttt{ImagesElm} (\ref{ImagesElm}), \texttt{ImagesRepresentative} (\ref{ImagesRepresentative}), \texttt{ImagesSet} (\ref{ImagesSet}), \texttt{ImageElm} (\ref{ImageElm}), \texttt{PreImagesElm} (\ref{PreImagesElm}), \texttt{PreImagesRepresentative} (\ref{PreImagesRepresentative}), \texttt{PreImagesSet} (\ref{PreImagesSet}), and \texttt{PreImageElm} (\ref{PreImageElm}) take two arguments, a general mapping \mbox{\texttt{\mdseries\slshape map}} and an element or collection of elements \mbox{\texttt{\mdseries\slshape elm}}. These methods must \emph{not} check whether \mbox{\texttt{\mdseries\slshape elm}} lies in the source or the range of \mbox{\texttt{\mdseries\slshape map}}. In the case that \mbox{\texttt{\mdseries\slshape elm}} does not, \texttt{fail} may be returned as well as any other \textsf{GAP} object, and even an error message is allowed. Checks of the arguments are done
only by the functions \texttt{Image} (\ref{Image:set of images of the source of a general mapping}), \texttt{Images} (\ref{Images:set of images of the source of a general mapping}), \texttt{PreImage} (\ref{PreImage:set of preimages of the range of a general mapping}), and \texttt{PreImages} (\ref{PreImages:set of preimages of the range of a general mapping}), which then delegate to the operations listed above. 

\subsection{\textcolor{Chapter }{IsSPGeneralMapping}}
\logpage{[ 32, 14, 1 ]}\nobreak
\hyperdef{L}{X7D28581F82481163}{}
{\noindent\textcolor{FuncColor}{$\triangleright$\enspace\texttt{IsSPGeneralMapping({\mdseries\slshape map})\index{IsSPGeneralMapping@\texttt{IsSPGeneralMapping}}
\label{IsSPGeneralMapping}
}\hfill{\scriptsize (Category)}}\\
\noindent\textcolor{FuncColor}{$\triangleright$\enspace\texttt{IsNonSPGeneralMapping({\mdseries\slshape map})\index{IsNonSPGeneralMapping@\texttt{IsNonSPGeneralMapping}}
\label{IsNonSPGeneralMapping}
}\hfill{\scriptsize (Category)}}\\


   }

 

\subsection{\textcolor{Chapter }{IsGeneralMappingFamily}}
\logpage{[ 32, 14, 2 ]}\nobreak
\hyperdef{L}{X80D02AD183E01F16}{}
{\noindent\textcolor{FuncColor}{$\triangleright$\enspace\texttt{IsGeneralMappingFamily({\mdseries\slshape obj})\index{IsGeneralMappingFamily@\texttt{IsGeneralMappingFamily}}
\label{IsGeneralMappingFamily}
}\hfill{\scriptsize (Category)}}\\


 The family category of the category of general mappings. }

 

\subsection{\textcolor{Chapter }{FamilyRange}}
\logpage{[ 32, 14, 3 ]}\nobreak
\hyperdef{L}{X86CFADBA7F2FE446}{}
{\noindent\textcolor{FuncColor}{$\triangleright$\enspace\texttt{FamilyRange({\mdseries\slshape Fam})\index{FamilyRange@\texttt{FamilyRange}}
\label{FamilyRange}
}\hfill{\scriptsize (attribute)}}\\


 is the elements family of the family of the range of each general mapping in
the family \mbox{\texttt{\mdseries\slshape Fam}}. }

 

\subsection{\textcolor{Chapter }{FamilySource}}
\logpage{[ 32, 14, 4 ]}\nobreak
\hyperdef{L}{X7C3736E281A9E505}{}
{\noindent\textcolor{FuncColor}{$\triangleright$\enspace\texttt{FamilySource({\mdseries\slshape Fam})\index{FamilySource@\texttt{FamilySource}}
\label{FamilySource}
}\hfill{\scriptsize (attribute)}}\\


 is the elements family of the family of the source of each general mapping in
the family \mbox{\texttt{\mdseries\slshape Fam}}. }

 

\subsection{\textcolor{Chapter }{FamiliesOfGeneralMappingsAndRanges}}
\logpage{[ 32, 14, 5 ]}\nobreak
\hyperdef{L}{X7AE54FB67E2E6374}{}
{\noindent\textcolor{FuncColor}{$\triangleright$\enspace\texttt{FamiliesOfGeneralMappingsAndRanges({\mdseries\slshape Fam})\index{FamiliesOfGeneralMappingsAndRanges@\texttt{FamiliesOfGeneralMappingsAndRanges}}
\label{FamiliesOfGeneralMappingsAndRanges}
}\hfill{\scriptsize (attribute)}}\\


 is a list that stores at the odd positions the families of general mappings
with source in the family \mbox{\texttt{\mdseries\slshape Fam}}, at the even positions the families of ranges of the general mappings. }

 

\subsection{\textcolor{Chapter }{GeneralMappingsFamily}}
\logpage{[ 32, 14, 6 ]}\nobreak
\hyperdef{L}{X7E1E26E37C413F6F}{}
{\noindent\textcolor{FuncColor}{$\triangleright$\enspace\texttt{GeneralMappingsFamily({\mdseries\slshape sourcefam, rangefam})\index{GeneralMappingsFamily@\texttt{GeneralMappingsFamily}}
\label{GeneralMappingsFamily}
}\hfill{\scriptsize (function)}}\\


 All general mappings with same source family \mbox{\texttt{\mdseries\slshape FS}} and same range family \mbox{\texttt{\mdseries\slshape FR}} lie in the family \texttt{GeneralMappingsFamily( \mbox{\texttt{\mdseries\slshape FS}}, \mbox{\texttt{\mdseries\slshape FR}} )}. }

 

\subsection{\textcolor{Chapter }{TypeOfDefaultGeneralMapping}}
\logpage{[ 32, 14, 7 ]}\nobreak
\hyperdef{L}{X7CF92CC37A6BBDA5}{}
{\noindent\textcolor{FuncColor}{$\triangleright$\enspace\texttt{TypeOfDefaultGeneralMapping({\mdseries\slshape source, range, filter})\index{TypeOfDefaultGeneralMapping@\texttt{TypeOfDefaultGeneralMapping}}
\label{TypeOfDefaultGeneralMapping}
}\hfill{\scriptsize (function)}}\\


 is the type of mappings with \texttt{IsDefaultGeneralMappingRep} with source \mbox{\texttt{\mdseries\slshape source}} and range \mbox{\texttt{\mdseries\slshape range}} and additional categories \mbox{\texttt{\mdseries\slshape filter}}. }

 }

 }

            
\chapter{\textcolor{Chapter }{Relations}}\label{Relations}
\logpage{[ 33, 0, 0 ]}
\hyperdef{L}{X838651287FCCEFD8}{}
{
  \index{binary relation} \index{IsBinaryRelation@\texttt{IsBinaryRelation}!same as IsEndoGeneralMapping} \index{IsEndoGeneralMapping@\texttt{IsEndoGeneralMapping}!same as IsBinaryRelation} A \emph{binary relation} $R$ on a set $X$ is a subset of $X \times X$. A binary relation can also be thought of as a (general) mapping from $X$ to itself or as a directed graph where each edge represents an element of $R$. 

 In \textsf{GAP}, a relation is conceptually represented as a general mapping from $X$ to itself. The category \texttt{IsBinaryRelation} (\ref{IsBinaryRelation}) is a synonym for \texttt{IsEndoGeneralMapping} (\ref{IsEndoGeneralMapping}). Attributes and properties of relations in \textsf{GAP} are supported for relations, via considering relations as a subset of $X \times X$, or as a directed graph; examples include finding the strongly connected
components of a relation, via \texttt{StronglyConnectedComponents} (\ref{StronglyConnectedComponents}), or enumerating the tuples of the relation.  
\section{\textcolor{Chapter }{General Binary Relations}}\label{General Binary Relations}
\logpage{[ 33, 1, 0 ]}
\hyperdef{L}{X7DED7F1F78D31785}{}
{
  This section lists general constructors of relations. 

\subsection{\textcolor{Chapter }{IsBinaryRelation}}
\logpage{[ 33, 1, 1 ]}\nobreak
\hyperdef{L}{X788D722F82165551}{}
{\noindent\textcolor{FuncColor}{$\triangleright$\enspace\texttt{IsBinaryRelation({\mdseries\slshape R})\index{IsBinaryRelation@\texttt{IsBinaryRelation}}
\label{IsBinaryRelation}
}\hfill{\scriptsize (property)}}\\


 is exactly the same category as (i.e. a synonym for) \texttt{IsEndoGeneralMapping} (\ref{IsEndoGeneralMapping}). }

 

\subsection{\textcolor{Chapter }{BinaryRelationByElements}}
\logpage{[ 33, 1, 2 ]}\nobreak
\hyperdef{L}{X7A1D8EEF8034B0B5}{}
{\noindent\textcolor{FuncColor}{$\triangleright$\enspace\texttt{BinaryRelationByElements({\mdseries\slshape domain, elms})\index{BinaryRelationByElements@\texttt{BinaryRelationByElements}}
\label{BinaryRelationByElements}
}\hfill{\scriptsize (function)}}\\


 is the binary relation on \mbox{\texttt{\mdseries\slshape domain}} and with underlying relation consisting of the tuples collection \mbox{\texttt{\mdseries\slshape elms}}. This construction is similar to \texttt{GeneralMappingByElements} (\ref{GeneralMappingByElements}) where the source and range are the same set. 
\begin{Verbatim}[commandchars=!@|,fontsize=\small,frame=single,label=Example]
  !gapprompt@gap>| !gapinput@r:=BinaryRelationByElements(Domain([1..3]),[DirectProductElement([1,2]),DirectProductElement([1,3])]);|
  <general mapping: Domain([ 1 .. 3 ]) -> Domain([ 1 .. 3 ]) >
\end{Verbatim}
 }

 
\subsection{\textcolor{Chapter }{IdentityBinaryRelation}}\logpage{[ 33, 1, 3 ]}
\hyperdef{L}{X81878EEF873B34D5}{}
{
\noindent\textcolor{FuncColor}{$\triangleright$\enspace\texttt{IdentityBinaryRelation({\mdseries\slshape degree})\index{IdentityBinaryRelation@\texttt{IdentityBinaryRelation}!for a degree}
\label{IdentityBinaryRelation:for a degree}
}\hfill{\scriptsize (function)}}\\
\noindent\textcolor{FuncColor}{$\triangleright$\enspace\texttt{IdentityBinaryRelation({\mdseries\slshape domain})\index{IdentityBinaryRelation@\texttt{IdentityBinaryRelation}!for a domain}
\label{IdentityBinaryRelation:for a domain}
}\hfill{\scriptsize (function)}}\\


 is the binary relation which consists of diagonal pairs, i.e., pairs of the
form $(x,x)$. In the first form if a positive integer \mbox{\texttt{\mdseries\slshape degree}} is given then the domain is the set of the integers $\{ 1, \ldots, \mbox{\texttt{\mdseries\slshape degree}} \}$. In the second form, the objects $x$ are from the domain \mbox{\texttt{\mdseries\slshape domain}}. 
\begin{Verbatim}[commandchars=!@|,fontsize=\small,frame=single,label=Example]
  !gapprompt@gap>| !gapinput@IdentityBinaryRelation(5);|
  <equivalence relation on Domain([ 1 .. 5 ]) >
  !gapprompt@gap>| !gapinput@s4:=SymmetricGroup(4);|
  Sym( [ 1 .. 4 ] )
  !gapprompt@gap>| !gapinput@IdentityBinaryRelation(s4);|
  IdentityMapping( Sym( [ 1 .. 4 ] ) )
\end{Verbatim}
 }

 

\subsection{\textcolor{Chapter }{EmptyBinaryRelation (for a degree)}}
\logpage{[ 33, 1, 4 ]}\nobreak
\hyperdef{L}{X80DDCDD387BA23F2}{}
{\noindent\textcolor{FuncColor}{$\triangleright$\enspace\texttt{EmptyBinaryRelation({\mdseries\slshape degree})\index{EmptyBinaryRelation@\texttt{EmptyBinaryRelation}!for a degree}
\label{EmptyBinaryRelation:for a degree}
}\hfill{\scriptsize (function)}}\\
\noindent\textcolor{FuncColor}{$\triangleright$\enspace\texttt{EmptyBinaryRelation({\mdseries\slshape domain})\index{EmptyBinaryRelation@\texttt{EmptyBinaryRelation}!for a domain}
\label{EmptyBinaryRelation:for a domain}
}\hfill{\scriptsize (function)}}\\


 is the relation with \mbox{\texttt{\mdseries\slshape R}} empty. In the first form of the command with \mbox{\texttt{\mdseries\slshape degree}} an integer, the domain is the set of points $\{ 1, \ldots, \mbox{\texttt{\mdseries\slshape degree}} \}$. In the second form, the domain is that given by the argument \mbox{\texttt{\mdseries\slshape domain}}. 
\begin{Verbatim}[commandchars=!@|,fontsize=\small,frame=single,label=Example]
  !gapprompt@gap>| !gapinput@EmptyBinaryRelation(3) = BinaryRelationOnPoints([ [], [], [] ]);|
  true
\end{Verbatim}
 }

 }

  
\section{\textcolor{Chapter }{Properties and Attributes of Binary Relations}}\label{Properties and Attributes of Binary Relations}
\logpage{[ 33, 2, 0 ]}
\hyperdef{L}{X7899E59181C46EBB}{}
{
  

\subsection{\textcolor{Chapter }{IsReflexiveBinaryRelation}}
\logpage{[ 33, 2, 1 ]}\nobreak
\hyperdef{L}{X79D69B667F5FE8FE}{}
{\noindent\textcolor{FuncColor}{$\triangleright$\enspace\texttt{IsReflexiveBinaryRelation({\mdseries\slshape R})\index{IsReflexiveBinaryRelation@\texttt{IsReflexiveBinaryRelation}}
\label{IsReflexiveBinaryRelation}
}\hfill{\scriptsize (property)}}\\


 returns \texttt{true} if the binary relation \mbox{\texttt{\mdseries\slshape R}} is reflexive, and \texttt{false} otherwise. 

 \index{reflexive relation} A binary relation $R$ (as a set of pairs) on a set $X$ is \emph{reflexive} if for all $x \in X$, $(x,x) \in R$. Alternatively, $R$ as a mapping is reflexive if for all $x \in X$, $x$ is an element of the image set $R(x)$. 

 A reflexive binary relation is necessarily a total endomorphic mapping (tested
via \texttt{IsTotal} (\ref{IsTotal})). 
\begin{Verbatim}[commandchars=!@|,fontsize=\small,frame=single,label=Example]
  !gapprompt@gap>| !gapinput@IsReflexiveBinaryRelation(BinaryRelationOnPoints([[1,3],[2],[3]]));|
  true
  !gapprompt@gap>| !gapinput@IsReflexiveBinaryRelation(BinaryRelationOnPoints([[2],[2]]));|
  false
\end{Verbatim}
 }

 

\subsection{\textcolor{Chapter }{IsSymmetricBinaryRelation}}
\logpage{[ 33, 2, 2 ]}\nobreak
\hyperdef{L}{X785916A181555368}{}
{\noindent\textcolor{FuncColor}{$\triangleright$\enspace\texttt{IsSymmetricBinaryRelation({\mdseries\slshape R})\index{IsSymmetricBinaryRelation@\texttt{IsSymmetricBinaryRelation}}
\label{IsSymmetricBinaryRelation}
}\hfill{\scriptsize (property)}}\\


 returns \texttt{true} if the binary relation \mbox{\texttt{\mdseries\slshape R}} is symmetric, and \texttt{false} otherwise. 

 \index{symmetric relation} A binary relation $R$ (as a set of pairs) on a set $X$ is \emph{symmetric} if $(x,y) \in R$ then $(y,x) \in R$. Alternatively, $R$ as a mapping is symmetric if for all $x \in X$, the preimage set of $x$ under $R$ equals the image set $R(x)$. 
\begin{Verbatim}[commandchars=!@|,fontsize=\small,frame=single,label=Example]
  !gapprompt@gap>| !gapinput@IsSymmetricBinaryRelation(BinaryRelationOnPoints([[2],[1]]));|
  true
  !gapprompt@gap>| !gapinput@IsSymmetricBinaryRelation(BinaryRelationOnPoints([[2],[2]]));|
  false
\end{Verbatim}
 }

 

\subsection{\textcolor{Chapter }{IsTransitiveBinaryRelation}}
\logpage{[ 33, 2, 3 ]}\nobreak
\hyperdef{L}{X7823942478124563}{}
{\noindent\textcolor{FuncColor}{$\triangleright$\enspace\texttt{IsTransitiveBinaryRelation({\mdseries\slshape R})\index{IsTransitiveBinaryRelation@\texttt{IsTransitiveBinaryRelation}}
\label{IsTransitiveBinaryRelation}
}\hfill{\scriptsize (property)}}\\


 returns \texttt{true} if the binary relation \mbox{\texttt{\mdseries\slshape R}} is transitive, and \texttt{false} otherwise. 

 \index{transitive relation} A binary relation \mbox{\texttt{\mdseries\slshape R}} (as a set of pairs) on a set $X$ is \emph{transitive} if $(x,y), (y,z) \in R$ implies $(x,z) \in R$. Alternatively, $R$ as a mapping is transitive if for all $x \in X$, the image set $R(R(x))$ of the image set $R(x)$ of $x$ is a subset of $R(x)$. 
\begin{Verbatim}[commandchars=!@|,fontsize=\small,frame=single,label=Example]
  !gapprompt@gap>| !gapinput@IsTransitiveBinaryRelation(BinaryRelationOnPoints([[1,2,3],[2,3],[]]));|
  true
  !gapprompt@gap>| !gapinput@IsTransitiveBinaryRelation(BinaryRelationOnPoints([[1,2],[2,3],[]]));|
  false
\end{Verbatim}
 }

 

\subsection{\textcolor{Chapter }{IsAntisymmetricBinaryRelation}}
\logpage{[ 33, 2, 4 ]}\nobreak
\hyperdef{L}{X870F72C38550A0A4}{}
{\noindent\textcolor{FuncColor}{$\triangleright$\enspace\texttt{IsAntisymmetricBinaryRelation({\mdseries\slshape rel})\index{IsAntisymmetricBinaryRelation@\texttt{IsAntisymmetricBinaryRelation}}
\label{IsAntisymmetricBinaryRelation}
}\hfill{\scriptsize (property)}}\\


 returns \texttt{true} if the binary relation \mbox{\texttt{\mdseries\slshape rel}} is antisymmetric, and \texttt{false} otherwise. 

 \index{antisymmetric relation} A binary relation \mbox{\texttt{\mdseries\slshape R}} (as a set of pairs) on a set $X$ is \emph{antisymmetric} if $(x,y), (y,x) \in R$ implies $x = y$. Alternatively, $R$ as a mapping is antisymmetric if for all $x \in X$, the intersection of the preimage set of $x$ under $R$ and the image set $R(x)$ is $\{ x \}$. }

 

\subsection{\textcolor{Chapter }{IsPreOrderBinaryRelation}}
\logpage{[ 33, 2, 5 ]}\nobreak
\hyperdef{L}{X782B7C8A8136532F}{}
{\noindent\textcolor{FuncColor}{$\triangleright$\enspace\texttt{IsPreOrderBinaryRelation({\mdseries\slshape rel})\index{IsPreOrderBinaryRelation@\texttt{IsPreOrderBinaryRelation}}
\label{IsPreOrderBinaryRelation}
}\hfill{\scriptsize (property)}}\\


 returns \texttt{true} if the binary relation \mbox{\texttt{\mdseries\slshape rel}} is a preorder, and \texttt{false} otherwise. 

 \index{preorder} A \emph{preorder} is a binary relation that is both reflexive and transitive. }

 

\subsection{\textcolor{Chapter }{IsPartialOrderBinaryRelation}}
\logpage{[ 33, 2, 6 ]}\nobreak
\hyperdef{L}{X7A1228207AB4FBA3}{}
{\noindent\textcolor{FuncColor}{$\triangleright$\enspace\texttt{IsPartialOrderBinaryRelation({\mdseries\slshape rel})\index{IsPartialOrderBinaryRelation@\texttt{IsPartialOrderBinaryRelation}}
\label{IsPartialOrderBinaryRelation}
}\hfill{\scriptsize (property)}}\\


 returns \texttt{true} if the binary relation \mbox{\texttt{\mdseries\slshape rel}} is a partial order, and \texttt{false} otherwise. 

 \index{partial order} A \emph{partial order} is a preorder which is also antisymmetric. }

 

\subsection{\textcolor{Chapter }{IsHasseDiagram}}
\logpage{[ 33, 2, 7 ]}\nobreak
\hyperdef{L}{X80D3735C84D1CDC2}{}
{\noindent\textcolor{FuncColor}{$\triangleright$\enspace\texttt{IsHasseDiagram({\mdseries\slshape rel})\index{IsHasseDiagram@\texttt{IsHasseDiagram}}
\label{IsHasseDiagram}
}\hfill{\scriptsize (property)}}\\


 returns \texttt{true} if the binary relation \mbox{\texttt{\mdseries\slshape rel}} is a Hasse Diagram of a partial order, i.e., was computed via \texttt{HasseDiagramBinaryRelation} (\ref{HasseDiagramBinaryRelation}). }

 

\subsection{\textcolor{Chapter }{IsEquivalenceRelation}}
\logpage{[ 33, 2, 8 ]}\nobreak
\hyperdef{L}{X82D6CB4B7CCE9E25}{}
{\noindent\textcolor{FuncColor}{$\triangleright$\enspace\texttt{IsEquivalenceRelation({\mdseries\slshape R})\index{IsEquivalenceRelation@\texttt{IsEquivalenceRelation}}
\label{IsEquivalenceRelation}
}\hfill{\scriptsize (property)}}\\


 returns \texttt{true} if the binary relation \mbox{\texttt{\mdseries\slshape R}} is an equivalence relation, and \texttt{false} otherwise. 

 \index{equivalence relation} Recall, that a relation \mbox{\texttt{\mdseries\slshape R}} is an \emph{equivalence relation} if it is symmetric, transitive, and reflexive. }

 

\subsection{\textcolor{Chapter }{Successors}}
\logpage{[ 33, 2, 9 ]}\nobreak
\hyperdef{L}{X85E2FD8B82652876}{}
{\noindent\textcolor{FuncColor}{$\triangleright$\enspace\texttt{Successors({\mdseries\slshape R})\index{Successors@\texttt{Successors}}
\label{Successors}
}\hfill{\scriptsize (attribute)}}\\


 returns the list of images of a binary relation \mbox{\texttt{\mdseries\slshape R}}. If the underlying domain of the relation is not $\{ 1, \ldots, n \}$, for some positive integer $n$, then an error is signalled. 

 The returned value of \texttt{Successors} is a list of lists where the lists are ordered as the elements according to
the sorted order of the underlying set of \mbox{\texttt{\mdseries\slshape R}}. Each list consists of the images of the element whose index is the same as
the list with the underlying set in sorted order. 

 The \texttt{Successors} of a relation is the adjacency list representation of the relation. 
\begin{Verbatim}[commandchars=!@|,fontsize=\small,frame=single,label=Example]
  !gapprompt@gap>| !gapinput@r1:=BinaryRelationOnPoints([[2],[3],[1]]);;|
  !gapprompt@gap>| !gapinput@Successors(r1);|
  [ [ 2 ], [ 3 ], [ 1 ] ]
\end{Verbatim}
 }

 

\subsection{\textcolor{Chapter }{DegreeOfBinaryRelation}}
\logpage{[ 33, 2, 10 ]}\nobreak
\hyperdef{L}{X7B4D22A17E752A91}{}
{\noindent\textcolor{FuncColor}{$\triangleright$\enspace\texttt{DegreeOfBinaryRelation({\mdseries\slshape R})\index{DegreeOfBinaryRelation@\texttt{DegreeOfBinaryRelation}}
\label{DegreeOfBinaryRelation}
}\hfill{\scriptsize (attribute)}}\\


 returns the size of the underlying domain of the binary relation \mbox{\texttt{\mdseries\slshape R}}. This is most natural when working with a binary relation on points. 
\begin{Verbatim}[commandchars=!@|,fontsize=\small,frame=single,label=Example]
  !gapprompt@gap>| !gapinput@DegreeOfBinaryRelation(r1);|
  3
\end{Verbatim}
 }

 

\subsection{\textcolor{Chapter }{PartialOrderOfHasseDiagram}}
\logpage{[ 33, 2, 11 ]}\nobreak
\hyperdef{L}{X8278E4457C3C3A0D}{}
{\noindent\textcolor{FuncColor}{$\triangleright$\enspace\texttt{PartialOrderOfHasseDiagram({\mdseries\slshape HD})\index{PartialOrderOfHasseDiagram@\texttt{PartialOrderOfHasseDiagram}}
\label{PartialOrderOfHasseDiagram}
}\hfill{\scriptsize (attribute)}}\\


 is the partial order associated with the Hasse Diagram \mbox{\texttt{\mdseries\slshape HD}} i.e. the partial order generated by the reflexive and transitive closure of \mbox{\texttt{\mdseries\slshape HD}}. }

 }

  
\section{\textcolor{Chapter }{Binary Relations on Points}}\label{Binary Relations on Points}
\logpage{[ 33, 3, 0 ]}
\hyperdef{L}{X78032F927F078E19}{}
{
  We have special construction methods when the underlying \mbox{\texttt{\mdseries\slshape X}} of our relation is the set of integers $\{ 1, \ldots, n \}$. 

\subsection{\textcolor{Chapter }{BinaryRelationOnPoints}}
\logpage{[ 33, 3, 1 ]}\nobreak
\hyperdef{L}{X79E40E9385274F89}{}
{\noindent\textcolor{FuncColor}{$\triangleright$\enspace\texttt{BinaryRelationOnPoints({\mdseries\slshape list})\index{BinaryRelationOnPoints@\texttt{BinaryRelationOnPoints}}
\label{BinaryRelationOnPoints}
}\hfill{\scriptsize (function)}}\\
\noindent\textcolor{FuncColor}{$\triangleright$\enspace\texttt{BinaryRelationOnPointsNC({\mdseries\slshape list})\index{BinaryRelationOnPointsNC@\texttt{BinaryRelationOnPointsNC}}
\label{BinaryRelationOnPointsNC}
}\hfill{\scriptsize (function)}}\\


 Given a list of $n$ lists, each containing elements from the set $\{ 1, \ldots, n \}$, this function constructs a binary relation such that $1$ is related to \mbox{\texttt{\mdseries\slshape list}}\texttt{[1]}, $2$ to \mbox{\texttt{\mdseries\slshape list}}\texttt{[2]} and so on. The first version checks whether the list supplied is valid. The \texttt{NC} version skips this check. 
\begin{Verbatim}[commandchars=!@|,fontsize=\small,frame=single,label=Example]
  !gapprompt@gap>| !gapinput@R:=BinaryRelationOnPoints([[1,2],[2],[3]]);|
  Binary Relation on 3 points
\end{Verbatim}
 }

 

\subsection{\textcolor{Chapter }{RandomBinaryRelationOnPoints}}
\logpage{[ 33, 3, 2 ]}\nobreak
\hyperdef{L}{X7D9323C283867515}{}
{\noindent\textcolor{FuncColor}{$\triangleright$\enspace\texttt{RandomBinaryRelationOnPoints({\mdseries\slshape degree})\index{RandomBinaryRelationOnPoints@\texttt{RandomBinaryRelationOnPoints}}
\label{RandomBinaryRelationOnPoints}
}\hfill{\scriptsize (function)}}\\


 creates a relation on points with degree \mbox{\texttt{\mdseries\slshape degree}}. }

 
\subsection{\textcolor{Chapter }{AsBinaryRelationOnPoints}}\logpage{[ 33, 3, 3 ]}
\hyperdef{L}{X8315C7A47CEB6BB3}{}
{
\noindent\textcolor{FuncColor}{$\triangleright$\enspace\texttt{AsBinaryRelationOnPoints({\mdseries\slshape trans})\index{AsBinaryRelationOnPoints@\texttt{AsBinaryRelationOnPoints}!for a transformation}
\label{AsBinaryRelationOnPoints:for a transformation}
}\hfill{\scriptsize (function)}}\\
\noindent\textcolor{FuncColor}{$\triangleright$\enspace\texttt{AsBinaryRelationOnPoints({\mdseries\slshape perm})\index{AsBinaryRelationOnPoints@\texttt{AsBinaryRelationOnPoints}!for a permutation}
\label{AsBinaryRelationOnPoints:for a permutation}
}\hfill{\scriptsize (function)}}\\
\noindent\textcolor{FuncColor}{$\triangleright$\enspace\texttt{AsBinaryRelationOnPoints({\mdseries\slshape rel})\index{AsBinaryRelationOnPoints@\texttt{AsBinaryRelationOnPoints}!for a binary relation}
\label{AsBinaryRelationOnPoints:for a binary relation}
}\hfill{\scriptsize (function)}}\\


 return the relation on points represented by general relation \mbox{\texttt{\mdseries\slshape rel}}, transformation \mbox{\texttt{\mdseries\slshape trans}} or permutation \mbox{\texttt{\mdseries\slshape perm}}. If \mbox{\texttt{\mdseries\slshape rel}} is already a binary relation on points then \mbox{\texttt{\mdseries\slshape rel}} is returned. 

 Transformations and permutations are special general endomorphic mappings and
have a natural representation as a binary relation on points. 

 In the last form, an isomorphic relation on points is constructed where the
points are indices of the elements of the underlying domain in sorted order. 
\begin{Verbatim}[commandchars=!@|,fontsize=\small,frame=single,label=Example]
  !gapprompt@gap>| !gapinput@t:=Transformation([2,3,1]);;|
  !gapprompt@gap>| !gapinput@r1:=AsBinaryRelationOnPoints(t);|
  Binary Relation on 3 points
  !gapprompt@gap>| !gapinput@r2:=AsBinaryRelationOnPoints((1,2,3));|
  Binary Relation on 3 points
  !gapprompt@gap>| !gapinput@r1=r2;|
  true
\end{Verbatim}
 }

 }

  
\section{\textcolor{Chapter }{Closure Operations and Other Constructors}}\label{Closure Operations and Other Constructors}
\logpage{[ 33, 4, 0 ]}
\hyperdef{L}{X7D9A14AE799142EF}{}
{
  

\subsection{\textcolor{Chapter }{ReflexiveClosureBinaryRelation}}
\logpage{[ 33, 4, 1 ]}\nobreak
\hyperdef{L}{X8252B17C864A4904}{}
{\noindent\textcolor{FuncColor}{$\triangleright$\enspace\texttt{ReflexiveClosureBinaryRelation({\mdseries\slshape R})\index{ReflexiveClosureBinaryRelation@\texttt{ReflexiveClosureBinaryRelation}}
\label{ReflexiveClosureBinaryRelation}
}\hfill{\scriptsize (operation)}}\\


 is the smallest binary relation containing the binary relation \mbox{\texttt{\mdseries\slshape R}} which is reflexive. This closure inherits the properties symmetric and
transitive from \mbox{\texttt{\mdseries\slshape R}}. E.g., if \mbox{\texttt{\mdseries\slshape R}} is symmetric then its reflexive closure is also. }

 

\subsection{\textcolor{Chapter }{SymmetricClosureBinaryRelation}}
\logpage{[ 33, 4, 2 ]}\nobreak
\hyperdef{L}{X820811E9785A7274}{}
{\noindent\textcolor{FuncColor}{$\triangleright$\enspace\texttt{SymmetricClosureBinaryRelation({\mdseries\slshape R})\index{SymmetricClosureBinaryRelation@\texttt{SymmetricClosureBinaryRelation}}
\label{SymmetricClosureBinaryRelation}
}\hfill{\scriptsize (operation)}}\\


 is the smallest binary relation containing the binary relation \mbox{\texttt{\mdseries\slshape R}} which is symmetric. This closure inherits the properties reflexive and
transitive from \mbox{\texttt{\mdseries\slshape R}}. E.g., if \mbox{\texttt{\mdseries\slshape R}} is reflexive then its symmetric closure is also. }

 

\subsection{\textcolor{Chapter }{TransitiveClosureBinaryRelation}}
\logpage{[ 33, 4, 3 ]}\nobreak
\hyperdef{L}{X853BFAD9858DCDF7}{}
{\noindent\textcolor{FuncColor}{$\triangleright$\enspace\texttt{TransitiveClosureBinaryRelation({\mdseries\slshape rel})\index{TransitiveClosureBinaryRelation@\texttt{TransitiveClosureBinaryRelation}}
\label{TransitiveClosureBinaryRelation}
}\hfill{\scriptsize (operation)}}\\


 is the smallest binary relation containing the binary relation \mbox{\texttt{\mdseries\slshape R}} which is transitive. This closure inherits the properties reflexive and
symmetric from \mbox{\texttt{\mdseries\slshape R}}. E.g., if \mbox{\texttt{\mdseries\slshape R}} is symmetric then its transitive closure is also. 

 \texttt{TransitiveClosureBinaryRelation} is a modified version of the Floyd\texttt{\symbol{45}}Warshall method of
solving the all\texttt{\symbol{45}}pairs shortest\texttt{\symbol{45}}paths
problem on a directed graph. Its asymptotic runtime is $O(n^3)$ where $n$ is the size of the vertex set. It only assumes there is an arbitrary (but
fixed) ordering of the vertex set. }

 

\subsection{\textcolor{Chapter }{HasseDiagramBinaryRelation}}
\logpage{[ 33, 4, 4 ]}\nobreak
\hyperdef{L}{X79672B3A7BCB6991}{}
{\noindent\textcolor{FuncColor}{$\triangleright$\enspace\texttt{HasseDiagramBinaryRelation({\mdseries\slshape partial\texttt{\symbol{45}}order})\index{HasseDiagramBinaryRelation@\texttt{HasseDiagramBinaryRelation}}
\label{HasseDiagramBinaryRelation}
}\hfill{\scriptsize (operation)}}\\


 is the smallest relation contained in the partial order \mbox{\texttt{\mdseries\slshape partial\texttt{\symbol{45}}order}} whose reflexive and transitive closure is equal to \mbox{\texttt{\mdseries\slshape partial\texttt{\symbol{45}}order}}. }

 

\subsection{\textcolor{Chapter }{StronglyConnectedComponents}}
\logpage{[ 33, 4, 5 ]}\nobreak
\hyperdef{L}{X85C22B3D812957C0}{}
{\noindent\textcolor{FuncColor}{$\triangleright$\enspace\texttt{StronglyConnectedComponents({\mdseries\slshape R})\index{StronglyConnectedComponents@\texttt{StronglyConnectedComponents}}
\label{StronglyConnectedComponents}
}\hfill{\scriptsize (operation)}}\\


 returns an equivalence relation on the vertices of the binary relation \mbox{\texttt{\mdseries\slshape R}}. }

 

\subsection{\textcolor{Chapter }{PartialOrderByOrderingFunction}}
\logpage{[ 33, 4, 6 ]}\nobreak
\hyperdef{L}{X86AAE6027A3AEF72}{}
{\noindent\textcolor{FuncColor}{$\triangleright$\enspace\texttt{PartialOrderByOrderingFunction({\mdseries\slshape dom, orderfunc})\index{PartialOrderByOrderingFunction@\texttt{PartialOrderByOrderingFunction}}
\label{PartialOrderByOrderingFunction}
}\hfill{\scriptsize (function)}}\\


 constructs a partial order whose elements are from the domain \mbox{\texttt{\mdseries\slshape dom}} and are ordered using the ordering function \mbox{\texttt{\mdseries\slshape orderfunc}}. The ordering function must be a binary function returning a boolean value.
If the ordering function does not describe a partial order then \texttt{fail} is returned. }

 }

  
\section{\textcolor{Chapter }{Equivalence Relations}}\label{Equivalence Relations}
\logpage{[ 33, 5, 0 ]}
\hyperdef{L}{X7DAA67338458BB64}{}
{
  \index{equivalence relation} An \emph{equivalence relation} \mbox{\texttt{\mdseries\slshape E}} over the set \mbox{\texttt{\mdseries\slshape X}} is a relation on \mbox{\texttt{\mdseries\slshape X}} which is reflexive, symmetric, and transitive. A \emph{partition} \mbox{\texttt{\mdseries\slshape P}} is a set of subsets of \mbox{\texttt{\mdseries\slshape X}} such that for all $R, S \in P$, $R \cap S$ is the empty set and $\cup P = X$. An equivalence relation induces a partition such that if $(x,y) \in E$ then $x, y$ are in the same element of \mbox{\texttt{\mdseries\slshape P}}. 

 Like all binary relations in \textsf{GAP} equivalence relations are regarded as general endomorphic mappings (and the
operations, properties and attributes of general mappings are available).
However, partitions provide an efficient way of representing equivalence
relations. Moreover, only the non\texttt{\symbol{45}}singleton classes or
blocks are listed allowing for small equivalence relations to be represented
on infinite sets. Hence the main attribute of equivalence relations is \texttt{EquivalenceRelationPartition} (\ref{EquivalenceRelationPartition}) which provides the partition induced by the given equivalence. 

\subsection{\textcolor{Chapter }{EquivalenceRelationByPartition}}
\logpage{[ 33, 5, 1 ]}\nobreak
\hyperdef{L}{X7A44D73D8150266A}{}
{\noindent\textcolor{FuncColor}{$\triangleright$\enspace\texttt{EquivalenceRelationByPartition({\mdseries\slshape domain, list})\index{EquivalenceRelationByPartition@\texttt{EquivalenceRelationByPartition}}
\label{EquivalenceRelationByPartition}
}\hfill{\scriptsize (function)}}\\
\noindent\textcolor{FuncColor}{$\triangleright$\enspace\texttt{EquivalenceRelationByPartitionNC({\mdseries\slshape domain, list})\index{EquivalenceRelationByPartitionNC@\texttt{EquivalenceRelationByPartitionNC}}
\label{EquivalenceRelationByPartitionNC}
}\hfill{\scriptsize (function)}}\\


 constructs the equivalence relation over the set \mbox{\texttt{\mdseries\slshape domain}} which induces the partition represented by \mbox{\texttt{\mdseries\slshape list}}. This representation includes only the non\texttt{\symbol{45}}trivial blocks
(or equivalent classes). \mbox{\texttt{\mdseries\slshape list}} is a list of lists, each of these lists contain elements of \mbox{\texttt{\mdseries\slshape domain}} and are pairwise mutually exclusive. 

 The list of lists do not need to be in any order nor do the elements in the
blocks (see \texttt{EquivalenceRelationPartition} (\ref{EquivalenceRelationPartition})). a list of elements of \mbox{\texttt{\mdseries\slshape domain}} The partition \mbox{\texttt{\mdseries\slshape list}} is a list of lists, each of these is a list of elements of \mbox{\texttt{\mdseries\slshape domain}} that makes up a block (or equivalent class). The \mbox{\texttt{\mdseries\slshape domain}} is the domain over which the relation is defined, and \mbox{\texttt{\mdseries\slshape list}} is a list of lists, each of these is a list of elements of \mbox{\texttt{\mdseries\slshape domain}} which are related to each other. \mbox{\texttt{\mdseries\slshape list}} need only contain the nontrivial blocks and singletons will be ignored. The \texttt{NC} version will not check to see if the lists are pairwise mutually exclusive or
that they contain only elements of the domain. 
\begin{Verbatim}[commandchars=!@|,fontsize=\small,frame=single,label=Example]
  !gapprompt@gap>| !gapinput@er:=EquivalenceRelationByPartition(Domain([1..10]),[[1,3,5,7,9],[2,4,6,8,10]]);|
  <equivalence relation on Domain([ 1 .. 10 ]) >
  !gapprompt@gap>| !gapinput@IsEquivalenceRelation(er);|
  true
\end{Verbatim}
 }

 

\subsection{\textcolor{Chapter }{EquivalenceRelationByRelation}}
\logpage{[ 33, 5, 2 ]}\nobreak
\hyperdef{L}{X82CD1C00810F6042}{}
{\noindent\textcolor{FuncColor}{$\triangleright$\enspace\texttt{EquivalenceRelationByRelation({\mdseries\slshape rel})\index{EquivalenceRelationByRelation@\texttt{EquivalenceRelationByRelation}}
\label{EquivalenceRelationByRelation}
}\hfill{\scriptsize (function)}}\\


 returns the smallest equivalence relation containing the binary relation \mbox{\texttt{\mdseries\slshape rel}}. }

 

\subsection{\textcolor{Chapter }{EquivalenceRelationByPairs}}
\logpage{[ 33, 5, 3 ]}\nobreak
\hyperdef{L}{X7B70215E7E3F9CA4}{}
{\noindent\textcolor{FuncColor}{$\triangleright$\enspace\texttt{EquivalenceRelationByPairs({\mdseries\slshape D, elms})\index{EquivalenceRelationByPairs@\texttt{EquivalenceRelationByPairs}}
\label{EquivalenceRelationByPairs}
}\hfill{\scriptsize (function)}}\\
\noindent\textcolor{FuncColor}{$\triangleright$\enspace\texttt{EquivalenceRelationByPairsNC({\mdseries\slshape D, elms})\index{EquivalenceRelationByPairsNC@\texttt{EquivalenceRelationByPairsNC}}
\label{EquivalenceRelationByPairsNC}
}\hfill{\scriptsize (function)}}\\


 return the smallest equivalence relation on the domain \mbox{\texttt{\mdseries\slshape D}} such that every pair in \mbox{\texttt{\mdseries\slshape elms}} is in the relation. 

 In the \texttt{NC} form, it is not checked that \mbox{\texttt{\mdseries\slshape elms}} are in the domain \mbox{\texttt{\mdseries\slshape D}}. }

 

\subsection{\textcolor{Chapter }{EquivalenceRelationByProperty}}
\logpage{[ 33, 5, 4 ]}\nobreak
\hyperdef{L}{X7C5AA9B97EE42DA5}{}
{\noindent\textcolor{FuncColor}{$\triangleright$\enspace\texttt{EquivalenceRelationByProperty({\mdseries\slshape domain, property})\index{EquivalenceRelationByProperty@\texttt{EquivalenceRelationByProperty}}
\label{EquivalenceRelationByProperty}
}\hfill{\scriptsize (function)}}\\


 creates an equivalence relation on \mbox{\texttt{\mdseries\slshape domain}} whose only defining datum is that of having the property \mbox{\texttt{\mdseries\slshape property}}. }

 }

  
\section{\textcolor{Chapter }{Attributes of and Operations on Equivalence Relations}}\label{Attributes of and Operations on Equivalence Relations}
\logpage{[ 33, 6, 0 ]}
\hyperdef{L}{X85A2A8E27AF52769}{}
{
  

\subsection{\textcolor{Chapter }{EquivalenceRelationPartition}}
\logpage{[ 33, 6, 1 ]}\nobreak
\hyperdef{L}{X877389B683DD8F1A}{}
{\noindent\textcolor{FuncColor}{$\triangleright$\enspace\texttt{EquivalenceRelationPartition({\mdseries\slshape equiv})\index{EquivalenceRelationPartition@\texttt{EquivalenceRelationPartition}}
\label{EquivalenceRelationPartition}
}\hfill{\scriptsize (attribute)}}\\


 returns a list of lists of elements of the underlying set of the equivalence
relation \mbox{\texttt{\mdseries\slshape equiv}}. The lists are precisely the nonsingleton equivalence classes of the
equivalence. This allows us to describe ``small'' equivalences on infinite sets. }

 

\subsection{\textcolor{Chapter }{GeneratorsOfEquivalenceRelationPartition}}
\logpage{[ 33, 6, 2 ]}\nobreak
\hyperdef{L}{X79DC914C82D7903B}{}
{\noindent\textcolor{FuncColor}{$\triangleright$\enspace\texttt{GeneratorsOfEquivalenceRelationPartition({\mdseries\slshape equiv})\index{GeneratorsOfEquivalenceRelationPartition@\texttt{Generators}\-\texttt{Of}\-\texttt{Equivalence}\-\texttt{Relation}\-\texttt{Partition}}
\label{GeneratorsOfEquivalenceRelationPartition}
}\hfill{\scriptsize (attribute)}}\\


 is a set of generating pairs for the equivalence relation \mbox{\texttt{\mdseries\slshape equiv}}. This set is not unique. The equivalence \mbox{\texttt{\mdseries\slshape equiv}} is the smallest equivalence relation over the underlying set which contains
the generating pairs. }

 

\subsection{\textcolor{Chapter }{JoinEquivalenceRelations}}
\logpage{[ 33, 6, 3 ]}\nobreak
\hyperdef{L}{X82BE360381476D92}{}
{\noindent\textcolor{FuncColor}{$\triangleright$\enspace\texttt{JoinEquivalenceRelations({\mdseries\slshape equiv1, equiv2})\index{JoinEquivalenceRelations@\texttt{JoinEquivalenceRelations}}
\label{JoinEquivalenceRelations}
}\hfill{\scriptsize (operation)}}\\
\noindent\textcolor{FuncColor}{$\triangleright$\enspace\texttt{MeetEquivalenceRelations({\mdseries\slshape equiv1, equiv2})\index{MeetEquivalenceRelations@\texttt{MeetEquivalenceRelations}}
\label{MeetEquivalenceRelations}
}\hfill{\scriptsize (operation)}}\\


 \texttt{JoinEquivalenceRelations} returns the smallest equivalence relation containing both the equivalence
relations \mbox{\texttt{\mdseries\slshape equiv1}} and \mbox{\texttt{\mdseries\slshape equiv2}}. 

 \texttt{MeetEquivalenceRelations} returns the intersection of the two equivalence relations \mbox{\texttt{\mdseries\slshape equiv1}} and \mbox{\texttt{\mdseries\slshape equiv2}}. }

 }

  
\section{\textcolor{Chapter }{Equivalence Classes}}\label{Equivalence Classes}
\logpage{[ 33, 7, 0 ]}
\hyperdef{L}{X79EE13287DEB11B1}{}
{
  

\subsection{\textcolor{Chapter }{IsEquivalenceClass}}
\logpage{[ 33, 7, 1 ]}\nobreak
\hyperdef{L}{X8424996186DB14EA}{}
{\noindent\textcolor{FuncColor}{$\triangleright$\enspace\texttt{IsEquivalenceClass({\mdseries\slshape obj})\index{IsEquivalenceClass@\texttt{IsEquivalenceClass}}
\label{IsEquivalenceClass}
}\hfill{\scriptsize (Category)}}\\


 returns \texttt{true} if the object \mbox{\texttt{\mdseries\slshape obj}} is an equivalence class, and \texttt{false} otherwise. 

 \index{equivalence class} An \emph{equivalence class} is a collection of elements which are mutually related to each other in the
associated equivalence relation. Note, this is a special category of objects
and not just a list of elements. }

 

\subsection{\textcolor{Chapter }{EquivalenceClassRelation}}
\logpage{[ 33, 7, 2 ]}\nobreak
\hyperdef{L}{X78F967E77EB16386}{}
{\noindent\textcolor{FuncColor}{$\triangleright$\enspace\texttt{EquivalenceClassRelation({\mdseries\slshape C})\index{EquivalenceClassRelation@\texttt{EquivalenceClassRelation}}
\label{EquivalenceClassRelation}
}\hfill{\scriptsize (attribute)}}\\


 returns the equivalence relation of which \mbox{\texttt{\mdseries\slshape C}} is a class. }

 

\subsection{\textcolor{Chapter }{EquivalenceClasses (attribute)}}
\logpage{[ 33, 7, 3 ]}\nobreak
\hyperdef{L}{X879439897EF4D728}{}
{\noindent\textcolor{FuncColor}{$\triangleright$\enspace\texttt{EquivalenceClasses({\mdseries\slshape rel})\index{EquivalenceClasses@\texttt{EquivalenceClasses}!attribute}
\label{EquivalenceClasses:attribute}
}\hfill{\scriptsize (attribute)}}\\


 returns a list of all equivalence classes of the equivalence relation \mbox{\texttt{\mdseries\slshape rel}}. Note that it is possible for different methods to yield the list in
different orders, so that for two equivalence relations $c1$ and $c2$ we may have $c1 = c2$ without having \texttt{EquivalenceClasses}$( c1 ) = $\texttt{EquivalenceClasses}$( c2 )$. 
\begin{Verbatim}[commandchars=!@|,fontsize=\small,frame=single,label=Example]
  !gapprompt@gap>| !gapinput@er:=EquivalenceRelationByPartition(Domain([1..10]),[[1,3,5,7,9],[2,4,6,8,10]]);|
  <equivalence relation on Domain([ 1 .. 10 ]) >
  !gapprompt@gap>| !gapinput@classes := EquivalenceClasses(er);|
  [ {1}, {2} ]
\end{Verbatim}
 }

 

\subsection{\textcolor{Chapter }{EquivalenceClassOfElement}}
\logpage{[ 33, 7, 4 ]}\nobreak
\hyperdef{L}{X7BB985BA7FD7A82E}{}
{\noindent\textcolor{FuncColor}{$\triangleright$\enspace\texttt{EquivalenceClassOfElement({\mdseries\slshape rel, elt})\index{EquivalenceClassOfElement@\texttt{EquivalenceClassOfElement}}
\label{EquivalenceClassOfElement}
}\hfill{\scriptsize (operation)}}\\
\noindent\textcolor{FuncColor}{$\triangleright$\enspace\texttt{EquivalenceClassOfElementNC({\mdseries\slshape rel, elt})\index{EquivalenceClassOfElementNC@\texttt{EquivalenceClassOfElementNC}}
\label{EquivalenceClassOfElementNC}
}\hfill{\scriptsize (operation)}}\\


 return the equivalence class of \mbox{\texttt{\mdseries\slshape elt}} in the binary relation \mbox{\texttt{\mdseries\slshape rel}}, where \mbox{\texttt{\mdseries\slshape elt}} is an element (i.e. a pair) of the domain of \mbox{\texttt{\mdseries\slshape rel}}. In the \texttt{NC} form, it is not checked that \mbox{\texttt{\mdseries\slshape elt}} is in the domain over which \mbox{\texttt{\mdseries\slshape rel}} is defined. 
\begin{Verbatim}[commandchars=!@|,fontsize=\small,frame=single,label=Example]
  !gapprompt@gap>| !gapinput@EquivalenceClassOfElement(er,3);|
  {3}
\end{Verbatim}
 }

 }

 }

            
\chapter{\textcolor{Chapter }{Orderings}}\label{Orderings}
\logpage{[ 34, 0, 0 ]}
\hyperdef{L}{X7E4AAA7382D42361}{}
{
  In \textsf{GAP} an ordering is a relation defined on a family, which is reflexive,
anti\texttt{\symbol{45}}symmetric and transitive.  
\section{\textcolor{Chapter }{IsOrdering (Filter)}}\label{sect:IsOrdering}
\logpage{[ 34, 1, 0 ]}
\hyperdef{L}{X79B1262585CE5427}{}
{
  

\subsection{\textcolor{Chapter }{IsOrdering}}
\logpage{[ 34, 1, 1 ]}\nobreak
\hyperdef{L}{X7EFDF115780934AF}{}
{\noindent\textcolor{FuncColor}{$\triangleright$\enspace\texttt{IsOrdering({\mdseries\slshape obj})\index{IsOrdering@\texttt{IsOrdering}}
\label{IsOrdering}
}\hfill{\scriptsize (Category)}}\\


 returns \texttt{true} if and only if the object \mbox{\texttt{\mdseries\slshape ord}} is an ordering. }

 

\subsection{\textcolor{Chapter }{OrderingsFamily}}
\logpage{[ 34, 1, 2 ]}\nobreak
\hyperdef{L}{X85E6445C87283BEC}{}
{\noindent\textcolor{FuncColor}{$\triangleright$\enspace\texttt{OrderingsFamily({\mdseries\slshape fam})\index{OrderingsFamily@\texttt{OrderingsFamily}}
\label{OrderingsFamily}
}\hfill{\scriptsize (attribute)}}\\


 for a family \mbox{\texttt{\mdseries\slshape fam}}, returns the family of all orderings on elements of \mbox{\texttt{\mdseries\slshape fam}}. }

 }

  
\section{\textcolor{Chapter }{Building new orderings}}\label{Building new orderings}
\logpage{[ 34, 2, 0 ]}
\hyperdef{L}{X85C4CAA784BD7F01}{}
{
  

\subsection{\textcolor{Chapter }{OrderingByLessThanFunctionNC}}
\logpage{[ 34, 2, 1 ]}\nobreak
\hyperdef{L}{X78B5D91278EFAFC9}{}
{\noindent\textcolor{FuncColor}{$\triangleright$\enspace\texttt{OrderingByLessThanFunctionNC({\mdseries\slshape fam, lt[, l]})\index{OrderingByLessThanFunctionNC@\texttt{OrderingByLessThanFunctionNC}}
\label{OrderingByLessThanFunctionNC}
}\hfill{\scriptsize (operation)}}\\


 Called with two arguments, \texttt{OrderingByLessThanFunctionNC} returns the ordering on the elements of the elements of the family \mbox{\texttt{\mdseries\slshape fam}}, according to the \texttt{LessThanFunction} (\ref{LessThanFunction}) value given by \mbox{\texttt{\mdseries\slshape lt}}, where \mbox{\texttt{\mdseries\slshape lt}} is a function that takes two arguments in \mbox{\texttt{\mdseries\slshape fam}} and returns \texttt{true} or \texttt{false}. 

 Called with three arguments, for a family \mbox{\texttt{\mdseries\slshape fam}}, a function \mbox{\texttt{\mdseries\slshape lt}} that takes two arguments in \mbox{\texttt{\mdseries\slshape fam}} and returns \texttt{true} or \texttt{false}, and a list \mbox{\texttt{\mdseries\slshape l}} of properties of orderings, \texttt{OrderingByLessThanFunctionNC} returns the ordering on the elements of \mbox{\texttt{\mdseries\slshape fam}} with \texttt{LessThanFunction} (\ref{LessThanFunction}) value given by \mbox{\texttt{\mdseries\slshape lt}} and with the properties from \mbox{\texttt{\mdseries\slshape l}} set to \texttt{true}. }

 

\subsection{\textcolor{Chapter }{OrderingByLessThanOrEqualFunctionNC}}
\logpage{[ 34, 2, 2 ]}\nobreak
\hyperdef{L}{X813D5BEB80506CE4}{}
{\noindent\textcolor{FuncColor}{$\triangleright$\enspace\texttt{OrderingByLessThanOrEqualFunctionNC({\mdseries\slshape fam, lteq[, l]})\index{OrderingByLessThanOrEqualFunctionNC@\texttt{OrderingByLessThanOrEqualFunctionNC}}
\label{OrderingByLessThanOrEqualFunctionNC}
}\hfill{\scriptsize (operation)}}\\


 Called with two arguments, \texttt{OrderingByLessThanOrEqualFunctionNC} returns the ordering on the elements of the elements of the family \mbox{\texttt{\mdseries\slshape fam}} according to the \texttt{LessThanOrEqualFunction} (\ref{LessThanOrEqualFunction}) value given by \mbox{\texttt{\mdseries\slshape lteq}}, where \mbox{\texttt{\mdseries\slshape lteq}} is a function that takes two arguments in \mbox{\texttt{\mdseries\slshape fam}} and returns \texttt{true} or \texttt{false}. 

 Called with three arguments, for a family \mbox{\texttt{\mdseries\slshape fam}}, a function \mbox{\texttt{\mdseries\slshape lteq}} that takes two arguments in \mbox{\texttt{\mdseries\slshape fam}} and returns \texttt{true} or \texttt{false}, and a list \mbox{\texttt{\mdseries\slshape l}} of properties of orderings, \texttt{OrderingByLessThanOrEqualFunctionNC} returns the ordering on the elements of \mbox{\texttt{\mdseries\slshape fam}} with \texttt{LessThanOrEqualFunction} (\ref{LessThanOrEqualFunction}) value given by \mbox{\texttt{\mdseries\slshape lteq}} and with the properties from \mbox{\texttt{\mdseries\slshape l}} set to \texttt{true}. 

 Notice that these functions do not check whether \mbox{\texttt{\mdseries\slshape fam}} and \mbox{\texttt{\mdseries\slshape lt}} or \mbox{\texttt{\mdseries\slshape lteq}} are compatible, and whether the properties listed in \mbox{\texttt{\mdseries\slshape l}} are indeed satisfied. 
\begin{Verbatim}[commandchars=!@|,fontsize=\small,frame=single,label=Example]
  !gapprompt@gap>| !gapinput@f := FreeSemigroup("a","b");;|
  !gapprompt@gap>| !gapinput@a := GeneratorsOfSemigroup(f)[1];;|
  !gapprompt@gap>| !gapinput@b := GeneratorsOfSemigroup(f)[2];;|
  !gapprompt@gap>| !gapinput@lt := function(x,y) return Length(x)<Length(y); end;|
  function( x, y ) ... end
  !gapprompt@gap>| !gapinput@fam := FamilyObj(a);;|
  !gapprompt@gap>| !gapinput@ord := OrderingByLessThanFunctionNC(fam,lt);|
  Ordering
\end{Verbatim}
 }

 }

  
\section{\textcolor{Chapter }{Properties and basic functionality}}\label{Properties and basic functionality}
\logpage{[ 34, 3, 0 ]}
\hyperdef{L}{X7F62235B87C20A54}{}
{
  

\subsection{\textcolor{Chapter }{IsWellFoundedOrdering}}
\logpage{[ 34, 3, 1 ]}\nobreak
\hyperdef{L}{X84FA448B7B4DDFDC}{}
{\noindent\textcolor{FuncColor}{$\triangleright$\enspace\texttt{IsWellFoundedOrdering({\mdseries\slshape ord})\index{IsWellFoundedOrdering@\texttt{IsWellFoundedOrdering}}
\label{IsWellFoundedOrdering}
}\hfill{\scriptsize (property)}}\\


 for an ordering \mbox{\texttt{\mdseries\slshape ord}}, returns \texttt{true} if and only if the ordering is well founded. An ordering \mbox{\texttt{\mdseries\slshape ord}} is well founded if it admits no infinite descending chains. Normally this
property is set at the time of creation of the ordering and there is no
general method to check whether a certain ordering is well founded. }

 

\subsection{\textcolor{Chapter }{IsTotalOrdering}}
\logpage{[ 34, 3, 2 ]}\nobreak
\hyperdef{L}{X867AC932843AD921}{}
{\noindent\textcolor{FuncColor}{$\triangleright$\enspace\texttt{IsTotalOrdering({\mdseries\slshape ord})\index{IsTotalOrdering@\texttt{IsTotalOrdering}}
\label{IsTotalOrdering}
}\hfill{\scriptsize (property)}}\\


 for an ordering \mbox{\texttt{\mdseries\slshape ord}}, returns true if and only if the ordering is total. An ordering \mbox{\texttt{\mdseries\slshape ord}} is total if any two elements of the family are comparable under \mbox{\texttt{\mdseries\slshape ord}}. Normally this property is set at the time of creation of the ordering and
there is no general method to check whether a certain ordering is total. }

 

\subsection{\textcolor{Chapter }{IsIncomparableUnder}}
\logpage{[ 34, 3, 3 ]}\nobreak
\hyperdef{L}{X814E5E7D85EDCAC7}{}
{\noindent\textcolor{FuncColor}{$\triangleright$\enspace\texttt{IsIncomparableUnder({\mdseries\slshape ord, el1, el2})\index{IsIncomparableUnder@\texttt{IsIncomparableUnder}}
\label{IsIncomparableUnder}
}\hfill{\scriptsize (operation)}}\\


 for an ordering \mbox{\texttt{\mdseries\slshape ord}} on the elements of the family of \mbox{\texttt{\mdseries\slshape el1}} and \mbox{\texttt{\mdseries\slshape el2}}, returns \texttt{true} if \mbox{\texttt{\mdseries\slshape el1}} $\neq$ \mbox{\texttt{\mdseries\slshape el2}} and \texttt{IsLessThanUnder}(\mbox{\texttt{\mdseries\slshape ord}},\mbox{\texttt{\mdseries\slshape el1}},\mbox{\texttt{\mdseries\slshape el2}}), \texttt{IsLessThanUnder}(\mbox{\texttt{\mdseries\slshape ord}},\mbox{\texttt{\mdseries\slshape el2}},\mbox{\texttt{\mdseries\slshape el1}}) are both \texttt{false}; and returns \texttt{false} otherwise. }

 

\subsection{\textcolor{Chapter }{FamilyForOrdering}}
\logpage{[ 34, 3, 4 ]}\nobreak
\hyperdef{L}{X872497B9782B97B4}{}
{\noindent\textcolor{FuncColor}{$\triangleright$\enspace\texttt{FamilyForOrdering({\mdseries\slshape ord})\index{FamilyForOrdering@\texttt{FamilyForOrdering}}
\label{FamilyForOrdering}
}\hfill{\scriptsize (attribute)}}\\


 for an ordering \mbox{\texttt{\mdseries\slshape ord}}, returns the family of elements that the ordering \mbox{\texttt{\mdseries\slshape ord}} compares. }

 

\subsection{\textcolor{Chapter }{LessThanFunction}}
\logpage{[ 34, 3, 5 ]}\nobreak
\hyperdef{L}{X7D08ED6882015BFB}{}
{\noindent\textcolor{FuncColor}{$\triangleright$\enspace\texttt{LessThanFunction({\mdseries\slshape ord})\index{LessThanFunction@\texttt{LessThanFunction}}
\label{LessThanFunction}
}\hfill{\scriptsize (attribute)}}\\


 for an ordering \mbox{\texttt{\mdseries\slshape ord}}, returns a function $f$ which takes two elements $el1$, $el2$ in \texttt{FamilyForOrdering}(\mbox{\texttt{\mdseries\slshape ord}}) and returns \texttt{true} if $el1$ is strictly less than $el2$ (with respect to \mbox{\texttt{\mdseries\slshape ord}}), and returns \texttt{false} otherwise. }

 

\subsection{\textcolor{Chapter }{LessThanOrEqualFunction}}
\logpage{[ 34, 3, 6 ]}\nobreak
\hyperdef{L}{X857E800583E9026D}{}
{\noindent\textcolor{FuncColor}{$\triangleright$\enspace\texttt{LessThanOrEqualFunction({\mdseries\slshape ord})\index{LessThanOrEqualFunction@\texttt{LessThanOrEqualFunction}}
\label{LessThanOrEqualFunction}
}\hfill{\scriptsize (attribute)}}\\


 for an ordering \mbox{\texttt{\mdseries\slshape ord}}, returns a function that takes two elements $el1$, $el2$ in \texttt{FamilyForOrdering}(\mbox{\texttt{\mdseries\slshape ord}}) and returns \texttt{true} if $el1$ is less than \emph{or equal to} $el2$ (with respect to \mbox{\texttt{\mdseries\slshape ord}}), and returns \texttt{false} otherwise. }

 

\subsection{\textcolor{Chapter }{IsLessThanUnder}}
\logpage{[ 34, 3, 7 ]}\nobreak
\hyperdef{L}{X87F51D737C695D41}{}
{\noindent\textcolor{FuncColor}{$\triangleright$\enspace\texttt{IsLessThanUnder({\mdseries\slshape ord, el1, el2})\index{IsLessThanUnder@\texttt{IsLessThanUnder}}
\label{IsLessThanUnder}
}\hfill{\scriptsize (operation)}}\\


 for an ordering \mbox{\texttt{\mdseries\slshape ord}} on the elements of the family of \mbox{\texttt{\mdseries\slshape el1}} and \mbox{\texttt{\mdseries\slshape el2}}, returns \texttt{true} if \mbox{\texttt{\mdseries\slshape el1}} is (strictly) less than \mbox{\texttt{\mdseries\slshape el2}} with respect to \mbox{\texttt{\mdseries\slshape ord}}, and \texttt{false} otherwise. }

 

\subsection{\textcolor{Chapter }{IsLessThanOrEqualUnder}}
\logpage{[ 34, 3, 8 ]}\nobreak
\hyperdef{L}{X8308B7DF7AAF6D9C}{}
{\noindent\textcolor{FuncColor}{$\triangleright$\enspace\texttt{IsLessThanOrEqualUnder({\mdseries\slshape ord, el1, el2})\index{IsLessThanOrEqualUnder@\texttt{IsLessThanOrEqualUnder}}
\label{IsLessThanOrEqualUnder}
}\hfill{\scriptsize (operation)}}\\


 for an ordering \mbox{\texttt{\mdseries\slshape ord}} on the elements of the family of \mbox{\texttt{\mdseries\slshape el1}} and \mbox{\texttt{\mdseries\slshape el2}}, returns \texttt{true} if \mbox{\texttt{\mdseries\slshape el1}} is less than or equal to \mbox{\texttt{\mdseries\slshape el2}} with respect to \mbox{\texttt{\mdseries\slshape ord}}, and \texttt{false} otherwise. 
\begin{Verbatim}[commandchars=!@|,fontsize=\small,frame=single,label=Example]
  !gapprompt@gap>| !gapinput@IsLessThanUnder(ord,a,a*b);|
  true
  !gapprompt@gap>| !gapinput@IsLessThanOrEqualUnder(ord,a*b,a*b);|
  true
  !gapprompt@gap>| !gapinput@IsIncomparableUnder(ord,a,b);|
  true
  !gapprompt@gap>| !gapinput@FamilyForOrdering(ord) = FamilyObj(a);|
  true
\end{Verbatim}
 }

 }

  
\section{\textcolor{Chapter }{Orderings on families of associative words}}\label{Orderings on families of associative words}
\logpage{[ 34, 4, 0 ]}
\hyperdef{L}{X834CD021878745BC}{}
{
  We now consider orderings on families of associative words. 

 Examples of families of associative words are the families of elements of a
free semigroup or a free monoid; these are the two cases that we consider
mostly. Associated with those families is an alphabet, which is the semigroup
(resp. monoid) generating set of the correspondent free semigroup (resp. free
monoid). For definitions of the orderings considered, see Sims \cite{Sims94}. 

\subsection{\textcolor{Chapter }{IsOrderingOnFamilyOfAssocWords}}
\logpage{[ 34, 4, 1 ]}\nobreak
\hyperdef{L}{X7C1808AE84B989AE}{}
{\noindent\textcolor{FuncColor}{$\triangleright$\enspace\texttt{IsOrderingOnFamilyOfAssocWords({\mdseries\slshape ord})\index{IsOrderingOnFamilyOfAssocWords@\texttt{IsOrderingOnFamilyOfAssocWords}}
\label{IsOrderingOnFamilyOfAssocWords}
}\hfill{\scriptsize (property)}}\\


 for an ordering \mbox{\texttt{\mdseries\slshape ord}}, returns true if \mbox{\texttt{\mdseries\slshape ord}} is an ordering over a family of associative words. }

 

\subsection{\textcolor{Chapter }{IsTranslationInvariantOrdering}}
\logpage{[ 34, 4, 2 ]}\nobreak
\hyperdef{L}{X8175B8887868F29A}{}
{\noindent\textcolor{FuncColor}{$\triangleright$\enspace\texttt{IsTranslationInvariantOrdering({\mdseries\slshape ord})\index{IsTranslationInvariantOrdering@\texttt{IsTranslationInvariantOrdering}}
\label{IsTranslationInvariantOrdering}
}\hfill{\scriptsize (property)}}\\


 for an ordering \mbox{\texttt{\mdseries\slshape ord}} on a family of associative words, returns \texttt{true} if and only if the ordering is translation invariant. 

 This is a property of orderings on families of associative words. An ordering \mbox{\texttt{\mdseries\slshape ord}} over a family $F$, with alphabet $X$ is translation invariant if \texttt{IsLessThanUnder(} \mbox{\texttt{\mdseries\slshape ord}}, $u$, $v$ \texttt{)} implies that for any $a, b \in X^*$, \texttt{IsLessThanUnder(} \mbox{\texttt{\mdseries\slshape ord}}, $a*u*b$, $a*v*b$ \texttt{)}. }

 

\subsection{\textcolor{Chapter }{IsReductionOrdering}}
\logpage{[ 34, 4, 3 ]}\nobreak
\hyperdef{L}{X816CD4BD82D41ED0}{}
{\noindent\textcolor{FuncColor}{$\triangleright$\enspace\texttt{IsReductionOrdering({\mdseries\slshape ord})\index{IsReductionOrdering@\texttt{IsReductionOrdering}}
\label{IsReductionOrdering}
}\hfill{\scriptsize (property)}}\\


 for an ordering \mbox{\texttt{\mdseries\slshape ord}} on a family of associative words, returns \texttt{true} if and only if the ordering is a reduction ordering. An ordering \mbox{\texttt{\mdseries\slshape ord}} is a reduction ordering if it is well founded and translation invariant. }

 

\subsection{\textcolor{Chapter }{OrderingOnGenerators}}
\logpage{[ 34, 4, 4 ]}\nobreak
\hyperdef{L}{X7B6051C282EA88D5}{}
{\noindent\textcolor{FuncColor}{$\triangleright$\enspace\texttt{OrderingOnGenerators({\mdseries\slshape ord})\index{OrderingOnGenerators@\texttt{OrderingOnGenerators}}
\label{OrderingOnGenerators}
}\hfill{\scriptsize (attribute)}}\\


 for an ordering \mbox{\texttt{\mdseries\slshape ord}} on a family of associative words, returns a list in which the generators are
considered. This could be indeed the ordering of the generators in the
ordering, but, for example, if a weight is associated to each generator then
this is not true anymore. See the example for \texttt{WeightLexOrdering} (\ref{WeightLexOrdering}). }

 

\subsection{\textcolor{Chapter }{LexicographicOrdering}}
\logpage{[ 34, 4, 5 ]}\nobreak
\hyperdef{L}{X79B2DEB786680F51}{}
{\noindent\textcolor{FuncColor}{$\triangleright$\enspace\texttt{LexicographicOrdering({\mdseries\slshape D[, gens]})\index{LexicographicOrdering@\texttt{LexicographicOrdering}}
\label{LexicographicOrdering}
}\hfill{\scriptsize (operation)}}\\


 Let \mbox{\texttt{\mdseries\slshape D}} be a free semigroup, a free monoid, or the elements family of such a domain.
Called with only argument \mbox{\texttt{\mdseries\slshape D}}, \texttt{LexicographicOrdering} returns the lexicographic ordering on the elements of \mbox{\texttt{\mdseries\slshape D}}. 

 The optional argument \mbox{\texttt{\mdseries\slshape gens}} can be either the list of free generators of \mbox{\texttt{\mdseries\slshape D}}, in the desired order, or a list of the positions of these generators, in the
desired order, and \texttt{LexicographicOrdering} returns the lexicographic ordering on the elements of \mbox{\texttt{\mdseries\slshape D}} with the ordering on the generators as given. 
\begin{Verbatim}[commandchars=!@|,fontsize=\small,frame=single,label=Example]
  !gapprompt@gap>| !gapinput@f := FreeSemigroup(3);|
  <free semigroup on the generators [ s1, s2, s3 ]>
  !gapprompt@gap>| !gapinput@lex := LexicographicOrdering(f,[2,3,1]);|
  Ordering
  !gapprompt@gap>| !gapinput@IsLessThanUnder(lex,f.2*f.3,f.3);|
  true
  !gapprompt@gap>| !gapinput@IsLessThanUnder(lex,f.3,f.2);|
  false
\end{Verbatim}
 }

 

\subsection{\textcolor{Chapter }{ShortLexOrdering}}
\logpage{[ 34, 4, 6 ]}\nobreak
\hyperdef{L}{X802EB44B7E7B1F57}{}
{\noindent\textcolor{FuncColor}{$\triangleright$\enspace\texttt{ShortLexOrdering({\mdseries\slshape D[, gens]})\index{ShortLexOrdering@\texttt{ShortLexOrdering}}
\label{ShortLexOrdering}
}\hfill{\scriptsize (operation)}}\\


 Let \mbox{\texttt{\mdseries\slshape D}} be a free semigroup, a free monoid, or the elements family of such a domain.
Called with only argument \mbox{\texttt{\mdseries\slshape D}}, \texttt{ShortLexOrdering} returns the shortlex ordering on the elements of \mbox{\texttt{\mdseries\slshape D}}. 

 The optional argument \mbox{\texttt{\mdseries\slshape gens}} can be either the list of free generators of \mbox{\texttt{\mdseries\slshape D}}, in the desired order, or a list of the positions of these generators, in the
desired order, and \texttt{ShortLexOrdering} returns the shortlex ordering on the elements of \mbox{\texttt{\mdseries\slshape D}} with the ordering on the generators as given. }

 

\subsection{\textcolor{Chapter }{IsShortLexOrdering}}
\logpage{[ 34, 4, 7 ]}\nobreak
\hyperdef{L}{X7B6ED9327E0A2099}{}
{\noindent\textcolor{FuncColor}{$\triangleright$\enspace\texttt{IsShortLexOrdering({\mdseries\slshape ord})\index{IsShortLexOrdering@\texttt{IsShortLexOrdering}}
\label{IsShortLexOrdering}
}\hfill{\scriptsize (property)}}\\


 for an ordering \mbox{\texttt{\mdseries\slshape ord}} of a family of associative words, returns \texttt{true} if and only if \mbox{\texttt{\mdseries\slshape ord}} is a shortlex ordering. 
\begin{Verbatim}[commandchars=!@|,fontsize=\small,frame=single,label=Example]
  !gapprompt@gap>| !gapinput@f := FreeSemigroup(3);|
  <free semigroup on the generators [ s1, s2, s3 ]>
  !gapprompt@gap>| !gapinput@sl := ShortLexOrdering(f,[2,3,1]);|
  Ordering
  !gapprompt@gap>| !gapinput@IsLessThanUnder(sl,f.1,f.2);|
  false
  !gapprompt@gap>| !gapinput@IsLessThanUnder(sl,f.3,f.2);|
  false
  !gapprompt@gap>| !gapinput@IsLessThanUnder(sl,f.3,f.1);|
  true
\end{Verbatim}
 }

 

\subsection{\textcolor{Chapter }{WeightLexOrdering}}
\logpage{[ 34, 4, 8 ]}\nobreak
\hyperdef{L}{X849DD7C6782333D5}{}
{\noindent\textcolor{FuncColor}{$\triangleright$\enspace\texttt{WeightLexOrdering({\mdseries\slshape D, gens, wt})\index{WeightLexOrdering@\texttt{WeightLexOrdering}}
\label{WeightLexOrdering}
}\hfill{\scriptsize (operation)}}\\


 Let \mbox{\texttt{\mdseries\slshape D}} be a free semigroup, a free monoid, or the elements family of such a domain. \mbox{\texttt{\mdseries\slshape gens}} can be either the list of free generators of \mbox{\texttt{\mdseries\slshape D}}, in the desired order, or a list of the positions of these generators, in the
desired order. Let \mbox{\texttt{\mdseries\slshape wt}} be a list of weights. \texttt{WeightLexOrdering} returns the weightlex ordering on the elements of \mbox{\texttt{\mdseries\slshape D}} with the ordering on the generators and weights of the generators as given. }

 

\subsection{\textcolor{Chapter }{IsWeightLexOrdering}}
\logpage{[ 34, 4, 9 ]}\nobreak
\hyperdef{L}{X7C7D7954784F5C73}{}
{\noindent\textcolor{FuncColor}{$\triangleright$\enspace\texttt{IsWeightLexOrdering({\mdseries\slshape ord})\index{IsWeightLexOrdering@\texttt{IsWeightLexOrdering}}
\label{IsWeightLexOrdering}
}\hfill{\scriptsize (property)}}\\


 for an ordering \mbox{\texttt{\mdseries\slshape ord}} on a family of associative words, returns \texttt{true} if and only if \mbox{\texttt{\mdseries\slshape ord}} is a weightlex ordering. }

 

\subsection{\textcolor{Chapter }{WeightOfGenerators}}
\logpage{[ 34, 4, 10 ]}\nobreak
\hyperdef{L}{X7E7FAEA484148947}{}
{\noindent\textcolor{FuncColor}{$\triangleright$\enspace\texttt{WeightOfGenerators({\mdseries\slshape ord})\index{WeightOfGenerators@\texttt{WeightOfGenerators}}
\label{WeightOfGenerators}
}\hfill{\scriptsize (attribute)}}\\


 for a weightlex ordering \mbox{\texttt{\mdseries\slshape ord}}, returns a list with length the size of the alphabet of the family. This list
gives the weight of each of the letters of the alphabet which are used for
weightlex orderings with respect to the ordering given by \texttt{OrderingOnGenerators} (\ref{OrderingOnGenerators}). 
\begin{Verbatim}[commandchars=!@|,fontsize=\small,frame=single,label=Example]
  !gapprompt@gap>| !gapinput@f := FreeSemigroup(3);|
  <free semigroup on the generators [ s1, s2, s3 ]>
  !gapprompt@gap>| !gapinput@wtlex := WeightLexOrdering(f,[f.2,f.3,f.1],[3,2,1]);|
  Ordering
  !gapprompt@gap>| !gapinput@IsLessThanUnder(wtlex,f.1,f.2);|
  true
  !gapprompt@gap>| !gapinput@IsLessThanUnder(wtlex,f.3,f.2);|
  true
  !gapprompt@gap>| !gapinput@IsLessThanUnder(wtlex,f.3,f.1);|
  false
  !gapprompt@gap>| !gapinput@OrderingOnGenerators(wtlex);|
  [ s2, s3, s1 ]
  !gapprompt@gap>| !gapinput@WeightOfGenerators(wtlex);|
  [ 3, 2, 1 ]
\end{Verbatim}
 }

 

\subsection{\textcolor{Chapter }{BasicWreathProductOrdering}}
\logpage{[ 34, 4, 11 ]}\nobreak
\hyperdef{L}{X79D1019E7C3E575E}{}
{\noindent\textcolor{FuncColor}{$\triangleright$\enspace\texttt{BasicWreathProductOrdering({\mdseries\slshape D[, gens]})\index{BasicWreathProductOrdering@\texttt{BasicWreathProductOrdering}}
\label{BasicWreathProductOrdering}
}\hfill{\scriptsize (operation)}}\\


 Let \mbox{\texttt{\mdseries\slshape D}} be a free semigroup, a free monoid, or the elements family of such a domain.
Called with only argument \mbox{\texttt{\mdseries\slshape D}}, \texttt{BasicWreathProductOrdering} returns the basic wreath product ordering on the elements of \mbox{\texttt{\mdseries\slshape D}}. 

 The optional argument \mbox{\texttt{\mdseries\slshape gens}} can be either the list of free generators of \mbox{\texttt{\mdseries\slshape D}}, in the desired order, or a list of the positions of these generators, in the
desired order, and \texttt{BasicWreathProductOrdering} returns the lexicographic ordering on the elements of \mbox{\texttt{\mdseries\slshape D}} with the ordering on the generators as given. }

 

\subsection{\textcolor{Chapter }{IsBasicWreathProductOrdering}}
\logpage{[ 34, 4, 12 ]}\nobreak
\hyperdef{L}{X7CB765477FBC3383}{}
{\noindent\textcolor{FuncColor}{$\triangleright$\enspace\texttt{IsBasicWreathProductOrdering({\mdseries\slshape ord})\index{IsBasicWreathProductOrdering@\texttt{IsBasicWreathProductOrdering}}
\label{IsBasicWreathProductOrdering}
}\hfill{\scriptsize (property)}}\\


 
\begin{Verbatim}[commandchars=!@|,fontsize=\small,frame=single,label=Example]
  !gapprompt@gap>| !gapinput@f := FreeSemigroup(3);|
  <free semigroup on the generators [ s1, s2, s3 ]>
  !gapprompt@gap>| !gapinput@basic := BasicWreathProductOrdering(f,[2,3,1]);|
  Ordering
  !gapprompt@gap>| !gapinput@IsLessThanUnder(basic,f.3,f.1);|
  true
  !gapprompt@gap>| !gapinput@IsLessThanUnder(basic,f.3*f.2,f.1);|
  true
  !gapprompt@gap>| !gapinput@IsLessThanUnder(basic,f.3*f.2*f.1,f.1*f.3);|
  false
\end{Verbatim}
 }

 

\subsection{\textcolor{Chapter }{WreathProductOrdering}}
\logpage{[ 34, 4, 13 ]}\nobreak
\hyperdef{L}{X7E6DF1B17F53642E}{}
{\noindent\textcolor{FuncColor}{$\triangleright$\enspace\texttt{WreathProductOrdering({\mdseries\slshape D[, gens], levels})\index{WreathProductOrdering@\texttt{WreathProductOrdering}}
\label{WreathProductOrdering}
}\hfill{\scriptsize (operation)}}\\


 Let \mbox{\texttt{\mdseries\slshape D}} be a free semigroup, a free monoid, or the elements family of such a domain,
let \mbox{\texttt{\mdseries\slshape gens}} be either the list of free generators of \mbox{\texttt{\mdseries\slshape D}}, in the desired order, or a list of the positions of these generators, in the
desired order, and let \mbox{\texttt{\mdseries\slshape levels}} be a list of levels for the generators. If \mbox{\texttt{\mdseries\slshape gens}} is omitted then the default ordering is taken. \texttt{WreathProductOrdering} returns the wreath product ordering on the elements of \mbox{\texttt{\mdseries\slshape D}} with the ordering on the generators as given. }

 

\subsection{\textcolor{Chapter }{IsWreathProductOrdering}}
\logpage{[ 34, 4, 14 ]}\nobreak
\hyperdef{L}{X7F0EE6E987148C96}{}
{\noindent\textcolor{FuncColor}{$\triangleright$\enspace\texttt{IsWreathProductOrdering({\mdseries\slshape ord})\index{IsWreathProductOrdering@\texttt{IsWreathProductOrdering}}
\label{IsWreathProductOrdering}
}\hfill{\scriptsize (property)}}\\


 specifies whether an ordering is a wreath product ordering (see \texttt{WreathProductOrdering} (\ref{WreathProductOrdering})). }

 

\subsection{\textcolor{Chapter }{LevelsOfGenerators}}
\logpage{[ 34, 4, 15 ]}\nobreak
\hyperdef{L}{X7901AA4479EDBE72}{}
{\noindent\textcolor{FuncColor}{$\triangleright$\enspace\texttt{LevelsOfGenerators({\mdseries\slshape ord})\index{LevelsOfGenerators@\texttt{LevelsOfGenerators}}
\label{LevelsOfGenerators}
}\hfill{\scriptsize (attribute)}}\\


 for a wreath product ordering \mbox{\texttt{\mdseries\slshape ord}}, returns the levels of the generators as given at creation (with respect to \texttt{OrderingOnGenerators} (\ref{OrderingOnGenerators})). 
\begin{Verbatim}[commandchars=!@|,fontsize=\small,frame=single,label=Example]
  !gapprompt@gap>| !gapinput@f := FreeSemigroup(3);|
  <free semigroup on the generators [ s1, s2, s3 ]>
  !gapprompt@gap>| !gapinput@wrp := WreathProductOrdering(f,[1,2,3],[1,1,2,]);|
  Ordering
  !gapprompt@gap>| !gapinput@IsLessThanUnder(wrp,f.3,f.1);|
  false
  !gapprompt@gap>| !gapinput@IsLessThanUnder(wrp,f.3,f.2);|
  false
  !gapprompt@gap>| !gapinput@IsLessThanUnder(wrp,f.1,f.2);|
  true
  !gapprompt@gap>| !gapinput@LevelsOfGenerators(wrp);|
  [ 1, 1, 2 ]
\end{Verbatim}
 }

 }

 }

            
\chapter{\textcolor{Chapter }{Magmas}}\label{Magmas}
\logpage{[ 35, 0, 0 ]}
\hyperdef{L}{X873E502F7D21C39C}{}
{
  This chapter deals with domains (see{\nobreakspace}\ref{Domains and their Elements}) that are closed under multiplication \texttt{*}. Following{\nobreakspace}\cite{Bourbaki70}, we call them \emph{magmas} in \textsf{GAP}. Together with the domains closed under addition \texttt{+} (see{\nobreakspace}\ref{Additive Magmas}), they are the basic algebraic structures; every semigroup, monoid
(see{\nobreakspace}\ref{Semigroups}), group (see{\nobreakspace}\ref{Groups}), ring (see{\nobreakspace}\ref{Rings}), or field (see{\nobreakspace}\ref{Fields and Division Rings}) is a magma. In the cases of a \emph{magma\texttt{\symbol{45}}with\texttt{\symbol{45}}one} or \emph{magma\texttt{\symbol{45}}with\texttt{\symbol{45}}inverses}, additional multiplicative structure is present, see{\nobreakspace}\ref{Magma Categories}. For functions to create free magmas, see{\nobreakspace}\ref{Free Magmas}.  
\section{\textcolor{Chapter }{Magma Categories}}\label{Magma Categories}
\logpage{[ 35, 1, 0 ]}
\hyperdef{L}{X7E1248B186E7BB44}{}
{
  

\subsection{\textcolor{Chapter }{IsMagma}}
\logpage{[ 35, 1, 1 ]}\nobreak
\hyperdef{L}{X87D3F38B7EAB13FA}{}
{\noindent\textcolor{FuncColor}{$\triangleright$\enspace\texttt{IsMagma({\mdseries\slshape obj})\index{IsMagma@\texttt{IsMagma}}
\label{IsMagma}
}\hfill{\scriptsize (Category)}}\\


 A \emph{magma} in \textsf{GAP} is a domain $M$ with (not necessarily associative) multiplication \texttt{*}$: M \times M \rightarrow M$. }

 

\subsection{\textcolor{Chapter }{IsMagmaWithOne}}
\logpage{[ 35, 1, 2 ]}\nobreak
\hyperdef{L}{X86071DE7835F1C7C}{}
{\noindent\textcolor{FuncColor}{$\triangleright$\enspace\texttt{IsMagmaWithOne({\mdseries\slshape obj})\index{IsMagmaWithOne@\texttt{IsMagmaWithOne}}
\label{IsMagmaWithOne}
}\hfill{\scriptsize (Category)}}\\


 A \emph{magma\texttt{\symbol{45}}with\texttt{\symbol{45}}one} in \textsf{GAP} is a magma $M$ with an operation \texttt{\texttt{\symbol{94}}0} (or \texttt{One} (\ref{One})) that yields the identity of $M$. 

 So a magma\texttt{\symbol{45}}with\texttt{\symbol{45}}one $M$ does always contain a unique multiplicatively neutral element $e$, i.e., $e$\texttt{ * }$m = m = m$\texttt{ * }$e$ holds for all $m \in M$ (see{\nobreakspace}\texttt{MultiplicativeNeutralElement} (\ref{MultiplicativeNeutralElement})). This element $e$ can be computed with the operation \texttt{One} (\ref{One}) as \texttt{One( }$M$\texttt{ )}, and $e$ is also equal to \texttt{One( }$m$\texttt{ )} and to $m$\texttt{\texttt{\symbol{94}}0} for each element $m \in M$. 

 \emph{Note} that a magma may contain a multiplicatively neutral element but \emph{not} an identity (see{\nobreakspace}\texttt{One} (\ref{One})), and a magma containing an identity may \emph{not} lie in the category \texttt{IsMagmaWithOne} (see Section{\nobreakspace}\ref{Domain Categories}). }

 

\subsection{\textcolor{Chapter }{IsMagmaWithInversesIfNonzero}}
\logpage{[ 35, 1, 3 ]}\nobreak
\hyperdef{L}{X83E4903D7FBB2E24}{}
{\noindent\textcolor{FuncColor}{$\triangleright$\enspace\texttt{IsMagmaWithInversesIfNonzero({\mdseries\slshape obj})\index{IsMagmaWithInversesIfNonzero@\texttt{IsMagmaWithInversesIfNonzero}}
\label{IsMagmaWithInversesIfNonzero}
}\hfill{\scriptsize (Category)}}\\


 An object in this \textsf{GAP} category is a magma\texttt{\symbol{45}}with\texttt{\symbol{45}}one $M$ with an operation \texttt{\texttt{\symbol{94}}\texttt{\symbol{45}}1}$: M \setminus Z \rightarrow M \setminus Z$ that maps each element $m$ of $M \setminus Z$ to its inverse $m$\texttt{\texttt{\symbol{94}}\texttt{\symbol{45}}1} (or \texttt{Inverse( }$m$\texttt{ )}, see{\nobreakspace}\texttt{Inverse} (\ref{Inverse})), where $Z$ is either empty or consists exactly of one element of $M$. 

 This category was introduced mainly to describe division rings, since the
nonzero elements in a division ring form a group; So an object $M$ in \texttt{IsMagmaWithInversesIfNonzero} will usually have both a multiplicative and an additive structure
(see{\nobreakspace}\ref{Additive Magmas}), and the set $Z$, if it is nonempty, contains exactly the zero element (see{\nobreakspace}\texttt{Zero} (\ref{Zero})) of $M$. }

 

\subsection{\textcolor{Chapter }{IsMagmaWithInverses}}
\logpage{[ 35, 1, 4 ]}\nobreak
\hyperdef{L}{X82CBFF648574B830}{}
{\noindent\textcolor{FuncColor}{$\triangleright$\enspace\texttt{IsMagmaWithInverses({\mdseries\slshape obj})\index{IsMagmaWithInverses@\texttt{IsMagmaWithInverses}}
\label{IsMagmaWithInverses}
}\hfill{\scriptsize (Category)}}\\


 A \emph{magma\texttt{\symbol{45}}with\texttt{\symbol{45}}inverses} in \textsf{GAP} is a magma\texttt{\symbol{45}}with\texttt{\symbol{45}}one $M$ with an operation \texttt{\texttt{\symbol{94}}\texttt{\symbol{45}}1}$: M \rightarrow M$ that maps each element $m$ of $M$ to its inverse $m$\texttt{\texttt{\symbol{94}}\texttt{\symbol{45}}1} (or \texttt{Inverse( }$m$\texttt{ )}, see{\nobreakspace}\texttt{Inverse} (\ref{Inverse})). 

 Note that not every trivial magma is a
magma\texttt{\symbol{45}}with\texttt{\symbol{45}}one, but every trivial
magma\texttt{\symbol{45}}with\texttt{\symbol{45}}one is a
magma\texttt{\symbol{45}}with\texttt{\symbol{45}}inverses. This holds also if
the identity of the magma\texttt{\symbol{45}}with\texttt{\symbol{45}}one is a
zero element. So a
magma\texttt{\symbol{45}}with\texttt{\symbol{45}}inverses\texttt{\symbol{45}}if\texttt{\symbol{45}}nonzero
can be a magma\texttt{\symbol{45}}with\texttt{\symbol{45}}inverses if either
it contains no zero element or consists of a zero element that has itself as
zero\texttt{\symbol{45}}th power. }

 }

  
\section{\textcolor{Chapter }{Magma Generation}}\label{Magma Generation}
\logpage{[ 35, 2, 0 ]}
\hyperdef{L}{X808F1A148398733D}{}
{
  This section describes functions that create magmas from generators (see \texttt{Magma} (\ref{Magma}), \texttt{MagmaWithOne} (\ref{MagmaWithOne}), \texttt{MagmaWithInverses} (\ref{MagmaWithInverses})), the underlying operations for which methods can be installed (see \texttt{MagmaByGenerators} (\ref{MagmaByGenerators}), \texttt{MagmaWithOneByGenerators} (\ref{MagmaWithOneByGenerators}), \texttt{MagmaWithInversesByGenerators} (\ref{MagmaWithInversesByGenerators})), functions for forming submagmas (see \texttt{Submagma} (\ref{Submagma}), \texttt{SubmagmaWithOne} (\ref{SubmagmaWithOne}), \texttt{SubmagmaWithInverses} (\ref{SubmagmaWithInverses})), and functions that form a magma equal to a given collection (see \texttt{AsMagma} (\ref{AsMagma}), \texttt{AsSubmagma} (\ref{AsSubmagma})). 

 \texttt{InjectionZeroMagma} (\ref{InjectionZeroMagma}) creates a new magma which is the original magma with a zero adjoined. 

\subsection{\textcolor{Chapter }{Magma}}
\logpage{[ 35, 2, 1 ]}\nobreak
\hyperdef{L}{X839147CF813312D6}{}
{\noindent\textcolor{FuncColor}{$\triangleright$\enspace\texttt{Magma({\mdseries\slshape [Fam, ]gens})\index{Magma@\texttt{Magma}}
\label{Magma}
}\hfill{\scriptsize (function)}}\\


 returns the magma $M$ that is generated by the elements in the list \mbox{\texttt{\mdseries\slshape gens}}, that is, the closure of \mbox{\texttt{\mdseries\slshape gens}} under multiplication \texttt{\texttt{\symbol{92}}*} (\ref{bSlash*}). The family \mbox{\texttt{\mdseries\slshape Fam}} of $M$ can be entered as the first argument; this is obligatory if \mbox{\texttt{\mdseries\slshape gens}} is empty (and hence also $M$ is empty). }

 

\subsection{\textcolor{Chapter }{MagmaWithOne}}
\logpage{[ 35, 2, 2 ]}\nobreak
\hyperdef{L}{X7854B23286B17321}{}
{\noindent\textcolor{FuncColor}{$\triangleright$\enspace\texttt{MagmaWithOne({\mdseries\slshape [Fam, ]gens})\index{MagmaWithOne@\texttt{MagmaWithOne}}
\label{MagmaWithOne}
}\hfill{\scriptsize (function)}}\\


 returns the magma\texttt{\symbol{45}}with\texttt{\symbol{45}}one $M$ that is generated by the elements in the list \mbox{\texttt{\mdseries\slshape gens}}, that is, the closure of \mbox{\texttt{\mdseries\slshape gens}} under multiplication \texttt{\texttt{\symbol{92}}*} (\ref{bSlash*}) and \texttt{One} (\ref{One}). The family \mbox{\texttt{\mdseries\slshape Fam}} of $M$ can be entered as first argument; this is obligatory if \mbox{\texttt{\mdseries\slshape gens}} is empty (and hence $M$ is trivial). }

 

\subsection{\textcolor{Chapter }{MagmaWithInverses}}
\logpage{[ 35, 2, 3 ]}\nobreak
\hyperdef{L}{X7A2B51F67EF4DA28}{}
{\noindent\textcolor{FuncColor}{$\triangleright$\enspace\texttt{MagmaWithInverses({\mdseries\slshape [Fam, ]gens})\index{MagmaWithInverses@\texttt{MagmaWithInverses}}
\label{MagmaWithInverses}
}\hfill{\scriptsize (function)}}\\


 returns the magma\texttt{\symbol{45}}with\texttt{\symbol{45}}inverses $M$ that is generated by the elements in the list \mbox{\texttt{\mdseries\slshape gens}}, that is, the closure of \mbox{\texttt{\mdseries\slshape gens}} under multiplication \texttt{\texttt{\symbol{92}}*} (\ref{bSlash*}), \texttt{One} (\ref{One}), and \texttt{Inverse} (\ref{Inverse}). The family \mbox{\texttt{\mdseries\slshape Fam}} of $M$ can be entered as first argument; this is obligatory if \mbox{\texttt{\mdseries\slshape gens}} is empty (and hence $M$ is trivial). }

 

\subsection{\textcolor{Chapter }{MagmaByGenerators}}
\logpage{[ 35, 2, 4 ]}\nobreak
\hyperdef{L}{X7F629A498383A0AD}{}
{\noindent\textcolor{FuncColor}{$\triangleright$\enspace\texttt{MagmaByGenerators({\mdseries\slshape [Fam, ]gens})\index{MagmaByGenerators@\texttt{MagmaByGenerators}}
\label{MagmaByGenerators}
}\hfill{\scriptsize (operation)}}\\


 An underlying operation for \texttt{Magma} (\ref{Magma}). }

 

\subsection{\textcolor{Chapter }{MagmaWithOneByGenerators}}
\logpage{[ 35, 2, 5 ]}\nobreak
\hyperdef{L}{X84DABBEB803107EB}{}
{\noindent\textcolor{FuncColor}{$\triangleright$\enspace\texttt{MagmaWithOneByGenerators({\mdseries\slshape [Fam, ]gens})\index{MagmaWithOneByGenerators@\texttt{MagmaWithOneByGenerators}}
\label{MagmaWithOneByGenerators}
}\hfill{\scriptsize (operation)}}\\


 An underlying operation for \texttt{MagmaWithOne} (\ref{MagmaWithOne}). }

 

\subsection{\textcolor{Chapter }{MagmaWithInversesByGenerators}}
\logpage{[ 35, 2, 6 ]}\nobreak
\hyperdef{L}{X82C08CFB854E3F1A}{}
{\noindent\textcolor{FuncColor}{$\triangleright$\enspace\texttt{MagmaWithInversesByGenerators({\mdseries\slshape [Fam, ]gens})\index{MagmaWithInversesByGenerators@\texttt{MagmaWithInversesByGenerators}}
\label{MagmaWithInversesByGenerators}
}\hfill{\scriptsize (operation)}}\\


 An underlying operation for \texttt{MagmaWithInverses} (\ref{MagmaWithInverses}). }

 

\subsection{\textcolor{Chapter }{Submagma}}
\logpage{[ 35, 2, 7 ]}\nobreak
\hyperdef{L}{X8268EAA47E4A3A64}{}
{\noindent\textcolor{FuncColor}{$\triangleright$\enspace\texttt{Submagma({\mdseries\slshape D, gens})\index{Submagma@\texttt{Submagma}}
\label{Submagma}
}\hfill{\scriptsize (function)}}\\
\noindent\textcolor{FuncColor}{$\triangleright$\enspace\texttt{SubmagmaNC({\mdseries\slshape D, gens})\index{SubmagmaNC@\texttt{SubmagmaNC}}
\label{SubmagmaNC}
}\hfill{\scriptsize (function)}}\\


 \texttt{Submagma} returns the magma generated by the elements in the list \mbox{\texttt{\mdseries\slshape gens}}, with parent the domain \mbox{\texttt{\mdseries\slshape D}}. \texttt{SubmagmaNC} does the same, except that it is not checked whether the elements of \mbox{\texttt{\mdseries\slshape gens}} lie in \mbox{\texttt{\mdseries\slshape D}}. }

 

\subsection{\textcolor{Chapter }{SubmagmaWithOne}}
\logpage{[ 35, 2, 8 ]}\nobreak
\hyperdef{L}{X7F295EBC7A9CE87E}{}
{\noindent\textcolor{FuncColor}{$\triangleright$\enspace\texttt{SubmagmaWithOne({\mdseries\slshape D, gens})\index{SubmagmaWithOne@\texttt{SubmagmaWithOne}}
\label{SubmagmaWithOne}
}\hfill{\scriptsize (function)}}\\
\noindent\textcolor{FuncColor}{$\triangleright$\enspace\texttt{SubmagmaWithOneNC({\mdseries\slshape D, gens})\index{SubmagmaWithOneNC@\texttt{SubmagmaWithOneNC}}
\label{SubmagmaWithOneNC}
}\hfill{\scriptsize (function)}}\\


 \texttt{SubmagmaWithOne} returns the magma\texttt{\symbol{45}}with\texttt{\symbol{45}}one generated by
the elements in the list \mbox{\texttt{\mdseries\slshape gens}}, with parent the domain \mbox{\texttt{\mdseries\slshape D}}. \texttt{SubmagmaWithOneNC} does the same, except that it is not checked whether the elements of \mbox{\texttt{\mdseries\slshape gens}} lie in \mbox{\texttt{\mdseries\slshape D}}. }

 

\subsection{\textcolor{Chapter }{SubmagmaWithInverses}}
\logpage{[ 35, 2, 9 ]}\nobreak
\hyperdef{L}{X79441F1F7A277E28}{}
{\noindent\textcolor{FuncColor}{$\triangleright$\enspace\texttt{SubmagmaWithInverses({\mdseries\slshape D, gens})\index{SubmagmaWithInverses@\texttt{SubmagmaWithInverses}}
\label{SubmagmaWithInverses}
}\hfill{\scriptsize (function)}}\\
\noindent\textcolor{FuncColor}{$\triangleright$\enspace\texttt{SubmagmaWithInversesNC({\mdseries\slshape D, gens})\index{SubmagmaWithInversesNC@\texttt{SubmagmaWithInversesNC}}
\label{SubmagmaWithInversesNC}
}\hfill{\scriptsize (function)}}\\


 \texttt{SubmagmaWithInverses} returns the magma\texttt{\symbol{45}}with\texttt{\symbol{45}}inverses
generated by the elements in the list \mbox{\texttt{\mdseries\slshape gens}}, with parent the domain \mbox{\texttt{\mdseries\slshape D}}. \texttt{SubmagmaWithInversesNC} does the same, except that it is not checked whether the elements of \mbox{\texttt{\mdseries\slshape gens}} lie in \mbox{\texttt{\mdseries\slshape D}}. }

 

\subsection{\textcolor{Chapter }{AsMagma}}
\logpage{[ 35, 2, 10 ]}\nobreak
\hyperdef{L}{X84ED076D7E46AB79}{}
{\noindent\textcolor{FuncColor}{$\triangleright$\enspace\texttt{AsMagma({\mdseries\slshape C})\index{AsMagma@\texttt{AsMagma}}
\label{AsMagma}
}\hfill{\scriptsize (attribute)}}\\


 For a collection \mbox{\texttt{\mdseries\slshape C}} whose elements form a magma, \texttt{AsMagma} returns this magma. Otherwise \texttt{fail} is returned. }

 

\subsection{\textcolor{Chapter }{AsSubmagma}}
\logpage{[ 35, 2, 11 ]}\nobreak
\hyperdef{L}{X87EEEC018129F0F4}{}
{\noindent\textcolor{FuncColor}{$\triangleright$\enspace\texttt{AsSubmagma({\mdseries\slshape D, C})\index{AsSubmagma@\texttt{AsSubmagma}}
\label{AsSubmagma}
}\hfill{\scriptsize (operation)}}\\


 Let \mbox{\texttt{\mdseries\slshape D}} be a domain and \mbox{\texttt{\mdseries\slshape C}} a collection. If \mbox{\texttt{\mdseries\slshape C}} is a subset of \mbox{\texttt{\mdseries\slshape D}} that forms a magma then \texttt{AsSubmagma} returns this magma, with parent \mbox{\texttt{\mdseries\slshape D}}. Otherwise \texttt{fail} is returned. }

 

\subsection{\textcolor{Chapter }{IsMagmaWithZeroAdjoined}}
\logpage{[ 35, 2, 12 ]}\nobreak
\hyperdef{L}{X8553F44D8123B2C6}{}
{\noindent\textcolor{FuncColor}{$\triangleright$\enspace\texttt{IsMagmaWithZeroAdjoined({\mdseries\slshape M})\index{IsMagmaWithZeroAdjoined@\texttt{IsMagmaWithZeroAdjoined}}
\label{IsMagmaWithZeroAdjoined}
}\hfill{\scriptsize (Category)}}\\
\textbf{\indent Returns:\ }
\texttt{true} or \texttt{false}.



 \texttt{IsMagmaWithZeroAdjoined} returns \texttt{true} if the magma \mbox{\texttt{\mdseries\slshape M}} was created using \texttt{InjectionZeroMagma} (\ref{InjectionZeroMagma}) or \texttt{MagmaWithZeroAdjoined} (\ref{MagmaWithZeroAdjoined}) and returns \texttt{false} if it was not. 
\begin{Verbatim}[commandchars=!@|,fontsize=\small,frame=single,label=Example]
  !gapprompt@gap>| !gapinput@S:=Semigroup(Transformation([1,1,1]), Transformation([1,3,2]));;|
  !gapprompt@gap>| !gapinput@IsMagmaWithZeroAdjoined(S);|
  false
  !gapprompt@gap>| !gapinput@M:=MagmaWithZeroAdjoined(S);|
  <<transformation semigroup of degree 3 with 2 generators>
    with 0 adjoined>
  !gapprompt@gap>| !gapinput@IsMagmaWithZeroAdjoined(M);|
  true
\end{Verbatim}
 }

 

\subsection{\textcolor{Chapter }{InjectionZeroMagma}}
\logpage{[ 35, 2, 13 ]}\nobreak
\hyperdef{L}{X8620878D7FD98823}{}
{\noindent\textcolor{FuncColor}{$\triangleright$\enspace\texttt{InjectionZeroMagma({\mdseries\slshape M})\index{InjectionZeroMagma@\texttt{InjectionZeroMagma}}
\label{InjectionZeroMagma}
}\hfill{\scriptsize (attribute)}}\\
\noindent\textcolor{FuncColor}{$\triangleright$\enspace\texttt{MagmaWithZeroAdjoined({\mdseries\slshape M})\index{MagmaWithZeroAdjoined@\texttt{MagmaWithZeroAdjoined}}
\label{MagmaWithZeroAdjoined}
}\hfill{\scriptsize (attribute)}}\\


 \texttt{InjectionZeroMagma} returns an embedding from the magma \mbox{\texttt{\mdseries\slshape M}} into a new magma formed from \mbox{\texttt{\mdseries\slshape M}} by adjoining a single new element which is the multiplicative zero of the
resulting magma. The elements of the new magma form a family of elements in
the category \texttt{IsMultiplicativeElementWithZero} (\ref{IsMultiplicativeElementWithZero}) and the magma itself satisfies \texttt{IsMagmaWithZeroAdjoined} (\ref{IsMagmaWithZeroAdjoined}).

 \texttt{MagmaWithZeroAdjoined} is just shorthand for \texttt{Range(InjectionZeroMagma(\mbox{\texttt{\mdseries\slshape M}})))}.

 If \texttt{N} is a magma with zero adjoined, then the embedding used to create \texttt{N} can be recovered using \texttt{UnderlyingInjectionZeroMagma} (\ref{UnderlyingInjectionZeroMagma}). 
\begin{Verbatim}[commandchars=!@|,fontsize=\small,frame=single,label=Example]
  !gapprompt@gap>| !gapinput@S:=Monoid(Transformation( [ 7, 7, 5, 3, 1, 3, 7 ] ),|
  !gapprompt@>| !gapinput@Transformation( [ 5, 1, 4, 1, 4, 4, 7 ] ));;|
  !gapprompt@gap>| !gapinput@MultiplicativeZero(S);|
  Transformation( [ 7, 7, 7, 7, 7, 7, 7 ] )
  !gapprompt@gap>| !gapinput@T:=MagmaWithZeroAdjoined(S);|
  <<transformation monoid of degree 7 with 2 generators>
    with 0 adjoined>
  !gapprompt@gap>| !gapinput@map:=UnderlyingInjectionZeroMagma(T);;|
  !gapprompt@gap>| !gapinput@x:=Transformation( [ 7, 7, 7, 3, 7, 3, 7 ] );;|
  !gapprompt@gap>| !gapinput@x^map;|
  <monoid with 0 adjoined elt: Transformation( [ 7, 7, 7, 3, 7, 3, 7 ]
   )>
  !gapprompt@gap>| !gapinput@PreImage(map, x^map)=x;|
  true
\end{Verbatim}
 }

 

\subsection{\textcolor{Chapter }{UnderlyingInjectionZeroMagma}}
\logpage{[ 35, 2, 14 ]}\nobreak
\hyperdef{L}{X7B353674859BF659}{}
{\noindent\textcolor{FuncColor}{$\triangleright$\enspace\texttt{UnderlyingInjectionZeroMagma({\mdseries\slshape M})\index{UnderlyingInjectionZeroMagma@\texttt{UnderlyingInjectionZeroMagma}}
\label{UnderlyingInjectionZeroMagma}
}\hfill{\scriptsize (attribute)}}\\


 \texttt{UnderlyingInjectionZeroMagma} returns the embedding used to create the magma with zero adjoined \mbox{\texttt{\mdseries\slshape M}}. 
\begin{Verbatim}[commandchars=!@|,fontsize=\small,frame=single,label=Example]
  !gapprompt@gap>| !gapinput@S:=Monoid(Transformation( [ 8, 7, 5, 3, 1, 3, 8, 8 ] ),|
  !gapprompt@>| !gapinput@Transformation( [ 5, 1, 4, 1, 4, 4, 7, 8 ] ));;|
  !gapprompt@gap>| !gapinput@MultiplicativeZero(S);|
  Transformation( [ 8, 8, 8, 8, 8, 8, 8, 8 ] )
  !gapprompt@gap>| !gapinput@T:=MagmaWithZeroAdjoined(S);|
  <<transformation monoid of degree 8 with 2 generators>
    with 0 adjoined>
  !gapprompt@gap>| !gapinput@UnderlyingInjectionZeroMagma(T);|
  MappingByFunction( <transformation monoid of degree 8 with 2
   generators>, <<transformation monoid of degree 8 with 2 generators>
    with 0 adjoined>, function( elt ) ... end, function( x ) ... end )
\end{Verbatim}
 }

 }

  
\section{\textcolor{Chapter }{Magmas Defined by Multiplication Tables}}\label{Magmas Defined by Multiplication Tables}
\logpage{[ 35, 3, 0 ]}
\hyperdef{L}{X782215B982F2F01C}{}
{
  The most elementary (but of course usually not recommended) way to implement a
magma with only few elements is via a multiplication table. 

\subsection{\textcolor{Chapter }{MagmaByMultiplicationTable}}
\logpage{[ 35, 3, 1 ]}\nobreak
\hyperdef{L}{X85CD1E7678295CA6}{}
{\noindent\textcolor{FuncColor}{$\triangleright$\enspace\texttt{MagmaByMultiplicationTable({\mdseries\slshape A})\index{MagmaByMultiplicationTable@\texttt{MagmaByMultiplicationTable}}
\label{MagmaByMultiplicationTable}
}\hfill{\scriptsize (function)}}\\


 For a square matrix \mbox{\texttt{\mdseries\slshape A}} with $n$ rows such that all entries of \mbox{\texttt{\mdseries\slshape A}} are in the range $[ 1 .. n ]$, \texttt{MagmaByMultiplicationTable} returns a magma $M$ with multiplication \texttt{*} defined by \mbox{\texttt{\mdseries\slshape A}}. That is, $M$ consists of the elements $m_1, m_2, \ldots, m_n$, and $m_i * m_j = m_k$, with $k =$ \mbox{\texttt{\mdseries\slshape A}}$[i][j]$. 

 The ordering of elements is defined by $m_1 < m_2 < \cdots < m_n$, so $m_i$ can be accessed as \texttt{MagmaElement( \mbox{\texttt{\mdseries\slshape M}}, \mbox{\texttt{\mdseries\slshape i}} )}, see{\nobreakspace}\texttt{MagmaElement} (\ref{MagmaElement}). 
\begin{Verbatim}[commandchars=!@|,fontsize=\small,frame=single,label=Example]
  !gapprompt@gap>| !gapinput@MagmaByMultiplicationTable([[1,2,3],[2,3,1],[1,1,1]]);|
  <magma with 3 generators>
\end{Verbatim}
 }

 

\subsection{\textcolor{Chapter }{MagmaWithOneByMultiplicationTable}}
\logpage{[ 35, 3, 2 ]}\nobreak
\hyperdef{L}{X865526C881645D65}{}
{\noindent\textcolor{FuncColor}{$\triangleright$\enspace\texttt{MagmaWithOneByMultiplicationTable({\mdseries\slshape A})\index{MagmaWithOneByMultiplicationTable@\texttt{MagmaWithOneByMultiplicationTable}}
\label{MagmaWithOneByMultiplicationTable}
}\hfill{\scriptsize (function)}}\\


 The only differences between \texttt{MagmaByMultiplicationTable} (\ref{MagmaByMultiplicationTable}) and \texttt{MagmaWithOneByMultiplicationTable} are that the latter returns a
magma\texttt{\symbol{45}}with\texttt{\symbol{45}}one (see{\nobreakspace}\texttt{MagmaWithOne} (\ref{MagmaWithOne})) if the magma described by the matrix \mbox{\texttt{\mdseries\slshape A}} has an identity, and returns \texttt{fail} if not. 
\begin{Verbatim}[commandchars=!@|,fontsize=\small,frame=single,label=Example]
  !gapprompt@gap>| !gapinput@MagmaWithOneByMultiplicationTable([[1,2,3],[2,3,1],[3,1,1]]);|
  <magma-with-one with 3 generators>
  !gapprompt@gap>| !gapinput@MagmaWithOneByMultiplicationTable([[1,2,3],[2,3,1],[1,1,1]]);|
  fail
\end{Verbatim}
 }

 

\subsection{\textcolor{Chapter }{MagmaWithInversesByMultiplicationTable}}
\logpage{[ 35, 3, 3 ]}\nobreak
\hyperdef{L}{X7EDAFB987EE8A770}{}
{\noindent\textcolor{FuncColor}{$\triangleright$\enspace\texttt{MagmaWithInversesByMultiplicationTable({\mdseries\slshape A})\index{MagmaWithInversesByMultiplicationTable@\texttt{Magma}\-\texttt{With}\-\texttt{Inverses}\-\texttt{By}\-\texttt{Multiplication}\-\texttt{Table}}
\label{MagmaWithInversesByMultiplicationTable}
}\hfill{\scriptsize (function)}}\\


 \texttt{MagmaByMultiplicationTable} (\ref{MagmaByMultiplicationTable}) and \texttt{MagmaWithInversesByMultiplicationTable} differ only in that the latter returns
magma\texttt{\symbol{45}}with\texttt{\symbol{45}}inverses (see{\nobreakspace}\texttt{MagmaWithInverses} (\ref{MagmaWithInverses})) if each element in the magma described by the matrix \mbox{\texttt{\mdseries\slshape A}} has an inverse, and returns \texttt{fail} if not. 
\begin{Verbatim}[commandchars=!@|,fontsize=\small,frame=single,label=Example]
  !gapprompt@gap>| !gapinput@MagmaWithInversesByMultiplicationTable([[1,2,3],[2,3,1],[3,1,2]]);|
  <magma-with-inverses with 3 generators>
  !gapprompt@gap>| !gapinput@MagmaWithInversesByMultiplicationTable([[1,2,3],[2,3,1],[3,2,1]]);|
  fail
\end{Verbatim}
 }

 

\subsection{\textcolor{Chapter }{MagmaElement}}
\logpage{[ 35, 3, 4 ]}\nobreak
\hyperdef{L}{X828BED4580D28FB8}{}
{\noindent\textcolor{FuncColor}{$\triangleright$\enspace\texttt{MagmaElement({\mdseries\slshape M, i})\index{MagmaElement@\texttt{MagmaElement}}
\label{MagmaElement}
}\hfill{\scriptsize (function)}}\\


 For a magma \mbox{\texttt{\mdseries\slshape M}} and a positive integer \mbox{\texttt{\mdseries\slshape i}}, \texttt{MagmaElement} returns the \mbox{\texttt{\mdseries\slshape i}}\texttt{\symbol{45}}th element of \mbox{\texttt{\mdseries\slshape M}}, w.r.t.{\nobreakspace}the ordering \texttt{{\textless}}. If \mbox{\texttt{\mdseries\slshape M}} has less than \mbox{\texttt{\mdseries\slshape i}} elements then \texttt{fail} is returned. }

 
\subsection{\textcolor{Chapter }{MultiplicationTable}}\logpage{[ 35, 3, 5 ]}
\hyperdef{L}{X849BDCC27C4C3191}{}
{
\noindent\textcolor{FuncColor}{$\triangleright$\enspace\texttt{MultiplicationTable({\mdseries\slshape elms})\index{MultiplicationTable@\texttt{MultiplicationTable}!for a list of elements}
\label{MultiplicationTable:for a list of elements}
}\hfill{\scriptsize (attribute)}}\\
\noindent\textcolor{FuncColor}{$\triangleright$\enspace\texttt{MultiplicationTable({\mdseries\slshape M})\index{MultiplicationTable@\texttt{MultiplicationTable}!for a magma}
\label{MultiplicationTable:for a magma}
}\hfill{\scriptsize (attribute)}}\\


 For a list \mbox{\texttt{\mdseries\slshape elms}} of elements that form a magma $M$, \texttt{MultiplicationTable} returns a square matrix $A$ of positive integers such that $A[i][j] = k$ holds if and only if \mbox{\texttt{\mdseries\slshape elms}}$[i] *$ \mbox{\texttt{\mdseries\slshape elms}}$[j] =$ \mbox{\texttt{\mdseries\slshape elms}}$[k]$. This matrix can be used to construct a magma isomorphic to $M$, using \texttt{MagmaByMultiplicationTable} (\ref{MagmaByMultiplicationTable}). 

 For a magma \mbox{\texttt{\mdseries\slshape M}}, \texttt{MultiplicationTable} returns the multiplication table w.r.t.{\nobreakspace}the sorted list of
elements of \mbox{\texttt{\mdseries\slshape M}}. 

 
\begin{Verbatim}[commandchars=!@|,fontsize=\small,frame=single,label=Example]
  !gapprompt@gap>| !gapinput@l:= [ (), (1,2)(3,4), (1,3)(2,4), (1,4)(2,3) ];;|
  !gapprompt@gap>| !gapinput@a:= MultiplicationTable( l );|
  [ [ 1, 2, 3, 4 ], [ 2, 1, 4, 3 ], [ 3, 4, 1, 2 ], [ 4, 3, 2, 1 ] ]
  !gapprompt@gap>| !gapinput@m:= MagmaByMultiplicationTable( a );|
  <magma with 4 generators>
  !gapprompt@gap>| !gapinput@One( m );|
  m1
  !gapprompt@gap>| !gapinput@elm:= MagmaElement( m, 2 );  One( elm );  elm^2;|
  m2
  m1
  m1
  !gapprompt@gap>| !gapinput@Inverse( elm );|
  m2
  !gapprompt@gap>| !gapinput@AsGroup( m );|
  <group of size 4 with 2 generators>
  !gapprompt@gap>| !gapinput@a:= [ [ 1, 2 ], [ 2, 2 ] ];|
  [ [ 1, 2 ], [ 2, 2 ] ]
  !gapprompt@gap>| !gapinput@m:= MagmaByMultiplicationTable( a );|
  <magma with 2 generators>
  !gapprompt@gap>| !gapinput@One( m );  Inverse( MagmaElement( m, 2 ) );|
  m1
  fail
\end{Verbatim}
 }

 }

  
\section{\textcolor{Chapter }{Attributes and Properties for Magmas}}\label{Attributes and Properties for Magmas}
\logpage{[ 35, 4, 0 ]}
\hyperdef{L}{X87036FCE868FFEE9}{}
{
  \emph{Note} that \texttt{IsAssociative} (\ref{IsAssociative}) and \texttt{IsCommutative} (\ref{IsCommutative}) always refer to the multiplication of a domain. If a magma \mbox{\texttt{\mdseries\slshape M}} has also an \emph{additive structure}, e.g., if \mbox{\texttt{\mdseries\slshape M}} is a ring (see{\nobreakspace}\ref{Rings}), then the addition \texttt{+} is always assumed to be associative and commutative, see{\nobreakspace}\ref{Arithmetic Operations for Elements}. 

\subsection{\textcolor{Chapter }{GeneratorsOfMagma}}
\logpage{[ 35, 4, 1 ]}\nobreak
\hyperdef{L}{X872E05B478EC20CA}{}
{\noindent\textcolor{FuncColor}{$\triangleright$\enspace\texttt{GeneratorsOfMagma({\mdseries\slshape M})\index{GeneratorsOfMagma@\texttt{GeneratorsOfMagma}}
\label{GeneratorsOfMagma}
}\hfill{\scriptsize (attribute)}}\\


 is a list \mbox{\texttt{\mdseries\slshape gens}} of elements of the magma \mbox{\texttt{\mdseries\slshape M}} that generates \mbox{\texttt{\mdseries\slshape M}} as a magma, that is, the closure of \mbox{\texttt{\mdseries\slshape gens}} under multiplication \texttt{\texttt{\symbol{92}}*} (\ref{bSlash*}) is \mbox{\texttt{\mdseries\slshape M}}. 

 For a free magma, each generator can also be accessed using the \texttt{.} operator (see \texttt{GeneratorsOfDomain} (\ref{GeneratorsOfDomain})). }

 

\subsection{\textcolor{Chapter }{GeneratorsOfMagmaWithOne}}
\logpage{[ 35, 4, 2 ]}\nobreak
\hyperdef{L}{X87DD93EC8061DD81}{}
{\noindent\textcolor{FuncColor}{$\triangleright$\enspace\texttt{GeneratorsOfMagmaWithOne({\mdseries\slshape M})\index{GeneratorsOfMagmaWithOne@\texttt{GeneratorsOfMagmaWithOne}}
\label{GeneratorsOfMagmaWithOne}
}\hfill{\scriptsize (attribute)}}\\


 is a list \mbox{\texttt{\mdseries\slshape gens}} of elements of the magma\texttt{\symbol{45}}with\texttt{\symbol{45}}one \mbox{\texttt{\mdseries\slshape M}} that generates \mbox{\texttt{\mdseries\slshape M}} as a magma\texttt{\symbol{45}}with\texttt{\symbol{45}}one, that is, the
closure of \mbox{\texttt{\mdseries\slshape gens}} under multiplication \texttt{\texttt{\symbol{92}}*} (\ref{bSlash*}) and \texttt{One} (\ref{One}) is \mbox{\texttt{\mdseries\slshape M}}. 

 For a free magma with one, each generator can also be accessed using the \texttt{.} operator (see \texttt{GeneratorsOfDomain} (\ref{GeneratorsOfDomain})). }

 

\subsection{\textcolor{Chapter }{GeneratorsOfMagmaWithInverses}}
\logpage{[ 35, 4, 3 ]}\nobreak
\hyperdef{L}{X83A901B1857C8489}{}
{\noindent\textcolor{FuncColor}{$\triangleright$\enspace\texttt{GeneratorsOfMagmaWithInverses({\mdseries\slshape M})\index{GeneratorsOfMagmaWithInverses@\texttt{GeneratorsOfMagmaWithInverses}}
\label{GeneratorsOfMagmaWithInverses}
}\hfill{\scriptsize (attribute)}}\\


 is a list \mbox{\texttt{\mdseries\slshape gens}} of elements of the magma\texttt{\symbol{45}}with\texttt{\symbol{45}}inverses \mbox{\texttt{\mdseries\slshape M}} that generates \mbox{\texttt{\mdseries\slshape M}} as a magma\texttt{\symbol{45}}with\texttt{\symbol{45}}inverses, that is, the
closure of \mbox{\texttt{\mdseries\slshape gens}} under multiplication \texttt{\texttt{\symbol{92}}*} (\ref{bSlash*}) and taking inverses (see{\nobreakspace}\texttt{Inverse} (\ref{Inverse})) is \mbox{\texttt{\mdseries\slshape M}}. }

 
\subsection{\textcolor{Chapter }{Centralizer}}\logpage{[ 35, 4, 4 ]}
\hyperdef{L}{X7A2BF4527E08803C}{}
{
\noindent\textcolor{FuncColor}{$\triangleright$\enspace\texttt{Centralizer({\mdseries\slshape M, elm})\index{Centralizer@\texttt{Centralizer}!for a magma and an element}
\label{Centralizer:for a magma and an element}
}\hfill{\scriptsize (operation)}}\\
\noindent\textcolor{FuncColor}{$\triangleright$\enspace\texttt{Centralizer({\mdseries\slshape M, S})\index{Centralizer@\texttt{Centralizer}!for a magma and a submagma}
\label{Centralizer:for a magma and a submagma}
}\hfill{\scriptsize (operation)}}\\
\noindent\textcolor{FuncColor}{$\triangleright$\enspace\texttt{Centralizer({\mdseries\slshape class})\index{Centralizer@\texttt{Centralizer}!for a class of objects in a magma}
\label{Centralizer:for a class of objects in a magma}
}\hfill{\scriptsize (attribute)}}\\


 \index{centraliser}\index{center} For an element \mbox{\texttt{\mdseries\slshape elm}} of the magma \mbox{\texttt{\mdseries\slshape M}} this operation returns the \emph{centralizer} of \mbox{\texttt{\mdseries\slshape elm}}. This is the domain of those elements \mbox{\texttt{\mdseries\slshape m}} $\in$ \mbox{\texttt{\mdseries\slshape M}} that commute with \mbox{\texttt{\mdseries\slshape elm}}. 

 For a submagma \mbox{\texttt{\mdseries\slshape S}} it returns the domain of those elements that commute with \emph{all} elements \mbox{\texttt{\mdseries\slshape s}} of \mbox{\texttt{\mdseries\slshape S}}. 

 If \mbox{\texttt{\mdseries\slshape class}} is a class of objects of a magma (this magma then is stored as the \texttt{ActingDomain} of \mbox{\texttt{\mdseries\slshape class}}) such as given by \texttt{ConjugacyClass} (\ref{ConjugacyClass}), \texttt{Centralizer} returns the centralizer of \texttt{Representative(\mbox{\texttt{\mdseries\slshape class}})} (which is a slight abuse of the notation).    
\begin{Verbatim}[commandchars=!@|,fontsize=\small,frame=single,label=Example]
  !gapprompt@gap>| !gapinput@g:=Group((1,2,3,4),(1,2));;|
  !gapprompt@gap>| !gapinput@Centralizer(g,(1,2,3));|
  Group([ (1,2,3) ])
  !gapprompt@gap>| !gapinput@Centralizer(g,Subgroup(g,[(1,2,3)]));|
  Group([ (1,2,3) ])
  !gapprompt@gap>| !gapinput@Centralizer(g,Subgroup(g,[(1,2,3),(1,2)]));|
  Group(())
\end{Verbatim}
 }

 

\subsection{\textcolor{Chapter }{Centre}}
\logpage{[ 35, 4, 5 ]}\nobreak
\hyperdef{L}{X847ABE6F781C7FE8}{}
{\noindent\textcolor{FuncColor}{$\triangleright$\enspace\texttt{Centre({\mdseries\slshape M})\index{Centre@\texttt{Centre}}
\label{Centre}
}\hfill{\scriptsize (attribute)}}\\
\noindent\textcolor{FuncColor}{$\triangleright$\enspace\texttt{Center({\mdseries\slshape M})\index{Center@\texttt{Center}}
\label{Center}
}\hfill{\scriptsize (attribute)}}\\


 \texttt{Centre} returns the \emph{centre} of the magma \mbox{\texttt{\mdseries\slshape M}}, i.e., the domain of those elements \mbox{\texttt{\mdseries\slshape m}} $\in$ \mbox{\texttt{\mdseries\slshape M}} that commute and associate with all elements of \mbox{\texttt{\mdseries\slshape M}}. That is, the set $\{ m \in M; \forall a, b \in M: ma = am, (ma)b = m(ab), (am)b = a(mb), (ab)m =
a(bm) \}$. 

 \texttt{Center} is just a synonym for \texttt{Centre}. 

 For associative magmas we have that \texttt{Centre( \mbox{\texttt{\mdseries\slshape M}} ) = Centralizer( \mbox{\texttt{\mdseries\slshape M}}, \mbox{\texttt{\mdseries\slshape M}} )}, see{\nobreakspace}\texttt{Centralizer} (\ref{Centralizer:for a magma and a submagma}). 

 The centre of a magma is always commutative (see{\nobreakspace}\texttt{IsCommutative} (\ref{IsCommutative})). (When one installs a new method for \texttt{Centre}, one should set the \texttt{IsCommutative} (\ref{IsCommutative}) value of the result to \texttt{true}, in order to make this information available.) }

 

\subsection{\textcolor{Chapter }{Idempotents}}
\logpage{[ 35, 4, 6 ]}\nobreak
\hyperdef{L}{X7C651C9C78398FFF}{}
{\noindent\textcolor{FuncColor}{$\triangleright$\enspace\texttt{Idempotents({\mdseries\slshape M})\index{Idempotents@\texttt{Idempotents}}
\label{Idempotents}
}\hfill{\scriptsize (attribute)}}\\


 The set of elements of \mbox{\texttt{\mdseries\slshape M}} which are their own squares. }

 

\subsection{\textcolor{Chapter }{IsAssociative}}
\logpage{[ 35, 4, 7 ]}\nobreak
\hyperdef{L}{X7C83B5A47FD18FB7}{}
{\noindent\textcolor{FuncColor}{$\triangleright$\enspace\texttt{IsAssociative({\mdseries\slshape M})\index{IsAssociative@\texttt{IsAssociative}}
\label{IsAssociative}
}\hfill{\scriptsize (property)}}\\


 A collection \mbox{\texttt{\mdseries\slshape M}} of elements that can be multiplied via \texttt{\texttt{\symbol{92}}*} (\ref{bSlash*}) is \emph{associative} if for all elements $a, b, c \in$ \mbox{\texttt{\mdseries\slshape M}} the equality $(a$\texttt{ * }$b)$\texttt{ * }$c = a$\texttt{ * }$(b$\texttt{ * }$c)$ holds. 

 An associative magma is called a \emph{semigroup} (see{\nobreakspace}\ref{Semigroups}), an associative magma\texttt{\symbol{45}}with\texttt{\symbol{45}}one is
called a \emph{monoid} (see{\nobreakspace}\ref{Semigroups}), and an associative
magma\texttt{\symbol{45}}with\texttt{\symbol{45}}inverses is called a \emph{group} (see{\nobreakspace}\ref{Groups}). }

 

\subsection{\textcolor{Chapter }{IsCentral}}
\logpage{[ 35, 4, 8 ]}\nobreak
\hyperdef{L}{X857B0E507D745ADB}{}
{\noindent\textcolor{FuncColor}{$\triangleright$\enspace\texttt{IsCentral({\mdseries\slshape M, obj})\index{IsCentral@\texttt{IsCentral}}
\label{IsCentral}
}\hfill{\scriptsize (operation)}}\\


 \texttt{IsCentral} returns \texttt{true} if the object \mbox{\texttt{\mdseries\slshape obj}}, which must either be an element or a magma, commutes with all elements in
the magma \mbox{\texttt{\mdseries\slshape M}}. }

 

\subsection{\textcolor{Chapter }{IsCommutative}}
\logpage{[ 35, 4, 9 ]}\nobreak
\hyperdef{L}{X830A4A4C795FBC2D}{}
{\noindent\textcolor{FuncColor}{$\triangleright$\enspace\texttt{IsCommutative({\mdseries\slshape M})\index{IsCommutative@\texttt{IsCommutative}}
\label{IsCommutative}
}\hfill{\scriptsize (property)}}\\
\noindent\textcolor{FuncColor}{$\triangleright$\enspace\texttt{IsAbelian({\mdseries\slshape M})\index{IsAbelian@\texttt{IsAbelian}}
\label{IsAbelian}
}\hfill{\scriptsize (property)}}\\


 A collection \mbox{\texttt{\mdseries\slshape M}} of elements that can be multiplied via \texttt{\texttt{\symbol{92}}*} (\ref{bSlash*}) is \emph{commutative} if for all elements $a, b \in$ \mbox{\texttt{\mdseries\slshape M}} the equality $a$\texttt{ * }$b = b$\texttt{ * }$a$ holds. \texttt{IsAbelian} is a synonym of \texttt{IsCommutative}. 

 Note that the commutativity of the \emph{addition} \texttt{\texttt{\symbol{92}}+} (\ref{bSlash+}) in an additive structure can be tested with \texttt{IsAdditivelyCommutative} (\ref{IsAdditivelyCommutative}). }

 

\subsection{\textcolor{Chapter }{MultiplicativeNeutralElement}}
\logpage{[ 35, 4, 10 ]}\nobreak
\hyperdef{L}{X7EE2EA5F7EB7FEC2}{}
{\noindent\textcolor{FuncColor}{$\triangleright$\enspace\texttt{MultiplicativeNeutralElement({\mdseries\slshape M})\index{MultiplicativeNeutralElement@\texttt{MultiplicativeNeutralElement}}
\label{MultiplicativeNeutralElement}
}\hfill{\scriptsize (attribute)}}\\


 returns the element $e$ in the magma \mbox{\texttt{\mdseries\slshape M}} with the property that $e$\texttt{ * }$m = m = m$\texttt{ * }$e$ holds for all $m \in$ \mbox{\texttt{\mdseries\slshape M}}, if such an element exists. Otherwise \texttt{fail} is returned. 

 A magma that is not a magma\texttt{\symbol{45}}with\texttt{\symbol{45}}one can
have a multiplicative neutral element $e$; in this case, $e$ \emph{cannot} be obtained as \texttt{One( \mbox{\texttt{\mdseries\slshape M}} )}, see{\nobreakspace}\texttt{One} (\ref{One}). }

 

\subsection{\textcolor{Chapter }{MultiplicativeZero}}
\logpage{[ 35, 4, 11 ]}\nobreak
\hyperdef{L}{X7B39F93C8136D642}{}
{\noindent\textcolor{FuncColor}{$\triangleright$\enspace\texttt{MultiplicativeZero({\mdseries\slshape M})\index{MultiplicativeZero@\texttt{MultiplicativeZero}}
\label{MultiplicativeZero}
}\hfill{\scriptsize (attribute)}}\\
\noindent\textcolor{FuncColor}{$\triangleright$\enspace\texttt{IsMultiplicativeZero({\mdseries\slshape M, z})\index{IsMultiplicativeZero@\texttt{IsMultiplicativeZero}}
\label{IsMultiplicativeZero}
}\hfill{\scriptsize (operation)}}\\


 \texttt{MultiplicativeZero} returns the multiplicative zero of the magma \mbox{\texttt{\mdseries\slshape M}} which is the element \texttt{z} in \mbox{\texttt{\mdseries\slshape M}} such that \texttt{\mbox{\texttt{\mdseries\slshape z}} * \mbox{\texttt{\mdseries\slshape m}} = \mbox{\texttt{\mdseries\slshape m}} * \mbox{\texttt{\mdseries\slshape z}} = \mbox{\texttt{\mdseries\slshape z}}} for all \mbox{\texttt{\mdseries\slshape m}} in \mbox{\texttt{\mdseries\slshape M}}.

 \texttt{IsMultiplicativeZero} returns \texttt{true} if the element \mbox{\texttt{\mdseries\slshape z}} of the magma \mbox{\texttt{\mdseries\slshape M}} equals the multiplicative zero of \mbox{\texttt{\mdseries\slshape M}}. 
\begin{Verbatim}[commandchars=!@|,fontsize=\small,frame=single,label=Example]
  !gapprompt@gap>| !gapinput@S:=Semigroup( Transformation( [ 1, 1, 1 ] ),|
  !gapprompt@>| !gapinput@Transformation( [ 2, 3, 1 ] ) );|
  <transformation semigroup of degree 3 with 2 generators>
  !gapprompt@gap>| !gapinput@MultiplicativeZero(S);|
  fail
  !gapprompt@gap>| !gapinput@S:=Semigroup( Transformation( [ 1, 1, 1 ] ),|
  !gapprompt@>| !gapinput@Transformation( [ 1, 3, 2 ] ) );|
  <transformation semigroup of degree 3 with 2 generators>
  !gapprompt@gap>| !gapinput@MultiplicativeZero(S);|
  Transformation( [ 1, 1, 1 ] )
\end{Verbatim}
 }

 

\subsection{\textcolor{Chapter }{SquareRoots}}
\logpage{[ 35, 4, 12 ]}\nobreak
\hyperdef{L}{X867DB05A8218FB1E}{}
{\noindent\textcolor{FuncColor}{$\triangleright$\enspace\texttt{SquareRoots({\mdseries\slshape M, elm})\index{SquareRoots@\texttt{SquareRoots}}
\label{SquareRoots}
}\hfill{\scriptsize (operation)}}\\


 is the proper set of all elements $r$ in the magma \mbox{\texttt{\mdseries\slshape M}} such that $r * r =$ \mbox{\texttt{\mdseries\slshape elm}} holds. }

 

\subsection{\textcolor{Chapter }{TrivialSubmagmaWithOne}}
\logpage{[ 35, 4, 13 ]}\nobreak
\hyperdef{L}{X837DA95883CFB985}{}
{\noindent\textcolor{FuncColor}{$\triangleright$\enspace\texttt{TrivialSubmagmaWithOne({\mdseries\slshape M})\index{TrivialSubmagmaWithOne@\texttt{TrivialSubmagmaWithOne}}
\label{TrivialSubmagmaWithOne}
}\hfill{\scriptsize (attribute)}}\\


 is the magma\texttt{\symbol{45}}with\texttt{\symbol{45}}one that has the
identity of the magma\texttt{\symbol{45}}with\texttt{\symbol{45}}one \mbox{\texttt{\mdseries\slshape M}} as only element. }

 }

 }

            
\chapter{\textcolor{Chapter }{Words}}\label{Words}
\logpage{[ 36, 0, 0 ]}
\hyperdef{L}{X7CB0D2F780D15136}{}
{
  This chapter describes categories of \emph{words} and \emph{nonassociative words}, and operations for them. For information about \emph{associative words}, which occur for example as elements in free groups, see
Chapter{\nobreakspace}\ref{Associative Words}.  
\section{\textcolor{Chapter }{Categories of Words and Nonassociative Words}}\label{Categories of Words and Nonassociative Words}
\logpage{[ 36, 1, 0 ]}
\hyperdef{L}{X79AEC832815B9317}{}
{
  

\subsection{\textcolor{Chapter }{IsWord}}
\logpage{[ 36, 1, 1 ]}\nobreak
\hyperdef{L}{X843F5C3A82239398}{}
{\noindent\textcolor{FuncColor}{$\triangleright$\enspace\texttt{IsWord({\mdseries\slshape obj})\index{IsWord@\texttt{IsWord}}
\label{IsWord}
}\hfill{\scriptsize (Category)}}\\
\noindent\textcolor{FuncColor}{$\triangleright$\enspace\texttt{IsWordWithOne({\mdseries\slshape obj})\index{IsWordWithOne@\texttt{IsWordWithOne}}
\label{IsWordWithOne}
}\hfill{\scriptsize (Category)}}\\
\noindent\textcolor{FuncColor}{$\triangleright$\enspace\texttt{IsWordWithInverse({\mdseries\slshape obj})\index{IsWordWithInverse@\texttt{IsWordWithInverse}}
\label{IsWordWithInverse}
}\hfill{\scriptsize (Category)}}\\


 \index{abstract word} Given a free multiplicative structure $M$ that is freely generated by a subset $X$, any expression of an element in $M$ as an iterated product of elements in $X$ is called a \emph{word} over $X$. 

 Interesting cases of free multiplicative structures are those of free
semigroups, free monoids, and free groups, where the multiplication is
associative (see{\nobreakspace}\texttt{IsAssociative} (\ref{IsAssociative})), which are described in Chapter{\nobreakspace}\ref{Associative Words}, and also the case of free magmas, where the multiplication is nonassociative
(see{\nobreakspace}\texttt{IsNonassocWord} (\ref{IsNonassocWord})). 

 Elements in free magmas (see{\nobreakspace}\texttt{FreeMagma} (\ref{FreeMagma:for given rank})) lie in the category \texttt{IsWord}; similarly, elements in free
magmas\texttt{\symbol{45}}with\texttt{\symbol{45}}one (see{\nobreakspace}\texttt{FreeMagmaWithOne} (\ref{FreeMagmaWithOne:for given rank})) lie in the category \texttt{IsWordWithOne}, and so on. 

 \texttt{IsWord} is mainly a ``common roof'' for the two \emph{disjoint} categories \texttt{IsAssocWord} (\ref{IsAssocWord}) and \texttt{IsNonassocWord} (\ref{IsNonassocWord}) of associative and nonassociative words. This means that associative words are \emph{not} regarded as special cases of nonassociative words. The main reason for this
setup is that we are interested in different external representations for
associative and nonassociative words (see{\nobreakspace}\ref{External Representation for Nonassociative Words} and \ref{The External Representation for Associative Words}). 

 Note that elements in finitely presented groups and also elements in
polycyclic groups in \textsf{GAP} are \emph{not} in \texttt{IsWord} although they are usually called words, see Chapters{\nobreakspace}\ref{Finitely Presented Groups} and{\nobreakspace}\ref{Pc Groups}. 

 Words are \emph{constants} (see{\nobreakspace}\ref{Mutability and Copyability}), that is, they are not copyable and not mutable. 

 The usual way to create words is to form them as products of known words,
starting from \emph{generators} of a free structure such as a free magma or a free group (see{\nobreakspace}\texttt{FreeMagma} (\ref{FreeMagma:for given rank}), \texttt{FreeGroup} (\ref{FreeGroup:for given rank})). 

 Words are also used to implement free algebras, in the same way as group
elements are used to implement group algebras (see{\nobreakspace}\ref{Constructing Algebras as Free Algebras} and Chapter{\nobreakspace}\ref{Magma Rings}). 

 
\begin{Verbatim}[commandchars=!@|,fontsize=\small,frame=single,label=Example]
  !gapprompt@gap>| !gapinput@m:= FreeMagmaWithOne( 2 );;  gens:= GeneratorsOfMagmaWithOne( m );|
  [ x1, x2 ]
  !gapprompt@gap>| !gapinput@w1:= gens[1] * gens[2] * gens[1];|
  ((x1*x2)*x1)
  !gapprompt@gap>| !gapinput@w2:= gens[1] * ( gens[2] * gens[1] );|
  (x1*(x2*x1))
  !gapprompt@gap>| !gapinput@w1 = w2;  IsAssociative( m );|
  false
  false
  !gapprompt@gap>| !gapinput@IsWord( w1 );  IsAssocWord( w1 );  IsNonassocWord( w1 );|
  true
  false
  true
  !gapprompt@gap>| !gapinput@s:= FreeMonoid( 2 );;  gens:= GeneratorsOfMagmaWithOne( s );|
  [ m1, m2 ]
  !gapprompt@gap>| !gapinput@u1:= ( gens[1] * gens[2] ) * gens[1];|
  m1*m2*m1
  !gapprompt@gap>| !gapinput@u2:= gens[1] * ( gens[2] * gens[1] );|
  m1*m2*m1
  !gapprompt@gap>| !gapinput@u1 = u2;  IsAssociative( s );|
  true
  true
  !gapprompt@gap>| !gapinput@IsWord( u1 );  IsAssocWord( u1 );  IsNonassocWord( u1 );|
  true
  true
  false
  !gapprompt@gap>| !gapinput@a:= (1,2,3);;  b:= (1,2);;|
  !gapprompt@gap>| !gapinput@w:= a*b*a;;  IsWord( w );|
  false
\end{Verbatim}
 }

 

\subsection{\textcolor{Chapter }{IsWordCollection}}
\logpage{[ 36, 1, 2 ]}\nobreak
\hyperdef{L}{X804B616579F223D8}{}
{\noindent\textcolor{FuncColor}{$\triangleright$\enspace\texttt{IsWordCollection({\mdseries\slshape obj})\index{IsWordCollection@\texttt{IsWordCollection}}
\label{IsWordCollection}
}\hfill{\scriptsize (Category)}}\\


 \texttt{IsWordCollection} is the collections category (see{\nobreakspace}\texttt{CategoryCollections} (\ref{CategoryCollections})) of \texttt{IsWord} (\ref{IsWord}). 
\begin{Verbatim}[commandchars=!@|,fontsize=\small,frame=single,label=Example]
  !gapprompt@gap>| !gapinput@IsWordCollection( m );  IsWordCollection( s );|
  true
  true
  !gapprompt@gap>| !gapinput@IsWordCollection( [ "a", "b" ] );|
  false
\end{Verbatim}
 }

 

\subsection{\textcolor{Chapter }{IsNonassocWord}}
\logpage{[ 36, 1, 3 ]}\nobreak
\hyperdef{L}{X808FA6F97E16502F}{}
{\noindent\textcolor{FuncColor}{$\triangleright$\enspace\texttt{IsNonassocWord({\mdseries\slshape obj})\index{IsNonassocWord@\texttt{IsNonassocWord}}
\label{IsNonassocWord}
}\hfill{\scriptsize (Category)}}\\
\noindent\textcolor{FuncColor}{$\triangleright$\enspace\texttt{IsNonassocWordWithOne({\mdseries\slshape obj})\index{IsNonassocWordWithOne@\texttt{IsNonassocWordWithOne}}
\label{IsNonassocWordWithOne}
}\hfill{\scriptsize (Category)}}\\


 A \emph{nonassociative word} in \textsf{GAP} is an element in a free magma or a free
magma\texttt{\symbol{45}}with\texttt{\symbol{45}}one (see{\nobreakspace}\ref{Free Magmas}). 

 The default methods for \texttt{ViewObj} (\ref{ViewObj}) and \texttt{PrintObj} (\ref{PrintObj}) show nonassociative words as products of letters, where the succession of
multiplications is determined by round brackets. 

 In this sense each nonassociative word describes a ``program'' to form a product of generators. (Also associative words can be interpreted as
such programs, except that the exact succession of multiplications is not
prescribed due to the associativity.) The function \texttt{MappedWord} (\ref{MappedWord}) implements a way to apply such a program. A more general way is provided by
straight line programs (see{\nobreakspace}\ref{Straight Line Programs}). 

 Note that associative words (see Chapter{\nobreakspace}\ref{Associative Words}) are \emph{not} regarded as special cases of nonassociative words (see{\nobreakspace}\texttt{IsWord} (\ref{IsWord})). }

 

\subsection{\textcolor{Chapter }{IsNonassocWordCollection}}
\logpage{[ 36, 1, 4 ]}\nobreak
\hyperdef{L}{X7F81276C80F690DC}{}
{\noindent\textcolor{FuncColor}{$\triangleright$\enspace\texttt{IsNonassocWordCollection({\mdseries\slshape obj})\index{IsNonassocWordCollection@\texttt{IsNonassocWordCollection}}
\label{IsNonassocWordCollection}
}\hfill{\scriptsize (Category)}}\\
\noindent\textcolor{FuncColor}{$\triangleright$\enspace\texttt{IsNonassocWordWithOneCollection({\mdseries\slshape obj})\index{IsNonassocWordWithOneCollection@\texttt{IsNonassocWordWithOneCollection}}
\label{IsNonassocWordWithOneCollection}
}\hfill{\scriptsize (Category)}}\\


 \texttt{IsNonassocWordCollection} is the collections category (see{\nobreakspace}\texttt{CategoryCollections} (\ref{CategoryCollections})) of \texttt{IsNonassocWord} (\ref{IsNonassocWord}), and \texttt{IsNonassocWordWithOneCollection} is the collections category of \texttt{IsNonassocWordWithOne} (\ref{IsNonassocWordWithOne}). }

 }

  
\section{\textcolor{Chapter }{Comparison of Words}}\label{Comparison of Words}
\logpage{[ 36, 2, 0 ]}
\hyperdef{L}{X852C815F85DBE4BD}{}
{
  

\subsection{\textcolor{Chapter }{\texttt{\symbol{92}}= (for nonassociative words)}}
\logpage{[ 36, 2, 1 ]}\nobreak
\hyperdef{L}{X7CA51DD7874115DF}{}
{\noindent\textcolor{FuncColor}{$\triangleright$\enspace\texttt{\texttt{\symbol{92}}=({\mdseries\slshape w1, w2})\index{\texttt{\symbol{92}}=@\texttt{\texttt{\symbol{92}}=}!for nonassociative words}
\label{bSlash=:for nonassociative words}
}\hfill{\scriptsize (operation)}}\\


 \index{equality!nonassociative words} 

 Two words are equal if and only if they are words over the same alphabet and
with equal external representations (see{\nobreakspace}\ref{External Representation for Nonassociative Words} and \ref{The External Representation for Associative Words}). For nonassociative words, the latter means that the words arise from the
letters of the alphabet by the same sequence of multiplications. }

 

\subsection{\textcolor{Chapter }{\texttt{\symbol{92}}{\textless} (for nonassociative words)}}
\logpage{[ 36, 2, 2 ]}\nobreak
\hyperdef{L}{X82D4C7BE803166D6}{}
{\noindent\textcolor{FuncColor}{$\triangleright$\enspace\texttt{\texttt{\symbol{92}}{\textless}({\mdseries\slshape w1, w2})\index{\texttt{\symbol{92}}{\textless}@\texttt{\texttt{\symbol{92}}{\textless}}!for nonassociative words}
\label{bSlash<:for nonassociative words}
}\hfill{\scriptsize (operation)}}\\


 \index{smaller!nonassociative words} Words are ordered according to their external representation. More precisely,
two words can be compared if they are words over the same alphabet, and the
word with smaller external representation is smaller. For nonassociative
words, the ordering is defined in{\nobreakspace}\ref{External Representation for Nonassociative Words}; associative words are ordered by the shortlex ordering via \texttt{{\textless}} (see{\nobreakspace}\ref{The External Representation for Associative Words}). 

 Note that the alphabet of a word is determined by its family
(see{\nobreakspace}\ref{Families}), and that the result of each call to \texttt{FreeMagma} (\ref{FreeMagma:for given rank}), \texttt{FreeGroup} (\ref{FreeGroup:for given rank}) etc. consists of words over a new alphabet. In particular, there is no ``universal'' empty word, every families of words in \texttt{IsWordWithOne} (\ref{IsWordWithOne}) has its own empty word. 

 
\begin{Verbatim}[commandchars=!@|,fontsize=\small,frame=single,label=Example]
  !gapprompt@gap>| !gapinput@m:= FreeMagma( "a", "b" );;|
  !gapprompt@gap>| !gapinput@x:= FreeMagma( "a", "b" );;|
  !gapprompt@gap>| !gapinput@mgens:= GeneratorsOfMagma( m );|
  [ a, b ]
  !gapprompt@gap>| !gapinput@xgens:= GeneratorsOfMagma( x );|
  [ a, b ]
  !gapprompt@gap>| !gapinput@a:= mgens[1];;  b:= mgens[2];;|
  !gapprompt@gap>| !gapinput@a = xgens[1];|
  false
  !gapprompt@gap>| !gapinput@a*(a*a) = (a*a)*a;  a*b = b*a;  a*a = a*a;|
  false
  false
  true
  !gapprompt@gap>| !gapinput@a < b;  b < a;  a < a*b;|
  true
  false
  true
\end{Verbatim}
 }

 }

  
\section{\textcolor{Chapter }{Operations for Words}}\label{Operations for Words}
\logpage{[ 36, 3, 0 ]}
\hyperdef{L}{X7A60A8E57AF13901}{}
{
  Two words can be multiplied via \texttt{*} only if they are words over the same alphabet (see{\nobreakspace}\ref{Comparison of Words}). 

\subsection{\textcolor{Chapter }{MappedWord}}
\logpage{[ 36, 3, 1 ]}\nobreak
\hyperdef{L}{X7EC17930781D104A}{}
{\noindent\textcolor{FuncColor}{$\triangleright$\enspace\texttt{MappedWord({\mdseries\slshape w, gens, imgs})\index{MappedWord@\texttt{MappedWord}}
\label{MappedWord}
}\hfill{\scriptsize (operation)}}\\


 \texttt{MappedWord} returns the object that is obtained by replacing each occurrence in the word \mbox{\texttt{\mdseries\slshape w}} of a generator in the list \mbox{\texttt{\mdseries\slshape gens}} by the corresponding object in the list \mbox{\texttt{\mdseries\slshape imgs}}. The lists \mbox{\texttt{\mdseries\slshape gens}} and \mbox{\texttt{\mdseries\slshape imgs}} must of course have the same length. 

 \texttt{MappedWord} needs to do some preprocessing to get internal generator numbers etc. When
mapping many (several thousand) words, a dedicated loop might be faster. 

 For example, if the elements in \mbox{\texttt{\mdseries\slshape imgs}} are all \emph{associative words} (see Chapter{\nobreakspace}\ref{Associative Words}) in the same family as the elements in \mbox{\texttt{\mdseries\slshape gens}}, and some of them are equal to the corresponding generators in \mbox{\texttt{\mdseries\slshape gens}}, then those may be omitted from \mbox{\texttt{\mdseries\slshape gens}} and \mbox{\texttt{\mdseries\slshape imgs}}. In this situation, the special case that the lists \mbox{\texttt{\mdseries\slshape gens}} and \mbox{\texttt{\mdseries\slshape imgs}} have only length $1$ is handled more efficiently by \texttt{EliminatedWord} (\ref{EliminatedWord}). 

 If the word is from a free group, it is permitted to give inverses of (some)
of the generators as extra generators. This can speed up the execution by
removing the need to calculate inverses anew. 
\begin{Verbatim}[commandchars=!@|,fontsize=\small,frame=single,label=Example]
  !gapprompt@gap>| !gapinput@m:= FreeMagma( "a", "b" );;  gens:= GeneratorsOfMagma( m );;|
  !gapprompt@gap>| !gapinput@a:= gens[1];  b:= gens[2];|
  a
  b
  !gapprompt@gap>| !gapinput@w:= (a*b)*((b*a)*a)*b;|
  (((a*b)*((b*a)*a))*b)
  !gapprompt@gap>| !gapinput@MappedWord( w, gens, [ (1,2), (1,2,3,4) ] );|
  (2,4,3)
  !gapprompt@gap>| !gapinput@a:= (1,2);; b:= (1,2,3,4);;  (a*b)*((b*a)*a)*b;|
  (2,4,3)
  !gapprompt@gap>| !gapinput@f:= FreeGroup( "a", "b" );;|
  !gapprompt@gap>| !gapinput@a:= GeneratorsOfGroup(f)[1];;  b:= GeneratorsOfGroup(f)[2];;|
  !gapprompt@gap>| !gapinput@w:= a^5*b*a^2/b^4*a;|
  a^5*b*a^2*b^-4*a
  !gapprompt@gap>| !gapinput@MappedWord( w, [ a, b ], [ (1,2), (1,2,3,4) ] );|
  (1,3,4,2)
  !gapprompt@gap>| !gapinput@MappedWord( w, [ a, b, b^-1 ], [ (1,2), (1,2,3,4), (1,4,3,2) ] );|
  (1,3,4,2)
  !gapprompt@gap>| !gapinput@(1,2)^5*(1,2,3,4)*(1,2)^2/(1,2,3,4)^4*(1,2);|
  (1,3,4,2)
  !gapprompt@gap>| !gapinput@MappedWord( w, [ a ], [ a^2 ] );|
  a^10*b*a^4*b^-4*a^2
\end{Verbatim}
 }

 }

  
\section{\textcolor{Chapter }{Free Magmas}}\label{Free Magmas}
\logpage{[ 36, 4, 0 ]}
\hyperdef{L}{X7F51B17983019D3E}{}
{
  The easiest way to create a family of words is to construct the free object
generated by these words. Each such free object defines a unique alphabet, and
its generators are simply the words of length one over this alphabet; These
generators can be accessed via \texttt{GeneratorsOfMagma} (\ref{GeneratorsOfMagma}) in the case of a free magma, and via \texttt{GeneratorsOfMagmaWithOne} (\ref{GeneratorsOfMagmaWithOne}) in the case of a free magma\texttt{\symbol{45}}with\texttt{\symbol{45}}one. 
\subsection{\textcolor{Chapter }{FreeMagma}}\logpage{[ 36, 4, 1 ]}
\hyperdef{L}{X7CFFD9027DDD1555}{}
{
\noindent\textcolor{FuncColor}{$\triangleright$\enspace\texttt{FreeMagma({\mdseries\slshape rank[, name]})\index{FreeMagma@\texttt{FreeMagma}!for given rank}
\label{FreeMagma:for given rank}
}\hfill{\scriptsize (function)}}\\
\noindent\textcolor{FuncColor}{$\triangleright$\enspace\texttt{FreeMagma({\mdseries\slshape name1[, name2[, ...]]})\index{FreeMagma@\texttt{FreeMagma}!for various names}
\label{FreeMagma:for various names}
}\hfill{\scriptsize (function)}}\\
\noindent\textcolor{FuncColor}{$\triangleright$\enspace\texttt{FreeMagma({\mdseries\slshape names})\index{FreeMagma@\texttt{FreeMagma}!for a list of names}
\label{FreeMagma:for a list of names}
}\hfill{\scriptsize (function)}}\\
\noindent\textcolor{FuncColor}{$\triangleright$\enspace\texttt{FreeMagma({\mdseries\slshape infinity[, name][, init]})\index{FreeMagma@\texttt{FreeMagma}!for infinitely many generators}
\label{FreeMagma:for infinitely many generators}
}\hfill{\scriptsize (function)}}\\


 \texttt{FreeMagma} returns a free magma. The number of generators, and the labels given to the
generators, can be specified in several different ways. Warning: the labels of
generators are only an aid for printing, and do not necessarily distinguish
generators; see the examples at the end of \texttt{FreeSemigroup} (\ref{FreeSemigroup:for given rank}) for more information. 
\begin{description}
\item[{ 1: For a given rank, and an optional generator name prefix }]  Called with a positive integer \mbox{\texttt{\mdseries\slshape rank}}, \texttt{FreeMagma} returns a free magma on \mbox{\texttt{\mdseries\slshape rank}} generators. The optional argument \mbox{\texttt{\mdseries\slshape name}} must be a string; its default value is \texttt{"x"}. 

 If \mbox{\texttt{\mdseries\slshape name}} is not given but the \texttt{generatorNames} option is, then this option is respected as described in
Section{\nobreakspace}\ref{Generator Names}. 

 Otherwise, the generators of the returned free magma are labelled \mbox{\texttt{\mdseries\slshape name}}\texttt{1}, ..., \mbox{\texttt{\mdseries\slshape name}}\texttt{k}, where \texttt{k} is the value of \mbox{\texttt{\mdseries\slshape rank}}. 

 
\item[{2: For given generator names}]  Called with various (at least one) nonempty strings, \texttt{FreeMagma} returns a free magma on as many generators as arguments, which are labelled \mbox{\texttt{\mdseries\slshape name1}}, \mbox{\texttt{\mdseries\slshape name2}}, etc. 
\item[{3: For a given list of generator names}]  Called with a finite nonempty list \mbox{\texttt{\mdseries\slshape names}} of nonempty strings, \texttt{FreeMagma} returns a free magma on \texttt{Length(\mbox{\texttt{\mdseries\slshape names}})} generators, whose \texttt{i}\texttt{\symbol{45}}th generator is labelled \mbox{\texttt{\mdseries\slshape names}}\texttt{[i]}. 
\item[{ 4: For the rank \texttt{infinity}, an optional default generator name prefix, and an optional finite list of
generator names }]  Called in the fourth form, \texttt{FreeMagma} returns a free magma on infinitely many generators. The optional argument \mbox{\texttt{\mdseries\slshape name}} must be a string; its default value is \texttt{"x"}, and the optional argument \mbox{\texttt{\mdseries\slshape init}} must be a finite list of nonempty strings; its default value is an empty list.
The generators are initially labelled according to the list \mbox{\texttt{\mdseries\slshape init}}, followed by \mbox{\texttt{\mdseries\slshape name}}\texttt{i} for each \texttt{i} in the range from \texttt{Length(\mbox{\texttt{\mdseries\slshape init}})+1} to \texttt{infinity}. 
\end{description}
 
\begin{Verbatim}[commandchars=!@|,fontsize=\small,frame=single,label=Example]
  !gapprompt@gap>| !gapinput@FreeMagma( 4 );|
  <free magma on the generators [ x1, x2, x3, x4 ]>
  !gapprompt@gap>| !gapinput@FreeMagma( 3, "a" );|
  <free magma on the generators [ a1, a2, a3 ]>
  !gapprompt@gap>| !gapinput@FreeMagma( "a", "b" );|
  <free magma on the generators [ a, b ]>
  !gapprompt@gap>| !gapinput@FreeMagma( [ "a", "b" ] );|
  <free magma on the generators [ a, b ]>
  !gapprompt@gap>| !gapinput@FreeMagma( infinity );|
  <free magma with infinity generators>
  !gapprompt@gap>| !gapinput@F := FreeMagma( infinity, "gen" );;|
  !gapprompt@gap>| !gapinput@GeneratorsOfMagma( F ){[ 1 .. 4 ]};|
  [ gen1, gen2, gen3, gen4 ]
  !gapprompt@gap>| !gapinput@F := FreeMagma( infinity, [ "z", "a" ] );;|
  !gapprompt@gap>| !gapinput@GeneratorsOfMagma( F ){[ 1 .. 3 ]};|
  [ z, a, x3 ]
  !gapprompt@gap>| !gapinput@F := FreeMagma( infinity, "y", [ "z", "a" ] );;|
  !gapprompt@gap>| !gapinput@GeneratorsOfMagma( F ){[ 1 .. 4 ]};|
  [ z, a, y3, y4 ]
  !gapprompt@gap>| !gapinput@FreeMagma( 3 : generatorNames := "elt" );|
  <free magma on the generators [ elt1, elt2, elt3 ]>
  !gapprompt@gap>| !gapinput@FreeMagma( 2 : generatorNames := [ "u", "v", "w" ] );|
  <free magma on the generators [ u, v ]>
\end{Verbatim}
 }

 
\subsection{\textcolor{Chapter }{FreeMagmaWithOne}}\logpage{[ 36, 4, 2 ]}
\hyperdef{L}{X86DB748080B4A9B9}{}
{
\noindent\textcolor{FuncColor}{$\triangleright$\enspace\texttt{FreeMagmaWithOne({\mdseries\slshape rank[, name]})\index{FreeMagmaWithOne@\texttt{FreeMagmaWithOne}!for given rank}
\label{FreeMagmaWithOne:for given rank}
}\hfill{\scriptsize (function)}}\\
\noindent\textcolor{FuncColor}{$\triangleright$\enspace\texttt{FreeMagmaWithOne({\mdseries\slshape [name1[, name2[, ...]]]})\index{FreeMagmaWithOne@\texttt{FreeMagmaWithOne}!for various names}
\label{FreeMagmaWithOne:for various names}
}\hfill{\scriptsize (function)}}\\
\noindent\textcolor{FuncColor}{$\triangleright$\enspace\texttt{FreeMagmaWithOne({\mdseries\slshape names})\index{FreeMagmaWithOne@\texttt{FreeMagmaWithOne}!for a list of names}
\label{FreeMagmaWithOne:for a list of names}
}\hfill{\scriptsize (function)}}\\
\noindent\textcolor{FuncColor}{$\triangleright$\enspace\texttt{FreeMagmaWithOne({\mdseries\slshape infinity[, name][, init]})\index{FreeMagmaWithOne@\texttt{FreeMagmaWithOne}!for infinitely many generators}
\label{FreeMagmaWithOne:for infinitely many generators}
}\hfill{\scriptsize (function)}}\\


 \texttt{FreeMagmaWithOne} returns a free magma\texttt{\symbol{45}}with\texttt{\symbol{45}}one. The
number of generators, and the labels given to the generators, can be specified
in several different ways. Warning: the labels of generators are only an aid
for printing, and do not necessarily distinguish generators; see the examples
at the end of \texttt{FreeSemigroup} (\ref{FreeSemigroup:for given rank}) for more information. 
\begin{description}
\item[{ 1: For a given rank, and an optional generator name prefix }]  Called with a nonnegative integer \mbox{\texttt{\mdseries\slshape rank}}, \texttt{FreeMagmaWithOne} returns a free magma\texttt{\symbol{45}}with\texttt{\symbol{45}}one on \mbox{\texttt{\mdseries\slshape rank}} generators. The optional argument \mbox{\texttt{\mdseries\slshape name}} must be a string; its default value is \texttt{"x"}. 

 If \mbox{\texttt{\mdseries\slshape name}} is not given but the \texttt{generatorNames} option is, then this option is respected as described in
Section{\nobreakspace}\ref{Generator Names}. 

 Otherwise, the generators of the returned free
magma\texttt{\symbol{45}}with\texttt{\symbol{45}}one are labelled \mbox{\texttt{\mdseries\slshape name}}\texttt{1}, ..., \mbox{\texttt{\mdseries\slshape name}}\texttt{k}, where \texttt{k} is the value of \mbox{\texttt{\mdseries\slshape rank}}. 

 
\item[{2: For given generator names}]  Called with various nonempty strings, \texttt{FreeMagmaWithOne} returns a free magma\texttt{\symbol{45}}with\texttt{\symbol{45}}one on as many
generators as arguments, which are labelled \mbox{\texttt{\mdseries\slshape name1}}, \mbox{\texttt{\mdseries\slshape name2}}, etc. 
\item[{3: For a given list of generator names}]  Called with a finite list \mbox{\texttt{\mdseries\slshape names}} of nonempty strings, \texttt{FreeMagmaWithOne} returns a free magma\texttt{\symbol{45}}with\texttt{\symbol{45}}one on \texttt{Length(\mbox{\texttt{\mdseries\slshape names}})} generators, whose \texttt{i}\texttt{\symbol{45}}th generator is labelled \mbox{\texttt{\mdseries\slshape names}}\texttt{[i]}. 
\item[{ 4: For the rank \texttt{infinity}, an optional default generator name prefix, and an optional finite list of
generator names }]  Called in the fourth form, \texttt{FreeMagmaWithOne} returns a free magma\texttt{\symbol{45}}with\texttt{\symbol{45}}one on
infinitely many generators. The optional argument \mbox{\texttt{\mdseries\slshape name}} must be a string; its default value is \texttt{"x"}, and the optional argument \mbox{\texttt{\mdseries\slshape init}} must be a finite list of nonempty strings; its default value is an empty list.
The generators are initially labelled according to the list \mbox{\texttt{\mdseries\slshape init}}, followed by \mbox{\texttt{\mdseries\slshape name}}\texttt{i} for each \texttt{i} in the range from \texttt{Length(\mbox{\texttt{\mdseries\slshape init}})+1} to \texttt{infinity}. 
\end{description}
 
\begin{Verbatim}[commandchars=!@|,fontsize=\small,frame=single,label=Example]
  !gapprompt@gap>| !gapinput@FreeMagmaWithOne( 4 );|
  <free magma-with-one on the generators [ x1, x2, x3, x4 ]>
  !gapprompt@gap>| !gapinput@FreeMagmaWithOne( 3, "a" );|
  <free magma-with-one on the generators [ a1, a2, a3 ]>
  !gapprompt@gap>| !gapinput@FreeMagmaWithOne( "a", "b" );|
  <free magma-with-one on the generators [ a, b ]>
  !gapprompt@gap>| !gapinput@FreeMagmaWithOne( [ "a", "b" ] );|
  <free magma-with-one on the generators [ a, b ]>
  !gapprompt@gap>| !gapinput@FreeMagmaWithOne( infinity );|
  <free magma-with-one with infinity generators>
  !gapprompt@gap>| !gapinput@F := FreeMagmaWithOne( infinity, "gen" );;|
  !gapprompt@gap>| !gapinput@GeneratorsOfMagmaWithOne( F ){[ 1 .. 4 ]};|
  [ gen1, gen2, gen3, gen4 ]
  !gapprompt@gap>| !gapinput@F := FreeMagmaWithOne( infinity, [ "z", "a" ] );;|
  !gapprompt@gap>| !gapinput@GeneratorsOfMagmaWithOne( F ){[ 1 .. 3 ]};|
  [ z, a, x3 ]
  !gapprompt@gap>| !gapinput@F := FreeMagmaWithOne( infinity, "y", [ "z", "a" ] );;|
  !gapprompt@gap>| !gapinput@GeneratorsOfMagmaWithOne( F ){[ 1 .. 4 ]};|
  [ z, a, y3, y4 ]
  !gapprompt@gap>| !gapinput@FreeMagmaWithOne( 0 );|
  <free group of rank zero>
  !gapprompt@gap>| !gapinput@FreeMagmaWithOne( 3 : generatorNames := "elt" );|
  <free magma-with-one on the generators [ elt1, elt2, elt3 ]>
  !gapprompt@gap>| !gapinput@FreeMagmaWithOne( 2 : generatorNames := [ "u", "v", "w" ] );|
  <free magma-with-one on the generators [ u, v ]>
\end{Verbatim}
 }

 }

  
\section{\textcolor{Chapter }{External Representation for Nonassociative Words}}\label{External Representation for Nonassociative Words}
\logpage{[ 36, 5, 0 ]}
\hyperdef{L}{X84C2F9037EEE9CED}{}
{
  The external representation of nonassociative words is defined as follows. The $i$\texttt{\symbol{45}}th generator of the family of elements in question has
external representation $i$, the identity (if exists) has external representation $0$, the inverse of the $i$\texttt{\symbol{45}}th generator (if exists) has external representation $-i$. If $v$ and $w$ are nonassociative words with external representations $e_v$ and $e_w$, respectively then the product $v * w$ has external representation $[ e_v, e_w ]$. So the external representation of any nonassociative word is either an
integer or a nested list of integers and lists, where each list has length
two. 

 One can create a nonassociative word from a family of words and the external
representation of a nonassociative word using \texttt{ObjByExtRep} (\ref{ObjByExtRep}). 

 
\begin{Verbatim}[commandchars=!@|,fontsize=\small,frame=single,label=Example]
  !gapprompt@gap>| !gapinput@m:= FreeMagma( 2 );;  gens:= GeneratorsOfMagma( m );|
  [ x1, x2 ]
  !gapprompt@gap>| !gapinput@w:= ( gens[1] * gens[2] ) * gens[1];|
  ((x1*x2)*x1)
  !gapprompt@gap>| !gapinput@ExtRepOfObj( w );  ExtRepOfObj( gens[1] );|
  [ [ 1, 2 ], 1 ]
  1
  !gapprompt@gap>| !gapinput@ ExtRepOfObj( w*w );|
  [ [ [ 1, 2 ], 1 ], [ [ 1, 2 ], 1 ] ]
  !gapprompt@gap>| !gapinput@ObjByExtRep( FamilyObj( w ), 2 );|
  x2
  !gapprompt@gap>| !gapinput@ObjByExtRep( FamilyObj( w ), [ 1, [ 2, 1 ] ] );|
  (x1*(x2*x1))
\end{Verbatim}
 }

 }

            
\chapter{\textcolor{Chapter }{Associative Words}}\label{Associative Words}
\logpage{[ 37, 0, 0 ]}
\hyperdef{L}{X78C56A0A87CE380E}{}
{
   
\section{\textcolor{Chapter }{Categories of Associative Words}}\label{Categories of Associative Words}
\logpage{[ 37, 1, 0 ]}
\hyperdef{L}{X7AB546CB7B929253}{}
{
  Associative words are used to represent elements in free groups, semigroups
and monoids in \textsf{GAP} (see{\nobreakspace}\ref{Free Groups, Monoids and Semigroups}). An associative word is just a sequence of letters, where each letter is an
element of an alphabet (in the following called a \emph{generator}) or its inverse. Associative words can be multiplied; in free monoids also
the computation of an identity is permitted, in free groups also the
computation of inverses (see{\nobreakspace}\ref{Operations for Associative Words}). 

 Different alphabets correspond to different families of associative words.
There is no relation whatsoever between words in different families. 

 
\begin{Verbatim}[commandchars=!@|,fontsize=\small,frame=single,label=Example]
  !gapprompt@gap>| !gapinput@f:= FreeGroup( "a", "b", "c" );|
  <free group on the generators [ a, b, c ]>
  !gapprompt@gap>| !gapinput@gens:= GeneratorsOfGroup(f);|
  [ a, b, c ]
  !gapprompt@gap>| !gapinput@w:= gens[1]*gens[2]/gens[3]*gens[2]*gens[1]/gens[1]*gens[3]/gens[2];|
  a*b*c^-1*b*c*b^-1
  !gapprompt@gap>| !gapinput@w^-1;|
  b*c^-1*b^-1*c*b^-1*a^-1
\end{Verbatim}
 

 Words are displayed as products of letters. The letters are usually printed
like \texttt{f1}, \texttt{f2}, $\ldots$, but it is possible to give user defined names to them, which can be
arbitrary strings. These names do not necessarily identify a unique letter
(generator), it is possible to have several letters {\textendash}even in the
same family{\textendash} that are displayed in the same way. Note also that \emph{there is no relation between the names of letters and variable names}. In the example above, we might have typed 

 
\begin{Verbatim}[commandchars=!@|,fontsize=\small,frame=single,label=Example]
  !gapprompt@gap>| !gapinput@a:= f.1;; b:= f.2;; c:= f.3;;|
\end{Verbatim}
 

 (\emph{Interactively}, the function \texttt{AssignGeneratorVariables} (\ref{AssignGeneratorVariables}) provides a shorthand for this.) This allows us to define \texttt{w} more conveniently: 

 
\begin{Verbatim}[commandchars=!@|,fontsize=\small,frame=single,label=Example]
  !gapprompt@gap>| !gapinput@w := a*b/c*b*a/a*c/b;|
  a*b*c^-1*b*c*b^-1
\end{Verbatim}
 

 Using homomorphisms it is possible to express elements of a group as words in
terms of generators, see{\nobreakspace}\ref{Expressing Group Elements as Words in Generators}. 

\subsection{\textcolor{Chapter }{IsAssocWord}}
\logpage{[ 37, 1, 1 ]}\nobreak
\hyperdef{L}{X7FA8DA728773BA89}{}
{\noindent\textcolor{FuncColor}{$\triangleright$\enspace\texttt{IsAssocWord({\mdseries\slshape obj})\index{IsAssocWord@\texttt{IsAssocWord}}
\label{IsAssocWord}
}\hfill{\scriptsize (Category)}}\\
\noindent\textcolor{FuncColor}{$\triangleright$\enspace\texttt{IsAssocWordWithOne({\mdseries\slshape obj})\index{IsAssocWordWithOne@\texttt{IsAssocWordWithOne}}
\label{IsAssocWordWithOne}
}\hfill{\scriptsize (Category)}}\\
\noindent\textcolor{FuncColor}{$\triangleright$\enspace\texttt{IsAssocWordWithInverse({\mdseries\slshape obj})\index{IsAssocWordWithInverse@\texttt{IsAssocWordWithInverse}}
\label{IsAssocWordWithInverse}
}\hfill{\scriptsize (Category)}}\\


 \texttt{IsAssocWord} is the category of associative words in free semigroups, \texttt{IsAssocWordWithOne} is the category of associative words in free monoids (which admit the
operation \texttt{One} (\ref{One}) to compute an identity), \texttt{IsAssocWordWithInverse} is the category of associative words in free groups (which have an inverse).
See{\nobreakspace}\texttt{IsWord} (\ref{IsWord}) for more general categories of words. }

 }

  
\section{\textcolor{Chapter }{Free Groups, Monoids and Semigroups}}\label{Free Groups, Monoids and Semigroups}
\logpage{[ 37, 2, 0 ]}
\hyperdef{L}{X82E7EA7F7FD31EC3}{}
{
  Usually a family of associative words will be generated by constructing the
free object generated by them. See \texttt{FreeMonoid} (\ref{FreeMonoid:for given rank}), \texttt{FreeSemigroup} (\ref{FreeSemigroup:for given rank}) for details. 
\subsection{\textcolor{Chapter }{FreeGroup}}\logpage{[ 37, 2, 1 ]}
\hyperdef{L}{X8215999E835290F0}{}
{
\noindent\textcolor{FuncColor}{$\triangleright$\enspace\texttt{FreeGroup({\mdseries\slshape [wfilt, ]rank[, name]})\index{FreeGroup@\texttt{FreeGroup}!for given rank}
\label{FreeGroup:for given rank}
}\hfill{\scriptsize (function)}}\\
\noindent\textcolor{FuncColor}{$\triangleright$\enspace\texttt{FreeGroup({\mdseries\slshape [wfilt][,] [name1[, name2[, ...]]]})\index{FreeGroup@\texttt{FreeGroup}!for various names}
\label{FreeGroup:for various names}
}\hfill{\scriptsize (function)}}\\
\noindent\textcolor{FuncColor}{$\triangleright$\enspace\texttt{FreeGroup({\mdseries\slshape [wfilt, ]names})\index{FreeGroup@\texttt{FreeGroup}!for a list of names}
\label{FreeGroup:for a list of names}
}\hfill{\scriptsize (function)}}\\
\noindent\textcolor{FuncColor}{$\triangleright$\enspace\texttt{FreeGroup({\mdseries\slshape [wfilt, ]infinity[, name][, init]})\index{FreeGroup@\texttt{FreeGroup}!for infinitely many generators}
\label{FreeGroup:for infinitely many generators}
}\hfill{\scriptsize (function)}}\\


 \texttt{FreeGroup} returns a free group. The number of generators, and the labels given to the
generators, can be specified in several different ways. Warning: the labels of
generators are only an aid for printing, and do not necessarily distinguish
generators; see the examples at the end of \texttt{FreeSemigroup} (\ref{FreeSemigroup:for given rank}) for more information. 
\begin{description}
\item[{ 1: For a given rank, and an optional generator name prefix }]  Called with a nonnegative integer \mbox{\texttt{\mdseries\slshape rank}}, \texttt{FreeGroup} returns a free group on \mbox{\texttt{\mdseries\slshape rank}} generators. The optional argument \mbox{\texttt{\mdseries\slshape name}} must be a string; its default value is \texttt{"f"}. 

 If \mbox{\texttt{\mdseries\slshape name}} is not given but the \texttt{generatorNames} option is, then this option is respected as described in
Section{\nobreakspace}\ref{Generator Names}. 

 Otherwise, the generators of the returned free group are labelled \mbox{\texttt{\mdseries\slshape name}}\texttt{1}, ..., \mbox{\texttt{\mdseries\slshape name}}\texttt{k}, where \texttt{k} is the value of \mbox{\texttt{\mdseries\slshape rank}}. 

 
\item[{2: For given generator names}]  Called with various nonempty strings, \texttt{FreeGroup} returns a free group on as many generators as arguments, which are labelled \mbox{\texttt{\mdseries\slshape name1}}, \mbox{\texttt{\mdseries\slshape name2}}, etc. 
\item[{3: For a given list of generator names}]  Called with a finite list \mbox{\texttt{\mdseries\slshape names}} of nonempty strings, \texttt{FreeGroup} returns a free group on \texttt{Length(\mbox{\texttt{\mdseries\slshape names}})} generators, whose \texttt{i}\texttt{\symbol{45}}th generator is labelled \mbox{\texttt{\mdseries\slshape names}}\texttt{[i]}. 
\item[{ 4: For the rank \texttt{infinity}, an optional default generator name prefix, and an optional finite list of
generator names }]  Called in the fourth form, \texttt{FreeGroup} returns a free group on infinitely many generators. The optional argument \mbox{\texttt{\mdseries\slshape name}} must be a string; its default value is \texttt{"f"}, and the optional argument \mbox{\texttt{\mdseries\slshape init}} must be a finite list of nonempty strings; its default value is an empty list.
The generators are initially labelled according to the list \mbox{\texttt{\mdseries\slshape init}}, followed by \mbox{\texttt{\mdseries\slshape name}}\texttt{i} for each \texttt{i} in the range from \texttt{Length(\mbox{\texttt{\mdseries\slshape init}})+1} to \texttt{infinity}. 
\end{description}
 If the optional first argument \mbox{\texttt{\mdseries\slshape wfilt}} is given, then it must be either \texttt{IsSyllableWordsFamily}, \texttt{IsLetterWordsFamily}, \texttt{IsWLetterWordsFamily}, or \texttt{IsBLetterWordsFamily}. This filter specifies the representation used for the elements of the free
group (see{\nobreakspace}\ref{Representations for Associative Words}). If no such filter is given, a letter representation is used. 

 (For interfacing to old code that omits the representation flag, use of the
syllable representation is also triggered by setting the option \texttt{FreeGroupFamilyType} to the string \texttt{"syllable"}; this is overwritten by the optional first argument if it is given.) 
\begin{Verbatim}[commandchars=!@|,fontsize=\small,frame=single,label=Example]
  !gapprompt@gap>| !gapinput@FreeGroup(5);|
  <free group on the generators [ f1, f2, f3, f4, f5 ]>
  !gapprompt@gap>| !gapinput@FreeGroup(4, "gen");|
  <free group on the generators [ gen1, gen2, gen3, gen4 ]>
  !gapprompt@gap>| !gapinput@FreeGroup(3 : generatorNames := "ack");|
  <free group on the generators [ ack1, ack2, ack3 ]>
  !gapprompt@gap>| !gapinput@FreeGroup(2 : generatorNames := ["u", "v", "w"]);|
  <free group on the generators [ u, v ]>
  !gapprompt@gap>| !gapinput@FreeGroup();|
  <free group of rank zero>
  !gapprompt@gap>| !gapinput@FreeGroup("a", "b", "c");|
  <free group on the generators [ a, b, c ]>
  !gapprompt@gap>| !gapinput@FreeGroup(["x", "y"]);|
  <free group on the generators [ x, y ]>
  !gapprompt@gap>| !gapinput@FreeGroup(infinity);|
  <free group with infinity generators>
  !gapprompt@gap>| !gapinput@F := FreeGroup(infinity, "g", ["a", "b"]);|
  <free group with infinity generators>
  !gapprompt@gap>| !gapinput@GeneratorsOfGroup(F){[1..4]};|
  [ a, b, g3, g4 ]
  !gapprompt@gap>| !gapinput@GeneratorsOfGroup(FreeGroup(infinity, "gen")){[1..3]};|
  [ gen1, gen2, gen3 ]
  !gapprompt@gap>| !gapinput@FreeGroup(IsSyllableWordsFamily, 50);|
  <free group with 50 generators>
\end{Verbatim}
 }

 

\subsection{\textcolor{Chapter }{IsFreeGroup}}
\logpage{[ 37, 2, 2 ]}\nobreak
\hyperdef{L}{X8601654A7C4AF1E7}{}
{\noindent\textcolor{FuncColor}{$\triangleright$\enspace\texttt{IsFreeGroup({\mdseries\slshape obj})\index{IsFreeGroup@\texttt{IsFreeGroup}}
\label{IsFreeGroup}
}\hfill{\scriptsize (Category)}}\\


 Any group consisting of elements in \texttt{IsAssocWordWithInverse} (\ref{IsAssocWordWithInverse}) lies in the filter \texttt{IsFreeGroup}; this holds in particular for any group created with \texttt{FreeGroup} (\ref{FreeGroup:for given rank}), or any subgroup of such a group. 

     Also see Chapter{\nobreakspace}\ref{Finitely Presented Groups}. }

 

\subsection{\textcolor{Chapter }{AssignGeneratorVariables}}
\logpage{[ 37, 2, 3 ]}\nobreak
\hyperdef{L}{X814203E281F3272E}{}
{\noindent\textcolor{FuncColor}{$\triangleright$\enspace\texttt{AssignGeneratorVariables({\mdseries\slshape G})\index{AssignGeneratorVariables@\texttt{AssignGeneratorVariables}}
\label{AssignGeneratorVariables}
}\hfill{\scriptsize (operation)}}\\


 If \mbox{\texttt{\mdseries\slshape G}} is a group, whose generators are represented by symbols (for example a free
group, a finitely presented group or a pc group) this function assigns these
generators to global variables with the same names. 

 The aim of this function is to make it easy in interactive use to work with
(for example) a free group. It is a shorthand for a sequence of assignments of
the form 

 
\begin{Verbatim}[commandchars=!@|,fontsize=\small,frame=single,label=Example]
  var1:=GeneratorsOfGroup(G)[1];
  var2:=GeneratorsOfGroup(G)[2];
  ...
  varn:=GeneratorsOfGroup(G)[n];
\end{Verbatim}
 

 However, since overwriting global variables can be very dangerous, \emph{it is not permitted to use this function within a function}. (If {\textendash}despite this warning{\textendash} this is done, the result
is undefined.) 

 If the assignment overwrites existing variables a warning is given, if any of
the variables is write protected, or any of the generator names would not be a
proper variable name, an error is raised. }

 }

  
\section{\textcolor{Chapter }{Comparison of Associative Words}}\label{Comparison of Associative Words}
\logpage{[ 37, 3, 0 ]}
\hyperdef{L}{X8405BECB7AC4EB61}{}
{
  

\subsection{\textcolor{Chapter }{\texttt{\symbol{92}}= (for associative words)}}
\logpage{[ 37, 3, 1 ]}\nobreak
\hyperdef{L}{X8206153078E97B90}{}
{\noindent\textcolor{FuncColor}{$\triangleright$\enspace\texttt{\texttt{\symbol{92}}=({\mdseries\slshape w1, w2})\index{\texttt{\symbol{92}}=@\texttt{\texttt{\symbol{92}}=}!for associative words}
\label{bSlash=:for associative words}
}\hfill{\scriptsize (operation)}}\\


 \index{equality!associative words} Two associative words are equal if they are words over the same alphabet and
if they are sequences of the same letters. This is equivalent to saying that
the external representations of the words are equal, see{\nobreakspace}\ref{The External Representation for Associative Words} and \ref{Comparison of Words}. 

 There is no ``universal'' empty word, every alphabet (that is, every family of words) has its own empty
word. 
\begin{Verbatim}[commandchars=!@|,fontsize=\small,frame=single,label=Example]
  !gapprompt@gap>| !gapinput@f:= FreeGroup( "a", "b", "b" );;|
  !gapprompt@gap>| !gapinput@gens:= GeneratorsOfGroup(f);|
  [ a, b, b ]
  !gapprompt@gap>| !gapinput@gens[2] = gens[3];|
  false
  !gapprompt@gap>| !gapinput@x:= gens[1]*gens[2];|
  a*b
  !gapprompt@gap>| !gapinput@y:= gens[2]/gens[2]*gens[1]*gens[2];|
  a*b
  !gapprompt@gap>| !gapinput@x = y;|
  true
  !gapprompt@gap>| !gapinput@z:= gens[2]/gens[2]*gens[1]*gens[3];|
  a*b
  !gapprompt@gap>| !gapinput@x = z;|
  false
\end{Verbatim}
 }

 

\subsection{\textcolor{Chapter }{\texttt{\symbol{92}}{\textless} (for associative words)}}
\logpage{[ 37, 3, 2 ]}\nobreak
\hyperdef{L}{X7BB12B9D7F990899}{}
{\noindent\textcolor{FuncColor}{$\triangleright$\enspace\texttt{\texttt{\symbol{92}}{\textless}({\mdseries\slshape w1, w2})\index{\texttt{\symbol{92}}{\textless}@\texttt{\texttt{\symbol{92}}{\textless}}!for associative words}
\label{bSlash<:for associative words}
}\hfill{\scriptsize (operation)}}\\


 \index{smaller!associative words} The ordering of associative words is defined by length and lexicography (this
ordering is called \emph{short\texttt{\symbol{45}}lex} ordering), that is, shorter words are smaller than longer words, and words of
the same length are compared w.r.t.{\nobreakspace}the lexicographical ordering
induced by the ordering of generators. Generators are sorted according to the
order in which they were created. If the generators are invertible then each
generator \mbox{\texttt{\mdseries\slshape g}} is larger than its inverse \mbox{\texttt{\mdseries\slshape g}}\texttt{\texttt{\symbol{94}}\texttt{\symbol{45}}1}, and \mbox{\texttt{\mdseries\slshape g}}\texttt{\texttt{\symbol{94}}\texttt{\symbol{45}}1} is larger than every generator that is smaller than \mbox{\texttt{\mdseries\slshape g}}. 
\begin{Verbatim}[commandchars=!@|,fontsize=\small,frame=single,label=Example]
  !gapprompt@gap>| !gapinput@f:= FreeGroup( 2 );;  gens:= GeneratorsOfGroup( f );;|
  !gapprompt@gap>| !gapinput@a:= gens[1];;  b:= gens[2];;|
  !gapprompt@gap>| !gapinput@One(f) < a^-1;  a^-1 < a;  a < b^-1;  b^-1 < b; b < a^2;  a^2 < a*b;|
  true
  true
  true
  true
  true
  true
\end{Verbatim}
 }

 

\subsection{\textcolor{Chapter }{IsShortLexLessThanOrEqual}}
\logpage{[ 37, 3, 3 ]}\nobreak
\hyperdef{L}{X805C519682B0A7ED}{}
{\noindent\textcolor{FuncColor}{$\triangleright$\enspace\texttt{IsShortLexLessThanOrEqual({\mdseries\slshape u, v})\index{IsShortLexLessThanOrEqual@\texttt{IsShortLexLessThanOrEqual}}
\label{IsShortLexLessThanOrEqual}
}\hfill{\scriptsize (function)}}\\


 returns \texttt{IsLessThanOrEqualUnder(\mbox{\texttt{\mdseries\slshape ord}}, \mbox{\texttt{\mdseries\slshape u}}, \mbox{\texttt{\mdseries\slshape v}})} where \mbox{\texttt{\mdseries\slshape ord}} is the short less ordering for the family of \mbox{\texttt{\mdseries\slshape u}} and \mbox{\texttt{\mdseries\slshape v}}. (This is here for compatibility with \textsf{GAP}{\nobreakspace}4.2.) }

 

\subsection{\textcolor{Chapter }{IsBasicWreathLessThanOrEqual}}
\logpage{[ 37, 3, 4 ]}\nobreak
\hyperdef{L}{X84875E08847B39E1}{}
{\noindent\textcolor{FuncColor}{$\triangleright$\enspace\texttt{IsBasicWreathLessThanOrEqual({\mdseries\slshape u, v})\index{IsBasicWreathLessThanOrEqual@\texttt{IsBasicWreathLessThanOrEqual}}
\label{IsBasicWreathLessThanOrEqual}
}\hfill{\scriptsize (function)}}\\


 returns \texttt{IsLessThanOrEqualUnder(\mbox{\texttt{\mdseries\slshape ord}}, \mbox{\texttt{\mdseries\slshape u}}, \mbox{\texttt{\mdseries\slshape v}})} where \mbox{\texttt{\mdseries\slshape ord}} is the basic wreath product ordering for the family of \mbox{\texttt{\mdseries\slshape u}} and \mbox{\texttt{\mdseries\slshape v}}. (This is here for compatibility with \textsf{GAP}{\nobreakspace}4.2.) }

 }

  
\section{\textcolor{Chapter }{Operations for Associative Words}}\label{Operations for Associative Words}
\logpage{[ 37, 4, 0 ]}
\hyperdef{L}{X79AF6C757A3547BD}{}
{
  The product of two given associative words is defined as the freely reduced
concatenation of the words. \index{product!of words} \index{quotient!of words} \index{power!of words} \index{conjugate!of a word} Besides the multiplication \texttt{\texttt{\symbol{92}}*} (\ref{bSlash*}), the arithmetical operators \texttt{One} (\ref{One}) (if the word lies in a family with identity) and (if the generators are
invertible) \texttt{Inverse} (\ref{Inverse}), \texttt{\texttt{\symbol{92}}/} (\ref{bSlash/}),\texttt{\texttt{\symbol{92}}\texttt{\symbol{94}}} (\ref{bSlash^}), \index{Comm@\texttt{Comm}!for words} \texttt{Comm} (\ref{Comm}), and \index{LeftQuotient@\texttt{LeftQuotient}!for words} \texttt{LeftQuotient} (\ref{LeftQuotient}) are applicable to associative words, see{\nobreakspace}\ref{Arithmetic Operations for Elements}. 

 See also \texttt{MappedWord} (\ref{MappedWord}), an operation that is applicable to arbitrary words. 

 See Section \ref{Representations for Associative Words} for a discussion of the internal representations of associative words that are
supported by \textsf{GAP}. Note that operations to extract or act on parts of words (letter or
syllables) can carry substantially different costs, depending on the
representation the words are in. 

 

\subsection{\textcolor{Chapter }{Length (for an associative word)}}
\logpage{[ 37, 4, 1 ]}\nobreak
\hyperdef{L}{X87CD4C6978A7936A}{}
{\noindent\textcolor{FuncColor}{$\triangleright$\enspace\texttt{Length({\mdseries\slshape w})\index{Length@\texttt{Length}!for an associative word}
\label{Length:for an associative word}
}\hfill{\scriptsize (attribute)}}\\


 \index{length!of a word} For an associative word \mbox{\texttt{\mdseries\slshape w}}, \texttt{Length} returns the number of letters in \mbox{\texttt{\mdseries\slshape w}}. 

 
\begin{Verbatim}[commandchars=!@|,fontsize=\small,frame=single,label=Example]
  !gapprompt@gap>| !gapinput@f := FreeGroup("a","b");; gens := GeneratorsOfGroup(f);;|
  !gapprompt@gap>| !gapinput@a := gens[1];; b := gens[2];;w := a^5*b*a^2*b^-4*a;;|
  !gapprompt@gap>| !gapinput@ w; Length( w );  Length( a^17 );  Length( w^0 );|
  a^5*b*a^2*b^-4*a
  13
  17
  0
\end{Verbatim}
 }

 

\subsection{\textcolor{Chapter }{ExponentSumWord}}
\logpage{[ 37, 4, 2 ]}\nobreak
\hyperdef{L}{X7F5ED4357A9C12E6}{}
{\noindent\textcolor{FuncColor}{$\triangleright$\enspace\texttt{ExponentSumWord({\mdseries\slshape w, gen})\index{ExponentSumWord@\texttt{ExponentSumWord}}
\label{ExponentSumWord}
}\hfill{\scriptsize (operation)}}\\


 For an associative word \mbox{\texttt{\mdseries\slshape w}} and a generator \mbox{\texttt{\mdseries\slshape gen}}, \texttt{ExponentSumWord} returns the number of times \mbox{\texttt{\mdseries\slshape gen}} appears in \mbox{\texttt{\mdseries\slshape w}} minus the number of times its inverse appears in \mbox{\texttt{\mdseries\slshape w}}. If both \mbox{\texttt{\mdseries\slshape gen}} and its inverse do not occur in \mbox{\texttt{\mdseries\slshape w}} then $0$ is returned. \mbox{\texttt{\mdseries\slshape gen}} may also be the inverse of a generator. 
\begin{Verbatim}[commandchars=!@|,fontsize=\small,frame=single,label=Example]
  !gapprompt@gap>| !gapinput@w;  ExponentSumWord( w, a );  ExponentSumWord( w, b );|
  a^5*b*a^2*b^-4*a
  8
  -3
  !gapprompt@gap>| !gapinput@ExponentSumWord( (a*b*a^-1)^3, a );  ExponentSumWord( w, b^-1 );|
  0
  3
\end{Verbatim}
 }

 

\subsection{\textcolor{Chapter }{Subword}}
\logpage{[ 37, 4, 3 ]}\nobreak
\hyperdef{L}{X82CC92C17AF6FFA0}{}
{\noindent\textcolor{FuncColor}{$\triangleright$\enspace\texttt{Subword({\mdseries\slshape w, from, to})\index{Subword@\texttt{Subword}}
\label{Subword}
}\hfill{\scriptsize (operation)}}\\


 For an associative word \mbox{\texttt{\mdseries\slshape w}} and two positive integers \mbox{\texttt{\mdseries\slshape from}} and \mbox{\texttt{\mdseries\slshape to}}, \texttt{Subword} returns the subword of \mbox{\texttt{\mdseries\slshape w}} that begins at position \mbox{\texttt{\mdseries\slshape from}} and ends at position \mbox{\texttt{\mdseries\slshape to}}. Indexing is done with origin 1. 
\begin{Verbatim}[commandchars=!@|,fontsize=\small,frame=single,label=Example]
  !gapprompt@gap>| !gapinput@w;  Subword( w, 3, 7 );|
  a^5*b*a^2*b^-4*a
  a^3*b*a
\end{Verbatim}
 }

 

\subsection{\textcolor{Chapter }{PositionWord}}
\logpage{[ 37, 4, 4 ]}\nobreak
\hyperdef{L}{X8509A0A4851981BB}{}
{\noindent\textcolor{FuncColor}{$\triangleright$\enspace\texttt{PositionWord({\mdseries\slshape w, sub, from})\index{PositionWord@\texttt{PositionWord}}
\label{PositionWord}
}\hfill{\scriptsize (operation)}}\\


 Let \mbox{\texttt{\mdseries\slshape w}} and \mbox{\texttt{\mdseries\slshape sub}} be associative words, and \mbox{\texttt{\mdseries\slshape from}} a positive integer. \texttt{PositionWord} returns the position of the first occurrence of \mbox{\texttt{\mdseries\slshape sub}} as a subword of \mbox{\texttt{\mdseries\slshape w}}, starting at position \mbox{\texttt{\mdseries\slshape from}}. If there is no such occurrence, \texttt{fail} is returned. Indexing is done with origin 1. 

 In other words, \texttt{PositionWord( \mbox{\texttt{\mdseries\slshape w}}, \mbox{\texttt{\mdseries\slshape sub}}, \mbox{\texttt{\mdseries\slshape from}} )} is the smallest integer $i$ larger than or equal to \mbox{\texttt{\mdseries\slshape from}} such that \texttt{Subword( \mbox{\texttt{\mdseries\slshape w}}, }$i$\texttt{,} $i$\texttt{+Length( \mbox{\texttt{\mdseries\slshape sub}} )\texttt{\symbol{45}}1 ) =} \mbox{\texttt{\mdseries\slshape sub}}, see{\nobreakspace}\texttt{Subword} (\ref{Subword}). 

 
\begin{Verbatim}[commandchars=!@|,fontsize=\small,frame=single,label=Example]
  !gapprompt@gap>| !gapinput@w;  PositionWord( w, a/b, 1 );|
  a^5*b*a^2*b^-4*a
  8
  !gapprompt@gap>| !gapinput@Subword( w, 8, 9 );|
  a*b^-1
  !gapprompt@gap>| !gapinput@PositionWord( w, a^2, 1 );|
  1
  !gapprompt@gap>| !gapinput@PositionWord( w, a^2, 2 );|
  2
  !gapprompt@gap>| !gapinput@PositionWord( w, a^2, 6 );|
  7
  !gapprompt@gap>| !gapinput@PositionWord( w, a^2, 8 );|
  fail
\end{Verbatim}
 }

 
\subsection{\textcolor{Chapter }{SubstitutedWord}}\logpage{[ 37, 4, 5 ]}
\hyperdef{L}{X79186218787C224A}{}
{
\noindent\textcolor{FuncColor}{$\triangleright$\enspace\texttt{SubstitutedWord({\mdseries\slshape w, from, to, by})\index{SubstitutedWord@\texttt{SubstitutedWord}!replace an interval by a given word}
\label{SubstitutedWord:replace an interval by a given word}
}\hfill{\scriptsize (operation)}}\\
\noindent\textcolor{FuncColor}{$\triangleright$\enspace\texttt{SubstitutedWord({\mdseries\slshape w, sub, from, by})\index{SubstitutedWord@\texttt{SubstitutedWord}!replace a subword by a given word}
\label{SubstitutedWord:replace a subword by a given word}
}\hfill{\scriptsize (operation)}}\\


 Let \mbox{\texttt{\mdseries\slshape w}} be an associative word. 

 In the first form, \texttt{SubstitutedWord} returns the associative word obtained by replacing the subword of \mbox{\texttt{\mdseries\slshape w}} that begins at position \mbox{\texttt{\mdseries\slshape from}} and ends at position \mbox{\texttt{\mdseries\slshape to}} by the associative word \mbox{\texttt{\mdseries\slshape by}}. \mbox{\texttt{\mdseries\slshape from}} and \mbox{\texttt{\mdseries\slshape to}} must be positive integers, indexing is done with origin 1. In other words, \texttt{SubstitutedWord( \mbox{\texttt{\mdseries\slshape w}}, \mbox{\texttt{\mdseries\slshape from}}, \mbox{\texttt{\mdseries\slshape to}}, \mbox{\texttt{\mdseries\slshape by}} )} is the product of the three words \texttt{Subword( \mbox{\texttt{\mdseries\slshape w}}, 1, \mbox{\texttt{\mdseries\slshape from}}\texttt{\symbol{45}}1 )}, \mbox{\texttt{\mdseries\slshape by}}, and \texttt{Subword( \mbox{\texttt{\mdseries\slshape w}}, \mbox{\texttt{\mdseries\slshape to}}+1, Length( \mbox{\texttt{\mdseries\slshape w}} ) )}, see{\nobreakspace}\texttt{Subword} (\ref{Subword}). 

 In the second form, \texttt{SubstitutedWord} returns the associative word obtained by replacing the first occurrence of the
associative word \mbox{\texttt{\mdseries\slshape sub}} of \mbox{\texttt{\mdseries\slshape w}}, starting at position \mbox{\texttt{\mdseries\slshape from}}, by the associative word \mbox{\texttt{\mdseries\slshape by}}; if there is no such occurrence, \texttt{fail} is returned. 
\begin{Verbatim}[commandchars=!@|,fontsize=\small,frame=single,label=Example]
  !gapprompt@gap>| !gapinput@w;  SubstitutedWord( w, 3, 7, a^19 );|
  a^5*b*a^2*b^-4*a
  a^22*b^-4*a
  !gapprompt@gap>| !gapinput@SubstitutedWord( w, a, 6, b^7 );|
  a^5*b^8*a*b^-4*a
  !gapprompt@gap>| !gapinput@SubstitutedWord( w, a*b, 6, b^7 );|
  fail
\end{Verbatim}
 }

 

\subsection{\textcolor{Chapter }{EliminatedWord}}
\logpage{[ 37, 4, 6 ]}\nobreak
\hyperdef{L}{X8486BFE1844CFE59}{}
{\noindent\textcolor{FuncColor}{$\triangleright$\enspace\texttt{EliminatedWord({\mdseries\slshape w, gen, by})\index{EliminatedWord@\texttt{EliminatedWord}}
\label{EliminatedWord}
}\hfill{\scriptsize (operation)}}\\


 For an associative word \mbox{\texttt{\mdseries\slshape w}}, a generator \mbox{\texttt{\mdseries\slshape gen}}, and an associative word \mbox{\texttt{\mdseries\slshape by}}, \texttt{EliminatedWord} returns the associative word obtained by replacing each occurrence of \mbox{\texttt{\mdseries\slshape gen}} in \mbox{\texttt{\mdseries\slshape w}} by \mbox{\texttt{\mdseries\slshape by}}. 
\begin{Verbatim}[commandchars=!@|,fontsize=\small,frame=single,label=Example]
  !gapprompt@gap>| !gapinput@w;  EliminatedWord( w, a, a^2 );  EliminatedWord( w, a, b^-1 );|
  a^5*b*a^2*b^-4*a
  a^10*b*a^4*b^-4*a^2
  b^-11
\end{Verbatim}
 }

 }

  
\section{\textcolor{Chapter }{Operations for Associative Words by their Syllables}}\label{Operations for Associative Words by their Syllables}
\logpage{[ 37, 5, 0 ]}
\hyperdef{L}{X7D357E047ABD2C6B}{}
{
  For an associative word \mbox{\texttt{\mdseries\slshape w}} $= x_1^{{n_1}} x_2^{{n_2}} \cdots x_k^{{n_k}}$ over an alphabet containing $x_1, x_2, \ldots, x_k$, such that $x_i \neq x_{{i+1}}^{{\pm 1}}$ for $1 \leq i \leq k-1$, the subwords $x_i^{{e_i}}$ are uniquely determined; these powers of generators are called the \emph{syllables} of $w$. 

\subsection{\textcolor{Chapter }{NumberSyllables}}
\logpage{[ 37, 5, 1 ]}\nobreak
\hyperdef{L}{X842D0B547CE93CF2}{}
{\noindent\textcolor{FuncColor}{$\triangleright$\enspace\texttt{NumberSyllables({\mdseries\slshape w})\index{NumberSyllables@\texttt{NumberSyllables}}
\label{NumberSyllables}
}\hfill{\scriptsize (attribute)}}\\


 \texttt{NumberSyllables} returns the number of syllables of the associative word \mbox{\texttt{\mdseries\slshape w}}. }

 

\subsection{\textcolor{Chapter }{ExponentSyllable}}
\logpage{[ 37, 5, 2 ]}\nobreak
\hyperdef{L}{X7E91575F848F4526}{}
{\noindent\textcolor{FuncColor}{$\triangleright$\enspace\texttt{ExponentSyllable({\mdseries\slshape w, i})\index{ExponentSyllable@\texttt{ExponentSyllable}}
\label{ExponentSyllable}
}\hfill{\scriptsize (operation)}}\\


 \texttt{ExponentSyllable} returns the exponent of the \mbox{\texttt{\mdseries\slshape i}}\texttt{\symbol{45}}th syllable of the associative word \mbox{\texttt{\mdseries\slshape w}}. }

 

\subsection{\textcolor{Chapter }{GeneratorSyllable}}
\logpage{[ 37, 5, 3 ]}\nobreak
\hyperdef{L}{X7F2A8CFD811C73B1}{}
{\noindent\textcolor{FuncColor}{$\triangleright$\enspace\texttt{GeneratorSyllable({\mdseries\slshape w, i})\index{GeneratorSyllable@\texttt{GeneratorSyllable}}
\label{GeneratorSyllable}
}\hfill{\scriptsize (operation)}}\\


 \texttt{GeneratorSyllable} returns the number of the generator that is involved in the \mbox{\texttt{\mdseries\slshape i}}\texttt{\symbol{45}}th syllable of the associative word \mbox{\texttt{\mdseries\slshape w}}. }

 

\subsection{\textcolor{Chapter }{SubSyllables}}
\logpage{[ 37, 5, 4 ]}\nobreak
\hyperdef{L}{X7B4F7A167E844FA5}{}
{\noindent\textcolor{FuncColor}{$\triangleright$\enspace\texttt{SubSyllables({\mdseries\slshape w, from, to})\index{SubSyllables@\texttt{SubSyllables}}
\label{SubSyllables}
}\hfill{\scriptsize (operation)}}\\


 \texttt{SubSyllables} returns the subword of the associative word \mbox{\texttt{\mdseries\slshape w}} that consists of the syllables from positions \mbox{\texttt{\mdseries\slshape from}} to \mbox{\texttt{\mdseries\slshape to}}, where \mbox{\texttt{\mdseries\slshape from}} and \mbox{\texttt{\mdseries\slshape to}} must be positive integers, and indexing is done with origin 1. 
\begin{Verbatim}[commandchars=!@|,fontsize=\small,frame=single,label=Example]
  !gapprompt@gap>| !gapinput@w;  NumberSyllables( w );|
  a^5*b*a^2*b^-4*a
  5
  !gapprompt@gap>| !gapinput@ExponentSyllable( w, 3 );|
  2
  !gapprompt@gap>| !gapinput@GeneratorSyllable( w, 3 );|
  1
  !gapprompt@gap>| !gapinput@SubSyllables( w, 2, 3 );|
  b*a^2
\end{Verbatim}
 }

 }

  
\section{\textcolor{Chapter }{Representations for Associative Words}}\label{Representations for Associative Words}
\logpage{[ 37, 6, 0 ]}
\hyperdef{L}{X80A9F39582ED296E}{}
{
  \textsf{GAP} provides two different internal kinds of representations of associative words.
The first one are ``syllable representations'' in which words are stored in syllable (i.e. generator,exponent) form. (Older
versions of \textsf{GAP} only used this representation.) The second kind are ``letter representations'' in which each letter in a word is represented by its index number. Negative
numbers are used for inverses. Unless the syllable representation is specified
explicitly when creating the free group/monoid or semigroup, a letter
representation is used by default. 

 Depending on the task in mind, either of these two representations will
perform better in time or in memory use and algorithms that are syllable or
letter based (for example \texttt{GeneratorSyllable} (\ref{GeneratorSyllable}) and \texttt{Subword} (\ref{Subword})) perform substantially better in the corresponding representation. For
example when creating pc groups (see{\nobreakspace}\ref{Pc Groups}), it is advantageous to use a syllable representation while calculations in
free groups usually benefit from using a letter representation. 

\subsection{\textcolor{Chapter }{IsLetterAssocWordRep}}
\logpage{[ 37, 6, 1 ]}\nobreak
\hyperdef{L}{X7E3612247B3E241B}{}
{\noindent\textcolor{FuncColor}{$\triangleright$\enspace\texttt{IsLetterAssocWordRep({\mdseries\slshape obj})\index{IsLetterAssocWordRep@\texttt{IsLetterAssocWordRep}}
\label{IsLetterAssocWordRep}
}\hfill{\scriptsize (Representation)}}\\


 A word in letter representation stores a list of generator/inverses numbers
(as given by \texttt{LetterRepAssocWord} (\ref{LetterRepAssocWord})). Letter access is fast, syllable access is slow for such words. }

 

\subsection{\textcolor{Chapter }{IsLetterWordsFamily}}
\logpage{[ 37, 6, 2 ]}\nobreak
\hyperdef{L}{X7E36F7897D82417F}{}
{\noindent\textcolor{FuncColor}{$\triangleright$\enspace\texttt{IsLetterWordsFamily({\mdseries\slshape obj})\index{IsLetterWordsFamily@\texttt{IsLetterWordsFamily}}
\label{IsLetterWordsFamily}
}\hfill{\scriptsize (Category)}}\\


 A letter word family stores words by default in letter form. 

 Internally, there are letter representations that use integers (4 Byte) to
represent a generator and letter representations that use single bytes to
represent a character. The latter are more memory efficient, but can only be
used if there are less than 128 generators (in which case they are used by
default). }

 

\subsection{\textcolor{Chapter }{IsBLetterAssocWordRep}}
\logpage{[ 37, 6, 3 ]}\nobreak
\hyperdef{L}{X7C84789D7BB161E9}{}
{\noindent\textcolor{FuncColor}{$\triangleright$\enspace\texttt{IsBLetterAssocWordRep({\mdseries\slshape obj})\index{IsBLetterAssocWordRep@\texttt{IsBLetterAssocWordRep}}
\label{IsBLetterAssocWordRep}
}\hfill{\scriptsize (Representation)}}\\
\noindent\textcolor{FuncColor}{$\triangleright$\enspace\texttt{IsWLetterAssocWordRep({\mdseries\slshape obj})\index{IsWLetterAssocWordRep@\texttt{IsWLetterAssocWordRep}}
\label{IsWLetterAssocWordRep}
}\hfill{\scriptsize (Representation)}}\\


 these two subrepresentations of \texttt{IsLetterAssocWordRep} (\ref{IsLetterAssocWordRep}) indicate whether the word is stored as a list of bytes (in a string) or as a
list of integers). }

 

\subsection{\textcolor{Chapter }{IsBLetterWordsFamily}}
\logpage{[ 37, 6, 4 ]}\nobreak
\hyperdef{L}{X8719E7F27CDA1995}{}
{\noindent\textcolor{FuncColor}{$\triangleright$\enspace\texttt{IsBLetterWordsFamily({\mdseries\slshape obj})\index{IsBLetterWordsFamily@\texttt{IsBLetterWordsFamily}}
\label{IsBLetterWordsFamily}
}\hfill{\scriptsize (Category)}}\\
\noindent\textcolor{FuncColor}{$\triangleright$\enspace\texttt{IsWLetterWordsFamily({\mdseries\slshape obj})\index{IsWLetterWordsFamily@\texttt{IsWLetterWordsFamily}}
\label{IsWLetterWordsFamily}
}\hfill{\scriptsize (Category)}}\\


 These two subcategories of \texttt{IsLetterWordsFamily} (\ref{IsLetterWordsFamily}) specify the type of letter representation to be used. }

 

\subsection{\textcolor{Chapter }{IsSyllableAssocWordRep}}
\logpage{[ 37, 6, 5 ]}\nobreak
\hyperdef{L}{X7886F8BD83CD8081}{}
{\noindent\textcolor{FuncColor}{$\triangleright$\enspace\texttt{IsSyllableAssocWordRep({\mdseries\slshape obj})\index{IsSyllableAssocWordRep@\texttt{IsSyllableAssocWordRep}}
\label{IsSyllableAssocWordRep}
}\hfill{\scriptsize (Representation)}}\\


 A word in syllable representation stores generator/exponents pairs (as given
by \texttt{ExtRepOfObj} (\ref{ExtRepOfObj}). Syllable access is fast, letter access is slow for such words. }

 

\subsection{\textcolor{Chapter }{IsSyllableWordsFamily}}
\logpage{[ 37, 6, 6 ]}\nobreak
\hyperdef{L}{X7869716C84EA9D81}{}
{\noindent\textcolor{FuncColor}{$\triangleright$\enspace\texttt{IsSyllableWordsFamily({\mdseries\slshape obj})\index{IsSyllableWordsFamily@\texttt{IsSyllableWordsFamily}}
\label{IsSyllableWordsFamily}
}\hfill{\scriptsize (Category)}}\\


 A syllable word family stores words by default in syllable form. There are
also different versions of syllable representations, which compress a
generator exponent pair in 8, 16 or 32 bits or use a pair of integers.
Internal mechanisms try to make this as memory efficient as possible. }

 

\subsection{\textcolor{Chapter }{Is16BitsFamily}}
\logpage{[ 37, 6, 7 ]}\nobreak
\hyperdef{L}{X83F669828481FC32}{}
{\noindent\textcolor{FuncColor}{$\triangleright$\enspace\texttt{Is16BitsFamily({\mdseries\slshape obj})\index{Is16BitsFamily@\texttt{Is16BitsFamily}}
\label{Is16BitsFamily}
}\hfill{\scriptsize (Category)}}\\
\noindent\textcolor{FuncColor}{$\triangleright$\enspace\texttt{Is32BitsFamily({\mdseries\slshape obj})\index{Is32BitsFamily@\texttt{Is32BitsFamily}}
\label{Is32BitsFamily}
}\hfill{\scriptsize (Category)}}\\
\noindent\textcolor{FuncColor}{$\triangleright$\enspace\texttt{IsInfBitsFamily({\mdseries\slshape obj})\index{IsInfBitsFamily@\texttt{IsInfBitsFamily}}
\label{IsInfBitsFamily}
}\hfill{\scriptsize (Category)}}\\


 Regardless of the internal representation used, it is possible to convert a
word in a list of numbers in letter or syllable representation and vice versa. }

 

\subsection{\textcolor{Chapter }{LetterRepAssocWord}}
\logpage{[ 37, 6, 8 ]}\nobreak
\hyperdef{L}{X7BD7647C7B088389}{}
{\noindent\textcolor{FuncColor}{$\triangleright$\enspace\texttt{LetterRepAssocWord({\mdseries\slshape w[, gens]})\index{LetterRepAssocWord@\texttt{LetterRepAssocWord}}
\label{LetterRepAssocWord}
}\hfill{\scriptsize (operation)}}\\


 The \emph{letter representation} of an associated word is as a list of integers, each entry corresponding to a
group generator. Inverses of the generators are represented by negative
numbers. The generator numbers are as associated to the family. 

 This operation returns the letter representation of the associative word \mbox{\texttt{\mdseries\slshape w}}. 

 In the call with two arguments, the generator numbers correspond to the
generator order given in the list \mbox{\texttt{\mdseries\slshape gens}}. 

 (For words stored in syllable form the letter representation has to be
computed.) }

 

\subsection{\textcolor{Chapter }{AssocWordByLetterRep}}
\logpage{[ 37, 6, 9 ]}\nobreak
\hyperdef{L}{X7AC8EC757CFB9A51}{}
{\noindent\textcolor{FuncColor}{$\triangleright$\enspace\texttt{AssocWordByLetterRep({\mdseries\slshape Fam, lrep[, gens]})\index{AssocWordByLetterRep@\texttt{AssocWordByLetterRep}}
\label{AssocWordByLetterRep}
}\hfill{\scriptsize (operation)}}\\


 takes a letter representation \mbox{\texttt{\mdseries\slshape lrep}} (see \texttt{LetterRepAssocWord} (\ref{LetterRepAssocWord})) and returns an associative word in family \mbox{\texttt{\mdseries\slshape fam}} corresponding to this letter representation. 

 If \mbox{\texttt{\mdseries\slshape gens}} is given, the numbers in the letter representation correspond to \mbox{\texttt{\mdseries\slshape gens}}. 
\begin{Verbatim}[commandchars=!@|,fontsize=\small,frame=single,label=Example]
  !gapprompt@gap>| !gapinput@w:=AssocWordByLetterRep( FamilyObj(a), [-1,2,1,-2,-2,-2,1,1,1,1]);|
  a^-1*b*a*b^-3*a^4
  !gapprompt@gap>| !gapinput@LetterRepAssocWord( w^2 );|
  [ -1, 2, 1, -2, -2, -2, 1, 1, 1, 2, 1, -2, -2, -2, 1, 1, 1, 1 ]
\end{Verbatim}
 

 The external representation (see section{\nobreakspace}\ref{The External Representation for Associative Words}) can be used if a syllable representation is needed. }

 }

  
\section{\textcolor{Chapter }{The External Representation for Associative Words}}\label{The External Representation for Associative Words}
\logpage{[ 37, 7, 0 ]}
\hyperdef{L}{X7934D3D5797102EC}{}
{
  The external representation of the associative word $w$ is defined as follows. If $w = g_{{i_1}}^{{e_1}} * g_{{i_2}}^{{e_2}} * \cdots * g_{{i_k}}^{{e_k}}$ is a word over the alphabet $g_1, g_2, \ldots$, i.e., $g_i$ denotes the $i$\texttt{\symbol{45}}th generator of the family of $w$, then $w$ has external representation $[ i_1, e_1, i_2, e_2, \ldots, i_k, e_k ]$. The empty list describes the identity element (if exists) of the family.
Exponents may be negative if the family allows inverses. The external
representation of an associative word is guaranteed to be freely reduced; for
example, $g_1 * g_2 * g_2^{{-1}} * g_1$ has the external representation \texttt{[ 1, 2 ]}. 

 Regardless of the family preference for letter or syllable representations
(see{\nobreakspace}\ref{Representations for Associative Words}), \texttt{ExtRepOfObj} and \texttt{ObjByExtRep} can be used and interface to this ``syllable''\texttt{\symbol{45}}like representation. 

 
\begin{Verbatim}[commandchars=!@|,fontsize=\small,frame=single,label=Example]
  !gapprompt@gap>| !gapinput@w:= ObjByExtRep( FamilyObj(a), [1,5,2,-7,1,3,2,4,1,-2] );|
  a^5*b^-7*a^3*b^4*a^-2
  !gapprompt@gap>| !gapinput@ExtRepOfObj( w^2 );|
  [ 1, 5, 2, -7, 1, 3, 2, 4, 1, 3, 2, -7, 1, 3, 2, 4, 1, -2 ]
\end{Verbatim}
 }

  
\section{\textcolor{Chapter }{Straight Line Programs}}\label{Straight Line Programs}
\logpage{[ 37, 8, 0 ]}
\hyperdef{L}{X7DC99E4284093FBB}{}
{
  \emph{Straight line programs} describe an efficient way for evaluating an abstract word at concrete
generators, in a more efficient way than with \texttt{MappedWord} (\ref{MappedWord}). For example, the associative word $ababbab$ of length $7$ can be computed from the generators $a$, $b$ with only four multiplications, by first computing $c = ab$, then $d = cb$, and then $cdc$; Alternatively, one can compute $c = ab$, $e = bc$, and $aee$. In each step of these computations, one forms words in terms of the words
computed in the previous steps. 

 A straight line program in \textsf{GAP} is represented by an object in the category \texttt{IsStraightLineProgram} (\ref{IsStraightLineProgram})) that stores a list of ``lines'' each of which has one of the following three forms. 
\begin{enumerate}
\item  a nonempty dense list $l$ of integers, 
\item  a pair $[ l, i ]$ where $l$ is a list of form 1. and $i$ is a positive integer, 
\item  a list $[ l_1, l_2, \ldots, l_k ]$ where each $l_i$ is a list of form 1.; this may occur only for the last line of the program. 
\end{enumerate}
 

 The lists of integers that occur are interpreted as external representations
of associative words (see Section{\nobreakspace} \ref{The External Representation for Associative Words}); for example, the list $[ 1, 3, 2, -1 ]$ represents the word $g_1^3 g_2^{{-1}}$, with $g_1$ and $g_2$ the first and second abstract generator, respectively. 

 For the meaning of the list of lines, see \texttt{ResultOfStraightLineProgram} (\ref{ResultOfStraightLineProgram}). 

 Straight line programs can be constructed using \texttt{StraightLineProgram} (\ref{StraightLineProgram:for a list of lines (and the number of generators)}). 

 Defining attributes for straight line programs are \texttt{NrInputsOfStraightLineProgram} (\ref{NrInputsOfStraightLineProgram}) and \texttt{LinesOfStraightLineProgram} (\ref{LinesOfStraightLineProgram}). Another operation for straight line programs is \texttt{ResultOfStraightLineProgram} (\ref{ResultOfStraightLineProgram}). 

 Special methods applicable to straight line programs are installed for the
operations \texttt{Display} (\ref{Display}), \texttt{IsInternallyConsistent} (\ref{IsInternallyConsistent}), \texttt{PrintObj} (\ref{PrintObj}), and \texttt{ViewObj} (\ref{ViewObj}). 

 For a straight line program \mbox{\texttt{\mdseries\slshape prog}}, the default \texttt{Display} (\ref{Display}) method prints the interpretation of \mbox{\texttt{\mdseries\slshape prog}} as a sequence of assignments of associative words; a record with components \texttt{gensnames} (with value a list of strings) and \texttt{listname} (a string) may be entered as second argument of \texttt{Display} (\ref{Display}), in this case these names are used, the default for \texttt{gensnames} is \texttt{[ g1, g2, }$\ldots$\texttt{ ]}, the default for \texttt{listname} is $r$. 

\subsection{\textcolor{Chapter }{IsStraightLineProgram}}
\logpage{[ 37, 8, 1 ]}\nobreak
\hyperdef{L}{X7F69FF3F7C6694CB}{}
{\noindent\textcolor{FuncColor}{$\triangleright$\enspace\texttt{IsStraightLineProgram({\mdseries\slshape obj})\index{IsStraightLineProgram@\texttt{IsStraightLineProgram}}
\label{IsStraightLineProgram}
}\hfill{\scriptsize (Category)}}\\


 Each straight line program in \textsf{GAP} lies in the category \texttt{IsStraightLineProgram}. }

 

\subsection{\textcolor{Chapter }{StraightLineProgram (for a list of lines (and the number of generators))}}
\logpage{[ 37, 8, 2 ]}\nobreak
\hyperdef{L}{X7AECA57280DA3195}{}
{\noindent\textcolor{FuncColor}{$\triangleright$\enspace\texttt{StraightLineProgram({\mdseries\slshape lines[, nrgens]})\index{StraightLineProgram@\texttt{StraightLineProgram}!for a list of lines (and the number of generators)}
\label{StraightLineProgram:for a list of lines (and the number of generators)}
}\hfill{\scriptsize (function)}}\\
\noindent\textcolor{FuncColor}{$\triangleright$\enspace\texttt{StraightLineProgram({\mdseries\slshape string, gens})\index{StraightLineProgram@\texttt{StraightLineProgram}!for a string and a list of generators names}
\label{StraightLineProgram:for a string and a list of generators names}
}\hfill{\scriptsize (function)}}\\
\noindent\textcolor{FuncColor}{$\triangleright$\enspace\texttt{StraightLineProgramNC({\mdseries\slshape lines[, nrgens]})\index{StraightLineProgramNC@\texttt{StraightLineProgramNC}!for a list of lines (and the number of generators)}
\label{StraightLineProgramNC:for a list of lines (and the number of generators)}
}\hfill{\scriptsize (function)}}\\
\noindent\textcolor{FuncColor}{$\triangleright$\enspace\texttt{StraightLineProgramNC({\mdseries\slshape string, gens})\index{StraightLineProgramNC@\texttt{StraightLineProgramNC}!for a string and a list of generators names}
\label{StraightLineProgramNC:for a string and a list of generators names}
}\hfill{\scriptsize (function)}}\\


 In the first form, \mbox{\texttt{\mdseries\slshape lines}} must be a nonempty list of lists that defines a unique straight line program
(see{\nobreakspace}\texttt{IsStraightLineProgram} (\ref{IsStraightLineProgram})); in this case \texttt{StraightLineProgram} returns this program, otherwise an error is signalled. The optional argument \mbox{\texttt{\mdseries\slshape nrgens}} specifies the number of input generators of the program; if a line of form 1.
(that is, a list of integers) occurs in \mbox{\texttt{\mdseries\slshape lines}} except in the last position, this number is not determined by \mbox{\texttt{\mdseries\slshape lines}} and therefore \emph{must} be specified by the argument \mbox{\texttt{\mdseries\slshape nrgens}}; if not then \texttt{StraightLineProgram} returns \texttt{fail}. 

 In the second form, \mbox{\texttt{\mdseries\slshape string}} must be a nonempty string describing an arithmetic expression in terms of the
strings in the list \mbox{\texttt{\mdseries\slshape gens}}, where multiplication is denoted by concatenation, powering is denoted by \texttt{\texttt{\symbol{94}}}, and round brackets \texttt{(}, \texttt{)} may be used. Each entry in \mbox{\texttt{\mdseries\slshape gens}} must consist only of uppercase or lowercase letters (i.e., letters in \texttt{IsAlphaChar} (\ref{IsAlphaChar})) such that no entry is an initial part of another one. Called with this
input, \texttt{StraightLineProgram} returns a straight line program that evaluates to the word corresponding to \mbox{\texttt{\mdseries\slshape string}} when called with generators corresponding to \mbox{\texttt{\mdseries\slshape gens}}. 

 The \texttt{NC} variant does the same, except that the internal consistency of the program is
not checked. }

 

\subsection{\textcolor{Chapter }{LinesOfStraightLineProgram}}
\logpage{[ 37, 8, 3 ]}\nobreak
\hyperdef{L}{X81A8AFC47F8E4B91}{}
{\noindent\textcolor{FuncColor}{$\triangleright$\enspace\texttt{LinesOfStraightLineProgram({\mdseries\slshape prog})\index{LinesOfStraightLineProgram@\texttt{LinesOfStraightLineProgram}}
\label{LinesOfStraightLineProgram}
}\hfill{\scriptsize (attribute)}}\\


 For a straight line program \mbox{\texttt{\mdseries\slshape prog}}, \texttt{LinesOfStraightLineProgram} returns the list of program lines. There is no default method to compute these
lines if they are not stored. }

 

\subsection{\textcolor{Chapter }{NrInputsOfStraightLineProgram}}
\logpage{[ 37, 8, 4 ]}\nobreak
\hyperdef{L}{X820A592881D57802}{}
{\noindent\textcolor{FuncColor}{$\triangleright$\enspace\texttt{NrInputsOfStraightLineProgram({\mdseries\slshape prog})\index{NrInputsOfStraightLineProgram@\texttt{NrInputsOfStraightLineProgram}}
\label{NrInputsOfStraightLineProgram}
}\hfill{\scriptsize (attribute)}}\\


 For a straight line program \mbox{\texttt{\mdseries\slshape prog}}, \texttt{NrInputsOfStraightLineProgram} returns the number of generators that are needed as input. 

 If a line of form 1. (that is, a list of integers) occurs in the lines of \mbox{\texttt{\mdseries\slshape prog}} except the last line then the number of generators is not determined by the
lines, and must be set in the construction of the straight line program
(see{\nobreakspace}\texttt{StraightLineProgram} (\ref{StraightLineProgram:for a list of lines (and the number of generators)})). So if \mbox{\texttt{\mdseries\slshape prog}} contains a line of form 1. other than the last line and does \emph{not} store the number of generators then \texttt{NrInputsOfStraightLineProgram} signals an error. }

 

\subsection{\textcolor{Chapter }{ResultOfStraightLineProgram}}
\logpage{[ 37, 8, 5 ]}\nobreak
\hyperdef{L}{X7847D32B863E822F}{}
{\noindent\textcolor{FuncColor}{$\triangleright$\enspace\texttt{ResultOfStraightLineProgram({\mdseries\slshape prog, gens})\index{ResultOfStraightLineProgram@\texttt{ResultOfStraightLineProgram}}
\label{ResultOfStraightLineProgram}
}\hfill{\scriptsize (operation)}}\\


 \texttt{ResultOfStraightLineProgram} evaluates the straight line program (see{\nobreakspace}\texttt{IsStraightLineProgram} (\ref{IsStraightLineProgram})) \mbox{\texttt{\mdseries\slshape prog}} at the group elements in the list \mbox{\texttt{\mdseries\slshape gens}}. 

 The \emph{result} of a straight line program with lines $p_1, p_2, \ldots, p_k$ when applied to \mbox{\texttt{\mdseries\slshape gens}} is defined as follows. 
\begin{description}
\item[{(a)}]  First a list $r$ of intermediate results is initialized with a shallow copy of \mbox{\texttt{\mdseries\slshape gens}}. 
\item[{(b)}]  For $i < k$, before the $i$\texttt{\symbol{45}}th step, let $r$ be of length $n$. If $p_i$ is the external representation of an associative word in the first $n$ generators then the image of this word under the homomorphism that is given by
mapping $r$ to these first $n$ generators is added to $r$; if $p_i$ is a pair $[ l, j ]$, for a list $l$, then the same element is computed, but instead of being added to $r$, it replaces the $j$\texttt{\symbol{45}}th entry of $r$. 
\item[{(c)}]  For $i = k$, if $p_k$ is the external representation of an associative word then the element
described in (b) is the result of the program, if $p_k$ is a pair $[ l, j ]$, for a list $l$, then the result is the element described by $l$, and if $p_k$ is a list $[ l_1, l_2, \ldots, l_k ]$ of lists then the result is a list of group elements, where each $l_i$ is treated as in (b). 
\end{description}
 

 
\begin{Verbatim}[commandchars=!@|,fontsize=\small,frame=single,label=Example]
  !gapprompt@gap>| !gapinput@f:= FreeGroup( "x", "y" );;  gens:= GeneratorsOfGroup( f );;|
  !gapprompt@gap>| !gapinput@x:= gens[1];;  y:= gens[2];;|
  !gapprompt@gap>| !gapinput@prg:= StraightLineProgram( [ [] ] );|
  <straight line program>
  !gapprompt@gap>| !gapinput@ResultOfStraightLineProgram( prg, [] );|
  [  ]
\end{Verbatim}
 The above straight line program \texttt{prg} returns {\textendash}for \emph{any} list of input generators{\textendash} an empty list. 
\begin{Verbatim}[commandchars=!@|,fontsize=\small,frame=single,label=Example]
  !gapprompt@gap>| !gapinput@StraightLineProgram( [ [1,2,2,3], [3,-1] ] );|
  fail
  !gapprompt@gap>| !gapinput@prg:= StraightLineProgram( [ [1,2,2,3], [3,-1] ], 2 );|
  <straight line program>
  !gapprompt@gap>| !gapinput@LinesOfStraightLineProgram( prg );|
  [ [ 1, 2, 2, 3 ], [ 3, -1 ] ]
  !gapprompt@gap>| !gapinput@prg:= StraightLineProgram( "(a^2b^3)^-1", [ "a", "b" ] );|
  <straight line program>
  !gapprompt@gap>| !gapinput@LinesOfStraightLineProgram( prg );|
  [ [ [ 1, 2, 2, 3 ], 3 ], [ [ 3, -1 ], 4 ] ]
  !gapprompt@gap>| !gapinput@res:= ResultOfStraightLineProgram( prg, gens );|
  y^-3*x^-2
  !gapprompt@gap>| !gapinput@res = (x^2 * y^3)^-1;|
  true
  !gapprompt@gap>| !gapinput@NrInputsOfStraightLineProgram( prg );|
  2
  !gapprompt@gap>| !gapinput@Print( prg, "\n" );|
  StraightLineProgram( [ [ [ 1, 2, 2, 3 ], 3 ], [ [ 3, -1 ], 4 ] ], 2 )
  !gapprompt@gap>| !gapinput@Display( prg );|
  # input:
  r:= [ g1, g2 ];
  # program:
  r[3]:= r[1]^2*r[2]^3;
  r[4]:= r[3]^-1;
  # return value:
  r[4]
  !gapprompt@gap>| !gapinput@IsInternallyConsistent( StraightLineProgramNC( [ [1,2] ] ) );|
  true
  !gapprompt@gap>| !gapinput@IsInternallyConsistent( StraightLineProgramNC( [ [1,2,3] ] ) );|
  false
  !gapprompt@gap>| !gapinput@prg1:= StraightLineProgram( [ [1,1,2,2], [3,3,1,1] ], 2 );;|
  !gapprompt@gap>| !gapinput@prg2:= StraightLineProgram( [ [ [1,1,2,2], 2 ], [2,3,1,1] ] );;|
  !gapprompt@gap>| !gapinput@res1:= ResultOfStraightLineProgram( prg1, gens );|
  (x*y^2)^3*x
  !gapprompt@gap>| !gapinput@res1 = (x*y^2)^3*x;|
  true
  !gapprompt@gap>| !gapinput@res2:= ResultOfStraightLineProgram( prg2, gens );|
  (x*y^2)^3*x
  !gapprompt@gap>| !gapinput@res2 = (x*y^2)^3*x;|
  true
  !gapprompt@gap>| !gapinput@prg:= StraightLineProgram( [ [2,3], [ [3,1,1,4], [1,2,3,1] ] ], 2 );;|
  !gapprompt@gap>| !gapinput@res:= ResultOfStraightLineProgram( prg, gens );|
  [ y^3*x^4, x^2*y^3 ]
\end{Verbatim}
 }

 \index{LaTeX!for the result of a straight line program} 

\subsection{\textcolor{Chapter }{StringOfResultOfStraightLineProgram}}
\logpage{[ 37, 8, 6 ]}\nobreak
\hyperdef{L}{X8098EAAF7D344466}{}
{\noindent\textcolor{FuncColor}{$\triangleright$\enspace\texttt{StringOfResultOfStraightLineProgram({\mdseries\slshape prog, gensnames[, "LaTeX"]})\index{StringOfResultOfStraightLineProgram@\texttt{StringOfResultOfStraightLineProgram}}
\label{StringOfResultOfStraightLineProgram}
}\hfill{\scriptsize (function)}}\\


 \texttt{StringOfResultOfStraightLineProgram} returns a string that describes the result of the straight line program
(see{\nobreakspace}\texttt{IsStraightLineProgram} (\ref{IsStraightLineProgram})) \mbox{\texttt{\mdseries\slshape prog}} as word(s) in terms of the strings in the list \mbox{\texttt{\mdseries\slshape gensnames}}. If the result of \mbox{\texttt{\mdseries\slshape prog}} is a single element then the return value of \texttt{StringOfResultOfStraightLineProgram} is a string consisting of the entries of \mbox{\texttt{\mdseries\slshape gensnames}}, opening and closing brackets \texttt{(} and \texttt{)}, and powering by integers via \texttt{\texttt{\symbol{94}}}. If the result of \mbox{\texttt{\mdseries\slshape prog}} is a list of elements then the return value of \texttt{StringOfResultOfStraightLineProgram} is a comma separated concatenation of the strings of the single elements,
enclosed in square brackets \texttt{[}, \texttt{]}. 
\begin{Verbatim}[commandchars=!@|,fontsize=\small,frame=single,label=Example]
  !gapprompt@gap>| !gapinput@prg:= StraightLineProgram( [ [ 1, 2, 2, 3 ], [ 3, -1 ] ], 2 );;|
  !gapprompt@gap>| !gapinput@StringOfResultOfStraightLineProgram( prg, [ "a", "b" ] );|
  "(a^2b^3)^-1"
  !gapprompt@gap>| !gapinput@StringOfResultOfStraightLineProgram( prg, [ "a", "b" ], "LaTeX" );|
  "(a^{2}b^{3})^{-1}"
\end{Verbatim}
 }

 

\subsection{\textcolor{Chapter }{CompositionOfStraightLinePrograms}}
\logpage{[ 37, 8, 7 ]}\nobreak
\hyperdef{L}{X8274C7948248C053}{}
{\noindent\textcolor{FuncColor}{$\triangleright$\enspace\texttt{CompositionOfStraightLinePrograms({\mdseries\slshape prog2, prog1})\index{CompositionOfStraightLinePrograms@\texttt{CompositionOfStraightLinePrograms}}
\label{CompositionOfStraightLinePrograms}
}\hfill{\scriptsize (function)}}\\


 For two straight line programs \mbox{\texttt{\mdseries\slshape prog1}} and \mbox{\texttt{\mdseries\slshape prog2}}, \texttt{CompositionOfStraightLinePrograms} returns a straight line program \mbox{\texttt{\mdseries\slshape prog}} with the properties that \mbox{\texttt{\mdseries\slshape prog1}} and \mbox{\texttt{\mdseries\slshape prog}} have the same number of inputs, and the result of \mbox{\texttt{\mdseries\slshape prog}} when applied to given generators \mbox{\texttt{\mdseries\slshape gens}} equals the result of \mbox{\texttt{\mdseries\slshape prog2}} when this is applied to the output of \mbox{\texttt{\mdseries\slshape prog1}} applied to \mbox{\texttt{\mdseries\slshape gens}}. 

 (Of course the number of outputs of \mbox{\texttt{\mdseries\slshape prog1}} must be the same as the number of inputs of \mbox{\texttt{\mdseries\slshape prog2}}.) 
\begin{Verbatim}[commandchars=!@|,fontsize=\small,frame=single,label=Example]
  !gapprompt@gap>| !gapinput@prg1:= StraightLineProgram( "a^2b", [ "a","b" ] );;|
  !gapprompt@gap>| !gapinput@prg2:= StraightLineProgram( "c^5", [ "c" ] );;|
  !gapprompt@gap>| !gapinput@comp:= CompositionOfStraightLinePrograms( prg2, prg1 );|
  <straight line program>
  !gapprompt@gap>| !gapinput@StringOfResultOfStraightLineProgram( comp, [ "a", "b" ] );|
  "(a^2b)^5"
  !gapprompt@gap>| !gapinput@prg:= StraightLineProgram( [ [2,3], [ [3,1,1,4], [1,2,3,1] ] ], 2 );;|
  !gapprompt@gap>| !gapinput@StringOfResultOfStraightLineProgram( prg, [ "a", "b" ] );|
  "[ b^3a^4, a^2b^3 ]"
  !gapprompt@gap>| !gapinput@comp:= CompositionOfStraightLinePrograms( prg, prg );|
  <straight line program>
  !gapprompt@gap>| !gapinput@StringOfResultOfStraightLineProgram( comp, [ "a", "b" ] );|
  "[ (a^2b^3)^3(b^3a^4)^4, (b^3a^4)^2(a^2b^3)^3 ]"
\end{Verbatim}
 }

 

\subsection{\textcolor{Chapter }{IntegratedStraightLineProgram}}
\logpage{[ 37, 8, 8 ]}\nobreak
\hyperdef{L}{X7A582FA97C786640}{}
{\noindent\textcolor{FuncColor}{$\triangleright$\enspace\texttt{IntegratedStraightLineProgram({\mdseries\slshape listofprogs})\index{IntegratedStraightLineProgram@\texttt{IntegratedStraightLineProgram}}
\label{IntegratedStraightLineProgram}
}\hfill{\scriptsize (function)}}\\


 For a nonempty dense list \mbox{\texttt{\mdseries\slshape listofprogs}} of straight line programs $p_1, p_2, \ldots, p_m$ that have the same number $n$ of inputs (see{\nobreakspace}\texttt{NrInputsOfStraightLineProgram} (\ref{NrInputsOfStraightLineProgram})), \texttt{IntegratedStraightLineProgram} returns a straight line program $prog$ with $n$ inputs such that for each $n$\texttt{\symbol{45}}tuple $gens$ of generators, \texttt{ResultOfStraightLineProgram( }$prog, gens$\texttt{ )} is the concatenation of the lists $r_1, r_2, \ldots, r_m$, where $r_i$ is equal to \texttt{ResultOfStraightLineProgram( }$p_i, gens$\texttt{ )} if this result is a list of elements, and otherwise $r_i$ is equal to the list of length one that contains this result. 
\begin{Verbatim}[commandchars=!@|,fontsize=\small,frame=single,label=Example]
  !gapprompt@gap>| !gapinput@f:= FreeGroup( "x", "y" );;  gens:= GeneratorsOfGroup( f );;|
  !gapprompt@gap>| !gapinput@prg1:= StraightLineProgram([ [ [ 1, 2 ], 1 ], [ 1, 2, 2, -1 ] ], 2);;|
  !gapprompt@gap>| !gapinput@prg2:= StraightLineProgram([ [ [ 2, 2 ], 3 ], [ 1, 3, 3, 2 ] ], 2);;|
  !gapprompt@gap>| !gapinput@prg3:= StraightLineProgram([ [ 2, 2 ], [ 1, 3, 3, 2 ] ], 2);;|
  !gapprompt@gap>| !gapinput@prg:= IntegratedStraightLineProgram( [ prg1, prg2, prg3 ] );;|
  !gapprompt@gap>| !gapinput@ResultOfStraightLineProgram( prg, gens );|
  [ x^4*y^-1, x^3*y^4, x^3*y^4 ]
  !gapprompt@gap>| !gapinput@prg:= IntegratedStraightLineProgram( [ prg2, prg3, prg1 ] );;|
  !gapprompt@gap>| !gapinput@ResultOfStraightLineProgram( prg, gens );|
  [ x^3*y^4, x^3*y^4, x^4*y^-1 ]
  !gapprompt@gap>| !gapinput@prg:= IntegratedStraightLineProgram( [ prg3, prg1, prg2 ] );;|
  !gapprompt@gap>| !gapinput@ResultOfStraightLineProgram( prg, gens );|
  [ x^3*y^4, x^4*y^-1, x^3*y^4 ]
  !gapprompt@gap>| !gapinput@prg:= IntegratedStraightLineProgram( [ prg, prg ] );;|
  !gapprompt@gap>| !gapinput@ResultOfStraightLineProgram( prg, gens );|
  [ x^3*y^4, x^4*y^-1, x^3*y^4, x^3*y^4, x^4*y^-1, x^3*y^4 ]
\end{Verbatim}
 }

 

\subsection{\textcolor{Chapter }{RestrictOutputsOfSLP}}
\logpage{[ 37, 8, 9 ]}\nobreak
\hyperdef{L}{X7C9CABD17BE4850F}{}
{\noindent\textcolor{FuncColor}{$\triangleright$\enspace\texttt{RestrictOutputsOfSLP({\mdseries\slshape slp, k})\index{RestrictOutputsOfSLP@\texttt{RestrictOutputsOfSLP}}
\label{RestrictOutputsOfSLP}
}\hfill{\scriptsize (function)}}\\


 \mbox{\texttt{\mdseries\slshape slp}} must be a straight line program returning a tuple of values. This function
returns a new slp that calculates only those outputs specified by \mbox{\texttt{\mdseries\slshape k}}. The argument \mbox{\texttt{\mdseries\slshape k}} may be an integer or a list of integers. If \mbox{\texttt{\mdseries\slshape k}} is an integer, the resulting slp calculates only the result with that number
in the original output tuple. If \mbox{\texttt{\mdseries\slshape k}} is a list of integers, the resulting slp calculates those results with indices \mbox{\texttt{\mdseries\slshape k}} in the original output tuple. In both cases the resulting slp does only what
is necessary. Obviously, the slp must have a line with enough expressions
(lists) for the supplied \mbox{\texttt{\mdseries\slshape k}} as its last line. \mbox{\texttt{\mdseries\slshape slp}} is either an slp or a pair where the first entry are the lines of the slp and
the second is the number of inputs. }

 

\subsection{\textcolor{Chapter }{IntermediateResultOfSLP}}
\logpage{[ 37, 8, 10 ]}\nobreak
\hyperdef{L}{X7EF202F17DCA5D1C}{}
{\noindent\textcolor{FuncColor}{$\triangleright$\enspace\texttt{IntermediateResultOfSLP({\mdseries\slshape slp, k})\index{IntermediateResultOfSLP@\texttt{IntermediateResultOfSLP}}
\label{IntermediateResultOfSLP}
}\hfill{\scriptsize (function)}}\\


 Returns a new slp that calculates only the value of slot \mbox{\texttt{\mdseries\slshape k}} at the end of \mbox{\texttt{\mdseries\slshape slp}} doing only what is necessary. slp is either an slp or a pair where the first
entry are the lines of the slp and the second is the number of inputs. Note
that this assumes a general SLP with possible overwriting. If you know that
your SLP does not overwrite slots, please use \texttt{IntermediateResultOfSLPWithoutOverwrite} (\ref{IntermediateResultOfSLPWithoutOverwrite}), which is much faster in this case. }

 

\subsection{\textcolor{Chapter }{IntermediateResultOfSLPWithoutOverwrite}}
\logpage{[ 37, 8, 11 ]}\nobreak
\hyperdef{L}{X8085CF79856B2889}{}
{\noindent\textcolor{FuncColor}{$\triangleright$\enspace\texttt{IntermediateResultOfSLPWithoutOverwrite({\mdseries\slshape slp, k})\index{IntermediateResultOfSLPWithoutOverwrite@\texttt{Intermediate}\-\texttt{Result}\-\texttt{Of}\-\texttt{S}\-\texttt{L}\-\texttt{P}\-\texttt{Without}\-\texttt{Overwrite}}
\label{IntermediateResultOfSLPWithoutOverwrite}
}\hfill{\scriptsize (function)}}\\


 Returns a new slp that calculates only the value of slot \mbox{\texttt{\mdseries\slshape k}}, which must be an integer. Note that \mbox{\texttt{\mdseries\slshape slp}} must not overwrite slots but only append!!! Use \texttt{IntermediateResultOfSLP} (\ref{IntermediateResultOfSLP}) in the other case! \mbox{\texttt{\mdseries\slshape slp}} is either an slp or a pair where the first entry is the list of lines of the
slp and the second is the number of its inputs. }

 

\subsection{\textcolor{Chapter }{IntermediateResultsOfSLPWithoutOverwrite}}
\logpage{[ 37, 8, 12 ]}\nobreak
\hyperdef{L}{X873244F37FAA717A}{}
{\noindent\textcolor{FuncColor}{$\triangleright$\enspace\texttt{IntermediateResultsOfSLPWithoutOverwrite({\mdseries\slshape slp, k})\index{IntermediateResultsOfSLPWithoutOverwrite@\texttt{Intermediate}\-\texttt{Results}\-\texttt{Of}\-\texttt{S}\-\texttt{L}\-\texttt{P}\-\texttt{Without}\-\texttt{Overwrite}}
\label{IntermediateResultsOfSLPWithoutOverwrite}
}\hfill{\scriptsize (function)}}\\


 Returns a new slp that calculates only the values of slots contained in the
list \mbox{\texttt{\mdseries\slshape k}}. Note that \mbox{\texttt{\mdseries\slshape slp}} must not overwrite slots but only append!!! Use \texttt{IntermediateResultOfSLP} (\ref{IntermediateResultOfSLP}) in the other case! \mbox{\texttt{\mdseries\slshape slp}} is either a slp or a pair where the first entry is the list of lines of the
slp and the second is the number of its inputs. }

 

\subsection{\textcolor{Chapter }{ProductOfStraightLinePrograms}}
\logpage{[ 37, 8, 13 ]}\nobreak
\hyperdef{L}{X837101F982C35035}{}
{\noindent\textcolor{FuncColor}{$\triangleright$\enspace\texttt{ProductOfStraightLinePrograms({\mdseries\slshape s1, s2})\index{ProductOfStraightLinePrograms@\texttt{ProductOfStraightLinePrograms}}
\label{ProductOfStraightLinePrograms}
}\hfill{\scriptsize (function)}}\\


 \mbox{\texttt{\mdseries\slshape s1}} and \mbox{\texttt{\mdseries\slshape s2}} must be two slps that return a single element with the same number of inputs.
This function constructs an slp that returns the product of the two results
the slps \mbox{\texttt{\mdseries\slshape s1}} and \mbox{\texttt{\mdseries\slshape s2}} would produce with the same input. }

 

\subsection{\textcolor{Chapter }{SlotUsagePattern}}
\logpage{[ 37, 8, 14 ]}\nobreak
\hyperdef{L}{X84C83CE98194FD03}{}
{\noindent\textcolor{FuncColor}{$\triangleright$\enspace\texttt{SlotUsagePattern({\mdseries\slshape s})\index{SlotUsagePattern@\texttt{SlotUsagePattern}}
\label{SlotUsagePattern}
}\hfill{\scriptsize (attribute)}}\\


 Analyses the straight line program \mbox{\texttt{\mdseries\slshape s}} for more efficient evaluation. This means in particular two things, when this
attribute is known: First of all, intermediate results which are not actually
needed later on are not computed at all, and once an intermediate result is
used for the last time in this SLP, it is discarded. The latter leads to the
fact that the evaluation of the SLP needs less memory. }

 }

  
\section{\textcolor{Chapter }{Straight Line Program Elements}}\label{Straight Line Program Elements}
\logpage{[ 37, 9, 0 ]}
\hyperdef{L}{X8188799182D82A92}{}
{
  When computing with very large (in terms of memory) elements, for example
permutations of degree a few hundred thousands, it can be helpful (in terms of
memory usage) to represent them via straight line programs in terms of an
original generator set. (So every element takes only small extra storage for
the straight line program.) 

 A straight line program element has a \emph{seed} (a list of group elements) and a straight line program on the same number of
generators as the length of this seed, its value is the value of the evaluated
straight line program. 

 At the moment, the entries of the straight line program have to be simple
lists (i.e. of the first form). 

 Straight line program elements are in the same categories and families as the
elements of the seed, so they should work together with existing algorithms. 

 Note however, that due to the different way of storage some normally very
cheap operations (such as testing for element equality) can become more
expensive when dealing with straight line program elements. This is
essentially the tradeoff for using less memory. 

 See also Section{\nobreakspace}\ref{Working with large degree permutation groups}. 

\subsection{\textcolor{Chapter }{IsStraightLineProgElm}}
\logpage{[ 37, 9, 1 ]}\nobreak
\hyperdef{L}{X85A5838482944FA5}{}
{\noindent\textcolor{FuncColor}{$\triangleright$\enspace\texttt{IsStraightLineProgElm({\mdseries\slshape obj})\index{IsStraightLineProgElm@\texttt{IsStraightLineProgElm}}
\label{IsStraightLineProgElm}
}\hfill{\scriptsize (Representation)}}\\


 A straight line program element is a group element given (for memory reasons)
as a straight line program. Straight line program elements are positional
objects, the first component is a record with a component \texttt{seeds}, the second component the straight line program. }

 

\subsection{\textcolor{Chapter }{StraightLineProgElm}}
\logpage{[ 37, 9, 2 ]}\nobreak
\hyperdef{L}{X78889E5B7E1B3BFF}{}
{\noindent\textcolor{FuncColor}{$\triangleright$\enspace\texttt{StraightLineProgElm({\mdseries\slshape seed, prog})\index{StraightLineProgElm@\texttt{StraightLineProgElm}}
\label{StraightLineProgElm}
}\hfill{\scriptsize (function)}}\\


 Creates a straight line program element for seed \mbox{\texttt{\mdseries\slshape seed}} and program \mbox{\texttt{\mdseries\slshape prog}}. }

 

\subsection{\textcolor{Chapter }{StraightLineProgGens}}
\logpage{[ 37, 9, 3 ]}\nobreak
\hyperdef{L}{X81BC263A7E45E775}{}
{\noindent\textcolor{FuncColor}{$\triangleright$\enspace\texttt{StraightLineProgGens({\mdseries\slshape gens[, base]})\index{StraightLineProgGens@\texttt{StraightLineProgGens}}
\label{StraightLineProgGens}
}\hfill{\scriptsize (function)}}\\


 returns a set of straight line program elements corresponding to the
generators in \mbox{\texttt{\mdseries\slshape gens}}. If \mbox{\texttt{\mdseries\slshape gens}} is a set of permutations then \mbox{\texttt{\mdseries\slshape base}} can be given which must be a base for the group generated by \mbox{\texttt{\mdseries\slshape gens}}. (Such a base will be used to speed up equality tests.) }

 

\subsection{\textcolor{Chapter }{EvalStraightLineProgElm}}
\logpage{[ 37, 9, 4 ]}\nobreak
\hyperdef{L}{X7BEAE8AC809B27DC}{}
{\noindent\textcolor{FuncColor}{$\triangleright$\enspace\texttt{EvalStraightLineProgElm({\mdseries\slshape slpel})\index{EvalStraightLineProgElm@\texttt{EvalStraightLineProgElm}}
\label{EvalStraightLineProgElm}
}\hfill{\scriptsize (function)}}\\


 evaluates a straight line program element \mbox{\texttt{\mdseries\slshape slpel}} from its seeds. }

 

\subsection{\textcolor{Chapter }{StretchImportantSLPElement}}
\logpage{[ 37, 9, 5 ]}\nobreak
\hyperdef{L}{X7D85D1DF84DC68E3}{}
{\noindent\textcolor{FuncColor}{$\triangleright$\enspace\texttt{StretchImportantSLPElement({\mdseries\slshape elm})\index{StretchImportantSLPElement@\texttt{StretchImportantSLPElement}}
\label{StretchImportantSLPElement}
}\hfill{\scriptsize (operation)}}\\


 If \mbox{\texttt{\mdseries\slshape elm}} is a straight line program element whose straight line representation is very
long, this operation changes the representation of \mbox{\texttt{\mdseries\slshape elm}} to a straight line program element, equal to \mbox{\texttt{\mdseries\slshape elm}}, whose seed contains the evaluation of \mbox{\texttt{\mdseries\slshape elm}} and whose straight line program has length 1. 

 For other objects nothing happens. 

 This operation permits to designate ``important'' elements within an algorithm (elements that will be referred to often), which
will be represented by guaranteed short straight line program elements. 
\begin{Verbatim}[commandchars=!@A,fontsize=\small,frame=single,label=Example]
  !gapprompt@gap>A !gapinput@gens:=StraightLineProgGens([(1,2,3,4),(1,2)]);A
  [ <[ [ 2, 1 ] ]|(1,2,3,4)>, <[ [ 1, 1 ] ]|(1,2)> ]
  !gapprompt@gap>A !gapinput@g:=Group(gens);;A
  !gapprompt@gap>A !gapinput@(gens[1]^3)^gens[2];A
  <[ [ 1, -1, 2, 3, 1, 1 ] ]|(1,2,4,3)>
  !gapprompt@gap>A !gapinput@Size(g);A
  24
\end{Verbatim}
 }

 }

  
\section{\textcolor{Chapter }{Objects with Memory}}\label{Objects with Memory}
\logpage{[ 37, 10, 0 ]}
\hyperdef{L}{X877AAF977A6C1F09}{}
{
  The idea behind objects with memory is as follows. One starts with a list of
seed objects and creates corresponding new objects for them that store how
they arise from the seed objects (see \texttt{GeneratorsWithMemory} (\ref{GeneratorsWithMemory})). Each product, inverse, commutator etc. of the new objects then also stores
how it was obtained from the seed objects. One can use \texttt{SLPOfElm} (\ref{SLPOfElm}) to create a straight line program for a given element $x$ such that evaluating this program in the seed objects yields $x$. 

 
\begin{Verbatim}[commandchars=!@|,fontsize=\small,frame=single,label=Example]
  !gapprompt@gap>| !gapinput@seeds:= [ (1,2,3), (1,2) ];;|
  !gapprompt@gap>| !gapinput@elms:= GeneratorsWithMemory( seeds );|
  [ <(1,2,3) with mem>, <(1,2) with mem> ]
  !gapprompt@gap>| !gapinput@x:= (elms[1] * elms[2])^-1;|
  <(2,3) with mem>
  !gapprompt@gap>| !gapinput@xslp:= SLPOfElm( x );|
  <straight line program>
  !gapprompt@gap>| !gapinput@Display( xslp );|
  # input:
  r:= [ g1, g2 ];
  # program:
  r[3]:= r[1]*r[2];
  r[4]:= r[3]^-1;
  # return value:
  r[4]
  !gapprompt@gap>| !gapinput@ResultOfStraightLineProgram( xslp, seeds );|
  (2,3)
\end{Verbatim}
 

 The objects with memory are intended to behave equally to the underlying
objects without memory, apart from the behaviour concerning the memory
information. In particular, the tests for equality and ordering are delegated
to the underlying objects. 

 
\begin{Verbatim}[commandchars=!@|,fontsize=\small,frame=single,label=Example]
  !gapprompt@gap>| !gapinput@seeds = elms;|
  true
  !gapprompt@gap>| !gapinput@( seeds[1] < seeds[2] ) = ( elms[1] < elms[2] );|
  true
\end{Verbatim}
 

 It may, however, happen that some \textsf{GAP} function runs into an error when it is called with elements with memory. In
such cases, probably a method for elements with memory is missing, and perhaps
some existing methods are installed with too weak requirements. As a solution,
one can then add the missing method, or ask \textsf{GAP} support for help. 

 An overview of the currently supported functionality for objects with memory
can be computed with \texttt{MethodsForObjWithMemory} (\ref{MethodsForObjWithMemory}), this may be helpful for finding code that can be adapted. 

\subsection{\textcolor{Chapter }{IsObjWithMemory}}
\logpage{[ 37, 10, 1 ]}\nobreak
\hyperdef{L}{X7E39534E826C6A8E}{}
{\noindent\textcolor{FuncColor}{$\triangleright$\enspace\texttt{IsObjWithMemory({\mdseries\slshape obj})\index{IsObjWithMemory@\texttt{IsObjWithMemory}}
\label{IsObjWithMemory}
}\hfill{\scriptsize (Representation)}}\\


 The filter \texttt{IsObjWithMemory} describes objects that can be multiplied and store information how they arise
from seed objects via arithmetic operations. 

 
\begin{Verbatim}[commandchars=!@|,fontsize=\small,frame=single,label=Example]
  !gapprompt@gap>| !gapinput@seeds:= [ (1,2,3), (1,2) ];;|
  !gapprompt@gap>| !gapinput@elms:= GeneratorsWithMemory( seeds );;|
  !gapprompt@gap>| !gapinput@IsObjWithMemory( elms[1] );|
  true
  !gapprompt@gap>| !gapinput@IsObjWithMemory( seeds[1] );|
  false
\end{Verbatim}
 }

 

\subsection{\textcolor{Chapter }{GeneratorsWithMemory}}
\logpage{[ 37, 10, 2 ]}\nobreak
\hyperdef{L}{X8310FFAC82CD7652}{}
{\noindent\textcolor{FuncColor}{$\triangleright$\enspace\texttt{GeneratorsWithMemory({\mdseries\slshape list})\index{GeneratorsWithMemory@\texttt{GeneratorsWithMemory}}
\label{GeneratorsWithMemory}
}\hfill{\scriptsize (function)}}\\
\textbf{\indent Returns:\ }
a list of objects with memory



 For a list \mbox{\texttt{\mdseries\slshape list}} of multiplicative elements, \texttt{GeneratorsWithMemory} returns a list of corresponding objects with memory. 

 
\begin{Verbatim}[commandchars=!@|,fontsize=\small,frame=single,label=Example]
  !gapprompt@gap>| !gapinput@seeds:= [ (1,2,3), (1,2) ];;|
  !gapprompt@gap>| !gapinput@elms:= GeneratorsWithMemory( seeds );|
  [ <(1,2,3) with mem>, <(1,2) with mem> ]
  !gapprompt@gap>| !gapinput@seeds = elms;|
  true
  !gapprompt@gap>| !gapinput@ForAll( elms, IsObjWithMemory );|
  true
\end{Verbatim}
 }

 
\subsection{\textcolor{Chapter }{GroupWithMemory}}\logpage{[ 37, 10, 3 ]}
\hyperdef{L}{X7A73055D83FFB54A}{}
{
\noindent\textcolor{FuncColor}{$\triangleright$\enspace\texttt{GroupWithMemory({\mdseries\slshape G})\index{GroupWithMemory@\texttt{GroupWithMemory}!for a group}
\label{GroupWithMemory:for a group}
}\hfill{\scriptsize (function)}}\\
\noindent\textcolor{FuncColor}{$\triangleright$\enspace\texttt{GroupWithMemory({\mdseries\slshape gens})\index{GroupWithMemory@\texttt{GroupWithMemory}!for a list of generators}
\label{GroupWithMemory:for a list of generators}
}\hfill{\scriptsize (function)}}\\


 Return a group whose \texttt{GeneratorsOfGroup} (\ref{GeneratorsOfGroup}) value is the \texttt{GeneratorsWithMemory} (\ref{GeneratorsWithMemory}) value of \texttt{GeneratorsOfGroup( }\mbox{\texttt{\mdseries\slshape G}}\texttt{ )} or of \mbox{\texttt{\mdseries\slshape gens}}, respectively. 

 
\begin{Verbatim}[commandchars=!@|,fontsize=\small,frame=single,label=Example]
  !gapprompt@gap>| !gapinput@G:= GroupWithMemory( SymmetricGroup( 3 ) );|
  Group([ (1,2,3), (1,2) ])
  !gapprompt@gap>| !gapinput@IsObjWithMemory( One( G ) );|
  true
  !gapprompt@gap>| !gapinput@G:= GroupWithMemory( [ (1,2,3) ] );|
  Group([ (1,2,3) ])
  !gapprompt@gap>| !gapinput@IsObjWithMemory( One( G ) );|
  true
\end{Verbatim}
 }

 

\subsection{\textcolor{Chapter }{SLPOfElm}}
\logpage{[ 37, 10, 4 ]}\nobreak
\hyperdef{L}{X86EE81217C438D32}{}
{\noindent\textcolor{FuncColor}{$\triangleright$\enspace\texttt{SLPOfElm({\mdseries\slshape g})\index{SLPOfElm@\texttt{SLPOfElm}}
\label{SLPOfElm}
}\hfill{\scriptsize (function)}}\\


 For \mbox{\texttt{\mdseries\slshape g}} in the filter \texttt{IsObjWithMemory} (\ref{IsObjWithMemory}), return a straight line program that describes \mbox{\texttt{\mdseries\slshape g}} in terms of the generators that were used to initialize the memory. 

 
\begin{Verbatim}[commandchars=!@|,fontsize=\small,frame=single,label=Example]
  !gapprompt@gap>| !gapinput@seeds:= [ (1,2,3), (1,2) ];;|
  !gapprompt@gap>| !gapinput@elms:= GeneratorsWithMemory( seeds );|
  [ <(1,2,3) with mem>, <(1,2) with mem> ]
  !gapprompt@gap>| !gapinput@g:= (elms[2] * elms[1] * elms[2])^-1;|
  <(1,2,3) with mem>
  !gapprompt@gap>| !gapinput@slp:= SLPOfElm( g );|
  <straight line program>
  !gapprompt@gap>| !gapinput@g = ResultOfStraightLineProgram( slp, seeds );|
  true
\end{Verbatim}
 }

 

\subsection{\textcolor{Chapter }{SLPOfElms}}
\logpage{[ 37, 10, 5 ]}\nobreak
\hyperdef{L}{X7D3B81CE87D6D429}{}
{\noindent\textcolor{FuncColor}{$\triangleright$\enspace\texttt{SLPOfElms({\mdseries\slshape list})\index{SLPOfElms@\texttt{SLPOfElms}}
\label{SLPOfElms}
}\hfill{\scriptsize (function)}}\\


 For a list \mbox{\texttt{\mdseries\slshape list}} of elements in the filter \texttt{IsObjWithMemory} (\ref{IsObjWithMemory}), return a straight line program that describes \mbox{\texttt{\mdseries\slshape list}} in terms of the generators that were used to initialize the memory. 

 
\begin{Verbatim}[commandchars=!@|,fontsize=\small,frame=single,label=Example]
  !gapprompt@gap>| !gapinput@seeds:= [ (1,2,3), (1,2) ];;|
  !gapprompt@gap>| !gapinput@elms:= GeneratorsWithMemory( seeds );|
  [ <(1,2,3) with mem>, <(1,2) with mem> ]
  !gapprompt@gap>| !gapinput@list:= [ elms[2] * elms[1], elms[1] * elms[2] ];|
  [ <(1,3) with mem>, <(2,3) with mem> ]
  !gapprompt@gap>| !gapinput@slp:= SLPOfElms( list );|
  <straight line program>
  !gapprompt@gap>| !gapinput@list = ResultOfStraightLineProgram( slp, seeds );|
  true
\end{Verbatim}
 }

 

\subsection{\textcolor{Chapter }{StripMemory}}
\logpage{[ 37, 10, 6 ]}\nobreak
\hyperdef{L}{X783FCE658057D273}{}
{\noindent\textcolor{FuncColor}{$\triangleright$\enspace\texttt{StripMemory({\mdseries\slshape obj})\index{StripMemory@\texttt{StripMemory}}
\label{StripMemory}
}\hfill{\scriptsize (operation)}}\\


 For a list \mbox{\texttt{\mdseries\slshape obj}}, \texttt{StripMemory} returns a new list of the \texttt{StripMemory} results for the entries of \mbox{\texttt{\mdseries\slshape obj}}. 

 If \mbox{\texttt{\mdseries\slshape obj}} is not a list and in the filter \texttt{IsObjWithMemory} (\ref{IsObjWithMemory}), \texttt{StripMemory} returns the underlying object without memory. 

 Otherwise \mbox{\texttt{\mdseries\slshape obj}} is returned. 

 
\begin{Verbatim}[commandchars=!@|,fontsize=\small,frame=single,label=Example]
  !gapprompt@gap>| !gapinput@seeds:= [ (1,2,3), (1,2) ];;|
  !gapprompt@gap>| !gapinput@elms:= GeneratorsWithMemory( seeds );|
  [ <(1,2,3) with mem>, <(1,2) with mem> ]
  !gapprompt@gap>| !gapinput@StripMemory( elms );|
  [ (1,2,3), (1,2) ]
  !gapprompt@gap>| !gapinput@StripMemory( elms[1] );|
  (1,2,3)
  !gapprompt@gap>| !gapinput@IsIdenticalObj( seeds[1], StripMemory( seeds[1] ) );|
  true
\end{Verbatim}
 }

 

\subsection{\textcolor{Chapter }{MethodsForObjWithMemory}}
\logpage{[ 37, 10, 7 ]}\nobreak
\hyperdef{L}{X81A6A0FB7FEBCCC4}{}
{\noindent\textcolor{FuncColor}{$\triangleright$\enspace\texttt{MethodsForObjWithMemory({\mdseries\slshape })\index{MethodsForObjWithMemory@\texttt{MethodsForObjWithMemory}}
\label{MethodsForObjWithMemory}
}\hfill{\scriptsize (function)}}\\


 This function is intended to give an overview of those methods in the current \textsf{GAP} session for which at least one argument is required to be in the filter \texttt{IsObjWithMemory} (\ref{IsObjWithMemory}). 

 The return value is a list of records, with the same components as the ones
returned by \texttt{MethodsOperation} (\ref{MethodsOperation}). In particular, the \texttt{func} component of each record is the method itself, and the \texttt{info} component is a string that describes its requirements. 

 
\begin{Verbatim}[commandchars=!@|,fontsize=\small,frame=single,label=Example]
  !gapprompt@gap>| !gapinput@res:= MethodsForObjWithMemory();;|
  !gapprompt@gap>| !gapinput@for r in res{ [ 1 .. 10 ] } do|
  !gapprompt@>| !gapinput@     Print( r.info, "\n" );|
  !gapprompt@>| !gapinput@   od;|
  ViewObj: [ IsObjWithMemory ]
  =: [ IsObjWithMemory, IsObjWithMemory ]
  =: [ IsMultiplicativeElement, IsObjWithMemory ]
  =: [ IsObjWithMemory, IsMultiplicativeElement ]
  <: [ IsObjWithMemory, IsObjWithMemory ]
  <: [ IsMultiplicativeElement, IsObjWithMemory ]
  <: [ IsObjWithMemory, IsMultiplicativeElement ]
  PrintObj: [ IsObjWithMemory ]
  Length: [ IsMatrix and IsObjWithMemory ]
  Length: [ IsObjWithMemory and IsWord ]
\end{Verbatim}
 

 Note that some functions may provide a special treatment for objects in \texttt{IsObjWithMemory} (\ref{IsObjWithMemory}) but no arguments of the function are objects with memory. For example, there
is a \texttt{One} (\ref{One}) method for a group that checks whether the generators are objects with memory,
and if yes returns an identity element with memory. Such functions do not
occur in the output of \texttt{MethodsForObjWithMemory}. }

 }

 }

            
\chapter{\textcolor{Chapter }{Rewriting Systems}}\label{Rewriting Systems}
\logpage{[ 38, 0, 0 ]}
\hyperdef{L}{X7CA8FCFD81AA1890}{}
{
  Rewriting systems in \textsf{GAP} are a framework for dealing with the very general task of rewriting elements
of a free (or \emph{term}) algebra in some normal form. Although most rewriting systems currently in
use are \emph{string rewriting systems} (where the algebra has only one binary operation which is associative) the
framework in \textsf{GAP} is general enough to encompass the task of rewriting algebras of any signature
from groups to semirings. 

 Rewriting systems are already implemented in \textsf{GAP} for finitely presented semigroups and for pc groups. The use of these
particular rewriting systems is described in the corresponding chapters. We
describe here only the general framework of rewriting systems with a
particular emphasis on material which would be helpful for a developer
implementing a rewriting system. 

 We fix some definitions and terminology for the rest of this chapter. Let $T$ be a term algebra in some signature. A \emph{term rewriting system} for $T$ is a set of ordered pairs of elements of $T$ of the form $(l, r)$. Viewed as a set of relations, the rewriting system determines a presentation
for a quotient algebra $A$ of $T$. 

 When we take into account the fact that the relations are expressed as \emph{ordered} pairs, we have a way of \emph{reducing} the elements of $T$. Suppose an element $u$ of $T$ has a subword $l$ and $(l, r)$ is a rule of the rewriting system, then we can replace the subterm $l$ of $u$ by the term $r$ and obtain a new word $v$. We say that we have \emph{rewritten} $u$ as $v$. Note that $u$ and $v$ represent the same element of $A$. If $u$ cannot be rewritten using any rule of the rewriting system we sat that $u$ is \emph{reduced}.  
\section{\textcolor{Chapter }{Operations on rewriting systems}}\label{Operations on rewriting systems}
\logpage{[ 38, 1, 0 ]}
\hyperdef{L}{X8287CBE183EBE5D7}{}
{
  

\subsection{\textcolor{Chapter }{IsRewritingSystem}}
\logpage{[ 38, 1, 1 ]}\nobreak
\hyperdef{L}{X842C0ED87986F7AA}{}
{\noindent\textcolor{FuncColor}{$\triangleright$\enspace\texttt{IsRewritingSystem({\mdseries\slshape obj})\index{IsRewritingSystem@\texttt{IsRewritingSystem}}
\label{IsRewritingSystem}
}\hfill{\scriptsize (Category)}}\\


 This is the category in which all rewriting systems lie. }

 

\subsection{\textcolor{Chapter }{Rules}}
\logpage{[ 38, 1, 2 ]}\nobreak
\hyperdef{L}{X833EAA8C86356F42}{}
{\noindent\textcolor{FuncColor}{$\triangleright$\enspace\texttt{Rules({\mdseries\slshape rws})\index{Rules@\texttt{Rules}}
\label{Rules}
}\hfill{\scriptsize (attribute)}}\\


 The rules comprising the rewriting system. Note that these may change through
the life of the rewriting system, however they will always be a set of
defining relations of the algebra described by the rewriting system. }

 

\subsection{\textcolor{Chapter }{OrderOfRewritingSystem}}
\logpage{[ 38, 1, 3 ]}\nobreak
\hyperdef{L}{X7C38C2EF817F9E0A}{}
{\noindent\textcolor{FuncColor}{$\triangleright$\enspace\texttt{OrderOfRewritingSystem({\mdseries\slshape rws})\index{OrderOfRewritingSystem@\texttt{OrderOfRewritingSystem}}
\label{OrderOfRewritingSystem}
}\hfill{\scriptsize (attribute)}}\\
\noindent\textcolor{FuncColor}{$\triangleright$\enspace\texttt{OrderingOfRewritingSystem({\mdseries\slshape rws})\index{OrderingOfRewritingSystem@\texttt{OrderingOfRewritingSystem}}
\label{OrderingOfRewritingSystem}
}\hfill{\scriptsize (attribute)}}\\


 return the ordering of the rewriting system \mbox{\texttt{\mdseries\slshape rws}}.  }

 

\subsection{\textcolor{Chapter }{ReducedForm}}
\logpage{[ 38, 1, 4 ]}\nobreak
\hyperdef{L}{X8340EB2280DE6CCC}{}
{\noindent\textcolor{FuncColor}{$\triangleright$\enspace\texttt{ReducedForm({\mdseries\slshape rws, u})\index{ReducedForm@\texttt{ReducedForm}}
\label{ReducedForm}
}\hfill{\scriptsize (operation)}}\\


 Given an element \mbox{\texttt{\mdseries\slshape u}} in the free (or term) algebra $T$ over which \mbox{\texttt{\mdseries\slshape rws}} is defined, rewrite \mbox{\texttt{\mdseries\slshape u}} by successive applications of the rules of \mbox{\texttt{\mdseries\slshape rws}} until no further rewriting is possible, and return the resulting element of $T$. }

 
\subsection{\textcolor{Chapter }{IsConfluent}}\logpage{[ 38, 1, 5 ]}
\hyperdef{L}{X8006790B86328CE8}{}
{
\noindent\textcolor{FuncColor}{$\triangleright$\enspace\texttt{IsConfluent({\mdseries\slshape rws})\index{IsConfluent@\texttt{IsConfluent}!for a rewriting system}
\label{IsConfluent:for a rewriting system}
}\hfill{\scriptsize (property)}}\\
\noindent\textcolor{FuncColor}{$\triangleright$\enspace\texttt{IsConfluent({\mdseries\slshape A})\index{IsConfluent@\texttt{IsConfluent}!for an algebra with canonical rewriting system}
\label{IsConfluent:for an algebra with canonical rewriting system}
}\hfill{\scriptsize (property)}}\\


 For a rewriting system \mbox{\texttt{\mdseries\slshape rws}}, \texttt{IsConfluent} returns \texttt{true} if and only if \mbox{\texttt{\mdseries\slshape rws}} is confluent. A rewriting system is \emph{confluent} if, for every two words $u$ and $v$ in the free algebra $T$ which represent the same element of the algebra $A$ defined by \mbox{\texttt{\mdseries\slshape rws}}, \texttt{ReducedForm( \mbox{\texttt{\mdseries\slshape rws}}, }$u$ \texttt{) = ReducedForm( \mbox{\texttt{\mdseries\slshape rws}}, }$v$\texttt{)} as words in the free algebra $T$. This element is the \emph{unique normal form} of the element represented by $u$. 

 For an algebra \mbox{\texttt{\mdseries\slshape A}} with a canonical rewriting system associated with it, \texttt{IsConfluent} checks whether that rewriting system is confluent. 

 Also see{\nobreakspace}\texttt{IsConfluent} (\ref{IsConfluent:for pc groups}). }

 

\subsection{\textcolor{Chapter }{ConfluentRws}}
\logpage{[ 38, 1, 6 ]}\nobreak
\hyperdef{L}{X870A1E1C7FB45A55}{}
{\noindent\textcolor{FuncColor}{$\triangleright$\enspace\texttt{ConfluentRws({\mdseries\slshape rws})\index{ConfluentRws@\texttt{ConfluentRws}}
\label{ConfluentRws}
}\hfill{\scriptsize (attribute)}}\\


 Return a new rewriting system defining the same algebra as \mbox{\texttt{\mdseries\slshape rws}} which is confluent. }

 

\subsection{\textcolor{Chapter }{IsReduced}}
\logpage{[ 38, 1, 7 ]}\nobreak
\hyperdef{L}{X8134689C7B576946}{}
{\noindent\textcolor{FuncColor}{$\triangleright$\enspace\texttt{IsReduced({\mdseries\slshape rws})\index{IsReduced@\texttt{IsReduced}}
\label{IsReduced}
}\hfill{\scriptsize (property)}}\\


 A rewriting system is reduced if for each rule $(l, r)$, $l$ and $r$ are both reduced. 

 }

 

\subsection{\textcolor{Chapter }{ReduceRules}}
\logpage{[ 38, 1, 8 ]}\nobreak
\hyperdef{L}{X864C82FD7FBA31A6}{}
{\noindent\textcolor{FuncColor}{$\triangleright$\enspace\texttt{ReduceRules({\mdseries\slshape rws})\index{ReduceRules@\texttt{ReduceRules}}
\label{ReduceRules}
}\hfill{\scriptsize (operation)}}\\


 Reduce rules and remove redundant rules to make \mbox{\texttt{\mdseries\slshape rws}} reduced. }

 

\subsection{\textcolor{Chapter }{AddRule}}
\logpage{[ 38, 1, 9 ]}\nobreak
\hyperdef{L}{X81E6B5CB789A7C3A}{}
{\noindent\textcolor{FuncColor}{$\triangleright$\enspace\texttt{AddRule({\mdseries\slshape rws, rule})\index{AddRule@\texttt{AddRule}}
\label{AddRule}
}\hfill{\scriptsize (operation)}}\\


 Add \mbox{\texttt{\mdseries\slshape rule}} to a rewriting system \mbox{\texttt{\mdseries\slshape rws}}. }

 

\subsection{\textcolor{Chapter }{AddRuleReduced}}
\logpage{[ 38, 1, 10 ]}\nobreak
\hyperdef{L}{X7FA0B54D7C533DDC}{}
{\noindent\textcolor{FuncColor}{$\triangleright$\enspace\texttt{AddRuleReduced({\mdseries\slshape rws, rule})\index{AddRuleReduced@\texttt{AddRuleReduced}}
\label{AddRuleReduced}
}\hfill{\scriptsize (operation)}}\\


 Add \mbox{\texttt{\mdseries\slshape rule}} to rewriting system \mbox{\texttt{\mdseries\slshape rws}}. Performs a reduction operation on the resulting system, so that if \mbox{\texttt{\mdseries\slshape rws}} is reduced it will remain reduced. }

 

\subsection{\textcolor{Chapter }{MakeConfluent}}
\logpage{[ 38, 1, 11 ]}\nobreak
\hyperdef{L}{X7BD6299E85561DC3}{}
{\noindent\textcolor{FuncColor}{$\triangleright$\enspace\texttt{MakeConfluent({\mdseries\slshape rws})\index{MakeConfluent@\texttt{MakeConfluent}}
\label{MakeConfluent}
}\hfill{\scriptsize (operation)}}\\


 Add rules (and perhaps reduce) in order to make \mbox{\texttt{\mdseries\slshape rws}} confluent }

 

\subsection{\textcolor{Chapter }{GeneratorsOfRws}}
\logpage{[ 38, 1, 12 ]}\nobreak
\hyperdef{L}{X795DC25886007DFE}{}
{\noindent\textcolor{FuncColor}{$\triangleright$\enspace\texttt{GeneratorsOfRws({\mdseries\slshape rws})\index{GeneratorsOfRws@\texttt{GeneratorsOfRws}}
\label{GeneratorsOfRws}
}\hfill{\scriptsize (attribute)}}\\


 Returns the list of generators of the rewriting system \mbox{\texttt{\mdseries\slshape rws}}. }

 }

  
\section{\textcolor{Chapter }{Operations on elements of the algebra}}\label{Operations on elements of the algebra}
\logpage{[ 38, 2, 0 ]}
\hyperdef{L}{X81B812C778CB1E4E}{}
{
  In this section let $u$ denote an element of the term algebra $T$ representing $[u]$ in the quotient algebra $A$. 

\subsection{\textcolor{Chapter }{ReducedProduct}}
\logpage{[ 38, 2, 1 ]}\nobreak
\hyperdef{L}{X81BB38CC793F7CE2}{}
{\noindent\textcolor{FuncColor}{$\triangleright$\enspace\texttt{ReducedProduct({\mdseries\slshape rws, u, v})\index{ReducedProduct@\texttt{ReducedProduct}}
\label{ReducedProduct}
}\hfill{\scriptsize (operation)}}\\
\noindent\textcolor{FuncColor}{$\triangleright$\enspace\texttt{ReducedSum({\mdseries\slshape rws, left, right})\index{ReducedSum@\texttt{ReducedSum}}
\label{ReducedSum}
}\hfill{\scriptsize (operation)}}\\
\noindent\textcolor{FuncColor}{$\triangleright$\enspace\texttt{ReducedOne({\mdseries\slshape rws})\index{ReducedOne@\texttt{ReducedOne}}
\label{ReducedOne}
}\hfill{\scriptsize (operation)}}\\
\noindent\textcolor{FuncColor}{$\triangleright$\enspace\texttt{ReducedAdditiveInverse({\mdseries\slshape rws, obj})\index{ReducedAdditiveInverse@\texttt{ReducedAdditiveInverse}}
\label{ReducedAdditiveInverse}
}\hfill{\scriptsize (operation)}}\\
\noindent\textcolor{FuncColor}{$\triangleright$\enspace\texttt{ReducedComm({\mdseries\slshape rws, left, right})\index{ReducedComm@\texttt{ReducedComm}}
\label{ReducedComm}
}\hfill{\scriptsize (operation)}}\\
\noindent\textcolor{FuncColor}{$\triangleright$\enspace\texttt{ReducedConjugate({\mdseries\slshape rws, left, right})\index{ReducedConjugate@\texttt{ReducedConjugate}}
\label{ReducedConjugate}
}\hfill{\scriptsize (operation)}}\\
\noindent\textcolor{FuncColor}{$\triangleright$\enspace\texttt{ReducedDifference({\mdseries\slshape rws, left, right})\index{ReducedDifference@\texttt{ReducedDifference}}
\label{ReducedDifference}
}\hfill{\scriptsize (operation)}}\\
\noindent\textcolor{FuncColor}{$\triangleright$\enspace\texttt{ReducedInverse({\mdseries\slshape rws, obj})\index{ReducedInverse@\texttt{ReducedInverse}}
\label{ReducedInverse}
}\hfill{\scriptsize (operation)}}\\
\noindent\textcolor{FuncColor}{$\triangleright$\enspace\texttt{ReducedLeftQuotient({\mdseries\slshape rws, left, right})\index{ReducedLeftQuotient@\texttt{ReducedLeftQuotient}}
\label{ReducedLeftQuotient}
}\hfill{\scriptsize (operation)}}\\
\noindent\textcolor{FuncColor}{$\triangleright$\enspace\texttt{ReducedPower({\mdseries\slshape rws, obj, pow})\index{ReducedPower@\texttt{ReducedPower}}
\label{ReducedPower}
}\hfill{\scriptsize (operation)}}\\
\noindent\textcolor{FuncColor}{$\triangleright$\enspace\texttt{ReducedQuotient({\mdseries\slshape rws, left, right})\index{ReducedQuotient@\texttt{ReducedQuotient}}
\label{ReducedQuotient}
}\hfill{\scriptsize (operation)}}\\
\noindent\textcolor{FuncColor}{$\triangleright$\enspace\texttt{ReducedScalarProduct({\mdseries\slshape rws, left, right})\index{ReducedScalarProduct@\texttt{ReducedScalarProduct}}
\label{ReducedScalarProduct}
}\hfill{\scriptsize (operation)}}\\
\noindent\textcolor{FuncColor}{$\triangleright$\enspace\texttt{ReducedZero({\mdseries\slshape rws})\index{ReducedZero@\texttt{ReducedZero}}
\label{ReducedZero}
}\hfill{\scriptsize (operation)}}\\


 The result of \texttt{ReducedProduct} is $w$ where $[w]$ equals [\mbox{\texttt{\mdseries\slshape u}}][\mbox{\texttt{\mdseries\slshape v}}] in $A$ and $w$ is in reduced form. 

 The remaining operations are defined similarly when they are defined (as
determined by the signature of the term algebra). }

 }

  
\section{\textcolor{Chapter }{Properties of rewriting systems}}\label{Properties of rewriting systems}
\logpage{[ 38, 3, 0 ]}
\hyperdef{L}{X8318649681DF783B}{}
{
  

\subsection{\textcolor{Chapter }{IsBuiltFromAdditiveMagmaWithInverses}}
\logpage{[ 38, 3, 1 ]}\nobreak
\hyperdef{L}{X7B647DB77D138A49}{}
{\noindent\textcolor{FuncColor}{$\triangleright$\enspace\texttt{IsBuiltFromAdditiveMagmaWithInverses({\mdseries\slshape obj})\index{IsBuiltFromAdditiveMagmaWithInverses@\texttt{IsBuilt}\-\texttt{From}\-\texttt{Additive}\-\texttt{Magma}\-\texttt{With}\-\texttt{Inverses}}
\label{IsBuiltFromAdditiveMagmaWithInverses}
}\hfill{\scriptsize (property)}}\\
\noindent\textcolor{FuncColor}{$\triangleright$\enspace\texttt{IsBuiltFromMagma({\mdseries\slshape obj})\index{IsBuiltFromMagma@\texttt{IsBuiltFromMagma}}
\label{IsBuiltFromMagma}
}\hfill{\scriptsize (property)}}\\
\noindent\textcolor{FuncColor}{$\triangleright$\enspace\texttt{IsBuiltFromMagmaWithOne({\mdseries\slshape obj})\index{IsBuiltFromMagmaWithOne@\texttt{IsBuiltFromMagmaWithOne}}
\label{IsBuiltFromMagmaWithOne}
}\hfill{\scriptsize (property)}}\\
\noindent\textcolor{FuncColor}{$\triangleright$\enspace\texttt{IsBuiltFromMagmaWithInverses({\mdseries\slshape obj})\index{IsBuiltFromMagmaWithInverses@\texttt{IsBuiltFromMagmaWithInverses}}
\label{IsBuiltFromMagmaWithInverses}
}\hfill{\scriptsize (property)}}\\
\noindent\textcolor{FuncColor}{$\triangleright$\enspace\texttt{IsBuiltFromSemigroup({\mdseries\slshape obj})\index{IsBuiltFromSemigroup@\texttt{IsBuiltFromSemigroup}}
\label{IsBuiltFromSemigroup}
}\hfill{\scriptsize (property)}}\\
\noindent\textcolor{FuncColor}{$\triangleright$\enspace\texttt{IsBuiltFromGroup({\mdseries\slshape obj})\index{IsBuiltFromGroup@\texttt{IsBuiltFromGroup}}
\label{IsBuiltFromGroup}
}\hfill{\scriptsize (property)}}\\


 These properties may be used to identify the type of term algebra over which
the rewriting system is defined. }

 }

  
\section{\textcolor{Chapter }{Rewriting in Groups and Monoids}}\label{Rewriting in Groups and Monoids}
\logpage{[ 38, 4, 0 ]}
\hyperdef{L}{X7F8B7848851784DF}{}
{
  One application of rewriting is to reduce words in finitely presented groups
and monoids. The rewriting system still has to be built for a finitely
presented monoid (using \texttt{IsomorphismFpMonoid} for conversion). Rewriting then can take place for words in the underlying
free monoid. (These can be obtained from monoid elements with the command \texttt{UnderlyingElement}.) 

 
\begin{Verbatim}[commandchars=!@|,fontsize=\small,frame=single,label=Example]
  !gapprompt@gap>| !gapinput@f:=FreeGroup(3);;|
  !gapprompt@gap>| !gapinput@rels:=[f.1*f.2^2/f.3,f.2*f.3^2/f.1,f.3*f.1^2/f.2];;|
  !gapprompt@gap>| !gapinput@g:=f/rels;|
  <fp group on the generators [ f1, f2, f3 ]>
  !gapprompt@gap>| !gapinput@mhom:=IsomorphismFpMonoid(g);|
  MappingByFunction( <fp group on the generators
  [ f1, f2, f3 ]>, <fp monoid on the generators
  [ f1, f1^-1, f2, f2^-1, f3, f3^-1
   ]>, function( x ) ... end, function( x ) ... end )
  !gapprompt@gap>| !gapinput@mon:=Image(mhom);|
  <fp monoid on the generators [ f1, f1^-1, f2, f2^-1, f3, f3^-1 ]>
  !gapprompt@gap>| !gapinput@k:=KnuthBendixRewritingSystem(mon);|
  Knuth Bendix Rewriting System for Monoid(
  [ f1, f1^-1, f2, f2^-1, f3, f3^-1 ] ) with rules
  [ [ f1*f1^-1, <identity ...> ], [ f1^-1*f1, <identity ...> ],
    [ f2*f2^-1, <identity ...> ], [ f2^-1*f2, <identity ...> ],
    [ f3*f3^-1, <identity ...> ], [ f3^-1*f3, <identity ...> ],
    [ f1*f2^2*f3^-1, <identity ...> ], [ f2*f3^2*f1^-1, <identity ...> ]
      , [ f3*f1^2*f2^-1, <identity ...> ] ]
  !gapprompt@gap>| !gapinput@MakeConfluent(k);|
  !gapprompt@gap>| !gapinput@a:=Product(GeneratorsOfMonoid(mon));|
  f1*f1^-1*f2*f2^-1*f3*f3^-1
  !gapprompt@gap>| !gapinput@ReducedForm(k,UnderlyingElement(a));|
  <identity ...>
\end{Verbatim}
 

 To rewrite a word in the finitely presented group, one has to convert it to a
word in the monoid first, rewrite in the underlying free monoid and convert
back (by forming first again an element of the fp monoid) to the finitely
presented group. 

   
\begin{Verbatim}[commandchars=!@|,fontsize=\small,frame=single,label=Example]
  !gapprompt@gap>| !gapinput@r:=PseudoRandom(g);;|
  !gapprompt@gap>| !gapinput@Length(r);|
  3704
  !gapprompt@gap>| !gapinput@melm:=Image(mhom,r);;|
  !gapprompt@gap>| !gapinput@red:=ReducedForm(k,UnderlyingElement(melm));|
  f1^-1^3*f2^-1*f1^2
  !gapprompt@gap>| !gapinput@melm:=ElementOfFpMonoid(FamilyObj(One(mon)),red);|
  f1^-1^3*f2^-1*f1^2
  !gapprompt@gap>| !gapinput@gpelm:=PreImagesRepresentative(mhom,melm);|
  f1^-3*f2^-1*f1^2
  !gapprompt@gap>| !gapinput@r=gpelm;|
  true
  !gapprompt@gap>| !gapinput@CategoriesOfObject(red);|
  [ "IsExtLElement", "IsExtRElement", "IsMultiplicativeElement",
    "IsMultiplicativeElementWithOne", "IsAssociativeElement", "IsWord" ]
  !gapprompt@gap>| !gapinput@CategoriesOfObject(melm);|
  [ "IsExtLElement", "IsExtRElement", "IsMultiplicativeElement",
    "IsMultiplicativeElementWithOne", "IsAssociativeElement",
    "IsElementOfFpMonoid" ]
  !gapprompt@gap>| !gapinput@CategoriesOfObject(gpelm);|
  [ "IsExtLElement", "IsExtRElement", "IsMultiplicativeElement",
    "IsMultiplicativeElementWithOne", "IsMultiplicativeElementWithInverse",
    "IsAssociativeElement", "IsElementOfFpGroup" ]
\end{Verbatim}
 

 Note, that the elements \texttt{red} (free monoid) \texttt{melm} (fp monoid) and \texttt{gpelm} (group) differ, though they are displayed identically. 

 Under Unix, it is possible to use the \textsf{kbmag} package to replace the built\texttt{\symbol{45}}in rewriting by this packages
efficient C implementation. You can do this (after loading the \textsf{kbmag} package) by assigning the variable \texttt{KB{\textunderscore}REW} (\ref{KBuScoreREW}) to \texttt{KBMAG{\textunderscore}REW}. Assignment to \texttt{GAPKB{\textunderscore}REW} reverts to the built\texttt{\symbol{45}}in implementation. 
\begin{Verbatim}[commandchars=!@|,fontsize=\small,frame=single,label=Example]
  !gapprompt@gap>| !gapinput@LoadPackage("kbmag");|
  true
  !gapprompt@gap>| !gapinput@KB_REW:=KBMAG_REW;;|
\end{Verbatim}
 }

  
\section{\textcolor{Chapter }{Developing rewriting systems}}\label{Developing rewriting systems}
\logpage{[ 38, 5, 0 ]}
\hyperdef{L}{X8751F8FA7DC989A2}{}
{
  The key point to note about rewriting systems is that they have properties
such as \texttt{IsConfluent} (\ref{IsConfluent:for a rewriting system}) and attributes such as \texttt{Rules} (\ref{Rules}), however they are rarely stored, but rather computed afresh each time they
are asked for, from data stored in the private members of the rewriting system
object. This is because a rewriting system often evolves through a session,
starting with some rules which define the algebra \mbox{\texttt{\mdseries\slshape A}} as relations, and then adding more rules to make the system confluent. For
example, in the case of Knuth\texttt{\symbol{45}}Bendix rewriting systems (see
Chapter{\nobreakspace}\ref{Finitely Presented Semigroups and Monoids}), the function \texttt{CreateKnuthBendixRewritingSystem} creating the rewriting system (in the file \texttt{lib/kbsemi.gi}) uses 

 
\begin{Verbatim}[commandchars=!@|,fontsize=\small,frame=single,label=Example]
  kbrws := Objectify(NewType(rwsfam,
    IsMutable and IsKnuthBendixRewritingSystem and
    IsKnuthBendixRewritingSystemRep),
    rec(family:= fam,
    reduced:=false,
    tzrules:=List(relwco,i->
     [LetterRepAssocWord(i[1]),LetterRepAssocWord(i[2])]),
    pairs2check:=CantorList(Length(r)),
    ordering:=wordord,
    freefam:=freefam));
\end{Verbatim}
 

 In particular, since we don't use the filter \texttt{IsAttributeStoringRep} in the \texttt{Objectify} (\ref{Objectify}), whenever \texttt{IsConfluent} (\ref{IsConfluent:for a rewriting system}) is called, the appropriate method to determine confluence is called. }

 }

            
\chapter{\textcolor{Chapter }{Groups}}\label{Groups}
\logpage{[ 39, 0, 0 ]}
\hyperdef{L}{X8716635F7951801B}{}
{
  This chapter explains how to create groups and defines operations for groups,
that is operations whose definition does not depend on the representation
used. However methods for these operations in most cases will make use of the
representation. 

 If not otherwise specified, in all examples in this chapter the group \texttt{g} will be the symmetric group $S_4$ acting on the letters $\{ 1, \ldots, 4 \}$.  
\section{\textcolor{Chapter }{Group Elements}}\label{Group Elements}
\logpage{[ 39, 1, 0 ]}
\hyperdef{L}{X822370B47DEA37B1}{}
{
  Groups in \textsf{GAP} are written multiplicatively. The elements from which a group can be generated
must permit multiplication and multiplicative inversion (see{\nobreakspace}\ref{Useful Categories of Elements}). 

 
\begin{Verbatim}[commandchars=!@|,fontsize=\small,frame=single,label=Example]
  !gapprompt@gap>| !gapinput@a:=(1,2,3);;b:=(2,3,4);;|
  !gapprompt@gap>| !gapinput@One(a);|
  ()
  !gapprompt@gap>| !gapinput@Inverse(b);|
  (2,4,3)
  !gapprompt@gap>| !gapinput@a*b;|
  (1,3)(2,4)
  !gapprompt@gap>| !gapinput@Order(a*b);|
  2
  !gapprompt@gap>| !gapinput@Order( [ [ 1, 1 ], [ 0, 1 ] ] );|
  infinity
\end{Verbatim}
 

 The next example may run into an infinite loop because the given matrix in
fact has infinite order. 

 
\begin{Verbatim}[commandchars=!@|,fontsize=\small,frame=single,label=Example]
  !gapprompt@gap>| !gapinput@Order( [ [ 1, 1 ], [ 0, 1 ] ] * Indeterminate( Rationals ) );|
  #I  Order: warning, order of <mat> might be infinite
\end{Verbatim}
 

 \index{order!of a group} Since groups are domains, the recommended command to compute the order of a
group is \texttt{Size} (\ref{Size}). For convenience, group orders can also be computed with \texttt{Order} (\ref{Order}). 

 The operation \texttt{Comm} (\ref{Comm}) can be used to compute the commutator of two elements, the operation \texttt{LeftQuotient} (\ref{LeftQuotient}) computes the product $x^{{-1}} y$. }

  
\section{\textcolor{Chapter }{Creating Groups}}\label{Creating Groups}
\logpage{[ 39, 2, 0 ]}
\hyperdef{L}{X86A022F9800121F8}{}
{
  When groups are created from generators, this means that the generators must
be elements that can be multiplied and inverted (see also{\nobreakspace}\ref{Constructing Domains}). For creating a free group on a set of symbols, see{\nobreakspace}\texttt{FreeGroup} (\ref{FreeGroup:for given rank}). 
\subsection{\textcolor{Chapter }{Group}}\logpage{[ 39, 2, 1 ]}
\hyperdef{L}{X7D7B075385435151}{}
{
\noindent\textcolor{FuncColor}{$\triangleright$\enspace\texttt{Group({\mdseries\slshape gen, ...})\index{Group@\texttt{Group}!for several generators}
\label{Group:for several generators}
}\hfill{\scriptsize (function)}}\\
\noindent\textcolor{FuncColor}{$\triangleright$\enspace\texttt{Group({\mdseries\slshape gens[, id]})\index{Group@\texttt{Group}!for a list of generators (and an identity element)}
\label{Group:for a list of generators (and an identity element)}
}\hfill{\scriptsize (function)}}\\


 \texttt{Group( \mbox{\texttt{\mdseries\slshape gen}}, ... )} is the group generated by the arguments \mbox{\texttt{\mdseries\slshape gen}}, ... 

 If the only argument \mbox{\texttt{\mdseries\slshape gens}} is a list that is not a matrix then \texttt{Group( \mbox{\texttt{\mdseries\slshape gens}} )} is the group generated by the elements of that list. 

 If there are two arguments, a list \mbox{\texttt{\mdseries\slshape gens}} and an element \mbox{\texttt{\mdseries\slshape id}}, then \texttt{Group( \mbox{\texttt{\mdseries\slshape gens}}, \mbox{\texttt{\mdseries\slshape id}} )} is the group generated by the elements of \mbox{\texttt{\mdseries\slshape gens}}, with identity \mbox{\texttt{\mdseries\slshape id}}. 

 Note that the value of the attribute \texttt{GeneratorsOfGroup} (\ref{GeneratorsOfGroup}) need not be equal to the list \mbox{\texttt{\mdseries\slshape gens}} of generators entered as argument. Use \texttt{GroupWithGenerators} (\ref{GroupWithGenerators}) if you want to be sure that the argument \mbox{\texttt{\mdseries\slshape gens}} is stored as value of \texttt{GeneratorsOfGroup} (\ref{GeneratorsOfGroup}). 
\begin{Verbatim}[commandchars=!@|,fontsize=\small,frame=single,label=Example]
  !gapprompt@gap>| !gapinput@g:=Group((1,2,3,4),(1,2));|
  Group([ (1,2,3,4), (1,2) ])
\end{Verbatim}
 }

 

\subsection{\textcolor{Chapter }{GroupByGenerators}}
\logpage{[ 39, 2, 2 ]}\nobreak
\hyperdef{L}{X7F81960287F3E32A}{}
{\noindent\textcolor{FuncColor}{$\triangleright$\enspace\texttt{GroupByGenerators({\mdseries\slshape gens})\index{GroupByGenerators@\texttt{GroupByGenerators}}
\label{GroupByGenerators}
}\hfill{\scriptsize (operation)}}\\
\noindent\textcolor{FuncColor}{$\triangleright$\enspace\texttt{GroupByGenerators({\mdseries\slshape gens, id})\index{GroupByGenerators@\texttt{GroupByGenerators}!with explicitly specified identity element}
\label{GroupByGenerators:with explicitly specified identity element}
}\hfill{\scriptsize (operation)}}\\


 \texttt{GroupByGenerators} returns the group $G$ generated by the list \mbox{\texttt{\mdseries\slshape gens}}. If a second argument \mbox{\texttt{\mdseries\slshape id}} is present then this is stored as the identity element of the group. 

 The value of the attribute \texttt{GeneratorsOfGroup} (\ref{GeneratorsOfGroup}) of $G$ need not be equal to \mbox{\texttt{\mdseries\slshape gens}}. \texttt{GroupByGenerators} is the underlying operation called by \texttt{Group} (\ref{Group:for several generators}). }

 

\subsection{\textcolor{Chapter }{GroupWithGenerators}}
\logpage{[ 39, 2, 3 ]}\nobreak
\hyperdef{L}{X8589EF9C7B658B94}{}
{\noindent\textcolor{FuncColor}{$\triangleright$\enspace\texttt{GroupWithGenerators({\mdseries\slshape gens[, id]})\index{GroupWithGenerators@\texttt{GroupWithGenerators}}
\label{GroupWithGenerators}
}\hfill{\scriptsize (operation)}}\\


 \texttt{GroupWithGenerators} returns the group $G$ generated by the list \mbox{\texttt{\mdseries\slshape gens}}. If a second argument \mbox{\texttt{\mdseries\slshape id}} is present then this is stored as the identity element of the group. The value
of the attribute \texttt{GeneratorsOfGroup} (\ref{GeneratorsOfGroup}) of $G$ is equal to \mbox{\texttt{\mdseries\slshape gens}}. }

 

\subsection{\textcolor{Chapter }{GeneratorsOfGroup}}
\logpage{[ 39, 2, 4 ]}\nobreak
\hyperdef{L}{X79C44528864044C5}{}
{\noindent\textcolor{FuncColor}{$\triangleright$\enspace\texttt{GeneratorsOfGroup({\mdseries\slshape G})\index{GeneratorsOfGroup@\texttt{GeneratorsOfGroup}}
\label{GeneratorsOfGroup}
}\hfill{\scriptsize (attribute)}}\\


 returns a list of generators of the group \mbox{\texttt{\mdseries\slshape G}}. If \mbox{\texttt{\mdseries\slshape G}} has been created by the command \texttt{GroupWithGenerators} (\ref{GroupWithGenerators}) with argument \mbox{\texttt{\mdseries\slshape gens}}, then the list returned by \texttt{GeneratorsOfGroup} will be equal to \mbox{\texttt{\mdseries\slshape gens}}. For such a group, each generator can also be accessed using the \texttt{.} operator (see \texttt{GeneratorsOfDomain} (\ref{GeneratorsOfDomain})): for a positive integer $i$, \texttt{\mbox{\texttt{\mdseries\slshape G}}.i} returns the $i$\texttt{\symbol{45}}th element of the list returned by \texttt{GeneratorsOfGroup}. Moreover, if \mbox{\texttt{\mdseries\slshape G}} is a free group, and \texttt{name} is the name of a generator of \mbox{\texttt{\mdseries\slshape G}} then \texttt{\mbox{\texttt{\mdseries\slshape G}}.name} also returns this generator. 
\begin{Verbatim}[commandchars=!@|,fontsize=\small,frame=single,label=Example]
  !gapprompt@gap>| !gapinput@g:=GroupWithGenerators([(1,2,3,4),(1,2)]);|
  Group([ (1,2,3,4), (1,2) ])
  !gapprompt@gap>| !gapinput@GeneratorsOfGroup(g);|
  [ (1,2,3,4), (1,2) ]
\end{Verbatim}
 

 While in this example \textsf{GAP} displays the group via the generating set stored in the attribute \texttt{GeneratorsOfGroup}, the methods installed for \texttt{View} (\ref{View}) will in general display only some information about the group which may even
be just the fact that it is a group. }

 

\subsection{\textcolor{Chapter }{AsGroup}}
\logpage{[ 39, 2, 5 ]}\nobreak
\hyperdef{L}{X7A0747F17B50D967}{}
{\noindent\textcolor{FuncColor}{$\triangleright$\enspace\texttt{AsGroup({\mdseries\slshape D})\index{AsGroup@\texttt{AsGroup}}
\label{AsGroup}
}\hfill{\scriptsize (attribute)}}\\


 if the elements of the collection \mbox{\texttt{\mdseries\slshape D}} form a group the command returns this group, otherwise it returns \texttt{fail}. 
\begin{Verbatim}[commandchars=!@|,fontsize=\small,frame=single,label=Example]
  !gapprompt@gap>| !gapinput@AsGroup([(1,2)]);|
  fail
  !gapprompt@gap>| !gapinput@AsGroup([(),(1,2)]);|
  Group([ (1,2) ])
\end{Verbatim}
 }

 

\subsection{\textcolor{Chapter }{ConjugateGroup}}
\logpage{[ 39, 2, 6 ]}\nobreak
\hyperdef{L}{X7E4143A08040BB47}{}
{\noindent\textcolor{FuncColor}{$\triangleright$\enspace\texttt{ConjugateGroup({\mdseries\slshape G, obj})\index{ConjugateGroup@\texttt{ConjugateGroup}}
\label{ConjugateGroup}
}\hfill{\scriptsize (operation)}}\\


 returns the conjugate group of \mbox{\texttt{\mdseries\slshape G}}, obtained by applying the conjugating element \mbox{\texttt{\mdseries\slshape obj}}. 

 To form a conjugate (group) by any object acting via \texttt{\texttt{\symbol{94}}}, one can also use the infix operator \texttt{\texttt{\symbol{94}}}. 
\begin{Verbatim}[commandchars=!@|,fontsize=\small,frame=single,label=Example]
  !gapprompt@gap>| !gapinput@ConjugateGroup(g,(1,5));|
  Group([ (2,3,4,5), (2,5) ])
\end{Verbatim}
 }

 

\subsection{\textcolor{Chapter }{IsGroup}}
\logpage{[ 39, 2, 7 ]}\nobreak
\hyperdef{L}{X7939B3177BBD61E4}{}
{\noindent\textcolor{FuncColor}{$\triangleright$\enspace\texttt{IsGroup({\mdseries\slshape obj})\index{IsGroup@\texttt{IsGroup}}
\label{IsGroup}
}\hfill{\scriptsize (Category)}}\\


 A group is a magma\texttt{\symbol{45}}with\texttt{\symbol{45}}inverses
(see{\nobreakspace}\texttt{IsMagmaWithInverses} (\ref{IsMagmaWithInverses})) and associative (see{\nobreakspace}\texttt{IsAssociative} (\ref{IsAssociative})) multiplication. 

 \texttt{IsGroup} tests whether the object \mbox{\texttt{\mdseries\slshape obj}} fulfills these conditions, it does \emph{not} test whether \mbox{\texttt{\mdseries\slshape obj}} is a set of elements that forms a group under multiplication; use \texttt{AsGroup} (\ref{AsGroup}) if you want to perform such a test. (See{\nobreakspace}\ref{Categories} for details about categories.) 
\begin{Verbatim}[commandchars=!@|,fontsize=\small,frame=single,label=Example]
  !gapprompt@gap>| !gapinput@IsGroup(g);|
  true
\end{Verbatim}
 }

 

\subsection{\textcolor{Chapter }{InfoGroup}}
\logpage{[ 39, 2, 8 ]}\nobreak
\hyperdef{L}{X845874BA82E1A11F}{}
{\noindent\textcolor{FuncColor}{$\triangleright$\enspace\texttt{InfoGroup\index{InfoGroup@\texttt{InfoGroup}}
\label{InfoGroup}
}\hfill{\scriptsize (info class)}}\\


 is the info class for the generic group theoretic functions
(see{\nobreakspace}\ref{Info Functions}). }

 }

  
\section{\textcolor{Chapter }{Subgroups}}\label{Subgroups}
\logpage{[ 39, 3, 0 ]}
\hyperdef{L}{X7BA181CA81D785BB}{}
{
  For the general concept of parents and subdomains, see{\nobreakspace}\ref{Parents} and{\nobreakspace}\ref{Constructing Subdomains}. More functions that construct certain subgroups can be found in the
sections{\nobreakspace}\ref{Normal Structure}, \ref{Specific and Parametrized Subgroups}, \ref{Sylow Subgroups and Hall Subgroups}, and{\nobreakspace}\ref{Subgroups characterized by prime powers}. 

 If a group $U$ is created as a subgroup of another group $G$, $G$ becomes the parent of $U$. There is no ``universal'' parent group, parent\texttt{\symbol{45}}child chains can be arbitrary long. \textsf{GAP} stores the result of some operations (such as \texttt{Normalizer} (\ref{Normalizer:for two groups})) with the parent as an attribute. 

\subsection{\textcolor{Chapter }{Subgroup}}
\logpage{[ 39, 3, 1 ]}\nobreak
\hyperdef{L}{X7C82AA387A42DCA0}{}
{\noindent\textcolor{FuncColor}{$\triangleright$\enspace\texttt{Subgroup({\mdseries\slshape G, gens})\index{Subgroup@\texttt{Subgroup}}
\label{Subgroup}
}\hfill{\scriptsize (function)}}\\
\noindent\textcolor{FuncColor}{$\triangleright$\enspace\texttt{SubgroupNC({\mdseries\slshape G, gens})\index{SubgroupNC@\texttt{SubgroupNC}}
\label{SubgroupNC}
}\hfill{\scriptsize (function)}}\\
\noindent\textcolor{FuncColor}{$\triangleright$\enspace\texttt{Subgroup({\mdseries\slshape G})\index{Subgroup@\texttt{Subgroup}!for a group}
\label{Subgroup:for a group}
}\hfill{\scriptsize (function)}}\\


 creates the subgroup \mbox{\texttt{\mdseries\slshape U}} of \mbox{\texttt{\mdseries\slshape G}} generated by \mbox{\texttt{\mdseries\slshape gens}}. The \texttt{Parent} (\ref{Parent}) value of \mbox{\texttt{\mdseries\slshape U}} will be \mbox{\texttt{\mdseries\slshape G}}. The \texttt{NC} version does not check, whether the elements in \mbox{\texttt{\mdseries\slshape gens}} actually lie in \mbox{\texttt{\mdseries\slshape G}}. 

 The unary version of \texttt{Subgroup} creates a (shell) subgroup that does not even know generators but can be used
to collect information about a particular subgroup over time. 
\begin{Verbatim}[commandchars=!@|,fontsize=\small,frame=single,label=Example]
  !gapprompt@gap>| !gapinput@u:=Subgroup(g,[(1,2,3),(1,2)]);|
  Group([ (1,2,3), (1,2) ])
\end{Verbatim}
 }

 
\subsection{\textcolor{Chapter }{Index (\textsf{GAP} operation)}}\logpage{[ 39, 3, 2 ]}
\hyperdef{L}{X842AD37E79CE953E}{}
{
\noindent\textcolor{FuncColor}{$\triangleright$\enspace\texttt{Index({\mdseries\slshape G, U})\index{Index@\texttt{Index}!for a group and its subgroup}
\label{Index:for a group and its subgroup}
}\hfill{\scriptsize (operation)}}\\
\noindent\textcolor{FuncColor}{$\triangleright$\enspace\texttt{IndexNC({\mdseries\slshape G, U})\index{IndexNC@\texttt{IndexNC}!for a group and its subgroup}
\label{IndexNC:for a group and its subgroup}
}\hfill{\scriptsize (operation)}}\\


 For a subgroup \mbox{\texttt{\mdseries\slshape U}} of the group \mbox{\texttt{\mdseries\slshape G}}, \texttt{Index} returns the index $[\mbox{\texttt{\mdseries\slshape G}}:\mbox{\texttt{\mdseries\slshape U}}] = |\mbox{\texttt{\mdseries\slshape G}}| / |\mbox{\texttt{\mdseries\slshape U}}|$ of \mbox{\texttt{\mdseries\slshape U}} in \mbox{\texttt{\mdseries\slshape G}}. The \texttt{NC} version does not test whether \mbox{\texttt{\mdseries\slshape U}} is contained in \mbox{\texttt{\mdseries\slshape G}}. 
\begin{Verbatim}[commandchars=!@|,fontsize=\small,frame=single,label=Example]
  !gapprompt@gap>| !gapinput@Index(g,u);|
  4
\end{Verbatim}
 }

 

\subsection{\textcolor{Chapter }{IndexInWholeGroup}}
\logpage{[ 39, 3, 3 ]}\nobreak
\hyperdef{L}{X8014135884DCC53E}{}
{\noindent\textcolor{FuncColor}{$\triangleright$\enspace\texttt{IndexInWholeGroup({\mdseries\slshape G})\index{IndexInWholeGroup@\texttt{IndexInWholeGroup}}
\label{IndexInWholeGroup}
}\hfill{\scriptsize (attribute)}}\\


 If the family of elements of \mbox{\texttt{\mdseries\slshape G}} itself forms a group \mbox{\texttt{\mdseries\slshape P}}, this attribute returns the index of \mbox{\texttt{\mdseries\slshape G}} in \mbox{\texttt{\mdseries\slshape P}}. It is used primarily for free groups or finitely presented groups. }

 
\begin{Verbatim}[commandchars=!@|,fontsize=\small,frame=single,label=Example]
  !gapprompt@gap>| !gapinput@freegp:=FreeGroup(1);;|
  !gapprompt@gap>| !gapinput@freesub:=Subgroup(freegp,[freegp.1^5]);;|
  !gapprompt@gap>| !gapinput@IndexInWholeGroup(freesub);|
  5
\end{Verbatim}
 

\subsection{\textcolor{Chapter }{AsSubgroup}}
\logpage{[ 39, 3, 4 ]}\nobreak
\hyperdef{L}{X7904AC9D7E9A3BB7}{}
{\noindent\textcolor{FuncColor}{$\triangleright$\enspace\texttt{AsSubgroup({\mdseries\slshape G, U})\index{AsSubgroup@\texttt{AsSubgroup}}
\label{AsSubgroup}
}\hfill{\scriptsize (operation)}}\\


 creates a subgroup of \mbox{\texttt{\mdseries\slshape G}} which contains the same elements as \mbox{\texttt{\mdseries\slshape U}} 
\begin{Verbatim}[commandchars=!@|,fontsize=\small,frame=single,label=Example]
  !gapprompt@gap>| !gapinput@v:=AsSubgroup(g,Group((1,2,3),(1,4)));|
  Group([ (1,2,3), (1,4) ])
  !gapprompt@gap>| !gapinput@Parent(v);|
  Group([ (1,2,3,4), (1,2) ])
\end{Verbatim}
 }

 

\subsection{\textcolor{Chapter }{IsSubgroup}}
\logpage{[ 39, 3, 5 ]}\nobreak
\hyperdef{L}{X7839D8927E778334}{}
{\noindent\textcolor{FuncColor}{$\triangleright$\enspace\texttt{IsSubgroup({\mdseries\slshape G, U})\index{IsSubgroup@\texttt{IsSubgroup}}
\label{IsSubgroup}
}\hfill{\scriptsize (function)}}\\


 \texttt{IsSubgroup} returns \texttt{true} if \mbox{\texttt{\mdseries\slshape U}} is a group that is a subset of the domain \mbox{\texttt{\mdseries\slshape G}}. This is actually checked by calling \texttt{IsGroup( \mbox{\texttt{\mdseries\slshape U}} )} and \texttt{IsSubset( \mbox{\texttt{\mdseries\slshape G}}, \mbox{\texttt{\mdseries\slshape U}} )}; note that special methods for \texttt{IsSubset} (\ref{IsSubset}) are available that test only generators of \mbox{\texttt{\mdseries\slshape U}} if \mbox{\texttt{\mdseries\slshape G}} is closed under the group operations. So in most cases, for example whenever
one knows already that \mbox{\texttt{\mdseries\slshape U}} is a group, it is better to call only \texttt{IsSubset} (\ref{IsSubset}). 
\begin{Verbatim}[commandchars=!@|,fontsize=\small,frame=single,label=Example]
  !gapprompt@gap>| !gapinput@IsSubgroup(g,u);|
  true
  !gapprompt@gap>| !gapinput@v:=Group((1,2,3),(1,2));|
  Group([ (1,2,3), (1,2) ])
  !gapprompt@gap>| !gapinput@u=v;|
  true
  !gapprompt@gap>| !gapinput@IsSubgroup(g,v);|
  true
\end{Verbatim}
 }

 

\subsection{\textcolor{Chapter }{IsNormal}}
\logpage{[ 39, 3, 6 ]}\nobreak
\hyperdef{L}{X838186F9836F678C}{}
{\noindent\textcolor{FuncColor}{$\triangleright$\enspace\texttt{IsNormal({\mdseries\slshape G, U})\index{IsNormal@\texttt{IsNormal}}
\label{IsNormal}
}\hfill{\scriptsize (operation)}}\\


 returns \texttt{true} if the group \mbox{\texttt{\mdseries\slshape G}} normalizes the group \mbox{\texttt{\mdseries\slshape U}} and \texttt{false} otherwise. 

 A group \mbox{\texttt{\mdseries\slshape G}} \emph{normalizes} a group \mbox{\texttt{\mdseries\slshape U}} if and only if for every $g \in \mbox{\texttt{\mdseries\slshape G}}$ and $u \in \mbox{\texttt{\mdseries\slshape U}}$ the element $u^g$ is a member of \mbox{\texttt{\mdseries\slshape U}}. Note that \mbox{\texttt{\mdseries\slshape U}} need not be a subgroup of \mbox{\texttt{\mdseries\slshape G}}. 
\begin{Verbatim}[commandchars=!@|,fontsize=\small,frame=single,label=Example]
  !gapprompt@gap>| !gapinput@IsNormal(g,u);|
  false
\end{Verbatim}
 }

 

\subsection{\textcolor{Chapter }{IsCharacteristicSubgroup}}
\logpage{[ 39, 3, 7 ]}\nobreak
\hyperdef{L}{X8390B5117A10CC52}{}
{\noindent\textcolor{FuncColor}{$\triangleright$\enspace\texttt{IsCharacteristicSubgroup({\mdseries\slshape G, N})\index{IsCharacteristicSubgroup@\texttt{IsCharacteristicSubgroup}}
\label{IsCharacteristicSubgroup}
}\hfill{\scriptsize (operation)}}\\


 tests whether \mbox{\texttt{\mdseries\slshape N}} is invariant under all automorphisms of \mbox{\texttt{\mdseries\slshape G}}. 
\begin{Verbatim}[commandchars=!@|,fontsize=\small,frame=single,label=Example]
  !gapprompt@gap>| !gapinput@IsCharacteristicSubgroup(g,u);|
  false
\end{Verbatim}
 }

 

\subsection{\textcolor{Chapter }{ConjugateSubgroup}}
\logpage{[ 39, 3, 8 ]}\nobreak
\hyperdef{L}{X84F5464983655590}{}
{\noindent\textcolor{FuncColor}{$\triangleright$\enspace\texttt{ConjugateSubgroup({\mdseries\slshape G, g})\index{ConjugateSubgroup@\texttt{ConjugateSubgroup}}
\label{ConjugateSubgroup}
}\hfill{\scriptsize (operation)}}\\


 For a group \mbox{\texttt{\mdseries\slshape G}} which has a parent group \texttt{P} (see \texttt{Parent} (\ref{Parent})), returns the subgroup of \texttt{P}, obtained by conjugating \mbox{\texttt{\mdseries\slshape G}} using the conjugating element \mbox{\texttt{\mdseries\slshape g}}. 

 If \mbox{\texttt{\mdseries\slshape G}} has no parent group, it just delegates to the call to \texttt{ConjugateGroup} (\ref{ConjugateGroup}) with the same arguments. 

 To form a conjugate (subgroup) by any object acting via \texttt{\texttt{\symbol{94}}}, one can also use the infix operator \texttt{\texttt{\symbol{94}}}. }

 

\subsection{\textcolor{Chapter }{ConjugateSubgroups}}
\logpage{[ 39, 3, 9 ]}\nobreak
\hyperdef{L}{X7D9990EB837075A4}{}
{\noindent\textcolor{FuncColor}{$\triangleright$\enspace\texttt{ConjugateSubgroups({\mdseries\slshape G, U})\index{ConjugateSubgroups@\texttt{ConjugateSubgroups}}
\label{ConjugateSubgroups}
}\hfill{\scriptsize (operation)}}\\


 returns a list of all images of the group \mbox{\texttt{\mdseries\slshape U}} under conjugation action by \mbox{\texttt{\mdseries\slshape G}}. }

 

\subsection{\textcolor{Chapter }{IsSubnormal}}
\logpage{[ 39, 3, 10 ]}\nobreak
\hyperdef{L}{X82ABF80780CC27AF}{}
{\noindent\textcolor{FuncColor}{$\triangleright$\enspace\texttt{IsSubnormal({\mdseries\slshape G, U})\index{IsSubnormal@\texttt{IsSubnormal}}
\label{IsSubnormal}
}\hfill{\scriptsize (operation)}}\\


 A subgroup \mbox{\texttt{\mdseries\slshape U}} of the group \mbox{\texttt{\mdseries\slshape G}} is subnormal if it is contained in a subnormal series of \mbox{\texttt{\mdseries\slshape G}}. 
\begin{Verbatim}[commandchars=!@|,fontsize=\small,frame=single,label=Example]
  !gapprompt@gap>| !gapinput@IsSubnormal(g,Group((1,2,3)));|
  false
  !gapprompt@gap>| !gapinput@IsSubnormal(g,Group((1,2)(3,4)));|
  true
\end{Verbatim}
 }

 

\subsection{\textcolor{Chapter }{SubgroupByProperty}}
\logpage{[ 39, 3, 11 ]}\nobreak
\hyperdef{L}{X829766158665FB54}{}
{\noindent\textcolor{FuncColor}{$\triangleright$\enspace\texttt{SubgroupByProperty({\mdseries\slshape G, prop})\index{SubgroupByProperty@\texttt{SubgroupByProperty}}
\label{SubgroupByProperty}
}\hfill{\scriptsize (function)}}\\


 creates a subgroup of \mbox{\texttt{\mdseries\slshape G}} consisting of those elements fulfilling \mbox{\texttt{\mdseries\slshape prop}} (which is a tester function). No test is done whether the property actually
defines a subgroup. 

 Note that currently very little functionality beyond an element test exists
for groups created this way. }

 

\subsection{\textcolor{Chapter }{SubgroupShell}}
\logpage{[ 39, 3, 12 ]}\nobreak
\hyperdef{L}{X7E95101F80583E77}{}
{\noindent\textcolor{FuncColor}{$\triangleright$\enspace\texttt{SubgroupShell({\mdseries\slshape G})\index{SubgroupShell@\texttt{SubgroupShell}}
\label{SubgroupShell}
}\hfill{\scriptsize (function)}}\\


 creates a subgroup of \mbox{\texttt{\mdseries\slshape G}} which at this point is not yet specified further (but will be later, for
example by assigning a generating set). 
\begin{Verbatim}[commandchars=!@|,fontsize=\small,frame=single,label=Example]
  !gapprompt@gap>| !gapinput@u:=SubgroupByProperty(g,i->3^i=3);|
  <subgrp of Group([ (1,2,3,4), (1,2) ]) by property>
  !gapprompt@gap>| !gapinput@(1,3) in u; (1,4) in u; (1,5) in u;|
  false
  true
  false
  !gapprompt@gap>| !gapinput@GeneratorsOfGroup(u);|
  [ (1,2), (1,4,2) ]
  !gapprompt@gap>| !gapinput@u:=SubgroupShell(g);|
  <group>
\end{Verbatim}
 }

 }

  
\section{\textcolor{Chapter }{Closures of (Sub)groups}}\label{Closures of (Sub)groups}
\logpage{[ 39, 4, 0 ]}
\hyperdef{L}{X7B855B0485C3C6C5}{}
{
  

\subsection{\textcolor{Chapter }{ClosureGroup}}
\logpage{[ 39, 4, 1 ]}\nobreak
\hyperdef{L}{X7D13FC1F8576FFD8}{}
{\noindent\textcolor{FuncColor}{$\triangleright$\enspace\texttt{ClosureGroup({\mdseries\slshape G, obj})\index{ClosureGroup@\texttt{ClosureGroup}}
\label{ClosureGroup}
}\hfill{\scriptsize (operation)}}\\


 creates the group generated by the elements of \mbox{\texttt{\mdseries\slshape G}} and \mbox{\texttt{\mdseries\slshape obj}}. \mbox{\texttt{\mdseries\slshape obj}} can be either an element or a collection of elements, in particular another
group. 
\begin{Verbatim}[commandchars=!@|,fontsize=\small,frame=single,label=Example]
  !gapprompt@gap>| !gapinput@g:=SmallGroup(24,12);;u:=Subgroup(g,[g.3,g.4]);|
  Group([ f3, f4 ])
  !gapprompt@gap>| !gapinput@ClosureGroup(u,g.2);|
  Group([ f2, f3, f4 ])
  !gapprompt@gap>| !gapinput@ClosureGroup(u,[g.1,g.2]);|
  Group([ f1, f2, f3, f4 ])
  !gapprompt@gap>| !gapinput@ClosureGroup(u,Group(g.2*g.1));|
  Group([ f1*f2^2, f3, f4 ])
\end{Verbatim}
 }

 

\subsection{\textcolor{Chapter }{ClosureGroupAddElm}}
\logpage{[ 39, 4, 2 ]}\nobreak
\hyperdef{L}{X81A20A397C308483}{}
{\noindent\textcolor{FuncColor}{$\triangleright$\enspace\texttt{ClosureGroupAddElm({\mdseries\slshape G, elm})\index{ClosureGroupAddElm@\texttt{ClosureGroupAddElm}}
\label{ClosureGroupAddElm}
}\hfill{\scriptsize (function)}}\\
\noindent\textcolor{FuncColor}{$\triangleright$\enspace\texttt{ClosureGroupCompare({\mdseries\slshape G, elm})\index{ClosureGroupCompare@\texttt{ClosureGroupCompare}}
\label{ClosureGroupCompare}
}\hfill{\scriptsize (function)}}\\
\noindent\textcolor{FuncColor}{$\triangleright$\enspace\texttt{ClosureGroupIntest({\mdseries\slshape G, elm})\index{ClosureGroupIntest@\texttt{ClosureGroupIntest}}
\label{ClosureGroupIntest}
}\hfill{\scriptsize (function)}}\\


 These three functions together with \texttt{ClosureGroupDefault} (\ref{ClosureGroupDefault}) implement the main methods for \texttt{ClosureGroup} (\ref{ClosureGroup}). In the ordering given, they just add \mbox{\texttt{\mdseries\slshape elm}} to the generators, remove duplicates and identity elements, and test whether \mbox{\texttt{\mdseries\slshape elm}} is already contained in \mbox{\texttt{\mdseries\slshape G}}. }

 

\subsection{\textcolor{Chapter }{ClosureGroupDefault}}
\logpage{[ 39, 4, 3 ]}\nobreak
\hyperdef{L}{X82F59F6680D1B0D5}{}
{\noindent\textcolor{FuncColor}{$\triangleright$\enspace\texttt{ClosureGroupDefault({\mdseries\slshape G, elm})\index{ClosureGroupDefault@\texttt{ClosureGroupDefault}}
\label{ClosureGroupDefault}
}\hfill{\scriptsize (function)}}\\


 This functions returns the closure of the group \mbox{\texttt{\mdseries\slshape G}} with the element \mbox{\texttt{\mdseries\slshape elm}}. If \mbox{\texttt{\mdseries\slshape G}} has the attribute \texttt{AsSSortedList} (\ref{AsSSortedList}) then also the result has this attribute. This is used to implement the default
method for \texttt{Enumerator} (\ref{Enumerator}) and \texttt{EnumeratorSorted} (\ref{EnumeratorSorted}). }

 

\subsection{\textcolor{Chapter }{ClosureSubgroup}}
\logpage{[ 39, 4, 4 ]}\nobreak
\hyperdef{L}{X7A7AF14A8052F055}{}
{\noindent\textcolor{FuncColor}{$\triangleright$\enspace\texttt{ClosureSubgroup({\mdseries\slshape G, obj})\index{ClosureSubgroup@\texttt{ClosureSubgroup}}
\label{ClosureSubgroup}
}\hfill{\scriptsize (function)}}\\
\noindent\textcolor{FuncColor}{$\triangleright$\enspace\texttt{ClosureSubgroupNC({\mdseries\slshape G, obj})\index{ClosureSubgroupNC@\texttt{ClosureSubgroupNC}}
\label{ClosureSubgroupNC}
}\hfill{\scriptsize (function)}}\\


 For a group \mbox{\texttt{\mdseries\slshape G}} that stores a parent group (see{\nobreakspace}\ref{Parents}), \texttt{ClosureSubgroup} calls \texttt{ClosureGroup} (\ref{ClosureGroup}) with the same arguments; if the result is a subgroup of the parent of \mbox{\texttt{\mdseries\slshape G}} then the parent of \mbox{\texttt{\mdseries\slshape G}} is set as parent of the result, otherwise an error is raised. The check
whether the result is contained in the parent of \mbox{\texttt{\mdseries\slshape G}} is omitted by the \texttt{NC} version. As a wrong parent might imply wrong properties this version should be
used with care. }

 }

  
\section{\textcolor{Chapter }{Expressing Group Elements as Words in Generators}}\label{Expressing Group Elements as Words in Generators}
\logpage{[ 39, 5, 0 ]}
\hyperdef{L}{X7E19F92284F6684E}{}
{
  \index{factorization} \index{words!in generators} Using homomorphisms (see chapter{\nobreakspace}\ref{Group Homomorphisms}) it is possible to express group elements as words in given generators:
Create a free group (see{\nobreakspace}\texttt{FreeGroup} (\ref{FreeGroup:for given rank})) on the correct number of generators and create a homomorphism from this free
group onto the group \mbox{\texttt{\mdseries\slshape G}} in whose generators you want to factorize. Then the preimage of an element of \mbox{\texttt{\mdseries\slshape G}} is a word in the free generators, that will map on this element again. 

 

\subsection{\textcolor{Chapter }{EpimorphismFromFreeGroup}}
\logpage{[ 39, 5, 1 ]}\nobreak
\hyperdef{L}{X7FE8A3B08458A1BF}{}
{\noindent\textcolor{FuncColor}{$\triangleright$\enspace\texttt{EpimorphismFromFreeGroup({\mdseries\slshape G})\index{EpimorphismFromFreeGroup@\texttt{EpimorphismFromFreeGroup}}
\label{EpimorphismFromFreeGroup}
}\hfill{\scriptsize (attribute)}}\\


 For a group \mbox{\texttt{\mdseries\slshape G}} with a known generating set, this attribute returns a homomorphism from a free
group that maps the free generators to the generators of \mbox{\texttt{\mdseries\slshape G}}. 

 The option \texttt{names} can be used to prescribe a (print) name for the free generators. 

 The following example shows how to decompose elements of $S_4$ in the generators \texttt{(1,2,3,4)} and \texttt{(1,2)}: 

 
\begin{Verbatim}[commandchars=!@|,fontsize=\small,frame=single,label=Example]
  !gapprompt@gap>| !gapinput@g:=Group((1,2,3,4),(1,2));|
  Group([ (1,2,3,4), (1,2) ])
  !gapprompt@gap>| !gapinput@hom:=EpimorphismFromFreeGroup(g:names:=["x","y"]);|
  [ x, y ] -> [ (1,2,3,4), (1,2) ]
  !gapprompt@gap>| !gapinput@PreImagesRepresentative(hom,(1,4));|
  y^-1*x^-1*(x^-1*y^-1)^2*x
\end{Verbatim}
 

 The following example stems from a real request to the \textsf{GAP} Forum. In September 2000 a \textsf{GAP} user working with puzzles wanted to express the permutation \texttt{(1,2)} as a word as short as possible in particular generators of the symmetric group $S_{16}$. 

 
\begin{Verbatim}[commandchars=!@|,fontsize=\small,frame=single,label=Example]
  !gapprompt@gap>| !gapinput@perms := [ (1,2,3,7,11,10,9,5), (2,3,4,8,12,11,10,6),|
  !gapprompt@>| !gapinput@  (5,6,7,11,15,14,13,9), (6,7,8,12,16,15,14,10) ];;|
  !gapprompt@gap>| !gapinput@puzzle := Group( perms );;Size( puzzle );|
  20922789888000
  !gapprompt@gap>| !gapinput@hom:=EpimorphismFromFreeGroup(puzzle:names:=["a", "b", "c", "d"]);;|
  !gapprompt@gap>| !gapinput@word := PreImagesRepresentative( hom, (1,2) );|
  a^-1*c*b*c^-1*a*b^-1*a^-2*c^-1*a*b^-1*c*b
  !gapprompt@gap>| !gapinput@Length( word );|
  13
\end{Verbatim}
 }

  

\subsection{\textcolor{Chapter }{Factorization}}
\logpage{[ 39, 5, 2 ]}\nobreak
\hyperdef{L}{X8357294D7B164106}{}
{\noindent\textcolor{FuncColor}{$\triangleright$\enspace\texttt{Factorization({\mdseries\slshape G, elm})\index{Factorization@\texttt{Factorization}}
\label{Factorization}
}\hfill{\scriptsize (operation)}}\\


 returns a factorization of \mbox{\texttt{\mdseries\slshape elm}} as word in the generators of the group \mbox{\texttt{\mdseries\slshape G}} given in the attribute \texttt{GeneratorsOfGroup} (\ref{GeneratorsOfGroup}). The attribute \texttt{EpimorphismFromFreeGroup} (\ref{EpimorphismFromFreeGroup}) of \mbox{\texttt{\mdseries\slshape G}} will contain a map from the free group in which the word is expressed to the
group \mbox{\texttt{\mdseries\slshape G}}. The attribute \texttt{MappingGeneratorsImages} (\ref{MappingGeneratorsImages}) of this map gives a list of generators and corresponding letters. 

 The algorithm used forms all elements of the group to ensure a short word is
found. Therefore this function should \emph{not} be used when the group \mbox{\texttt{\mdseries\slshape G}} has more than a few million elements. Because of this, one should not call
this function within algorithms, but use homomorphisms instead. 
\begin{Verbatim}[commandchars=!@|,fontsize=\small,frame=single,label=Example]
  !gapprompt@gap>| !gapinput@G:=SymmetricGroup( 6 );;|
  !gapprompt@gap>| !gapinput@r:=(3,4);; s:=(1,2,3,4,5,6);;|
  !gapprompt@gap>| !gapinput@# create subgroup to force the system to use the generators r and s:|
  !gapprompt@gap>| !gapinput@H:= Subgroup(G, [ r, s ] );|
  Group([ (3,4), (1,2,3,4,5,6) ])
  !gapprompt@gap>| !gapinput@Factorization( H, (1,2,3) );|
  (x2*x1)^2*x2^-2
  !gapprompt@gap>| !gapinput@s*r*s*r*s^-2;|
  (1,2,3)
  !gapprompt@gap>| !gapinput@MappingGeneratorsImages(EpimorphismFromFreeGroup(H));|
  [ [ x1, x2 ], [ (3,4), (1,2,3,4,5,6) ] ]
\end{Verbatim}
 }

 

\subsection{\textcolor{Chapter }{GrowthFunctionOfGroup}}
\logpage{[ 39, 5, 3 ]}\nobreak
\hyperdef{L}{X871508DD808EB487}{}
{\noindent\textcolor{FuncColor}{$\triangleright$\enspace\texttt{GrowthFunctionOfGroup({\mdseries\slshape G})\index{GrowthFunctionOfGroup@\texttt{GrowthFunctionOfGroup}}
\label{GrowthFunctionOfGroup}
}\hfill{\scriptsize (attribute)}}\\
\noindent\textcolor{FuncColor}{$\triangleright$\enspace\texttt{GrowthFunctionOfGroup({\mdseries\slshape G, radius})\index{GrowthFunctionOfGroup@\texttt{GrowthFunctionOfGroup}!with word length limit}
\label{GrowthFunctionOfGroup:with word length limit}
}\hfill{\scriptsize (operation)}}\\


 For a group \mbox{\texttt{\mdseries\slshape G}} with a generating set given in \texttt{GeneratorsOfGroup} (\ref{GeneratorsOfGroup}), this function calculates the number of elements whose shortest expression as
words in the generating set is of a particular length. It returns a list \mbox{\texttt{\mdseries\slshape L}}, whose $i+1$ entry counts the number of elements whose shortest word expression has length $i$. If a maximal length \mbox{\texttt{\mdseries\slshape radius}} is given, only words up to length \mbox{\texttt{\mdseries\slshape radius}} are counted. Otherwise the group must be finite and all elements are
enumerated. 
\begin{Verbatim}[commandchars=!@|,fontsize=\small,frame=single,label=Example]
  !gapprompt@gap>| !gapinput@GrowthFunctionOfGroup(MathieuGroup(12));|
  [ 1, 5, 19, 70, 255, 903, 3134, 9870, 25511, 38532, 16358, 382 ]
  !gapprompt@gap>| !gapinput@GrowthFunctionOfGroup(MathieuGroup(12),2);|
  [ 1, 5, 19 ]
  !gapprompt@gap>| !gapinput@GrowthFunctionOfGroup(MathieuGroup(12),99);|
  [ 1, 5, 19, 70, 255, 903, 3134, 9870, 25511, 38532, 16358, 382 ]
  !gapprompt@gap>| !gapinput@free:=FreeGroup("a","b");|
  <free group on the generators [ a, b ]>
  !gapprompt@gap>| !gapinput@product:=free/ParseRelators(free,"a2,b3");|
  <fp group on the generators [ a, b ]>
  !gapprompt@gap>| !gapinput@SetIsFinite(product,false);|
  !gapprompt@gap>| !gapinput@GrowthFunctionOfGroup(product,10);|
  [ 1, 3, 4, 6, 8, 12, 16, 24, 32, 48, 64 ]
\end{Verbatim}
 }

 }

  
\section{\textcolor{Chapter }{Structure Descriptions}}\label{Structure Descriptions}
\logpage{[ 39, 6, 0 ]}
\hyperdef{L}{X87BF1B887C91CA2E}{}
{
  

\subsection{\textcolor{Chapter }{StructureDescription}}
\logpage{[ 39, 6, 1 ]}\nobreak
\hyperdef{L}{X8199B74B84446971}{}
{\noindent\textcolor{FuncColor}{$\triangleright$\enspace\texttt{StructureDescription({\mdseries\slshape G})\index{StructureDescription@\texttt{StructureDescription}}
\label{StructureDescription}
}\hfill{\scriptsize (attribute)}}\\


 The method for \texttt{StructureDescription} exhibits a structure of the given group \mbox{\texttt{\mdseries\slshape G}} to some extent, using the strategy outlined below. The idea is to return a
possibly short string which gives some insight in the structure of the
considered group. It is intended primarily for small groups (order less than
100) or groups with few normal subgroups, in other cases, in particular large $p$\texttt{\symbol{45}}groups, it can be very costly. Furthermore, the string
returned is \texttt{\symbol{45}}\texttt{\symbol{45}} as the action on chief
factors is not described \texttt{\symbol{45}}\texttt{\symbol{45}} often not
the most useful way to describe a group.

 The string returned by \texttt{StructureDescription} is \textsc{not} an isomorphism invariant: non\texttt{\symbol{45}}isomorphic groups can have
the same string value, and two isomorphic groups in different representations
can produce different strings. The value returned by \texttt{StructureDescription} is a string of the following form: 

 
\begin{Verbatim}[commandchars=!@I,fontsize=\small,frame=single,label=]
    StructureDescription(<G>) ::=
        1                                 ; trivial group
      | C<size>                           ; finite cyclic group
      | Z                                 ; infinite cyclic group
      | A<degree>                         ; alternating group
      | S<degree>                         ; symmetric group
      | D<size>                           ; dihedral group
      | Q<size>                           ; quaternion group
      | QD<size>                          ; quasidihedral group
      | PSL(<n>,<q>)                      ; projective special linear group
      | SL(<n>,<q>)                       ; special linear group
      | GL(<n>,<q>)                       ; general linear group
      | PSU(<n>,<q>)                      ; proj. special unitary group
      | O(2<n>+1,<q>)                     ; orthogonal group, type B
      | O+(2<n>,<q>)                      ; orthogonal group, type D
      | O-(2<n>,<q>)                      ; orthogonal group, type 2D
      | PSp(2<n>,<q>)                     ; proj. special symplectic group
      | Sz(<q>)                           ; Suzuki group
      | Ree(<q>)                          ; Ree group (type 2F or 2G)
      | E(6,<q>) | E(7,<q>) | E(8,<q>)    ; Lie group of exceptional type
      | 2E(6,<q>) | F(4,<q>) | G(2,<q>)
      | 3D(4,<q>)                         ; Steinberg triality group
      | M11 | M12 | M22 | M23 | M24
      | J1 | J2 | J3 | J4 | Co1 | Co2
      | Co3 | Fi22 | Fi23 | Fi24' | Suz
      | HS | McL | He | HN | Th | B
      | M | ON | Ly | Ru                  ; sporadic simple group
      | 2F(4,2)'                          ; Tits group
      | PerfectGroup(<size>,<id>)         ; the indicated group from the
                                          ; library of perfect groups
      | A x B                             ; direct product
      | N : H                             ; semidirect product
      | C(G) . G/C(G) = G' . G/G'         ; non-split extension
                                          ; (equal alternatives and
                                          ; trivial extensions omitted)
      | Phi(G) . G/Phi(G)                 ; non-split extension:
                                          ; Frattini subgroup and
                                          ; Frattini factor group
\end{Verbatim}
 

 Note that the \texttt{StructureDescription} is only \emph{one} possible way of building up the given group from smaller pieces. 

 The option ``short'' is recognized \texttt{\symbol{45}} if this option is set, an abbreviated
output format is used (e.g. \texttt{"6x3"} instead of \texttt{"C6 x C3"}). 

 If the \texttt{Name} (\ref{Name}) attribute is not bound, but \texttt{StructureDescription} is, \texttt{View} (\ref{View}) prints the value of the attribute \texttt{StructureDescription}. The \texttt{Print} (\ref{Print})ed representation of a group is not affected by computing a \texttt{StructureDescription}. 

 The strategy used to compute a \texttt{StructureDescription} is as follows: 

 
\begin{description}
\item[{1.}]  Lookup in a precomputed list, if the order of \mbox{\texttt{\mdseries\slshape G}} is not larger than 100 and not equal to 64 or 96. 
\item[{2.}]  If \mbox{\texttt{\mdseries\slshape G}} is abelian, then decompose it into cyclic factors in ``elementary divisors style''. For example, \texttt{"C2 x C3 x C3"} is \texttt{"C6 x C3"}. For infinite abelian groups, \texttt{"Z"} denotes the group of integers. 
\item[{3.}]  Recognize alternating groups, symmetric groups, dihedral groups, quasidihedral
groups, quaternion groups, PSL's, SL's, GL's and simple groups not listed so
far as basic building blocks. 
\item[{4.}]  Decompose \mbox{\texttt{\mdseries\slshape G}} into a direct product of irreducible factors. 
\item[{5.}]  Recognize semidirect products \mbox{\texttt{\mdseries\slshape G}}=$N$:$H$, where $N$ is normal. Select a pair $N$, $H$ with the following preferences: 
\begin{description}
\item[{1.}]  if \mbox{\texttt{\mdseries\slshape G}} is defined as a semidirect product of $N$, $H$ then select $N$, $H$, 
\item[{2.}]  if \mbox{\texttt{\mdseries\slshape G}} is solvable, then select a solvable normal Hall subgroup $N$, if exists, and consider the semidirect decomposition of $N$ and $\mbox{\texttt{\mdseries\slshape G}}/N$, 
\item[{3.}]  find any nontrivial normal subgroup $N$ which has a complement $H$. 
\end{description}
 The option ``nice'' is recognized. If this option is set, then all semidirect products are
computed in order to find a possibly nicer presentation. Note, that this may
take a very long time if \mbox{\texttt{\mdseries\slshape G}} has many normal subgroups, e.g. if $\mbox{\texttt{\mdseries\slshape G}}/\mbox{\texttt{\mdseries\slshape G}}'$ has many cyclic factors. If the option ``nice'' is set, then GAP would select a pair $N$, $H$ with the following preferences: 
\begin{description}
\item[{1.}]  $H$ is abelian 
\item[{2.}]  $N$ is abelian 
\item[{2a.}]  $N$ has many abelian invariants 
\item[{3.}]  $N$ is a direct product 
\item[{3a.}]  $N$ has many direct factors 
\item[{4.}]  $\phi: H \rightarrow$ Aut($N$), $h \mapsto (n \mapsto n^h)$ is injective. 
\end{description}
 
\item[{6.}]  Fall back to non\texttt{\symbol{45}}splitting extensions: If the centre or the
commutator factor group is non\texttt{\symbol{45}}trivial, write \mbox{\texttt{\mdseries\slshape G}} as $Z(\mbox{\texttt{\mdseries\slshape G}})$.$\mbox{\texttt{\mdseries\slshape G}}/Z(\mbox{\texttt{\mdseries\slshape G}})$ or $\mbox{\texttt{\mdseries\slshape G}}'$.$\mbox{\texttt{\mdseries\slshape G}}/\mbox{\texttt{\mdseries\slshape G}}'$, respectively. Otherwise if the Frattini subgroup is
non\texttt{\symbol{45}}trivial, write \mbox{\texttt{\mdseries\slshape G}} as $\Phi$(\mbox{\texttt{\mdseries\slshape G}}).\mbox{\texttt{\mdseries\slshape G}}/$\Phi$(\mbox{\texttt{\mdseries\slshape G}}). 
\item[{7.}]  If no decomposition is found (maybe this is not the case for any finite
group), try to identify \mbox{\texttt{\mdseries\slshape G}} in the perfect groups library. If this fails also, then return a string
describing this situation. 
\end{description}
 Note that \texttt{StructureDescription} is \emph{not} intended to be a research tool, but rather an educational tool. The reasons
for this are as follows: 
\begin{description}
\item[{1.}]  ``Most'' groups do not have ``nice'' decompositions. This is in some contrast to what is often taught in elementary
courses on group theory, where it is sometimes suggested that basically every
group can be written as iterated direct or semidirect product of cyclic groups
and nonabelian simple groups. 
\item[{2.}]  In particular many $p$\texttt{\symbol{45}}groups have very ``similar'' structure, and \texttt{StructureDescription} can only exhibit a little of it. Changing this would likely make the output
not essentially easier to read than a pc presentation. 
\end{description}
 
\begin{Verbatim}[commandchars=!@|,fontsize=\small,frame=single,label=Example]
  !gapprompt@gap>| !gapinput@l := AllSmallGroups(12);;|
  !gapprompt@gap>| !gapinput@List(l,StructureDescription);; l;|
  [ C3 : C4, C12, A4, D12, C6 x C2 ]
  !gapprompt@gap>| !gapinput@List(AllSmallGroups(40),G->StructureDescription(G:short));|
  [ "5:8", "40", "5:8", "5:Q8", "4xD10", "D40", "2x(5:4)", "(10x2):2",
    "20x2", "5xD8", "5xQ8", "2x(5:4)", "2^2xD10", "10x2^2" ]
  !gapprompt@gap>| !gapinput@List(AllTransitiveGroups(DegreeAction,6),|
  !gapprompt@>| !gapinput@        G->StructureDescription(G:short));|
  [ "6", "S3", "D12", "A4", "3xS3", "2xA4", "S4", "S4", "S3xS3",
    "(3^2):4", "2xS4", "A5", "(S3xS3):2", "S5", "A6", "S6" ]
  !gapprompt@gap>| !gapinput@StructureDescription(SmallGroup(504,7));|
  "C7 : (C9 x Q8)"
  !gapprompt@gap>| !gapinput@StructureDescription(SmallGroup(504,7):nice);|
  "(C7 : Q8) : C9"
  !gapprompt@gap>| !gapinput@StructureDescription(AbelianGroup([0,2,3]));|
  "Z x C6"
  !gapprompt@gap>| !gapinput@StructureDescription(AbelianGroup([0,0,0,2,3,6]):short);|
  "Z^3x6^2"
  !gapprompt@gap>| !gapinput@StructureDescription(PSL(4,2));|
  "A8"
\end{Verbatim}
 }

 }

  
\section{\textcolor{Chapter }{Cosets}}\label{Cosets}
\logpage{[ 39, 7, 0 ]}
\hyperdef{L}{X81002AA87DDBC02F}{}
{
  \index{right cosets} \index{coset} 

\subsection{\textcolor{Chapter }{RightCoset}}
\logpage{[ 39, 7, 1 ]}\nobreak
\hyperdef{L}{X8412ABD57986B9FC}{}
{\noindent\textcolor{FuncColor}{$\triangleright$\enspace\texttt{RightCoset({\mdseries\slshape U, g})\index{RightCoset@\texttt{RightCoset}}
\label{RightCoset}
}\hfill{\scriptsize (operation)}}\\


 returns the right coset of \mbox{\texttt{\mdseries\slshape U}} with representative \mbox{\texttt{\mdseries\slshape g}}, which is the set of all elements of the form $ug$ for all $u \in \mbox{\texttt{\mdseries\slshape U}}$. \mbox{\texttt{\mdseries\slshape g}} must be an element of a larger group \mbox{\texttt{\mdseries\slshape G}} which contains \mbox{\texttt{\mdseries\slshape U}}. For element operations such as \texttt{in} a right coset behaves like a set of group elements. 

 Right cosets are external orbits for the action of \mbox{\texttt{\mdseries\slshape U}} which acts via \texttt{OnLeftInverse} (\ref{OnLeftInverse}). Of course the action of a larger group \mbox{\texttt{\mdseries\slshape G}} on right cosets is via \texttt{OnRight} (\ref{OnRight}). 
\begin{Verbatim}[commandchars=!@|,fontsize=\small,frame=single,label=Example]
  !gapprompt@gap>| !gapinput@u:=Group((1,2,3), (1,2));;|
  !gapprompt@gap>| !gapinput@c:=RightCoset(u,(2,3,4));|
  RightCoset(Group( [ (1,2,3), (1,2) ] ),(2,3,4))
  !gapprompt@gap>| !gapinput@ActingDomain(c);|
  Group([ (1,2,3), (1,2) ])
  !gapprompt@gap>| !gapinput@Representative(c);|
  (2,3,4)
  !gapprompt@gap>| !gapinput@Size(c);|
  6
  !gapprompt@gap>| !gapinput@AsList(c);|
  [ (2,3,4), (1,4,2), (1,3,4,2), (1,3)(2,4), (2,4), (1,4,2,3) ]
  !gapprompt@gap>| !gapinput@IsBiCoset(c);|
  false
\end{Verbatim}
 }

 

\subsection{\textcolor{Chapter }{RightCosets}}
\logpage{[ 39, 7, 2 ]}\nobreak
\hyperdef{L}{X835F48248571364F}{}
{\noindent\textcolor{FuncColor}{$\triangleright$\enspace\texttt{RightCosets({\mdseries\slshape G, U})\index{RightCosets@\texttt{RightCosets}}
\label{RightCosets}
}\hfill{\scriptsize (function)}}\\
\noindent\textcolor{FuncColor}{$\triangleright$\enspace\texttt{RightCosetsNC({\mdseries\slshape G, U})\index{RightCosetsNC@\texttt{RightCosetsNC}}
\label{RightCosetsNC}
}\hfill{\scriptsize (operation)}}\\


 computes a duplicate free list of right cosets \mbox{\texttt{\mdseries\slshape U}} $g$ for $g \in$ \mbox{\texttt{\mdseries\slshape G}}. A set of representatives for the elements in this list forms a right
transversal of \mbox{\texttt{\mdseries\slshape U}} in \mbox{\texttt{\mdseries\slshape G}}. (By inverting the representatives one obtains a list of representatives of
the left cosets of \mbox{\texttt{\mdseries\slshape U}}.) The \texttt{NC} version does not check whether \mbox{\texttt{\mdseries\slshape U}} is a subgroup of \mbox{\texttt{\mdseries\slshape G}}. 
\begin{Verbatim}[commandchars=!@|,fontsize=\small,frame=single,label=Example]
  !gapprompt@gap>| !gapinput@RightCosets(g,u);|
  [ RightCoset(Group( [ (1,2,3), (1,2) ] ),()),
    RightCoset(Group( [ (1,2,3), (1,2) ] ),(1,3)(2,4)),
    RightCoset(Group( [ (1,2,3), (1,2) ] ),(1,4)(2,3)),
    RightCoset(Group( [ (1,2,3), (1,2) ] ),(1,2)(3,4)) ]
\end{Verbatim}
 }

 

\subsection{\textcolor{Chapter }{CanonicalRightCosetElement}}
\logpage{[ 39, 7, 3 ]}\nobreak
\hyperdef{L}{X85884F177B5D98AE}{}
{\noindent\textcolor{FuncColor}{$\triangleright$\enspace\texttt{CanonicalRightCosetElement({\mdseries\slshape U, g})\index{CanonicalRightCosetElement@\texttt{CanonicalRightCosetElement}}
\label{CanonicalRightCosetElement}
}\hfill{\scriptsize (operation)}}\\


 returns a ``canonical'' representative of the right coset \mbox{\texttt{\mdseries\slshape U}} \mbox{\texttt{\mdseries\slshape g}} which is independent of the given representative \mbox{\texttt{\mdseries\slshape g}}. This can be used to compare cosets by comparing their canonical
representatives. 

 The representative chosen to be the ``canonical'' one is representation dependent and only guaranteed to remain the same within
one \textsf{GAP} session. 
\begin{Verbatim}[commandchars=!@|,fontsize=\small,frame=single,label=Example]
  !gapprompt@gap>| !gapinput@CanonicalRightCosetElement(u,(2,4,3));|
  (3,4)
\end{Verbatim}
 }

 

\subsection{\textcolor{Chapter }{IsRightCoset}}
\logpage{[ 39, 7, 4 ]}\nobreak
\hyperdef{L}{X7D7625A1861D9DAB}{}
{\noindent\textcolor{FuncColor}{$\triangleright$\enspace\texttt{IsRightCoset({\mdseries\slshape obj})\index{IsRightCoset@\texttt{IsRightCoset}}
\label{IsRightCoset}
}\hfill{\scriptsize (Category)}}\\


 The category of right cosets. 

 \index{left cosets} \textsf{GAP} does not provide left cosets as a separate data type, but as the left coset $gU$ consists of exactly the inverses of the elements of the right coset $Ug^{{-1}}$ calculations with left cosets can be emulated using right cosets by inverting
the representatives. }

 

\subsection{\textcolor{Chapter }{IsBiCoset}}
\logpage{[ 39, 7, 5 ]}\nobreak
\hyperdef{L}{X78F4F0D8838F5ABF}{}
{\noindent\textcolor{FuncColor}{$\triangleright$\enspace\texttt{IsBiCoset({\mdseries\slshape C})\index{IsBiCoset@\texttt{IsBiCoset}}
\label{IsBiCoset}
}\hfill{\scriptsize (property)}}\\


 \index{bicoset} A (right) coset $Ug$ is considered a \emph{bicoset} if its set of elements simultaneously forms a left coset for the same
subgroup. This is the case if and only if the coset representative $g$ normalizes the subgroup $U$. }

 

\subsection{\textcolor{Chapter }{CosetDecomposition}}
\logpage{[ 39, 7, 6 ]}\nobreak
\hyperdef{L}{X82F6ABE378B928D1}{}
{\noindent\textcolor{FuncColor}{$\triangleright$\enspace\texttt{CosetDecomposition({\mdseries\slshape G, S})\index{CosetDecomposition@\texttt{CosetDecomposition}}
\label{CosetDecomposition}
}\hfill{\scriptsize (function)}}\\


 For a finite group \mbox{\texttt{\mdseries\slshape G}} and a subgroup $\mbox{\texttt{\mdseries\slshape S}} \leq \mbox{\texttt{\mdseries\slshape G}}$ this function returns a partition of the elements of \mbox{\texttt{\mdseries\slshape G}} according to the (right) cosets of \mbox{\texttt{\mdseries\slshape S}}. The result is a list of lists, each sublist corresponding to one coset. The
first sublist is the elements list of the subgroup, the other lists are
arranged accordingly. 
\begin{Verbatim}[commandchars=!@|,fontsize=\small,frame=single,label=Example]
  !gapprompt@gap>| !gapinput@CosetDecomposition(SymmetricGroup(4),SymmetricGroup(3));|
  [ [ (), (2,3), (1,2), (1,2,3), (1,3,2), (1,3) ],
    [ (1,4), (1,4)(2,3), (1,2,4), (1,2,3,4), (1,3,2,4), (1,3,4) ],
    [ (1,4,2), (1,4,2,3), (2,4), (2,3,4), (1,3)(2,4), (1,3,4,2) ],
    [ (1,4,3), (1,4,3,2), (1,2,4,3), (1,2)(3,4), (2,4,3), (3,4) ] ]
\end{Verbatim}
 }

 }

  
\section{\textcolor{Chapter }{Transversals}}\label{Transversals}
\logpage{[ 39, 8, 0 ]}
\hyperdef{L}{X83C723878230D616}{}
{
  

\subsection{\textcolor{Chapter }{RightTransversal}}
\logpage{[ 39, 8, 1 ]}\nobreak
\hyperdef{L}{X85C65D06822E716F}{}
{\noindent\textcolor{FuncColor}{$\triangleright$\enspace\texttt{RightTransversal({\mdseries\slshape G, U})\index{RightTransversal@\texttt{RightTransversal}}
\label{RightTransversal}
}\hfill{\scriptsize (operation)}}\\


 A right transversal $t$ is a list of representatives for the set $\mbox{\texttt{\mdseries\slshape U}} \setminus \mbox{\texttt{\mdseries\slshape G}}$ of right cosets (consisting of cosets $Ug$) of $U$ in $G$. 

 The object returned by \texttt{RightTransversal} is not a plain list, but an object that behaves like an immutable list of
length $[\mbox{\texttt{\mdseries\slshape G}}:\mbox{\texttt{\mdseries\slshape U}}]$, except if \mbox{\texttt{\mdseries\slshape U}} is the trivial subgroup of \mbox{\texttt{\mdseries\slshape G}} in which case \texttt{RightTransversal} may return the sorted plain list of coset representatives. 

 The operation \texttt{PositionCanonical} (\ref{PositionCanonical}), called for a transversal $t$ and an element $g$ of \mbox{\texttt{\mdseries\slshape G}}, will return the position of the representative in $t$ that lies in the same coset of \mbox{\texttt{\mdseries\slshape U}} as the element $g$ does. (In comparison, \texttt{Position} (\ref{Position}) will return \texttt{fail} if the element is not equal to the representative.) Functions that implement
group actions such as \texttt{Action} (\ref{Action:for a group, an action domain, etc.}) or \texttt{Permutation} (\ref{Permutation:for a group, an action domain, etc.}) (see Chapter{\nobreakspace}\ref{Group Actions}) use \texttt{PositionCanonical} (\ref{PositionCanonical}), therefore it is possible to ``act'' on a right transversal to implement the action on the cosets. This is often
much more efficient than acting on cosets. 
\begin{Verbatim}[commandchars=!@|,fontsize=\small,frame=single,label=Example]
  !gapprompt@gap>| !gapinput@g:=Group((1,2,3,4),(1,2));;|
  !gapprompt@gap>| !gapinput@u:=Subgroup(g,[(1,2,3),(1,2)]);;|
  !gapprompt@gap>| !gapinput@rt:=RightTransversal(g,u);|
  RightTransversal(Group([ (1,2,3,4), (1,2) ]),Group([ (1,2,3), (1,2) ]))
  !gapprompt@gap>| !gapinput@Length(rt);|
  4
  !gapprompt@gap>| !gapinput@Position(rt,(1,2,3));|
  fail
\end{Verbatim}
 

 Note that the elements of a right transversal are not necessarily ``canonical'' in the sense of \texttt{CanonicalRightCosetElement} (\ref{CanonicalRightCosetElement}), but we may compute a list of canonical coset representatives by calling that
function. (See also \texttt{PositionCanonical} (\ref{PositionCanonical}).) 

 
\begin{Verbatim}[commandchars=!@|,fontsize=\small,frame=single,label=Example]
  !gapprompt@gap>| !gapinput@List(RightTransversal(g,u),i->CanonicalRightCosetElement(u,i));|
  [ (), (2,3,4), (1,2,3,4), (3,4) ]
  !gapprompt@gap>| !gapinput@PositionCanonical(rt,(1,2,3));|
  1
  !gapprompt@gap>| !gapinput@rt[1];|
  ()
\end{Verbatim}
 }

 }

  
\section{\textcolor{Chapter }{Double Cosets}}\label{Double Cosets}
\logpage{[ 39, 9, 0 ]}
\hyperdef{L}{X78B98B257E981046}{}
{
  

\subsection{\textcolor{Chapter }{DoubleCoset}}
\logpage{[ 39, 9, 1 ]}\nobreak
\hyperdef{L}{X7E51ED757D17254B}{}
{\noindent\textcolor{FuncColor}{$\triangleright$\enspace\texttt{DoubleCoset({\mdseries\slshape U, g, V})\index{DoubleCoset@\texttt{DoubleCoset}}
\label{DoubleCoset}
}\hfill{\scriptsize (operation)}}\\


 The groups \mbox{\texttt{\mdseries\slshape U}} and \mbox{\texttt{\mdseries\slshape V}} must be subgroups of a common supergroup \mbox{\texttt{\mdseries\slshape G}} of which \mbox{\texttt{\mdseries\slshape g}} is an element. This command constructs the double coset \mbox{\texttt{\mdseries\slshape U}} \mbox{\texttt{\mdseries\slshape g}} \mbox{\texttt{\mdseries\slshape V}} which is the set of all elements of the form $ugv$ for any $u \in \mbox{\texttt{\mdseries\slshape U}}$, $v \in \mbox{\texttt{\mdseries\slshape V}}$. For element operations such as \texttt{in}, a double coset behaves like a set of group elements. The double coset stores \mbox{\texttt{\mdseries\slshape U}} in the attribute \texttt{LeftActingGroup}, \mbox{\texttt{\mdseries\slshape g}} as \texttt{Representative} (\ref{Representative}), and \mbox{\texttt{\mdseries\slshape V}} as \texttt{RightActingGroup}. }

 

\subsection{\textcolor{Chapter }{RepresentativesContainedRightCosets}}
\logpage{[ 39, 9, 2 ]}\nobreak
\hyperdef{L}{X7F53DABD79BA4F72}{}
{\noindent\textcolor{FuncColor}{$\triangleright$\enspace\texttt{RepresentativesContainedRightCosets({\mdseries\slshape D})\index{RepresentativesContainedRightCosets@\texttt{RepresentativesContainedRightCosets}}
\label{RepresentativesContainedRightCosets}
}\hfill{\scriptsize (attribute)}}\\


 A double coset $\mbox{\texttt{\mdseries\slshape D}} = U g V$ can be considered as a union of right cosets $U h_i$. (It is the union of the orbit of $U g$ under right multiplication by $V$.) For a double coset \mbox{\texttt{\mdseries\slshape D}} this function returns a set of representatives $h_i$ such that \mbox{\texttt{\mdseries\slshape D}} $= \bigcup_{{h_i}} U h_i$. The representatives returned are canonical for $U$ (see \texttt{CanonicalRightCosetElement} (\ref{CanonicalRightCosetElement})) and form a set. 
\begin{Verbatim}[commandchars=!@|,fontsize=\small,frame=single,label=Example]
  !gapprompt@gap>| !gapinput@u:=Subgroup(g,[(1,2,3),(1,2)]);;v:=Subgroup(g,[(3,4)]);;|
  !gapprompt@gap>| !gapinput@c:=DoubleCoset(u,(2,4),v);|
  DoubleCoset(Group( [ (1,2,3), (1,2) ] ),(2,4),Group( [ (3,4) ] ))
  !gapprompt@gap>| !gapinput@(1,2,3) in c;|
  false
  !gapprompt@gap>| !gapinput@(2,3,4) in c;|
  true
  !gapprompt@gap>| !gapinput@LeftActingGroup(c);|
  Group([ (1,2,3), (1,2) ])
  !gapprompt@gap>| !gapinput@RightActingGroup(c);|
  Group([ (3,4) ])
  !gapprompt@gap>| !gapinput@RepresentativesContainedRightCosets(c);|
  [ (2,3,4) ]
\end{Verbatim}
 }

 

\subsection{\textcolor{Chapter }{DoubleCosets}}
\logpage{[ 39, 9, 3 ]}\nobreak
\hyperdef{L}{X7A5EFABB86E6D4D5}{}
{\noindent\textcolor{FuncColor}{$\triangleright$\enspace\texttt{DoubleCosets({\mdseries\slshape G, U, V})\index{DoubleCosets@\texttt{DoubleCosets}}
\label{DoubleCosets}
}\hfill{\scriptsize (function)}}\\
\noindent\textcolor{FuncColor}{$\triangleright$\enspace\texttt{DoubleCosetsNC({\mdseries\slshape G, U, V})\index{DoubleCosetsNC@\texttt{DoubleCosetsNC}}
\label{DoubleCosetsNC}
}\hfill{\scriptsize (operation)}}\\


 computes a duplicate free list of all double cosets \mbox{\texttt{\mdseries\slshape U}} $g$ \mbox{\texttt{\mdseries\slshape V}} for $g \in \mbox{\texttt{\mdseries\slshape G}}$. The groups \mbox{\texttt{\mdseries\slshape U}} and \mbox{\texttt{\mdseries\slshape V}} must be subgroups of the group \mbox{\texttt{\mdseries\slshape G}}. The \texttt{NC} version does not check whether \mbox{\texttt{\mdseries\slshape U}} and \mbox{\texttt{\mdseries\slshape V}} are subgroups of \mbox{\texttt{\mdseries\slshape G}}. 
\begin{Verbatim}[commandchars=!@|,fontsize=\small,frame=single,label=Example]
  !gapprompt@gap>| !gapinput@dc:=DoubleCosets(g,u,v);|
  [ DoubleCoset(Group( [ (1,2,3), (1,2) ] ),(),Group( [ (3,4) ] )),
    DoubleCoset(Group( [ (1,2,3), (1,2) ] ),(1,3)(2,4),Group(
      [ (3,4) ] )), DoubleCoset(Group( [ (1,2,3), (1,2) ] ),(1,4)
      (2,3),Group( [ (3,4) ] )) ]
  !gapprompt@gap>| !gapinput@List(dc,Representative);|
  [ (), (1,3)(2,4), (1,4)(2,3) ]
\end{Verbatim}
 }

 

\subsection{\textcolor{Chapter }{IsDoubleCoset (operation)}}
\logpage{[ 39, 9, 4 ]}\nobreak
\hyperdef{L}{X85ED464F878EF24C}{}
{\noindent\textcolor{FuncColor}{$\triangleright$\enspace\texttt{IsDoubleCoset({\mdseries\slshape obj})\index{IsDoubleCoset@\texttt{IsDoubleCoset}!operation}
\label{IsDoubleCoset:operation}
}\hfill{\scriptsize (Category)}}\\


 The category of double cosets. }

 

\subsection{\textcolor{Chapter }{DoubleCosetRepsAndSizes}}
\logpage{[ 39, 9, 5 ]}\nobreak
\hyperdef{L}{X7A25B1C886CF8C6A}{}
{\noindent\textcolor{FuncColor}{$\triangleright$\enspace\texttt{DoubleCosetRepsAndSizes({\mdseries\slshape G, U, V})\index{DoubleCosetRepsAndSizes@\texttt{DoubleCosetRepsAndSizes}}
\label{DoubleCosetRepsAndSizes}
}\hfill{\scriptsize (operation)}}\\


 returns a list of double coset representatives and their sizes, the entries
are lists of the form $[ r, n ]$ where $r$ and $n$ are an element of the double coset and the size of the coset, respectively.
This operation is faster than \texttt{DoubleCosetsNC} (\ref{DoubleCosetsNC}) because no double coset objects have to be created. 
\begin{Verbatim}[commandchars=!@|,fontsize=\small,frame=single,label=Example]
  !gapprompt@gap>| !gapinput@dc:=DoubleCosetRepsAndSizes(g,u,v);|
  [ [ (), 12 ], [ (1,3)(2,4), 6 ], [ (1,4)(2,3), 6 ] ]
\end{Verbatim}
 }

 

\subsection{\textcolor{Chapter }{InfoCoset}}
\logpage{[ 39, 9, 6 ]}\nobreak
\hyperdef{L}{X84AE7EE77E5FB30E}{}
{\noindent\textcolor{FuncColor}{$\triangleright$\enspace\texttt{InfoCoset\index{InfoCoset@\texttt{InfoCoset}}
\label{InfoCoset}
}\hfill{\scriptsize (info class)}}\\


 The information function for coset and double coset operations is \texttt{InfoCoset}. }

 }

  
\section{\textcolor{Chapter }{Conjugacy Classes}}\label{Conjugacy Classes}
\logpage{[ 39, 10, 0 ]}
\hyperdef{L}{X7D474F8F87E4E5D9}{}
{
  

\subsection{\textcolor{Chapter }{ConjugacyClass}}
\logpage{[ 39, 10, 1 ]}\nobreak
\hyperdef{L}{X7B2F207F7F85F5B8}{}
{\noindent\textcolor{FuncColor}{$\triangleright$\enspace\texttt{ConjugacyClass({\mdseries\slshape G, g})\index{ConjugacyClass@\texttt{ConjugacyClass}}
\label{ConjugacyClass}
}\hfill{\scriptsize (operation)}}\\


 creates the conjugacy class in \mbox{\texttt{\mdseries\slshape G}} with representative \mbox{\texttt{\mdseries\slshape g}}. This class is an external set, so functions such as \texttt{Representative} (\ref{Representative}) (which returns \mbox{\texttt{\mdseries\slshape g}}), \texttt{ActingDomain} (\ref{ActingDomain}) (which returns \mbox{\texttt{\mdseries\slshape G}}), \texttt{StabilizerOfExternalSet} (\ref{StabilizerOfExternalSet}) (which returns the centralizer of \mbox{\texttt{\mdseries\slshape g}}) and \texttt{AsList} (\ref{AsList}) work for it. 

 A conjugacy class is an external orbit (see \texttt{ExternalOrbit} (\ref{ExternalOrbit})) of group elements with the group acting by conjugation on it. Thus element
tests or operation representatives can be computed. The attribute \texttt{Centralizer} (\ref{Centralizer:for a class of objects in a magma}) gives the centralizer of the representative (which is the same result as \texttt{StabilizerOfExternalSet} (\ref{StabilizerOfExternalSet})). (This is a slight abuse of notation: This is \emph{not} the centralizer of the class as a \emph{set} which would be the standard behaviour of \texttt{Centralizer} (\ref{Centralizer:for a class of objects in a magma}).) }

 

\subsection{\textcolor{Chapter }{ConjugacyClasses (attribute)}}
\logpage{[ 39, 10, 2 ]}\nobreak
\hyperdef{L}{X871B570284BBA685}{}
{\noindent\textcolor{FuncColor}{$\triangleright$\enspace\texttt{ConjugacyClasses({\mdseries\slshape G})\index{ConjugacyClasses@\texttt{ConjugacyClasses}!attribute}
\label{ConjugacyClasses:attribute}
}\hfill{\scriptsize (attribute)}}\\


 returns the conjugacy classes of elements of \mbox{\texttt{\mdseries\slshape G}} as a list of class objects of \mbox{\texttt{\mdseries\slshape G}} (see{\nobreakspace}\texttt{ConjugacyClass} (\ref{ConjugacyClass}) for details). It is guaranteed that the class of the identity is in the first
position, the further arrangement depends on the method chosen (and might be
different for equal but not identical groups). 

 For very small groups (of size up to 500) the classes will be computed by the
conjugation action of \mbox{\texttt{\mdseries\slshape G}} on itself (see{\nobreakspace}\texttt{ConjugacyClassesByOrbits} (\ref{ConjugacyClassesByOrbits})). This can be deliberately switched off using the ``\texttt{noaction}'' option shown below. 

 For solvable groups, the default method to compute the classes is by
homomorphic lift (see section{\nobreakspace}\ref{Conjugacy Classes in Solvable Groups}). 

 For other groups the method of \cite{HulpkeClasses} is employed. 

 \texttt{ConjugacyClasses} supports the following options that can be used to modify this strategy: 
\begin{description}
\item[{\texttt{random}}]  The classes are computed by random search. See \texttt{ConjugacyClassesByRandomSearch} (\ref{ConjugacyClassesByRandomSearch}) below. 
\item[{\texttt{action}}]  The classes are computed by action of \mbox{\texttt{\mdseries\slshape G}} on itself. See \texttt{ConjugacyClassesByOrbits} (\ref{ConjugacyClassesByOrbits}) below. 
\item[{\texttt{noaction}}]  Even for small groups \texttt{ConjugacyClassesByOrbits} (\ref{ConjugacyClassesByOrbits}) is not used as a default. This can be useful if the elements of the group use
a lot of memory. 
\end{description}
 
\begin{Verbatim}[commandchars=!@|,fontsize=\small,frame=single,label=Example]
  !gapprompt@gap>| !gapinput@g:=SymmetricGroup(4);;|
  !gapprompt@gap>| !gapinput@cl:=ConjugacyClasses(g);|
  [ ()^G, (1,2)^G, (1,2)(3,4)^G, (1,2,3)^G, (1,2,3,4)^G ]
  !gapprompt@gap>| !gapinput@Representative(cl[3]);Centralizer(cl[3]);|
  (1,2)(3,4)
  Group([ (1,2), (1,3)(2,4), (3,4) ])
  !gapprompt@gap>| !gapinput@Size(Centralizer(cl[5]));|
  4
  !gapprompt@gap>| !gapinput@Size(cl[2]);|
  6
\end{Verbatim}
 

 In general, you will not need to have to influence the method, but simply call \texttt{ConjugacyClasses} {\textendash}\textsf{GAP} will try to select a suitable method on its own. The method specifications are
provided here mainly for expert use. }

 

\subsection{\textcolor{Chapter }{ConjugacyClassesByRandomSearch}}
\logpage{[ 39, 10, 3 ]}\nobreak
\hyperdef{L}{X7D6ED84C86C2979B}{}
{\noindent\textcolor{FuncColor}{$\triangleright$\enspace\texttt{ConjugacyClassesByRandomSearch({\mdseries\slshape G})\index{ConjugacyClassesByRandomSearch@\texttt{ConjugacyClassesByRandomSearch}}
\label{ConjugacyClassesByRandomSearch}
}\hfill{\scriptsize (function)}}\\


 computes the classes of the group \mbox{\texttt{\mdseries\slshape G}} by random search. This works very efficiently for almost simple groups. 

 This function is also accessible via the option \texttt{random} to the function \texttt{ConjugacyClass} (\ref{ConjugacyClass}). }

 

\subsection{\textcolor{Chapter }{ConjugacyClassesByOrbits}}
\logpage{[ 39, 10, 4 ]}\nobreak
\hyperdef{L}{X852B3634789D770E}{}
{\noindent\textcolor{FuncColor}{$\triangleright$\enspace\texttt{ConjugacyClassesByOrbits({\mdseries\slshape G})\index{ConjugacyClassesByOrbits@\texttt{ConjugacyClassesByOrbits}}
\label{ConjugacyClassesByOrbits}
}\hfill{\scriptsize (function)}}\\


 computes the classes of the group \mbox{\texttt{\mdseries\slshape G}} as orbits of \mbox{\texttt{\mdseries\slshape G}} on its elements. This can be quick but unsurprisingly may also take a lot of
memory if \mbox{\texttt{\mdseries\slshape G}} becomes larger. All the classes will store their element list and thus a
membership test will be quick as well. 

 This function is also accessible via the option \texttt{action} to the function \texttt{ConjugacyClass} (\ref{ConjugacyClass}). 

 Typically, for small groups (roughly of order up to $10^3$) the computation of classes as orbits under the action is fastest; memory
restrictions (and the increasing cost of eliminating duplicates) make this
less efficient for larger groups. 

 Calculation by random search has the smallest memory requirement, but in
generally performs worse, the more classes are there. 

 The following example shows the effect of this for a small group with many
classes: 

  
\begin{Verbatim}[commandchars=!@|,fontsize=\small,frame=single,label=Example]
  !gapprompt@gap>| !gapinput@h:=Group((4,5)(6,7,8),(1,2,3)(5,6,9));;ConjugacyClasses(h:noaction);;time;|
  110
  !gapprompt@gap>| !gapinput@h:=Group((4,5)(6,7,8),(1,2,3)(5,6,9));;ConjugacyClasses(h:random);;time;|
  300
  !gapprompt@gap>| !gapinput@h:=Group((4,5)(6,7,8),(1,2,3)(5,6,9));;ConjugacyClasses(h:action);;time;|
  30
\end{Verbatim}
 }

 

\subsection{\textcolor{Chapter }{NrConjugacyClasses}}
\logpage{[ 39, 10, 5 ]}\nobreak
\hyperdef{L}{X8733F87B7E4C9903}{}
{\noindent\textcolor{FuncColor}{$\triangleright$\enspace\texttt{NrConjugacyClasses({\mdseries\slshape G})\index{NrConjugacyClasses@\texttt{NrConjugacyClasses}}
\label{NrConjugacyClasses}
}\hfill{\scriptsize (attribute)}}\\


 returns the number of conjugacy classes of \mbox{\texttt{\mdseries\slshape G}}. 
\begin{Verbatim}[commandchars=!@|,fontsize=\small,frame=single,label=Example]
  !gapprompt@gap>| !gapinput@g:=Group((1,2,3,4),(1,2));;|
  !gapprompt@gap>| !gapinput@NrConjugacyClasses(g);|
  5
\end{Verbatim}
 }

 

\subsection{\textcolor{Chapter }{RationalClass}}
\logpage{[ 39, 10, 6 ]}\nobreak
\hyperdef{L}{X7BD2A4427B7FE248}{}
{\noindent\textcolor{FuncColor}{$\triangleright$\enspace\texttt{RationalClass({\mdseries\slshape G, g})\index{RationalClass@\texttt{RationalClass}}
\label{RationalClass}
}\hfill{\scriptsize (operation)}}\\


 creates the rational class in \mbox{\texttt{\mdseries\slshape G}} with representative \mbox{\texttt{\mdseries\slshape g}}. A rational class consists of all elements that are conjugate to \mbox{\texttt{\mdseries\slshape g}} or to an $i$\texttt{\symbol{45}}th power of \mbox{\texttt{\mdseries\slshape g}} where $i$ is coprime to the order of $g$. Thus a rational class can be interpreted as a conjugacy class of cyclic
subgroups. A rational class is an external set (\texttt{IsExternalSet} (\ref{IsExternalSet})) of group elements with the group acting by conjugation on it, but not an
external orbit. }

 

\subsection{\textcolor{Chapter }{RationalClasses}}
\logpage{[ 39, 10, 7 ]}\nobreak
\hyperdef{L}{X81E9EF0A811072E8}{}
{\noindent\textcolor{FuncColor}{$\triangleright$\enspace\texttt{RationalClasses({\mdseries\slshape G})\index{RationalClasses@\texttt{RationalClasses}}
\label{RationalClasses}
}\hfill{\scriptsize (attribute)}}\\


 returns a list of the rational classes of the group \mbox{\texttt{\mdseries\slshape G}}. (See \texttt{RationalClass} (\ref{RationalClass}).) 
\begin{Verbatim}[commandchars=!@|,fontsize=\small,frame=single,label=Example]
  !gapprompt@gap>| !gapinput@RationalClasses(DerivedSubgroup(g));|
  [ RationalClass( AlternatingGroup( [ 1 .. 4 ] ), () ),
    RationalClass( AlternatingGroup( [ 1 .. 4 ] ), (1,2)(3,4) ),
    RationalClass( AlternatingGroup( [ 1 .. 4 ] ), (1,2,3) ) ]
\end{Verbatim}
 }

 

\subsection{\textcolor{Chapter }{GaloisGroup (of rational class of a group)}}
\logpage{[ 39, 10, 8 ]}\nobreak
\hyperdef{L}{X877691247DE23386}{}
{\noindent\textcolor{FuncColor}{$\triangleright$\enspace\texttt{GaloisGroup({\mdseries\slshape ratcl})\index{GaloisGroup@\texttt{GaloisGroup}!of rational class of a group}
\label{GaloisGroup:of rational class of a group}
}\hfill{\scriptsize (attribute)}}\\


 Suppose that \mbox{\texttt{\mdseries\slshape ratcl}} is a rational class of a group $G$ with representative $g$. The exponents $i$ for which $g^i$ lies already in the ordinary conjugacy class of $g$, form a subgroup of the \emph{prime residue class group} $P_n$ (see \texttt{PrimitiveRootMod} (\ref{PrimitiveRootMod})), the so\texttt{\symbol{45}}called \emph{Galois group} of the rational class. The prime residue class group $P_n$ is obtained in \textsf{GAP} as \texttt{Units( Integers mod \mbox{\texttt{\mdseries\slshape n}} )}, the unit group of a residue class ring. The Galois group of a rational class \mbox{\texttt{\mdseries\slshape ratcl}} is stored in the attribute \texttt{GaloisGroup} as a subgroup of this group. }

 
\subsection{\textcolor{Chapter }{IsConjugate}}\logpage{[ 39, 10, 9 ]}
\hyperdef{L}{X83DD148D7DA2ABA9}{}
{
\noindent\textcolor{FuncColor}{$\triangleright$\enspace\texttt{IsConjugate({\mdseries\slshape G, x, y})\index{IsConjugate@\texttt{IsConjugate}!for a group and two elements}
\label{IsConjugate:for a group and two elements}
}\hfill{\scriptsize (operation)}}\\
\noindent\textcolor{FuncColor}{$\triangleright$\enspace\texttt{IsConjugate({\mdseries\slshape G, U, V})\index{IsConjugate@\texttt{IsConjugate}!for a group and two groups}
\label{IsConjugate:for a group and two groups}
}\hfill{\scriptsize (operation)}}\\


 tests whether the elements \mbox{\texttt{\mdseries\slshape x}} and \mbox{\texttt{\mdseries\slshape y}} or the subgroups \mbox{\texttt{\mdseries\slshape U}} and \mbox{\texttt{\mdseries\slshape V}} are conjugate under the action of \mbox{\texttt{\mdseries\slshape G}}. (They do not need to be \emph{contained in} \mbox{\texttt{\mdseries\slshape G}}.) This command is only a shortcut to \texttt{RepresentativeAction} (\ref{RepresentativeAction}). 
\begin{Verbatim}[commandchars=!@|,fontsize=\small,frame=single,label=Example]
  !gapprompt@gap>| !gapinput@IsConjugate(g,Group((1,2,3,4),(1,3)),Group((1,3,2,4),(1,2)));|
  true
\end{Verbatim}
 

 \texttt{RepresentativeAction} (\ref{RepresentativeAction}) can be used to obtain conjugating elements. 
\begin{Verbatim}[commandchars=!@|,fontsize=\small,frame=single,label=Example]
  !gapprompt@gap>| !gapinput@RepresentativeAction(g,(1,2),(3,4));|
  (1,3)(2,4)
\end{Verbatim}
 }

 

\subsection{\textcolor{Chapter }{NthRootsInGroup}}
\logpage{[ 39, 10, 10 ]}\nobreak
\hyperdef{L}{X81A92F828400FC8A}{}
{\noindent\textcolor{FuncColor}{$\triangleright$\enspace\texttt{NthRootsInGroup({\mdseries\slshape G, e, n})\index{NthRootsInGroup@\texttt{NthRootsInGroup}}
\label{NthRootsInGroup}
}\hfill{\scriptsize (function)}}\\


 Let \mbox{\texttt{\mdseries\slshape e}} be an element in the group \mbox{\texttt{\mdseries\slshape G}}. This function returns a list of all those elements in \mbox{\texttt{\mdseries\slshape G}} whose \mbox{\texttt{\mdseries\slshape n}}\texttt{\symbol{45}}th power is \mbox{\texttt{\mdseries\slshape e}}. 

 
\begin{Verbatim}[commandchars=!@|,fontsize=\small,frame=single,label=Example]
  !gapprompt@gap>| !gapinput@NthRootsInGroup(g,(1,2)(3,4),2);|
  [ (1,3,2,4), (1,4,2,3) ]
\end{Verbatim}
 }

 }

  
\section{\textcolor{Chapter }{Normal Structure}}\label{Normal Structure}
\logpage{[ 39, 11, 0 ]}
\hyperdef{L}{X804F0F037F06E25E}{}
{
  For the operations \texttt{Centralizer} (\ref{Centralizer:for a magma and a submagma}) and \texttt{Centre} (\ref{Centre}), see Chapter{\nobreakspace}\ref{Magmas}. \index{normalizer} 
\subsection{\textcolor{Chapter }{Normalizer}}\logpage{[ 39, 11, 1 ]}
\hyperdef{L}{X87B5370C7DFD401D}{}
{
\noindent\textcolor{FuncColor}{$\triangleright$\enspace\texttt{Normalizer({\mdseries\slshape G, U})\index{Normalizer@\texttt{Normalizer}!for two groups}
\label{Normalizer:for two groups}
}\hfill{\scriptsize (operation)}}\\
\noindent\textcolor{FuncColor}{$\triangleright$\enspace\texttt{Normalizer({\mdseries\slshape G, g})\index{Normalizer@\texttt{Normalizer}!for a group and a group element}
\label{Normalizer:for a group and a group element}
}\hfill{\scriptsize (operation)}}\\


 For two groups \mbox{\texttt{\mdseries\slshape G}}, \mbox{\texttt{\mdseries\slshape U}}, \texttt{Normalizer} computes the normalizer $N_{\mbox{\texttt{\mdseries\slshape G}}}(\mbox{\texttt{\mdseries\slshape U}})$, that is, the stabilizer of \mbox{\texttt{\mdseries\slshape U}} under the conjugation action of \mbox{\texttt{\mdseries\slshape G}}. 

 For a group \mbox{\texttt{\mdseries\slshape G}} and a group element \mbox{\texttt{\mdseries\slshape g}}, \texttt{Normalizer} computes $N_{\mbox{\texttt{\mdseries\slshape G}}}(\langle \mbox{\texttt{\mdseries\slshape g}} \rangle)$. 

 
\begin{Verbatim}[commandchars=!@|,fontsize=\small,frame=single,label=Example]
  !gapprompt@gap>| !gapinput@Normalizer(g,Subgroup(g,[(1,2,3)]));|
  Group([ (1,2,3), (2,3) ])
\end{Verbatim}
 }

 

\subsection{\textcolor{Chapter }{Core}}
\logpage{[ 39, 11, 2 ]}\nobreak
\hyperdef{L}{X7C4E00297E37AA44}{}
{\noindent\textcolor{FuncColor}{$\triangleright$\enspace\texttt{Core({\mdseries\slshape S, U})\index{Core@\texttt{Core}}
\label{Core}
}\hfill{\scriptsize (operation)}}\\


 If \mbox{\texttt{\mdseries\slshape S}} and \mbox{\texttt{\mdseries\slshape U}} are groups of elements in the same family, this operation returns the core of \mbox{\texttt{\mdseries\slshape U}} in \mbox{\texttt{\mdseries\slshape S}}, that is the intersection of all \mbox{\texttt{\mdseries\slshape S}}\texttt{\symbol{45}}conjugates of \mbox{\texttt{\mdseries\slshape U}}. 
\begin{Verbatim}[commandchars=!@|,fontsize=\small,frame=single,label=Example]
  !gapprompt@gap>| !gapinput@g:=Group((1,2,3,4),(1,2));;|
  !gapprompt@gap>| !gapinput@Core(g,Subgroup(g,[(1,2,3,4)]));|
  Group(())
\end{Verbatim}
 }

 

\subsection{\textcolor{Chapter }{PCore}}
\logpage{[ 39, 11, 3 ]}\nobreak
\hyperdef{L}{X7CF497C77B1E8938}{}
{\noindent\textcolor{FuncColor}{$\triangleright$\enspace\texttt{PCore({\mdseries\slshape G, p})\index{PCore@\texttt{PCore}}
\label{PCore}
}\hfill{\scriptsize (operation)}}\\


 \index{Op(G)@$O_p(G)$!see PCore} The \emph{\mbox{\texttt{\mdseries\slshape p}}\texttt{\symbol{45}}core} of \mbox{\texttt{\mdseries\slshape G}} is the largest normal \mbox{\texttt{\mdseries\slshape p}}\texttt{\symbol{45}}subgroup of \mbox{\texttt{\mdseries\slshape G}}. It is the core of a Sylow \mbox{\texttt{\mdseries\slshape p}}\texttt{\symbol{45}}subgroup of \mbox{\texttt{\mdseries\slshape G}}, see \texttt{Core} (\ref{Core}). 
\begin{Verbatim}[commandchars=!@|,fontsize=\small,frame=single,label=Example]
  !gapprompt@gap>| !gapinput@g:=DicyclicGroup(12);;|
  !gapprompt@gap>| !gapinput@PCore(g,2);|
  Group([ y3 ])
  !gapprompt@gap>| !gapinput@PCore(g,2) = Core(g,SylowSubgroup(g,2));|
  true
  !gapprompt@gap>| !gapinput@PCore(g,3);|
  Group([ y*y3 ])
  !gapprompt@gap>| !gapinput@PCore(g,5);|
  Group([  ])
  !gapprompt@gap>| !gapinput@g:=SymmetricGroup(4);;|
  !gapprompt@gap>| !gapinput@PCore(g,2);|
  Group([ (1,4)(2,3), (1,2)(3,4) ])
  !gapprompt@gap>| !gapinput@PCore(g,2) = Core(g,SylowSubgroup(g,2));|
  true
\end{Verbatim}
 }

 

\subsection{\textcolor{Chapter }{NormalClosure}}
\logpage{[ 39, 11, 4 ]}\nobreak
\hyperdef{L}{X7BDEA0A98720D1BB}{}
{\noindent\textcolor{FuncColor}{$\triangleright$\enspace\texttt{NormalClosure({\mdseries\slshape G, U})\index{NormalClosure@\texttt{NormalClosure}}
\label{NormalClosure}
}\hfill{\scriptsize (operation)}}\\
\noindent\textcolor{FuncColor}{$\triangleright$\enspace\texttt{NormalClosure({\mdseries\slshape G, list})\index{NormalClosure@\texttt{NormalClosure}!for group and a list}
\label{NormalClosure:for group and a list}
}\hfill{\scriptsize (operation)}}\\


 The normal closure of \mbox{\texttt{\mdseries\slshape U}} in \mbox{\texttt{\mdseries\slshape G}} is the smallest normal subgroup of the closure of \mbox{\texttt{\mdseries\slshape G}} and \mbox{\texttt{\mdseries\slshape U}} which contains \mbox{\texttt{\mdseries\slshape U}}. 

 The second argument may also be a list of group elements, in which case the
normal closure of the group generated by these elements is computed. 
\begin{Verbatim}[commandchars=!@|,fontsize=\small,frame=single,label=Example]
  !gapprompt@gap>| !gapinput@NormalClosure(g,Subgroup(g,[(1,2,3)])) = Group([ (1,2,3), (2,3,4) ]);|
  true
  !gapprompt@gap>| !gapinput@NormalClosure(g,[(1,2,3)]) = Group([ (1,2,3), (2,3,4) ]);|
  true
  !gapprompt@gap>| !gapinput@NormalClosure(g,Group((3,4,5))) = Group([ (3,4,5), (1,5,4), (1,2,5) ]);|
  true
\end{Verbatim}
 }

 

\subsection{\textcolor{Chapter }{NormalIntersection}}
\logpage{[ 39, 11, 5 ]}\nobreak
\hyperdef{L}{X7D25E7DC7834A703}{}
{\noindent\textcolor{FuncColor}{$\triangleright$\enspace\texttt{NormalIntersection({\mdseries\slshape G, U})\index{NormalIntersection@\texttt{NormalIntersection}}
\label{NormalIntersection}
}\hfill{\scriptsize (operation)}}\\


 computes the intersection of \mbox{\texttt{\mdseries\slshape G}} and \mbox{\texttt{\mdseries\slshape U}}, assuming that \mbox{\texttt{\mdseries\slshape G}} is normalized by \mbox{\texttt{\mdseries\slshape U}}. This works faster than \texttt{Intersection}, but will not produce the intersection if \mbox{\texttt{\mdseries\slshape G}} is not normalized by \mbox{\texttt{\mdseries\slshape U}}. 
\begin{Verbatim}[commandchars=!@|,fontsize=\small,frame=single,label=Example]
  !gapprompt@gap>| !gapinput@NormalIntersection(Group((1,2)(3,4),(1,3)(2,4)),Group((1,2,3,4)));|
  Group([ (1,3)(2,4) ])
\end{Verbatim}
 }

 

\subsection{\textcolor{Chapter }{ComplementClassesRepresentatives}}
\logpage{[ 39, 11, 6 ]}\nobreak
\hyperdef{L}{X811B8A4683DDE1F9}{}
{\noindent\textcolor{FuncColor}{$\triangleright$\enspace\texttt{ComplementClassesRepresentatives({\mdseries\slshape G, N})\index{ComplementClassesRepresentatives@\texttt{ComplementClassesRepresentatives}}
\label{ComplementClassesRepresentatives}
}\hfill{\scriptsize (operation)}}\\


 Let \mbox{\texttt{\mdseries\slshape N}} be a normal subgroup of \mbox{\texttt{\mdseries\slshape G}}. This command returns a set of representatives for the conjugacy classes of
complements of \mbox{\texttt{\mdseries\slshape N}} in \mbox{\texttt{\mdseries\slshape G}}. Complements are subgroups of \mbox{\texttt{\mdseries\slshape G}} which intersect trivially with \mbox{\texttt{\mdseries\slshape N}} and together with \mbox{\texttt{\mdseries\slshape N}} generate \mbox{\texttt{\mdseries\slshape G}}. 

 At the moment methods are available only for the case that \mbox{\texttt{\mdseries\slshape N}} or \mbox{\texttt{\mdseries\slshape G}}\texttt{/}\mbox{\texttt{\mdseries\slshape N}} is solvable. 
\begin{Verbatim}[commandchars=!@|,fontsize=\small,frame=single,label=Example]
  !gapprompt@gap>| !gapinput@ComplementClassesRepresentatives(g,Group((1,2)(3,4),(1,3)(2,4)));|
  [ Group([ (3,4), (2,4,3) ]) ]
\end{Verbatim}
 }

 

\subsection{\textcolor{Chapter }{InfoComplement}}
\logpage{[ 39, 11, 7 ]}\nobreak
\hyperdef{L}{X8581F4E77B11C610}{}
{\noindent\textcolor{FuncColor}{$\triangleright$\enspace\texttt{InfoComplement\index{InfoComplement@\texttt{InfoComplement}}
\label{InfoComplement}
}\hfill{\scriptsize (info class)}}\\


 Info class for the complement routines. }

 }

  
\section{\textcolor{Chapter }{Specific and Parametrized Subgroups}}\label{Specific and Parametrized Subgroups}
\logpage{[ 39, 12, 0 ]}
\hyperdef{L}{X7C39EE3E836D6BC6}{}
{
  The centre of a group (the subgroup of those elements that commute with all
other elements of the group) can be computed by the operation \texttt{Centre} (\ref{Centre}). 

\subsection{\textcolor{Chapter }{TrivialSubgroup}}
\logpage{[ 39, 12, 1 ]}\nobreak
\hyperdef{L}{X829759F67D4247CA}{}
{\noindent\textcolor{FuncColor}{$\triangleright$\enspace\texttt{TrivialSubgroup({\mdseries\slshape G})\index{TrivialSubgroup@\texttt{TrivialSubgroup}}
\label{TrivialSubgroup}
}\hfill{\scriptsize (attribute)}}\\


 
\begin{Verbatim}[commandchars=!@|,fontsize=\small,frame=single,label=Example]
  !gapprompt@gap>| !gapinput@TrivialSubgroup(g);|
  Group(())
\end{Verbatim}
 }

 

\subsection{\textcolor{Chapter }{CommutatorSubgroup}}
\logpage{[ 39, 12, 2 ]}\nobreak
\hyperdef{L}{X7A9A3D5578CE33A0}{}
{\noindent\textcolor{FuncColor}{$\triangleright$\enspace\texttt{CommutatorSubgroup({\mdseries\slshape G, H})\index{CommutatorSubgroup@\texttt{CommutatorSubgroup}}
\label{CommutatorSubgroup}
}\hfill{\scriptsize (operation)}}\\


 If \mbox{\texttt{\mdseries\slshape G}} and \mbox{\texttt{\mdseries\slshape H}} are two groups of elements in the same family, this operation returns the
group generated by all commutators $[ g, h ] = g^{{-1}} h^{{-1}} g h$ (see{\nobreakspace}\texttt{Comm} (\ref{Comm})) of elements $g \in \mbox{\texttt{\mdseries\slshape G}}$ and $h \in \mbox{\texttt{\mdseries\slshape H}}$, that is the group $\left \langle [ g, h ] \mid g \in \mbox{\texttt{\mdseries\slshape G}}, h \in \mbox{\texttt{\mdseries\slshape H}} \right \rangle$. 
\begin{Verbatim}[commandchars=!@|,fontsize=\small,frame=single,label=Example]
  !gapprompt@gap>| !gapinput@CommutatorSubgroup(Group((1,2,3),(1,2)),Group((2,3,4),(3,4)));|
  Group([ (1,4)(2,3), (1,3,4) ])
  !gapprompt@gap>| !gapinput@Size(last);|
  12
\end{Verbatim}
 }

 

\subsection{\textcolor{Chapter }{DerivedSubgroup}}
\logpage{[ 39, 12, 3 ]}\nobreak
\hyperdef{L}{X7CC17CF179ED7EF2}{}
{\noindent\textcolor{FuncColor}{$\triangleright$\enspace\texttt{DerivedSubgroup({\mdseries\slshape G})\index{DerivedSubgroup@\texttt{DerivedSubgroup}}
\label{DerivedSubgroup}
}\hfill{\scriptsize (attribute)}}\\


 The derived subgroup $\mbox{\texttt{\mdseries\slshape G}}'$ of \mbox{\texttt{\mdseries\slshape G}} is the subgroup generated by all commutators of pairs of elements of \mbox{\texttt{\mdseries\slshape G}}. It is normal in \mbox{\texttt{\mdseries\slshape G}} and the factor group $\mbox{\texttt{\mdseries\slshape G}}/\mbox{\texttt{\mdseries\slshape G}}'$ is the largest abelian factor group of \mbox{\texttt{\mdseries\slshape G}}. 
\begin{Verbatim}[commandchars=!@|,fontsize=\small,frame=single,label=Example]
  !gapprompt@gap>| !gapinput@g:=Group((1,2,3,4),(1,2));;|
  !gapprompt@gap>| !gapinput@DerivedSubgroup(g) = Group([ (1,3,2), (2,4,3) ]);|
  true
\end{Verbatim}
 }

 

\subsection{\textcolor{Chapter }{CommutatorLength}}
\logpage{[ 39, 12, 4 ]}\nobreak
\hyperdef{L}{X7B10B58F83DDE56E}{}
{\noindent\textcolor{FuncColor}{$\triangleright$\enspace\texttt{CommutatorLength({\mdseries\slshape G})\index{CommutatorLength@\texttt{CommutatorLength}}
\label{CommutatorLength}
}\hfill{\scriptsize (attribute)}}\\


 returns the minimal number $n$ such that each element in the derived subgroup (see{\nobreakspace}\texttt{DerivedSubgroup} (\ref{DerivedSubgroup})) of the group \mbox{\texttt{\mdseries\slshape G}} can be written as a product of (at most) $n$ commutators of elements in \mbox{\texttt{\mdseries\slshape G}}. 
\begin{Verbatim}[commandchars=!@|,fontsize=\small,frame=single,label=Example]
  !gapprompt@gap>| !gapinput@CommutatorLength( g );|
  1
\end{Verbatim}
 }

 

\subsection{\textcolor{Chapter }{FittingSubgroup}}
\logpage{[ 39, 12, 5 ]}\nobreak
\hyperdef{L}{X780552B57C30DD8F}{}
{\noindent\textcolor{FuncColor}{$\triangleright$\enspace\texttt{FittingSubgroup({\mdseries\slshape G})\index{FittingSubgroup@\texttt{FittingSubgroup}}
\label{FittingSubgroup}
}\hfill{\scriptsize (attribute)}}\\


 The Fitting subgroup of a group \mbox{\texttt{\mdseries\slshape G}} is its largest nilpotent normal subgroup. 
\begin{Verbatim}[commandchars=!@|,fontsize=\small,frame=single,label=Example]
  !gapprompt@gap>| !gapinput@FittingSubgroup(g);|
  Group([ (1,2)(3,4), (1,4)(2,3) ])
\end{Verbatim}
 }

 

\subsection{\textcolor{Chapter }{FrattiniSubgroup}}
\logpage{[ 39, 12, 6 ]}\nobreak
\hyperdef{L}{X788C856C82243274}{}
{\noindent\textcolor{FuncColor}{$\triangleright$\enspace\texttt{FrattiniSubgroup({\mdseries\slshape G})\index{FrattiniSubgroup@\texttt{FrattiniSubgroup}}
\label{FrattiniSubgroup}
}\hfill{\scriptsize (attribute)}}\\


 The Frattini subgroup of a group \mbox{\texttt{\mdseries\slshape G}} is the intersection of all maximal subgroups of \mbox{\texttt{\mdseries\slshape G}}. 
\begin{Verbatim}[commandchars=!@|,fontsize=\small,frame=single,label=Example]
  !gapprompt@gap>| !gapinput@FrattiniSubgroup(g);|
  Group(())
\end{Verbatim}
 }

 

\subsection{\textcolor{Chapter }{PrefrattiniSubgroup}}
\logpage{[ 39, 12, 7 ]}\nobreak
\hyperdef{L}{X81D86CCE84193E4F}{}
{\noindent\textcolor{FuncColor}{$\triangleright$\enspace\texttt{PrefrattiniSubgroup({\mdseries\slshape G})\index{PrefrattiniSubgroup@\texttt{PrefrattiniSubgroup}}
\label{PrefrattiniSubgroup}
}\hfill{\scriptsize (attribute)}}\\


 returns a Prefrattini subgroup of the finite solvable group \mbox{\texttt{\mdseries\slshape G}}. 

 A factor $M/N$ of \mbox{\texttt{\mdseries\slshape G}} is called a Frattini factor if $M/N$ is contained in the Frattini subgroup of $\mbox{\texttt{\mdseries\slshape G}}/N$. A subgroup $P$ is a Prefrattini subgroup of \mbox{\texttt{\mdseries\slshape G}} if $P$ covers each Frattini chief factor of \mbox{\texttt{\mdseries\slshape G}}, and if for each maximal subgroup of \mbox{\texttt{\mdseries\slshape G}} there exists a conjugate maximal subgroup, which contains $P$. In a finite solvable group \mbox{\texttt{\mdseries\slshape G}} the Prefrattini subgroups form a characteristic conjugacy class of subgroups
and the intersection of all these subgroups is the Frattini subgroup of \mbox{\texttt{\mdseries\slshape G}}. 
\begin{Verbatim}[commandchars=!@|,fontsize=\small,frame=single,label=Example]
  !gapprompt@gap>| !gapinput@G := SmallGroup( 60, 7 );|
  <pc group of size 60 with 4 generators>
  !gapprompt@gap>| !gapinput@P := PrefrattiniSubgroup(G);|
  Group([ f2 ])
  !gapprompt@gap>| !gapinput@Size(P);|
  2
  !gapprompt@gap>| !gapinput@IsNilpotent(P);|
  true
  !gapprompt@gap>| !gapinput@Core(G,P);|
  Group([  ])
  !gapprompt@gap>| !gapinput@FrattiniSubgroup(G);|
  Group([  ])
\end{Verbatim}
 }

 

\subsection{\textcolor{Chapter }{PerfectResiduum}}
\logpage{[ 39, 12, 8 ]}\nobreak
\hyperdef{L}{X83D5C8B8865C85F1}{}
{\noindent\textcolor{FuncColor}{$\triangleright$\enspace\texttt{PerfectResiduum({\mdseries\slshape G})\index{PerfectResiduum@\texttt{PerfectResiduum}}
\label{PerfectResiduum}
}\hfill{\scriptsize (attribute)}}\\


 is the smallest normal subgroup of \mbox{\texttt{\mdseries\slshape G}} that has a solvable factor group. 
\begin{Verbatim}[commandchars=!@|,fontsize=\small,frame=single,label=Example]
  !gapprompt@gap>| !gapinput@PerfectResiduum(SymmetricGroup(5));|
  Alt( [ 1 .. 5 ] )
\end{Verbatim}
 }

 

\subsection{\textcolor{Chapter }{SolvableRadical}}
\logpage{[ 39, 12, 9 ]}\nobreak
\hyperdef{L}{X8250D99A830DA832}{}
{\noindent\textcolor{FuncColor}{$\triangleright$\enspace\texttt{SolvableRadical({\mdseries\slshape G})\index{SolvableRadical@\texttt{SolvableRadical}}
\label{SolvableRadical}
}\hfill{\scriptsize (attribute)}}\\


 is the solvable radical of the group \mbox{\texttt{\mdseries\slshape G}}, i.e., the largest solvable normal subgroup of \mbox{\texttt{\mdseries\slshape G}}. 
\begin{Verbatim}[commandchars=!@|,fontsize=\small,frame=single,label=Example]
  !gapprompt@gap>| !gapinput@rad:= SolvableRadical( SL(2,5) );|
  <group of 2x2 matrices of size 2 over GF(5)>
  !gapprompt@gap>| !gapinput@Size( rad );|
  2
\end{Verbatim}
 }

 

\subsection{\textcolor{Chapter }{Socle}}
\logpage{[ 39, 12, 10 ]}\nobreak
\hyperdef{L}{X81F647FA83D8854F}{}
{\noindent\textcolor{FuncColor}{$\triangleright$\enspace\texttt{Socle({\mdseries\slshape G})\index{Socle@\texttt{Socle}}
\label{Socle}
}\hfill{\scriptsize (attribute)}}\\


 The socle of the group \mbox{\texttt{\mdseries\slshape G}} is the subgroup generated by all minimal normal subgroups. 
\begin{Verbatim}[commandchars=!@|,fontsize=\small,frame=single,label=Example]
  !gapprompt@gap>| !gapinput@Socle(g);|
  Group([ (1,4)(2,3), (1,2)(3,4) ])
\end{Verbatim}
 }

 

\subsection{\textcolor{Chapter }{SupersolvableResiduum}}
\logpage{[ 39, 12, 11 ]}\nobreak
\hyperdef{L}{X8440C61080CDAA14}{}
{\noindent\textcolor{FuncColor}{$\triangleright$\enspace\texttt{SupersolvableResiduum({\mdseries\slshape G})\index{SupersolvableResiduum@\texttt{SupersolvableResiduum}}
\label{SupersolvableResiduum}
}\hfill{\scriptsize (attribute)}}\\


 is the supersolvable residuum of the group \mbox{\texttt{\mdseries\slshape G}}, that is, its smallest normal subgroup $N$ such that the factor group $\mbox{\texttt{\mdseries\slshape G}} / N$ is supersolvable. 
\begin{Verbatim}[commandchars=!@|,fontsize=\small,frame=single,label=Example]
  !gapprompt@gap>| !gapinput@SupersolvableResiduum(g) = Group([ (1,3)(2,4), (1,4)(2,3) ]);|
  true
\end{Verbatim}
 }

 

\subsection{\textcolor{Chapter }{PRump}}
\logpage{[ 39, 12, 12 ]}\nobreak
\hyperdef{L}{X796DA805853FAC90}{}
{\noindent\textcolor{FuncColor}{$\triangleright$\enspace\texttt{PRump({\mdseries\slshape G, p})\index{PRump@\texttt{PRump}}
\label{PRump}
}\hfill{\scriptsize (operation)}}\\


 For a prime $p$, the \emph{\mbox{\texttt{\mdseries\slshape p}}\texttt{\symbol{45}}rump} of a group \mbox{\texttt{\mdseries\slshape G}} is the subgroup $\mbox{\texttt{\mdseries\slshape G}}' \mbox{\texttt{\mdseries\slshape G}}^{\mbox{\texttt{\mdseries\slshape p}}}$. Unless it equals \mbox{\texttt{\mdseries\slshape G}} itself (which is the e.g. the case if \mbox{\texttt{\mdseries\slshape G}} is perfect), it is equal to the second term of the \mbox{\texttt{\mdseries\slshape p}}\texttt{\symbol{45}}central series of \mbox{\texttt{\mdseries\slshape G}}, see \texttt{PCentralSeries} (\ref{PCentralSeries}). 

 
\begin{Verbatim}[commandchars=!@|,fontsize=\small,frame=single,label=Example]
  !gapprompt@gap>| !gapinput@g:=DicyclicGroup(12);;|
  !gapprompt@gap>| !gapinput@PRump(g,2) = PCentralSeries(g,2)[2];|
  true
  !gapprompt@gap>| !gapinput@g:=SymmetricGroup(4);;|
  !gapprompt@gap>| !gapinput@PRump(g,2) = AlternatingGroup(4);|
  true
\end{Verbatim}
 }

 }

  
\section{\textcolor{Chapter }{Sylow Subgroups and Hall Subgroups}}\label{Sylow Subgroups and Hall Subgroups}
\logpage{[ 39, 13, 0 ]}
\hyperdef{L}{X7FF0BBDD80E8F6BF}{}
{
  With respect to the following \textsf{GAP} functions, please note that by theorems of P.{\nobreakspace}Hall, a group $G$ is solvable if and only if one of the following conditions holds. 
\begin{enumerate}
\item  For each prime $p$ dividing the order of $G$, there exists a $p$\texttt{\symbol{45}}complement (see{\nobreakspace}\texttt{SylowComplement} (\ref{SylowComplement})). 
\item  For each set $P$ of primes dividing the order of $G$, there exists a $P$\texttt{\symbol{45}}Hall subgroup (see{\nobreakspace}\texttt{HallSubgroup} (\ref{HallSubgroup})). 
\item  $G$ has a Sylow system (see{\nobreakspace}\texttt{SylowSystem} (\ref{SylowSystem})). 
\item  $G$ has a complement system (see{\nobreakspace}\texttt{ComplementSystem} (\ref{ComplementSystem})). 
\end{enumerate}
 

\subsection{\textcolor{Chapter }{SylowSubgroup}}
\logpage{[ 39, 13, 1 ]}\nobreak
\hyperdef{L}{X7AA351308787544C}{}
{\noindent\textcolor{FuncColor}{$\triangleright$\enspace\texttt{SylowSubgroup({\mdseries\slshape G, p})\index{SylowSubgroup@\texttt{SylowSubgroup}}
\label{SylowSubgroup}
}\hfill{\scriptsize (operation)}}\\


 returns a Sylow \mbox{\texttt{\mdseries\slshape p}}\texttt{\symbol{45}}subgroup of the finite group \mbox{\texttt{\mdseries\slshape G}}. This is a \mbox{\texttt{\mdseries\slshape p}}\texttt{\symbol{45}}subgroup of \mbox{\texttt{\mdseries\slshape G}} whose index in \mbox{\texttt{\mdseries\slshape G}} is coprime to \mbox{\texttt{\mdseries\slshape p}}. \texttt{SylowSubgroup} computes Sylow subgroups via the operation \texttt{SylowSubgroupOp}. 
\begin{Verbatim}[commandchars=!@|,fontsize=\small,frame=single,label=Example]
  !gapprompt@gap>| !gapinput@g:=SymmetricGroup(4);;|
  !gapprompt@gap>| !gapinput@SylowSubgroup(g,2);|
  Group([ (1,2), (3,4), (1,3)(2,4) ])
\end{Verbatim}
 }

 

\subsection{\textcolor{Chapter }{SylowComplement}}
\logpage{[ 39, 13, 2 ]}\nobreak
\hyperdef{L}{X8605F3FE7A3B8E12}{}
{\noindent\textcolor{FuncColor}{$\triangleright$\enspace\texttt{SylowComplement({\mdseries\slshape G, p})\index{SylowComplement@\texttt{SylowComplement}}
\label{SylowComplement}
}\hfill{\scriptsize (operation)}}\\


 returns a Sylow \mbox{\texttt{\mdseries\slshape p}}\texttt{\symbol{45}}complement of the finite group \mbox{\texttt{\mdseries\slshape G}}. This is a subgroup $U$ of order coprime to \mbox{\texttt{\mdseries\slshape p}} such that the index $[\mbox{\texttt{\mdseries\slshape G}}:U]$ is a \mbox{\texttt{\mdseries\slshape p}}\texttt{\symbol{45}}power. 

 At the moment methods exist only if \mbox{\texttt{\mdseries\slshape G}} is solvable and \textsf{GAP} will issue an error if \mbox{\texttt{\mdseries\slshape G}} is not solvable. 

 
\begin{Verbatim}[commandchars=!@|,fontsize=\small,frame=single,label=Example]
  !gapprompt@gap>| !gapinput@SylowComplement(g,3);|
  Group([ (1,2), (3,4), (1,3)(2,4) ])
\end{Verbatim}
 }

 

\subsection{\textcolor{Chapter }{HallSubgroup}}
\logpage{[ 39, 13, 3 ]}\nobreak
\hyperdef{L}{X7EDBA19E828CD584}{}
{\noindent\textcolor{FuncColor}{$\triangleright$\enspace\texttt{HallSubgroup({\mdseries\slshape G, P})\index{HallSubgroup@\texttt{HallSubgroup}}
\label{HallSubgroup}
}\hfill{\scriptsize (operation)}}\\


 computes a \mbox{\texttt{\mdseries\slshape P}}\texttt{\symbol{45}}Hall subgroup for a set \mbox{\texttt{\mdseries\slshape P}} of primes. This is a subgroup the order of which is only divisible by primes
in \mbox{\texttt{\mdseries\slshape P}} and whose index is coprime to all primes in \mbox{\texttt{\mdseries\slshape P}}. Such a subgroup is unique up to conjugacy if \mbox{\texttt{\mdseries\slshape G}} is solvable. The function computes Hall subgroups via the operation \texttt{HallSubgroupOp}. 

 If \mbox{\texttt{\mdseries\slshape G}} is solvable this function always returns a subgroup. If \mbox{\texttt{\mdseries\slshape G}} is not solvable this function might return a subgroup (if it is unique up to
conjugacy), a list of subgroups (which are representatives of the conjugacy
classes in case there are several such classes) or \texttt{fail} if no such subgroup exists. 
\begin{Verbatim}[commandchars=!@|,fontsize=\small,frame=single,label=Example]
  !gapprompt@gap>| !gapinput@h:=SmallGroup(60,10);;|
  !gapprompt@gap>| !gapinput@u:=HallSubgroup(h,[2,3]);|
  Group([ f1, f2, f3 ])
  !gapprompt@gap>| !gapinput@Size(u);|
  12
  !gapprompt@gap>| !gapinput@h:=PSL(3,5);;|
  !gapprompt@gap>| !gapinput@HallSubgroup(h,[2,3]);|
  [ <permutation group of size 96 with 6 generators>,
    <permutation group of size 96 with 6 generators> ]
  !gapprompt@gap>| !gapinput@u := HallSubgroup(h,[3,31]);;|
  !gapprompt@gap>| !gapinput@Size(u); StructureDescription(u);|
  93
  "C31 : C3"
  !gapprompt@gap>| !gapinput@HallSubgroup(h,[5,31]);|
  fail
\end{Verbatim}
 }

 

\subsection{\textcolor{Chapter }{SylowSystem}}
\logpage{[ 39, 13, 4 ]}\nobreak
\hyperdef{L}{X832E8E6B8347B13F}{}
{\noindent\textcolor{FuncColor}{$\triangleright$\enspace\texttt{SylowSystem({\mdseries\slshape G})\index{SylowSystem@\texttt{SylowSystem}}
\label{SylowSystem}
}\hfill{\scriptsize (attribute)}}\\


 A Sylow system of a group \mbox{\texttt{\mdseries\slshape G}} is a set of Sylow subgroups of \mbox{\texttt{\mdseries\slshape G}} such that every pair of subgroups from this set commutes as subgroups. Sylow
systems exist only for solvable groups. The operation returns \texttt{fail} if the group \mbox{\texttt{\mdseries\slshape G}} is not solvable. 
\begin{Verbatim}[commandchars=!@|,fontsize=\small,frame=single,label=Example]
  !gapprompt@gap>| !gapinput@h:=SmallGroup(60,10);;|
  !gapprompt@gap>| !gapinput@SylowSystem(h);|
  [ Group([ f1, f2 ]), Group([ f3 ]), Group([ f4 ]) ]
  !gapprompt@gap>| !gapinput@List(last,Size);|
  [ 4, 3, 5 ]
\end{Verbatim}
 }

 

\subsection{\textcolor{Chapter }{ComplementSystem}}
\logpage{[ 39, 13, 5 ]}\nobreak
\hyperdef{L}{X87A245E180D27147}{}
{\noindent\textcolor{FuncColor}{$\triangleright$\enspace\texttt{ComplementSystem({\mdseries\slshape G})\index{ComplementSystem@\texttt{ComplementSystem}}
\label{ComplementSystem}
}\hfill{\scriptsize (attribute)}}\\


 A complement system of a group \mbox{\texttt{\mdseries\slshape G}} is a set of Hall $p'$\texttt{\symbol{45}}subgroups of \mbox{\texttt{\mdseries\slshape G}}, where $p'$ runs through the subsets of prime factors of $|\mbox{\texttt{\mdseries\slshape G}}|$ that omit exactly one prime. Every pair of subgroups from this set commutes as
subgroups. Complement systems exist only for solvable groups, therefore \texttt{ComplementSystem} returns \texttt{fail} if the group \mbox{\texttt{\mdseries\slshape G}} is not solvable. 
\begin{Verbatim}[commandchars=!@|,fontsize=\small,frame=single,label=Example]
  !gapprompt@gap>| !gapinput@ComplementSystem(h);|
  [ Group([ f3, f4 ]), Group([ f1, f2, f4 ]), Group([ f1, f2, f3 ]) ]
  !gapprompt@gap>| !gapinput@List(last,Size);|
  [ 15, 20, 12 ]
\end{Verbatim}
 }

 

\subsection{\textcolor{Chapter }{HallSystem}}
\logpage{[ 39, 13, 6 ]}\nobreak
\hyperdef{L}{X82FE5DFD84F8A3C6}{}
{\noindent\textcolor{FuncColor}{$\triangleright$\enspace\texttt{HallSystem({\mdseries\slshape G})\index{HallSystem@\texttt{HallSystem}}
\label{HallSystem}
}\hfill{\scriptsize (attribute)}}\\


 returns a list containing one Hall $P$\texttt{\symbol{45}}subgroup for each set $P$ of prime divisors of the order of \mbox{\texttt{\mdseries\slshape G}}. Hall systems exist only for solvable groups. The operation returns \texttt{fail} if the group \mbox{\texttt{\mdseries\slshape G}} is not solvable. 
\begin{Verbatim}[commandchars=!@|,fontsize=\small,frame=single,label=Example]
  !gapprompt@gap>| !gapinput@HallSystem(h);|
  [ Group([  ]), Group([ f1, f2 ]), Group([ f1, f2, f3 ]),
    Group([ f1, f2, f3, f4 ]), Group([ f1, f2, f4 ]), Group([ f3 ]),
    Group([ f3, f4 ]), Group([ f4 ]) ]
  !gapprompt@gap>| !gapinput@List(last,Size);|
  [ 1, 4, 12, 60, 20, 3, 15, 5 ]
\end{Verbatim}
 }

   }

  
\section{\textcolor{Chapter }{Subgroups characterized by prime powers}}\label{Subgroups characterized by prime powers}
\logpage{[ 39, 14, 0 ]}
\hyperdef{L}{X87AF37E980382499}{}
{
  

\subsection{\textcolor{Chapter }{Omega}}
\logpage{[ 39, 14, 1 ]}\nobreak
\hyperdef{L}{X7F069ACC83DB3374}{}
{\noindent\textcolor{FuncColor}{$\triangleright$\enspace\texttt{Omega({\mdseries\slshape G, p[, n]})\index{Omega@\texttt{Omega}}
\label{Omega}
}\hfill{\scriptsize (operation)}}\\


 For a \mbox{\texttt{\mdseries\slshape p}}\texttt{\symbol{45}}group \mbox{\texttt{\mdseries\slshape G}}, one defines $\Omega_{\mbox{\texttt{\mdseries\slshape n}}}(\mbox{\texttt{\mdseries\slshape G}}) = \langle g \in \mbox{\texttt{\mdseries\slshape G}} \mid g^{{\mbox{\texttt{\mdseries\slshape p}}^{\mbox{\texttt{\mdseries\slshape n}}}}} = 1 \rangle$. The default value for \mbox{\texttt{\mdseries\slshape n}} is \texttt{1}. 

 
\begin{Verbatim}[commandchars=!@|,fontsize=\small,frame=single,label=Example]
  !gapprompt@gap>| !gapinput@h:=SmallGroup(16,10);|
  <pc group of size 16 with 4 generators>
  !gapprompt@gap>| !gapinput@Omega(h,2);|
  Group([ f2, f3, f4 ])
\end{Verbatim}
 }

 

\subsection{\textcolor{Chapter }{Agemo}}
\logpage{[ 39, 14, 2 ]}\nobreak
\hyperdef{L}{X83DB33747F069ACC}{}
{\noindent\textcolor{FuncColor}{$\triangleright$\enspace\texttt{Agemo({\mdseries\slshape G, p[, n]})\index{Agemo@\texttt{Agemo}}
\label{Agemo}
}\hfill{\scriptsize (function)}}\\


 For a \mbox{\texttt{\mdseries\slshape p}}\texttt{\symbol{45}}group \mbox{\texttt{\mdseries\slshape G}}, one defines $\mho_{\mbox{\texttt{\mdseries\slshape n}}}(G) = \langle g^{{\mbox{\texttt{\mdseries\slshape p}}^{\mbox{\texttt{\mdseries\slshape n}}}}} \mid g \in \mbox{\texttt{\mdseries\slshape G}} \rangle$. The default value for \mbox{\texttt{\mdseries\slshape n}} is \texttt{1}. 
\begin{Verbatim}[commandchars=!@|,fontsize=\small,frame=single,label=Example]
  !gapprompt@gap>| !gapinput@Agemo(h,2);Agemo(h,2,2);|
  Group([ f4 ])
  Group([  ])
\end{Verbatim}
 }

 }

  
\section{\textcolor{Chapter }{Group Properties}}\label{Group Properties}
\logpage{[ 39, 15, 0 ]}
\hyperdef{L}{X7B75879B8085120A}{}
{
  Some properties of groups can be defined not only for groups but also for
other structures. For example, nilpotency and solvability make sense also for
algebras. Note that these names refer to different definitions for groups and
algebras, contrary to the situation with finiteness or commutativity. In such
cases, the name of the function for groups got a suffix \texttt{Group} to distinguish different meanings for different structures. 

 Some functions, such as \texttt{IsPSolvable} (\ref{IsPSolvable}) and \texttt{IsPNilpotent} (\ref{IsPNilpotent}), although they are mathematical properties, are not properties in the sense
of \textsf{GAP} (see{\nobreakspace}\ref{Attributes} and \ref{Properties}), as they depend on a parameter. 

\subsection{\textcolor{Chapter }{IsCyclic}}
\logpage{[ 39, 15, 1 ]}\nobreak
\hyperdef{L}{X7DA27D338374FD28}{}
{\noindent\textcolor{FuncColor}{$\triangleright$\enspace\texttt{IsCyclic({\mdseries\slshape G})\index{IsCyclic@\texttt{IsCyclic}}
\label{IsCyclic}
}\hfill{\scriptsize (property)}}\\


 A group is \emph{cyclic} if it can be generated by one element. For a cyclic group, one can compute a
generating set consisting of only one element using \texttt{MinimalGeneratingSet} (\ref{MinimalGeneratingSet}). }

 

\subsection{\textcolor{Chapter }{IsElementaryAbelian}}
\logpage{[ 39, 15, 2 ]}\nobreak
\hyperdef{L}{X813C952F80E775D4}{}
{\noindent\textcolor{FuncColor}{$\triangleright$\enspace\texttt{IsElementaryAbelian({\mdseries\slshape G})\index{IsElementaryAbelian@\texttt{IsElementaryAbelian}}
\label{IsElementaryAbelian}
}\hfill{\scriptsize (property)}}\\


 A group \mbox{\texttt{\mdseries\slshape G}} is elementary abelian if it is commutative and if there is a prime $p$ such that the order of each element in \mbox{\texttt{\mdseries\slshape G}} divides $p$. }

 

\subsection{\textcolor{Chapter }{IsNilpotentGroup}}
\logpage{[ 39, 15, 3 ]}\nobreak
\hyperdef{L}{X87D062608719F2CD}{}
{\noindent\textcolor{FuncColor}{$\triangleright$\enspace\texttt{IsNilpotentGroup({\mdseries\slshape G})\index{IsNilpotentGroup@\texttt{IsNilpotentGroup}}
\label{IsNilpotentGroup}
}\hfill{\scriptsize (property)}}\\


 A group is \emph{nilpotent} if the lower central series (see{\nobreakspace}\texttt{LowerCentralSeriesOfGroup} (\ref{LowerCentralSeriesOfGroup}) for a definition) reaches the trivial subgroup in a finite number of steps. }

 

\subsection{\textcolor{Chapter }{NilpotencyClassOfGroup}}
\logpage{[ 39, 15, 4 ]}\nobreak
\hyperdef{L}{X7E3056237C6A5D43}{}
{\noindent\textcolor{FuncColor}{$\triangleright$\enspace\texttt{NilpotencyClassOfGroup({\mdseries\slshape G})\index{NilpotencyClassOfGroup@\texttt{NilpotencyClassOfGroup}}
\label{NilpotencyClassOfGroup}
}\hfill{\scriptsize (attribute)}}\\


 The nilpotency class of a nilpotent group \mbox{\texttt{\mdseries\slshape G}} is the number of steps in the lower central series of \mbox{\texttt{\mdseries\slshape G}} (see \texttt{LowerCentralSeriesOfGroup} (\ref{LowerCentralSeriesOfGroup})); 

 If \mbox{\texttt{\mdseries\slshape G}} is not nilpotent an error is issued. }

 

\subsection{\textcolor{Chapter }{IsPerfectGroup}}
\logpage{[ 39, 15, 5 ]}\nobreak
\hyperdef{L}{X8755147280C84DBB}{}
{\noindent\textcolor{FuncColor}{$\triangleright$\enspace\texttt{IsPerfectGroup({\mdseries\slshape G})\index{IsPerfectGroup@\texttt{IsPerfectGroup}}
\label{IsPerfectGroup}
}\hfill{\scriptsize (property)}}\\


 A group is \emph{perfect} if it equals its derived subgroup (see{\nobreakspace}\texttt{DerivedSubgroup} (\ref{DerivedSubgroup})). }

 

\subsection{\textcolor{Chapter }{IsSolvableGroup}}
\logpage{[ 39, 15, 6 ]}\nobreak
\hyperdef{L}{X809C78D5877D31DF}{}
{\noindent\textcolor{FuncColor}{$\triangleright$\enspace\texttt{IsSolvableGroup({\mdseries\slshape G})\index{IsSolvableGroup@\texttt{IsSolvableGroup}}
\label{IsSolvableGroup}
}\hfill{\scriptsize (property)}}\\


 A group is \emph{solvable} if the derived series (see{\nobreakspace}\texttt{DerivedSeriesOfGroup} (\ref{DerivedSeriesOfGroup}) for a definition) reaches the trivial subgroup in a finite number of steps. 

 For finite groups this is the same as being polycyclic (see{\nobreakspace}\texttt{IsPolycyclicGroup} (\ref{IsPolycyclicGroup})), and each polycyclic group is solvable, but there are infinite solvable
groups that are not polycyclic. }

 

\subsection{\textcolor{Chapter }{IsPolycyclicGroup}}
\logpage{[ 39, 15, 7 ]}\nobreak
\hyperdef{L}{X7D7456077D3D1B86}{}
{\noindent\textcolor{FuncColor}{$\triangleright$\enspace\texttt{IsPolycyclicGroup({\mdseries\slshape G})\index{IsPolycyclicGroup@\texttt{IsPolycyclicGroup}}
\label{IsPolycyclicGroup}
}\hfill{\scriptsize (property)}}\\


 A group is polycyclic if it has a subnormal series with cyclic factors. For
finite groups this is the same as if the group is solvable (see{\nobreakspace}\texttt{IsSolvableGroup} (\ref{IsSolvableGroup})). }

 

\subsection{\textcolor{Chapter }{IsSupersolvableGroup}}
\logpage{[ 39, 15, 8 ]}\nobreak
\hyperdef{L}{X7AADF2E88501B9FF}{}
{\noindent\textcolor{FuncColor}{$\triangleright$\enspace\texttt{IsSupersolvableGroup({\mdseries\slshape G})\index{IsSupersolvableGroup@\texttt{IsSupersolvableGroup}}
\label{IsSupersolvableGroup}
}\hfill{\scriptsize (property)}}\\


 A finite group is \emph{supersolvable} if it has a normal series with cyclic factors. }

 

\subsection{\textcolor{Chapter }{IsMonomialGroup}}
\logpage{[ 39, 15, 9 ]}\nobreak
\hyperdef{L}{X83977EB97A8E2290}{}
{\noindent\textcolor{FuncColor}{$\triangleright$\enspace\texttt{IsMonomialGroup({\mdseries\slshape G})\index{IsMonomialGroup@\texttt{IsMonomialGroup}}
\label{IsMonomialGroup}
}\hfill{\scriptsize (property)}}\\


 A finite group is \emph{monomial} if every irreducible complex character is induced from a linear character of a
subgroup. }

 

\subsection{\textcolor{Chapter }{IsSimpleGroup}}
\logpage{[ 39, 15, 10 ]}\nobreak
\hyperdef{L}{X7A6685D7819AEC32}{}
{\noindent\textcolor{FuncColor}{$\triangleright$\enspace\texttt{IsSimpleGroup({\mdseries\slshape G})\index{IsSimpleGroup@\texttt{IsSimpleGroup}}
\label{IsSimpleGroup}
}\hfill{\scriptsize (property)}}\\
\noindent\textcolor{FuncColor}{$\triangleright$\enspace\texttt{IsNonabelianSimpleGroup({\mdseries\slshape G})\index{IsNonabelianSimpleGroup@\texttt{IsNonabelianSimpleGroup}}
\label{IsNonabelianSimpleGroup}
}\hfill{\scriptsize (property)}}\\


 A group is \emph{simple} if it is nontrivial and has no nontrivial normal subgroups. A \emph{nonabelian simple} group is simple and not abelian. }

 

\subsection{\textcolor{Chapter }{IsAlmostSimpleGroup}}
\logpage{[ 39, 15, 11 ]}\nobreak
\hyperdef{L}{X78CC9764803601E7}{}
{\noindent\textcolor{FuncColor}{$\triangleright$\enspace\texttt{IsAlmostSimpleGroup({\mdseries\slshape G})\index{IsAlmostSimpleGroup@\texttt{IsAlmostSimpleGroup}}
\label{IsAlmostSimpleGroup}
}\hfill{\scriptsize (property)}}\\


 A group \mbox{\texttt{\mdseries\slshape G}} is \emph{almost simple} if a nonabelian simple group $S$ exists such that \mbox{\texttt{\mdseries\slshape G}} is isomorphic to a subgroup of the automorphism group of $S$ that contains all inner automorphisms of $S$. 

 Equivalently, \mbox{\texttt{\mdseries\slshape G}} is almost simple if and only if it has a unique minimal normal subgroup $N$ and if $N$ is a nonabelian simple group. 

  Note that an almost simple group is \emph{not} defined as an extension of a simple group by outer automorphisms, since we
want to exclude extensions of groups of prime order. In particular, a \emph{simple} group is \emph{almost simple} if and only if it is nonabelian. 

 
\begin{Verbatim}[commandchars=!@|,fontsize=\small,frame=single,label=Example]
  !gapprompt@gap>| !gapinput@IsAlmostSimpleGroup( AlternatingGroup( 5 ) );|
  true
  !gapprompt@gap>| !gapinput@IsAlmostSimpleGroup( SymmetricGroup( 5 ) );|
  true
  !gapprompt@gap>| !gapinput@IsAlmostSimpleGroup( SymmetricGroup( 3 ) );|
  false
  !gapprompt@gap>| !gapinput@IsAlmostSimpleGroup( SL( 2, 5 ) );|
  false
\end{Verbatim}
 }

 

\subsection{\textcolor{Chapter }{IsQuasisimpleGroup}}
\logpage{[ 39, 15, 12 ]}\nobreak
\hyperdef{L}{X7C1709A986B00F97}{}
{\noindent\textcolor{FuncColor}{$\triangleright$\enspace\texttt{IsQuasisimpleGroup({\mdseries\slshape G})\index{IsQuasisimpleGroup@\texttt{IsQuasisimpleGroup}}
\label{IsQuasisimpleGroup}
}\hfill{\scriptsize (property)}}\\


 A group \mbox{\texttt{\mdseries\slshape G}} is \emph{quasisimple} if \mbox{\texttt{\mdseries\slshape G}} is perfect (see \texttt{IsPerfectGroup} (\ref{IsPerfectGroup})) and if \mbox{\texttt{\mdseries\slshape G}}$/Z($\mbox{\texttt{\mdseries\slshape G}}$)$ is simple (see \texttt{IsSimpleGroup} (\ref{IsSimpleGroup})), where $Z($\mbox{\texttt{\mdseries\slshape G}}$)$ is the centre of \mbox{\texttt{\mdseries\slshape G}} (see \texttt{Centre} (\ref{Centre})). 

 
\begin{Verbatim}[commandchars=!@|,fontsize=\small,frame=single,label=Example]
  !gapprompt@gap>| !gapinput@IsQuasisimpleGroup( AlternatingGroup( 5 ) );|
  true
  !gapprompt@gap>| !gapinput@IsQuasisimpleGroup( SymmetricGroup( 5 ) );|
  false
  !gapprompt@gap>| !gapinput@IsQuasisimpleGroup( SL( 2, 5 ) );|
  true
\end{Verbatim}
 }

 
\subsection{\textcolor{Chapter }{IsomorphismTypeInfoFiniteSimpleGroup}}\logpage{[ 39, 15, 13 ]}
\hyperdef{L}{X7C6AA6897C4409AC}{}
{
\noindent\textcolor{FuncColor}{$\triangleright$\enspace\texttt{IsomorphismTypeInfoFiniteSimpleGroup({\mdseries\slshape G})\index{IsomorphismTypeInfoFiniteSimpleGroup@\texttt{Isomorphism}\-\texttt{Type}\-\texttt{Info}\-\texttt{Finite}\-\texttt{Simple}\-\texttt{Group}!for a group}
\label{IsomorphismTypeInfoFiniteSimpleGroup:for a group}
}\hfill{\scriptsize (attribute)}}\\
\noindent\textcolor{FuncColor}{$\triangleright$\enspace\texttt{IsomorphismTypeInfoFiniteSimpleGroup({\mdseries\slshape n})\index{IsomorphismTypeInfoFiniteSimpleGroup@\texttt{Isomorphism}\-\texttt{Type}\-\texttt{Info}\-\texttt{Finite}\-\texttt{Simple}\-\texttt{Group}!for a group order}
\label{IsomorphismTypeInfoFiniteSimpleGroup:for a group order}
}\hfill{\scriptsize (attribute)}}\\


 For a finite simple group \mbox{\texttt{\mdseries\slshape G}}, \texttt{IsomorphismTypeInfoFiniteSimpleGroup} returns a record with the components \texttt{name}, \texttt{shortname}, \texttt{series}, and possibly \texttt{parameter}, describing the isomorphism type of \mbox{\texttt{\mdseries\slshape G}}. 

 The values of the components \texttt{name}, \texttt{shortname}, and \texttt{series} are strings, \texttt{name} gives name(s) for \mbox{\texttt{\mdseries\slshape G}}, \texttt{shortname} gives one name for \mbox{\texttt{\mdseries\slshape G}} that is compatible with the naming scheme used in the \textsf{GAP} packages \textsf{CTblLib} and \textsf{AtlasRep} (and in the \textsf{Atlas} of Finite Groups{\nobreakspace}\cite{CCN85}), and \texttt{series} describes the following series. 

 (If different characterizations of \mbox{\texttt{\mdseries\slshape G}} are possible only one is given by \texttt{series} and \texttt{parameter}, while \texttt{name} may give several names.) 
\begin{description}
\item[{\texttt{"A"}}]  Alternating groups, \texttt{parameter} gives the natural degree. 
\item[{\texttt{"L"}}]  Linear groups (Chevalley type $A$), \texttt{parameter} is a list $[ n, q ]$ that indicates $L(n,q)$. 
\item[{\texttt{"2A"}}]  Twisted Chevalley type ${}^2A$, \texttt{parameter} is a list $[ n, q ]$ that indicates ${}^2A(n,q)$. 
\item[{\texttt{"B"}}]  Chevalley type $B$, \texttt{parameter} is a list $[n, q ]$ that indicates $B(n,q)$. 
\item[{\texttt{"2B"}}]  Twisted Chevalley type ${}^2B$, \texttt{parameter} is a value $q$ that indicates ${}^2B(2,q)$. 
\item[{\texttt{"C"}}]  Chevalley type $C$, \texttt{parameter} is a list $[ n, q ]$ that indicates $C(n,q)$. 
\item[{\texttt{"D"}}]  Chevalley type $D$, \texttt{parameter} is a list $[ n, q ]$ that indicates $D(n,q)$. 
\item[{\texttt{"2D"}}]  Twisted Chevalley type ${}^2D$, \texttt{parameter} is a list $[ n, q ]$ that indicates ${}^2D(n,q)$. 
\item[{\texttt{"3D"}}]  Twisted Chevalley type ${}^3D$, \texttt{parameter} is a value $q$ that indicates ${}^3D(4,q)$. 
\item[{\texttt{"E"}}]  Exceptional Chevalley type $E$, \texttt{parameter} is a list $[ n, q ]$ that indicates $E_n(q)$. The value of \mbox{\texttt{\mdseries\slshape n}} is 6, 7, or 8. 
\item[{\texttt{"2E"}}]  Twisted exceptional Chevalley type $E_6$, \texttt{parameter} is a value $q$ that indicates ${}^2E_6(q)$. 
\item[{\texttt{"F"}}]  Exceptional Chevalley type $F$, \texttt{parameter} is a value $q$ that indicates $F(4,q)$. 
\item[{\texttt{"2F"}}]  Twisted exceptional Chevalley type ${}^2F$ (Ree groups), \texttt{parameter} is a value $q$ that indicates ${}^2F(4,q)$. 
\item[{\texttt{"G"}}]  Exceptional Chevalley type $G$, \texttt{parameter} is a value $q$ that indicates $G(2,q)$. 
\item[{\texttt{"2G"}}]  Twisted exceptional Chevalley type ${}^2G$ (Ree groups), \texttt{parameter} is a value $q$ that indicates ${}^2G(2,q)$. 
\item[{\texttt{"Spor"}}]  Sporadic simple groups, \texttt{name} gives the name. 
\item[{\texttt{"Z"}}]  Cyclic groups of prime size, \texttt{parameter} gives the size. 
\end{description}
 

 An equal sign in the name denotes different naming schemes for the same group,
a tilde sign abstract isomorphisms between groups constructed in a different
way. 

 
\begin{Verbatim}[commandchars=!@|,fontsize=\small,frame=single,label=Example]
  !gapprompt@gap>| !gapinput@IsomorphismTypeInfoFiniteSimpleGroup(|
  !gapprompt@>| !gapinput@                            Group((4,5)(6,7),(1,2,4)(3,5,6)));|
  rec(
    name := "A(1,7) = L(2,7) ~ B(1,7) = O(3,7) ~ C(1,7) = S(2,7) ~ 2A(1,\
  7) = U(2,7) ~ A(2,2) = L(3,2)", parameter := [ 2, 7 ], series := "L",
    shortname := "L3(2)" )
\end{Verbatim}
 

 For a positive integer \mbox{\texttt{\mdseries\slshape n}}, \texttt{IsomorphismTypeInfoFiniteSimpleGroup} returns \texttt{fail} if \mbox{\texttt{\mdseries\slshape n}} is not the order of a finite simple group, and a record as described for the
case of a group \mbox{\texttt{\mdseries\slshape G}} otherwise. If more than one simple group of order \mbox{\texttt{\mdseries\slshape n}} exists then the result record contains only the \texttt{name} component, a string that lists the two possible isomorphism types of simple
groups of this order. 

 
\begin{Verbatim}[commandchars=!@|,fontsize=\small,frame=single,label=Example]
  !gapprompt@gap>| !gapinput@IsomorphismTypeInfoFiniteSimpleGroup( 5 );|
  rec( name := "Z(5)", parameter := 5, series := "Z", shortname := "C5"
   )
  !gapprompt@gap>| !gapinput@IsomorphismTypeInfoFiniteSimpleGroup( 6 );|
  fail
  !gapprompt@gap>| !gapinput@IsomorphismTypeInfoFiniteSimpleGroup(Size(SymplecticGroup(6,3))/2);|
  rec(
    name := "cannot decide from size alone between B(3,3) = O(7,3) and C\
  (3,3) = S(6,3)", parameter := [ 3, 3 ] )
\end{Verbatim}
 }

 

\subsection{\textcolor{Chapter }{SimpleGroup}}
\logpage{[ 39, 15, 14 ]}\nobreak
\hyperdef{L}{X8492B05B822AC58C}{}
{\noindent\textcolor{FuncColor}{$\triangleright$\enspace\texttt{SimpleGroup({\mdseries\slshape id[, param]})\index{SimpleGroup@\texttt{SimpleGroup}}
\label{SimpleGroup}
}\hfill{\scriptsize (function)}}\\


 This function will construct \textsc{an} instance of the specified nonabelian simple group. Groups are specified via
their name in ATLAS style notation, with parameters added if necessary. The
intelligence applied to parsing the name is limited, and at the moment no
proper extensions can be constructed. For groups which do not have a
permutation representation of small degree the \textsf{AtlasRep} package might need to be installed to construct these groups. 
\begin{Verbatim}[commandchars=!@|,fontsize=\small,frame=single,label=Example]
  !gapprompt@gap>| !gapinput@g:=SimpleGroup("M(23)");|
  M23
  !gapprompt@gap>| !gapinput@Size(g);|
  10200960
  !gapprompt@gap>| !gapinput@g:=SimpleGroup("PSL",3,5);|
  PSL(3,5)
  !gapprompt@gap>| !gapinput@Size(g);|
  372000
  !gapprompt@gap>| !gapinput@g:=SimpleGroup("PSp6",2);|
  PSp(6,2)
\end{Verbatim}
 }

 

\subsection{\textcolor{Chapter }{SimpleGroupsIterator}}
\logpage{[ 39, 15, 15 ]}\nobreak
\hyperdef{L}{X839CDD8C7AE39FD6}{}
{\noindent\textcolor{FuncColor}{$\triangleright$\enspace\texttt{SimpleGroupsIterator({\mdseries\slshape [start[, end]]})\index{SimpleGroupsIterator@\texttt{SimpleGroupsIterator}}
\label{SimpleGroupsIterator}
}\hfill{\scriptsize (function)}}\\


 This function returns an iterator that will run over all nonabelian simple
groups, starting at order \mbox{\texttt{\mdseries\slshape start}} if specified, up to order $10^{27}$ (or \texttt{\symbol{45}}\texttt{\symbol{45}} if specified
\texttt{\symbol{45}}\texttt{\symbol{45}} order \mbox{\texttt{\mdseries\slshape end}}). If the option \mbox{\texttt{\mdseries\slshape NOPSL2}} is given, groups of type $PSL_2(q)$ are omitted. 
\begin{Verbatim}[commandchars=!@|,fontsize=\small,frame=single,label=Example]
  !gapprompt@gap>| !gapinput@it:=SimpleGroupsIterator(20000);|
  <iterator>
  !gapprompt@gap>| !gapinput@List([1..8],x->NextIterator(it));|
  [ A8, PSL(3,4), PSL(2,37), PSp(4,3), Sz(8), PSL(2,32), PSL(2,41),
    PSL(2,43) ]
  !gapprompt@gap>| !gapinput@it:=SimpleGroupsIterator(1,2000);;|
  !gapprompt@gap>| !gapinput@l:=[];;for i in it do Add(l,i);od;l;|
  [ A5, PSL(2,7), A6, PSL(2,8), PSL(2,11), PSL(2,13) ]
  !gapprompt@gap>| !gapinput@it:=SimpleGroupsIterator(20000,100000:NOPSL2);;|
  !gapprompt@gap>| !gapinput@l:=[];;for i in it do Add(l,i);od;l;|
  [ A8, PSL(3,4), PSp(4,3), Sz(8), PSU(3,4), M12 ]
\end{Verbatim}
 }

 

\subsection{\textcolor{Chapter }{SmallSimpleGroup}}
\logpage{[ 39, 15, 16 ]}\nobreak
\hyperdef{L}{X872E93F586F54FCE}{}
{\noindent\textcolor{FuncColor}{$\triangleright$\enspace\texttt{SmallSimpleGroup({\mdseries\slshape order[, i]})\index{SmallSimpleGroup@\texttt{SmallSimpleGroup}}
\label{SmallSimpleGroup}
}\hfill{\scriptsize (function)}}\\
\textbf{\indent Returns:\ }
 The \mbox{\texttt{\mdseries\slshape i}}th simple group of order \mbox{\texttt{\mdseries\slshape order}} in the stored list, given in a small\texttt{\symbol{45}}degree permutation
representation, or \texttt{fail} (\ref{fail}) if no such simple group exists. 



 If \mbox{\texttt{\mdseries\slshape i}} is not given, it defaults to{\nobreakspace}1. Currently, all simple groups of
order less than $10^6$ are available via this function. 
\begin{Verbatim}[commandchars=!@|,fontsize=\small,frame=single,label=Example]
  !gapprompt@gap>| !gapinput@SmallSimpleGroup(60);|
  A5
  !gapprompt@gap>| !gapinput@SmallSimpleGroup(20160,1);|
  A8
  !gapprompt@gap>| !gapinput@SmallSimpleGroup(20160,2);|
  PSL(3,4)
\end{Verbatim}
 }

 

\subsection{\textcolor{Chapter }{AllSmallNonabelianSimpleGroups}}
\logpage{[ 39, 15, 17 ]}\nobreak
\hyperdef{L}{X7EB47BF27D8CBF72}{}
{\noindent\textcolor{FuncColor}{$\triangleright$\enspace\texttt{AllSmallNonabelianSimpleGroups({\mdseries\slshape orders})\index{AllSmallNonabelianSimpleGroups@\texttt{AllSmallNonabelianSimpleGroups}}
\label{AllSmallNonabelianSimpleGroups}
}\hfill{\scriptsize (function)}}\\
\textbf{\indent Returns:\ }
 A list of all nonabelian simple groups whose order lies in the range \mbox{\texttt{\mdseries\slshape orders}}. 



 The groups are given in small\texttt{\symbol{45}}degree permutation
representations. The returned list is sorted by ascending group order.
Currently, all simple groups of order less than $10^6$ are available via this function. 
\begin{Verbatim}[commandchars=!@|,fontsize=\small,frame=single,label=Example]
  !gapprompt@gap>| !gapinput@List(AllSmallNonabelianSimpleGroups([1..1000000]),|
  !gapprompt@>| !gapinput@        StructureDescription);|
  [ "A5", "PSL(3,2)", "A6", "PSL(2,8)", "PSL(2,11)", "PSL(2,13)",
    "PSL(2,17)", "A7", "PSL(2,19)", "PSL(2,16)", "PSL(3,3)",
    "PSU(3,3)", "PSL(2,23)", "PSL(2,25)", "M11", "PSL(2,27)",
    "PSL(2,29)", "PSL(2,31)", "A8", "PSL(3,4)", "PSL(2,37)", "O(5,3)",
    "Sz(8)", "PSL(2,32)", "PSL(2,41)", "PSL(2,43)", "PSL(2,47)",
    "PSL(2,49)", "PSU(3,4)", "PSL(2,53)", "M12", "PSL(2,59)",
    "PSL(2,61)", "PSU(3,5)", "PSL(2,67)", "J1", "PSL(2,71)", "A9",
    "PSL(2,73)", "PSL(2,79)", "PSL(2,64)", "PSL(2,81)", "PSL(2,83)",
    "PSL(2,89)", "PSL(3,5)", "M22", "PSL(2,97)", "PSL(2,101)",
    "PSL(2,103)", "HJ", "PSL(2,107)", "PSL(2,109)", "PSL(2,113)",
    "PSL(2,121)", "PSL(2,125)", "O(5,4)" ]
\end{Verbatim}
 }

 

\subsection{\textcolor{Chapter }{IsFinitelyGeneratedGroup}}
\logpage{[ 39, 15, 18 ]}\nobreak
\hyperdef{L}{X81E22D07871DF37E}{}
{\noindent\textcolor{FuncColor}{$\triangleright$\enspace\texttt{IsFinitelyGeneratedGroup({\mdseries\slshape G})\index{IsFinitelyGeneratedGroup@\texttt{IsFinitelyGeneratedGroup}}
\label{IsFinitelyGeneratedGroup}
}\hfill{\scriptsize (property)}}\\


 tests whether the group \mbox{\texttt{\mdseries\slshape G}} can be generated by a finite number of generators. (This property is mainly
used to obtain finiteness conditions.) 

 Note that this is a pure existence statement. Even if a group is known to be
generated by a finite number of elements, it can be very hard or even
impossible to obtain such a generating set if it is not known. }

 

\subsection{\textcolor{Chapter }{IsSubsetLocallyFiniteGroup}}
\logpage{[ 39, 15, 19 ]}\nobreak
\hyperdef{L}{X8648EDA287829755}{}
{\noindent\textcolor{FuncColor}{$\triangleright$\enspace\texttt{IsSubsetLocallyFiniteGroup({\mdseries\slshape U})\index{IsSubsetLocallyFiniteGroup@\texttt{IsSubsetLocallyFiniteGroup}}
\label{IsSubsetLocallyFiniteGroup}
}\hfill{\scriptsize (property)}}\\


 A group is called locally finite if every finitely generated subgroup is
finite. This property checks whether the group \mbox{\texttt{\mdseries\slshape U}} is a subset of a locally finite group. This is used to check whether finite
generation will imply finiteness, as it does for example for permutation
groups. }

 

\subsection{\textcolor{Chapter }{IsPGroup}}
\logpage{[ 39, 15, 20 ]}\nobreak
\hyperdef{L}{X8089F18C810B7E3E}{}
{\noindent\textcolor{FuncColor}{$\triangleright$\enspace\texttt{IsPGroup({\mdseries\slshape G})\index{IsPGroup@\texttt{IsPGroup}}
\label{IsPGroup}
}\hfill{\scriptsize (property)}}\\


 \index{p-group@$p$\texttt{\symbol{45}}group} A \emph{$p$\texttt{\symbol{45}}group} is a group in which the order (see{\nobreakspace}\texttt{Order} (\ref{Order})) of every element is of the form $p^n$ for a prime integer $p$ and a nonnegative integer $n$. \texttt{IsPGroup} returns \texttt{true} if \mbox{\texttt{\mdseries\slshape G}} is a $p$\texttt{\symbol{45}}group, and \texttt{false} otherwise. 

 Finite $p$\texttt{\symbol{45}}groups are precisely those groups whose order
(see{\nobreakspace}\texttt{Size} (\ref{Size})) is $1$ or a prime power (see{\nobreakspace}\texttt{IsPrimePowerInt} (\ref{IsPrimePowerInt}), and are always nilpotent. 

 Note that $p$\texttt{\symbol{45}}groups can also be infinite, and in that case, need not be
nilpotent. 

 
\begin{Verbatim}[commandchars=!@|,fontsize=\small,frame=single,label=Example]
  !gapprompt@gap>| !gapinput@IsPGroup( DihedralGroup( 8 ) );|
  true
  !gapprompt@gap>| !gapinput@IsPGroup( TrivialGroup() );|
  true
  !gapprompt@gap>| !gapinput@IsPGroup( DihedralGroup( 10 ) );|
  false
\end{Verbatim}
 }

 

\subsection{\textcolor{Chapter }{IsPowerfulPGroup}}
\logpage{[ 39, 15, 21 ]}\nobreak
\hyperdef{L}{X7F232B3F8261CE25}{}
{\noindent\textcolor{FuncColor}{$\triangleright$\enspace\texttt{IsPowerfulPGroup({\mdseries\slshape G})\index{IsPowerfulPGroup@\texttt{IsPowerfulPGroup}}
\label{IsPowerfulPGroup}
}\hfill{\scriptsize (property)}}\\


 \index{PowerfulPGroup@Powerful $p$\texttt{\symbol{45}}group} A finite $p$\texttt{\symbol{45}}group \mbox{\texttt{\mdseries\slshape G}} is said to be a \emph{powerful $p$\texttt{\symbol{45}}group} if the commutator subgroup $[\mbox{\texttt{\mdseries\slshape G}},\mbox{\texttt{\mdseries\slshape G}}]$ is contained in $\mbox{\texttt{\mdseries\slshape G}}^{p}$ if the prime $p$ is odd, or if $[\mbox{\texttt{\mdseries\slshape G}},\mbox{\texttt{\mdseries\slshape G}}]$ is contained in $\mbox{\texttt{\mdseries\slshape G}}^{4}$ if $p = 2$. The subgroup $\mbox{\texttt{\mdseries\slshape G}}^{p}$ is called the first Agemo subgroup, (see{\nobreakspace}\texttt{Agemo} (\ref{Agemo})). \texttt{IsPowerfulPGroup} returns \texttt{true} if \mbox{\texttt{\mdseries\slshape G}} is a powerful $p$\texttt{\symbol{45}}group, and \texttt{false} otherwise. \emph{Note: }This function returns \texttt{true} if \mbox{\texttt{\mdseries\slshape G}} is the trivial group. }

 

\subsection{\textcolor{Chapter }{IsRegularPGroup}}
\logpage{[ 39, 15, 22 ]}\nobreak
\hyperdef{L}{X7ED4A14F7A235617}{}
{\noindent\textcolor{FuncColor}{$\triangleright$\enspace\texttt{IsRegularPGroup({\mdseries\slshape G})\index{IsRegularPGroup@\texttt{IsRegularPGroup}}
\label{IsRegularPGroup}
}\hfill{\scriptsize (property)}}\\


 \index{RegularPGroup@Regular $p$\texttt{\symbol{45}}group} A finite $p$\texttt{\symbol{45}}group \mbox{\texttt{\mdseries\slshape G}} is a \emph{regular $p$\texttt{\symbol{45}}group} if for all $a,b$ in \mbox{\texttt{\mdseries\slshape G}}, one has $a^p b^p = (a b)^p c^p$ where $c$ is an element of the derived subgroup of the group generated by $a$ and $b$ (see{\nobreakspace}\cite{Hal34}). \texttt{IsRegularPGroup} returns \texttt{true} if \mbox{\texttt{\mdseries\slshape G}} is a regular $p$\texttt{\symbol{45}}group, and \texttt{false} otherwise. \emph{Note: }This function returns \texttt{true} if \mbox{\texttt{\mdseries\slshape G}} is the trivial group. }

 

\subsection{\textcolor{Chapter }{PrimePGroup}}
\logpage{[ 39, 15, 23 ]}\nobreak
\hyperdef{L}{X87356BAA7E9E2142}{}
{\noindent\textcolor{FuncColor}{$\triangleright$\enspace\texttt{PrimePGroup({\mdseries\slshape G})\index{PrimePGroup@\texttt{PrimePGroup}}
\label{PrimePGroup}
}\hfill{\scriptsize (attribute)}}\\


 If \mbox{\texttt{\mdseries\slshape G}} is a nontrivial $p$\texttt{\symbol{45}}group (see{\nobreakspace}\texttt{IsPGroup} (\ref{IsPGroup})), \texttt{PrimePGroup} returns the prime integer $p$; if \mbox{\texttt{\mdseries\slshape G}} is trivial then \texttt{PrimePGroup} returns \texttt{fail}. Otherwise an error is issued. 

 (One should avoid a common error of writing \texttt{if IsPGroup(g) then ... PrimePGroup(g) ...} where the code represented by dots assumes that \texttt{PrimePGroup(g)} is an integer.) }

 

\subsection{\textcolor{Chapter }{PClassPGroup}}
\logpage{[ 39, 15, 24 ]}\nobreak
\hyperdef{L}{X863434AD7DDE514B}{}
{\noindent\textcolor{FuncColor}{$\triangleright$\enspace\texttt{PClassPGroup({\mdseries\slshape G})\index{PClassPGroup@\texttt{PClassPGroup}}
\label{PClassPGroup}
}\hfill{\scriptsize (attribute)}}\\


 The $p$\texttt{\symbol{45}}class of a $p$\texttt{\symbol{45}}group \mbox{\texttt{\mdseries\slshape G}} (see{\nobreakspace}\texttt{IsPGroup} (\ref{IsPGroup})) is the number of steps in the lower $p$\texttt{\symbol{45}}central series (see{\nobreakspace}\texttt{PCentralSeries} (\ref{PCentralSeries})) of \mbox{\texttt{\mdseries\slshape G}}. If \mbox{\texttt{\mdseries\slshape G}} is not a $p$\texttt{\symbol{45}}group then an error is issued. }

 

\subsection{\textcolor{Chapter }{RankPGroup}}
\logpage{[ 39, 15, 25 ]}\nobreak
\hyperdef{L}{X840A4F937ABF15E1}{}
{\noindent\textcolor{FuncColor}{$\triangleright$\enspace\texttt{RankPGroup({\mdseries\slshape G})\index{RankPGroup@\texttt{RankPGroup}}
\label{RankPGroup}
}\hfill{\scriptsize (attribute)}}\\


 For a $p$\texttt{\symbol{45}}group \mbox{\texttt{\mdseries\slshape G}} (see{\nobreakspace}\texttt{IsPGroup} (\ref{IsPGroup})), \texttt{RankPGroup} returns the \emph{rank} of \mbox{\texttt{\mdseries\slshape G}}, which is defined as the minimal size of a generating system of \mbox{\texttt{\mdseries\slshape G}}. If \mbox{\texttt{\mdseries\slshape G}} is not a $p$\texttt{\symbol{45}}group then an error is issued. 
\begin{Verbatim}[commandchars=!@|,fontsize=\small,frame=single,label=Example]
  !gapprompt@gap>| !gapinput@h:= Group((1,2,3,4), (1,3));;|
  !gapprompt@gap>| !gapinput@Length(PCentralSeries(h));|
  3
  !gapprompt@gap>| !gapinput@PClassPGroup(h);|
  2
  !gapprompt@gap>| !gapinput@RankPGroup(h);|
  2
\end{Verbatim}
 }

 

\subsection{\textcolor{Chapter }{IsPSolvable}}
\logpage{[ 39, 15, 26 ]}\nobreak
\hyperdef{L}{X81130F9A7CFCF6BF}{}
{\noindent\textcolor{FuncColor}{$\triangleright$\enspace\texttt{IsPSolvable({\mdseries\slshape G, p})\index{IsPSolvable@\texttt{IsPSolvable}}
\label{IsPSolvable}
}\hfill{\scriptsize (operation)}}\\


 A finite group is $p$\texttt{\symbol{45}}solvable if every chief factor either has order not
divisible by $p$, or is solvable. 

 }

 

\subsection{\textcolor{Chapter }{IsPNilpotent}}
\logpage{[ 39, 15, 27 ]}\nobreak
\hyperdef{L}{X87415A8485FCF510}{}
{\noindent\textcolor{FuncColor}{$\triangleright$\enspace\texttt{IsPNilpotent({\mdseries\slshape G, p})\index{IsPNilpotent@\texttt{IsPNilpotent}}
\label{IsPNilpotent}
}\hfill{\scriptsize (operation)}}\\


 A group is $p$\texttt{\symbol{45}}nilpotent if it possesses a normal $p$\texttt{\symbol{45}}complement. }

 }

  
\section{\textcolor{Chapter }{Numerical Group Attributes}}\label{Numerical Group Attributes}
\logpage{[ 39, 16, 0 ]}
\hyperdef{L}{X7F8264FA796B2B7D}{}
{
  This section gives only some examples of numerical group attributes, so it
should not serve as a collection of all numerical group attributes. The manual
contains more such attributes documented in this manual, for example, \texttt{NrConjugacyClasses} (\ref{NrConjugacyClasses}), \texttt{NilpotencyClassOfGroup} (\ref{NilpotencyClassOfGroup}) and others. 

 Note also that some functions, such as \texttt{EulerianFunction} (\ref{EulerianFunction}), are mathematical attributes, but not \textsf{GAP} attributes (see \ref{Attributes}) as they are depending on a parameter. 

\subsection{\textcolor{Chapter }{AbelianInvariants}}
\logpage{[ 39, 16, 1 ]}\nobreak
\hyperdef{L}{X812827937F403300}{}
{\noindent\textcolor{FuncColor}{$\triangleright$\enspace\texttt{AbelianInvariants({\mdseries\slshape G})\index{AbelianInvariants@\texttt{AbelianInvariants}}
\label{AbelianInvariants}
}\hfill{\scriptsize (attribute)}}\\


 \index{AbelianInvariants@\texttt{AbelianInvariants}!for groups} returns the abelian invariants (also sometimes called primary decomposition)
of the commutator factor group of the group \mbox{\texttt{\mdseries\slshape G}}. These are given as a list of prime\texttt{\symbol{45}}powers or zeroes and
describe the structure of $\mbox{\texttt{\mdseries\slshape G}}/\mbox{\texttt{\mdseries\slshape G}}'$ as a direct product of cyclic groups of prime power (or infinite) order. 

 (See \texttt{IndependentGeneratorsOfAbelianGroup} (\ref{IndependentGeneratorsOfAbelianGroup}) to obtain actual generators). 
\begin{Verbatim}[commandchars=!@|,fontsize=\small,frame=single,label=Example]
  !gapprompt@gap>| !gapinput@g:=Group((1,2,3,4),(1,2),(5,6));;|
  !gapprompt@gap>| !gapinput@AbelianInvariants(g);|
  [ 2, 2 ]
  !gapprompt@gap>| !gapinput@h:=FreeGroup(2);;h:=h/[h.1^3];;|
  !gapprompt@gap>| !gapinput@AbelianInvariants(h);|
  [ 0, 3 ]
\end{Verbatim}
 }

 

\subsection{\textcolor{Chapter }{Exponent}}
\logpage{[ 39, 16, 2 ]}\nobreak
\hyperdef{L}{X7D44470C7DA59C1C}{}
{\noindent\textcolor{FuncColor}{$\triangleright$\enspace\texttt{Exponent({\mdseries\slshape G})\index{Exponent@\texttt{Exponent}}
\label{Exponent}
}\hfill{\scriptsize (attribute)}}\\


 The exponent $e$ of a group \mbox{\texttt{\mdseries\slshape G}} is the lcm of the orders of its elements, that is, $e$ is the smallest integer such that $g^e = 1$ for all $g \in \mbox{\texttt{\mdseries\slshape G}}$. 
\begin{Verbatim}[commandchars=!@|,fontsize=\small,frame=single,label=Example]
  !gapprompt@gap>| !gapinput@Exponent(g);|
  12
\end{Verbatim}
 }

 

\subsection{\textcolor{Chapter }{EulerianFunction}}
\logpage{[ 39, 16, 3 ]}\nobreak
\hyperdef{L}{X843E0CCA8351FDF4}{}
{\noindent\textcolor{FuncColor}{$\triangleright$\enspace\texttt{EulerianFunction({\mdseries\slshape G, n})\index{EulerianFunction@\texttt{EulerianFunction}}
\label{EulerianFunction}
}\hfill{\scriptsize (operation)}}\\


 returns the number of \mbox{\texttt{\mdseries\slshape n}}\texttt{\symbol{45}}tuples $(g_1, g_2, \ldots, g_n)$ of elements of the group \mbox{\texttt{\mdseries\slshape G}} that generate the whole group \mbox{\texttt{\mdseries\slshape G}}. The elements of such an \mbox{\texttt{\mdseries\slshape n}}\texttt{\symbol{45}}tuple need not be different. 

 In \cite{Hal36}, the notation $\phi_{\mbox{\texttt{\mdseries\slshape n}}}(\mbox{\texttt{\mdseries\slshape G}})$ is used for the value returned by \texttt{EulerianFunction}, and the quotient of $\phi_{\mbox{\texttt{\mdseries\slshape n}}}(\mbox{\texttt{\mdseries\slshape G}})$ by the order of the automorphism group of \mbox{\texttt{\mdseries\slshape G}} is called $d_{\mbox{\texttt{\mdseries\slshape n}}}(\mbox{\texttt{\mdseries\slshape G}})$. If \mbox{\texttt{\mdseries\slshape G}} is a nonabelian simple group then $d_{\mbox{\texttt{\mdseries\slshape n}}}(\mbox{\texttt{\mdseries\slshape G}})$ is the greatest number $d$ for which the direct product of $d$ groups isomorphic with \mbox{\texttt{\mdseries\slshape G}} can be generated by \mbox{\texttt{\mdseries\slshape n}} elements. 

 If the Library of Tables of Marks (see Chapter \ref{Tables of Marks}) covers the group \mbox{\texttt{\mdseries\slshape G}}, you may also use \texttt{EulerianFunctionByTom} (\ref{EulerianFunctionByTom}). 

 
\begin{Verbatim}[commandchars=!@|,fontsize=\small,frame=single,label=Example]
  !gapprompt@gap>| !gapinput@EulerianFunction( g, 2 );|
  432
\end{Verbatim}
 }

 }

  
\section{\textcolor{Chapter }{Subgroup Series}}\label{Subgroup Series}
\logpage{[ 39, 17, 0 ]}
\hyperdef{L}{X7AEDEDF67CFED672}{}
{
  In group theory many subgroup series are considered, and \textsf{GAP} provides commands to compute them. In the following sections, there is always
a series $G = U_1 > U_2 > \cdots > U_m = \langle 1 \rangle$ of subgroups considered. A series also may stop without reaching $G$ or $\langle 1 \rangle$. 

 A series is called \emph{subnormal} if every $U_{{i+1}}$ is normal in $U_i$. 

 A series is called \emph{normal} if every $U_i$ is normal in $G$. 

 A series of normal subgroups is called \emph{central} if $U_i/U_{{i+1}}$ is central in $G / U_{{i+1}}$. 

 We call a series \emph{refinable} if intermediate subgroups can be added to the series without destroying the
properties of the series. 

 Unless explicitly declared otherwise, all subgroup series are descending. That
is they are stored in decreasing order. 

\subsection{\textcolor{Chapter }{ChiefSeries}}
\logpage{[ 39, 17, 1 ]}\nobreak
\hyperdef{L}{X7BDD116F7833800F}{}
{\noindent\textcolor{FuncColor}{$\triangleright$\enspace\texttt{ChiefSeries({\mdseries\slshape G})\index{ChiefSeries@\texttt{ChiefSeries}}
\label{ChiefSeries}
}\hfill{\scriptsize (attribute)}}\\


 is a series of normal subgroups of \mbox{\texttt{\mdseries\slshape G}} which cannot be refined further. That is there is no normal subgroup $N$ of \mbox{\texttt{\mdseries\slshape G}} with $U_i > N > U_{{i+1}}$. This attribute returns \emph{one} chief series (of potentially many possibilities). 
\begin{Verbatim}[commandchars=!@|,fontsize=\small,frame=single,label=Example]
  !gapprompt@gap>| !gapinput@g:=Group((1,2,3,4),(1,2));;|
  !gapprompt@gap>| !gapinput@ChiefSeries(g);|
  [ Group([ (1,2,3,4), (1,2) ]),
    Group([ (2,4,3), (1,4)(2,3), (1,3)(2,4) ]),
    Group([ (1,4)(2,3), (1,3)(2,4) ]), Group(()) ]
\end{Verbatim}
 }

 

\subsection{\textcolor{Chapter }{ChiefSeriesThrough}}
\logpage{[ 39, 17, 2 ]}\nobreak
\hyperdef{L}{X7AC93E977AC9ED58}{}
{\noindent\textcolor{FuncColor}{$\triangleright$\enspace\texttt{ChiefSeriesThrough({\mdseries\slshape G, l})\index{ChiefSeriesThrough@\texttt{ChiefSeriesThrough}}
\label{ChiefSeriesThrough}
}\hfill{\scriptsize (operation)}}\\


 is a chief series of the group \mbox{\texttt{\mdseries\slshape G}} going through the normal subgroups in the list \mbox{\texttt{\mdseries\slshape l}}, which must be a list of normal subgroups of \mbox{\texttt{\mdseries\slshape G}} contained in each other, sorted by descending size. This attribute returns \emph{one} chief series (of potentially many possibilities). }

 

\subsection{\textcolor{Chapter }{ChiefSeriesUnderAction}}
\logpage{[ 39, 17, 3 ]}\nobreak
\hyperdef{L}{X8724E15F81B51173}{}
{\noindent\textcolor{FuncColor}{$\triangleright$\enspace\texttt{ChiefSeriesUnderAction({\mdseries\slshape H, G})\index{ChiefSeriesUnderAction@\texttt{ChiefSeriesUnderAction}}
\label{ChiefSeriesUnderAction}
}\hfill{\scriptsize (operation)}}\\


 returns a series of normal subgroups of \mbox{\texttt{\mdseries\slshape G}} which are invariant under \mbox{\texttt{\mdseries\slshape H}} such that the series cannot be refined any further. \mbox{\texttt{\mdseries\slshape G}} must be a subgroup of \mbox{\texttt{\mdseries\slshape H}}. This attribute returns \emph{one} such series (of potentially many possibilities). }

 

\subsection{\textcolor{Chapter }{SubnormalSeries}}
\logpage{[ 39, 17, 4 ]}\nobreak
\hyperdef{L}{X7A0E7A8B8495B79D}{}
{\noindent\textcolor{FuncColor}{$\triangleright$\enspace\texttt{SubnormalSeries({\mdseries\slshape G, U})\index{SubnormalSeries@\texttt{SubnormalSeries}}
\label{SubnormalSeries}
}\hfill{\scriptsize (operation)}}\\


 If \mbox{\texttt{\mdseries\slshape U}} is a subgroup of \mbox{\texttt{\mdseries\slshape G}} this operation returns a subnormal series that descends from \mbox{\texttt{\mdseries\slshape G}} to a subnormal subgroup $V \geq $\mbox{\texttt{\mdseries\slshape U}}. If \mbox{\texttt{\mdseries\slshape U}} is subnormal, $V =$ \mbox{\texttt{\mdseries\slshape U}}. 
\begin{Verbatim}[commandchars=!@|,fontsize=\small,frame=single,label=Example]
  !gapprompt@gap>| !gapinput@s:=SubnormalSeries(g,Group((1,2)(3,4))) =|
  !gapprompt@>| !gapinput@[ Group([ (1,2,3,4), (1,2) ]),|
  !gapprompt@>| !gapinput@  Group([ (1,2)(3,4), (1,3)(2,4) ]),|
  !gapprompt@>| !gapinput@  Group([ (1,2)(3,4) ]) ];|
  true
\end{Verbatim}
 }

 

\subsection{\textcolor{Chapter }{CompositionSeries}}
\logpage{[ 39, 17, 5 ]}\nobreak
\hyperdef{L}{X81CDCBD67BC98A5A}{}
{\noindent\textcolor{FuncColor}{$\triangleright$\enspace\texttt{CompositionSeries({\mdseries\slshape G})\index{CompositionSeries@\texttt{CompositionSeries}}
\label{CompositionSeries}
}\hfill{\scriptsize (attribute)}}\\
\noindent\textcolor{FuncColor}{$\triangleright$\enspace\texttt{CompositionSeriesThrough({\mdseries\slshape G, normals})\index{CompositionSeriesThrough@\texttt{CompositionSeriesThrough}}
\label{CompositionSeriesThrough}
}\hfill{\scriptsize (operation)}}\\


 A composition series is a subnormal series which cannot be refined. This
attribute returns \emph{one} composition series (of potentially many possibilities). The variant \texttt{CompositionSeriesThrough} takes as second argument a list \mbox{\texttt{\mdseries\slshape normals}} of normal subgroups of the group, and returns a composition series that
incorporates these normal subgroups. }

 

\subsection{\textcolor{Chapter }{DisplayCompositionSeries}}
\logpage{[ 39, 17, 6 ]}\nobreak
\hyperdef{L}{X82C0D0217ACB2042}{}
{\noindent\textcolor{FuncColor}{$\triangleright$\enspace\texttt{DisplayCompositionSeries({\mdseries\slshape G})\index{DisplayCompositionSeries@\texttt{DisplayCompositionSeries}}
\label{DisplayCompositionSeries}
}\hfill{\scriptsize (function)}}\\


 Displays a composition series of \mbox{\texttt{\mdseries\slshape G}} in a nice way, identifying the simple factors. 
\begin{Verbatim}[commandchars=!@B,fontsize=\small,frame=single,label=Example]
  !gapprompt@gap>B !gapinput@CompositionSeries(g);B
  [ Group([ (3,4), (2,4,3), (1,4)(2,3), (1,3)(2,4) ]),
    Group([ (2,4,3), (1,4)(2,3), (1,3)(2,4) ]),
    Group([ (1,4)(2,3), (1,3)(2,4) ]), Group([ (1,3)(2,4) ]), Group(())
   ]
  !gapprompt@gap>B !gapinput@DisplayCompositionSeries(Group((1,2,3,4,5,6,7),(1,2)));B
  G (2 gens, size 5040)
   | C2
  S (5 gens, size 2520)
   | A7
  1 (0 gens, size 1)
\end{Verbatim}
 }

 

\subsection{\textcolor{Chapter }{DerivedSeriesOfGroup}}
\logpage{[ 39, 17, 7 ]}\nobreak
\hyperdef{L}{X7A879948834BD889}{}
{\noindent\textcolor{FuncColor}{$\triangleright$\enspace\texttt{DerivedSeriesOfGroup({\mdseries\slshape G})\index{DerivedSeriesOfGroup@\texttt{DerivedSeriesOfGroup}}
\label{DerivedSeriesOfGroup}
}\hfill{\scriptsize (attribute)}}\\


 The derived series of a group is obtained by $U_{{i+1}} = U_i'$. It stops if $U_i$ is perfect. }

 

\subsection{\textcolor{Chapter }{DerivedLength}}
\logpage{[ 39, 17, 8 ]}\nobreak
\hyperdef{L}{X7A9AA1577CEC891F}{}
{\noindent\textcolor{FuncColor}{$\triangleright$\enspace\texttt{DerivedLength({\mdseries\slshape G})\index{DerivedLength@\texttt{DerivedLength}}
\label{DerivedLength}
}\hfill{\scriptsize (attribute)}}\\


 The derived length of a group is the number of steps in the derived series.
(As there is always the group, it is the series length minus 1.) 
\begin{Verbatim}[commandchars=!@|,fontsize=\small,frame=single,label=Example]
  !gapprompt@gap>| !gapinput@List(DerivedSeriesOfGroup(g),Size);|
  [ 24, 12, 4, 1 ]
  !gapprompt@gap>| !gapinput@DerivedLength(g);|
  3
\end{Verbatim}
 }

 
\subsection{\textcolor{Chapter }{ElementaryAbelianSeries}}\logpage{[ 39, 17, 9 ]}
\hyperdef{L}{X83F057E5791944D6}{}
{
\noindent\textcolor{FuncColor}{$\triangleright$\enspace\texttt{ElementaryAbelianSeries({\mdseries\slshape G})\index{ElementaryAbelianSeries@\texttt{ElementaryAbelianSeries}!for a group}
\label{ElementaryAbelianSeries:for a group}
}\hfill{\scriptsize (attribute)}}\\
\noindent\textcolor{FuncColor}{$\triangleright$\enspace\texttt{ElementaryAbelianSeriesLargeSteps({\mdseries\slshape G})\index{ElementaryAbelianSeriesLargeSteps@\texttt{ElementaryAbelianSeriesLargeSteps}}
\label{ElementaryAbelianSeriesLargeSteps}
}\hfill{\scriptsize (attribute)}}\\
\noindent\textcolor{FuncColor}{$\triangleright$\enspace\texttt{ElementaryAbelianSeries({\mdseries\slshape list})\index{ElementaryAbelianSeries@\texttt{ElementaryAbelianSeries}!for a list}
\label{ElementaryAbelianSeries:for a list}
}\hfill{\scriptsize (attribute)}}\\


 returns a series of normal subgroups of $G$ such that all factors are elementary abelian. If the group is not solvable
(and thus no such series exists) it returns \texttt{fail}. 

 The variant \texttt{ElementaryAbelianSeriesLargeSteps} tries to make the steps in this series large (by eliminating intermediate
subgroups if possible) at a small additional cost. 

 In the third variant, an elementary abelian series through the given series of
normal subgroups in the list \mbox{\texttt{\mdseries\slshape list}} is constructed. 
\begin{Verbatim}[commandchars=!@|,fontsize=\small,frame=single,label=Example]
  !gapprompt@gap>| !gapinput@List(ElementaryAbelianSeries(g),Size);|
  [ 24, 12, 4, 1 ]
\end{Verbatim}
 }

 

\subsection{\textcolor{Chapter }{InvariantElementaryAbelianSeries}}
\logpage{[ 39, 17, 10 ]}\nobreak
\hyperdef{L}{X782BD7A47D6B6503}{}
{\noindent\textcolor{FuncColor}{$\triangleright$\enspace\texttt{InvariantElementaryAbelianSeries({\mdseries\slshape G, morph[, N[, fine]]})\index{InvariantElementaryAbelianSeries@\texttt{InvariantElementaryAbelianSeries}}
\label{InvariantElementaryAbelianSeries}
}\hfill{\scriptsize (function)}}\\


 For a (solvable) group \mbox{\texttt{\mdseries\slshape G}} and a list of automorphisms \mbox{\texttt{\mdseries\slshape morph}} of \mbox{\texttt{\mdseries\slshape G}}, this command finds a normal series of \mbox{\texttt{\mdseries\slshape G}} with elementary abelian factors such that every group in this series is
invariant under every automorphism in \mbox{\texttt{\mdseries\slshape morph}}. 

 If a normal subgroup \mbox{\texttt{\mdseries\slshape N}} of \mbox{\texttt{\mdseries\slshape G}} which is invariant under \mbox{\texttt{\mdseries\slshape morph}} is given, this series is chosen to contain \mbox{\texttt{\mdseries\slshape N}}. No tests are performed to check the validity of the arguments. 

 The series obtained will be constructed to prefer large steps unless \mbox{\texttt{\mdseries\slshape fine}} is given as \texttt{true}. 
\begin{Verbatim}[commandchars=!@|,fontsize=\small,frame=single,label=Example]
  !gapprompt@gap>| !gapinput@g:=Group((1,2,3,4),(1,3));|
  Group([ (1,2,3,4), (1,3) ])
  !gapprompt@gap>| !gapinput@hom:=GroupHomomorphismByImages(g,g,GeneratorsOfGroup(g),|
  !gapprompt@>| !gapinput@[(1,4,3,2),(1,4)(2,3)]);|
  [ (1,2,3,4), (1,3) ] -> [ (1,4,3,2), (1,4)(2,3) ]
  !gapprompt@gap>| !gapinput@InvariantElementaryAbelianSeries(g,[hom]);|
  [ Group([ (1,2,3,4), (1,3) ]), Group([ (1,3)(2,4) ]), Group(()) ]
\end{Verbatim}
 }

 

\subsection{\textcolor{Chapter }{LowerCentralSeriesOfGroup}}
\logpage{[ 39, 17, 11 ]}\nobreak
\hyperdef{L}{X879D55A67DB42676}{}
{\noindent\textcolor{FuncColor}{$\triangleright$\enspace\texttt{LowerCentralSeriesOfGroup({\mdseries\slshape G})\index{LowerCentralSeriesOfGroup@\texttt{LowerCentralSeriesOfGroup}}
\label{LowerCentralSeriesOfGroup}
}\hfill{\scriptsize (attribute)}}\\


 The lower central series of a group \mbox{\texttt{\mdseries\slshape G}} is defined as $U_{{i+1}}:= [\mbox{\texttt{\mdseries\slshape G}}, U_i]$. It is a central series of normal subgroups. The name derives from the fact
that $U_i$ is contained in the $i$\texttt{\symbol{45}}th step subgroup of any central series. }

 

\subsection{\textcolor{Chapter }{UpperCentralSeriesOfGroup}}
\logpage{[ 39, 17, 12 ]}\nobreak
\hyperdef{L}{X8428592E8773CD7B}{}
{\noindent\textcolor{FuncColor}{$\triangleright$\enspace\texttt{UpperCentralSeriesOfGroup({\mdseries\slshape G})\index{UpperCentralSeriesOfGroup@\texttt{UpperCentralSeriesOfGroup}}
\label{UpperCentralSeriesOfGroup}
}\hfill{\scriptsize (attribute)}}\\


 The upper central series of a group \mbox{\texttt{\mdseries\slshape G}} is defined as an ending series $U_i / U_{{i+1}}:= Z(\mbox{\texttt{\mdseries\slshape G}}/U_{{i+1}})$. It is a central series of normal subgroups. The name derives from the fact
that $U_i$ contains every $i$\texttt{\symbol{45}}th step subgroup of a central series. }

 

\subsection{\textcolor{Chapter }{PCentralSeries}}
\logpage{[ 39, 17, 13 ]}\nobreak
\hyperdef{L}{X7809B7ED792669F3}{}
{\noindent\textcolor{FuncColor}{$\triangleright$\enspace\texttt{PCentralSeries({\mdseries\slshape G, p})\index{PCentralSeries@\texttt{PCentralSeries}}
\label{PCentralSeries}
}\hfill{\scriptsize (operation)}}\\


 The \mbox{\texttt{\mdseries\slshape p}}\texttt{\symbol{45}}central series of \mbox{\texttt{\mdseries\slshape G}} is defined by $U_1:= \mbox{\texttt{\mdseries\slshape G}}$, $U_i:= [\mbox{\texttt{\mdseries\slshape G}}, U_{{i-1}}] U_{{i-1}}^{\mbox{\texttt{\mdseries\slshape p}}}$. 
\begin{Verbatim}[commandchars=!@|,fontsize=\small,frame=single,label=Example]
  !gapprompt@gap>| !gapinput@g:=DicyclicGroup(12);;|
  !gapprompt@gap>| !gapinput@PCentralSeries(g,2);|
  [ <pc group of size 12 with 3 generators>, Group([ y3, y*y3 ]), Group([ y*y3 ]) ]
  !gapprompt@gap>| !gapinput@g:=SymmetricGroup(4);;|
  !gapprompt@gap>| !gapinput@List(PCentralSeries(g,2), StructureDescription);|
  [ "S4", "A4" ]
\end{Verbatim}
 }

 

\subsection{\textcolor{Chapter }{JenningsSeries}}
\logpage{[ 39, 17, 14 ]}\nobreak
\hyperdef{L}{X82A34BD681F24A94}{}
{\noindent\textcolor{FuncColor}{$\triangleright$\enspace\texttt{JenningsSeries({\mdseries\slshape G})\index{JenningsSeries@\texttt{JenningsSeries}}
\label{JenningsSeries}
}\hfill{\scriptsize (attribute)}}\\


 For a $p$\texttt{\symbol{45}}group \mbox{\texttt{\mdseries\slshape G}}, this function returns its Jennings series. This series is defined by setting $G_1 = \mbox{\texttt{\mdseries\slshape G}}$ and for $i \geq 0$, $G_{{i+1}} = [G_i,\mbox{\texttt{\mdseries\slshape G}}] G_j^p$, where $j$ is the smallest integer $> i/p$. }

 

\subsection{\textcolor{Chapter }{DimensionsLoewyFactors}}
\logpage{[ 39, 17, 15 ]}\nobreak
\hyperdef{L}{X7C08A8B77EC09CFF}{}
{\noindent\textcolor{FuncColor}{$\triangleright$\enspace\texttt{DimensionsLoewyFactors({\mdseries\slshape G})\index{DimensionsLoewyFactors@\texttt{DimensionsLoewyFactors}}
\label{DimensionsLoewyFactors}
}\hfill{\scriptsize (attribute)}}\\


 This operation computes the dimensions of the factors of the Loewy series of \mbox{\texttt{\mdseries\slshape G}}. (See \cite[p. 157]{Hup82} for the slightly complicated definition of the Loewy Series.) 

 The dimensions are computed via the \texttt{JenningsSeries} (\ref{JenningsSeries}) without computing the Loewy series itself. 
\begin{Verbatim}[commandchars=!@|,fontsize=\small,frame=single,label=Example]
  !gapprompt@gap>| !gapinput@G:= SmallGroup( 3^6, 100 );|
  <pc group of size 729 with 6 generators>
  !gapprompt@gap>| !gapinput@JenningsSeries( G );|
  [ <pc group of size 729 with 6 generators>, Group([ f3, f4, f5, f6 ]),
    Group([ f4, f5, f6 ]), Group([ f5, f6 ]), Group([ f5, f6 ]),
    Group([ f5, f6 ]), Group([ f6 ]), Group([ f6 ]), Group([ f6 ]),
    Group([ <identity> of ... ]) ]
  !gapprompt@gap>| !gapinput@DimensionsLoewyFactors(G);|
  [ 1, 2, 4, 5, 7, 8, 10, 11, 13, 14, 16, 17, 19, 20, 22, 23, 25, 26,
    27, 27, 27, 27, 27, 27, 27, 27, 27, 26, 25, 23, 22, 20, 19, 17, 16,
    14, 13, 11, 10, 8, 7, 5, 4, 2, 1 ]
\end{Verbatim}
 }

 

\subsection{\textcolor{Chapter }{AscendingChain}}
\logpage{[ 39, 17, 16 ]}\nobreak
\hyperdef{L}{X84112774812180DD}{}
{\noindent\textcolor{FuncColor}{$\triangleright$\enspace\texttt{AscendingChain({\mdseries\slshape G, U})\index{AscendingChain@\texttt{AscendingChain}}
\label{AscendingChain}
}\hfill{\scriptsize (function)}}\\


 This function computes an ascending chain of subgroups from \mbox{\texttt{\mdseries\slshape U}} to \mbox{\texttt{\mdseries\slshape G}}. This chain is given as a list whose first entry is \mbox{\texttt{\mdseries\slshape U}} and the last entry is \mbox{\texttt{\mdseries\slshape G}}. The function tries to make the links in this chain small. 

 The option \texttt{refineIndex} can be used to give a bound for refinements of steps to avoid \textsf{GAP} trying to enforce too small steps. The option \texttt{cheap} (if set to \texttt{true}) will overall limit the amount of heuristic searches. }

 

\subsection{\textcolor{Chapter }{IntermediateGroup}}
\logpage{[ 39, 17, 17 ]}\nobreak
\hyperdef{L}{X7C5029EE86D7FC96}{}
{\noindent\textcolor{FuncColor}{$\triangleright$\enspace\texttt{IntermediateGroup({\mdseries\slshape G, U})\index{IntermediateGroup@\texttt{IntermediateGroup}}
\label{IntermediateGroup}
}\hfill{\scriptsize (function)}}\\


 This routine tries to find a subgroup $E$ of \mbox{\texttt{\mdseries\slshape G}}, such that $\mbox{\texttt{\mdseries\slshape G}} > E > \mbox{\texttt{\mdseries\slshape U}}$ holds. If \mbox{\texttt{\mdseries\slshape U}} is maximal in \mbox{\texttt{\mdseries\slshape G}}, the function returns \texttt{fail}. This is done by finding minimal blocks for the operation of \mbox{\texttt{\mdseries\slshape G}} on the right cosets of \mbox{\texttt{\mdseries\slshape U}}. }

 

\subsection{\textcolor{Chapter }{IntermediateSubgroups}}
\logpage{[ 39, 17, 18 ]}\nobreak
\hyperdef{L}{X781661FB78DC83B5}{}
{\noindent\textcolor{FuncColor}{$\triangleright$\enspace\texttt{IntermediateSubgroups({\mdseries\slshape G, U})\index{IntermediateSubgroups@\texttt{IntermediateSubgroups}}
\label{IntermediateSubgroups}
}\hfill{\scriptsize (operation)}}\\


 returns a list of all subgroups of \mbox{\texttt{\mdseries\slshape G}} that properly contain \mbox{\texttt{\mdseries\slshape U}}; that is all subgroups between \mbox{\texttt{\mdseries\slshape G}} and \mbox{\texttt{\mdseries\slshape U}}. It returns a record with a component \texttt{subgroups}, which is a list of these subgroups, as well as a component \texttt{inclusions}, which lists all maximality inclusions among these subgroups. A maximality
inclusion is given as a list $[i, j]$ indicating that the subgroup number $i$ is a maximal subgroup of the subgroup number $j$, the numbers $0$ and $1 +$ \texttt{Length(subgroups)} are used to denote \mbox{\texttt{\mdseries\slshape U}} and \mbox{\texttt{\mdseries\slshape G}}, respectively. }

 

\subsection{\textcolor{Chapter }{StructuralSeriesOfGroup}}
\logpage{[ 39, 17, 19 ]}\nobreak
\hyperdef{L}{X783CDAA67BDD8195}{}
{\noindent\textcolor{FuncColor}{$\triangleright$\enspace\texttt{StructuralSeriesOfGroup({\mdseries\slshape G})\index{StructuralSeriesOfGroup@\texttt{StructuralSeriesOfGroup}}
\label{StructuralSeriesOfGroup}
}\hfill{\scriptsize (attribute)}}\\


 The structural series of a finite group \mbox{\texttt{\mdseries\slshape G}} is a descending series of characteristic subgroups which goes through the
derived series of the solvable radical of \$G\$, refined into elementary
abelian factors, as well as the socle and the \mbox{\texttt{\mdseries\slshape Pker}} (kernel of the action on socle components) of the radical factor 
\begin{Verbatim}[commandchars=!@|,fontsize=\small,frame=single,label=Example]
  !gapprompt@gap>| !gapinput@gp:=WreathProduct(SymmetricGroup(5),SymmetricGroup(3));;|
  !gapprompt@gap>| !gapinput@gp:=WreathProduct(Group((1,2,3,4,5,6)),gp);;|
  !gapprompt@gap>| !gapinput@List(StructuralSeriesOfGroup(gp),Size);|
  [ 4874877920083968000, 812479653347328000, 101559956668416000, 470184984576, 32768, 1 ]
\end{Verbatim}
 }

 }

  
\section{\textcolor{Chapter }{Factor Groups}}\label{Factor Groups}
\logpage{[ 39, 18, 0 ]}
\hyperdef{L}{X84091B0A7E401E2B}{}
{
  

\subsection{\textcolor{Chapter }{NaturalHomomorphismByNormalSubgroup}}
\logpage{[ 39, 18, 1 ]}\nobreak
\hyperdef{L}{X80FC390C7F38A13F}{}
{\noindent\textcolor{FuncColor}{$\triangleright$\enspace\texttt{NaturalHomomorphismByNormalSubgroup({\mdseries\slshape G, N})\index{NaturalHomomorphismByNormalSubgroup@\texttt{NaturalHomomorphismByNormalSubgroup}}
\label{NaturalHomomorphismByNormalSubgroup}
}\hfill{\scriptsize (function)}}\\
\noindent\textcolor{FuncColor}{$\triangleright$\enspace\texttt{NaturalHomomorphismByNormalSubgroupNC({\mdseries\slshape G, N})\index{NaturalHomomorphismByNormalSubgroupNC@\texttt{Natural}\-\texttt{Homomorphism}\-\texttt{By}\-\texttt{Normal}\-\texttt{SubgroupNC}}
\label{NaturalHomomorphismByNormalSubgroupNC}
}\hfill{\scriptsize (function)}}\\


 returns a homomorphism from \mbox{\texttt{\mdseries\slshape G}} to another group whose kernel is \mbox{\texttt{\mdseries\slshape N}}. \textsf{GAP} will try to select the image group as to make computations in it as efficient
as possible. As the factor group $\mbox{\texttt{\mdseries\slshape G}}/\mbox{\texttt{\mdseries\slshape N}}$ can be identified with the image of \mbox{\texttt{\mdseries\slshape G}} this permits efficient computations in the factor group. The homomorphism
returned is not necessarily surjective, so \texttt{ImagesSource} (\ref{ImagesSource}) should be used instead of \texttt{Range} (\ref{Range:of a general mapping}) to get a group isomorphic to the factor group. The \texttt{NC} variant does not check whether \mbox{\texttt{\mdseries\slshape N}} is normal in \mbox{\texttt{\mdseries\slshape G}}. }

 

\subsection{\textcolor{Chapter }{FactorGroup}}
\logpage{[ 39, 18, 2 ]}\nobreak
\hyperdef{L}{X7E6EED0185B27C48}{}
{\noindent\textcolor{FuncColor}{$\triangleright$\enspace\texttt{FactorGroup({\mdseries\slshape G, N})\index{FactorGroup@\texttt{FactorGroup}}
\label{FactorGroup}
}\hfill{\scriptsize (function)}}\\
\noindent\textcolor{FuncColor}{$\triangleright$\enspace\texttt{FactorGroupNC({\mdseries\slshape G, N})\index{FactorGroupNC@\texttt{FactorGroupNC}}
\label{FactorGroupNC}
}\hfill{\scriptsize (operation)}}\\


 returns the image of the \texttt{NaturalHomomorphismByNormalSubgroup(\mbox{\texttt{\mdseries\slshape G}},\mbox{\texttt{\mdseries\slshape N}})}. This function is provided for compatibility with older code, but if a
connection between group and factor is desired, users need to start by
obtaining the \texttt{NaturalHomomorphismByNormalSubgroup} in the first place. The \texttt{NC} version does not test whether \mbox{\texttt{\mdseries\slshape N}} is normal in \mbox{\texttt{\mdseries\slshape G}}. 
\begin{Verbatim}[commandchars=!@|,fontsize=\small,frame=single,label=Example]
  !gapprompt@gap>| !gapinput@g:=Group((1,2,3,4),(1,2));;n:=Subgroup(g,[(1,2)(3,4),(1,3)(2,4)]);;|
  !gapprompt@gap>| !gapinput@hom:=NaturalHomomorphismByNormalSubgroup(g,n);|
  [ (1,2,3,4), (1,2) ] -> [ f1*f2, f1 ]
  !gapprompt@gap>| !gapinput@Size(ImagesSource(hom));|
  6
  !gapprompt@gap>| !gapinput@StructureDescription(Image(hom,g));|
  "S3"
\end{Verbatim}
 }

 

\subsection{\textcolor{Chapter }{CommutatorFactorGroup}}
\logpage{[ 39, 18, 3 ]}\nobreak
\hyperdef{L}{X7816FA867BF1B8ED}{}
{\noindent\textcolor{FuncColor}{$\triangleright$\enspace\texttt{CommutatorFactorGroup({\mdseries\slshape G})\index{CommutatorFactorGroup@\texttt{CommutatorFactorGroup}}
\label{CommutatorFactorGroup}
}\hfill{\scriptsize (attribute)}}\\


 computes the commutator factor group $\mbox{\texttt{\mdseries\slshape G}}/\mbox{\texttt{\mdseries\slshape G}}'$ of the group \mbox{\texttt{\mdseries\slshape G}}. 
\begin{Verbatim}[commandchars=!@|,fontsize=\small,frame=single,label=Example]
  !gapprompt@gap>| !gapinput@CommutatorFactorGroup(g);|
  Group([ f1 ])
\end{Verbatim}
 }

 

\subsection{\textcolor{Chapter }{MaximalAbelianQuotient}}
\logpage{[ 39, 18, 4 ]}\nobreak
\hyperdef{L}{X7BB93B9778C5A0B2}{}
{\noindent\textcolor{FuncColor}{$\triangleright$\enspace\texttt{MaximalAbelianQuotient({\mdseries\slshape G})\index{MaximalAbelianQuotient@\texttt{MaximalAbelianQuotient}}
\label{MaximalAbelianQuotient}
}\hfill{\scriptsize (attribute)}}\\


 returns an epimorphism from \mbox{\texttt{\mdseries\slshape G}} onto the maximal abelian quotient of \mbox{\texttt{\mdseries\slshape G}}. The kernel of this epimorphism is the derived subgroup of \mbox{\texttt{\mdseries\slshape G}}, see \texttt{DerivedSubgroup} (\ref{DerivedSubgroup}). }

 

\subsection{\textcolor{Chapter }{HasAbelianFactorGroup}}
\logpage{[ 39, 18, 5 ]}\nobreak
\hyperdef{L}{X7FC83E4C783572E7}{}
{\noindent\textcolor{FuncColor}{$\triangleright$\enspace\texttt{HasAbelianFactorGroup({\mdseries\slshape G, N})\index{HasAbelianFactorGroup@\texttt{HasAbelianFactorGroup}}
\label{HasAbelianFactorGroup}
}\hfill{\scriptsize (function)}}\\


 tests whether \mbox{\texttt{\mdseries\slshape G}} $/$ \mbox{\texttt{\mdseries\slshape N}} is abelian (without explicitly constructing the factor group and without
testing whether \mbox{\texttt{\mdseries\slshape N}} is in fact a normal subgroup). 
\begin{Verbatim}[commandchars=!@|,fontsize=\small,frame=single,label=Example]
  !gapprompt@gap>| !gapinput@HasAbelianFactorGroup(g,n);|
  false
  !gapprompt@gap>| !gapinput@HasAbelianFactorGroup(DerivedSubgroup(g),n);|
  true
\end{Verbatim}
 }

 

\subsection{\textcolor{Chapter }{HasElementaryAbelianFactorGroup}}
\logpage{[ 39, 18, 6 ]}\nobreak
\hyperdef{L}{X7FAC018680B766B7}{}
{\noindent\textcolor{FuncColor}{$\triangleright$\enspace\texttt{HasElementaryAbelianFactorGroup({\mdseries\slshape G, N})\index{HasElementaryAbelianFactorGroup@\texttt{HasElementaryAbelianFactorGroup}}
\label{HasElementaryAbelianFactorGroup}
}\hfill{\scriptsize (function)}}\\


 tests whether \mbox{\texttt{\mdseries\slshape G}} $/$ \mbox{\texttt{\mdseries\slshape N}} is elementary abelian (without explicitly constructing the factor group and
without testing whether \mbox{\texttt{\mdseries\slshape N}} is in fact a normal subgroup). }

 

\subsection{\textcolor{Chapter }{CentralizerModulo}}
\logpage{[ 39, 18, 7 ]}\nobreak
\hyperdef{L}{X822A3AB27919BC1E}{}
{\noindent\textcolor{FuncColor}{$\triangleright$\enspace\texttt{CentralizerModulo({\mdseries\slshape G, N, elm})\index{CentralizerModulo@\texttt{CentralizerModulo}}
\label{CentralizerModulo}
}\hfill{\scriptsize (operation)}}\\


 Computes the full preimage of the centralizer $C_{{\mbox{\texttt{\mdseries\slshape G}}/\mbox{\texttt{\mdseries\slshape N}}}}(\mbox{\texttt{\mdseries\slshape elm}} \cdot \mbox{\texttt{\mdseries\slshape N}})$ in \mbox{\texttt{\mdseries\slshape G}} (without necessarily constructing the factor group). 
\begin{Verbatim}[commandchars=!@|,fontsize=\small,frame=single,label=Example]
  !gapprompt@gap>| !gapinput@CentralizerModulo(g,n,(1,2));|
  Group([ (3,4), (1,3)(2,4), (1,4)(2,3) ])
\end{Verbatim}
 }

  }

  
\section{\textcolor{Chapter }{Sets of Subgroups}}\label{Sets of Subgroups}
\logpage{[ 39, 19, 0 ]}
\hyperdef{L}{X7D8EFB2F85AA24EE}{}
{
  

\subsection{\textcolor{Chapter }{ConjugacyClassSubgroups}}
\logpage{[ 39, 19, 1 ]}\nobreak
\hyperdef{L}{X7DDE67C67E871336}{}
{\noindent\textcolor{FuncColor}{$\triangleright$\enspace\texttt{ConjugacyClassSubgroups({\mdseries\slshape G, U})\index{ConjugacyClassSubgroups@\texttt{ConjugacyClassSubgroups}}
\label{ConjugacyClassSubgroups}
}\hfill{\scriptsize (operation)}}\\


 generates the conjugacy class of subgroups of \mbox{\texttt{\mdseries\slshape G}} with representative \mbox{\texttt{\mdseries\slshape U}}. This class is an external set, so functions such as \texttt{Representative} (\ref{Representative}), (which returns \mbox{\texttt{\mdseries\slshape U}}), \texttt{ActingDomain} (\ref{ActingDomain}) (which returns \mbox{\texttt{\mdseries\slshape G}}), \texttt{StabilizerOfExternalSet} (\ref{StabilizerOfExternalSet}) (which returns the normalizer of \mbox{\texttt{\mdseries\slshape U}}), and \texttt{AsList} (\ref{AsList}) work for it. 

 (The use of the \texttt{[]} list access to select elements of the class is considered obsolescent and will
be removed in future versions. Use \texttt{ClassElementLattice} (\ref{ClassElementLattice}) instead.) 

 
\begin{Verbatim}[commandchars=!@|,fontsize=\small,frame=single,label=Example]
  !gapprompt@gap>| !gapinput@g:=Group((1,2,3,4),(1,2));;IsNaturalSymmetricGroup(g);;|
  !gapprompt@gap>| !gapinput@cl:=ConjugacyClassSubgroups(g,Subgroup(g,[(1,2)]));|
  Group( [ (1,2) ] )^G
  !gapprompt@gap>| !gapinput@Size(cl);|
  6
  !gapprompt@gap>| !gapinput@ClassElementLattice(cl,4);|
  Group([ (2,3) ])
\end{Verbatim}
 }

 

\subsection{\textcolor{Chapter }{IsConjugacyClassSubgroupsRep}}
\logpage{[ 39, 19, 2 ]}\nobreak
\hyperdef{L}{X7C5BBF487977B8CD}{}
{\noindent\textcolor{FuncColor}{$\triangleright$\enspace\texttt{IsConjugacyClassSubgroupsRep({\mdseries\slshape obj})\index{IsConjugacyClassSubgroupsRep@\texttt{IsConjugacyClassSubgroupsRep}}
\label{IsConjugacyClassSubgroupsRep}
}\hfill{\scriptsize (Representation)}}\\
\noindent\textcolor{FuncColor}{$\triangleright$\enspace\texttt{IsConjugacyClassSubgroupsByStabilizerRep({\mdseries\slshape obj})\index{IsConjugacyClassSubgroupsByStabilizerRep@\texttt{IsConjugacy}\-\texttt{Class}\-\texttt{Subgroups}\-\texttt{By}\-\texttt{StabilizerRep}}
\label{IsConjugacyClassSubgroupsByStabilizerRep}
}\hfill{\scriptsize (Representation)}}\\


 Is the representation \textsf{GAP} uses for conjugacy classes of subgroups. It can be used to check whether an
object is a class of subgroups. The second representation \texttt{IsConjugacyClassSubgroupsByStabilizerRep} in addition is an external orbit by stabilizer and will compute its elements
via a transversal of the stabilizer. }

 

\subsection{\textcolor{Chapter }{ConjugacyClassesSubgroups}}
\logpage{[ 39, 19, 3 ]}\nobreak
\hyperdef{L}{X7E986BF48393113A}{}
{\noindent\textcolor{FuncColor}{$\triangleright$\enspace\texttt{ConjugacyClassesSubgroups({\mdseries\slshape G})\index{ConjugacyClassesSubgroups@\texttt{ConjugacyClassesSubgroups}}
\label{ConjugacyClassesSubgroups}
}\hfill{\scriptsize (attribute)}}\\


 This attribute returns a list of all conjugacy classes of subgroups of the
group \mbox{\texttt{\mdseries\slshape G}}. It also is applicable for lattices of subgroups (see{\nobreakspace}\texttt{LatticeSubgroups} (\ref{LatticeSubgroups})). The order in which the classes are listed depends on the method chosen by \textsf{GAP}. For each class of subgroups, a representative can be accessed using \texttt{Representative} (\ref{Representative}). 
\begin{Verbatim}[commandchars=!@|,fontsize=\small,frame=single,label=Example]
  !gapprompt@gap>| !gapinput@ConjugacyClassesSubgroups(g);|
  [ Group( () )^G, Group( [ (1,3)(2,4) ] )^G, Group( [ (3,4) ] )^G,
    Group( [ (2,4,3) ] )^G, Group( [ (1,4)(2,3), (1,3)(2,4) ] )^G,
    Group( [ (3,4), (1,2)(3,4) ] )^G,
    Group( [ (1,3,2,4), (1,2)(3,4) ] )^G, Group( [ (3,4), (2,4,3) ] )^G,
    Group( [ (1,4)(2,3), (1,3)(2,4), (3,4) ] )^G,
    Group( [ (1,4)(2,3), (1,3)(2,4), (2,4,3) ] )^G,
    Group( [ (1,4)(2,3), (1,3)(2,4), (2,4,3), (3,4) ] )^G ]
\end{Verbatim}
 }

 

\subsection{\textcolor{Chapter }{ConjugacyClassesMaximalSubgroups}}
\logpage{[ 39, 19, 4 ]}\nobreak
\hyperdef{L}{X8486C25380853F9B}{}
{\noindent\textcolor{FuncColor}{$\triangleright$\enspace\texttt{ConjugacyClassesMaximalSubgroups({\mdseries\slshape G})\index{ConjugacyClassesMaximalSubgroups@\texttt{ConjugacyClassesMaximalSubgroups}}
\label{ConjugacyClassesMaximalSubgroups}
}\hfill{\scriptsize (attribute)}}\\


 returns the conjugacy classes of maximal subgroups of \mbox{\texttt{\mdseries\slshape G}}. Representatives of the classes can be computed directly by \texttt{MaximalSubgroupClassReps} (\ref{MaximalSubgroupClassReps}). 
\begin{Verbatim}[commandchars=!@|,fontsize=\small,frame=single,label=Example]
  !gapprompt@gap>| !gapinput@ConjugacyClassesMaximalSubgroups(g);|
  [ Group( [ (2,4,3), (1,4)(2,3), (1,3)(2,4) ] )^G,
    Group( [ (3,4), (1,3)(2,4) ] )^G, Group( [ (3,4), (2,4,3) ] )^G ]
\end{Verbatim}
 }

 

\subsection{\textcolor{Chapter }{MaximalSubgroupClassReps}}
\logpage{[ 39, 19, 5 ]}\nobreak
\hyperdef{L}{X798BF55C837DB188}{}
{\noindent\textcolor{FuncColor}{$\triangleright$\enspace\texttt{MaximalSubgroupClassReps({\mdseries\slshape G})\index{MaximalSubgroupClassReps@\texttt{MaximalSubgroupClassReps}}
\label{MaximalSubgroupClassReps}
}\hfill{\scriptsize (attribute)}}\\


 returns a list of conjugacy representatives of the maximal subgroups of \mbox{\texttt{\mdseries\slshape G}}. 
\begin{Verbatim}[commandchars=!@|,fontsize=\small,frame=single,label=Example]
  !gapprompt@gap>| !gapinput@MaximalSubgroupClassReps(g);|
  [ Group([ (2,4,3), (1,4)(2,3), (1,3)(2,4) ]),
    Group([ (3,4), (1,3)(2,4) ]), Group([ (3,4), (2,4,3) ]) ]
\end{Verbatim}
 }

 

\subsection{\textcolor{Chapter }{LowIndexSubgroups}}
\logpage{[ 39, 19, 6 ]}\nobreak
\hyperdef{L}{X85DAFB7582A88463}{}
{\noindent\textcolor{FuncColor}{$\triangleright$\enspace\texttt{LowIndexSubgroups({\mdseries\slshape G, index})\index{LowIndexSubgroups@\texttt{LowIndexSubgroups}}
\label{LowIndexSubgroups}
}\hfill{\scriptsize (operation)}}\\


 The operation \texttt{LowIndexSubgroups} computes representatives of the conjugacy classes of subgroups of the group \mbox{\texttt{\mdseries\slshape G}} that have index less than or equal to \mbox{\texttt{\mdseries\slshape index}}. 

 For finitely presented groups this operation simply defaults to \texttt{LowIndexSubgroupsFpGroup} (\ref{LowIndexSubgroupsFpGroup}). In other cases, it uses repeated calculation of maximal subgroups. 

 The function \texttt{LowLayerSubgroups} (\ref{LowLayerSubgroups}) works similar but does not bound the index, but instead considers up to \mbox{\texttt{\mdseries\slshape layer}}\texttt{\symbol{45}}th maximal subgroups. 
\begin{Verbatim}[commandchars=!@|,fontsize=\small,frame=single,label=Example]
  !gapprompt@gap>| !gapinput@g:=TransitiveGroup(18,950);;|
  !gapprompt@gap>| !gapinput@l:=LowIndexSubgroups(g,20);;Collected(List(l,x->Index(g,x)));|
  [ [ 1, 1 ], [ 2, 1 ], [ 5, 1 ], [ 6, 1 ], [ 10, 2 ], [ 12, 3 ], [ 15, 1 ],
    [ 16, 2 ], [ 18, 1 ], [ 20, 9 ] ]
\end{Verbatim}
 }

 

\subsection{\textcolor{Chapter }{AllSubgroups}}
\logpage{[ 39, 19, 7 ]}\nobreak
\hyperdef{L}{X80399CD4870FFC4B}{}
{\noindent\textcolor{FuncColor}{$\triangleright$\enspace\texttt{AllSubgroups({\mdseries\slshape G})\index{AllSubgroups@\texttt{AllSubgroups}}
\label{AllSubgroups}
}\hfill{\scriptsize (function)}}\\


 For a finite group \mbox{\texttt{\mdseries\slshape G}} \texttt{AllSubgroups} returns a list of all subgroups of \mbox{\texttt{\mdseries\slshape G}}, intended primarily for use in class for small examples. This list will
quickly get very long and in general use of \texttt{ConjugacyClassesSubgroups} (\ref{ConjugacyClassesSubgroups}) is recommended. 
\begin{Verbatim}[commandchars=!@|,fontsize=\small,frame=single,label=Example]
  !gapprompt@gap>| !gapinput@AllSubgroups(SymmetricGroup(3));|
  [ Group(()), Group([ (2,3) ]), Group([ (1,2) ]), Group([ (1,3) ]),
    Group([ (1,2,3) ]), Group([ (1,2,3), (2,3) ]) ]
\end{Verbatim}
 }

 

\subsection{\textcolor{Chapter }{MaximalSubgroups}}
\logpage{[ 39, 19, 8 ]}\nobreak
\hyperdef{L}{X861CD8DA790D81C2}{}
{\noindent\textcolor{FuncColor}{$\triangleright$\enspace\texttt{MaximalSubgroups({\mdseries\slshape G})\index{MaximalSubgroups@\texttt{MaximalSubgroups}}
\label{MaximalSubgroups}
}\hfill{\scriptsize (attribute)}}\\


 returns a list of all maximal subgroups of \mbox{\texttt{\mdseries\slshape G}}. This may take up much space, therefore the command should be avoided if
possible. See \texttt{ConjugacyClassesMaximalSubgroups} (\ref{ConjugacyClassesMaximalSubgroups}). 
\begin{Verbatim}[commandchars=!@|,fontsize=\small,frame=single,label=Example]
  !gapprompt@gap>| !gapinput@MaximalSubgroups(Group((1,2,3),(1,2)));|
  [ Group([ (1,2,3) ]), Group([ (2,3) ]), Group([ (1,2) ]),
    Group([ (1,3) ]) ]
\end{Verbatim}
 }

 

\subsection{\textcolor{Chapter }{NormalSubgroups}}
\logpage{[ 39, 19, 9 ]}\nobreak
\hyperdef{L}{X80237A847E24E6CF}{}
{\noindent\textcolor{FuncColor}{$\triangleright$\enspace\texttt{NormalSubgroups({\mdseries\slshape G})\index{NormalSubgroups@\texttt{NormalSubgroups}}
\label{NormalSubgroups}
}\hfill{\scriptsize (attribute)}}\\


 returns a list of all normal subgroups of \mbox{\texttt{\mdseries\slshape G}}. 
\begin{Verbatim}[commandchars=!@|,fontsize=\small,frame=single,label=Example]
  !gapprompt@gap>| !gapinput@g:=SymmetricGroup(4);;|
  !gapprompt@gap>| !gapinput@List( NormalSubgroups(g), StructureDescription );|
  [ "S4", "A4", "C2 x C2", "1" ]
  !gapprompt@gap>| !gapinput@g:=AbelianGroup([2,2]);; NormalSubgroups(g);|
  [ <pc group of size 4 with 2 generators>, Group([ f2 ]),
    Group([ f1*f2 ]), Group([ f1 ]), Group([  ]) ]
\end{Verbatim}
 

 The algorithm for the computation of normal subgroups is described in \cite{Hulpke98}. }

 

\subsection{\textcolor{Chapter }{MaximalNormalSubgroups}}
\logpage{[ 39, 19, 10 ]}\nobreak
\hyperdef{L}{X82ECAA427C987318}{}
{\noindent\textcolor{FuncColor}{$\triangleright$\enspace\texttt{MaximalNormalSubgroups({\mdseries\slshape G})\index{MaximalNormalSubgroups@\texttt{MaximalNormalSubgroups}}
\label{MaximalNormalSubgroups}
}\hfill{\scriptsize (attribute)}}\\


 is a list containing those proper normal subgroups of the group \mbox{\texttt{\mdseries\slshape G}} that are maximal among the proper normal subgroups. Gives error if \mbox{\texttt{\mdseries\slshape G}}/\mbox{\texttt{\mdseries\slshape G'}} is infinite, yielding infinitely many maximal normal subgroups. Note, that the
maximal normal subgroups of a group \mbox{\texttt{\mdseries\slshape G}} can be computed more efficiently if the character table of \mbox{\texttt{\mdseries\slshape G}} is known or if \mbox{\texttt{\mdseries\slshape G}} is known to be abelian or solvable (even if infinite). So if the character
table is needed, anyhow, or \mbox{\texttt{\mdseries\slshape G}} is suspected to be abelian or solvable, then these should be computed before
computing the maximal normal subgroups. 
\begin{Verbatim}[commandchars=!@|,fontsize=\small,frame=single,label=Example]
  !gapprompt@gap>| !gapinput@g:=SymmetricGroup(4);; MaximalNormalSubgroups( g );|
  [ Alt( [ 1 .. 4 ] ) ]
  !gapprompt@gap>| !gapinput@f := FreeGroup("x", "y");; x := f.1;; y := f.2;;|
  !gapprompt@gap>| !gapinput@List(MaximalNormalSubgroups(f/[x^2, y^2]), GeneratorsOfGroup);|
  [ [ x, y*x*y^-1 ], [ y, x*y*x^-1 ], [ y*x^-1 ] ]
\end{Verbatim}
 }

 

\subsection{\textcolor{Chapter }{MinimalNormalSubgroups}}
\logpage{[ 39, 19, 11 ]}\nobreak
\hyperdef{L}{X86FDD9BA819F5644}{}
{\noindent\textcolor{FuncColor}{$\triangleright$\enspace\texttt{MinimalNormalSubgroups({\mdseries\slshape G})\index{MinimalNormalSubgroups@\texttt{MinimalNormalSubgroups}}
\label{MinimalNormalSubgroups}
}\hfill{\scriptsize (attribute)}}\\


 is a list containing those nontrivial normal subgroups of the group \mbox{\texttt{\mdseries\slshape G}} that are minimal among the nontrivial normal subgroups. 
\begin{Verbatim}[commandchars=!@|,fontsize=\small,frame=single,label=Example]
  !gapprompt@gap>| !gapinput@g:=SymmetricGroup(4);; MinimalNormalSubgroups( g );|
  [ Group([ (1,4)(2,3), (1,3)(2,4) ]) ]
\end{Verbatim}
 }

 

\subsection{\textcolor{Chapter }{CharacteristicSubgroups}}
\logpage{[ 39, 19, 12 ]}\nobreak
\hyperdef{L}{X7A823C5A810910C3}{}
{\noindent\textcolor{FuncColor}{$\triangleright$\enspace\texttt{CharacteristicSubgroups({\mdseries\slshape G})\index{CharacteristicSubgroups@\texttt{CharacteristicSubgroups}}
\label{CharacteristicSubgroups}
}\hfill{\scriptsize (attribute)}}\\


 returns a list of all characteristic subgroups of \mbox{\texttt{\mdseries\slshape G}}, that is subgroups that are invariant under all automorphisms. 
\begin{Verbatim}[commandchars=!@|,fontsize=\small,frame=single,label=Example]
  !gapprompt@gap>| !gapinput@g:=SymmetricGroup(4);;|
  !gapprompt@gap>| !gapinput@List( CharacteristicSubgroups(g), StructureDescription );|
  [ "S4", "A4", "C2 x C2", "1" ]
  !gapprompt@gap>| !gapinput@g:=AbelianGroup([2,2]);; CharacteristicSubgroups(g);|
  [ <pc group of size 4 with 2 generators>, Group([  ]) ]
\end{Verbatim}
 

 }

    }

  
\section{\textcolor{Chapter }{Subgroup Lattice}}\label{Subgroup Lattice}
\logpage{[ 39, 20, 0 ]}
\hyperdef{L}{X7FA267497CFC0550}{}
{
  

\subsection{\textcolor{Chapter }{LatticeSubgroups}}
\logpage{[ 39, 20, 1 ]}\nobreak
\hyperdef{L}{X7B104E2C86166188}{}
{\noindent\textcolor{FuncColor}{$\triangleright$\enspace\texttt{LatticeSubgroups({\mdseries\slshape G})\index{LatticeSubgroups@\texttt{LatticeSubgroups}}
\label{LatticeSubgroups}
}\hfill{\scriptsize (attribute)}}\\


 computes the lattice of subgroups of the group \mbox{\texttt{\mdseries\slshape G}}. This lattice has the conjugacy classes of subgroups as attribute \texttt{ConjugacyClassesSubgroups} (\ref{ConjugacyClassesSubgroups}) and permits one to test maximality/minimality relations. 
\begin{Verbatim}[commandchars=!@|,fontsize=\small,frame=single,label=Example]
  !gapprompt@gap>| !gapinput@g:=SymmetricGroup(4);;|
  !gapprompt@gap>| !gapinput@l:=LatticeSubgroups(g);|
  <subgroup lattice of Sym( [ 1 .. 4 ] ), 11 classes, 30 subgroups>
  !gapprompt@gap>| !gapinput@ConjugacyClassesSubgroups(l);|
  [ Group( () )^G, Group( [ (1,3)(2,4) ] )^G, Group( [ (3,4) ] )^G,
    Group( [ (2,4,3) ] )^G, Group( [ (1,4)(2,3), (1,3)(2,4) ] )^G,
    Group( [ (3,4), (1,2)(3,4) ] )^G,
    Group( [ (1,3,2,4), (1,2)(3,4) ] )^G, Group( [ (3,4), (2,4,3) ] )^G,
    Group( [ (1,4)(2,3), (1,3)(2,4), (3,4) ] )^G,
    Group( [ (1,4)(2,3), (1,3)(2,4), (2,4,3) ] )^G,
    Group( [ (1,4)(2,3), (1,3)(2,4), (2,4,3), (3,4) ] )^G ]
\end{Verbatim}
 }

 

\subsection{\textcolor{Chapter }{ClassElementLattice}}
\logpage{[ 39, 20, 2 ]}\nobreak
\hyperdef{L}{X78928A3582882BFD}{}
{\noindent\textcolor{FuncColor}{$\triangleright$\enspace\texttt{ClassElementLattice({\mdseries\slshape C, n})\index{ClassElementLattice@\texttt{ClassElementLattice}}
\label{ClassElementLattice}
}\hfill{\scriptsize (operation)}}\\


 For a class \mbox{\texttt{\mdseries\slshape C}} of subgroups, obtained by a lattice computation, this operation returns the \mbox{\texttt{\mdseries\slshape n}}\texttt{\symbol{45}}th conjugate subgroup in the class. 

 \emph{Because of other methods installed, calling \texttt{AsList} (\ref{AsList}) with \mbox{\texttt{\mdseries\slshape C}} can give a different arrangement of the class elements!} 

 The \textsf{GAP} package \textsf{XGAP} permits a graphical display of the lattice of subgroups in a nice way. }

 

\subsection{\textcolor{Chapter }{DotFileLatticeSubgroups}}
\logpage{[ 39, 20, 3 ]}\nobreak
\hyperdef{L}{X7E5DF287825EE7BA}{}
{\noindent\textcolor{FuncColor}{$\triangleright$\enspace\texttt{DotFileLatticeSubgroups({\mdseries\slshape L, file})\index{DotFileLatticeSubgroups@\texttt{DotFileLatticeSubgroups}}
\label{DotFileLatticeSubgroups}
}\hfill{\scriptsize (function)}}\\


 \index{dot\texttt{\symbol{45}}file} \index{graphviz} \index{OmniGraffle} This function produces a graphical representation of the subgroup lattice \mbox{\texttt{\mdseries\slshape L}} in file \mbox{\texttt{\mdseries\slshape file}}. The output is in \texttt{.dot} (also known as \texttt{GraphViz} format). For details on the format, and information about how to display or
edit this format see \href{https://www.graphviz.org} {\texttt{https://www.graphviz.org}}. (On the Macintosh, the program \texttt{OmniGraffle} is also able to read this format.) 

 Subgroups are labelled in the form \texttt{\mbox{\texttt{\mdseries\slshape i}}\texttt{\symbol{45}}\mbox{\texttt{\mdseries\slshape j}}} where \mbox{\texttt{\mdseries\slshape i}} is the number of the class of subgroups and \mbox{\texttt{\mdseries\slshape j}} the number within this class. Normal subgroups are represented by a box. 

 
\begin{Verbatim}[commandchars=!@|,fontsize=\small,frame=single,label=Example]
  !gapprompt@gap>| !gapinput@DotFileLatticeSubgroups(l,"s4lat.dot");|
\end{Verbatim}
 }

 

\subsection{\textcolor{Chapter }{MaximalSubgroupsLattice}}
\logpage{[ 39, 20, 4 ]}\nobreak
\hyperdef{L}{X815CDA447C5DB285}{}
{\noindent\textcolor{FuncColor}{$\triangleright$\enspace\texttt{MaximalSubgroupsLattice({\mdseries\slshape lat})\index{MaximalSubgroupsLattice@\texttt{MaximalSubgroupsLattice}}
\label{MaximalSubgroupsLattice}
}\hfill{\scriptsize (attribute)}}\\


 For a lattice \mbox{\texttt{\mdseries\slshape lat}} of subgroups this attribute contains the maximal subgroup relations among the
subgroups of the lattice. It is a list corresponding to the \texttt{ConjugacyClassesSubgroups} (\ref{ConjugacyClassesSubgroups}) value of the lattice, each entry giving a list of the maximal subgroups of the
representative of this class. Every maximal subgroup is indicated by a list of
the form $[ c, n ]$ which means that the $n$\texttt{\symbol{45}}th subgroup in class number $c$ is a maximal subgroup of the representative. 

 The number $n$ corresponds to access via \texttt{ClassElementLattice} (\ref{ClassElementLattice}) and \emph{not} necessarily the \texttt{AsList} (\ref{AsList}) arrangement! See also \texttt{MinimalSupergroupsLattice} (\ref{MinimalSupergroupsLattice}). 
\begin{Verbatim}[commandchars=!@|,fontsize=\small,frame=single,label=Example]
  !gapprompt@gap>| !gapinput@MaximalSubgroupsLattice(l);|
  [ [  ], [ [ 1, 1 ] ], [ [ 1, 1 ] ], [ [ 1, 1 ] ],
    [ [ 2, 1 ], [ 2, 2 ], [ 2, 3 ] ], [ [ 3, 1 ], [ 3, 6 ], [ 2, 3 ] ],
    [ [ 2, 3 ] ], [ [ 4, 1 ], [ 3, 1 ], [ 3, 2 ], [ 3, 3 ] ],
    [ [ 7, 1 ], [ 6, 1 ], [ 5, 1 ] ],
    [ [ 5, 1 ], [ 4, 1 ], [ 4, 2 ], [ 4, 3 ], [ 4, 4 ] ],
    [ [ 10, 1 ], [ 9, 1 ], [ 9, 2 ], [ 9, 3 ], [ 8, 1 ], [ 8, 2 ],
        [ 8, 3 ], [ 8, 4 ] ] ]
  !gapprompt@gap>| !gapinput@last[6];|
  [ [ 3, 1 ], [ 3, 6 ], [ 2, 3 ] ]
  !gapprompt@gap>| !gapinput@u1:=Representative(ConjugacyClassesSubgroups(l)[6]);|
  Group([ (3,4), (1,2)(3,4) ])
  !gapprompt@gap>| !gapinput@u2:=ClassElementLattice(ConjugacyClassesSubgroups(l)[3],1);;|
  !gapprompt@gap>| !gapinput@u3:=ClassElementLattice(ConjugacyClassesSubgroups(l)[3],6);;|
  !gapprompt@gap>| !gapinput@u4:=ClassElementLattice(ConjugacyClassesSubgroups(l)[2],3);;|
  !gapprompt@gap>| !gapinput@IsSubgroup(u1,u2);IsSubgroup(u1,u3);IsSubgroup(u1,u4);|
  true
  true
  true
\end{Verbatim}
 }

 

\subsection{\textcolor{Chapter }{MinimalSupergroupsLattice}}
\logpage{[ 39, 20, 5 ]}\nobreak
\hyperdef{L}{X8138997C871EDF96}{}
{\noindent\textcolor{FuncColor}{$\triangleright$\enspace\texttt{MinimalSupergroupsLattice({\mdseries\slshape lat})\index{MinimalSupergroupsLattice@\texttt{MinimalSupergroupsLattice}}
\label{MinimalSupergroupsLattice}
}\hfill{\scriptsize (attribute)}}\\


 For a lattice \mbox{\texttt{\mdseries\slshape lat}} of subgroups this attribute contains the minimal supergroup relations among
the subgroups of the lattice. It is a list corresponding to the \texttt{ConjugacyClassesSubgroups} (\ref{ConjugacyClassesSubgroups}) value of the lattice, each entry giving a list of the minimal supergroups of
the representative of this class. Every minimal supergroup is indicated by a
list of the form $[ c, n ]$, which means that the $n$\texttt{\symbol{45}}th subgroup in class number $c$ is a minimal supergroup of the representative. 

 The number $n$ corresponds to access via \texttt{ClassElementLattice} (\ref{ClassElementLattice}) and \emph{not} necessarily the \texttt{AsList} (\ref{AsList}) arrangement! See also \texttt{MaximalSubgroupsLattice} (\ref{MaximalSubgroupsLattice}). 
\begin{Verbatim}[commandchars=!@|,fontsize=\small,frame=single,label=Example]
  !gapprompt@gap>| !gapinput@MinimalSupergroupsLattice(l);|
  [ [ [ 2, 1 ], [ 2, 2 ], [ 2, 3 ], [ 3, 1 ], [ 3, 2 ], [ 3, 3 ],
        [ 3, 4 ], [ 3, 5 ], [ 3, 6 ], [ 4, 1 ], [ 4, 2 ], [ 4, 3 ],
        [ 4, 4 ] ], [ [ 5, 1 ], [ 6, 2 ], [ 7, 2 ] ],
    [ [ 6, 1 ], [ 8, 1 ], [ 8, 3 ] ], [ [ 8, 1 ], [ 10, 1 ] ],
    [ [ 9, 1 ], [ 9, 2 ], [ 9, 3 ], [ 10, 1 ] ], [ [ 9, 1 ] ],
    [ [ 9, 1 ] ], [ [ 11, 1 ] ], [ [ 11, 1 ] ], [ [ 11, 1 ] ], [  ] ]
  !gapprompt@gap>| !gapinput@last[3];|
  [ [ 6, 1 ], [ 8, 1 ], [ 8, 3 ] ]
  !gapprompt@gap>| !gapinput@u5:=ClassElementLattice(ConjugacyClassesSubgroups(l)[8],1);|
  Group([ (3,4), (2,4,3) ])
  !gapprompt@gap>| !gapinput@u6:=ClassElementLattice(ConjugacyClassesSubgroups(l)[8],3);|
  Group([ (1,3), (1,3,4) ])
  !gapprompt@gap>| !gapinput@IsSubgroup(u5,u2);|
  true
  !gapprompt@gap>| !gapinput@IsSubgroup(u6,u2);|
  true
\end{Verbatim}
 }

 

\subsection{\textcolor{Chapter }{LowLayerSubgroups}}
\logpage{[ 39, 20, 6 ]}\nobreak
\hyperdef{L}{X87BE970D7B18E2C5}{}
{\noindent\textcolor{FuncColor}{$\triangleright$\enspace\texttt{LowLayerSubgroups({\mdseries\slshape [act, ]G, lim[, cond, dosub]})\index{LowLayerSubgroups@\texttt{LowLayerSubgroups}}
\label{LowLayerSubgroups}
}\hfill{\scriptsize (function)}}\\


 This function computes representatives of the conjugacy classes of subgroups
of the finite group \mbox{\texttt{\mdseries\slshape G}} such that the subgroups can be obtained as \mbox{\texttt{\mdseries\slshape lim}}\texttt{\symbol{45}}fold iterated maximal subgroups. If a function \mbox{\texttt{\mdseries\slshape cond}} is given, only subgroups for which this function returns true (also for their
intermediate overgroups) is returned. If also a function \mbox{\texttt{\mdseries\slshape dosub}} is given, maximal subgroups are only attempted if this function returns true
(this is separated for performance reasons). In the example below, the result
would be the same with leaving out the fourth function, but calculation this
way is slightly faster. If an initial argument \mbox{\texttt{\mdseries\slshape act}} is given, it must be a group containing and normalizing \mbox{\texttt{\mdseries\slshape G}}, and representatives for classes under the action of this group are chosen. 
\begin{Verbatim}[commandchars=!@|,fontsize=\small,frame=single,label=Example]
  !gapprompt@gap>| !gapinput@g:=SymmetricGroup(12);;|
  !gapprompt@gap>| !gapinput@l:=LowLayerSubgroups(g,2,x->Size(x)>100000,x->Size(x)>200000);;|
  !gapprompt@gap>| !gapinput@Collected(List(l,Size));|
  [ [ 100800, 1 ], [ 120960, 1 ], [ 161280, 1 ], [ 241920, 1 ], [ 302400, 3 ],
    [ 322560, 1 ], [ 483840, 3 ], [ 518400, 3 ], [ 604800, 1 ], [ 725760, 1 ],
    [ 967680, 1 ], [ 1036800, 1 ], [ 1088640, 3 ], [ 2177280, 1 ],
    [ 3628800, 3 ], [ 7257600, 1 ], [ 19958400, 1 ], [ 39916800, 1 ],
    [ 239500800, 1 ], [ 479001600, 1 ] ]
\end{Verbatim}
 }

 

\subsection{\textcolor{Chapter }{ContainedConjugates}}
\logpage{[ 39, 20, 7 ]}\nobreak
\hyperdef{L}{X87FABD5F87AD2568}{}
{\noindent\textcolor{FuncColor}{$\triangleright$\enspace\texttt{ContainedConjugates({\mdseries\slshape G, A, B[, onlyone]})\index{ContainedConjugates@\texttt{ContainedConjugates}}
\label{ContainedConjugates}
}\hfill{\scriptsize (operation)}}\\


 For $A,B \leq G$ this operation returns representatives of the \mbox{\texttt{\mdseries\slshape A}} conjugacy classes of subgroups that are conjugate to \mbox{\texttt{\mdseries\slshape B}} under \mbox{\texttt{\mdseries\slshape G}}. The function returns a list of pairs of subgroup and conjugating element. If
the optional fourth argument \mbox{\texttt{\mdseries\slshape onlyone}} is given as \mbox{\texttt{\mdseries\slshape true}}, then only one pair (or \mbox{\texttt{\mdseries\slshape fail}} if none exists) is returned. 
\begin{Verbatim}[commandchars=!@|,fontsize=\small,frame=single,label=Example]
  !gapprompt@gap>| !gapinput@g:=SymmetricGroup(8);;|
  !gapprompt@gap>| !gapinput@a:=TransitiveGroup(8,47);;b:=TransitiveGroup(8,9);;|
  !gapprompt@gap>| !gapinput@ContainedConjugates(g,a,b);|
  [ [ Group([ (1,8)(2,3)(4,5)(6,7), (1,3)(2,8)(4,6)(5,7), (1,5)(2,6)(3,7)(4,8),
          (4,5)(6,7) ]), () ],
    [ Group([ (1,8)(2,3)(4,5)(6,7), (1,5)(2,6)(3,7)(4,8), (1,3)(2,8)(4,6)(5,7),
          (2,3)(6,7) ]), (2,4)(3,5) ] ]
  !gapprompt@gap>| !gapinput@ContainedConjugates(g,a,b,true);|
  [ Group([ (1,8)(2,3)(4,5)(6,7), (1,3)(2,8)(4,6)(5,7), (1,5)(2,6)(3,7)(4,8),
      (4,5)(6,7) ]), () ]
\end{Verbatim}
 }

 

\subsection{\textcolor{Chapter }{ContainingConjugates}}
\logpage{[ 39, 20, 8 ]}\nobreak
\hyperdef{L}{X79C3619C849F97B8}{}
{\noindent\textcolor{FuncColor}{$\triangleright$\enspace\texttt{ContainingConjugates({\mdseries\slshape G, A, B})\index{ContainingConjugates@\texttt{ContainingConjugates}}
\label{ContainingConjugates}
}\hfill{\scriptsize (operation)}}\\


 For $A,B \leq G$ this operation returns all \mbox{\texttt{\mdseries\slshape G}} conjugates of \mbox{\texttt{\mdseries\slshape A}} that contain \mbox{\texttt{\mdseries\slshape B}}. The function returns a list of pairs of subgroup and conjugating element. 
\begin{Verbatim}[commandchars=!@|,fontsize=\small,frame=single,label=Example]
  !gapprompt@gap>| !gapinput@g:=SymmetricGroup(8);;|
  !gapprompt@gap>| !gapinput@a:=TransitiveGroup(8,47);;b:=TransitiveGroup(8,7);;|
  !gapprompt@gap>| !gapinput@ContainingConjugates(g,a,b);|
  [ [ Group([ (1,3,5,7), (3,5), (1,4)(2,7)(3,6)(5,8) ]), (2,3,5,4)(7,8) ] ]
\end{Verbatim}
 }

 

\subsection{\textcolor{Chapter }{MinimalFaithfulPermutationDegree}}
\logpage{[ 39, 20, 9 ]}\nobreak
\hyperdef{L}{X8111F50C798B0D76}{}
{\noindent\textcolor{FuncColor}{$\triangleright$\enspace\texttt{MinimalFaithfulPermutationDegree({\mdseries\slshape G})\index{MinimalFaithfulPermutationDegree@\texttt{MinimalFaithfulPermutationDegree}}
\label{MinimalFaithfulPermutationDegree}
}\hfill{\scriptsize (operation)}}\\
\noindent\textcolor{FuncColor}{$\triangleright$\enspace\texttt{MinimalFaithfulPermutationRepresentation({\mdseries\slshape G})\index{MinimalFaithfulPermutationRepresentation@\texttt{Minimal}\-\texttt{Faithful}\-\texttt{Permutation}\-\texttt{Representation}}
\label{MinimalFaithfulPermutationRepresentation}
}\hfill{\scriptsize (operation)}}\\


 For a finite group \mbox{\texttt{\mdseries\slshape G}}, \texttt{MinimalFaithfulPermutationDegree} calculates the least positive integer $n=\mu(G)$ such that \mbox{\texttt{\mdseries\slshape G}} is isomorphic to a subgroup of the symmetric group of degree $n$. This can require calculating the whole subgroup lattice. The operation \texttt{MinimalFaithfulPermutationRepresentation} returns a corresponding isomorphism. 
\begin{Verbatim}[commandchars=!@|,fontsize=\small,frame=single,label=Example]
  !gapprompt@gap>| !gapinput@MinimalFaithfulPermutationDegree(SmallGroup(96,3));|
  12
  !gapprompt@gap>| !gapinput@g:=TransitiveGroup(10,32);;|
  !gapprompt@gap>| !gapinput@MinimalFaithfulPermutationDegree(g);|
  6
  !gapprompt@gap>| !gapinput@map:=MinimalFaithfulPermutationRepresentation(g);;|
  !gapprompt@gap>| !gapinput@Size(Image(map));|
  720
\end{Verbatim}
 }

 

\subsection{\textcolor{Chapter }{RepresentativesPerfectSubgroups}}
\logpage{[ 39, 20, 10 ]}\nobreak
\hyperdef{L}{X7BA3484E7AE0A0E1}{}
{\noindent\textcolor{FuncColor}{$\triangleright$\enspace\texttt{RepresentativesPerfectSubgroups({\mdseries\slshape G})\index{RepresentativesPerfectSubgroups@\texttt{RepresentativesPerfectSubgroups}}
\label{RepresentativesPerfectSubgroups}
}\hfill{\scriptsize (attribute)}}\\
\noindent\textcolor{FuncColor}{$\triangleright$\enspace\texttt{RepresentativesSimpleSubgroups({\mdseries\slshape G})\index{RepresentativesSimpleSubgroups@\texttt{RepresentativesSimpleSubgroups}}
\label{RepresentativesSimpleSubgroups}
}\hfill{\scriptsize (attribute)}}\\


 returns a list of conjugacy representatives of perfect (respectively simple)
subgroups of \mbox{\texttt{\mdseries\slshape G}}. This uses the library of perfect groups (see \texttt{PerfectGroup} (\ref{PerfectGroup:for group order (and index)})), thus it will issue an error if the library is insufficient to determine all
perfect subgroups. 
\begin{Verbatim}[commandchars=!@|,fontsize=\small,frame=single,label=Example]
  !gapprompt@gap>| !gapinput@m11:=TransitiveGroup(11,6);|
  M(11)
  !gapprompt@gap>| !gapinput@r:=RepresentativesPerfectSubgroups(m11);;|
  !gapprompt@gap>| !gapinput@List(r,Size);|
  [ 60, 60, 360, 660, 7920, 1 ]
  !gapprompt@gap>| !gapinput@List(r,StructureDescription);|
  [ "A5", "A5", "A6", "PSL(2,11)", "M11", "1" ]
\end{Verbatim}
 }

 

\subsection{\textcolor{Chapter }{ConjugacyClassesPerfectSubgroups}}
\logpage{[ 39, 20, 11 ]}\nobreak
\hyperdef{L}{X7B2233D180DF77A1}{}
{\noindent\textcolor{FuncColor}{$\triangleright$\enspace\texttt{ConjugacyClassesPerfectSubgroups({\mdseries\slshape G})\index{ConjugacyClassesPerfectSubgroups@\texttt{ConjugacyClassesPerfectSubgroups}}
\label{ConjugacyClassesPerfectSubgroups}
}\hfill{\scriptsize (attribute)}}\\


 returns a list of the conjugacy classes of perfect subgroups of \mbox{\texttt{\mdseries\slshape G}}. (see \texttt{RepresentativesPerfectSubgroups} (\ref{RepresentativesPerfectSubgroups}).) 
\begin{Verbatim}[commandchars=!@|,fontsize=\small,frame=single,label=Example]
  !gapprompt@gap>| !gapinput@r := ConjugacyClassesPerfectSubgroups(m11);;|
  !gapprompt@gap>| !gapinput@List(r, x -> StructureDescription(Representative(x)));|
  [ "A5", "A5", "A6", "PSL(2,11)", "M11", "1" ]
  !gapprompt@gap>| !gapinput@SortedList( List(r,Size) );|
  [ 1, 1, 11, 12, 66, 132 ]
\end{Verbatim}
 }

 

\subsection{\textcolor{Chapter }{Zuppos}}
\logpage{[ 39, 20, 12 ]}\nobreak
\hyperdef{L}{X7BFE573187B4BEF8}{}
{\noindent\textcolor{FuncColor}{$\triangleright$\enspace\texttt{Zuppos({\mdseries\slshape G})\index{Zuppos@\texttt{Zuppos}}
\label{Zuppos}
}\hfill{\scriptsize (attribute)}}\\


 The \emph{Zuppos} of a group are the cyclic subgroups of prime power order. (The name ``Zuppo'' derives from the German abbreviation for ``zyklische Untergruppen von Primzahlpotenzordnung''.) This attribute gives generators of all such subgroups of a group \mbox{\texttt{\mdseries\slshape G}}. That is all elements of \mbox{\texttt{\mdseries\slshape G}} of prime power order up to the equivalence that they generate the same cyclic
subgroup. }

 

\subsection{\textcolor{Chapter }{InfoLattice}}
\logpage{[ 39, 20, 13 ]}\nobreak
\hyperdef{L}{X82C12E2C81963B23}{}
{\noindent\textcolor{FuncColor}{$\triangleright$\enspace\texttt{InfoLattice\index{InfoLattice@\texttt{InfoLattice}}
\label{InfoLattice}
}\hfill{\scriptsize (info class)}}\\


 is the information class used by the cyclic extension methods for subgroup
lattice calculations. }

 }

  
\section{\textcolor{Chapter }{Specific Methods for Subgroup Lattice Computations}}\label{Specific Methods for Subgroup Lattice Computations}
\logpage{[ 39, 21, 0 ]}
\hyperdef{L}{X85E613D57F28AEFF}{}
{
  

\subsection{\textcolor{Chapter }{LatticeByCyclicExtension}}
\logpage{[ 39, 21, 1 ]}\nobreak
\hyperdef{L}{X86462A567DDBA6BC}{}
{\noindent\textcolor{FuncColor}{$\triangleright$\enspace\texttt{LatticeByCyclicExtension({\mdseries\slshape G[, func[, noperf]]})\index{LatticeByCyclicExtension@\texttt{LatticeByCyclicExtension}}
\label{LatticeByCyclicExtension}
}\hfill{\scriptsize (function)}}\\


 computes the lattice of \mbox{\texttt{\mdseries\slshape G}} using the cyclic extension algorithm. If the function \mbox{\texttt{\mdseries\slshape func}} is given, the algorithm will discard all subgroups not fulfilling \mbox{\texttt{\mdseries\slshape func}} (and will also not extend them), returning a partial lattice. This can be
useful to compute only subgroups with certain properties. Note however that
this will \emph{not} necessarily yield all subgroups that fulfill \mbox{\texttt{\mdseries\slshape func}}, but the subgroups whose subgroups are used for the construction must also
fulfill \mbox{\texttt{\mdseries\slshape func}} as well. (In fact the filter \mbox{\texttt{\mdseries\slshape func}} will simply discard subgroups in the cyclic extension algorithm. Therefore the
trivial subgroup will always be included.) Also note, that for such a partial
lattice maximality/minimality inclusion relations cannot be computed. (If \mbox{\texttt{\mdseries\slshape func}} is a list of length 2, its first entry is such a discarding function, the
second a function for discarding zuppos.) 

 The cyclic extension algorithm requires the perfect subgroups of \mbox{\texttt{\mdseries\slshape G}}. However \textsf{GAP} cannot analyze the function \mbox{\texttt{\mdseries\slshape func}} for its implication but can only apply it. If it is known that \mbox{\texttt{\mdseries\slshape func}} implies solvability, the computation of the perfect subgroups can be avoided
by giving a third parameter \mbox{\texttt{\mdseries\slshape noperf}} set to \texttt{true}. 

 
\begin{Verbatim}[commandchars=@|A,fontsize=\small,frame=single,label=Example]
  @gapprompt|gap>A @gapinput|g:=WreathProduct(Group((1,2,3),(1,2)),Group((1,2,3,4)));;A
  @gapprompt|gap>A @gapinput|l:=LatticeByCyclicExtension(g,function(G)A
  @gapprompt|>A @gapinput|return Size(G) in [1,2,3,6];end);A
  <subgroup lattice of <permutation group of size 5184 with
  9 generators>, 47 classes,
  2628 subgroups, restricted under further condition l!.func>
\end{Verbatim}
 

 The total number of classes in this example is much bigger, as the following
example shows: 
\begin{Verbatim}[commandchars=!@|,fontsize=\small,frame=single,label=Example]
  !gapprompt@gap>| !gapinput@LatticeSubgroups(g);|
  <subgroup lattice of <permutation group of size 5184 with
  9 generators>, 566 classes, 27134 subgroups>
\end{Verbatim}
 }

 

\subsection{\textcolor{Chapter }{InvariantSubgroupsElementaryAbelianGroup}}
\logpage{[ 39, 21, 2 ]}\nobreak
\hyperdef{L}{X78918D83835A0EDF}{}
{\noindent\textcolor{FuncColor}{$\triangleright$\enspace\texttt{InvariantSubgroupsElementaryAbelianGroup({\mdseries\slshape G, homs[, dims]})\index{InvariantSubgroupsElementaryAbelianGroup@\texttt{Invariant}\-\texttt{Subgroups}\-\texttt{Elementary}\-\texttt{Abelian}\-\texttt{Group}}
\label{InvariantSubgroupsElementaryAbelianGroup}
}\hfill{\scriptsize (function)}}\\


 Let \mbox{\texttt{\mdseries\slshape G}} be an elementary abelian group and \mbox{\texttt{\mdseries\slshape homs}} be a set of automorphisms of \mbox{\texttt{\mdseries\slshape G}}. Then this function computes all subspaces of \mbox{\texttt{\mdseries\slshape G}} which are invariant under all automorphisms in \mbox{\texttt{\mdseries\slshape homs}}. When considering \mbox{\texttt{\mdseries\slshape G}} as a module for the algebra generated by \mbox{\texttt{\mdseries\slshape homs}}, these are all submodules. If \mbox{\texttt{\mdseries\slshape homs}} is empty, it computes all subgroups. If the optional parameter \mbox{\texttt{\mdseries\slshape dims}} is given, only submodules of this dimension are computed. 
\begin{Verbatim}[commandchars=!@|,fontsize=\small,frame=single,label=Example]
  !gapprompt@gap>| !gapinput@g:=Group((1,2,3),(4,5,6),(7,8,9));|
  Group([ (1,2,3), (4,5,6), (7,8,9) ])
  !gapprompt@gap>| !gapinput@hom:=GroupHomomorphismByImages(g,g,[(1,2,3),(4,5,6),(7,8,9)],|
  !gapprompt@>| !gapinput@[(7,8,9),(1,2,3),(4,5,6)]);|
  [ (1,2,3), (4,5,6), (7,8,9) ] -> [ (7,8,9), (1,2,3), (4,5,6) ]
  !gapprompt@gap>| !gapinput@u:=InvariantSubgroupsElementaryAbelianGroup(g,[hom]);|
  [ Group([ (7,8,9), (4,5,6), (1,2,3) ]),
    Group([ (1,3,2)(7,8,9), (1,3,2)(4,5,6) ]),
    Group([ (1,2,3)(4,5,6)(7,8,9) ]), Group(()) ]
\end{Verbatim}
 }

 

\subsection{\textcolor{Chapter }{SubgroupsSolvableGroup}}
\logpage{[ 39, 21, 3 ]}\nobreak
\hyperdef{L}{X7AD7804A803910AC}{}
{\noindent\textcolor{FuncColor}{$\triangleright$\enspace\texttt{SubgroupsSolvableGroup({\mdseries\slshape G[, opt]})\index{SubgroupsSolvableGroup@\texttt{SubgroupsSolvableGroup}}
\label{SubgroupsSolvableGroup}
}\hfill{\scriptsize (function)}}\\


 This function (implementing the algorithm published in \cite{Hulpke99}) computes subgroups of a solvable group \mbox{\texttt{\mdseries\slshape G}}, using the homomorphism principle. It returns a list of representatives up to \mbox{\texttt{\mdseries\slshape G}}\texttt{\symbol{45}}conjugacy. 

 The optional argument \mbox{\texttt{\mdseries\slshape opt}} is a record, which may be used to suggest restrictions on the subgroups
computed. The following record components of \mbox{\texttt{\mdseries\slshape opt}} are recognized and have the following effects. Note that all of the following
restrictions to subgroups with particular properties are only used to speed up
the calculation, but the result might still contain subgroups (that had to be
computed in any case) that do not satisfy the properties. If this is not
desired, the calculation must be followed by an explicit test for the desired
properties (which is not done by default, as it would be a general slowdown).
The function guarantees that representatives of all subgroups that satisfy the
properties are found, i.e. there can be only false positives. 
\begin{description}
\item[{\texttt{actions}}]  must be a list of automorphisms of \mbox{\texttt{\mdseries\slshape G}}. If given, only groups which are invariant under all these automorphisms are
computed. The algorithm must know the normalizer in \mbox{\texttt{\mdseries\slshape G}} of the group generated by \texttt{actions} (defined formally by embedding in the semidirect product of \mbox{\texttt{\mdseries\slshape G}} with \mbox{\texttt{\mdseries\slshape actions}}). This can be given in the component \texttt{funcnorm} and will be computed if this component is not given. 
\item[{\texttt{normal}}]  if set to \texttt{true} only normal subgroups are guaranteed to be returned (though some of the
returned subgroups might still be not normal). 
\item[{\texttt{consider}}]  a function to restrict the groups computed. This must be a function of five
parameters, $C$, $A$, $N$, $B$, $M$, that are interpreted as follows: The arguments are subgroups of a factor $F$ of \mbox{\texttt{\mdseries\slshape G}} in the relation $F \geq C > A > N > B > M$. $N$ and $M$ are normal subgroups. $C$ is the full preimage of the normalizer of $A/N$ in $F/N$. When computing modulo $M$ and looking for subgroups $U$ such that $U \cap N = B$ and $\langle U, N \rangle = A$, this function is called. If it returns \texttt{false} then all potential groups $U$ (and therefore all groups later arising from them) are disregarded. This can
be used for example to compute only subgroups of certain sizes. 

 (\emph{This is just a restriction to speed up computations. The function may still
return (invariant) subgroups which don't fulfill this condition!}) This parameter is used to permit calculations of some subgroups if the set
of all subgroups would be too large to handle. 

 The actual groups $C$, $A$, $N$ and $B$ which are passed to this function are not necessarily subgroups of \mbox{\texttt{\mdseries\slshape G}} but might be subgroups of a proper factor group $F = \mbox{\texttt{\mdseries\slshape G}}/H$. Therefore the \texttt{consider} function may not relate the parameter groups to \mbox{\texttt{\mdseries\slshape G}}. 
\item[{\texttt{retnorm}}]  if set to \texttt{true} the function not only returns a list \texttt{subs} of subgroups but also a corresponding list \texttt{norms} of normalizers in the form \texttt{[ subs, norms ]}. 
\item[{\texttt{series}}]  is an elementary abelian series of \mbox{\texttt{\mdseries\slshape G}} which will be used for the computation. 
\item[{\texttt{groups}}]  is a list of groups to seed the calculation. Only subgroups of these groups
are constructed. 
\end{description}
 
\begin{Verbatim}[commandchars=!@|,fontsize=\small,frame=single,label=Example]
  !gapprompt@gap>| !gapinput@g:=Group((1,2,3),(1,2),(4,5,6),(4,5),(7,8,9),(7,8));|
  Group([ (1,2,3), (1,2), (4,5,6), (4,5), (7,8,9), (7,8) ])
  !gapprompt@gap>| !gapinput@hom:=GroupHomomorphismByImages(g,g,|
  !gapprompt@>| !gapinput@[(1,2,3),(1,2),(4,5,6),(4,5),(7,8,9),(7,8)],|
  !gapprompt@>| !gapinput@[(4,5,6),(4,5),(7,8,9),(7,8),(1,2,3),(1,2)]);|
  [ (1,2,3), (1,2), (4,5,6), (4,5), (7,8,9), (7,8) ] ->
  [ (4,5,6), (4,5), (7,8,9), (7,8), (1,2,3), (1,2) ]
  !gapprompt@gap>| !gapinput@l:=SubgroupsSolvableGroup(g,rec(actions:=[hom]));;|
  !gapprompt@gap>| !gapinput@SortedList(List(l,Size));|
  [ 1, 2, 3, 4, 6, 8, 9, 18, 27, 54, 108, 216 ]
  !gapprompt@gap>| !gapinput@Length(ConjugacyClassesSubgroups(g)); # to compare|
  162
\end{Verbatim}
 }

 

\subsection{\textcolor{Chapter }{SizeConsiderFunction}}
\logpage{[ 39, 21, 4 ]}\nobreak
\hyperdef{L}{X7F60BBB8874DFE40}{}
{\noindent\textcolor{FuncColor}{$\triangleright$\enspace\texttt{SizeConsiderFunction({\mdseries\slshape size})\index{SizeConsiderFunction@\texttt{SizeConsiderFunction}}
\label{SizeConsiderFunction}
}\hfill{\scriptsize (function)}}\\


 This function returns a function \texttt{consider} of five arguments that can be used in \texttt{SubgroupsSolvableGroup} (\ref{SubgroupsSolvableGroup}) for the option \texttt{consider} to compute subgroups whose sizes are divisible by \mbox{\texttt{\mdseries\slshape size}}. 
\begin{Verbatim}[commandchars=!@|,fontsize=\small,frame=single,label=Example]
  !gapprompt@gap>| !gapinput@l:=SubgroupsSolvableGroup(g,rec(actions:=[hom],|
  !gapprompt@>| !gapinput@consider:=SizeConsiderFunction(6)));;|
  !gapprompt@gap>| !gapinput@SortedList(List(l,Size));|
  [ 1, 3, 6, 9, 18, 27, 54, 108, 216 ]
\end{Verbatim}
 

 This example shows that in general the \texttt{consider} function does not provide a perfect filter. It is guaranteed that all
subgroups fulfilling the condition are returned, but not all subgroups
returned necessarily fulfill the condition. }

 

\subsection{\textcolor{Chapter }{ExactSizeConsiderFunction}}
\logpage{[ 39, 21, 5 ]}\nobreak
\hyperdef{L}{X833C51BD7E7812C4}{}
{\noindent\textcolor{FuncColor}{$\triangleright$\enspace\texttt{ExactSizeConsiderFunction({\mdseries\slshape size})\index{ExactSizeConsiderFunction@\texttt{ExactSizeConsiderFunction}}
\label{ExactSizeConsiderFunction}
}\hfill{\scriptsize (function)}}\\


 This function returns a function \texttt{consider} of five arguments that can be used in \texttt{SubgroupsSolvableGroup} (\ref{SubgroupsSolvableGroup}) for the option \texttt{consider} to compute subgroups whose sizes are exactly \mbox{\texttt{\mdseries\slshape size}}. 
\begin{Verbatim}[commandchars=!@|,fontsize=\small,frame=single,label=Example]
  !gapprompt@gap>| !gapinput@l:=SubgroupsSolvableGroup(g,rec(actions:=[hom],|
  !gapprompt@>| !gapinput@consider:=ExactSizeConsiderFunction(6)));;|
  !gapprompt@gap>| !gapinput@SortedList(List(l,Size));|
  [ 1, 3, 6, 9, 27, 54, 108, 216 ]
\end{Verbatim}
 

 Again, the \texttt{consider} function does not provide a perfect filter. It is guaranteed that all
subgroups fulfilling the condition are returned, but not all subgroups
returned necessarily fulfill the condition. }

 

\subsection{\textcolor{Chapter }{InfoPcSubgroup}}
\logpage{[ 39, 21, 6 ]}\nobreak
\hyperdef{L}{X7A2C774B7CFF3E07}{}
{\noindent\textcolor{FuncColor}{$\triangleright$\enspace\texttt{InfoPcSubgroup\index{InfoPcSubgroup@\texttt{InfoPcSubgroup}}
\label{InfoPcSubgroup}
}\hfill{\scriptsize (info class)}}\\


 Information function for the subgroup lattice functions using pcgs. }

 }

  
\section{\textcolor{Chapter }{Special Generating Sets}}\label{Special Generating Sets}
\logpage{[ 39, 22, 0 ]}
\hyperdef{L}{X79F894537D526B61}{}
{
  

\subsection{\textcolor{Chapter }{GeneratorsSmallest}}
\logpage{[ 39, 22, 1 ]}\nobreak
\hyperdef{L}{X82FD78AF7F80A0E2}{}
{\noindent\textcolor{FuncColor}{$\triangleright$\enspace\texttt{GeneratorsSmallest({\mdseries\slshape G})\index{GeneratorsSmallest@\texttt{GeneratorsSmallest}}
\label{GeneratorsSmallest}
}\hfill{\scriptsize (attribute)}}\\


 returns a ``smallest'' generating set for the group \mbox{\texttt{\mdseries\slshape G}}. This is the lexicographically (using \textsf{GAP}s order of group elements) smallest list $l$ of elements of \mbox{\texttt{\mdseries\slshape G}} such that $G = \langle l \rangle$ and $l_i \not \in \langle l_1, \ldots, l_{{i-1}} \rangle$ (in particular $l_1$ is not the identity element of the group). The comparison of two groups via
lexicographic comparison of their sorted element lists yields the same
relation as lexicographic comparison of their smallest generating sets. 
\begin{Verbatim}[commandchars=!@|,fontsize=\small,frame=single,label=Example]
  !gapprompt@gap>| !gapinput@g:=SymmetricGroup(4);;|
  !gapprompt@gap>| !gapinput@GeneratorsSmallest(g);|
  [ (3,4), (2,3), (1,2) ]
\end{Verbatim}
 }

 

\subsection{\textcolor{Chapter }{LargestElementGroup}}
\logpage{[ 39, 22, 2 ]}\nobreak
\hyperdef{L}{X7A258CCF79552198}{}
{\noindent\textcolor{FuncColor}{$\triangleright$\enspace\texttt{LargestElementGroup({\mdseries\slshape G})\index{LargestElementGroup@\texttt{LargestElementGroup}}
\label{LargestElementGroup}
}\hfill{\scriptsize (attribute)}}\\


 returns the largest element of \mbox{\texttt{\mdseries\slshape G}} with respect to the ordering \texttt{{\textless}} of the elements family. }

 

\subsection{\textcolor{Chapter }{MinimalGeneratingSet}}
\logpage{[ 39, 22, 3 ]}\nobreak
\hyperdef{L}{X81D15723804771E2}{}
{\noindent\textcolor{FuncColor}{$\triangleright$\enspace\texttt{MinimalGeneratingSet({\mdseries\slshape G})\index{MinimalGeneratingSet@\texttt{MinimalGeneratingSet}}
\label{MinimalGeneratingSet}
}\hfill{\scriptsize (attribute)}}\\


 returns a generating set of \mbox{\texttt{\mdseries\slshape G}} of minimal possible length. 

 Note that only methods for finite groups, solvable groups, or finitely
generated nilpotent groups are available (the latter through the \textsf{Polycyclic} package) and that calculations for nonsolvable finite groups of higher rank
can be expensive. 

 If you do not really need a minimal generating set, but are satisfied with
getting a reasonably small set of generators, you better use \texttt{SmallGeneratingSet} (\ref{SmallGeneratingSet}). 

 Information about the minimal generating sets of the finite simple groups of
order less than $10^6$ can be found in \cite{MY79}. See also the package \textsf{AtlasRep}. 
\begin{Verbatim}[commandchars=!@|,fontsize=\small,frame=single,label=Example]
  !gapprompt@gap>| !gapinput@MinimalGeneratingSet(g);|
  [ (2,4,3), (1,4,2,3) ]
\end{Verbatim}
 }

 

\subsection{\textcolor{Chapter }{SmallGeneratingSet}}
\logpage{[ 39, 22, 4 ]}\nobreak
\hyperdef{L}{X814DBABC878D5232}{}
{\noindent\textcolor{FuncColor}{$\triangleright$\enspace\texttt{SmallGeneratingSet({\mdseries\slshape G})\index{SmallGeneratingSet@\texttt{SmallGeneratingSet}}
\label{SmallGeneratingSet}
}\hfill{\scriptsize (attribute)}}\\


 returns a generating set of \mbox{\texttt{\mdseries\slshape G}} which has few elements. As neither irredundancy, nor minimal length is proven
it runs much faster than \texttt{MinimalGeneratingSet} (\ref{MinimalGeneratingSet}). It can be used whenever a short generating set is desired which not
necessarily needs to be optimal. 
\begin{Verbatim}[commandchars=!@|,fontsize=\small,frame=single,label=Example]
  !gapprompt@gap>| !gapinput@SmallGeneratingSet(g);|
  [ (1,2,3,4), (1,2) ]
\end{Verbatim}
 }

 

\subsection{\textcolor{Chapter }{IndependentGeneratorsOfAbelianGroup}}
\logpage{[ 39, 22, 5 ]}\nobreak
\hyperdef{L}{X7D1574457B152333}{}
{\noindent\textcolor{FuncColor}{$\triangleright$\enspace\texttt{IndependentGeneratorsOfAbelianGroup({\mdseries\slshape A})\index{IndependentGeneratorsOfAbelianGroup@\texttt{IndependentGeneratorsOfAbelianGroup}}
\label{IndependentGeneratorsOfAbelianGroup}
}\hfill{\scriptsize (attribute)}}\\


 returns a list of generators $a_1, a_2, \ldots$ of prime power order or infinite order of the abelian group \mbox{\texttt{\mdseries\slshape A}} such that \mbox{\texttt{\mdseries\slshape A}} is the direct product of the cyclic groups generated by the $a_i$. The list of orders of the returned generators must match the result of \texttt{AbelianInvariants} (\ref{AbelianInvariants}) (taking into account that zero and \texttt{infinity} (\ref{infinity}) are identified). 
\begin{Verbatim}[commandchars=!@|,fontsize=\small,frame=single,label=Example]
  !gapprompt@gap>| !gapinput@g:=AbelianGroup(IsPermGroup,[15,14,22,78]);;|
  !gapprompt@gap>| !gapinput@List(IndependentGeneratorsOfAbelianGroup(g),Order);|
  [ 2, 2, 2, 3, 3, 5, 7, 11, 13 ]
  !gapprompt@gap>| !gapinput@AbelianInvariants(g);|
  [ 2, 2, 2, 3, 3, 5, 7, 11, 13 ]
\end{Verbatim}
 }

 

\subsection{\textcolor{Chapter }{IndependentGeneratorExponents}}
\logpage{[ 39, 22, 6 ]}\nobreak
\hyperdef{L}{X86F835DA8264A0CE}{}
{\noindent\textcolor{FuncColor}{$\triangleright$\enspace\texttt{IndependentGeneratorExponents({\mdseries\slshape G, g})\index{IndependentGeneratorExponents@\texttt{IndependentGeneratorExponents}}
\label{IndependentGeneratorExponents}
}\hfill{\scriptsize (operation)}}\\


 For an abelian group \mbox{\texttt{\mdseries\slshape G}}, with \texttt{IndependentGeneratorsOfAbelianGroup} (\ref{IndependentGeneratorsOfAbelianGroup}) value the list $[ a_1, \ldots, a_n ]$, this operation returns the exponent vector $[ e_1, \ldots, e_n ]$ to represent $\mbox{\texttt{\mdseries\slshape g}} = \prod_i a_i^{{e_i}}$. 
\begin{Verbatim}[commandchars=!@|,fontsize=\small,frame=single,label=Example]
  !gapprompt@gap>| !gapinput@g := AbelianGroup([16,9,625]);;|
  !gapprompt@gap>| !gapinput@gens := IndependentGeneratorsOfAbelianGroup(g);;|
  !gapprompt@gap>| !gapinput@List(gens, Order);|
  [ 9, 16, 625 ]
  !gapprompt@gap>| !gapinput@AbelianInvariants(g);|
  [ 9, 16, 625 ]
  !gapprompt@gap>| !gapinput@r:=gens[1]^4*gens[2]^12*gens[3]^128;;|
  !gapprompt@gap>| !gapinput@IndependentGeneratorExponents(g,r);|
  [ 4, 12, 128 ]
\end{Verbatim}
 }

 }

  
\section{\textcolor{Chapter }{1\texttt{\symbol{45}}Cohomology}}\label{1-Cohomology}
\logpage{[ 39, 23, 0 ]}
\hyperdef{L}{X7CA0B6A27E0BE6B8}{}
{
  \index{one cohomology} \index{cohomology} \index{cocycles} Let $G$ be a finite group and $M$ an elementary abelian normal $p$\texttt{\symbol{45}}subgroup of $G$. Then the group of 1\texttt{\symbol{45}}cocycles $Z^1( G/M, M )$ is defined as 
\[ Z^1(G/M, M) = \{ \gamma: G/M \rightarrow M \mid \forall g_1, g_2 \in G :
\gamma(g_1 M \cdot g_2 M ) = \gamma(g_1 M)^{{g_2}} \cdot \gamma(g_2 M) \} \]
 and is a $GF(p)$\texttt{\symbol{45}}vector space. 

 The group of 1\texttt{\symbol{45}}coboundaries $B^1( G/M, M )$ is defined as 
\[ B^1(G/M, M) = \{ \gamma : G/M \rightarrow M \mid \exists m \in M \forall g \in
G : \gamma(gM) = (m^{{-1}})^g \cdot m \} \]
 It also is a $GF(p)$\texttt{\symbol{45}}vector space. 

 Let $\alpha$ be the isomorphism of $M$ into a row vector space ${\cal W}$ and $(g_1, \ldots, g_l)$ representatives for a generating set of $G/M$. Then there exists a monomorphism $\beta$ of $Z^1( G/M, M )$ in the $l$\texttt{\symbol{45}}fold direct sum of ${\cal W}$, such that $\beta( \gamma ) = ( \alpha( \gamma(g_1 M) ),\ldots, \alpha( \gamma(g_l M) ) )$ for every $\gamma \in Z^1( G/M, M )$. 
\subsection{\textcolor{Chapter }{OneCocycles}}\logpage{[ 39, 23, 1 ]}
\hyperdef{L}{X847BEC137A49BAF4}{}
{
\noindent\textcolor{FuncColor}{$\triangleright$\enspace\texttt{OneCocycles({\mdseries\slshape G, M})\index{OneCocycles@\texttt{OneCocycles}!for two groups}
\label{OneCocycles:for two groups}
}\hfill{\scriptsize (function)}}\\
\noindent\textcolor{FuncColor}{$\triangleright$\enspace\texttt{OneCocycles({\mdseries\slshape G, mpcgs})\index{OneCocycles@\texttt{OneCocycles}!for a group and a pcgs}
\label{OneCocycles:for a group and a pcgs}
}\hfill{\scriptsize (function)}}\\
\noindent\textcolor{FuncColor}{$\triangleright$\enspace\texttt{OneCocycles({\mdseries\slshape gens, M})\index{OneCocycles@\texttt{OneCocycles}!for generators and a group}
\label{OneCocycles:for generators and a group}
}\hfill{\scriptsize (function)}}\\
\noindent\textcolor{FuncColor}{$\triangleright$\enspace\texttt{OneCocycles({\mdseries\slshape gens, mpcgs})\index{OneCocycles@\texttt{OneCocycles}!for generators and a pcgs}
\label{OneCocycles:for generators and a pcgs}
}\hfill{\scriptsize (function)}}\\


 Computes the group of 1\texttt{\symbol{45}}cocycles $Z^1(\mbox{\texttt{\mdseries\slshape G}}/\mbox{\texttt{\mdseries\slshape M}},\mbox{\texttt{\mdseries\slshape M}})$. The normal subgroup \mbox{\texttt{\mdseries\slshape M}} may be given by a (Modulo)Pcgs \mbox{\texttt{\mdseries\slshape mpcgs}}. In this case the whole calculation is performed modulo the normal subgroup
defined by \texttt{DenominatorOfModuloPcgs(\mbox{\texttt{\mdseries\slshape mpcgs}})} (see{\nobreakspace}\ref{Polycyclic Generating Systems}). Similarly the group \mbox{\texttt{\mdseries\slshape G}} may instead be specified by a set of elements \mbox{\texttt{\mdseries\slshape gens}} that are representatives for a generating system for the factor group \mbox{\texttt{\mdseries\slshape G}}/\mbox{\texttt{\mdseries\slshape M}}. If this is done the 1\texttt{\symbol{45}}cocycles are computed with respect
to these generators (otherwise the routines try to select suitable generators
themselves). The current version of the code assumes that \mbox{\texttt{\mdseries\slshape G}} is a permutation group or a pc group. }

 

\subsection{\textcolor{Chapter }{OneCoboundaries}}
\logpage{[ 39, 23, 2 ]}\nobreak
\hyperdef{L}{X7E6438D5834ACCDA}{}
{\noindent\textcolor{FuncColor}{$\triangleright$\enspace\texttt{OneCoboundaries({\mdseries\slshape G, M})\index{OneCoboundaries@\texttt{OneCoboundaries}}
\label{OneCoboundaries}
}\hfill{\scriptsize (function)}}\\


 computes the group of 1\texttt{\symbol{45}}coboundaries. Syntax of input and
output otherwise is the same as with \texttt{OneCocycles} (\ref{OneCocycles:for two groups}) except that entries that refer to cocycles are not computed. 

 The operations \texttt{OneCocycles} (\ref{OneCocycles:for two groups}) and \texttt{OneCoboundaries} return a record with (at least) the components: 

 
\begin{description}
\item[{\texttt{generators}}]  Is a list of representatives for a generating set of \mbox{\texttt{\mdseries\slshape G}}/\mbox{\texttt{\mdseries\slshape M}}. Cocycles are represented with respect to these generators. 
\item[{\texttt{oneCocycles}}]  A space of row vectors over GF($p$), representing $Z^1$. The vectors are represented in dimension $a \cdot b$ where $a$ is the length of \texttt{generators} and $p^b$ the size of \mbox{\texttt{\mdseries\slshape M}}. 
\item[{\texttt{oneCoboundaries}}]  A space of row vectors that represents $B^1$. 
\item[{\texttt{cocycleToList}}]  is a function to convert a cocycle (a row vector in \texttt{oneCocycles}) to a corresponding list of elements of \mbox{\texttt{\mdseries\slshape M}}. 
\item[{\texttt{listToCocycle}}]  is a function to convert a list of elements of \mbox{\texttt{\mdseries\slshape M}} to a cocycle. 
\item[{\texttt{isSplitExtension}}]  indicates whether \mbox{\texttt{\mdseries\slshape G}} splits over \mbox{\texttt{\mdseries\slshape M}}. The following components are only bound if the extension splits. Note that
if \mbox{\texttt{\mdseries\slshape M}} is given by a modulo pcgs all subgroups are given as subgroups of \mbox{\texttt{\mdseries\slshape G}} by generators corresponding to \texttt{generators} and thus may not contain the denominator of the modulo pcgs. In this case
taking the closure with this denominator will give the full preimage of the
complement in the factor group. 
\item[{\texttt{complement}}]  One complement to \mbox{\texttt{\mdseries\slshape M}} in \mbox{\texttt{\mdseries\slshape G}}. 
\item[{\texttt{cocycleToComplement( cyc )}}]  is a function that takes a cocycle from \texttt{oneCocycles} and returns the corresponding complement to \mbox{\texttt{\mdseries\slshape M}} in \mbox{\texttt{\mdseries\slshape G}} (with respect to the fixed complement \texttt{complement}). 
\item[{\texttt{complementToCocycle(\mbox{\texttt{\mdseries\slshape U}})}}]  is a function that takes a complement and returns the corresponding cocycle. 
\end{description}
 

 If the factor \mbox{\texttt{\mdseries\slshape G}}/\mbox{\texttt{\mdseries\slshape M}} is given by a (modulo) pcgs \mbox{\texttt{\mdseries\slshape gens}} then special methods are used that compute a presentation for the factor
implicitly from the pcgs. 

 Note that the groups of 1\texttt{\symbol{45}}cocycles and
1\texttt{\symbol{45}}coboundaries are not groups in the sense of \texttt{Group} (\ref{Group:for several generators}) for \textsf{GAP} but vector spaces. 

 
\begin{Verbatim}[commandchars=!@|,fontsize=\small,frame=single,label=Example]
  !gapprompt@gap>| !gapinput@g:=Group((1,2,3,4),(1,2));;|
  !gapprompt@gap>| !gapinput@n:=Group((1,2)(3,4),(1,3)(2,4));;|
  !gapprompt@gap>| !gapinput@oc:=OneCocycles(g,n);|
  rec( cocycleToComplement := function( c ) ... end,
    cocycleToList := function( c ) ... end,
    complement := Group([ (3,4), (2,4,3) ]),
    complementGens := [ (3,4), (2,4,3) ],
    complementToCocycle := function( K ) ... end,
    factorGens := [ (3,4), (2,4,3) ], generators := [ (3,4), (2,4,3) ],
    isSplitExtension := true, listToCocycle := function( L ) ... end,
    oneCoboundaries := <vector space over GF(2), with 2 generators>,
    oneCocycles := <vector space over GF(2), with 2 generators> )
  !gapprompt@gap>| !gapinput@oc.cocycleToList([ 0*Z(2), Z(2)^0, 0*Z(2), Z(2)^0 ]);|
  [ (1,2)(3,4), (1,2)(3,4) ]
  !gapprompt@gap>| !gapinput@oc.listToCocycle([(),(1,3)(2,4)]) = Z(2) * [ 0, 0, 1, 0];|
  true
  !gapprompt@gap>| !gapinput@oc.cocycleToComplement([ 0*Z(2), 0*Z(2), Z(2)^0, 0*Z(2) ]);|
  Group([ (3,4), (1,3,4) ])
  !gapprompt@gap>| !gapinput@oc.complementToCocycle(Group((1,2,4),(1,4))) = Z(2) * [ 0, 1, 1, 1 ];|
  true
\end{Verbatim}
 

 The factor group $H^1(\mbox{\texttt{\mdseries\slshape G}}/\mbox{\texttt{\mdseries\slshape M}}, \mbox{\texttt{\mdseries\slshape M}}) = Z^1(\mbox{\texttt{\mdseries\slshape G}}/\mbox{\texttt{\mdseries\slshape M}}, \mbox{\texttt{\mdseries\slshape M}}) / B^1(\mbox{\texttt{\mdseries\slshape G}}/\mbox{\texttt{\mdseries\slshape M}}, \mbox{\texttt{\mdseries\slshape M}})$ is called the first cohomology group. Currently there is no function which
explicitly computes this group. The easiest way to represent it is as a vector
space complement to $B^1$ in $Z^1$. 

 If the only purpose of the calculation of $H^1$ is the determination of complements it might be desirable to stop calculations
once it is known that the extension cannot split. This can be achieved via the
more technical function \texttt{OCOneCocycles} (\ref{OCOneCocycles}). }

 

\subsection{\textcolor{Chapter }{OCOneCocycles}}
\logpage{[ 39, 23, 3 ]}\nobreak
\hyperdef{L}{X80400ABD7F40FAA0}{}
{\noindent\textcolor{FuncColor}{$\triangleright$\enspace\texttt{OCOneCocycles({\mdseries\slshape ocr, onlySplit})\index{OCOneCocycles@\texttt{OCOneCocycles}}
\label{OCOneCocycles}
}\hfill{\scriptsize (function)}}\\


 is the more technical function to compute 1\texttt{\symbol{45}}cocycles. It
takes a record \mbox{\texttt{\mdseries\slshape ocr}} as first argument which must contain at least the components \texttt{group} for the group and \texttt{modulePcgs} for a (modulo) pcgs of the module. This record will also be returned with
components as described under \texttt{OneCocycles} (\ref{OneCocycles:for two groups}) (with the exception of \texttt{isSplitExtension} which is indicated by the existence of a \texttt{complement}) but components such as \texttt{oneCoboundaries} will only be computed if not already present. 

 If \mbox{\texttt{\mdseries\slshape onlySplit}} is \texttt{true}, \texttt{OCOneCocycles} returns \texttt{false} as soon as possible if the extension does not split. }

 

\subsection{\textcolor{Chapter }{ComplementClassesRepresentativesEA}}
\logpage{[ 39, 23, 4 ]}\nobreak
\hyperdef{L}{X811E1CF07DABE924}{}
{\noindent\textcolor{FuncColor}{$\triangleright$\enspace\texttt{ComplementClassesRepresentativesEA({\mdseries\slshape G, N})\index{ComplementClassesRepresentativesEA@\texttt{ComplementClassesRepresentativesEA}}
\label{ComplementClassesRepresentativesEA}
}\hfill{\scriptsize (function)}}\\


 computes complement classes to an elementary abelian normal subgroup \mbox{\texttt{\mdseries\slshape N}} via 1\texttt{\symbol{45}}Cohomology. Normally, a user program should call \texttt{ComplementClassesRepresentatives} (\ref{ComplementClassesRepresentatives}) instead, which also works for a solvable (not necessarily elementary abelian) \mbox{\texttt{\mdseries\slshape N}}. }

 

\subsection{\textcolor{Chapter }{InfoCoh}}
\logpage{[ 39, 23, 5 ]}\nobreak
\hyperdef{L}{X8199B1D27D487897}{}
{\noindent\textcolor{FuncColor}{$\triangleright$\enspace\texttt{InfoCoh\index{InfoCoh@\texttt{InfoCoh}}
\label{InfoCoh}
}\hfill{\scriptsize (info class)}}\\


 The info class for the cohomology calculations is \texttt{InfoCoh}. }

     }

  
\section{\textcolor{Chapter }{Schur Covers and Multipliers}}\label{Schur Covers and Multipliers}
\logpage{[ 39, 24, 0 ]}
\hyperdef{L}{X80A4B0F282977074}{}
{
  Additional attributes and properties of a group can be derived from computing
its Schur cover. For example, if $G$ is a finitely presented group, the derived subgroup of a Schur cover of $G$ is invariant and isomorphic to the \texttt{NonabelianExteriorSquare} (\ref{NonabelianExteriorSquare}) value of $G$, see \cite{BJR87}. \index{Darstellungsgruppe!see EpimorphismSchurCover} 

\subsection{\textcolor{Chapter }{EpimorphismSchurCover}}
\logpage{[ 39, 24, 1 ]}\nobreak
\hyperdef{L}{X7F619DDA7DD6C43B}{}
{\noindent\textcolor{FuncColor}{$\triangleright$\enspace\texttt{EpimorphismSchurCover({\mdseries\slshape G[, pl]})\index{EpimorphismSchurCover@\texttt{EpimorphismSchurCover}}
\label{EpimorphismSchurCover}
}\hfill{\scriptsize (attribute)}}\\


 returns an epimorphism $epi$ from a group $D$ onto \mbox{\texttt{\mdseries\slshape G}}. The group $D$ is one (of possibly several) Schur covers of \mbox{\texttt{\mdseries\slshape G}}. The group $D$ can be obtained as the \texttt{Source} (\ref{Source}) value of \mbox{\texttt{\mdseries\slshape epi}}. The kernel of $epi$ is the Schur multiplier of \mbox{\texttt{\mdseries\slshape G}}. If \mbox{\texttt{\mdseries\slshape pl}} is given as a list of primes, only the multiplier part for these primes is
realized. At the moment, $D$ is represented as a finitely presented group. }

 

\subsection{\textcolor{Chapter }{SchurCover}}
\logpage{[ 39, 24, 2 ]}\nobreak
\hyperdef{L}{X7DD1E37987612042}{}
{\noindent\textcolor{FuncColor}{$\triangleright$\enspace\texttt{SchurCover({\mdseries\slshape G})\index{SchurCover@\texttt{SchurCover}}
\label{SchurCover}
}\hfill{\scriptsize (attribute)}}\\


 returns one (of possibly several) Schur covers of the group \mbox{\texttt{\mdseries\slshape G}}. 

 At the moment this cover is represented as a finitely presented group and \texttt{IsomorphismPermGroup} (\ref{IsomorphismPermGroup}) would be needed to convert it to a permutation group. 

 If also the relation to \mbox{\texttt{\mdseries\slshape G}} is needed, \texttt{EpimorphismSchurCover} (\ref{EpimorphismSchurCover}) should be used. 
\begin{Verbatim}[commandchars=!@|,fontsize=\small,frame=single,label=Example]
  !gapprompt@gap>| !gapinput@g:=Group((1,2,3,4),(1,2));;|
  !gapprompt@gap>| !gapinput@epi:=EpimorphismSchurCover(g);|
  [ F1, F2, F3 ] -> [ (1,2), (2,3), (3,4) ]
  !gapprompt@gap>| !gapinput@Size(Source(epi));|
  48
  !gapprompt@gap>| !gapinput@f:=FreeGroup("a","b");;|
  !gapprompt@gap>| !gapinput@g:=f/ParseRelators(f,"a2,b3,(ab)5");;|
  !gapprompt@gap>| !gapinput@epi:=EpimorphismSchurCover(g);|
  [ a, b ] -> [ a, b ]
  !gapprompt@gap>| !gapinput@Size(Kernel(epi));|
  2
\end{Verbatim}
 

 If the group becomes bigger, Schur Cover calculations might become unfeasible. 

 There is another operation, \texttt{AbelianInvariantsMultiplier} (\ref{AbelianInvariantsMultiplier}), which only returns the structure of the Schur Multiplier, and which should
work for larger groups as well. }

 

\subsection{\textcolor{Chapter }{AbelianInvariantsMultiplier}}
\logpage{[ 39, 24, 3 ]}\nobreak
\hyperdef{L}{X792BC39D7CEB1D27}{}
{\noindent\textcolor{FuncColor}{$\triangleright$\enspace\texttt{AbelianInvariantsMultiplier({\mdseries\slshape G})\index{AbelianInvariantsMultiplier@\texttt{AbelianInvariantsMultiplier}}
\label{AbelianInvariantsMultiplier}
}\hfill{\scriptsize (attribute)}}\\


 \index{Multiplier} \index{Schur multiplier} returns a list of the abelian invariants of the Schur multiplier of \mbox{\texttt{\mdseries\slshape G}}. 

 At the moment, this operation will not give any information about how to
extend the multiplier to a Schur Cover. 
\begin{Verbatim}[commandchars=!@|,fontsize=\small,frame=single,label=Example]
  !gapprompt@gap>| !gapinput@AbelianInvariantsMultiplier(g);|
  [ 2 ]
  !gapprompt@gap>| !gapinput@AbelianInvariantsMultiplier(AlternatingGroup(6));|
  [ 2, 3 ]
  !gapprompt@gap>| !gapinput@AbelianInvariantsMultiplier(SL(2,3));|
  [  ]
  !gapprompt@gap>| !gapinput@AbelianInvariantsMultiplier(SL(3,2));|
  [ 2 ]
  !gapprompt@gap>| !gapinput@AbelianInvariantsMultiplier(PSU(4,2));|
  [ 2 ]
\end{Verbatim}
 (Note that the last command from the example will take some time.) 

 The \textsf{GAP}{\nobreakspace}4.4.12 manual contained examples for larger groups e.g. $M_{22}$. However, some issues that may very rarely (and not easily reproducibly) lead
to wrong results were discovered in the code capable of handling larger
groups, and in \textsf{GAP}{\nobreakspace}4.5 it was replaced by a more reliable basic method. To deal
with larger groups, one can use the function \texttt{SchurMultiplier} (\textbf{cohomolo: SchurMultiplier}) from the \textsf{cohomolo} package. Also, additional methods for \texttt{AbelianInvariantsMultiplier} are installed in the \textsf{Polycyclic} package for pcp\texttt{\symbol{45}}groups. }

 

\subsection{\textcolor{Chapter }{Epicentre}}
\logpage{[ 39, 24, 4 ]}\nobreak
\hyperdef{L}{X819E8AEC835F8CD1}{}
{\noindent\textcolor{FuncColor}{$\triangleright$\enspace\texttt{Epicentre({\mdseries\slshape G})\index{Epicentre@\texttt{Epicentre}}
\label{Epicentre}
}\hfill{\scriptsize (attribute)}}\\
\noindent\textcolor{FuncColor}{$\triangleright$\enspace\texttt{ExteriorCentre({\mdseries\slshape G})\index{ExteriorCentre@\texttt{ExteriorCentre}}
\label{ExteriorCentre}
}\hfill{\scriptsize (attribute)}}\\


 There are various ways of describing the epicentre of a group \mbox{\texttt{\mdseries\slshape G}}. It is the smallest normal subgroup $N$ of \mbox{\texttt{\mdseries\slshape G}} such that $\mbox{\texttt{\mdseries\slshape G}}/N$ is a central quotient of a group. It is also equal to the Exterior Center of \mbox{\texttt{\mdseries\slshape G}}, see \cite{Ellis98}. }

 

\subsection{\textcolor{Chapter }{NonabelianExteriorSquare}}
\logpage{[ 39, 24, 5 ]}\nobreak
\hyperdef{L}{X8739CD4686301A0E}{}
{\noindent\textcolor{FuncColor}{$\triangleright$\enspace\texttt{NonabelianExteriorSquare({\mdseries\slshape G})\index{NonabelianExteriorSquare@\texttt{NonabelianExteriorSquare}}
\label{NonabelianExteriorSquare}
}\hfill{\scriptsize (operation)}}\\


 Computes the nonabelian exterior square $\mbox{\texttt{\mdseries\slshape G}} \wedge \mbox{\texttt{\mdseries\slshape G}}$ of the group \mbox{\texttt{\mdseries\slshape G}}, which for a finitely presented group is the derived subgroup of any Schur
cover of \mbox{\texttt{\mdseries\slshape G}} (see \cite{BJR87}). }

 

\subsection{\textcolor{Chapter }{EpimorphismNonabelianExteriorSquare}}
\logpage{[ 39, 24, 6 ]}\nobreak
\hyperdef{L}{X7E1C8CD77CDB9F71}{}
{\noindent\textcolor{FuncColor}{$\triangleright$\enspace\texttt{EpimorphismNonabelianExteriorSquare({\mdseries\slshape G})\index{EpimorphismNonabelianExteriorSquare@\texttt{EpimorphismNonabelianExteriorSquare}}
\label{EpimorphismNonabelianExteriorSquare}
}\hfill{\scriptsize (operation)}}\\


 Computes the mapping $\mbox{\texttt{\mdseries\slshape G}} \wedge \mbox{\texttt{\mdseries\slshape G}} \rightarrow \mbox{\texttt{\mdseries\slshape G}}$. The kernel of this mapping is equal to the Schur multiplier of \mbox{\texttt{\mdseries\slshape G}}. }

 

\subsection{\textcolor{Chapter }{IsCentralFactor}}
\logpage{[ 39, 24, 7 ]}\nobreak
\hyperdef{L}{X7BF8DB3D8300BB3F}{}
{\noindent\textcolor{FuncColor}{$\triangleright$\enspace\texttt{IsCentralFactor({\mdseries\slshape G})\index{IsCentralFactor@\texttt{IsCentralFactor}}
\label{IsCentralFactor}
}\hfill{\scriptsize (property)}}\\


 This function determines if there exists a group $H$ such that \mbox{\texttt{\mdseries\slshape G}} is isomorphic to the quotient $H/Z(H)$. A group with this property is called in literature \emph{capable}. A group being capable is equivalent to the epicentre of \mbox{\texttt{\mdseries\slshape G}} being trivial, see \cite{BFS79}. }

 
\subsection{\textcolor{Chapter }{Covering groups of symmetric groups}}\logpage{[ 39, 24, 8 ]}
\hyperdef{L}{X7F4240CD782B6032}{}
{
 The covering groups of symmetric groups were classified in \cite{Schur1911}; an inductive procedure to construct faithful, irreducible representations of
minimal degree over all fields was presented in \cite{Maas2010}. Methods for \texttt{EpimorphismSchurCover} (\ref{EpimorphismSchurCover}) are provided for natural symmetric groups which use these representations. For
alternating groups, the restriction of these representations are provided, but
they may not be irreducible. In the case of degree $6$ and $7$, they are not the full covering groups and so matrix representations are just
stored explicitly for the six\texttt{\symbol{45}}fold covers. 
\begin{Verbatim}[commandchars=!@|,fontsize=\small,frame=single,label=Example]
  !gapprompt@gap>| !gapinput@EpimorphismSchurCover(SymmetricGroup(15));|
  [ < immutable compressed matrix 64x64 over GF(9) >,
    < immutable compressed matrix 64x64 over GF(9) > ] ->
  [ (1,2,3,4,5,6,7,8,9,10,11,12,13,14,15), (1,2) ]
  !gapprompt@gap>| !gapinput@EpimorphismSchurCover(AlternatingGroup(15));|
  [ < immutable compressed matrix 64x64 over GF(9) >,
    < immutable compressed matrix 64x64 over GF(9) > ] ->
  [ (1,2,3,4,5,6,7,8,9,10,11,12,13,14,15), (13,14,15) ]
  !gapprompt@gap>| !gapinput@SchurCoverOfSymmetricGroup(12);|
  <matrix group of size 958003200 with 2 generators>
  !gapprompt@gap>| !gapinput@DoubleCoverOfAlternatingGroup(12);|
  <matrix group of size 479001600 with 2 generators>
  !gapprompt@gap>| !gapinput@BasicSpinRepresentationOfSymmetricGroup( 10, 3, -1 );|
  [ < immutable compressed matrix 16x16 over GF(9) >,
    < immutable compressed matrix 16x16 over GF(9) >,
    < immutable compressed matrix 16x16 over GF(9) >,
    < immutable compressed matrix 16x16 over GF(9) >,
    < immutable compressed matrix 16x16 over GF(9) >,
    < immutable compressed matrix 16x16 over GF(9) >,
    < immutable compressed matrix 16x16 over GF(9) >,
    < immutable compressed matrix 16x16 over GF(9) >,
    < immutable compressed matrix 16x16 over GF(9) > ]
\end{Verbatim}
 }

 

\subsection{\textcolor{Chapter }{BasicSpinRepresentationOfSymmetricGroup}}
\logpage{[ 39, 24, 9 ]}\nobreak
\hyperdef{L}{X7DDA6BC1824F78FD}{}
{\noindent\textcolor{FuncColor}{$\triangleright$\enspace\texttt{BasicSpinRepresentationOfSymmetricGroup({\mdseries\slshape n, p, sign})\index{BasicSpinRepresentationOfSymmetricGroup@\texttt{Basic}\-\texttt{Spin}\-\texttt{Representation}\-\texttt{Of}\-\texttt{Symmetric}\-\texttt{Group}}
\label{BasicSpinRepresentationOfSymmetricGroup}
}\hfill{\scriptsize (function)}}\\


 Constructs the image of the Coxeter generators in the basic spin (projective)
representation of the symmetric group of degree \mbox{\texttt{\mdseries\slshape n}} over a field of characteristic $\mbox{\texttt{\mdseries\slshape p}} \geq 0$. There are two such representations and \mbox{\texttt{\mdseries\slshape sign}} controls which is returned: +1 gives a group where the preimage of an adjacent
transposition $(i,i+1)$ has order 4, \texttt{\symbol{45}}1 gives a group where the preimage of an
adjacent transposition $(i,i+1)$ has order 2. If no \mbox{\texttt{\mdseries\slshape sign}} is specified, +1 is used by default. If no \mbox{\texttt{\mdseries\slshape p}} is specified, 3 is used by default. (Note that the convention of which cover
is labelled as +1 is inconsistent in the literature.)}

 

\subsection{\textcolor{Chapter }{SchurCoverOfSymmetricGroup}}
\logpage{[ 39, 24, 10 ]}\nobreak
\hyperdef{L}{X844CFFDE80F6AD15}{}
{\noindent\textcolor{FuncColor}{$\triangleright$\enspace\texttt{SchurCoverOfSymmetricGroup({\mdseries\slshape n, p, sign})\index{SchurCoverOfSymmetricGroup@\texttt{SchurCoverOfSymmetricGroup}}
\label{SchurCoverOfSymmetricGroup}
}\hfill{\scriptsize (operation)}}\\


 Constructs a Schur cover of \texttt{SymmetricGroup(\mbox{\texttt{\mdseries\slshape n}})} as a faithful, irreducible matrix group in characteristic \mbox{\texttt{\mdseries\slshape p}} ($\mbox{\texttt{\mdseries\slshape p}} \neq 2$). For $\mbox{\texttt{\mdseries\slshape n}} \geq 4$, there are two such covers, and \mbox{\texttt{\mdseries\slshape sign}} determines which is returned: +1 gives a group where the preimage of an
adjacent transposition $(i,i+1)$ has order 4, \texttt{\symbol{45}}1 gives a group where the preimage of an
adjacent transposition $(i,i+1)$ has order 2. If no \mbox{\texttt{\mdseries\slshape sign}} is specified, +1 is used by default. If no \mbox{\texttt{\mdseries\slshape p}} is specified, 3 is used by default. (Note that the convention of which cover
is labelled as +1 is inconsistent in the literature.) For $\mbox{\texttt{\mdseries\slshape n}} \leq 3$, the symmetric group is its own Schur cover and \mbox{\texttt{\mdseries\slshape sign}} is ignored. For $\mbox{\texttt{\mdseries\slshape p}} = 2$, there is no faithful, irreducible representation of the Schur cover unless $\mbox{\texttt{\mdseries\slshape n}} = 1$ or $\mbox{\texttt{\mdseries\slshape n}} = 3$, so \texttt{fail} is returned if $\mbox{\texttt{\mdseries\slshape p}} = 2$. For $\mbox{\texttt{\mdseries\slshape p}} = 3$, $\mbox{\texttt{\mdseries\slshape n}} = 3$, the representation is indecomposable, but reducible. The field of the matrix
group is generally \texttt{GF(\mbox{\texttt{\mdseries\slshape p}}\texttt{\symbol{94}}2)} if $\mbox{\texttt{\mdseries\slshape p}} > 0$, and an abelian number field if $\mbox{\texttt{\mdseries\slshape p}} = 0$. }

 

\subsection{\textcolor{Chapter }{DoubleCoverOfAlternatingGroup}}
\logpage{[ 39, 24, 11 ]}\nobreak
\hyperdef{L}{X7E0F4896795E34FC}{}
{\noindent\textcolor{FuncColor}{$\triangleright$\enspace\texttt{DoubleCoverOfAlternatingGroup({\mdseries\slshape n, p})\index{DoubleCoverOfAlternatingGroup@\texttt{DoubleCoverOfAlternatingGroup}}
\label{DoubleCoverOfAlternatingGroup}
}\hfill{\scriptsize (operation)}}\\


 Constructs a double cover of \texttt{AlternatingGroup(\mbox{\texttt{\mdseries\slshape n}})} as a faithful, completely reducible matrix group in characteristic \mbox{\texttt{\mdseries\slshape p}} ($p \neq 2$) for $n \geq 4$. For $n \leq 3$, the alternating group is its own Schur cover, and \texttt{fail} is returned. For $p = 2$, there is no faithful, completely reducible representation of the double
cover, so \texttt{fail} is returned. The field of the matrix group is generally \texttt{GF(p\texttt{\symbol{94}}2)} if $p>0$, and an abelian number field if $p=0$. If \mbox{\texttt{\mdseries\slshape p}} is omitted, the default is 3. }

 }

  
\section{\textcolor{Chapter }{2\texttt{\symbol{45}}Cohomology}}\label{2-Cohomology}
\logpage{[ 39, 25, 0 ]}
\hyperdef{L}{X7BD95B8D879B73A3}{}
{
  

\subsection{\textcolor{Chapter }{TwoCohomologyGeneric}}
\logpage{[ 39, 25, 1 ]}\nobreak
\hyperdef{L}{X7A1EBC3A7AB0D614}{}
{\noindent\textcolor{FuncColor}{$\triangleright$\enspace\texttt{TwoCohomologyGeneric({\mdseries\slshape G, M})\index{TwoCohomologyGeneric@\texttt{TwoCohomologyGeneric}}
\label{TwoCohomologyGeneric}
}\hfill{\scriptsize (operation)}}\\


 This function computes the second cohomology group for an arbitrary finite
group \mbox{\texttt{\mdseries\slshape G}}. The generators of the module \mbox{\texttt{\mdseries\slshape M}} must correspond to the generators of \mbox{\texttt{\mdseries\slshape G}}. It returns a record with components \texttt{coboundaries}, \texttt{cocycles} and \texttt{cohomology} which are lists of vectors that form a basis of the respective group. \texttt{cohomology} is chosen as a vector space complement to \texttt{coboundaries} in the \texttt{cocycles}. These vectors are representing tails in \mbox{\texttt{\mdseries\slshape M}} with respect to the \texttt{relators} of the presentation \texttt{presentation} of \mbox{\texttt{\mdseries\slshape G}}. (Note that this presentation is on a generating set chosen by the routine,
this generating system corresponds to the components \texttt{group} and \texttt{module} of the record returned. The extension corresponding to a cocyle \texttt{c} can be constructed as \texttt{FpGroupCocycle(r,c)} where \texttt{r} is the cohomology record. This is currently done as a finitely presented
group. 
\begin{Verbatim}[commandchars=!@|,fontsize=\small,frame=single,label=Example]
  !gapprompt@gap>| !gapinput@g:=Group((1,2,3,4,5),(1,2,3));;|
  !gapprompt@gap>| !gapinput@mats:=[[[2,0,0,1],[1,2,1,0],[2,1,1,1],[2,1,1,0]],|
  !gapprompt@>| !gapinput@[[0,2,0,0],[1,2,1,0],[0,0,1,0],[0,0,0,1]]]*Z(3)^0;;|
  !gapprompt@gap>| !gapinput@mo:=GModuleByMats(mats,GF(3));;|
  !gapprompt@gap>| !gapinput@coh:=TwoCohomologyGeneric(g,mo);;|
  !gapprompt@gap>| !gapinput@coh.cocycles;|
  < immutable compressed matrix 8x44 over GF(3) >
  !gapprompt@gap>| !gapinput@coh.coboundaries;|
  [ < immutable compressed vector length 44 over GF(3) >,
    < immutable compressed vector length 44 over GF(3) >,
    < immutable compressed vector length 44 over GF(3) >,
    < immutable compressed vector length 44 over GF(3) >,
    < immutable compressed vector length 44 over GF(3) >,
    < immutable compressed vector length 44 over GF(3) >,
    < immutable compressed vector length 44 over GF(3) > ]
  !gapprompt@gap>| !gapinput@coh.cohomology;|
  [ < immutable compressed vector length 44 over GF(3) > ]
  !gapprompt@gap>| !gapinput@g1:=FpGroupCocycle(coh,coh.zero,true);|
  <fp group of size 4860 on the generators [ F1, F2, F3, m1, m2, m3, m4 ]>
  !gapprompt@gap>| !gapinput@g2:=FpGroupCocycle(coh,coh.cohomology[1],true);|
  <fp group of size 4860 on the generators [ F1, F2, F3, m1, m2, m3, m4 ]>
  !gapprompt@gap>| !gapinput@g1:=Image(IsomorphismPermGroup(g1));|
  <permutation group with 7 generators>
  !gapprompt@gap>| !gapinput@Length(ComplementClassesRepresentatives(g1,SolvableRadical(g1)));|
  3
  !gapprompt@gap>| !gapinput@g2:=Image(IsomorphismPermGroup(g2));|
  <permutation group with 7 generators>
  !gapprompt@gap>| !gapinput@Length(ComplementClassesRepresentatives(g2,SolvableRadical(g2)));|
  0
\end{Verbatim}
 }

 

\subsection{\textcolor{Chapter }{FpGroupCocycle}}
\logpage{[ 39, 25, 2 ]}\nobreak
\hyperdef{L}{X7A65366879BB3977}{}
{\noindent\textcolor{FuncColor}{$\triangleright$\enspace\texttt{FpGroupCocycle({\mdseries\slshape r, c[, doperm]})\index{FpGroupCocycle@\texttt{FpGroupCocycle}}
\label{FpGroupCocycle}
}\hfill{\scriptsize (function)}}\\


 For a record \mbox{\texttt{\mdseries\slshape r}} as returned by \texttt{TwoCohomologyGeneric} (\ref{TwoCohomologyGeneric}) and a vector \mbox{\texttt{\mdseries\slshape c}} in the space of two\texttt{\symbol{45}}cocycles, this operation returns a
finitely presented group that is an extension corresponding to the cocycle \mbox{\texttt{\mdseries\slshape c}}. If the underlying module has dimension $d$, the last $d$ generators generate the normal subgroup. If the optional parameter \mbox{\texttt{\mdseries\slshape doperm}} is given as \mbox{\texttt{\mdseries\slshape true}}, a faithful permutation representation is computed and stored in the
attribute \texttt{IsomorphismPermGroup} (\ref{IsomorphismPermGroup}) of the computed group. If the option \texttt{normalform} is given as \texttt{true}, arithmetic in the resulting finitely presented group will bring words into
normal form. 
\begin{Verbatim}[commandchars=!@|,fontsize=\small,frame=single,label=Example]
  !gapprompt@gap>| !gapinput@g:=Group((2,15,8,16)(3,17,14,21)(4,23,20,6)(5,9,22,11)(7,13,19,25),|
  !gapprompt@>| !gapinput@(2,12,7,17)(3,18,13,23)(4,24,19,9)(5,10,25,15)(6,11,16,21));;|
  !gapprompt@gap>| !gapinput@StructureDescription(g);|
  "GL(2,5)"
  !gapprompt@gap>| !gapinput@mats:=[[[1,1,0,2],[2,0,0,0],[0,2,2,0],[0,1,0,0]],|
  !gapprompt@>| !gapinput@[[0,0,0,1],[1,1,2,0],[1,0,2,1],[1,0,1,0]]]*Z(3)^0;;|
  !gapprompt@gap>| !gapinput@mo:=GModuleByMats(mats,GF(3));;|
  !gapprompt@gap>| !gapinput@coh:=TwoCohomologyGeneric(g,mo);;|
  !gapprompt@gap>| !gapinput@coh.cohomology;|
  [ < immutable compressed vector length 116 over GF(3) > ]
  !gapprompt@gap>| !gapinput@p:=FpGroupCocycle(coh,coh.zero,true);|
  <fp group of size 38880 on the generators
  [ F1, F2, F3, F4, F5, F6, m1, m2, m3, m4 ]>
  !gapprompt@gap>| !gapinput@g1:=Image(IsomorphismPermGroup(p));|
  <permutation group with 10 generators>
  !gapprompt@gap>| !gapinput@p:=FpGroupCocycle(coh,coh.cohomology[1],true);|
  <fp group of size 38880 on the generators
  [ F1, F2, F3, F4, F5, F6, m1, m2, m3, m4 ]>
  !gapprompt@gap>| !gapinput@g2:=Image(IsomorphismPermGroup(p));|
  <permutation group with 10 generators>
  !gapprompt@gap>| !gapinput@Collected(List(MaximalSubgroupClassReps(g1),Size));|
  [ [ 480, 3 ], [ 3888, 1 ], [ 6480, 1 ], [ 7776, 1 ], [ 19440, 1 ] ]
  !gapprompt@gap>| !gapinput@Collected(List(MaximalSubgroupClassReps(g2),Size));|
  [ [ 3888, 1 ], [ 6480, 1 ], [ 7776, 1 ], [ 19440, 1 ] ]
  !gapprompt@gap>| !gapinput@p:=FpGroupCocycle(coh,coh.cohomology[1],true:normalform);;|
\end{Verbatim}
 
\begin{Verbatim}[commandchars=!@|,fontsize=\small,frame=single,label=Example]
  !gapprompt@gap>| !gapinput@p.7*p.1; # i.e. m1*F1, but in normal form|
  F1*m4
\end{Verbatim}
 }

 Also see Section \ref{2-Cohomology and Extensions} for operations and methods specific for Pc groups. }

  
\section{\textcolor{Chapter }{Tests for the Availability of Methods}}\label{Tests for the Availability of Methods}
\logpage{[ 39, 26, 0 ]}
\hyperdef{L}{X865722987E0E19B6}{}
{
  The following filters and operations indicate capabilities of \textsf{GAP}. They can be used in the method selection or algorithms to check whether it
is feasible to compute certain operations for a given group. In general, they
return \texttt{true} if good algorithms for the given arguments are available in \textsf{GAP}. An answer \texttt{false} indicates that no method for this group may exist, or that the existing
methods might run into problems. 

 Typical examples when this might happen is with finitely presented groups, for
which many of the methods cannot be guaranteed to succeed in all situations. 

 The willingness of \textsf{GAP} to perform certain operations may change, depending on which further
information is known about the arguments. Therefore the filters used are not
implemented as properties but as ``other filters'' (see{\nobreakspace}\ref{Properties} and{\nobreakspace}\ref{Other Filters}). 

\subsection{\textcolor{Chapter }{CanEasilyTestMembership}}
\logpage{[ 39, 26, 1 ]}\nobreak
\hyperdef{L}{X798F13EA810FB215}{}
{\noindent\textcolor{FuncColor}{$\triangleright$\enspace\texttt{CanEasilyTestMembership({\mdseries\slshape G})\index{CanEasilyTestMembership@\texttt{CanEasilyTestMembership}}
\label{CanEasilyTestMembership}
}\hfill{\scriptsize (filter)}}\\


 This filter indicates whether \textsf{GAP} can test membership of elements in the group \mbox{\texttt{\mdseries\slshape G}} (via the operation \texttt{\texttt{\symbol{92}}in} (\ref{bSlashin:for a collection})) in reasonable time. It is used by the method selection to decide whether an
algorithm that relies on membership tests may be used. }

 

\subsection{\textcolor{Chapter }{CanEasilyComputeWithIndependentGensAbelianGroup}}
\logpage{[ 39, 26, 2 ]}\nobreak
\hyperdef{L}{X7C2A89607BDFD920}{}
{\noindent\textcolor{FuncColor}{$\triangleright$\enspace\texttt{CanEasilyComputeWithIndependentGensAbelianGroup({\mdseries\slshape G})\index{CanEasilyComputeWithIndependentGensAbelianGroup@\texttt{Can}\-\texttt{Easily}\-\texttt{Compute}\-\texttt{With}\-\texttt{Independent}\-\texttt{Gens}\-\texttt{Abelian}\-\texttt{Group}}
\label{CanEasilyComputeWithIndependentGensAbelianGroup}
}\hfill{\scriptsize (filter)}}\\


 This filter indicates whether \textsf{GAP} can in reasonable time compute independent abelian generators of the group \mbox{\texttt{\mdseries\slshape G}} (via \texttt{IndependentGeneratorsOfAbelianGroup} (\ref{IndependentGeneratorsOfAbelianGroup})) and then can decompose arbitrary group elements with respect to these
generators using \texttt{IndependentGeneratorExponents} (\ref{IndependentGeneratorExponents}). It is used by the method selection to decide whether an algorithm that
relies on these two operations may be used. }

 

\subsection{\textcolor{Chapter }{CanComputeSize}}
\logpage{[ 39, 26, 3 ]}\nobreak
\hyperdef{L}{X83245C82835D496C}{}
{\noindent\textcolor{FuncColor}{$\triangleright$\enspace\texttt{CanComputeSize({\mdseries\slshape dom})\index{CanComputeSize@\texttt{CanComputeSize}}
\label{CanComputeSize}
}\hfill{\scriptsize (filter)}}\\


 This filter indicates that we know that the size of the domain \mbox{\texttt{\mdseries\slshape dom}} (which might be \texttt{infinity} (\ref{infinity})) can be computed reasonably easily. It doesn't imply as quick a computation
as \texttt{HasSize} would but its absence does not imply that the size cannot be computed. }

 

\subsection{\textcolor{Chapter }{CanComputeSizeAnySubgroup}}
\logpage{[ 39, 26, 4 ]}\nobreak
\hyperdef{L}{X8268965487364912}{}
{\noindent\textcolor{FuncColor}{$\triangleright$\enspace\texttt{CanComputeSizeAnySubgroup({\mdseries\slshape G})\index{CanComputeSizeAnySubgroup@\texttt{CanComputeSizeAnySubgroup}}
\label{CanComputeSizeAnySubgroup}
}\hfill{\scriptsize (filter)}}\\


 This filter indicates whether \textsf{GAP} can easily compute the size of any subgroup of the group \mbox{\texttt{\mdseries\slshape G}}. (This is for example advantageous if one can test that a stabilizer index
equals the length of the orbit computed so far to stop early.) }

 

\subsection{\textcolor{Chapter }{CanComputeIndex}}
\logpage{[ 39, 26, 5 ]}\nobreak
\hyperdef{L}{X82DDE00D82A32083}{}
{\noindent\textcolor{FuncColor}{$\triangleright$\enspace\texttt{CanComputeIndex({\mdseries\slshape G, H})\index{CanComputeIndex@\texttt{CanComputeIndex}}
\label{CanComputeIndex}
}\hfill{\scriptsize (operation)}}\\


 This function indicates whether the index $[\mbox{\texttt{\mdseries\slshape G}}:\mbox{\texttt{\mdseries\slshape H}}]$ (which might be \texttt{infinity} (\ref{infinity})) can be computed. It assumes that $\mbox{\texttt{\mdseries\slshape H}} \leq \mbox{\texttt{\mdseries\slshape G}}$ (see \texttt{CanComputeIsSubset} (\ref{CanComputeIsSubset})). }

 

\subsection{\textcolor{Chapter }{CanComputeIsSubset}}
\logpage{[ 39, 26, 6 ]}\nobreak
\hyperdef{L}{X7BE7C36B84C23511}{}
{\noindent\textcolor{FuncColor}{$\triangleright$\enspace\texttt{CanComputeIsSubset({\mdseries\slshape A, B})\index{CanComputeIsSubset@\texttt{CanComputeIsSubset}}
\label{CanComputeIsSubset}
}\hfill{\scriptsize (operation)}}\\


 This filter indicates that \textsf{GAP} can test (via \texttt{IsSubset} (\ref{IsSubset})) whether \mbox{\texttt{\mdseries\slshape B}} is a subset of \mbox{\texttt{\mdseries\slshape A}}. }

 

\subsection{\textcolor{Chapter }{KnowsHowToDecompose}}
\logpage{[ 39, 26, 7 ]}\nobreak
\hyperdef{L}{X87D62C2C7C375E2D}{}
{\noindent\textcolor{FuncColor}{$\triangleright$\enspace\texttt{KnowsHowToDecompose({\mdseries\slshape G[, gens]})\index{KnowsHowToDecompose@\texttt{KnowsHowToDecompose}}
\label{KnowsHowToDecompose}
}\hfill{\scriptsize (property)}}\\


 Tests whether the group \mbox{\texttt{\mdseries\slshape G}} can decompose elements in the generators \mbox{\texttt{\mdseries\slshape gens}}. If \mbox{\texttt{\mdseries\slshape gens}} is not given it tests, whether it can decompose in the generators given in the \texttt{GeneratorsOfGroup} (\ref{GeneratorsOfGroup}) value of \mbox{\texttt{\mdseries\slshape G}}. 

 This property can be used for example to check whether a group homomorphism by
images (see \texttt{GroupHomomorphismByImages} (\ref{GroupHomomorphismByImages})) can be reasonably defined from this group. }

 }

  
\section{\textcolor{Chapter }{Specific functions for Normalizer calculation}}\label{Specific functions for Normalizer calculation}
\logpage{[ 39, 27, 0 ]}
\hyperdef{L}{X83A9997586694DC0}{}
{
  

\subsection{\textcolor{Chapter }{NormalizerViaRadical}}
\logpage{[ 39, 27, 1 ]}\nobreak
\hyperdef{L}{X84ABCA997D294B36}{}
{\noindent\textcolor{FuncColor}{$\triangleright$\enspace\texttt{NormalizerViaRadical({\mdseries\slshape G, S})\index{NormalizerViaRadical@\texttt{NormalizerViaRadical}}
\label{NormalizerViaRadical}
}\hfill{\scriptsize (function)}}\\


 This function implements a particular approach, following the SolvableRadical
paradigm, for calculating the normalizer of a subgroup \mbox{\texttt{\mdseries\slshape S}} in \mbox{\texttt{\mdseries\slshape G}}. It is at the moment provided only as a separate function, and not as method
for the operation \texttt{Normalizer}, as it can often be slower than other built\texttt{\symbol{45}}in routines.
In certain hard cases (non\texttt{\symbol{45}}solvable groups with nontrivial
radical), however its performance is substantially superior. The function thus
is provided as a non\texttt{\symbol{45}}automated tool for advanced users. 
\begin{Verbatim}[commandchars=!@|,fontsize=\small,frame=single,label=Example]
  !gapprompt@gap>| !gapinput@g:=TransitiveGroup(30,2030);;|
  !gapprompt@gap>| !gapinput@s:=SylowSubgroup(g,5);;|
  !gapprompt@gap>| !gapinput@Size(NormalizerViaRadical(g,s));|
  28800
\end{Verbatim}
 Note that this example only demonstrates usage, but that in this case in fact
the ordinary \texttt{Normalizer} routine performs faster. }

 }

 }

            
\chapter{\textcolor{Chapter }{Group Homomorphisms}}\label{Group Homomorphisms}
\logpage{[ 40, 0, 0 ]}
\hyperdef{L}{X83702FC27B3C3098}{}
{
  A group homomorphism is a mapping from one group to another that respects
multiplication and inverses. They are implemented as a special class of
mappings, so in particular all operations for mappings, such as \texttt{Image} (\ref{Image:set of images of the source of a general mapping}), \texttt{PreImage} (\ref{PreImage:set of preimages of the range of a general mapping}), \texttt{PreImagesRepresentative} (\ref{PreImagesRepresentative}), \texttt{KernelOfMultiplicativeGeneralMapping} (\ref{KernelOfMultiplicativeGeneralMapping}), \texttt{Source} (\ref{Source}), \texttt{Range} (\ref{Range:of a general mapping}), \texttt{IsInjective} (\ref{IsInjective}) and \texttt{IsSurjective} (\ref{IsSurjective}) (see chapter{\nobreakspace}\ref{Mappings}, in particular section{\nobreakspace}\ref{Mappings that Respect Multiplication}) are applicable to them. 

 Homomorphisms can be used to transfer calculations into isomorphic groups in
another representation, for which better algorithms are available.
Section{\nobreakspace}\ref{Nice Monomorphisms} explains a technique how to enforce this automatically. 

 Homomorphisms are also used to represent group automorphisms, and section \ref{Group Automorphisms} explains \textsf{GAP}'s facilities to work with automorphism groups. 

 Section \ref{Searching for Homomorphisms} explains how to make \textsf{GAP} to search for all homomorphisms between two groups which fulfill certain
specifications.  
\section{\textcolor{Chapter }{Creating Group Homomorphisms}}\label{Creating Group Homomorphisms}
\logpage{[ 40, 1, 0 ]}
\hyperdef{L}{X81A7BB0F7D81B247}{}
{
  The most important way of creating group homomorphisms is to give images for a
set of group generators and to extend it to the group generated by them by the
homomorphism property. 

 \emph{A second} way to create homomorphisms is to give functions that compute image and
preimage. (A similar case are homomorphisms that are induced by conjugation.
Special constructors for such mappings are described in section{\nobreakspace}\ref{Group Automorphisms}). 

 \emph{The third} class are epimorphisms from a group onto its factor group. Such homomorphisms
can be constructed by \texttt{NaturalHomomorphismByNormalSubgroup} (\ref{NaturalHomomorphismByNormalSubgroup}). 

 \emph{The fourth} class is homomorphisms in a permutation group that are induced by an action on
a set. Such homomorphisms are described in the context of group actions, see
chapter{\nobreakspace}\ref{Group Actions} and in particular \texttt{ActionHomomorphism} (\ref{ActionHomomorphism:for a group, an action domain, etc.}). 

\subsection{\textcolor{Chapter }{GroupHomomorphismByImages}}
\logpage{[ 40, 1, 1 ]}\nobreak
\hyperdef{L}{X7F348F497C813BE0}{}
{\noindent\textcolor{FuncColor}{$\triangleright$\enspace\texttt{GroupHomomorphismByImages({\mdseries\slshape G, H[[, gens], imgs]})\index{GroupHomomorphismByImages@\texttt{GroupHomomorphismByImages}}
\label{GroupHomomorphismByImages}
}\hfill{\scriptsize (function)}}\\


 \texttt{GroupHomomorphismByImages} returns the group homomorphism with source \mbox{\texttt{\mdseries\slshape G}} and range \mbox{\texttt{\mdseries\slshape H}} that is defined by mapping the list \mbox{\texttt{\mdseries\slshape gens}} of generators of \mbox{\texttt{\mdseries\slshape G}} to the list \mbox{\texttt{\mdseries\slshape imgs}} of images in \mbox{\texttt{\mdseries\slshape H}}. 

 If omitted, the arguments \mbox{\texttt{\mdseries\slshape gens}} and \mbox{\texttt{\mdseries\slshape imgs}} default to the \texttt{GeneratorsOfGroup} (\ref{GeneratorsOfGroup}) value of \mbox{\texttt{\mdseries\slshape G}} and \mbox{\texttt{\mdseries\slshape H}}, respectively. If \mbox{\texttt{\mdseries\slshape H}} is not given the mapping is automatically considered as surjective. 

 If \mbox{\texttt{\mdseries\slshape gens}} does not generate \mbox{\texttt{\mdseries\slshape G}} or if the mapping of the generators does not extend to a homomorphism (i.e.,
if mapping the generators describes only a multi\texttt{\symbol{45}}valued
mapping) then \texttt{fail} is returned. 

 This test can be quite expensive. If one is certain that the mapping of the
generators extends to a homomorphism, one can avoid the checks by calling \texttt{GroupHomomorphismByImagesNC} (\ref{GroupHomomorphismByImagesNC}). (There also is the possibility to construct potentially
multi\texttt{\symbol{45}}valued mappings with \texttt{GroupGeneralMappingByImages} (\ref{GroupGeneralMappingByImages}) and to test with \texttt{IsMapping} (\ref{IsMapping}) whether they are indeed homomorphisms.) }

 

\subsection{\textcolor{Chapter }{GroupHomomorphismByImagesNC}}
\logpage{[ 40, 1, 2 ]}\nobreak
\hyperdef{L}{X7AB15AF5830F2A6B}{}
{\noindent\textcolor{FuncColor}{$\triangleright$\enspace\texttt{GroupHomomorphismByImagesNC({\mdseries\slshape G, H[[, gens], imgs]})\index{GroupHomomorphismByImagesNC@\texttt{GroupHomomorphismByImagesNC}}
\label{GroupHomomorphismByImagesNC}
}\hfill{\scriptsize (operation)}}\\


 \texttt{GroupHomomorphismByImagesNC} creates a homomorphism as \texttt{GroupHomomorphismByImages} (\ref{GroupHomomorphismByImages}) does, however it does not test whether \mbox{\texttt{\mdseries\slshape gens}} generates \mbox{\texttt{\mdseries\slshape G}} and that the mapping of \mbox{\texttt{\mdseries\slshape gens}} to \mbox{\texttt{\mdseries\slshape imgs}} indeed defines a group homomorphism. Because these tests can be expensive it
can be substantially faster than \texttt{GroupHomomorphismByImages} (\ref{GroupHomomorphismByImages}). Results are unpredictable if the conditions do not hold. 

 If omitted, the arguments \mbox{\texttt{\mdseries\slshape gens}} and \mbox{\texttt{\mdseries\slshape imgs}} default to the \texttt{GeneratorsOfGroup} (\ref{GeneratorsOfGroup}) value of \mbox{\texttt{\mdseries\slshape G}} and \mbox{\texttt{\mdseries\slshape H}}, respectively. 

 (For creating a possibly multi\texttt{\symbol{45}}valued mapping from \mbox{\texttt{\mdseries\slshape G}} to \mbox{\texttt{\mdseries\slshape H}} that respects multiplication and inverses, \texttt{GroupGeneralMappingByImages} (\ref{GroupGeneralMappingByImages}) can be used.)          

 
\begin{Verbatim}[commandchars=!@|,fontsize=\small,frame=single,label=Example]
  !gapprompt@gap>| !gapinput@gens:=[(1,2,3,4),(1,2)];|
  [ (1,2,3,4), (1,2) ]
  !gapprompt@gap>| !gapinput@g:=Group(gens);|
  Group([ (1,2,3,4), (1,2) ])
  !gapprompt@gap>| !gapinput@h:=Group((1,2,3),(1,2));|
  Group([ (1,2,3), (1,2) ])
  !gapprompt@gap>| !gapinput@hom:=GroupHomomorphismByImages(g,h,gens,[(1,2),(1,3)]);|
  [ (1,2,3,4), (1,2) ] -> [ (1,2), (1,3) ]
  !gapprompt@gap>| !gapinput@Image(hom,(1,4));|
  (2,3)
  !gapprompt@gap>| !gapinput@map:=GroupHomomorphismByImages(g,h,gens,[(1,2,3),(1,2)]);|
  fail
\end{Verbatim}
 }

 

\subsection{\textcolor{Chapter }{GroupGeneralMappingByImages}}
\logpage{[ 40, 1, 3 ]}\nobreak
\hyperdef{L}{X7A59F2C47BD41DC8}{}
{\noindent\textcolor{FuncColor}{$\triangleright$\enspace\texttt{GroupGeneralMappingByImages({\mdseries\slshape G, H, gens, imgs})\index{GroupGeneralMappingByImages@\texttt{GroupGeneralMappingByImages}}
\label{GroupGeneralMappingByImages}
}\hfill{\scriptsize (operation)}}\\
\noindent\textcolor{FuncColor}{$\triangleright$\enspace\texttt{GroupGeneralMappingByImages({\mdseries\slshape G, gens, imgs})\index{GroupGeneralMappingByImages@\texttt{GroupGeneralMappingByImages}!from group to itself}
\label{GroupGeneralMappingByImages:from group to itself}
}\hfill{\scriptsize (operation)}}\\
\noindent\textcolor{FuncColor}{$\triangleright$\enspace\texttt{GroupGeneralMappingByImagesNC({\mdseries\slshape G, H, gens, imgs})\index{GroupGeneralMappingByImagesNC@\texttt{GroupGeneralMappingByImagesNC}}
\label{GroupGeneralMappingByImagesNC}
}\hfill{\scriptsize (operation)}}\\
\noindent\textcolor{FuncColor}{$\triangleright$\enspace\texttt{GroupGeneralMappingByImagesNC({\mdseries\slshape G, gens, imgs})\index{GroupGeneralMappingByImagesNC@\texttt{GroupGeneralMappingByImagesNC}!from group to itself}
\label{GroupGeneralMappingByImagesNC:from group to itself}
}\hfill{\scriptsize (operation)}}\\


 returns a general mapping defined by extending the mapping from \mbox{\texttt{\mdseries\slshape gens}} to \mbox{\texttt{\mdseries\slshape imgs}} homomorphically. If the range \mbox{\texttt{\mdseries\slshape H}} is not given the mapping will be made automatically surjective. The NC version
does not test whether \mbox{\texttt{\mdseries\slshape gens}} are contained in \mbox{\texttt{\mdseries\slshape G}} or \mbox{\texttt{\mdseries\slshape imgs}} are contained in \mbox{\texttt{\mdseries\slshape H}}. (\texttt{GroupHomomorphismByImages} (\ref{GroupHomomorphismByImages}) creates a group general mapping by images and tests whether it is in \texttt{IsMapping} (\ref{IsMapping}).) 
\begin{Verbatim}[commandchars=!@|,fontsize=\small,frame=single,label=Example]
  !gapprompt@gap>| !gapinput@map:=GroupGeneralMappingByImages(g,h,gens,[(1,2,3),(1,2)]);|
  [ (1,2,3,4), (1,2) ] -> [ (1,2,3), (1,2) ]
  !gapprompt@gap>| !gapinput@IsMapping(map);|
  false
\end{Verbatim}
 }

 
\subsection{\textcolor{Chapter }{GroupHomomorphismByFunction}}\logpage{[ 40, 1, 4 ]}
\hyperdef{L}{X7BC6C20E7CEDBFC5}{}
{
\noindent\textcolor{FuncColor}{$\triangleright$\enspace\texttt{GroupHomomorphismByFunction({\mdseries\slshape S, R, fun[, invfun]})\index{GroupHomomorphismByFunction@\texttt{GroupHomomorphismByFunction}!by function (and inverse function) between two domains}
\label{GroupHomomorphismByFunction:by function (and inverse function) between two domains}
}\hfill{\scriptsize (function)}}\\
\noindent\textcolor{FuncColor}{$\triangleright$\enspace\texttt{GroupHomomorphismByFunction({\mdseries\slshape S, R, fun, false, prefun})\index{GroupHomomorphismByFunction@\texttt{GroupHomomorphismByFunction}!by function and function that computes one preimage}
\label{GroupHomomorphismByFunction:by function and function that computes one preimage}
}\hfill{\scriptsize (function)}}\\


 \texttt{GroupHomomorphismByFunction} returns a group homomorphism \texttt{hom} with source \mbox{\texttt{\mdseries\slshape S}} and range \mbox{\texttt{\mdseries\slshape R}}, such that each element \texttt{s} of \mbox{\texttt{\mdseries\slshape S}} is mapped to the element \mbox{\texttt{\mdseries\slshape fun}}\texttt{( s )}, where \mbox{\texttt{\mdseries\slshape fun}} is a \textsf{GAP} function. 

 If the argument \mbox{\texttt{\mdseries\slshape invfun}} is bound then \mbox{\texttt{\mdseries\slshape hom}} is a bijection between \mbox{\texttt{\mdseries\slshape S}} and \mbox{\texttt{\mdseries\slshape R}}, and the preimage of each element \texttt{r} of \mbox{\texttt{\mdseries\slshape R}} is given by \mbox{\texttt{\mdseries\slshape invfun}}\texttt{( r )}, where \mbox{\texttt{\mdseries\slshape invfun}} is a \textsf{GAP} function. 

 If five arguments are given and the fourth argument is \texttt{false} then the \textsf{GAP} function \mbox{\texttt{\mdseries\slshape prefun}} can be used to compute a single preimage also if \texttt{hom} is not bijective. 

 No test is performed on whether the functions actually give an homomorphism
between both groups because this would require testing the full multiplication
table. 

 \texttt{GroupHomomorphismByFunction} creates a mapping which lies in \texttt{IsSPGeneralMapping} (\ref{IsSPGeneralMapping}). 

 
\begin{Verbatim}[commandchars=!@|,fontsize=\small,frame=single,label=Example]
  !gapprompt@gap>| !gapinput@hom:=GroupHomomorphismByFunction(g,h,|
  !gapprompt@>| !gapinput@function(x) if SignPerm(x)=-1 then return (1,2); else return ();fi;end);|
  MappingByFunction( Group([ (1,2,3,4), (1,2) ]), Group(
  [ (1,2,3), (1,2) ]), function( x ) ... end )
  !gapprompt@gap>| !gapinput@ImagesSource(hom);|
  Group([ (1,2), (1,2) ])
  !gapprompt@gap>| !gapinput@Image(hom,(1,2,3,4));|
  (1,2)
\end{Verbatim}
 }

 

\subsection{\textcolor{Chapter }{AsGroupGeneralMappingByImages}}
\logpage{[ 40, 1, 5 ]}\nobreak
\hyperdef{L}{X785AB6057F736344}{}
{\noindent\textcolor{FuncColor}{$\triangleright$\enspace\texttt{AsGroupGeneralMappingByImages({\mdseries\slshape map})\index{AsGroupGeneralMappingByImages@\texttt{AsGroupGeneralMappingByImages}}
\label{AsGroupGeneralMappingByImages}
}\hfill{\scriptsize (attribute)}}\\


 If \mbox{\texttt{\mdseries\slshape map}} is a mapping from one group to another this attribute returns a group general
mapping that which implements the same abstract mapping. (Some operations can
be performed more effective in this representation, see also{\nobreakspace}\texttt{IsGroupGeneralMappingByAsGroupGeneralMappingByImages} (\ref{IsGroupGeneralMappingByAsGroupGeneralMappingByImages}).) 
\begin{Verbatim}[commandchars=!@|,fontsize=\small,frame=single,label=Example]
  !gapprompt@gap>| !gapinput@AsGroupGeneralMappingByImages(hom);|
  [ (1,2,3,4), (1,2) ] -> [ (1,2), (1,2) ]
\end{Verbatim}
 }

 }

  
\section{\textcolor{Chapter }{Operations for Group Homomorphisms}}\label{Operations for Group Homomorphisms}
\logpage{[ 40, 2, 0 ]}
\hyperdef{L}{X794043AC7E4FAF9E}{}
{
  \index{kernel!group homomorphism} Group homomorphisms are mappings, so all the operations and properties for
mappings described in chapter{\nobreakspace}\ref{Mappings} are applicable to them. (However often much better methods, than for general
mappings are available.) 

 Group homomorphisms will map groups to groups by just mapping the set of
generators. 

 \texttt{KernelOfMultiplicativeGeneralMapping} (\ref{KernelOfMultiplicativeGeneralMapping}) can be used to compute the kernel of a group homomorphism. 

 
\begin{Verbatim}[commandchars=!@|,fontsize=\small,frame=single,label=Example]
  !gapprompt@gap>| !gapinput@hom:=GroupHomomorphismByImages(g,h,gens,[(1,2),(1,3)]);;|
  !gapprompt@gap>| !gapinput@Kernel(hom);|
  Group([ (1,4)(2,3), (1,2)(3,4) ])
\end{Verbatim}
 

 Homomorphisms can map between groups in different representations and are also
used to get isomorphic groups in a different representation. 

 
\begin{Verbatim}[commandchars=!@|,fontsize=\small,frame=single,label=Example]
  !gapprompt@gap>| !gapinput@m1:=[[0,-1],[1,0]];;m2:=[[0,-1],[1,1]];;|
  !gapprompt@gap>| !gapinput@sl2z:=Group(m1,m2);; # SL(2,Integers) as matrix group|
  !gapprompt@gap>| !gapinput@F:=FreeGroup(2);;|
  !gapprompt@gap>| !gapinput@psl2z:=F/[F.1^2,F.2^3]; #PSL(2,Z) as FP group|
  <fp group on the generators [ f1, f2 ]>
  !gapprompt@gap>| !gapinput@phom:=GroupHomomorphismByImagesNC(sl2z,psl2z,[m1,m2],|
  !gapprompt@>| !gapinput@GeneratorsOfGroup(psl2z)); # the non NC-version would be expensive|
  [ [ [ 0, -1 ], [ 1, 0 ] ], [ [ 0, -1 ], [ 1, 1 ] ] ] -> [ f1, f2 ]
  !gapprompt@gap>| !gapinput@Kernel(phom); # the diagonal matrices|
  Group([ [ [ -1, 0 ], [ 0, -1 ] ], [ [ -1, 0 ], [ 0, -1 ] ] ])
  !gapprompt@gap>| !gapinput@p1:=(1,2)(3,4);;p2:=(2,4,5);;a5:=Group(p1,p2);;|
  !gapprompt@gap>| !gapinput@ahom:=GroupHomomorphismByImages(psl2z,a5,|
  !gapprompt@>| !gapinput@GeneratorsOfGroup(psl2z),[p1,p2]); # here homomorphism test is cheap.|
  [ f1, f2 ] -> [ (1,2)(3,4), (2,4,5) ]
  !gapprompt@gap>| !gapinput@u:=PreImage(ahom,Group((1,2,3),(1,2)(4,5)));|
  Group(<fp, no generators known>)
  !gapprompt@gap>| !gapinput@Index(psl2z,u);|
  10
  !gapprompt@gap>| !gapinput@isofp:=IsomorphismFpGroup(u);; Image(isofp);|
  <fp group of size infinity on the generators [ F1, F2, F3, F4 ]>
  !gapprompt@gap>| !gapinput@RelatorsOfFpGroup(Image(isofp));|
  [ F1^2, F4^2, F3^3 ]
  !gapprompt@gap>| !gapinput@up:=PreImage(phom,u);;|
  !gapprompt@gap>| !gapinput@List(GeneratorsOfGroup(up),TraceMat);|
  [ -2, -2, 0, -4, 1, 0 ]
\end{Verbatim}
 

 For an automorphism \mbox{\texttt{\mdseries\slshape aut}}, \texttt{Inverse} (\ref{Inverse}) \index{Inverse!group homomorphism} returns the inverse automorphism $\mbox{\texttt{\mdseries\slshape aut}}^{{-1}}$. However if \mbox{\texttt{\mdseries\slshape hom}} is a bijective homomorphism between different groups, or if \mbox{\texttt{\mdseries\slshape hom}} is injective and considered to be a bijection to its image, the operation \texttt{InverseGeneralMapping} (\ref{InverseGeneralMapping}) should be used instead. (See{\nobreakspace}\texttt{Inverse} (\ref{Inverse}) for a further discussion of this problem.) 

 
\begin{Verbatim}[commandchars=!@|,fontsize=\small,frame=single,label=Example]
  !gapprompt@gap>| !gapinput@iso:=IsomorphismPcGroup(g);|
  Pcgs([ (3,4), (2,4,3), (1,4)(2,3), (1,3)(2,4) ]) -> [ f1, f2, f3, f4 ]
  !gapprompt@gap>| !gapinput@Inverse(iso);|
  #I  The mapping must be bijective and have source=range
  #I  You might want to use `InverseGeneralMapping'
  fail
  !gapprompt@gap>| !gapinput@InverseGeneralMapping(iso);|
  [ f1, f2, f3, f4 ] -> Pcgs([ (3,4), (2,4,3), (1,4)(2,3), (1,3)(2,4) ])
\end{Verbatim}
 }

  
\section{\textcolor{Chapter }{Efficiency of Homomorphisms}}\label{Efficiency of Homomorphisms}
\logpage{[ 40, 3, 0 ]}
\hyperdef{L}{X7A121B9E7F78138A}{}
{
  \textsf{GAP} permits to create homomorphisms between arbitrary groups. This section
considers the efficiency of the implementation and shows ways how to choose
suitable representations. For permutation groups (see{\nobreakspace}\ref{Permutation Groups}) or Pc groups (see{\nobreakspace}\ref{Pc Groups}) this is normally nothing to worry about, unless the groups get extremely
large. For other groups however certain calculations might be expensive and
some precaution might be needed to avoid unnecessarily expensive calculations. 

 In short, it is always worth to tell a mapping that it is a homomorphism (this
can be done by calling \texttt{SetIsMapping}) (or to create it directly with \texttt{GroupHomomorphismByImagesNC} (\ref{GroupHomomorphismByImagesNC})). 

 The basic operations required are to compute image and preimage of elements
and to test whether a mapping is a homomorphism. Their cost will differ
depending on the type of the mapping.  
\subsection{\textcolor{Chapter }{Mappings given on generators}}\label{Mappings given on generators}
\logpage{[ 40, 3, 1 ]}
\hyperdef{L}{X84CFBB577BAFFD4D}{}
{
  See \texttt{GroupHomomorphismByImages} (\ref{GroupHomomorphismByImages}) and \texttt{GroupGeneralMappingByImages} (\ref{GroupGeneralMappingByImages}). 

 Computing images requires to express an element of the source as word in the
generators. If it cannot be done effectively (this is determined by \texttt{KnowsHowToDecompose} (\ref{KnowsHowToDecompose}) which returns \texttt{true} for example for arbitrary permutation groups, for Pc groups or for finitely
presented groups with the images of the free generators) the span of the
generators has to be computed elementwise which can be very expensive and
memory consuming. 

 Computing preimages adheres to the same rules with swapped r{\^o}les of
generators and their images. 

 The test whether a mapping is a homomorphism requires the computation of a
presentation for the source and evaluation of its relators in the images of
its generators. For larger groups this can be expensive and \texttt{GroupHomomorphismByImagesNC} (\ref{GroupHomomorphismByImagesNC}) should be used if the mapping is known to be a homomorphism. }

  
\subsection{\textcolor{Chapter }{Action homomorphisms}}\label{Action homomorphisms}
\logpage{[ 40, 3, 2 ]}
\hyperdef{L}{X86C2BE2481FDC8EE}{}
{
  See \texttt{ActionHomomorphism} (\ref{ActionHomomorphism:for a group, an action domain, etc.}). 

 The calculation of images is determined by the acting function used and
{\textendash}for large domains{\textendash} is often dominated by the search
for the position of an image in a list of the domain elements. This can be
improved by sorting this list if an efficient method for \texttt{\texttt{\symbol{92}}{\textless}} (\ref{bSlash<}) to compare elements of the domain is available. 

 Once the images of a generating set are computed, computing preimages (which
is done via \texttt{AsGroupGeneralMappingByImages} (\ref{AsGroupGeneralMappingByImages})) and computing the kernel behaves the same as for a homomorphism created with \texttt{GroupHomomorphismByImages} (\ref{GroupHomomorphismByImages}) from a permutation group. 

 \textsf{GAP} will always assume that the acting function provided implements a proper group
action and thus that the mapping is indeed a homomorphism. }

  
\subsection{\textcolor{Chapter }{Mappings given by functions}}\label{Mappings given by functions}
\logpage{[ 40, 3, 3 ]}
\hyperdef{L}{X802C5A887D8A7CC4}{}
{
  See \texttt{GroupHomomorphismByFunction} (\ref{GroupHomomorphismByFunction:by function (and inverse function) between two domains}).   

 Computing images is wholly determined by the function that performs the image
calculation. If no function to compute preimages is given, computing preimages
requires mapping every element of the source to find an element that maps to
the requested image. This is time and memory consuming. 

   }

  
\subsection{\textcolor{Chapter }{Other operations}}\label{Other operations}
\logpage{[ 40, 3, 4 ]}
\hyperdef{L}{X87497C207B7D7511}{}
{
  To compute the kernel of a homomorphism (unless the mapping is known to be
injective) requires the capability to compute a presentation of the image and
to evaluate the relators of this presentation in preimages of the
presentations generators. 

 The calculation of the \texttt{Image} (\ref{Image:set of images of the source of a general mapping}) (respectively \texttt{ImagesSource} (\ref{ImagesSource})) value requires to map a generating set of the source, testing surjectivity
is a comparison for equality with the range. 

 Testing injectivity is a test for triviality of the kernel. 

 The comparison of mappings is based on a lexicographic comparison of a sorted
element list of the source. For group homomorphisms, this can be simplified,
using \texttt{ImagesSmallestGenerators} (\ref{ImagesSmallestGenerators}) }

 

\subsection{\textcolor{Chapter }{ImagesSmallestGenerators}}
\logpage{[ 40, 3, 5 ]}\nobreak
\hyperdef{L}{X80B8ABEC7CC20DFB}{}
{\noindent\textcolor{FuncColor}{$\triangleright$\enspace\texttt{ImagesSmallestGenerators({\mdseries\slshape map})\index{ImagesSmallestGenerators@\texttt{ImagesSmallestGenerators}}
\label{ImagesSmallestGenerators}
}\hfill{\scriptsize (attribute)}}\\


 returns the list of images of \texttt{GeneratorsSmallest(Source(\mbox{\texttt{\mdseries\slshape map}}))}. This list can be used to compare group homomorphisms. (The standard
comparison is to compare the image lists on the set of elements of the source.
If however x and y have the same images under a and b, certainly all their
products have. Therefore it is sufficient to test this on the images of the
smallest generators.) }

 }

  
\section{\textcolor{Chapter }{Homomorphism for very large groups}}\label{Homomorphism for very large groups}
\logpage{[ 40, 4, 0 ]}
\hyperdef{L}{X7BA90DA481A1C6D6}{}
{
  Some homomorphisms (notably particular actions) transfer known information
about the source group (such as a stabilizer chain) to the image group if this
is substantially cheaper than to compute the information in the image group
anew. In most cases this is no problem and in fact speeds up further
calculations notably. 

 For a huge source group, however this can be time consuming or take a large
amount of extra memory for storage. In this case it can be helpful to avoid as
much automatism as possible. 

 The following list of tricks might be useful in such a case. (However you will
lose much automatic deduction. So please restrict the use of these to cases
where the standard approach does not work.) 
\begin{itemize}
\item  Compute only images (or the \texttt{PreImagesRepresentative} (\ref{PreImagesRepresentative})) of group elements. Do not compute the images of (sub)groups or the full
preimage of a subgroup. 
\item  Create action homomorphisms as ``surjective'' (see \texttt{ActionHomomorphism} (\ref{ActionHomomorphism:for a group, an action domain, etc.})), otherwise the range is set to be the full symmetric group. However do not
compute \texttt{Range} (\ref{Range:of a general mapping}) or \texttt{Image} (\ref{Image:set of images of the source of a general mapping}) values, but only the images of a generator set. 
\item  If you suspect an action homomorphism to do too much internally, replace the
action function with a function that does the same; i.e. replace \texttt{OnPoints} (\ref{OnPoints}) by \texttt{function( p, g ) return p\texttt{\symbol{94}}g; end;}. The action will be the same, but as the action function is not \texttt{OnPoints} (\ref{OnPoints}), the extra processing for special cases is not triggered. 
\end{itemize}
 }

  
\section{\textcolor{Chapter }{Nice Monomorphisms}}\label{Nice Monomorphisms}
\logpage{[ 40, 5, 0 ]}
\hyperdef{L}{X7FFD731684606BC6}{}
{
  \textsf{GAP} contains very efficient algorithms for some special representations of groups
(for example pc groups or permutation groups) while for other representations
only slow generic methods are available. In this case it can be worthwhile to
do all calculations rather in an isomorphic image of the group, which is in a ``better'' representation. The way to achieve this in \textsf{GAP} is via \emph{nice monomorphisms}. 

 For this mechanism to work, of course there must be effective methods to
evaluate the \texttt{NiceMonomorphism} (\ref{NiceMonomorphism}) value on elements and to take preimages under it. As by definition no good
algorithms exist for the source group, normally this can only be achieved by
using the result of a call to \texttt{ActionHomomorphism} (\ref{ActionHomomorphism:for a group, an action domain, etc.}) or \texttt{GroupHomomorphismByFunction} (\ref{GroupHomomorphismByFunction:by function (and inverse function) between two domains}) (see also section{\nobreakspace}\ref{Efficiency of Homomorphisms}). 

 It may happen that one knows a monomorphism \texttt{map} that is suitable as a nice monomorphism of a given group \texttt{G}, for example if \texttt{G} is a matrix group and \texttt{map} describes the faithful action of \texttt{G} on a small set of vectors. In such a case, one can prescribe this monomorphism
via \texttt{SetNiceMonomorphism( G, map )} if certain technical requirements are satisfied. 

 
\begin{itemize}
\item  \texttt{G} does not store a nice monomorphism already (i.e., \texttt{HasNiceMonomorphism( G )} returns \texttt{false}). 
\item  \texttt{map} stores that it is injective. This can be achieved by calling \texttt{IsInjective( map )} or, if one is sure that \texttt{map} is injective and one wants to avoid the overhead of the test, by calling \texttt{SetIsInjective( map, true )}. Note that the \texttt{IsInjective} call for \texttt{map} may compute and store some different nice monomorphism in the source of \texttt{map}, which is probably \texttt{G}. 
\item  \texttt{map} knows that it is a group homomorphism, calling \texttt{ImagesSet} (\ref{ImagesSet}) with \texttt{map} and a subgroup of \texttt{G} returns a group object, and analogously \texttt{PreImagesSet} (\ref{PreImagesSet}) for an image group returns a group object. For example, mappings constructed
with \texttt{GeneralMappingByElements} (\ref{GeneralMappingByElements}) will \emph{not} be ``nice enough'' to serve as nice monomorphisms. 
\end{itemize}
 

\subsection{\textcolor{Chapter }{IsHandledByNiceMonomorphism}}
\logpage{[ 40, 5, 1 ]}\nobreak
\hyperdef{L}{X78849F81804C44B3}{}
{\noindent\textcolor{FuncColor}{$\triangleright$\enspace\texttt{IsHandledByNiceMonomorphism({\mdseries\slshape obj})\index{IsHandledByNiceMonomorphism@\texttt{IsHandledByNiceMonomorphism}}
\label{IsHandledByNiceMonomorphism}
}\hfill{\scriptsize (property)}}\\


 If this property is \texttt{true}, high\texttt{\symbol{45}}valued methods that translate all calculations in \mbox{\texttt{\mdseries\slshape obj}} in the image under the \texttt{NiceMonomorphism} (\ref{NiceMonomorphism}) value of \mbox{\texttt{\mdseries\slshape obj}} become available for \mbox{\texttt{\mdseries\slshape obj}}. }

 

\subsection{\textcolor{Chapter }{NiceMonomorphism}}
\logpage{[ 40, 5, 2 ]}\nobreak
\hyperdef{L}{X7965086E82ABCF41}{}
{\noindent\textcolor{FuncColor}{$\triangleright$\enspace\texttt{NiceMonomorphism({\mdseries\slshape obj})\index{NiceMonomorphism@\texttt{NiceMonomorphism}}
\label{NiceMonomorphism}
}\hfill{\scriptsize (attribute)}}\\


 is a homomorphism that is defined (at least) on the whole of \mbox{\texttt{\mdseries\slshape obj}} and whose restriction to \mbox{\texttt{\mdseries\slshape obj}} is injective. The concrete morphism (and also the image group) will depend on
the representation of \mbox{\texttt{\mdseries\slshape obj}}. 

 WARNING: The domain of the homomorphism may be larger than \mbox{\texttt{\mdseries\slshape obj}}. To obtain the image of \mbox{\texttt{\mdseries\slshape obj}} under the homomorphism, use \texttt{NiceObject} (\ref{NiceObject}); see there for an example where it matters. }

 

\subsection{\textcolor{Chapter }{NiceObject}}
\logpage{[ 40, 5, 3 ]}\nobreak
\hyperdef{L}{X7B47BE0983E93A83}{}
{\noindent\textcolor{FuncColor}{$\triangleright$\enspace\texttt{NiceObject({\mdseries\slshape obj})\index{NiceObject@\texttt{NiceObject}}
\label{NiceObject}
}\hfill{\scriptsize (attribute)}}\\


 The \texttt{NiceObject} value of \mbox{\texttt{\mdseries\slshape obj}} is the image of \mbox{\texttt{\mdseries\slshape obj}} under the mapping stored as the value of \texttt{NiceMonomorphism} (\ref{NiceMonomorphism}) for \mbox{\texttt{\mdseries\slshape obj}}. 

 A typical example are finite matrix groups, which use a faithful action on
vectors to translate all calculations in a permutation group. 

 
\begin{Verbatim}[commandchars=!@|,fontsize=\small,frame=single,label=Example]
  !gapprompt@gap>| !gapinput@gl:=GL(3,2);|
  SL(3,2)
  !gapprompt@gap>| !gapinput@IsHandledByNiceMonomorphism(gl);|
  true
  !gapprompt@gap>| !gapinput@NiceObject(gl);|
  Group([ (5,7)(6,8), (2,3,5)(4,7,6) ])
  !gapprompt@gap>| !gapinput@Image(NiceMonomorphism(gl),Z(2)*[[1,0,0],[0,1,1],[1,0,1]]);|
  (2,6)(3,4,7,8)
\end{Verbatim}
 }

 

\subsection{\textcolor{Chapter }{IsCanonicalNiceMonomorphism}}
\logpage{[ 40, 5, 4 ]}\nobreak
\hyperdef{L}{X8652149F7F291EE3}{}
{\noindent\textcolor{FuncColor}{$\triangleright$\enspace\texttt{IsCanonicalNiceMonomorphism({\mdseries\slshape nhom})\index{IsCanonicalNiceMonomorphism@\texttt{IsCanonicalNiceMonomorphism}}
\label{IsCanonicalNiceMonomorphism}
}\hfill{\scriptsize (property)}}\\


 A nice monomorphism (see \texttt{NiceMonomorphism} (\ref{NiceMonomorphism}) \mbox{\texttt{\mdseries\slshape nhom}} is canonical if the image set will only depend on the set of group elements
but not on the generating set and \texttt{\texttt{\symbol{92}}{\textless}} (\ref{bSlash<}) comparison of group elements translates through the nice monomorphism. This
implies that equal objects will always have equal \texttt{NiceObject} (\ref{NiceObject}) values. In some situations however this condition would be expensive to
achieve, therefore it is not guaranteed for every nice monomorphism. }

 }

  
\section{\textcolor{Chapter }{Group Automorphisms}}\label{Group Automorphisms}
\logpage{[ 40, 6, 0 ]}
\hyperdef{L}{X783030917CB43959}{}
{
  Group automorphisms are bijective homomorphism from a group onto itself. An
important subclass are automorphisms which are induced by conjugation of the
group itself or a supergroup. 

\subsection{\textcolor{Chapter }{ConjugatorIsomorphism}}
\logpage{[ 40, 6, 1 ]}\nobreak
\hyperdef{L}{X7E52E99487562F3A}{}
{\noindent\textcolor{FuncColor}{$\triangleright$\enspace\texttt{ConjugatorIsomorphism({\mdseries\slshape G, g})\index{ConjugatorIsomorphism@\texttt{ConjugatorIsomorphism}}
\label{ConjugatorIsomorphism}
}\hfill{\scriptsize (operation)}}\\


 Let \mbox{\texttt{\mdseries\slshape G}} be a group, and \mbox{\texttt{\mdseries\slshape g}} an element in the same family as the elements of \mbox{\texttt{\mdseries\slshape G}}. \texttt{ConjugatorIsomorphism} returns the isomorphism from \mbox{\texttt{\mdseries\slshape G}} to \texttt{\mbox{\texttt{\mdseries\slshape G}}\texttt{\symbol{94}}\mbox{\texttt{\mdseries\slshape g}}} defined by $h \mapsto h^{\mbox{\texttt{\mdseries\slshape g}}}$ for all $h \in \mbox{\texttt{\mdseries\slshape G}}$. 

 If \mbox{\texttt{\mdseries\slshape g}} normalizes \mbox{\texttt{\mdseries\slshape G}} then \texttt{ConjugatorIsomorphism} does the same as \texttt{ConjugatorAutomorphismNC} (\ref{ConjugatorAutomorphismNC}). }

 

\subsection{\textcolor{Chapter }{ConjugatorAutomorphism}}
\logpage{[ 40, 6, 2 ]}\nobreak
\hyperdef{L}{X79ED68CF8139F08A}{}
{\noindent\textcolor{FuncColor}{$\triangleright$\enspace\texttt{ConjugatorAutomorphism({\mdseries\slshape G, g})\index{ConjugatorAutomorphism@\texttt{ConjugatorAutomorphism}}
\label{ConjugatorAutomorphism}
}\hfill{\scriptsize (function)}}\\
\noindent\textcolor{FuncColor}{$\triangleright$\enspace\texttt{ConjugatorAutomorphismNC({\mdseries\slshape G, g})\index{ConjugatorAutomorphismNC@\texttt{ConjugatorAutomorphismNC}}
\label{ConjugatorAutomorphismNC}
}\hfill{\scriptsize (operation)}}\\


 Let \mbox{\texttt{\mdseries\slshape G}} be a group, and \mbox{\texttt{\mdseries\slshape g}} an element in the same family as the elements of \mbox{\texttt{\mdseries\slshape G}} such that \mbox{\texttt{\mdseries\slshape g}} normalizes \mbox{\texttt{\mdseries\slshape G}}. \texttt{ConjugatorAutomorphism} returns the automorphism of \mbox{\texttt{\mdseries\slshape G}} defined by $h \mapsto h^{\mbox{\texttt{\mdseries\slshape g}}}$ for all $h \in \mbox{\texttt{\mdseries\slshape G}}$. 

 If conjugation by \mbox{\texttt{\mdseries\slshape g}} does \emph{not} leave \mbox{\texttt{\mdseries\slshape G}} invariant, \texttt{ConjugatorAutomorphism} returns \texttt{fail}; in this case, the isomorphism from \mbox{\texttt{\mdseries\slshape G}} to \texttt{\mbox{\texttt{\mdseries\slshape G}}\texttt{\symbol{94}}\mbox{\texttt{\mdseries\slshape g}}} induced by conjugation with \mbox{\texttt{\mdseries\slshape g}} can be constructed with \texttt{ConjugatorIsomorphism} (\ref{ConjugatorIsomorphism}). 

 \texttt{ConjugatorAutomorphismNC} does the same as \texttt{ConjugatorAutomorphism}, except that the check is omitted whether \mbox{\texttt{\mdseries\slshape g}} normalizes \mbox{\texttt{\mdseries\slshape G}} and it is assumed that \mbox{\texttt{\mdseries\slshape g}} is chosen to be in \mbox{\texttt{\mdseries\slshape G}} if possible. }

 

\subsection{\textcolor{Chapter }{InnerAutomorphism}}
\logpage{[ 40, 6, 3 ]}\nobreak
\hyperdef{L}{X7E937A947856D9DA}{}
{\noindent\textcolor{FuncColor}{$\triangleright$\enspace\texttt{InnerAutomorphism({\mdseries\slshape G, g})\index{InnerAutomorphism@\texttt{InnerAutomorphism}}
\label{InnerAutomorphism}
}\hfill{\scriptsize (function)}}\\
\noindent\textcolor{FuncColor}{$\triangleright$\enspace\texttt{InnerAutomorphismNC({\mdseries\slshape G, g})\index{InnerAutomorphismNC@\texttt{InnerAutomorphismNC}}
\label{InnerAutomorphismNC}
}\hfill{\scriptsize (operation)}}\\


 Let \mbox{\texttt{\mdseries\slshape G}} be a group, and $\mbox{\texttt{\mdseries\slshape g}} \in \mbox{\texttt{\mdseries\slshape G}}$. \texttt{InnerAutomorphism} returns the automorphism of \mbox{\texttt{\mdseries\slshape G}} defined by $h \mapsto h^{\mbox{\texttt{\mdseries\slshape g}}}$ for all $h \in \mbox{\texttt{\mdseries\slshape G}}$. 

 If \mbox{\texttt{\mdseries\slshape g}} is \emph{not} an element of \mbox{\texttt{\mdseries\slshape G}}, \texttt{InnerAutomorphism} returns \texttt{fail}; in this case, the isomorphism from \mbox{\texttt{\mdseries\slshape G}} to \texttt{\mbox{\texttt{\mdseries\slshape G}}\texttt{\symbol{94}}\mbox{\texttt{\mdseries\slshape g}}} induced by conjugation with \mbox{\texttt{\mdseries\slshape g}} can be constructed with \texttt{ConjugatorIsomorphism} (\ref{ConjugatorIsomorphism}) or with \texttt{ConjugatorAutomorphism} (\ref{ConjugatorAutomorphism}). 

 \texttt{InnerAutomorphismNC} does the same as \texttt{InnerAutomorphism}, except that the check is omitted whether $\mbox{\texttt{\mdseries\slshape g}} \in \mbox{\texttt{\mdseries\slshape G}}$. }

 

\subsection{\textcolor{Chapter }{IsConjugatorIsomorphism}}
\logpage{[ 40, 6, 4 ]}\nobreak
\hyperdef{L}{X7F31FECC7A3D4A8A}{}
{\noindent\textcolor{FuncColor}{$\triangleright$\enspace\texttt{IsConjugatorIsomorphism({\mdseries\slshape hom})\index{IsConjugatorIsomorphism@\texttt{IsConjugatorIsomorphism}}
\label{IsConjugatorIsomorphism}
}\hfill{\scriptsize (property)}}\\
\noindent\textcolor{FuncColor}{$\triangleright$\enspace\texttt{IsConjugatorAutomorphism({\mdseries\slshape hom})\index{IsConjugatorAutomorphism@\texttt{IsConjugatorAutomorphism}}
\label{IsConjugatorAutomorphism}
}\hfill{\scriptsize (property)}}\\
\noindent\textcolor{FuncColor}{$\triangleright$\enspace\texttt{IsInnerAutomorphism({\mdseries\slshape hom})\index{IsInnerAutomorphism@\texttt{IsInnerAutomorphism}}
\label{IsInnerAutomorphism}
}\hfill{\scriptsize (property)}}\\


 Let \mbox{\texttt{\mdseries\slshape hom}} be a group general mapping (see{\nobreakspace}\texttt{IsGroupGeneralMapping} (\ref{IsGroupGeneralMapping})) with source $G$. \texttt{IsConjugatorIsomorphism} returns \texttt{true} if \mbox{\texttt{\mdseries\slshape hom}} is induced by conjugation of $G$ by an element $g$ that lies in $G$ or in a group into which $G$ is naturally embedded in the sense described below, and \texttt{false} otherwise. 

 Natural embeddings are dealt with in the case that $G$ is a permutation group (see Chapter{\nobreakspace}\ref{Permutation Groups}), a matrix group (see Chapter{\nobreakspace}\ref{Matrix Groups}), a finitely presented group (see Chapter{\nobreakspace}\ref{Finitely Presented Groups}), or a group given w.r.t.{\nobreakspace}a polycyclic presentation (see
Chapter{\nobreakspace}\ref{Pc Groups}). In all other cases, \texttt{IsConjugatorIsomorphism} may return \texttt{false} if \mbox{\texttt{\mdseries\slshape hom}} is induced by conjugation but is not an inner automorphism. 

 If \texttt{IsConjugatorIsomorphism} returns \texttt{true} for \mbox{\texttt{\mdseries\slshape hom}} then an element $g$ that induces \mbox{\texttt{\mdseries\slshape hom}} can be accessed as value of the attribute \texttt{ConjugatorOfConjugatorIsomorphism} (\ref{ConjugatorOfConjugatorIsomorphism}). 

 \texttt{IsConjugatorAutomorphism} returns \texttt{true} if \mbox{\texttt{\mdseries\slshape hom}} is an automorphism (see{\nobreakspace}\texttt{IsEndoGeneralMapping} (\ref{IsEndoGeneralMapping})) that is regarded as a conjugator isomorphism by \texttt{IsConjugatorIsomorphism}, and \texttt{false} otherwise. 

 \texttt{IsInnerAutomorphism} returns \texttt{true} if \mbox{\texttt{\mdseries\slshape hom}} is a conjugator automorphism such that an element $g$ inducing \mbox{\texttt{\mdseries\slshape hom}} can be chosen in $G$, and \texttt{false} otherwise. }

 

\subsection{\textcolor{Chapter }{ConjugatorOfConjugatorIsomorphism}}
\logpage{[ 40, 6, 5 ]}\nobreak
\hyperdef{L}{X78FE7E307E86525A}{}
{\noindent\textcolor{FuncColor}{$\triangleright$\enspace\texttt{ConjugatorOfConjugatorIsomorphism({\mdseries\slshape hom})\index{ConjugatorOfConjugatorIsomorphism@\texttt{ConjugatorOfConjugatorIsomorphism}}
\label{ConjugatorOfConjugatorIsomorphism}
}\hfill{\scriptsize (attribute)}}\\


 For a conjugator isomorphism \mbox{\texttt{\mdseries\slshape hom}} (see{\nobreakspace}\texttt{ConjugatorIsomorphism} (\ref{ConjugatorIsomorphism})), \texttt{ConjugatorOfConjugatorIsomorphism} returns an element $g$ such that mapping under \mbox{\texttt{\mdseries\slshape hom}} is induced by conjugation with $g$. 

 To avoid problems with \texttt{IsInnerAutomorphism} (\ref{IsInnerAutomorphism}), it is guaranteed that the conjugator is taken from the source of \mbox{\texttt{\mdseries\slshape hom}} if possible. 

 
\begin{Verbatim}[commandchars=!@|,fontsize=\small,frame=single,label=Example]
  !gapprompt@gap>| !gapinput@hgens:=[(1,2,3),(1,2,4)];;h:=Group(hgens);;|
  !gapprompt@gap>| !gapinput@hom:=GroupHomomorphismByImages(h,h,hgens,[(1,2,3),(2,3,4)]);;|
  !gapprompt@gap>| !gapinput@IsInnerAutomorphism(hom);|
  true
  !gapprompt@gap>| !gapinput@ConjugatorOfConjugatorIsomorphism(hom);|
  (1,2,3)
  !gapprompt@gap>| !gapinput@hom:=GroupHomomorphismByImages(h,h,hgens,[(1,3,2),(1,4,2)]);|
  [ (1,2,3), (1,2,4) ] -> [ (1,3,2), (1,4,2) ]
  !gapprompt@gap>| !gapinput@IsInnerAutomorphism(hom);|
  false
  !gapprompt@gap>| !gapinput@IsConjugatorAutomorphism(hom);|
  true
  !gapprompt@gap>| !gapinput@ConjugatorOfConjugatorIsomorphism(hom);|
  (1,2)
\end{Verbatim}
 }

 }

  
\section{\textcolor{Chapter }{Groups of Automorphisms}}\label{Groups of Automorphisms}
\logpage{[ 40, 7, 0 ]}
\hyperdef{L}{X79640F3682BDBFC1}{}
{
  Group automorphism can be multiplied and inverted and thus it is possible to
form groups of automorphisms. 

\subsection{\textcolor{Chapter }{AutomorphismGroup}}
\logpage{[ 40, 7, 1 ]}\nobreak
\hyperdef{L}{X87677B0787B4461A}{}
{\noindent\textcolor{FuncColor}{$\triangleright$\enspace\texttt{AutomorphismGroup({\mdseries\slshape G})\index{AutomorphismGroup@\texttt{AutomorphismGroup}}
\label{AutomorphismGroup}
}\hfill{\scriptsize (attribute)}}\\


 returns the full automorphism group of the group \mbox{\texttt{\mdseries\slshape G}}. The automorphisms act on \mbox{\texttt{\mdseries\slshape G}} by the caret operator \texttt{\texttt{\symbol{94}}}. The automorphism group often stores a \texttt{NiceMonomorphism} (\ref{NiceMonomorphism}) value whose image is a permutation group, obtained by the action on a subset
of \mbox{\texttt{\mdseries\slshape G}}. 

 Note that current methods for the calculation of the automorphism group of a
group $G$ require $G$ to be a permutation group or a pc group to be efficient. For groups in other
representations the calculation is likely very slow. 

 Also, the \textsf{AutPGrp} package installs enhanced methods for \texttt{AutomorphismGroup} for finite $p$\texttt{\symbol{45}}groups, and the \textsf{FGA} package \texttt{\symbol{45}} for finitely generated subgroups of free groups. 

 Methods may be installed for \texttt{AutomorphismGroup} for other domains, such as e.g. for linear codes in the \textsf{GUAVA} package, loops in the \textsf{loops} package and nilpotent Lie algebras in the \textsf{Sophus} package (see package manuals for their descriptions). }

 

\subsection{\textcolor{Chapter }{IsGroupOfAutomorphisms}}
\logpage{[ 40, 7, 2 ]}\nobreak
\hyperdef{L}{X7FC631B786C1DC8B}{}
{\noindent\textcolor{FuncColor}{$\triangleright$\enspace\texttt{IsGroupOfAutomorphisms({\mdseries\slshape G})\index{IsGroupOfAutomorphisms@\texttt{IsGroupOfAutomorphisms}}
\label{IsGroupOfAutomorphisms}
}\hfill{\scriptsize (property)}}\\


 indicates whether \mbox{\texttt{\mdseries\slshape G}} consists of automorphisms of another group $H$. The group $H$ can be obtained from \mbox{\texttt{\mdseries\slshape G}} via the attribute \texttt{AutomorphismDomain} (\ref{AutomorphismDomain}). }

 

\subsection{\textcolor{Chapter }{AutomorphismDomain}}
\logpage{[ 40, 7, 3 ]}\nobreak
\hyperdef{L}{X7B767B9D827EB0FC}{}
{\noindent\textcolor{FuncColor}{$\triangleright$\enspace\texttt{AutomorphismDomain({\mdseries\slshape G})\index{AutomorphismDomain@\texttt{AutomorphismDomain}}
\label{AutomorphismDomain}
}\hfill{\scriptsize (attribute)}}\\


 If \mbox{\texttt{\mdseries\slshape G}} consists of automorphisms of $H$, this attribute returns $H$. }

 

\subsection{\textcolor{Chapter }{IsAutomorphismGroup}}
\logpage{[ 40, 7, 4 ]}\nobreak
\hyperdef{L}{X7F87D5957D9B991E}{}
{\noindent\textcolor{FuncColor}{$\triangleright$\enspace\texttt{IsAutomorphismGroup({\mdseries\slshape G})\index{IsAutomorphismGroup@\texttt{IsAutomorphismGroup}}
\label{IsAutomorphismGroup}
}\hfill{\scriptsize (property)}}\\


 indicates whether \mbox{\texttt{\mdseries\slshape G}}, which must be \texttt{IsGroupOfAutomorphisms} (\ref{IsGroupOfAutomorphisms}), is the full automorphism group of another group $H$, this group is given as \texttt{AutomorphismDomain} (\ref{AutomorphismDomain}) value of \mbox{\texttt{\mdseries\slshape G}}. 
\begin{Verbatim}[commandchars=!@|,fontsize=\small,frame=single,label=Example]
  !gapprompt@gap>| !gapinput@g:=Group((1,2,3,4),(1,3));|
  Group([ (1,2,3,4), (1,3) ])
  !gapprompt@gap>| !gapinput@au:=AutomorphismGroup(g);|
  <group of size 8 with 3 generators>
  !gapprompt@gap>| !gapinput@GeneratorsOfGroup(au);|
  [ Pcgs([ (2,4), (1,2,3,4), (1,3)(2,4) ]) ->
      [ (1,2)(3,4), (1,2,3,4), (1,3)(2,4) ],
    Pcgs([ (2,4), (1,2,3,4), (1,3)(2,4) ]) ->
      [ (1,3), (1,2,3,4), (1,3)(2,4) ],
    Pcgs([ (2,4), (1,2,3,4), (1,3)(2,4) ]) ->
      [ (2,4), (1,4,3,2), (1,3)(2,4) ] ]
  !gapprompt@gap>| !gapinput@NiceObject(au);|
  Group([ (1,2,4,6), (1,4)(2,6), (2,6)(3,5) ])
\end{Verbatim}
 }

 

\subsection{\textcolor{Chapter }{InnerAutomorphismsAutomorphismGroup}}
\logpage{[ 40, 7, 5 ]}\nobreak
\hyperdef{L}{X8476738A7BF9BADA}{}
{\noindent\textcolor{FuncColor}{$\triangleright$\enspace\texttt{InnerAutomorphismsAutomorphismGroup({\mdseries\slshape autgroup})\index{InnerAutomorphismsAutomorphismGroup@\texttt{InnerAutomorphismsAutomorphismGroup}}
\label{InnerAutomorphismsAutomorphismGroup}
}\hfill{\scriptsize (attribute)}}\\


 For an automorphism group \mbox{\texttt{\mdseries\slshape autgroup}} of a group this attribute stores the subgroup of inner automorphisms
(automorphisms induced by conjugation) of the original group. 
\begin{Verbatim}[commandchars=!@|,fontsize=\small,frame=single,label=Example]
  !gapprompt@gap>| !gapinput@InnerAutomorphismsAutomorphismGroup(au);|
  <group with 2 generators>
\end{Verbatim}
 }

 

\subsection{\textcolor{Chapter }{InnerAutomorphismGroup}}
\logpage{[ 40, 7, 6 ]}\nobreak
\hyperdef{L}{X7957AC21782B6C8C}{}
{\noindent\textcolor{FuncColor}{$\triangleright$\enspace\texttt{InnerAutomorphismGroup({\mdseries\slshape G})\index{InnerAutomorphismGroup@\texttt{InnerAutomorphismGroup}}
\label{InnerAutomorphismGroup}
}\hfill{\scriptsize (attribute)}}\\


 For a group \mbox{\texttt{\mdseries\slshape G}} this attribute stores the group of inner automorphisms (automorphisms induced
by conjugation) of the original group. 
\begin{Verbatim}[commandchars=!@|,fontsize=\small,frame=single,label=Example]
  !gapprompt@gap>| !gapinput@InnerAutomorphismGroup(SymmetricGroup(5));|
  <group with 2 generators>
\end{Verbatim}
 }

 

\subsection{\textcolor{Chapter }{InducedAutomorphism}}
\logpage{[ 40, 7, 7 ]}\nobreak
\hyperdef{L}{X7FC9B6EA7CAADC0A}{}
{\noindent\textcolor{FuncColor}{$\triangleright$\enspace\texttt{InducedAutomorphism({\mdseries\slshape epi, aut})\index{InducedAutomorphism@\texttt{InducedAutomorphism}}
\label{InducedAutomorphism}
}\hfill{\scriptsize (function)}}\\


 Let \mbox{\texttt{\mdseries\slshape aut}} be an automorphism of a group $G$ and \mbox{\texttt{\mdseries\slshape epi}} be a homomorphism from $G$ to a group $H$ such that the kernel of \mbox{\texttt{\mdseries\slshape epi}} is fixed under \mbox{\texttt{\mdseries\slshape aut}}. Let $U$ be the image of \mbox{\texttt{\mdseries\slshape epi}}. This command returns the automorphism of $U$ induced by \mbox{\texttt{\mdseries\slshape aut}} via \mbox{\texttt{\mdseries\slshape epi}}, that is, the automorphism of $U$ which maps $g$\texttt{\texttt{\symbol{94}}\mbox{\texttt{\mdseries\slshape epi}}} to \texttt{(}$g$\texttt{\texttt{\symbol{94}}\mbox{\texttt{\mdseries\slshape aut}})\texttt{\symbol{94}}\mbox{\texttt{\mdseries\slshape epi}}}, for $g \in G$. 
\begin{Verbatim}[commandchars=!@|,fontsize=\small,frame=single,label=Example]
  !gapprompt@gap>| !gapinput@g:=Group((1,2,3,4),(1,2));|
  Group([ (1,2,3,4), (1,2) ])
  !gapprompt@gap>| !gapinput@n:=Subgroup(g,[(1,2)(3,4),(1,3)(2,4)]);|
  Group([ (1,2)(3,4), (1,3)(2,4) ])
  !gapprompt@gap>| !gapinput@epi:=NaturalHomomorphismByNormalSubgroup(g,n);|
  [ (1,2,3,4), (1,2) ] -> [ f1*f2, f1 ]
  !gapprompt@gap>| !gapinput@aut:=InnerAutomorphism(g,(1,2,3));|
  ^(1,2,3)
  !gapprompt@gap>| !gapinput@InducedAutomorphism(epi,aut);|
  ^f2
\end{Verbatim}
 }

   }

  
\section{\textcolor{Chapter }{Calculating with Group Automorphisms}}\label{Calculating with Group Automorphisms}
\logpage{[ 40, 8, 0 ]}
\hyperdef{L}{X7A8E961C7F1A57B3}{}
{
  Usually the best way to calculate in a group of automorphisms is to translate
all calculations to an isomorphic group in a representation for which better
algorithms are available, such as a permutation group. This translation can be
done automatically using \texttt{NiceMonomorphism} (\ref{NiceMonomorphism}). 

 Once a group knows to be a group of automorphisms (this can be achieved by
testing or setting the property \texttt{IsGroupOfAutomorphisms} (\ref{IsGroupOfAutomorphisms})), \textsf{GAP} will try itself to find such a nice monomorphism once calculations in the
automorphism group are done. 

 Note that nice homomorphisms inherit down to subgroups, but cannot necessarily
be extended from a subgroup to the whole group. Thus when working with a group
of automorphisms, it can be beneficial to enforce calculation of the nice
monomorphism for the whole group (for example by explicitly calling \texttt{Random} (\ref{Random:for a list or collection}) and ignoring the result {\textendash}it will be stored internally) at the
start of the calculation. Otherwise \textsf{GAP} might first calculate a nice monomorphism for the subgroup, only to be forced
to calculate a new nice monomorphism for the whole group later on. 

 If a good domain for a faithful permutation action is known already, a
homomorphism for the action on it can be created using \texttt{NiceMonomorphismAutomGroup} (\ref{NiceMonomorphismAutomGroup}). It might be stored by \texttt{SetNiceMonomorphism} (see \texttt{NiceMonomorphism} (\ref{NiceMonomorphism})). 

 Another nice way of representing automorphisms as permutations has been
described in \cite{Sims97}. It is not yet available in \textsf{GAP}, a description however can be found in section {\nobreakspace}\ref{Stabilizer Chains for Automorphisms Acting on Enumerators}. 

\subsection{\textcolor{Chapter }{AssignNiceMonomorphismAutomorphismGroup}}
\logpage{[ 40, 8, 1 ]}\nobreak
\hyperdef{L}{X85691E8386107403}{}
{\noindent\textcolor{FuncColor}{$\triangleright$\enspace\texttt{AssignNiceMonomorphismAutomorphismGroup({\mdseries\slshape autgrp, group})\index{AssignNiceMonomorphismAutomorphismGroup@\texttt{Assign}\-\texttt{Nice}\-\texttt{Monomorphism}\-\texttt{Automorphism}\-\texttt{Group}}
\label{AssignNiceMonomorphismAutomorphismGroup}
}\hfill{\scriptsize (function)}}\\


 computes a nice monomorphism for \mbox{\texttt{\mdseries\slshape autgroup}} acting on \mbox{\texttt{\mdseries\slshape group}} and stores it as \texttt{NiceMonomorphism} (\ref{NiceMonomorphism}) value of \mbox{\texttt{\mdseries\slshape autgrp}}. 

 If the centre of \texttt{AutomorphismDomain} (\ref{AutomorphismDomain}) of \mbox{\texttt{\mdseries\slshape autgrp}} is trivial, the operation will first try to represent all automorphisms by
conjugation (in \mbox{\texttt{\mdseries\slshape group}} or in a natural parent of \mbox{\texttt{\mdseries\slshape group}}). 

 If this fails the operation tries to find a small subset of \mbox{\texttt{\mdseries\slshape group}} on which the action will be faithful. 

 The operation sets the attribute \texttt{NiceMonomorphism} (\ref{NiceMonomorphism}) and does not return a value. }

 

\subsection{\textcolor{Chapter }{NiceMonomorphismAutomGroup}}
\logpage{[ 40, 8, 2 ]}\nobreak
\hyperdef{L}{X7C9FB0A57BFF6CC0}{}
{\noindent\textcolor{FuncColor}{$\triangleright$\enspace\texttt{NiceMonomorphismAutomGroup({\mdseries\slshape autgrp, elms, elmsgens})\index{NiceMonomorphismAutomGroup@\texttt{NiceMonomorphismAutomGroup}}
\label{NiceMonomorphismAutomGroup}
}\hfill{\scriptsize (function)}}\\


 This function creates a monomorphism for an automorphism group \mbox{\texttt{\mdseries\slshape autgrp}} of a group by permuting the group elements in the list \mbox{\texttt{\mdseries\slshape elms}}. This list must be chosen to yield a faithful representation. \mbox{\texttt{\mdseries\slshape elmsgens}} is a list of generators which are a subset of \mbox{\texttt{\mdseries\slshape elms}}. (They can differ from the group's original generators.) It does not yet
assign it as \texttt{NiceMonomorphism} (\ref{NiceMonomorphism}) value. }

 }

  
\section{\textcolor{Chapter }{Searching for Homomorphisms}}\label{Searching for Homomorphisms}
\logpage{[ 40, 9, 0 ]}
\hyperdef{L}{X81B79CC27F47D429}{}
{
  \index{homomorphisms!find all} 

\subsection{\textcolor{Chapter }{IsomorphismGroups}}
\logpage{[ 40, 9, 1 ]}\nobreak
\hyperdef{L}{X7B536A32827788C6}{}
{\noindent\textcolor{FuncColor}{$\triangleright$\enspace\texttt{IsomorphismGroups({\mdseries\slshape G, H})\index{IsomorphismGroups@\texttt{IsomorphismGroups}}
\label{IsomorphismGroups}
}\hfill{\scriptsize (function)}}\\


 computes an isomorphism between the groups \mbox{\texttt{\mdseries\slshape G}} and \mbox{\texttt{\mdseries\slshape H}} if they are isomorphic and returns \texttt{fail} otherwise. 

 With the existing methods the amount of time needed grows with the size of a
generating system of \mbox{\texttt{\mdseries\slshape G}}. (Thus in particular for $p$\texttt{\symbol{45}}groups calculations can be slow.) If you do only need to
know whether groups are isomorphic, you might want to consider \texttt{IdGroup} (\textbf{smallgrp: IdGroup}) or the random isomorphism test (see{\nobreakspace}\texttt{RandomIsomorphismTest} (\ref{RandomIsomorphismTest})). 

 \index{isomorphisms!find all} 
\begin{Verbatim}[commandchars=!@|,fontsize=\small,frame=single,label=Example]
  !gapprompt@gap>| !gapinput@g:=Group((1,2,3,4),(1,3));;|
  !gapprompt@gap>| !gapinput@h:=Group((1,4,6,7)(2,3,5,8), (1,5)(2,6)(3,4)(7,8));;|
  !gapprompt@gap>| !gapinput@IsomorphismGroups(g,h);|
  [ (1,2,3,4), (1,3) ] -> [ (1,4,6,7)(2,3,5,8), (1,2)(3,7)(4,8)(5,6) ]
  !gapprompt@gap>| !gapinput@IsomorphismGroups(g,Group((1,2,3,4),(1,2)));|
  fail
\end{Verbatim}
 }

 

\subsection{\textcolor{Chapter }{AllHomomorphismClasses}}
\logpage{[ 40, 9, 2 ]}\nobreak
\hyperdef{L}{X7D0C3D5E864CE954}{}
{\noindent\textcolor{FuncColor}{$\triangleright$\enspace\texttt{AllHomomorphismClasses({\mdseries\slshape G, H})\index{AllHomomorphismClasses@\texttt{AllHomomorphismClasses}}
\label{AllHomomorphismClasses}
}\hfill{\scriptsize (function)}}\\


 For two finite groups \mbox{\texttt{\mdseries\slshape G}} and \mbox{\texttt{\mdseries\slshape H}}, this function returns representatives of all homomorphisms $\mbox{\texttt{\mdseries\slshape G}} to \mbox{\texttt{\mdseries\slshape H}}$ up to \mbox{\texttt{\mdseries\slshape H}}\texttt{\symbol{45}}conjugacy. 
\begin{Verbatim}[commandchars=!@|,fontsize=\small,frame=single,label=Example]
  !gapprompt@gap>| !gapinput@AllHomomorphismClasses(SymmetricGroup(4),SymmetricGroup(3));|
  [ [ (2,4,3), (1,4,2,3) ] -> [ (), () ],
    [ (2,4,3), (1,4,2,3) ] -> [ (), (1,2) ],
    [ (2,4,3), (1,4,2,3) ] -> [ (1,2,3), (1,2) ] ]
\end{Verbatim}
 }

 

\subsection{\textcolor{Chapter }{AllHomomorphisms}}
\logpage{[ 40, 9, 3 ]}\nobreak
\hyperdef{L}{X791D12B7845610CE}{}
{\noindent\textcolor{FuncColor}{$\triangleright$\enspace\texttt{AllHomomorphisms({\mdseries\slshape G, H})\index{AllHomomorphisms@\texttt{AllHomomorphisms}}
\label{AllHomomorphisms}
}\hfill{\scriptsize (function)}}\\
\noindent\textcolor{FuncColor}{$\triangleright$\enspace\texttt{AllEndomorphisms({\mdseries\slshape G})\index{AllEndomorphisms@\texttt{AllEndomorphisms}}
\label{AllEndomorphisms}
}\hfill{\scriptsize (function)}}\\
\noindent\textcolor{FuncColor}{$\triangleright$\enspace\texttt{AllAutomorphisms({\mdseries\slshape G})\index{AllAutomorphisms@\texttt{AllAutomorphisms}}
\label{AllAutomorphisms}
}\hfill{\scriptsize (function)}}\\


 For two finite groups \mbox{\texttt{\mdseries\slshape G}} and \mbox{\texttt{\mdseries\slshape H}}, this function returns all homomorphisms $\mbox{\texttt{\mdseries\slshape G}} to \mbox{\texttt{\mdseries\slshape H}}$. Since this number will grow quickly, \texttt{AllHomomorphismClasses} (\ref{AllHomomorphismClasses}) should be used in most cases. \texttt{AllEndomorphisms} returns all homomorphisms from \mbox{\texttt{\mdseries\slshape G}} to itself, \texttt{AllAutomorphisms} returns all bijective endomorphisms. 
\begin{Verbatim}[commandchars=!@|,fontsize=\small,frame=single,label=Example]
  !gapprompt@gap>| !gapinput@AllHomomorphisms(SymmetricGroup(3),SymmetricGroup(3));|
  [ [ (2,3), (1,2,3) ] -> [ (), () ],
    [ (2,3), (1,2,3) ] -> [ (1,2), () ],
    [ (2,3), (1,2,3) ] -> [ (2,3), () ],
    [ (2,3), (1,2,3) ] -> [ (1,3), () ],
    [ (2,3), (1,2,3) ] -> [ (2,3), (1,2,3) ],
    [ (2,3), (1,2,3) ] -> [ (1,3), (1,2,3) ],
    [ (2,3), (1,2,3) ] -> [ (1,3), (1,3,2) ],
    [ (2,3), (1,2,3) ] -> [ (1,2), (1,2,3) ],
    [ (2,3), (1,2,3) ] -> [ (2,3), (1,3,2) ],
    [ (2,3), (1,2,3) ] -> [ (1,2), (1,3,2) ] ]
\end{Verbatim}
 }

 

\subsection{\textcolor{Chapter }{GQuotients}}
\logpage{[ 40, 9, 4 ]}\nobreak
\hyperdef{L}{X790C261184EEAB94}{}
{\noindent\textcolor{FuncColor}{$\triangleright$\enspace\texttt{GQuotients({\mdseries\slshape F, G})\index{GQuotients@\texttt{GQuotients}}
\label{GQuotients}
}\hfill{\scriptsize (operation)}}\\


 computes all epimorphisms from \mbox{\texttt{\mdseries\slshape F}} onto \mbox{\texttt{\mdseries\slshape G}} up to automorphisms of \mbox{\texttt{\mdseries\slshape G}}. This classifies all factor groups of \mbox{\texttt{\mdseries\slshape F}} which are isomorphic to \mbox{\texttt{\mdseries\slshape G}}. 

 With the existing methods the amount of time needed grows with the size of a
generating system of \mbox{\texttt{\mdseries\slshape G}}. (Thus in particular for $p$\texttt{\symbol{45}}groups calculations can be slow.) 

 If the \texttt{findall} option is set to \texttt{false}, the algorithm will stop once one homomorphism has been found (this can be
faster and might be sufficient if not all homomorphisms are needed). 

 \index{epimorphisms!find all} \index{projections!find all} 
\begin{Verbatim}[commandchars=!@|,fontsize=\small,frame=single,label=Example]
  !gapprompt@gap>| !gapinput@g:=Group((1,2,3,4),(1,2));|
  Group([ (1,2,3,4), (1,2) ])
  !gapprompt@gap>| !gapinput@h:=Group((1,2,3),(1,2));|
  Group([ (1,2,3), (1,2) ])
  !gapprompt@gap>| !gapinput@quo:=GQuotients(g,h);|
  [ [ (2,4), (1,4,3) ] -> [ (2,3), (1,2,3) ] ]
\end{Verbatim}
 }

 

\subsection{\textcolor{Chapter }{IsomorphicSubgroups}}
\logpage{[ 40, 9, 5 ]}\nobreak
\hyperdef{L}{X83B417BE7C508DC4}{}
{\noindent\textcolor{FuncColor}{$\triangleright$\enspace\texttt{IsomorphicSubgroups({\mdseries\slshape G, H})\index{IsomorphicSubgroups@\texttt{IsomorphicSubgroups}}
\label{IsomorphicSubgroups}
}\hfill{\scriptsize (operation)}}\\


 computes all monomorphisms from \mbox{\texttt{\mdseries\slshape H}} into \mbox{\texttt{\mdseries\slshape G}} up to \mbox{\texttt{\mdseries\slshape G}}\texttt{\symbol{45}}conjugacy of the image groups. This classifies all \mbox{\texttt{\mdseries\slshape G}}\texttt{\symbol{45}}classes of subgroups of \mbox{\texttt{\mdseries\slshape G}} which are isomorphic to \mbox{\texttt{\mdseries\slshape H}}. 

 With the existing methods, the amount of time needed grows with the size of a
generating system of \mbox{\texttt{\mdseries\slshape G}}. (Thus in particular for $p$\texttt{\symbol{45}}groups calculations can be slow.) A main use of \texttt{IsomorphicSubgroups} therefore is to find nonsolvable subgroups (which often can be generated by 2
elements). 

 (To find $p$\texttt{\symbol{45}}subgroups it is often faster to compute the subgroup
lattice of the Sylow subgroup and to use \texttt{IdGroup} (\textbf{smallgrp: IdGroup}) to identify the type of the subgroups.) 

 If the \texttt{findall} option is set to \texttt{false}, the algorithm will stop once one homomorphism has been found (this can be
faster and might be sufficient if not all homomorphisms are needed). 

 \index{embeddings!find all} \index{monomorphisms!find all} 
\begin{Verbatim}[commandchars=!@|,fontsize=\small,frame=single,label=Example]
  !gapprompt@gap>| !gapinput@g:=Group((1,2,3,4),(1,2));|
  Group([ (1,2,3,4), (1,2) ])
  !gapprompt@gap>| !gapinput@h:=Group((3,4),(1,2));;|
  !gapprompt@gap>| !gapinput@emb:=IsomorphicSubgroups(g,h);|
  [ [ (3,4), (1,2) ] -> [ (1,2), (3,4) ],
    [ (3,4), (1,2) ] -> [ (1,3)(2,4), (1,2)(3,4) ] ]
  !gapprompt@gap>| !gapinput@g1:=PSO(-1,8,2);;|
  !gapprompt@gap>| !gapinput@Length(IsomorphicSubgroups(g1,PSL(2,7)));|
  3
  !gapprompt@gap>| !gapinput@Length(IsomorphicSubgroups(g1,PSL(2,7):findall:=false));|
  1
\end{Verbatim}
 }

 

\subsection{\textcolor{Chapter }{MorClassLoop}}
\logpage{[ 40, 9, 6 ]}\nobreak
\hyperdef{L}{X7AABA9A27E30BF2B}{}
{\noindent\textcolor{FuncColor}{$\triangleright$\enspace\texttt{MorClassLoop({\mdseries\slshape range, classes, params, action})\index{MorClassLoop@\texttt{MorClassLoop}}
\label{MorClassLoop}
}\hfill{\scriptsize (function)}}\\


 This function loops over element tuples taken from \mbox{\texttt{\mdseries\slshape classes}} and checks these for properties such as generating a given group, or
fulfilling relations. This can be used to find small generating sets or all
types of Morphisms. The element tuples are used only up to inner automorphisms
as all images can be obtained easily from them by conjugation while running
through all of them usually would take too long. 

 \mbox{\texttt{\mdseries\slshape range}} is a group from which these elements are taken. The classes are given in a
list \mbox{\texttt{\mdseries\slshape classes}} which is a list of records with the following components. 
\begin{description}
\item[{\texttt{classes}}]  list of conjugacy classes 
\item[{\texttt{representative}}]  One element in the union of these classes 
\item[{\texttt{size}}]  The sum of the sizes of these classes 
\end{description}
 

 \mbox{\texttt{\mdseries\slshape params}} is a record containing the following optional components. 
\begin{description}
\item[{\texttt{gens}}]  generators that are to be mapped (for testing morphisms). The length of this
list determines the length of element tuples considered. 
\item[{\texttt{from}}]  a preimage group (that contains \texttt{gens}) 
\item[{\texttt{to}}]  image group (which might be smaller than \texttt{range}) 
\item[{\texttt{free}}]  free generators, a list of the same length than the generators \texttt{gens}. 
\item[{\texttt{rels}}]  some relations that hold among the generators \texttt{gens}. They are given as a list \texttt{[ word, order ]} where \texttt{word} is a word in the free generators \texttt{free}. 
\item[{\texttt{dom}}]  a set of elements on which automorphisms act faithfully (used to do element
tests in partial automorphism groups). 
\item[{\texttt{aut}}]  Subgroup of already known automorphisms. 
\item[{\texttt{condition}}]  A function that will be applied to the homomorphism and must return \texttt{true} for the homomorphism to be accepted. 
\end{description}
 

 \mbox{\texttt{\mdseries\slshape action}} is a number whose bit\texttt{\symbol{45}}representation indicates the
requirements which are enforced on the element tuples found, as follows. 
\begin{description}
\item[{1}]  homomorphism 
\item[{2}]  injective 
\item[{4}]  surjective 
\item[{8}]  find all (otherwise stops after the first find) 
\end{description}
 If the search is for homomorphisms, the function returns homomorphisms
obtained by mapping the given generators \texttt{gens} instead of element tuples. 

 The ``Morpheus'' algorithm used to find homomorphisms is described in \cite[Section V.5]{Hulpke96}. }

 }

  
\section{\textcolor{Chapter }{Representations for Group Homomorphisms}}\label{Representations for Group Homomorphisms}
\logpage{[ 40, 10, 0 ]}
\hyperdef{L}{X81FC3CEF85CED3DC}{}
{
  The different representations of group homomorphisms are used to indicate from
what type of group to what type of group they map and thus determine which
methods are used to compute images and preimages. 

 The information in this section is mainly relevant for implementing new
methods and not for using homomorphisms. 

\subsection{\textcolor{Chapter }{IsGroupGeneralMappingByImages}}
\logpage{[ 40, 10, 1 ]}\nobreak
\hyperdef{L}{X82B77A5F7F9EDBC7}{}
{\noindent\textcolor{FuncColor}{$\triangleright$\enspace\texttt{IsGroupGeneralMappingByImages({\mdseries\slshape map})\index{IsGroupGeneralMappingByImages@\texttt{IsGroupGeneralMappingByImages}}
\label{IsGroupGeneralMappingByImages}
}\hfill{\scriptsize (Representation)}}\\


 Representation for mappings from one group to another that are defined by
extending a mapping of group generators homomorphically. Instead of record
components, the attribute \texttt{MappingGeneratorsImages} (\ref{MappingGeneratorsImages}) is used to store generators and their images. }

 

\subsection{\textcolor{Chapter }{MappingGeneratorsImages}}
\logpage{[ 40, 10, 2 ]}\nobreak
\hyperdef{L}{X863805187A24B5E3}{}
{\noindent\textcolor{FuncColor}{$\triangleright$\enspace\texttt{MappingGeneratorsImages({\mdseries\slshape map})\index{MappingGeneratorsImages@\texttt{MappingGeneratorsImages}}
\label{MappingGeneratorsImages}
}\hfill{\scriptsize (attribute)}}\\


 This attribute contains a list of length 2, the first entry being a list of
generators of the source of \mbox{\texttt{\mdseries\slshape map}} and the second entry a list of their images. This attribute is used, for
example, by \texttt{GroupHomomorphismByImages} (\ref{GroupHomomorphismByImages}) to store generators and images.   }

 

\subsection{\textcolor{Chapter }{IsGroupGeneralMappingByAsGroupGeneralMappingByImages}}
\logpage{[ 40, 10, 3 ]}\nobreak
\hyperdef{L}{X7DFBBAB18126B4D9}{}
{\noindent\textcolor{FuncColor}{$\triangleright$\enspace\texttt{IsGroupGeneralMappingByAsGroupGeneralMappingByImages({\mdseries\slshape map})\index{IsGroupGeneralMappingByAsGroupGeneralMappingByImages@\texttt{IsGroup}\-\texttt{General}\-\texttt{Mapping}\-\texttt{By}\-\texttt{As}\-\texttt{Group}\-\texttt{General}\-\texttt{Mapping}\-\texttt{By}\-\texttt{Images}}
\label{IsGroupGeneralMappingByAsGroupGeneralMappingByImages}
}\hfill{\scriptsize (Representation)}}\\


 Representation for mappings that delegate work on a \texttt{GroupHomomorphismByImages} (\ref{GroupHomomorphismByImages}). }

 

\subsection{\textcolor{Chapter }{IsPreimagesByAsGroupGeneralMappingByImages}}
\logpage{[ 40, 10, 4 ]}\nobreak
\hyperdef{L}{X78707DF57C3927EB}{}
{\noindent\textcolor{FuncColor}{$\triangleright$\enspace\texttt{IsPreimagesByAsGroupGeneralMappingByImages({\mdseries\slshape map})\index{IsPreimagesByAsGroupGeneralMappingByImages@\texttt{IsPreimages}\-\texttt{By}\-\texttt{As}\-\texttt{Group}\-\texttt{General}\-\texttt{Mapping}\-\texttt{By}\-\texttt{Images}}
\label{IsPreimagesByAsGroupGeneralMappingByImages}
}\hfill{\scriptsize (Representation)}}\\


 Representation for mappings that delegate work for preimages to a mapping
created with \texttt{GroupHomomorphismByImages} (\ref{GroupHomomorphismByImages}). }

 

\subsection{\textcolor{Chapter }{IsPermGroupGeneralMapping}}
\logpage{[ 40, 10, 5 ]}\nobreak
\hyperdef{L}{X83E10338798F552B}{}
{\noindent\textcolor{FuncColor}{$\triangleright$\enspace\texttt{IsPermGroupGeneralMapping({\mdseries\slshape map})\index{IsPermGroupGeneralMapping@\texttt{IsPermGroupGeneralMapping}}
\label{IsPermGroupGeneralMapping}
}\hfill{\scriptsize (Representation)}}\\
\noindent\textcolor{FuncColor}{$\triangleright$\enspace\texttt{IsPermGroupGeneralMappingByImages({\mdseries\slshape map})\index{IsPermGroupGeneralMappingByImages@\texttt{IsPermGroupGeneralMappingByImages}}
\label{IsPermGroupGeneralMappingByImages}
}\hfill{\scriptsize (Representation)}}\\
\noindent\textcolor{FuncColor}{$\triangleright$\enspace\texttt{IsPermGroupHomomorphism({\mdseries\slshape map})\index{IsPermGroupHomomorphism@\texttt{IsPermGroupHomomorphism}}
\label{IsPermGroupHomomorphism}
}\hfill{\scriptsize (Representation)}}\\
\noindent\textcolor{FuncColor}{$\triangleright$\enspace\texttt{IsPermGroupHomomorphismByImages({\mdseries\slshape map})\index{IsPermGroupHomomorphismByImages@\texttt{IsPermGroupHomomorphismByImages}}
\label{IsPermGroupHomomorphismByImages}
}\hfill{\scriptsize (Representation)}}\\


 are the representations for mappings that map from a perm group }

 

\subsection{\textcolor{Chapter }{IsToPermGroupGeneralMappingByImages}}
\logpage{[ 40, 10, 6 ]}\nobreak
\hyperdef{L}{X83DADD9F7CAD829B}{}
{\noindent\textcolor{FuncColor}{$\triangleright$\enspace\texttt{IsToPermGroupGeneralMappingByImages({\mdseries\slshape map})\index{IsToPermGroupGeneralMappingByImages@\texttt{IsToPermGroupGeneralMappingByImages}}
\label{IsToPermGroupGeneralMappingByImages}
}\hfill{\scriptsize (Representation)}}\\
\noindent\textcolor{FuncColor}{$\triangleright$\enspace\texttt{IsToPermGroupHomomorphismByImages({\mdseries\slshape map})\index{IsToPermGroupHomomorphismByImages@\texttt{IsToPermGroupHomomorphismByImages}}
\label{IsToPermGroupHomomorphismByImages}
}\hfill{\scriptsize (Representation)}}\\


 is the representation for mappings that map to a perm group }

 

\subsection{\textcolor{Chapter }{IsGroupGeneralMappingByPcgs}}
\logpage{[ 40, 10, 7 ]}\nobreak
\hyperdef{L}{X798E72E77EC85D4A}{}
{\noindent\textcolor{FuncColor}{$\triangleright$\enspace\texttt{IsGroupGeneralMappingByPcgs({\mdseries\slshape map})\index{IsGroupGeneralMappingByPcgs@\texttt{IsGroupGeneralMappingByPcgs}}
\label{IsGroupGeneralMappingByPcgs}
}\hfill{\scriptsize (Representation)}}\\


 is the representations for mappings that map a pcgs to images and thus may use
exponents to decompose generators. }

 

\subsection{\textcolor{Chapter }{IsPcGroupGeneralMappingByImages}}
\logpage{[ 40, 10, 8 ]}\nobreak
\hyperdef{L}{X86FF63B784FB8F85}{}
{\noindent\textcolor{FuncColor}{$\triangleright$\enspace\texttt{IsPcGroupGeneralMappingByImages({\mdseries\slshape map})\index{IsPcGroupGeneralMappingByImages@\texttt{IsPcGroupGeneralMappingByImages}}
\label{IsPcGroupGeneralMappingByImages}
}\hfill{\scriptsize (Representation)}}\\
\noindent\textcolor{FuncColor}{$\triangleright$\enspace\texttt{IsPcGroupHomomorphismByImages({\mdseries\slshape map})\index{IsPcGroupHomomorphismByImages@\texttt{IsPcGroupHomomorphismByImages}}
\label{IsPcGroupHomomorphismByImages}
}\hfill{\scriptsize (Representation)}}\\


 is the representation for mappings from a pc group }

 

\subsection{\textcolor{Chapter }{IsToPcGroupGeneralMappingByImages}}
\logpage{[ 40, 10, 9 ]}\nobreak
\hyperdef{L}{X79A853B579B250C0}{}
{\noindent\textcolor{FuncColor}{$\triangleright$\enspace\texttt{IsToPcGroupGeneralMappingByImages({\mdseries\slshape map})\index{IsToPcGroupGeneralMappingByImages@\texttt{IsToPcGroupGeneralMappingByImages}}
\label{IsToPcGroupGeneralMappingByImages}
}\hfill{\scriptsize (Representation)}}\\
\noindent\textcolor{FuncColor}{$\triangleright$\enspace\texttt{IsToPcGroupHomomorphismByImages({\mdseries\slshape map})\index{IsToPcGroupHomomorphismByImages@\texttt{IsToPcGroupHomomorphismByImages}}
\label{IsToPcGroupHomomorphismByImages}
}\hfill{\scriptsize (Representation)}}\\


 is the representation for mappings to a pc group }

 

\subsection{\textcolor{Chapter }{IsFromFpGroupGeneralMappingByImages}}
\logpage{[ 40, 10, 10 ]}\nobreak
\hyperdef{L}{X7BE2A2EB80DC5CFF}{}
{\noindent\textcolor{FuncColor}{$\triangleright$\enspace\texttt{IsFromFpGroupGeneralMappingByImages({\mdseries\slshape map})\index{IsFromFpGroupGeneralMappingByImages@\texttt{IsFromFpGroupGeneralMappingByImages}}
\label{IsFromFpGroupGeneralMappingByImages}
}\hfill{\scriptsize (Representation)}}\\
\noindent\textcolor{FuncColor}{$\triangleright$\enspace\texttt{IsFromFpGroupHomomorphismByImages({\mdseries\slshape map})\index{IsFromFpGroupHomomorphismByImages@\texttt{IsFromFpGroupHomomorphismByImages}}
\label{IsFromFpGroupHomomorphismByImages}
}\hfill{\scriptsize (Representation)}}\\


 is the representation of mappings from an fp group. }

 

\subsection{\textcolor{Chapter }{IsFromFpGroupStdGensGeneralMappingByImages}}
\logpage{[ 40, 10, 11 ]}\nobreak
\hyperdef{L}{X81090C207F4F6423}{}
{\noindent\textcolor{FuncColor}{$\triangleright$\enspace\texttt{IsFromFpGroupStdGensGeneralMappingByImages({\mdseries\slshape map})\index{IsFromFpGroupStdGensGeneralMappingByImages@\texttt{IsFrom}\-\texttt{Fp}\-\texttt{Group}\-\texttt{Std}\-\texttt{Gens}\-\texttt{General}\-\texttt{Mapping}\-\texttt{By}\-\texttt{Images}}
\label{IsFromFpGroupStdGensGeneralMappingByImages}
}\hfill{\scriptsize (Representation)}}\\
\noindent\textcolor{FuncColor}{$\triangleright$\enspace\texttt{IsFromFpGroupStdGensHomomorphismByImages({\mdseries\slshape map})\index{IsFromFpGroupStdGensHomomorphismByImages@\texttt{IsFrom}\-\texttt{Fp}\-\texttt{Group}\-\texttt{Std}\-\texttt{Gens}\-\texttt{Homomorphism}\-\texttt{By}\-\texttt{Images}}
\label{IsFromFpGroupStdGensHomomorphismByImages}
}\hfill{\scriptsize (Representation)}}\\


 is the representation of total mappings from an fp group that give images of
the standard generators. }

 }

 }

             
\chapter{\textcolor{Chapter }{Group Actions}}\label{Group Actions}
\logpage{[ 41, 0, 0 ]}
\hyperdef{L}{X87115591851FB7F4}{}
{
  \index{group actions} A \emph{group action} is a triple $(G, \Omega, \mu)$, where $G$ is a group, $\Omega$ a set and $\mu \colon \Omega \times G \rightarrow \Omega$ a function that is compatible with the group arithmetic. We call $\Omega$ the \emph{domain} of the action. 

 In \textsf{GAP}, $\Omega$ can be a duplicate\texttt{\symbol{45}}free collection (an object that permits
access to its elements via the $\Omega[n]$ operation, for example a list), it does not need to be sorted
(see{\nobreakspace}\texttt{IsSet} (\ref{IsSet})). 

 The acting function $\mu$ is a binary \textsf{GAP} function that returns the image $\mu( x, g )$ for a point $x \in \Omega$ and a group element $g \in G$. 

 In \textsf{GAP}, groups always act from the right, that is $\mu( \mu( x, g ), h ) = \mu( x, gh )$. 

 \textsf{GAP} does not test whether the acting function $\mu$ satisfies the conditions for a group operation but silently assumes that is
does. (If it does not, results are unpredictable.) 

 The first section of this chapter, \ref{About Group Actions}, describes the various ways how operations for group actions can be called. 

 Functions for several commonly used action are already built into \textsf{GAP}. These are listed in section{\nobreakspace}\ref{Basic Actions}. 

 The sections \ref{The Permutation Image of an Action} and \ref{Action of a group on itself} describe homomorphisms and mappings associated to group actions as well as the
permutation group image of an action. 

 The other sections then describe operations to compute orbits, stabilizers, as
well as properties of actions. 

 Finally section{\nobreakspace}\ref{External Sets} describes the concept of ``external sets'' which represent the concept of a \emph{$G$\texttt{\symbol{45}}set} and underly the actions mechanism.  
\section{\textcolor{Chapter }{About Group Actions}}\label{About Group Actions}
\logpage{[ 41, 1, 0 ]}
\hyperdef{L}{X83661AFD7B7BD1D9}{}
{
  \index{group actions!operations syntax} The syntax which is used by the operations for group actions is quite
flexible. For example we can call the operation \texttt{OrbitsDomain} (\ref{OrbitsDomain:for a group and an action domain}) for the orbits of the group \mbox{\texttt{\mdseries\slshape G}} on the domain \mbox{\texttt{\mdseries\slshape Omega}} in the following ways: 

 
\begin{description}
\item[{\texttt{OrbitsDomain}$( G, \Omega[, \mu] )$}]  The acting function $\mu$ is optional. If it is not given, the built\texttt{\symbol{45}}in action \texttt{OnPoints} (\ref{OnPoints}) (which defines an action via the caret operator \texttt{\texttt{\symbol{94}}}) is used as a default. 
\item[{\texttt{OrbitsDomain}$( G, \Omega, gens, acts[, \mu] )$}]  This second version of \texttt{OrbitsDomain} (\ref{OrbitsDomain:for a group and an action domain}) permits one to implement an action induced by a homomorphism: If the group $H$ acts on $\Omega$ via $\mu$ and $\varphi \colon G \rightarrow H$ is a homomorphism, $G$ acts on $\Omega$ via the induced action $\mu'( x, g ) = \mu( x, g^{\varphi} )$. 

 Here $gens$ must be a set of generators of $G$ and $acts$ the images of $gens$ under $\varphi$. $\mu$ is the acting function for $H$. Again, the function $\mu$ is optional and \texttt{OnPoints} (\ref{OnPoints}) is used as a default. 

 The advantage of this notation is that \textsf{GAP} does not need to construct this homomorphism $\varphi$ and the range group \mbox{\texttt{\mdseries\slshape H}} as \textsf{GAP} objects. (If a small group $G$ acts via complicated objects $acts$ this otherwise could lead to performance problems.) 

 \textsf{GAP} does not test whether the mapping $gens \mapsto acts$ actually induces a homomorphism and the results are unpredictable if this is
not the case. 
\item[{\texttt{OrbitsDomain}$( xset )$}]  A third variant is to call the operation with an external set, which then
provides $G$, $\Omega$ and $\mu$. You will find more about external sets in Section{\nobreakspace}\ref{External Sets}. 
\end{description}
 

 For operations like \texttt{Stabilizer} (\ref{Stabilizer}) of course the domain must be replaced by an element of the domain of the
action. }

  
\section{\textcolor{Chapter }{Basic Actions}}\label{Basic Actions}
\logpage{[ 41, 2, 0 ]}
\hyperdef{L}{X81B8F9CD868CD953}{}
{
  \index{group actions} \index{actions} \index{group operations} \textsf{GAP} already provides acting functions for the more common actions of a group. For
built\texttt{\symbol{45}}in operations such as \texttt{Stabilizer} (\ref{Stabilizer}) special methods are available for many of these actions. 

 If one needs an action for which no acting function is provided by the library
it can be implemented via a \textsf{GAP} function that conforms to the syntax 

 \texttt{actfun( omega, g )} 

 where \texttt{omega} is an element of the action domain, \texttt{g} is an element of the acting group, and the return value is the image of \texttt{omega} under \texttt{g}. 

 For example one could define the following function that acts on pairs of
polynomials via \texttt{OnIndeterminates} (\ref{OnIndeterminates:as a permutation action}): 
\begin{Verbatim}[commandchars=!@|,fontsize=\small,frame=single,label=Example]
  OnIndeterminatesPairs:= function( polypair, g )
    return [ OnIndeterminates( polypair[1], g ),
             OnIndeterminates( polypair[2], g ) ];
  end;
\end{Verbatim}
 

 Note that this function \emph{must} implement a group action from the \emph{right}. This is not verified by \textsf{GAP} and results are unpredictable otherwise. 

\subsection{\textcolor{Chapter }{OnPoints}}
\logpage{[ 41, 2, 1 ]}\nobreak
\hyperdef{L}{X7FE417DD837987B4}{}
{\noindent\textcolor{FuncColor}{$\triangleright$\enspace\texttt{OnPoints({\mdseries\slshape pnt, g})\index{OnPoints@\texttt{OnPoints}}
\label{OnPoints}
}\hfill{\scriptsize (function)}}\\


 \index{conjugation} \index{action!by conjugation} returns \texttt{\mbox{\texttt{\mdseries\slshape pnt}} \texttt{\symbol{94}} \mbox{\texttt{\mdseries\slshape g}}}. This is for example the action of a permutation group on points, or the
action of a group on its elements via conjugation, that is, if both \mbox{\texttt{\mdseries\slshape pnt}} and \mbox{\texttt{\mdseries\slshape g}} are elements from a common group then \texttt{\mbox{\texttt{\mdseries\slshape pnt}} \texttt{\symbol{94}} \mbox{\texttt{\mdseries\slshape g}}} is equal to \mbox{\texttt{\mdseries\slshape g}}$^{{-1}}$\texttt{*}\mbox{\texttt{\mdseries\slshape pnt}}\texttt{*}\mbox{\texttt{\mdseries\slshape g}}. The action of a matrix group on vectors from the right is described by both \texttt{OnPoints} and \texttt{OnRight} (\ref{OnRight}). 
\begin{Verbatim}[commandchars=!@|,fontsize=\small,frame=single,label=Example]
  !gapprompt@gap>| !gapinput@OnPoints( 1, (1,2,3) );|
  2
  !gapprompt@gap>| !gapinput@OnPoints( (1,2), (1,2,3) );|
  (2,3)
  !gapprompt@gap>| !gapinput@g:= Group( (1,2,3), (2,3,4) );;|
  !gapprompt@gap>| !gapinput@Orbit( g, 1, OnPoints );|
  [ 1, 2, 3, 4 ]
\end{Verbatim}
 }

 

\subsection{\textcolor{Chapter }{OnRight}}
\logpage{[ 41, 2, 2 ]}\nobreak
\hyperdef{L}{X7960924D84B5B18F}{}
{\noindent\textcolor{FuncColor}{$\triangleright$\enspace\texttt{OnRight({\mdseries\slshape pnt, g})\index{OnRight@\texttt{OnRight}}
\label{OnRight}
}\hfill{\scriptsize (function)}}\\


 returns \texttt{\mbox{\texttt{\mdseries\slshape pnt}} * \mbox{\texttt{\mdseries\slshape g}}}. This is for example the action of a group on its elements via right
multiplication, or the action of a group on the cosets of a subgroup. The
action of a matrix group on vectors from the right is described by both \texttt{OnPoints} (\ref{OnPoints}) and \texttt{OnRight}. 
\begin{Verbatim}[commandchars=!@|,fontsize=\small,frame=single,label=Example]
  !gapprompt@gap>| !gapinput@OnRight( [ 1, 2 ], [ [ 1, 2 ], [ 3, 4 ] ] );|
  [ 7, 10 ]
  !gapprompt@gap>| !gapinput@OnRight( (1,2,3), (2,3,4) );|
  (1,3)(2,4)
  !gapprompt@gap>| !gapinput@g:= Group( (1,2,3), (2,3,4) );;|
  !gapprompt@gap>| !gapinput@Orbit( g, (), OnRight );|
  [ (), (1,2,3), (2,3,4), (1,3,2), (1,3)(2,4), (1,2)(3,4), (2,4,3),
    (1,4,2), (1,4,3), (1,3,4), (1,2,4), (1,4)(2,3) ]
\end{Verbatim}
 }

 

\subsection{\textcolor{Chapter }{OnLeftInverse}}
\logpage{[ 41, 2, 3 ]}\nobreak
\hyperdef{L}{X832DF5327ECA0E44}{}
{\noindent\textcolor{FuncColor}{$\triangleright$\enspace\texttt{OnLeftInverse({\mdseries\slshape pnt, g})\index{OnLeftInverse@\texttt{OnLeftInverse}}
\label{OnLeftInverse}
}\hfill{\scriptsize (function)}}\\


 returns $\mbox{\texttt{\mdseries\slshape g}}^{{-1}}$ \texttt{* \mbox{\texttt{\mdseries\slshape pnt}}}. Forming the inverse is necessary to make this a proper action, as in \textsf{GAP} groups always act from the right. 

 \texttt{OnLeftInverse} is used for example in the representation of a right coset as an external set
(see{\nobreakspace}\ref{External Sets}), that is, a right coset $Ug$ is an external set for the group $U$ acting on it via \texttt{OnLeftInverse}.) 
\begin{Verbatim}[commandchars=!@|,fontsize=\small,frame=single,label=Example]
  !gapprompt@gap>| !gapinput@OnLeftInverse( [ 1, 2 ], [ [ 1, 2 ], [ 3, 4 ] ] );|
  [ 0, 1/2 ]
  !gapprompt@gap>| !gapinput@OnLeftInverse( (1,2,3), (2,3,4) );|
  (1,2,4)
  !gapprompt@gap>| !gapinput@g:= Group( (1,2,3), (2,3,4) );;|
  !gapprompt@gap>| !gapinput@Orbit( g, (), OnLeftInverse );|
  [ (), (1,3,2), (2,4,3), (1,2,3), (1,3)(2,4), (1,2)(3,4), (2,3,4),
    (1,2,4), (1,3,4), (1,4,3), (1,4,2), (1,4)(2,3) ]
\end{Verbatim}
 }

 

\subsection{\textcolor{Chapter }{OnSets}}
\logpage{[ 41, 2, 4 ]}\nobreak
\hyperdef{L}{X85AA04347CD117F9}{}
{\noindent\textcolor{FuncColor}{$\triangleright$\enspace\texttt{OnSets({\mdseries\slshape set, g})\index{OnSets@\texttt{OnSets}}
\label{OnSets}
}\hfill{\scriptsize (function)}}\\


 \index{action!on sets} \index{action!on blocks} Let \mbox{\texttt{\mdseries\slshape set}} be a proper set (see{\nobreakspace}\ref{Sorted Lists and Sets}). \texttt{OnSets} returns the proper set formed by the images of all points $x$ of \mbox{\texttt{\mdseries\slshape set}} via the action function \texttt{OnPoints} (\ref{OnPoints}), applied to $x$ and \mbox{\texttt{\mdseries\slshape g}}. 

 \texttt{OnSets} is for example used to compute the action of a permutation group on blocks. 

 (\texttt{OnTuples} (\ref{OnTuples}) is an action on lists that preserves the ordering of entries.) 
\begin{Verbatim}[commandchars=!@|,fontsize=\small,frame=single,label=Example]
  !gapprompt@gap>| !gapinput@OnSets( [ 1, 3 ], (1,2,3) );|
  [ 1, 2 ]
  !gapprompt@gap>| !gapinput@OnSets( [ (2,3), (1,2) ], (1,2,3) );|
  [ (2,3), (1,3) ]
  !gapprompt@gap>| !gapinput@g:= Group( (1,2,3), (2,3,4) );;|
  !gapprompt@gap>| !gapinput@Orbit( g, [ 1, 2 ], OnSets );|
  [ [ 1, 2 ], [ 2, 3 ], [ 1, 3 ], [ 3, 4 ], [ 1, 4 ], [ 2, 4 ] ]
\end{Verbatim}
 }

 

\subsection{\textcolor{Chapter }{OnTuples}}
\logpage{[ 41, 2, 5 ]}\nobreak
\hyperdef{L}{X832CC5F87EEA4A7E}{}
{\noindent\textcolor{FuncColor}{$\triangleright$\enspace\texttt{OnTuples({\mdseries\slshape tup, g})\index{OnTuples@\texttt{OnTuples}}
\label{OnTuples}
}\hfill{\scriptsize (function)}}\\


 Let \mbox{\texttt{\mdseries\slshape tup}} be a list. \texttt{OnTuples} returns the list formed by the images of all points $x$ of \mbox{\texttt{\mdseries\slshape tup}} via the action function \texttt{OnPoints} (\ref{OnPoints}), applied to $x$ and \mbox{\texttt{\mdseries\slshape g}}. 

 (\texttt{OnSets} (\ref{OnSets}) is an action on lists that additionally sorts the entries of the result.) 
\begin{Verbatim}[commandchars=!@|,fontsize=\small,frame=single,label=Example]
  !gapprompt@gap>| !gapinput@OnTuples( [ 1, 3 ], (1,2,3) );|
  [ 2, 1 ]
  !gapprompt@gap>| !gapinput@OnTuples( [ (2,3), (1,2) ], (1,2,3) );|
  [ (1,3), (2,3) ]
  !gapprompt@gap>| !gapinput@g:= Group( (1,2,3), (2,3,4) );;|
  !gapprompt@gap>| !gapinput@Orbit( g, [ 1, 2 ], OnTuples );|
  [ [ 1, 2 ], [ 2, 3 ], [ 1, 3 ], [ 3, 1 ], [ 3, 4 ], [ 2, 1 ],
    [ 1, 4 ], [ 4, 1 ], [ 4, 2 ], [ 3, 2 ], [ 2, 4 ], [ 4, 3 ] ]
\end{Verbatim}
 }

 

\subsection{\textcolor{Chapter }{OnPairs}}
\logpage{[ 41, 2, 6 ]}\nobreak
\hyperdef{L}{X80DAA1D2855B1456}{}
{\noindent\textcolor{FuncColor}{$\triangleright$\enspace\texttt{OnPairs({\mdseries\slshape tup, g})\index{OnPairs@\texttt{OnPairs}}
\label{OnPairs}
}\hfill{\scriptsize (function)}}\\


 is a special case of \texttt{OnTuples} (\ref{OnTuples}) for lists \mbox{\texttt{\mdseries\slshape tup}} of length 2. }

 

\subsection{\textcolor{Chapter }{OnSetsSets}}
\logpage{[ 41, 2, 7 ]}\nobreak
\hyperdef{L}{X7C10492081D72376}{}
{\noindent\textcolor{FuncColor}{$\triangleright$\enspace\texttt{OnSetsSets({\mdseries\slshape set, g})\index{OnSetsSets@\texttt{OnSetsSets}}
\label{OnSetsSets}
}\hfill{\scriptsize (function)}}\\


 implements the action on sets of sets. \mbox{\texttt{\mdseries\slshape set}} must be a proper set whose entries are again proper sets (see{\nobreakspace}\ref{Sorted Lists and Sets}). 

 \texttt{OnSetsSets} returns the proper set whose elements are the images obtained by applying the
action function \texttt{OnSets} (\ref{OnSets}) to $x$ and \mbox{\texttt{\mdseries\slshape g}}, for each element $x$ of \mbox{\texttt{\mdseries\slshape set}}. 

 For the special case that the elements of \mbox{\texttt{\mdseries\slshape set}} are pairwise disjoint, it is possible to use \texttt{OnSetsDisjointSets} (\ref{OnSetsDisjointSets}). 
\begin{Verbatim}[commandchars=!@|,fontsize=\small,frame=single,label=Example]
  !gapprompt@gap>| !gapinput@OnSetsSets( [ [ 1, 2 ], [ 3, 4 ] ], (1,2,3) );|
  [ [ 1, 4 ], [ 2, 3 ] ]
  !gapprompt@gap>| !gapinput@g:= Group( (1,2,3), (2,3,4) );;|
  !gapprompt@gap>| !gapinput@Orbit( g, [ [ 1, 2 ], [ 3, 4 ] ], OnSetsSets );|
  [ [ [ 1, 2 ], [ 3, 4 ] ], [ [ 1, 4 ], [ 2, 3 ] ],
    [ [ 1, 3 ], [ 2, 4 ] ] ]
\end{Verbatim}
 }

 

\subsection{\textcolor{Chapter }{OnSetsDisjointSets}}
\logpage{[ 41, 2, 8 ]}\nobreak
\hyperdef{L}{X7E23686E7A9D3A20}{}
{\noindent\textcolor{FuncColor}{$\triangleright$\enspace\texttt{OnSetsDisjointSets({\mdseries\slshape set, g})\index{OnSetsDisjointSets@\texttt{OnSetsDisjointSets}}
\label{OnSetsDisjointSets}
}\hfill{\scriptsize (function)}}\\


 implements the action on sets of pairwise disjoint sets (see
also{\nobreakspace}\texttt{OnSetsSets} (\ref{OnSetsSets})). \mbox{\texttt{\mdseries\slshape set}} must be a proper set whose entries are pairwise disjoint proper sets. However,
it is \emph{not} checked that the entries of \mbox{\texttt{\mdseries\slshape set}} are indeed pairwise disjoint. 

 \texttt{OnSetsDisjointSets} returns the proper set whose elements are the images obtained by applying the
action function \texttt{OnSets} (\ref{OnSets}) to $x$ and \mbox{\texttt{\mdseries\slshape g}}, for each element $x$ of \mbox{\texttt{\mdseries\slshape set}}. }

 

\subsection{\textcolor{Chapter }{OnSetsTuples}}
\logpage{[ 41, 2, 9 ]}\nobreak
\hyperdef{L}{X7ADE244E819035FF}{}
{\noindent\textcolor{FuncColor}{$\triangleright$\enspace\texttt{OnSetsTuples({\mdseries\slshape set, g})\index{OnSetsTuples@\texttt{OnSetsTuples}}
\label{OnSetsTuples}
}\hfill{\scriptsize (function)}}\\


 implements the action on sets of tuples. \mbox{\texttt{\mdseries\slshape set}} must be a proper set of lists without holes (see{\nobreakspace}\ref{Sorted Lists and Sets}). 

 \texttt{OnSetsTuples} returns the proper set whose elements are the images obtained by applying the
action function \texttt{OnTuples} (\ref{OnTuples}) to $x$ and \mbox{\texttt{\mdseries\slshape g}}, for each element $x$ of \mbox{\texttt{\mdseries\slshape set}}. 
\begin{Verbatim}[commandchars=!@|,fontsize=\small,frame=single,label=Example]
  !gapprompt@gap>| !gapinput@OnSetsTuples( [ [ 1, 2 ], [ 3, 4 ] ], (1,2,3) );|
  [ [ 1, 4 ], [ 2, 3 ] ]
  !gapprompt@gap>| !gapinput@g:= Group( (1,2,3), (2,3,4) );;|
  !gapprompt@gap>| !gapinput@Orbit( g, [ [ 1, 2 ], [ 3, 4 ] ], OnSetsTuples );|
  [ [ [ 1, 2 ], [ 3, 4 ] ], [ [ 1, 4 ], [ 2, 3 ] ],
    [ [ 1, 3 ], [ 4, 2 ] ], [ [ 2, 4 ], [ 3, 1 ] ],
    [ [ 2, 1 ], [ 4, 3 ] ], [ [ 3, 2 ], [ 4, 1 ] ] ]
\end{Verbatim}
 }

 

\subsection{\textcolor{Chapter }{OnTuplesSets}}
\logpage{[ 41, 2, 10 ]}\nobreak
\hyperdef{L}{X7FF556CD7E6739A9}{}
{\noindent\textcolor{FuncColor}{$\triangleright$\enspace\texttt{OnTuplesSets({\mdseries\slshape tup, g})\index{OnTuplesSets@\texttt{OnTuplesSets}}
\label{OnTuplesSets}
}\hfill{\scriptsize (function)}}\\


 implements the action on tuples of sets. \mbox{\texttt{\mdseries\slshape tup}} must be a list without holes whose entries are proper sets (see{\nobreakspace}\ref{Sorted Lists and Sets}). 

 \texttt{OnTuplesSets} returns the list of the corresponding images obtained by applying the action
function \texttt{OnSets} (\ref{OnSets}) to $x$ and \mbox{\texttt{\mdseries\slshape g}}, for each entry $x$ of \mbox{\texttt{\mdseries\slshape tup}}. 
\begin{Verbatim}[commandchars=!@|,fontsize=\small,frame=single,label=Example]
  !gapprompt@gap>| !gapinput@OnTuplesSets( [ [ 2, 3 ], [ 3, 4 ] ], (1,2,3) );|
  [ [ 1, 3 ], [ 1, 4 ] ]
  !gapprompt@gap>| !gapinput@g:= Group( (1,2,3), (2,3,4) );;|
  !gapprompt@gap>| !gapinput@Orbit( g, [ [ 1, 2 ], [ 3, 4 ] ], OnTuplesSets );|
  [ [ [ 1, 2 ], [ 3, 4 ] ], [ [ 2, 3 ], [ 1, 4 ] ],
    [ [ 1, 3 ], [ 2, 4 ] ], [ [ 3, 4 ], [ 1, 2 ] ],
    [ [ 1, 4 ], [ 2, 3 ] ], [ [ 2, 4 ], [ 1, 3 ] ] ]
\end{Verbatim}
 }

 

\subsection{\textcolor{Chapter }{OnTuplesTuples}}
\logpage{[ 41, 2, 11 ]}\nobreak
\hyperdef{L}{X844E902382EB4151}{}
{\noindent\textcolor{FuncColor}{$\triangleright$\enspace\texttt{OnTuplesTuples({\mdseries\slshape tup, g})\index{OnTuplesTuples@\texttt{OnTuplesTuples}}
\label{OnTuplesTuples}
}\hfill{\scriptsize (function)}}\\


 implements the action on tuples of tuples. \mbox{\texttt{\mdseries\slshape tup}} must be a list without holes whose entries are again lists without holes. 

 \texttt{OnTuplesTuples} returns the list of the corresponding images obtained by applying the action
function \texttt{OnTuples} (\ref{OnTuples}) to $x$ and \mbox{\texttt{\mdseries\slshape g}}, for each entry $x$ of the list \mbox{\texttt{\mdseries\slshape tup}}. 
\begin{Verbatim}[commandchars=!@|,fontsize=\small,frame=single,label=Example]
  !gapprompt@gap>| !gapinput@OnTuplesTuples( [ [ 2, 3 ], [ 3, 4 ] ], (1,2,3) );|
  [ [ 3, 1 ], [ 1, 4 ] ]
  !gapprompt@gap>| !gapinput@g:=Group((1,2,3),(2,3,4));;|
  !gapprompt@gap>| !gapinput@Orbit(g,[[1,2],[3,4]],OnTuplesTuples);|
  [ [ [ 1, 2 ], [ 3, 4 ] ], [ [ 2, 3 ], [ 1, 4 ] ],
    [ [ 1, 3 ], [ 4, 2 ] ], [ [ 3, 1 ], [ 2, 4 ] ],
    [ [ 3, 4 ], [ 1, 2 ] ], [ [ 2, 1 ], [ 4, 3 ] ],
    [ [ 1, 4 ], [ 2, 3 ] ], [ [ 4, 1 ], [ 3, 2 ] ],
    [ [ 4, 2 ], [ 1, 3 ] ], [ [ 3, 2 ], [ 4, 1 ] ],
    [ [ 2, 4 ], [ 3, 1 ] ], [ [ 4, 3 ], [ 2, 1 ] ] ]
\end{Verbatim}
 }

 

\subsection{\textcolor{Chapter }{OnLines}}
\logpage{[ 41, 2, 12 ]}\nobreak
\hyperdef{L}{X86DC2DD5829CAD9A}{}
{\noindent\textcolor{FuncColor}{$\triangleright$\enspace\texttt{OnLines({\mdseries\slshape vec, g})\index{OnLines@\texttt{OnLines}}
\label{OnLines}
}\hfill{\scriptsize (function)}}\\


 Let \mbox{\texttt{\mdseries\slshape vec}} be a \emph{normed} row vector, that is, its first nonzero entry is normed to the identity of the
relevant field, see \texttt{NormedRowVector} (\ref{NormedRowVector}). The function \texttt{OnLines} returns the row vector obtained from first multiplying \mbox{\texttt{\mdseries\slshape vec}} from the right with \mbox{\texttt{\mdseries\slshape g}} (via \texttt{OnRight} (\ref{OnRight})) and then normalizing the resulting row vector by scalar multiplication from
the left. 

 This action corresponds to the projective action of a matrix group on
one\texttt{\symbol{45}}dimensional subspaces. 

 If \mbox{\texttt{\mdseries\slshape vec}} is a zero vector or is not normed then an error is triggered
(see{\nobreakspace}\ref{Action on canonical representatives}). 

 
\begin{Verbatim}[commandchars=!@|,fontsize=\small,frame=single,label=Example]
  !gapprompt@gap>| !gapinput@OnLines( [ 1, 2 ], [ [ 1, 2 ], [ 3, 4 ] ] );|
  [ 1, 10/7 ]
  !gapprompt@gap>| !gapinput@gl:=GL(2,5);;v:=[1,0]*Z(5)^0;|
  [ Z(5)^0, 0*Z(5) ]
  !gapprompt@gap>| !gapinput@h:=Action(gl,Orbit(gl,v,OnLines),OnLines);|
  Group([ (2,3,5,6), (1,2,4)(3,6,5) ])
\end{Verbatim}
 }

 

\subsection{\textcolor{Chapter }{OnIndeterminates (as a permutation action)}}
\logpage{[ 41, 2, 13 ]}\nobreak
\hyperdef{L}{X7FA394D27E721E2B}{}
{\noindent\textcolor{FuncColor}{$\triangleright$\enspace\texttt{OnIndeterminates({\mdseries\slshape poly, perm})\index{OnIndeterminates@\texttt{OnIndeterminates}!as a permutation action}
\label{OnIndeterminates:as a permutation action}
}\hfill{\scriptsize (function)}}\\


 A permutation \mbox{\texttt{\mdseries\slshape perm}} acts on the multivariate polynomial \mbox{\texttt{\mdseries\slshape poly}} by permuting the indeterminates as it permutes points. 
\begin{Verbatim}[commandchars=!@|,fontsize=\small,frame=single,label=Example]
  !gapprompt@gap>| !gapinput@x:=Indeterminate(Rationals,1);; y:=Indeterminate(Rationals,2);;|
  !gapprompt@gap>| !gapinput@OnIndeterminates(x^7*y+x*y^4,(1,17)(2,28));|
  x_17^7*x_28+x_17*x_28^4
  !gapprompt@gap>| !gapinput@Stabilizer(Group((1,2,3,4),(1,2)),x*y,OnIndeterminates);|
  Group([ (1,2), (3,4) ])
\end{Verbatim}
 }

 

\subsection{\textcolor{Chapter }{Permuted (as a permutation action)}}
\logpage{[ 41, 2, 14 ]}\nobreak
\hyperdef{L}{X7BA8D76586F1F06E}{}
{\noindent\textcolor{FuncColor}{$\triangleright$\enspace\texttt{Permuted({\mdseries\slshape list, perm})\index{Permuted@\texttt{Permuted}!as a permutation action}
\label{Permuted:as a permutation action}
}\hfill{\scriptsize (method)}}\\


 The following example demonstrates \texttt{Permuted} (\ref{Permuted}) being used to implement a permutation action on a domain: 

 
\begin{Verbatim}[commandchars=!@|,fontsize=\small,frame=single,label=Example]
  !gapprompt@gap>| !gapinput@g:=Group((1,2,3),(1,2));;|
  !gapprompt@gap>| !gapinput@dom:=[ "a", "b", "c" ];;|
  !gapprompt@gap>| !gapinput@Orbit(g,dom,Permuted);|
  [ [ "a", "b", "c" ], [ "c", "a", "b" ], [ "b", "a", "c" ],
    [ "b", "c", "a" ], [ "a", "c", "b" ], [ "c", "b", "a" ] ]
\end{Verbatim}
 }

 

\subsection{\textcolor{Chapter }{OnSubspacesByCanonicalBasis}}
\logpage{[ 41, 2, 15 ]}\nobreak
\hyperdef{L}{X85124D197F0F9C4D}{}
{\noindent\textcolor{FuncColor}{$\triangleright$\enspace\texttt{OnSubspacesByCanonicalBasis({\mdseries\slshape bas, mat})\index{OnSubspacesByCanonicalBasis@\texttt{OnSubspacesByCanonicalBasis}}
\label{OnSubspacesByCanonicalBasis}
}\hfill{\scriptsize (function)}}\\
\noindent\textcolor{FuncColor}{$\triangleright$\enspace\texttt{OnSubspacesByCanonicalBasisConcatenations({\mdseries\slshape basvec, mat})\index{OnSubspacesByCanonicalBasisConcatenations@\texttt{OnSubspaces}\-\texttt{By}\-\texttt{Canonical}\-\texttt{Basis}\-\texttt{Concatenations}}
\label{OnSubspacesByCanonicalBasisConcatenations}
}\hfill{\scriptsize (function)}}\\


 implements the operation of a matrix group on subspaces of a vector space. \mbox{\texttt{\mdseries\slshape bas}} must be a list of (linearly independent) vectors which forms a basis of the
subspace in Hermite normal form. \mbox{\texttt{\mdseries\slshape mat}} is an element of the acting matrix group. The function returns a mutable
matrix which gives the basis of the image of the subspace in Hermite normal
form. (In other words: it triangulizes the product of \mbox{\texttt{\mdseries\slshape bas}} with \mbox{\texttt{\mdseries\slshape mat}}.) 

 \mbox{\texttt{\mdseries\slshape bas}} must be given in Hermite normal form, otherwise an error is triggered
(see{\nobreakspace}\ref{Action on canonical representatives}). }

 }

  
\section{\textcolor{Chapter }{Action on canonical representatives}}\label{Action on canonical representatives}
\logpage{[ 41, 3, 0 ]}
\hyperdef{L}{X82181CA07A5B2056}{}
{
  A variety of action functions assumes that the objects on which it acts are
given in a particular form, for example canonical representatives. Affected
actions are for example \texttt{OnSetsSets} (\ref{OnSetsSets}), \texttt{OnSetsDisjointSets} (\ref{OnSetsDisjointSets}), \texttt{OnSetsTuples} (\ref{OnSetsTuples}), \texttt{OnTuplesSets} (\ref{OnTuplesSets}), \texttt{OnLines} (\ref{OnLines}) and \texttt{OnSubspacesByCanonicalBasis} (\ref{OnSubspacesByCanonicalBasis}). 

 If orbit seeds or domain elements are not given in the required form \textsf{GAP} will issue an error message: 
\begin{Verbatim}[commandchars=!@|,fontsize=\small,frame=single,label=Example]
  !gapprompt@gap>| !gapinput@Orbit(SymmetricGroup(5),[[2,4],[1,3]],OnSetsSets);|
  Error, Action not well-defined. See the manual section
  ``Action on canonical representatives''.
\end{Verbatim}
 In this case the affected domain elements have to be brought in canonical
form, as documented for the respective action function. For interactive use
this is most easily done by acting with the identity element of the group. 

 (A similar error could arise if a user\texttt{\symbol{45}}defined action
function is used which actually does not implement an action from the right.) }

  
\section{\textcolor{Chapter }{Orbits}}\label{Orbits}
\logpage{[ 41, 4, 0 ]}
\hyperdef{L}{X81E0FF0587C54543}{}
{
  If a group $G$ acts on a set $\Omega$, the set of all images of $x \in \Omega$ under elements of $G$ is called the \emph{orbit} of $x$. The set of orbits of $G$ is a partition of $\Omega$. 

\subsection{\textcolor{Chapter }{Orbit}}
\logpage{[ 41, 4, 1 ]}\nobreak
\hyperdef{L}{X80E0234E7BD79409}{}
{\noindent\textcolor{FuncColor}{$\triangleright$\enspace\texttt{Orbit({\mdseries\slshape G[, Omega], pnt[, gens, acts][, act]})\index{Orbit@\texttt{Orbit}}
\label{Orbit}
}\hfill{\scriptsize (operation)}}\\


 The orbit of the point \mbox{\texttt{\mdseries\slshape pnt}} is the list of all images of \mbox{\texttt{\mdseries\slshape pnt}} under the action of the group \mbox{\texttt{\mdseries\slshape G}} w.r.t. the action function \mbox{\texttt{\mdseries\slshape act}} or \texttt{OnPoints} (\ref{OnPoints}) if no action function is given. 

 (Note that the arrangement of points in this list is not defined by the
operation.) 

 The orbit of \mbox{\texttt{\mdseries\slshape pnt}} will always contain one element that is \emph{equal} to \mbox{\texttt{\mdseries\slshape pnt}}, however for performance reasons this element is not necessarily \emph{identical} to \mbox{\texttt{\mdseries\slshape pnt}}, in particular if \mbox{\texttt{\mdseries\slshape pnt}} is mutable. 
\begin{Verbatim}[commandchars=!@|,fontsize=\small,frame=single,label=Example]
  !gapprompt@gap>| !gapinput@g:=Group((1,3,2),(2,4,3));;|
  !gapprompt@gap>| !gapinput@Orbit(g,1);|
  [ 1, 3, 2, 4 ]
  !gapprompt@gap>| !gapinput@Orbit(g,[1,2],OnSets);|
  [ [ 1, 2 ], [ 1, 3 ], [ 1, 4 ], [ 2, 3 ], [ 3, 4 ], [ 2, 4 ] ]
\end{Verbatim}
 

 (See Section{\nobreakspace}\ref{Basic Actions} for information about specific actions.) }

 

\subsection{\textcolor{Chapter }{Orbits (operation)}}
\logpage{[ 41, 4, 2 ]}\nobreak
\hyperdef{L}{X86BCAE17869BBEAA}{}
{\noindent\textcolor{FuncColor}{$\triangleright$\enspace\texttt{Orbits({\mdseries\slshape G, seeds[, gens, acts][, act]})\index{Orbits@\texttt{Orbits}!operation}
\label{Orbits:operation}
}\hfill{\scriptsize (operation)}}\\
\noindent\textcolor{FuncColor}{$\triangleright$\enspace\texttt{Orbits({\mdseries\slshape G})\index{Orbits@\texttt{Orbits}!for a permutation group}
\label{Orbits:for a permutation group}
}\hfill{\scriptsize (attribute)}}\\
\noindent\textcolor{FuncColor}{$\triangleright$\enspace\texttt{Orbits({\mdseries\slshape xset})\index{Orbits@\texttt{Orbits}!attribute}
\label{Orbits:attribute}
}\hfill{\scriptsize (attribute)}}\\


 returns a duplicate\texttt{\symbol{45}}free list of the orbits of the elements
in \mbox{\texttt{\mdseries\slshape seeds}} under the action \mbox{\texttt{\mdseries\slshape act}} of \mbox{\texttt{\mdseries\slshape G}} or under \texttt{OnPoints} (\ref{OnPoints}) if no action function is given. 

 For a permutation group \mbox{\texttt{\mdseries\slshape G}}, one may also invoke this as \texttt{Orbits(\mbox{\texttt{\mdseries\slshape G}})}, which returns all the orbits of its natural action on the set of points
moved by it. For example the group $\langle (1,2,3), (4,5) \rangle$ has the orbits $\{1,2,3\}$ and $\{4,5\}$. 

 (Note that the arrangement of orbits or of points within one orbit is not
defined by the operation.) }

 
\subsection{\textcolor{Chapter }{OrbitsDomain}}\logpage{[ 41, 4, 3 ]}
\hyperdef{L}{X86BC8B958123F953}{}
{
\noindent\textcolor{FuncColor}{$\triangleright$\enspace\texttt{OrbitsDomain({\mdseries\slshape G, Omega[, gens, acts][, act]})\index{OrbitsDomain@\texttt{OrbitsDomain}!for a group and an action domain}
\label{OrbitsDomain:for a group and an action domain}
}\hfill{\scriptsize (operation)}}\\
\noindent\textcolor{FuncColor}{$\triangleright$\enspace\texttt{OrbitsDomain({\mdseries\slshape G})\index{OrbitsDomain@\texttt{OrbitsDomain}!for a permutation group}
\label{OrbitsDomain:for a permutation group}
}\hfill{\scriptsize (attribute)}}\\
\noindent\textcolor{FuncColor}{$\triangleright$\enspace\texttt{OrbitsDomain({\mdseries\slshape xset})\index{OrbitsDomain@\texttt{OrbitsDomain}!of an external set}
\label{OrbitsDomain:of an external set}
}\hfill{\scriptsize (attribute)}}\\


 returns a list of the orbits of \mbox{\texttt{\mdseries\slshape G}} on the domain \mbox{\texttt{\mdseries\slshape Omega}} (given as lists) under the action \mbox{\texttt{\mdseries\slshape act}} or under \texttt{OnPoints} (\ref{OnPoints}) if no action function is given. 

 This operation is often faster than \texttt{Orbits} (\ref{Orbits:operation}). The domain \mbox{\texttt{\mdseries\slshape Omega}} must be closed under the action of \mbox{\texttt{\mdseries\slshape G}}, otherwise an error can occur. 

 For a permutation group \mbox{\texttt{\mdseries\slshape G}}, one may also invoke this as \texttt{OrbitsDomain(\mbox{\texttt{\mdseries\slshape G}})}, which returns all the orbits of its natural action on the set of points
moved by it. 

 (Note that the arrangement of orbits or of points within one orbit is not
defined by the operation.) 
\begin{Verbatim}[commandchars=!@|,fontsize=\small,frame=single,label=Example]
  !gapprompt@gap>| !gapinput@g:=Group((1,3,2),(2,4,3));;|
  !gapprompt@gap>| !gapinput@Orbits(g,[1..5]);|
  [ [ 1, 3, 2, 4 ], [ 5 ] ]
  !gapprompt@gap>| !gapinput@OrbitsDomain(g,Arrangements([1..4],3),OnTuples);|
  [ [ [ 1, 2, 3 ], [ 3, 1, 2 ], [ 1, 4, 2 ], [ 2, 3, 1 ], [ 2, 1, 4 ],
        [ 3, 4, 1 ], [ 1, 3, 4 ], [ 4, 2, 1 ], [ 4, 1, 3 ],
        [ 2, 4, 3 ], [ 3, 2, 4 ], [ 4, 3, 2 ] ],
    [ [ 1, 2, 4 ], [ 3, 1, 4 ], [ 1, 4, 3 ], [ 2, 3, 4 ], [ 2, 1, 3 ],
        [ 3, 4, 2 ], [ 1, 3, 2 ], [ 4, 2, 3 ], [ 4, 1, 2 ],
        [ 2, 4, 1 ], [ 3, 2, 1 ], [ 4, 3, 1 ] ] ]
  !gapprompt@gap>| !gapinput@OrbitsDomain(g,GF(2)^2,[(1,2,3),(1,4)(2,3)],|
  !gapprompt@>| !gapinput@[[[Z(2)^0,Z(2)^0],[Z(2)^0,0*Z(2)]],[[Z(2)^0,0*Z(2)],[0*Z(2),Z(2)^0]]]);|
  [ [ <an immutable GF2 vector of length 2> ],
    [ <an immutable GF2 vector of length 2>,
        <an immutable GF2 vector of length 2>,
        <an immutable GF2 vector of length 2> ] ]
\end{Verbatim}
 

 (See Section{\nobreakspace}\ref{Basic Actions} for information about specific actions.) }

 

\subsection{\textcolor{Chapter }{OrbitLength}}
\logpage{[ 41, 4, 4 ]}\nobreak
\hyperdef{L}{X799910CF832EDC45}{}
{\noindent\textcolor{FuncColor}{$\triangleright$\enspace\texttt{OrbitLength({\mdseries\slshape G[, Omega], pnt[, gens, acts][, act]})\index{OrbitLength@\texttt{OrbitLength}}
\label{OrbitLength}
}\hfill{\scriptsize (operation)}}\\


 computes the length of the orbit of \mbox{\texttt{\mdseries\slshape pnt}} under the action function \mbox{\texttt{\mdseries\slshape act}} or \texttt{OnPoints} (\ref{OnPoints}) if no action function is given. }

 
\subsection{\textcolor{Chapter }{OrbitLengths}}\logpage{[ 41, 4, 5 ]}
\hyperdef{L}{X8032F73078DF2DDB}{}
{
\noindent\textcolor{FuncColor}{$\triangleright$\enspace\texttt{OrbitLengths({\mdseries\slshape G, seeds[, gens, acts][, act]})\index{OrbitLengths@\texttt{OrbitLengths}!for a group, a set of seeds, etc.}
\label{OrbitLengths:for a group, a set of seeds, etc.}
}\hfill{\scriptsize (operation)}}\\
\noindent\textcolor{FuncColor}{$\triangleright$\enspace\texttt{OrbitLengths({\mdseries\slshape G})\index{OrbitLengths@\texttt{OrbitLengths}!for a permutation group}
\label{OrbitLengths:for a permutation group}
}\hfill{\scriptsize (attribute)}}\\
\noindent\textcolor{FuncColor}{$\triangleright$\enspace\texttt{OrbitLengths({\mdseries\slshape xset})\index{OrbitLengths@\texttt{OrbitLengths}!for an external set}
\label{OrbitLengths:for an external set}
}\hfill{\scriptsize (attribute)}}\\


 computes the lengths of all the orbits of the elements in \mbox{\texttt{\mdseries\slshape seeds}} under the action \mbox{\texttt{\mdseries\slshape act}} of \mbox{\texttt{\mdseries\slshape G}}. 

 For a permutation group \mbox{\texttt{\mdseries\slshape G}}, one may also invoke this as \texttt{OrbitLengths(\mbox{\texttt{\mdseries\slshape G}})}, which returns the lengths of all the orbits of its natural action on the set
of points moved by it. For example the group $\langle (1,2,3), (5,6) \rangle$ has the orbit lengths 2 and 3. }

 
\subsection{\textcolor{Chapter }{OrbitLengthsDomain}}\logpage{[ 41, 4, 6 ]}
\hyperdef{L}{X8520E2487F7E98AF}{}
{
\noindent\textcolor{FuncColor}{$\triangleright$\enspace\texttt{OrbitLengthsDomain({\mdseries\slshape G, Omega[, gens, acts][, act]})\index{OrbitLengthsDomain@\texttt{OrbitLengthsDomain}!for a group and a set of seeds}
\label{OrbitLengthsDomain:for a group and a set of seeds}
}\hfill{\scriptsize (operation)}}\\
\noindent\textcolor{FuncColor}{$\triangleright$\enspace\texttt{OrbitLengthsDomain({\mdseries\slshape G})\index{OrbitLengthsDomain@\texttt{OrbitLengthsDomain}!for a permutation group}
\label{OrbitLengthsDomain:for a permutation group}
}\hfill{\scriptsize (attribute)}}\\
\noindent\textcolor{FuncColor}{$\triangleright$\enspace\texttt{OrbitLengthsDomain({\mdseries\slshape xset})\index{OrbitLengthsDomain@\texttt{OrbitLengthsDomain}!of an external set}
\label{OrbitLengthsDomain:of an external set}
}\hfill{\scriptsize (attribute)}}\\


 computes the lengths of all the orbits of \mbox{\texttt{\mdseries\slshape G}} on \mbox{\texttt{\mdseries\slshape Omega}}. 

 This operation is often faster than \texttt{OrbitLengths} (\ref{OrbitLengths:for a group, a set of seeds, etc.}). The domain \mbox{\texttt{\mdseries\slshape Omega}} must be closed under the action of \mbox{\texttt{\mdseries\slshape G}}, otherwise an error can occur. 

 For a permutation group \mbox{\texttt{\mdseries\slshape G}}, one may also invoke this as \texttt{OrbitLengthsDomain(\mbox{\texttt{\mdseries\slshape G}})}, which returns the length of all the orbits of its natural action on the set
of points moved by it. 
\begin{Verbatim}[commandchars=!@|,fontsize=\small,frame=single,label=Example]
  !gapprompt@gap>| !gapinput@g:=Group((1,3,2),(2,4,3));;|
  !gapprompt@gap>| !gapinput@OrbitLength(g,[1,2,3,4],OnTuples);|
  12
  !gapprompt@gap>| !gapinput@OrbitLengths(g,Arrangements([1..4],4),OnTuples);|
  [ 12, 12 ]
  !gapprompt@gap>| !gapinput@g:=Group((1,2,3),(5,6,7));;|
  !gapprompt@gap>| !gapinput@OrbitLengthsDomain(g,[1,2,3]);|
  [ 3 ]
  !gapprompt@gap>| !gapinput@OrbitLengthsDomain(g);|
  [ 3, 3 ]
\end{Verbatim}
 }

 }

  
\section{\textcolor{Chapter }{Stabilizers}}\label{Stabilizers}
\logpage{[ 41, 5, 0 ]}
\hyperdef{L}{X797BD60E7ACEF1B1}{}
{
  \index{point stabilizer}\index{set stabilizer} \index{tuple stabilizer} The \emph{stabilizer} of a point $x$ under the action of a group $G$ is the set of all those elements in $G$ which fix $x$. 

\subsection{\textcolor{Chapter }{OrbitStabilizer}}
\logpage{[ 41, 5, 1 ]}\nobreak
\hyperdef{L}{X7C34EC437EF598BF}{}
{\noindent\textcolor{FuncColor}{$\triangleright$\enspace\texttt{OrbitStabilizer({\mdseries\slshape G[, Omega], pnt[, gens, acts][, act]})\index{OrbitStabilizer@\texttt{OrbitStabilizer}}
\label{OrbitStabilizer}
}\hfill{\scriptsize (operation)}}\\


 computes the orbit and the stabilizer of \mbox{\texttt{\mdseries\slshape pnt}} simultaneously in a single orbit\texttt{\symbol{45}}stabilizer algorithm.
Returns an immutable record with components \texttt{orbit} (the orbit as a list) and \texttt{stabilizer} (the point stabilizer, a subgroup of \mbox{\texttt{\mdseries\slshape G}}). 

 The stabilizer will have \mbox{\texttt{\mdseries\slshape G}} as its parent. 
\begin{Verbatim}[commandchars=!@|,fontsize=\small,frame=single,label=Example]
  !gapprompt@gap>| !gapinput@g:=Group((1,3,2),(2,4,3));;|
  !gapprompt@gap>| !gapinput@orbstab:=OrbitStabilizer(g,[1,2],OnSets);|
  rec(
    orbit := [ [ 1, 2 ], [ 1, 3 ], [ 1, 4 ], [ 2, 3 ], [ 3, 4 ],
        [ 2, 4 ] ], stabilizer := Group([ (1,2)(3,4) ]) )
  !gapprompt@gap>| !gapinput@orbstab.orbit;|
  [ [ 1, 2 ], [ 1, 3 ], [ 1, 4 ], [ 2, 3 ], [ 3, 4 ], [ 2, 4 ] ]
  !gapprompt@gap>| !gapinput@orbstab.stabilizer;|
  Group([ (1,2)(3,4) ])
\end{Verbatim}
 }

 

\subsection{\textcolor{Chapter }{Stabilizer}}
\logpage{[ 41, 5, 2 ]}\nobreak
\hyperdef{L}{X86FB962786397E02}{}
{\noindent\textcolor{FuncColor}{$\triangleright$\enspace\texttt{Stabilizer({\mdseries\slshape G[, Omega], pnt[, gens, acts][, act]})\index{Stabilizer@\texttt{Stabilizer}}
\label{Stabilizer}
}\hfill{\scriptsize (function)}}\\


 computes the stabilizer in \mbox{\texttt{\mdseries\slshape G}} of the point \mbox{\texttt{\mdseries\slshape pnt}}, that is the subgroup of those elements of \mbox{\texttt{\mdseries\slshape G}} that fix \mbox{\texttt{\mdseries\slshape pnt}}. The stabilizer will have \mbox{\texttt{\mdseries\slshape G}} as its parent. 
\begin{Verbatim}[commandchars=!@|,fontsize=\small,frame=single,label=Example]
  !gapprompt@gap>| !gapinput@g:=Group((1,3,2),(2,4,3));;|
  !gapprompt@gap>| !gapinput@stab:=Stabilizer(g,4);|
  Group([ (1,3,2) ])
  !gapprompt@gap>| !gapinput@Parent(stab);|
  Group([ (1,3,2), (2,4,3) ])
\end{Verbatim}
 

 The stabilizer of a set or tuple of points can be computed by specifying an
action of sets or tuples of points. 
\begin{Verbatim}[commandchars=!@|,fontsize=\small,frame=single,label=Example]
  !gapprompt@gap>| !gapinput@Stabilizer(g,[1,2],OnSets);|
  Group([ (1,2)(3,4) ])
  !gapprompt@gap>| !gapinput@Stabilizer(g,[1,2],OnTuples);|
  Group(())
  !gapprompt@gap>| !gapinput@orbstab:=OrbitStabilizer(g,[1,2],OnSets);|
  rec(
    orbit := [ [ 1, 2 ], [ 1, 3 ], [ 1, 4 ], [ 2, 3 ], [ 3, 4 ],
        [ 2, 4 ] ], stabilizer := Group([ (1,2)(3,4) ]) )
  !gapprompt@gap>| !gapinput@Parent(orbstab.stabilizer);|
  Group([ (1,3,2), (2,4,3) ])
\end{Verbatim}
 

 (See Section{\nobreakspace}\ref{Basic Actions} for information about specific actions.) 

 The standard methods for all these actions are an
orbit\texttt{\symbol{45}}stabilizer algorithm. For permutation groups
backtrack algorithms are used. For solvable groups an
orbit\texttt{\symbol{45}}stabilizer algorithm for solvable groups, which uses
the fact that the orbits of a normal subgroup form a block system (see \cite{SOGOS}) is used. }

 

\subsection{\textcolor{Chapter }{OrbitStabilizerAlgorithm}}
\logpage{[ 41, 5, 3 ]}\nobreak
\hyperdef{L}{X78C3A8568414BC44}{}
{\noindent\textcolor{FuncColor}{$\triangleright$\enspace\texttt{OrbitStabilizerAlgorithm({\mdseries\slshape G, Omega, blist, gens, acts, pntact})\index{OrbitStabilizerAlgorithm@\texttt{OrbitStabilizerAlgorithm}}
\label{OrbitStabilizerAlgorithm}
}\hfill{\scriptsize (operation)}}\\


 This operation should not be called by a user. It is documented however for
purposes to extend or maintain the group actions package (the word ``package'' here refers to the \textsf{GAP} functionality for group actions, not to a \textsf{GAP} package). 

 \texttt{OrbitStabilizerAlgorithm} performs an orbit stabilizer algorithm for the group \mbox{\texttt{\mdseries\slshape G}} acting with the generators \mbox{\texttt{\mdseries\slshape gens}} via the generator images \mbox{\texttt{\mdseries\slshape gens}} and the group action \mbox{\texttt{\mdseries\slshape act}} on the element \mbox{\texttt{\mdseries\slshape pnt}}. (For technical reasons \mbox{\texttt{\mdseries\slshape pnt}} and \mbox{\texttt{\mdseries\slshape act}} are put in one record with components \texttt{pnt} and \texttt{act} respectively.) 

 The \mbox{\texttt{\mdseries\slshape pntact}} record may carry a component \mbox{\texttt{\mdseries\slshape stabsub}}. If given, this must be a subgroup stabilizing \emph{all} points in the domain and can be used to abbreviate stabilizer calculations. 

 The \mbox{\texttt{\mdseries\slshape pntact}} component also may contain the boolean entry \texttt{onlystab} set to \texttt{true}. In this case the \texttt{orbit} component may be omitted from the result. 

 The argument \mbox{\texttt{\mdseries\slshape Omega}} (which may be replaced by \texttt{false} to be ignored) is the set within which the orbit is computed (once the orbit
is the full domain, the orbit calculation may stop). If \mbox{\texttt{\mdseries\slshape blist}} is given it must be a bit list corresponding to \mbox{\texttt{\mdseries\slshape Omega}} in which elements which have been found already will be ``ticked off'' with \texttt{true}. (In particular, the entries for the orbit of \mbox{\texttt{\mdseries\slshape pnt}} still must be all set to \texttt{false}). Again the remaining action domain (the bits set initially to \texttt{false}) can be used to stop if the orbit cannot grow any longer. Another use of the
bit list is if \mbox{\texttt{\mdseries\slshape Omega}} is an enumerator which can determine \texttt{PositionCanonical} (\ref{PositionCanonical}) values very quickly. In this situation it can be worth to search images not in
the orbit found so far, but via their position in \mbox{\texttt{\mdseries\slshape Omega}} and use a the bit list to keep track whether the element is in the orbit found
so far. }

 }

  
\section{\textcolor{Chapter }{Elements with Prescribed Images}}\label{Elements with Prescribed Images}
\logpage{[ 41, 6, 0 ]}
\hyperdef{L}{X7A9389097BAF670D}{}
{
  \index{transporter} 

\subsection{\textcolor{Chapter }{RepresentativeAction}}
\logpage{[ 41, 6, 1 ]}\nobreak
\hyperdef{L}{X857DC7B085EB0539}{}
{\noindent\textcolor{FuncColor}{$\triangleright$\enspace\texttt{RepresentativeAction({\mdseries\slshape G[, Omega], d, e[, gens, acts][, act]})\index{RepresentativeAction@\texttt{RepresentativeAction}}
\label{RepresentativeAction}
}\hfill{\scriptsize (function)}}\\


 computes an element of \mbox{\texttt{\mdseries\slshape G}} that maps \mbox{\texttt{\mdseries\slshape d}} to \mbox{\texttt{\mdseries\slshape e}} under the given action and returns \texttt{fail} if no such element exists. 
\begin{Verbatim}[commandchars=!@|,fontsize=\small,frame=single,label=Example]
  !gapprompt@gap>| !gapinput@g:=Group((1,3,2),(2,4,3));;|
  !gapprompt@gap>| !gapinput@RepresentativeAction(g,1,3);|
  (1,3)(2,4)
  !gapprompt@gap>| !gapinput@RepresentativeAction(g,1,3,OnPoints);|
  (1,3)(2,4)
  !gapprompt@gap>| !gapinput@RepresentativeAction(g,(1,2,3),(2,4,3));|
  (1,2,4)
  !gapprompt@gap>| !gapinput@RepresentativeAction(g,(1,2,3),(2,3,4));|
  fail
  !gapprompt@gap>| !gapinput@RepresentativeAction(g,Group((1,2,3)),Group((2,3,4)));|
  (1,2,4)
  !gapprompt@gap>| !gapinput@ RepresentativeAction(g,[1,2,3],[1,2,4],OnSets);|
  (2,4,3)
  !gapprompt@gap>| !gapinput@ RepresentativeAction(g,[1,2,3],[1,2,4],OnTuples);|
  fail
\end{Verbatim}
 

 (See Section{\nobreakspace}\ref{Basic Actions} for information about specific actions.) 

 Again the standard method for \texttt{RepresentativeAction} is an orbit\texttt{\symbol{45}}stabilizer algorithm, for permutation groups
and standard actions a backtrack algorithm is used. }

 }

  
\section{\textcolor{Chapter }{The Permutation Image of an Action}}\label{The Permutation Image of an Action}
\logpage{[ 41, 7, 0 ]}
\hyperdef{L}{X87F73CCA7921DE65}{}
{
  When a group $G$ acts on a domain $\Omega$, an enumeration of $Omega$ yields a homomorphism from $G$ into the symmetric group on $\{ 1, \ldots, |\Omega| \}$. In \textsf{GAP}, the enumeration of $\Omega$ is provided by the \texttt{Enumerator} (\ref{Enumerator}) value of $\Omega$ which of course is $\Omega$ itself if it is a list. 

 For an action homomorphism, the operation \texttt{UnderlyingExternalSet} (\ref{UnderlyingExternalSet}) will return the external set on $\Omega$ which affords the action. 
\subsection{\textcolor{Chapter }{ActionHomomorphism}}\logpage{[ 41, 7, 1 ]}
\hyperdef{L}{X78E6A002835288A4}{}
{
\noindent\textcolor{FuncColor}{$\triangleright$\enspace\texttt{ActionHomomorphism({\mdseries\slshape G, Omega[, gens, acts][, act][, "surjective"]})\index{ActionHomomorphism@\texttt{ActionHomomorphism}!for a group, an action domain, etc.}
\label{ActionHomomorphism:for a group, an action domain, etc.}
}\hfill{\scriptsize (function)}}\\
\noindent\textcolor{FuncColor}{$\triangleright$\enspace\texttt{ActionHomomorphism({\mdseries\slshape xset[, "surjective"]})\index{ActionHomomorphism@\texttt{ActionHomomorphism}!for an external set}
\label{ActionHomomorphism:for an external set}
}\hfill{\scriptsize (function)}}\\
\noindent\textcolor{FuncColor}{$\triangleright$\enspace\texttt{ActionHomomorphism({\mdseries\slshape action})\index{ActionHomomorphism@\texttt{ActionHomomorphism}!for an action image}
\label{ActionHomomorphism:for an action image}
}\hfill{\scriptsize (function)}}\\


 computes a homomorphism from \mbox{\texttt{\mdseries\slshape G}} into the symmetric group on $|\mbox{\texttt{\mdseries\slshape Omega}}|$ points that gives the permutation action of \mbox{\texttt{\mdseries\slshape G}} on \mbox{\texttt{\mdseries\slshape Omega}}. (In particular, this homomorphism is a permutation equivalence, that is the
permutation image of a group element is given by the positions of points in \mbox{\texttt{\mdseries\slshape Omega}}.) 

 The result is undefined if \mbox{\texttt{\mdseries\slshape G}} does not act on \mbox{\texttt{\mdseries\slshape Omega}}. 

 By default the homomorphism returned by \texttt{ActionHomomorphism} is not necessarily surjective (its \texttt{Range} (\ref{Range:of a general mapping}) value is the full symmetric group) to avoid unnecessary computation of the
image. If the optional string argument \texttt{"surjective"} is given, a surjective homomorphism is created. 

 The third version (which is supported only for \textsf{GAP}3 compatibility) returns the action homomorphism that belongs to the image
obtained via \texttt{Action} (\ref{Action:for a group, an action domain, etc.}). 

 (See Section{\nobreakspace}\ref{Basic Actions} for information about specific actions.) 

 
\begin{Verbatim}[commandchars=!@|,fontsize=\small,frame=single,label=Example]
  !gapprompt@gap>| !gapinput@g:=Group((1,2,3),(1,2));;|
  !gapprompt@gap>| !gapinput@hom:=ActionHomomorphism(g,Arrangements([1..4],3),OnTuples);|
  <action homomorphism>
  !gapprompt@gap>| !gapinput@Image(hom);|
  Group(
  [ (1,9,13)(2,10,14)(3,7,15)(4,8,16)(5,12,17)(6,11,18)(19,22,23)(20,21,
      24), (1,7)(2,8)(3,9)(4,10)(5,11)(6,12)(13,15)(14,16)(17,18)(19,
      21)(20,22)(23,24) ])
  !gapprompt@gap>| !gapinput@Size(Range(hom));Size(Image(hom));|
  620448401733239439360000
  6
  !gapprompt@gap>| !gapinput@hom:=ActionHomomorphism(g,Arrangements([1..4],3),OnTuples,|
  !gapprompt@>| !gapinput@"surjective");;|
  !gapprompt@gap>| !gapinput@Size(Range(hom));|
  6
\end{Verbatim}
 

 When acting on a domain, the operation \texttt{PositionCanonical} (\ref{PositionCanonical}) is used to determine the position of elements in the domain. This can be used
to act on a domain given by a list of representatives for which \texttt{PositionCanonical} (\ref{PositionCanonical}) is implemented, for example the return value of \texttt{RightTransversal} (\ref{RightTransversal}). }

 

\subsection{\textcolor{Chapter }{Action (for a group, an action domain, etc.)}}
\logpage{[ 41, 7, 2 ]}\nobreak
\hyperdef{L}{X85A8E93D786C3C9C}{}
{\noindent\textcolor{FuncColor}{$\triangleright$\enspace\texttt{Action({\mdseries\slshape G, Omega[, gens, acts][, act]})\index{Action@\texttt{Action}!for a group, an action domain, etc.}
\label{Action:for a group, an action domain, etc.}
}\hfill{\scriptsize (function)}}\\
\noindent\textcolor{FuncColor}{$\triangleright$\enspace\texttt{Action({\mdseries\slshape xset})\index{Action@\texttt{Action}!for an external set}
\label{Action:for an external set}
}\hfill{\scriptsize (function)}}\\


 returns the image group of \texttt{ActionHomomorphism} (\ref{ActionHomomorphism:for a group, an action domain, etc.}) called with the same parameters. 

 Note that (for compatibility reasons to be able to get the action
homomorphism) this image group internally stores the action homomorphism. If \mbox{\texttt{\mdseries\slshape G}} or \mbox{\texttt{\mdseries\slshape Omega}} are extremely big, this can cause memory problems. In this case compute only
generator images and form the image group yourself. 

 (See Section{\nobreakspace}\ref{Basic Actions} for information about specific actions.) 

 \index{regular action} The following code shows for example how to create the regular action of a
group. 

 
\begin{Verbatim}[commandchars=!@|,fontsize=\small,frame=single,label=Example]
  !gapprompt@gap>| !gapinput@g:=Group((1,2,3),(1,2));;|
  !gapprompt@gap>| !gapinput@Action(g,AsList(g),OnRight);|
  Group([ (1,5,3)(2,6,4), (1,6)(2,5)(3,4) ])
\end{Verbatim}
 }

 

\subsection{\textcolor{Chapter }{SparseActionHomomorphism}}
\logpage{[ 41, 7, 3 ]}\nobreak
\hyperdef{L}{X86FF54A383B73967}{}
{\noindent\textcolor{FuncColor}{$\triangleright$\enspace\texttt{SparseActionHomomorphism({\mdseries\slshape G, start[, gens, acts][, act]})\index{SparseActionHomomorphism@\texttt{SparseActionHomomorphism}}
\label{SparseActionHomomorphism}
}\hfill{\scriptsize (operation)}}\\
\noindent\textcolor{FuncColor}{$\triangleright$\enspace\texttt{SortedSparseActionHomomorphism({\mdseries\slshape G, start[, gens, acts][, act]})\index{SortedSparseActionHomomorphism@\texttt{SortedSparseActionHomomorphism}}
\label{SortedSparseActionHomomorphism}
}\hfill{\scriptsize (operation)}}\\


 \texttt{SparseActionHomomorphism} computes the action homomorphism (see \texttt{ActionHomomorphism} (\ref{ActionHomomorphism:for a group, an action domain, etc.})) with arguments \mbox{\texttt{\mdseries\slshape G}}, $D$, and the optional arguments given, where $D$ is the union of the \mbox{\texttt{\mdseries\slshape G}}\texttt{\symbol{45}}orbits of all points in \mbox{\texttt{\mdseries\slshape start}}. In the \texttt{Orbit} (\ref{Orbit}) calls that are used to create $D$, again the optional arguments given are entered.) 

 If \mbox{\texttt{\mdseries\slshape G}} acts on a very large domain not surjectively this may yield a permutation
image of substantially smaller degree than by action on the whole domain. 

 The operation \texttt{SparseActionHomomorphism} will only use \texttt{\texttt{\symbol{92}}=} (\ref{bSlash=}) comparisons of points in the orbit. Therefore it can be used even if no good \texttt{\texttt{\symbol{92}}{\textless}} (\ref{bSlash<}) comparison method for these points is available. However the image group will
depend on the generators \mbox{\texttt{\mdseries\slshape gens}} of \mbox{\texttt{\mdseries\slshape G}}. 

 The operation \texttt{SortedSparseActionHomomorphism} in contrast will sort the orbit and thus produce an image group which does not
depend on these generators. 

 
\begin{Verbatim}[commandchars=!@|,fontsize=\small,frame=single,label=Example]
  !gapprompt@gap>| !gapinput@h:=Group(Z(3)*[[[1,1],[0,1]]]);|
  Group([ [ [ Z(3), Z(3) ], [ 0*Z(3), Z(3) ] ] ])
  !gapprompt@gap>| !gapinput@hom:=ActionHomomorphism(h,GF(3)^2,OnRight);;|
  !gapprompt@gap>| !gapinput@Image(hom);|
  Group([ (2,3)(4,9,6,7,5,8) ])
  !gapprompt@gap>| !gapinput@hom:=SparseActionHomomorphism(h,[Z(3)*[1,0]],OnRight);;|
  !gapprompt@gap>| !gapinput@Image(hom);|
  Group([ (1,2,3,4,5,6) ])
\end{Verbatim}
 }

 }

  
\section{\textcolor{Chapter }{Action of a group on itself}}\label{Action of a group on itself}
\logpage{[ 41, 8, 0 ]}
\hyperdef{L}{X7FED50ED7ACA5FB2}{}
{
  Of particular importance is the action of a group on its elements or cosets of
a subgroup. These actions can be obtained by using \texttt{ActionHomomorphism} (\ref{ActionHomomorphism:for a group, an action domain, etc.}) for a suitable domain (for example a list of subgroups). For the following
(frequently used) types of actions however special (often particularly
efficient) functions are provided. A special case is the regular action on all
elements. 

\subsection{\textcolor{Chapter }{FactorCosetAction (for a group and subgroup)}}
\logpage{[ 41, 8, 1 ]}\nobreak
\hyperdef{L}{X784D417D87F4E58D}{}
{\noindent\textcolor{FuncColor}{$\triangleright$\enspace\texttt{FactorCosetAction({\mdseries\slshape G, U[, N]})\index{FactorCosetAction@\texttt{FactorCosetAction}!for a group and subgroup}
\label{FactorCosetAction:for a group and subgroup}
}\hfill{\scriptsize (operation)}}\\
\noindent\textcolor{FuncColor}{$\triangleright$\enspace\texttt{FactorCosetAction({\mdseries\slshape G, L})\index{FactorCosetAction@\texttt{FactorCosetAction}!for a group and list of subgroups}
\label{FactorCosetAction:for a group and list of subgroups}
}\hfill{\scriptsize (operation)}}\\


 This command computes the action of the group \mbox{\texttt{\mdseries\slshape G}} on the right cosets of the subgroup \mbox{\texttt{\mdseries\slshape U}}. If a normal subgroup \mbox{\texttt{\mdseries\slshape N}} of \mbox{\texttt{\mdseries\slshape G}} is given, it is stored as kernel of this action. When calling \texttt{FactorCosetAction} with a list of subgroups as the second argument, an action with image
isomorphic to the subdirect product of the coset actions of all subgroups is
computed. (However a degree reduction may take place if some of the actions
are redundant, i.e. there is no guarantee that every subgroup in the list is
represented by an orbit.) 
\begin{Verbatim}[commandchars=!@|,fontsize=\small,frame=single,label=Example]
  !gapprompt@gap>| !gapinput@g:=Group((1,2,3,4,5),(1,2));;u:=SylowSubgroup(g,2);;Index(g,u);|
  15
  !gapprompt@gap>| !gapinput@FactorCosetAction(g,u);|
  [ (1,2,3,4,5), (1,2) ] -> [ (1,4,7,10,13)(2,5,8,11,14)(3,6,9,12,15),\
    (1,4)(2,6)(3,5)(7,8)(10,12)(13,14) ]
  !gapprompt@gap>| !gapinput@StructureDescription(Range(last));|
  "S5"
  !gapprompt@gap>| !gapinput@FactorCosetAction(g,[u,SylowSubgroup(g,3)]);;|
  !gapprompt@gap>| !gapinput@Size(Image(last));|
  120
\end{Verbatim}
 The correspondence of points with cosets will, for performance reasons, depend
on the method used. It is not guaranteed that it will be the same as used by \texttt{RightTransversal} or \texttt{RightCosets}. }

 

\subsection{\textcolor{Chapter }{RegularActionHomomorphism}}
\logpage{[ 41, 8, 2 ]}\nobreak
\hyperdef{L}{X8561DEBA79E01ABD}{}
{\noindent\textcolor{FuncColor}{$\triangleright$\enspace\texttt{RegularActionHomomorphism({\mdseries\slshape G})\index{RegularActionHomomorphism@\texttt{RegularActionHomomorphism}}
\label{RegularActionHomomorphism}
}\hfill{\scriptsize (attribute)}}\\


 returns an isomorphism from \mbox{\texttt{\mdseries\slshape G}} onto the regular permutation representation of \mbox{\texttt{\mdseries\slshape G}}. }

 

\subsection{\textcolor{Chapter }{AbelianSubfactorAction}}
\logpage{[ 41, 8, 3 ]}\nobreak
\hyperdef{L}{X835317A7847477D4}{}
{\noindent\textcolor{FuncColor}{$\triangleright$\enspace\texttt{AbelianSubfactorAction({\mdseries\slshape G, M, N})\index{AbelianSubfactorAction@\texttt{AbelianSubfactorAction}}
\label{AbelianSubfactorAction}
}\hfill{\scriptsize (operation)}}\\


 Let \mbox{\texttt{\mdseries\slshape G}} be a group and $\mbox{\texttt{\mdseries\slshape M}} \geq \mbox{\texttt{\mdseries\slshape N}}$ be subgroups of a common parent that are normal under \mbox{\texttt{\mdseries\slshape G}}, such that the subfactor $\mbox{\texttt{\mdseries\slshape M}}/\mbox{\texttt{\mdseries\slshape N}}$ is elementary abelian. The operation \texttt{AbelianSubfactorAction} returns a list \texttt{[ \mbox{\texttt{\mdseries\slshape phi}}, \mbox{\texttt{\mdseries\slshape alpha}}, \mbox{\texttt{\mdseries\slshape bas}} ]} where \mbox{\texttt{\mdseries\slshape bas}} is a list of elements of \mbox{\texttt{\mdseries\slshape M}} which are representatives for a basis of $\mbox{\texttt{\mdseries\slshape M}}/\mbox{\texttt{\mdseries\slshape N}}$, \mbox{\texttt{\mdseries\slshape alpha}} is a map from \mbox{\texttt{\mdseries\slshape M}} into a $n$\texttt{\symbol{45}}dimensional row space over $GF(p)$ where $[\mbox{\texttt{\mdseries\slshape M}}:\mbox{\texttt{\mdseries\slshape N}}] = p^n$ that is the natural homomorphism of \mbox{\texttt{\mdseries\slshape M}} by \mbox{\texttt{\mdseries\slshape N}} with the quotient represented as an additive group. Finally \mbox{\texttt{\mdseries\slshape phi}} is a homomorphism from \mbox{\texttt{\mdseries\slshape G}} into $GL_n(p)$ that represents the action of \mbox{\texttt{\mdseries\slshape G}} on the factor $\mbox{\texttt{\mdseries\slshape M}}/\mbox{\texttt{\mdseries\slshape N}}$. 

 Note: If only matrices for the action are needed, \texttt{LinearActionLayer} (\ref{LinearActionLayer}) might be faster. 
\begin{Verbatim}[commandchars=!@|,fontsize=\small,frame=single,label=Example]
  !gapprompt@gap>| !gapinput@g:=Group((1,8,10,7,3,5)(2,4,12,9,11,6),|
  !gapprompt@>| !gapinput@            (1,9,5,6,3,10)(2,11,12,8,4,7));;|
  !gapprompt@gap>| !gapinput@c:=ChiefSeries(g);;List(c,Size);|
  [ 96, 48, 16, 4, 1 ]
  !gapprompt@gap>| !gapinput@HasElementaryAbelianFactorGroup(c[3],c[4]);|
  true
  !gapprompt@gap>| !gapinput@SetName(c[3],"my_group");;|
  !gapprompt@gap>| !gapinput@a:=AbelianSubfactorAction(g,c[3],c[4]);|
  [ [ (1,8,10,7,3,5)(2,4,12,9,11,6), (1,9,5,6,3,10)(2,11,12,8,4,7) ] ->
      [ <an immutable 2x2 matrix over GF2>,
        <an immutable 2x2 matrix over GF2> ],
    MappingByFunction( my_group, ( GF(2)^
      2 ), function( e ) ... end, function( r ) ... end ),
    Pcgs([ (2,9,3,8)(4,11,5,10), (1,6,12,7)(4,10,5,11) ]) ]
  !gapprompt@gap>| !gapinput@mat:=Image(a[1],g);|
  Group([ <an immutable 2x2 matrix over GF2>,
    <an immutable 2x2 matrix over GF2> ])
  !gapprompt@gap>| !gapinput@Size(mat);|
  3
  !gapprompt@gap>| !gapinput@e:=PreImagesRepresentative(a[2],[Z(2),0*Z(2)]);|
  (2,9,3,8)(4,11,5,10)
  !gapprompt@gap>| !gapinput@e in c[3];e in c[4];|
  true
  false
\end{Verbatim}
 }

 }

  
\section{\textcolor{Chapter }{Permutations Induced by Elements and Cycles}}\label{Permutations Induced by Elements and Cycles}
\logpage{[ 41, 9, 0 ]}
\hyperdef{L}{X807AA91E841D132B}{}
{
  If only the permutation image of a single element is needed, it might not be
worth to create the action homomorphism, the following operations yield the
permutation image and cycles of a single element. 
\subsection{\textcolor{Chapter }{Permutation}}\logpage{[ 41, 9, 1 ]}
\hyperdef{L}{X7807A33381DCAB26}{}
{
\noindent\textcolor{FuncColor}{$\triangleright$\enspace\texttt{Permutation({\mdseries\slshape g, Omega[, gens, acts][, act]})\index{Permutation@\texttt{Permutation}!for a group, an action domain, etc.}
\label{Permutation:for a group, an action domain, etc.}
}\hfill{\scriptsize (function)}}\\
\noindent\textcolor{FuncColor}{$\triangleright$\enspace\texttt{Permutation({\mdseries\slshape g, xset})\index{Permutation@\texttt{Permutation}!for an external set}
\label{Permutation:for an external set}
}\hfill{\scriptsize (function)}}\\


 computes the permutation that corresponds to the action of \mbox{\texttt{\mdseries\slshape g}} on the permutation domain \mbox{\texttt{\mdseries\slshape Omega}} (a list of objects that are permuted). If an external set \mbox{\texttt{\mdseries\slshape xset}} is given, the permutation domain is the \texttt{HomeEnumerator} (\ref{HomeEnumerator}) value of this external set (see Section{\nobreakspace}\ref{External Sets}). Note that the points of the returned permutation refer to the positions in \mbox{\texttt{\mdseries\slshape Omega}}, even if \mbox{\texttt{\mdseries\slshape Omega}} itself consists of integers. 

 If \mbox{\texttt{\mdseries\slshape g}} does not leave the domain invariant, or does not map the domain injectively
then \texttt{fail} is returned. }

 

\subsection{\textcolor{Chapter }{PermutationCycle}}
\logpage{[ 41, 9, 2 ]}\nobreak
\hyperdef{L}{X81D4EA42810974A0}{}
{\noindent\textcolor{FuncColor}{$\triangleright$\enspace\texttt{PermutationCycle({\mdseries\slshape g, Omega, pnt[, act]})\index{PermutationCycle@\texttt{PermutationCycle}}
\label{PermutationCycle}
}\hfill{\scriptsize (function)}}\\


 computes the permutation that represents the cycle of \mbox{\texttt{\mdseries\slshape pnt}} under the action of the element \mbox{\texttt{\mdseries\slshape g}}. 
\begin{Verbatim}[commandchars=!@|,fontsize=\small,frame=single,label=Example]
  !gapprompt@gap>| !gapinput@Permutation([[Z(3),-Z(3)],[Z(3),0*Z(3)]],AsList(GF(3)^2));|
  (2,7,6)(3,4,8)
  !gapprompt@gap>| !gapinput@Permutation((1,2,3)(4,5)(6,7),[4..7]);|
  (1,2)(3,4)
  !gapprompt@gap>| !gapinput@PermutationCycle((1,2,3)(4,5)(6,7),[4..7],4);|
  (1,2)
\end{Verbatim}
 }

 

\subsection{\textcolor{Chapter }{Cycle}}
\logpage{[ 41, 9, 3 ]}\nobreak
\hyperdef{L}{X80AF6E0683CA7F14}{}
{\noindent\textcolor{FuncColor}{$\triangleright$\enspace\texttt{Cycle({\mdseries\slshape g, Omega, pnt[, act]})\index{Cycle@\texttt{Cycle}}
\label{Cycle}
}\hfill{\scriptsize (function)}}\\


 returns a list of the points in the cycle of \mbox{\texttt{\mdseries\slshape pnt}} under the action of the element \mbox{\texttt{\mdseries\slshape g}}. }

 

\subsection{\textcolor{Chapter }{CycleLength}}
\logpage{[ 41, 9, 4 ]}\nobreak
\hyperdef{L}{X7F559E897B333758}{}
{\noindent\textcolor{FuncColor}{$\triangleright$\enspace\texttt{CycleLength({\mdseries\slshape g, Omega, pnt[, act]})\index{CycleLength@\texttt{CycleLength}}
\label{CycleLength}
}\hfill{\scriptsize (function)}}\\


 returns the length of the cycle of \mbox{\texttt{\mdseries\slshape pnt}} under the action of the element \mbox{\texttt{\mdseries\slshape g}}. }

 

\subsection{\textcolor{Chapter }{Cycles}}
\logpage{[ 41, 9, 5 ]}\nobreak
\hyperdef{L}{X7F3B387A7FD8AE5E}{}
{\noindent\textcolor{FuncColor}{$\triangleright$\enspace\texttt{Cycles({\mdseries\slshape g, Omega[, act]})\index{Cycles@\texttt{Cycles}}
\label{Cycles}
}\hfill{\scriptsize (function)}}\\


 returns a list of the cycles (as lists of points) of the action of the element \mbox{\texttt{\mdseries\slshape g}}. }

 

\subsection{\textcolor{Chapter }{CycleLengths}}
\logpage{[ 41, 9, 6 ]}\nobreak
\hyperdef{L}{X83040A6080C2C6C6}{}
{\noindent\textcolor{FuncColor}{$\triangleright$\enspace\texttt{CycleLengths({\mdseries\slshape g, Omega[, act]})\index{CycleLengths@\texttt{CycleLengths}}
\label{CycleLengths}
}\hfill{\scriptsize (operation)}}\\


 returns the lengths of all the cycles under the action of the element \mbox{\texttt{\mdseries\slshape g}} on \mbox{\texttt{\mdseries\slshape Omega}}. 
\begin{Verbatim}[commandchars=!@|,fontsize=\small,frame=single,label=Example]
  !gapprompt@gap>| !gapinput@Cycle((1,2,3)(4,5)(6,7),[4..7],4);|
  [ 4, 5 ]
  !gapprompt@gap>| !gapinput@CycleLength((1,2,3)(4,5)(6,7),[4..7],4);|
  2
  !gapprompt@gap>| !gapinput@Cycles((1,2,3)(4,5)(6,7),[4..7]);|
  [ [ 4, 5 ], [ 6, 7 ] ]
  !gapprompt@gap>| !gapinput@CycleLengths((1,2,3)(4,5)(6,7),[4..7]);|
  [ 2, 2 ]
\end{Verbatim}
 }

 
\subsection{\textcolor{Chapter }{CycleIndex}}\logpage{[ 41, 9, 7 ]}
\hyperdef{L}{X87FDA6838065CDCB}{}
{
\noindent\textcolor{FuncColor}{$\triangleright$\enspace\texttt{CycleIndex({\mdseries\slshape g, Omega[, act]})\index{CycleIndex@\texttt{CycleIndex}!for a permutation and an action domain}
\label{CycleIndex:for a permutation and an action domain}
}\hfill{\scriptsize (function)}}\\
\noindent\textcolor{FuncColor}{$\triangleright$\enspace\texttt{CycleIndex({\mdseries\slshape G, Omega[, act]})\index{CycleIndex@\texttt{CycleIndex}!for a permutation group and an action domain}
\label{CycleIndex:for a permutation group and an action domain}
}\hfill{\scriptsize (function)}}\\


 The \emph{cycle index} of a permutation \mbox{\texttt{\mdseries\slshape g}} acting on \mbox{\texttt{\mdseries\slshape Omega}} is defined as 
\[ z(\mbox{\texttt{\mdseries\slshape g}}) = s_1^{{c_1}} s_2^{{c_2}} \cdots s_n^{{c_n}} \]
 where $c_k$ is the number of $k$\texttt{\symbol{45}}cycles in the cycle decomposition of \mbox{\texttt{\mdseries\slshape g}} and the $s_i$ are indeterminates. 

 The \emph{cycle index} of a group \mbox{\texttt{\mdseries\slshape G}} is defined as 
\[ Z(\mbox{\texttt{\mdseries\slshape G}}) = \left( \sum_{{g \in \mbox{\texttt{\mdseries\slshape G}}}} z(g) \right) / |\mbox{\texttt{\mdseries\slshape G}}| . \]
 

 The indeterminates used by \texttt{CycleIndex} are the indeterminates $1$ to $n$ over the rationals (see{\nobreakspace}\texttt{Indeterminate} (\ref{Indeterminate:for a ring (and a number)})). 

 
\begin{Verbatim}[commandchars=!@|,fontsize=\small,frame=single,label=Example]
  !gapprompt@gap>| !gapinput@g:=TransitiveGroup(6,8);|
  S_4(6c) = 1/2[2^3]S(3)
  !gapprompt@gap>| !gapinput@CycleIndex(g);|
  1/24*x_1^6+1/8*x_1^2*x_2^2+1/4*x_1^2*x_4+1/4*x_2^3+1/3*x_3^2
\end{Verbatim}
 }

 }

  
\section{\textcolor{Chapter }{Tests for Actions}}\label{Tests for Actions}
\logpage{[ 41, 10, 0 ]}
\hyperdef{L}{X850A84618421392A}{}
{
  
\subsection{\textcolor{Chapter }{IsTransitive}}\logpage{[ 41, 10, 1 ]}
\hyperdef{L}{X79B15750851828CB}{}
{
\noindent\textcolor{FuncColor}{$\triangleright$\enspace\texttt{IsTransitive({\mdseries\slshape G, Omega[, gens, acts][, act]})\index{IsTransitive@\texttt{IsTransitive}!for a group, an action domain, etc.}
\label{IsTransitive:for a group, an action domain, etc.}
}\hfill{\scriptsize (operation)}}\\
\noindent\textcolor{FuncColor}{$\triangleright$\enspace\texttt{IsTransitive({\mdseries\slshape G})\index{IsTransitive@\texttt{IsTransitive}!for a permutation group}
\label{IsTransitive:for a permutation group}
}\hfill{\scriptsize (property)}}\\
\noindent\textcolor{FuncColor}{$\triangleright$\enspace\texttt{IsTransitive({\mdseries\slshape xset})\index{IsTransitive@\texttt{IsTransitive}!for an external set}
\label{IsTransitive:for an external set}
}\hfill{\scriptsize (property)}}\\


 returns \texttt{true} if the action implied by the arguments is transitive, or \texttt{false} otherwise. 

 \index{transitive} We say that a group \mbox{\texttt{\mdseries\slshape G}} acts \emph{transitively} on a domain $D$ if and only if \mbox{\texttt{\mdseries\slshape G}} acts on $D$ and for every pair of points $d, e \in D$ there is an element $g$ in \mbox{\texttt{\mdseries\slshape G}} such that $d^g = e$. 

 For a permutation group \mbox{\texttt{\mdseries\slshape G}}, one may also invoke this as \texttt{IsTransitive(\mbox{\texttt{\mdseries\slshape G}})}, which tests whether the group is transitive with respect to its natural
action on the set of points moved by it. For example the group $\langle (2,3,4),(2,3) \rangle$ is transitive on the set $\{2, 3, 4\}$. 
\begin{Verbatim}[commandchars=!@|,fontsize=\small,frame=single,label=Example]
  !gapprompt@gap>| !gapinput@G:= Group( (2,3,4), (2,3) );;|
  !gapprompt@gap>| !gapinput@IsTransitive( G, [ 2 .. 4 ] );|
  true
  !gapprompt@gap>| !gapinput@IsTransitive( G, [ 2, 3 ] );   # G does not act on [ 2, 3 ]|
  false
  !gapprompt@gap>| !gapinput@IsTransitive( G, [ 1 .. 4 ] );  # G has two orbits on [ 1 .. 4 ]|
  false
  !gapprompt@gap>| !gapinput@IsTransitive( G );  # G is transitive on [ 2 .. 4 ]|
  true
  !gapprompt@gap>| !gapinput@IsTransitive( SL(2, 3), NormedRowVectors( GF(3)^2 ) );|
  false
  !gapprompt@gap>| !gapinput@IsTransitive( SL(2, 3), NormedRowVectors( GF(3)^2 ), OnLines );|
  true
\end{Verbatim}
 }

 
\subsection{\textcolor{Chapter }{Transitivity}}\logpage{[ 41, 10, 2 ]}
\hyperdef{L}{X8295D733796B7A37}{}
{
\noindent\textcolor{FuncColor}{$\triangleright$\enspace\texttt{Transitivity({\mdseries\slshape G, Omega[, gens, acts][, act]})\index{Transitivity@\texttt{Transitivity}!for a group and an action domain}
\label{Transitivity:for a group and an action domain}
}\hfill{\scriptsize (operation)}}\\
\noindent\textcolor{FuncColor}{$\triangleright$\enspace\texttt{Transitivity({\mdseries\slshape G})\index{Transitivity@\texttt{Transitivity}!for a permutation group}
\label{Transitivity:for a permutation group}
}\hfill{\scriptsize (attribute)}}\\
\noindent\textcolor{FuncColor}{$\triangleright$\enspace\texttt{Transitivity({\mdseries\slshape xset})\index{Transitivity@\texttt{Transitivity}!for an external set}
\label{Transitivity:for an external set}
}\hfill{\scriptsize (attribute)}}\\


 returns the degree $k$ (a non\texttt{\symbol{45}}negative integer) of transitivity of the action
implied by the arguments, i.e. the largest integer $k$ such that the action is $k$\texttt{\symbol{45}}transitive. If the action is not transitive \texttt{0} is returned. 

 An action is \emph{$k$\texttt{\symbol{45}}transitive} if every $k$\texttt{\symbol{45}}tuple of points can be mapped simultaneously to every
other $k$\texttt{\symbol{45}}tuple. 

 For a permutation group \mbox{\texttt{\mdseries\slshape G}}, one may also invoke this as \texttt{Transitivity(\mbox{\texttt{\mdseries\slshape G}})}, which returns the degree of transitivity of the group with respect to its
natural action on the set of points moved by it. For example the group $\langle (2,3,4),(2,3) \rangle$ is 3\texttt{\symbol{45}}transitive on the set $\{2, 3, 4\}$. 
\begin{Verbatim}[commandchars=!@|,fontsize=\small,frame=single,label=Example]
  !gapprompt@gap>| !gapinput@g:=Group((1,3,2),(2,4,3));;|
  !gapprompt@gap>| !gapinput@IsTransitive(g,[1..5]);|
  false
  !gapprompt@gap>| !gapinput@Transitivity(g,[1..4]);|
  2
  !gapprompt@gap>| !gapinput@Transitivity(g);|
  2
\end{Verbatim}
 }

 
\subsection{\textcolor{Chapter }{RankAction}}\logpage{[ 41, 10, 3 ]}
\hyperdef{L}{X8166A6A17C8D6E73}{}
{
\noindent\textcolor{FuncColor}{$\triangleright$\enspace\texttt{RankAction({\mdseries\slshape G, Omega[, gens, acts][, act]})\index{RankAction@\texttt{RankAction}!for a group, an action domain, etc.}
\label{RankAction:for a group, an action domain, etc.}
}\hfill{\scriptsize (operation)}}\\
\noindent\textcolor{FuncColor}{$\triangleright$\enspace\texttt{RankAction({\mdseries\slshape xset})\index{RankAction@\texttt{RankAction}!for an external set}
\label{RankAction:for an external set}
}\hfill{\scriptsize (attribute)}}\\


 returns the rank of the transitive (see \texttt{IsTransitive} (\ref{IsTransitive:for a group, an action domain, etc.})) action of \mbox{\texttt{\mdseries\slshape G}} on \mbox{\texttt{\mdseries\slshape Omega}}, i. e., the number of orbits of any point stabilizer. 
\begin{Verbatim}[commandchars=!@|,fontsize=\small,frame=single,label=Example]
  !gapprompt@gap>| !gapinput@RankAction(g,Combinations([1..4],2),OnSets);|
  4
\end{Verbatim}
 }

 
\subsection{\textcolor{Chapter }{IsSemiRegular}}\logpage{[ 41, 10, 4 ]}
\hyperdef{L}{X7B77040F8543CD6E}{}
{
\noindent\textcolor{FuncColor}{$\triangleright$\enspace\texttt{IsSemiRegular({\mdseries\slshape G, Omega[, gens, acts][, act]})\index{IsSemiRegular@\texttt{IsSemiRegular}!for a group, an action domain, etc.}
\label{IsSemiRegular:for a group, an action domain, etc.}
}\hfill{\scriptsize (operation)}}\\
\noindent\textcolor{FuncColor}{$\triangleright$\enspace\texttt{IsSemiRegular({\mdseries\slshape G})\index{IsSemiRegular@\texttt{IsSemiRegular}!for a permutation group}
\label{IsSemiRegular:for a permutation group}
}\hfill{\scriptsize (property)}}\\
\noindent\textcolor{FuncColor}{$\triangleright$\enspace\texttt{IsSemiRegular({\mdseries\slshape xset})\index{IsSemiRegular@\texttt{IsSemiRegular}!for an external set}
\label{IsSemiRegular:for an external set}
}\hfill{\scriptsize (property)}}\\


 returns \texttt{true} if the action implied by the arguments is semiregular, or \texttt{false} otherwise. 

 \index{semiregular} An action is \emph{semiregular} if the stabilizer of each point is the identity. 

 For a permutation group \mbox{\texttt{\mdseries\slshape G}}, one may also invoke this as \texttt{IsSemiRegular(\mbox{\texttt{\mdseries\slshape G}})}, which tests whether the group is semiregular with respect to its natural
action on the set of points moved by it. For example the group $\langle (2,3,4) (5,6,7) \rangle$ is semiregular on the set $\{2, 3, 4, 5, 6, 7\}$. }

 
\subsection{\textcolor{Chapter }{IsRegular}}\logpage{[ 41, 10, 5 ]}
\hyperdef{L}{X7CF02C4785F0EAB5}{}
{
\noindent\textcolor{FuncColor}{$\triangleright$\enspace\texttt{IsRegular({\mdseries\slshape G, Omega[, gens, acts][, act]})\index{IsRegular@\texttt{IsRegular}!for a group, an action domain, etc.}
\label{IsRegular:for a group, an action domain, etc.}
}\hfill{\scriptsize (operation)}}\\
\noindent\textcolor{FuncColor}{$\triangleright$\enspace\texttt{IsRegular({\mdseries\slshape G})\index{IsRegular@\texttt{IsRegular}!for a permutation group}
\label{IsRegular:for a permutation group}
}\hfill{\scriptsize (property)}}\\
\noindent\textcolor{FuncColor}{$\triangleright$\enspace\texttt{IsRegular({\mdseries\slshape xset})\index{IsRegular@\texttt{IsRegular}!for an external set}
\label{IsRegular:for an external set}
}\hfill{\scriptsize (property)}}\\


 returns \texttt{true} if the action implied by the arguments is regular, or \texttt{false} otherwise. 

 \index{regular} An action is \emph{regular} if it is both semiregular (see{\nobreakspace}\texttt{IsSemiRegular} (\ref{IsSemiRegular:for a group, an action domain, etc.})) and transitive (see{\nobreakspace}\texttt{IsTransitive} (\ref{IsTransitive:for a group, an action domain, etc.})). In this case every point \mbox{\texttt{\mdseries\slshape pnt}} of \mbox{\texttt{\mdseries\slshape Omega}} defines a one\texttt{\symbol{45}}to\texttt{\symbol{45}}one correspondence
between \mbox{\texttt{\mdseries\slshape G}} and \mbox{\texttt{\mdseries\slshape Omega}}. 

 For a permutation group \mbox{\texttt{\mdseries\slshape G}}, one may also invoke this as \texttt{IsRegular(\mbox{\texttt{\mdseries\slshape G}})}, which tests whether the group is regular with respect to its natural action
on the set of points moved by it. For example the group $\langle (2,3,4) \rangle$ is regular on the set $\{2, 3, 4\}$. 
\begin{Verbatim}[commandchars=!@|,fontsize=\small,frame=single,label=Example]
  !gapprompt@gap>| !gapinput@IsSemiRegular(g,Arrangements([1..4],3),OnTuples);|
  true
  !gapprompt@gap>| !gapinput@IsRegular(g,Arrangements([1..4],3),OnTuples);|
  false
\end{Verbatim}
 }

 
\subsection{\textcolor{Chapter }{Earns}}\logpage{[ 41, 10, 6 ]}
\hyperdef{L}{X7CB1D74280F92AFC}{}
{
\noindent\textcolor{FuncColor}{$\triangleright$\enspace\texttt{Earns({\mdseries\slshape G, Omega[, gens, acts][, act]})\index{Earns@\texttt{Earns}!for a group, an action domain, etc.}
\label{Earns:for a group, an action domain, etc.}
}\hfill{\scriptsize (operation)}}\\
\noindent\textcolor{FuncColor}{$\triangleright$\enspace\texttt{Earns({\mdseries\slshape xset})\index{Earns@\texttt{Earns}!for an external set}
\label{Earns:for an external set}
}\hfill{\scriptsize (attribute)}}\\


 returns a list of the elementary abelian regular (when acting on \mbox{\texttt{\mdseries\slshape Omega}}) normal subgroups of \mbox{\texttt{\mdseries\slshape G}}. 

 At the moment only methods for a primitive group \mbox{\texttt{\mdseries\slshape G}} are implemented. }

 
\subsection{\textcolor{Chapter }{IsPrimitive}}\logpage{[ 41, 10, 7 ]}
\hyperdef{L}{X84C19AD68247B760}{}
{
\noindent\textcolor{FuncColor}{$\triangleright$\enspace\texttt{IsPrimitive({\mdseries\slshape G, Omega[, gens, acts][, act]})\index{IsPrimitive@\texttt{IsPrimitive}!for a group, an action domain, etc.}
\label{IsPrimitive:for a group, an action domain, etc.}
}\hfill{\scriptsize (operation)}}\\
\noindent\textcolor{FuncColor}{$\triangleright$\enspace\texttt{IsPrimitive({\mdseries\slshape G})\index{IsPrimitive@\texttt{IsPrimitive}!for a permutation group}
\label{IsPrimitive:for a permutation group}
}\hfill{\scriptsize (property)}}\\
\noindent\textcolor{FuncColor}{$\triangleright$\enspace\texttt{IsPrimitive({\mdseries\slshape xset})\index{IsPrimitive@\texttt{IsPrimitive}!for an external set}
\label{IsPrimitive:for an external set}
}\hfill{\scriptsize (property)}}\\


 returns \texttt{true} if the action implied by the arguments is primitive, or \texttt{false} otherwise. 

 \index{primitive} An action is \emph{primitive} if it is transitive (see \texttt{IsTransitive} (\ref{IsTransitive:for a group, an action domain, etc.})) and the action admits no nontrivial block systems. See{\nobreakspace}\ref{Block Systems} for the definition of block systems. 

 For a permutation group \mbox{\texttt{\mdseries\slshape G}}, one may also invoke this as \texttt{IsPrimitive(\mbox{\texttt{\mdseries\slshape G}})}, which tests whether the group is primitive with respect to its natural
action on the set of points moved by it. For example the group $\langle (2,3,4),(2,3) \rangle$ is primitive on the set $\{2, 3, 4\}$. 

 For an explanation of the meaning of all the inputs, please refer to
{\nobreakspace}\ref{About Group Actions}. 

 \emph{Note:} This operation does not tell whether a matrix group is primitive in the sense
of preserving a direct sum of vector spaces. To do this use \texttt{IsPrimitiveMatrixGroup} or \texttt{IsPrimitive} from the package \textsf{IRREDSOL}. 
\begin{Verbatim}[commandchars=!@|,fontsize=\small,frame=single,label=Example]
  !gapprompt@gap>| !gapinput@IsPrimitive(g,Orbit(g,(1,2)(3,4)));|
  true
\end{Verbatim}
 }

  }

  
\section{\textcolor{Chapter }{Block Systems}}\label{Block Systems}
\logpage{[ 41, 11, 0 ]}
\hyperdef{L}{X7E9D3D0B7A9A8572}{}
{
  A \emph{block system} (system of imprimitivity) for the action of a group $G$ on an action domain $\Omega$ is a partition of $\Omega$ which {\textendash}as a partition{\textendash} remains invariant under the
action of $G$. For operations concerning block systems, \textsf{GAP} assumes that $G$ acts transitively on $\Omega$ (see \texttt{IsTransitive} (\ref{IsTransitive:for a group, an action domain, etc.})). One may get wrong results or error messages (perhaps at a much later stage)
if this condition is not satisfied. 
\subsection{\textcolor{Chapter }{Blocks}}\logpage{[ 41, 11, 1 ]}
\hyperdef{L}{X84FE699F85371643}{}
{
\noindent\textcolor{FuncColor}{$\triangleright$\enspace\texttt{Blocks({\mdseries\slshape G, Omega[, seed][, gens, acts][, act]})\index{Blocks@\texttt{Blocks}!for a group, an action domain, etc.}
\label{Blocks:for a group, an action domain, etc.}
}\hfill{\scriptsize (operation)}}\\
\noindent\textcolor{FuncColor}{$\triangleright$\enspace\texttt{Blocks({\mdseries\slshape xset[, seed]})\index{Blocks@\texttt{Blocks}!for an external set}
\label{Blocks:for an external set}
}\hfill{\scriptsize (attribute)}}\\


 computes a block system for the transitive (see \texttt{IsTransitive} (\ref{IsTransitive:for a group, an action domain, etc.})) action of \mbox{\texttt{\mdseries\slshape G}} on \mbox{\texttt{\mdseries\slshape Omega}}. If \mbox{\texttt{\mdseries\slshape seed}} is not given and the action is imprimitive, a minimal nontrivial block system
will be found. If \mbox{\texttt{\mdseries\slshape seed}} is given, a block system in which \mbox{\texttt{\mdseries\slshape seed}} is the subset of one block is computed. 

 The result is undefined if the action is not transitive. 
\begin{Verbatim}[commandchars=!@|,fontsize=\small,frame=single,label=Example]
  !gapprompt@gap>| !gapinput@g:=TransitiveGroup(8,3);|
  E(8)=2[x]2[x]2
  !gapprompt@gap>| !gapinput@Blocks(g,[1..8]);|
  [ [ 1, 8 ], [ 2, 3 ], [ 4, 5 ], [ 6, 7 ] ]
  !gapprompt@gap>| !gapinput@Blocks(g,[1..8],[1,4]);|
  [ [ 1, 4 ], [ 2, 7 ], [ 3, 6 ], [ 5, 8 ] ]
\end{Verbatim}
 

 (See Section{\nobreakspace}\ref{Basic Actions} for information about specific actions.) }

 
\subsection{\textcolor{Chapter }{MaximalBlocks}}\logpage{[ 41, 11, 2 ]}
\hyperdef{L}{X79936EB97AAD1144}{}
{
\noindent\textcolor{FuncColor}{$\triangleright$\enspace\texttt{MaximalBlocks({\mdseries\slshape G, Omega[, seed][, gens, acts][, act]})\index{MaximalBlocks@\texttt{MaximalBlocks}!for a group, an action domain, etc.}
\label{MaximalBlocks:for a group, an action domain, etc.}
}\hfill{\scriptsize (operation)}}\\
\noindent\textcolor{FuncColor}{$\triangleright$\enspace\texttt{MaximalBlocks({\mdseries\slshape xset[, seed]})\index{MaximalBlocks@\texttt{MaximalBlocks}!for an external set}
\label{MaximalBlocks:for an external set}
}\hfill{\scriptsize (attribute)}}\\


 returns a block system that is maximal (i.e., blocks are maximal with respect
to inclusion) for the transitive (see \texttt{IsTransitive} (\ref{IsTransitive:for a group, an action domain, etc.})) action of \mbox{\texttt{\mdseries\slshape G}} on \mbox{\texttt{\mdseries\slshape Omega}}. If \mbox{\texttt{\mdseries\slshape seed}} is given, a block system is computed in which \mbox{\texttt{\mdseries\slshape seed}} is a subset of one block. 

 The result is undefined if the action is not transitive. 
\begin{Verbatim}[commandchars=!@|,fontsize=\small,frame=single,label=Example]
  !gapprompt@gap>| !gapinput@MaximalBlocks(g,[1..8]);|
  [ [ 1, 2, 3, 8 ], [ 4 .. 7 ] ]
\end{Verbatim}
 }

 
\subsection{\textcolor{Chapter }{RepresentativesMinimalBlocks}}\logpage{[ 41, 11, 3 ]}
\hyperdef{L}{X7941DB6380B74510}{}
{
\noindent\textcolor{FuncColor}{$\triangleright$\enspace\texttt{RepresentativesMinimalBlocks({\mdseries\slshape G, Omega[, gens, acts][, act]})\index{RepresentativesMinimalBlocks@\texttt{RepresentativesMinimalBlocks}!for a group, an action domain, etc.}
\label{RepresentativesMinimalBlocks:for a group, an action domain, etc.}
}\hfill{\scriptsize (operation)}}\\
\noindent\textcolor{FuncColor}{$\triangleright$\enspace\texttt{RepresentativesMinimalBlocks({\mdseries\slshape xset})\index{RepresentativesMinimalBlocks@\texttt{RepresentativesMinimalBlocks}!for an external set}
\label{RepresentativesMinimalBlocks:for an external set}
}\hfill{\scriptsize (attribute)}}\\


 computes a list of block representatives for all minimal (i.e blocks are
minimal with respect to inclusion) nontrivial block systems for the transitive
(see \texttt{IsTransitive} (\ref{IsTransitive:for a group, an action domain, etc.})) action of \mbox{\texttt{\mdseries\slshape G}} on \mbox{\texttt{\mdseries\slshape Omega}}. 

 The result is undefined if the action is not transitive. 
\begin{Verbatim}[commandchars=!@|,fontsize=\small,frame=single,label=Example]
  !gapprompt@gap>| !gapinput@RepresentativesMinimalBlocks(g,[1..8]);|
  [ [ 1, 2 ], [ 1, 3 ], [ 1, 4 ], [ 1, 5 ], [ 1, 6 ], [ 1, 7 ],
    [ 1, 8 ] ]
\end{Verbatim}
 }

 

\subsection{\textcolor{Chapter }{AllBlocks}}
\logpage{[ 41, 11, 4 ]}\nobreak
\hyperdef{L}{X835658B07B28EF3B}{}
{\noindent\textcolor{FuncColor}{$\triangleright$\enspace\texttt{AllBlocks({\mdseries\slshape G})\index{AllBlocks@\texttt{AllBlocks}}
\label{AllBlocks}
}\hfill{\scriptsize (attribute)}}\\


 computes a list of representatives of all block systems for a permutation
group \mbox{\texttt{\mdseries\slshape G}} acting transitively on the points moved by the group. 

 Each representative in the returned list is sorted and contains the smallest
point moved by \mbox{\texttt{\mdseries\slshape G}}. 

 
\begin{Verbatim}[commandchars=!@|,fontsize=\small,frame=single,label=Example]
  !gapprompt@gap>| !gapinput@AllBlocks(g);|
  [ [ 1, 8 ], [ 1, 2, 3, 8 ], [ 1, 4, 5, 8 ], [ 1, 6, 7, 8 ], [ 1, 3 ],
    [ 1, 3, 5, 7 ], [ 1, 3, 4, 6 ], [ 1, 5 ], [ 1, 2, 5, 6 ], [ 1, 2 ],
    [ 1, 2, 4, 7 ], [ 1, 4 ], [ 1, 7 ], [ 1, 6 ] ]
\end{Verbatim}
 

 The stabilizer of a block can be computed via the action \texttt{OnSets} (\ref{OnSets}): 

 
\begin{Verbatim}[commandchars=!@|,fontsize=\small,frame=single,label=Example]
  !gapprompt@gap>| !gapinput@Stabilizer(g,[1,8],OnSets);|
  Group([ (1,8)(2,3)(4,5)(6,7) ])
\end{Verbatim}
 

 If \texttt{bs} is a partition of the action domain, given as a set of sets, the stabilizer
under the action \texttt{OnSetsDisjointSets} (\ref{OnSetsDisjointSets}) returns the largest subgroup which preserves \texttt{bs} as a block system. 

 
\begin{Verbatim}[commandchars=!@|,fontsize=\small,frame=single,label=Example]
  !gapprompt@gap>| !gapinput@g:=Group((1,2,3,4,5,6,7,8),(1,2));;|
  !gapprompt@gap>| !gapinput@bs:=[[1,2,3,4],[5,6,7,8]];;|
  !gapprompt@gap>| !gapinput@Stabilizer(g,bs,OnSetsDisjointSets);|
  Group([ (6,7), (5,6), (5,8), (2,3), (3,4)(5,7), (1,4),
    (1,5,4,8)(2,6,3,7) ])
\end{Verbatim}
 }

 }

  
\section{\textcolor{Chapter }{External Sets}}\label{External Sets}
\logpage{[ 41, 12, 0 ]}
\hyperdef{L}{X7FD3D2D2788709B7}{}
{
  \index{G-sets@$G$\texttt{\symbol{45}}sets} When considering group actions, sometimes the concept of a \emph{$G$\texttt{\symbol{45}}set} is used. This is a set $\Omega$ endowed with an action of $G$. The elements of the $G$\texttt{\symbol{45}}set are the same as those of $\Omega$, however concepts like equality and equivalence of $G$\texttt{\symbol{45}}sets do not only consider the underlying domain $\Omega$ but the group action as well. 

 This concept is implemented in \textsf{GAP} via \emph{external sets}. 

 The constituents of an external set are stored in the attributes \texttt{ActingDomain} (\ref{ActingDomain}), \texttt{FunctionAction} (\ref{FunctionAction}) and \texttt{HomeEnumerator} (\ref{HomeEnumerator}). 

 Most operations for actions are applicable as an attribute for an external
set. 

 The most prominent external subsets are orbits, see \texttt{ExternalOrbit} (\ref{ExternalOrbit}). 

 Many subsets of a group, such as conjugacy classes or cosets
(see{\nobreakspace}\texttt{ConjugacyClass} (\ref{ConjugacyClass}) and \texttt{RightCoset} (\ref{RightCoset})) are implemented as external orbits. 

 External sets also are implicitly underlying action homomorphisms, see \texttt{UnderlyingExternalSet} (\ref{UnderlyingExternalSet}) and \texttt{SurjectiveActionHomomorphismAttr} (\ref{SurjectiveActionHomomorphismAttr}). 

\subsection{\textcolor{Chapter }{IsExternalSet}}
\logpage{[ 41, 12, 1 ]}\nobreak
\hyperdef{L}{X8264C3C479FF0A8B}{}
{\noindent\textcolor{FuncColor}{$\triangleright$\enspace\texttt{IsExternalSet({\mdseries\slshape obj})\index{IsExternalSet@\texttt{IsExternalSet}}
\label{IsExternalSet}
}\hfill{\scriptsize (Category)}}\\


 An \emph{external set} specifies a group action $\mu: \Omega \times G \mapsto \Omega$ of a group $G$ on a domain $\Omega$. The external set knows the group, the domain and the actual acting function.
Mathematically, an external set is the set{\nobreakspace}$\Omega$, which is endowed with the action of a group $G$ via the group action $\mu$. For this reason \textsf{GAP} treats an external set as a domain whose elements are the elements of $\Omega$. An external set is always a union of orbits. Currently the
domain{\nobreakspace}$\Omega$ must always be finite. If $\Omega$ is not a list, an enumerator for $\Omega$ is automatically chosen, see \texttt{Enumerator} (\ref{Enumerator}). }

 

\subsection{\textcolor{Chapter }{ExternalSet}}
\logpage{[ 41, 12, 2 ]}\nobreak
\hyperdef{L}{X7C90F648793E47DD}{}
{\noindent\textcolor{FuncColor}{$\triangleright$\enspace\texttt{ExternalSet({\mdseries\slshape G, Omega[, gens, acts][, act]})\index{ExternalSet@\texttt{ExternalSet}}
\label{ExternalSet}
}\hfill{\scriptsize (operation)}}\\


 creates the external set for the action \mbox{\texttt{\mdseries\slshape act}} of \mbox{\texttt{\mdseries\slshape G}} on \mbox{\texttt{\mdseries\slshape Omega}}. \mbox{\texttt{\mdseries\slshape Omega}} can be either a proper set, or a domain which is represented as described in \ref{Domains} and \ref{Collections}, or (to use less memory but with a slower performance) an enumerator (see \texttt{Enumerator} (\ref{Enumerator}) ) of this domain. 

 The result is undefined if \mbox{\texttt{\mdseries\slshape G}} does not act on \mbox{\texttt{\mdseries\slshape Omega}}. 
\begin{Verbatim}[commandchars=!@|,fontsize=\small,frame=single,label=Example]
  !gapprompt@gap>| !gapinput@g:=Group((1,2,3),(2,3,4));;|
  !gapprompt@gap>| !gapinput@e:=ExternalSet(g,[1..4]);|
  <xset:[ 1, 2, 3, 4 ]>
  !gapprompt@gap>| !gapinput@e:=ExternalSet(g,g,OnRight);|
  <xset:[ (), (2,3,4), (2,4,3), (1,2)(3,4), (1,2,3), (1,2,4), (1,3,2),
    (1,3,4), (1,3)(2,4), (1,4,2), (1,4,3), (1,4)(2,3) ]>
  !gapprompt@gap>| !gapinput@Orbits(e);|
  [ [ (), (1,2)(3,4), (1,3)(2,4), (1,4)(2,3), (2,4,3), (1,4,2),
        (1,2,3), (1,3,4), (2,3,4), (1,3,2), (1,4,3), (1,2,4) ] ]
\end{Verbatim}
 }

 

\subsection{\textcolor{Chapter }{ActingDomain}}
\logpage{[ 41, 12, 3 ]}\nobreak
\hyperdef{L}{X7B9DB15D80CE28B4}{}
{\noindent\textcolor{FuncColor}{$\triangleright$\enspace\texttt{ActingDomain({\mdseries\slshape xset})\index{ActingDomain@\texttt{ActingDomain}}
\label{ActingDomain}
}\hfill{\scriptsize (attribute)}}\\


 This attribute returns the group with which the external set \mbox{\texttt{\mdseries\slshape xset}} was defined. }

 

\subsection{\textcolor{Chapter }{FunctionAction}}
\logpage{[ 41, 12, 4 ]}\nobreak
\hyperdef{L}{X86153CB087394DC1}{}
{\noindent\textcolor{FuncColor}{$\triangleright$\enspace\texttt{FunctionAction({\mdseries\slshape xset})\index{FunctionAction@\texttt{FunctionAction}}
\label{FunctionAction}
}\hfill{\scriptsize (attribute)}}\\


 is the acting function with which the external set \mbox{\texttt{\mdseries\slshape xset}} was defined. }

 

\subsection{\textcolor{Chapter }{HomeEnumerator}}
\logpage{[ 41, 12, 5 ]}\nobreak
\hyperdef{L}{X86A0CC1479A5932A}{}
{\noindent\textcolor{FuncColor}{$\triangleright$\enspace\texttt{HomeEnumerator({\mdseries\slshape xset})\index{HomeEnumerator@\texttt{HomeEnumerator}}
\label{HomeEnumerator}
}\hfill{\scriptsize (attribute)}}\\


 returns an enumerator of the action domain with which the external set \mbox{\texttt{\mdseries\slshape xset}} was defined. For external subsets, this is in general different from the \texttt{Enumerator} (\ref{Enumerator}) value of \mbox{\texttt{\mdseries\slshape xset}}, which enumerates only the subset. 
\begin{Verbatim}[commandchars=!@|,fontsize=\small,frame=single,label=Example]
  !gapprompt@gap>| !gapinput@ActingDomain(e);|
  Group([ (1,2,3), (2,3,4) ])
  !gapprompt@gap>| !gapinput@FunctionAction(e)=OnRight;|
  true
  !gapprompt@gap>| !gapinput@HomeEnumerator(e);|
  [ (), (2,3,4), (2,4,3), (1,2)(3,4), (1,2,3), (1,2,4), (1,3,2),
    (1,3,4), (1,3)(2,4), (1,4,2), (1,4,3), (1,4)(2,3) ]
\end{Verbatim}
 }

 

\subsection{\textcolor{Chapter }{IsExternalSubset}}
\logpage{[ 41, 12, 6 ]}\nobreak
\hyperdef{L}{X879DE63C7858453C}{}
{\noindent\textcolor{FuncColor}{$\triangleright$\enspace\texttt{IsExternalSubset({\mdseries\slshape obj})\index{IsExternalSubset@\texttt{IsExternalSubset}}
\label{IsExternalSubset}
}\hfill{\scriptsize (Representation)}}\\


 An external subset is the restriction of an external set to a subset of the
domain (which must be invariant under the action). It is again an external
set. }

 

\subsection{\textcolor{Chapter }{ExternalSubset}}
\logpage{[ 41, 12, 7 ]}\nobreak
\hyperdef{L}{X87D1EA1486D86233}{}
{\noindent\textcolor{FuncColor}{$\triangleright$\enspace\texttt{ExternalSubset({\mdseries\slshape G, Omega, start[, gens, acts], act})\index{ExternalSubset@\texttt{ExternalSubset}}
\label{ExternalSubset}
}\hfill{\scriptsize (operation)}}\\


 constructs the external subset of \mbox{\texttt{\mdseries\slshape Omega}} on the union of orbits of the points in \mbox{\texttt{\mdseries\slshape start}}. 

 The result is undefined if \mbox{\texttt{\mdseries\slshape G}} does not act on \mbox{\texttt{\mdseries\slshape Omega}}. }

 

\subsection{\textcolor{Chapter }{IsExternalOrbit}}
\logpage{[ 41, 12, 8 ]}\nobreak
\hyperdef{L}{X7E081F568407317F}{}
{\noindent\textcolor{FuncColor}{$\triangleright$\enspace\texttt{IsExternalOrbit({\mdseries\slshape obj})\index{IsExternalOrbit@\texttt{IsExternalOrbit}}
\label{IsExternalOrbit}
}\hfill{\scriptsize (Representation)}}\\


 An external orbit is an external subset consisting of one orbit. }

 

\subsection{\textcolor{Chapter }{ExternalOrbit}}
\logpage{[ 41, 12, 9 ]}\nobreak
\hyperdef{L}{X7FB656AE7A066C35}{}
{\noindent\textcolor{FuncColor}{$\triangleright$\enspace\texttt{ExternalOrbit({\mdseries\slshape G, Omega, pnt[, gens, acts], act})\index{ExternalOrbit@\texttt{ExternalOrbit}}
\label{ExternalOrbit}
}\hfill{\scriptsize (operation)}}\\


 constructs the external subset on the orbit of \mbox{\texttt{\mdseries\slshape pnt}}. The \texttt{Representative} (\ref{Representative}) value of this external set is \mbox{\texttt{\mdseries\slshape pnt}}. 

 The result is undefined if \mbox{\texttt{\mdseries\slshape G}} does not act on \mbox{\texttt{\mdseries\slshape Omega}}. 
\begin{Verbatim}[commandchars=!@|,fontsize=\small,frame=single,label=Example]
  !gapprompt@gap>| !gapinput@e:=ExternalOrbit(g,g,(1,2,3));|
  (1,2,3)^G
\end{Verbatim}
 }

 

\subsection{\textcolor{Chapter }{StabilizerOfExternalSet}}
\logpage{[ 41, 12, 10 ]}\nobreak
\hyperdef{L}{X7BAFF02B7D6DF9F2}{}
{\noindent\textcolor{FuncColor}{$\triangleright$\enspace\texttt{StabilizerOfExternalSet({\mdseries\slshape xset})\index{StabilizerOfExternalSet@\texttt{StabilizerOfExternalSet}}
\label{StabilizerOfExternalSet}
}\hfill{\scriptsize (attribute)}}\\


 computes the stabilizer of the \texttt{Representative} (\ref{Representative}) value of the external set \mbox{\texttt{\mdseries\slshape xset}}. The stabilizer will have the acting group of \mbox{\texttt{\mdseries\slshape xset}} as its parent. 
\begin{Verbatim}[commandchars=!@|,fontsize=\small,frame=single,label=Example]
  !gapprompt@gap>| !gapinput@Representative(e);|
  (1,2,3)
  !gapprompt@gap>| !gapinput@StabilizerOfExternalSet(e);|
  Group([ (1,2,3) ])
\end{Verbatim}
 }

 
\subsection{\textcolor{Chapter }{ExternalOrbits}}\logpage{[ 41, 12, 11 ]}
\hyperdef{L}{X867262FA82FDD592}{}
{
\noindent\textcolor{FuncColor}{$\triangleright$\enspace\texttt{ExternalOrbits({\mdseries\slshape G, Omega[, gens, acts][, act]})\index{ExternalOrbits@\texttt{ExternalOrbits}!for a group, an action domain, etc.}
\label{ExternalOrbits:for a group, an action domain, etc.}
}\hfill{\scriptsize (operation)}}\\
\noindent\textcolor{FuncColor}{$\triangleright$\enspace\texttt{ExternalOrbits({\mdseries\slshape xset})\index{ExternalOrbits@\texttt{ExternalOrbits}!for an external set}
\label{ExternalOrbits:for an external set}
}\hfill{\scriptsize (attribute)}}\\


 computes a list of external orbits that give the orbits of \mbox{\texttt{\mdseries\slshape G}}. 
\begin{Verbatim}[commandchars=!@|,fontsize=\small,frame=single,label=Example]
  !gapprompt@gap>| !gapinput@ExternalOrbits(g,AsList(g));|
  [ ()^G, (2,3,4)^G, (2,4,3)^G, (1,2)(3,4)^G ]
\end{Verbatim}
 }

 
\subsection{\textcolor{Chapter }{ExternalOrbitsStabilizers}}\logpage{[ 41, 12, 12 ]}
\hyperdef{L}{X7A64EF807CE8893E}{}
{
\noindent\textcolor{FuncColor}{$\triangleright$\enspace\texttt{ExternalOrbitsStabilizers({\mdseries\slshape G, Omega[, gens, acts][, act]})\index{ExternalOrbitsStabilizers@\texttt{ExternalOrbitsStabilizers}!for a group, an action domain, etc.}
\label{ExternalOrbitsStabilizers:for a group, an action domain, etc.}
}\hfill{\scriptsize (operation)}}\\
\noindent\textcolor{FuncColor}{$\triangleright$\enspace\texttt{ExternalOrbitsStabilizers({\mdseries\slshape xset})\index{ExternalOrbitsStabilizers@\texttt{ExternalOrbitsStabilizers}!for an external set}
\label{ExternalOrbitsStabilizers:for an external set}
}\hfill{\scriptsize (attribute)}}\\


 In addition to \texttt{ExternalOrbits} (\ref{ExternalOrbits:for a group, an action domain, etc.}), this operation also computes the stabilizers of the representatives of the
external orbits at the same time. (This can be quicker than computing the \texttt{ExternalOrbits} (\ref{ExternalOrbits:for a group, an action domain, etc.}) value first and the stabilizers afterwards.) 
\begin{Verbatim}[commandchars=!@|,fontsize=\small,frame=single,label=Example]
  !gapprompt@gap>| !gapinput@e:=ExternalOrbitsStabilizers(g,AsList(g));|
  [ ()^G, (2,3,4)^G, (2,4,3)^G, (1,2)(3,4)^G ]
  !gapprompt@gap>| !gapinput@HasStabilizerOfExternalSet(e[3]);|
  true
  !gapprompt@gap>| !gapinput@StabilizerOfExternalSet(e[3]);|
  Group([ (2,4,3) ])
\end{Verbatim}
 }

 

\subsection{\textcolor{Chapter }{CanonicalRepresentativeOfExternalSet}}
\logpage{[ 41, 12, 13 ]}\nobreak
\hyperdef{L}{X8048AE727A7F1A2F}{}
{\noindent\textcolor{FuncColor}{$\triangleright$\enspace\texttt{CanonicalRepresentativeOfExternalSet({\mdseries\slshape xset})\index{CanonicalRepresentativeOfExternalSet@\texttt{Canonical}\-\texttt{Representative}\-\texttt{Of}\-\texttt{ExternalSet}}
\label{CanonicalRepresentativeOfExternalSet}
}\hfill{\scriptsize (attribute)}}\\


 The canonical representative of an external set \mbox{\texttt{\mdseries\slshape xset}} may only depend on the defining attributes \mbox{\texttt{\mdseries\slshape G}}, \mbox{\texttt{\mdseries\slshape Omega}}, \mbox{\texttt{\mdseries\slshape act}} of \mbox{\texttt{\mdseries\slshape xset}} and (in the case of external subsets) \texttt{Enumerator( \mbox{\texttt{\mdseries\slshape xset}} )}. It must \emph{not} depend, e.g., on the representative of an external orbit. \textsf{GAP} does not know methods for arbitrary external sets to compute a canonical
representative, see \texttt{CanonicalRepresentativeDeterminatorOfExternalSet} (\ref{CanonicalRepresentativeDeterminatorOfExternalSet}). }

 

\subsection{\textcolor{Chapter }{CanonicalRepresentativeDeterminatorOfExternalSet}}
\logpage{[ 41, 12, 14 ]}\nobreak
\hyperdef{L}{X8071A8D784DC8325}{}
{\noindent\textcolor{FuncColor}{$\triangleright$\enspace\texttt{CanonicalRepresentativeDeterminatorOfExternalSet({\mdseries\slshape xset})\index{CanonicalRepresentativeDeterminatorOfExternalSet@\texttt{Canonical}\-\texttt{Representative}\-\texttt{Determinator}\-\texttt{Of}\-\texttt{ExternalSet}}
\label{CanonicalRepresentativeDeterminatorOfExternalSet}
}\hfill{\scriptsize (attribute)}}\\


 returns a function that takes as its arguments the acting group and a point.
This function returns a list of length 1 or 3, the first entry being the
canonical representative and the other entries (if bound) being the stabilizer
of the canonical representative and a conjugating element, respectively. An
external set is only guaranteed to be able to compute a canonical
representative if it has a \texttt{CanonicalRepresentativeDeterminatorOfExternalSet}. }

 

\subsection{\textcolor{Chapter }{ActorOfExternalSet}}
\logpage{[ 41, 12, 15 ]}\nobreak
\hyperdef{L}{X85E9A6A77B8D00B8}{}
{\noindent\textcolor{FuncColor}{$\triangleright$\enspace\texttt{ActorOfExternalSet({\mdseries\slshape xset})\index{ActorOfExternalSet@\texttt{ActorOfExternalSet}}
\label{ActorOfExternalSet}
}\hfill{\scriptsize (attribute)}}\\


 returns an element mapping \texttt{Representative(\mbox{\texttt{\mdseries\slshape xset}})} to \texttt{CanonicalRepresentativeOfExternalSet(\mbox{\texttt{\mdseries\slshape xset}})} under the given action. 
\begin{Verbatim}[commandchars=!@|,fontsize=\small,frame=single,label=Example]
  !gapprompt@gap>| !gapinput@u:=Subgroup(g,[(1,2,3)]);;|
  !gapprompt@gap>| !gapinput@e:=RightCoset(u,(1,2)(3,4));;|
  !gapprompt@gap>| !gapinput@CanonicalRepresentativeOfExternalSet(e);|
  (2,4,3)
  !gapprompt@gap>| !gapinput@ActorOfExternalSet(e);|
  (1,3,2)
  !gapprompt@gap>| !gapinput@FunctionAction(e)((1,2)(3,4),last);|
  (2,4,3)
\end{Verbatim}
 }

 

\subsection{\textcolor{Chapter }{UnderlyingExternalSet}}
\logpage{[ 41, 12, 16 ]}\nobreak
\hyperdef{L}{X8190A8247F29A5C7}{}
{\noindent\textcolor{FuncColor}{$\triangleright$\enspace\texttt{UnderlyingExternalSet({\mdseries\slshape acthom})\index{UnderlyingExternalSet@\texttt{UnderlyingExternalSet}}
\label{UnderlyingExternalSet}
}\hfill{\scriptsize (attribute)}}\\


 The underlying set of an action homomorphism \mbox{\texttt{\mdseries\slshape acthom}} is the external set on which it was defined. 
\begin{Verbatim}[commandchars=!@|,fontsize=\small,frame=single,label=Example]
  !gapprompt@gap>| !gapinput@g:=Group((1,2,3),(1,2));;|
  !gapprompt@gap>| !gapinput@hom:=ActionHomomorphism(g,Arrangements([1..4],3),OnTuples);;|
  !gapprompt@gap>| !gapinput@s:=UnderlyingExternalSet(hom);|
  <xset:[[ 1, 2, 3 ],[ 1, 2, 4 ],[ 1, 3, 2 ],[ 1, 3, 4 ],[ 1, 4, 2 ],
  [ 1, 4, 3 ],[ 2, 1, 3 ],[ 2, 1, 4 ],[ 2, 3, 1 ],[ 2, 3, 4 ],
  [ 2, 4, 1 ],[ 2, 4, 3 ],[ 3, 1, 2 ],[ 3, 1, 4 ],[ 3, 2, 1 ], ...]>
  !gapprompt@gap>| !gapinput@Print(s,"\n");|
  [ [ 1, 2, 3 ], [ 1, 2, 4 ], [ 1, 3, 2 ], [ 1, 3, 4 ], [ 1, 4, 2 ],
    [ 1, 4, 3 ], [ 2, 1, 3 ], [ 2, 1, 4 ], [ 2, 3, 1 ], [ 2, 3, 4 ],
    [ 2, 4, 1 ], [ 2, 4, 3 ], [ 3, 1, 2 ], [ 3, 1, 4 ], [ 3, 2, 1 ],
    [ 3, 2, 4 ], [ 3, 4, 1 ], [ 3, 4, 2 ], [ 4, 1, 2 ], [ 4, 1, 3 ],
    [ 4, 2, 1 ], [ 4, 2, 3 ], [ 4, 3, 1 ], [ 4, 3, 2 ] ]
\end{Verbatim}
 }

 

\subsection{\textcolor{Chapter }{SurjectiveActionHomomorphismAttr}}
\logpage{[ 41, 12, 17 ]}\nobreak
\hyperdef{L}{X7A3D87DE809FBFD4}{}
{\noindent\textcolor{FuncColor}{$\triangleright$\enspace\texttt{SurjectiveActionHomomorphismAttr({\mdseries\slshape xset})\index{SurjectiveActionHomomorphismAttr@\texttt{SurjectiveActionHomomorphismAttr}}
\label{SurjectiveActionHomomorphismAttr}
}\hfill{\scriptsize (attribute)}}\\


 returns an action homomorphism for the external set \mbox{\texttt{\mdseries\slshape xset}} which is surjective. (As the \texttt{Image} (\ref{Image:set of images of the source of a general mapping}) value of this homomorphism has to be computed to obtain the range, this may
take substantially longer than \texttt{ActionHomomorphism} (\ref{ActionHomomorphism:for an external set}).) }

 }

 }

             
\chapter{\textcolor{Chapter }{Permutations}}\label{Permutations}
\logpage{[ 42, 0, 0 ]}
\hyperdef{L}{X80F808307A2D5AB8}{}
{
  \textsf{GAP} offers a data type \emph{permutation} to describe the elements of permutation groups. 

 The points on which permutations in \textsf{GAP} act are the positive integers up to a certain architecture dependent limit,
and the image of a point $i$ under a permutation $p$ is written $i^p$, which is expressed as $i$\texttt{\texttt{\symbol{94}}}$p$ in \textsf{GAP}. (This action is also implemented by the function \texttt{OnPoints} (\ref{OnPoints}).) If $i$\texttt{\texttt{\symbol{94}}}$p$ is different from $i$, we say that $i$ is \emph{moved} by{\nobreakspace}$p$, otherwise it is \emph{fixed}. Permutations in \textsf{GAP} are entered and displayed in cycle notation, such as \texttt{(1,2,3)(4,5)}. 

 The preimage of the point $i$ under the permutation $p$ can be computed as $i$\texttt{/}$p$, see \texttt{PERM{\textunderscore}INVERSE{\textunderscore}THRESHOLD} (\ref{PERMuScoreINVERSEuScoreTHRESHOLD}). 

 For arithmetic operations for permutations and their precedence,
see{\nobreakspace}\ref{Arithmetic Operations for Elements}. 

 In the names of the \textsf{GAP} functions that deal with permutations, the word ``Permutation'' is usually abbreviated to ``Perm'', to save typing. For example, the category test function for permutations is \texttt{IsPerm} (\ref{IsPerm}).  
\section{\textcolor{Chapter }{IsPerm (Filter)}}\label{sect:IsPerm}
\logpage{[ 42, 1, 0 ]}
\hyperdef{L}{X80F07BE2811D4BAC}{}
{
  Internally, \textsf{GAP} stores a permutation as a list of the $d$ images of the integers $1, \ldots, d$, where the ``internal degree'' $d$ is the largest integer moved by the permutation or bigger. When a permutation
is read in cycle notation, $d$ is always set to the largest moved integer, but a bigger $d$ can result from a multiplication of two permutations, because the product is
not shortened if it fixes{\nobreakspace}$d$. The images are stored all as $16$\texttt{\symbol{45}}bit integers or all as $32$\texttt{\symbol{45}}bit integers, depending on whether $d \leq 65536$ or not. For example, if $m\geq 65536$, the permutation $(1, 2, \ldots, m)$ has internal degree $d=m$ and takes $4m$ bytes of memory for storage. But
\texttt{\symbol{45}}\texttt{\symbol{45}}\texttt{\symbol{45}} since the
internal degree is not reduced
\texttt{\symbol{45}}\texttt{\symbol{45}}\texttt{\symbol{45}} this means that
the identity permutation \texttt{()} calculated as $(1, 2, \ldots, m) * (1, 2, \ldots, m)^{{-1}}$ also takes $4m$ bytes of storage. It can take even more because the internal list has
sometimes room for more than $d$ images. 

 On 32\texttt{\symbol{45}}bit systems, the limit on the degree of permutations
is, for technical reasons, $2^{28}-1$. On 64\texttt{\symbol{45}}bit systems, it is $2^{32}-1$ because only a 32\texttt{\symbol{45}}bit integer is used to represent each
image internally. Error messages should be given if any command would require
creating a permutation exceeding this limit. 

 The operation \texttt{RestrictedPerm} (\ref{RestrictedPerm}) reduces the storage degree of its result and therefore can be used to save
memory if intermediate calculations in large degree result in a small degree
result. 

 Permutations do not belong to a specific group. That means that one can work
with permutations without defining a permutation group that contains them. 

 
\begin{Verbatim}[commandchars=!@|,fontsize=\small,frame=single,label=Example]
  !gapprompt@gap>| !gapinput@(1,2,3);|
  (1,2,3)
  !gapprompt@gap>| !gapinput@(1,2,3) * (2,3,4);|
  (1,3)(2,4)
  !gapprompt@gap>| !gapinput@17^(2,5,17,9,8);|
  9
  !gapprompt@gap>| !gapinput@OnPoints(17,(2,5,17,9,8));|
  9
\end{Verbatim}
 

 The operation \texttt{Permuted} (\ref{Permuted}) can be used to permute the entries of a list according to a permutation. 

\subsection{\textcolor{Chapter }{IsPerm}}
\logpage{[ 42, 1, 1 ]}\nobreak
\hyperdef{L}{X7AA69C6686FC49EA}{}
{\noindent\textcolor{FuncColor}{$\triangleright$\enspace\texttt{IsPerm({\mdseries\slshape obj})\index{IsPerm@\texttt{IsPerm}}
\label{IsPerm}
}\hfill{\scriptsize (Category)}}\\


 Each \emph{permutation} in \textsf{GAP} lies in the category \texttt{IsPerm}. Basic operations for permutations are \texttt{LargestMovedPoint} (\ref{LargestMovedPoint:for a permutation}), multiplication of two permutations via \texttt{*}, and exponentiation \texttt{\texttt{\symbol{94}}} with first argument a positive integer $i$ and second argument a permutation $\pi$, the result being the image $i^{\pi}$ of the point $i$ under $\pi$.  }

 

\subsection{\textcolor{Chapter }{IsPermCollection}}
\logpage{[ 42, 1, 2 ]}\nobreak
\hyperdef{L}{X82069E437D2DF9AA}{}
{\noindent\textcolor{FuncColor}{$\triangleright$\enspace\texttt{IsPermCollection({\mdseries\slshape obj})\index{IsPermCollection@\texttt{IsPermCollection}}
\label{IsPermCollection}
}\hfill{\scriptsize (Category)}}\\
\noindent\textcolor{FuncColor}{$\triangleright$\enspace\texttt{IsPermCollColl({\mdseries\slshape obj})\index{IsPermCollColl@\texttt{IsPermCollColl}}
\label{IsPermCollColl}
}\hfill{\scriptsize (Category)}}\\


 are the categories for collections of permutations and collections of
collections of permutations, respectively. }

 

\subsection{\textcolor{Chapter }{PermutationsFamily}}
\logpage{[ 42, 1, 3 ]}\nobreak
\hyperdef{L}{X819628B083B3939B}{}
{\noindent\textcolor{FuncColor}{$\triangleright$\enspace\texttt{PermutationsFamily\index{PermutationsFamily@\texttt{PermutationsFamily}}
\label{PermutationsFamily}
}\hfill{\scriptsize (family)}}\\


 is the family of all permutations. }

 

\subsection{\textcolor{Chapter }{PERM{\textunderscore}INVERSE{\textunderscore}THRESHOLD}}
\logpage{[ 42, 1, 4 ]}\nobreak
\hyperdef{L}{X83C711557DEB4B36}{}
{\noindent\textcolor{FuncColor}{$\triangleright$\enspace\texttt{PERM{\textunderscore}INVERSE{\textunderscore}THRESHOLD\index{PERM{\textunderscore}INVERSE{\textunderscore}THRESHOLD@\texttt{PER}\-\texttt{M{\textunderscore}}\-\texttt{I}\-\texttt{N}\-\texttt{V}\-\texttt{E}\-\texttt{R}\-\texttt{S}\-\texttt{E{\textunderscore}}\-\texttt{T}\-\texttt{H}\-\texttt{R}\-\texttt{E}\-\texttt{S}\-\texttt{HOLD}}
\label{PERMuScoreINVERSEuScoreTHRESHOLD}
}\hfill{\scriptsize (global variable)}}\\


 For permutations of degree up to \texttt{PERM{\textunderscore}INVERSE{\textunderscore}THRESHOLD} whenever the inverse image of a point under a permutations is needed, the
entire inverse is computed and stored. Otherwise, if the inverse is not
stored, the point is traced around the cycle it is part of to find the inverse
image. This takes time when it happens, and uses memory, but saves time on a
variety of subsequent computations. This threshold can be adjusted by simply
assigning to the variable. The default is 10000. }

 }

  
\section{\textcolor{Chapter }{Comparison of Permutations}}\label{Comparison of Permutations}
\logpage{[ 42, 2, 0 ]}
\hyperdef{L}{X7A21DE5779D21A6D}{}
{
  

\subsection{\textcolor{Chapter }{\texttt{\symbol{92}}= (for permutations)}}
\logpage{[ 42, 2, 1 ]}\nobreak
\hyperdef{L}{X7CEC03A9808E2E7C}{}
{\noindent\textcolor{FuncColor}{$\triangleright$\enspace\texttt{\texttt{\symbol{92}}=({\mdseries\slshape p1, p2})\index{\texttt{\symbol{92}}=@\texttt{\texttt{\symbol{92}}=}!for permutations}
\label{bSlash=:for permutations}
}\hfill{\scriptsize (method)}}\\
\noindent\textcolor{FuncColor}{$\triangleright$\enspace\texttt{\texttt{\symbol{92}}{\textless}({\mdseries\slshape p1, p2})\index{\texttt{\symbol{92}}{\textless}@\texttt{\texttt{\symbol{92}}{\textless}}!for permutations}
\label{bSlash<:for permutations}
}\hfill{\scriptsize (method)}}\\


 \index{equality test!for permutations} \index{precedence test!for permutations} Two permutations are equal if they move the same points and all these points
have the same images under both permutations. 

 The permutation \mbox{\texttt{\mdseries\slshape p1}} is smaller than \mbox{\texttt{\mdseries\slshape p2}} if \mbox{\texttt{\mdseries\slshape p1}} $\neq$ \mbox{\texttt{\mdseries\slshape p2}} and $i^{{\mbox{\texttt{\mdseries\slshape p1}}}} < i^{{\mbox{\texttt{\mdseries\slshape p2}}}}$, where $i$ is the smallest point with $i^{{\mbox{\texttt{\mdseries\slshape p1}}}} \neq i^{{\mbox{\texttt{\mdseries\slshape p2}}}}$. Therefore the identity permutation is the smallest permutation, see also
Section{\nobreakspace}\ref{Comparison Operations for Elements}. 

 Permutations can be compared with certain other \textsf{GAP} objects, see{\nobreakspace}\ref{Comparisons} for the details. 

 
\begin{Verbatim}[commandchars=!@|,fontsize=\small,frame=single,label=Example]
  !gapprompt@gap>| !gapinput@(1,2,3) = (2,3,1);|
  true
  !gapprompt@gap>| !gapinput@(1,2,3) * (2,3,4) = (1,3)(2,4);|
  true
  !gapprompt@gap>| !gapinput@(1,2,3) < (1,3,2);      # 1^(1,2,3) = 2 < 3 = 1^(1,3,2)|
  true
  !gapprompt@gap>| !gapinput@(1,3,2,4) < (1,3,4,2);  # 2^(1,3,2,4) = 4 > 1 = 2^(1,3,4,2)|
  false
\end{Verbatim}
 }

 

\subsection{\textcolor{Chapter }{DistancePerms}}
\logpage{[ 42, 2, 2 ]}\nobreak
\hyperdef{L}{X7BC944F57A04AFF2}{}
{\noindent\textcolor{FuncColor}{$\triangleright$\enspace\texttt{DistancePerms({\mdseries\slshape perm1, perm2})\index{DistancePerms@\texttt{DistancePerms}}
\label{DistancePerms}
}\hfill{\scriptsize (operation)}}\\


 returns the number of points for which \mbox{\texttt{\mdseries\slshape perm1}} and \mbox{\texttt{\mdseries\slshape perm2}} have different images. This should always produce the same result as \texttt{NrMovedPoints(\mbox{\texttt{\mdseries\slshape perm1}}/\mbox{\texttt{\mdseries\slshape perm2}})} but some methods may be much faster than this form, since no new permutation
needs to be created. }

 

\subsection{\textcolor{Chapter }{SmallestGeneratorPerm}}
\logpage{[ 42, 2, 3 ]}\nobreak
\hyperdef{L}{X83A917F67D45C7EA}{}
{\noindent\textcolor{FuncColor}{$\triangleright$\enspace\texttt{SmallestGeneratorPerm({\mdseries\slshape perm})\index{SmallestGeneratorPerm@\texttt{SmallestGeneratorPerm}}
\label{SmallestGeneratorPerm}
}\hfill{\scriptsize (attribute)}}\\


 is the smallest permutation that generates the same cyclic group as the
permutation \mbox{\texttt{\mdseries\slshape perm}}. This is very efficient, even when \mbox{\texttt{\mdseries\slshape perm}} has large order. 
\begin{Verbatim}[commandchars=!@|,fontsize=\small,frame=single,label=Example]
  !gapprompt@gap>| !gapinput@SmallestGeneratorPerm( (1,4,3,2) );|
  (1,2,3,4)
\end{Verbatim}
 }

 }

  
\section{\textcolor{Chapter }{Moved Points of Permutations}}\label{Moved Points of Permutations}
\logpage{[ 42, 3, 0 ]}
\hyperdef{L}{X82C255E2821C0721}{}
{
  

\subsection{\textcolor{Chapter }{SmallestMovedPoint (for a permutation)}}
\logpage{[ 42, 3, 1 ]}\nobreak
\hyperdef{L}{X84EF0A697F7A87DC}{}
{\noindent\textcolor{FuncColor}{$\triangleright$\enspace\texttt{SmallestMovedPoint({\mdseries\slshape perm})\index{SmallestMovedPoint@\texttt{SmallestMovedPoint}!for a permutation}
\label{SmallestMovedPoint:for a permutation}
}\hfill{\scriptsize (attribute)}}\\
\noindent\textcolor{FuncColor}{$\triangleright$\enspace\texttt{SmallestMovedPoint({\mdseries\slshape C})\index{SmallestMovedPoint@\texttt{SmallestMovedPoint}!for a list or collection of permutations}
\label{SmallestMovedPoint:for a list or collection of permutations}
}\hfill{\scriptsize (attribute)}}\\


 is the smallest positive integer that is moved by \mbox{\texttt{\mdseries\slshape perm}} if such an integer exists, and \texttt{infinity} (\ref{infinity}) if \mbox{\texttt{\mdseries\slshape perm}} is the identity. For \mbox{\texttt{\mdseries\slshape C}} a collection or list of permutations, the smallest value of \texttt{SmallestMovedPoint} for the elements of \mbox{\texttt{\mdseries\slshape C}} is returned (and \texttt{infinity} (\ref{infinity}) if \mbox{\texttt{\mdseries\slshape C}} is empty). }

 

\subsection{\textcolor{Chapter }{LargestMovedPoint (for a permutation)}}
\logpage{[ 42, 3, 2 ]}\nobreak
\hyperdef{L}{X84AA603987C94AC0}{}
{\noindent\textcolor{FuncColor}{$\triangleright$\enspace\texttt{LargestMovedPoint({\mdseries\slshape perm})\index{LargestMovedPoint@\texttt{LargestMovedPoint}!for a permutation}
\label{LargestMovedPoint:for a permutation}
}\hfill{\scriptsize (attribute)}}\\
\noindent\textcolor{FuncColor}{$\triangleright$\enspace\texttt{LargestMovedPoint({\mdseries\slshape C})\index{LargestMovedPoint@\texttt{LargestMovedPoint}!for a list or collection of permutations}
\label{LargestMovedPoint:for a list or collection of permutations}
}\hfill{\scriptsize (attribute)}}\\


 For a permutation \mbox{\texttt{\mdseries\slshape perm}}, this attribute contains the largest positive integer which is moved by \mbox{\texttt{\mdseries\slshape perm}} if such an integer exists, and \texttt{0} if \mbox{\texttt{\mdseries\slshape perm}} is the identity. For \mbox{\texttt{\mdseries\slshape C}} a collection or list of permutations, the largest value of \texttt{LargestMovedPoint} for the elements of \mbox{\texttt{\mdseries\slshape C}} is returned (and \texttt{0} if \mbox{\texttt{\mdseries\slshape C}} is empty). }

 

\subsection{\textcolor{Chapter }{MovedPoints (for a permutation)}}
\logpage{[ 42, 3, 3 ]}\nobreak
\hyperdef{L}{X85E61B9C7A6B0CCA}{}
{\noindent\textcolor{FuncColor}{$\triangleright$\enspace\texttt{MovedPoints({\mdseries\slshape perm})\index{MovedPoints@\texttt{MovedPoints}!for a permutation}
\label{MovedPoints:for a permutation}
}\hfill{\scriptsize (attribute)}}\\
\noindent\textcolor{FuncColor}{$\triangleright$\enspace\texttt{MovedPoints({\mdseries\slshape C})\index{MovedPoints@\texttt{MovedPoints}!for a list or collection of permutations}
\label{MovedPoints:for a list or collection of permutations}
}\hfill{\scriptsize (attribute)}}\\


 is the proper set of the positive integers moved by at least one permutation
in the collection \mbox{\texttt{\mdseries\slshape C}}, respectively by the permutation \mbox{\texttt{\mdseries\slshape perm}}. }

 

\subsection{\textcolor{Chapter }{NrMovedPoints (for a permutation)}}
\logpage{[ 42, 3, 4 ]}\nobreak
\hyperdef{L}{X85E7B1E28430F49E}{}
{\noindent\textcolor{FuncColor}{$\triangleright$\enspace\texttt{NrMovedPoints({\mdseries\slshape perm})\index{NrMovedPoints@\texttt{NrMovedPoints}!for a permutation}
\label{NrMovedPoints:for a permutation}
}\hfill{\scriptsize (attribute)}}\\
\noindent\textcolor{FuncColor}{$\triangleright$\enspace\texttt{NrMovedPoints({\mdseries\slshape C})\index{NrMovedPoints@\texttt{NrMovedPoints}!for a list or collection of permutations}
\label{NrMovedPoints:for a list or collection of permutations}
}\hfill{\scriptsize (attribute)}}\\


 is the number of positive integers that are moved by \mbox{\texttt{\mdseries\slshape perm}}, respectively by at least one element in the collection \mbox{\texttt{\mdseries\slshape C}}. (The actual moved points are returned by \texttt{MovedPoints} (\ref{MovedPoints:for a permutation}).) 
\begin{Verbatim}[commandchars=!@|,fontsize=\small,frame=single,label=Example]
  !gapprompt@gap>| !gapinput@SmallestMovedPointPerm((4,5,6)(7,2,8));|
  2
  !gapprompt@gap>| !gapinput@LargestMovedPointPerm((4,5,6)(7,2,8));|
  8
  !gapprompt@gap>| !gapinput@NrMovedPointsPerm((4,5,6)(7,2,8));|
  6
  !gapprompt@gap>| !gapinput@MovedPoints([(2,3,4),(7,6,3),(5,47)]);|
  [ 2, 3, 4, 5, 6, 7, 47 ]
  !gapprompt@gap>| !gapinput@NrMovedPoints([(2,3,4),(7,6,3),(5,47)]);|
  7
  !gapprompt@gap>| !gapinput@SmallestMovedPoint([(2,3,4),(7,6,3),(5,47)]);|
  2
  !gapprompt@gap>| !gapinput@LargestMovedPoint([(2,3,4),(7,6,3),(5,47)]);|
  47
  !gapprompt@gap>| !gapinput@LargestMovedPoint([()]);|
  0
\end{Verbatim}
 }

 }

  
\section{\textcolor{Chapter }{Sign and Cycle Structure}}\label{Sign and Cycle Structure}
\logpage{[ 42, 4, 0 ]}
\hyperdef{L}{X79BE80267F4AA2B0}{}
{
  

\subsection{\textcolor{Chapter }{SignPerm}}
\logpage{[ 42, 4, 1 ]}\nobreak
\hyperdef{L}{X7BE5011B7C0DB704}{}
{\noindent\textcolor{FuncColor}{$\triangleright$\enspace\texttt{SignPerm({\mdseries\slshape perm})\index{SignPerm@\texttt{SignPerm}}
\label{SignPerm}
}\hfill{\scriptsize (attribute)}}\\


 The \emph{sign} of a permutation \mbox{\texttt{\mdseries\slshape perm}} is defined as $(-1)^k$ where $k$ is the number of cycles of \mbox{\texttt{\mdseries\slshape perm}} of even length. 

 The sign is a homomorphism from the symmetric group onto the multiplicative
group $\{ +1, -1 \}$, the kernel of which is the alternating group. }

 

\subsection{\textcolor{Chapter }{CycleStructurePerm}}
\logpage{[ 42, 4, 2 ]}\nobreak
\hyperdef{L}{X7944D1447804A69A}{}
{\noindent\textcolor{FuncColor}{$\triangleright$\enspace\texttt{CycleStructurePerm({\mdseries\slshape perm})\index{CycleStructurePerm@\texttt{CycleStructurePerm}}
\label{CycleStructurePerm}
}\hfill{\scriptsize (attribute)}}\\


 is the cycle structure (i.e. the numbers of cycles of different lengths) of
the permutation \mbox{\texttt{\mdseries\slshape perm}}. This is encoded in a list $l$ in the following form: The $i$\texttt{\symbol{45}}th entry of $l$ contains the number of cycles of \mbox{\texttt{\mdseries\slshape perm}} of length $i+1$. If \mbox{\texttt{\mdseries\slshape perm}} contains no cycles of length $i+1$ it is not bound. Cycles of length 1 are ignored. 
\begin{Verbatim}[commandchars=!@|,fontsize=\small,frame=single,label=Example]
  !gapprompt@gap>| !gapinput@SignPerm((1,2,3)(4,5));|
  -1
  !gapprompt@gap>| !gapinput@CycleStructurePerm((1,2,3)(4,5,9,7,8));|
  [ , 1,, 1 ]
\end{Verbatim}
 }

 }

  
\section{\textcolor{Chapter }{Creating Permutations}}\label{Creating Permutations}
\logpage{[ 42, 5, 0 ]}
\hyperdef{L}{X7B3194EC869D936D}{}
{
  

\subsection{\textcolor{Chapter }{ListPerm}}
\logpage{[ 42, 5, 1 ]}\nobreak
\hyperdef{L}{X7A9DCFD986958C1E}{}
{\noindent\textcolor{FuncColor}{$\triangleright$\enspace\texttt{ListPerm({\mdseries\slshape perm[, n]})\index{ListPerm@\texttt{ListPerm}}
\label{ListPerm}
}\hfill{\scriptsize (function)}}\\


 is a list $l$ that contains the images of the positive integers from 1 to \mbox{\texttt{\mdseries\slshape n}} under the permutation \mbox{\texttt{\mdseries\slshape perm}}. That means that $l$\texttt{[}$i$\texttt{]} $= i$\texttt{\texttt{\symbol{94}}}\mbox{\texttt{\mdseries\slshape perm}}, where $i$ lies between 1 and \mbox{\texttt{\mdseries\slshape n}}. 

 If the optional second argument \mbox{\texttt{\mdseries\slshape n}} is omitted then the largest point moved by \mbox{\texttt{\mdseries\slshape perm}} is used (see{\nobreakspace}\texttt{LargestMovedPoint} (\ref{LargestMovedPoint:for a permutation})). }

 

\subsection{\textcolor{Chapter }{PermList}}
\logpage{[ 42, 5, 2 ]}\nobreak
\hyperdef{L}{X78D611D17EA6E3BC}{}
{\noindent\textcolor{FuncColor}{$\triangleright$\enspace\texttt{PermList({\mdseries\slshape list})\index{PermList@\texttt{PermList}}
\label{PermList}
}\hfill{\scriptsize (function)}}\\


 is the permutation $\pi$ that moves points as described by the list \mbox{\texttt{\mdseries\slshape list}}. That means that $i^{\pi} =$ \mbox{\texttt{\mdseries\slshape list}}\texttt{[}$i$\texttt{]} if $i$ lies between $1$ and the length of \mbox{\texttt{\mdseries\slshape list}}, and $i^{\pi} = i$ if $i$ is larger than the length of the list \mbox{\texttt{\mdseries\slshape list}}. \texttt{PermList} will return \texttt{fail} if \mbox{\texttt{\mdseries\slshape list}} does not define a permutation, i.e., if \mbox{\texttt{\mdseries\slshape list}} is not dense, or if \mbox{\texttt{\mdseries\slshape list}} contains a positive integer twice, or if \mbox{\texttt{\mdseries\slshape list}} contains an integer not in the range \texttt{[ 1 .. Length( \mbox{\texttt{\mdseries\slshape list}} ) ]}, or if \mbox{\texttt{\mdseries\slshape list}} contains non\texttt{\symbol{45}}integer entries, etc. }

 

\subsection{\textcolor{Chapter }{MappingPermListList}}
\logpage{[ 42, 5, 3 ]}\nobreak
\hyperdef{L}{X8087DCC780B9656A}{}
{\noindent\textcolor{FuncColor}{$\triangleright$\enspace\texttt{MappingPermListList({\mdseries\slshape src, dst})\index{MappingPermListList@\texttt{MappingPermListList}}
\label{MappingPermListList}
}\hfill{\scriptsize (function)}}\\


 Let \mbox{\texttt{\mdseries\slshape src}} and \mbox{\texttt{\mdseries\slshape dst}} be lists of positive integers of the same length, such that there is a
permutation $\pi$ such that \texttt{OnTuples(\mbox{\texttt{\mdseries\slshape src}},} $\pi$\texttt{) = \mbox{\texttt{\mdseries\slshape dst}}}. \texttt{MappingPermListList} returns such a permutation (i. e., \mbox{\texttt{\mdseries\slshape src}}\texttt{[}$i$\texttt{]\texttt{\symbol{94}}}$\pi =$ \mbox{\texttt{\mdseries\slshape dst}}\texttt{[}$i$\texttt{]} holds), with the property that $\pi$ fixes any point which is not in \mbox{\texttt{\mdseries\slshape src}} or \mbox{\texttt{\mdseries\slshape dst}}. If there are several such permutations, it is not specified which of them \texttt{MappingPermListList} returns. If there is no such permutation, then \texttt{MappingPermListList} returns \texttt{fail}. }

 

\subsection{\textcolor{Chapter }{RestrictedPerm}}
\logpage{[ 42, 5, 4 ]}\nobreak
\hyperdef{L}{X7EF8388E7DA8E600}{}
{\noindent\textcolor{FuncColor}{$\triangleright$\enspace\texttt{RestrictedPerm({\mdseries\slshape perm, list})\index{RestrictedPerm@\texttt{RestrictedPerm}}
\label{RestrictedPerm}
}\hfill{\scriptsize (operation)}}\\
\noindent\textcolor{FuncColor}{$\triangleright$\enspace\texttt{RestrictedPermNC({\mdseries\slshape perm, list})\index{RestrictedPermNC@\texttt{RestrictedPermNC}}
\label{RestrictedPermNC}
}\hfill{\scriptsize (operation)}}\\


 \texttt{RestrictedPerm} returns the new permutation that acts on the points in the list \mbox{\texttt{\mdseries\slshape list}} in the same way as the permutation \mbox{\texttt{\mdseries\slshape perm}}, and that fixes those points that are not in \mbox{\texttt{\mdseries\slshape list}}. The resulting permutation is stored internally of degree given by the
maximal entry of \mbox{\texttt{\mdseries\slshape list}}. \mbox{\texttt{\mdseries\slshape list}} must be a list of positive integers such that for each $i$ in \mbox{\texttt{\mdseries\slshape list}} the image $i$\texttt{\texttt{\symbol{94}}}\mbox{\texttt{\mdseries\slshape perm}} is also in \mbox{\texttt{\mdseries\slshape list}}, i.e., \mbox{\texttt{\mdseries\slshape list}} must be the union of cycles of \mbox{\texttt{\mdseries\slshape perm}}. 

 \texttt{RestrictedPermNC} does not check whether \mbox{\texttt{\mdseries\slshape list}} is a union of cycles. 

 
\begin{Verbatim}[commandchars=!@|,fontsize=\small,frame=single,label=Example]
  !gapprompt@gap>| !gapinput@ListPerm((3,4,5));|
  [ 1, 2, 4, 5, 3 ]
  !gapprompt@gap>| !gapinput@PermList([1,2,4,5,3]);|
  (3,4,5)
  !gapprompt@gap>| !gapinput@MappingPermListList([2,5,1,6],[7,12,8,2]);|
  (1,8,5,12,6,2,7)
  !gapprompt@gap>| !gapinput@RestrictedPerm((1,2)(3,4),[3..5]);|
  (3,4)
\end{Verbatim}
 }

 

\subsection{\textcolor{Chapter }{CycleFromList}}
\logpage{[ 42, 5, 5 ]}\nobreak
\hyperdef{L}{X80665A5D800CAFE1}{}
{\noindent\textcolor{FuncColor}{$\triangleright$\enspace\texttt{CycleFromList({\mdseries\slshape list})\index{CycleFromList@\texttt{CycleFromList}}
\label{CycleFromList}
}\hfill{\scriptsize (function)}}\\


 For the given dense, duplicate\texttt{\symbol{45}}free list of positive
integers $[a_1, a_2, ..., a_n]$ return the $n$\texttt{\symbol{45}}cycle $(a_1,a_2,...,a_n)$. For the empty list the trivial permutation $()$ is returned. 

 If the given \mbox{\texttt{\mdseries\slshape list}} contains duplicates or holes, return \texttt{fail}. 

 
\begin{Verbatim}[commandchars=!@|,fontsize=\small,frame=single,label=Example]
  !gapprompt@gap>| !gapinput@CycleFromList( [1,2,3,4] );|
  (1,2,3,4)
  !gapprompt@gap>| !gapinput@CycleFromList( [3,2,6,4,5] );|
  (2,6,4,5,3)
  !gapprompt@gap>| !gapinput@CycleFromList( [2,3,2] );|
  fail
  !gapprompt@gap>| !gapinput@CycleFromList( [1,,3] );|
  fail
\end{Verbatim}
 }

 

\subsection{\textcolor{Chapter }{AsPermutation}}
\logpage{[ 42, 5, 6 ]}\nobreak
\hyperdef{L}{X8353AB8987E35DF3}{}
{\noindent\textcolor{FuncColor}{$\triangleright$\enspace\texttt{AsPermutation({\mdseries\slshape f})\index{AsPermutation@\texttt{AsPermutation}}
\label{AsPermutation}
}\hfill{\scriptsize (attribute)}}\\
\textbf{\indent Returns:\ }
A permutation or \texttt{fail}.



 Partial permutations and transformations which define permutations
(mathematically) can be converted into \textsf{GAP} permutations using \texttt{AsPermutation}; see Chapters \ref{Transformations} and \ref{Partial permutations} for more details about transformations and partial permutations. 
\begin{description}
\item[{for partial permutations}]  If the partial permutation \mbox{\texttt{\mdseries\slshape f}} is a permutation of its image, then \texttt{AsPermutation} returns this permutation. If \mbox{\texttt{\mdseries\slshape f}} does not permute its image, then \texttt{fail} is returned. 

 
\item[{for transformations}]  A transformation is a permutation if and only if its rank equals its degree.
If a transformation in \textsf{GAP} is a permutation, then \texttt{AsPermutation} returns this permutation. If \mbox{\texttt{\mdseries\slshape f}} is not a permutation, then \texttt{fail} is returned. 
\end{description}
 The function \texttt{Permutation} (\ref{Permutation:for a group, an action domain, etc.}) can also be used to convert partial permutations and transformations into
permutations where appropriate. 
\begin{Verbatim}[commandchars=!@|,fontsize=\small,frame=single,label=Example]
  !gapprompt@gap>| !gapinput@f:=PartialPerm( [ 1, 2, 3, 4, 5, 6, 7, 8, 9, 10 ],|
  !gapprompt@>| !gapinput@[ 2, 7, 9, 4, 1, 10, 5, 6, 3, 8 ] );|
  (1,2,7,5)(3,9)(4)(6,10,8)
  !gapprompt@gap>| !gapinput@AsPermutation(f);|
  (1,2,7,5)(3,9)(6,10,8)
  !gapprompt@gap>| !gapinput@f:= PartialPerm( [ 1, 2, 3, 4, 5, 7, 8 ], [ 5, 3, 8, 1, 9, 4, 10 ] );|
  [2,3,8,10][7,4,1,5,9]
  !gapprompt@gap>| !gapinput@AsPermutation(f);|
  fail
  !gapprompt@gap>| !gapinput@f:=Transformation( [ 5, 8, 3, 5, 8, 6, 2, 2, 7, 8 ] );;|
  !gapprompt@gap>| !gapinput@AsPermutation(f);|
  fail
  !gapprompt@gap>| !gapinput@f:=Transformation( [ 1, 3, 6, 6, 2, 10, 2, 3, 10, 5 ] );;|
  !gapprompt@gap>| !gapinput@AsPermutation(f);|
  fail
  !gapprompt@gap>| !gapinput@f:=Transformation( [ 2, 7, 9, 4, 1, 10, 5, 6, 3, 8 ] );|
  Transformation( [ 2, 7, 9, 4, 1, 10, 5, 6, 3, 8 ] )
  !gapprompt@gap>| !gapinput@AsPermutation(f);|
  (1,2,7,5)(3,9)(6,10,8)
\end{Verbatim}
 }

 }

 }

              
\chapter{\textcolor{Chapter }{Permutation Groups}}\label{Permutation Groups}
\logpage{[ 43, 0, 0 ]}
\hyperdef{L}{X85ED46007CED6191}{}
{
       
\section{\textcolor{Chapter }{IsPermGroup (Filter)}}\label{sect:IsPermGroup}
\logpage{[ 43, 1, 0 ]}
\hyperdef{L}{X7F38777E7BBE12AE}{}
{
  

\subsection{\textcolor{Chapter }{IsPermGroup}}
\logpage{[ 43, 1, 1 ]}\nobreak
\hyperdef{L}{X7879877482F59676}{}
{\noindent\textcolor{FuncColor}{$\triangleright$\enspace\texttt{IsPermGroup({\mdseries\slshape obj})\index{IsPermGroup@\texttt{IsPermGroup}}
\label{IsPermGroup}
}\hfill{\scriptsize (Category)}}\\


 A permutation group is a group of permutations on a finite set $\Omega$ of positive integers. \textsf{GAP} does \emph{not} require the user to specify the operation domain $\Omega$ when a permutation group is defined. 

 
\begin{Verbatim}[commandchars=!@|,fontsize=\small,frame=single,label=Example]
  !gapprompt@gap>| !gapinput@g:=Group((1,2,3,4),(1,2));|
  Group([ (1,2,3,4), (1,2) ])
\end{Verbatim}
 

 Permutation groups are groups and therefore all operations for groups (see
Chapter{\nobreakspace}\ref{Groups}) can be applied to them. In many cases special methods are installed for
permutation groups that make computations more effective. }

 }

  
\section{\textcolor{Chapter }{The Natural Action}}\label{The Natural Action}
\logpage{[ 43, 2, 0 ]}
\hyperdef{L}{X85D769FF85545AAB}{}
{
  The functions \texttt{MovedPoints} (\ref{MovedPoints:for a list or collection of permutations}), \texttt{NrMovedPoints} (\ref{NrMovedPoints:for a list or collection of permutations}), \texttt{LargestMovedPoint} (\ref{LargestMovedPoint:for a list or collection of permutations}), and \texttt{SmallestMovedPoint} (\ref{SmallestMovedPoint:for a list or collection of permutations}) are defined for arbitrary collections of permutations (see{\nobreakspace}\ref{Moved Points of Permutations}), in particular they can be applied to permutation groups. 
\begin{Verbatim}[commandchars=!@|,fontsize=\small,frame=single,label=Example]
  !gapprompt@gap>| !gapinput@g:= Group( (2,3,5,6), (2,3) );;|
  !gapprompt@gap>| !gapinput@MovedPoints( g );  NrMovedPoints( g );|
  [ 2, 3, 5, 6 ]
  4
  !gapprompt@gap>| !gapinput@LargestMovedPoint( g );  SmallestMovedPoint( g );|
  6
  2
\end{Verbatim}
 

 The action of a permutation group on the positive integers is a group action
(via the acting function \texttt{OnPoints} (\ref{OnPoints})). Therefore all action functions can be applied (see the
Chapter{\nobreakspace}\ref{Group Actions}), for example \texttt{Orbit} (\ref{Orbit}), \texttt{Stabilizer} (\ref{Stabilizer}), \texttt{Blocks} (\ref{Blocks:for a group, an action domain, etc.}), \texttt{IsTransitive} (\ref{IsTransitive:for a group, an action domain, etc.}), \texttt{IsPrimitive} (\ref{IsPrimitive:for a group, an action domain, etc.}). 

 If one has a list of group generators and is interested in the moved points
(see above) or orbits, it may be useful to avoid the explicit construction of
the group for efficiency reasons. For the special case of the action of
permutations on positive integers via \texttt{\texttt{\symbol{94}}}, the functions \texttt{OrbitPerms} (\ref{OrbitPerms}) and \texttt{OrbitsPerms} (\ref{OrbitsPerms}) are provided for this purpose. 

 Similarly, several functions concerning the natural action of permutation
groups address stabilizer chains (see{\nobreakspace}\ref{Stabilizer Chains}) rather than permutation groups themselves, for example \texttt{BaseStabChain} (\ref{BaseStabChain}). 

\subsection{\textcolor{Chapter }{OrbitPerms}}
\logpage{[ 43, 2, 1 ]}\nobreak
\hyperdef{L}{X84CFA16D858B00B8}{}
{\noindent\textcolor{FuncColor}{$\triangleright$\enspace\texttt{OrbitPerms({\mdseries\slshape perms, pnt})\index{OrbitPerms@\texttt{OrbitPerms}}
\label{OrbitPerms}
}\hfill{\scriptsize (function)}}\\


 returns the orbit of the positive integer \mbox{\texttt{\mdseries\slshape pnt}} under the group generated by the permutations in the list \mbox{\texttt{\mdseries\slshape perms}}. }

 

\subsection{\textcolor{Chapter }{OrbitsPerms}}
\logpage{[ 43, 2, 2 ]}\nobreak
\hyperdef{L}{X81F98222818DA35B}{}
{\noindent\textcolor{FuncColor}{$\triangleright$\enspace\texttt{OrbitsPerms({\mdseries\slshape perms, D})\index{OrbitsPerms@\texttt{OrbitsPerms}}
\label{OrbitsPerms}
}\hfill{\scriptsize (function)}}\\


 returns the list of orbits of the positive integers in the list \mbox{\texttt{\mdseries\slshape D}} under the group generated by the permutations in the list \mbox{\texttt{\mdseries\slshape perms}}. 
\begin{Verbatim}[commandchars=!@|,fontsize=\small,frame=single,label=Example]
  !gapprompt@gap>| !gapinput@OrbitPerms( [ (1,2,3)(4,5), (3,6) ], 1 );|
  [ 1, 2, 3, 6 ]
  !gapprompt@gap>| !gapinput@OrbitsPerms( [ (1,2,3)(4,5), (3,6) ], [ 1 .. 6 ] );|
  [ [ 1, 2, 3, 6 ], [ 4, 5 ] ]
\end{Verbatim}
 }

 }

  
\section{\textcolor{Chapter }{Computing a Permutation Representation}}\label{Computing a Permutation Representation}
\logpage{[ 43, 3, 0 ]}
\hyperdef{L}{X7E468B64860D5604}{}
{
  

\subsection{\textcolor{Chapter }{IsomorphismPermGroup}}
\logpage{[ 43, 3, 1 ]}\nobreak
\hyperdef{L}{X80B7B1C783AA1567}{}
{\noindent\textcolor{FuncColor}{$\triangleright$\enspace\texttt{IsomorphismPermGroup({\mdseries\slshape G})\index{IsomorphismPermGroup@\texttt{IsomorphismPermGroup}}
\label{IsomorphismPermGroup}
}\hfill{\scriptsize (attribute)}}\\


 returns an isomorphism from the group \mbox{\texttt{\mdseries\slshape G}} onto a permutation group which is isomorphic to \mbox{\texttt{\mdseries\slshape G}}. The method will select a suitable permutation representation. 
\begin{Verbatim}[commandchars=!@|,fontsize=\small,frame=single,label=Example]
  !gapprompt@gap>| !gapinput@g:=SmallGroup(24,12);|
  <pc group of size 24 with 4 generators>
  !gapprompt@gap>| !gapinput@iso:=IsomorphismPermGroup(g);|
  [ f1, f2, f3, f4 ] -> [ (2,3), (2,3,4), (1,2)(3,4), (1,3)(2,4) ]
  !gapprompt@gap>| !gapinput@Image(iso,g.3*g.4);|
  (1,4)(2,3)
\end{Verbatim}
 

 In many cases the permutation representation constructed by \texttt{IsomorphismPermGroup} is regular. }

 

\subsection{\textcolor{Chapter }{SmallerDegreePermutationRepresentation}}
\logpage{[ 43, 3, 2 ]}\nobreak
\hyperdef{L}{X8086628878AFD3EA}{}
{\noindent\textcolor{FuncColor}{$\triangleright$\enspace\texttt{SmallerDegreePermutationRepresentation({\mdseries\slshape G})\index{SmallerDegreePermutationRepresentation@\texttt{Smaller}\-\texttt{Degree}\-\texttt{Permutation}\-\texttt{Representation}}
\label{SmallerDegreePermutationRepresentation}
}\hfill{\scriptsize (function)}}\\


 Let \mbox{\texttt{\mdseries\slshape G}} be a permutation group. \texttt{SmallerDegreePermutationRepresentation} tries to find a faithful permutation representation of smaller degree. The
result is a group homomorphism onto a permutation group, in the worst case
this is the identity mapping on \mbox{\texttt{\mdseries\slshape G}}. 

 If the \texttt{cheap} option is given, the function only tries to reduce to orbits or actions on
blocks, otherwise also actions on cosets of random subgroups are tried. 

 Note that the result is not guaranteed to be a faithful permutation
representation of smallest degree, or of smallest degree among the transitive
permutation representations of \mbox{\texttt{\mdseries\slshape G}}. Using \textsf{GAP} interactively, one might be able to choose subgroups of small index for which
the cores intersect trivially; in this case, the actions on the cosets of
these subgroups give rise to an intransitive permutation representation the
degree of which may be smaller than the original degree. 

 The methods used might involve the use of random elements and the permutation
representation (or even the degree of the representation) is not guaranteed to
be the same for different calls of \texttt{SmallerDegreePermutationRepresentation}. 

 If the option cheap is given less work is spent on trying to get a small
degree representation, if the value of this option is set to the string "skip"
the identity mapping is returned. (This is useful if a function called
internally might try a degree reduction.) 

 
\begin{Verbatim}[commandchars=!@|,fontsize=\small,frame=single,label=Example]
  !gapprompt@gap>| !gapinput@iso:=RegularActionHomomorphism(SymmetricGroup(4));;|
  !gapprompt@gap>| !gapinput@image:= Image( iso );;  NrMovedPoints( image );|
  24
  !gapprompt@gap>| !gapinput@small:= SmallerDegreePermutationRepresentation( image );;|
  !gapprompt@gap>| !gapinput@Image( small );|
  Group([ (2,5,4,3), (1,4)(2,6)(3,5) ])
  !gapprompt@gap>| !gapinput@g:=Image(IsomorphismPermGroup(GL(4,5)));;|
  !gapprompt@gap>| !gapinput@sm:=SmallerDegreePermutationRepresentation(g:cheap);;|
  !gapprompt@gap>| !gapinput@NrMovedPoints(Range(sm));|
  624
\end{Verbatim}
 }

 
\begin{Verbatim}[commandchars=!@|,fontsize=\small,frame=single,label=Example]
  !gapprompt@gap>| !gapinput@p:=Group((1,2,3,4,5,6),(1,2));;p:=Action(p,AsList(p),OnRight);;|
  !gapprompt@gap>| !gapinput@Length(MovedPoints(p));|
  720
  !gapprompt@gap>| !gapinput@q:=SmallerDegreePermutationRepresentation(p);;|
  !gapprompt@gap>| !gapinput@NrMovedPoints(Image(q));|
  6
\end{Verbatim}
 }

  
\section{\textcolor{Chapter }{Symmetric and Alternating Groups}}\label{Symmetric and Alternating Groups}
\logpage{[ 43, 4, 0 ]}
\hyperdef{L}{X834208CD7C2956A3}{}
{
  The commands \texttt{SymmetricGroup} (\ref{SymmetricGroup:for a degree}) and \texttt{AlternatingGroup} (\ref{AlternatingGroup:for a degree}) (see Section{\nobreakspace}\ref{Basic Groups}) construct symmetric and alternating permutation groups. \textsf{GAP} can also detect whether a given permutation group is a symmetric or
alternating group on the set of its moved points; if so then the group is
called a \emph{natural} symmetric or alternating group, respectively. 

 The functions \texttt{IsSymmetricGroup} (\ref{IsSymmetricGroup}) and \texttt{IsAlternatingGroup} (\ref{IsAlternatingGroup}) can be used to check whether a given group (not necessarily a permutation
group) is isomorphic to a symmetric or alternating group. 

 

\subsection{\textcolor{Chapter }{IsNaturalSymmetricGroup}}
\logpage{[ 43, 4, 1 ]}\nobreak
\hyperdef{L}{X8129BE59781478E1}{}
{\noindent\textcolor{FuncColor}{$\triangleright$\enspace\texttt{IsNaturalSymmetricGroup({\mdseries\slshape group})\index{IsNaturalSymmetricGroup@\texttt{IsNaturalSymmetricGroup}}
\label{IsNaturalSymmetricGroup}
}\hfill{\scriptsize (property)}}\\
\noindent\textcolor{FuncColor}{$\triangleright$\enspace\texttt{IsNaturalAlternatingGroup({\mdseries\slshape group})\index{IsNaturalAlternatingGroup@\texttt{IsNaturalAlternatingGroup}}
\label{IsNaturalAlternatingGroup}
}\hfill{\scriptsize (property)}}\\


 A group is a natural symmetric or alternating group if it is a permutation
group acting as symmetric or alternating group, respectively, on its moved
points. 

 For groups that are known to be natural symmetric or natural alternating
groups, very efficient methods for computing membership, conjugacy classes,
Sylow subgroups etc.{\nobreakspace}are used. 

 
\begin{Verbatim}[commandchars=!@|,fontsize=\small,frame=single,label=Example]
  !gapprompt@gap>| !gapinput@g:=Group((1,5,7,8,99),(1,99,13,72));;|
  !gapprompt@gap>| !gapinput@IsNaturalSymmetricGroup(g);|
  true
  !gapprompt@gap>| !gapinput@g;|
  Sym( [ 1, 5, 7, 8, 13, 72, 99 ] )
  !gapprompt@gap>| !gapinput@IsNaturalSymmetricGroup( Group( (1,2)(4,5), (1,2,3)(4,5,6) ) );|
  false
\end{Verbatim}
 }

 

\subsection{\textcolor{Chapter }{IsSymmetricGroup}}
\logpage{[ 43, 4, 2 ]}\nobreak
\hyperdef{L}{X85CA6AD17BE90C95}{}
{\noindent\textcolor{FuncColor}{$\triangleright$\enspace\texttt{IsSymmetricGroup({\mdseries\slshape group})\index{IsSymmetricGroup@\texttt{IsSymmetricGroup}}
\label{IsSymmetricGroup}
}\hfill{\scriptsize (property)}}\\


 is \texttt{true} if the group \mbox{\texttt{\mdseries\slshape group}} is isomorphic to a symmetric group. }

 

\subsection{\textcolor{Chapter }{IsAlternatingGroup}}
\logpage{[ 43, 4, 3 ]}\nobreak
\hyperdef{L}{X8514BE9E79C608E0}{}
{\noindent\textcolor{FuncColor}{$\triangleright$\enspace\texttt{IsAlternatingGroup({\mdseries\slshape group})\index{IsAlternatingGroup@\texttt{IsAlternatingGroup}}
\label{IsAlternatingGroup}
}\hfill{\scriptsize (property)}}\\


 is \texttt{true} if the group \mbox{\texttt{\mdseries\slshape group}} is isomorphic to an alternating group. }

 

\subsection{\textcolor{Chapter }{SymmetricParentGroup}}
\logpage{[ 43, 4, 4 ]}\nobreak
\hyperdef{L}{X7ED60F7E81F1B614}{}
{\noindent\textcolor{FuncColor}{$\triangleright$\enspace\texttt{SymmetricParentGroup({\mdseries\slshape grp})\index{SymmetricParentGroup@\texttt{SymmetricParentGroup}}
\label{SymmetricParentGroup}
}\hfill{\scriptsize (attribute)}}\\


 For a permutation group \mbox{\texttt{\mdseries\slshape grp}} this function returns the symmetric group that moves the same points as \mbox{\texttt{\mdseries\slshape grp}} does. 
\begin{Verbatim}[commandchars=!@|,fontsize=\small,frame=single,label=Example]
  !gapprompt@gap>| !gapinput@SymmetricParentGroup( Group( (1,2), (4,5), (7,8,9) ) );|
  Sym( [ 1, 2, 4, 5, 7, 8, 9 ] )
\end{Verbatim}
 }

 }

  
\section{\textcolor{Chapter }{Primitive Groups}}\label{Primitive Groups}
\logpage{[ 43, 5, 0 ]}
\hyperdef{L}{X83F8D3B578A7BEEB}{}
{
  

\subsection{\textcolor{Chapter }{ONanScottType}}
\logpage{[ 43, 5, 1 ]}\nobreak
\hyperdef{L}{X7E50211A7B92455F}{}
{\noindent\textcolor{FuncColor}{$\triangleright$\enspace\texttt{ONanScottType({\mdseries\slshape G})\index{ONanScottType@\texttt{ONanScottType}}
\label{ONanScottType}
}\hfill{\scriptsize (attribute)}}\\


 returns the type of a primitive permutation group \mbox{\texttt{\mdseries\slshape G}}, according to the O'Nan\texttt{\symbol{45}}Scott classification. The
labelling of the different types is not consistent in the literature, we use
the following identifications. The two\texttt{\symbol{45}}letter code given is
the name of the type as used by Praeger. 
\begin{description}
\item[{1}]  Affine. (HA) 
\item[{2}]  Almost simple. (AS) 
\item[{3a}]  Diagonal, Socle consists of two normal subgroups. (HS) 
\item[{3b}]  Diagonal, Socle is minimal normal. (SD) 
\item[{4a}]  Product action with the first factor primitive of type 3a. (HC) 
\item[{4b}]  Product action with the first factor primitive of type 3b. (CD) 
\item[{4c}]  Product action with the first factor primitive of type 2. (PA) 
\item[{5}]  Twisted wreath product (TW) 
\end{description}
 See \cite{EickHulpke01} for correspondence to other labellings used in the literature. As it can
contain letters, the type is returned as a string. 

 If \mbox{\texttt{\mdseries\slshape G}} is not a permutation group or does not act primitively on the points moved by
it, the result is undefined. }

 

\subsection{\textcolor{Chapter }{SocleTypePrimitiveGroup}}
\logpage{[ 43, 5, 2 ]}\nobreak
\hyperdef{L}{X7E89A46A86A3F4A2}{}
{\noindent\textcolor{FuncColor}{$\triangleright$\enspace\texttt{SocleTypePrimitiveGroup({\mdseries\slshape G})\index{SocleTypePrimitiveGroup@\texttt{SocleTypePrimitiveGroup}}
\label{SocleTypePrimitiveGroup}
}\hfill{\scriptsize (attribute)}}\\


 returns the socle type of the primitive permutation group \mbox{\texttt{\mdseries\slshape G}}. The socle of a primitive group is the direct product of isomorphic simple
groups, therefore the type is indicated by a record with components \texttt{series}, \texttt{parameter} (both as described under \texttt{IsomorphismTypeInfoFiniteSimpleGroup} (\ref{IsomorphismTypeInfoFiniteSimpleGroup:for a group})), and \texttt{width} for the number of direct factors. 

 If \mbox{\texttt{\mdseries\slshape G}} does not act primitively on its moved points, an error is returned. 
\begin{Verbatim}[commandchars=!@|,fontsize=\small,frame=single,label=Example]
  !gapprompt@gap>| !gapinput@g:=AlternatingGroup(5);;|
  !gapprompt@gap>| !gapinput@h:=DirectProduct(g,g);;|
  !gapprompt@gap>| !gapinput@p:=List([1,2],i->Projection(h,i));;|
  !gapprompt@gap>| !gapinput@ac:=Action(h,AsList(g),|
  !gapprompt@>| !gapinput@function(g,h) return Image(p[1],h)^-1*g*Image(p[2],h);end);;|
  !gapprompt@gap>| !gapinput@Size(ac);NrMovedPoints(ac);IsPrimitive(ac,[1..60]);|
  3600
  60
  true
  !gapprompt@gap>| !gapinput@ONanScottType(ac);|
  "3a"
  !gapprompt@gap>| !gapinput@SocleTypePrimitiveGroup(ac);|
  rec(
    name := "A(5) ~ A(1,4) = L(2,4) ~ B(1,4) = O(3,4) ~ C(1,4) = S(2,4) \
  ~ 2A(1,4) = U(2,4) ~ A(1,5) = L(2,5) ~ B(1,5) = O(3,5) ~ C(1,5) = S(2,\
  5) ~ 2A(1,5) = U(2,5)", parameter := 5, series := "A", width := 2 )
\end{Verbatim}
 }

 }

  
\section{\textcolor{Chapter }{Stabilizer Chains}}\label{Stabilizer Chains}
\logpage{[ 43, 6, 0 ]}
\hyperdef{L}{X7FA58C3A8283F3BD}{}
{
  Many of the algorithms for permutation groups use a \emph{stabilizer chain} of the group. The concepts of stabilizer chains, \emph{bases}, and \emph{strong generating sets} were introduced by Charles Sims in{\nobreakspace}\cite{Sim70}. An extensive account of basic algorithms together with asymptotic runtime
analysis can be found in reference{\nobreakspace}\cite[Chapter 4]{Seress2003}. A further discussion of base change is given in section{\nobreakspace}\ref{Generalized Conjugation Technique}. 

 Let $B = [ b_1, \ldots, b_n ]$ be a list of points, $G^{(1)} = G$ and $G^{{(i+1)}} = Stab_{{G^{(i)}}}(b_i)$, such that $G^{(n+1)} = \{ () \}$. Then the list $[ b_1, \ldots, b_n ]$ is called a \emph{base} of $G$, the points $b_i$ are called \emph{base points}. A set $S$ of generators for $G$ satisfying the condition $\langle S \cap G^{(i)} \rangle = G^{(i)}$ for each $1 \leq i \leq n$, is called a \emph{strong generating set} (SGS) of $G$. (More precisely we ought to say that it is a SGS of $G$ \emph{relative} to $B$). The chain of subgroups $G^{(i)}$ of $G$ itself is called the \emph{stabilizer chain} of $G$ relative to $B$. 

 Since $[ b_1, \ldots, b_n ]$, where $n$ is the degree of $G$ and $b_i$ are the moved points of $G$, certainly is a base for $G$ there exists a base for each permutation group. The number of points in a base
is called the \emph{length} of the base. A base $B$ is called \emph{reduced} if there exists no $i$ such that $G^{(i)} = G^{(i+1)}$. (This however does not imply that no subset of $B$ could also serve as a base.) Note that different reduced bases for one
permutation group $G$ may have different lengths. For example, the irreducible degree $416$ permutation representation of the Chevalley Group $G_2(4)$ possesses reduced bases of lengths $5$ and $7$. 

 Let $R^{(i)}$ be a right transversal of $G^{(i+1)}$ in $G^{(i)}$, i.e. a set of right coset representatives of the cosets of $G^{(i+1)}$ in $G^{(i)}$. Then each element $g$ of $G$ has a unique representation as a product of the form $g = r_n \ldots r_1$ with $r_i \in R^{(i)}$. The cosets of $G^{(i+1)}$ in $G^{(i)}$ are in bijective correspondence with the points in $O^{(i)} := b_i^{{G^{(i)}}}$. So we could represent a transversal as a list $T$ such that $T[p]$ is a representative of the coset corresponding to the point $p \in O^{(i)}$, i.e., an element of $G^{(i)}$ that takes $b_i$ to $p$. (Note that such a list has holes in all positions corresponding to points
not contained in $O^{(i)}$.) 

 This approach however will store many different permutations as coset
representatives which can be a problem if the degree $n$ gets bigger. Our goal therefore is to store as few different permutations as
possible such that we can still reconstruct each representative in $R^{(i)}$, and from them the elements in $G$. A \emph{factorized inverse transversal} $T$ is a list where $T[p]$ is a generator of $G^{(i)}$ such that $p^{{T[p]}}$ is a point that lies earlier in $O^{(i)}$ than $p$ (note that we consider $O^{(i)}$ as a list, not as a set). If we assume inductively that we know an element $r \in G^{(i)}$ that takes $b_i$ to $p^{{T[p]}}$, then $r T[p]^{{-1}}$ is an element in $G^{(i)}$ that takes $b_i$ to $p$. \textsf{GAP} uses such factorized inverse transversals. 

 Another name for a factorized inverse transversal is a \emph{Schreier tree}. The vertices of the tree are the points in $O^{(i)}$, and the root of the tree is $b_i$. The edges are defined as the ordered pairs $(p, p^{{T[p]}})$, for $p \in O^{(i)} \setminus \{ b_i \}$. The edge $(p, p^{{T[p]}})$ is labelled with the generator $T[p]$, and the product of edge labels along the unique path from $p$ to $b_i$ is the inverse of the transversal element carrying $b_i$ to $p$. 

 Before we describe the construction of stabilizer chains in{\nobreakspace}\ref{Construction of Stabilizer Chains}, we explain in{\nobreakspace}\ref{Randomized Methods for Permutation Groups} the idea of using non\texttt{\symbol{45}}deterministic algorithms; this is
necessary for understanding the options available for the construction of
stabilizer chains. After that, in{\nobreakspace}\ref{Stabilizer Chain Records} it is explained how a stabilizer chain is stored in \textsf{GAP}, \ref{Operations for Stabilizer Chains} lists operations for stabilizer chains, and{\nobreakspace}\ref{Low Level Routines to Modify and Create Stabilizer Chains} lists low level routines for manipulating stabilizer chains. }

  
\section{\textcolor{Chapter }{Randomized Methods for Permutation Groups}}\label{Randomized Methods for Permutation Groups}
\logpage{[ 43, 7, 0 ]}
\hyperdef{L}{X7C2406B97E057196}{}
{
  \index{Schreier\texttt{\symbol{45}}Sims!random} For most computations with permutation groups, it is crucial to construct
stabilizer chains efficiently. Sims's original construction in \cite{Sim70} is deterministic, and is called the Schreier\texttt{\symbol{45}}Sims
algorithm, because it is based on Schreier's Lemma (\cite[p. 96]{Hall}): given $K = \langle S \rangle$ and a transversal $T$ for $K$ mod $L$, one can obtain $|S||T|$ generators for $L$. This lemma is applied recursively, with consecutive point stabilizers $G^{(i)}$ and $G^{(i+1)}$ playing the role of $K$ and $L$. 

 In permutation groups of large degree, the number of Schreier generators to be
processed becomes too large, and the deterministic
Schreier\texttt{\symbol{45}}Sims algorithm becomes impractical. Therefore, \textsf{GAP} uses randomized algorithms. The method selection process, which is quite
different from Version 3, works the following way. 

 If a group acts on not more than a hundred points, Sims's original
deterministic algorithm is applied. In groups of degree greater than hundred,
a heuristic algorithm based on ideas in \cite{BCFS91} constructs a stabilizer chain. This construction is complemented by a
verify\texttt{\symbol{45}}routine that either proves the correctness of the
stabilizer chain or causes the extension of the chain to a correct one. The
user can influence the verification process by setting the value of the record
component \texttt{random} (cf.{\nobreakspace}\ref{Construction of Stabilizer Chains}). 

 If the \texttt{random} value equals $1000$ then a slight extension of an unpublished method of Sims is used. The outcome
of this verification process is always correct. The user also can prescribe
any integer $x$, $1 \leq x \leq 999$ as the value of \texttt{random}. In this case, a randomized verification process from \cite{BCFS91} is applied, and the result of the stabilizer chain construction is guaranteed
to be correct with probability at least $x/1000$. The practical performance of the algorithm is much better than the
theoretical guarantee. 

 If the stabilizer chain is not correct then the elements in the product of
transversals $R^{(m)} R^{(m-1)} \cdots R^{(1)}$ constitute a proper subset of the group $G$ in question. This means that a membership test with this stabilizer chain
returns \texttt{false} for all elements that are not in $G$, but it may also return \texttt{false} for some elements of $G$; in other words, the result \texttt{true} of a membership test is always correct, whereas the result \texttt{false} may be incorrect. 

 The construction and verification phases are separated because there are
situations where the verification step can be omitted; if one happens to know
the order of the group in advance then the randomized construction of the
stabilizer chain stops as soon as the product of the lengths of the basic
orbits of the chain equals the group order, and the chain will be correct (see
the \texttt{size} option of the \texttt{StabChain} (\ref{StabChain:for a group (and a record)}) command). 

 Although the worst case running time is roughly quadratic for Sims's
verification and roughly linear for the randomized one, in most examples the
running time of the stabilizer chain construction with \texttt{random} value $1000$ (i.e., guaranteed correct output) is about the same as the running time of
randomized verification with guarantee of at least $90$ percent correctness. Therefore, we suggest to use the default value \texttt{random} $= 1000$. Possible uses of \texttt{random} values less than $1000$ are when one has to run through a large collection of subgroups, and a low
value of random is used to choose quickly a candidate for more thorough
examination; another use is when the user suspects that the quadratic
bottleneck of the guaranteed correct verification is hit. 

 We will give two examples to illustrate these ideas. 

 
\begin{Verbatim}[commandchars=!@|,fontsize=\small,frame=single,label=Example]
  !gapprompt@gap>| !gapinput@h:= SL(4,7);;|
  !gapprompt@gap>| !gapinput@o:= Orbit( h, [1,0,0,0]*Z(7)^0, OnLines );;|
  !gapprompt@gap>| !gapinput@op:= Action( h, o, OnLines );;|
  !gapprompt@gap>| !gapinput@NrMovedPoints( op );|
  400
\end{Verbatim}
 

 We created a permutation group on $400$ points. First we compute a guaranteed correct stabilizer chain
(see{\nobreakspace}\texttt{StabChain} (\ref{StabChain:for a group (and a record)})). 

 
\begin{Verbatim}[commandchars=!@|,fontsize=\small,frame=single,label=Example]
  !gapprompt@gap>| !gapinput@h:= Group( GeneratorsOfGroup( op ) );;|
  !gapprompt@gap>| !gapinput@StabChain( h );;  time;|
  1120
  !gapprompt@gap>| !gapinput@Size( h );|
  2317591180800
\end{Verbatim}
 

 Now randomized verification will be used. We require that the result is
guaranteed correct with probability $90$ percent. This means that if we would do this calculation many times over, \textsf{GAP} would \emph{guarantee} that in least $90$ percent of all calculations the result is correct. In fact the results are
much better than the guarantee, but we cannot promise that this will really
happen. (For the meaning of the \texttt{random} component in the second argument of \texttt{StabChain} (\ref{StabChain:for a group (and a record)}).) 

 First the group is created anew. 

 
\begin{Verbatim}[commandchars=!@|,fontsize=\small,frame=single,label=Example]
  !gapprompt@gap>| !gapinput@h:= Group( GeneratorsOfGroup( op ) );;|
  !gapprompt@gap>| !gapinput@StabChain( h, rec( random:= 900 ) );;  time;|
  1410
  !gapprompt@gap>| !gapinput@Size( h );|
  2317591180800
\end{Verbatim}
 

 The result is still correct, and the running time is actually somewhat slower.
If you give the algorithm the order of the group, then it can check its
result, and so things become faster and the result is guaranteed to be
correct. This can be done with the \texttt{size} option (see{\nobreakspace}\texttt{StabChain} (\ref{StabChain:for a group (and a record)})), or by setting the size of the group beforehand with \texttt{SetSize}. 

 
\begin{Verbatim}[commandchars=!@|,fontsize=\small,frame=single,label=Example]
  !gapprompt@gap>| !gapinput@h:=Group( GeneratorsOfGroup( op ) );;|
  !gapprompt@gap>| !gapinput@SetSize( h, 2317591180800 );|
  !gapprompt@gap>| !gapinput@StabChain( h );;  time;|
  170
\end{Verbatim}
 

 The second example gives a typical group when the verification with \texttt{random} value $1000$ is slow. The problem is that the group has a stabilizer subgroup $G^{(i)}$ such that the fundamental orbit $O^{(i)}$ is split into a lot of orbits when we stabilize $b_i$ and one additional point of $O^{(i)}$. 

 
\begin{Verbatim}[commandchars=!@|,fontsize=\small,frame=single,label=Example]
  !gapprompt@gap>| !gapinput@p1:=PermList(Concatenation([401],[1..400]));;|
  !gapprompt@gap>| !gapinput@p2:=PermList(List([1..400],i->(i*20 mod 401)));;|
  !gapprompt@gap>| !gapinput@d:=DirectProduct(Group(p1,p2),SymmetricGroup(5));;|
  !gapprompt@gap>| !gapinput@h:=Group(GeneratorsOfGroup(d));;|
  !gapprompt@gap>| !gapinput@StabChain(h);;time;Size(h);|
  1030
  192480
  !gapprompt@gap>| !gapinput@h:=Group(GeneratorsOfGroup(d));;|
  !gapprompt@gap>| !gapinput@StabChain(h,rec(random:=900));;time;Size(h);|
  570
  192480
\end{Verbatim}
 

 When stabilizer chains of a group $G$ are created with \texttt{random} value less than $1000$, this is noted in the group $G$, by setting of the record component \texttt{random} in the value of the attribute \texttt{StabChainOptions} (\ref{StabChainOptions}) for $G$. As errors induced by the random methods might propagate, any group or
homomorphism created from $G$ inherits a \texttt{random} component in its \texttt{StabChainOptions} (\ref{StabChainOptions}) value from the corresponding component for $G$. 

 A lot of algorithms dealing with permutation groups use randomized methods;
however, if the initial stabilizer chain construction for a group is correct,
these further methods will provide guaranteed correct output. }

  
\section{\textcolor{Chapter }{Construction of Stabilizer Chains}}\label{Construction of Stabilizer Chains}
\logpage{[ 43, 8, 0 ]}
\hyperdef{L}{X7C7EA55C80E457FA}{}
{
  

\subsection{\textcolor{Chapter }{StabChain (for a group (and a record))}}
\logpage{[ 43, 8, 1 ]}\nobreak
\hyperdef{L}{X80B5CF78829495C2}{}
{\noindent\textcolor{FuncColor}{$\triangleright$\enspace\texttt{StabChain({\mdseries\slshape G[, options]})\index{StabChain@\texttt{StabChain}!for a group (and a record)}
\label{StabChain:for a group (and a record)}
}\hfill{\scriptsize (function)}}\\
\noindent\textcolor{FuncColor}{$\triangleright$\enspace\texttt{StabChain({\mdseries\slshape G, base})\index{StabChain@\texttt{StabChain}!for a group and a base}
\label{StabChain:for a group and a base}
}\hfill{\scriptsize (function)}}\\
\noindent\textcolor{FuncColor}{$\triangleright$\enspace\texttt{StabChainOp({\mdseries\slshape G, options})\index{StabChainOp@\texttt{StabChainOp}}
\label{StabChainOp}
}\hfill{\scriptsize (operation)}}\\
\noindent\textcolor{FuncColor}{$\triangleright$\enspace\texttt{StabChainMutable({\mdseries\slshape G})\index{StabChainMutable@\texttt{StabChainMutable}!for a group}
\label{StabChainMutable:for a group}
}\hfill{\scriptsize (attribute)}}\\
\noindent\textcolor{FuncColor}{$\triangleright$\enspace\texttt{StabChainMutable({\mdseries\slshape permhomom})\index{StabChainMutable@\texttt{StabChainMutable}!for a homomorphism}
\label{StabChainMutable:for a homomorphism}
}\hfill{\scriptsize (attribute)}}\\
\noindent\textcolor{FuncColor}{$\triangleright$\enspace\texttt{StabChainImmutable({\mdseries\slshape G})\index{StabChainImmutable@\texttt{StabChainImmutable}}
\label{StabChainImmutable}
}\hfill{\scriptsize (attribute)}}\\


 These commands compute a stabilizer chain for the permutation group \mbox{\texttt{\mdseries\slshape G}}; additionally, \texttt{StabChainMutable} is also an attribute for the group homomorphism \mbox{\texttt{\mdseries\slshape permhomom}} whose source is a permutation group. 

 (The mathematical background of stabilizer chains is sketched
in{\nobreakspace}\ref{Stabilizer Chains}, more information about the objects representing stabilizer chains in \textsf{GAP} can be found in{\nobreakspace}\ref{Stabilizer Chain Records}.) 

 \texttt{StabChainOp} is an operation with two arguments \mbox{\texttt{\mdseries\slshape G}} and \mbox{\texttt{\mdseries\slshape options}}, the latter being a record which controls some aspects of the computation of
a stabilizer chain (see below); \texttt{StabChainOp} returns a \emph{mutable} stabilizer chain. \texttt{StabChainMutable} is a \emph{mutable} attribute for groups or homomorphisms, its default method for groups is to
call \texttt{StabChainOp} with empty options record. \texttt{StabChainImmutable} is an attribute with \emph{immutable} values; its default method dispatches to \texttt{StabChainMutable}. 

 \texttt{StabChain} is a function with first argument a permutation group \mbox{\texttt{\mdseries\slshape G}}, and optionally a record \mbox{\texttt{\mdseries\slshape options}} as second argument. If the value of \texttt{StabChainImmutable} for \mbox{\texttt{\mdseries\slshape G}} is already known and if this stabilizer chain matches the requirements of \mbox{\texttt{\mdseries\slshape options}}, \texttt{StabChain} simply returns this stored stabilizer chain. Otherwise \texttt{StabChain} calls \texttt{StabChainOp} and returns an immutable copy of the result; additionally, this chain is
stored as \texttt{StabChainImmutable} value for \mbox{\texttt{\mdseries\slshape G}}. If no \mbox{\texttt{\mdseries\slshape options}} argument is given, its components default to the global variable \texttt{DefaultStabChainOptions} (\ref{DefaultStabChainOptions}). If \mbox{\texttt{\mdseries\slshape base}} is a list of positive integers, the version \texttt{StabChain( \mbox{\texttt{\mdseries\slshape G}}, \mbox{\texttt{\mdseries\slshape base}} )} defaults to \texttt{StabChain( \mbox{\texttt{\mdseries\slshape G}}, rec( base:= \mbox{\texttt{\mdseries\slshape base}} ) )}. 

 If given, \mbox{\texttt{\mdseries\slshape options}} is a record whose components specify properties of the desired stabilizer
chain or which may help the algorithm. Default values for all of them can be
given in the global variable \texttt{DefaultStabChainOptions} (\ref{DefaultStabChainOptions}). The following options are supported. 
\begin{description}
\item[{\texttt{base} (default an empty list)}]  A list of points, through which the resulting stabilizer chain shall run. For
the base $B$ of the resulting stabilizer chain \mbox{\texttt{\mdseries\slshape S}} this means the following. If the \texttt{reduced} component of \mbox{\texttt{\mdseries\slshape options}} is \texttt{true} then those points of \texttt{base} with nontrivial basic orbits form the initial segment of $B$, if the \texttt{reduced} component is \texttt{false} then \texttt{base} itself is the initial segment of $B$. Repeated occurrences of points in \texttt{base} are ignored. If a stabilizer chain for \mbox{\texttt{\mdseries\slshape G}} is already known then the stabilizer chain is computed via a base change. 
\item[{\texttt{knownBase} (no default value)}]  A list of points which is known to be a base for the group. Such a known base
makes it easier to test whether a permutation given as a word in terms of a
set of generators is the identity, since it suffices to map the known base
with each factor consecutively, rather than multiplying the whole permutations
(which would mean to map every point). This speeds up the
Schreier\texttt{\symbol{45}}Sims algorithm which is used when a new stabilizer
chain is constructed; it will not affect a base change, however. The component \texttt{knownBase} bears no relation to the \texttt{base} component, you may specify a known base \texttt{knownBase} and a desired base \texttt{base} independently. 
\item[{\texttt{reduced} (default \texttt{true})}]  If this is \texttt{true} the resulting stabilizer chain \mbox{\texttt{\mdseries\slshape S}} is reduced, i.e., the case $G^{(i)} = G^{(i+1)}$ does not occur. Setting \texttt{reduced} to \texttt{false} makes sense only if the component \texttt{base} (see above) is also set; in this case all points of \texttt{base} will occur in the base $B$ of \mbox{\texttt{\mdseries\slshape S}}, even if they have trivial basic orbits. Note that if \texttt{base} is just an initial segment of $B$, the basic orbits of the points in $B \setminus $\texttt{base} are always nontrivial. 
\item[{\texttt{tryPcgs} (default \texttt{true})}]  If this is \texttt{true} and either the degree is at most $100$ or the group is known to be solvable, \textsf{GAP} will first try to construct a pcgs (see Chapter{\nobreakspace}\ref{Polycyclic Groups}) for \mbox{\texttt{\mdseries\slshape G}} which will succeed and implicitly construct a stabilizer chain if \mbox{\texttt{\mdseries\slshape G}} is solvable. If \mbox{\texttt{\mdseries\slshape G}} turns out non\texttt{\symbol{45}}solvable, one of the other methods will be
used. This solvability check is comparatively fast, even if it fails, and it
can save a lot of time if \mbox{\texttt{\mdseries\slshape G}} is solvable. 
\item[{\texttt{random} (default \texttt{1000})}]  If the value is less than{\nobreakspace}$1000$, the resulting chain is correct with probability at least{\nobreakspace}\texttt{random}$ / 1000$. The \texttt{random} option is explained in more detail in{\nobreakspace}\ref{Randomized Methods for Permutation Groups}. 
\item[{\texttt{size} (default \texttt{Size(\mbox{\texttt{\mdseries\slshape G}})} if this is known, i.e., if \texttt{HasSize(\mbox{\texttt{\mdseries\slshape G}})} is \texttt{true})}]  If this component is present, its value is assumed to be the order of the
group \mbox{\texttt{\mdseries\slshape G}}. This information can be used to prove that a
non\texttt{\symbol{45}}deterministically constructed stabilizer chain is
correct. In this case, \textsf{GAP} does a non\texttt{\symbol{45}}deterministic construction until the size is
correct. 
\item[{\texttt{limit} (default \texttt{Size(Parent(\mbox{\texttt{\mdseries\slshape G}}))} or \texttt{StabChainOptions(Parent(\mbox{\texttt{\mdseries\slshape G}})).limit} if it is present)}]  If this component is present, it must be greater than or equal to the order of \mbox{\texttt{\mdseries\slshape G}}. The stabilizer chain construction stops if size \texttt{limit} is reached. 
\end{description}
 }

 

\subsection{\textcolor{Chapter }{StabChainOptions}}
\logpage{[ 43, 8, 2 ]}\nobreak
\hyperdef{L}{X790C27B8783EDE68}{}
{\noindent\textcolor{FuncColor}{$\triangleright$\enspace\texttt{StabChainOptions({\mdseries\slshape G})\index{StabChainOptions@\texttt{StabChainOptions}}
\label{StabChainOptions}
}\hfill{\scriptsize (attribute)}}\\


 is a record that stores the options with which the stabilizer chain stored in \texttt{StabChainImmutable} (\ref{StabChainImmutable}) has been computed (see{\nobreakspace}\texttt{StabChain} (\ref{StabChain:for a group (and a record)}) for the options that are supported). }

 

\subsection{\textcolor{Chapter }{DefaultStabChainOptions}}
\logpage{[ 43, 8, 3 ]}\nobreak
\hyperdef{L}{X87E1292E85A5D31C}{}
{\noindent\textcolor{FuncColor}{$\triangleright$\enspace\texttt{DefaultStabChainOptions\index{DefaultStabChainOptions@\texttt{DefaultStabChainOptions}}
\label{DefaultStabChainOptions}
}\hfill{\scriptsize (global variable)}}\\


 are the options for \texttt{StabChain} (\ref{StabChain:for a group (and a record)}) which are set as default. }

 

\subsection{\textcolor{Chapter }{StabChainBaseStrongGenerators}}
\logpage{[ 43, 8, 4 ]}\nobreak
\hyperdef{L}{X86D64D2B81D58431}{}
{\noindent\textcolor{FuncColor}{$\triangleright$\enspace\texttt{StabChainBaseStrongGenerators({\mdseries\slshape base, sgs[, one]})\index{StabChainBaseStrongGenerators@\texttt{StabChainBaseStrongGenerators}}
\label{StabChainBaseStrongGenerators}
}\hfill{\scriptsize (function)}}\\


 Let \mbox{\texttt{\mdseries\slshape base}} be a base for a permutation group $G$, and let \mbox{\texttt{\mdseries\slshape sgs}} be a strong generating set for $G$ with respect to \mbox{\texttt{\mdseries\slshape base}}; \mbox{\texttt{\mdseries\slshape one}} must be the appropriate identity element of $G$ (see \texttt{One} (\ref{One}), in most cases this will be \texttt{()}). This function constructs a stabilizer chain corresponding to the given base
and strong generating set without the need to find Schreier generators; so
this is much faster than the other algorithms. 

 If \mbox{\texttt{\mdseries\slshape sgs}} is nonempty, then the argument \mbox{\texttt{\mdseries\slshape one}} is optional; if not given, then the \texttt{One} (\ref{One}) of \texttt{\mbox{\texttt{\mdseries\slshape sgs}}[1]} is taken as the identity element. 
\begin{Verbatim}[commandchars=!@|,fontsize=\small,frame=single,label=Example]
  !gapprompt@gap>| !gapinput@sc := StabChainBaseStrongGenerators([1,2], [(1,3,4), (2,3,4)], ());|
  <stabilizer chain record, Base [ 1, 2 ], Orbit length 4, Size: 12>
  !gapprompt@gap>| !gapinput@GroupStabChain(sc) = AlternatingGroup(4);|
  true
  !gapprompt@gap>| !gapinput@StabChainBaseStrongGenerators([1,3], [(1,2),(3,4)]);|
  <stabilizer chain record, Base [ 1, 3 ], Orbit length 2, Size: 4>
\end{Verbatim}
 }

 

\subsection{\textcolor{Chapter }{MinimalStabChain}}
\logpage{[ 43, 8, 5 ]}\nobreak
\hyperdef{L}{X7BEC5F5A7851CAAB}{}
{\noindent\textcolor{FuncColor}{$\triangleright$\enspace\texttt{MinimalStabChain({\mdseries\slshape G})\index{MinimalStabChain@\texttt{MinimalStabChain}}
\label{MinimalStabChain}
}\hfill{\scriptsize (attribute)}}\\


 returns the reduced stabilizer chain corresponding to the base $[ 1, 2, 3, 4, \ldots ]$. }

 }

  
\section{\textcolor{Chapter }{Stabilizer Chain Records}}\label{Stabilizer Chain Records}
\logpage{[ 43, 9, 0 ]}
\hyperdef{L}{X81D7FCE47AC7F942}{}
{
  If a permutation group has a stabilizer chain, this is stored as a recursive
structure. This structure is itself a record{\nobreakspace}\mbox{\texttt{\mdseries\slshape S}} and it has 
\begin{description}
\item[{(1)}]  components that provide information about one level{\nobreakspace}$G^{(i)}$ of the stabilizer chain (which we call the ``current stabilizer'') and 
\item[{(2)}]  a component \texttt{stabilizer} that holds another such record, namely the stabilizer chain of the next
stabilizer{\nobreakspace}$G^{(i+1)}$. 
\end{description}
 

 This gives a recursive structure where the ``outermost'' record representing the ``topmost'' stabilizer is bound to the group record component \texttt{stabChain} and has the components explained below. Note: Since the structure is
recursive, \emph{never print a stabilizer chain!} (Unless you want to exercise the scrolling capabilities of your terminal.) 
\begin{description}
\item[{\texttt{identity} }]  the identity element of the current stabilizer. 
\item[{\texttt{labels} }]  a list of permutations which contains labels for the Schreier tree of the
current stabilizer, i.e., it contains elements for the factorized inverse
transversal. The first entry in this list is always the \texttt{identity}. Note that \textsf{GAP} tries to arrange things so that the \texttt{labels} components are identical (i.e., the same \textsf{GAP} object) in every stabilizer of the chain; thus the \texttt{labels} of a stabilizer do not necessarily all lie in the this stabilizer (but see \texttt{genlabels} below). 
\item[{\texttt{genlabels} }]  a list of integers indexing some of the permutations in the \texttt{labels} component. The \texttt{labels} addressed in this way form a generating set for the current stabilizer. If the \texttt{genlabels} component is empty, the rest of the stabilizer chain represents the trivial
subgroup, and can be ignored, e.g., when calculating the size. 
\item[{\texttt{generators} }]  a list of generators for the current stabilizer. Usually, it is \texttt{labels\texttt{\symbol{123}} genlabels \texttt{\symbol{125}}}. 
\item[{\texttt{orbit} }]  the vertices of the Schreier tree, which form the basic orbit $b_i^{{G^{(i)}}}$, ordered in such a way that the base point $b_i$ is in the first position in the orbit. 
\item[{\texttt{transversal}}]  The factorized inverse transversal found during the orbit algorithm. The
element $g$ stored at \texttt{transversal}$[i]$ will map $i$ to another point $j$ that in the Schreier tree is closer to the base point. By iterated application
(\texttt{transversal}$[j]$ and so on) eventually the base point is reached and an element that maps $i$ to the base point found as product. 
\item[{\texttt{translabels}}]  An index list such that \texttt{transversal}$[j] =$ \texttt{labels}$[$ \texttt{translabels}$[j] ]$. This list takes up comparatively little memory and is used to speed up base
changes. 
\item[{\texttt{stabilizer} }]  If the current stabilizer is not yet the trivial group, the stabilizer chain
continues with the stabilizer of the current base point, which is again
represented as a record with components \texttt{labels}, \texttt{identity}, \texttt{genlabels}, \texttt{generators}, \texttt{orbit}, \texttt{translabels}, \texttt{transversal} (and perhaps \texttt{stabilizer}). This record is bound to the \texttt{stabilizer} component of the current stabilizer. The last member of a stabilizer chain is
recognized by the fact that it has no \texttt{stabilizer} component bound. 
\end{description}
 It is possible that different stabilizer chains share the same record as one
of their iterated \texttt{stabilizer} components. 

 
\begin{Verbatim}[commandchars=!@|,fontsize=\small,frame=single,label=Example]
  !gapprompt@gap>| !gapinput@g:=Group((1,2,3,4),(1,2));;|
  !gapprompt@gap>| !gapinput@StabChain(g);|
  <stabilizer chain record, Base [ 1, 2, 3 ], Orbit length 4, Size: 24>
  !gapprompt@gap>| !gapinput@BaseOfGroup(g);|
  [ 1, 2, 3 ]
  !gapprompt@gap>| !gapinput@StabChainOptions(g);|
  rec( random := 1000 )
  !gapprompt@gap>| !gapinput@DefaultStabChainOptions;|
  rec( random := 1000, reduced := true, tryPcgs := true )
\end{Verbatim}
 }

  
\section{\textcolor{Chapter }{Operations for Stabilizer Chains}}\label{Operations for Stabilizer Chains}
\logpage{[ 43, 10, 0 ]}
\hyperdef{L}{X7ECF8A4586346FD4}{}
{
  

\subsection{\textcolor{Chapter }{BaseStabChain}}
\logpage{[ 43, 10, 1 ]}\nobreak
\hyperdef{L}{X7FBE6EB57EBE8B7D}{}
{\noindent\textcolor{FuncColor}{$\triangleright$\enspace\texttt{BaseStabChain({\mdseries\slshape S})\index{BaseStabChain@\texttt{BaseStabChain}}
\label{BaseStabChain}
}\hfill{\scriptsize (function)}}\\


 returns the base belonging to the stabilizer chain \mbox{\texttt{\mdseries\slshape S}}. }

 

\subsection{\textcolor{Chapter }{BaseOfGroup}}
\logpage{[ 43, 10, 2 ]}\nobreak
\hyperdef{L}{X7D2A190D8308ED39}{}
{\noindent\textcolor{FuncColor}{$\triangleright$\enspace\texttt{BaseOfGroup({\mdseries\slshape G})\index{BaseOfGroup@\texttt{BaseOfGroup}}
\label{BaseOfGroup}
}\hfill{\scriptsize (attribute)}}\\


 returns a base of the permutation group \mbox{\texttt{\mdseries\slshape G}}. There is \emph{no} guarantee that a stabilizer chain stored in \mbox{\texttt{\mdseries\slshape G}} corresponds to this base! }

 

\subsection{\textcolor{Chapter }{SizeStabChain}}
\logpage{[ 43, 10, 3 ]}\nobreak
\hyperdef{L}{X7EF36DC78465026A}{}
{\noindent\textcolor{FuncColor}{$\triangleright$\enspace\texttt{SizeStabChain({\mdseries\slshape S})\index{SizeStabChain@\texttt{SizeStabChain}}
\label{SizeStabChain}
}\hfill{\scriptsize (function)}}\\


 returns the product of the orbit lengths in the stabilizer chain \mbox{\texttt{\mdseries\slshape S}}, that is, the order of the group described by \mbox{\texttt{\mdseries\slshape S}}. }

 

\subsection{\textcolor{Chapter }{StrongGeneratorsStabChain}}
\logpage{[ 43, 10, 4 ]}\nobreak
\hyperdef{L}{X8384170881B9B531}{}
{\noindent\textcolor{FuncColor}{$\triangleright$\enspace\texttt{StrongGeneratorsStabChain({\mdseries\slshape S})\index{StrongGeneratorsStabChain@\texttt{StrongGeneratorsStabChain}}
\label{StrongGeneratorsStabChain}
}\hfill{\scriptsize (function)}}\\


 returns a strong generating set corresponding to the stabilizer chain \mbox{\texttt{\mdseries\slshape S}}. }

 

\subsection{\textcolor{Chapter }{GroupStabChain}}
\logpage{[ 43, 10, 5 ]}\nobreak
\hyperdef{L}{X87F473777EFDE867}{}
{\noindent\textcolor{FuncColor}{$\triangleright$\enspace\texttt{GroupStabChain({\mdseries\slshape [G, ]S})\index{GroupStabChain@\texttt{GroupStabChain}}
\label{GroupStabChain}
}\hfill{\scriptsize (function)}}\\


 constructs a permutation group with stabilizer chain \mbox{\texttt{\mdseries\slshape S}}, i.e., a group with generators \texttt{Generators( \mbox{\texttt{\mdseries\slshape S}} )} to which \mbox{\texttt{\mdseries\slshape S}} is assigned as component \texttt{stabChain}. If the optional argument \mbox{\texttt{\mdseries\slshape G}} is given, the result will have the parent \mbox{\texttt{\mdseries\slshape G}}. }

 

\subsection{\textcolor{Chapter }{OrbitStabChain}}
\logpage{[ 43, 10, 6 ]}\nobreak
\hyperdef{L}{X87FB6DED80692D3F}{}
{\noindent\textcolor{FuncColor}{$\triangleright$\enspace\texttt{OrbitStabChain({\mdseries\slshape S, pnt})\index{OrbitStabChain@\texttt{OrbitStabChain}}
\label{OrbitStabChain}
}\hfill{\scriptsize (function)}}\\


 returns the orbit of \mbox{\texttt{\mdseries\slshape pnt}} under the group described by the stabilizer chain \mbox{\texttt{\mdseries\slshape S}}. }

 

\subsection{\textcolor{Chapter }{IndicesStabChain}}
\logpage{[ 43, 10, 7 ]}\nobreak
\hyperdef{L}{X7AC8F165875906DE}{}
{\noindent\textcolor{FuncColor}{$\triangleright$\enspace\texttt{IndicesStabChain({\mdseries\slshape S})\index{IndicesStabChain@\texttt{IndicesStabChain}}
\label{IndicesStabChain}
}\hfill{\scriptsize (function)}}\\


 returns a list of the indices of the stabilizers in the stabilizer chain \mbox{\texttt{\mdseries\slshape S}}. }

 

\subsection{\textcolor{Chapter }{ListStabChain}}
\logpage{[ 43, 10, 8 ]}\nobreak
\hyperdef{L}{X7CF607BC82C2C202}{}
{\noindent\textcolor{FuncColor}{$\triangleright$\enspace\texttt{ListStabChain({\mdseries\slshape S})\index{ListStabChain@\texttt{ListStabChain}}
\label{ListStabChain}
}\hfill{\scriptsize (function)}}\\


 returns a list that contains at position $i$ the stabilizer of the first $i-1$ base points in the stabilizer chain \mbox{\texttt{\mdseries\slshape S}}. }

 

\subsection{\textcolor{Chapter }{ElementsStabChain}}
\logpage{[ 43, 10, 9 ]}\nobreak
\hyperdef{L}{X7F40E52D7B0438BF}{}
{\noindent\textcolor{FuncColor}{$\triangleright$\enspace\texttt{ElementsStabChain({\mdseries\slshape S})\index{ElementsStabChain@\texttt{ElementsStabChain}}
\label{ElementsStabChain}
}\hfill{\scriptsize (function)}}\\


 returns a list of all elements of the group described by the stabilizer chain \mbox{\texttt{\mdseries\slshape S}}. The list is duplicate free but may be unsorted. }

 

\subsection{\textcolor{Chapter }{IteratorStabChain}}
\logpage{[ 43, 10, 10 ]}\nobreak
\hyperdef{L}{X780875477CD2A57D}{}
{\noindent\textcolor{FuncColor}{$\triangleright$\enspace\texttt{IteratorStabChain({\mdseries\slshape S})\index{IteratorStabChain@\texttt{IteratorStabChain}}
\label{IteratorStabChain}
}\hfill{\scriptsize (function)}}\\


 returns an iterator for the elements of the group described by the stabilizer
chain \mbox{\texttt{\mdseries\slshape S}}. The elements of the group \mbox{\texttt{\mdseries\slshape G}} are produced by iterating through all base images in turn, and in the ordering
induced by the base. For more details see{\nobreakspace}\ref{Stabilizer Chains} }

 

\subsection{\textcolor{Chapter }{InverseRepresentative}}
\logpage{[ 43, 10, 11 ]}\nobreak
\hyperdef{L}{X861062AE87ACF340}{}
{\noindent\textcolor{FuncColor}{$\triangleright$\enspace\texttt{InverseRepresentative({\mdseries\slshape S, pnt})\index{InverseRepresentative@\texttt{InverseRepresentative}}
\label{InverseRepresentative}
}\hfill{\scriptsize (function)}}\\


 calculates the transversal element which maps \mbox{\texttt{\mdseries\slshape pnt}} back to the base point of \mbox{\texttt{\mdseries\slshape S}}. It just runs back through the Schreier tree from \mbox{\texttt{\mdseries\slshape pnt}} to the root and multiplies the labels along the way. }

 

\subsection{\textcolor{Chapter }{SiftedPermutation}}
\logpage{[ 43, 10, 12 ]}\nobreak
\hyperdef{L}{X79D2248C8787EAF2}{}
{\noindent\textcolor{FuncColor}{$\triangleright$\enspace\texttt{SiftedPermutation({\mdseries\slshape S, g})\index{SiftedPermutation@\texttt{SiftedPermutation}}
\label{SiftedPermutation}
}\hfill{\scriptsize (function)}}\\


 sifts the permutation \mbox{\texttt{\mdseries\slshape g}} through the stabilizer chain \mbox{\texttt{\mdseries\slshape S}} and returns the result after the last step. 

 The element \mbox{\texttt{\mdseries\slshape g}} is sifted as follows: \mbox{\texttt{\mdseries\slshape g}} is replaced by \texttt{\mbox{\texttt{\mdseries\slshape g}} * InverseRepresentative( \mbox{\texttt{\mdseries\slshape S}}, \mbox{\texttt{\mdseries\slshape S}}.orbit[1]\texttt{\symbol{94}}\mbox{\texttt{\mdseries\slshape g}} )}, then \mbox{\texttt{\mdseries\slshape S}} is replaced by \texttt{\mbox{\texttt{\mdseries\slshape S}}.stabilizer} and this process is repeated until \mbox{\texttt{\mdseries\slshape S}} is trivial or \texttt{\mbox{\texttt{\mdseries\slshape S}}.orbit[1]\texttt{\symbol{94}}\mbox{\texttt{\mdseries\slshape g}}} is not in the basic orbit \texttt{\mbox{\texttt{\mdseries\slshape S}}.orbit}. The remainder \mbox{\texttt{\mdseries\slshape g}} is returned, it is the identity permutation if and only if the original \mbox{\texttt{\mdseries\slshape g}} is in the group $G$ described by the original{\nobreakspace}\mbox{\texttt{\mdseries\slshape S}}. }

 

\subsection{\textcolor{Chapter }{MinimalElementCosetStabChain}}
\logpage{[ 43, 10, 13 ]}\nobreak
\hyperdef{L}{X7B870C217D0B9997}{}
{\noindent\textcolor{FuncColor}{$\triangleright$\enspace\texttt{MinimalElementCosetStabChain({\mdseries\slshape S, g})\index{MinimalElementCosetStabChain@\texttt{MinimalElementCosetStabChain}}
\label{MinimalElementCosetStabChain}
}\hfill{\scriptsize (function)}}\\


 Let $G$ be the group described by the stabilizer chain \mbox{\texttt{\mdseries\slshape S}}. This function returns a permutation $h$ such that $G \mbox{\texttt{\mdseries\slshape g}} = G h$ (that is, $\mbox{\texttt{\mdseries\slshape g}} / h \in G$) and with the additional property that the list of images under $h$ of the base belonging to \mbox{\texttt{\mdseries\slshape S}} is minimal w.r.t.{\nobreakspace}lexicographical ordering. }

 

\subsection{\textcolor{Chapter }{LargestElementStabChain}}
\logpage{[ 43, 10, 14 ]}\nobreak
\hyperdef{L}{X87435B7884D9B353}{}
{\noindent\textcolor{FuncColor}{$\triangleright$\enspace\texttt{LargestElementStabChain({\mdseries\slshape S, id})\index{LargestElementStabChain@\texttt{LargestElementStabChain}}
\label{LargestElementStabChain}
}\hfill{\scriptsize (function)}}\\


 Let $G$ be the group described by the stabilizer chain \mbox{\texttt{\mdseries\slshape S}}. This function returns the element $h \in G$ with the property that the list of images under $h$ of the base belonging to \mbox{\texttt{\mdseries\slshape S}} is maximal w.r.t.{\nobreakspace}lexicographical ordering. The second argument
must be an identity element (used to start the recursion). }

 

\subsection{\textcolor{Chapter }{ApproximateSuborbitsStabilizerPermGroup}}
\logpage{[ 43, 10, 15 ]}\nobreak
\hyperdef{L}{X809B2C3B7C5F77AB}{}
{\noindent\textcolor{FuncColor}{$\triangleright$\enspace\texttt{ApproximateSuborbitsStabilizerPermGroup({\mdseries\slshape G, pnt})\index{ApproximateSuborbitsStabilizerPermGroup@\texttt{Approximate}\-\texttt{Suborbits}\-\texttt{Stabilizer}\-\texttt{Perm}\-\texttt{Group}}
\label{ApproximateSuborbitsStabilizerPermGroup}
}\hfill{\scriptsize (function)}}\\


 returns an approximation of the orbits of \texttt{Stabilizer( \mbox{\texttt{\mdseries\slshape G}}, \mbox{\texttt{\mdseries\slshape pnt}} )} on all points of the orbit \texttt{Orbit( \mbox{\texttt{\mdseries\slshape G}}, \mbox{\texttt{\mdseries\slshape pnt}} )}, without computing the full point stabilizer; As not all Schreier generators
are used, the result may represent the orbits of only a subgroup of the point
stabilizer. }

 }

  
\section{\textcolor{Chapter }{Low Level Routines to Modify and Create Stabilizer Chains}}\label{Low Level Routines to Modify and Create Stabilizer Chains}
\logpage{[ 43, 11, 0 ]}
\hyperdef{L}{X8188051F79E72A95}{}
{
  These operations modify a stabilizer chain or obtain new chains with specific
properties. They are rather technical and should only be used if such
low\texttt{\symbol{45}}level routines are deliberately required. (For all
functions in this section the parameter \mbox{\texttt{\mdseries\slshape S}} is a stabilizer chain.) 

\subsection{\textcolor{Chapter }{CopyStabChain}}
\logpage{[ 43, 11, 1 ]}\nobreak
\hyperdef{L}{X86B31E6A81AE5FCB}{}
{\noindent\textcolor{FuncColor}{$\triangleright$\enspace\texttt{CopyStabChain({\mdseries\slshape S})\index{CopyStabChain@\texttt{CopyStabChain}}
\label{CopyStabChain}
}\hfill{\scriptsize (function)}}\\


 This function returns a mutable copy of the stabilizer chain \mbox{\texttt{\mdseries\slshape S}} that has no mutable object (list or record) in common with \mbox{\texttt{\mdseries\slshape S}}. The \texttt{labels} components of the result are possibly shared by several levels, but
superfluous labels are removed. (An entry in \texttt{labels} is superfluous if it does not occur among the \texttt{genlabels} or \texttt{translabels} on any of the levels which share that \texttt{labels} component.) 

 This is useful for stabiliser sub\texttt{\symbol{45}}chains that have been
obtained as the (iterated) \texttt{stabilizer} component of a bigger chain. }

 

\subsection{\textcolor{Chapter }{CopyOptionsDefaults}}
\logpage{[ 43, 11, 2 ]}\nobreak
\hyperdef{L}{X7E167E557B567C6A}{}
{\noindent\textcolor{FuncColor}{$\triangleright$\enspace\texttt{CopyOptionsDefaults({\mdseries\slshape G, options})\index{CopyOptionsDefaults@\texttt{CopyOptionsDefaults}}
\label{CopyOptionsDefaults}
}\hfill{\scriptsize (function)}}\\


 sets components in a stabilizer chain options record \mbox{\texttt{\mdseries\slshape options}} according to what is known about the group \mbox{\texttt{\mdseries\slshape G}}. This can be used to obtain a new stabilizer chain for \mbox{\texttt{\mdseries\slshape G}} quickly. }

 

\subsection{\textcolor{Chapter }{ChangeStabChain}}
\logpage{[ 43, 11, 3 ]}\nobreak
\hyperdef{L}{X87FF64AB87BFC779}{}
{\noindent\textcolor{FuncColor}{$\triangleright$\enspace\texttt{ChangeStabChain({\mdseries\slshape S, base[, reduced]})\index{ChangeStabChain@\texttt{ChangeStabChain}}
\label{ChangeStabChain}
}\hfill{\scriptsize (function)}}\\


 changes or reduces a stabilizer chain \mbox{\texttt{\mdseries\slshape S}} to be adapted to the base \mbox{\texttt{\mdseries\slshape base}}. The optional argument \mbox{\texttt{\mdseries\slshape reduced}} is interpreted as follows. 
\begin{description}
\item[{\texttt{reduced = }\texttt{false} : }]  change the stabilizer chain, do not reduce it, 
\item[{\texttt{reduced = }\texttt{true} : }]  change the stabilizer chain, reduce it. 
\end{description}
 }

 

\subsection{\textcolor{Chapter }{ExtendStabChain}}
\logpage{[ 43, 11, 4 ]}\nobreak
\hyperdef{L}{X8778B4657D3FD97B}{}
{\noindent\textcolor{FuncColor}{$\triangleright$\enspace\texttt{ExtendStabChain({\mdseries\slshape S, base})\index{ExtendStabChain@\texttt{ExtendStabChain}}
\label{ExtendStabChain}
}\hfill{\scriptsize (function)}}\\


 extends the stabilizer chain \mbox{\texttt{\mdseries\slshape S}} so that it corresponds to base \mbox{\texttt{\mdseries\slshape base}}. The original base of \mbox{\texttt{\mdseries\slshape S}} must be a subset of \mbox{\texttt{\mdseries\slshape base}}. }

 

\subsection{\textcolor{Chapter }{ReduceStabChain}}
\logpage{[ 43, 11, 5 ]}\nobreak
\hyperdef{L}{X7E5E9F727D0B19D9}{}
{\noindent\textcolor{FuncColor}{$\triangleright$\enspace\texttt{ReduceStabChain({\mdseries\slshape S})\index{ReduceStabChain@\texttt{ReduceStabChain}}
\label{ReduceStabChain}
}\hfill{\scriptsize (function)}}\\


 changes the stabilizer chain \mbox{\texttt{\mdseries\slshape S}} to a reduced stabilizer chain by eliminating trivial steps. }

  

\subsection{\textcolor{Chapter }{RemoveStabChain}}
\logpage{[ 43, 11, 6 ]}\nobreak
\hyperdef{L}{X85BF290D848C4091}{}
{\noindent\textcolor{FuncColor}{$\triangleright$\enspace\texttt{RemoveStabChain({\mdseries\slshape S})\index{RemoveStabChain@\texttt{RemoveStabChain}}
\label{RemoveStabChain}
}\hfill{\scriptsize (function)}}\\


 \mbox{\texttt{\mdseries\slshape S}} must be a stabilizer record in a stabilizer chain. This chain then is cut off
at \mbox{\texttt{\mdseries\slshape S}} by changing the entries in \mbox{\texttt{\mdseries\slshape S}}. This can be used to remove trailing trivial steps. }

 

\subsection{\textcolor{Chapter }{EmptyStabChain}}
\logpage{[ 43, 11, 7 ]}\nobreak
\hyperdef{L}{X84E4906B86E5C089}{}
{\noindent\textcolor{FuncColor}{$\triangleright$\enspace\texttt{EmptyStabChain({\mdseries\slshape labels, id[, pnt]})\index{EmptyStabChain@\texttt{EmptyStabChain}}
\label{EmptyStabChain}
}\hfill{\scriptsize (function)}}\\


 constructs a stabilizer chain for the trivial group with \texttt{identity} value equal to\mbox{\texttt{\mdseries\slshape id}} and \texttt{labels = }$\{ \mbox{\texttt{\mdseries\slshape id}} \} \cup$ \mbox{\texttt{\mdseries\slshape labels}} (but of course with \texttt{genlabels} and \texttt{generators} values an empty list). If the optional third argument \mbox{\texttt{\mdseries\slshape pnt}} is present, the only stabilizer of the chain is initialized with the
one\texttt{\symbol{45}}point basic orbit \texttt{[ \mbox{\texttt{\mdseries\slshape pnt}} ]} and with \texttt{translabels} and \texttt{transversal} components. }

 

\subsection{\textcolor{Chapter }{InsertTrivialStabilizer}}
\logpage{[ 43, 11, 8 ]}\nobreak
\hyperdef{L}{X80C7D2E87E6EE357}{}
{\noindent\textcolor{FuncColor}{$\triangleright$\enspace\texttt{InsertTrivialStabilizer({\mdseries\slshape S, pnt})\index{InsertTrivialStabilizer@\texttt{InsertTrivialStabilizer}}
\label{InsertTrivialStabilizer}
}\hfill{\scriptsize (function)}}\\


 \texttt{InsertTrivialStabilizer} initializes the current stabilizer with \mbox{\texttt{\mdseries\slshape pnt}} as \texttt{EmptyStabChain} (\ref{EmptyStabChain}) did, but assigns the original \mbox{\texttt{\mdseries\slshape S}} to the new \texttt{\mbox{\texttt{\mdseries\slshape S}}.stabilizer} component, such that a new level with trivial basic orbit (but identical \texttt{labels} and \texttt{ShallowCopy}ed \texttt{genlabels} and \texttt{generators}) is inserted. This function should be used only if \mbox{\texttt{\mdseries\slshape pnt}} really is fixed by the generators of \mbox{\texttt{\mdseries\slshape S}}, because then new generators can be added and the orbit and transversal at
the same time extended with \texttt{AddGeneratorsExtendSchreierTree} (\ref{AddGeneratorsExtendSchreierTree}). }

 

\subsection{\textcolor{Chapter }{IsFixedStabilizer}}
\logpage{[ 43, 11, 9 ]}\nobreak
\hyperdef{L}{X7B47B379824F6150}{}
{\noindent\textcolor{FuncColor}{$\triangleright$\enspace\texttt{IsFixedStabilizer({\mdseries\slshape S, pnt})\index{IsFixedStabilizer@\texttt{IsFixedStabilizer}}
\label{IsFixedStabilizer}
}\hfill{\scriptsize (function)}}\\


 returns \texttt{true} if \mbox{\texttt{\mdseries\slshape pnt}} is fixed by all generators of \mbox{\texttt{\mdseries\slshape S}} and \texttt{false} otherwise. }

 

\subsection{\textcolor{Chapter }{AddGeneratorsExtendSchreierTree}}
\logpage{[ 43, 11, 10 ]}\nobreak
\hyperdef{L}{X8373007880EBF736}{}
{\noindent\textcolor{FuncColor}{$\triangleright$\enspace\texttt{AddGeneratorsExtendSchreierTree({\mdseries\slshape S, new})\index{AddGeneratorsExtendSchreierTree@\texttt{AddGeneratorsExtendSchreierTree}}
\label{AddGeneratorsExtendSchreierTree}
}\hfill{\scriptsize (function)}}\\


 adds the elements in \mbox{\texttt{\mdseries\slshape new}} to the list of generators of \mbox{\texttt{\mdseries\slshape S}} and at the same time extends the orbit and transversal. This is the only legal
way to extend a Schreier tree (because this involves careful handling of the
tree components). }

 }

  
\section{\textcolor{Chapter }{Backtrack}}\label{Backtrack}
\logpage{[ 43, 12, 0 ]}
\hyperdef{L}{X86C78160854C7F30}{}
{
  A main use for stabilizer chains is in backtrack algorithms for permutation
groups. \textsf{GAP} implements a partition\texttt{\symbol{45}}backtrack algorithm as described in \cite{Leon91} and refined in \cite{Theissen97}. 

 

\subsection{\textcolor{Chapter }{SubgroupProperty}}
\logpage{[ 43, 12, 1 ]}\nobreak
\hyperdef{L}{X7BE3F03C80BF8B08}{}
{\noindent\textcolor{FuncColor}{$\triangleright$\enspace\texttt{SubgroupProperty({\mdseries\slshape G, Pr[, L]})\index{SubgroupProperty@\texttt{SubgroupProperty}}
\label{SubgroupProperty}
}\hfill{\scriptsize (function)}}\\


 \mbox{\texttt{\mdseries\slshape Pr}} must be a one\texttt{\symbol{45}}argument function that returns \texttt{true} or \texttt{false} for elements of the permutation group \mbox{\texttt{\mdseries\slshape G}}, and the subset of elements of \mbox{\texttt{\mdseries\slshape G}} that fulfill \mbox{\texttt{\mdseries\slshape Pr}} must be a subgroup. (\emph{If the latter is not true the result of this operation is unpredictable!}) This command computes this subgroup. The optional argument \mbox{\texttt{\mdseries\slshape L}} must be a subgroup of the set of all elements in \mbox{\texttt{\mdseries\slshape G}} fulfilling \mbox{\texttt{\mdseries\slshape Pr}} and can be given if known in order to speed up the calculation. }

 

\subsection{\textcolor{Chapter }{ElementProperty}}
\logpage{[ 43, 12, 2 ]}\nobreak
\hyperdef{L}{X7EE7DDCC87C4BC31}{}
{\noindent\textcolor{FuncColor}{$\triangleright$\enspace\texttt{ElementProperty({\mdseries\slshape G, Pr[, L[, R]]})\index{ElementProperty@\texttt{ElementProperty}}
\label{ElementProperty}
}\hfill{\scriptsize (function)}}\\


 \texttt{ElementProperty} returns an element $\pi$ of the permutation group \mbox{\texttt{\mdseries\slshape G}} such that the one\texttt{\symbol{45}}argument function \mbox{\texttt{\mdseries\slshape Pr}} returns \texttt{true} for $\pi$. It returns \texttt{fail} if no such element exists in \mbox{\texttt{\mdseries\slshape G}}. The optional arguments \mbox{\texttt{\mdseries\slshape L}} and \mbox{\texttt{\mdseries\slshape R}} are subgroups of \mbox{\texttt{\mdseries\slshape G}} such that the property \mbox{\texttt{\mdseries\slshape Pr}} has the same value for all elements in the cosets \mbox{\texttt{\mdseries\slshape L}} $g$ and $g$ \mbox{\texttt{\mdseries\slshape R}}, respectively, with $g \in \mbox{\texttt{\mdseries\slshape G}}$. 

 A typical example of using the optional subgroups \mbox{\texttt{\mdseries\slshape L}} and \mbox{\texttt{\mdseries\slshape R}} is the conjugacy test for elements $a$ and $b$ for which one can set \mbox{\texttt{\mdseries\slshape L}}$:= C_{\mbox{\texttt{\mdseries\slshape G}}}(a)$ and \mbox{\texttt{\mdseries\slshape R}}$:= C_{\mbox{\texttt{\mdseries\slshape G}}}(b)$. 

 
\begin{Verbatim}[commandchars=!@|,fontsize=\small,frame=single,label=Example]
  !gapprompt@gap>| !gapinput@propfun:= el -> (1,2,3)^el in [ (1,2,3), (1,3,2) ];;|
  !gapprompt@gap>| !gapinput@SubgroupProperty( g, propfun, Subgroup( g, [ (1,2,3) ] ) );|
  Group([ (1,2,3), (2,3) ])
  !gapprompt@gap>| !gapinput@ElementProperty( g, el -> Order( el ) = 2 );|
  (2,4)
\end{Verbatim}
 

 Chapter{\nobreakspace}\ref{Permutations} describes special operations to construct permutations in the symmetric group
without using backtrack constructions. 

 Backtrack routines are also called by the methods for permutation groups that
compute centralizers, normalizers, intersections, conjugating elements as well
as stabilizers for the operations of a permutation group via \texttt{OnPoints} (\ref{OnPoints}), \texttt{OnSets} (\ref{OnSets}), \texttt{OnTuples} (\ref{OnTuples}) and \texttt{OnSetsSets} (\ref{OnSetsSets}). Some of these methods use more specific refinements than \texttt{SubgroupProperty} (\ref{SubgroupProperty}) or \texttt{ElementProperty}. For the definition of refinements, and how one can define refinements, see
Section{\nobreakspace}\ref{The General Backtrack Algorithm with Ordered Partitions}. }

 

\subsection{\textcolor{Chapter }{TwoClosure}}
\logpage{[ 43, 12, 3 ]}\nobreak
\hyperdef{L}{X7A2D046B83DD5F5F}{}
{\noindent\textcolor{FuncColor}{$\triangleright$\enspace\texttt{TwoClosure({\mdseries\slshape G})\index{TwoClosure@\texttt{TwoClosure}}
\label{TwoClosure}
}\hfill{\scriptsize (attribute)}}\\


 The \emph{2\texttt{\symbol{45}}closure} of a transitive permutation group \mbox{\texttt{\mdseries\slshape G}} on $n$ points is the largest subgroup of the symmetric group $S_n$ which has the same orbits on sets of ordered pairs of points as the group \mbox{\texttt{\mdseries\slshape G}} has. It also can be interpreted as the stabilizer of the orbital graphs of \mbox{\texttt{\mdseries\slshape G}}. 
\begin{Verbatim}[commandchars=!@|,fontsize=\small,frame=single,label=Example]
  !gapprompt@gap>| !gapinput@TwoClosure(Group((1,2,3),(2,3,4)));|
  Sym( [ 1 .. 4 ] )
\end{Verbatim}
 }

 

\subsection{\textcolor{Chapter }{InfoBckt}}
\logpage{[ 43, 12, 4 ]}\nobreak
\hyperdef{L}{X861461AB7964DC64}{}
{\noindent\textcolor{FuncColor}{$\triangleright$\enspace\texttt{InfoBckt\index{InfoBckt@\texttt{InfoBckt}}
\label{InfoBckt}
}\hfill{\scriptsize (info class)}}\\


 is the info class for the partition backtrack routines. }

 }

  
\section{\textcolor{Chapter }{Working with large degree permutation groups}}\label{Working with large degree permutation groups}
\logpage{[ 43, 13, 0 ]}
\hyperdef{L}{X78A68F5A80ADD1B6}{}
{
  Permutation groups of large degree (usually at least a few $10000$) can pose a challenge to the heuristics used in the algorithms for
permutation groups. This section lists a few useful tricks that may speed up
calculations with such large groups enormously. 

 The first aspect concerns solvable groups: A lot of calculations (including an
initial stabilizer chain computation thanks to the algorithm from \cite{Sims90b}) are faster if a permutation group is known to be solvable. On the other
hand, proving nonsolvability can be expensive for higher degrees. Therefore \textsf{GAP} will automatically test a permutation group for solvability, only if the
degree is not exceeding $100$. (See also the \texttt{tryPcgs} component of \texttt{StabChainOptions} (\ref{StabChainOptions}).) It is therefore beneficial to tell a group of larger degree, which is known
to be solvable, that it is, using \texttt{SetIsSolvableGroup(\mbox{\texttt{\mdseries\slshape G}},true)}. 

 The second aspect concerns memory usage. A permutation on more than $65536$ points requires $4$ bytes per point for storing. So permutations on $256000$ points require roughly 1MB of storage per permutation. Just storing the
permutations required for a stabilizer chain might already go beyond the
available memory, in particular if the base is not very short. In such a
situation it can be useful, to replace the permutations by straight line
program elements (see{\nobreakspace}\ref{Straight Line Program Elements}). 

 The following code gives an example of usage: We create a group of degree $231000$. Using straight line program elements, one can compute a stabilizer chain in
about $200$ MB of memory. 

    
\begin{Verbatim}[commandchars=!@|,fontsize=\small,frame=single,label=Example]
  !gapprompt@gap>| !gapinput@Read("largeperms"); # read generators from file|
  !gapprompt@gap>| !gapinput@gens:=StraightLineProgGens(permutationlist);;|
  !gapprompt@gap>| !gapinput@g:=Group(gens);|
  <permutation group with 5 generators>
  !gapprompt@gap>| !gapinput@# use random algorithm (faster, but result is monte carlo)|
  !gapprompt@gap>| !gapinput@StabChainOptions(g).random:=1;;|
  !gapprompt@gap>| !gapinput@Size(g); # enforce computation of a stabilizer chain|
  3529698298145066075557232833758234188056080273649172207877011796336000
\end{Verbatim}
 

 Without straight line program elements, the same calculation runs into memory
problems after a while even with 512MB of workspace: 
\begin{Verbatim}[commandchars=!@|,fontsize=\small,frame=single,label=Example]
  !gapprompt@gap>| !gapinput@h:=Group(permutationlist);|
  <permutation group with 5 generators>
  !gapprompt@gap>| !gapinput@StabChainOptions(h).random:=1;;|
  !gapprompt@gap>| !gapinput@Size(h);|
  exceeded the permitted memory (`-o' command line option)
  ...
\end{Verbatim}
 

 The advantage in memory usage however is paid for in runtime: Comparisons of
elements become much more expensive. One can avoid some of the related
problems by registering a known base with the straight line program elements
(see{\nobreakspace}\texttt{StraightLineProgGens} (\ref{StraightLineProgGens})). In this case element comparison will only compare the images of the given
base points. If we are planning to do extensive calculations with the group,
it can even be worth to recreate it with straight line program elements
knowing a previously computed base: 

 
\begin{Verbatim}[commandchars=!@|,fontsize=\small,frame=single,label=Example]
  !gapprompt@gap>| !gapinput@# get the base we computed already|
  !gapprompt@gap>| !gapinput@bas:=BaseStabChain(StabChainMutable(g));|
  [ 1, 4, 7, 10, 13, 16, 19, 22, 25, 28, 31, 34, 37, 40, 43, 46, 49, 52, 55,
  ...
    2530, 2533, 2554, 2563, 2569 ]
  !gapprompt@gap>| !gapinput@gens:=StraightLineProgGens(permutationlist,bas);;|
  !gapprompt@gap>| !gapinput@g:=Group(gens);;|
  !gapprompt@gap>| !gapinput@SetSize(g,|
  !gapprompt@>| !gapinput@3529698298145066075557232833758234188056080273649172207877011796336000);|
  !gapprompt@gap>| !gapinput@Random(g);; # enforce computation of a stabilizer chain|
\end{Verbatim}
 

 As we know already base and size, this second stabilizer chain calculation is
much faster than the first one and takes less memory. }

 }

            
\chapter{\textcolor{Chapter }{Matrix Groups}}\label{Matrix Groups}
\logpage{[ 44, 0, 0 ]}
\hyperdef{L}{X7CF51CB48610A07D}{}
{
  Matrix groups are groups generated by invertible square matrices. 

 In the following example we temporarily increase the line length limit from
its default value 80 to 83 in order to get a nicer output format. 

 
\begin{Verbatim}[commandchars=!@|,fontsize=\small,frame=single,label=Example]
  !gapprompt@gap>| !gapinput@m1 := [ [ Z(3)^0, Z(3)^0,   Z(3) ],|
  !gapprompt@>| !gapinput@           [   Z(3), 0*Z(3),   Z(3) ],|
  !gapprompt@>| !gapinput@           [ 0*Z(3),   Z(3), 0*Z(3) ] ];;|
  !gapprompt@gap>| !gapinput@m2 := [ [   Z(3),   Z(3), Z(3)^0 ],|
  !gapprompt@>| !gapinput@           [   Z(3), 0*Z(3),   Z(3) ],|
  !gapprompt@>| !gapinput@           [ Z(3)^0, 0*Z(3),   Z(3) ] ];;|
  !gapprompt@gap>| !gapinput@m := Group( m1, m2 );|
  Group(
  [
    [ [ Z(3)^0, Z(3)^0, Z(3) ], [ Z(3), 0*Z(3), Z(3) ],
        [ 0*Z(3), Z(3), 0*Z(3) ] ],
    [ [ Z(3), Z(3), Z(3)^0 ], [ Z(3), 0*Z(3), Z(3) ],
        [ Z(3)^0, 0*Z(3), Z(3) ] ] ])
\end{Verbatim}
  
\section{\textcolor{Chapter }{IsMatrixGroup (Filter)}}\label{sect:IsMatrixGroup}
\logpage{[ 44, 1, 0 ]}
\hyperdef{L}{X86CEA60E7C04744C}{}
{
  For most operations, \textsf{GAP} only provides methods for finite matrix groups. Many calculations in finite
matrix groups are done via so\texttt{\symbol{45}}called ``nice monomorphisms'' (see Section \ref{Nice Monomorphisms}) that represent a faithful action on vectors. 

\subsection{\textcolor{Chapter }{IsMatrixGroup}}
\logpage{[ 44, 1, 1 ]}\nobreak
\hyperdef{L}{X7E6093FF85F1C3A1}{}
{\noindent\textcolor{FuncColor}{$\triangleright$\enspace\texttt{IsMatrixGroup({\mdseries\slshape grp})\index{IsMatrixGroup@\texttt{IsMatrixGroup}}
\label{IsMatrixGroup}
}\hfill{\scriptsize (Category)}}\\


 The category of matrix groups. }

 }

  
\section{\textcolor{Chapter }{Attributes and Properties for Matrix Groups}}\label{Attributes and Properties for Matrix Groups}
\logpage{[ 44, 2, 0 ]}
\hyperdef{L}{X7FD808E386FAF9B0}{}
{
  

\subsection{\textcolor{Chapter }{DimensionOfMatrixGroup}}
\logpage{[ 44, 2, 1 ]}\nobreak
\hyperdef{L}{X7E55258C783C50CA}{}
{\noindent\textcolor{FuncColor}{$\triangleright$\enspace\texttt{DimensionOfMatrixGroup({\mdseries\slshape mat\texttt{\symbol{45}}grp})\index{DimensionOfMatrixGroup@\texttt{DimensionOfMatrixGroup}}
\label{DimensionOfMatrixGroup}
}\hfill{\scriptsize (attribute)}}\\


 The dimension of the matrix group. }

 

\subsection{\textcolor{Chapter }{DefaultFieldOfMatrixGroup}}
\logpage{[ 44, 2, 2 ]}\nobreak
\hyperdef{L}{X7D540083793CD496}{}
{\noindent\textcolor{FuncColor}{$\triangleright$\enspace\texttt{DefaultFieldOfMatrixGroup({\mdseries\slshape mat\texttt{\symbol{45}}grp})\index{DefaultFieldOfMatrixGroup@\texttt{DefaultFieldOfMatrixGroup}}
\label{DefaultFieldOfMatrixGroup}
}\hfill{\scriptsize (attribute)}}\\


 Is a field containing all the matrix entries. It is not guaranteed to be the
smallest field with this property. }

 

\subsection{\textcolor{Chapter }{FieldOfMatrixGroup}}
\logpage{[ 44, 2, 3 ]}\nobreak
\hyperdef{L}{X78A9F0E580DA613A}{}
{\noindent\textcolor{FuncColor}{$\triangleright$\enspace\texttt{FieldOfMatrixGroup({\mdseries\slshape matgrp})\index{FieldOfMatrixGroup@\texttt{FieldOfMatrixGroup}}
\label{FieldOfMatrixGroup}
}\hfill{\scriptsize (attribute)}}\\


 The smallest field containing all the matrix entries of all elements of the
matrix group \mbox{\texttt{\mdseries\slshape matgrp}}. As the calculation of this can be hard, this should only be used if one \emph{really} needs the smallest field, use \texttt{DefaultFieldOfMatrixGroup} (\ref{DefaultFieldOfMatrixGroup}) to get (for example) the characteristic. 
\begin{Verbatim}[commandchars=!@|,fontsize=\small,frame=single,label=Example]
  !gapprompt@gap>| !gapinput@DimensionOfMatrixGroup(m);|
  3
  !gapprompt@gap>| !gapinput@DefaultFieldOfMatrixGroup(m);|
  GF(3)
\end{Verbatim}
 }

 

\subsection{\textcolor{Chapter }{TransposedMatrixGroup}}
\logpage{[ 44, 2, 4 ]}\nobreak
\hyperdef{L}{X832D18C77ED608DE}{}
{\noindent\textcolor{FuncColor}{$\triangleright$\enspace\texttt{TransposedMatrixGroup({\mdseries\slshape matgrp})\index{TransposedMatrixGroup@\texttt{TransposedMatrixGroup}}
\label{TransposedMatrixGroup}
}\hfill{\scriptsize (attribute)}}\\


 returns the transpose of the matrix group \mbox{\texttt{\mdseries\slshape matgrp}}. The transpose of the transpose of \mbox{\texttt{\mdseries\slshape matgrp}} is identical to \mbox{\texttt{\mdseries\slshape matgrp}}. 
\begin{Verbatim}[commandchars=!@|,fontsize=\small,frame=single,label=Example]
  !gapprompt@gap>| !gapinput@G := Group( [[0,-1],[1,0]] );|
  Group([ [ [ 0, -1 ], [ 1, 0 ] ] ])
  !gapprompt@gap>| !gapinput@T := TransposedMatrixGroup( G );|
  Group([ [ [ 0, 1 ], [ -1, 0 ] ] ])
  !gapprompt@gap>| !gapinput@IsIdenticalObj( G, TransposedMatrixGroup( T ) );|
  true
\end{Verbatim}
 }

 

\subsection{\textcolor{Chapter }{IsFFEMatrixGroup}}
\logpage{[ 44, 2, 5 ]}\nobreak
\hyperdef{L}{X84B36A827E5EFC35}{}
{\noindent\textcolor{FuncColor}{$\triangleright$\enspace\texttt{IsFFEMatrixGroup({\mdseries\slshape G})\index{IsFFEMatrixGroup@\texttt{IsFFEMatrixGroup}}
\label{IsFFEMatrixGroup}
}\hfill{\scriptsize (Category)}}\\


 tests whether all matrices in \mbox{\texttt{\mdseries\slshape G}} have finite field element entries. }

 }

  
\section{\textcolor{Chapter }{Actions of Matrix Groups}}\label{Actions of Matrix Groups}
\logpage{[ 44, 3, 0 ]}
\hyperdef{L}{X7F4B0B397AAC7659}{}
{
  The basic operations for groups are described in Chapter{\nobreakspace}\ref{Group Actions}, special actions for \emph{matrix} groups mentioned there are \texttt{OnLines} (\ref{OnLines}), \texttt{OnRight} (\ref{OnRight}), and \texttt{OnSubspacesByCanonicalBasis} (\ref{OnSubspacesByCanonicalBasis}).  

 For subtleties concerning multiplication from the left or from the right,
see{\nobreakspace}\ref{Acting OnRight and OnLeft}. 

\subsection{\textcolor{Chapter }{ProjectiveActionOnFullSpace}}
\logpage{[ 44, 3, 1 ]}\nobreak
\hyperdef{L}{X7BD4F38E8624735D}{}
{\noindent\textcolor{FuncColor}{$\triangleright$\enspace\texttt{ProjectiveActionOnFullSpace({\mdseries\slshape G, F, n})\index{ProjectiveActionOnFullSpace@\texttt{ProjectiveActionOnFullSpace}}
\label{ProjectiveActionOnFullSpace}
}\hfill{\scriptsize (function)}}\\


 Let \mbox{\texttt{\mdseries\slshape G}} be a group of \mbox{\texttt{\mdseries\slshape n}} by \mbox{\texttt{\mdseries\slshape n}} matrices over a field contained in the finite field \mbox{\texttt{\mdseries\slshape F}}.   \texttt{ProjectiveActionOnFullSpace} returns the image of the projective action of \mbox{\texttt{\mdseries\slshape G}} on the full row space $\mbox{\texttt{\mdseries\slshape F}}^{\mbox{\texttt{\mdseries\slshape n}}}$. }

 

\subsection{\textcolor{Chapter }{ProjectiveActionHomomorphismMatrixGroup}}
\logpage{[ 44, 3, 2 ]}\nobreak
\hyperdef{L}{X7F8EA8D583C1E9B2}{}
{\noindent\textcolor{FuncColor}{$\triangleright$\enspace\texttt{ProjectiveActionHomomorphismMatrixGroup({\mdseries\slshape G})\index{ProjectiveActionHomomorphismMatrixGroup@\texttt{Projective}\-\texttt{Action}\-\texttt{Homomorphism}\-\texttt{Matrix}\-\texttt{Group}}
\label{ProjectiveActionHomomorphismMatrixGroup}
}\hfill{\scriptsize (function)}}\\


 returns an action homomorphism for a faithful projective action of \mbox{\texttt{\mdseries\slshape G}} on the underlying vector space. (Note: The action is not necessarily on the
full space, if a smaller subset can be found on which the action is faithful.) }

 

\subsection{\textcolor{Chapter }{BlowUpIsomorphism}}
\logpage{[ 44, 3, 3 ]}\nobreak
\hyperdef{L}{X849C451A80B4A210}{}
{\noindent\textcolor{FuncColor}{$\triangleright$\enspace\texttt{BlowUpIsomorphism({\mdseries\slshape matgrp, B})\index{BlowUpIsomorphism@\texttt{BlowUpIsomorphism}}
\label{BlowUpIsomorphism}
}\hfill{\scriptsize (function)}}\\


 For a matrix group \mbox{\texttt{\mdseries\slshape matgrp}} and a basis \mbox{\texttt{\mdseries\slshape B}} of a field extension $L / K$ such that the entries of all matrices in \mbox{\texttt{\mdseries\slshape matgrp}} lie in $L$, \texttt{BlowUpIsomorphism} returns the isomorphism with source \mbox{\texttt{\mdseries\slshape matgrp}} that is defined by mapping the matrix $A$ to \texttt{BlownUpMat}$( A, \mbox{\texttt{\mdseries\slshape B}} )$, see{\nobreakspace}\texttt{BlownUpMat} (\ref{BlownUpMat}). 
\begin{Verbatim}[commandchars=!@|,fontsize=\small,frame=single,label=Example]
  !gapprompt@gap>| !gapinput@g:= GL(2,4);;|
  !gapprompt@gap>| !gapinput@B:= CanonicalBasis( GF(4) );;  BasisVectors( B );|
  [ Z(2)^0, Z(2^2) ]
  !gapprompt@gap>| !gapinput@iso:= BlowUpIsomorphism( g, B );;|
  !gapprompt@gap>| !gapinput@Display( Image( iso, [ [ Z(4), Z(2) ], [ 0*Z(2), Z(4)^2 ] ] ) );|
   . 1 1 .
   1 1 . 1
   . . 1 1
   . . 1 .
  !gapprompt@gap>| !gapinput@img:= Image( iso, g );|
  <matrix group with 2 generators>
  !gapprompt@gap>| !gapinput@Index( GL(4,2), img );|
  112
\end{Verbatim}
 }

 }

  
\section{\textcolor{Chapter }{GL and SL}}\label{GL and SL}
\logpage{[ 44, 4, 0 ]}
\hyperdef{L}{X7934EED77891BE6B}{}
{
  (See also section{\nobreakspace}\ref{Classical Groups}.) 

\subsection{\textcolor{Chapter }{IsGeneralLinearGroup}}
\logpage{[ 44, 4, 1 ]}\nobreak
\hyperdef{L}{X781387AF7999EA99}{}
{\noindent\textcolor{FuncColor}{$\triangleright$\enspace\texttt{IsGeneralLinearGroup({\mdseries\slshape grp})\index{IsGeneralLinearGroup@\texttt{IsGeneralLinearGroup}}
\label{IsGeneralLinearGroup}
}\hfill{\scriptsize (property)}}\\
\noindent\textcolor{FuncColor}{$\triangleright$\enspace\texttt{IsGL({\mdseries\slshape grp})\index{IsGL@\texttt{IsGL}}
\label{IsGL}
}\hfill{\scriptsize (property)}}\\


 The General Linear group is the group of all invertible matrices over a ring.
This property tests, whether a group is isomorphic to a General Linear group.
(Note that currently only a few trivial methods are available for this
operation. We hope to improve this in the future.) }

 

\subsection{\textcolor{Chapter }{IsNaturalGL}}
\logpage{[ 44, 4, 2 ]}\nobreak
\hyperdef{L}{X86F9A27D7AFAEB5A}{}
{\noindent\textcolor{FuncColor}{$\triangleright$\enspace\texttt{IsNaturalGL({\mdseries\slshape matgrp})\index{IsNaturalGL@\texttt{IsNaturalGL}}
\label{IsNaturalGL}
}\hfill{\scriptsize (property)}}\\


 This property tests, whether a matrix group is the General Linear group in the
right dimension over the (smallest) ring which contains all entries of its
elements. (Currently, only a trivial test that computes the order of the group
is available.) }

 

\subsection{\textcolor{Chapter }{IsSpecialLinearGroup}}
\logpage{[ 44, 4, 3 ]}\nobreak
\hyperdef{L}{X816677CD821261FA}{}
{\noindent\textcolor{FuncColor}{$\triangleright$\enspace\texttt{IsSpecialLinearGroup({\mdseries\slshape grp})\index{IsSpecialLinearGroup@\texttt{IsSpecialLinearGroup}}
\label{IsSpecialLinearGroup}
}\hfill{\scriptsize (property)}}\\
\noindent\textcolor{FuncColor}{$\triangleright$\enspace\texttt{IsSL({\mdseries\slshape grp})\index{IsSL@\texttt{IsSL}}
\label{IsSL}
}\hfill{\scriptsize (property)}}\\


 The Special Linear group is the group of all invertible matrices over a ring,
whose determinant is equal to 1. This property tests, whether a group is
isomorphic to a Special Linear group. (Note that currently only a few trivial
methods are available for this operation. We hope to improve this in the
future.) }

 

\subsection{\textcolor{Chapter }{IsNaturalSL}}
\logpage{[ 44, 4, 4 ]}\nobreak
\hyperdef{L}{X84134F08781EB943}{}
{\noindent\textcolor{FuncColor}{$\triangleright$\enspace\texttt{IsNaturalSL({\mdseries\slshape matgrp})\index{IsNaturalSL@\texttt{IsNaturalSL}}
\label{IsNaturalSL}
}\hfill{\scriptsize (property)}}\\


 This property tests, whether a matrix group is the Special Linear group in the
right dimension over the (smallest) ring which contains all entries of its
elements. (Currently, only a trivial test that computes the order of the group
is available.) 
\begin{Verbatim}[commandchars=!@|,fontsize=\small,frame=single,label=Example]
  !gapprompt@gap>| !gapinput@IsNaturalGL(m);|
  false
\end{Verbatim}
 }

 

\subsection{\textcolor{Chapter }{IsSubgroupSL}}
\logpage{[ 44, 4, 5 ]}\nobreak
\hyperdef{L}{X7ED43D4F7E993A31}{}
{\noindent\textcolor{FuncColor}{$\triangleright$\enspace\texttt{IsSubgroupSL({\mdseries\slshape matgrp})\index{IsSubgroupSL@\texttt{IsSubgroupSL}}
\label{IsSubgroupSL}
}\hfill{\scriptsize (property)}}\\


 This property tests, whether a matrix group is a subgroup of the Special
Linear group in the right dimension over the (smallest) ring which contains
all entries of its elements. }

 }

  
\section{\textcolor{Chapter }{Invariant Forms}}\label{Invariant Forms}
\logpage{[ 44, 5, 0 ]}
\hyperdef{L}{X7CA4097C79F5BD90}{}
{
  

\subsection{\textcolor{Chapter }{InvariantBilinearForm}}
\logpage{[ 44, 5, 1 ]}\nobreak
\hyperdef{L}{X7C08385A81AB05E1}{}
{\noindent\textcolor{FuncColor}{$\triangleright$\enspace\texttt{InvariantBilinearForm({\mdseries\slshape matgrp})\index{InvariantBilinearForm@\texttt{InvariantBilinearForm}}
\label{InvariantBilinearForm}
}\hfill{\scriptsize (attribute)}}\\


 For a matrix group \mbox{\texttt{\mdseries\slshape matgrp}}, \texttt{InvariantBilinearForm} returns a record containing at least the components \texttt{matrix}, whose value is a matrix $M$, and \texttt{baseDomain}, whose value is a field. 

 This record describes a bilinear form that is invariant under \mbox{\texttt{\mdseries\slshape matgrp}}. For every generator $g$ of \mbox{\texttt{\mdseries\slshape matgrp}}, the equation $g \cdot M \cdot g^{tr} = M$ holds. 

 The attribute value is set for example in groups constructed with \texttt{SymplecticGroup} (\ref{SymplecticGroup:for dimension and field size}). Together with \texttt{IsFullSubgroupGLorSLRespectingBilinearForm} (\ref{IsFullSubgroupGLorSLRespectingBilinearForm}), it can be used for efficient membership tests in \mbox{\texttt{\mdseries\slshape matgrp}}. }

 

\subsection{\textcolor{Chapter }{IsFullSubgroupGLorSLRespectingBilinearForm}}
\logpage{[ 44, 5, 2 ]}\nobreak
\hyperdef{L}{X8652FBF781940AC3}{}
{\noindent\textcolor{FuncColor}{$\triangleright$\enspace\texttt{IsFullSubgroupGLorSLRespectingBilinearForm({\mdseries\slshape matgrp})\index{IsFullSubgroupGLorSLRespectingBilinearForm@\texttt{IsFull}\-\texttt{Subgroup}\-\texttt{G}\-\texttt{Lor}\-\texttt{S}\-\texttt{L}\-\texttt{Respecting}\-\texttt{Bilinear}\-\texttt{Form}}
\label{IsFullSubgroupGLorSLRespectingBilinearForm}
}\hfill{\scriptsize (property)}}\\


 If this property is set then the matrix group \mbox{\texttt{\mdseries\slshape matgrp}} is the full subgroup of \texttt{GL}$(d, q)$ or \texttt{SL}$(d, q)$ respecting the form stored as the value of \texttt{InvariantBilinearForm} (\ref{InvariantBilinearForm}) for \mbox{\texttt{\mdseries\slshape matgrp}}. 

 The value $d$ is the \texttt{DimensionOfMatrixGroup} (\ref{DimensionOfMatrixGroup}) value of \mbox{\texttt{\mdseries\slshape matgrp}}, and $q$ is the size of the \texttt{baseDomain} component of the record stored as \texttt{InvariantBilinearForm} (\ref{InvariantBilinearForm}). 

 The property \texttt{IsSubgroupSL} (\ref{IsSubgroupSL}) of \mbox{\texttt{\mdseries\slshape matgrp}} determines whether \texttt{GL} or \texttt{SL} is meant for \texttt{IsFullSubgroupGLorSLRespectingBilinearForm}. }

 

\subsection{\textcolor{Chapter }{InvariantBilinearFormUpToScalars}}
\logpage{[ 44, 5, 3 ]}\nobreak
\hyperdef{L}{X78B2705B87ED5384}{}
{\noindent\textcolor{FuncColor}{$\triangleright$\enspace\texttt{InvariantBilinearFormUpToScalars({\mdseries\slshape matgrp})\index{InvariantBilinearFormUpToScalars@\texttt{InvariantBilinearFormUpToScalars}}
\label{InvariantBilinearFormUpToScalars}
}\hfill{\scriptsize (attribute)}}\\


 This attribute describes a bilinear form that is invariant up to scalars under
the matrix group \mbox{\texttt{\mdseries\slshape matgrp}}. The form is given by a record with the component \texttt{matrix} which is a matrix $J$ such that for every generator $g$ of \mbox{\texttt{\mdseries\slshape matgrp}} the equation $g \cdot J \cdot g^{tr} = \lambda(g) J$ holds, for $\lambda(g)$ in the \texttt{FieldOfMatrixGroup} (\ref{FieldOfMatrixGroup}) value of \mbox{\texttt{\mdseries\slshape matgrp}}. }

 

\subsection{\textcolor{Chapter }{IsFullSubgroupGLRespectingBilinearFormUpToScalars}}
\logpage{[ 44, 5, 4 ]}\nobreak
\hyperdef{L}{X8461A2D87E4208EF}{}
{\noindent\textcolor{FuncColor}{$\triangleright$\enspace\texttt{IsFullSubgroupGLRespectingBilinearFormUpToScalars({\mdseries\slshape matgrp})\index{IsFullSubgroupGLRespectingBilinearFormUpToScalars@\texttt{IsFull}\-\texttt{Subgroup}\-\texttt{G}\-\texttt{L}\-\texttt{Respecting}\-\texttt{Bilinear}\-\texttt{Form}\-\texttt{Up}\-\texttt{To}\-\texttt{Scalars}}
\label{IsFullSubgroupGLRespectingBilinearFormUpToScalars}
}\hfill{\scriptsize (property)}}\\


 This property tests whether the matrix group \mbox{\texttt{\mdseries\slshape matgrp}} is the full subgroup of GL respecting, up to scalars, the form stored as the
value of \texttt{InvariantBilinearFormUpToScalars} (\ref{InvariantBilinearFormUpToScalars}) for \mbox{\texttt{\mdseries\slshape matgrp}}. }

 

\subsection{\textcolor{Chapter }{InvariantSesquilinearForm}}
\logpage{[ 44, 5, 5 ]}\nobreak
\hyperdef{L}{X82F22079852130C9}{}
{\noindent\textcolor{FuncColor}{$\triangleright$\enspace\texttt{InvariantSesquilinearForm({\mdseries\slshape matgrp})\index{InvariantSesquilinearForm@\texttt{InvariantSesquilinearForm}}
\label{InvariantSesquilinearForm}
}\hfill{\scriptsize (attribute)}}\\


 For a matrix group \mbox{\texttt{\mdseries\slshape matgrp}}, \texttt{InvariantSesquilinearForm} returns a record containing at least the components \texttt{matrix}, whose value is a matrix $M$, and \texttt{baseDomain}, whose value is a field $F$. 

 This record describes a sesquilinear form that is invariant under \mbox{\texttt{\mdseries\slshape matgrp}}. For every generator $g$ of \mbox{\texttt{\mdseries\slshape matgrp}}, the equation $g \cdot M \cdot (g^{tr})^f = M$ holds, where $f$ is an automorphism of $F$. If $F$ has $q^2$ elements then $f$ raises each element of $F$ to its $q$\texttt{\symbol{45}}th power. ($f$ can be obtained as a power of the \texttt{FrobeniusAutomorphism} (\ref{FrobeniusAutomorphism}) value of $F$.) 

 The attribute value is set for example in groups constructed with \texttt{GeneralUnitaryGroup} (\ref{GeneralUnitaryGroup}). Together with \texttt{IsFullSubgroupGLorSLRespectingSesquilinearForm} (\ref{IsFullSubgroupGLorSLRespectingSesquilinearForm}), it can be used for efficient membership tests in \mbox{\texttt{\mdseries\slshape matgrp}}. }

 

\subsection{\textcolor{Chapter }{IsFullSubgroupGLorSLRespectingSesquilinearForm}}
\logpage{[ 44, 5, 6 ]}\nobreak
\hyperdef{L}{X7B35A8AF7D8F0313}{}
{\noindent\textcolor{FuncColor}{$\triangleright$\enspace\texttt{IsFullSubgroupGLorSLRespectingSesquilinearForm({\mdseries\slshape matgrp})\index{IsFullSubgroupGLorSLRespectingSesquilinearForm@\texttt{IsFull}\-\texttt{Subgroup}\-\texttt{G}\-\texttt{Lor}\-\texttt{S}\-\texttt{L}\-\texttt{Respecting}\-\texttt{Sesquilinear}\-\texttt{Form}}
\label{IsFullSubgroupGLorSLRespectingSesquilinearForm}
}\hfill{\scriptsize (property)}}\\


 If this property is set then the matrix group \mbox{\texttt{\mdseries\slshape matgrp}} is the full subgroup of \texttt{GL}$(d, q)$ or \texttt{SL}$(d, q)$ respecting the form stored as the value of \texttt{InvariantSesquilinearForm} (\ref{InvariantSesquilinearForm}) for \mbox{\texttt{\mdseries\slshape matgrp}}. 

 The value $d$ is the \texttt{DimensionOfMatrixGroup} (\ref{DimensionOfMatrixGroup}) value of \mbox{\texttt{\mdseries\slshape matgrp}}, and $q$ is the size of the \texttt{baseDomain} component of the record stored as \texttt{InvariantSesquilinearForm} (\ref{InvariantSesquilinearForm}). 

 The property \texttt{IsSubgroupSL} (\ref{IsSubgroupSL}) of \mbox{\texttt{\mdseries\slshape matgrp}} determines whether \texttt{GL} or \texttt{SL} is meant for \texttt{IsFullSubgroupGLorSLRespectingSesquilinearForm}. }

 

\subsection{\textcolor{Chapter }{InvariantQuadraticForm}}
\logpage{[ 44, 5, 7 ]}\nobreak
\hyperdef{L}{X7BCACC007EB9B613}{}
{\noindent\textcolor{FuncColor}{$\triangleright$\enspace\texttt{InvariantQuadraticForm({\mdseries\slshape matgrp})\index{InvariantQuadraticForm@\texttt{InvariantQuadraticForm}}
\label{InvariantQuadraticForm}
}\hfill{\scriptsize (attribute)}}\\


 For a matrix group \mbox{\texttt{\mdseries\slshape matgrp}}, \texttt{InvariantQuadraticForm} returns a record containing at least the components \texttt{matrix}, whose value is a matrix $Q$, and \texttt{baseDomain}, whose value is a field $F$. 

 The quadratic form $q$ on the natural $F$\texttt{\symbol{45}}vector space $V$ on which \mbox{\texttt{\mdseries\slshape matgrp}} acts is given by $q(v) = v Q v^{tr}$, and the invariance under \mbox{\texttt{\mdseries\slshape matgrp}} is given by the equation $q(v) = q(v M)$ for all $v \in V$ and $M$ in \mbox{\texttt{\mdseries\slshape matgrp}}. (Note that the invariance of $q$ does \emph{not} imply that the matrix $Q$ is invariant under \mbox{\texttt{\mdseries\slshape matgrp}}.) 

 The function $q$ is defined relative to an invariant symmetric bilinear form $b$ (see{\nobreakspace}\texttt{InvariantBilinearForm} (\ref{InvariantBilinearForm})), via the equation $q(\lambda x + \mu y) = \lambda^2 q(x) + \lambda \mu b(x,y) + \mu^2 q(y)$, see \cite[Chapter 3.4]{CCN85}. If $b$ is represented by the matrix $B$ then this implies $B = Q + Q^{tr}$. In characteristic different from $2$, we have $q(x) = b(x,x)/2$, so $Q$ can be chosen as the strictly upper triangular part of $B$ plus half of the diagonal part of $B$. In characteristic $2$, $B$ does not determine $Q$ but still $Q$ can be chosen as an upper (or lower) triangular matrix. 

 Whenever the \texttt{InvariantQuadraticForm} value is set in a matrix group then also the \texttt{InvariantBilinearForm} (\ref{InvariantBilinearForm}) value can be accessed, and the two values are compatible in the above sense. 

 The attribute value is set for example in groups constructed with \texttt{GeneralOrthogonalGroup} (\ref{GeneralOrthogonalGroup}). Together with \texttt{IsFullSubgroupGLorSLRespectingQuadraticForm} (\ref{IsFullSubgroupGLorSLRespectingQuadraticForm}), it can be used for efficient membership tests in \mbox{\texttt{\mdseries\slshape matgrp}}. }

 

\subsection{\textcolor{Chapter }{IsFullSubgroupGLorSLRespectingQuadraticForm}}
\logpage{[ 44, 5, 8 ]}\nobreak
\hyperdef{L}{X84AB04A67DFC0274}{}
{\noindent\textcolor{FuncColor}{$\triangleright$\enspace\texttt{IsFullSubgroupGLorSLRespectingQuadraticForm({\mdseries\slshape matgrp})\index{IsFullSubgroupGLorSLRespectingQuadraticForm@\texttt{IsFull}\-\texttt{Subgroup}\-\texttt{G}\-\texttt{Lor}\-\texttt{S}\-\texttt{L}\-\texttt{Respecting}\-\texttt{Quadratic}\-\texttt{Form}}
\label{IsFullSubgroupGLorSLRespectingQuadraticForm}
}\hfill{\scriptsize (property)}}\\


 If this property is set then the matrix group \mbox{\texttt{\mdseries\slshape matgrp}} is the full subgroup of \texttt{GL}$(d, q)$ or \texttt{SL}$(d, q)$ respecting the form stored as the value of \texttt{InvariantQuadraticForm} (\ref{InvariantQuadraticForm}) for \mbox{\texttt{\mdseries\slshape matgrp}}. 

 The value $d$ is the \texttt{DimensionOfMatrixGroup} (\ref{DimensionOfMatrixGroup}) value of \mbox{\texttt{\mdseries\slshape matgrp}}, and $q$ is the size of the \texttt{baseDomain} component of the record stored as \texttt{InvariantQuadraticForm} (\ref{InvariantQuadraticForm}). 

 The property \texttt{IsSubgroupSL} (\ref{IsSubgroupSL}) of \mbox{\texttt{\mdseries\slshape matgrp}} determines whether \texttt{GL} or \texttt{SL} is meant for \texttt{IsFullSubgroupGLorSLRespectingQuadraticForm}. 

 
\begin{Verbatim}[commandchars=!@|,fontsize=\small,frame=single,label=Example]
  !gapprompt@gap>| !gapinput@g:= Sp( 2, 3 );;|
  !gapprompt@gap>| !gapinput@m:= InvariantBilinearForm( g ).matrix;|
  [ [ 0*Z(3), Z(3)^0 ], [ Z(3), 0*Z(3) ] ]
  !gapprompt@gap>| !gapinput@[ 0, 1 ] * m * [ 1, -1 ];       # evaluate the bilinear form|
  Z(3)
  !gapprompt@gap>| !gapinput@IsFullSubgroupGLorSLRespectingBilinearForm( g );|
  true
  !gapprompt@gap>| !gapinput@g:= SU( 2, 4 );;|
  !gapprompt@gap>| !gapinput@m:= InvariantSesquilinearForm( g ).matrix;|
  [ [ 0*Z(2), Z(2)^0 ], [ Z(2)^0, 0*Z(2) ] ]
  !gapprompt@gap>| !gapinput@[ 0, 1 ] * m * [ 1, 1 ];        # evaluate the sesquilinear form|
  Z(2)^0
  !gapprompt@gap>| !gapinput@IsFullSubgroupGLorSLRespectingSesquilinearForm( g );|
  true
  !gapprompt@gap>| !gapinput@IsSubgroupSL( g );|
  true
  !gapprompt@gap>| !gapinput@g:= GO( 1, 2, 3 );;|
  !gapprompt@gap>| !gapinput@m:= InvariantBilinearForm( g ).matrix;|
  [ [ 0*Z(3), Z(3)^0 ], [ Z(3)^0, 0*Z(3) ] ]
  !gapprompt@gap>| !gapinput@[ 0, 1 ] * m * [ 1, 1 ];        # evaluate the bilinear form|
  Z(3)^0
  !gapprompt@gap>| !gapinput@q:= InvariantQuadraticForm( g ).matrix;|
  [ [ 0*Z(3), Z(3)^0 ], [ 0*Z(3), 0*Z(3) ] ]
  !gapprompt@gap>| !gapinput@[ 0, 1 ] * q * [ 0, 1 ];        # evaluate the quadratic form|
  0*Z(3)
  !gapprompt@gap>| !gapinput@IsFullSubgroupGLorSLRespectingQuadraticForm( g );|
  true
\end{Verbatim}
 }

 }

  
\section{\textcolor{Chapter }{Matrix Groups in Characteristic 0}}\label{Matrix Groups in Characteristic 0}
\logpage{[ 44, 6, 0 ]}
\hyperdef{L}{X7FB0138F79E8C5E7}{}
{
  Most of the functions described in this and the following section have
implementations which use functions from the \textsf{GAP} package \textsf{CaratInterface}. If \textsf{CaratInterface} is not installed or not compiled, no suitable methods are available. 

\subsection{\textcolor{Chapter }{IsCyclotomicMatrixGroup}}
\logpage{[ 44, 6, 1 ]}\nobreak
\hyperdef{L}{X850821F78558C829}{}
{\noindent\textcolor{FuncColor}{$\triangleright$\enspace\texttt{IsCyclotomicMatrixGroup({\mdseries\slshape G})\index{IsCyclotomicMatrixGroup@\texttt{IsCyclotomicMatrixGroup}}
\label{IsCyclotomicMatrixGroup}
}\hfill{\scriptsize (Category)}}\\


 tests whether all matrices in \mbox{\texttt{\mdseries\slshape G}} have cyclotomic entries. }

 

\subsection{\textcolor{Chapter }{IsRationalMatrixGroup}}
\logpage{[ 44, 6, 2 ]}\nobreak
\hyperdef{L}{X7FEDB2E17EE02674}{}
{\noindent\textcolor{FuncColor}{$\triangleright$\enspace\texttt{IsRationalMatrixGroup({\mdseries\slshape G})\index{IsRationalMatrixGroup@\texttt{IsRationalMatrixGroup}}
\label{IsRationalMatrixGroup}
}\hfill{\scriptsize (property)}}\\


 tests whether all matrices in \mbox{\texttt{\mdseries\slshape G}} have rational entries. }

 

\subsection{\textcolor{Chapter }{IsIntegerMatrixGroup}}
\logpage{[ 44, 6, 3 ]}\nobreak
\hyperdef{L}{X7F737FC4795F3E48}{}
{\noindent\textcolor{FuncColor}{$\triangleright$\enspace\texttt{IsIntegerMatrixGroup({\mdseries\slshape G})\index{IsIntegerMatrixGroup@\texttt{IsIntegerMatrixGroup}}
\label{IsIntegerMatrixGroup}
}\hfill{\scriptsize (property)}}\\


 tests whether all matrices in \mbox{\texttt{\mdseries\slshape G}} have integer entries. 
\begin{Verbatim}[commandchars=!@|,fontsize=\small,frame=single,label=Example]
  !gapprompt@gap>| !gapinput@A:=[[0,1,0],[0,0,1],[1,0,0]];;|
  !gapprompt@gap>| !gapinput@B:=[[0,0,1],[0,1,0],[-1,0,0]];;|
  !gapprompt@gap>| !gapinput@C:=[[E(4),0,0],[0,E(4)^(-1),0],[0,0,1]];;|
  !gapprompt@gap>| !gapinput@IsIntegerMatrixGroup(Group(A, B));|
  true
  !gapprompt@gap>| !gapinput@IsIntegerMatrixGroup(Group(A, C));|
  false
\end{Verbatim}
   }

 

\subsection{\textcolor{Chapter }{IsNaturalGLnZ}}
\logpage{[ 44, 6, 4 ]}\nobreak
\hyperdef{L}{X86F9CC1E7DB97CB6}{}
{\noindent\textcolor{FuncColor}{$\triangleright$\enspace\texttt{IsNaturalGLnZ({\mdseries\slshape G})\index{IsNaturalGLnZ@\texttt{IsNaturalGLnZ}}
\label{IsNaturalGLnZ}
}\hfill{\scriptsize (property)}}\\


 tests whether \mbox{\texttt{\mdseries\slshape G}} is $GL_n({\ensuremath{\mathbb Z}})$ in its natural representation by $n \times n$ integer matrices. (The dimension $n$ will be read off the generating matrices.) 
\begin{Verbatim}[commandchars=!@|,fontsize=\small,frame=single,label=Example]
  !gapprompt@gap>| !gapinput@IsNaturalGLnZ( GL( 2, Integers ) );|
  true
\end{Verbatim}
 }

 

\subsection{\textcolor{Chapter }{IsNaturalSLnZ}}
\logpage{[ 44, 6, 5 ]}\nobreak
\hyperdef{L}{X7B0E70127F5D2EAF}{}
{\noindent\textcolor{FuncColor}{$\triangleright$\enspace\texttt{IsNaturalSLnZ({\mdseries\slshape G})\index{IsNaturalSLnZ@\texttt{IsNaturalSLnZ}}
\label{IsNaturalSLnZ}
}\hfill{\scriptsize (property)}}\\


 tests whether \mbox{\texttt{\mdseries\slshape G}} is $SL_n({\ensuremath{\mathbb Z}})$ in its natural representation by $n \times n$ integer matrices. (The dimension $n$ will be read off the generating matrices.) 
\begin{Verbatim}[commandchars=!@|,fontsize=\small,frame=single,label=Example]
  !gapprompt@gap>| !gapinput@IsNaturalSLnZ( SL( 2, Integers ) );|
  true
\end{Verbatim}
 }

 

\subsection{\textcolor{Chapter }{InvariantLattice}}
\logpage{[ 44, 6, 6 ]}\nobreak
\hyperdef{L}{X7DE412A37A6975B3}{}
{\noindent\textcolor{FuncColor}{$\triangleright$\enspace\texttt{InvariantLattice({\mdseries\slshape G})\index{InvariantLattice@\texttt{InvariantLattice}}
\label{InvariantLattice}
}\hfill{\scriptsize (attribute)}}\\


 returns a matrix $B$, whose rows form a basis of a ${\ensuremath{\mathbb Z}}$\texttt{\symbol{45}}lattice that is invariant under the rational matrix group \mbox{\texttt{\mdseries\slshape G}} acting from the right. It returns \texttt{fail} if the group is not unimodular. The columns of the inverse of $B$ span a ${\ensuremath{\mathbb Z}}$\texttt{\symbol{45}}lattice invariant under \mbox{\texttt{\mdseries\slshape G}} acting from the left. }

 

\subsection{\textcolor{Chapter }{NormalizerInGLnZ}}
\logpage{[ 44, 6, 7 ]}\nobreak
\hyperdef{L}{X7CC4D6DC81739698}{}
{\noindent\textcolor{FuncColor}{$\triangleright$\enspace\texttt{NormalizerInGLnZ({\mdseries\slshape G})\index{NormalizerInGLnZ@\texttt{NormalizerInGLnZ}}
\label{NormalizerInGLnZ}
}\hfill{\scriptsize (attribute)}}\\


 is an attribute used to store the normalizer of \mbox{\texttt{\mdseries\slshape G}} in $GL_n({\ensuremath{\mathbb Z}})$, where \mbox{\texttt{\mdseries\slshape G}} is an integer matrix group of dimension \mbox{\texttt{\mdseries\slshape n}}. This attribute is used by \texttt{Normalizer( GL( n, Integers ), G )}. }

 

\subsection{\textcolor{Chapter }{CentralizerInGLnZ}}
\logpage{[ 44, 6, 8 ]}\nobreak
\hyperdef{L}{X7DAFB71F86525DE7}{}
{\noindent\textcolor{FuncColor}{$\triangleright$\enspace\texttt{CentralizerInGLnZ({\mdseries\slshape G})\index{CentralizerInGLnZ@\texttt{CentralizerInGLnZ}}
\label{CentralizerInGLnZ}
}\hfill{\scriptsize (attribute)}}\\


 is an attribute used to store the centralizer of \mbox{\texttt{\mdseries\slshape G}} in $GL_n({\ensuremath{\mathbb Z}})$, where \mbox{\texttt{\mdseries\slshape G}} is an integer matrix group of dimension \mbox{\texttt{\mdseries\slshape n}}. This attribute is used by \texttt{Centralizer( GL( n, Integers ), G )}. }

 

\subsection{\textcolor{Chapter }{ZClassRepsQClass}}
\logpage{[ 44, 6, 9 ]}\nobreak
\hyperdef{L}{X8217762A863F1382}{}
{\noindent\textcolor{FuncColor}{$\triangleright$\enspace\texttt{ZClassRepsQClass({\mdseries\slshape G})\index{ZClassRepsQClass@\texttt{ZClassRepsQClass}}
\label{ZClassRepsQClass}
}\hfill{\scriptsize (attribute)}}\\


 The conjugacy class in $GL_n({\ensuremath{\mathbb Q}})$ of the finite integer matrix group \mbox{\texttt{\mdseries\slshape G}} splits into finitely many conjugacy classes in $GL_n({\ensuremath{\mathbb Z}})$. \texttt{ZClassRepsQClass( \mbox{\texttt{\mdseries\slshape G}} )} returns representative groups for these. }

 

\subsection{\textcolor{Chapter }{IsBravaisGroup}}
\logpage{[ 44, 6, 10 ]}\nobreak
\hyperdef{L}{X84FD9FC97FB90795}{}
{\noindent\textcolor{FuncColor}{$\triangleright$\enspace\texttt{IsBravaisGroup({\mdseries\slshape G})\index{IsBravaisGroup@\texttt{IsBravaisGroup}}
\label{IsBravaisGroup}
}\hfill{\scriptsize (property)}}\\


 test whether \mbox{\texttt{\mdseries\slshape G}} coincides with its Bravais group (see \texttt{BravaisGroup} (\ref{BravaisGroup})). }

 

\subsection{\textcolor{Chapter }{BravaisGroup}}
\logpage{[ 44, 6, 11 ]}\nobreak
\hyperdef{L}{X7AAE301C83116451}{}
{\noindent\textcolor{FuncColor}{$\triangleright$\enspace\texttt{BravaisGroup({\mdseries\slshape G})\index{BravaisGroup@\texttt{BravaisGroup}}
\label{BravaisGroup}
}\hfill{\scriptsize (attribute)}}\\


 returns the Bravais group of a finite integer matrix group \mbox{\texttt{\mdseries\slshape G}}. If $C$ is the cone of positive definite quadratic forms $Q$ invariant under $g \mapsto g Q g^{tr}$ for all $g \in \mbox{\texttt{\mdseries\slshape G}}$, then the Bravais group of \mbox{\texttt{\mdseries\slshape G}} is the maximal subgroup of $GL_n({\ensuremath{\mathbb Z}})$ leaving the forms in that same cone invariant. Alternatively, the Bravais
group of \mbox{\texttt{\mdseries\slshape G}} can also be defined with respect to the action $g \mapsto g^{tr} Q g$ on positive definite quadratic forms $Q$. This latter definition is appropriate for groups \mbox{\texttt{\mdseries\slshape G}} acting from the right on row vectors, whereas the former definition is
appropriate for groups acting from the left on column vectors. Both
definitions yield the same Bravais group. }

 

\subsection{\textcolor{Chapter }{BravaisSubgroups}}
\logpage{[ 44, 6, 12 ]}\nobreak
\hyperdef{L}{X788C7D9C7C2301C5}{}
{\noindent\textcolor{FuncColor}{$\triangleright$\enspace\texttt{BravaisSubgroups({\mdseries\slshape G})\index{BravaisSubgroups@\texttt{BravaisSubgroups}}
\label{BravaisSubgroups}
}\hfill{\scriptsize (attribute)}}\\


 returns the subgroups of the Bravais group of \mbox{\texttt{\mdseries\slshape G}}, which are themselves Bravais groups. }

 

\subsection{\textcolor{Chapter }{BravaisSupergroups}}
\logpage{[ 44, 6, 13 ]}\nobreak
\hyperdef{L}{X7F5FF1A481E08AD5}{}
{\noindent\textcolor{FuncColor}{$\triangleright$\enspace\texttt{BravaisSupergroups({\mdseries\slshape G})\index{BravaisSupergroups@\texttt{BravaisSupergroups}}
\label{BravaisSupergroups}
}\hfill{\scriptsize (attribute)}}\\


 returns the subgroups of $GL_n({\ensuremath{\mathbb Z}})$ that contain the Bravais group of \mbox{\texttt{\mdseries\slshape G}} and are Bravais groups themselves. }

 

\subsection{\textcolor{Chapter }{NormalizerInGLnZBravaisGroup}}
\logpage{[ 44, 6, 14 ]}\nobreak
\hyperdef{L}{X79B7CD797A420720}{}
{\noindent\textcolor{FuncColor}{$\triangleright$\enspace\texttt{NormalizerInGLnZBravaisGroup({\mdseries\slshape G})\index{NormalizerInGLnZBravaisGroup@\texttt{NormalizerInGLnZBravaisGroup}}
\label{NormalizerInGLnZBravaisGroup}
}\hfill{\scriptsize (attribute)}}\\


 returns the normalizer of the Bravais group of \mbox{\texttt{\mdseries\slshape G}} in the appropriate $GL_n({\ensuremath{\mathbb Z}})$. }

 }

  
\section{\textcolor{Chapter }{Acting OnRight and OnLeft}}\label{Acting OnRight and OnLeft}
\logpage{[ 44, 7, 0 ]}
\hyperdef{L}{X868288377CFA8D1B}{}
{
  In \textsf{GAP}, matrices by convention act on row vectors from the right, whereas in
crystallography the convention is to act on column vectors from the left. The
definition of certain algebraic objects important in crystallography
implicitly depends on which action is assumed. This holds true in particular
for quadratic forms invariant under a matrix group. In a similar way, the
representation of affine crystallographic groups, as they are provided by the \textsf{GAP} package \textsf{CrystGap}, depends on which action is assumed. Crystallographers are used to the action
from the left, whereas the action from the right is the natural one for \textsf{GAP}. For this reason, a number of functions which are important in
crystallography, and whose result depends on which action is assumed, are
provided in two versions, one for the usual action from the right, and one for
the crystallographic action from the left. 

 For every such function, this fact is explicitly mentioned. The naming scheme
is as follows: If \texttt{SomeThing} is such a function, there will be functions \texttt{SomeThingOnRight} and \texttt{SomeThingOnLeft}, assuming action from the right and from the left, respectively. In addition,
there is a generic function \texttt{SomeThing}, which returns either the result of \texttt{SomeThingOnRight} or \texttt{SomeThingOnLeft}, depending on the global variable \texttt{CrystGroupDefaultAction} (\ref{CrystGroupDefaultAction}). 

\subsection{\textcolor{Chapter }{CrystGroupDefaultAction}}
\logpage{[ 44, 7, 1 ]}\nobreak
\hyperdef{L}{X7D1318A6780CD88B}{}
{\noindent\textcolor{FuncColor}{$\triangleright$\enspace\texttt{CrystGroupDefaultAction\index{CrystGroupDefaultAction@\texttt{CrystGroupDefaultAction}}
\label{CrystGroupDefaultAction}
}\hfill{\scriptsize (global variable)}}\\


 can have either of the two values \texttt{RightAction} and \texttt{LeftAction}. The initial value is \texttt{RightAction}. For functions which have variants OnRight and OnLeft, this variable
determines which variant is returned by the generic form. The value of \texttt{CrystGroupDefaultAction} can be changed with the function \texttt{SetCrystGroupDefaultAction} (\ref{SetCrystGroupDefaultAction}). }

 

\subsection{\textcolor{Chapter }{SetCrystGroupDefaultAction}}
\logpage{[ 44, 7, 2 ]}\nobreak
\hyperdef{L}{X792D237385977BE6}{}
{\noindent\textcolor{FuncColor}{$\triangleright$\enspace\texttt{SetCrystGroupDefaultAction({\mdseries\slshape action})\index{SetCrystGroupDefaultAction@\texttt{SetCrystGroupDefaultAction}}
\label{SetCrystGroupDefaultAction}
}\hfill{\scriptsize (function)}}\\


 allows one to set the value of the global variable \texttt{CrystGroupDefaultAction} (\ref{CrystGroupDefaultAction}). Only the arguments \texttt{RightAction} and \texttt{LeftAction} are allowed. Initially, the value of \texttt{CrystGroupDefaultAction} (\ref{CrystGroupDefaultAction}) is \texttt{RightAction}. }

 }

 }

             
\chapter{\textcolor{Chapter }{Polycyclic Groups}}\label{Polycyclic Groups}
\logpage{[ 45, 0, 0 ]}
\hyperdef{L}{X86007B0083F60470}{}
{
  A group \mbox{\texttt{\mdseries\slshape G}} is \emph{polycyclic} if there exists a subnormal series $G = C_1 > C_2 > \ldots > C_n > C_{{n+1}} = \{ 1 \}$ with cyclic factors. Such a series is called \emph{pc series} of \mbox{\texttt{\mdseries\slshape G}}. 

 Every polycyclic group is solvable and every finite solvable group is
polycyclic. However, there are infinite solvable groups which are not
polycyclic. 

 In \textsf{GAP} there exists a large number of methods for polycyclic groups which are based
upon the polycyclic structure of these groups. These methods are usually very
efficient, especially for groups which are given by a
pc\texttt{\symbol{45}}presentation (see chapter \ref{Pc Groups}), and can be applied to many types of groups. Hence \textsf{GAP} tries to use them whenever possible, for example, for permutation groups and
matrix groups over finite fields that are known to be polycyclic (the only
exception is the representation as finitely presented group for which the
polycyclic methods cannot be used in general). 

 At the current state of implementations the \textsf{GAP} library contains methods to compute with finite polycyclic groups, while the \textsf{GAP} package \textsf{Polycyclic} by Bettina Eick and Werner Nickel allows also computations with infinite
polycyclic groups which are given by a pc\texttt{\symbol{45}}presentation. 

     
\section{\textcolor{Chapter }{Polycyclic Generating Systems}}\label{Polycyclic Generating Systems}
\logpage{[ 45, 1, 0 ]}
\hyperdef{L}{X7F18A01785DBAC4E}{}
{
  Let \mbox{\texttt{\mdseries\slshape G}} be a polycyclic group with a pc series as above. A \emph{polycyclic generating sequence} (\emph{pcgs} for short) of \mbox{\texttt{\mdseries\slshape G}} is a sequence $P := (g_1, \ldots, g_n)$ of elements of \mbox{\texttt{\mdseries\slshape G}} such that $C_i = \langle C_{{i+1}}, g_i \rangle$ for $1 \leq i \leq n$. Note that each polycyclic group has a pcgs, but except for very small
groups, a pcgs is not unique. 

 For each index $i$ the subsequence of elements $(g_i, \ldots, g_n)$ forms a pcgs of the subgroup $C_i$. In particular, these \emph{tails} generate the subgroups of the pc series and hence we say that the pc series is \emph{determined} by $P$. 

 Let $r_i$ be the index of $C_{{i+1}}$ in $C_i$ which is either a finite positive number or infinity. Then $r_i$ is the order of $g_i C_{{i+1}}$ and we call the resulting list of indices the \emph{relative orders} of the pcgs \mbox{\texttt{\mdseries\slshape P}}. 

 Moreover, with respect to a given pcgs $(g_1, \ldots, g_n)$ each element \mbox{\texttt{\mdseries\slshape g}} of \mbox{\texttt{\mdseries\slshape G}} can be represented in a unique way as a product $g = g_1^{{e_1}} \cdot g_2^{{e_2}} \cdots g_n^{{e_n}}$ with exponents $e_i \in \{0, \ldots, r_i-1\}$, if $r_i$ is finite, and $e_i \in {\ensuremath{\mathbb Z}}$ otherwise. Words of this form are called \emph{normal words} or \emph{words in normal form}. Then the integer vector $[ e_1, \ldots, e_n ]$ is called the \emph{exponent vector} of the element $g$. Furthermore, the smallest index $k$ such that $e_k \neq 0$ is called the \emph{depth} of \mbox{\texttt{\mdseries\slshape g}} and $e_k$ is the \emph{leading exponent} of \mbox{\texttt{\mdseries\slshape g}}. 

 For many applications we have to assume that each of the relative orders $r_i$ is either a prime or infinity. This is equivalent to saying that there are no
trivial factors in the pc series and the finite factors of the pc series are
maximal refined. Then we obtain that $r_i$ is the order of $g C_{{i+1}}$ for all elements $g$ in $C_i \setminus C_{{i+1}}$ and we call $r_i$ the \emph{relative order} of the element $g$. }

  
\section{\textcolor{Chapter }{Computing a Pcgs}}\label{Computing a Pcgs}
\logpage{[ 45, 2, 0 ]}
\hyperdef{L}{X87F7E31879AFA06C}{}
{
  Suppose a group \mbox{\texttt{\mdseries\slshape G}} is given; for example, let \mbox{\texttt{\mdseries\slshape G}} be a permutation or matrix group. Then we can ask \textsf{GAP} to compute a pcgs of this group. If \mbox{\texttt{\mdseries\slshape G}} is not polycyclic, the result will be \texttt{fail}. 

 Note that these methods can only be applied if \mbox{\texttt{\mdseries\slshape G}} is not given as finitely presented group. For finitely presented groups one
can try to compute a pcgs via the polycyclic quotient methods, see \ref{Quotient Methods}. 

 Note also that a pcgs behaves like a list. 

\subsection{\textcolor{Chapter }{Pcgs}}
\logpage{[ 45, 2, 1 ]}\nobreak
\hyperdef{L}{X84C3750C7A4EEC34}{}
{\noindent\textcolor{FuncColor}{$\triangleright$\enspace\texttt{Pcgs({\mdseries\slshape G})\index{Pcgs@\texttt{Pcgs}}
\label{Pcgs}
}\hfill{\scriptsize (attribute)}}\\


 returns a pcgs for the group \mbox{\texttt{\mdseries\slshape G}}. If \mbox{\texttt{\mdseries\slshape grp}} is not polycyclic it returns \texttt{fail} \emph{and this result is not stored as attribute value}, in particular in this case the filter \texttt{HasPcgs} is \emph{not} set for \mbox{\texttt{\mdseries\slshape G}}! }

 

\subsection{\textcolor{Chapter }{IsPcgs}}
\logpage{[ 45, 2, 2 ]}\nobreak
\hyperdef{L}{X8635E61A7DB73BA6}{}
{\noindent\textcolor{FuncColor}{$\triangleright$\enspace\texttt{IsPcgs({\mdseries\slshape obj})\index{IsPcgs@\texttt{IsPcgs}}
\label{IsPcgs}
}\hfill{\scriptsize (Category)}}\\


 The category of pcgs. 
\begin{Verbatim}[commandchars=!@|,fontsize=\small,frame=single,label=Example]
  !gapprompt@gap>| !gapinput@G := Group((1,2,3,4),(1,2));;|
  !gapprompt@gap>| !gapinput@p := Pcgs(G);|
  Pcgs([ (3,4), (2,4,3), (1,4)(2,3), (1,3)(2,4) ])
  !gapprompt@gap>| !gapinput@IsPcgs( p );|
  true
  !gapprompt@gap>| !gapinput@p[1];|
  (3,4)
  !gapprompt@gap>| !gapinput@G := Group((1,2,3,4,5),(1,2));;|
  !gapprompt@gap>| !gapinput@Pcgs(G);|
  fail
\end{Verbatim}
 }

 

\subsection{\textcolor{Chapter }{CanEasilyComputePcgs}}
\logpage{[ 45, 2, 3 ]}\nobreak
\hyperdef{L}{X7B561B1685CEC2AB}{}
{\noindent\textcolor{FuncColor}{$\triangleright$\enspace\texttt{CanEasilyComputePcgs({\mdseries\slshape grp})\index{CanEasilyComputePcgs@\texttt{CanEasilyComputePcgs}}
\label{CanEasilyComputePcgs}
}\hfill{\scriptsize (filter)}}\\


 This filter indicates whether it is possible to compute a pcgs for \mbox{\texttt{\mdseries\slshape grp}} cheaply. Clearly, \mbox{\texttt{\mdseries\slshape grp}} must be polycyclic in this case. However, not for every polycyclic group there
is a method to compute a pcgs at low costs. This filter is used in the method
selection mainly. Note that this filter may change its value from \texttt{false} to \texttt{true}. 
\begin{Verbatim}[commandchars=!@|,fontsize=\small,frame=single,label=Example]
  !gapprompt@gap>| !gapinput@G := Group( (1,2,3,4),(1,2) );|
  Group([ (1,2,3,4), (1,2) ])
  !gapprompt@gap>| !gapinput@CanEasilyComputePcgs(G);|
  false
  !gapprompt@gap>| !gapinput@Pcgs(G);|
  Pcgs([ (3,4), (2,4,3), (1,4)(2,3), (1,3)(2,4) ])
  !gapprompt@gap>| !gapinput@CanEasilyComputePcgs(G);|
  true
\end{Verbatim}
 }

 }

  
\section{\textcolor{Chapter }{Defining a Pcgs Yourself}}\label{Defining a Pcgs Yourself}
\logpage{[ 45, 3, 0 ]}
\hyperdef{L}{X7CAAD6D2838354D9}{}
{
  In a number of situations it might be useful to supply a pcgs to a group. 

 Note that the elementary operations for such a pcgs might be rather
inefficient, since \textsf{GAP} has to use generic methods in this case. It might be helpful to supply the
relative orders of the self\texttt{\symbol{45}}defined pcgs as well by \texttt{SetRelativeOrder}. See also{\nobreakspace}\texttt{IsPrimeOrdersPcgs} (\ref{IsPrimeOrdersPcgs}). 

\subsection{\textcolor{Chapter }{PcgsByPcSequence}}
\logpage{[ 45, 3, 1 ]}\nobreak
\hyperdef{L}{X7E139C3D80847D76}{}
{\noindent\textcolor{FuncColor}{$\triangleright$\enspace\texttt{PcgsByPcSequence({\mdseries\slshape fam, pcs})\index{PcgsByPcSequence@\texttt{PcgsByPcSequence}}
\label{PcgsByPcSequence}
}\hfill{\scriptsize (operation)}}\\
\noindent\textcolor{FuncColor}{$\triangleright$\enspace\texttt{PcgsByPcSequenceNC({\mdseries\slshape fam, pcs})\index{PcgsByPcSequenceNC@\texttt{PcgsByPcSequenceNC}}
\label{PcgsByPcSequenceNC}
}\hfill{\scriptsize (operation)}}\\


 constructs a pcgs for the elements family \mbox{\texttt{\mdseries\slshape fam}} from the elements in the list \mbox{\texttt{\mdseries\slshape pcs}}. The elements must lie in the family \mbox{\texttt{\mdseries\slshape fam}}. \texttt{PcgsByPcSequence} and its \texttt{NC} variant will always create a new pcgs which is not induced by any other pcgs
(cf. \texttt{InducedPcgsByPcSequence} (\ref{InducedPcgsByPcSequence})). 
\begin{Verbatim}[commandchars=!@|,fontsize=\small,frame=single,label=Example]
  !gapprompt@gap>| !gapinput@fam := FamilyObj( (1,2) );; # the family of permutations|
  !gapprompt@gap>| !gapinput@p := PcgsByPcSequence( fam, [(1,2),(1,2,3)] );|
  Pcgs([ (1,2), (1,2,3) ])
  !gapprompt@gap>| !gapinput@RelativeOrders(p);|
  [ 2, 3 ]
  !gapprompt@gap>| !gapinput@ExponentsOfPcElement( p, (1,3,2) );|
  [ 0, 2 ]
\end{Verbatim}
 }

 }

  
\section{\textcolor{Chapter }{Elementary Operations for a Pcgs}}\label{Elementary Operations for a Pcgs}
\logpage{[ 45, 4, 0 ]}
\hyperdef{L}{X816C5E8E7F71C9D8}{}
{
  

\subsection{\textcolor{Chapter }{RelativeOrders}}
\logpage{[ 45, 4, 1 ]}\nobreak
\hyperdef{L}{X7DD0DF677AC1CF10}{}
{\noindent\textcolor{FuncColor}{$\triangleright$\enspace\texttt{RelativeOrders({\mdseries\slshape pcgs})\index{RelativeOrders@\texttt{RelativeOrders}}
\label{RelativeOrders}
}\hfill{\scriptsize (attribute)}}\\


 \index{RelativeOrders@\texttt{RelativeOrders}!of a pcgs} returns the list of relative orders of the pcgs \mbox{\texttt{\mdseries\slshape pcgs}}. }

 

\subsection{\textcolor{Chapter }{IsFiniteOrdersPcgs}}
\logpage{[ 45, 4, 2 ]}\nobreak
\hyperdef{L}{X80D526848427A5C6}{}
{\noindent\textcolor{FuncColor}{$\triangleright$\enspace\texttt{IsFiniteOrdersPcgs({\mdseries\slshape pcgs})\index{IsFiniteOrdersPcgs@\texttt{IsFiniteOrdersPcgs}}
\label{IsFiniteOrdersPcgs}
}\hfill{\scriptsize (property)}}\\


 tests whether the relative orders of \mbox{\texttt{\mdseries\slshape pcgs}} are all finite. }

 

\subsection{\textcolor{Chapter }{IsPrimeOrdersPcgs}}
\logpage{[ 45, 4, 3 ]}\nobreak
\hyperdef{L}{X866C3A5382FF231A}{}
{\noindent\textcolor{FuncColor}{$\triangleright$\enspace\texttt{IsPrimeOrdersPcgs({\mdseries\slshape pcgs})\index{IsPrimeOrdersPcgs@\texttt{IsPrimeOrdersPcgs}}
\label{IsPrimeOrdersPcgs}
}\hfill{\scriptsize (property)}}\\


 tests whether the relative orders of \mbox{\texttt{\mdseries\slshape pcgs}} are prime numbers. Many algorithms require a pcgs to have this property. The
operation{\nobreakspace}\texttt{IsomorphismRefinedPcGroup} (\ref{IsomorphismRefinedPcGroup}) can be of help here. }

 

\subsection{\textcolor{Chapter }{PcSeries}}
\logpage{[ 45, 4, 4 ]}\nobreak
\hyperdef{L}{X827A7B097A959579}{}
{\noindent\textcolor{FuncColor}{$\triangleright$\enspace\texttt{PcSeries({\mdseries\slshape pcgs})\index{PcSeries@\texttt{PcSeries}}
\label{PcSeries}
}\hfill{\scriptsize (attribute)}}\\


 returns the subnormal series determined by \mbox{\texttt{\mdseries\slshape pcgs}}. }

 

\subsection{\textcolor{Chapter }{GroupOfPcgs}}
\logpage{[ 45, 4, 5 ]}\nobreak
\hyperdef{L}{X7903702E8194EF29}{}
{\noindent\textcolor{FuncColor}{$\triangleright$\enspace\texttt{GroupOfPcgs({\mdseries\slshape pcgs})\index{GroupOfPcgs@\texttt{GroupOfPcgs}}
\label{GroupOfPcgs}
}\hfill{\scriptsize (attribute)}}\\


 The group generated by \mbox{\texttt{\mdseries\slshape pcgs}}. }

 

\subsection{\textcolor{Chapter }{OneOfPcgs}}
\logpage{[ 45, 4, 6 ]}\nobreak
\hyperdef{L}{X878FB11887524E2C}{}
{\noindent\textcolor{FuncColor}{$\triangleright$\enspace\texttt{OneOfPcgs({\mdseries\slshape pcgs})\index{OneOfPcgs@\texttt{OneOfPcgs}}
\label{OneOfPcgs}
}\hfill{\scriptsize (attribute)}}\\


 The identity of the group generated by \mbox{\texttt{\mdseries\slshape pcgs}}. 
\begin{Verbatim}[commandchars=!@|,fontsize=\small,frame=single,label=Example]
  !gapprompt@gap>| !gapinput@G := Group( (1,2,3,4),(1,2) );; p := Pcgs(G);;|
  !gapprompt@gap>| !gapinput@RelativeOrders(p);|
  [ 2, 3, 2, 2 ]
  !gapprompt@gap>| !gapinput@IsFiniteOrdersPcgs(p);|
  true
  !gapprompt@gap>| !gapinput@IsPrimeOrdersPcgs(p);|
  true
  !gapprompt@gap>| !gapinput@PcSeries(p);|
  [ Group([ (3,4), (2,4,3), (1,4)(2,3), (1,3)(2,4) ]),
    Group([ (2,4,3), (1,4)(2,3), (1,3)(2,4) ]),
    Group([ (1,4)(2,3), (1,3)(2,4) ]), Group([ (1,3)(2,4) ]), Group(())
   ]
\end{Verbatim}
 }

 }

  
\section{\textcolor{Chapter }{Elementary Operations for a Pcgs and an Element}}\label{Elementary Operations for a Pcgs and an Element}
\logpage{[ 45, 5, 0 ]}
\hyperdef{L}{X84243AA07DA5A827}{}
{
  

\subsection{\textcolor{Chapter }{RelativeOrderOfPcElement}}
\logpage{[ 45, 5, 1 ]}\nobreak
\hyperdef{L}{X7B941D4A7CAFCD73}{}
{\noindent\textcolor{FuncColor}{$\triangleright$\enspace\texttt{RelativeOrderOfPcElement({\mdseries\slshape pcgs, elm})\index{RelativeOrderOfPcElement@\texttt{RelativeOrderOfPcElement}}
\label{RelativeOrderOfPcElement}
}\hfill{\scriptsize (operation)}}\\


 The relative order of \mbox{\texttt{\mdseries\slshape elm}} with respect to the prime order pcgs \mbox{\texttt{\mdseries\slshape pcgs}}. }

 

\subsection{\textcolor{Chapter }{ExponentOfPcElement}}
\logpage{[ 45, 5, 2 ]}\nobreak
\hyperdef{L}{X78134914842E2F5F}{}
{\noindent\textcolor{FuncColor}{$\triangleright$\enspace\texttt{ExponentOfPcElement({\mdseries\slshape pcgs, elm, pos})\index{ExponentOfPcElement@\texttt{ExponentOfPcElement}}
\label{ExponentOfPcElement}
}\hfill{\scriptsize (operation)}}\\


 returns the \mbox{\texttt{\mdseries\slshape pos}}\texttt{\symbol{45}}th exponent of \mbox{\texttt{\mdseries\slshape elm}} with respect to \mbox{\texttt{\mdseries\slshape pcgs}}. }

 

\subsection{\textcolor{Chapter }{ExponentsOfPcElement}}
\logpage{[ 45, 5, 3 ]}\nobreak
\hyperdef{L}{X848DAEBF7DC448A5}{}
{\noindent\textcolor{FuncColor}{$\triangleright$\enspace\texttt{ExponentsOfPcElement({\mdseries\slshape pcgs, elm[, posran]})\index{ExponentsOfPcElement@\texttt{ExponentsOfPcElement}}
\label{ExponentsOfPcElement}
}\hfill{\scriptsize (operation)}}\\


 returns the exponents of \mbox{\texttt{\mdseries\slshape elm}} with respect to \mbox{\texttt{\mdseries\slshape pcgs}}. The three argument version returns the exponents in the positions given in \mbox{\texttt{\mdseries\slshape posran}}. }

 

\subsection{\textcolor{Chapter }{DepthOfPcElement}}
\logpage{[ 45, 5, 4 ]}\nobreak
\hyperdef{L}{X829BCB267CDBC5C0}{}
{\noindent\textcolor{FuncColor}{$\triangleright$\enspace\texttt{DepthOfPcElement({\mdseries\slshape pcgs, elm})\index{DepthOfPcElement@\texttt{DepthOfPcElement}}
\label{DepthOfPcElement}
}\hfill{\scriptsize (operation)}}\\


 returns the depth of the element \mbox{\texttt{\mdseries\slshape elm}} with respect to \mbox{\texttt{\mdseries\slshape pcgs}}. }

 

\subsection{\textcolor{Chapter }{LeadingExponentOfPcElement}}
\logpage{[ 45, 5, 5 ]}\nobreak
\hyperdef{L}{X7D47966479EA2890}{}
{\noindent\textcolor{FuncColor}{$\triangleright$\enspace\texttt{LeadingExponentOfPcElement({\mdseries\slshape pcgs, elm})\index{LeadingExponentOfPcElement@\texttt{LeadingExponentOfPcElement}}
\label{LeadingExponentOfPcElement}
}\hfill{\scriptsize (operation)}}\\


 returns the leading exponent of \mbox{\texttt{\mdseries\slshape elm}} with respect to \mbox{\texttt{\mdseries\slshape pcgs}}. }

 

\subsection{\textcolor{Chapter }{PcElementByExponents}}
\logpage{[ 45, 5, 6 ]}\nobreak
\hyperdef{L}{X87AF746B8328F5D0}{}
{\noindent\textcolor{FuncColor}{$\triangleright$\enspace\texttt{PcElementByExponents({\mdseries\slshape pcgs, list})\index{PcElementByExponents@\texttt{PcElementByExponents}}
\label{PcElementByExponents}
}\hfill{\scriptsize (function)}}\\
\noindent\textcolor{FuncColor}{$\triangleright$\enspace\texttt{PcElementByExponentsNC({\mdseries\slshape pcgs[, basisind], list})\index{PcElementByExponentsNC@\texttt{PcElementByExponentsNC}}
\label{PcElementByExponentsNC}
}\hfill{\scriptsize (operation)}}\\


 returns the element corresponding to the exponent vector \mbox{\texttt{\mdseries\slshape list}} with respect to \mbox{\texttt{\mdseries\slshape pcgs}}. The exponents in \mbox{\texttt{\mdseries\slshape list}} must be in the range of permissible exponents for \mbox{\texttt{\mdseries\slshape pcgs}}. \emph{It is not guaranteed that \texttt{PcElementByExponents} will reduce the exponents modulo the relative orders}. (You should use the operation \texttt{LinearCombinationPcgs} (\ref{LinearCombinationPcgs}) for this purpose.) The \texttt{NC} version does not check that the lengths of the lists fit together and does not
check the exponent range. 

 The three argument version gives exponents only w.r.t. the generators in \mbox{\texttt{\mdseries\slshape pcgs}} indexed by \mbox{\texttt{\mdseries\slshape basisind}}. }

 

\subsection{\textcolor{Chapter }{LinearCombinationPcgs}}
\logpage{[ 45, 5, 7 ]}\nobreak
\hyperdef{L}{X7F8BD7A87DB3933A}{}
{\noindent\textcolor{FuncColor}{$\triangleright$\enspace\texttt{LinearCombinationPcgs({\mdseries\slshape pcgs, list[, one]})\index{LinearCombinationPcgs@\texttt{LinearCombinationPcgs}}
\label{LinearCombinationPcgs}
}\hfill{\scriptsize (function)}}\\


 returns the product $\prod_i \mbox{\texttt{\mdseries\slshape pcgs}}[i]^{{\mbox{\texttt{\mdseries\slshape list}}[i]}}$. In contrast to \texttt{PcElementByExponents} (\ref{PcElementByExponents}) this permits negative exponents. \mbox{\texttt{\mdseries\slshape pcgs}} might be a list of group elements. In this case, an appropriate identity
element \mbox{\texttt{\mdseries\slshape one}} must be given. \mbox{\texttt{\mdseries\slshape list}} can be empty. 
\begin{Verbatim}[commandchars=!@|,fontsize=\small,frame=single,label=Example]
  !gapprompt@gap>| !gapinput@G := Group( (1,2,3,4),(1,2) );; P := Pcgs(G);;|
  !gapprompt@gap>| !gapinput@g := PcElementByExponents(P, [0,1,1,1]);|
  (1,2,3)
  !gapprompt@gap>| !gapinput@ExponentsOfPcElement(P, g);|
  [ 0, 1, 1, 1 ]
\end{Verbatim}
 }

 

\subsection{\textcolor{Chapter }{SiftedPcElement}}
\logpage{[ 45, 5, 8 ]}\nobreak
\hyperdef{L}{X8066B66D8069BAB4}{}
{\noindent\textcolor{FuncColor}{$\triangleright$\enspace\texttt{SiftedPcElement({\mdseries\slshape pcgs, elm})\index{SiftedPcElement@\texttt{SiftedPcElement}}
\label{SiftedPcElement}
}\hfill{\scriptsize (operation)}}\\


 sifts \mbox{\texttt{\mdseries\slshape elm}} through \mbox{\texttt{\mdseries\slshape pcgs}}, reducing it if the depth is the same as the depth of one of the generators
in \mbox{\texttt{\mdseries\slshape pcgs}}. Thus the identity is returned if \mbox{\texttt{\mdseries\slshape elm}} lies in the group generated by \mbox{\texttt{\mdseries\slshape pcgs}}. \mbox{\texttt{\mdseries\slshape pcgs}} must be an induced pcgs (see section \ref{Subgroups of Polycyclic Groups - Induced Pcgs}) and \mbox{\texttt{\mdseries\slshape elm}} must lie in the span of the parent of \mbox{\texttt{\mdseries\slshape pcgs}}. }

 

\subsection{\textcolor{Chapter }{CanonicalPcElement}}
\logpage{[ 45, 5, 9 ]}\nobreak
\hyperdef{L}{X7B52ADE7878A749A}{}
{\noindent\textcolor{FuncColor}{$\triangleright$\enspace\texttt{CanonicalPcElement({\mdseries\slshape ipcgs, elm})\index{CanonicalPcElement@\texttt{CanonicalPcElement}}
\label{CanonicalPcElement}
}\hfill{\scriptsize (operation)}}\\


 reduces \mbox{\texttt{\mdseries\slshape elm}} at the induces pcgs \mbox{\texttt{\mdseries\slshape ipcgs}} such that the exponents of the reduced result \mbox{\texttt{\mdseries\slshape r}} are zero at the depths for which there are generators in \mbox{\texttt{\mdseries\slshape ipcgs}}. Elements, whose quotient lies in the group generated by \mbox{\texttt{\mdseries\slshape ipcgs}} yield the same canonical element. }

 

\subsection{\textcolor{Chapter }{ReducedPcElement}}
\logpage{[ 45, 5, 10 ]}\nobreak
\hyperdef{L}{X7A94AA357DB2F86C}{}
{\noindent\textcolor{FuncColor}{$\triangleright$\enspace\texttt{ReducedPcElement({\mdseries\slshape pcgs, x, y})\index{ReducedPcElement@\texttt{ReducedPcElement}}
\label{ReducedPcElement}
}\hfill{\scriptsize (operation)}}\\


 reduces the element \mbox{\texttt{\mdseries\slshape x}} by dividing off (from the left) a power of \mbox{\texttt{\mdseries\slshape y}} such that the leading coefficient of the result with respect to \mbox{\texttt{\mdseries\slshape pcgs}} becomes zero. The elements \mbox{\texttt{\mdseries\slshape x}} and \mbox{\texttt{\mdseries\slshape y}} therefore have to have the same depth. }

 

\subsection{\textcolor{Chapter }{CleanedTailPcElement}}
\logpage{[ 45, 5, 11 ]}\nobreak
\hyperdef{L}{X8702D76D8284CF3E}{}
{\noindent\textcolor{FuncColor}{$\triangleright$\enspace\texttt{CleanedTailPcElement({\mdseries\slshape pcgs, elm, dep})\index{CleanedTailPcElement@\texttt{CleanedTailPcElement}}
\label{CleanedTailPcElement}
}\hfill{\scriptsize (operation)}}\\


 returns an element in the span of \mbox{\texttt{\mdseries\slshape pcgs}} whose exponents for indices $1$ to $\mbox{\texttt{\mdseries\slshape dep}}-1$ with respect to \mbox{\texttt{\mdseries\slshape pcgs}} are the same as those of \mbox{\texttt{\mdseries\slshape elm}}, the remaining exponents are undefined. This can be used to obtain more ``simple'' elements if only representatives in a factor are required, see{\nobreakspace}\ref{Factor Groups of Polycyclic Groups - Modulo Pcgs}. 

 The difference to \texttt{HeadPcElementByNumber} (\ref{HeadPcElementByNumber}) is that this function is guaranteed to zero out trailing coefficients while \texttt{CleanedTailPcElement} will only do this if it can be done cheaply. }

 

\subsection{\textcolor{Chapter }{HeadPcElementByNumber}}
\logpage{[ 45, 5, 12 ]}\nobreak
\hyperdef{L}{X830A0D037DBEAA97}{}
{\noindent\textcolor{FuncColor}{$\triangleright$\enspace\texttt{HeadPcElementByNumber({\mdseries\slshape pcgs, elm, dep})\index{HeadPcElementByNumber@\texttt{HeadPcElementByNumber}}
\label{HeadPcElementByNumber}
}\hfill{\scriptsize (operation)}}\\


 returns an element in the span of \mbox{\texttt{\mdseries\slshape pcgs}} whose exponents for indices $1$ to \mbox{\texttt{\mdseries\slshape dep}}$-1$ with respect to \mbox{\texttt{\mdseries\slshape pcgs}} are the same as those of \mbox{\texttt{\mdseries\slshape elm}}, the remaining exponents are zero. This can be used to obtain more ``simple'' elements if only representatives in a factor are required. }

 }

  
\section{\textcolor{Chapter }{Exponents of Special Products}}\label{Exponents of Special Products}
\logpage{[ 45, 6, 0 ]}
\hyperdef{L}{X7EF61EA4822870E7}{}
{
  There are certain products of elements whose exponents are used often within
algorithms, and which might be obtained more easily than by computing the
product first and to obtain its exponents afterwards. The operations in this
section provide a way to obtain such exponent vectors directly. 

 (The circumstances under which these operations give a speedup depend very
much on the pcgs and the representation of elements that is used. So the
following operations are not guaranteed to give a speedup in every case,
however the default methods are not slower than to compute the exponents of a
product and thus these operations should \emph{always} be used if applicable.) 

 The second class are exponents of products of the generators which make up the
pcgs. If the pcgs used is a family pcgs (see \texttt{FamilyPcgs} (\ref{FamilyPcgs})) then these exponents can be looked up and do not need to be computed. 

\subsection{\textcolor{Chapter }{ExponentsConjugateLayer}}
\logpage{[ 45, 6, 1 ]}\nobreak
\hyperdef{L}{X868D6DB07D349460}{}
{\noindent\textcolor{FuncColor}{$\triangleright$\enspace\texttt{ExponentsConjugateLayer({\mdseries\slshape mpcgs, elm, e})\index{ExponentsConjugateLayer@\texttt{ExponentsConjugateLayer}}
\label{ExponentsConjugateLayer}
}\hfill{\scriptsize (operation)}}\\


 Computes the exponents of \mbox{\texttt{\mdseries\slshape elm}}\texttt{\texttt{\symbol{94}}}\mbox{\texttt{\mdseries\slshape e}} with respect to \mbox{\texttt{\mdseries\slshape mpcgs}}; \mbox{\texttt{\mdseries\slshape elm}} must be in the span of \mbox{\texttt{\mdseries\slshape mpcgs}}, \mbox{\texttt{\mdseries\slshape e}} a pc element in the span of the parent pcgs of \mbox{\texttt{\mdseries\slshape mpcgs}} and \mbox{\texttt{\mdseries\slshape mpcgs}} must be the modulo pcgs for an abelian layer. (This is the usual case when
acting on a chief factor). In this case if \mbox{\texttt{\mdseries\slshape mpcgs}} is induced by the family pcgs (see section \ref{Subgroups of Polycyclic Groups - Induced Pcgs}), the exponents can be computed directly by looking up exponents without
having to compute in the group and having to collect a potential tail. }

 

\subsection{\textcolor{Chapter }{ExponentsOfRelativePower}}
\logpage{[ 45, 6, 2 ]}\nobreak
\hyperdef{L}{X874F70697FE7B6DF}{}
{\noindent\textcolor{FuncColor}{$\triangleright$\enspace\texttt{ExponentsOfRelativePower({\mdseries\slshape pcgs, i})\index{ExponentsOfRelativePower@\texttt{ExponentsOfRelativePower}}
\label{ExponentsOfRelativePower}
}\hfill{\scriptsize (operation)}}\\


 For $p = \mbox{\texttt{\mdseries\slshape pcgs}}[\mbox{\texttt{\mdseries\slshape i}}]$ this function returns the exponent vector with respect to \mbox{\texttt{\mdseries\slshape pcgs}} of the element $p^e$ where $e$ is the relative order of \mbox{\texttt{\mdseries\slshape p}} in \mbox{\texttt{\mdseries\slshape pcgs}}. For the family pcgs or pcgs induced by it (see section \ref{Subgroups of Polycyclic Groups - Induced Pcgs}), this might be faster than computing the element and computing its exponent
vector. }

 

\subsection{\textcolor{Chapter }{ExponentsOfConjugate}}
\logpage{[ 45, 6, 3 ]}\nobreak
\hyperdef{L}{X78CAF32F864EF656}{}
{\noindent\textcolor{FuncColor}{$\triangleright$\enspace\texttt{ExponentsOfConjugate({\mdseries\slshape pcgs, i, j})\index{ExponentsOfConjugate@\texttt{ExponentsOfConjugate}}
\label{ExponentsOfConjugate}
}\hfill{\scriptsize (operation)}}\\


 returns the exponents of \texttt{\mbox{\texttt{\mdseries\slshape pcgs}}[\mbox{\texttt{\mdseries\slshape i}}]\texttt{\symbol{94}}\mbox{\texttt{\mdseries\slshape pcgs}}[\mbox{\texttt{\mdseries\slshape j}}]} with respect to \mbox{\texttt{\mdseries\slshape pcgs}}. For the family pcgs or pcgs induced by it (see section \ref{Subgroups of Polycyclic Groups - Induced Pcgs}), this might be faster than computing the element and computing its exponent
vector. }

 

\subsection{\textcolor{Chapter }{ExponentsOfCommutator}}
\logpage{[ 45, 6, 4 ]}\nobreak
\hyperdef{L}{X875689897DD0CAFC}{}
{\noindent\textcolor{FuncColor}{$\triangleright$\enspace\texttt{ExponentsOfCommutator({\mdseries\slshape pcgs, i, j})\index{ExponentsOfCommutator@\texttt{ExponentsOfCommutator}}
\label{ExponentsOfCommutator}
}\hfill{\scriptsize (operation)}}\\


 returns the exponents of the commutator \texttt{Comm( }$\mbox{\texttt{\mdseries\slshape pcgs}}[\mbox{\texttt{\mdseries\slshape i}}], \mbox{\texttt{\mdseries\slshape pcgs}}[\mbox{\texttt{\mdseries\slshape j}}]$\texttt{ )} with respect to \mbox{\texttt{\mdseries\slshape pcgs}}. For the family pcgs or pcgs induced by it, (see section \ref{Subgroups of Polycyclic Groups - Induced Pcgs}), this might be faster than computing the element and computing its exponent
vector. }

 }

  
\section{\textcolor{Chapter }{Subgroups of Polycyclic Groups {\textendash} Induced Pcgs}}\label{Subgroups of Polycyclic Groups - Induced Pcgs}
\logpage{[ 45, 7, 0 ]}
\hyperdef{L}{X8676397383093D1E}{}
{
  Let \mbox{\texttt{\mdseries\slshape U}} be a subgroup of \mbox{\texttt{\mdseries\slshape G}} and let \mbox{\texttt{\mdseries\slshape P}} be a pcgs of \mbox{\texttt{\mdseries\slshape G}} as above such that \mbox{\texttt{\mdseries\slshape P}} determines the subnormal series $G = C_1 > \ldots > C_{{n+1}} = \{ 1 \}$. Then the series of subgroups $U \cap C_i$ is a subnormal series of \mbox{\texttt{\mdseries\slshape U}} with cyclic or trivial factors. Hence, if we choose an element $u_{{i_j}} \in (U \cap C_{{i_j}}) \setminus (U \cap C_{{i_j+1}})$ whenever this factor is non\texttt{\symbol{45}}trivial, then we obtain a pcgs $Q = (u_{{i_1}}, \ldots, u_{{i_m}})$ of $U$. We say that $Q$ is an \emph{induced pcgs} with respect to \mbox{\texttt{\mdseries\slshape P}}. The pcgs \mbox{\texttt{\mdseries\slshape P}} is the \emph{parent pcgs} to the induced pcgs \mbox{\texttt{\mdseries\slshape Q}}. 

 Note that the pcgs $Q$ is induced with respect to \mbox{\texttt{\mdseries\slshape P}} if and only if the matrix of exponent vectors of the elements $u_{{i_j}}$ with respect to \mbox{\texttt{\mdseries\slshape P}} is in upper triangular form. Thus $Q$ is not unique in general. 

 In particular, the elements of an induced pcgs do \emph{not necessarily} have leading coefficient 1 relative to the inducing pcgs. The attribute \texttt{LeadCoeffsIGS} (\ref{LeadCoeffsIGS}) holds the leading coefficients in case they have to be renormed in an
algorithm. 

 Each induced pcgs is a pcgs and hence allows all elementary operations for
pcgs. On the other hand each pcgs could be transformed into an induced pcgs
for the group defined by the pcgs, but note that an arbitrary pcgs is in
general not an induced pcgs for technical reasons. 

 An induced pcgs is ``compatible'' with its parent, see \texttt{ParentPcgs} (\ref{ParentPcgs}). 

 In \cite{SOGOS} a ``non\texttt{\symbol{45}}commutative Gauss'' algorithm is described to compute an induced pcgs of a subgroup \mbox{\texttt{\mdseries\slshape U}} from a generating set of \mbox{\texttt{\mdseries\slshape U}}. For calling this in \textsf{GAP}, see \ref{InducedPcgs} to \ref{ExtendedPcgs}. 

 To create a subgroup generated by an induced pcgs such that the induced pcgs
gets stored automatically, use \texttt{SubgroupByPcgs} (\ref{SubgroupByPcgs}). 

\subsection{\textcolor{Chapter }{IsInducedPcgs}}
\logpage{[ 45, 7, 1 ]}\nobreak
\hyperdef{L}{X81FA878C854D63F8}{}
{\noindent\textcolor{FuncColor}{$\triangleright$\enspace\texttt{IsInducedPcgs({\mdseries\slshape pcgs})\index{IsInducedPcgs@\texttt{IsInducedPcgs}}
\label{IsInducedPcgs}
}\hfill{\scriptsize (Category)}}\\


 The category of induced pcgs. This a subcategory of pcgs. }

 

\subsection{\textcolor{Chapter }{InducedPcgsByPcSequence}}
\logpage{[ 45, 7, 2 ]}\nobreak
\hyperdef{L}{X83F6759184937F1B}{}
{\noindent\textcolor{FuncColor}{$\triangleright$\enspace\texttt{InducedPcgsByPcSequence({\mdseries\slshape pcgs, pcs})\index{InducedPcgsByPcSequence@\texttt{InducedPcgsByPcSequence}}
\label{InducedPcgsByPcSequence}
}\hfill{\scriptsize (operation)}}\\
\noindent\textcolor{FuncColor}{$\triangleright$\enspace\texttt{InducedPcgsByPcSequenceNC({\mdseries\slshape pcgs, pcs[, depths]})\index{InducedPcgsByPcSequenceNC@\texttt{InducedPcgsByPcSequenceNC}}
\label{InducedPcgsByPcSequenceNC}
}\hfill{\scriptsize (operation)}}\\


 If \mbox{\texttt{\mdseries\slshape pcs}} is a list of elements that form an induced pcgs with respect to \mbox{\texttt{\mdseries\slshape pcgs}} this operation returns an induced pcgs with these elements. 

 In the third version, the depths of \mbox{\texttt{\mdseries\slshape pcs}} with respect to \mbox{\texttt{\mdseries\slshape pcgs}} can be given (they are computed anew otherwise). }

 

\subsection{\textcolor{Chapter }{ParentPcgs}}
\logpage{[ 45, 7, 3 ]}\nobreak
\hyperdef{L}{X86308E80843BF9E5}{}
{\noindent\textcolor{FuncColor}{$\triangleright$\enspace\texttt{ParentPcgs({\mdseries\slshape pcgs})\index{ParentPcgs@\texttt{ParentPcgs}}
\label{ParentPcgs}
}\hfill{\scriptsize (attribute)}}\\


 returns the pcgs by which \mbox{\texttt{\mdseries\slshape pcgs}} was induced. If \mbox{\texttt{\mdseries\slshape pcgs}} was not induced, it simply returns \mbox{\texttt{\mdseries\slshape pcgs}}. 
\begin{Verbatim}[commandchars=!@|,fontsize=\small,frame=single,label=Example]
  !gapprompt@gap>| !gapinput@G := Group( (1,2,3,4),(1,2) );;|
  !gapprompt@gap>| !gapinput@P := Pcgs(G);;|
  !gapprompt@gap>| !gapinput@K := InducedPcgsByPcSequence( P, [(1,2,3,4),(1,3)(2,4)] );|
  Pcgs([ (1,2,3,4), (1,3)(2,4) ])
  !gapprompt@gap>| !gapinput@ParentPcgs( K );|
  Pcgs([ (3,4), (2,4,3), (1,4)(2,3), (1,3)(2,4) ])
  !gapprompt@gap>| !gapinput@IsInducedPcgs( K );|
  true
\end{Verbatim}
 }

 

\subsection{\textcolor{Chapter }{InducedPcgs}}
\logpage{[ 45, 7, 4 ]}\nobreak
\hyperdef{L}{X7F0EB20080590B23}{}
{\noindent\textcolor{FuncColor}{$\triangleright$\enspace\texttt{InducedPcgs({\mdseries\slshape pcgs, grp})\index{InducedPcgs@\texttt{InducedPcgs}}
\label{InducedPcgs}
}\hfill{\scriptsize (function)}}\\


 computes a pcgs for \mbox{\texttt{\mdseries\slshape grp}} which is induced by \mbox{\texttt{\mdseries\slshape pcgs}}. If \mbox{\texttt{\mdseries\slshape pcgs}} has a parent pcgs, then the result is induced with respect to this parent
pcgs. 

 \texttt{InducedPcgs} is a wrapper function only. Therefore, methods for computing an induced pcgs
should be installed for the operation \texttt{InducedPcgsOp}. }

 

\subsection{\textcolor{Chapter }{InducedPcgsByGenerators}}
\logpage{[ 45, 7, 5 ]}\nobreak
\hyperdef{L}{X8332F1197DF6FEDE}{}
{\noindent\textcolor{FuncColor}{$\triangleright$\enspace\texttt{InducedPcgsByGenerators({\mdseries\slshape pcgs, gens})\index{InducedPcgsByGenerators@\texttt{InducedPcgsByGenerators}}
\label{InducedPcgsByGenerators}
}\hfill{\scriptsize (operation)}}\\
\noindent\textcolor{FuncColor}{$\triangleright$\enspace\texttt{InducedPcgsByGeneratorsNC({\mdseries\slshape pcgs, gens})\index{InducedPcgsByGeneratorsNC@\texttt{InducedPcgsByGeneratorsNC}}
\label{InducedPcgsByGeneratorsNC}
}\hfill{\scriptsize (operation)}}\\


 returns an induced pcgs with respect to \mbox{\texttt{\mdseries\slshape pcgs}} for the subgroup generated by \mbox{\texttt{\mdseries\slshape gens}}. }

 

\subsection{\textcolor{Chapter }{InducedPcgsByPcSequenceAndGenerators}}
\logpage{[ 45, 7, 6 ]}\nobreak
\hyperdef{L}{X7AF82BD079D811E5}{}
{\noindent\textcolor{FuncColor}{$\triangleright$\enspace\texttt{InducedPcgsByPcSequenceAndGenerators({\mdseries\slshape pcgs, ind, gens})\index{InducedPcgsByPcSequenceAndGenerators@\texttt{Induced}\-\texttt{Pcgs}\-\texttt{By}\-\texttt{Pc}\-\texttt{Sequence}\-\texttt{And}\-\texttt{Generators}}
\label{InducedPcgsByPcSequenceAndGenerators}
}\hfill{\scriptsize (operation)}}\\


 returns an induced pcgs with respect to \mbox{\texttt{\mdseries\slshape pcgs}} of the subgroup generated by \mbox{\texttt{\mdseries\slshape ind}} and \mbox{\texttt{\mdseries\slshape gens}}. Here \mbox{\texttt{\mdseries\slshape ind}} must be an induced pcgs with respect to \mbox{\texttt{\mdseries\slshape pcgs}} (or a list of group elements that form such an igs) and it will be used as
initial sequence for the computation. 
\begin{Verbatim}[commandchars=!@|,fontsize=\small,frame=single,label=Example]
  !gapprompt@gap>| !gapinput@G := Group( (1,2,3,4),(1,2) );;  P := Pcgs(G);;|
  !gapprompt@gap>| !gapinput@I := InducedPcgsByGenerators( P, [(1,2,3,4)] );|
  Pcgs([ (1,2,3,4), (1,3)(2,4) ])
  !gapprompt@gap>| !gapinput@J := InducedPcgsByPcSequenceAndGenerators( P, I, [(1,2)] );|
  Pcgs([ (1,2,3,4), (2,4,3), (1,4)(2,3), (1,3)(2,4) ])
\end{Verbatim}
 }

 

\subsection{\textcolor{Chapter }{LeadCoeffsIGS}}
\logpage{[ 45, 7, 7 ]}\nobreak
\hyperdef{L}{X845FF8CA783D6CB3}{}
{\noindent\textcolor{FuncColor}{$\triangleright$\enspace\texttt{LeadCoeffsIGS({\mdseries\slshape igs})\index{LeadCoeffsIGS@\texttt{LeadCoeffsIGS}}
\label{LeadCoeffsIGS}
}\hfill{\scriptsize (attribute)}}\\


 This attribute is used to store leading coefficients with respect to the
parent pcgs. the \mbox{\texttt{\mdseries\slshape i}}\texttt{\symbol{45}}th entry {\textendash}if bound{\textendash} is the leading
exponent of the element of \mbox{\texttt{\mdseries\slshape igs}} that has depth \mbox{\texttt{\mdseries\slshape i}} in the parent. (It cannot be assigned to a component in the object created by \texttt{InducedPcgsByPcSequenceNC} (\ref{InducedPcgsByPcSequenceNC}) as the permutation group methods call it from within the postprocessing,
before this postprocessing however no coefficients may be computed.) }

 

\subsection{\textcolor{Chapter }{ExtendedPcgs}}
\logpage{[ 45, 7, 8 ]}\nobreak
\hyperdef{L}{X800287C680C5DEC3}{}
{\noindent\textcolor{FuncColor}{$\triangleright$\enspace\texttt{ExtendedPcgs({\mdseries\slshape N, gens})\index{ExtendedPcgs@\texttt{ExtendedPcgs}}
\label{ExtendedPcgs}
}\hfill{\scriptsize (operation)}}\\


 extends the pcgs \mbox{\texttt{\mdseries\slshape N}} (induced w.r.t. \mbox{\texttt{\mdseries\slshape home}}) to a new induced pcgs by prepending \mbox{\texttt{\mdseries\slshape gens}}. No checks are performed that this really yields an induced pcgs. }

 

\subsection{\textcolor{Chapter }{SubgroupByPcgs}}
\logpage{[ 45, 7, 9 ]}\nobreak
\hyperdef{L}{X817E16D67B31389B}{}
{\noindent\textcolor{FuncColor}{$\triangleright$\enspace\texttt{SubgroupByPcgs({\mdseries\slshape G, pcgs})\index{SubgroupByPcgs@\texttt{SubgroupByPcgs}}
\label{SubgroupByPcgs}
}\hfill{\scriptsize (operation)}}\\


 returns a subgroup of \mbox{\texttt{\mdseries\slshape G}} generated by the elements of \mbox{\texttt{\mdseries\slshape pcgs}}. }

 }

  
\section{\textcolor{Chapter }{Subgroups of Polycyclic Groups {\textendash} Canonical Pcgs}}\label{Subgroups of Polycyclic Groups - Canonical Pcgs}
\logpage{[ 45, 8, 0 ]}
\hyperdef{L}{X84068D2478C134C1}{}
{
  The induced pcgs \mbox{\texttt{\mdseries\slshape Q}} of \mbox{\texttt{\mdseries\slshape U}} is called \emph{canonical} if the matrix of exponent vectors contains normed vectors only and above each
leading entry in the matrix there are 0's only. The canonical pcgs of \mbox{\texttt{\mdseries\slshape U}} with respect to \mbox{\texttt{\mdseries\slshape P}} is unique and hence such pcgs can be used to compare subgroups. 

\subsection{\textcolor{Chapter }{IsCanonicalPcgs}}
\logpage{[ 45, 8, 1 ]}\nobreak
\hyperdef{L}{X80D122B986B42F80}{}
{\noindent\textcolor{FuncColor}{$\triangleright$\enspace\texttt{IsCanonicalPcgs({\mdseries\slshape pcgs})\index{IsCanonicalPcgs@\texttt{IsCanonicalPcgs}}
\label{IsCanonicalPcgs}
}\hfill{\scriptsize (property)}}\\


 An induced pcgs is canonical if the matrix of the exponent vectors of the
elements of \mbox{\texttt{\mdseries\slshape pcgs}} with respect to the \texttt{ParentPcgs} (\ref{ParentPcgs}) value of \mbox{\texttt{\mdseries\slshape pcgs}} is in Hermite normal form (see \cite{SOGOS}). While a subgroup can have various induced pcgs with respect to a parent
pcgs a canonical pcgs is unique. }

 

\subsection{\textcolor{Chapter }{CanonicalPcgs}}
\logpage{[ 45, 8, 2 ]}\nobreak
\hyperdef{L}{X816F6B4187032A10}{}
{\noindent\textcolor{FuncColor}{$\triangleright$\enspace\texttt{CanonicalPcgs({\mdseries\slshape pcgs})\index{CanonicalPcgs@\texttt{CanonicalPcgs}}
\label{CanonicalPcgs}
}\hfill{\scriptsize (attribute)}}\\


 returns the canonical pcgs corresponding to the induced pcgs \mbox{\texttt{\mdseries\slshape pcgs}}. 
\begin{Verbatim}[commandchars=!@|,fontsize=\small,frame=single,label=Example]
  !gapprompt@gap>| !gapinput@G := Group((1,2,3,4),(5,6,7));|
  Group([ (1,2,3,4), (5,6,7) ])
  !gapprompt@gap>| !gapinput@P := Pcgs(G);|
  Pcgs([ (5,6,7), (1,2,3,4), (1,3)(2,4) ])
  !gapprompt@gap>| !gapinput@I := InducedPcgsByPcSequence(P, [(5,6,7)*(1,3)(2,4),(1,3)(2,4)] );|
  Pcgs([ (1,3)(2,4)(5,6,7), (1,3)(2,4) ])
  !gapprompt@gap>| !gapinput@CanonicalPcgs(I);|
  Pcgs([ (5,6,7), (1,3)(2,4) ])
\end{Verbatim}
 }

 }

  
\section{\textcolor{Chapter }{Factor Groups of Polycyclic Groups {\textendash} Modulo Pcgs}}\label{Factor Groups of Polycyclic Groups - Modulo Pcgs}
\logpage{[ 45, 9, 0 ]}
\hyperdef{L}{X8294F5EF81B7ABA0}{}
{
  Let \mbox{\texttt{\mdseries\slshape N}} be a normal subgroup of \mbox{\texttt{\mdseries\slshape G}} such that \mbox{\texttt{\mdseries\slshape G/N}} is polycyclic with pcgs $(h_1 N, \ldots, h_r N)$. Then we call the sequence of preimages $(h_1, \ldots h_r)$ a \emph{modulo pcgs} of \mbox{\texttt{\mdseries\slshape G/N}}. \mbox{\texttt{\mdseries\slshape G}} is called the \emph{numerator} of the modulo pcgs and \mbox{\texttt{\mdseries\slshape N}} is the \emph{denominator} of the modulo pcgs. 

 Modulo pcgs are often used to facilitate efficient computations with factor
groups, since they allow computations with factor groups without formally
defining the factor group at all. 

 All elementary operations of pcgs, see Sections \ref{Elementary Operations for a Pcgs} and \ref{Elementary Operations for a Pcgs and an Element}, apply to modulo pcgs as well. However, it is in general not possible to
compute induced pcgs with respect to a modulo pcgs. 

 Two more elementary operations for modulo pcgs are \texttt{NumeratorOfModuloPcgs} (\ref{NumeratorOfModuloPcgs}) and \texttt{DenominatorOfModuloPcgs} (\ref{DenominatorOfModuloPcgs}). 

\subsection{\textcolor{Chapter }{ModuloPcgs}}
\logpage{[ 45, 9, 1 ]}\nobreak
\hyperdef{L}{X7FE689A37E559F66}{}
{\noindent\textcolor{FuncColor}{$\triangleright$\enspace\texttt{ModuloPcgs({\mdseries\slshape G, N})\index{ModuloPcgs@\texttt{ModuloPcgs}}
\label{ModuloPcgs}
}\hfill{\scriptsize (operation)}}\\


 returns a modulo pcgs for the factor $\mbox{\texttt{\mdseries\slshape G}}/\mbox{\texttt{\mdseries\slshape N}}$ which must be solvable, while \mbox{\texttt{\mdseries\slshape N}} may be non\texttt{\symbol{45}}solvable. \texttt{ModuloPcgs} will return \emph{a} pcgs for the factor, there is no guarantee that it will be ``compatible'' with any other pcgs. If this is required, the \texttt{mod} operator must be used on induced pcgs, see \texttt{\texttt{\symbol{92}}mod} (\ref{bSlashmod:for two pcgs}). }

 

\subsection{\textcolor{Chapter }{IsModuloPcgs}}
\logpage{[ 45, 9, 2 ]}\nobreak
\hyperdef{L}{X868207D77D09D915}{}
{\noindent\textcolor{FuncColor}{$\triangleright$\enspace\texttt{IsModuloPcgs({\mdseries\slshape obj})\index{IsModuloPcgs@\texttt{IsModuloPcgs}}
\label{IsModuloPcgs}
}\hfill{\scriptsize (Category)}}\\


 The category of modulo pcgs. Note that each pcgs is a modulo pcgs for the
trivial subgroup. }

 

\subsection{\textcolor{Chapter }{NumeratorOfModuloPcgs}}
\logpage{[ 45, 9, 3 ]}\nobreak
\hyperdef{L}{X8027CC9878031D74}{}
{\noindent\textcolor{FuncColor}{$\triangleright$\enspace\texttt{NumeratorOfModuloPcgs({\mdseries\slshape pcgs})\index{NumeratorOfModuloPcgs@\texttt{NumeratorOfModuloPcgs}}
\label{NumeratorOfModuloPcgs}
}\hfill{\scriptsize (attribute)}}\\


 returns a generating set for the numerator of the modulo pcgs \mbox{\texttt{\mdseries\slshape pcgs}}. }

 

\subsection{\textcolor{Chapter }{DenominatorOfModuloPcgs}}
\logpage{[ 45, 9, 4 ]}\nobreak
\hyperdef{L}{X87DBE2797D51B2F1}{}
{\noindent\textcolor{FuncColor}{$\triangleright$\enspace\texttt{DenominatorOfModuloPcgs({\mdseries\slshape pcgs})\index{DenominatorOfModuloPcgs@\texttt{DenominatorOfModuloPcgs}}
\label{DenominatorOfModuloPcgs}
}\hfill{\scriptsize (attribute)}}\\


 returns a generating set for the denominator of the modulo pcgs \mbox{\texttt{\mdseries\slshape pcgs}}. 
\begin{Verbatim}[commandchars=!@|,fontsize=\small,frame=single,label=Example]
  !gapprompt@gap>| !gapinput@G := Group( (1,2,3,4,5),(1,2) );|
  Group([ (1,2,3,4,5), (1,2) ])
  !gapprompt@gap>| !gapinput@P := ModuloPcgs(G, DerivedSubgroup(G) );|
  Pcgs([ (4,5) ])
  !gapprompt@gap>| !gapinput@NumeratorOfModuloPcgs(P);|
  [ (1,2,3,4,5), (1,2) ]
  !gapprompt@gap>| !gapinput@DenominatorOfModuloPcgs(P);|
  [ (1,3,2), (1,4,3), (2,5,4) ]
  !gapprompt@gap>| !gapinput@RelativeOrders(P);|
  [ 2 ]
  !gapprompt@gap>| !gapinput@ExponentsOfPcElement( P, (1,2,3,4,5) );|
  [ 0 ]
  !gapprompt@gap>| !gapinput@ExponentsOfPcElement( P, (4,5) );|
  [ 1 ]
\end{Verbatim}
 }

 

\subsection{\textcolor{Chapter }{\texttt{\symbol{92}}mod (for two pcgs)}}
\logpage{[ 45, 9, 5 ]}\nobreak
\hyperdef{L}{X7E923519832F2D08}{}
{\noindent\textcolor{FuncColor}{$\triangleright$\enspace\texttt{\texttt{\symbol{92}}mod({\mdseries\slshape P, I})\index{\texttt{\symbol{92}}mod@\texttt{\texttt{\symbol{92}}mod}!for two pcgs}
\label{bSlashmod:for two pcgs}
}\hfill{\scriptsize (method)}}\\


 Modulo Pcgs can also be built from compatible induced pcgs. Let $G$ be a group with pcgs \mbox{\texttt{\mdseries\slshape P}} and let \mbox{\texttt{\mdseries\slshape I}} be an induced pcgs of a normal subgroup $N$ of $G$. (Respectively: \mbox{\texttt{\mdseries\slshape P}} and \mbox{\texttt{\mdseries\slshape I}} are both induced with respect to the \emph{same} Pcgs.) Then we can compute a modulo pcgs of $G$ mod $N$ by 

 \mbox{\texttt{\mdseries\slshape P}} \texttt{mod} \mbox{\texttt{\mdseries\slshape I}} 

 Note that in this case we obtain the advantage that the values of \texttt{NumeratorOfModuloPcgs} (\ref{NumeratorOfModuloPcgs}) and \texttt{DenominatorOfModuloPcgs} (\ref{DenominatorOfModuloPcgs}) are just \mbox{\texttt{\mdseries\slshape P}} and \mbox{\texttt{\mdseries\slshape I}}, respectively, and hence are unique. 

 The resulting modulo pcgs will consist of a subset of \mbox{\texttt{\mdseries\slshape P}} and will be ``compatible'' with \mbox{\texttt{\mdseries\slshape P}} (or its parent). 

 
\begin{Verbatim}[commandchars=!@|,fontsize=\small,frame=single,label=Example]
  !gapprompt@gap>| !gapinput@G := Group((1,2,3,4));;|
  !gapprompt@gap>| !gapinput@P := Pcgs(G);|
  Pcgs([ (1,2,3,4), (1,3)(2,4) ])
  !gapprompt@gap>| !gapinput@I := InducedPcgsByGenerators(P, [(1,3)(2,4)]);|
  Pcgs([ (1,3)(2,4) ])
  !gapprompt@gap>| !gapinput@M := P mod I;|
  [ (1,2,3,4) ]
  !gapprompt@gap>| !gapinput@NumeratorOfModuloPcgs(M);|
  Pcgs([ (1,2,3,4), (1,3)(2,4) ])
  !gapprompt@gap>| !gapinput@DenominatorOfModuloPcgs(M);|
  Pcgs([ (1,3)(2,4) ])
\end{Verbatim}
 }

 

\subsection{\textcolor{Chapter }{CorrespondingGeneratorsByModuloPcgs}}
\logpage{[ 45, 9, 6 ]}\nobreak
\hyperdef{L}{X876A41F97FBA7754}{}
{\noindent\textcolor{FuncColor}{$\triangleright$\enspace\texttt{CorrespondingGeneratorsByModuloPcgs({\mdseries\slshape mpcgs, imgs})\index{CorrespondingGeneratorsByModuloPcgs@\texttt{CorrespondingGeneratorsByModuloPcgs}}
\label{CorrespondingGeneratorsByModuloPcgs}
}\hfill{\scriptsize (function)}}\\


 Let \mbox{\texttt{\mdseries\slshape mpcgs}} be a modulo pcgs for a factor of a group $G$ and let $U$ be a subgroup of $G$ generated by \mbox{\texttt{\mdseries\slshape imgs}} such that $U$ covers the factor for the modulo pcgs. Then this function computes elements in $U$ corresponding to the generators of the modulo pcgs. 

 Note that the computation of induced generating sets is not possible for some
modulo pcgs. }

 

\subsection{\textcolor{Chapter }{CanonicalPcgsByGeneratorsWithImages}}
\logpage{[ 45, 9, 7 ]}\nobreak
\hyperdef{L}{X8480852A7D49BC3F}{}
{\noindent\textcolor{FuncColor}{$\triangleright$\enspace\texttt{CanonicalPcgsByGeneratorsWithImages({\mdseries\slshape pcgs, gens, imgs})\index{CanonicalPcgsByGeneratorsWithImages@\texttt{CanonicalPcgsByGeneratorsWithImages}}
\label{CanonicalPcgsByGeneratorsWithImages}
}\hfill{\scriptsize (operation)}}\\


 computes a canonical, \mbox{\texttt{\mdseries\slshape pcgs}}\texttt{\symbol{45}}induced pcgs for the span of \mbox{\texttt{\mdseries\slshape gens}} and simultaneously does the same transformations on \mbox{\texttt{\mdseries\slshape imgs}}, preserving thus a correspondence between \mbox{\texttt{\mdseries\slshape gens}} and \mbox{\texttt{\mdseries\slshape imgs}}. This operation is used to represent homomorphisms from a pc group. }

 }

  
\section{\textcolor{Chapter }{Factor Groups of Polycyclic Groups in their Own Representation}}\label{Factor Groups of Polycyclic Groups in their Own Representation}
\logpage{[ 45, 10, 0 ]}
\hyperdef{L}{X8254C0F485F945BD}{}
{
  If substantial calculations are done in a factor it might be worth still to
construct the factor group in its own representation (for example by calling \texttt{PcGroupWithPcgs} (\ref{PcGroupWithPcgs}) on a modulo pcgs. 

 The following functions are intended for working with factor groups obtained
by factoring out the tail of a pcgs. They provide a way to map elements or
induced pcgs quickly in the factor (respectively to take preimages) without
the need to construct a homomorphism. 

 The setup is always a pcgs \mbox{\texttt{\mdseries\slshape pcgs}} of \mbox{\texttt{\mdseries\slshape G}} and a pcgs \mbox{\texttt{\mdseries\slshape fpcgs}} of a factor group $H = \mbox{\texttt{\mdseries\slshape G}}/\mbox{\texttt{\mdseries\slshape N}}$ which corresponds to a head of \mbox{\texttt{\mdseries\slshape pcgs}}. 

 No tests for validity of the input are performed. 

\subsection{\textcolor{Chapter }{ProjectedPcElement}}
\logpage{[ 45, 10, 1 ]}\nobreak
\hyperdef{L}{X806C2D827E04ACF3}{}
{\noindent\textcolor{FuncColor}{$\triangleright$\enspace\texttt{ProjectedPcElement({\mdseries\slshape pcgs, fpcgs, elm})\index{ProjectedPcElement@\texttt{ProjectedPcElement}}
\label{ProjectedPcElement}
}\hfill{\scriptsize (function)}}\\


 returns the image in \mbox{\texttt{\mdseries\slshape H}} of an element \mbox{\texttt{\mdseries\slshape elm}} of \mbox{\texttt{\mdseries\slshape G}}. }

 

\subsection{\textcolor{Chapter }{ProjectedInducedPcgs}}
\logpage{[ 45, 10, 2 ]}\nobreak
\hyperdef{L}{X82F39CCE7C928D3A}{}
{\noindent\textcolor{FuncColor}{$\triangleright$\enspace\texttt{ProjectedInducedPcgs({\mdseries\slshape pcgs, fpcgs, ipcgs})\index{ProjectedInducedPcgs@\texttt{ProjectedInducedPcgs}}
\label{ProjectedInducedPcgs}
}\hfill{\scriptsize (function)}}\\


 \mbox{\texttt{\mdseries\slshape ipcgs}} must be an induced pcgs with respect to \mbox{\texttt{\mdseries\slshape pcgs}}. This operation returns an induced pcgs with respect to \mbox{\texttt{\mdseries\slshape fpcgs}} consisting of the nontrivial images of \mbox{\texttt{\mdseries\slshape ipcgs}}. }

 

\subsection{\textcolor{Chapter }{LiftedPcElement}}
\logpage{[ 45, 10, 3 ]}\nobreak
\hyperdef{L}{X816813A078B93A6B}{}
{\noindent\textcolor{FuncColor}{$\triangleright$\enspace\texttt{LiftedPcElement({\mdseries\slshape pcgs, fpcgs, elm})\index{LiftedPcElement@\texttt{LiftedPcElement}}
\label{LiftedPcElement}
}\hfill{\scriptsize (function)}}\\


 returns a preimage in \mbox{\texttt{\mdseries\slshape G}} of an element \mbox{\texttt{\mdseries\slshape elm}} of \mbox{\texttt{\mdseries\slshape H}}. }

 

\subsection{\textcolor{Chapter }{LiftedInducedPcgs}}
\logpage{[ 45, 10, 4 ]}\nobreak
\hyperdef{L}{X83C60F1587577D65}{}
{\noindent\textcolor{FuncColor}{$\triangleright$\enspace\texttt{LiftedInducedPcgs({\mdseries\slshape pcgs, fpcgs, ipcgs, ker})\index{LiftedInducedPcgs@\texttt{LiftedInducedPcgs}}
\label{LiftedInducedPcgs}
}\hfill{\scriptsize (function)}}\\


 \mbox{\texttt{\mdseries\slshape ipcgs}} must be an induced pcgs with respect to \mbox{\texttt{\mdseries\slshape fpcgs}}. This operation returns an induced pcgs with respect to \mbox{\texttt{\mdseries\slshape pcgs}} consisting of the preimages of \mbox{\texttt{\mdseries\slshape ipcgs}}, appended by the elements in \mbox{\texttt{\mdseries\slshape ker}} (assuming there is a bijection of \mbox{\texttt{\mdseries\slshape pcgs}} mod \mbox{\texttt{\mdseries\slshape ker}} to \mbox{\texttt{\mdseries\slshape fpcgs}}). \mbox{\texttt{\mdseries\slshape ker}} might be a simple element list. }

 }

  
\section{\textcolor{Chapter }{Pcgs and Normal Series}}\label{Pcgs and Normal Series}
\logpage{[ 45, 11, 0 ]}
\hyperdef{L}{X83FE235E7B208EC0}{}
{
  By definition, a pcgs determines a pc series of its underlying group. However,
in many applications it will be necessary that this pc series refines a normal
series with certain properties; for example, a normal series with abelian
factors. 

 There are functions in \textsf{GAP} to compute a pcgs through a normal series with elementary abelian factors, a
central series or the lower p\texttt{\symbol{45}}central series. See also
Section \ref{Special Pcgs} for a more explicit possibility. 

\subsection{\textcolor{Chapter }{IsPcgsElementaryAbelianSeries}}
\logpage{[ 45, 11, 1 ]}\nobreak
\hyperdef{L}{X7E7E89C278DDE20D}{}
{\noindent\textcolor{FuncColor}{$\triangleright$\enspace\texttt{IsPcgsElementaryAbelianSeries({\mdseries\slshape pcgs})\index{IsPcgsElementaryAbelianSeries@\texttt{IsPcgsElementaryAbelianSeries}}
\label{IsPcgsElementaryAbelianSeries}
}\hfill{\scriptsize (property)}}\\


 returns \texttt{true} if the pcgs \mbox{\texttt{\mdseries\slshape pcgs}} refines an elementary abelian series. \texttt{IndicesEANormalSteps} (\ref{IndicesEANormalSteps}) then gives the indices in the Pcgs, at which the subgroups of this series
start. }

 

\subsection{\textcolor{Chapter }{PcgsElementaryAbelianSeries (for a group)}}
\logpage{[ 45, 11, 2 ]}\nobreak
\hyperdef{L}{X863A20B57EA37BAC}{}
{\noindent\textcolor{FuncColor}{$\triangleright$\enspace\texttt{PcgsElementaryAbelianSeries({\mdseries\slshape G})\index{PcgsElementaryAbelianSeries@\texttt{PcgsElementaryAbelianSeries}!for a group}
\label{PcgsElementaryAbelianSeries:for a group}
}\hfill{\scriptsize (attribute)}}\\
\noindent\textcolor{FuncColor}{$\triangleright$\enspace\texttt{PcgsElementaryAbelianSeries({\mdseries\slshape list})\index{PcgsElementaryAbelianSeries@\texttt{PcgsElementaryAbelianSeries}!for a list of normal subgroups}
\label{PcgsElementaryAbelianSeries:for a list of normal subgroups}
}\hfill{\scriptsize (attribute)}}\\


 computes a pcgs for \mbox{\texttt{\mdseries\slshape G}} that refines an elementary abelian series. \texttt{IndicesEANormalSteps} (\ref{IndicesEANormalSteps}) gives the indices in the pcgs, at which the normal subgroups of this series
start. The second variant returns a pcgs that runs through the normal
subgroups in the list \mbox{\texttt{\mdseries\slshape list}}. }

 

\subsection{\textcolor{Chapter }{IndicesEANormalSteps}}
\logpage{[ 45, 11, 3 ]}\nobreak
\hyperdef{L}{X7BCC1E2A80544CC7}{}
{\noindent\textcolor{FuncColor}{$\triangleright$\enspace\texttt{IndicesEANormalSteps({\mdseries\slshape pcgs})\index{IndicesEANormalSteps@\texttt{IndicesEANormalSteps}}
\label{IndicesEANormalSteps}
}\hfill{\scriptsize (attribute)}}\\
\noindent\textcolor{FuncColor}{$\triangleright$\enspace\texttt{IndicesEANormalStepsBounded({\mdseries\slshape pcgs, bound})\index{IndicesEANormalStepsBounded@\texttt{IndicesEANormalStepsBounded}}
\label{IndicesEANormalStepsBounded}
}\hfill{\scriptsize (function)}}\\


 Let \mbox{\texttt{\mdseries\slshape pcgs}} be a pcgs obtained as corresponding to a series of normal subgroups with
elementary abelian factors (for example from calling \texttt{PcgsElementaryAbelianSeries} (\ref{PcgsElementaryAbelianSeries:for a group})) Then \texttt{IndicesEANormalSteps} returns a sorted list of integers, indicating the tails of \mbox{\texttt{\mdseries\slshape pcgs}} which generate these normal subgroup of \mbox{\texttt{\mdseries\slshape G}}. If $i$ is an element of this list, $(g_i, \ldots, g_n)$ is a normal subgroup of \mbox{\texttt{\mdseries\slshape G}}. The list always starts with $1$ and ends with $n+1$. (These indices form \emph{one} series with elementary abelian subfactors, not necessarily the most refined
one.) 

 The attribute \texttt{EANormalSeriesByPcgs} (\ref{EANormalSeriesByPcgs}) returns the actual series of subgroups. 

 For arbitrary pcgs not obtained as belonging to a special series such a set of
indices not necessarily exists, and \texttt{IndicesEANormalSteps} is not guaranteed to work in this situation. 

 Typically, \texttt{IndicesEANormalSteps} is set by \texttt{PcgsElementaryAbelianSeries} (\ref{PcgsElementaryAbelianSeries:for a group}). 

 The variant \texttt{IndicesEANormalStepsBounded} will aim to ensure that no factor will be larger than the given bound. }

 

\subsection{\textcolor{Chapter }{EANormalSeriesByPcgs}}
\logpage{[ 45, 11, 4 ]}\nobreak
\hyperdef{L}{X7FCE308887F621FC}{}
{\noindent\textcolor{FuncColor}{$\triangleright$\enspace\texttt{EANormalSeriesByPcgs({\mdseries\slshape pcgs})\index{EANormalSeriesByPcgs@\texttt{EANormalSeriesByPcgs}}
\label{EANormalSeriesByPcgs}
}\hfill{\scriptsize (attribute)}}\\


 Let \mbox{\texttt{\mdseries\slshape pcgs}} be a pcgs obtained as corresponding to a series of normal subgroups with
elementary abelian factors (for example from calling \texttt{PcgsElementaryAbelianSeries} (\ref{PcgsElementaryAbelianSeries:for a group})). This attribute returns the actual series of normal subgroups, corresponding
to \texttt{IndicesEANormalSteps} (\ref{IndicesEANormalSteps}). }

 

\subsection{\textcolor{Chapter }{IsPcgsCentralSeries}}
\logpage{[ 45, 11, 5 ]}\nobreak
\hyperdef{L}{X79675266796D7254}{}
{\noindent\textcolor{FuncColor}{$\triangleright$\enspace\texttt{IsPcgsCentralSeries({\mdseries\slshape pcgs})\index{IsPcgsCentralSeries@\texttt{IsPcgsCentralSeries}}
\label{IsPcgsCentralSeries}
}\hfill{\scriptsize (property)}}\\


 returns \texttt{true} if the pcgs \mbox{\texttt{\mdseries\slshape pcgs}} refines an central elementary abelian series. \texttt{IndicesCentralNormalSteps} (\ref{IndicesCentralNormalSteps}) then gives the indices in the pcgs, at which the subgroups of this series
start. }

 

\subsection{\textcolor{Chapter }{PcgsCentralSeries}}
\logpage{[ 45, 11, 6 ]}\nobreak
\hyperdef{L}{X8187FCF483659E69}{}
{\noindent\textcolor{FuncColor}{$\triangleright$\enspace\texttt{PcgsCentralSeries({\mdseries\slshape G})\index{PcgsCentralSeries@\texttt{PcgsCentralSeries}}
\label{PcgsCentralSeries}
}\hfill{\scriptsize (attribute)}}\\


 computes a pcgs for \mbox{\texttt{\mdseries\slshape G}} that refines a central elementary abelian series. \texttt{IndicesCentralNormalSteps} (\ref{IndicesCentralNormalSteps}) gives the indices in the pcgs, at which the normal subgroups of this series
start. }

 

\subsection{\textcolor{Chapter }{IndicesCentralNormalSteps}}
\logpage{[ 45, 11, 7 ]}\nobreak
\hyperdef{L}{X7FB73FEB7BED5BFA}{}
{\noindent\textcolor{FuncColor}{$\triangleright$\enspace\texttt{IndicesCentralNormalSteps({\mdseries\slshape pcgs})\index{IndicesCentralNormalSteps@\texttt{IndicesCentralNormalSteps}}
\label{IndicesCentralNormalSteps}
}\hfill{\scriptsize (attribute)}}\\


 Let \mbox{\texttt{\mdseries\slshape pcgs}} be a pcgs obtained as corresponding to a series of normal subgroups with
central elementary abelian factors (for example from calling \texttt{PcgsCentralSeries} (\ref{PcgsCentralSeries})). Then \texttt{IndicesCentralNormalSteps} returns a sorted list of integers, indicating the tails of \mbox{\texttt{\mdseries\slshape pcgs}} which generate these normal subgroups of \mbox{\texttt{\mdseries\slshape G}}. If $i$ is an element of this list, $(g_i, \ldots, g_n)$ is a normal subgroup of \mbox{\texttt{\mdseries\slshape G}}. The list always starts with $1$ and ends with $n+1$. (These indices form \emph{one} series with central elementary abelian subfactors, not necessarily the most
refined one.) 

 The attribute \texttt{CentralNormalSeriesByPcgs} (\ref{CentralNormalSeriesByPcgs}) returns the actual series of subgroups. 

 For arbitrary pcgs not obtained as belonging to a special series such a set of
indices not necessarily exists, and \texttt{IndicesCentralNormalSteps} is not guaranteed to work in this situation. 

 Typically, \texttt{IndicesCentralNormalSteps} is set by \texttt{PcgsCentralSeries} (\ref{PcgsCentralSeries}). }

 

\subsection{\textcolor{Chapter }{CentralNormalSeriesByPcgs}}
\logpage{[ 45, 11, 8 ]}\nobreak
\hyperdef{L}{X82266ADA86B2A689}{}
{\noindent\textcolor{FuncColor}{$\triangleright$\enspace\texttt{CentralNormalSeriesByPcgs({\mdseries\slshape pcgs})\index{CentralNormalSeriesByPcgs@\texttt{CentralNormalSeriesByPcgs}}
\label{CentralNormalSeriesByPcgs}
}\hfill{\scriptsize (attribute)}}\\


 Let \mbox{\texttt{\mdseries\slshape pcgs}} be a pcgs obtained as corresponding to a series of normal subgroups with
central elementary abelian factors (for example from calling \texttt{PcgsCentralSeries} (\ref{PcgsCentralSeries})). This attribute returns the actual series of normal subgroups, corresponding
to \texttt{IndicesCentralNormalSteps} (\ref{IndicesCentralNormalSteps}). }

 

\subsection{\textcolor{Chapter }{IsPcgsPCentralSeriesPGroup}}
\logpage{[ 45, 11, 9 ]}\nobreak
\hyperdef{L}{X786E60AF7B61BF9E}{}
{\noindent\textcolor{FuncColor}{$\triangleright$\enspace\texttt{IsPcgsPCentralSeriesPGroup({\mdseries\slshape pcgs})\index{IsPcgsPCentralSeriesPGroup@\texttt{IsPcgsPCentralSeriesPGroup}}
\label{IsPcgsPCentralSeriesPGroup}
}\hfill{\scriptsize (property)}}\\


 returns \texttt{true} if the pcgs \mbox{\texttt{\mdseries\slshape pcgs}} refines a $p$\texttt{\symbol{45}}central elementary abelian series for a $p$\texttt{\symbol{45}}group. \texttt{IndicesPCentralNormalStepsPGroup} (\ref{IndicesPCentralNormalStepsPGroup}) then gives the indices in the pcgs, at which the subgroups of this series
start. }

 

\subsection{\textcolor{Chapter }{PcgsPCentralSeriesPGroup}}
\logpage{[ 45, 11, 10 ]}\nobreak
\hyperdef{L}{X86F19DBD7D346E7F}{}
{\noindent\textcolor{FuncColor}{$\triangleright$\enspace\texttt{PcgsPCentralSeriesPGroup({\mdseries\slshape G})\index{PcgsPCentralSeriesPGroup@\texttt{PcgsPCentralSeriesPGroup}}
\label{PcgsPCentralSeriesPGroup}
}\hfill{\scriptsize (attribute)}}\\


 computes a pcgs for the $p$\texttt{\symbol{45}}group \mbox{\texttt{\mdseries\slshape G}} that refines a $p$\texttt{\symbol{45}}central elementary abelian series. \texttt{IndicesPCentralNormalStepsPGroup} (\ref{IndicesPCentralNormalStepsPGroup}) gives the indices in the pcgs, at which the normal subgroups of this series
start. }

 

\subsection{\textcolor{Chapter }{IndicesPCentralNormalStepsPGroup}}
\logpage{[ 45, 11, 11 ]}\nobreak
\hyperdef{L}{X863968F08509E7D4}{}
{\noindent\textcolor{FuncColor}{$\triangleright$\enspace\texttt{IndicesPCentralNormalStepsPGroup({\mdseries\slshape pcgs})\index{IndicesPCentralNormalStepsPGroup@\texttt{IndicesPCentralNormalStepsPGroup}}
\label{IndicesPCentralNormalStepsPGroup}
}\hfill{\scriptsize (attribute)}}\\


 Let \mbox{\texttt{\mdseries\slshape pcgs}} be a pcgs obtained as corresponding to a series of normal subgroups with $p$\texttt{\symbol{45}}central elementary abelian factors (for example from
calling \texttt{PcgsPCentralSeriesPGroup} (\ref{PcgsPCentralSeriesPGroup})). Then \texttt{IndicesPCentralNormalStepsPGroup} returns a sorted list of integers, indicating the tails of \mbox{\texttt{\mdseries\slshape pcgs}} which generate these normal subgroups of \mbox{\texttt{\mdseries\slshape G}}. If $i$ is an element of this list, $(g_i, \ldots, g_n)$ is a normal subgroup of \mbox{\texttt{\mdseries\slshape G}}. The list always starts with $1$ and ends with $n+1$. (These indices form \emph{one} series with central elementary abelian subfactors, not necessarily the most
refined one.) 

 The attribute \texttt{PCentralNormalSeriesByPcgsPGroup} (\ref{PCentralNormalSeriesByPcgsPGroup}) returns the actual series of subgroups. 

 For arbitrary pcgs not obtained as belonging to a special series such a set of
indices not necessarily exists, and \texttt{IndicesPCentralNormalStepsPGroup} is not guaranteed to work in this situation. 

 Typically, \texttt{IndicesPCentralNormalStepsPGroup} is set by \texttt{PcgsPCentralSeriesPGroup} (\ref{PcgsPCentralSeriesPGroup}). }

 

\subsection{\textcolor{Chapter }{PCentralNormalSeriesByPcgsPGroup}}
\logpage{[ 45, 11, 12 ]}\nobreak
\hyperdef{L}{X7A92C9EA7BAF60CA}{}
{\noindent\textcolor{FuncColor}{$\triangleright$\enspace\texttt{PCentralNormalSeriesByPcgsPGroup({\mdseries\slshape pcgs})\index{PCentralNormalSeriesByPcgsPGroup@\texttt{PCentralNormalSeriesByPcgsPGroup}}
\label{PCentralNormalSeriesByPcgsPGroup}
}\hfill{\scriptsize (attribute)}}\\


 Let \mbox{\texttt{\mdseries\slshape pcgs}} be a pcgs obtained as corresponding to a series of normal subgroups with $p$\texttt{\symbol{45}}central elementary abelian factors (for example from
calling \texttt{PcgsPCentralSeriesPGroup} (\ref{PcgsPCentralSeriesPGroup})). This attribute returns the actual series of normal subgroups, corresponding
to \texttt{IndicesPCentralNormalStepsPGroup} (\ref{IndicesPCentralNormalStepsPGroup}). }

 

\subsection{\textcolor{Chapter }{IsPcgsChiefSeries}}
\logpage{[ 45, 11, 13 ]}\nobreak
\hyperdef{L}{X7EA5BC3B7FE9D98D}{}
{\noindent\textcolor{FuncColor}{$\triangleright$\enspace\texttt{IsPcgsChiefSeries({\mdseries\slshape pcgs})\index{IsPcgsChiefSeries@\texttt{IsPcgsChiefSeries}}
\label{IsPcgsChiefSeries}
}\hfill{\scriptsize (property)}}\\


 returns \texttt{true} if the pcgs \mbox{\texttt{\mdseries\slshape pcgs}} refines a chief series. \texttt{IndicesChiefNormalSteps} (\ref{IndicesChiefNormalSteps}) then gives the indices in the pcgs, at which the subgroups of this series
start. }

 

\subsection{\textcolor{Chapter }{PcgsChiefSeries}}
\logpage{[ 45, 11, 14 ]}\nobreak
\hyperdef{L}{X7E7326947EAE4BC9}{}
{\noindent\textcolor{FuncColor}{$\triangleright$\enspace\texttt{PcgsChiefSeries({\mdseries\slshape G})\index{PcgsChiefSeries@\texttt{PcgsChiefSeries}}
\label{PcgsChiefSeries}
}\hfill{\scriptsize (attribute)}}\\


 computes a pcgs for \mbox{\texttt{\mdseries\slshape G}} that refines a chief series. \texttt{IndicesChiefNormalSteps} (\ref{IndicesChiefNormalSteps}) gives the indices in the pcgs, at which the normal subgroups of this series
start. }

 

\subsection{\textcolor{Chapter }{IndicesChiefNormalSteps}}
\logpage{[ 45, 11, 15 ]}\nobreak
\hyperdef{L}{X7C05E84A78CA405E}{}
{\noindent\textcolor{FuncColor}{$\triangleright$\enspace\texttt{IndicesChiefNormalSteps({\mdseries\slshape pcgs})\index{IndicesChiefNormalSteps@\texttt{IndicesChiefNormalSteps}}
\label{IndicesChiefNormalSteps}
}\hfill{\scriptsize (attribute)}}\\


 Let \mbox{\texttt{\mdseries\slshape pcgs}} be a pcgs obtained as corresponding to a chief series for example from calling \texttt{PcgsChiefSeries} (\ref{PcgsChiefSeries})). Then \texttt{IndicesChiefNormalSteps} returns a sorted list of integers, indicating the tails of \mbox{\texttt{\mdseries\slshape pcgs}} which generate these normal subgroups of \mbox{\texttt{\mdseries\slshape G}}. If $i$ is an element of this list, $(g_i, \ldots, g_n)$ is a normal subgroup of \mbox{\texttt{\mdseries\slshape G}}. The list always starts with $1$ and ends with $n+1$. (These indices form \emph{one} series with elementary abelian subfactors, not necessarily the most refined
one.) 

 The attribute \texttt{ChiefNormalSeriesByPcgs} (\ref{ChiefNormalSeriesByPcgs}) returns the actual series of subgroups. 

 For arbitrary pcgs not obtained as belonging to a special series such a set of
indices not necessarily exists, and \texttt{IndicesChiefNormalSteps} is not guaranteed to work in this situation. 

 Typically, \texttt{IndicesChiefNormalSteps} is set by \texttt{PcgsChiefSeries} (\ref{PcgsChiefSeries}). }

 

\subsection{\textcolor{Chapter }{ChiefNormalSeriesByPcgs}}
\logpage{[ 45, 11, 16 ]}\nobreak
\hyperdef{L}{X83C5ABC587074B14}{}
{\noindent\textcolor{FuncColor}{$\triangleright$\enspace\texttt{ChiefNormalSeriesByPcgs({\mdseries\slshape pcgs})\index{ChiefNormalSeriesByPcgs@\texttt{ChiefNormalSeriesByPcgs}}
\label{ChiefNormalSeriesByPcgs}
}\hfill{\scriptsize (attribute)}}\\


 Let \mbox{\texttt{\mdseries\slshape pcgs}} be a pcgs obtained as corresponding to a chief series (for example from
calling \texttt{PcgsChiefSeries} (\ref{PcgsChiefSeries})). This attribute returns the actual series of normal subgroups, corresponding
to \texttt{IndicesChiefNormalSteps} (\ref{IndicesChiefNormalSteps}). 
\begin{Verbatim}[commandchars=!@|,fontsize=\small,frame=single,label=Example]
  !gapprompt@gap>| !gapinput@g:=Group((1,2,3,4),(1,2));;|
  !gapprompt@gap>| !gapinput@p:=PcgsElementaryAbelianSeries(g);|
  Pcgs([ (3,4), (2,4,3), (1,4)(2,3), (1,3)(2,4) ])
  !gapprompt@gap>| !gapinput@IndicesEANormalSteps(p);|
  [ 1, 2, 3, 5 ]
  !gapprompt@gap>| !gapinput@g:=Group((1,2,3,4),(1,5)(2,6)(3,7)(4,8));;|
  !gapprompt@gap>| !gapinput@p:=PcgsCentralSeries(g);|
  Pcgs([ (1,5)(2,6)(3,7)(4,8), (5,6,7,8), (5,7)(6,8),
    (1,4,3,2)(5,6,7,8), (1,3)(2,4)(5,7)(6,8) ])
  !gapprompt@gap>| !gapinput@IndicesCentralNormalSteps(p);|
  [ 1, 2, 4, 5, 6 ]
  !gapprompt@gap>| !gapinput@q:=PcgsPCentralSeriesPGroup(g);|
  Pcgs([ (1,5)(2,6)(3,7)(4,8), (5,6,7,8), (5,7)(6,8),
    (1,4,3,2)(5,6,7,8), (1,3)(2,4)(5,7)(6,8) ])
  !gapprompt@gap>| !gapinput@IndicesPCentralNormalStepsPGroup(q);|
  [ 1, 3, 5, 6 ]
\end{Verbatim}
 }

 

\subsection{\textcolor{Chapter }{IndicesNormalSteps}}
\logpage{[ 45, 11, 17 ]}\nobreak
\hyperdef{L}{X7A954E3887189842}{}
{\noindent\textcolor{FuncColor}{$\triangleright$\enspace\texttt{IndicesNormalSteps({\mdseries\slshape pcgs})\index{IndicesNormalSteps@\texttt{IndicesNormalSteps}}
\label{IndicesNormalSteps}
}\hfill{\scriptsize (attribute)}}\\


 returns the indices of \emph{all} steps in the pc series, which are normal in the group defined by the pcgs. 

 (In general, this function yields a slower performance than the more
specialized index functions for elementary abelian series etc.) }

 

\subsection{\textcolor{Chapter }{NormalSeriesByPcgs}}
\logpage{[ 45, 11, 18 ]}\nobreak
\hyperdef{L}{X7947B0FB87F44042}{}
{\noindent\textcolor{FuncColor}{$\triangleright$\enspace\texttt{NormalSeriesByPcgs({\mdseries\slshape pcgs})\index{NormalSeriesByPcgs@\texttt{NormalSeriesByPcgs}}
\label{NormalSeriesByPcgs}
}\hfill{\scriptsize (attribute)}}\\


 returns the subgroups the pc series, which are normal in the group defined by
the pcgs. 

 (In general, this function yields a slower performance than the more
specialized index functions for elementary abelian series etc.) }

 }

  
\section{\textcolor{Chapter }{Sum and Intersection of Pcgs}}\label{Sum and Intersection of Pcgs}
\logpage{[ 45, 12, 0 ]}
\hyperdef{L}{X7E624B4E8224DE2D}{}
{
     

\subsection{\textcolor{Chapter }{SumFactorizationFunctionPcgs}}
\logpage{[ 45, 12, 1 ]}\nobreak
\hyperdef{L}{X7833DAAA7C07CFD7}{}
{\noindent\textcolor{FuncColor}{$\triangleright$\enspace\texttt{SumFactorizationFunctionPcgs({\mdseries\slshape parentpcgs, n, u, kerpcgs})\index{SumFactorizationFunctionPcgs@\texttt{SumFactorizationFunctionPcgs}}
\label{SumFactorizationFunctionPcgs}
}\hfill{\scriptsize (operation)}}\\


 computes the sum and intersection of the lists \mbox{\texttt{\mdseries\slshape n}} and \mbox{\texttt{\mdseries\slshape u}} whose elements form modulo pcgs induced by \mbox{\texttt{\mdseries\slshape parentpcgs}} for two subgroups modulo a kernel given by \mbox{\texttt{\mdseries\slshape kerpcgs}}. If \mbox{\texttt{\mdseries\slshape kerpcgs}} is a tail if the \mbox{\texttt{\mdseries\slshape parent\texttt{\symbol{45}}pcgs}} it is sufficient to give the starting depth, this can be more efficient than
to construct an explicit pcgs. The factor group modulo \mbox{\texttt{\mdseries\slshape kerpcgs}} generated by \mbox{\texttt{\mdseries\slshape n}} must be elementary abelian and normal under \mbox{\texttt{\mdseries\slshape u}}. 

 The function returns a record with components 
\begin{description}
\item[{\texttt{sum}}]  elements that form a modulo pcgs for the span of both subgroups. 
\item[{\texttt{intersection}}]  elements that form a modulo pcgs for the intersection of both subgroups. 
\item[{\texttt{factorization}}]  a function that returns for an element \mbox{\texttt{\mdseries\slshape x}} in the span of \texttt{sum} a record with components \texttt{u} and \texttt{n} that give its decomposition. 
\end{description}
 

 The record components \texttt{sum} and \texttt{intersection} are \emph{not} pcgs but only lists of pc elements (to avoid unnecessary creation of induced
pcgs). }

 }

  
\section{\textcolor{Chapter }{Special Pcgs}}\label{Special Pcgs}
\logpage{[ 45, 13, 0 ]}
\hyperdef{L}{X83039CF97D27D819}{}
{
  In short, a special pcgs is a pcgs which has particularly nice properties, for
example it always refines an elementary abelian series, for $p$\texttt{\symbol{45}}groups it even refines a central series. These nice
properties permit particularly efficient algorithms. 

 Let \mbox{\texttt{\mdseries\slshape G}} be a finite polycyclic group. A \emph{special pcgs} of \mbox{\texttt{\mdseries\slshape G}} is a pcgs which is closely related to a Hall system and the maximal subgroups
of \mbox{\texttt{\mdseries\slshape G}}. These pcgs have been introduced by C. R. Leedham\texttt{\symbol{45}}Green
who also gave an algorithm to compute them. Improvements to this algorithm  are due to Bettina Eick. For a more detailed account of their definition the
reader is referred to \cite{Eick97} 

 To introduce the definition of special pcgs we first need to define the \emph{LG\texttt{\symbol{45}}series} and \emph{head complements} of a finite polycyclic group \mbox{\texttt{\mdseries\slshape G}}. Let $G = G_1 > G_2 > \ldots G_m > G_{{m+1}} = \{ 1 \}$ be the lower nilpotent series of $G$; that is, $G_i$ is the smallest normal subgroup of $G_{{i-1}}$ with nilpotent factor. To obtain the LG\texttt{\symbol{45}}series of \mbox{\texttt{\mdseries\slshape G}} we need to refine this series. Thus consider a factor $F_i := G_i / G_{{i+1}}$. Since $F_i$ is finite nilpotent, it is a direct product of its Sylow subgroups $F_i = P_{{i,1}} \cdots P_{{i,r_i}}$. For each Sylow $p_j$\texttt{\symbol{45}}subgroup $P_{{i,j}}$ we can consider its lower $p_j$\texttt{\symbol{45}}central series. To obtain a characteristic central series
with elementary abelian factors of $F_i$ we loop over its Sylow subgroups. Each time we consider $P_{{i,j}}$ in this process we take the next step of its lower $p_j$\texttt{\symbol{45}}central series into the series of $F_i$. If there is no next step, then we just skip the consideration of $P_{{i,j}}$. Note that the second term of the lower $p$\texttt{\symbol{45}}central series of a $p$\texttt{\symbol{45}}group is in fact its Frattini subgroup. Thus the Frattini
subgroup of $F_i$ is contained in the computed series of this group. We denote the Frattini
subgroup of $F_i = G_i / G_{{i+1}}$ by $G_i^* / G_{{i+1}}$. 

 The factors $G_i / G_i^*$ are called the heads of $G$, while the (possibly trivial) factors $G_i^* / G_{{i+1}}$ are the tails of $G$. A head complement of $G$ is a subgroup $U$ of $G$ such that $U / G_i^*$ is a complement to the head $G_i / G_i^*$ in $G / G_i^*$ for some $i$. 

 Now we are able to define a special pcgs of \mbox{\texttt{\mdseries\slshape G}}. It is a pcgs of \mbox{\texttt{\mdseries\slshape G}} with three additional properties. First, the pc series determined by the pcgs
refines the LG\texttt{\symbol{45}}series of \mbox{\texttt{\mdseries\slshape G}}. Second, a special pcgs \emph{exhibits} a Hall system of the group \mbox{\texttt{\mdseries\slshape G}}; that is, for each set of primes $\pi$ the elements of the pcgs with relative order in $\pi$ form a pcgs of a Hall $\pi$\texttt{\symbol{45}}subgroup in a Hall system of \mbox{\texttt{\mdseries\slshape G}}. Third, a special pcgs exhibits a head complement for each head of \mbox{\texttt{\mdseries\slshape G}}. 

 To record information about the LG\texttt{\symbol{45}}series with the special
pcgs we define the \emph{LGWeights} of the special pcgs. These weights are a list which contains a weight $w$ for each elements $g$ of the special pcgs. Such a weight $w$ represents the smallest subgroup of the LG\texttt{\symbol{45}}series
containing $g$. 

 Since the LG\texttt{\symbol{45}}series is defined in terms of the lower
nilpotent series, Sylow subgroups of the factors and lower $p$\texttt{\symbol{45}}central series of the Sylow subgroup, the weight $w$ is a triple. More precisely, $g$ is contained in the $w[1]$th term $U$ of the lower nilpotent series of \mbox{\texttt{\mdseries\slshape G}}, but not in the next smaller one $V$. Then $w[3]$ is a prime such that $g V$ is contained in the Sylow $w[3]$\texttt{\symbol{45}}subgroup $P/V$ of $U/V$. Moreover, $gV$ is contained in the $w[2]$th term of the lower $p$\texttt{\symbol{45}}central series of $P/V$. 

 There are two more attributes of a special pcgs containing information about
the LG\texttt{\symbol{45}}series: the list \emph{LGLayers} and the list \emph{LGFirst}. The list of layers corresponds to the elements of the special pcgs and
denotes the layer of the LG\texttt{\symbol{45}}series in which an element
lies. The list LGFirst corresponds to the LG\texttt{\symbol{45}}series and
gives the number of the first element in the special pcgs of the corresponding
subgroup. 

\subsection{\textcolor{Chapter }{IsSpecialPcgs}}
\logpage{[ 45, 13, 1 ]}\nobreak
\hyperdef{L}{X7C8A82FA786AC021}{}
{\noindent\textcolor{FuncColor}{$\triangleright$\enspace\texttt{IsSpecialPcgs({\mdseries\slshape obj})\index{IsSpecialPcgs@\texttt{IsSpecialPcgs}}
\label{IsSpecialPcgs}
}\hfill{\scriptsize (property)}}\\


 tests whether \mbox{\texttt{\mdseries\slshape obj}} is a special pcgs. }

 
\subsection{\textcolor{Chapter }{SpecialPcgs}}\logpage{[ 45, 13, 2 ]}
\hyperdef{L}{X827EB7767BACD023}{}
{
\noindent\textcolor{FuncColor}{$\triangleright$\enspace\texttt{SpecialPcgs({\mdseries\slshape pcgs})\index{SpecialPcgs@\texttt{SpecialPcgs}!for a pcgs}
\label{SpecialPcgs:for a pcgs}
}\hfill{\scriptsize (attribute)}}\\
\noindent\textcolor{FuncColor}{$\triangleright$\enspace\texttt{SpecialPcgs({\mdseries\slshape G})\index{SpecialPcgs@\texttt{SpecialPcgs}!for a group}
\label{SpecialPcgs:for a group}
}\hfill{\scriptsize (attribute)}}\\


 computes a special pcgs for the group defined by \mbox{\texttt{\mdseries\slshape pcgs}} or for \mbox{\texttt{\mdseries\slshape G}}. }

 

\subsection{\textcolor{Chapter }{LGWeights}}
\logpage{[ 45, 13, 3 ]}\nobreak
\hyperdef{L}{X82DC7CE682140588}{}
{\noindent\textcolor{FuncColor}{$\triangleright$\enspace\texttt{LGWeights({\mdseries\slshape pcgs})\index{LGWeights@\texttt{LGWeights}}
\label{LGWeights}
}\hfill{\scriptsize (attribute)}}\\


 returns the LGWeights of the special pcgs \mbox{\texttt{\mdseries\slshape pcgs}}. }

 

\subsection{\textcolor{Chapter }{LGLayers}}
\logpage{[ 45, 13, 4 ]}\nobreak
\hyperdef{L}{X824645C97E347EEE}{}
{\noindent\textcolor{FuncColor}{$\triangleright$\enspace\texttt{LGLayers({\mdseries\slshape pcgs})\index{LGLayers@\texttt{LGLayers}}
\label{LGLayers}
}\hfill{\scriptsize (attribute)}}\\


 returns the layers of the special pcgs \mbox{\texttt{\mdseries\slshape pcgs}}. }

 

\subsection{\textcolor{Chapter }{LGFirst}}
\logpage{[ 45, 13, 5 ]}\nobreak
\hyperdef{L}{X7A655F4C7D9AE130}{}
{\noindent\textcolor{FuncColor}{$\triangleright$\enspace\texttt{LGFirst({\mdseries\slshape pcgs})\index{LGFirst@\texttt{LGFirst}}
\label{LGFirst}
}\hfill{\scriptsize (attribute)}}\\


 returns the first indices for each layer of the special pcgs \mbox{\texttt{\mdseries\slshape pcgs}}. }

 

\subsection{\textcolor{Chapter }{LGLength}}
\logpage{[ 45, 13, 6 ]}\nobreak
\hyperdef{L}{X7C3912F77B12C8B6}{}
{\noindent\textcolor{FuncColor}{$\triangleright$\enspace\texttt{LGLength({\mdseries\slshape G})\index{LGLength@\texttt{LGLength}}
\label{LGLength}
}\hfill{\scriptsize (attribute)}}\\


 returns the length of the LG\texttt{\symbol{45}}series of the group \mbox{\texttt{\mdseries\slshape G}}, if \mbox{\texttt{\mdseries\slshape G}} is solvable, and \texttt{fail} otherwise. 
\begin{Verbatim}[commandchars=!@|,fontsize=\small,frame=single,label=Example]
  !gapprompt@gap>| !gapinput@G := SmallGroup( 96, 220 );|
  <pc group of size 96 with 6 generators>
  !gapprompt@gap>| !gapinput@spec := SpecialPcgs( G );|
  Pcgs([ f1, f2, f3, f4, f5, f6 ])
  !gapprompt@gap>| !gapinput@LGWeights(spec);|
  [ [ 1, 1, 2 ], [ 1, 1, 2 ], [ 1, 1, 2 ], [ 1, 1, 2 ], [ 1, 1, 3 ],
    [ 1, 2, 2 ] ]
  !gapprompt@gap>| !gapinput@LGLayers(spec);|
  [ 1, 1, 1, 1, 2, 3 ]
  !gapprompt@gap>| !gapinput@LGFirst(spec);|
  [ 1, 5, 6, 7 ]
  !gapprompt@gap>| !gapinput@LGLength( G );|
  3
  !gapprompt@gap>| !gapinput@p := SpecialPcgs( Pcgs( SmallGroup( 96, 120 ) ) );|
  Pcgs([ f1, f2, f3, f4, f5, f6 ])
  !gapprompt@gap>| !gapinput@LGWeights(p);|
  [ [ 1, 1, 2 ], [ 1, 1, 2 ], [ 1, 1, 2 ], [ 1, 2, 2 ], [ 1, 3, 2 ],
    [ 2, 1, 3 ] ]
\end{Verbatim}
 

 Thus the first group, \texttt{SmallGroup(96, 220)}, has a lower nilpotent series of length $1$; that is, the group is nilpotent. It is a direct product of its Sylow
subgroups. Moreover the Sylow $2$\texttt{\symbol{45}}subgroup is generated by the elements \texttt{f1, f2, f3, f4, f6}, and the Sylow $3$\texttt{\symbol{45}}subgroup is generated by \texttt{f5}. The lower $2$\texttt{\symbol{45}}central series of the Sylow $2$\texttt{\symbol{45}}subgroup has length $2$ and the second subgroup in this series is generated by \texttt{f6}. 

 The second group, \texttt{SmallGroup(96, 120)}, has a lower nilpotent series of length $2$ and hence is not nilpotent. The second subgroup in this series is just the
Sylow $3$\texttt{\symbol{45}}subgroup and it is generated by \texttt{f6}. The subgroup generated by \texttt{f1}, $\ldots$, \texttt{f5} is a Sylow $2$\texttt{\symbol{45}}subgroup of the group and also a head complement to the
second head of the group. Its lower $2$\texttt{\symbol{45}}central series has length $2$. 

 In this example the \texttt{FamilyPcgs} (\ref{FamilyPcgs}) value of the groups used was a special pcgs, but this is not necessarily the
case. For performance reasons it can be worth to enforce this,
see{\nobreakspace}\texttt{IsomorphismSpecialPcGroup} (\ref{IsomorphismSpecialPcGroup}). }

 

\subsection{\textcolor{Chapter }{IsInducedPcgsWrtSpecialPcgs}}
\logpage{[ 45, 13, 7 ]}\nobreak
\hyperdef{L}{X814C35BF7C9A8DEF}{}
{\noindent\textcolor{FuncColor}{$\triangleright$\enspace\texttt{IsInducedPcgsWrtSpecialPcgs({\mdseries\slshape pcgs})\index{IsInducedPcgsWrtSpecialPcgs@\texttt{IsInducedPcgsWrtSpecialPcgs}}
\label{IsInducedPcgsWrtSpecialPcgs}
}\hfill{\scriptsize (property)}}\\


 tests whether \mbox{\texttt{\mdseries\slshape pcgs}} is induced with respect to a special pcgs. }

 

\subsection{\textcolor{Chapter }{InducedPcgsWrtSpecialPcgs}}
\logpage{[ 45, 13, 8 ]}\nobreak
\hyperdef{L}{X7C14AE5C82FB0771}{}
{\noindent\textcolor{FuncColor}{$\triangleright$\enspace\texttt{InducedPcgsWrtSpecialPcgs({\mdseries\slshape G})\index{InducedPcgsWrtSpecialPcgs@\texttt{InducedPcgsWrtSpecialPcgs}}
\label{InducedPcgsWrtSpecialPcgs}
}\hfill{\scriptsize (attribute)}}\\


 computes an induced pcgs with respect to the special pcgs of the parent of \mbox{\texttt{\mdseries\slshape G}}. 

 \texttt{InducedPcgsWrtSpecialPcgs} will return a pcgs induced by \emph{a} special pcgs (which might differ from the one you had in mind). If you need an
induced pcgs compatible with a \emph{given} special pcgs use \texttt{InducedPcgs} (\ref{InducedPcgs}) for this special pcgs. }

 }

  
\section{\textcolor{Chapter }{Action on Subfactors Defined by a Pcgs}}\label{Action on Subfactors Defined by a Pcgs}
\logpage{[ 45, 14, 0 ]}
\hyperdef{L}{X7E86EB517DC08809}{}
{
  When working with a polycyclic group, one often needs to compute matrix
operations of the group on a factor of the group. For this purpose there are
the functions described in \ref{VectorSpaceByPcgsOfElementaryAbelianGroup} to \ref{LinearActionLayer}. 

 In certain situations, for example within the computation of conjugacy classes
of finite soluble groups as described in \cite{MeckyNeubueser89}, affine actions of groups are required. For this purpose we introduce the
functions \texttt{AffineAction} (\ref{AffineAction}) and \texttt{AffineActionLayer} (\ref{AffineActionLayer}). 

\subsection{\textcolor{Chapter }{VectorSpaceByPcgsOfElementaryAbelianGroup}}
\logpage{[ 45, 14, 1 ]}\nobreak
\hyperdef{L}{X7A9BB9D0817CA949}{}
{\noindent\textcolor{FuncColor}{$\triangleright$\enspace\texttt{VectorSpaceByPcgsOfElementaryAbelianGroup({\mdseries\slshape mpcgs, fld})\index{VectorSpaceByPcgsOfElementaryAbelianGroup@\texttt{Vector}\-\texttt{Space}\-\texttt{By}\-\texttt{Pcgs}\-\texttt{Of}\-\texttt{Elementary}\-\texttt{Abelian}\-\texttt{Group}}
\label{VectorSpaceByPcgsOfElementaryAbelianGroup}
}\hfill{\scriptsize (function)}}\\


 returns the vector space over \mbox{\texttt{\mdseries\slshape fld}} corresponding to the modulo pcgs \mbox{\texttt{\mdseries\slshape mpcgs}}. Note that \mbox{\texttt{\mdseries\slshape mpcgs}} has to define an elementary abelian $p$\texttt{\symbol{45}}group where $p$ is the characteristic of \mbox{\texttt{\mdseries\slshape fld}}. }

 

\subsection{\textcolor{Chapter }{LinearAction}}
\logpage{[ 45, 14, 2 ]}\nobreak
\hyperdef{L}{X81FC09DD7FC06C6E}{}
{\noindent\textcolor{FuncColor}{$\triangleright$\enspace\texttt{LinearAction({\mdseries\slshape gens, basisvectors, linear})\index{LinearAction@\texttt{LinearAction}}
\label{LinearAction}
}\hfill{\scriptsize (operation)}}\\
\noindent\textcolor{FuncColor}{$\triangleright$\enspace\texttt{LinearOperation({\mdseries\slshape gens, basisvectors, linear})\index{LinearOperation@\texttt{LinearOperation}}
\label{LinearOperation}
}\hfill{\scriptsize (operation)}}\\


 returns a list of matrices, one for each element of \mbox{\texttt{\mdseries\slshape gens}}, which corresponds to the matrix action of the elements in \mbox{\texttt{\mdseries\slshape gens}} on the basis \mbox{\texttt{\mdseries\slshape basisvectors}} via \mbox{\texttt{\mdseries\slshape linear}}. }

 

\subsection{\textcolor{Chapter }{LinearActionLayer}}
\logpage{[ 45, 14, 3 ]}\nobreak
\hyperdef{L}{X7C2135B98732BBC3}{}
{\noindent\textcolor{FuncColor}{$\triangleright$\enspace\texttt{LinearActionLayer({\mdseries\slshape G, gens, pcgs})\index{LinearActionLayer@\texttt{LinearActionLayer}}
\label{LinearActionLayer}
}\hfill{\scriptsize (function)}}\\
\noindent\textcolor{FuncColor}{$\triangleright$\enspace\texttt{LinearOperationLayer({\mdseries\slshape G, gens, pcgs})\index{LinearOperationLayer@\texttt{LinearOperationLayer}}
\label{LinearOperationLayer}
}\hfill{\scriptsize (function)}}\\


 returns a list of matrices, one for each element of \mbox{\texttt{\mdseries\slshape gens}}, which corresponds to the matrix action of \mbox{\texttt{\mdseries\slshape G}} on the vector space corresponding to the modulo pcgs \mbox{\texttt{\mdseries\slshape pcgs}}. }

 

\subsection{\textcolor{Chapter }{AffineAction}}
\logpage{[ 45, 14, 4 ]}\nobreak
\hyperdef{L}{X79C2D6BF7DD69ED6}{}
{\noindent\textcolor{FuncColor}{$\triangleright$\enspace\texttt{AffineAction({\mdseries\slshape gens, basisvectors, linear, transl})\index{AffineAction@\texttt{AffineAction}}
\label{AffineAction}
}\hfill{\scriptsize (operation)}}\\


 return a list of matrices, one for each element of \mbox{\texttt{\mdseries\slshape gens}}, which corresponds to the affine action of the elements in \mbox{\texttt{\mdseries\slshape gens}} on the basis \mbox{\texttt{\mdseries\slshape basisvectors}} via \mbox{\texttt{\mdseries\slshape linear}} with translation \mbox{\texttt{\mdseries\slshape transl}}. }

 

\subsection{\textcolor{Chapter }{AffineActionLayer}}
\logpage{[ 45, 14, 5 ]}\nobreak
\hyperdef{L}{X7E4CB1358524497B}{}
{\noindent\textcolor{FuncColor}{$\triangleright$\enspace\texttt{AffineActionLayer({\mdseries\slshape G, gens, pcgs, transl})\index{AffineActionLayer@\texttt{AffineActionLayer}}
\label{AffineActionLayer}
}\hfill{\scriptsize (function)}}\\


 returns a list of matrices, one for each element of \mbox{\texttt{\mdseries\slshape gens}}, which corresponds to the affine action of \mbox{\texttt{\mdseries\slshape G}} on the vector space corresponding to the modulo pcgs \mbox{\texttt{\mdseries\slshape pcgs}} with translation \mbox{\texttt{\mdseries\slshape transl}}. 
\begin{Verbatim}[commandchars=!@|,fontsize=\small,frame=single,label=Example]
  !gapprompt@gap>| !gapinput@G := SmallGroup( 96, 51 );|
  <pc group of size 96 with 6 generators>
  !gapprompt@gap>| !gapinput@spec := SpecialPcgs( G );|
  Pcgs([ f1, f2, f3, f4, f5, f6 ])
  !gapprompt@gap>| !gapinput@LGWeights( spec );|
  [ [ 1, 1, 2 ], [ 1, 1, 2 ], [ 1, 1, 3 ], [ 1, 2, 2 ], [ 1, 2, 2 ],
    [ 1, 3, 2 ] ]
  !gapprompt@gap>| !gapinput@mpcgs := InducedPcgsByPcSequence( spec, spec{[4,5,6]} );|
  Pcgs([ f4, f5, f6 ])
  !gapprompt@gap>| !gapinput@npcgs := InducedPcgsByPcSequence( spec, spec{[6]} );|
  Pcgs([ f6 ])
  !gapprompt@gap>| !gapinput@modu := mpcgs mod npcgs;|
  [ f4, f5 ]
  !gapprompt@gap>| !gapinput@mat:=LinearActionLayer( G, spec{[1,2,3]}, modu );|
  [ <an immutable 2x2 matrix over GF2>,
    <an immutable 2x2 matrix over GF2>,
    <an immutable 2x2 matrix over GF2> ]
  !gapprompt@gap>| !gapinput@Print( mat, "\n" );|
  [ [ [ Z(2)^0, 0*Z(2) ], [ 0*Z(2), Z(2)^0 ] ],
    [ [ Z(2)^0, 0*Z(2) ], [ 0*Z(2), Z(2)^0 ] ],
    [ [ Z(2)^0, 0*Z(2) ], [ 0*Z(2), Z(2)^0 ] ] ]
\end{Verbatim}
 }

 }

  
\section{\textcolor{Chapter }{Orbit Stabilizer Methods for Polycyclic Groups}}\label{Orbit Stabilizer Methods for Polycyclic Groups}
\logpage{[ 45, 15, 0 ]}
\hyperdef{L}{X7EEA8D638492F432}{}
{
  If a pcgs \mbox{\texttt{\mdseries\slshape pcgs}} is known for a group \mbox{\texttt{\mdseries\slshape G}}, then orbits and stabilizers can be computed by a special method which is
particularly efficient. Note that within this function only the elements in \mbox{\texttt{\mdseries\slshape pcgs}} and the relative orders of \mbox{\texttt{\mdseries\slshape pcgs}} are needed. Hence this function works effectively even if the elementary
operations for \mbox{\texttt{\mdseries\slshape pcgs}} are slow. 

\subsection{\textcolor{Chapter }{StabilizerPcgs}}
\logpage{[ 45, 15, 1 ]}\nobreak
\hyperdef{L}{X7CFCCF607A30B5EE}{}
{\noindent\textcolor{FuncColor}{$\triangleright$\enspace\texttt{StabilizerPcgs({\mdseries\slshape pcgs, pnt[, acts][, act]})\index{StabilizerPcgs@\texttt{StabilizerPcgs}}
\label{StabilizerPcgs}
}\hfill{\scriptsize (function)}}\\


 computes the stabilizer in the group generated by \mbox{\texttt{\mdseries\slshape pcgs}} of the point \mbox{\texttt{\mdseries\slshape pnt}}. If given, \mbox{\texttt{\mdseries\slshape acts}} are elements by which \mbox{\texttt{\mdseries\slshape pcgs}} acts, \mbox{\texttt{\mdseries\slshape act}} is the acting function. This function returns a pcgs for the stabilizer which
is induced by the \texttt{ParentPcgs} of \mbox{\texttt{\mdseries\slshape pcgs}}, that is it is compatible with \mbox{\texttt{\mdseries\slshape pcgs}}. }

 

\subsection{\textcolor{Chapter }{Pcgs{\textunderscore}OrbitStabilizer}}
\logpage{[ 45, 15, 2 ]}\nobreak
\hyperdef{L}{X7A87E72F86813132}{}
{\noindent\textcolor{FuncColor}{$\triangleright$\enspace\texttt{Pcgs{\textunderscore}OrbitStabilizer({\mdseries\slshape pcgs, domain, pnt, oprs, opr})\index{Pcgs{\textunderscore}OrbitStabilizer@\texttt{Pcgs{\textunderscore}}\-\texttt{Orbit}\-\texttt{Stabilizer}}
\label{PcgsuScoreOrbitStabilizer}
}\hfill{\scriptsize (function)}}\\


 runs a solvable group orbit\texttt{\symbol{45}}stabilizer algorithm on \mbox{\texttt{\mdseries\slshape pnt}} with \mbox{\texttt{\mdseries\slshape pcgs}} acting via the images \mbox{\texttt{\mdseries\slshape oprs}} and the operation function \mbox{\texttt{\mdseries\slshape opr}}. The domain \mbox{\texttt{\mdseries\slshape domain}} can be used to speed up search, if it is not known, \texttt{false} can be given instead. The function returns a record with components \texttt{orbit}, \texttt{stabpcgs} and \texttt{lengths}, the latter indicating the lengths of the orbit whenever it got extended.
This can be used to recompute transversal elements. This function should be
used only inside algorithms when speed is essential. }

 }

  
\section{\textcolor{Chapter }{Operations which have Special Methods for Groups with Pcgs}}\label{Operations which have Special Methods for Groups with Pcgs}
\logpage{[ 45, 16, 0 ]}
\hyperdef{L}{X7A19DF1E7E841074}{}
{
  \index{IsNilpotent@\texttt{IsNilpotent}!for groups with pcgs} \index{IsSupersolvable@\texttt{IsSupersolvable}!for groups with pcgs} \index{Size@\texttt{Size}!for groups with pcgs} \index{CompositionSeries@\texttt{CompositionSeries}!for groups with pcgs} \index{ConjugacyClasses@\texttt{ConjugacyClasses}!for groups with pcgs} \index{Centralizer@\texttt{Centralizer}!for groups with pcgs} \index{FrattiniSubgroup@\texttt{FrattiniSubgroup}!for groups with pcgs} \index{PrefrattiniSubgroup@\texttt{PrefrattiniSubgroup}!for groups with pcgs} \index{MaximalSubgroups@\texttt{MaximalSubgroups}!for groups with pcgs} \index{HallSystem@\texttt{HallSystem}!for groups with pcgs} \index{MinimalGeneratingSet@\texttt{MinimalGeneratingSet}!for groups with pcgs} \index{Centre@\texttt{Centre}!for groups with pcgs} \index{Intersection@\texttt{Intersection}!for groups with pcgs} \index{AutomorphismGroup@\texttt{AutomorphismGroup}!for groups with pcgs} \index{IrreducibleModules@\texttt{IrreducibleModules}!for groups with pcgs} For the following operations there are special methods for groups with pcgs
installed: 

 \texttt{IsNilpotentGroup} (\ref{IsNilpotentGroup}), \texttt{IsSupersolvableGroup} (\ref{IsSupersolvableGroup}), \texttt{Size} (\ref{Size}), \texttt{CompositionSeries} (\ref{CompositionSeries}), \texttt{ConjugacyClasses} (\ref{ConjugacyClasses:attribute}), \texttt{Centralizer} (\ref{Centralizer:for a magma and an element}), \texttt{FrattiniSubgroup} (\ref{FrattiniSubgroup}), \texttt{PrefrattiniSubgroup} (\ref{PrefrattiniSubgroup}), \texttt{MaximalSubgroups} (\ref{MaximalSubgroups}) and related operations, \texttt{HallSystem} (\ref{HallSystem}) and related operations, \texttt{MinimalGeneratingSet} (\ref{MinimalGeneratingSet}), \texttt{Centre} (\ref{Centre}), \texttt{Intersection} (\ref{Intersection:for various collections}), \texttt{AutomorphismGroup} (\ref{AutomorphismGroup}), \texttt{IrreducibleModules} (\ref{IrreducibleModules}). }

  
\section{\textcolor{Chapter }{Conjugacy Classes in Solvable Groups}}\label{Conjugacy Classes in Solvable Groups}
\logpage{[ 45, 17, 0 ]}
\hyperdef{L}{X79DCCF6D80351859}{}
{
  There are a variety of algorithms to compute conjugacy classes and
centralizers in solvable groups via epimorphic images (\cite{FelschNeubueser79}, \cite{MeckyNeubueser89}, \cite{Theissen93}). Usually these are only invoked as methods, but it is possible to access the
algorithm directly. 

 

\subsection{\textcolor{Chapter }{ClassesSolvableGroup}}
\logpage{[ 45, 17, 1 ]}\nobreak
\hyperdef{L}{X79593F667A68A21D}{}
{\noindent\textcolor{FuncColor}{$\triangleright$\enspace\texttt{ClassesSolvableGroup({\mdseries\slshape G, mode[, opt]})\index{ClassesSolvableGroup@\texttt{ClassesSolvableGroup}}
\label{ClassesSolvableGroup}
}\hfill{\scriptsize (function)}}\\


 computes conjugacy classes and centralizers in solvable groups. \mbox{\texttt{\mdseries\slshape G}} is the acting group. \mbox{\texttt{\mdseries\slshape mode}} indicates the type of the calculation: 

 0 Conjugacy classes 

 4 Conjugacy test for the two elements in \mbox{\texttt{\mdseries\slshape opt}}\texttt{.candidates} 

 In mode 0 the function returns a list of records containing components \mbox{\texttt{\mdseries\slshape representative}} and \mbox{\texttt{\mdseries\slshape centralizer}}. In mode 4 it returns a conjugating element. 

 The optional record \mbox{\texttt{\mdseries\slshape opt}} may contain the following components that will affect the algorithm's
behaviour: 

 
\begin{description}
\item[{\texttt{pcgs}}]  is a pcgs that will be used for the calculation. The attribute \texttt{EANormalSeriesByPcgs} (\ref{EANormalSeriesByPcgs}) must return an appropriate series of normal subgroups with elementary abelian
factors among them. The algorithm will step down this series. In the case of
the calculation of rational classes, it must be a pcgs refining a central
series. 
\item[{\texttt{candidates}}]  is a list of elements for which canonical representatives are to be computed
or for which a conjugacy test is performed. Both elements must lie in \mbox{\texttt{\mdseries\slshape G}}, but this is not tested. In mode 4 these elements must be given. In mode 0 a
list of classes corresponding to \texttt{candidates} is returned (which may contain duplicates). The \texttt{representative}s chosen are canonical with respect to \texttt{pcgs}. The records returned also contain components \texttt{operator} such that \texttt{candidate \texttt{\symbol{94}} operator = representative}. 
\item[{\texttt{consider}}]  is a function \texttt{consider( fhome, rep, cenp, K, L )}. Here \texttt{fhome} is a home pcgs for the factor group $F$ in which the calculation currently takes place, \texttt{rep} is an element of the factor and \texttt{cenp} is a pcgs for the centralizer of \texttt{rep} modulo \texttt{K}. In mode 0, when lifting from $F$/\texttt{K} to $F$/\texttt{L} (note: for efficiency reasons, $F$ can be different from \mbox{\texttt{\mdseries\slshape G}} or \texttt{L} might be not trivial) this function is called before performing the actual
lifting and only those representatives for which it returns \texttt{true} are passed to the next level. This permits for example the calculation of only
those classes with small centralizers or classes of restricted orders. 
\end{description}
 }

 

\subsection{\textcolor{Chapter }{CentralizerSizeLimitConsiderFunction}}
\logpage{[ 45, 17, 2 ]}\nobreak
\hyperdef{L}{X7B358D3B7E236973}{}
{\noindent\textcolor{FuncColor}{$\triangleright$\enspace\texttt{CentralizerSizeLimitConsiderFunction({\mdseries\slshape sz})\index{CentralizerSizeLimitConsiderFunction@\texttt{Centralizer}\-\texttt{Size}\-\texttt{Limit}\-\texttt{Consider}\-\texttt{Function}}
\label{CentralizerSizeLimitConsiderFunction}
}\hfill{\scriptsize (function)}}\\


 returns a function (with arguments \texttt{fhome}, \texttt{rep}, \texttt{cen}, \texttt{K}, \texttt{L}) that can be used in \texttt{ClassesSolvableGroup} (\ref{ClassesSolvableGroup}) as the \texttt{consider} component of the options record. It will restrict the lifting to those
classes, for which the size of the centralizer (in the factor) is at most \mbox{\texttt{\mdseries\slshape sz}}. 

 See also \texttt{SubgroupsSolvableGroup} (\ref{SubgroupsSolvableGroup}). }

 }

 }

          
\chapter{\textcolor{Chapter }{Pc Groups}}\label{Pc Groups}
\logpage{[ 46, 0, 0 ]}
\hyperdef{L}{X7EAD57C97EBF7E67}{}
{
     Pc groups are polycyclic groups that use the polycyclic presentation for
element arithmetic. This presentation gives them a ``natural'' pcgs, the \texttt{FamilyPcgs} (\ref{FamilyPcgs}) with respect to which pcgs operations as described in chapter{\nobreakspace}\ref{Polycyclic Groups} are particularly efficient. 

 Let $G$ be a polycyclic group with pcgs $P = (g_1, \ldots, g_n)$ and corresponding relative orders $(r_1, \ldots, r_n)$. Recall that the $r_i$ are positive integers or infinity and let $I$ be the set of indices $i$ with $r_i$ a positive integer. Then $G$ has a finite presentation on the generators $g_1, \ldots, g_n$ with relations of the following form. \begin{center}
\begin{tabular}{lcl}$g_i^{{r_i}}$&
=&
$g_{{i+1}}^{a(i,i,i+1)} \cdots g_n^{a(i,i,n)}$\\
&
&
for $1 \leq i \leq n$ and $i \in I$\\
$g_i^{{-1}} g_j g_i$&
=&
$g_{{i+1}}^{a(i,j,i+1)} \cdots g_n^{a(i,j,n)}$\\
&
&
for $1 \leq i < j \leq n$\\
\end{tabular}\\[2mm]
\end{center}

 For infinite groups we need additionally \begin{center}
\begin{tabular}{lcl}$g_i^{{-1}} g_j^{{-1}} g_i$&
=&
$g_{{i+1}}^{b(i,j,i+1)} \cdots g_n^{b(i,j,n)}$\\
&
&
for $1 \leq i < j \leq n$ and $j \not \in I$\\
$g_i g_j g_i^{{-1}}$&
=&
$g_{{i+1}}^{c(i,j,i+1)} \cdots g_n^{c(i,j,n)}$\\
&
&
for $1 \leq i < j \leq n$ and $i \not \in I$\\
$g_i g_j^{{-1}} g_i^{{-1}}$&
=&
$g_{{i+1}}^{d(i,j,i+1)} \cdots g_n^{d(i,j,n)}$\\
&
&
for $1 \leq i < j \leq n$ and $i, j, \not \in I$\\
\end{tabular}\\[2mm]
\end{center}

 Here the right hand sides are assumed to be words in normal form; that is, for $k \in I$ we have for all exponents $0 \leq a(i,j,k), b(i,j,k), c(i,j,k), d(i,j,k) < r_k$. 

 A finite presentation of this type is called a \emph{power\texttt{\symbol{45}}conjugate presentation} and a \emph{pc group} is a polycyclic group defined by a power\texttt{\symbol{45}}conjugate
presentation. Instead of conjugates we could just as well work with
commutators and then the presentation would be called a \emph{power\texttt{\symbol{45}}commutator} presentation. Both types of presentation are abbreviated as \emph{pc presentation}. Note that a pc presentation is a rewriting system. 

 Clearly, whenever a group $G$ with pcgs $P$ is given, then we can write down the corresponding pc presentation. On the
other hand, one may just write down a presentation on $n$ abstract generators $g_1, \ldots, g_n$ with relations of the above form and define a group $H$ by this. Then the subgroups $C_i = \langle g_i, \ldots, g_n \rangle$ of $H$ form a subnormal series whose factors are cyclic or trivial. In the case that
all factors are non\texttt{\symbol{45}}trivial, we say that the pc
presentation of $H$ is \emph{confluent}. Note that \textsf{GAP}{\nobreakspace}4 can only work correctly with pc groups defined by a confluent
pc presentation. 

 At the current state of implementations the \textsf{GAP} library contains methods to compute with finite polycyclic groups, while the \textsf{GAP} package \textsf{Polycyclic} by Bettina Eick and Werner Nickel allows also computations with infinite
polycyclic groups which are given by a pc\texttt{\symbol{45}}presentation. 

     Algorithms for pc groups use the methods for polycyclic groups described in
chapter{\nobreakspace}\ref{Polycyclic Groups}.  
\section{\textcolor{Chapter }{The Family Pcgs}}\label{The Family Pcgs}
\logpage{[ 46, 1, 0 ]}
\hyperdef{L}{X78E9E4D778A57A96}{}
{
  Clearly, the generators of a power\texttt{\symbol{45}}conjugate presentation
of a pc group $G$ form a pcgs of the pc group. This pcgs is called the \emph{family pcgs}. 

\subsection{\textcolor{Chapter }{FamilyPcgs}}
\logpage{[ 46, 1, 1 ]}\nobreak
\hyperdef{L}{X79EDB35E82C99304}{}
{\noindent\textcolor{FuncColor}{$\triangleright$\enspace\texttt{FamilyPcgs({\mdseries\slshape grp})\index{FamilyPcgs@\texttt{FamilyPcgs}}
\label{FamilyPcgs}
}\hfill{\scriptsize (attribute)}}\\


 returns, for a pc group \mbox{\texttt{\mdseries\slshape grp}}, a ``natural'' pcgs of some group $G$ which contains \mbox{\texttt{\mdseries\slshape grp}} and is maximal with this property. 

 The pcgs operations described in Chapter{\nobreakspace}\ref{Polycyclic Groups} are particularly efficient with respect to this pcgs. }

 

\subsection{\textcolor{Chapter }{IsFamilyPcgs}}
\logpage{[ 46, 1, 2 ]}\nobreak
\hyperdef{L}{X80893D2A7FFC791B}{}
{\noindent\textcolor{FuncColor}{$\triangleright$\enspace\texttt{IsFamilyPcgs({\mdseries\slshape pcgs})\index{IsFamilyPcgs@\texttt{IsFamilyPcgs}}
\label{IsFamilyPcgs}
}\hfill{\scriptsize (property)}}\\


 specifies whether the pcgs is a \texttt{FamilyPcgs} (\ref{FamilyPcgs}) of a pc group. }

 

\subsection{\textcolor{Chapter }{InducedPcgsWrtFamilyPcgs}}
\logpage{[ 46, 1, 3 ]}\nobreak
\hyperdef{L}{X85C1596A867BE93D}{}
{\noindent\textcolor{FuncColor}{$\triangleright$\enspace\texttt{InducedPcgsWrtFamilyPcgs({\mdseries\slshape grp})\index{InducedPcgsWrtFamilyPcgs@\texttt{InducedPcgsWrtFamilyPcgs}}
\label{InducedPcgsWrtFamilyPcgs}
}\hfill{\scriptsize (attribute)}}\\


 returns the pcgs which induced with respect to a family pcgs (see \texttt{IsParentPcgsFamilyPcgs} (\ref{IsParentPcgsFamilyPcgs}) for further details). }

 

\subsection{\textcolor{Chapter }{IsParentPcgsFamilyPcgs}}
\logpage{[ 46, 1, 4 ]}\nobreak
\hyperdef{L}{X8333ACCB7F530406}{}
{\noindent\textcolor{FuncColor}{$\triangleright$\enspace\texttt{IsParentPcgsFamilyPcgs({\mdseries\slshape pcgs})\index{IsParentPcgsFamilyPcgs@\texttt{IsParentPcgsFamilyPcgs}}
\label{IsParentPcgsFamilyPcgs}
}\hfill{\scriptsize (property)}}\\


 This property indicates that the pcgs \mbox{\texttt{\mdseries\slshape pcgs}} is induced with respect to a family pcgs. 

 This property is needed to distinguish between different independent
polycyclic generating sequences which a pc group may have, since the
elementary operations for a non\texttt{\symbol{45}}family pcgs may not be as
efficient as the elementary operations for the family pcgs. 

 This can have a significant influence on the performance of algorithms for
polycyclic groups. Many algorithms require a pcgs that corresponds to an
elementary abelian series (see{\nobreakspace}\texttt{PcgsElementaryAbelianSeries} (\ref{PcgsElementaryAbelianSeries:for a group})) or even a special pcgs (see{\nobreakspace}\ref{Special Pcgs}). If the family pcgs has the required properties, it will be used for these
purposes, if not \textsf{GAP} has to work with respect to a new pcgs which is \emph{not} the family pcgs and thus takes longer for elementary calculations like \texttt{ExponentsOfPcElement} (\ref{ExponentsOfPcElement}). 

 Therefore, if the family pcgs chosen for arithmetic is not of importance it
might be worth to \emph{change} to another, nicer, pcgs to speed up calculations. This can be achieved, for
example, by using the \texttt{Range} (\ref{Range:of a general mapping}) value of the isomorphism obtained by \texttt{IsomorphismSpecialPcGroup} (\ref{IsomorphismSpecialPcGroup}). }

 }

  
\section{\textcolor{Chapter }{Elements of Pc Groups}}\label{Elements of Pc Groups}
\logpage{[ 46, 2, 0 ]}
\hyperdef{L}{X842526BE7FEFE8BD}{}
{
  
\subsection{\textcolor{Chapter }{Comparison of elements of pc groups}}\logpage{[ 46, 2, 1 ]}
\hyperdef{L}{X869DCE7D86E32337}{}
{
\noindent\textcolor{FuncColor}{$\triangleright$\enspace\texttt{\texttt{\symbol{92}}=({\mdseries\slshape pcword1, pcword2})\index{\texttt{\symbol{92}}=@\texttt{\texttt{\symbol{92}}=}!for pcwords}
\label{bSlash=:for pcwords}
}\hfill{\scriptsize (method)}}\\
\noindent\textcolor{FuncColor}{$\triangleright$\enspace\texttt{\texttt{\symbol{92}}{\textless}({\mdseries\slshape pcword1, pcword2})\index{\texttt{\symbol{92}}{\textless}@\texttt{\texttt{\symbol{92}}{\textless}}!for pcwords}
\label{bSlash<:for pcwords}
}\hfill{\scriptsize (method)}}\\


 \index{equality!for pcwords} \index{smaller!for pcwords} The elements of a pc group $G$ are always represented as words in normal form with respect to the family pcgs
of $G$. Thus it is straightforward to compare elements of a pc group, since this
boils down to a mere comparison of exponent vectors with respect to the family
pcgs. In particular, the word problem is efficiently solvable in pc groups. }

 
\subsection{\textcolor{Chapter }{Arithmetic operations for elements of pc groups}}\logpage{[ 46, 2, 2 ]}
\hyperdef{L}{X7D1B700882FC6C78}{}
{
\noindent\textcolor{FuncColor}{$\triangleright$\enspace\texttt{\texttt{\symbol{92}}*({\mdseries\slshape pcword1, pcword2})\index{\texttt{\symbol{92}}*@\texttt{\texttt{\symbol{92}}*}!for pcwords}
\label{bSlash*:for pcwords}
}\hfill{\scriptsize (method)}}\\
\noindent\textcolor{FuncColor}{$\triangleright$\enspace\texttt{Inverse({\mdseries\slshape pcword})\index{Inverse@\texttt{Inverse}!for a pcword}
\label{Inverse:for a pcword}
}\hfill{\scriptsize (attribute)}}\\


 However, multiplication and inversion of elements in pc groups is not as
straightforward as in arbitrary finitely presented groups where a simple
concatenation or reversion of the corresponding words is sufficient (but one
cannot solve the word problem). 

 To multiply two elements in a pc group, we first concatenate the corresponding
words and then use an algorithm called \emph{collection} to transform the new word into a word in normal form. 

 
\begin{Verbatim}[commandchars=!@|,fontsize=\small,frame=single,label=Example]
  !gapprompt@gap>| !gapinput@g := FamilyPcgs( SmallGroup( 24, 12 ) );|
  Pcgs([ f1, f2, f3, f4 ])
  !gapprompt@gap>| !gapinput@g[4] * g[1];|
  f1*f3
  !gapprompt@gap>| !gapinput@(g[2] * g[3])^-1;|
  f2^2*f3*f4
\end{Verbatim}
 }

 }

  
\section{\textcolor{Chapter }{Pc Groups versus Fp Groups}}\label{Pc Groups versus Fp Groups}
\logpage{[ 46, 3, 0 ]}
\hyperdef{L}{X87B866C386B386E4}{}
{
  In theory pc groups are finitely presented groups. In practice the arithmetic
in pc groups is different from the arithmetic in fp groups. Thus for technical
reasons the pc groups in \textsf{GAP} do not form a subcategory of the fp groups and hence the methods for fp groups
cannot be applied to pc groups in general. 

\subsection{\textcolor{Chapter }{IsPcGroup}}
\logpage{[ 46, 3, 1 ]}\nobreak
\hyperdef{L}{X7D1F506D7830B1D9}{}
{\noindent\textcolor{FuncColor}{$\triangleright$\enspace\texttt{IsPcGroup({\mdseries\slshape G})\index{IsPcGroup@\texttt{IsPcGroup}}
\label{IsPcGroup}
}\hfill{\scriptsize (Category)}}\\


 tests whether \mbox{\texttt{\mdseries\slshape G}} is a pc group. 
\begin{Verbatim}[commandchars=!@|,fontsize=\small,frame=single,label=Example]
  !gapprompt@gap>| !gapinput@G := SmallGroup( 24, 12 );|
  <pc group of size 24 with 4 generators>
  !gapprompt@gap>| !gapinput@IsPcGroup( G );|
  true
  !gapprompt@gap>| !gapinput@IsFpGroup( G );|
  false
\end{Verbatim}
 }

 

\subsection{\textcolor{Chapter }{IsomorphismFpGroupByPcgs}}
\logpage{[ 46, 3, 2 ]}\nobreak
\hyperdef{L}{X7D2735A18111FE39}{}
{\noindent\textcolor{FuncColor}{$\triangleright$\enspace\texttt{IsomorphismFpGroupByPcgs({\mdseries\slshape pcgs, str})\index{IsomorphismFpGroupByPcgs@\texttt{IsomorphismFpGroupByPcgs}}
\label{IsomorphismFpGroupByPcgs}
}\hfill{\scriptsize (function)}}\\


 It is possible to convert a pc group to a fp group in \textsf{GAP}. The function \texttt{IsomorphismFpGroupByPcgs} computes the power\texttt{\symbol{45}}commutator presentation defined by \mbox{\texttt{\mdseries\slshape pcgs}}. The string \mbox{\texttt{\mdseries\slshape str}} can be used to give a name to the generators of the fp group. 

 
\begin{Verbatim}[commandchars=!@|,fontsize=\small,frame=single,label=Example]
  !gapprompt@gap>| !gapinput@p := FamilyPcgs( SmallGroup( 24, 12 ) );|
  Pcgs([ f1, f2, f3, f4 ])
  !gapprompt@gap>| !gapinput@iso := IsomorphismFpGroupByPcgs( p, "g" );|
  [ f1, f2, f3, f4 ] -> [ g1, g2, g3, g4 ]
  !gapprompt@gap>| !gapinput@F := Image( iso );|
  <fp group of size 24 on the generators [ g1, g2, g3, g4 ]>
  !gapprompt@gap>| !gapinput@RelatorsOfFpGroup( F );|
  [ g1^2, g2^-1*g1^-1*g2*g1*g2^-1, g3^-1*g1^-1*g3*g1*g4^-1*g3^-1,
    g4^-1*g1^-1*g4*g1*g4^-1*g3^-1, g2^3, g3^-1*g2^-1*g3*g2*g4^-1*g3^-1,
    g4^-1*g2^-1*g4*g2*g3^-1, g3^2, g4^-1*g3^-1*g4*g3, g4^2 ]
\end{Verbatim}
 }

 }

  
\section{\textcolor{Chapter }{Constructing Pc Groups}}\label{Constructing Pc Groups}
\logpage{[ 46, 4, 0 ]}
\hyperdef{L}{X8581887880556E0C}{}
{
  If necessary, you can supply \textsf{GAP} with a pc presentation by hand. (Although this is the most tedious way to
input a pc group.) Note that the pc presentation has to be confluent in order
to work with the pc group in \textsf{GAP}. 

 (If you have already a suitable pcgs in another representation, use \texttt{PcGroupWithPcgs} (\ref{PcGroupWithPcgs}), see below.) 

 One way is to define a finitely presented group with a pc presentation in \textsf{GAP} and then convert this presentation into a pc group, see \texttt{PcGroupFpGroup} (\ref{PcGroupFpGroup}). Note that this does not work for arbitrary presentations of polycyclic
groups, see Chapter \ref{Quotient Methods} for further information. 

 Another way is to create and manipulate a collector of a pc group by hand and
to use it to define a pc group. This is the most technical way and has little
error checking and thus is intended mostly for experts who want to create a pc
presentation in a particular way. \textsf{GAP} provides different collectors for different collecting strategies; at the
moment, there are two collectors to choose from: the single collector for
finite pc groups (see \texttt{SingleCollector} (\ref{SingleCollector})) and the combinatorial collector for finite $p$\texttt{\symbol{45}}groups. See \cite{Sims94} for further information on collecting strategies. 

 A collector is initialized with an underlying free group and the relative
orders of the pc series. Then one adds the right hand sides of the power and
the commutator or conjugate relations one by one. Note that omitted relators
are assumed to be trivial. 

 For performance reasons it is beneficial to enforce a ``syllable'' representation in the free group (see{\nobreakspace}\ref{Representations for Associative Words}). 

 Note that in the end, the collector has to be converted to a group, see \texttt{GroupByRws} (\ref{GroupByRws}). 

 With these methods a pc group with arbitrary defining pcgs can be constructed.
However, for almost all applications within \textsf{GAP} we need to have a pc group whose defining pcgs is a prime order pcgs, see \texttt{IsomorphismRefinedPcGroup} (\ref{IsomorphismRefinedPcGroup}) and \texttt{RefinedPcGroup} (\ref{RefinedPcGroup}). 

\subsection{\textcolor{Chapter }{PcGroupFpGroup}}
\logpage{[ 46, 4, 1 ]}\nobreak
\hyperdef{L}{X84C10D1F7CB5274F}{}
{\noindent\textcolor{FuncColor}{$\triangleright$\enspace\texttt{PcGroupFpGroup({\mdseries\slshape G})\index{PcGroupFpGroup@\texttt{PcGroupFpGroup}}
\label{PcGroupFpGroup}
}\hfill{\scriptsize (function)}}\\


 creates a pc group \mbox{\texttt{\mdseries\slshape P}} from an fp group (see Chapter \ref{Finitely Presented Groups}) \mbox{\texttt{\mdseries\slshape G}} whose presentation is polycyclic. The resulting group \mbox{\texttt{\mdseries\slshape P}} has generators corresponding to the generators of \mbox{\texttt{\mdseries\slshape G}}. They are printed in the same way as generators of \mbox{\texttt{\mdseries\slshape G}}, but they lie in a different family. If the pc presentation of \mbox{\texttt{\mdseries\slshape G}} is not confluent, an error message occurs. 

 
\begin{Verbatim}[commandchars=!@|,fontsize=\small,frame=single,label=Example]
  !gapprompt@gap>| !gapinput@F := FreeGroup(IsSyllableWordsFamily,"a","b","c","d");;|
  !gapprompt@gap>| !gapinput@a := F.1;; b := F.2;; c := F.3;; d := F.4;;|
  !gapprompt@gap>| !gapinput@rels := [a^2, b^3, c^2, d^2, Comm(b,a)/b, Comm(c,a)/d, Comm(d,a),|
  !gapprompt@>| !gapinput@            Comm(c,b)/(c*d), Comm(d,b)/c, Comm(d,c)];|
  [ a^2, b^3, c^2, d^2, b^-1*a^-1*b*a*b^-1, c^-1*a^-1*c*a*d^-1,
    d^-1*a^-1*d*a, c^-1*b^-1*c*b*d^-1*c^-1, d^-1*b^-1*d*b*c^-1,
    d^-1*c^-1*d*c ]
  !gapprompt@gap>| !gapinput@G := F / rels;|
  <fp group on the generators [ a, b, c, d ]>
  !gapprompt@gap>| !gapinput@H := PcGroupFpGroup( G );|
  <pc group of size 24 with 4 generators>
\end{Verbatim}
 }

 

\subsection{\textcolor{Chapter }{SingleCollector}}
\logpage{[ 46, 4, 2 ]}\nobreak
\hyperdef{L}{X7E958DB281E070FD}{}
{\noindent\textcolor{FuncColor}{$\triangleright$\enspace\texttt{SingleCollector({\mdseries\slshape fgrp, relorders})\index{SingleCollector@\texttt{SingleCollector}}
\label{SingleCollector}
}\hfill{\scriptsize (operation)}}\\
\noindent\textcolor{FuncColor}{$\triangleright$\enspace\texttt{CombinatorialCollector({\mdseries\slshape fgrp, relorders})\index{CombinatorialCollector@\texttt{CombinatorialCollector}}
\label{CombinatorialCollector}
}\hfill{\scriptsize (operation)}}\\


 initializes a single collector or a combinatorial collector, where \mbox{\texttt{\mdseries\slshape fgrp}} must be a free group and \mbox{\texttt{\mdseries\slshape relorders}} must be a list of the relative orders of the pc series. 

 A combinatorial collector can only be set up for a finite $p$\texttt{\symbol{45}}group. Here, the relative orders \mbox{\texttt{\mdseries\slshape relorders}} must all be equal and a prime. }

 

\subsection{\textcolor{Chapter }{SetConjugate}}
\logpage{[ 46, 4, 3 ]}\nobreak
\hyperdef{L}{X86A08D887E049347}{}
{\noindent\textcolor{FuncColor}{$\triangleright$\enspace\texttt{SetConjugate({\mdseries\slshape coll, j, i, w})\index{SetConjugate@\texttt{SetConjugate}}
\label{SetConjugate}
}\hfill{\scriptsize (operation)}}\\


 Let $f_1, \ldots, f_n$ be the generators of the underlying free group of the collector \mbox{\texttt{\mdseries\slshape coll}}. 

 For \mbox{\texttt{\mdseries\slshape i}} $<$ \mbox{\texttt{\mdseries\slshape j}}, \texttt{SetConjugate} sets the conjugate $f_j^{{f_i}}$ to equal \mbox{\texttt{\mdseries\slshape w}}, which is assumed to be a canonical word in $f_{{i+1}}, \ldots, f_n$. No check of the arguments is performed. }

 

\subsection{\textcolor{Chapter }{SetCommutator}}
\logpage{[ 46, 4, 4 ]}\nobreak
\hyperdef{L}{X7B25997C7DF92B6D}{}
{\noindent\textcolor{FuncColor}{$\triangleright$\enspace\texttt{SetCommutator({\mdseries\slshape coll, j, i, w})\index{SetCommutator@\texttt{SetCommutator}}
\label{SetCommutator}
}\hfill{\scriptsize (operation)}}\\


 Let $f_1, \ldots, f_n$ be the generators of the underlying free group of the collector \mbox{\texttt{\mdseries\slshape coll}}. 

 For \mbox{\texttt{\mdseries\slshape i}} $<$ \mbox{\texttt{\mdseries\slshape j}}, \texttt{SetCommutator} sets the commutator of $f_j$ and $f_i$ to equal \mbox{\texttt{\mdseries\slshape w}}, which is assumed to be a canonical word in $f_{{i+1}}, \ldots, f_n$. No check of the arguments is performed. }

 

\subsection{\textcolor{Chapter }{SetPower}}
\logpage{[ 46, 4, 5 ]}\nobreak
\hyperdef{L}{X7BC319BA8698420C}{}
{\noindent\textcolor{FuncColor}{$\triangleright$\enspace\texttt{SetPower({\mdseries\slshape coll, i, w})\index{SetPower@\texttt{SetPower}}
\label{SetPower}
}\hfill{\scriptsize (operation)}}\\


 Let $f_1, \ldots, f_n$ be the generators of the underlying free group of the collector \mbox{\texttt{\mdseries\slshape coll}}, and let $r_i$ be the corresponding relative orders. 

 \texttt{SetPower} sets the power $f_i^{{r_i}}$ to equal \mbox{\texttt{\mdseries\slshape w}}, which is assumed to be a canonical word in $f_{{i+1}}, \ldots, f_n$. No check of the arguments is performed. }

 

\subsection{\textcolor{Chapter }{GroupByRws}}
\logpage{[ 46, 4, 6 ]}\nobreak
\hyperdef{L}{X84F0521486672C3C}{}
{\noindent\textcolor{FuncColor}{$\triangleright$\enspace\texttt{GroupByRws({\mdseries\slshape coll})\index{GroupByRws@\texttt{GroupByRws}}
\label{GroupByRws}
}\hfill{\scriptsize (operation)}}\\
\noindent\textcolor{FuncColor}{$\triangleright$\enspace\texttt{GroupByRwsNC({\mdseries\slshape coll})\index{GroupByRwsNC@\texttt{GroupByRwsNC}}
\label{GroupByRwsNC}
}\hfill{\scriptsize (operation)}}\\


 creates a group from a rewriting system. In the first version it is checked
whether the rewriting system is confluent, in the second version this is
assumed to be true. }

 

\subsection{\textcolor{Chapter }{IsConfluent (for pc groups)}}
\logpage{[ 46, 4, 7 ]}\nobreak
\hyperdef{L}{X7DF4835F79667099}{}
{\noindent\textcolor{FuncColor}{$\triangleright$\enspace\texttt{IsConfluent({\mdseries\slshape G})\index{IsConfluent@\texttt{IsConfluent}!for pc groups}
\label{IsConfluent:for pc groups}
}\hfill{\scriptsize (property)}}\\


 checks whether the pc group \mbox{\texttt{\mdseries\slshape G}} has been built from a collector with a confluent
power\texttt{\symbol{45}}commutator presentation. 

 
\begin{Verbatim}[commandchars=!@|,fontsize=\small,frame=single,label=Example]
  !gapprompt@gap>| !gapinput@F := FreeGroup(IsSyllableWordsFamily, 2 );;|
  !gapprompt@gap>| !gapinput@coll1 := SingleCollector( F, [2,3] );|
  <<single collector, 8 Bits>>
  !gapprompt@gap>| !gapinput@SetConjugate( coll1, 2, 1, F.2 );|
  !gapprompt@gap>| !gapinput@SetPower( coll1, 1, F.2 );|
  !gapprompt@gap>| !gapinput@G1 := GroupByRws( coll1 );|
  <pc group of size 6 with 2 generators>
  !gapprompt@gap>| !gapinput@IsConfluent(G1);|
  true
  !gapprompt@gap>| !gapinput@IsAbelian(G1);|
  true
  !gapprompt@gap>| !gapinput@coll2 := SingleCollector( F, [2,3] );|
  <<single collector, 8 Bits>>
  !gapprompt@gap>| !gapinput@SetConjugate( coll2, 2, 1, F.2^2 );|
  !gapprompt@gap>| !gapinput@G2 := GroupByRws( coll2 );|
  <pc group of size 6 with 2 generators>
  !gapprompt@gap>| !gapinput@IsAbelian(G2);|
  false
\end{Verbatim}
 }

 

\subsection{\textcolor{Chapter }{IsomorphismRefinedPcGroup}}
\logpage{[ 46, 4, 8 ]}\nobreak
\hyperdef{L}{X7E6226597DFE5F8F}{}
{\noindent\textcolor{FuncColor}{$\triangleright$\enspace\texttt{IsomorphismRefinedPcGroup({\mdseries\slshape G})\index{IsomorphismRefinedPcGroup@\texttt{IsomorphismRefinedPcGroup}}
\label{IsomorphismRefinedPcGroup}
}\hfill{\scriptsize (attribute)}}\\


 \index{isomorphic!pc group} returns an isomorphism from a pc group \mbox{\texttt{\mdseries\slshape G}} onto an isomorphic pc group whose family pcgs is a prime order pcgs. }

 

\subsection{\textcolor{Chapter }{RefinedPcGroup}}
\logpage{[ 46, 4, 9 ]}\nobreak
\hyperdef{L}{X821560A387762DD1}{}
{\noindent\textcolor{FuncColor}{$\triangleright$\enspace\texttt{RefinedPcGroup({\mdseries\slshape G})\index{RefinedPcGroup@\texttt{RefinedPcGroup}}
\label{RefinedPcGroup}
}\hfill{\scriptsize (attribute)}}\\


 returns the range of the \texttt{IsomorphismRefinedPcGroup} (\ref{IsomorphismRefinedPcGroup}) value of \mbox{\texttt{\mdseries\slshape G}}. }

 }

  
\section{\textcolor{Chapter }{Computing Pc Groups}}\label{Computing Pc Groups}
\logpage{[ 46, 5, 0 ]}
\hyperdef{L}{X83F69FE27B024E24}{}
{
  Another possibility to get a pc group in \textsf{GAP} is to convert a polycyclic group given by some other representation to a pc
group. For finitely presented groups there are various quotient methods
available. For all other types of groups one can use the following functions. 

\subsection{\textcolor{Chapter }{PcGroupWithPcgs}}
\logpage{[ 46, 5, 1 ]}\nobreak
\hyperdef{L}{X81C55D4F825C36D4}{}
{\noindent\textcolor{FuncColor}{$\triangleright$\enspace\texttt{PcGroupWithPcgs({\mdseries\slshape mpcgs})\index{PcGroupWithPcgs@\texttt{PcGroupWithPcgs}}
\label{PcGroupWithPcgs}
}\hfill{\scriptsize (attribute)}}\\


 creates a new pc group \mbox{\texttt{\mdseries\slshape G}} whose family pcgs is isomorphic to the (modulo) pcgs \mbox{\texttt{\mdseries\slshape mpcgs}}. 
\begin{Verbatim}[commandchars=!@|,fontsize=\small,frame=single,label=Example]
  !gapprompt@gap>| !gapinput@G := Group( (1,2,3), (3,4,1) );;|
  !gapprompt@gap>| !gapinput@PcGroupWithPcgs( Pcgs(G) );|
  <pc group of size 12 with 3 generators>
\end{Verbatim}
 

 If a pcgs is only given by a list of pc elements, \texttt{PcgsByPcSequence} (\ref{PcgsByPcSequence}) can be used: 
\begin{Verbatim}[commandchars=!@|,fontsize=\small,frame=single,label=Example]
  !gapprompt@gap>| !gapinput@G:=Group((1,2,3,4),(1,2));;|
  !gapprompt@gap>| !gapinput@p:=PcgsByPcSequence(FamilyObj(One(G)),|
  !gapprompt@>| !gapinput@[ (3,4), (2,4,3), (1,4)(2,3), (1,3)(2,4) ]);|
  Pcgs([ (3,4), (2,4,3), (1,4)(2,3), (1,3)(2,4) ])
  !gapprompt@gap>| !gapinput@PcGroupWithPcgs(p);|
  <pc group of size 24 with 4 generators>
  !gapprompt@gap>| !gapinput@G := SymmetricGroup( 5 );|
  Sym( [ 1 .. 5 ] )
  !gapprompt@gap>| !gapinput@H := Subgroup( G, [(1,2,3,4,5), (3,4,5)] );|
  Group([ (1,2,3,4,5), (3,4,5) ])
  !gapprompt@gap>| !gapinput@modu := ModuloPcgs( G, H );|
  Pcgs([ (4,5) ])
  !gapprompt@gap>| !gapinput@PcGroupWithPcgs(modu);|
  <pc group of size 2 with 1 generator>
\end{Verbatim}
 }

 

\subsection{\textcolor{Chapter }{IsomorphismPcGroup}}
\logpage{[ 46, 5, 2 ]}\nobreak
\hyperdef{L}{X873CEB137BA1CD6E}{}
{\noindent\textcolor{FuncColor}{$\triangleright$\enspace\texttt{IsomorphismPcGroup({\mdseries\slshape G})\index{IsomorphismPcGroup@\texttt{IsomorphismPcGroup}}
\label{IsomorphismPcGroup}
}\hfill{\scriptsize (attribute)}}\\


 \index{isomorphic!pc group} returns an isomorphism from \mbox{\texttt{\mdseries\slshape G}} onto an isomorphic pc group. The series chosen for this pc representation
depends on the method chosen. \mbox{\texttt{\mdseries\slshape G}} must be a finite polycyclic group of any kind, for example a solvable
permutation group. If \mbox{\texttt{\mdseries\slshape G}} is finite but not solvable, \texttt{fail} is returned. 
\begin{Verbatim}[commandchars=!@|,fontsize=\small,frame=single,label=Example]
  !gapprompt@gap>| !gapinput@G := Group( (1,2,3), (3,4,1) );;|
  !gapprompt@gap>| !gapinput@iso := IsomorphismPcGroup( G );|
  Pcgs([ (2,4,3), (1,2)(3,4), (1,3)(2,4) ]) -> [ f1, f2, f3 ]
  !gapprompt@gap>| !gapinput@H := Image( iso );|
  Group([ f1, f2, f3 ])
\end{Verbatim}
 }

 

\subsection{\textcolor{Chapter }{IsomorphismSpecialPcGroup}}
\logpage{[ 46, 5, 3 ]}\nobreak
\hyperdef{L}{X82BE14A986FA6882}{}
{\noindent\textcolor{FuncColor}{$\triangleright$\enspace\texttt{IsomorphismSpecialPcGroup({\mdseries\slshape G})\index{IsomorphismSpecialPcGroup@\texttt{IsomorphismSpecialPcGroup}}
\label{IsomorphismSpecialPcGroup}
}\hfill{\scriptsize (attribute)}}\\


 returns an isomorphism from \mbox{\texttt{\mdseries\slshape G}} onto an isomorphic pc group whose family pcgs is a special pcgs. (This can be
beneficial to the runtime of calculations.) \mbox{\texttt{\mdseries\slshape G}} must be a finite polycyclic group of any kind, for example a solvable
permutation group. If \mbox{\texttt{\mdseries\slshape G}} is finite but not solvable, \texttt{fail} is returned. }

 }

  
\section{\textcolor{Chapter }{Saving a Pc Group}}\label{Saving a Pc Group}
\logpage{[ 46, 6, 0 ]}
\hyperdef{L}{X85696AB9791DF047}{}
{
  As printing a polycyclic group does not display the presentation, one cannot
simply print a pc group to a file to save it. For this purpose we need the
following function. 

\subsection{\textcolor{Chapter }{GapInputPcGroup}}
\logpage{[ 46, 6, 1 ]}\nobreak
\hyperdef{L}{X8593253380D84508}{}
{\noindent\textcolor{FuncColor}{$\triangleright$\enspace\texttt{GapInputPcGroup({\mdseries\slshape grp, string})\index{GapInputPcGroup@\texttt{GapInputPcGroup}}
\label{GapInputPcGroup}
}\hfill{\scriptsize (function)}}\\


 
\begin{Verbatim}[commandchars=!@|,fontsize=\small,frame=single,label=Example]
  !gapprompt@gap>| !gapinput@G := SmallGroup( 24, 12 );|
  <pc group of size 24 with 4 generators>
  !gapprompt@gap>| !gapinput@PrintTo( "save", GapInputPcGroup( G, "H" ) );|
  !gapprompt@gap>| !gapinput@Read( "save" );|
  #I A group of order 24 has been defined.
  #I It is called H
  !gapprompt@gap>| !gapinput@H = G;|
  false
  !gapprompt@gap>| !gapinput@IdSmallGroup( H ) = IdSmallGroup( G );|
  true
  !gapprompt@gap>| !gapinput@RemoveFile( "save" );;|
\end{Verbatim}
 }

 }

  
\section{\textcolor{Chapter }{Operations for Pc Groups}}\label{Operations for Pc Groups}
\logpage{[ 46, 7, 0 ]}
\hyperdef{L}{X8391EE8D782D0C9E}{}
{
  All the operations described in Chapters \ref{Groups} and \ref{Polycyclic Groups} apply to a pc group. Nearly all methods for pc groups are methods for groups
with pcgs as described in Chapter \ref{Polycyclic Groups}. The only method with is special for pc groups is a method to compute
intersections of subgroups, since here a pcgs of a parent group is needed and
this can only by guaranteed within pc groups. Section \ref{2-Cohomology} describes operations and methods for arbitrary finite groups. }

  
\section{\textcolor{Chapter }{$2$\texttt{\symbol{45}}Cohomology and Extensions}}\label{2-Cohomology and Extensions}
\logpage{[ 46, 8, 0 ]}
\hyperdef{L}{X877AAB887D4507E3}{}
{
  One of the most interesting applications of pc groups is the possibility to
compute with extensions of these groups by elementary abelian groups; that is, $H$ is an extension of $G$ by $M$, if there exists a normal subgroup $N$ in $H$ which is isomorphic to $M$ such that $H/N$ is isomorphic to $G$. 

 Pc groups are particularly suited for such applications, since the $2$\texttt{\symbol{45}}cohomology can be computed efficiently for such groups
and, moreover, extensions of pc groups by elementary abelian groups can be
represented as pc groups again. 

 To define the elementary abelian group $M$ together with an action of $G$ on $M$ we consider $M$ as a MeatAxe module for $G$ over a finite field (section{\nobreakspace}\texttt{IrreducibleModules} (\ref{IrreducibleModules}) describes functions that can be used to obtain certain modules). For further
information on meataxe modules see Chapter \ref{The MeatAxe}. Note that the matrices defining the module must correspond to the pcgs of
the group \mbox{\texttt{\mdseries\slshape G}}. 

 There exists an action of the subgroup of \emph{compatible pairs} in $Aut(G) \times Aut(M)$ which acts on the second cohomology group, see \texttt{CompatiblePairs} (\ref{CompatiblePairs}). $2$\texttt{\symbol{45}}cocycles which lie in the same orbit under this action
define isomorphic extensions of $G$. However, there may be isomorphic extensions of $G$ corresponding to cocycles in different orbits. 

 See also the \textsf{GAP} package \textsf{GrpConst} by Hans Ulrich Besche and Bettina Eick that contains methods to construct up
to isomorphism the groups of a given order. 

 Finally we note that for the computation of split extensions it is not
necessary that \mbox{\texttt{\mdseries\slshape M}} must correspond to an elementary abelian group. Here it is possible to
construct split extensions of arbitrary pc groups, see \texttt{SplitExtension} (\ref{SplitExtension}). 

\subsection{\textcolor{Chapter }{TwoCoboundaries}}
\logpage{[ 46, 8, 1 ]}\nobreak
\hyperdef{L}{X78E6E11E8285E288}{}
{\noindent\textcolor{FuncColor}{$\triangleright$\enspace\texttt{TwoCoboundaries({\mdseries\slshape G, M})\index{TwoCoboundaries@\texttt{TwoCoboundaries}}
\label{TwoCoboundaries}
}\hfill{\scriptsize (operation)}}\\


 returns the group of $2$\texttt{\symbol{45}}coboundaries of a pc group \mbox{\texttt{\mdseries\slshape G}} by the \mbox{\texttt{\mdseries\slshape G}}\texttt{\symbol{45}}module \mbox{\texttt{\mdseries\slshape M}}. The generators of \mbox{\texttt{\mdseries\slshape M}} must correspond to the \texttt{Pcgs} (\ref{Pcgs}) value of \mbox{\texttt{\mdseries\slshape G}}. The group of coboundaries is given as vector space over the field underlying \mbox{\texttt{\mdseries\slshape M}}. }

 

\subsection{\textcolor{Chapter }{TwoCocycles}}
\logpage{[ 46, 8, 2 ]}\nobreak
\hyperdef{L}{X784FCA207B8694A6}{}
{\noindent\textcolor{FuncColor}{$\triangleright$\enspace\texttt{TwoCocycles({\mdseries\slshape G, M})\index{TwoCocycles@\texttt{TwoCocycles}}
\label{TwoCocycles}
}\hfill{\scriptsize (operation)}}\\


 returns the $2$\texttt{\symbol{45}}cocycles of a pc group \mbox{\texttt{\mdseries\slshape G}} by the \mbox{\texttt{\mdseries\slshape G}}\texttt{\symbol{45}}module \mbox{\texttt{\mdseries\slshape M}}. The generators of \mbox{\texttt{\mdseries\slshape M}} must correspond to the \texttt{Pcgs} (\ref{Pcgs}) value of \mbox{\texttt{\mdseries\slshape G}}. The operation returns a list of vectors over the field underlying \mbox{\texttt{\mdseries\slshape M}} and the additive group generated by these vectors is the group of $2$\texttt{\symbol{45}}cocyles. }

 

\subsection{\textcolor{Chapter }{TwoCohomology}}
\logpage{[ 46, 8, 3 ]}\nobreak
\hyperdef{L}{X838065F97F60468F}{}
{\noindent\textcolor{FuncColor}{$\triangleright$\enspace\texttt{TwoCohomology({\mdseries\slshape G, M})\index{TwoCohomology@\texttt{TwoCohomology}}
\label{TwoCohomology}
}\hfill{\scriptsize (operation)}}\\


 This operation computes the second cohomology group for the special case of a
Pc Group. It returns a record defining the second cohomology group as factor
space of the space of cocycles by the space of coboundaries. \mbox{\texttt{\mdseries\slshape G}} must be a pc group and the generators of \mbox{\texttt{\mdseries\slshape M}} must correspond to the pcgs of \mbox{\texttt{\mdseries\slshape G}}. 
\begin{Verbatim}[commandchars=!@|,fontsize=\small,frame=single,label=Example]
  !gapprompt@gap>| !gapinput@G := SmallGroup( 4, 2 );|
  <pc group of size 4 with 2 generators>
  !gapprompt@gap>| !gapinput@mats := List( Pcgs( G ), x -> IdentityMat( 1, GF(2) ) );|
  [ [ <a GF2 vector of length 1> ], [ <a GF2 vector of length 1> ] ]
  !gapprompt@gap>| !gapinput@M := GModuleByMats( mats, GF(2) );|
  rec( IsOverFiniteField := true, dimension := 1, field := GF(2),
    generators := [ <an immutable 1x1 matrix over GF2>,
        <an immutable 1x1 matrix over GF2> ], isMTXModule := true )
  !gapprompt@gap>| !gapinput@TwoCoboundaries( G, M );|
  [  ]
  !gapprompt@gap>| !gapinput@TwoCocycles( G, M );|
  [ [ Z(2)^0, 0*Z(2), 0*Z(2) ], [ 0*Z(2), Z(2)^0, 0*Z(2) ],
    [ 0*Z(2), 0*Z(2), Z(2)^0 ] ]
  !gapprompt@gap>| !gapinput@cc := TwoCohomology( G, M );;|
  !gapprompt@gap>| !gapinput@cc.cohom;|
  <linear mapping by matrix, <vector space of dimension 3 over GF(
  2)> -> ( GF(2)^3 )>
\end{Verbatim}
 }

 

\subsection{\textcolor{Chapter }{Extensions}}
\logpage{[ 46, 8, 4 ]}\nobreak
\hyperdef{L}{X8236AD927A5A0E5A}{}
{\noindent\textcolor{FuncColor}{$\triangleright$\enspace\texttt{Extensions({\mdseries\slshape G, M})\index{Extensions@\texttt{Extensions}}
\label{Extensions}
}\hfill{\scriptsize (operation)}}\\


 returns all extensions of \mbox{\texttt{\mdseries\slshape G}} by the \mbox{\texttt{\mdseries\slshape G}}\texttt{\symbol{45}}module \mbox{\texttt{\mdseries\slshape M}} up to equivalence as pc groups. }

 

\subsection{\textcolor{Chapter }{Extension}}
\logpage{[ 46, 8, 5 ]}\nobreak
\hyperdef{L}{X7B3BE908867CE4F9}{}
{\noindent\textcolor{FuncColor}{$\triangleright$\enspace\texttt{Extension({\mdseries\slshape G, M, c})\index{Extension@\texttt{Extension}}
\label{Extension}
}\hfill{\scriptsize (operation)}}\\
\noindent\textcolor{FuncColor}{$\triangleright$\enspace\texttt{ExtensionNC({\mdseries\slshape G, M, c})\index{ExtensionNC@\texttt{ExtensionNC}}
\label{ExtensionNC}
}\hfill{\scriptsize (operation)}}\\


 returns the extension of \mbox{\texttt{\mdseries\slshape G}} by the \mbox{\texttt{\mdseries\slshape G}}\texttt{\symbol{45}}module \mbox{\texttt{\mdseries\slshape M}} via the cocycle \mbox{\texttt{\mdseries\slshape c}} as pc groups. The \texttt{NC} version does not check the resulting group for consistence. }

 

\subsection{\textcolor{Chapter }{SplitExtension}}
\logpage{[ 46, 8, 6 ]}\nobreak
\hyperdef{L}{X83DCB5AB7B6EE785}{}
{\noindent\textcolor{FuncColor}{$\triangleright$\enspace\texttt{SplitExtension({\mdseries\slshape G, M})\index{SplitExtension@\texttt{SplitExtension}}
\label{SplitExtension}
}\hfill{\scriptsize (operation)}}\\


 returns the split extension of \mbox{\texttt{\mdseries\slshape G}} by the \mbox{\texttt{\mdseries\slshape G}}\texttt{\symbol{45}}module \mbox{\texttt{\mdseries\slshape M}}. See also \texttt{SplitExtension} (\ref{SplitExtension:with specified homomorphism}) for its 3\texttt{\symbol{45}}argument version. }

 

\subsection{\textcolor{Chapter }{ModuleOfExtension}}
\logpage{[ 46, 8, 7 ]}\nobreak
\hyperdef{L}{X7EAC6B8B7ABEEB86}{}
{\noindent\textcolor{FuncColor}{$\triangleright$\enspace\texttt{ModuleOfExtension({\mdseries\slshape E})\index{ModuleOfExtension@\texttt{ModuleOfExtension}}
\label{ModuleOfExtension}
}\hfill{\scriptsize (attribute)}}\\


 returns the module of an extension \mbox{\texttt{\mdseries\slshape E}} of \mbox{\texttt{\mdseries\slshape G}} by \mbox{\texttt{\mdseries\slshape M}}. This is the normal subgroup of \mbox{\texttt{\mdseries\slshape E}} which corresponds to \mbox{\texttt{\mdseries\slshape M}}. 

 
\begin{Verbatim}[commandchars=!@|,fontsize=\small,frame=single,label=Example]
  !gapprompt@gap>| !gapinput@G := SmallGroup( 4, 2 );;|
  !gapprompt@gap>| !gapinput@mats := List( Pcgs( G ), x -> IdentityMat( 1, GF(2) ) );;|
  !gapprompt@gap>| !gapinput@M := GModuleByMats( mats, GF(2) );;|
  !gapprompt@gap>| !gapinput@co := TwoCocycles( G, M );;|
  !gapprompt@gap>| !gapinput@Extension( G, M, co[2] );|
  <pc group of size 8 with 3 generators>
  !gapprompt@gap>| !gapinput@SplitExtension( G, M );|
  <pc group of size 8 with 3 generators>
  !gapprompt@gap>| !gapinput@Extensions( G, M );|
  [ <pc group of size 8 with 3 generators>,
    <pc group of size 8 with 3 generators>,
    <pc group of size 8 with 3 generators>,
    <pc group of size 8 with 3 generators>,
    <pc group of size 8 with 3 generators>,
    <pc group of size 8 with 3 generators>,
    <pc group of size 8 with 3 generators>,
    <pc group of size 8 with 3 generators> ]
  !gapprompt@gap>| !gapinput@List(last, IdSmallGroup);|
  [ [ 8, 5 ], [ 8, 2 ], [ 8, 3 ], [ 8, 3 ], [ 8, 2 ], [ 8, 2 ],
    [ 8, 3 ], [ 8, 4 ] ]
\end{Verbatim}
 

 Note that the extensions returned by \texttt{Extensions} (\ref{Extensions}) are computed up to equivalence, but not up to isomorphism. }

 

\subsection{\textcolor{Chapter }{CompatiblePairs}}
\logpage{[ 46, 8, 8 ]}\nobreak
\hyperdef{L}{X824F2B2E7C11ABAF}{}
{\noindent\textcolor{FuncColor}{$\triangleright$\enspace\texttt{CompatiblePairs({\mdseries\slshape [A, ]G, M[, D]})\index{CompatiblePairs@\texttt{CompatiblePairs}}
\label{CompatiblePairs}
}\hfill{\scriptsize (function)}}\\


 returns the group of compatible pairs of the group \mbox{\texttt{\mdseries\slshape G}} with the \mbox{\texttt{\mdseries\slshape G}}\texttt{\symbol{45}}module \mbox{\texttt{\mdseries\slshape M}} as subgroup of the direct product Aut(\mbox{\texttt{\mdseries\slshape G}}) $\times$ Aut(\mbox{\texttt{\mdseries\slshape M}}). Here Aut(\mbox{\texttt{\mdseries\slshape M}}) is considered as subgroup of a general linear group. If \mbox{\texttt{\mdseries\slshape G}} is solvable, then the generators of \mbox{\texttt{\mdseries\slshape M}} must correspond to the \texttt{Pcgs} (\ref{Pcgs}) value of \mbox{\texttt{\mdseries\slshape G}}, otherwise the generators of \mbox{\texttt{\mdseries\slshape M}} must correspond to the \texttt{GeneratorsOfGroup} (\ref{GeneratorsOfGroup}) value of \mbox{\texttt{\mdseries\slshape G}}. 

 The optional argument \mbox{\texttt{\mdseries\slshape D}} should be a subgroup of Aut(\mbox{\texttt{\mdseries\slshape G}}) $\times$ Aut(\mbox{\texttt{\mdseries\slshape M}}). If it is given, then only the compatible pairs in \mbox{\texttt{\mdseries\slshape D}} are computed. If a group \mbox{\texttt{\mdseries\slshape A}} of automorphisms of \mbox{\texttt{\mdseries\slshape G}} is given as optional first argument, it is used in place of the full
automorphism group of \mbox{\texttt{\mdseries\slshape G}}, avoiding the need to compute this automorphism group. }

 

\subsection{\textcolor{Chapter }{ExtensionRepresentatives}}
\logpage{[ 46, 8, 9 ]}\nobreak
\hyperdef{L}{X854FFEF187C4AAB9}{}
{\noindent\textcolor{FuncColor}{$\triangleright$\enspace\texttt{ExtensionRepresentatives({\mdseries\slshape G, M, P})\index{ExtensionRepresentatives@\texttt{ExtensionRepresentatives}}
\label{ExtensionRepresentatives}
}\hfill{\scriptsize (operation)}}\\


 returns all extensions of \mbox{\texttt{\mdseries\slshape G}} by the \mbox{\texttt{\mdseries\slshape G}}\texttt{\symbol{45}}module \mbox{\texttt{\mdseries\slshape M}} up to equivalence under action of \mbox{\texttt{\mdseries\slshape P}} where \mbox{\texttt{\mdseries\slshape P}} has to be a subgroup of the group of compatible pairs of \mbox{\texttt{\mdseries\slshape G}} with \mbox{\texttt{\mdseries\slshape M}}. 
\begin{Verbatim}[commandchars=!@|,fontsize=\small,frame=single,label=Example]
  !gapprompt@gap>| !gapinput@G := SmallGroup( 4, 2 );;|
  !gapprompt@gap>| !gapinput@mats := List( Pcgs( G ), x -> IdentityMat( 1, GF(2) ) );;|
  !gapprompt@gap>| !gapinput@M := GModuleByMats( mats, GF(2) );;|
  !gapprompt@gap>| !gapinput@A := AutomorphismGroup( G );;|
  !gapprompt@gap>| !gapinput@B := GL( 1, 2 );;|
  !gapprompt@gap>| !gapinput@D := DirectProduct( A, B );; Size(D);|
  6
  !gapprompt@gap>| !gapinput@P := CompatiblePairs( G, M, D );|
  <group of size 6 with 2 generators>
  !gapprompt@gap>| !gapinput@ExtensionRepresentatives( G, M, P );|
  [ <pc group of size 8 with 3 generators>,
    <pc group of size 8 with 3 generators>,
    <pc group of size 8 with 3 generators>,
    <pc group of size 8 with 3 generators> ]
  !gapprompt@gap>| !gapinput@Extensions( G, M );|
  [ <pc group of size 8 with 3 generators>,
    <pc group of size 8 with 3 generators>,
    <pc group of size 8 with 3 generators>,
    <pc group of size 8 with 3 generators>,
    <pc group of size 8 with 3 generators>,
    <pc group of size 8 with 3 generators>,
    <pc group of size 8 with 3 generators>,
    <pc group of size 8 with 3 generators> ]
\end{Verbatim}
 }

 

\subsection{\textcolor{Chapter }{SplitExtension (with specified homomorphism)}}
\logpage{[ 46, 8, 10 ]}\nobreak
\hyperdef{L}{X84E2DA897FAAF6D8}{}
{\noindent\textcolor{FuncColor}{$\triangleright$\enspace\texttt{SplitExtension({\mdseries\slshape G, aut, N})\index{SplitExtension@\texttt{SplitExtension}!with specified homomorphism}
\label{SplitExtension:with specified homomorphism}
}\hfill{\scriptsize (operation)}}\\


 returns the split extensions of the pc group \mbox{\texttt{\mdseries\slshape G}} by the pc group \mbox{\texttt{\mdseries\slshape N}}. \mbox{\texttt{\mdseries\slshape aut}} should be a homomorphism from \mbox{\texttt{\mdseries\slshape G}} into Aut(\mbox{\texttt{\mdseries\slshape N}}). 

 In the following example we construct the holomorph of $Q_8$ as split extension of $Q_8$ by $S_4$. 

 
\begin{Verbatim}[commandchars=!@|,fontsize=\small,frame=single,label=Example]
  !gapprompt@gap>| !gapinput@N := SmallGroup( 8, 4 );|
  <pc group of size 8 with 3 generators>
  !gapprompt@gap>| !gapinput@IsAbelian( N );|
  false
  !gapprompt@gap>| !gapinput@A := AutomorphismGroup( N );|
  <group of size 24 with 4 generators>
  !gapprompt@gap>| !gapinput@iso := IsomorphismPcGroup( A );|
  CompositionMapping( Pcgs([ (2,6,5,3), (1,3,5)(2,4,6), (2,5)(3,6),
    (1,4)(3,6) ]) -> [ f1, f2, f3, f4 ], <action isomorphism> )
  !gapprompt@gap>| !gapinput@H := Image( iso );|
  Group([ f1, f2, f3, f4 ])
  !gapprompt@gap>| !gapinput@G := Subgroup( H, Pcgs(H){[1,2]} );|
  Group([ f1, f2 ])
  !gapprompt@gap>| !gapinput@inv := InverseGeneralMapping( iso );|
  [ f1*f2, f2^2*f3, f4, f3 ] ->
  [ Pcgs([ f1, f2, f3 ]) -> [ f1*f2, f2, f3 ],
    Pcgs([ f1, f2, f3 ]) -> [ f2, f1*f2, f3 ],
    Pcgs([ f1, f2, f3 ]) -> [ f1*f3, f2, f3 ],
    Pcgs([ f1, f2, f3 ]) -> [ f1, f2*f3, f3 ] ]
  !gapprompt@gap>| !gapinput@K := SplitExtension( G, inv, N );|
  <pc group of size 192 with 7 generators>
\end{Verbatim}
 }

 }

  
\section{\textcolor{Chapter }{Coding a Pc Presentation}}\label{Coding a Pc Presentation}
\logpage{[ 46, 9, 0 ]}
\hyperdef{L}{X874E4B107BD78F5A}{}
{
  If one wants to store a large number of pc groups, then it can be useful to
store them in a compressed format, since pc presentations can be space
consuming. Here we introduce a method to code and decode pc presentations by
integers. To decode a given code the size of the underlying pc group is needed
as well. For the full definition and the coding and decoding procedures see \cite{BescheEick98}. This method is used with the small groups library (see  (\textbf{smallgrp: The Small Groups Library}). 

\subsection{\textcolor{Chapter }{CodePcgs}}
\logpage{[ 46, 9, 1 ]}\nobreak
\hyperdef{L}{X79948F1D7D4FF8D9}{}
{\noindent\textcolor{FuncColor}{$\triangleright$\enspace\texttt{CodePcgs({\mdseries\slshape pcgs})\index{CodePcgs@\texttt{CodePcgs}}
\label{CodePcgs}
}\hfill{\scriptsize (function)}}\\


 returns the code corresponding to \mbox{\texttt{\mdseries\slshape pcgs}}. 
\begin{Verbatim}[commandchars=!@|,fontsize=\small,frame=single,label=Example]
  !gapprompt@gap>| !gapinput@G := CyclicGroup(512);;|
  !gapprompt@gap>| !gapinput@p := Pcgs( G );;|
  !gapprompt@gap>| !gapinput@CodePcgs( p );|
  162895587718739690298008513020159
\end{Verbatim}
 }

 

\subsection{\textcolor{Chapter }{CodePcGroup}}
\logpage{[ 46, 9, 2 ]}\nobreak
\hyperdef{L}{X8041C2D88721EEA9}{}
{\noindent\textcolor{FuncColor}{$\triangleright$\enspace\texttt{CodePcGroup({\mdseries\slshape G})\index{CodePcGroup@\texttt{CodePcGroup}}
\label{CodePcGroup}
}\hfill{\scriptsize (function)}}\\


 returns the code for a pcgs of \mbox{\texttt{\mdseries\slshape G}}. 
\begin{Verbatim}[commandchars=!@|,fontsize=\small,frame=single,label=Example]
  !gapprompt@gap>| !gapinput@G := DihedralGroup(512);;|
  !gapprompt@gap>| !gapinput@CodePcGroup( G );|
  2940208627577393070560341803949986912431725641726
\end{Verbatim}
 }

 

\subsection{\textcolor{Chapter }{PcGroupCode}}
\logpage{[ 46, 9, 3 ]}\nobreak
\hyperdef{L}{X826BFDA07A707C54}{}
{\noindent\textcolor{FuncColor}{$\triangleright$\enspace\texttt{PcGroupCode({\mdseries\slshape code, size})\index{PcGroupCode@\texttt{PcGroupCode}}
\label{PcGroupCode}
}\hfill{\scriptsize (function)}}\\


 returns a pc group of size \mbox{\texttt{\mdseries\slshape size}} corresponding to \mbox{\texttt{\mdseries\slshape code}}. The argument \mbox{\texttt{\mdseries\slshape code}} must be a valid code for a pcgs, otherwise anything may happen. Valid codes
are usually obtained by one of the functions \texttt{CodePcgs} (\ref{CodePcgs}) or \texttt{CodePcGroup} (\ref{CodePcGroup}). 
\begin{Verbatim}[commandchars=!@|,fontsize=\small,frame=single,label=Example]
  !gapprompt@gap>| !gapinput@G := SmallGroup( 24, 12 );;|
  !gapprompt@gap>| !gapinput@p := Pcgs( G );;|
  !gapprompt@gap>| !gapinput@code := CodePcgs( p );|
  5790338948
  !gapprompt@gap>| !gapinput@H := PcGroupCode( code, 24 );|
  <pc group of size 24 with 4 generators>
  !gapprompt@gap>| !gapinput@map := GroupHomomorphismByImages( G, H, p, FamilyPcgs(H) );|
  Pcgs([ f1, f2, f3, f4 ]) -> Pcgs([ f1, f2, f3, f4 ])
  !gapprompt@gap>| !gapinput@IsBijective(map);|
  true
\end{Verbatim}
 }

 }

  
\section{\textcolor{Chapter }{Random Isomorphism Testing}}\label{Random Isomorphism Testing}
\logpage{[ 46, 10, 0 ]}
\hyperdef{L}{X81D211D8838B875C}{}
{
  The generic isomorphism test for groups may be applied to pc groups as well.
However, this test is often quite time consuming. Here we describe another
method to test isomorphism by a probabilistic approach. 

\subsection{\textcolor{Chapter }{RandomIsomorphismTest}}
\logpage{[ 46, 10, 1 ]}\nobreak
\hyperdef{L}{X84F6F9787CB2CF16}{}
{\noindent\textcolor{FuncColor}{$\triangleright$\enspace\texttt{RandomIsomorphismTest({\mdseries\slshape coderecs, n})\index{RandomIsomorphismTest@\texttt{RandomIsomorphismTest}}
\label{RandomIsomorphismTest}
}\hfill{\scriptsize (function)}}\\


 The first argument is a list \mbox{\texttt{\mdseries\slshape coderecs}} containing records describing groups, and the second argument is a
non\texttt{\symbol{45}}negative integer \mbox{\texttt{\mdseries\slshape n}}. 

 The test returns a sublist of \mbox{\texttt{\mdseries\slshape coderecs}} where isomorphic copies detected by the probabilistic test have been removed. 

 The list \mbox{\texttt{\mdseries\slshape coderecs}} should contain records with two components, \texttt{code} and \texttt{order}, describing a group via \texttt{PcGroupCode( code, order )} (see \texttt{PcGroupCode} (\ref{PcGroupCode})). 

 The integer \mbox{\texttt{\mdseries\slshape n}} gives a certain amount of control over the probability to detect all
isomorphisms. If it is $0$, then nothing will be done at all. The larger \mbox{\texttt{\mdseries\slshape n}} is, the larger is the probability of finding all isomorphisms. However, due to
the underlying method we cannot guarantee that the algorithm finds all
isomorphisms, no matter how large \mbox{\texttt{\mdseries\slshape n}} is. }

 }

 }

          
\chapter{\textcolor{Chapter }{Finitely Presented Groups}}\label{Finitely Presented Groups}
\logpage{[ 47, 0, 0 ]}
\hyperdef{L}{X7AA982637E90B35A}{}
{
  A \emph{finitely presented group} (in short: FpGroup) is a group generated by a finite set of \emph{abstract generators} subject to a finite set of \emph{relations} that these generators satisfy. Every finite group can be represented as a
finitely presented group, though in almost all cases it is computationally
much more efficient to work in another representation (even the regular
permutation representation). 

 Finitely presented groups are obtained by factoring a free group by a set of
relators. Their elements know about this presentation and compare accordingly. 

 So to create a finitely presented group you first have to generate a free
group (see{\nobreakspace}\texttt{FreeGroup} (\ref{FreeGroup:for given rank}) for details). There are two ways to specify a quotient of the free group:
either by giving a list of relators or by giving a list of equations. Relators
are just words in the generators of the free group. Equations are represented
as pairs of words in the generators of the free group. In either case the
generators of the quotient are \emph{the images} of the free generators under the canonical homomorphism from the free group
onto the quotient. So for example to create the group 
\[ \langle a, b \mid a^2, b^3, (a b)^5 \rangle \]
 you can use the following commands: 
\begin{Verbatim}[commandchars=!@|,fontsize=\small,frame=single,label=Example]
  !gapprompt@gap>| !gapinput@f := FreeGroup( "a", "b" );;|
  !gapprompt@gap>| !gapinput@g := f / [ f.1^2, f.2^3, (f.1*f.2)^5 ];|
  <fp group on the generators [ a, b ]>
  !gapprompt@gap>| !gapinput@h := f / [ [f.1^2, f.1^0], [f.2^3, f.1^0], [(f.1*f.2)^4, f.2^-1*f.1^-1] ];|
  <fp group on the generators [ a, b ]>
\end{Verbatim}
 

 Note that you cannot call the generators by their names. These names are not
variables, but just display figures. So, if you want to access the generators
by their names, you first have to introduce the respective variables and to
assign the generators to them. 

 
\begin{Verbatim}[commandchars=!@|,fontsize=\small,frame=single,label=Example]
  !gapprompt@gap>| !gapinput@Unbind(a);|
  !gapprompt@gap>| !gapinput@GeneratorsOfGroup( g );|
  [ a, b ]
  !gapprompt@gap>| !gapinput@a;|
  Error, Variable: 'a' must have a value
  !gapprompt@gap>| !gapinput@a := g.1;; b := g.2;; # assign variables|
  !gapprompt@gap>| !gapinput@GeneratorsOfGroup( g );|
  [ a, b ]
  !gapprompt@gap>| !gapinput@a in f;|
  false
  !gapprompt@gap>| !gapinput@a in g;|
  true
\end{Verbatim}
 

 To relieve you of the tedium of typing the above assignments, \emph{when working interactively}, there is the function \texttt{AssignGeneratorVariables} (\ref{AssignGeneratorVariables}). 

 Note that the generators of the free group are different from the generators
of the FpGroup (even though they are displayed by the same names). That means
that words in the generators of the free group are not elements of the
finitely presented group. Vice versa elements of the FpGroup are not words. 

 
\begin{Verbatim}[commandchars=!@|,fontsize=\small,frame=single,label=Example]
  !gapprompt@gap>| !gapinput@a*b = b*a;|
  false
  !gapprompt@gap>| !gapinput@(b^2*a*b)^2 = a^0;|
  true
\end{Verbatim}
 

 Such calculations comparing elements of an FpGroup may run into problems:
There exist finitely presented groups for which no algorithm exists (it is
known that no such algorithm can exist) that will tell for two arbitrary words
in the generators whether the corresponding elements in the FpGroup are equal. 

 Therefore the methods used by \textsf{GAP} to compute in finitely presented groups may run into warning errors, run out
of memory or run forever. If the FpGroup is (by theory) known to be finite the
algorithms are guaranteed to terminate (if there is sufficient memory
available), but the time needed for the calculation cannot be bounded a
priori. See \ref{Coset Tables and Coset Enumeration} and \ref{Testing Finiteness of Finitely Presented Groups}. 

 
\begin{Verbatim}[commandchars=!@|,fontsize=\small,frame=single,label=Example]
  !gapprompt@gap>| !gapinput@(b^2*a*b)^2;|
  (b^2*a*b)^2
  !gapprompt@gap>| !gapinput@a^0;|
  <identity ...>
\end{Verbatim}
 

 A consequence of our convention is that elements of finitely presented groups
are not printed in a unique way. See also \texttt{SetReducedMultiplication} (\ref{SetReducedMultiplication}). For many ``higher level'' computations for a finite group, such as conjugacy classes or character table,
it is advisable not to use an FpGroup but an isomorphic group in a better
internal representation, such as a pc group (see \texttt{IsomorphismPcGroup} (\ref{IsomorphismPcGroup})) or a permutation group (see \texttt{IsomorphismPermGroup} (\ref{IsomorphismPermGroup})). In fact, calling a function such as \texttt{ConjugacyClasses} (\ref{ConjugacyClasses:attribute}) with an FpGroup for which \texttt{HasIsFinite} returns \texttt{false} may result in an error message, because the developers of the FpGroup code do
not want an \emph{automatic} \texttt{IsFinite} (\ref{IsFinite}) test (which may not terminate) in the function. (Once the \texttt{IsFinite} (\ref{IsFinite}) value is known, the computation in question can be done with the FpGroup, but
still it is advisable to use a better internal representation.)  
\section{\textcolor{Chapter }{IsSubgroupFpGroup and IsFpGroup}}\label{sect:IsSubgroupFpGroup}
\logpage{[ 47, 1, 0 ]}
\hyperdef{L}{X7824C8167B3CFAB1}{}
{
  

\subsection{\textcolor{Chapter }{IsSubgroupFpGroup}}
\logpage{[ 47, 1, 1 ]}\nobreak
\hyperdef{L}{X7AF7E2B48199452C}{}
{\noindent\textcolor{FuncColor}{$\triangleright$\enspace\texttt{IsSubgroupFpGroup({\mdseries\slshape H})\index{IsSubgroupFpGroup@\texttt{IsSubgroupFpGroup}}
\label{IsSubgroupFpGroup}
}\hfill{\scriptsize (Category)}}\\


 is the category for finitely presented groups or subgroups of a finitely
presented group. }

 

\subsection{\textcolor{Chapter }{IsFpGroup}}
\logpage{[ 47, 1, 2 ]}\nobreak
\hyperdef{L}{X850B9DF17D90C3A2}{}
{\noindent\textcolor{FuncColor}{$\triangleright$\enspace\texttt{IsFpGroup({\mdseries\slshape G})\index{IsFpGroup@\texttt{IsFpGroup}}
\label{IsFpGroup}
}\hfill{\scriptsize (filter)}}\\


 is a synonym for \texttt{IsSubgroupFpGroup(\mbox{\texttt{\mdseries\slshape G}}) and IsGroupOfFamily(\mbox{\texttt{\mdseries\slshape G}})}. 

 Free groups are a special case of finitely presented groups, namely finitely
presented groups with no relators. 

 Note that \texttt{FreeGroup(infinity)} (which exists e.g. for purposes of rewriting presentations with further
generators) satisfies this filter, though of course it is not finitely
generated (and thus not finitely presented). \texttt{IsFpGroup} thus is not a proper property test and slightly misnamed for the sake of its
most prominent uses. 

 Another special case are groups given by polycyclic presentations. \textsf{GAP} uses a special representation for these groups which is created in a different
way. See chapter \ref{Pc Groups} for details. 
\begin{Verbatim}[commandchars=!@|,fontsize=\small,frame=single,label=Example]
  !gapprompt@gap>| !gapinput@g:=FreeGroup(2);|
  <free group on the generators [ f1, f2 ]>
  !gapprompt@gap>| !gapinput@IsFpGroup(g);|
  true
  !gapprompt@gap>| !gapinput@h:=CyclicGroup(2);|
  <pc group of size 2 with 1 generator>
  !gapprompt@gap>| !gapinput@IsFpGroup(h);|
  false
\end{Verbatim}
 }

 

\subsection{\textcolor{Chapter }{InfoFpGroup}}
\logpage{[ 47, 1, 3 ]}\nobreak
\hyperdef{L}{X8370BF3B78D0B14D}{}
{\noindent\textcolor{FuncColor}{$\triangleright$\enspace\texttt{InfoFpGroup\index{InfoFpGroup@\texttt{InfoFpGroup}}
\label{InfoFpGroup}
}\hfill{\scriptsize (info class)}}\\


 The info class for functions dealing with finitely presented groups is \texttt{InfoFpGroup}. }

 }

  
\section{\textcolor{Chapter }{Creating Finitely Presented Groups}}\label{Creating Finitely Presented Groups}
\logpage{[ 47, 2, 0 ]}
\hyperdef{L}{X7D55E56E790F85FD}{}
{
  

\subsection{\textcolor{Chapter }{\texttt{\symbol{92}}/ (for a free group and a list of elements)}}
\logpage{[ 47, 2, 1 ]}\nobreak
\hyperdef{L}{X7EF4179E78BC7313}{}
{\noindent\textcolor{FuncColor}{$\triangleright$\enspace\texttt{\texttt{\symbol{92}}/({\mdseries\slshape F, rels})\index{\texttt{\symbol{92}}/@\texttt{\texttt{\symbol{92}}/}!for a free group and a list of elements}
\label{bSlash/:for a free group and a list of elements}
}\hfill{\scriptsize (method)}}\\
\noindent\textcolor{FuncColor}{$\triangleright$\enspace\texttt{\texttt{\symbol{92}}/({\mdseries\slshape F, eqns})\index{\texttt{\symbol{92}}/@\texttt{\texttt{\symbol{92}}/}!for a free group and a list of pairs of elements}
\label{bSlash/:for a free group and a list of pairs of elements}
}\hfill{\scriptsize (method)}}\\


 \index{quotient!for finitely presented groups} creates a finitely presented group given by the presentation $\langle gens \mid \mbox{\texttt{\mdseries\slshape rels}} \rangle$ or $\langle gens \mid \mbox{\texttt{\mdseries\slshape eqns}} \rangle$, respectively where $gens$ are the free generators of the free group \mbox{\texttt{\mdseries\slshape F}}. Relations can be entered either as words or as pairs of words in the
generators of \mbox{\texttt{\mdseries\slshape F}}. In the former case we refer to the words given as \emph{relators}, in the latter we refer to the pairs of words as \emph{equations}. The two methods can currently not be mixed. 

 The same result is obtained with the infix operator \texttt{/}, i.e., as \mbox{\texttt{\mdseries\slshape F}} \texttt{/} \mbox{\texttt{\mdseries\slshape rels}}. 

 
\begin{Verbatim}[commandchars=!@|,fontsize=\small,frame=single,label=Example]
  !gapprompt@gap>| !gapinput@f := FreeGroup( 3 );;|
  !gapprompt@gap>| !gapinput@f / [ f.1^4, f.2^3, f.3^5, f.1*f.2*f.3 ];|
  <fp group on the generators [ f1, f2, f3 ]>
  !gapprompt@gap>| !gapinput@f / [ [ f.1^4, f.1^0 ], [ f.2^3, f.1^0 ], [ f.1, f.2^-1*f.3^-1 ] ];|
  <fp group on the generators [ f1, f2, f3 ]>
\end{Verbatim}
 }

 

\subsection{\textcolor{Chapter }{FactorGroupFpGroupByRels}}
\logpage{[ 47, 2, 2 ]}\nobreak
\hyperdef{L}{X7CE0FA5F8695241E}{}
{\noindent\textcolor{FuncColor}{$\triangleright$\enspace\texttt{FactorGroupFpGroupByRels({\mdseries\slshape G, elts})\index{FactorGroupFpGroupByRels@\texttt{FactorGroupFpGroupByRels}}
\label{FactorGroupFpGroupByRels}
}\hfill{\scriptsize (function)}}\\


 returns the factor group \mbox{\texttt{\mdseries\slshape G}}/$N$ of \mbox{\texttt{\mdseries\slshape G}} by the normal closure $N$ of \mbox{\texttt{\mdseries\slshape elts}} where \mbox{\texttt{\mdseries\slshape elts}} is expected to be a list of elements of \mbox{\texttt{\mdseries\slshape G}}. }

 

\subsection{\textcolor{Chapter }{ParseRelators}}
\logpage{[ 47, 2, 3 ]}\nobreak
\hyperdef{L}{X7B3D290B87B6EFE4}{}
{\noindent\textcolor{FuncColor}{$\triangleright$\enspace\texttt{ParseRelators({\mdseries\slshape gens, rels})\index{ParseRelators@\texttt{ParseRelators}}
\label{ParseRelators}
}\hfill{\scriptsize (function)}}\\


 Will translate a list of relations as given in print, e.g. $x y^2 = (x y^3 x)^2 xy = yzx$ into relators. \mbox{\texttt{\mdseries\slshape gens}} must be a list of generators of a free group, each being displayed by a single
letter. \mbox{\texttt{\mdseries\slshape rels}} is a string that lists a sequence of equalities. These must be written in the
letters which are the names of the generators in \mbox{\texttt{\mdseries\slshape gens}}. Change of upper/lower case is interpreted to indicate inverses. 

 
\begin{Verbatim}[commandchars=!@|,fontsize=\small,frame=single,label=Example]
  !gapprompt@gap>| !gapinput@f:=FreeGroup("x","y","z");;|
  !gapprompt@gap>| !gapinput@AssignGeneratorVariables(f);|
  #I  Assigned the global variables [ x, y, z ]
  !gapprompt@gap>| !gapinput@r:=ParseRelators([x,y,z],|
  !gapprompt@>| !gapinput@"x^2 = y^5 = z^3 = (xyxyxy^4)^2 = (xz)^2 = (y^2z)^2 = 1");|
  [ x^2, y^5, z^3, (x*z)^2, (y^2*z)^2, ((x*y)^3*y^3)^2 ]
  !gapprompt@gap>| !gapinput@g:=f/r;|
  <fp group on the generators [ x, y, z ]>
\end{Verbatim}
 }

 

\subsection{\textcolor{Chapter }{StringFactorizationWord}}
\logpage{[ 47, 2, 4 ]}\nobreak
\hyperdef{L}{X85EAA789848B528E}{}
{\noindent\textcolor{FuncColor}{$\triangleright$\enspace\texttt{StringFactorizationWord({\mdseries\slshape w})\index{StringFactorizationWord@\texttt{StringFactorizationWord}}
\label{StringFactorizationWord}
}\hfill{\scriptsize (function)}}\\


 returns a string that expresses a given word \mbox{\texttt{\mdseries\slshape w}} in compact form written as a string. Inverses are expressed by changing the
upper/lower case of the generators, recurring expressions are written as
products. 
\begin{Verbatim}[commandchars=!@|,fontsize=\small,frame=single,label=Example]
  !gapprompt@gap>| !gapinput@StringFactorizationWord(z^-1*x*y*y*y*x*x*y*y*y*x*y^-1*x);|
  "Z(xy3x)2Yx"
\end{Verbatim}
 }

 }

  
\section{\textcolor{Chapter }{Comparison of Elements of Finitely Presented Groups}}\label{Comparison of Elements of Finitely Presented Groups}
\logpage{[ 47, 3, 0 ]}
\hyperdef{L}{X84D693EC872DAA55}{}
{
  

\subsection{\textcolor{Chapter }{\texttt{\symbol{92}}= (for two elements in a f.p. group)}}
\logpage{[ 47, 3, 1 ]}\nobreak
\hyperdef{L}{X797D29628203CBD6}{}
{\noindent\textcolor{FuncColor}{$\triangleright$\enspace\texttt{\texttt{\symbol{92}}=({\mdseries\slshape a, b})\index{\texttt{\symbol{92}}=@\texttt{\texttt{\symbol{92}}=}!for two elements in a f.p. group}
\label{bSlash=:for two elements in a f.p. group}
}\hfill{\scriptsize (method)}}\\


 \index{equality!elements of finitely presented groups} Two elements of a finitely presented group are equal if they are equal in this
group. Nevertheless they may be represented as different words in the
generators. Because of the fundamental problems mentioned in the introduction
to this chapter such a test may take very long and cannot be guaranteed to
finish. 

 The method employed by \textsf{GAP} for such an equality test use the underlying finitely presented group. First
(unless this group is known to be infinite) \textsf{GAP} tries to find a faithful permutation representation by a bounded
Todd\texttt{\symbol{45}}Coxeter. If this fails, a
Knuth\texttt{\symbol{45}}Bendix (see \ref{Rewriting Systems and the Knuth-Bendix Procedure}) is attempted and the words are compared via their normal form. 

 If only elements in a subgroup are to be tested for equality it thus can be
useful to translate the problem in a new finitely presented group by rewriting
(see \texttt{IsomorphismFpGroup} (\ref{IsomorphismFpGroup})); 

 The equality test of elements underlies many ``basic'' calculations, such as the order of an element, and the same type of problems
can arise there. In some cases, working with rewriting systems can still help
to solve the problem. The \textsf{kbmag} package provides such functionality, see the package manual for further
details. }

 

\subsection{\textcolor{Chapter }{\texttt{\symbol{92}}{\textless} (for two elements in a f.p. group)}}
\logpage{[ 47, 3, 2 ]}\nobreak
\hyperdef{L}{X7B350C718573B8DF}{}
{\noindent\textcolor{FuncColor}{$\triangleright$\enspace\texttt{\texttt{\symbol{92}}{\textless}({\mdseries\slshape a, b})\index{\texttt{\symbol{92}}{\textless}@\texttt{\texttt{\symbol{92}}{\textless}}!for two elements in a f.p. group}
\label{bSlash<:for two elements in a f.p. group}
}\hfill{\scriptsize (method)}}\\


 \index{smaller!elements of finitely presented groups} Compared with equality testing, problems get even worse when trying to compute
a total ordering on the elements of a finitely presented group. As any
ordering that is guaranteed to be reproducible in different runs of \textsf{GAP} or even with different groups given by syntactically equal presentations would
be prohibitively expensive to implement, the ordering of elements is depending
on a method chosen by \textsf{GAP} and not guaranteed to stay the same when repeating the construction of an
FpGroup. The only guarantee given for the \texttt{{\textless}} ordering for such elements is that it will stay the same for one family during
its lifetime. The attribute \texttt{FpElmComparisonMethod} (\ref{FpElmComparisonMethod}) is used to obtain a comparison function for a family of FpGroup elements. }

 

\subsection{\textcolor{Chapter }{FpElmComparisonMethod}}
\logpage{[ 47, 3, 3 ]}\nobreak
\hyperdef{L}{X87512CF485CC4128}{}
{\noindent\textcolor{FuncColor}{$\triangleright$\enspace\texttt{FpElmComparisonMethod({\mdseries\slshape fam})\index{FpElmComparisonMethod@\texttt{FpElmComparisonMethod}}
\label{FpElmComparisonMethod}
}\hfill{\scriptsize (attribute)}}\\


 If \mbox{\texttt{\mdseries\slshape fam}} is the elements family of a finitely presented group this attribute returns a
function \texttt{smaller(\mbox{\texttt{\mdseries\slshape left}}, \mbox{\texttt{\mdseries\slshape right}})} that will be used to compare elements in \mbox{\texttt{\mdseries\slshape fam}}. }

 

\subsection{\textcolor{Chapter }{SetReducedMultiplication}}
\logpage{[ 47, 3, 4 ]}\nobreak
\hyperdef{L}{X82CB9EC982CDAEAC}{}
{\noindent\textcolor{FuncColor}{$\triangleright$\enspace\texttt{SetReducedMultiplication({\mdseries\slshape obj})\index{SetReducedMultiplication@\texttt{SetReducedMultiplication}}
\label{SetReducedMultiplication}
}\hfill{\scriptsize (function)}}\\


 For an FpGroup \mbox{\texttt{\mdseries\slshape obj}}, an element \mbox{\texttt{\mdseries\slshape obj}} of it or the family \mbox{\texttt{\mdseries\slshape obj}} of its elements, this function will force immediate reduction when
multiplying, keeping words short at extra cost per multiplication. }

 }

  
\section{\textcolor{Chapter }{Preimages in the Free Group}}\label{Preimages in the Free Group}
\logpage{[ 47, 4, 0 ]}
\hyperdef{L}{X7B0B2781796800AD}{}
{
  

\subsection{\textcolor{Chapter }{FreeGroupOfFpGroup}}
\logpage{[ 47, 4, 1 ]}\nobreak
\hyperdef{L}{X85CF3931849FB441}{}
{\noindent\textcolor{FuncColor}{$\triangleright$\enspace\texttt{FreeGroupOfFpGroup({\mdseries\slshape G})\index{FreeGroupOfFpGroup@\texttt{FreeGroupOfFpGroup}}
\label{FreeGroupOfFpGroup}
}\hfill{\scriptsize (attribute)}}\\


 returns the underlying free group for the finitely presented group \mbox{\texttt{\mdseries\slshape G}}. This is the group generated by the free generators provided by the \texttt{FreeGeneratorsOfFpGroup} (\ref{FreeGeneratorsOfFpGroup}) value of \mbox{\texttt{\mdseries\slshape G}}. }

 

\subsection{\textcolor{Chapter }{FreeGeneratorsOfFpGroup}}
\logpage{[ 47, 4, 2 ]}\nobreak
\hyperdef{L}{X79C77C5184CA02B6}{}
{\noindent\textcolor{FuncColor}{$\triangleright$\enspace\texttt{FreeGeneratorsOfFpGroup({\mdseries\slshape G})\index{FreeGeneratorsOfFpGroup@\texttt{FreeGeneratorsOfFpGroup}}
\label{FreeGeneratorsOfFpGroup}
}\hfill{\scriptsize (attribute)}}\\
\noindent\textcolor{FuncColor}{$\triangleright$\enspace\texttt{FreeGeneratorsOfWholeGroup({\mdseries\slshape U})\index{FreeGeneratorsOfWholeGroup@\texttt{FreeGeneratorsOfWholeGroup}}
\label{FreeGeneratorsOfWholeGroup}
}\hfill{\scriptsize (operation)}}\\


 \texttt{FreeGeneratorsOfFpGroup} returns the underlying free generators corresponding to the generators of the
finitely presented group \mbox{\texttt{\mdseries\slshape G}} which must be a full FpGroup. 

 \texttt{FreeGeneratorsOfWholeGroup} also works for subgroups of an FpGroup and returns the free generators of the
full group that defines the family. }

 

\subsection{\textcolor{Chapter }{RelatorsOfFpGroup}}
\logpage{[ 47, 4, 3 ]}\nobreak
\hyperdef{L}{X87BA180287CD1F71}{}
{\noindent\textcolor{FuncColor}{$\triangleright$\enspace\texttt{RelatorsOfFpGroup({\mdseries\slshape G})\index{RelatorsOfFpGroup@\texttt{RelatorsOfFpGroup}}
\label{RelatorsOfFpGroup}
}\hfill{\scriptsize (attribute)}}\\


 returns the relators of the finitely presented group \mbox{\texttt{\mdseries\slshape G}} as words in the free generators provided by the \texttt{FreeGeneratorsOfFpGroup} (\ref{FreeGeneratorsOfFpGroup}) value of \mbox{\texttt{\mdseries\slshape G}}. 
\begin{Verbatim}[commandchars=!@|,fontsize=\small,frame=single,label=Example]
  !gapprompt@gap>| !gapinput@f := FreeGroup( "a", "b" );;|
  !gapprompt@gap>| !gapinput@g := f / [ f.1^5, f.2^2, f.1^f.2*f.1 ];|
  <fp group on the generators [ a, b ]>
  !gapprompt@gap>| !gapinput@Size( g );|
  10
  !gapprompt@gap>| !gapinput@FreeGroupOfFpGroup( g ) = f;|
  true
  !gapprompt@gap>| !gapinput@FreeGeneratorsOfFpGroup( g );|
  [ a, b ]
  !gapprompt@gap>| !gapinput@RelatorsOfFpGroup( g );|
  [ a^5, b^2, b^-1*a*b*a ]
\end{Verbatim}
 

 Note that these attributes are only available for the \emph{full} finitely presented group. It is possible (for example by using \texttt{Subgroup} (\ref{Subgroup})) to construct a subgroup of index $1$ which is not identical to the whole group. The latter one can be obtained in
this situation via \texttt{Parent} (\ref{Parent}). 

 Elements of a finitely presented group are not words, but are represented
using a word from the free group as representative. The following two commands
obtain this representative, respectively create an element in the finitely
presented group. }

 

\subsection{\textcolor{Chapter }{UnderlyingElement (fp group elements)}}
\logpage{[ 47, 4, 4 ]}\nobreak
\hyperdef{L}{X8447A2397A1E524B}{}
{\noindent\textcolor{FuncColor}{$\triangleright$\enspace\texttt{UnderlyingElement({\mdseries\slshape elm})\index{UnderlyingElement@\texttt{UnderlyingElement}!fp group elements}
\label{UnderlyingElement:fp group elements}
}\hfill{\scriptsize (operation)}}\\


 Let \mbox{\texttt{\mdseries\slshape elm}} be an element of a group whose elements are represented as words with further
properties. Then \texttt{UnderlyingElement} returns the word from the free group that is used as a representative for \mbox{\texttt{\mdseries\slshape elm}}. 

 
\begin{Verbatim}[commandchars=!@|,fontsize=\small,frame=single,label=Example]
  !gapprompt@gap>| !gapinput@w := g.1*g.2;|
  a*b
  !gapprompt@gap>| !gapinput@IsWord( w );|
  false
  !gapprompt@gap>| !gapinput@ue := UnderlyingElement( w );|
  a*b
  !gapprompt@gap>| !gapinput@IsWord( ue );|
  true
\end{Verbatim}
 }

 

\subsection{\textcolor{Chapter }{ElementOfFpGroup}}
\logpage{[ 47, 4, 5 ]}\nobreak
\hyperdef{L}{X7F34C8017DC03FDB}{}
{\noindent\textcolor{FuncColor}{$\triangleright$\enspace\texttt{ElementOfFpGroup({\mdseries\slshape fam, word})\index{ElementOfFpGroup@\texttt{ElementOfFpGroup}}
\label{ElementOfFpGroup}
}\hfill{\scriptsize (operation)}}\\


 If \mbox{\texttt{\mdseries\slshape fam}} is the elements family of a finitely presented group and \mbox{\texttt{\mdseries\slshape word}} is a word in the free generators underlying this finitely presented group,
this operation creates the element with the representative \mbox{\texttt{\mdseries\slshape word}} in the free group. 
\begin{Verbatim}[commandchars=!@|,fontsize=\small,frame=single,label=Example]
  !gapprompt@gap>| !gapinput@ge := ElementOfFpGroup( FamilyObj( g.1 ), f.1*f.2 );|
  a*b
  !gapprompt@gap>| !gapinput@ge in f;|
  false
  !gapprompt@gap>| !gapinput@ge in g;|
  true
\end{Verbatim}
 }

 }

  
\section{\textcolor{Chapter }{Operations for Finitely Presented Groups}}\label{Operations for Finitely Presented Groups}
\logpage{[ 47, 5, 0 ]}
\hyperdef{L}{X869143D284F3379D}{}
{
  Finitely presented groups are groups and so all operations for groups should
be applicable to them (though not necessarily efficient methods are
available). However, one may run into errors if the FpGroup does not store
that it is finite; see the introduction to this chapter. Most methods for
finitely presented groups rely on coset enumeration. See{\nobreakspace}\ref{Coset Tables and Coset Enumeration} for details. 

 The command \texttt{IsomorphismPermGroup} (\ref{IsomorphismPermGroup}) can be used to obtain a faithful permutation representation, if such a
representation of small degree exists. (Otherwise it might run very long or
fail.) 
\begin{Verbatim}[commandchars=!@|,fontsize=\small,frame=single,label=Example]
  !gapprompt@gap>| !gapinput@f := FreeGroup( "a", "b" );|
  <free group on the generators [ a, b ]>
  !gapprompt@gap>| !gapinput@g := f / [ f.1^2, f.2^3, (f.1*f.2)^5 ];|
  <fp group on the generators [ a, b ]>
  !gapprompt@gap>| !gapinput@h := IsomorphismPermGroup( g );|
  [ a, b ] -> [ (2,4)(5,6), (1,2,3)(4,5,6) ]
  !gapprompt@gap>| !gapinput@u:=Subgroup(g,[g.1*g.2]);;rt:=RightTransversal(g,u);|
  RightTransversal(<fp group of size 60 on the generators
  [ a, b ]>,Group([ a*b ]))
  !gapprompt@gap>| !gapinput@Image(ActionHomomorphism(g,rt,OnRight));|
  Group([ (1,2)(3,4)(5,7)(6,8)(9,10)(11,12),
    (1,3,2)(4,5,6)(7,8,9)(10,11,12) ])
\end{Verbatim}
 

\subsection{\textcolor{Chapter }{PseudoRandom (for finitely presented groups)}}
\logpage{[ 47, 5, 1 ]}\nobreak
\hyperdef{L}{X7AB7187779EDC9BA}{}
{\noindent\textcolor{FuncColor}{$\triangleright$\enspace\texttt{PseudoRandom({\mdseries\slshape F: radius := l})\index{PseudoRandom@\texttt{PseudoRandom}!for finitely presented groups}
\label{PseudoRandom:for finitely presented groups}
}\hfill{\scriptsize (method)}}\\


 The default algorithm for \texttt{PseudoRandom} (\ref{PseudoRandom}) makes little sense for finitely presented or free groups, as it produces words
that are extremely long. 

 By specifying the option \texttt{radius}, instead elements are taken as words in the generators of \mbox{\texttt{\mdseries\slshape F}} in the ball of radius \mbox{\texttt{\mdseries\slshape l}} with equal distribution in the free group. 

 
\begin{Verbatim}[commandchars=!@|,fontsize=\small,frame=single,label=Example]
  !gapprompt@gap>| !gapinput@PseudoRandom(g:radius:=20);|
  a^3*b^2*a^-2*b^-1*a*b^-4*a*b^-1*a*b^-4
\end{Verbatim}
 }

 }

  
\section{\textcolor{Chapter }{Coset Tables and Coset Enumeration}}\label{Coset Tables and Coset Enumeration}
\logpage{[ 47, 6, 0 ]}
\hyperdef{L}{X7BD0CEBA7B225416}{}
{
  Coset enumeration (see \cite{Neu82} for an explanation) is one of the fundamental tools for the examination of
finitely presented groups. This section describes \textsf{GAP} functions that can be used to invoke a coset enumeration. 

 Note that in addition to the built\texttt{\symbol{45}}in coset enumerator
there is the \textsf{GAP} package \textsf{ACE}. Moreover, \textsf{GAP} provides an interactive Todd\texttt{\symbol{45}}Coxeter in the \textsf{GAP} package \textsf{ITC} which is based on the \textsf{XGAP} package. 

\subsection{\textcolor{Chapter }{CosetTable}}
\logpage{[ 47, 6, 1 ]}\nobreak
\hyperdef{L}{X7F7F31E47D7F6EF8}{}
{\noindent\textcolor{FuncColor}{$\triangleright$\enspace\texttt{CosetTable({\mdseries\slshape G, H})\index{CosetTable@\texttt{CosetTable}}
\label{CosetTable}
}\hfill{\scriptsize (operation)}}\\


 returns the coset table of the finitely presented group \mbox{\texttt{\mdseries\slshape G}} on the cosets of the subgroup \mbox{\texttt{\mdseries\slshape H}}. 

 Basically a coset table is the permutation representation of the finitely
presented group on the cosets of a subgroup (which need not be faithful if the
subgroup has a nontrivial core). Most of the set theoretic and group functions
use the regular representation of \mbox{\texttt{\mdseries\slshape G}}, i.e., the coset table of \mbox{\texttt{\mdseries\slshape G}} over the trivial subgroup. 

 The coset table is returned as a list of lists. For each generator of \mbox{\texttt{\mdseries\slshape G}} and its inverse the table contains a generator list. A generator list is
simply a list of integers. If $l$ is the generator list for the generator $g$ and if $l[i] = j$ then generator $g$ takes the coset $i$ to the coset $j$ by multiplication from the right. Thus the permutation representation of \mbox{\texttt{\mdseries\slshape G}} on the cosets of \mbox{\texttt{\mdseries\slshape H}} is obtained by applying \texttt{PermList} (\ref{PermList}) to each generator list. 

 The coset table is standard (see below). 

 For finitely presented groups, a coset table is computed by a
Todd\texttt{\symbol{45}}Coxeter coset enumeration. Note that you may influence
the performance of that enumeration by changing the values of the global
variables \texttt{CosetTableDefaultLimit} (\ref{CosetTableDefaultLimit}) and \texttt{CosetTableDefaultMaxLimit} (\ref{CosetTableDefaultMaxLimit}) described below and that the options described under \texttt{CosetTableFromGensAndRels} (\ref{CosetTableFromGensAndRels}) are recognized. 

 
\begin{Verbatim}[commandchars=!@|,fontsize=\small,frame=single,label=Example]
  !gapprompt@gap>| !gapinput@tab := CosetTable(g, Subgroup(g, [ g.1, g.2*g.1*g.2*g.1*g.2^-1 ]));|
  [ [ 1, 4, 5, 2, 3 ], [ 1, 4, 5, 2, 3 ], [ 2, 3, 1, 4, 5 ],
    [ 3, 1, 2, 4, 5 ] ]
  !gapprompt@gap>| !gapinput@List( last, PermList );|
  [ (2,4)(3,5), (2,4)(3,5), (1,2,3), (1,3,2) ]
  !gapprompt@gap>| !gapinput@PrintArray( TransposedMat( tab ) );|
  [ [  1,  1,  2,  3 ],
    [  4,  4,  3,  1 ],
    [  5,  5,  1,  2 ],
    [  2,  2,  4,  4 ],
    [  3,  3,  5,  5 ] ]
\end{Verbatim}
 

 The last printout in the preceding example provides the coset table in the
form in which it is usually used in hand calculations: The rows correspond to
the cosets, the columns correspond to the generators and their inverses in the
ordering $g_1, g_1^{{-1}}, g_2, g_2^{{-1}}$. (See section{\nobreakspace}\ref{Standardization of coset tables} for a description on the way the numbers are assigned.) }

 

\subsection{\textcolor{Chapter }{TracedCosetFpGroup}}
\logpage{[ 47, 6, 2 ]}\nobreak
\hyperdef{L}{X87D175757C581E62}{}
{\noindent\textcolor{FuncColor}{$\triangleright$\enspace\texttt{TracedCosetFpGroup({\mdseries\slshape tab, word, pt})\index{TracedCosetFpGroup@\texttt{TracedCosetFpGroup}}
\label{TracedCosetFpGroup}
}\hfill{\scriptsize (function)}}\\


 Traces the coset number \mbox{\texttt{\mdseries\slshape pt}} under the word \mbox{\texttt{\mdseries\slshape word}} through the coset table \mbox{\texttt{\mdseries\slshape tab}}. (Note: \mbox{\texttt{\mdseries\slshape word}} must be in the free group, use \texttt{UnderlyingElement} (\ref{UnderlyingElement:fp group elements}) if in doubt.) 
\begin{Verbatim}[commandchars=!@|,fontsize=\small,frame=single,label=Example]
  !gapprompt@gap>| !gapinput@TracedCosetFpGroup(tab,UnderlyingElement(g.1),2);|
  4
\end{Verbatim}
 }

 

\subsection{\textcolor{Chapter }{FactorCosetAction (for fp groups)}}
\logpage{[ 47, 6, 3 ]}\nobreak
\hyperdef{L}{X7EC1B0EE876E478A}{}
{\noindent\textcolor{FuncColor}{$\triangleright$\enspace\texttt{FactorCosetAction({\mdseries\slshape G, H})\index{FactorCosetAction@\texttt{FactorCosetAction}!for fp groups}
\label{FactorCosetAction:for fp groups}
}\hfill{\scriptsize (operation)}}\\


 returns the action of \mbox{\texttt{\mdseries\slshape G}} on the cosets of its subgroup \mbox{\texttt{\mdseries\slshape H}}. 

 
\begin{Verbatim}[commandchars=!@|,fontsize=\small,frame=single,label=Example]
  !gapprompt@gap>| !gapinput@u := Subgroup( g, [ g.1, g.1^g.2 ] );|
  Group([ a, b^-1*a*b ])
  !gapprompt@gap>| !gapinput@FactorCosetAction( g, u );|
  [ a, b ] -> [ (2,4)(5,6), (1,2,3)(4,5,6) ]
\end{Verbatim}
 }

 

\subsection{\textcolor{Chapter }{CosetTableBySubgroup}}
\logpage{[ 47, 6, 4 ]}\nobreak
\hyperdef{L}{X82926A7F8365A341}{}
{\noindent\textcolor{FuncColor}{$\triangleright$\enspace\texttt{CosetTableBySubgroup({\mdseries\slshape G, H})\index{CosetTableBySubgroup@\texttt{CosetTableBySubgroup}}
\label{CosetTableBySubgroup}
}\hfill{\scriptsize (operation)}}\\


 returns a coset table for the action of \mbox{\texttt{\mdseries\slshape G}} on the cosets of \mbox{\texttt{\mdseries\slshape H}}. The columns of the table correspond to the \texttt{GeneratorsOfGroup} (\ref{GeneratorsOfGroup}) value of \mbox{\texttt{\mdseries\slshape G}}. }

 

\subsection{\textcolor{Chapter }{CosetTableFromGensAndRels}}
\logpage{[ 47, 6, 5 ]}\nobreak
\hyperdef{L}{X7DE601F179E6FD09}{}
{\noindent\textcolor{FuncColor}{$\triangleright$\enspace\texttt{CosetTableFromGensAndRels({\mdseries\slshape fgens, grels, fsgens})\index{CosetTableFromGensAndRels@\texttt{CosetTableFromGensAndRels}}
\label{CosetTableFromGensAndRels}
}\hfill{\scriptsize (function)}}\\


 \index{TCENUM@\texttt{TCENUM}} \index{GAPTCENUM@\texttt{GAPTCENUM}} is an internal function which is called by the functions \texttt{CosetTable} (\ref{CosetTable}), \texttt{CosetTableInWholeGroup} (\ref{CosetTableInWholeGroup}) and others. It is, in fact, the workhorse that performs a
Todd\texttt{\symbol{45}}Coxeter coset enumeration. \mbox{\texttt{\mdseries\slshape fgens}} must be a set of free generators and \mbox{\texttt{\mdseries\slshape grels}} a set of relators in these generators. \mbox{\texttt{\mdseries\slshape fsgens}} are subgroup generators expressed as words in these generators. The function
returns a coset table with respect to \mbox{\texttt{\mdseries\slshape fgens}}. 

 \texttt{CosetTableFromGensAndRels} will call \texttt{TCENUM.CosetTableFromGensAndRels}. This makes it possible to replace the built\texttt{\symbol{45}}in coset
enumerator with another one by assigning \texttt{TCENUM} to another record. 

 The library version which is used by default performs a standard Felsch
strategy coset enumeration. You can call this function explicitly as \texttt{GAPTCENUM.CosetTableFromGensAndRels} even if other coset enumerators are installed. 

 The expected parameters are 
\begin{description}
\item[{\mbox{\texttt{\mdseries\slshape fgens}}}]  generators of the free group \mbox{\texttt{\mdseries\slshape F}} 
\item[{\mbox{\texttt{\mdseries\slshape grels}}}]  relators as words in \mbox{\texttt{\mdseries\slshape F}} 
\item[{\mbox{\texttt{\mdseries\slshape fsgens}}}]  subgroup generators as words in \mbox{\texttt{\mdseries\slshape F}}. 
\end{description}
 

 \texttt{CosetTableFromGensAndRels} processes two options (see chapter{\nobreakspace}\ref{Options Stack}): 
\begin{description}
\item[{\texttt{max}}]  The limit of the number of cosets to be defined. If the enumeration does not
finish with this number of cosets, an error is raised and the user is asked
whether she wants to continue. The default value is the value given in the
variable \texttt{CosetTableDefaultMaxLimit}. (Due to the algorithm the actual limit used can be a bit higher than the
number given.) 
\item[{\texttt{silent}}]  If set to \texttt{true} the algorithm will not raise the error mentioned under option \texttt{max} but silently return \texttt{fail}. This can be useful if an enumeration is only wanted unless it becomes too
big. 
\end{description}
 }

 

\subsection{\textcolor{Chapter }{CosetTableDefaultMaxLimit}}
\logpage{[ 47, 6, 6 ]}\nobreak
\hyperdef{L}{X822B188F87E9E642}{}
{\noindent\textcolor{FuncColor}{$\triangleright$\enspace\texttt{CosetTableDefaultMaxLimit\index{CosetTableDefaultMaxLimit@\texttt{CosetTableDefaultMaxLimit}}
\label{CosetTableDefaultMaxLimit}
}\hfill{\scriptsize (global variable)}}\\


 is the default limit for the number of cosets allowed in a coset enumeration. 

 A coset enumeration will not finish if the subgroup does not have finite
index, and even if it has it may take many more intermediate cosets than the
actual index of the subgroup is. To avoid a coset enumeration ``running away'' therefore \textsf{GAP} has a ``safety stop'' built in. This is controlled by the global variable \texttt{CosetTableDefaultMaxLimit}. 

 If this number of cosets is reached, \textsf{GAP} will issue an error message and prompt the user to either continue the
calculation or to stop it. The default value is $4096000$. 

 See also the description of the options to \texttt{CosetTableFromGensAndRels} (\ref{CosetTableFromGensAndRels}). 

 
\begin{Verbatim}[commandchars=!|B,fontsize=\small,frame=single,label=Example]
  !gapprompt|gap>B !gapinput|f := FreeGroup( "a", "b" );;B
  !gapprompt|gap>B !gapinput|u := Subgroup( f, [ f.2 ] );B
  Group([ b ])
  !gapprompt|gap>B !gapinput|Index( f, u );B
  Error, the coset enumeration has defined more than 4096000 cosets
  Stack trace:
  *[1] Error( "the coset enumeration has defined more ", "than ", limit, " cosets\n" );
     @ GAPROOT/lib/grpfp.gi:1230
   [2] TCENUM.CosetTableFromGensAndRels( fgens, grels, fsgens )
     @ GAPROOT/lib/grpfp.gi:1068
   [3] CosetTableFromGensAndRels( fgens, grels, fsgens )
     @ GAPROOT/lib/grpfp.gi:1346
   [4] TryCosetTableInWholeGroup( H )
     @ GAPROOT/lib/grpfp.gi:1359
   [5] CosetTableInWholeGroup( H )
     @ GAPROOT/lib/grpfp.gi:1733
  ...  at *stdin*:3
  you can enter 'return;' to continue with a new limit of 8192000 cosets,
  you can enter 'quit;' to abort the coset enumeration,
  you can enter 'maxlimit := 0; return;' in order to continue without a limit
  !gapbrkprompt|brk>B !gapinput|quit;B
\end{Verbatim}
 

 At this point, a \texttt{break}\texttt{\symbol{45}}loop (see Section{\nobreakspace}\ref{Break Loops}) has been entered. The line beginning \texttt{Error} tells you why this occurred. The next seven lines occur if \texttt{OnBreak} (\ref{OnBreak}) has its default value \texttt{Where} (\ref{Where}). They explain, in this case, how \textsf{GAP} came to be doing a coset enumeration. Then you are given a number of options
of how to escape the \texttt{break}\texttt{\symbol{45}}loop: you can either continue the calculation with a
larger number of permitted cosets, stop the calculation if you don't expect
the enumeration to finish (like in the example above), or continue without a
limit on the number of cosets. (Choosing the first option will, of course,
land you back in a \texttt{break}\texttt{\symbol{45}}loop. Try it!) 

 Setting \texttt{CosetTableDefaultMaxLimit} (or the \texttt{max} option value, for any function that invokes a coset enumeration) to \texttt{infinity} (\ref{infinity}) (or to $0$) will force all coset enumerations to continue until they either get a result
or exhaust the whole available space. For example, each of the following two
inputs 

 
\begin{Verbatim}[commandchars=!@|,fontsize=\small,frame=single,label=]
  !gapprompt@gap>| !gapinput@CosetTableDefaultMaxLimit := 0;;|
  !gapprompt@gap>| !gapinput@Index( f, u );|
\end{Verbatim}
 

 or 

 
\begin{Verbatim}[commandchars=!@|,fontsize=\small,frame=single,label=]
  !gapprompt@gap>| !gapinput@Index( f, u : max := 0 );|
\end{Verbatim}
 

 have essentially the same effect as choosing the third option (typing: \texttt{maxlimit := 0; return;}) at the \texttt{brk{\textgreater}} prompt above (instead of \texttt{quit;}). }

 

\subsection{\textcolor{Chapter }{CosetTableDefaultLimit}}
\logpage{[ 47, 6, 7 ]}\nobreak
\hyperdef{L}{X7A80A00E7E088E44}{}
{\noindent\textcolor{FuncColor}{$\triangleright$\enspace\texttt{CosetTableDefaultLimit\index{CosetTableDefaultLimit@\texttt{CosetTableDefaultLimit}}
\label{CosetTableDefaultLimit}
}\hfill{\scriptsize (global variable)}}\\


 is the default number of cosets with which any coset table is initialized
before doing a coset enumeration. 

 The function performing this coset enumeration will automatically extend the
table whenever necessary (as long as the number of cosets does not exceed the
value of \texttt{CosetTableDefaultMaxLimit} (\ref{CosetTableDefaultMaxLimit})), but this is an expensive operation. Thus, if you change the value of \texttt{CosetTableDefaultLimit}, you should set it to a number of cosets that you expect to be sufficient for
your subsequent coset enumerations. On the other hand, if you make it too
large, your job will unnecessarily waste a lot of space. 

 The default value of \texttt{CosetTableDefaultLimit} is $1000$. }

 

\subsection{\textcolor{Chapter }{MostFrequentGeneratorFpGroup}}
\logpage{[ 47, 6, 8 ]}\nobreak
\hyperdef{L}{X829D31A981CB2AF4}{}
{\noindent\textcolor{FuncColor}{$\triangleright$\enspace\texttt{MostFrequentGeneratorFpGroup({\mdseries\slshape G})\index{MostFrequentGeneratorFpGroup@\texttt{MostFrequentGeneratorFpGroup}}
\label{MostFrequentGeneratorFpGroup}
}\hfill{\scriptsize (function)}}\\


 is an internal function which is used in some applications of coset table
methods. It returns the first of those generators of the given finitely
presented group \mbox{\texttt{\mdseries\slshape G}} which occur most frequently in the relators. }

 

\subsection{\textcolor{Chapter }{IndicesInvolutaryGenerators}}
\logpage{[ 47, 6, 9 ]}\nobreak
\hyperdef{L}{X7912E6577B577A5C}{}
{\noindent\textcolor{FuncColor}{$\triangleright$\enspace\texttt{IndicesInvolutaryGenerators({\mdseries\slshape G})\index{IndicesInvolutaryGenerators@\texttt{IndicesInvolutaryGenerators}}
\label{IndicesInvolutaryGenerators}
}\hfill{\scriptsize (attribute)}}\\


 returns the indices of those generators of the finitely presented group \mbox{\texttt{\mdseries\slshape G}} which are known to be involutions. This knowledge is used by internal
functions to improve the performance of coset enumerations. }

 }

  
\section{\textcolor{Chapter }{Standardization of coset tables}}\label{Standardization of coset tables}
\logpage{[ 47, 7, 0 ]}
\hyperdef{L}{X85B882F782D7AFD0}{}
{
  For any two coset numbers $i$ and $j$ with $i < j$ the first occurrence of $i$ in a coset table precedes the first occurrence of $j$ with respect to the usual row\texttt{\symbol{45}}wise ordering of the table
entries. Following the notation of Charles Sims' book on computation with
finitely presented groups \cite{Sims94} we call such a table a \emph{standard coset table}. 

 The table entries which contain the first occurrences of the coset numbers $i > 1$ recursively provide for each $i$ a representative of the corresponding coset in form of a unique word $w_i$ in the generators and inverse generators of $G$. The first coset (which is $H$ itself) can be represented by the empty word $w_1$. A coset table is standard if and only if the words $w_1, w_2, \ldots$ are length\texttt{\symbol{45}}plus\texttt{\symbol{45}}lexicographic ordered
(as defined in \cite{Sims94}), for short: \emph{lenlex}. 

 This standardization of coset tables is different from that used in \textsf{GAP} versions 4.2 and earlier. Before that, we ignored the columns that correspond
to inverse generators and hence only considered words in the generators of $G$. We call this older ordering the \emph{semilenlex} standard as it also applies to the case of semigroups where no inverses of the
generators are known. 

 We changed our default from the semilenlex standard to the lenlex standard to
be consistent with \cite{Sims94}. However, the semilenlex standardisation remains available and the convention
used for all implicit standardisations can be selected by setting the value of
the global variable \texttt{CosetTableStandard} (\ref{CosetTableStandard}) to either \texttt{"lenlex"} or \texttt{"semilenlex"}. Independent of the current value of \texttt{CosetTableStandard} (\ref{CosetTableStandard}) you can standardize (or restandardize) a coset table at any time using \texttt{StandardizeTable} (\ref{StandardizeTable}). 

\subsection{\textcolor{Chapter }{CosetTableStandard}}
\logpage{[ 47, 7, 1 ]}\nobreak
\hyperdef{L}{X85FD1D637EF1EBE7}{}
{\noindent\textcolor{FuncColor}{$\triangleright$\enspace\texttt{CosetTableStandard\index{CosetTableStandard@\texttt{CosetTableStandard}}
\label{CosetTableStandard}
}\hfill{\scriptsize (global variable)}}\\


 specifies the definition of a \emph{standard coset table}. It is used whenever coset tables or augmented coset tables are created. Its
value may be \texttt{"lenlex"} or \texttt{"semilenlex"}. If it is \texttt{"lenlex"} coset tables will be standardized using all their columns as defined in
Charles Sims' book (this is the new default standard of \textsf{GAP}). If it is \texttt{"semilenlex"} they will be standardized using only their generator columns (this was the
original \textsf{GAP} standard). The default value of \texttt{CosetTableStandard} is \texttt{"lenlex"}. }

 

\subsection{\textcolor{Chapter }{StandardizeTable}}
\logpage{[ 47, 7, 2 ]}\nobreak
\hyperdef{L}{X85FCD8DF81BA94D5}{}
{\noindent\textcolor{FuncColor}{$\triangleright$\enspace\texttt{StandardizeTable({\mdseries\slshape table, standard})\index{StandardizeTable@\texttt{StandardizeTable}}
\label{StandardizeTable}
}\hfill{\scriptsize (function)}}\\


 standardizes the given coset table \mbox{\texttt{\mdseries\slshape table}}. The second argument is optional. It defines the standard to be used, its
values may be \texttt{"lenlex"} or \texttt{"semilenlex"} specifying the new or the old convention, respectively. If no value for the
parameter \mbox{\texttt{\mdseries\slshape standard}} is provided the function will use the global variable \texttt{CosetTableStandard} (\ref{CosetTableStandard}) instead. Note that the function alters the given table, it does not create a
copy. 
\begin{Verbatim}[commandchars=!@|,fontsize=\small,frame=single,label=Example]
  !gapprompt@gap>| !gapinput@StandardizeTable( tab, "semilenlex" );|
  !gapprompt@gap>| !gapinput@PrintArray( TransposedMat( tab ) );|
  [ [  1,  1,  2,  4 ],
    [  3,  3,  4,  1 ],
    [  2,  2,  3,  3 ],
    [  5,  5,  1,  2 ],
    [  4,  4,  5,  5 ] ]
\end{Verbatim}
 }

 }

  
\section{\textcolor{Chapter }{Coset tables for subgroups in the whole group}}\label{Coset tables for subgroups in the whole group}
\logpage{[ 47, 8, 0 ]}
\hyperdef{L}{X87C3FA0784A85309}{}
{
  

\subsection{\textcolor{Chapter }{CosetTableInWholeGroup}}
\logpage{[ 47, 8, 1 ]}\nobreak
\hyperdef{L}{X846EC8AB7803114D}{}
{\noindent\textcolor{FuncColor}{$\triangleright$\enspace\texttt{CosetTableInWholeGroup({\mdseries\slshape H})\index{CosetTableInWholeGroup@\texttt{CosetTableInWholeGroup}}
\label{CosetTableInWholeGroup}
}\hfill{\scriptsize (attribute)}}\\
\noindent\textcolor{FuncColor}{$\triangleright$\enspace\texttt{TryCosetTableInWholeGroup({\mdseries\slshape H})\index{TryCosetTableInWholeGroup@\texttt{TryCosetTableInWholeGroup}}
\label{TryCosetTableInWholeGroup}
}\hfill{\scriptsize (operation)}}\\


 is equivalent to \texttt{CosetTable(\mbox{\texttt{\mdseries\slshape G}},\mbox{\texttt{\mdseries\slshape H}})} where \mbox{\texttt{\mdseries\slshape G}} is the (unique) finitely presented group such that \mbox{\texttt{\mdseries\slshape H}} is a subgroup of \mbox{\texttt{\mdseries\slshape G}}. It overrides a \texttt{silent} option (see{\nobreakspace}\texttt{CosetTableFromGensAndRels} (\ref{CosetTableFromGensAndRels})) with \texttt{false}. 

 The variant \texttt{TryCosetTableInWholeGroup} does not override the \texttt{silent} option with \texttt{false} in case a coset table is only wanted if not too expensive. It will store a
result that is not \texttt{fail} in the attribute \texttt{CosetTableInWholeGroup}. }

 

\subsection{\textcolor{Chapter }{SubgroupOfWholeGroupByCosetTable}}
\logpage{[ 47, 8, 2 ]}\nobreak
\hyperdef{L}{X857F239583AFE0B7}{}
{\noindent\textcolor{FuncColor}{$\triangleright$\enspace\texttt{SubgroupOfWholeGroupByCosetTable({\mdseries\slshape fpfam, tab})\index{SubgroupOfWholeGroupByCosetTable@\texttt{SubgroupOfWholeGroupByCosetTable}}
\label{SubgroupOfWholeGroupByCosetTable}
}\hfill{\scriptsize (function)}}\\


 takes a family \mbox{\texttt{\mdseries\slshape fpfam}} of an FpGroup and a standardized coset table \mbox{\texttt{\mdseries\slshape tab}} and returns the subgroup of \mbox{\texttt{\mdseries\slshape fpfam}}\texttt{!.wholeGroup} defined by this coset table. The function will not check whether the coset
table is standardized. See also{\nobreakspace}\texttt{CosetTableBySubgroup} (\ref{CosetTableBySubgroup}). }

 }

  
\section{\textcolor{Chapter }{Augmented Coset Tables and Rewriting}}\label{Augmented Coset Tables and Rewriting}
\logpage{[ 47, 9, 0 ]}
\hyperdef{L}{X7E17A14E823F953D}{}
{
  

\subsection{\textcolor{Chapter }{AugmentedCosetTableInWholeGroup}}
\logpage{[ 47, 9, 1 ]}\nobreak
\hyperdef{L}{X80F8BF1D867DA7C1}{}
{\noindent\textcolor{FuncColor}{$\triangleright$\enspace\texttt{AugmentedCosetTableInWholeGroup({\mdseries\slshape H[, gens]})\index{AugmentedCosetTableInWholeGroup@\texttt{AugmentedCosetTableInWholeGroup}}
\label{AugmentedCosetTableInWholeGroup}
}\hfill{\scriptsize (function)}}\\


 For a subgroup \mbox{\texttt{\mdseries\slshape H}} of a finitely presented group, this function returns an augmented coset table.
If a generator set \mbox{\texttt{\mdseries\slshape gens}} is given, it is guaranteed that \mbox{\texttt{\mdseries\slshape gens}} will be a subset of the primary and secondary subgroup generators of this
coset table. 

 It is mutable so we are permitted to add further entries. However existing
entries may not be changed. Any entries added however should correspond to the
subgroup only and not to a homomorphism. }

 

\subsection{\textcolor{Chapter }{AugmentedCosetTableMtc}}
\logpage{[ 47, 9, 2 ]}\nobreak
\hyperdef{L}{X7AF67CFD846C1159}{}
{\noindent\textcolor{FuncColor}{$\triangleright$\enspace\texttt{AugmentedCosetTableMtc({\mdseries\slshape G, H, type, string})\index{AugmentedCosetTableMtc@\texttt{AugmentedCosetTableMtc}}
\label{AugmentedCosetTableMtc}
}\hfill{\scriptsize (function)}}\\


 is an internal function used by the subgroup presentation functions described
in \ref{Subgroup Presentations}. It applies a Modified Todd\texttt{\symbol{45}}Coxeter coset representative
enumeration to construct an augmented coset table (see \ref{Subgroup Presentations}) for the given subgroup \mbox{\texttt{\mdseries\slshape H}} of \mbox{\texttt{\mdseries\slshape G}}. The subgroup generators will be named \mbox{\texttt{\mdseries\slshape string}}\texttt{1}, \mbox{\texttt{\mdseries\slshape string}}\texttt{2}, $\ldots$. 

 The function accepts the options \texttt{max} and \texttt{silent} as described for the function \texttt{CosetTableFromGensAndRels} (\ref{CosetTableFromGensAndRels}). }

 

\subsection{\textcolor{Chapter }{AugmentedCosetTableRrs}}
\logpage{[ 47, 9, 3 ]}\nobreak
\hyperdef{L}{X7F3F09C778552811}{}
{\noindent\textcolor{FuncColor}{$\triangleright$\enspace\texttt{AugmentedCosetTableRrs({\mdseries\slshape G, table, type, string})\index{AugmentedCosetTableRrs@\texttt{AugmentedCosetTableRrs}}
\label{AugmentedCosetTableRrs}
}\hfill{\scriptsize (function)}}\\


 is an internal function used by the subgroup presentation functions described
in \ref{Subgroup Presentations}. It applies the Reduced Reidemeister\texttt{\symbol{45}}Schreier method to
construct an augmented coset table for the subgroup of \mbox{\texttt{\mdseries\slshape G}} which is defined by the given coset table \mbox{\texttt{\mdseries\slshape table}}. The new subgroup generators will be named \mbox{\texttt{\mdseries\slshape string}}\texttt{1}, \mbox{\texttt{\mdseries\slshape string}}\texttt{2}, $\ldots$. }

 

\subsection{\textcolor{Chapter }{RewriteWord}}
\logpage{[ 47, 9, 4 ]}\nobreak
\hyperdef{L}{X86B65EA186140244}{}
{\noindent\textcolor{FuncColor}{$\triangleright$\enspace\texttt{RewriteWord({\mdseries\slshape aug, word})\index{RewriteWord@\texttt{RewriteWord}}
\label{RewriteWord}
}\hfill{\scriptsize (function)}}\\


 \texttt{RewriteWord} rewrites \mbox{\texttt{\mdseries\slshape word}} (which must be a word in the underlying free group with respect to which the
augmented coset table \mbox{\texttt{\mdseries\slshape aug}} is given) in the subgroup generators given by the augmented coset table \mbox{\texttt{\mdseries\slshape aug}}. It returns a Tietze\texttt{\symbol{45}}type word (i.e.{\nobreakspace}a list
of integers), referring to the primary and secondary generators of \mbox{\texttt{\mdseries\slshape aug}}. 

 If \mbox{\texttt{\mdseries\slshape word}} is not contained in the subgroup, \texttt{fail} is returned. }

 }

  
\section{\textcolor{Chapter }{Low Index Subgroups}}\label{Low Index Subgroups}
\logpage{[ 47, 10, 0 ]}
\hyperdef{L}{X87FBDA2B815A8776}{}
{
  

\subsection{\textcolor{Chapter }{LowIndexSubgroupsFpGroupIterator}}
\logpage{[ 47, 10, 1 ]}\nobreak
\hyperdef{L}{X85C5151380E19122}{}
{\noindent\textcolor{FuncColor}{$\triangleright$\enspace\texttt{LowIndexSubgroupsFpGroupIterator({\mdseries\slshape G[, H], index[, excluded]})\index{LowIndexSubgroupsFpGroupIterator@\texttt{LowIndexSubgroupsFpGroupIterator}}
\label{LowIndexSubgroupsFpGroupIterator}
}\hfill{\scriptsize (operation)}}\\
\noindent\textcolor{FuncColor}{$\triangleright$\enspace\texttt{LowIndexSubgroupsFpGroup({\mdseries\slshape G[, H], index[, excluded]})\index{LowIndexSubgroupsFpGroup@\texttt{LowIndexSubgroupsFpGroup}}
\label{LowIndexSubgroupsFpGroup}
}\hfill{\scriptsize (operation)}}\\


 \index{iterator!for low index subgroups} These functions compute representatives of the conjugacy classes of subgroups
of the finitely presented group \mbox{\texttt{\mdseries\slshape G}} that contain the subgroup \mbox{\texttt{\mdseries\slshape H}} of \mbox{\texttt{\mdseries\slshape G}} and that have index less than or equal to \mbox{\texttt{\mdseries\slshape index}}. 

 \texttt{LowIndexSubgroupsFpGroupIterator} returns an iterator (see{\nobreakspace}\ref{Iterators}) that can be used to run over these subgroups, and \texttt{LowIndexSubgroupsFpGroup} returns the list of these subgroups. If one is interested only in one or a few
subgroups up to a given index then preferably the iterator should be used. 

 If the optional argument \mbox{\texttt{\mdseries\slshape excluded}} has been specified, then it is expected to be a list of words in the free
generators of the underlying free group of \mbox{\texttt{\mdseries\slshape G}}, and \texttt{LowIndexSubgroupsFpGroup} returns only those subgroups of index at most \mbox{\texttt{\mdseries\slshape index}} that contain \mbox{\texttt{\mdseries\slshape H}}, but do not contain any conjugate of any of the group elements defined by
these words. 

 If not given, \mbox{\texttt{\mdseries\slshape H}} defaults to the trivial subgroup. 

 The algorithm used finds the requested subgroups by systematically running
through a tree of all potential coset tables of \mbox{\texttt{\mdseries\slshape G}} of length at most \mbox{\texttt{\mdseries\slshape index}} (where it skips all branches of that tree for which it knows in advance that
they cannot provide new classes of such subgroups). The time required to do
this depends, of course, on the presentation of \mbox{\texttt{\mdseries\slshape G}}, but in general it will grow exponentially with the value of \mbox{\texttt{\mdseries\slshape index}}. So you should be careful with the choice of \mbox{\texttt{\mdseries\slshape index}}. 
\begin{Verbatim}[commandchars=!@|,fontsize=\small,frame=single,label=Example]
  !gapprompt@gap>| !gapinput@li:=LowIndexSubgroupsFpGroup( g, TrivialSubgroup( g ), 10 );|
  [ Group(<fp, no generators known>), Group(<fp, no generators known>),
    Group(<fp, no generators known>), Group(<fp, no generators known>) ]
\end{Verbatim}
 

 By default, the algorithm computes no generating sets for the subgroups. This
can be enforced with \texttt{GeneratorsOfGroup} (\ref{GeneratorsOfGroup}): 
\begin{Verbatim}[commandchars=!@|,fontsize=\small,frame=single,label=Example]
  !gapprompt@gap>| !gapinput@GeneratorsOfGroup(li[2]);|
  [ a, b*a*b^-1 ]
\end{Verbatim}
 

 If we are interested just in one (proper) subgroup of index at most $10$, we can use the function that returns an iterator. The first subgroup found
is the group itself, except if a list of excluded elements is entered (see
below), so we look at the second subgroup. 

 
\begin{Verbatim}[commandchars=!@|,fontsize=\small,frame=single,label=Example]
  !gapprompt@gap>| !gapinput@iter:= LowIndexSubgroupsFpGroupIterator( g, 10 );;|
  !gapprompt@gap>| !gapinput@s1:= NextIterator( iter );;  Index( g, s1 );|
  1
  !gapprompt@gap>| !gapinput@IsDoneIterator( iter );|
  false
  !gapprompt@gap>| !gapinput@s2:= NextIterator( iter );;  s2 = li[2];|
  true
\end{Verbatim}
 

 As an example for an application of the optional parameter \mbox{\texttt{\mdseries\slshape excluded}}, we compute all conjugacy classes of torsion free subgroups of index at most $24$ in the group $G = \langle x,y,z \mid x^2, y^4, z^3, (xy)^3, (yz)^2, (xz)^3 \rangle$. It is know from theory that each torsion element of this group is conjugate
to a power of $x$, $y$, $z$, $xy$, $xz$, or $yz$. (Note that this includes conjugates of $y^2$.) 

 
\begin{Verbatim}[commandchars=!@|,fontsize=\small,frame=single,label=Example]
  !gapprompt@gap>| !gapinput@F := FreeGroup( "x", "y", "z" );;|
  !gapprompt@gap>| !gapinput@x := F.1;; y := F.2;; z := F.3;;|
  !gapprompt@gap>| !gapinput@G := F / [ x^2, y^4, z^3, (x*y)^3, (y*z)^2, (x*z)^3 ];;|
  !gapprompt@gap>| !gapinput@torsion := [ x, y, y^2, z, x*y, x*z, y*z ];;|
  !gapprompt@gap>| !gapinput@SetInfoLevel( InfoFpGroup, 2 );|
  !gapprompt@gap>| !gapinput@lis := LowIndexSubgroupsFpGroup(G, TrivialSubgroup(G), 24, torsion);;|
  #I  LowIndexSubgroupsFpGroup called
  #I   class 1 of index 24 and length 8
  #I   class 2 of index 24 and length 24
  #I   class 3 of index 24 and length 24
  #I   class 4 of index 24 and length 24
  #I   class 5 of index 24 and length 24
  #I  LowIndexSubgroupsFpGroup done. Found 5 classes
  !gapprompt@gap>| !gapinput@SetInfoLevel( InfoFpGroup, 0 );|
\end{Verbatim}
 

 If a particular image group is desired, the operation \texttt{GQuotients} (\ref{GQuotients}) (see{\nobreakspace}\ref{Quotient Methods}) can be useful as well. }

 }

  
\section{\textcolor{Chapter }{Converting Groups to Finitely Presented Groups}}\label{Converting Groups to Finitely Presented Groups}
\logpage{[ 47, 11, 0 ]}
\hyperdef{L}{X81003D217D92E342}{}
{
  

\subsection{\textcolor{Chapter }{IsomorphismFpGroup}}
\logpage{[ 47, 11, 1 ]}\nobreak
\hyperdef{L}{X7F28268F850F454E}{}
{\noindent\textcolor{FuncColor}{$\triangleright$\enspace\texttt{IsomorphismFpGroup({\mdseries\slshape G})\index{IsomorphismFpGroup@\texttt{IsomorphismFpGroup}}
\label{IsomorphismFpGroup}
}\hfill{\scriptsize (attribute)}}\\


 returns an isomorphism from the given finite group \mbox{\texttt{\mdseries\slshape G}} to a finitely presented group isomorphic to \mbox{\texttt{\mdseries\slshape G}}. The function first \emph{chooses a set of generators of \mbox{\texttt{\mdseries\slshape G}}} and then computes a presentation in terms of these generators. 
\begin{Verbatim}[commandchars=!@|,fontsize=\small,frame=single,label=Example]
  !gapprompt@gap>| !gapinput@g := Group( (2,3,4,5), (1,2,5) );;|
  !gapprompt@gap>| !gapinput@iso := IsomorphismFpGroup( g );|
  [ (1,2), (2,3), (3,4), (4,5) ] -> [ F1, F2, F3, F4 ]
  !gapprompt@gap>| !gapinput@fp := Image( iso );|
  <fp group of size 120 on the generators [ F1, F2, F3, F4 ]>
  !gapprompt@gap>| !gapinput@RelatorsOfFpGroup( fp );|
  [ F1^2, F2^2, F3^2, F4^2, (F1*F2)^3, (F1*F3)^2, (F1*F4)^2, (F2*F3)^3,
    (F2*F4)^2, (F3*F4)^3 ]
\end{Verbatim}
 }

 

\subsection{\textcolor{Chapter }{IsomorphismFpGroupByGenerators}}
\logpage{[ 47, 11, 2 ]}\nobreak
\hyperdef{L}{X81B2B3B6812FD62D}{}
{\noindent\textcolor{FuncColor}{$\triangleright$\enspace\texttt{IsomorphismFpGroupByGenerators({\mdseries\slshape G, gens[, string]})\index{IsomorphismFpGroupByGenerators@\texttt{IsomorphismFpGroupByGenerators}}
\label{IsomorphismFpGroupByGenerators}
}\hfill{\scriptsize (function)}}\\
\noindent\textcolor{FuncColor}{$\triangleright$\enspace\texttt{IsomorphismFpGroupByGeneratorsNC({\mdseries\slshape G, gens, string})\index{IsomorphismFpGroupByGeneratorsNC@\texttt{IsomorphismFpGroupByGeneratorsNC}}
\label{IsomorphismFpGroupByGeneratorsNC}
}\hfill{\scriptsize (operation)}}\\


 returns an isomorphism from a finite group \mbox{\texttt{\mdseries\slshape G}} to a finitely presented group \mbox{\texttt{\mdseries\slshape F}} isomorphic to \mbox{\texttt{\mdseries\slshape G}}. The generators of \mbox{\texttt{\mdseries\slshape F}} correspond to the \emph{generators of \mbox{\texttt{\mdseries\slshape G}} given in the list \mbox{\texttt{\mdseries\slshape gens}}}. If \mbox{\texttt{\mdseries\slshape string}} is given it is used to name the generators of the finitely presented group. 

 The \texttt{NC} version will avoid testing whether the elements in \mbox{\texttt{\mdseries\slshape gens}} generate \mbox{\texttt{\mdseries\slshape G}}. 
\begin{Verbatim}[commandchars=!@|,fontsize=\small,frame=single,label=Example]
  !gapprompt@gap>| !gapinput@SetInfoLevel( InfoFpGroup, 1 );|
  !gapprompt@gap>| !gapinput@iso := IsomorphismFpGroupByGenerators( g, [ (1,2), (1,2,3,4,5) ] );|
  #I  the image group has 2 gens and 5 rels of total length 39
  [ (1,2), (1,2,3,4,5) ] -> [ F1, F2 ]
  !gapprompt@gap>| !gapinput@fp := Image( iso );|
  <fp group of size 120 on the generators [ F1, F2 ]>
\end{Verbatim}
 

 The main task of the function \texttt{IsomorphismFpGroupByGenerators} is to find a presentation of \mbox{\texttt{\mdseries\slshape G}} in the provided generators \mbox{\texttt{\mdseries\slshape gens}}. In the case of a permutation group \mbox{\texttt{\mdseries\slshape G}} it does this by first constructing a stabilizer chain of \mbox{\texttt{\mdseries\slshape G}} and then it works through that chain from the bottom to the top, recursively
computing a presentation for each of the involved stabilizers. The method used
is essentially an implementation of John Cannon's
multi\texttt{\symbol{45}}stage relations\texttt{\symbol{45}}finding algorithm
as described in \cite{Neu82} (see also \cite{Can73} for a more graph theoretical description). Moreover, it makes heavy use of
Tietze transformations in each stage to avoid an explosion of the total length
of the relators. 

 Note that because of the random methods involved in the construction of the
stabilizer chain the resulting presentations of \mbox{\texttt{\mdseries\slshape G}} will in general be different for repeated calls with the same arguments. 

 
\begin{Verbatim}[commandchars=!@|,fontsize=\small,frame=single,label=Example]
  !gapprompt@gap>| !gapinput@M12 := MathieuGroup( 12 );|
  Group([ (1,2,3,4,5,6,7,8,9,10,11), (3,7,11,8)(4,10,5,6),
    (1,12)(2,11)(3,6)(4,8)(5,9)(7,10) ])
  !gapprompt@gap>| !gapinput@gens := GeneratorsOfGroup( M12 );;|
  !gapprompt@gap>| !gapinput@iso := IsomorphismFpGroupByGenerators( M12, gens );;|
  #I  the image group has 3 gens and 24 rels of total length 669
  !gapprompt@gap>| !gapinput@iso := IsomorphismFpGroupByGenerators( M12, gens );;|
  #I  the image group has 3 gens and 20 rels of total length 414
\end{Verbatim}
 

 Also in the case of a permutation group \mbox{\texttt{\mdseries\slshape G}}, the function \texttt{IsomorphismFpGroupByGenerators} supports the option \texttt{method} that can be used to modify the strategy. The option \texttt{method} may take the following values. 

 
\begin{description}
\item[{\texttt{method := "regular"}}]  This may be specified for groups of small size, up to $10^5$ say. It implies that the function first constructs a regular representation \mbox{\texttt{\mdseries\slshape R}} of \mbox{\texttt{\mdseries\slshape G}} and then a presentation of \mbox{\texttt{\mdseries\slshape R}}. In general, this presentation will be much more concise than the default
one, but the price is the time needed for the construction of \mbox{\texttt{\mdseries\slshape R}}. 
\item[{\texttt{method := [ "regular", bound ]}}]  This is a refinement of the previous possibility. In this case, \texttt{bound} should be an integer, and if so the method \texttt{"regular"} as described above is applied to the largest stabilizer in the stabilizer
chain of \mbox{\texttt{\mdseries\slshape G}} whose size does not exceed the given bound and then the
multi\texttt{\symbol{45}}stage algorithm is used to work through the chain
from that subgroup to the top. 
\item[{\texttt{method := "fast"}}]  This chooses an alternative method which essentially is a kind of
multi\texttt{\symbol{45}}stage algorithm for a stabilizer chain of \mbox{\texttt{\mdseries\slshape G}} but does not make any attempt do reduce the number of relators as it is done
in Cannon's algorithm or to reduce their total length. Hence it is often much
faster than the default method, but the total length of the resulting
presentation may be huge. 
\item[{\texttt{method := "default"}}]  This simply means that the default method shall be used, which is the case if
the option \texttt{method} is not given a value. 
\end{description}
 

 
\begin{Verbatim}[commandchars=!@|,fontsize=\small,frame=single,label=Example]
  !gapprompt@gap>| !gapinput@iso := IsomorphismFpGroupByGenerators( M12, gens :|
  !gapprompt@>| !gapinput@                                          method := "regular" );;|
  #I  the image group has 3 gens and 11 rels of total length 92
  !gapprompt@gap>| !gapinput@iso := IsomorphismFpGroupByGenerators( M12, gens :|
  !gapprompt@>| !gapinput@                                          method := "fast" );;|
  #I  the image group has 3 gens and 136 rels of total length 3170
\end{Verbatim}
 

 Though the option \texttt{method := "regular"} is only checked in the case of a permutation group it also affects the
performance and the results of the function \texttt{IsomorphismFpGroupByGenerators} for other groups, e. g. for matrix groups. This happens because, for these
groups, the function first calls the function \texttt{NiceMonomorphism} (\ref{NiceMonomorphism}) to get a bijective action homomorphism from \mbox{\texttt{\mdseries\slshape G}} to a suitable permutation group, $P$ say, and then, recursively, calls itself for the group $P$ so that now the option becomes relevant. 

 
\begin{Verbatim}[commandchars=!@|,fontsize=\small,frame=single,label=Example]
  !gapprompt@gap>| !gapinput@G := ImfMatrixGroup( 5, 1, 3 );|
  ImfMatrixGroup(5,1,3)
  !gapprompt@gap>| !gapinput@gens := GeneratorsOfGroup( G );|
  [ [ [ -1, 0, 0, 0, 0 ], [ 0, 1, 0, 0, 0 ], [ 0, 0, 0, 1, 0 ],
        [ -1, -1, -1, -1, 2 ], [ -1, 0, 0, 0, 1 ] ],
    [ [ 0, 1, 0, 0, 0 ], [ 0, 0, 1, 0, 0 ], [ 0, 0, 0, 1, 0 ],
        [ 1, 0, 0, 0, 0 ], [ 0, 0, 0, 0, 1 ] ] ]
  !gapprompt@gap>| !gapinput@iso := IsomorphismFpGroupByGenerators( G, gens );;|
  #I  the image group has 2 gens and 11 rels of total length 120
  !gapprompt@gap>| !gapinput@iso := IsomorphismFpGroupByGenerators( G, gens :|
  !gapprompt@>| !gapinput@                                          method := "regular");;|
  #I  the image group has 2 gens and 6 rels of total length 56
  !gapprompt@gap>| !gapinput@SetInfoLevel( InfoFpGroup, 0 );|
  !gapprompt@gap>| !gapinput@iso;|
  <composed isomorphism:[ [ [ -1, 0, 0, 0, 0 ], [ 0, 1, 0, 0, 0 ], [ 0, \
  0, 0, 1, 0 ], [ -1, -1, -1, -1, 2 ], [ -1, 0, 0, 0, 1 ] ], [ [ 0, 1, 0\
  , 0, 0 ], [ 0, 0, 1, 0, 0 ], [ 0, 0, 0, 1, 0 ], [ 1, 0, 0, 0, 0 ], [ 0\
  , 0, 0, 0, 1 ] ] ]->[ F1, F2 ]>
  !gapprompt@gap>| !gapinput@ConstituentsCompositionMapping(iso);|
  [ <action isomorphism>,
    [ (2,3,4)(5,6)(8,9,10), (1,2,3,5)(6,7,8,9) ] -> [ F1, F2 ] ]
\end{Verbatim}
 

 Since \textsf{GAP} cannot decompose elements of a matrix group into generators, the resulting
isomorphism is stored as a composition of a (faithful) permutation action on
vectors and a homomorphism from the permutation image to the finitely
presented group. In such a situation the constituent mappings can be obtained
via \texttt{ConstituentsCompositionMapping} (\ref{ConstituentsCompositionMapping}) as separate \textsf{GAP} objects. }

 }

  
\section{\textcolor{Chapter }{New Presentations and Presentations for Subgroups}}\label{New Presentations and Presentations for Subgroups}
\logpage{[ 47, 12, 0 ]}
\hyperdef{L}{X826604AA7F18BFA3}{}
{
  \index{IsomorphismFpGroup@\texttt{IsomorphismFpGroup}!for subgroups of fp groups} \texttt{IsomorphismFpGroup} (\ref{IsomorphismFpGroup}) is also used to compute a new finitely presented group that is isomorphic to
the given subgroup of a finitely presented group. (This is typically the only
method to compute with subgroups of a finitely presented group.) 

 
\begin{Verbatim}[commandchars=!@A,fontsize=\small,frame=single,label=Example]
  !gapprompt@gap>A !gapinput@f:=FreeGroup(2);;A
  !gapprompt@gap>A !gapinput@g:=f/[f.1^2,f.2^3,(f.1*f.2)^5];A
  <fp group on the generators [ f1, f2 ]>
  !gapprompt@gap>A !gapinput@u:=Subgroup(g,[g.1*g.2]);A
  Group([ f1*f2 ])
  !gapprompt@gap>A !gapinput@hom:=IsomorphismFpGroup(u);A
  [ <[ [ 1, 1 ] ]|f2^-1*f1^-1> ] -> [ F1 ]
  !gapprompt@gap>A !gapinput@new:=Range(hom);A
  <fp group on the generators [ F1 ]>
  !gapprompt@gap>A !gapinput@List(GeneratorsOfGroup(new),i->PreImagesRepresentative(hom,i));A
  [ <[ [ 1, 1 ] ]|f2^-1*f1^-1> ]
\end{Verbatim}
 

 When working with such homomorphisms, some subgroup elements are expressed as
extremely long words in the group generators. Therefore the underlying words
of subgroup generators stored in the isomorphism (as obtained by \texttt{MappingGeneratorsImages} (\ref{MappingGeneratorsImages}) and displayed when \texttt{View} (\ref{View})ing the homomorphism) as well as preimages under the homomorphism are stored
in the form of straight line program elements (see{\nobreakspace}\ref{Straight Line Program Elements}). These will behave like ordinary words and no extra treatment should be
necessary. 

 
\begin{Verbatim}[commandchars=!@A,fontsize=\small,frame=single,label=Example]
  !gapprompt@gap>A !gapinput@r:=Range(hom).1^10;A
  F1^10
  !gapprompt@gap>A !gapinput@p:=PreImagesRepresentative(hom,r);A
  <[ [ 1, 10 ] ]|(f2^-1*f1^-1)^10>
\end{Verbatim}
 If desired, it also is possible to convert these underlying words using \texttt{EvalStraightLineProgElm} (\ref{EvalStraightLineProgElm}): 
\begin{Verbatim}[commandchars=!@|,fontsize=\small,frame=single,label=Example]
  !gapprompt@gap>| !gapinput@r:=EvalStraightLineProgElm(UnderlyingElement(p));|
  (f2^-1*f1^-1)^10
  !gapprompt@gap>| !gapinput@p:=ElementOfFpGroup(FamilyObj(p),r);|
  (f2^-1*f1^-1)^10
\end{Verbatim}
 

 (If you are only interested in a finitely presented group isomorphic to the
given subgroup but not in the isomorphism, you may also use the functions \texttt{PresentationViaCosetTable} (\ref{PresentationViaCosetTable}) and \texttt{FpGroupPresentation} (\ref{FpGroupPresentation}) (see \ref{Creating Presentations}).) 

 Homomorphisms can also be used to obtain an isomorphic finitely presented
group with a (hopefully) simpler presentation. 

 

\subsection{\textcolor{Chapter }{IsomorphismSimplifiedFpGroup}}
\logpage{[ 47, 12, 1 ]}\nobreak
\hyperdef{L}{X78D87FA68233C401}{}
{\noindent\textcolor{FuncColor}{$\triangleright$\enspace\texttt{IsomorphismSimplifiedFpGroup({\mdseries\slshape G})\index{IsomorphismSimplifiedFpGroup@\texttt{IsomorphismSimplifiedFpGroup}}
\label{IsomorphismSimplifiedFpGroup}
}\hfill{\scriptsize (attribute)}}\\


 applies Tietze transformations to a copy of the presentation of the given
finitely presented group \mbox{\texttt{\mdseries\slshape G}} in order to reduce it with respect to the number of generators, the number of
relators, and the relator lengths. 

 The operation returns an isomorphism with source \mbox{\texttt{\mdseries\slshape G}}, range a group \mbox{\texttt{\mdseries\slshape H}} isomorphic to \mbox{\texttt{\mdseries\slshape G}}, so that the presentation of \mbox{\texttt{\mdseries\slshape H}} has been simplified using Tietze transformations. 
\begin{Verbatim}[commandchars=!@|,fontsize=\small,frame=single,label=Example]
  !gapprompt@gap>| !gapinput@f:=FreeGroup(3);;|
  !gapprompt@gap>| !gapinput@g:=f/[f.1^2,f.2^3,(f.1*f.2)^5,f.1/f.3];|
  <fp group on the generators [ f1, f2, f3 ]>
  !gapprompt@gap>| !gapinput@hom:=IsomorphismSimplifiedFpGroup(g);|
  [ f1, f2, f3 ] -> [ f1, f2, f1 ]
  !gapprompt@gap>| !gapinput@Range(hom);|
  <fp group on the generators [ f1, f2 ]>
  !gapprompt@gap>| !gapinput@RelatorsOfFpGroup(Range(hom));|
  [ f1^2, f2^3, (f1*f2)^5 ]
  !gapprompt@gap>| !gapinput@RelatorsOfFpGroup(g);|
  [ f1^2, f2^3, (f1*f2)^5, f1*f3^-1 ]
\end{Verbatim}
 

 \texttt{IsomorphismSimplifiedFpGroup} uses Tietze transformations to simplify the presentation, see \ref{SimplifiedFpGroup}. }

 }

  
\section{\textcolor{Chapter }{Preimages under Homomorphisms from an FpGroup}}\label{Preimages under Homomorphisms from an FpGroup}
\logpage{[ 47, 13, 0 ]}
\hyperdef{L}{X86E7CE077D82133D}{}
{
  For some subgroups of a finitely presented group the number of subgroup
generators increases with the index of the subgroup. However often these
generators are not needed at all for further calculations, but what is needed
is the action of the cosets of the subgroup. This gives the image of the
subgroup in a finite quotient and this finite quotient can be used to
calculate normalizers, closures, intersections and so forth{\nobreakspace}\cite{HulpkeQuot}. 

 The same applies for subgroups that are obtained as preimages under
homomorphisms. 

\subsection{\textcolor{Chapter }{SubgroupOfWholeGroupByQuotientSubgroup}}
\logpage{[ 47, 13, 1 ]}\nobreak
\hyperdef{L}{X7ABC3C917D41A74B}{}
{\noindent\textcolor{FuncColor}{$\triangleright$\enspace\texttt{SubgroupOfWholeGroupByQuotientSubgroup({\mdseries\slshape fpfam, Q, U})\index{SubgroupOfWholeGroupByQuotientSubgroup@\texttt{Subgroup}\-\texttt{Of}\-\texttt{Whole}\-\texttt{Group}\-\texttt{By}\-\texttt{Quotient}\-\texttt{Subgroup}}
\label{SubgroupOfWholeGroupByQuotientSubgroup}
}\hfill{\scriptsize (function)}}\\


 takes a FpGroup family \mbox{\texttt{\mdseries\slshape fpfam}}, a finitely generated group \mbox{\texttt{\mdseries\slshape Q}} such that the fp generators of \mbox{\texttt{\mdseries\slshape fpfam}} can be mapped by an epimorphism $phi$ onto the \texttt{GeneratorsOfGroup} (\ref{GeneratorsOfGroup}) value of \mbox{\texttt{\mdseries\slshape Q}}, and a subgroup \mbox{\texttt{\mdseries\slshape U}} of \mbox{\texttt{\mdseries\slshape Q}}. It returns the subgroup of \mbox{\texttt{\mdseries\slshape fpfam}}\texttt{!.wholeGroup} which is the full preimage of \mbox{\texttt{\mdseries\slshape U}} under $phi$. }

 

\subsection{\textcolor{Chapter }{IsSubgroupOfWholeGroupByQuotientRep}}
\logpage{[ 47, 13, 2 ]}\nobreak
\hyperdef{L}{X8047D7A37B27FEEA}{}
{\noindent\textcolor{FuncColor}{$\triangleright$\enspace\texttt{IsSubgroupOfWholeGroupByQuotientRep({\mdseries\slshape G})\index{IsSubgroupOfWholeGroupByQuotientRep@\texttt{IsSubgroupOfWholeGroupByQuotientRep}}
\label{IsSubgroupOfWholeGroupByQuotientRep}
}\hfill{\scriptsize (Representation)}}\\


 is the representation for subgroups of an FpGroup, given by a quotient
subgroup. The components \mbox{\texttt{\mdseries\slshape G}}\texttt{!.quot} and \mbox{\texttt{\mdseries\slshape G}}\texttt{!.sub} hold quotient, respectively subgroup. }

 

\subsection{\textcolor{Chapter }{AsSubgroupOfWholeGroupByQuotient}}
\logpage{[ 47, 13, 3 ]}\nobreak
\hyperdef{L}{X84E6CEA28611C112}{}
{\noindent\textcolor{FuncColor}{$\triangleright$\enspace\texttt{AsSubgroupOfWholeGroupByQuotient({\mdseries\slshape U})\index{AsSubgroupOfWholeGroupByQuotient@\texttt{AsSubgroupOfWholeGroupByQuotient}}
\label{AsSubgroupOfWholeGroupByQuotient}
}\hfill{\scriptsize (attribute)}}\\


 returns the same subgroup in the representation \texttt{AsSubgroupOfWholeGroupByQuotient}. 

 See also \texttt{SubgroupOfWholeGroupByCosetTable} (\ref{SubgroupOfWholeGroupByCosetTable}) and \texttt{CosetTableBySubgroup} (\ref{CosetTableBySubgroup}). 

 This technique is used by \textsf{GAP} for example to represent the derived subgroup, which is obtained from the
quotient $G/G'$. 
\begin{Verbatim}[commandchars=!@|,fontsize=\small,frame=single,label=Example]
  !gapprompt@gap>| !gapinput@f:=FreeGroup(2);;g:=f/[f.1^6,f.2^6,(f.1*f.2)^6];;|
  !gapprompt@gap>| !gapinput@d:=DerivedSubgroup(g);|
  Group(<fp, no generators known>)
  !gapprompt@gap>| !gapinput@Index(g,d);|
  36
\end{Verbatim}
 }

 

\subsection{\textcolor{Chapter }{DefiningQuotientHomomorphism}}
\logpage{[ 47, 13, 4 ]}\nobreak
\hyperdef{L}{X7DA1151D84289FC9}{}
{\noindent\textcolor{FuncColor}{$\triangleright$\enspace\texttt{DefiningQuotientHomomorphism({\mdseries\slshape U})\index{DefiningQuotientHomomorphism@\texttt{DefiningQuotientHomomorphism}}
\label{DefiningQuotientHomomorphism}
}\hfill{\scriptsize (function)}}\\


 if \mbox{\texttt{\mdseries\slshape U}} is a subgroup in quotient representation (\texttt{IsSubgroupOfWholeGroupByQuotientRep} (\ref{IsSubgroupOfWholeGroupByQuotientRep})), this function returns the defining homomorphism from the whole group to \mbox{\texttt{\mdseries\slshape U}}\texttt{!.quot}. }

 }

  
\section{\textcolor{Chapter }{Quotient Methods}}\label{Quotient Methods}
\logpage{[ 47, 14, 0 ]}
\hyperdef{L}{X846072F779B51087}{}
{
  An important class of algorithms for finitely presented groups are the \emph{quotient algorithms} which compute quotient groups of a given finitely presented group. There are
algorithms for epimorphisms onto abelian groups, $p$\texttt{\symbol{45}}groups and solvable groups. (The ``low index'' algorithm {\textendash}\texttt{LowIndexSubgroupsFpGroup} (\ref{LowIndexSubgroupsFpGroup}){\textendash} can be considered as well as an algorithm that produces
permutation group quotients.) 

 \texttt{MaximalAbelianQuotient} (\ref{MaximalAbelianQuotient}), as defined for general groups, returns the largest abelian quotient of the
given group. 
\begin{Verbatim}[commandchars=!@|,fontsize=\small,frame=single,label=Example]
  !gapprompt@gap>| !gapinput@f:=FreeGroup(2);;fp:=f/[f.1^6,f.2^6,(f.1*f.2)^12];|
  <fp group on the generators [ f1, f2 ]>
  !gapprompt@gap>| !gapinput@hom:=MaximalAbelianQuotient(fp);|
  [ f1, f2 ] -> [ f1, f3 ]
  !gapprompt@gap>| !gapinput@Size(Image(hom));|
  36
\end{Verbatim}
 

\subsection{\textcolor{Chapter }{PQuotient}}
\logpage{[ 47, 14, 1 ]}\nobreak
\hyperdef{L}{X7B5DDADC80F5796B}{}
{\noindent\textcolor{FuncColor}{$\triangleright$\enspace\texttt{PQuotient({\mdseries\slshape F, p[, c][, logord][, ctype]})\index{PQuotient@\texttt{PQuotient}}
\label{PQuotient}
}\hfill{\scriptsize (function)}}\\


 computes a factor \mbox{\texttt{\mdseries\slshape p}}\texttt{\symbol{45}}group of a finitely presented group \mbox{\texttt{\mdseries\slshape F}} in form of a quotient system. The quotient system can be converted into an
epimorphism from \mbox{\texttt{\mdseries\slshape F}} onto the \mbox{\texttt{\mdseries\slshape p}}\texttt{\symbol{45}}group computed by the function \texttt{EpimorphismQuotientSystem} (\ref{EpimorphismQuotientSystem}). 

 For a group $G$ define the exponent\texttt{\symbol{45}}$p$ central series of $G$ inductively by ${\cal P}_1(G) = G$ and ${\cal P}_{{i+1}}(G) = [{\cal P}_i(G),G]{\cal P}_{{i+1}}(G)^p$. The factor groups modulo the terms of the lower exponent\texttt{\symbol{45}}$p$ central series are $p$\texttt{\symbol{45}}groups. The group $G$ has $p$\texttt{\symbol{45}}class $c$ if ${\cal P}_c(G) \neq {\cal P}_{{c+1}}(G) = 1$. 

 The algorithm computes successive quotients modulo the terms of the
exponent\texttt{\symbol{45}}$p$ central series of \mbox{\texttt{\mdseries\slshape F}}. If the parameter \mbox{\texttt{\mdseries\slshape c}} is present, then the factor group modulo the $(c+1)$\texttt{\symbol{45}}th term of the exponent\texttt{\symbol{45}}$p$ central series of \mbox{\texttt{\mdseries\slshape F}} is returned. If \mbox{\texttt{\mdseries\slshape c}} is not present, then the algorithm attempts to compute the largest factor \mbox{\texttt{\mdseries\slshape p}}\texttt{\symbol{45}}group of \mbox{\texttt{\mdseries\slshape F}}. In case \mbox{\texttt{\mdseries\slshape F}} does not have a largest factor \mbox{\texttt{\mdseries\slshape p}}\texttt{\symbol{45}}group, the algorithm will not terminate. 

 By default the algorithm computes only with factor groups of order at most $p^{256}$. If the parameter \mbox{\texttt{\mdseries\slshape logord}} is present, it will compute with factor groups of order at most $p^{\mbox{\texttt{\mdseries\slshape logord}}}$. If this parameter is specified, then the parameter \mbox{\texttt{\mdseries\slshape c}} must also be given. The present implementation produces an error message if
the order of a $p$\texttt{\symbol{45}}quotient exceeds $p^{256}$ or $p^{\mbox{\texttt{\mdseries\slshape logord}}}$, respectively. Note that the order of intermediate $p$\texttt{\symbol{45}}groups may be larger than the final order of a $p$\texttt{\symbol{45}}quotient. 

 The parameter \mbox{\texttt{\mdseries\slshape ctype}} determines the type of collector that is used for computations within the
factor \mbox{\texttt{\mdseries\slshape p}}\texttt{\symbol{45}}group. \mbox{\texttt{\mdseries\slshape ctype}} must either be \texttt{"single"} in which case a simple collector from the left is used or \texttt{"combinatorial"} in which case a combinatorial collector from the left is used. }

 

\subsection{\textcolor{Chapter }{EpimorphismQuotientSystem}}
\logpage{[ 47, 14, 2 ]}\nobreak
\hyperdef{L}{X86EB30A7867EEF16}{}
{\noindent\textcolor{FuncColor}{$\triangleright$\enspace\texttt{EpimorphismQuotientSystem({\mdseries\slshape quotsys})\index{EpimorphismQuotientSystem@\texttt{EpimorphismQuotientSystem}}
\label{EpimorphismQuotientSystem}
}\hfill{\scriptsize (operation)}}\\


 For a quotient system \mbox{\texttt{\mdseries\slshape quotsys}} obtained from the function \texttt{PQuotient} (\ref{PQuotient}), this operation returns an epimorphism $\mbox{\texttt{\mdseries\slshape F}} \rightarrow \mbox{\texttt{\mdseries\slshape P}}$ where $\mbox{\texttt{\mdseries\slshape F}}$ is the finitely presented group of which \mbox{\texttt{\mdseries\slshape quotsys}} is a quotient system and $\mbox{\texttt{\mdseries\slshape P}}$ is a pc group isomorphic to the quotient of \mbox{\texttt{\mdseries\slshape F}} determined by \mbox{\texttt{\mdseries\slshape quotsys}}. 

 Different calls to this operation will create different groups \mbox{\texttt{\mdseries\slshape P}}, each with its own family. 

 
\begin{Verbatim}[commandchars=!@|,fontsize=\small,frame=single,label=Example]
  !gapprompt@gap>| !gapinput@PQuotient( FreeGroup(2), 5, 10, 1024, "combinatorial" );|
  <5-quotient system of 5-class 10 with 520 generators>
  !gapprompt@gap>| !gapinput@phi := EpimorphismQuotientSystem( last );|
  [ f1, f2 ] -> [ a1, a2 ]
  !gapprompt@gap>| !gapinput@Collected( Factors( Size( Image( phi ) ) ) );|
  [ [ 5, 520 ] ]
\end{Verbatim}
 }

 

\subsection{\textcolor{Chapter }{EpimorphismPGroup}}
\logpage{[ 47, 14, 3 ]}\nobreak
\hyperdef{L}{X7CA738DB80B20D67}{}
{\noindent\textcolor{FuncColor}{$\triangleright$\enspace\texttt{EpimorphismPGroup({\mdseries\slshape fpgrp, p[, cl]})\index{EpimorphismPGroup@\texttt{EpimorphismPGroup}}
\label{EpimorphismPGroup}
}\hfill{\scriptsize (operation)}}\\


 computes an epimorphism from the finitely presented group \mbox{\texttt{\mdseries\slshape fpgrp}} to the largest $p$\texttt{\symbol{45}}group of $p$\texttt{\symbol{45}}class \mbox{\texttt{\mdseries\slshape cl}} which is a quotient of \mbox{\texttt{\mdseries\slshape fpgrp}}. If \mbox{\texttt{\mdseries\slshape cl}} is omitted, the largest finite $p$\texttt{\symbol{45}}group quotient (of $p$\texttt{\symbol{45}}class up to $1000$) is determined. 

 
\begin{Verbatim}[commandchars=!@|,fontsize=\small,frame=single,label=Example]
  !gapprompt@gap>| !gapinput@hom:=EpimorphismPGroup(fp,2);|
  [ f1, f2 ] -> [ a1, a2 ]
  !gapprompt@gap>| !gapinput@Size(Image(hom));|
  8
  !gapprompt@gap>| !gapinput@hom:=EpimorphismPGroup(fp,3,7);|
  [ f1, f2 ] -> [ a1, a2 ]
  !gapprompt@gap>| !gapinput@Size(Image(hom));|
  6561
\end{Verbatim}
 }

 

\subsection{\textcolor{Chapter }{EpimorphismNilpotentQuotient}}
\logpage{[ 47, 14, 4 ]}\nobreak
\hyperdef{L}{X7CA20E2582DC45FD}{}
{\noindent\textcolor{FuncColor}{$\triangleright$\enspace\texttt{EpimorphismNilpotentQuotient({\mdseries\slshape fpgrp[, n]})\index{EpimorphismNilpotentQuotient@\texttt{EpimorphismNilpotentQuotient}}
\label{EpimorphismNilpotentQuotient}
}\hfill{\scriptsize (function)}}\\


 returns an epimorphism on the class \mbox{\texttt{\mdseries\slshape n}} finite nilpotent quotient of the finitely presented group \mbox{\texttt{\mdseries\slshape fpgrp}}. If \mbox{\texttt{\mdseries\slshape n}} is omitted, the largest finite nilpotent quotient (of $p$\texttt{\symbol{45}}class up to $1000$) is taken. 

 
\begin{Verbatim}[commandchars=!@|,fontsize=\small,frame=single,label=Example]
  !gapprompt@gap>| !gapinput@hom:=EpimorphismNilpotentQuotient(fp,7);|
  [ f1, f2 ] -> [ f1*f4, f2*f5 ]
  !gapprompt@gap>| !gapinput@Size(Image(hom));|
  52488
\end{Verbatim}
 

 A related operation which is also applicable to finitely presented groups is \texttt{GQuotients} (\ref{GQuotients}), which computes all epimorphisms from a (finitely presented) group \mbox{\texttt{\mdseries\slshape F}} onto a given (finite) group \mbox{\texttt{\mdseries\slshape G}}. 

 
\begin{Verbatim}[commandchars=!@|,fontsize=\small,frame=single,label=Example]
  !gapprompt@gap>| !gapinput@GQuotients(fp,Group((1,2,3),(1,2)));|
  [ [ f1, f2 ] -> [ (1,2), (2,3) ], [ f1, f2 ] -> [ (2,3), (1,2,3) ],
    [ f1, f2 ] -> [ (1,2,3), (2,3) ] ]
\end{Verbatim}
 }

  

\subsection{\textcolor{Chapter }{SolvableQuotient (for a f.p. group and a size)}}
\logpage{[ 47, 14, 5 ]}\nobreak
\hyperdef{L}{X869F70CC818C946D}{}
{\noindent\textcolor{FuncColor}{$\triangleright$\enspace\texttt{SolvableQuotient({\mdseries\slshape F, size})\index{SolvableQuotient@\texttt{SolvableQuotient}!for a f.p. group and a size}
\label{SolvableQuotient:for a f.p. group and a size}
}\hfill{\scriptsize (function)}}\\
\noindent\textcolor{FuncColor}{$\triangleright$\enspace\texttt{SolvableQuotient({\mdseries\slshape F, primes})\index{SolvableQuotient@\texttt{SolvableQuotient}!for a f.p. group and a list of primes}
\label{SolvableQuotient:for a f.p. group and a list of primes}
}\hfill{\scriptsize (function)}}\\
\noindent\textcolor{FuncColor}{$\triangleright$\enspace\texttt{SolvableQuotient({\mdseries\slshape F, tuples})\index{SolvableQuotient@\texttt{SolvableQuotient}!for a f.p. group and a list of tuples}
\label{SolvableQuotient:for a f.p. group and a list of tuples}
}\hfill{\scriptsize (function)}}\\
\noindent\textcolor{FuncColor}{$\triangleright$\enspace\texttt{SQ({\mdseries\slshape F, ...})\index{SQ@\texttt{SQ}!synonym of SolvableQuotient}
\label{SQ:synonym of SolvableQuotient}
}\hfill{\scriptsize (function)}}\\


 This routine calls the solvable quotient algorithm for a finitely presented
group \mbox{\texttt{\mdseries\slshape F}}. The quotient to be found can be specified in the following ways: Specifying
an integer \mbox{\texttt{\mdseries\slshape size}} finds a quotient of size up to \mbox{\texttt{\mdseries\slshape size}} (if such large quotients exist). Specifying a list of primes in \mbox{\texttt{\mdseries\slshape primes}} finds the largest quotient involving the given primes. Finally \mbox{\texttt{\mdseries\slshape tuples}} can be used to prescribe a chief series. 

 \texttt{SQ} can be used as a synonym for \texttt{SolvableQuotient}. }

 

\subsection{\textcolor{Chapter }{EpimorphismSolvableQuotient}}
\logpage{[ 47, 14, 6 ]}\nobreak
\hyperdef{L}{X79A4D3B68110F48A}{}
{\noindent\textcolor{FuncColor}{$\triangleright$\enspace\texttt{EpimorphismSolvableQuotient({\mdseries\slshape F, param})\index{EpimorphismSolvableQuotient@\texttt{EpimorphismSolvableQuotient}}
\label{EpimorphismSolvableQuotient}
}\hfill{\scriptsize (function)}}\\


 computes an epimorphism from the finitely presented group \mbox{\texttt{\mdseries\slshape fpgrp}} to the largest solvable quotient given by \mbox{\texttt{\mdseries\slshape param}} (specified as in \texttt{SolvableQuotient} (\ref{SolvableQuotient:for a f.p. group and a size})). 

 
\begin{Verbatim}[commandchars=!@|,fontsize=\small,frame=single,label=Example]
  !gapprompt@gap>| !gapinput@f := FreeGroup( "a", "b", "c", "d" );;|
  !gapprompt@gap>| !gapinput@fp := f / [ f.1^2, f.2^2, f.3^2, f.4^2, f.1*f.2*f.1*f.2*f.1*f.2,|
  !gapprompt@>| !gapinput@ f.2*f.3*f.2*f.3*f.2*f.3*f.2*f.3, f.3*f.4*f.3*f.4*f.3*f.4,|
  !gapprompt@>| !gapinput@f.1^-1*f.3^-1*f.1*f.3, f.1^-1*f.4^-1*f.1*f.4,|
  !gapprompt@>| !gapinput@f.2^-1*f.4^-1*f.2*f.4 ];;|
  !gapprompt@gap>| !gapinput@hom:=EpimorphismSolvableQuotient(fp,300);Size(Image(hom));|
  [ a, b, c, d ] -> [ f1*f2, f1*f2, f2*f3, f2 ]
  12
  !gapprompt@gap>| !gapinput@hom:=EpimorphismSolvableQuotient(fp,[2,3]);Size(Image(hom));|
  [ a, b, c, d ] -> [ f1*f2*f4, f1*f2*f6*f8, f2*f3, f2 ]
  1152
\end{Verbatim}
 }

 

\subsection{\textcolor{Chapter }{LargerQuotientBySubgroupAbelianization}}
\logpage{[ 47, 14, 7 ]}\nobreak
\hyperdef{L}{X81167847832DD3B1}{}
{\noindent\textcolor{FuncColor}{$\triangleright$\enspace\texttt{LargerQuotientBySubgroupAbelianization({\mdseries\slshape hom, U})\index{LargerQuotientBySubgroupAbelianization@\texttt{Larger}\-\texttt{Quotient}\-\texttt{By}\-\texttt{Subgroup}\-\texttt{Abelianization}}
\label{LargerQuotientBySubgroupAbelianization}
}\hfill{\scriptsize (function)}}\\


 Let \mbox{\texttt{\mdseries\slshape hom}} a homomorphism from a finitely presented group $G$ to a finite group $H$ and $\mbox{\texttt{\mdseries\slshape U}} \leq H$. This function will {\textendash}if it exists{\textendash} return a subgroup $S \leq \mbox{\texttt{\mdseries\slshape G}}$, such that the core of $S$ is properly contained in the kernel of \mbox{\texttt{\mdseries\slshape hom}} as well as in the derived subgroup of $V$, where $V$ is the pre\texttt{\symbol{45}}image of \mbox{\texttt{\mdseries\slshape U}} under \mbox{\texttt{\mdseries\slshape hom}}. Thus $S$ exposes a larger quotient of $G$. If no such subgroup exists, \texttt{fail} is returned. 
\begin{Verbatim}[commandchars=!@|,fontsize=\small,frame=single,label=Example]
  !gapprompt@gap>| !gapinput@f:=FreeGroup("x","y","z");;|
  !gapprompt@gap>| !gapinput@g:=f/ParseRelators(f,"x^3=y^3=z^5=(xyx^2y^2)^2=(xz)^2=(yz^3)^2=1");|
  <fp group on the generators [ x, y, z ]>
  !gapprompt@gap>| !gapinput@l:=LowIndexSubgroupsFpGroup(g,6);;|
  !gapprompt@gap>| !gapinput@List(l,IndexInWholeGroup);|
  [ 1, 5, 5, 5, 5, 6, 6, 6, 6, 6, 6 ]
  !gapprompt@gap>| !gapinput@q:=DefiningQuotientHomomorphism(l[6]);;p:=Image(q);Size(p);|
  Group([ (4,5,6), (1,2,3)(4,6,5), (2,4,6,3,5) ])
  360
  !gapprompt@gap>| !gapinput@s:=LargerQuotientBySubgroupAbelianization(q,SylowSubgroup(p,3));|
  Group(<fp, no generators known>)
  !gapprompt@gap>| !gapinput@Size(Image(DefiningQuotientHomomorphism(s)));|
  193273528320
\end{Verbatim}
 }

 }

  
\section{\textcolor{Chapter }{Abelian Invariants for Subgroups}}\label{Abelian Invariants for Subgroups}
\logpage{[ 47, 15, 0 ]}
\hyperdef{L}{X81451C4B8463B848}{}
{
  Using variations of coset enumeration it is possible to compute the abelian
invariants of a subgroup of a finitely presented group without computing a
complete presentation for the subgroup in the first place. Typically, the
operation \texttt{AbelianInvariants} (\ref{AbelianInvariants}) when called for subgroups should automatically take care of this, but in case
you want to have further control about the methods used, the following
operations might be of use. 

\subsection{\textcolor{Chapter }{AbelianInvariantsSubgroupFpGroup}}
\logpage{[ 47, 15, 1 ]}\nobreak
\hyperdef{L}{X83B63ED8826F4268}{}
{\noindent\textcolor{FuncColor}{$\triangleright$\enspace\texttt{AbelianInvariantsSubgroupFpGroup({\mdseries\slshape G, H})\index{AbelianInvariantsSubgroupFpGroup@\texttt{AbelianInvariantsSubgroupFpGroup}}
\label{AbelianInvariantsSubgroupFpGroup}
}\hfill{\scriptsize (function)}}\\


 \texttt{AbelianInvariantsSubgroupFpGroup} is a synonym for \texttt{AbelianInvariantsSubgroupFpGroupRrs} (\ref{AbelianInvariantsSubgroupFpGroupRrs:for a group and a coset table}). }

 

\subsection{\textcolor{Chapter }{AbelianInvariantsSubgroupFpGroupMtc}}
\logpage{[ 47, 15, 2 ]}\nobreak
\hyperdef{L}{X804F664180BA2134}{}
{\noindent\textcolor{FuncColor}{$\triangleright$\enspace\texttt{AbelianInvariantsSubgroupFpGroupMtc({\mdseries\slshape G, H})\index{AbelianInvariantsSubgroupFpGroupMtc@\texttt{AbelianInvariantsSubgroupFpGroupMtc}}
\label{AbelianInvariantsSubgroupFpGroupMtc}
}\hfill{\scriptsize (function)}}\\


 uses the Modified Todd\texttt{\symbol{45}}Coxeter method to compute the
abelian invariants of a subgroup \mbox{\texttt{\mdseries\slshape H}} of a finitely presented group \mbox{\texttt{\mdseries\slshape G}}. }

 
\subsection{\textcolor{Chapter }{AbelianInvariantsSubgroupFpGroupRrs}}\logpage{[ 47, 15, 3 ]}
\hyperdef{L}{X8586137B7AAA6C10}{}
{
\noindent\textcolor{FuncColor}{$\triangleright$\enspace\texttt{AbelianInvariantsSubgroupFpGroupRrs({\mdseries\slshape G, H})\index{AbelianInvariantsSubgroupFpGroupRrs@\texttt{AbelianInvariantsSubgroupFpGroupRrs}!for two groups}
\label{AbelianInvariantsSubgroupFpGroupRrs:for two groups}
}\hfill{\scriptsize (function)}}\\
\noindent\textcolor{FuncColor}{$\triangleright$\enspace\texttt{AbelianInvariantsSubgroupFpGroupRrs({\mdseries\slshape G, table})\index{AbelianInvariantsSubgroupFpGroupRrs@\texttt{AbelianInvariantsSubgroupFpGroupRrs}!for a group and a coset table}
\label{AbelianInvariantsSubgroupFpGroupRrs:for a group and a coset table}
}\hfill{\scriptsize (function)}}\\


 uses the Reduced Reidemeister\texttt{\symbol{45}}Schreier method to compute
the abelian invariants of a subgroup \mbox{\texttt{\mdseries\slshape H}} of a finitely presented group \mbox{\texttt{\mdseries\slshape G}}. 

 Alternatively to the subgroup \mbox{\texttt{\mdseries\slshape H}}, its coset table \mbox{\texttt{\mdseries\slshape table}} in \mbox{\texttt{\mdseries\slshape G}} may be given as second argument. }

 

\subsection{\textcolor{Chapter }{AbelianInvariantsNormalClosureFpGroup}}
\logpage{[ 47, 15, 4 ]}\nobreak
\hyperdef{L}{X850E4CD784F6EAA8}{}
{\noindent\textcolor{FuncColor}{$\triangleright$\enspace\texttt{AbelianInvariantsNormalClosureFpGroup({\mdseries\slshape G, H})\index{AbelianInvariantsNormalClosureFpGroup@\texttt{Abelian}\-\texttt{Invariants}\-\texttt{Normal}\-\texttt{Closure}\-\texttt{Fp}\-\texttt{Group}}
\label{AbelianInvariantsNormalClosureFpGroup}
}\hfill{\scriptsize (function)}}\\


 \texttt{AbelianInvariantsNormalClosureFpGroup} is a synonym for \texttt{AbelianInvariantsNormalClosureFpGroupRrs} (\ref{AbelianInvariantsNormalClosureFpGroupRrs}). }

 

\subsection{\textcolor{Chapter }{AbelianInvariantsNormalClosureFpGroupRrs}}
\logpage{[ 47, 15, 5 ]}\nobreak
\hyperdef{L}{X801635B28079E56A}{}
{\noindent\textcolor{FuncColor}{$\triangleright$\enspace\texttt{AbelianInvariantsNormalClosureFpGroupRrs({\mdseries\slshape G, H})\index{AbelianInvariantsNormalClosureFpGroupRrs@\texttt{Abelian}\-\texttt{Invariants}\-\texttt{Normal}\-\texttt{Closure}\-\texttt{Fp}\-\texttt{GroupRrs}}
\label{AbelianInvariantsNormalClosureFpGroupRrs}
}\hfill{\scriptsize (function)}}\\


 uses the Reduced Reidemeister\texttt{\symbol{45}}Schreier method to compute
the abelian invariants of the normal closure of a subgroup \mbox{\texttt{\mdseries\slshape H}} of a finitely presented group \mbox{\texttt{\mdseries\slshape G}}. See \ref{Subgroup Presentations} for details on the different strategies. 

 The following example shows a calculation for the Coxeter group $B_1$. This calculation and a similar one for $B_0$ have been used to prove that $B_1' / B_1'' \cong Z_2^9 \times Z^3$ and $B_0' / B_0'' \cong Z_2^{91} \times Z^{27}$ as stated in in \cite[Proposition 5]{FJNT95}. 

 
\begin{Verbatim}[commandchars=!@|,fontsize=\small,frame=single,label=Example]
  !gapprompt@gap>| !gapinput@# Define the Coxeter group E1.|
  !gapprompt@gap>| !gapinput@F := FreeGroup( "x1", "x2", "x3", "x4", "x5" );|
  <free group on the generators [ x1, x2, x3, x4, x5 ]>
  !gapprompt@gap>| !gapinput@x1 := F.1;; x2 := F.2;; x3 := F.3;; x4 := F.4;; x5 := F.5;;|
  !gapprompt@gap>| !gapinput@rels := [ x1^2, x2^2, x3^2, x4^2, x5^2,|
  !gapprompt@>| !gapinput@ (x1 * x3)^2, (x2 * x4)^2, (x1 * x2)^3, (x2 * x3)^3, (x3 * x4)^3,|
  !gapprompt@>| !gapinput@ (x4 * x1)^3, (x1 * x5)^3, (x2 * x5)^2, (x3 * x5)^3, (x4 * x5)^2,|
  !gapprompt@>| !gapinput@ (x1 * x2 * x3 * x4 * x3 * x2)^2 ];;|
  !gapprompt@gap>| !gapinput@E1 := F / rels;|
  <fp group on the generators [ x1, x2, x3, x4, x5 ]>
  !gapprompt@gap>| !gapinput@x1 := E1.1;; x2 := E1.2;; x3 := E1.3;; x4 := E1.4;; x5 := E1.5;;|
  !gapprompt@gap>| !gapinput@# Get normal subgroup generators for B1.|
  !gapprompt@gap>| !gapinput@H := Subgroup( E1, [ x5 * x2^-1, x5 * x4^-1 ] );;|
  !gapprompt@gap>| !gapinput@# Compute the abelian invariants of B1/B1'.|
  !gapprompt@gap>| !gapinput@A := AbelianInvariantsNormalClosureFpGroup( E1, H );|
  [ 2, 2, 2, 2, 2, 2, 2, 2 ]
  !gapprompt@gap>| !gapinput@# Compute a presentation for B1.|
  !gapprompt@gap>| !gapinput@P := PresentationNormalClosure( E1, H );|
  <presentation with 18 gens and 46 rels of total length 132>
  !gapprompt@gap>| !gapinput@SimplifyPresentation( P );|
  #I  there are 8 generators and 30 relators of total length 148
  !gapprompt@gap>| !gapinput@B1 := FpGroupPresentation( P );|
  <fp group on the generators [ _x1, _x2, _x3, _x4, _x6, _x7, _x8, _x11
   ]>
  !gapprompt@gap>| !gapinput@# Compute normal subgroup generators for B1'.|
  !gapprompt@gap>| !gapinput@gens := GeneratorsOfGroup( B1 );;|
  !gapprompt@gap>| !gapinput@numgens := Length( gens );;|
  !gapprompt@gap>| !gapinput@comms := [ ];;|
  !gapprompt@gap>| !gapinput@for i in [ 1 .. numgens - 1 ] do|
  !gapprompt@>| !gapinput@    for j in [i+1 .. numgens ] do|
  !gapprompt@>| !gapinput@        Add( comms, Comm( gens[i], gens[j] ) );|
  !gapprompt@>| !gapinput@    od;|
  !gapprompt@>| !gapinput@od;|
  !gapprompt@gap>| !gapinput@# Compute the abelian invariants of B1'/B1".|
  !gapprompt@gap>| !gapinput@K := Subgroup( B1, comms );;|
  !gapprompt@gap>| !gapinput@A := AbelianInvariantsNormalClosureFpGroup( B1, K );|
  [ 0, 0, 0, 2, 2, 2, 2, 2, 2, 2, 2, 2 ]
\end{Verbatim}
 }

 }

  
\section{\textcolor{Chapter }{Testing Finiteness of Finitely Presented Groups}}\label{Testing Finiteness of Finitely Presented Groups}
\logpage{[ 47, 16, 0 ]}
\hyperdef{L}{X86C43E3B81ED25DC}{}
{
  As a consequence of the algorithmic insolvabilities mentioned in the
introduction to this chapter, there cannot be a general method that will test
whether a given finitely presented group is actually finite. 

 Therefore testing the finiteness of a finitely presented group can be
problematic. What \textsf{GAP} actually does upon a call of \texttt{IsFinite} (\ref{IsFinite}) (or if it is {\textendash}probably implicitly{\textendash} asked for a
faithful permutation representation) is to test whether it can find (via coset
enumeration) a cyclic subgroup of finite index. If it can, it rewrites the
presentation to this subgroup. Since the subgroup is cyclic, its size can be
checked easily from the resulting presentation, the size of the whole group is
the product of the index and the subgroup size. Since however no bound for the
index of such a subgroup (if any exist) is known, such a test might continue
unsuccessfully until memory is exhausted. 

 On the other hand, a couple of methods exist, that might prove that a group is
infinite. Again, none is guaranteed to work in every case: 

 The first method is to find (for example via the low index algorithm,
see{\nobreakspace}\texttt{LowIndexSubgroupsFpGroup} (\ref{LowIndexSubgroupsFpGroup})) a subgroup $U$ such that $[U:U']$ is infinite. If $U$ has finite index, this can be checked by \texttt{IsInfiniteAbelianizationGroup} (\ref{IsInfiniteAbelianizationGroup}). 

 Note that this test has been done traditionally by checking the \texttt{AbelianInvariants} (\ref{AbelianInvariants}) (see section{\nobreakspace}\ref{Abelian Invariants for Subgroups}) of $U$, \texttt{IsInfiniteAbelianizationGroup} (\ref{IsInfiniteAbelianizationGroup}) does a similar calculation but stops as soon as it is known whether $0$ is an invariant without computing the actual values. This can be notably
faster. 

 Another method is based on $p$\texttt{\symbol{45}}group quotients, see \texttt{NewmanInfinityCriterion} (\ref{NewmanInfinityCriterion}). 

\subsection{\textcolor{Chapter }{IsInfiniteAbelianizationGroup}}
\logpage{[ 47, 16, 1 ]}\nobreak
\hyperdef{L}{X82F444F67BE0E4FE}{}
{\noindent\textcolor{FuncColor}{$\triangleright$\enspace\texttt{IsInfiniteAbelianizationGroup({\mdseries\slshape G})\index{IsInfiniteAbelianizationGroup@\texttt{IsInfiniteAbelianizationGroup}}
\label{IsInfiniteAbelianizationGroup}
}\hfill{\scriptsize (property)}}\\


 \index{IsInfiniteAbelianizationGroup@\texttt{IsInfiniteAbelianizationGroup}!for groups} returns true if the commutator factor group $\mbox{\texttt{\mdseries\slshape G}}/\mbox{\texttt{\mdseries\slshape G}}'$ is infinite. This might be done without computing the full structure of the
commutator factor group. }

 

\subsection{\textcolor{Chapter }{NewmanInfinityCriterion}}
\logpage{[ 47, 16, 2 ]}\nobreak
\hyperdef{L}{X85C9FD548394C1E2}{}
{\noindent\textcolor{FuncColor}{$\triangleright$\enspace\texttt{NewmanInfinityCriterion({\mdseries\slshape G, p})\index{NewmanInfinityCriterion@\texttt{NewmanInfinityCriterion}}
\label{NewmanInfinityCriterion}
}\hfill{\scriptsize (function)}}\\


 Let \mbox{\texttt{\mdseries\slshape G}} be a finitely presented group and \mbox{\texttt{\mdseries\slshape p}} a prime that divides the order of the commutator factor group of \mbox{\texttt{\mdseries\slshape G}}. This function applies an infinity criterion due to M.
F.{\nobreakspace}Newman \cite{New90} to \mbox{\texttt{\mdseries\slshape G}}. (See \cite[chapter 16]{Joh97} for a more explicit description.) It returns \texttt{true} if the criterion succeeds in proving that \mbox{\texttt{\mdseries\slshape G}} is infinite and \texttt{fail} otherwise. 

 Note that the criterion uses the number of generators and relations in the
presentation of \mbox{\texttt{\mdseries\slshape G}}. Reduction of the presentation via Tietze transformations (\texttt{IsomorphismSimplifiedFpGroup} (\ref{IsomorphismSimplifiedFpGroup})) therefore might produce an isomorphic group, for which the criterion will
work better. 
\begin{Verbatim}[commandchars=!@|,fontsize=\small,frame=single,label=Example]
  !gapprompt@gap>| !gapinput@g:=FibonacciGroup(2,9);|
  <fp group on the generators [ f1, f2, f3, f4, f5, f6, f7, f8, f9 ]>
  !gapprompt@gap>| !gapinput@hom:=EpimorphismNilpotentQuotient(g,2);;|
  !gapprompt@gap>| !gapinput@k:=Kernel(hom);;|
  !gapprompt@gap>| !gapinput@Index(g,k);|
  152
  !gapprompt@gap>| !gapinput@AbelianInvariants(k);|
  [ 5, 5, 5, 5, 5, 5, 5, 5, 5, 5, 5, 5, 5, 5, 5, 5, 5, 5 ]
  !gapprompt@gap>| !gapinput@NewmanInfinityCriterion(Kernel(hom),5);|
  true
\end{Verbatim}
 

 This proves that the subgroup \texttt{k} (and thus the whole group \texttt{g}) is infinite. (This is the original example from{\nobreakspace}\cite{New90}.) }

 }

 }

          
\chapter{\textcolor{Chapter }{Presentations and Tietze Transformations}}\label{Presentations and Tietze Transformations}
\logpage{[ 48, 0, 0 ]}
\hyperdef{L}{X782985197BE809BF}{}
{
  A finite presentation describes a group, but usually there is a multitude of
presentations that describe isomorphic groups. Therefore a presentation in \textsf{GAP} is different from a finitely presented group though there are ways to
translate between both. 

 An important feature of presentations is that they can be modified (see
sections \ref{Changing Presentations} to \ref{Tietze Transformations that introduce new Generators}). 

  If you only want to get new presentations for subgroups of a finitely
presented group (and do not want to manipulate presentations yourself),
chances are that the operation \texttt{IsomorphismFpGroup} (\ref{IsomorphismFpGroup}) already does what you want (see{\nobreakspace}\ref{New Presentations and Presentations for Subgroups}).  
\section{\textcolor{Chapter }{Creating Presentations}}\label{Creating Presentations}
\logpage{[ 48, 1, 0 ]}
\hyperdef{L}{X867D00387957450F}{}
{
  Most of the functions creating presentations and all functions performing
Tietze transformations on them sort the relators by increasing lengths. The
function \texttt{PresentationFpGroup} (\ref{PresentationFpGroup}) is an exception because it is intended to reflect the relators that were used
to define the involved f. p. group. You may use the command \texttt{TzSort} (\ref{TzSort}) to sort the presentation. 

\subsection{\textcolor{Chapter }{PresentationFpGroup}}
\logpage{[ 48, 1, 1 ]}\nobreak
\hyperdef{L}{X797867B287AD92F8}{}
{\noindent\textcolor{FuncColor}{$\triangleright$\enspace\texttt{PresentationFpGroup({\mdseries\slshape G[, printlevel]})\index{PresentationFpGroup@\texttt{PresentationFpGroup}}
\label{PresentationFpGroup}
}\hfill{\scriptsize (function)}}\\


 creates a presentation, i. e., a Tietze object, for the given finitely
presented group \mbox{\texttt{\mdseries\slshape G}}. This presentation will be exactly as the presentation of \mbox{\texttt{\mdseries\slshape G}} and \emph{no} initial Tietze transformations are applied to it. 

 The optional \mbox{\texttt{\mdseries\slshape printlevel}} parameter can be used to restrict or to extend the amount of output provided
by Tietze transformation commands when being applied to the created
presentation. The default value 1 is designed for interactive use and implies
explicit messages to be displayed by most of these commands. A \mbox{\texttt{\mdseries\slshape printlevel}} value of 0 will suppress these messages, whereas a \mbox{\texttt{\mdseries\slshape printlevel}} value of 2 will enforce some additional output. 
\begin{Verbatim}[commandchars=!@|,fontsize=\small,frame=single,label=Example]
  !gapprompt@gap>| !gapinput@f := FreeGroup( "a", "b" );|
  <free group on the generators [ a, b ]>
  !gapprompt@gap>| !gapinput@g := f / [ f.1^3, f.2^2, (f.1*f.2)^3 ];|
  <fp group on the generators [ a, b ]>
  !gapprompt@gap>| !gapinput@p := PresentationFpGroup( g );|
  <presentation with 2 gens and 3 rels of total length 11>
\end{Verbatim}
 }

 

\subsection{\textcolor{Chapter }{TzSort}}
\logpage{[ 48, 1, 2 ]}\nobreak
\hyperdef{L}{X8637837A79422445}{}
{\noindent\textcolor{FuncColor}{$\triangleright$\enspace\texttt{TzSort({\mdseries\slshape P})\index{TzSort@\texttt{TzSort}}
\label{TzSort}
}\hfill{\scriptsize (function)}}\\


 sorts the relators of the given presentation \mbox{\texttt{\mdseries\slshape P}} by increasing lengths. There is no particular ordering defined for the
relators of equal length. Note that \texttt{TzSort} does not return a new object. It changes the given presentation. }

 

\subsection{\textcolor{Chapter }{GeneratorsOfPresentation}}
\logpage{[ 48, 1, 3 ]}\nobreak
\hyperdef{L}{X849429BC7D435F77}{}
{\noindent\textcolor{FuncColor}{$\triangleright$\enspace\texttt{GeneratorsOfPresentation({\mdseries\slshape P})\index{GeneratorsOfPresentation@\texttt{GeneratorsOfPresentation}}
\label{GeneratorsOfPresentation}
}\hfill{\scriptsize (function)}}\\


 returns a list of free generators that is a shallow copy (see \texttt{ShallowCopy} (\ref{ShallowCopy})) of the current generators of the presentation \mbox{\texttt{\mdseries\slshape P}}. }

 

\subsection{\textcolor{Chapter }{FpGroupPresentation}}
\logpage{[ 48, 1, 4 ]}\nobreak
\hyperdef{L}{X7D6F40A87F24D3D6}{}
{\noindent\textcolor{FuncColor}{$\triangleright$\enspace\texttt{FpGroupPresentation({\mdseries\slshape P[, nam]})\index{FpGroupPresentation@\texttt{FpGroupPresentation}}
\label{FpGroupPresentation}
}\hfill{\scriptsize (function)}}\\


 constructs an f. p. group as defined by the given Tietze presentation \mbox{\texttt{\mdseries\slshape P}}. 
\begin{Verbatim}[commandchars=!@|,fontsize=\small,frame=single,label=Example]
  !gapprompt@gap>| !gapinput@h := FpGroupPresentation( p );|
  <fp group on the generators [ a, b ]>
  !gapprompt@gap>| !gapinput@h = g;|
  false
\end{Verbatim}
 }

 

\subsection{\textcolor{Chapter }{PresentationViaCosetTable}}
\logpage{[ 48, 1, 5 ]}\nobreak
\hyperdef{L}{X84E056C57AFEDEA8}{}
{\noindent\textcolor{FuncColor}{$\triangleright$\enspace\texttt{PresentationViaCosetTable({\mdseries\slshape G[, F, words]})\index{PresentationViaCosetTable@\texttt{PresentationViaCosetTable}}
\label{PresentationViaCosetTable}
}\hfill{\scriptsize (function)}}\\


 constructs a presentation for a given concrete finite group. It applies the
relations finding algorithm which has been described in \cite{Can73} and \cite{Neu82}. It automatically applies Tietze transformations to the presentation found. 

 If only a group \mbox{\texttt{\mdseries\slshape G}} has been specified, the single stage algorithm is applied. 

 The operation \texttt{IsomorphismFpGroup} (\ref{IsomorphismFpGroup}) in contrast uses a multiple\texttt{\symbol{45}}stage algorithm using a chief
series and stabilizer chains. It usually should be used rather than \texttt{PresentationViaCosetTable}. (It does not apply Tietze transformations automatically.) 

 If the two stage algorithm is to be used, \texttt{PresentationViaCosetTable} expects a subgroup \mbox{\texttt{\mdseries\slshape H}} of \mbox{\texttt{\mdseries\slshape G}} to be provided in form of two additional arguments \mbox{\texttt{\mdseries\slshape F}} and \mbox{\texttt{\mdseries\slshape words}}, where \mbox{\texttt{\mdseries\slshape F}} is a free group with the same number of generators as \mbox{\texttt{\mdseries\slshape G}}, and \mbox{\texttt{\mdseries\slshape words}} is a list of words in the generators of \mbox{\texttt{\mdseries\slshape F}} which supply a list of generators of \mbox{\texttt{\mdseries\slshape H}} if they are evaluated as words in the corresponding generators of \mbox{\texttt{\mdseries\slshape G}}. 
\begin{Verbatim}[commandchars=!@|,fontsize=\small,frame=single,label=Example]
  !gapprompt@gap>| !gapinput@G := GeneralLinearGroup( 2, 7 );|
  GL(2,7)
  !gapprompt@gap>| !gapinput@GeneratorsOfGroup( G );|
  [ [ [ Z(7), 0*Z(7) ], [ 0*Z(7), Z(7)^0 ] ],
    [ [ Z(7)^3, Z(7)^0 ], [ Z(7)^3, 0*Z(7) ] ] ]
  !gapprompt@gap>| !gapinput@Size( G );|
  2016
  !gapprompt@gap>| !gapinput@P := PresentationViaCosetTable( G );|
  <presentation with 2 gens and 5 rels of total length 46>
  !gapprompt@gap>| !gapinput@TzPrintRelators( P );|
  #I  1. f2^3
  #I  2. f1^6
  #I  3. (f1*f2)^6
  #I  4. f1*f2*f1^-1*f2*f1*f2^-1*f1^-1*f2*f1*f2*f1^-1*f2^-1
  #I  5. f1^-3*f2*f1*f2*(f1^-1*f2^-1)^2*f1^-2*f2
\end{Verbatim}
 

 The two stage algorithm saves an essential amount of space by constructing two
coset tables of lengths $|H|$ and $|G|/|H|$ instead of just one coset table of length $|G|$. The next example shows an application of this option in the case of a
subgroup of size 7920 and index 12 in a permutation group of size 95040. 

 
\begin{Verbatim}[commandchars=!@|,fontsize=\small,frame=single,label=Example]
  !gapprompt@gap>| !gapinput@M12 := Group( [ (1,2,3,4,5,6,7,8,9,10,11), (3,7,11,8)(4,10,5,6),|
  !gapprompt@>| !gapinput@(1,12)(2,11)(3,6)(4,8)(5,9)(7,10) ], () );;|
  !gapprompt@gap>| !gapinput@F := FreeGroup( "a", "b", "c" );|
  <free group on the generators [ a, b, c ]>
  !gapprompt@gap>| !gapinput@words := [ F.1, F.2 ];|
  [ a, b ]
  !gapprompt@gap>| !gapinput@P := PresentationViaCosetTable( M12, F, words );|
  <presentation with 3 gens and 10 rels of total length 97>
  !gapprompt@gap>| !gapinput@G := FpGroupPresentation( P );|
  <fp group on the generators [ a, b, c ]>
  !gapprompt@gap>| !gapinput@RelatorsOfFpGroup( G );|
  [ c^2, b^4, (a*c)^3, (a*b^-2)^3, a^11,
    a^2*b*a^-2*b^-1*(b^-1*a)^2*a*b^-1, (a*(b*a^-1)^2*b^-1)^2,
    a^2*b*a^2*b^-2*a^-1*b*(a^-1*b^-1)^2,
    a^2*b^-1*a^-1*b^-1*a*c*b*c*(a*b)^2, a^2*(a^2*b)^2*a^-2*c*a*b*a^-1*c
   ]
\end{Verbatim}
 

 Before it is returned, the resulting presentation is being simplified by
appropriate calls of the function \texttt{SimplifyPresentation} (\ref{SimplifyPresentation}) (see \ref{Tietze Transformations}), but without allowing any eliminations of generators. This restriction
guarantees that we get a bijection between the list of generators of \mbox{\texttt{\mdseries\slshape G}} and the list of generators in the presentation. Hence, if the generators of \mbox{\texttt{\mdseries\slshape G}} are redundant and if you don't care for the bijection, you may get a shorter
presentation by calling the function \texttt{SimplifyPresentation} (\ref{SimplifyPresentation}), now without this restriction, once more yourself. 

 
\begin{Verbatim}[commandchars=!@|,fontsize=\small,frame=single,label=Example]
  !gapprompt@gap>| !gapinput@H := Group(|
  !gapprompt@>| !gapinput@[ (2,5,3), (2,7,5), (1,8,4), (1,8,6), (4,8,6), (3,5,7) ], () );;|
  !gapprompt@gap>| !gapinput@P := PresentationViaCosetTable( H );|
  <presentation with 6 gens and 12 rels of total length 42>
  !gapprompt@gap>| !gapinput@SimplifyPresentation( P );|
  #I  there are 4 generators and 10 relators of total length 36
\end{Verbatim}
 

 If you apply the function \texttt{FpGroupPresentation} (\ref{FpGroupPresentation}) to the resulting presentation you will get a finitely presented group
isomorphic to \mbox{\texttt{\mdseries\slshape G}}. Note, however, that the function \texttt{IsomorphismFpGroup} (\ref{IsomorphismFpGroup}) is recommended for this purpose. }

 

\subsection{\textcolor{Chapter }{SimplifiedFpGroup}}
\logpage{[ 48, 1, 6 ]}\nobreak
\hyperdef{L}{X7E1F2658827FC228}{}
{\noindent\textcolor{FuncColor}{$\triangleright$\enspace\texttt{SimplifiedFpGroup({\mdseries\slshape G})\index{SimplifiedFpGroup@\texttt{SimplifiedFpGroup}}
\label{SimplifiedFpGroup}
}\hfill{\scriptsize (function)}}\\


 applies Tietze transformations to a copy of the presentation of the given
finitely presented group \mbox{\texttt{\mdseries\slshape G}} in order to reduce it with respect to the number of generators, the number of
relators, and the relator lengths. 

 \texttt{SimplifiedFpGroup} returns a group isomorphic to the given one with a presentation which has been
tried to simplify via Tietze transformations. 

 If the connection to the original group is important, then the operation \texttt{IsomorphismSimplifiedFpGroup} (\ref{IsomorphismSimplifiedFpGroup}) should be used instead. 

 
\begin{Verbatim}[commandchars=!@|,fontsize=\small,frame=single,label=Example]
  !gapprompt@gap>| !gapinput@F6 := FreeGroup( 6, "G" );;|
  !gapprompt@gap>| !gapinput@G := F6 / [ F6.1^2, F6.2^2, F6.4*F6.6^-1, F6.5^2, F6.6^2,|
  !gapprompt@>| !gapinput@F6.1*F6.2^-1*F6.3, F6.1*F6.5*F6.3^-1, F6.2*F6.4^-1*F6.3,|
  !gapprompt@>| !gapinput@F6.3*F6.4*F6.5^-1, F6.1*F6.6*F6.3^-2, F6.3^4 ];;|
  !gapprompt@gap>| !gapinput@H := SimplifiedFpGroup( G );|
  <fp group on the generators [ G1, G3 ]>
  !gapprompt@gap>| !gapinput@RelatorsOfFpGroup( H );|
  [ G1^2, (G1*G3^-1)^2, G3^4 ]
\end{Verbatim}
 

 In fact, the command 

 
\begin{Verbatim}[commandchars=!@|,fontsize=\small,frame=single,label=Example]
  H := SimplifiedFpGroup( G );
\end{Verbatim}
 

 is an abbreviation of the command sequence 

 
\begin{Verbatim}[commandchars=!@|,fontsize=\small,frame=single,label=Example]
  P := PresentationFpGroup( G, 0 );;
  SimplifyPresentation( P );
  H := FpGroupPresentation( P );
\end{Verbatim}
 

 which applies a rather simple\texttt{\symbol{45}}minded strategy of Tietze
transformations to the intermediate presentation \mbox{\texttt{\mdseries\slshape P}}. If, for some concrete group, the resulting presentation is unsatisfying,
then you should try a more sophisticated, interactive use of the available
Tietze transformation commands (see \ref{Tietze Transformations}). }

 }

  
\section{\textcolor{Chapter }{Subgroup Presentations}}\label{Subgroup Presentations}
\logpage{[ 48, 2, 0 ]}
\hyperdef{L}{X8118FECE7AD1879B}{}
{
  \index{Schreier} 

\subsection{\textcolor{Chapter }{PresentationSubgroup}}
\logpage{[ 48, 2, 1 ]}\nobreak
\hyperdef{L}{X7DB32FA97DAC5AC8}{}
{\noindent\textcolor{FuncColor}{$\triangleright$\enspace\texttt{PresentationSubgroup({\mdseries\slshape G, H[, string]})\index{PresentationSubgroup@\texttt{PresentationSubgroup}}
\label{PresentationSubgroup}
}\hfill{\scriptsize (function)}}\\


 \texttt{PresentationSubgroup} is a synonym for \texttt{PresentationSubgroupRrs} (\ref{PresentationSubgroupRrs:for a group and a coset table (and a string)}). }

 
\subsection{\textcolor{Chapter }{PresentationSubgroupRrs}}\logpage{[ 48, 2, 2 ]}
\hyperdef{L}{X857365CD87ADC29E}{}
{
\noindent\textcolor{FuncColor}{$\triangleright$\enspace\texttt{PresentationSubgroupRrs({\mdseries\slshape G, H[, string]})\index{PresentationSubgroupRrs@\texttt{PresentationSubgroupRrs}!for two groups (and a string)}
\label{PresentationSubgroupRrs:for two groups (and a string)}
}\hfill{\scriptsize (function)}}\\
\noindent\textcolor{FuncColor}{$\triangleright$\enspace\texttt{PresentationSubgroupRrs({\mdseries\slshape G, table[, string]})\index{PresentationSubgroupRrs@\texttt{PresentationSubgroupRrs}!for a group and a coset table (and a string)}
\label{PresentationSubgroupRrs:for a group and a coset table (and a string)}
}\hfill{\scriptsize (function)}}\\


 uses the Reduced Reidemeister\texttt{\symbol{45}}Schreier method to compute a
presentation \mbox{\texttt{\mdseries\slshape P}} for a subgroup \mbox{\texttt{\mdseries\slshape H}} of a finitely presented group \mbox{\texttt{\mdseries\slshape G}}. The generators in the resulting presentation will be named \mbox{\texttt{\mdseries\slshape string}}\texttt{1}, \mbox{\texttt{\mdseries\slshape string}}\texttt{2}, $\ldots$, the default string is \texttt{"{\textunderscore}x"}. You may access the $i$\texttt{\symbol{45}}th of these generators by \mbox{\texttt{\mdseries\slshape P}}\texttt{!.}$i$. 

 Alternatively to the subgroup \mbox{\texttt{\mdseries\slshape H}}, its coset table \mbox{\texttt{\mdseries\slshape table}} in \mbox{\texttt{\mdseries\slshape G}} may be given as second argument. 
\begin{Verbatim}[commandchars=!@|,fontsize=\small,frame=single,label=Example]
  !gapprompt@gap>| !gapinput@f := FreeGroup( "a", "b" );;|
  !gapprompt@gap>| !gapinput@g := f / [ f.1^2, f.2^3, (f.1*f.2)^5 ];|
  <fp group on the generators [ a, b ]>
  !gapprompt@gap>| !gapinput@g1 := Size( g );|
  60
  !gapprompt@gap>| !gapinput@u := Subgroup( g, [ g.1, g.1^g.2 ] );|
  Group([ a, b^-1*a*b ])
  !gapprompt@gap>| !gapinput@p := PresentationSubgroup( g, u, "g" );|
  <presentation with 3 gens and 4 rels of total length 12>
  !gapprompt@gap>| !gapinput@gens := GeneratorsOfPresentation( p );|
  [ g1, g2, g3 ]
  !gapprompt@gap>| !gapinput@TzPrintRelators( p );|
  #I  1. g1^2
  #I  2. g2^2
  #I  3. g3*g2*g1
  #I  4. g3^5
\end{Verbatim}
 

 Note that you cannot call the generators by their names. These names are not
variables, but just display figures. So, if you want to access the generators
by their names, you first will have to introduce the respective variables and
to assign the generators to them. 

 
\begin{Verbatim}[commandchars=!@|,fontsize=\small,frame=single,label=Example]
  !gapprompt@gap>| !gapinput@gens[1] = g1;|
  false
  !gapprompt@gap>| !gapinput@g1;|
  60
  !gapprompt@gap>| !gapinput@g1 := gens[1];; g2 := gens[2];; g3 := gens[3];;|
  !gapprompt@gap>| !gapinput@g1;|
  g1
\end{Verbatim}
 

 The Reduced Reidemeister\texttt{\symbol{45}}Schreier algorithm is a
modification of the Reidemeister\texttt{\symbol{45}}Schreier algorithm of
George Havas \cite{Hav74b}. It was proposed by Joachim Neub{\"u}ser and first implemented in 1986 by
Andrea Lucchini and Volkmar Felsch in the SPAS system \cite{Spa89}. Like the Reidemeister\texttt{\symbol{45}}Schreier algorithm of George Havas,
it needs only the presentation of \mbox{\texttt{\mdseries\slshape G}} and a coset table of \mbox{\texttt{\mdseries\slshape H}} in \mbox{\texttt{\mdseries\slshape G}} to construct a presentation of \mbox{\texttt{\mdseries\slshape H}}. 

 Whenever you call the command \texttt{PresentationSubgroupRrs}, it first obtains a coset table of \mbox{\texttt{\mdseries\slshape H}} in \mbox{\texttt{\mdseries\slshape G}} if not given. Next, a set of generators of \mbox{\texttt{\mdseries\slshape H}} is determined by reconstructing the coset table and introducing in that
process as many Schreier generators of \mbox{\texttt{\mdseries\slshape H}} in \mbox{\texttt{\mdseries\slshape G}} as are needed to do a Felsch strategy coset enumeration without any
coincidences. (In general, though containing redundant generators, this set
will be much smaller than the set of all Schreier generators. That is why we
call the method the \emph{Reduced} Reidemeister\texttt{\symbol{45}}Schreier.) 

 After having constructed this set of \emph{primary subgroup generators}, say, the coset table is extended to an \emph{augmented coset table} which describes the action of the group generators on coset representatives,
i.e., on elements instead of cosets. For this purpose, suitable words in the
(primary) subgroup generators have to be associated to the coset table
entries. In order to keep the lengths of these words short, additional \emph{secondary subgroup generators} are introduced as abbreviations of subwords. Their number may be large. 

 Finally, a Reidemeister rewriting process is used to get defining relators for \mbox{\texttt{\mdseries\slshape H}} from the relators of \mbox{\texttt{\mdseries\slshape G}}. As the resulting presentation of \mbox{\texttt{\mdseries\slshape H}} is a presentation on primary \emph{and} secondary generators, in general you will have to simplify it by appropriate
Tietze transformations (see \ref{Tietze Transformations}) or by the command \texttt{DecodeTree} (\ref{DecodeTree}) before you can use it. Therefore it is returned in the form of a presentation, \mbox{\texttt{\mdseries\slshape P}} say. 

 Compared with the Modified Todd\texttt{\symbol{45}}Coxeter method described
below, the Reduced Reidemeister\texttt{\symbol{45}}Schreier method (as well as
Havas' original Reidemeister\texttt{\symbol{45}}Schreier program) has the
advantage that it does not require generators of \mbox{\texttt{\mdseries\slshape H}} to be given if a coset table of \mbox{\texttt{\mdseries\slshape H}} in \mbox{\texttt{\mdseries\slshape G}} is known. This provides a possibility to compute a presentation of the normal
closure of a given subgroup (see \texttt{PresentationNormalClosureRrs} (\ref{PresentationNormalClosureRrs})). 

 For certain applications you may be interested in getting not only just a
presentation for \mbox{\texttt{\mdseries\slshape H}}, but also a relation between the involved generators of \mbox{\texttt{\mdseries\slshape H}} and the generators of \mbox{\texttt{\mdseries\slshape G}}. The subgroup generators in the presentation are sorted such that the primary
generators precede the secondary ones. Moreover, for each secondary subgroup
generator there is a relator in the presentation which expresses this
generator as a word in preceding ones. Hence, all we need in addition is a
list of words in the generators of \mbox{\texttt{\mdseries\slshape G}} which express the primary subgroup generators. In fact, such a list is
provided in the attribute \texttt{PrimaryGeneratorWords} (\ref{PrimaryGeneratorWords}) of the resulting presentation. }

 

\subsection{\textcolor{Chapter }{PrimaryGeneratorWords}}
\logpage{[ 48, 2, 3 ]}\nobreak
\hyperdef{L}{X7FCE7ED581CF7897}{}
{\noindent\textcolor{FuncColor}{$\triangleright$\enspace\texttt{PrimaryGeneratorWords({\mdseries\slshape P})\index{PrimaryGeneratorWords@\texttt{PrimaryGeneratorWords}}
\label{PrimaryGeneratorWords}
}\hfill{\scriptsize (attribute)}}\\


 is an attribute of the presentation \mbox{\texttt{\mdseries\slshape P}} which holds a list of words in the associated group generators (of the
underlying free group) which express the primary subgroup generators of \mbox{\texttt{\mdseries\slshape P}}. 
\begin{Verbatim}[commandchars=!@|,fontsize=\small,frame=single,label=Example]
  !gapprompt@gap>| !gapinput@PrimaryGeneratorWords( p );|
  [ a, b^-1*a*b ]
\end{Verbatim}
 }

 

\subsection{\textcolor{Chapter }{PresentationSubgroupMtc}}
\logpage{[ 48, 2, 4 ]}\nobreak
\hyperdef{L}{X80BA10F780EAE68E}{}
{\noindent\textcolor{FuncColor}{$\triangleright$\enspace\texttt{PresentationSubgroupMtc({\mdseries\slshape G, H[, string][, print, level]})\index{PresentationSubgroupMtc@\texttt{PresentationSubgroupMtc}}
\label{PresentationSubgroupMtc}
}\hfill{\scriptsize (function)}}\\


 uses the Modified Todd\texttt{\symbol{45}}Coxeter coset representative
enumeration method to compute a presentation $P$ for a subgroup \mbox{\texttt{\mdseries\slshape H}} of a finitely presented group \mbox{\texttt{\mdseries\slshape G}}. The presentation returned is in generators corresponding to the generators
of \mbox{\texttt{\mdseries\slshape H}}. The generators in the resulting presentation will be named \mbox{\texttt{\mdseries\slshape string}}\texttt{1}, \mbox{\texttt{\mdseries\slshape string}}\texttt{2}, $\ldots$, the default string is \texttt{"{\textunderscore}x"}. You may access the $i$\texttt{\symbol{45}}th of these generators by $P$\texttt{!.}$i$. 

 The default print level is $1$. If the print level is set to $0$, then the printout of the implicitly called function \texttt{DecodeTree} (\ref{DecodeTree}) will be suppressed. 
\begin{Verbatim}[commandchars=!@|,fontsize=\small,frame=single,label=Example]
  !gapprompt@gap>| !gapinput@p := PresentationSubgroupMtc( g, u );|
  <presentation with 2 gens and 3 rels of total length 14>
\end{Verbatim}
 

 The so called Modified Todd\texttt{\symbol{45}}Coxeter method was proposed, in
slightly different forms, by Nathan S.{\nobreakspace}Mendelsohn and William
O.{\nobreakspace}J.{\nobreakspace}Moser in 1966. Moser's method was proved in \cite{BC76}. It has been generalized to cover a broad spectrum of different versions (see
the survey \cite{Neu82}). 

 The \emph{Modified Todd\texttt{\symbol{45}}Coxeter} method performs an enumeration of coset representatives. It proceeds like an
ordinary coset enumeration (see \ref{Coset Tables and Coset Enumeration}), but as the product of a coset representative by a group generator or its
inverse need not be a coset representative itself, the Modified
Todd\texttt{\symbol{45}}Coxeter has to store a kind of correction element for
each coset table entry. Hence it builds up a so called \emph{augmented coset table} of \mbox{\texttt{\mdseries\slshape H}} in \mbox{\texttt{\mdseries\slshape G}} consisting of the ordinary coset table and a second table in parallel which
contains the associated subgroup elements. 

 Theoretically, these subgroup elements could be expressed as words in the
given generators of \mbox{\texttt{\mdseries\slshape H}}, but in general these words tend to become unmanageable because of their
enormous lengths. Therefore, a highly redundant list of subgroup generators is
built up starting from the given (``primary'') generators of \mbox{\texttt{\mdseries\slshape H}} and adding additional (``secondary'') generators which are defined as abbreviations of suitable words of length
two in the preceding generators such that each of the subgroup elements in the
augmented coset table can be expressed as a word of length at most one in the
resulting (primary \emph{and} secondary) subgroup generators. 

 Then a rewriting process (which is essentially a kind of Reidemeister
rewriting process) is used to get relators for \mbox{\texttt{\mdseries\slshape H}} from the defining relators of \mbox{\texttt{\mdseries\slshape G}}. 

 The resulting presentation involves all the primary, but not all the secondary
generators of \mbox{\texttt{\mdseries\slshape H}}. In fact, it contains only those secondary generators which explicitly occur
in the augmented coset table. If we extended this presentation by those
secondary generators which are not yet contained in it as additional
generators, and by the definitions of all secondary generators as additional
relators, we would get a presentation of \mbox{\texttt{\mdseries\slshape H}}, but, in general, we would end up with a large number of generators and
relators. 

 On the other hand, if we avoid this extension, the current presentation will
not necessarily define \mbox{\texttt{\mdseries\slshape H}} although we have used the same rewriting process which in the case of the \texttt{PresentationSubgroupRrs} (\ref{PresentationSubgroupRrs:for a group and a coset table (and a string)}) command computes a defining set of relators for \mbox{\texttt{\mdseries\slshape H}} from an augmented coset table and defining relators of \mbox{\texttt{\mdseries\slshape G}}. The different behaviour here is caused by the fact that coincidences may
have occurred in the Modified Todd\texttt{\symbol{45}}Coxeter coset
enumeration. 

 To overcome this problem without extending the presentation by all secondary
generators, the \texttt{PresentationSubgroupMtc} command applies the so called \emph{decoding tree} algorithm which provides a more economical approach. The reader is strongly
recommended to carefully read section \ref{sect:DecodeTree} where this algorithm is described in more detail. Here we will only mention
that this procedure may add a lot of intermediate generators and relators (and
even change the isomorphism type) in a process which in fact eliminates all
secondary generators from the presentation and hence finally provides a
presentation of \mbox{\texttt{\mdseries\slshape H}} on the primary, i.e., the originally given, generators of \mbox{\texttt{\mdseries\slshape H}}. This is a remarkable advantage of the command \texttt{PresentationSubgroupMtc} compared to the command \texttt{PresentationSubgroupRrs} (\ref{PresentationSubgroupRrs:for a group and a coset table (and a string)}). But note that, for some particular subgroup \mbox{\texttt{\mdseries\slshape H}}, the Reduced Reidemeister\texttt{\symbol{45}}Schreier method might quite well
produce a more concise presentation. 

 The resulting presentation is returned in the form of a presentation, $P$ say. 

 As the function \texttt{PresentationSubgroupRrs} (\ref{PresentationSubgroupRrs:for a group and a coset table (and a string)}) described above (see there for details), the function \texttt{PresentationSubgroupMtc} returns a list of the primary subgroup generators of \mbox{\texttt{\mdseries\slshape H}} in the attribute \texttt{PrimaryGeneratorWords} (\ref{PrimaryGeneratorWords}) of $P$. In fact, this list is not very exciting here because it is just a shallow
copy of the value of \texttt{GeneratorsOfPresentation} (\ref{GeneratorsOfPresentation}) of \mbox{\texttt{\mdseries\slshape H}}, however it is needed to guarantee a certain consistency between the results
of the different functions for computing subgroup presentations. 

 Though the decoding tree routine already involves a lot of Tietze
transformations, we recommend that you try to further simplify the resulting
presentation by appropriate Tietze transformations (see \ref{Tietze Transformations}). }

 

\subsection{\textcolor{Chapter }{PresentationNormalClosureRrs}}
\logpage{[ 48, 2, 5 ]}\nobreak
\hyperdef{L}{X7D6A52837BEE5C3D}{}
{\noindent\textcolor{FuncColor}{$\triangleright$\enspace\texttt{PresentationNormalClosureRrs({\mdseries\slshape G, H[, string]})\index{PresentationNormalClosureRrs@\texttt{PresentationNormalClosureRrs}}
\label{PresentationNormalClosureRrs}
}\hfill{\scriptsize (function)}}\\


 uses the Reduced Reidemeister\texttt{\symbol{45}}Schreier method to compute a
presentation $P$ for the normal closure of a subgroup \mbox{\texttt{\mdseries\slshape H}} of a finitely presented group \mbox{\texttt{\mdseries\slshape G}}. The generators in the resulting presentation will be named \mbox{\texttt{\mdseries\slshape string}}\texttt{1}, \mbox{\texttt{\mdseries\slshape string}}\texttt{2}, $\ldots$, the default string is \texttt{"{\textunderscore}x"}. You may access the $i$\texttt{\symbol{45}}th of these generators by $P$\texttt{!.}$i$. }

 

\subsection{\textcolor{Chapter }{PresentationNormalClosure}}
\logpage{[ 48, 2, 6 ]}\nobreak
\hyperdef{L}{X7A7E5D0084DB0B4F}{}
{\noindent\textcolor{FuncColor}{$\triangleright$\enspace\texttt{PresentationNormalClosure({\mdseries\slshape G, H[, string]})\index{PresentationNormalClosure@\texttt{PresentationNormalClosure}}
\label{PresentationNormalClosure}
}\hfill{\scriptsize (function)}}\\


 \texttt{PresentationNormalClosure} is a synonym for \texttt{PresentationNormalClosureRrs} (\ref{PresentationNormalClosureRrs}). }

 }

  
\section{\textcolor{Chapter }{Relators in a Presentation}}\label{Relators in a Presentation}
\logpage{[ 48, 3, 0 ]}
\hyperdef{L}{X7BC960AB7E8DE419}{}
{
  In order to speed up the Tietze transformation routines, each relator in a
presentation is internally represented by a list of positive or negative
generator numbers, i.e., each factor of the proper \textsf{GAP} word is represented by the position number of the corresponding generator with
respect to the current list of generators, or by the respective negative
number, if the factor is the inverse of a generator. Note that the numbering
of the generators in Tietze words is always relative to a generator list and
bears no relation to the internal numbering of generators in a family of
associative words. 

\subsection{\textcolor{Chapter }{TietzeWordAbstractWord}}
\logpage{[ 48, 3, 1 ]}\nobreak
\hyperdef{L}{X8365BAFA785FCD8D}{}
{\noindent\textcolor{FuncColor}{$\triangleright$\enspace\texttt{TietzeWordAbstractWord({\mdseries\slshape word, fgens})\index{TietzeWordAbstractWord@\texttt{TietzeWordAbstractWord}}
\label{TietzeWordAbstractWord}
}\hfill{\scriptsize (operation)}}\\


 assumes \mbox{\texttt{\mdseries\slshape fgens}} to be a list of free group generators and \mbox{\texttt{\mdseries\slshape word}} to be an abstract word in these generators. It converts \mbox{\texttt{\mdseries\slshape word}} into a Tietze word, i. e., a list of positive or negative generator numbers. 

 This function simply calls \texttt{LetterRepAssocWord} (\ref{LetterRepAssocWord}). }

 

\subsection{\textcolor{Chapter }{AbstractWordTietzeWord}}
\logpage{[ 48, 3, 2 ]}\nobreak
\hyperdef{L}{X8573E91C838B1D13}{}
{\noindent\textcolor{FuncColor}{$\triangleright$\enspace\texttt{AbstractWordTietzeWord({\mdseries\slshape word, fgens})\index{AbstractWordTietzeWord@\texttt{AbstractWordTietzeWord}}
\label{AbstractWordTietzeWord}
}\hfill{\scriptsize (function)}}\\


 assumes \mbox{\texttt{\mdseries\slshape fgens}} to be a list of free group generators and \mbox{\texttt{\mdseries\slshape word}} to be a Tietze word in these generators, i. e., a list of positive or negative
generator numbers. It converts \mbox{\texttt{\mdseries\slshape word}} to an abstract word. 

 This function simply calls \texttt{AssocWordByLetterRep} (\ref{AssocWordByLetterRep}). 
\begin{Verbatim}[commandchars=!@|,fontsize=\small,frame=single,label=Example]
  !gapprompt@gap>| !gapinput@F := FreeGroup( "a", "b", "c" ,"d");|
  <free group on the generators [ a, b, c, d ]>
  !gapprompt@gap>| !gapinput@tzword := TietzeWordAbstractWord(|
  !gapprompt@>| !gapinput@Comm(F.4,F.2) * (F.3^2 * F.2)^-1, GeneratorsOfGroup( F ){[2,3,4]} );|
  [ -3, -1, 3, -2, -2 ]
  !gapprompt@gap>| !gapinput@AbstractWordTietzeWord( tzword, GeneratorsOfGroup( F ){[2,3,4]} );|
  d^-1*b^-1*d*c^-2
\end{Verbatim}
 }

 }

  
\section{\textcolor{Chapter }{Printing Presentations}}\label{Printing Presentations}
\logpage{[ 48, 4, 0 ]}
\hyperdef{L}{X867F64FA840B3F81}{}
{
  Whenever you create a presentation $P$, or assign it to a variable, \textsf{GAP} will respond by printing $P$. However, as $P$ may contain a lot of generators and many relators of large length, it would be
annoying if the standard print facilities displayed all this information in
detail. So they restrict the printout to just one line of text containing the
number of generators, the number of relators, and the total length of all
relators of $P$. As compensation, \textsf{GAP} offers some special print commands which display various details of a
presentation. Note that there is also a function \texttt{TzPrintOptions} (\ref{TzPrintOptions}). It is described in Section \ref{Tietze Options}. 

\subsection{\textcolor{Chapter }{TzPrintGenerators}}
\logpage{[ 48, 4, 1 ]}\nobreak
\hyperdef{L}{X847EA6737C21171C}{}
{\noindent\textcolor{FuncColor}{$\triangleright$\enspace\texttt{TzPrintGenerators({\mdseries\slshape P[, list]})\index{TzPrintGenerators@\texttt{TzPrintGenerators}}
\label{TzPrintGenerators}
}\hfill{\scriptsize (function)}}\\


 prints the generators of the given Tietze presentation \mbox{\texttt{\mdseries\slshape P}} together with the number of their occurrences in the relators. The optional
second argument can be used to specify the numbers of the generators to be
printed. Default: all generators are printed. 
\begin{Verbatim}[commandchars=!@|,fontsize=\small,frame=single,label=Example]
  !gapprompt@gap>| !gapinput@G := Group( [ (1,2,3,4,5), (2,3,5,4), (1,6)(3,4) ], () );|
  Group([ (1,2,3,4,5), (2,3,5,4), (1,6)(3,4) ])
  !gapprompt@gap>| !gapinput@P := PresentationViaCosetTable( G );|
  <presentation with 3 gens and 6 rels of total length 28>
  !gapprompt@gap>| !gapinput@TzPrintGenerators( P );|
  #I  1.  f1   11 occurrences
  #I  2.  f2   10 occurrences
  #I  3.  f3   7 occurrences   involution
\end{Verbatim}
 }

 

\subsection{\textcolor{Chapter }{TzPrintRelators}}
\logpage{[ 48, 4, 2 ]}\nobreak
\hyperdef{L}{X821B63DD82894443}{}
{\noindent\textcolor{FuncColor}{$\triangleright$\enspace\texttt{TzPrintRelators({\mdseries\slshape P[, list]})\index{TzPrintRelators@\texttt{TzPrintRelators}}
\label{TzPrintRelators}
}\hfill{\scriptsize (function)}}\\


 prints the relators of the given Tietze presentation \mbox{\texttt{\mdseries\slshape P}}. The optional second argument \mbox{\texttt{\mdseries\slshape list}} can be used to specify the numbers of the relators to be printed. Default: all
relators are printed. 
\begin{Verbatim}[commandchars=!@|,fontsize=\small,frame=single,label=Example]
  !gapprompt@gap>| !gapinput@TzPrintRelators( P );|
  #I  1. f3^2
  #I  2. f2^4
  #I  3. (f2*f3)^2
  #I  4. f1^5
  #I  5. f1^2*f2*f1*f2^-1
  #I  6. f1*f2^-2*f3*f1*f3*f1^-1*f3
\end{Verbatim}
 }

 

\subsection{\textcolor{Chapter }{TzPrintLengths}}
\logpage{[ 48, 4, 3 ]}\nobreak
\hyperdef{L}{X852C52C37FAAB7DD}{}
{\noindent\textcolor{FuncColor}{$\triangleright$\enspace\texttt{TzPrintLengths({\mdseries\slshape P})\index{TzPrintLengths@\texttt{TzPrintLengths}}
\label{TzPrintLengths}
}\hfill{\scriptsize (function)}}\\


 prints just a list of all relator lengths of the given presentation \mbox{\texttt{\mdseries\slshape P}}. 
\begin{Verbatim}[commandchars=!@|,fontsize=\small,frame=single,label=Example]
  !gapprompt@gap>| !gapinput@TzPrintLengths( P );|
  [ 2, 4, 4, 5, 5, 8 ]
\end{Verbatim}
 }

 

\subsection{\textcolor{Chapter }{TzPrintStatus}}
\logpage{[ 48, 4, 4 ]}\nobreak
\hyperdef{L}{X7D7B3F46865443E4}{}
{\noindent\textcolor{FuncColor}{$\triangleright$\enspace\texttt{TzPrintStatus({\mdseries\slshape P[, norepeat]})\index{TzPrintStatus@\texttt{TzPrintStatus}}
\label{TzPrintStatus}
}\hfill{\scriptsize (function)}}\\


 is an internal function which is used by the Tietze transformation routines to
print the number of generators, the number of relators, and the total length
of all relators in the given Tietze presentation \mbox{\texttt{\mdseries\slshape P}}. If \mbox{\texttt{\mdseries\slshape norepeat}} is specified as \texttt{true}, the printing is suppressed if none of the three values has changed since the
last call. 
\begin{Verbatim}[commandchars=!@|,fontsize=\small,frame=single,label=Example]
  !gapprompt@gap>| !gapinput@TzPrintStatus( P );|
  #I  there are 3 generators and 6 relators of total length 28
\end{Verbatim}
 }

 

\subsection{\textcolor{Chapter }{TzPrintPresentation}}
\logpage{[ 48, 4, 5 ]}\nobreak
\hyperdef{L}{X85F8DAE27F06C32B}{}
{\noindent\textcolor{FuncColor}{$\triangleright$\enspace\texttt{TzPrintPresentation({\mdseries\slshape P})\index{TzPrintPresentation@\texttt{TzPrintPresentation}}
\label{TzPrintPresentation}
}\hfill{\scriptsize (function)}}\\


 prints the generators and the relators of a Tietze presentation. In fact, it
is an abbreviation for the successive call of the three commands \texttt{TzPrintGenerators} (\ref{TzPrintGenerators}), \texttt{TzPrintRelators} (\ref{TzPrintRelators}), and \texttt{TzPrintStatus} (\ref{TzPrintStatus}), each with the presentation \mbox{\texttt{\mdseries\slshape P}} as only argument. }

 

\subsection{\textcolor{Chapter }{TzPrint}}
\logpage{[ 48, 4, 6 ]}\nobreak
\hyperdef{L}{X7CA8BA51802655FC}{}
{\noindent\textcolor{FuncColor}{$\triangleright$\enspace\texttt{TzPrint({\mdseries\slshape P[, list]})\index{TzPrint@\texttt{TzPrint}}
\label{TzPrint}
}\hfill{\scriptsize (function)}}\\


 prints the current generators of the given presentation \mbox{\texttt{\mdseries\slshape P}}, and prints the relators of \mbox{\texttt{\mdseries\slshape P}} as Tietze words (without converting them back to abstract words as the
functions \texttt{TzPrintRelators} (\ref{TzPrintRelators}) and \texttt{TzPrintPresentation} (\ref{TzPrintPresentation}) do). The optional second argument can be used to specify the numbers of the
relators to be printed. Default: all relators are printed. 
\begin{Verbatim}[commandchars=!@|,fontsize=\small,frame=single,label=Example]
  !gapprompt@gap>| !gapinput@TzPrint( P );|
  #I  generators: [ f1, f2, f3 ]
  #I  relators:
  #I  1.  2  [ 3, 3 ]
  #I  2.  4  [ 2, 2, 2, 2 ]
  #I  3.  4  [ 2, 3, 2, 3 ]
  #I  4.  5  [ 1, 1, 1, 1, 1 ]
  #I  5.  5  [ 1, 1, 2, 1, -2 ]
  #I  6.  8  [ 1, -2, -2, 3, 1, 3, -1, 3 ]
\end{Verbatim}
 }

 

\subsection{\textcolor{Chapter }{TzPrintPairs}}
\logpage{[ 48, 4, 7 ]}\nobreak
\hyperdef{L}{X82F6B0EE7C7C7901}{}
{\noindent\textcolor{FuncColor}{$\triangleright$\enspace\texttt{TzPrintPairs({\mdseries\slshape P[, n]})\index{TzPrintPairs@\texttt{TzPrintPairs}}
\label{TzPrintPairs}
}\hfill{\scriptsize (function)}}\\


 prints the \mbox{\texttt{\mdseries\slshape n}} most often occurring relator subwords of the form $a b$, where $a$ and $b$ are different generators or inverses of generators, together with the number
of their occurrences. The default value of \mbox{\texttt{\mdseries\slshape n}} is 10. A value \mbox{\texttt{\mdseries\slshape n}} = 0 is interpreted as \texttt{infinity} (\ref{infinity}). 

 The function \texttt{TzPrintPairs} is useful in the context of Tietze transformations which introduce new
generators by substituting words in the current generators (see \ref{Tietze Transformations that introduce new Generators}). It gives some evidence for an appropriate choice of a word of length 2 to
be substituted. 
\begin{Verbatim}[commandchars=!@|,fontsize=\small,frame=single,label=Example]
  !gapprompt@gap>| !gapinput@TzPrintPairs( P, 3 );|
  #I  1.  3  occurrences of  f2^-1 * f3
  #I  2.  2  occurrences of  f2 * f3
  #I  3.  2  occurrences of  f1^-1 * f3
\end{Verbatim}
 }

 }

  
\section{\textcolor{Chapter }{Changing Presentations}}\label{Changing Presentations}
\logpage{[ 48, 5, 0 ]}
\hyperdef{L}{X82455E5885D73FFF}{}
{
  The functions described in this section may be used to change a presentation.
Note, however, that in general they do not perform Tietze transformations
because they change or may change the isomorphism type of the group defined by
the presentation. 

\subsection{\textcolor{Chapter }{AddGenerator}}
\logpage{[ 48, 5, 1 ]}\nobreak
\hyperdef{L}{X7F632A6D8685855D}{}
{\noindent\textcolor{FuncColor}{$\triangleright$\enspace\texttt{AddGenerator({\mdseries\slshape P})\index{AddGenerator@\texttt{AddGenerator}}
\label{AddGenerator}
}\hfill{\scriptsize (function)}}\\


 extends the presentation \mbox{\texttt{\mdseries\slshape P}} by a new generator. 

 Let $i$ be the smallest positive integer which has not yet been used as a generator
number in the given presentation. \texttt{AddGenerator} defines a new abstract generator $x_i$ with the name \texttt{"{\textunderscore}x}$i$\texttt{"} and adds it to the list of generators of \mbox{\texttt{\mdseries\slshape P}}. 

 You may access the generator $x_i$ by typing \mbox{\texttt{\mdseries\slshape P}}\texttt{!.}$i$. However, this is only practicable if you are running an interactive job
because you have to know the value of $i$. Hence the proper way to access the new generator is to write \texttt{Last(GeneratorsOfPresentation(P))}. 
\begin{Verbatim}[commandchars=@|B,fontsize=\small,frame=single,label=Example]
  @gapprompt|gap>B @gapinput|G := PerfectGroup(IsFpGroup, 120 );;B
  @gapprompt|gap>B @gapinput|H := Subgroup( G, [ G.1^G.2, G.3 ] );;B
  @gapprompt|gap>B @gapinput|P := PresentationSubgroup( G, H );B
  <presentation with 4 gens and 7 rels of total length 21>
  @gapprompt|gap>B @gapinput|AddGenerator( P );B
  #I  now the presentation has 5 generators, the new generator is _x7
  @gapprompt|gap>B @gapinput|gens := GeneratorsOfPresentation( P );B
  [ _x1, _x2, _x4, _x5, _x7 ]
  @gapprompt|gap>B @gapinput|gen := Last(gens);B
  _x7
  @gapprompt|gap>B @gapinput|gen = P!.7;B
  true
\end{Verbatim}
 }

 

\subsection{\textcolor{Chapter }{TzNewGenerator}}
\logpage{[ 48, 5, 2 ]}\nobreak
\hyperdef{L}{X83A5667086FD538A}{}
{\noindent\textcolor{FuncColor}{$\triangleright$\enspace\texttt{TzNewGenerator({\mdseries\slshape P})\index{TzNewGenerator@\texttt{TzNewGenerator}}
\label{TzNewGenerator}
}\hfill{\scriptsize (function)}}\\


 is an internal function which defines a new abstract generator and adds it to
the presentation \mbox{\texttt{\mdseries\slshape P}}. It is called by \texttt{AddGenerator} (\ref{AddGenerator}) and by several Tietze transformation commands. As it does not know which
global lists have to be kept consistent, you should not call it. Instead, you
should call the function \texttt{AddGenerator} (\ref{AddGenerator}), if needed. }

 

\subsection{\textcolor{Chapter }{AddRelator}}
\logpage{[ 48, 5, 3 ]}\nobreak
\hyperdef{L}{X78D1BCE67FA852D8}{}
{\noindent\textcolor{FuncColor}{$\triangleright$\enspace\texttt{AddRelator({\mdseries\slshape P, word})\index{AddRelator@\texttt{AddRelator}}
\label{AddRelator}
}\hfill{\scriptsize (function)}}\\


 adds the relator \mbox{\texttt{\mdseries\slshape word}} to the presentation \mbox{\texttt{\mdseries\slshape P}}, probably changing the group defined by \mbox{\texttt{\mdseries\slshape P}}. \mbox{\texttt{\mdseries\slshape word}} must be an abstract word in the generators of \mbox{\texttt{\mdseries\slshape P}}. }

 

\subsection{\textcolor{Chapter }{RemoveRelator}}
\logpage{[ 48, 5, 4 ]}\nobreak
\hyperdef{L}{X7B11E89E78A22EBF}{}
{\noindent\textcolor{FuncColor}{$\triangleright$\enspace\texttt{RemoveRelator({\mdseries\slshape P, n})\index{RemoveRelator@\texttt{RemoveRelator}}
\label{RemoveRelator}
}\hfill{\scriptsize (function)}}\\


 removes the \mbox{\texttt{\mdseries\slshape n}}\texttt{\symbol{45}}th relator from the presentation \mbox{\texttt{\mdseries\slshape P}}, probably changing the group defined by \mbox{\texttt{\mdseries\slshape P}}. }

 }

  
\section{\textcolor{Chapter }{Tietze Transformations}}\label{Tietze Transformations}
\logpage{[ 48, 6, 0 ]}
\hyperdef{L}{X829B634286471AB7}{}
{
  The commands in this section can be used to modify a presentation by Tietze
transformations. 

 In general, the aim of such modifications will be to \emph{simplify} the given presentation, i.e., to reduce the number of generators and the
number of relators without increasing too much the sum of all relator lengths
which we will call the \emph{total length} of the presentation. Depending on the concrete presentation under
investigation one may end up with a nice, short presentation or with a very
huge one. 

 Unfortunately there is no algorithm which could be applied to find the
shortest presentation which can be obtained by Tietze transformations from a
given one. Therefore, what \textsf{GAP} offers are some lower\texttt{\symbol{45}}level Tietze transformation commands
and, in addition, some higher\texttt{\symbol{45}}level commands which apply
the lower\texttt{\symbol{45}}level ones in a kind of default strategy which of
course cannot be the optimal choice for all presentations. 

 The design of these commands follows closely the concept of the ANU Tietze
transformation program \cite{Hav69} and its later revisions (see \cite{HKRR84}, \cite{Rob88}). 

\subsection{\textcolor{Chapter }{TzGo}}
\logpage{[ 48, 6, 1 ]}\nobreak
\hyperdef{L}{X7C4A30328224C466}{}
{\noindent\textcolor{FuncColor}{$\triangleright$\enspace\texttt{TzGo({\mdseries\slshape P[, silent]})\index{TzGo@\texttt{TzGo}}
\label{TzGo}
}\hfill{\scriptsize (function)}}\\


 automatically performs suitable Tietze transformations of the given
presentation \mbox{\texttt{\mdseries\slshape P}}. It is perhaps the most convenient one among the interactive Tietze
transformation commands. It offers a kind of default strategy which, in
general, saves you from explicitly calling the lower\texttt{\symbol{45}}level
commands it involves. 

 If \mbox{\texttt{\mdseries\slshape silent}} is specified as \texttt{true}, the printing of the status line by \texttt{TzGo} is suppressed if the Tietze option \texttt{printLevel} (see \ref{Tietze Options}) has a value less than $2$. }

 

\subsection{\textcolor{Chapter }{SimplifyPresentation}}
\logpage{[ 48, 6, 2 ]}\nobreak
\hyperdef{L}{X78C3D23387DAC35A}{}
{\noindent\textcolor{FuncColor}{$\triangleright$\enspace\texttt{SimplifyPresentation({\mdseries\slshape P})\index{SimplifyPresentation@\texttt{SimplifyPresentation}}
\label{SimplifyPresentation}
}\hfill{\scriptsize (function)}}\\


 \texttt{SimplifyPresentation} is a synonym for \texttt{TzGo} (\ref{TzGo}). 
\begin{Verbatim}[commandchars=!@|,fontsize=\small,frame=single,label=Example]
  !gapprompt@gap>| !gapinput@F2 := FreeGroup( "a", "b" );;|
  !gapprompt@gap>| !gapinput@G := F2 / [ F2.1^9, F2.2^2, (F2.1*F2.2)^4, (F2.1^2*F2.2)^3 ];;|
  !gapprompt@gap>| !gapinput@a := G.1;; b := G.2;;|
  !gapprompt@gap>| !gapinput@H := Subgroup( G, [ (a*b)^2, (a^-1*b)^2 ] );;|
  !gapprompt@gap>| !gapinput@Index( G, H );|
  408
  !gapprompt@gap>| !gapinput@P := PresentationSubgroup( G, H );|
  <presentation with 8 gens and 36 rels of total length 111>
  !gapprompt@gap>| !gapinput@PrimaryGeneratorWords( P );|
  [ b, a*b*a ]
  !gapprompt@gap>| !gapinput@TzOptions( P ).protected := 2;|
  2
  !gapprompt@gap>| !gapinput@TzOptions( P ).printLevel := 2;|
  2
  !gapprompt@gap>| !gapinput@SimplifyPresentation( P );|
  #I  eliminating _x7 = _x5^-1
  #I  eliminating _x5 = _x4
  #I  eliminating _x18 = _x3
  #I  eliminating _x8 = _x3
  #I  there are 4 generators and 8 relators of total length 21
  #I  there are 4 generators and 7 relators of total length 18
  #I  eliminating _x4 = _x3^-1*_x2^-1
  #I  eliminating _x3 = _x2*_x1^-1
  #I  there are 2 generators and 4 relators of total length 14
  #I  there are 2 generators and 4 relators of total length 13
  #I  there are 2 generators and 3 relators of total length 9
  !gapprompt@gap>| !gapinput@TzPrintRelators( P );|
  #I  1. _x1^2
  #I  2. _x2^3
  #I  3. (_x2*_x1)^2
\end{Verbatim}
 

 Roughly speaking, \texttt{TzGo} (\ref{TzGo}) consists of a loop over a procedure which involves two phases: In the \emph{search phase} it calls \texttt{TzSearch} (\ref{TzSearch}) and \texttt{TzSearchEqual} (\ref{TzSearchEqual}) described below which try to reduce the relator lengths by substituting common
subwords of relators, in the \emph{elimination phase} it calls the command \texttt{TzEliminate} (\ref{TzEliminate:for a presentation (and a generator)}) described below (or, more precisely, a subroutine of \texttt{TzEliminate} (\ref{TzEliminate:for a presentation (and a generator)}) in order to save some administrative overhead) which tries to eliminate
generators that can be expressed as words in the remaining generators. 

 If \texttt{TzGo} (\ref{TzGo}) succeeds in reducing the number of generators, the number of relators, or the
total length of all relators, it displays the new status before returning
(provided that you did not set the print level to zero). However, it does not
provide any output if all these three values have remained unchanged, even if
the command \texttt{TzSearchEqual} (\ref{TzSearchEqual}) involved has changed the presentation such that another call of \texttt{TzGo} (\ref{TzGo}) might provide further progress. Hence, in such a case it makes sense to repeat
the call of the command for several times (or to call the command \texttt{TzGoGo} (\ref{TzGoGo}) instead). }

 

\subsection{\textcolor{Chapter }{TzGoGo}}
\logpage{[ 48, 6, 3 ]}\nobreak
\hyperdef{L}{X801D3D8984E1CA55}{}
{\noindent\textcolor{FuncColor}{$\triangleright$\enspace\texttt{TzGoGo({\mdseries\slshape P})\index{TzGoGo@\texttt{TzGoGo}}
\label{TzGoGo}
}\hfill{\scriptsize (function)}}\\


 calls the command \texttt{TzGo} (\ref{TzGo}) again and again until it does not reduce the presentation any more. 

 The result of the Tietze transformations can be affected substantially by the
options parameters (see \ref{Tietze Options}). To demonstrate the effect of the \texttt{eliminationsLimit} parameter, we will give an example in which we handle a subgroup of index 240
in a group of order 40320 given by a presentation due to B.{\nobreakspace}H.
Neumann. First we construct a presentation of the subgroup, and then we apply
to it the command \texttt{TzGoGo} for different values of the parameter \texttt{eliminationsLimit} (including the default value 100). In fact, we also alter the \texttt{printLevel} parameter, but this is only done in order to suppress most of the output. In
all cases the resulting presentations cannot be improved any more by applying
the command \texttt{TzGoGo} again, i.e., they are the best results which we can get without substituting
new generators. 

 
\begin{Verbatim}[commandchars=!@|,fontsize=\small,frame=single,label=Example]
  !gapprompt@gap>| !gapinput@F3 := FreeGroup( "a", "b", "c" );;|
  !gapprompt@gap>| !gapinput@G := F3 / [ F3.1^3, F3.2^3, F3.3^3, (F3.1*F3.2)^5,|
  !gapprompt@>| !gapinput@(F3.1^-1*F3.2)^5, (F3.1*F3.3)^4, (F3.1*F3.3^-1)^4,|
  !gapprompt@>| !gapinput@F3.1*F3.2^-1*F3.1*F3.2*F3.3^-1*F3.1*F3.3*F3.1*F3.3^-1,|
  !gapprompt@>| !gapinput@(F3.2*F3.3)^3, (F3.2^-1*F3.3)^4 ];;|
  !gapprompt@gap>| !gapinput@a := G.1;; b := G.2;; c := G.3;;|
  !gapprompt@gap>| !gapinput@H := Subgroup( G, [ a, c ] );;|
  !gapprompt@gap>| !gapinput@for i in [ 61, 62, 63, 90, 97 ] do|
  !gapprompt@>| !gapinput@Pi := PresentationSubgroup( G, H );|
  !gapprompt@>| !gapinput@TzOptions( Pi ).eliminationsLimit := i;|
  !gapprompt@>| !gapinput@Print("#I eliminationsLimit set to ",i,"\n");|
  !gapprompt@>| !gapinput@TzOptions( Pi ).printLevel := 0;|
  !gapprompt@>| !gapinput@TzGoGo( Pi );|
  !gapprompt@>| !gapinput@TzPrintStatus( Pi );|
  !gapprompt@>| !gapinput@od;|
  #I eliminationsLimit set to 61
  #I  there are 2 generators and 104 relators of total length 7012
  #I eliminationsLimit set to 62
  #I  there are 2 generators and 7 relators of total length 56
  #I eliminationsLimit set to 63
  #I  there are 3 generators and 97 relators of total length 5998
  #I eliminationsLimit set to 90
  #I  there are 3 generators and 11 relators of total length 68
  #I eliminationsLimit set to 97
  #I  there are 4 generators and 109 relators of total length 3813
\end{Verbatim}
 

 Similarly, we demonstrate the influence of the \texttt{saveLimit} parameter by just continuing the preceding example for some different values
of the \texttt{saveLimit} parameter (including its default value 10), but without changing the \texttt{eliminationsLimit} parameter which keeps its default value 100. 

 
\begin{Verbatim}[commandchars=!@|,fontsize=\small,frame=single,label=Example]
  !gapprompt@gap>| !gapinput@for i in [ 7 .. 11 ] do|
  !gapprompt@>| !gapinput@Pi := PresentationSubgroup( G, H );|
  !gapprompt@>| !gapinput@TzOptions( Pi ).saveLimit := i;|
  !gapprompt@>| !gapinput@Print( "#I saveLimit set to ", i, "\n" );|
  !gapprompt@>| !gapinput@TzOptions( Pi ).printLevel := 0;|
  !gapprompt@>| !gapinput@TzGoGo( Pi );|
  !gapprompt@>| !gapinput@TzPrintStatus( Pi );|
  !gapprompt@>| !gapinput@od;|
  #I saveLimit set to 7
  #I  there are 3 generators and 99 relators of total length 2713
  #I saveLimit set to 8
  #I  there are 2 generators and 103 relators of total length 11982
  #I saveLimit set to 9
  #I  there are 2 generators and 6 relators of total length 41
  #I saveLimit set to 10
  #I  there are 3 generators and 118 relators of total length 13713
  #I saveLimit set to 11
  #I  there are 3 generators and 11 relators of total length 58
\end{Verbatim}
 }

 }

  
\section{\textcolor{Chapter }{Elementary Tietze Transformations}}\label{Elementary Tietze Transformations}
\logpage{[ 48, 7, 0 ]}
\hyperdef{L}{X7D19E30080290FC7}{}
{
  
\subsection{\textcolor{Chapter }{TzEliminate}}\logpage{[ 48, 7, 1 ]}
\hyperdef{L}{X85989AF886EC2BF6}{}
{
\noindent\textcolor{FuncColor}{$\triangleright$\enspace\texttt{TzEliminate({\mdseries\slshape P[, gen]})\index{TzEliminate@\texttt{TzEliminate}!for a presentation (and a generator)}
\label{TzEliminate:for a presentation (and a generator)}
}\hfill{\scriptsize (function)}}\\
\noindent\textcolor{FuncColor}{$\triangleright$\enspace\texttt{TzEliminate({\mdseries\slshape P[, n]})\index{TzEliminate@\texttt{TzEliminate}!for a presentation (and an integer)}
\label{TzEliminate:for a presentation (and an integer)}
}\hfill{\scriptsize (function)}}\\


 tries to eliminate a generator from a presentation \mbox{\texttt{\mdseries\slshape P}} via Tietze transformations. 

 Any relator which contains some generator just once can be used to substitute
that generator by a word in the remaining generators. If such generators and
relators exist, then \texttt{TzEliminate} chooses a generator for which the product of its number of occurrences and the
length of the substituting word is minimal, and then it eliminates this
generator from the presentation, provided that the resulting total length of
the relators does not exceed the associated Tietze option parameter \texttt{spaceLimit} (see \ref{Tietze Options}). The default value of that parameter is \texttt{infinity} (\ref{infinity}), but you may alter it appropriately. 

 If a generator \mbox{\texttt{\mdseries\slshape gen}} has been specified, \texttt{TzEliminate} eliminates it if possible, i. e. if there is a relator in which \mbox{\texttt{\mdseries\slshape gen}} occurs just once. If no second argument has been specified, \texttt{TzEliminate} eliminates some appropriate generator if possible and if the resulting total
length of the relators will not exceed the Tietze options parameter \texttt{lengthLimit}. 

 If an integer \mbox{\texttt{\mdseries\slshape n}} has been specified, \texttt{TzEliminate} tries to eliminate up to \mbox{\texttt{\mdseries\slshape n}} generators. Note that the calls \texttt{TzEliminate(\mbox{\texttt{\mdseries\slshape P}})} and \texttt{TzEliminate(\mbox{\texttt{\mdseries\slshape P}},1)} are equivalent. }

 

\subsection{\textcolor{Chapter }{TzSearch}}
\logpage{[ 48, 7, 2 ]}\nobreak
\hyperdef{L}{X7DF4BBDF839643DD}{}
{\noindent\textcolor{FuncColor}{$\triangleright$\enspace\texttt{TzSearch({\mdseries\slshape P})\index{TzSearch@\texttt{TzSearch}}
\label{TzSearch}
}\hfill{\scriptsize (function)}}\\


 searches for relator subwords which, in some relator, have a complement of
shorter length and which occur in other relators, too, and uses them to reduce
these other relators. 

 The idea is to find pairs of relators $r_1$ and $r_2$ of length $l_1$ and $l_2$, respectively, such that $l_1 \leq l_2$ and $r_1$ and $r_2$ coincide (possibly after inverting or conjugating one of them) in some maximal
subword $w$ of length greater than $l_1/2$, and then to substitute each copy of $w$ in $r_2$ by the inverse complement of $w$ in $r_1$. 

 Two of the Tietze option parameters which are listed in section \ref{Tietze Options} may strongly influence the performance and the results of the command \texttt{TzSearch}. These are the parameters \texttt{saveLimit} and \texttt{searchSimultaneous}. The first of them has the following effect: 

 When \texttt{TzSearch} has finished its main loop over all relators, then, in general, there are
relators which have changed and hence should be handled again in another run
through the whole procedure. However, experience shows that it really does not
pay to continue this way until no more relators change. Therefore, \texttt{TzSearch} starts a new loop only if the loop just finished has reduced the total length
of the relators by at least \texttt{saveLimit} per cent. 

 The default value of \texttt{saveLimit} is 10 per cent. 

 To understand the effect of the option \texttt{searchSimultaneous}, we have to look in more detail at how \texttt{TzSearch} proceeds: 

 First, it sorts the list of relators by increasing lengths. Then it performs a
loop over this list. In each step of this loop, the current relator is treated
as \emph{short relator} $r_1$, and a subroutine is called which loops over the succeeding relators,
treating them as \emph{long relators} $r_2$ and performing the respective comparisons and substitutions. 

 As this subroutine performs a very expensive process, it has been implemented
as a C routine in the \textsf{GAP} kernel. For the given relator $r_1$ of length $l_1$ it first determines the \emph{minimal match length} $l$ which is $l_1/2+1$, if $l_1$ is even, or $(l_1+1)/2$, otherwise. Then it builds up a hash list for all subwords of length $l$ occurring in the conjugates of $r_1$ or $r_1^{{-1}}$, and finally it loops over all long relators $r_2$ and compares the hash values of their subwords of length $l$ against this list. A comparison of subwords which is much more expensive is
only done if a hash match has been found. 

 To improve the efficiency of this process we allow the subroutine to handle
several short relators simultaneously provided that they have the same minimal
match length. If, for example, it handles $n$ short relators simultaneously, then you save $n - 1$ loops over the long relators $r_2$, but you pay for it by additional fruitless subword comparisons. In general,
you will not get the best performance by always choosing the maximal possible
number of short relators to be handled simultaneously. In fact, the optimal
choice of the number will depend on the concrete presentation under
investigation. You can use the parameter \texttt{searchSimultaneous} to prescribe an upper bound for the number of short relators to be handled
simultaneously. 

 The default value of \texttt{searchSimultaneous} is 20. }

 

\subsection{\textcolor{Chapter }{TzSearchEqual}}
\logpage{[ 48, 7, 3 ]}\nobreak
\hyperdef{L}{X87F7A87A7ACF2445}{}
{\noindent\textcolor{FuncColor}{$\triangleright$\enspace\texttt{TzSearchEqual({\mdseries\slshape P})\index{TzSearchEqual@\texttt{TzSearchEqual}}
\label{TzSearchEqual}
}\hfill{\scriptsize (function)}}\\


 searches for Tietze relator subwords which, in some relator, have a complement
of equal length and which occur in other relators, too, and uses them to
modify these other relators. 

 The idea is to find pairs of relators $r_1$ and $r_2$ of length $l_1$ and $l_2$, respectively, such that $l_1$ is even, $l_1 \leq l_2$, and $r_1$ and $r_2$ coincide (possibly after inverting or conjugating one of them) in some maximal
subword $w$ of length at least $l_1/2$. Let $l$ be the length of $w$. Then, if $l > l_1/2$, the pair is handled as in \texttt{TzSearch} (\ref{TzSearch}). Otherwise, if $l = l_1/2$, then \texttt{TzSearchEqual} substitutes each copy of $w$ in $r_2$ by the inverse complement of $w$ in $r_1$. 

 The Tietze option parameter \texttt{searchSimultaneous} is used by \texttt{TzSearchEqual} in the same way as described for \texttt{TzSearch} (\ref{TzSearch}). However, \texttt{TzSearchEqual} does not use the parameter \texttt{saveLimit}: The loop over the relators is executed exactly once. }

 

\subsection{\textcolor{Chapter }{TzFindCyclicJoins}}
\logpage{[ 48, 7, 4 ]}\nobreak
\hyperdef{L}{X80D31A0F7C2A51BD}{}
{\noindent\textcolor{FuncColor}{$\triangleright$\enspace\texttt{TzFindCyclicJoins({\mdseries\slshape P})\index{TzFindCyclicJoins@\texttt{TzFindCyclicJoins}}
\label{TzFindCyclicJoins}
}\hfill{\scriptsize (function)}}\\


 searches for power and commutator relators in order to find pairs of
generators which generate a common cyclic subgroup. It uses these pairs to
introduce new relators, but it does not introduce any new generators as is
done by \texttt{TzSubstituteCyclicJoins} (\ref{TzSubstituteCyclicJoins}). 

 More precisely: \texttt{TzFindCyclicJoins} searches for pairs of generators $a$ and $b$ such that (possibly after inverting or conjugating some relators) the set of
relators contains the commutator $[a,b]$, a power $a^n$, and a product of the form $a^s b^t$ with $s$ prime to $n$. For each such pair, \texttt{TzFindCyclicJoins} uses the Euclidean algorithm to express $a$ as a power of $b$, and then it eliminates $a$. }

 }

  
\section{\textcolor{Chapter }{Tietze Transformations that introduce new Generators}}\label{Tietze Transformations that introduce new Generators}
\logpage{[ 48, 8, 0 ]}
\hyperdef{L}{X7D2FACCF79F57040}{}
{
  Some of the Tietze transformation commands listed so far may eliminate
generators and hence change the given presentation to a presentation on a
subset of the given set of generators, but they all do \emph{not} introduce new generators. However, sometimes there will be the need to
substitute certain words as new generators in order to improve a presentation.
Therefore \textsf{GAP} offers the two commands \texttt{TzSubstitute} (\ref{TzSubstitute:for a presentation and a word}) and \texttt{TzSubstituteCyclicJoins} (\ref{TzSubstituteCyclicJoins}) which introduce new generators. 
\subsection{\textcolor{Chapter }{TzSubstitute}}\logpage{[ 48, 8, 1 ]}
\hyperdef{L}{X846DB23E8236FF8A}{}
{
\noindent\textcolor{FuncColor}{$\triangleright$\enspace\texttt{TzSubstitute({\mdseries\slshape P, word})\index{TzSubstitute@\texttt{TzSubstitute}!for a presentation and a word}
\label{TzSubstitute:for a presentation and a word}
}\hfill{\scriptsize (function)}}\\
\noindent\textcolor{FuncColor}{$\triangleright$\enspace\texttt{TzSubstitute({\mdseries\slshape P[, n[, eliminate]]})\index{TzSubstitute@\texttt{TzSubstitute}!for a presentation (and an integer and 0/1/2)}
\label{TzSubstitute:for a presentation (and an integer and 0/1/2)}
}\hfill{\scriptsize (function)}}\\


 In the first form \texttt{TzSubstitute} expects \mbox{\texttt{\mdseries\slshape P}} to be a presentation and \mbox{\texttt{\mdseries\slshape word}} to be either an abstract word or a Tietze word in the generators of \mbox{\texttt{\mdseries\slshape P}}. It substitutes the given word as a new generator of \mbox{\texttt{\mdseries\slshape P}}. This is done as follows: First, \texttt{TzSubstitute} creates a new abstract generator, $g$ say, and adds it to the presentation, then it adds a new relator $g^{{-1}} \cdot \mbox{\texttt{\mdseries\slshape word}}$. 

 In its second form, \texttt{TzSubstitute} substitutes a squarefree word of length 2 as a new generator and then
eliminates a generator from the extended generator list. We will describe this
process in more detail below. 

 The parameters \mbox{\texttt{\mdseries\slshape n}} and \mbox{\texttt{\mdseries\slshape eliminate}} are optional. If you specify arguments for them, then \mbox{\texttt{\mdseries\slshape n}} is expected to be a positive integer, and \mbox{\texttt{\mdseries\slshape eliminate}} is expected to be 0, 1, or 2. The default values are \mbox{\texttt{\mdseries\slshape n}} $= 1$ and \mbox{\texttt{\mdseries\slshape eliminate}} $= 0$. 

 \texttt{TzSubstitute} first determines the \mbox{\texttt{\mdseries\slshape n}} most frequently occurring relator subwords of the form $g_1 g_2$, where $g_1$ and $g_2$ are different generators or their inverses, and sorts them by decreasing
numbers of occurrences. 

 Let $a b$ be the last word in that list, and let $i$ be the smallest positive integer which has not yet been used as a generator
number in the presentation \mbox{\texttt{\mdseries\slshape P}} so far. \texttt{TzSubstitute} defines a new abstract generator $x_i$ named \texttt{"{\textunderscore}x\mbox{\texttt{\mdseries\slshape i}}"} and adds it to \mbox{\texttt{\mdseries\slshape P}} (see \texttt{AddGenerator} (\ref{AddGenerator})). Then it adds the word $x_i^{{-1}} a b$ as a new relator to \mbox{\texttt{\mdseries\slshape P}} and replaces all occurrences of $a b$ in the relators by $x_i$. Finally, it eliminates some suitable generator from \mbox{\texttt{\mdseries\slshape P}}. 

 The choice of the generator to be eliminated depends on the actual value of
the parameter \mbox{\texttt{\mdseries\slshape eliminate}}: 

 If \mbox{\texttt{\mdseries\slshape eliminate}} is zero, \texttt{TzSubstitute} just calls the function \texttt{TzEliminate} (\ref{TzEliminate:for a presentation (and a generator)}). So it may happen that it is the just introduced generator $x_i$ which now is deleted again so that you don't get any remarkable progress in
simplifying your presentation. On the first glance this does not look
reasonable, but it is a consequence of the request that a call of \texttt{TzSubstitute} with \mbox{\texttt{\mdseries\slshape eliminate}} = 0 must not increase the total length of the relators. 

 Otherwise, if \mbox{\texttt{\mdseries\slshape eliminate}} is 1 or 2, \texttt{TzSubstitute} eliminates the respective factor of the substituted word $a b$, i. e., it eliminates $a$ if \mbox{\texttt{\mdseries\slshape eliminate}} = 1 or $b$ if \mbox{\texttt{\mdseries\slshape eliminate}} = 2. In this case, it may happen that the total length of the relators
increases, but sometimes such an intermediate extension is the only way to
finally reduce a given presentation. 

 There is still another property of the command \texttt{TzSubstitute} which should be mentioned. If, for instance, \texttt{word} is an abstract word, a call 

 
\begin{Verbatim}[commandchars=!@|,fontsize=\small,frame=single,label=Example]
  TzSubstitute( P, word );
\end{Verbatim}
 

 is more or less equivalent to 

 
\begin{Verbatim}[commandchars=!@|,fontsize=\small,frame=single,label=Example]
  AddGenerator( P );
  g := Last(GeneratorsOfPresentation(P));
  AddRelator( P, g^-1 * word );
\end{Verbatim}
 

 However, there is a difference: If you are tracing generator images and
preimages of \mbox{\texttt{\mdseries\slshape P}} through the Tietze transformations applied to \mbox{\texttt{\mdseries\slshape P}} (see \ref{Tracing generator images through Tietze transformations}), then \texttt{TzSubstitute}, as a Tietze transformation of \mbox{\texttt{\mdseries\slshape P}}, will update and save the respective lists, whereas a call of the function \texttt{AddGenerator} (\ref{AddGenerator}) (which does not perform a Tietze transformation) will delete these lists and
hence terminate the tracing. 

 
\begin{Verbatim}[commandchars=!@|,fontsize=\small,frame=single,label=Example]
  !gapprompt@gap>| !gapinput@G := PerfectGroup( IsSubgroupFpGroup, 960, 1 );|
  A5 2^4
  !gapprompt@gap>| !gapinput@P := PresentationFpGroup( G );|
  <presentation with 6 gens and 21 rels of total length 84>
  !gapprompt@gap>| !gapinput@GeneratorsOfPresentation( P );|
  [ a, b, s, t, u, v ]
  !gapprompt@gap>| !gapinput@TzGoGo( P );|
  #I  there are 3 generators and 10 relators of total length 81
  #I  there are 3 generators and 10 relators of total length 80
  !gapprompt@gap>| !gapinput@TzPrintGenerators( P );|
  #I  1.  a   31 occurrences   involution
  #I  2.  b   26 occurrences
  #I  3.  t   23 occurrences   involution
  !gapprompt@gap>| !gapinput@a := GeneratorsOfPresentation( P )[1];;|
  !gapprompt@gap>| !gapinput@b := GeneratorsOfPresentation( P )[2];;|
  !gapprompt@gap>| !gapinput@TzSubstitute( P, a*b );|
  #I  now the presentation has 4 generators, the new generator is _x7
  #I  substituting new generator _x7 defined by a*b
  #I  there are 4 generators and 11 relators of total length 83
  !gapprompt@gap>| !gapinput@TzGo( P );|
  #I  there are 3 generators and 10 relators of total length 74
  !gapprompt@gap>| !gapinput@TzPrintGenerators( P );|
  #I  1.  a   23 occurrences   involution
  #I  2.  t   23 occurrences   involution
  #I  3.  _x7   28 occurrences
\end{Verbatim}
 

 As an example of an application of the command \texttt{TzSubstitute} in its second form we handle a subgroup of index 266 in the Janko group $J_1$. 

 
\begin{Verbatim}[commandchars=!@|,fontsize=\small,frame=single,label=Example]
  !gapprompt@gap>| !gapinput@F2 := FreeGroup( "a", "b" );;|
  !gapprompt@gap>| !gapinput@J1 := F2 / [ F2.1^2, F2.2^3, (F2.1*F2.2)^7,|
  !gapprompt@>| !gapinput@Comm(F2.1,F2.2)^10, Comm(F2.1,F2.2^-1*(F2.1*F2.2)^2)^6 ];;|
  !gapprompt@gap>| !gapinput@a := J1.1;; b := J1.2;;|
  !gapprompt@gap>| !gapinput@H := Subgroup ( J1, [ a, b^(a*b*(a*b^-1)^2) ] );;|
  !gapprompt@gap>| !gapinput@P := PresentationSubgroup( J1, H );|
  <presentation with 23 gens and 82 rels of total length 530>
  !gapprompt@gap>| !gapinput@TzGoGo( P );|
  #I  there are 3 generators and 47 relators of total length 1368
  #I  there are 2 generators and 46 relators of total length 3773
  #I  there are 2 generators and 46 relators of total length 2570
  !gapprompt@gap>| !gapinput@TzGoGo( P );|
  #I  there are 2 generators and 46 relators of total length 2568
  !gapprompt@gap>| !gapinput@TzGoGo( P );|
\end{Verbatim}
 

 Here we do not get any more progress without substituting a new generator. 

 
\begin{Verbatim}[commandchars=!@|,fontsize=\small,frame=single,label=Example]
  !gapprompt@gap>| !gapinput@TzSubstitute( P );|
  #I  substituting new generator _x28 defined by _x6*_x23^-1
  #I  eliminating _x28 = _x6*_x23^-1
\end{Verbatim}
 

 \textsf{GAP} cannot substitute a new generator without extending the total length, so we
have to explicitly ask for it by using the second form of the command \texttt{TzSubstitute}. Our problem is to choose appropriate values for the arguments \mbox{\texttt{\mdseries\slshape n}} and \mbox{\texttt{\mdseries\slshape eliminate}}. For this purpose it may be helpful to print out a list of the most
frequently occurring squarefree relator subwords of length 2. 

 
\begin{Verbatim}[commandchars=!@|,fontsize=\small,frame=single,label=Example]
  !gapprompt@gap>| !gapinput@TzPrintPairs( P );|
  #I  1.  504  occurrences of  _x6 * _x23^-1
  #I  2.  504  occurrences of  _x6^-1 * _x23
  #I  3.  448  occurrences of  _x6 * _x23
  #I  4.  448  occurrences of  _x6^-1 * _x23^-1
  !gapprompt@gap>| !gapinput@TzSubstitute( P, 2, 1 );|
  #I  substituting new generator _x29 defined by _x6^-1*_x23
  #I  eliminating _x6 = _x23*_x29^-1
  #I  there are 2 generators and 46 relators of total length 2867
  !gapprompt@gap>| !gapinput@TzGoGo( P );|
  #I  there are 2 generators and 45 relators of total length 2417
  #I  there are 2 generators and 45 relators of total length 2122
  !gapprompt@gap>| !gapinput@TzSubstitute( P, 1, 2 );|
  #I  substituting new generator _x30 defined by _x23*_x29^-1
  #I  eliminating _x29 = _x30^-1*_x23
  #I  there are 2 generators and 45 relators of total length 2192
  !gapprompt@gap>| !gapinput@TzGoGo( P );|
  #I  there are 2 generators and 42 relators of total length 1637
  #I  there are 2 generators and 40 relators of total length 1286
  #I  there are 2 generators and 36 relators of total length 807
  #I  there are 2 generators and 32 relators of total length 625
  #I  there are 2 generators and 22 relators of total length 369
  #I  there are 2 generators and 18 relators of total length 213
  #I  there are 2 generators and 13 relators of total length 141
  #I  there are 2 generators and 12 relators of total length 121
  #I  there are 2 generators and 10 relators of total length 101
  !gapprompt@gap>| !gapinput@TzPrintPairs( P );|
  #I  1.  19  occurrences of  _x23 * _x30^-1
  #I  2.  19  occurrences of  _x23^-1 * _x30
  #I  3.  14  occurrences of  _x23 * _x30
  #I  4.  14  occurrences of  _x23^-1 * _x30^-1
\end{Verbatim}
 

 If we save a copy of the current presentation, then later we will be able to
restart the computation from the current state. 

 
\begin{Verbatim}[commandchars=!@|,fontsize=\small,frame=single,label=Example]
  !gapprompt@gap>| !gapinput@P1 := ShallowCopy( P );|
  <presentation with 2 gens and 10 rels of total length 101>
\end{Verbatim}
 

 Just for demonstration we make an inconvenient choice: 

 
\begin{Verbatim}[commandchars=!@|,fontsize=\small,frame=single,label=Example]
  !gapprompt@gap>| !gapinput@TzSubstitute( P, 3, 1 );|
  #I  substituting new generator _x31 defined by _x23*_x30
  #I  eliminating _x23 = _x31*_x30^-1
  #I  there are 2 generators and 10 relators of total length 122
  !gapprompt@gap>| !gapinput@TzGoGo( P );|
  #I  there are 2 generators and 9 relators of total length 105
\end{Verbatim}
 

 This presentation is worse than the one we have saved, so we restart from that
presentation again. 

 
\begin{Verbatim}[commandchars=!@|,fontsize=\small,frame=single,label=Example]
  !gapprompt@gap>| !gapinput@P := ShallowCopy( P1 );|
  <presentation with 2 gens and 10 rels of total length 101>
  !gapprompt@gap>| !gapinput@TzSubstitute( P, 2, 1);|
  #I  substituting new generator _x31 defined by _x23^-1*_x30
  #I  eliminating _x23 = _x30*_x31^-1
  #I  there are 2 generators and 10 relators of total length 107
  !gapprompt@gap>| !gapinput@TzGoGo( P );|
  #I  there are 2 generators and 9 relators of total length 84
  #I  there are 2 generators and 8 relators of total length 75
  !gapprompt@gap>| !gapinput@TzSubstitute( P, 2, 1);|
  #I  substituting new generator _x32 defined by _x30^-1*_x31
  #I  eliminating _x30 = _x31*_x32^-1
  #I  there are 2 generators and 8 relators of total length 71
  !gapprompt@gap>| !gapinput@TzGoGo( P );|
  #I  there are 2 generators and 7 relators of total length 56
  #I  there are 2 generators and 5 relators of total length 36
  !gapprompt@gap>| !gapinput@TzPrintRelators( P );|
  #I  1. _x32^5
  #I  2. _x31^5
  #I  3. (_x31^-1*_x32^-1)^3
  #I  4. _x31*(_x32*_x31^-1)^2*_x32*_x31*_x32^-2
  #I  5. _x31^-1*_x32^2*(_x31*_x32^-1*_x31)^2*_x32^2
\end{Verbatim}
 }

 

\subsection{\textcolor{Chapter }{TzSubstituteCyclicJoins}}
\logpage{[ 48, 8, 2 ]}\nobreak
\hyperdef{L}{X7ADE3B437C19B94D}{}
{\noindent\textcolor{FuncColor}{$\triangleright$\enspace\texttt{TzSubstituteCyclicJoins({\mdseries\slshape P})\index{TzSubstituteCyclicJoins@\texttt{TzSubstituteCyclicJoins}}
\label{TzSubstituteCyclicJoins}
}\hfill{\scriptsize (function)}}\\


 tries to find pairs of commuting generators $a$ and $b$ such that the exponent of $a$ (i. e. the least currently known positive integer $n$ such that $a^n$ is a relator in \mbox{\texttt{\mdseries\slshape P}}) is prime to the exponent of $b$. For each such pair, their product $a b$ is substituted as a new generator, and $a$ and $b$ are eliminated. }

 }

  
\section{\textcolor{Chapter }{Tracing generator images through Tietze transformations}}\label{Tracing generator images through Tietze transformations}
\logpage{[ 48, 9, 0 ]}
\hyperdef{L}{X85E703997A0212EE}{}
{
  Any sequence of Tietze transformations applied to a presentation, starting
from some presentation $P_1$ and ending up with some presentation $P_2$, defines an isomorphism, $\varphi$ say, between the groups defined by $P_1$ and $P_2$, respectively. Sometimes it is desirable to know the images of the (old)
generators of $P_1$ or the preimages of the (new) generators of $P_2$ under $\varphi$. The \textsf{GAP} Tietze transformation functions are able to trace these images. This is not
automatically done because the involved words may grow to tremendous length,
but it will be done if you explicitly request for it by calling the function \texttt{TzInitGeneratorImages} (\ref{TzInitGeneratorImages}). 

\subsection{\textcolor{Chapter }{TzInitGeneratorImages}}
\logpage{[ 48, 9, 1 ]}\nobreak
\hyperdef{L}{X7D855FA08242898A}{}
{\noindent\textcolor{FuncColor}{$\triangleright$\enspace\texttt{TzInitGeneratorImages({\mdseries\slshape P})\index{TzInitGeneratorImages@\texttt{TzInitGeneratorImages}}
\label{TzInitGeneratorImages}
}\hfill{\scriptsize (function)}}\\


 expects \mbox{\texttt{\mdseries\slshape P}} to be a presentation. It defines the current generators to be the ``old generators'' of \mbox{\texttt{\mdseries\slshape P}} and initializes the (pre)image tracing. See \texttt{TzImagesOldGens} (\ref{TzImagesOldGens}) and \texttt{TzPreImagesNewGens} (\ref{TzPreImagesNewGens}) for details. 

 You can reinitialize the tracing of the generator images at any later state by
just calling the function \texttt{TzInitGeneratorImages} again. 

 Note: A subsequent call of the function \texttt{DecodeTree} (\ref{DecodeTree}) will imply that the images and preimages are deleted and reinitialized after
decoding the tree. 

 Moreover, if you introduce a new generator by calling the function \texttt{AddGenerator} (\ref{AddGenerator}) described in Section \ref{Changing Presentations}, this new generator cannot be traced in the old generators. Therefore \texttt{AddGenerator} (\ref{AddGenerator}) will terminate the tracing of the generator images and preimages and delete
the respective lists whenever it is called. }

 

\subsection{\textcolor{Chapter }{OldGeneratorsOfPresentation}}
\logpage{[ 48, 9, 2 ]}\nobreak
\hyperdef{L}{X7AB9A06F80FB3659}{}
{\noindent\textcolor{FuncColor}{$\triangleright$\enspace\texttt{OldGeneratorsOfPresentation({\mdseries\slshape P})\index{OldGeneratorsOfPresentation@\texttt{OldGeneratorsOfPresentation}}
\label{OldGeneratorsOfPresentation}
}\hfill{\scriptsize (function)}}\\


 assumes that \mbox{\texttt{\mdseries\slshape P}} is a presentation for which the generator images and preimages are being
traced under Tietze transformations. It returns the list of old generators of \mbox{\texttt{\mdseries\slshape P}}. }

 

\subsection{\textcolor{Chapter }{TzImagesOldGens}}
\logpage{[ 48, 9, 3 ]}\nobreak
\hyperdef{L}{X798B38F87C082C45}{}
{\noindent\textcolor{FuncColor}{$\triangleright$\enspace\texttt{TzImagesOldGens({\mdseries\slshape P})\index{TzImagesOldGens@\texttt{TzImagesOldGens}}
\label{TzImagesOldGens}
}\hfill{\scriptsize (function)}}\\


 assumes that \mbox{\texttt{\mdseries\slshape P}} is a presentation for which the generator images and preimages are being
traced under Tietze transformations. It returns a list $l$ of words in the (current) \texttt{GeneratorsOfPresentation} (\ref{GeneratorsOfPresentation}) value of \mbox{\texttt{\mdseries\slshape P}} such that the $i$\texttt{\symbol{45}}th word $l[i]$ represents the $i$\texttt{\symbol{45}}th old generator of \mbox{\texttt{\mdseries\slshape P}}, i. e., the $i$\texttt{\symbol{45}}th entry of the \texttt{OldGeneratorsOfPresentation} (\ref{OldGeneratorsOfPresentation}) value of \mbox{\texttt{\mdseries\slshape P}}. }

 

\subsection{\textcolor{Chapter }{TzPreImagesNewGens}}
\logpage{[ 48, 9, 4 ]}\nobreak
\hyperdef{L}{X7AC41B117DBB87D6}{}
{\noindent\textcolor{FuncColor}{$\triangleright$\enspace\texttt{TzPreImagesNewGens({\mdseries\slshape P})\index{TzPreImagesNewGens@\texttt{TzPreImagesNewGens}}
\label{TzPreImagesNewGens}
}\hfill{\scriptsize (function)}}\\


 assumes that \mbox{\texttt{\mdseries\slshape P}} is a presentation for which the generator images and preimages are being
traced under Tietze transformations. It returns a list $l$ of words in the old generators of \mbox{\texttt{\mdseries\slshape P}} (the \texttt{OldGeneratorsOfPresentation} (\ref{OldGeneratorsOfPresentation}) value of \mbox{\texttt{\mdseries\slshape P}}) such that the $i$\texttt{\symbol{45}}th entry of $l$ represents the $i$\texttt{\symbol{45}}th (current) generator of \mbox{\texttt{\mdseries\slshape P}} (the \texttt{GeneratorsOfPresentation} (\ref{GeneratorsOfPresentation}) value of \mbox{\texttt{\mdseries\slshape P}}). }

 

\subsection{\textcolor{Chapter }{TzPrintGeneratorImages}}
\logpage{[ 48, 9, 5 ]}\nobreak
\hyperdef{L}{X7F086D0E7AD6173B}{}
{\noindent\textcolor{FuncColor}{$\triangleright$\enspace\texttt{TzPrintGeneratorImages({\mdseries\slshape P})\index{TzPrintGeneratorImages@\texttt{TzPrintGeneratorImages}}
\label{TzPrintGeneratorImages}
}\hfill{\scriptsize (function)}}\\


 assumes that \mbox{\texttt{\mdseries\slshape P}} is a presentation for which the generator images and preimages are being
traced under Tietze transformations. It displays the preimages of the current
generators as Tietze words in the old generators, and the images of the old
generators as Tietze words in the current generators. 
\begin{Verbatim}[commandchars=!@|,fontsize=\small,frame=single,label=Example]
  !gapprompt@gap>| !gapinput@G := PerfectGroup( IsSubgroupFpGroup, 960, 1 );|
  A5 2^4
  !gapprompt@gap>| !gapinput@P := PresentationFpGroup( G );|
  <presentation with 6 gens and 21 rels of total length 84>
  !gapprompt@gap>| !gapinput@TzInitGeneratorImages( P );|
  !gapprompt@gap>| !gapinput@TzGo( P );|
  #I  there are 3 generators and 11 relators of total length 96
  #I  there are 3 generators and 10 relators of total length 81
  !gapprompt@gap>| !gapinput@TzPrintGeneratorImages( P );|
  #I  preimages of current generators as Tietze words in the old ones:
  #I  1. [ 1 ]
  #I  2. [ 2 ]
  #I  3. [ 4 ]
  #I  images of old generators as Tietze words in the current ones:
  #I  1. [ 1 ]
  #I  2. [ 2 ]
  #I  3. [ 1, -2, 1, 3, 1, 2, 1 ]
  #I  4. [ 3 ]
  #I  5. [ -2, 1, 3, 1, 2 ]
  #I  6. [ 1, 3, 1 ]
  !gapprompt@gap>| !gapinput@gens := GeneratorsOfPresentation( P );|
  [ a, b, t ]
  !gapprompt@gap>| !gapinput@oldgens := OldGeneratorsOfPresentation( P );|
  [ a, b, s, t, u, v ]
  !gapprompt@gap>| !gapinput@TzImagesOldGens( P );|
  [ a, b, a*b^-1*a*t*a*b*a, t, b^-1*a*t*a*b, a*t*a ]
  !gapprompt@gap>| !gapinput@for i in [ 1 .. Length( oldgens ) ] do|
  !gapprompt@>| !gapinput@Print( oldgens[i], " = ", TzImagesOldGens( P )[i], "\n" );|
  !gapprompt@>| !gapinput@od;|
  a = a
  b = b
  s = a*b^-1*a*t*a*b*a
  t = t
  u = b^-1*a*t*a*b
  v = a*t*a
\end{Verbatim}
 }

 }

  
\section{\textcolor{Chapter }{The Decoding Tree Procedure}}\label{sect:DecodeTree}
\logpage{[ 48, 10, 0 ]}
\hyperdef{L}{X7D9E283D81CCCF1A}{}
{
  

\subsection{\textcolor{Chapter }{DecodeTree}}
\logpage{[ 48, 10, 1 ]}\nobreak
\hyperdef{L}{X7ACBFE2F78D72A31}{}
{\noindent\textcolor{FuncColor}{$\triangleright$\enspace\texttt{DecodeTree({\mdseries\slshape P})\index{DecodeTree@\texttt{DecodeTree}}
\label{DecodeTree}
}\hfill{\scriptsize (function)}}\\


 assumes that \mbox{\texttt{\mdseries\slshape P}} is a subgroup presentation provided by the Reduced
Reidemeister\texttt{\symbol{45}}Schreier or by the Modified
Todd\texttt{\symbol{45}}Coxeter method (see \texttt{PresentationSubgroupRrs} (\ref{PresentationSubgroupRrs:for two groups (and a string)}), \texttt{PresentationNormalClosureRrs} (\ref{PresentationNormalClosureRrs}), \texttt{PresentationSubgroupMtc} (\ref{PresentationSubgroupMtc})). It eliminates the secondary generators of \mbox{\texttt{\mdseries\slshape P}} (see Section \ref{Subgroup Presentations}) by applying the so called ``decoding tree'' procedure. 

 \texttt{DecodeTree} is called automatically by the command \texttt{PresentationSubgroupMtc} (\ref{PresentationSubgroupMtc}) where it reduces \mbox{\texttt{\mdseries\slshape P}} to a presentation on the given (primary) subgroup generators. \index{secondary subgroup generators} 

 In order to explain the effect of this command we need to insert a few remarks
on the subgroup presentation commands described in section \ref{Subgroup Presentations}. All these commands have the common property that in the process of
constructing a presentation for a given subgroup \mbox{\texttt{\mdseries\slshape H}} of a finitely presented group \mbox{\texttt{\mdseries\slshape G}} they first build up a highly redundant list of generators of \mbox{\texttt{\mdseries\slshape H}} which consists of an (in general small) list of ``primary'' generators, followed by an (in general large) list of ``secondary'' generators, and then construct a presentation $P_0$ \emph{on a sublist of these generators} by rewriting the defining relators of \mbox{\texttt{\mdseries\slshape G}}. This sublist contains all primary, but, at least in general, by far not all
secondary generators. \index{primary subgroup generators} 

 The role of the primary generators depends on the concrete choice of the
subgroup presentation command. If the Modified Todd\texttt{\symbol{45}}Coxeter
method is used, they are just the given generators of \mbox{\texttt{\mdseries\slshape H}}, whereas in the case of the Reduced Reidemeister\texttt{\symbol{45}}Schreier
algorithm they are constructed by the program. 

 Each of the secondary generators is defined by a word of length two in the
preceding generators and their inverses. By historical reasons, the list of
these definitions is called the \emph{subgroup generators tree} though in fact it is not a tree but rather a kind of bush. \index{subgroup generators tree} 

 Now we have to distinguish two cases. If $P_0$ has been constructed by the Reduced Reidemeister\texttt{\symbol{45}}Schreier
routines, it is a presentation of \mbox{\texttt{\mdseries\slshape H}}. However, if the Modified Todd\texttt{\symbol{45}}Coxeter routines have been
used instead, then the relators in $P_0$ are valid relators of \mbox{\texttt{\mdseries\slshape H}}, but they do not necessarily define \mbox{\texttt{\mdseries\slshape H}}. We handle these cases in turn, starting with the latter one. 

 In fact, we could easily receive a presentation of \mbox{\texttt{\mdseries\slshape H}} also in this case if we extended $P_0$ by adding to it all the secondary generators which are not yet contained in it
and all the definitions from the generators tree as additional generators and
relators. Then we could recursively eliminate all secondary generators by
Tietze transformations using the new relators. However, this procedure turns
out to be too inefficient to be of interest. 

 Instead, we use the so called \emph{decoding tree} procedure (see \cite{AMW82}, \cite{AR84}). It proceeds as follows. 

 Starting from $P = P_0$, it runs through a number of steps in each of which it eliminates the current ``last'' generator (with respect to the list of all primary and secondary generators).
If the last generator \mbox{\texttt{\mdseries\slshape g}} is a primary generator, then the procedure terminates. Otherwise it checks
whether there is a relator in the current presentation which can be used to
substitute \mbox{\texttt{\mdseries\slshape g}} by a Tietze transformation. If so, this is done. Otherwise, and only then, the
tree definition of \mbox{\texttt{\mdseries\slshape g}} is added to \mbox{\texttt{\mdseries\slshape P}} as a new relator, and the generators involved are added as new generators if
they have not yet been contained in \mbox{\texttt{\mdseries\slshape P}}. Subsequently, \mbox{\texttt{\mdseries\slshape g}} is eliminated. 

 Note that the extension of \mbox{\texttt{\mdseries\slshape P}} by one or two new generators is \emph{not} a Tietze transformation. In general, it will change the isomorphism type of
the group defined by \mbox{\texttt{\mdseries\slshape P}}. However, it is a remarkable property of this procedure, that at the end,
i.e., as soon as all secondary generators have been eliminated, it provides a
presentation $P = P_1$, say, which defines a group isomorphic to \mbox{\texttt{\mdseries\slshape H}}. In fact, it is this presentation which is returned by the command \texttt{DecodeTree} and hence by the command \texttt{PresentationSubgroupMtc} (\ref{PresentationSubgroupMtc}). 

 If, in the other case, the presentation $P_0$ has been constructed by the Reduced Reidemeister\texttt{\symbol{45}}Schreier
algorithm, then $P_0$ itself is a presentation of \mbox{\texttt{\mdseries\slshape H}}, and the corresponding subgroup presentation command (\texttt{PresentationSubgroupRrs} (\ref{PresentationSubgroupRrs:for two groups (and a string)}) or \texttt{PresentationNormalClosureRrs} (\ref{PresentationNormalClosureRrs})) just returns $P_0$. 

 As mentioned in section \ref{Subgroup Presentations}, we recommend to further simplify this presentation before you use it. The
standard way to do this is to start from $P_0$ and to apply suitable Tietze transformations, e. g., by calling the commands \texttt{TzGo} (\ref{TzGo}) or \texttt{TzGoGo} (\ref{TzGoGo}). This is probably the most efficient approach, but you will end up with a
presentation on some unpredictable set of generators. As an alternative, \textsf{GAP} offers you the \texttt{DecodeTree} command which you can use to eliminate all secondary generators (provided that
there are no space or time problems). For this purpose, the subgroup
presentation commands do not only return the resulting presentation, but also
the tree (together with some associated lists) as a kind of side result in a
component \mbox{\texttt{\mdseries\slshape P}}\texttt{!.tree} of the resulting presentation \mbox{\texttt{\mdseries\slshape P}}. 

 Note, however, that the decoding tree routines will not work correctly any
more on a presentation from which generators have already been eliminated by
Tietze transformations. Therefore, to prevent you from getting wrong results
by calling \texttt{DecodeTree} in such a situation, \textsf{GAP} will automatically remove the subgroup generators tree from a presentation as
soon as one of the generators is substituted by a Tietze transformation. 

 Nevertheless, a certain misuse of the command is still possible, and we want
to explicitly warn you from this. The reason is that the Tietze option
parameters described in Section \ref{Tietze Options} apply to \texttt{DecodeTree} as well. Hence, in case of inadequate values of these parameters, it may
happen that \texttt{DecodeTree} stops before all the secondary generators have vanished. In this case \textsf{GAP} will display an appropriate warning. Then you should change the respective
parameters and continue the process by calling \texttt{DecodeTree} again. Otherwise, if you would apply Tietze transformations, it might happen
because of the convention described above that the tree is removed and that
you end up with a wrong presentation. 

 After a successful run of \texttt{DecodeTree} it is convenient to further simplify the resulting presentation by suitable
Tietze transformations. 

 As an example of an explicit call of \texttt{DecodeTree} we compute two presentations of a subgroup of order $384$ in a group of order $6912$. In both cases we use the Reduced Reidemeister\texttt{\symbol{45}}Schreier
algorithm, but in the first run we just apply the Tietze transformations
offered by the \texttt{TzGoGo} (\ref{TzGoGo}) command with its default parameters, whereas in the second run we call the \texttt{DecodeTree} command before. 

 
\begin{Verbatim}[commandchars=!@|,fontsize=\small,frame=single,label=Example]
  !gapprompt@gap>| !gapinput@F2 := FreeGroup( "a", "b" );;|
  !gapprompt@gap>| !gapinput@G := F2 / [ F2.1*F2.2^2*F2.1^-1*F2.2^-1*F2.1^3*F2.2^-1,|
  !gapprompt@>| !gapinput@               F2.2*F2.1^2*F2.2^-1*F2.1^-1*F2.2^3*F2.1^-1 ];;|
  !gapprompt@gap>| !gapinput@a := G.1;;  b := G.2;;|
  !gapprompt@gap>| !gapinput@H := Subgroup( G, [ Comm(a^-1,b^-1), Comm(a^-1,b), Comm(a,b) ] );;|
\end{Verbatim}
 

 We use the Reduced Reidemeister Schreier method and default Tietze
transformations to get a presentation for \mbox{\texttt{\mdseries\slshape H}}. 

 
\begin{Verbatim}[commandchars=!@|,fontsize=\small,frame=single,label=Example]
  !gapprompt@gap>| !gapinput@P := PresentationSubgroupRrs( G, H );|
  <presentation with 18 gens and 35 rels of total length 169>
  !gapprompt@gap>| !gapinput@TzGoGo( P );|
  #I  there are 3 generators and 20 relators of total length 488
  #I  there are 3 generators and 20 relators of total length 466
\end{Verbatim}
 

 We end up with 20 relators of total length 466. Now we repeat the procedure,
but we call the decoding tree algorithm before doing the Tietze
transformations. 

 
\begin{Verbatim}[commandchars=!@|,fontsize=\small,frame=single,label=Example]
  !gapprompt@gap>| !gapinput@P := PresentationSubgroupRrs( G, H );|
  <presentation with 18 gens and 35 rels of total length 169>
  !gapprompt@gap>| !gapinput@DecodeTree( P );|
  #I  there are 9 generators and 26 relators of total length 185
  #I  there are 6 generators and 23 relators of total length 213
  #I  there are 3 generators and 20 relators of total length 252
  #I  there are 3 generators and 20 relators of total length 244
  !gapprompt@gap>| !gapinput@TzGoGo( P );|
  #I  there are 3 generators and 19 relators of total length 168
  #I  there are 3 generators and 17 relators of total length 138
  #I  there are 3 generators and 15 relators of total length 114
  #I  there are 3 generators and 13 relators of total length 96
  #I  there are 3 generators and 12 relators of total length 84
\end{Verbatim}
 

 This time we end up with a shorter presentation. 

 

 }

 }

  
\section{\textcolor{Chapter }{Tietze Options}}\label{Tietze Options}
\logpage{[ 48, 11, 0 ]}
\hyperdef{L}{X856F37537E9927EE}{}
{
  Several of the Tietze transformation commands described above are controlled
by certain parameters, the \emph{Tietze options}, which often have a tremendous influence on their performance and results.
However, in each application of the commands, an appropriate choice of these
option parameters will depend on the concrete presentation under
investigation. Therefore we have implemented the Tietze options in such a way
that they are associated to the presentation: Each presentation keeps its own
set of Tietze option parameters as an attribute. 

\subsection{\textcolor{Chapter }{TzOptions}}
\logpage{[ 48, 11, 1 ]}\nobreak
\hyperdef{L}{X8178683283214D88}{}
{\noindent\textcolor{FuncColor}{$\triangleright$\enspace\texttt{TzOptions({\mdseries\slshape P})\index{TzOptions@\texttt{TzOptions}}
\label{TzOptions}
}\hfill{\scriptsize (attribute)}}\\


 is a record whose components direct the heuristics applied by the Tietze
transformation functions. 

 You may alter the value of any of these Tietze options by just assigning a new
value to the respective record component. 

 The following Tietze options are recognized by \textsf{GAP}: 

 
\begin{description}
\item[{\texttt{protected}:}]  The first \texttt{protected} generators in a presentation \mbox{\texttt{\mdseries\slshape P}} are protected from being eliminated by the Tietze transformations functions.
There are only two exceptions: The option \texttt{protected} is ignored by the functions \texttt{TzEliminate} (\ref{TzEliminate:for a presentation (and a generator)}) and \texttt{TzSubstitute} (\ref{TzSubstitute:for a presentation and a word}) because they explicitly specify the generator to be eliminated. The default
value of \texttt{protected} is 0. 
\item[{\texttt{eliminationsLimit}:}]  Whenever the elimination phase of the \texttt{TzGo} (\ref{TzGo}) command is entered for a presentation \mbox{\texttt{\mdseries\slshape P}}, then it will eliminate at most \texttt{eliminationsLimit} generators (except for further ones which have turned out to be trivial).
Hence you may use the \texttt{eliminationsLimit} parameter as a break criterion for the \texttt{TzGo} (\ref{TzGo}) command. Note, however, that it is ignored by the \texttt{TzEliminate} (\ref{TzEliminate:for a presentation (and a generator)}) command. The default value of \texttt{eliminationsLimit} is 100. 
\item[{\texttt{expandLimit}:}]  Whenever the routine for eliminating more than 1 generator is called for a
presentation \mbox{\texttt{\mdseries\slshape P}} by the \texttt{TzEliminate} (\ref{TzEliminate:for a presentation (and a generator)}) command or the elimination phase of the \texttt{TzGo} (\ref{TzGo}) command, then it saves the given total length of the relators, and
subsequently it checks the current total length against its value before each
elimination. If the total length has increased to more than \texttt{expandLimit} per cent of its original value, then the routine returns instead of
eliminating another generator. Hence you may use the \texttt{expandLimit} parameter as a break criterion for the \texttt{TzGo} (\ref{TzGo}) command. The default value of \texttt{expandLimit} is 150. 
\item[{\texttt{generatorsLimit}:}]  Whenever the elimination phase of the \texttt{TzGo} (\ref{TzGo}) command is entered for a presentation \mbox{\texttt{\mdseries\slshape P}} with $n$ generators, then it will eliminate at most $n - $\texttt{generatorsLimit} generators (except for generators which turn out to be trivial). Hence you may
use the \texttt{generatorsLimit} parameter as a break criterion for the \texttt{TzGo} (\ref{TzGo}) command. The default value of \texttt{generatorsLimit} is 0. 
\item[{\texttt{lengthLimit}:}]  The Tietze transformation commands will never eliminate a generator of a
presentation \mbox{\texttt{\mdseries\slshape P}}, if they cannot exclude the possibility that the resulting total length of
the relators exceeds the maximal \textsf{GAP} list length of $2^{31}-1$ or the value of the option \texttt{lengthLimit}. The default value of \texttt{lengthLimit} is $2^{31}-1$. 
\item[{\texttt{loopLimit}:}]  Whenever the \texttt{TzGo} (\ref{TzGo}) command is called for a presentation \mbox{\texttt{\mdseries\slshape P}}, then it will loop over at most \texttt{loopLimit} of its basic steps. Hence you may use the \texttt{loopLimit} parameter as a break criterion for the \texttt{TzGo} (\ref{TzGo}) command. The default value of \texttt{loopLimit} is \texttt{infinity} (\ref{infinity}). 
\item[{\texttt{printLevel}:}]  Whenever Tietze transformation commands are called for a presentation \mbox{\texttt{\mdseries\slshape P}} with \texttt{printLevel} $= 0$, they will not provide any output except for error messages. If \texttt{printLevel} $= 1$, they will display some reasonable amount of output which allows you to watch
the progress of the computation and to decide about your next commands. In the
case \texttt{printLevel} $= 2$, you will get a much more generous amount of output. Finally, if \texttt{printLevel} $= 3$, various messages on internal details will be added. The default value of \texttt{printLevel} is 1. 
\item[{\texttt{saveLimit}:}]  Whenever the \texttt{TzSearch} (\ref{TzSearch}) command has finished its main loop over all relators of a presentation \mbox{\texttt{\mdseries\slshape P}}, then it checks whether during this loop the total length of the relators has
been reduced by at least \texttt{saveLimit} per cent. If this is the case, then \texttt{TzSearch} (\ref{TzSearch}) repeats its procedure instead of returning. Hence you may use the \texttt{saveLimit} parameter as a break criterion for the \texttt{TzSearch} (\ref{TzSearch}) command and, in particular, for the search phase of the \texttt{TzGo} (\ref{TzGo}) command. The default value of \texttt{saveLimit} is 10. 
\item[{\texttt{searchSimultaneous}:}]  Whenever the \texttt{TzSearch} (\ref{TzSearch}) or the \texttt{TzSearchEqual} (\ref{TzSearchEqual}) command is called for a presentation \mbox{\texttt{\mdseries\slshape P}}, then it is allowed to handle up to \texttt{searchSimultaneous} short relators simultaneously (see the description of the \texttt{TzSearch} (\ref{TzSearch}) command for more details). The choice of this parameter may heavily influence
the performance as well as the result of the \texttt{TzSearch} (\ref{TzSearch}) and the \texttt{TzSearchEqual} (\ref{TzSearchEqual}) commands and hence also of the search phase of the \texttt{TzGo} (\ref{TzGo}) command. The default value of \texttt{searchSimultaneous} is 20. 
\end{description}
 }

 

\subsection{\textcolor{Chapter }{TzPrintOptions}}
\logpage{[ 48, 11, 2 ]}\nobreak
\hyperdef{L}{X7BC90B6882DE6D10}{}
{\noindent\textcolor{FuncColor}{$\triangleright$\enspace\texttt{TzPrintOptions({\mdseries\slshape P})\index{TzPrintOptions@\texttt{TzPrintOptions}}
\label{TzPrintOptions}
}\hfill{\scriptsize (function)}}\\


 prints the current values of the Tietze options of the presentation \mbox{\texttt{\mdseries\slshape P}}. 
\begin{Verbatim}[commandchars=!@|,fontsize=\small,frame=single,label=Example]
  !gapprompt@gap>| !gapinput@TzPrintOptions( P );|
  #I  protected          = 0
  #I  eliminationsLimit  = 100
  #I  expandLimit        = 150
  #I  generatorsLimit    = 0
  #I  lengthLimit        = 2147483647
  #I  loopLimit          = infinity
  #I  printLevel         = 1
  #I  saveLimit          = 10
  #I  searchSimultaneous = 20
\end{Verbatim}
 }

 }

 }

         
\chapter{\textcolor{Chapter }{Group Products}}\label{Group Products}
\logpage{[ 49, 0, 0 ]}
\hyperdef{L}{X7D5C75647DB168F1}{}
{
  This chapter describes the various group product constructions that are
possible in \textsf{GAP}. 

 At the moment for some of the products methods are available only if both
factors are given in the same representation or only for certain types of
groups such as permutation groups and pc groups when the product can be
naturally represented as a group of the same kind. 

 \textsf{GAP} does not guarantee that a product of two groups will be in a particular
representation. (Exceptions are \texttt{WreathProductImprimitiveAction} (\ref{WreathProductImprimitiveAction}) and \texttt{WreathProductProductAction} (\ref{WreathProductProductAction}) which are construction that makes sense only for permutation groups,
see{\nobreakspace}\texttt{WreathProduct} (\ref{WreathProduct})). 

 \textsf{GAP} however will try to choose an efficient representation, so products of
permutation groups or pc groups often will be represented as a group of the
same kind again. 

 Therefore the only guaranteed way to relate a product to its factors is via
the embedding and projection homomorphisms, see{\nobreakspace}\ref{Embeddings and Projections for Group Products}.  
\section{\textcolor{Chapter }{Direct Products}}\label{Direct Products}
\logpage{[ 49, 1, 0 ]}
\hyperdef{L}{X7D39232A84CD8DBD}{}
{
  The direct product of groups is the cartesian product of the groups
(considered as element sets) with component\texttt{\symbol{45}}wise
multiplication. 

 

\subsection{\textcolor{Chapter }{DirectProduct}}
\logpage{[ 49, 1, 1 ]}\nobreak
\hyperdef{L}{X861BA02C7902A4F4}{}
{\noindent\textcolor{FuncColor}{$\triangleright$\enspace\texttt{DirectProduct({\mdseries\slshape G[, H, ...]})\index{DirectProduct@\texttt{DirectProduct}}
\label{DirectProduct}
}\hfill{\scriptsize (function)}}\\
\noindent\textcolor{FuncColor}{$\triangleright$\enspace\texttt{DirectProductOp({\mdseries\slshape list, expl})\index{DirectProductOp@\texttt{DirectProductOp}}
\label{DirectProductOp}
}\hfill{\scriptsize (operation)}}\\


 These functions construct the direct product of the groups given as arguments. \texttt{DirectProduct} takes an arbitrary positive number of arguments and calls the operation \texttt{DirectProductOp}, which takes exactly two arguments, namely a nonempty list \mbox{\texttt{\mdseries\slshape list}} of groups and one of these groups, \mbox{\texttt{\mdseries\slshape expl}}. (This somewhat strange syntax allows the method selection to choose a
reasonable method for special cases, e.g., if all groups are permutation
groups or pc groups.) 

 \textsf{GAP} will try to choose an efficient representation for the direct product. For
example the direct product of permutation groups will be a permutation group
again and the direct product of pc groups will be a pc group. 

 If the groups are in different representations a generic direct product will
be formed which may not be particularly efficient for many calculations.
Instead it may be worth to convert all factors to a common representation
first, before forming the product. 

 \index{Embedding@\texttt{Embedding}!example for direct products} \index{Projection@\texttt{Projection}!example for direct products} For a direct product $P$, calling \texttt{Embedding} (\ref{Embedding:for a domain and a positive integer}) with $P$ and $n$ yields the homomorphism embedding the $n$\texttt{\symbol{45}}th factor into $P$; calling \texttt{Projection} (\ref{Projection:for a domain and a positive integer}) with \mbox{\texttt{\mdseries\slshape P}} and \mbox{\texttt{\mdseries\slshape n}} yields the projection of $P$ onto the $n$\texttt{\symbol{45}}th factor, see{\nobreakspace}\ref{Embeddings and Projections for Group Products}. 

 
\begin{Verbatim}[commandchars=!@|,fontsize=\small,frame=single,label=Example]
  !gapprompt@gap>| !gapinput@g:=Group((1,2,3),(1,2));;|
  !gapprompt@gap>| !gapinput@d:=DirectProduct(g,g,g);|
  Group([ (1,2,3), (1,2), (4,5,6), (4,5), (7,8,9), (7,8) ])
  !gapprompt@gap>| !gapinput@Size(d);|
  216
  !gapprompt@gap>| !gapinput@e:=Embedding(d,2);|
  2nd embedding into Group([ (1,2,3), (1,2), (4,5,6), (4,5), (7,8,9),
    (7,8) ])
  !gapprompt@gap>| !gapinput@Image(e,(1,2));|
  (4,5)
  !gapprompt@gap>| !gapinput@Image(Projection(d,2),(1,2,3)(4,5)(8,9));|
  (1,2)
  !gapprompt@gap>| !gapinput@f:=FreeGroup("a","b");;|
  !gapprompt@gap>| !gapinput@g:=f/ParseRelators(f,"a2,b3,(ab)5");|
  <fp group on the generators [ a, b ]>
  !gapprompt@gap>| !gapinput@f2:=FreeGroup("x","y");;|
  !gapprompt@gap>| !gapinput@h:=f2/ParseRelators(f2,"x2,y4,xy=Yx");|
  <fp group on the generators [ x, y ]>
  !gapprompt@gap>| !gapinput@d:=DirectProduct(g,h);|
  <fp group on the generators [ a, b, x, y ]>
\end{Verbatim}
 }

 }

  
\section{\textcolor{Chapter }{Semidirect Products}}\label{Semidirect Products}
\logpage{[ 49, 2, 0 ]}
\hyperdef{L}{X87FE512E7DB7346C}{}
{
  The semidirect product of a group $N$ with a group $G$ acting on $N$ via a homomorphism $\alpha$ from $G$ into the automorphism group of $N$ is the cartesian product $G \times N$ with the multiplication $(g, n) \cdot (h, m) = (gh, n^{{h^\alpha}}m)$. 
\subsection{\textcolor{Chapter }{SemidirectProduct}}\logpage{[ 49, 2, 1 ]}
\hyperdef{L}{X7D905A5778D7ACDE}{}
{
\noindent\textcolor{FuncColor}{$\triangleright$\enspace\texttt{SemidirectProduct({\mdseries\slshape G, alpha, N})\index{SemidirectProduct@\texttt{SemidirectProduct}!for acting group, action, and a group}
\label{SemidirectProduct:for acting group, action, and a group}
}\hfill{\scriptsize (operation)}}\\
\noindent\textcolor{FuncColor}{$\triangleright$\enspace\texttt{SemidirectProduct({\mdseries\slshape autgp, N})\index{SemidirectProduct@\texttt{SemidirectProduct}!for a group of automorphisms and a group}
\label{SemidirectProduct:for a group of automorphisms and a group}
}\hfill{\scriptsize (operation)}}\\


 constructs the semidirect product of \mbox{\texttt{\mdseries\slshape N}} with \mbox{\texttt{\mdseries\slshape G}} acting via \mbox{\texttt{\mdseries\slshape alpha}}, which must be a homomorphism from \mbox{\texttt{\mdseries\slshape G}} into a group of automorphisms of \mbox{\texttt{\mdseries\slshape N}}. 

 If \mbox{\texttt{\mdseries\slshape N}} is a group, \mbox{\texttt{\mdseries\slshape alpha}} must be a homomorphism from \mbox{\texttt{\mdseries\slshape G}} into a group of automorphisms of \mbox{\texttt{\mdseries\slshape N}}. 

 If \mbox{\texttt{\mdseries\slshape N}} is a full row space over a field \mbox{\texttt{\mdseries\slshape F}}, \mbox{\texttt{\mdseries\slshape alpha}} must be a homomorphism from \mbox{\texttt{\mdseries\slshape G}} into a matrix group of the right dimension over a subfield of \mbox{\texttt{\mdseries\slshape F}}, or into a permutation group (in this case permutation matrices are taken). 

 In the second variant, \mbox{\texttt{\mdseries\slshape autgp}} must be a group of automorphism of \mbox{\texttt{\mdseries\slshape N}}, it is a shorthand for \texttt{SemidirectProduct(\mbox{\texttt{\mdseries\slshape autgp}},IdentityMapping(\mbox{\texttt{\mdseries\slshape autgp}}),\mbox{\texttt{\mdseries\slshape N}})}. Note that (unless \mbox{\texttt{\mdseries\slshape autgrp}} has been obtained by the operation \texttt{AutomorphismGroup} (\ref{AutomorphismGroup})) you have to test \texttt{IsGroupOfAutomorphisms} (\ref{IsGroupOfAutomorphisms}) for \mbox{\texttt{\mdseries\slshape autgrp}} to ensure that \textsf{GAP} knows that \mbox{\texttt{\mdseries\slshape autgrp}} consists of group automorphisms. 
\begin{Verbatim}[commandchars=!@|,fontsize=\small,frame=single,label=Example]
  !gapprompt@gap>| !gapinput@n:=AbelianGroup(IsPcGroup,[5,5]);|
  <pc group of size 25 with 2 generators>
  !gapprompt@gap>| !gapinput@au:=DerivedSubgroup(AutomorphismGroup(n));;|
  !gapprompt@gap>| !gapinput@Size(au);|
  120
  !gapprompt@gap>| !gapinput@p:=SemidirectProduct(au,n);;|
  !gapprompt@gap>| !gapinput@Size(p);|
  3000
  !gapprompt@gap>| !gapinput@n:=Group((1,2),(3,4));;|
  !gapprompt@gap>| !gapinput@au:=AutomorphismGroup(n);;|
  !gapprompt@gap>| !gapinput@au:=First(AsSet(au),i->Order(i)=3);;|
  !gapprompt@gap>| !gapinput@au:=Group(au);|
  <group with 1 generator>
  !gapprompt@gap>| !gapinput@IsGroupOfAutomorphisms(au);|
  true
  !gapprompt@gap>| !gapinput@SemidirectProduct(au,n);|
  <pc group with 3 generators>
  !gapprompt@gap>| !gapinput@n:=AbelianGroup(IsPcGroup,[2,2]);|
  <pc group of size 4 with 2 generators>
  !gapprompt@gap>| !gapinput@au:=AutomorphismGroup(n);;|
  !gapprompt@gap>| !gapinput@apc:=IsomorphismPcGroup(au);;|
  !gapprompt@gap>| !gapinput@g:=Image(apc);|
  Group([ f1, f2 ])
  !gapprompt@gap>| !gapinput@apci:=InverseGeneralMapping(apc);;|
  !gapprompt@gap>| !gapinput@IsGroupHomomorphism(apci);|
  true
  !gapprompt@gap>| !gapinput@p:=SemidirectProduct(g,apci,n);|
  <pc group of size 24 with 4 generators>
  !gapprompt@gap>| !gapinput@IsomorphismGroups(p,Group((1,2,3,4),(1,2))) <> fail;|
  true
  !gapprompt@gap>| !gapinput@SemidirectProduct(SU(3,3),GF(9)^3);|
  <matrix group of size 4408992 with 3 generators>
  !gapprompt@gap>| !gapinput@SemidirectProduct(Group((1,2,3),(2,3,4)),GF(5)^4);|
  <matrix group of size 7500 with 3 generators>
  !gapprompt@gap>| !gapinput@g:=Group((3,4,5),(1,2,3));;|
  !gapprompt@gap>| !gapinput@mats:=[[[Z(2^2),0*Z(2)],[0*Z(2),Z(2^2)^2]],|
  !gapprompt@>| !gapinput@         [[Z(2)^0,Z(2)^0], [Z(2)^0,0*Z(2)]]];;|
  !gapprompt@gap>| !gapinput@hom:=GroupHomomorphismByImages(g,Group(mats),[g.1,g.2],mats);;|
  !gapprompt@gap>| !gapinput@SemidirectProduct(g,hom,GF(4)^2);|
  <matrix group of size 960 with 3 generators>
  !gapprompt@gap>| !gapinput@SemidirectProduct(g,hom,GF(16)^2);|
  <matrix group of size 15360 with 4 generators>
\end{Verbatim}
 

 \index{Embedding@\texttt{Embedding}!example for semidirect products} \index{Projection@\texttt{Projection}!example for semidirect products} For a semidirect product $P$ of \mbox{\texttt{\mdseries\slshape G}} with \mbox{\texttt{\mdseries\slshape N}}, calling \texttt{Embedding} (\ref{Embedding:for a domain and a positive integer}) with $P$ and \texttt{1} yields the embedding of \mbox{\texttt{\mdseries\slshape G}}, calling \texttt{Embedding} (\ref{Embedding:for a domain and a positive integer}) with $P$ and \texttt{2} yields the embedding of \mbox{\texttt{\mdseries\slshape N}}; calling \texttt{Projection} (\ref{Projection:for a domain and a positive integer}) with \mbox{\texttt{\mdseries\slshape P}} yields the projection of $P$ onto \mbox{\texttt{\mdseries\slshape G}}, see{\nobreakspace}\ref{Embeddings and Projections for Group Products}. 

 
\begin{Verbatim}[commandchars=!@|,fontsize=\small,frame=single,label=Example]
  !gapprompt@gap>| !gapinput@Size(Image(Embedding(p,1)));|
  6
  !gapprompt@gap>| !gapinput@Embedding(p,2);|
  [ f1, f2 ] -> [ f3, f4 ]
  !gapprompt@gap>| !gapinput@Projection(p);|
  [ f1, f2, f3, f4 ] -> [ f1, f2, <identity> of ..., <identity> of ... ]
\end{Verbatim}
 }

 }

  
\section{\textcolor{Chapter }{Subdirect Products}}\label{Subdirect Products}
\logpage{[ 49, 3, 0 ]}
\hyperdef{L}{X815AFC537B215D7B}{}
{
  The subdirect product of the groups $G$ and $H$ with respect to the epimorphisms $\varphi\colon G \rightarrow A$ and $\psi\colon H \rightarrow A$ (for a common group $A$) is the subgroup of the direct product $G \times H$ consisting of the elements $(g,h)$ for which $g^{\varphi} = h^{\psi}$. It is the pull\texttt{\symbol{45}}back of the following diagram. 

 
\[ \matrix{ &&G&\cr &&\Big\downarrow&\varphi\cr
H&\mathop{\longrightarrow}\limits^{\psi}&A&\cr } \]
  

\subsection{\textcolor{Chapter }{SubdirectProduct}}
\logpage{[ 49, 3, 1 ]}\nobreak
\hyperdef{L}{X82112D768085AD98}{}
{\noindent\textcolor{FuncColor}{$\triangleright$\enspace\texttt{SubdirectProduct({\mdseries\slshape G, H, Ghom, Hhom})\index{SubdirectProduct@\texttt{SubdirectProduct}}
\label{SubdirectProduct}
}\hfill{\scriptsize (function)}}\\


 constructs the subdirect product of \mbox{\texttt{\mdseries\slshape G}} and \mbox{\texttt{\mdseries\slshape H}} with respect to the epimorphisms \mbox{\texttt{\mdseries\slshape Ghom}} from \mbox{\texttt{\mdseries\slshape G}} onto a group $A$ and \mbox{\texttt{\mdseries\slshape Hhom}} from \mbox{\texttt{\mdseries\slshape H}} onto the same group $A$. 

 \index{Projection@\texttt{Projection}!example for subdirect products} For a subdirect product $P$, calling \texttt{Projection} (\ref{Projection:for a domain and a positive integer}) with $P$ and $n$ yields the projection on the $n$\texttt{\symbol{45}}th factor. (In general the factors do not embed into a
subdirect product.) 

 
\begin{Verbatim}[commandchars=!@|,fontsize=\small,frame=single,label=Example]
  !gapprompt@gap>| !gapinput@g:=Group((1,2,3),(1,2));|
  Group([ (1,2,3), (1,2) ])
  !gapprompt@gap>| !gapinput@hom:=GroupHomomorphismByImagesNC(g,g,[(1,2,3),(1,2)],[(),(1,2)]);|
  [ (1,2,3), (1,2) ] -> [ (), (1,2) ]
  !gapprompt@gap>| !gapinput@s:=SubdirectProduct(g,g,hom,hom);|
  Group([ (1,2,3), (1,2)(4,5), (4,5,6) ])
  !gapprompt@gap>| !gapinput@Size(s);|
  18
  !gapprompt@gap>| !gapinput@p:=Projection(s,2);|
  2nd projection of Group([ (1,2,3), (1,2)(4,5), (4,5,6) ])
  !gapprompt@gap>| !gapinput@Image(p,(1,3,2)(4,5,6));|
  (1,2,3)
\end{Verbatim}
 }

 

\subsection{\textcolor{Chapter }{SubdirectProducts}}
\logpage{[ 49, 3, 2 ]}\nobreak
\hyperdef{L}{X814204E97812894C}{}
{\noindent\textcolor{FuncColor}{$\triangleright$\enspace\texttt{SubdirectProducts({\mdseries\slshape G, H})\index{SubdirectProducts@\texttt{SubdirectProducts}}
\label{SubdirectProducts}
}\hfill{\scriptsize (function)}}\\


 this function computes all subdirect products of \mbox{\texttt{\mdseries\slshape G}} and \mbox{\texttt{\mdseries\slshape H}} up to conjugacy in the direct product of Parent(\mbox{\texttt{\mdseries\slshape G}}) and Parent(\mbox{\texttt{\mdseries\slshape H}}). The subdirect products are returned as subgroups of this direct product. }

 }

  
\section{\textcolor{Chapter }{Wreath Products}}\label{Wreath Products}
\logpage{[ 49, 4, 0 ]}
\hyperdef{L}{X7DF2AEBC8518FFA4}{}
{
  The wreath product of a group $G$ with a permutation group $P$ acting on $n$ points is the semidirect product of the normal subgroup $\mbox{\texttt{\mdseries\slshape G}}^n$ with the group $P$ which acts on $\mbox{\texttt{\mdseries\slshape G}}^n$ by permuting the components. 

 Note that \textsf{GAP} always considers the domain of a permutation group to be the points moved by
elements of the group as returned by \texttt{MovedPoints} (\ref{MovedPoints:for a list or collection of permutations}), i.e. it is not possible to have a domain to include fixed points, I.e. $P = \langle (1,2,3) \rangle$ and $P = \langle (1,3,5) \rangle$ result in isomorphic wreath products. (If fixed points are desired the wreath
product $G \wr T$ has to be formed with a transitive overgroup $T$ of $P$ and then the pre\texttt{\symbol{45}}image of $P$ under the projection $G \wr T \rightarrow T$ has to be taken.) 

 

\subsection{\textcolor{Chapter }{WreathProduct}}
\logpage{[ 49, 4, 1 ]}\nobreak
\hyperdef{L}{X8786EFBC78D7D6ED}{}
{\noindent\textcolor{FuncColor}{$\triangleright$\enspace\texttt{WreathProduct({\mdseries\slshape G, H[, hom]})\index{WreathProduct@\texttt{WreathProduct}}
\label{WreathProduct}
}\hfill{\scriptsize (operation)}}\\
\noindent\textcolor{FuncColor}{$\triangleright$\enspace\texttt{StandardWreathProduct({\mdseries\slshape G, H})\index{StandardWreathProduct@\texttt{StandardWreathProduct}}
\label{StandardWreathProduct}
}\hfill{\scriptsize (operation)}}\\


 \texttt{WreathProduct} constructs the wreath product of the group \mbox{\texttt{\mdseries\slshape G}} with the group \mbox{\texttt{\mdseries\slshape H}}, acting as a permutation group. 

 If a third argument \mbox{\texttt{\mdseries\slshape hom}} is given, it must be a homomorphism from \mbox{\texttt{\mdseries\slshape H}} into a permutation group, and the action of this group on its moved points is
considered. 

 If only two arguments are given, \mbox{\texttt{\mdseries\slshape H}} must be a permutation group. 

 \texttt{StandardWreathProduct} returns the wreath product for the (right regular) permutation action of \mbox{\texttt{\mdseries\slshape H}} on its elements. 

 \index{Embedding@\texttt{Embedding}!example for wreath products} \index{Projection@\texttt{Projection}!example for wreath products} For a wreath product $W$ of \mbox{\texttt{\mdseries\slshape G}} with a permutation group $P$ of degree $n$ and $1 \leq i \leq n$ calling \texttt{Embedding} (\ref{Embedding:for a domain and a positive integer}) with $W$ and $i$ yields the embedding of \mbox{\texttt{\mdseries\slshape G}} in the $i$\texttt{\symbol{45}}th component of the direct product of the base group $\mbox{\texttt{\mdseries\slshape G}}^n$ of $W$. For $i = n+1$, \texttt{Embedding} (\ref{Embedding:for a domain and a positive integer}) yields the embedding of $P$ into $W$. Calling \texttt{Projection} (\ref{Projection:for a domain and a positive integer}) with $W$ yields the projection onto the acting group $P$, see{\nobreakspace}\ref{Embeddings and Projections for Group Products}. 

 
\begin{Verbatim}[commandchars=!@|,fontsize=\small,frame=single,label=Example]
  !gapprompt@gap>| !gapinput@g:=Group((1,2,3),(1,2));|
  Group([ (1,2,3), (1,2) ])
  !gapprompt@gap>| !gapinput@p:=Group((1,2,3));|
  Group([ (1,2,3) ])
  !gapprompt@gap>| !gapinput@w:=WreathProduct(g,p);|
  Group([ (1,2,3), (1,2), (4,5,6), (4,5), (7,8,9), (7,8),
    (1,4,7)(2,5,8)(3,6,9) ])
  !gapprompt@gap>| !gapinput@Size(w);|
  648
  !gapprompt@gap>| !gapinput@Embedding(w,1);|
  1st embedding into Group( [ (1,2,3), (1,2), (4,5,6), (4,5), (7,8,9),
    (7,8), (1,4,7)(2,5,8)(3,6,9) ] )
  !gapprompt@gap>| !gapinput@Image(Embedding(w,3));|
  Group([ (7,8,9), (7,8) ])
  !gapprompt@gap>| !gapinput@Image(Embedding(w,4));|
  Group([ (1,4,7)(2,5,8)(3,6,9) ])
  !gapprompt@gap>| !gapinput@Image(Projection(w),(1,4,8,2,6,7,3,5,9));|
  (1,2,3)
\end{Verbatim}
 }

 

\subsection{\textcolor{Chapter }{WreathProductImprimitiveAction}}
\logpage{[ 49, 4, 2 ]}\nobreak
\hyperdef{L}{X8589DCFA7C2E5FAA}{}
{\noindent\textcolor{FuncColor}{$\triangleright$\enspace\texttt{WreathProductImprimitiveAction({\mdseries\slshape G, H})\index{WreathProductImprimitiveAction@\texttt{WreathProductImprimitiveAction}}
\label{WreathProductImprimitiveAction}
}\hfill{\scriptsize (function)}}\\


 For two permutation groups \mbox{\texttt{\mdseries\slshape G}} and \mbox{\texttt{\mdseries\slshape H}}, this function constructs the wreath product of \mbox{\texttt{\mdseries\slshape G}} and \mbox{\texttt{\mdseries\slshape H}} in the imprimitive action. If \mbox{\texttt{\mdseries\slshape G}} acts on $l$ points and \mbox{\texttt{\mdseries\slshape H}} on $m$ points this action will be on $l \cdot m$ points, it will be imprimitive with $m$ blocks of size $l$ each. 

 The operations \texttt{Embedding} (\ref{Embedding:for two domains}) and \texttt{Projection} (\ref{Projection:for two domains}) operate on this product as described for general wreath products. 

 
\begin{Verbatim}[commandchars=!@|,fontsize=\small,frame=single,label=Example]
  !gapprompt@gap>| !gapinput@w:=WreathProductImprimitiveAction(g,p);;|
  !gapprompt@gap>| !gapinput@LargestMovedPoint(w);|
  9
\end{Verbatim}
 }

 

\subsection{\textcolor{Chapter }{WreathProductProductAction}}
\logpage{[ 49, 4, 3 ]}\nobreak
\hyperdef{L}{X82B8DD1C868A3726}{}
{\noindent\textcolor{FuncColor}{$\triangleright$\enspace\texttt{WreathProductProductAction({\mdseries\slshape G, H})\index{WreathProductProductAction@\texttt{WreathProductProductAction}}
\label{WreathProductProductAction}
}\hfill{\scriptsize (function)}}\\


 For two permutation groups \mbox{\texttt{\mdseries\slshape G}} and \mbox{\texttt{\mdseries\slshape H}}, this function constructs the wreath product in product action. If \mbox{\texttt{\mdseries\slshape G}} acts on $l$ points and \mbox{\texttt{\mdseries\slshape H}} on $m$ points this action will be on $l^m$ points. 

 The operations \texttt{Embedding} (\ref{Embedding:for two domains}) and \texttt{Projection} (\ref{Projection:for two domains}) operate on this product as described for general wreath products. 
\begin{Verbatim}[commandchars=!@|,fontsize=\small,frame=single,label=Example]
  !gapprompt@gap>| !gapinput@w:=WreathProductProductAction(g,p);|
  <permutation group of size 648 with 7 generators>
  !gapprompt@gap>| !gapinput@LargestMovedPoint(w);|
  27
\end{Verbatim}
 }

 

\subsection{\textcolor{Chapter }{KuKGenerators}}
\logpage{[ 49, 4, 4 ]}\nobreak
\hyperdef{L}{X80634C3180E0C593}{}
{\noindent\textcolor{FuncColor}{$\triangleright$\enspace\texttt{KuKGenerators({\mdseries\slshape G, beta, alpha})\index{KuKGenerators@\texttt{KuKGenerators}}
\label{KuKGenerators}
}\hfill{\scriptsize (function)}}\\


 \index{Krasner\texttt{\symbol{45}}Kaloujnine theorem} \index{Wreath product embedding} If \mbox{\texttt{\mdseries\slshape beta}} is a homomorphism from \mbox{\texttt{\mdseries\slshape G}} into a transitive permutation group, $U$ the full preimage of the point stabilizer and \mbox{\texttt{\mdseries\slshape alpha}} a homomorphism defined on (a superset) of $U$, this function returns images of the generators of \mbox{\texttt{\mdseries\slshape G}} when mapping to the wreath product $(U \mbox{\texttt{\mdseries\slshape alpha}}) \wr (\mbox{\texttt{\mdseries\slshape G}} \mbox{\texttt{\mdseries\slshape beta}})$. (This is the Krasner\texttt{\symbol{45}}Kaloujnine embedding theorem.) 
\begin{Verbatim}[commandchars=!@|,fontsize=\small,frame=single,label=Example]
  !gapprompt@gap>| !gapinput@g:=Group((1,2,3,4),(1,2));;|
  !gapprompt@gap>| !gapinput@hom:=GroupHomomorphismByImages(g,Group((1,2)),|
  !gapprompt@>| !gapinput@GeneratorsOfGroup(g),[(1,2),(1,2)]);;|
  !gapprompt@gap>| !gapinput@u:=PreImage(hom,Stabilizer(Image(hom),1));|
  Group([ (2,3,4), (1,2,4) ])
  !gapprompt@gap>| !gapinput@hom2:=GroupHomomorphismByImages(u,Group((1,2,3)),|
  !gapprompt@>| !gapinput@GeneratorsOfGroup(u),[ (1,2,3), (1,2,3) ]);;|
  !gapprompt@gap>| !gapinput@KuKGenerators(g,hom,hom2);|
  [ (1,4)(2,5)(3,6), (1,6)(2,4)(3,5) ]
\end{Verbatim}
 }

 

\subsection{\textcolor{Chapter }{ListWreathProductElement}}
\logpage{[ 49, 4, 5 ]}\nobreak
\hyperdef{L}{X801A358E879A0FF0}{}
{\noindent\textcolor{FuncColor}{$\triangleright$\enspace\texttt{ListWreathProductElement({\mdseries\slshape G, x[, testDecomposition]})\index{ListWreathProductElement@\texttt{ListWreathProductElement}}
\label{ListWreathProductElement}
}\hfill{\scriptsize (function)}}\\
\noindent\textcolor{FuncColor}{$\triangleright$\enspace\texttt{ListWreathProductElementNC({\mdseries\slshape G, x, testDecomposition})\index{ListWreathProductElementNC@\texttt{ListWreathProductElementNC}}
\label{ListWreathProductElementNC}
}\hfill{\scriptsize (operation)}}\\


 Let \mbox{\texttt{\mdseries\slshape x}} be an element of a wreath product \mbox{\texttt{\mdseries\slshape G}} where $G = K \wr H$ and $H$ acts as a finite permutation group of degree $m$. We can identify the element \mbox{\texttt{\mdseries\slshape x}} with a tuple $(f_1, \ldots, f_m; h)$, where $f_i \in K$ is the $i$\texttt{\symbol{45}}th base component of \mbox{\texttt{\mdseries\slshape x}} and $h \in H$ is the top component of \mbox{\texttt{\mdseries\slshape x}}. 

 \texttt{ListWreathProductElement} returns a list $[f_1, \ldots, f_m, h]$ containing the components of \mbox{\texttt{\mdseries\slshape x}} or \texttt{fail} if \mbox{\texttt{\mdseries\slshape x}} cannot be decomposed in the wreath product. 

 If omitted, the argument \mbox{\texttt{\mdseries\slshape testDecomposition}} defaults to true. If \mbox{\texttt{\mdseries\slshape testDecomposition}} is true, \texttt{ListWreathProductElement} makes additional tests to ensure that the computed decomposition of \mbox{\texttt{\mdseries\slshape x}} is correct, i.e. it checks that \mbox{\texttt{\mdseries\slshape x}} is an element of the parent wreath product of \mbox{\texttt{\mdseries\slshape G}}: 

 If $K \leq \mathop{Sym}(l)$, this ensures that $x \in \mathop{Sym}(l) \wr \mathop{Sym}(m)$ where the parent wreath product is considered in the same action as \mbox{\texttt{\mdseries\slshape G}}, i.e. either in imprimitive action or product action. 

 If $K \leq \mathop{GL}(n,q)$, this ensures that $x \in \mathop{GL}(n,q) \wr \mathop{Sym}(m)$. }

 

\subsection{\textcolor{Chapter }{WreathProductElementList}}
\logpage{[ 49, 4, 6 ]}\nobreak
\hyperdef{L}{X7ECB076E81D8D402}{}
{\noindent\textcolor{FuncColor}{$\triangleright$\enspace\texttt{WreathProductElementList({\mdseries\slshape G, list})\index{WreathProductElementList@\texttt{WreathProductElementList}}
\label{WreathProductElementList}
}\hfill{\scriptsize (function)}}\\
\noindent\textcolor{FuncColor}{$\triangleright$\enspace\texttt{WreathProductElementListNC({\mdseries\slshape G, list})\index{WreathProductElementListNC@\texttt{WreathProductElementListNC}}
\label{WreathProductElementListNC}
}\hfill{\scriptsize (operation)}}\\


 Let \mbox{\texttt{\mdseries\slshape list}} be equal to $[f_1, \ldots, f_m, h]$ and \mbox{\texttt{\mdseries\slshape G}} be a wreath product where $G = K \wr H$, $H$ acts as a finite permutation group of degree $m$, $f_i \in K$ and $h \in H$. 

 \texttt{WreathProductElementList} returns the element $x \in G$ identified by the tuple $(f_1, \ldots, f_m; h)$. }

 }

  
\section{\textcolor{Chapter }{Free Products}}\label{Free Products}
\logpage{[ 49, 5, 0 ]}
\hyperdef{L}{X7AC1AD17833117DF}{}
{
  Let $G$ and $H$ be groups with presentations $\langle X \mid R \rangle$ and $\langle Y \mid S \rangle$, respectively. Then the free product $G*H$ is the group with presentation $\langle X \cup Y \mid R \cup S \rangle$. This construction can be generalized to an arbitrary number of groups. 
\subsection{\textcolor{Chapter }{FreeProduct}}\logpage{[ 49, 5, 1 ]}
\hyperdef{L}{X837AC5A081EECF50}{}
{
\noindent\textcolor{FuncColor}{$\triangleright$\enspace\texttt{FreeProduct({\mdseries\slshape G[, H, ...]})\index{FreeProduct@\texttt{FreeProduct}!for several groups}
\label{FreeProduct:for several groups}
}\hfill{\scriptsize (function)}}\\
\noindent\textcolor{FuncColor}{$\triangleright$\enspace\texttt{FreeProduct({\mdseries\slshape list})\index{FreeProduct@\texttt{FreeProduct}!for a list}
\label{FreeProduct:for a list}
}\hfill{\scriptsize (function)}}\\


 constructs a finitely presented group which is the free product of the groups
given as arguments. If the group arguments are not finitely presented groups,
then \texttt{IsomorphismFpGroup} (\ref{IsomorphismFpGroup}) must be defined for them. 

 The operation \texttt{Embedding} (\ref{Embedding:for two domains}) operates on this product. 
\begin{Verbatim}[commandchars=!@|,fontsize=\small,frame=single,label=Example]
  !gapprompt@gap>| !gapinput@g := DihedralGroup(8);;|
  !gapprompt@gap>| !gapinput@h := CyclicGroup(5);;|
  !gapprompt@gap>| !gapinput@fp := FreeProduct(g,h,h);|
  <fp group on the generators [ f1, f2, f3, f4, f5 ]>
  !gapprompt@gap>| !gapinput@fp := FreeProduct([g,h,h]);|
  <fp group on the generators [ f1, f2, f3, f4, f5 ]>
  !gapprompt@gap>| !gapinput@Embedding(fp,2);|
  [ f1 ] -> [ f4 ]
\end{Verbatim}
 }

 }

  
\section{\textcolor{Chapter }{Embeddings and Projections for Group Products}}\label{Embeddings and Projections for Group Products}
\logpage{[ 49, 6, 0 ]}
\hyperdef{L}{X798FDA1386A0EAC6}{}
{
  The relation between a group product and its factors is provided via
homomorphisms, the embeddings in the product and the projections from the
product. Depending on the kind of product only some of these are defined. 

 

\subsection{\textcolor{Chapter }{Embedding (for group products)}}
\logpage{[ 49, 6, 1 ]}\nobreak
\hyperdef{L}{X784149B8847B20FF}{}
{\noindent\textcolor{FuncColor}{$\triangleright$\enspace\texttt{Embedding({\mdseries\slshape P, nr})\index{Embedding@\texttt{Embedding}!for group products}
\label{Embedding:for group products}
}\hfill{\scriptsize (operation)}}\\


 returns the \mbox{\texttt{\mdseries\slshape nr}}\texttt{\symbol{45}}th embedding in the group product \mbox{\texttt{\mdseries\slshape P}}. The actual meaning of this embedding is described in the manual section for
the appropriate product. }

 

\subsection{\textcolor{Chapter }{Projection (for group products)}}
\logpage{[ 49, 6, 2 ]}\nobreak
\hyperdef{L}{X86F275AC7C625626}{}
{\noindent\textcolor{FuncColor}{$\triangleright$\enspace\texttt{Projection({\mdseries\slshape P, nr})\index{Projection@\texttt{Projection}!for group products}
\label{Projection:for group products}
}\hfill{\scriptsize (operation)}}\\


 returns the (\mbox{\texttt{\mdseries\slshape nr}}\texttt{\symbol{45}}th) projection of the group product \mbox{\texttt{\mdseries\slshape P}}. The actual meaning of the projection returned is described in the manual
section for the appropriate product. }

 }

 }

             
\chapter{\textcolor{Chapter }{Group Libraries}}\label{Group Libraries}
\logpage{[ 50, 0, 0 ]}
\hyperdef{L}{X81B00B667D2BD022}{}
{
  When you start \textsf{GAP}, it already knows several groups. Currently \textsf{GAP} initially knows the following groups: 
\begin{itemize}
\item  some basic groups, such as cyclic groups or symmetric groups
(see{\nobreakspace}\ref{Basic Groups}), 
\item  classical matrix groups (see{\nobreakspace}\ref{Classical Groups}), 
\item  a library of transitive permutation groups, provided by the \textsf{TransGrp} package (see{\nobreakspace} (\textbf{transgrp: Transitive Permutation Groups})), 
\item  a library of groups of small order, provided by the \textsf{SmallGrp} package (see{\nobreakspace} (\textbf{smallgrp: The Small Groups Library})), 
\item  a libary of finite perfect groups, (see{\nobreakspace}\ref{Finite Perfect Groups}), 
\item  a library of primitive permutations groups, provided by the \textsf{PrimGrp} package (see{\nobreakspace} (\textbf{primgrp: Primitive Permutation Groups})), 
\item  the irreducible solvable subgroups of $GL(n,q)$ for $n>1$, $q$ a prime power and $q^n < 2^{24}$, provided by the \textsf{IRREDSOL} package (see  (\textbf{irredsol: Overview})), 
\item  the irreducible maximal finite integral matrix groups of dimension at most 31
(see{\nobreakspace}\ref{Irreducible Maximal Finite Integral Matrix Groups}), 
\item  the crystallographic groups of dimension at most 4, provided by the \textsf{CrystCat} package (see  (\textbf{CrystCat: The Crystallographic Groups Catalog})). 
\end{itemize}
 

 There is usually no relation between the groups in the different libraries and
a group may occur in different libraries in different incarnations. 

 Note that a system administrator may choose to install all, or only a few, or
even none of the libraries. So some of the libraries mentioned below may not
be available on your installation. 

 \textsf{GAP} might use data libraries that are available to speed up calculations, for
example in using a classification to determine that groups must be isomorphic,
based on agreement of properties; or to determine maximal subgroups or
subgroup maximality. This will be indicated by an info message of level 2 in
the info class \texttt{InfoPerformance}. If the calculation is to be independent of such data library use, for
example if it is used to verify the data library, functions can be called with
the option \texttt{NoPrecomputedData}, to turn these features off. Doing so might cause significantly longer
calculations, or even failure of certain calculations.  
\section{\textcolor{Chapter }{Basic Groups}}\label{Basic Groups}
\logpage{[ 50, 1, 0 ]}
\hyperdef{L}{X839981CC7D9B671B}{}
{
  There are several infinite families of groups which are parametrized by
numbers. \textsf{GAP} provides various functions to construct these groups. The functions always
permit (but do not require) one to indicate a filter (see{\nobreakspace}\ref{Filters}), for example \texttt{IsPermGroup} (\ref{IsPermGroup}), \texttt{IsMatrixGroup} (\ref{IsMatrixGroup}) or \texttt{IsPcGroup} (\ref{IsPcGroup}), in which the group shall be constructed. There always is a default filter
corresponding to a ``natural'' way to describe the group in question. Note that not every group can be
constructed in every filter, there may be theoretical restrictions (\texttt{IsPcGroup} (\ref{IsPcGroup}) only works for solvable groups) or methods may be available only for a few
filters. 

 Certain filters may admit additional hints. For example, groups constructed in \texttt{IsMatrixGroup} (\ref{IsMatrixGroup}) may be constructed over a specified field, which can be given as second
argument of the function that constructs the group; The default field is \texttt{Rationals} (\ref{Rationals}). 

 In the case of matrix groups, also the internal representation of the group
elements can be prescribed. For convenience, one can enter one of the filters
in the list \texttt{ConstructingFiltersForMatrixGroupElements} (\ref{ConstructingFiltersForMatrixGroupElements}), which then implies \texttt{IsMatrixGroup} (\ref{IsMatrixGroup}). 

\subsection{\textcolor{Chapter }{TrivialGroup}}
\logpage{[ 50, 1, 1 ]}\nobreak
\hyperdef{L}{X8489BECB78664847}{}
{\noindent\textcolor{FuncColor}{$\triangleright$\enspace\texttt{TrivialGroup({\mdseries\slshape [filter]})\index{TrivialGroup@\texttt{TrivialGroup}}
\label{TrivialGroup}
}\hfill{\scriptsize (function)}}\\


 constructs a trivial group in the category given by the filter \mbox{\texttt{\mdseries\slshape filter}}. If \mbox{\texttt{\mdseries\slshape filter}} is not given it defaults to \texttt{IsPcGroup} (\ref{IsPcGroup}). For more information on possible values of \mbox{\texttt{\mdseries\slshape filt}} see section (\ref{Basic Groups}). 

 
\begin{Verbatim}[commandchars=!@|,fontsize=\small,frame=single,label=Example]
  !gapprompt@gap>| !gapinput@TrivialGroup();|
  <pc group of size 1 with 0 generators>
  !gapprompt@gap>| !gapinput@TrivialGroup( IsPermGroup );|
  Group(())
\end{Verbatim}
 }

 

\subsection{\textcolor{Chapter }{CyclicGroup}}
\logpage{[ 50, 1, 2 ]}\nobreak
\hyperdef{L}{X7A7C473D87B31F3B}{}
{\noindent\textcolor{FuncColor}{$\triangleright$\enspace\texttt{CyclicGroup({\mdseries\slshape [filt, ]n})\index{CyclicGroup@\texttt{CyclicGroup}}
\label{CyclicGroup}
}\hfill{\scriptsize (function)}}\\


 constructs the cyclic group of size \mbox{\texttt{\mdseries\slshape n}} in the category given by the filter \mbox{\texttt{\mdseries\slshape filt}}. If \mbox{\texttt{\mdseries\slshape filt}} is not given it defaults to \texttt{IsPcGroup} (\ref{IsPcGroup}), unless \mbox{\texttt{\mdseries\slshape n}} equals \texttt{infinity} (\ref{infinity}), in which case the default filter is switched to \texttt{IsFpGroup} (\ref{IsFpGroup}). For more information on possible values of \mbox{\texttt{\mdseries\slshape filt}} see section (\ref{Basic Groups}). 

 
\begin{Verbatim}[commandchars=!@|,fontsize=\small,frame=single,label=Example]
  !gapprompt@gap>| !gapinput@CyclicGroup(12);|
  <pc group of size 12 with 3 generators>
  !gapprompt@gap>| !gapinput@CyclicGroup(infinity);|
  <free group on the generators [ a ]>
  !gapprompt@gap>| !gapinput@CyclicGroup(IsPermGroup,12);|
  Group([ (1,2,3,4,5,6,7,8,9,10,11,12) ])
  !gapprompt@gap>| !gapinput@matgrp1:= CyclicGroup( IsMatrixGroup, 12 );|
  <matrix group of size 12 with 1 generator>
  !gapprompt@gap>| !gapinput@FieldOfMatrixGroup( matgrp1 );|
  Rationals
  !gapprompt@gap>| !gapinput@matgrp2:= CyclicGroup( IsMatrixGroup, GF(2), 12 );|
  <matrix group of size 12 with 1 generator>
  !gapprompt@gap>| !gapinput@FieldOfMatrixGroup( matgrp2 );|
  GF(2)
\end{Verbatim}
 }

 

\subsection{\textcolor{Chapter }{AbelianGroup}}
\logpage{[ 50, 1, 3 ]}\nobreak
\hyperdef{L}{X81CCC3BF8005A2D7}{}
{\noindent\textcolor{FuncColor}{$\triangleright$\enspace\texttt{AbelianGroup({\mdseries\slshape [filt, ]ints})\index{AbelianGroup@\texttt{AbelianGroup}}
\label{AbelianGroup}
}\hfill{\scriptsize (function)}}\\


 constructs an abelian group in the category given by the filter \mbox{\texttt{\mdseries\slshape filt}} which is of isomorphism type $C_{{\mbox{\texttt{\mdseries\slshape ints}}[1]}} \times C_{{\mbox{\texttt{\mdseries\slshape ints}}[2]}} \times \ldots \times C_{{\mbox{\texttt{\mdseries\slshape ints}}[n]}}$, where \mbox{\texttt{\mdseries\slshape ints}} must be a list of non\texttt{\symbol{45}}negative integers or \texttt{infinity} (\ref{infinity}); for the latter value or 0, $C_{{\mbox{\texttt{\mdseries\slshape ints}}[i]}}$ is taken as an infinite cyclic group, otherwise as a cyclic group of order \mbox{\texttt{\mdseries\slshape ints}}[i]. If \mbox{\texttt{\mdseries\slshape filt}} is not given it defaults to \texttt{IsPcGroup} (\ref{IsPcGroup}), unless any 0 or \texttt{infinity} is contained in \mbox{\texttt{\mdseries\slshape ints}}, in which the default filter is switched to \texttt{IsFpGroup} (\ref{IsFpGroup}). The generators of the group returned are the elements corresponding to the
factors $C_{{\mbox{\texttt{\mdseries\slshape ints}}[i]}}$ and hence the integers in \mbox{\texttt{\mdseries\slshape ints}}. For more information on possible values of \mbox{\texttt{\mdseries\slshape filt}} see section (\ref{Basic Groups}). 

 
\begin{Verbatim}[commandchars=!@|,fontsize=\small,frame=single,label=Example]
  !gapprompt@gap>| !gapinput@AbelianGroup([1,2,3]);|
  <pc group of size 6 with 3 generators>
  !gapprompt@gap>| !gapinput@G:=AbelianGroup([0,3]);|
  <fp group of size infinity on the generators [ f1, f2 ]>
  !gapprompt@gap>| !gapinput@AbelianInvariants(G);|
  [ 0, 3 ]
\end{Verbatim}
 }

 

\subsection{\textcolor{Chapter }{ElementaryAbelianGroup}}
\logpage{[ 50, 1, 4 ]}\nobreak
\hyperdef{L}{X8778256286E50743}{}
{\noindent\textcolor{FuncColor}{$\triangleright$\enspace\texttt{ElementaryAbelianGroup({\mdseries\slshape [filt, ]n})\index{ElementaryAbelianGroup@\texttt{ElementaryAbelianGroup}}
\label{ElementaryAbelianGroup}
}\hfill{\scriptsize (function)}}\\


 constructs the elementary abelian group of size \mbox{\texttt{\mdseries\slshape n}} in the category given by the filter \mbox{\texttt{\mdseries\slshape filt}}. If \mbox{\texttt{\mdseries\slshape filt}} is not given it defaults to \texttt{IsPcGroup} (\ref{IsPcGroup}). For more information on possible values of \mbox{\texttt{\mdseries\slshape filt}} see section (\ref{Basic Groups}). 

 
\begin{Verbatim}[commandchars=!@|,fontsize=\small,frame=single,label=Example]
  !gapprompt@gap>| !gapinput@ElementaryAbelianGroup(8192);|
  <pc group of size 8192 with 13 generators>
\end{Verbatim}
 }

 

\subsection{\textcolor{Chapter }{FreeAbelianGroup}}
\logpage{[ 50, 1, 5 ]}\nobreak
\hyperdef{L}{X7F43050D8587E767}{}
{\noindent\textcolor{FuncColor}{$\triangleright$\enspace\texttt{FreeAbelianGroup({\mdseries\slshape [filt, ]rank})\index{FreeAbelianGroup@\texttt{FreeAbelianGroup}}
\label{FreeAbelianGroup}
}\hfill{\scriptsize (function)}}\\


 constructs the free abelian group of rank \mbox{\texttt{\mdseries\slshape n}} in the category given by the filter \mbox{\texttt{\mdseries\slshape filt}}. If \mbox{\texttt{\mdseries\slshape filt}} is not given it defaults to \texttt{IsFpGroup} (\ref{IsFpGroup}). For more information on possible values of \mbox{\texttt{\mdseries\slshape filt}} see section (\ref{Basic Groups}). 

 
\begin{Verbatim}[commandchars=!@|,fontsize=\small,frame=single,label=Example]
  !gapprompt@gap>| !gapinput@FreeAbelianGroup(4);|
  <fp group of size infinity on the generators [ f1, f2, f3, f4 ]>
\end{Verbatim}
 }

 

\subsection{\textcolor{Chapter }{DihedralGroup}}
\logpage{[ 50, 1, 6 ]}\nobreak
\hyperdef{L}{X838DE1AB7B3D70FF}{}
{\noindent\textcolor{FuncColor}{$\triangleright$\enspace\texttt{DihedralGroup({\mdseries\slshape [filt, ]n})\index{DihedralGroup@\texttt{DihedralGroup}}
\label{DihedralGroup}
}\hfill{\scriptsize (function)}}\\


 constructs the dihedral group of size \mbox{\texttt{\mdseries\slshape n}} in the category given by the filter \mbox{\texttt{\mdseries\slshape filt}}. If \mbox{\texttt{\mdseries\slshape filt}} is not given it defaults to \texttt{IsPcGroup} (\ref{IsPcGroup}), unless \mbox{\texttt{\mdseries\slshape n}} equals \texttt{infinity} (\ref{infinity}), in which case the default filter is switched to \texttt{IsFpGroup} (\ref{IsFpGroup}). For more information on possible values of \mbox{\texttt{\mdseries\slshape filt}} see section (\ref{Basic Groups}). 

 
\begin{Verbatim}[commandchars=!@|,fontsize=\small,frame=single,label=Example]
  !gapprompt@gap>| !gapinput@DihedralGroup(8);|
  <pc group of size 8 with 3 generators>
  !gapprompt@gap>| !gapinput@DihedralGroup( IsPermGroup, 8 );|
  Group([ (1,2,3,4), (2,4) ])
  !gapprompt@gap>| !gapinput@DihedralGroup(infinity);|
  <fp group of size infinity on the generators [ r, s ]>
\end{Verbatim}
 }

 

\subsection{\textcolor{Chapter }{IsDihedralGroup}}
\logpage{[ 50, 1, 7 ]}\nobreak
\hyperdef{L}{X8233A853818CAF33}{}
{\noindent\textcolor{FuncColor}{$\triangleright$\enspace\texttt{IsDihedralGroup({\mdseries\slshape G})\index{IsDihedralGroup@\texttt{IsDihedralGroup}}
\label{IsDihedralGroup}
}\hfill{\scriptsize (property)}}\\
\noindent\textcolor{FuncColor}{$\triangleright$\enspace\texttt{DihedralGenerators({\mdseries\slshape G})\index{DihedralGenerators@\texttt{DihedralGenerators}}
\label{DihedralGenerators}
}\hfill{\scriptsize (attribute)}}\\


 \texttt{IsDihedralGroup} indicates whether the group \mbox{\texttt{\mdseries\slshape G}} is a dihedral group. If it is, methods may set the attribute \texttt{DihedralGenerators} to [\mbox{\texttt{\mdseries\slshape t}},\mbox{\texttt{\mdseries\slshape s}}], where \mbox{\texttt{\mdseries\slshape t}} and \mbox{\texttt{\mdseries\slshape s}} are two elements such that \mbox{\texttt{\mdseries\slshape G}} = $\langle t, s | t^2 = s^n = 1, s^t = s^{-1} \rangle$. }

 

\subsection{\textcolor{Chapter }{DicyclicGroup}}
\logpage{[ 50, 1, 8 ]}\nobreak
\hyperdef{L}{X7E9844EF7C47EEB0}{}
{\noindent\textcolor{FuncColor}{$\triangleright$\enspace\texttt{DicyclicGroup({\mdseries\slshape [filt, [field, ]]n})\index{DicyclicGroup@\texttt{DicyclicGroup}}
\label{DicyclicGroup}
}\hfill{\scriptsize (function)}}\\
\noindent\textcolor{FuncColor}{$\triangleright$\enspace\texttt{QuaternionGroup({\mdseries\slshape [filt, [field, ]]n})\index{QuaternionGroup@\texttt{QuaternionGroup}}
\label{QuaternionGroup}
}\hfill{\scriptsize (function)}}\\


 \texttt{DicyclicGroup} constructs the dicyclic group of size \mbox{\texttt{\mdseries\slshape n}} in the category given by the filter \mbox{\texttt{\mdseries\slshape filt}}. Here, \mbox{\texttt{\mdseries\slshape n}} must be a multiple of 4. The synonym \texttt{QuaternionGroup} for \texttt{DicyclicGroup} is provided for backward compatibility, but will print a warning if \mbox{\texttt{\mdseries\slshape n}} is not a power of $2$. If \mbox{\texttt{\mdseries\slshape filt}} is not given it defaults to \texttt{IsPcGroup} (\ref{IsPcGroup}). For more information on possible values of \mbox{\texttt{\mdseries\slshape filt}} see section (\ref{Basic Groups}). Methods are also available for permutation and matrix groups (of minimal
degree and minimal dimension in coprime characteristic). 

 
\begin{Verbatim}[commandchars=!@|,fontsize=\small,frame=single,label=Example]
  !gapprompt@gap>| !gapinput@DicyclicGroup(24);|
  <pc group of size 24 with 4 generators>
  !gapprompt@gap>| !gapinput@g:=QuaternionGroup(IsMatrixGroup,CF(16),32);|
  Group([ [ [ 0, 1 ], [ -1, 0 ] ], [ [ E(16), 0 ], [ 0, -E(16)^7 ] ] ])
\end{Verbatim}
 }

 

\subsection{\textcolor{Chapter }{IsDicyclicGroup}}
\logpage{[ 50, 1, 9 ]}\nobreak
\hyperdef{L}{X7F260D177FD4BE4C}{}
{\noindent\textcolor{FuncColor}{$\triangleright$\enspace\texttt{IsDicyclicGroup({\mdseries\slshape G})\index{IsDicyclicGroup@\texttt{IsDicyclicGroup}}
\label{IsDicyclicGroup}
}\hfill{\scriptsize (property)}}\\
\noindent\textcolor{FuncColor}{$\triangleright$\enspace\texttt{DicyclicGenerators({\mdseries\slshape G})\index{DicyclicGenerators@\texttt{DicyclicGenerators}}
\label{DicyclicGenerators}
}\hfill{\scriptsize (attribute)}}\\
\noindent\textcolor{FuncColor}{$\triangleright$\enspace\texttt{IsGeneralisedQuaternionGroup({\mdseries\slshape G})\index{IsGeneralisedQuaternionGroup@\texttt{IsGeneralisedQuaternionGroup}}
\label{IsGeneralisedQuaternionGroup}
}\hfill{\scriptsize (property)}}\\
\noindent\textcolor{FuncColor}{$\triangleright$\enspace\texttt{GeneralisedQuaternionGenerators({\mdseries\slshape G})\index{GeneralisedQuaternionGenerators@\texttt{GeneralisedQuaternionGenerators}}
\label{GeneralisedQuaternionGenerators}
}\hfill{\scriptsize (attribute)}}\\
\noindent\textcolor{FuncColor}{$\triangleright$\enspace\texttt{IsQuaternionGroup({\mdseries\slshape G})\index{IsQuaternionGroup@\texttt{IsQuaternionGroup}}
\label{IsQuaternionGroup}
}\hfill{\scriptsize (property)}}\\
\noindent\textcolor{FuncColor}{$\triangleright$\enspace\texttt{QuaternionGenerators({\mdseries\slshape G})\index{QuaternionGenerators@\texttt{QuaternionGenerators}}
\label{QuaternionGenerators}
}\hfill{\scriptsize (attribute)}}\\


 \texttt{IsDicyclicGroup} indicates whether the group \mbox{\texttt{\mdseries\slshape G}} is a dicyclic group of order $N = 4n$. If it is, methods may set the attribute \texttt{DicyclicGenerators} to $[ t, s ]$, where $t$ and $s$ are two elements such that \mbox{\texttt{\mdseries\slshape G}} = $\langle t, s | s^{2n} = 1, t^2 = s^n, s^t = s^{-1} \rangle$ holds. 

 \texttt{IsGeneralisedQuaternionGroup} indicates whether \mbox{\texttt{\mdseries\slshape G}} is a generalized quaternion group, i. e., a dicyclic group of $2$\texttt{\symbol{45}}power order. If it is, methods may set the attribute \texttt{GeneralisedQuaternionGenerators} to the value of \texttt{DicyclicGenerators} for \mbox{\texttt{\mdseries\slshape G}}. 

 \texttt{IsQuaternionGroup} and \texttt{QuaternionGenerators} are provided for backwards compatibility with existing code. }

 

\subsection{\textcolor{Chapter }{ExtraspecialGroup}}
\logpage{[ 50, 1, 10 ]}\nobreak
\hyperdef{L}{X86E76B3A796BEFA8}{}
{\noindent\textcolor{FuncColor}{$\triangleright$\enspace\texttt{ExtraspecialGroup({\mdseries\slshape [filt, ]order, exp})\index{ExtraspecialGroup@\texttt{ExtraspecialGroup}}
\label{ExtraspecialGroup}
}\hfill{\scriptsize (function)}}\\


 Let \mbox{\texttt{\mdseries\slshape order}} be of the form $p^{{2n+1}}$, for a prime integer $p$ and a positive integer $n$. \texttt{ExtraspecialGroup} returns the extraspecial group of order \mbox{\texttt{\mdseries\slshape order}} that is determined by \mbox{\texttt{\mdseries\slshape exp}}, in the category given by the filter \mbox{\texttt{\mdseries\slshape filt}}. 

 If $p$ is odd then admissible values of \mbox{\texttt{\mdseries\slshape exp}} are the exponent of the group (either $p$ or $p^2$) or one of \texttt{'+'}, \texttt{"+"}, \texttt{'\texttt{\symbol{45}}'}, \texttt{"\texttt{\symbol{45}}"}. For $p = 2$, only the above plus or minus signs are admissible. 

 If \mbox{\texttt{\mdseries\slshape filt}} is not given it defaults to \texttt{IsPcGroup} (\ref{IsPcGroup}). For more information on possible values of \mbox{\texttt{\mdseries\slshape filt}} see section (\ref{Basic Groups}). 

 
\begin{Verbatim}[commandchars=!@|,fontsize=\small,frame=single,label=Example]
  !gapprompt@gap>| !gapinput@ExtraspecialGroup( 27, 3 );|
  <pc group of size 27 with 3 generators>
  !gapprompt@gap>| !gapinput@ExtraspecialGroup( 27, '+' );|
  <pc group of size 27 with 3 generators>
  !gapprompt@gap>| !gapinput@ExtraspecialGroup( 8, "-" );|
  <pc group of size 8 with 3 generators>
\end{Verbatim}
 }

 
\subsection{\textcolor{Chapter }{AlternatingGroup}}\logpage{[ 50, 1, 11 ]}
\hyperdef{L}{X7E54D3E778E6A53E}{}
{
\noindent\textcolor{FuncColor}{$\triangleright$\enspace\texttt{AlternatingGroup({\mdseries\slshape [filt, ]deg})\index{AlternatingGroup@\texttt{AlternatingGroup}!for a degree}
\label{AlternatingGroup:for a degree}
}\hfill{\scriptsize (function)}}\\
\noindent\textcolor{FuncColor}{$\triangleright$\enspace\texttt{AlternatingGroup({\mdseries\slshape [filt, ]dom})\index{AlternatingGroup@\texttt{AlternatingGroup}!for a domain}
\label{AlternatingGroup:for a domain}
}\hfill{\scriptsize (function)}}\\


 constructs the alternating group of degree \mbox{\texttt{\mdseries\slshape deg}} in the category given by the filter \mbox{\texttt{\mdseries\slshape filt}}. If \mbox{\texttt{\mdseries\slshape filt}} is not given it defaults to \texttt{IsPermGroup} (\ref{IsPermGroup}). For more information on possible values of \mbox{\texttt{\mdseries\slshape filt}} see section (\ref{Basic Groups}). In the second version, the function constructs the alternating group on the
points given in the set \mbox{\texttt{\mdseries\slshape dom}} which must be a set of positive integers. 
\begin{Verbatim}[commandchars=!@|,fontsize=\small,frame=single,label=Example]
  !gapprompt@gap>| !gapinput@AlternatingGroup(5);|
  Alt( [ 1 .. 5 ] )
\end{Verbatim}
 }

 
\subsection{\textcolor{Chapter }{SymmetricGroup}}\logpage{[ 50, 1, 12 ]}
\hyperdef{L}{X858666F97BD85ABB}{}
{
\noindent\textcolor{FuncColor}{$\triangleright$\enspace\texttt{SymmetricGroup({\mdseries\slshape [filt, ]deg})\index{SymmetricGroup@\texttt{SymmetricGroup}!for a degree}
\label{SymmetricGroup:for a degree}
}\hfill{\scriptsize (function)}}\\
\noindent\textcolor{FuncColor}{$\triangleright$\enspace\texttt{SymmetricGroup({\mdseries\slshape [filt, ]dom})\index{SymmetricGroup@\texttt{SymmetricGroup}!for a domain}
\label{SymmetricGroup:for a domain}
}\hfill{\scriptsize (function)}}\\


 constructs the symmetric group of degree \mbox{\texttt{\mdseries\slshape deg}} in the category given by the filter \mbox{\texttt{\mdseries\slshape filt}}. If \mbox{\texttt{\mdseries\slshape filt}} is not given it defaults to \texttt{IsPermGroup} (\ref{IsPermGroup}). For more information on possible values of \mbox{\texttt{\mdseries\slshape filt}} see section (\ref{Basic Groups}). In the second version, the function constructs the symmetric group on the
points given in the set \mbox{\texttt{\mdseries\slshape dom}} which must be a set of positive integers. 

 
\begin{Verbatim}[commandchars=!@|,fontsize=\small,frame=single,label=Example]
  !gapprompt@gap>| !gapinput@SymmetricGroup(10);|
  Sym( [ 1 .. 10 ] )
\end{Verbatim}
 

 Note that permutation groups provide special treatment of symmetric and
alternating groups, see{\nobreakspace}\ref{Symmetric and Alternating Groups}. }

 

\subsection{\textcolor{Chapter }{MathieuGroup}}
\logpage{[ 50, 1, 13 ]}\nobreak
\hyperdef{L}{X788FA7DE84E0FE6A}{}
{\noindent\textcolor{FuncColor}{$\triangleright$\enspace\texttt{MathieuGroup({\mdseries\slshape [filt, ]degree})\index{MathieuGroup@\texttt{MathieuGroup}}
\label{MathieuGroup}
}\hfill{\scriptsize (function)}}\\


 constructs the Mathieu group of degree \mbox{\texttt{\mdseries\slshape degree}} in the category given by the filter \mbox{\texttt{\mdseries\slshape filt}}, where \mbox{\texttt{\mdseries\slshape degree}} must be in the set $\{ 9, 10, 11, 12, 21, 22, 23, 24 \}$. If \mbox{\texttt{\mdseries\slshape filt}} is not given it defaults to \texttt{IsPermGroup} (\ref{IsPermGroup}). For more information on possible values of \mbox{\texttt{\mdseries\slshape filt}} see section (\ref{Basic Groups}). 

 
\begin{Verbatim}[commandchars=!@|,fontsize=\small,frame=single,label=Example]
  !gapprompt@gap>| !gapinput@MathieuGroup( 11 );|
  Group([ (1,2,3,4,5,6,7,8,9,10,11), (3,7,11,8)(4,10,5,6) ])
\end{Verbatim}
 }

 

\subsection{\textcolor{Chapter }{SuzukiGroup}}
\logpage{[ 50, 1, 14 ]}\nobreak
\hyperdef{L}{X8469DBBF82F8E5C3}{}
{\noindent\textcolor{FuncColor}{$\triangleright$\enspace\texttt{SuzukiGroup({\mdseries\slshape [filt, ]q})\index{SuzukiGroup@\texttt{SuzukiGroup}}
\label{SuzukiGroup}
}\hfill{\scriptsize (function)}}\\
\noindent\textcolor{FuncColor}{$\triangleright$\enspace\texttt{Sz({\mdseries\slshape [filt, ]q})\index{Sz@\texttt{Sz}}
\label{Sz}
}\hfill{\scriptsize (function)}}\\


 Constructs a group isomorphic to the Suzuki group Sz( \mbox{\texttt{\mdseries\slshape q}} ) over the field with \mbox{\texttt{\mdseries\slshape q}} elements, where \mbox{\texttt{\mdseries\slshape q}} is a non\texttt{\symbol{45}}square power of $2$. 

 If \mbox{\texttt{\mdseries\slshape filt}} is not given it defaults to \texttt{IsMatrixGroup} (\ref{IsMatrixGroup}), and the returned group is the Suzuki group itself. For more information on
possible values of \mbox{\texttt{\mdseries\slshape filt}} see section (\ref{Basic Groups}). 
\begin{Verbatim}[commandchars=!@|,fontsize=\small,frame=single,label=Example]
  !gapprompt@gap>| !gapinput@SuzukiGroup( 32 );|
  Sz(32)
\end{Verbatim}
 }

 

\subsection{\textcolor{Chapter }{ReeGroup}}
\logpage{[ 50, 1, 15 ]}\nobreak
\hyperdef{L}{X87E5B0F679CA7FE4}{}
{\noindent\textcolor{FuncColor}{$\triangleright$\enspace\texttt{ReeGroup({\mdseries\slshape [filt, ]q})\index{ReeGroup@\texttt{ReeGroup}}
\label{ReeGroup}
}\hfill{\scriptsize (function)}}\\
\noindent\textcolor{FuncColor}{$\triangleright$\enspace\texttt{Ree({\mdseries\slshape [filt, ]q})\index{Ree@\texttt{Ree}}
\label{Ree}
}\hfill{\scriptsize (function)}}\\


 Constructs a group isomorphic to the Ree group $^2G_2(q)$ where $q = 3^{{1+2m}}$ for $m$ a non\texttt{\symbol{45}}negative integer. 

 If \mbox{\texttt{\mdseries\slshape filt}} is not given it defaults to \texttt{IsMatrixGroup} (\ref{IsMatrixGroup}) and the generating matrices are based on{\nobreakspace}\cite{KLM01}. (No particular choice of a generating set is guaranteed.) For more
information on possible values of \mbox{\texttt{\mdseries\slshape filt}} see section (\ref{Basic Groups}). 

 
\begin{Verbatim}[commandchars=!@|,fontsize=\small,frame=single,label=Example]
  !gapprompt@gap>| !gapinput@ReeGroup( 27 );|
  Ree(27)
\end{Verbatim}
 }

 
\subsection{\textcolor{Chapter }{Generator Names}}\label{Generator Names}
\logpage{[ 50, 1, 16 ]}
\hyperdef{L}{X7D0FFDA4793995FC}{}
{
  For groups created as finitely presented groups, including polycyclic groups,
the generators are labelled, by default, with a letter and a number. It is
possible to influence this naming with the option \texttt{generatorNames}, see Section{\nobreakspace}\ref{Function Call With Options}. If this option holds a string, then the generators are named with this
string and sequential numbers starting with \texttt{1}. If this option holds a list of sufficient length consisting of nonempty
strings, then the generator names are taken from this list, in order. 

 
\begin{Verbatim}[commandchars=!@|,fontsize=\small,frame=single,label=Example]
  !gapprompt@gap>| !gapinput@GeneratorsOfGroup(AbelianGroup([5,7]));|
  [ f1, f2 ]
  !gapprompt@gap>| !gapinput@GeneratorsOfGroup(AbelianGroup([5,7]:generatorNames:="a"));|
  [ a1, a2 ]
  !gapprompt@gap>| !gapinput@GeneratorsOfGroup(AbelianGroup([5,7]:generatorNames:=["u","v","w"]));|
  [ u, v ]
  !gapprompt@gap>| !gapinput@AsSSortedList(DihedralGroup(12:generatorNames:="a"));|
  [ <identity> of ..., a1, a2, a3, a1*a2, a1*a3, a2*a3, a3^2, a1*a2*a3,
    a1*a3^2, a2*a3^2, a1*a2*a3^2 ]
  !gapprompt@gap>| !gapinput@AsSSortedList(DihedralGroup(12:generatorNames:=["a","b","c"]));|
  [ <identity> of ..., a, b, c, a*b, a*c, b*c, c^2, a*b*c, a*c^2, b*c^2,
    a*b*c^2 ]
\end{Verbatim}
 }

 }

  
\section{\textcolor{Chapter }{Classical Groups}}\label{Classical Groups}
\logpage{[ 50, 2, 0 ]}
\hyperdef{L}{X8674AAA578FE4AEE}{}
{
  The following functions return classical groups. 

 \emph{Generators} 

 
\begin{itemize}
\item  For the linear, symplectic, and unitary groups (the latter in dimension at
least $3$), the generators are taken from{\nobreakspace}\cite{Tay87}. 
\item  For the unitary groups in dimension $2$, the isomorphism of SU$(2,q)$ and SL$(2,q)$ is used, see for example{\nobreakspace}\cite{Hup67}. 
\item  The generators of the general and special orthogonal groups are taken
from{\nobreakspace}\cite{IshibashiEarnest94} and \cite{KleidmanLiebeck90}, except that the generators of the groups in odd dimension in even
characteristic are constructed via the isomorphism to a symplectic group, see
for example{\nobreakspace}\cite{Car72a}. 
\item  The generators of the groups $\Omega^\epsilon(d, q)$ are taken from{\nobreakspace}\cite{RylandsTalor98}, except that in odd dimension and even characteristic, the generators of SO$(d, q)$ are taken for $\Omega(d, q)$. Note that the generators claimed in{\nobreakspace}\cite[Section 4.5 and 4.6]{RylandsTalor98} do not describe orthogonal groups, one would have to transpose these matrices
in order to get groups that respect the required forms. The matrices
from{\nobreakspace}\cite{RylandsTalor98} generate groups of the right isomorphism types but not orthogonal groups,
except in the case $(d,q) = (5,2)$, where the matrices from{\nobreakspace}\cite{RylandsTalor98} generate the simple group $S_4(2)'$ and not the group $S_4(2)$. 
\item  The generators for the semilinear groups are constructed from the generators
of the corresponding linear groups plus one additional generator that
describes the action of the group of field automorphisms; for prime integers $p$ and positive integers $f$, this yields the matrix groups $Gamma$L$(d, p^f)$ and $Sigma$L$(d, p^f)$ as groups of $d f \times df$ matrices over the field with $p$ elements. 
\item  The generators for the conformal symplectic groups are taken
from{\nobreakspace}\cite{Tay21}. 
\end{itemize}
 

 \emph{Invariant forms} 

 
\begin{itemize}
\item  For symplectic and orthogonal matrix groups returned by the functions
described below, the invariant bilinear form is stored as the value of the
attribute \texttt{InvariantBilinearForm} (\ref{InvariantBilinearForm}). 
\item  The invariant sesquilinear form defining the unitary groups is stored as the
value of the attribute \texttt{InvariantSesquilinearForm} (\ref{InvariantSesquilinearForm})). 
\item  The defining quadratic form of orthogonal groups is stored as the value of the
attribute \texttt{InvariantQuadraticForm} (\ref{InvariantQuadraticForm}). 
\item  The bilinear form that is invariant up to scalars under the conformal
symplectic group is stored as the value of the attribute \texttt{InvariantBilinearFormUpToScalars} (\ref{InvariantBilinearFormUpToScalars}). 
\end{itemize}
 

 Note that due to the different sources for the generators, the invariant forms
for the groups $\Omega(e,d,q)$ are in general different from the forms for SO$(e,d,q)$ and GO$(e,d,q)$. If version at least 1.2.6 of the \textsf{Forms} package is loaded then compatible groups can be created by specifying the
desired form, see the sections below. 
\subsection{\textcolor{Chapter }{GeneralLinearGroup}}\logpage{[ 50, 2, 1 ]}
\hyperdef{L}{X85D607DD82AF3E27}{}
{
\noindent\textcolor{FuncColor}{$\triangleright$\enspace\texttt{GeneralLinearGroup({\mdseries\slshape [filt, ]d, R})\index{GeneralLinearGroup@\texttt{GeneralLinearGroup}!for dimension and a ring}
\label{GeneralLinearGroup:for dimension and a ring}
}\hfill{\scriptsize (function)}}\\
\noindent\textcolor{FuncColor}{$\triangleright$\enspace\texttt{GL({\mdseries\slshape [filt, ]d, R})\index{GL@\texttt{GL}!for dimension and a ring}
\label{GL:for dimension and a ring}
}\hfill{\scriptsize (function)}}\\
\noindent\textcolor{FuncColor}{$\triangleright$\enspace\texttt{GeneralLinearGroup({\mdseries\slshape [filt, ]d, q})\index{GeneralLinearGroup@\texttt{GeneralLinearGroup}!for dimension and field size}
\label{GeneralLinearGroup:for dimension and field size}
}\hfill{\scriptsize (function)}}\\
\noindent\textcolor{FuncColor}{$\triangleright$\enspace\texttt{GL({\mdseries\slshape [filt, ]d, q})\index{GL@\texttt{GL}!for dimension and field size}
\label{GL:for dimension and field size}
}\hfill{\scriptsize (function)}}\\


 The first two forms construct a group isomorphic to the general linear group
GL( \mbox{\texttt{\mdseries\slshape d}}, \mbox{\texttt{\mdseries\slshape R}} ) of all $\mbox{\texttt{\mdseries\slshape d}} \times \mbox{\texttt{\mdseries\slshape d}}$ matrices that are invertible over the ring \mbox{\texttt{\mdseries\slshape R}}, in the category given by the filter \mbox{\texttt{\mdseries\slshape filt}}. 

 The third and the fourth form construct the general linear group over the
finite field with \mbox{\texttt{\mdseries\slshape q}} elements. 

 If \mbox{\texttt{\mdseries\slshape filt}} is not given it defaults to \texttt{IsMatrixGroup} (\ref{IsMatrixGroup}), and the returned group is the general linear group as a matrix group in its
natural action (see also{\nobreakspace}\texttt{IsNaturalGL} (\ref{IsNaturalGL}), \texttt{IsNaturalGLnZ} (\ref{IsNaturalGLnZ})). 

 Currently supported rings \mbox{\texttt{\mdseries\slshape R}} are finite fields, the ring \texttt{Integers} (\ref{Integers:global variable}), and residue class rings \texttt{Integers mod \mbox{\texttt{\mdseries\slshape m}}}, see \ref{Residue Class Rings}. 

 
\begin{Verbatim}[commandchars=!@|,fontsize=\small,frame=single,label=Example]
  !gapprompt@gap>| !gapinput@GL(4,3);|
  GL(4,3)
  !gapprompt@gap>| !gapinput@GL(2,Integers);|
  GL(2,Integers)
  !gapprompt@gap>| !gapinput@GL(3,Integers mod 12);|
  GL(3,Z/12Z)
\end{Verbatim}
 

 \index{OnLines@\texttt{OnLines}!example} Using the \texttt{OnLines} (\ref{OnLines}) operation it is possible to obtain the corresponding projective groups in a
permutation action: 

 
\begin{Verbatim}[commandchars=!@|,fontsize=\small,frame=single,label=Example]
  !gapprompt@gap>| !gapinput@g:=GL(4,3);;Size(g);|
  24261120
  !gapprompt@gap>| !gapinput@pgl:=Action(g,Orbit(g,Z(3)^0*[1,0,0,0],OnLines),OnLines);;|
  !gapprompt@gap>| !gapinput@Size(pgl);|
  12130560
\end{Verbatim}
 

 If you are interested only in the projective group as a permutation group and
not in the correspondence between its moved points and the points in the
projective space, you can also use \texttt{PGL} (\ref{PGL}). }

 
\subsection{\textcolor{Chapter }{SpecialLinearGroup}}\logpage{[ 50, 2, 2 ]}
\hyperdef{L}{X7CA3F7BF83992C6B}{}
{
\noindent\textcolor{FuncColor}{$\triangleright$\enspace\texttt{SpecialLinearGroup({\mdseries\slshape [filt, ]d, R})\index{SpecialLinearGroup@\texttt{SpecialLinearGroup}!for dimension and a ring}
\label{SpecialLinearGroup:for dimension and a ring}
}\hfill{\scriptsize (function)}}\\
\noindent\textcolor{FuncColor}{$\triangleright$\enspace\texttt{SL({\mdseries\slshape [filt, ]d, R})\index{SL@\texttt{SL}!for dimension and a ring}
\label{SL:for dimension and a ring}
}\hfill{\scriptsize (function)}}\\
\noindent\textcolor{FuncColor}{$\triangleright$\enspace\texttt{SpecialLinearGroup({\mdseries\slshape [filt, ]d, q})\index{SpecialLinearGroup@\texttt{SpecialLinearGroup}!for dimension and a field size}
\label{SpecialLinearGroup:for dimension and a field size}
}\hfill{\scriptsize (function)}}\\
\noindent\textcolor{FuncColor}{$\triangleright$\enspace\texttt{SL({\mdseries\slshape [filt, ]d, q})\index{SL@\texttt{SL}!for dimension and a field size}
\label{SL:for dimension and a field size}
}\hfill{\scriptsize (function)}}\\


 The first two forms construct a group isomorphic to the special linear group
SL( \mbox{\texttt{\mdseries\slshape d}}, \mbox{\texttt{\mdseries\slshape R}} ) of all those $\mbox{\texttt{\mdseries\slshape d}} \times \mbox{\texttt{\mdseries\slshape d}}$ matrices over the ring \mbox{\texttt{\mdseries\slshape R}} whose determinant is the identity of \mbox{\texttt{\mdseries\slshape R}}, in the category given by the filter \mbox{\texttt{\mdseries\slshape filt}}. 

 The third and the fourth form construct the special linear group over the
finite field with \mbox{\texttt{\mdseries\slshape q}} elements. 

 If \mbox{\texttt{\mdseries\slshape filt}} is not given it defaults to \texttt{IsMatrixGroup} (\ref{IsMatrixGroup}), and the returned group is the special linear group as a matrix group in its
natural action (see also{\nobreakspace}\texttt{IsNaturalSL} (\ref{IsNaturalSL}), \texttt{IsNaturalSLnZ} (\ref{IsNaturalSLnZ})). 

 Currently supported rings \mbox{\texttt{\mdseries\slshape R}} are finite fields, the ring \texttt{Integers} (\ref{Integers}), and residue class rings \texttt{Integers mod \mbox{\texttt{\mdseries\slshape m}}}, see \ref{Residue Class Rings}. 

 
\begin{Verbatim}[commandchars=!@|,fontsize=\small,frame=single,label=Example]
  !gapprompt@gap>| !gapinput@SpecialLinearGroup(2,2);|
  SL(2,2)
  !gapprompt@gap>| !gapinput@SL(3,Integers);|
  SL(3,Integers)
  !gapprompt@gap>| !gapinput@SL(4,Integers mod 4);|
  SL(4,Z/4Z)
\end{Verbatim}
 }

 

\subsection{\textcolor{Chapter }{GeneralUnitaryGroup}}
\logpage{[ 50, 2, 3 ]}\nobreak
\hyperdef{L}{X866D4E2B816BDFA5}{}
{\noindent\textcolor{FuncColor}{$\triangleright$\enspace\texttt{GeneralUnitaryGroup({\mdseries\slshape [filt, ]d, q[, form]})\index{GeneralUnitaryGroup@\texttt{GeneralUnitaryGroup}}
\label{GeneralUnitaryGroup}
}\hfill{\scriptsize (function)}}\\
\noindent\textcolor{FuncColor}{$\triangleright$\enspace\texttt{GeneralUnitaryGroup({\mdseries\slshape [filt, ]form})\index{GeneralUnitaryGroup@\texttt{GeneralUnitaryGroup}!for a form}
\label{GeneralUnitaryGroup:for a form}
}\hfill{\scriptsize (function)}}\\
\noindent\textcolor{FuncColor}{$\triangleright$\enspace\texttt{GU({\mdseries\slshape [filt, ]d, q[, form]})\index{GU@\texttt{GU}}
\label{GU}
}\hfill{\scriptsize (function)}}\\
\noindent\textcolor{FuncColor}{$\triangleright$\enspace\texttt{GU({\mdseries\slshape [filt, ]form})\index{GU@\texttt{GU}!for a form}
\label{GU:for a form}
}\hfill{\scriptsize (function)}}\\


 constructs a group isomorphic to the general unitary group GU( \mbox{\texttt{\mdseries\slshape d}}, \mbox{\texttt{\mdseries\slshape q}} ) of those $\mbox{\texttt{\mdseries\slshape d}} \times \mbox{\texttt{\mdseries\slshape d}}$ matrices over the field with $\mbox{\texttt{\mdseries\slshape q}}^2$ elements that respect a fixed nondegenerate sesquilinear form, in the category
given by the filter \mbox{\texttt{\mdseries\slshape filt}}. 

 If \mbox{\texttt{\mdseries\slshape filt}} is not given it defaults to \texttt{IsMatrixGroup} (\ref{IsMatrixGroup}), and the returned group is the general unitary group itself. 

 If version at least 1.2.6 of the \textsf{Forms} package is loaded then the desired sesquilinear form can be specified via \mbox{\texttt{\mdseries\slshape form}}, which can be either a matrix or a form object in \texttt{IsHermitianForm} (\textbf{Forms: IsHermitianForm}) or a group with stored \texttt{InvariantSesquilinearForm} (\ref{InvariantSesquilinearForm}) value (and then this form is taken). 

 A given \mbox{\texttt{\mdseries\slshape form}} determines \mbox{\texttt{\mdseries\slshape d}}, and also \mbox{\texttt{\mdseries\slshape q}} except if \mbox{\texttt{\mdseries\slshape form}} is a matrix that does not store its \texttt{BaseDomain} (\ref{BaseDomain:for a matrix object}) value. These parameters can be entered, and an error is signalled if they do
not fit to the given \mbox{\texttt{\mdseries\slshape form}}. 

 If \mbox{\texttt{\mdseries\slshape form}} is not given then a default is chosen as described in the introduction to
Section \ref{Classical Groups}. 

 
\begin{Verbatim}[commandchars=!@|,fontsize=\small,frame=single,label=Example]
  !gapprompt@gap>| !gapinput@GeneralUnitaryGroup( 3, 5 );|
  GU(3,5)
  !gapprompt@gap>| !gapinput@GeneralUnitaryGroup( IsPermGroup, 3, 5 );|
  Perm_GU(3,5)
\end{Verbatim}
 }

 

\subsection{\textcolor{Chapter }{SpecialUnitaryGroup}}
\logpage{[ 50, 2, 4 ]}\nobreak
\hyperdef{L}{X82A2AADE805DCDE9}{}
{\noindent\textcolor{FuncColor}{$\triangleright$\enspace\texttt{SpecialUnitaryGroup({\mdseries\slshape [filt, ]d, q[, form]})\index{SpecialUnitaryGroup@\texttt{SpecialUnitaryGroup}}
\label{SpecialUnitaryGroup}
}\hfill{\scriptsize (function)}}\\
\noindent\textcolor{FuncColor}{$\triangleright$\enspace\texttt{SpecialUnitaryGroup({\mdseries\slshape [filt, ]form})\index{SpecialUnitaryGroup@\texttt{SpecialUnitaryGroup}!for a form}
\label{SpecialUnitaryGroup:for a form}
}\hfill{\scriptsize (function)}}\\
\noindent\textcolor{FuncColor}{$\triangleright$\enspace\texttt{SU({\mdseries\slshape [filt, ]d, q[, form]})\index{SU@\texttt{SU}}
\label{SU}
}\hfill{\scriptsize (function)}}\\
\noindent\textcolor{FuncColor}{$\triangleright$\enspace\texttt{SU({\mdseries\slshape [filt, ]form})\index{SU@\texttt{SU}!for a form}
\label{SU:for a form}
}\hfill{\scriptsize (function)}}\\


 constructs a group isomorphic to the special unitary group SU( \mbox{\texttt{\mdseries\slshape d}}, \mbox{\texttt{\mdseries\slshape q}} ) of those $\mbox{\texttt{\mdseries\slshape d}} \times \mbox{\texttt{\mdseries\slshape d}}$ matrices over the field with $\mbox{\texttt{\mdseries\slshape q}}^2$ elements whose determinant is the identity of the field and that respect a
fixed nondegenerate sesquilinear form, in the category given by the filter \mbox{\texttt{\mdseries\slshape filt}}. 

 If \mbox{\texttt{\mdseries\slshape filt}} is not given it defaults to \texttt{IsMatrixGroup} (\ref{IsMatrixGroup}), and the returned group is the special unitary group itself. 

 If version at least 1.2.6 of the \textsf{Forms} package is loaded then the desired sesquilinear form can be specified via \mbox{\texttt{\mdseries\slshape form}}, which can be either a matrix or a form object in \texttt{IsHermitianForm} (\textbf{Forms: IsHermitianForm}) or a group with stored \texttt{InvariantSesquilinearForm} (\ref{InvariantSesquilinearForm}) value (and then this form is taken). 

 A given \mbox{\texttt{\mdseries\slshape form}} determines \mbox{\texttt{\mdseries\slshape d}}, and also \mbox{\texttt{\mdseries\slshape q}} except if \mbox{\texttt{\mdseries\slshape form}} is a matrix that does not store its \texttt{BaseDomain} (\ref{BaseDomain:for a matrix object}) value. These parameters can be entered, and an error is signalled if they do
not fit to the given \mbox{\texttt{\mdseries\slshape form}}. 

 If \mbox{\texttt{\mdseries\slshape form}} is not given then a default is chosen as described in the introduction to
Section \ref{Classical Groups}. 

 
\begin{Verbatim}[commandchars=!@|,fontsize=\small,frame=single,label=Example]
  !gapprompt@gap>| !gapinput@SpecialUnitaryGroup( 3, 5 );|
  SU(3,5)
  !gapprompt@gap>| !gapinput@SpecialUnitaryGroup( IsPermGroup, 3, 5 );|
  Perm_SU(3,5)
\end{Verbatim}
 }

 
\subsection{\textcolor{Chapter }{SymplecticGroup}}\logpage{[ 50, 2, 5 ]}
\hyperdef{L}{X8142A8B07811CA90}{}
{
\noindent\textcolor{FuncColor}{$\triangleright$\enspace\texttt{SymplecticGroup({\mdseries\slshape [filt, ]d, R[, form]})\index{SymplecticGroup@\texttt{SymplecticGroup}!for dimension and a ring}
\label{SymplecticGroup:for dimension and a ring}
}\hfill{\scriptsize (function)}}\\
\noindent\textcolor{FuncColor}{$\triangleright$\enspace\texttt{SymplecticGroup({\mdseries\slshape [filt, ]d, q[, form]})\index{SymplecticGroup@\texttt{SymplecticGroup}!for dimension and field size}
\label{SymplecticGroup:for dimension and field size}
}\hfill{\scriptsize (function)}}\\
\noindent\textcolor{FuncColor}{$\triangleright$\enspace\texttt{SymplecticGroup({\mdseries\slshape [filt, ]form})\index{SymplecticGroup@\texttt{SymplecticGroup}!for form}
\label{SymplecticGroup:for form}
}\hfill{\scriptsize (function)}}\\
\noindent\textcolor{FuncColor}{$\triangleright$\enspace\texttt{Sp({\mdseries\slshape [filt, ]d, R[, form]})\index{Sp@\texttt{Sp}!for dimension and a ring}
\label{Sp:for dimension and a ring}
}\hfill{\scriptsize (function)}}\\
\noindent\textcolor{FuncColor}{$\triangleright$\enspace\texttt{Sp({\mdseries\slshape [filt, ]d, q[, form]})\index{Sp@\texttt{Sp}!for dimension and field size}
\label{Sp:for dimension and field size}
}\hfill{\scriptsize (function)}}\\
\noindent\textcolor{FuncColor}{$\triangleright$\enspace\texttt{Sp({\mdseries\slshape [filt, ]form})\index{Sp@\texttt{Sp}!for form}
\label{Sp:for form}
}\hfill{\scriptsize (function)}}\\
\noindent\textcolor{FuncColor}{$\triangleright$\enspace\texttt{SP({\mdseries\slshape [filt, ]d, R[, form]})\index{SP@\texttt{SP}!for dimension and a ring}
\label{SP:for dimension and a ring}
}\hfill{\scriptsize (function)}}\\
\noindent\textcolor{FuncColor}{$\triangleright$\enspace\texttt{SP({\mdseries\slshape [filt, ]d, q[, form]})\index{SP@\texttt{SP}!for dimension and field size}
\label{SP:for dimension and field size}
}\hfill{\scriptsize (function)}}\\
\noindent\textcolor{FuncColor}{$\triangleright$\enspace\texttt{SP({\mdseries\slshape [filt, ]form})\index{SP@\texttt{SP}!for form}
\label{SP:for form}
}\hfill{\scriptsize (function)}}\\


 constructs a group isomorphic to the symplectic group Sp( \mbox{\texttt{\mdseries\slshape d}}, \mbox{\texttt{\mdseries\slshape R}} ) of those $\mbox{\texttt{\mdseries\slshape d}} \times \mbox{\texttt{\mdseries\slshape d}}$ matrices over the ring \mbox{\texttt{\mdseries\slshape R}} or the field with \mbox{\texttt{\mdseries\slshape q}} elements, respectively, that respect a fixed nondegenerate symplectic form, in
the category given by the filter \mbox{\texttt{\mdseries\slshape filt}}. 

 If \mbox{\texttt{\mdseries\slshape filt}} is not given it defaults to \texttt{IsMatrixGroup} (\ref{IsMatrixGroup}), and the returned group is the symplectic group itself. Other supported
values for \mbox{\texttt{\mdseries\slshape filt}} are \texttt{IsPermGroup} (\ref{IsPermGroup}) (in this case, the argument \mbox{\texttt{\mdseries\slshape form}} is not supported) and the entries of \texttt{ConstructingFiltersForMatrixGroupElements} (\ref{ConstructingFiltersForMatrixGroupElements}). 

 At the moment finite fields or residue class rings \texttt{Integers mod \mbox{\texttt{\mdseries\slshape q}}}, with \mbox{\texttt{\mdseries\slshape q}} an odd prime power, are supported. 

 If version at least 1.2.6 of the \textsf{Forms} package is loaded and the arguments describe a matrix group over a finite
field then the desired bilinear form can be specified via \mbox{\texttt{\mdseries\slshape form}}, which can be either a matrix or a form object in \texttt{IsBilinearForm} (\textbf{Forms: IsBilinearForm}) or a group with stored \texttt{InvariantBilinearForm} (\ref{InvariantBilinearForm}) value (and then this form is taken). 

 A given \mbox{\texttt{\mdseries\slshape form}} determines \mbox{\texttt{\mdseries\slshape d}}, and also \mbox{\texttt{\mdseries\slshape R}} except if \mbox{\texttt{\mdseries\slshape form}} is a matrix that does not store its \texttt{BaseDomain} (\ref{BaseDomain:for a matrix object}) value. These parameters can be entered, and an error is signalled if they do
not fit to the given \mbox{\texttt{\mdseries\slshape form}}. 

 If \mbox{\texttt{\mdseries\slshape form}} is not given then a default is chosen as described in the introduction to
Section \ref{Classical Groups}. 

 
\begin{Verbatim}[commandchars=!@|,fontsize=\small,frame=single,label=Example]
  !gapprompt@gap>| !gapinput@SymplecticGroup( 4, 2 );|
  Sp(4,2)
  !gapprompt@gap>| !gapinput@g:=SymplecticGroup(6,Integers mod 9);|
  Sp(6,Z/9Z)
  !gapprompt@gap>| !gapinput@Size(g);|
  95928796265538862080
  !gapprompt@gap>| !gapinput@SymplecticGroup( IsPermGroup, 4, 2 );|
  Perm_Sp(4,2)
\end{Verbatim}
 }

 

\subsection{\textcolor{Chapter }{GeneralOrthogonalGroup}}
\logpage{[ 50, 2, 6 ]}\nobreak
\hyperdef{L}{X7C2051CB7B94CEB1}{}
{\noindent\textcolor{FuncColor}{$\triangleright$\enspace\texttt{GeneralOrthogonalGroup({\mdseries\slshape [filt, ][e, ]d, q[, form]})\index{GeneralOrthogonalGroup@\texttt{GeneralOrthogonalGroup}}
\label{GeneralOrthogonalGroup}
}\hfill{\scriptsize (function)}}\\
\noindent\textcolor{FuncColor}{$\triangleright$\enspace\texttt{GeneralOrthogonalGroup({\mdseries\slshape [filt, ]form})\index{GeneralOrthogonalGroup@\texttt{GeneralOrthogonalGroup}!for a form}
\label{GeneralOrthogonalGroup:for a form}
}\hfill{\scriptsize (function)}}\\
\noindent\textcolor{FuncColor}{$\triangleright$\enspace\texttt{GO({\mdseries\slshape [filt, ][e, ]d, q[, form]})\index{GO@\texttt{GO}}
\label{GO}
}\hfill{\scriptsize (function)}}\\
\noindent\textcolor{FuncColor}{$\triangleright$\enspace\texttt{GO({\mdseries\slshape [filt, ]form})\index{GO@\texttt{GO}!for a form}
\label{GO:for a form}
}\hfill{\scriptsize (function)}}\\


 constructs a group isomorphic to the general orthogonal group GO( \mbox{\texttt{\mdseries\slshape e}}, \mbox{\texttt{\mdseries\slshape d}}, \mbox{\texttt{\mdseries\slshape q}} ) of those $\mbox{\texttt{\mdseries\slshape d}} \times \mbox{\texttt{\mdseries\slshape d}}$ matrices over the field with \mbox{\texttt{\mdseries\slshape q}} elements that respect a non\texttt{\symbol{45}}singular quadratic form
(see{\nobreakspace}\texttt{InvariantQuadraticForm} (\ref{InvariantQuadraticForm})) specified by \mbox{\texttt{\mdseries\slshape e}}, in the category given by the filter \mbox{\texttt{\mdseries\slshape filt}}. 

 The value of \mbox{\texttt{\mdseries\slshape e}} must be $0$ for odd \mbox{\texttt{\mdseries\slshape d}} (and can optionally be omitted in this case), respectively one of $1$ or $-1$ for even \mbox{\texttt{\mdseries\slshape d}}. If \mbox{\texttt{\mdseries\slshape filt}} is not given it defaults to \texttt{IsMatrixGroup} (\ref{IsMatrixGroup}), and the returned group is the general orthogonal group itself. 

 If version at least 1.2.6 of the \textsf{Forms} package is loaded then the desired quadratic form can be specified via \mbox{\texttt{\mdseries\slshape form}}, which can be either a matrix or a form object in \texttt{IsQuadraticForm} (\textbf{Forms: IsQuadraticForm}) or a group with stored \texttt{InvariantQuadraticForm} (\ref{InvariantQuadraticForm}) value (and then this form is taken). 

 A given \mbox{\texttt{\mdseries\slshape form}} determines \mbox{\texttt{\mdseries\slshape e}} and \mbox{\texttt{\mdseries\slshape d}}, and also \mbox{\texttt{\mdseries\slshape q}} except if \mbox{\texttt{\mdseries\slshape form}} is a matrix that does not store its \texttt{BaseDomain} (\ref{BaseDomain:for a matrix object}) value. These parameters can be entered, and an error is signalled if they do
not fit to the given \mbox{\texttt{\mdseries\slshape form}}. 

 If \mbox{\texttt{\mdseries\slshape form}} is not given then a default is chosen as described in the introduction to
Section \ref{Classical Groups}. 

 Note that in{\nobreakspace}\cite{KleidmanLiebeck90}, GO is defined as the stabilizer $\Delta(V, F, \kappa)$ of the quadratic form, up to scalars, whereas our GO is called $I(V, F, \kappa)$ there. 

 
\begin{Verbatim}[commandchars=!@|,fontsize=\small,frame=single,label=Example]
  !gapprompt@gap>| !gapinput@GeneralOrthogonalGroup( 5, 3 );|
  GO(0,5,3)
  !gapprompt@gap>| !gapinput@GeneralOrthogonalGroup( -1, 8, 2 );|
  GO(-1,8,2)
  !gapprompt@gap>| !gapinput@GeneralOrthogonalGroup( IsPermGroup, -1, 8, 2 );|
  Perm_GO(-1,8,2)
\end{Verbatim}
 }

 

\subsection{\textcolor{Chapter }{SpecialOrthogonalGroup}}
\logpage{[ 50, 2, 7 ]}\nobreak
\hyperdef{L}{X78D4EEF27AA2DCFD}{}
{\noindent\textcolor{FuncColor}{$\triangleright$\enspace\texttt{SpecialOrthogonalGroup({\mdseries\slshape [filt, ][e, ]d, q[, form]})\index{SpecialOrthogonalGroup@\texttt{SpecialOrthogonalGroup}}
\label{SpecialOrthogonalGroup}
}\hfill{\scriptsize (function)}}\\
\noindent\textcolor{FuncColor}{$\triangleright$\enspace\texttt{SpecialOrthogonalGroup({\mdseries\slshape [filt, ]form})\index{SpecialOrthogonalGroup@\texttt{SpecialOrthogonalGroup}!for a form}
\label{SpecialOrthogonalGroup:for a form}
}\hfill{\scriptsize (function)}}\\
\noindent\textcolor{FuncColor}{$\triangleright$\enspace\texttt{SO({\mdseries\slshape [filt, ][e, ]d, q[, form]})\index{SO@\texttt{SO}}
\label{SO}
}\hfill{\scriptsize (function)}}\\
\noindent\textcolor{FuncColor}{$\triangleright$\enspace\texttt{SO({\mdseries\slshape [filt, ]form})\index{SO@\texttt{SO}!for a form}
\label{SO:for a form}
}\hfill{\scriptsize (function)}}\\


 constructs a group isomorphic to the special orthogonal group SO( \mbox{\texttt{\mdseries\slshape e}}, \mbox{\texttt{\mdseries\slshape d}}, \mbox{\texttt{\mdseries\slshape q}} ), which is the subgroup of all those matrices in the general orthogonal group
(see{\nobreakspace}\texttt{GeneralOrthogonalGroup} (\ref{GeneralOrthogonalGroup})) that have determinant one, in the category given by the filter \mbox{\texttt{\mdseries\slshape filt}}. (The index of SO( \mbox{\texttt{\mdseries\slshape e}}, \mbox{\texttt{\mdseries\slshape d}}, \mbox{\texttt{\mdseries\slshape q}} ) in GO( \mbox{\texttt{\mdseries\slshape e}}, \mbox{\texttt{\mdseries\slshape d}}, \mbox{\texttt{\mdseries\slshape q}} ) is $2$ if \mbox{\texttt{\mdseries\slshape q}} is odd, and $1$ if \mbox{\texttt{\mdseries\slshape q}} is even.) Also interesting is the group Omega( \mbox{\texttt{\mdseries\slshape e}}, \mbox{\texttt{\mdseries\slshape d}}, \mbox{\texttt{\mdseries\slshape q}} ), see \texttt{Omega} (\ref{Omega:construct an orthogonal group}), which is of index $2$ in SO( \mbox{\texttt{\mdseries\slshape e}}, \mbox{\texttt{\mdseries\slshape d}}, \mbox{\texttt{\mdseries\slshape q}} ), except in the case $\mbox{\texttt{\mdseries\slshape d}} = 1$. 

 If \mbox{\texttt{\mdseries\slshape filt}} is not given it defaults to \texttt{IsMatrixGroup} (\ref{IsMatrixGroup}), and the returned group is the special orthogonal group itself. 

 If version at least 1.2.6 of the \textsf{Forms} package is loaded then the desired quadratic form can be specified via \mbox{\texttt{\mdseries\slshape form}}, which can be either a matrix or a form object in \texttt{IsQuadraticForm} (\textbf{Forms: IsQuadraticForm}) or a group with stored \texttt{InvariantQuadraticForm} (\ref{InvariantQuadraticForm}) value (and then this form is taken). 

 A given \mbox{\texttt{\mdseries\slshape form}} determines \mbox{\texttt{\mdseries\slshape e}} and \mbox{\texttt{\mdseries\slshape d}}, and also \mbox{\texttt{\mdseries\slshape q}} except if \mbox{\texttt{\mdseries\slshape form}} is a matrix that does not store its \texttt{BaseDomain} (\ref{BaseDomain:for a matrix object}) value. These parameters can be entered, and an error is signalled if they do
not fit to the given \mbox{\texttt{\mdseries\slshape form}}. 

 If \mbox{\texttt{\mdseries\slshape form}} is not given then a default is chosen as described in the introduction to
Section \ref{Classical Groups}. 

 
\begin{Verbatim}[commandchars=!@|,fontsize=\small,frame=single,label=Example]
  !gapprompt@gap>| !gapinput@SpecialOrthogonalGroup( 5, 3 );|
  SO(0,5,3)
  !gapprompt@gap>| !gapinput@SpecialOrthogonalGroup( -1, 8, 2 );  # here SO and GO coincide|
  GO(-1,8,2)
  !gapprompt@gap>| !gapinput@SpecialOrthogonalGroup( IsPermGroup, 5, 3 );|
  Perm_SO(0,5,3)
\end{Verbatim}
 }

 

\subsection{\textcolor{Chapter }{Omega (construct an orthogonal group)}}
\logpage{[ 50, 2, 8 ]}\nobreak
\hyperdef{L}{X8365E0AB8338DA3F}{}
{\noindent\textcolor{FuncColor}{$\triangleright$\enspace\texttt{Omega({\mdseries\slshape [filt, ][e, ]d, q[, form]})\index{Omega@\texttt{Omega}!construct an orthogonal group}
\label{Omega:construct an orthogonal group}
}\hfill{\scriptsize (operation)}}\\
\noindent\textcolor{FuncColor}{$\triangleright$\enspace\texttt{Omega({\mdseries\slshape [filt, ]form})\index{Omega@\texttt{Omega}!construct an orthogonal group for a given quadratic form}
\label{Omega:construct an orthogonal group for a given quadratic form}
}\hfill{\scriptsize (operation)}}\\


 constructs a group isomorphic to the group $\Omega$( \mbox{\texttt{\mdseries\slshape e}}, \mbox{\texttt{\mdseries\slshape d}}, \mbox{\texttt{\mdseries\slshape q}} ) of those $\mbox{\texttt{\mdseries\slshape d}} \times \mbox{\texttt{\mdseries\slshape d}}$ matrices over the field with \mbox{\texttt{\mdseries\slshape q}} elements that respect a non\texttt{\symbol{45}}singular quadratic form
(see{\nobreakspace}\texttt{InvariantQuadraticForm} (\ref{InvariantQuadraticForm})) specified by \mbox{\texttt{\mdseries\slshape e}}, and that have square spinor norm in odd characteristic or Dickson invariant $0$ in even characteristic, respectively, in the category given by the filter \mbox{\texttt{\mdseries\slshape filt}}. 

 For odd \mbox{\texttt{\mdseries\slshape q}} and $\mbox{\texttt{\mdseries\slshape d}} \geq 2$, this group has always index two in the corresponding special orthogonal
group, which will be conjugate in $GL(d,q)$ to the group returned by SO( \mbox{\texttt{\mdseries\slshape e}}, \mbox{\texttt{\mdseries\slshape d}}, \mbox{\texttt{\mdseries\slshape q}} ), see \texttt{SpecialOrthogonalGroup} (\ref{SpecialOrthogonalGroup}), but may fix a different form (see \ref{Classical Groups}). 

 The value of \mbox{\texttt{\mdseries\slshape e}} must be $0$ for odd \mbox{\texttt{\mdseries\slshape d}} (and can optionally be omitted in this case), respectively one of $1$ or $-1$ for even \mbox{\texttt{\mdseries\slshape d}}. If \mbox{\texttt{\mdseries\slshape filt}} is not given it defaults to \texttt{IsMatrixGroup} (\ref{IsMatrixGroup}), and the returned group is the group $\Omega$( \mbox{\texttt{\mdseries\slshape e}}, \mbox{\texttt{\mdseries\slshape d}}, \mbox{\texttt{\mdseries\slshape q}} ) itself. 

 If version at least 1.2.6 of the \textsf{Forms} package is loaded then the desired quadratic form can be specified via \mbox{\texttt{\mdseries\slshape form}}, which can be either a matrix or a form object in \texttt{IsQuadraticForm} (\textbf{Forms: IsQuadraticForm}) or a group with stored \texttt{InvariantQuadraticForm} (\ref{InvariantQuadraticForm}) value (and then this form is taken). 

 A given \mbox{\texttt{\mdseries\slshape form}} determines \mbox{\texttt{\mdseries\slshape e}} and \mbox{\texttt{\mdseries\slshape d}}, and also \mbox{\texttt{\mdseries\slshape q}} except if \mbox{\texttt{\mdseries\slshape form}} is a matrix that does not store its \texttt{BaseDomain} (\ref{BaseDomain:for a matrix object}) value. These parameters can be entered, and an error is signalled if they do
not fit to the given \mbox{\texttt{\mdseries\slshape form}}. 

 If \mbox{\texttt{\mdseries\slshape form}} is not given then a default is chosen as described in the introduction to
Section \ref{Classical Groups}. 

 
\begin{Verbatim}[commandchars=@|B,fontsize=\small,frame=single,label=Example]
  @gapprompt|gap>B @gapinput|g:= Omega( 3, 5 );  StructureDescription( g );B
  Omega(0,3,5)
  "A5"
  @gapprompt|gap>B @gapinput|g:= Omega( 1, 4, 4 );  StructureDescription( g );B
  Omega(+1,4,4)
  "A5 x A5"
  @gapprompt|gap>B @gapinput|g:= Omega( -1, 4, 3 );  StructureDescription( g );B
  Omega(-1,4,3)
  "A6"
  @gapprompt|gap>B @gapinput|g:= Omega( IsPermGroup, 1, 6, 2 );  StructureDescription( g );B
  Perm_Omega(+1,6,2)
  "A8"
  @gapprompt|gap>B @gapinput|IsSubset( GO( 3, 5 ), Omega( 3, 5 ) );  # different forms!B
  false
\end{Verbatim}
 }

 

\subsection{\textcolor{Chapter }{GeneralSemilinearGroup}}
\logpage{[ 50, 2, 9 ]}\nobreak
\hyperdef{L}{X79C3C61A7D83A6D0}{}
{\noindent\textcolor{FuncColor}{$\triangleright$\enspace\texttt{GeneralSemilinearGroup({\mdseries\slshape [filt, ]d, q})\index{GeneralSemilinearGroup@\texttt{GeneralSemilinearGroup}}
\label{GeneralSemilinearGroup}
}\hfill{\scriptsize (function)}}\\
\noindent\textcolor{FuncColor}{$\triangleright$\enspace\texttt{GammaL({\mdseries\slshape [filt, ]d, q})\index{GammaL@\texttt{GammaL}}
\label{GammaL}
}\hfill{\scriptsize (function)}}\\


 \texttt{GeneralSemilinearGroup} returns a group isomorphic to the general semilinear group $\Gamma$L( \mbox{\texttt{\mdseries\slshape d}}, \mbox{\texttt{\mdseries\slshape q}} ) of semilinear mappings of the vector space \texttt{GF( }\mbox{\texttt{\mdseries\slshape q}}\texttt{ )\texttt{\symbol{94}}}\mbox{\texttt{\mdseries\slshape d}}. 

 If \mbox{\texttt{\mdseries\slshape filt}} is not given it defaults to \texttt{IsMatrixGroup} (\ref{IsMatrixGroup}), and the returned group consists of matrices of dimension \mbox{\texttt{\mdseries\slshape d}} $f$ over the field with $p$ elements, where \mbox{\texttt{\mdseries\slshape q}} $= p^f$, for a prime integer $p$. }

 

\subsection{\textcolor{Chapter }{SpecialSemilinearGroup}}
\logpage{[ 50, 2, 10 ]}\nobreak
\hyperdef{L}{X7D3779237CB5B49C}{}
{\noindent\textcolor{FuncColor}{$\triangleright$\enspace\texttt{SpecialSemilinearGroup({\mdseries\slshape [filt, ]d, q})\index{SpecialSemilinearGroup@\texttt{SpecialSemilinearGroup}}
\label{SpecialSemilinearGroup}
}\hfill{\scriptsize (function)}}\\
\noindent\textcolor{FuncColor}{$\triangleright$\enspace\texttt{SigmaL({\mdseries\slshape [filt, ]d, q})\index{SigmaL@\texttt{SigmaL}}
\label{SigmaL}
}\hfill{\scriptsize (function)}}\\


 \texttt{SpecialSemilinearGroup} returns a group isomorphic to the special semilinear group $\Sigma$L( \mbox{\texttt{\mdseries\slshape d}}, \mbox{\texttt{\mdseries\slshape q}} ) of those semilinear mappings of the vector space \texttt{GF( }\mbox{\texttt{\mdseries\slshape q}}\texttt{ )\texttt{\symbol{94}}}\mbox{\texttt{\mdseries\slshape d}} (see \texttt{GeneralSemilinearGroup} (\ref{GeneralSemilinearGroup})) whose linear part has determinant one. 

 If \mbox{\texttt{\mdseries\slshape filt}} is not given it defaults to \texttt{IsMatrixGroup} (\ref{IsMatrixGroup}), and the returned group consists of matrices of dimension \mbox{\texttt{\mdseries\slshape d}} $f$ over the field with $p$ elements, where \mbox{\texttt{\mdseries\slshape q}} $= p^f$, for a prime integer $p$. }

 

\subsection{\textcolor{Chapter }{ProjectiveGeneralLinearGroup}}
\logpage{[ 50, 2, 11 ]}\nobreak
\hyperdef{L}{X7F0DBEB880D2D574}{}
{\noindent\textcolor{FuncColor}{$\triangleright$\enspace\texttt{ProjectiveGeneralLinearGroup({\mdseries\slshape [filt, ]d, q})\index{ProjectiveGeneralLinearGroup@\texttt{ProjectiveGeneralLinearGroup}}
\label{ProjectiveGeneralLinearGroup}
}\hfill{\scriptsize (function)}}\\
\noindent\textcolor{FuncColor}{$\triangleright$\enspace\texttt{PGL({\mdseries\slshape [filt, ]d, q})\index{PGL@\texttt{PGL}}
\label{PGL}
}\hfill{\scriptsize (function)}}\\


 constructs a group isomorphic to the projective general linear group PGL( \mbox{\texttt{\mdseries\slshape d}}, \mbox{\texttt{\mdseries\slshape q}} ) of those $\mbox{\texttt{\mdseries\slshape d}} \times \mbox{\texttt{\mdseries\slshape d}}$ matrices over the field with \mbox{\texttt{\mdseries\slshape q}} elements, modulo the centre, in the category given by the filter \mbox{\texttt{\mdseries\slshape filt}}. 

 If \mbox{\texttt{\mdseries\slshape filt}} is not given it defaults to \texttt{IsPermGroup} (\ref{IsPermGroup}), and the returned group is the action on lines of the underlying vector
space. 

 }

 

\subsection{\textcolor{Chapter }{ProjectiveSpecialLinearGroup}}
\logpage{[ 50, 2, 12 ]}\nobreak
\hyperdef{L}{X86784EDA80224B74}{}
{\noindent\textcolor{FuncColor}{$\triangleright$\enspace\texttt{ProjectiveSpecialLinearGroup({\mdseries\slshape [filt, ]d, q})\index{ProjectiveSpecialLinearGroup@\texttt{ProjectiveSpecialLinearGroup}}
\label{ProjectiveSpecialLinearGroup}
}\hfill{\scriptsize (function)}}\\
\noindent\textcolor{FuncColor}{$\triangleright$\enspace\texttt{PSL({\mdseries\slshape [filt, ]d, q})\index{PSL@\texttt{PSL}}
\label{PSL}
}\hfill{\scriptsize (function)}}\\


 constructs a group isomorphic to the projective special linear group PSL( \mbox{\texttt{\mdseries\slshape d}}, \mbox{\texttt{\mdseries\slshape q}} ) of those $\mbox{\texttt{\mdseries\slshape d}} \times \mbox{\texttt{\mdseries\slshape d}}$ matrices over the field with \mbox{\texttt{\mdseries\slshape q}} elements whose determinant is the identity of the field, modulo the centre, in
the category given by the filter \mbox{\texttt{\mdseries\slshape filt}}. 

 If \mbox{\texttt{\mdseries\slshape filt}} is not given it defaults to \texttt{IsPermGroup} (\ref{IsPermGroup}), and the returned group is the action on lines of the underlying vector
space. }

 

\subsection{\textcolor{Chapter }{ProjectiveGeneralUnitaryGroup}}
\logpage{[ 50, 2, 13 ]}\nobreak
\hyperdef{L}{X7E471ADE7E095604}{}
{\noindent\textcolor{FuncColor}{$\triangleright$\enspace\texttt{ProjectiveGeneralUnitaryGroup({\mdseries\slshape [filt, ]d, q})\index{ProjectiveGeneralUnitaryGroup@\texttt{ProjectiveGeneralUnitaryGroup}}
\label{ProjectiveGeneralUnitaryGroup}
}\hfill{\scriptsize (function)}}\\
\noindent\textcolor{FuncColor}{$\triangleright$\enspace\texttt{PGU({\mdseries\slshape [filt, ]d, q})\index{PGU@\texttt{PGU}}
\label{PGU}
}\hfill{\scriptsize (function)}}\\


 constructs a group isomorphic to the projective general unitary group PGU( \mbox{\texttt{\mdseries\slshape d}}, \mbox{\texttt{\mdseries\slshape q}} ) of those $\mbox{\texttt{\mdseries\slshape d}} \times \mbox{\texttt{\mdseries\slshape d}}$ matrices over the field with $\mbox{\texttt{\mdseries\slshape q}}^2$ elements that respect a fixed nondegenerate sesquilinear form, modulo the
centre, in the category given by the filter \mbox{\texttt{\mdseries\slshape filt}}. 

 If \mbox{\texttt{\mdseries\slshape filt}} is not given it defaults to \texttt{IsPermGroup} (\ref{IsPermGroup}), and the returned group is the action on lines of the underlying vector
space. }

 

\subsection{\textcolor{Chapter }{ProjectiveSpecialUnitaryGroup}}
\logpage{[ 50, 2, 14 ]}\nobreak
\hyperdef{L}{X7A88FE2B7EF9C804}{}
{\noindent\textcolor{FuncColor}{$\triangleright$\enspace\texttt{ProjectiveSpecialUnitaryGroup({\mdseries\slshape [filt, ]d, q})\index{ProjectiveSpecialUnitaryGroup@\texttt{ProjectiveSpecialUnitaryGroup}}
\label{ProjectiveSpecialUnitaryGroup}
}\hfill{\scriptsize (function)}}\\
\noindent\textcolor{FuncColor}{$\triangleright$\enspace\texttt{PSU({\mdseries\slshape [filt, ]d, q})\index{PSU@\texttt{PSU}}
\label{PSU}
}\hfill{\scriptsize (function)}}\\


 constructs a group isomorphic to the projective special unitary group PSU( \mbox{\texttt{\mdseries\slshape d}}, \mbox{\texttt{\mdseries\slshape q}} ) of those $\mbox{\texttt{\mdseries\slshape d}} \times \mbox{\texttt{\mdseries\slshape d}}$ matrices over the field with $\mbox{\texttt{\mdseries\slshape q}}^2$ elements that respect a fixed nondegenerate sesquilinear form and have
determinant 1, modulo the centre, in the category given by the filter \mbox{\texttt{\mdseries\slshape filt}}. 

 If \mbox{\texttt{\mdseries\slshape filt}} is not given it defaults to \texttt{IsPermGroup} (\ref{IsPermGroup}), and the returned group is the action on lines of the underlying vector
space. }

 

\subsection{\textcolor{Chapter }{ProjectiveSymplecticGroup}}
\logpage{[ 50, 2, 15 ]}\nobreak
\hyperdef{L}{X7DEDE2537B8FFFF5}{}
{\noindent\textcolor{FuncColor}{$\triangleright$\enspace\texttt{ProjectiveSymplecticGroup({\mdseries\slshape [filt, ]d, q})\index{ProjectiveSymplecticGroup@\texttt{ProjectiveSymplecticGroup}}
\label{ProjectiveSymplecticGroup}
}\hfill{\scriptsize (function)}}\\
\noindent\textcolor{FuncColor}{$\triangleright$\enspace\texttt{PSP({\mdseries\slshape [filt, ]d, q})\index{PSP@\texttt{PSP}}
\label{PSP}
}\hfill{\scriptsize (function)}}\\
\noindent\textcolor{FuncColor}{$\triangleright$\enspace\texttt{PSp({\mdseries\slshape [filt, ]d, q})\index{PSp@\texttt{PSp}}
\label{PSp}
}\hfill{\scriptsize (function)}}\\


 constructs a group isomorphic to the projective symplectic group PSp(\mbox{\texttt{\mdseries\slshape d}},\mbox{\texttt{\mdseries\slshape q}}) of those $\mbox{\texttt{\mdseries\slshape d}} \times \mbox{\texttt{\mdseries\slshape d}}$ matrices over the field with \mbox{\texttt{\mdseries\slshape q}} elements that respect a fixed nondegenerate symplectic form, modulo the
centre, in the category given by the filter \mbox{\texttt{\mdseries\slshape filt}}. 

 If \mbox{\texttt{\mdseries\slshape filt}} is not given it defaults to \texttt{IsPermGroup} (\ref{IsPermGroup}), and the returned group is the action on lines of the underlying vector
space. }

 

\subsection{\textcolor{Chapter }{ProjectiveGeneralOrthogonalGroup}}
\logpage{[ 50, 2, 16 ]}\nobreak
\hyperdef{L}{X87AAB20A8434356B}{}
{\noindent\textcolor{FuncColor}{$\triangleright$\enspace\texttt{ProjectiveGeneralOrthogonalGroup({\mdseries\slshape [filt, ][e, ]d, q})\index{ProjectiveGeneralOrthogonalGroup@\texttt{ProjectiveGeneralOrthogonalGroup}}
\label{ProjectiveGeneralOrthogonalGroup}
}\hfill{\scriptsize (function)}}\\
\noindent\textcolor{FuncColor}{$\triangleright$\enspace\texttt{PGO({\mdseries\slshape [filt, ][e, ]d, q})\index{PGO@\texttt{PGO}}
\label{PGO}
}\hfill{\scriptsize (function)}}\\


 constructs a group isomorphic to the projective group PGO( \mbox{\texttt{\mdseries\slshape e}}, \mbox{\texttt{\mdseries\slshape d}}, \mbox{\texttt{\mdseries\slshape q}} ) of GO( \mbox{\texttt{\mdseries\slshape e}}, \mbox{\texttt{\mdseries\slshape d}}, \mbox{\texttt{\mdseries\slshape q}} ), modulo the centre (see \texttt{GeneralOrthogonalGroup} (\ref{GeneralOrthogonalGroup})), in the category given by the filter \mbox{\texttt{\mdseries\slshape filt}}. 

 If \mbox{\texttt{\mdseries\slshape filt}} is not given it defaults to \texttt{IsPermGroup} (\ref{IsPermGroup}), and the returned group is the action on lines of the underlying vector
space. }

 

\subsection{\textcolor{Chapter }{ProjectiveSpecialOrthogonalGroup}}
\logpage{[ 50, 2, 17 ]}\nobreak
\hyperdef{L}{X835E0D3384C4AB6B}{}
{\noindent\textcolor{FuncColor}{$\triangleright$\enspace\texttt{ProjectiveSpecialOrthogonalGroup({\mdseries\slshape [filt, ][e, ]d, q})\index{ProjectiveSpecialOrthogonalGroup@\texttt{ProjectiveSpecialOrthogonalGroup}}
\label{ProjectiveSpecialOrthogonalGroup}
}\hfill{\scriptsize (function)}}\\
\noindent\textcolor{FuncColor}{$\triangleright$\enspace\texttt{PSO({\mdseries\slshape [filt, ][e, ]d, q})\index{PSO@\texttt{PSO}}
\label{PSO}
}\hfill{\scriptsize (function)}}\\


 constructs a group isomorphic to the projective group PSO( \mbox{\texttt{\mdseries\slshape e}}, \mbox{\texttt{\mdseries\slshape d}}, \mbox{\texttt{\mdseries\slshape q}} ) of SO( \mbox{\texttt{\mdseries\slshape e}}, \mbox{\texttt{\mdseries\slshape d}}, \mbox{\texttt{\mdseries\slshape q}} ), modulo the centre (see \texttt{SpecialOrthogonalGroup} (\ref{SpecialOrthogonalGroup})), in the category given by the filter \mbox{\texttt{\mdseries\slshape filt}}. 

 If \mbox{\texttt{\mdseries\slshape filt}} is not given it defaults to \texttt{IsPermGroup} (\ref{IsPermGroup}), and the returned group is the action on lines of the underlying vector
space. }

 

\subsection{\textcolor{Chapter }{ProjectiveOmega}}
\logpage{[ 50, 2, 18 ]}\nobreak
\hyperdef{L}{X7F546F907A37DF15}{}
{\noindent\textcolor{FuncColor}{$\triangleright$\enspace\texttt{ProjectiveOmega({\mdseries\slshape [filt, ][e, ]d, q})\index{ProjectiveOmega@\texttt{ProjectiveOmega}}
\label{ProjectiveOmega}
}\hfill{\scriptsize (function)}}\\
\noindent\textcolor{FuncColor}{$\triangleright$\enspace\texttt{POmega({\mdseries\slshape [filt, ][e, ]d, q})\index{POmega@\texttt{POmega}}
\label{POmega}
}\hfill{\scriptsize (function)}}\\


 constructs a group isomorphic to the projective group P$\Omega$( \mbox{\texttt{\mdseries\slshape e}}, \mbox{\texttt{\mdseries\slshape d}}, \mbox{\texttt{\mdseries\slshape q}} ) of $\Omega$( \mbox{\texttt{\mdseries\slshape e}}, \mbox{\texttt{\mdseries\slshape d}}, \mbox{\texttt{\mdseries\slshape q}} ), modulo the centre (see \texttt{Omega} (\ref{Omega:construct an orthogonal group})), in the category given by the filter \mbox{\texttt{\mdseries\slshape filt}}. 

 If \mbox{\texttt{\mdseries\slshape filt}} is not given it defaults to \texttt{IsPermGroup} (\ref{IsPermGroup}), and the returned group is the action on lines of the underlying vector
space. }

 

\subsection{\textcolor{Chapter }{ProjectiveGeneralSemilinearGroup}}
\logpage{[ 50, 2, 19 ]}\nobreak
\hyperdef{L}{X824925DB7C3A2FA6}{}
{\noindent\textcolor{FuncColor}{$\triangleright$\enspace\texttt{ProjectiveGeneralSemilinearGroup({\mdseries\slshape [filt, ]d, q})\index{ProjectiveGeneralSemilinearGroup@\texttt{ProjectiveGeneralSemilinearGroup}}
\label{ProjectiveGeneralSemilinearGroup}
}\hfill{\scriptsize (function)}}\\
\noindent\textcolor{FuncColor}{$\triangleright$\enspace\texttt{PGammaL({\mdseries\slshape [filt, ]d, q})\index{PGammaL@\texttt{PGammaL}}
\label{PGammaL}
}\hfill{\scriptsize (function)}}\\


 \texttt{ProjectiveGeneralSemilinearGroup} returns a group isomorphic to the factor group of the general semilinear group \texttt{GammaL(} \mbox{\texttt{\mdseries\slshape d}}, \mbox{\texttt{\mdseries\slshape q}} \texttt{)} modulo the center of its normal subgroup \texttt{GL(} \mbox{\texttt{\mdseries\slshape d}}, \mbox{\texttt{\mdseries\slshape q}} \texttt{)}. 

 If \mbox{\texttt{\mdseries\slshape filt}} is not given it defaults to \texttt{IsPermGroup} (\ref{IsPermGroup}), and the returned group is the action on lines of the underlying vector space \texttt{GF(}\mbox{\texttt{\mdseries\slshape q}}\texttt{)\texttt{\symbol{94}}}\mbox{\texttt{\mdseries\slshape d}}. }

 

\subsection{\textcolor{Chapter }{ProjectiveSpecialSemilinearGroup}}
\logpage{[ 50, 2, 20 ]}\nobreak
\hyperdef{L}{X86BD9AE27CCAB1A6}{}
{\noindent\textcolor{FuncColor}{$\triangleright$\enspace\texttt{ProjectiveSpecialSemilinearGroup({\mdseries\slshape [filt, ]d, q})\index{ProjectiveSpecialSemilinearGroup@\texttt{ProjectiveSpecialSemilinearGroup}}
\label{ProjectiveSpecialSemilinearGroup}
}\hfill{\scriptsize (function)}}\\
\noindent\textcolor{FuncColor}{$\triangleright$\enspace\texttt{PSigmaL({\mdseries\slshape [filt, ]d, q})\index{PSigmaL@\texttt{PSigmaL}}
\label{PSigmaL}
}\hfill{\scriptsize (function)}}\\


 \texttt{ProjectiveSpecialSemilinearGroup} returns a group isomorphic to the factor group of the special semilinear group \texttt{SigmaL(} \mbox{\texttt{\mdseries\slshape d}}, \mbox{\texttt{\mdseries\slshape q}} \texttt{)} modulo the center of its normal subgroup \texttt{SL(} \mbox{\texttt{\mdseries\slshape d}}, \mbox{\texttt{\mdseries\slshape q}} \texttt{)}. 

 If \mbox{\texttt{\mdseries\slshape filt}} is not given it defaults to \texttt{IsPermGroup} (\ref{IsPermGroup}), and the returned group is the action on lines of the underlying vector space \texttt{GF(}\mbox{\texttt{\mdseries\slshape q}}\texttt{)\texttt{\symbol{94}}}\mbox{\texttt{\mdseries\slshape d}}. }

 }

  
\section{\textcolor{Chapter }{Conjugacy Classes in Classical Groups}}\label{Conjugacy Classes in Classical Groups}
\logpage{[ 50, 3, 0 ]}
\hyperdef{L}{X85B9F2D379616C35}{}
{
  \index{ConjugacyClasses@\texttt{ConjugacyClasses}!for linear groups} For general and special linear groups (see{\nobreakspace}\texttt{GeneralLinearGroup} (\ref{GeneralLinearGroup:for dimension and a ring}) and{\nobreakspace}\texttt{SpecialLinearGroup} (\ref{SpecialLinearGroup:for dimension and a ring})) \textsf{GAP} has an efficient method to generate representatives of the conjugacy classes.
This uses results from linear algebra on normal forms of matrices. If you know
how to do this for other types of classical groups, please, tell us. 

 
\begin{Verbatim}[commandchars=!@|,fontsize=\small,frame=single,label=Example]
  !gapprompt@gap>| !gapinput@g := SL(4,9);|
  SL(4,9)
  !gapprompt@gap>| !gapinput@NrConjugacyClasses(g);|
  861
  !gapprompt@gap>| !gapinput@cl := ConjugacyClasses(g);;|
  !gapprompt@gap>| !gapinput@Length(cl);|
  861
\end{Verbatim}
 

 

\subsection{\textcolor{Chapter }{NrConjugacyClassesGL}}
\logpage{[ 50, 3, 1 ]}\nobreak
\hyperdef{L}{X831789117E93171E}{}
{\noindent\textcolor{FuncColor}{$\triangleright$\enspace\texttt{NrConjugacyClassesGL({\mdseries\slshape n, q})\index{NrConjugacyClassesGL@\texttt{NrConjugacyClassesGL}}
\label{NrConjugacyClassesGL}
}\hfill{\scriptsize (function)}}\\
\noindent\textcolor{FuncColor}{$\triangleright$\enspace\texttt{NrConjugacyClassesGU({\mdseries\slshape n, q})\index{NrConjugacyClassesGU@\texttt{NrConjugacyClassesGU}}
\label{NrConjugacyClassesGU}
}\hfill{\scriptsize (function)}}\\
\noindent\textcolor{FuncColor}{$\triangleright$\enspace\texttt{NrConjugacyClassesSL({\mdseries\slshape n, q})\index{NrConjugacyClassesSL@\texttt{NrConjugacyClassesSL}}
\label{NrConjugacyClassesSL}
}\hfill{\scriptsize (function)}}\\
\noindent\textcolor{FuncColor}{$\triangleright$\enspace\texttt{NrConjugacyClassesSU({\mdseries\slshape n, q})\index{NrConjugacyClassesSU@\texttt{NrConjugacyClassesSU}}
\label{NrConjugacyClassesSU}
}\hfill{\scriptsize (function)}}\\
\noindent\textcolor{FuncColor}{$\triangleright$\enspace\texttt{NrConjugacyClassesPGL({\mdseries\slshape n, q})\index{NrConjugacyClassesPGL@\texttt{NrConjugacyClassesPGL}}
\label{NrConjugacyClassesPGL}
}\hfill{\scriptsize (function)}}\\
\noindent\textcolor{FuncColor}{$\triangleright$\enspace\texttt{NrConjugacyClassesPGU({\mdseries\slshape n, q})\index{NrConjugacyClassesPGU@\texttt{NrConjugacyClassesPGU}}
\label{NrConjugacyClassesPGU}
}\hfill{\scriptsize (function)}}\\
\noindent\textcolor{FuncColor}{$\triangleright$\enspace\texttt{NrConjugacyClassesPSL({\mdseries\slshape n, q})\index{NrConjugacyClassesPSL@\texttt{NrConjugacyClassesPSL}}
\label{NrConjugacyClassesPSL}
}\hfill{\scriptsize (function)}}\\
\noindent\textcolor{FuncColor}{$\triangleright$\enspace\texttt{NrConjugacyClassesPSU({\mdseries\slshape n, q})\index{NrConjugacyClassesPSU@\texttt{NrConjugacyClassesPSU}}
\label{NrConjugacyClassesPSU}
}\hfill{\scriptsize (function)}}\\
\noindent\textcolor{FuncColor}{$\triangleright$\enspace\texttt{NrConjugacyClassesSLIsogeneous({\mdseries\slshape n, q, f})\index{NrConjugacyClassesSLIsogeneous@\texttt{NrConjugacyClassesSLIsogeneous}}
\label{NrConjugacyClassesSLIsogeneous}
}\hfill{\scriptsize (function)}}\\
\noindent\textcolor{FuncColor}{$\triangleright$\enspace\texttt{NrConjugacyClassesSUIsogeneous({\mdseries\slshape n, q, f})\index{NrConjugacyClassesSUIsogeneous@\texttt{NrConjugacyClassesSUIsogeneous}}
\label{NrConjugacyClassesSUIsogeneous}
}\hfill{\scriptsize (function)}}\\


 The first of these functions compute for given positive integer \mbox{\texttt{\mdseries\slshape n}} and prime power \mbox{\texttt{\mdseries\slshape q}} the number of conjugacy classes in the classical groups GL( \mbox{\texttt{\mdseries\slshape n}}, \mbox{\texttt{\mdseries\slshape q}} ), GU( \mbox{\texttt{\mdseries\slshape n}}, \mbox{\texttt{\mdseries\slshape q}} ), SL( \mbox{\texttt{\mdseries\slshape n}}, \mbox{\texttt{\mdseries\slshape q}} ), SU( \mbox{\texttt{\mdseries\slshape n}}, \mbox{\texttt{\mdseries\slshape q}} ), PGL( \mbox{\texttt{\mdseries\slshape n}}, \mbox{\texttt{\mdseries\slshape q}} ), PGU( \mbox{\texttt{\mdseries\slshape n}}, \mbox{\texttt{\mdseries\slshape q}} ), PSL( \mbox{\texttt{\mdseries\slshape n}}, \mbox{\texttt{\mdseries\slshape q}} ), PSL( \mbox{\texttt{\mdseries\slshape n}}, \mbox{\texttt{\mdseries\slshape q}} ), respectively. (See also \texttt{ConjugacyClasses} (\ref{ConjugacyClasses:attribute}) and Section{\nobreakspace}\ref{Classical Groups}.) 

 For each divisor \mbox{\texttt{\mdseries\slshape f}} of \mbox{\texttt{\mdseries\slshape n}} there is a group of Lie type with the same order as SL( \mbox{\texttt{\mdseries\slshape n}}, \mbox{\texttt{\mdseries\slshape q}} ), such that its derived subgroup modulo its center is isomorphic to PSL( \mbox{\texttt{\mdseries\slshape n}}, \mbox{\texttt{\mdseries\slshape q}} ). The various such groups with fixed \mbox{\texttt{\mdseries\slshape n}} and \mbox{\texttt{\mdseries\slshape q}} are called \emph{isogeneous}. (Depending on congruence conditions on \mbox{\texttt{\mdseries\slshape q}} and \mbox{\texttt{\mdseries\slshape n}} several of these groups may actually be isomorphic.) The function \texttt{NrConjugacyClassesSLIsogeneous} computes the number of conjugacy classes in this group. The extreme cases \mbox{\texttt{\mdseries\slshape f}} $= 1$ and \mbox{\texttt{\mdseries\slshape f}} $= n$ lead to the groups SL( \mbox{\texttt{\mdseries\slshape n}}, \mbox{\texttt{\mdseries\slshape q}} ) and PGL( \mbox{\texttt{\mdseries\slshape n}}, \mbox{\texttt{\mdseries\slshape q}} ), respectively. 

 The function \texttt{NrConjugacyClassesSUIsogeneous} is the analogous one for the corresponding unitary groups. 

 The formulae for the number of conjugacy classes are taken from{\nobreakspace}\cite{Mac81}. 

 
\begin{Verbatim}[commandchars=!@|,fontsize=\small,frame=single,label=Example]
  !gapprompt@gap>| !gapinput@NrConjugacyClassesGL(24,27);|
  22528399544939174406067288580609952
  !gapprompt@gap>| !gapinput@NrConjugacyClassesPSU(19,17);|
  15052300411163848367708
  !gapprompt@gap>| !gapinput@NrConjugacyClasses(SL(16,16));|
  1229782938228219920
\end{Verbatim}
 }

 }

  
\section{\textcolor{Chapter }{Constructors for Basic Groups}}\label{Constructors for Basic Groups}
\logpage{[ 50, 4, 0 ]}
\hyperdef{L}{X817EBD6E841285CD}{}
{
  All functions described in the previous sections call constructor operations
to do the work. The names of the constructors are obtained from the names of
the functions by appending \texttt{"Cons"}, so for example \texttt{CyclicGroup} (\ref{CyclicGroup}) calls the constructor 

 \texttt{CyclicGroupCons( \mbox{\texttt{\mdseries\slshape cat}}, \mbox{\texttt{\mdseries\slshape n}} )} 

 The first argument \mbox{\texttt{\mdseries\slshape cat}} for each method of this constructor must be the category for which the method
is installed. For example the method for constructing a cyclic permutation
group is installed as follows (see{\nobreakspace}\texttt{InstallMethod} (\ref{InstallMethod}) for the meaning of the arguments. 

 
\begin{Verbatim}[commandchars=!@|,fontsize=\small,frame=single,label=Example]
  InstallMethod( CyclicGroupCons,
      "regular perm group",
      true,
      [ IsPermGroup and IsRegularProp and IsFinite, IsInt and IsPosRat ], 0,
      function( filter, n )
  
      ...
  
      end );
\end{Verbatim}
 }

  
\section{\textcolor{Chapter }{Selection Functions}}\label{Selection Functions}
\logpage{[ 50, 5, 0 ]}
\hyperdef{L}{X82676ED5826E9E2E}{}
{
  \index{AllPrimitiveGroups@\texttt{AllPrimitiveGroups}} \index{AllTransitiveGroups@\texttt{AllTransitiveGroups}} \index{AllLibraryGroups@\texttt{All\mbox{\texttt{\mdseries\slshape Library}}Groups}} \texttt{All\mbox{\texttt{\mdseries\slshape Library}}Groups( \mbox{\texttt{\mdseries\slshape fun1}}, \mbox{\texttt{\mdseries\slshape val1}}, ... )} 

 For a number of the group libraries two \emph{selection functions} are provided. Each \texttt{All\mbox{\texttt{\mdseries\slshape Library}}Groups} selection function permits one to select \emph{all} groups from the library \mbox{\texttt{\mdseries\slshape Library}} that have a given set of properties.    Currently, the library selection functions provided, of this type, are \texttt{AllSmallGroups} (\textbf{smallgrp: AllSmallGroups}), \texttt{AllIrreducibleSolvableGroups} (\textbf{primgrp: AllIrreducibleSolvableGroups}), \texttt{AllTransitiveGroups} (\textbf{transgrp: AllTransitiveGroups}), and \texttt{AllPrimitiveGroups} (\textbf{primgrp: AllPrimitiveGroups}). Corresponding to each of these there is a \texttt{One\mbox{\texttt{\mdseries\slshape Library}}Group} function (see below) which returns at most one group. 

 These functions take an arbitrary number of pairs (but at least one pair) of
arguments. The first argument in such a pair is a function that can be applied
to the groups in the library, and the second argument is either a single value
that this function must return in order to have this group included in the
selection, or a list of such values. For the function \texttt{AllSmallGroups} (\textbf{smallgrp: AllSmallGroups}) the first such function must be \texttt{Size} (\ref{Size}), and, unlike the other library selection functions, it supports an
alternative syntax where \texttt{Size} (\ref{Size}) is omitted (see{\nobreakspace}\texttt{AllSmallGroups} (\textbf{smallgrp: AllSmallGroups})). Also, see \texttt{AllIrreducibleSolvableGroups} (\textbf{primgrp: AllIrreducibleSolvableGroups}), for details pertaining to this function. 

 For an example, let us consider the selection function for the library of
transitive groups (also see{\nobreakspace} (\textbf{transgrp: Transitive Permutation Groups})). The command 

 
\begin{Verbatim}[commandchars=!@|,fontsize=\small,frame=single,label=Example]
  !gapprompt@gap>| !gapinput@AllTransitiveGroups(NrMovedPoints,[10..15],|
  !gapprompt@>| !gapinput@                       Size,         [1..100],|
  !gapprompt@>| !gapinput@                       IsAbelian,    false    );|
\end{Verbatim}
 

 returns a list of all transitive groups with degree between 10 and 15 and size
less than 100 that are not abelian. 

 Thus \texttt{AllTransitiveGroups} behaves as if it was implemented by a function similar to the one defined
below, where \texttt{TransitiveGroupsList} is a list of all transitive groups. (Note that in the definition below we
assume for simplicity that \texttt{AllTransitiveGroups} accepts exactly 4 arguments. It is of course obvious how to change this
definition so that the function would accept a variable number of arguments.) 

 
\begin{Verbatim}[commandchars=!@|,fontsize=\small,frame=single,label=Example]
  AllTransitiveGroups := function( fun1, val1, fun2, val2 )
  local    groups, g, i;
    groups := [];
    for i  in [ 1 .. Length( TransitiveGroupsList ) ] do
      g := TransitiveGroupsList[i];
      if      fun1(g) = val1  or IsList(val1) and fun1(g) in val1
          and fun2(g) = val2  or IsList(val2) and fun2(g) in val2
       then
        Add( groups, g );
      fi;
    od;
    return groups;
  end;
\end{Verbatim}
 

 Note that the real selection functions are considerably more difficult, to
improve the efficiency. Most important, each recognizes a certain set of
properties which are precomputed for the library without having to compute
them anew for each group. This will substantially speed up the selection
process. In the description of each library we will list the properties that
are stored for this library. 

 \index{OnePrimitiveGroup@\texttt{OnePrimitiveGroup}} \index{OneTransitiveGroup@\texttt{OneTransitiveGroup}} \index{OneLibraryGroup@\texttt{One\mbox{\texttt{\mdseries\slshape Library}}Group}} \texttt{One\mbox{\texttt{\mdseries\slshape Library}}Group( \mbox{\texttt{\mdseries\slshape fun1}}, \mbox{\texttt{\mdseries\slshape val1}}, ... )} 

 For each \texttt{All\mbox{\texttt{\mdseries\slshape Library}}Groups} function (see above) there is a corresponding function \texttt{One\mbox{\texttt{\mdseries\slshape Library}}Group} on exactly the same arguments, i.e., there are \texttt{OneSmallGroup} (\textbf{smallgrp: OneSmallGroup}), \texttt{OneIrreducibleSolvableGroup} (\textbf{primgrp: OneIrreducibleSolvableGroup}), \texttt{OneTransitiveGroup} (\textbf{transgrp: OneTransitiveGroup}), and \texttt{OnePrimitiveGroup} (\textbf{primgrp: OnePrimitiveGroup}). Each function simply returns \emph{one} group in the library that has the prescribed properties, instead of \emph{all} such groups. It returns \texttt{fail} if no such group exists in the library. }

  
\section{\textcolor{Chapter }{Finite Perfect Groups}}\label{Finite Perfect Groups}
\logpage{[ 50, 6, 0 ]}
\hyperdef{L}{X7A884ECF813C2026}{}
{
  \index{perfect groups} The \textsf{GAP} library of finite perfect groups provides, up to isomorphism, a list of all
perfect groups whose sizes are less than $2\cdot 10^6$. The groups of orders up to $10^6$ have been enumerated by Derek{\nobreakspace}F. Holt and Wilhelm Plesken and
published in their book ``Perfect Groups'' \cite{HP89}. For orders $n = 86016$, 368640, or 737280 this work only counted the groups (but did not explicitly
list them), the groups of orders $n = 61440$, 122880, 172032, 245760, 344064, 491520, 688128, or 983040 were omitted. 

 We are grateful to Derek Holt and Wilhelm Plesken for making their groups
available to the \textsf{GAP} community by contributing their files. It should be noted that their book
contains a lot of further information for many of the library groups. So we
would like to recommend it to any \textsf{GAP} user who is interested in the groups. The library of these has been brought
into \textsf{GAP} format by Volkmar Felsch. 

 Several additional groups omitted from the book ``Perfect Groups'' have also been included. Two groups \texttt{\symbol{45}}\texttt{\symbol{45}}
one of order 450000 with a factor group of type $A_6$ and the one of order 962280 \texttt{\symbol{45}}\texttt{\symbol{45}} were
found by Jack Schmidt in 2005. Two groups of order 243000 and one each of
orders 729000, 871200, 878460 were found in 2020 by Alexander Hulpke. 

 The perfect groups of size less than $2\cdot 10^6$ which had not been classified in the work of Holt and Plesken have been
enumerated by Alexander Hulpke. They are stored directly and provide less
construction information in their names. 

 

 As all groups are stored by presentations, a permutation representation is
obtained by coset enumeration. Note that some of the library groups do not
have a faithful permutation representation of small degree. Computations in
these groups may be rather time consuming. 

\subsection{\textcolor{Chapter }{SizesPerfectGroups}}
\logpage{[ 50, 6, 1 ]}\nobreak
\hyperdef{L}{X866A25F882A4E97B}{}
{\noindent\textcolor{FuncColor}{$\triangleright$\enspace\texttt{SizesPerfectGroups({\mdseries\slshape })\index{SizesPerfectGroups@\texttt{SizesPerfectGroups}}
\label{SizesPerfectGroups}
}\hfill{\scriptsize (function)}}\\


 This is the ordered list of all numbers up to $2\cdot 10^6$ that occur as sizes of perfect groups. One can iterate over part of the
perfect groups library with: 
\begin{Verbatim}[commandchars=!@|,fontsize=\small,frame=single,label=Example]
  !gapprompt@gap>| !gapinput@for n in Intersection([100..500],SizesPerfectGroups()) do|
  !gapprompt@>| !gapinput@     for k in [1..NrPerfectGroups(n)] do|
  !gapprompt@>| !gapinput@       pg := PerfectGroup(n,k);|
  !gapprompt@>| !gapinput@     od;|
  !gapprompt@>| !gapinput@   od;|
\end{Verbatim}
 }

 
\subsection{\textcolor{Chapter }{PerfectGroup}}\logpage{[ 50, 6, 2 ]}
\hyperdef{L}{X7906BBA7818E9415}{}
{
\noindent\textcolor{FuncColor}{$\triangleright$\enspace\texttt{PerfectGroup({\mdseries\slshape [filt, ]size[, n]})\index{PerfectGroup@\texttt{PerfectGroup}!for group order (and index)}
\label{PerfectGroup:for group order (and index)}
}\hfill{\scriptsize (function)}}\\
\noindent\textcolor{FuncColor}{$\triangleright$\enspace\texttt{PerfectGroup({\mdseries\slshape [filt, ]sizenumberpair})\index{PerfectGroup@\texttt{PerfectGroup}!for a pair [ order, index ]}
\label{PerfectGroup:for a pair [ order, index ]}
}\hfill{\scriptsize (function)}}\\


 returns a group which is isomorphic to the library group specified by the size
number \texttt{[ \mbox{\texttt{\mdseries\slshape size}}, \mbox{\texttt{\mdseries\slshape n}} ]} or by the two separate arguments \mbox{\texttt{\mdseries\slshape size}} and \mbox{\texttt{\mdseries\slshape n}}, assuming a default value of $\mbox{\texttt{\mdseries\slshape n}} = 1$. The optional argument \mbox{\texttt{\mdseries\slshape filt}} defines the filter in which the group is returned. Possible filters so far are \texttt{IsPermGroup} (\ref{IsPermGroup}) and \texttt{IsSubgroupFpGroup} (\ref{IsSubgroupFpGroup}). In the latter case, the generators and relators used coincide with those
given in{\nobreakspace}\cite{HP89}. The default filter is \texttt{IsPermGroup} (\ref{IsPermGroup}). 
\begin{Verbatim}[commandchars=!@|,fontsize=\small,frame=single,label=Example]
  !gapprompt@gap>| !gapinput@G := PerfectGroup(IsPermGroup,6048,1);|
  U3(3)
  !gapprompt@gap>| !gapinput@G:=PerfectGroup(IsPermGroup,823080,2);|
  A5 2^1 19^2 C 19^1
  !gapprompt@gap>| !gapinput@NrMovedPoints(G);|
  6859
  !gapprompt@gap>| !gapinput@G:=PerfectGroup(1866240,12);|
  PG1866240.12
  !gapprompt@gap>| !gapinput@NrMovedPoints(G);|
  270
\end{Verbatim}
 }

 

\subsection{\textcolor{Chapter }{PerfectIdentification}}
\logpage{[ 50, 6, 3 ]}\nobreak
\hyperdef{L}{X7E1CB2D18085FF9D}{}
{\noindent\textcolor{FuncColor}{$\triangleright$\enspace\texttt{PerfectIdentification({\mdseries\slshape G})\index{PerfectIdentification@\texttt{PerfectIdentification}}
\label{PerfectIdentification}
}\hfill{\scriptsize (attribute)}}\\


 This attribute is set for all groups obtained from the perfect groups library
and has the value \texttt{[\mbox{\texttt{\mdseries\slshape size}},\mbox{\texttt{\mdseries\slshape nr}}]} if the group is obtained with these parameters from the library. }

 

\subsection{\textcolor{Chapter }{NumberPerfectGroups}}
\logpage{[ 50, 6, 4 ]}\nobreak
\hyperdef{L}{X7D68BE547FE5C0F5}{}
{\noindent\textcolor{FuncColor}{$\triangleright$\enspace\texttt{NumberPerfectGroups({\mdseries\slshape size})\index{NumberPerfectGroups@\texttt{NumberPerfectGroups}}
\label{NumberPerfectGroups}
}\hfill{\scriptsize (function)}}\\
\noindent\textcolor{FuncColor}{$\triangleright$\enspace\texttt{NrPerfectGroups({\mdseries\slshape size})\index{NrPerfectGroups@\texttt{NrPerfectGroups}}
\label{NrPerfectGroups}
}\hfill{\scriptsize (function)}}\\
\noindent\textcolor{FuncColor}{$\triangleright$\enspace\texttt{NumberPerfectLibraryGroups({\mdseries\slshape size})\index{NumberPerfectLibraryGroups@\texttt{NumberPerfectLibraryGroups}}
\label{NumberPerfectLibraryGroups}
}\hfill{\scriptsize (function)}}\\
\noindent\textcolor{FuncColor}{$\triangleright$\enspace\texttt{NrPerfectLibraryGroups({\mdseries\slshape size})\index{NrPerfectLibraryGroups@\texttt{NrPerfectLibraryGroups}}
\label{NrPerfectLibraryGroups}
}\hfill{\scriptsize (function)}}\\


 returns the number of non\texttt{\symbol{45}}isomorphic perfect groups of size \mbox{\texttt{\mdseries\slshape size}} for each positive integer \mbox{\texttt{\mdseries\slshape size}} up to $2\cdot10^6$. Additionally, for odd \mbox{\texttt{\mdseries\slshape size}} an answer is returned (odd order groups are solvable). For any other argument
out of range it returns \texttt{fail}. \mbox{\texttt{\mdseries\slshape NrPerfectGroups}} is a synonym for \texttt{NumberPerfectGroups}. Moreover \mbox{\texttt{\mdseries\slshape NumberPerfectLibraryGroups}} (and its synonym \mbox{\texttt{\mdseries\slshape NrPerfectLibraryGroups}}) exist for historical reasons, and return 0 instead of fail for arguments
outside the library scope. }

 

\subsection{\textcolor{Chapter }{SizeNumbersPerfectGroups}}
\logpage{[ 50, 6, 5 ]}\nobreak
\hyperdef{L}{X866356A684F6B15E}{}
{\noindent\textcolor{FuncColor}{$\triangleright$\enspace\texttt{SizeNumbersPerfectGroups({\mdseries\slshape factor1, factor2, ...})\index{SizeNumbersPerfectGroups@\texttt{SizeNumbersPerfectGroups}}
\label{SizeNumbersPerfectGroups}
}\hfill{\scriptsize (function)}}\\


 \texttt{SizeNumbersPerfectGroups} returns a list of pairs, each entry consisting of a group order and the number
of those groups in the library of perfect groups that contain the specified
factors \mbox{\texttt{\mdseries\slshape factor1}}, \mbox{\texttt{\mdseries\slshape factor2}}, ... among their composition factors. 

 Each argument must either be the name of a nonabelian simple group or an
integer which stands for the product of the sizes of one or more cyclic
factors. (In fact, the function replaces all integers among the arguments by
their product.) 

 The following text strings are accepted as simple group names. 
\begin{itemize}
\item  \texttt{A\mbox{\texttt{\mdseries\slshape n}}} or \texttt{A(\mbox{\texttt{\mdseries\slshape n}})} for the alternating groups $A_{\mbox{\texttt{\mdseries\slshape n}}}$, $5 \leq n \leq 9$, for example \texttt{A5} or \texttt{A(6)}. 
\item  \texttt{L\mbox{\texttt{\mdseries\slshape n}}(\mbox{\texttt{\mdseries\slshape q}})} or \texttt{L(\mbox{\texttt{\mdseries\slshape n}},\mbox{\texttt{\mdseries\slshape q}})} for PSL$(n,q)$, where $n \in \{ 2, 3 \}$ and $q$ a prime power, ranging 
\begin{itemize}
\item  for $n = 2$ from 4 to 125 
\item  for $n = 3$ from 2 to 5 
\end{itemize}
 
\item  \texttt{U\mbox{\texttt{\mdseries\slshape n}}(\mbox{\texttt{\mdseries\slshape q}})} or \texttt{U(\mbox{\texttt{\mdseries\slshape n}},\mbox{\texttt{\mdseries\slshape q}})} for PSU$(n,q)$, where $n \in \{ 3, 4 \}$ and $q$ a prime power, ranging 
\begin{itemize}
\item  for $n = 3$ from 3 to 5 
\item  for $n = 4$ from 2 to 2 
\end{itemize}
 
\item  \texttt{Sp4(4)} or \texttt{S(4,4)} for the symplectic group Sp$(4,4)$, 
\item  \texttt{Sz(8)} for the Suzuki group Sz$(8)$, 
\item  \texttt{M\mbox{\texttt{\mdseries\slshape n}}} or \texttt{M(\mbox{\texttt{\mdseries\slshape n}})} for the Mathieu groups $M_{11}$, $M_{12}$, and $M_{22}$, and 
\item  \texttt{J\mbox{\texttt{\mdseries\slshape n}}} or \texttt{J(\mbox{\texttt{\mdseries\slshape n}})} for the Janko groups $J_1$ and $J_2$. 
\end{itemize}
 

 Note that, for most of the groups, the preceding list offers two different
names in order to be consistent with the notation used in \cite{HP89} as well as with the notation used in the \texttt{DisplayCompositionSeries} (\ref{DisplayCompositionSeries}) command of \textsf{GAP}. However, as the names are compared as text strings, you are restricted to
the above choice. Even expressions like \texttt{L2(2\texttt{\symbol{94}}5)} are not accepted. 

 As the use of the term PSU$(n,q)$ is not unique in the literature, we mention that in this library it denotes
the factor group of SU$(n,q)$ by its centre, where SU$(n,q)$ is the group of all $n \times n$ unitary matrices with entries in $GF(q^2)$ and determinant 1. 

 The purpose of the function is to provide a simple way to formulate a loop
over all library groups which contain certain composition factors. }

 
\subsection{\textcolor{Chapter }{DisplayInformationPerfectGroups}}\logpage{[ 50, 6, 6 ]}
\hyperdef{L}{X845419F07BB92867}{}
{
\noindent\textcolor{FuncColor}{$\triangleright$\enspace\texttt{DisplayInformationPerfectGroups({\mdseries\slshape size[, n]})\index{DisplayInformationPerfectGroups@\texttt{DisplayInformationPerfectGroups}!for group order (and index)}
\label{DisplayInformationPerfectGroups:for group order (and index)}
}\hfill{\scriptsize (function)}}\\
\noindent\textcolor{FuncColor}{$\triangleright$\enspace\texttt{DisplayInformationPerfectGroups({\mdseries\slshape sizenumberpair})\index{DisplayInformationPerfectGroups@\texttt{DisplayInformationPerfectGroups}!for a pair [ order, index ]}
\label{DisplayInformationPerfectGroups:for a pair [ order, index ]}
}\hfill{\scriptsize (function)}}\\


 \texttt{DisplayInformationPerfectGroups} displays some invariants of the \mbox{\texttt{\mdseries\slshape n}}\texttt{\symbol{45}}th group of order \mbox{\texttt{\mdseries\slshape size}} from the perfect groups library. 

 If no value of \mbox{\texttt{\mdseries\slshape n}} has been specified, the invariants will be displayed for all groups of size \mbox{\texttt{\mdseries\slshape size}} available in the library. 

 Alternatively, also a list of length two may be entered as the only argument,
with entries \mbox{\texttt{\mdseries\slshape size}} and \mbox{\texttt{\mdseries\slshape n}}. 

 The information provided for $G$ includes the following items: 
\begin{itemize}
\item  a headline containing the size number \texttt{[ \mbox{\texttt{\mdseries\slshape size}}, \mbox{\texttt{\mdseries\slshape n}} ]} of $G$ in the form \texttt{\mbox{\texttt{\mdseries\slshape size}}.\mbox{\texttt{\mdseries\slshape n}}} (the suffix \texttt{.\mbox{\texttt{\mdseries\slshape n}}} will be suppressed if, up to isomorphism, $G$ is the only perfect group of order \mbox{\texttt{\mdseries\slshape size}}), 
\item  a message if $G$ is simple or quasisimple, i.e., if the factor group of $G$ by its centre is simple, 
\item  the ``description'' of the structure of $G$ as it is given by Holt and Plesken in{\nobreakspace}\cite{HP89} (see below), 
\item  the size of the centre of $G$ (suppressed, if $G$ is simple), 
\item  the prime decomposition of the size of $G$, 
\item  orbit sizes for a faithful permutation representation of $G$ which is provided by the library (see below), 
\item  a reference to each occurrence of $G$ in the tables of section 5.3 of \cite{HP89}. Each of these references consists of a class number and an internal number $(i,j)$ under which $G$ is listed in that class. For some groups, there is more than one reference
because these groups belong to more than one of the classes in the book. 
\end{itemize}
 
\begin{Verbatim}[commandchars=!@|,fontsize=\small,frame=single,label=Example]
  !gapprompt@gap>| !gapinput@DisplayInformationPerfectGroups( 30720, 3 );|
  #I Perfect group 30720:  A5 ( 2^4 E N 2^1 E 2^4 ) A
  #I   size = 2^11*3*5  orbit size = 240
  #I   Holt-Plesken class 1 (9,3)
  !gapprompt@gap>| !gapinput@DisplayInformationPerfectGroups( 30720, 6 );|
  #I Perfect group 30720:  A5 ( 2^4 x 2^4 ) C N 2^1
  #I   centre = 2  size = 2^11*3*5  orbit size = 384
  #I   Holt-Plesken class 1 (9,6)
  !gapprompt@gap>| !gapinput@DisplayInformationPerfectGroups( Factorial( 8 ) / 2 );|
  #I Perfect group 20160.1:  A5 x L3(2) 2^1
  #I   centre = 2  size = 2^6*3^2*5*7  orbit sizes = 5 + 16
  #I   Holt-Plesken class 31 (1,1) (occurs also in class 32)
  #I Perfect group 20160.2:  A5 2^1 x L3(2)
  #I   centre = 2  size = 2^6*3^2*5*7  orbit sizes = 7 + 24
  #I   Holt-Plesken class 31 (1,2) (occurs also in class 32)
  #I Perfect group 20160.3:  ( A5 x L3(2) ) 2^1
  #I   centre = 2  size = 2^6*3^2*5*7  orbit size = 192
  #I   Holt-Plesken class 31 (1,3)
  #I Perfect group 20160.4:  simple group  A8
  #I   size = 2^6*3^2*5*7  orbit size = 8
  #I   Holt-Plesken class 26 (0,1)
  #I Perfect group 20160.5:  simple group  L3(4)
  #I   size = 2^6*3^2*5*7  orbit size = 21
  #I   Holt-Plesken class 27 (0,1)
\end{Verbatim}
 }

 
\subsection{\textcolor{Chapter }{More about the Perfect Groups Library}}\label{More about the Perfect Groups Library}
\logpage{[ 50, 6, 7 ]}
\hyperdef{L}{X875C5BE67BAB7F71}{}
{
  For any library group $G$, the library files do not only provide a presentation, but, in addition, a
list of one or more subgroups $S_1, \ldots, S_r$ of $G$ such that there is a faithful permutation representation of $G$ of degree $\sum_{{i = 1}}^r [G:S_i]$ on the set $\{ S_i g \mid 1 \leq i \leq r, g \in G \}$ of the cosets of the $S_i$. This allows one to construct the groups as permutation groups. The function \texttt{DisplayInformationPerfectGroups} (\ref{DisplayInformationPerfectGroups:for group order (and index)}) displays only the available degree. The message 
\begin{Verbatim}[commandchars=!@|,fontsize=\small,frame=single,label=Example]
  orbit size = 8
\end{Verbatim}
 

 in the above example means that the available permutation representation is
transitive and of degree 8, whereas the message 
\begin{Verbatim}[commandchars=!@|,fontsize=\small,frame=single,label=Example]
  orbit sizes = 5 + 16
\end{Verbatim}
 means that a nontransitive permutation representation is available which acts
on two orbits of size 5 and 16 respectively. 

 The notation used in the ``description'' of a group is explained in section{\nobreakspace}5.1.2 of \cite{HP89}. We quote the respective page from there: 

 Within a class $Q\,\#\,p$, an isomorphism type of groups will be denoted by an ordered pair of integers $(r,n)$, where $r \geq 0$ and $n > 0$. More precisely, the isomorphism types in $Q \# p$ of order $p^r |Q|$ will be denoted by $(r,1), (r,2), (r,3), \ldots\,$. Thus $Q$ will always get the size number $(0,1)$. 

 In addition to the symbol $(r,n)$, the groups in $Q\,\#\,p$ will also be given a more descriptive name. The purpose of this is to provide
a very rough idea of the structure of the group. The names are derived in the
following manner. First of all, the isomorphism classes of irreducible $F_pQ$\texttt{\symbol{45}}modules $M$ with $|Q|.|M| \leq 10^6$, where $F_p$ is the field of order $p$, are assigned symbols. These will either be simply $p^x$, where $x$ is the dimension of the module, or, if there is more than one isomorphism
class of irreducible modules having the same dimension, they will be denoted
by$p^x$, $p^{{x'}}$, etc. The one\texttt{\symbol{45}}dimensional module with trivial $Q$\texttt{\symbol{45}}action will therefore be denoted by $p^1$. These symbols will be listed under the description of $Q$. The group name consists essentially of a list of the composition factors
working from the top of the group downwards; hence it always starts with the
name of $Q$ itself. (This convention is the most convenient in our context, but it is
different from that adopted in the ATLAS \cite{CCN85}, for example, where composition factors are listed in the reverse order. For
example, we denote a group isomorphic to $SL(2,5)$ by $A_5 2^1$ rather than $2.A_5$.) 

 Some other symbols are used in the name, in order to give some idea of the
relationship between these composition factors, and splitting properties. We
shall now list these additional symbols. 

 
\begin{description}
\item[{$\times$}]  between two factors denotes a direct product of $F_pQ$\texttt{\symbol{45}}modules or groups. 
\item[{C}]  (for ``commutator'') between two factors means that the second lies in the commutator subgroup of
the first. Similarly, a segment of the form $(f_1 \! \times \! f_2) C f_3$ would mean that the factors $f_1$ and $f_2$ commute modulo $f_3$ and $f_3$ lies in $[f_1,f_2]$. 
\item[{A}]  (for ``abelian'') between two factors indicates that the second is in the $p$th power (but not the commutator subgroup) of the first. ``A'' may also follow the factors, if bracketed. 
\item[{E}]  (for ``elementary abelian'') between two factors indicates that together they generate an elementary
abelian group (modulo subsequent factors), but that the resulting $F_p Q$\texttt{\symbol{45}}module extension does not split. 
\item[{N}]  (for ``nonsplit'') before a factor indicates that $Q$ (or possibly its covering group) splits down as far at this factor but not
over the factor itself. So ``$Q f_1 N f_2$'' means that the normal subgroup $f_1 f_2$ of the group has no complement but, modulo $f_2$, $f_1$, does have a complement. 
\end{description}
 

 Brackets have their obvious meaning. Summarizing, we have: 

 
\begin{description}
\item[{$\times$}]  = direct product; 
\item[{C}]  = commutator subgroup; 
\item[{A}]  = abelian; 
\item[{E}]  = elementary abelian; and 
\item[{N}]  = nonsplit. 
\end{description}
 

 Here are some examples. 

 
\begin{description}
\item[{(i)}]  $A_5 (2^4 E 2^1 E 2^4) A$ means that the pairs $2^4 E 2^1$ and $2^1 E 2^4$ are both elementary abelian of exponent 4. 
\item[{(ii)}]  $A_5 (2^4 E 2^1 A) C 2^1$ means that $O_2(G)$ is of symplectic type $2^{{1+5}}$, with Frattini factor group of type $2^4 E 2^1$. The ``A'' after the $2^1$ indicates that $G$ has a central cyclic subgroup $2^1 A 2^1$ of order 4. 
\item[{(iii)}]  $L_3(2) ((2^1 E) \! \times \! ( N 2^3 E 2^{{3'}} A) C) 2^{{3'}}$ means that the $2^{{3'}}$ factor at the bottom lies in the commutator subgroup of the pair $2^3 E 2^{{3'}}$ in the middle, but the lower pair $2^{{3'}} A 2^{{3'}}$ is abelian of exponent 4. There is also a submodule $2^1 E 2^{{3'}}$, and the covering group $L_3(2) 2^1$ of $L_3(2)$ does not split over the $2^3$ factor. (Since $G$ is perfect, it goes without saying that the extension $L_3(2) 2^1$ cannot split itself.) 
\end{description}
 

 We must stress that this notation does not always succeed in being precise or
even unambiguous, and the reader is free to ignore it if it does not seem
helpful. 

 If such a group description has been given in the book for $G$ (and, in fact, this is the case for most of the library groups), it is
displayed by \texttt{DisplayInformationPerfectGroups} (\ref{DisplayInformationPerfectGroups:for group order (and index)}). Otherwise the function provides a less explicit description of the (in these
cases unique) Holt\texttt{\symbol{45}}Plesken class to which $G$ belongs, together with a serial number if this is necessary to make it unique. }

 }

  
\section{\textcolor{Chapter }{Irreducible Maximal Finite Integral Matrix Groups}}\label{Irreducible Maximal Finite Integral Matrix Groups}
\logpage{[ 50, 7, 0 ]}
\hyperdef{L}{X7873506D873EDB95}{}
{
  A library of irreducible maximal finite integral matrix groups is provided
with \textsf{GAP}. It contains ${\ensuremath{\mathbb Q}}$\texttt{\symbol{45}}class representatives for all of these groups of dimension
at most 31, and ${\ensuremath{\mathbb Z}}$\texttt{\symbol{45}}class representatives for those of dimension at most 11 or
of dimension 13, 17, 19, or 23. 

 The groups provided in this library have been determined by Wilhelm Plesken,
partially as joint work with Michael Pohst, or by members of his institute
(Lehrstuhl B f{\"u}r Mathematik, RWTH Aachen). In particular, the data for the
groups of dimensions 2 to 9 have been taken from the output of computer
calculations which they performed in 1979 (see \cite{PP77}, \cite{PP80}). The ${\ensuremath{\mathbb Z}}$\texttt{\symbol{45}}class representatives of the groups of dimension 10 have
been determined and computed by Bernd Souvignier (\cite{Sou94}), and those of dimensions 11, 13, and 17 have been recomputed for this
library from the circulant Gram matrices given in \cite{Ple85}, using the stand\texttt{\symbol{45}}alone programs for the computation of
short vectors and Bravais groups which have been developed in Plesken's
institute. The ${\ensuremath{\mathbb Z}}$\texttt{\symbol{45}}class representatives of the groups of dimensions 19 and
23 had already been determined in \cite{Ple85}. Gabriele Nebe has recomputed them for us. Her main contribution to this
library, however, is that she has determined and computed the ${\ensuremath{\mathbb Q}}$\texttt{\symbol{45}}class representatives of the groups of
non\texttt{\symbol{45}}prime dimensions between 12 and 24 and the groups of
dimensions 25 to 31 (see \cite{PN95}, \cite{NP95}, \cite{Neb95}, \cite{Neb96}). 

 The library has been brought into \textsf{GAP} format by Volkmar Felsch. He has applied several \textsf{GAP} routines to check certain consistency of the data. However, the credit and
responsibility for the lists remain with the authors. We are grateful to
Wilhelm Plesken, Gabriele Nebe, and Bernd Souvignier for supplying their
results to \textsf{GAP}. 

 In the preceding acknowledgement, we used some notations that will also be
needed in the sequel. We first define these. 

 Any integral matrix group $G$ of dimension $n$ is a subgroup of $GL_n({\ensuremath{\mathbb Z}})$ as well as of $GL_n({\ensuremath{\mathbb Q}})$ and hence lies in some conjugacy class of integral matrix groups under $GL_n({\ensuremath{\mathbb Z}})$ and also in some conjugacy class of rational matrix groups under $GL_n({\ensuremath{\mathbb Q}})$. As usual, we call these classes the ${\ensuremath{\mathbb Z}}$\texttt{\symbol{45}}class and the ${\ensuremath{\mathbb Q}}$\texttt{\symbol{45}}class of $G$, respectively. Note that any conjugacy class of subgroups of $GL_n({\ensuremath{\mathbb Q}})$ contains at least one ${\ensuremath{\mathbb Z}}$\texttt{\symbol{45}}class of subgroups of $GL_n({\ensuremath{\mathbb Z}})$ and hence can be considered as the ${\ensuremath{\mathbb Q}}$\texttt{\symbol{45}}class of some integral matrix group. 

 In the context of this library we are only concerned with ${\ensuremath{\mathbb Z}}$\texttt{\symbol{45}}classes and ${\ensuremath{\mathbb Q}}$\texttt{\symbol{45}}classes of subgroups of $GL_n({\ensuremath{\mathbb Z}})$ which are irreducible and maximal finite in $GL_n({\ensuremath{\mathbb Z}})$ (we will call them \emph{i.m.f.}{\nobreakspace}subgroups of $GL_n({\ensuremath{\mathbb Z}})$). We can distinguish two types of these groups: 

 First, there are those i.m.f.{\nobreakspace}subgroups of $GL_n({\ensuremath{\mathbb Z}})$ which are also maximal finite subgroups of $GL_n({\ensuremath{\mathbb Q}})$. Let us denote the set of their ${\ensuremath{\mathbb Q}}$\texttt{\symbol{45}}classes by $Q_1(n)$. It is clear from the above remark that $Q_1(n)$ just consists of the ${\ensuremath{\mathbb Q}}$\texttt{\symbol{45}}classes of i.m.f.{\nobreakspace}subgroups of $GL_n({\ensuremath{\mathbb Q}})$. 

 Secondly, there is the set $Q_2(n)$ of the ${\ensuremath{\mathbb Q}}$\texttt{\symbol{45}}classes of the remaining i.m.f.{\nobreakspace}subgroups of $GL_n({\ensuremath{\mathbb Z}})$, i.e., of those which are not maximal finite subgroups of $GL_n({\ensuremath{\mathbb Q}})$. For any such group $G$, there is at least one class $C \in Q_1(n)$ such that $G$ is conjugate under ${\ensuremath{\mathbb Q}}$ to a proper subgroup of some group $H \in C$. In fact, the class $C$ is uniquely determined for any group $G$ occurring in the library (though there seems to be no reason to assume that
this property should hold in general). Hence we may call $C$ the \emph{rational i.m.f.{\nobreakspace}class} of $G$. Finally, we will denote the number of classes in $Q_1(n)$ and $Q_2(n)$ by $q_1(n)$ and $q_2(n)$, respectively. 

 As an example, let us consider the case $n = 4$. There are 6 ${\ensuremath{\mathbb Z}}$\texttt{\symbol{45}}classes of i.m.f.{\nobreakspace}subgroups of $GL_4({\ensuremath{\mathbb Z}})$ with representative subgroups $G_1, \ldots, G_6$ of isomorphism types $G_1 \cong W(F_4)$, $G_2 \cong D_{12} \wr C_2$, $G_3 \cong G_4 \cong C_2 \times S_5$, $G_5 \cong W(B_4)$, and $G_6 \cong (D_{12} $\texttt{Y}$ D_{12}) \!:\! C_2$. The corresponding ${\ensuremath{\mathbb Q}}$\texttt{\symbol{45}}classes, which we denote $R_1, \ldots, R_6$, are pairwise different except that $R_3$ coincides with $R_4$. The groups $G_1$, $G_2$, and $G_3$ are i.m.f.{\nobreakspace}subgroups of $GL_4({\ensuremath{\mathbb Q}})$, but $G_5$ and $G_6$ are not because they are conjugate under $GL_4({\ensuremath{\mathbb Q}})$ to proper subgroups of $G_1$ and $G_2$, respectively. So we have $Q_1(4) = \{ R_1, R_2, R_3 \}$, $Q_2(4) = \{ R_5, R_6 \}$, $q_1(4) = 3$, and $q_2(4) = 2$. 

 The $q_1(n)$ ${\ensuremath{\mathbb Q}}$\texttt{\symbol{45}}classes of i.m.f.{\nobreakspace}subgroups of $GL_n({\ensuremath{\mathbb Q}})$ have been determined for each dimension $n \leq 31$. The current \textsf{GAP} library provides integral representative groups for all these classes.
Moreover, all ${\ensuremath{\mathbb Z}}$\texttt{\symbol{45}}classes of i.m.f.{\nobreakspace}subgroups of $GL_n({\ensuremath{\mathbb Z}})$ are known for $n \leq 11$ and for $n \in \{13,17,19,23\}$. For these dimensions, the library offers integral representative groups for
all ${\ensuremath{\mathbb Q}}$\texttt{\symbol{45}}classes in $Q_1(n)$ and $Q_2(n)$ as well as for all ${\ensuremath{\mathbb Z}}$\texttt{\symbol{45}}classes of i.m.f.{\nobreakspace}subgroups of $GL_n({\ensuremath{\mathbb Z}})$. 

 Any group $G$ of dimension $n$ given in the library is represented as the automorphism group $G = Aut(F,L) = \{ g \in GL_n({\ensuremath{\mathbb Z}}) \mid Lg = L, g F g^{tr}
= F \}$ of a positive definite symmetric $n \times n$ matrix $F \in {\ensuremath{\mathbb Z}}^{{n \times n}}$ on an $n$\texttt{\symbol{45}}dimensional lattice $L \cong {\ensuremath{\mathbb Z}}^{{1 \times n}}$ (for details see e.g. \cite{PN95}). \textsf{GAP} provides for $G$ a list of matrix generators and the \emph{Gram matrix} $F$. 

 The positive definite quadratic form defined by $F$ defines a \emph{norm} $v F v^{tr}$ for each vector $v \in L$, and there is only a finite set of vectors of minimal norm. These vectors are
often simply called the \emph{short vectors}. Their set splits into orbits under $G$, and $G$ being irreducible acts faithfully on each of these orbits by multiplication
from the right. \textsf{GAP} provides for each of these orbits the orbit size and a representative vector. 

 Like most of the other \textsf{GAP} libraries, the library of i.m.f.{\nobreakspace}integral matrix groups supplies
an extraction function, \texttt{ImfMatrixGroup}. However, as the library involves only 525 different groups, there is no need
for a selection or an example function. Instead, there are two functions, \texttt{ImfInvariants} (\ref{ImfInvariants}) and \texttt{DisplayImfInvariants} (\ref{DisplayImfInvariants}), which provide some ${\ensuremath{\mathbb Z}}$\texttt{\symbol{45}}class invariants that can be extracted from the library
without actually constructing the representative groups themselves. The
difference between these two functions is that the latter one displays the
resulting data in some easily readable format, whereas the first one returns
them as record components so that you can properly access them. 

 We shall give an individual description of each of the library functions, but
first we would like to insert a short remark concerning their names: Any
self\texttt{\symbol{45}}explaining name of a function handling \emph{irreducible maximal finite integral matrix groups} would have to include this term in full length and hence would grow extremely
long. Therefore we have decided to use the abbreviation \texttt{Imf} instead in order to restrict the names to some reasonable length. 

 The first three functions can be used to formulate loops over the classes. 

 

\subsection{\textcolor{Chapter }{ImfNumberQQClasses}}
\logpage{[ 50, 7, 1 ]}\nobreak
\hyperdef{L}{X8693FD647EF3C53B}{}
{\noindent\textcolor{FuncColor}{$\triangleright$\enspace\texttt{ImfNumberQQClasses({\mdseries\slshape dim})\index{ImfNumberQQClasses@\texttt{ImfNumberQQClasses}}
\label{ImfNumberQQClasses}
}\hfill{\scriptsize (function)}}\\
\noindent\textcolor{FuncColor}{$\triangleright$\enspace\texttt{ImfNumberQClasses({\mdseries\slshape dim})\index{ImfNumberQClasses@\texttt{ImfNumberQClasses}}
\label{ImfNumberQClasses}
}\hfill{\scriptsize (function)}}\\
\noindent\textcolor{FuncColor}{$\triangleright$\enspace\texttt{ImfNumberZClasses({\mdseries\slshape dim, q})\index{ImfNumberZClasses@\texttt{ImfNumberZClasses}}
\label{ImfNumberZClasses}
}\hfill{\scriptsize (function)}}\\


 \texttt{ImfNumberQQClasses} returns the number $q_1($\mbox{\texttt{\mdseries\slshape dim}}$)$ of ${\ensuremath{\mathbb Q}}$\texttt{\symbol{45}}classes of i.m.f.{\nobreakspace}rational matrix groups of
dimension \mbox{\texttt{\mdseries\slshape dim}}. Valid values of \mbox{\texttt{\mdseries\slshape dim}} are all positive integers up to 31. 

 Note: In order to enable you to loop just over the classes belonging to $Q_1($\mbox{\texttt{\mdseries\slshape dim}}$)$, we have arranged the list of ${\ensuremath{\mathbb Q}}$\texttt{\symbol{45}}classes of dimension \mbox{\texttt{\mdseries\slshape dim}} for any dimension \mbox{\texttt{\mdseries\slshape dim}} in the library such that, whenever the classes of $Q_2($\mbox{\texttt{\mdseries\slshape dim}}$)$ are known, too, i.e., in the cases $dim \leq 11$ or $dim \in \{13,17,19,23\}$, the classes of $Q_1($\mbox{\texttt{\mdseries\slshape dim}}$)$ precede those of $Q_2($\mbox{\texttt{\mdseries\slshape dim}}$)$ and hence are numbered from 1 to $q_1($\mbox{\texttt{\mdseries\slshape dim}}$)$. 

 \texttt{ImfNumberQClasses} returns the number of ${\ensuremath{\mathbb Q}}$\texttt{\symbol{45}}classes of groups of dimension \mbox{\texttt{\mdseries\slshape dim}} which are available in the library. If $dim \leq 11$ or $dim \in \{13,17,19,23\}$, this is the number $q_1($\mbox{\texttt{\mdseries\slshape dim}}$) + q_2($\mbox{\texttt{\mdseries\slshape dim}}$)$ of ${\ensuremath{\mathbb Q}}$\texttt{\symbol{45}}classes of i.m.f.{\nobreakspace}subgroups of $GL_{dim}({\ensuremath{\mathbb Z}})$. Otherwise, it is just the number $q_1($\mbox{\texttt{\mdseries\slshape dim}}$)$ of ${\ensuremath{\mathbb Q}}$\texttt{\symbol{45}}classes of i.m.f.{\nobreakspace}subgroups of $GL_{dim}({\ensuremath{\mathbb Q}})$. Valid values of \mbox{\texttt{\mdseries\slshape dim}} are all positive integers up to 31. 

 \texttt{ImfNumberZClasses} returns the number of ${\ensuremath{\mathbb Z}}$\texttt{\symbol{45}}classes in the \mbox{\texttt{\mdseries\slshape q}}\texttt{\symbol{45}}th ${\ensuremath{\mathbb Q}}$\texttt{\symbol{45}}class of i.m.f.{\nobreakspace}integral matrix groups of
dimension \mbox{\texttt{\mdseries\slshape dim}}. Valid values of \mbox{\texttt{\mdseries\slshape dim}} are all positive integers up to 11 and all primes up to 23. }

 

 

\subsection{\textcolor{Chapter }{DisplayImfInvariants}}
\logpage{[ 50, 7, 2 ]}\nobreak
\hyperdef{L}{X8705F64B7E19DDC7}{}
{\noindent\textcolor{FuncColor}{$\triangleright$\enspace\texttt{DisplayImfInvariants({\mdseries\slshape dim, q[, z]})\index{DisplayImfInvariants@\texttt{DisplayImfInvariants}}
\label{DisplayImfInvariants}
}\hfill{\scriptsize (function)}}\\


 \texttt{DisplayImfInvariants} displays the following ${\ensuremath{\mathbb Z}}$\texttt{\symbol{45}}class invariants of the groups in the \mbox{\texttt{\mdseries\slshape z}}\texttt{\symbol{45}}th ${\ensuremath{\mathbb Z}}$\texttt{\symbol{45}}class in the \mbox{\texttt{\mdseries\slshape q}}\texttt{\symbol{45}}th ${\ensuremath{\mathbb Q}}$\texttt{\symbol{45}}class of i.m.f.{\nobreakspace}integral matrix groups of
dimension \mbox{\texttt{\mdseries\slshape dim}}: 

 
\begin{itemize}
\item  its ${\ensuremath{\mathbb Z}}$\texttt{\symbol{45}}class number in the form \mbox{\texttt{\mdseries\slshape dim}}.\mbox{\texttt{\mdseries\slshape q}}.\mbox{\texttt{\mdseries\slshape z}}, if \mbox{\texttt{\mdseries\slshape dim}} is at most 11 or a prime at most 23, or its ${\ensuremath{\mathbb Q}}$\texttt{\symbol{45}}class number in the form \mbox{\texttt{\mdseries\slshape dim}}.\mbox{\texttt{\mdseries\slshape q}}, else, 
\item  a message if the group is solvable, 
\item  the size of the group, 
\item  the isomorphism type of the group, 
\item  the elementary divisors of the associated quadratic form, 
\item  the sizes of the orbits of short vectors (these sizes are the degrees of the
faithful permutation representations which you may construct using the
functions \texttt{IsomorphismPermGroup} (\ref{IsomorphismPermGroup:for Imf matrix groups}) or \texttt{IsomorphismPermGroupImfGroup} (\ref{IsomorphismPermGroupImfGroup}) below), 
\item  the norm of the associated short vectors, 
\item  only in case that the group is not an i.m.f.{\nobreakspace}group in $GL_n({\ensuremath{\mathbb Q}})$: an appropriate message, including the ${\ensuremath{\mathbb Q}}$\texttt{\symbol{45}}class number of the corresponding rational
i.m.f.{\nobreakspace}class. 
\end{itemize}
 

 If you specify the value 0 for any of the parameters \mbox{\texttt{\mdseries\slshape dim}}, \mbox{\texttt{\mdseries\slshape q}}, or \mbox{\texttt{\mdseries\slshape z}}, the command will loop over all available dimensions, ${\ensuremath{\mathbb Q}}$\texttt{\symbol{45}}classes of given dimension, or ${\ensuremath{\mathbb Z}}$\texttt{\symbol{45}}classes within the given ${\ensuremath{\mathbb Q}}$\texttt{\symbol{45}}class, respectively. Otherwise, the values of the
arguments must be in range. A value \mbox{\texttt{\mdseries\slshape z}} $\neq 1$ must not be specified if the ${\ensuremath{\mathbb Z}}$\texttt{\symbol{45}}classes are not known for the given dimension, i.e., if \mbox{\texttt{\mdseries\slshape dim}} $> 11$ and \mbox{\texttt{\mdseries\slshape dim}} $\not \in \{ 13, 17, 19, 23 \}$. The default value of \mbox{\texttt{\mdseries\slshape z}} is{\nobreakspace}1. This value of \mbox{\texttt{\mdseries\slshape z}} will be accepted even if the ${\ensuremath{\mathbb Z}}$\texttt{\symbol{45}}classes are not known. Then it specifies the only
representative group which is available for the \mbox{\texttt{\mdseries\slshape q}}\texttt{\symbol{45}}th ${\ensuremath{\mathbb Q}}$\texttt{\symbol{45}}class. The greatest legal value of \mbox{\texttt{\mdseries\slshape dim}} is 31. 

 
\begin{Verbatim}[commandchars=!@|,fontsize=\small,frame=single,label=Example]
  !gapprompt@gap>| !gapinput@DisplayImfInvariants( 3, 1, 0 );|
  #I Z-class 3.1.1:  Solvable, size = 2^4*3
  #I   isomorphism type = C2 wr S3 = C2 x S4 = W(B3)
  #I   elementary divisors = 1^3
  #I   orbit size = 6, minimal norm = 1
  #I Z-class 3.1.2:  Solvable, size = 2^4*3
  #I   isomorphism type = C2 wr S3 = C2 x S4 = C2 x W(A3)
  #I   elementary divisors = 1*4^2
  #I   orbit size = 8, minimal norm = 3
  #I Z-class 3.1.3:  Solvable, size = 2^4*3
  #I   isomorphism type = C2 wr S3 = C2 x S4 = C2 x W(A3)
  #I   elementary divisors = 1^2*4
  #I   orbit size = 12, minimal norm = 2
  !gapprompt@gap>| !gapinput@DisplayImfInvariants( 8, 15, 1 );|
  #I Z-class 8.15.1:  Solvable, size = 2^5*3^4
  #I   isomorphism type = C2 x (S3 wr S3)
  #I   elementary divisors = 1*3^3*9^3*27
  #I   orbit size = 54, minimal norm = 8
  #I   not maximal finite in GL(8,Q), rational imf class is 8.5
  !gapprompt@gap>| !gapinput@DisplayImfInvariants( 20, 23 );|
  #I Q-class 20.23:  Size = 2^5*3^2*5*11
  #I   isomorphism type = (PSL(2,11) x D12).C2
  #I   elementary divisors = 1^18*11^2
  #I   orbit size = 3*660 + 2*1980 + 2640 + 3960, minimal norm = 4
\end{Verbatim}
 

 Note that the function \texttt{DisplayImfInvariants} uses a kind of shorthand to display the elementary divisors.
E.{\nobreakspace}g., the expression \texttt{1*3\texttt{\symbol{94}}3*9\texttt{\symbol{94}}3*27} in the preceding example stands for the elementary divisors $1,3,3,3,9,9,9,27$. (See also the next example which shows that the function \texttt{ImfInvariants} (\ref{ImfInvariants}) provides the elementary divisors in form of an ordinary \textsf{GAP} list.) 

 In the description of the isomorphism types the following notations are used: 
\begin{description}
\item[{$A$ \texttt{x} $B$ }]  denotes a direct product of a group $A$ by a group $B$, 
\item[{$A$ \texttt{subd} $B$ }]  denotes a subdirect product of $A$ by $B$, 
\item[{$A$ \texttt{Y} $B$ }]  denotes a central product of $A$ by $B$, 
\item[{$A$ \texttt{wr} $B$ }]  denotes a wreath product of $A$ by $B$, 
\item[{$A$\texttt{:}$B$ }]  denotes a split extension of $A$ by $B$, 
\item[{$A$\texttt{.}$B$ }]  denotes just an extension of $A$ by $B$ (split or nonsplit). 
\end{description}
 

 The groups involved are 
\begin{itemize}
\item  the cyclic groups $C_n$, dihedral groups $D_n$, and generalized quaternion groups $Q_n$ of order $n$, denoted by \texttt{C}\mbox{\texttt{\mdseries\slshape n}}, \texttt{D}\mbox{\texttt{\mdseries\slshape n}}, and \texttt{Q}\mbox{\texttt{\mdseries\slshape n}}, respectively, 
\item  the alternating groups $A_n$ and symmetric groups $S_n$ of degree $n$, denoted by \texttt{A}\mbox{\texttt{\mdseries\slshape n}} and \texttt{S}\mbox{\texttt{\mdseries\slshape n}}, respectively, 
\item  the linear groups $GL_n(q)$, $PGL_n(q)$, $SL_n(q)$, and $PSL_n(q)$, denoted by \texttt{GL}(\mbox{\texttt{\mdseries\slshape n}},\mbox{\texttt{\mdseries\slshape q}}), \texttt{PGL}(\mbox{\texttt{\mdseries\slshape n}},\mbox{\texttt{\mdseries\slshape q}}), \texttt{SL}(\mbox{\texttt{\mdseries\slshape n}},\mbox{\texttt{\mdseries\slshape q}}), and \texttt{PSL}(\mbox{\texttt{\mdseries\slshape n}},\mbox{\texttt{\mdseries\slshape q}}), respectively, 
\item  the unitary groups $SU_n(q)$ and $PSU_n(q)$, denoted by \texttt{SU}(\mbox{\texttt{\mdseries\slshape n}},\mbox{\texttt{\mdseries\slshape q}}) and \texttt{PSU}(\mbox{\texttt{\mdseries\slshape n}},\mbox{\texttt{\mdseries\slshape q}}), respectively, 
\item  the symplectic groups $Sp(n,q)$ and $PSp(n,q)$, denoted by \texttt{Sp}(\mbox{\texttt{\mdseries\slshape n}},\mbox{\texttt{\mdseries\slshape q}}) and \texttt{PSp}(\mbox{\texttt{\mdseries\slshape n}},\mbox{\texttt{\mdseries\slshape q}}), respectively, 
\item  the orthogonal groups $O_8^+(2)$ and $PO_8^+(2)$, denoted by \texttt{O+}(8,2) and \texttt{PO+}(8,2), respectively, 
\item  the extraspecial groups $2_+^{{1+8}}$, $3_+^{{1+2}}$, $3_+^{{1+4}}$, and $5_+^{{1+2}}$, denoted by \texttt{2+\texttt{\symbol{94}}(1+8)}, \texttt{3+\texttt{\symbol{94}}(1+2)}, \texttt{3+\texttt{\symbol{94}}(1+4)}, and \texttt{5+\texttt{\symbol{94}}(1+2)}, respectively, 
\item  the Chevalley group $G_2(3)$, denoted by \texttt{G2(3)}, 
\item  the twisted Chevalley group ${^3}D_4(2)$, denoted by \texttt{3D4(2)}, 
\item  the Suzuki group $Sz(8)$, denoted by \texttt{Sz(8)}, 
\item  the Weyl groups $W(A_n)$, $W(B_n)$, $W(D_n)$, $W(E_n)$, and $W(F_4)$, denoted by \texttt{W(A\mbox{\texttt{\mdseries\slshape n}})}, \texttt{W(B\mbox{\texttt{\mdseries\slshape n}})}, \texttt{W(D\mbox{\texttt{\mdseries\slshape n}})}, \texttt{W(E\mbox{\texttt{\mdseries\slshape n}})}, and \texttt{W(F4)}, respectively, 
\item  the sporadic simple groups $Co_1$, $Co_2$, $Co_3$, $HS$, $J_2$, $M_{12}$, $M_{22}$, $M_{23}$, $M_{24}$, and $Mc$, denoted by \texttt{Co1}, \texttt{Co2}, \texttt{Co3}, \texttt{HS}, \texttt{J2}, \texttt{M12}, \texttt{M22}, \texttt{M23}, \texttt{M24}, and \texttt{Mc}, respectively, 
\item  a point stabilizer of index 11 in $M_{11}$, denoted by \texttt{M10}. 
\end{itemize}
 

 As mentioned above, the data assembled by the function \texttt{DisplayImfInvariants} are ``cheap data'' in the sense that they can be provided by the library without loading any of
its large matrix files or performing any matrix calculations. The following
function allows you to get proper access to these cheap data instead of just
displaying them. }

 

\subsection{\textcolor{Chapter }{ImfInvariants}}
\logpage{[ 50, 7, 3 ]}\nobreak
\hyperdef{L}{X8604A2167B2E8434}{}
{\noindent\textcolor{FuncColor}{$\triangleright$\enspace\texttt{ImfInvariants({\mdseries\slshape dim, q[, z]})\index{ImfInvariants@\texttt{ImfInvariants}}
\label{ImfInvariants}
}\hfill{\scriptsize (function)}}\\


 \texttt{ImfInvariants} returns a record which provides some ${\ensuremath{\mathbb Z}}$\texttt{\symbol{45}}class invariants of the groups in the \mbox{\texttt{\mdseries\slshape z}}\texttt{\symbol{45}}th ${\ensuremath{\mathbb Z}}$\texttt{\symbol{45}}class in the \mbox{\texttt{\mdseries\slshape q}}\texttt{\symbol{45}}th ${\ensuremath{\mathbb Q}}$\texttt{\symbol{45}}class of i.m.f.{\nobreakspace}integral matrix groups of
dimension \mbox{\texttt{\mdseries\slshape dim}}. A value \mbox{\texttt{\mdseries\slshape z}} $\neq 1$ must not be specified if the ${\ensuremath{\mathbb Z}}$\texttt{\symbol{45}}classes are not known for the given dimension, i.e., if \mbox{\texttt{\mdseries\slshape dim}} $> 11$ and \mbox{\texttt{\mdseries\slshape dim}} $\not \in \{ 13, 17, 19, 23 \}$. The default value of \mbox{\texttt{\mdseries\slshape z}} is{\nobreakspace}1. This value of \mbox{\texttt{\mdseries\slshape z}} will be accepted even if the ${\ensuremath{\mathbb Z}}$\texttt{\symbol{45}}classes are not known. Then it specifies the only
representative group which is available for the \mbox{\texttt{\mdseries\slshape q}}\texttt{\symbol{45}}th ${\ensuremath{\mathbb Q}}$\texttt{\symbol{45}}class. The greatest legal value of \mbox{\texttt{\mdseries\slshape dim}} is 31. 

 The resulting record contains six or seven components: 

 
\begin{description}
\item[{\texttt{size} }]  the size of any representative group \mbox{\texttt{\mdseries\slshape G}}, 
\item[{\texttt{isSolvable} }]  is \texttt{true} if \mbox{\texttt{\mdseries\slshape G}} is solvable, 
\item[{\texttt{isomorphismType} }]  a text string describing the isomorphism type of \mbox{\texttt{\mdseries\slshape G}} (in the same notation as used by the function \texttt{DisplayImfInvariants} above), 
\item[{\texttt{elementaryDivisors} }]  the elementary divisors of the associated Gram matrix \mbox{\texttt{\mdseries\slshape F}} (in the same format as the result of the function \texttt{ElementaryDivisorsMat} (\ref{ElementaryDivisorsMat}), 
\item[{\texttt{minimalNorm} }]  the norm of the associated short vectors, 
\item[{\texttt{sizesOrbitsShortVectors} }]  the sizes of the orbits of short vectors under \mbox{\texttt{\mdseries\slshape F}}, 
\item[{\texttt{maximalQClass} }]  the ${\ensuremath{\mathbb Q}}$\texttt{\symbol{45}}class number of an i.m.f.{\nobreakspace}group in $GL_n({\ensuremath{\mathbb Q}})$ that contains \mbox{\texttt{\mdseries\slshape G}} as a subgroup (only in case that not \mbox{\texttt{\mdseries\slshape G}} itself is an i.m.f.{\nobreakspace}subgroup of $GL_n({\ensuremath{\mathbb Q}})$). 
\end{description}
 

 Note that four of these data, namely the group size, the solvability, the
isomorphism type, and the corresponding rational i.m.f.{\nobreakspace}class,
are not only ${\ensuremath{\mathbb Z}}$\texttt{\symbol{45}}class invariants, but also ${\ensuremath{\mathbb Q}}$\texttt{\symbol{45}}class invariants. 

 Note further that, though the isomorphism type is a ${\ensuremath{\mathbb Q}}$\texttt{\symbol{45}}class invariant, you will sometimes get different
descriptions for different ${\ensuremath{\mathbb Z}}$\texttt{\symbol{45}}classes of the same ${\ensuremath{\mathbb Q}}$\texttt{\symbol{45}}class (as, e.g., for the classes 3.1.1 and 3.1.2 in the
last example above). The purpose of this behaviour is to provide some more
information about the underlying lattices. 

 
\begin{Verbatim}[commandchars=!@|,fontsize=\small,frame=single,label=Example]
  !gapprompt@gap>| !gapinput@ImfInvariants( 8, 15, 1 );|
  rec( elementaryDivisors := [ 1, 3, 3, 3, 9, 9, 9, 27 ],
    isSolvable := true, isomorphismType := "C2 x (S3 wr S3)",
    maximalQClass := 5, minimalNorm := 8, size := 2592,
    sizesOrbitsShortVectors := [ 54 ] )
  !gapprompt@gap>| !gapinput@ImfInvariants( 24, 1 ).size;|
  10409396852733332453861621760000
  !gapprompt@gap>| !gapinput@ImfInvariants( 23, 5, 2 ).sizesOrbitsShortVectors;|
  [ 552, 53130 ]
  !gapprompt@gap>| !gapinput@for i in [ 1 .. ImfNumberQClasses( 22 ) ] do|
  !gapprompt@>| !gapinput@   Print( ImfInvariants( 22, i ).isomorphismType, "\n" ); od;|
  C2 wr S22 = W(B22)
  (C2 x PSU(6,2)).S3
  (C2 x S3) wr S11 = (C2 x W(A2)) wr S11
  (C2 x S12) wr C2 = (C2 x W(A11)) wr C2
  C2 x S3 x S12 = C2 x W(A2) x W(A11)
  (C2 x HS).C2
  (C2 x Mc).C2
  C2 x S23 = C2 x W(A22)
  C2 x PSL(2,23)
  C2 x PSL(2,23)
  C2 x PGL(2,23)
  C2 x PGL(2,23)
\end{Verbatim}
 }

 

\subsection{\textcolor{Chapter }{ImfMatrixGroup}}
\logpage{[ 50, 7, 4 ]}\nobreak
\hyperdef{L}{X78935B307B909101}{}
{\noindent\textcolor{FuncColor}{$\triangleright$\enspace\texttt{ImfMatrixGroup({\mdseries\slshape dim, q[, z]})\index{ImfMatrixGroup@\texttt{ImfMatrixGroup}}
\label{ImfMatrixGroup}
}\hfill{\scriptsize (function)}}\\


 \texttt{ImfMatrixGroup} is the essential extraction function of this library (note that its name has
been changed from \texttt{ImfMatGroup} in \textsf{GAP}{\nobreakspace}3 to \texttt{ImfMatrixGroup} in \textsf{GAP}{\nobreakspace}4). It returns a representative group, $G$ say, of the \mbox{\texttt{\mdseries\slshape z}}\texttt{\symbol{45}}th ${\ensuremath{\mathbb Z}}$\texttt{\symbol{45}}class in the \mbox{\texttt{\mdseries\slshape q}}\texttt{\symbol{45}}th ${\ensuremath{\mathbb Q}}$\texttt{\symbol{45}}class of i.m.f.{\nobreakspace}integral matrix groups of
dimension \mbox{\texttt{\mdseries\slshape dim}}. A value \mbox{\texttt{\mdseries\slshape z}}$ \neq 1$ must not be specified if the ${\ensuremath{\mathbb Z}}$\texttt{\symbol{45}}classes are not known for the given dimension, i.e., if \mbox{\texttt{\mdseries\slshape dim}} $> 11$ and \mbox{\texttt{\mdseries\slshape dim}} $\not \in \{ 13, 17, 19, 23 \}$. The default value of \mbox{\texttt{\mdseries\slshape z}} is{\nobreakspace}1. This value of \mbox{\texttt{\mdseries\slshape z}} will be accepted even if the ${\ensuremath{\mathbb Z}}$\texttt{\symbol{45}}classes are not known. Then it specifies the only
representative group which is available for the \mbox{\texttt{\mdseries\slshape q}}\texttt{\symbol{45}}th ${\ensuremath{\mathbb Q}}$\texttt{\symbol{45}}class. The greatest legal value of \mbox{\texttt{\mdseries\slshape dim}} is 31. 

 
\begin{Verbatim}[commandchars=!@|,fontsize=\small,frame=single,label=Example]
  !gapprompt@gap>| !gapinput@G := ImfMatrixGroup( 5, 1, 3 );|
  ImfMatrixGroup(5,1,3)
  !gapprompt@gap>| !gapinput@for m in GeneratorsOfGroup( G ) do PrintArray( m ); od;|
  [ [  -1,   0,   0,   0,   0 ],
    [   0,   1,   0,   0,   0 ],
    [   0,   0,   0,   1,   0 ],
    [  -1,  -1,  -1,  -1,   2 ],
    [  -1,   0,   0,   0,   1 ] ]
  [ [  0,  1,  0,  0,  0 ],
    [  0,  0,  1,  0,  0 ],
    [  0,  0,  0,  1,  0 ],
    [  1,  0,  0,  0,  0 ],
    [  0,  0,  0,  0,  1 ] ]
\end{Verbatim}
 

 The attributes \texttt{Size} (\ref{Size}) and \texttt{IsSolvable} will be properly set in the resulting matrix group $G$. In addition, it has two attributes \texttt{IsImfMatrixGroup} and \texttt{ImfRecord} where the first one is just a logical flag set to \texttt{true} and the latter one is a record. Except for the group size and the solvability
flag, this record contains the same components as the resulting record of the
function \texttt{ImfInvariants} (\ref{ImfInvariants}) described above, namely the components \texttt{isomorphismType}, \texttt{elementaryDivisors}, \texttt{minimalNorm}, and \texttt{sizesOrbitsShortVectors} and, if $G$ is not a rational i.m.f.{\nobreakspace}group, \texttt{maximalQClass}. Moreover, it has the two components 

 
\begin{description}
\item[{\texttt{form}}]  the associated Gram matrix $F$, and 
\item[{\texttt{repsOrbitsShortVectors}}]  representatives of the orbits of short vectors under $F$. 
\end{description}
 

 The last one of these components will be required by the function \texttt{IsomorphismPermGroup} (\ref{IsomorphismPermGroup:for Imf matrix groups}) below. 

 
\begin{Verbatim}[commandchars=!@|,fontsize=\small,frame=single,label=Example]
  !gapprompt@gap>| !gapinput@Size( G );|
  3840
  !gapprompt@gap>| !gapinput@imf := ImfRecord( G );;|
  !gapprompt@gap>| !gapinput@imf.isomorphismType;|
  "C2 wr S5 = C2 x W(D5)"
  !gapprompt@gap>| !gapinput@PrintArray( imf.form );|
  [ [  4,  0,  0,  0,  2 ],
    [  0,  4,  0,  0,  2 ],
    [  0,  0,  4,  0,  2 ],
    [  0,  0,  0,  4,  2 ],
    [  2,  2,  2,  2,  5 ] ]
  !gapprompt@gap>| !gapinput@imf.elementaryDivisors;|
  [ 1, 4, 4, 4, 4 ]
  !gapprompt@gap>| !gapinput@imf.minimalNorm;|
  4
\end{Verbatim}
 

 If you want to perform calculations in such a matrix group $G$ you should be aware of the fact that the permutation group routines of \textsf{GAP} are much more efficient than the matrix group routines. Hence we recommend
that you do your computations, whenever possible, in the isomorphic
permutation group which is induced by the action of $G$ on one of the orbits of the associated short vectors. You may call one of the
following functions \texttt{IsomorphismPermGroup} (\ref{IsomorphismPermGroup:for Imf matrix groups}) or \texttt{IsomorphismPermGroupImfGroup} (\ref{IsomorphismPermGroupImfGroup}) to get an isomorphism to such a permutation group (note that these \textsf{GAP}{\nobreakspace}4 functions have replaced the \textsf{GAP}{\nobreakspace}3 functions \texttt{PermGroup} and \texttt{PermGroupImfGroup}). }

 

\subsection{\textcolor{Chapter }{IsomorphismPermGroup (for Imf matrix groups)}}
\logpage{[ 50, 7, 5 ]}\nobreak
\hyperdef{L}{X84BF34B27CD5E85C}{}
{\noindent\textcolor{FuncColor}{$\triangleright$\enspace\texttt{IsomorphismPermGroup({\mdseries\slshape G})\index{IsomorphismPermGroup@\texttt{IsomorphismPermGroup}!for Imf matrix groups}
\label{IsomorphismPermGroup:for Imf matrix groups}
}\hfill{\scriptsize (method)}}\\


 returns an isomorphism, $\varphi$ say, from the given i.m.f.{\nobreakspace}integral matrix group $G$ to a permutation group $P := \varphi(G)$ acting on a minimal orbit, $S$ say, of short vectors of $G$ such that each matrix $m \in G$ is mapped to the permutation induced by its action on $S$. 

 Note that in case of a large orbit the construction of $\varphi$ may be space and time consuming. Fortunately, there are only six ${\ensuremath{\mathbb Q}}$\texttt{\symbol{45}}classes in the library for which the smallest orbit of
short vectors is of size greater than $20000$, the worst case being the orbit of size $196560$ for the Leech lattice (\mbox{\texttt{\mdseries\slshape dim}} $= 24$, \mbox{\texttt{\mdseries\slshape q}} $= 3$). 

 The inverse isomorphism $\varphi^{{-1}}$ from $P$ to $G$ is constructed by determining a ${\ensuremath{\mathbb Q}}$\texttt{\symbol{45}}base $B \subset S$ of ${\ensuremath{\mathbb Q}}^{{1 \times dim}}$ in $S$ and, in addition, the associated base change matrix $M$ which transforms $B$ into the standard base of ${\ensuremath{\mathbb Z}}^{{1 \times dim}}$. This allows a simple computation of the preimage $\varphi^{{-1}}(p)$ of any permutation $p \in P$, as follows. If, for $1 \leq i \leq$ \mbox{\texttt{\mdseries\slshape dim}}, $b_i$ is the position number in $S$ of the $i$\texttt{\symbol{45}}th base vector in $B$, it suffices to look up the vector whose position number in $S$ is the image of $b_i$ under $p$ and to multiply this vector by $M$ to get the $i$\texttt{\symbol{45}}th row of $\varphi^{{-1}}(p)$. 

 You may use the functions \texttt{Image} (\ref{Image:set of images of the source of a general mapping}) and \texttt{PreImage} (\ref{PreImage:set of preimages of the range of a general mapping}) to switch from $G$ to $P$ and back from $P$ to $G$. 

 As an example, let us continue the preceding example and compute the solvable
residuum of the group $G$. 

 
\begin{Verbatim}[commandchars=!@|,fontsize=\small,frame=single,label=Example]
  !gapprompt@gap>| !gapinput@# Perform the computations in an isomorphic permutation group.|
  !gapprompt@gap>| !gapinput@phi := IsomorphismPermGroup( G );;|
  !gapprompt@gap>| !gapinput@P := Image( phi );|
  Group([ (1,7,6)(2,9)(4,5,10), (2,3,4,5)(6,9,8,7) ])
  !gapprompt@gap>| !gapinput@D := DerivedSubgroup( P );;|
  !gapprompt@gap>| !gapinput@Size( D );|
  960
  !gapprompt@gap>| !gapinput@IsPerfectGroup( D );|
  true
  !gapprompt@gap>| !gapinput@# We have found the solvable residuum of P,|
  !gapprompt@gap>| !gapinput@# now move the results back to the matrix group G.|
  !gapprompt@gap>| !gapinput@R := PreImage( phi, D );;|
  !gapprompt@gap>| !gapinput@StructureDescription(R);|
  "(C2 x C2 x C2 x C2) : A5"
  !gapprompt@gap>| !gapinput@IdGroup(D)=IdGroup(R);|
  true
\end{Verbatim}
 }

 

\subsection{\textcolor{Chapter }{IsomorphismPermGroupImfGroup}}
\logpage{[ 50, 7, 6 ]}\nobreak
\hyperdef{L}{X7CEDB6CE7BAC4518}{}
{\noindent\textcolor{FuncColor}{$\triangleright$\enspace\texttt{IsomorphismPermGroupImfGroup({\mdseries\slshape G, n})\index{IsomorphismPermGroupImfGroup@\texttt{IsomorphismPermGroupImfGroup}}
\label{IsomorphismPermGroupImfGroup}
}\hfill{\scriptsize (function)}}\\


 \texttt{IsomorphismPermGroupImfGroup} returns an isomorphism, $\varphi$ say, from the given i.m.f.{\nobreakspace}integral matrix group \mbox{\texttt{\mdseries\slshape G}} to a permutation group $P$ acting on the \mbox{\texttt{\mdseries\slshape n}}\texttt{\symbol{45}}th orbit, $S$ say, of short vectors of \mbox{\texttt{\mdseries\slshape G}} such that each matrix $m \in$ \mbox{\texttt{\mdseries\slshape G}} is mapped to the permutation induced by its action on $S$. 

 The only difference to the above function \texttt{IsomorphismPermGroup} (\ref{IsomorphismPermGroup:for Imf matrix groups}) is that you can specify the orbit to be used. In fact, as the orbits of short
vectors are sorted by increasing sizes, the function \texttt{IsomorphismPermGroup( \mbox{\texttt{\mdseries\slshape G}} )} has been implemented such that it is equivalent to \texttt{IsomorphismPermGroupImfGroup( \mbox{\texttt{\mdseries\slshape G}}, 1 )}. 

 
\begin{Verbatim}[commandchars=!@|,fontsize=\small,frame=single,label=Example]
  !gapprompt@gap>| !gapinput@ImfInvariants( 12, 9 ).sizesOrbitsShortVectors;|
  [ 120, 300 ]
  !gapprompt@gap>| !gapinput@G := ImfMatrixGroup( 12, 9 );|
  ImfMatrixGroup(12,9)
  !gapprompt@gap>| !gapinput@phi1 := IsomorphismPermGroupImfGroup( G, 1 );;|
  !gapprompt@gap>| !gapinput@P1 := Image( phi1 );|
  <permutation group of size 2400 with 2 generators>
  !gapprompt@gap>| !gapinput@LargestMovedPoint( P1 );|
  120
  !gapprompt@gap>| !gapinput@phi2 := IsomorphismPermGroupImfGroup( G, 2 );;|
  !gapprompt@gap>| !gapinput@P2 := Image( phi2 );|
  <permutation group of size 2400 with 2 generators>
  !gapprompt@gap>| !gapinput@LargestMovedPoint( P2 );|
  300
\end{Verbatim}
 }

 }

 }

            
\chapter{\textcolor{Chapter }{Semigroups and Monoids}}\label{Semigroups}
\logpage{[ 51, 0, 0 ]}
\hyperdef{L}{X8665D8737FDD5B10}{}
{
  This chapter describes functions for creating semigroups and monoids and
determining information about them.  
\section{\textcolor{Chapter }{Semigroups}}\label{sect:IsSemigroup}
\logpage{[ 51, 1, 0 ]}
\hyperdef{L}{X80AF5F307DBDC2B4}{}
{
  

\subsection{\textcolor{Chapter }{IsSemigroup}}
\logpage{[ 51, 1, 1 ]}\nobreak
\hyperdef{L}{X7B412E5B8543E9B7}{}
{\noindent\textcolor{FuncColor}{$\triangleright$\enspace\texttt{IsSemigroup({\mdseries\slshape D})\index{IsSemigroup@\texttt{IsSemigroup}}
\label{IsSemigroup}
}\hfill{\scriptsize (Synonym)}}\\


 returns \texttt{true} if the object \mbox{\texttt{\mdseries\slshape D}} is a semigroup. \index{semigroup} A \emph{semigroup} is a magma (see{\nobreakspace}\ref{Magmas}) with associative multiplication. }

 
\subsection{\textcolor{Chapter }{Semigroup}}\logpage{[ 51, 1, 2 ]}
\hyperdef{L}{X7F55D28F819B2817}{}
{
\noindent\textcolor{FuncColor}{$\triangleright$\enspace\texttt{Semigroup({\mdseries\slshape gen1, gen2, ...})\index{Semigroup@\texttt{Semigroup}!for various generators}
\label{Semigroup:for various generators}
}\hfill{\scriptsize (function)}}\\
\noindent\textcolor{FuncColor}{$\triangleright$\enspace\texttt{Semigroup({\mdseries\slshape gens})\index{Semigroup@\texttt{Semigroup}!for a list}
\label{Semigroup:for a list}
}\hfill{\scriptsize (function)}}\\


 In the first form, \texttt{Semigroup} returns the semigroup generated by the arguments \mbox{\texttt{\mdseries\slshape gen1}}, \mbox{\texttt{\mdseries\slshape gen2}}, $\ldots$, that is, the closure of these elements under multiplication. In the second
form, \texttt{Semigroup} returns the semigroup generated by the elements in the homogeneous list \mbox{\texttt{\mdseries\slshape gens}}; a square matrix as only argument is treated as one generator, not as a list
of generators. 

 It is \emph{not} checked whether the underlying multiplication is associative, use \texttt{Magma} (\ref{Magma}) and \texttt{IsAssociative} (\ref{IsAssociative}) if you want to check whether a magma is in fact a semigroup. 

 
\begin{Verbatim}[commandchars=!@|,fontsize=\small,frame=single,label=Example]
  !gapprompt@gap>| !gapinput@a:= Transformation( [ 2, 3, 4, 1 ] );|
  Transformation( [ 2, 3, 4, 1 ] )
  !gapprompt@gap>| !gapinput@b:= Transformation( [ 2, 2, 3, 4 ] );|
  Transformation( [ 2, 2 ] )
  !gapprompt@gap>| !gapinput@s:= Semigroup(a, b);|
  <transformation semigroup of degree 4 with 2 generators>
\end{Verbatim}
 }

 

\subsection{\textcolor{Chapter }{Subsemigroup}}
\logpage{[ 51, 1, 3 ]}\nobreak
\hyperdef{L}{X8678D40878CC09A1}{}
{\noindent\textcolor{FuncColor}{$\triangleright$\enspace\texttt{Subsemigroup({\mdseries\slshape S, gens})\index{Subsemigroup@\texttt{Subsemigroup}}
\label{Subsemigroup}
}\hfill{\scriptsize (function)}}\\
\noindent\textcolor{FuncColor}{$\triangleright$\enspace\texttt{SubsemigroupNC({\mdseries\slshape S, gens})\index{SubsemigroupNC@\texttt{SubsemigroupNC}}
\label{SubsemigroupNC}
}\hfill{\scriptsize (function)}}\\


 are just synonyms of \texttt{Submagma} (\ref{Submagma}) and \texttt{SubmagmaNC} (\ref{SubmagmaNC}), respectively. 

 
\begin{Verbatim}[commandchars=!@|,fontsize=\small,frame=single,label=Example]
  !gapprompt@gap>| !gapinput@a:=GeneratorsOfSemigroup(s)[1];|
  Transformation( [ 2, 3, 4, 1 ] )
  !gapprompt@gap>| !gapinput@t:=Subsemigroup(s,[a]);|
  <commutative transformation semigroup of degree 4 with 1 generator>
\end{Verbatim}
 }

 

\subsection{\textcolor{Chapter }{IsSubsemigroup}}
\logpage{[ 51, 1, 4 ]}\nobreak
\hyperdef{L}{X782B7BDD8252581C}{}
{\noindent\textcolor{FuncColor}{$\triangleright$\enspace\texttt{IsSubsemigroup({\mdseries\slshape S, T})\index{IsSubsemigroup@\texttt{IsSubsemigroup}}
\label{IsSubsemigroup}
}\hfill{\scriptsize (operation)}}\\
\textbf{\indent Returns:\ }
\texttt{true} or \texttt{false}.



 This operation returns \texttt{true} if the semigroup \mbox{\texttt{\mdseries\slshape T}} is a subsemigroup of the semigroup \mbox{\texttt{\mdseries\slshape S}} and \texttt{false} if it is not. 
\begin{Verbatim}[commandchars=!@|,fontsize=\small,frame=single,label=Example]
  !gapprompt@gap>| !gapinput@f := Transformation([5, 6, 7, 1, 4, 3, 2, 7]);|
  Transformation( [ 5, 6, 7, 1, 4, 3, 2, 7 ] )
  !gapprompt@gap>| !gapinput@T := Semigroup(f);;|
  !gapprompt@gap>| !gapinput@IsSubsemigroup(FullTransformationSemigroup(4), T);|
  false
  !gapprompt@gap>| !gapinput@S := Semigroup(f);;|
  !gapprompt@gap>| !gapinput@T := Semigroup(f ^ 2);;|
  !gapprompt@gap>| !gapinput@IsSubsemigroup(S, T);|
  true
\end{Verbatim}
 }

 

\subsection{\textcolor{Chapter }{SemigroupByGenerators}}
\logpage{[ 51, 1, 5 ]}\nobreak
\hyperdef{L}{X79FBBEC9841544F3}{}
{\noindent\textcolor{FuncColor}{$\triangleright$\enspace\texttt{SemigroupByGenerators({\mdseries\slshape gens})\index{SemigroupByGenerators@\texttt{SemigroupByGenerators}}
\label{SemigroupByGenerators}
}\hfill{\scriptsize (operation)}}\\


 is the underlying operation of{\nobreakspace}\texttt{Semigroup} (\ref{Semigroup:for various generators}). }

 

\subsection{\textcolor{Chapter }{AsSemigroup}}
\logpage{[ 51, 1, 6 ]}\nobreak
\hyperdef{L}{X80ED104F85AE5134}{}
{\noindent\textcolor{FuncColor}{$\triangleright$\enspace\texttt{AsSemigroup({\mdseries\slshape C})\index{AsSemigroup@\texttt{AsSemigroup}}
\label{AsSemigroup}
}\hfill{\scriptsize (operation)}}\\


 If \mbox{\texttt{\mdseries\slshape C}} is a collection whose elements form a semigroup under \texttt{\texttt{\symbol{92}}*} (\ref{bSlash*}) (see{\nobreakspace}\texttt{IsSemigroup} (\ref{IsSemigroup})) then \texttt{AsSemigroup} returns this semigroup. Otherwise \texttt{fail} is returned. }

 

\subsection{\textcolor{Chapter }{AsSubsemigroup}}
\logpage{[ 51, 1, 7 ]}\nobreak
\hyperdef{L}{X7B1EEA3E82BFE09F}{}
{\noindent\textcolor{FuncColor}{$\triangleright$\enspace\texttt{AsSubsemigroup({\mdseries\slshape D, C})\index{AsSubsemigroup@\texttt{AsSubsemigroup}}
\label{AsSubsemigroup}
}\hfill{\scriptsize (operation)}}\\


 Let \mbox{\texttt{\mdseries\slshape D}} be a domain and \mbox{\texttt{\mdseries\slshape C}} a collection. If \mbox{\texttt{\mdseries\slshape C}} is a subset of \mbox{\texttt{\mdseries\slshape D}} that forms a semigroup then \texttt{AsSubsemigroup} returns this semigroup, with parent \mbox{\texttt{\mdseries\slshape D}}. Otherwise \texttt{fail} is returned. }

 

\subsection{\textcolor{Chapter }{GeneratorsOfSemigroup}}
\logpage{[ 51, 1, 8 ]}\nobreak
\hyperdef{L}{X78147A247963F23B}{}
{\noindent\textcolor{FuncColor}{$\triangleright$\enspace\texttt{GeneratorsOfSemigroup({\mdseries\slshape S})\index{GeneratorsOfSemigroup@\texttt{GeneratorsOfSemigroup}}
\label{GeneratorsOfSemigroup}
}\hfill{\scriptsize (attribute)}}\\


 Semigroup generators of a semigroup \mbox{\texttt{\mdseries\slshape D}} are the same as magma generators, see{\nobreakspace}\texttt{GeneratorsOfMagma} (\ref{GeneratorsOfMagma}). 
\begin{Verbatim}[commandchars=!@|,fontsize=\small,frame=single,label=Example]
  !gapprompt@gap>| !gapinput@GeneratorsOfSemigroup(s);|
  [ Transformation( [ 2, 3, 4, 1 ] ), Transformation( [ 2, 2 ] ) ]
  !gapprompt@gap>| !gapinput@GeneratorsOfSemigroup(t);|
  [ Transformation( [ 2, 3, 4, 1 ] ) ]
\end{Verbatim}
 }

 

\subsection{\textcolor{Chapter }{IsGeneratorsOfSemigroup}}
\logpage{[ 51, 1, 9 ]}\nobreak
\hyperdef{L}{X79776D7C8399F2CF}{}
{\noindent\textcolor{FuncColor}{$\triangleright$\enspace\texttt{IsGeneratorsOfSemigroup({\mdseries\slshape C})\index{IsGeneratorsOfSemigroup@\texttt{IsGeneratorsOfSemigroup}}
\label{IsGeneratorsOfSemigroup}
}\hfill{\scriptsize (property)}}\\


 This property reflects whether the list or collection \mbox{\texttt{\mdseries\slshape C}} generates a semigroup. \texttt{IsAssociativeElementCollection} (\ref{IsAssociativeElementCollection}) implies {\nobreakspace}\texttt{IsGeneratorsOfSemigroup}, but is not used directly in semigroup code, because of conflicts with
matrices. 
\begin{Verbatim}[commandchars=!@|,fontsize=\small,frame=single,label=Example]
  !gapprompt@gap>| !gapinput@IsGeneratorsOfSemigroup([Transformation([2,3,1])]);|
  true
\end{Verbatim}
 }

 
\subsection{\textcolor{Chapter }{FreeSemigroup}}\logpage{[ 51, 1, 10 ]}
\hyperdef{L}{X7C72E4747BF642BB}{}
{
\noindent\textcolor{FuncColor}{$\triangleright$\enspace\texttt{FreeSemigroup({\mdseries\slshape [wfilt, ]rank[, name]})\index{FreeSemigroup@\texttt{FreeSemigroup}!for given rank}
\label{FreeSemigroup:for given rank}
}\hfill{\scriptsize (function)}}\\
\noindent\textcolor{FuncColor}{$\triangleright$\enspace\texttt{FreeSemigroup({\mdseries\slshape [wfilt, ]name1[, name2[, ...]]})\index{FreeSemigroup@\texttt{FreeSemigroup}!for various names}
\label{FreeSemigroup:for various names}
}\hfill{\scriptsize (function)}}\\
\noindent\textcolor{FuncColor}{$\triangleright$\enspace\texttt{FreeSemigroup({\mdseries\slshape [wfilt, ]names})\index{FreeSemigroup@\texttt{FreeSemigroup}!for a list of names}
\label{FreeSemigroup:for a list of names}
}\hfill{\scriptsize (function)}}\\
\noindent\textcolor{FuncColor}{$\triangleright$\enspace\texttt{FreeSemigroup({\mdseries\slshape [wfilt, ]infinity[, name][, init]})\index{FreeSemigroup@\texttt{FreeSemigroup}!for infinitely many generators}
\label{FreeSemigroup:for infinitely many generators}
}\hfill{\scriptsize (function)}}\\


 \texttt{FreeSemigroup} returns a free semigroup. The number of generators, and the labels given to
the generators, can be specified in several different ways. Warning: the
labels of generators are only an aid for printing, and do not necessarily
distinguish generators; see the examples at the end for more information. 
\begin{description}
\item[{ 1: For a given rank, and an optional generator name prefix }]  Called with a positive integer \mbox{\texttt{\mdseries\slshape rank}}, \texttt{FreeSemigroup} returns a free semigroup on \mbox{\texttt{\mdseries\slshape rank}} generators. The optional argument \mbox{\texttt{\mdseries\slshape name}} must be a string; its default value is \texttt{"s"}. 

 If \mbox{\texttt{\mdseries\slshape name}} is not given but the \texttt{generatorNames} option is, then this option is respected as described in
Section{\nobreakspace}\ref{Generator Names}. 

 Otherwise, the generators of the returned free semigroup are labelled \mbox{\texttt{\mdseries\slshape name}}\texttt{1}, ..., \mbox{\texttt{\mdseries\slshape name}}\texttt{k}, where \texttt{k} is the value of \mbox{\texttt{\mdseries\slshape rank}}. 

 
\item[{2: For given generator names}]  Called with various (at least one) nonempty strings, \texttt{FreeSemigroup} returns a free semigroup on as many generators as arguments, which are
labelled \mbox{\texttt{\mdseries\slshape name1}}, \mbox{\texttt{\mdseries\slshape name2}}, etc. 
\item[{3: For a given list of generator names}]  Called with a nonempty finite list \mbox{\texttt{\mdseries\slshape names}} of nonempty strings, \texttt{FreeSemigroup} returns a free semigroup on \texttt{Length(\mbox{\texttt{\mdseries\slshape names}})} generators, whose \texttt{i}\texttt{\symbol{45}}th generator is labelled \mbox{\texttt{\mdseries\slshape names}}\texttt{[i]}. 
\item[{ 4: For the rank \texttt{infinity}, an optional default generator name prefix, and an optional finite list of
generator names }]  Called in the fourth form, \texttt{FreeSemigroup} returns a free semigroup on infinitely many generators. The optional argument \mbox{\texttt{\mdseries\slshape name}} must be a string; its default value is \texttt{"s"}, and the optional argument \mbox{\texttt{\mdseries\slshape init}} must be a finite list of nonempty strings; its default value is an empty list.
The generators are initially labelled according to the list \mbox{\texttt{\mdseries\slshape init}}, followed by \mbox{\texttt{\mdseries\slshape name}}\texttt{i} for each \texttt{i} in the range from \texttt{Length(\mbox{\texttt{\mdseries\slshape init}})+1} to \texttt{infinity}; such a label is not allowed to appear in \mbox{\texttt{\mdseries\slshape init}}. 
\end{description}
 If the optional first argument \mbox{\texttt{\mdseries\slshape wfilt}} is given, then it must be either \texttt{IsSyllableWordsFamily}, \texttt{IsLetterWordsFamily}, \texttt{IsWLetterWordsFamily}, or \texttt{IsBLetterWordsFamily}. This filter specifies the representation used for the elements of the free
semigroup (see{\nobreakspace}\ref{Representations for Associative Words}). If no such filter is given, a letter representation is used. 

 For more on associative words see Chapter{\nobreakspace}\ref{Associative Words}. 

 
\begin{Verbatim}[commandchars=!@|,fontsize=\small,frame=single,label=Example]
  !gapprompt@gap>| !gapinput@f1 := FreeSemigroup( 3 );|
  <free semigroup on the generators [ s1, s2, s3 ]>
  !gapprompt@gap>| !gapinput@f2 := FreeSemigroup( 3 , "generator" );|
  <free semigroup on the generators
  [ generator1, generator2, generator3 ]>
  !gapprompt@gap>| !gapinput@f3 := FreeSemigroup( "gen1" , "gen2" );|
  <free semigroup on the generators [ gen1, gen2 ]>
  !gapprompt@gap>| !gapinput@f4 := FreeSemigroup( ["gen1" , "gen2"] );|
  <free semigroup on the generators [ gen1, gen2 ]>
  !gapprompt@gap>| !gapinput@FreeSemigroup( 3 : generatorNames := "boom" );|
  <free semigroup on the generators [ boom1, boom2, boom3 ]>
  !gapprompt@gap>| !gapinput@FreeSemigroup( 2 : generatorNames := [ "u", "v", "w" ] );|
  <free semigroup on the generators [ u, v ]>
  !gapprompt@gap>| !gapinput@FreeSemigroup( infinity ) ;|
  <free semigroup on the generators [ s1, s2, ... ]>
  !gapprompt@gap>| !gapinput@F := FreeSemigroup( infinity, "g", [ "a", "b" ]);|
  <free semigroup on the generators [ a, b, ... ]>
  !gapprompt@gap>| !gapinput@GeneratorsOfSemigroup( F ){[1..4]};|
  [ a, b, g3, g4 ]
  !gapprompt@gap>| !gapinput@GeneratorsOfSemigroup( FreeSemigroup( infinity, "gen" ) ){[1..3]};|
  [ gen1, gen2, gen3 ]
  !gapprompt@gap>| !gapinput@GeneratorsOfSemigroup( FreeSemigroup( infinity, [ "f" ] ) ){[1..3]};|
  [ f, s2, s3 ]
  !gapprompt@gap>| !gapinput@FreeSemigroup(IsSyllableWordsFamily, 5);|
  <free semigroup on the generators [ s1, s2, s3, s4, s5 ]>
\end{Verbatim}
 

 Each free object defines a unique alphabet (and a unique family of words). Its
generators are the letters of this alphabet, thus words of length one. 

 
\begin{Verbatim}[commandchars=!@|,fontsize=\small,frame=single,label=Example]
  !gapprompt@gap>| !gapinput@FreeSemigroup( 5 );|
  <free semigroup on the generators [ s1, s2, s3, s4, s5 ]>
  !gapprompt@gap>| !gapinput@FreeMonoid( "a", "b" );|
  <free monoid on the generators [ a, b ]>
  !gapprompt@gap>| !gapinput@FreeGroup( infinity );|
  <free group with infinity generators>
  !gapprompt@gap>| !gapinput@FreeSemigroup( "x", "y" );|
  <free semigroup on the generators [ x, y ]>
  !gapprompt@gap>| !gapinput@FreeMonoid( 7 );|
  <free monoid on the generators [ m1, m2, m3, m4, m5, m6, m7 ]>
\end{Verbatim}
 

 Remember that names are just a help for printing and do not necessarily
distinguish letters. It is possible to create arbitrarily weird situations by
choosing strange names for the letters. 

 
\begin{Verbatim}[commandchars=!@|,fontsize=\small,frame=single,label=Example]
  !gapprompt@gap>| !gapinput@f := FreeGroup( "x", "x" );|
  <free group on the generators [ x, x ]>
  !gapprompt@gap>| !gapinput@gens := GeneratorsOfGroup( f );|
  [ x, x ]
  !gapprompt@gap>| !gapinput@gens[1] = gens[2];|
  false
  !gapprompt@gap>| !gapinput@f:= FreeGroup( "f1*f2", "f2^-1", "Group( [ f1, f2 ] )" );|
  <free group on the generators [ f1*f2, f2^-1, Group( [ f1, f2 ] ) ]>
  !gapprompt@gap>| !gapinput@gens:= GeneratorsOfGroup( f );;|
  !gapprompt@gap>| !gapinput@gens[1] * gens[2];|
  f1*f2*f2^-1
  !gapprompt@gap>| !gapinput@gens[1] / gens[3];|
  f1*f2*Group( [ f1, f2 ] )^-1
  !gapprompt@gap>| !gapinput@gens[3] / gens[1] / gens[2];|
  Group( [ f1, f2 ] )*f1*f2^-1*f2^-1^-1
\end{Verbatim}
 }

 

\subsection{\textcolor{Chapter }{SemigroupByMultiplicationTable}}
\logpage{[ 51, 1, 11 ]}\nobreak
\hyperdef{L}{X7E67E13F7A01F8D3}{}
{\noindent\textcolor{FuncColor}{$\triangleright$\enspace\texttt{SemigroupByMultiplicationTable({\mdseries\slshape A})\index{SemigroupByMultiplicationTable@\texttt{SemigroupByMultiplicationTable}}
\label{SemigroupByMultiplicationTable}
}\hfill{\scriptsize (function)}}\\


 returns the semigroup whose multiplication is defined by the square matrix \mbox{\texttt{\mdseries\slshape A}} (see{\nobreakspace}\texttt{MagmaByMultiplicationTable} (\ref{MagmaByMultiplicationTable})) if such a semigroup exists. Otherwise \texttt{fail} is returned. 
\begin{Verbatim}[commandchars=!@|,fontsize=\small,frame=single,label=Example]
  !gapprompt@gap>| !gapinput@SemigroupByMultiplicationTable([[1,2,3],[2,3,1],[3,1,2]]);|
  <semigroup of size 3, with 3 generators>
  !gapprompt@gap>| !gapinput@SemigroupByMultiplicationTable([[1,2,3],[2,3,1],[3,2,1]]);|
  fail
\end{Verbatim}
 }

 }

  
\section{\textcolor{Chapter }{Monoids}}\label{sect:IsMonoid}
\logpage{[ 51, 2, 0 ]}
\hyperdef{L}{X872FE34A7814C0DC}{}
{
  

\subsection{\textcolor{Chapter }{IsMonoid}}
\logpage{[ 51, 2, 1 ]}\nobreak
\hyperdef{L}{X861C523483C6248C}{}
{\noindent\textcolor{FuncColor}{$\triangleright$\enspace\texttt{IsMonoid({\mdseries\slshape D})\index{IsMonoid@\texttt{IsMonoid}}
\label{IsMonoid}
}\hfill{\scriptsize (Synonym)}}\\


 A \emph{monoid} is a magma\texttt{\symbol{45}}with\texttt{\symbol{45}}one (see{\nobreakspace}\ref{Magmas}) with associative multiplication. }

 
\subsection{\textcolor{Chapter }{Monoid}}\logpage{[ 51, 2, 2 ]}
\hyperdef{L}{X7F95328B7C7E49EA}{}
{
\noindent\textcolor{FuncColor}{$\triangleright$\enspace\texttt{Monoid({\mdseries\slshape gen1, gen2, ...})\index{Monoid@\texttt{Monoid}!for various generators}
\label{Monoid:for various generators}
}\hfill{\scriptsize (function)}}\\
\noindent\textcolor{FuncColor}{$\triangleright$\enspace\texttt{Monoid({\mdseries\slshape gens[, id]})\index{Monoid@\texttt{Monoid}!for a list}
\label{Monoid:for a list}
}\hfill{\scriptsize (function)}}\\


 In the first form, \texttt{Monoid} returns the monoid generated by the arguments \mbox{\texttt{\mdseries\slshape gen1}}, \mbox{\texttt{\mdseries\slshape gen2}}, $\ldots$, that is, the closure of these elements under multiplication and taking the
0\texttt{\symbol{45}}th power. In the second form, \texttt{Monoid} returns the monoid generated by the elements in the homogeneous list \mbox{\texttt{\mdseries\slshape gens}}; a square matrix as only argument is treated as one generator, not as a list
of generators. In the second form, the identity element \mbox{\texttt{\mdseries\slshape id}} may be given as the second argument. 

 It is \emph{not} checked whether the underlying multiplication is associative, use \texttt{MagmaWithOne} (\ref{MagmaWithOne}) and \texttt{IsAssociative} (\ref{IsAssociative}) if you want to check whether a
magma\texttt{\symbol{45}}with\texttt{\symbol{45}}one is in fact a monoid. }

 

\subsection{\textcolor{Chapter }{Submonoid}}
\logpage{[ 51, 2, 3 ]}\nobreak
\hyperdef{L}{X8322D01E84912FD7}{}
{\noindent\textcolor{FuncColor}{$\triangleright$\enspace\texttt{Submonoid({\mdseries\slshape M, gens})\index{Submonoid@\texttt{Submonoid}}
\label{Submonoid}
}\hfill{\scriptsize (function)}}\\
\noindent\textcolor{FuncColor}{$\triangleright$\enspace\texttt{SubmonoidNC({\mdseries\slshape M, gens})\index{SubmonoidNC@\texttt{SubmonoidNC}}
\label{SubmonoidNC}
}\hfill{\scriptsize (function)}}\\


 are just synonyms of \texttt{SubmagmaWithOne} (\ref{SubmagmaWithOne}) and \texttt{SubmagmaWithOneNC} (\ref{SubmagmaWithOneNC}), respectively. }

 

\subsection{\textcolor{Chapter }{MonoidByGenerators}}
\logpage{[ 51, 2, 4 ]}\nobreak
\hyperdef{L}{X85129EE387CC4D28}{}
{\noindent\textcolor{FuncColor}{$\triangleright$\enspace\texttt{MonoidByGenerators({\mdseries\slshape gens[, one]})\index{MonoidByGenerators@\texttt{MonoidByGenerators}}
\label{MonoidByGenerators}
}\hfill{\scriptsize (operation)}}\\


 is the underlying operation of \texttt{Monoid} (\ref{Monoid:for a list}). }

 

\subsection{\textcolor{Chapter }{AsMonoid}}
\logpage{[ 51, 2, 5 ]}\nobreak
\hyperdef{L}{X7B22038F832B9C0F}{}
{\noindent\textcolor{FuncColor}{$\triangleright$\enspace\texttt{AsMonoid({\mdseries\slshape C})\index{AsMonoid@\texttt{AsMonoid}}
\label{AsMonoid}
}\hfill{\scriptsize (operation)}}\\


 If \mbox{\texttt{\mdseries\slshape C}} is a collection whose elements form a monoid, then \texttt{AsMonoid} returns this monoid. Otherwise \texttt{fail} is returned. }

 

\subsection{\textcolor{Chapter }{AsSubmonoid}}
\logpage{[ 51, 2, 6 ]}\nobreak
\hyperdef{L}{X7C9A12DE8287B2D3}{}
{\noindent\textcolor{FuncColor}{$\triangleright$\enspace\texttt{AsSubmonoid({\mdseries\slshape D, C})\index{AsSubmonoid@\texttt{AsSubmonoid}}
\label{AsSubmonoid}
}\hfill{\scriptsize (operation)}}\\


 Let \mbox{\texttt{\mdseries\slshape D}} be a domain and \mbox{\texttt{\mdseries\slshape C}} a collection. If \mbox{\texttt{\mdseries\slshape C}} is a subset of \mbox{\texttt{\mdseries\slshape D}} that forms a monoid then \texttt{AsSubmonoid} returns this monoid, with parent \mbox{\texttt{\mdseries\slshape D}}. Otherwise \texttt{fail} is returned. }

 

\subsection{\textcolor{Chapter }{GeneratorsOfMonoid}}
\logpage{[ 51, 2, 7 ]}\nobreak
\hyperdef{L}{X83CA2E7279C44718}{}
{\noindent\textcolor{FuncColor}{$\triangleright$\enspace\texttt{GeneratorsOfMonoid({\mdseries\slshape M})\index{GeneratorsOfMonoid@\texttt{GeneratorsOfMonoid}}
\label{GeneratorsOfMonoid}
}\hfill{\scriptsize (attribute)}}\\


 Monoid generators of a monoid \mbox{\texttt{\mdseries\slshape M}} are the same as magma\texttt{\symbol{45}}with\texttt{\symbol{45}}one
generators (see{\nobreakspace}\texttt{GeneratorsOfMagmaWithOne} (\ref{GeneratorsOfMagmaWithOne})). }

 

\subsection{\textcolor{Chapter }{TrivialSubmonoid}}
\logpage{[ 51, 2, 8 ]}\nobreak
\hyperdef{L}{X7EC77C0184587181}{}
{\noindent\textcolor{FuncColor}{$\triangleright$\enspace\texttt{TrivialSubmonoid({\mdseries\slshape M})\index{TrivialSubmonoid@\texttt{TrivialSubmonoid}}
\label{TrivialSubmonoid}
}\hfill{\scriptsize (attribute)}}\\


 is just a synonym for \texttt{TrivialSubmagmaWithOne} (\ref{TrivialSubmagmaWithOne}). }

 
\subsection{\textcolor{Chapter }{FreeMonoid}}\logpage{[ 51, 2, 9 ]}
\hyperdef{L}{X79FA3FA978CA2E43}{}
{
\noindent\textcolor{FuncColor}{$\triangleright$\enspace\texttt{FreeMonoid({\mdseries\slshape [wfilt, ]rank[, name]})\index{FreeMonoid@\texttt{FreeMonoid}!for given rank}
\label{FreeMonoid:for given rank}
}\hfill{\scriptsize (function)}}\\
\noindent\textcolor{FuncColor}{$\triangleright$\enspace\texttt{FreeMonoid({\mdseries\slshape [wfilt][,] [name1[, name2[, ...]]]})\index{FreeMonoid@\texttt{FreeMonoid}!for various names}
\label{FreeMonoid:for various names}
}\hfill{\scriptsize (function)}}\\
\noindent\textcolor{FuncColor}{$\triangleright$\enspace\texttt{FreeMonoid({\mdseries\slshape [wfilt, ]names})\index{FreeMonoid@\texttt{FreeMonoid}!for a list of names}
\label{FreeMonoid:for a list of names}
}\hfill{\scriptsize (function)}}\\
\noindent\textcolor{FuncColor}{$\triangleright$\enspace\texttt{FreeMonoid({\mdseries\slshape [wfilt, ]infinity[, name][, init]})\index{FreeMonoid@\texttt{FreeMonoid}!for infinitely many generators}
\label{FreeMonoid:for infinitely many generators}
}\hfill{\scriptsize (function)}}\\


 \texttt{FreeMonoid} returns a free monoid. The number of generators, and the labels given to the
generators, can be specified in several different ways. Warning: the labels of
generators are only an aid for printing, and do not necessarily distinguish
generators; see the examples at the end of \texttt{FreeSemigroup} (\ref{FreeSemigroup:for given rank}) for more information. 
\begin{description}
\item[{ 1: For a given rank, and an optional generator name prefix }]  Called with a nonnegative integer \mbox{\texttt{\mdseries\slshape rank}}, \texttt{FreeMonoid} returns a free monoid on \mbox{\texttt{\mdseries\slshape rank}} generators. The optional argument \mbox{\texttt{\mdseries\slshape name}} must be a string; its default value is \texttt{"m"}. 

 If \mbox{\texttt{\mdseries\slshape name}} is not given but the \texttt{generatorNames} option is, then this option is respected as described in
Section{\nobreakspace}\ref{Generator Names}. 

 Otherwise, the generators of the returned free monoid are labelled \mbox{\texttt{\mdseries\slshape name}}\texttt{1}, ..., \mbox{\texttt{\mdseries\slshape name}}\texttt{k}, where \texttt{k} is the value of \mbox{\texttt{\mdseries\slshape rank}}. 

 
\item[{2: For given generator names}]  Called with various nonempty strings, \texttt{FreeMonoid} returns a free monoid on as many generators as arguments, which are labelled \mbox{\texttt{\mdseries\slshape name1}}, \mbox{\texttt{\mdseries\slshape name2}}, etc. 
\item[{3: For a given list of generator names}]  Called with a finite list \mbox{\texttt{\mdseries\slshape names}} of nonempty strings, \texttt{FreeMonoid} returns a free monoid on \texttt{Length(\mbox{\texttt{\mdseries\slshape names}})} generators, whose \texttt{i}\texttt{\symbol{45}}th generator is labelled \mbox{\texttt{\mdseries\slshape names}}\texttt{[i]}. 
\item[{ 4: For the rank \texttt{infinity}, an optional default generator name prefix, and an optional finite list of
generator names }]  Called in the fourth form, \texttt{FreeMonoid} returns a free monoid on infinitely many generators. The optional argument \mbox{\texttt{\mdseries\slshape name}} must be a string; its default value is \texttt{"m"}, and the optional argument \mbox{\texttt{\mdseries\slshape init}} must be a finite list of nonempty strings; its default value is an empty list.
The generators are initially labelled according to the list \mbox{\texttt{\mdseries\slshape init}}, followed by \mbox{\texttt{\mdseries\slshape name}}\texttt{i} for each \texttt{i} in the range from \texttt{Length(\mbox{\texttt{\mdseries\slshape init}})+1} to \texttt{infinity}. 
\end{description}
 If the optional first argument \mbox{\texttt{\mdseries\slshape wfilt}} is given, then it must be either \texttt{IsSyllableWordsFamily}, \texttt{IsLetterWordsFamily}, \texttt{IsWLetterWordsFamily}, or \texttt{IsBLetterWordsFamily}. This filter specifies the representation used for the elements of the free
monoid (see{\nobreakspace}\ref{Representations for Associative Words}). If no such filter is given, a letter representation is used. 

 For more on associative words see Chapter{\nobreakspace}\ref{Associative Words}. 
\begin{Verbatim}[commandchars=!@|,fontsize=\small,frame=single,label=Example]
  !gapprompt@gap>| !gapinput@FreeMonoid(5);|
  <free monoid on the generators [ m1, m2, m3, m4, m5 ]>
  !gapprompt@gap>| !gapinput@FreeMonoid(4, "gen");|
  <free monoid on the generators [ gen1, gen2, gen3, gen4 ]>
  !gapprompt@gap>| !gapinput@FreeMonoid(3 : generatorNames := "turbo");|
  <free monoid on the generators [ turbo1, turbo2, turbo3 ]>
  !gapprompt@gap>| !gapinput@FreeMonoid(2 : generatorNames := ["u", "v", "w"]);|
  <free monoid on the generators [ u, v ]>
  !gapprompt@gap>| !gapinput@FreeMonoid(); FreeMonoid(0); FreeMonoid([]);|
  <free monoid of rank zero>
  <free monoid of rank zero>
  <free monoid of rank zero>
  !gapprompt@gap>| !gapinput@FreeMonoid("a", "b", "c");|
  <free monoid on the generators [ a, b, c ]>
  !gapprompt@gap>| !gapinput@FreeMonoid(["x", "y"]);|
  <free monoid on the generators [ x, y ]>
  !gapprompt@gap>| !gapinput@FreeMonoid(infinity);|
  <free monoid with infinity generators>
  !gapprompt@gap>| !gapinput@F := FreeMonoid(infinity, "g", ["a", "b"]);|
  <free monoid with infinity generators>
  !gapprompt@gap>| !gapinput@GeneratorsOfMonoid(F){[1..4]};|
  [ a, b, g3, g4 ]
  !gapprompt@gap>| !gapinput@GeneratorsOfMonoid(FreeMonoid(infinity, "gen")){[1..3]};|
  [ gen1, gen2, gen3 ]
  !gapprompt@gap>| !gapinput@GeneratorsOfMonoid(FreeMonoid(infinity, [ "f", "g" ])){[1..3]};|
  [ f, g, m3 ]
  !gapprompt@gap>| !gapinput@FreeMonoid(IsSyllableWordsFamily, 50);|
  <free monoid with 50 generators>
\end{Verbatim}
 }

 

\subsection{\textcolor{Chapter }{MonoidByMultiplicationTable}}
\logpage{[ 51, 2, 10 ]}\nobreak
\hyperdef{L}{X7BFE938E857CA27D}{}
{\noindent\textcolor{FuncColor}{$\triangleright$\enspace\texttt{MonoidByMultiplicationTable({\mdseries\slshape A})\index{MonoidByMultiplicationTable@\texttt{MonoidByMultiplicationTable}}
\label{MonoidByMultiplicationTable}
}\hfill{\scriptsize (function)}}\\


 returns the monoid whose multiplication is defined by the square matrix \mbox{\texttt{\mdseries\slshape A}} (see{\nobreakspace}\texttt{MagmaByMultiplicationTable} (\ref{MagmaByMultiplicationTable})) if such a monoid exists. Otherwise \texttt{fail} is returned. 
\begin{Verbatim}[commandchars=!@|,fontsize=\small,frame=single,label=Example]
  !gapprompt@gap>| !gapinput@MonoidByMultiplicationTable([[1,2,3],[2,3,1],[3,1,2]]);|
  <monoid of size 3, with 3 generators>
  !gapprompt@gap>| !gapinput@MonoidByMultiplicationTable([[1,2,3],[2,3,1],[1,3,2]]);|
  fail
\end{Verbatim}
 }

 }

  
\section{\textcolor{Chapter }{Inverse semigroups and monoids}}\label{sect:IsInverseSemigroup}
\logpage{[ 51, 3, 0 ]}
\hyperdef{L}{X840847B6810BD0E1}{}
{
  

\subsection{\textcolor{Chapter }{InverseSemigroup}}
\logpage{[ 51, 3, 1 ]}\nobreak
\hyperdef{L}{X78B13FED7AFB4326}{}
{\noindent\textcolor{FuncColor}{$\triangleright$\enspace\texttt{InverseSemigroup({\mdseries\slshape obj1, obj2, ...})\index{InverseSemigroup@\texttt{InverseSemigroup}}
\label{InverseSemigroup}
}\hfill{\scriptsize (function)}}\\
\textbf{\indent Returns:\ }
An inverse semigroup.



 If \mbox{\texttt{\mdseries\slshape obj1}}, \mbox{\texttt{\mdseries\slshape obj2}}, ... are (any combination) of associative elements with unique semigroup
inverses, semigroups of such elements, or collections of such elements, then \texttt{InverseSemigroup} returns the inverse semigroup generated by the union of \mbox{\texttt{\mdseries\slshape obj1}}, \mbox{\texttt{\mdseries\slshape obj2}}, .... This equals the semigroup generated by the union of \mbox{\texttt{\mdseries\slshape obj1}}, \mbox{\texttt{\mdseries\slshape obj2}}, ... and their inverses.

 For example if \texttt{S} and \texttt{T} are inverse semigroups, then \texttt{InverseSemigroup(S, f, Idempotents(T));} is the inverse semigroup generated by \texttt{Union(GeneratorsOfInverseSemigroup(S), [f], Idempotents(T)));}.

 As present, the only associative elements with unique semigroup inverses,
which do not always generate a group, are partial permutations; see Chapter \ref{Partial permutations}. 
\begin{Verbatim}[commandchars=!@|,fontsize=\small,frame=single,label=Example]
  !gapprompt@gap>| !gapinput@S := InverseSemigroup(|
  !gapprompt@>| !gapinput@PartialPerm( [ 1, 2, 3, 6, 8, 10 ], [ 2, 6, 7, 9, 1, 5 ] ) );;|
  !gapprompt@gap>| !gapinput@f := PartialPerm( [ 1, 2, 3, 4, 5, 8, 10 ],|
  !gapprompt@>| !gapinput@[ 7, 1, 4, 3, 2, 6, 5 ] );;|
  !gapprompt@gap>| !gapinput@S := InverseSemigroup(S, f, Idempotents(SymmetricInverseSemigroup(5)));|
  <inverse partial perm semigroup of rank 10 with 34 generators>
  !gapprompt@gap>| !gapinput@Size(S);|
  1233
\end{Verbatim}
 }

 

\subsection{\textcolor{Chapter }{InverseMonoid}}
\logpage{[ 51, 3, 2 ]}\nobreak
\hyperdef{L}{X80D9B9A98736051B}{}
{\noindent\textcolor{FuncColor}{$\triangleright$\enspace\texttt{InverseMonoid({\mdseries\slshape obj1, obj2, ...})\index{InverseMonoid@\texttt{InverseMonoid}}
\label{InverseMonoid}
}\hfill{\scriptsize (function)}}\\
\textbf{\indent Returns:\ }
An inverse monoid.



 If \mbox{\texttt{\mdseries\slshape obj1}}, \mbox{\texttt{\mdseries\slshape obj2}}, ... are (any combination) of associative elements with unique semigroup
inverses, semigroups of such elements, or collections of such elements, then \texttt{InverseMonoid} returns the inverse monoid generated by the union of \mbox{\texttt{\mdseries\slshape obj1}}, \mbox{\texttt{\mdseries\slshape obj2}}, .... This equals the monoid generated by the union of \mbox{\texttt{\mdseries\slshape obj1}}, \mbox{\texttt{\mdseries\slshape obj2}}, ... and their inverses.

 As present, the only associative elements with unique semigroup inverses are
partial permutations; see Chapter \ref{Partial permutations}.

 For example if \texttt{S} and \texttt{T} are inverse monoids, then \texttt{InverseMonoid(S, f, Idempotents(T));} is the inverse monoid generated by \texttt{Union(GeneratorsOfInverseMonoid(S), [f], Idempotents(T)));}. 
\begin{Verbatim}[commandchars=!@|,fontsize=\small,frame=single,label=Example]
  !gapprompt@gap>| !gapinput@S := InverseMonoid(|
  !gapprompt@>| !gapinput@PartialPerm( [ 1, 2, 3, 6, 8, 10 ], [ 2, 6, 7, 9, 1, 5 ] ) );;|
  !gapprompt@gap>| !gapinput@f := PartialPerm( [ 1, 2, 3, 4, 5, 8, 10 ],|
  !gapprompt@>| !gapinput@[ 7, 1, 4, 3, 2, 6, 5 ] );;|
  !gapprompt@gap>| !gapinput@S := InverseMonoid(S, f, Idempotents(SymmetricInverseSemigroup(5)));|
  <inverse partial perm monoid of rank 10 with 35 generators>
  !gapprompt@gap>| !gapinput@Size(S);|
  1243
\end{Verbatim}
 }

 

\subsection{\textcolor{Chapter }{GeneratorsOfInverseSemigroup}}
\logpage{[ 51, 3, 3 ]}\nobreak
\hyperdef{L}{X87C373597F787250}{}
{\noindent\textcolor{FuncColor}{$\triangleright$\enspace\texttt{GeneratorsOfInverseSemigroup({\mdseries\slshape S})\index{GeneratorsOfInverseSemigroup@\texttt{GeneratorsOfInverseSemigroup}}
\label{GeneratorsOfInverseSemigroup}
}\hfill{\scriptsize (attribute)}}\\
\textbf{\indent Returns:\ }
The generators of an inverse semigroup.



 If \mbox{\texttt{\mdseries\slshape S}} is an inverse semigroup, then \texttt{GeneratorsOfInverseSemigroup} returns the generators used to define \mbox{\texttt{\mdseries\slshape S}}, i.e. an inverse semigroup generating set for \mbox{\texttt{\mdseries\slshape S}}. 

 The value of \texttt{GeneratorsOfSemigroup(\mbox{\texttt{\mdseries\slshape S}})}, for an inverse semigroup \mbox{\texttt{\mdseries\slshape S}}, is the union of inverse semigroup generator and their inverses. So, \mbox{\texttt{\mdseries\slshape S}} is the semigroup, as opposed to inverse semigroup, generated by the elements
of \texttt{GeneratorsOfInverseSemigroup(\mbox{\texttt{\mdseries\slshape S}})} and their inverses. 

 If \mbox{\texttt{\mdseries\slshape S}} is an inverse monoid, then \texttt{GeneratorsOfInverseSemigroup} returns the generators used to define \mbox{\texttt{\mdseries\slshape S}}, as described above, and the identity of \mbox{\texttt{\mdseries\slshape S}}. 
\begin{Verbatim}[commandchars=!@|,fontsize=\small,frame=single,label=Example]
  !gapprompt@gap>| !gapinput@S:=InverseMonoid(|
  !gapprompt@>| !gapinput@ PartialPerm( [ 1, 2 ], [ 1, 4 ] ),|
  !gapprompt@>| !gapinput@ PartialPerm( [ 1, 2, 4 ], [ 3, 4, 1 ] ) );;|
  !gapprompt@gap>| !gapinput@GeneratorsOfSemigroup(S);|
  [ <identity partial perm on [ 1, 2, 3, 4 ]>, [2,4](1), [2,4,1,3],
    [4,2](1), [3,1,4,2] ]
  !gapprompt@gap>| !gapinput@GeneratorsOfInverseSemigroup(S);|
  [ [2,4](1), [2,4,1,3], <identity partial perm on [ 1, 2, 3, 4 ]> ]
  !gapprompt@gap>| !gapinput@GeneratorsOfMonoid(S);|
  [ [2,4](1), [2,4,1,3], [4,2](1), [3,1,4,2] ]
\end{Verbatim}
 }

 

\subsection{\textcolor{Chapter }{GeneratorsOfInverseMonoid}}
\logpage{[ 51, 3, 4 ]}\nobreak
\hyperdef{L}{X7A3B262C85B6D475}{}
{\noindent\textcolor{FuncColor}{$\triangleright$\enspace\texttt{GeneratorsOfInverseMonoid({\mdseries\slshape S})\index{GeneratorsOfInverseMonoid@\texttt{GeneratorsOfInverseMonoid}}
\label{GeneratorsOfInverseMonoid}
}\hfill{\scriptsize (attribute)}}\\
\textbf{\indent Returns:\ }
The generators of an inverse monoid.



 If \mbox{\texttt{\mdseries\slshape S}} is an inverse monoid, then \texttt{GeneratorsOfInverseMonoid} returns the generators used to define \mbox{\texttt{\mdseries\slshape S}}, i.e. an inverse monoid generating set for \mbox{\texttt{\mdseries\slshape S}}. 

 There are four different possible generating sets which define an inverse
monoid. More precisely, an inverse monoid can be generated as an inverse
monoid, inverse semigroup, monoid, or semigroup. The different generating sets
in each case can be obtained using \texttt{GeneratorsOfInverseMonoid}, \texttt{GeneratorsOfInverseSemigroup} (\ref{GeneratorsOfInverseSemigroup}), \texttt{GeneratorsOfMonoid} (\ref{GeneratorsOfMonoid}), and \texttt{GeneratorsOfSemigroup} (\ref{GeneratorsOfSemigroup}), respectively. 
\begin{Verbatim}[commandchars=!@|,fontsize=\small,frame=single,label=Example]
  !gapprompt@gap>| !gapinput@S:=InverseMonoid(|
  !gapprompt@>| !gapinput@ PartialPerm( [ 1, 2 ], [ 1, 4 ] ),|
  !gapprompt@>| !gapinput@ PartialPerm( [ 1, 2, 4 ], [ 3, 4, 1 ] ) );;|
  !gapprompt@gap>| !gapinput@GeneratorsOfInverseMonoid(S);|
  [ [2,4](1), [2,4,1,3] ]
\end{Verbatim}
 }

 

\subsection{\textcolor{Chapter }{IsInverseSubsemigroup}}
\logpage{[ 51, 3, 5 ]}\nobreak
\hyperdef{L}{X7C4C6EE681E7A57E}{}
{\noindent\textcolor{FuncColor}{$\triangleright$\enspace\texttt{IsInverseSubsemigroup({\mdseries\slshape S, T})\index{IsInverseSubsemigroup@\texttt{IsInverseSubsemigroup}}
\label{IsInverseSubsemigroup}
}\hfill{\scriptsize (operation)}}\\
\textbf{\indent Returns:\ }
\texttt{true} or \texttt{false}.



 If the semigroup \mbox{\texttt{\mdseries\slshape T}} is an inverse subsemigroup of the semigroup \mbox{\texttt{\mdseries\slshape S}}, then this operation returns \texttt{true}. 
\begin{Verbatim}[commandchars=!@|,fontsize=\small,frame=single,label=Example]
  !gapprompt@gap>| !gapinput@T:=InverseSemigroup(RandomPartialPerm(4));;|
  !gapprompt@gap>| !gapinput@IsInverseSubsemigroup(SymmetricInverseSemigroup(4), T);|
  true
  !gapprompt@gap>| !gapinput@T:=Semigroup(Transformation( [ 1, 2, 4, 5, 6, 3, 7, 8 ] ),|
  !gapprompt@>| !gapinput@Transformation( [ 3, 3, 4, 5, 6, 2, 7, 8 ] ),|
  !gapprompt@>| !gapinput@Transformation([ 1, 2, 5, 3, 6, 8, 4, 4 ] ));;|
  !gapprompt@gap>| !gapinput@IsInverseSubsemigroup(FullTransformationSemigroup(8), T);|
  true
\end{Verbatim}
 }

 }

  
\section{\textcolor{Chapter }{Properties of Semigroups}}\label{Properties of Semigroups}
\logpage{[ 51, 4, 0 ]}
\hyperdef{L}{X78274024827F306D}{}
{
  The following functions determine information about semigroups. 

\subsection{\textcolor{Chapter }{IsRegularSemigroup}}
\logpage{[ 51, 4, 1 ]}\nobreak
\hyperdef{L}{X7C4663827C5ACEF1}{}
{\noindent\textcolor{FuncColor}{$\triangleright$\enspace\texttt{IsRegularSemigroup({\mdseries\slshape S})\index{IsRegularSemigroup@\texttt{IsRegularSemigroup}}
\label{IsRegularSemigroup}
}\hfill{\scriptsize (property)}}\\


 returns \texttt{true} if \mbox{\texttt{\mdseries\slshape S}} is regular, i.e., if every $\mathcal{D}$\texttt{\symbol{45}}class of \mbox{\texttt{\mdseries\slshape S}} is regular. }

 

\subsection{\textcolor{Chapter }{IsRegularSemigroupElement}}
\logpage{[ 51, 4, 2 ]}\nobreak
\hyperdef{L}{X87532A76854347E0}{}
{\noindent\textcolor{FuncColor}{$\triangleright$\enspace\texttt{IsRegularSemigroupElement({\mdseries\slshape S, x})\index{IsRegularSemigroupElement@\texttt{IsRegularSemigroupElement}}
\label{IsRegularSemigroupElement}
}\hfill{\scriptsize (operation)}}\\


 returns \texttt{true} if \mbox{\texttt{\mdseries\slshape x}} has a general inverse in \mbox{\texttt{\mdseries\slshape S}}, i.e., there is an element $y \in \mbox{\texttt{\mdseries\slshape S}}$ such that $\mbox{\texttt{\mdseries\slshape x}} y \mbox{\texttt{\mdseries\slshape x}} = \mbox{\texttt{\mdseries\slshape x}}$ and $y \mbox{\texttt{\mdseries\slshape x}} y = y$. }

 

\subsection{\textcolor{Chapter }{InversesOfSemigroupElement}}
\logpage{[ 51, 4, 3 ]}\nobreak
\hyperdef{L}{X7AFDE0F17AE516C5}{}
{\noindent\textcolor{FuncColor}{$\triangleright$\enspace\texttt{InversesOfSemigroupElement({\mdseries\slshape S, x})\index{InversesOfSemigroupElement@\texttt{InversesOfSemigroupElement}}
\label{InversesOfSemigroupElement}
}\hfill{\scriptsize (operation)}}\\
\textbf{\indent Returns:\ }
A list of the inverses of an element of a semigroup.



 \texttt{InversesOfSemigroupElement} returns a list of the inverses of the element \mbox{\texttt{\mdseries\slshape x}} in the semigroup \mbox{\texttt{\mdseries\slshape S}}.

 An element \mbox{\texttt{\mdseries\slshape y}} in \mbox{\texttt{\mdseries\slshape S}} is an \emph{inverse} of \mbox{\texttt{\mdseries\slshape x}} if \texttt{\mbox{\texttt{\mdseries\slshape x}}*y*\mbox{\texttt{\mdseries\slshape x}}=\mbox{\texttt{\mdseries\slshape x}}} and \texttt{y*\mbox{\texttt{\mdseries\slshape x}}*y=y}. The element \mbox{\texttt{\mdseries\slshape x}} has an inverse if and only if \mbox{\texttt{\mdseries\slshape x}} is a regular element of \mbox{\texttt{\mdseries\slshape S}}. 
\begin{Verbatim}[commandchars=!@|,fontsize=\small,frame=single,label=Example]
  !gapprompt@gap>| !gapinput@S := Semigroup([|
  !gapprompt@>| !gapinput@ Transformation([3, 1, 4, 2, 5, 2, 1, 6, 1]),|
  !gapprompt@>| !gapinput@ Transformation([5, 7, 8, 8, 7, 5, 9, 1, 9]),|
  !gapprompt@>| !gapinput@ Transformation([7, 6, 2, 8, 4, 7, 5, 8, 3])]);|
  <transformation semigroup of degree 9 with 3 generators>
  !gapprompt@gap>| !gapinput@x := Transformation([3, 1, 4, 2, 5, 2, 1, 6, 1]);;|
  !gapprompt@gap>| !gapinput@InversesOfSemigroupElement(S, x);|
  [  ]
  !gapprompt@gap>| !gapinput@IsRegularSemigroupElement(S, x);|
  false
  !gapprompt@gap>| !gapinput@x := Transformation([1, 9, 7, 5, 5, 1, 9, 5, 1]);;|
  !gapprompt@gap>| !gapinput@Set(InversesOfSemigroupElement(S, x));|
  [ Transformation( [ 1, 2, 3, 5, 5, 1, 3, 5, 2 ] ),
    Transformation( [ 1, 5, 1, 1, 5, 1, 3, 1, 2 ] ),
    Transformation( [ 1, 5, 1, 2, 5, 1, 3, 2, 2 ] ) ]
  !gapprompt@gap>| !gapinput@IsRegularSemigroupElement(S, x);|
  true
  !gapprompt@gap>| !gapinput@S := ReesZeroMatrixSemigroup(Group((1,2,3)),|
  !gapprompt@>| !gapinput@ [[(), ()], [(), 0], [(), (1,2,3)]]);;|
  !gapprompt@gap>| !gapinput@x := ReesZeroMatrixSemigroupElement(S, 2, (1,2,3), 3);;|
  !gapprompt@gap>| !gapinput@InversesOfSemigroupElement(S, x);|
  [ (1,(1,2,3),3), (1,(1,3,2),1), (2,(),3), (2,(1,2,3),1) ]
\end{Verbatim}
 }

 

\subsection{\textcolor{Chapter }{IsSimpleSemigroup}}
\logpage{[ 51, 4, 4 ]}\nobreak
\hyperdef{L}{X836F4692839F4874}{}
{\noindent\textcolor{FuncColor}{$\triangleright$\enspace\texttt{IsSimpleSemigroup({\mdseries\slshape S})\index{IsSimpleSemigroup@\texttt{IsSimpleSemigroup}}
\label{IsSimpleSemigroup}
}\hfill{\scriptsize (property)}}\\


 is \texttt{true} if and only if the semigroup \mbox{\texttt{\mdseries\slshape S}} has no proper ideals. }

 

\subsection{\textcolor{Chapter }{IsZeroSimpleSemigroup}}
\logpage{[ 51, 4, 5 ]}\nobreak
\hyperdef{L}{X8193A60F839C064E}{}
{\noindent\textcolor{FuncColor}{$\triangleright$\enspace\texttt{IsZeroSimpleSemigroup({\mdseries\slshape S})\index{IsZeroSimpleSemigroup@\texttt{IsZeroSimpleSemigroup}}
\label{IsZeroSimpleSemigroup}
}\hfill{\scriptsize (property)}}\\


 is \texttt{true} if and only if the semigroup has no proper ideals except for 0, where \mbox{\texttt{\mdseries\slshape S}} is a semigroup with zero. If the semigroup does not find its zero, then a
break\texttt{\symbol{45}}loop is entered. }

 

\subsection{\textcolor{Chapter }{IsZeroGroup}}
\logpage{[ 51, 4, 6 ]}\nobreak
\hyperdef{L}{X85F7E5CD86F0643B}{}
{\noindent\textcolor{FuncColor}{$\triangleright$\enspace\texttt{IsZeroGroup({\mdseries\slshape S})\index{IsZeroGroup@\texttt{IsZeroGroup}}
\label{IsZeroGroup}
}\hfill{\scriptsize (property)}}\\


 is \texttt{true} if and only if the semigroup \mbox{\texttt{\mdseries\slshape S}} is a group with zero adjoined. }

 

\subsection{\textcolor{Chapter }{IsReesCongruenceSemigroup}}
\logpage{[ 51, 4, 7 ]}\nobreak
\hyperdef{L}{X7FFEC81F7F2C4EAA}{}
{\noindent\textcolor{FuncColor}{$\triangleright$\enspace\texttt{IsReesCongruenceSemigroup({\mdseries\slshape S})\index{IsReesCongruenceSemigroup@\texttt{IsReesCongruenceSemigroup}}
\label{IsReesCongruenceSemigroup}
}\hfill{\scriptsize (property)}}\\


 returns \texttt{true} if \mbox{\texttt{\mdseries\slshape S}} is a Rees Congruence semigroup, that is, if all congruences of \mbox{\texttt{\mdseries\slshape S}} are Rees Congruences. }

 

\subsection{\textcolor{Chapter }{IsInverseSemigroup}}
\logpage{[ 51, 4, 8 ]}\nobreak
\hyperdef{L}{X83F1529479D56665}{}
{\noindent\textcolor{FuncColor}{$\triangleright$\enspace\texttt{IsInverseSemigroup({\mdseries\slshape S})\index{IsInverseSemigroup@\texttt{IsInverseSemigroup}}
\label{IsInverseSemigroup}
}\hfill{\scriptsize (property)}}\\
\noindent\textcolor{FuncColor}{$\triangleright$\enspace\texttt{IsInverseMonoid({\mdseries\slshape S})\index{IsInverseMonoid@\texttt{IsInverseMonoid}}
\label{IsInverseMonoid}
}\hfill{\scriptsize (Category)}}\\
\textbf{\indent Returns:\ }
\texttt{true} or \texttt{false}.



 A semigroup \mbox{\texttt{\mdseries\slshape S}} is an \emph{inverse semigroup} if every element \texttt{x} in \mbox{\texttt{\mdseries\slshape S}} has a unique semigroup inverse, that is, a unique element \texttt{y} in \mbox{\texttt{\mdseries\slshape S}} such that \texttt{x*y*x=x} and \texttt{y*x*y=y}.

 A monoid that happens to be an inverse semigroup is called an \emph{inverse monoid}; see \texttt{IsMonoid} (\ref{IsMonoid}). 
\begin{Verbatim}[commandchars=!@|,fontsize=\small,frame=single,label=Example]
  !gapprompt@gap>| !gapinput@S := Semigroup([|
  !gapprompt@>| !gapinput@ Transformation([1, 2, 4, 5, 6, 3, 7, 8]),|
  !gapprompt@>| !gapinput@ Transformation([3, 3, 4, 5, 6, 2, 7, 8]),|
  !gapprompt@>| !gapinput@ Transformation([1, 2, 5, 3, 6, 8, 4, 4])]);;|
  !gapprompt@gap>| !gapinput@IsInverseSemigroup(S);|
  true
\end{Verbatim}
 }

 }

  
\section{\textcolor{Chapter }{Ideals of semigroups}}\label{Ideals of semigroups}
\logpage{[ 51, 5, 0 ]}
\hyperdef{L}{X7BB32D508183C0F1}{}
{
  Ideals of semigroups are the same as ideals of the semigroup when considered
as a magma. For documentation on ideals for magmas, see \texttt{Magma} (\ref{Magma}). 

\subsection{\textcolor{Chapter }{SemigroupIdealByGenerators}}
\logpage{[ 51, 5, 1 ]}\nobreak
\hyperdef{L}{X7D5CEE4D7D4318ED}{}
{\noindent\textcolor{FuncColor}{$\triangleright$\enspace\texttt{SemigroupIdealByGenerators({\mdseries\slshape S, gens})\index{SemigroupIdealByGenerators@\texttt{SemigroupIdealByGenerators}}
\label{SemigroupIdealByGenerators}
}\hfill{\scriptsize (operation)}}\\


 \mbox{\texttt{\mdseries\slshape S}} is a semigroup, \mbox{\texttt{\mdseries\slshape gens}} is a list of elements of \mbox{\texttt{\mdseries\slshape S}}. Returns the two\texttt{\symbol{45}}sided ideal of \mbox{\texttt{\mdseries\slshape S}} generated by \mbox{\texttt{\mdseries\slshape gens}}. }

 

\subsection{\textcolor{Chapter }{ReesCongruenceOfSemigroupIdeal}}
\logpage{[ 51, 5, 2 ]}\nobreak
\hyperdef{L}{X7F01FFB18125DED5}{}
{\noindent\textcolor{FuncColor}{$\triangleright$\enspace\texttt{ReesCongruenceOfSemigroupIdeal({\mdseries\slshape I})\index{ReesCongruenceOfSemigroupIdeal@\texttt{ReesCongruenceOfSemigroupIdeal}}
\label{ReesCongruenceOfSemigroupIdeal}
}\hfill{\scriptsize (attribute)}}\\


 A two sided ideal \mbox{\texttt{\mdseries\slshape I}} of a semigroup \mbox{\texttt{\mdseries\slshape S}} defines a congruence on \mbox{\texttt{\mdseries\slshape S}} given by $\Delta \cup I \times I$. }

 

\subsection{\textcolor{Chapter }{IsLeftSemigroupIdeal}}
\logpage{[ 51, 5, 3 ]}\nobreak
\hyperdef{L}{X7A3FF85984345540}{}
{\noindent\textcolor{FuncColor}{$\triangleright$\enspace\texttt{IsLeftSemigroupIdeal({\mdseries\slshape I})\index{IsLeftSemigroupIdeal@\texttt{IsLeftSemigroupIdeal}}
\label{IsLeftSemigroupIdeal}
}\hfill{\scriptsize (property)}}\\
\noindent\textcolor{FuncColor}{$\triangleright$\enspace\texttt{IsRightSemigroupIdeal({\mdseries\slshape I})\index{IsRightSemigroupIdeal@\texttt{IsRightSemigroupIdeal}}
\label{IsRightSemigroupIdeal}
}\hfill{\scriptsize (property)}}\\
\noindent\textcolor{FuncColor}{$\triangleright$\enspace\texttt{IsSemigroupIdeal({\mdseries\slshape I})\index{IsSemigroupIdeal@\texttt{IsSemigroupIdeal}}
\label{IsSemigroupIdeal}
}\hfill{\scriptsize (property)}}\\


 Categories of semigroup ideals. }

 }

  
\section{\textcolor{Chapter }{Congruences on semigroups}}\label{Congruences on semigroups}
\logpage{[ 51, 6, 0 ]}
\hyperdef{L}{X7914691E7DFFE27A}{}
{
  An equivalence or a congruence on a semigroup is the equivalence or congruence
on the semigroup considered as a magma. So, to deal with equivalences and
congruences on semigroups, magma functions are used. For documentation on
equivalences and congruences on magmas, see \texttt{Magma} (\ref{Magma}). 

\subsection{\textcolor{Chapter }{IsSemigroupCongruence}}
\logpage{[ 51, 6, 1 ]}\nobreak
\hyperdef{L}{X78E34B737F0E009F}{}
{\noindent\textcolor{FuncColor}{$\triangleright$\enspace\texttt{IsSemigroupCongruence({\mdseries\slshape c})\index{IsSemigroupCongruence@\texttt{IsSemigroupCongruence}}
\label{IsSemigroupCongruence}
}\hfill{\scriptsize (property)}}\\


 a magma congruence \mbox{\texttt{\mdseries\slshape c}} on a semigroup. }

 

\subsection{\textcolor{Chapter }{IsReesCongruence}}
\logpage{[ 51, 6, 2 ]}\nobreak
\hyperdef{L}{X822DB78579BCB7B5}{}
{\noindent\textcolor{FuncColor}{$\triangleright$\enspace\texttt{IsReesCongruence({\mdseries\slshape c})\index{IsReesCongruence@\texttt{IsReesCongruence}}
\label{IsReesCongruence}
}\hfill{\scriptsize (property)}}\\


 returns \texttt{true} if and only if the congruence \mbox{\texttt{\mdseries\slshape c}} has at most one nonsingleton congruence class. }

 }

  
\section{\textcolor{Chapter }{Quotients}}\label{Quotients}
\logpage{[ 51, 7, 0 ]}
\hyperdef{L}{X87CE9EAB7EE3A128}{}
{
  Given a semigroup and a congruence on the semigroup, one can construct a new
semigroup: the quotient semigroup. The following functions deal with quotient
semigroups in \textsf{GAP}. For a semigroup $S$, elements of a quotient semigroup are equivalence classes of elements of the \texttt{QuotientSemigroupPreimage} (\ref{QuotientSemigroupPreimage}) value under the congruence given by the value of \texttt{QuotientSemigroupCongruence} (\ref{QuotientSemigroupCongruence}). 

 It is probably most useful for calculating the elements of the equivalence
classes by using \texttt{Elements} (\ref{Elements}) or by looking at the images of elements of \texttt{QuotientSemigroupPreimage} (\ref{QuotientSemigroupPreimage}) under the map returned by \texttt{QuotientSemigroupHomomorphism} (\ref{QuotientSemigroupHomomorphism}), which maps the \texttt{QuotientSemigroupPreimage} (\ref{QuotientSemigroupPreimage}) value to \mbox{\texttt{\mdseries\slshape S}}. 

 For intensive computations in a quotient semigroup, it is probably worthwhile
finding another representation as the equality test could involve enumeration
of the elements of the congruence classes being compared. 

\subsection{\textcolor{Chapter }{IsQuotientSemigroup}}
\logpage{[ 51, 7, 1 ]}\nobreak
\hyperdef{L}{X80EF3E6F842BE64E}{}
{\noindent\textcolor{FuncColor}{$\triangleright$\enspace\texttt{IsQuotientSemigroup({\mdseries\slshape S})\index{IsQuotientSemigroup@\texttt{IsQuotientSemigroup}}
\label{IsQuotientSemigroup}
}\hfill{\scriptsize (Category)}}\\


 is the category of semigroups constructed from another semigroup and a
congruence on it. }

 

\subsection{\textcolor{Chapter }{HomomorphismQuotientSemigroup}}
\logpage{[ 51, 7, 2 ]}\nobreak
\hyperdef{L}{X7CAD3D1687956F7F}{}
{\noindent\textcolor{FuncColor}{$\triangleright$\enspace\texttt{HomomorphismQuotientSemigroup({\mdseries\slshape cong})\index{HomomorphismQuotientSemigroup@\texttt{HomomorphismQuotientSemigroup}}
\label{HomomorphismQuotientSemigroup}
}\hfill{\scriptsize (function)}}\\


 for a congruence \mbox{\texttt{\mdseries\slshape cong}} and a semigroup \mbox{\texttt{\mdseries\slshape S}}. Returns the homomorphism from \mbox{\texttt{\mdseries\slshape S}} to the quotient of \mbox{\texttt{\mdseries\slshape S}} by \mbox{\texttt{\mdseries\slshape cong}}. }

 

\subsection{\textcolor{Chapter }{QuotientSemigroupPreimage}}
\logpage{[ 51, 7, 3 ]}\nobreak
\hyperdef{L}{X87120C46808F7289}{}
{\noindent\textcolor{FuncColor}{$\triangleright$\enspace\texttt{QuotientSemigroupPreimage({\mdseries\slshape S})\index{QuotientSemigroupPreimage@\texttt{QuotientSemigroupPreimage}}
\label{QuotientSemigroupPreimage}
}\hfill{\scriptsize (attribute)}}\\
\noindent\textcolor{FuncColor}{$\triangleright$\enspace\texttt{QuotientSemigroupCongruence({\mdseries\slshape S})\index{QuotientSemigroupCongruence@\texttt{QuotientSemigroupCongruence}}
\label{QuotientSemigroupCongruence}
}\hfill{\scriptsize (attribute)}}\\
\noindent\textcolor{FuncColor}{$\triangleright$\enspace\texttt{QuotientSemigroupHomomorphism({\mdseries\slshape S})\index{QuotientSemigroupHomomorphism@\texttt{QuotientSemigroupHomomorphism}}
\label{QuotientSemigroupHomomorphism}
}\hfill{\scriptsize (attribute)}}\\


 for a quotient semigroup \mbox{\texttt{\mdseries\slshape S}}. }

 }

  
\section{\textcolor{Chapter }{Green's Relations}}\label{Green's Relations}
\logpage{[ 51, 8, 0 ]}
\hyperdef{L}{X80C6C718801855E9}{}
{
  Green's equivalence relations play a very important role in semigroup theory.
In this section we describe how they can be used in \textsf{GAP}. 

 The five Green's relations are $R$, $L$, $J$, $H$, $D$: two elements $x$, $y$ from a semigroup $S$ are $R$\texttt{\symbol{45}}related if and only if $xS^1 = yS^1$, $L$\texttt{\symbol{45}}related if and only if $S^1 x = S^1 y$ and $J$\texttt{\symbol{45}}related if and only if $S^1 xS^1 = S^1 yS^1$; finally, $H = R \wedge L$, and $D = R \circ L$. 

 Recall that relations $R$, $L$ and $J$ induce a partial order among the elements of the semigroup $S$: for two elements $x$, $y$ from $S$, we say that $x$ is less than or equal to $y$ in the order on $R$ if $xS^1 \subseteq yS^1$; similarly, $x$ is less than or equal to $y$ under $L$ if $S^1x \subseteq S^1y$; finally $x$ is less than or equal to $y$ under $J$ if $S^1 xS^1 \subseteq S^1 tS^1$. We extend this preorder to a partial order on equivalence classes in the
natural way. 

\subsection{\textcolor{Chapter }{GreensRRelation}}
\logpage{[ 51, 8, 1 ]}\nobreak
\hyperdef{L}{X786CEDD4814A9079}{}
{\noindent\textcolor{FuncColor}{$\triangleright$\enspace\texttt{GreensRRelation({\mdseries\slshape semigroup})\index{GreensRRelation@\texttt{GreensRRelation}}
\label{GreensRRelation}
}\hfill{\scriptsize (attribute)}}\\
\noindent\textcolor{FuncColor}{$\triangleright$\enspace\texttt{GreensLRelation({\mdseries\slshape semigroup})\index{GreensLRelation@\texttt{GreensLRelation}}
\label{GreensLRelation}
}\hfill{\scriptsize (attribute)}}\\
\noindent\textcolor{FuncColor}{$\triangleright$\enspace\texttt{GreensJRelation({\mdseries\slshape semigroup})\index{GreensJRelation@\texttt{GreensJRelation}}
\label{GreensJRelation}
}\hfill{\scriptsize (attribute)}}\\
\noindent\textcolor{FuncColor}{$\triangleright$\enspace\texttt{GreensDRelation({\mdseries\slshape semigroup})\index{GreensDRelation@\texttt{GreensDRelation}}
\label{GreensDRelation}
}\hfill{\scriptsize (attribute)}}\\
\noindent\textcolor{FuncColor}{$\triangleright$\enspace\texttt{GreensHRelation({\mdseries\slshape semigroup})\index{GreensHRelation@\texttt{GreensHRelation}}
\label{GreensHRelation}
}\hfill{\scriptsize (attribute)}}\\


 The Green's relations (which are equivalence relations) are attributes of the
semigroup \mbox{\texttt{\mdseries\slshape semigroup}}. }

 

\subsection{\textcolor{Chapter }{IsGreensRelation}}
\logpage{[ 51, 8, 2 ]}\nobreak
\hyperdef{L}{X8364D69987D49DE1}{}
{\noindent\textcolor{FuncColor}{$\triangleright$\enspace\texttt{IsGreensRelation({\mdseries\slshape bin\texttt{\symbol{45}}relation})\index{IsGreensRelation@\texttt{IsGreensRelation}}
\label{IsGreensRelation}
}\hfill{\scriptsize (filter)}}\\
\noindent\textcolor{FuncColor}{$\triangleright$\enspace\texttt{IsGreensRRelation({\mdseries\slshape equiv\texttt{\symbol{45}}relation})\index{IsGreensRRelation@\texttt{IsGreensRRelation}}
\label{IsGreensRRelation}
}\hfill{\scriptsize (filter)}}\\
\noindent\textcolor{FuncColor}{$\triangleright$\enspace\texttt{IsGreensLRelation({\mdseries\slshape equiv\texttt{\symbol{45}}relation})\index{IsGreensLRelation@\texttt{IsGreensLRelation}}
\label{IsGreensLRelation}
}\hfill{\scriptsize (filter)}}\\
\noindent\textcolor{FuncColor}{$\triangleright$\enspace\texttt{IsGreensJRelation({\mdseries\slshape equiv\texttt{\symbol{45}}relation})\index{IsGreensJRelation@\texttt{IsGreensJRelation}}
\label{IsGreensJRelation}
}\hfill{\scriptsize (filter)}}\\
\noindent\textcolor{FuncColor}{$\triangleright$\enspace\texttt{IsGreensHRelation({\mdseries\slshape equiv\texttt{\symbol{45}}relation})\index{IsGreensHRelation@\texttt{IsGreensHRelation}}
\label{IsGreensHRelation}
}\hfill{\scriptsize (filter)}}\\
\noindent\textcolor{FuncColor}{$\triangleright$\enspace\texttt{IsGreensDRelation({\mdseries\slshape equiv\texttt{\symbol{45}}relation})\index{IsGreensDRelation@\texttt{IsGreensDRelation}}
\label{IsGreensDRelation}
}\hfill{\scriptsize (filter)}}\\


 Categories for the Green's relations. }

 

\subsection{\textcolor{Chapter }{IsGreensClass}}
\logpage{[ 51, 8, 3 ]}\nobreak
\hyperdef{L}{X82A11A087AFB3EB0}{}
{\noindent\textcolor{FuncColor}{$\triangleright$\enspace\texttt{IsGreensClass({\mdseries\slshape equiv\texttt{\symbol{45}}class})\index{IsGreensClass@\texttt{IsGreensClass}}
\label{IsGreensClass}
}\hfill{\scriptsize (filter)}}\\
\noindent\textcolor{FuncColor}{$\triangleright$\enspace\texttt{IsGreensRClass({\mdseries\slshape equiv\texttt{\symbol{45}}class})\index{IsGreensRClass@\texttt{IsGreensRClass}}
\label{IsGreensRClass}
}\hfill{\scriptsize (filter)}}\\
\noindent\textcolor{FuncColor}{$\triangleright$\enspace\texttt{IsGreensLClass({\mdseries\slshape equiv\texttt{\symbol{45}}class})\index{IsGreensLClass@\texttt{IsGreensLClass}}
\label{IsGreensLClass}
}\hfill{\scriptsize (filter)}}\\
\noindent\textcolor{FuncColor}{$\triangleright$\enspace\texttt{IsGreensJClass({\mdseries\slshape equiv\texttt{\symbol{45}}class})\index{IsGreensJClass@\texttt{IsGreensJClass}}
\label{IsGreensJClass}
}\hfill{\scriptsize (filter)}}\\
\noindent\textcolor{FuncColor}{$\triangleright$\enspace\texttt{IsGreensHClass({\mdseries\slshape equiv\texttt{\symbol{45}}class})\index{IsGreensHClass@\texttt{IsGreensHClass}}
\label{IsGreensHClass}
}\hfill{\scriptsize (filter)}}\\
\noindent\textcolor{FuncColor}{$\triangleright$\enspace\texttt{IsGreensDClass({\mdseries\slshape equiv\texttt{\symbol{45}}class})\index{IsGreensDClass@\texttt{IsGreensDClass}}
\label{IsGreensDClass}
}\hfill{\scriptsize (filter)}}\\


 return \texttt{true} if the equivalence class \mbox{\texttt{\mdseries\slshape equiv\texttt{\symbol{45}}class}} is a Green's class of any type, or of $R$, $L$, $J$, $H$, $D$ type, respectively, or \texttt{false} otherwise. }

 

\subsection{\textcolor{Chapter }{IsGreensLessThanOrEqual}}
\logpage{[ 51, 8, 4 ]}\nobreak
\hyperdef{L}{X7AA204C8850F9070}{}
{\noindent\textcolor{FuncColor}{$\triangleright$\enspace\texttt{IsGreensLessThanOrEqual({\mdseries\slshape C1, C2})\index{IsGreensLessThanOrEqual@\texttt{IsGreensLessThanOrEqual}}
\label{IsGreensLessThanOrEqual}
}\hfill{\scriptsize (operation)}}\\


 returns \texttt{true} if the Green's class \mbox{\texttt{\mdseries\slshape C1}} is less than or equal to \mbox{\texttt{\mdseries\slshape C2}} under the respective ordering (as defined above), and \texttt{false} otherwise. 

 Only defined for $R$, $L$ and $J$ classes. }

 

\subsection{\textcolor{Chapter }{RClassOfHClass}}
\logpage{[ 51, 8, 5 ]}\nobreak
\hyperdef{L}{X86FE5F5585EBCF13}{}
{\noindent\textcolor{FuncColor}{$\triangleright$\enspace\texttt{RClassOfHClass({\mdseries\slshape H})\index{RClassOfHClass@\texttt{RClassOfHClass}}
\label{RClassOfHClass}
}\hfill{\scriptsize (attribute)}}\\
\noindent\textcolor{FuncColor}{$\triangleright$\enspace\texttt{LClassOfHClass({\mdseries\slshape H})\index{LClassOfHClass@\texttt{LClassOfHClass}}
\label{LClassOfHClass}
}\hfill{\scriptsize (attribute)}}\\


 are attributes reflecting the natural ordering over the various Green's
classes. \texttt{RClassOfHClass} and \texttt{LClassOfHClass} return the $R$ and $L$ classes, respectively, in which an $H$ class is contained. }

 

\subsection{\textcolor{Chapter }{EggBoxOfDClass}}
\logpage{[ 51, 8, 6 ]}\nobreak
\hyperdef{L}{X78C56F4A78E0088A}{}
{\noindent\textcolor{FuncColor}{$\triangleright$\enspace\texttt{EggBoxOfDClass({\mdseries\slshape Dclass})\index{EggBoxOfDClass@\texttt{EggBoxOfDClass}}
\label{EggBoxOfDClass}
}\hfill{\scriptsize (attribute)}}\\


 returns for a Green's $D$ class \mbox{\texttt{\mdseries\slshape Dclass}} a matrix whose rows represent $R$ classes and columns represent $L$ classes. The entries are the $H$ classes. }

 

\subsection{\textcolor{Chapter }{DisplayEggBoxOfDClass}}
\logpage{[ 51, 8, 7 ]}\nobreak
\hyperdef{L}{X803237F17ACD44E3}{}
{\noindent\textcolor{FuncColor}{$\triangleright$\enspace\texttt{DisplayEggBoxOfDClass({\mdseries\slshape Dclass})\index{DisplayEggBoxOfDClass@\texttt{DisplayEggBoxOfDClass}}
\label{DisplayEggBoxOfDClass}
}\hfill{\scriptsize (function)}}\\


 displays a ``picture'' of the $D$ class \mbox{\texttt{\mdseries\slshape Dclass}}, as an array of 1s and 0s. A 1 represents a group $H$ class. }

 

\subsection{\textcolor{Chapter }{GreensRClassOfElement}}
\logpage{[ 51, 8, 8 ]}\nobreak
\hyperdef{L}{X87C75A9D86122D93}{}
{\noindent\textcolor{FuncColor}{$\triangleright$\enspace\texttt{GreensRClassOfElement({\mdseries\slshape S, a})\index{GreensRClassOfElement@\texttt{GreensRClassOfElement}}
\label{GreensRClassOfElement}
}\hfill{\scriptsize (operation)}}\\
\noindent\textcolor{FuncColor}{$\triangleright$\enspace\texttt{GreensLClassOfElement({\mdseries\slshape S, a})\index{GreensLClassOfElement@\texttt{GreensLClassOfElement}}
\label{GreensLClassOfElement}
}\hfill{\scriptsize (operation)}}\\
\noindent\textcolor{FuncColor}{$\triangleright$\enspace\texttt{GreensDClassOfElement({\mdseries\slshape S, a})\index{GreensDClassOfElement@\texttt{GreensDClassOfElement}}
\label{GreensDClassOfElement}
}\hfill{\scriptsize (operation)}}\\
\noindent\textcolor{FuncColor}{$\triangleright$\enspace\texttt{GreensJClassOfElement({\mdseries\slshape S, a})\index{GreensJClassOfElement@\texttt{GreensJClassOfElement}}
\label{GreensJClassOfElement}
}\hfill{\scriptsize (operation)}}\\
\noindent\textcolor{FuncColor}{$\triangleright$\enspace\texttt{GreensHClassOfElement({\mdseries\slshape S, a})\index{GreensHClassOfElement@\texttt{GreensHClassOfElement}}
\label{GreensHClassOfElement}
}\hfill{\scriptsize (operation)}}\\


 Creates the $X$ class of the element \mbox{\texttt{\mdseries\slshape a}} in the semigroup \mbox{\texttt{\mdseries\slshape S}} where $X$ is one of $L$, $R$, $D$, $J$, or $H$. }

 

\subsection{\textcolor{Chapter }{GreensRClasses}}
\logpage{[ 51, 8, 9 ]}\nobreak
\hyperdef{L}{X844D20467A644811}{}
{\noindent\textcolor{FuncColor}{$\triangleright$\enspace\texttt{GreensRClasses({\mdseries\slshape S})\index{GreensRClasses@\texttt{GreensRClasses}}
\label{GreensRClasses}
}\hfill{\scriptsize (attribute)}}\\
\noindent\textcolor{FuncColor}{$\triangleright$\enspace\texttt{GreensLClasses({\mdseries\slshape S})\index{GreensLClasses@\texttt{GreensLClasses}}
\label{GreensLClasses}
}\hfill{\scriptsize (attribute)}}\\
\noindent\textcolor{FuncColor}{$\triangleright$\enspace\texttt{GreensHClasses({\mdseries\slshape S})\index{GreensHClasses@\texttt{GreensHClasses}}
\label{GreensHClasses}
}\hfill{\scriptsize (attribute)}}\\
\noindent\textcolor{FuncColor}{$\triangleright$\enspace\texttt{GreensJClasses({\mdseries\slshape S})\index{GreensJClasses@\texttt{GreensJClasses}}
\label{GreensJClasses}
}\hfill{\scriptsize (attribute)}}\\
\noindent\textcolor{FuncColor}{$\triangleright$\enspace\texttt{GreensDClasses({\mdseries\slshape S})\index{GreensDClasses@\texttt{GreensDClasses}}
\label{GreensDClasses}
}\hfill{\scriptsize (attribute)}}\\


 If \mbox{\texttt{\mdseries\slshape S}} is a semigroup, then these attributes return the Green's $R$\texttt{\symbol{45}}, $L$\texttt{\symbol{45}}, $H$\texttt{\symbol{45}}, $J$\texttt{\symbol{45}}, or $D$\texttt{\symbol{45}}classes, respectively for the semigroup \mbox{\texttt{\mdseries\slshape S}}. 

 Additionally, if \mbox{\texttt{\mdseries\slshape S}} is a Green's $D$\texttt{\symbol{45}}class of a semigroup, then \texttt{GreensRClasses} and \texttt{GreensLClasses} return the Green's $R$\texttt{\symbol{45}} or $L-$classes of the semigroup, respectively, contained in the $D$\texttt{\symbol{45}}class \mbox{\texttt{\mdseries\slshape S}}; if \mbox{\texttt{\mdseries\slshape S}} is a Green's $D$\texttt{\symbol{45}}, $R$\texttt{\symbol{45}}, or $L$\texttt{\symbol{45}}class of a semigroup, then \texttt{GreensHClasses} returns the Green's $H$\texttt{\symbol{45}}classes of the semigroup contained in the Green's class \mbox{\texttt{\mdseries\slshape S}}. 

 \texttt{EquivalenceClasses} (\ref{EquivalenceClasses:attribute}) for a Green's relation lead to one of these functions. }

 

\subsection{\textcolor{Chapter }{GroupHClassOfGreensDClass}}
\logpage{[ 51, 8, 10 ]}\nobreak
\hyperdef{L}{X7CB4A18685B850E2}{}
{\noindent\textcolor{FuncColor}{$\triangleright$\enspace\texttt{GroupHClassOfGreensDClass({\mdseries\slshape Dclass})\index{GroupHClassOfGreensDClass@\texttt{GroupHClassOfGreensDClass}}
\label{GroupHClassOfGreensDClass}
}\hfill{\scriptsize (attribute)}}\\


 for a $D$ class \mbox{\texttt{\mdseries\slshape Dclass}} of a semigroup, returns a group $H$ class of the $D$ class, or \texttt{fail} if there is no group $H$ class. }

 

\subsection{\textcolor{Chapter }{IsGroupHClass}}
\logpage{[ 51, 8, 11 ]}\nobreak
\hyperdef{L}{X79D740EF7F0E53BD}{}
{\noindent\textcolor{FuncColor}{$\triangleright$\enspace\texttt{IsGroupHClass({\mdseries\slshape Hclass})\index{IsGroupHClass@\texttt{IsGroupHClass}}
\label{IsGroupHClass}
}\hfill{\scriptsize (property)}}\\


 returns \texttt{true} if the Green's $H$ class \mbox{\texttt{\mdseries\slshape Hclass}} is a group, which in turn is true if and only if \mbox{\texttt{\mdseries\slshape Hclass}}$^2$ intersects \mbox{\texttt{\mdseries\slshape Hclass}}. }

 

\subsection{\textcolor{Chapter }{IsRegularDClass}}
\logpage{[ 51, 8, 12 ]}\nobreak
\hyperdef{L}{X7F5860927CAD920F}{}
{\noindent\textcolor{FuncColor}{$\triangleright$\enspace\texttt{IsRegularDClass({\mdseries\slshape Dclass})\index{IsRegularDClass@\texttt{IsRegularDClass}}
\label{IsRegularDClass}
}\hfill{\scriptsize (property)}}\\


 returns \texttt{true} if the Greens $D$ class \mbox{\texttt{\mdseries\slshape Dclass}} is regular. A $D$ class is regular if and only if each of its elements is regular, which in turn
is true if and only if any one element of \mbox{\texttt{\mdseries\slshape Dclass}} is regular. Idempotents are regular since $eee = e$ so it follows that a Green's $D$ class containing an idempotent is regular. Conversely, it is true that a
regular $D$ class must contain at least one idempotent. (See{\nobreakspace}\cite[Prop. 3.2]{Howie76}.) }

 

\subsection{\textcolor{Chapter }{DisplaySemigroup}}
\logpage{[ 51, 8, 13 ]}\nobreak
\hyperdef{L}{X81AF2EAB7CEF8C19}{}
{\noindent\textcolor{FuncColor}{$\triangleright$\enspace\texttt{DisplaySemigroup({\mdseries\slshape S})\index{DisplaySemigroup@\texttt{DisplaySemigroup}}
\label{DisplaySemigroup}
}\hfill{\scriptsize (operation)}}\\


 Produces a convenient display of a transformation semigroup's
D\texttt{\symbol{45}}Class structure. Let \mbox{\texttt{\mdseries\slshape S}} be a transformation semigroup of degree $n$. Then for each $r\leq n$, we show all D\texttt{\symbol{45}}classes of rank $r$. 

 A regular D\texttt{\symbol{45}}class with a single H\texttt{\symbol{45}}class
of size 120 appears as 
\begin{Verbatim}[commandchars=!@|,fontsize=\small,frame=single,label=Example]
  *[H size = 120, 1 L-class, 1 R-class]
\end{Verbatim}
 (the \texttt{*} denoting regularity). }

 }

 
\section{\textcolor{Chapter }{Rees Matrix Semigroups}}\label{Rees Matrix Semigroups}
\logpage{[ 51, 9, 0 ]}
\hyperdef{L}{X8225A9EC87A255E6}{}
{
  In this section, we describe the functions in \textsf{GAP} for Rees matrix and 0\texttt{\symbol{45}}matrix semigroups and their
subsemigroups. The importance of these semigroups lies in the fact that Rees
matrix semigroups over groups are exactly the completely simple semigroups,
and Rees 0\texttt{\symbol{45}}matrix semigroups over groups are the completely
0\texttt{\symbol{45}}simple semigroups. 

 Let $I$ and $J$ be sets, let $S$ be a semigroup, and let $P=(p_{ji})_{j\in J, i\in I}$ be a $|J|\times |I|$ matrix with entries in $S$. Then the \emph{Rees matrix semigroup} with underlying semigroup $S$ and matrix $P$ is just the direct product $I\times S \times J$ with multiplication defined by 
\[(i, s, j)(k, t, l)=(i,s\cdot p_{j,k}\cdot t, l).\]
 Rees 0\texttt{\symbol{45}}matrix semigroups are defined as follows. If $I$, $J$, $S$, and $P$ are as above and $0$ denotes a new element, then the \emph{Rees 0\texttt{\symbol{45}}matrix semigroup} with underlying semigroup $S$ and matrix $P$ is $(I\times S\times J)\cup \{0\}$ with multiplication defined by 
\[(i, s, j)(k, t, l)=(i, s\cdot p_{j,k}\cdot t, l)\]
 when $p_{j,k}$ is not $0$ and $0$ if $p_{j,k}$ is 0.

 If $R$ is a Rees matrix or 0\texttt{\symbol{45}}matrix semigroup, then the \emph{rows} of $R$ is the index set $I$, the \emph{columns} of $R$ is the index set $J$, the semigroup $S$ is the \emph{underlying semigroup} of $R$, and the \emph{matrix} $P$ is the matrix of $S$. 

 Thoroughout this section, wherever the distinction is unimportant, we will
refer to Rees matrix or 0\texttt{\symbol{45}}matrix semigroups collectively as
Rees matrix semigroups.

 Multiplication of elements of a Rees matrix semigroup obviously depends on the
matrix used to create the semigroup. Hence elements of a Rees matrix semigroup
can only be created with reference to the semigroup to which they belong. More
specifically, every collection or semigroup of Rees matrix semigroup elements
is created from a specific Rees matrix semigroup, which contains the whole
family of its elements. So, it is not possible to multiply or compare elements
belonging to distinct Rees matrix semigroups, since they belong to different
families. For example, this situation is similar to free groups, but it is
different to permutations, which belong to a single family, and where
arbitrary permutations can be compared and multiplied without reference to any
group containing them. 

 A subsemigroup of a Rees matrix semigroup is not necessarily a Rees matrix
semigroup. Every semigroup consisting of elements of a Rees matrix semigroup
satisfies the property \texttt{IsReesMatrixSubsemigroup} (\ref{IsReesMatrixSubsemigroup}) and every semigroup of Rees 0\texttt{\symbol{45}}matrix semigroup elements
satisfies \texttt{IsReesZeroMatrixSubsemigroup} (\ref{IsReesZeroMatrixSubsemigroup}). 

 Rees matrix and 0\texttt{\symbol{45}}matrix semigroups can be created using
the operations \texttt{ReesMatrixSemigroup} (\ref{ReesMatrixSemigroup}) and \texttt{ReesZeroMatrixSemigroup} (\ref{ReesZeroMatrixSemigroup}), respectively, from an underlying semigroup and a matrix. Rees matrix
semigroups created in this way contain the whole family of their elements.
Every element of a Rees matrix semigroup belongs to a unique semigroup created
in this way; every subsemigroup of a Rees matrix semigroup is a subsemigroup
of a unique semigroup created in this way.

 Subsemigroups of Rees matrix semigroups can also be created by specifying
generators. A subsemigroup of a Rees matrix semigroup $I\times U\times J$ satisfies \texttt{IsReesMatrixSemigroup} (\ref{IsReesMatrixSemigroup}) if and only if it is equal to $I'\times U'\times J'$ where $I'\subseteq I$, $J'\subseteq J$, and $U'$ is a subsemigroup of $U$. The analogous statements holds for Rees 0\texttt{\symbol{45}}matrix
semigroups. 

 It is not necessarily the case that a simple subsemigroups of Rees matrix
semigroups satisfies \texttt{IsReesMatrixSemigroup} (\ref{IsReesMatrixSemigroup}). A Rees matrix semigroup is simple if and only if its underlying semigroup is
simple. A finite semigroup is simple if and only if it is isomorphic to a Rees
matrix semigroup over a group; this isomorphism can be obtained explicitly
using \texttt{IsomorphismReesMatrixSemigroup} (\ref{IsomorphismReesMatrixSemigroup}).

 Similarly, 0\texttt{\symbol{45}}simple subsemigroups of Rees
0\texttt{\symbol{45}}matrix semigroups do not have to satisfy \texttt{IsReesZeroMatrixSemigroup} (\ref{IsReesZeroMatrixSemigroup}). A Rees 0\texttt{\symbol{45}}matrix semigroup with more than 2 elements is
0\texttt{\symbol{45}}simple if and only if every row and every column of its
matrix contains a non\texttt{\symbol{45}}zero entry, and its underlying
semigroup is simple. A finite semigroup is 0\texttt{\symbol{45}}simple if and
only if it is isomorphic to a Rees 0\texttt{\symbol{45}}matrix semigroup over
a group; again this isomorphism can be found by using \texttt{IsomorphismReesZeroMatrixSemigroup} (\ref{IsomorphismReesZeroMatrixSemigroup}).

 Elements of a Rees matrix or 0\texttt{\symbol{45}}matrix semigroup belong to
the categories \texttt{IsReesMatrixSemigroupElement} (\ref{IsReesMatrixSemigroupElement}) and \texttt{IsReesZeroMatrixSemigroupElement} (\ref{IsReesZeroMatrixSemigroupElement}), respectively. Such elements can be created directly using the functions \texttt{ReesMatrixSemigroupElement} (\ref{ReesMatrixSemigroupElement}) and \texttt{ReesZeroMatrixSemigroupElement} (\ref{ReesZeroMatrixSemigroupElement}). 

 A semigroup in \textsf{GAP} can either satisfies \texttt{IsReesMatrixSemigroup} (\ref{IsReesMatrixSemigroup}) or \texttt{IsReesZeroMatrixSemigroup} (\ref{IsReesZeroMatrixSemigroup}) but not both. 

\subsection{\textcolor{Chapter }{ReesMatrixSemigroup}}
\logpage{[ 51, 9, 1 ]}\nobreak
\hyperdef{L}{X8526AA557CDF6C49}{}
{\noindent\textcolor{FuncColor}{$\triangleright$\enspace\texttt{ReesMatrixSemigroup({\mdseries\slshape S, mat})\index{ReesMatrixSemigroup@\texttt{ReesMatrixSemigroup}}
\label{ReesMatrixSemigroup}
}\hfill{\scriptsize (operation)}}\\
\noindent\textcolor{FuncColor}{$\triangleright$\enspace\texttt{ReesZeroMatrixSemigroup({\mdseries\slshape S, mat})\index{ReesZeroMatrixSemigroup@\texttt{ReesZeroMatrixSemigroup}}
\label{ReesZeroMatrixSemigroup}
}\hfill{\scriptsize (operation)}}\\
\textbf{\indent Returns:\ }
A Rees matrix or 0\texttt{\symbol{45}}matrix semigroup.



 When \mbox{\texttt{\mdseries\slshape S}} is a semigroup and \mbox{\texttt{\mdseries\slshape mat}} is an \texttt{m} by \texttt{n} matrix with entries in \mbox{\texttt{\mdseries\slshape S}}, the function \texttt{ReesMatrixSemigroup} returns the \texttt{n} by \texttt{m} Rees matrix semigroup over \mbox{\texttt{\mdseries\slshape S}} with multiplication defined by \mbox{\texttt{\mdseries\slshape mat}}. 

 The arguments of \texttt{ReesZeroMatrixSemigroup} should be a semigroup \mbox{\texttt{\mdseries\slshape S}} and an \texttt{m} by \texttt{n} matrix \mbox{\texttt{\mdseries\slshape mat}} with entries in \mbox{\texttt{\mdseries\slshape S}} or equal to the integer \texttt{0}. \texttt{ReesZeroMatrixSemigroup} returns the \texttt{n} by \texttt{m} Rees 0\texttt{\symbol{45}}matrix semigroup over \mbox{\texttt{\mdseries\slshape S}} with multiplication defined by \mbox{\texttt{\mdseries\slshape mat}}. In \textsf{GAP} a Rees 0\texttt{\symbol{45}}matrix semigroup always contains a multiplicative
zero element, regardless of whether there are any entries in \mbox{\texttt{\mdseries\slshape mat}} which are equal to \texttt{0}. 
\begin{Verbatim}[commandchars=!@|,fontsize=\small,frame=single,label=Example]
  !gapprompt@gap>| !gapinput@G:=Random(AllSmallGroups(Size, 32));;|
  !gapprompt@gap>| !gapinput@mat:=List([1..5], x-> List([1..3], y-> Random(G)));;|
  !gapprompt@gap>| !gapinput@S:=ReesMatrixSemigroup(G, mat);|
  <Rees matrix semigroup 3x5 over <pc group of size 32 with
   5 generators>>
  !gapprompt@gap>| !gapinput@mat:=[[(), 0, (), ()], [0, 0, 0, 0]];;|
  !gapprompt@gap>| !gapinput@S:=ReesZeroMatrixSemigroup(DihedralGroup(IsPermGroup, 8), mat);|
  <Rees 0-matrix semigroup 4x2 over Group([ (1,2,3,4), (2,4) ])>
\end{Verbatim}
 }

 

\subsection{\textcolor{Chapter }{ReesMatrixSubsemigroup}}
\logpage{[ 51, 9, 2 ]}\nobreak
\hyperdef{L}{X78D2A48C87FC8E38}{}
{\noindent\textcolor{FuncColor}{$\triangleright$\enspace\texttt{ReesMatrixSubsemigroup({\mdseries\slshape R, I, U, J})\index{ReesMatrixSubsemigroup@\texttt{ReesMatrixSubsemigroup}}
\label{ReesMatrixSubsemigroup}
}\hfill{\scriptsize (operation)}}\\
\noindent\textcolor{FuncColor}{$\triangleright$\enspace\texttt{ReesZeroMatrixSubsemigroup({\mdseries\slshape R, I, U, J})\index{ReesZeroMatrixSubsemigroup@\texttt{ReesZeroMatrixSubsemigroup}}
\label{ReesZeroMatrixSubsemigroup}
}\hfill{\scriptsize (operation)}}\\
\textbf{\indent Returns:\ }
A Rees matrix or 0\texttt{\symbol{45}}matrix subsemigroup.



 The arguments of \texttt{ReesMatrixSubsemigroup} should be a Rees matrix semigroup \mbox{\texttt{\mdseries\slshape R}}, subsets \mbox{\texttt{\mdseries\slshape I}} and \mbox{\texttt{\mdseries\slshape J}} of the rows and columns of \mbox{\texttt{\mdseries\slshape R}}, respectively, and a subsemigroup \mbox{\texttt{\mdseries\slshape U}} of the underlying semigroup of \mbox{\texttt{\mdseries\slshape R}}. \texttt{ReesMatrixSubsemigroup} returns the subsemigroup of \mbox{\texttt{\mdseries\slshape R}} generated by the direct product of \mbox{\texttt{\mdseries\slshape I}}, \mbox{\texttt{\mdseries\slshape U}}, and \mbox{\texttt{\mdseries\slshape J}}. 

 The usage and returned value of \texttt{ReesZeroMatrixSubsemigroup} is analogous when \mbox{\texttt{\mdseries\slshape R}} is a Rees 0\texttt{\symbol{45}}matrix semigroup. 
\begin{Verbatim}[commandchars=!@|,fontsize=\small,frame=single,label=Example]
  !gapprompt@gap>| !gapinput@G:=CyclicGroup(IsPermGroup, 1007);;|
  !gapprompt@gap>| !gapinput@mat:=[[(), 0, 0], [0, (), 0], [0, 0, ()],|
  !gapprompt@>| !gapinput@[(), (), ()], [0, 0, ()]];;|
  !gapprompt@gap>| !gapinput@R:=ReesZeroMatrixSemigroup(G, mat);|
  <Rees 0-matrix semigroup 3x5 over
    <permutation group of size 1007 with 1 generator>>
  !gapprompt@gap>| !gapinput@ReesZeroMatrixSubsemigroup(R, [1,3], G, [1..5]);|
  <Rees 0-matrix semigroup 2x5 over
    <permutation group of size 1007 with 1 generator>>
\end{Verbatim}
 }

 

\subsection{\textcolor{Chapter }{IsomorphismReesMatrixSemigroup}}
\logpage{[ 51, 9, 3 ]}\nobreak
\hyperdef{L}{X7964B5C97FB9C07D}{}
{\noindent\textcolor{FuncColor}{$\triangleright$\enspace\texttt{IsomorphismReesMatrixSemigroup({\mdseries\slshape S})\index{IsomorphismReesMatrixSemigroup@\texttt{IsomorphismReesMatrixSemigroup}}
\label{IsomorphismReesMatrixSemigroup}
}\hfill{\scriptsize (attribute)}}\\
\noindent\textcolor{FuncColor}{$\triangleright$\enspace\texttt{IsomorphismReesZeroMatrixSemigroup({\mdseries\slshape S})\index{IsomorphismReesZeroMatrixSemigroup@\texttt{IsomorphismReesZeroMatrixSemigroup}}
\label{IsomorphismReesZeroMatrixSemigroup}
}\hfill{\scriptsize (attribute)}}\\
\textbf{\indent Returns:\ }
An isomorphism.



 Every finite simple semigroup is isomorphic to a Rees matrix semigroup over a
group, and every finite 0\texttt{\symbol{45}}simple semigroup is isomorphic to
a Rees 0\texttt{\symbol{45}}matrix semigroup over a group. 

 If the argument \mbox{\texttt{\mdseries\slshape S}} is a simple semigroup, then \texttt{IsomorphismReesMatrixSemigroup} returns an isomorphism to a Rees matrix semigroup over a group. If \mbox{\texttt{\mdseries\slshape S}} is not simple, then \texttt{IsomorphismReesMatrixSemigroup} returns an error. 

 If the argument \mbox{\texttt{\mdseries\slshape S}} is a 0\texttt{\symbol{45}}simple semigroup, then \texttt{IsomorphismReesZeroMatrixSemigroup} returns an isomorphism to a Rees 0\texttt{\symbol{45}}matrix semigroup over a
group. If \mbox{\texttt{\mdseries\slshape S}} is not 0\texttt{\symbol{45}}simple, then \texttt{IsomorphismReesZeroMatrixSemigroup} returns an error. 

 See \texttt{IsSimpleSemigroup} (\ref{IsSimpleSemigroup}) and \texttt{IsZeroSimpleSemigroup} (\ref{IsZeroSimpleSemigroup}). 
\begin{Verbatim}[commandchars=!@|,fontsize=\small,frame=single,label=Example]
  !gapprompt@gap>| !gapinput@S := Semigroup(Transformation([2, 1, 1, 2, 1]),|
  !gapprompt@>| !gapinput@                  Transformation([3, 4, 3, 4, 4]),|
  !gapprompt@>| !gapinput@                  Transformation([3, 4, 3, 4, 3]),|
  !gapprompt@>| !gapinput@                  Transformation([4, 3, 3, 4, 4]));;|
  !gapprompt@gap>| !gapinput@IsSimpleSemigroup(S);|
  true
  !gapprompt@gap>| !gapinput@Range(IsomorphismReesMatrixSemigroup(S));|
  <Rees matrix semigroup 4x2 over Group([ (1,2) ])>
  !gapprompt@gap>| !gapinput@mat := [[(), 0, 0],|
  !gapprompt@>| !gapinput@           [0, (), 0],|
  !gapprompt@>| !gapinput@           [0, 0, ()]];;|
  !gapprompt@gap>| !gapinput@R := ReesZeroMatrixSemigroup(Group((1,2,4,5,6)), mat);|
  <Rees 0-matrix semigroup 3x3 over Group([ (1,2,4,5,6) ])>
  !gapprompt@gap>| !gapinput@U := ReesZeroMatrixSubsemigroup(R, [1, 2], Group(()), [2, 3]);|
  <subsemigroup of 3x3 Rees 0-matrix semigroup with 4 generators>
  !gapprompt@gap>| !gapinput@IsZeroSimpleSemigroup(U);|
  false
  !gapprompt@gap>| !gapinput@U := ReesZeroMatrixSubsemigroup(R, [2, 3], Group(()), [2, 3]);|
  <subsemigroup of 3x3 Rees 0-matrix semigroup with 3 generators>
  !gapprompt@gap>| !gapinput@IsZeroSimpleSemigroup(U);|
  true
  !gapprompt@gap>| !gapinput@Rows(U); Columns(U);|
  [ 2, 3 ]
  [ 2, 3 ]
  !gapprompt@gap>| !gapinput@V := Range(IsomorphismReesZeroMatrixSemigroup(U));|
  <Rees 0-matrix semigroup 2x2 over Group(())>
  !gapprompt@gap>| !gapinput@Rows(V); Columns(V);|
  [ 1, 2 ]
  [ 1, 2 ]
\end{Verbatim}
 }

 

\subsection{\textcolor{Chapter }{IsReesMatrixSemigroupElement}}
\logpage{[ 51, 9, 4 ]}\nobreak
\hyperdef{L}{X7F6B852B81488C86}{}
{\noindent\textcolor{FuncColor}{$\triangleright$\enspace\texttt{IsReesMatrixSemigroupElement({\mdseries\slshape elt})\index{IsReesMatrixSemigroupElement@\texttt{IsReesMatrixSemigroupElement}}
\label{IsReesMatrixSemigroupElement}
}\hfill{\scriptsize (Category)}}\\
\noindent\textcolor{FuncColor}{$\triangleright$\enspace\texttt{IsReesZeroMatrixSemigroupElement({\mdseries\slshape elt})\index{IsReesZeroMatrixSemigroupElement@\texttt{IsReesZeroMatrixSemigroupElement}}
\label{IsReesZeroMatrixSemigroupElement}
}\hfill{\scriptsize (Category)}}\\
\textbf{\indent Returns:\ }
\texttt{true} or \texttt{false}.



 Every element of a Rees matrix semigroup belongs to the category \texttt{IsReesMatrixSemigroupElement}, and every element of a Rees 0\texttt{\symbol{45}}matrix semigroup belongs to
the category \texttt{IsReesZeroMatrixSemigroupElement}. 
\begin{Verbatim}[commandchars=!@|,fontsize=\small,frame=single,label=Example]
  !gapprompt@gap>| !gapinput@G:=Group((1,2,3));;|
  !gapprompt@gap>| !gapinput@mat:=[ [ (), (1,3,2) ], [ (1,3,2), () ] ];;|
  !gapprompt@gap>| !gapinput@R:=ReesMatrixSemigroup(G, mat);|
  <Rees matrix semigroup 2x2 over Group([ (1,2,3) ])>
  !gapprompt@gap>| !gapinput@GeneratorsOfSemigroup(R);|
  [ (1,(1,2,3),1), (2,(),2) ]
  !gapprompt@gap>| !gapinput@IsReesMatrixSemigroupElement(last[1]);|
  true
  !gapprompt@gap>| !gapinput@IsReesZeroMatrixSemigroupElement(last2[1]);|
  false
\end{Verbatim}
 }

 

\subsection{\textcolor{Chapter }{ReesMatrixSemigroupElement}}
\logpage{[ 51, 9, 5 ]}\nobreak
\hyperdef{L}{X7A0DE1F28470295E}{}
{\noindent\textcolor{FuncColor}{$\triangleright$\enspace\texttt{ReesMatrixSemigroupElement({\mdseries\slshape R, i, x, j})\index{ReesMatrixSemigroupElement@\texttt{ReesMatrixSemigroupElement}}
\label{ReesMatrixSemigroupElement}
}\hfill{\scriptsize (function)}}\\
\noindent\textcolor{FuncColor}{$\triangleright$\enspace\texttt{ReesZeroMatrixSemigroupElement({\mdseries\slshape R, i, x, j})\index{ReesZeroMatrixSemigroupElement@\texttt{ReesZeroMatrixSemigroupElement}}
\label{ReesZeroMatrixSemigroupElement}
}\hfill{\scriptsize (function)}}\\
\textbf{\indent Returns:\ }
An element of a Rees matrix or \texttt{0}\texttt{\symbol{45}}matrix semigroup.



 The arguments of \mbox{\texttt{\mdseries\slshape ReesMatrixSemigroupElement}} should be a Rees matrix subsemigroup \mbox{\texttt{\mdseries\slshape R}}, elements \mbox{\texttt{\mdseries\slshape i}} and \mbox{\texttt{\mdseries\slshape j}} of the the rows and columns of \mbox{\texttt{\mdseries\slshape R}}, respectively, and an element \mbox{\texttt{\mdseries\slshape x}} of the underlying semigroup of \mbox{\texttt{\mdseries\slshape R}}. \texttt{ReesMatrixSemigroupElement} returns the element of \mbox{\texttt{\mdseries\slshape R}} with row index \mbox{\texttt{\mdseries\slshape i}}, underlying element \mbox{\texttt{\mdseries\slshape x}} in the underlying semigroup of \mbox{\texttt{\mdseries\slshape R}}, and column index \mbox{\texttt{\mdseries\slshape j}}, if such an element exist, if such an element exists.

 The usage of \texttt{ReesZeroMatrixSemigroupElement} is analogous to that of \texttt{ReesMatrixSemigroupElement}, when \mbox{\texttt{\mdseries\slshape R}} is a Rees 0\texttt{\symbol{45}}matrix semigroup. 

 The row \mbox{\texttt{\mdseries\slshape i}}, underlying element \mbox{\texttt{\mdseries\slshape x}}, and column \mbox{\texttt{\mdseries\slshape j}} of an element \texttt{y} of a Rees matrix (or 0\texttt{\symbol{45}}matrix) semigroup can be recovered
from \texttt{y} using \texttt{y[1]}, \texttt{y[2]}, and \texttt{y[3]}, respectively.

 
\begin{Verbatim}[commandchars=!@|,fontsize=\small,frame=single,label=Example]
  !gapprompt@gap>| !gapinput@G:=Group((1,2,3));;|
  !gapprompt@gap>| !gapinput@mat:=[ [ 0, () ], [ (1,3,2), (1,3,2) ] ];;|
  !gapprompt@gap>| !gapinput@R:=ReesZeroMatrixSemigroup(G, mat);|
  <Rees 0-matrix semigroup 2x2 over Group([ (1,2,3) ])>
  !gapprompt@gap>| !gapinput@ReesZeroMatrixSemigroupElement(R, 1, (1,2,3), 2);|
  (1,(1,2,3),2)
  !gapprompt@gap>| !gapinput@MultiplicativeZero(R);|
  0
\end{Verbatim}
 }

 

\subsection{\textcolor{Chapter }{IsReesMatrixSubsemigroup}}
\logpage{[ 51, 9, 6 ]}\nobreak
\hyperdef{L}{X7F03BE707AC7F8A0}{}
{\noindent\textcolor{FuncColor}{$\triangleright$\enspace\texttt{IsReesMatrixSubsemigroup({\mdseries\slshape R})\index{IsReesMatrixSubsemigroup@\texttt{IsReesMatrixSubsemigroup}}
\label{IsReesMatrixSubsemigroup}
}\hfill{\scriptsize (Synonym)}}\\
\noindent\textcolor{FuncColor}{$\triangleright$\enspace\texttt{IsReesZeroMatrixSubsemigroup({\mdseries\slshape R})\index{IsReesZeroMatrixSubsemigroup@\texttt{IsReesZeroMatrixSubsemigroup}}
\label{IsReesZeroMatrixSubsemigroup}
}\hfill{\scriptsize (Synonym)}}\\
\textbf{\indent Returns:\ }
\texttt{true} or \texttt{false}.



 Every semigroup consisting of elements of a Rees matrix semigroup satisfies
the property \texttt{IsReesMatrixSubsemigroup} and every semigroup of Rees 0\texttt{\symbol{45}}matrix semigroup elements
satisfies \texttt{IsReesZeroMatrixSubsemigroup}. 

 Note that a subsemigroup of a Rees matrix semigroup is not necessarily a Rees
matrix semigroup. 
\begin{Verbatim}[commandchars=!@|,fontsize=\small,frame=single,label=Example]
  !gapprompt@gap>| !gapinput@G:=DihedralGroup(32);;|
  !gapprompt@gap>| !gapinput@mat:=List([1..2], x-> List([1..10], x-> Random(G)));;|
  !gapprompt@gap>| !gapinput@R:=ReesMatrixSemigroup(G, mat);|
  <Rees matrix semigroup 10x2 over <pc group of size 32 with
   5 generators>>
  !gapprompt@gap>| !gapinput@S:=Semigroup(GeneratorsOfSemigroup(R));|
  <subsemigroup of 10x2 Rees matrix semigroup with 14 generators>
  !gapprompt@gap>| !gapinput@IsReesMatrixSubsemigroup(S);|
  true
  !gapprompt@gap>| !gapinput@S:=Semigroup(GeneratorsOfSemigroup(R)[1]);|
  <subsemigroup of 10x2 Rees matrix semigroup with 1 generator>
  !gapprompt@gap>| !gapinput@IsReesMatrixSubsemigroup(S);|
  true
\end{Verbatim}
 }

 

\subsection{\textcolor{Chapter }{IsReesMatrixSemigroup}}
\logpage{[ 51, 9, 7 ]}\nobreak
\hyperdef{L}{X780BB78A79275244}{}
{\noindent\textcolor{FuncColor}{$\triangleright$\enspace\texttt{IsReesMatrixSemigroup({\mdseries\slshape R})\index{IsReesMatrixSemigroup@\texttt{IsReesMatrixSemigroup}}
\label{IsReesMatrixSemigroup}
}\hfill{\scriptsize (property)}}\\
\noindent\textcolor{FuncColor}{$\triangleright$\enspace\texttt{IsReesZeroMatrixSemigroup({\mdseries\slshape R})\index{IsReesZeroMatrixSemigroup@\texttt{IsReesZeroMatrixSemigroup}}
\label{IsReesZeroMatrixSemigroup}
}\hfill{\scriptsize (property)}}\\
\textbf{\indent Returns:\ }
\texttt{true} or \texttt{false}.



 A subsemigroup of a Rees matrix semigroup $I\times U\times J$ satisfies \texttt{IsReesMatrixSemigroup} if and only if it is equal to $I'\times U'\times J'$ where $I'\subseteq I$, $J'\subseteq J$, and $U'$ is a subsemigroup of $U$. It can be costly to check that a subsemigroup defined by generators
satisfies \texttt{IsReesMatrixSemigroup}. The analogous statements holds for Rees 0\texttt{\symbol{45}}matrix
semigroups. 

 It is not necessarily the case that a simple subsemigroups of Rees matrix
semigroups satisfies \texttt{IsReesMatrixSemigroup}. A Rees matrix semigroup is simple if and only if its underlying semigroup is
simple. A finite semigroup is simple if and only if it is isomorphic to a Rees
matrix semigroup over a group; this isomorphism can be obtained explicitly
using \texttt{IsomorphismReesMatrixSemigroup} (\ref{IsomorphismReesMatrixSemigroup}).

 Similarly, 0\texttt{\symbol{45}}simple subsemigroups of Rees
0\texttt{\symbol{45}}matrix semigroups do not have to satisfy \texttt{IsReesZeroMatrixSemigroup}. A Rees 0\texttt{\symbol{45}}matrix semigroup with more than 2 elements is
0\texttt{\symbol{45}}simple if and only if every row and every column of its
matrix contains a non\texttt{\symbol{45}}zero entry, and its underlying
semigroup is simple. A finite semigroup is 0\texttt{\symbol{45}}simple if and
only if it is isomorphic to a Rees 0\texttt{\symbol{45}}matrix semigroup over
a group; again this isomorphism can be found by using \texttt{IsomorphismReesMatrixSemigroup} (\ref{IsomorphismReesMatrixSemigroup}). 
\begin{Verbatim}[commandchars=!@|,fontsize=\small,frame=single,label=Example]
  !gapprompt@gap>| !gapinput@G:=PSL(2,5);;|
  !gapprompt@gap>| !gapinput@mat:=[ [ 0, (), 0, (2,6,3,5,4) ],|
  !gapprompt@>| !gapinput@[ (), 0, (), 0 ], [ 0, 0, 0, () ] ];;|
  !gapprompt@gap>| !gapinput@R:=ReesZeroMatrixSemigroup(G, mat);|
  <Rees 0-matrix semigroup 4x3 over Group([ (3,5)(4,6), (1,2,5)
  (3,4,6) ])>
  !gapprompt@gap>| !gapinput@IsReesZeroMatrixSemigroup(R);|
  true
  !gapprompt@gap>| !gapinput@U:=ReesZeroMatrixSubsemigroup(R, [1..3], Group(()), [1..2]);|
  <subsemigroup of 4x3 Rees 0-matrix semigroup with 4 generators>
  !gapprompt@gap>| !gapinput@IsReesZeroMatrixSemigroup(U);|
  true
  !gapprompt@gap>| !gapinput@V:=Semigroup(GeneratorsOfSemigroup(U));|
  <subsemigroup of 4x3 Rees 0-matrix semigroup with 4 generators>
  !gapprompt@gap>| !gapinput@IsReesZeroMatrixSemigroup(V);|
  true
  !gapprompt@gap>| !gapinput@S:=Semigroup(Transformation([1,1]), Transformation([1,2]));|
  <commutative transformation monoid of degree 2 with 1 generator>
  !gapprompt@gap>| !gapinput@IsSimpleSemigroup(S);|
  false
  !gapprompt@gap>| !gapinput@mat:=[[0, One(S), 0, One(S)], [One(S), 0, One(S), 0],|
  !gapprompt@>| !gapinput@[0, 0, 0, One(S)]];;|
  !gapprompt@gap>| !gapinput@R:=ReesZeroMatrixSemigroup(S, mat);;|
  !gapprompt@gap>| !gapinput@U:=ReesZeroMatrixSubsemigroup(R, [1..3],|
  !gapprompt@>| !gapinput@Semigroup(Transformation([1,1])), [1..2]);|
  <subsemigroup of 4x3 Rees 0-matrix semigroup with 6 generators>
  !gapprompt@gap>| !gapinput@V:=Semigroup(GeneratorsOfSemigroup(U));|
  <subsemigroup of 4x3 Rees 0-matrix semigroup with 6 generators>
  !gapprompt@gap>| !gapinput@IsReesZeroMatrixSemigroup(V);|
  true
  !gapprompt@gap>| !gapinput@T:=Semigroup(|
  !gapprompt@>| !gapinput@ReesZeroMatrixSemigroupElement(R, 3, Transformation( [ 1, 1 ] ), 3),|
  !gapprompt@>| !gapinput@ReesZeroMatrixSemigroupElement(R, 2, Transformation( [ 1, 1 ] ), 2));|
  <subsemigroup of 4x3 Rees 0-matrix semigroup with 2 generators>
  !gapprompt@gap>| !gapinput@IsReesZeroMatrixSemigroup(T);|
  false
\end{Verbatim}
 }

 

\subsection{\textcolor{Chapter }{Matrix (for Rees matrix semigroups)}}
\logpage{[ 51, 9, 8 ]}\nobreak
\hyperdef{L}{X7CACF4D686AF1D19}{}
{\noindent\textcolor{FuncColor}{$\triangleright$\enspace\texttt{Matrix({\mdseries\slshape R})\index{Matrix@\texttt{Matrix}!for Rees matrix semigroups}
\label{Matrix:for Rees matrix semigroups}
}\hfill{\scriptsize (operation)}}\\
\noindent\textcolor{FuncColor}{$\triangleright$\enspace\texttt{MatrixOfReesMatrixSemigroup({\mdseries\slshape R})\index{MatrixOfReesMatrixSemigroup@\texttt{MatrixOfReesMatrixSemigroup}}
\label{MatrixOfReesMatrixSemigroup}
}\hfill{\scriptsize (attribute)}}\\
\noindent\textcolor{FuncColor}{$\triangleright$\enspace\texttt{MatrixOfReesZeroMatrixSemigroup({\mdseries\slshape R})\index{MatrixOfReesZeroMatrixSemigroup@\texttt{MatrixOfReesZeroMatrixSemigroup}}
\label{MatrixOfReesZeroMatrixSemigroup}
}\hfill{\scriptsize (attribute)}}\\
\textbf{\indent Returns:\ }
A matrix.



 If \mbox{\texttt{\mdseries\slshape R}} is a Rees matrix or 0\texttt{\symbol{45}}matrix semigroup, then \texttt{MatrixOfReesMatrixSemigroup} respectively \texttt{MatrixOfReesZeroMatrixSemigroup} return the matrix used to define multiplication in \mbox{\texttt{\mdseries\slshape R}}. For convenience, one may also abbreviate either to \texttt{Matrix}. 

 More specifically, if \mbox{\texttt{\mdseries\slshape R}} is a Rees matrix or 0\texttt{\symbol{45}}matrix semigroup, which is a proper
subsemigroup of another such semigroup, then \texttt{Matrix} returns the matrix used to define the Rees matrix (or
0\texttt{\symbol{45}}matrix) semigroup consisting of the whole family to which
the elements of \mbox{\texttt{\mdseries\slshape R}} belong. Thus, for example, a \texttt{1} by \texttt{1} Rees matrix semigroup can have a \texttt{65} by \texttt{15} matrix. 

 Arbitrary subsemigroups of Rees matrix or 0\texttt{\symbol{45}}matrix
semigroups do not have a matrix. Such a subsemigroup \mbox{\texttt{\mdseries\slshape R}} has a matrix if and only if it satisfies \texttt{IsReesMatrixSemigroup} (\ref{IsReesMatrixSemigroup}) or \texttt{IsReesZeroMatrixSemigroup} (\ref{IsReesZeroMatrixSemigroup}). 
\begin{Verbatim}[commandchars=!@|,fontsize=\small,frame=single,label=Example]
  !gapprompt@gap>| !gapinput@G:=AlternatingGroup(5);;|
  !gapprompt@gap>| !gapinput@mat:=[[(), (), ()], [(), (), ()]];;|
  !gapprompt@gap>| !gapinput@R:=ReesMatrixSemigroup(G, mat);|
  <Rees matrix semigroup 3x2 over Alt( [ 1 .. 5 ] )>
  !gapprompt@gap>| !gapinput@Matrix(R);|
  [ [ (), (), () ], [ (), (), () ] ]
  !gapprompt@gap>| !gapinput@R:=ReesMatrixSubsemigroup(R, [1,2], Group(()), [2]);|
  <subsemigroup of 3x2 Rees matrix semigroup with 2 generators>
  !gapprompt@gap>| !gapinput@Matrix(R);|
  [ [ (), (), () ], [ (), (), () ] ]
\end{Verbatim}
 }

 
\subsection{\textcolor{Chapter }{Rows and columns}}\logpage{[ 51, 9, 9 ]}
\hyperdef{L}{X82FC5D6980C66AC4}{}
{
\noindent\textcolor{FuncColor}{$\triangleright$\enspace\texttt{Rows({\mdseries\slshape R})\index{Rows@\texttt{Rows}}
\label{Rows}
}\hfill{\scriptsize (attribute)}}\\
\noindent\textcolor{FuncColor}{$\triangleright$\enspace\texttt{Columns({\mdseries\slshape R})\index{Columns@\texttt{Columns}}
\label{Columns}
}\hfill{\scriptsize (attribute)}}\\
\textbf{\indent Returns:\ }
The rows or columns of \mbox{\texttt{\mdseries\slshape R}}.



 \texttt{Rows} returns the rows of the Rees matrix or 0\texttt{\symbol{45}}matrix semigroup \mbox{\texttt{\mdseries\slshape R}}. Note that the rows of the semigroup correspond to the columns of the matrix
used to define multiplication in \mbox{\texttt{\mdseries\slshape R}}.

 \texttt{Columns} returns the columns of the Rees matrix or 0\texttt{\symbol{45}}matrix
semigroup \mbox{\texttt{\mdseries\slshape R}}. Note that the columns of the semigroup correspond to the rows of the matrix
used to define multiplication in \mbox{\texttt{\mdseries\slshape R}}.

 Arbitrary subsemigroups of Rees matrix or 0\texttt{\symbol{45}}matrix
semigroups do not have rows or columns. Such a subsemigroup \mbox{\texttt{\mdseries\slshape R}} has rows and columns if and only if it satisfies \texttt{IsReesMatrixSemigroup} (\ref{IsReesMatrixSemigroup}) or \texttt{IsReesZeroMatrixSemigroup} (\ref{IsReesZeroMatrixSemigroup}). 
\begin{Verbatim}[commandchars=!@|,fontsize=\small,frame=single,label=Example]
  !gapprompt@gap>| !gapinput@G:=Group((1,2,3));;|
  !gapprompt@gap>| !gapinput@mat:=List([1..100], x-> List([1..200], x->Random(G)));;|
  !gapprompt@gap>| !gapinput@R:=ReesZeroMatrixSemigroup(G, mat);|
  <Rees 0-matrix semigroup 200x100 over Group([ (1,2,3) ])>
  !gapprompt@gap>| !gapinput@Rows(R);|
  [ 1 .. 200 ]
  !gapprompt@gap>| !gapinput@Columns(R);|
  [ 1 .. 100 ]
\end{Verbatim}
 }

 

\subsection{\textcolor{Chapter }{UnderlyingSemigroup (for a Rees matrix semigroup)}}
\logpage{[ 51, 9, 10 ]}\nobreak
\hyperdef{L}{X7D9719F887AFCF8F}{}
{\noindent\textcolor{FuncColor}{$\triangleright$\enspace\texttt{UnderlyingSemigroup({\mdseries\slshape R})\index{UnderlyingSemigroup@\texttt{UnderlyingSemigroup}!for a Rees matrix semigroup}
\label{UnderlyingSemigroup:for a Rees matrix semigroup}
}\hfill{\scriptsize (attribute)}}\\
\noindent\textcolor{FuncColor}{$\triangleright$\enspace\texttt{UnderlyingSemigroup({\mdseries\slshape R})\index{UnderlyingSemigroup@\texttt{UnderlyingSemigroup}!for a Rees 0-matrix semigroup}
\label{UnderlyingSemigroup:for a Rees 0-matrix semigroup}
}\hfill{\scriptsize (attribute)}}\\
\textbf{\indent Returns:\ }
A semigroup.



 \texttt{UnderlyingSemigroup} returns the underlying semigroup of the Rees matrix or
0\texttt{\symbol{45}}matrix semigroup \mbox{\texttt{\mdseries\slshape R}}. 

 Arbitrary subsemigroups of Rees matrix or 0\texttt{\symbol{45}}matrix
semigroups do not have an underlying semigroup. Such a subsemigroup \mbox{\texttt{\mdseries\slshape R}} has an underlying semigroup if and only if it satisfies \texttt{IsReesMatrixSemigroup} (\ref{IsReesMatrixSemigroup}) or \texttt{IsReesZeroMatrixSemigroup} (\ref{IsReesZeroMatrixSemigroup}). 
\begin{Verbatim}[commandchars=!@|,fontsize=\small,frame=single,label=Example]
  !gapprompt@gap>| !gapinput@S:=Semigroup(Transformation( [ 2, 1, 1, 2, 1 ] ),|
  !gapprompt@>| !gapinput@Transformation( [ 3, 4, 3, 4, 4 ] ), Transformation([ 3, 4, 3, 4, 3 ] ),|
  !gapprompt@>| !gapinput@Transformation([ 4, 3, 3, 4, 4 ] ) );;|
  !gapprompt@gap>| !gapinput@R:=Range(IsomorphismReesMatrixSemigroup(S));|
  <Rees matrix semigroup 4x2 over Group([ (1,2) ])>
  !gapprompt@gap>| !gapinput@UnderlyingSemigroup(R);|
  Group([ (1,2) ])
\end{Verbatim}
 }

 

\subsection{\textcolor{Chapter }{AssociatedReesMatrixSemigroupOfDClass}}
\logpage{[ 51, 9, 11 ]}\nobreak
\hyperdef{L}{X7D1D9A0382064B8F}{}
{\noindent\textcolor{FuncColor}{$\triangleright$\enspace\texttt{AssociatedReesMatrixSemigroupOfDClass({\mdseries\slshape D})\index{AssociatedReesMatrixSemigroupOfDClass@\texttt{Associated}\-\texttt{Rees}\-\texttt{Matrix}\-\texttt{Semigroup}\-\texttt{Of}\-\texttt{D}\-\texttt{Class}}
\label{AssociatedReesMatrixSemigroupOfDClass}
}\hfill{\scriptsize (attribute)}}\\
\textbf{\indent Returns:\ }
A Rees matrix or 0\texttt{\symbol{45}}matrix semigroup.



 If \mbox{\texttt{\mdseries\slshape D}} is a regular $\mathcal{D}$\texttt{\symbol{45}}class of a finite semigroup \texttt{S}, then there is a standard way of associating a Rees matrix semigroup to \mbox{\texttt{\mdseries\slshape D}}. If \mbox{\texttt{\mdseries\slshape D}} is a subsemigroup of \texttt{S}, then \mbox{\texttt{\mdseries\slshape D}} is simple and hence is isomorphic to a Rees matrix semigroup. In this case,
the associated Rees matrix semigroup of \mbox{\texttt{\mdseries\slshape D}} is just the Rees matrix semigroup isomorphic to \mbox{\texttt{\mdseries\slshape D}}. 

 If \mbox{\texttt{\mdseries\slshape D}} is not a subsemigroup of \texttt{S}, then we define a semigroup with elements \mbox{\texttt{\mdseries\slshape D}} and a new element \texttt{0} with multiplication of $x,y\in D$ defined by:  
\[ xy=\left\{\begin{array}{ll} x*y\ (\textrm{in }S)&\textrm{if }x*y\in D\\
0&\textrm{if }xy\not\in D. \end{array}\right. \]
   The semigroup thus defined is 0\texttt{\symbol{45}}simple and hence is
isomorphic to a Rees 0\texttt{\symbol{45}}matrix semigroup. This semigroup can
also be described as the Rees quotient of the ideal generated by \mbox{\texttt{\mdseries\slshape D}} by it maximal subideal. The associated Rees matrix semigroup of \mbox{\texttt{\mdseries\slshape D}} is just the Rees 0\texttt{\symbol{45}}matrix semigroup isomorphic to the
semigroup defined above. 
\begin{Verbatim}[commandchars=!@|,fontsize=\small,frame=single,label=Example]
  !gapprompt@gap>| !gapinput@S:=FullTransformationSemigroup(5);;|
  !gapprompt@gap>| !gapinput@D:=GreensDClasses(S)[3];|
  {Transformation( [ 1, 1, 1, 2, 3 ] )}
  !gapprompt@gap>| !gapinput@AssociatedReesMatrixSemigroupOfDClass(D);|
  <Rees 0-matrix semigroup 25x10 over Group([ (1,2)(3,5)(4,6), (1,3)
  (2,4)(5,6) ])>
\end{Verbatim}
 }

 }

  }

            
\chapter{\textcolor{Chapter }{Finitely Presented Semigroups and Monoids}}\label{Finitely Presented Semigroups and Monoids}
\logpage{[ 52, 0, 0 ]}
\hyperdef{L}{X7DE7C52A7C4BDADE}{}
{
  A \emph{finitely presented semigroup} (resp. \emph{finitely presented monoid}) is a quotient of a free semigroup (resp. free monoid) on a finite number of
generators over a finitely generated congruence on the free semigroup (resp.
free monoid). 

 Finitely presented semigroups are obtained by factoring a free semigroup by a
set of relations (a generating set for the congruence), i.e., a set of pairs
of words in the free semigroup. 

 
\begin{Verbatim}[commandchars=!@|,fontsize=\small,frame=single,label=Example]
  !gapprompt@gap>| !gapinput@f := FreeSemigroup( "a", "b" );;|
  !gapprompt@gap>| !gapinput@x := GeneratorsOfSemigroup( f );;|
  !gapprompt@gap>| !gapinput@s := f / [ [ x[1] * x[2], x[2] * x[1] ] ];|
  <fp semigroup on the generators [ a, b ]>
  !gapprompt@gap>| !gapinput@GeneratorsOfSemigroup( s );|
  [ a, b ]
  !gapprompt@gap>| !gapinput@RelationsOfFpSemigroup( s );|
  [ [ a*b, b*a ] ]
\end{Verbatim}
 

 Finitely presented monoids are obtained by factoring a free monoid by a set of
relations, i.e. a set of pairs of words in the free monoid. 

 
\begin{Verbatim}[commandchars=!@|,fontsize=\small,frame=single,label=Example]
  !gapprompt@gap>| !gapinput@f := FreeMonoid( "a", "b" );;|
  !gapprompt@gap>| !gapinput@x := GeneratorsOfMonoid( f );|
  [ a, b ]
  !gapprompt@gap>| !gapinput@e := Identity( f );|
  <identity ...>
  !gapprompt@gap>| !gapinput@m := f / [ [ x[1] * x[2], e ] ];|
  <fp monoid on the generators [ a, b ]>
  !gapprompt@gap>| !gapinput@RelationsOfFpMonoid( m );|
  [ [ a*b, <identity ...> ] ]
\end{Verbatim}
 

 Notice that for \textsf{GAP} a finitely presented monoid is not a finitely presented semigroup. 

 
\begin{Verbatim}[commandchars=!@|,fontsize=\small,frame=single,label=Example]
  !gapprompt@gap>| !gapinput@IsFpSemigroup( m );|
  false
\end{Verbatim}
 

 However, one can build a finitely presented semigroup isomorphic to that
finitely presented monoid (see \texttt{IsomorphismFpSemigroup} (\ref{IsomorphismFpSemigroup})). 

 Also note that is not possible to refer to the generators by their names.
These names are not variables, but just display figures. So, if one wants to
access the generators by their names, one first has to introduce the
respective variables and to assign the generators to them. 

 
\begin{Verbatim}[commandchars=!@|,fontsize=\small,frame=single,label=Example]
  !gapprompt@gap>| !gapinput@Unbind( a );|
  !gapprompt@gap>| !gapinput@f := FreeSemigroup( "a", "b" );;|
  !gapprompt@gap>| !gapinput@s := f / [ [ f.1 * f.2, f.2 * f.1 ] ];;|
  !gapprompt@gap>| !gapinput@a;|
  Error, Variable: 'a' must have a value
  !gapprompt@gap>| !gapinput@a := s.1;|
  a
  !gapprompt@gap>| !gapinput@a in f;|
  false
  !gapprompt@gap>| !gapinput@a in s;|
  true
\end{Verbatim}
 

 The generators of the free semigroup (resp. free monoid) are different from
the generators of the finitely presented semigroup (resp. finitely presented
monoid) (even though they are displayed by the same names). This means that
words in the generators of the free semigroup (resp. free monoid) are not
elements of the finitely presented semigroup (resp. finitely presented
monoid). Conversely elements of the finitely presented semigroup (resp.
finitely presented monoid) are not words of the free semigroup (resp. free
monoid). 

 Calculations comparing elements of a finitely presented semigroup may run into
problems: there are finitely presented semigroups for which no algorithm
exists (it is known that no such algorithm can exist) that will tell for two
arbitrary words in the generators whether the corresponding elements in the
finitely presented semigroup are equal. Therefore the methods used by \textsf{GAP} to compute in finitely presented semigroups may run into warning errors, run
out of memory or run forever. If the finitely presented semigroup is (by
theory) known to be finite the algorithms are guaranteed to terminate (if
there is sufficient memory available), but the time needed for the calculation
cannot be bounded a priori. The same can be said for monoids. (See \ref{Rewriting Systems and the Knuth-Bendix Procedure}.) 

 
\begin{Verbatim}[commandchars=!@|,fontsize=\small,frame=single,label=Example]
  !gapprompt@gap>| !gapinput@b := s.2;;|
  !gapprompt@gap>| !gapinput@a*b = a^5;|
  false
  !gapprompt@gap>| !gapinput@a^5 * b^2 * a = a^6 * b^2;|
  true
\end{Verbatim}
 

 Note that elements of a finitely presented semigroup (or monoid) are not
printed in a unique way: 

 
\begin{Verbatim}[commandchars=!@|,fontsize=\small,frame=single,label=Example]
  !gapprompt@gap>| !gapinput@a^5 * b^2 * a;|
  a^5*b^2*a
  !gapprompt@gap>| !gapinput@a^6 * b^2;|
  a^6*b^2
\end{Verbatim}
  
\section{\textcolor{Chapter }{IsSubsemigroupFpSemigroup (Filter)}}\label{sect:IsSubsemigroupFpSemigroup}
\logpage{[ 52, 1, 0 ]}
\hyperdef{L}{X78C80F1A84C58E1E}{}
{
  

\subsection{\textcolor{Chapter }{IsSubsemigroupFpSemigroup}}
\logpage{[ 52, 1, 1 ]}\nobreak
\hyperdef{L}{X8496E23C80453C33}{}
{\noindent\textcolor{FuncColor}{$\triangleright$\enspace\texttt{IsSubsemigroupFpSemigroup({\mdseries\slshape t})\index{IsSubsemigroupFpSemigroup@\texttt{IsSubsemigroupFpSemigroup}}
\label{IsSubsemigroupFpSemigroup}
}\hfill{\scriptsize (filter)}}\\
\noindent\textcolor{FuncColor}{$\triangleright$\enspace\texttt{IsSubmonoidFpMonoid({\mdseries\slshape t})\index{IsSubmonoidFpMonoid@\texttt{IsSubmonoidFpMonoid}}
\label{IsSubmonoidFpMonoid}
}\hfill{\scriptsize (filter)}}\\


 The first function returns true if \mbox{\texttt{\mdseries\slshape t}} is a finitely presented semigroup or a subsemigroup of a finitely presented
semigroup. The second function does the equivalent thing for monoids.
(Generally speaking, such a subsemigroup or monoid can be constructed with \texttt{Semigroup(\mbox{\texttt{\mdseries\slshape gens}})} or \texttt{Monoid(\mbox{\texttt{\mdseries\slshape gens}})}, where \mbox{\texttt{\mdseries\slshape gens}} is a list of elements of a finitely presented semigroup or monoid.) 

 A submonoid of a monoid has the same identity as the monoid. }

 

\subsection{\textcolor{Chapter }{IsFpSemigroup}}
\logpage{[ 52, 1, 2 ]}\nobreak
\hyperdef{L}{X8239EF2B853411E9}{}
{\noindent\textcolor{FuncColor}{$\triangleright$\enspace\texttt{IsFpSemigroup({\mdseries\slshape s})\index{IsFpSemigroup@\texttt{IsFpSemigroup}}
\label{IsFpSemigroup}
}\hfill{\scriptsize (filter)}}\\
\noindent\textcolor{FuncColor}{$\triangleright$\enspace\texttt{IsFpMonoid({\mdseries\slshape m})\index{IsFpMonoid@\texttt{IsFpMonoid}}
\label{IsFpMonoid}
}\hfill{\scriptsize (filter)}}\\


 The first function is a synonym for \texttt{IsSubsemigroupFpSemigroup(\mbox{\texttt{\mdseries\slshape s}})} and \texttt{IsWholeFamily(\mbox{\texttt{\mdseries\slshape s}})} (this is because a subsemigroup of a finitely presented semigroup is not
necessarily finitely presented). 

 Similarly, the second function is a synonym for \texttt{IsSubmonoidFpMonoid(\mbox{\texttt{\mdseries\slshape m}})} and \texttt{IsWholeFamily(\mbox{\texttt{\mdseries\slshape m}})}. }

 

\subsection{\textcolor{Chapter }{IsElementOfFpSemigroup}}
\logpage{[ 52, 1, 3 ]}\nobreak
\hyperdef{L}{X81ABBE997A4C19B7}{}
{\noindent\textcolor{FuncColor}{$\triangleright$\enspace\texttt{IsElementOfFpSemigroup({\mdseries\slshape elm})\index{IsElementOfFpSemigroup@\texttt{IsElementOfFpSemigroup}}
\label{IsElementOfFpSemigroup}
}\hfill{\scriptsize (Category)}}\\
\noindent\textcolor{FuncColor}{$\triangleright$\enspace\texttt{IsElementOfFpMonoid({\mdseries\slshape elm})\index{IsElementOfFpMonoid@\texttt{IsElementOfFpMonoid}}
\label{IsElementOfFpMonoid}
}\hfill{\scriptsize (Category)}}\\


 returns true if \mbox{\texttt{\mdseries\slshape elm}} is an element of a finitely presented semigroup or monoid. 

 
\begin{Verbatim}[commandchars=!@|,fontsize=\small,frame=single,label=Example]
  !gapprompt@gap>| !gapinput@f := FreeSemigroup( "a", "b" );;|
  !gapprompt@gap>| !gapinput@IsFpSemigroup( f );|
  false
  !gapprompt@gap>| !gapinput@s := f / [ [ f.1^2, f.2^2 ] ];;|
  !gapprompt@gap>| !gapinput@IsFpSemigroup( s );|
  true
  !gapprompt@gap>| !gapinput@t := Semigroup( [ s.1^2 ] );|
  <commutative semigroup with 1 generator>
  !gapprompt@gap>| !gapinput@IsSubsemigroupFpSemigroup( t );|
  true
  !gapprompt@gap>| !gapinput@IsSubsemigroupFpSemigroup( s );|
  true
  !gapprompt@gap>| !gapinput@IsSubsemigroupFpSemigroup( f );|
  false
  !gapprompt@gap>| !gapinput@IsElementOfFpSemigroup( t.1^3 );|
  true
\end{Verbatim}
 }

 

\subsection{\textcolor{Chapter }{FpGrpMonSmgOfFpGrpMonSmgElement}}
\logpage{[ 52, 1, 4 ]}\nobreak
\hyperdef{L}{X7DC8A5D380AFE5DB}{}
{\noindent\textcolor{FuncColor}{$\triangleright$\enspace\texttt{FpGrpMonSmgOfFpGrpMonSmgElement({\mdseries\slshape elm})\index{FpGrpMonSmgOfFpGrpMonSmgElement@\texttt{FpGrpMonSmgOfFpGrpMonSmgElement}}
\label{FpGrpMonSmgOfFpGrpMonSmgElement}
}\hfill{\scriptsize (operation)}}\\


 returns the finitely presented group, monoid or semigroup to which \mbox{\texttt{\mdseries\slshape elm}} belongs. 

 
\begin{Verbatim}[commandchars=!@|,fontsize=\small,frame=single,label=Example]
  !gapprompt@gap>| !gapinput@s = FpGrpMonSmgOfFpGrpMonSmgElement( s.1 );|
  true
  !gapprompt@gap>| !gapinput@s = FpGrpMonSmgOfFpGrpMonSmgElement( t.1 );|
  true
  !gapprompt@gap>| !gapinput@f := FreeMonoid( 2 );;|
  !gapprompt@gap>| !gapinput@m := f / [ [ f.1^2, f.2^2 ] ];|
  <fp monoid on the generators [ m1, m2 ]>
  !gapprompt@gap>| !gapinput@m = FpGrpMonSmgOfFpGrpMonSmgElement( m.1 * m.2 );|
  true
\end{Verbatim}
 }

 }

  
\section{\textcolor{Chapter }{Creating Finitely Presented Semigroups and Monoids}}\label{Creating Finitely Presented Semigroups and Monoids}
\logpage{[ 52, 2, 0 ]}
\hyperdef{L}{X865E230B83982E66}{}
{
  

\subsection{\textcolor{Chapter }{\texttt{\symbol{92}}/ (for a free semigroup or monoid and a list of pairs of elements)}}
\logpage{[ 52, 2, 1 ]}\nobreak
\hyperdef{L}{X84745EC6789FEB4C}{}
{\noindent\textcolor{FuncColor}{$\triangleright$\enspace\texttt{\texttt{\symbol{92}}/({\mdseries\slshape F, rels})\index{\texttt{\symbol{92}}/@\texttt{\texttt{\symbol{92}}/}!for a free semigroup or monoid and a list of pairs of elements}
\label{bSlash/:for a free semigroup or monoid and a list of pairs of elements}
}\hfill{\scriptsize (method)}}\\


 \index{quotient!of free semigroup} \index{quotient!of free monoid} creates a finitely presented semigroup or monoid given by the presentation $\langle \mbox{\texttt{\mdseries\slshape gens}} \mid \mbox{\texttt{\mdseries\slshape rels}} \rangle$ where \mbox{\texttt{\mdseries\slshape gens}} are the generators of the free semigroup or free monoid \mbox{\texttt{\mdseries\slshape F}}, and the relations \mbox{\texttt{\mdseries\slshape rels}} are entered as pairs of words in the generators of the free semigroup or free
monoid. 

 The same result is obtained with the infix operator \texttt{/}, i.e. as \mbox{\texttt{\mdseries\slshape F}} \texttt{/} \mbox{\texttt{\mdseries\slshape rels}}. 

 
\begin{Verbatim}[commandchars=!@|,fontsize=\small,frame=single,label=Example]
  !gapprompt@gap>| !gapinput@fs := FreeSemigroup( 3 );;|
  !gapprompt@gap>| !gapinput@x := GeneratorsOfSemigroup( fs );;|
  !gapprompt@gap>| !gapinput@s := fs / [ [ x[1] * x[2] * x[1], x[1] ], [ x[2]^4, x[1] ] ];|
  <fp semigroup on the generators [ s1, s2, s3 ]>
\end{Verbatim}
 }

 

\subsection{\textcolor{Chapter }{FactorFreeSemigroupByRelations}}
\logpage{[ 52, 2, 2 ]}\nobreak
\hyperdef{L}{X822F04B2833BE254}{}
{\noindent\textcolor{FuncColor}{$\triangleright$\enspace\texttt{FactorFreeSemigroupByRelations({\mdseries\slshape f, rels})\index{FactorFreeSemigroupByRelations@\texttt{FactorFreeSemigroupByRelations}}
\label{FactorFreeSemigroupByRelations}
}\hfill{\scriptsize (function)}}\\
\noindent\textcolor{FuncColor}{$\triangleright$\enspace\texttt{FactorFreeMonoidByRelations({\mdseries\slshape f, rels})\index{FactorFreeMonoidByRelations@\texttt{FactorFreeMonoidByRelations}}
\label{FactorFreeMonoidByRelations}
}\hfill{\scriptsize (function)}}\\


 for a free semigroup or free monoid \mbox{\texttt{\mdseries\slshape f}} and a list \mbox{\texttt{\mdseries\slshape rels}} of pairs of elements of \mbox{\texttt{\mdseries\slshape f}}. Returns the finitely presented semigroup or monoid which is the quotient of \mbox{\texttt{\mdseries\slshape f}} by the least congruence on \mbox{\texttt{\mdseries\slshape f}} generated by the pairs in \mbox{\texttt{\mdseries\slshape rels}}. 

 Users should be aware that much of the code described in this chapter is in
need of substantial revision. In particular, the two functions described here
are \emph{not} called by the operation \texttt{\texttt{\symbol{92}}/} of the previous subsection, and so are liable to be removed in due course. 

 
\begin{Verbatim}[commandchars=!@|,fontsize=\small,frame=single,label=Example]
  !gapprompt@gap>| !gapinput@fm := FreeMonoid( 3 );;|
  !gapprompt@gap>| !gapinput@y := GeneratorsOfMonoid( fm );;|
  !gapprompt@gap>| !gapinput@m := FactorFreeMonoidByRelations( fm,|
  !gapprompt@>| !gapinput@          [ [ y[1] * y[2] * y[1], y[1] ],[ y[2]^4, y[1] ] ] );|
  <fp monoid on the generators [ m1, m2, m3 ]>
\end{Verbatim}
 }

 

\subsection{\textcolor{Chapter }{IsomorphismFpSemigroup}}
\logpage{[ 52, 2, 3 ]}\nobreak
\hyperdef{L}{X869F966B8196F28C}{}
{\noindent\textcolor{FuncColor}{$\triangleright$\enspace\texttt{IsomorphismFpSemigroup({\mdseries\slshape m})\index{IsomorphismFpSemigroup@\texttt{IsomorphismFpSemigroup}}
\label{IsomorphismFpSemigroup}
}\hfill{\scriptsize (attribute)}}\\
\noindent\textcolor{FuncColor}{$\triangleright$\enspace\texttt{IsomorphismFpMonoid({\mdseries\slshape g})\index{IsomorphismFpMonoid@\texttt{IsomorphismFpMonoid}}
\label{IsomorphismFpMonoid}
}\hfill{\scriptsize (attribute)}}\\


 for a finitely presented monoid \mbox{\texttt{\mdseries\slshape m}} or a finitely presented group \mbox{\texttt{\mdseries\slshape g}}. Returns an isomorphism from \mbox{\texttt{\mdseries\slshape m}} or \mbox{\texttt{\mdseries\slshape g}} to a finitely presented semigroup or monoid. 

 
\begin{Verbatim}[commandchars=!@|,fontsize=\small,frame=single,label=Example]
  !gapprompt@gap>| !gapinput@phis := IsomorphismFpSemigroup( m );|
  MappingByFunction( <fp monoid on the generators
  [ m1, m2, m3 ]>, <fp semigroup on the generators [ <identity ...>, m1, m2, m3
   ]>, function( x ) ... end, function( x ) ... end )
  !gapprompt@gap>| !gapinput@fg := FreeGroup( 2 );;|
  !gapprompt@gap>| !gapinput@g := fg / [ fg.1^4, fg.2^5 ];|
  <fp group on the generators [ f1, f2 ]>
  !gapprompt@gap>| !gapinput@phim := IsomorphismFpMonoid( g );|
  MappingByFunction( <fp group on the generators
  [ f1, f2 ]>, <fp monoid on the generators [ f1, f1^-1, f2, f2^-1
   ]>, function( x ) ... end, function( x ) ... end )
\end{Verbatim}
 }

 }

  
\section{\textcolor{Chapter }{Comparison of Elements of Finitely Presented Semigroups}}\label{Comparison of Elements of Finitely Presented Semigroups}
\logpage{[ 52, 3, 0 ]}
\hyperdef{L}{X85E7C8407C9D5FBE}{}
{
  

\subsection{\textcolor{Chapter }{\texttt{\symbol{92}}= (for two elements in a f.p. semigroup)}}
\logpage{[ 52, 3, 1 ]}\nobreak
\hyperdef{L}{X7DD9D81F863EBE31}{}
{\noindent\textcolor{FuncColor}{$\triangleright$\enspace\texttt{\texttt{\symbol{92}}=({\mdseries\slshape a, b})\index{\texttt{\symbol{92}}=@\texttt{\texttt{\symbol{92}}=}!for two elements in a f.p. semigroup}
\label{bSlash=:for two elements in a f.p. semigroup}
}\hfill{\scriptsize (method)}}\\


 \index{comparison!fp semigroup elements} Two elements \mbox{\texttt{\mdseries\slshape a}}, \mbox{\texttt{\mdseries\slshape b}} of a finitely presented semigroup are equal if they are equal in the
semigroup. Nevertheless they may be represented as different words in the
generators. Because of the fundamental problems mentioned in the introduction
to this chapter such a test may take a very long time and cannot be guaranteed
to finish (see \ref{Rewriting Systems and the Knuth-Bendix Procedure}). }

 }

  
\section{\textcolor{Chapter }{Preimages in the Free Semigroup or Monoid}}\label{Preimages in the Free Semigroup or Monoid}
\logpage{[ 52, 4, 0 ]}
\hyperdef{L}{X7CD806CA7E0A1438}{}
{
  Elements of a finitely presented semigroup or monoid are not words, but are
represented using a word from the free semigroup or free monoid as
representative. 

 

\subsection{\textcolor{Chapter }{UnderlyingElement (of an element in a fp semigroup or monoid)}}
\logpage{[ 52, 4, 1 ]}\nobreak
\hyperdef{L}{X784B3DB686E7080C}{}
{\noindent\textcolor{FuncColor}{$\triangleright$\enspace\texttt{UnderlyingElement({\mdseries\slshape elm})\index{UnderlyingElement@\texttt{UnderlyingElement}!of an element in a fp semigroup or monoid}
\label{UnderlyingElement:of an element in a fp semigroup or monoid}
}\hfill{\scriptsize (operation)}}\\


 for an element \mbox{\texttt{\mdseries\slshape elm}} of a finitely presented semigroup or monoid. Returns the word from the free
semigroup or free monoid that is used as a representative for \mbox{\texttt{\mdseries\slshape elm}}. 

 
\begin{Verbatim}[commandchars=!@|,fontsize=\small,frame=single,label=Example]
  !gapprompt@gap>| !gapinput@genm := GeneratorsOfMonoid( m );;|
  !gapprompt@gap>| !gapinput@e := genm[2]^4;|
  m2^4
  !gapprompt@gap>| !gapinput@IsWord( e );|
  false
  !gapprompt@gap>| !gapinput@ue := UnderlyingElement( e );|
  m2^4
  !gapprompt@gap>| !gapinput@IsWord( ue );|
  true
\end{Verbatim}
 }

 

\subsection{\textcolor{Chapter }{ElementOfFpSemigroup}}
\logpage{[ 52, 4, 2 ]}\nobreak
\hyperdef{L}{X847012347856C55E}{}
{\noindent\textcolor{FuncColor}{$\triangleright$\enspace\texttt{ElementOfFpSemigroup({\mdseries\slshape fam, word})\index{ElementOfFpSemigroup@\texttt{ElementOfFpSemigroup}}
\label{ElementOfFpSemigroup}
}\hfill{\scriptsize (operation)}}\\
\noindent\textcolor{FuncColor}{$\triangleright$\enspace\texttt{ElementOfFpMonoid({\mdseries\slshape fam, word})\index{ElementOfFpMonoid@\texttt{ElementOfFpMonoid}}
\label{ElementOfFpMonoid}
}\hfill{\scriptsize (operation)}}\\


 for a family \mbox{\texttt{\mdseries\slshape fam}} of elements of a finitely presented semigroup or monoid and a word \mbox{\texttt{\mdseries\slshape word}} in the free generators underlying this finitely presented semigroup or monoid.
Returns the element of the finitely presented semigroup or monoid with the
representative \mbox{\texttt{\mdseries\slshape word}} in the free semigroup or free monoid. These operations are inverse to \texttt{UnderlyingElement}. 

 
\begin{Verbatim}[commandchars=!@|,fontsize=\small,frame=single,label=Example]
  !gapprompt@gap>| !gapinput@fam := FamilyObj( genm[1] );;|
  !gapprompt@gap>| !gapinput@w := y[1]^3 * y[2]^4 * y[3]^5;|
  m1^3*m2^4*m3^5
  !gapprompt@gap>| !gapinput@ew := ElementOfFpMonoid( fam, w );|
  m1^3*m2^4*m3^5
  !gapprompt@gap>| !gapinput@ew in fm;|
  false
  !gapprompt@gap>| !gapinput@ew in m;|
  true
  !gapprompt@gap>| !gapinput@w = UnderlyingElement( ew );|
  true
\end{Verbatim}
 }

 

\subsection{\textcolor{Chapter }{FreeSemigroupOfFpSemigroup}}
\logpage{[ 52, 4, 3 ]}\nobreak
\hyperdef{L}{X8726523779601873}{}
{\noindent\textcolor{FuncColor}{$\triangleright$\enspace\texttt{FreeSemigroupOfFpSemigroup({\mdseries\slshape s})\index{FreeSemigroupOfFpSemigroup@\texttt{FreeSemigroupOfFpSemigroup}}
\label{FreeSemigroupOfFpSemigroup}
}\hfill{\scriptsize (attribute)}}\\
\noindent\textcolor{FuncColor}{$\triangleright$\enspace\texttt{FreeMonoidOfFpMonoid({\mdseries\slshape m})\index{FreeMonoidOfFpMonoid@\texttt{FreeMonoidOfFpMonoid}}
\label{FreeMonoidOfFpMonoid}
}\hfill{\scriptsize (attribute)}}\\


 returns the underlying free semigroup or free monoid for the finitely
presented semigroup \mbox{\texttt{\mdseries\slshape s}} or monoid \mbox{\texttt{\mdseries\slshape m}}, i.e. the free semigroup or free monoid over which \mbox{\texttt{\mdseries\slshape s}} or \mbox{\texttt{\mdseries\slshape m}} is defined as a quotient. (This is the free semigroup or free monoid generated
by the free generators provided by \texttt{FreeGeneratorsOfFpSemigroup(\mbox{\texttt{\mdseries\slshape s}})} or \texttt{FreeGeneratorsOfFpMonoid(\mbox{\texttt{\mdseries\slshape m}})}). }

 

\subsection{\textcolor{Chapter }{FreeGeneratorsOfFpSemigroup}}
\logpage{[ 52, 4, 4 ]}\nobreak
\hyperdef{L}{X79A39402806B5EB7}{}
{\noindent\textcolor{FuncColor}{$\triangleright$\enspace\texttt{FreeGeneratorsOfFpSemigroup({\mdseries\slshape s})\index{FreeGeneratorsOfFpSemigroup@\texttt{FreeGeneratorsOfFpSemigroup}}
\label{FreeGeneratorsOfFpSemigroup}
}\hfill{\scriptsize (attribute)}}\\
\noindent\textcolor{FuncColor}{$\triangleright$\enspace\texttt{FreeGeneratorsOfFpMonoid({\mdseries\slshape m})\index{FreeGeneratorsOfFpMonoid@\texttt{FreeGeneratorsOfFpMonoid}}
\label{FreeGeneratorsOfFpMonoid}
}\hfill{\scriptsize (attribute)}}\\


 returns the underlying free generators corresponding to the generators of the
finitely presented semigroup \mbox{\texttt{\mdseries\slshape s}} or monoid \mbox{\texttt{\mdseries\slshape m}}. }

 

\subsection{\textcolor{Chapter }{RelationsOfFpSemigroup}}
\logpage{[ 52, 4, 5 ]}\nobreak
\hyperdef{L}{X862BE9FA7C987CAB}{}
{\noindent\textcolor{FuncColor}{$\triangleright$\enspace\texttt{RelationsOfFpSemigroup({\mdseries\slshape s})\index{RelationsOfFpSemigroup@\texttt{RelationsOfFpSemigroup}}
\label{RelationsOfFpSemigroup}
}\hfill{\scriptsize (attribute)}}\\
\noindent\textcolor{FuncColor}{$\triangleright$\enspace\texttt{RelationsOfFpMonoid({\mdseries\slshape m})\index{RelationsOfFpMonoid@\texttt{RelationsOfFpMonoid}}
\label{RelationsOfFpMonoid}
}\hfill{\scriptsize (attribute)}}\\


 returns the relations of the finitely presented semigroup \mbox{\texttt{\mdseries\slshape s}} or monoid \mbox{\texttt{\mdseries\slshape m}} as pairs of words in the free generators provided by \texttt{FreeGeneratorsOfFpSemigroup(\mbox{\texttt{\mdseries\slshape s}})} or \texttt{FreeGeneratorsOfFpMonoid(\mbox{\texttt{\mdseries\slshape m}})}. 

 
\begin{Verbatim}[commandchars=!@|,fontsize=\small,frame=single,label=Example]
  !gapprompt@gap>| !gapinput@fs = FreeSemigroupOfFpSemigroup( s );|
  true
  !gapprompt@gap>| !gapinput@FreeGeneratorsOfFpMonoid( m );|
  [ m1, m2, m3 ]
  !gapprompt@gap>| !gapinput@RelationsOfFpSemigroup( s );|
  [ [ s1*s2*s1, s1 ], [ s2^4, s1 ] ]
\end{Verbatim}
 }

 }

  
\section{\textcolor{Chapter }{Rewriting Systems and the Knuth\texttt{\symbol{45}}Bendix Procedure}}\label{Rewriting Systems and the Knuth-Bendix Procedure}
\logpage{[ 52, 5, 0 ]}
\hyperdef{L}{X87693BDC79DC6EBF}{}
{
  If a finitely presented semigroup (or monoid) has a confluent rewriting system
then it has a solvable word problem, that is, there is an algorithm to decide
when two words in the free underlying semigroup represent the same element of
the finitely presented semigroup. Indeed, once we have a confluent rewriting
system, it is possible to successfully test that two words represent the same
element in the semigroup, by reducing both words using the rewriting system
rules. This is, at the moment, the method that \textsf{GAP} uses to check equality in finitely presented semigroups and monoids. 

 

\subsection{\textcolor{Chapter }{ReducedConfluentRewritingSystem}}
\logpage{[ 52, 5, 1 ]}\nobreak
\hyperdef{L}{X7D8F804E814D894D}{}
{\noindent\textcolor{FuncColor}{$\triangleright$\enspace\texttt{ReducedConfluentRewritingSystem({\mdseries\slshape S[, ordering]})\index{ReducedConfluentRewritingSystem@\texttt{ReducedConfluentRewritingSystem}}
\label{ReducedConfluentRewritingSystem}
}\hfill{\scriptsize (attribute)}}\\


 returns a reduced confluent rewriting system of the finitely presented
semigroup or monoid \mbox{\texttt{\mdseries\slshape S}} with respect to the reduction ordering \mbox{\texttt{\mdseries\slshape ordering}} (see \ref{Orderings}). 

 The default for \mbox{\texttt{\mdseries\slshape ordering}} is the length plus lexicographic ordering on words, also called the shortlex
ordering; for the definition see for example \cite{Sims94}. 

 Notice that this might not terminate. In particular, if the semigroup or
monoid \mbox{\texttt{\mdseries\slshape S}} does not have a solvable word problem then it this will certainly never end.
Also, in this case, the object returned is an immutable rewriting system,
because once we have a confluent rewriting system for a finitely presented
semigroup or monoid we do not want to allow it to change (as it was most
probably very time consuming to get it in the first place). Furthermore, this
is also an attribute storing object (see \ref{Representation}). 

 
\begin{Verbatim}[commandchars=@|A,fontsize=\small,frame=single,label=Example]
  @gapprompt|gap>A @gapinput|f := FreeSemigroup( "a", "b" );;A
  @gapprompt|gap>A @gapinput|a := f.1;;  b := f.2;;A
  @gapprompt|gap>A @gapinput|s := f / [ [ a*b*a, b ], [ b*a*b, a ] ];;A
  @gapprompt|gap>A @gapinput|rws := ReducedConfluentRewritingSystem( s );A
  Rewriting System for Semigroup( [ a, b ] ) with rules
  [ [ a*b*a, b ], [ b*a*b, a ], [ b*a^2, a^2*b ], [ b^2, a^2 ],
    [ a^5, a ], [ a^3*b, b*a ] ]
  @gapprompt|gap>A @gapinput|c := s.1;;  d := s.2;;A
  @gapprompt|gap>A @gapinput|e := (c*d^2)^3;A
  (a*b^2)^3
  @gapprompt|gap>A @gapinput|## ReducedForm( rws, e );  gives an error!A
  @gapprompt|gap>A @gapinput|w := UnderlyingElement( e );A
  (a*b^2)^3
  @gapprompt|gap>A @gapinput|ReducedForm( rws, w );A
  a
\end{Verbatim}
 

 The creation of a reduced confluent rewriting system for a semigroup or for a
monoid, in \textsf{GAP}, uses the Knuth\texttt{\symbol{45}}Bendix procedure for strings, which
manipulates a rewriting system of the semigroup or monoid and attempts to make
it confluent, (see Chapter \ref{Rewriting Systems} and also Sims \cite{Sims94}). (Since the word problem for semigroups/monoids is not solvable in general,
the Knuth\texttt{\symbol{45}}Bendix procedure cannot always terminate). 

 In order to apply this procedure we will build a rewriting system for the
semigroup or monoid, which we will call a \emph{Knuth\texttt{\symbol{45}}Bendix Rewriting System} (we need to define this because we need the rewriting system to store some
information needed for the implementation of the
Knuth\texttt{\symbol{45}}Bendix procedure). 

 Actually, Knuth\texttt{\symbol{45}}Bendix Rewriting Systems do not only serve
this purpose. Indeed these are objects which are mutable and which can be
manipulated (see \ref{Rewriting Systems}). 

 Note that the implemented version of the Knuth\texttt{\symbol{45}}Bendix
procedure, in \textsf{GAP} returns, if it terminates, a confluent rewriting system which is reduced.
Also, a reduction ordering has to be specified when building a rewriting
system. If none is specified, the shortlex ordering is assumed (note that the
procedure may terminate with a certain ordering and not with another one). 

 On Unix systems it is possible to replace the built\texttt{\symbol{45}}in
Knuth\texttt{\symbol{45}}Bendix by other routines, for example the package \textsf{kbmag} offers such a possibility. }

 

\subsection{\textcolor{Chapter }{KB{\textunderscore}REW}}
\logpage{[ 52, 5, 2 ]}\nobreak
\hyperdef{L}{X7A3F8AE285C41D80}{}
{\noindent\textcolor{FuncColor}{$\triangleright$\enspace\texttt{KB{\textunderscore}REW\index{KB{\textunderscore}REW@\texttt{KB{\textunderscore}REW}}
\label{KBuScoreREW}
}\hfill{\scriptsize (global variable)}}\\
\noindent\textcolor{FuncColor}{$\triangleright$\enspace\texttt{GAPKB{\textunderscore}REW\index{GAPKB{\textunderscore}REW@\texttt{GAPKB{\textunderscore}REW}}
\label{GAPKBuScoreREW}
}\hfill{\scriptsize (global variable)}}\\


 \texttt{KB{\textunderscore}REW} is a global record variable whose components contain functions used for
Knuth\texttt{\symbol{45}}Bendix. By default \texttt{KB{\textunderscore}REW} is assigned to \texttt{GAPKB{\textunderscore}REW}, which contains the KB functions provided by the GAP library. }

 
\subsection{\textcolor{Chapter }{KnuthBendixRewritingSystem}}\logpage{[ 52, 5, 3 ]}
\hyperdef{L}{X87A3823483E4FF86}{}
{
\noindent\textcolor{FuncColor}{$\triangleright$\enspace\texttt{KnuthBendixRewritingSystem({\mdseries\slshape s, wordord})\index{KnuthBendixRewritingSystem@\texttt{KnuthBendixRewritingSystem}!for a semigroup and a reduction ordering}
\label{KnuthBendixRewritingSystem:for a semigroup and a reduction ordering}
}\hfill{\scriptsize (operation)}}\\
\noindent\textcolor{FuncColor}{$\triangleright$\enspace\texttt{KnuthBendixRewritingSystem({\mdseries\slshape m, wordord})\index{KnuthBendixRewritingSystem@\texttt{KnuthBendixRewritingSystem}!for a monoid and a reduction ordering}
\label{KnuthBendixRewritingSystem:for a monoid and a reduction ordering}
}\hfill{\scriptsize (operation)}}\\


 in the first form, for a semigroup \mbox{\texttt{\mdseries\slshape s}} and a reduction ordering for the underlying free semigroup, it returns the
Knuth\texttt{\symbol{45}}Bendix rewriting system of the finitely presented
semigroup \mbox{\texttt{\mdseries\slshape s}} using the reduction ordering \mbox{\texttt{\mdseries\slshape wordord}}. In the second form, for a monoid \mbox{\texttt{\mdseries\slshape m}} and a reduction ordering for the underlying free monoid, it returns the
Knuth\texttt{\symbol{45}}Bendix rewriting system of the finitely presented
monoid \mbox{\texttt{\mdseries\slshape m}} using the reduction ordering \mbox{\texttt{\mdseries\slshape wordord}}. }

 

\subsection{\textcolor{Chapter }{SemigroupOfRewritingSystem}}
\logpage{[ 52, 5, 4 ]}\nobreak
\hyperdef{L}{X7966343587A04AFF}{}
{\noindent\textcolor{FuncColor}{$\triangleright$\enspace\texttt{SemigroupOfRewritingSystem({\mdseries\slshape rws})\index{SemigroupOfRewritingSystem@\texttt{SemigroupOfRewritingSystem}}
\label{SemigroupOfRewritingSystem}
}\hfill{\scriptsize (attribute)}}\\
\noindent\textcolor{FuncColor}{$\triangleright$\enspace\texttt{MonoidOfRewritingSystem({\mdseries\slshape rws})\index{MonoidOfRewritingSystem@\texttt{MonoidOfRewritingSystem}}
\label{MonoidOfRewritingSystem}
}\hfill{\scriptsize (attribute)}}\\


 returns the semigroup or monoid over which \mbox{\texttt{\mdseries\slshape rws}} is a rewriting system. }

 

\subsection{\textcolor{Chapter }{FreeSemigroupOfRewritingSystem}}
\logpage{[ 52, 5, 5 ]}\nobreak
\hyperdef{L}{X80B8115C8147F605}{}
{\noindent\textcolor{FuncColor}{$\triangleright$\enspace\texttt{FreeSemigroupOfRewritingSystem({\mdseries\slshape rws})\index{FreeSemigroupOfRewritingSystem@\texttt{FreeSemigroupOfRewritingSystem}}
\label{FreeSemigroupOfRewritingSystem}
}\hfill{\scriptsize (attribute)}}\\
\noindent\textcolor{FuncColor}{$\triangleright$\enspace\texttt{FreeMonoidOfRewritingSystem({\mdseries\slshape rws})\index{FreeMonoidOfRewritingSystem@\texttt{FreeMonoidOfRewritingSystem}}
\label{FreeMonoidOfRewritingSystem}
}\hfill{\scriptsize (attribute)}}\\


 returns the free semigroup or monoid over which \mbox{\texttt{\mdseries\slshape rws}} is a rewriting system. 

 
\begin{Verbatim}[commandchars=!@|,fontsize=\small,frame=single,label=Example]
  !gapprompt@gap>| !gapinput@f1 := FreeSemigroupOfRewritingSystem( rws );|
  <free semigroup on the generators [ a, b ]>
  !gapprompt@gap>| !gapinput@f1 = f;|
  true
  !gapprompt@gap>| !gapinput@s1 := SemigroupOfRewritingSystem( rws );|
  <fp semigroup on the generators [ a, b ]>
  !gapprompt@gap>| !gapinput@s1 = s;|
  true
\end{Verbatim}
 

 As mentioned before, having a confluent rewriting system, one can decide
whether two words represent the same element of a finitely presented semigroup
(or finitely presented monoid). 

 
\begin{Verbatim}[commandchars=!@|,fontsize=\small,frame=single,label=Example]
  !gapprompt@gap>| !gapinput@d^6 = c^2;|
  true
  !gapprompt@gap>| !gapinput@ReducedForm( rws, UnderlyingElement( d^6 ) );|
  a^2
  !gapprompt@gap>| !gapinput@ReducedForm( rws, UnderlyingElement( c^2 ) );|
  a^2
\end{Verbatim}
 }

 }

  
\section{\textcolor{Chapter }{Todd\texttt{\symbol{45}}Coxeter Procedure}}\label{Todd-Coxeter Procedure}
\logpage{[ 52, 6, 0 ]}
\hyperdef{L}{X812C28217F3E6720}{}
{
  This procedure gives a standard way of finding a transformation representation
of a finitely presented semigroup. Usually one does not explicitly call this
procedure but uses \texttt{IsomorphismTransformationSemigroup} (\ref{IsomorphismTransformationSemigroup}).  

\subsection{\textcolor{Chapter }{CosetTableOfFpSemigroup}}
\logpage{[ 52, 6, 1 ]}\nobreak
\hyperdef{L}{X7C24508A7F677520}{}
{\noindent\textcolor{FuncColor}{$\triangleright$\enspace\texttt{CosetTableOfFpSemigroup({\mdseries\slshape r})\index{CosetTableOfFpSemigroup@\texttt{CosetTableOfFpSemigroup}}
\label{CosetTableOfFpSemigroup}
}\hfill{\scriptsize (attribute)}}\\


 \mbox{\texttt{\mdseries\slshape r}} is a right congruence of an fp\texttt{\symbol{45}}semigroup \mbox{\texttt{\mdseries\slshape S}}. This attribute is the coset table of FP semigroup \mbox{\texttt{\mdseries\slshape S}} on a right congruence \mbox{\texttt{\mdseries\slshape r}}. Given a right congruence \mbox{\texttt{\mdseries\slshape r}} we represent \mbox{\texttt{\mdseries\slshape S}} as a set of transformations of the congruence classes of \mbox{\texttt{\mdseries\slshape r}}. 

 The images of the cosets under the generators are compiled in a list \mbox{\texttt{\mdseries\slshape table}} such that \mbox{\texttt{\mdseries\slshape table[i][s]}} contains the image of coset \mbox{\texttt{\mdseries\slshape s}} under generator \mbox{\texttt{\mdseries\slshape i}}. }

 }

 }

    
\chapter{\textcolor{Chapter }{Transformations}}\label{Transformations}
\logpage{[ 53, 0, 0 ]}
\hyperdef{L}{X860026B880BCB2A5}{}
{
  This chapter describes the functions in \textsf{GAP} for transformations. 

 A \emph{transformation} in \textsf{GAP} is simply a function from the positive integers to the positive integers.
Transformations are to semigroup theory what permutations are to group theory,
in the sense that every semigroup can be realised as a semigroup of
transformations. In \textsf{GAP} transformation semigroups are always finite, and so only finite semigroups can
be realised in this way. 

 A transformation in \textsf{GAP} acts on the positive integers (up to some architecture dependent limit) on the
right. The image of a point \texttt{i} under a transformation \texttt{f} is expressed as \texttt{i \texttt{\symbol{94}} f} in \textsf{GAP}. This action is also implemented by the function \texttt{OnPoints} (\ref{OnPoints}). If \texttt{i \texttt{\symbol{94}} f} is different from \texttt{i}, then \texttt{i} is \emph{moved} by \emph{f} and otherwise it is \emph{fixed} by \texttt{f}. Transformations in \textsf{GAP} are created using the operations described in Section \ref{sect:CreatingTransformations}. 

 The \emph{degree} of a transformation \texttt{f} is usually defined as the largest positive integer where \texttt{f} is defined. In previous versions of \textsf{GAP}, transformations were only defined on positive integers less than their
degree, it was only possible to multiply transformations of equal degree, and
a transformation did not act on any point exceeding its degree. Starting with
version 4.7 of \textsf{GAP}, transformations behave more like permutations, in that they fix unspecified
points and it is possible to multiply arbitrary transformations; see Chapter \ref{Permutations}. The definition of the degree of a transformation \texttt{f} in the current version of \textsf{GAP} is the largest value \texttt{n} such that \texttt{n \texttt{\symbol{94}} f {\textless}{\textgreater} n} or \texttt{i \texttt{\symbol{94}} f = n} for some \texttt{i {\textless}{\textgreater} n}. Equivalently, the degree of a transformation is the least value \texttt{n} such that \texttt{[ n + 1, n + 2, ... ]} is fixed pointwise by \texttt{f}. 

 The transformations of a given degree belong to the full transformation
semigroup of that degree; see \texttt{FullTransformationSemigroup} (\ref{FullTransformationSemigroup}). Transformation semigroups are hence subsemigroups of the full transformation
semigroup. 

 It is possible to use transformations in \textsf{GAP} without reference to the degree, much as it is possible to use permutations in
this way. However, for backwards compatibility, and because it is sometimes
useful, it is possible to access the degree of a transformation using \texttt{DegreeOfTransformation} (\ref{DegreeOfTransformation}). Certain attributes of transformations are also calculated with respect to
the degree, such as the rank, image set, or kernel (these values can also be
calculated with respect to any positive integer). So, it is possible to ignore
the degree of a transformation if you prefer to think of transformations as
acting on the positive integers in a similar way to permutations. For example,
this approach is used in the \textsf{FR} package. It is also possible to think of transformations as only acting on the
positive integers not exceeding their degree. For example, this was the
approach formerly used in \textsf{GAP} and it is also useful in the \textsf{Semigroups} package. 

 Transformations are displayed, by default, using the list \texttt{[ 1 \texttt{\symbol{94}} f .. n \texttt{\symbol{94}} f ]} where \texttt{n} is the degree of \texttt{f}. This behaviour differs from that of versions of \textsf{GAP} earlier than 4.7. See Section \ref{sect:DisplayingTransformations} for more information. 

 The \emph{rank} of a transformation on the positive integers up to \texttt{n} is the number of distinct points in \texttt{[ 1 \texttt{\symbol{94}} f .. n \texttt{\symbol{94}} f ]}. The \emph{kernel} of a transformation \texttt{f} on \texttt{[ 1 .. n ]} is the equivalence relation on \texttt{[ 1 .. n ]} consisting of those pairs \texttt{(i, j)} of positive integers such that \texttt{i \texttt{\symbol{94}} f = j \texttt{\symbol{94}} f}. The kernel of a transformation is represented in two ways: as a partition of \texttt{[ 1 .. n ]} or as the image list of a transformation \texttt{g} such that the kernel of \texttt{g} on \texttt{[ 1 .. n ]} equals the kernel of \texttt{f} and \texttt{j \texttt{\symbol{94}} g = i} for all \texttt{j} in \texttt{i}th class. The latter is referred to as the \emph{flat kernel} of \texttt{f}. For any given transformation \texttt{f} and value \texttt{n}, there is a unique transformation \texttt{g} with this property. 

 A \emph{functional digraph} is a directed graph where every vertex has out\texttt{\symbol{45}}degree $1$. A transformation \mbox{\texttt{\mdseries\slshape f}} can be thought of as a functional digraph with vertices the positive integers
and edges from \texttt{i} to \texttt{i \texttt{\symbol{94}} f} for every \texttt{i}. A \emph{component} of a transformation is defined as a component of the corresponding functional
digraph. More specifically, \texttt{i} and \texttt{j} are in the same component if and only if there are $i = v_0, v_1, \ldots, v_n = j$ such that either $v_{k+1}=v_{k}^f$ or $v_{k}=v_{k+1}^f$ for all $k$. A \emph{cycle} of a transformation is defined as a cycle (or strongly connected component) of
the corresponding functional digraph. More specifically, \texttt{i} belongs to a cycle of \mbox{\texttt{\mdseries\slshape f}} if there are $i=v_0, v_1, \ldots, v_n=i$ such that either $v_{k+1}=v_{k}^f$ or $v_{k}=v_{k+1}^f$ for all $k$. 

 Internally, \textsf{GAP} stores a transformation \texttt{f} as a list consisting of the images \texttt{i \texttt{\symbol{94}} f} for all values of \texttt{i} less than a value which is at least the degree of \texttt{f} and which is determined at the time of the creation of \texttt{f}. When the degree of a transformation \texttt{f} is at most 65536, the images of points under \texttt{f} are stored as 16\texttt{\symbol{45}}bit integers, the kernel and image set are
subobjects of \texttt{f} which are plain lists of \textsf{GAP} integers. When the degree of \texttt{f} is greater than 65536, the images of points under \texttt{f} are stored as 32\texttt{\symbol{45}}bit integers; the kernel and image set are
stored in the same way as before. A transformation belongs to \texttt{IsTrans2Rep} if it is stored using 16\texttt{\symbol{45}}bit integers and to \texttt{IsTrans4Rep} if it is stored using 32\texttt{\symbol{45}}bit integers. 

 
\section{\textcolor{Chapter }{The family and categories of transformations}}\logpage{[ 53, 1, 0 ]}
\hyperdef{L}{X7CF9291C7CC42340}{}
{
  

\subsection{\textcolor{Chapter }{IsTransformation}}
\logpage{[ 53, 1, 1 ]}\nobreak
\hyperdef{L}{X7B6259467974FB70}{}
{\noindent\textcolor{FuncColor}{$\triangleright$\enspace\texttt{IsTransformation({\mdseries\slshape obj})\index{IsTransformation@\texttt{IsTransformation}}
\label{IsTransformation}
}\hfill{\scriptsize (Category)}}\\


 Every transformation in \textsf{GAP} belongs to the category \texttt{IsTransformation}. Basic operations for transformations are \texttt{ImageListOfTransformation} (\ref{ImageListOfTransformation}), \texttt{ImageSetOfTransformation} (\ref{ImageSetOfTransformation}), \texttt{KernelOfTransformation} (\ref{KernelOfTransformation}), \texttt{FlatKernelOfTransformation} (\ref{FlatKernelOfTransformation}), \texttt{RankOfTransformation} (\ref{RankOfTransformation:for a transformation and a list}), \texttt{DegreeOfTransformation} (\ref{DegreeOfTransformation}), multiplication of two transformations via \texttt{*}, and exponentiation with the first argument a positive integer \texttt{i} and second argument a transformation \texttt{f} where the result is the image \texttt{i \texttt{\symbol{94}} f} of the point \texttt{i} under \texttt{f}. }

 

\subsection{\textcolor{Chapter }{IsTransformationCollection}}
\logpage{[ 53, 1, 2 ]}\nobreak
\hyperdef{L}{X7A6747CE85F2E6EA}{}
{\noindent\textcolor{FuncColor}{$\triangleright$\enspace\texttt{IsTransformationCollection({\mdseries\slshape obj})\index{IsTransformationCollection@\texttt{IsTransformationCollection}}
\label{IsTransformationCollection}
}\hfill{\scriptsize (Category)}}\\


 Every collection of transformations belongs to the category \texttt{IsTransformationCollection}. For example, transformation semigroups belong to \texttt{IsTransformationCollection}. }

 

\subsection{\textcolor{Chapter }{TransformationFamily}}
\logpage{[ 53, 1, 3 ]}\nobreak
\hyperdef{L}{X7E58AFA1832FF064}{}
{\noindent\textcolor{FuncColor}{$\triangleright$\enspace\texttt{TransformationFamily\index{TransformationFamily@\texttt{TransformationFamily}}
\label{TransformationFamily}
}\hfill{\scriptsize (family)}}\\


 The family of all transformations is \texttt{TransformationFamily}. }

 }

  
\section{\textcolor{Chapter }{Creating transformations}}\label{sect:CreatingTransformations}
\logpage{[ 53, 2, 0 ]}
\hyperdef{L}{X80F3086F87E93DF8}{}
{
  There are several ways of creating transformations in \textsf{GAP}, which are described in this section.  

\subsection{\textcolor{Chapter }{Transformation (for an image list)}}
\logpage{[ 53, 2, 1 ]}\nobreak
\hyperdef{L}{X86ADBDE57A20E323}{}
{\noindent\textcolor{FuncColor}{$\triangleright$\enspace\texttt{Transformation({\mdseries\slshape list})\index{Transformation@\texttt{Transformation}!for an image list}
\label{Transformation:for an image list}
}\hfill{\scriptsize (operation)}}\\
\noindent\textcolor{FuncColor}{$\triangleright$\enspace\texttt{Transformation({\mdseries\slshape list, func})\index{Transformation@\texttt{Transformation}!for a list and function}
\label{Transformation:for a list and function}
}\hfill{\scriptsize (operation)}}\\
\noindent\textcolor{FuncColor}{$\triangleright$\enspace\texttt{TransformationList({\mdseries\slshape list})\index{TransformationList@\texttt{TransformationList}!for an image list}
\label{TransformationList:for an image list}
}\hfill{\scriptsize (operation)}}\\
\textbf{\indent Returns:\ }
A transformation.



 \texttt{TransformationList} returns the transformation \texttt{f} such that \texttt{i \texttt{\symbol{94}} \mbox{\texttt{\mdseries\slshape f}} = \mbox{\texttt{\mdseries\slshape list}}[i]} if \texttt{i} is between \texttt{1} and the length of \mbox{\texttt{\mdseries\slshape list}} and \texttt{i \texttt{\symbol{94}} \mbox{\texttt{\mdseries\slshape f}} = i} if \texttt{i} is larger than the length of \mbox{\texttt{\mdseries\slshape list}}. An error will occur in \texttt{TransformationList} if \mbox{\texttt{\mdseries\slshape list}} is not dense, if \mbox{\texttt{\mdseries\slshape list}} contains an element which is not a positive integer, or if \mbox{\texttt{\mdseries\slshape list}} contains an integer not in \texttt{[ 1 .. Length( \mbox{\texttt{\mdseries\slshape list}} ) ]}. 

 \texttt{TransformationList} is the analogue in the context of transformations of \texttt{PermList} (\ref{PermList}). \texttt{Transformation} is a synonym of \texttt{TransformationList} when the argument is a list. 

 When the arguments are a list of positive integers \mbox{\texttt{\mdseries\slshape list}} and a function \mbox{\texttt{\mdseries\slshape func}}, \texttt{Transformation} returns the transformation \texttt{f} such that \texttt{\mbox{\texttt{\mdseries\slshape list}}[i] \texttt{\symbol{94}} f = \mbox{\texttt{\mdseries\slshape func}}( \mbox{\texttt{\mdseries\slshape list}}[i] )} if \texttt{i} is in the range \texttt{[ 1 .. Length( \mbox{\texttt{\mdseries\slshape list}} ) ]} and \texttt{f} fixes all other points. 
\begin{Verbatim}[commandchars=!@|,fontsize=\small,frame=single,label=Example]
  !gapprompt@gap>| !gapinput@SetUserPreference( "NotationForTransformations", "input" );|
  !gapprompt@gap>| !gapinput@f := Transformation( [ 11, 10, 2, 11, 4, 4, 7, 6, 9, 10, 1, 11 ] );|
  Transformation( [ 11, 10, 2, 11, 4, 4, 7, 6, 9, 10, 1, 11 ] )
  !gapprompt@gap>| !gapinput@f := TransformationList( [ 2, 3, 3, 1 ] );|
  Transformation( [ 2, 3, 3, 1 ] )
  !gapprompt@gap>| !gapinput@SetUserPreference( "NotationForTransformations", "fr" );|
  !gapprompt@gap>| !gapinput@f := Transformation( [ 10, 11 ], x -> x ^ 2 );|
  <transformation: 1,2,3,4,5,6,7,8,9,100,121>
  !gapprompt@gap>| !gapinput@SetUserPreference( "NotationForTransformations", "input" );|
\end{Verbatim}
 }

  

\subsection{\textcolor{Chapter }{Transformation (for a source and destination)}}
\logpage{[ 53, 2, 2 ]}\nobreak
\hyperdef{L}{X8040642687531E7F}{}
{\noindent\textcolor{FuncColor}{$\triangleright$\enspace\texttt{Transformation({\mdseries\slshape src, dst})\index{Transformation@\texttt{Transformation}!for a source and destination}
\label{Transformation:for a source and destination}
}\hfill{\scriptsize (operation)}}\\
\noindent\textcolor{FuncColor}{$\triangleright$\enspace\texttt{TransformationListList({\mdseries\slshape src, dst})\index{TransformationListList@\texttt{TransformationListList}!for a source and destination}
\label{TransformationListList:for a source and destination}
}\hfill{\scriptsize (operation)}}\\
\textbf{\indent Returns:\ }
A transformation.



 If \mbox{\texttt{\mdseries\slshape src}} and \mbox{\texttt{\mdseries\slshape dst}} are lists of positive integers of the same length, such that \mbox{\texttt{\mdseries\slshape src}} contains no element twice, then \texttt{TransformationListList( \mbox{\texttt{\mdseries\slshape src}}, \mbox{\texttt{\mdseries\slshape dst}} )} returns a transformation \texttt{f} such that \texttt{src[i] \texttt{\symbol{94}} \mbox{\texttt{\mdseries\slshape f}} = dst[i]}. The transformation \mbox{\texttt{\mdseries\slshape f}} fixes all points larger than the maximum of the entries in \mbox{\texttt{\mdseries\slshape src}} and \mbox{\texttt{\mdseries\slshape dst}}. 

 \texttt{TransformationListList} is the analogue in the context of transformations of \texttt{MappingPermListList} (\ref{MappingPermListList}). \texttt{Transformation} is a synonym of \texttt{TransformationListList} when its arguments are two lists of positive integers. 
\begin{Verbatim}[commandchars=!@|,fontsize=\small,frame=single,label=Example]
  !gapprompt@gap>| !gapinput@Transformation( [ 10, 11 ],[ 11, 12 ] );|
  Transformation( [ 1, 2, 3, 4, 5, 6, 7, 8, 9, 11, 12, 12 ] )
  !gapprompt@gap>| !gapinput@TransformationListList( [ 1, 2, 3 ], [ 4, 5, 6 ] );|
  Transformation( [ 4, 5, 6, 4, 5, 6 ] )
\end{Verbatim}
 }

  

\subsection{\textcolor{Chapter }{TransformationByImageAndKernel (for an image and kernel)}}
\logpage{[ 53, 2, 3 ]}\nobreak
\hyperdef{L}{X7E82EBD68455EE4A}{}
{\noindent\textcolor{FuncColor}{$\triangleright$\enspace\texttt{TransformationByImageAndKernel({\mdseries\slshape im, ker})\index{TransformationByImageAndKernel@\texttt{TransformationByImageAndKernel}!for an image and kernel}
\label{TransformationByImageAndKernel:for an image and kernel}
}\hfill{\scriptsize (operation)}}\\
\textbf{\indent Returns:\ }
A transformation or \texttt{fail}.



 This operation returns the transformation \texttt{f} where \texttt{i \texttt{\symbol{94}} f = \mbox{\texttt{\mdseries\slshape im}}[\mbox{\texttt{\mdseries\slshape ker}}[i]]} for \texttt{i} in the range \texttt{[ 1 .. Length( \mbox{\texttt{\mdseries\slshape ker}} ) ]}. This transformation has flat kernel equal to \mbox{\texttt{\mdseries\slshape ker}} and image set equal to \texttt{Set( \mbox{\texttt{\mdseries\slshape im}} )}. 

 The argument \mbox{\texttt{\mdseries\slshape im}} should be a duplicate free list of positive integers and \mbox{\texttt{\mdseries\slshape ker}} should be the flat kernel of a transformation with rank equal to the length of \mbox{\texttt{\mdseries\slshape im}}. If the arguments do not fulfil these conditions, then \texttt{fail} is returned. 
\begin{Verbatim}[commandchars=!@|,fontsize=\small,frame=single,label=Example]
  !gapprompt@gap>| !gapinput@TransformationByImageAndKernel( [ 8, 1, 3, 4 ],|
  !gapprompt@>| !gapinput@                                   [ 1, 2, 3, 1, 2, 1, 2, 4 ] );|
  Transformation( [ 8, 1, 3, 8, 1, 8, 1, 4 ] )
  !gapprompt@gap>| !gapinput@TransformationByImageAndKernel( [ 1, 3, 8, 4 ],|
  !gapprompt@>| !gapinput@                                   [ 1, 2, 3, 1, 2, 1, 2, 4 ] );|
  Transformation( [ 1, 3, 8, 1, 3, 1, 3, 4 ] )
\end{Verbatim}
 }

  

\subsection{\textcolor{Chapter }{Idempotent}}
\logpage{[ 53, 2, 4 ]}\nobreak
\hyperdef{L}{X85D1071484CE004C}{}
{\noindent\textcolor{FuncColor}{$\triangleright$\enspace\texttt{Idempotent({\mdseries\slshape im, ker})\index{Idempotent@\texttt{Idempotent}}
\label{Idempotent}
}\hfill{\scriptsize (operation)}}\\
\textbf{\indent Returns:\ }
A transformation or \texttt{fail}.



 \texttt{Idempotent} returns the idempotent transformation with image set \mbox{\texttt{\mdseries\slshape im}} and flat kernel \mbox{\texttt{\mdseries\slshape ker}} if such a transformation exists and \texttt{fail} if it does not. More specifically, a transformation is returned when the
argument \mbox{\texttt{\mdseries\slshape im}} is a set of positive integers and \mbox{\texttt{\mdseries\slshape ker}} is the flat kernel of a transformation with rank equal to the length of \mbox{\texttt{\mdseries\slshape im}} and where \mbox{\texttt{\mdseries\slshape im}} has one element in every class of the kernel corresponding to \mbox{\texttt{\mdseries\slshape ker}}. 

 Note that this is function does not always return the same transformation as \texttt{TransformationByImageAndKernel} with the same arguments. 
\begin{Verbatim}[commandchars=!@|,fontsize=\small,frame=single,label=Example]
  !gapprompt@gap>| !gapinput@Idempotent( [ 2, 4, 6, 7, 8, 10, 11 ],|
  !gapprompt@>| !gapinput@               [ 1, 2, 1, 3, 3, 4, 5, 1, 6, 6, 7, 5 ] );|
  Transformation( [ 8, 2, 8, 4, 4, 6, 7, 8, 10, 10, 11, 7 ] )
  !gapprompt@gap>| !gapinput@TransformationByImageAndKernel( [ 2, 4, 6, 7, 8, 10, 11 ],|
  !gapprompt@>| !gapinput@                     [ 1, 2, 1, 3, 3, 4, 5, 1, 6, 6, 7, 5 ] );|
  Transformation( [ 2, 4, 2, 6, 6, 7, 8, 2, 10, 10, 11, 8 ] )
\end{Verbatim}
 }

  

\subsection{\textcolor{Chapter }{TransformationOp}}
\logpage{[ 53, 2, 5 ]}\nobreak
\hyperdef{L}{X7C2A3FC9782F2099}{}
{\noindent\textcolor{FuncColor}{$\triangleright$\enspace\texttt{TransformationOp({\mdseries\slshape obj, list[, func]})\index{TransformationOp@\texttt{TransformationOp}}
\label{TransformationOp}
}\hfill{\scriptsize (operation)}}\\
\noindent\textcolor{FuncColor}{$\triangleright$\enspace\texttt{TransformationOpNC({\mdseries\slshape obj, list[, func]})\index{TransformationOpNC@\texttt{TransformationOpNC}}
\label{TransformationOpNC}
}\hfill{\scriptsize (operation)}}\\
\textbf{\indent Returns:\ }
A transformation or \texttt{fail}.



 \texttt{TransformationOp} returns the transformation that corresponds to the action of the object \mbox{\texttt{\mdseries\slshape obj}} on the domain or list \mbox{\texttt{\mdseries\slshape list}} via the function \mbox{\texttt{\mdseries\slshape func}}. If the optional third argument \mbox{\texttt{\mdseries\slshape func}} is not specified, then the action \texttt{OnPoints} (\ref{OnPoints}) is used by default. Note that the returned transformation refers to the
positions in \mbox{\texttt{\mdseries\slshape list}} even if \mbox{\texttt{\mdseries\slshape list}} itself consists of integers. 

 This function is the analogue in the context of transformations of \texttt{Permutation} (\ref{Permutation:for a group, an action domain, etc.}). 

 If \mbox{\texttt{\mdseries\slshape obj}} does not map elements of \mbox{\texttt{\mdseries\slshape list}} into \mbox{\texttt{\mdseries\slshape list}}, then \texttt{fail} is returned. 

 \texttt{TransformationOpNC} does not check that \mbox{\texttt{\mdseries\slshape obj}} maps elements of \mbox{\texttt{\mdseries\slshape list}} to elements of \mbox{\texttt{\mdseries\slshape list}} or that a transformation is defined by the action of \mbox{\texttt{\mdseries\slshape obj}} on \mbox{\texttt{\mdseries\slshape list}} via \mbox{\texttt{\mdseries\slshape func}}. This function should be used only with caution, and in situations where it
is guaranteed that the arguments have the required properties. 
\begin{Verbatim}[commandchars=!@|,fontsize=\small,frame=single,label=Example]
  !gapprompt@gap>| !gapinput@f := Transformation( [ 10, 2, 3, 10, 5, 10, 7, 2, 5, 6 ] );;|
  !gapprompt@gap>| !gapinput@TransformationOp( f, [ 2, 3 ] );|
  IdentityTransformation
  !gapprompt@gap>| !gapinput@TransformationOp( f, [ 1, 2, 3 ] );|
  fail
  !gapprompt@gap>| !gapinput@S := SemigroupByMultiplicationTable( [ [ 1, 1, 1 ],|
  !gapprompt@>| !gapinput@                                          [ 1, 1, 1 ],|
  !gapprompt@>| !gapinput@                                          [ 1, 1, 2 ] ] );;|
  !gapprompt@gap>| !gapinput@TransformationOp( Elements( S )[1], S, OnRight );|
  Transformation( [ 1, 1, 1 ] )
  !gapprompt@gap>| !gapinput@TransformationOp( Elements( S )[3], S, OnRight );|
  Transformation( [ 1, 1, 2 ] )
\end{Verbatim}
 }

 

\subsection{\textcolor{Chapter }{TransformationNumber}}
\logpage{[ 53, 2, 6 ]}\nobreak
\hyperdef{L}{X7D6FCC417DE86CD1}{}
{\noindent\textcolor{FuncColor}{$\triangleright$\enspace\texttt{TransformationNumber({\mdseries\slshape m, n})\index{TransformationNumber@\texttt{TransformationNumber}}
\label{TransformationNumber}
}\hfill{\scriptsize (operation)}}\\
\noindent\textcolor{FuncColor}{$\triangleright$\enspace\texttt{NumberTransformation({\mdseries\slshape f[, n]})\index{NumberTransformation@\texttt{NumberTransformation}}
\label{NumberTransformation}
}\hfill{\scriptsize (operation)}}\\
\textbf{\indent Returns:\ }
A transformation or a number.



 These functions implement a bijection from the transformations with degree at
most \mbox{\texttt{\mdseries\slshape n}} to the numbers \texttt{[ 1 .. \mbox{\texttt{\mdseries\slshape n}} \texttt{\symbol{94}} \mbox{\texttt{\mdseries\slshape n}} ]}. 

 More precisely, if \mbox{\texttt{\mdseries\slshape m}} and \mbox{\texttt{\mdseries\slshape n}} are positive integers such that \mbox{\texttt{\mdseries\slshape m}} is at most \texttt{\mbox{\texttt{\mdseries\slshape n}} \texttt{\symbol{94}} \mbox{\texttt{\mdseries\slshape n}}}, then \texttt{TransformationNumber} returns the \mbox{\texttt{\mdseries\slshape m}}th transformation with degree at most \mbox{\texttt{\mdseries\slshape n}}. 

 If \mbox{\texttt{\mdseries\slshape f}} is a transformation and \mbox{\texttt{\mdseries\slshape n}} is a positive integer, which is greater than or equal to the degree of \mbox{\texttt{\mdseries\slshape f}}, then \texttt{NumberTransformation} returns the number in \texttt{[ 1 .. \mbox{\texttt{\mdseries\slshape n}} \texttt{\symbol{94}} \mbox{\texttt{\mdseries\slshape n}} ]} that corresponds to \mbox{\texttt{\mdseries\slshape f}}. If the optional second argument \mbox{\texttt{\mdseries\slshape n}} is not specified, then the degree of \mbox{\texttt{\mdseries\slshape f}} is used by default. 
\begin{Verbatim}[commandchars=!@|,fontsize=\small,frame=single,label=Example]
  !gapprompt@gap>| !gapinput@f := Transformation( [ 3, 3, 5, 3, 3 ] );;|
  !gapprompt@gap>| !gapinput@NumberTransformation( f, 5 );|
  1613
  !gapprompt@gap>| !gapinput@NumberTransformation( f, 10 );|
  2242256790
  !gapprompt@gap>| !gapinput@TransformationNumber( 2242256790, 10 );|
  Transformation( [ 3, 3, 5, 3, 3 ] )
  !gapprompt@gap>| !gapinput@TransformationNumber( 1613, 5 );|
  Transformation( [ 3, 3, 5, 3, 3 ] )
\end{Verbatim}
 }

  
\subsection{\textcolor{Chapter }{RandomTransformation}}\logpage{[ 53, 2, 7 ]}
\hyperdef{L}{X8475448F87E8CB8A}{}
{
\noindent\textcolor{FuncColor}{$\triangleright$\enspace\texttt{RandomTransformation({\mdseries\slshape n})\index{RandomTransformation@\texttt{RandomTransformation}}
\label{RandomTransformation}
}\hfill{\scriptsize (operation)}}\\
\textbf{\indent Returns:\ }
A random transformation.



 If \mbox{\texttt{\mdseries\slshape n}} is a positive integer, then \texttt{RandomTransformation} returns a random transformation with degree at most \mbox{\texttt{\mdseries\slshape n}}. 
\begin{Verbatim}[commandchars=!@|,fontsize=\small,frame=single,label=Example]
  !gapprompt@gap>| !gapinput@RandomTransformation( 6 );|
  Transformation( [ 2, 1, 2, 1, 1, 2 ] )
\end{Verbatim}
 }

  

\subsection{\textcolor{Chapter }{IdentityTransformation}}
\logpage{[ 53, 2, 8 ]}\nobreak
\hyperdef{L}{X8268A58685BEFD6F}{}
{\noindent\textcolor{FuncColor}{$\triangleright$\enspace\texttt{IdentityTransformation\index{IdentityTransformation@\texttt{IdentityTransformation}}
\label{IdentityTransformation}
}\hfill{\scriptsize (global variable)}}\\


 This variable is bound to the identity transformation, which has degree \texttt{0}. 
\begin{Verbatim}[commandchars=!@|,fontsize=\small,frame=single,label=Example]
  !gapprompt@gap>| !gapinput@IdentityTransformation;|
  IdentityTransformation
\end{Verbatim}
 }

  

\subsection{\textcolor{Chapter }{ConstantTransformation}}
\logpage{[ 53, 2, 9 ]}\nobreak
\hyperdef{L}{X7F1E4B5184210D2B}{}
{\noindent\textcolor{FuncColor}{$\triangleright$\enspace\texttt{ConstantTransformation({\mdseries\slshape m, n})\index{ConstantTransformation@\texttt{ConstantTransformation}}
\label{ConstantTransformation}
}\hfill{\scriptsize (operation)}}\\
\textbf{\indent Returns:\ }
A transformation.



 This function returns a constant transformation \texttt{f} such that \texttt{i \texttt{\symbol{94}} f = \mbox{\texttt{\mdseries\slshape n}}} for all \texttt{i} less than or equal to \mbox{\texttt{\mdseries\slshape m}}, when \mbox{\texttt{\mdseries\slshape n}} and \mbox{\texttt{\mdseries\slshape m}} are positive integers. 
\begin{Verbatim}[commandchars=!@|,fontsize=\small,frame=single,label=Example]
  !gapprompt@gap>| !gapinput@ConstantTransformation( 5, 1 );|
  Transformation( [ 1, 1, 1, 1, 1 ] )
  !gapprompt@gap>| !gapinput@ConstantTransformation( 6, 4 );|
  Transformation( [ 4, 4, 4, 4, 4, 4 ] )
\end{Verbatim}
 }

 }

  
\section{\textcolor{Chapter }{Changing the representation of a transformation}}\label{sect:ChangingRepTransformations}
\logpage{[ 53, 3, 0 ]}
\hyperdef{L}{X7F81A18B813C9DF0}{}
{
  It is possible that a transformation in \textsf{GAP} can be represented as another type of object, or that another type of \textsf{GAP} object can be represented as a transformation. 

 The operations \texttt{AsPermutation} (\ref{AsPermutation}) and \texttt{AsPartialPerm} (\ref{AsPartialPerm:for a transformation and a positive integer}) can be used to convert transformations into permutations or partial
permutations, where appropriate. In this section we describe functions for
converting other types of objects into transformations. 

\subsection{\textcolor{Chapter }{AsTransformation}}
\logpage{[ 53, 3, 1 ]}\nobreak
\hyperdef{L}{X7C5360B2799943F3}{}
{\noindent\textcolor{FuncColor}{$\triangleright$\enspace\texttt{AsTransformation({\mdseries\slshape f[, n]})\index{AsTransformation@\texttt{AsTransformation}}
\label{AsTransformation}
}\hfill{\scriptsize (attribute)}}\\
\textbf{\indent Returns:\ }
A transformation.



 \texttt{AsTransformation} returns the permutation, transformation, partial permutation or binary
relation \mbox{\texttt{\mdseries\slshape f}} as a transformation. 
\begin{description}
\item[{for permutations}]  If \mbox{\texttt{\mdseries\slshape f}} is a permutation and \mbox{\texttt{\mdseries\slshape n}} is a non\texttt{\symbol{45}}negative integer, then \texttt{AsTransformation( \mbox{\texttt{\mdseries\slshape f}}, \mbox{\texttt{\mdseries\slshape n}} )} returns the transformation \texttt{g} such that \texttt{i \texttt{\symbol{94}} g = i \texttt{\symbol{94}} f} for all \texttt{i} in the range \texttt{[ 1 .. \mbox{\texttt{\mdseries\slshape n}} ]}. 

 If no non\texttt{\symbol{45}}negative integer \mbox{\texttt{\mdseries\slshape n}} is specified, then the largest moved point of \mbox{\texttt{\mdseries\slshape f}} is used as the value for \mbox{\texttt{\mdseries\slshape n}}; see \texttt{LargestMovedPoint} (\ref{LargestMovedPoint:for a permutation}). 
\item[{for transformations}]  If \mbox{\texttt{\mdseries\slshape f}} is a transformation and \mbox{\texttt{\mdseries\slshape n}} is a non\texttt{\symbol{45}}negative integer less than the degree of \mbox{\texttt{\mdseries\slshape f}} such that \mbox{\texttt{\mdseries\slshape f}} is a transformation of \texttt{[ 1 .. \mbox{\texttt{\mdseries\slshape n}} ]}, then \texttt{AsTransformation} returns the restriction of \mbox{\texttt{\mdseries\slshape f}} to \texttt{[ 1 .. \mbox{\texttt{\mdseries\slshape n}} ]}. 

 If \mbox{\texttt{\mdseries\slshape f}} is a transformation and \mbox{\texttt{\mdseries\slshape n}} is not specified or is greater than or equal to the degree of \mbox{\texttt{\mdseries\slshape f}}, then \mbox{\texttt{\mdseries\slshape f}} is returned. 
\item[{for partial permutations}]  A partial permutation \mbox{\texttt{\mdseries\slshape f}} can be converted into a transformation \texttt{g} as follows. The degree \texttt{m} of \texttt{g} is equal to the maximum of \mbox{\texttt{\mdseries\slshape n}}, the largest moved point of \mbox{\texttt{\mdseries\slshape f}} plus \texttt{1}, and the largest image of a moved point plus \texttt{1}. The transformation \texttt{g} agrees with \mbox{\texttt{\mdseries\slshape f}} on the domain of \mbox{\texttt{\mdseries\slshape f}} and maps the points in \texttt{[ 1 .. m ]}, which are not in the domain of \mbox{\texttt{\mdseries\slshape f}} to \texttt{n}, i.e. \texttt{i \texttt{\symbol{94}} g = i \texttt{\symbol{94}} \mbox{\texttt{\mdseries\slshape f}}} for all \texttt{i} in the domain of \mbox{\texttt{\mdseries\slshape f}}, \texttt{i \texttt{\symbol{94}} g = n} for all \texttt{i} in \texttt{[ 1 .. n ]}, and \texttt{i \texttt{\symbol{94}} g = i} for all \texttt{i} greater than \mbox{\texttt{\mdseries\slshape n}}. \texttt{AsTransformation( \mbox{\texttt{\mdseries\slshape f}} )} returns the transformation \texttt{g} defined in the previous sentences. 

 If the optional argument \mbox{\texttt{\mdseries\slshape n}} is not present, then the default value of the maximum of the largest moved
point and the largest image of a moved point of \mbox{\texttt{\mdseries\slshape f}} plus \texttt{1} is used. 
\item[{for binary relations}]  In the case that \mbox{\texttt{\mdseries\slshape f}} is a binary relation, which defines a transformation, \texttt{AsTransformation} returns that transformation. 
\end{description}
 
\begin{Verbatim}[commandchars=!@|,fontsize=\small,frame=single,label=Example]
  !gapprompt@gap>| !gapinput@f := Transformation( [ 3, 5, 3, 4, 1, 2 ] );;|
  !gapprompt@gap>| !gapinput@AsTransformation( f, 5 );|
  Transformation( [ 3, 5, 3, 4, 1 ] )
  !gapprompt@gap>| !gapinput@AsTransformation( f, 10 );|
  Transformation( [ 3, 5, 3, 4, 1, 2 ] )
  !gapprompt@gap>| !gapinput@AsTransformation( (1,3)(2,4) );|
  Transformation( [ 3, 4, 1, 2 ] )
  !gapprompt@gap>| !gapinput@AsTransformation( (1,3)(2,4), 10 );|
  Transformation( [ 3, 4, 1, 2 ] )
  !gapprompt@gap>| !gapinput@f := PartialPerm( [ 1, 2, 3, 4, 5, 6 ], [ 6, 7, 1, 4, 3, 2 ] );|
  [5,3,1,6,2,7](4)
  !gapprompt@gap>| !gapinput@AsTransformation( f, 11 );|
  Transformation( [ 6, 7, 1, 4, 3, 2, 11, 11, 11, 11, 11 ] )
  !gapprompt@gap>| !gapinput@AsPartialPerm( last, DomainOfPartialPerm( f ) );|
  [5,3,1,6,2,7](4)
  !gapprompt@gap>| !gapinput@AsTransformation( f, 14 );|
  Transformation( [ 6, 7, 1, 4, 3, 2, 14, 14, 14, 14, 14, 14, 14, 14 ] )
  !gapprompt@gap>| !gapinput@AsPartialPerm( last, DomainOfPartialPerm( f ) );|
  [5,3,1,6,2,7](4)
  !gapprompt@gap>| !gapinput@AsTransformation( f );|
  Transformation( [ 6, 7, 1, 4, 3, 2, 8, 8 ] )
  !gapprompt@gap>| !gapinput@AsTransformation( Transformation( [ 1, 1, 2 ] ), 0 );|
  IdentityTransformation
\end{Verbatim}
 }

 

\subsection{\textcolor{Chapter }{RestrictedTransformation}}
\logpage{[ 53, 3, 2 ]}\nobreak
\hyperdef{L}{X846A6F6B7B715188}{}
{\noindent\textcolor{FuncColor}{$\triangleright$\enspace\texttt{RestrictedTransformation({\mdseries\slshape f, list})\index{RestrictedTransformation@\texttt{RestrictedTransformation}}
\label{RestrictedTransformation}
}\hfill{\scriptsize (function)}}\\
\textbf{\indent Returns:\ }
A transformation.



 \texttt{RestrictedTransformation} returns the new transformation \texttt{g} such that \texttt{ i \texttt{\symbol{94}} g = i \texttt{\symbol{94}} \mbox{\texttt{\mdseries\slshape f}}} for all \texttt{i} in \mbox{\texttt{\mdseries\slshape list}} and such that \texttt{i \texttt{\symbol{94}} g = i} for all \texttt{i} not in \mbox{\texttt{\mdseries\slshape list}}. 

 
\begin{Verbatim}[commandchars=!@|,fontsize=\small,frame=single,label=Example]
  !gapprompt@gap>| !gapinput@f := Transformation( [ 2, 10, 5, 9, 10, 9, 6, 3, 8, 4, 6, 5 ] );;|
  !gapprompt@gap>| !gapinput@RestrictedTransformation( f, [ 1, 2, 3, 10, 11, 12 ] );|
  Transformation( [ 2, 10, 5, 4, 5, 6, 7, 8, 9, 4, 6, 5 ] )
\end{Verbatim}
 }

 

\subsection{\textcolor{Chapter }{PermutationOfImage}}
\logpage{[ 53, 3, 3 ]}\nobreak
\hyperdef{L}{X8708AE247F5B129B}{}
{\noindent\textcolor{FuncColor}{$\triangleright$\enspace\texttt{PermutationOfImage({\mdseries\slshape f})\index{PermutationOfImage@\texttt{PermutationOfImage}}
\label{PermutationOfImage}
}\hfill{\scriptsize (function)}}\\
\textbf{\indent Returns:\ }
A permutation or \texttt{fail}.



 If the transformation \mbox{\texttt{\mdseries\slshape f}} is a permutation of the points in its image, then \texttt{PermutationOfImage} returns this permutation. If \mbox{\texttt{\mdseries\slshape f}} does not permute its image, then \texttt{fail} is returned. 

 If \mbox{\texttt{\mdseries\slshape f}} happens to be a permutation, then \texttt{PermutationOfImage} with argument \mbox{\texttt{\mdseries\slshape f}} returns the same value as \texttt{AsPermutation} with argument \mbox{\texttt{\mdseries\slshape f}}. 
\begin{Verbatim}[commandchars=!@|,fontsize=\small,frame=single,label=Example]
  !gapprompt@gap>| !gapinput@f := Transformation( [ 5, 8, 3, 5, 8, 6, 2, 2, 7, 8 ] );;|
  !gapprompt@gap>| !gapinput@PermutationOfImage( f );|
  fail
  !gapprompt@gap>| !gapinput@f := Transformation( [ 8, 2, 10, 2, 4, 4, 7, 6, 9, 10 ] );;|
  !gapprompt@gap>| !gapinput@PermutationOfImage( f );|
  fail
  !gapprompt@gap>| !gapinput@f := Transformation( [ 1, 3, 6, 6, 2, 10, 2, 3, 10, 5 ] );;|
  !gapprompt@gap>| !gapinput@PermutationOfImage( f );|
  (2,3,6,10,5)
  !gapprompt@gap>| !gapinput@f := Transformation( [ 5, 2, 8, 4, 1, 8, 10, 3, 5, 7 ] );;|
  !gapprompt@gap>| !gapinput@PermutationOfImage( f );|
  (1,5)(3,8)(7,10)
\end{Verbatim}
 }

 }

   
\section{\textcolor{Chapter }{Operators for transformations}}\label{sect:OperatorsTransformations}
\logpage{[ 53, 4, 0 ]}
\hyperdef{L}{X812CEC008609A8A2}{}
{
  

\subsection{\textcolor{Chapter }{\texttt{\symbol{92}}\texttt{\symbol{94}} (for a positive integer and a transformation)}}
\logpage{[ 53, 4, 1 ]}\nobreak
\hyperdef{L}{X853031E37F214D10}{}
{\noindent\textcolor{FuncColor}{$\triangleright$\enspace\texttt{\texttt{\symbol{92}}\texttt{\symbol{94}}({\mdseries\slshape i, f})\index{\texttt{\symbol{92}}\texttt{\symbol{94}}@\texttt{\texttt{\symbol{92}}\texttt{\symbol{94}}}!for a positive integer and a transformation}
\label{bSlash^:for a positive integer and a transformation}
}\hfill{\scriptsize (method)}}\\


 \texttt{\mbox{\texttt{\mdseries\slshape i}} \texttt{\symbol{94}} \mbox{\texttt{\mdseries\slshape f}}} returns the image of the positive integer \mbox{\texttt{\mdseries\slshape i}} under the transformation \mbox{\texttt{\mdseries\slshape f}}. }

 

\subsection{\textcolor{Chapter }{\texttt{\symbol{92}}\texttt{\symbol{94}} (for a transformation and a permutation)}}
\logpage{[ 53, 4, 2 ]}\nobreak
\hyperdef{L}{X824A8E247DE7E53E}{}
{\noindent\textcolor{FuncColor}{$\triangleright$\enspace\texttt{\texttt{\symbol{92}}\texttt{\symbol{94}}({\mdseries\slshape f, g})\index{\texttt{\symbol{92}}\texttt{\symbol{94}}@\texttt{\texttt{\symbol{92}}\texttt{\symbol{94}}}!for a transformation and a permutation}
\label{bSlash^:for a transformation and a permutation}
}\hfill{\scriptsize (method)}}\\


 \texttt{\mbox{\texttt{\mdseries\slshape f}} \texttt{\symbol{94}} \mbox{\texttt{\mdseries\slshape g}}} returns \texttt{\mbox{\texttt{\mdseries\slshape g}} \texttt{\symbol{94}} \texttt{\symbol{45}}1 * \mbox{\texttt{\mdseries\slshape f}} * \mbox{\texttt{\mdseries\slshape g}}} when \mbox{\texttt{\mdseries\slshape f}} is a transformation and \mbox{\texttt{\mdseries\slshape g}} is a permutation \texttt{\texttt{\symbol{92}}\texttt{\symbol{94}}} (\ref{bSlash^}). This operation requires essentially the same number of steps as multiplying
a transformation by a permutation, which is approximately one third of the
number required to first invert \mbox{\texttt{\mdseries\slshape g}}, take the product with \mbox{\texttt{\mdseries\slshape f}}, and then the product with \mbox{\texttt{\mdseries\slshape g}}. }

 

\subsection{\textcolor{Chapter }{\texttt{\symbol{92}}* (for transformations)}}
\logpage{[ 53, 4, 3 ]}\nobreak
\hyperdef{L}{X7F6573E67D27D822}{}
{\noindent\textcolor{FuncColor}{$\triangleright$\enspace\texttt{\texttt{\symbol{92}}*({\mdseries\slshape f, g})\index{\texttt{\symbol{92}}*@\texttt{\texttt{\symbol{92}}*}!for transformations}
\label{bSlash*:for transformations}
}\hfill{\scriptsize (method)}}\\


 \texttt{\mbox{\texttt{\mdseries\slshape f}} * \mbox{\texttt{\mdseries\slshape g}}} returns the composition of \mbox{\texttt{\mdseries\slshape f}} and \mbox{\texttt{\mdseries\slshape g}} when \mbox{\texttt{\mdseries\slshape f}} and \mbox{\texttt{\mdseries\slshape g}} are transformations or permutations. The product of a permutation and a
transformation is returned as a transformation. }

 

\subsection{\textcolor{Chapter }{\texttt{\symbol{92}}/ (for a transformation and a permutation)}}
\logpage{[ 53, 4, 4 ]}\nobreak
\hyperdef{L}{X828CA7137F97C124}{}
{\noindent\textcolor{FuncColor}{$\triangleright$\enspace\texttt{\texttt{\symbol{92}}/({\mdseries\slshape f, g})\index{\texttt{\symbol{92}}/@\texttt{\texttt{\symbol{92}}/}!for a transformation and a permutation}
\label{bSlash/:for a transformation and a permutation}
}\hfill{\scriptsize (method)}}\\


 \texttt{\mbox{\texttt{\mdseries\slshape f}} / \mbox{\texttt{\mdseries\slshape g}}} returns \texttt{\mbox{\texttt{\mdseries\slshape f}} * \mbox{\texttt{\mdseries\slshape g}} \texttt{\symbol{94}} \texttt{\symbol{45}}1} when \mbox{\texttt{\mdseries\slshape f}} is a transformation and \mbox{\texttt{\mdseries\slshape g}} is a permutation. This operation requires essentially the same number of steps
as multiplying a transformation by a permutation, which is approximately half
the number required to first invert \mbox{\texttt{\mdseries\slshape g}} and then take the product with \mbox{\texttt{\mdseries\slshape f}}. }

 

\subsection{\textcolor{Chapter }{LeftQuotient (for a permutation and transformation)}}
\logpage{[ 53, 4, 5 ]}\nobreak
\hyperdef{L}{X7856D91E8709EF5B}{}
{\noindent\textcolor{FuncColor}{$\triangleright$\enspace\texttt{LeftQuotient({\mdseries\slshape g, f})\index{LeftQuotient@\texttt{LeftQuotient}!for a permutation and transformation}
\label{LeftQuotient:for a permutation and transformation}
}\hfill{\scriptsize (method)}}\\


 returns \texttt{\mbox{\texttt{\mdseries\slshape g}} \texttt{\symbol{94}} \texttt{\symbol{45}}1 * \mbox{\texttt{\mdseries\slshape f}}} when \mbox{\texttt{\mdseries\slshape f}} is a transformation and \mbox{\texttt{\mdseries\slshape g}} is a permutation. This operation uses essentially the same number of steps as
multiplying a transformation by a permutation, which is approximately half the
number required to first invert \mbox{\texttt{\mdseries\slshape g}} and then take the product with \mbox{\texttt{\mdseries\slshape f}}. }

 

\subsection{\textcolor{Chapter }{\texttt{\symbol{92}}{\textless} (for transformations)}}
\logpage{[ 53, 4, 6 ]}\nobreak
\hyperdef{L}{X84B8E294826A9377}{}
{\noindent\textcolor{FuncColor}{$\triangleright$\enspace\texttt{\texttt{\symbol{92}}{\textless}({\mdseries\slshape i, f})\index{\texttt{\symbol{92}}{\textless}@\texttt{\texttt{\symbol{92}}{\textless}}!for transformations}
\label{bSlash<:for transformations}
}\hfill{\scriptsize (method)}}\\


 \index{smaller!for transformations} \texttt{\mbox{\texttt{\mdseries\slshape f}} {\textless} \mbox{\texttt{\mdseries\slshape g}}} returns \texttt{true} if the image list of \mbox{\texttt{\mdseries\slshape f}} is lexicographically less than the image list of \mbox{\texttt{\mdseries\slshape g}} and \texttt{false} if it is not. }

 

\subsection{\textcolor{Chapter }{\texttt{\symbol{92}}= (for transformations)}}
\logpage{[ 53, 4, 7 ]}\nobreak
\hyperdef{L}{X7D454AAD851AE07E}{}
{\noindent\textcolor{FuncColor}{$\triangleright$\enspace\texttt{\texttt{\symbol{92}}=({\mdseries\slshape f, g})\index{\texttt{\symbol{92}}=@\texttt{\texttt{\symbol{92}}=}!for transformations}
\label{bSlash=:for transformations}
}\hfill{\scriptsize (method)}}\\


 \index{equality!for transformations} \texttt{\mbox{\texttt{\mdseries\slshape f}} = \mbox{\texttt{\mdseries\slshape g}}} returns \texttt{true} if the transformation \mbox{\texttt{\mdseries\slshape f}} equals the transformation \mbox{\texttt{\mdseries\slshape g}} and returns \texttt{false} if it does not. }

  

\subsection{\textcolor{Chapter }{PermLeftQuoTransformation}}
\logpage{[ 53, 4, 8 ]}\nobreak
\hyperdef{L}{X83DBA2A18719EFA8}{}
{\noindent\textcolor{FuncColor}{$\triangleright$\enspace\texttt{PermLeftQuoTransformation({\mdseries\slshape f, g})\index{PermLeftQuoTransformation@\texttt{PermLeftQuoTransformation}}
\label{PermLeftQuoTransformation}
}\hfill{\scriptsize (operation)}}\\
\noindent\textcolor{FuncColor}{$\triangleright$\enspace\texttt{PermLeftQuoTransformationNC({\mdseries\slshape f, g})\index{PermLeftQuoTransformationNC@\texttt{PermLeftQuoTransformationNC}}
\label{PermLeftQuoTransformationNC}
}\hfill{\scriptsize (function)}}\\
\textbf{\indent Returns:\ }
A permutation.



 Returns the permutation on the image set of \mbox{\texttt{\mdseries\slshape f}} induced by \texttt{\mbox{\texttt{\mdseries\slshape f}} \texttt{\symbol{94}} \texttt{\symbol{45}}1 * \mbox{\texttt{\mdseries\slshape g}}} when the transformations \mbox{\texttt{\mdseries\slshape f}} and \mbox{\texttt{\mdseries\slshape g}} have equal kernel and image set. 

 \texttt{PermLeftQuoTransformation} verifies that \mbox{\texttt{\mdseries\slshape f}} and \mbox{\texttt{\mdseries\slshape g}} have equal kernels and image sets, and returns an error if they do not. \texttt{PermLeftQuoTransformationNC} does no checks. 
\begin{Verbatim}[commandchars=!@|,fontsize=\small,frame=single,label=Example]
  !gapprompt@gap>| !gapinput@f := Transformation( [ 5, 6, 7, 1, 4, 3, 2, 7 ] );;|
  !gapprompt@gap>| !gapinput@g := Transformation( [ 5, 7, 1, 6, 4, 3, 2, 1 ] );;|
  !gapprompt@gap>| !gapinput@PermLeftQuoTransformation( f, g );|
  (1,6,7)
  !gapprompt@gap>| !gapinput@PermLeftQuoTransformation( g, f );|
  (1,7,6)
\end{Verbatim}
 }

  

\subsection{\textcolor{Chapter }{IsInjectiveListTrans}}
\logpage{[ 53, 4, 9 ]}\nobreak
\hyperdef{L}{X8275DFAA8270BB59}{}
{\noindent\textcolor{FuncColor}{$\triangleright$\enspace\texttt{IsInjectiveListTrans({\mdseries\slshape list, obj})\index{IsInjectiveListTrans@\texttt{IsInjectiveListTrans}}
\label{IsInjectiveListTrans}
}\hfill{\scriptsize (function)}}\\
\textbf{\indent Returns:\ }
\texttt{true} or \texttt{false}.



 The argument \mbox{\texttt{\mdseries\slshape obj}} should be a transformation or the list of images of a transformation and \mbox{\texttt{\mdseries\slshape list}} should be a list of positive integers. \texttt{IsInjectiveListTrans} checks if \mbox{\texttt{\mdseries\slshape obj}} is injective on \mbox{\texttt{\mdseries\slshape list}}. 

 More precisely, if \mbox{\texttt{\mdseries\slshape obj}} is a transformation, then we define \texttt{f := \mbox{\texttt{\mdseries\slshape obj}}} and if \mbox{\texttt{\mdseries\slshape obj}} is the image list of a transformation we define \texttt{f := Transformation( \mbox{\texttt{\mdseries\slshape obj}} )}. \texttt{IsInjectiveListTrans} returns \texttt{true} if \texttt{f} is injective on \mbox{\texttt{\mdseries\slshape list}} and \texttt{false} if it is not. If \mbox{\texttt{\mdseries\slshape list}} is not duplicate free, then \texttt{false} is returned. 

 
\begin{Verbatim}[commandchars=!@|,fontsize=\small,frame=single,label=Example]
  !gapprompt@gap>| !gapinput@f := Transformation( [ 2, 6, 7, 2, 6, 9, 9, 1, 1, 5 ] );;|
  !gapprompt@gap>| !gapinput@IsInjectiveListTrans( [ 1, 5 ], f );|
  true
  !gapprompt@gap>| !gapinput@IsInjectiveListTrans( [ 5, 1 ], f );|
  true
  !gapprompt@gap>| !gapinput@IsInjectiveListTrans( [ 5, 1, 5, 1, 1, ], f );|
  false
  !gapprompt@gap>| !gapinput@IsInjectiveListTrans( [ 5, 1, 2, 3 ], [ 1, 2, 3, 4, 5 ] );|
  true
\end{Verbatim}
 }

  

\subsection{\textcolor{Chapter }{ComponentTransformationInt}}
\logpage{[ 53, 4, 10 ]}\nobreak
\hyperdef{L}{X834A313B7DAF06D5}{}
{\noindent\textcolor{FuncColor}{$\triangleright$\enspace\texttt{ComponentTransformationInt({\mdseries\slshape f, n})\index{ComponentTransformationInt@\texttt{ComponentTransformationInt}}
\label{ComponentTransformationInt}
}\hfill{\scriptsize (operation)}}\\
\textbf{\indent Returns:\ }
A list of positive integers.



 If \mbox{\texttt{\mdseries\slshape f}} is a transformation and \mbox{\texttt{\mdseries\slshape n}} is a positive integer, then \texttt{ComponentTransformationInt} returns those elements \texttt{i} such that \texttt{\mbox{\texttt{\mdseries\slshape n}} \texttt{\symbol{94}} \mbox{\texttt{\mdseries\slshape f}} \texttt{\symbol{94}} j = i} for some positive integer \texttt{j}, i.e. the elements of the component of \mbox{\texttt{\mdseries\slshape f}} containing \mbox{\texttt{\mdseries\slshape n}} that can be obtained by applying powers of \mbox{\texttt{\mdseries\slshape f}} to \mbox{\texttt{\mdseries\slshape n}}. 
\begin{Verbatim}[commandchars=!@|,fontsize=\small,frame=single,label=Example]
  !gapprompt@gap>| !gapinput@f := Transformation( [ 6, 2, 8, 4, 7, 5, 8, 3, 5, 8 ] );;|
  !gapprompt@gap>| !gapinput@ComponentTransformationInt( f, 1 );|
  [ 1, 6, 5, 7, 8, 3 ]
  !gapprompt@gap>| !gapinput@ComponentTransformationInt( f, 12 );|
  [ 12 ]
  !gapprompt@gap>| !gapinput@ComponentTransformationInt( f, 5 );|
  [ 5, 7, 8, 3 ]
\end{Verbatim}
 }

  

\subsection{\textcolor{Chapter }{PreImagesOfTransformation}}
\logpage{[ 53, 4, 11 ]}\nobreak
\hyperdef{L}{X82F5DEEC837B60A3}{}
{\noindent\textcolor{FuncColor}{$\triangleright$\enspace\texttt{PreImagesOfTransformation({\mdseries\slshape f, n})\index{PreImagesOfTransformation@\texttt{PreImagesOfTransformation}}
\label{PreImagesOfTransformation}
}\hfill{\scriptsize (operation)}}\\
\textbf{\indent Returns:\ }
A set of positive integers.



 Returns the preimages of the positive integer \mbox{\texttt{\mdseries\slshape n}} under the transformation \mbox{\texttt{\mdseries\slshape f}}, i.e. the positive integers \texttt{i} such that \texttt{i \texttt{\symbol{94}} \mbox{\texttt{\mdseries\slshape f}} = n}. 
\begin{Verbatim}[commandchars=!@|,fontsize=\small,frame=single,label=Example]
  !gapprompt@gap>| !gapinput@f := Transformation( [ 2, 6, 7, 2, 6, 9, 9, 1, 1, 5 ] );;|
  !gapprompt@gap>| !gapinput@PreImagesOfTransformation( f, 1 );|
  [ 8, 9 ]
  !gapprompt@gap>| !gapinput@PreImagesOfTransformation( f, 3 );|
  [  ]
  !gapprompt@gap>| !gapinput@PreImagesOfTransformation( f, 100 );|
  [ 100 ]
\end{Verbatim}
 }

 }

  
\section{\textcolor{Chapter }{Attributes for transformations}}\label{sect:AttributesTransformations}
\logpage{[ 53, 5, 0 ]}
\hyperdef{L}{X86DE4F7A7C535820}{}
{
  In this section we describe the functions available in \textsf{GAP} for finding various properties and attributes of transformations. 

\subsection{\textcolor{Chapter }{DegreeOfTransformation}}
\logpage{[ 53, 5, 1 ]}\nobreak
\hyperdef{L}{X78A209C87CF0E32B}{}
{\noindent\textcolor{FuncColor}{$\triangleright$\enspace\texttt{DegreeOfTransformation({\mdseries\slshape f})\index{DegreeOfTransformation@\texttt{DegreeOfTransformation}}
\label{DegreeOfTransformation}
}\hfill{\scriptsize (function)}}\\
\noindent\textcolor{FuncColor}{$\triangleright$\enspace\texttt{DegreeOfTransformationCollection({\mdseries\slshape coll})\index{DegreeOfTransformationCollection@\texttt{DegreeOfTransformationCollection}}
\label{DegreeOfTransformationCollection}
}\hfill{\scriptsize (attribute)}}\\
\textbf{\indent Returns:\ }
A positive integer.



 The \emph{degree} of a transformation \mbox{\texttt{\mdseries\slshape f}} is the largest value such that \texttt{n \texttt{\symbol{94}} \mbox{\texttt{\mdseries\slshape f}} {\textless}{\textgreater} n} or \texttt{i \texttt{\symbol{94}} \mbox{\texttt{\mdseries\slshape f}} = n} for some \texttt{i {\textless}{\textgreater} n}. Equivalently, the degree of a transformation is the least value \texttt{n} such that \texttt{[ n + 1, n + 2, ... ]} is fixed pointwise by \mbox{\texttt{\mdseries\slshape f}}. 

 The degree of a collection of transformations \mbox{\texttt{\mdseries\slshape coll}} is the maximum degree of any transformation in \mbox{\texttt{\mdseries\slshape coll}}. 
\begin{Verbatim}[commandchars=!@|,fontsize=\small,frame=single,label=Example]
  !gapprompt@gap>| !gapinput@DegreeOfTransformation( IdentityTransformation );|
  0
  !gapprompt@gap>| !gapinput@DegreeOfTransformationCollection(|
  !gapprompt@>| !gapinput@[ Transformation( [ 1, 3, 4, 1 ] ),|
  !gapprompt@>| !gapinput@  Transformation( [ 3, 1, 1, 3, 4 ] ),|
  !gapprompt@>| !gapinput@  Transformation( [ 2, 4, 1, 2 ] ) ] );|
  5
\end{Verbatim}
 }

  

\subsection{\textcolor{Chapter }{ImageListOfTransformation}}
\logpage{[ 53, 5, 2 ]}\nobreak
\hyperdef{L}{X7AEC9E6687B3505A}{}
{\noindent\textcolor{FuncColor}{$\triangleright$\enspace\texttt{ImageListOfTransformation({\mdseries\slshape f[, n]})\index{ImageListOfTransformation@\texttt{ImageListOfTransformation}}
\label{ImageListOfTransformation}
}\hfill{\scriptsize (operation)}}\\
\noindent\textcolor{FuncColor}{$\triangleright$\enspace\texttt{ListTransformation({\mdseries\slshape f[, n]})\index{ListTransformation@\texttt{ListTransformation}}
\label{ListTransformation}
}\hfill{\scriptsize (operation)}}\\
\textbf{\indent Returns:\ }
The list of images of a transformation.



 Returns the list of images of \texttt{[ 1 .. \mbox{\texttt{\mdseries\slshape n}} ]} under the transformation \mbox{\texttt{\mdseries\slshape f}}, which is \texttt{[ 1 \texttt{\symbol{94}} \mbox{\texttt{\mdseries\slshape f}} .. \mbox{\texttt{\mdseries\slshape n}} \texttt{\symbol{94}} \mbox{\texttt{\mdseries\slshape f}} ]}. If the optional second argument \mbox{\texttt{\mdseries\slshape n}} is not present, then the degree of \mbox{\texttt{\mdseries\slshape f}} is used by default. 

 This is the analogue for transformations of \texttt{ListPerm} (\ref{ListPerm}) for permutations. 
\begin{Verbatim}[commandchars=!@|,fontsize=\small,frame=single,label=Example]
  !gapprompt@gap>| !gapinput@f := Transformation( [ 2 ,3, 4, 2, 4 ] );;|
  !gapprompt@gap>| !gapinput@ImageListOfTransformation( f );|
  [ 2, 3, 4, 2, 4 ]
  !gapprompt@gap>| !gapinput@ImageListOfTransformation( f, 10 );|
  [ 2, 3, 4, 2, 4, 6, 7, 8, 9, 10 ]
\end{Verbatim}
 }

  

\subsection{\textcolor{Chapter }{ImageSetOfTransformation}}
\logpage{[ 53, 5, 3 ]}\nobreak
\hyperdef{L}{X839A6D6082A21D1F}{}
{\noindent\textcolor{FuncColor}{$\triangleright$\enspace\texttt{ImageSetOfTransformation({\mdseries\slshape f[, n]})\index{ImageSetOfTransformation@\texttt{ImageSetOfTransformation}}
\label{ImageSetOfTransformation}
}\hfill{\scriptsize (attribute)}}\\
\textbf{\indent Returns:\ }
The set of images of the transformation.



 Returns the set of points in the list of images of \texttt{[ 1 .. \mbox{\texttt{\mdseries\slshape n}} ]} under \mbox{\texttt{\mdseries\slshape f}}, i.e. the sorted list of images with duplicates removed. If the optional
second argument \mbox{\texttt{\mdseries\slshape n}} is not given, then the degree of \mbox{\texttt{\mdseries\slshape f}} is used. 

 
\begin{Verbatim}[commandchars=!@|,fontsize=\small,frame=single,label=Example]
  !gapprompt@gap>| !gapinput@f := Transformation( [ 5, 6, 7, 1, 4, 3, 2, 7 ] );;|
  !gapprompt@gap>| !gapinput@ImageSetOfTransformation( f );|
  [ 1, 2, 3, 4, 5, 6, 7 ]
  !gapprompt@gap>| !gapinput@ImageSetOfTransformation( f, 10 );|
  [ 1, 2, 3, 4, 5, 6, 7, 9, 10 ]
\end{Verbatim}
 }

  

\subsection{\textcolor{Chapter }{RankOfTransformation (for a transformation and a positive integer)}}
\logpage{[ 53, 5, 4 ]}\nobreak
\hyperdef{L}{X818EBB167C7EA37B}{}
{\noindent\textcolor{FuncColor}{$\triangleright$\enspace\texttt{RankOfTransformation({\mdseries\slshape f[, n]})\index{RankOfTransformation@\texttt{RankOfTransformation}!for a transformation and a positive integer}
\label{RankOfTransformation:for a transformation and a positive integer}
}\hfill{\scriptsize (attribute)}}\\
\noindent\textcolor{FuncColor}{$\triangleright$\enspace\texttt{RankOfTransformation({\mdseries\slshape f[, list]})\index{RankOfTransformation@\texttt{RankOfTransformation}!for a transformation and a list}
\label{RankOfTransformation:for a transformation and a list}
}\hfill{\scriptsize (attribute)}}\\
\textbf{\indent Returns:\ }
The rank of a transformation.



 When the arguments are a transformation \mbox{\texttt{\mdseries\slshape f}} and a positive integer \mbox{\texttt{\mdseries\slshape n}}, \texttt{RankOfTransformation} returns the size of the set of images of the transformation \mbox{\texttt{\mdseries\slshape f}} in the range \texttt{[ 1 .. \mbox{\texttt{\mdseries\slshape n}} ]}. If the optional second argument \mbox{\texttt{\mdseries\slshape n}} is not specified, then the degree of \mbox{\texttt{\mdseries\slshape f}} is used. 

 When the arguments are a transformation \mbox{\texttt{\mdseries\slshape f}} and a list \mbox{\texttt{\mdseries\slshape list}} of positive integers, this function returns the size of the set of images of
the transformation \mbox{\texttt{\mdseries\slshape f}} on \mbox{\texttt{\mdseries\slshape list}}. 
\begin{Verbatim}[commandchars=!@|,fontsize=\small,frame=single,label=Example]
  !gapprompt@gap>| !gapinput@f := Transformation( [ 8, 5, 8, 2, 2, 8, 4, 7, 3, 1 ] );;|
  !gapprompt@gap>| !gapinput@ImageSetOfTransformation( f );|
  [ 1, 2, 3, 4, 5, 7, 8 ]
  !gapprompt@gap>| !gapinput@RankOfTransformation( f );|
  7
  !gapprompt@gap>| !gapinput@RankOfTransformation( f, 100 );|
  97
  !gapprompt@gap>| !gapinput@RankOfTransformation( f, [ 2, 5, 8 ] );|
  3
\end{Verbatim}
 }

  

\subsection{\textcolor{Chapter }{MovedPoints (for a transformation)}}
\logpage{[ 53, 5, 5 ]}\nobreak
\hyperdef{L}{X844F00F982D5BD3C}{}
{\noindent\textcolor{FuncColor}{$\triangleright$\enspace\texttt{MovedPoints({\mdseries\slshape f})\index{MovedPoints@\texttt{MovedPoints}!for a transformation}
\label{MovedPoints:for a transformation}
}\hfill{\scriptsize (attribute)}}\\
\noindent\textcolor{FuncColor}{$\triangleright$\enspace\texttt{MovedPoints({\mdseries\slshape coll})\index{MovedPoints@\texttt{MovedPoints}!for a transformation coll}
\label{MovedPoints:for a transformation coll}
}\hfill{\scriptsize (attribute)}}\\
\textbf{\indent Returns:\ }
A set of positive integers.



 When the argument is a transformation, \texttt{MovedPoints} returns the set of positive integers \texttt{i} such that \texttt{i \texttt{\symbol{94}} \mbox{\texttt{\mdseries\slshape f}} {\textless}{\textgreater} i}. 

 \texttt{MovedPoints} returns the set of points moved by some element of the collection of
transformations \mbox{\texttt{\mdseries\slshape coll}}. 
\begin{Verbatim}[commandchars=!@|,fontsize=\small,frame=single,label=Example]
  !gapprompt@gap>| !gapinput@f := Transformation( [ 6, 10, 1, 4, 6, 5, 1, 2, 3, 3 ] );;|
  !gapprompt@gap>| !gapinput@MovedPoints( f );|
  [ 1, 2, 3, 5, 6, 7, 8, 9, 10 ]
  !gapprompt@gap>| !gapinput@f := IdentityTransformation;|
  IdentityTransformation
  !gapprompt@gap>| !gapinput@MovedPoints( f );|
  [  ]
\end{Verbatim}
 }

  

\subsection{\textcolor{Chapter }{NrMovedPoints (for a transformation)}}
\logpage{[ 53, 5, 6 ]}\nobreak
\hyperdef{L}{X7FA6A4B57FDA003D}{}
{\noindent\textcolor{FuncColor}{$\triangleright$\enspace\texttt{NrMovedPoints({\mdseries\slshape f})\index{NrMovedPoints@\texttt{NrMovedPoints}!for a transformation}
\label{NrMovedPoints:for a transformation}
}\hfill{\scriptsize (attribute)}}\\
\noindent\textcolor{FuncColor}{$\triangleright$\enspace\texttt{NrMovedPoints({\mdseries\slshape coll})\index{NrMovedPoints@\texttt{NrMovedPoints}!for a transformation coll}
\label{NrMovedPoints:for a transformation coll}
}\hfill{\scriptsize (attribute)}}\\
\textbf{\indent Returns:\ }
A positive integer.



 When the argument is a transformation,\texttt{NrMovedPoints} returns the number of positive integers \texttt{i} such that \texttt{i \texttt{\symbol{94}} \mbox{\texttt{\mdseries\slshape f}} {\textless}{\textgreater} i}. 

 \texttt{MovedPoints} returns the number of points which are moved by at least one element of the
collection of transformations \mbox{\texttt{\mdseries\slshape coll}}. 
\begin{Verbatim}[commandchars=!@|,fontsize=\small,frame=single,label=Example]
  !gapprompt@gap>| !gapinput@f := Transformation( [ 7, 1, 4, 3, 2, 7, 7, 6, 6, 5 ] );;|
  !gapprompt@gap>| !gapinput@NrMovedPoints( f );|
  9
  !gapprompt@gap>| !gapinput@NrMovedPoints( IdentityTransformation );|
  0
\end{Verbatim}
 }

  

\subsection{\textcolor{Chapter }{SmallestMovedPoint (for a transformation)}}
\logpage{[ 53, 5, 7 ]}\nobreak
\hyperdef{L}{X86C0DDDC7881273A}{}
{\noindent\textcolor{FuncColor}{$\triangleright$\enspace\texttt{SmallestMovedPoint({\mdseries\slshape f})\index{SmallestMovedPoint@\texttt{SmallestMovedPoint}!for a transformation}
\label{SmallestMovedPoint:for a transformation}
}\hfill{\scriptsize (attribute)}}\\
\noindent\textcolor{FuncColor}{$\triangleright$\enspace\texttt{SmallestMovedPoint({\mdseries\slshape coll})\index{SmallestMovedPoint@\texttt{SmallestMovedPoint}!for a transformation coll}
\label{SmallestMovedPoint:for a transformation coll}
}\hfill{\scriptsize (method)}}\\
\textbf{\indent Returns:\ }
A positive integer or \texttt{infinity}.



 \texttt{SmallestMovedPoint} returns the smallest positive integer \texttt{i} such that \texttt{i \texttt{\symbol{94}} \mbox{\texttt{\mdseries\slshape f}} {\textless}{\textgreater} i} if such an \texttt{i} exists. If \mbox{\texttt{\mdseries\slshape f}} is the identity transformation, then \texttt{infinity} is returned. 

 If the argument is a collection of transformations \mbox{\texttt{\mdseries\slshape coll}}, then the smallest point which is moved by at least one element of \mbox{\texttt{\mdseries\slshape coll}} is returned, if such a point exists. If \mbox{\texttt{\mdseries\slshape coll}} only contains identity transformations, then \texttt{SmallestMovedPoint} returns \texttt{infinity}. 
\begin{Verbatim}[commandchars=!@|,fontsize=\small,frame=single,label=Example]
  !gapprompt@gap>| !gapinput@S := FullTransformationSemigroup( 5 );|
  <full transformation monoid of degree 5>
  !gapprompt@gap>| !gapinput@SmallestMovedPoint( S );|
  1
  !gapprompt@gap>| !gapinput@S := Semigroup( IdentityTransformation );|
  <trivial transformation group of degree 0 with 1 generator>
  !gapprompt@gap>| !gapinput@SmallestMovedPoint( S );|
  infinity
  !gapprompt@gap>| !gapinput@f := Transformation( [ 1, 2, 3, 6, 6, 6 ] );;|
  !gapprompt@gap>| !gapinput@SmallestMovedPoint( f );|
  4
\end{Verbatim}
 }

  

\subsection{\textcolor{Chapter }{LargestMovedPoint (for a transformation)}}
\logpage{[ 53, 5, 8 ]}\nobreak
\hyperdef{L}{X8383A7727AC97724}{}
{\noindent\textcolor{FuncColor}{$\triangleright$\enspace\texttt{LargestMovedPoint({\mdseries\slshape f})\index{LargestMovedPoint@\texttt{LargestMovedPoint}!for a transformation}
\label{LargestMovedPoint:for a transformation}
}\hfill{\scriptsize (attribute)}}\\
\noindent\textcolor{FuncColor}{$\triangleright$\enspace\texttt{LargestMovedPoint({\mdseries\slshape coll})\index{LargestMovedPoint@\texttt{LargestMovedPoint}!for a transformation coll}
\label{LargestMovedPoint:for a transformation coll}
}\hfill{\scriptsize (method)}}\\
\textbf{\indent Returns:\ }
A positive integer.



 \texttt{LargestMovedPoint} returns the largest positive integers \texttt{i} such that \texttt{i \texttt{\symbol{94}} \mbox{\texttt{\mdseries\slshape f}} {\textless}{\textgreater} i} if such an \texttt{i} exists. If \mbox{\texttt{\mdseries\slshape f}} is the identity transformation, then \texttt{0} is returned. 

 If the argument is a collection of transformations \mbox{\texttt{\mdseries\slshape coll}}, then the largest point which is moved by at least one element of \mbox{\texttt{\mdseries\slshape coll}} is returned, if such a point exists. If \mbox{\texttt{\mdseries\slshape coll}} only contains identity transformations, then \texttt{LargestMovedPoint} returns \texttt{0}. 
\begin{Verbatim}[commandchars=!@|,fontsize=\small,frame=single,label=Example]
  !gapprompt@gap>| !gapinput@S := FullTransformationSemigroup( 5 );|
  <full transformation monoid of degree 5>
  !gapprompt@gap>| !gapinput@LargestMovedPoint( S );|
  5
  !gapprompt@gap>| !gapinput@S := Semigroup( IdentityTransformation );|
  <trivial transformation group of degree 0 with 1 generator>
  !gapprompt@gap>| !gapinput@LargestMovedPoint( S );|
  0
  !gapprompt@gap>| !gapinput@f := Transformation( [ 1, 2, 3, 6, 6, 6 ] );;|
  !gapprompt@gap>| !gapinput@LargestMovedPoint( f );|
  5
\end{Verbatim}
 }

  

\subsection{\textcolor{Chapter }{SmallestImageOfMovedPoint (for a transformation)}}
\logpage{[ 53, 5, 9 ]}\nobreak
\hyperdef{L}{X7CCFE27E83676572}{}
{\noindent\textcolor{FuncColor}{$\triangleright$\enspace\texttt{SmallestImageOfMovedPoint({\mdseries\slshape f})\index{SmallestImageOfMovedPoint@\texttt{SmallestImageOfMovedPoint}!for a transformation}
\label{SmallestImageOfMovedPoint:for a transformation}
}\hfill{\scriptsize (attribute)}}\\
\noindent\textcolor{FuncColor}{$\triangleright$\enspace\texttt{SmallestImageOfMovedPoint({\mdseries\slshape coll})\index{SmallestImageOfMovedPoint@\texttt{SmallestImageOfMovedPoint}!for a transformation coll}
\label{SmallestImageOfMovedPoint:for a transformation coll}
}\hfill{\scriptsize (method)}}\\
\textbf{\indent Returns:\ }
A positive integer or \texttt{infinity}.



 \texttt{SmallestImageOfMovedPoint} returns the smallest positive integer \texttt{i \texttt{\symbol{94}} \mbox{\texttt{\mdseries\slshape f}}} such that \texttt{i \texttt{\symbol{94}} \mbox{\texttt{\mdseries\slshape f}} {\textless}{\textgreater} i} if such an \texttt{i} exists. If \mbox{\texttt{\mdseries\slshape f}} is the identity transformation, then \texttt{infinity} is returned.

 If the argument is a collection of transformations \mbox{\texttt{\mdseries\slshape coll}}, then the smallest integer which is the image a point moved by at least one
element of \mbox{\texttt{\mdseries\slshape coll}} is returned, if such a point exists. If \mbox{\texttt{\mdseries\slshape coll}} only contains identity transformations, then \texttt{SmallestImageOfMovedPoint} returns \texttt{infinity}. 
\begin{Verbatim}[commandchars=!@|,fontsize=\small,frame=single,label=Example]
  !gapprompt@gap>| !gapinput@S := FullTransformationSemigroup( 5 );|
  <full transformation monoid of degree 5>
  !gapprompt@gap>| !gapinput@SmallestImageOfMovedPoint( S );|
  1
  !gapprompt@gap>| !gapinput@S := Semigroup( IdentityTransformation );|
  <trivial transformation group of degree 0 with 1 generator>
  !gapprompt@gap>| !gapinput@SmallestImageOfMovedPoint( S );|
  infinity
  !gapprompt@gap>| !gapinput@f := Transformation( [ 1, 2, 3, 6, 6, 6 ] );;|
  !gapprompt@gap>| !gapinput@SmallestImageOfMovedPoint( f );|
  6
\end{Verbatim}
 }

  

\subsection{\textcolor{Chapter }{LargestImageOfMovedPoint (for a transformation)}}
\logpage{[ 53, 5, 10 ]}\nobreak
\hyperdef{L}{X7E7172567C3A3E63}{}
{\noindent\textcolor{FuncColor}{$\triangleright$\enspace\texttt{LargestImageOfMovedPoint({\mdseries\slshape f})\index{LargestImageOfMovedPoint@\texttt{LargestImageOfMovedPoint}!for a transformation}
\label{LargestImageOfMovedPoint:for a transformation}
}\hfill{\scriptsize (attribute)}}\\
\noindent\textcolor{FuncColor}{$\triangleright$\enspace\texttt{LargestImageOfMovedPoint({\mdseries\slshape coll})\index{LargestImageOfMovedPoint@\texttt{LargestImageOfMovedPoint}!for a transformation coll}
\label{LargestImageOfMovedPoint:for a transformation coll}
}\hfill{\scriptsize (method)}}\\
\textbf{\indent Returns:\ }
A positive integer.



 \texttt{LargestImageOfMovedPoint} returns the largest positive integer \texttt{i \texttt{\symbol{94}} \mbox{\texttt{\mdseries\slshape f}}} such that \texttt{i \texttt{\symbol{94}} \mbox{\texttt{\mdseries\slshape f}} {\textless}{\textgreater} i} if such an \texttt{i} exists. If \mbox{\texttt{\mdseries\slshape f}} is the identity transformation, then \texttt{0} is returned. 

 If the argument is a collection of transformations \mbox{\texttt{\mdseries\slshape coll}}, then the largest integer which is the image a point moved by at least one
element of \mbox{\texttt{\mdseries\slshape coll}} is returned, if such a point exists. If \mbox{\texttt{\mdseries\slshape coll}} only contains identity transformations, then \texttt{LargestImageOfMovedPoint} returns \texttt{0}. 
\begin{Verbatim}[commandchars=!@|,fontsize=\small,frame=single,label=Example]
  !gapprompt@gap>| !gapinput@S := FullTransformationSemigroup( 5 );|
  <full transformation monoid of degree 5>
  !gapprompt@gap>| !gapinput@LargestImageOfMovedPoint( S );|
  5
  !gapprompt@gap>| !gapinput@S := Semigroup( IdentityTransformation );;|
  !gapprompt@gap>| !gapinput@LargestImageOfMovedPoint( S );|
  0
  !gapprompt@gap>| !gapinput@f := Transformation( [ 1, 2, 3, 6, 6, 6 ] );;|
  !gapprompt@gap>| !gapinput@LargestImageOfMovedPoint( f );|
  6
\end{Verbatim}
 }

  

\subsection{\textcolor{Chapter }{FlatKernelOfTransformation}}
\logpage{[ 53, 5, 11 ]}\nobreak
\hyperdef{L}{X8083794579274E87}{}
{\noindent\textcolor{FuncColor}{$\triangleright$\enspace\texttt{FlatKernelOfTransformation({\mdseries\slshape f[, n]})\index{FlatKernelOfTransformation@\texttt{FlatKernelOfTransformation}}
\label{FlatKernelOfTransformation}
}\hfill{\scriptsize (attribute)}}\\
\textbf{\indent Returns:\ }
The flat kernel of a transformation.



 If the kernel classes of the transformation \mbox{\texttt{\mdseries\slshape f}} on \texttt{[ 1 .. \mbox{\texttt{\mdseries\slshape n}} ]} are $K_1, \dots, K_r$, then \texttt{FlatKernelOfTransformation} returns a list \texttt{L} such that \texttt{L[i] = j} for all \texttt{i} in $K_j$. For a given transformation and positive integer \mbox{\texttt{\mdseries\slshape n}}, there is a unique such list. 

 If the optional second argument \mbox{\texttt{\mdseries\slshape n}} is not present, then the degree of \mbox{\texttt{\mdseries\slshape f}} is used by default. 
\begin{Verbatim}[commandchars=!@|,fontsize=\small,frame=single,label=Example]
  !gapprompt@gap>| !gapinput@f := Transformation( [ 10, 3, 7, 10, 1, 5, 9, 2, 6, 10 ] );;|
  !gapprompt@gap>| !gapinput@FlatKernelOfTransformation( f );|
  [ 1, 2, 3, 1, 4, 5, 6, 7, 8, 1 ]
\end{Verbatim}
 }

  

\subsection{\textcolor{Chapter }{KernelOfTransformation}}
\logpage{[ 53, 5, 12 ]}\nobreak
\hyperdef{L}{X80FCB5048789CF75}{}
{\noindent\textcolor{FuncColor}{$\triangleright$\enspace\texttt{KernelOfTransformation({\mdseries\slshape f[, n, bool]})\index{KernelOfTransformation@\texttt{KernelOfTransformation}}
\label{KernelOfTransformation}
}\hfill{\scriptsize (attribute)}}\\
\textbf{\indent Returns:\ }
The kernel of a transformation.



 When the arguments are a transformation \mbox{\texttt{\mdseries\slshape f}}, a positive integer \mbox{\texttt{\mdseries\slshape n}}, and \texttt{true}, \texttt{KernelOfTransformation} returns the kernel of the transformation \mbox{\texttt{\mdseries\slshape f}} on \texttt{[ 1 .. \mbox{\texttt{\mdseries\slshape n}} ]} as a set of sets of positive integers. If the argument \mbox{\texttt{\mdseries\slshape bool}} is \texttt{false}, then only the non\texttt{\symbol{45}}singleton classes are returned. 

 The second and third arguments are optional, the default values are the degree
of \mbox{\texttt{\mdseries\slshape f}} and \texttt{true}. 
\begin{Verbatim}[commandchars=!@|,fontsize=\small,frame=single,label=Example]
  !gapprompt@gap>| !gapinput@f := Transformation( [ 2, 6, 7, 2, 6, 9, 9, 1, 11, 1, 12, 5 ] );;|
  !gapprompt@gap>| !gapinput@KernelOfTransformation( f );|
  [ [ 1, 4 ], [ 2, 5 ], [ 3 ], [ 6, 7 ], [ 8, 10 ], [ 9 ], [ 11 ],
    [ 12 ] ]
  !gapprompt@gap>| !gapinput@KernelOfTransformation( f, 5 );|
  [ [ 1, 4 ], [ 2, 5 ], [ 3 ] ]
  !gapprompt@gap>| !gapinput@KernelOfTransformation( f, 5, false );|
  [ [ 1, 4 ], [ 2, 5 ] ]
  !gapprompt@gap>| !gapinput@KernelOfTransformation( f, 15 );|
  [ [ 1, 4 ], [ 2, 5 ], [ 3 ], [ 6, 7 ], [ 8, 10 ], [ 9 ], [ 11 ],
    [ 12 ], [ 13 ], [ 14 ], [ 15 ] ]
  !gapprompt@gap>| !gapinput@KernelOfTransformation( f, false );|
  [ [ 1, 4 ], [ 2, 5 ], [ 6, 7 ], [ 8, 10 ] ]
\end{Verbatim}
 }

  

\subsection{\textcolor{Chapter }{InverseOfTransformation}}
\logpage{[ 53, 5, 13 ]}\nobreak
\hyperdef{L}{X860306EB7FAAD2D4}{}
{\noindent\textcolor{FuncColor}{$\triangleright$\enspace\texttt{InverseOfTransformation({\mdseries\slshape f})\index{InverseOfTransformation@\texttt{InverseOfTransformation}}
\label{InverseOfTransformation}
}\hfill{\scriptsize (function)}}\\
\textbf{\indent Returns:\ }
A transformation.



 \texttt{InverseOfTransformation} returns a semigroup inverse of the transformation \mbox{\texttt{\mdseries\slshape f}} in the full transformation semigroup. An \emph{inverse} of \mbox{\texttt{\mdseries\slshape f}} is any transformation \texttt{g} such that \texttt{\mbox{\texttt{\mdseries\slshape f}} * g * \mbox{\texttt{\mdseries\slshape f}} = \mbox{\texttt{\mdseries\slshape f}}} and \texttt{g * \mbox{\texttt{\mdseries\slshape f}} * g = g}. Every transformation has at least one inverse. 
\begin{Verbatim}[commandchars=!@|,fontsize=\small,frame=single,label=Example]
  !gapprompt@gap>| !gapinput@f := Transformation( [ 2, 6, 7, 2, 6, 9, 9, 1, 1, 5 ] );;|
  !gapprompt@gap>| !gapinput@g := InverseOfTransformation( f );|
  Transformation( [ 8, 1, 1, 1, 10, 2, 3, 1, 6, 1 ] )
  !gapprompt@gap>| !gapinput@f * g * f;|
  Transformation( [ 2, 6, 7, 2, 6, 9, 9, 1, 1, 5 ] )
  !gapprompt@gap>| !gapinput@g * f * g;|
  Transformation( [ 8, 1, 1, 1, 10, 2, 3, 1, 6, 1 ] )
\end{Verbatim}
 }

 

\subsection{\textcolor{Chapter }{Inverse (for a transformation)}}
\logpage{[ 53, 5, 14 ]}\nobreak
\hyperdef{L}{X7BB9DB6E8558356D}{}
{\noindent\textcolor{FuncColor}{$\triangleright$\enspace\texttt{Inverse({\mdseries\slshape f})\index{Inverse@\texttt{Inverse}!for a transformation}
\label{Inverse:for a transformation}
}\hfill{\scriptsize (attribute)}}\\
\textbf{\indent Returns:\ }
A transformation.



 If the transformation \mbox{\texttt{\mdseries\slshape f}} is a bijection, then \texttt{Inverse} or \texttt{\mbox{\texttt{\mdseries\slshape f}} \texttt{\symbol{94}} \texttt{\symbol{45}}1} returns the inverse of \mbox{\texttt{\mdseries\slshape f}}. If \mbox{\texttt{\mdseries\slshape f}} is not a bijection, then \texttt{fail} is returned. 
\begin{Verbatim}[commandchars=!@|,fontsize=\small,frame=single,label=Example]
  !gapprompt@gap>| !gapinput@Transformation( [ 3, 8, 12, 1, 11, 9, 9, 4, 10, 5, 10, 6 ] ) ^ -1;|
  fail
  !gapprompt@gap>| !gapinput@Transformation( [ 2, 3, 1 ] ) ^ -1;|
  Transformation( [ 3, 1, 2 ] )
\end{Verbatim}
 }

  

\subsection{\textcolor{Chapter }{IndexPeriodOfTransformation}}
\logpage{[ 53, 5, 15 ]}\nobreak
\hyperdef{L}{X863216CB7AF88BED}{}
{\noindent\textcolor{FuncColor}{$\triangleright$\enspace\texttt{IndexPeriodOfTransformation({\mdseries\slshape f})\index{IndexPeriodOfTransformation@\texttt{IndexPeriodOfTransformation}}
\label{IndexPeriodOfTransformation}
}\hfill{\scriptsize (function)}}\\
\textbf{\indent Returns:\ }
A pair of positive integers.



 Returns the least positive integers \texttt{m} and \texttt{r} such that \texttt{\mbox{\texttt{\mdseries\slshape f}} \texttt{\symbol{94}} (m + r) = \mbox{\texttt{\mdseries\slshape f}} \texttt{\symbol{94}} m}, which are known as the \emph{index} and \emph{period} of the transformation \mbox{\texttt{\mdseries\slshape f}}. 
\begin{Verbatim}[commandchars=!@|,fontsize=\small,frame=single,label=Example]
  !gapprompt@gap>| !gapinput@f := Transformation( [ 3, 4, 4, 6, 1, 3, 3, 7, 1 ] );;|
  !gapprompt@gap>| !gapinput@IndexPeriodOfTransformation( f );|
  [ 2, 3 ]
  !gapprompt@gap>| !gapinput@f ^ 2 = f ^ 5;|
  true
  !gapprompt@gap>| !gapinput@IndexPeriodOfTransformation( IdentityTransformation );|
  [ 1, 1 ]
  !gapprompt@gap>| !gapinput@IndexPeriodOfTransformation( Transformation( [ 1, 2, 1 ] ) );|
  [ 1, 1 ]
  !gapprompt@gap>| !gapinput@IndexPeriodOfTransformation( Transformation( [ 1, 2, 3 ] ) );|
  [ 1, 1 ]
  !gapprompt@gap>| !gapinput@IndexPeriodOfTransformation( Transformation( [ 1, 3, 2 ] ) );|
  [ 1, 2 ]
\end{Verbatim}
 }

  

\subsection{\textcolor{Chapter }{SmallestIdempotentPower (for a transformation)}}
\logpage{[ 53, 5, 16 ]}\nobreak
\hyperdef{L}{X85FE9F20810BCC70}{}
{\noindent\textcolor{FuncColor}{$\triangleright$\enspace\texttt{SmallestIdempotentPower({\mdseries\slshape f})\index{SmallestIdempotentPower@\texttt{SmallestIdempotentPower}!for a transformation}
\label{SmallestIdempotentPower:for a transformation}
}\hfill{\scriptsize (attribute)}}\\
\textbf{\indent Returns:\ }
A positive integer.



 This function returns the least positive integer \texttt{n} such that the transformation \texttt{\mbox{\texttt{\mdseries\slshape f}} \texttt{\symbol{94}} n} is an idempotent. The smallest idempotent power of \mbox{\texttt{\mdseries\slshape f}} is the least multiple of the period of \mbox{\texttt{\mdseries\slshape f}} that is greater than or equal to the index of \mbox{\texttt{\mdseries\slshape f}}; see \texttt{IndexPeriodOfTransformation} (\ref{IndexPeriodOfTransformation}). 
\begin{Verbatim}[commandchars=!@|,fontsize=\small,frame=single,label=Example]
  !gapprompt@gap>| !gapinput@f := Transformation( [ 6, 7, 4, 1, 7, 4, 6, 1, 3, 4 ] );;|
  !gapprompt@gap>| !gapinput@SmallestIdempotentPower( f );|
  3
  !gapprompt@gap>| !gapinput@f := Transformation( [ 6, 6, 6, 2, 7, 1, 5, 3, 10, 6 ] );;|
  !gapprompt@gap>| !gapinput@SmallestIdempotentPower( f );|
  2
\end{Verbatim}
 }

  

\subsection{\textcolor{Chapter }{ComponentsOfTransformation}}
\logpage{[ 53, 5, 17 ]}\nobreak
\hyperdef{L}{X858E944481F6B591}{}
{\noindent\textcolor{FuncColor}{$\triangleright$\enspace\texttt{ComponentsOfTransformation({\mdseries\slshape f})\index{ComponentsOfTransformation@\texttt{ComponentsOfTransformation}}
\label{ComponentsOfTransformation}
}\hfill{\scriptsize (attribute)}}\\
\textbf{\indent Returns:\ }
A list of lists of positive integers.



 \texttt{ComponentsOfTransformation} returns a list of the components of the transformation \mbox{\texttt{\mdseries\slshape f}}. Each component is a subset of \texttt{[ 1 .. DegreeOfTransformation( f ) ]}, and the union of the components is \texttt{[ 1 .. DegreeOfTransformation( f ) ]}. 
\begin{Verbatim}[commandchars=!@|,fontsize=\small,frame=single,label=Example]
  !gapprompt@gap>| !gapinput@f := Transformation( [ 6, 12, 11, 1, 7, 6, 2, 8, 4, 7, 5, 12 ] );|
  Transformation( [ 6, 12, 11, 1, 7, 6, 2, 8, 4, 7, 5, 12 ] )
  !gapprompt@gap>| !gapinput@ComponentsOfTransformation( f );|
  [ [ 1, 6, 4, 9 ], [ 2, 12, 3, 11, 5, 7, 10 ], [ 8 ] ]
  !gapprompt@gap>| !gapinput@f := AsTransformation( (1,8,2,4,11,5,10)(3,7)(9,12) );|
  Transformation( [ 8, 4, 7, 11, 10, 6, 3, 2, 12, 1, 5, 9 ] )
  !gapprompt@gap>| !gapinput@ComponentsOfTransformation( f );|
  [ [ 1, 8, 2, 4, 11, 5, 10 ], [ 3, 7 ], [ 6 ], [ 9, 12 ] ]
\end{Verbatim}
 }

  

\subsection{\textcolor{Chapter }{NrComponentsOfTransformation}}
\logpage{[ 53, 5, 18 ]}\nobreak
\hyperdef{L}{X8640AE1C79201470}{}
{\noindent\textcolor{FuncColor}{$\triangleright$\enspace\texttt{NrComponentsOfTransformation({\mdseries\slshape f})\index{NrComponentsOfTransformation@\texttt{NrComponentsOfTransformation}}
\label{NrComponentsOfTransformation}
}\hfill{\scriptsize (attribute)}}\\
\textbf{\indent Returns:\ }
A positive integer.



 \texttt{NrComponentsOfTransformation} returns the number of components of the transformation \mbox{\texttt{\mdseries\slshape f}} on the range \texttt{[ 1 .. DegreeOfTransformation( \mbox{\texttt{\mdseries\slshape f}} ) ]}. 
\begin{Verbatim}[commandchars=!@|,fontsize=\small,frame=single,label=Example]
  !gapprompt@gap>| !gapinput@f := Transformation( [ 6, 12, 11, 1, 7, 6, 2, 8, 4, 7, 5, 12 ] );|
  Transformation( [ 6, 12, 11, 1, 7, 6, 2, 8, 4, 7, 5, 12 ] )
  !gapprompt@gap>| !gapinput@NrComponentsOfTransformation( f );|
  3
  !gapprompt@gap>| !gapinput@f := AsTransformation( (1,8,2,4,11,5,10)(3,7)(9,12) );|
  Transformation( [ 8, 4, 7, 11, 10, 6, 3, 2, 12, 1, 5, 9 ] )
  !gapprompt@gap>| !gapinput@NrComponentsOfTransformation( f );|
  4
\end{Verbatim}
 }

  

\subsection{\textcolor{Chapter }{ComponentRepsOfTransformation}}
\logpage{[ 53, 5, 19 ]}\nobreak
\hyperdef{L}{X784650B583CEAF7D}{}
{\noindent\textcolor{FuncColor}{$\triangleright$\enspace\texttt{ComponentRepsOfTransformation({\mdseries\slshape f})\index{ComponentRepsOfTransformation@\texttt{ComponentRepsOfTransformation}}
\label{ComponentRepsOfTransformation}
}\hfill{\scriptsize (attribute)}}\\
\textbf{\indent Returns:\ }
A list of lists of positive integers.



 \texttt{ComponentRepsOfTransformation} returns the representatives, in the following sense, of the components of the
transformation \mbox{\texttt{\mdseries\slshape f}}. For every \texttt{i} in \texttt{[ 1 .. DegreeOfTransformation( f ) ]} there exists a representative \texttt{j} and a positive integer \texttt{k} such that \texttt{i \texttt{\symbol{94}} (\mbox{\texttt{\mdseries\slshape f}} \texttt{\symbol{94}} k) = j}. The representatives returned by \texttt{ComponentRepsOfTransformation} are partitioned according to the component they belong to. \texttt{ComponentRepsOfTransformation} returns the least number of representatives. 
\begin{Verbatim}[commandchars=!@|,fontsize=\small,frame=single,label=Example]
  !gapprompt@gap>| !gapinput@f := Transformation( [ 6, 12, 11, 1, 7, 6, 2, 8, 4, 7, 5, 12 ] );|
  Transformation( [ 6, 12, 11, 1, 7, 6, 2, 8, 4, 7, 5, 12 ] )
  !gapprompt@gap>| !gapinput@ComponentRepsOfTransformation( f );|
  [ [ 3, 10 ], [ 9 ], [ 8 ] ]
  !gapprompt@gap>| !gapinput@f := AsTransformation( (1,8,2,4,11,5,10)(3,7)(9,12) );|
  Transformation( [ 8, 4, 7, 11, 10, 6, 3, 2, 12, 1, 5, 9 ] )
  !gapprompt@gap>| !gapinput@ComponentRepsOfTransformation( f );|
  [ [ 1 ], [ 3 ], [ 6 ], [ 9 ] ]
\end{Verbatim}
 }

  

\subsection{\textcolor{Chapter }{CyclesOfTransformation}}
\logpage{[ 53, 5, 20 ]}\nobreak
\hyperdef{L}{X7EAA15557D55D93B}{}
{\noindent\textcolor{FuncColor}{$\triangleright$\enspace\texttt{CyclesOfTransformation({\mdseries\slshape f[, list]})\index{CyclesOfTransformation@\texttt{CyclesOfTransformation}}
\label{CyclesOfTransformation}
}\hfill{\scriptsize (attribute)}}\\
\textbf{\indent Returns:\ }
A list of lists of positive integers.



 When the arguments of this function are a transformation \mbox{\texttt{\mdseries\slshape f}} and a list \mbox{\texttt{\mdseries\slshape list}}, it returns a list of the cycles of the components of \mbox{\texttt{\mdseries\slshape f}} containing any element of \mbox{\texttt{\mdseries\slshape list}}.

 If the optional second argument is not present, then the range \texttt{[ 1 .. DegreeOfTransformation( \mbox{\texttt{\mdseries\slshape f}} ) ]} is used as the default value for \mbox{\texttt{\mdseries\slshape list}}. 
\begin{Verbatim}[commandchars=!@|,fontsize=\small,frame=single,label=Example]
  !gapprompt@gap>| !gapinput@f := Transformation( [ 6, 12, 11, 1, 7, 6, 2, 8, 4, 7, 5, 12 ] );|
  Transformation( [ 6, 12, 11, 1, 7, 6, 2, 8, 4, 7, 5, 12 ] )
  !gapprompt@gap>| !gapinput@CyclesOfTransformation( f );|
  [ [ 6 ], [ 12 ], [ 8 ] ]
  !gapprompt@gap>| !gapinput@CyclesOfTransformation( f, [ 1, 2, 4 ] );|
  [ [ 6 ], [ 12 ] ]
  !gapprompt@gap>| !gapinput@CyclesOfTransformation( f, [ 1 .. 17 ] );|
  [ [ 6 ], [ 12 ], [ 8 ], [ 13 ], [ 14 ], [ 15 ], [ 16 ], [ 17 ] ]
\end{Verbatim}
 }

  

\subsection{\textcolor{Chapter }{CycleTransformationInt}}
\logpage{[ 53, 5, 21 ]}\nobreak
\hyperdef{L}{X786EB02A829260DB}{}
{\noindent\textcolor{FuncColor}{$\triangleright$\enspace\texttt{CycleTransformationInt({\mdseries\slshape f, n})\index{CycleTransformationInt@\texttt{CycleTransformationInt}}
\label{CycleTransformationInt}
}\hfill{\scriptsize (operation)}}\\
\textbf{\indent Returns:\ }
A list of positive integers.



 If \mbox{\texttt{\mdseries\slshape f}} is a transformation and \mbox{\texttt{\mdseries\slshape n}} is a positive integer, then \texttt{CycleTransformationInt} returns the cycle of the component of \mbox{\texttt{\mdseries\slshape f}} containing \mbox{\texttt{\mdseries\slshape n}}. 
\begin{Verbatim}[commandchars=!@|,fontsize=\small,frame=single,label=Example]
  !gapprompt@gap>| !gapinput@f := Transformation( [ 6, 2, 8, 4, 7, 5, 8, 3, 5, 8 ] );;|
  !gapprompt@gap>| !gapinput@CycleTransformationInt( f, 1 );|
  [ 8, 3 ]
  !gapprompt@gap>| !gapinput@CycleTransformationInt( f, 12 );|
  [ 12 ]
  !gapprompt@gap>| !gapinput@CycleTransformationInt( f, 5 );|
  [ 8, 3 ]
\end{Verbatim}
 }

  

\subsection{\textcolor{Chapter }{LeftOne (for a transformation)}}
\logpage{[ 53, 5, 22 ]}\nobreak
\hyperdef{L}{X845869E0815A6AA6}{}
{\noindent\textcolor{FuncColor}{$\triangleright$\enspace\texttt{LeftOne({\mdseries\slshape f})\index{LeftOne@\texttt{LeftOne}!for a transformation}
\label{LeftOne:for a transformation}
}\hfill{\scriptsize (attribute)}}\\
\noindent\textcolor{FuncColor}{$\triangleright$\enspace\texttt{RightOne({\mdseries\slshape f})\index{RightOne@\texttt{RightOne}!for a transformation}
\label{RightOne:for a transformation}
}\hfill{\scriptsize (attribute)}}\\
\textbf{\indent Returns:\ }
A transformation.



 \texttt{LeftOne} returns an idempotent transformation \texttt{e} such that the kernel (with respect to the degree of \mbox{\texttt{\mdseries\slshape f}}) of \texttt{e} equals the kernel of the transformation \mbox{\texttt{\mdseries\slshape f}} and \texttt{e * \mbox{\texttt{\mdseries\slshape f}} = f}. 

 \texttt{RightOne} returns an idempotent transformation \texttt{e} such that the image set (with respect to the degree of \mbox{\texttt{\mdseries\slshape f}}) of \texttt{e} equals the image set of \mbox{\texttt{\mdseries\slshape f}} and \texttt{\mbox{\texttt{\mdseries\slshape f}} * e = f}. 
\begin{Verbatim}[commandchars=!@|,fontsize=\small,frame=single,label=Example]
  !gapprompt@gap>| !gapinput@f := Transformation( [ 11, 10, 2, 11, 4, 4, 7, 6, 9, 10, 1, 11 ] );;|
  !gapprompt@gap>| !gapinput@e := RightOne( f );|
  Transformation( [ 1, 2, 2, 4, 4, 6, 7, 7, 9, 10, 11, 11 ] )
  !gapprompt@gap>| !gapinput@IsIdempotent( e );|
  true
  !gapprompt@gap>| !gapinput@f * e = f;|
  true
  !gapprompt@gap>| !gapinput@e := LeftOne( f );|
  Transformation( [ 1, 2, 3, 1, 5, 5, 7, 8, 9, 2, 11, 1 ] )
  !gapprompt@gap>| !gapinput@e * f = f;|
  true
  !gapprompt@gap>| !gapinput@IsIdempotent( e );|
  true
\end{Verbatim}
 }

 

\subsection{\textcolor{Chapter }{TrimTransformation}}
\logpage{[ 53, 5, 23 ]}\nobreak
\hyperdef{L}{X7F19C9C77F9F8981}{}
{\noindent\textcolor{FuncColor}{$\triangleright$\enspace\texttt{TrimTransformation({\mdseries\slshape f[, n]})\index{TrimTransformation@\texttt{TrimTransformation}}
\label{TrimTransformation}
}\hfill{\scriptsize (operation)}}\\
\textbf{\indent Returns:\ }
Nothing.



 It can happen that the internal representation of a transformation uses more
memory than necessary. For example, this can happen when composing
transformations where it is possible that the resulting transformation \mbox{\texttt{\mdseries\slshape f}} belongs to \texttt{IsTrans4Rep} and stores its images as 32\texttt{\symbol{45}}bit integers, while none of its
moved points exceeds 65536. The purpose of \texttt{TrimTransformation} is to change the internal representation of such an \mbox{\texttt{\mdseries\slshape f}} to remove the trailing fixed points in the internal representation of \mbox{\texttt{\mdseries\slshape f}}. 

 If the optional second argument \mbox{\texttt{\mdseries\slshape n}} is provided, then the internal representation of \mbox{\texttt{\mdseries\slshape f}} is reduced to the images of the first \mbox{\texttt{\mdseries\slshape n}} positive integers. Please note that it must be the case that \texttt{i \texttt{\symbol{94}} \mbox{\texttt{\mdseries\slshape f}} {\textless}= n} for all \texttt{i} in the range \texttt{[ 1 .. \mbox{\texttt{\mdseries\slshape n}} ]} otherwise the resulting object will not define a transformation. 

 If the optional second argument is not included, then the degree of \mbox{\texttt{\mdseries\slshape f}} is used by default. 

 The transformation \mbox{\texttt{\mdseries\slshape f}} is changed in\texttt{\symbol{45}}place, and nothing is returned by this
function. 
\begin{Verbatim}[commandchars=!@|,fontsize=\small,frame=single,label=Example]
  !gapprompt@gap>| !gapinput@f := Transformation( [ 1 .. 2 ^ 16 ], x -> x + 1 );|
  <transformation on 65537 pts with rank 65536>
  !gapprompt@gap>| !gapinput@g := Transformation( [ 1 .. 2 ^ 16 + 1 ],|
  !gapprompt@>| !gapinput@function( x )|
  !gapprompt@>| !gapinput@  if x = 1 or x = 65537 then|
  !gapprompt@>| !gapinput@    return x;|
  !gapprompt@>| !gapinput@  else|
  !gapprompt@>| !gapinput@    return x - 1;|
  !gapprompt@>| !gapinput@  fi;|
  !gapprompt@>| !gapinput@end );|
  <transformation on 65536 pts with rank 65535>
  !gapprompt@gap>| !gapinput@h := g * f;|
  Transformation( [ 2, 2 ] )
  !gapprompt@gap>| !gapinput@DegreeOfTransformation( h ); IsTrans4Rep( h ); MemoryUsage( h );|
  65537
  true
  262188
  !gapprompt@gap>| !gapinput@TrimTransformation( h ); h;|
  Transformation( [ 2, 2 ] )
  !gapprompt@gap>| !gapinput@DegreeOfTransformation( h ); IsTrans4Rep( h ); MemoryUsage( h );|
  2
  false
  44
\end{Verbatim}
 }

 }

  
\section{\textcolor{Chapter }{Displaying transformations}}\label{sect:DisplayingTransformations}
\logpage{[ 53, 6, 0 ]}
\hyperdef{L}{X810D23017A5527B7}{}
{
  It is possible to change the way that \textsf{GAP} displays transformations using the user preferences \texttt{TransformationDisplayLimit} and \texttt{NotationForTransformations}; see Section \texttt{UserPreference} (\ref{UserPreference}) for more information about user preferences. 

 If \texttt{f} is a transformation where the degree \texttt{n} of \texttt{f} exceeds the value of the user preference \texttt{TransformationDisplayLimit}, then \texttt{f} is displayed as: 
\begin{Verbatim}[commandchars=!@|,fontsize=\small,frame=single,label=Example]
  <transformation on n pts with rank r>
\end{Verbatim}
 where \texttt{r} is the rank of \texttt{f} relative to \texttt{n}. The idea is to abbreviate the display of transformations defined on many
points. The default value for the \texttt{TransformationDisplayLimit} is \texttt{100}. 

 If the degree of \texttt{f} does not exceed the value of \texttt{TransformationDisplayLimit}, then how \texttt{f} is displayed depends on the value of the user preference \texttt{NotationForTransformations}. 

 There are two possible values for \texttt{NotationForTransformations}: 
\begin{description}
\item[{input}]  With this option a transformation \mbox{\texttt{\mdseries\slshape f}} is displayed in as: \texttt{Transformation( ImageListOfTransformation( \mbox{\texttt{\mdseries\slshape f}}, n ) )} where \texttt{n} is the degree of \mbox{\texttt{\mdseries\slshape f}}. The only exception is the identity transformation, which is displayed as: \texttt{IdentityTransformation}. 
\item[{fr}] With this option a transformation \mbox{\texttt{\mdseries\slshape f}} is displayed in as: \texttt{{\textless}transformation: ImageListOfTransformation( \mbox{\texttt{\mdseries\slshape f}}, n ){\textgreater}} where \texttt{n} is the largest moved point of \mbox{\texttt{\mdseries\slshape f}}. The only exception is the identity transformation, which is displayed as: \texttt{{\textless}identity transformation{\textgreater}}. 
\end{description}
 
\begin{Verbatim}[commandchars=!@|,fontsize=\small,frame=single,label=Example]
  !gapprompt@gap>| !gapinput@SetUserPreference( "TransformationDisplayLimit", 12 );|
  !gapprompt@gap>| !gapinput@f := Transformation( [ 3, 8, 12, 1, 11, 9, 9, 4, 10, 5, 10, 6 ] );|
  <transformation on 12 pts with rank 10>
  !gapprompt@gap>| !gapinput@SetUserPreference( "TransformationDisplayLimit", 100 );|
  !gapprompt@gap>| !gapinput@f;|
  Transformation( [ 3, 8, 12, 1, 11, 9, 9, 4, 10, 5, 10, 6 ] )
  !gapprompt@gap>| !gapinput@SetUserPreference( "NotationForTransformations", "fr" );|
  !gapprompt@gap>| !gapinput@f;|
  <transformation: 3,8,12,1,11,9,9,4,10,5,10,6>
\end{Verbatim}
 }

  
\section{\textcolor{Chapter }{Semigroups of transformations}}\label{Making transformation semigroups}
\logpage{[ 53, 7, 0 ]}
\hyperdef{L}{X7B51CE257B814B09}{}
{
  As mentioned at the start of the chapter, every semigroup is isomorphic to a
semigroup of transformations, and in this section we describe the functions in \textsf{GAP} specific to transformation semigroups. For more information about semigroups
in general see Chapter \ref{Semigroups}. 

 The \textsf{Semigroups} package contains many additional functions and methods for computing with
semigroups of transformations. In particular, \textsf{Semigroups} contains more efficient methods than those available in the \textsf{GAP} library (and in many cases more efficient than any other software) for
creating semigroups of transformations, calculating their Green's classes,
size, elements, group of units, minimal ideal, small generating sets, testing
membership, finding the inverses of a regular element, factorizing elements
over the generators, and more. Since a transformation semigroup is also a
transformation collection, there are special methods for \texttt{MovedPoints} (\ref{MovedPoints:for a transformation coll}), \texttt{NrMovedPoints} (\ref{NrMovedPoints:for a transformation coll}), \texttt{LargestMovedPoint} (\ref{LargestMovedPoint:for a transformation coll}), \texttt{SmallestMovedPoint} (\ref{SmallestMovedPoint:for a transformation coll}), \texttt{LargestImageOfMovedPoint} (\ref{LargestImageOfMovedPoint:for a transformation coll}), and \texttt{SmallestImageOfMovedPoint} (\ref{SmallestImageOfMovedPoint:for a transformation coll}), when applied to a transformation semigroup.  

\subsection{\textcolor{Chapter }{IsTransformationSemigroup}}
\logpage{[ 53, 7, 1 ]}\nobreak
\hyperdef{L}{X7EAF835D7FE4026F}{}
{\noindent\textcolor{FuncColor}{$\triangleright$\enspace\texttt{IsTransformationSemigroup({\mdseries\slshape obj})\index{IsTransformationSemigroup@\texttt{IsTransformationSemigroup}}
\label{IsTransformationSemigroup}
}\hfill{\scriptsize (Synonym)}}\\
\noindent\textcolor{FuncColor}{$\triangleright$\enspace\texttt{IsTransformationMonoid({\mdseries\slshape obj})\index{IsTransformationMonoid@\texttt{IsTransformationMonoid}}
\label{IsTransformationMonoid}
}\hfill{\scriptsize (Synonym)}}\\
\textbf{\indent Returns:\ }
\texttt{true} or \texttt{false}.



 A \emph{transformation semigroup} is simply a semigroup consisting of transformations. An object \mbox{\texttt{\mdseries\slshape obj}} is a transformation semigroup in \textsf{GAP} if it satisfies \texttt{IsSemigroup} (\ref{IsSemigroup}) and \texttt{IsTransformationCollection} (\ref{IsTransformationCollection}). 

 A \emph{transformation monoid} is a monoid consisting of transformations. An object \mbox{\texttt{\mdseries\slshape obj}} is a transformation monoid in \textsf{GAP} if it satisfies \texttt{IsMonoid} (\ref{IsMonoid}) and \texttt{IsTransformationCollection} (\ref{IsTransformationCollection}). 

 Note that it is possible for a transformation semigroup to have a
multiplicative neutral element (i.e. an identity element) but not to satisfy \texttt{IsTransformationMonoid}. For example, 
\begin{Verbatim}[commandchars=!@|,fontsize=\small,frame=single,label=Example]
  !gapprompt@gap>| !gapinput@f := Transformation( [ 2, 6, 7, 2, 6, 9, 9, 1, 1, 5 ] );;|
  !gapprompt@gap>| !gapinput@S := Semigroup( f, One( f ) );|
  <commutative transformation monoid of degree 10 with 1 generator>
  !gapprompt@gap>| !gapinput@IsMonoid( S );|
  true
  !gapprompt@gap>| !gapinput@IsTransformationMonoid( S );|
  true
  !gapprompt@gap>| !gapinput@S := Semigroup(|
  !gapprompt@>| !gapinput@Transformation( [ 3, 8, 1, 4, 5, 6, 7, 1, 10, 10 ] ),|
  !gapprompt@>| !gapinput@Transformation( [ 1, 2, 3, 4, 5, 6, 7, 8, 10, 10 ] ) );|
  <transformation semigroup of degree 10 with 2 generators>
  !gapprompt@gap>| !gapinput@One( S );|
  fail
  !gapprompt@gap>| !gapinput@MultiplicativeNeutralElement( S );|
  Transformation( [ 1, 2, 3, 4, 5, 6, 7, 8, 10, 10 ] )
  !gapprompt@gap>| !gapinput@IsMonoid( S );|
  false
\end{Verbatim}
 In this example \texttt{S} cannot be converted into a monoid using \texttt{AsMonoid} (\ref{AsMonoid}) since the \texttt{One} (\ref{One}) of any element in \texttt{S} differs from the multiplicative neutral element. 

 For more details see \texttt{IsMagmaWithOne} (\ref{IsMagmaWithOne}). }

 

\subsection{\textcolor{Chapter }{DegreeOfTransformationSemigroup}}
\logpage{[ 53, 7, 2 ]}\nobreak
\hyperdef{L}{X7EA699C687952544}{}
{\noindent\textcolor{FuncColor}{$\triangleright$\enspace\texttt{DegreeOfTransformationSemigroup({\mdseries\slshape S})\index{DegreeOfTransformationSemigroup@\texttt{DegreeOfTransformationSemigroup}}
\label{DegreeOfTransformationSemigroup}
}\hfill{\scriptsize (attribute)}}\\
\textbf{\indent Returns:\ }
A non\texttt{\symbol{45}}negative integer.



 The \emph{degree} of a transformation semigroup \mbox{\texttt{\mdseries\slshape S}} is just the maximum of the degrees of the elements of \mbox{\texttt{\mdseries\slshape S}}. 
\begin{Verbatim}[commandchars=!@|,fontsize=\small,frame=single,label=Example]
  !gapprompt@gap>| !gapinput@S := Semigroup(|
  !gapprompt@>| !gapinput@Transformation( [ 3, 8, 1, 4, 5, 6, 7, 1, 10, 10, 11 ] ),|
  !gapprompt@>| !gapinput@Transformation( [ 1, 2, 3, 4, 5, 6, 7, 8, 1, 1, 11 ] ) );|
  <transformation semigroup of degree 10 with 2 generators>
  !gapprompt@gap>| !gapinput@DegreeOfTransformationSemigroup( S );|
  10
\end{Verbatim}
 }

  

\subsection{\textcolor{Chapter }{FullTransformationSemigroup}}
\logpage{[ 53, 7, 3 ]}\nobreak
\hyperdef{L}{X7D2B0685815B4053}{}
{\noindent\textcolor{FuncColor}{$\triangleright$\enspace\texttt{FullTransformationSemigroup({\mdseries\slshape n})\index{FullTransformationSemigroup@\texttt{FullTransformationSemigroup}}
\label{FullTransformationSemigroup}
}\hfill{\scriptsize (function)}}\\
\noindent\textcolor{FuncColor}{$\triangleright$\enspace\texttt{FullTransformationMonoid({\mdseries\slshape n})\index{FullTransformationMonoid@\texttt{FullTransformationMonoid}}
\label{FullTransformationMonoid}
}\hfill{\scriptsize (function)}}\\
\textbf{\indent Returns:\ }
The full transformation semigroup of degree \mbox{\texttt{\mdseries\slshape n}}.



 If \mbox{\texttt{\mdseries\slshape n}} is a positive integer, then \texttt{FullTransformationSemigroup} returns the monoid consisting of all transformations with degree at most \mbox{\texttt{\mdseries\slshape n}}, called the \emph{full transformation semigroup}. 

 The full transformation semigroup is regular, has \texttt{\mbox{\texttt{\mdseries\slshape n}} \texttt{\symbol{94}} \mbox{\texttt{\mdseries\slshape n}}} elements, and is generated by any set containing transformations that generate
the symmetric group on \mbox{\texttt{\mdseries\slshape n}} points and any transformation of rank \texttt{\mbox{\texttt{\mdseries\slshape n}} \texttt{\symbol{45}} 1}. 

 \texttt{FulTransformationMonoid} is a synonym for \texttt{FullTransformationSemigroup}. 
\begin{Verbatim}[commandchars=!@|,fontsize=\small,frame=single,label=Example]
  !gapprompt@gap>| !gapinput@FullTransformationSemigroup( 1234 );|
  <full transformation monoid of degree 1234>
\end{Verbatim}
 }

  

\subsection{\textcolor{Chapter }{IsFullTransformationSemigroup}}
\logpage{[ 53, 7, 4 ]}\nobreak
\hyperdef{L}{X85C58E1E818C838C}{}
{\noindent\textcolor{FuncColor}{$\triangleright$\enspace\texttt{IsFullTransformationSemigroup({\mdseries\slshape S})\index{IsFullTransformationSemigroup@\texttt{IsFullTransformationSemigroup}}
\label{IsFullTransformationSemigroup}
}\hfill{\scriptsize (property)}}\\
\noindent\textcolor{FuncColor}{$\triangleright$\enspace\texttt{IsFullTransformationMonoid({\mdseries\slshape S})\index{IsFullTransformationMonoid@\texttt{IsFullTransformationMonoid}}
\label{IsFullTransformationMonoid}
}\hfill{\scriptsize (property)}}\\
\textbf{\indent Returns:\ }
\texttt{true} or \texttt{false}.



 If the transformation semigroup \mbox{\texttt{\mdseries\slshape S}} of degree \texttt{n} contains every transformation of degree at most \texttt{n}, then \texttt{IsFullTransformationSemigroup} returns \texttt{true} and otherwise it returns \texttt{false}. 

 \texttt{IsFullTransformationMonoid} is a synonym of \texttt{IsFullTransformationSemigroup}. It is common in the literature for the full transformation monoid to be
referred to as the full transformation semigroup. 
\begin{Verbatim}[commandchars=!@|,fontsize=\small,frame=single,label=Example]
  !gapprompt@gap>| !gapinput@S := Semigroup( AsTransformation( (1,3,4,2), 5 ),|
  !gapprompt@>| !gapinput@                   AsTransformation( (1,3,5), 5 ),|
  !gapprompt@>| !gapinput@                   Transformation( [ 1, 1, 2, 3, 4 ] ) );|
  <transformation semigroup of degree 5 with 3 generators>
  !gapprompt@gap>| !gapinput@IsFullTransformationSemigroup( S );|
  true
  !gapprompt@gap>| !gapinput@S;|
  <full transformation monoid of degree 5>
  !gapprompt@gap>| !gapinput@IsFullTransformationMonoid( S );|
  true
  !gapprompt@gap>| !gapinput@S := FullTransformationSemigroup( 5 );;|
  !gapprompt@gap>| !gapinput@IsFullTransformationSemigroup( S );|
  true
\end{Verbatim}
 }

  

\subsection{\textcolor{Chapter }{IsomorphismTransformationSemigroup}}
\logpage{[ 53, 7, 5 ]}\nobreak
\hyperdef{L}{X78F29C817CF6827F}{}
{\noindent\textcolor{FuncColor}{$\triangleright$\enspace\texttt{IsomorphismTransformationSemigroup({\mdseries\slshape S})\index{IsomorphismTransformationSemigroup@\texttt{IsomorphismTransformationSemigroup}}
\label{IsomorphismTransformationSemigroup}
}\hfill{\scriptsize (attribute)}}\\
\noindent\textcolor{FuncColor}{$\triangleright$\enspace\texttt{IsomorphismTransformationMonoid({\mdseries\slshape S})\index{IsomorphismTransformationMonoid@\texttt{IsomorphismTransformationMonoid}}
\label{IsomorphismTransformationMonoid}
}\hfill{\scriptsize (attribute)}}\\
\textbf{\indent Returns:\ }
An isomorphism to a transformation semigroup or monoid.



 Returns an isomorphism from the finite semigroup \mbox{\texttt{\mdseries\slshape S}} to a transformation semigroup. For most types of objects in \textsf{GAP} the degree of this transformation semigroup will be equal to the size of \mbox{\texttt{\mdseries\slshape S}} plus \texttt{1}. 

 Let \texttt{\mbox{\texttt{\mdseries\slshape S}} \texttt{\symbol{94}} 1} denote the monoid obtained from \mbox{\texttt{\mdseries\slshape S}} by adjoining an identity element. Then \mbox{\texttt{\mdseries\slshape S}} acts faithfully on \texttt{\mbox{\texttt{\mdseries\slshape S}} \texttt{\symbol{94}} 1} by right multiplication, i.e. every element of \mbox{\texttt{\mdseries\slshape S}} describes a transformation on \texttt{1, .. , |S| + 1}. The isomorphism from \mbox{\texttt{\mdseries\slshape S}} to the transformation semigroup described in this way is called the \emph{right regular representation} of \mbox{\texttt{\mdseries\slshape S}}. In most cases, \texttt{IsomorphismTransformationSemigroup} will return the right regular representation of \mbox{\texttt{\mdseries\slshape S}}. 

 As exceptions, if \mbox{\texttt{\mdseries\slshape S}} is a permutation group or a partial perm semigroup, then the elements of \mbox{\texttt{\mdseries\slshape S}} act naturally and faithfully by transformations on the values from \texttt{1} to the largest moved point of \mbox{\texttt{\mdseries\slshape S}}. 

 If \mbox{\texttt{\mdseries\slshape S}} is a finitely presented semigroup, then the Todd\texttt{\symbol{45}}Coxeter
approach will be attempted.

 \texttt{IsomorphismTransformationMonoid} differs from \texttt{IsomorphismTransformationSemigroup} only in that its range is a transformation monoid, and not only a semigroup,
when the semigroup \mbox{\texttt{\mdseries\slshape S}} is a monoid. 
\begin{Verbatim}[commandchars=!@|,fontsize=\small,frame=single,label=Example]
  !gapprompt@gap>| !gapinput@S := Semigroup( [ [ [ Z(3), 0*Z(3) ], [ 0*Z(3), Z(3) ^ 0 ] ],|
  !gapprompt@>| !gapinput@ [ [ Z(3), Z(3)^0 ], [ Z(3), 0*Z(3) ] ],|
  !gapprompt@>| !gapinput@ [ [ Z(3)^0, 0*Z(3) ], [ 0*Z(3), 0*Z(3) ] ] ] );;|
  !gapprompt@gap>| !gapinput@Size( S );|
  81
  !gapprompt@gap>| !gapinput@IsomorphismTransformationSemigroup( S );;|
  !gapprompt@gap>| !gapinput@S := SymmetricInverseSemigroup( 4 );|
  <symmetric inverse monoid of degree 4>
  !gapprompt@gap>| !gapinput@IsomorphismTransformationMonoid( S );|
  MappingByFunction( <symmetric inverse monoid of degree 4>,
  <transformation monoid of degree 5 with 4 generators>
   , function( x ) ... end, <Operation "AsPartialPerm"> )
  !gapprompt@gap>| !gapinput@G := Group( (1,2,3) );|
  Group([ (1,2,3) ])
  !gapprompt@gap>| !gapinput@IsomorphismTransformationMonoid( G );|
  MappingByFunction( Group([ (1,2,3) ]), <commutative transformation
   monoid of degree 3 with 1 generator>
   , function( x ) ... end, function( x ) ... end )
\end{Verbatim}
 }

 

\subsection{\textcolor{Chapter }{AntiIsomorphismTransformationSemigroup}}
\logpage{[ 53, 7, 6 ]}\nobreak
\hyperdef{L}{X820ECE00846E480F}{}
{\noindent\textcolor{FuncColor}{$\triangleright$\enspace\texttt{AntiIsomorphismTransformationSemigroup({\mdseries\slshape S})\index{AntiIsomorphismTransformationSemigroup@\texttt{Anti}\-\texttt{Isomorphism}\-\texttt{Transformation}\-\texttt{Semigroup}}
\label{AntiIsomorphismTransformationSemigroup}
}\hfill{\scriptsize (attribute)}}\\
\textbf{\indent Returns:\ }
An anti\texttt{\symbol{45}}isomorphism.



 If \mbox{\texttt{\mdseries\slshape S}} is a semigroup, then \texttt{AntiIsomorphismTransformationSemigroup} returns an anti\texttt{\symbol{45}}isomorphism from \mbox{\texttt{\mdseries\slshape S}} to a transformation semigroup. At present, the degree of the resulting
transformation semigroup equals the size of \mbox{\texttt{\mdseries\slshape S}} plus $1$, and, consequently, this function is of limited use. 

 
\begin{Verbatim}[commandchars=!@|,fontsize=\small,frame=single,label=Example]
  !gapprompt@gap>| !gapinput@S := Semigroup( Transformation( [ 5, 5, 1, 1, 3 ] ),|
  !gapprompt@>| !gapinput@                   Transformation( [ 2, 4, 1, 5, 5 ] ) );|
  <transformation semigroup of degree 5 with 2 generators>
  !gapprompt@gap>| !gapinput@Size( S );|
  172
  !gapprompt@gap>| !gapinput@AntiIsomorphismTransformationSemigroup( S );|
  MappingByFunction( <transformation semigroup of size 172, degree 5
   with 2 generators>, <transformation semigroup of degree 173 with 2
   generators>, function( x ) ... end, function( x ) ... end )
\end{Verbatim}
 }

 }

 }

 
\chapter{\textcolor{Chapter }{Partial permutations}}\label{Partial permutations}
\logpage{[ 54, 0, 0 ]}
\hyperdef{L}{X7D6495F77B8A77BD}{}
{
  This chapter describes the functions in \textsf{GAP} for partial permutations. 

 A \emph{partial permutation} in \textsf{GAP} is simply an injective function from any finite set of positive integers to
any other finite set of positive integers. The largest point on which a
partial permutation can be defined, and the largest value that the image of
such a point can have, are defined by certain architecture dependent limits. 

 Every inverse semigroup is isomorphic to an inverse semigroup of partial
permutations and, as such, partial permutations are to inverse semigroup
theory what permutations are to group theory and transformations are to
semigroup theory. In this way, partial permutations are the elements of
inverse partial permutation semigroups. 

 A partial permutations in \textsf{GAP} acts on a finite set of positive integers on the right. The image of a point \texttt{i} under a partial permutation \texttt{f} is expressed as \texttt{i\texttt{\symbol{94}}f} in \textsf{GAP}. This action is also implemented by the function \texttt{OnPoints} (\ref{OnPoints}). The preimage of a point \texttt{i} under the partial permutation \texttt{f} can be computed using \texttt{i/f} without constructing the inverse of \texttt{f}. Partial permutations in \textsf{GAP} are created using the operations described in Section \ref{sect:CreatingPartialPerms}. Partial permutations are, by default, displayed in component notation, which
is described in Section \ref{sect:DisplayingPartialPerms}. 

 The fundamental attributes of a partial permutation are: 
\begin{description}
\item[{Domain}] The \emph{domain} of a partial permutation is just the set of positive integers where it is
defined; see \texttt{DomainOfPartialPerm} (\ref{DomainOfPartialPerm}). We will denote the domain of a partial permutation \texttt{f} by dom(\texttt{f}). 
\item[{Degree}]  The \emph{degree} of a partial permutation \texttt{f} is just the largest positive integer where \texttt{f} is defined. In other words, the degree of \texttt{f} is the largest element in the domain of \texttt{f}; see \texttt{DegreeOfPartialPerm} (\ref{DegreeOfPartialPerm}). 
\item[{Image list}]  The \emph{image list} of a partial permutation \texttt{f} is the list \texttt{[i{\textunderscore}1\texttt{\symbol{94}}f,
i{\textunderscore}2\texttt{\symbol{94}}f, .. ,
i{\textunderscore}n\texttt{\symbol{94}}f]} where the domain of \texttt{f} is \texttt{[i{\textunderscore}1, i{\textunderscore}2, .., i{\textunderscore}n]} see \texttt{ImageListOfPartialPerm} (\ref{ImageListOfPartialPerm}). For example, the partial perm sending \texttt{1} to \texttt{5} and \texttt{2} to \texttt{4} has image list \texttt{[ 5, 4 ]}. 
\item[{Image set}]  The \emph{image set} of a partial permutation \texttt{f} is just the set of points in the image list (i.e. the image list after it has
been sorted into increasing order); see \texttt{ImageSetOfPartialPerm} (\ref{ImageSetOfPartialPerm}). We will denote the image set of a partial permutation \texttt{f} by im(\texttt{f}). 
\item[{Codegree}]  The \emph{codegree} of a partial permutation \texttt{f} is just the largest positive integer of the form \texttt{i\texttt{\symbol{94}}f} for any \texttt{i} in the domain of \texttt{f}. In other words, the codegree of \texttt{f} is the largest element in the image of \texttt{f}; see \texttt{CodegreeOfPartialPerm} (\ref{CodegreeOfPartialPerm}). 
\item[{Rank}] The \emph{rank} of a partial permutation \texttt{f} is the size of its domain, or equivalently the size of its image set or image
list; see \texttt{RankOfPartialPerm} (\ref{RankOfPartialPerm}). 
\end{description}
 A \emph{functional digraph} is a directed graph where every vertex has out\texttt{\symbol{45}}degree \texttt{1}. A partial permutation \mbox{\texttt{\mdseries\slshape f}} can be thought of as a functional digraph with vertices \texttt{[1..DegreeOfPartialPerm(f)]} and edges from \texttt{i} to \texttt{i\texttt{\symbol{94}}f} for every \texttt{i}. A \emph{component} of a partial permutation is defined as a component of the corresponding
functional digraph. More specifically, \texttt{i} and \texttt{j} are in the same component if and only if there are $i=v_0, v_1, \ldots, v_n=j$ such that either $v_{k+1}=v_{k}^f$ or $v_{k}=v_{k+1}^f$ for all \texttt{k}. 

 If \texttt{S} is a semigroup and \texttt{s} is an element of \texttt{S}, then an element \texttt{t} in \texttt{S} is a \emph{semigroup inverse} for \texttt{s} if \texttt{s*t*s=s} and \texttt{t*s*t=t}; see, for example, \texttt{InverseOfTransformation} (\ref{InverseOfTransformation}). A semigroup in which every element has a unique semigroup inverse is called
an \emph{inverse semigroup}.

 Every partial permutation belongs to a symmetric inverse monoid; see \texttt{SymmetricInverseSemigroup} (\ref{SymmetricInverseSemigroup}). Inverse semigroups of partial permutations are hence inverse subsemigroups
of the symmetric inverse monoids. 

 The inverse \texttt{f\texttt{\symbol{94}}\texttt{\symbol{45}}1} of a partial permutation \texttt{f} is simply the partial permutation that maps \texttt{i\texttt{\symbol{94}}f} to \texttt{i} for all \texttt{i} in the image of \texttt{f}. It follows that the domain of \texttt{f\texttt{\symbol{94}}\texttt{\symbol{45}}1} equals the image of \texttt{f} and that the image of \texttt{f\texttt{\symbol{94}}\texttt{\symbol{45}}1} equals the domain of \texttt{f}. The inverse \texttt{f\texttt{\symbol{94}}\texttt{\symbol{45}}1} is the unique partial permutation with the property that \texttt{f*f\texttt{\symbol{94}}\texttt{\symbol{45}}1*f=f} and \texttt{f\texttt{\symbol{94}}\texttt{\symbol{45}}1*f*f\texttt{\symbol{94}}\texttt{\symbol{45}}1=f\texttt{\symbol{94}}\texttt{\symbol{45}}1}. In other words, \texttt{f\texttt{\symbol{94}}\texttt{\symbol{45}}1} is the unique semigroup inverse of \texttt{f} in the symmetric inverse monoid. 

 If \texttt{f} and \texttt{g} are partial permutations, then the domain and image of the product are:   
\[ \textrm{dom}(fg)=(\textrm{im}(f)\cap \textrm{dom}(g))f^{-1}\textrm{ and }
\textrm{im}(fg)=(\textrm{im}(f)\cap \textrm{dom}(g))g \]
  A partial permutation is an idempotent if and only if it is the identity
function on its domain. The products \texttt{f*f\texttt{\symbol{94}}\texttt{\symbol{45}}1} and \texttt{f\texttt{\symbol{94}}\texttt{\symbol{45}}1*f} are just the identity functions on the domain and image of \texttt{f}, respectively. It follows that \texttt{f*f\texttt{\symbol{94}}\texttt{\symbol{45}}1} is a left identity for \texttt{f} and \texttt{f\texttt{\symbol{94}}\texttt{\symbol{45}}1*f} is a right identity. These products will be referred to here as the \emph{left one} and \emph{right one} of the partial permutation \texttt{f}; see \texttt{LeftOne} (\ref{LeftOne:for a partial perm}). The \emph{one} of a partial permutation is just the identity on the union of its domain and
its image, and the \emph{zero} of a partial permutation is just the empty partial permutation; see \texttt{One} (\ref{One:for a partial perm}) and \texttt{MultiplicativeZero} (\ref{MultiplicativeZero:for a partial perm}). 

 If \texttt{S} is an arbitrary inverse semigroup, the \emph{natural partial order} on \texttt{S} is defined as follows: for elements \texttt{x} and \texttt{y} of \texttt{S} we say \texttt{x}$\leq$\texttt{y} if there exists an idempotent element \texttt{e} in \texttt{S} such that \texttt{x=ey}. In the context of the symmetric inverse monoid, a partial permutation \texttt{f} is less than or equal to a partial permutation \texttt{g} in the natural partial order if and only if \texttt{f} is a restriction of \texttt{g}. The natural partial order is a meet semilattice, in other words, every pair
of elements has a greatest lower bound; see \texttt{MeetOfPartialPerms} (\ref{MeetOfPartialPerms}).

 Note that unlike permutations, partial permutations do not fix unspecified
points but are simply undefined on such points; see Chapter \ref{Permutations}. Similar to permutations, and unlike transformations, it is possible to
multiply any two partial permutations in \textsf{GAP}.

 Internally, \textsf{GAP} stores a partial permutation \texttt{f} as a list consisting of the codegree of \texttt{f} and the images \texttt{i\texttt{\symbol{94}}f} of the points \texttt{i} that are less than or equal to the degree of \texttt{f}; the value \texttt{0} is stored where \texttt{i\texttt{\symbol{94}}f} is undefined. The domain and image set of \texttt{f} are also stored after either of these values is computed. When the codegree of
a partial permutation \texttt{f} is less than 65536, the codegree and images \texttt{i\texttt{\symbol{94}}f} are stored as 16\texttt{\symbol{45}}bit integers, the domain and image set are
subobjects of \texttt{f} which are immutable plain lists of \textsf{GAP} integers. When the codegree of \texttt{f} is greater than or equal to 65536, the codegree and images are stored as
32\texttt{\symbol{45}}bit integers; the domain and image set are stored in the
same way as before. A partial permutation belongs to \texttt{IsPPerm2Rep} if it is stored using 16\texttt{\symbol{45}}bit integers and to \texttt{IsPPerm4Rep} otherwise. 

 In the names of the \textsf{GAP} functions that deal with partial permutations, the word ``Permutation'' is usually abbreviated to ``Perm'', to save typing. For example, the category test function for partial
permutations is \texttt{IsPartialPerm} (\ref{IsPartialPerm}). 
\section{\textcolor{Chapter }{The family and categories of partial permutations}}\logpage{[ 54, 1, 0 ]}
\hyperdef{L}{X87B0D6657A3F2B0E}{}
{
  

\subsection{\textcolor{Chapter }{IsPartialPerm}}
\logpage{[ 54, 1, 1 ]}\nobreak
\hyperdef{L}{X7EECE133792B30FC}{}
{\noindent\textcolor{FuncColor}{$\triangleright$\enspace\texttt{IsPartialPerm({\mdseries\slshape obj})\index{IsPartialPerm@\texttt{IsPartialPerm}}
\label{IsPartialPerm}
}\hfill{\scriptsize (Category)}}\\
\textbf{\indent Returns:\ }
\texttt{true} or \texttt{false}.



 Every partial permutation in \textsf{GAP} belongs to the category \texttt{IsPartialPerm}. Basic operations for partial permutations are \texttt{DomainOfPartialPerm} (\ref{DomainOfPartialPerm}), \texttt{ImageListOfPartialPerm} (\ref{ImageListOfPartialPerm}), \texttt{ImageSetOfPartialPerm} (\ref{ImageSetOfPartialPerm}), \texttt{RankOfPartialPerm} (\ref{RankOfPartialPerm}), \texttt{DegreeOfPartialPerm} (\ref{DegreeOfPartialPerm}), multiplication of two partial permutations is via \texttt{*}, and exponentiation with the first argument a positive integer \texttt{i} and second argument a partial permutation \texttt{f} where the result is the image \texttt{i\texttt{\symbol{94}}f} of the point \texttt{i} under \texttt{f}. The inverse of a partial permutation \texttt{f} can be obtains using \texttt{f\texttt{\symbol{94}}\texttt{\symbol{45}}1}. }

 

\subsection{\textcolor{Chapter }{IsPartialPermCollection}}
\logpage{[ 54, 1, 2 ]}\nobreak
\hyperdef{L}{X8262A827790DD1CC}{}
{\noindent\textcolor{FuncColor}{$\triangleright$\enspace\texttt{IsPartialPermCollection({\mdseries\slshape obj})\index{IsPartialPermCollection@\texttt{IsPartialPermCollection}}
\label{IsPartialPermCollection}
}\hfill{\scriptsize (Category)}}\\


 Every collection of partial permutations belongs to the category \texttt{IsPartialPermCollection}. For example, a semigroup of partial permutations belongs in \texttt{IsPartialPermCollection}. }

 

\subsection{\textcolor{Chapter }{PartialPermFamily}}
\logpage{[ 54, 1, 3 ]}\nobreak
\hyperdef{L}{X7E63D17780F64FBA}{}
{\noindent\textcolor{FuncColor}{$\triangleright$\enspace\texttt{PartialPermFamily\index{PartialPermFamily@\texttt{PartialPermFamily}}
\label{PartialPermFamily}
}\hfill{\scriptsize (family)}}\\


 The family of all partial permutations is \texttt{PartialPermFamily} }

 }

   
\section{\textcolor{Chapter }{Creating partial permutations}}\label{sect:CreatingPartialPerms}
\logpage{[ 54, 2, 0 ]}
\hyperdef{L}{X7B9D451D7FDA1DD8}{}
{
  There are several ways of creating partial permutations in \textsf{GAP}, which are described in this section.  

\subsection{\textcolor{Chapter }{PartialPerm (for a domain and image)}}
\logpage{[ 54, 2, 1 ]}\nobreak
\hyperdef{L}{X8538BAE77F2FB2F8}{}
{\noindent\textcolor{FuncColor}{$\triangleright$\enspace\texttt{PartialPerm({\mdseries\slshape dom, img})\index{PartialPerm@\texttt{PartialPerm}!for a domain and image}
\label{PartialPerm:for a domain and image}
}\hfill{\scriptsize (function)}}\\
\noindent\textcolor{FuncColor}{$\triangleright$\enspace\texttt{PartialPerm({\mdseries\slshape list})\index{PartialPerm@\texttt{PartialPerm}!for a dense image}
\label{PartialPerm:for a dense image}
}\hfill{\scriptsize (function)}}\\
\textbf{\indent Returns:\ }
A partial permutation.



 Partial permutations can be created in two ways: by giving the domain and the
image, or the dense image list.

 
\begin{description}
\item[{Domain and image}]  The partial permutation defined by a domain \mbox{\texttt{\mdseries\slshape dom}} and image \mbox{\texttt{\mdseries\slshape img}}, where \mbox{\texttt{\mdseries\slshape dom}} is a set of positive integers and \mbox{\texttt{\mdseries\slshape img}} is a duplicate free list of positive integers, maps \mbox{\texttt{\mdseries\slshape dom}}\texttt{[i]} to \mbox{\texttt{\mdseries\slshape img}}\texttt{[i]}. For example, the partial permutation mapping \texttt{1} and \texttt{5} to \texttt{20} and \texttt{2} can be created using: 
\begin{Verbatim}[commandchars=!@|,fontsize=\small,frame=single,label=Example]
  PartialPerm([1,5],[20,2]); 
\end{Verbatim}
 In this setting, \texttt{PartialPerm} is the analogue in the context of partial permutations of \texttt{MappingPermListList} (\ref{MappingPermListList}). 
\item[{Dense image list}]  The partial permutation defined by a dense image list \mbox{\texttt{\mdseries\slshape list}}, maps the positive integer \texttt{i} to \mbox{\texttt{\mdseries\slshape list}}\texttt{[i]} if \mbox{\texttt{\mdseries\slshape list}}\texttt{[i]{\textless}{\textgreater}0} and is undefined at \texttt{i} if \mbox{\texttt{\mdseries\slshape list}}\texttt{[i]=0}. For example, the partial permutation mapping \texttt{1} and \texttt{5} to \texttt{20} and \texttt{2} can be created using: 
\begin{Verbatim}[commandchars=!@|,fontsize=\small,frame=single,label=Example]
  PartialPerm([20,0,0,0,2]);
\end{Verbatim}
 In this setting, \texttt{PartialPerm} is the analogue in the context of partial permutations of \texttt{PermList} (\ref{PermList}). 
\end{description}
 Regardless of which of these two methods are used to create a partial
permutation in \textsf{GAP} the internal representation is the same. 

 If the largest point in the domain is larger than the rank of the partial
permutation, then using the dense image list to define the partial permutation
will require less typing; otherwise using the domain and the image will
require less typing. For example, the partial permutation mapping \texttt{10000} to \texttt{1} can be defined using: 
\begin{Verbatim}[commandchars=!@|,fontsize=\small,frame=single,label=Example]
  PartialPerm([10000], [1]);
\end{Verbatim}
 but using the dense image list would require a list with \texttt{9999} entries equal to \texttt{0} and the final entry equal to \texttt{1}. On the other hand, the identity on \texttt{[1,2,3,4,6]} can be defined using: 
\begin{Verbatim}[commandchars=!@|,fontsize=\small,frame=single,label=Example]
  PartialPerm([1,2,3,4,0,6]);
\end{Verbatim}
 

 Please note that a partial permutation in \textsf{GAP} is never a permutation nor is a permutation ever a partial permutation. For
example, the permutation \texttt{(1,4,2)} fixes \texttt{3} but the partial permutation \texttt{PartialPerm([4,1,0,2]);} is not defined on \texttt{3}. }

  

\subsection{\textcolor{Chapter }{PartialPermOp}}
\logpage{[ 54, 2, 2 ]}\nobreak
\hyperdef{L}{X81188D9F83F64222}{}
{\noindent\textcolor{FuncColor}{$\triangleright$\enspace\texttt{PartialPermOp({\mdseries\slshape obj, list[, func]})\index{PartialPermOp@\texttt{PartialPermOp}}
\label{PartialPermOp}
}\hfill{\scriptsize (operation)}}\\
\noindent\textcolor{FuncColor}{$\triangleright$\enspace\texttt{PartialPermOpNC({\mdseries\slshape obj, list[, func]})\index{PartialPermOpNC@\texttt{PartialPermOpNC}}
\label{PartialPermOpNC}
}\hfill{\scriptsize (operation)}}\\
\textbf{\indent Returns:\ }
A partial permutation or \texttt{fail}.



 \texttt{PartialPermOp} returns the partial permutation that corresponds to the action of the object \mbox{\texttt{\mdseries\slshape obj}} on the domain or list \mbox{\texttt{\mdseries\slshape list}} via the function \mbox{\texttt{\mdseries\slshape func}}. If the optional third argument \mbox{\texttt{\mdseries\slshape func}} is not specified, then the action \texttt{OnPoints} (\ref{OnPoints}) is used by default. Note that the returned partial permutation refers to the
positions in \mbox{\texttt{\mdseries\slshape list}} even if \mbox{\texttt{\mdseries\slshape list}} itself consists of integers. 

 This function is the analogue in the context of partial permutations of \texttt{Permutation} (\ref{Permutation:for a group, an action domain, etc.}) or \texttt{TransformationOp} (\ref{TransformationOp}).

 If \mbox{\texttt{\mdseries\slshape obj}} does not map the elements of \mbox{\texttt{\mdseries\slshape list}} injectively, then \texttt{fail} is returned.

 \texttt{PartialPermOpNC} does not check that \mbox{\texttt{\mdseries\slshape obj}} maps elements of \mbox{\texttt{\mdseries\slshape list}} injectively or that a partial permutation is defined by the action of \mbox{\texttt{\mdseries\slshape obj}} on \mbox{\texttt{\mdseries\slshape list}} via \mbox{\texttt{\mdseries\slshape func}}. This function should be used only with caution, in situations where it is
guaranteed that the arguments have the required properties. 
\begin{Verbatim}[commandchars=!@|,fontsize=\small,frame=single,label=Example]
  !gapprompt@gap>| !gapinput@f:=Transformation( [ 9, 10, 4, 2, 10, 5, 9, 10, 9, 6 ] );;|
  !gapprompt@gap>| !gapinput@PartialPermOp(f, [ 6 .. 8 ], OnPoints);|
  [1,4][2,5][3,6]
\end{Verbatim}
 }

  

\subsection{\textcolor{Chapter }{RestrictedPartialPerm}}
\logpage{[ 54, 2, 3 ]}\nobreak
\hyperdef{L}{X80ABBF4883C79060}{}
{\noindent\textcolor{FuncColor}{$\triangleright$\enspace\texttt{RestrictedPartialPerm({\mdseries\slshape f, set})\index{RestrictedPartialPerm@\texttt{RestrictedPartialPerm}}
\label{RestrictedPartialPerm}
}\hfill{\scriptsize (operation)}}\\
\textbf{\indent Returns:\ }
A partial permutation.



 \texttt{RestrictedPartialPerm} returns a new partial permutation that acts on the points in the set of
positive integers \mbox{\texttt{\mdseries\slshape set}} in the same way as the partial permutation \mbox{\texttt{\mdseries\slshape f}}, and that is undefined on those points that are not in \mbox{\texttt{\mdseries\slshape set}}. 
\begin{Verbatim}[commandchars=!@|,fontsize=\small,frame=single,label=Example]
  !gapprompt@gap>| !gapinput@f:=PartialPerm( [ 1, 3, 4, 7, 8, 9 ], [ 9, 4, 1, 6, 2, 8 ] );;|
  !gapprompt@gap>| !gapinput@RestrictedPartialPerm(f, [ 2, 3, 6, 10 ] );|
  [3,4]
\end{Verbatim}
 }

  

\subsection{\textcolor{Chapter }{JoinOfPartialPerms}}
\logpage{[ 54, 2, 4 ]}\nobreak
\hyperdef{L}{X849668DD7B0B9E3B}{}
{\noindent\textcolor{FuncColor}{$\triangleright$\enspace\texttt{JoinOfPartialPerms({\mdseries\slshape arg})\index{JoinOfPartialPerms@\texttt{JoinOfPartialPerms}}
\label{JoinOfPartialPerms}
}\hfill{\scriptsize (function)}}\\
\noindent\textcolor{FuncColor}{$\triangleright$\enspace\texttt{JoinOfIdempotentPartialPermsNC({\mdseries\slshape arg})\index{JoinOfIdempotentPartialPermsNC@\texttt{JoinOfIdempotentPartialPermsNC}}
\label{JoinOfIdempotentPartialPermsNC}
}\hfill{\scriptsize (function)}}\\
\textbf{\indent Returns:\ }
A partial permutation or \texttt{fail}.



 The join of partial permutations \mbox{\texttt{\mdseries\slshape f}} and \mbox{\texttt{\mdseries\slshape g}} is just the join, or supremum, of \mbox{\texttt{\mdseries\slshape f}} and \mbox{\texttt{\mdseries\slshape g}} under the natural partial ordering of partial permutations. 

 \texttt{JoinOfPartialPerms} returns the union of the partial permutations in its argument if this defines
a partial permutation, and \texttt{fail} if it is not. The argument \mbox{\texttt{\mdseries\slshape arg}} can be a partial permutation collection or a number of partial permutations. 

 The function \texttt{JoinOfIdempotentPartialPermsNC} returns the join of its argument which is assumed to be a collection of
idempotent partial permutations or a number of idempotent partial
permutations. It is not checked that the arguments are idempotents. The
performance of this function is higher than \texttt{JoinOfPartialPerms} when it is known \emph{a priori} that the argument consists of idempotents.

 The union of \mbox{\texttt{\mdseries\slshape f}} and \mbox{\texttt{\mdseries\slshape g}} is a partial permutation if and only if \mbox{\texttt{\mdseries\slshape f}} and \mbox{\texttt{\mdseries\slshape g}} agree on the intersection dom(\mbox{\texttt{\mdseries\slshape f}})$\cap$ dom(\mbox{\texttt{\mdseries\slshape g}}) of their domains and the images of dom(\mbox{\texttt{\mdseries\slshape f}})$\setminus$ dom(\mbox{\texttt{\mdseries\slshape g}}) and dom(\mbox{\texttt{\mdseries\slshape g}})$\setminus$ dom(\mbox{\texttt{\mdseries\slshape f}}) under \mbox{\texttt{\mdseries\slshape f}} and \mbox{\texttt{\mdseries\slshape g}}, respectively, are disjoint. 
\begin{Verbatim}[commandchars=!@|,fontsize=\small,frame=single,label=Example]
  !gapprompt@gap>| !gapinput@f:=PartialPerm( [ 1, 2, 3, 6, 8, 10 ], [ 2, 6, 7, 9, 1, 5 ] );|
  [3,7][8,1,2,6,9][10,5]
  !gapprompt@gap>| !gapinput@g:=PartialPerm( [ 11, 12, 14, 16, 18, 19 ],|
  !gapprompt@>| !gapinput@[ 17, 20, 11, 19, 14, 12 ] );|
  [16,19,12,20][18,14,11,17]
  !gapprompt@gap>| !gapinput@JoinOfPartialPerms(f, g);|
  [3,7][8,1,2,6,9][10,5][16,19,12,20][18,14,11,17]
  !gapprompt@gap>| !gapinput@f:=PartialPerm( [ 1, 4, 5, 6, 7 ], [ 5, 7, 3, 1, 4 ] );|
  [6,1,5,3](4,7)
  !gapprompt@gap>| !gapinput@g:=PartialPerm( [ 100 ], [ 1 ] );|
  [100,1]
  !gapprompt@gap>| !gapinput@JoinOfPartialPerms(f, g);|
  fail
  !gapprompt@gap>| !gapinput@f:=PartialPerm( [ 1, 3, 4 ], [ 3, 2, 4 ] );|
  [1,3,2](4)
  !gapprompt@gap>| !gapinput@g:=PartialPerm( [ 1, 2, 4 ], [ 2, 3, 4 ] );|
  [1,2,3](4)
  !gapprompt@gap>| !gapinput@JoinOfPartialPerms(f, g);|
  fail
  !gapprompt@gap>| !gapinput@f:=PartialPerm( [ 1 ], [ 2 ] );|
  [1,2]
  !gapprompt@gap>| !gapinput@JoinOfPartialPerms(f, f^-1);|
  (1,2)
\end{Verbatim}
 }

  

\subsection{\textcolor{Chapter }{MeetOfPartialPerms}}
\logpage{[ 54, 2, 5 ]}\nobreak
\hyperdef{L}{X81E2B6977E28CD00}{}
{\noindent\textcolor{FuncColor}{$\triangleright$\enspace\texttt{MeetOfPartialPerms({\mdseries\slshape arg})\index{MeetOfPartialPerms@\texttt{MeetOfPartialPerms}}
\label{MeetOfPartialPerms}
}\hfill{\scriptsize (function)}}\\
\textbf{\indent Returns:\ }
A partial permutation.



 The meet of partial permutations \mbox{\texttt{\mdseries\slshape f}} and \mbox{\texttt{\mdseries\slshape g}} is just the meet, or infimum, of \mbox{\texttt{\mdseries\slshape f}} and \mbox{\texttt{\mdseries\slshape g}} under the natural partial ordering of partial permutations. In other words,
the meet is the greatest partial permutation which is a restriction of both \mbox{\texttt{\mdseries\slshape f}} and \mbox{\texttt{\mdseries\slshape g}}. 

 Note that unlike the join of partial permutations, the meet always exists. 

 \texttt{MeetOfPartialPerms} returns the meet of the partial permutations in its argument. The argument \mbox{\texttt{\mdseries\slshape arg}} can be a partial permutation collection or a number of partial permutations. 
\begin{Verbatim}[commandchars=!@|,fontsize=\small,frame=single,label=Example]
  !gapprompt@gap>| !gapinput@f:=PartialPerm( [ 1, 2, 3, 6, 100000 ], [ 2, 6, 7, 1, 5 ] );|
  [3,7][100000,5](1,2,6)
  !gapprompt@gap>| !gapinput@g:=PartialPerm( [ 1, 2, 3, 4, 6 ], [ 2, 4, 6, 1, 5 ] );|
  [3,6,5](1,2,4)
  !gapprompt@gap>| !gapinput@MeetOfPartialPerms(f, g);|
  [1,2]
  !gapprompt@gap>| !gapinput@g:=PartialPerm( [ 1, 2, 3, 5, 6, 7, 9, 10 ],|
  !gapprompt@>| !gapinput@[ 4, 10, 5, 6, 7, 1, 3, 2 ] );|
  [9,3,5,6,7,1,4](2,10)
  !gapprompt@gap>| !gapinput@MeetOfPartialPerms(f, g);|
  <empty partial perm>
\end{Verbatim}
 }

  

\subsection{\textcolor{Chapter }{EmptyPartialPerm}}
\logpage{[ 54, 2, 6 ]}\nobreak
\hyperdef{L}{X80EFB142817A0A9F}{}
{\noindent\textcolor{FuncColor}{$\triangleright$\enspace\texttt{EmptyPartialPerm({\mdseries\slshape })\index{EmptyPartialPerm@\texttt{EmptyPartialPerm}}
\label{EmptyPartialPerm}
}\hfill{\scriptsize (function)}}\\
\textbf{\indent Returns:\ }
The empty partial permutation.



 The empty partial permutation is returned by this function when it is called
with no arguments. This is just short hand for \texttt{PartialPerm([]);}. 
\begin{Verbatim}[commandchars=!@|,fontsize=\small,frame=single,label=Example]
  !gapprompt@gap>| !gapinput@EmptyPartialPerm();|
  <empty partial perm>
\end{Verbatim}
 }

  
\subsection{\textcolor{Chapter }{RandomPartialPerm}}\logpage{[ 54, 2, 7 ]}
\hyperdef{L}{X7E6ADC8583C31530}{}
{
\noindent\textcolor{FuncColor}{$\triangleright$\enspace\texttt{RandomPartialPerm({\mdseries\slshape n})\index{RandomPartialPerm@\texttt{RandomPartialPerm}!for a positive integer}
\label{RandomPartialPerm:for a positive integer}
}\hfill{\scriptsize (function)}}\\
\noindent\textcolor{FuncColor}{$\triangleright$\enspace\texttt{RandomPartialPerm({\mdseries\slshape set})\index{RandomPartialPerm@\texttt{RandomPartialPerm}!for a set of positive
      integers}
\label{RandomPartialPerm:for a set of positive
      integers}
}\hfill{\scriptsize (function)}}\\
\noindent\textcolor{FuncColor}{$\triangleright$\enspace\texttt{RandomPartialPerm({\mdseries\slshape dom, img})\index{RandomPartialPerm@\texttt{RandomPartialPerm}!for domain and image}
\label{RandomPartialPerm:for domain and image}
}\hfill{\scriptsize (function)}}\\
\textbf{\indent Returns:\ }
A random partial permutation.



 In its first form, \texttt{RandomPartialPerm} returns a randomly chosen partial permutation where points in the domain and
image are bounded above by the positive integer \mbox{\texttt{\mdseries\slshape n}}. 
\begin{Verbatim}[commandchars=!@|,fontsize=\small,frame=single,label=Example]
  !gapprompt@gap>| !gapinput@RandomPartialPerm(10);|
  [2,9][4,1,6,5][7,3](8)
\end{Verbatim}
 In its second form, \texttt{RandomPartialPerm} returns a randomly chosen partial permutation with points in the domain and
image contained in the set of positive integers \mbox{\texttt{\mdseries\slshape set}}. 
\begin{Verbatim}[commandchars=!@|,fontsize=\small,frame=single,label=Example]
  !gapprompt@gap>| !gapinput@RandomPartialPerm([1,2,3,1000]);|
  [2,3,1000](1)
\end{Verbatim}
 In its third form, \texttt{RandomPartialPerm} creates a randomly chosen partial permutation with domain contained in the set
of positive integers \mbox{\texttt{\mdseries\slshape dom}} and image contained in the set of positive integers \mbox{\texttt{\mdseries\slshape img}}. The arguments \mbox{\texttt{\mdseries\slshape dom}} and \mbox{\texttt{\mdseries\slshape img}} do not have to have equal length.

 Note that it is not guaranteed in either of these cases that partial
permutations are chosen with a uniform distribution. }

 }

  
\section{\textcolor{Chapter }{Attributes for partial permutations}}\label{sect:AttributesPartialPerms}
\logpage{[ 54, 3, 0 ]}
\hyperdef{L}{X8779F0997D0FDA78}{}
{
  In this section we describe the functions available in \textsf{GAP} for finding various attributes of partial permutations. 

 

\subsection{\textcolor{Chapter }{DegreeOfPartialPerm}}
\logpage{[ 54, 3, 1 ]}\nobreak
\hyperdef{L}{X8612A4DC864E7959}{}
{\noindent\textcolor{FuncColor}{$\triangleright$\enspace\texttt{DegreeOfPartialPerm({\mdseries\slshape f})\index{DegreeOfPartialPerm@\texttt{DegreeOfPartialPerm}}
\label{DegreeOfPartialPerm}
}\hfill{\scriptsize (function)}}\\
\noindent\textcolor{FuncColor}{$\triangleright$\enspace\texttt{DegreeOfPartialPermCollection({\mdseries\slshape coll})\index{DegreeOfPartialPermCollection@\texttt{DegreeOfPartialPermCollection}}
\label{DegreeOfPartialPermCollection}
}\hfill{\scriptsize (attribute)}}\\
\textbf{\indent Returns:\ }
A non\texttt{\symbol{45}}negative integer.



 The \emph{degree} of a partial permutation \mbox{\texttt{\mdseries\slshape f}} is the largest positive integer where it is defined, i.e. the maximum element
in the domain of \mbox{\texttt{\mdseries\slshape f}}. 

 The degree a collection of partial permutations \mbox{\texttt{\mdseries\slshape coll}} is the largest degree of any partial permutation in \mbox{\texttt{\mdseries\slshape coll}}. 
\begin{Verbatim}[commandchars=!@|,fontsize=\small,frame=single,label=Example]
  !gapprompt@gap>| !gapinput@f:=PartialPerm( [ 1, 2, 3, 6, 8, 10 ], [ 2, 6, 7, 9, 1, 5 ] );|
  [3,7][8,1,2,6,9][10,5]
  !gapprompt@gap>| !gapinput@DegreeOfPartialPerm(f);|
  10
\end{Verbatim}
 }

  

\subsection{\textcolor{Chapter }{CodegreeOfPartialPerm}}
\logpage{[ 54, 3, 2 ]}\nobreak
\hyperdef{L}{X8413D0EF7DEE1FFF}{}
{\noindent\textcolor{FuncColor}{$\triangleright$\enspace\texttt{CodegreeOfPartialPerm({\mdseries\slshape f})\index{CodegreeOfPartialPerm@\texttt{CodegreeOfPartialPerm}}
\label{CodegreeOfPartialPerm}
}\hfill{\scriptsize (function)}}\\
\noindent\textcolor{FuncColor}{$\triangleright$\enspace\texttt{CodegreeOfPartialPermCollection({\mdseries\slshape coll})\index{CodegreeOfPartialPermCollection@\texttt{CodegreeOfPartialPermCollection}}
\label{CodegreeOfPartialPermCollection}
}\hfill{\scriptsize (attribute)}}\\
\textbf{\indent Returns:\ }
A non\texttt{\symbol{45}}negative integer.



 The \emph{codegree} of a partial permutation \mbox{\texttt{\mdseries\slshape f}} is the largest positive integer in its image. 

 The codegree a collection of partial permutations \mbox{\texttt{\mdseries\slshape coll}} is the largest codegree of any partial permutation in \mbox{\texttt{\mdseries\slshape coll}}. 
\begin{Verbatim}[commandchars=!@|,fontsize=\small,frame=single,label=Example]
  !gapprompt@gap>| !gapinput@f:=PartialPerm( [ 1, 2, 3, 4, 5, 8, 10 ], [ 7, 1, 4, 3, 2, 6, 5 ] );|
  [8,6][10,5,2,1,7](3,4)
  !gapprompt@gap>| !gapinput@CodegreeOfPartialPerm(f);|
  7
\end{Verbatim}
 }

  

\subsection{\textcolor{Chapter }{RankOfPartialPerm}}
\logpage{[ 54, 3, 3 ]}\nobreak
\hyperdef{L}{X7C1ABD8A80E95B39}{}
{\noindent\textcolor{FuncColor}{$\triangleright$\enspace\texttt{RankOfPartialPerm({\mdseries\slshape f})\index{RankOfPartialPerm@\texttt{RankOfPartialPerm}}
\label{RankOfPartialPerm}
}\hfill{\scriptsize (function)}}\\
\noindent\textcolor{FuncColor}{$\triangleright$\enspace\texttt{RankOfPartialPermCollection({\mdseries\slshape coll})\index{RankOfPartialPermCollection@\texttt{RankOfPartialPermCollection}}
\label{RankOfPartialPermCollection}
}\hfill{\scriptsize (attribute)}}\\
\textbf{\indent Returns:\ }
A non\texttt{\symbol{45}}negative integer.



 The \emph{rank} of a partial permutation \mbox{\texttt{\mdseries\slshape f}} is the size of its domain, or equivalently the size of its image set or image
list.

 The rank of a partial permutation collection \mbox{\texttt{\mdseries\slshape coll}} is the size of the union of the domains of the elements of \mbox{\texttt{\mdseries\slshape coll}}, or equivalently, the total number of points on which the elements of \mbox{\texttt{\mdseries\slshape coll}} act. Note that this is value may not the same as the size of the union of the
images of the elements in \mbox{\texttt{\mdseries\slshape coll}}. 
\begin{Verbatim}[commandchars=!@|,fontsize=\small,frame=single,label=Example]
  !gapprompt@gap>| !gapinput@f:=PartialPerm( [ 1, 2, 4, 6, 8, 9 ], [ 7, 10, 1, 9, 4, 2 ] );|
  [6,9,2,10][8,4,1,7]
  !gapprompt@gap>| !gapinput@RankOfPartialPerm(f);|
  6
\end{Verbatim}
 }

  

\subsection{\textcolor{Chapter }{DomainOfPartialPerm}}
\logpage{[ 54, 3, 4 ]}\nobreak
\hyperdef{L}{X784A14F787E041D7}{}
{\noindent\textcolor{FuncColor}{$\triangleright$\enspace\texttt{DomainOfPartialPerm({\mdseries\slshape f})\index{DomainOfPartialPerm@\texttt{DomainOfPartialPerm}}
\label{DomainOfPartialPerm}
}\hfill{\scriptsize (attribute)}}\\
\noindent\textcolor{FuncColor}{$\triangleright$\enspace\texttt{DomainOfPartialPermCollection({\mdseries\slshape f})\index{DomainOfPartialPermCollection@\texttt{DomainOfPartialPermCollection}}
\label{DomainOfPartialPermCollection}
}\hfill{\scriptsize (attribute)}}\\
\textbf{\indent Returns:\ }
A set of positive integers (maybe empty).



 The \emph{domain} of a partial permutation \mbox{\texttt{\mdseries\slshape f}} is the set of positive integers where \mbox{\texttt{\mdseries\slshape f}} is defined. 

 The domain of a partial permutation collection \mbox{\texttt{\mdseries\slshape coll}} is the union of the domains of its elements. 
\begin{Verbatim}[commandchars=!@|,fontsize=\small,frame=single,label=Example]
  !gapprompt@gap>| !gapinput@f:=PartialPerm( [ 1, 2, 3, 6, 8, 10 ], [ 2, 6, 7, 9, 1, 5 ] );|
  [3,7][8,1,2,6,9][10,5]
  !gapprompt@gap>| !gapinput@DomainOfPartialPerm(f);|
  [ 1, 2, 3, 6, 8, 10 ]
\end{Verbatim}
 }

  

\subsection{\textcolor{Chapter }{ImageOfPartialPermCollection}}
\logpage{[ 54, 3, 5 ]}\nobreak
\hyperdef{L}{X7CD84B107831E0FC}{}
{\noindent\textcolor{FuncColor}{$\triangleright$\enspace\texttt{ImageOfPartialPermCollection({\mdseries\slshape coll})\index{ImageOfPartialPermCollection@\texttt{ImageOfPartialPermCollection}}
\label{ImageOfPartialPermCollection}
}\hfill{\scriptsize (attribute)}}\\
\textbf{\indent Returns:\ }
A set of positive integers (maybe empty).



 The \emph{image} of a partial permutation collection \mbox{\texttt{\mdseries\slshape coll}} is the union of the images of its elements. 
\begin{Verbatim}[commandchars=!@|,fontsize=\small,frame=single,label=Example]
  !gapprompt@gap>| !gapinput@S := SymmetricInverseSemigroup(5);|
  <symmetric inverse monoid of degree 5>
  !gapprompt@gap>| !gapinput@ImageOfPartialPermCollection(GeneratorsOfInverseSemigroup(S));|
  [ 1 .. 5 ]
\end{Verbatim}
 }

  

\subsection{\textcolor{Chapter }{ImageListOfPartialPerm}}
\logpage{[ 54, 3, 6 ]}\nobreak
\hyperdef{L}{X8333293F87F654FA}{}
{\noindent\textcolor{FuncColor}{$\triangleright$\enspace\texttt{ImageListOfPartialPerm({\mdseries\slshape f})\index{ImageListOfPartialPerm@\texttt{ImageListOfPartialPerm}}
\label{ImageListOfPartialPerm}
}\hfill{\scriptsize (attribute)}}\\
\textbf{\indent Returns:\ }
The list of images of a partial permutation.



 The \emph{image list} of a partial permutation \mbox{\texttt{\mdseries\slshape f}} is the list of images of the elements of the domain \mbox{\texttt{\mdseries\slshape f}} where \texttt{ImageListOfPartialPerm(\mbox{\texttt{\mdseries\slshape f}})[i]=DomainOfPartialPerm(\mbox{\texttt{\mdseries\slshape f}})[i]\texttt{\symbol{94}}\mbox{\texttt{\mdseries\slshape f}}} for any \texttt{i} in the range from \texttt{1} to the rank of \mbox{\texttt{\mdseries\slshape f}}. 
\begin{Verbatim}[commandchars=!@|,fontsize=\small,frame=single,label=Example]
  !gapprompt@gap>| !gapinput@f:=PartialPerm( [ 1, 2, 3, 4, 5, 8, 10 ], [ 7, 1, 4, 3, 2, 6, 5 ] );|
  [8,6][10,5,2,1,7](3,4)
  !gapprompt@gap>| !gapinput@ImageListOfPartialPerm(f);|
  [ 7, 1, 4, 3, 2, 6, 5 ]
\end{Verbatim}
 }

  

\subsection{\textcolor{Chapter }{ImageSetOfPartialPerm}}
\logpage{[ 54, 3, 7 ]}\nobreak
\hyperdef{L}{X7F0724A07A14DCF7}{}
{\noindent\textcolor{FuncColor}{$\triangleright$\enspace\texttt{ImageSetOfPartialPerm({\mdseries\slshape f})\index{ImageSetOfPartialPerm@\texttt{ImageSetOfPartialPerm}}
\label{ImageSetOfPartialPerm}
}\hfill{\scriptsize (attribute)}}\\
\textbf{\indent Returns:\ }
The image set of a partial permutation.



 The \emph{image set} of a partial permutation \texttt{f} is just the set of points in the image list (i.e. the image list after it has
been sorted into increasing order). 
\begin{Verbatim}[commandchars=!@|,fontsize=\small,frame=single,label=Example]
  !gapprompt@gap>| !gapinput@f:=PartialPerm( [ 1, 2, 3, 5, 7, 10 ], [ 10, 2, 3, 5, 7, 6 ] );|
  [1,10,6](2)(3)(5)(7)
  !gapprompt@gap>| !gapinput@ImageSetOfPartialPerm(f);|
  [ 2, 3, 5, 6, 7, 10 ]
\end{Verbatim}
 }

  

\subsection{\textcolor{Chapter }{FixedPointsOfPartialPerm (for a partial perm)}}
\logpage{[ 54, 3, 8 ]}\nobreak
\hyperdef{L}{X82AAFF938623422E}{}
{\noindent\textcolor{FuncColor}{$\triangleright$\enspace\texttt{FixedPointsOfPartialPerm({\mdseries\slshape f})\index{FixedPointsOfPartialPerm@\texttt{FixedPointsOfPartialPerm}!for a partial perm}
\label{FixedPointsOfPartialPerm:for a partial perm}
}\hfill{\scriptsize (attribute)}}\\
\noindent\textcolor{FuncColor}{$\triangleright$\enspace\texttt{FixedPointsOfPartialPerm({\mdseries\slshape coll})\index{FixedPointsOfPartialPerm@\texttt{FixedPointsOfPartialPerm}!for a partial perm coll}
\label{FixedPointsOfPartialPerm:for a partial perm coll}
}\hfill{\scriptsize (method)}}\\
\textbf{\indent Returns:\ }
A set of positive integers.



 \texttt{FixedPointsOfPartialPerm} returns the set of points \texttt{i} in the domain of the partial permutation \mbox{\texttt{\mdseries\slshape f}} such that \texttt{i\texttt{\symbol{94}}\mbox{\texttt{\mdseries\slshape f}}=i}. 

 When the argument is a collection of partial permutations \mbox{\texttt{\mdseries\slshape coll}}, \texttt{FixedPointsOfPartialPerm} returns the set of points fixed by every element of the collection of partial
permutations \mbox{\texttt{\mdseries\slshape coll}}. 
\begin{Verbatim}[commandchars=!@|,fontsize=\small,frame=single,label=Example]
  !gapprompt@gap>| !gapinput@f := PartialPerm( [ 1, 2, 3, 6, 7 ], [ 1, 3, 4, 7, 5 ] );|
  [2,3,4][6,7,5](1)
  !gapprompt@gap>| !gapinput@FixedPointsOfPartialPerm(f);|
  [ 1 ]
  !gapprompt@gap>| !gapinput@f := PartialPerm([1 .. 10]);;|
  !gapprompt@gap>| !gapinput@FixedPointsOfPartialPerm(f);|
  [ 1, 2, 3, 4, 5, 6, 7, 8, 9, 10 ]
\end{Verbatim}
 }

 

\subsection{\textcolor{Chapter }{MovedPoints (for a partial perm)}}
\logpage{[ 54, 3, 9 ]}\nobreak
\hyperdef{L}{X82FE981A87FAA2DC}{}
{\noindent\textcolor{FuncColor}{$\triangleright$\enspace\texttt{MovedPoints({\mdseries\slshape f})\index{MovedPoints@\texttt{MovedPoints}!for a partial perm}
\label{MovedPoints:for a partial perm}
}\hfill{\scriptsize (attribute)}}\\
\noindent\textcolor{FuncColor}{$\triangleright$\enspace\texttt{MovedPoints({\mdseries\slshape coll})\index{MovedPoints@\texttt{MovedPoints}!for a partial perm coll}
\label{MovedPoints:for a partial perm coll}
}\hfill{\scriptsize (method)}}\\
\textbf{\indent Returns:\ }
A set of positive integers.



 \texttt{MovedPoints} returns the set of points \texttt{i} in the domain of the partial permutation \mbox{\texttt{\mdseries\slshape f}} such that \texttt{i\texttt{\symbol{94}}\mbox{\texttt{\mdseries\slshape f}}{\textless}{\textgreater}i}. 

 When the argument is a collection of partial permutations \mbox{\texttt{\mdseries\slshape coll}}, \texttt{MovedPoints} returns the set of points moved by some element of the collection of partial
permutations \mbox{\texttt{\mdseries\slshape coll}}. 
\begin{Verbatim}[commandchars=!@|,fontsize=\small,frame=single,label=Example]
  !gapprompt@gap>| !gapinput@f := PartialPerm( [ 1, 2, 3, 4 ], [ 5, 7, 1, 6 ] );|
  [2,7][3,1,5][4,6]
  !gapprompt@gap>| !gapinput@MovedPoints(f);|
  [ 1, 2, 3, 4 ]
  !gapprompt@gap>| !gapinput@FixedPointsOfPartialPerm(f);|
  [  ]
  !gapprompt@gap>| !gapinput@FixedPointsOfPartialPerm(PartialPerm([1 .. 4]));|
  [ 1, 2, 3, 4 ]
\end{Verbatim}
 }

  

\subsection{\textcolor{Chapter }{NrFixedPoints (for a partial perm)}}
\logpage{[ 54, 3, 10 ]}\nobreak
\hyperdef{L}{X7FAF969C84CDC742}{}
{\noindent\textcolor{FuncColor}{$\triangleright$\enspace\texttt{NrFixedPoints({\mdseries\slshape f})\index{NrFixedPoints@\texttt{NrFixedPoints}!for a partial perm}
\label{NrFixedPoints:for a partial perm}
}\hfill{\scriptsize (attribute)}}\\
\noindent\textcolor{FuncColor}{$\triangleright$\enspace\texttt{NrFixedPoints({\mdseries\slshape coll})\index{NrFixedPoints@\texttt{NrFixedPoints}!for a partial perm coll}
\label{NrFixedPoints:for a partial perm coll}
}\hfill{\scriptsize (method)}}\\
\textbf{\indent Returns:\ }
A positive integer.



 \texttt{NrFixedPoints} returns the number of points \texttt{i} in the domain of the partial permutation \mbox{\texttt{\mdseries\slshape f}} such that \texttt{i\texttt{\symbol{94}}\mbox{\texttt{\mdseries\slshape f}}=i}. 

 When the argument is a collection of partial permutations \mbox{\texttt{\mdseries\slshape coll}}, \texttt{NrFixedPoints} returns the number of points fixed by every element of the collection of
partial permutations \mbox{\texttt{\mdseries\slshape coll}}. 
\begin{Verbatim}[commandchars=!@|,fontsize=\small,frame=single,label=Example]
  !gapprompt@gap>| !gapinput@f := PartialPerm( [ 1, 2, 3, 4, 5 ], [ 3, 2, 4, 6, 1 ] );|
  [5,1,3,4,6](2)
  !gapprompt@gap>| !gapinput@NrFixedPoints(f);|
  1
  !gapprompt@gap>| !gapinput@NrFixedPoints(PartialPerm([1 .. 10]));|
  10
\end{Verbatim}
 }

 

\subsection{\textcolor{Chapter }{NrMovedPoints (for a partial perm)}}
\logpage{[ 54, 3, 11 ]}\nobreak
\hyperdef{L}{X81F5C64E7DAD27A7}{}
{\noindent\textcolor{FuncColor}{$\triangleright$\enspace\texttt{NrMovedPoints({\mdseries\slshape f})\index{NrMovedPoints@\texttt{NrMovedPoints}!for a partial perm}
\label{NrMovedPoints:for a partial perm}
}\hfill{\scriptsize (attribute)}}\\
\noindent\textcolor{FuncColor}{$\triangleright$\enspace\texttt{NrMovedPoints({\mdseries\slshape coll})\index{NrMovedPoints@\texttt{NrMovedPoints}!for a partial perm coll}
\label{NrMovedPoints:for a partial perm coll}
}\hfill{\scriptsize (method)}}\\
\textbf{\indent Returns:\ }
A positive integer.



 \texttt{NrMovedPoints} returns the number of points \texttt{i} in the domain of the partial permutation \mbox{\texttt{\mdseries\slshape f}} such that \texttt{i\texttt{\symbol{94}}\mbox{\texttt{\mdseries\slshape f}}{\textless}{\textgreater}i}. 

 When the argument is a collection of partial permutations \mbox{\texttt{\mdseries\slshape coll}}, \texttt{NrMovedPoints} returns the number of points moved by some element of the collection of
partial permutations \mbox{\texttt{\mdseries\slshape coll}}. 
\begin{Verbatim}[commandchars=!@|,fontsize=\small,frame=single,label=Example]
  !gapprompt@gap>| !gapinput@f := PartialPerm( [ 1, 2, 3, 4, 5, 7, 8 ], [ 4, 5, 6, 7, 1, 3, 2 ] );|
  [8,2,5,1,4,7,3,6]
  !gapprompt@gap>| !gapinput@NrMovedPoints(f);|
  7
  !gapprompt@gap>| !gapinput@NrMovedPoints(PartialPerm([1 .. 4]));|
  0
\end{Verbatim}
 }

  

\subsection{\textcolor{Chapter }{SmallestMovedPoint (for a partial perm)}}
\logpage{[ 54, 3, 12 ]}\nobreak
\hyperdef{L}{X84A49C977E1E29AA}{}
{\noindent\textcolor{FuncColor}{$\triangleright$\enspace\texttt{SmallestMovedPoint({\mdseries\slshape f})\index{SmallestMovedPoint@\texttt{SmallestMovedPoint}!for a partial perm}
\label{SmallestMovedPoint:for a partial perm}
}\hfill{\scriptsize (attribute)}}\\
\noindent\textcolor{FuncColor}{$\triangleright$\enspace\texttt{SmallestMovedPoint({\mdseries\slshape coll})\index{SmallestMovedPoint@\texttt{SmallestMovedPoint}!for a partial perm coll}
\label{SmallestMovedPoint:for a partial perm coll}
}\hfill{\scriptsize (method)}}\\
\textbf{\indent Returns:\ }
A positive integer or \texttt{infinity}.



 \texttt{SmallestMovedPoint} returns the smallest positive integer \texttt{i} such that \texttt{i\texttt{\symbol{94}}\mbox{\texttt{\mdseries\slshape f}}{\textless}{\textgreater}i} if such an \texttt{i} exists. If \mbox{\texttt{\mdseries\slshape f}} is an identity partial permutation, then \texttt{infinity} is returned.

 If the argument is a collection of partial permutations \mbox{\texttt{\mdseries\slshape coll}}, then the smallest point which is moved by at least one element of \mbox{\texttt{\mdseries\slshape coll}} is returned, if such a point exists. If \mbox{\texttt{\mdseries\slshape coll}} only contains identity partial permutations, then \texttt{SmallestMovedPoint} returns \texttt{infinity}. 
\begin{Verbatim}[commandchars=!@|,fontsize=\small,frame=single,label=Example]
  !gapprompt@gap>| !gapinput@f := PartialPerm( [ 1, 3 ], [ 4, 3 ] );|
  [1,4](3)
  !gapprompt@gap>| !gapinput@SmallestMovedPoint(f);|
  1
  !gapprompt@gap>| !gapinput@SmallestMovedPoint(PartialPerm([1 .. 10]));|
  infinity
\end{Verbatim}
 }

 

\subsection{\textcolor{Chapter }{LargestMovedPoint (for a partial perm)}}
\logpage{[ 54, 3, 13 ]}\nobreak
\hyperdef{L}{X7D4290A785ABC86D}{}
{\noindent\textcolor{FuncColor}{$\triangleright$\enspace\texttt{LargestMovedPoint({\mdseries\slshape f})\index{LargestMovedPoint@\texttt{LargestMovedPoint}!for a partial perm}
\label{LargestMovedPoint:for a partial perm}
}\hfill{\scriptsize (attribute)}}\\
\noindent\textcolor{FuncColor}{$\triangleright$\enspace\texttt{LargestMovedPoint({\mdseries\slshape coll})\index{LargestMovedPoint@\texttt{LargestMovedPoint}!for a partial perm coll}
\label{LargestMovedPoint:for a partial perm coll}
}\hfill{\scriptsize (method)}}\\
\textbf{\indent Returns:\ }
A positive integer or \texttt{infinity}.



 \texttt{LargestMovedPoint} returns the largest positive integers \texttt{i} such that \texttt{i\texttt{\symbol{94}}\mbox{\texttt{\mdseries\slshape f}}{\textless}{\textgreater}i} if such an \texttt{i} exists. If \mbox{\texttt{\mdseries\slshape f}} is the identity partial permutation, then \texttt{0} is returned.

 If the argument is a collection of partial permutations \mbox{\texttt{\mdseries\slshape coll}}, then the largest point which is moved by at least one element of \mbox{\texttt{\mdseries\slshape coll}} is returned, if such a point exists. If \mbox{\texttt{\mdseries\slshape coll}} only contains identity partial permutations, then \texttt{LargestMovedPoint} returns \texttt{0}. 
\begin{Verbatim}[commandchars=!@|,fontsize=\small,frame=single,label=Example]
  !gapprompt@gap>| !gapinput@f := PartialPerm( [ 1, 3, 4, 5 ], [ 5, 1, 6, 4 ] );|
  [3,1,5,4,6]
  !gapprompt@gap>| !gapinput@LargestMovedPoint(f);|
  5
  !gapprompt@gap>| !gapinput@LargestMovedPoint(PartialPerm([1 .. 10]));|
  0
\end{Verbatim}
 }

 

\subsection{\textcolor{Chapter }{SmallestImageOfMovedPoint (for a partial permutation)}}
\logpage{[ 54, 3, 14 ]}\nobreak
\hyperdef{L}{X85280F1A7B1014BA}{}
{\noindent\textcolor{FuncColor}{$\triangleright$\enspace\texttt{SmallestImageOfMovedPoint({\mdseries\slshape f})\index{SmallestImageOfMovedPoint@\texttt{SmallestImageOfMovedPoint}!for a partial permutation}
\label{SmallestImageOfMovedPoint:for a partial permutation}
}\hfill{\scriptsize (attribute)}}\\
\noindent\textcolor{FuncColor}{$\triangleright$\enspace\texttt{SmallestImageOfMovedPoint({\mdseries\slshape coll})\index{SmallestImageOfMovedPoint@\texttt{SmallestImageOfMovedPoint}!for a partial permutation coll}
\label{SmallestImageOfMovedPoint:for a partial permutation coll}
}\hfill{\scriptsize (method)}}\\
\textbf{\indent Returns:\ }
A positive integer or \texttt{infinity}.



 \texttt{SmallestImageOfMovedPoint} returns the smallest positive integer \texttt{i\texttt{\symbol{94}}\mbox{\texttt{\mdseries\slshape f}}} such that \texttt{i\texttt{\symbol{94}}\mbox{\texttt{\mdseries\slshape f}}{\textless}{\textgreater}i} if such an \texttt{i} exists. If \mbox{\texttt{\mdseries\slshape f}} is the identity partial permutation, then \texttt{infinity} is returned.

 If the argument is a collection of partial permutations \mbox{\texttt{\mdseries\slshape coll}}, then the smallest integer which is the image a point moved by at least one
element of \mbox{\texttt{\mdseries\slshape coll}} is returned, if such a point exists. If \mbox{\texttt{\mdseries\slshape coll}} only contains identity partial permutations, then \texttt{SmallestImageOfMovedPoint} returns \texttt{infinity}. 
\begin{Verbatim}[commandchars=!@|,fontsize=\small,frame=single,label=Example]
  !gapprompt@gap>| !gapinput@S := SymmetricInverseSemigroup(5);|
  <symmetric inverse monoid of degree 5>
  !gapprompt@gap>| !gapinput@SmallestImageOfMovedPoint(S);|
  1
  !gapprompt@gap>| !gapinput@S := Semigroup(PartialPerm([10 .. 100], [10 .. 100]));;|
  !gapprompt@gap>| !gapinput@SmallestImageOfMovedPoint(S);|
  infinity
  !gapprompt@gap>| !gapinput@f := PartialPerm( [ 1, 2, 3, 6 ] );|
  [4,6](1)(2)(3)
  !gapprompt@gap>| !gapinput@SmallestImageOfMovedPoint(f);|
  6
\end{Verbatim}
 }

  

\subsection{\textcolor{Chapter }{LargestImageOfMovedPoint (for a partial permutation)}}
\logpage{[ 54, 3, 15 ]}\nobreak
\hyperdef{L}{X7A95CD437BC1CB1A}{}
{\noindent\textcolor{FuncColor}{$\triangleright$\enspace\texttt{LargestImageOfMovedPoint({\mdseries\slshape f})\index{LargestImageOfMovedPoint@\texttt{LargestImageOfMovedPoint}!for a partial permutation}
\label{LargestImageOfMovedPoint:for a partial permutation}
}\hfill{\scriptsize (attribute)}}\\
\noindent\textcolor{FuncColor}{$\triangleright$\enspace\texttt{LargestImageOfMovedPoint({\mdseries\slshape coll})\index{LargestImageOfMovedPoint@\texttt{LargestImageOfMovedPoint}!for a partial permutation coll}
\label{LargestImageOfMovedPoint:for a partial permutation coll}
}\hfill{\scriptsize (method)}}\\
\textbf{\indent Returns:\ }
A positive integer.



 \texttt{LargestImageOfMovedPoint} returns the largest positive integer \texttt{i\texttt{\symbol{94}}\mbox{\texttt{\mdseries\slshape f}}} such that \texttt{i\texttt{\symbol{94}}\mbox{\texttt{\mdseries\slshape f}}{\textless}{\textgreater}i} if such an \texttt{i} exists. If \mbox{\texttt{\mdseries\slshape f}} is an identity partial permutation, then \texttt{0} is returned.

 If the argument is a collection of partial permutations \mbox{\texttt{\mdseries\slshape coll}}, then the largest integer which is the image of a point moved by at least one
element of \mbox{\texttt{\mdseries\slshape coll}} is returned, if such a point exists. If \mbox{\texttt{\mdseries\slshape coll}} only contains identity partial permutations, then \texttt{LargestImageOfMovedPoint} returns \texttt{0}. 
\begin{Verbatim}[commandchars=!@|,fontsize=\small,frame=single,label=Example]
  !gapprompt@gap>| !gapinput@S := SymmetricInverseSemigroup(5);|
  <symmetric inverse monoid of degree 5>
  !gapprompt@gap>| !gapinput@LargestImageOfMovedPoint(S);|
  5
  !gapprompt@gap>| !gapinput@S := Semigroup(PartialPerm([10 .. 100], [10 .. 100]));;|
  !gapprompt@gap>| !gapinput@LargestImageOfMovedPoint(S);|
  0
  !gapprompt@gap>| !gapinput@f := PartialPerm( [ 1, 2, 3, 6 ] );;|
  !gapprompt@gap>| !gapinput@LargestImageOfMovedPoint(f);|
  6
\end{Verbatim}
 }

  

\subsection{\textcolor{Chapter }{IndexPeriodOfPartialPerm}}
\logpage{[ 54, 3, 16 ]}\nobreak
\hyperdef{L}{X873A9F717DA75CBC}{}
{\noindent\textcolor{FuncColor}{$\triangleright$\enspace\texttt{IndexPeriodOfPartialPerm({\mdseries\slshape f})\index{IndexPeriodOfPartialPerm@\texttt{IndexPeriodOfPartialPerm}}
\label{IndexPeriodOfPartialPerm}
}\hfill{\scriptsize (attribute)}}\\
\textbf{\indent Returns:\ }
A pair of positive integers.



 Returns the least positive integers \texttt{m, r} such that \texttt{\mbox{\texttt{\mdseries\slshape f}}\texttt{\symbol{94}}(m+r)=\mbox{\texttt{\mdseries\slshape f}}\texttt{\symbol{94}}m}, which are known as the \emph{index} and \emph{period} of the partial permutation \mbox{\texttt{\mdseries\slshape f}}. 
\begin{Verbatim}[commandchars=!@|,fontsize=\small,frame=single,label=Example]
  !gapprompt@gap>| !gapinput@f:=PartialPerm( [ 1, 2, 3, 5, 6, 7, 8, 11, 12, 16, 19 ],|
  !gapprompt@>| !gapinput@[ 9, 18, 20, 11, 5, 16, 8, 19, 14, 13, 1 ] );|
  [2,18][3,20][6,5,11,19,1,9][7,16,13][12,14](8)
  !gapprompt@gap>| !gapinput@IndexPeriodOfPartialPerm(f);|
  [ 6, 1 ]
  !gapprompt@gap>| !gapinput@f^6=f^7;|
  true
\end{Verbatim}
 }

  

\subsection{\textcolor{Chapter }{SmallestIdempotentPower (for a partial perm)}}
\logpage{[ 54, 3, 17 ]}\nobreak
\hyperdef{L}{X7C04AA377F080722}{}
{\noindent\textcolor{FuncColor}{$\triangleright$\enspace\texttt{SmallestIdempotentPower({\mdseries\slshape f})\index{SmallestIdempotentPower@\texttt{SmallestIdempotentPower}!for a partial perm}
\label{SmallestIdempotentPower:for a partial perm}
}\hfill{\scriptsize (attribute)}}\\
\textbf{\indent Returns:\ }
A positive integer.



 This function returns the least positive integer \texttt{n} such that the partial permutation \texttt{\mbox{\texttt{\mdseries\slshape f}}\texttt{\symbol{94}}n} is an idempotent. The smallest idempotent power of \mbox{\texttt{\mdseries\slshape f}} is the least multiple of the period of \mbox{\texttt{\mdseries\slshape f}} that is greater than or equal to the index of \mbox{\texttt{\mdseries\slshape f}}; see \texttt{IndexPeriodOfPartialPerm} (\ref{IndexPeriodOfPartialPerm}). 
\begin{Verbatim}[commandchars=!@|,fontsize=\small,frame=single,label=Example]
  !gapprompt@gap>| !gapinput@f:=PartialPerm( [ 1, 2, 3, 4, 5, 7, 8, 10, 11, 13, 18, 19, 20 ],|
  !gapprompt@>| !gapinput@[ 5, 1, 7, 3, 10, 2, 12, 14, 11, 16, 6, 9, 15 ] );|
  [4,3,7,2,1,5,10,14][8,12][13,16][18,6][19,9][20,15](11)
  !gapprompt@gap>| !gapinput@SmallestIdempotentPower(f);|
  8
  !gapprompt@gap>| !gapinput@f^8;|
  <identity partial perm on [ 11 ]>
\end{Verbatim}
 }

  

\subsection{\textcolor{Chapter }{ComponentsOfPartialPerm}}
\logpage{[ 54, 3, 18 ]}\nobreak
\hyperdef{L}{X8185065E788BDD0D}{}
{\noindent\textcolor{FuncColor}{$\triangleright$\enspace\texttt{ComponentsOfPartialPerm({\mdseries\slshape f})\index{ComponentsOfPartialPerm@\texttt{ComponentsOfPartialPerm}}
\label{ComponentsOfPartialPerm}
}\hfill{\scriptsize (attribute)}}\\
\textbf{\indent Returns:\ }
A list of lists of positive integer.



 \texttt{ComponentsOfPartialPerm} returns a list of the components of the partial permutation \mbox{\texttt{\mdseries\slshape f}}. Each component is a subset of the domain of \mbox{\texttt{\mdseries\slshape f}}, and the union of the components equals the domain. 
\begin{Verbatim}[commandchars=!@|,fontsize=\small,frame=single,label=Example]
  !gapprompt@gap>| !gapinput@f:=PartialPerm( [ 1, 2, 3, 4, 5, 7, 8, 10, 11, 12, 13, 19 ],|
  !gapprompt@>| !gapinput@[ 20, 4, 6, 19, 9, 14, 3, 12, 17, 5, 15, 13 ] );|
  [1,20][2,4,19,13,15][7,14][8,3,6][10,12,5,9][11,17]
  !gapprompt@gap>| !gapinput@ComponentsOfPartialPerm(f);|
  [ [ 1, 20 ], [ 2, 4, 19, 13, 15 ], [ 7, 14 ], [ 8, 3, 6 ],
    [ 10, 12, 5, 9 ], [ 11, 17 ] ]
\end{Verbatim}
 }

  

\subsection{\textcolor{Chapter }{NrComponentsOfPartialPerm}}
\logpage{[ 54, 3, 19 ]}\nobreak
\hyperdef{L}{X7CB51EB67FFA95E9}{}
{\noindent\textcolor{FuncColor}{$\triangleright$\enspace\texttt{NrComponentsOfPartialPerm({\mdseries\slshape f})\index{NrComponentsOfPartialPerm@\texttt{NrComponentsOfPartialPerm}}
\label{NrComponentsOfPartialPerm}
}\hfill{\scriptsize (attribute)}}\\
\textbf{\indent Returns:\ }
A positive integer.



 \texttt{NrComponentsOfPartialPerm} returns the number of components of the partial permutation \mbox{\texttt{\mdseries\slshape f}} on its domain. 
\begin{Verbatim}[commandchars=!@|,fontsize=\small,frame=single,label=Example]
  !gapprompt@gap>| !gapinput@f:=PartialPerm( [ 1, 2, 3, 4, 5, 7, 8, 10, 11, 12, 13, 19 ],|
  !gapprompt@>| !gapinput@[ 20, 4, 6, 19, 9, 14, 3, 12, 17, 5, 15, 13 ] );|
  [1,20][2,4,19,13,15][7,14][8,3,6][10,12,5,9][11,17]
  !gapprompt@gap>| !gapinput@NrComponentsOfPartialPerm(f);|
  6
\end{Verbatim}
 }

  

\subsection{\textcolor{Chapter }{ComponentRepsOfPartialPerm}}
\logpage{[ 54, 3, 20 ]}\nobreak
\hyperdef{L}{X7AAAAE4082B30E18}{}
{\noindent\textcolor{FuncColor}{$\triangleright$\enspace\texttt{ComponentRepsOfPartialPerm({\mdseries\slshape f})\index{ComponentRepsOfPartialPerm@\texttt{ComponentRepsOfPartialPerm}}
\label{ComponentRepsOfPartialPerm}
}\hfill{\scriptsize (attribute)}}\\
\textbf{\indent Returns:\ }
A list of positive integers.



 \texttt{ComponentRepsOfPartialPerm} returns the representatives, in the following sense, of the components of the
partial permutation \mbox{\texttt{\mdseries\slshape f}}. Every component of \mbox{\texttt{\mdseries\slshape f}} contains a unique element in the domain but not the image of \mbox{\texttt{\mdseries\slshape f}}; this element is called the \emph{representative} of the component. If \texttt{i} is a representative of a component of \mbox{\texttt{\mdseries\slshape f}}, then for every \texttt{j}$\not=$\texttt{i} in the component of \texttt{i}, there exists a positive integer \texttt{k} such that \texttt{i \texttt{\symbol{94}} (\mbox{\texttt{\mdseries\slshape f}} \texttt{\symbol{94}} k) = j}. Unlike transformations, there is exactly one representative for every
component of \mbox{\texttt{\mdseries\slshape f}}. \texttt{ComponentRepsOfPartialPerm} returns the least number of representatives. 
\begin{Verbatim}[commandchars=!@|,fontsize=\small,frame=single,label=Example]
  !gapprompt@gap>| !gapinput@f:=PartialPerm( [ 1, 2, 3, 4, 5, 7, 8, 10, 11, 12, 13, 19 ],|
  !gapprompt@>| !gapinput@[ 20, 4, 6, 19, 9, 14, 3, 12, 17, 5, 15, 13 ] );|
  [1,20][2,4,19,13,15][7,14][8,3,6][10,12,5,9][11,17]
  !gapprompt@gap>| !gapinput@ComponentRepsOfPartialPerm(f);|
  [ 1, 2, 7, 8, 10, 11 ]
\end{Verbatim}
 }

  

\subsection{\textcolor{Chapter }{LeftOne (for a partial perm)}}
\logpage{[ 54, 3, 21 ]}\nobreak
\hyperdef{L}{X7A8FB86C78C49F85}{}
{\noindent\textcolor{FuncColor}{$\triangleright$\enspace\texttt{LeftOne({\mdseries\slshape f})\index{LeftOne@\texttt{LeftOne}!for a partial perm}
\label{LeftOne:for a partial perm}
}\hfill{\scriptsize (attribute)}}\\
\noindent\textcolor{FuncColor}{$\triangleright$\enspace\texttt{RightOne({\mdseries\slshape f})\index{RightOne@\texttt{RightOne}!for a partial perm}
\label{RightOne:for a partial perm}
}\hfill{\scriptsize (attribute)}}\\
\textbf{\indent Returns:\ }
A partial permutation.



 \texttt{LeftOne} returns the identity partial permutation \texttt{e} such that the domain and image of \texttt{e} equal the domain of the partial permutation \mbox{\texttt{\mdseries\slshape f}} and such that \texttt{e*\mbox{\texttt{\mdseries\slshape f}}=f}. 

 \texttt{RightOne} returns the identity partial permutation \texttt{e} such that the domain and image of \texttt{e} equal the image of \mbox{\texttt{\mdseries\slshape f}} and such that \texttt{\mbox{\texttt{\mdseries\slshape f}}*e=f}. 
\begin{Verbatim}[commandchars=!@|,fontsize=\small,frame=single,label=Example]
  !gapprompt@gap>| !gapinput@f:=PartialPerm( [ 1, 2, 4, 5, 6, 7 ], [ 10, 1, 6, 5, 8, 7 ] );|
  [2,1,10][4,6,8](5)(7)
  !gapprompt@gap>| !gapinput@RightOne(f);|
  <identity partial perm on [ 1, 5, 6, 7, 8, 10 ]>
  !gapprompt@gap>| !gapinput@LeftOne(f);|
  <identity partial perm on [ 1, 2, 4, 5, 6, 7 ]>
\end{Verbatim}
 }

  

\subsection{\textcolor{Chapter }{One (for a partial perm)}}
\logpage{[ 54, 3, 22 ]}\nobreak
\hyperdef{L}{X857FC10C81507E8B}{}
{\noindent\textcolor{FuncColor}{$\triangleright$\enspace\texttt{One({\mdseries\slshape f})\index{One@\texttt{One}!for a partial perm}
\label{One:for a partial perm}
}\hfill{\scriptsize (method)}}\\
\textbf{\indent Returns:\ }
A partial permutation.



 As described in \texttt{OneImmutable} (\ref{OneImmutable}), \texttt{One} returns the multiplicative neutral element of the partial permutation \mbox{\texttt{\mdseries\slshape f}}, which is the identity partial permutation on the union of the domain and
image of \mbox{\texttt{\mdseries\slshape f}}. Equivalently, the one of \mbox{\texttt{\mdseries\slshape f}} is the join of the right one and left one of \mbox{\texttt{\mdseries\slshape f}}. 
\begin{Verbatim}[commandchars=!@|,fontsize=\small,frame=single,label=Example]
  !gapprompt@gap>| !gapinput@f:=PartialPerm([ 1, 2, 3, 4, 5, 7, 10 ], [ 3, 7, 9, 6, 1, 10, 2 ]);;|
  !gapprompt@gap>| !gapinput@One(f);|
  <identity partial perm on [ 1, 2, 3, 4, 5, 6, 7, 9, 10 ]>
\end{Verbatim}
 }

  

\subsection{\textcolor{Chapter }{MultiplicativeZero (for a partial perm)}}
\logpage{[ 54, 3, 23 ]}\nobreak
\hyperdef{L}{X7D90CF497D58D759}{}
{\noindent\textcolor{FuncColor}{$\triangleright$\enspace\texttt{MultiplicativeZero({\mdseries\slshape f})\index{MultiplicativeZero@\texttt{MultiplicativeZero}!for a partial perm}
\label{MultiplicativeZero:for a partial perm}
}\hfill{\scriptsize (method)}}\\
\textbf{\indent Returns:\ }
The empty partial permutation.



 As described in \texttt{MultiplicativeZero} (\ref{MultiplicativeZero}), \texttt{Zero} returns the multiplicative zero element of the partial permutation \mbox{\texttt{\mdseries\slshape f}}, which is the empty partial permutation. 
\begin{Verbatim}[commandchars=!@|,fontsize=\small,frame=single,label=Example]
  !gapprompt@gap>| !gapinput@f := PartialPerm([ 1, 2, 3, 4, 5, 7, 10 ], [ 3, 7, 9, 6, 1, 10, 2 ]);;|
  !gapprompt@gap>| !gapinput@MultiplicativeZero(f);|
  <empty partial perm>
\end{Verbatim}
 }

  }

  
\section{\textcolor{Chapter }{Changing the representation of a partial permutation}}\label{sect:ChangingRepPartialPerms}
\logpage{[ 54, 4, 0 ]}
\hyperdef{L}{X8585AA8B78E9CDFB}{}
{
  It is possible that a partial permutation in \textsf{GAP} can be represented by other types of objects, or that other types of \textsf{GAP} objects can be represented by partial permutations. Partial permutations which
are mathematically permutations can be converted into permutations in \textsf{GAP} using the function \texttt{AsPermutation} (\ref{AsPermutation}). Similarly, a partial permutation can be converted into a transformation
using \texttt{AsTransformation} (\ref{AsTransformation}). 

 In this section we describe functions for converting other types of objects in \textsf{GAP} into partial permutations. 

\subsection{\textcolor{Chapter }{AsPartialPerm (for a permutation and a set of
    positive integers)}}
\logpage{[ 54, 4, 1 ]}\nobreak
\hyperdef{L}{X81B32CB182489ACA}{}
{\noindent\textcolor{FuncColor}{$\triangleright$\enspace\texttt{AsPartialPerm({\mdseries\slshape f, set})\index{AsPartialPerm@\texttt{AsPartialPerm}!for a permutation and a set of
    positive integers}
\label{AsPartialPerm:for a permutation and a set of
    positive integers}
}\hfill{\scriptsize (operation)}}\\
\noindent\textcolor{FuncColor}{$\triangleright$\enspace\texttt{AsPartialPerm({\mdseries\slshape f})\index{AsPartialPerm@\texttt{AsPartialPerm}!for a permutation}
\label{AsPartialPerm:for a permutation}
}\hfill{\scriptsize (method)}}\\
\noindent\textcolor{FuncColor}{$\triangleright$\enspace\texttt{AsPartialPerm({\mdseries\slshape f, n})\index{AsPartialPerm@\texttt{AsPartialPerm}!for a permutation and a positive integer}
\label{AsPartialPerm:for a permutation and a positive integer}
}\hfill{\scriptsize (method)}}\\
\textbf{\indent Returns:\ }
A partial permutation.



 A permutation \mbox{\texttt{\mdseries\slshape f}} defines a partial permutation when it is restricted to any finite set of
positive integers. \texttt{AsPartialPerm} can be used to obtain this partial permutation.

 There are several possible arguments for \texttt{AsPartialPerm}: 
\begin{description}
\item[{for a permutation and set of positive integers}] \texttt{AsPartialPerm} returns the partial permutation that equals \mbox{\texttt{\mdseries\slshape f}} on the set of positive integers \mbox{\texttt{\mdseries\slshape set}} and that is undefined on every other positive integer. 

 Note that as explained in \texttt{PartialPerm} (\ref{PartialPerm:for a domain and image}) \emph{a permutation is never a partial permutation} in \textsf{GAP}, please keep this in mind when using \texttt{AsPartialPerm}. 
\item[{for a permutation}] \texttt{AsPartialPerm} returns the partial permutation that agrees with \mbox{\texttt{\mdseries\slshape f}} on \texttt{[1..LargestMovedPoint(\mbox{\texttt{\mdseries\slshape f}})]} and that is not defined on any other positive integer. 
\item[{for a permutation and a positive integer}] \texttt{AsPartialPerm} returns the partial permutation that agrees with \mbox{\texttt{\mdseries\slshape f}} on \texttt{[1..\mbox{\texttt{\mdseries\slshape n}}]}, when \mbox{\texttt{\mdseries\slshape n}} is a positive integer, and that is not defined on any other positive integer. 
\end{description}
 The operation \texttt{PartialPermOp} (\ref{PartialPermOp}) can also be used to convert permutations into partial permutations. 
\begin{Verbatim}[commandchars=!@|,fontsize=\small,frame=single,label=Example]
  !gapprompt@gap>| !gapinput@f:=(2,8,19,9,14,10,20,17,4,13,12,3,5,7,18,16);;|
  !gapprompt@gap>| !gapinput@AsPartialPerm(f);|
  (1)(2,8,19,9,14,10,20,17,4,13,12,3,5,7,18,16)(6)(11)(15)
  !gapprompt@gap>| !gapinput@AsPartialPerm(f, [ 1, 2, 3 ] );|
  [2,8][3,5](1)
\end{Verbatim}
 }

  

\subsection{\textcolor{Chapter }{AsPartialPerm (for a transformation and a set of positive integer)}}
\logpage{[ 54, 4, 2 ]}\nobreak
\hyperdef{L}{X87EC67747B260E98}{}
{\noindent\textcolor{FuncColor}{$\triangleright$\enspace\texttt{AsPartialPerm({\mdseries\slshape f, set})\index{AsPartialPerm@\texttt{AsPartialPerm}!for a transformation and a set of positive integer}
\label{AsPartialPerm:for a transformation and a set of positive integer}
}\hfill{\scriptsize (operation)}}\\
\noindent\textcolor{FuncColor}{$\triangleright$\enspace\texttt{AsPartialPerm({\mdseries\slshape f, n})\index{AsPartialPerm@\texttt{AsPartialPerm}!for a transformation and a positive integer}
\label{AsPartialPerm:for a transformation and a positive integer}
}\hfill{\scriptsize (method)}}\\
\textbf{\indent Returns:\ }
A partial permutation or \texttt{fail}.



 A transformation \mbox{\texttt{\mdseries\slshape f}} defines a partial permutation when it is restricted to a set of positive
integers where it is injective. \texttt{AsPartialPerm} can be used to obtain this partial permutation. 

 There are several possible arguments for \texttt{AsPartialPerm}: 
\begin{description}
\item[{for a transformation and set of positive integers}]  \texttt{AsPartialPerm} returns the partial permutation obtained by restricting \mbox{\texttt{\mdseries\slshape f}} to the set of positive integers \mbox{\texttt{\mdseries\slshape set}} when: 
\begin{itemize}
\item \mbox{\texttt{\mdseries\slshape set}} contains no elements exceeding the degree of \mbox{\texttt{\mdseries\slshape f}}; 
\item \mbox{\texttt{\mdseries\slshape f}} is injective on \mbox{\texttt{\mdseries\slshape set}}. 
\end{itemize}
 
\item[{for a transformation and a positive integer}] \texttt{AsPartialPerm} returns the partial permutation that agrees with \mbox{\texttt{\mdseries\slshape f}} on \texttt{[1..\mbox{\texttt{\mdseries\slshape n}}]} when \mbox{\texttt{\mdseries\slshape A}} is a positive integer and this set satisfies the conditions given above. 
\end{description}
 The operation \texttt{PartialPermOp} (\ref{PartialPermOp}) can also be used to convert transformations into partial permutations. 
\begin{Verbatim}[commandchars=!@|,fontsize=\small,frame=single,label=Example]
  !gapprompt@gap>| !gapinput@f:=Transformation( [ 8, 3, 5, 9, 6, 2, 9, 7, 9 ] );;|
  !gapprompt@gap>| !gapinput@AsPartialPerm(f, [1, 2, 3, 5, 8]);|
  [1,8,7][2,3,5,6]
  !gapprompt@gap>| !gapinput@AsPartialPerm(f, 3);|
  [1,8][2,3,5]
  !gapprompt@gap>| !gapinput@AsPartialPerm(f, [ 2 .. 4 ] );|
  [2,3,5][4,9]
\end{Verbatim}
 }

 }

   
\section{\textcolor{Chapter }{Operators and operations for partial permutations}}\label{sect:OperatorsPartialPerms}
\logpage{[ 54, 5, 0 ]}
\hyperdef{L}{X848CD855802C6CE1}{}
{
  

\subsection{\textcolor{Chapter }{Inverse (for a partial permutation)}}
\logpage{[ 54, 5, 1 ]}\nobreak
\hyperdef{L}{X7B8630027B7F0BCC}{}
{\noindent\textcolor{FuncColor}{$\triangleright$\enspace\texttt{Inverse({\mdseries\slshape f})\index{Inverse@\texttt{Inverse}!for a partial permutation}
\label{Inverse:for a partial permutation}
}\hfill{\scriptsize (method)}}\\


 returns the inverse of the partial permutation \mbox{\texttt{\mdseries\slshape f}}. }

 

\subsection{\textcolor{Chapter }{\texttt{\symbol{92}}\texttt{\symbol{94}} (for a positive integer and a partial permutation)}}
\logpage{[ 54, 5, 2 ]}\nobreak
\hyperdef{L}{X792D3BA278DAB869}{}
{\noindent\textcolor{FuncColor}{$\triangleright$\enspace\texttt{\texttt{\symbol{92}}\texttt{\symbol{94}}({\mdseries\slshape i, f})\index{\texttt{\symbol{92}}\texttt{\symbol{94}}@\texttt{\texttt{\symbol{92}}\texttt{\symbol{94}}}!for a positive integer and a partial permutation}
\label{bSlash^:for a positive integer and a partial permutation}
}\hfill{\scriptsize (method)}}\\


 returns the image of the positive integer \mbox{\texttt{\mdseries\slshape i}} under the partial permutation \mbox{\texttt{\mdseries\slshape f}} if it is defined and \texttt{0} if it is not. }

 

\subsection{\textcolor{Chapter }{\texttt{\symbol{92}}/ (for a positive integer and a partial permutation)}}
\logpage{[ 54, 5, 3 ]}\nobreak
\hyperdef{L}{X7DC25BC47AAA9C73}{}
{\noindent\textcolor{FuncColor}{$\triangleright$\enspace\texttt{\texttt{\symbol{92}}/({\mdseries\slshape i, f})\index{\texttt{\symbol{92}}/@\texttt{\texttt{\symbol{92}}/}!for a positive integer and a partial permutation}
\label{bSlash/:for a positive integer and a partial permutation}
}\hfill{\scriptsize (method)}}\\


 returns the preimage of the positive integer \mbox{\texttt{\mdseries\slshape i}} under the partial permutation \mbox{\texttt{\mdseries\slshape f}} if it is defined and \texttt{0} if it is not. Note that the inverse of \mbox{\texttt{\mdseries\slshape f}} is not calculated to find the preimage of \mbox{\texttt{\mdseries\slshape i}}. }

 

\subsection{\textcolor{Chapter }{\texttt{\symbol{92}}\texttt{\symbol{94}} (for a partial permutation and a permutation or partial permutation)}}
\logpage{[ 54, 5, 4 ]}\nobreak
\hyperdef{L}{X8213CD6E7C461169}{}
{\noindent\textcolor{FuncColor}{$\triangleright$\enspace\texttt{\texttt{\symbol{92}}\texttt{\symbol{94}}({\mdseries\slshape f, g})\index{\texttt{\symbol{92}}\texttt{\symbol{94}}@\texttt{\texttt{\symbol{92}}\texttt{\symbol{94}}}!for a partial permutation and a permutation or partial permutation}
\label{bSlash^:for a partial permutation and a permutation or partial permutation}
}\hfill{\scriptsize (method)}}\\


 \texttt{\mbox{\texttt{\mdseries\slshape f}} \texttt{\symbol{94}} \mbox{\texttt{\mdseries\slshape g}}} returns \texttt{\mbox{\texttt{\mdseries\slshape g}}\texttt{\symbol{94}}\texttt{\symbol{45}}1*\mbox{\texttt{\mdseries\slshape f}}*\mbox{\texttt{\mdseries\slshape g}}} when \mbox{\texttt{\mdseries\slshape f}} is a partial permutation and \mbox{\texttt{\mdseries\slshape g}} is a permutation or partial permutation; see \texttt{\texttt{\symbol{92}}\texttt{\symbol{94}}} (\ref{bSlash^}). This operation requires essentially the same number of steps as multiplying
partial permutations, which is around one third as many as inverting and
multiplying twice. }

 

\subsection{\textcolor{Chapter }{\texttt{\symbol{92}}* (for permutations and partial permutations)}}
\logpage{[ 54, 5, 5 ]}\nobreak
\hyperdef{L}{X86469D597F8BC7CE}{}
{\noindent\textcolor{FuncColor}{$\triangleright$\enspace\texttt{\texttt{\symbol{92}}*({\mdseries\slshape f, g})\index{\texttt{\symbol{92}}*@\texttt{\texttt{\symbol{92}}*}!for permutations and partial permutations}
\label{bSlash*:for permutations and partial permutations}
}\hfill{\scriptsize (method)}}\\


 \texttt{\mbox{\texttt{\mdseries\slshape f}} * \mbox{\texttt{\mdseries\slshape g}}} returns the composition of \mbox{\texttt{\mdseries\slshape f}} and \mbox{\texttt{\mdseries\slshape g}} when \mbox{\texttt{\mdseries\slshape f}} and \mbox{\texttt{\mdseries\slshape g}} are partial permutations or permutations. The product of a permutation and a
partial permutation is returned as a partial permutation. }

 

\subsection{\textcolor{Chapter }{\texttt{\symbol{92}}/ (for a partial permutation and permutation or partial permutation)}}
\logpage{[ 54, 5, 6 ]}\nobreak
\hyperdef{L}{X869DBDF67FA3817B}{}
{\noindent\textcolor{FuncColor}{$\triangleright$\enspace\texttt{\texttt{\symbol{92}}/({\mdseries\slshape f, g})\index{\texttt{\symbol{92}}/@\texttt{\texttt{\symbol{92}}/}!for a partial permutation and permutation or partial permutation}
\label{bSlash/:for a partial permutation and permutation or partial permutation}
}\hfill{\scriptsize (method)}}\\


 \texttt{\mbox{\texttt{\mdseries\slshape f}} / \mbox{\texttt{\mdseries\slshape g}}} returns \texttt{\mbox{\texttt{\mdseries\slshape f}}*\mbox{\texttt{\mdseries\slshape g}}\texttt{\symbol{94}}\texttt{\symbol{45}}1} when \mbox{\texttt{\mdseries\slshape f}} is a partial permutation and \mbox{\texttt{\mdseries\slshape g}} is a permutation or partial permutation. This operation requires essentially
the same number of steps as multiplying partial permutations, which is
approximately half that required to first invert \mbox{\texttt{\mdseries\slshape g}} and then take the product with \mbox{\texttt{\mdseries\slshape f}}. }

 

\subsection{\textcolor{Chapter }{LeftQuotient (for a permutation or partial permutation
        and a partial permutation)}}
\logpage{[ 54, 5, 7 ]}\nobreak
\hyperdef{L}{X82E3A3E186A4F2D2}{}
{\noindent\textcolor{FuncColor}{$\triangleright$\enspace\texttt{LeftQuotient({\mdseries\slshape g, f})\index{LeftQuotient@\texttt{LeftQuotient}!for a permutation or partial permutation
        and a partial permutation}
\label{LeftQuotient:for a permutation or partial permutation
        and a partial permutation}
}\hfill{\scriptsize (method)}}\\


 returns \texttt{\mbox{\texttt{\mdseries\slshape g}}\texttt{\symbol{94}}\texttt{\symbol{45}}1*\mbox{\texttt{\mdseries\slshape f}}} when \mbox{\texttt{\mdseries\slshape f}} is a partial permutation and \mbox{\texttt{\mdseries\slshape g}} is a permutation or partial permutation. This operation requires essentially
the same number of steps as multiplying partial permutations, which is
approximately half that required to first invert \mbox{\texttt{\mdseries\slshape g}} and then take the product with \mbox{\texttt{\mdseries\slshape f}}. }

 

\subsection{\textcolor{Chapter }{\texttt{\symbol{92}}{\textless} (for partial permutations)}}
\logpage{[ 54, 5, 8 ]}\nobreak
\hyperdef{L}{X8659E9E57AC8D9CE}{}
{\noindent\textcolor{FuncColor}{$\triangleright$\enspace\texttt{\texttt{\symbol{92}}{\textless}({\mdseries\slshape f, g})\index{\texttt{\symbol{92}}{\textless}@\texttt{\texttt{\symbol{92}}{\textless}}!for partial permutations}
\label{bSlash<:for partial permutations}
}\hfill{\scriptsize (method)}}\\


 \texttt{\mbox{\texttt{\mdseries\slshape f}} {\textless} \mbox{\texttt{\mdseries\slshape g}}} returns \texttt{true} if the image of \mbox{\texttt{\mdseries\slshape f}} on the range from 1 to the degree of \mbox{\texttt{\mdseries\slshape f}} is lexicographically less than the corresponding image for \mbox{\texttt{\mdseries\slshape g}} and \texttt{false} if it is not. See \texttt{NaturalLeqPartialPerm} (\ref{NaturalLeqPartialPerm}) and \texttt{ShortLexLeqPartialPerm} (\ref{ShortLexLeqPartialPerm}) for additional orders for partial permutations. }

 

\subsection{\textcolor{Chapter }{\texttt{\symbol{92}}= (for partial permutations)}}
\logpage{[ 54, 5, 9 ]}\nobreak
\hyperdef{L}{X7828338C7DB8AAC7}{}
{\noindent\textcolor{FuncColor}{$\triangleright$\enspace\texttt{\texttt{\symbol{92}}=({\mdseries\slshape f, g})\index{\texttt{\symbol{92}}=@\texttt{\texttt{\symbol{92}}=}!for partial permutations}
\label{bSlash=:for partial permutations}
}\hfill{\scriptsize (method)}}\\


 \texttt{\mbox{\texttt{\mdseries\slshape f}} = \mbox{\texttt{\mdseries\slshape g}}} returns \texttt{true} if the partial permutation \mbox{\texttt{\mdseries\slshape f}} equals the partial permutation \mbox{\texttt{\mdseries\slshape g}} and returns \texttt{false} if it does not. }

  

\subsection{\textcolor{Chapter }{PermLeftQuoPartialPerm}}
\logpage{[ 54, 5, 10 ]}\nobreak
\hyperdef{L}{X8382A0F8875CEB08}{}
{\noindent\textcolor{FuncColor}{$\triangleright$\enspace\texttt{PermLeftQuoPartialPerm({\mdseries\slshape f, g})\index{PermLeftQuoPartialPerm@\texttt{PermLeftQuoPartialPerm}}
\label{PermLeftQuoPartialPerm}
}\hfill{\scriptsize (operation)}}\\
\noindent\textcolor{FuncColor}{$\triangleright$\enspace\texttt{PermLeftQuoPartialPermNC({\mdseries\slshape f, g})\index{PermLeftQuoPartialPermNC@\texttt{PermLeftQuoPartialPermNC}}
\label{PermLeftQuoPartialPermNC}
}\hfill{\scriptsize (operation)}}\\
\textbf{\indent Returns:\ }
A permutation.



 Returns the permutation on the image set of \mbox{\texttt{\mdseries\slshape f}} induced by \texttt{\mbox{\texttt{\mdseries\slshape f}}\texttt{\symbol{94}}\texttt{\symbol{45}}1*\mbox{\texttt{\mdseries\slshape g}}} when the partial permutations \mbox{\texttt{\mdseries\slshape f}} and \mbox{\texttt{\mdseries\slshape g}} have equal domain and image set. 

 \texttt{PermLeftQuoPartialPerm} verifies that \mbox{\texttt{\mdseries\slshape f}} and \mbox{\texttt{\mdseries\slshape g}} have equal domains and image sets, and returns an error if they do not. \texttt{PermLeftQuoPartialPermNC} does no checks. 
\begin{Verbatim}[commandchars=!@|,fontsize=\small,frame=single,label=Example]
  !gapprompt@gap>| !gapinput@f:=PartialPerm( [ 1, 2, 3, 4, 5, 7 ], [ 7, 9, 10, 4, 2, 5 ] );|
  [1,7,5,2,9][3,10](4)
  !gapprompt@gap>| !gapinput@g:=PartialPerm( [ 1, 2, 3, 4, 5, 7 ], [ 7, 4, 9, 2, 5, 10 ] );|
  [1,7,10][3,9](2,4)(5)
  !gapprompt@gap>| !gapinput@PermLeftQuoPartialPerm(f, g);|
  (2,5,10,9,4)
\end{Verbatim}
 }

  

\subsection{\textcolor{Chapter }{PreImagePartialPerm}}
\logpage{[ 54, 5, 11 ]}\nobreak
\hyperdef{L}{X7C7F5EAB7E9A381D}{}
{\noindent\textcolor{FuncColor}{$\triangleright$\enspace\texttt{PreImagePartialPerm({\mdseries\slshape f, i})\index{PreImagePartialPerm@\texttt{PreImagePartialPerm}}
\label{PreImagePartialPerm}
}\hfill{\scriptsize (operation)}}\\
\textbf{\indent Returns:\ }
A positive integer or \texttt{fail}.



 \texttt{PreImagePartialPerm} returns the preimage of the positive integer \mbox{\texttt{\mdseries\slshape i}} under the partial permutation \mbox{\texttt{\mdseries\slshape f}} if \mbox{\texttt{\mdseries\slshape i}} belongs to the image of \mbox{\texttt{\mdseries\slshape f}}. If \mbox{\texttt{\mdseries\slshape i}} does not belong to the image of \mbox{\texttt{\mdseries\slshape f}}, then \texttt{fail} is returned. 

 The same result can be obtained by using \texttt{\mbox{\texttt{\mdseries\slshape i}}/\mbox{\texttt{\mdseries\slshape f}}} as described in Section \ref{sect:OperatorsPartialPerms}. 
\begin{Verbatim}[commandchars=!@|,fontsize=\small,frame=single,label=Example]
  !gapprompt@gap>| !gapinput@f:=PartialPerm( [ 1, 2, 3, 5, 9, 10 ], [ 5, 10, 7, 8, 9, 1 ] );|
  [2,10,1,5,8][3,7](9)
  !gapprompt@gap>| !gapinput@PreImagePartialPerm(f, 8);|
  5
  !gapprompt@gap>| !gapinput@PreImagePartialPerm(f, 5);|
  1
  !gapprompt@gap>| !gapinput@PreImagePartialPerm(f, 1);|
  10
  !gapprompt@gap>| !gapinput@PreImagePartialPerm(f, 10);|
  2
  !gapprompt@gap>| !gapinput@PreImagePartialPerm(f, 2);|
  fail
\end{Verbatim}
 }

  

\subsection{\textcolor{Chapter }{ComponentPartialPermInt}}
\logpage{[ 54, 5, 12 ]}\nobreak
\hyperdef{L}{X797A6CC084068219}{}
{\noindent\textcolor{FuncColor}{$\triangleright$\enspace\texttt{ComponentPartialPermInt({\mdseries\slshape f, i})\index{ComponentPartialPermInt@\texttt{ComponentPartialPermInt}}
\label{ComponentPartialPermInt}
}\hfill{\scriptsize (operation)}}\\
\textbf{\indent Returns:\ }
A set of positive integers.



 \texttt{ComponentPartialPermInt} returns the elements of the component of \mbox{\texttt{\mdseries\slshape f}} containing \mbox{\texttt{\mdseries\slshape i}} that can be obtained by repeatedly applying \mbox{\texttt{\mdseries\slshape f}} to \mbox{\texttt{\mdseries\slshape i}}. 
\begin{Verbatim}[commandchars=!@|,fontsize=\small,frame=single,label=Example]
  !gapprompt@gap>| !gapinput@f:=PartialPerm( [ 1, 2, 4, 5, 6, 7, 8, 10, 14, 15, 16, 17, 18 ],|
  !gapprompt@>| !gapinput@[ 11, 4, 14, 16, 15, 3, 20, 8, 17, 19, 1, 6, 12 ] );|
  [2,4,14,17,6,15,19][5,16,1,11][7,3][10,8,20][18,12]
  !gapprompt@gap>| !gapinput@ComponentPartialPermInt(f, 4);|
  [ 4, 14, 17, 6, 15, 19 ]
  !gapprompt@gap>| !gapinput@ComponentPartialPermInt(f, 3);|
  [  ]
  !gapprompt@gap>| !gapinput@ComponentPartialPermInt(f, 10);|
  [ 10, 8, 20 ]
  !gapprompt@gap>| !gapinput@ComponentPartialPermInt(f, 100);|
  [  ]
\end{Verbatim}
 }

  

\subsection{\textcolor{Chapter }{NaturalLeqPartialPerm}}
\logpage{[ 54, 5, 13 ]}\nobreak
\hyperdef{L}{X87B1ED93785257C1}{}
{\noindent\textcolor{FuncColor}{$\triangleright$\enspace\texttt{NaturalLeqPartialPerm({\mdseries\slshape f, g})\index{NaturalLeqPartialPerm@\texttt{NaturalLeqPartialPerm}}
\label{NaturalLeqPartialPerm}
}\hfill{\scriptsize (function)}}\\
\textbf{\indent Returns:\ }
\texttt{true} or \texttt{false}.



 The \emph{natural partial order} $\leq$ on an inverse semigroup \texttt{S} is defined by \texttt{s}$\leq$\texttt{t} if there exists an idempotent \texttt{e} in \texttt{S} such that \texttt{s=et}. Hence if \mbox{\texttt{\mdseries\slshape f}} and \mbox{\texttt{\mdseries\slshape g}} are partial permutations, then \mbox{\texttt{\mdseries\slshape f}}$\leq$\mbox{\texttt{\mdseries\slshape g}} if and only if \mbox{\texttt{\mdseries\slshape f}} is a restriction of \mbox{\texttt{\mdseries\slshape g}}; see \texttt{RestrictedPartialPerm} (\ref{RestrictedPartialPerm}). 

 \texttt{NaturalLeqPartialPerm} returns \texttt{true} if \mbox{\texttt{\mdseries\slshape f}} is a restriction of \mbox{\texttt{\mdseries\slshape g}} and \texttt{false} if it is not. Note that since this is a partial order and not a total order,
it is possible that \mbox{\texttt{\mdseries\slshape f}} and \mbox{\texttt{\mdseries\slshape g}} are incomparable with respect to the natural partial order. 
\begin{Verbatim}[commandchars=!@|,fontsize=\small,frame=single,label=Example]
  !gapprompt@gap>| !gapinput@f:=PartialPerm(|
  !gapprompt@>| !gapinput@[ 1, 2, 3, 4, 5, 6, 7, 8, 10, 11, 12, 16, 17, 18, 19 ],|
  !gapprompt@>| !gapinput@[ 3, 12, 14, 4, 11, 18, 17, 2, 9, 5, 15, 8, 20, 10, 19 ] );;|
  !gapprompt@gap>| !gapinput@g:=RestrictedPartialPerm(f, [ 1, 2, 3, 9, 13, 20 ] );|
  [1,3,14][2,12]
  !gapprompt@gap>| !gapinput@NaturalLeqPartialPerm(g,f);|
  true
  !gapprompt@gap>| !gapinput@NaturalLeqPartialPerm(f,g);|
  false
  !gapprompt@gap>| !gapinput@g:=PartialPerm( [ 1, 2, 3, 4, 5, 8, 10 ],|
  !gapprompt@>| !gapinput@[ 7, 1, 4, 3, 2, 6, 5 ] );;|
  !gapprompt@gap>| !gapinput@NaturalLeqPartialPerm(f, g);|
  false
  !gapprompt@gap>| !gapinput@NaturalLeqPartialPerm(g, f);|
  false
\end{Verbatim}
 }

  

\subsection{\textcolor{Chapter }{ShortLexLeqPartialPerm}}
\logpage{[ 54, 5, 14 ]}\nobreak
\hyperdef{L}{X81BD69307D294A1C}{}
{\noindent\textcolor{FuncColor}{$\triangleright$\enspace\texttt{ShortLexLeqPartialPerm({\mdseries\slshape f, g})\index{ShortLexLeqPartialPerm@\texttt{ShortLexLeqPartialPerm}}
\label{ShortLexLeqPartialPerm}
}\hfill{\scriptsize (function)}}\\
\textbf{\indent Returns:\ }
\texttt{true} or \texttt{false}.



 \texttt{ShortLexLeqPartialPerm} returns \texttt{true} if the concatenation of the domain and image list of \mbox{\texttt{\mdseries\slshape f}} is short\texttt{\symbol{45}}lex less than the corresponding concatenation for \mbox{\texttt{\mdseries\slshape g}} and \texttt{false} otherwise.

 Note that this is not the natural partial order on partial permutation or the
same as comparing \mbox{\texttt{\mdseries\slshape f}} and \mbox{\texttt{\mdseries\slshape g}} using \texttt{\texttt{\symbol{92}}{\textless}}. 
\begin{Verbatim}[commandchars=!@|,fontsize=\small,frame=single,label=Example]
  !gapprompt@gap>| !gapinput@f:=PartialPerm( [ 1, 2, 3, 4, 6, 7, 8, 10 ],|
  !gapprompt@>| !gapinput@[ 3, 8, 1, 9, 4, 10, 5, 6 ] );|
  [2,8,5][7,10,6,4,9](1,3)
  !gapprompt@gap>| !gapinput@g:=PartialPerm( [ 1, 2, 3, 4, 5, 8, 10 ],|
  !gapprompt@>| !gapinput@[ 7, 1, 4, 3, 2, 6, 5 ] );|
  [8,6][10,5,2,1,7](3,4)
  !gapprompt@gap>| !gapinput@f<g;|
  true
  !gapprompt@gap>| !gapinput@g<f;|
  false
  !gapprompt@gap>| !gapinput@ShortLexLeqPartialPerm(f, g);|
  false
  !gapprompt@gap>| !gapinput@ShortLexLeqPartialPerm(g, f);|
  true
  !gapprompt@gap>| !gapinput@NaturalLeqPartialPerm(f, g);|
  false
  !gapprompt@gap>| !gapinput@NaturalLeqPartialPerm(g, f);|
  false
\end{Verbatim}
 }

 

\subsection{\textcolor{Chapter }{TrimPartialPerm}}
\logpage{[ 54, 5, 15 ]}\nobreak
\hyperdef{L}{X83560BE678ACF855}{}
{\noindent\textcolor{FuncColor}{$\triangleright$\enspace\texttt{TrimPartialPerm({\mdseries\slshape f})\index{TrimPartialPerm@\texttt{TrimPartialPerm}}
\label{TrimPartialPerm}
}\hfill{\scriptsize (operation)}}\\
\textbf{\indent Returns:\ }
Nothing.



 It can happen that the internal representation of a partial permutation uses
more memory than necessary. For example, by composing a partial permutation
with codegree less than 65536 with a partial permutation with codegree greater
than 65535. It is possible that the resulting partial permutation \mbox{\texttt{\mdseries\slshape f}} has its codegree and images stored as 32\texttt{\symbol{45}}bit integers,
while none of its image points exceeds 65536. The purpose of this function is
to change the internal representation of such an \mbox{\texttt{\mdseries\slshape f}} from using 32\texttt{\symbol{45}}bit to using 16\texttt{\symbol{45}}bit
integers. 

 Note that the partial permutation \mbox{\texttt{\mdseries\slshape f}} is changed in\texttt{\symbol{45}}place, and nothing is returned by this
function. 
\begin{Verbatim}[commandchars=!@|,fontsize=\small,frame=single,label=Example]
  !gapprompt@gap>| !gapinput@f:=PartialPerm( [ 1, 2 ], [ 3, 4 ] )|
  !gapprompt@>| !gapinput@*PartialPerm( [ 3, 5 ], [ 3, 100000 ] );|
  [1,3]
  !gapprompt@gap>| !gapinput@IsPPerm4Rep(f);|
  true
  !gapprompt@gap>| !gapinput@TrimPartialPerm(f); f;|
  [1,3]
  !gapprompt@gap>| !gapinput@IsPPerm4Rep(f);|
  false
\end{Verbatim}
 }

  }

  
\section{\textcolor{Chapter }{Displaying partial permutations}}\label{sect:DisplayingPartialPerms}
\logpage{[ 54, 6, 0 ]}
\hyperdef{L}{X7849595B81D063EE}{}
{
  It is possible to change the way that \textsf{GAP} displays partial permutations using the user preferences \texttt{PartialPermDisplayLimit} and \texttt{NotationForPartialPerms}; see Section \texttt{UserPreference} (\ref{UserPreference}) for more information about user preferences. 

 If \texttt{f} is a partial permutation of rank \texttt{r} exceeding the value of the user preference \texttt{PartialPermDisplayLimit}, then \texttt{f} is displayed as: 
\begin{Verbatim}[commandchars=!@|,fontsize=\small,frame=single,label=Example]
  <partial perm on r pts with degree m, codegree n>
\end{Verbatim}
 where the degree and codegree are \texttt{m} and \texttt{n}, respectively. The idea is to abbreviate the display of partial permutations
defined on many points. The default value for the \texttt{PartialPermDisplayLimit} is \texttt{100}. 

 If the rank of \texttt{f} does not exceed the value of \texttt{PartialPermDisplayLimit}, then how \texttt{f} is displayed depends on the value of the user preference \texttt{NotationForPartialPerms} except in the case that \texttt{f} is the empty partial permutation or an identity partial permutation. 

 There are three possible values for \texttt{NotationForPartialPerms} user preference, which are described below. 
\begin{description}
\item[{component}]  Similar to permutations, and unlike transformations, partial permutations can
be expressed as products of disjoint permutations and chains. A \emph{chain} is a list \texttt{c} of some length \texttt{n} such that: 
\begin{itemize}
\item  \texttt{c[1]} is an element of the domain of \mbox{\texttt{\mdseries\slshape f}} but not the image 
\item \texttt{c[i]\texttt{\symbol{94}}\mbox{\texttt{\mdseries\slshape f}}=c[i+1]} for all \texttt{i} in the range from \texttt{1} to \texttt{n\texttt{\symbol{45}}1}. 
\item \texttt{c[n]} is in the image of \mbox{\texttt{\mdseries\slshape f}} but not the domain.
\end{itemize}
 In the display, permutations are displayed as they usually are in \textsf{GAP}, except that fixed points are displayed enclosed in round brackets, and
chains are displayed enclosed in square brackets. 
\begin{Verbatim}[commandchars=!@|,fontsize=\small,frame=single,label=Example]
  !gapprompt@gap>| !gapinput@f := PartialPerm([ 1, 2, 3, 4, 5, 6, 7, 8, 10, 11, 12, 16, 17, 18, 19 ],|
  !gapprompt@>| !gapinput@[ 3, 12, 14, 4, 11, 18, 17, 2, 9, 5, 15, 8, 20, 10, 19 ]);|
  [1,3,14][16,8,2,12,15](4)(5,11)[6,18,10,9][7,17,20](19)
\end{Verbatim}
 This option is the most compact way to display a partial permutation and is
the default value of the user preference \texttt{NotationForPartialPerms}. 
\item[{domainimage}]  With this option a partial permutation \texttt{f} is displayed in the format: \texttt{DomainOfPartialPerm(\mbox{\texttt{\mdseries\slshape f}})\texttt{\symbol{45}}{\textgreater} ImageListOfPartialPerm(\mbox{\texttt{\mdseries\slshape f}})}. 
\begin{Verbatim}[commandchars=!@|,fontsize=\small,frame=single,label=Example]
  !gapprompt@gap>| !gapinput@f:=PartialPerm( [ 1, 2, 4, 5, 6, 7 ], [ 10, 1, 6, 5, 8, 7 ]);|
  [ 1, 2, 4, 5, 6, 7 ] -> [ 10, 1, 6, 5, 8, 7 ]
\end{Verbatim}
 
\item[{input}]  With this option a partial permutation \mbox{\texttt{\mdseries\slshape f}} is displayed as: \texttt{PartialPerm(DomainOfPartialPerm(\mbox{\texttt{\mdseries\slshape f}}), ImageListOfPartialPerm(\mbox{\texttt{\mdseries\slshape f}}))} which corresponds to the input (of the first type described in \texttt{PartialPerm} (\ref{PartialPerm:for a domain and image})). 
\begin{Verbatim}[commandchars=!@|,fontsize=\small,frame=single,label=Example]
  !gapprompt@gap>| !gapinput@f:=PartialPerm( [ 1, 2, 3, 5, 6, 9, 10 ],|
  !gapprompt@>| !gapinput@[ 4, 7, 3, 8, 2, 1, 6 ] );|
  PartialPerm( [ 1, 2, 3, 5, 6, 9, 10 ], [ 4, 7, 3, 8, 2, 1, 6 ] )
\end{Verbatim}
 
\end{description}
 
\begin{Verbatim}[commandchars=!@|,fontsize=\small,frame=single,label=Example]
  !gapprompt@gap>| !gapinput@SetUserPreference("PartialPermDisplayLimit", 12);|
  !gapprompt@gap>| !gapinput@UserPreference("PartialPermDisplayLimit");|
  12
  !gapprompt@gap>| !gapinput@f:=PartialPerm([1,2,3,4,5,6], [6,7,1,4,3,2]);|
  [5,3,1,6,2,7](4)
  !gapprompt@gap>| !gapinput@f:=PartialPerm(|
  !gapprompt@>| !gapinput@[ 1, 2, 3, 4, 5, 6, 7, 8, 10, 11, 12, 16, 17, 18, 19 ],|
  !gapprompt@>| !gapinput@[ 3, 12, 14, 4, 11, 18, 17, 2, 9, 5, 15, 8, 20, 10, 19 ] );|
  <partial perm on 15 pts with degree 19, codegree 20>
  !gapprompt@gap>| !gapinput@SetUserPreference("PartialPermDisplayLimit", 100);|
  !gapprompt@gap>| !gapinput@f;|
  [1,3,14][6,18,10,9][7,17,20][16,8,2,12,15](4)(5,11)(19)
  !gapprompt@gap>| !gapinput@UserPreference("NotationForPartialPerms");|
  "component"
  !gapprompt@gap>| !gapinput@SetUserPreference("NotationForPartialPerms", "domainimage");|
  !gapprompt@gap>| !gapinput@f;|
  [ 1, 2, 3, 4, 5, 6, 7, 8, 10, 11, 12, 16, 17, 18, 19 ] ->
  [ 3, 12, 14, 4, 11, 18, 17, 2, 9, 5, 15, 8, 20, 10, 19 ]
  !gapprompt@gap>| !gapinput@SetUserPreference("NotationForPartialPerms", "input");|
  !gapprompt@gap>| !gapinput@f;|
  PartialPerm(
  [ 1, 2, 3, 4, 5, 6, 7, 8, 10, 11, 12, 16, 17, 18, 19 ],
  [ 3, 12, 14, 4, 11, 18, 17, 2, 9, 5, 15, 8, 20, 10, 19 ] )
\end{Verbatim}
 }

 
\section{\textcolor{Chapter }{Semigroups and inverse semigroups of partial permutations}}\label{Partial perm semigroups}
\logpage{[ 54, 7, 0 ]}
\hyperdef{L}{X7CCC82E07A73EB55}{}
{
  As mentioned at the start of the chapter, every inverse semigroup is
isomorphic to a semigroup of partial permutations, and in this section we
describe the functions in \textsf{GAP} specific to partial permutation semigroups. For more information about
semigroups and inverse semigroups see Chapter \ref{Semigroups}. 

 The \textsf{Semigroups} package contains many additional functions and methods for computing with
semigroups of partial permutations. In particular, \textsf{Semigroups} contains more efficient methods than those available in the \textsf{GAP} library (and in many cases more efficient than any other software) for
creating semigroups of transformations, calculating their Green's classes,
size, elements, group of units, minimal ideal, small generating sets, testing
membership, finding the inverses of a regular element, factorizing elements
over the generators, and more. 

 Since a partial permutation semigroup is also a partial permutation
collection, there are special methods for \texttt{DomainOfPartialPermCollection} (\ref{DomainOfPartialPermCollection}), \texttt{ImageOfPartialPermCollection} (\ref{ImageOfPartialPermCollection}), \texttt{FixedPointsOfPartialPerm} (\ref{FixedPointsOfPartialPerm:for a partial perm coll}), \texttt{MovedPoints} (\ref{MovedPoints:for a partial perm coll}), \texttt{NrFixedPoints} (\ref{NrFixedPoints:for a partial perm coll}), \texttt{NrMovedPoints} (\ref{NrMovedPoints:for a partial perm coll}), \texttt{LargestMovedPoint} (\ref{LargestMovedPoint:for a partial perm coll}), and \texttt{SmallestMovedPoint} (\ref{SmallestMovedPoint:for a partial perm coll}) when applied to a partial permutation semigroup. 

\subsection{\textcolor{Chapter }{IsPartialPermSemigroup}}
\logpage{[ 54, 7, 1 ]}\nobreak
\hyperdef{L}{X7D161674800B50E0}{}
{\noindent\textcolor{FuncColor}{$\triangleright$\enspace\texttt{IsPartialPermSemigroup({\mdseries\slshape obj})\index{IsPartialPermSemigroup@\texttt{IsPartialPermSemigroup}}
\label{IsPartialPermSemigroup}
}\hfill{\scriptsize (filter)}}\\
\noindent\textcolor{FuncColor}{$\triangleright$\enspace\texttt{IsPartialPermMonoid({\mdseries\slshape obj})\index{IsPartialPermMonoid@\texttt{IsPartialPermMonoid}}
\label{IsPartialPermMonoid}
}\hfill{\scriptsize (filter)}}\\
\textbf{\indent Returns:\ }
\texttt{true} or \texttt{false}.



 A \emph{partial perm semigroup} is simply a semigroup consisting of partial permutations, which may or may not
be an inverse semigroup. An object \mbox{\texttt{\mdseries\slshape obj}} in \textsf{GAP} is a partial perm semigroup if and only if it satisfies \texttt{IsSemigroup} (\ref{IsSemigroup}) and \texttt{IsPartialPermCollection} (\ref{IsPartialPermCollection}).

 A \emph{partial perm monoid} is a monoid consisting of partial permutations. An object in \textsf{GAP} is a partial perm monoid if it satisfies \texttt{IsMonoid} (\ref{IsMonoid}) and \texttt{IsPartialPermCollection} (\ref{IsPartialPermCollection}).

 Note that it is possible for a partial perm semigroup to have a multiplicative
neutral element (i.e. an identity element) but not to satisfy \texttt{IsPartialPermMonoid}. For example, 
\begin{Verbatim}[commandchars=!@|,fontsize=\small,frame=single,label=Example]
  !gapprompt@gap>| !gapinput@f := PartialPerm( [ 1, 2, 3, 6, 8, 10 ], [ 2, 6, 7, 9, 1, 5 ] );;|
  !gapprompt@gap>| !gapinput@S := Semigroup(f, One(f));|
  <commutative partial perm monoid of rank 9 with 1 generator>
  !gapprompt@gap>| !gapinput@IsMonoid(S);|
  true
  !gapprompt@gap>| !gapinput@IsPartialPermMonoid(S);|
  true
\end{Verbatim}
 Note that unlike transformation semigroups, the \texttt{One} (\ref{One}) of a partial permutation semigroup must coincide with the multiplicative
neutral element, if either exists.

 For more details see \texttt{IsMagmaWithOne} (\ref{IsMagmaWithOne}). }

 

\subsection{\textcolor{Chapter }{DegreeOfPartialPermSemigroup}}
\logpage{[ 54, 7, 2 ]}\nobreak
\hyperdef{L}{X7D7F0BAB82F0D820}{}
{\noindent\textcolor{FuncColor}{$\triangleright$\enspace\texttt{DegreeOfPartialPermSemigroup({\mdseries\slshape S})\index{DegreeOfPartialPermSemigroup@\texttt{DegreeOfPartialPermSemigroup}}
\label{DegreeOfPartialPermSemigroup}
}\hfill{\scriptsize (attribute)}}\\
\noindent\textcolor{FuncColor}{$\triangleright$\enspace\texttt{CodegreeOfPartialPermSemigroup({\mdseries\slshape S})\index{CodegreeOfPartialPermSemigroup@\texttt{CodegreeOfPartialPermSemigroup}}
\label{CodegreeOfPartialPermSemigroup}
}\hfill{\scriptsize (attribute)}}\\
\noindent\textcolor{FuncColor}{$\triangleright$\enspace\texttt{RankOfPartialPermSemigroup({\mdseries\slshape S})\index{RankOfPartialPermSemigroup@\texttt{RankOfPartialPermSemigroup}}
\label{RankOfPartialPermSemigroup}
}\hfill{\scriptsize (attribute)}}\\
\textbf{\indent Returns:\ }
A non\texttt{\symbol{45}}negative integer.



 The \emph{degree} of a partial permutation semigroup \mbox{\texttt{\mdseries\slshape S}} is the largest degree of any partial permutation in \mbox{\texttt{\mdseries\slshape S}}. 

 The \emph{codegree} of a partial permutation semigroup \mbox{\texttt{\mdseries\slshape S}} is the largest positive integer in its image.

 The \emph{rank} of a partial permutation semigroup \mbox{\texttt{\mdseries\slshape S}} is the number of points on which it acts. 
\begin{Verbatim}[commandchars=!@|,fontsize=\small,frame=single,label=Example]
  !gapprompt@gap>| !gapinput@S := Semigroup( PartialPerm( [ 1, 5 ], [ 10000, 3 ] ) );|
  <commutative partial perm semigroup of rank 2 with 1 generator>
  !gapprompt@gap>| !gapinput@DegreeOfPartialPermSemigroup(S);|
  5
  !gapprompt@gap>| !gapinput@CodegreeOfPartialPermSemigroup(S);|
  10000
  !gapprompt@gap>| !gapinput@RankOfPartialPermSemigroup(S);|
  2
\end{Verbatim}
 }

 

\subsection{\textcolor{Chapter }{SymmetricInverseSemigroup}}
\logpage{[ 54, 7, 3 ]}\nobreak
\hyperdef{L}{X81D271B380995F8A}{}
{\noindent\textcolor{FuncColor}{$\triangleright$\enspace\texttt{SymmetricInverseSemigroup({\mdseries\slshape n})\index{SymmetricInverseSemigroup@\texttt{SymmetricInverseSemigroup}}
\label{SymmetricInverseSemigroup}
}\hfill{\scriptsize (operation)}}\\
\noindent\textcolor{FuncColor}{$\triangleright$\enspace\texttt{SymmetricInverseMonoid({\mdseries\slshape n})\index{SymmetricInverseMonoid@\texttt{SymmetricInverseMonoid}}
\label{SymmetricInverseMonoid}
}\hfill{\scriptsize (operation)}}\\
\textbf{\indent Returns:\ }
The symmetric inverse semigroup of degree \mbox{\texttt{\mdseries\slshape n}}.



 If \mbox{\texttt{\mdseries\slshape n}} is a non\texttt{\symbol{45}}negative integer, then \texttt{SymmetricInverseSemigroup} returns the inverse semigroup consisting of all partial permutations with
degree and codegree at most \mbox{\texttt{\mdseries\slshape n}}. Note that \mbox{\texttt{\mdseries\slshape n}} must be non\texttt{\symbol{45}}negative, but in particular, can equal \texttt{0}. 

 The symmetric inverse semigroup  has $\sum_{r=0}^n{n\choose r}^2\cdot r!$ elements and  is generated by any set that of partial permutations that generate the
symmetric group on \mbox{\texttt{\mdseries\slshape n}} points and any partial permutation of rank \texttt{\mbox{\texttt{\mdseries\slshape n}}\texttt{\symbol{45}}1}.

 \texttt{SymmetricInverseMonoid} is a synonym for \texttt{SymmetricInverseSemigroup}. 
\begin{Verbatim}[commandchars=!@|,fontsize=\small,frame=single,label=Example]
  !gapprompt@gap>| !gapinput@S := SymmetricInverseSemigroup(5);|
  <symmetric inverse monoid of degree 5>
  !gapprompt@gap>| !gapinput@Size(S);|
  1546
  !gapprompt@gap>| !gapinput@GeneratorsOfInverseMonoid(S);|
  [ (1,2,3,4,5), (1,2)(3)(4)(5), [5,4,3,2,1] ]
\end{Verbatim}
 }

 

\subsection{\textcolor{Chapter }{IsSymmetricInverseSemigroup}}
\logpage{[ 54, 7, 4 ]}\nobreak
\hyperdef{L}{X7C8AEA50834060DD}{}
{\noindent\textcolor{FuncColor}{$\triangleright$\enspace\texttt{IsSymmetricInverseSemigroup({\mdseries\slshape S})\index{IsSymmetricInverseSemigroup@\texttt{IsSymmetricInverseSemigroup}}
\label{IsSymmetricInverseSemigroup}
}\hfill{\scriptsize (property)}}\\
\noindent\textcolor{FuncColor}{$\triangleright$\enspace\texttt{IsSymmetricInverseMonoid({\mdseries\slshape S})\index{IsSymmetricInverseMonoid@\texttt{IsSymmetricInverseMonoid}}
\label{IsSymmetricInverseMonoid}
}\hfill{\scriptsize (property)}}\\
\textbf{\indent Returns:\ }
\texttt{true} or \texttt{false}.



 If the partial perm semigroup \mbox{\texttt{\mdseries\slshape S}} of degree and codegree \mbox{\texttt{\mdseries\slshape n}} equals the symmetric inverse semigroup on \mbox{\texttt{\mdseries\slshape n}} points, then \texttt{IsSymmetricInverseSemigroup} return \texttt{true} and otherwise it returns \texttt{false}.

 \texttt{IsSymmetricInverseMonoid} is a synonym of \texttt{IsSymmetricInverseSemigroup}. It is common in the literature for the symmetric inverse monoid to be
referred to as the symmetric inverse semigroup. 
\begin{Verbatim}[commandchars=!@|,fontsize=\small,frame=single,label=Example]
  !gapprompt@gap>| !gapinput@S := Semigroup(AsPartialPerm((1, 3, 4, 2), 5), AsPartialPerm((1, 3, 5), 5),|
  !gapprompt@>| !gapinput@PartialPerm( [ 1, 2, 3, 4 ] ) );|
  <partial perm semigroup of rank 5 with 3 generators>
  !gapprompt@gap>| !gapinput@IsSymmetricInverseSemigroup(S);|
  true
  !gapprompt@gap>| !gapinput@S;|
  <symmetric inverse monoid of degree 5>
\end{Verbatim}
 }

 

\subsection{\textcolor{Chapter }{NaturalPartialOrder}}
\logpage{[ 54, 7, 5 ]}\nobreak
\hyperdef{L}{X7EA51F087CF7621F}{}
{\noindent\textcolor{FuncColor}{$\triangleright$\enspace\texttt{NaturalPartialOrder({\mdseries\slshape S})\index{NaturalPartialOrder@\texttt{NaturalPartialOrder}}
\label{NaturalPartialOrder}
}\hfill{\scriptsize (attribute)}}\\
\noindent\textcolor{FuncColor}{$\triangleright$\enspace\texttt{ReverseNaturalPartialOrder({\mdseries\slshape S})\index{ReverseNaturalPartialOrder@\texttt{ReverseNaturalPartialOrder}}
\label{ReverseNaturalPartialOrder}
}\hfill{\scriptsize (attribute)}}\\
\textbf{\indent Returns:\ }
The natural partial order on an inverse semigroup.



 The \emph{natural partial order} $\leq$ on an inverse semigroup \mbox{\texttt{\mdseries\slshape S}} is defined by \texttt{s}$\leq$\texttt{t} if there exists an idempotent \texttt{e} in \mbox{\texttt{\mdseries\slshape S}} such that \texttt{s=et}. Hence if \texttt{f} and \texttt{g} are partial permutations, then \texttt{f}$\leq$\texttt{g} if and only if \texttt{f} is a restriction of \texttt{g}; see \texttt{RestrictedPartialPerm} (\ref{RestrictedPartialPerm}).

 \texttt{NaturalPartialOrder} returns the natural partial order on the inverse semigroup of partial
permutations \mbox{\texttt{\mdseries\slshape S}} as a list of sets of positive integers where entry \texttt{i} in \texttt{NaturalPartialOrder(\mbox{\texttt{\mdseries\slshape S}})} is the set of positions in \texttt{Elements(\mbox{\texttt{\mdseries\slshape S}})} of elements which are less than \texttt{Elements(\mbox{\texttt{\mdseries\slshape S}})[i]}. See also \texttt{NaturalLeqPartialPerm} (\ref{NaturalLeqPartialPerm}).

 \texttt{ReverseNaturalPartialOrder} returns the reverse of the natural partial order on the inverse semigroup of
partial permutations \mbox{\texttt{\mdseries\slshape S}} as a list of sets of positive integers where entry \texttt{i} in \texttt{ReverseNaturalPartialOrder(\mbox{\texttt{\mdseries\slshape S}})} is the set of positions in \texttt{Elements(\mbox{\texttt{\mdseries\slshape S}})} of elements which are greater than \texttt{Elements(\mbox{\texttt{\mdseries\slshape S}})[i]}. See also \texttt{NaturalLeqPartialPerm} (\ref{NaturalLeqPartialPerm}). 
\begin{Verbatim}[commandchars=!@|,fontsize=\small,frame=single,label=Example]
  !gapprompt@gap>| !gapinput@S := InverseSemigroup([ PartialPerm( [ 1, 3 ], [ 1, 3 ] ),|
  !gapprompt@>| !gapinput@PartialPerm( [ 1, 2 ], [ 3, 2 ] ) ] );|
  <inverse partial perm semigroup of rank 3 with 2 generators>
  !gapprompt@gap>| !gapinput@Size(S);|
  11
  !gapprompt@gap>| !gapinput@NaturalPartialOrder(S);|
  [ [  ], [ 1 ], [ 1 ], [ 1 ], [ 1, 2, 4 ], [ 1, 3, 4 ], [ 1 ], [ 1 ],
    [ 1, 4, 7 ], [ 1, 4, 8 ], [ 1, 2, 8 ] ]
  !gapprompt@gap>| !gapinput@NaturalLeqPartialPerm(Elements(S)[4], Elements(S)[10]);|
  true
  !gapprompt@gap>| !gapinput@NaturalLeqPartialPerm(Elements(S)[4], Elements(S)[1]);|
  false
\end{Verbatim}
 }

 

\subsection{\textcolor{Chapter }{IsomorphismPartialPermSemigroup}}
\logpage{[ 54, 7, 6 ]}\nobreak
\hyperdef{L}{X7FE18EBE79B9C17C}{}
{\noindent\textcolor{FuncColor}{$\triangleright$\enspace\texttt{IsomorphismPartialPermSemigroup({\mdseries\slshape S})\index{IsomorphismPartialPermSemigroup@\texttt{IsomorphismPartialPermSemigroup}}
\label{IsomorphismPartialPermSemigroup}
}\hfill{\scriptsize (attribute)}}\\
\noindent\textcolor{FuncColor}{$\triangleright$\enspace\texttt{IsomorphismPartialPermMonoid({\mdseries\slshape S})\index{IsomorphismPartialPermMonoid@\texttt{IsomorphismPartialPermMonoid}}
\label{IsomorphismPartialPermMonoid}
}\hfill{\scriptsize (attribute)}}\\
\textbf{\indent Returns:\ }
An isomorphism.



 \texttt{IsomorphismPartialPermSemigroup(\mbox{\texttt{\mdseries\slshape S}})} returns an isomorphism from the inverse semigroup \mbox{\texttt{\mdseries\slshape S}} to an inverse semigroup of partial permutations.

 \texttt{IsomorphismPartialPermMonoid(\mbox{\texttt{\mdseries\slshape S}})} returns an isomorphism from the inverse semigroup \mbox{\texttt{\mdseries\slshape S}} to an inverse monoid of partial permutations, if possible.

 We only describe \texttt{IsomorphismPartialPermMonoid}, the corresponding statements for \texttt{IsomorphismPartialPermSemigroup} also hold. 
\begin{description}
\item[{Partial permutation semigroups}]  If \mbox{\texttt{\mdseries\slshape S}} is a partial permutation semigroup that does not satisfy \texttt{IsMonoid} (\ref{IsMonoid}) but where \texttt{MultiplicativeNeutralElement(\mbox{\texttt{\mdseries\slshape S}}){\textless}{\textgreater}fail}, then \texttt{IsomorphismPartialPermMonoid(\mbox{\texttt{\mdseries\slshape S}})} returns an isomorphism from \mbox{\texttt{\mdseries\slshape S}} to an inverse monoid of partial permutations. 
\item[{Permutation groups}]  If \mbox{\texttt{\mdseries\slshape S}} is a permutation group, then \texttt{IsomorphismPartialPermMonoid} returns an isomorphism from \mbox{\texttt{\mdseries\slshape S}} to an inverse monoid of partial permutations on the set \texttt{MovedPoints(\mbox{\texttt{\mdseries\slshape S}})} obtained using \texttt{AsPartialPerm} (\ref{AsPartialPerm:for a permutation}). The inverse of this isomorphism is obtained using \texttt{AsPermutation} (\ref{AsPermutation}). 
\item[{Transformation semigroups}]  If \mbox{\texttt{\mdseries\slshape S}} is a transformation semigroup which is mathematically a monoid but which does
not necessarily belong to the category \texttt{IsMonoid} (\ref{IsMonoid}), then \texttt{IsomorphismPartialPermMonoid} returns an isomorphism from \mbox{\texttt{\mdseries\slshape S}} to an inverse monoid of partial permutations. 
\end{description}
 
\begin{Verbatim}[commandchars=!@|,fontsize=\small,frame=single,label=Example]
  !gapprompt@gap>| !gapinput@S := InverseSemigroup(|
  !gapprompt@>| !gapinput@PartialPerm( [ 1, 2, 3, 4, 5 ], [ 4, 2, 3, 1, 5 ] ),|
  !gapprompt@>| !gapinput@PartialPerm( [ 1, 2, 4, 5 ], [ 3, 1, 4, 2 ] ) );;|
  !gapprompt@gap>| !gapinput@IsMonoid(S);|
  false
  !gapprompt@gap>| !gapinput@Size(S);|
  508
  !gapprompt@gap>| !gapinput@iso := IsomorphismPartialPermMonoid(S);|
  MappingByFunction( <inverse partial perm semigroup of size 508,
   rank 5 with 2 generators>, <inverse partial perm monoid of size 508,
   rank 5 with 2 generators>
   , function( object ) ... end, function( object ) ... end )
  !gapprompt@gap>| !gapinput@Size(S);|
  508
  !gapprompt@gap>| !gapinput@Size(Range(iso));|
  508
  !gapprompt@gap>| !gapinput@G := Group((1,2)(3,8)(4,6)(5,7), (1,3,4,7)(2,5,6,8), (1,4)(2,6)(3,7)(5,8));;|
  !gapprompt@gap>| !gapinput@IsomorphismPartialPermSemigroup(G);|
  MappingByFunction( Group([ (1,2)(3,8)(4,6)(5,7), (1,3,4,7)
  (2,5,6,8), (1,4)(2,6)(3,7)
  (5,8) ]), <partial perm group of rank 8 with 3 generators>
  , function( p ) ... end, <Attribute "AsPermutation"> )
  !gapprompt@gap>| !gapinput@S := Semigroup(Transformation( [ 2, 5, 1, 7, 3, 7, 7 ] ),|
  !gapprompt@>| !gapinput@Transformation( [ 3, 6, 5, 7, 2, 1, 7 ] ) );;|
  !gapprompt@gap>| !gapinput@iso := IsomorphismPartialPermMonoid(S);;|
  !gapprompt@gap>| !gapinput@MultiplicativeNeutralElement(S) ^ iso;|
  <identity partial perm on [ 1, 2, 3, 4, 5, 6 ]>
  !gapprompt@gap>| !gapinput@One(Range(iso));|
  <identity partial perm on [ 1, 2, 3, 4, 5, 6 ]>
  !gapprompt@gap>| !gapinput@MovedPoints(Range(iso));|
  [ 1 .. 5 ]
\end{Verbatim}
 }

 }

 }

         
\chapter{\textcolor{Chapter }{Additive Magmas}}\label{Additive Magmas}
\logpage{[ 55, 0, 0 ]}
\hyperdef{L}{X7D0D096B81365B02}{}
{
  This chapter deals with domains that are closed under addition \texttt{+}, which are called \emph{near\texttt{\symbol{45}}additive magmas} in \textsf{GAP}. Together with the domains closed under multiplication \texttt{*} (see{\nobreakspace}\ref{Magmas}), they are the basic algebraic structures. In many cases, the addition is
commutative (see{\nobreakspace}\texttt{IsAdditivelyCommutative} (\ref{IsAdditivelyCommutative})), the domain is called an \emph{additive magma} then. Every module (see{\nobreakspace}\ref{Modules}), vector space (see{\nobreakspace}\ref{Vector Spaces}), ring (see{\nobreakspace}\ref{Rings}), or field (see{\nobreakspace}\ref{Fields and Division Rings}) is an additive magma. In the cases of all \emph{(near\texttt{\symbol{45}})additive
magma\texttt{\symbol{45}}with\texttt{\symbol{45}}zero} or \emph{(near\texttt{\symbol{45}})additive
magma\texttt{\symbol{45}}with\texttt{\symbol{45}}inverses}, additional additive structure is present (see{\nobreakspace}\ref{(Near-)Additive Magma Categories}).  
\section{\textcolor{Chapter }{(Near\texttt{\symbol{45}})Additive Magma Categories}}\label{(Near-)Additive Magma Categories}
\logpage{[ 55, 1, 0 ]}
\hyperdef{L}{X82A4AB7B812B063B}{}
{
  

\subsection{\textcolor{Chapter }{IsNearAdditiveMagma}}
\logpage{[ 55, 1, 1 ]}\nobreak
\hyperdef{L}{X8129E95D83227658}{}
{\noindent\textcolor{FuncColor}{$\triangleright$\enspace\texttt{IsNearAdditiveMagma({\mdseries\slshape obj})\index{IsNearAdditiveMagma@\texttt{IsNearAdditiveMagma}}
\label{IsNearAdditiveMagma}
}\hfill{\scriptsize (Category)}}\\


 A \emph{near\texttt{\symbol{45}}additive magma} in \textsf{GAP} is a domain $A$ with an associative but not necessarily commutative addition \texttt{+}$: A \times A \rightarrow A$. }

 

\subsection{\textcolor{Chapter }{IsNearAdditiveMagmaWithZero}}
\logpage{[ 55, 1, 2 ]}\nobreak
\hyperdef{L}{X7DADE4577D0A7208}{}
{\noindent\textcolor{FuncColor}{$\triangleright$\enspace\texttt{IsNearAdditiveMagmaWithZero({\mdseries\slshape obj})\index{IsNearAdditiveMagmaWithZero@\texttt{IsNearAdditiveMagmaWithZero}}
\label{IsNearAdditiveMagmaWithZero}
}\hfill{\scriptsize (Category)}}\\


 A \emph{near\texttt{\symbol{45}}additive
magma\texttt{\symbol{45}}with\texttt{\symbol{45}}zero} in \textsf{GAP} is a near\texttt{\symbol{45}}additive magma $A$ with an operation \texttt{0*} (or \texttt{Zero} (\ref{Zero})) that yields the zero element of $A$. 

 So a near\texttt{\symbol{45}}additive
magma\texttt{\symbol{45}}with\texttt{\symbol{45}}zero \mbox{\texttt{\mdseries\slshape A}} does always contain a unique additively neutral element $z$, i.e., $z + a = a = a + z$ holds for all $a \in A$ (see{\nobreakspace}\texttt{AdditiveNeutralElement} (\ref{AdditiveNeutralElement})). This zero element $z$ can be computed with the operation \texttt{Zero} (\ref{Zero}), by applying this function to $A$ or to any element $a$ in $A$. The zero element can be computed also as \texttt{0 * }$a$, for any $a$ in $A$. 

 \emph{Note} that it may happen that a near\texttt{\symbol{45}}additive magma containing a
zero does \emph{not} lie in the category \texttt{IsNearAdditiveMagmaWithZero} (see{\nobreakspace}\ref{Domain Categories}). }

 

\subsection{\textcolor{Chapter }{IsNearAdditiveGroup}}
\logpage{[ 55, 1, 3 ]}\nobreak
\hyperdef{L}{X7FC3A9C178185942}{}
{\noindent\textcolor{FuncColor}{$\triangleright$\enspace\texttt{IsNearAdditiveGroup({\mdseries\slshape obj})\index{IsNearAdditiveGroup@\texttt{IsNearAdditiveGroup}}
\label{IsNearAdditiveGroup}
}\hfill{\scriptsize (Category)}}\\
\noindent\textcolor{FuncColor}{$\triangleright$\enspace\texttt{IsNearAdditiveMagmaWithInverses({\mdseries\slshape obj})\index{IsNearAdditiveMagmaWithInverses@\texttt{IsNearAdditiveMagmaWithInverses}}
\label{IsNearAdditiveMagmaWithInverses}
}\hfill{\scriptsize (Category)}}\\


 A \emph{near\texttt{\symbol{45}}additive group} in \textsf{GAP} is a near\texttt{\symbol{45}}additive
magma\texttt{\symbol{45}}with\texttt{\symbol{45}}zero $A$ with an operation \texttt{\texttt{\symbol{45}}1*}$: A \rightarrow A$ that maps each element $a$ of $A$ to its additive inverse \texttt{\texttt{\symbol{45}}1*}$a$ (or \texttt{AdditiveInverse( }\mbox{\texttt{\mdseries\slshape a}}\texttt{ )}, see{\nobreakspace}\texttt{AdditiveInverse} (\ref{AdditiveInverse})). 

 The addition \texttt{+} of $A$ is assumed to be associative, so a near\texttt{\symbol{45}}additive group is
not more than a \emph{near\texttt{\symbol{45}}additive
magma\texttt{\symbol{45}}with\texttt{\symbol{45}}inverses}. \texttt{IsNearAdditiveMagmaWithInverses} is just a synonym for \texttt{IsNearAdditiveGroup}, and can be used alternatively in all function names involving the string \texttt{"NearAdditiveGroup"}. 

 Note that not every trivial near\texttt{\symbol{45}}additive magma is a
near\texttt{\symbol{45}}additive
magma\texttt{\symbol{45}}with\texttt{\symbol{45}}zero, but every trivial
near\texttt{\symbol{45}}additive
magma\texttt{\symbol{45}}with\texttt{\symbol{45}}zero is a
near\texttt{\symbol{45}}additive group. }

 

\subsection{\textcolor{Chapter }{IsAdditiveMagma}}
\logpage{[ 55, 1, 4 ]}\nobreak
\hyperdef{L}{X8565FD0C847BAA3A}{}
{\noindent\textcolor{FuncColor}{$\triangleright$\enspace\texttt{IsAdditiveMagma({\mdseries\slshape obj})\index{IsAdditiveMagma@\texttt{IsAdditiveMagma}}
\label{IsAdditiveMagma}
}\hfill{\scriptsize (Category)}}\\


 An \emph{additive magma} in \textsf{GAP} is a domain $A$ with an associative and commutative addition \texttt{+}$: A \times A \rightarrow A$, see{\nobreakspace}\texttt{IsNearAdditiveMagma} (\ref{IsNearAdditiveMagma}) and \texttt{IsAdditivelyCommutative} (\ref{IsAdditivelyCommutative}). }

 

\subsection{\textcolor{Chapter }{IsAdditiveMagmaWithZero}}
\logpage{[ 55, 1, 5 ]}\nobreak
\hyperdef{L}{X785B41A67D791783}{}
{\noindent\textcolor{FuncColor}{$\triangleright$\enspace\texttt{IsAdditiveMagmaWithZero({\mdseries\slshape obj})\index{IsAdditiveMagmaWithZero@\texttt{IsAdditiveMagmaWithZero}}
\label{IsAdditiveMagmaWithZero}
}\hfill{\scriptsize (Category)}}\\


 An \emph{additive magma\texttt{\symbol{45}}with\texttt{\symbol{45}}zero} in \textsf{GAP} is an additive magma $A$ (see \texttt{IsAdditiveMagma} (\ref{IsAdditiveMagma}) with an operation \texttt{0*} (or \texttt{Zero} (\ref{Zero})) that yields the zero of $A$. 

 So an additive magma\texttt{\symbol{45}}with\texttt{\symbol{45}}zero $A$ does always contain a unique additively neutral element $z$, i.e., $z + a = a = a + z$ holds for all $a \in A$ (see{\nobreakspace}\texttt{AdditiveNeutralElement} (\ref{AdditiveNeutralElement})). This element $z$ can be computed with the operation \texttt{Zero} (\ref{Zero}) as \texttt{Zero( }$A$\texttt{ )}, and $z$ is also equal to \texttt{Zero( }$a$\texttt{ )} and to \texttt{0*}$a$ for each element $a$ in $A$. 

 \emph{Note} that it may happen that an additive magma containing a zero does \emph{not} lie in the category \texttt{IsAdditiveMagmaWithZero} (see{\nobreakspace}\ref{Domain Categories}). }

 

\subsection{\textcolor{Chapter }{IsAdditiveGroup}}
\logpage{[ 55, 1, 6 ]}\nobreak
\hyperdef{L}{X7B8FBD9082CE271B}{}
{\noindent\textcolor{FuncColor}{$\triangleright$\enspace\texttt{IsAdditiveGroup({\mdseries\slshape obj})\index{IsAdditiveGroup@\texttt{IsAdditiveGroup}}
\label{IsAdditiveGroup}
}\hfill{\scriptsize (Category)}}\\
\noindent\textcolor{FuncColor}{$\triangleright$\enspace\texttt{IsAdditiveMagmaWithInverses({\mdseries\slshape obj})\index{IsAdditiveMagmaWithInverses@\texttt{IsAdditiveMagmaWithInverses}}
\label{IsAdditiveMagmaWithInverses}
}\hfill{\scriptsize (Category)}}\\


 An \emph{additive group} in \textsf{GAP} is an additive magma\texttt{\symbol{45}}with\texttt{\symbol{45}}zero $A$ with an operation \texttt{\texttt{\symbol{45}}1*}$: A \rightarrow A$ that maps each element $a$ of $A$ to its additive inverse \texttt{\texttt{\symbol{45}}1*}$a$ (or \texttt{AdditiveInverse( }$a$\texttt{ )}, see{\nobreakspace}\texttt{AdditiveInverse} (\ref{AdditiveInverse})). 

 The addition \texttt{+} of $A$ is assumed to be commutative and associative, so an additive group is not more
than an \emph{additive magma\texttt{\symbol{45}}with\texttt{\symbol{45}}inverses}. \texttt{IsAdditiveMagmaWithInverses} is just a synonym for \texttt{IsAdditiveGroup}, and can be used alternatively in all function names involving the string \texttt{"AdditiveGroup"}. 

 Note that not every trivial additive magma is an additive
magma\texttt{\symbol{45}}with\texttt{\symbol{45}}zero, but every trivial
additive magma\texttt{\symbol{45}}with\texttt{\symbol{45}}zero is an additive
group. }

 }

  
\section{\textcolor{Chapter }{(Near\texttt{\symbol{45}})Additive Magma Generation}}\label{(Near-)Additive Magma Generation}
\logpage{[ 55, 2, 0 ]}
\hyperdef{L}{X7C39F9DE7CA22688}{}
{
  This section describes functions that create additive magmas from generators
(see \texttt{NearAdditiveMagma} (\ref{NearAdditiveMagma}), \texttt{NearAdditiveMagmaWithZero} (\ref{NearAdditiveMagmaWithZero}), \texttt{NearAdditiveGroup} (\ref{NearAdditiveGroup})), the underlying operations for which methods can be installed (see \texttt{NearAdditiveMagmaByGenerators} (\ref{NearAdditiveMagmaByGenerators}), \texttt{NearAdditiveMagmaWithZeroByGenerators} (\ref{NearAdditiveMagmaWithZeroByGenerators}), \texttt{NearAdditiveGroupByGenerators} (\ref{NearAdditiveGroupByGenerators})) and functions for forming additive submagmas (see \texttt{SubnearAdditiveMagma} (\ref{SubnearAdditiveMagma}), \texttt{SubnearAdditiveMagmaWithZero} (\ref{SubnearAdditiveMagmaWithZero}), \texttt{SubnearAdditiveGroup} (\ref{SubnearAdditiveGroup})). 

\subsection{\textcolor{Chapter }{NearAdditiveMagma}}
\logpage{[ 55, 2, 1 ]}\nobreak
\hyperdef{L}{X79C947CF8060335A}{}
{\noindent\textcolor{FuncColor}{$\triangleright$\enspace\texttt{NearAdditiveMagma({\mdseries\slshape [Fam, ]gens})\index{NearAdditiveMagma@\texttt{NearAdditiveMagma}}
\label{NearAdditiveMagma}
}\hfill{\scriptsize (function)}}\\


 returns the (near\texttt{\symbol{45}})additive magma $A$ that is generated by the elements in the list \mbox{\texttt{\mdseries\slshape gens}}, that is, the closure of \mbox{\texttt{\mdseries\slshape gens}} under addition \texttt{+}. The family \mbox{\texttt{\mdseries\slshape Fam}} of $A$ can be entered as first argument; this is obligatory if \mbox{\texttt{\mdseries\slshape gens}} is empty (and hence also $A$ is empty). }

 

\subsection{\textcolor{Chapter }{NearAdditiveMagmaWithZero}}
\logpage{[ 55, 2, 2 ]}\nobreak
\hyperdef{L}{X80F57FB47E1DB380}{}
{\noindent\textcolor{FuncColor}{$\triangleright$\enspace\texttt{NearAdditiveMagmaWithZero({\mdseries\slshape [Fam, ]gens})\index{NearAdditiveMagmaWithZero@\texttt{NearAdditiveMagmaWithZero}}
\label{NearAdditiveMagmaWithZero}
}\hfill{\scriptsize (function)}}\\


 returns the (near\texttt{\symbol{45}})additive
magma\texttt{\symbol{45}}with\texttt{\symbol{45}}zero $A$ that is generated by the elements in the list \mbox{\texttt{\mdseries\slshape gens}}, that is, the closure of \mbox{\texttt{\mdseries\slshape gens}} under addition \texttt{+} and \texttt{Zero} (\ref{Zero}). The family \mbox{\texttt{\mdseries\slshape Fam}} of $A$ can be entered as first argument; this is obligatory if \mbox{\texttt{\mdseries\slshape gens}} is empty (and hence $A$ is trivial). }

 

\subsection{\textcolor{Chapter }{NearAdditiveGroup}}
\logpage{[ 55, 2, 3 ]}\nobreak
\hyperdef{L}{X872307537ECC5755}{}
{\noindent\textcolor{FuncColor}{$\triangleright$\enspace\texttt{NearAdditiveGroup({\mdseries\slshape [Fam, ]gens})\index{NearAdditiveGroup@\texttt{NearAdditiveGroup}}
\label{NearAdditiveGroup}
}\hfill{\scriptsize (function)}}\\


 returns the (near\texttt{\symbol{45}})additive group $A$ that is generated by the elements in the list \mbox{\texttt{\mdseries\slshape gens}}, that is, the closure of \mbox{\texttt{\mdseries\slshape gens}} under addition \texttt{+}, \texttt{Zero} (\ref{Zero}), and \texttt{AdditiveInverse} (\ref{AdditiveInverse}). The family \mbox{\texttt{\mdseries\slshape Fam}} of $A$ can be entered as first argument; this is obligatory if \mbox{\texttt{\mdseries\slshape gens}} is empty (and hence $A$ is trivial). }

 

\subsection{\textcolor{Chapter }{NearAdditiveMagmaByGenerators}}
\logpage{[ 55, 2, 4 ]}\nobreak
\hyperdef{L}{X85122CFD7BDAD668}{}
{\noindent\textcolor{FuncColor}{$\triangleright$\enspace\texttt{NearAdditiveMagmaByGenerators({\mdseries\slshape [Fam, ]gens})\index{NearAdditiveMagmaByGenerators@\texttt{NearAdditiveMagmaByGenerators}}
\label{NearAdditiveMagmaByGenerators}
}\hfill{\scriptsize (operation)}}\\


 An underlying operation for \texttt{NearAdditiveMagma} (\ref{NearAdditiveMagma}). }

 

\subsection{\textcolor{Chapter }{NearAdditiveMagmaWithZeroByGenerators}}
\logpage{[ 55, 2, 5 ]}\nobreak
\hyperdef{L}{X81880460851DEFBC}{}
{\noindent\textcolor{FuncColor}{$\triangleright$\enspace\texttt{NearAdditiveMagmaWithZeroByGenerators({\mdseries\slshape [Fam, ]gens})\index{NearAdditiveMagmaWithZeroByGenerators@\texttt{Near}\-\texttt{Additive}\-\texttt{Magma}\-\texttt{With}\-\texttt{Zero}\-\texttt{By}\-\texttt{Generators}}
\label{NearAdditiveMagmaWithZeroByGenerators}
}\hfill{\scriptsize (operation)}}\\


 An underlying operation for \texttt{NearAdditiveMagmaWithZero} (\ref{NearAdditiveMagmaWithZero}). }

 

\subsection{\textcolor{Chapter }{NearAdditiveGroupByGenerators}}
\logpage{[ 55, 2, 6 ]}\nobreak
\hyperdef{L}{X85F120B68576B267}{}
{\noindent\textcolor{FuncColor}{$\triangleright$\enspace\texttt{NearAdditiveGroupByGenerators({\mdseries\slshape [Fam, ]gens})\index{NearAdditiveGroupByGenerators@\texttt{NearAdditiveGroupByGenerators}}
\label{NearAdditiveGroupByGenerators}
}\hfill{\scriptsize (operation)}}\\


 An underlying operation for \texttt{NearAdditiveGroup} (\ref{NearAdditiveGroup}). }

 

\subsection{\textcolor{Chapter }{SubnearAdditiveMagma}}
\logpage{[ 55, 2, 7 ]}\nobreak
\hyperdef{L}{X7AA6092683FC0F9C}{}
{\noindent\textcolor{FuncColor}{$\triangleright$\enspace\texttt{SubnearAdditiveMagma({\mdseries\slshape D, gens})\index{SubnearAdditiveMagma@\texttt{SubnearAdditiveMagma}}
\label{SubnearAdditiveMagma}
}\hfill{\scriptsize (function)}}\\
\noindent\textcolor{FuncColor}{$\triangleright$\enspace\texttt{SubadditiveMagma({\mdseries\slshape D, gens})\index{SubadditiveMagma@\texttt{SubadditiveMagma}}
\label{SubadditiveMagma}
}\hfill{\scriptsize (function)}}\\
\noindent\textcolor{FuncColor}{$\triangleright$\enspace\texttt{SubnearAdditiveMagmaNC({\mdseries\slshape D, gens})\index{SubnearAdditiveMagmaNC@\texttt{SubnearAdditiveMagmaNC}}
\label{SubnearAdditiveMagmaNC}
}\hfill{\scriptsize (function)}}\\
\noindent\textcolor{FuncColor}{$\triangleright$\enspace\texttt{SubadditiveMagmaNC({\mdseries\slshape D, gens})\index{SubadditiveMagmaNC@\texttt{SubadditiveMagmaNC}}
\label{SubadditiveMagmaNC}
}\hfill{\scriptsize (function)}}\\


 \texttt{SubnearAdditiveMagma} returns the near\texttt{\symbol{45}}additive magma generated by the elements
in the list \mbox{\texttt{\mdseries\slshape gens}}, with parent the domain \mbox{\texttt{\mdseries\slshape D}}. \texttt{SubnearAdditiveMagmaNC} does the same, except that it does not check whether the elements of \mbox{\texttt{\mdseries\slshape gens}} lie in \mbox{\texttt{\mdseries\slshape D}}. 

 \texttt{SubadditiveMagma} and \texttt{SubadditiveMagmaNC} are just synonyms of these functions. }

 

\subsection{\textcolor{Chapter }{SubnearAdditiveMagmaWithZero}}
\logpage{[ 55, 2, 8 ]}\nobreak
\hyperdef{L}{X784859197D89A548}{}
{\noindent\textcolor{FuncColor}{$\triangleright$\enspace\texttt{SubnearAdditiveMagmaWithZero({\mdseries\slshape D, gens})\index{SubnearAdditiveMagmaWithZero@\texttt{SubnearAdditiveMagmaWithZero}}
\label{SubnearAdditiveMagmaWithZero}
}\hfill{\scriptsize (function)}}\\
\noindent\textcolor{FuncColor}{$\triangleright$\enspace\texttt{SubadditiveMagmaWithZero({\mdseries\slshape D, gens})\index{SubadditiveMagmaWithZero@\texttt{SubadditiveMagmaWithZero}}
\label{SubadditiveMagmaWithZero}
}\hfill{\scriptsize (function)}}\\
\noindent\textcolor{FuncColor}{$\triangleright$\enspace\texttt{SubnearAdditiveMagmaWithZeroNC({\mdseries\slshape D, gens})\index{SubnearAdditiveMagmaWithZeroNC@\texttt{SubnearAdditiveMagmaWithZeroNC}}
\label{SubnearAdditiveMagmaWithZeroNC}
}\hfill{\scriptsize (function)}}\\
\noindent\textcolor{FuncColor}{$\triangleright$\enspace\texttt{SubadditiveMagmaWithZeroNC({\mdseries\slshape D, gens})\index{SubadditiveMagmaWithZeroNC@\texttt{SubadditiveMagmaWithZeroNC}}
\label{SubadditiveMagmaWithZeroNC}
}\hfill{\scriptsize (function)}}\\


 \texttt{SubnearAdditiveMagmaWithZero} returns the near\texttt{\symbol{45}}additive
magma\texttt{\symbol{45}}with\texttt{\symbol{45}}zero generated by the
elements in the list \mbox{\texttt{\mdseries\slshape gens}}, with parent the domain \mbox{\texttt{\mdseries\slshape D}}. \texttt{SubnearAdditiveMagmaWithZeroNC} does the same, except that it does not check whether the elements of \mbox{\texttt{\mdseries\slshape gens}} lie in \mbox{\texttt{\mdseries\slshape D}}. 

 \texttt{SubadditiveMagmaWithZero} and \texttt{SubadditiveMagmaWithZeroNC} are just synonyms of these functions. }

 

\subsection{\textcolor{Chapter }{SubnearAdditiveGroup}}
\logpage{[ 55, 2, 9 ]}\nobreak
\hyperdef{L}{X844C49BA807AB99F}{}
{\noindent\textcolor{FuncColor}{$\triangleright$\enspace\texttt{SubnearAdditiveGroup({\mdseries\slshape D, gens})\index{SubnearAdditiveGroup@\texttt{SubnearAdditiveGroup}}
\label{SubnearAdditiveGroup}
}\hfill{\scriptsize (function)}}\\
\noindent\textcolor{FuncColor}{$\triangleright$\enspace\texttt{SubadditiveGroup({\mdseries\slshape D, gens})\index{SubadditiveGroup@\texttt{SubadditiveGroup}}
\label{SubadditiveGroup}
}\hfill{\scriptsize (function)}}\\
\noindent\textcolor{FuncColor}{$\triangleright$\enspace\texttt{SubnearAdditiveGroupNC({\mdseries\slshape D, gens})\index{SubnearAdditiveGroupNC@\texttt{SubnearAdditiveGroupNC}}
\label{SubnearAdditiveGroupNC}
}\hfill{\scriptsize (function)}}\\
\noindent\textcolor{FuncColor}{$\triangleright$\enspace\texttt{SubadditiveGroupNC({\mdseries\slshape D, gens})\index{SubadditiveGroupNC@\texttt{SubadditiveGroupNC}}
\label{SubadditiveGroupNC}
}\hfill{\scriptsize (function)}}\\


 \texttt{SubnearAdditiveGroup} returns the near\texttt{\symbol{45}}additive group generated by the elements
in the list \mbox{\texttt{\mdseries\slshape gens}}, with parent the domain \mbox{\texttt{\mdseries\slshape D}}. \texttt{SubadditiveGroupNC} does the same, except that it does not check whether the elements of \mbox{\texttt{\mdseries\slshape gens}} lie in \mbox{\texttt{\mdseries\slshape D}}. 

 \texttt{SubadditiveGroup} and \texttt{SubadditiveGroupNC} are just synonyms of these functions. }

      }

  
\section{\textcolor{Chapter }{Attributes and Properties for (Near\texttt{\symbol{45}})Additive Magmas}}\label{Attributes and Properties for (Near-)Additive Magmas}
\logpage{[ 55, 3, 0 ]}
\hyperdef{L}{X799E6CC28737BF1B}{}
{
  

\subsection{\textcolor{Chapter }{IsAdditivelyCommutative}}
\logpage{[ 55, 3, 1 ]}\nobreak
\hyperdef{L}{X82D471327A9CA960}{}
{\noindent\textcolor{FuncColor}{$\triangleright$\enspace\texttt{IsAdditivelyCommutative({\mdseries\slshape A})\index{IsAdditivelyCommutative@\texttt{IsAdditivelyCommutative}}
\label{IsAdditivelyCommutative}
}\hfill{\scriptsize (property)}}\\


 A near\texttt{\symbol{45}}additive magma \mbox{\texttt{\mdseries\slshape A}} in \textsf{GAP} is \emph{additively commutative} if for all elements $a, b \in \mbox{\texttt{\mdseries\slshape A}}$ the equality $a + b = b + a$ holds. 

 Note that the commutativity of the \emph{multiplication} \texttt{*} in a multiplicative structure can be tested with \texttt{IsCommutative} (\ref{IsCommutative}). }

 

\subsection{\textcolor{Chapter }{GeneratorsOfNearAdditiveMagma}}
\logpage{[ 55, 3, 2 ]}\nobreak
\hyperdef{L}{X804B178884002A40}{}
{\noindent\textcolor{FuncColor}{$\triangleright$\enspace\texttt{GeneratorsOfNearAdditiveMagma({\mdseries\slshape A})\index{GeneratorsOfNearAdditiveMagma@\texttt{GeneratorsOfNearAdditiveMagma}}
\label{GeneratorsOfNearAdditiveMagma}
}\hfill{\scriptsize (attribute)}}\\
\noindent\textcolor{FuncColor}{$\triangleright$\enspace\texttt{GeneratorsOfAdditiveMagma({\mdseries\slshape A})\index{GeneratorsOfAdditiveMagma@\texttt{GeneratorsOfAdditiveMagma}}
\label{GeneratorsOfAdditiveMagma}
}\hfill{\scriptsize (attribute)}}\\


 is a list of elements of the near\texttt{\symbol{45}}additive magma \mbox{\texttt{\mdseries\slshape A}} that generates \mbox{\texttt{\mdseries\slshape A}} as a near\texttt{\symbol{45}}additive magma, that is, the closure of this list
under addition is \mbox{\texttt{\mdseries\slshape A}}. }

 

\subsection{\textcolor{Chapter }{GeneratorsOfNearAdditiveMagmaWithZero}}
\logpage{[ 55, 3, 3 ]}\nobreak
\hyperdef{L}{X7EB9ABF880DCAE01}{}
{\noindent\textcolor{FuncColor}{$\triangleright$\enspace\texttt{GeneratorsOfNearAdditiveMagmaWithZero({\mdseries\slshape A})\index{GeneratorsOfNearAdditiveMagmaWithZero@\texttt{Generators}\-\texttt{Of}\-\texttt{Near}\-\texttt{Additive}\-\texttt{Magma}\-\texttt{With}\-\texttt{Zero}}
\label{GeneratorsOfNearAdditiveMagmaWithZero}
}\hfill{\scriptsize (attribute)}}\\
\noindent\textcolor{FuncColor}{$\triangleright$\enspace\texttt{GeneratorsOfAdditiveMagmaWithZero({\mdseries\slshape A})\index{GeneratorsOfAdditiveMagmaWithZero@\texttt{GeneratorsOfAdditiveMagmaWithZero}}
\label{GeneratorsOfAdditiveMagmaWithZero}
}\hfill{\scriptsize (attribute)}}\\


 is a list of elements of the near\texttt{\symbol{45}}additive
magma\texttt{\symbol{45}}with\texttt{\symbol{45}}zero \mbox{\texttt{\mdseries\slshape A}} that generates \mbox{\texttt{\mdseries\slshape A}} as a near\texttt{\symbol{45}}additive
magma\texttt{\symbol{45}}with\texttt{\symbol{45}}zero, that is, the closure of
this list under addition and \texttt{Zero} (\ref{Zero}) is \mbox{\texttt{\mdseries\slshape A}}. }

 

\subsection{\textcolor{Chapter }{GeneratorsOfNearAdditiveGroup}}
\logpage{[ 55, 3, 4 ]}\nobreak
\hyperdef{L}{X7EA15714795D71CF}{}
{\noindent\textcolor{FuncColor}{$\triangleright$\enspace\texttt{GeneratorsOfNearAdditiveGroup({\mdseries\slshape A})\index{GeneratorsOfNearAdditiveGroup@\texttt{GeneratorsOfNearAdditiveGroup}}
\label{GeneratorsOfNearAdditiveGroup}
}\hfill{\scriptsize (attribute)}}\\
\noindent\textcolor{FuncColor}{$\triangleright$\enspace\texttt{GeneratorsOfAdditiveGroup({\mdseries\slshape A})\index{GeneratorsOfAdditiveGroup@\texttt{GeneratorsOfAdditiveGroup}}
\label{GeneratorsOfAdditiveGroup}
}\hfill{\scriptsize (attribute)}}\\


 is a list of elements of the near\texttt{\symbol{45}}additive group \mbox{\texttt{\mdseries\slshape A}} that generates \mbox{\texttt{\mdseries\slshape A}} as a near\texttt{\symbol{45}}additive group, that is, the closure of this list
under addition, taking the zero element, and taking additive inverses
(see{\nobreakspace}\texttt{AdditiveInverse} (\ref{AdditiveInverse})) is \mbox{\texttt{\mdseries\slshape A}}. }

 

\subsection{\textcolor{Chapter }{AdditiveNeutralElement}}
\logpage{[ 55, 3, 5 ]}\nobreak
\hyperdef{L}{X851EA2E67F0C9A75}{}
{\noindent\textcolor{FuncColor}{$\triangleright$\enspace\texttt{AdditiveNeutralElement({\mdseries\slshape A})\index{AdditiveNeutralElement@\texttt{AdditiveNeutralElement}}
\label{AdditiveNeutralElement}
}\hfill{\scriptsize (attribute)}}\\


 returns the element $z$ in the near\texttt{\symbol{45}}additive magma \mbox{\texttt{\mdseries\slshape A}} with the property that $z + a = a = a + z$ holds for all $a \in$ \mbox{\texttt{\mdseries\slshape A}}, if such an element exists. Otherwise \texttt{fail} is returned. 

 A near\texttt{\symbol{45}}additive magma that is not a
near\texttt{\symbol{45}}additive
magma\texttt{\symbol{45}}with\texttt{\symbol{45}}zero can have an additive
neutral element $z$; in this case, $z$ \emph{cannot} be obtained as \texttt{Zero( \mbox{\texttt{\mdseries\slshape A}} )} or as \texttt{0*}$a$ for an element $a$ in \mbox{\texttt{\mdseries\slshape A}}, see{\nobreakspace}\texttt{Zero} (\ref{Zero}). }

 

\subsection{\textcolor{Chapter }{TrivialSubnearAdditiveMagmaWithZero}}
\logpage{[ 55, 3, 6 ]}\nobreak
\hyperdef{L}{X78FB0A5C86DC86F9}{}
{\noindent\textcolor{FuncColor}{$\triangleright$\enspace\texttt{TrivialSubnearAdditiveMagmaWithZero({\mdseries\slshape A})\index{TrivialSubnearAdditiveMagmaWithZero@\texttt{TrivialSubnearAdditiveMagmaWithZero}}
\label{TrivialSubnearAdditiveMagmaWithZero}
}\hfill{\scriptsize (attribute)}}\\


 is the additive magma\texttt{\symbol{45}}with\texttt{\symbol{45}}zero that has
the zero of the near\texttt{\symbol{45}}additive
magma\texttt{\symbol{45}}with\texttt{\symbol{45}}zero \mbox{\texttt{\mdseries\slshape A}} as its only element. }

 }

  
\section{\textcolor{Chapter }{Operations for (Near\texttt{\symbol{45}})Additive Magmas}}\label{Operations for (Near-)Additive Magmas}
\logpage{[ 55, 4, 0 ]}
\hyperdef{L}{X7BB03781863BE4EB}{}
{
  
\subsection{\textcolor{Chapter }{ClosureNearAdditiveGroup}}\logpage{[ 55, 4, 1 ]}
\hyperdef{L}{X845E915B87D2AC16}{}
{
\noindent\textcolor{FuncColor}{$\triangleright$\enspace\texttt{ClosureNearAdditiveGroup({\mdseries\slshape A, a})\index{ClosureNearAdditiveGroup@\texttt{ClosureNearAdditiveGroup}!for a near-additive group and an element}
\label{ClosureNearAdditiveGroup:for a near-additive group and an element}
}\hfill{\scriptsize (operation)}}\\
\noindent\textcolor{FuncColor}{$\triangleright$\enspace\texttt{ClosureNearAdditiveGroup({\mdseries\slshape A, B})\index{ClosureNearAdditiveGroup@\texttt{ClosureNearAdditiveGroup}!for two near-additive groups}
\label{ClosureNearAdditiveGroup:for two near-additive groups}
}\hfill{\scriptsize (operation)}}\\


 returns the closure of the near\texttt{\symbol{45}}additive magma \mbox{\texttt{\mdseries\slshape A}} with the element \mbox{\texttt{\mdseries\slshape a}} or with the near\texttt{\symbol{45}}additive magma \mbox{\texttt{\mdseries\slshape B}}, w.r.t.{\nobreakspace}addition, taking the zero element, and taking additive
inverses. }

 

\subsection{\textcolor{Chapter }{ShowAdditionTable}}
\logpage{[ 55, 4, 2 ]}\nobreak
\hyperdef{L}{X8142D994794B700A}{}
{\noindent\textcolor{FuncColor}{$\triangleright$\enspace\texttt{ShowAdditionTable({\mdseries\slshape R})\index{ShowAdditionTable@\texttt{ShowAdditionTable}}
\label{ShowAdditionTable}
}\hfill{\scriptsize (function)}}\\
\noindent\textcolor{FuncColor}{$\triangleright$\enspace\texttt{ShowMultiplicationTable({\mdseries\slshape M})\index{ShowMultiplicationTable@\texttt{ShowMultiplicationTable}}
\label{ShowMultiplicationTable}
}\hfill{\scriptsize (function)}}\\


 For a structure \mbox{\texttt{\mdseries\slshape R}} with an addition given by \texttt{+}, respectively a structure \mbox{\texttt{\mdseries\slshape M}} with a multiplication given by \texttt{*}, this command displays the addition (multiplication) table of the structure
in a pretty way. 
\begin{Verbatim}[commandchars=!@B,fontsize=\small,frame=single,label=Example]
  !gapprompt@gap>B !gapinput@ShowAdditionTable(GF(4));B
  +        | 0*Z(2)   Z(2)^0   Z(2^2)   Z(2^2)^2
  ---------+------------------------------------
  0*Z(2)   | 0*Z(2)   Z(2)^0   Z(2^2)   Z(2^2)^2
  Z(2)^0   | Z(2)^0   0*Z(2)   Z(2^2)^2 Z(2^2)
  Z(2^2)   | Z(2^2)   Z(2^2)^2 0*Z(2)   Z(2)^0
  Z(2^2)^2 | Z(2^2)^2 Z(2^2)   Z(2)^0   0*Z(2)
  
  !gapprompt@gap>B !gapinput@ShowMultiplicationTable(GF(4));B
  *        | 0*Z(2)   Z(2)^0   Z(2^2)   Z(2^2)^2
  ---------+------------------------------------
  0*Z(2)   | 0*Z(2)   0*Z(2)   0*Z(2)   0*Z(2)
  Z(2)^0   | 0*Z(2)   Z(2)^0   Z(2^2)   Z(2^2)^2
  Z(2^2)   | 0*Z(2)   Z(2^2)   Z(2^2)^2 Z(2)^0
  Z(2^2)^2 | 0*Z(2)   Z(2^2)^2 Z(2)^0   Z(2^2)
  
\end{Verbatim}
 }

 }

 }

            
\chapter{\textcolor{Chapter }{Rings}}\label{Rings}
\logpage{[ 56, 0, 0 ]}
\hyperdef{L}{X81897F6082CACB59}{}
{
  This chapter deals with domains that are additive groups (see \texttt{IsAdditiveGroup} (\ref{IsAdditiveGroup}) closed under multiplication \texttt{*}. Such a domain, if \texttt{*} and \texttt{+} are distributive, is called a \emph{ring} in \textsf{GAP}. Each division ring, field (see{\nobreakspace}\ref{Fields and Division Rings}), or algebra (see{\nobreakspace}\ref{Algebras}) is a ring. Important examples of rings are the integers (see{\nobreakspace}\ref{Integers}) and matrix rings.   

 In the case of a \emph{ring\texttt{\symbol{45}}with\texttt{\symbol{45}}one}, additional multiplicative structure is present, see{\nobreakspace}\texttt{IsRingWithOne} (\ref{IsRingWithOne}). There is a little support in \textsf{GAP} for rings that have no additional structure: it is possible to perform some
computations for small finite rings; infinite rings are handled by \textsf{GAP} in an acceptable way in the case that they are algebras. 

 Also, the \textsf{SONATA} package provides support for near\texttt{\symbol{45}}rings, and a related
functionality for multiplicative semigroups of near\texttt{\symbol{45}}rings
is available in the \textsf{Smallsemi} package. 

 Several functions for ring elements, such as \texttt{IsPrime} (\ref{IsPrime}) and \texttt{Factors} (\ref{Factors}), are defined only relative to a ring \mbox{\texttt{\mdseries\slshape R}}, which can be entered as an optional argument; if \mbox{\texttt{\mdseries\slshape R}} is omitted then a \emph{default ring} is formed from the ring elements given as arguments, see{\nobreakspace}\texttt{DefaultRing} (\ref{DefaultRing:for ring elements}).  
\section{\textcolor{Chapter }{Generating Rings}}\label{Generating Rings}
\logpage{[ 56, 1, 0 ]}
\hyperdef{L}{X839FC48687C25FCD}{}
{
  

\subsection{\textcolor{Chapter }{IsRing}}
\logpage{[ 56, 1, 1 ]}\nobreak
\hyperdef{L}{X80FD843C8221DAC9}{}
{\noindent\textcolor{FuncColor}{$\triangleright$\enspace\texttt{IsRing({\mdseries\slshape R})\index{IsRing@\texttt{IsRing}}
\label{IsRing}
}\hfill{\scriptsize (filter)}}\\


 A \emph{ring} in \textsf{GAP} is an additive group (see{\nobreakspace}\texttt{IsAdditiveGroup} (\ref{IsAdditiveGroup})) that is also a magma (see{\nobreakspace}\texttt{IsMagma} (\ref{IsMagma})), such that addition \texttt{+} and multiplication \texttt{*} are distributive, see \texttt{IsDistributive} (\ref{IsDistributive}). 

 The multiplication need \emph{not} be associative (see{\nobreakspace}\texttt{IsAssociative} (\ref{IsAssociative})). For example, a Lie algebra (see{\nobreakspace}\ref{Lie Algebras}) is regarded as a ring in \textsf{GAP}. }

 
\subsection{\textcolor{Chapter }{Ring}}\logpage{[ 56, 1, 2 ]}
\hyperdef{L}{X820B172A860A5B1A}{}
{
\noindent\textcolor{FuncColor}{$\triangleright$\enspace\texttt{Ring({\mdseries\slshape r, s, ...})\index{Ring@\texttt{Ring}!for ring elements}
\label{Ring:for ring elements}
}\hfill{\scriptsize (function)}}\\
\noindent\textcolor{FuncColor}{$\triangleright$\enspace\texttt{Ring({\mdseries\slshape coll})\index{Ring@\texttt{Ring}!for a collection}
\label{Ring:for a collection}
}\hfill{\scriptsize (function)}}\\


 In the first form \texttt{Ring} returns the smallest ring that contains all the elements \mbox{\texttt{\mdseries\slshape r}}, \mbox{\texttt{\mdseries\slshape s}}, $\ldots$ In the second form \texttt{Ring} returns the smallest ring that contains all the elements in the collection \mbox{\texttt{\mdseries\slshape coll}}. If any element is not an element of a ring or if the elements lie in no
common ring an error is raised. 

 \texttt{Ring} differs from \texttt{DefaultRing} (\ref{DefaultRing:for ring elements}) in that it returns the smallest ring in which the elements lie, while \texttt{DefaultRing} (\ref{DefaultRing:for ring elements}) may return a larger ring if that makes sense. 
\begin{Verbatim}[commandchars=!@|,fontsize=\small,frame=single,label=Example]
  !gapprompt@gap>| !gapinput@Ring( 2, E(4) );|
  <ring with 2 generators>
\end{Verbatim}
 }

 
\subsection{\textcolor{Chapter }{DefaultRing}}\logpage{[ 56, 1, 3 ]}
\hyperdef{L}{X83AFFCC77DE6ABDA}{}
{
\noindent\textcolor{FuncColor}{$\triangleright$\enspace\texttt{DefaultRing({\mdseries\slshape r, s, ...})\index{DefaultRing@\texttt{DefaultRing}!for ring elements}
\label{DefaultRing:for ring elements}
}\hfill{\scriptsize (function)}}\\
\noindent\textcolor{FuncColor}{$\triangleright$\enspace\texttt{DefaultRing({\mdseries\slshape coll})\index{DefaultRing@\texttt{DefaultRing}!for a collection}
\label{DefaultRing:for a collection}
}\hfill{\scriptsize (function)}}\\


 In the first form \texttt{DefaultRing} returns a ring that contains all the elements \mbox{\texttt{\mdseries\slshape r}}, \mbox{\texttt{\mdseries\slshape s}}, $\ldots$ etc. In the second form \texttt{DefaultRing} returns a ring that contains all the elements in the collection \mbox{\texttt{\mdseries\slshape coll}}. If any element is not an element of a ring or if the elements lie in no
common ring an error is raised. 

 The ring returned by \texttt{DefaultRing} need not be the smallest ring in which the elements lie. For example for
elements from cyclotomic fields, \texttt{DefaultRing} may return the ring of integers of the smallest cyclotomic field in which the
elements lie, which need not be the smallest ring overall, because the
elements may in fact lie in a smaller number field which is itself not a
cyclotomic field. 

 (For the exact definition of the default ring of a certain type of elements,
look at the corresponding method installation.) 

 \texttt{DefaultRing} is used by ring functions such as \texttt{Quotient} (\ref{Quotient}), \texttt{IsPrime} (\ref{IsPrime}), \texttt{Factors} (\ref{Factors}), or \texttt{Gcd} (\ref{Gcd:for (a ring and) several elements}) if no explicit ring is given. 

 \texttt{Ring} (\ref{Ring:for ring elements}) differs from \texttt{DefaultRing} in that it returns the smallest ring in which the elements lie, while \texttt{DefaultRing} may return a larger ring if that makes sense. 

 
\begin{Verbatim}[commandchars=!@|,fontsize=\small,frame=single,label=Example]
  !gapprompt@gap>| !gapinput@DefaultRing( 2, E(4) );|
  GaussianIntegers
\end{Verbatim}
 }

 

\subsection{\textcolor{Chapter }{RingByGenerators}}
\logpage{[ 56, 1, 4 ]}\nobreak
\hyperdef{L}{X7D736E027DFD8961}{}
{\noindent\textcolor{FuncColor}{$\triangleright$\enspace\texttt{RingByGenerators({\mdseries\slshape C})\index{RingByGenerators@\texttt{RingByGenerators}}
\label{RingByGenerators}
}\hfill{\scriptsize (operation)}}\\


 \texttt{RingByGenerators} returns the ring generated by the elements in the collection \mbox{\texttt{\mdseries\slshape C}}, i.{\nobreakspace}e., the closure of \mbox{\texttt{\mdseries\slshape C}} under addition, multiplication, and taking additive inverses. 
\begin{Verbatim}[commandchars=!@|,fontsize=\small,frame=single,label=Example]
  !gapprompt@gap>| !gapinput@RingByGenerators([ 2, E(4) ]);|
  <ring with 2 generators>
\end{Verbatim}
 }

 

\subsection{\textcolor{Chapter }{DefaultRingByGenerators}}
\logpage{[ 56, 1, 5 ]}\nobreak
\hyperdef{L}{X839E609480495E27}{}
{\noindent\textcolor{FuncColor}{$\triangleright$\enspace\texttt{DefaultRingByGenerators({\mdseries\slshape coll})\index{DefaultRingByGenerators@\texttt{DefaultRingByGenerators}}
\label{DefaultRingByGenerators}
}\hfill{\scriptsize (operation)}}\\


 For a collection \mbox{\texttt{\mdseries\slshape coll}}, returns a default ring in which \mbox{\texttt{\mdseries\slshape coll}} is contained. 
\begin{Verbatim}[commandchars=!@|,fontsize=\small,frame=single,label=Example]
  !gapprompt@gap>| !gapinput@DefaultRingByGenerators([ 2, E(4) ]);|
  GaussianIntegers
\end{Verbatim}
 }

 

\subsection{\textcolor{Chapter }{GeneratorsOfRing}}
\logpage{[ 56, 1, 6 ]}\nobreak
\hyperdef{L}{X7D0428D87E63288C}{}
{\noindent\textcolor{FuncColor}{$\triangleright$\enspace\texttt{GeneratorsOfRing({\mdseries\slshape R})\index{GeneratorsOfRing@\texttt{GeneratorsOfRing}}
\label{GeneratorsOfRing}
}\hfill{\scriptsize (attribute)}}\\


 \texttt{GeneratorsOfRing} returns a list of elements such that the ring \mbox{\texttt{\mdseries\slshape R}} is the closure of these elements under addition, multiplication, and taking
additive inverses. 
\begin{Verbatim}[commandchars=!@|,fontsize=\small,frame=single,label=Example]
  !gapprompt@gap>| !gapinput@R:=Ring( 2, 1/2 );|
  <ring with 2 generators>
  !gapprompt@gap>| !gapinput@GeneratorsOfRing( R );|
  [ 2, 1/2 ]
\end{Verbatim}
 }

  

\subsection{\textcolor{Chapter }{Subring}}
\logpage{[ 56, 1, 7 ]}\nobreak
\hyperdef{L}{X860E4AC78520D27E}{}
{\noindent\textcolor{FuncColor}{$\triangleright$\enspace\texttt{Subring({\mdseries\slshape R, gens})\index{Subring@\texttt{Subring}}
\label{Subring}
}\hfill{\scriptsize (function)}}\\
\noindent\textcolor{FuncColor}{$\triangleright$\enspace\texttt{SubringNC({\mdseries\slshape R, gens})\index{SubringNC@\texttt{SubringNC}}
\label{SubringNC}
}\hfill{\scriptsize (function)}}\\


 returns the ring with parent \mbox{\texttt{\mdseries\slshape R}} generated by the elements in \mbox{\texttt{\mdseries\slshape gens}}. When the second form, \texttt{SubringNC} is used, it is \emph{not} checked whether all elements in \mbox{\texttt{\mdseries\slshape gens}} lie in \mbox{\texttt{\mdseries\slshape R}}. 

 
\begin{Verbatim}[commandchars=!@|,fontsize=\small,frame=single,label=Example]
  !gapprompt@gap>| !gapinput@R:= Integers;|
  Integers
  !gapprompt@gap>| !gapinput@S:= Subring( R, [ 4, 6 ] );|
  <ring with 1 generator>
  !gapprompt@gap>| !gapinput@Parent( S );|
  Integers
\end{Verbatim}
 }

 
\subsection{\textcolor{Chapter }{ClosureRing}}\logpage{[ 56, 1, 8 ]}
\hyperdef{L}{X819B0AFE79C78C34}{}
{
\noindent\textcolor{FuncColor}{$\triangleright$\enspace\texttt{ClosureRing({\mdseries\slshape R, r})\index{ClosureRing@\texttt{ClosureRing}!for a ring and a ring element}
\label{ClosureRing:for a ring and a ring element}
}\hfill{\scriptsize (operation)}}\\
\noindent\textcolor{FuncColor}{$\triangleright$\enspace\texttt{ClosureRing({\mdseries\slshape R, S})\index{ClosureRing@\texttt{ClosureRing}!for two rings}
\label{ClosureRing:for two rings}
}\hfill{\scriptsize (operation)}}\\


 For a ring \mbox{\texttt{\mdseries\slshape R}} and either an element \mbox{\texttt{\mdseries\slshape r}} of its elements family or a ring \mbox{\texttt{\mdseries\slshape S}}, \texttt{ClosureRing} returns the ring generated by both arguments. 

 
\begin{Verbatim}[commandchars=!@|,fontsize=\small,frame=single,label=Example]
  !gapprompt@gap>| !gapinput@ClosureRing( Integers, E(4) );|
  <ring-with-one, with 2 generators>
\end{Verbatim}
 }

 

\subsection{\textcolor{Chapter }{Quotient}}
\logpage{[ 56, 1, 9 ]}\nobreak
\hyperdef{L}{X8350500B8576F833}{}
{\noindent\textcolor{FuncColor}{$\triangleright$\enspace\texttt{Quotient({\mdseries\slshape [R, ]r, s})\index{Quotient@\texttt{Quotient}}
\label{Quotient}
}\hfill{\scriptsize (operation)}}\\


 \texttt{Quotient} returns a (right) quotient of the two ring elements \mbox{\texttt{\mdseries\slshape r}} and \mbox{\texttt{\mdseries\slshape s}} in the ring \mbox{\texttt{\mdseries\slshape R}}, if given, and otherwise in their default ring (see \texttt{DefaultRing} (\ref{DefaultRing:for ring elements})). More specifically, it returns a ring element $q$ such that $r = q * s$ holds, or \texttt{fail} if no such elements exists in the respective ring. 

 The result may not be unique if the ring contains zero divisors. 

 (To perform the division in the quotient field of a ring, use the quotient
operator \texttt{/}.) 
\begin{Verbatim}[commandchars=!@|,fontsize=\small,frame=single,label=Example]
  !gapprompt@gap>| !gapinput@Quotient( 2, 3 );|
  fail
  !gapprompt@gap>| !gapinput@Quotient( 6, 3 );|
  2
\end{Verbatim}
 }

 }

  
\section{\textcolor{Chapter }{Ideals of Rings}}\label{Ideals of Rings}
\logpage{[ 56, 2, 0 ]}
\hyperdef{L}{X8776C3F97A731E70}{}
{
  A \emph{left ideal} in a ring $R$ is a subring of $R$ that is closed under multiplication with elements of $R$ from the left. 

 A \emph{right ideal} in a ring $R$ is a subring of $R$ that is closed under multiplication with elements of $R$ from the right. 

 A \emph{two\texttt{\symbol{45}}sided ideal} or simply \emph{ideal} in a ring $R$ is both a left ideal and a right ideal in $R$. 

 So being a (left/right/two\texttt{\symbol{45}}sided) ideal is not a property
of a domain but refers to the acting ring(s). Hence we must ask,
e.{\nobreakspace}g., \texttt{IsIdeal( }$R, I$\texttt{ )} if we want to know whether the ring $I$ is an ideal in the ring $R$. The property \texttt{IsTwoSidedIdealInParent} (\ref{IsTwoSidedIdealInParent}) can be used to store whether a ring is an ideal in its parent. 

 (Whenever the term \texttt{"Ideal"} occurs in an identifier without a specifying prefix \texttt{"Left"} or \texttt{"Right"}, this means the same as \texttt{"TwoSidedIdeal"}. Conversely, any occurrence of \texttt{"TwoSidedIdeal"} can be substituted by \texttt{"Ideal"}.) 

 For any of the above kinds of ideals, there is a notion of generators, namely \texttt{GeneratorsOfLeftIdeal} (\ref{GeneratorsOfLeftIdeal}), \texttt{GeneratorsOfRightIdeal} (\ref{GeneratorsOfRightIdeal}), and \texttt{GeneratorsOfTwoSidedIdeal} (\ref{GeneratorsOfTwoSidedIdeal}). The acting rings can be accessed as \texttt{LeftActingRingOfIdeal} (\ref{LeftActingRingOfIdeal}) and \texttt{RightActingRingOfIdeal} (\ref{RightActingRingOfIdeal}), respectively. Note that ideals are detected from known values of these
attributes, especially it is assumed that whenever a domain has both a left
and a right acting ring then these two are equal. 

 Note that we cannot use \texttt{LeftActingDomain} (\ref{LeftActingDomain}) and \texttt{RightActingDomain} here, since ideals in algebras are themselves vector spaces, and such a space
can of course also be a module for an action from the right. In order to make
the usual vector space functionality automatically available for ideals, we
have to distinguish the left and right module structure from the additional
closure properties of the ideal. 

 Further note that the attributes denoting ideal generators and acting ring are
used to create ideals if this is explicitly wanted, but the ideal relation in
the sense of \texttt{IsTwoSidedIdeal} (\ref{IsTwoSidedIdeal}) is of course independent of the presence of the attribute values. 

 Ideals are constructed with \texttt{LeftIdeal} (\ref{LeftIdeal}), \texttt{RightIdeal} (\ref{RightIdeal}), \texttt{TwoSidedIdeal} (\ref{TwoSidedIdeal}). Principal ideals of the form $x * R$, $R * x$, $R * x * R$ can also be constructed with a simple multiplication. 

 Currently many methods for dealing with ideals need linear algebra to work, so
they are mainly applicable to ideals in algebras. 

 

\subsection{\textcolor{Chapter }{TwoSidedIdeal}}
\logpage{[ 56, 2, 1 ]}\nobreak
\hyperdef{L}{X7C486A7C821D79F0}{}
{\noindent\textcolor{FuncColor}{$\triangleright$\enspace\texttt{TwoSidedIdeal({\mdseries\slshape R, gens[, "basis"]})\index{TwoSidedIdeal@\texttt{TwoSidedIdeal}}
\label{TwoSidedIdeal}
}\hfill{\scriptsize (function)}}\\
\noindent\textcolor{FuncColor}{$\triangleright$\enspace\texttt{Ideal({\mdseries\slshape R, gens[, "basis"]})\index{Ideal@\texttt{Ideal}}
\label{Ideal}
}\hfill{\scriptsize (function)}}\\
\noindent\textcolor{FuncColor}{$\triangleright$\enspace\texttt{LeftIdeal({\mdseries\slshape R, gens[, "basis"]})\index{LeftIdeal@\texttt{LeftIdeal}}
\label{LeftIdeal}
}\hfill{\scriptsize (function)}}\\
\noindent\textcolor{FuncColor}{$\triangleright$\enspace\texttt{RightIdeal({\mdseries\slshape R, gens[, "basis"]})\index{RightIdeal@\texttt{RightIdeal}}
\label{RightIdeal}
}\hfill{\scriptsize (function)}}\\


 Let \mbox{\texttt{\mdseries\slshape R}} be a ring, and \mbox{\texttt{\mdseries\slshape gens}} a list of collection of elements in \mbox{\texttt{\mdseries\slshape R}}. \texttt{TwoSidedIdeal}, \texttt{LeftIdeal}, and \texttt{RightIdeal} return the two\texttt{\symbol{45}}sided, left, or right ideal, respectively, $I$ in \mbox{\texttt{\mdseries\slshape R}} that is generated by \mbox{\texttt{\mdseries\slshape gens}}. The ring \mbox{\texttt{\mdseries\slshape R}} can be accessed as \texttt{LeftActingRingOfIdeal} (\ref{LeftActingRingOfIdeal}) or \texttt{RightActingRingOfIdeal} (\ref{RightActingRingOfIdeal}) (or both) of $I$. 

 If \mbox{\texttt{\mdseries\slshape R}} is a left $F$\texttt{\symbol{45}}module then also $I$ is a left $F$\texttt{\symbol{45}}module, in particular the \texttt{LeftActingDomain} (\ref{LeftActingDomain}) values of \mbox{\texttt{\mdseries\slshape R}} and $I$ are equal. 

 If the optional argument \texttt{"basis"} is given then \mbox{\texttt{\mdseries\slshape gens}} are assumed to be a list of basis vectors of $I$ viewed as a free $F$\texttt{\symbol{45}}module. (This is mainly applicable to ideals in algebras.)
In this case, it is \emph{not} checked whether \mbox{\texttt{\mdseries\slshape gens}} really is linearly independent and whether \mbox{\texttt{\mdseries\slshape gens}} is a subset of \mbox{\texttt{\mdseries\slshape R}}. 

 \texttt{Ideal} is simply a synonym of \texttt{TwoSidedIdeal}. 

 
\begin{Verbatim}[commandchars=!@|,fontsize=\small,frame=single,label=Example]
  !gapprompt@gap>| !gapinput@R:= Integers;;|
  !gapprompt@gap>| !gapinput@I:= Ideal( R, [ 2 ] );|
  <two-sided ideal in Integers, (1 generator)>
\end{Verbatim}
 }

 

\subsection{\textcolor{Chapter }{TwoSidedIdealNC}}
\logpage{[ 56, 2, 2 ]}\nobreak
\hyperdef{L}{X7C8E196478C7431A}{}
{\noindent\textcolor{FuncColor}{$\triangleright$\enspace\texttt{TwoSidedIdealNC({\mdseries\slshape R, gens[, "basis"]})\index{TwoSidedIdealNC@\texttt{TwoSidedIdealNC}}
\label{TwoSidedIdealNC}
}\hfill{\scriptsize (function)}}\\
\noindent\textcolor{FuncColor}{$\triangleright$\enspace\texttt{IdealNC({\mdseries\slshape R, gens[, "basis"]})\index{IdealNC@\texttt{IdealNC}}
\label{IdealNC}
}\hfill{\scriptsize (function)}}\\
\noindent\textcolor{FuncColor}{$\triangleright$\enspace\texttt{LeftIdealNC({\mdseries\slshape R, gens[, "basis"]})\index{LeftIdealNC@\texttt{LeftIdealNC}}
\label{LeftIdealNC}
}\hfill{\scriptsize (function)}}\\
\noindent\textcolor{FuncColor}{$\triangleright$\enspace\texttt{RightIdealNC({\mdseries\slshape R, gens[, "basis"]})\index{RightIdealNC@\texttt{RightIdealNC}}
\label{RightIdealNC}
}\hfill{\scriptsize (function)}}\\


 The effects of \texttt{TwoSidedIdealNC}, \texttt{LeftIdealNC}, and \texttt{RightIdealNC} are the same as \texttt{TwoSidedIdeal} (\ref{TwoSidedIdeal}), \texttt{LeftIdeal} (\ref{LeftIdeal}), and \texttt{RightIdeal} (\ref{RightIdeal}), respectively, but they do not check whether all entries of \mbox{\texttt{\mdseries\slshape gens}} lie in \mbox{\texttt{\mdseries\slshape R}}. }

 

\subsection{\textcolor{Chapter }{IsTwoSidedIdeal}}
\logpage{[ 56, 2, 3 ]}\nobreak
\hyperdef{L}{X7DF623847B338850}{}
{\noindent\textcolor{FuncColor}{$\triangleright$\enspace\texttt{IsTwoSidedIdeal({\mdseries\slshape R, I})\index{IsTwoSidedIdeal@\texttt{IsTwoSidedIdeal}}
\label{IsTwoSidedIdeal}
}\hfill{\scriptsize (operation)}}\\
\noindent\textcolor{FuncColor}{$\triangleright$\enspace\texttt{IsLeftIdeal({\mdseries\slshape R, I})\index{IsLeftIdeal@\texttt{IsLeftIdeal}}
\label{IsLeftIdeal}
}\hfill{\scriptsize (operation)}}\\
\noindent\textcolor{FuncColor}{$\triangleright$\enspace\texttt{IsRightIdeal({\mdseries\slshape R, I})\index{IsRightIdeal@\texttt{IsRightIdeal}}
\label{IsRightIdeal}
}\hfill{\scriptsize (operation)}}\\
\noindent\textcolor{FuncColor}{$\triangleright$\enspace\texttt{IsTwoSidedIdealInParent({\mdseries\slshape I})\index{IsTwoSidedIdealInParent@\texttt{IsTwoSidedIdealInParent}}
\label{IsTwoSidedIdealInParent}
}\hfill{\scriptsize (property)}}\\
\noindent\textcolor{FuncColor}{$\triangleright$\enspace\texttt{IsLeftIdealInParent({\mdseries\slshape I})\index{IsLeftIdealInParent@\texttt{IsLeftIdealInParent}}
\label{IsLeftIdealInParent}
}\hfill{\scriptsize (property)}}\\
\noindent\textcolor{FuncColor}{$\triangleright$\enspace\texttt{IsRightIdealInParent({\mdseries\slshape I})\index{IsRightIdealInParent@\texttt{IsRightIdealInParent}}
\label{IsRightIdealInParent}
}\hfill{\scriptsize (property)}}\\


 The properties \texttt{IsTwoSidedIdealInParent} etc., are attributes of the ideal, and once known they are stored in the
ideal. 
\begin{Verbatim}[commandchars=!@|,fontsize=\small,frame=single,label=Example]
  !gapprompt@gap>| !gapinput@A:= FullMatrixAlgebra( Rationals, 3 );|
  ( Rationals^[ 3, 3 ] )
  !gapprompt@gap>| !gapinput@I:= Ideal( A, [ Random( A ) ] );|
  <two-sided ideal in ( Rationals^[ 3, 3 ] ), (1 generator)>
  !gapprompt@gap>| !gapinput@IsTwoSidedIdeal( A, I );|
  true
\end{Verbatim}
 }

 

\subsection{\textcolor{Chapter }{TwoSidedIdealByGenerators}}
\logpage{[ 56, 2, 4 ]}\nobreak
\hyperdef{L}{X86C998178690DAE0}{}
{\noindent\textcolor{FuncColor}{$\triangleright$\enspace\texttt{TwoSidedIdealByGenerators({\mdseries\slshape R, gens})\index{TwoSidedIdealByGenerators@\texttt{TwoSidedIdealByGenerators}}
\label{TwoSidedIdealByGenerators}
}\hfill{\scriptsize (operation)}}\\
\noindent\textcolor{FuncColor}{$\triangleright$\enspace\texttt{IdealByGenerators({\mdseries\slshape R, gens})\index{IdealByGenerators@\texttt{IdealByGenerators}}
\label{IdealByGenerators}
}\hfill{\scriptsize (operation)}}\\


 \texttt{TwoSidedIdealByGenerators} returns the ring that is generated by the elements of the collection \mbox{\texttt{\mdseries\slshape gens}} under addition, multiplication, and multiplication with elements of the ring \mbox{\texttt{\mdseries\slshape R}} from the left and from the right. 

 \mbox{\texttt{\mdseries\slshape R}} can be accessed by \texttt{LeftActingRingOfIdeal} (\ref{LeftActingRingOfIdeal}) or \texttt{RightActingRingOfIdeal} (\ref{RightActingRingOfIdeal}), \mbox{\texttt{\mdseries\slshape gens}} can be accessed by \texttt{GeneratorsOfTwoSidedIdeal} (\ref{GeneratorsOfTwoSidedIdeal}). }

 

\subsection{\textcolor{Chapter }{LeftIdealByGenerators}}
\logpage{[ 56, 2, 5 ]}\nobreak
\hyperdef{L}{X82D8B07281EB0AC7}{}
{\noindent\textcolor{FuncColor}{$\triangleright$\enspace\texttt{LeftIdealByGenerators({\mdseries\slshape R, gens})\index{LeftIdealByGenerators@\texttt{LeftIdealByGenerators}}
\label{LeftIdealByGenerators}
}\hfill{\scriptsize (operation)}}\\


 \texttt{LeftIdealByGenerators} returns the ring that is generated by the elements of the collection \mbox{\texttt{\mdseries\slshape gens}} under addition, multiplication, and multiplication with elements of the ring \mbox{\texttt{\mdseries\slshape R}} from the left. 

 \mbox{\texttt{\mdseries\slshape R}} can be accessed by \texttt{LeftActingRingOfIdeal} (\ref{LeftActingRingOfIdeal}), \mbox{\texttt{\mdseries\slshape gens}} can be accessed by \texttt{GeneratorsOfLeftIdeal} (\ref{GeneratorsOfLeftIdeal}). }

 

\subsection{\textcolor{Chapter }{RightIdealByGenerators}}
\logpage{[ 56, 2, 6 ]}\nobreak
\hyperdef{L}{X858EAEAF87751428}{}
{\noindent\textcolor{FuncColor}{$\triangleright$\enspace\texttt{RightIdealByGenerators({\mdseries\slshape R, gens})\index{RightIdealByGenerators@\texttt{RightIdealByGenerators}}
\label{RightIdealByGenerators}
}\hfill{\scriptsize (operation)}}\\


 \texttt{RightIdealByGenerators} returns the ring that is generated by the elements of the collection \mbox{\texttt{\mdseries\slshape gens}} under addition, multiplication, and multiplication with elements of the ring \mbox{\texttt{\mdseries\slshape R}} from the right. 

 \mbox{\texttt{\mdseries\slshape R}} can be accessed by \texttt{RightActingRingOfIdeal} (\ref{RightActingRingOfIdeal}), \mbox{\texttt{\mdseries\slshape gens}} can be accessed by \texttt{GeneratorsOfRightIdeal} (\ref{GeneratorsOfRightIdeal}). }

 

\subsection{\textcolor{Chapter }{GeneratorsOfTwoSidedIdeal}}
\logpage{[ 56, 2, 7 ]}\nobreak
\hyperdef{L}{X86AAF5F9800E97EE}{}
{\noindent\textcolor{FuncColor}{$\triangleright$\enspace\texttt{GeneratorsOfTwoSidedIdeal({\mdseries\slshape I})\index{GeneratorsOfTwoSidedIdeal@\texttt{GeneratorsOfTwoSidedIdeal}}
\label{GeneratorsOfTwoSidedIdeal}
}\hfill{\scriptsize (attribute)}}\\
\noindent\textcolor{FuncColor}{$\triangleright$\enspace\texttt{GeneratorsOfIdeal({\mdseries\slshape I})\index{GeneratorsOfIdeal@\texttt{GeneratorsOfIdeal}}
\label{GeneratorsOfIdeal}
}\hfill{\scriptsize (attribute)}}\\


 is a list of generators for the ideal \mbox{\texttt{\mdseries\slshape I}}, with respect to the action of the rings that are stored as the values of \texttt{LeftActingRingOfIdeal} (\ref{LeftActingRingOfIdeal}) and \texttt{RightActingRingOfIdeal} (\ref{RightActingRingOfIdeal}), from the left and from the right, respectively. 

 
\begin{Verbatim}[commandchars=!@|,fontsize=\small,frame=single,label=Example]
  !gapprompt@gap>| !gapinput@A:= FullMatrixAlgebra( Rationals, 3 );;|
  !gapprompt@gap>| !gapinput@I:= Ideal( A, [ One( A ) ] );;|
  !gapprompt@gap>| !gapinput@GeneratorsOfIdeal( I );|
  [ [ [ 1, 0, 0 ], [ 0, 1, 0 ], [ 0, 0, 1 ] ] ]
\end{Verbatim}
 }

 

\subsection{\textcolor{Chapter }{GeneratorsOfLeftIdeal}}
\logpage{[ 56, 2, 8 ]}\nobreak
\hyperdef{L}{X7B20BD2B7FAFBD64}{}
{\noindent\textcolor{FuncColor}{$\triangleright$\enspace\texttt{GeneratorsOfLeftIdeal({\mdseries\slshape I})\index{GeneratorsOfLeftIdeal@\texttt{GeneratorsOfLeftIdeal}}
\label{GeneratorsOfLeftIdeal}
}\hfill{\scriptsize (attribute)}}\\


 is a list of generators for the left ideal \mbox{\texttt{\mdseries\slshape I}}, with respect to the action from the left of the ring that is stored as the
value of \texttt{LeftActingRingOfIdeal} (\ref{LeftActingRingOfIdeal}). }

 

\subsection{\textcolor{Chapter }{GeneratorsOfRightIdeal}}
\logpage{[ 56, 2, 9 ]}\nobreak
\hyperdef{L}{X80F2239F8653FF74}{}
{\noindent\textcolor{FuncColor}{$\triangleright$\enspace\texttt{GeneratorsOfRightIdeal({\mdseries\slshape I})\index{GeneratorsOfRightIdeal@\texttt{GeneratorsOfRightIdeal}}
\label{GeneratorsOfRightIdeal}
}\hfill{\scriptsize (attribute)}}\\


 is a list of generators for the right ideal \mbox{\texttt{\mdseries\slshape I}}, with respect to the action from the right of the ring that is stored as the
value of \texttt{RightActingRingOfIdeal} (\ref{RightActingRingOfIdeal}). }

 

\subsection{\textcolor{Chapter }{LeftActingRingOfIdeal}}
\logpage{[ 56, 2, 10 ]}\nobreak
\hyperdef{L}{X81D81D027C2F8D06}{}
{\noindent\textcolor{FuncColor}{$\triangleright$\enspace\texttt{LeftActingRingOfIdeal({\mdseries\slshape I})\index{LeftActingRingOfIdeal@\texttt{LeftActingRingOfIdeal}}
\label{LeftActingRingOfIdeal}
}\hfill{\scriptsize (attribute)}}\\
\noindent\textcolor{FuncColor}{$\triangleright$\enspace\texttt{RightActingRingOfIdeal({\mdseries\slshape I})\index{RightActingRingOfIdeal@\texttt{RightActingRingOfIdeal}}
\label{RightActingRingOfIdeal}
}\hfill{\scriptsize (attribute)}}\\


 returns the left (resp. right) acting ring of an ideal \mbox{\texttt{\mdseries\slshape I}}. }

 

\subsection{\textcolor{Chapter }{AsLeftIdeal}}
\logpage{[ 56, 2, 11 ]}\nobreak
\hyperdef{L}{X83D9D7408706B69A}{}
{\noindent\textcolor{FuncColor}{$\triangleright$\enspace\texttt{AsLeftIdeal({\mdseries\slshape R, S})\index{AsLeftIdeal@\texttt{AsLeftIdeal}}
\label{AsLeftIdeal}
}\hfill{\scriptsize (operation)}}\\
\noindent\textcolor{FuncColor}{$\triangleright$\enspace\texttt{AsRightIdeal({\mdseries\slshape R, S})\index{AsRightIdeal@\texttt{AsRightIdeal}}
\label{AsRightIdeal}
}\hfill{\scriptsize (operation)}}\\
\noindent\textcolor{FuncColor}{$\triangleright$\enspace\texttt{AsTwoSidedIdeal({\mdseries\slshape R, S})\index{AsTwoSidedIdeal@\texttt{AsTwoSidedIdeal}}
\label{AsTwoSidedIdeal}
}\hfill{\scriptsize (operation)}}\\


 Let \mbox{\texttt{\mdseries\slshape S}} be a subring of the ring \mbox{\texttt{\mdseries\slshape R}}. 

 If \mbox{\texttt{\mdseries\slshape S}} is a left ideal in \mbox{\texttt{\mdseries\slshape R}} then \texttt{AsLeftIdeal} returns this left ideal, otherwise \texttt{fail} is returned. 

 If \mbox{\texttt{\mdseries\slshape S}} is a right ideal in \mbox{\texttt{\mdseries\slshape R}} then \texttt{AsRightIdeal} returns this right ideal, otherwise \texttt{fail} is returned. 

 If \mbox{\texttt{\mdseries\slshape S}} is a two\texttt{\symbol{45}}sided ideal in \mbox{\texttt{\mdseries\slshape R}} then \texttt{AsTwoSidedIdeal} returns this two\texttt{\symbol{45}}sided ideal, otherwise \texttt{fail} is returned. 

 
\begin{Verbatim}[commandchars=!@|,fontsize=\small,frame=single,label=Example]
  !gapprompt@gap>| !gapinput@A:= FullMatrixAlgebra( Rationals, 3 );;|
  !gapprompt@gap>| !gapinput@B:= DirectSumOfAlgebras( A, A );|
  <algebra over Rationals, with 6 generators>
  !gapprompt@gap>| !gapinput@C:= Subalgebra( B, Basis( B ){[1..9]} );|
  <algebra over Rationals, with 9 generators>
  !gapprompt@gap>| !gapinput@I:= AsTwoSidedIdeal( B, C );|
  <two-sided ideal in <algebra of dimension 18 over Rationals>,
    (9 generators)>
\end{Verbatim}
 }

 }

  
\section{\textcolor{Chapter }{Rings With One}}\label{Rings With One}
\logpage{[ 56, 3, 0 ]}
\hyperdef{L}{X790DD00586F9B8B8}{}
{
  

\subsection{\textcolor{Chapter }{IsRingWithOne}}
\logpage{[ 56, 3, 1 ]}\nobreak
\hyperdef{L}{X7E601FBD8020A0F3}{}
{\noindent\textcolor{FuncColor}{$\triangleright$\enspace\texttt{IsRingWithOne({\mdseries\slshape R})\index{IsRingWithOne@\texttt{IsRingWithOne}}
\label{IsRingWithOne}
}\hfill{\scriptsize (filter)}}\\


 A \emph{ring\texttt{\symbol{45}}with\texttt{\symbol{45}}one} in \textsf{GAP} is a ring (see{\nobreakspace}\texttt{IsRing} (\ref{IsRing})) that is also a magma\texttt{\symbol{45}}with\texttt{\symbol{45}}one
(see{\nobreakspace}\texttt{IsMagmaWithOne} (\ref{IsMagmaWithOne})). 

 Note that the identity and the zero of a
ring\texttt{\symbol{45}}with\texttt{\symbol{45}}one need \emph{not} be distinct. This means that a ring that consists only of its zero element can
be regarded as a ring\texttt{\symbol{45}}with\texttt{\symbol{45}}one.   

 This is especially useful in the case of finitely presented rings, in the
sense that each factor of a
ring\texttt{\symbol{45}}with\texttt{\symbol{45}}one is again a
ring\texttt{\symbol{45}}with\texttt{\symbol{45}}one. }

 
\subsection{\textcolor{Chapter }{RingWithOne}}\logpage{[ 56, 3, 2 ]}
\hyperdef{L}{X80942A318417366E}{}
{
\noindent\textcolor{FuncColor}{$\triangleright$\enspace\texttt{RingWithOne({\mdseries\slshape r, s, ...})\index{RingWithOne@\texttt{RingWithOne}!for ring elements}
\label{RingWithOne:for ring elements}
}\hfill{\scriptsize (function)}}\\
\noindent\textcolor{FuncColor}{$\triangleright$\enspace\texttt{RingWithOne({\mdseries\slshape coll})\index{RingWithOne@\texttt{RingWithOne}!for a collection}
\label{RingWithOne:for a collection}
}\hfill{\scriptsize (function)}}\\


 In the first form \texttt{RingWithOne} returns the smallest ring with one that contains all the elements \mbox{\texttt{\mdseries\slshape r}}, \mbox{\texttt{\mdseries\slshape s}}, $\ldots$ In the second form \texttt{RingWithOne} returns the smallest ring with one that contains all the elements in the
collection \mbox{\texttt{\mdseries\slshape C}}. If any element is not an element of a ring or if the elements lie in no
common ring an error is raised. 
\begin{Verbatim}[commandchars=!@|,fontsize=\small,frame=single,label=Example]
  !gapprompt@gap>| !gapinput@RingWithOne( [ 4, 6 ] );|
  Integers
\end{Verbatim}
 }

 

\subsection{\textcolor{Chapter }{RingWithOneByGenerators}}
\logpage{[ 56, 3, 3 ]}\nobreak
\hyperdef{L}{X851115EC79B8C393}{}
{\noindent\textcolor{FuncColor}{$\triangleright$\enspace\texttt{RingWithOneByGenerators({\mdseries\slshape coll})\index{RingWithOneByGenerators@\texttt{RingWithOneByGenerators}}
\label{RingWithOneByGenerators}
}\hfill{\scriptsize (operation)}}\\


 \texttt{RingWithOneByGenerators} returns the ring\texttt{\symbol{45}}with\texttt{\symbol{45}}one generated by
the elements in the collection \mbox{\texttt{\mdseries\slshape coll}}, i.{\nobreakspace}e., the closure of \mbox{\texttt{\mdseries\slshape coll}} under addition, multiplication, taking additive inverses, and taking the
identity of an element. }

 

\subsection{\textcolor{Chapter }{GeneratorsOfRingWithOne}}
\logpage{[ 56, 3, 4 ]}\nobreak
\hyperdef{L}{X7F9F122C831BCDD1}{}
{\noindent\textcolor{FuncColor}{$\triangleright$\enspace\texttt{GeneratorsOfRingWithOne({\mdseries\slshape R})\index{GeneratorsOfRingWithOne@\texttt{GeneratorsOfRingWithOne}}
\label{GeneratorsOfRingWithOne}
}\hfill{\scriptsize (attribute)}}\\


 \texttt{GeneratorsOfRingWithOne} returns a list of elements such that the ring \mbox{\texttt{\mdseries\slshape R}} is the closure of these elements under addition, multiplication, taking
additive inverses, and taking the identity element \texttt{One( \mbox{\texttt{\mdseries\slshape R}} )}. 

 \mbox{\texttt{\mdseries\slshape R}} itself need \emph{not} be known to be a ring\texttt{\symbol{45}}with\texttt{\symbol{45}}one. 

 
\begin{Verbatim}[commandchars=!@|,fontsize=\small,frame=single,label=Example]
  !gapprompt@gap>| !gapinput@R:= RingWithOne( [ 4, 6 ] );|
  Integers
  !gapprompt@gap>| !gapinput@GeneratorsOfRingWithOne( R );|
  [ 1 ]
\end{Verbatim}
 }

 

\subsection{\textcolor{Chapter }{SubringWithOne}}
\logpage{[ 56, 3, 5 ]}\nobreak
\hyperdef{L}{X7D0BADF178D4DDF8}{}
{\noindent\textcolor{FuncColor}{$\triangleright$\enspace\texttt{SubringWithOne({\mdseries\slshape R, gens})\index{SubringWithOne@\texttt{SubringWithOne}}
\label{SubringWithOne}
}\hfill{\scriptsize (function)}}\\
\noindent\textcolor{FuncColor}{$\triangleright$\enspace\texttt{SubringWithOneNC({\mdseries\slshape R, gens})\index{SubringWithOneNC@\texttt{SubringWithOneNC}}
\label{SubringWithOneNC}
}\hfill{\scriptsize (function)}}\\


 returns the ring with one with parent \mbox{\texttt{\mdseries\slshape R}} generated by the elements in \mbox{\texttt{\mdseries\slshape gens}}. When the second form, \texttt{SubringWithOneNC} is used, it is \emph{not} checked whether all elements in \mbox{\texttt{\mdseries\slshape gens}} lie in \mbox{\texttt{\mdseries\slshape R}}. 

 
\begin{Verbatim}[commandchars=!@|,fontsize=\small,frame=single,label=Example]
  !gapprompt@gap>| !gapinput@R:= SubringWithOne( Integers, [ 4, 6 ] );|
  Integers
  !gapprompt@gap>| !gapinput@Parent( R );|
  Integers
\end{Verbatim}
 }

 }

  
\section{\textcolor{Chapter }{Properties of Rings}}\label{Properties of Rings}
\logpage{[ 56, 4, 0 ]}
\hyperdef{L}{X797F5869874BDBFB}{}
{
  

\subsection{\textcolor{Chapter }{IsIntegralRing}}
\logpage{[ 56, 4, 1 ]}\nobreak
\hyperdef{L}{X87A7D5B584713B52}{}
{\noindent\textcolor{FuncColor}{$\triangleright$\enspace\texttt{IsIntegralRing({\mdseries\slshape R})\index{IsIntegralRing@\texttt{IsIntegralRing}}
\label{IsIntegralRing}
}\hfill{\scriptsize (property)}}\\


 A ring\texttt{\symbol{45}}with\texttt{\symbol{45}}one \mbox{\texttt{\mdseries\slshape R}} is integral if it is commutative, contains no nontrivial zero divisors, and if
its identity is distinct from its zero. 
\begin{Verbatim}[commandchars=!@|,fontsize=\small,frame=single,label=Example]
  !gapprompt@gap>| !gapinput@IsIntegralRing( Integers );|
  true
\end{Verbatim}
 }

 

\subsection{\textcolor{Chapter }{IsUniqueFactorizationRing}}
\logpage{[ 56, 4, 2 ]}\nobreak
\hyperdef{L}{X789A917085DB7527}{}
{\noindent\textcolor{FuncColor}{$\triangleright$\enspace\texttt{IsUniqueFactorizationRing({\mdseries\slshape R})\index{IsUniqueFactorizationRing@\texttt{IsUniqueFactorizationRing}}
\label{IsUniqueFactorizationRing}
}\hfill{\scriptsize (Category)}}\\


 A ring \mbox{\texttt{\mdseries\slshape R}} is called a \emph{unique factorization ring} if every nonzero element has a unique factorization into irreducible elements,
i.e., a unique representation as product of irreducibles (see \texttt{IsIrreducibleRingElement} (\ref{IsIrreducibleRingElement})). Unique in this context means unique up to permutations of the factors and
up to multiplication of the factors by units (see{\nobreakspace}\texttt{Units} (\ref{Units})). 

 (Note that we cannot install a subset maintained method for this filter since
the factorization of an element needs not exist in a subring. As an example,
consider the subring $4 {\ensuremath{\mathbb N}} + 1$ of the ring $4 {\ensuremath{\mathbb Z}} + 1$; in the subring, the element $3 \cdot 3 \cdot 11 \cdot 7$ has the two factorizations $33 \cdot 21 = 9 \cdot 77$, but in the large ring there is the unique factorization $(-3) \cdot (-3) \cdot (-11) \cdot (-7)$, and it is easy to see that every element in $4 {\ensuremath{\mathbb Z}} + 1$ has a unique factorization.) 

 
\begin{Verbatim}[commandchars=!@|,fontsize=\small,frame=single,label=Example]
  !gapprompt@gap>| !gapinput@IsUniqueFactorizationRing( PolynomialRing( Rationals, 1 ) );|
  true
\end{Verbatim}
 }

 

\subsection{\textcolor{Chapter }{IsLDistributive}}
\logpage{[ 56, 4, 3 ]}\nobreak
\hyperdef{L}{X7D4BB44187C55BF2}{}
{\noindent\textcolor{FuncColor}{$\triangleright$\enspace\texttt{IsLDistributive({\mdseries\slshape C})\index{IsLDistributive@\texttt{IsLDistributive}}
\label{IsLDistributive}
}\hfill{\scriptsize (property)}}\\


 is \texttt{true} if the relation $a * ( b + c ) = ( a * b ) + ( a * c )$ holds for all elements $a$, $b$, $c$ in the collection \mbox{\texttt{\mdseries\slshape C}}, and \texttt{false} otherwise. }

 

\subsection{\textcolor{Chapter }{IsRDistributive}}
\logpage{[ 56, 4, 4 ]}\nobreak
\hyperdef{L}{X79A5AEE786AED315}{}
{\noindent\textcolor{FuncColor}{$\triangleright$\enspace\texttt{IsRDistributive({\mdseries\slshape C})\index{IsRDistributive@\texttt{IsRDistributive}}
\label{IsRDistributive}
}\hfill{\scriptsize (property)}}\\


 is \texttt{true} if the relation $( a + b ) * c = ( a * c ) + ( b * c )$ holds for all elements $a$, $b$, $c$ in the collection \mbox{\texttt{\mdseries\slshape C}}, and \texttt{false} otherwise. }

 

\subsection{\textcolor{Chapter }{IsDistributive}}
\logpage{[ 56, 4, 5 ]}\nobreak
\hyperdef{L}{X86716D4F7B968604}{}
{\noindent\textcolor{FuncColor}{$\triangleright$\enspace\texttt{IsDistributive({\mdseries\slshape C})\index{IsDistributive@\texttt{IsDistributive}}
\label{IsDistributive}
}\hfill{\scriptsize (property)}}\\


 is \texttt{true} if the collection \mbox{\texttt{\mdseries\slshape C}} is both left and right distributive (see \texttt{IsLDistributive} (\ref{IsLDistributive}), \texttt{IsRDistributive} (\ref{IsRDistributive})), and \texttt{false} otherwise. 
\begin{Verbatim}[commandchars=!@|,fontsize=\small,frame=single,label=Example]
  !gapprompt@gap>| !gapinput@IsDistributive( Integers );|
  true
\end{Verbatim}
 }

 

\subsection{\textcolor{Chapter }{IsAnticommutative}}
\logpage{[ 56, 4, 6 ]}\nobreak
\hyperdef{L}{X82DECD237D49D937}{}
{\noindent\textcolor{FuncColor}{$\triangleright$\enspace\texttt{IsAnticommutative({\mdseries\slshape R})\index{IsAnticommutative@\texttt{IsAnticommutative}}
\label{IsAnticommutative}
}\hfill{\scriptsize (property)}}\\


 is \texttt{true} if the relation $a * b = - b * a$ holds for all elements $a$, $b$ in the ring \mbox{\texttt{\mdseries\slshape R}}, and \texttt{false} otherwise. }

 

\subsection{\textcolor{Chapter }{IsZeroSquaredRing}}
\logpage{[ 56, 4, 7 ]}\nobreak
\hyperdef{L}{X7EC0FEC88535E8CC}{}
{\noindent\textcolor{FuncColor}{$\triangleright$\enspace\texttt{IsZeroSquaredRing({\mdseries\slshape R})\index{IsZeroSquaredRing@\texttt{IsZeroSquaredRing}}
\label{IsZeroSquaredRing}
}\hfill{\scriptsize (property)}}\\


 is \texttt{true} if $a * a$ is the zero element of the ring \mbox{\texttt{\mdseries\slshape R}} for all $a$ in \mbox{\texttt{\mdseries\slshape R}}, and \texttt{false} otherwise. }

 

\subsection{\textcolor{Chapter }{IsJacobianRing}}
\logpage{[ 56, 4, 8 ]}\nobreak
\hyperdef{L}{X799BEF8581971A13}{}
{\noindent\textcolor{FuncColor}{$\triangleright$\enspace\texttt{IsJacobianRing({\mdseries\slshape R})\index{IsJacobianRing@\texttt{IsJacobianRing}}
\label{IsJacobianRing}
}\hfill{\scriptsize (property)}}\\


 is \texttt{true} if the Jacobi identity holds in the ring \mbox{\texttt{\mdseries\slshape R}}, and \texttt{false} otherwise. The Jacobi identity means that $x * (y * z) + z * (x * y) + y * (z * x)$ is the zero element of \mbox{\texttt{\mdseries\slshape R}}, for all elements $x$, $y$, $z$ in \mbox{\texttt{\mdseries\slshape R}}. 
\begin{Verbatim}[commandchars=!@|,fontsize=\small,frame=single,label=Example]
  !gapprompt@gap>| !gapinput@L:= FullMatrixLieAlgebra( GF( 5 ), 7 );|
  <Lie algebra over GF(5), with 13 generators>
  !gapprompt@gap>| !gapinput@IsJacobianRing( L );|
  true
\end{Verbatim}
 }

 }

  
\section{\textcolor{Chapter }{Units and Factorizations}}\label{Units and Factorizations}
\logpage{[ 56, 5, 0 ]}
\hyperdef{L}{X8130085978A9B3C4}{}
{
  

\subsection{\textcolor{Chapter }{IsUnit}}
\logpage{[ 56, 5, 1 ]}\nobreak
\hyperdef{L}{X85CBFBAE78DE72E8}{}
{\noindent\textcolor{FuncColor}{$\triangleright$\enspace\texttt{IsUnit({\mdseries\slshape [R, ]r})\index{IsUnit@\texttt{IsUnit}}
\label{IsUnit}
}\hfill{\scriptsize (operation)}}\\


 \texttt{IsUnit} returns \texttt{true} if \mbox{\texttt{\mdseries\slshape r}} is a unit in the ring \mbox{\texttt{\mdseries\slshape R}}, if given, and otherwise in its default ring (see \texttt{DefaultRing} (\ref{DefaultRing:for ring elements})). If \mbox{\texttt{\mdseries\slshape r}} is not a unit then \texttt{false} is returned. 

 An element \mbox{\texttt{\mdseries\slshape r}} is called a \emph{unit} in a ring \mbox{\texttt{\mdseries\slshape R}}, if \mbox{\texttt{\mdseries\slshape r}} has an inverse in \mbox{\texttt{\mdseries\slshape R}}. 

 \texttt{IsUnit} may call \texttt{Quotient} (\ref{Quotient}).  }

 

\subsection{\textcolor{Chapter }{Units}}
\logpage{[ 56, 5, 2 ]}\nobreak
\hyperdef{L}{X853C045B7BA6A580}{}
{\noindent\textcolor{FuncColor}{$\triangleright$\enspace\texttt{Units({\mdseries\slshape R})\index{Units@\texttt{Units}}
\label{Units}
}\hfill{\scriptsize (attribute)}}\\


 \texttt{Units} returns the group of units of the ring \mbox{\texttt{\mdseries\slshape R}}. This may either be returned as a list or as a group. 

 An element $r$ is called a \emph{unit} of a ring $R$ if $r$ has an inverse in $R$. It is easy to see that the set of units forms a multiplicative group. 

 
\begin{Verbatim}[commandchars=!@|,fontsize=\small,frame=single,label=Example]
  !gapprompt@gap>| !gapinput@Units( GaussianIntegers );|
  [ -1, 1, -E(4), E(4) ]
  !gapprompt@gap>| !gapinput@Units( GF( 16 ) );|
  <group of size 15 with 1 generator>
\end{Verbatim}
 }

 

\subsection{\textcolor{Chapter }{IsAssociated}}
\logpage{[ 56, 5, 3 ]}\nobreak
\hyperdef{L}{X7B307F217DDC7E20}{}
{\noindent\textcolor{FuncColor}{$\triangleright$\enspace\texttt{IsAssociated({\mdseries\slshape [R, ]r, s})\index{IsAssociated@\texttt{IsAssociated}}
\label{IsAssociated}
}\hfill{\scriptsize (operation)}}\\


 \texttt{IsAssociated} returns \texttt{true} if the two ring elements \mbox{\texttt{\mdseries\slshape r}} and \mbox{\texttt{\mdseries\slshape s}} are associated in the ring \mbox{\texttt{\mdseries\slshape R}}, if given, and otherwise in their default ring (see \texttt{DefaultRing} (\ref{DefaultRing:for ring elements})). If the two elements are not associated then \texttt{false} is returned. 

 Two elements \mbox{\texttt{\mdseries\slshape r}} and \mbox{\texttt{\mdseries\slshape s}} of a ring \mbox{\texttt{\mdseries\slshape R}} are called \emph{associated} if there is a unit $u$ of \mbox{\texttt{\mdseries\slshape R}} such that \mbox{\texttt{\mdseries\slshape r}} $u = $\mbox{\texttt{\mdseries\slshape s}}. }

 

\subsection{\textcolor{Chapter }{Associates}}
\logpage{[ 56, 5, 4 ]}\nobreak
\hyperdef{L}{X7A69C9097E17D161}{}
{\noindent\textcolor{FuncColor}{$\triangleright$\enspace\texttt{Associates({\mdseries\slshape [R, ]r})\index{Associates@\texttt{Associates}}
\label{Associates}
}\hfill{\scriptsize (operation)}}\\


 \texttt{Associates} returns the set of associates of \mbox{\texttt{\mdseries\slshape r}} in the ring \mbox{\texttt{\mdseries\slshape R}}, if given, and otherwise in its default ring (see \texttt{DefaultRing} (\ref{DefaultRing:for ring elements})). 

 Two elements \mbox{\texttt{\mdseries\slshape r}} and $s$ of a ring $R$ are called \emph{associated} if there is a unit $u$ of $R$ such that $\mbox{\texttt{\mdseries\slshape r}} u = s$. 

 
\begin{Verbatim}[commandchars=!@|,fontsize=\small,frame=single,label=Example]
  !gapprompt@gap>| !gapinput@Associates( Integers, 2 );|
  [ -2, 2 ]
  !gapprompt@gap>| !gapinput@Associates( GaussianIntegers, 2 );|
  [ -2, 2, -2*E(4), 2*E(4) ]
\end{Verbatim}
 }

 

\subsection{\textcolor{Chapter }{StandardAssociate}}
\logpage{[ 56, 5, 5 ]}\nobreak
\hyperdef{L}{X7B1A9A4C7C59FB36}{}
{\noindent\textcolor{FuncColor}{$\triangleright$\enspace\texttt{StandardAssociate({\mdseries\slshape [R, ]r})\index{StandardAssociate@\texttt{StandardAssociate}}
\label{StandardAssociate}
}\hfill{\scriptsize (operation)}}\\


 \texttt{StandardAssociate} returns the standard associate of the ring element \mbox{\texttt{\mdseries\slshape r}} in the ring \mbox{\texttt{\mdseries\slshape R}}, if given, and otherwise in its default ring (see \texttt{DefaultRing} (\ref{DefaultRing:for ring elements})). 

 The \emph{standard associate} of a ring element \mbox{\texttt{\mdseries\slshape r}} of \mbox{\texttt{\mdseries\slshape R}} is an associated element of \mbox{\texttt{\mdseries\slshape r}} which is, in a ring dependent way, distinguished among the set of associates
of \mbox{\texttt{\mdseries\slshape r}}. For example, in the ring of integers the standard associate is the absolute
value. 

 
\begin{Verbatim}[commandchars=!@|,fontsize=\small,frame=single,label=Example]
  !gapprompt@gap>| !gapinput@x:= Indeterminate( Rationals, "x" );;|
  !gapprompt@gap>| !gapinput@StandardAssociate( -x^2-x+1 );|
  x^2+x-1
\end{Verbatim}
 }

 

\subsection{\textcolor{Chapter }{StandardAssociateUnit}}
\logpage{[ 56, 5, 6 ]}\nobreak
\hyperdef{L}{X7EB6803C789E027D}{}
{\noindent\textcolor{FuncColor}{$\triangleright$\enspace\texttt{StandardAssociateUnit({\mdseries\slshape [R, ]r})\index{StandardAssociateUnit@\texttt{StandardAssociateUnit}}
\label{StandardAssociateUnit}
}\hfill{\scriptsize (operation)}}\\


 \texttt{StandardAssociateUnit} returns a unit in the ring \mbox{\texttt{\mdseries\slshape R}} such that the ring element \mbox{\texttt{\mdseries\slshape r}} times this unit equals the standard associate of \mbox{\texttt{\mdseries\slshape r}} in \mbox{\texttt{\mdseries\slshape R}}. 

 If \mbox{\texttt{\mdseries\slshape R}} is not given, the default ring of \mbox{\texttt{\mdseries\slshape r}} is used instead. (see \texttt{DefaultRing} (\ref{DefaultRing:for ring elements})). 

 

 
\begin{Verbatim}[commandchars=!@|,fontsize=\small,frame=single,label=Example]
  !gapprompt@gap>| !gapinput@y:= Indeterminate( Rationals, "y" );;|
  !gapprompt@gap>| !gapinput@r:= -y^2-y+1;|
  -y^2-y+1
  !gapprompt@gap>| !gapinput@StandardAssociateUnit( r );|
  -1
  !gapprompt@gap>| !gapinput@StandardAssociateUnit( r ) * r = StandardAssociate( r );|
  true
\end{Verbatim}
 }

 

\subsection{\textcolor{Chapter }{IsIrreducibleRingElement}}
\logpage{[ 56, 5, 7 ]}\nobreak
\hyperdef{L}{X7CD7C64A7D961A18}{}
{\noindent\textcolor{FuncColor}{$\triangleright$\enspace\texttt{IsIrreducibleRingElement({\mdseries\slshape [R, ]r})\index{IsIrreducibleRingElement@\texttt{IsIrreducibleRingElement}}
\label{IsIrreducibleRingElement}
}\hfill{\scriptsize (operation)}}\\


 \texttt{IsIrreducibleRingElement} returns \texttt{true} if the ring element \mbox{\texttt{\mdseries\slshape r}} is irreducible in the ring \mbox{\texttt{\mdseries\slshape R}}, if given, and otherwise in its default ring (see \texttt{DefaultRing} (\ref{DefaultRing:for ring elements})). If \mbox{\texttt{\mdseries\slshape r}} is not irreducible then \texttt{false} is returned. 

 An element \mbox{\texttt{\mdseries\slshape r}} of a ring \mbox{\texttt{\mdseries\slshape R}} is called \emph{irreducible} if \mbox{\texttt{\mdseries\slshape r}} is not a unit in \mbox{\texttt{\mdseries\slshape R}} and if there is no nontrivial factorization of \mbox{\texttt{\mdseries\slshape r}} in \mbox{\texttt{\mdseries\slshape R}}, i.e., if there is no representation of \mbox{\texttt{\mdseries\slshape r}} as product $s t$ such that neither $s$ nor $t$ is a unit (see \texttt{IsUnit} (\ref{IsUnit})). Each prime element (see \texttt{IsPrime} (\ref{IsPrime})) is irreducible. 
\begin{Verbatim}[commandchars=!@|,fontsize=\small,frame=single,label=Example]
  !gapprompt@gap>| !gapinput@IsIrreducibleRingElement( Integers, 2 );|
  true
\end{Verbatim}
 }

 

\subsection{\textcolor{Chapter }{IsPrime}}
\logpage{[ 56, 5, 8 ]}\nobreak
\hyperdef{L}{X7AA107AE7F79C6D8}{}
{\noindent\textcolor{FuncColor}{$\triangleright$\enspace\texttt{IsPrime({\mdseries\slshape [R, ]r})\index{IsPrime@\texttt{IsPrime}}
\label{IsPrime}
}\hfill{\scriptsize (operation)}}\\


 \texttt{IsPrime} returns \texttt{true} if the ring element \mbox{\texttt{\mdseries\slshape r}} is a prime in the ring \mbox{\texttt{\mdseries\slshape R}}, if given, and otherwise in its default ring (see \texttt{DefaultRing} (\ref{DefaultRing:for ring elements})). If \mbox{\texttt{\mdseries\slshape r}} is not a prime then \texttt{false} is returned. 

 An element \mbox{\texttt{\mdseries\slshape r}} of a ring \mbox{\texttt{\mdseries\slshape R}} is called \emph{prime} if for each pair $s$ and $t$ such that \mbox{\texttt{\mdseries\slshape r}} divides $s t$ the element \mbox{\texttt{\mdseries\slshape r}} divides either $s$ or $t$. Note that there are rings where not every irreducible element (see \texttt{IsIrreducibleRingElement} (\ref{IsIrreducibleRingElement})) is a prime. }

 

\subsection{\textcolor{Chapter }{Factors}}
\logpage{[ 56, 5, 9 ]}\nobreak
\hyperdef{L}{X82D6EDC685D12AE2}{}
{\noindent\textcolor{FuncColor}{$\triangleright$\enspace\texttt{Factors({\mdseries\slshape [R, ]r})\index{Factors@\texttt{Factors}}
\label{Factors}
}\hfill{\scriptsize (operation)}}\\


 \texttt{Factors} returns the factorization of the ring element \mbox{\texttt{\mdseries\slshape r}} in the ring \mbox{\texttt{\mdseries\slshape R}}, if given, and otherwise in its default ring (see \texttt{DefaultRing} (\ref{DefaultRing:for ring elements})). The factorization is returned as a list of primes (see \texttt{IsPrime} (\ref{IsPrime})). Each element in the list is a standard associate (see \texttt{StandardAssociate} (\ref{StandardAssociate})) except the first one, which is multiplied by a unit as necessary to have \texttt{Product( Factors( \mbox{\texttt{\mdseries\slshape R}}, \mbox{\texttt{\mdseries\slshape r}} ) ) = \mbox{\texttt{\mdseries\slshape r}}}. This list is usually also sorted, thus smallest prime factors come first. If \mbox{\texttt{\mdseries\slshape r}} is a unit or zero, \texttt{Factors( \mbox{\texttt{\mdseries\slshape R}}, \mbox{\texttt{\mdseries\slshape r}} ) = [ \mbox{\texttt{\mdseries\slshape r}} ]}. 

  
\begin{Verbatim}[commandchars=!@|,fontsize=\small,frame=single,label=Example]
  !gapprompt@gap>| !gapinput@x:= Indeterminate( GF(2), "x" );;|
  !gapprompt@gap>| !gapinput@pol:= x^2+x+1;|
  x^2+x+Z(2)^0
  !gapprompt@gap>| !gapinput@Factors( pol );|
  [ x^2+x+Z(2)^0 ]
  !gapprompt@gap>| !gapinput@Factors( PolynomialRing( GF(4) ), pol );|
  [ x+Z(2^2), x+Z(2^2)^2 ]
\end{Verbatim}
 }

 

\subsection{\textcolor{Chapter }{PadicValuation}}
\logpage{[ 56, 5, 10 ]}\nobreak
\hyperdef{L}{X8559CC7B80C479F1}{}
{\noindent\textcolor{FuncColor}{$\triangleright$\enspace\texttt{PadicValuation({\mdseries\slshape r, p})\index{PadicValuation@\texttt{PadicValuation}}
\label{PadicValuation}
}\hfill{\scriptsize (operation)}}\\


 \texttt{PadicValuation} is the operation to compute the \mbox{\texttt{\mdseries\slshape p}}\texttt{\symbol{45}}adic valuation of a ring element \mbox{\texttt{\mdseries\slshape r}}. }

 }

  
\section{\textcolor{Chapter }{Euclidean Rings}}\label{Euclidean Rings}
\logpage{[ 56, 6, 0 ]}
\hyperdef{L}{X7F12BB99865EB7BF}{}
{
  

\subsection{\textcolor{Chapter }{IsEuclideanRing}}
\logpage{[ 56, 6, 1 ]}\nobreak
\hyperdef{L}{X808B8E8E80D48E4A}{}
{\noindent\textcolor{FuncColor}{$\triangleright$\enspace\texttt{IsEuclideanRing({\mdseries\slshape R})\index{IsEuclideanRing@\texttt{IsEuclideanRing}}
\label{IsEuclideanRing}
}\hfill{\scriptsize (Category)}}\\


 A ring $R$ is called a Euclidean ring if it is a non\texttt{\symbol{45}}trivial
commutative ring and there exists a function $\delta$, called the Euclidean degree, from $R-\{0_R\}$ into a well\texttt{\symbol{45}}ordered set (such as the nonnegative integers),
such that for every pair $r \in R$ and $s \in R-\{0_R\}$ there exists an element $q$ such that either $r - q s = 0_R$ or $\delta(r - q s) < \delta( s )$. In \textsf{GAP} the Euclidean degree $\delta$ is implicitly built into a ring and cannot be changed. The existence of this
division with remainder implies that the Euclidean algorithm can be applied to
compute a greatest common divisor of two elements, which in turn implies that $R$ is a unique factorization ring. 

  
\begin{Verbatim}[commandchars=!@|,fontsize=\small,frame=single,label=Example]
  !gapprompt@gap>| !gapinput@IsEuclideanRing( GaussianIntegers );|
  true
\end{Verbatim}
 }

 

\subsection{\textcolor{Chapter }{EuclideanDegree}}
\logpage{[ 56, 6, 2 ]}\nobreak
\hyperdef{L}{X784234088350D4E4}{}
{\noindent\textcolor{FuncColor}{$\triangleright$\enspace\texttt{EuclideanDegree({\mdseries\slshape [R, ]r})\index{EuclideanDegree@\texttt{EuclideanDegree}}
\label{EuclideanDegree}
}\hfill{\scriptsize (operation)}}\\


 \texttt{EuclideanDegree} returns the Euclidean degree of the ring element \mbox{\texttt{\mdseries\slshape r}} in the ring \mbox{\texttt{\mdseries\slshape R}}, if given, and otherwise in its default ring (see \texttt{DefaultRing} (\ref{DefaultRing:for ring elements})). 

 The ring \mbox{\texttt{\mdseries\slshape R}} must be a Euclidean ring (see \texttt{IsEuclideanRing} (\ref{IsEuclideanRing})). 
\begin{Verbatim}[commandchars=!@|,fontsize=\small,frame=single,label=Example]
  !gapprompt@gap>| !gapinput@EuclideanDegree( GaussianIntegers, 3 );|
  9
\end{Verbatim}
 }

 

\subsection{\textcolor{Chapter }{EuclideanQuotient}}
\logpage{[ 56, 6, 3 ]}\nobreak
\hyperdef{L}{X7A93FA788318B147}{}
{\noindent\textcolor{FuncColor}{$\triangleright$\enspace\texttt{EuclideanQuotient({\mdseries\slshape [R, ]r, m})\index{EuclideanQuotient@\texttt{EuclideanQuotient}}
\label{EuclideanQuotient}
}\hfill{\scriptsize (operation)}}\\


 \texttt{EuclideanQuotient} returns the Euclidean quotient of the ring elements \mbox{\texttt{\mdseries\slshape r}} and \mbox{\texttt{\mdseries\slshape m}} in the ring \mbox{\texttt{\mdseries\slshape R}}, if given, and otherwise in their default ring (see \texttt{DefaultRing} (\ref{DefaultRing:for ring elements})). 

 The ring \mbox{\texttt{\mdseries\slshape R}} must be a Euclidean ring (see \texttt{IsEuclideanRing} (\ref{IsEuclideanRing})), otherwise an error is signalled. 
\begin{Verbatim}[commandchars=!@|,fontsize=\small,frame=single,label=Example]
  !gapprompt@gap>| !gapinput@EuclideanQuotient( 8, 3 );|
  2
\end{Verbatim}
 }

 

\subsection{\textcolor{Chapter }{EuclideanRemainder}}
\logpage{[ 56, 6, 4 ]}\nobreak
\hyperdef{L}{X7B5E9639865E91BA}{}
{\noindent\textcolor{FuncColor}{$\triangleright$\enspace\texttt{EuclideanRemainder({\mdseries\slshape [R, ]r, m})\index{EuclideanRemainder@\texttt{EuclideanRemainder}}
\label{EuclideanRemainder}
}\hfill{\scriptsize (operation)}}\\


 \texttt{EuclideanRemainder} returns the Euclidean remainder of the ring element \mbox{\texttt{\mdseries\slshape r}} modulo the ring element \mbox{\texttt{\mdseries\slshape m}} in the ring \mbox{\texttt{\mdseries\slshape R}}, if given, and otherwise in their default ring (see \texttt{DefaultRing} (\ref{DefaultRing:for ring elements})). 

 The ring \mbox{\texttt{\mdseries\slshape R}} must be a Euclidean ring (see \texttt{IsEuclideanRing} (\ref{IsEuclideanRing})), otherwise an error is signalled. 
\begin{Verbatim}[commandchars=!@|,fontsize=\small,frame=single,label=Example]
  !gapprompt@gap>| !gapinput@EuclideanRemainder( 8, 3 );|
  2
\end{Verbatim}
 }

 

\subsection{\textcolor{Chapter }{QuotientRemainder}}
\logpage{[ 56, 6, 5 ]}\nobreak
\hyperdef{L}{X876B7532801A1B35}{}
{\noindent\textcolor{FuncColor}{$\triangleright$\enspace\texttt{QuotientRemainder({\mdseries\slshape [R, ]r, m})\index{QuotientRemainder@\texttt{QuotientRemainder}}
\label{QuotientRemainder}
}\hfill{\scriptsize (operation)}}\\


 \texttt{QuotientRemainder} returns the Euclidean quotient and the Euclidean remainder of the ring
elements \mbox{\texttt{\mdseries\slshape r}} and \mbox{\texttt{\mdseries\slshape m}} in the ring \mbox{\texttt{\mdseries\slshape R}}, if given, and otherwise in their default ring (see \texttt{DefaultRing} (\ref{DefaultRing:for ring elements})). The result is a pair of ring elements. 

 The ring \mbox{\texttt{\mdseries\slshape R}} must be a Euclidean ring (see \texttt{IsEuclideanRing} (\ref{IsEuclideanRing})), otherwise an error is signalled. 
\begin{Verbatim}[commandchars=!@|,fontsize=\small,frame=single,label=Example]
  !gapprompt@gap>| !gapinput@QuotientRemainder( GaussianIntegers, 8, 3 );|
  [ 3, -1 ]
\end{Verbatim}
 }

 }

  
\section{\textcolor{Chapter }{Gcd and Lcm}}\label{Gcd and Lcm}
\logpage{[ 56, 7, 0 ]}
\hyperdef{L}{X7E9CF2C07C4A6CEE}{}
{
  
\subsection{\textcolor{Chapter }{Gcd}}\logpage{[ 56, 7, 1 ]}
\hyperdef{L}{X7DE207718456F98F}{}
{
\noindent\textcolor{FuncColor}{$\triangleright$\enspace\texttt{Gcd({\mdseries\slshape [R, ]r1, r2, ...})\index{Gcd@\texttt{Gcd}!for (a ring and) several elements}
\label{Gcd:for (a ring and) several elements}
}\hfill{\scriptsize (function)}}\\
\noindent\textcolor{FuncColor}{$\triangleright$\enspace\texttt{Gcd({\mdseries\slshape [R, ]list})\index{Gcd@\texttt{Gcd}!for (a ring and) a list of elements}
\label{Gcd:for (a ring and) a list of elements}
}\hfill{\scriptsize (function)}}\\


 \texttt{Gcd} returns the greatest common divisor of the ring elements \mbox{\texttt{\mdseries\slshape r1}}, \mbox{\texttt{\mdseries\slshape r2}}, $\ldots$ resp. of the ring elements in the list \mbox{\texttt{\mdseries\slshape list}} in the ring \mbox{\texttt{\mdseries\slshape R}}, if given, and otherwise in their default ring, see \texttt{DefaultRing} (\ref{DefaultRing:for ring elements}). 

 \texttt{Gcd} returns the standard associate (see \texttt{StandardAssociate} (\ref{StandardAssociate})) of the greatest common divisors. 

 A divisor of an element $r$ in the ring $R$ is an element $d\in R$ such that $r$ is a multiple of $d$. A common divisor of the elements $r_1, r_2, \ldots$ in the ring $R$ is an element $d\in R$ which is a divisor of each $r_1, r_2, \ldots$. A greatest common divisor $d$ in addition has the property that every other common divisor of $r_1, r_2, \ldots$ is a divisor of $d$. 

 Note that this in particular implies the following: For the zero element $z$ of \mbox{\texttt{\mdseries\slshape R}}, we have \texttt{Gcd( \mbox{\texttt{\mdseries\slshape r}}, }$z$\texttt{ ) = Gcd( }$z$\texttt{, \mbox{\texttt{\mdseries\slshape r}} ) = StandardAssociate( \mbox{\texttt{\mdseries\slshape r}} )} and \texttt{Gcd( }$z$\texttt{, }$z$\texttt{ ) = }$z$. 
\begin{Verbatim}[commandchars=!@|,fontsize=\small,frame=single,label=Example]
  !gapprompt@gap>| !gapinput@Gcd( Integers, [ 10, 15 ] );|
  5
\end{Verbatim}
 }

 

\subsection{\textcolor{Chapter }{GcdOp}}
\logpage{[ 56, 7, 2 ]}\nobreak
\hyperdef{L}{X7836D50F8341D6E1}{}
{\noindent\textcolor{FuncColor}{$\triangleright$\enspace\texttt{GcdOp({\mdseries\slshape [R, ]r, s})\index{GcdOp@\texttt{GcdOp}}
\label{GcdOp}
}\hfill{\scriptsize (operation)}}\\


 \texttt{GcdOp} is the operation to compute the greatest common divisor of two ring elements \mbox{\texttt{\mdseries\slshape r}}, \mbox{\texttt{\mdseries\slshape s}} in the ring \mbox{\texttt{\mdseries\slshape R}} or in their default ring. }

 
\subsection{\textcolor{Chapter }{GcdRepresentation}}\logpage{[ 56, 7, 3 ]}
\hyperdef{L}{X7ABB91EF838075EF}{}
{
\noindent\textcolor{FuncColor}{$\triangleright$\enspace\texttt{GcdRepresentation({\mdseries\slshape [R, ]r1, r2, ...})\index{GcdRepresentation@\texttt{GcdRepresentation}!for (a ring and) several elements}
\label{GcdRepresentation:for (a ring and) several elements}
}\hfill{\scriptsize (function)}}\\
\noindent\textcolor{FuncColor}{$\triangleright$\enspace\texttt{GcdRepresentation({\mdseries\slshape [R, ]list})\index{GcdRepresentation@\texttt{GcdRepresentation}!for (a ring and) a list of elements}
\label{GcdRepresentation:for (a ring and) a list of elements}
}\hfill{\scriptsize (function)}}\\


 \texttt{GcdRepresentation} returns a representation of the greatest common divisor of the ring elements \mbox{\texttt{\mdseries\slshape r1}}, \mbox{\texttt{\mdseries\slshape r2}}, $\ldots$ resp. of the ring elements in the list \mbox{\texttt{\mdseries\slshape list}} in the Euclidean ring \mbox{\texttt{\mdseries\slshape R}}, if given, and otherwise in their default ring, see \texttt{DefaultRing} (\ref{DefaultRing:for ring elements}). 

 A representation of the gcd $g$ of the elements $r_1, r_2, \ldots$ of a ring $R$ is a list of ring elements $s_1, s_2, \ldots$ of $R$, such that $g = s_1 r_1 + s_2 r_2 + \cdots$. Such representations do not exist in all rings, but they do exist in
Euclidean rings (see \texttt{IsEuclideanRing} (\ref{IsEuclideanRing})), which can be shown using the Euclidean algorithm, which in fact can compute
those coefficients. 
\begin{Verbatim}[commandchars=!@|,fontsize=\small,frame=single,label=Example]
  !gapprompt@gap>| !gapinput@a:= Indeterminate( Rationals, "a" );;|
  !gapprompt@gap>| !gapinput@GcdRepresentation( a^2+1, a^3+1 );|
  [ -1/2*a^2-1/2*a+1/2, 1/2*a+1/2 ]
\end{Verbatim}
 

 \texttt{Gcdex} (\ref{Gcdex}) provides similar functionality over the integers. }

 

\subsection{\textcolor{Chapter }{GcdRepresentationOp}}
\logpage{[ 56, 7, 4 ]}\nobreak
\hyperdef{L}{X81392E7F84956341}{}
{\noindent\textcolor{FuncColor}{$\triangleright$\enspace\texttt{GcdRepresentationOp({\mdseries\slshape [R, ]r, s})\index{GcdRepresentationOp@\texttt{GcdRepresentationOp}}
\label{GcdRepresentationOp}
}\hfill{\scriptsize (operation)}}\\


 \texttt{GcdRepresentationOp} is the operation to compute the representation of the greatest common divisor
of two ring elements \mbox{\texttt{\mdseries\slshape r}}, \mbox{\texttt{\mdseries\slshape s}} in the Euclidean ring \mbox{\texttt{\mdseries\slshape R}} or in their default ring, respectively. }

 

\subsection{\textcolor{Chapter }{ShowGcd}}
\logpage{[ 56, 7, 5 ]}\nobreak
\hyperdef{L}{X836DB8B47A0219FB}{}
{\noindent\textcolor{FuncColor}{$\triangleright$\enspace\texttt{ShowGcd({\mdseries\slshape a, b})\index{ShowGcd@\texttt{ShowGcd}}
\label{ShowGcd}
}\hfill{\scriptsize (function)}}\\


 This function takes two elements \mbox{\texttt{\mdseries\slshape a}} and \mbox{\texttt{\mdseries\slshape b}} of an Euclidean ring and returns their greatest common divisor. It will print
out the steps performed by the Euclidean algorithm, as well as the
rearrangement of these steps to express the gcd as a ring combination of \mbox{\texttt{\mdseries\slshape a}} and \mbox{\texttt{\mdseries\slshape b}}. 
\begin{Verbatim}[commandchars=!@|,fontsize=\small,frame=single,label=Example]
  !gapprompt@gap>| !gapinput@ShowGcd(192,42);|
  192=4*42 + 24
  42=1*24 + 18
  24=1*18 + 6
  18=3*6 + 0
  The Gcd is 6
   = 1*24 -1*18
   = -1*42 + 2*24
   = 2*192 -9*42
  6
\end{Verbatim}
 }

 
\subsection{\textcolor{Chapter }{Lcm}}\logpage{[ 56, 7, 6 ]}
\hyperdef{L}{X7ABA92057DD6C7AF}{}
{
\noindent\textcolor{FuncColor}{$\triangleright$\enspace\texttt{Lcm({\mdseries\slshape [R, ]r1, r2, ...})\index{Lcm@\texttt{Lcm}!for (a ring and) several elements}
\label{Lcm:for (a ring and) several elements}
}\hfill{\scriptsize (function)}}\\
\noindent\textcolor{FuncColor}{$\triangleright$\enspace\texttt{Lcm({\mdseries\slshape [R, ]list})\index{Lcm@\texttt{Lcm}!for (a ring and) a list of elements}
\label{Lcm:for (a ring and) a list of elements}
}\hfill{\scriptsize (function)}}\\


 \texttt{Lcm} returns the least common multiple of the ring elements \mbox{\texttt{\mdseries\slshape r1}}, \mbox{\texttt{\mdseries\slshape r2}}, $\ldots$ resp. of the ring elements in the list \mbox{\texttt{\mdseries\slshape list}} in the ring \mbox{\texttt{\mdseries\slshape R}}, if given, and otherwise in their default ring, see \texttt{DefaultRing} (\ref{DefaultRing:for ring elements}). 

 \texttt{Lcm} returns the standard associate (see{\nobreakspace}\texttt{StandardAssociate} (\ref{StandardAssociate})) of the least common multiples. 

 A least common multiple of the elements $r_1, r_2, \ldots$ of the ring $R$ is an element $m$ that is a multiple of $r_1, r_2, \ldots$, and every other multiple of these elements is a multiple of $m$. 

 Note that this in particular implies the following: For the zero element $z$ of \mbox{\texttt{\mdseries\slshape R}}, we have \texttt{Lcm( \mbox{\texttt{\mdseries\slshape r}}, }$z$\texttt{ ) = Lcm( }$z$\texttt{, \mbox{\texttt{\mdseries\slshape r}} ) = StandardAssociate( \mbox{\texttt{\mdseries\slshape r}} )} and \texttt{Lcm( }$z$\texttt{, }$z$\texttt{ ) = }$z$. }

 

\subsection{\textcolor{Chapter }{LcmOp}}
\logpage{[ 56, 7, 7 ]}\nobreak
\hyperdef{L}{X7FB6C5A67AC1E8C1}{}
{\noindent\textcolor{FuncColor}{$\triangleright$\enspace\texttt{LcmOp({\mdseries\slshape [R, ]r, s})\index{LcmOp@\texttt{LcmOp}}
\label{LcmOp}
}\hfill{\scriptsize (operation)}}\\


 \texttt{LcmOp} is the operation to compute the least common multiple of two ring elements \mbox{\texttt{\mdseries\slshape r}}, \mbox{\texttt{\mdseries\slshape s}} in the ring \mbox{\texttt{\mdseries\slshape R}} or in their default ring, respectively. 

 The default methods for this uses the equality $lcm( m, n ) = m*n / gcd( m, n )$ (see{\nobreakspace}\texttt{GcdOp} (\ref{GcdOp})). }

 

\subsection{\textcolor{Chapter }{QuotientMod}}
\logpage{[ 56, 7, 8 ]}\nobreak
\hyperdef{L}{X8555913A83D716A4}{}
{\noindent\textcolor{FuncColor}{$\triangleright$\enspace\texttt{QuotientMod({\mdseries\slshape [R, ]r, s, m})\index{QuotientMod@\texttt{QuotientMod}}
\label{QuotientMod}
}\hfill{\scriptsize (operation)}}\\


 \texttt{QuotientMod} returns a quotient of the ring elements \mbox{\texttt{\mdseries\slshape r}} and \mbox{\texttt{\mdseries\slshape s}} modulo the ring element \mbox{\texttt{\mdseries\slshape m}} in the ring \mbox{\texttt{\mdseries\slshape R}}, if given, and otherwise in their default ring, see \texttt{DefaultRing} (\ref{DefaultRing:for ring elements}). 

 \mbox{\texttt{\mdseries\slshape R}} must be a Euclidean ring (see \texttt{IsEuclideanRing} (\ref{IsEuclideanRing})) so that \texttt{EuclideanRemainder} (\ref{EuclideanRemainder}) can be applied. If no modular quotient exists, \texttt{fail} is returned. 

 A quotient $q$ of \mbox{\texttt{\mdseries\slshape r}} and \mbox{\texttt{\mdseries\slshape s}} modulo \mbox{\texttt{\mdseries\slshape m}} is an element of \mbox{\texttt{\mdseries\slshape R}} such that $q \mbox{\texttt{\mdseries\slshape s}} = \mbox{\texttt{\mdseries\slshape r}}$ modulo $m$, i.e., such that $q \mbox{\texttt{\mdseries\slshape s}} - \mbox{\texttt{\mdseries\slshape r}}$ is divisible by \mbox{\texttt{\mdseries\slshape m}} in \mbox{\texttt{\mdseries\slshape R}} and that $q$ is either zero (if \mbox{\texttt{\mdseries\slshape r}} is divisible by \mbox{\texttt{\mdseries\slshape m}}) or the Euclidean degree of $q$ is strictly smaller than the Euclidean degree of \mbox{\texttt{\mdseries\slshape m}}. 
\begin{Verbatim}[commandchars=!@|,fontsize=\small,frame=single,label=Example]
  !gapprompt@gap>| !gapinput@QuotientMod( 7, 2, 3 );|
  2
\end{Verbatim}
 }

 

\subsection{\textcolor{Chapter }{PowerMod}}
\logpage{[ 56, 7, 9 ]}\nobreak
\hyperdef{L}{X805A35D684B7A952}{}
{\noindent\textcolor{FuncColor}{$\triangleright$\enspace\texttt{PowerMod({\mdseries\slshape [R, ]r, e, m})\index{PowerMod@\texttt{PowerMod}}
\label{PowerMod}
}\hfill{\scriptsize (operation)}}\\


 \texttt{PowerMod} returns the \mbox{\texttt{\mdseries\slshape e}}\texttt{\symbol{45}}th power of the ring element \mbox{\texttt{\mdseries\slshape r}} modulo the ring element \mbox{\texttt{\mdseries\slshape m}} in the ring \mbox{\texttt{\mdseries\slshape R}}, if given, and otherwise in their default ring, see \texttt{DefaultRing} (\ref{DefaultRing:for ring elements}). \mbox{\texttt{\mdseries\slshape e}} must be an integer. 

 \mbox{\texttt{\mdseries\slshape R}} must be a Euclidean ring (see \texttt{IsEuclideanRing} (\ref{IsEuclideanRing})) so that \texttt{EuclideanRemainder} (\ref{EuclideanRemainder}) can be applied to its elements. 

 If \mbox{\texttt{\mdseries\slshape e}} is positive the result is \mbox{\texttt{\mdseries\slshape r}}\texttt{\texttt{\symbol{94}}}\mbox{\texttt{\mdseries\slshape e}} modulo \mbox{\texttt{\mdseries\slshape m}}. If \mbox{\texttt{\mdseries\slshape e}} is negative then \texttt{PowerMod} first tries to find the inverse of \mbox{\texttt{\mdseries\slshape r}} modulo \mbox{\texttt{\mdseries\slshape m}}, i.e., $i$ such that $i \mbox{\texttt{\mdseries\slshape r}} = 1$ modulo \mbox{\texttt{\mdseries\slshape m}}. If the inverse does not exist an error is signalled. If the inverse does
exist \texttt{PowerMod} returns \texttt{PowerMod( \mbox{\texttt{\mdseries\slshape R}}, \mbox{\texttt{\mdseries\slshape i}}, \texttt{\symbol{45}}\mbox{\texttt{\mdseries\slshape e}}, \mbox{\texttt{\mdseries\slshape m}} )}. 

 \texttt{PowerMod} reduces the intermediate values modulo \mbox{\texttt{\mdseries\slshape m}}, improving performance drastically when \mbox{\texttt{\mdseries\slshape e}} is large and \mbox{\texttt{\mdseries\slshape m}} small. 
\begin{Verbatim}[commandchars=!@|,fontsize=\small,frame=single,label=Example]
  !gapprompt@gap>| !gapinput@PowerMod( 12, 100000, 7 );|
  2
\end{Verbatim}
 }

 

\subsection{\textcolor{Chapter }{InterpolatedPolynomial}}
\logpage{[ 56, 7, 10 ]}\nobreak
\hyperdef{L}{X87711E6F8024A358}{}
{\noindent\textcolor{FuncColor}{$\triangleright$\enspace\texttt{InterpolatedPolynomial({\mdseries\slshape R, x, y})\index{InterpolatedPolynomial@\texttt{InterpolatedPolynomial}}
\label{InterpolatedPolynomial}
}\hfill{\scriptsize (operation)}}\\


 \texttt{InterpolatedPolynomial} returns, for given lists \mbox{\texttt{\mdseries\slshape x}}, \mbox{\texttt{\mdseries\slshape y}} of elements in a ring \mbox{\texttt{\mdseries\slshape R}} of the same length $n$ the unique polynomial of degree less than $n$ which has value \mbox{\texttt{\mdseries\slshape y}}[$i$] at \mbox{\texttt{\mdseries\slshape x}}$[i]$, for all $i \in \{ 1, \ldots, n \}$. Note that the elements in \mbox{\texttt{\mdseries\slshape x}} must be distinct. 
\begin{Verbatim}[commandchars=!@|,fontsize=\small,frame=single,label=Example]
  !gapprompt@gap>| !gapinput@InterpolatedPolynomial( Integers, [ 1, 2, 3 ], [ 5, 7, 0 ] );|
  -9/2*x^2+31/2*x-6
\end{Verbatim}
 }

 }

  
\section{\textcolor{Chapter }{Homomorphisms of Rings}}\label{Homomorphisms of Rings}
\logpage{[ 56, 8, 0 ]}
\hyperdef{L}{X7B13484581169439}{}
{
  A \emph{ring homomorphism} is a mapping between two rings that respects addition and multiplication. 

 Currently \textsf{GAP} supports ring homomorphisms between finite rings (using straightforward
methods) and ring homomorphisms with additional structures, where source and
range are in fact algebras and where also the linear structure is respected,
see \ref{Homomorphisms of Algebras}. 

\subsection{\textcolor{Chapter }{RingGeneralMappingByImages}}
\logpage{[ 56, 8, 1 ]}\nobreak
\hyperdef{L}{X7DE9CC5B877C91DA}{}
{\noindent\textcolor{FuncColor}{$\triangleright$\enspace\texttt{RingGeneralMappingByImages({\mdseries\slshape R, S, gens, imgs})\index{RingGeneralMappingByImages@\texttt{RingGeneralMappingByImages}}
\label{RingGeneralMappingByImages}
}\hfill{\scriptsize (operation)}}\\


 is a general mapping from the ring \mbox{\texttt{\mdseries\slshape A}} to the ring \mbox{\texttt{\mdseries\slshape S}}. This general mapping is defined by mapping the entries in the list \mbox{\texttt{\mdseries\slshape gens}} (elements of \mbox{\texttt{\mdseries\slshape R}}) to the entries in the list \mbox{\texttt{\mdseries\slshape imgs}} (elements of \mbox{\texttt{\mdseries\slshape S}}), and taking the additive and multiplicative closure. 

 \mbox{\texttt{\mdseries\slshape gens}} need not generate \mbox{\texttt{\mdseries\slshape R}} as a ring, and if the specification does not define an additive and
multiplicative mapping then the result will be multivalued. Hence, in general
it is not a mapping. }

 

\subsection{\textcolor{Chapter }{RingHomomorphismByImages}}
\logpage{[ 56, 8, 2 ]}\nobreak
\hyperdef{L}{X78C1016284F08026}{}
{\noindent\textcolor{FuncColor}{$\triangleright$\enspace\texttt{RingHomomorphismByImages({\mdseries\slshape R, S, gens, imgs})\index{RingHomomorphismByImages@\texttt{RingHomomorphismByImages}}
\label{RingHomomorphismByImages}
}\hfill{\scriptsize (function)}}\\


 \texttt{RingHomomorphismByImages} returns the ring homomorphism with source \mbox{\texttt{\mdseries\slshape R}} and range \mbox{\texttt{\mdseries\slshape S}} that is defined by mapping the list \mbox{\texttt{\mdseries\slshape gens}} of generators of \mbox{\texttt{\mdseries\slshape R}} to the list \mbox{\texttt{\mdseries\slshape imgs}} of images in \mbox{\texttt{\mdseries\slshape S}}. 

 If \mbox{\texttt{\mdseries\slshape gens}} does not generate \mbox{\texttt{\mdseries\slshape R}} or if the homomorphism does not exist (i.e., if mapping the generators
describes only a multi\texttt{\symbol{45}}valued mapping) then \texttt{fail} is returned. 

 One can avoid the checks by calling \texttt{RingHomomorphismByImagesNC} (\ref{RingHomomorphismByImagesNC}), and one can construct multi\texttt{\symbol{45}}valued mappings with \texttt{RingGeneralMappingByImages} (\ref{RingGeneralMappingByImages}). }

 

\subsection{\textcolor{Chapter }{RingHomomorphismByImagesNC}}
\logpage{[ 56, 8, 3 ]}\nobreak
\hyperdef{L}{X7D01646A7CCBEDBB}{}
{\noindent\textcolor{FuncColor}{$\triangleright$\enspace\texttt{RingHomomorphismByImagesNC({\mdseries\slshape R, S, gens, imgs})\index{RingHomomorphismByImagesNC@\texttt{RingHomomorphismByImagesNC}}
\label{RingHomomorphismByImagesNC}
}\hfill{\scriptsize (operation)}}\\


 \texttt{RingHomomorphismByImagesNC} is the operation that is called by the function \texttt{RingHomomorphismByImages} (\ref{RingHomomorphismByImages}). Its methods may assume that \mbox{\texttt{\mdseries\slshape gens}} generates \mbox{\texttt{\mdseries\slshape R}} as a ring and that the mapping of \mbox{\texttt{\mdseries\slshape gens}} to \mbox{\texttt{\mdseries\slshape imgs}} defines a ring homomorphism. Results are unpredictable if these conditions do
not hold. 

 For creating a possibly multi\texttt{\symbol{45}}valued mapping from \mbox{\texttt{\mdseries\slshape R}} to \mbox{\texttt{\mdseries\slshape S}} that respects addition and multiplication, \texttt{RingGeneralMappingByImages} (\ref{RingGeneralMappingByImages}) can be used. }

  

\subsection{\textcolor{Chapter }{NaturalHomomorphismByIdeal}}
\logpage{[ 56, 8, 4 ]}\nobreak
\hyperdef{L}{X83D53D98809EC461}{}
{\noindent\textcolor{FuncColor}{$\triangleright$\enspace\texttt{NaturalHomomorphismByIdeal({\mdseries\slshape R, I})\index{NaturalHomomorphismByIdeal@\texttt{NaturalHomomorphismByIdeal}}
\label{NaturalHomomorphismByIdeal}
}\hfill{\scriptsize (operation)}}\\


 is the homomorphism of rings provided by the natural projection map of \mbox{\texttt{\mdseries\slshape R}} onto the quotient ring \mbox{\texttt{\mdseries\slshape R}}/\mbox{\texttt{\mdseries\slshape I}}. This map can be used to take pre\texttt{\symbol{45}}images in the original
ring from elements in the quotient. }

 }

  
\section{\textcolor{Chapter }{Small Rings}}\label{Small Rings}
\logpage{[ 56, 9, 0 ]}
\hyperdef{L}{X81D526A57B375AAD}{}
{
  \textsf{GAP} contains a library of small (order up to 15) rings. 

\subsection{\textcolor{Chapter }{SmallRing}}
\logpage{[ 56, 9, 1 ]}\nobreak
\hyperdef{L}{X7E86DCB7812DF04C}{}
{\noindent\textcolor{FuncColor}{$\triangleright$\enspace\texttt{SmallRing({\mdseries\slshape s, n})\index{SmallRing@\texttt{SmallRing}}
\label{SmallRing}
}\hfill{\scriptsize (function)}}\\


 returns the $n$\texttt{\symbol{45}}th ring of order $s$ from a library of rings of small order (up to isomorphism). 
\begin{Verbatim}[commandchars=!@A,fontsize=\small,frame=single,label=Example]
  !gapprompt@gap>A !gapinput@R:=SmallRing(8,37);A
  <ring with 3 generators>
  !gapprompt@gap>A !gapinput@ShowMultiplicationTable(R);A
  *     | 0*a   c     b     b+c   a     a+c   a+b   a+b+c
  ------+------------------------------------------------
  0*a   | 0*a   0*a   0*a   0*a   0*a   0*a   0*a   0*a
  c     | 0*a   0*a   0*a   0*a   0*a   0*a   0*a   0*a
  b     | 0*a   0*a   0*a   0*a   b     b     b     b
  b+c   | 0*a   0*a   0*a   0*a   b     b     b     b
  a     | 0*a   c     b     b+c   a+b   a+b+c a     a+c
  a+c   | 0*a   c     b     b+c   a+b   a+b+c a     a+c
  a+b   | 0*a   c     b     b+c   a     a+c   a+b   a+b+c
  a+b+c | 0*a   c     b     b+c   a     a+c   a+b   a+b+c
\end{Verbatim}
 }

 

\subsection{\textcolor{Chapter }{NumberSmallRings}}
\logpage{[ 56, 9, 2 ]}\nobreak
\hyperdef{L}{X7F2EE9AF83DCE641}{}
{\noindent\textcolor{FuncColor}{$\triangleright$\enspace\texttt{NumberSmallRings({\mdseries\slshape s})\index{NumberSmallRings@\texttt{NumberSmallRings}}
\label{NumberSmallRings}
}\hfill{\scriptsize (function)}}\\


 returns the number of (nonisomorphic) rings of order $s$ stored in the library of small rings. 
\begin{Verbatim}[commandchars=!@|,fontsize=\small,frame=single,label=Example]
  !gapprompt@gap>| !gapinput@List([1..15],NumberSmallRings);|
  [ 1, 2, 2, 11, 2, 4, 2, 52, 11, 4, 2, 22, 2, 4, 4 ]
\end{Verbatim}
 }

 

\subsection{\textcolor{Chapter }{Subrings}}
\logpage{[ 56, 9, 3 ]}\nobreak
\hyperdef{L}{X8070D20B86148929}{}
{\noindent\textcolor{FuncColor}{$\triangleright$\enspace\texttt{Subrings({\mdseries\slshape R})\index{Subrings@\texttt{Subrings}}
\label{Subrings}
}\hfill{\scriptsize (attribute)}}\\


 for a finite ring \mbox{\texttt{\mdseries\slshape R}} this function returns a list of all subrings of \mbox{\texttt{\mdseries\slshape R}}. 
\begin{Verbatim}[commandchars=!@|,fontsize=\small,frame=single,label=Example]
  !gapprompt@gap>| !gapinput@Subrings(SmallRing(8,37));|
  [ <ring with 1 generator>, <ring with 1 generator>,
    <ring with 1 generator>, <ring with 1 generator>,
    <ring with 1 generator>, <ring with 1 generator>,
    <ring with 2 generators>, <ring with 2 generators>,
    <ring with 2 generators>, <ring with 2 generators>,
    <ring with 3 generators> ]
\end{Verbatim}
 }

 

\subsection{\textcolor{Chapter }{Ideals}}
\logpage{[ 56, 9, 4 ]}\nobreak
\hyperdef{L}{X83629803819C4A6F}{}
{\noindent\textcolor{FuncColor}{$\triangleright$\enspace\texttt{Ideals({\mdseries\slshape R})\index{Ideals@\texttt{Ideals}}
\label{Ideals}
}\hfill{\scriptsize (attribute)}}\\


 for a finite ring \mbox{\texttt{\mdseries\slshape R}} this function returns a list of all ideals of \mbox{\texttt{\mdseries\slshape R}}. 
\begin{Verbatim}[commandchars=!@|,fontsize=\small,frame=single,label=Example]
  !gapprompt@gap>| !gapinput@Ideals(SmallRing(8,37));|
  [ <ring with 1 generator>, <ring with 1 generator>,
    <ring with 1 generator>, <ring with 2 generators>,
    <ring with 3 generators> ]
\end{Verbatim}
 }

 

\subsection{\textcolor{Chapter }{DirectSum}}
\logpage{[ 56, 9, 5 ]}\nobreak
\hyperdef{L}{X82AD6F187B550060}{}
{\noindent\textcolor{FuncColor}{$\triangleright$\enspace\texttt{DirectSum({\mdseries\slshape R\texttt{\symbol{123}}, S\texttt{\symbol{125}}})\index{DirectSum@\texttt{DirectSum}}
\label{DirectSum}
}\hfill{\scriptsize (function)}}\\
\noindent\textcolor{FuncColor}{$\triangleright$\enspace\texttt{DirectSumOp({\mdseries\slshape list, expl})\index{DirectSumOp@\texttt{DirectSumOp}}
\label{DirectSumOp}
}\hfill{\scriptsize (operation)}}\\


 These functions construct the direct sum of the rings given as arguments. \texttt{DirectSum} takes an arbitrary positive number of arguments and calls the operation \texttt{DirectSumOp}, which takes exactly two arguments, namely a nonempty list of rings and one
of these rings. (This somewhat strange syntax allows the method selection to
choose a reasonable method for special cases.) 
\begin{Verbatim}[commandchars=!@|,fontsize=\small,frame=single,label=Example]
  !gapprompt@gap>| !gapinput@DirectSum(SmallRing(5,1),SmallRing(5,1));|
  <ring with 2 generators>
\end{Verbatim}
 }

 

\subsection{\textcolor{Chapter }{RingByStructureConstants}}
\logpage{[ 56, 9, 6 ]}\nobreak
\hyperdef{L}{X7E7B1B727EA434CF}{}
{\noindent\textcolor{FuncColor}{$\triangleright$\enspace\texttt{RingByStructureConstants({\mdseries\slshape moduli, sctable[, nameinfo]})\index{RingByStructureConstants@\texttt{RingByStructureConstants}}
\label{RingByStructureConstants}
}\hfill{\scriptsize (function)}}\\


 returns a ring $R$ whose additive group is described by the list \mbox{\texttt{\mdseries\slshape moduli}}, with multiplication defined by the structure constants table \mbox{\texttt{\mdseries\slshape sctable}}. The optional argument \mbox{\texttt{\mdseries\slshape nameinfo}} can be used to prescribe names for the elements of the canonical generators of $R$; it can be either a string \mbox{\texttt{\mdseries\slshape name}} (then \mbox{\texttt{\mdseries\slshape name}}\texttt{1}, \mbox{\texttt{\mdseries\slshape name}}\texttt{2} etc. are chosen) or a list of strings which are then chosen. 
\begin{Verbatim}[commandchars=!@|,fontsize=\small,frame=single,label=Example]
  !gapprompt@gap>| !gapinput@T:=EmptySCTable( 1, 0 );;|
  !gapprompt@gap>| !gapinput@SetEntrySCTable( T, 1, 1, [ 1, 1 ] );|
  !gapprompt@gap>| !gapinput@R:=RingByStructureConstants([9], T);   # Z/9Z|
  <ring with 1 generator>
  !gapprompt@gap>| !gapinput@Elements(R);|
  [ 0*r.1, r.1, 2*r.1, 3*r.1, 4*r.1, 5*r.1, 6*r.1, 7*r.1, -r.1 ]
\end{Verbatim}
 }

 }

 }

            
\chapter{\textcolor{Chapter }{Modules}}\label{Modules}
\logpage{[ 57, 0, 0 ]}
\hyperdef{L}{X8183A6857B0C3633}{}
{
   
\section{\textcolor{Chapter }{Generating modules}}\label{Generating modules}
\logpage{[ 57, 1, 0 ]}
\hyperdef{L}{X87A33EFD7CC179C1}{}
{
  

\subsection{\textcolor{Chapter }{IsLeftOperatorAdditiveGroup}}
\logpage{[ 57, 1, 1 ]}\nobreak
\hyperdef{L}{X7C62FE5282E9C505}{}
{\noindent\textcolor{FuncColor}{$\triangleright$\enspace\texttt{IsLeftOperatorAdditiveGroup({\mdseries\slshape D})\index{IsLeftOperatorAdditiveGroup@\texttt{IsLeftOperatorAdditiveGroup}}
\label{IsLeftOperatorAdditiveGroup}
}\hfill{\scriptsize (Category)}}\\


 A domain \mbox{\texttt{\mdseries\slshape D}} lies in \texttt{IsLeftOperatorAdditiveGroup} if it is an additive group that is closed under scalar multiplication from the
left, and such that $\lambda * ( x + y ) = \lambda * x + \lambda * y$ for all scalars $\lambda$ and elements $x, y \in D$ (here and below by scalars we mean elements of a domain acting on \mbox{\texttt{\mdseries\slshape D}} from left or right as appropriate). }

 

\subsection{\textcolor{Chapter }{IsLeftModule}}
\logpage{[ 57, 1, 2 ]}\nobreak
\hyperdef{L}{X7ED323027B291BDF}{}
{\noindent\textcolor{FuncColor}{$\triangleright$\enspace\texttt{IsLeftModule({\mdseries\slshape M})\index{IsLeftModule@\texttt{IsLeftModule}}
\label{IsLeftModule}
}\hfill{\scriptsize (Category)}}\\


 A domain \mbox{\texttt{\mdseries\slshape M}} lies in \texttt{IsLeftModule} if it lies in \texttt{IsLeftOperatorAdditiveGroup}, \emph{and} the set of scalars forms a ring, \emph{and} $(\lambda + \mu) * x = \lambda * x + \mu * x$ for scalars $\lambda, \mu$ and $x \in M$, \emph{and} scalar multiplication satisfies $\lambda * (\mu * x) = (\lambda * \mu) * x$ for scalars $\lambda, \mu$ and $x \in M$. 
\begin{Verbatim}[commandchars=!@|,fontsize=\small,frame=single,label=Example]
  !gapprompt@gap>| !gapinput@V:= FullRowSpace( Rationals, 3 );|
  ( Rationals^3 )
  !gapprompt@gap>| !gapinput@IsLeftModule( V );|
  true
\end{Verbatim}
 }

 

\subsection{\textcolor{Chapter }{GeneratorsOfLeftOperatorAdditiveGroup}}
\logpage{[ 57, 1, 3 ]}\nobreak
\hyperdef{L}{X7F76B1FD84775025}{}
{\noindent\textcolor{FuncColor}{$\triangleright$\enspace\texttt{GeneratorsOfLeftOperatorAdditiveGroup({\mdseries\slshape D})\index{GeneratorsOfLeftOperatorAdditiveGroup@\texttt{Generators}\-\texttt{Of}\-\texttt{Left}\-\texttt{Operator}\-\texttt{Additive}\-\texttt{Group}}
\label{GeneratorsOfLeftOperatorAdditiveGroup}
}\hfill{\scriptsize (attribute)}}\\


 returns a list of elements of \mbox{\texttt{\mdseries\slshape D}} that generates \mbox{\texttt{\mdseries\slshape D}} as a left operator additive group. }

 

\subsection{\textcolor{Chapter }{GeneratorsOfLeftModule}}
\logpage{[ 57, 1, 4 ]}\nobreak
\hyperdef{L}{X7C7684EF867323C2}{}
{\noindent\textcolor{FuncColor}{$\triangleright$\enspace\texttt{GeneratorsOfLeftModule({\mdseries\slshape M})\index{GeneratorsOfLeftModule@\texttt{GeneratorsOfLeftModule}}
\label{GeneratorsOfLeftModule}
}\hfill{\scriptsize (attribute)}}\\


 returns a list of elements of \mbox{\texttt{\mdseries\slshape M}} that generate \mbox{\texttt{\mdseries\slshape M}} as a left module. 
\begin{Verbatim}[commandchars=!@|,fontsize=\small,frame=single,label=Example]
  !gapprompt@gap>| !gapinput@V:= FullRowSpace( Rationals, 3 );;|
  !gapprompt@gap>| !gapinput@GeneratorsOfLeftModule( V );|
  [ [ 1, 0, 0 ], [ 0, 1, 0 ], [ 0, 0, 1 ] ]
\end{Verbatim}
 }

 

\subsection{\textcolor{Chapter }{AsLeftModule}}
\logpage{[ 57, 1, 5 ]}\nobreak
\hyperdef{L}{X7EB3E46D7BC4A35C}{}
{\noindent\textcolor{FuncColor}{$\triangleright$\enspace\texttt{AsLeftModule({\mdseries\slshape R, D})\index{AsLeftModule@\texttt{AsLeftModule}}
\label{AsLeftModule}
}\hfill{\scriptsize (operation)}}\\


 if the domain \mbox{\texttt{\mdseries\slshape D}} forms an additive group and is closed under left multiplication by the
elements of \mbox{\texttt{\mdseries\slshape R}}, then \texttt{AsLeftModule( \mbox{\texttt{\mdseries\slshape R}}, \mbox{\texttt{\mdseries\slshape D}} )} returns the domain \mbox{\texttt{\mdseries\slshape D}} viewed as a left module. 
\begin{Verbatim}[commandchars=!@|,fontsize=\small,frame=single,label=Example]
  !gapprompt@gap>| !gapinput@coll:= [[0*Z(2),0*Z(2)], [Z(2),0*Z(2)], [0*Z(2),Z(2)], [Z(2),Z(2)]];|
  [ [ 0*Z(2), 0*Z(2) ], [ Z(2)^0, 0*Z(2) ], [ 0*Z(2), Z(2)^0 ],
    [ Z(2)^0, Z(2)^0 ] ]
  !gapprompt@gap>| !gapinput@AsLeftModule( GF(2), coll );|
  <vector space of dimension 2 over GF(2)>
\end{Verbatim}
 }

 

\subsection{\textcolor{Chapter }{IsRightOperatorAdditiveGroup}}
\logpage{[ 57, 1, 6 ]}\nobreak
\hyperdef{L}{X7F19AD3D799D0469}{}
{\noindent\textcolor{FuncColor}{$\triangleright$\enspace\texttt{IsRightOperatorAdditiveGroup({\mdseries\slshape D})\index{IsRightOperatorAdditiveGroup@\texttt{IsRightOperatorAdditiveGroup}}
\label{IsRightOperatorAdditiveGroup}
}\hfill{\scriptsize (Category)}}\\


 A domain \mbox{\texttt{\mdseries\slshape D}} lies in \texttt{IsRightOperatorAdditiveGroup} if it is an additive group that is closed under scalar multiplication from the
right, and such that $( x + y ) * \lambda = x * \lambda + y * \lambda$ for all scalars $\lambda$ and elements $x, y \in D$. }

 

\subsection{\textcolor{Chapter }{IsRightModule}}
\logpage{[ 57, 1, 7 ]}\nobreak
\hyperdef{L}{X8479A5AA7DF25F50}{}
{\noindent\textcolor{FuncColor}{$\triangleright$\enspace\texttt{IsRightModule({\mdseries\slshape M})\index{IsRightModule@\texttt{IsRightModule}}
\label{IsRightModule}
}\hfill{\scriptsize (Category)}}\\


 A domain \mbox{\texttt{\mdseries\slshape M}} lies in \texttt{IsRightModule} if it lies in \texttt{IsRightOperatorAdditiveGroup}, \emph{and} the set of scalars forms a ring, \emph{and} $x * (\lambda + \mu) = x * \lambda + x * \mu$ for scalars $\lambda, \mu$ and $x \in M$, \emph{and} scalar multiplication satisfies $(x * \mu) * \lambda = x * (\mu * \lambda)$ for scalars $\lambda, \mu$ and $x \in M$. }

 

\subsection{\textcolor{Chapter }{GeneratorsOfRightOperatorAdditiveGroup}}
\logpage{[ 57, 1, 8 ]}\nobreak
\hyperdef{L}{X7DBC4BCB876EEE1C}{}
{\noindent\textcolor{FuncColor}{$\triangleright$\enspace\texttt{GeneratorsOfRightOperatorAdditiveGroup({\mdseries\slshape D})\index{GeneratorsOfRightOperatorAdditiveGroup@\texttt{Generators}\-\texttt{Of}\-\texttt{Right}\-\texttt{Operator}\-\texttt{Additive}\-\texttt{Group}}
\label{GeneratorsOfRightOperatorAdditiveGroup}
}\hfill{\scriptsize (attribute)}}\\


 returns a list of elements of \mbox{\texttt{\mdseries\slshape D}} that generates \mbox{\texttt{\mdseries\slshape D}} as a right operator additive group. }

 

\subsection{\textcolor{Chapter }{GeneratorsOfRightModule}}
\logpage{[ 57, 1, 9 ]}\nobreak
\hyperdef{L}{X8586A83B85F176F6}{}
{\noindent\textcolor{FuncColor}{$\triangleright$\enspace\texttt{GeneratorsOfRightModule({\mdseries\slshape M})\index{GeneratorsOfRightModule@\texttt{GeneratorsOfRightModule}}
\label{GeneratorsOfRightModule}
}\hfill{\scriptsize (attribute)}}\\


 returns a list of elements of \mbox{\texttt{\mdseries\slshape M}} that generate \mbox{\texttt{\mdseries\slshape M}} as a left module. }

 

\subsection{\textcolor{Chapter }{LeftModuleByGenerators}}
\logpage{[ 57, 1, 10 ]}\nobreak
\hyperdef{L}{X79ED1D7D7F0AE59A}{}
{\noindent\textcolor{FuncColor}{$\triangleright$\enspace\texttt{LeftModuleByGenerators({\mdseries\slshape R, gens[, zero]})\index{LeftModuleByGenerators@\texttt{LeftModuleByGenerators}}
\label{LeftModuleByGenerators}
}\hfill{\scriptsize (operation)}}\\


 returns the left module over \mbox{\texttt{\mdseries\slshape R}} generated by \mbox{\texttt{\mdseries\slshape gens}}. 
\begin{Verbatim}[commandchars=!@|,fontsize=\small,frame=single,label=Example]
  !gapprompt@gap>| !gapinput@coll:= [ [Z(2),0*Z(2)], [0*Z(2),Z(2)], [Z(2),Z(2)] ];;|
  !gapprompt@gap>| !gapinput@V:= LeftModuleByGenerators( GF(16), coll );|
  <vector space over GF(2^4), with 3 generators>
\end{Verbatim}
 }

 

\subsection{\textcolor{Chapter }{LeftActingDomain}}
\logpage{[ 57, 1, 11 ]}\nobreak
\hyperdef{L}{X86F070E0807DC34E}{}
{\noindent\textcolor{FuncColor}{$\triangleright$\enspace\texttt{LeftActingDomain({\mdseries\slshape D})\index{LeftActingDomain@\texttt{LeftActingDomain}}
\label{LeftActingDomain}
}\hfill{\scriptsize (attribute)}}\\


 Let \mbox{\texttt{\mdseries\slshape D}} be an external left set, that is, \mbox{\texttt{\mdseries\slshape D}} is closed under the action of a domain $L$ by multiplication from the left. Then $L$ can be accessed as value of \texttt{LeftActingDomain} for \mbox{\texttt{\mdseries\slshape D}}. }

 }

  
\section{\textcolor{Chapter }{Submodules}}\label{Submodules}
\logpage{[ 57, 2, 0 ]}
\hyperdef{L}{X7934FAE97B6D2AD8}{}
{
  

\subsection{\textcolor{Chapter }{Submodule}}
\logpage{[ 57, 2, 1 ]}\nobreak
\hyperdef{L}{X8465103F874BC07B}{}
{\noindent\textcolor{FuncColor}{$\triangleright$\enspace\texttt{Submodule({\mdseries\slshape M, gens[, "basis"]})\index{Submodule@\texttt{Submodule}}
\label{Submodule}
}\hfill{\scriptsize (function)}}\\


 is the left module generated by the collection \mbox{\texttt{\mdseries\slshape gens}}, with parent module \mbox{\texttt{\mdseries\slshape M}}. If the string \texttt{"basis"} is entered as the third argument then the submodule of \mbox{\texttt{\mdseries\slshape M}} is created for which the list \mbox{\texttt{\mdseries\slshape gens}} is known to be a list of basis vectors; in this case, it is \emph{not} checked whether \mbox{\texttt{\mdseries\slshape gens}} really is linearly independent and whether all in \mbox{\texttt{\mdseries\slshape gens}} lie in \mbox{\texttt{\mdseries\slshape M}}. 
\begin{Verbatim}[commandchars=!@|,fontsize=\small,frame=single,label=Example]
  !gapprompt@gap>| !gapinput@coll:= [ [Z(2),0*Z(2)], [0*Z(2),Z(2)], [Z(2),Z(2)] ];;|
  !gapprompt@gap>| !gapinput@V:= LeftModuleByGenerators( GF(16), coll );;|
  !gapprompt@gap>| !gapinput@W:= Submodule( V, [ coll[1], coll[2] ] );|
  <vector space over GF(2^4), with 2 generators>
  !gapprompt@gap>| !gapinput@Parent( W ) = V;|
  true
\end{Verbatim}
 }

 

\subsection{\textcolor{Chapter }{SubmoduleNC}}
\logpage{[ 57, 2, 2 ]}\nobreak
\hyperdef{L}{X83CF3AD18050C982}{}
{\noindent\textcolor{FuncColor}{$\triangleright$\enspace\texttt{SubmoduleNC({\mdseries\slshape M, gens[, "basis"]})\index{SubmoduleNC@\texttt{SubmoduleNC}}
\label{SubmoduleNC}
}\hfill{\scriptsize (function)}}\\


 \texttt{SubmoduleNC} does the same as \texttt{Submodule} (\ref{Submodule}), except that it does not check whether all in \mbox{\texttt{\mdseries\slshape gens}} lie in \mbox{\texttt{\mdseries\slshape M}}. }

 

\subsection{\textcolor{Chapter }{ClosureLeftModule}}
\logpage{[ 57, 2, 3 ]}\nobreak
\hyperdef{L}{X7C68C4E287481EC0}{}
{\noindent\textcolor{FuncColor}{$\triangleright$\enspace\texttt{ClosureLeftModule({\mdseries\slshape M, m})\index{ClosureLeftModule@\texttt{ClosureLeftModule}}
\label{ClosureLeftModule}
}\hfill{\scriptsize (operation)}}\\


 is the left module generated by the left module generators of \mbox{\texttt{\mdseries\slshape M}} and the element \mbox{\texttt{\mdseries\slshape m}}. 
\begin{Verbatim}[commandchars=!@|,fontsize=\small,frame=single,label=Example]
  !gapprompt@gap>| !gapinput@V:= LeftModuleByGenerators(Rationals, [ [ 1, 0, 0 ], [ 0, 1, 0 ] ]);|
  <vector space over Rationals, with 2 generators>
  !gapprompt@gap>| !gapinput@ClosureLeftModule( V, [ 1, 1, 1 ] );|
  <vector space over Rationals, with 3 generators>
\end{Verbatim}
 }

 

\subsection{\textcolor{Chapter }{TrivialSubmodule}}
\logpage{[ 57, 2, 4 ]}\nobreak
\hyperdef{L}{X7980BC20856B2B7D}{}
{\noindent\textcolor{FuncColor}{$\triangleright$\enspace\texttt{TrivialSubmodule({\mdseries\slshape M})\index{TrivialSubmodule@\texttt{TrivialSubmodule}}
\label{TrivialSubmodule}
}\hfill{\scriptsize (attribute)}}\\


 returns the zero submodule of \mbox{\texttt{\mdseries\slshape M}}. 
\begin{Verbatim}[commandchars=!@|,fontsize=\small,frame=single,label=Example]
  !gapprompt@gap>| !gapinput@V:= LeftModuleByGenerators(Rationals, [[ 1, 0, 0 ], [ 0, 1, 0 ]]);;|
  !gapprompt@gap>| !gapinput@TrivialSubmodule( V );|
  <vector space of dimension 0 over Rationals>
\end{Verbatim}
 }

 }

  
\section{\textcolor{Chapter }{Free Modules}}\label{Free Modules}
\logpage{[ 57, 3, 0 ]}
\hyperdef{L}{X85BD57F27F513D3E}{}
{
  

\subsection{\textcolor{Chapter }{IsFreeLeftModule}}
\logpage{[ 57, 3, 1 ]}\nobreak
\hyperdef{L}{X7C4832187F3D9228}{}
{\noindent\textcolor{FuncColor}{$\triangleright$\enspace\texttt{IsFreeLeftModule({\mdseries\slshape M})\index{IsFreeLeftModule@\texttt{IsFreeLeftModule}}
\label{IsFreeLeftModule}
}\hfill{\scriptsize (Category)}}\\


 A left module is free as module if it is isomorphic to a direct sum of copies
of its left acting domain. 

 Free left modules can have bases. 

 The characteristic (see{\nobreakspace}\texttt{Characteristic} (\ref{Characteristic})) of a free left module is defined as the characteristic of its left acting
domain (see{\nobreakspace}\texttt{LeftActingDomain} (\ref{LeftActingDomain})). }

 

\subsection{\textcolor{Chapter }{FreeLeftModule}}
\logpage{[ 57, 3, 2 ]}\nobreak
\hyperdef{L}{X7C043E307E344AEE}{}
{\noindent\textcolor{FuncColor}{$\triangleright$\enspace\texttt{FreeLeftModule({\mdseries\slshape R, gens[, zero][, "basis"]})\index{FreeLeftModule@\texttt{FreeLeftModule}}
\label{FreeLeftModule}
}\hfill{\scriptsize (function)}}\\


 \texttt{FreeLeftModule( \mbox{\texttt{\mdseries\slshape R}}, \mbox{\texttt{\mdseries\slshape gens}} )} is the free left module over the ring \mbox{\texttt{\mdseries\slshape R}}, generated by the vectors in the collection \mbox{\texttt{\mdseries\slshape gens}}. 

 If there are three arguments, a ring \mbox{\texttt{\mdseries\slshape R}} and a collection \mbox{\texttt{\mdseries\slshape gens}} and an element \mbox{\texttt{\mdseries\slshape zero}}, then \texttt{FreeLeftModule( \mbox{\texttt{\mdseries\slshape R}}, \mbox{\texttt{\mdseries\slshape gens}}, \mbox{\texttt{\mdseries\slshape zero}} )} is the \mbox{\texttt{\mdseries\slshape R}}\texttt{\symbol{45}}free left module generated by \mbox{\texttt{\mdseries\slshape gens}}, with zero element \mbox{\texttt{\mdseries\slshape zero}}. 

 If the last argument is the string \texttt{"basis"} then the vectors in \mbox{\texttt{\mdseries\slshape gens}} are known to form a basis of the free module. 

 It should be noted that the generators \mbox{\texttt{\mdseries\slshape gens}} must be vectors, that is, they must support an addition and a scalar action of \mbox{\texttt{\mdseries\slshape R}} via left multiplication. (See also Section{\nobreakspace}\ref{Constructing Domains} for the general meaning of ``generators'' in \textsf{GAP}.) In particular, \texttt{FreeLeftModule} is \emph{not} an equivalent of commands such as \texttt{FreeGroup} (\ref{FreeGroup:for given rank}) in the sense of a constructor of a free group on abstract generators. Such a
construction seems to be unnecessary for vector spaces, for that one can use
for example row spaces (see{\nobreakspace}\texttt{FullRowSpace} (\ref{FullRowSpace})) in the finite dimensional case and polynomial rings (see{\nobreakspace}\texttt{PolynomialRing} (\ref{PolynomialRing:for a ring and a rank (and an exclusion list)})) in the infinite dimensional case. Moreover, the definition of a ``natural'' addition for elements of a given magma (for example a permutation group) is
possible via the construction of magma rings (see Chapter \ref{Magma Rings}). 
\begin{Verbatim}[commandchars=!@|,fontsize=\small,frame=single,label=Example]
  !gapprompt@gap>| !gapinput@V:= FreeLeftModule(Rationals, [[ 1, 0, 0 ], [ 0, 1, 0 ]], "basis");|
  <vector space of dimension 2 over Rationals>
\end{Verbatim}
 }

 

\subsection{\textcolor{Chapter }{Dimension}}
\logpage{[ 57, 3, 3 ]}\nobreak
\hyperdef{L}{X7E6926C6850E7C4E}{}
{\noindent\textcolor{FuncColor}{$\triangleright$\enspace\texttt{Dimension({\mdseries\slshape M})\index{Dimension@\texttt{Dimension}}
\label{Dimension}
}\hfill{\scriptsize (attribute)}}\\


 A free left module has dimension $n$ if it is isomorphic to a direct sum of $n$ copies of its left acting domain. 

 (We do \emph{not} mark \texttt{Dimension} as invariant under isomorphisms since we want to call \texttt{UseIsomorphismRelation} (\ref{UseIsomorphismRelation}) also for free left modules over different left acting domains.) 
\begin{Verbatim}[commandchars=!@|,fontsize=\small,frame=single,label=Example]
  !gapprompt@gap>| !gapinput@V:= FreeLeftModule( Rationals, [ [ 1, 0 ], [ 0, 1 ], [ 1, 1 ] ] );;|
  !gapprompt@gap>| !gapinput@Dimension( V );|
  2
\end{Verbatim}
 }

 

\subsection{\textcolor{Chapter }{IsFiniteDimensional}}
\logpage{[ 57, 3, 4 ]}\nobreak
\hyperdef{L}{X802DB9FB824B0167}{}
{\noindent\textcolor{FuncColor}{$\triangleright$\enspace\texttt{IsFiniteDimensional({\mdseries\slshape M})\index{IsFiniteDimensional@\texttt{IsFiniteDimensional}}
\label{IsFiniteDimensional}
}\hfill{\scriptsize (property)}}\\


 is \texttt{true} if \mbox{\texttt{\mdseries\slshape M}} is a free left module that is finite dimensional over its left acting domain,
and \texttt{false} otherwise. 
\begin{Verbatim}[commandchars=!@|,fontsize=\small,frame=single,label=Example]
  !gapprompt@gap>| !gapinput@V:= FreeLeftModule( Rationals, [ [ 1, 0 ], [ 0, 1 ], [ 1, 1 ] ] );;|
  !gapprompt@gap>| !gapinput@IsFiniteDimensional( V );|
  true
\end{Verbatim}
 }

 

\subsection{\textcolor{Chapter }{UseBasis}}
\logpage{[ 57, 3, 5 ]}\nobreak
\hyperdef{L}{X7909E8E785420F0E}{}
{\noindent\textcolor{FuncColor}{$\triangleright$\enspace\texttt{UseBasis({\mdseries\slshape V, gens})\index{UseBasis@\texttt{UseBasis}}
\label{UseBasis}
}\hfill{\scriptsize (operation)}}\\


 The vectors in the list \mbox{\texttt{\mdseries\slshape gens}} are known to form a basis of the free left module \mbox{\texttt{\mdseries\slshape V}}. \texttt{UseBasis} stores information in \mbox{\texttt{\mdseries\slshape V}} that can be derived form this fact, namely 
\begin{itemize}
\item  \mbox{\texttt{\mdseries\slshape gens}} are stored as left module generators if no such generators were bound (this is
useful especially if \mbox{\texttt{\mdseries\slshape V}} is an algebra), 
\item  the dimension of \mbox{\texttt{\mdseries\slshape V}} is stored. 
\end{itemize}
 
\begin{Verbatim}[commandchars=!@|,fontsize=\small,frame=single,label=Example]
  !gapprompt@gap>| !gapinput@V:= FreeLeftModule( Rationals, [ [ 1, 0 ], [ 0, 1 ], [ 1, 1 ] ] );;|
  !gapprompt@gap>| !gapinput@UseBasis( V, [ [ 1, 0 ], [ 1, 1 ] ] );|
  !gapprompt@gap>| !gapinput@V;  # now V knows its dimension|
  <vector space of dimension 2 over Rationals>
\end{Verbatim}
 }

 

\subsection{\textcolor{Chapter }{IsRowModule}}
\logpage{[ 57, 3, 6 ]}\nobreak
\hyperdef{L}{X7C8F844783F4FA09}{}
{\noindent\textcolor{FuncColor}{$\triangleright$\enspace\texttt{IsRowModule({\mdseries\slshape V})\index{IsRowModule@\texttt{IsRowModule}}
\label{IsRowModule}
}\hfill{\scriptsize (property)}}\\


 A \emph{row module} is a free left module whose elements are row vectors. }

 

\subsection{\textcolor{Chapter }{IsMatrixModule}}
\logpage{[ 57, 3, 7 ]}\nobreak
\hyperdef{L}{X81FCC1D780435CF1}{}
{\noindent\textcolor{FuncColor}{$\triangleright$\enspace\texttt{IsMatrixModule({\mdseries\slshape V})\index{IsMatrixModule@\texttt{IsMatrixModule}}
\label{IsMatrixModule}
}\hfill{\scriptsize (property)}}\\


 A \emph{matrix module} is a free left module whose elements are matrices. }

 

\subsection{\textcolor{Chapter }{IsFullRowModule}}
\logpage{[ 57, 3, 8 ]}\nobreak
\hyperdef{L}{X853E085C868196EF}{}
{\noindent\textcolor{FuncColor}{$\triangleright$\enspace\texttt{IsFullRowModule({\mdseries\slshape M})\index{IsFullRowModule@\texttt{IsFullRowModule}}
\label{IsFullRowModule}
}\hfill{\scriptsize (property)}}\\


 A \emph{full row module} is a module $R^n$, for a ring $R$ and a nonnegative integer $n$. 

 More precisely, a full row module is a free left module over a ring $R$ such that the elements are row vectors of the same length $n$ and with entries in $R$ and such that the dimension is equal to $n$. 

 Several functions delegate their tasks to full row modules, for example \texttt{Iterator} (\ref{Iterator}) and \texttt{Enumerator} (\ref{Enumerator}). }

 

\subsection{\textcolor{Chapter }{FullRowModule}}
\logpage{[ 57, 3, 9 ]}\nobreak
\hyperdef{L}{X848041A47BC4B038}{}
{\noindent\textcolor{FuncColor}{$\triangleright$\enspace\texttt{FullRowModule({\mdseries\slshape R, n})\index{FullRowModule@\texttt{FullRowModule}}
\label{FullRowModule}
}\hfill{\scriptsize (function)}}\\


 is the row module \texttt{\mbox{\texttt{\mdseries\slshape R}}\texttt{\symbol{94}}\mbox{\texttt{\mdseries\slshape n}}}, for a ring \mbox{\texttt{\mdseries\slshape R}} and a nonnegative integer \mbox{\texttt{\mdseries\slshape n}}. 
\begin{Verbatim}[commandchars=!@|,fontsize=\small,frame=single,label=Example]
  !gapprompt@gap>| !gapinput@V:= FullRowModule( Integers, 5 );|
  ( Integers^5 )
\end{Verbatim}
 }

 

\subsection{\textcolor{Chapter }{IsFullMatrixModule}}
\logpage{[ 57, 3, 10 ]}\nobreak
\hyperdef{L}{X814CEA62842CF5BB}{}
{\noindent\textcolor{FuncColor}{$\triangleright$\enspace\texttt{IsFullMatrixModule({\mdseries\slshape M})\index{IsFullMatrixModule@\texttt{IsFullMatrixModule}}
\label{IsFullMatrixModule}
}\hfill{\scriptsize (property)}}\\


 A \emph{full matrix module} is a module $R^{{[m,n]}}$, for a ring $R$ and two nonnegative integers $m$, $n$. 

 More precisely, a full matrix module is a free left module over a ring $R$ such that the elements are $m$ by $n$ matrices with entries in $R$ and such that the dimension is equal to $m n$. }

 

\subsection{\textcolor{Chapter }{FullMatrixModule}}
\logpage{[ 57, 3, 11 ]}\nobreak
\hyperdef{L}{X7A0C871B7C446F1F}{}
{\noindent\textcolor{FuncColor}{$\triangleright$\enspace\texttt{FullMatrixModule({\mdseries\slshape R, m, n})\index{FullMatrixModule@\texttt{FullMatrixModule}}
\label{FullMatrixModule}
}\hfill{\scriptsize (function)}}\\


 is the matrix module \texttt{\mbox{\texttt{\mdseries\slshape R}}\texttt{\symbol{94}}[\mbox{\texttt{\mdseries\slshape m}},\mbox{\texttt{\mdseries\slshape n}}]}, for a ring \mbox{\texttt{\mdseries\slshape R}} and nonnegative integers \mbox{\texttt{\mdseries\slshape m}} and \mbox{\texttt{\mdseries\slshape n}}. 
\begin{Verbatim}[commandchars=!@|,fontsize=\small,frame=single,label=Example]
  !gapprompt@gap>| !gapinput@FullMatrixModule( GaussianIntegers, 3, 6 );|
  ( GaussianIntegers^[ 3, 6 ] )
\end{Verbatim}
 }

 }

 }

         
\chapter{\textcolor{Chapter }{Fields and Division Rings}}\label{Fields and Division Rings}
\logpage{[ 58, 0, 0 ]}
\hyperdef{L}{X80A8E676814A19FD}{}
{
  \index{fields} \index{division rings} A \emph{division ring} is a ring (see Chapter{\nobreakspace}\ref{Rings}) in which every non\texttt{\symbol{45}}zero element has an inverse. The most
important class of division rings are the commutative ones, which are called \emph{fields}. 

 \textsf{GAP} supports finite fields (see Chapter{\nobreakspace}\ref{Finite Fields}) and abelian number fields (see Chapter{\nobreakspace}\ref{Abelian Number Fields}), in particular the field of rationals (see Chapter{\nobreakspace}\ref{Rational Numbers}). 

 This chapter describes the general \textsf{GAP} functions for fields and division rings. 

 If a field \mbox{\texttt{\mdseries\slshape F}} is a subfield of a commutative ring \mbox{\texttt{\mdseries\slshape C}}, \mbox{\texttt{\mdseries\slshape C}} can be considered as a vector space over the (left) acting domain \mbox{\texttt{\mdseries\slshape F}} (see Chapter{\nobreakspace}\ref{Vector Spaces}). In this situation, we call \mbox{\texttt{\mdseries\slshape F}} the \emph{field of definition} of \mbox{\texttt{\mdseries\slshape C}}. 

 Each field in \textsf{GAP} is represented as a vector space over a subfield (see{\nobreakspace}\texttt{IsField} (\ref{IsField})), thus each field is in fact a field extension in a natural way, which is
used by functions such as \texttt{Norm} (\ref{Norm}) and \texttt{Trace} (\ref{Trace:for a field element}) (see{\nobreakspace}\ref{Galois Action}).  
\section{\textcolor{Chapter }{Generating Fields}}\label{Generating Fields}
\logpage{[ 58, 1, 0 ]}
\hyperdef{L}{X82B74B458705B3CE}{}
{
  

\subsection{\textcolor{Chapter }{IsDivisionRing}}
\logpage{[ 58, 1, 1 ]}\nobreak
\hyperdef{L}{X7F2CAA9E7A16913D}{}
{\noindent\textcolor{FuncColor}{$\triangleright$\enspace\texttt{IsDivisionRing({\mdseries\slshape D})\index{IsDivisionRing@\texttt{IsDivisionRing}}
\label{IsDivisionRing}
}\hfill{\scriptsize (Category)}}\\


 A \emph{division ring} in \textsf{GAP} is a nontrivial associative algebra \mbox{\texttt{\mdseries\slshape D}} with a multiplicative inverse for each nonzero element. In \textsf{GAP} every division ring is a vector space over a division ring (possibly over
itself). Note that being a division ring is thus not a property that a ring
can get, because a ring is usually not represented as a vector space. 

 The field of coefficients is stored as the value of the attribute \texttt{LeftActingDomain} (\ref{LeftActingDomain}) of \mbox{\texttt{\mdseries\slshape D}}. }

 

\subsection{\textcolor{Chapter }{IsField}}
\logpage{[ 58, 1, 2 ]}\nobreak
\hyperdef{L}{X7A5AE30E7C0F457C}{}
{\noindent\textcolor{FuncColor}{$\triangleright$\enspace\texttt{IsField({\mdseries\slshape D})\index{IsField@\texttt{IsField}}
\label{IsField}
}\hfill{\scriptsize (filter)}}\\


 A \emph{field} is a commutative division ring (see{\nobreakspace}\texttt{IsDivisionRing} (\ref{IsDivisionRing}) and{\nobreakspace}\texttt{IsCommutative} (\ref{IsCommutative})). 
\begin{Verbatim}[commandchars=!@|,fontsize=\small,frame=single,label=Example]
  !gapprompt@gap>| !gapinput@IsField( GaloisField(16) );           # the field with 16 elements|
  true
  !gapprompt@gap>| !gapinput@IsField( Rationals );                 # the field of rationals|
  true
  !gapprompt@gap>| !gapinput@q:= QuaternionAlgebra( Rationals );;  # noncommutative division ring|
  !gapprompt@gap>| !gapinput@IsField( q );  IsDivisionRing( q );|
  false
  true
  !gapprompt@gap>| !gapinput@mat:= [ [ 1 ] ];;  a:= Algebra( Rationals, [ mat ] );;|
  !gapprompt@gap>| !gapinput@IsDivisionRing( a );   # algebra not constructed as a division ring|
  false
\end{Verbatim}
 }

 

\subsection{\textcolor{Chapter }{Field (for several generators)}}
\logpage{[ 58, 1, 3 ]}\nobreak
\hyperdef{L}{X871AA7D58263E9AC}{}
{\noindent\textcolor{FuncColor}{$\triangleright$\enspace\texttt{Field({\mdseries\slshape z, ...})\index{Field@\texttt{Field}!for several generators}
\label{Field:for several generators}
}\hfill{\scriptsize (function)}}\\
\noindent\textcolor{FuncColor}{$\triangleright$\enspace\texttt{Field({\mdseries\slshape [F, ]list})\index{Field@\texttt{Field}!for (a field and) a list of generators}
\label{Field:for (a field and) a list of generators}
}\hfill{\scriptsize (function)}}\\


 \texttt{Field} returns the smallest field $K$ that contains all the elements $\mbox{\texttt{\mdseries\slshape z}}, \ldots$, or the smallest field $K$ that contains all elements in the list \mbox{\texttt{\mdseries\slshape list}}. If no subfield \mbox{\texttt{\mdseries\slshape F}} is given, $K$ is constructed as a field over itself, i.e. the left acting domain of $K$ is $K$. Called with a field \mbox{\texttt{\mdseries\slshape F}} and a list \mbox{\texttt{\mdseries\slshape list}}, \texttt{Field} constructs the field generated by \mbox{\texttt{\mdseries\slshape F}} and the elements in \mbox{\texttt{\mdseries\slshape list}}, as a vector space over \mbox{\texttt{\mdseries\slshape F}}. }

 

\subsection{\textcolor{Chapter }{DefaultField (for several generators)}}
\logpage{[ 58, 1, 4 ]}\nobreak
\hyperdef{L}{X7D9F7FD4786691EE}{}
{\noindent\textcolor{FuncColor}{$\triangleright$\enspace\texttt{DefaultField({\mdseries\slshape z, ...})\index{DefaultField@\texttt{DefaultField}!for several generators}
\label{DefaultField:for several generators}
}\hfill{\scriptsize (function)}}\\
\noindent\textcolor{FuncColor}{$\triangleright$\enspace\texttt{DefaultField({\mdseries\slshape list})\index{DefaultField@\texttt{DefaultField}!for a list of generators}
\label{DefaultField:for a list of generators}
}\hfill{\scriptsize (function)}}\\


 \texttt{DefaultField} returns a field $K$ that contains all the elements $\mbox{\texttt{\mdseries\slshape z}}, \ldots$, or a field $K$ that contains all elements in the list \mbox{\texttt{\mdseries\slshape list}}. 

 This field need not be the smallest field in which the elements lie,
cf.{\nobreakspace}\texttt{Field} (\ref{Field:for several generators}). For example, for elements from cyclotomic fields \texttt{DefaultField} returns the smallest cyclotomic field in which the elements lie, but the
elements may lie in a smaller number field which is not a cyclotomic field. 
\begin{Verbatim}[commandchars=!@|,fontsize=\small,frame=single,label=Example]
  !gapprompt@gap>| !gapinput@Field( Z(4) );  Field( [ Z(4), Z(8) ] );  # finite fields|
  GF(2^2)
  GF(2^6)
  !gapprompt@gap>| !gapinput@Field( E(9) );  Field( CF(4), [ E(9) ] ); # abelian number fields|
  CF(9)
  AsField( GaussianRationals, CF(36) )
  !gapprompt@gap>| !gapinput@f1:= Field( EB(5) );  f2:= DefaultField( EB(5) );|
  NF(5,[ 1, 4 ])
  CF(5)
  !gapprompt@gap>| !gapinput@f1 = f2;  IsSubset( f2, f1 );|
  false
  true
\end{Verbatim}
 }

 

\subsection{\textcolor{Chapter }{DefaultFieldByGenerators}}
\logpage{[ 58, 1, 5 ]}\nobreak
\hyperdef{L}{X7C298A40852C2AFF}{}
{\noindent\textcolor{FuncColor}{$\triangleright$\enspace\texttt{DefaultFieldByGenerators({\mdseries\slshape [z, ...]})\index{DefaultFieldByGenerators@\texttt{DefaultFieldByGenerators}}
\label{DefaultFieldByGenerators}
}\hfill{\scriptsize (operation)}}\\


 returns the default field containing the elements \mbox{\texttt{\mdseries\slshape z}}, $\ldots$. This field may be bigger than the smallest field containing these elements. }

 

\subsection{\textcolor{Chapter }{GeneratorsOfDivisionRing}}
\logpage{[ 58, 1, 6 ]}\nobreak
\hyperdef{L}{X7EF624958648D0FA}{}
{\noindent\textcolor{FuncColor}{$\triangleright$\enspace\texttt{GeneratorsOfDivisionRing({\mdseries\slshape D})\index{GeneratorsOfDivisionRing@\texttt{GeneratorsOfDivisionRing}}
\label{GeneratorsOfDivisionRing}
}\hfill{\scriptsize (attribute)}}\\


 generators with respect to addition, multiplication, and taking inverses (the
identity cannot be omitted ...) }

 

\subsection{\textcolor{Chapter }{GeneratorsOfField}}
\logpage{[ 58, 1, 7 ]}\nobreak
\hyperdef{L}{X7AA715317A81261B}{}
{\noindent\textcolor{FuncColor}{$\triangleright$\enspace\texttt{GeneratorsOfField({\mdseries\slshape F})\index{GeneratorsOfField@\texttt{GeneratorsOfField}}
\label{GeneratorsOfField}
}\hfill{\scriptsize (attribute)}}\\


 generators with respect to addition, multiplication, and taking inverses. This
attribute is the same as \texttt{GeneratorsOfDivisionRing} (\ref{GeneratorsOfDivisionRing}). }

 

\subsection{\textcolor{Chapter }{DivisionRingByGenerators}}
\logpage{[ 58, 1, 8 ]}\nobreak
\hyperdef{L}{X8641861A8550F8BE}{}
{\noindent\textcolor{FuncColor}{$\triangleright$\enspace\texttt{DivisionRingByGenerators({\mdseries\slshape [F, ]gens})\index{DivisionRingByGenerators@\texttt{DivisionRingByGenerators}}
\label{DivisionRingByGenerators}
}\hfill{\scriptsize (operation)}}\\
\noindent\textcolor{FuncColor}{$\triangleright$\enspace\texttt{FieldByGenerators({\mdseries\slshape [F, ]gens})\index{FieldByGenerators@\texttt{FieldByGenerators}}
\label{FieldByGenerators}
}\hfill{\scriptsize (operation)}}\\


 Called with a field \mbox{\texttt{\mdseries\slshape F}} and a list \mbox{\texttt{\mdseries\slshape gens}} of scalars, \texttt{DivisionRingByGenerators} returns the division ring over \mbox{\texttt{\mdseries\slshape F}} generated by \mbox{\texttt{\mdseries\slshape gens}}. The unary version returns the division ring as vector space over \texttt{FieldOverItselfByGenerators( \mbox{\texttt{\mdseries\slshape gens}} )}. 

 \texttt{FieldByGenerators} is just a synonym for \texttt{DivisionRingByGenerators}. }

 

\subsection{\textcolor{Chapter }{AsDivisionRing}}
\logpage{[ 58, 1, 9 ]}\nobreak
\hyperdef{L}{X7C193B7D7AFB29BE}{}
{\noindent\textcolor{FuncColor}{$\triangleright$\enspace\texttt{AsDivisionRing({\mdseries\slshape [F, ]C})\index{AsDivisionRing@\texttt{AsDivisionRing}}
\label{AsDivisionRing}
}\hfill{\scriptsize (operation)}}\\
\noindent\textcolor{FuncColor}{$\triangleright$\enspace\texttt{AsField({\mdseries\slshape [F, ]C})\index{AsField@\texttt{AsField}}
\label{AsField}
}\hfill{\scriptsize (operation)}}\\


 If the collection \mbox{\texttt{\mdseries\slshape C}} can be regarded as a division ring then \texttt{AsDivisionRing( \mbox{\texttt{\mdseries\slshape C}} )} is the division ring that consists of the elements of \mbox{\texttt{\mdseries\slshape C}}, viewed as a vector space over its prime field; otherwise \texttt{fail} is returned. 

 In the second form, if \mbox{\texttt{\mdseries\slshape F}} is a division ring contained in \mbox{\texttt{\mdseries\slshape C}} then the returned division ring is viewed as a vector space over \mbox{\texttt{\mdseries\slshape F}}. 

 \texttt{AsField} is just a synonym for \texttt{AsDivisionRing}. }

 }

  
\section{\textcolor{Chapter }{Subfields of Fields}}\label{Subfields of Fields}
\logpage{[ 58, 2, 0 ]}
\hyperdef{L}{X7C53566A839B57F6}{}
{
  

\subsection{\textcolor{Chapter }{Subfield}}
\logpage{[ 58, 2, 1 ]}\nobreak
\hyperdef{L}{X7FE1FA217A08DCE5}{}
{\noindent\textcolor{FuncColor}{$\triangleright$\enspace\texttt{Subfield({\mdseries\slshape F, gens})\index{Subfield@\texttt{Subfield}}
\label{Subfield}
}\hfill{\scriptsize (function)}}\\
\noindent\textcolor{FuncColor}{$\triangleright$\enspace\texttt{SubfieldNC({\mdseries\slshape F, gens})\index{SubfieldNC@\texttt{SubfieldNC}}
\label{SubfieldNC}
}\hfill{\scriptsize (function)}}\\


 Constructs the subfield of \mbox{\texttt{\mdseries\slshape F}} generated by \mbox{\texttt{\mdseries\slshape gens}}. }

 

\subsection{\textcolor{Chapter }{FieldOverItselfByGenerators}}
\logpage{[ 58, 2, 2 ]}\nobreak
\hyperdef{L}{X82A0E79A7B9799E0}{}
{\noindent\textcolor{FuncColor}{$\triangleright$\enspace\texttt{FieldOverItselfByGenerators({\mdseries\slshape [z, ...]})\index{FieldOverItselfByGenerators@\texttt{FieldOverItselfByGenerators}}
\label{FieldOverItselfByGenerators}
}\hfill{\scriptsize (operation)}}\\


 This operation is needed for the call of \texttt{Field} (\ref{Field:for several generators}) or \texttt{FieldByGenerators} (\ref{FieldByGenerators}) without explicitly given subfield, in order to construct a left acting domain
for such a field. }

 

\subsection{\textcolor{Chapter }{PrimitiveElement}}
\logpage{[ 58, 2, 3 ]}\nobreak
\hyperdef{L}{X86DB31B57FB4F570}{}
{\noindent\textcolor{FuncColor}{$\triangleright$\enspace\texttt{PrimitiveElement({\mdseries\slshape D})\index{PrimitiveElement@\texttt{PrimitiveElement}}
\label{PrimitiveElement}
}\hfill{\scriptsize (attribute)}}\\


 is an element of \mbox{\texttt{\mdseries\slshape D}} that generates \mbox{\texttt{\mdseries\slshape D}} as a division ring together with the left acting domain. }

 

\subsection{\textcolor{Chapter }{PrimeField}}
\logpage{[ 58, 2, 4 ]}\nobreak
\hyperdef{L}{X7DD27F927BD57FDE}{}
{\noindent\textcolor{FuncColor}{$\triangleright$\enspace\texttt{PrimeField({\mdseries\slshape D})\index{PrimeField@\texttt{PrimeField}}
\label{PrimeField}
}\hfill{\scriptsize (attribute)}}\\


 The \emph{prime field} of a division ring \mbox{\texttt{\mdseries\slshape D}} is the smallest field which is contained in \mbox{\texttt{\mdseries\slshape D}}. For example, the prime field of any field in characteristic zero is
isomorphic to the field of rational numbers. }

 

\subsection{\textcolor{Chapter }{IsPrimeField}}
\logpage{[ 58, 2, 5 ]}\nobreak
\hyperdef{L}{X84B6F1E67AD0E33D}{}
{\noindent\textcolor{FuncColor}{$\triangleright$\enspace\texttt{IsPrimeField({\mdseries\slshape D})\index{IsPrimeField@\texttt{IsPrimeField}}
\label{IsPrimeField}
}\hfill{\scriptsize (property)}}\\


 A division ring is a prime field if it is equal to its prime field
(see{\nobreakspace}\texttt{PrimeField} (\ref{PrimeField})). }

 

\subsection{\textcolor{Chapter }{DegreeOverPrimeField}}
\logpage{[ 58, 2, 6 ]}\nobreak
\hyperdef{L}{X7845CECE86A83219}{}
{\noindent\textcolor{FuncColor}{$\triangleright$\enspace\texttt{DegreeOverPrimeField({\mdseries\slshape F})\index{DegreeOverPrimeField@\texttt{DegreeOverPrimeField}}
\label{DegreeOverPrimeField}
}\hfill{\scriptsize (attribute)}}\\


 is the degree of the field \mbox{\texttt{\mdseries\slshape F}} over its prime field (see{\nobreakspace}\texttt{PrimeField} (\ref{PrimeField})). }

 

\subsection{\textcolor{Chapter }{DefiningPolynomial}}
\logpage{[ 58, 2, 7 ]}\nobreak
\hyperdef{L}{X7ADDCBF47E2ED3D4}{}
{\noindent\textcolor{FuncColor}{$\triangleright$\enspace\texttt{DefiningPolynomial({\mdseries\slshape F})\index{DefiningPolynomial@\texttt{DefiningPolynomial}}
\label{DefiningPolynomial}
}\hfill{\scriptsize (attribute)}}\\


 is the defining polynomial of the field \mbox{\texttt{\mdseries\slshape F}} as a field extension over the left acting domain of \mbox{\texttt{\mdseries\slshape F}}. A root of the defining polynomial can be computed with \texttt{RootOfDefiningPolynomial} (\ref{RootOfDefiningPolynomial}). }

 

\subsection{\textcolor{Chapter }{RootOfDefiningPolynomial}}
\logpage{[ 58, 2, 8 ]}\nobreak
\hyperdef{L}{X8173DA4982DB1E8A}{}
{\noindent\textcolor{FuncColor}{$\triangleright$\enspace\texttt{RootOfDefiningPolynomial({\mdseries\slshape F})\index{RootOfDefiningPolynomial@\texttt{RootOfDefiningPolynomial}}
\label{RootOfDefiningPolynomial}
}\hfill{\scriptsize (attribute)}}\\


 is a root in the field \mbox{\texttt{\mdseries\slshape F}} of its defining polynomial as a field extension over the left acting domain of \mbox{\texttt{\mdseries\slshape F}}. The defining polynomial can be computed with \texttt{DefiningPolynomial} (\ref{DefiningPolynomial}). }

 

\subsection{\textcolor{Chapter }{FieldExtension}}
\logpage{[ 58, 2, 9 ]}\nobreak
\hyperdef{L}{X82718B3B818DC699}{}
{\noindent\textcolor{FuncColor}{$\triangleright$\enspace\texttt{FieldExtension({\mdseries\slshape F, poly})\index{FieldExtension@\texttt{FieldExtension}}
\label{FieldExtension}
}\hfill{\scriptsize (operation)}}\\


 is the field obtained on adjoining a root of the irreducible polynomial \mbox{\texttt{\mdseries\slshape poly}} to the field \mbox{\texttt{\mdseries\slshape F}}. }

 

\subsection{\textcolor{Chapter }{Subfields}}
\logpage{[ 58, 2, 10 ]}\nobreak
\hyperdef{L}{X83490C65819D85FE}{}
{\noindent\textcolor{FuncColor}{$\triangleright$\enspace\texttt{Subfields({\mdseries\slshape F})\index{Subfields@\texttt{Subfields}}
\label{Subfields}
}\hfill{\scriptsize (attribute)}}\\


 is the set of all subfields of the field \mbox{\texttt{\mdseries\slshape F}}.  }

 }

  
\section{\textcolor{Chapter }{Galois Action}}\label{Galois Action}
\logpage{[ 58, 3, 0 ]}
\hyperdef{L}{X7D9A02B07D08FA40}{}
{
  Let $L > K$ be a field extension of finite degree. Then to each element $\alpha \in L$, we can associate a $K$\texttt{\symbol{45}}linear mapping $\varphi_{\alpha}$ on $L$, and for a fixed $K$\texttt{\symbol{45}}basis of $L$, we can associate to $\alpha$ the matrix $M_{\alpha}$ (over $K$) of this mapping. 

 The \emph{norm} of $\alpha$ is defined as the determinant of $M_{\alpha}$, the \emph{trace} of $\alpha$ is defined as the trace of $M_{\alpha}$, the \emph{minimal polynomial} $\mu_{\alpha}$ and the \emph{trace polynomial} $\chi_{\alpha}$ of $\alpha$ are defined as the minimal polynomial (see{\nobreakspace}\ref{MinimalPolynomial}) and the characteristic polynomial (see{\nobreakspace}\texttt{CharacteristicPolynomial} (\ref{CharacteristicPolynomial}) and \texttt{TracePolynomial} (\ref{TracePolynomial})) of $M_{\alpha}$. (Note that $\mu_{\alpha}$ depends only on $K$ whereas $\chi_{\alpha}$ depends on both $L$ and $K$.) 

 Thus norm and trace of $\alpha$ are elements of $K$, and $\mu_{\alpha}$ and $\chi_{\alpha}$ are polynomials over $K$, $\chi_{\alpha}$ being a power of $\mu_{\alpha}$, and the degree of $\chi_{\alpha}$ equals the degree of the field extension $L > K$. 

 The \emph{conjugates} of $\alpha$ in $L$ are those roots of $\chi_{\alpha}$ (with multiplicity) that lie in $L$; note that if only $L$ is given, there is in general no way to access the roots outside $L$. 

 Analogously, the \emph{Galois group} of the extension $L > K$ is defined as the group of all those field automorphisms of $L$ that fix $K$ pointwise. 

 If $L > K$ is a Galois extension then the conjugates of $\alpha$ are all roots of $\chi_{\alpha}$ (with multiplicity), the set of conjugates equals the roots of $\mu_{\alpha}$, the norm of $\alpha$ equals the product and the trace of $\alpha$ equals the sum of the conjugates of $\alpha$, and the Galois group in the sense of the above definition equals the usual
Galois group, 

 Note that \texttt{MinimalPolynomial( \mbox{\texttt{\mdseries\slshape F}}, \mbox{\texttt{\mdseries\slshape z}} )} is a polynomial \emph{over} \mbox{\texttt{\mdseries\slshape F}}, whereas \texttt{Norm( \mbox{\texttt{\mdseries\slshape F}}, \mbox{\texttt{\mdseries\slshape z}} )} is the norm of the element \mbox{\texttt{\mdseries\slshape z}} \emph{in} \mbox{\texttt{\mdseries\slshape F}} w.r.t.{\nobreakspace}the field extension \texttt{\mbox{\texttt{\mdseries\slshape F}} {\textgreater} LeftActingDomain( \mbox{\texttt{\mdseries\slshape F}} )}. 

 The default methods for field elements are as follows. \texttt{MinimalPolynomial} (\ref{MinimalPolynomial}) solves a system of linear equations, \texttt{TracePolynomial} (\ref{TracePolynomial}) computes the appropriate power of the minimal polynomial, \texttt{Norm} (\ref{Norm}) and \texttt{Trace} (\ref{Trace:for a field element}) values are obtained as coefficients of the characteristic polynomial, and \texttt{Conjugates} (\ref{Conjugates}) uses the factorization of the characteristic polynomial. 

 For elements in finite fields and cyclotomic fields, one wants to do the
computations in a different way since the field extensions in question are
Galois extensions, and the Galois groups are well\texttt{\symbol{45}}known in
these cases. More general, if a field is in the category \texttt{IsFieldControlledByGaloisGroup} then the default methods are the following. \texttt{Conjugates} (\ref{Conjugates}) returns the sorted list of images (with multiplicity) of the element under the
Galois group, \texttt{Norm} (\ref{Norm}) computes the product of the conjugates, \texttt{Trace} (\ref{Trace:for a field element}) computes the sum of the conjugates, \texttt{TracePolynomial} (\ref{TracePolynomial}) and \texttt{MinimalPolynomial} (\ref{MinimalPolynomial}) compute the product of linear factors $x - c$ with $c$ ranging over the conjugates and the set of conjugates, respectively. \index{IsFieldControlledByGaloisGroup@\texttt{IsFieldControlledByGaloisGroup}} 

\subsection{\textcolor{Chapter }{GaloisGroup (of field)}}
\logpage{[ 58, 3, 1 ]}\nobreak
\hyperdef{L}{X80CAA5BA82F09ED2}{}
{\noindent\textcolor{FuncColor}{$\triangleright$\enspace\texttt{GaloisGroup({\mdseries\slshape F})\index{GaloisGroup@\texttt{GaloisGroup}!of field}
\label{GaloisGroup:of field}
}\hfill{\scriptsize (attribute)}}\\


 The \emph{Galois group} of a field \mbox{\texttt{\mdseries\slshape F}} is the group of all field automorphisms of \mbox{\texttt{\mdseries\slshape F}} that fix the subfield $K = $\texttt{LeftActingDomain( \mbox{\texttt{\mdseries\slshape F}} )} pointwise. 

 Note that the field extension $\mbox{\texttt{\mdseries\slshape F}} > K$ need \emph{not} be a Galois extension. 
\begin{Verbatim}[commandchars=!@|,fontsize=\small,frame=single,label=Example]
  !gapprompt@gap>| !gapinput@g:= GaloisGroup( AsField( GF(2^2), GF(2^12) ) );;|
  !gapprompt@gap>| !gapinput@Size( g );  IsCyclic( g );|
  6
  true
  !gapprompt@gap>| !gapinput@h:= GaloisGroup( CF(60) );;|
  !gapprompt@gap>| !gapinput@Size( h );  IsAbelian( h );|
  16
  true
\end{Verbatim}
 }

 

\subsection{\textcolor{Chapter }{MinimalPolynomial (over a field)}}
\logpage{[ 58, 3, 2 ]}\nobreak
\hyperdef{L}{X8738C6687D784BB5}{}
{\noindent\textcolor{FuncColor}{$\triangleright$\enspace\texttt{MinimalPolynomial({\mdseries\slshape F, z[, ind]})\index{MinimalPolynomial@\texttt{MinimalPolynomial}!over a field}
\label{MinimalPolynomial:over a field}
}\hfill{\scriptsize (operation)}}\\


 returns the minimal polynomial of \mbox{\texttt{\mdseries\slshape z}} over the field \mbox{\texttt{\mdseries\slshape F}}. This is a generator of the ideal in $\mbox{\texttt{\mdseries\slshape F}}[x]$ of all polynomials which vanish on \mbox{\texttt{\mdseries\slshape z}}. (This definition is consistent with the general definition of \texttt{MinimalPolynomial} (\ref{MinimalPolynomial}) for rings.) 

 
\begin{Verbatim}[commandchars=!@|,fontsize=\small,frame=single,label=Example]
  !gapprompt@gap>| !gapinput@MinimalPolynomial( Rationals, E(8) );|
  x_1^4+1
  !gapprompt@gap>| !gapinput@MinimalPolynomial( CF(4), E(8) );|
  x_1^2+(-E(4))
  !gapprompt@gap>| !gapinput@MinimalPolynomial( CF(8), E(8) );|
  x_1+(-E(8))
\end{Verbatim}
 }

 

\subsection{\textcolor{Chapter }{TracePolynomial}}
\logpage{[ 58, 3, 3 ]}\nobreak
\hyperdef{L}{X80FE7E017C2D255C}{}
{\noindent\textcolor{FuncColor}{$\triangleright$\enspace\texttt{TracePolynomial({\mdseries\slshape L, K, z[, inum]})\index{TracePolynomial@\texttt{TracePolynomial}}
\label{TracePolynomial}
}\hfill{\scriptsize (operation)}}\\


 \index{characteristic polynomial!for field elements} returns the polynomial that is the product of $(X - c)$ where $c$ runs over the conjugates of \mbox{\texttt{\mdseries\slshape z}} in the field extension \mbox{\texttt{\mdseries\slshape L}} over \mbox{\texttt{\mdseries\slshape K}}. The polynomial is returned as a univariate polynomial over \mbox{\texttt{\mdseries\slshape K}} in the indeterminate number \mbox{\texttt{\mdseries\slshape inum}} (defaulting to 1). 

 This polynomial is sometimes also called the \emph{characteristic polynomial} of \mbox{\texttt{\mdseries\slshape z}} w.r.t.{\nobreakspace}the field extension $\mbox{\texttt{\mdseries\slshape L}} > \mbox{\texttt{\mdseries\slshape K}}$. Therefore methods are installed for \texttt{CharacteristicPolynomial} (\ref{CharacteristicPolynomial}) that call \texttt{TracePolynomial} in the case of field extensions. 

 
\begin{Verbatim}[commandchars=!@|,fontsize=\small,frame=single,label=Example]
  !gapprompt@gap>| !gapinput@TracePolynomial( CF(8), Rationals, E(8) );|
  x_1^4+1
  !gapprompt@gap>| !gapinput@TracePolynomial( CF(16), Rationals, E(8) );|
  x_1^8+2*x_1^4+1
\end{Verbatim}
 }

 

\subsection{\textcolor{Chapter }{Norm}}
\logpage{[ 58, 3, 4 ]}\nobreak
\hyperdef{L}{X838515278587FF01}{}
{\noindent\textcolor{FuncColor}{$\triangleright$\enspace\texttt{Norm({\mdseries\slshape [L, [K, ]]z})\index{Norm@\texttt{Norm}}
\label{Norm}
}\hfill{\scriptsize (attribute)}}\\


 \texttt{Norm} returns the norm of the field element \mbox{\texttt{\mdseries\slshape z}}. If two fields \mbox{\texttt{\mdseries\slshape L}} and \mbox{\texttt{\mdseries\slshape K}} are given then the norm is computed w.r.t.{\nobreakspace}the field extension \mbox{\texttt{\mdseries\slshape L}}$ > $\mbox{\texttt{\mdseries\slshape K}}, if only one field \mbox{\texttt{\mdseries\slshape L}} is given then \texttt{LeftActingDomain( \mbox{\texttt{\mdseries\slshape L}} )} is taken as default for the subfield \mbox{\texttt{\mdseries\slshape K}}, and if no field is given then \texttt{DefaultField( \mbox{\texttt{\mdseries\slshape z}} )} is taken as default for \mbox{\texttt{\mdseries\slshape L}}. }

 
\subsection{\textcolor{Chapter }{Traces of field elements and matrices}}\logpage{[ 58, 3, 5 ]}
\hyperdef{L}{X7DD17EB581200AD6}{}
{
\noindent\textcolor{FuncColor}{$\triangleright$\enspace\texttt{Trace({\mdseries\slshape [L, [K, ]]z})\index{Trace@\texttt{Trace}!for a field element}
\label{Trace:for a field element}
}\hfill{\scriptsize (attribute)}}\\
\noindent\textcolor{FuncColor}{$\triangleright$\enspace\texttt{Trace({\mdseries\slshape mat})\index{Trace@\texttt{Trace}!for a matrix}
\label{Trace:for a matrix}
}\hfill{\scriptsize (attribute)}}\\


 \texttt{Trace} returns the trace of the field element \mbox{\texttt{\mdseries\slshape z}}. If two fields \mbox{\texttt{\mdseries\slshape L}} and \mbox{\texttt{\mdseries\slshape K}} are given then the trace is computed w.r.t.{\nobreakspace}the field extension $\mbox{\texttt{\mdseries\slshape L}} > \mbox{\texttt{\mdseries\slshape K}}$, if only one field \mbox{\texttt{\mdseries\slshape L}} is given then \texttt{LeftActingDomain( \mbox{\texttt{\mdseries\slshape L}} )} is taken as default for the subfield \mbox{\texttt{\mdseries\slshape K}}, and if no field is given then \texttt{DefaultField( \mbox{\texttt{\mdseries\slshape z}} )} is taken as default for \mbox{\texttt{\mdseries\slshape L}}. 

 The \emph{trace of a matrix} is the sum of its diagonal entries. Note that this is \emph{not} compatible with the definition of \texttt{Trace} for field elements, so the one\texttt{\symbol{45}}argument version is not
suitable when matrices shall be regarded as field elements.   }

 

\subsection{\textcolor{Chapter }{Conjugates}}
\logpage{[ 58, 3, 6 ]}\nobreak
\hyperdef{L}{X837A4A5781F8EE92}{}
{\noindent\textcolor{FuncColor}{$\triangleright$\enspace\texttt{Conjugates({\mdseries\slshape [L, [K, ]]z})\index{Conjugates@\texttt{Conjugates}}
\label{Conjugates}
}\hfill{\scriptsize (attribute)}}\\


 \texttt{Conjugates} returns the list of \emph{conjugates} of the field element \mbox{\texttt{\mdseries\slshape z}}. If two fields \mbox{\texttt{\mdseries\slshape L}} and \mbox{\texttt{\mdseries\slshape K}} are given then the conjugates are computed w.r.t.{\nobreakspace}the field
extension \mbox{\texttt{\mdseries\slshape L}}$ > $\mbox{\texttt{\mdseries\slshape K}}, if only one field \mbox{\texttt{\mdseries\slshape L}} is given then \texttt{LeftActingDomain( \mbox{\texttt{\mdseries\slshape L}} )} is taken as default for the subfield \mbox{\texttt{\mdseries\slshape K}}, and if no field is given then \texttt{DefaultField( \mbox{\texttt{\mdseries\slshape z}} )} is taken as default for \mbox{\texttt{\mdseries\slshape L}}. 

 The result list will contain duplicates if \mbox{\texttt{\mdseries\slshape z}} lies in a proper subfield of \mbox{\texttt{\mdseries\slshape L}}, or of the default field of \mbox{\texttt{\mdseries\slshape z}}, respectively. The result list need not be sorted. 

 
\begin{Verbatim}[commandchars=!@|,fontsize=\small,frame=single,label=Example]
  !gapprompt@gap>| !gapinput@Norm( E(8) );  Norm( CF(8), E(8) );|
  1
  1
  !gapprompt@gap>| !gapinput@Norm( CF(8), CF(4), E(8) );|
  -E(4)
  !gapprompt@gap>| !gapinput@Norm( AsField( CF(4), CF(8) ), E(8) );|
  -E(4)
  !gapprompt@gap>| !gapinput@Trace( E(8) );  Trace( CF(8), CF(8), E(8) );|
  0
  E(8)
  !gapprompt@gap>| !gapinput@Conjugates( CF(8), E(8) );|
  [ E(8), E(8)^3, -E(8), -E(8)^3 ]
  !gapprompt@gap>| !gapinput@Conjugates( CF(8), CF(4), E(8) );|
  [ E(8), -E(8) ]
  !gapprompt@gap>| !gapinput@Conjugates( CF(16), E(8) );|
  [ E(8), E(8)^3, -E(8), -E(8)^3, E(8), E(8)^3, -E(8), -E(8)^3 ]
\end{Verbatim}
 }

 

\subsection{\textcolor{Chapter }{NormalBase}}
\logpage{[ 58, 3, 7 ]}\nobreak
\hyperdef{L}{X8236A8B47E6AAD93}{}
{\noindent\textcolor{FuncColor}{$\triangleright$\enspace\texttt{NormalBase({\mdseries\slshape F[, elm]})\index{NormalBase@\texttt{NormalBase}}
\label{NormalBase}
}\hfill{\scriptsize (attribute)}}\\


 Let \mbox{\texttt{\mdseries\slshape F}} be a field that is a Galois extension of its subfield \texttt{LeftActingDomain( \mbox{\texttt{\mdseries\slshape F}} )}. Then \texttt{NormalBase} returns a list of elements in \mbox{\texttt{\mdseries\slshape F}} that form a normal basis of \mbox{\texttt{\mdseries\slshape F}}, that is, a vector space basis that is closed under the action of the Galois
group (see{\nobreakspace}\texttt{GaloisGroup} (\ref{GaloisGroup:of field})) of \mbox{\texttt{\mdseries\slshape F}}. 

 If a second argument \mbox{\texttt{\mdseries\slshape elm}} is given, it is used as a hint for the algorithm to find a normal basis with
the algorithm described in{\nobreakspace}\cite{Art68}. 
\begin{Verbatim}[commandchars=!@|,fontsize=\small,frame=single,label=Example]
  !gapprompt@gap>| !gapinput@NormalBase( CF(5) );|
  [ -E(5), -E(5)^2, -E(5)^3, -E(5)^4 ]
  !gapprompt@gap>| !gapinput@NormalBase( CF(4) );|
  [ 1/2-1/2*E(4), 1/2+1/2*E(4) ]
  !gapprompt@gap>| !gapinput@NormalBase( GF(3^6) );|
  [ Z(3^6)^2, Z(3^6)^6, Z(3^6)^18, Z(3^6)^54, Z(3^6)^162, Z(3^6)^486 ]
  !gapprompt@gap>| !gapinput@NormalBase( GF( GF(8), 2 ) );|
  [ Z(2^6), Z(2^6)^8 ]
\end{Verbatim}
 }

 }

     }

            
\chapter{\textcolor{Chapter }{Finite Fields}}\label{Finite Fields}
\logpage{[ 59, 0, 0 ]}
\hyperdef{L}{X7893ABF67A028802}{}
{
  This chapter describes the special functionality which exists in \textsf{GAP} for finite fields and their elements. Of course the general functionality for
fields (see Chapter{\nobreakspace}\ref{Fields and Division Rings}) also applies to finite fields. 

 In the following, the term \emph{finite field element} is used to denote \textsf{GAP} objects in the category \texttt{IsFFE} (\ref{IsFFE}), and \emph{finite field} means a field consisting of such elements. Note that in principle we must
distinguish these fields from (abstract) finite fields. For example, the image
of the embedding of a finite field into a field of rational functions in the
same characteristic is of course a finite field but its elements are not in \texttt{IsFFE} (\ref{IsFFE}), and in fact \textsf{GAP} does currently not support such fields. 

 Special representations exist for row vectors and matrices over small finite
fields (see sections{\nobreakspace}\ref{Row Vectors over Finite Fields} and{\nobreakspace}\ref{Matrices over Finite Fields}).  
\section{\textcolor{Chapter }{Finite Field Elements}}\label{Finite Field Elements}
\logpage{[ 59, 1, 0 ]}
\hyperdef{L}{X7B9DCCCC83400B47}{}
{
  

\subsection{\textcolor{Chapter }{IsFFE}}
\logpage{[ 59, 1, 1 ]}\nobreak
\hyperdef{L}{X7D3DF32C84FEBD25}{}
{\noindent\textcolor{FuncColor}{$\triangleright$\enspace\texttt{IsFFE({\mdseries\slshape obj})\index{IsFFE@\texttt{IsFFE}}
\label{IsFFE}
}\hfill{\scriptsize (Category)}}\\
\noindent\textcolor{FuncColor}{$\triangleright$\enspace\texttt{IsFFECollection({\mdseries\slshape obj})\index{IsFFECollection@\texttt{IsFFECollection}}
\label{IsFFECollection}
}\hfill{\scriptsize (Category)}}\\
\noindent\textcolor{FuncColor}{$\triangleright$\enspace\texttt{IsFFECollColl({\mdseries\slshape obj})\index{IsFFECollColl@\texttt{IsFFECollColl}}
\label{IsFFECollColl}
}\hfill{\scriptsize (Category)}}\\
\noindent\textcolor{FuncColor}{$\triangleright$\enspace\texttt{IsFFECollCollColl({\mdseries\slshape obj})\index{IsFFECollCollColl@\texttt{IsFFECollCollColl}}
\label{IsFFECollCollColl}
}\hfill{\scriptsize (Category)}}\\


 Objects in the category \texttt{IsFFE} are used to implement elements of finite fields. In this manual, the term \emph{finite field element} always means an object in \texttt{IsFFE}. All finite field elements of the same characteristic form a family in \textsf{GAP} (see{\nobreakspace}\ref{Families}). Any collection of finite field elements of the same characteristic
(see{\nobreakspace}\texttt{IsCollection} (\ref{IsCollection})) lies in \texttt{IsFFECollection}, and a collection of such collections (e.g., a matrix of finite field
elements) lies in \texttt{IsFFECollColl}. }

 

\subsection{\textcolor{Chapter }{Z (for field size)}}
\logpage{[ 59, 1, 2 ]}\nobreak
\hyperdef{L}{X7AA52FAF7EDEDD56}{}
{\noindent\textcolor{FuncColor}{$\triangleright$\enspace\texttt{Z({\mdseries\slshape p\texttt{\symbol{94}}d})\index{Z@\texttt{Z}!for field size}
\label{Z:for field size}
}\hfill{\scriptsize (function)}}\\
\noindent\textcolor{FuncColor}{$\triangleright$\enspace\texttt{Z({\mdseries\slshape p, d})\index{Z@\texttt{Z}!for prime and degree}
\label{Z:for prime and degree}
}\hfill{\scriptsize (function)}}\\


 For creating elements of a finite field, the function \texttt{Z} can be used. The call \texttt{Z(\mbox{\texttt{\mdseries\slshape p}},\mbox{\texttt{\mdseries\slshape d}})} (alternatively \texttt{Z(\mbox{\texttt{\mdseries\slshape p}}\texttt{\symbol{94}}\mbox{\texttt{\mdseries\slshape d}})}) returns the designated generator of the multiplicative group of the finite
field with \mbox{\texttt{\mdseries\slshape p\texttt{\symbol{94}}d}} elements. \mbox{\texttt{\mdseries\slshape p}} must be a prime integer. 

 \textsf{GAP} can represent elements of all finite fields \texttt{GF(\mbox{\texttt{\mdseries\slshape p\texttt{\symbol{94}}d}})} such that either (1) \mbox{\texttt{\mdseries\slshape p\texttt{\symbol{94}}d}} $<= 65536$ (in which case an extremely efficient internal representation is used); (2) d
= 1, (in which case, for large \mbox{\texttt{\mdseries\slshape p}}, the field is represented using the machinery of residue class rings (see
section{\nobreakspace}\ref{Residue Class Rings}) or (3) if the Conway polynomial of degree \mbox{\texttt{\mdseries\slshape d}} over the field with \mbox{\texttt{\mdseries\slshape p}} elements is known, or can be computed (see \texttt{ConwayPolynomial} (\ref{ConwayPolynomial})). 

 If you attempt to construct an element of \texttt{GF(\mbox{\texttt{\mdseries\slshape p\texttt{\symbol{94}}d}})} for which \mbox{\texttt{\mdseries\slshape d}} $> 1$ and the relevant Conway polynomial is not known, and not necessarily easy to
find (see \texttt{IsCheapConwayPolynomial} (\ref{IsCheapConwayPolynomial})), then \textsf{GAP} will stop with an error and enter the break loop. If you leave this break loop
by entering \texttt{return;} \textsf{GAP} will attempt to compute the Conway polynomial, which may take a very long
time. 

 The root returned by \texttt{Z} is a generator of the multiplicative group of the finite field with \mbox{\texttt{\mdseries\slshape p\texttt{\symbol{94}}d}} elements, which is cyclic. The order of the element is of course \mbox{\texttt{\mdseries\slshape p\texttt{\symbol{94}}d}} $-1$. The \mbox{\texttt{\mdseries\slshape p\texttt{\symbol{94}}d}} $-1$ different powers of the root are exactly the nonzero elements of the finite
field. 

 Thus all nonzero elements of the finite field with \mbox{\texttt{\mdseries\slshape p\texttt{\symbol{94}}d}} elements can be entered as \texttt{Z(\mbox{\texttt{\mdseries\slshape p\texttt{\symbol{94}}d}})\texttt{\symbol{94}}}$i$. Note that this is also the form that \textsf{GAP} uses to output those elements when they are stored in the internal
representation. In larger fields, it is more convenient to enter and print
elements as linear combinations of powers of the primitive element, see
section \ref{Printing, Viewing and Displaying Finite Field Elements}. 

 The additive neutral element is \texttt{0 * Z(\mbox{\texttt{\mdseries\slshape p}})}. It is different from the integer \texttt{0} in subtle ways. First \texttt{IsInt( 0 * Z(\mbox{\texttt{\mdseries\slshape p}}) )} (see \texttt{IsInt} (\ref{IsInt})) is \texttt{false} and \texttt{IsFFE( 0 * Z(\mbox{\texttt{\mdseries\slshape p}}) )} (see \texttt{IsFFE} (\ref{IsFFE})) is \texttt{true}, whereas it is just the other way around for the integer \texttt{0}. 

 The multiplicative neutral element is \texttt{Z(\mbox{\texttt{\mdseries\slshape p}})\texttt{\symbol{94}}0}. It is different from the integer \texttt{1} in subtle ways. First \texttt{IsInt( Z(\mbox{\texttt{\mdseries\slshape p}})\texttt{\symbol{94}}0 )} (see \texttt{IsInt} (\ref{IsInt})) is \texttt{false} and \texttt{IsFFE( Z(\mbox{\texttt{\mdseries\slshape p}})\texttt{\symbol{94}}0 )} (see \texttt{IsFFE} (\ref{IsFFE})) is \texttt{true}, whereas it is just the other way around for the integer \texttt{1}. Also \texttt{1+1} is \texttt{2}, whereas, e.g., \texttt{Z(2)\texttt{\symbol{94}}0 + Z(2)\texttt{\symbol{94}}0} is \texttt{0 * Z(2)}. 

 The various roots returned by \texttt{Z} for finite fields of the same characteristic are compatible in the following
sense. If the field \texttt{GF(\mbox{\texttt{\mdseries\slshape p}},}$n$\texttt{)} is a subfield of the field \texttt{GF(\mbox{\texttt{\mdseries\slshape p}},}$m$\texttt{)}, i.e., $n$ divides $m$, then \texttt{Z}$(\mbox{\texttt{\mdseries\slshape p}}^n) = $\texttt{Z}$(\mbox{\texttt{\mdseries\slshape p}}^m)^{{(\mbox{\texttt{\mdseries\slshape p}}^m-1)/(\mbox{\texttt{\mdseries\slshape p}}^n-1)}}$. Note that this is the simplest relation that may hold between a generator of \texttt{GF(\mbox{\texttt{\mdseries\slshape p}},}$n$\texttt{)} and \texttt{GF(\mbox{\texttt{\mdseries\slshape p}},}$m$\texttt{)}, since \texttt{Z}$(\mbox{\texttt{\mdseries\slshape p}}^n)$ is an element of order $\mbox{\texttt{\mdseries\slshape p}}^m-1$ and \texttt{Z}$(\mbox{\texttt{\mdseries\slshape p}}^m)$ is an element of order $\mbox{\texttt{\mdseries\slshape p}}^n-1$. This is achieved by choosing \texttt{Z(\mbox{\texttt{\mdseries\slshape p}})} as the smallest primitive root modulo \mbox{\texttt{\mdseries\slshape p}} and \texttt{Z(}\mbox{\texttt{\mdseries\slshape p\texttt{\symbol{94}}n}}\texttt{)} as a root of the $n$\texttt{\symbol{45}}th \emph{Conway polynomial} (see{\nobreakspace}\texttt{ConwayPolynomial} (\ref{ConwayPolynomial})) of characteristic \mbox{\texttt{\mdseries\slshape p}}. Those polynomials were defined by J.{\nobreakspace}H.{\nobreakspace}Conway,
and many of them were computed by R.{\nobreakspace}A.{\nobreakspace}Parker. 

 
\begin{Verbatim}[commandchars=!@|,fontsize=\small,frame=single,label=Example]
  !gapprompt@gap>| !gapinput@a:= Z( 32 );|
  Z(2^5)
  !gapprompt@gap>| !gapinput@a+a;|
  0*Z(2)
  !gapprompt@gap>| !gapinput@a*a;|
  Z(2^5)^2
  !gapprompt@gap>| !gapinput@b := Z(3,12);|
  z
  !gapprompt@gap>| !gapinput@b*b;|
  z2
  !gapprompt@gap>| !gapinput@b+b;|
  2z
  !gapprompt@gap>| !gapinput@Print(b^100,"\n");|
  Z(3)^0+Z(3,12)^5+Z(3,12)^6+2*Z(3,12)^8+Z(3,12)^10+Z(3,12)^11
\end{Verbatim}
 
\begin{Verbatim}[commandchars=!|B,fontsize=\small,frame=single,label=Example]
  !gapprompt|gap>B !gapinput|Z(11,40);B
  Error, Conway Polynomial 11^40 will need to be computed and might be slow
  Stack trace:
  *[1] Error( "Conway Polynomial ", p, "^", d, " will need to be computed and might be slow" );
     @ GAPROOT/lib/ffeconway.gi:81
   [2] FFECONWAY.SetUpConwayStuff( p, d );
     @ GAPROOT/lib/ffeconway.gi:140
   [3] FFECONWAY.ZNC( p, d )
     @ GAPROOT/lib/ffeconway.gi:167
  <function "ZOp method">( <arguments> )
   called from read-eval loop at *stdin*:2
  you can enter 'quit;' to quit to outer loop, or
  you can enter 'return;' to continue
  !gapbrkprompt|brk>B !gapinput|B
\end{Verbatim}
 }

 

\subsection{\textcolor{Chapter }{IsLexOrderedFFE}}
\logpage{[ 59, 1, 3 ]}\nobreak
\hyperdef{L}{X8612BCEA816CF1B9}{}
{\noindent\textcolor{FuncColor}{$\triangleright$\enspace\texttt{IsLexOrderedFFE({\mdseries\slshape ffe})\index{IsLexOrderedFFE@\texttt{IsLexOrderedFFE}}
\label{IsLexOrderedFFE}
}\hfill{\scriptsize (Category)}}\\
\noindent\textcolor{FuncColor}{$\triangleright$\enspace\texttt{IsLogOrderedFFE({\mdseries\slshape ffe})\index{IsLogOrderedFFE@\texttt{IsLogOrderedFFE}}
\label{IsLogOrderedFFE}
}\hfill{\scriptsize (Category)}}\\


 Elements of finite fields can be compared using the operators \texttt{=} and \texttt{{\textless}}. The call \texttt{\mbox{\texttt{\mdseries\slshape a}} = \mbox{\texttt{\mdseries\slshape b}}} returns \texttt{true} if and only if the finite field elements \mbox{\texttt{\mdseries\slshape a}} and \mbox{\texttt{\mdseries\slshape b}} are equal. Furthermore \texttt{\mbox{\texttt{\mdseries\slshape a}} {\textless} \mbox{\texttt{\mdseries\slshape b}}} tests whether \mbox{\texttt{\mdseries\slshape a}} is smaller than \mbox{\texttt{\mdseries\slshape b}}. The exact behaviour of this comparison depends on which of two categories
the field elements belong to: 

 Finite field elements are ordered in \textsf{GAP} (by \texttt{\texttt{\symbol{92}}{\textless}} (\ref{bSlash<})) first by characteristic and then by their degree (i.e. the sizes of the
smallest fields containing them). Amongst irreducible elements of a given
field, the ordering depends on which of these categories the elements of the
field belong to (all irreducible elements of a given field should belong to
the same one) 

 Elements in \texttt{IsLexOrderedFFE} are ordered lexicographically by their coefficients with respect to the
canonical basis of the field. 

 Elements in \texttt{IsLogOrderedFFE} are ordered according to their discrete logarithms with respect to the \texttt{PrimitiveRoot} (\ref{PrimitiveRoot}) attribute of the field. For the comparison of finite field elements with other \textsf{GAP} objects, see{\nobreakspace}\ref{Comparisons}. 

 
\begin{Verbatim}[commandchars=!@|,fontsize=\small,frame=single,label=Example]
  !gapprompt@gap>| !gapinput@Z( 16 )^10 = Z( 4 )^2;  # illustrates embedding of GF(4) in GF(16)|
  true
  !gapprompt@gap>| !gapinput@0 < 0*Z(101);|
  true
  !gapprompt@gap>| !gapinput@Z(256) > Z(101);|
  false
  !gapprompt@gap>| !gapinput@Z(2,20) < Z(2,20)^2; # this illustrates the lexicographic ordering|
  false
\end{Verbatim}
 }

 }

  
\section{\textcolor{Chapter }{Operations for Finite Field Elements}}\label{Operations for Finite Field Elements}
\logpage{[ 59, 2, 0 ]}
\hyperdef{L}{X7A79399283EF78D0}{}
{
  Since finite field elements are scalars, the operations \texttt{Characteristic} (\ref{Characteristic}), \texttt{One} (\ref{One}), \texttt{Zero} (\ref{Zero}), \texttt{Inverse} (\ref{Inverse}), \texttt{AdditiveInverse} (\ref{AdditiveInverse}), \texttt{Order} (\ref{Order}) can be applied to them (see{\nobreakspace}\ref{Attributes and Properties of Elements}). Contrary to the situation with other scalars, \texttt{Order} (\ref{Order}) is defined also for the zero element in a finite field, with value \texttt{0}. 

 
\begin{Verbatim}[commandchars=!@|,fontsize=\small,frame=single,label=Example]
  !gapprompt@gap>| !gapinput@Characteristic( Z( 16 )^10 );  Characteristic( Z( 9 )^2 );|
  2
  3
  !gapprompt@gap>| !gapinput@Characteristic( [ Z(4), Z(8) ] );|
  2
  !gapprompt@gap>| !gapinput@One( Z(9) );  One( 0*Z(4) );|
  Z(3)^0
  Z(2)^0
  !gapprompt@gap>| !gapinput@Inverse( Z(9) );  AdditiveInverse( Z(9) );|
  Z(3^2)^7
  Z(3^2)^5
  !gapprompt@gap>| !gapinput@Order( Z(9)^7 );|
  8
\end{Verbatim}
 

\subsection{\textcolor{Chapter }{DegreeFFE (for a FFE)}}
\logpage{[ 59, 2, 1 ]}\nobreak
\hyperdef{L}{X828E846E7C1EA3DD}{}
{\noindent\textcolor{FuncColor}{$\triangleright$\enspace\texttt{DegreeFFE({\mdseries\slshape z})\index{DegreeFFE@\texttt{DegreeFFE}!for a FFE}
\label{DegreeFFE:for a FFE}
}\hfill{\scriptsize (attribute)}}\\
\noindent\textcolor{FuncColor}{$\triangleright$\enspace\texttt{DegreeFFE({\mdseries\slshape vec})\index{DegreeFFE@\texttt{DegreeFFE}!for a vector of FFEs}
\label{DegreeFFE:for a vector of FFEs}
}\hfill{\scriptsize (method)}}\\
\noindent\textcolor{FuncColor}{$\triangleright$\enspace\texttt{DegreeFFE({\mdseries\slshape mat})\index{DegreeFFE@\texttt{DegreeFFE}!for a matrix of FFEs}
\label{DegreeFFE:for a matrix of FFEs}
}\hfill{\scriptsize (method)}}\\


 \texttt{DegreeFFE} returns the degree of the smallest finite field \mbox{\texttt{\mdseries\slshape F}} containing the element \mbox{\texttt{\mdseries\slshape z}}, respectively all elements of the row vector \mbox{\texttt{\mdseries\slshape vec}} over a finite field (see{\nobreakspace}\ref{Row Vectors}), or the matrix \mbox{\texttt{\mdseries\slshape mat}} over a finite field (see{\nobreakspace}\ref{Matrices}). 

 
\begin{Verbatim}[commandchars=!@|,fontsize=\small,frame=single,label=Example]
  !gapprompt@gap>| !gapinput@DegreeFFE( Z( 16 )^10 );|
  2
  !gapprompt@gap>| !gapinput@DegreeFFE( Z( 16 )^11 );|
  4
  !gapprompt@gap>| !gapinput@DegreeFFE( [ Z(2^13), Z(2^10) ] );|
  130
\end{Verbatim}
 }

 

\subsection{\textcolor{Chapter }{LogFFE}}
\logpage{[ 59, 2, 2 ]}\nobreak
\hyperdef{L}{X7B049A3478B369E4}{}
{\noindent\textcolor{FuncColor}{$\triangleright$\enspace\texttt{LogFFE({\mdseries\slshape z, r})\index{LogFFE@\texttt{LogFFE}}
\label{LogFFE}
}\hfill{\scriptsize (operation)}}\\


 \texttt{LogFFE} returns the discrete logarithm of the element \mbox{\texttt{\mdseries\slshape z}} in a finite field with respect to the root \mbox{\texttt{\mdseries\slshape r}}. An error is signalled if \mbox{\texttt{\mdseries\slshape z}} is zero. \texttt{fail} is returned if \mbox{\texttt{\mdseries\slshape z}} is not a power of \mbox{\texttt{\mdseries\slshape r}}. 

 The \emph{discrete logarithm} of the element \mbox{\texttt{\mdseries\slshape z}} with respect to the root \mbox{\texttt{\mdseries\slshape r}} is the smallest nonnegative integer $i$ such that $\mbox{\texttt{\mdseries\slshape r}}^i = \mbox{\texttt{\mdseries\slshape z}}$ holds. 

 
\begin{Verbatim}[commandchars=!@|,fontsize=\small,frame=single,label=Example]
  !gapprompt@gap>| !gapinput@LogFFE( Z(409)^116, Z(409) );  LogFFE( Z(409)^116, Z(409)^2 );|
  116
  58
\end{Verbatim}
 }

 

\subsection{\textcolor{Chapter }{IntFFE}}
\logpage{[ 59, 2, 3 ]}\nobreak
\hyperdef{L}{X79F48E337FC2746A}{}
{\noindent\textcolor{FuncColor}{$\triangleright$\enspace\texttt{IntFFE({\mdseries\slshape z})\index{IntFFE@\texttt{IntFFE}}
\label{IntFFE}
}\hfill{\scriptsize (attribute)}}\\
\noindent\textcolor{FuncColor}{$\triangleright$\enspace\texttt{Int({\mdseries\slshape z})\index{Int@\texttt{Int}!for a FFE}
\label{Int:for a FFE}
}\hfill{\scriptsize (method)}}\\


 \texttt{IntFFE} returns the integer corresponding to the element \mbox{\texttt{\mdseries\slshape z}}, which must lie in a finite prime field. That is, \texttt{IntFFE} returns the smallest nonnegative integer $i$ such that $i$\texttt{ * One( }\mbox{\texttt{\mdseries\slshape z}}\texttt{ ) = }\mbox{\texttt{\mdseries\slshape z}}. 

 The correspondence between elements from a finite prime field of
characteristic $p$ (for $p < 2^{16}$) and the integers between $0$ and $p-1$ is defined by choosing \texttt{Z(}$p$\texttt{)} the element corresponding to the smallest primitive root mod $p$ (see{\nobreakspace}\texttt{PrimitiveRootMod} (\ref{PrimitiveRootMod})). 

 \texttt{IntFFE} is installed as a method for the operation \texttt{Int} (\ref{Int}) with argument a finite field element. 

 
\begin{Verbatim}[commandchars=!@|,fontsize=\small,frame=single,label=Example]
  !gapprompt@gap>| !gapinput@IntFFE( Z(13) );  PrimitiveRootMod( 13 );|
  2
  2
  !gapprompt@gap>| !gapinput@IntFFE( Z(409) );|
  21
  !gapprompt@gap>| !gapinput@IntFFE( Z(409)^116 );  21^116 mod 409;|
  311
  311
\end{Verbatim}
 See also \texttt{IntFFESymm} (\ref{IntFFESymm:for a FFE}). }

 

\subsection{\textcolor{Chapter }{IntFFESymm (for a FFE)}}
\logpage{[ 59, 2, 4 ]}\nobreak
\hyperdef{L}{X7DABD827848BCC2A}{}
{\noindent\textcolor{FuncColor}{$\triangleright$\enspace\texttt{IntFFESymm({\mdseries\slshape z})\index{IntFFESymm@\texttt{IntFFESymm}!for a FFE}
\label{IntFFESymm:for a FFE}
}\hfill{\scriptsize (attribute)}}\\
\noindent\textcolor{FuncColor}{$\triangleright$\enspace\texttt{IntFFESymm({\mdseries\slshape vec})\index{IntFFESymm@\texttt{IntFFESymm}!for a vector of FFEs}
\label{IntFFESymm:for a vector of FFEs}
}\hfill{\scriptsize (attribute)}}\\


 For a finite prime field element \mbox{\texttt{\mdseries\slshape z}}, \texttt{IntFFESymm} returns the corresponding integer of smallest absolute value. That is, \texttt{IntFFESymm} returns the integer $i$ of smallest absolute value such that $i$\texttt{ * One( }\mbox{\texttt{\mdseries\slshape z}}\texttt{ ) = }\mbox{\texttt{\mdseries\slshape z}} holds. 

 For a vector \mbox{\texttt{\mdseries\slshape vec}} of FFEs, the operation returns the result of applying \texttt{IntFFESymm} to every entry of the vector. 

 The correspondence between elements from a finite prime field of
characteristic $p$ (for $p < 2^{16}$) and the integers between $-p/2$ and $p/2$ is defined by choosing \texttt{Z(}$p$\texttt{)} the element corresponding to the smallest positive primitive root mod $p$ (see{\nobreakspace}\texttt{PrimitiveRootMod} (\ref{PrimitiveRootMod})) and reducing results to the $-p/2 .. p/2$ range. 

 
\begin{Verbatim}[commandchars=!@|,fontsize=\small,frame=single,label=Example]
  !gapprompt@gap>| !gapinput@IntFFE(Z(13)^2);IntFFE(Z(13)^3);|
  4
  8
  !gapprompt@gap>| !gapinput@IntFFESymm(Z(13)^2);IntFFESymm(Z(13)^3);|
  4
  -5
\end{Verbatim}
 See also \texttt{IntFFE} (\ref{IntFFE}) }

 

\subsection{\textcolor{Chapter }{IntVecFFE}}
\logpage{[ 59, 2, 5 ]}\nobreak
\hyperdef{L}{X8009968782F18888}{}
{\noindent\textcolor{FuncColor}{$\triangleright$\enspace\texttt{IntVecFFE({\mdseries\slshape vecffe})\index{IntVecFFE@\texttt{IntVecFFE}}
\label{IntVecFFE}
}\hfill{\scriptsize (operation)}}\\


 is the list of integers corresponding to the vector \mbox{\texttt{\mdseries\slshape vecffe}} of finite field elements in a prime field (see{\nobreakspace}\texttt{IntFFE} (\ref{IntFFE})). }

 

\subsection{\textcolor{Chapter }{AsInternalFFE}}
\logpage{[ 59, 2, 6 ]}\nobreak
\hyperdef{L}{X807959EE82CED148}{}
{\noindent\textcolor{FuncColor}{$\triangleright$\enspace\texttt{AsInternalFFE({\mdseries\slshape ffe})\index{AsInternalFFE@\texttt{AsInternalFFE}}
\label{AsInternalFFE}
}\hfill{\scriptsize (attribute)}}\\


 return an internal FFE equal to \mbox{\texttt{\mdseries\slshape ffe}} if one exists, otherwise \texttt{fail} }

 

\subsection{\textcolor{Chapter }{RootFFE}}
\logpage{[ 59, 2, 7 ]}\nobreak
\hyperdef{L}{X794AEB148410825B}{}
{\noindent\textcolor{FuncColor}{$\triangleright$\enspace\texttt{RootFFE({\mdseries\slshape F, z, k})\index{RootFFE@\texttt{RootFFE}}
\label{RootFFE}
}\hfill{\scriptsize (operation)}}\\


 \texttt{RootFFE} returns a finite field element \mbox{\texttt{\mdseries\slshape r}} from \mbox{\texttt{\mdseries\slshape F}} whose \mbox{\texttt{\mdseries\slshape k}}\texttt{\symbol{45}}th power is \mbox{\texttt{\mdseries\slshape z}}. If no such element exists then \texttt{fail} is returned. }

 }

  
\section{\textcolor{Chapter }{Creating Finite Fields}}\label{Creating Finite Fields}
\logpage{[ 59, 3, 0 ]}
\hyperdef{L}{X81B54A8378734C33}{}
{
  

\subsection{\textcolor{Chapter }{DefaultField (for finite field elements)}}
\logpage{[ 59, 3, 1 ]}\nobreak
\hyperdef{L}{X7979F51D7C43AB05}{}
{\noindent\textcolor{FuncColor}{$\triangleright$\enspace\texttt{DefaultField({\mdseries\slshape list})\index{DefaultField@\texttt{DefaultField}!for finite field elements}
\label{DefaultField:for finite field elements}
}\hfill{\scriptsize (function)}}\\
\noindent\textcolor{FuncColor}{$\triangleright$\enspace\texttt{DefaultRing({\mdseries\slshape list})\index{DefaultRing@\texttt{DefaultRing}!for finite field elements}
\label{DefaultRing:for finite field elements}
}\hfill{\scriptsize (function)}}\\


 \texttt{DefaultField} and \texttt{DefaultRing} for finite field elements are defined to return the \emph{smallest} field containing the given elements. 

 
\begin{Verbatim}[commandchars=!@|,fontsize=\small,frame=single,label=Example]
  !gapprompt@gap>| !gapinput@DefaultField( [ Z(4), Z(4)^2 ] );  DefaultField( [ Z(4), Z(8) ] );|
  GF(2^2)
  GF(2^6)
\end{Verbatim}
 }

 

\subsection{\textcolor{Chapter }{GaloisField (for field size)}}
\logpage{[ 59, 3, 2 ]}\nobreak
\hyperdef{L}{X8592DBB086A8A9BE}{}
{\noindent\textcolor{FuncColor}{$\triangleright$\enspace\texttt{GaloisField({\mdseries\slshape p\texttt{\symbol{94}}d})\index{GaloisField@\texttt{GaloisField}!for field size}
\label{GaloisField:for field size}
}\hfill{\scriptsize (function)}}\\
\noindent\textcolor{FuncColor}{$\triangleright$\enspace\texttt{GF({\mdseries\slshape p\texttt{\symbol{94}}d})\index{GF@\texttt{GF}!for field size}
\label{GF:for field size}
}\hfill{\scriptsize (function)}}\\
\noindent\textcolor{FuncColor}{$\triangleright$\enspace\texttt{GaloisField({\mdseries\slshape p, d})\index{GaloisField@\texttt{GaloisField}!for characteristic and degree}
\label{GaloisField:for characteristic and degree}
}\hfill{\scriptsize (function)}}\\
\noindent\textcolor{FuncColor}{$\triangleright$\enspace\texttt{GF({\mdseries\slshape p, d})\index{GF@\texttt{GF}!for characteristic and degree}
\label{GF:for characteristic and degree}
}\hfill{\scriptsize (function)}}\\
\noindent\textcolor{FuncColor}{$\triangleright$\enspace\texttt{GaloisField({\mdseries\slshape subfield, d})\index{GaloisField@\texttt{GaloisField}!for subfield and degree}
\label{GaloisField:for subfield and degree}
}\hfill{\scriptsize (function)}}\\
\noindent\textcolor{FuncColor}{$\triangleright$\enspace\texttt{GF({\mdseries\slshape subfield, d})\index{GF@\texttt{GF}!for subfield and degree}
\label{GF:for subfield and degree}
}\hfill{\scriptsize (function)}}\\
\noindent\textcolor{FuncColor}{$\triangleright$\enspace\texttt{GaloisField({\mdseries\slshape p, pol})\index{GaloisField@\texttt{GaloisField}!for characteristic and polynomial}
\label{GaloisField:for characteristic and polynomial}
}\hfill{\scriptsize (function)}}\\
\noindent\textcolor{FuncColor}{$\triangleright$\enspace\texttt{GF({\mdseries\slshape p, pol})\index{GF@\texttt{GF}!for characteristic and polynomial}
\label{GF:for characteristic and polynomial}
}\hfill{\scriptsize (function)}}\\
\noindent\textcolor{FuncColor}{$\triangleright$\enspace\texttt{GaloisField({\mdseries\slshape subfield, pol})\index{GaloisField@\texttt{GaloisField}!for subfield and polynomial}
\label{GaloisField:for subfield and polynomial}
}\hfill{\scriptsize (function)}}\\
\noindent\textcolor{FuncColor}{$\triangleright$\enspace\texttt{GF({\mdseries\slshape subfield, pol})\index{GF@\texttt{GF}!for subfield and polynomial}
\label{GF:for subfield and polynomial}
}\hfill{\scriptsize (function)}}\\


 \texttt{GaloisField} returns a finite field. It takes two arguments. The form \texttt{GaloisField( \mbox{\texttt{\mdseries\slshape p}}, \mbox{\texttt{\mdseries\slshape d}} )}, where \mbox{\texttt{\mdseries\slshape p}}, \mbox{\texttt{\mdseries\slshape d}} are integers, can also be given as \texttt{GaloisField( \mbox{\texttt{\mdseries\slshape p}}\texttt{\symbol{94}}\mbox{\texttt{\mdseries\slshape d}} )}. \texttt{GF} is an abbreviation for \texttt{GaloisField}. 

 The first argument specifies the subfield $S$ over which the new field is to be taken. It can be a prime integer or a finite
field. If it is a prime \mbox{\texttt{\mdseries\slshape p}}, the subfield is the prime field of this characteristic. 

 The second argument specifies the extension. It can be an integer or an
irreducible polynomial over the field $S$. If it is an integer \mbox{\texttt{\mdseries\slshape d}}, the new field is constructed as the polynomial extension w.r.t. the Conway
polynomial (see{\nobreakspace}\texttt{ConwayPolynomial} (\ref{ConwayPolynomial})) of degree \mbox{\texttt{\mdseries\slshape d}} over $S$. If it is an irreducible polynomial \mbox{\texttt{\mdseries\slshape pol}} over $S$, the new field is constructed as polynomial extension of $S$ with this polynomial; in this case, \mbox{\texttt{\mdseries\slshape pol}} is accessible as the value of \texttt{DefiningPolynomial} (\ref{DefiningPolynomial}) for the new field, and a root of \mbox{\texttt{\mdseries\slshape pol}} in the new field is accessible as the value of \texttt{RootOfDefiningPolynomial} (\ref{RootOfDefiningPolynomial}). 

 Note that the subfield over which a field was constructed determines over
which field the Galois group, conjugates, norm, trace, minimal polynomial, and
trace polynomial are computed (see{\nobreakspace}\texttt{GaloisGroup} (\ref{GaloisGroup:of field}), \texttt{Conjugates} (\ref{Conjugates}), \texttt{Norm} (\ref{Norm}), \texttt{Trace} (\ref{Trace:for a field element}), \texttt{MinimalPolynomial} (\ref{MinimalPolynomial:over a field}), \texttt{TracePolynomial} (\ref{TracePolynomial})). 

 The field is regarded as a vector space (see{\nobreakspace}\ref{Vector Spaces}) over the given subfield, so this determines the dimension and the canonical
basis of the field. 

 
\begin{Verbatim}[commandchars=!@|,fontsize=\small,frame=single,label=Example]
  !gapprompt@gap>| !gapinput@f1:= GF( 2^4 );|
  GF(2^4)
  !gapprompt@gap>| !gapinput@Size( GaloisGroup ( f1 ) );|
  4
  !gapprompt@gap>| !gapinput@BasisVectors( Basis( f1 ) );|
  [ Z(2)^0, Z(2^4), Z(2^4)^2, Z(2^4)^3 ]
  !gapprompt@gap>| !gapinput@f2:= GF( GF(4), 2 );|
  AsField( GF(2^2), GF(2^4) )
  !gapprompt@gap>| !gapinput@Size( GaloisGroup( f2 ) );|
  2
  !gapprompt@gap>| !gapinput@BasisVectors( Basis( f2 ) );|
  [ Z(2)^0, Z(2^4) ]
\end{Verbatim}
 }

 

\subsection{\textcolor{Chapter }{PrimitiveRoot}}
\logpage{[ 59, 3, 3 ]}\nobreak
\hyperdef{L}{X788B1ECD83C70516}{}
{\noindent\textcolor{FuncColor}{$\triangleright$\enspace\texttt{PrimitiveRoot({\mdseries\slshape F})\index{PrimitiveRoot@\texttt{PrimitiveRoot}}
\label{PrimitiveRoot}
}\hfill{\scriptsize (attribute)}}\\


 A \emph{primitive root} of a finite field is a generator of its multiplicative group. A primitive root
is always a primitive element (see{\nobreakspace}\texttt{PrimitiveElement} (\ref{PrimitiveElement})), the converse is in general not true.   
\begin{Verbatim}[commandchars=!@|,fontsize=\small,frame=single,label=Example]
  !gapprompt@gap>| !gapinput@f:= GF( 3^5 );|
  GF(3^5)
  !gapprompt@gap>| !gapinput@PrimitiveRoot( f );|
  Z(3^5)
\end{Verbatim}
 }

 }

  
\section{\textcolor{Chapter }{Frobenius Automorphisms}}\label{Frobenius Automorphisms}
\logpage{[ 59, 4, 0 ]}
\hyperdef{L}{X7A5F075185CE5B06}{}
{
  

\subsection{\textcolor{Chapter }{FrobeniusAutomorphism}}
\logpage{[ 59, 4, 1 ]}\nobreak
\hyperdef{L}{X8758E4AB7D0A1955}{}
{\noindent\textcolor{FuncColor}{$\triangleright$\enspace\texttt{FrobeniusAutomorphism({\mdseries\slshape F})\index{FrobeniusAutomorphism@\texttt{FrobeniusAutomorphism}}
\label{FrobeniusAutomorphism}
}\hfill{\scriptsize (attribute)}}\\


 \index{homomorphisms!Frobenius, field} \index{field homomorphisms!Frobenius} \index{CompositionMapping@\texttt{CompositionMapping}!for Frobenius automorphisms} returns the Frobenius automorphism of the finite field \mbox{\texttt{\mdseries\slshape F}} as a field homomorphism (see{\nobreakspace}\ref{Ring Homomorphisms}). 

 \index{Frobenius automorphism} The \emph{Frobenius automorphism} $f$ of a finite field $F$ of characteristic $p$ is the function that takes each element $z$ of $F$ to its $p$\texttt{\symbol{45}}th power. Each field automorphism of $F$ is a power of $f$. Thus $f$ is a generator for the Galois group of $F$ relative to the prime field of $F$, and an appropriate power of $f$ is a generator of the Galois group of $F$ over a subfield (see{\nobreakspace}\texttt{GaloisGroup} (\ref{GaloisGroup:of field})). 

 
\begin{Verbatim}[commandchars=!@|,fontsize=\small,frame=single,label=Example]
  !gapprompt@gap>| !gapinput@f := GF(16);|
  GF(2^4)
  !gapprompt@gap>| !gapinput@x := FrobeniusAutomorphism( f );|
  FrobeniusAutomorphism( GF(2^4) )
  !gapprompt@gap>| !gapinput@Z(16) ^ x;|
  Z(2^4)^2
  !gapprompt@gap>| !gapinput@x^2;|
  FrobeniusAutomorphism( GF(2^4) )^2
\end{Verbatim}
 

 \index{Image@\texttt{Image}!for Frobenius automorphisms} The image of an element $z$ under the $i$\texttt{\symbol{45}}th power of $f$ is computed as the $p^i$\texttt{\symbol{45}}th power of $z$. The product of the $i$\texttt{\symbol{45}}th power and the $j$\texttt{\symbol{45}}th power of $f$ is the $k$\texttt{\symbol{45}}th power of $f$, where $k$ is $i j \bmod $ \texttt{Size(\mbox{\texttt{\mdseries\slshape F}})}$-1$. The zeroth power of $f$ is \texttt{IdentityMapping( \mbox{\texttt{\mdseries\slshape F}} )}. }

 }

  
\section{\textcolor{Chapter }{Conway Polynomials}}\label{Conway Polynomials}
\logpage{[ 59, 5, 0 ]}
\hyperdef{L}{X869919BB7EBE5741}{}
{
  

\subsection{\textcolor{Chapter }{ConwayPolynomial}}
\logpage{[ 59, 5, 1 ]}\nobreak
\hyperdef{L}{X7C2425A786F09054}{}
{\noindent\textcolor{FuncColor}{$\triangleright$\enspace\texttt{ConwayPolynomial({\mdseries\slshape p, n})\index{ConwayPolynomial@\texttt{ConwayPolynomial}}
\label{ConwayPolynomial}
}\hfill{\scriptsize (function)}}\\


 is the Conway polynomial of the finite field $GF(p^n)$ as polynomial over the prime field in characteristic \mbox{\texttt{\mdseries\slshape p}}. 

 The \emph{Conway polynomial} $\Phi_{{n,p}}$ of $GF(p^n)$ is defined by the following properties. 

 First define an ordering of polynomials of degree $n$ over $GF(p)$, as follows. $f = \sum_{{i = 0}}^n (-1)^i f_i x^i$ is smaller than $g = \sum_{{i = 0}}^n (-1)^i g_i x^i$ if and only if there is an index $m \leq n$ such that $f_i = g_i$ for all $i > m$, and $\tilde{{f_m}} < \tilde{{g_m}}$, where $\tilde{{c}}$ denotes the integer value in $\{ 0, 1, \ldots, p-1 \}$ that is mapped to $c \in GF(p)$ under the canonical epimorphism that maps the integers onto $GF(p)$. 

 $\Phi_{{n,p}}$ is \emph{primitive} over $GF(p)$ (see{\nobreakspace}\texttt{IsPrimitivePolynomial} (\ref{IsPrimitivePolynomial})). That is, $\Phi_{{n,p}}$ is irreducible, monic, and is the minimal polynomial of a primitive root of $GF(p^n)$. 

 For all divisors $d$ of $n$ the compatibility condition $\Phi_{{d,p}}( x^{{\frac{{p^n-1}}{{p^m-1}}}} ) \equiv 0
\pmod{{\Phi_{{n,p}}(x)}}$ holds. (That is, the appropriate power of a zero of $\Phi_{{n,p}}$ is a zero of the Conway polynomial $\Phi_{{d,p}}$.) 

 With respect to the ordering defined above, $\Phi_{{n,p}}$ shall be minimal. 

 The computation of Conway polynomials can be time consuming. Therefore, \textsf{GAP} comes with a list of precomputed polynomials. If a requested polynomial is not
stored then \textsf{GAP} prints a warning and computes it by checking all polynomials in the order
defined above for the defining conditions. If $n$ is not a prime this is probably a very long computation. (Some previously
known polynomials with prime $n$ are not stored in \textsf{GAP} because they are quickly recomputed.) Use the function \texttt{IsCheapConwayPolynomial} (\ref{IsCheapConwayPolynomial}) to check in advance if \texttt{ConwayPolynomial} will give a result after a short time. 

 Note that primitivity of a polynomial can only be checked if \textsf{GAP} can factorize $p^n-1$. A sufficiently new version of the \textsf{FactInt} package contains many precomputed factors of such numbers from various
factorization projects. 

 See{\nobreakspace}\cite{L03} for further information on known Conway polynomials. 

 An interactive overview of the Conway polynomials known to \textsf{GAP} is provided by the function \texttt{BrowseConwayPolynomials} from the \textsf{GAP} package \textsf{Browse}, see \texttt{BrowseGapData} (\textbf{Browse: BrowseGapData}). 

 \index{InfoText@\texttt{InfoText}!(for Conway polynomials)} If \mbox{\texttt{\mdseries\slshape pol}} is a result returned by \texttt{ConwayPolynomial} the command \texttt{Print( InfoText( \mbox{\texttt{\mdseries\slshape pol}} ) );} will print some info on the origin of that particular polynomial. 

 For some purposes it may be enough to have any primitive polynomial for an
extension of a finite field instead of the Conway polynomial,
see{\nobreakspace}\texttt{RandomPrimitivePolynomial} (\ref{RandomPrimitivePolynomial}) below. 
\begin{Verbatim}[commandchars=!@|,fontsize=\small,frame=single,label=Example]
  !gapprompt@gap>| !gapinput@ConwayPolynomial( 2, 5 );  ConwayPolynomial( 3, 7 );|
  x_1^5+x_1^2+Z(2)^0
  x_1^7-x_1^2+Z(3)^0
\end{Verbatim}
 }

 

\subsection{\textcolor{Chapter }{IsCheapConwayPolynomial}}
\logpage{[ 59, 5, 2 ]}\nobreak
\hyperdef{L}{X78A7C1247E129AD9}{}
{\noindent\textcolor{FuncColor}{$\triangleright$\enspace\texttt{IsCheapConwayPolynomial({\mdseries\slshape p, n})\index{IsCheapConwayPolynomial@\texttt{IsCheapConwayPolynomial}}
\label{IsCheapConwayPolynomial}
}\hfill{\scriptsize (function)}}\\


 Returns \texttt{true} if \texttt{ConwayPolynomial( \mbox{\texttt{\mdseries\slshape p}}, \mbox{\texttt{\mdseries\slshape n}} )} will give a result in \emph{reasonable} time. This is either the case when this polynomial is
pre\texttt{\symbol{45}}computed, or if \mbox{\texttt{\mdseries\slshape n}} is a not too big prime. }

 

\subsection{\textcolor{Chapter }{RandomPrimitivePolynomial}}
\logpage{[ 59, 5, 3 ]}\nobreak
\hyperdef{L}{X7ECC593583E68A6C}{}
{\noindent\textcolor{FuncColor}{$\triangleright$\enspace\texttt{RandomPrimitivePolynomial({\mdseries\slshape F, n[, i]})\index{RandomPrimitivePolynomial@\texttt{RandomPrimitivePolynomial}}
\label{RandomPrimitivePolynomial}
}\hfill{\scriptsize (function)}}\\


 For a finite field \mbox{\texttt{\mdseries\slshape F}} and a positive integer \mbox{\texttt{\mdseries\slshape n}} this function returns a primitive polynomial of degree \mbox{\texttt{\mdseries\slshape n}} over \mbox{\texttt{\mdseries\slshape F}}, that is a zero of this polynomial has maximal multiplicative order $|\mbox{\texttt{\mdseries\slshape F}}|^n-1$. If \mbox{\texttt{\mdseries\slshape i}} is given then the polynomial is written in variable number \mbox{\texttt{\mdseries\slshape i}} over \mbox{\texttt{\mdseries\slshape F}} (see{\nobreakspace}\texttt{Indeterminate} (\ref{Indeterminate:for a ring (and a number)})), the default for \mbox{\texttt{\mdseries\slshape i}} is 1. 

 Alternatively, \mbox{\texttt{\mdseries\slshape F}} can be a prime power q, then \mbox{\texttt{\mdseries\slshape F}} = GF(q) is assumed. And \mbox{\texttt{\mdseries\slshape i}} can be a univariate polynomial over \mbox{\texttt{\mdseries\slshape F}}, then the result is a polynomial in the same variable. 

 This function can work for much larger fields than those for which Conway
polynomials are available, of course \textsf{GAP} must be able to factorize $|\mbox{\texttt{\mdseries\slshape F}}|^n-1$. }

 }

  
\section{\textcolor{Chapter }{Printing, Viewing and Displaying Finite Field Elements}}\label{Printing, Viewing and Displaying Finite Field Elements}
\logpage{[ 59, 6, 0 ]}
\hyperdef{L}{X78EE3656879C3B88}{}
{
  

\subsection{\textcolor{Chapter }{ViewObj (for a ffe)}}
\logpage{[ 59, 6, 1 ]}\nobreak
\hyperdef{L}{X80DAAA5E7C79C94C}{}
{\noindent\textcolor{FuncColor}{$\triangleright$\enspace\texttt{ViewObj({\mdseries\slshape z})\index{ViewObj@\texttt{ViewObj}!for a ffe}
\label{ViewObj:for a ffe}
}\hfill{\scriptsize (method)}}\\
\noindent\textcolor{FuncColor}{$\triangleright$\enspace\texttt{PrintObj({\mdseries\slshape z})\index{PrintObj@\texttt{PrintObj}!for a ffe}
\label{PrintObj:for a ffe}
}\hfill{\scriptsize (method)}}\\
\noindent\textcolor{FuncColor}{$\triangleright$\enspace\texttt{Display({\mdseries\slshape z})\index{Display@\texttt{Display}!for a ffe}
\label{Display:for a ffe}
}\hfill{\scriptsize (method)}}\\


 Internal finite field elements are viewed, printed and displayed (see section \ref{View and Print} for the distinctions between these operations) as powers of the primitive root
(except for the zero element, which is displayed as 0 times the primitive
root). Thus: 

 
\begin{Verbatim}[commandchars=!@|,fontsize=\small,frame=single,label=Example]
  !gapprompt@gap>| !gapinput@Z(2);|
  Z(2)^0
  !gapprompt@gap>| !gapinput@Z(5)+Z(5);|
  Z(5)^2
  !gapprompt@gap>| !gapinput@Z(256);|
  Z(2^8)
  !gapprompt@gap>| !gapinput@Zero(Z(125));|
  0*Z(5)
\end{Verbatim}
 

 Note also that each element is displayed as an element of the field it
generates, and that the size of the field is printed as a power of the
characteristic. 

 Elements of larger fields are printed as \textsf{GAP} expressions which represent them as sums of low powers of the primitive root: 

 
\begin{Verbatim}[commandchars=!@|,fontsize=\small,frame=single,label=Example]
  !gapprompt@gap>| !gapinput@Print( Z(3,20)^100, "\n" );|
  2*Z(3,20)^2+Z(3,20)^4+Z(3,20)^6+Z(3,20)^7+2*Z(3,20)^9+2*Z(3,20)^10+2*Z\
  (3,20)^12+2*Z(3,20)^15+2*Z(3,20)^17+Z(3,20)^18+Z(3,20)^19
  !gapprompt@gap>| !gapinput@Print( Z(3,20)^((3^20-1)/(3^10-1)), "\n" );|
  Z(3,20)^3+2*Z(3,20)^4+2*Z(3,20)^7+Z(3,20)^8+2*Z(3,20)^10+Z(3,20)^11+2*\
  Z(3,20)^12+Z(3,20)^13+Z(3,20)^14+Z(3,20)^15+Z(3,20)^17+Z(3,20)^18+2*Z(\
  3,20)^19
  !gapprompt@gap>| !gapinput@Z(3,20)^((3^20-1)/(3^10-1)) = Z(3,10);|
  true
\end{Verbatim}
 

 Note from the second example above, that these elements are not always written
over the smallest possible field before being output. 

 The \texttt{ViewObj} and \texttt{Display} methods for these large finite field elements use a slightly more compact, but
mathematically equivalent representation. The primitive root is represented by \texttt{z}; its $i$\texttt{\symbol{45}}th power by \texttt{z}$i$ and $k$ times this power by $k$\texttt{z}$i$. 

 
\begin{Verbatim}[commandchars=!@|,fontsize=\small,frame=single,label=Example]
  !gapprompt@gap>| !gapinput@Z(5,20)^100;|
  z2+z4+4z5+2z6+z8+3z9+4z10+3z12+z13+2z14+4z16+3z17+2z18+2z19
\end{Verbatim}
 

 This output format is always used for \texttt{Display}. For \texttt{ViewObj} it is used only if its length would not exceed the number of lines specified
in the user preference \texttt{ViewLength} (see \texttt{SetUserPreference} (\ref{SetUserPreference}). Longer output is replaced by \texttt{{\textless}{\textless}an element of GF(\mbox{\texttt{\mdseries\slshape p}}, \mbox{\texttt{\mdseries\slshape d}}){\textgreater}{\textgreater}}. 

 
\begin{Verbatim}[commandchars=!@|,fontsize=\small,frame=single,label=Example]
  !gapprompt@gap>| !gapinput@Z(2,409)^100000;|
  <<an element of GF(2, 409)>>
  !gapprompt@gap>| !gapinput@Display(Z(2,409)^100000);|
  z2+z3+z4+z5+z6+z7+z8+z10+z11+z13+z17+z19+z20+z29+z32+z34+z35+z37+z40+z\
  45+z46+z48+z50+z52+z54+z55+z58+z59+z60+z66+z67+z68+z70+z74+z79+z80+z81\
  +z82+z83+z86+z91+z93+z94+z95+z96+z98+z99+z100+z101+z102+z104+z106+z109\
  +z110+z112+z114+z115+z118+z119+z123+z126+z127+z135+z138+z140+z142+z143\
  +z146+z147+z154+z159+z161+z162+z168+z170+z171+z173+z174+z181+z182+z183\
  +z186+z188+z189+z192+z193+z194+z195+z196+z199+z202+z204+z205+z207+z208\
  +z209+z211+z212+z213+z214+z215+z216+z218+z219+z220+z222+z223+z229+z232\
  +z235+z236+z237+z238+z240+z243+z244+z248+z250+z251+z256+z258+z262+z263\
  +z268+z270+z271+z272+z274+z276+z282+z286+z288+z289+z294+z295+z299+z300\
  +z301+z302+z303+z304+z305+z306+z307+z308+z309+z310+z312+z314+z315+z316\
  +z320+z321+z322+z324+z325+z326+z327+z330+z332+z335+z337+z338+z341+z344\
  +z348+z350+z352+z353+z356+z357+z358+z360+z362+z364+z366+z368+z372+z373\
  +z374+z375+z378+z379+z380+z381+z383+z384+z386+z387+z390+z395+z401+z402\
  +z406+z408
\end{Verbatim}
 

 Finally note that elements of large prime fields are stored and displayed as
residue class objects. So 

 
\begin{Verbatim}[commandchars=!@|,fontsize=\small,frame=single,label=Example]
  !gapprompt@gap>| !gapinput@Z(65537);|
  ZmodpZObj( 3, 65537 )
\end{Verbatim}
 }

 }

 }

            
\chapter{\textcolor{Chapter }{Abelian Number Fields}}\label{Abelian Number Fields}
\logpage{[ 60, 0, 0 ]}
\hyperdef{L}{X80510B5880521FDC}{}
{
  An \emph{abelian number field} is a field in characteristic zero that is a finite dimensional normal
extension of its prime field such that the Galois group is abelian. In \textsf{GAP}, one implementation of abelian number fields is given by fields of cyclotomic
numbers (see Chapter{\nobreakspace}\ref{Cyclotomic Numbers}). Note that abelian number fields can also be constructed with the more
general \texttt{AlgebraicExtension} (\ref{AlgebraicExtension}), a discussion of advantages and disadvantages can be found in{\nobreakspace}\ref{Internally Represented Cyclotomics}. The functions described in this chapter have been developed for fields whose
elements are in the filter \texttt{IsCyclotomic} (\ref{IsCyclotomic}), they may or may not work well for abelian number fields consisting of other
kinds of elements. 

 Throughout this chapter, ${\ensuremath{\mathbb Q}}_n$ will denote the cyclotomic field generated by the field ${\ensuremath{\mathbb Q}}$ of rationals together with $n$\texttt{\symbol{45}}th roots of unity. 

 In{\nobreakspace}\ref{Construction of Abelian Number Fields}, constructors for abelian number fields are described, \ref{Operations for Abelian Number Fields} introduces operations for abelian number fields, \ref{Integral Bases of Abelian Number Fields} deals with the vector space structure of abelian number fields, and \ref{Galois Groups of Abelian Number Fields} describes field automorphisms of abelian number fields,   
\section{\textcolor{Chapter }{Construction of Abelian Number Fields}}\label{Construction of Abelian Number Fields}
\logpage{[ 60, 1, 0 ]}
\hyperdef{L}{X7D4E43E5799753B5}{}
{
  Besides the usual construction using \texttt{Field} (\ref{Field:for several generators}) or \texttt{DefaultField} (\ref{DefaultField:for cyclotomics}) (see{\nobreakspace}\texttt{DefaultField} (\ref{DefaultField:for cyclotomics})), abelian number fields consisting of cyclotomics can be created with \texttt{CyclotomicField} (\ref{CyclotomicField:for (subfield and) conductor}) and \texttt{AbelianNumberField} (\ref{AbelianNumberField}). 

\subsection{\textcolor{Chapter }{CyclotomicField (for (subfield and) conductor)}}
\logpage{[ 60, 1, 1 ]}\nobreak
\hyperdef{L}{X80D21D80850EFA4B}{}
{\noindent\textcolor{FuncColor}{$\triangleright$\enspace\texttt{CyclotomicField({\mdseries\slshape [subfield, ]n})\index{CyclotomicField@\texttt{CyclotomicField}!for (subfield and) conductor}
\label{CyclotomicField:for (subfield and) conductor}
}\hfill{\scriptsize (function)}}\\
\noindent\textcolor{FuncColor}{$\triangleright$\enspace\texttt{CyclotomicField({\mdseries\slshape [subfield, ]gens})\index{CyclotomicField@\texttt{CyclotomicField}!for (subfield and) generators}
\label{CyclotomicField:for (subfield and) generators}
}\hfill{\scriptsize (function)}}\\
\noindent\textcolor{FuncColor}{$\triangleright$\enspace\texttt{CF({\mdseries\slshape [subfield, ]n})\index{CF@\texttt{CF}!for (subfield and) conductor}
\label{CF:for (subfield and) conductor}
}\hfill{\scriptsize (function)}}\\
\noindent\textcolor{FuncColor}{$\triangleright$\enspace\texttt{CF({\mdseries\slshape [subfield, ]gens})\index{CF@\texttt{CF}!for (subfield and) generators}
\label{CF:for (subfield and) generators}
}\hfill{\scriptsize (function)}}\\


 The first version creates the \mbox{\texttt{\mdseries\slshape n}}\texttt{\symbol{45}}th cyclotomic field ${\ensuremath{\mathbb Q}}_n$. The second version creates the smallest cyclotomic field containing the
elements in the list \mbox{\texttt{\mdseries\slshape gens}}. In both cases the field can be generated as an extension of a designated
subfield \mbox{\texttt{\mdseries\slshape subfield}} (cf.{\nobreakspace}\ref{Integral Bases of Abelian Number Fields}). 

 \texttt{CyclotomicField} can be abbreviated to \texttt{CF}, this form is used also when \textsf{GAP} prints cyclotomic fields. 

 Fields constructed with the one argument version of \texttt{CF} are stored in the global list \texttt{CYCLOTOMIC{\textunderscore}FIELDS}, so repeated calls of \texttt{CF} just fetch these field objects after they have been created once.  

 
\begin{Verbatim}[commandchars=!@|,fontsize=\small,frame=single,label=Example]
  !gapprompt@gap>| !gapinput@CyclotomicField( 5 );  CyclotomicField( [ Sqrt(3) ] );|
  CF(5)
  CF(12)
  !gapprompt@gap>| !gapinput@CF( CF(3), 12 );  CF( CF(4), [ Sqrt(7) ] );|
  AsField( CF(3), CF(12) )
  AsField( GaussianRationals, CF(28) )
\end{Verbatim}
 }

 

\subsection{\textcolor{Chapter }{AbelianNumberField}}
\logpage{[ 60, 1, 2 ]}\nobreak
\hyperdef{L}{X80E5AD028143E11E}{}
{\noindent\textcolor{FuncColor}{$\triangleright$\enspace\texttt{AbelianNumberField({\mdseries\slshape n, stab})\index{AbelianNumberField@\texttt{AbelianNumberField}}
\label{AbelianNumberField}
}\hfill{\scriptsize (function)}}\\
\noindent\textcolor{FuncColor}{$\triangleright$\enspace\texttt{NF({\mdseries\slshape n, stab})\index{NF@\texttt{NF}}
\label{NF}
}\hfill{\scriptsize (function)}}\\


 For a positive integer \mbox{\texttt{\mdseries\slshape n}} and a list \mbox{\texttt{\mdseries\slshape stab}} of prime residues modulo \mbox{\texttt{\mdseries\slshape n}}, \texttt{AbelianNumberField} returns the fixed field of the group described by \mbox{\texttt{\mdseries\slshape stab}} (cf.{\nobreakspace}\texttt{GaloisStabilizer} (\ref{GaloisStabilizer})), in the \mbox{\texttt{\mdseries\slshape n}}\texttt{\symbol{45}}th cyclotomic field. \texttt{AbelianNumberField} is mainly thought for internal use and for printing fields in a standard way; \texttt{Field} (\ref{Field:for several generators}) (cf.{\nobreakspace}also{\nobreakspace}\ref{Operations for Abelian Number Fields}) is probably more suitable if one knows generators of the field in question. 

 \texttt{AbelianNumberField} can be abbreviated to \texttt{NF}, this form is used also when \textsf{GAP} prints abelian number fields. 

 Fields constructed with \texttt{NF} are stored in the global list \texttt{ABELIAN{\textunderscore}NUMBER{\textunderscore}FIELDS}, so repeated calls of \texttt{NF} just fetch these field objects after they have been created once.  

 
\begin{Verbatim}[commandchars=!@|,fontsize=\small,frame=single,label=Example]
  !gapprompt@gap>| !gapinput@NF( 7, [ 1 ] );|
  CF(7)
  !gapprompt@gap>| !gapinput@f:= NF( 7, [ 1, 2 ] );  Sqrt(-7); Sqrt(-7) in f;|
  NF(7,[ 1, 2, 4 ])
  E(7)+E(7)^2-E(7)^3+E(7)^4-E(7)^5-E(7)^6
  true
\end{Verbatim}
 }

 

\subsection{\textcolor{Chapter }{GaussianRationals}}
\logpage{[ 60, 1, 3 ]}\nobreak
\hyperdef{L}{X82F53C65802FF551}{}
{\noindent\textcolor{FuncColor}{$\triangleright$\enspace\texttt{GaussianRationals\index{GaussianRationals@\texttt{GaussianRationals}}
\label{GaussianRationals}
}\hfill{\scriptsize (global variable)}}\\
\noindent\textcolor{FuncColor}{$\triangleright$\enspace\texttt{IsGaussianRationals({\mdseries\slshape obj})\index{IsGaussianRationals@\texttt{IsGaussianRationals}}
\label{IsGaussianRationals}
}\hfill{\scriptsize (Category)}}\\


 \texttt{GaussianRationals} is the field ${\ensuremath{\mathbb Q}}_4 = {\ensuremath{\mathbb Q}}(\sqrt{{-1}})$ of Gaussian rationals, as a set of cyclotomic numbers, see
Chapter{\nobreakspace}\ref{Cyclotomic Numbers} for basic operations. This field can also be obtained as \texttt{CF(4)} (see \texttt{CyclotomicField} (\ref{CyclotomicField:for (subfield and) conductor})). 

 The filter \texttt{IsGaussianRationals} returns \texttt{true} for the \textsf{GAP} object \texttt{GaussianRationals}, and \texttt{false} for all other \textsf{GAP} objects. 

 (For details about the field of rationals, see Chapter{\nobreakspace}\texttt{Rationals} (\ref{Rationals}).) 

 
\begin{Verbatim}[commandchars=!@|,fontsize=\small,frame=single,label=Example]
  !gapprompt@gap>| !gapinput@CF(4) = GaussianRationals;|
  true
  !gapprompt@gap>| !gapinput@Sqrt(-1) in GaussianRationals;|
  true
\end{Verbatim}
 }

 }

  
\section{\textcolor{Chapter }{Operations for Abelian Number Fields}}\label{Operations for Abelian Number Fields}
\logpage{[ 60, 2, 0 ]}
\hyperdef{L}{X81B5FE06781DB824}{}
{
  For operations for elements of abelian number fields, e.g., \texttt{Conductor} (\ref{Conductor:for a cyclotomic}) or \texttt{ComplexConjugate} (\ref{ComplexConjugate}), see Chapter{\nobreakspace}\ref{Cyclotomic Numbers}. 

\subsection{\textcolor{Chapter }{Factors (for polynomials over abelian number fields)}}
\logpage{[ 60, 2, 1 ]}\nobreak
\hyperdef{L}{X7B0AB0FB7A4136C4}{}
{\noindent\textcolor{FuncColor}{$\triangleright$\enspace\texttt{Factors({\mdseries\slshape F})\index{Factors@\texttt{Factors}!for polynomials over abelian number fields}
\label{Factors:for polynomials over abelian number fields}
}\hfill{\scriptsize (method)}}\\


 Factoring of polynomials over abelian number fields consisting of cyclotomics
works in principle but is not very efficient if the degree of the field
extension is large. 

 
\begin{Verbatim}[commandchars=!@|,fontsize=\small,frame=single,label=Example]
  !gapprompt@gap>| !gapinput@x:= Indeterminate( CF(5) );|
  x_1
  !gapprompt@gap>| !gapinput@Factors( PolynomialRing( Rationals ), x^5-1 );|
  [ x_1-1, x_1^4+x_1^3+x_1^2+x_1+1 ]
  !gapprompt@gap>| !gapinput@Factors( PolynomialRing( CF(5) ), x^5-1 );|
  [ x_1-1, x_1+(-E(5)), x_1+(-E(5)^2), x_1+(-E(5)^3), x_1+(-E(5)^4) ]
\end{Verbatim}
 }

 

\subsection{\textcolor{Chapter }{IsNumberField}}
\logpage{[ 60, 2, 2 ]}\nobreak
\hyperdef{L}{X87D78F5E875F2E8A}{}
{\noindent\textcolor{FuncColor}{$\triangleright$\enspace\texttt{IsNumberField({\mdseries\slshape F})\index{IsNumberField@\texttt{IsNumberField}}
\label{IsNumberField}
}\hfill{\scriptsize (property)}}\\


 \index{number field} returns \texttt{true} if the field \mbox{\texttt{\mdseries\slshape F}} is a finite dimensional extension of a prime field in characteristic zero, and \texttt{false} otherwise. }

 

\subsection{\textcolor{Chapter }{IsAbelianNumberField}}
\logpage{[ 60, 2, 3 ]}\nobreak
\hyperdef{L}{X7D202D707D5708FA}{}
{\noindent\textcolor{FuncColor}{$\triangleright$\enspace\texttt{IsAbelianNumberField({\mdseries\slshape F})\index{IsAbelianNumberField@\texttt{IsAbelianNumberField}}
\label{IsAbelianNumberField}
}\hfill{\scriptsize (property)}}\\


 \index{abelian number field} returns \texttt{true} if the field \mbox{\texttt{\mdseries\slshape F}} is a number field (see{\nobreakspace}\texttt{IsNumberField} (\ref{IsNumberField})) that is a Galois extension of the prime field, with abelian Galois group
(see{\nobreakspace}\texttt{GaloisGroup} (\ref{GaloisGroup:of field})). }

 

\subsection{\textcolor{Chapter }{IsCyclotomicField}}
\logpage{[ 60, 2, 4 ]}\nobreak
\hyperdef{L}{X84CAE4627F0CD639}{}
{\noindent\textcolor{FuncColor}{$\triangleright$\enspace\texttt{IsCyclotomicField({\mdseries\slshape F})\index{IsCyclotomicField@\texttt{IsCyclotomicField}}
\label{IsCyclotomicField}
}\hfill{\scriptsize (property)}}\\


 returns \texttt{true} if the field \mbox{\texttt{\mdseries\slshape F}} is a \emph{cyclotomic field}, i.e., an abelian number field (see{\nobreakspace}\texttt{IsAbelianNumberField} (\ref{IsAbelianNumberField})) that can be generated by roots of unity. 

 
\begin{Verbatim}[commandchars=!@|,fontsize=\small,frame=single,label=Example]
  !gapprompt@gap>| !gapinput@IsNumberField( CF(9) ); IsAbelianNumberField( Field( [ ER(3) ] ) );|
  true
  true
  !gapprompt@gap>| !gapinput@IsNumberField( GF(2) );|
  false
  !gapprompt@gap>| !gapinput@IsCyclotomicField( CF(9) );|
  true
  !gapprompt@gap>| !gapinput@IsCyclotomicField( Field( [ Sqrt(-3) ] ) );|
  true
  !gapprompt@gap>| !gapinput@IsCyclotomicField( Field( [ Sqrt(3) ] ) );|
  false
\end{Verbatim}
 }

 

\subsection{\textcolor{Chapter }{GaloisStabilizer}}
\logpage{[ 60, 2, 5 ]}\nobreak
\hyperdef{L}{X87E7313D8070B9CC}{}
{\noindent\textcolor{FuncColor}{$\triangleright$\enspace\texttt{GaloisStabilizer({\mdseries\slshape F})\index{GaloisStabilizer@\texttt{GaloisStabilizer}}
\label{GaloisStabilizer}
}\hfill{\scriptsize (attribute)}}\\


 Let \mbox{\texttt{\mdseries\slshape F}} be an abelian number field (see{\nobreakspace}\texttt{IsAbelianNumberField} (\ref{IsAbelianNumberField})) with conductor $n$, say. (This means that the $n$\texttt{\symbol{45}}th cyclotomic field is the smallest cyclotomic field
containing \mbox{\texttt{\mdseries\slshape F}}, see{\nobreakspace}\texttt{Conductor} (\ref{Conductor:for a cyclotomic}).) \texttt{GaloisStabilizer} returns the set of all those integers $k$ in the range $[ 1 .. n ]$ such that the field automorphism induced by raising $n$\texttt{\symbol{45}}th roots of unity to the $k$\texttt{\symbol{45}}th power acts trivially on \mbox{\texttt{\mdseries\slshape F}}. 

 
\begin{Verbatim}[commandchars=!@|,fontsize=\small,frame=single,label=Example]
  !gapprompt@gap>| !gapinput@r5:= Sqrt(5);|
  E(5)-E(5)^2-E(5)^3+E(5)^4
  !gapprompt@gap>| !gapinput@GaloisCyc( r5, 4 ) = r5;  GaloisCyc( r5, 2 ) = r5;|
  true
  false
  !gapprompt@gap>| !gapinput@GaloisStabilizer( Field( [ r5 ] ) );|
  [ 1, 4 ]
\end{Verbatim}
 }

 }

  
\section{\textcolor{Chapter }{Integral Bases of Abelian Number Fields}}\label{Integral Bases of Abelian Number Fields}
\logpage{[ 60, 3, 0 ]}
\hyperdef{L}{X7D2421AC8491D2BE}{}
{
  Each abelian number field is naturally a vector space over ${\ensuremath{\mathbb Q}}$. Moreover, if the abelian number field $F$ contains the $n$\texttt{\symbol{45}}th cyclotomic field ${\ensuremath{\mathbb Q}}_n$ then $F$ is a vector space over ${\ensuremath{\mathbb Q}}_n$. In \textsf{GAP}, each field object represents a vector space object over a certain subfield $S$, which depends on the way $F$ was constructed. The subfield $S$ can be accessed as the value of the attribute \texttt{LeftActingDomain} (\ref{LeftActingDomain}). 

 The return values of \texttt{NF} (\ref{NF}) and of the one argument versions of \texttt{CF} (\ref{CF:for (subfield and) conductor}) represent vector spaces over ${\ensuremath{\mathbb Q}}$, and the return values of the two argument version of \texttt{CF} (\ref{CF:for (subfield and) conductor}) represent vector spaces over the field that is given as the first argument.
For an abelian number field \mbox{\texttt{\mdseries\slshape F}} and a subfield \mbox{\texttt{\mdseries\slshape S}} of \mbox{\texttt{\mdseries\slshape F}}, a \textsf{GAP} object representing \mbox{\texttt{\mdseries\slshape F}} as a vector space over \mbox{\texttt{\mdseries\slshape S}} can be constructed using \texttt{AsField} (\ref{AsField}). 

 \index{cyclotomic fields!CanonicalBasis} Let \mbox{\texttt{\mdseries\slshape F}} be the cyclotomic field ${\ensuremath{\mathbb Q}}_n$, represented as a vector space over the subfield \mbox{\texttt{\mdseries\slshape S}}. If \mbox{\texttt{\mdseries\slshape S}} is the cyclotomic field ${\ensuremath{\mathbb Q}}_m$, with $m$ a divisor of $n$, then \texttt{CanonicalBasis( \mbox{\texttt{\mdseries\slshape F}} )} returns the Zumbroich basis of \mbox{\texttt{\mdseries\slshape F}} relative to \mbox{\texttt{\mdseries\slshape S}}, which consists of the roots of unity \texttt{E(\mbox{\texttt{\mdseries\slshape n}})}\texttt{\symbol{94}}\mbox{\texttt{\mdseries\slshape i}} where \mbox{\texttt{\mdseries\slshape i}} is an element of the list \texttt{ZumbroichBase( \mbox{\texttt{\mdseries\slshape n}}, \mbox{\texttt{\mdseries\slshape m}} )} (see{\nobreakspace}\texttt{ZumbroichBase} (\ref{ZumbroichBase})). If \mbox{\texttt{\mdseries\slshape S}} is an abelian number field that is not a cyclotomic field then \texttt{CanonicalBasis( \mbox{\texttt{\mdseries\slshape F}} )} returns a normal \mbox{\texttt{\mdseries\slshape S}}\texttt{\symbol{45}}basis of \mbox{\texttt{\mdseries\slshape F}}, i.e., a basis that is closed under the field automorphisms of \mbox{\texttt{\mdseries\slshape F}}. 

 \index{abelian number fields!CanonicalBasis} Let \mbox{\texttt{\mdseries\slshape F}} be the abelian number field \texttt{NF( \mbox{\texttt{\mdseries\slshape n}}, \mbox{\texttt{\mdseries\slshape stab}} )}, with conductor \mbox{\texttt{\mdseries\slshape n}}, that is itself not a cyclotomic field, represented as a vector space over
the subfield \mbox{\texttt{\mdseries\slshape S}}. If \mbox{\texttt{\mdseries\slshape S}} is the cyclotomic field ${\ensuremath{\mathbb Q}}_m$, with $m$ a divisor of $n$, then \texttt{CanonicalBasis( \mbox{\texttt{\mdseries\slshape F}} )} returns the Lenstra basis of \mbox{\texttt{\mdseries\slshape F}} relative to \mbox{\texttt{\mdseries\slshape S}} that consists of the sums of roots of unity described by \texttt{LenstraBase( \mbox{\texttt{\mdseries\slshape n}}, \mbox{\texttt{\mdseries\slshape stab}}, \mbox{\texttt{\mdseries\slshape stab}}, \mbox{\texttt{\mdseries\slshape m}} )} (see{\nobreakspace}\texttt{LenstraBase} (\ref{LenstraBase})). If \mbox{\texttt{\mdseries\slshape S}} is an abelian number field that is not a cyclotomic field then \texttt{CanonicalBasis( \mbox{\texttt{\mdseries\slshape F}} )} returns a normal \mbox{\texttt{\mdseries\slshape S}}\texttt{\symbol{45}}basis of \mbox{\texttt{\mdseries\slshape F}}. 

 
\begin{Verbatim}[commandchars=!@|,fontsize=\small,frame=single,label=Example]
  !gapprompt@gap>| !gapinput@f:= CF(8);;   # a cycl. field over the rationals|
  !gapprompt@gap>| !gapinput@b:= CanonicalBasis( f );;  BasisVectors( b );|
  [ 1, E(8), E(4), E(8)^3 ]
  !gapprompt@gap>| !gapinput@Coefficients( b, Sqrt(-2) );|
  [ 0, 1, 0, 1 ]
  !gapprompt@gap>| !gapinput@f:= AsField( CF(4), CF(8) );;  # a cycl. field over a cycl. field|
  !gapprompt@gap>| !gapinput@b:= CanonicalBasis( f );;  BasisVectors( b );|
  [ 1, E(8) ]
  !gapprompt@gap>| !gapinput@Coefficients( b, Sqrt(-2) );|
  [ 0, 1+E(4) ]
  !gapprompt@gap>| !gapinput@f:= AsField( Field( [ Sqrt(-2) ] ), CF(8) );;|
  !gapprompt@gap>| !gapinput@# a cycl. field over a non-cycl. field|
  !gapprompt@gap>| !gapinput@b:= CanonicalBasis( f );;  BasisVectors( b );|
  [ 1/2+1/2*E(8)-1/2*E(8)^2-1/2*E(8)^3,
    1/2-1/2*E(8)+1/2*E(8)^2+1/2*E(8)^3 ]
  !gapprompt@gap>| !gapinput@Coefficients( b, Sqrt(-2) );|
  [ E(8)+E(8)^3, E(8)+E(8)^3 ]
  !gapprompt@gap>| !gapinput@f:= Field( [ Sqrt(-2) ] );   # a non-cycl. field over the rationals|
  NF(8,[ 1, 3 ])
  !gapprompt@gap>| !gapinput@b:= CanonicalBasis( f );;  BasisVectors( b );|
  [ 1, E(8)+E(8)^3 ]
  !gapprompt@gap>| !gapinput@Coefficients( b, Sqrt(-2) );|
  [ 0, 1 ]
\end{Verbatim}
 

\subsection{\textcolor{Chapter }{ZumbroichBase}}
\logpage{[ 60, 3, 1 ]}\nobreak
\hyperdef{L}{X7F52BEA0862E06F2}{}
{\noindent\textcolor{FuncColor}{$\triangleright$\enspace\texttt{ZumbroichBase({\mdseries\slshape n, m})\index{ZumbroichBase@\texttt{ZumbroichBase}}
\label{ZumbroichBase}
}\hfill{\scriptsize (function)}}\\


 Let \mbox{\texttt{\mdseries\slshape n}} and \mbox{\texttt{\mdseries\slshape m}} be positive integers, such that \mbox{\texttt{\mdseries\slshape m}} divides \mbox{\texttt{\mdseries\slshape n}}. \texttt{ZumbroichBase} returns the set of exponents $i$ for which \texttt{E(\mbox{\texttt{\mdseries\slshape n}})\texttt{\symbol{94}}}$i$ belongs to the (generalized) Zumbroich basis of the cyclotomic field ${\ensuremath{\mathbb Q}}_n$, viewed as a vector space over ${\ensuremath{\mathbb Q}}_m$. 

 This basis is defined as follows. Let $P$ denote the set of prime divisors of \mbox{\texttt{\mdseries\slshape n}}, $\mbox{\texttt{\mdseries\slshape n}} = \prod_{{p \in P}} p^{{\nu_p}}$, and $\mbox{\texttt{\mdseries\slshape m}} = \prod_{{p \in P}} p^{{\mu_p}}$ with $\mu_p \leq \nu_p$. Let $e_l =$ \texttt{E}$(l)$ for any positive integer $l$, and $\{ e_{{n_1}}^j \}_{{j \in J}} \otimes \{ e_{{n_2}}^k \}_{{k \in K}} = \{
e_{{n_1}}^j \cdot e_{{n_2}}^k \}_{{j \in J, k \in K}}$. 

 Then the basis is 
\[ B_{{n,m}} = \bigotimes_{{p \in P}} \bigotimes_{{k = \mu_p}}^{{\nu_p-1}} \{
e_{{p^{{k+1}}}}^j \}_{{j \in J_{{k,p}}}} \]
 where $J_{{k,p}} =$ \begin{center}
\begin{tabular}{lcl}$\{ 0 \}$&
;&
$k = 0, p = 2$\\
$\{ 0, 1 \}$&
;&
$k > 0, p = 2$\\
$\{ 1, \ldots, p-1 \}$&
;&
$k = 0, p \neq 2$\\
$\{ -(p-1)/2, \ldots, (p-1)/2 \}$&
;&
$k > 0, p \neq 2$\\
\end{tabular}\\[2mm]
\end{center}

 

 $B_{{n,1}}$ is equal to the basis of ${\ensuremath{\mathbb Q}}_n$ over the rationals which is introduced in{\nobreakspace}\cite{Zum89}. Also the conversion of arbitrary sums of roots of unity into its basis
representation, and the reduction to the minimal cyclotomic field are
described in this thesis. (Note that the notation here is slightly different
from that there.) 

 $B_{{n,m}}$ consists of roots of unity, it is an integral basis (that is, exactly the
integral elements in ${\ensuremath{\mathbb Q}}_n$ have integral coefficients w.r.t.{\nobreakspace}$B_{{n,m}}$, cf.{\nobreakspace}\texttt{IsIntegralCyclotomic} (\ref{IsIntegralCyclotomic})), it is a normal basis for squarefree $n$ and closed under complex conjugation for odd $n$. 

 \emph{Note:} For $\mbox{\texttt{\mdseries\slshape n}} \equiv 2 \pmod 4$, we have \texttt{ZumbroichBase(\mbox{\texttt{\mdseries\slshape n}}, 1) = 2 * ZumbroichBase(\mbox{\texttt{\mdseries\slshape n}}/2, 1)} and \texttt{List( ZumbroichBase(\mbox{\texttt{\mdseries\slshape n}}, 1), x \texttt{\symbol{45}}{\textgreater} E(\mbox{\texttt{\mdseries\slshape n}})\texttt{\symbol{94}}x ) = List( ZumbroichBase(\mbox{\texttt{\mdseries\slshape n}}/2, 1), x \texttt{\symbol{45}}{\textgreater} E(\mbox{\texttt{\mdseries\slshape n}}/2)\texttt{\symbol{94}}x )}. 

 
\begin{Verbatim}[commandchars=!@|,fontsize=\small,frame=single,label=Example]
  !gapprompt@gap>| !gapinput@ZumbroichBase( 15, 1 ); ZumbroichBase( 12, 3 );|
  [ 1, 2, 4, 7, 8, 11, 13, 14 ]
  [ 0, 3 ]
  !gapprompt@gap>| !gapinput@ZumbroichBase( 10, 2 ); ZumbroichBase( 32, 4 );|
  [ 2, 4, 6, 8 ]
  [ 0, 1, 2, 3, 4, 5, 6, 7 ]
\end{Verbatim}
 }

 

\subsection{\textcolor{Chapter }{LenstraBase}}
\logpage{[ 60, 3, 2 ]}\nobreak
\hyperdef{L}{X87DB9C2C858B722A}{}
{\noindent\textcolor{FuncColor}{$\triangleright$\enspace\texttt{LenstraBase({\mdseries\slshape n, stabilizer, super, m})\index{LenstraBase@\texttt{LenstraBase}}
\label{LenstraBase}
}\hfill{\scriptsize (function)}}\\


 Let \mbox{\texttt{\mdseries\slshape n}} and \mbox{\texttt{\mdseries\slshape m}} be positive integers such that \mbox{\texttt{\mdseries\slshape m}} divides \mbox{\texttt{\mdseries\slshape n}}, \mbox{\texttt{\mdseries\slshape stabilizer}} be a list of prime residues modulo \mbox{\texttt{\mdseries\slshape n}}, which describes a subfield of the \mbox{\texttt{\mdseries\slshape n}}\texttt{\symbol{45}}th cyclotomic field (see{\nobreakspace}\texttt{GaloisStabilizer} (\ref{GaloisStabilizer})), and \mbox{\texttt{\mdseries\slshape super}} be a list representing a supergroup of the group given by \mbox{\texttt{\mdseries\slshape stabilizer}}. 

 \texttt{LenstraBase} returns a list $[ b_1, b_2, \ldots, b_k ]$ of lists, each $b_i$ consisting of integers such that the elements $\sum_{{j \in b_i}} $\texttt{E(n)}$^j$ form a basis of the abelian number field \texttt{NF( \mbox{\texttt{\mdseries\slshape n}}, \mbox{\texttt{\mdseries\slshape stabilizer}} )}, as a vector space over the \mbox{\texttt{\mdseries\slshape m}}\texttt{\symbol{45}}th cyclotomic field (see{\nobreakspace}\texttt{AbelianNumberField} (\ref{AbelianNumberField})). 

 This basis is an integral basis, that is, exactly the integral elements in \texttt{NF( \mbox{\texttt{\mdseries\slshape n}}, \mbox{\texttt{\mdseries\slshape stabilizer}} )} have integral coefficients. (For details about this basis, see{\nobreakspace}\cite{Bre97}.) 

 If possible then the result is chosen such that the group described by \mbox{\texttt{\mdseries\slshape super}} acts on it, consistently with the action of \mbox{\texttt{\mdseries\slshape stabilizer}}, i.e., each orbit of \mbox{\texttt{\mdseries\slshape super}} is a union of orbits of \mbox{\texttt{\mdseries\slshape stabilizer}}. (A usual case is \mbox{\texttt{\mdseries\slshape super}}\texttt{ = }\mbox{\texttt{\mdseries\slshape stabilizer}}, so there is no additional condition. 

 \emph{Note:} The $b_i$ are in general not sets, since for \texttt{\mbox{\texttt{\mdseries\slshape stabilizer}} = \mbox{\texttt{\mdseries\slshape super}}}, the first entry is always an element of \texttt{ZumbroichBase( \mbox{\texttt{\mdseries\slshape n}}, \mbox{\texttt{\mdseries\slshape m}} )}; this property is used by \texttt{NF} (\ref{NF}) and \texttt{Coefficients} (\ref{Coefficients}) (see{\nobreakspace}\ref{Integral Bases of Abelian Number Fields}). 

 \mbox{\texttt{\mdseries\slshape stabilizer}} must not contain the stabilizer of a proper cyclotomic subfield of the \mbox{\texttt{\mdseries\slshape n}}\texttt{\symbol{45}}th cyclotomic field, i.e., the result must describe a
basis for a field with conductor \mbox{\texttt{\mdseries\slshape n}}. 

 
\begin{Verbatim}[commandchars=!@|,fontsize=\small,frame=single,label=Example]
  !gapprompt@gap>| !gapinput@LenstraBase( 24, [ 1, 19 ], [ 1, 19 ], 1 );|
  [ [ 1, 19 ], [ 8 ], [ 11, 17 ], [ 16 ] ]
  !gapprompt@gap>| !gapinput@LenstraBase( 24, [ 1, 19 ], [ 1, 5, 19, 23 ], 1 );|
  [ [ 1, 19 ], [ 5, 23 ], [ 8 ], [ 16 ] ]
  !gapprompt@gap>| !gapinput@LenstraBase( 15, [ 1, 4 ], PrimeResidues( 15 ), 1 );|
  [ [ 1, 4 ], [ 2, 8 ], [ 7, 13 ], [ 11, 14 ] ]
\end{Verbatim}
 

 The first two results describe two bases of the field ${\ensuremath{\mathbb Q}}_3(\sqrt{{6}})$, the third result describes a normal basis of ${\ensuremath{\mathbb Q}}_3(\sqrt{{5}})$. }

   }

  
\section{\textcolor{Chapter }{Galois Groups of Abelian Number Fields}}\label{Galois Groups of Abelian Number Fields}
\logpage{[ 60, 4, 0 ]}
\hyperdef{L}{X7E4AB4B17C7BA10C}{}
{
  \index{abelian number fields!Galois group} \index{number fields!Galois group} \index{automorphism group!of number fields} The field automorphisms of the cyclotomic field ${\ensuremath{\mathbb Q}}_n$ (see Chapter{\nobreakspace}\ref{Cyclotomic Numbers}) are given by the linear maps $*k$ on ${\ensuremath{\mathbb Q}}_n$ that are defined by \texttt{E}$(n)^{{*k}} = $\texttt{E}$(n)^k$, where $1 \leq k < n$ and \texttt{Gcd}$( n, k ) = 1$ hold (see{\nobreakspace}\texttt{GaloisCyc} (\ref{GaloisCyc:for a cyclotomic})). Note that this action is \emph{not} equal to exponentiation of cyclotomics, i.e., for general cyclotomics $z$, $z^{{*k}}$ is different from $z^k$. 

 (In \textsf{GAP}, the image of a cyclotomic $z$ under $*k$ can be computed as \texttt{GaloisCyc( }$z, k$\texttt{ )}.) 

 
\begin{Verbatim}[commandchars=!@|,fontsize=\small,frame=single,label=Example]
  !gapprompt@gap>| !gapinput@( E(5) + E(5)^4 )^2; GaloisCyc( E(5) + E(5)^4, 2 );|
  -2*E(5)-E(5)^2-E(5)^3-2*E(5)^4
  E(5)^2+E(5)^3
\end{Verbatim}
 

 For \texttt{Gcd}$( n, k ) \neq 1$, the map \texttt{E}$(n) \mapsto$ \texttt{E}$(n)^k$ does \emph{not} define a field automorphism of ${\ensuremath{\mathbb Q}}_n$ but only a ${\ensuremath{\mathbb Q}}$\texttt{\symbol{45}}linear map. 

 
\begin{Verbatim}[commandchars=!@|,fontsize=\small,frame=single,label=Example]
  !gapprompt@gap>| !gapinput@GaloisCyc( E(5)+E(5)^4, 5 ); GaloisCyc( ( E(5)+E(5)^4 )^2, 5 );|
  2
  -6
\end{Verbatim}
 

\subsection{\textcolor{Chapter }{GaloisGroup (for abelian number fields)}}
\logpage{[ 60, 4, 1 ]}\nobreak
\hyperdef{L}{X7B55A90582E818F3}{}
{\noindent\textcolor{FuncColor}{$\triangleright$\enspace\texttt{GaloisGroup({\mdseries\slshape F})\index{GaloisGroup@\texttt{GaloisGroup}!for abelian number fields}
\label{GaloisGroup:for abelian number fields}
}\hfill{\scriptsize (method)}}\\


 The Galois group $Gal( {\ensuremath{\mathbb Q}}_n, {\ensuremath{\mathbb Q}} )$ of the field extension ${\ensuremath{\mathbb Q}}_n / {\ensuremath{\mathbb Q}}$ is isomorphic to the group $({\ensuremath{\mathbb Z}} / n {\ensuremath{\mathbb Z}})^{*}$ of prime residues modulo $n$, via the isomorphism $({\ensuremath{\mathbb Z}} / n {\ensuremath{\mathbb Z}})^{*} \rightarrow Gal(
{\ensuremath{\mathbb Q}}_n, {\ensuremath{\mathbb Q}} )$ that is defined by $k + n {\ensuremath{\mathbb Z}} \mapsto ( z \mapsto z^{*k} )$. 

 The Galois group of the field extension ${\ensuremath{\mathbb Q}}_n / L$ with any abelian number field $L \subseteq {\ensuremath{\mathbb Q}}_n$ is simply the factor group of $Gal( {\ensuremath{\mathbb Q}}_n, {\ensuremath{\mathbb Q}} )$ modulo the stabilizer of $L$, and the Galois group of $L / L'$, with $L'$ an abelian number field contained in $L$, is the subgroup in this group that stabilizes $L'$. These groups are easily described in terms of $({\ensuremath{\mathbb Z}} / n {\ensuremath{\mathbb Z}})^{*}$. Generators of $({\ensuremath{\mathbb Z}} / n {\ensuremath{\mathbb Z}})^{*}$ can be computed using \texttt{GeneratorsPrimeResidues} (\ref{GeneratorsPrimeResidues}). 

 In \textsf{GAP}, a field extension $L / L'$ is given by the field object $L$ with \texttt{LeftActingDomain} (\ref{LeftActingDomain}) value $L'$ (see{\nobreakspace}\ref{Integral Bases of Abelian Number Fields}). 

 
\begin{Verbatim}[commandchars=!@|,fontsize=\small,frame=single,label=Example]
  !gapprompt@gap>| !gapinput@f:= CF(15);|
  CF(15)
  !gapprompt@gap>| !gapinput@g:= GaloisGroup( f );|
  <group with 2 generators>
  !gapprompt@gap>| !gapinput@Size( g ); IsCyclic( g ); IsAbelian( g );|
  8
  false
  true
  !gapprompt@gap>| !gapinput@Action( g, NormalBase( f ), OnPoints );|
  Group([ (1,6)(2,4)(3,8)(5,7), (1,4,3,7)(2,8,5,6) ])
\end{Verbatim}
 

 The following example shows Galois groups of a cyclotomic field and of a
proper subfield that is not a cyclotomic field. 

 
\begin{Verbatim}[commandchars=!@|,fontsize=\small,frame=single,label=Example]
  !gapprompt@gap>| !gapinput@gens1:= GeneratorsOfGroup( GaloisGroup( CF(5) ) );|
  [ ANFAutomorphism( CF(5), 2 ) ]
  !gapprompt@gap>| !gapinput@gens2:= GeneratorsOfGroup( GaloisGroup( Field( Sqrt(5) ) ) );|
  [ ANFAutomorphism( NF(5,[ 1, 4 ]), 2 ) ]
  !gapprompt@gap>| !gapinput@Order( gens1[1] );  Order( gens2[1] );|
  4
  2
  !gapprompt@gap>| !gapinput@Sqrt(5)^gens1[1] = Sqrt(5)^gens2[1];|
  true
\end{Verbatim}
 

 The following example shows the Galois group of a cyclotomic field over a
non\texttt{\symbol{45}}cyclotomic field. 

 
\begin{Verbatim}[commandchars=!@|,fontsize=\small,frame=single,label=Example]
  !gapprompt@gap>| !gapinput@g:= GaloisGroup( AsField( Field( [ Sqrt(5) ] ), CF(5) ) );|
  <group of size 2 with 1 generator>
  !gapprompt@gap>| !gapinput@gens:= GeneratorsOfGroup( g );|
  [ ANFAutomorphism( AsField( NF(5,[ 1, 4 ]), CF(5) ), 4 ) ]
  !gapprompt@gap>| !gapinput@x:= last[1];;  x^2;|
  IdentityMapping( AsField( NF(5,[ 1, 4 ]), CF(5) ) )
\end{Verbatim}
 }

 

\subsection{\textcolor{Chapter }{ANFAutomorphism}}
\logpage{[ 60, 4, 2 ]}\nobreak
\hyperdef{L}{X8643D4B47A827D9D}{}
{\noindent\textcolor{FuncColor}{$\triangleright$\enspace\texttt{ANFAutomorphism({\mdseries\slshape F, k})\index{ANFAutomorphism@\texttt{ANFAutomorphism}}
\label{ANFAutomorphism}
}\hfill{\scriptsize (function)}}\\


 Let \mbox{\texttt{\mdseries\slshape F}} be an abelian number field and \mbox{\texttt{\mdseries\slshape k}} be an integer that is coprime to the conductor (see \texttt{Conductor} (\ref{Conductor:for a collection of cyclotomics})) of \mbox{\texttt{\mdseries\slshape F}}. Then \texttt{ANFAutomorphism} returns the automorphism of \mbox{\texttt{\mdseries\slshape F}} that is defined as the linear extension of the map that raises each root of
unity in \mbox{\texttt{\mdseries\slshape F}} to its \mbox{\texttt{\mdseries\slshape k}}\texttt{\symbol{45}}th power. 

 
\begin{Verbatim}[commandchars=!@|,fontsize=\small,frame=single,label=Example]
  !gapprompt@gap>| !gapinput@f:= CF(25);|
  CF(25)
  !gapprompt@gap>| !gapinput@alpha:= ANFAutomorphism( f, 2 );|
  ANFAutomorphism( CF(25), 2 )
  !gapprompt@gap>| !gapinput@alpha^2;|
  ANFAutomorphism( CF(25), 4 )
  !gapprompt@gap>| !gapinput@Order( alpha );|
  20
  !gapprompt@gap>| !gapinput@E(5)^alpha;|
  E(5)^2
\end{Verbatim}
 }

 }

  
\section{\textcolor{Chapter }{Gaussians}}\label{Gaussians}
\logpage{[ 60, 5, 0 ]}
\hyperdef{L}{X85E9E90D7FE877CC}{}
{
  

\subsection{\textcolor{Chapter }{GaussianIntegers}}
\logpage{[ 60, 5, 1 ]}\nobreak
\hyperdef{L}{X80BD5EAB879F096E}{}
{\noindent\textcolor{FuncColor}{$\triangleright$\enspace\texttt{GaussianIntegers\index{GaussianIntegers@\texttt{GaussianIntegers}}
\label{GaussianIntegers}
}\hfill{\scriptsize (global variable)}}\\


 \texttt{GaussianIntegers} is the ring ${\ensuremath{\mathbb Z}}[\sqrt{{-1}}]$ of Gaussian integers. This is a subring of the cyclotomic field \texttt{GaussianRationals} (\ref{GaussianRationals}). }

 

\subsection{\textcolor{Chapter }{IsGaussianIntegers}}
\logpage{[ 60, 5, 2 ]}\nobreak
\hyperdef{L}{X7BFD33D27BFB7C5A}{}
{\noindent\textcolor{FuncColor}{$\triangleright$\enspace\texttt{IsGaussianIntegers({\mdseries\slshape obj})\index{IsGaussianIntegers@\texttt{IsGaussianIntegers}}
\label{IsGaussianIntegers}
}\hfill{\scriptsize (Category)}}\\


 is the defining category for the domain \texttt{GaussianIntegers} (\ref{GaussianIntegers}). }

  }

 }

             
\chapter{\textcolor{Chapter }{Vector Spaces}}\label{Vector Spaces}
\logpage{[ 61, 0, 0 ]}
\hyperdef{L}{X7DAD6700787EC845}{}
{
   
\section{\textcolor{Chapter }{IsLeftVectorSpace (Filter)}}\label{sect:IsLeftVectorSpace}
\logpage{[ 61, 1, 0 ]}
\hyperdef{L}{X8754F7207CFDA38B}{}
{
  

\subsection{\textcolor{Chapter }{IsLeftVectorSpace}}
\logpage{[ 61, 1, 1 ]}\nobreak
\hyperdef{L}{X80290A908241706B}{}
{\noindent\textcolor{FuncColor}{$\triangleright$\enspace\texttt{IsLeftVectorSpace({\mdseries\slshape V})\index{IsLeftVectorSpace@\texttt{IsLeftVectorSpace}}
\label{IsLeftVectorSpace}
}\hfill{\scriptsize (Category)}}\\
\noindent\textcolor{FuncColor}{$\triangleright$\enspace\texttt{IsVectorSpace({\mdseries\slshape V})\index{IsVectorSpace@\texttt{IsVectorSpace}}
\label{IsVectorSpace}
}\hfill{\scriptsize (Category)}}\\


 A \emph{vector space} in \textsf{GAP} is a free left module (see{\nobreakspace}\texttt{IsFreeLeftModule} (\ref{IsFreeLeftModule})) over a division ring (see Chapter{\nobreakspace}\ref{Fields and Division Rings}). 

 Whenever we talk about an $F$\texttt{\symbol{45}}vector space \mbox{\texttt{\mdseries\slshape V}} then \mbox{\texttt{\mdseries\slshape V}} is an additive group (see{\nobreakspace}\texttt{IsAdditiveGroup} (\ref{IsAdditiveGroup})) on which the division ring $F$ acts via multiplication from the left such that this action and the addition
in \mbox{\texttt{\mdseries\slshape V}} are left and right distributive. The division ring $F$ can be accessed as value of the attribute \texttt{LeftActingDomain} (\ref{LeftActingDomain}). 

 Vector spaces in \textsf{GAP} are always \emph{left} vector spaces, \texttt{IsLeftVectorSpace} and \texttt{IsVectorSpace} are synonyms. }

 }

  
\section{\textcolor{Chapter }{Constructing Vector Spaces}}\label{Constructing Vector Spaces}
\logpage{[ 61, 2, 0 ]}
\hyperdef{L}{X87AD06FE873619EA}{}
{
  

\subsection{\textcolor{Chapter }{VectorSpace}}
\logpage{[ 61, 2, 1 ]}\nobreak
\hyperdef{L}{X805413157CE9BECF}{}
{\noindent\textcolor{FuncColor}{$\triangleright$\enspace\texttt{VectorSpace({\mdseries\slshape F, gens[, zero][, "basis"]})\index{VectorSpace@\texttt{VectorSpace}}
\label{VectorSpace}
}\hfill{\scriptsize (function)}}\\


 For a field \mbox{\texttt{\mdseries\slshape F}} and a collection \mbox{\texttt{\mdseries\slshape gens}} of vectors, \texttt{VectorSpace} returns the \mbox{\texttt{\mdseries\slshape F}}\texttt{\symbol{45}}vector space spanned by the elements in \mbox{\texttt{\mdseries\slshape gens}}. 

 The optional argument \mbox{\texttt{\mdseries\slshape zero}} can be used to specify the zero element of the space; \mbox{\texttt{\mdseries\slshape zero}} \emph{must} be given if \mbox{\texttt{\mdseries\slshape gens}} is empty. The optional string \texttt{"basis"} indicates that \mbox{\texttt{\mdseries\slshape gens}} is known to be linearly independent over \mbox{\texttt{\mdseries\slshape F}}, in particular the dimension of the vector space is immediately set; note
that \texttt{Basis} (\ref{Basis}) need \emph{not} return the basis formed by \mbox{\texttt{\mdseries\slshape gens}} if the string \texttt{"basis"} is given as an argument.  
\begin{Verbatim}[commandchars=!@|,fontsize=\small,frame=single,label=Example]
  !gapprompt@gap>| !gapinput@V:= VectorSpace( Rationals, [ [ 1, 2, 3 ], [ 1, 1, 1 ] ] );|
  <vector space over Rationals, with 2 generators>
\end{Verbatim}
 }

 

\subsection{\textcolor{Chapter }{Subspace}}
\logpage{[ 61, 2, 2 ]}\nobreak
\hyperdef{L}{X78C9826780BC9AE6}{}
{\noindent\textcolor{FuncColor}{$\triangleright$\enspace\texttt{Subspace({\mdseries\slshape V, gens[, "basis"]})\index{Subspace@\texttt{Subspace}}
\label{Subspace}
}\hfill{\scriptsize (function)}}\\
\noindent\textcolor{FuncColor}{$\triangleright$\enspace\texttt{SubspaceNC({\mdseries\slshape V, gens[, "basis"]})\index{SubspaceNC@\texttt{SubspaceNC}}
\label{SubspaceNC}
}\hfill{\scriptsize (function)}}\\


 For an $F$\texttt{\symbol{45}}vector space \mbox{\texttt{\mdseries\slshape V}} and a list or collection \mbox{\texttt{\mdseries\slshape gens}} that is a subset of \mbox{\texttt{\mdseries\slshape V}}, \texttt{Subspace} returns the $F$\texttt{\symbol{45}}vector space spanned by \mbox{\texttt{\mdseries\slshape gens}}; if \mbox{\texttt{\mdseries\slshape gens}} is empty then the trivial subspace (see{\nobreakspace}\texttt{TrivialSubspace} (\ref{TrivialSubspace})) of \mbox{\texttt{\mdseries\slshape V}} is returned. The parent (see{\nobreakspace}\ref{Parents}) of the returned vector space is set to \mbox{\texttt{\mdseries\slshape V}}. 

 \texttt{SubspaceNC} does the same as \texttt{Subspace}, except that it omits the check whether \mbox{\texttt{\mdseries\slshape gens}} is a subset of \mbox{\texttt{\mdseries\slshape V}}. 

 The optional string \mbox{\texttt{\mdseries\slshape "basis"}} indicates that \mbox{\texttt{\mdseries\slshape gens}} is known to be linearly independent over $F$. In this case the dimension of the subspace is immediately set, and both \texttt{Subspace} and \texttt{SubspaceNC} do \emph{not} check whether \mbox{\texttt{\mdseries\slshape gens}} really is linearly independent and whether \mbox{\texttt{\mdseries\slshape gens}} is a subset of \mbox{\texttt{\mdseries\slshape V}}.  
\begin{Verbatim}[commandchars=!@|,fontsize=\small,frame=single,label=Example]
  !gapprompt@gap>| !gapinput@V:= VectorSpace( Rationals, [ [ 1, 2, 3 ], [ 1, 1, 1 ] ] );;|
  !gapprompt@gap>| !gapinput@W:= Subspace( V, [ [ 0, 1, 2 ] ] );|
  <vector space over Rationals, with 1 generator>
\end{Verbatim}
 }

 

\subsection{\textcolor{Chapter }{AsVectorSpace}}
\logpage{[ 61, 2, 3 ]}\nobreak
\hyperdef{L}{X7B001BAF7D5FD5D0}{}
{\noindent\textcolor{FuncColor}{$\triangleright$\enspace\texttt{AsVectorSpace({\mdseries\slshape F, D})\index{AsVectorSpace@\texttt{AsVectorSpace}}
\label{AsVectorSpace}
}\hfill{\scriptsize (operation)}}\\


 Let \mbox{\texttt{\mdseries\slshape F}} be a division ring and \mbox{\texttt{\mdseries\slshape D}} a domain. If the elements in \mbox{\texttt{\mdseries\slshape D}} form an \mbox{\texttt{\mdseries\slshape F}}\texttt{\symbol{45}}vector space then \texttt{AsVectorSpace} returns this \mbox{\texttt{\mdseries\slshape F}}\texttt{\symbol{45}}vector space, otherwise \texttt{fail} is returned. 

 \texttt{AsVectorSpace} can be used for example to view a given vector space as a vector space over a
smaller or larger division ring. 
\begin{Verbatim}[commandchars=!@|,fontsize=\small,frame=single,label=Example]
  !gapprompt@gap>| !gapinput@V:= FullRowSpace( GF( 27 ), 3 );|
  ( GF(3^3)^3 )
  !gapprompt@gap>| !gapinput@Dimension( V );  LeftActingDomain( V );|
  3
  GF(3^3)
  !gapprompt@gap>| !gapinput@W:= AsVectorSpace( GF( 3 ), V );|
  <vector space over GF(3), with 9 generators>
  !gapprompt@gap>| !gapinput@Dimension( W );  LeftActingDomain( W );|
  9
  GF(3)
  !gapprompt@gap>| !gapinput@AsVectorSpace( GF( 9 ), V );|
  fail
\end{Verbatim}
 }

 

\subsection{\textcolor{Chapter }{AsSubspace}}
\logpage{[ 61, 2, 4 ]}\nobreak
\hyperdef{L}{X7D4F84C27EDAC89B}{}
{\noindent\textcolor{FuncColor}{$\triangleright$\enspace\texttt{AsSubspace({\mdseries\slshape V, U})\index{AsSubspace@\texttt{AsSubspace}}
\label{AsSubspace}
}\hfill{\scriptsize (operation)}}\\


 Let \mbox{\texttt{\mdseries\slshape V}} be an $F$\texttt{\symbol{45}}vector space, and \mbox{\texttt{\mdseries\slshape U}} a collection. If \mbox{\texttt{\mdseries\slshape U}} is a subset of \mbox{\texttt{\mdseries\slshape V}} such that the elements of \mbox{\texttt{\mdseries\slshape U}} form an $F$\texttt{\symbol{45}}vector space then \texttt{AsSubspace} returns this vector space, with parent set to \mbox{\texttt{\mdseries\slshape V}} (see{\nobreakspace}\texttt{AsVectorSpace} (\ref{AsVectorSpace})). Otherwise \texttt{fail} is returned. 
\begin{Verbatim}[commandchars=!@|,fontsize=\small,frame=single,label=Example]
  !gapprompt@gap>| !gapinput@V:= VectorSpace( Rationals, [ [ 1, 2, 3 ], [ 1, 1, 1 ] ] );;|
  !gapprompt@gap>| !gapinput@W:= VectorSpace( Rationals, [ [ 1/2, 1/2, 1/2 ] ] );;|
  !gapprompt@gap>| !gapinput@U:= AsSubspace( V, W );|
  <vector space over Rationals, with 1 generator>
  !gapprompt@gap>| !gapinput@Parent( U ) = V;|
  true
  !gapprompt@gap>| !gapinput@AsSubspace( V, [ [ 1, 1, 1 ] ] );|
  fail
\end{Verbatim}
 }

 }

  
\section{\textcolor{Chapter }{Operations and Attributes for Vector Spaces}}\label{Operations and Attributes for Vector Spaces}
\logpage{[ 61, 3, 0 ]}
\hyperdef{L}{X789FB2D883E53662}{}
{
  

\subsection{\textcolor{Chapter }{GeneratorsOfLeftVectorSpace}}
\logpage{[ 61, 3, 1 ]}\nobreak
\hyperdef{L}{X849651C6830C94A1}{}
{\noindent\textcolor{FuncColor}{$\triangleright$\enspace\texttt{GeneratorsOfLeftVectorSpace({\mdseries\slshape V})\index{GeneratorsOfLeftVectorSpace@\texttt{GeneratorsOfLeftVectorSpace}}
\label{GeneratorsOfLeftVectorSpace}
}\hfill{\scriptsize (attribute)}}\\
\noindent\textcolor{FuncColor}{$\triangleright$\enspace\texttt{GeneratorsOfVectorSpace({\mdseries\slshape V})\index{GeneratorsOfVectorSpace@\texttt{GeneratorsOfVectorSpace}}
\label{GeneratorsOfVectorSpace}
}\hfill{\scriptsize (attribute)}}\\


 For an $F$\texttt{\symbol{45}}vector space \mbox{\texttt{\mdseries\slshape V}}, \texttt{GeneratorsOfLeftVectorSpace} returns a list of vectors in \mbox{\texttt{\mdseries\slshape V}} that generate \mbox{\texttt{\mdseries\slshape V}} as an $F$\texttt{\symbol{45}}vector space. 
\begin{Verbatim}[commandchars=!@|,fontsize=\small,frame=single,label=Example]
  !gapprompt@gap>| !gapinput@GeneratorsOfVectorSpace( FullRowSpace( Rationals, 3 ) );|
  [ [ 1, 0, 0 ], [ 0, 1, 0 ], [ 0, 0, 1 ] ]
\end{Verbatim}
 }

 

\subsection{\textcolor{Chapter }{TrivialSubspace}}
\logpage{[ 61, 3, 2 ]}\nobreak
\hyperdef{L}{X86DC71A9835430FD}{}
{\noindent\textcolor{FuncColor}{$\triangleright$\enspace\texttt{TrivialSubspace({\mdseries\slshape V})\index{TrivialSubspace@\texttt{TrivialSubspace}}
\label{TrivialSubspace}
}\hfill{\scriptsize (attribute)}}\\


 For a vector space \mbox{\texttt{\mdseries\slshape V}}, \texttt{TrivialSubspace} returns the subspace of \mbox{\texttt{\mdseries\slshape V}} that consists of the zero vector in \mbox{\texttt{\mdseries\slshape V}}. 
\begin{Verbatim}[commandchars=!@|,fontsize=\small,frame=single,label=Example]
  !gapprompt@gap>| !gapinput@V:= GF(3)^3;;|
  !gapprompt@gap>| !gapinput@triv:= TrivialSubspace( V );|
  <vector space of dimension 0 over GF(3)>
  !gapprompt@gap>| !gapinput@AsSet( triv );|
  [ [ 0*Z(3), 0*Z(3), 0*Z(3) ] ]
\end{Verbatim}
 }

 }

  
\section{\textcolor{Chapter }{Domains of Subspaces of Vector Spaces}}\label{Domains of Subspaces of Vector Spaces}
\logpage{[ 61, 4, 0 ]}
\hyperdef{L}{X8125675583357131}{}
{
  

\subsection{\textcolor{Chapter }{Subspaces}}
\logpage{[ 61, 4, 1 ]}\nobreak
\hyperdef{L}{X7975E41A7B29C3FD}{}
{\noindent\textcolor{FuncColor}{$\triangleright$\enspace\texttt{Subspaces({\mdseries\slshape V[, k]})\index{Subspaces@\texttt{Subspaces}}
\label{Subspaces}
}\hfill{\scriptsize (attribute)}}\\


 Called with a finite vector space \mbox{\texttt{\mdseries\slshape V}}, \texttt{Subspaces} returns the domain of all subspaces of \mbox{\texttt{\mdseries\slshape V}}. 

 Called with \mbox{\texttt{\mdseries\slshape V}} and a nonnegative integer \mbox{\texttt{\mdseries\slshape k}}, \texttt{Subspaces} returns the domain of all \mbox{\texttt{\mdseries\slshape k}}\texttt{\symbol{45}}dimensional subspaces of \mbox{\texttt{\mdseries\slshape V}}. 

 Special \texttt{Size} (\ref{Size}) and \texttt{Iterator} (\ref{Iterator}) methods are provided for these domains.  }

 

\subsection{\textcolor{Chapter }{IsSubspacesVectorSpace}}
\logpage{[ 61, 4, 2 ]}\nobreak
\hyperdef{L}{X7A8F5C367FAE3D1B}{}
{\noindent\textcolor{FuncColor}{$\triangleright$\enspace\texttt{IsSubspacesVectorSpace({\mdseries\slshape D})\index{IsSubspacesVectorSpace@\texttt{IsSubspacesVectorSpace}}
\label{IsSubspacesVectorSpace}
}\hfill{\scriptsize (Category)}}\\


 The domain of all subspaces of a (finite) vector space or of all subspaces of
fixed dimension, as returned by \texttt{Subspaces} (\ref{Subspaces}) (see{\nobreakspace}\texttt{Subspaces} (\ref{Subspaces})) lies in the category \texttt{IsSubspacesVectorSpace}. 
\begin{Verbatim}[commandchars=!@|,fontsize=\small,frame=single,label=Example]
  !gapprompt@gap>| !gapinput@D:= Subspaces( GF(3)^3 );|
  Subspaces( ( GF(3)^3 ) )
  !gapprompt@gap>| !gapinput@Size( D );|
  28
  !gapprompt@gap>| !gapinput@iter:= Iterator( D );;|
  !gapprompt@gap>| !gapinput@NextIterator( iter );|
  <vector space of dimension 0 over GF(3)>
  !gapprompt@gap>| !gapinput@NextIterator( iter );|
  <vector space of dimension 1 over GF(3)>
  !gapprompt@gap>| !gapinput@IsSubspacesVectorSpace( D );|
  true
\end{Verbatim}
 }

    }

  
\section{\textcolor{Chapter }{Bases of Vector Spaces}}\label{Bases of Vector Spaces}
\logpage{[ 61, 5, 0 ]}
\hyperdef{L}{X828AA09B87F14FAD}{}
{
  In \textsf{GAP}, a \emph{basis} of a free left $F$\texttt{\symbol{45}}module $V$ is a list of vectors $B = [ v_1, v_2, \ldots, v_n ]$ in $V$ such that $V$ is generated as a left $F$\texttt{\symbol{45}}module by these vectors and such that $B$ is linearly independent over $F$. The integer $n$ is the dimension of $V$ (see{\nobreakspace}\texttt{Dimension} (\ref{Dimension})). In particular, as each basis is a list (see Chapter{\nobreakspace}\ref{Lists}), it has a length (see{\nobreakspace}\texttt{Length} (\ref{Length})), and the $i$\texttt{\symbol{45}}th vector of $B$ can be accessed as $B[i]$. 
\begin{Verbatim}[commandchars=!@|,fontsize=\small,frame=single,label=Example]
  !gapprompt@gap>| !gapinput@V:= Rationals^3;|
  ( Rationals^3 )
  !gapprompt@gap>| !gapinput@B:= Basis( V );|
  CanonicalBasis( ( Rationals^3 ) )
  !gapprompt@gap>| !gapinput@Length( B );|
  3
  !gapprompt@gap>| !gapinput@B[1];|
  [ 1, 0, 0 ]
\end{Verbatim}
 

 The operations described below make sense only for bases of \emph{finite} dimensional vector spaces. (In practice this means that the vector spaces must
be \emph{low} dimensional, that is, the dimension should not exceed a few hundred.) 

 Besides the basic operations for lists (see{\nobreakspace}\ref{Basic Operations for Lists}), the \emph{basic operations for bases} are \texttt{BasisVectors} (\ref{BasisVectors}), \texttt{Coefficients} (\ref{Coefficients}), \texttt{LinearCombination} (\ref{LinearCombination}), and \texttt{UnderlyingLeftModule} (\ref{UnderlyingLeftModule}). These and other operations for arbitrary bases are described
in{\nobreakspace}\ref{Operations for Vector Space Bases}. 

 For special kinds of bases, further operations are defined (see{\nobreakspace}\ref{Operations for Special Kinds of Bases}). 

 \textsf{GAP} supports the following three kinds of bases. 

 \emph{Relative bases} delegate the work to other bases of the same free left module, via basechange
matrices (see{\nobreakspace}\texttt{RelativeBasis} (\ref{RelativeBasis})). 

 \emph{Bases handled by nice bases} delegate the work to bases of isomorphic left modules over the same left
acting domain (see{\nobreakspace}\ref{Vector Spaces Handled By Nice Bases}). 

 Finally, of course there must be bases in \textsf{GAP} that really do the work. 

 For example, in the case of a Gaussian row or matrix space \mbox{\texttt{\mdseries\slshape V}} (see{\nobreakspace}\ref{Row and Matrix Spaces}), \texttt{Basis( \mbox{\texttt{\mdseries\slshape V}} )} is a semi\texttt{\symbol{45}}echelonized basis (see{\nobreakspace}\texttt{IsSemiEchelonized} (\ref{IsSemiEchelonized})) that uses Gaussian elimination; such a basis is of the third kind. \texttt{Basis( \mbox{\texttt{\mdseries\slshape V}}, \mbox{\texttt{\mdseries\slshape vectors}} )} is either semi\texttt{\symbol{45}}echelonized or a relative basis. Other
examples of bases of the third kind are canonical bases of finite fields and
of abelian number fields. 

 Bases handled by nice bases are described in{\nobreakspace}\ref{Vector Spaces Handled By Nice Bases}. Examples are non\texttt{\symbol{45}}Gaussian row and matrix spaces, and
subspaces of finite fields and abelian number fields that are themselves not
fields. 

\subsection{\textcolor{Chapter }{IsBasis}}
\logpage{[ 61, 5, 1 ]}\nobreak
\hyperdef{L}{X8739510881F5D862}{}
{\noindent\textcolor{FuncColor}{$\triangleright$\enspace\texttt{IsBasis({\mdseries\slshape obj})\index{IsBasis@\texttt{IsBasis}}
\label{IsBasis}
}\hfill{\scriptsize (Category)}}\\


 In \textsf{GAP}, a \emph{basis} of a free left module is an object that knows how to compute coefficients
w.r.t.{\nobreakspace}its basis vectors (see{\nobreakspace}\texttt{Coefficients} (\ref{Coefficients})). Bases are constructed by \texttt{Basis} (\ref{Basis}). Each basis is an immutable list, the $i$\texttt{\symbol{45}}th entry being the $i$\texttt{\symbol{45}}th basis vector. 

 (See{\nobreakspace}\ref{Mutable Bases} for mutable bases.) 

 
\begin{Verbatim}[commandchars=!@|,fontsize=\small,frame=single,label=Example]
  !gapprompt@gap>| !gapinput@V:= GF(2)^2;;|
  !gapprompt@gap>| !gapinput@B:= Basis( V );;|
  !gapprompt@gap>| !gapinput@IsBasis( B );|
  true
  !gapprompt@gap>| !gapinput@IsBasis( [ [ 1, 0 ], [ 0, 1 ] ] );|
  false
  !gapprompt@gap>| !gapinput@IsBasis( Basis( Rationals^2, [ [ 1, 0 ], [ 0, 1 ] ] ) );|
  true
\end{Verbatim}
 }

 

\subsection{\textcolor{Chapter }{Basis}}
\logpage{[ 61, 5, 2 ]}\nobreak
\hyperdef{L}{X837BE54C80DE368E}{}
{\noindent\textcolor{FuncColor}{$\triangleright$\enspace\texttt{Basis({\mdseries\slshape V[, vectors]})\index{Basis@\texttt{Basis}}
\label{Basis}
}\hfill{\scriptsize (attribute)}}\\
\noindent\textcolor{FuncColor}{$\triangleright$\enspace\texttt{BasisNC({\mdseries\slshape V, vectors})\index{BasisNC@\texttt{BasisNC}}
\label{BasisNC}
}\hfill{\scriptsize (operation)}}\\


 Called with a free left $F$\texttt{\symbol{45}}module \mbox{\texttt{\mdseries\slshape V}} as the only argument, \texttt{Basis} returns an $F$\texttt{\symbol{45}}basis of \mbox{\texttt{\mdseries\slshape V}} whose vectors are not further specified. 

 If additionally a list \mbox{\texttt{\mdseries\slshape vectors}} of vectors in \mbox{\texttt{\mdseries\slshape V}} is given that forms an $F$\texttt{\symbol{45}}basis of \mbox{\texttt{\mdseries\slshape V}} then \texttt{Basis} returns this basis; if \mbox{\texttt{\mdseries\slshape vectors}} is not linearly independent over $F$ or does not generate \mbox{\texttt{\mdseries\slshape V}} as a free left $F$\texttt{\symbol{45}}module then \texttt{fail} is returned. 

 \texttt{BasisNC} does the same as the two argument version of \texttt{Basis}, except that it does not check whether \mbox{\texttt{\mdseries\slshape vectors}} form a basis. 

 If no basis vectors are prescribed then \texttt{Basis} need not compute basis vectors; in this case, the vectors are computed in the
first call to \texttt{BasisVectors} (\ref{BasisVectors}). 
\begin{Verbatim}[commandchars=!@|,fontsize=\small,frame=single,label=Example]
  !gapprompt@gap>| !gapinput@V:= VectorSpace( Rationals, [ [ 1, 2, 7 ], [ 1/2, 1/3, 5 ] ] );;|
  !gapprompt@gap>| !gapinput@B:= Basis( V );|
  SemiEchelonBasis( <vector space over Rationals, with
  2 generators>, ... )
  !gapprompt@gap>| !gapinput@BasisVectors( B );|
  [ [ 1, 2, 7 ], [ 0, 1, -9/4 ] ]
  !gapprompt@gap>| !gapinput@B:= Basis( V, [ [ 1, 2, 7 ], [ 3, 2, 30 ] ] );|
  Basis( <vector space over Rationals, with 2 generators>,
  [ [ 1, 2, 7 ], [ 3, 2, 30 ] ] )
  !gapprompt@gap>| !gapinput@Basis( V, [ [ 1, 2, 3 ] ] );|
  fail
\end{Verbatim}
 }

  

\subsection{\textcolor{Chapter }{CanonicalBasis}}
\logpage{[ 61, 5, 3 ]}\nobreak
\hyperdef{L}{X7C8EBFF5805F8C51}{}
{\noindent\textcolor{FuncColor}{$\triangleright$\enspace\texttt{CanonicalBasis({\mdseries\slshape V})\index{CanonicalBasis@\texttt{CanonicalBasis}}
\label{CanonicalBasis}
}\hfill{\scriptsize (attribute)}}\\


 If the vector space \mbox{\texttt{\mdseries\slshape V}} supports a \emph{canonical basis} then \texttt{CanonicalBasis} returns this basis, otherwise \texttt{fail} is returned. 

 The defining property of a canonical basis is that its vectors are uniquely
determined by the vector space. If canonical bases exist for two vector spaces
over the same left acting domain (see{\nobreakspace}\texttt{LeftActingDomain} (\ref{LeftActingDomain})) then the equality of these vector spaces can be decided by comparing the
canonical bases. 

 The exact meaning of a canonical basis depends on the type of \mbox{\texttt{\mdseries\slshape V}}. Canonical bases are defined for example for Gaussian row and matrix spaces
(see{\nobreakspace}\ref{Row and Matrix Spaces}). 

 If one designs a new kind of vector spaces (see{\nobreakspace}\ref{How to Implement New Kinds of Vector Spaces}) and defines a canonical basis for these spaces then the \texttt{CanonicalBasis} method one installs (see{\nobreakspace}\texttt{InstallMethod} (\ref{InstallMethod})) must \emph{not} call \texttt{Basis} (\ref{Basis}). On the other hand, one probably should install a \texttt{Basis} (\ref{Basis}) method that simply calls \texttt{CanonicalBasis}, the value of the method (see{\nobreakspace}\ref{Method Installation} and \ref{Applicable Methods and Method Selection}) being \texttt{CANONICAL{\textunderscore}BASIS{\textunderscore}FLAGS}. 
\begin{Verbatim}[commandchars=!@|,fontsize=\small,frame=single,label=Example]
  !gapprompt@gap>| !gapinput@vecs:= [ [ 1, 2, 3 ], [ 1, 1, 1 ], [ 1, 1, 1 ] ];;|
  !gapprompt@gap>| !gapinput@V:= VectorSpace( Rationals, vecs );;|
  !gapprompt@gap>| !gapinput@B:= CanonicalBasis( V );|
  CanonicalBasis( <vector space over Rationals, with 3 generators> )
  !gapprompt@gap>| !gapinput@BasisVectors( B );|
  [ [ 1, 0, -1 ], [ 0, 1, 2 ] ]
\end{Verbatim}
 }

 

\subsection{\textcolor{Chapter }{RelativeBasis}}
\logpage{[ 61, 5, 4 ]}\nobreak
\hyperdef{L}{X8786D40B84120F38}{}
{\noindent\textcolor{FuncColor}{$\triangleright$\enspace\texttt{RelativeBasis({\mdseries\slshape B, vectors})\index{RelativeBasis@\texttt{RelativeBasis}}
\label{RelativeBasis}
}\hfill{\scriptsize (operation)}}\\
\noindent\textcolor{FuncColor}{$\triangleright$\enspace\texttt{RelativeBasisNC({\mdseries\slshape B, vectors})\index{RelativeBasisNC@\texttt{RelativeBasisNC}}
\label{RelativeBasisNC}
}\hfill{\scriptsize (operation)}}\\


 A relative basis is a basis of the free left module \mbox{\texttt{\mdseries\slshape V}} that delegates the computation of coefficients etc. to another basis of \mbox{\texttt{\mdseries\slshape V}} via a basechange matrix. 

 Let \mbox{\texttt{\mdseries\slshape B}} be a basis of the free left module \mbox{\texttt{\mdseries\slshape V}}, and \mbox{\texttt{\mdseries\slshape vectors}} a list of vectors in \mbox{\texttt{\mdseries\slshape V}}. 

 \texttt{RelativeBasis} checks whether \mbox{\texttt{\mdseries\slshape vectors}} form a basis of \mbox{\texttt{\mdseries\slshape V}}, and in this case a basis is returned in which \mbox{\texttt{\mdseries\slshape vectors}} are the basis vectors; otherwise \texttt{fail} is returned. 

 \texttt{RelativeBasisNC} does the same, except that it omits the check. }

 }

  
\section{\textcolor{Chapter }{Operations for Vector Space Bases}}\label{Operations for Vector Space Bases}
\logpage{[ 61, 6, 0 ]}
\hyperdef{L}{X839B9C4880EBFB5F}{}
{
  

\subsection{\textcolor{Chapter }{BasisVectors}}
\logpage{[ 61, 6, 1 ]}\nobreak
\hyperdef{L}{X7B1F17AE8027A590}{}
{\noindent\textcolor{FuncColor}{$\triangleright$\enspace\texttt{BasisVectors({\mdseries\slshape B})\index{BasisVectors@\texttt{BasisVectors}}
\label{BasisVectors}
}\hfill{\scriptsize (attribute)}}\\


 For a vector space basis \mbox{\texttt{\mdseries\slshape B}}, \texttt{BasisVectors} returns the list of basis vectors of \mbox{\texttt{\mdseries\slshape B}}. The lists \mbox{\texttt{\mdseries\slshape B}} and \texttt{BasisVectors( \mbox{\texttt{\mdseries\slshape B}} )} are equal; the main purpose of \texttt{BasisVectors} is to provide access to a list of vectors that does \emph{not} know about an underlying vector space. 
\begin{Verbatim}[commandchars=!@|,fontsize=\small,frame=single,label=Example]
  !gapprompt@gap>| !gapinput@V:= VectorSpace( Rationals, [ [ 1, 2, 7 ], [ 1/2, 1/3, 5 ] ] );;|
  !gapprompt@gap>| !gapinput@B:= Basis( V, [ [ 1, 2, 7 ], [ 0, 1, -9/4 ] ] );;|
  !gapprompt@gap>| !gapinput@BasisVectors( B );|
  [ [ 1, 2, 7 ], [ 0, 1, -9/4 ] ]
\end{Verbatim}
 }

 

\subsection{\textcolor{Chapter }{UnderlyingLeftModule}}
\logpage{[ 61, 6, 2 ]}\nobreak
\hyperdef{L}{X81E8AE88843B70FF}{}
{\noindent\textcolor{FuncColor}{$\triangleright$\enspace\texttt{UnderlyingLeftModule({\mdseries\slshape B})\index{UnderlyingLeftModule@\texttt{UnderlyingLeftModule}}
\label{UnderlyingLeftModule}
}\hfill{\scriptsize (attribute)}}\\


 For a basis \mbox{\texttt{\mdseries\slshape B}} of a free left module $V$, \texttt{UnderlyingLeftModule} returns $V$. 

 The reason why a basis stores a free left module is that otherwise one would
have to store the basis vectors and the coefficient domain separately. Storing
the module allows one for example to deal with bases whose basis vectors have
not yet been computed yet (see{\nobreakspace}\texttt{Basis} (\ref{Basis})); furthermore, in some cases it is convenient to test membership of a vector
in the module before computing coefficients w.r.t.{\nobreakspace}a basis.  
\begin{Verbatim}[commandchars=!@|,fontsize=\small,frame=single,label=Example]
  !gapprompt@gap>| !gapinput@B:= Basis( GF(2)^6 );;  UnderlyingLeftModule( B );|
  ( GF(2)^6 )
\end{Verbatim}
 }

 

\subsection{\textcolor{Chapter }{Coefficients}}
\logpage{[ 61, 6, 3 ]}\nobreak
\hyperdef{L}{X80B32F667BF6AFD8}{}
{\noindent\textcolor{FuncColor}{$\triangleright$\enspace\texttt{Coefficients({\mdseries\slshape B, v})\index{Coefficients@\texttt{Coefficients}}
\label{Coefficients}
}\hfill{\scriptsize (operation)}}\\


 Let $V$ be the underlying left module of the basis \mbox{\texttt{\mdseries\slshape B}}, and \mbox{\texttt{\mdseries\slshape v}} a vector such that the family of \mbox{\texttt{\mdseries\slshape v}} is the elements family of the family of $V$. Then \texttt{Coefficients( \mbox{\texttt{\mdseries\slshape B}}, \mbox{\texttt{\mdseries\slshape v}} )} is the list of coefficients of \mbox{\texttt{\mdseries\slshape v}} w.r.t. \mbox{\texttt{\mdseries\slshape B}} if \mbox{\texttt{\mdseries\slshape v}} lies in $V$, and \texttt{fail} otherwise. 
\begin{Verbatim}[commandchars=!@|,fontsize=\small,frame=single,label=Example]
  !gapprompt@gap>| !gapinput@V:= VectorSpace( Rationals, [ [ 1, 2, 7 ], [ 1/2, 1/3, 5 ] ] );;|
  !gapprompt@gap>| !gapinput@B:= Basis( V, [ [ 1, 2, 7 ], [ 0, 1, -9/4 ] ] );;|
  !gapprompt@gap>| !gapinput@Coefficients( B, [ 1/2, 1/3, 5 ] );|
  [ 1/2, -2/3 ]
  !gapprompt@gap>| !gapinput@Coefficients( B, [ 1, 0, 0 ] );|
  fail
\end{Verbatim}
 }

 

\subsection{\textcolor{Chapter }{LinearCombination}}
\logpage{[ 61, 6, 4 ]}\nobreak
\hyperdef{L}{X7D305AB3834889BF}{}
{\noindent\textcolor{FuncColor}{$\triangleright$\enspace\texttt{LinearCombination({\mdseries\slshape B, coeff})\index{LinearCombination@\texttt{LinearCombination}}
\label{LinearCombination}
}\hfill{\scriptsize (operation)}}\\


 If \mbox{\texttt{\mdseries\slshape B}} is a basis object (see \texttt{IsBasis} (\ref{IsBasis})) or a homogeneous list of length $n$, and \mbox{\texttt{\mdseries\slshape coeff}} is a row vector of the same length, \texttt{LinearCombination} returns the vector $\sum_{{i = 1}}^n \mbox{\texttt{\mdseries\slshape coeff}}[i] * \mbox{\texttt{\mdseries\slshape B}}[i]$. 

 Perhaps the most important usage is the case where \mbox{\texttt{\mdseries\slshape B}} forms a basis. 
\begin{Verbatim}[commandchars=!@|,fontsize=\small,frame=single,label=Example]
  !gapprompt@gap>| !gapinput@V:= VectorSpace( Rationals, [ [ 1, 2, 7 ], [ 1/2, 1/3, 5 ] ] );;|
  !gapprompt@gap>| !gapinput@B:= Basis( V, [ [ 1, 2, 7 ], [ 0, 1, -9/4 ] ] );;|
  !gapprompt@gap>| !gapinput@LinearCombination( B, [ 1/2, -2/3 ] );|
  [ 1/2, 1/3, 5 ]
\end{Verbatim}
 }

 

\subsection{\textcolor{Chapter }{EnumeratorByBasis}}
\logpage{[ 61, 6, 5 ]}\nobreak
\hyperdef{L}{X7EB0D16A7EC2DEE3}{}
{\noindent\textcolor{FuncColor}{$\triangleright$\enspace\texttt{EnumeratorByBasis({\mdseries\slshape B})\index{EnumeratorByBasis@\texttt{EnumeratorByBasis}}
\label{EnumeratorByBasis}
}\hfill{\scriptsize (attribute)}}\\


 For a basis \mbox{\texttt{\mdseries\slshape B}} of the free left $F$\texttt{\symbol{45}}module $V$ of dimension $n$, \texttt{EnumeratorByBasis} returns an enumerator that loops over the elements of $V$ as linear combinations of the vectors of \mbox{\texttt{\mdseries\slshape B}} with coefficients the row vectors in the full row space (see{\nobreakspace}\texttt{FullRowSpace} (\ref{FullRowSpace})) of dimension $n$ over $F$, in the succession given by the default enumerator of this row space. 
\begin{Verbatim}[commandchars=!@|,fontsize=\small,frame=single,label=Example]
  !gapprompt@gap>| !gapinput@V:= GF(2)^3;;|
  !gapprompt@gap>| !gapinput@enum:= EnumeratorByBasis( CanonicalBasis( V ) );;|
  !gapprompt@gap>| !gapinput@Print( enum{ [ 1 .. 4 ] }, "\n" );|
  [ [ 0*Z(2), 0*Z(2), 0*Z(2) ], [ 0*Z(2), 0*Z(2), Z(2)^0 ],
    [ 0*Z(2), Z(2)^0, 0*Z(2) ], [ 0*Z(2), Z(2)^0, Z(2)^0 ] ]
  !gapprompt@gap>| !gapinput@B:= Basis( V, [ [ 1, 1, 1 ], [ 1, 1, 0 ], [ 1, 0, 0 ] ] * Z(2) );;|
  !gapprompt@gap>| !gapinput@enum:= EnumeratorByBasis( B );;|
  !gapprompt@gap>| !gapinput@Print( enum{ [ 1 .. 4 ] }, "\n" );|
  [ [ 0*Z(2), 0*Z(2), 0*Z(2) ], [ Z(2)^0, 0*Z(2), 0*Z(2) ],
    [ Z(2)^0, Z(2)^0, 0*Z(2) ], [ 0*Z(2), Z(2)^0, 0*Z(2) ] ]
\end{Verbatim}
 }

 

\subsection{\textcolor{Chapter }{IteratorByBasis}}
\logpage{[ 61, 6, 6 ]}\nobreak
\hyperdef{L}{X855625D47979005D}{}
{\noindent\textcolor{FuncColor}{$\triangleright$\enspace\texttt{IteratorByBasis({\mdseries\slshape B})\index{IteratorByBasis@\texttt{IteratorByBasis}}
\label{IteratorByBasis}
}\hfill{\scriptsize (operation)}}\\


 For a basis \mbox{\texttt{\mdseries\slshape B}} of the free left $F$\texttt{\symbol{45}}module $V$ of dimension $n$, \texttt{IteratorByBasis} returns an iterator that loops over the elements of $V$ as linear combinations of the vectors of \mbox{\texttt{\mdseries\slshape B}} with coefficients the row vectors in the full row space (see{\nobreakspace}\texttt{FullRowSpace} (\ref{FullRowSpace})) of dimension $n$ over $F$, in the succession given by the default enumerator of this row space. 
\begin{Verbatim}[commandchars=!@|,fontsize=\small,frame=single,label=Example]
  !gapprompt@gap>| !gapinput@V:= GF(2)^3;;|
  !gapprompt@gap>| !gapinput@iter:= IteratorByBasis( CanonicalBasis( V ) );;|
  !gapprompt@gap>| !gapinput@for i in [ 1 .. 4 ] do Print( NextIterator( iter ), "\n" ); od;|
  [ 0*Z(2), 0*Z(2), 0*Z(2) ]
  [ 0*Z(2), 0*Z(2), Z(2)^0 ]
  [ 0*Z(2), Z(2)^0, 0*Z(2) ]
  [ 0*Z(2), Z(2)^0, Z(2)^0 ]
  !gapprompt@gap>| !gapinput@B:= Basis( V, [ [ 1, 1, 1 ], [ 1, 1, 0 ], [ 1, 0, 0 ] ] * Z(2) );;|
  !gapprompt@gap>| !gapinput@iter:= IteratorByBasis( B );;|
  !gapprompt@gap>| !gapinput@for i in [ 1 .. 4 ] do Print( NextIterator( iter ), "\n" ); od;|
  [ 0*Z(2), 0*Z(2), 0*Z(2) ]
  [ Z(2)^0, 0*Z(2), 0*Z(2) ]
  [ Z(2)^0, Z(2)^0, 0*Z(2) ]
  [ 0*Z(2), Z(2)^0, 0*Z(2) ]
\end{Verbatim}
 }

 }

  
\section{\textcolor{Chapter }{Operations for Special Kinds of Bases}}\label{Operations for Special Kinds of Bases}
\logpage{[ 61, 7, 0 ]}
\hyperdef{L}{X82809D6C82DE4EC2}{}
{
  

\subsection{\textcolor{Chapter }{IsCanonicalBasis}}
\logpage{[ 61, 7, 1 ]}\nobreak
\hyperdef{L}{X7CC2B3DD81628CE9}{}
{\noindent\textcolor{FuncColor}{$\triangleright$\enspace\texttt{IsCanonicalBasis({\mdseries\slshape B})\index{IsCanonicalBasis@\texttt{IsCanonicalBasis}}
\label{IsCanonicalBasis}
}\hfill{\scriptsize (property)}}\\


 If the underlying free left module $V$ of the basis \mbox{\texttt{\mdseries\slshape B}} supports a canonical basis (see{\nobreakspace}\texttt{CanonicalBasis} (\ref{CanonicalBasis})) then \texttt{IsCanonicalBasis} returns \texttt{true} if \mbox{\texttt{\mdseries\slshape B}} is equal to the canonical basis of $V$, and \texttt{false} otherwise. }

 

\subsection{\textcolor{Chapter }{IsIntegralBasis}}
\logpage{[ 61, 7, 2 ]}\nobreak
\hyperdef{L}{X86DE147F8606B739}{}
{\noindent\textcolor{FuncColor}{$\triangleright$\enspace\texttt{IsIntegralBasis({\mdseries\slshape B})\index{IsIntegralBasis@\texttt{IsIntegralBasis}}
\label{IsIntegralBasis}
}\hfill{\scriptsize (property)}}\\


 Let \mbox{\texttt{\mdseries\slshape B}} be an $S$\texttt{\symbol{45}}basis of a \emph{field} $F$ for a subfield $S$ of $F$, and let $R$ and $M$ be the rings of algebraic integers in $S$ and $F$, respectively. \texttt{IsIntegralBasis} returns \texttt{true} if \mbox{\texttt{\mdseries\slshape B}} is also an $R$\texttt{\symbol{45}}basis of $M$, and \texttt{false} otherwise. }

 

\subsection{\textcolor{Chapter }{IsNormalBasis}}
\logpage{[ 61, 7, 3 ]}\nobreak
\hyperdef{L}{X7FC051C579D61223}{}
{\noindent\textcolor{FuncColor}{$\triangleright$\enspace\texttt{IsNormalBasis({\mdseries\slshape B})\index{IsNormalBasis@\texttt{IsNormalBasis}}
\label{IsNormalBasis}
}\hfill{\scriptsize (property)}}\\


 Let \mbox{\texttt{\mdseries\slshape B}} be an $S$\texttt{\symbol{45}}basis of a \emph{field} $F$ for a subfield $S$ of $F$. \texttt{IsNormalBasis} returns \texttt{true} if \mbox{\texttt{\mdseries\slshape B}} is invariant under the Galois group (see{\nobreakspace}\texttt{GaloisGroup} (\ref{GaloisGroup:of field})) of the field extension $F / S$, and \texttt{false} otherwise. 
\begin{Verbatim}[commandchars=!@|,fontsize=\small,frame=single,label=Example]
  !gapprompt@gap>| !gapinput@B:= CanonicalBasis( GaussianRationals );|
  CanonicalBasis( GaussianRationals )
  !gapprompt@gap>| !gapinput@IsIntegralBasis( B );  IsNormalBasis( B );|
  true
  false
\end{Verbatim}
 }

    }

  
\section{\textcolor{Chapter }{Mutable Bases}}\label{Mutable Bases}
\logpage{[ 61, 8, 0 ]}
\hyperdef{L}{X7C11B9C3819F3EA2}{}
{
  It is useful to have a \emph{mutable basis} of a free module when successively closures with new vectors are formed, since
one does not want to create a new module and a corresponding basis for each
step. 

 Note that the situation here is different from the situation with stabilizer
chains, which are (mutable or immutable) records that do not need to know
about the groups they describe, whereas each (immutable) basis stores the
underlying left module (see{\nobreakspace}\texttt{UnderlyingLeftModule} (\ref{UnderlyingLeftModule})). 

 So immutable bases and mutable bases are different categories of objects. The
only thing they have in common is that one can ask both for their basis
vectors and for the coefficients of a given vector. 

 Since \texttt{Immutable} produces an immutable copy of any \textsf{GAP} object, it would in principle be possible to construct a mutable basis that is
in fact immutable. In the sequel, we will deal only with mutable bases that
are in fact \emph{mutable} \textsf{GAP} objects, hence these objects are unable to store attribute values. 

 Basic operations for immutable bases are \texttt{NrBasisVectors} (\ref{NrBasisVectors}), \texttt{IsContainedInSpan} (\ref{IsContainedInSpan}), \texttt{CloseMutableBasis} (\ref{CloseMutableBasis}), \texttt{ImmutableBasis} (\ref{ImmutableBasis}), \texttt{Coefficients} (\ref{Coefficients}), and \texttt{BasisVectors} (\ref{BasisVectors}). \texttt{ShallowCopy} (\ref{ShallowCopy}) for a mutable basis returns a mutable plain list containing the current basis
vectors.    

 Since mutable bases do not admit arbitrary changes of their lists of basis
vectors, a mutable basis is \emph{not} a list. It is, however, a collection, more precisely its family
(see{\nobreakspace}\ref{Families}) equals the family of its collection of basis vectors. 

 Mutable bases can be constructed with \texttt{MutableBasis}. 

 Similar to the situation with bases (cf.{\nobreakspace}\ref{Bases of Vector Spaces}), \textsf{GAP} supports the following three kinds of mutable bases. 

 The \emph{generic method} of \texttt{MutableBasis} returns a mutable basis that simply stores an immutable basis; clearly one
wants to avoid this whenever possible with reasonable effort. 

 There are mutable bases that store a mutable basis for a nicer module.     Note that this is meaningful only if the mechanism of computing nice and ugly
vectors (see{\nobreakspace}\ref{Vector Spaces Handled By Nice Bases}) is invariant under closures of the basis; this is the case for example if
the vectors are matrices, Lie objects, or elements of structure constants
algebras. 

 There are mutable bases that use special information to perform their tasks;
examples are mutable bases of Gaussian row and matrix spaces. 

\subsection{\textcolor{Chapter }{IsMutableBasis}}
\logpage{[ 61, 8, 1 ]}\nobreak
\hyperdef{L}{X7F466FB47F7E9F00}{}
{\noindent\textcolor{FuncColor}{$\triangleright$\enspace\texttt{IsMutableBasis({\mdseries\slshape MB})\index{IsMutableBasis@\texttt{IsMutableBasis}}
\label{IsMutableBasis}
}\hfill{\scriptsize (Category)}}\\


 Every mutable basis lies in the category \texttt{IsMutableBasis}. }

 

\subsection{\textcolor{Chapter }{MutableBasis}}
\logpage{[ 61, 8, 2 ]}\nobreak
\hyperdef{L}{X8115C061819E5172}{}
{\noindent\textcolor{FuncColor}{$\triangleright$\enspace\texttt{MutableBasis({\mdseries\slshape R, vectors[, zero]})\index{MutableBasis@\texttt{MutableBasis}}
\label{MutableBasis}
}\hfill{\scriptsize (operation)}}\\


 \texttt{MutableBasis} returns a mutable basis for the \mbox{\texttt{\mdseries\slshape R}}\texttt{\symbol{45}}free module generated by the vectors in the list \mbox{\texttt{\mdseries\slshape vectors}}. The optional argument \mbox{\texttt{\mdseries\slshape zero}} is the zero vector of the module; it must be given if \mbox{\texttt{\mdseries\slshape vectors}} is empty. 

 \emph{Note} that \mbox{\texttt{\mdseries\slshape vectors}} will in general \emph{not} be the basis vectors of the mutable basis!  
\begin{Verbatim}[commandchars=!@|,fontsize=\small,frame=single,label=Example]
  !gapprompt@gap>| !gapinput@MB:= MutableBasis( Rationals, [ [ 1, 2, 3 ], [ 0, 1, 0 ] ] );|
  <mutable basis over Rationals, 2 vectors>
\end{Verbatim}
 }

 

\subsection{\textcolor{Chapter }{NrBasisVectors}}
\logpage{[ 61, 8, 3 ]}\nobreak
\hyperdef{L}{X7EC90F4F7BCAF8D4}{}
{\noindent\textcolor{FuncColor}{$\triangleright$\enspace\texttt{NrBasisVectors({\mdseries\slshape MB})\index{NrBasisVectors@\texttt{NrBasisVectors}}
\label{NrBasisVectors}
}\hfill{\scriptsize (operation)}}\\


 For a mutable basis \mbox{\texttt{\mdseries\slshape MB}}, \texttt{NrBasisVectors} returns the current number of basis vectors of \mbox{\texttt{\mdseries\slshape MB}}. Note that this operation is \emph{not} an attribute, as it makes no sense to store the value. \texttt{NrBasisVectors} is used mainly as an equivalent of \texttt{Dimension} for the underlying left module in the case of immutable bases. 
\begin{Verbatim}[commandchars=!@|,fontsize=\small,frame=single,label=Example]
  !gapprompt@gap>| !gapinput@MB:= MutableBasis( Rationals, [ [ 1, 1], [ 2, 2 ] ] );;|
  !gapprompt@gap>| !gapinput@NrBasisVectors( MB );|
  1
\end{Verbatim}
 }

 

\subsection{\textcolor{Chapter }{ImmutableBasis}}
\logpage{[ 61, 8, 4 ]}\nobreak
\hyperdef{L}{X7BA87512823A8CFD}{}
{\noindent\textcolor{FuncColor}{$\triangleright$\enspace\texttt{ImmutableBasis({\mdseries\slshape MB[, V]})\index{ImmutableBasis@\texttt{ImmutableBasis}}
\label{ImmutableBasis}
}\hfill{\scriptsize (operation)}}\\


 \texttt{ImmutableBasis} returns the immutable basis $B$ with the same basis vectors as in the mutable basis \mbox{\texttt{\mdseries\slshape MB}}. 

 If the second argument \mbox{\texttt{\mdseries\slshape V}} is present then \mbox{\texttt{\mdseries\slshape V}} is the value of \texttt{UnderlyingLeftModule} (\ref{UnderlyingLeftModule}) for $B$. The second variant is used mainly for the case that one knows the module for
the desired basis in advance, and if it has a nicer structure than the module
known to \mbox{\texttt{\mdseries\slshape MB}}, for example if it is an algebra.    
\begin{Verbatim}[commandchars=!@|,fontsize=\small,frame=single,label=Example]
  !gapprompt@gap>| !gapinput@MB:= MutableBasis( Rationals, [ [ 1, 1 ], [ 2, 2 ] ] );;|
  !gapprompt@gap>| !gapinput@B:= ImmutableBasis( MB );|
  SemiEchelonBasis( <vector space of dimension 1 over Rationals>,
  [ [ 1, 1 ] ] )
  !gapprompt@gap>| !gapinput@UnderlyingLeftModule( B );|
  <vector space of dimension 1 over Rationals>
\end{Verbatim}
 }

 

\subsection{\textcolor{Chapter }{IsContainedInSpan}}
\logpage{[ 61, 8, 5 ]}\nobreak
\hyperdef{L}{X85B50AC77A22108B}{}
{\noindent\textcolor{FuncColor}{$\triangleright$\enspace\texttt{IsContainedInSpan({\mdseries\slshape MB, v})\index{IsContainedInSpan@\texttt{IsContainedInSpan}}
\label{IsContainedInSpan}
}\hfill{\scriptsize (operation)}}\\


 For a mutable basis \mbox{\texttt{\mdseries\slshape MB}} over the coefficient ring $R$ and a vector \mbox{\texttt{\mdseries\slshape v}}, \texttt{IsContainedInSpan} returns \texttt{true} is \mbox{\texttt{\mdseries\slshape v}} lies in the $R$\texttt{\symbol{45}}span of the current basis vectors of \mbox{\texttt{\mdseries\slshape MB}}, and \texttt{false} otherwise. }

 

\subsection{\textcolor{Chapter }{CloseMutableBasis}}
\logpage{[ 61, 8, 6 ]}\nobreak
\hyperdef{L}{X7B52C99B84316F61}{}
{\noindent\textcolor{FuncColor}{$\triangleright$\enspace\texttt{CloseMutableBasis({\mdseries\slshape MB, v})\index{CloseMutableBasis@\texttt{CloseMutableBasis}}
\label{CloseMutableBasis}
}\hfill{\scriptsize (operation)}}\\


 For a mutable basis \mbox{\texttt{\mdseries\slshape MB}} over the coefficient ring $R$ and a vector \mbox{\texttt{\mdseries\slshape v}}, \texttt{CloseMutableBasis} changes \mbox{\texttt{\mdseries\slshape MB}} such that afterwards it describes the $R$\texttt{\symbol{45}}span of the former basis vectors together with \mbox{\texttt{\mdseries\slshape v}}. 

 \emph{Note} that if \mbox{\texttt{\mdseries\slshape v}} enlarges the dimension then this does in general \emph{not} mean that \mbox{\texttt{\mdseries\slshape v}} is simply added to the basis vectors of \mbox{\texttt{\mdseries\slshape MB}}. Usually a linear combination of \mbox{\texttt{\mdseries\slshape v}} and the other basis vectors is added, and also the old basis vectors may be
modified, for example in order to keep the list of basis vectors echelonized
(see{\nobreakspace}\texttt{IsSemiEchelonized} (\ref{IsSemiEchelonized})). 

 \texttt{CloseMutableBasis} returns \texttt{false} if \mbox{\texttt{\mdseries\slshape v}} was already in the $R$\texttt{\symbol{45}}span described by \mbox{\texttt{\mdseries\slshape MB}}, and \texttt{true} if \mbox{\texttt{\mdseries\slshape MB}} got extended. 
\begin{Verbatim}[commandchars=!@|,fontsize=\small,frame=single,label=Example]
  !gapprompt@gap>| !gapinput@MB:= MutableBasis( Rationals, [ [ 1, 1, 3 ], [ 2, 2, 1 ] ] );|
  <mutable basis over Rationals, 2 vectors>
  !gapprompt@gap>| !gapinput@IsContainedInSpan( MB, [ 1, 0, 0 ] );|
  false
  !gapprompt@gap>| !gapinput@CloseMutableBasis( MB, [ 1, 0, 0 ] );|
  true
  !gapprompt@gap>| !gapinput@MB;|
  <mutable basis over Rationals, 3 vectors>
  !gapprompt@gap>| !gapinput@IsContainedInSpan( MB, [ 1, 0, 0 ] );|
  true
  !gapprompt@gap>| !gapinput@CloseMutableBasis( MB, [ 1, 0, 0 ] );|
  false
\end{Verbatim}
 }

 }

  
\section{\textcolor{Chapter }{Row and Matrix Spaces}}\label{Row and Matrix Spaces}
\logpage{[ 61, 9, 0 ]}
\hyperdef{L}{X7D937EBC7DE2819B}{}
{
  \index{row spaces} \index{matrix spaces} 

\subsection{\textcolor{Chapter }{IsRowSpace}}
\logpage{[ 61, 9, 1 ]}\nobreak
\hyperdef{L}{X79B305CE87511C4B}{}
{\noindent\textcolor{FuncColor}{$\triangleright$\enspace\texttt{IsRowSpace({\mdseries\slshape V})\index{IsRowSpace@\texttt{IsRowSpace}}
\label{IsRowSpace}
}\hfill{\scriptsize (filter)}}\\


 A \emph{row space} in \textsf{GAP} is a vector space that consists of row vectors (see Chapter{\nobreakspace}\ref{Row Vectors}). }

 

\subsection{\textcolor{Chapter }{IsMatrixSpace}}
\logpage{[ 61, 9, 2 ]}\nobreak
\hyperdef{L}{X7A2BBBA07B2BE8F8}{}
{\noindent\textcolor{FuncColor}{$\triangleright$\enspace\texttt{IsMatrixSpace({\mdseries\slshape V})\index{IsMatrixSpace@\texttt{IsMatrixSpace}}
\label{IsMatrixSpace}
}\hfill{\scriptsize (filter)}}\\


 A \emph{matrix space} in \textsf{GAP} is a vector space that consists of matrices (see Chapter{\nobreakspace}\ref{Matrices}). }

 

\subsection{\textcolor{Chapter }{IsGaussianSpace}}
\logpage{[ 61, 9, 3 ]}\nobreak
\hyperdef{L}{X83724C157F4FDFB4}{}
{\noindent\textcolor{FuncColor}{$\triangleright$\enspace\texttt{IsGaussianSpace({\mdseries\slshape V})\index{IsGaussianSpace@\texttt{IsGaussianSpace}}
\label{IsGaussianSpace}
}\hfill{\scriptsize (filter)}}\\


 The filter \texttt{IsGaussianSpace} (see{\nobreakspace}\ref{Filters}) for the row space (see{\nobreakspace}\texttt{IsRowSpace} (\ref{IsRowSpace})) or matrix space (see{\nobreakspace}\texttt{IsMatrixSpace} (\ref{IsMatrixSpace})) \mbox{\texttt{\mdseries\slshape V}} over a field $F$ indicates that the entries of all row vectors or matrices in \mbox{\texttt{\mdseries\slshape V}}, respectively, are all contained in $F$. In this case, \mbox{\texttt{\mdseries\slshape V}} is called a \emph{Gaussian} vector space. Bases for Gaussian spaces can be computed using Gaussian
elimination for a given list of vector space generators. 
\begin{Verbatim}[commandchars=!@|,fontsize=\small,frame=single,label=Example]
  !gapprompt@gap>| !gapinput@mats:= [ [[1,1],[2,2]], [[3,4],[0,1]] ];;|
  !gapprompt@gap>| !gapinput@V:= VectorSpace( Rationals, mats );;|
  !gapprompt@gap>| !gapinput@IsGaussianSpace( V );|
  true
  !gapprompt@gap>| !gapinput@mats[1][1][1]:= E(4);;   # an element in an extension field|
  !gapprompt@gap>| !gapinput@V:= VectorSpace( Rationals, mats );;|
  !gapprompt@gap>| !gapinput@IsGaussianSpace( V );|
  false
  !gapprompt@gap>| !gapinput@V:= VectorSpace( Field( Rationals, [ E(4) ] ), mats );;|
  !gapprompt@gap>| !gapinput@IsGaussianSpace( V );|
  true
\end{Verbatim}
 }

 

\subsection{\textcolor{Chapter }{FullRowSpace}}
\logpage{[ 61, 9, 4 ]}\nobreak
\hyperdef{L}{X80209A8785126AAB}{}
{\noindent\textcolor{FuncColor}{$\triangleright$\enspace\texttt{FullRowSpace({\mdseries\slshape F, n})\index{FullRowSpace@\texttt{FullRowSpace}}
\label{FullRowSpace}
}\hfill{\scriptsize (function)}}\\
\noindent\textcolor{FuncColor}{$\triangleright$\enspace\texttt{\texttt{\symbol{92}}\texttt{\symbol{94}}({\mdseries\slshape F, n})\index{\texttt{\symbol{92}}\texttt{\symbol{94}}@\texttt{\texttt{\symbol{92}}\texttt{\symbol{94}}}!for a field and an integer}
\label{bSlash^:for a field and an integer}
}\hfill{\scriptsize (method)}}\\


 For a field \mbox{\texttt{\mdseries\slshape F}} and a nonnegative integer \mbox{\texttt{\mdseries\slshape n}}, \texttt{FullRowSpace} returns the \mbox{\texttt{\mdseries\slshape F}}\texttt{\symbol{45}}vector space that consists of all row vectors
(see{\nobreakspace}\texttt{IsRowVector} (\ref{IsRowVector})) of length \mbox{\texttt{\mdseries\slshape n}} with entries in \mbox{\texttt{\mdseries\slshape F}}. 

 An alternative to construct this vector space is via \mbox{\texttt{\mdseries\slshape F}}\texttt{\texttt{\symbol{94}}}\mbox{\texttt{\mdseries\slshape n}}. 
\begin{Verbatim}[commandchars=!@|,fontsize=\small,frame=single,label=Example]
  !gapprompt@gap>| !gapinput@FullRowSpace( GF( 9 ), 3 );|
  ( GF(3^2)^3 )
  !gapprompt@gap>| !gapinput@GF(9)^3;           # the same as above|
  ( GF(3^2)^3 )
\end{Verbatim}
 }

 

\subsection{\textcolor{Chapter }{FullMatrixSpace}}
\logpage{[ 61, 9, 5 ]}\nobreak
\hyperdef{L}{X876B66C37A7B749F}{}
{\noindent\textcolor{FuncColor}{$\triangleright$\enspace\texttt{FullMatrixSpace({\mdseries\slshape F, m, n})\index{FullMatrixSpace@\texttt{FullMatrixSpace}}
\label{FullMatrixSpace}
}\hfill{\scriptsize (function)}}\\
\noindent\textcolor{FuncColor}{$\triangleright$\enspace\texttt{\texttt{\symbol{92}}\texttt{\symbol{94}}({\mdseries\slshape F, dims})\index{\texttt{\symbol{92}}\texttt{\symbol{94}}@\texttt{\texttt{\symbol{92}}\texttt{\symbol{94}}}!for a field and a pair of integers}
\label{bSlash^:for a field and a pair of integers}
}\hfill{\scriptsize (method)}}\\


 For a field \mbox{\texttt{\mdseries\slshape F}} and two positive integers \mbox{\texttt{\mdseries\slshape m}} and \mbox{\texttt{\mdseries\slshape n}}, \texttt{FullMatrixSpace} returns the \mbox{\texttt{\mdseries\slshape F}}\texttt{\symbol{45}}vector space that consists of all \mbox{\texttt{\mdseries\slshape m}} by \mbox{\texttt{\mdseries\slshape n}} matrices (see{\nobreakspace}\texttt{IsMatrix} (\ref{IsMatrix})) with entries in \mbox{\texttt{\mdseries\slshape F}}. 

 If \mbox{\texttt{\mdseries\slshape m}}\texttt{ = }\mbox{\texttt{\mdseries\slshape n}} then the result is in fact an algebra (see{\nobreakspace}\texttt{FullMatrixAlgebra} (\ref{FullMatrixAlgebra})). 

 An alternative to construct this vector space is via \mbox{\texttt{\mdseries\slshape F}}\texttt{\texttt{\symbol{94}}[}\mbox{\texttt{\mdseries\slshape m}},\mbox{\texttt{\mdseries\slshape n}}\texttt{]}. 
\begin{Verbatim}[commandchars=!@|,fontsize=\small,frame=single,label=Example]
  !gapprompt@gap>| !gapinput@FullMatrixSpace( GF(2), 4, 5 );|
  ( GF(2)^[ 4, 5 ] )
  !gapprompt@gap>| !gapinput@GF(2)^[ 4, 5 ];    # the same as above|
  ( GF(2)^[ 4, 5 ] )
\end{Verbatim}
 }

 

\subsection{\textcolor{Chapter }{DimensionOfVectors}}
\logpage{[ 61, 9, 6 ]}\nobreak
\hyperdef{L}{X8534A750878478D0}{}
{\noindent\textcolor{FuncColor}{$\triangleright$\enspace\texttt{DimensionOfVectors({\mdseries\slshape M})\index{DimensionOfVectors@\texttt{DimensionOfVectors}}
\label{DimensionOfVectors}
}\hfill{\scriptsize (attribute)}}\\


 For a left module \mbox{\texttt{\mdseries\slshape M}} that consists of row vectors (see{\nobreakspace}\texttt{IsRowModule} (\ref{IsRowModule})), \texttt{DimensionOfVectors} returns the common length of all row vectors in \mbox{\texttt{\mdseries\slshape M}}. For a left module \mbox{\texttt{\mdseries\slshape M}} that consists of matrices (see{\nobreakspace}\texttt{IsMatrixModule} (\ref{IsMatrixModule})), \texttt{DimensionOfVectors} returns the common matrix dimensions (see{\nobreakspace}\texttt{DimensionsMat} (\ref{DimensionsMat})) of all matrices in \mbox{\texttt{\mdseries\slshape M}}. 
\begin{Verbatim}[commandchars=!@|,fontsize=\small,frame=single,label=Example]
  !gapprompt@gap>| !gapinput@DimensionOfVectors( GF(2)^5 );|
  5
  !gapprompt@gap>| !gapinput@DimensionOfVectors( GF(2)^[2,3] );|
  [ 2, 3 ]
\end{Verbatim}
 }

 

\subsection{\textcolor{Chapter }{IsSemiEchelonized}}
\logpage{[ 61, 9, 7 ]}\nobreak
\hyperdef{L}{X865A540F85FAE2DF}{}
{\noindent\textcolor{FuncColor}{$\triangleright$\enspace\texttt{IsSemiEchelonized({\mdseries\slshape B})\index{IsSemiEchelonized@\texttt{IsSemiEchelonized}}
\label{IsSemiEchelonized}
}\hfill{\scriptsize (property)}}\\


 Let \mbox{\texttt{\mdseries\slshape B}} be a basis of a Gaussian row or matrix space $V$ (see{\nobreakspace}\texttt{IsGaussianSpace} (\ref{IsGaussianSpace})) over the field $F$. 

 If $V$ is a row space then \mbox{\texttt{\mdseries\slshape B}} is semi\texttt{\symbol{45}}echelonized if the matrix formed by its basis
vectors has the property that the first nonzero element in each row is the
identity of $F$, and all values exactly below these pivot elements are the zero of $F$ (cf.{\nobreakspace}\texttt{SemiEchelonMat} (\ref{SemiEchelonMat})). 

 If $V$ is a matrix space then \mbox{\texttt{\mdseries\slshape B}} is semi\texttt{\symbol{45}}echelonized if the matrix obtained by replacing
each basis vector by the concatenation of its rows is
semi\texttt{\symbol{45}}echelonized (see above, cf.{\nobreakspace}\texttt{SemiEchelonMats} (\ref{SemiEchelonMats})). 
\begin{Verbatim}[commandchars=!@|,fontsize=\small,frame=single,label=Example]
  !gapprompt@gap>| !gapinput@V:= GF(2)^2;;|
  !gapprompt@gap>| !gapinput@B1:= Basis( V, [ [ 0, 1 ], [ 1, 0 ] ] * Z(2) );;|
  !gapprompt@gap>| !gapinput@IsSemiEchelonized( B1 );|
  true
  !gapprompt@gap>| !gapinput@B2:= Basis( V, [ [ 0, 1 ], [ 1, 1 ] ] * Z(2) );;|
  !gapprompt@gap>| !gapinput@IsSemiEchelonized( B2 );|
  false
\end{Verbatim}
 }

 

\subsection{\textcolor{Chapter }{SemiEchelonBasis}}
\logpage{[ 61, 9, 8 ]}\nobreak
\hyperdef{L}{X87DCA09579589106}{}
{\noindent\textcolor{FuncColor}{$\triangleright$\enspace\texttt{SemiEchelonBasis({\mdseries\slshape V[, vectors]})\index{SemiEchelonBasis@\texttt{SemiEchelonBasis}}
\label{SemiEchelonBasis}
}\hfill{\scriptsize (attribute)}}\\
\noindent\textcolor{FuncColor}{$\triangleright$\enspace\texttt{SemiEchelonBasisNC({\mdseries\slshape V, vectors})\index{SemiEchelonBasisNC@\texttt{SemiEchelonBasisNC}}
\label{SemiEchelonBasisNC}
}\hfill{\scriptsize (operation)}}\\


 Let \mbox{\texttt{\mdseries\slshape V}} be a Gaussian row or matrix vector space over the field $F$ (see{\nobreakspace}\texttt{IsGaussianSpace} (\ref{IsGaussianSpace}), \texttt{IsRowSpace} (\ref{IsRowSpace}), \texttt{IsMatrixSpace} (\ref{IsMatrixSpace})). 

 Called with \mbox{\texttt{\mdseries\slshape V}} as the only argument, \texttt{SemiEchelonBasis} returns a basis of \mbox{\texttt{\mdseries\slshape V}} that has the property \texttt{IsSemiEchelonized} (\ref{IsSemiEchelonized}). 

 If additionally a list \mbox{\texttt{\mdseries\slshape vectors}} of vectors in \mbox{\texttt{\mdseries\slshape V}} is given that forms a semi\texttt{\symbol{45}}echelonized basis of \mbox{\texttt{\mdseries\slshape V}} then \texttt{SemiEchelonBasis} returns this basis; if \mbox{\texttt{\mdseries\slshape vectors}} do not form a basis of \mbox{\texttt{\mdseries\slshape V}} then \texttt{fail} is returned. 

 \texttt{SemiEchelonBasisNC} does the same as the two argument version of \texttt{SemiEchelonBasis}, except that it is not checked whether \mbox{\texttt{\mdseries\slshape vectors}} form a semi\texttt{\symbol{45}}echelonized basis. 
\begin{Verbatim}[commandchars=!@|,fontsize=\small,frame=single,label=Example]
  !gapprompt@gap>| !gapinput@V:= GF(2)^2;;|
  !gapprompt@gap>| !gapinput@B:= SemiEchelonBasis( V );|
  SemiEchelonBasis( ( GF(2)^2 ), ... )
  !gapprompt@gap>| !gapinput@Print( BasisVectors( B ), "\n" );|
  [ [ Z(2)^0, 0*Z(2) ], [ 0*Z(2), Z(2)^0 ] ]
  !gapprompt@gap>| !gapinput@B:= SemiEchelonBasis( V, [ [ 1, 1 ], [ 0, 1 ] ] * Z(2) );|
  SemiEchelonBasis( ( GF(2)^2 ), <an immutable 2x2 matrix over GF2> )
  !gapprompt@gap>| !gapinput@Print( BasisVectors( B ), "\n" );|
  [ [ Z(2)^0, Z(2)^0 ], [ 0*Z(2), Z(2)^0 ] ]
  !gapprompt@gap>| !gapinput@Coefficients( B, [ 0, 1 ] * Z(2) );|
  [ 0*Z(2), Z(2)^0 ]
  !gapprompt@gap>| !gapinput@Coefficients( B, [ 1, 0 ] * Z(2) );|
  [ Z(2)^0, Z(2)^0 ]
  !gapprompt@gap>| !gapinput@SemiEchelonBasis( V, [ [ 0, 1 ], [ 1, 1 ] ] * Z(2) );|
  fail
\end{Verbatim}
 }

 

\subsection{\textcolor{Chapter }{IsCanonicalBasisFullRowModule}}
\logpage{[ 61, 9, 9 ]}\nobreak
\hyperdef{L}{X7C3CC5F97FA048A4}{}
{\noindent\textcolor{FuncColor}{$\triangleright$\enspace\texttt{IsCanonicalBasisFullRowModule({\mdseries\slshape B})\index{IsCanonicalBasisFullRowModule@\texttt{IsCanonicalBasisFullRowModule}}
\label{IsCanonicalBasisFullRowModule}
}\hfill{\scriptsize (property)}}\\


 \index{canonical basis!for row spaces} \texttt{IsCanonicalBasisFullRowModule} returns \texttt{true} if \mbox{\texttt{\mdseries\slshape B}} is the canonical basis (see{\nobreakspace}\texttt{IsCanonicalBasis} (\ref{IsCanonicalBasis})) of a full row module (see{\nobreakspace}\texttt{IsFullRowModule} (\ref{IsFullRowModule})), and \texttt{false} otherwise. 

 The \emph{canonical basis} of a Gaussian row space is defined as the unique
semi\texttt{\symbol{45}}echelonized (see{\nobreakspace}\texttt{IsSemiEchelonized} (\ref{IsSemiEchelonized})) basis with the additional property that for $j > i$ the position of the pivot of row $j$ is bigger than the position of the pivot of row $i$, and that each pivot column contains exactly one nonzero entry. }

 

\subsection{\textcolor{Chapter }{IsCanonicalBasisFullMatrixModule}}
\logpage{[ 61, 9, 10 ]}\nobreak
\hyperdef{L}{X83D282697C1A3148}{}
{\noindent\textcolor{FuncColor}{$\triangleright$\enspace\texttt{IsCanonicalBasisFullMatrixModule({\mdseries\slshape B})\index{IsCanonicalBasisFullMatrixModule@\texttt{IsCanonicalBasisFullMatrixModule}}
\label{IsCanonicalBasisFullMatrixModule}
}\hfill{\scriptsize (property)}}\\


 \index{canonical basis!for matrix spaces} \texttt{IsCanonicalBasisFullMatrixModule} returns \texttt{true} if \mbox{\texttt{\mdseries\slshape B}} is the canonical basis (see{\nobreakspace}\texttt{IsCanonicalBasis} (\ref{IsCanonicalBasis})) of a full matrix module (see{\nobreakspace}\texttt{IsFullMatrixModule} (\ref{IsFullMatrixModule})), and \texttt{false} otherwise. 

 The \emph{canonical basis} of a Gaussian matrix space is defined as the unique
semi\texttt{\symbol{45}}echelonized (see{\nobreakspace}\texttt{IsSemiEchelonized} (\ref{IsSemiEchelonized})) basis for which the list of concatenations of the basis vectors forms the
canonical basis of the corresponding Gaussian row space. }

 

\subsection{\textcolor{Chapter }{NormedRowVectors}}
\logpage{[ 61, 9, 11 ]}\nobreak
\hyperdef{L}{X7D6537F87E940344}{}
{\noindent\textcolor{FuncColor}{$\triangleright$\enspace\texttt{NormedRowVectors({\mdseries\slshape V})\index{NormedRowVectors@\texttt{NormedRowVectors}}
\label{NormedRowVectors}
}\hfill{\scriptsize (attribute)}}\\


 For a finite Gaussian row space \mbox{\texttt{\mdseries\slshape V}} (see{\nobreakspace}\texttt{IsRowSpace} (\ref{IsRowSpace}), \texttt{IsGaussianSpace} (\ref{IsGaussianSpace})), \texttt{NormedRowVectors} returns a list of those nonzero vectors in \mbox{\texttt{\mdseries\slshape V}} that have a one in the first nonzero component. 

 The result list can be used as action domain for the action of a matrix group
via \texttt{OnLines} (\ref{OnLines}), which yields the natural action on one\texttt{\symbol{45}}dimensional
subspaces of \mbox{\texttt{\mdseries\slshape V}} (see also{\nobreakspace}\texttt{Subspaces} (\ref{Subspaces})). 
\begin{Verbatim}[commandchars=!@|,fontsize=\small,frame=single,label=Example]
  !gapprompt@gap>| !gapinput@vecs:= NormedRowVectors( GF(3)^2 );|
  [ [ 0*Z(3), Z(3)^0 ], [ Z(3)^0, 0*Z(3) ], [ Z(3)^0, Z(3)^0 ],
    [ Z(3)^0, Z(3) ] ]
  !gapprompt@gap>| !gapinput@Action( GL(2,3), vecs, OnLines );|
  Group([ (3,4), (1,2,4) ])
\end{Verbatim}
 }

 

\subsection{\textcolor{Chapter }{SiftedVector}}
\logpage{[ 61, 9, 12 ]}\nobreak
\hyperdef{L}{X815C69A57C042C34}{}
{\noindent\textcolor{FuncColor}{$\triangleright$\enspace\texttt{SiftedVector({\mdseries\slshape B, v})\index{SiftedVector@\texttt{SiftedVector}}
\label{SiftedVector}
}\hfill{\scriptsize (operation)}}\\


 Let \mbox{\texttt{\mdseries\slshape B}} be a semi\texttt{\symbol{45}}echelonized basis (see{\nobreakspace}\texttt{IsSemiEchelonized} (\ref{IsSemiEchelonized})) of a Gaussian row or matrix space $V$ (see{\nobreakspace}\texttt{IsGaussianSpace} (\ref{IsGaussianSpace})), and \mbox{\texttt{\mdseries\slshape v}} a row vector or matrix, respectively, of the same dimension as the elements in $V$. \texttt{SiftedVector} returns the \emph{residuum} of \mbox{\texttt{\mdseries\slshape v}} with respect to \mbox{\texttt{\mdseries\slshape B}}, which is obtained by successively cleaning the pivot positions in \mbox{\texttt{\mdseries\slshape v}} by subtracting multiples of the basis vectors in \mbox{\texttt{\mdseries\slshape B}}. So the result is the zero vector in $V$ if and only if \mbox{\texttt{\mdseries\slshape v}} lies in $V$. 

 \mbox{\texttt{\mdseries\slshape B}} may also be a mutable basis (see{\nobreakspace}\ref{Mutable Bases}) of a Gaussian row or matrix space. 
\begin{Verbatim}[commandchars=!@|,fontsize=\small,frame=single,label=Example]
  !gapprompt@gap>| !gapinput@V:= VectorSpace( Rationals, [ [ 1, 2, 7 ], [ 1/2, 1/3, 5 ] ] );;|
  !gapprompt@gap>| !gapinput@B:= Basis( V );;|
  !gapprompt@gap>| !gapinput@SiftedVector( B, [ 1, 2, 8 ] );|
  [ 0, 0, 1 ]
\end{Verbatim}
 }

 }

  
\section{\textcolor{Chapter }{Vector Space Homomorphisms}}\label{Vector Space Homomorphisms}
\logpage{[ 61, 10, 0 ]}
\hyperdef{L}{X7F61CECA84CEF39D}{}
{
  \emph{Vector space homomorphisms} (or \emph{linear mappings}) are defined in Section{\nobreakspace}\ref{Linear Mappings}. \textsf{GAP} provides special functions to construct a particular linear mapping from
images of given elements in the source, from a matrix of coefficients, or as a
natural epimorphism. 

 $F$\texttt{\symbol{45}}linear mappings with same source and same range can be
added, so one can form vector spaces of linear mappings. 

\subsection{\textcolor{Chapter }{LeftModuleGeneralMappingByImages}}
\logpage{[ 61, 10, 1 ]}\nobreak
\hyperdef{L}{X82013D328645E370}{}
{\noindent\textcolor{FuncColor}{$\triangleright$\enspace\texttt{LeftModuleGeneralMappingByImages({\mdseries\slshape V, W, gens, imgs})\index{LeftModuleGeneralMappingByImages@\texttt{LeftModuleGeneralMappingByImages}}
\label{LeftModuleGeneralMappingByImages}
}\hfill{\scriptsize (operation)}}\\


 Let \mbox{\texttt{\mdseries\slshape V}} and \mbox{\texttt{\mdseries\slshape W}} be two left modules over the same left acting domain $R$ and \mbox{\texttt{\mdseries\slshape gens}} and \mbox{\texttt{\mdseries\slshape imgs}} lists (of the same length) of elements in \mbox{\texttt{\mdseries\slshape V}} and \mbox{\texttt{\mdseries\slshape W}}, respectively. \texttt{LeftModuleGeneralMappingByImages} returns the general mapping with source \mbox{\texttt{\mdseries\slshape V}} and range \mbox{\texttt{\mdseries\slshape W}} that is defined by mapping the elements in \mbox{\texttt{\mdseries\slshape gens}} to the corresponding elements in \mbox{\texttt{\mdseries\slshape imgs}}, and taking the $R$\texttt{\symbol{45}}linear closure. 

 \mbox{\texttt{\mdseries\slshape gens}} need not generate \mbox{\texttt{\mdseries\slshape V}} as a left $R$\texttt{\symbol{45}}module, and if the specification does not define a linear
mapping then the result will be multi\texttt{\symbol{45}}valued; hence in
general it is not a mapping (see{\nobreakspace}\texttt{IsMapping} (\ref{IsMapping})). 
\begin{Verbatim}[commandchars=!@|,fontsize=\small,frame=single,label=Example]
  !gapprompt@gap>| !gapinput@V:= Rationals^2;;|
  !gapprompt@gap>| !gapinput@W:= VectorSpace( Rationals, [ [1,2,3], [1,0,1] ] );;|
  !gapprompt@gap>| !gapinput@f:= LeftModuleGeneralMappingByImages( V, W,|
  !gapprompt@>| !gapinput@                               [[1,0],[2,0]], [[1,0,1],[1,0,1] ] );|
  [ [ 1, 0 ], [ 2, 0 ] ] -> [ [ 1, 0, 1 ], [ 1, 0, 1 ] ]
  !gapprompt@gap>| !gapinput@IsMapping( f );|
  false
\end{Verbatim}
 }

 

\subsection{\textcolor{Chapter }{LeftModuleHomomorphismByImages}}
\logpage{[ 61, 10, 2 ]}\nobreak
\hyperdef{L}{X85F5293983E47B5A}{}
{\noindent\textcolor{FuncColor}{$\triangleright$\enspace\texttt{LeftModuleHomomorphismByImages({\mdseries\slshape V, W, gens, imgs})\index{LeftModuleHomomorphismByImages@\texttt{LeftModuleHomomorphismByImages}}
\label{LeftModuleHomomorphismByImages}
}\hfill{\scriptsize (function)}}\\
\noindent\textcolor{FuncColor}{$\triangleright$\enspace\texttt{LeftModuleHomomorphismByImagesNC({\mdseries\slshape V, W, gens, imgs})\index{LeftModuleHomomorphismByImagesNC@\texttt{LeftModuleHomomorphismByImagesNC}}
\label{LeftModuleHomomorphismByImagesNC}
}\hfill{\scriptsize (operation)}}\\


 Let \mbox{\texttt{\mdseries\slshape V}} and \mbox{\texttt{\mdseries\slshape W}} be two left modules over the same left acting domain $R$ and \mbox{\texttt{\mdseries\slshape gens}} and \mbox{\texttt{\mdseries\slshape imgs}} lists (of the same length) of elements in \mbox{\texttt{\mdseries\slshape V}} and \mbox{\texttt{\mdseries\slshape W}}, respectively. \texttt{LeftModuleHomomorphismByImages} returns the left $R$\texttt{\symbol{45}}module homomorphism with source \mbox{\texttt{\mdseries\slshape V}} and range \mbox{\texttt{\mdseries\slshape W}} that is defined by mapping the elements in \mbox{\texttt{\mdseries\slshape gens}} to the corresponding elements in \mbox{\texttt{\mdseries\slshape imgs}}. 

 If \mbox{\texttt{\mdseries\slshape gens}} does not generate \mbox{\texttt{\mdseries\slshape V}} or if the homomorphism does not exist (i.e., if mapping the generators
describes only a multi\texttt{\symbol{45}}valued mapping) then \texttt{fail} is returned. For creating a possibly multi\texttt{\symbol{45}}valued mapping
from \mbox{\texttt{\mdseries\slshape V}} to \mbox{\texttt{\mdseries\slshape W}} that respects addition, multiplication, and scalar multiplication, \texttt{LeftModuleGeneralMappingByImages} (\ref{LeftModuleGeneralMappingByImages}) can be used. 

 \texttt{LeftModuleHomomorphismByImagesNC} does the same as \texttt{LeftModuleHomomorphismByImages}, except that it omits all checks. 
\begin{Verbatim}[commandchars=!@|,fontsize=\small,frame=single,label=Example]
  !gapprompt@gap>| !gapinput@V:=Rationals^2;;|
  !gapprompt@gap>| !gapinput@W:=VectorSpace( Rationals, [ [ 1, 0, 1 ], [ 1, 2, 3 ] ] );;|
  !gapprompt@gap>| !gapinput@f:=LeftModuleHomomorphismByImages( V, W,|
  !gapprompt@>| !gapinput@[ [ 1, 0 ], [ 0, 1 ] ], [ [ 1, 0, 1 ], [ 1, 2, 3 ] ] );|
  [ [ 1, 0 ], [ 0, 1 ] ] -> [ [ 1, 0, 1 ], [ 1, 2, 3 ] ]
  !gapprompt@gap>| !gapinput@Image( f, [1,1] );|
  [ 2, 2, 4 ]
\end{Verbatim}
 }

  

\subsection{\textcolor{Chapter }{LeftModuleHomomorphismByMatrix}}
\logpage{[ 61, 10, 3 ]}\nobreak
\hyperdef{L}{X8477E6C3872A6DBB}{}
{\noindent\textcolor{FuncColor}{$\triangleright$\enspace\texttt{LeftModuleHomomorphismByMatrix({\mdseries\slshape BS, matrix, BR})\index{LeftModuleHomomorphismByMatrix@\texttt{LeftModuleHomomorphismByMatrix}}
\label{LeftModuleHomomorphismByMatrix}
}\hfill{\scriptsize (operation)}}\\


 Let \mbox{\texttt{\mdseries\slshape BS}} and \mbox{\texttt{\mdseries\slshape BR}} be bases of the left $R$\texttt{\symbol{45}}modules $V$ and $W$, respectively. \texttt{LeftModuleHomomorphismByMatrix} returns the $R$\texttt{\symbol{45}}linear mapping from $V$ to $W$ that is defined by the matrix \mbox{\texttt{\mdseries\slshape matrix}}, as follows. The image of the $i$\texttt{\symbol{45}}th basis vector of \mbox{\texttt{\mdseries\slshape BS}} is the linear combination of the basis vectors of \mbox{\texttt{\mdseries\slshape BR}} with coefficients the $i$\texttt{\symbol{45}}th row of \mbox{\texttt{\mdseries\slshape matrix}}. 
\begin{Verbatim}[commandchars=!@|,fontsize=\small,frame=single,label=Example]
  !gapprompt@gap>| !gapinput@V:= Rationals^2;;|
  !gapprompt@gap>| !gapinput@W:= VectorSpace( Rationals, [ [ 1, 0, 1 ], [ 1, 2, 3 ] ] );;|
  !gapprompt@gap>| !gapinput@f:= LeftModuleHomomorphismByMatrix( Basis( V ),|
  !gapprompt@>| !gapinput@[ [ 1, 2 ], [ 3, 1 ] ], Basis( W ) );|
  <linear mapping by matrix, ( Rationals^
  2 ) -> <vector space over Rationals, with 2 generators>>
\end{Verbatim}
 }

  

\subsection{\textcolor{Chapter }{NaturalHomomorphismBySubspace}}
\logpage{[ 61, 10, 4 ]}\nobreak
\hyperdef{L}{X8494AA5D7C3B88AD}{}
{\noindent\textcolor{FuncColor}{$\triangleright$\enspace\texttt{NaturalHomomorphismBySubspace({\mdseries\slshape V, W})\index{NaturalHomomorphismBySubspace@\texttt{NaturalHomomorphismBySubspace}}
\label{NaturalHomomorphismBySubspace}
}\hfill{\scriptsize (operation)}}\\


 For an $R$\texttt{\symbol{45}}vector space \mbox{\texttt{\mdseries\slshape V}} and a subspace \mbox{\texttt{\mdseries\slshape W}} of \mbox{\texttt{\mdseries\slshape V}}, \texttt{NaturalHomomorphismBySubspace} returns the $R$\texttt{\symbol{45}}linear mapping that is the natural projection of \mbox{\texttt{\mdseries\slshape V}} onto the factor space \texttt{\mbox{\texttt{\mdseries\slshape V}} / \mbox{\texttt{\mdseries\slshape W}}}. 
\begin{Verbatim}[commandchars=!@|,fontsize=\small,frame=single,label=Example]
  !gapprompt@gap>| !gapinput@V:= Rationals^3;;|
  !gapprompt@gap>| !gapinput@W:= VectorSpace( Rationals, [ [ 1, 1, 1 ] ] );;|
  !gapprompt@gap>| !gapinput@f:= NaturalHomomorphismBySubspace( V, W );|
  <linear mapping by matrix, ( Rationals^3 ) -> ( Rationals^2 )>
\end{Verbatim}
 }

  

\subsection{\textcolor{Chapter }{Hom}}
\logpage{[ 61, 10, 5 ]}\nobreak
\hyperdef{L}{X80015C78876B4F1E}{}
{\noindent\textcolor{FuncColor}{$\triangleright$\enspace\texttt{Hom({\mdseries\slshape F, V, W})\index{Hom@\texttt{Hom}}
\label{Hom}
}\hfill{\scriptsize (operation)}}\\


 For a field \mbox{\texttt{\mdseries\slshape F}} and two vector spaces \mbox{\texttt{\mdseries\slshape V}} and \mbox{\texttt{\mdseries\slshape W}} that can be regarded as \mbox{\texttt{\mdseries\slshape F}}\texttt{\symbol{45}}modules (see{\nobreakspace}\texttt{AsLeftModule} (\ref{AsLeftModule})), \texttt{Hom} returns the \mbox{\texttt{\mdseries\slshape F}}\texttt{\symbol{45}}vector space of all \mbox{\texttt{\mdseries\slshape F}}\texttt{\symbol{45}}linear mappings from \mbox{\texttt{\mdseries\slshape V}} to \mbox{\texttt{\mdseries\slshape W}}. 
\begin{Verbatim}[commandchars=!@|,fontsize=\small,frame=single,label=Example]
  !gapprompt@gap>| !gapinput@V:= Rationals^2;;|
  !gapprompt@gap>| !gapinput@W:= VectorSpace( Rationals, [ [ 1, 0, 1 ], [ 1, 2, 3 ] ] );;|
  !gapprompt@gap>| !gapinput@H:= Hom( Rationals, V, W );|
  Hom( Rationals, ( Rationals^2 ), <vector space over Rationals, with
  2 generators> )
  !gapprompt@gap>| !gapinput@Dimension( H );|
  4
\end{Verbatim}
 }

 

\subsection{\textcolor{Chapter }{End}}
\logpage{[ 61, 10, 6 ]}\nobreak
\hyperdef{L}{X8680ADD381ECF879}{}
{\noindent\textcolor{FuncColor}{$\triangleright$\enspace\texttt{End({\mdseries\slshape F, V})\index{End@\texttt{End}}
\label{End}
}\hfill{\scriptsize (operation)}}\\


 For a field \mbox{\texttt{\mdseries\slshape F}} and a vector space \mbox{\texttt{\mdseries\slshape V}} that can be regarded as an \mbox{\texttt{\mdseries\slshape F}}\texttt{\symbol{45}}module (see{\nobreakspace}\texttt{AsLeftModule} (\ref{AsLeftModule})), \texttt{End} returns the \mbox{\texttt{\mdseries\slshape F}}\texttt{\symbol{45}}algebra of all \mbox{\texttt{\mdseries\slshape F}}\texttt{\symbol{45}}linear mappings from \mbox{\texttt{\mdseries\slshape V}} to \mbox{\texttt{\mdseries\slshape V}}. 
\begin{Verbatim}[commandchars=!@|,fontsize=\small,frame=single,label=Example]
  !gapprompt@gap>| !gapinput@A:= End( Rationals, Rationals^2 );|
  End( Rationals, ( Rationals^2 ) )
  !gapprompt@gap>| !gapinput@Dimension( A );|
  4
\end{Verbatim}
 }

 

\subsection{\textcolor{Chapter }{IsFullHomModule}}
\logpage{[ 61, 10, 7 ]}\nobreak
\hyperdef{L}{X7A9A08EA79259659}{}
{\noindent\textcolor{FuncColor}{$\triangleright$\enspace\texttt{IsFullHomModule({\mdseries\slshape M})\index{IsFullHomModule@\texttt{IsFullHomModule}}
\label{IsFullHomModule}
}\hfill{\scriptsize (property)}}\\


 A \emph{full hom module} is a module of all $R$\texttt{\symbol{45}}linear mappings between two left $R$\texttt{\symbol{45}}modules. The function \texttt{Hom} (\ref{Hom}) can be used to construct a full hom module. 
\begin{Verbatim}[commandchars=!@|,fontsize=\small,frame=single,label=Example]
  !gapprompt@gap>| !gapinput@V:= Rationals^2;;|
  !gapprompt@gap>| !gapinput@W:= VectorSpace( Rationals, [ [ 1, 0, 1 ], [ 1, 2, 3 ] ] );;|
  !gapprompt@gap>| !gapinput@H:= Hom( Rationals, V, W );;|
  !gapprompt@gap>| !gapinput@IsFullHomModule( H );|
  true
\end{Verbatim}
 }

 

\subsection{\textcolor{Chapter }{IsPseudoCanonicalBasisFullHomModule}}
\logpage{[ 61, 10, 8 ]}\nobreak
\hyperdef{L}{X7C4737687E76A24A}{}
{\noindent\textcolor{FuncColor}{$\triangleright$\enspace\texttt{IsPseudoCanonicalBasisFullHomModule({\mdseries\slshape B})\index{IsPseudoCanonicalBasisFullHomModule@\texttt{IsPseudoCanonicalBasisFullHomModule}}
\label{IsPseudoCanonicalBasisFullHomModule}
}\hfill{\scriptsize (property)}}\\


 A basis of a full hom module is called pseudo canonical basis if the matrices
of its basis vectors w.r.t. the stored bases of source and range contain
exactly one identity entry and otherwise zeros. 

 Note that this is not a canonical basis (see{\nobreakspace}\texttt{CanonicalBasis} (\ref{CanonicalBasis})) because it depends on the stored bases of source and range. 
\begin{Verbatim}[commandchars=!@|,fontsize=\small,frame=single,label=Example]
  !gapprompt@gap>| !gapinput@IsPseudoCanonicalBasisFullHomModule( Basis( H ) );|
  true
\end{Verbatim}
 }

 

\subsection{\textcolor{Chapter }{IsLinearMappingsModule}}
\logpage{[ 61, 10, 9 ]}\nobreak
\hyperdef{L}{X84F87C327A1856F2}{}
{\noindent\textcolor{FuncColor}{$\triangleright$\enspace\texttt{IsLinearMappingsModule({\mdseries\slshape V})\index{IsLinearMappingsModule@\texttt{IsLinearMappingsModule}}
\label{IsLinearMappingsModule}
}\hfill{\scriptsize (filter)}}\\


 If an $F$\texttt{\symbol{45}}vector space \mbox{\texttt{\mdseries\slshape V}} is in the filter \texttt{IsLinearMappingsModule} then this expresses that \mbox{\texttt{\mdseries\slshape V}} consists of linear mappings, and that \mbox{\texttt{\mdseries\slshape V}} is handled via the mechanism of nice bases (see{\nobreakspace}\ref{Vector Spaces Handled By Nice Bases}), in the following way. Let $S$ and $R$ be the source and the range, respectively, of each mapping in $V$. Then the \texttt{NiceFreeLeftModuleInfo} (\ref{NiceFreeLeftModuleInfo}) value of \mbox{\texttt{\mdseries\slshape V}} is a record with the components \texttt{basissource} (a basis $B_S$ of $S$) and \texttt{basisrange} (a basis $B_R$ of $R$), and the \texttt{NiceVector} (\ref{NiceVector}) value of $v \in \mbox{\texttt{\mdseries\slshape V}}$ is defined as the matrix of the $F$\texttt{\symbol{45}}linear mapping $v$ w.r.t.{\nobreakspace}the bases $B_S$ and $B_R$. }

  }

  
\section{\textcolor{Chapter }{Vector Spaces Handled By Nice Bases}}\label{Vector Spaces Handled By Nice Bases}
\logpage{[ 61, 11, 0 ]}
\hyperdef{L}{X81503EB77FCE648D}{}
{
  There are kinds of free $R$\texttt{\symbol{45}}modules for which efficient computations are possible
because the elements are ``nice'', for example subspaces of full row modules or of full matrix modules. In
other cases, a ``nice'' canonical basis is known that allows one to do the necessary computations in
the corresponding row module, for example algebras given by structure
constants. 

 In many other situations, one knows at least an isomorphism from the given
module $V$ to a ``nicer'' free left module $W$, in the sense that for each vector in $V$, the image in $W$ can easily be computed, and analogously for each vector in $W$, one can compute the preimage in $V$. 

 This allows one to delegate computations w.r.t.{\nobreakspace}a basis $B$ of $V$ to the corresponding basis $C$ of $W$. We call $W$ the \emph{nice free left module} of $V$, and $C$ the \emph{nice basis} of $B$. (Note that it may happen that also $C$ delegates questions to a ``nicer'' basis.) The basis $B$ indicates the intended behaviour by the filter \texttt{IsBasisByNiceBasis} (\ref{IsBasisByNiceBasis}), and stores $C$ as value of the attribute \texttt{NiceBasis} (\ref{NiceBasis}). $V$ indicates the intended behaviour by the filter \texttt{IsHandledByNiceBasis} (\ref{IsHandledByNiceBasis}), and stores $W$ as value of the attribute \texttt{NiceFreeLeftModule} (\ref{NiceFreeLeftModule}). 

 The bijection between $V$ and $W$ is implemented by the functions \texttt{NiceVector} (\ref{NiceVector}) and \texttt{UglyVector} (\ref{UglyVector}); additional data needed to compute images and preimages can be stored as
value of \texttt{NiceFreeLeftModuleInfo} (\ref{NiceFreeLeftModuleInfo}). 

\subsection{\textcolor{Chapter }{NiceFreeLeftModule}}
\logpage{[ 61, 11, 1 ]}\nobreak
\hyperdef{L}{X826FD4BC7BA0559D}{}
{\noindent\textcolor{FuncColor}{$\triangleright$\enspace\texttt{NiceFreeLeftModule({\mdseries\slshape V})\index{NiceFreeLeftModule@\texttt{NiceFreeLeftModule}}
\label{NiceFreeLeftModule}
}\hfill{\scriptsize (attribute)}}\\


 For a free left module \mbox{\texttt{\mdseries\slshape V}} that is handled via the mechanism of nice bases, this attribute stores the
associated free left module to which the tasks are delegated. }

 

\subsection{\textcolor{Chapter }{NiceVector}}
\logpage{[ 61, 11, 2 ]}\nobreak
\hyperdef{L}{X807B8032780C59A4}{}
{\noindent\textcolor{FuncColor}{$\triangleright$\enspace\texttt{NiceVector({\mdseries\slshape V, v})\index{NiceVector@\texttt{NiceVector}}
\label{NiceVector}
}\hfill{\scriptsize (operation)}}\\
\noindent\textcolor{FuncColor}{$\triangleright$\enspace\texttt{UglyVector({\mdseries\slshape V, r})\index{UglyVector@\texttt{UglyVector}}
\label{UglyVector}
}\hfill{\scriptsize (operation)}}\\


 \texttt{NiceVector} and \texttt{UglyVector} provide the linear bijection between the free left module \mbox{\texttt{\mdseries\slshape V}} and \texttt{\mbox{\texttt{\mdseries\slshape W}}:= NiceFreeLeftModule( \mbox{\texttt{\mdseries\slshape V}} )}. 

 If \mbox{\texttt{\mdseries\slshape v}} lies in the elements family of the family of \mbox{\texttt{\mdseries\slshape V}} then \texttt{NiceVector( \mbox{\texttt{\mdseries\slshape v}} )} is either \texttt{fail} or an element in the elements family of the family of \mbox{\texttt{\mdseries\slshape W}}. 

 If \mbox{\texttt{\mdseries\slshape r}} lies in the elements family of the family of \mbox{\texttt{\mdseries\slshape W}} then \texttt{UglyVector( \mbox{\texttt{\mdseries\slshape r}} )} is either \texttt{fail} or an element in the elements family of the family of \mbox{\texttt{\mdseries\slshape V}}. 

 If \mbox{\texttt{\mdseries\slshape v}} lies in \mbox{\texttt{\mdseries\slshape V}} (which usually \emph{cannot} be checked without using \mbox{\texttt{\mdseries\slshape W}}) then \texttt{UglyVector( \mbox{\texttt{\mdseries\slshape V}}, NiceVector( \mbox{\texttt{\mdseries\slshape V}}, \mbox{\texttt{\mdseries\slshape v}} ) ) = \mbox{\texttt{\mdseries\slshape v}}}. If \mbox{\texttt{\mdseries\slshape r}} lies in \mbox{\texttt{\mdseries\slshape W}} (which usually \emph{can} be checked) then \texttt{NiceVector( \mbox{\texttt{\mdseries\slshape V}}, UglyVector( \mbox{\texttt{\mdseries\slshape V}}, \mbox{\texttt{\mdseries\slshape r}} ) ) = \mbox{\texttt{\mdseries\slshape r}}}. 

 (This allows one to implement for example a membership test for \mbox{\texttt{\mdseries\slshape V}} using the membership test in \mbox{\texttt{\mdseries\slshape W}}.) }

 

\subsection{\textcolor{Chapter }{NiceFreeLeftModuleInfo}}
\logpage{[ 61, 11, 3 ]}\nobreak
\hyperdef{L}{X79350786800C2DD8}{}
{\noindent\textcolor{FuncColor}{$\triangleright$\enspace\texttt{NiceFreeLeftModuleInfo({\mdseries\slshape V})\index{NiceFreeLeftModuleInfo@\texttt{NiceFreeLeftModuleInfo}}
\label{NiceFreeLeftModuleInfo}
}\hfill{\scriptsize (attribute)}}\\


 For a free left module \mbox{\texttt{\mdseries\slshape V}} that is handled via the mechanism of nice bases, this operation has to provide
the necessary information (if any) for calls of \texttt{NiceVector} (\ref{NiceVector}) and \texttt{UglyVector} (\ref{UglyVector}). }

 

\subsection{\textcolor{Chapter }{NiceBasis}}
\logpage{[ 61, 11, 4 ]}\nobreak
\hyperdef{L}{X8388E0248690C214}{}
{\noindent\textcolor{FuncColor}{$\triangleright$\enspace\texttt{NiceBasis({\mdseries\slshape B})\index{NiceBasis@\texttt{NiceBasis}}
\label{NiceBasis}
}\hfill{\scriptsize (attribute)}}\\


 Let \mbox{\texttt{\mdseries\slshape B}} be a basis of a free left module \mbox{\texttt{\mdseries\slshape V}} that is handled via nice bases. If \mbox{\texttt{\mdseries\slshape B}} has no basis vectors stored at the time of the first call to \texttt{NiceBasis} then \texttt{NiceBasis( \mbox{\texttt{\mdseries\slshape B}} )} is obtained as \texttt{Basis( NiceFreeLeftModule( \mbox{\texttt{\mdseries\slshape V}} ) )}. If basis vectors are stored then \texttt{NiceBasis( \mbox{\texttt{\mdseries\slshape B}} )} is the result of the call of \texttt{Basis} with arguments \texttt{NiceFreeLeftModule( \mbox{\texttt{\mdseries\slshape V}} )} and the \texttt{NiceVector} values of the basis vectors of \mbox{\texttt{\mdseries\slshape B}}. 

 Note that the result is \texttt{fail} if and only if the ``basis vectors'' stored in \mbox{\texttt{\mdseries\slshape B}} are in fact not basis vectors. 

 The attributes \texttt{GeneratorsOfLeftModule} of the underlying left modules of \mbox{\texttt{\mdseries\slshape B}} and the result of \texttt{NiceBasis} correspond via \texttt{NiceVector} (\ref{NiceVector}) and \texttt{UglyVector} (\ref{UglyVector}). }

 

\subsection{\textcolor{Chapter }{IsBasisByNiceBasis}}
\logpage{[ 61, 11, 5 ]}\nobreak
\hyperdef{L}{X82BC30A487967F5B}{}
{\noindent\textcolor{FuncColor}{$\triangleright$\enspace\texttt{IsBasisByNiceBasis({\mdseries\slshape B})\index{IsBasisByNiceBasis@\texttt{IsBasisByNiceBasis}}
\label{IsBasisByNiceBasis}
}\hfill{\scriptsize (Category)}}\\


 This filter indicates that the basis \mbox{\texttt{\mdseries\slshape B}} delegates tasks such as the computation of coefficients (see{\nobreakspace}\texttt{Coefficients} (\ref{Coefficients})) to a basis of an isomorphic ``nicer'' free left module.   }

 

\subsection{\textcolor{Chapter }{IsHandledByNiceBasis}}
\logpage{[ 61, 11, 6 ]}\nobreak
\hyperdef{L}{X79D1DEA679AEDA40}{}
{\noindent\textcolor{FuncColor}{$\triangleright$\enspace\texttt{IsHandledByNiceBasis({\mdseries\slshape M})\index{IsHandledByNiceBasis@\texttt{IsHandledByNiceBasis}}
\label{IsHandledByNiceBasis}
}\hfill{\scriptsize (Category)}}\\


 For a free left module \mbox{\texttt{\mdseries\slshape M}} in this category, essentially all operations are performed using a ``nicer'' free left module, which is usually a row module. }

 }

  
\section{\textcolor{Chapter }{How to Implement New Kinds of Vector Spaces}}\label{How to Implement New Kinds of Vector Spaces}
\logpage{[ 61, 12, 0 ]}
\hyperdef{L}{X8238195B851D3C44}{}
{
   

\subsection{\textcolor{Chapter }{DeclareHandlingByNiceBasis}}
\logpage{[ 61, 12, 1 ]}\nobreak
\hyperdef{L}{X7DE34C3E837FCBC3}{}
{\noindent\textcolor{FuncColor}{$\triangleright$\enspace\texttt{DeclareHandlingByNiceBasis({\mdseries\slshape name, info})\index{DeclareHandlingByNiceBasis@\texttt{DeclareHandlingByNiceBasis}}
\label{DeclareHandlingByNiceBasis}
}\hfill{\scriptsize (function)}}\\
\noindent\textcolor{FuncColor}{$\triangleright$\enspace\texttt{InstallHandlingByNiceBasis({\mdseries\slshape name, record})\index{InstallHandlingByNiceBasis@\texttt{InstallHandlingByNiceBasis}}
\label{InstallHandlingByNiceBasis}
}\hfill{\scriptsize (function)}}\\


 These functions are used to implement a new kind of free left modules that
shall be handled via the mechanism of nice bases (see{\nobreakspace}\ref{Vector Spaces Handled By Nice Bases}). 

 \mbox{\texttt{\mdseries\slshape name}} must be a string, a filter $f$ with this name is created which implies \texttt{IsFreeLeftModule} (\ref{IsFreeLeftModule}), and a logical implication from the join of $f$ with \texttt{IsAttributeStoringRep} (\ref{IsAttributeStoringRep}) to \texttt{IsHandledByNiceBasis} (\ref{IsHandledByNiceBasis}) is installed. 

 \mbox{\texttt{\mdseries\slshape record}} must be a record with the following components. 
\begin{description}
\item[{\texttt{detect} }]  a function of four arguments $R$, $l$, $V$, and $z$, where $V$ is a free left module over the ring $R$ with generators the list or collection $l$, and $z$ is either the zero element of $V$ or \texttt{false} (then $l$ is nonempty); the function returns \texttt{true} if $V$ shall lie in the filter $f$, and \texttt{false} otherwise; the return value may also be \texttt{fail}, which indicates that $V$ is \emph{not} to be handled via the mechanism of nice bases at all, 
\item[{\texttt{NiceFreeLeftModuleInfo} }]  the \texttt{NiceFreeLeftModuleInfo} (\ref{NiceFreeLeftModuleInfo}) method for left modules in $f$, 
\item[{\texttt{NiceVector} }]  the \texttt{NiceVector} (\ref{NiceVector}) method for left modules $V$ in $f$; called with $V$ and a vector $v \in V$, this function returns the nice vector $r$ associated with $v$, and 
\item[{\texttt{UglyVector}}]  the \texttt{UglyVector} (\ref{UglyVector}) method for left modules $V$ in $f$; called with $V$ and a vector $r$ in the \texttt{NiceFreeLeftModule} (\ref{NiceFreeLeftModule}) value of $V$, this function returns the vector $v \in V$ to which $r$ is associated. 
\end{description}
 

 The idea is that all one has to do for implementing a new kind of free left
modules handled by the mechanism of nice bases is to call \texttt{DeclareHandlingByNiceBasis} and \texttt{InstallHandlingByNiceBasis}, which causes the installation of the necessary methods and adds the pair $[ f, $\mbox{\texttt{\mdseries\slshape record}}\texttt{.detect}$ ]$ to the global list \texttt{NiceBasisFiltersInfo} (\ref{NiceBasisFiltersInfo}). The \texttt{LeftModuleByGenerators} (\ref{LeftModuleByGenerators}) methods call \texttt{CheckForHandlingByNiceBasis} (\ref{CheckForHandlingByNiceBasis}), which sets the appropriate filter for the desired left module if applicable. }

 

\subsection{\textcolor{Chapter }{NiceBasisFiltersInfo}}
\logpage{[ 61, 12, 2 ]}\nobreak
\hyperdef{L}{X7E6077F0830A28DA}{}
{\noindent\textcolor{FuncColor}{$\triangleright$\enspace\texttt{NiceBasisFiltersInfo\index{NiceBasisFiltersInfo@\texttt{NiceBasisFiltersInfo}}
\label{NiceBasisFiltersInfo}
}\hfill{\scriptsize (global variable)}}\\


 An overview of all kinds of vector spaces that are currently handled by nice
bases is given by the global list \texttt{NiceBasisFiltersInfo}. Examples of such vector spaces are vector spaces of field elements (but not
the fields themselves) and non\texttt{\symbol{45}}Gaussian row and matrix
spaces (see{\nobreakspace}\texttt{IsGaussianSpace} (\ref{IsGaussianSpace})). }

 

\subsection{\textcolor{Chapter }{CheckForHandlingByNiceBasis}}
\logpage{[ 61, 12, 3 ]}\nobreak
\hyperdef{L}{X7A374553786DF5E7}{}
{\noindent\textcolor{FuncColor}{$\triangleright$\enspace\texttt{CheckForHandlingByNiceBasis({\mdseries\slshape R, gens, M, zero})\index{CheckForHandlingByNiceBasis@\texttt{CheckForHandlingByNiceBasis}}
\label{CheckForHandlingByNiceBasis}
}\hfill{\scriptsize (function)}}\\


 Whenever a free left module is constructed for which the filter \texttt{IsHandledByNiceBasis} may be useful, \texttt{CheckForHandlingByNiceBasis} should be called. (This is done in the methods for \texttt{VectorSpaceByGenerators}, \texttt{AlgebraByGenerators}, \texttt{IdealByGenerators} etc.{\nobreakspace}in the \textsf{GAP} library.) 

 The arguments of this function are the coefficient ring \mbox{\texttt{\mdseries\slshape R}}, the list \mbox{\texttt{\mdseries\slshape gens}} of generators, the constructed module \mbox{\texttt{\mdseries\slshape M}} itself, and the zero element \mbox{\texttt{\mdseries\slshape zero}} of \mbox{\texttt{\mdseries\slshape M}}; if \mbox{\texttt{\mdseries\slshape gens}} is nonempty then the \mbox{\texttt{\mdseries\slshape zero}} value may also be \texttt{false}. }

 }

 
\section{\textcolor{Chapter }{Tensor Products and Exterior and Symmetric Powers}}\label{Tensor Products and Exterior and Symmetric Powers}
\logpage{[ 61, 13, 0 ]}
\hyperdef{L}{X78515F448644204E}{}
{
  

\subsection{\textcolor{Chapter }{TensorProduct (for a list of vector spaces)}}
\logpage{[ 61, 13, 1 ]}\nobreak
\hyperdef{L}{X81B2276A7EBA8ED1}{}
{\noindent\textcolor{FuncColor}{$\triangleright$\enspace\texttt{TensorProduct({\mdseries\slshape list})\index{TensorProduct@\texttt{TensorProduct}!for a list of vector spaces}
\label{TensorProduct:for a list of vector spaces}
}\hfill{\scriptsize (operation)}}\\
\noindent\textcolor{FuncColor}{$\triangleright$\enspace\texttt{TensorProduct({\mdseries\slshape V, W, ...})\index{TensorProduct@\texttt{TensorProduct}!for vector spaces}
\label{TensorProduct:for vector spaces}
}\hfill{\scriptsize (operation)}}\\


 Here \mbox{\texttt{\mdseries\slshape list}} must be a list of vector spaces. This function returns the tensor product of
the elements in the list. The vector spaces must be defined over the same
field. 

 In the second form, the vector spaces are given individually. 

 Elements of the tensor product $V_1\otimes \cdots \otimes V_k$ are linear combinations of $v_1\otimes\cdots \otimes v_k$, where the $v_i$ are arbitrary basis elements of $V_i$. In \textsf{GAP} a tensor element like that is printed as 
\begin{Verbatim}[commandchars=!@|,fontsize=\small,frame=single,label=Example]
     v_1<x> ... <x>v_k
\end{Verbatim}
 Furthermore, the zero of a tensor product is printed as 
\begin{Verbatim}[commandchars=!@|,fontsize=\small,frame=single,label=Example]
   <0-tensor>
\end{Verbatim}
 This does not mean that all tensor products have the same zero element: zeros
of different tensor products have different families. 
\begin{Verbatim}[commandchars=!@|,fontsize=\small,frame=single,label=Example]
  !gapprompt@gap>| !gapinput@V:=TensorProduct(Rationals^2, Rationals^3);|
  <vector space over Rationals, with 6 generators>
  !gapprompt@gap>| !gapinput@Basis(V);|
  Basis( <vector space over Rationals, with 6 generators>,
  [ 1*([ 0, 1 ]<x>[ 0, 0, 1 ]), 1*([ 0, 1 ]<x>[ 0, 1, 0 ]),
    1*([ 0, 1 ]<x>[ 1, 0, 0 ]), 1*([ 1, 0 ]<x>[ 0, 0, 1 ]),
    1*([ 1, 0 ]<x>[ 0, 1, 0 ]), 1*([ 1, 0 ]<x>[ 1, 0, 0 ]) ] )
\end{Verbatim}
 See also \texttt{KroneckerProduct} (\ref{KroneckerProduct}). }

 

\subsection{\textcolor{Chapter }{ExteriorPower}}
\logpage{[ 61, 13, 2 ]}\nobreak
\hyperdef{L}{X787BB7FF85F0AD68}{}
{\noindent\textcolor{FuncColor}{$\triangleright$\enspace\texttt{ExteriorPower({\mdseries\slshape V, k})\index{ExteriorPower@\texttt{ExteriorPower}}
\label{ExteriorPower}
}\hfill{\scriptsize (operation)}}\\


 Here \mbox{\texttt{\mdseries\slshape V}} must be a vector space. This function returns the \mbox{\texttt{\mdseries\slshape k}}\texttt{\symbol{45}}th exterior power of \mbox{\texttt{\mdseries\slshape V}}. 

 Elements of the exterior power $\bigwedge^k V$ are linear combinations of $v_{i_1}\wedge\cdots \wedge v_{i_k}$, where the $v_{i_j}$ are basis elements of $V$, and $1 \leq i_1 < i_2 \cdots < i_k$. In \textsf{GAP} a wedge element like that is printed as 
\begin{Verbatim}[commandchars=!@|,fontsize=\small,frame=single,label=Example]
     v_1/\ ... /\v_k
\end{Verbatim}
 Furthermore, the zero of an exterior power is printed as 
\begin{Verbatim}[commandchars=!@|,fontsize=\small,frame=single,label=Example]
   <0-wedge>
\end{Verbatim}
 This does not mean that all exterior powers have the same zero element: zeros
of different exterior powers have different families. 
\begin{Verbatim}[commandchars=!@|,fontsize=\small,frame=single,label=Example]
  !gapprompt@gap>| !gapinput@V:=ExteriorPower(Rationals^3, 2);|
  <vector space of dimension 3 over Rationals>
  !gapprompt@gap>| !gapinput@Basis(V);|
  Basis( <vector space of dimension 3 over Rationals>, [
    1*([ 0, 1, 0 ]/\[ 0, 0, 1 ]), 1*([ 1, 0, 0 ]/\[ 0, 0, 1 ]),
    1*([ 1, 0, 0 ]/\[ 0, 1, 0 ]) ] )
\end{Verbatim}
 }

 

\subsection{\textcolor{Chapter }{SymmetricPower}}
\logpage{[ 61, 13, 3 ]}\nobreak
\hyperdef{L}{X79E2C2AF842E8419}{}
{\noindent\textcolor{FuncColor}{$\triangleright$\enspace\texttt{SymmetricPower({\mdseries\slshape V, k})\index{SymmetricPower@\texttt{SymmetricPower}}
\label{SymmetricPower}
}\hfill{\scriptsize (operation)}}\\


 Here \mbox{\texttt{\mdseries\slshape V}} must be a vector space. This function returns the \mbox{\texttt{\mdseries\slshape k}}\texttt{\symbol{45}}th symmetric power of \mbox{\texttt{\mdseries\slshape V}}. 
\begin{Verbatim}[commandchars=!@|,fontsize=\small,frame=single,label=Example]
  !gapprompt@gap>| !gapinput@V:=SymmetricPower(Rationals^3, 2);|
  <vector space over Rationals, with 6 generators>
  !gapprompt@gap>| !gapinput@Basis(V);|
  Basis( <vector space over Rationals, with 6 generators>,
  [ 1*([ 0, 0, 1 ].[ 0, 0, 1 ]), 1*([ 0, 1, 0 ].[ 0, 0, 1 ]),
    1*([ 0, 1, 0 ].[ 0, 1, 0 ]), 1*([ 1, 0, 0 ].[ 0, 0, 1 ]),
    1*([ 1, 0, 0 ].[ 0, 1, 0 ]), 1*([ 1, 0, 0 ].[ 1, 0, 0 ])
   ] )
\end{Verbatim}
 }

 }

 }

            
\chapter{\textcolor{Chapter }{Algebras}}\label{Algebras}
\logpage{[ 62, 0, 0 ]}
\hyperdef{L}{X7DDBF6F47A2E021C}{}
{
  An algebra is a vector space equipped with a bilinear map (multiplication).
This chapter describes the functions in \textsf{GAP} that deal with general algebras and associative algebras. 

 Algebras in \textsf{GAP} are vector spaces in a natural way. So all the functionality for vector spaces
(see Chapter{\nobreakspace}\ref{Vector Spaces}) is also applicable to algebras.    
\section{\textcolor{Chapter }{InfoAlgebra (Info Class)}}\label{sect:InfoAlgebra}
\logpage{[ 62, 1, 0 ]}
\hyperdef{L}{X830EDB5F85645FFB}{}
{
  

\subsection{\textcolor{Chapter }{InfoAlgebra}}
\logpage{[ 62, 1, 1 ]}\nobreak
\hyperdef{L}{X8665F459841AAD53}{}
{\noindent\textcolor{FuncColor}{$\triangleright$\enspace\texttt{InfoAlgebra\index{InfoAlgebra@\texttt{InfoAlgebra}}
\label{InfoAlgebra}
}\hfill{\scriptsize (info class)}}\\


 is the info class for the functions dealing with algebras (see{\nobreakspace}\ref{Info Functions}). }

 }

  
\section{\textcolor{Chapter }{Constructing Algebras by Generators}}\label{Constructing Algebras by Generators}
\logpage{[ 62, 2, 0 ]}
\hyperdef{L}{X8686DEBA85D3F3B6}{}
{
   

\subsection{\textcolor{Chapter }{Algebra}}
\logpage{[ 62, 2, 1 ]}\nobreak
\hyperdef{L}{X7B213851791A594B}{}
{\noindent\textcolor{FuncColor}{$\triangleright$\enspace\texttt{Algebra({\mdseries\slshape F, gens[, zero][, "basis"]})\index{Algebra@\texttt{Algebra}}
\label{Algebra}
}\hfill{\scriptsize (function)}}\\


 \texttt{Algebra( \mbox{\texttt{\mdseries\slshape F}}, \mbox{\texttt{\mdseries\slshape gens}} )} is the algebra over the division ring \mbox{\texttt{\mdseries\slshape F}}, generated by the vectors in the list \mbox{\texttt{\mdseries\slshape gens}}. 

 If there are three arguments, a division ring \mbox{\texttt{\mdseries\slshape F}} and a list \mbox{\texttt{\mdseries\slshape gens}} and an element \mbox{\texttt{\mdseries\slshape zero}}, then \texttt{Algebra( \mbox{\texttt{\mdseries\slshape F}}, \mbox{\texttt{\mdseries\slshape gens}}, \mbox{\texttt{\mdseries\slshape zero}} )} is the \mbox{\texttt{\mdseries\slshape F}}\texttt{\symbol{45}}algebra generated by \mbox{\texttt{\mdseries\slshape gens}}, with zero element \mbox{\texttt{\mdseries\slshape zero}}. 

 If the last argument is the string \texttt{"basis"} then the vectors in \mbox{\texttt{\mdseries\slshape gens}} are known to form a basis of the algebra (as an \mbox{\texttt{\mdseries\slshape F}}\texttt{\symbol{45}}vector space). 
\begin{Verbatim}[commandchars=!@|,fontsize=\small,frame=single,label=Example]
  !gapprompt@gap>| !gapinput@m:= [ [ 0, 1, 2 ], [ 0, 0, 3], [ 0, 0, 0 ] ];;|
  !gapprompt@gap>| !gapinput@A:= Algebra( Rationals, [ m ] );|
  <algebra over Rationals, with 1 generator>
  !gapprompt@gap>| !gapinput@Dimension( A );|
  2
\end{Verbatim}
 }

  

\subsection{\textcolor{Chapter }{AlgebraWithOne}}
\logpage{[ 62, 2, 2 ]}\nobreak
\hyperdef{L}{X80FE16EA84EE56CD}{}
{\noindent\textcolor{FuncColor}{$\triangleright$\enspace\texttt{AlgebraWithOne({\mdseries\slshape F, gens[, zero][, "basis"]})\index{AlgebraWithOne@\texttt{AlgebraWithOne}}
\label{AlgebraWithOne}
}\hfill{\scriptsize (function)}}\\


 \texttt{AlgebraWithOne( \mbox{\texttt{\mdseries\slshape F}}, \mbox{\texttt{\mdseries\slshape gens}} )} is the algebra\texttt{\symbol{45}}with\texttt{\symbol{45}}one over the
division ring \mbox{\texttt{\mdseries\slshape F}}, generated by the vectors in the list \mbox{\texttt{\mdseries\slshape gens}}. 

 If there are three arguments, a division ring \mbox{\texttt{\mdseries\slshape F}} and a list \mbox{\texttt{\mdseries\slshape gens}} and an element \mbox{\texttt{\mdseries\slshape zero}}, then \texttt{AlgebraWithOne( \mbox{\texttt{\mdseries\slshape F}}, \mbox{\texttt{\mdseries\slshape gens}}, \mbox{\texttt{\mdseries\slshape zero}} )} is the \mbox{\texttt{\mdseries\slshape F}}\texttt{\symbol{45}}algebra\texttt{\symbol{45}}with\texttt{\symbol{45}}one
generated by \mbox{\texttt{\mdseries\slshape gens}}, with zero element \mbox{\texttt{\mdseries\slshape zero}}. 

 If the last argument is the string \texttt{"basis"} then the vectors in \mbox{\texttt{\mdseries\slshape gens}} are known to form a basis of the algebra (as an \mbox{\texttt{\mdseries\slshape F}}\texttt{\symbol{45}}vector space). 
\begin{Verbatim}[commandchars=!@|,fontsize=\small,frame=single,label=Example]
  !gapprompt@gap>| !gapinput@m:= [ [ 0, 1, 2 ], [ 0, 0, 3], [ 0, 0, 0 ] ];;|
  !gapprompt@gap>| !gapinput@A:= AlgebraWithOne( Rationals, [ m ] );|
  <algebra-with-one over Rationals, with 1 generator>
  !gapprompt@gap>| !gapinput@Dimension( A );|
  3
  !gapprompt@gap>| !gapinput@One(A);|
  [ [ 1, 0, 0 ], [ 0, 1, 0 ], [ 0, 0, 1 ] ]
\end{Verbatim}
 }

 }

  
\section{\textcolor{Chapter }{Constructing Algebras as Free Algebras}}\label{Constructing Algebras as Free Algebras}
\logpage{[ 62, 3, 0 ]}
\hyperdef{L}{X7A7B00127DC9DD40}{}
{
  

\subsection{\textcolor{Chapter }{FreeAlgebra (for ring, rank (and name))}}
\logpage{[ 62, 3, 1 ]}\nobreak
\hyperdef{L}{X83484C917D8F7A1A}{}
{\noindent\textcolor{FuncColor}{$\triangleright$\enspace\texttt{FreeAlgebra({\mdseries\slshape R, rank[, name]})\index{FreeAlgebra@\texttt{FreeAlgebra}!for ring, rank (and name)}
\label{FreeAlgebra:for ring, rank (and name)}
}\hfill{\scriptsize (function)}}\\
\noindent\textcolor{FuncColor}{$\triangleright$\enspace\texttt{FreeAlgebra({\mdseries\slshape R, name1, name2, ...})\index{FreeAlgebra@\texttt{FreeAlgebra}!for ring and several names}
\label{FreeAlgebra:for ring and several names}
}\hfill{\scriptsize (function)}}\\


 is a free (nonassociative) algebra of rank \mbox{\texttt{\mdseries\slshape rank}} over the division ring \mbox{\texttt{\mdseries\slshape R}}. Here \mbox{\texttt{\mdseries\slshape name}}, and \mbox{\texttt{\mdseries\slshape name1}}, \mbox{\texttt{\mdseries\slshape name2}}, ... are optional strings that can be used to provide names for the
generators. 
\begin{Verbatim}[commandchars=!@|,fontsize=\small,frame=single,label=Example]
  !gapprompt@gap>| !gapinput@A:= FreeAlgebra( Rationals, "a", "b" );|
  <algebra over Rationals, with 2 generators>
  !gapprompt@gap>| !gapinput@g:= GeneratorsOfAlgebra( A );|
  [ (1)*a, (1)*b ]
  !gapprompt@gap>| !gapinput@(g[1]*g[2])*((g[2]*g[1])*g[1]);|
  (1)*((a*b)*((b*a)*a))
\end{Verbatim}
 }

 

\subsection{\textcolor{Chapter }{FreeAlgebraWithOne (for ring, rank (and name))}}
\logpage{[ 62, 3, 2 ]}\nobreak
\hyperdef{L}{X7FBD04B07B85623D}{}
{\noindent\textcolor{FuncColor}{$\triangleright$\enspace\texttt{FreeAlgebraWithOne({\mdseries\slshape R, rank[, name]})\index{FreeAlgebraWithOne@\texttt{FreeAlgebraWithOne}!for ring, rank (and name)}
\label{FreeAlgebraWithOne:for ring, rank (and name)}
}\hfill{\scriptsize (function)}}\\
\noindent\textcolor{FuncColor}{$\triangleright$\enspace\texttt{FreeAlgebraWithOne({\mdseries\slshape R, name1, name2, ...})\index{FreeAlgebraWithOne@\texttt{FreeAlgebraWithOne}!for ring and several names}
\label{FreeAlgebraWithOne:for ring and several names}
}\hfill{\scriptsize (function)}}\\


 is a free (nonassociative)
algebra\texttt{\symbol{45}}with\texttt{\symbol{45}}one of rank \mbox{\texttt{\mdseries\slshape rank}} over the division ring \mbox{\texttt{\mdseries\slshape R}}. Here \mbox{\texttt{\mdseries\slshape name}}, and \mbox{\texttt{\mdseries\slshape name1}}, \mbox{\texttt{\mdseries\slshape name2}}, ... are optional strings that can be used to provide names for the
generators. 
\begin{Verbatim}[commandchars=!@|,fontsize=\small,frame=single,label=Example]
  !gapprompt@gap>| !gapinput@A:= FreeAlgebraWithOne( Rationals, 4, "q" );|
  <algebra-with-one over Rationals, with 4 generators>
  !gapprompt@gap>| !gapinput@GeneratorsOfAlgebra( A );|
  [ (1)*<identity ...>, (1)*q.1, (1)*q.2, (1)*q.3, (1)*q.4 ]
  !gapprompt@gap>| !gapinput@One( A );|
  (1)*<identity ...>
\end{Verbatim}
 }

 

\subsection{\textcolor{Chapter }{FreeAssociativeAlgebra (for ring, rank (and name))}}
\logpage{[ 62, 3, 3 ]}\nobreak
\hyperdef{L}{X87835FFE79D2E068}{}
{\noindent\textcolor{FuncColor}{$\triangleright$\enspace\texttt{FreeAssociativeAlgebra({\mdseries\slshape R, rank[, name]})\index{FreeAssociativeAlgebra@\texttt{FreeAssociativeAlgebra}!for ring, rank (and name)}
\label{FreeAssociativeAlgebra:for ring, rank (and name)}
}\hfill{\scriptsize (function)}}\\
\noindent\textcolor{FuncColor}{$\triangleright$\enspace\texttt{FreeAssociativeAlgebra({\mdseries\slshape R, name1, name2, ...})\index{FreeAssociativeAlgebra@\texttt{FreeAssociativeAlgebra}!for ring and several names}
\label{FreeAssociativeAlgebra:for ring and several names}
}\hfill{\scriptsize (function)}}\\


 is a free associative algebra of rank \mbox{\texttt{\mdseries\slshape rank}} over the division ring \mbox{\texttt{\mdseries\slshape R}}. Here \mbox{\texttt{\mdseries\slshape name}}, and \mbox{\texttt{\mdseries\slshape name1}}, \mbox{\texttt{\mdseries\slshape name2}}, ... are optional strings that can be used to provide names for the
generators. 
\begin{Verbatim}[commandchars=!@|,fontsize=\small,frame=single,label=Example]
  !gapprompt@gap>| !gapinput@A:= FreeAssociativeAlgebra( GF( 5 ), 4, "a" );|
  <algebra over GF(5), with 4 generators>
\end{Verbatim}
 }

 

\subsection{\textcolor{Chapter }{FreeAssociativeAlgebraWithOne (for ring, rank (and name))}}
\logpage{[ 62, 3, 4 ]}\nobreak
\hyperdef{L}{X845A777584A7D711}{}
{\noindent\textcolor{FuncColor}{$\triangleright$\enspace\texttt{FreeAssociativeAlgebraWithOne({\mdseries\slshape R, rank[, name]})\index{FreeAssociativeAlgebraWithOne@\texttt{FreeAssociativeAlgebraWithOne}!for ring, rank (and name)}
\label{FreeAssociativeAlgebraWithOne:for ring, rank (and name)}
}\hfill{\scriptsize (function)}}\\
\noindent\textcolor{FuncColor}{$\triangleright$\enspace\texttt{FreeAssociativeAlgebraWithOne({\mdseries\slshape R, name1, name2, ...})\index{FreeAssociativeAlgebraWithOne@\texttt{FreeAssociativeAlgebraWithOne}!for ring and several names}
\label{FreeAssociativeAlgebraWithOne:for ring and several names}
}\hfill{\scriptsize (function)}}\\


 is a free associative algebra\texttt{\symbol{45}}with\texttt{\symbol{45}}one
of rank \mbox{\texttt{\mdseries\slshape rank}} over the division ring \mbox{\texttt{\mdseries\slshape R}}. Here \mbox{\texttt{\mdseries\slshape name}}, and \mbox{\texttt{\mdseries\slshape name1}}, \mbox{\texttt{\mdseries\slshape name2}}, ... are optional strings that can be used to provide names for the
generators. 
\begin{Verbatim}[commandchars=!@|,fontsize=\small,frame=single,label=Example]
  !gapprompt@gap>| !gapinput@A:= FreeAssociativeAlgebraWithOne( Rationals, "a", "b", "c" );|
  <algebra-with-one over Rationals, with 3 generators>
  !gapprompt@gap>| !gapinput@GeneratorsOfAlgebra( A );|
  [ (1)*<identity ...>, (1)*a, (1)*b, (1)*c ]
  !gapprompt@gap>| !gapinput@One( A );|
  (1)*<identity ...>
\end{Verbatim}
 }

 }

  
\section{\textcolor{Chapter }{Constructing Algebras by Structure Constants}}\label{Constructing Algebras by Structure Constants}
\logpage{[ 62, 4, 0 ]}
\hyperdef{L}{X7E8F45547CC07CE5}{}
{
  For an introduction into structure constants and how they are handled by \textsf{GAP}, we refer to Section  (\textbf{Tutorial: Algebras}) of the user's tutorial. 

\subsection{\textcolor{Chapter }{AlgebraByStructureConstants}}
\logpage{[ 62, 4, 1 ]}\nobreak
\hyperdef{L}{X7CC58DFD816E6B65}{}
{\noindent\textcolor{FuncColor}{$\triangleright$\enspace\texttt{AlgebraByStructureConstants({\mdseries\slshape R, sctable[, nameinfo]})\index{AlgebraByStructureConstants@\texttt{AlgebraByStructureConstants}}
\label{AlgebraByStructureConstants}
}\hfill{\scriptsize (function)}}\\


 returns a free left module $A$ over the division ring \mbox{\texttt{\mdseries\slshape R}}, with multiplication defined by the structure constants table \mbox{\texttt{\mdseries\slshape sctable}}. The optional argument \mbox{\texttt{\mdseries\slshape nameinfo}} can be used to prescribe names for the elements of the canonical basis of $A$; it can be either a string \mbox{\texttt{\mdseries\slshape name}} (then \mbox{\texttt{\mdseries\slshape name}}\texttt{1}, \mbox{\texttt{\mdseries\slshape name}}\texttt{2} etc. are chosen) or a list of strings which are then chosen. The vectors of
the canonical basis of $A$ correspond to the vectors of the basis given by \mbox{\texttt{\mdseries\slshape sctable}}. 

  It is \emph{not} checked whether the coefficients in \mbox{\texttt{\mdseries\slshape sctable}} are really elements in \mbox{\texttt{\mdseries\slshape R}}. 
\begin{Verbatim}[commandchars=!@|,fontsize=\small,frame=single,label=Example]
  !gapprompt@gap>| !gapinput@T:= EmptySCTable( 2, 0 );;|
  !gapprompt@gap>| !gapinput@SetEntrySCTable( T, 1, 1, [ 1/2, 1, 2/3, 2 ] );|
  !gapprompt@gap>| !gapinput@A:= AlgebraByStructureConstants( Rationals, T );|
  <algebra of dimension 2 over Rationals>
  !gapprompt@gap>| !gapinput@b:= BasisVectors( Basis( A ) );;|
  !gapprompt@gap>| !gapinput@b[1]^2;|
  (1/2)*v.1+(2/3)*v.2
  !gapprompt@gap>| !gapinput@b[1]*b[2];|
  0*v.1
\end{Verbatim}
 }

 

\subsection{\textcolor{Chapter }{AlgebraWithOneByStructureConstants}}
\logpage{[ 62, 4, 2 ]}\nobreak
\hyperdef{L}{X80D528A57FD64873}{}
{\noindent\textcolor{FuncColor}{$\triangleright$\enspace\texttt{AlgebraWithOneByStructureConstants({\mdseries\slshape R, sctable[, nameinfo], onecoeffs})\index{AlgebraWithOneByStructureConstants@\texttt{AlgebraWithOneByStructureConstants}}
\label{AlgebraWithOneByStructureConstants}
}\hfill{\scriptsize (function)}}\\


 The only differences between this function and \texttt{AlgebraByStructureConstants} (\ref{AlgebraByStructureConstants}) are that \texttt{AlgebraWithOneByStructureConstants} takes an additional argument \mbox{\texttt{\mdseries\slshape onecoeffs}}, the coefficients vector over the ring \mbox{\texttt{\mdseries\slshape R}} that describes the unique multiplicative identity element of the returned
algebra w. r. t. the defining basis of this algebra, and that the returned
algebra is an algebra\texttt{\symbol{45}}with\texttt{\symbol{45}}one (see \texttt{IsAlgebraWithOne} (\ref{IsAlgebraWithOne})). 

 
\begin{Verbatim}[commandchars=!@|,fontsize=\small,frame=single,label=Example]
  !gapprompt@gap>| !gapinput@A:= GF(2)^[2,2];;|
  !gapprompt@gap>| !gapinput@B:= Basis( A );;|
  !gapprompt@gap>| !gapinput@onecoeffs:= Coefficients( B, One( A ) );|
  [ Z(2)^0, 0*Z(2), 0*Z(2), Z(2)^0 ]
  !gapprompt@gap>| !gapinput@T:= StructureConstantsTable( B );;|
  !gapprompt@gap>| !gapinput@sc1:= AlgebraByStructureConstants( GF(2), T );|
  <algebra of dimension 4 over GF(2)>
  !gapprompt@gap>| !gapinput@HasOne( sc1 );|
  false
  !gapprompt@gap>| !gapinput@One( sc1 );|
  v.1+v.4
  !gapprompt@gap>| !gapinput@sc2:= AlgebraWithOneByStructureConstants( GF(2), T, onecoeffs );|
  <algebra-with-one of dimension 4 over GF(2)>
  !gapprompt@gap>| !gapinput@HasOne( sc2 );|
  true
  !gapprompt@gap>| !gapinput@One( sc2 );|
  v.1+v.4
\end{Verbatim}
 }

 

\subsection{\textcolor{Chapter }{StructureConstantsTable}}
\logpage{[ 62, 4, 3 ]}\nobreak
\hyperdef{L}{X804ADF0280F67CDC}{}
{\noindent\textcolor{FuncColor}{$\triangleright$\enspace\texttt{StructureConstantsTable({\mdseries\slshape B})\index{StructureConstantsTable@\texttt{StructureConstantsTable}}
\label{StructureConstantsTable}
}\hfill{\scriptsize (attribute)}}\\


 Let \mbox{\texttt{\mdseries\slshape B}} be a basis of a free left module $R$ that is also a ring. In this case \texttt{StructureConstantsTable} returns a structure constants table $T$ in sparse representation, as used for structure constants algebras (see
Section{\nobreakspace} (\textbf{Tutorial: Algebras}) of the \textsf{GAP} User's Tutorial). 

 If \mbox{\texttt{\mdseries\slshape B}} has length $n$ then $T$ is a list of length $n+2$. The first $n$ entries of $T$ are lists of length $n$. $T[ n+1 ]$ is one of $1$, $-1$, or $0$; in the case of $1$ the table is known to be symmetric, in the case of $-1$ it is known to be antisymmetric, and $0$ occurs in all other cases. $T[ n+2 ]$ is the zero element of the coefficient domain. 

 The coefficients w.r.t.{\nobreakspace}\mbox{\texttt{\mdseries\slshape B}} of the product of the $i$\texttt{\symbol{45}}th and $j$\texttt{\symbol{45}}th basis vector of \mbox{\texttt{\mdseries\slshape B}} are stored in $T[i][j]$ as a list of length $2$; its first entry is the list of positions of nonzero coefficients, the second
entry is the list of these coefficients themselves. 

 The multiplication in an algebra $A$ with vector space basis \mbox{\texttt{\mdseries\slshape B}} with basis vectors $[ v_1, \ldots, v_n ]$ is determined by the so\texttt{\symbol{45}}called structure matrices $M_k = [ m_{ijk} ]_{ij}$, $1 \leq k \leq n$. The $M_k$ are defined by $v_i v_j = \sum_k m_{ijk} v_k$. Let $a = [ a_1, \ldots, a_n ]$ and $b = [ b_1, \ldots, b_n ]$. Then 
\[ \left( \sum_i a_i v_i \right) \left( \sum_j b_j v_j \right) = \sum_{{i,j}} a_i
b_j \left( v_i v_j \right) = \sum_k \left( \sum_j \left( \sum_i a_i m_{ijk}
\right) b_j \right) v_k = \sum_k \left( a M_k b^{tr} \right) v_k. \]
 

 
\begin{Verbatim}[commandchars=!@|,fontsize=\small,frame=single,label=Example]
  !gapprompt@gap>| !gapinput@A:= QuaternionAlgebra( Rationals );;|
  !gapprompt@gap>| !gapinput@StructureConstantsTable( Basis( A ) );|
  [ [ [ [ 1 ], [ 1 ] ], [ [ 2 ], [ 1 ] ], [ [ 3 ], [ 1 ] ],
        [ [ 4 ], [ 1 ] ] ],
    [ [ [ 2 ], [ 1 ] ], [ [ 1 ], [ -1 ] ], [ [ 4 ], [ 1 ] ],
        [ [ 3 ], [ -1 ] ] ],
    [ [ [ 3 ], [ 1 ] ], [ [ 4 ], [ -1 ] ], [ [ 1 ], [ -1 ] ],
        [ [ 2 ], [ 1 ] ] ],
    [ [ [ 4 ], [ 1 ] ], [ [ 3 ], [ 1 ] ], [ [ 2 ], [ -1 ] ],
        [ [ 1 ], [ -1 ] ] ], 0, 0 ]
\end{Verbatim}
 }

 

\subsection{\textcolor{Chapter }{EmptySCTable}}
\logpage{[ 62, 4, 4 ]}\nobreak
\hyperdef{L}{X7F1203A1793411DF}{}
{\noindent\textcolor{FuncColor}{$\triangleright$\enspace\texttt{EmptySCTable({\mdseries\slshape dim, zero[, flag]})\index{EmptySCTable@\texttt{EmptySCTable}}
\label{EmptySCTable}
}\hfill{\scriptsize (function)}}\\


 \texttt{EmptySCTable} returns a structure constants table for an algebra of dimension \mbox{\texttt{\mdseries\slshape dim}}, describing trivial multiplication. \mbox{\texttt{\mdseries\slshape zero}} must be the zero of the coefficients domain. If the multiplication is known to
be (anti)commutative then this can be indicated by the optional third argument \mbox{\texttt{\mdseries\slshape flag}}, which must be one of the strings \texttt{"symmetric"}, \texttt{"antisymmetric"}. 

 For filling up the structure constants table, see \texttt{SetEntrySCTable} (\ref{SetEntrySCTable}). 
\begin{Verbatim}[commandchars=!@|,fontsize=\small,frame=single,label=Example]
  !gapprompt@gap>| !gapinput@EmptySCTable( 2, Zero( GF(5) ), "antisymmetric" );|
  [ [ [ [  ], [  ] ], [ [  ], [  ] ] ],
    [ [ [  ], [  ] ], [ [  ], [  ] ] ], -1, 0*Z(5) ]
\end{Verbatim}
 }

 

\subsection{\textcolor{Chapter }{SetEntrySCTable}}
\logpage{[ 62, 4, 5 ]}\nobreak
\hyperdef{L}{X817BD086876EC1C4}{}
{\noindent\textcolor{FuncColor}{$\triangleright$\enspace\texttt{SetEntrySCTable({\mdseries\slshape T, i, j, list})\index{SetEntrySCTable@\texttt{SetEntrySCTable}}
\label{SetEntrySCTable}
}\hfill{\scriptsize (function)}}\\


 sets the entry of the structure constants table \mbox{\texttt{\mdseries\slshape T}} that describes the product of the \mbox{\texttt{\mdseries\slshape i}}\texttt{\symbol{45}}th basis element with the \mbox{\texttt{\mdseries\slshape j}}\texttt{\symbol{45}}th basis element to the value given by the list \mbox{\texttt{\mdseries\slshape list}}. 

 If \mbox{\texttt{\mdseries\slshape T}} is known to be antisymmetric or symmetric then also the value \texttt{\mbox{\texttt{\mdseries\slshape T}}[\mbox{\texttt{\mdseries\slshape j}}][\mbox{\texttt{\mdseries\slshape i}}]} is set. 

 \mbox{\texttt{\mdseries\slshape list}} must be of the form $[ c_{ij}^{{k_1}}, k_1, c_{ij}^{{k_2}}, k_2, \ldots ]$. 

 The entries at the odd positions of \mbox{\texttt{\mdseries\slshape list}} must be compatible with the zero element stored in \mbox{\texttt{\mdseries\slshape T}}. For convenience, these entries may also be rational numbers that are
automatically replaced by the corresponding elements in the appropriate prime
field in finite characteristic if necessary. 
\begin{Verbatim}[commandchars=!@|,fontsize=\small,frame=single,label=Example]
  !gapprompt@gap>| !gapinput@T:= EmptySCTable( 2, 0 );;|
  !gapprompt@gap>| !gapinput@SetEntrySCTable( T, 1, 1, [ 1/2, 1, 2/3, 2 ] );|
  !gapprompt@gap>| !gapinput@T;|
  [ [ [ [ 1, 2 ], [ 1/2, 2/3 ] ], [ [  ], [  ] ] ],
    [ [ [  ], [  ] ], [ [  ], [  ] ] ], 0, 0 ]
\end{Verbatim}
 }

 

\subsection{\textcolor{Chapter }{GapInputSCTable}}
\logpage{[ 62, 4, 6 ]}\nobreak
\hyperdef{L}{X7F333822780B6731}{}
{\noindent\textcolor{FuncColor}{$\triangleright$\enspace\texttt{GapInputSCTable({\mdseries\slshape T, varname})\index{GapInputSCTable@\texttt{GapInputSCTable}}
\label{GapInputSCTable}
}\hfill{\scriptsize (function)}}\\


 is a string that describes the structure constants table \mbox{\texttt{\mdseries\slshape T}} in terms of \texttt{EmptySCTable} (\ref{EmptySCTable}) and \texttt{SetEntrySCTable} (\ref{SetEntrySCTable}). The assignments are made to the variable \mbox{\texttt{\mdseries\slshape varname}}. 
\begin{Verbatim}[commandchars=!@|,fontsize=\small,frame=single,label=Example]
  !gapprompt@gap>| !gapinput@T:= EmptySCTable( 2, 0 );;|
  !gapprompt@gap>| !gapinput@SetEntrySCTable( T, 1, 2, [ 1, 2 ] );|
  !gapprompt@gap>| !gapinput@SetEntrySCTable( T, 2, 1, [ 1, 2 ] );|
  !gapprompt@gap>| !gapinput@GapInputSCTable( T, "T" );|
  "T:= EmptySCTable( 2, 0 );\nSetEntrySCTable( T, 1, 2, [1,2] );\nSetEnt\
  rySCTable( T, 2, 1, [1,2] );\n"
\end{Verbatim}
 }

 

\subsection{\textcolor{Chapter }{TestJacobi}}
\logpage{[ 62, 4, 7 ]}\nobreak
\hyperdef{L}{X7C23ED85814C0371}{}
{\noindent\textcolor{FuncColor}{$\triangleright$\enspace\texttt{TestJacobi({\mdseries\slshape T})\index{TestJacobi@\texttt{TestJacobi}}
\label{TestJacobi}
}\hfill{\scriptsize (function)}}\\


 tests whether the structure constants table \mbox{\texttt{\mdseries\slshape T}} satisfies the Jacobi identity $v_i * (v_j * v_k) + v_j * (v_k * v_i) + v_k * (v_i * v_j) = 0$ for all basis vectors $v_i$ of the underlying algebra, where $i \leq j \leq k$. (Thus antisymmetry is assumed.) 

 The function returns \texttt{true} if the Jacobi identity is satisfied, and a failing triple $[ i, j, k ]$ otherwise. 
\begin{Verbatim}[commandchars=!@|,fontsize=\small,frame=single,label=Example]
  !gapprompt@gap>| !gapinput@T:= EmptySCTable( 2, 0, "antisymmetric" );;|
  !gapprompt@gap>| !gapinput@SetEntrySCTable( T, 1, 2, [ 1, 2 ] );;|
  !gapprompt@gap>| !gapinput@TestJacobi( T );|
  true
\end{Verbatim}
 }

 

\subsection{\textcolor{Chapter }{IdentityFromSCTable}}
\logpage{[ 62, 4, 8 ]}\nobreak
\hyperdef{L}{X78B633CE7A5B9F9A}{}
{\noindent\textcolor{FuncColor}{$\triangleright$\enspace\texttt{IdentityFromSCTable({\mdseries\slshape T})\index{IdentityFromSCTable@\texttt{IdentityFromSCTable}}
\label{IdentityFromSCTable}
}\hfill{\scriptsize (function)}}\\


 Let \mbox{\texttt{\mdseries\slshape T}} be a structure constants table of an algebra $A$ of dimension $n$. \texttt{IdentityFromSCTable( \mbox{\texttt{\mdseries\slshape T}} )} is either \texttt{fail} or the vector of length $n$ that contains the coefficients of the multiplicative identity of $A$ with respect to the basis that belongs to \mbox{\texttt{\mdseries\slshape T}}. 
\begin{Verbatim}[commandchars=!@|,fontsize=\small,frame=single,label=Example]
  !gapprompt@gap>| !gapinput@T:= EmptySCTable( 2, 0 );;|
  !gapprompt@gap>| !gapinput@SetEntrySCTable( T, 1, 1, [ 1, 1 ] );;|
  !gapprompt@gap>| !gapinput@SetEntrySCTable( T, 1, 2, [ 1, 2 ] );;|
  !gapprompt@gap>| !gapinput@SetEntrySCTable( T, 2, 1, [ 1, 2 ] );;|
  !gapprompt@gap>| !gapinput@IdentityFromSCTable( T );|
  [ 1, 0 ]
\end{Verbatim}
 }

 

\subsection{\textcolor{Chapter }{QuotientFromSCTable}}
\logpage{[ 62, 4, 9 ]}\nobreak
\hyperdef{L}{X7F2A71608602635D}{}
{\noindent\textcolor{FuncColor}{$\triangleright$\enspace\texttt{QuotientFromSCTable({\mdseries\slshape T, num, den})\index{QuotientFromSCTable@\texttt{QuotientFromSCTable}}
\label{QuotientFromSCTable}
}\hfill{\scriptsize (function)}}\\


 Let \mbox{\texttt{\mdseries\slshape T}} be a structure constants table of an algebra $A$ of dimension $n$. \texttt{QuotientFromSCTable( \mbox{\texttt{\mdseries\slshape T}} )} is either \texttt{fail} or the vector of length $n$ that contains the coefficients of the quotient of \mbox{\texttt{\mdseries\slshape num}} and \mbox{\texttt{\mdseries\slshape den}} with respect to the basis that belongs to \mbox{\texttt{\mdseries\slshape T}}. 

 We solve the equation system \mbox{\texttt{\mdseries\slshape num}}$ = x *$ \mbox{\texttt{\mdseries\slshape den}}. If no solution exists, \texttt{fail} is returned. 

 In terms of the basis $B$ with vectors $b_1, \ldots, b_n$ this means for $\mbox{\texttt{\mdseries\slshape num}} = \sum_{{i = 1}}^n a_i b_i$, $\mbox{\texttt{\mdseries\slshape den}} = \sum_{{i = 1}}^n c_i b_i$, $x = \sum_{{i = 1}}^n x_i b_i$ that $a_k = \sum_{{i,j}} c_i x_j c_{ijk}$ for all $k$. Here $c_{ijk}$ denotes the structure constants with respect to $B$. This means that (as a vector) $a = x M$ with $M_{jk} = \sum_{{i = 1}}^n c_{ijk} c_i$. 

 
\begin{Verbatim}[commandchars=!@|,fontsize=\small,frame=single,label=Example]
  !gapprompt@gap>| !gapinput@T:= EmptySCTable( 2, 0 );;|
  !gapprompt@gap>| !gapinput@SetEntrySCTable( T, 1, 1, [ 1, 1 ] );;|
  !gapprompt@gap>| !gapinput@SetEntrySCTable( T, 2, 1, [ 1, 2 ] );;|
  !gapprompt@gap>| !gapinput@SetEntrySCTable( T, 1, 2, [ 1, 2 ] );;|
  !gapprompt@gap>| !gapinput@QuotientFromSCTable( T, [0,1], [1,0] );|
  [ 0, 1 ]
\end{Verbatim}
 }

 }

  
\section{\textcolor{Chapter }{Some Special Algebras}}\label{Some Special Algebras}
\logpage{[ 62, 5, 0 ]}
\hyperdef{L}{X79B7C3078112E7E1}{}
{
  

\subsection{\textcolor{Chapter }{QuaternionAlgebra}}
\logpage{[ 62, 5, 1 ]}\nobreak
\hyperdef{L}{X83DF4BCC7CE494FC}{}
{\noindent\textcolor{FuncColor}{$\triangleright$\enspace\texttt{QuaternionAlgebra({\mdseries\slshape F[, a, b]})\index{QuaternionAlgebra@\texttt{QuaternionAlgebra}}
\label{QuaternionAlgebra}
}\hfill{\scriptsize (function)}}\\
\textbf{\indent Returns:\ }
 a quaternion algebra over \mbox{\texttt{\mdseries\slshape F}}, with parameters \mbox{\texttt{\mdseries\slshape a}} and \mbox{\texttt{\mdseries\slshape b}}. 



 Let \mbox{\texttt{\mdseries\slshape F}} be a field or a list of field elements, let $F$ be the field generated by \mbox{\texttt{\mdseries\slshape F}}, and let \mbox{\texttt{\mdseries\slshape a}} and \mbox{\texttt{\mdseries\slshape b}} two elements in $F$. \texttt{QuaternionAlgebra} returns a quaternion algebra over $F$, with parameters \mbox{\texttt{\mdseries\slshape a}} and \mbox{\texttt{\mdseries\slshape b}}, i.e., a four\texttt{\symbol{45}}dimensional associative $F$\texttt{\symbol{45}}algebra with basis $(e,i,j,k)$ and multiplication defined by $e e = e$, $e i = i e = i$, $e j = j e = j$, $e k = k e = k$, $i i = \mbox{\texttt{\mdseries\slshape a}} e$, $i j = - j i = k$, $i k = - k i = \mbox{\texttt{\mdseries\slshape a}} j$, $j j = \mbox{\texttt{\mdseries\slshape b}} e$, $k j = - j k = \mbox{\texttt{\mdseries\slshape b}} i$, $k k = - \mbox{\texttt{\mdseries\slshape a}} \mbox{\texttt{\mdseries\slshape b}} e$. The default value for both \mbox{\texttt{\mdseries\slshape a}} and \mbox{\texttt{\mdseries\slshape b}} is $-1 \in F$. 

 The \texttt{GeneratorsOfAlgebra} (\ref{GeneratorsOfAlgebra}) and \texttt{CanonicalBasis} (\ref{CanonicalBasis}) value of an algebra constructed with \texttt{QuaternionAlgebra} is the list $[ e, i, j, k ]$. 

 Two quaternion algebras with the same parameters \mbox{\texttt{\mdseries\slshape a}}, \mbox{\texttt{\mdseries\slshape b}} lie in the same family, so it makes sense to consider their intersection or to
ask whether they are contained in each other. (This is due to the fact that
the results of \texttt{QuaternionAlgebra} are cached, in the global variable \texttt{QuaternionAlgebraData}.) 

 The embedding of the field \texttt{GaussianRationals} (\ref{GaussianRationals}) into a quaternion algebra $A$ over \texttt{Rationals} (\ref{Rationals}) is not uniquely determined. One can specify one embedding as a vector space
homomorphism that maps \texttt{1} to the first algebra generator of $A$, and \texttt{E(4)} to one of the others. 
\begin{Verbatim}[commandchars=!@|,fontsize=\small,frame=single,label=Example]
  !gapprompt@gap>| !gapinput@QuaternionAlgebra( Rationals );|
  <algebra-with-one of dimension 4 over Rationals>
\end{Verbatim}
 }

 

\subsection{\textcolor{Chapter }{ComplexificationQuat (for a vector)}}
\logpage{[ 62, 5, 2 ]}\nobreak
\hyperdef{L}{X7B807702782F56FF}{}
{\noindent\textcolor{FuncColor}{$\triangleright$\enspace\texttt{ComplexificationQuat({\mdseries\slshape vector})\index{ComplexificationQuat@\texttt{ComplexificationQuat}!for a vector}
\label{ComplexificationQuat:for a vector}
}\hfill{\scriptsize (function)}}\\
\noindent\textcolor{FuncColor}{$\triangleright$\enspace\texttt{ComplexificationQuat({\mdseries\slshape matrix})\index{ComplexificationQuat@\texttt{ComplexificationQuat}!for a matrix}
\label{ComplexificationQuat:for a matrix}
}\hfill{\scriptsize (function)}}\\


 Let $A = e F \oplus i F \oplus j F \oplus k F$ be a quaternion algebra over the field $F$ of cyclotomics, with basis $(e,i,j,k)$. 

 If $v = v_1 + v_2 j$ is a row vector over $A$ with $v_1 = e w_1 + i w_2$ and $v_2 = e w_3 + i w_4$ then \texttt{ComplexificationQuat} called with argument $v$ returns the concatenation of $w_1 + $\texttt{E(4)}$ w_2$ and $w_3 + $\texttt{E(4)}$ w_4$. 

 If $M = M_1 + M_2 j$ is a matrix over $A$ with $M_1 = e N_1 + i N_2$ and $M_2 = e N_3 + i N_4$ then \texttt{ComplexificationQuat} called with argument $M$ returns the block matrix $A$ over $e F \oplus i F$ such that $A(1,1) = N_1 + $\texttt{E(4)}$ N_2$, $A(2,2) = N_1 - $\texttt{E(4)}$ N_2$, $A(1,2) = N_3 + $\texttt{E(4)}$ N_4$, and $A(2,1) = - N_3 + $\texttt{E(4)}$ N_4$. 

 Then \texttt{ComplexificationQuat(\mbox{\texttt{\mdseries\slshape v}}) * ComplexificationQuat(\mbox{\texttt{\mdseries\slshape M}})= ComplexificationQuat(\mbox{\texttt{\mdseries\slshape v}} * \mbox{\texttt{\mdseries\slshape M}})}, since 
\[ v M = v_1 M_1 + v_2 j M_1 + v_1 M_2 j + v_2 j M_2 j = ( v_1 M_1 - v_2
\overline{{M_2}} ) + ( v_1 M_2 + v_2 \overline{{M_1}} ) j. \]
 }

 

\subsection{\textcolor{Chapter }{OctaveAlgebra}}
\logpage{[ 62, 5, 3 ]}\nobreak
\hyperdef{L}{X78C88A38853A8443}{}
{\noindent\textcolor{FuncColor}{$\triangleright$\enspace\texttt{OctaveAlgebra({\mdseries\slshape F})\index{OctaveAlgebra@\texttt{OctaveAlgebra}}
\label{OctaveAlgebra}
}\hfill{\scriptsize (function)}}\\


 The algebra of octonions over \mbox{\texttt{\mdseries\slshape F}}. 
\begin{Verbatim}[commandchars=!@|,fontsize=\small,frame=single,label=Example]
  !gapprompt@gap>| !gapinput@OctaveAlgebra( Rationals );|
  <algebra of dimension 8 over Rationals>
\end{Verbatim}
 }

 

\subsection{\textcolor{Chapter }{FullMatrixAlgebra}}
\logpage{[ 62, 5, 4 ]}\nobreak
\hyperdef{L}{X7D88E42B7DE087B0}{}
{\noindent\textcolor{FuncColor}{$\triangleright$\enspace\texttt{FullMatrixAlgebra({\mdseries\slshape R, n})\index{FullMatrixAlgebra@\texttt{FullMatrixAlgebra}}
\label{FullMatrixAlgebra}
}\hfill{\scriptsize (function)}}\\
\noindent\textcolor{FuncColor}{$\triangleright$\enspace\texttt{MatrixAlgebra({\mdseries\slshape R, n})\index{MatrixAlgebra@\texttt{MatrixAlgebra}}
\label{MatrixAlgebra}
}\hfill{\scriptsize (function)}}\\
\noindent\textcolor{FuncColor}{$\triangleright$\enspace\texttt{MatAlgebra({\mdseries\slshape R, n})\index{MatAlgebra@\texttt{MatAlgebra}}
\label{MatAlgebra}
}\hfill{\scriptsize (function)}}\\


 is the full matrix algebra of $\mbox{\texttt{\mdseries\slshape n}} \times \mbox{\texttt{\mdseries\slshape n}}$ matrices over the ring \mbox{\texttt{\mdseries\slshape R}}, for a nonnegative integer \mbox{\texttt{\mdseries\slshape n}}. 
\begin{Verbatim}[commandchars=!@|,fontsize=\small,frame=single,label=Example]
  !gapprompt@gap>| !gapinput@A:=FullMatrixAlgebra( Rationals, 20 );|
  ( Rationals^[ 20, 20 ] )
  !gapprompt@gap>| !gapinput@Dimension( A );|
  400
\end{Verbatim}
 }

 

\subsection{\textcolor{Chapter }{NullAlgebra}}
\logpage{[ 62, 5, 5 ]}\nobreak
\hyperdef{L}{X78B8BA77869DAA13}{}
{\noindent\textcolor{FuncColor}{$\triangleright$\enspace\texttt{NullAlgebra({\mdseries\slshape R})\index{NullAlgebra@\texttt{NullAlgebra}}
\label{NullAlgebra}
}\hfill{\scriptsize (attribute)}}\\


 The zero\texttt{\symbol{45}}dimensional algebra over \mbox{\texttt{\mdseries\slshape R}}.  
\begin{Verbatim}[commandchars=!@|,fontsize=\small,frame=single,label=Example]
  !gapprompt@gap>| !gapinput@A:= NullAlgebra( Rationals );|
  <algebra of dimension 0 over Rationals>
  !gapprompt@gap>| !gapinput@Dimension( A );|
  0
\end{Verbatim}
 }

 }

  
\section{\textcolor{Chapter }{Subalgebras}}\label{Subalgebras}
\logpage{[ 62, 6, 0 ]}
\hyperdef{L}{X7DF5989886BE611E}{}
{
  

\subsection{\textcolor{Chapter }{Subalgebra}}
\logpage{[ 62, 6, 1 ]}\nobreak
\hyperdef{L}{X8396643D7A49EEAD}{}
{\noindent\textcolor{FuncColor}{$\triangleright$\enspace\texttt{Subalgebra({\mdseries\slshape A, gens[, "basis"]})\index{Subalgebra@\texttt{Subalgebra}}
\label{Subalgebra}
}\hfill{\scriptsize (function)}}\\


 is the $F$\texttt{\symbol{45}}algebra generated by \mbox{\texttt{\mdseries\slshape gens}}, with parent algebra \mbox{\texttt{\mdseries\slshape A}}, where $F$ is the left acting domain of \mbox{\texttt{\mdseries\slshape A}}. 

 \emph{Note} that being a subalgebra of \mbox{\texttt{\mdseries\slshape A}} means to be an algebra, to be contained in \mbox{\texttt{\mdseries\slshape A}}, \emph{and} to have the same left acting domain as \mbox{\texttt{\mdseries\slshape A}}. 

 An optional argument \texttt{"basis"} may be added if it is known that the generators already form a basis of the
algebra. Then it is \emph{not} checked whether \mbox{\texttt{\mdseries\slshape gens}} really are linearly independent and whether all elements in \mbox{\texttt{\mdseries\slshape gens}} lie in \mbox{\texttt{\mdseries\slshape A}}. 
\begin{Verbatim}[commandchars=!@|,fontsize=\small,frame=single,label=Example]
  !gapprompt@gap>| !gapinput@m:= [ [ 0, 1, 2 ], [ 0, 0, 3], [ 0, 0, 0 ] ];;|
  !gapprompt@gap>| !gapinput@A:= Algebra( Rationals, [ m ] );|
  <algebra over Rationals, with 1 generator>
  !gapprompt@gap>| !gapinput@B:= Subalgebra( A, [ m^2 ] );|
  <algebra over Rationals, with 1 generator>
\end{Verbatim}
 }

 

\subsection{\textcolor{Chapter }{SubalgebraNC}}
\logpage{[ 62, 6, 2 ]}\nobreak
\hyperdef{L}{X7C6B08657BD836C3}{}
{\noindent\textcolor{FuncColor}{$\triangleright$\enspace\texttt{SubalgebraNC({\mdseries\slshape A, gens[, "basis"]})\index{SubalgebraNC@\texttt{SubalgebraNC}}
\label{SubalgebraNC}
}\hfill{\scriptsize (function)}}\\


 \texttt{SubalgebraNC} does the same as \texttt{Subalgebra} (\ref{Subalgebra}), except that it does not check whether all elements in \mbox{\texttt{\mdseries\slshape gens}} lie in \mbox{\texttt{\mdseries\slshape A}}. 
\begin{Verbatim}[commandchars=!@|,fontsize=\small,frame=single,label=Example]
  !gapprompt@gap>| !gapinput@m:= RandomMat( 3, 3 );;|
  !gapprompt@gap>| !gapinput@A:= Algebra( Rationals, [ m ] );|
  <algebra over Rationals, with 1 generator>
  !gapprompt@gap>| !gapinput@SubalgebraNC( A, [ IdentityMat( 3, 3 ) ], "basis" );|
  <algebra of dimension 1 over Rationals>
\end{Verbatim}
 }

 

\subsection{\textcolor{Chapter }{SubalgebraWithOne}}
\logpage{[ 62, 6, 3 ]}\nobreak
\hyperdef{L}{X83ECF489846F00B0}{}
{\noindent\textcolor{FuncColor}{$\triangleright$\enspace\texttt{SubalgebraWithOne({\mdseries\slshape A, gens[, "basis"]})\index{SubalgebraWithOne@\texttt{SubalgebraWithOne}}
\label{SubalgebraWithOne}
}\hfill{\scriptsize (function)}}\\


 is the algebra\texttt{\symbol{45}}with\texttt{\symbol{45}}one generated by \mbox{\texttt{\mdseries\slshape gens}}, with parent algebra \mbox{\texttt{\mdseries\slshape A}}. 

 The optional third argument, the string \texttt{"basis"}, may be added if it is known that the elements from \mbox{\texttt{\mdseries\slshape gens}} are linearly independent. Then it is \emph{not} checked whether \mbox{\texttt{\mdseries\slshape gens}} really are linearly independent and whether all elements in \mbox{\texttt{\mdseries\slshape gens}} lie in \mbox{\texttt{\mdseries\slshape A}}. 
\begin{Verbatim}[commandchars=!@|,fontsize=\small,frame=single,label=Example]
  !gapprompt@gap>| !gapinput@m:= [ [ 0, 1, 2 ], [ 0, 0, 3], [ 0, 0, 0 ] ];;|
  !gapprompt@gap>| !gapinput@A:= AlgebraWithOne( Rationals, [ m ] );|
  <algebra-with-one over Rationals, with 1 generator>
  !gapprompt@gap>| !gapinput@B1:= SubalgebraWithOne( A, [ m ] );;|
  !gapprompt@gap>| !gapinput@B2:= Subalgebra( A, [ m ] );;|
  !gapprompt@gap>| !gapinput@Dimension( B1 );|
  3
  !gapprompt@gap>| !gapinput@Dimension( B2 );|
  2
\end{Verbatim}
 }

 

\subsection{\textcolor{Chapter }{SubalgebraWithOneNC}}
\logpage{[ 62, 6, 4 ]}\nobreak
\hyperdef{L}{X7A11B177868E76AA}{}
{\noindent\textcolor{FuncColor}{$\triangleright$\enspace\texttt{SubalgebraWithOneNC({\mdseries\slshape A, gens[, "basis"]})\index{SubalgebraWithOneNC@\texttt{SubalgebraWithOneNC}}
\label{SubalgebraWithOneNC}
}\hfill{\scriptsize (function)}}\\


 \texttt{SubalgebraWithOneNC} does the same as \texttt{SubalgebraWithOne} (\ref{SubalgebraWithOne}), except that it does not check whether all elements in \mbox{\texttt{\mdseries\slshape gens}} lie in \mbox{\texttt{\mdseries\slshape A}}. 
\begin{Verbatim}[commandchars=!@|,fontsize=\small,frame=single,label=Example]
  !gapprompt@gap>| !gapinput@m:= RandomMat( 3, 3 );; A:= Algebra( Rationals, [ m ] );;|
  !gapprompt@gap>| !gapinput@SubalgebraWithOneNC( A, [ m ] );|
  <algebra-with-one over Rationals, with 1 generator>
\end{Verbatim}
 }

 

\subsection{\textcolor{Chapter }{TrivialSubalgebra}}
\logpage{[ 62, 6, 5 ]}\nobreak
\hyperdef{L}{X7FDD953A84CFC3D2}{}
{\noindent\textcolor{FuncColor}{$\triangleright$\enspace\texttt{TrivialSubalgebra({\mdseries\slshape A})\index{TrivialSubalgebra@\texttt{TrivialSubalgebra}}
\label{TrivialSubalgebra}
}\hfill{\scriptsize (attribute)}}\\


 The zero dimensional subalgebra of the algebra \mbox{\texttt{\mdseries\slshape A}}. 
\begin{Verbatim}[commandchars=!@|,fontsize=\small,frame=single,label=Example]
  !gapprompt@gap>| !gapinput@A:= QuaternionAlgebra( Rationals );;|
  !gapprompt@gap>| !gapinput@B:= TrivialSubalgebra( A );|
  <algebra of dimension 0 over Rationals>
  !gapprompt@gap>| !gapinput@Dimension( B );|
  0
\end{Verbatim}
 }

 }

  
\section{\textcolor{Chapter }{Ideals of Algebras}}\label{Ideals of Algebras}
\logpage{[ 62, 7, 0 ]}
\hyperdef{L}{X81EE8C1F7D7A7CF8}{}
{
  For constructing and working with ideals in algebras the same functions are
available as for ideals in rings. So for the precise description of these
functions we refer to Chapter \ref{Rings}. Here we give examples demonstrating the use of ideals in algebras. For an
introduction into the construction of quotient algebras we refer to Chapter  (\textbf{Tutorial: Algebras}) of the user's tutorial. 

 
\begin{Verbatim}[commandchars=!@|,fontsize=\small,frame=single,label=Example]
  !gapprompt@gap>| !gapinput@m:= [ [ 0, 2, 3 ], [ 0, 0, 4 ], [ 0, 0, 0] ];;|
  !gapprompt@gap>| !gapinput@A:= AlgebraWithOne( Rationals, [ m ] );;|
  !gapprompt@gap>| !gapinput@I:= Ideal( A, [ m ] );  # the two-sided ideal of `A' generated by `m'|
  <two-sided ideal in <algebra-with-one of dimension 3 over Rationals>,
    (1 generator)>
  !gapprompt@gap>| !gapinput@Dimension( I );|
  2
  !gapprompt@gap>| !gapinput@GeneratorsOfIdeal( I );|
  [ [ [ 0, 2, 3 ], [ 0, 0, 4 ], [ 0, 0, 0 ] ] ]
  !gapprompt@gap>| !gapinput@BasisVectors( Basis( I ) );|
  [ [ [ 0, 1, 3/2 ], [ 0, 0, 2 ], [ 0, 0, 0 ] ],
    [ [ 0, 0, 1 ], [ 0, 0, 0 ], [ 0, 0, 0 ] ] ]
  !gapprompt@gap>| !gapinput@A:= FullMatrixAlgebra( Rationals, 4 );;|
  !gapprompt@gap>| !gapinput@m:= NullMat( 4, 4 );; m[1][4]:=1;;|
  !gapprompt@gap>| !gapinput@I:= LeftIdeal( A, [ m ] );|
  <left ideal in ( Rationals^[ 4, 4 ] ), (1 generator)>
  !gapprompt@gap>| !gapinput@Dimension( I );|
  4
  !gapprompt@gap>| !gapinput@GeneratorsOfLeftIdeal( I );|
  [ [ [ 0, 0, 0, 1 ], [ 0, 0, 0, 0 ], [ 0, 0, 0, 0 ], [ 0, 0, 0, 0 ] ] ]
  !gapprompt@gap>| !gapinput@mats:= [ [[1,0],[0,0]], [[0,1],[0,0]], [[0,0],[0,1]] ];;|
  !gapprompt@gap>| !gapinput@A:= Algebra( Rationals, mats );;|
  !gapprompt@gap>| !gapinput@# Form the two-sided ideal for which `mats[2]' is known to be|
  !gapprompt@gap>| !gapinput@# the unique basis element.|
  !gapprompt@gap>| !gapinput@I:= Ideal( A, [ mats[2] ], "basis" );|
  <two-sided ideal in <algebra of dimension 3 over Rationals>,
    (dimension 1)>
\end{Verbatim}
 }

  
\section{\textcolor{Chapter }{Categories and Properties of Algebras}}\label{Categories and Properties of Algebras}
\logpage{[ 62, 8, 0 ]}
\hyperdef{L}{X7DC95D6982C9D7B0}{}
{
  

\subsection{\textcolor{Chapter }{IsFLMLOR}}
\logpage{[ 62, 8, 1 ]}\nobreak
\hyperdef{L}{X7FEDFAA383AB20D2}{}
{\noindent\textcolor{FuncColor}{$\triangleright$\enspace\texttt{IsFLMLOR({\mdseries\slshape obj})\index{IsFLMLOR@\texttt{IsFLMLOR}}
\label{IsFLMLOR}
}\hfill{\scriptsize (Category)}}\\


 A FLMLOR (``free left module left operator ring'') in \textsf{GAP} is a ring that is also a free left module. 

 Note that this means that being a FLMLOR is not a property a ring can get,
since a ring is usually not represented as an external left set. 

 Examples are magma rings (e.g. over the integers) or algebras. 
\begin{Verbatim}[commandchars=!@|,fontsize=\small,frame=single,label=Example]
  !gapprompt@gap>| !gapinput@A:= FullMatrixAlgebra( Rationals, 2 );;|
  !gapprompt@gap>| !gapinput@IsFLMLOR ( A );|
  true
\end{Verbatim}
 }

 

\subsection{\textcolor{Chapter }{IsFLMLORWithOne}}
\logpage{[ 62, 8, 2 ]}\nobreak
\hyperdef{L}{X85C1E13A877DF2C8}{}
{\noindent\textcolor{FuncColor}{$\triangleright$\enspace\texttt{IsFLMLORWithOne({\mdseries\slshape obj})\index{IsFLMLORWithOne@\texttt{IsFLMLORWithOne}}
\label{IsFLMLORWithOne}
}\hfill{\scriptsize (Category)}}\\


 A FLMLOR\texttt{\symbol{45}}with\texttt{\symbol{45}}one in \textsf{GAP} is a ring\texttt{\symbol{45}}with\texttt{\symbol{45}}one that is also a free
left module. 

 Note that this means that being a
FLMLOR\texttt{\symbol{45}}with\texttt{\symbol{45}}one is not a property a
ring\texttt{\symbol{45}}with\texttt{\symbol{45}}one can get, since a
ring\texttt{\symbol{45}}with\texttt{\symbol{45}}one is usually not represented
as an external left set. 

 Examples are magma rings\texttt{\symbol{45}}with\texttt{\symbol{45}}one or
algebras\texttt{\symbol{45}}with\texttt{\symbol{45}}one (but also over the
integers). 
\begin{Verbatim}[commandchars=!@|,fontsize=\small,frame=single,label=Example]
  !gapprompt@gap>| !gapinput@A:= FullMatrixAlgebra( Rationals, 2 );;|
  !gapprompt@gap>| !gapinput@IsFLMLORWithOne ( A );|
  true
\end{Verbatim}
 }

 

\subsection{\textcolor{Chapter }{IsAlgebra}}
\logpage{[ 62, 8, 3 ]}\nobreak
\hyperdef{L}{X801ED693808F6C84}{}
{\noindent\textcolor{FuncColor}{$\triangleright$\enspace\texttt{IsAlgebra({\mdseries\slshape obj})\index{IsAlgebra@\texttt{IsAlgebra}}
\label{IsAlgebra}
}\hfill{\scriptsize (Category)}}\\


 An algebra in \textsf{GAP} is a ring that is also a left vector space. Note that this means that being an
algebra is not a property a ring can get, since a ring is usually not
represented as an external left set. 
\begin{Verbatim}[commandchars=!@|,fontsize=\small,frame=single,label=Example]
  !gapprompt@gap>| !gapinput@A:= MatAlgebra( Rationals, 3 );;|
  !gapprompt@gap>| !gapinput@IsAlgebra( A );|
  true
\end{Verbatim}
 }

 

\subsection{\textcolor{Chapter }{IsAlgebraWithOne}}
\logpage{[ 62, 8, 4 ]}\nobreak
\hyperdef{L}{X80B21AC27DE6D068}{}
{\noindent\textcolor{FuncColor}{$\triangleright$\enspace\texttt{IsAlgebraWithOne({\mdseries\slshape obj})\index{IsAlgebraWithOne@\texttt{IsAlgebraWithOne}}
\label{IsAlgebraWithOne}
}\hfill{\scriptsize (Category)}}\\


 An algebra\texttt{\symbol{45}}with\texttt{\symbol{45}}one in \textsf{GAP} is a ring\texttt{\symbol{45}}with\texttt{\symbol{45}}one that is also a left
vector space. Note that this means that being an
algebra\texttt{\symbol{45}}with\texttt{\symbol{45}}one is not a property a
ring\texttt{\symbol{45}}with\texttt{\symbol{45}}one can get, since a
ring\texttt{\symbol{45}}with\texttt{\symbol{45}}one is usually not represented
as an external left set. 
\begin{Verbatim}[commandchars=!@|,fontsize=\small,frame=single,label=Example]
  !gapprompt@gap>| !gapinput@A:= MatAlgebra( Rationals, 3 );;|
  !gapprompt@gap>| !gapinput@IsAlgebraWithOne( A );|
  true
\end{Verbatim}
 }

 

\subsection{\textcolor{Chapter }{IsLieAlgebra}}
\logpage{[ 62, 8, 5 ]}\nobreak
\hyperdef{L}{X839BAC687B4E1A1D}{}
{\noindent\textcolor{FuncColor}{$\triangleright$\enspace\texttt{IsLieAlgebra({\mdseries\slshape A})\index{IsLieAlgebra@\texttt{IsLieAlgebra}}
\label{IsLieAlgebra}
}\hfill{\scriptsize (filter)}}\\


 An algebra \mbox{\texttt{\mdseries\slshape A}} is called Lie algebra if $a * a = 0$ for all $a$ in \mbox{\texttt{\mdseries\slshape A}} and $( a * ( b * c ) ) + ( b * ( c * a ) ) + ( c * ( a * b ) ) = 0$ for all $a, b, c \in $\mbox{\texttt{\mdseries\slshape A}} (Jacobi identity). 
\begin{Verbatim}[commandchars=!@|,fontsize=\small,frame=single,label=Example]
  !gapprompt@gap>| !gapinput@A:= FullMatrixLieAlgebra( Rationals, 3 );;|
  !gapprompt@gap>| !gapinput@IsLieAlgebra( A );|
  true
\end{Verbatim}
 }

 

\subsection{\textcolor{Chapter }{IsSimpleAlgebra}}
\logpage{[ 62, 8, 6 ]}\nobreak
\hyperdef{L}{X877DF13387831A6A}{}
{\noindent\textcolor{FuncColor}{$\triangleright$\enspace\texttt{IsSimpleAlgebra({\mdseries\slshape A})\index{IsSimpleAlgebra@\texttt{IsSimpleAlgebra}}
\label{IsSimpleAlgebra}
}\hfill{\scriptsize (property)}}\\


 is \texttt{true} if the algebra \mbox{\texttt{\mdseries\slshape A}} is simple, and \texttt{false} otherwise. This function is only implemented for the cases where \mbox{\texttt{\mdseries\slshape A}} is an associative or a Lie algebra. And for Lie algebras it is only
implemented for the case where the ground field is of characteristic zero. 
\begin{Verbatim}[commandchars=!@|,fontsize=\small,frame=single,label=Example]
  !gapprompt@gap>| !gapinput@A:= FullMatrixLieAlgebra( Rationals, 3 );;|
  !gapprompt@gap>| !gapinput@IsSimpleAlgebra( A );|
  false
  !gapprompt@gap>| !gapinput@A:= MatAlgebra( Rationals, 3 );;|
  !gapprompt@gap>| !gapinput@IsSimpleAlgebra( A );|
  true
\end{Verbatim}
 }

  

\subsection{\textcolor{Chapter }{IsFiniteDimensional (for matrix algebras)}}
\logpage{[ 62, 8, 7 ]}\nobreak
\hyperdef{L}{X7C5AECE87D79D075}{}
{\noindent\textcolor{FuncColor}{$\triangleright$\enspace\texttt{IsFiniteDimensional({\mdseries\slshape matalg})\index{IsFiniteDimensional@\texttt{IsFiniteDimensional}!for matrix algebras}
\label{IsFiniteDimensional:for matrix algebras}
}\hfill{\scriptsize (method)}}\\


 returns \texttt{true} (always) for a matrix algebra \mbox{\texttt{\mdseries\slshape matalg}}, since matrix algebras are always finite dimensional. 

 
\begin{Verbatim}[commandchars=!@|,fontsize=\small,frame=single,label=Example]
  !gapprompt@gap>| !gapinput@A:= MatAlgebra( Rationals, 3 );;|
  !gapprompt@gap>| !gapinput@IsFiniteDimensional( A );|
  true
\end{Verbatim}
 }

 

\subsection{\textcolor{Chapter }{IsQuaternion}}
\logpage{[ 62, 8, 8 ]}\nobreak
\hyperdef{L}{X82B3A9077D0CB453}{}
{\noindent\textcolor{FuncColor}{$\triangleright$\enspace\texttt{IsQuaternion({\mdseries\slshape obj})\index{IsQuaternion@\texttt{IsQuaternion}}
\label{IsQuaternion}
}\hfill{\scriptsize (Category)}}\\
\noindent\textcolor{FuncColor}{$\triangleright$\enspace\texttt{IsQuaternionCollection({\mdseries\slshape obj})\index{IsQuaternionCollection@\texttt{IsQuaternionCollection}}
\label{IsQuaternionCollection}
}\hfill{\scriptsize (Category)}}\\
\noindent\textcolor{FuncColor}{$\triangleright$\enspace\texttt{IsQuaternionCollColl({\mdseries\slshape obj})\index{IsQuaternionCollColl@\texttt{IsQuaternionCollColl}}
\label{IsQuaternionCollColl}
}\hfill{\scriptsize (Category)}}\\


 \texttt{IsQuaternion} is the category of elements in an algebra constructed by \texttt{QuaternionAlgebra} (\ref{QuaternionAlgebra}). A collection of quaternions lies in the category \texttt{IsQuaternionCollection}. Finally, a collection of quaternion collections (e.g., a matrix of
quaternions) lies in the category \texttt{IsQuaternionCollColl}. 
\begin{Verbatim}[commandchars=!@|,fontsize=\small,frame=single,label=Example]
  !gapprompt@gap>| !gapinput@A:= QuaternionAlgebra( Rationals );;|
  !gapprompt@gap>| !gapinput@b:= BasisVectors( Basis( A ) );|
  [ e, i, j, k ]
  !gapprompt@gap>| !gapinput@IsQuaternion( b[1] );|
  true
  !gapprompt@gap>| !gapinput@IsQuaternionCollColl( [ [ b[1], b[2] ], [ b[3], b[4] ] ] );|
  true
\end{Verbatim}
 }

 }

  
\section{\textcolor{Chapter }{Attributes and Operations for Algebras}}\label{Attributes and Operations for Algebras}
\logpage{[ 62, 9, 0 ]}
\hyperdef{L}{X7E9273E47CF38CF1}{}
{
    

\subsection{\textcolor{Chapter }{GeneratorsOfAlgebra}}
\logpage{[ 62, 9, 1 ]}\nobreak
\hyperdef{L}{X83B055F37EBF2438}{}
{\noindent\textcolor{FuncColor}{$\triangleright$\enspace\texttt{GeneratorsOfAlgebra({\mdseries\slshape A})\index{GeneratorsOfAlgebra@\texttt{GeneratorsOfAlgebra}}
\label{GeneratorsOfAlgebra}
}\hfill{\scriptsize (attribute)}}\\


 returns a list of elements that generate \mbox{\texttt{\mdseries\slshape A}} as an algebra. 

 For a free algebra, each generator can also be accessed using the \texttt{.} operator (see \texttt{GeneratorsOfDomain} (\ref{GeneratorsOfDomain})). 
\begin{Verbatim}[commandchars=!@|,fontsize=\small,frame=single,label=Example]
  !gapprompt@gap>| !gapinput@m:= [ [ 0, 1, 2 ], [ 0, 0, 3 ], [ 0, 0, 0 ] ];;|
  !gapprompt@gap>| !gapinput@A:= AlgebraWithOne( Rationals, [ m ] );|
  <algebra-with-one over Rationals, with 1 generator>
  !gapprompt@gap>| !gapinput@GeneratorsOfAlgebra( A );|
  [ [ [ 1, 0, 0 ], [ 0, 1, 0 ], [ 0, 0, 1 ] ],
    [ [ 0, 1, 2 ], [ 0, 0, 3 ], [ 0, 0, 0 ] ] ]
\end{Verbatim}
 }

 

\subsection{\textcolor{Chapter }{GeneratorsOfAlgebraWithOne}}
\logpage{[ 62, 9, 2 ]}\nobreak
\hyperdef{L}{X7FA408307A5A420E}{}
{\noindent\textcolor{FuncColor}{$\triangleright$\enspace\texttt{GeneratorsOfAlgebraWithOne({\mdseries\slshape A})\index{GeneratorsOfAlgebraWithOne@\texttt{GeneratorsOfAlgebraWithOne}}
\label{GeneratorsOfAlgebraWithOne}
}\hfill{\scriptsize (attribute)}}\\


 returns a list of elements of \mbox{\texttt{\mdseries\slshape A}} that generate \mbox{\texttt{\mdseries\slshape A}} as an algebra with one. 

 For a free algebra with one, each generator can also be accessed using the \texttt{.} operator (see \texttt{GeneratorsOfDomain} (\ref{GeneratorsOfDomain})). 
\begin{Verbatim}[commandchars=!@|,fontsize=\small,frame=single,label=Example]
  !gapprompt@gap>| !gapinput@m:= [ [ 0, 1, 2 ], [ 0, 0, 3 ], [ 0, 0, 0 ] ];;|
  !gapprompt@gap>| !gapinput@A:= AlgebraWithOne( Rationals, [ m ] );|
  <algebra-with-one over Rationals, with 1 generator>
  !gapprompt@gap>| !gapinput@GeneratorsOfAlgebraWithOne( A );|
  [ [ [ 0, 1, 2 ], [ 0, 0, 3 ], [ 0, 0, 0 ] ] ]
\end{Verbatim}
 }

 

\subsection{\textcolor{Chapter }{ProductSpace}}
\logpage{[ 62, 9, 3 ]}\nobreak
\hyperdef{L}{X7D309FD37D94B196}{}
{\noindent\textcolor{FuncColor}{$\triangleright$\enspace\texttt{ProductSpace({\mdseries\slshape U, V})\index{ProductSpace@\texttt{ProductSpace}}
\label{ProductSpace}
}\hfill{\scriptsize (operation)}}\\


 is the vector space $\langle u * v ; u \in U, v \in V \rangle$, where $U$ and $V$ are subspaces of the same algebra. 

 If $\mbox{\texttt{\mdseries\slshape U}} = \mbox{\texttt{\mdseries\slshape V}}$ is known to be an algebra then the product space is also an algebra, moreover
it is an ideal in \mbox{\texttt{\mdseries\slshape U}}. If \mbox{\texttt{\mdseries\slshape U}} and \mbox{\texttt{\mdseries\slshape V}} are known to be ideals in an algebra $A$ then the product space is known to be an algebra and an ideal in $A$. 
\begin{Verbatim}[commandchars=!@|,fontsize=\small,frame=single,label=Example]
  !gapprompt@gap>| !gapinput@A:= QuaternionAlgebra( Rationals );;|
  !gapprompt@gap>| !gapinput@b:= BasisVectors( Basis( A ) );;|
  !gapprompt@gap>| !gapinput@B:= Subalgebra( A, [ b[4] ] );|
  <algebra over Rationals, with 1 generator>
  !gapprompt@gap>| !gapinput@ProductSpace( A, B );|
  <vector space of dimension 4 over Rationals>
\end{Verbatim}
 }

 

\subsection{\textcolor{Chapter }{PowerSubalgebraSeries}}
\logpage{[ 62, 9, 4 ]}\nobreak
\hyperdef{L}{X875CD2B37EE9A8A2}{}
{\noindent\textcolor{FuncColor}{$\triangleright$\enspace\texttt{PowerSubalgebraSeries({\mdseries\slshape A})\index{PowerSubalgebraSeries@\texttt{PowerSubalgebraSeries}}
\label{PowerSubalgebraSeries}
}\hfill{\scriptsize (attribute)}}\\


 returns a list of subalgebras of \mbox{\texttt{\mdseries\slshape A}}, the first term of which is \mbox{\texttt{\mdseries\slshape A}}; and every next term is the product space of the previous term with itself. 
\begin{Verbatim}[commandchars=!@|,fontsize=\small,frame=single,label=Example]
  !gapprompt@gap>| !gapinput@A:= QuaternionAlgebra( Rationals );|
  <algebra-with-one of dimension 4 over Rationals>
  !gapprompt@gap>| !gapinput@PowerSubalgebraSeries( A );|
  [ <algebra-with-one of dimension 4 over Rationals> ]
\end{Verbatim}
 }

 

\subsection{\textcolor{Chapter }{AdjointBasis}}
\logpage{[ 62, 9, 5 ]}\nobreak
\hyperdef{L}{X788F4E6184E5C863}{}
{\noindent\textcolor{FuncColor}{$\triangleright$\enspace\texttt{AdjointBasis({\mdseries\slshape B})\index{AdjointBasis@\texttt{AdjointBasis}}
\label{AdjointBasis}
}\hfill{\scriptsize (attribute)}}\\


 The \emph{adjoint map} $ad(x)$ of an element $x$ in an $F$\texttt{\symbol{45}}algebra $A$ is the left multiplication by $x$. This map is $F$\texttt{\symbol{45}}linear and thus, w.r.t. the given basis \mbox{\texttt{\mdseries\slshape B}}$ = (x_1, x_2, \ldots, x_n)$ of $A$, $ad(x)$ can be represented by a matrix over $F$. Let $V$ denote the $F$\texttt{\symbol{45}}vector space of the matrices corresponding to $ad(x)$, for $x \in A$. Then \texttt{AdjointBasis} returns the basis of $V$ that consists of the matrices for $ad(x_1), \ldots, ad(x_n)$. 

 
\begin{Verbatim}[commandchars=!@|,fontsize=\small,frame=single,label=Example]
  !gapprompt@gap>| !gapinput@A:= QuaternionAlgebra( Rationals );;|
  !gapprompt@gap>| !gapinput@AdjointBasis( Basis( A ) );|
  Basis( <vector space over Rationals, with 4 generators>,
  [ [ [ 1, 0, 0, 0 ], [ 0, 1, 0, 0 ], [ 0, 0, 1, 0 ], [ 0, 0, 0, 1 ] ],
    [ [ 0, -1, 0, 0 ], [ 1, 0, 0, 0 ], [ 0, 0, 0, -1 ], [ 0, 0, 1, 0 ] ]
      ,
    [ [ 0, 0, -1, 0 ], [ 0, 0, 0, 1 ], [ 1, 0, 0, 0 ], [ 0, -1, 0, 0 ] ]
      ,
    [ [ 0, 0, 0, -1 ], [ 0, 0, -1, 0 ], [ 0, 1, 0, 0 ], [ 1, 0, 0, 0 ]
       ] ] )
\end{Verbatim}
 }

 

\subsection{\textcolor{Chapter }{IndicesOfAdjointBasis}}
\logpage{[ 62, 9, 6 ]}\nobreak
\hyperdef{L}{X800A410B8536E6DD}{}
{\noindent\textcolor{FuncColor}{$\triangleright$\enspace\texttt{IndicesOfAdjointBasis({\mdseries\slshape B})\index{IndicesOfAdjointBasis@\texttt{IndicesOfAdjointBasis}}
\label{IndicesOfAdjointBasis}
}\hfill{\scriptsize (attribute)}}\\


 Let \mbox{\texttt{\mdseries\slshape A}} be an algebra and let \mbox{\texttt{\mdseries\slshape B}} be the basis that is output by \texttt{AdjointBasis( Basis( \mbox{\texttt{\mdseries\slshape A}} ) )}. This function returns a list of indices. If $i$ is an index belonging to this list, then $ad x_i$ is a basis vector of the matrix space spanned by $ad A$, where $x_i$ is the $i$\texttt{\symbol{45}}th basis vector of the basis \mbox{\texttt{\mdseries\slshape B}}. 
\begin{Verbatim}[commandchars=!@|,fontsize=\small,frame=single,label=Example]
  !gapprompt@gap>| !gapinput@L:= FullMatrixLieAlgebra( Rationals, 3 );;|
  !gapprompt@gap>| !gapinput@B:= AdjointBasis( Basis( L ) );;|
  !gapprompt@gap>| !gapinput@IndicesOfAdjointBasis( B );|
  [ 1, 2, 3, 4, 5, 6, 7, 8 ]
\end{Verbatim}
 }

 

\subsection{\textcolor{Chapter }{AsAlgebra}}
\logpage{[ 62, 9, 7 ]}\nobreak
\hyperdef{L}{X7BA35CB28062D407}{}
{\noindent\textcolor{FuncColor}{$\triangleright$\enspace\texttt{AsAlgebra({\mdseries\slshape F, A})\index{AsAlgebra@\texttt{AsAlgebra}}
\label{AsAlgebra}
}\hfill{\scriptsize (operation)}}\\


 Returns the algebra over \mbox{\texttt{\mdseries\slshape F}} generated by \mbox{\texttt{\mdseries\slshape A}}. 
\begin{Verbatim}[commandchars=!@|,fontsize=\small,frame=single,label=Example]
  !gapprompt@gap>| !gapinput@V:= VectorSpace( Rationals, [ IdentityMat( 2 ) ] );;|
  !gapprompt@gap>| !gapinput@AsAlgebra( Rationals, V );|
  <algebra of dimension 1 over Rationals>
\end{Verbatim}
 }

 

\subsection{\textcolor{Chapter }{AsAlgebraWithOne}}
\logpage{[ 62, 9, 8 ]}\nobreak
\hyperdef{L}{X878323367D0B68EB}{}
{\noindent\textcolor{FuncColor}{$\triangleright$\enspace\texttt{AsAlgebraWithOne({\mdseries\slshape F, A})\index{AsAlgebraWithOne@\texttt{AsAlgebraWithOne}}
\label{AsAlgebraWithOne}
}\hfill{\scriptsize (operation)}}\\


 If the algebra \mbox{\texttt{\mdseries\slshape A}} has an identity, then it can be viewed as an algebra with one over \mbox{\texttt{\mdseries\slshape F}}. This function returns this algebra with one. 
\begin{Verbatim}[commandchars=!@|,fontsize=\small,frame=single,label=Example]
  !gapprompt@gap>| !gapinput@V:= VectorSpace( Rationals, [ IdentityMat( 2 ) ] );;|
  !gapprompt@gap>| !gapinput@A:= AsAlgebra( Rationals, V );;|
  !gapprompt@gap>| !gapinput@AsAlgebraWithOne( Rationals, A );|
  <algebra-with-one over Rationals, with 1 generator>
\end{Verbatim}
 }

 

\subsection{\textcolor{Chapter }{AsSubalgebra}}
\logpage{[ 62, 9, 9 ]}\nobreak
\hyperdef{L}{X7A922D26805AFF99}{}
{\noindent\textcolor{FuncColor}{$\triangleright$\enspace\texttt{AsSubalgebra({\mdseries\slshape A, B})\index{AsSubalgebra@\texttt{AsSubalgebra}}
\label{AsSubalgebra}
}\hfill{\scriptsize (operation)}}\\


 If all elements of the algebra \mbox{\texttt{\mdseries\slshape B}} happen to be contained in the algebra \mbox{\texttt{\mdseries\slshape A}}, then \mbox{\texttt{\mdseries\slshape B}} can be viewed as a subalgebra of \mbox{\texttt{\mdseries\slshape A}}. This function returns this subalgebra. 
\begin{Verbatim}[commandchars=!@|,fontsize=\small,frame=single,label=Example]
  !gapprompt@gap>| !gapinput@A:= FullMatrixAlgebra( Rationals, 2 );;|
  !gapprompt@gap>| !gapinput@V:= VectorSpace( Rationals, [ IdentityMat( 2 ) ] );;|
  !gapprompt@gap>| !gapinput@B:= AsAlgebra( Rationals, V );;|
  !gapprompt@gap>| !gapinput@BA:= AsSubalgebra( A, B );|
  <algebra of dimension 1 over Rationals>
\end{Verbatim}
 }

 

\subsection{\textcolor{Chapter }{AsSubalgebraWithOne}}
\logpage{[ 62, 9, 10 ]}\nobreak
\hyperdef{L}{X7B964BC37A975E48}{}
{\noindent\textcolor{FuncColor}{$\triangleright$\enspace\texttt{AsSubalgebraWithOne({\mdseries\slshape A, B})\index{AsSubalgebraWithOne@\texttt{AsSubalgebraWithOne}}
\label{AsSubalgebraWithOne}
}\hfill{\scriptsize (operation)}}\\


 If \mbox{\texttt{\mdseries\slshape B}} is an algebra with one, all elements of which happen to be contained in the
algebra with one \mbox{\texttt{\mdseries\slshape A}}, then \mbox{\texttt{\mdseries\slshape B}} can be viewed as a subalgebra with one of \mbox{\texttt{\mdseries\slshape A}}. This function returns this subalgebra with one. 
\begin{Verbatim}[commandchars=!@|,fontsize=\small,frame=single,label=Example]
  !gapprompt@gap>| !gapinput@A:= FullMatrixAlgebra( Rationals, 2 );;|
  !gapprompt@gap>| !gapinput@V:= VectorSpace( Rationals, [ IdentityMat( 2 ) ] );;|
  !gapprompt@gap>| !gapinput@B:= AsAlgebra( Rationals, V );;|
  !gapprompt@gap>| !gapinput@C:= AsAlgebraWithOne( Rationals, B );;|
  !gapprompt@gap>| !gapinput@AC:= AsSubalgebraWithOne( A, C );|
  <algebra-with-one over Rationals, with 1 generator>
\end{Verbatim}
 }

 

\subsection{\textcolor{Chapter }{MutableBasisOfClosureUnderAction}}
\logpage{[ 62, 9, 11 ]}\nobreak
\hyperdef{L}{X7C280DAC7F840B60}{}
{\noindent\textcolor{FuncColor}{$\triangleright$\enspace\texttt{MutableBasisOfClosureUnderAction({\mdseries\slshape F, Agens, from, init, opr, zero, maxdim})\index{MutableBasisOfClosureUnderAction@\texttt{MutableBasisOfClosureUnderAction}}
\label{MutableBasisOfClosureUnderAction}
}\hfill{\scriptsize (function)}}\\


 Let \mbox{\texttt{\mdseries\slshape F}} be a ring, \mbox{\texttt{\mdseries\slshape Agens}} a list of generators for an \mbox{\texttt{\mdseries\slshape F}}\texttt{\symbol{45}}algebra $A$, and \mbox{\texttt{\mdseries\slshape from}} one of \texttt{"left"}, \texttt{"right"}, \texttt{"both"}; this means that elements of $A$ act via multiplication from the respective side(s). \mbox{\texttt{\mdseries\slshape init}} must be a list of initial generating vectors, and \mbox{\texttt{\mdseries\slshape opr}} the operation (a function of two arguments). 

 \texttt{MutableBasisOfClosureUnderAction} returns a mutable basis of the \mbox{\texttt{\mdseries\slshape F}}\texttt{\symbol{45}}free left module generated by the vectors in \mbox{\texttt{\mdseries\slshape init}} and their images under the action of \mbox{\texttt{\mdseries\slshape Agens}} from the respective side(s). 

 \mbox{\texttt{\mdseries\slshape zero}} is the zero element of the desired module. \mbox{\texttt{\mdseries\slshape maxdim}} is an upper bound for the dimension of the closure; if no such upper bound is
known then the value of \mbox{\texttt{\mdseries\slshape maxdim}} must be \texttt{infinity} (\ref{infinity}). 

 \texttt{MutableBasisOfClosureUnderAction} can be used to compute a basis of an \emph{associative} algebra generated by the elements in \mbox{\texttt{\mdseries\slshape Agens}}. In this case \mbox{\texttt{\mdseries\slshape from}} may be \texttt{"left"} or \texttt{"right"}, \mbox{\texttt{\mdseries\slshape opr}} is the multiplication \texttt{*}, and \mbox{\texttt{\mdseries\slshape init}} is a list containing either the identity of the algebra or a list of algebra
generators. (Note that if the algebra has an identity then it is in general
not sufficient to take algebra\texttt{\symbol{45}}with\texttt{\symbol{45}}one
generators as \mbox{\texttt{\mdseries\slshape init}}, whereas of course \mbox{\texttt{\mdseries\slshape Agens}} need not contain the identity.) 

 (Note that bases of \emph{not} necessarily associative algebras can be computed using \texttt{MutableBasisOfNonassociativeAlgebra} (\ref{MutableBasisOfNonassociativeAlgebra}).) 

 Other applications of \texttt{MutableBasisOfClosureUnderAction} are the computations of bases for (left/ right/ two\texttt{\symbol{45}}sided)
ideals $I$ in an \emph{associative} algebra $A$ from ideal generators of $I$; in these cases \mbox{\texttt{\mdseries\slshape Agens}} is a list of algebra generators of $A$, \mbox{\texttt{\mdseries\slshape from}} denotes the appropriate side(s), \mbox{\texttt{\mdseries\slshape init}} is a list of ideal generators of $I$, and \mbox{\texttt{\mdseries\slshape opr}} is again \texttt{*}. 

 (Note that bases of ideals in \emph{not} necessarily associative algebras can be computed using \texttt{MutableBasisOfIdealInNonassociativeAlgebra} (\ref{MutableBasisOfIdealInNonassociativeAlgebra}).) 

 Finally, bases of right $A$\texttt{\symbol{45}}modules also can be computed using \texttt{MutableBasisOfClosureUnderAction}. The only difference to the ideal case is that \mbox{\texttt{\mdseries\slshape init}} is now a list of right module generators, and \mbox{\texttt{\mdseries\slshape opr}} is the operation of the module. 

  
\begin{Verbatim}[commandchars=!@|,fontsize=\small,frame=single,label=Example]
  !gapprompt@gap>| !gapinput@A:= QuaternionAlgebra( Rationals );;|
  !gapprompt@gap>| !gapinput@g:= GeneratorsOfAlgebra( A );;|
  !gapprompt@gap>| !gapinput@B:= MutableBasisOfClosureUnderAction( Rationals,|
  !gapprompt@>| !gapinput@                               g, "left", [ g[1] ], \*, Zero(A), 4 );|
  <mutable basis over Rationals, 4 vectors>
  !gapprompt@gap>| !gapinput@BasisVectors( B );|
  [ e, i, j, k ]
\end{Verbatim}
 }

 

\subsection{\textcolor{Chapter }{MutableBasisOfNonassociativeAlgebra}}
\logpage{[ 62, 9, 12 ]}\nobreak
\hyperdef{L}{X7BA1739D7F8B3A2B}{}
{\noindent\textcolor{FuncColor}{$\triangleright$\enspace\texttt{MutableBasisOfNonassociativeAlgebra({\mdseries\slshape F, Agens, zero, maxdim})\index{MutableBasisOfNonassociativeAlgebra@\texttt{MutableBasisOfNonassociativeAlgebra}}
\label{MutableBasisOfNonassociativeAlgebra}
}\hfill{\scriptsize (function)}}\\


 is a mutable basis of the (not necessarily associative) \mbox{\texttt{\mdseries\slshape F}}\texttt{\symbol{45}}algebra that is generated by \mbox{\texttt{\mdseries\slshape Agens}}, has zero element \mbox{\texttt{\mdseries\slshape zero}}, and has dimension at most \mbox{\texttt{\mdseries\slshape maxdim}}. If no finite bound for the dimension is known then \texttt{infinity} (\ref{infinity}) must be the value of \mbox{\texttt{\mdseries\slshape maxdim}}. 

 The difference to \texttt{MutableBasisOfClosureUnderAction} (\ref{MutableBasisOfClosureUnderAction}) is that in general it is not sufficient to multiply just with algebra
generators. (For special cases of nonassociative algebras, especially for Lie
algebras, multiplying with algebra generators suffices.) 
\begin{Verbatim}[commandchars=!@|,fontsize=\small,frame=single,label=Example]
  !gapprompt@gap>| !gapinput@L:= FullMatrixLieAlgebra( Rationals, 4 );;|
  !gapprompt@gap>| !gapinput@m1:= Random( L );;|
  !gapprompt@gap>| !gapinput@m2:= Random( L );;|
  !gapprompt@gap>| !gapinput@MutableBasisOfNonassociativeAlgebra( Rationals, [ m1, m2 ],|
  !gapprompt@>| !gapinput@Zero( L ), 16 );|
  <mutable basis over Rationals, 16 vectors>
\end{Verbatim}
 }

 

\subsection{\textcolor{Chapter }{MutableBasisOfIdealInNonassociativeAlgebra}}
\logpage{[ 62, 9, 13 ]}\nobreak
\hyperdef{L}{X8467B687823371F9}{}
{\noindent\textcolor{FuncColor}{$\triangleright$\enspace\texttt{MutableBasisOfIdealInNonassociativeAlgebra({\mdseries\slshape F, Vgens, Igens, zero, from, maxdim})\index{MutableBasisOfIdealInNonassociativeAlgebra@\texttt{Mutable}\-\texttt{Basis}\-\texttt{Of}\-\texttt{Ideal}\-\texttt{In}\-\texttt{Nonassociative}\-\texttt{Algebra}}
\label{MutableBasisOfIdealInNonassociativeAlgebra}
}\hfill{\scriptsize (function)}}\\


 is a mutable basis of the ideal generated by \mbox{\texttt{\mdseries\slshape Igens}} under the action of the (not necessarily associative) \mbox{\texttt{\mdseries\slshape F}}\texttt{\symbol{45}}algebra with vector space generators \mbox{\texttt{\mdseries\slshape Vgens}}. The zero element of the ideal is \mbox{\texttt{\mdseries\slshape zero}}, \mbox{\texttt{\mdseries\slshape from}} is one of \texttt{"left"}, \texttt{"right"}, \texttt{"both"} (with the same meaning as in \texttt{MutableBasisOfClosureUnderAction} (\ref{MutableBasisOfClosureUnderAction})), and \mbox{\texttt{\mdseries\slshape maxdim}} is a known upper bound on the dimension of the ideal; if no finite bound for
the dimension is known then \texttt{infinity} (\ref{infinity}) must be the value of \mbox{\texttt{\mdseries\slshape maxdim}}. 

 The difference to \texttt{MutableBasisOfClosureUnderAction} (\ref{MutableBasisOfClosureUnderAction}) is that in general it is not sufficient to multiply just with algebra
generators. (For special cases of nonassociative algebras, especially for Lie
algebras, multiplying with algebra generators suffices.) 
\begin{Verbatim}[commandchars=!@|,fontsize=\small,frame=single,label=Example]
  !gapprompt@gap>| !gapinput@mats:= [  [[ 1, 0 ], [ 0, -1 ]], [[0,1],[0,0]] ];;|
  !gapprompt@gap>| !gapinput@A:= Algebra( Rationals, mats );;|
  !gapprompt@gap>| !gapinput@basA:= BasisVectors( Basis( A ) );;|
  !gapprompt@gap>| !gapinput@B:= MutableBasisOfIdealInNonassociativeAlgebra( Rationals, basA,|
  !gapprompt@>| !gapinput@[ mats[2] ], 0*mats[1], "both", infinity );|
  <mutable basis over Rationals, 1 vector>
  !gapprompt@gap>| !gapinput@BasisVectors( B );|
  [ [ [ 0, 1 ], [ 0, 0 ] ] ]
\end{Verbatim}
 }

 

\subsection{\textcolor{Chapter }{DirectSumOfAlgebras (for two algebras)}}
\logpage{[ 62, 9, 14 ]}\nobreak
\hyperdef{L}{X7C591B7C7DEA7EEB}{}
{\noindent\textcolor{FuncColor}{$\triangleright$\enspace\texttt{DirectSumOfAlgebras({\mdseries\slshape A1, A2})\index{DirectSumOfAlgebras@\texttt{DirectSumOfAlgebras}!for two algebras}
\label{DirectSumOfAlgebras:for two algebras}
}\hfill{\scriptsize (operation)}}\\
\noindent\textcolor{FuncColor}{$\triangleright$\enspace\texttt{DirectSumOfAlgebras({\mdseries\slshape list})\index{DirectSumOfAlgebras@\texttt{DirectSumOfAlgebras}!for a list of algebras}
\label{DirectSumOfAlgebras:for a list of algebras}
}\hfill{\scriptsize (operation)}}\\


 is the direct sum of the two algebras \mbox{\texttt{\mdseries\slshape A1}} and \mbox{\texttt{\mdseries\slshape A2}} respectively of the algebras in the list \mbox{\texttt{\mdseries\slshape list}}. 

 If all involved algebras are associative algebras then the result is also
known to be associative. If all involved algebras are Lie algebras then the
result is also known to be a Lie algebra. 

 All involved algebras must have the same left acting domain. 

 The default case is that the result is a structure constants algebra. If all
involved algebras are matrix algebras, and either both are Lie algebras or
both are associative then the result is again a matrix algebra of the
appropriate type. 
\begin{Verbatim}[commandchars=!@|,fontsize=\small,frame=single,label=Example]
  !gapprompt@gap>| !gapinput@A:= QuaternionAlgebra( Rationals );;|
  !gapprompt@gap>| !gapinput@DirectSumOfAlgebras( [A, A, A] );|
  <algebra of dimension 12 over Rationals>
\end{Verbatim}
 }

 

\subsection{\textcolor{Chapter }{FullMatrixAlgebraCentralizer}}
\logpage{[ 62, 9, 15 ]}\nobreak
\hyperdef{L}{X7D0EB1437D3D9495}{}
{\noindent\textcolor{FuncColor}{$\triangleright$\enspace\texttt{FullMatrixAlgebraCentralizer({\mdseries\slshape F, lst})\index{FullMatrixAlgebraCentralizer@\texttt{FullMatrixAlgebraCentralizer}}
\label{FullMatrixAlgebraCentralizer}
}\hfill{\scriptsize (function)}}\\


 Let \mbox{\texttt{\mdseries\slshape lst}} be a nonempty list of square matrices of the same dimension $n$ with entries in the field \mbox{\texttt{\mdseries\slshape F}}. \texttt{FullMatrixAlgebraCentralizer} returns the (pointwise) centralizer of all matrices in \mbox{\texttt{\mdseries\slshape lst}}, inside the full matrix algebra of $n \times n$ matrices over \mbox{\texttt{\mdseries\slshape F}}. 

 
\begin{Verbatim}[commandchars=!@|,fontsize=\small,frame=single,label=Example]
  !gapprompt@gap>| !gapinput@A:= QuaternionAlgebra( Rationals );;|
  !gapprompt@gap>| !gapinput@b:= Basis( A );;|
  !gapprompt@gap>| !gapinput@mats:= List( BasisVectors( b ), x -> AdjointMatrix( b, x ) );;|
  !gapprompt@gap>| !gapinput@FullMatrixAlgebraCentralizer( Rationals, mats );|
  <algebra-with-one of dimension 4 over Rationals>
\end{Verbatim}
 }

 

\subsection{\textcolor{Chapter }{RadicalOfAlgebra}}
\logpage{[ 62, 9, 16 ]}\nobreak
\hyperdef{L}{X850C29907A509533}{}
{\noindent\textcolor{FuncColor}{$\triangleright$\enspace\texttt{RadicalOfAlgebra({\mdseries\slshape A})\index{RadicalOfAlgebra@\texttt{RadicalOfAlgebra}}
\label{RadicalOfAlgebra}
}\hfill{\scriptsize (attribute)}}\\


 is the maximal nilpotent ideal of \mbox{\texttt{\mdseries\slshape A}}, where \mbox{\texttt{\mdseries\slshape A}} is an associative algebra. 
\begin{Verbatim}[commandchars=!@|,fontsize=\small,frame=single,label=Example]
  !gapprompt@gap>| !gapinput@m:= [ [ 0, 1, 2 ], [ 0, 0, 3 ], [ 0, 0, 0 ] ];;|
  !gapprompt@gap>| !gapinput@A:= AlgebraWithOneByGenerators( Rationals, [ m ] );|
  <algebra-with-one over Rationals, with 1 generator>
  !gapprompt@gap>| !gapinput@RadicalOfAlgebra( A );|
  <algebra of dimension 2 over Rationals>
\end{Verbatim}
 }

 

\subsection{\textcolor{Chapter }{CentralIdempotentsOfAlgebra}}
\logpage{[ 62, 9, 17 ]}\nobreak
\hyperdef{L}{X82571785846CF05C}{}
{\noindent\textcolor{FuncColor}{$\triangleright$\enspace\texttt{CentralIdempotentsOfAlgebra({\mdseries\slshape A})\index{CentralIdempotentsOfAlgebra@\texttt{CentralIdempotentsOfAlgebra}}
\label{CentralIdempotentsOfAlgebra}
}\hfill{\scriptsize (attribute)}}\\


 For an associative algebra \mbox{\texttt{\mdseries\slshape A}}, this function returns a list of central primitive idempotents such that
their sum is the identity element of \mbox{\texttt{\mdseries\slshape A}}. Therefore \mbox{\texttt{\mdseries\slshape A}} is required to have an identity. 

 (This is a synonym of \texttt{CentralIdempotentsOfSemiring}.)  
\begin{Verbatim}[commandchars=!@|,fontsize=\small,frame=single,label=Example]
  !gapprompt@gap>| !gapinput@A:= QuaternionAlgebra( Rationals );;|
  !gapprompt@gap>| !gapinput@B:= DirectSumOfAlgebras( [A, A, A] );|
  <algebra of dimension 12 over Rationals>
  !gapprompt@gap>| !gapinput@CentralIdempotentsOfAlgebra( B );|
  [ v.9, v.5, v.1 ]
\end{Verbatim}
 }

 

\subsection{\textcolor{Chapter }{DirectSumDecomposition (for Lie algebras)}}
\logpage{[ 62, 9, 18 ]}\nobreak
\hyperdef{L}{X7CFB230582C26DAA}{}
{\noindent\textcolor{FuncColor}{$\triangleright$\enspace\texttt{DirectSumDecomposition({\mdseries\slshape L})\index{DirectSumDecomposition@\texttt{DirectSumDecomposition}!for Lie algebras}
\label{DirectSumDecomposition:for Lie algebras}
}\hfill{\scriptsize (attribute)}}\\


 This function calculates a list of ideals of the algebra \mbox{\texttt{\mdseries\slshape L}} such that \mbox{\texttt{\mdseries\slshape L}} is equal to their direct sum. Currently this is only implemented for
semisimple associative algebras, and for Lie algebras (semisimple or not). 
\begin{Verbatim}[commandchars=!@|,fontsize=\small,frame=single,label=Example]
  !gapprompt@gap>| !gapinput@G:= SymmetricGroup( 4 );;|
  !gapprompt@gap>| !gapinput@A:= GroupRing( Rationals, G );|
  <algebra-with-one over Rationals, with 2 generators>
  !gapprompt@gap>| !gapinput@dd:= DirectSumDecomposition( A );|
  [ <two-sided ideal in
        <algebra-with-one of dimension 24 over Rationals>,
        (1 generator)>,
    <two-sided ideal in
        <algebra-with-one of dimension 24 over Rationals>,
        (1 generator)>,
    <two-sided ideal in
        <algebra-with-one of dimension 24 over Rationals>,
        (1 generator)>,
    <two-sided ideal in
        <algebra-with-one of dimension 24 over Rationals>,
        (1 generator)>,
    <two-sided ideal in
        <algebra-with-one of dimension 24 over Rationals>,
        (1 generator)> ]
  !gapprompt@gap>| !gapinput@List( dd, Dimension );|
  [ 1, 1, 4, 9, 9 ]
\end{Verbatim}
 
\begin{Verbatim}[commandchars=!@|,fontsize=\small,frame=single,label=Example]
  !gapprompt@gap>| !gapinput@L:= FullMatrixLieAlgebra( Rationals, 5 );;|
  !gapprompt@gap>| !gapinput@DirectSumDecomposition( L );|
  [ <two-sided ideal in
        <two-sided ideal in <Lie algebra of dimension 25 over Rationals>
              , (dimension 1)>, (dimension 1)>,
    <two-sided ideal in
        <two-sided ideal in <Lie algebra of dimension 25 over Rationals>
              , (dimension 24)>, (dimension 24)> ]
\end{Verbatim}
 }

 

\subsection{\textcolor{Chapter }{LeviMalcevDecomposition (for Lie algebras)}}
\logpage{[ 62, 9, 19 ]}\nobreak
\hyperdef{L}{X85C58364833E014C}{}
{\noindent\textcolor{FuncColor}{$\triangleright$\enspace\texttt{LeviMalcevDecomposition({\mdseries\slshape L})\index{LeviMalcevDecomposition@\texttt{LeviMalcevDecomposition}!for Lie algebras}
\label{LeviMalcevDecomposition:for Lie algebras}
}\hfill{\scriptsize (attribute)}}\\


 A Levi\texttt{\symbol{45}}Malcev subalgebra of the algebra \mbox{\texttt{\mdseries\slshape L}} is a semisimple subalgebra complementary to the radical of \mbox{\texttt{\mdseries\slshape L}}. This function returns a list with two components. The first component is a
Levi\texttt{\symbol{45}}Malcev subalgebra, the second the radical. This
function is implemented for associative and Lie algebras. 
\begin{Verbatim}[commandchars=!@|,fontsize=\small,frame=single,label=Example]
  !gapprompt@gap>| !gapinput@m:= [ [ 1, 2, 0 ], [ 0, 1, 3 ], [ 0, 0, 1] ];;|
  !gapprompt@gap>| !gapinput@A:= Algebra( Rationals, [ m ] );;|
  !gapprompt@gap>| !gapinput@LeviMalcevDecomposition( A );|
  [ <algebra of dimension 1 over Rationals>,
    <algebra of dimension 2 over Rationals> ]
\end{Verbatim}
 
\begin{Verbatim}[commandchars=!@|,fontsize=\small,frame=single,label=Example]
  !gapprompt@gap>| !gapinput@L:= FullMatrixLieAlgebra( Rationals, 5 );;|
  !gapprompt@gap>| !gapinput@LeviMalcevDecomposition( L );|
  [ <Lie algebra of dimension 24 over Rationals>,
    <two-sided ideal in <Lie algebra of dimension 25 over Rationals>,
        (dimension 1)> ]
\end{Verbatim}
 }

 

\subsection{\textcolor{Chapter }{Grading}}
\logpage{[ 62, 9, 20 ]}\nobreak
\hyperdef{L}{X7DCA2568870A2D34}{}
{\noindent\textcolor{FuncColor}{$\triangleright$\enspace\texttt{Grading({\mdseries\slshape A})\index{Grading@\texttt{Grading}}
\label{Grading}
}\hfill{\scriptsize (attribute)}}\\


 Let $G$ be an Abelian group and $A$ an algebra. Then $A$ is said to be graded over $G$ if for every $g \in G$ there is a subspace $A_g$ of $A$ such that $A_g \cdot A_h \subset A_{{g+h}}$ for $g, h \in G$. In \textsf{GAP}{\nobreakspace}4 a \emph{grading} of an algebra is a record containing the following components. 
\begin{description}
\item[{\texttt{source}}]  the Abelian group over which the algebra is graded. 
\item[{\texttt{hom{\textunderscore}components}}]  a function assigning to each element from the source a subspace of the
algebra. 
\item[{\texttt{min{\textunderscore}degree}}]  in the case where the algebra is graded over the integers this is the minimum
number for which \texttt{hom{\textunderscore}components} returns a nonzero subspace. 
\item[{\texttt{max{\textunderscore}degree}}]  is analogous to \texttt{min{\textunderscore}degree}. 
\end{description}
 We note that there are no methods to compute a grading of an arbitrary
algebra; however some algebras get a natural grading when they are constructed
(see \texttt{JenningsLieAlgebra} (\ref{JenningsLieAlgebra}), \texttt{NilpotentQuotientOfFpLieAlgebra} (\ref{NilpotentQuotientOfFpLieAlgebra})). 

 We note also that these components may be not enough to handle the grading
efficiently, and another record component may be needed. For instance in a Lie
algebra $L$ constructed by \texttt{JenningsLieAlgebra} (\ref{JenningsLieAlgebra}), the length of the of the range \texttt{[ Grading(L)!.min{\textunderscore}degree ..
Grading(L)!.max{\textunderscore}degree ]} may be non\texttt{\symbol{45}}polynomial in the dimension of $L$. To handle efficiently this situation, an optional component can be used: 
\begin{description}
\item[{\texttt{non{\textunderscore}zero{\textunderscore}hom{\textunderscore}components}}]  the subset of \texttt{source} for which \texttt{hom{\textunderscore}components} returns a nonzero subspace. 
\end{description}
 
\begin{Verbatim}[commandchars=!@|,fontsize=\small,frame=single,label=Example]
  !gapprompt@gap>| !gapinput@G:= SmallGroup(3^6, 100 );|
  <pc group of size 729 with 6 generators>
  !gapprompt@gap>| !gapinput@L:= JenningsLieAlgebra( G );|
  <Lie algebra of dimension 6 over GF(3)>
  !gapprompt@gap>| !gapinput@g:= Grading( L );|
  rec( hom_components := function( d ) ... end, max_degree := 9,
    min_degree := 1, source := Integers )
  !gapprompt@gap>| !gapinput@g.hom_components( 3 );|
  <vector space over GF(3), with 1 generator>
  !gapprompt@gap>| !gapinput@g.hom_components( 14 );|
  <vector space of dimension 0 over GF(3)>
\end{Verbatim}
 }

 }

  
\section{\textcolor{Chapter }{Homomorphisms of Algebras}}\label{Homomorphisms of Algebras}
\logpage{[ 62, 10, 0 ]}
\hyperdef{L}{X7E94B857847F95C1}{}
{
  Algebra homomorphisms are vector space homomorphisms that preserve the
multiplication. So the default methods for vector space homomorphisms work,
and in fact there is not much use of the fact that source and range are
algebras, except that preimages and images are algebras (or even ideals) in
certain cases. 

\subsection{\textcolor{Chapter }{AlgebraGeneralMappingByImages}}
\logpage{[ 62, 10, 1 ]}\nobreak
\hyperdef{L}{X83CE798C7D39E368}{}
{\noindent\textcolor{FuncColor}{$\triangleright$\enspace\texttt{AlgebraGeneralMappingByImages({\mdseries\slshape A, B, gens, imgs})\index{AlgebraGeneralMappingByImages@\texttt{AlgebraGeneralMappingByImages}}
\label{AlgebraGeneralMappingByImages}
}\hfill{\scriptsize (operation)}}\\


 is a general mapping from the $F$\texttt{\symbol{45}}algebra \mbox{\texttt{\mdseries\slshape A}} to the $F$\texttt{\symbol{45}}algebra \mbox{\texttt{\mdseries\slshape B}}. This general mapping is defined by mapping the entries in the list \mbox{\texttt{\mdseries\slshape gens}} (elements of \mbox{\texttt{\mdseries\slshape A}}) to the entries in the list \mbox{\texttt{\mdseries\slshape imgs}} (elements of \mbox{\texttt{\mdseries\slshape B}}), and taking the $F$\texttt{\symbol{45}}linear and multiplicative closure. 

 \mbox{\texttt{\mdseries\slshape gens}} need not generate \mbox{\texttt{\mdseries\slshape A}} as an $F$\texttt{\symbol{45}}algebra, and if the specification does not define a linear
and multiplicative mapping then the result will be multivalued. Hence, in
general it is not a mapping. For constructing a linear map that is not
necessarily multiplicative, we refer to \texttt{LeftModuleHomomorphismByImages} (\ref{LeftModuleHomomorphismByImages}). 
\begin{Verbatim}[commandchars=!@|,fontsize=\small,frame=single,label=Example]
  !gapprompt@gap>| !gapinput@A:= QuaternionAlgebra( Rationals );;|
  !gapprompt@gap>| !gapinput@B:= FullMatrixAlgebra( Rationals, 2 );;|
  !gapprompt@gap>| !gapinput@bA:= BasisVectors( Basis( A ) );; bB:= BasisVectors( Basis( B ) );;|
  !gapprompt@gap>| !gapinput@f:= AlgebraGeneralMappingByImages( A, B, bA, bB );|
  [ e, i, j, k ] -> [ [ [ 1, 0 ], [ 0, 0 ] ], [ [ 0, 1 ], [ 0, 0 ] ],
    [ [ 0, 0 ], [ 1, 0 ] ], [ [ 0, 0 ], [ 0, 1 ] ] ]
  !gapprompt@gap>| !gapinput@Images( f, bA[1] );|
  <add. coset of <algebra over Rationals, with 16 generators>>
\end{Verbatim}
 }

 

\subsection{\textcolor{Chapter }{AlgebraHomomorphismByImages}}
\logpage{[ 62, 10, 2 ]}\nobreak
\hyperdef{L}{X7A7F97ED8608C882}{}
{\noindent\textcolor{FuncColor}{$\triangleright$\enspace\texttt{AlgebraHomomorphismByImages({\mdseries\slshape A, B, gens, imgs})\index{AlgebraHomomorphismByImages@\texttt{AlgebraHomomorphismByImages}}
\label{AlgebraHomomorphismByImages}
}\hfill{\scriptsize (function)}}\\


 \texttt{AlgebraHomomorphismByImages} returns the algebra homomorphism with source \mbox{\texttt{\mdseries\slshape A}} and range \mbox{\texttt{\mdseries\slshape B}} that is defined by mapping the list \mbox{\texttt{\mdseries\slshape gens}} of generators of \mbox{\texttt{\mdseries\slshape A}} to the list \mbox{\texttt{\mdseries\slshape imgs}} of images in \mbox{\texttt{\mdseries\slshape B}}. 

 If \mbox{\texttt{\mdseries\slshape gens}} does not generate \mbox{\texttt{\mdseries\slshape A}} or if the homomorphism does not exist (i.e., if mapping the generators
describes only a multi\texttt{\symbol{45}}valued mapping) then \texttt{fail} is returned. 

 One can avoid the checks by calling \texttt{AlgebraHomomorphismByImagesNC} (\ref{AlgebraHomomorphismByImagesNC}), and one can construct multi\texttt{\symbol{45}}valued mappings with \texttt{AlgebraGeneralMappingByImages} (\ref{AlgebraGeneralMappingByImages}). 
\begin{Verbatim}[commandchars=!@|,fontsize=\small,frame=single,label=Example]
  !gapprompt@gap>| !gapinput@T:= EmptySCTable( 2, 0 );;|
  !gapprompt@gap>| !gapinput@SetEntrySCTable( T, 1, 1, [1,1] ); SetEntrySCTable( T, 2, 2, [1,2] );|
  !gapprompt@gap>| !gapinput@A:= AlgebraByStructureConstants( Rationals, T );;|
  !gapprompt@gap>| !gapinput@m1:= NullMat( 2, 2 );; m1[1][1]:= 1;;|
  !gapprompt@gap>| !gapinput@m2:= NullMat( 2, 2 );; m2[2][2]:= 1;;|
  !gapprompt@gap>| !gapinput@B:= AlgebraByGenerators( Rationals, [ m1, m2 ] );;|
  !gapprompt@gap>| !gapinput@bA:= BasisVectors( Basis( A ) );; bB:= BasisVectors( Basis( B ) );;|
  !gapprompt@gap>| !gapinput@f:= AlgebraHomomorphismByImages( A, B, bA, bB );|
  [ v.1, v.2 ] -> [ [ [ 1, 0 ], [ 0, 0 ] ], [ [ 0, 0 ], [ 0, 1 ] ] ]
  !gapprompt@gap>| !gapinput@Image( f, bA[1]+bA[2] );|
  [ [ 1, 0 ], [ 0, 1 ] ]
\end{Verbatim}
 }

 

\subsection{\textcolor{Chapter }{AlgebraHomomorphismByImagesNC}}
\logpage{[ 62, 10, 3 ]}\nobreak
\hyperdef{L}{X8326D1BD79725462}{}
{\noindent\textcolor{FuncColor}{$\triangleright$\enspace\texttt{AlgebraHomomorphismByImagesNC({\mdseries\slshape A, B, gens, imgs})\index{AlgebraHomomorphismByImagesNC@\texttt{AlgebraHomomorphismByImagesNC}}
\label{AlgebraHomomorphismByImagesNC}
}\hfill{\scriptsize (operation)}}\\


 \texttt{AlgebraHomomorphismByImagesNC} is the operation that is called by the function \texttt{AlgebraHomomorphismByImages} (\ref{AlgebraHomomorphismByImages}). Its methods may assume that \mbox{\texttt{\mdseries\slshape gens}} generates \mbox{\texttt{\mdseries\slshape A}} and that the mapping of \mbox{\texttt{\mdseries\slshape gens}} to \mbox{\texttt{\mdseries\slshape imgs}} defines an algebra homomorphism. Results are unpredictable if these conditions
do not hold. 

 For creating a possibly multi\texttt{\symbol{45}}valued mapping from \mbox{\texttt{\mdseries\slshape A}} to \mbox{\texttt{\mdseries\slshape B}} that respects addition, multiplication, and scalar multiplication, \texttt{AlgebraGeneralMappingByImages} (\ref{AlgebraGeneralMappingByImages}) can be used.  

 For the definitions of the algebras \texttt{A} and \texttt{B} in the next example we refer to the previous example. 

 
\begin{Verbatim}[commandchars=!@|,fontsize=\small,frame=single,label=Example]
  !gapprompt@gap>| !gapinput@f:= AlgebraHomomorphismByImagesNC( A, B, bA, bB );|
  [ v.1, v.2 ] -> [ [ [ 1, 0 ], [ 0, 0 ] ], [ [ 0, 0 ], [ 0, 1 ] ] ]
\end{Verbatim}
 }

 

\subsection{\textcolor{Chapter }{AlgebraWithOneGeneralMappingByImages}}
\logpage{[ 62, 10, 4 ]}\nobreak
\hyperdef{L}{X8057E55B864567AD}{}
{\noindent\textcolor{FuncColor}{$\triangleright$\enspace\texttt{AlgebraWithOneGeneralMappingByImages({\mdseries\slshape A, B, gens, imgs})\index{AlgebraWithOneGeneralMappingByImages@\texttt{Algebra}\-\texttt{With}\-\texttt{One}\-\texttt{General}\-\texttt{Mapping}\-\texttt{By}\-\texttt{Images}}
\label{AlgebraWithOneGeneralMappingByImages}
}\hfill{\scriptsize (operation)}}\\


 This function is analogous to \texttt{AlgebraGeneralMappingByImages} (\ref{AlgebraGeneralMappingByImages}); the only difference being that the identity of \mbox{\texttt{\mdseries\slshape A}} is automatically mapped to the identity of \mbox{\texttt{\mdseries\slshape B}}. 
\begin{Verbatim}[commandchars=!@|,fontsize=\small,frame=single,label=Example]
  !gapprompt@gap>| !gapinput@A:= QuaternionAlgebra( Rationals );;|
  !gapprompt@gap>| !gapinput@B:= FullMatrixAlgebra( Rationals, 2 );;|
  !gapprompt@gap>| !gapinput@bA:= BasisVectors( Basis( A ) );; bB:= BasisVectors( Basis( B ) );;|
  !gapprompt@gap>| !gapinput@f:=AlgebraWithOneGeneralMappingByImages(A,B,bA{[2,3,4]},bB{[1,2,3]});|
  [ i, j, k, e ] -> [ [ [ 1, 0 ], [ 0, 0 ] ], [ [ 0, 1 ], [ 0, 0 ] ],
    [ [ 0, 0 ], [ 1, 0 ] ], [ [ 1, 0 ], [ 0, 1 ] ] ]
\end{Verbatim}
 }

 

\subsection{\textcolor{Chapter }{AlgebraWithOneHomomorphismByImages}}
\logpage{[ 62, 10, 5 ]}\nobreak
\hyperdef{L}{X866F32B5846E5857}{}
{\noindent\textcolor{FuncColor}{$\triangleright$\enspace\texttt{AlgebraWithOneHomomorphismByImages({\mdseries\slshape A, B, gens, imgs})\index{AlgebraWithOneHomomorphismByImages@\texttt{AlgebraWithOneHomomorphismByImages}}
\label{AlgebraWithOneHomomorphismByImages}
}\hfill{\scriptsize (function)}}\\


 \texttt{AlgebraWithOneHomomorphismByImages} returns the algebra\texttt{\symbol{45}}with\texttt{\symbol{45}}one
homomorphism with source \mbox{\texttt{\mdseries\slshape A}} and range \mbox{\texttt{\mdseries\slshape B}} that is defined by mapping the list \mbox{\texttt{\mdseries\slshape gens}} of generators of \mbox{\texttt{\mdseries\slshape A}} to the list \mbox{\texttt{\mdseries\slshape imgs}} of images in \mbox{\texttt{\mdseries\slshape B}}. 

 The difference between an algebra homomorphism and an
algebra\texttt{\symbol{45}}with\texttt{\symbol{45}}one homomorphism is that in
the latter case, it is assumed that the identity of \mbox{\texttt{\mdseries\slshape A}} is mapped to the identity of \mbox{\texttt{\mdseries\slshape B}}, and therefore \mbox{\texttt{\mdseries\slshape gens}} needs to generate \mbox{\texttt{\mdseries\slshape A}} only as an algebra\texttt{\symbol{45}}with\texttt{\symbol{45}}one. 

 If \mbox{\texttt{\mdseries\slshape gens}} does not generate \mbox{\texttt{\mdseries\slshape A}} or if the homomorphism does not exist (i.e., if mapping the generators
describes only a multi\texttt{\symbol{45}}valued mapping) then \texttt{fail} is returned. 

 One can avoid the checks by calling \texttt{AlgebraWithOneHomomorphismByImagesNC} (\ref{AlgebraWithOneHomomorphismByImagesNC}), and one can construct multi\texttt{\symbol{45}}valued mappings with \texttt{AlgebraWithOneGeneralMappingByImages} (\ref{AlgebraWithOneGeneralMappingByImages}). 
\begin{Verbatim}[commandchars=!@|,fontsize=\small,frame=single,label=Example]
  !gapprompt@gap>| !gapinput@m1:= NullMat( 2, 2 );; m1[1][1]:=1;;|
  !gapprompt@gap>| !gapinput@m2:= NullMat( 2, 2 );; m2[2][2]:=1;;|
  !gapprompt@gap>| !gapinput@A:= AlgebraByGenerators( Rationals, [m1,m2] );;|
  !gapprompt@gap>| !gapinput@T:= EmptySCTable( 2, 0 );;|
  !gapprompt@gap>| !gapinput@SetEntrySCTable( T, 1, 1, [1,1] );|
  !gapprompt@gap>| !gapinput@SetEntrySCTable( T, 2, 2, [1,2] );|
  !gapprompt@gap>| !gapinput@B:= AlgebraByStructureConstants(Rationals, T);;|
  !gapprompt@gap>| !gapinput@bA:= BasisVectors( Basis( A ) );; bB:= BasisVectors( Basis( B ) );;|
  !gapprompt@gap>| !gapinput@f:= AlgebraWithOneHomomorphismByImages( A, B, bA{[1]}, bB{[1]} );|
  [ [ [ 1, 0 ], [ 0, 0 ] ], [ [ 1, 0 ], [ 0, 1 ] ] ] -> [ v.1, v.1+v.2 ]
\end{Verbatim}
 }

 

\subsection{\textcolor{Chapter }{AlgebraWithOneHomomorphismByImagesNC}}
\logpage{[ 62, 10, 6 ]}\nobreak
\hyperdef{L}{X80BF4D6A7FDC959A}{}
{\noindent\textcolor{FuncColor}{$\triangleright$\enspace\texttt{AlgebraWithOneHomomorphismByImagesNC({\mdseries\slshape A, B, gens, imgs})\index{AlgebraWithOneHomomorphismByImagesNC@\texttt{Algebra}\-\texttt{With}\-\texttt{One}\-\texttt{Homomorphism}\-\texttt{By}\-\texttt{ImagesNC}}
\label{AlgebraWithOneHomomorphismByImagesNC}
}\hfill{\scriptsize (operation)}}\\


 \texttt{AlgebraWithOneHomomorphismByImagesNC} is the operation that is called by the function \texttt{AlgebraWithOneHomomorphismByImages} (\ref{AlgebraWithOneHomomorphismByImages}). Its methods may assume that \mbox{\texttt{\mdseries\slshape gens}} generates \mbox{\texttt{\mdseries\slshape A}} and that the mapping of \mbox{\texttt{\mdseries\slshape gens}} to \mbox{\texttt{\mdseries\slshape imgs}} defines an algebra\texttt{\symbol{45}}with\texttt{\symbol{45}}one
homomorphism. Results are unpredictable if these conditions do not hold. 

 For creating a possibly multi\texttt{\symbol{45}}valued mapping from \mbox{\texttt{\mdseries\slshape A}} to \mbox{\texttt{\mdseries\slshape B}} that respects addition, multiplication, identity, and scalar multiplication, \texttt{AlgebraWithOneGeneralMappingByImages} (\ref{AlgebraWithOneGeneralMappingByImages}) can be used. 

  
\begin{Verbatim}[commandchars=!@|,fontsize=\small,frame=single,label=Example]
  !gapprompt@gap>| !gapinput@m1:= NullMat( 2, 2 );; m1[1][1]:=1;;|
  !gapprompt@gap>| !gapinput@m2:= NullMat( 2, 2 );; m2[2][2]:=1;;|
  !gapprompt@gap>| !gapinput@A:= AlgebraByGenerators( Rationals, [m1,m2] );;|
  !gapprompt@gap>| !gapinput@T:= EmptySCTable( 2, 0 );;|
  !gapprompt@gap>| !gapinput@SetEntrySCTable( T, 1, 1, [1,1] );|
  !gapprompt@gap>| !gapinput@SetEntrySCTable( T, 2, 2, [1,2] );|
  !gapprompt@gap>| !gapinput@B:= AlgebraByStructureConstants( Rationals, T);;|
  !gapprompt@gap>| !gapinput@bA:= BasisVectors( Basis( A ) );; bB:= BasisVectors( Basis( B ) );;|
  !gapprompt@gap>| !gapinput@f:= AlgebraWithOneHomomorphismByImagesNC( A, B, bA{[1]}, bB{[1]} );|
  [ [ [ 1, 0 ], [ 0, 0 ] ], [ [ 1, 0 ], [ 0, 1 ] ] ] -> [ v.1, v.1+v.2 ]
\end{Verbatim}
 }

 

\subsection{\textcolor{Chapter }{AlgebraHomomorphismByFunction}}
\logpage{[ 62, 10, 7 ]}\nobreak
\hyperdef{L}{X825149467C57DEFC}{}
{\noindent\textcolor{FuncColor}{$\triangleright$\enspace\texttt{AlgebraHomomorphismByFunction({\mdseries\slshape A, B, f})\index{AlgebraHomomorphismByFunction@\texttt{AlgebraHomomorphismByFunction}}
\label{AlgebraHomomorphismByFunction}
}\hfill{\scriptsize (operation)}}\\
\noindent\textcolor{FuncColor}{$\triangleright$\enspace\texttt{AlgebraWithOneHomomorphismByFunction({\mdseries\slshape A, B, f})\index{AlgebraWithOneHomomorphismByFunction@\texttt{Algebra}\-\texttt{With}\-\texttt{One}\-\texttt{Homomorphism}\-\texttt{By}\-\texttt{Function}}
\label{AlgebraWithOneHomomorphismByFunction}
}\hfill{\scriptsize (operation)}}\\


 These functions construct an algebra homomorphism from the algebra \mbox{\texttt{\mdseries\slshape A}} to the algebra \mbox{\texttt{\mdseries\slshape B}} using a one\texttt{\symbol{45}}argument function \mbox{\texttt{\mdseries\slshape f}}. They do not check that the function actually defines a homomorphism. 
\begin{Verbatim}[commandchars=!@|,fontsize=\small,frame=single,label=Example]
  !gapprompt@gap>| !gapinput@A := MatrixAlgebra( Rationals, 2 );;|
  !gapprompt@gap>| !gapinput@f := AlgebraHomomorphismByFunction( Rationals, A, q->[[q,0],[0,0]] );|
  MappingByFunction( Rationals, ( Rationals^[ 2, 2 ] ), function( q ) ... end )
  !gapprompt@gap>| !gapinput@11^f;|
  [ [ 11, 0 ], [ 0, 0 ] ]
\end{Verbatim}
 }

 

\subsection{\textcolor{Chapter }{NaturalHomomorphismByIdeal (for an algebra and an ideal)}}
\logpage{[ 62, 10, 8 ]}\nobreak
\hyperdef{L}{X8712E5C1861CC32C}{}
{\noindent\textcolor{FuncColor}{$\triangleright$\enspace\texttt{NaturalHomomorphismByIdeal({\mdseries\slshape A, I})\index{NaturalHomomorphismByIdeal@\texttt{NaturalHomomorphismByIdeal}!for an algebra and an ideal}
\label{NaturalHomomorphismByIdeal:for an algebra and an ideal}
}\hfill{\scriptsize (method)}}\\


 For an algebra \mbox{\texttt{\mdseries\slshape A}} and an ideal \mbox{\texttt{\mdseries\slshape I}} in \mbox{\texttt{\mdseries\slshape A}}, the return value of \texttt{NaturalHomomorphismByIdeal} (\ref{NaturalHomomorphismByIdeal}) is a homomorphism of algebras, in particular the range of this mapping is also
an algebra. 

 
\begin{Verbatim}[commandchars=!@|,fontsize=\small,frame=single,label=Example]
  !gapprompt@gap>| !gapinput@L:= FullMatrixLieAlgebra( Rationals, 3 );;|
  !gapprompt@gap>| !gapinput@C:= LieCentre( L );|
  <two-sided ideal in <Lie algebra of dimension 9 over Rationals>,
    (dimension 1)>
  !gapprompt@gap>| !gapinput@hom:= NaturalHomomorphismByIdeal( L, C );|
  <linear mapping by matrix, <Lie algebra of dimension
  9 over Rationals> -> <Lie algebra of dimension 8 over Rationals>>
  !gapprompt@gap>| !gapinput@ImagesSource( hom );|
  <Lie algebra of dimension 8 over Rationals>
\end{Verbatim}
 }

 

\subsection{\textcolor{Chapter }{OperationAlgebraHomomorphism (action w.r.t. a basis of the module)}}
\logpage{[ 62, 10, 9 ]}\nobreak
\hyperdef{L}{X8705A9C68102FEA3}{}
{\noindent\textcolor{FuncColor}{$\triangleright$\enspace\texttt{OperationAlgebraHomomorphism({\mdseries\slshape A, B[, opr]})\index{OperationAlgebraHomomorphism@\texttt{OperationAlgebraHomomorphism}!action w.r.t. a basis of the module}
\label{OperationAlgebraHomomorphism:action w.r.t. a basis of the module}
}\hfill{\scriptsize (operation)}}\\
\noindent\textcolor{FuncColor}{$\triangleright$\enspace\texttt{OperationAlgebraHomomorphism({\mdseries\slshape A, V[, opr]})\index{OperationAlgebraHomomorphism@\texttt{OperationAlgebraHomomorphism}!action on a free left module}
\label{OperationAlgebraHomomorphism:action on a free left module}
}\hfill{\scriptsize (operation)}}\\


 \texttt{OperationAlgebraHomomorphism} returns an algebra homomorphism from the $F$\texttt{\symbol{45}}algebra \mbox{\texttt{\mdseries\slshape A}} into a matrix algebra over $F$ that describes the $F$\texttt{\symbol{45}}linear action of \mbox{\texttt{\mdseries\slshape A}} on the basis \mbox{\texttt{\mdseries\slshape B}} of a free left module respectively on the free left module \mbox{\texttt{\mdseries\slshape V}} (in which case some basis of \mbox{\texttt{\mdseries\slshape V}} is chosen), via the operation \mbox{\texttt{\mdseries\slshape opr}}. 

 The homomorphism need not be surjective. The default value for \mbox{\texttt{\mdseries\slshape opr}} is \texttt{OnRight} (\ref{OnRight}). 

 If \mbox{\texttt{\mdseries\slshape A}} is an algebra\texttt{\symbol{45}}with\texttt{\symbol{45}}one then the
operation homomorphism is an
algebra\texttt{\symbol{45}}with\texttt{\symbol{45}}one homomorphism because
the identity of \mbox{\texttt{\mdseries\slshape A}} must act as the identity. 

  
\begin{Verbatim}[commandchars=!@|,fontsize=\small,frame=single,label=Example]
  !gapprompt@gap>| !gapinput@m1:= NullMat( 2, 2 );; m1[1][1]:= 1;;|
  !gapprompt@gap>| !gapinput@m2:= NullMat( 2, 2 );; m2[2][2]:= 1;;|
  !gapprompt@gap>| !gapinput@B:= AlgebraByGenerators( Rationals, [ m1, m2 ] );;|
  !gapprompt@gap>| !gapinput@V:= FullRowSpace( Rationals, 2 );|
  ( Rationals^2 )
  !gapprompt@gap>| !gapinput@f:=OperationAlgebraHomomorphism( B, Basis( V ), OnRight );|
  <op. hom. Algebra( Rationals,
  [ [ [ 1, 0 ], [ 0, 0 ] ], [ [ 0, 0 ], [ 0, 1 ] ]
   ] ) -> matrices of dim. 2>
  !gapprompt@gap>| !gapinput@Image( f, m1 );|
  [ [ 1, 0 ], [ 0, 0 ] ]
\end{Verbatim}
 }

 

\subsection{\textcolor{Chapter }{NiceAlgebraMonomorphism}}
\logpage{[ 62, 10, 10 ]}\nobreak
\hyperdef{L}{X7B249E8E86D895F0}{}
{\noindent\textcolor{FuncColor}{$\triangleright$\enspace\texttt{NiceAlgebraMonomorphism({\mdseries\slshape A})\index{NiceAlgebraMonomorphism@\texttt{NiceAlgebraMonomorphism}}
\label{NiceAlgebraMonomorphism}
}\hfill{\scriptsize (attribute)}}\\


 If \mbox{\texttt{\mdseries\slshape A}} is an associative algebra with one, returns an isomorphism from \mbox{\texttt{\mdseries\slshape A}} onto a matrix algebra (see \texttt{IsomorphismMatrixAlgebra} (\ref{IsomorphismMatrixAlgebra}) for an example). If \mbox{\texttt{\mdseries\slshape A}} is a finitely presented Lie algebra, returns an isomorphism from \mbox{\texttt{\mdseries\slshape A}} onto a Lie algebra defined by a structure constants table (see \ref{Finitely Presented Lie Algebras} for an example). }

 

\subsection{\textcolor{Chapter }{IsomorphismFpAlgebra}}
\logpage{[ 62, 10, 11 ]}\nobreak
\hyperdef{L}{X79D770777D873F80}{}
{\noindent\textcolor{FuncColor}{$\triangleright$\enspace\texttt{IsomorphismFpAlgebra({\mdseries\slshape A})\index{IsomorphismFpAlgebra@\texttt{IsomorphismFpAlgebra}}
\label{IsomorphismFpAlgebra}
}\hfill{\scriptsize (attribute)}}\\


 isomorphism from the algebra \mbox{\texttt{\mdseries\slshape A}} onto a finitely presented algebra. Currently this is only implemented for
associative algebras with one. 
\begin{Verbatim}[commandchars=!@|,fontsize=\small,frame=single,label=Example]
  !gapprompt@gap>| !gapinput@A:= QuaternionAlgebra( Rationals );|
  <algebra-with-one of dimension 4 over Rationals>
  !gapprompt@gap>| !gapinput@f:= IsomorphismFpAlgebra( A );|
  [ e, i, j, k, e ] -> [ [(1)*x.1], [(1)*x.2], [(1)*x.3], [(1)*x.4],
    [(1)*<identity ...>] ]
\end{Verbatim}
 }

 

\subsection{\textcolor{Chapter }{IsomorphismMatrixAlgebra}}
\logpage{[ 62, 10, 12 ]}\nobreak
\hyperdef{L}{X7FB760F9813B0789}{}
{\noindent\textcolor{FuncColor}{$\triangleright$\enspace\texttt{IsomorphismMatrixAlgebra({\mdseries\slshape A})\index{IsomorphismMatrixAlgebra@\texttt{IsomorphismMatrixAlgebra}}
\label{IsomorphismMatrixAlgebra}
}\hfill{\scriptsize (attribute)}}\\


 isomorphism from the algebra \mbox{\texttt{\mdseries\slshape A}} onto a matrix algebra. Currently this is only implemented for associative
algebras with one. 
\begin{Verbatim}[commandchars=!@|,fontsize=\small,frame=single,label=Example]
  !gapprompt@gap>| !gapinput@T:= EmptySCTable( 2, 0 );;|
  !gapprompt@gap>| !gapinput@SetEntrySCTable( T, 1, 1, [1,1] ); SetEntrySCTable( T, 2, 2, [1,2] );|
  !gapprompt@gap>| !gapinput@A:= AlgebraByStructureConstants( Rationals, T );;|
  !gapprompt@gap>| !gapinput@A:= AsAlgebraWithOne( Rationals, A );;|
  !gapprompt@gap>| !gapinput@f:=IsomorphismMatrixAlgebra( A );|
  <op. hom. AlgebraWithOne( Rationals, ... ) -> matrices of dim. 2>
  !gapprompt@gap>| !gapinput@Image( f, BasisVectors( Basis( A ) )[1] );|
  [ [ 1, 0 ], [ 0, 0 ] ]
\end{Verbatim}
 }

 

\subsection{\textcolor{Chapter }{IsomorphismSCAlgebra (w.r.t. a given basis)}}
\logpage{[ 62, 10, 13 ]}\nobreak
\hyperdef{L}{X7F8D3DF2863EC50D}{}
{\noindent\textcolor{FuncColor}{$\triangleright$\enspace\texttt{IsomorphismSCAlgebra({\mdseries\slshape B})\index{IsomorphismSCAlgebra@\texttt{IsomorphismSCAlgebra}!w.r.t. a given basis}
\label{IsomorphismSCAlgebra:w.r.t. a given basis}
}\hfill{\scriptsize (attribute)}}\\
\noindent\textcolor{FuncColor}{$\triangleright$\enspace\texttt{IsomorphismSCAlgebra({\mdseries\slshape A})\index{IsomorphismSCAlgebra@\texttt{IsomorphismSCAlgebra}!for an algebra}
\label{IsomorphismSCAlgebra:for an algebra}
}\hfill{\scriptsize (attribute)}}\\


 For a basis \mbox{\texttt{\mdseries\slshape B}} of an algebra $A$, \texttt{IsomorphismSCAlgebra} returns an algebra isomorphism from $A$ to an algebra $S$ given by structure constants (see{\nobreakspace}\ref{Constructing Algebras by Structure Constants}), such that the canonical basis of $S$ is the image of \mbox{\texttt{\mdseries\slshape B}}. 

 For an algebra \mbox{\texttt{\mdseries\slshape A}}, \texttt{IsomorphismSCAlgebra} chooses a basis of \mbox{\texttt{\mdseries\slshape A}} and returns the \texttt{IsomorphismSCAlgebra} value for that basis. 

 
\begin{Verbatim}[commandchars=!@|,fontsize=\small,frame=single,label=Example]
  !gapprompt@gap>| !gapinput@IsomorphismSCAlgebra( GF(8) );|
  CanonicalBasis( GF(2^3) ) -> CanonicalBasis( <algebra of dimension
  3 over GF(2)> )
  !gapprompt@gap>| !gapinput@IsomorphismSCAlgebra( GF(2)^[2,2] );|
  CanonicalBasis( ( GF(2)^
  [ 2, 2 ] ) ) -> CanonicalBasis( <algebra of dimension 4 over GF(2)> )
\end{Verbatim}
 }

 

\subsection{\textcolor{Chapter }{RepresentativeLinearOperation}}
\logpage{[ 62, 10, 14 ]}\nobreak
\hyperdef{L}{X7F34244B81979696}{}
{\noindent\textcolor{FuncColor}{$\triangleright$\enspace\texttt{RepresentativeLinearOperation({\mdseries\slshape A, v, w, opr})\index{RepresentativeLinearOperation@\texttt{RepresentativeLinearOperation}}
\label{RepresentativeLinearOperation}
}\hfill{\scriptsize (operation)}}\\


 is an element of the algebra \mbox{\texttt{\mdseries\slshape A}} that maps the vector \mbox{\texttt{\mdseries\slshape v}} to the vector \mbox{\texttt{\mdseries\slshape w}} under the linear operation described by the function \mbox{\texttt{\mdseries\slshape opr}}. If no such element exists then \texttt{fail} is returned. 

   
\begin{Verbatim}[commandchars=!@|,fontsize=\small,frame=single,label=Example]
  !gapprompt@gap>| !gapinput@m1:= NullMat( 2, 2 );; m1[1][1]:= 1;;|
  !gapprompt@gap>| !gapinput@m2:= NullMat( 2, 2 );; m2[2][2]:= 1;;|
  !gapprompt@gap>| !gapinput@B:= AlgebraByGenerators( Rationals, [ m1, m2 ] );;|
  !gapprompt@gap>| !gapinput@RepresentativeLinearOperation( B, [1,0], [1,0], OnRight );|
  [ [ 1, 0 ], [ 0, 0 ] ]
  !gapprompt@gap>| !gapinput@RepresentativeLinearOperation( B, [1,0], [0,1], OnRight );|
  fail
\end{Verbatim}
 }

 }

  
\section{\textcolor{Chapter }{Representations of Algebras}}\label{Representations of Algebras}
\logpage{[ 62, 11, 0 ]}
\hyperdef{L}{X818DE6C57D1A4B33}{}
{
  An algebra module is a vector space together with an action of an algebra. So
a module over an algebra is constructed by giving generators of a vector
space, and a function for calculating the action of algebra elements on
elements of the vector space. When creating an algebra module, the generators
of the vector space are wrapped up and given the category \texttt{IsLeftAlgebraModuleElement} or \texttt{IsRightModuleElement} if the algebra acts from the left, or right respectively. (So in the case of a
bi\texttt{\symbol{45}}module the elements get both categories.) Most linear
algebra computations are delegated to the original vector space. 

 The transition between the original vector space and the corresponding algebra
module is handled by \texttt{ExtRepOfObj} and \texttt{ObjByExtRep}. For an element \texttt{v} of the algebra module, \texttt{ExtRepOfObj( v )} returns the underlying element of the original vector space. Furthermore, if \texttt{vec} is an element of the original vector space, and \texttt{fam} the elements family of the corresponding algebra module, then \texttt{ObjByExtRep( fam, vec )} returns the corresponding element of the algebra module. Below is an example
of this. 

 The action of the algebra on elements of the algebra module is constructed by
using the operator \texttt{\texttt{\symbol{94}}}. If \texttt{x} is an element of an algebra \texttt{A}, and \texttt{v} an element of a left \texttt{A}\texttt{\symbol{45}}module, then \texttt{x\texttt{\symbol{94}}v} calculates the result of the action of \texttt{x} on \texttt{v}. Similarly, if \texttt{v} is an element of a right \texttt{A}\texttt{\symbol{45}}module, then \texttt{v\texttt{\symbol{94}}x} calculates the action of \texttt{x} on \texttt{v}. 

\subsection{\textcolor{Chapter }{LeftAlgebraModuleByGenerators}}
\logpage{[ 62, 11, 1 ]}\nobreak
\hyperdef{L}{X8055B87F7ADBD66B}{}
{\noindent\textcolor{FuncColor}{$\triangleright$\enspace\texttt{LeftAlgebraModuleByGenerators({\mdseries\slshape A, op, gens})\index{LeftAlgebraModuleByGenerators@\texttt{LeftAlgebraModuleByGenerators}}
\label{LeftAlgebraModuleByGenerators}
}\hfill{\scriptsize (operation)}}\\


 Constructs the left algebra module over \mbox{\texttt{\mdseries\slshape A}} generated by the list of vectors \mbox{\texttt{\mdseries\slshape gens}}. The action of \mbox{\texttt{\mdseries\slshape A}} is described by the function \mbox{\texttt{\mdseries\slshape op}}. This must be a function of two arguments; the first argument is the algebra
element, and the second argument is a vector; it outputs the result of
applying the algebra element to the vector. }

 

\subsection{\textcolor{Chapter }{RightAlgebraModuleByGenerators}}
\logpage{[ 62, 11, 2 ]}\nobreak
\hyperdef{L}{X8026B99B7955A355}{}
{\noindent\textcolor{FuncColor}{$\triangleright$\enspace\texttt{RightAlgebraModuleByGenerators({\mdseries\slshape A, op, gens})\index{RightAlgebraModuleByGenerators@\texttt{RightAlgebraModuleByGenerators}}
\label{RightAlgebraModuleByGenerators}
}\hfill{\scriptsize (operation)}}\\


 Constructs the right algebra module over \mbox{\texttt{\mdseries\slshape A}} generated by the list of vectors \mbox{\texttt{\mdseries\slshape gens}}. The action of \mbox{\texttt{\mdseries\slshape A}} is described by the function \mbox{\texttt{\mdseries\slshape op}}. This must be a function of two arguments; the first argument is a vector,
and the second argument is the algebra element; it outputs the result of
applying the algebra element to the vector. }

 

\subsection{\textcolor{Chapter }{BiAlgebraModuleByGenerators}}
\logpage{[ 62, 11, 3 ]}\nobreak
\hyperdef{L}{X7F28A47E876427E0}{}
{\noindent\textcolor{FuncColor}{$\triangleright$\enspace\texttt{BiAlgebraModuleByGenerators({\mdseries\slshape A, B, opl, opr, gens})\index{BiAlgebraModuleByGenerators@\texttt{BiAlgebraModuleByGenerators}}
\label{BiAlgebraModuleByGenerators}
}\hfill{\scriptsize (operation)}}\\


 Constructs the algebra bi\texttt{\symbol{45}}module over \mbox{\texttt{\mdseries\slshape A}} and \mbox{\texttt{\mdseries\slshape B}} generated by the list of vectors \mbox{\texttt{\mdseries\slshape gens}}. The left action of \mbox{\texttt{\mdseries\slshape A}} is described by the function \mbox{\texttt{\mdseries\slshape opl}}, and the right action of \mbox{\texttt{\mdseries\slshape B}} by the function \mbox{\texttt{\mdseries\slshape opr}}. \mbox{\texttt{\mdseries\slshape opl}} must be a function of two arguments; the first argument is the algebra
element, and the second argument is a vector; it outputs the result of
applying the algebra element on the left to the vector. \mbox{\texttt{\mdseries\slshape opr}} must be a function of two arguments; the first argument is a vector, and the
second argument is the algebra element; it outputs the result of applying the
algebra element on the right to the vector. 
\begin{Verbatim}[commandchars=!@|,fontsize=\small,frame=single,label=Example]
  !gapprompt@gap>| !gapinput@A:= Rationals^[3,3];|
  ( Rationals^[ 3, 3 ] )
  !gapprompt@gap>| !gapinput@V:= LeftAlgebraModuleByGenerators( A, \*, [ [ 1, 0, 0 ] ] );|
  <left-module over ( Rationals^[ 3, 3 ] )>
  !gapprompt@gap>| !gapinput@W:= RightAlgebraModuleByGenerators( A, \*, [ [ 1, 0, 0 ] ] );|
  <right-module over ( Rationals^[ 3, 3 ] )>
  !gapprompt@gap>| !gapinput@M:= BiAlgebraModuleByGenerators( A, A, \*, \*, [ [ 1, 0, 0 ] ] );|
  <bi-module over ( Rationals^[ 3, 3 ] ) (left) and ( Rationals^
  [ 3, 3 ] ) (right)>
\end{Verbatim}
 

 In the above examples, the modules \texttt{V}, \texttt{W}, and \texttt{M} are $3$\texttt{\symbol{45}}dimensional vector spaces over the rationals. The algebra \texttt{A} acts from the left on \texttt{V}, from the right on \texttt{W}, and from the left and from the right on \texttt{M}. }

 

\subsection{\textcolor{Chapter }{LeftAlgebraModule}}
\logpage{[ 62, 11, 4 ]}\nobreak
\hyperdef{L}{X852524F581613359}{}
{\noindent\textcolor{FuncColor}{$\triangleright$\enspace\texttt{LeftAlgebraModule({\mdseries\slshape A, op, V})\index{LeftAlgebraModule@\texttt{LeftAlgebraModule}}
\label{LeftAlgebraModule}
}\hfill{\scriptsize (operation)}}\\


 Constructs the left algebra module over \mbox{\texttt{\mdseries\slshape A}} with underlying space \mbox{\texttt{\mdseries\slshape V}}. The action of \mbox{\texttt{\mdseries\slshape A}} is described by the function \mbox{\texttt{\mdseries\slshape op}}. This must be a function of two arguments; the first argument is the algebra
element, and the second argument is a vector from \mbox{\texttt{\mdseries\slshape V}}; it outputs the result of applying the algebra element to the vector. }

 

\subsection{\textcolor{Chapter }{RightAlgebraModule}}
\logpage{[ 62, 11, 5 ]}\nobreak
\hyperdef{L}{X8222F2B67D753036}{}
{\noindent\textcolor{FuncColor}{$\triangleright$\enspace\texttt{RightAlgebraModule({\mdseries\slshape A, op, V})\index{RightAlgebraModule@\texttt{RightAlgebraModule}}
\label{RightAlgebraModule}
}\hfill{\scriptsize (operation)}}\\


 Constructs the right algebra module over \mbox{\texttt{\mdseries\slshape A}} with underlying space \mbox{\texttt{\mdseries\slshape V}}. The action of \mbox{\texttt{\mdseries\slshape A}} is described by the function \mbox{\texttt{\mdseries\slshape op}}. This must be a function of two arguments; the first argument is a vector,
from \mbox{\texttt{\mdseries\slshape V}} and the second argument is the algebra element; it outputs the result of
applying the algebra element to the vector. }

 

\subsection{\textcolor{Chapter }{BiAlgebraModule}}
\logpage{[ 62, 11, 6 ]}\nobreak
\hyperdef{L}{X84517770868DDA02}{}
{\noindent\textcolor{FuncColor}{$\triangleright$\enspace\texttt{BiAlgebraModule({\mdseries\slshape A, B, opl, opr, V})\index{BiAlgebraModule@\texttt{BiAlgebraModule}}
\label{BiAlgebraModule}
}\hfill{\scriptsize (operation)}}\\


 Constructs the algebra bi\texttt{\symbol{45}}module over \mbox{\texttt{\mdseries\slshape A}} and \mbox{\texttt{\mdseries\slshape B}} with underlying space \mbox{\texttt{\mdseries\slshape V}}. The left action of \mbox{\texttt{\mdseries\slshape A}} is described by the function \mbox{\texttt{\mdseries\slshape opl}}, and the right action of \mbox{\texttt{\mdseries\slshape B}} by the function \mbox{\texttt{\mdseries\slshape opr}}. \mbox{\texttt{\mdseries\slshape opl}} must be a function of two arguments; the first argument is the algebra
element, and the second argument is a vector from \mbox{\texttt{\mdseries\slshape V}}; it outputs the result of applying the algebra element on the left to the
vector. \mbox{\texttt{\mdseries\slshape opr}} must be a function of two arguments; the first argument is a vector from \mbox{\texttt{\mdseries\slshape V}}, and the second argument is the algebra element; it outputs the result of
applying the algebra element on the right to the vector. 
\begin{Verbatim}[commandchars=!@|,fontsize=\small,frame=single,label=Example]
  !gapprompt@gap>| !gapinput@A:= Rationals^[3,3];;|
  !gapprompt@gap>| !gapinput@V:= Rationals^3;|
  ( Rationals^3 )
  !gapprompt@gap>| !gapinput@V:= Rationals^3;;|
  !gapprompt@gap>| !gapinput@M:= BiAlgebraModule( A, A, \*, \*, V );|
  <bi-module over ( Rationals^[ 3, 3 ] ) (left) and ( Rationals^
  [ 3, 3 ] ) (right)>
  !gapprompt@gap>| !gapinput@Dimension( M );|
  3
\end{Verbatim}
 }

 

\subsection{\textcolor{Chapter }{GeneratorsOfAlgebraModule}}
\logpage{[ 62, 11, 7 ]}\nobreak
\hyperdef{L}{X79AAB50D83A14A43}{}
{\noindent\textcolor{FuncColor}{$\triangleright$\enspace\texttt{GeneratorsOfAlgebraModule({\mdseries\slshape M})\index{GeneratorsOfAlgebraModule@\texttt{GeneratorsOfAlgebraModule}}
\label{GeneratorsOfAlgebraModule}
}\hfill{\scriptsize (attribute)}}\\


 A list of elements of \mbox{\texttt{\mdseries\slshape M}} that generate \mbox{\texttt{\mdseries\slshape M}} as an algebra module. 
\begin{Verbatim}[commandchars=!@|,fontsize=\small,frame=single,label=Example]
  !gapprompt@gap>| !gapinput@A:= Rationals^[3,3];;|
  !gapprompt@gap>| !gapinput@V:= LeftAlgebraModuleByGenerators( A, \*, [ [ 1, 0, 0 ] ] );;|
  !gapprompt@gap>| !gapinput@GeneratorsOfAlgebraModule( V );|
  [ [ 1, 0, 0 ] ]
\end{Verbatim}
 }

 

\subsection{\textcolor{Chapter }{IsAlgebraModuleElement}}
\logpage{[ 62, 11, 8 ]}\nobreak
\hyperdef{L}{X82B708BD84F3DAB1}{}
{\noindent\textcolor{FuncColor}{$\triangleright$\enspace\texttt{IsAlgebraModuleElement({\mdseries\slshape obj})\index{IsAlgebraModuleElement@\texttt{IsAlgebraModuleElement}}
\label{IsAlgebraModuleElement}
}\hfill{\scriptsize (Category)}}\\
\noindent\textcolor{FuncColor}{$\triangleright$\enspace\texttt{IsAlgebraModuleElementCollection({\mdseries\slshape obj})\index{IsAlgebraModuleElementCollection@\texttt{IsAlgebraModuleElementCollection}}
\label{IsAlgebraModuleElementCollection}
}\hfill{\scriptsize (Category)}}\\
\noindent\textcolor{FuncColor}{$\triangleright$\enspace\texttt{IsAlgebraModuleElementFamily({\mdseries\slshape fam})\index{IsAlgebraModuleElementFamily@\texttt{IsAlgebraModuleElementFamily}}
\label{IsAlgebraModuleElementFamily}
}\hfill{\scriptsize (Category)}}\\


 Category of algebra module elements. If an object has \texttt{IsAlgebraModuleElementCollection}, then it is an algebra module. If a family has \texttt{IsAlgebraModuleElementFamily}, then it is a family of algebra module elements (every algebra module has its
own elements family). }

 

\subsection{\textcolor{Chapter }{IsLeftAlgebraModuleElement}}
\logpage{[ 62, 11, 9 ]}\nobreak
\hyperdef{L}{X80E786467F9163F9}{}
{\noindent\textcolor{FuncColor}{$\triangleright$\enspace\texttt{IsLeftAlgebraModuleElement({\mdseries\slshape obj})\index{IsLeftAlgebraModuleElement@\texttt{IsLeftAlgebraModuleElement}}
\label{IsLeftAlgebraModuleElement}
}\hfill{\scriptsize (Category)}}\\
\noindent\textcolor{FuncColor}{$\triangleright$\enspace\texttt{IsLeftAlgebraModuleElementCollection({\mdseries\slshape obj})\index{IsLeftAlgebraModuleElementCollection@\texttt{IsLeft}\-\texttt{Algebra}\-\texttt{Module}\-\texttt{Element}\-\texttt{Collection}}
\label{IsLeftAlgebraModuleElementCollection}
}\hfill{\scriptsize (Category)}}\\


 Category of left algebra module elements. If an object has \texttt{IsLeftAlgebraModuleElementCollection}, then it is a left\texttt{\symbol{45}}algebra module. }

 

\subsection{\textcolor{Chapter }{IsRightAlgebraModuleElement}}
\logpage{[ 62, 11, 10 ]}\nobreak
\hyperdef{L}{X863756787E2B6E75}{}
{\noindent\textcolor{FuncColor}{$\triangleright$\enspace\texttt{IsRightAlgebraModuleElement({\mdseries\slshape obj})\index{IsRightAlgebraModuleElement@\texttt{IsRightAlgebraModuleElement}}
\label{IsRightAlgebraModuleElement}
}\hfill{\scriptsize (Category)}}\\
\noindent\textcolor{FuncColor}{$\triangleright$\enspace\texttt{IsRightAlgebraModuleElementCollection({\mdseries\slshape obj})\index{IsRightAlgebraModuleElementCollection@\texttt{IsRight}\-\texttt{Algebra}\-\texttt{Module}\-\texttt{Element}\-\texttt{Collection}}
\label{IsRightAlgebraModuleElementCollection}
}\hfill{\scriptsize (Category)}}\\


 Category of right algebra module elements. If an object has \texttt{IsRightAlgebraModuleElementCollection}, then it is a right\texttt{\symbol{45}}algebra module. 
\begin{Verbatim}[commandchars=!@|,fontsize=\small,frame=single,label=Example]
  !gapprompt@gap>| !gapinput@A:= Rationals^[3,3];|
  ( Rationals^[ 3, 3 ] )
  !gapprompt@gap>| !gapinput@M:= BiAlgebraModuleByGenerators( A, A, \*, \*, [ [ 1, 0, 0 ] ] );|
  <bi-module over ( Rationals^[ 3, 3 ] ) (left) and ( Rationals^
  [ 3, 3 ] ) (right)>
  !gapprompt@gap>| !gapinput@vv:= BasisVectors( Basis( M ) );|
  [ [ 1, 0, 0 ], [ 0, 1, 0 ], [ 0, 0, 1 ] ]
  !gapprompt@gap>| !gapinput@IsLeftAlgebraModuleElement( vv[1] );|
  true
  !gapprompt@gap>| !gapinput@IsRightAlgebraModuleElement( vv[1] );|
  true
  !gapprompt@gap>| !gapinput@vv[1] = [ 1, 0, 0 ];|
  false
  !gapprompt@gap>| !gapinput@ExtRepOfObj( vv[1] ) = [ 1, 0, 0 ];|
  true
  !gapprompt@gap>| !gapinput@ObjByExtRep( ElementsFamily( FamilyObj( M ) ), [ 1, 0, 0 ] ) in M;|
  true
  !gapprompt@gap>| !gapinput@xx:= BasisVectors( Basis( A ) );;|
  !gapprompt@gap>| !gapinput@xx[4]^vv[1];  # left action|
  [ 0, 1, 0 ]
  !gapprompt@gap>| !gapinput@vv[1]^xx[2];  # right action|
  [ 0, 1, 0 ]
\end{Verbatim}
 }

 

\subsection{\textcolor{Chapter }{LeftActingAlgebra}}
\logpage{[ 62, 11, 11 ]}\nobreak
\hyperdef{L}{X85654EF07F708AC3}{}
{\noindent\textcolor{FuncColor}{$\triangleright$\enspace\texttt{LeftActingAlgebra({\mdseries\slshape V})\index{LeftActingAlgebra@\texttt{LeftActingAlgebra}}
\label{LeftActingAlgebra}
}\hfill{\scriptsize (attribute)}}\\


 Here \mbox{\texttt{\mdseries\slshape V}} is a left\texttt{\symbol{45}}algebra module; this function returns the algebra
that acts from the left on \mbox{\texttt{\mdseries\slshape V}}. }

 

\subsection{\textcolor{Chapter }{RightActingAlgebra}}
\logpage{[ 62, 11, 12 ]}\nobreak
\hyperdef{L}{X826298B37E1B1520}{}
{\noindent\textcolor{FuncColor}{$\triangleright$\enspace\texttt{RightActingAlgebra({\mdseries\slshape V})\index{RightActingAlgebra@\texttt{RightActingAlgebra}}
\label{RightActingAlgebra}
}\hfill{\scriptsize (attribute)}}\\


 Here \mbox{\texttt{\mdseries\slshape V}} is a right\texttt{\symbol{45}}algebra module; this function returns the
algebra that acts from the right on \mbox{\texttt{\mdseries\slshape V}}. }

 

\subsection{\textcolor{Chapter }{ActingAlgebra}}
\logpage{[ 62, 11, 13 ]}\nobreak
\hyperdef{L}{X8308408D86CFC3C9}{}
{\noindent\textcolor{FuncColor}{$\triangleright$\enspace\texttt{ActingAlgebra({\mdseries\slshape V})\index{ActingAlgebra@\texttt{ActingAlgebra}}
\label{ActingAlgebra}
}\hfill{\scriptsize (operation)}}\\


 Here \mbox{\texttt{\mdseries\slshape V}} is an algebra module; this function returns the algebra that acts on \mbox{\texttt{\mdseries\slshape V}} (this is the same as \texttt{LeftActingAlgebra( \mbox{\texttt{\mdseries\slshape V}} )} if \mbox{\texttt{\mdseries\slshape V}} is a left module, and \texttt{RightActingAlgebra( \mbox{\texttt{\mdseries\slshape V}} )} if \mbox{\texttt{\mdseries\slshape V}} is a right module; it will signal an error if \mbox{\texttt{\mdseries\slshape V}} is a bi\texttt{\symbol{45}}module). 
\begin{Verbatim}[commandchars=!@|,fontsize=\small,frame=single,label=Example]
  !gapprompt@gap>| !gapinput@A:= Rationals^[3,3];;|
  !gapprompt@gap>| !gapinput@M:= BiAlgebraModuleByGenerators( A, A, \*, \*, [ [ 1, 0, 0 ] ] );;|
  !gapprompt@gap>| !gapinput@LeftActingAlgebra( M );|
  ( Rationals^[ 3, 3 ] )
  !gapprompt@gap>| !gapinput@RightActingAlgebra( M );|
  ( Rationals^[ 3, 3 ] )
  !gapprompt@gap>| !gapinput@V:= RightAlgebraModuleByGenerators( A, \*, [ [ 1, 0, 0 ] ] );;|
  !gapprompt@gap>| !gapinput@ActingAlgebra( V );|
  ( Rationals^[ 3, 3 ] )
\end{Verbatim}
 }

 

\subsection{\textcolor{Chapter }{IsBasisOfAlgebraModuleElementSpace}}
\logpage{[ 62, 11, 14 ]}\nobreak
\hyperdef{L}{X7C325A507EC9BA18}{}
{\noindent\textcolor{FuncColor}{$\triangleright$\enspace\texttt{IsBasisOfAlgebraModuleElementSpace({\mdseries\slshape B})\index{IsBasisOfAlgebraModuleElementSpace@\texttt{IsBasisOfAlgebraModuleElementSpace}}
\label{IsBasisOfAlgebraModuleElementSpace}
}\hfill{\scriptsize (Category)}}\\


 If a basis \mbox{\texttt{\mdseries\slshape B}} lies in the category \texttt{IsBasisOfAlgebraModuleElementSpace}, then \mbox{\texttt{\mdseries\slshape B}} is a basis of a subspace of an algebra module. This means that \mbox{\texttt{\mdseries\slshape B}} has the record field \texttt{\mbox{\texttt{\mdseries\slshape B}}!.delegateBasis} set. This last object is a basis of the corresponding subspace of the vector
space underlying the algebra module (i.e., the vector space spanned by all \texttt{ExtRepOfObj( v )} for \texttt{v} in the algebra module). 
\begin{Verbatim}[commandchars=@|C,fontsize=\small,frame=single,label=Example]
  @gapprompt|gap>C @gapinput|A:= Rationals^[3,3];;C
  @gapprompt|gap>C @gapinput|M:= BiAlgebraModuleByGenerators( A, A, \*, \*, [ [ 1, 0, 0 ] ] );;C
  @gapprompt|gap>C @gapinput|B:= Basis( M );C
  Basis( <3-dimensional bi-module over ( Rationals^
  [ 3, 3 ] ) (left) and ( Rationals^[ 3, 3 ] ) (right)>,
  [ [ 1, 0, 0 ], [ 0, 1, 0 ], [ 0, 0, 1 ] ] )
  @gapprompt|gap>C @gapinput|IsBasisOfAlgebraModuleElementSpace( B );C
  true
  @gapprompt|gap>C @gapinput|B!.delegateBasis;C
  SemiEchelonBasis( <vector space of dimension 3 over Rationals>,
  [ [ 1, 0, 0 ], [ 0, 1, 0 ], [ 0, 0, 1 ] ] )
\end{Verbatim}
 }

 

\subsection{\textcolor{Chapter }{MatrixOfAction}}
\logpage{[ 62, 11, 15 ]}\nobreak
\hyperdef{L}{X789863037B0E35D2}{}
{\noindent\textcolor{FuncColor}{$\triangleright$\enspace\texttt{MatrixOfAction({\mdseries\slshape B, x[, side]})\index{MatrixOfAction@\texttt{MatrixOfAction}}
\label{MatrixOfAction}
}\hfill{\scriptsize (operation)}}\\


 Here \mbox{\texttt{\mdseries\slshape B}} is a basis of an algebra module and \mbox{\texttt{\mdseries\slshape x}} is an element of the algebra that acts on this module. This function returns
the matrix of the action of \mbox{\texttt{\mdseries\slshape x}} with respect to \mbox{\texttt{\mdseries\slshape B}}. If \mbox{\texttt{\mdseries\slshape x}} acts from the left, then the coefficients of the images of the basis elements
of \mbox{\texttt{\mdseries\slshape B}} (under the action of \mbox{\texttt{\mdseries\slshape x}}) are the columns of the output. If \mbox{\texttt{\mdseries\slshape x}} acts from the right, then they are the rows of the output. 

 If the module is a bi\texttt{\symbol{45}}module, then the third parameter \mbox{\texttt{\mdseries\slshape side}} must be specified. This is the string \texttt{"left"}, or \texttt{"right"} depending whether \mbox{\texttt{\mdseries\slshape x}} acts from the left or the right. 
\begin{Verbatim}[commandchars=!@|,fontsize=\small,frame=single,label=Example]
  !gapprompt@gap>| !gapinput@M:= LeftAlgebraModuleByGenerators( A, \*, [ [ 1, 0, 0 ] ] );;|
  !gapprompt@gap>| !gapinput@x:= Basis(A)[3];|
  [ [ 0, 0, 1 ], [ 0, 0, 0 ], [ 0, 0, 0 ] ]
  !gapprompt@gap>| !gapinput@MatrixOfAction( Basis( M ), x );|
  [ [ 0, 0, 1 ], [ 0, 0, 0 ], [ 0, 0, 0 ] ]
\end{Verbatim}
 }

 

\subsection{\textcolor{Chapter }{SubAlgebraModule}}
\logpage{[ 62, 11, 16 ]}\nobreak
\hyperdef{L}{X8742A7D27F26AFAB}{}
{\noindent\textcolor{FuncColor}{$\triangleright$\enspace\texttt{SubAlgebraModule({\mdseries\slshape M, gens[, "basis"]})\index{SubAlgebraModule@\texttt{SubAlgebraModule}}
\label{SubAlgebraModule}
}\hfill{\scriptsize (operation)}}\\


 is the sub\texttt{\symbol{45}}module of the algebra module \mbox{\texttt{\mdseries\slshape M}}, generated by the vectors in \mbox{\texttt{\mdseries\slshape gens}}. If as an optional argument the string \texttt{basis} is added, then it is assumed that the vectors in \mbox{\texttt{\mdseries\slshape gens}} form a basis of the submodule. 
\begin{Verbatim}[commandchars=!@|,fontsize=\small,frame=single,label=Example]
  !gapprompt@gap>| !gapinput@m1:= NullMat( 2, 2 );; m1[1][1]:= 1;;|
  !gapprompt@gap>| !gapinput@m2:= NullMat( 2, 2 );; m2[2][2]:= 1;;|
  !gapprompt@gap>| !gapinput@A:= Algebra( Rationals, [ m1, m2 ] );;|
  !gapprompt@gap>| !gapinput@M:= LeftAlgebraModuleByGenerators( A, \*, [ [ 1, 0 ], [ 0, 1 ] ] );|
  <left-module over <algebra over Rationals, with 2 generators>>
  !gapprompt@gap>| !gapinput@bb:= BasisVectors( Basis( M ) );|
  [ [ 1, 0 ], [ 0, 1 ] ]
  !gapprompt@gap>| !gapinput@V:= SubAlgebraModule( M, [ bb[1] ] );|
  <left-module over <algebra over Rationals, with 2 generators>>
  !gapprompt@gap>| !gapinput@Dimension( V );|
  1
\end{Verbatim}
 }

 

\subsection{\textcolor{Chapter }{LeftModuleByHomomorphismToMatAlg}}
\logpage{[ 62, 11, 17 ]}\nobreak
\hyperdef{L}{X86E0515987192F0E}{}
{\noindent\textcolor{FuncColor}{$\triangleright$\enspace\texttt{LeftModuleByHomomorphismToMatAlg({\mdseries\slshape A, hom})\index{LeftModuleByHomomorphismToMatAlg@\texttt{LeftModuleByHomomorphismToMatAlg}}
\label{LeftModuleByHomomorphismToMatAlg}
}\hfill{\scriptsize (operation)}}\\


 Here \mbox{\texttt{\mdseries\slshape A}} is an algebra and \mbox{\texttt{\mdseries\slshape hom}} a homomorphism from \mbox{\texttt{\mdseries\slshape A}} into a matrix algebra. This function returns the left \mbox{\texttt{\mdseries\slshape A}}\texttt{\symbol{45}}module defined by the homomorphism \mbox{\texttt{\mdseries\slshape hom}}. }

 

\subsection{\textcolor{Chapter }{RightModuleByHomomorphismToMatAlg}}
\logpage{[ 62, 11, 18 ]}\nobreak
\hyperdef{L}{X7EE41297867E41A8}{}
{\noindent\textcolor{FuncColor}{$\triangleright$\enspace\texttt{RightModuleByHomomorphismToMatAlg({\mdseries\slshape A, hom})\index{RightModuleByHomomorphismToMatAlg@\texttt{RightModuleByHomomorphismToMatAlg}}
\label{RightModuleByHomomorphismToMatAlg}
}\hfill{\scriptsize (operation)}}\\


 Here \mbox{\texttt{\mdseries\slshape A}} is an algebra and \mbox{\texttt{\mdseries\slshape hom}} a homomorphism from \mbox{\texttt{\mdseries\slshape A}} into a matrix algebra. This function returns the right \mbox{\texttt{\mdseries\slshape A}}\texttt{\symbol{45}}module defined by the homomorphism \mbox{\texttt{\mdseries\slshape hom}}. 

 First we produce a structure constants algebra with basis elements $x$, $y$, $z$ such that $x^2 = x$, $y^2 = y$, $xz = z$, $zy = z$ and all other products are zero. 

 
\begin{Verbatim}[commandchars=!@|,fontsize=\small,frame=single,label=Example]
  !gapprompt@gap>| !gapinput@T:= EmptySCTable( 3, 0 );;|
  !gapprompt@gap>| !gapinput@SetEntrySCTable( T, 1, 1, [ 1, 1 ]);|
  !gapprompt@gap>| !gapinput@SetEntrySCTable( T, 2, 2, [ 1, 2 ]);|
  !gapprompt@gap>| !gapinput@SetEntrySCTable( T, 1, 3, [ 1, 3 ]);|
  !gapprompt@gap>| !gapinput@SetEntrySCTable( T, 3, 2, [ 1, 3 ]);|
  !gapprompt@gap>| !gapinput@A:= AlgebraByStructureConstants( Rationals, T );|
  <algebra of dimension 3 over Rationals>
\end{Verbatim}
 

 Now we construct an isomorphic matrix algebra. 

 
\begin{Verbatim}[commandchars=!@|,fontsize=\small,frame=single,label=Example]
  !gapprompt@gap>| !gapinput@m1:= NullMat( 2, 2 );; m1[1][1]:= 1;;|
  !gapprompt@gap>| !gapinput@m2:= NullMat( 2, 2 );; m2[2][2]:= 1;;|
  !gapprompt@gap>| !gapinput@m3:= NullMat( 2, 2 );; m3[1][2]:= 1;;|
  !gapprompt@gap>| !gapinput@B:= Algebra( Rationals, [ m1, m2, m3 ] );|
  <algebra over Rationals, with 3 generators>
\end{Verbatim}
 

 Finally we construct the homomorphism and the corresponding right module. 

 
\begin{Verbatim}[commandchars=!@|,fontsize=\small,frame=single,label=Example]
  !gapprompt@gap>| !gapinput@f:= AlgebraHomomorphismByImages( A, B, Basis(A), [ m1, m2, m3 ] );;|
  !gapprompt@gap>| !gapinput@RightModuleByHomomorphismToMatAlg( A, f );|
  <right-module over <algebra of dimension 3 over Rationals>>
\end{Verbatim}
 }

 

\subsection{\textcolor{Chapter }{AdjointModule}}
\logpage{[ 62, 11, 19 ]}\nobreak
\hyperdef{L}{X8729F0A678A4A09C}{}
{\noindent\textcolor{FuncColor}{$\triangleright$\enspace\texttt{AdjointModule({\mdseries\slshape A})\index{AdjointModule@\texttt{AdjointModule}}
\label{AdjointModule}
}\hfill{\scriptsize (attribute)}}\\


 returns the \mbox{\texttt{\mdseries\slshape A}}\texttt{\symbol{45}}module defined by the left action of \mbox{\texttt{\mdseries\slshape A}} on itself. 
\begin{Verbatim}[commandchars=!@|,fontsize=\small,frame=single,label=Example]
  !gapprompt@gap>| !gapinput@m1:= NullMat( 2, 2 );; m1[1][1]:= 1;;|
  !gapprompt@gap>| !gapinput@m2:= NullMat( 2, 2 );; m2[2][2]:= 1;;|
  !gapprompt@gap>| !gapinput@m3:= NullMat( 2, 2 );; m3[1][2]:= 1;;|
  !gapprompt@gap>| !gapinput@A:= Algebra( Rationals, [ m1, m2, m3 ] );|
  <algebra over Rationals, with 3 generators>
  !gapprompt@gap>| !gapinput@V:= AdjointModule( A );|
  <3-dimensional left-module over <algebra of dimension
  3 over Rationals>>
  !gapprompt@gap>| !gapinput@v:= Basis( V )[3];|
  [ [ 0, 1 ], [ 0, 0 ] ]
  !gapprompt@gap>| !gapinput@W:= SubAlgebraModule( V, [ v ] );|
  <left-module over <algebra of dimension 3 over Rationals>>
  !gapprompt@gap>| !gapinput@Dimension( W );|
  1
\end{Verbatim}
 }

     

\subsection{\textcolor{Chapter }{FaithfulModule (for Lie algebras)}}
\logpage{[ 62, 11, 20 ]}\nobreak
\hyperdef{L}{X84813BCD80BDF3C4}{}
{\noindent\textcolor{FuncColor}{$\triangleright$\enspace\texttt{FaithfulModule({\mdseries\slshape A})\index{FaithfulModule@\texttt{FaithfulModule}!for Lie algebras}
\label{FaithfulModule:for Lie algebras}
}\hfill{\scriptsize (attribute)}}\\


 returns a faithful finite\texttt{\symbol{45}}dimensional
left\texttt{\symbol{45}}module over the algebra \mbox{\texttt{\mdseries\slshape A}}. This is only implemented for associative algebras, and for Lie algebras of
characteristic $0$. (It may also work for certain Lie algebras of characteristic $p > 0$.) 
\begin{Verbatim}[commandchars=!@|,fontsize=\small,frame=single,label=Example]
  !gapprompt@gap>| !gapinput@T:= EmptySCTable( 2, 0 );;|
  !gapprompt@gap>| !gapinput@A:= AlgebraByStructureConstants( Rationals, T );|
  <algebra of dimension 2 over Rationals>
\end{Verbatim}
 
\begin{Verbatim}[commandchars=!@|,fontsize=\small,frame=single,label=Example]
  !gapprompt@gap>| !gapinput@T:= EmptySCTable( 3, 0, "antisymmetric" );;|
  !gapprompt@gap>| !gapinput@SetEntrySCTable( T, 1, 2, [ 1, 3 ]);|
  !gapprompt@gap>| !gapinput@L:= LieAlgebraByStructureConstants( Rationals, T );|
  <Lie algebra of dimension 3 over Rationals>
  !gapprompt@gap>| !gapinput@V:= FaithfulModule( L );|
  <left-module over <Lie algebra of dimension 3 over Rationals>>
  !gapprompt@gap>| !gapinput@vv:= BasisVectors( Basis( V ) );|
  [ [ 1, 0, 0 ], [ 0, 1, 0 ], [ 0, 0, 1 ] ]
  !gapprompt@gap>| !gapinput@x:= Basis( L )[3];|
  v.3
  !gapprompt@gap>| !gapinput@List( vv, v -> x^v );|
  [ [ 0, 0, 0 ], [ 1, 0, 0 ], [ 0, 0, 0 ] ]
\end{Verbatim}
 

 \texttt{A} is a $2$\texttt{\symbol{45}}dimensional algebra where all products are zero. 

 
\begin{Verbatim}[commandchars=!@|,fontsize=\small,frame=single,label=Example]
  !gapprompt@gap>| !gapinput@V:= FaithfulModule( A );|
  <left-module over <algebra of dimension 2 over Rationals>>
  !gapprompt@gap>| !gapinput@vv:= BasisVectors( Basis( V ) );|
  [ [ 1, 0, 0 ], [ 0, 1, 0 ], [ 0, 0, 1 ] ]
  !gapprompt@gap>| !gapinput@xx:= BasisVectors( Basis( A ) );|
  [ v.1, v.2 ]
  !gapprompt@gap>| !gapinput@xx[1]^vv[3];|
  [ 1, 0, 0 ]
\end{Verbatim}
 }

 

\subsection{\textcolor{Chapter }{ModuleByRestriction}}
\logpage{[ 62, 11, 21 ]}\nobreak
\hyperdef{L}{X7E16630185CE2C10}{}
{\noindent\textcolor{FuncColor}{$\triangleright$\enspace\texttt{ModuleByRestriction({\mdseries\slshape V, sub1[, sub2]})\index{ModuleByRestriction@\texttt{ModuleByRestriction}}
\label{ModuleByRestriction}
}\hfill{\scriptsize (operation)}}\\


 Here \mbox{\texttt{\mdseries\slshape V}} is an algebra module and \mbox{\texttt{\mdseries\slshape sub1}} is a subalgebra of the acting algebra of \mbox{\texttt{\mdseries\slshape V}}. This function returns the module that is the restriction of \mbox{\texttt{\mdseries\slshape V}} to \mbox{\texttt{\mdseries\slshape sub1}}. So it has the same underlying vector space as \mbox{\texttt{\mdseries\slshape V}}, but the acting algebra is \mbox{\texttt{\mdseries\slshape sub}}. If two subalgebras \mbox{\texttt{\mdseries\slshape sub1}}, \mbox{\texttt{\mdseries\slshape sub2}} are given then \mbox{\texttt{\mdseries\slshape V}} is assumed to be a bi\texttt{\symbol{45}}module, and \mbox{\texttt{\mdseries\slshape sub1}} a subalgebra of the algebra acting on the left, and \mbox{\texttt{\mdseries\slshape sub2}} a subalgebra of the algebra acting on the right. 
\begin{Verbatim}[commandchars=!@|,fontsize=\small,frame=single,label=Example]
  !gapprompt@gap>| !gapinput@A:= Rationals^[3,3];;|
  !gapprompt@gap>| !gapinput@V:= LeftAlgebraModuleByGenerators( A, \*, [ [ 1, 0, 0 ] ] );;|
  !gapprompt@gap>| !gapinput@B:= Subalgebra( A, [ Basis(A)[1] ] );|
  <algebra over Rationals, with 1 generator>
  !gapprompt@gap>| !gapinput@W:= ModuleByRestriction( V, B );|
  <left-module over <algebra over Rationals, with 1 generator>>
\end{Verbatim}
 }

 

\subsection{\textcolor{Chapter }{NaturalHomomorphismBySubAlgebraModule}}
\logpage{[ 62, 11, 22 ]}\nobreak
\hyperdef{L}{X7885AAC87FDCF649}{}
{\noindent\textcolor{FuncColor}{$\triangleright$\enspace\texttt{NaturalHomomorphismBySubAlgebraModule({\mdseries\slshape V, W})\index{NaturalHomomorphismBySubAlgebraModule@\texttt{Natural}\-\texttt{Homomorphism}\-\texttt{By}\-\texttt{Sub}\-\texttt{Algebra}\-\texttt{Module}}
\label{NaturalHomomorphismBySubAlgebraModule}
}\hfill{\scriptsize (operation)}}\\


 Here \mbox{\texttt{\mdseries\slshape V}} must be a sub\texttt{\symbol{45}}algebra module of \mbox{\texttt{\mdseries\slshape V}}. This function returns the projection from \mbox{\texttt{\mdseries\slshape V}} onto \texttt{\mbox{\texttt{\mdseries\slshape V}}/\mbox{\texttt{\mdseries\slshape W}}}. It is a linear map, that is also a module homomorphism. As usual images can
be formed with \texttt{Image( f, v )} and pre\texttt{\symbol{45}}images with \texttt{PreImagesRepresentative( f, u )}. 

 The quotient module can also be formed by entering \texttt{\mbox{\texttt{\mdseries\slshape V}}/\mbox{\texttt{\mdseries\slshape W}}}. 
\begin{Verbatim}[commandchars=!@|,fontsize=\small,frame=single,label=Example]
  !gapprompt@gap>| !gapinput@A:= Rationals^[3,3];;|
  !gapprompt@gap>| !gapinput@B:= DirectSumOfAlgebras( A, A );|
  <algebra over Rationals, with 6 generators>
  !gapprompt@gap>| !gapinput@T:= StructureConstantsTable( Basis( B ) );;|
  !gapprompt@gap>| !gapinput@C:= AlgebraByStructureConstants( Rationals, T );|
  <algebra of dimension 18 over Rationals>
  !gapprompt@gap>| !gapinput@V:= AdjointModule( C );|
  <left-module over <algebra of dimension 18 over Rationals>>
  !gapprompt@gap>| !gapinput@W:= SubAlgebraModule( V, [ Basis(V)[1] ] );|
  <left-module over <algebra of dimension 18 over Rationals>>
  !gapprompt@gap>| !gapinput@f:= NaturalHomomorphismBySubAlgebraModule( V, W );|
  <linear mapping by matrix, <
  18-dimensional left-module over <algebra of dimension
  18 over Rationals>> -> <
  9-dimensional left-module over <algebra of dimension
  18 over Rationals>>>
  !gapprompt@gap>| !gapinput@quo:= ImagesSource( f );  # i.e., the quotient module|
  <9-dimensional left-module over <algebra of dimension
  18 over Rationals>>
  !gapprompt@gap>| !gapinput@v:= Basis( quo )[1];|
  [ 1, 0, 0, 0, 0, 0, 0, 0, 0 ]
  !gapprompt@gap>| !gapinput@PreImagesRepresentative( f, v );|
  v.4
  !gapprompt@gap>| !gapinput@Basis( C )[4]^v;|
  [ 1, 0, 0, 0, 0, 0, 0, 0, 0 ]
\end{Verbatim}
 }

 

\subsection{\textcolor{Chapter }{DirectSumOfAlgebraModules (for a list of Lie algebra modules)}}
\logpage{[ 62, 11, 23 ]}\nobreak
\hyperdef{L}{X85D0F3758551DADC}{}
{\noindent\textcolor{FuncColor}{$\triangleright$\enspace\texttt{DirectSumOfAlgebraModules({\mdseries\slshape list})\index{DirectSumOfAlgebraModules@\texttt{DirectSumOfAlgebraModules}!for a list of Lie algebra modules}
\label{DirectSumOfAlgebraModules:for a list of Lie algebra modules}
}\hfill{\scriptsize (operation)}}\\
\noindent\textcolor{FuncColor}{$\triangleright$\enspace\texttt{DirectSumOfAlgebraModules({\mdseries\slshape V, W})\index{DirectSumOfAlgebraModules@\texttt{DirectSumOfAlgebraModules}!for two Lie algebra modules}
\label{DirectSumOfAlgebraModules:for two Lie algebra modules}
}\hfill{\scriptsize (operation)}}\\


 Here \mbox{\texttt{\mdseries\slshape list}} must be a list of algebra modules. This function returns the direct sum of the
elements in the list (as an algebra module). The modules must be defined over
the same algebras. 

 In the second form is short for \texttt{DirectSumOfAlgebraModules( [ \mbox{\texttt{\mdseries\slshape V}}, \mbox{\texttt{\mdseries\slshape W}} ] )} 
\begin{Verbatim}[commandchars=!@|,fontsize=\small,frame=single,label=Example]
  !gapprompt@gap>| !gapinput@A:= FullMatrixAlgebra( Rationals, 3 );;|
  !gapprompt@gap>| !gapinput@V:= BiAlgebraModuleByGenerators( A, A, \*, \*, [ [1,0,0] ] );;|
  !gapprompt@gap>| !gapinput@W:= DirectSumOfAlgebraModules( V, V );|
  <6-dimensional left-module over ( Rationals^[ 3, 3 ] )>
  !gapprompt@gap>| !gapinput@BasisVectors( Basis( W ) );|
  [ ( [ 1, 0, 0 ] )(+)( [ 0, 0, 0 ] ), ( [ 0, 1, 0 ] )(+)( [ 0, 0, 0 ] )
      , ( [ 0, 0, 1 ] )(+)( [ 0, 0, 0 ] ),
    ( [ 0, 0, 0 ] )(+)( [ 1, 0, 0 ] ), ( [ 0, 0, 0 ] )(+)( [ 0, 1, 0 ] )
      , ( [ 0, 0, 0 ] )(+)( [ 0, 0, 1 ] ) ]
\end{Verbatim}
 
\begin{Verbatim}[commandchars=!@|,fontsize=\small,frame=single,label=Example]
  !gapprompt@gap>| !gapinput@L:= SimpleLieAlgebra( "C", 3, Rationals );;|
  !gapprompt@gap>| !gapinput@V:= HighestWeightModule( L, [ 1, 1, 0 ] );|
  <64-dimensional left-module over <Lie algebra of dimension
  21 over Rationals>>
  !gapprompt@gap>| !gapinput@W:= HighestWeightModule( L, [ 0, 0, 2 ] );|
  <84-dimensional left-module over <Lie algebra of dimension
  21 over Rationals>>
  !gapprompt@gap>| !gapinput@U:= DirectSumOfAlgebraModules( V, W );|
  <148-dimensional left-module over <Lie algebra of dimension
  21 over Rationals>>
\end{Verbatim}
 }

 

\subsection{\textcolor{Chapter }{TranslatorSubalgebra}}
\logpage{[ 62, 11, 24 ]}\nobreak
\hyperdef{L}{X7D7A6486803B15CE}{}
{\noindent\textcolor{FuncColor}{$\triangleright$\enspace\texttt{TranslatorSubalgebra({\mdseries\slshape M, U, W})\index{TranslatorSubalgebra@\texttt{TranslatorSubalgebra}}
\label{TranslatorSubalgebra}
}\hfill{\scriptsize (operation)}}\\


 Here \mbox{\texttt{\mdseries\slshape M}} is an algebra module, and \mbox{\texttt{\mdseries\slshape U}} and \mbox{\texttt{\mdseries\slshape W}} are two subspaces of \mbox{\texttt{\mdseries\slshape M}}. Let \mbox{\texttt{\mdseries\slshape A}} be the algebra acting on \mbox{\texttt{\mdseries\slshape M}}. This function returns the subspace of elements of \mbox{\texttt{\mdseries\slshape A}} that map \mbox{\texttt{\mdseries\slshape U}} into \mbox{\texttt{\mdseries\slshape W}}. If \mbox{\texttt{\mdseries\slshape W}} is a sub\texttt{\symbol{45}}algebra\texttt{\symbol{45}}module (i.e., closed
under the action of \mbox{\texttt{\mdseries\slshape A}}), then this space is a subalgebra of \mbox{\texttt{\mdseries\slshape A}}. 

 This function works for left, or right modules over a
finite\texttt{\symbol{45}}dimensional algebra. We stress that it is not
checked whether \mbox{\texttt{\mdseries\slshape U}} and \mbox{\texttt{\mdseries\slshape W}} are indeed subspaces of \mbox{\texttt{\mdseries\slshape M}}. If this is not the case nothing is guaranteed about the behaviour of the
function. 
\begin{Verbatim}[commandchars=!@|,fontsize=\small,frame=single,label=Example]
  !gapprompt@gap>| !gapinput@A:= FullMatrixAlgebra( Rationals, 3 );|
  ( Rationals^[ 3, 3 ] )
  !gapprompt@gap>| !gapinput@V:= Rationals^[3,2];|
  ( Rationals^[ 3, 2 ] )
  !gapprompt@gap>| !gapinput@M:= LeftAlgebraModule( A, \*, V );|
  <left-module over ( Rationals^[ 3, 3 ] )>
  !gapprompt@gap>| !gapinput@bm:= Basis(M);;|
  !gapprompt@gap>| !gapinput@U:= SubAlgebraModule( M, [ bm[1] ] );|
  <left-module over ( Rationals^[ 3, 3 ] )>
  !gapprompt@gap>| !gapinput@TranslatorSubalgebra( M, U, M );|
  <algebra of dimension 9 over Rationals>
  !gapprompt@gap>| !gapinput@W:= SubAlgebraModule( M, [ bm[4] ] );|
  <left-module over ( Rationals^[ 3, 3 ] )>
  !gapprompt@gap>| !gapinput@T:=TranslatorSubalgebra( M, U, W );|
  <algebra of dimension 0 over Rationals>
\end{Verbatim}
 }

 }

 }

            
\chapter{\textcolor{Chapter }{Finitely Presented Algebras}}\label{Finitely Presented Algebras}
\logpage{[ 63, 0, 0 ]}
\hyperdef{L}{X85A22A8286596D02}{}
{
  Currently the \textsf{GAP} library contains only few functions dealing with general finitely presented
algebras, so this chapter is merely a placeholder. 

 The special case of finitely presented \emph{Lie} algebras is described in{\nobreakspace}\ref{Finitely Presented Lie Algebras}, and there is also a \textsf{GAP} package \texttt{fplsa} for computing structure constants of \emph{f}initely \emph{p}resented \emph{L}ie (\emph{s}uper)\emph{a}lgebras. 

   }

            
\chapter{\textcolor{Chapter }{Lie Algebras}}\label{Lie Algebras}
\logpage{[ 64, 0, 0 ]}
\hyperdef{L}{X78559D4C800AF58A}{}
{
  A Lie algebra $L$ is an algebra such that $x x = 0$ and $x(yz) + y(zx) + z(xy) = 0$ for all $x, y, z \in L$. A common way of creating a Lie algebra is by taking an associative algebra
together with the commutator as product. Therefore the product of two elements $x, y$ of a Lie algebra is usually denoted by $[x,y]$, but in \textsf{GAP} this denotes the list of the elements $x$ and $y$; hence the product of elements is made by the usual \texttt{*}. This gives no problems when dealing with Lie algebras given by a table of
structure constants. However, for matrix Lie algebras the situation is not so
easy as \texttt{*} denotes the ordinary (associative) matrix multiplication. In \textsf{GAP} this problem is solved by wrapping elements of a matrix Lie algebra up as \texttt{LieObject}s, and then define the \texttt{*} for \texttt{LieObject}s to be the commutator (see \ref{Lie Objects}).   
\section{\textcolor{Chapter }{Lie Objects}}\label{Lie Objects}
\logpage{[ 64, 1, 0 ]}
\hyperdef{L}{X80A607C47B7A2E69}{}
{
  Let \texttt{x} be a ring element, then \texttt{LieObject(x)} (see{\nobreakspace}\texttt{LieObject} (\ref{LieObject})) wraps \texttt{x} up into an object that contains the same data (namely \texttt{x}). The multiplication \texttt{*} for Lie objects is formed by taking the commutator. More exactly, if \texttt{l1} and \texttt{l2} are the Lie objects corresponding to the ring elements \texttt{r1} and \texttt{r2}, then \texttt{l1 * l2} is equal to the Lie object corresponding to \texttt{r1 * r2 \texttt{\symbol{45}} r2 * r1}. Two rules for Lie objects are worth noting: 

 
\begin{itemize}
\item  An element is \emph{not} equal to its Lie element. 
\item  If we take the Lie object of an ordinary (associative) matrix then this is
again a matrix; it is therefore a collection (of its rows) and a list. But it
is \emph{not} a collection of collections of its entries, and its family is \emph{not} a collections family. 
\end{itemize}
 

 Given a family \texttt{F} of ring elements, we can form its Lie family \texttt{L}. The elements of \texttt{F} and \texttt{L} are in bijection, only the multiplications via \texttt{*} differ for both families. More exactly, if \texttt{l1} and \texttt{l2} are the Lie elements corresponding to the elements \texttt{f1} and \texttt{f2} in \texttt{F}, we have \texttt{l1 * l2} equal to the Lie element corresponding to \texttt{f1 * f2 \texttt{\symbol{45}} f2 * f1}. Furthermore, the product of Lie elements \texttt{l1}, \texttt{l2} and \texttt{l3} is left\texttt{\symbol{45}}normed, that is \texttt{l1*l2*l3} is equal to \texttt{(l1*l2)*l3}. 

 The main reason to distinguish elements and Lie elements on the family level
is that this helps to avoid forming domains that contain elements of both
types. For example, if we could form vector spaces of matrices then at first
glance it would be no problem to have both ordinary and Lie matrices in it,
but as soon as we find out that the space is in fact an algebra (e.g., because
its dimension is that of the full matrix algebra), we would run into strange
problems. 

 Note that the family situation with Lie families may be not familiar. 

 
\begin{itemize}
\item  We have to be careful when installing methods for certain types of domains
that may involve Lie elements. For example, the zero element of a matrix space
is either an ordinary matrix or its Lie element, depending on the space. So
either the method must be aware of both cases, or the method selection must
distinguish the two cases. In the latter situation, only one method may be
applicable to each case; this means that it is not sufficient to treat the Lie
case with the additional requirement \texttt{IsLieObjectCollection} but that we must explicitly require non\texttt{\symbol{45}}Lie elements for
the non\texttt{\symbol{45}}Lie case. 
\item  Being a full matrix space is a property that may hold for a space of ordinary
matrices or a space of Lie matrices. So methods for full matrix spaces must
also be aware of Lie matrices. 
\end{itemize}
 

\subsection{\textcolor{Chapter }{LieObject}}
\logpage{[ 64, 1, 1 ]}\nobreak
\hyperdef{L}{X87F121978775AF48}{}
{\noindent\textcolor{FuncColor}{$\triangleright$\enspace\texttt{LieObject({\mdseries\slshape obj})\index{LieObject@\texttt{LieObject}}
\label{LieObject}
}\hfill{\scriptsize (attribute)}}\\


 Let \mbox{\texttt{\mdseries\slshape obj}} be a ring element. Then \texttt{LieObject( \mbox{\texttt{\mdseries\slshape obj}} )} is the corresponding Lie object. If \mbox{\texttt{\mdseries\slshape obj}} lies in the family \texttt{F}, then \texttt{LieObject( \mbox{\texttt{\mdseries\slshape obj}} )} lies in the family \texttt{LieFamily( F )} (see{\nobreakspace}\texttt{LieFamily} (\ref{LieFamily})). 
\begin{Verbatim}[commandchars=!@|,fontsize=\small,frame=single,label=Example]
  !gapprompt@gap>| !gapinput@m:= [ [ 1, 0 ], [ 0, 1 ] ];;|
  !gapprompt@gap>| !gapinput@lo:= LieObject( m );|
  LieObject( [ [ 1, 0 ], [ 0, 1 ] ] )
  !gapprompt@gap>| !gapinput@m*m;|
  [ [ 1, 0 ], [ 0, 1 ] ]
  !gapprompt@gap>| !gapinput@lo*lo;|
  LieObject( [ [ 0, 0 ], [ 0, 0 ] ] )
\end{Verbatim}
 }

 

\subsection{\textcolor{Chapter }{IsLieObject}}
\logpage{[ 64, 1, 2 ]}\nobreak
\hyperdef{L}{X83E5DD4381D9A65D}{}
{\noindent\textcolor{FuncColor}{$\triangleright$\enspace\texttt{IsLieObject({\mdseries\slshape obj})\index{IsLieObject@\texttt{IsLieObject}}
\label{IsLieObject}
}\hfill{\scriptsize (Category)}}\\
\noindent\textcolor{FuncColor}{$\triangleright$\enspace\texttt{IsLieObjectCollection({\mdseries\slshape obj})\index{IsLieObjectCollection@\texttt{IsLieObjectCollection}}
\label{IsLieObjectCollection}
}\hfill{\scriptsize (Category)}}\\
\noindent\textcolor{FuncColor}{$\triangleright$\enspace\texttt{IsRestrictedLieObject({\mdseries\slshape obj})\index{IsRestrictedLieObject@\texttt{IsRestrictedLieObject}}
\label{IsRestrictedLieObject}
}\hfill{\scriptsize (Category)}}\\
\noindent\textcolor{FuncColor}{$\triangleright$\enspace\texttt{IsRestrictedLieObjectCollection({\mdseries\slshape obj})\index{IsRestrictedLieObjectCollection@\texttt{IsRestrictedLieObjectCollection}}
\label{IsRestrictedLieObjectCollection}
}\hfill{\scriptsize (Category)}}\\


 An object lies in \texttt{IsLieObject} if and only if it lies in a family constructed by \texttt{LieFamily} (\ref{LieFamily}). 
\begin{Verbatim}[commandchars=!@|,fontsize=\small,frame=single,label=Example]
  !gapprompt@gap>| !gapinput@m:= [ [ 1, 0 ], [ 0, 1 ] ];;|
  !gapprompt@gap>| !gapinput@lo:= LieObject( m );|
  LieObject( [ [ 1, 0 ], [ 0, 1 ] ] )
  !gapprompt@gap>| !gapinput@IsLieObject( m );|
  false
  !gapprompt@gap>| !gapinput@IsLieObject( lo );|
  true
\end{Verbatim}
 }

 

\subsection{\textcolor{Chapter }{LieFamily}}
\logpage{[ 64, 1, 3 ]}\nobreak
\hyperdef{L}{X8725993C7BF386EE}{}
{\noindent\textcolor{FuncColor}{$\triangleright$\enspace\texttt{LieFamily({\mdseries\slshape Fam})\index{LieFamily@\texttt{LieFamily}}
\label{LieFamily}
}\hfill{\scriptsize (attribute)}}\\


 is a family \texttt{F} in bijection with the family \mbox{\texttt{\mdseries\slshape Fam}}, but with the Lie bracket as infix multiplication. That is, for \texttt{x}, \texttt{y} in \mbox{\texttt{\mdseries\slshape Fam}}, the product of the images in \texttt{F} will be the image of \texttt{x * y \texttt{\symbol{45}} y * x}. 

 The standard type of objects in a Lie family \texttt{F} is \texttt{\mbox{\texttt{\mdseries\slshape F}}!.packedType}. 

 \index{Embedding@\texttt{Embedding}!for Lie algebras} The bijection from \mbox{\texttt{\mdseries\slshape Fam}} to \texttt{F} is given by \texttt{Embedding( \mbox{\texttt{\mdseries\slshape Fam}}, F )} (see{\nobreakspace}\texttt{Embedding} (\ref{Embedding:for two domains})); this bijection respects addition and additive inverses. }

 

\subsection{\textcolor{Chapter }{UnderlyingFamily}}
\logpage{[ 64, 1, 4 ]}\nobreak
\hyperdef{L}{X81D9F5C6876FE93B}{}
{\noindent\textcolor{FuncColor}{$\triangleright$\enspace\texttt{UnderlyingFamily({\mdseries\slshape Fam})\index{UnderlyingFamily@\texttt{UnderlyingFamily}}
\label{UnderlyingFamily}
}\hfill{\scriptsize (attribute)}}\\


 If \mbox{\texttt{\mdseries\slshape Fam}} is a Lie family then \texttt{UnderlyingFamily( \mbox{\texttt{\mdseries\slshape Fam}} )} is a family \texttt{F} such that \texttt{\mbox{\texttt{\mdseries\slshape Fam}} = LieFamily( F )}. }

 

\subsection{\textcolor{Chapter }{UnderlyingRingElement}}
\logpage{[ 64, 1, 5 ]}\nobreak
\hyperdef{L}{X874B2B2A7F5A9A78}{}
{\noindent\textcolor{FuncColor}{$\triangleright$\enspace\texttt{UnderlyingRingElement({\mdseries\slshape obj})\index{UnderlyingRingElement@\texttt{UnderlyingRingElement}}
\label{UnderlyingRingElement}
}\hfill{\scriptsize (attribute)}}\\


 Let \mbox{\texttt{\mdseries\slshape obj}} be a Lie object constructed from a ring element \texttt{r} by calling \texttt{LieObject( r )}. Then \texttt{UnderlyingRingElement( \mbox{\texttt{\mdseries\slshape obj}} )} returns the ring element \texttt{r} used to construct \mbox{\texttt{\mdseries\slshape obj}}. If \texttt{r} lies in the family \texttt{F}, then \mbox{\texttt{\mdseries\slshape obj}} lies in the family \texttt{LieFamily( F )} (see{\nobreakspace}\texttt{LieFamily} (\ref{LieFamily})). 
\begin{Verbatim}[commandchars=!@|,fontsize=\small,frame=single,label=Example]
  !gapprompt@gap>| !gapinput@lo:= LieObject( [ [ 1, 0 ], [ 0, 1 ] ] );|
  LieObject( [ [ 1, 0 ], [ 0, 1 ] ] )
  !gapprompt@gap>| !gapinput@m:=UnderlyingRingElement(lo);|
  [ [ 1, 0 ], [ 0, 1 ] ]
  !gapprompt@gap>| !gapinput@lo*lo;|
  LieObject( [ [ 0, 0 ], [ 0, 0 ] ] )
  !gapprompt@gap>| !gapinput@m*m;|
  [ [ 1, 0 ], [ 0, 1 ] ]
\end{Verbatim}
 }

 }

  
\section{\textcolor{Chapter }{Constructing Lie algebras}}\label{Constructing Lie algebras}
\logpage{[ 64, 2, 0 ]}
\hyperdef{L}{X789A44F283C16B2B}{}
{
  In this section we describe functions that create Lie algebras. Creating and
working with subalgebras goes exactly in the same way as for general algebras;
so for that we refer to Chapter \ref{Algebras}. 

\subsection{\textcolor{Chapter }{LieAlgebraByStructureConstants}}
\logpage{[ 64, 2, 1 ]}\nobreak
\hyperdef{L}{X7D362350824FA115}{}
{\noindent\textcolor{FuncColor}{$\triangleright$\enspace\texttt{LieAlgebraByStructureConstants({\mdseries\slshape R, sct[, nameinfo]})\index{LieAlgebraByStructureConstants@\texttt{LieAlgebraByStructureConstants}}
\label{LieAlgebraByStructureConstants}
}\hfill{\scriptsize (function)}}\\


 \texttt{LieAlgebraByStructureConstants} does the same as \texttt{AlgebraByStructureConstants} (\ref{AlgebraByStructureConstants}), and has the same meaning of arguments, except that the result is assumed to
be a Lie algebra. Note that the function does not check whether \mbox{\texttt{\mdseries\slshape sct}} satisfies the Jacobi identity. (So if one creates a Lie algebra this way with
a table that does not satisfy the Jacobi identity, errors may occur later on.) 
\begin{Verbatim}[commandchars=!@|,fontsize=\small,frame=single,label=Example]
  !gapprompt@gap>| !gapinput@T:= EmptySCTable( 2, 0, "antisymmetric" );;|
  !gapprompt@gap>| !gapinput@SetEntrySCTable( T, 1, 2, [ 1/2, 1 ] );|
  !gapprompt@gap>| !gapinput@L:= LieAlgebraByStructureConstants( Rationals, T );|
  <Lie algebra of dimension 2 over Rationals>
\end{Verbatim}
 }

 

\subsection{\textcolor{Chapter }{RestrictedLieAlgebraByStructureConstants}}
\logpage{[ 64, 2, 2 ]}\nobreak
\hyperdef{L}{X7EEB79EE855E124C}{}
{\noindent\textcolor{FuncColor}{$\triangleright$\enspace\texttt{RestrictedLieAlgebraByStructureConstants({\mdseries\slshape R, sct[, nameinfo], pmapping})\index{RestrictedLieAlgebraByStructureConstants@\texttt{Restricted}\-\texttt{Lie}\-\texttt{Algebra}\-\texttt{By}\-\texttt{Structure}\-\texttt{Constants}}
\label{RestrictedLieAlgebraByStructureConstants}
}\hfill{\scriptsize (function)}}\\


 \texttt{RestrictedLieAlgebraByStructureConstants} does the same as \texttt{LieAlgebraByStructureConstants} (\ref{LieAlgebraByStructureConstants}), and has the same meaning of all arguments, except that the result is assumed
to be a restricted Lie algebra (see \ref{Restricted Lie algebras}) with the $p$\texttt{\symbol{45}}map given by the additional argument \mbox{\texttt{\mdseries\slshape pmapping}}. This last argument is a list of the length equal to the dimension of the
algebra; its $i$\texttt{\symbol{45}}th entry specifies the $p$\texttt{\symbol{45}}th power of the $i$\texttt{\symbol{45}}th basis vector in the same format \texttt{[ coeff1, position1, coeff2, position2, ... ]} as \texttt{SetEntrySCTable} (\ref{SetEntrySCTable}) uses to specify entries of the structure constants table. 

 Note that the function does not check whether \mbox{\texttt{\mdseries\slshape sct}} satisfies the Jacobi identity, of whether \mbox{\texttt{\mdseries\slshape pmapping}} specifies a legitimate $p$\texttt{\symbol{45}}mapping. 

 The following example creates a commutative restricted Lie algebra of
dimension 3, in which the $p$\texttt{\symbol{45}}th power of the $i$\texttt{\symbol{45}}th basis element is the $i+1$\texttt{\symbol{45}}th basis element (except for the 3rd basis element which
goes to zero). 
\begin{Verbatim}[commandchars=!@|,fontsize=\small,frame=single,label=Example]
  !gapprompt@gap>| !gapinput@T:= EmptySCTable( 3, Zero(GF(5)), "antisymmetric" );;|
  !gapprompt@gap>| !gapinput@L:= RestrictedLieAlgebraByStructureConstants(|
  !gapprompt@>| !gapinput@                                    GF(5), T, [[1,2],[1,3],[]] );|
  <Lie algebra of dimension 3 over GF(5)>
  !gapprompt@gap>| !gapinput@List(Basis(L),PthPowerImage);|
  [ v.2, v.3, 0*v.1 ]
  !gapprompt@gap>| !gapinput@PthPowerImage(L.1+L.2);|
  v.2+v.3
\end{Verbatim}
 }

 

\subsection{\textcolor{Chapter }{LieAlgebra (for an associative algebra)}}
\logpage{[ 64, 2, 3 ]}\nobreak
\hyperdef{L}{X7C840A9F85D28C81}{}
{\noindent\textcolor{FuncColor}{$\triangleright$\enspace\texttt{LieAlgebra({\mdseries\slshape L})\index{LieAlgebra@\texttt{LieAlgebra}!for an associative algebra}
\label{LieAlgebra:for an associative algebra}
}\hfill{\scriptsize (function)}}\\
\noindent\textcolor{FuncColor}{$\triangleright$\enspace\texttt{LieAlgebra({\mdseries\slshape F, gens[, zero][, "basis"]})\index{LieAlgebra@\texttt{LieAlgebra}!for field and generators}
\label{LieAlgebra:for field and generators}
}\hfill{\scriptsize (function)}}\\


 For an associative algebra \mbox{\texttt{\mdseries\slshape L}}, \texttt{LieAlgebra( \mbox{\texttt{\mdseries\slshape L}} )} is the Lie algebra isomorphic to \mbox{\texttt{\mdseries\slshape L}} as a vector space but with the Lie bracket as product. 

 \texttt{LieAlgebra( \mbox{\texttt{\mdseries\slshape F}}, \mbox{\texttt{\mdseries\slshape gens}} )} is the Lie algebra over the division ring \mbox{\texttt{\mdseries\slshape F}}, generated \emph{as Lie algebra} by the Lie objects corresponding to the vectors in the list \mbox{\texttt{\mdseries\slshape gens}}. 

 \emph{Note} that the algebra returned by \texttt{LieAlgebra} does not contain the vectors in \mbox{\texttt{\mdseries\slshape gens}}. The elements in \mbox{\texttt{\mdseries\slshape gens}} are wrapped up as Lie objects (see \ref{Lie Objects}). This allows one to create Lie algebras from ring elements with respect to
the Lie bracket as product. But of course the product in the Lie algebra is
the usual \texttt{*}. 

 If there are three arguments, a division ring \mbox{\texttt{\mdseries\slshape F}} and a list \mbox{\texttt{\mdseries\slshape gens}} and an element \mbox{\texttt{\mdseries\slshape zero}}, then \texttt{LieAlgebra( \mbox{\texttt{\mdseries\slshape F}}, \mbox{\texttt{\mdseries\slshape gens}}, \mbox{\texttt{\mdseries\slshape zero}} )} is the corresponding \mbox{\texttt{\mdseries\slshape F}}\texttt{\symbol{45}}Lie algebra with zero element the Lie object corresponding
to \mbox{\texttt{\mdseries\slshape zero}}. 

 If the last argument is the string \texttt{"basis"} then the vectors in \mbox{\texttt{\mdseries\slshape gens}} are known to form a basis of the algebra (as an \mbox{\texttt{\mdseries\slshape F}}\texttt{\symbol{45}}vector space). 

 \emph{Note} that even if each element in \mbox{\texttt{\mdseries\slshape gens}} is already a Lie element, i.e., is of the form \texttt{LieElement( \mbox{\texttt{\mdseries\slshape elm}} )} for an object \mbox{\texttt{\mdseries\slshape elm}}, the elements of the result lie in the Lie family of the family that contains \mbox{\texttt{\mdseries\slshape gens}} as a subset. 
\begin{Verbatim}[commandchars=!@|,fontsize=\small,frame=single,label=Example]
  !gapprompt@gap>| !gapinput@A:= FullMatrixAlgebra( GF( 7 ), 4 );;|
  !gapprompt@gap>| !gapinput@L:= LieAlgebra( A );|
  <Lie algebra of dimension 16 over GF(7)>
  !gapprompt@gap>| !gapinput@mats:= [ [ [ 1, 0 ], [ 0, -1 ] ], [ [ 0, 1 ], [ 0, 0 ] ],|
  !gapprompt@>| !gapinput@            [ [ 0, 0 ], [ 1, 0] ] ];;|
  !gapprompt@gap>| !gapinput@L:= LieAlgebra( Rationals, mats );|
  <Lie algebra over Rationals, with 3 generators>
\end{Verbatim}
 }

 

\subsection{\textcolor{Chapter }{FreeLieAlgebra (for ring, rank (and name))}}
\logpage{[ 64, 2, 4 ]}\nobreak
\hyperdef{L}{X7F7B34BD80F0F1C8}{}
{\noindent\textcolor{FuncColor}{$\triangleright$\enspace\texttt{FreeLieAlgebra({\mdseries\slshape R, rank[, name]})\index{FreeLieAlgebra@\texttt{FreeLieAlgebra}!for ring, rank (and name)}
\label{FreeLieAlgebra:for ring, rank (and name)}
}\hfill{\scriptsize (function)}}\\
\noindent\textcolor{FuncColor}{$\triangleright$\enspace\texttt{FreeLieAlgebra({\mdseries\slshape R, name1, name2, ...})\index{FreeLieAlgebra@\texttt{FreeLieAlgebra}!for ring and several names}
\label{FreeLieAlgebra:for ring and several names}
}\hfill{\scriptsize (function)}}\\


 Returns a free Lie algebra of rank \mbox{\texttt{\mdseries\slshape rank}} over the ring \mbox{\texttt{\mdseries\slshape R}}. \texttt{FreeLieAlgebra( \mbox{\texttt{\mdseries\slshape R}}, \mbox{\texttt{\mdseries\slshape name1}}, \mbox{\texttt{\mdseries\slshape name2}},...)} returns a free Lie algebra over \mbox{\texttt{\mdseries\slshape R}} with generators named \mbox{\texttt{\mdseries\slshape name1}}, \mbox{\texttt{\mdseries\slshape name2}}, and so on. The elements of a free Lie algebra are written on the
Hall\texttt{\symbol{45}}Lyndon basis. 
\begin{Verbatim}[commandchars=!@|,fontsize=\small,frame=single,label=Example]
  !gapprompt@gap>| !gapinput@L:= FreeLieAlgebra( Rationals, "x", "y", "z" );|
  <Lie algebra over Rationals, with 3 generators>
  !gapprompt@gap>| !gapinput@g:= GeneratorsOfAlgebra( L );; x:= g[1];; y:=g[2];; z:= g[3];;|
  !gapprompt@gap>| !gapinput@z*(y*(x*(z*y)));|
  (-1)*((x*(y*z))*(y*z))+(-1)*((x*((y*z)*z))*y)+(-1)*(((x*z)*(y*z))*y)
\end{Verbatim}
 }

 

\subsection{\textcolor{Chapter }{FullMatrixLieAlgebra}}
\logpage{[ 64, 2, 5 ]}\nobreak
\hyperdef{L}{X8735EE937A0081F0}{}
{\noindent\textcolor{FuncColor}{$\triangleright$\enspace\texttt{FullMatrixLieAlgebra({\mdseries\slshape R, n})\index{FullMatrixLieAlgebra@\texttt{FullMatrixLieAlgebra}}
\label{FullMatrixLieAlgebra}
}\hfill{\scriptsize (function)}}\\
\noindent\textcolor{FuncColor}{$\triangleright$\enspace\texttt{MatrixLieAlgebra({\mdseries\slshape R, n})\index{MatrixLieAlgebra@\texttt{MatrixLieAlgebra}}
\label{MatrixLieAlgebra}
}\hfill{\scriptsize (function)}}\\
\noindent\textcolor{FuncColor}{$\triangleright$\enspace\texttt{MatLieAlgebra({\mdseries\slshape R, n})\index{MatLieAlgebra@\texttt{MatLieAlgebra}}
\label{MatLieAlgebra}
}\hfill{\scriptsize (function)}}\\


 is the full matrix Lie algebra of $\mbox{\texttt{\mdseries\slshape n}} \times \mbox{\texttt{\mdseries\slshape n}}$ matrices over the ring \mbox{\texttt{\mdseries\slshape R}}, for a nonnegative integer \mbox{\texttt{\mdseries\slshape n}}. 
\begin{Verbatim}[commandchars=!@|,fontsize=\small,frame=single,label=Example]
  !gapprompt@gap>| !gapinput@FullMatrixLieAlgebra( GF(9), 10 );|
  <Lie algebra over GF(3^2), with 19 generators>
\end{Verbatim}
 }

 

\subsection{\textcolor{Chapter }{RightDerivations}}
\logpage{[ 64, 2, 6 ]}\nobreak
\hyperdef{L}{X821B6C197C08878B}{}
{\noindent\textcolor{FuncColor}{$\triangleright$\enspace\texttt{RightDerivations({\mdseries\slshape B})\index{RightDerivations@\texttt{RightDerivations}}
\label{RightDerivations}
}\hfill{\scriptsize (attribute)}}\\
\noindent\textcolor{FuncColor}{$\triangleright$\enspace\texttt{LeftDerivations({\mdseries\slshape B})\index{LeftDerivations@\texttt{LeftDerivations}}
\label{LeftDerivations}
}\hfill{\scriptsize (attribute)}}\\
\noindent\textcolor{FuncColor}{$\triangleright$\enspace\texttt{Derivations({\mdseries\slshape B})\index{Derivations@\texttt{Derivations}}
\label{Derivations}
}\hfill{\scriptsize (attribute)}}\\


 These functions all return the matrix Lie algebra of derivations of the
algebra $A$ with basis \mbox{\texttt{\mdseries\slshape B}}. 

 \texttt{RightDerivations( \mbox{\texttt{\mdseries\slshape B}} )} returns the algebra of derivations represented by their right action on the
algebra $A$. This means that with respect to the basis $B$ of $A$, the derivation $D$ is described by the matrix $[ d_{{i,j}} ]$ which means that $D$ maps the $i$\texttt{\symbol{45}}th basis element $b_i$ to $\sum_{{j = 1}}^n d_{{i,j}} b_j$. 

 \texttt{LeftDerivations( \mbox{\texttt{\mdseries\slshape B}} )} returns the Lie algebra of derivations represented by their left action on the
algebra $A$. So the matrices contained in the algebra output by \texttt{LeftDerivations( \mbox{\texttt{\mdseries\slshape B}} )} are the transposes of the matrices contained in the output of \texttt{RightDerivations( \mbox{\texttt{\mdseries\slshape B}} )}. 

 \texttt{Derivations} is just a synonym for \texttt{RightDerivations}. 
\begin{Verbatim}[commandchars=!@|,fontsize=\small,frame=single,label=Example]
  !gapprompt@gap>| !gapinput@A:= OctaveAlgebra( Rationals );|
  <algebra of dimension 8 over Rationals>
  !gapprompt@gap>| !gapinput@L:= Derivations( Basis( A ) );|
  <Lie algebra of dimension 14 over Rationals>
\end{Verbatim}
 }

 

\subsection{\textcolor{Chapter }{SimpleLieAlgebra}}
\logpage{[ 64, 2, 7 ]}\nobreak
\hyperdef{L}{X7933F05F7DE342AB}{}
{\noindent\textcolor{FuncColor}{$\triangleright$\enspace\texttt{SimpleLieAlgebra({\mdseries\slshape type, n, F})\index{SimpleLieAlgebra@\texttt{SimpleLieAlgebra}}
\label{SimpleLieAlgebra}
}\hfill{\scriptsize (function)}}\\


 This function constructs the simple Lie algebra of type given by the string \mbox{\texttt{\mdseries\slshape type}} and rank \mbox{\texttt{\mdseries\slshape n}} over the field \mbox{\texttt{\mdseries\slshape F}}. The string \mbox{\texttt{\mdseries\slshape type}} must be one of \texttt{"A"}, \texttt{"B"}, \texttt{"C"}, \texttt{"D"}, \texttt{"E"}, \texttt{"F"}, \texttt{"G"}, \texttt{"H"}, \texttt{"K"}, \texttt{"S"}, \texttt{"W"} or \texttt{"M"}. For the types \texttt{A} to \texttt{G}, \mbox{\texttt{\mdseries\slshape n}} must be a positive integer. The last five types only exist over fields of
characteristic $p>0$. If the type is \texttt{H}, then \mbox{\texttt{\mdseries\slshape n}} must be a list of positive integers of even length. If the type is \texttt{K}, then \mbox{\texttt{\mdseries\slshape n}} must be a list of positive integers of odd length. For the types \texttt{S} and \texttt{W}, \mbox{\texttt{\mdseries\slshape n}} must be a list of positive integers of any length. If the type is \texttt{M}, then the Melikyan algebra is constructed. In this case \mbox{\texttt{\mdseries\slshape n}} must be a list of two positive integers. This Lie algebra only exists over
fields of characteristic $5$. This Lie algebra is ${\ensuremath{\mathbb Z}} \times {\ensuremath{\mathbb Z}}$ graded; and the grading can be accessed via the attribute \texttt{Grading(L)} (see{\nobreakspace}\texttt{Grading} (\ref{Grading})). In some cases the Lie algebra returned by this function is not simple.
Examples are the Lie algebras of type $A_n$ over a field of characteristic $p>0$ where $p$ divides $n+1$, and the Lie algebras of type $K_n$ where $n$ is a list of length 1. 

 If \mbox{\texttt{\mdseries\slshape type}} is one of \texttt{A}, \texttt{B}, \texttt{C}, \texttt{D}, \texttt{E}, \texttt{F}, \texttt{G}, and \mbox{\texttt{\mdseries\slshape F}} is a field of characteristic zero, then the basis of the returned Lie algebra
is a Chevalley basis. 

 
\begin{Verbatim}[commandchars=!@|,fontsize=\small,frame=single,label=Example]
  !gapprompt@gap>| !gapinput@SimpleLieAlgebra( "E", 6, Rationals );|
  <Lie algebra of dimension 78 over Rationals>
  !gapprompt@gap>| !gapinput@SimpleLieAlgebra( "A", 6, GF(5) );|
  <Lie algebra of dimension 48 over GF(5)>
  !gapprompt@gap>| !gapinput@SimpleLieAlgebra( "W", [1,2], GF(5) );|
  <Lie algebra of dimension 250 over GF(5)>
  !gapprompt@gap>| !gapinput@SimpleLieAlgebra( "H", [1,2], GF(5) );|
  <Lie algebra of dimension 123 over GF(5)>
  !gapprompt@gap>| !gapinput@L:= SimpleLieAlgebra( "M", [1,1], GF(5) );|
  <Lie algebra of dimension 125 over GF(5)>
\end{Verbatim}
 }

 }

  
\section{\textcolor{Chapter }{Distinguished Subalgebras}}\label{Distinguished Subalgebras}
\logpage{[ 64, 3, 0 ]}
\hyperdef{L}{X798391F47E835F85}{}
{
  Here we describe functions that calculate well\texttt{\symbol{45}}known
subalgebras and ideals of a Lie algebra (such as the centre, the centralizer
of a subalgebra, etc.). 

\subsection{\textcolor{Chapter }{LieCentre}}
\logpage{[ 64, 3, 1 ]}\nobreak
\hyperdef{L}{X8111F58E7DE3E25C}{}
{\noindent\textcolor{FuncColor}{$\triangleright$\enspace\texttt{LieCentre({\mdseries\slshape L})\index{LieCentre@\texttt{LieCentre}}
\label{LieCentre}
}\hfill{\scriptsize (attribute)}}\\
\noindent\textcolor{FuncColor}{$\triangleright$\enspace\texttt{LieCenter({\mdseries\slshape L})\index{LieCenter@\texttt{LieCenter}}
\label{LieCenter}
}\hfill{\scriptsize (attribute)}}\\


 The \emph{Lie} centre of the Lie algebra \mbox{\texttt{\mdseries\slshape L}} is the kernel of the adjoint mapping, that is, the set $\{ a \in L : \forall x \in L: a x = 0 \}$. 

 In characteristic $2$ this may differ from the usual centre (that is the set of all $a \in L$ such that $a x = x a$ for all $x \in L$). Therefore, this operation is named \texttt{LieCentre} and not \texttt{Centre} (\ref{Centre}). 
\begin{Verbatim}[commandchars=!@|,fontsize=\small,frame=single,label=Example]
  !gapprompt@gap>| !gapinput@L:= FullMatrixLieAlgebra( GF(3), 3 );|
  <Lie algebra over GF(3), with 5 generators>
  !gapprompt@gap>| !gapinput@LieCentre( L );|
  <two-sided ideal in <Lie algebra of dimension 9 over GF(3)>,
    (dimension 1)>
\end{Verbatim}
 }

 

\subsection{\textcolor{Chapter }{LieCentralizer}}
\logpage{[ 64, 3, 2 ]}\nobreak
\hyperdef{L}{X811444717EEDCC34}{}
{\noindent\textcolor{FuncColor}{$\triangleright$\enspace\texttt{LieCentralizer({\mdseries\slshape L, S})\index{LieCentralizer@\texttt{LieCentralizer}}
\label{LieCentralizer}
}\hfill{\scriptsize (operation)}}\\


 is the annihilator of \mbox{\texttt{\mdseries\slshape S}} in the Lie algebra \mbox{\texttt{\mdseries\slshape L}}, that is, the set $\{ a \in L : \forall s \in S: a*s = 0 \}$. Here \mbox{\texttt{\mdseries\slshape S}} may be a subspace or a subalgebra of \mbox{\texttt{\mdseries\slshape L}}. 
\begin{Verbatim}[commandchars=!@|,fontsize=\small,frame=single,label=Example]
  !gapprompt@gap>| !gapinput@L:= SimpleLieAlgebra( "G", 2, Rationals );|
  <Lie algebra of dimension 14 over Rationals>
  !gapprompt@gap>| !gapinput@b:= BasisVectors( Basis( L ) );;|
  !gapprompt@gap>| !gapinput@LieCentralizer( L, Subalgebra( L, [ b[1], b[2] ] ) );|
  <Lie algebra of dimension 1 over Rationals>
\end{Verbatim}
 }

 

\subsection{\textcolor{Chapter }{LieNormalizer}}
\logpage{[ 64, 3, 3 ]}\nobreak
\hyperdef{L}{X7E62B6B37A75E09D}{}
{\noindent\textcolor{FuncColor}{$\triangleright$\enspace\texttt{LieNormalizer({\mdseries\slshape L, U})\index{LieNormalizer@\texttt{LieNormalizer}}
\label{LieNormalizer}
}\hfill{\scriptsize (operation)}}\\


 is the normalizer of the subspace \mbox{\texttt{\mdseries\slshape U}} in the Lie algebra \mbox{\texttt{\mdseries\slshape L}}, that is, the set $N_L(U) = \{ x \in L : [x,U] \subset U \}$. 
\begin{Verbatim}[commandchars=!@|,fontsize=\small,frame=single,label=Example]
  !gapprompt@gap>| !gapinput@L:= SimpleLieAlgebra( "G", 2, Rationals );|
  <Lie algebra of dimension 14 over Rationals>
  !gapprompt@gap>| !gapinput@b:= BasisVectors( Basis( L ) );;|
  !gapprompt@gap>| !gapinput@LieNormalizer( L, Subalgebra( L, [ b[1], b[2] ] ) );|
  <Lie algebra of dimension 8 over Rationals>
\end{Verbatim}
 }

 

\subsection{\textcolor{Chapter }{LieDerivedSubalgebra}}
\logpage{[ 64, 3, 4 ]}\nobreak
\hyperdef{L}{X7C95C0057C977747}{}
{\noindent\textcolor{FuncColor}{$\triangleright$\enspace\texttt{LieDerivedSubalgebra({\mdseries\slshape L})\index{LieDerivedSubalgebra@\texttt{LieDerivedSubalgebra}}
\label{LieDerivedSubalgebra}
}\hfill{\scriptsize (attribute)}}\\


 is the (Lie) derived subalgebra of the Lie algebra \mbox{\texttt{\mdseries\slshape L}}. 
\begin{Verbatim}[commandchars=!@|,fontsize=\small,frame=single,label=Example]
  !gapprompt@gap>| !gapinput@ L:= FullMatrixLieAlgebra( GF( 3 ), 3 );|
  <Lie algebra over GF(3), with 5 generators>
  !gapprompt@gap>| !gapinput@LieDerivedSubalgebra( L );|
  <Lie algebra of dimension 8 over GF(3)>
\end{Verbatim}
 }

 

\subsection{\textcolor{Chapter }{LieNilRadical}}
\logpage{[ 64, 3, 5 ]}\nobreak
\hyperdef{L}{X7D072F6D7A3D0BAF}{}
{\noindent\textcolor{FuncColor}{$\triangleright$\enspace\texttt{LieNilRadical({\mdseries\slshape L})\index{LieNilRadical@\texttt{LieNilRadical}}
\label{LieNilRadical}
}\hfill{\scriptsize (attribute)}}\\


 This function calculates the (Lie) nil radical of the Lie algebra \mbox{\texttt{\mdseries\slshape L}}. 

 
\begin{Verbatim}[commandchars=!@|,fontsize=\small,frame=single,label=Example]
  !gapprompt@gap>| !gapinput@mats:= [ [[1,0],[0,0]], [[0,1],[0,0]], [[0,0],[0,1]] ];;|
  !gapprompt@gap>| !gapinput@L:= LieAlgebra( Rationals, mats );;|
  !gapprompt@gap>| !gapinput@LieNilRadical( L );|
  <two-sided ideal in <Lie algebra of dimension 3 over Rationals>,
    (dimension 2)>
\end{Verbatim}
 }

 

\subsection{\textcolor{Chapter }{LieSolvableRadical}}
\logpage{[ 64, 3, 6 ]}\nobreak
\hyperdef{L}{X8445C9F17F7CBEA1}{}
{\noindent\textcolor{FuncColor}{$\triangleright$\enspace\texttt{LieSolvableRadical({\mdseries\slshape L})\index{LieSolvableRadical@\texttt{LieSolvableRadical}}
\label{LieSolvableRadical}
}\hfill{\scriptsize (attribute)}}\\


 Returns the (Lie) solvable radical of the Lie algebra \mbox{\texttt{\mdseries\slshape L}}. 
\begin{Verbatim}[commandchars=!@|,fontsize=\small,frame=single,label=Example]
  !gapprompt@gap>| !gapinput@L:= FullMatrixLieAlgebra( Rationals, 3 );;|
  !gapprompt@gap>| !gapinput@LieSolvableRadical( L );|
  <two-sided ideal in <Lie algebra of dimension 9 over Rationals>,
    (dimension 1)>
\end{Verbatim}
 }

 

\subsection{\textcolor{Chapter }{CartanSubalgebra}}
\logpage{[ 64, 3, 7 ]}\nobreak
\hyperdef{L}{X86114F157DFF6523}{}
{\noindent\textcolor{FuncColor}{$\triangleright$\enspace\texttt{CartanSubalgebra({\mdseries\slshape L})\index{CartanSubalgebra@\texttt{CartanSubalgebra}}
\label{CartanSubalgebra}
}\hfill{\scriptsize (attribute)}}\\


 A Cartan subalgebra of a Lie algebra \mbox{\texttt{\mdseries\slshape L}} is defined as a nilpotent subalgebra of \mbox{\texttt{\mdseries\slshape L}} equal to its own Lie normalizer in \mbox{\texttt{\mdseries\slshape L}}. 
\begin{Verbatim}[commandchars=!@|,fontsize=\small,frame=single,label=Example]
  !gapprompt@gap>| !gapinput@L:= SimpleLieAlgebra( "G", 2, Rationals );;|
  !gapprompt@gap>| !gapinput@CartanSubalgebra( L );|
  <Lie algebra of dimension 2 over Rationals>
\end{Verbatim}
 }

 }

  
\section{\textcolor{Chapter }{Series of Ideals}}\label{Series of Ideals}
\logpage{[ 64, 4, 0 ]}
\hyperdef{L}{X7A72840882F7A9B6}{}
{
  

\subsection{\textcolor{Chapter }{LieDerivedSeries}}
\logpage{[ 64, 4, 1 ]}\nobreak
\hyperdef{L}{X7DEF89A8869809F5}{}
{\noindent\textcolor{FuncColor}{$\triangleright$\enspace\texttt{LieDerivedSeries({\mdseries\slshape L})\index{LieDerivedSeries@\texttt{LieDerivedSeries}}
\label{LieDerivedSeries}
}\hfill{\scriptsize (attribute)}}\\


 is the (Lie) derived series of the Lie algebra \mbox{\texttt{\mdseries\slshape L}}. 
\begin{Verbatim}[commandchars=!@|,fontsize=\small,frame=single,label=Example]
  !gapprompt@gap>| !gapinput@mats:= [ [[1,0],[0,0]], [[0,1],[0,0]], [[0,0],[0,1]] ];;|
  !gapprompt@gap>| !gapinput@L:= LieAlgebra( Rationals, mats );;|
  !gapprompt@gap>| !gapinput@LieDerivedSeries( L );|
  [ <Lie algebra of dimension 3 over Rationals>,
    <Lie algebra of dimension 1 over Rationals>,
    <Lie algebra of dimension 0 over Rationals> ]
\end{Verbatim}
 }

 

\subsection{\textcolor{Chapter }{LieLowerCentralSeries}}
\logpage{[ 64, 4, 2 ]}\nobreak
\hyperdef{L}{X7900D17E7BA26A48}{}
{\noindent\textcolor{FuncColor}{$\triangleright$\enspace\texttt{LieLowerCentralSeries({\mdseries\slshape L})\index{LieLowerCentralSeries@\texttt{LieLowerCentralSeries}}
\label{LieLowerCentralSeries}
}\hfill{\scriptsize (attribute)}}\\


 is the (Lie) lower central series of the Lie algebra \mbox{\texttt{\mdseries\slshape L}}. 
\begin{Verbatim}[commandchars=!@|,fontsize=\small,frame=single,label=Example]
  !gapprompt@gap>| !gapinput@mats:= [ [[ 1, 0 ], [ 0, 0 ]], [[0,1],[0,0]], [[0,0],[0,1]] ];;|
  !gapprompt@gap>| !gapinput@L:=LieAlgebra( Rationals, mats );;|
  !gapprompt@gap>| !gapinput@LieLowerCentralSeries( L );|
  [ <Lie algebra of dimension 3 over Rationals>,
    <Lie algebra of dimension 1 over Rationals> ]
\end{Verbatim}
 }

 

\subsection{\textcolor{Chapter }{LieUpperCentralSeries}}
\logpage{[ 64, 4, 3 ]}\nobreak
\hyperdef{L}{X86A8701C868828C7}{}
{\noindent\textcolor{FuncColor}{$\triangleright$\enspace\texttt{LieUpperCentralSeries({\mdseries\slshape L})\index{LieUpperCentralSeries@\texttt{LieUpperCentralSeries}}
\label{LieUpperCentralSeries}
}\hfill{\scriptsize (attribute)}}\\


 is the (Lie) upper central series of the Lie algebra \mbox{\texttt{\mdseries\slshape L}}. 
\begin{Verbatim}[commandchars=!@|,fontsize=\small,frame=single,label=Example]
  !gapprompt@gap>| !gapinput@mats:= [ [[ 1, 0 ], [ 0, 0 ]], [[0,1],[0,0]], [[0,0],[0,1]] ];;|
  !gapprompt@gap>| !gapinput@L:=LieAlgebra( Rationals, mats );;|
  !gapprompt@gap>| !gapinput@LieUpperCentralSeries( L );|
  [ <two-sided ideal in <Lie algebra of dimension 3 over Rationals>,
        (dimension 1)>, <Lie algebra of dimension 0 over Rationals>
   ]
\end{Verbatim}
 }

 }

  
\section{\textcolor{Chapter }{Properties of a Lie Algebra}}\label{Properties of a Lie Algebra}
\logpage{[ 64, 5, 0 ]}
\hyperdef{L}{X8208CE5F8286155F}{}
{
  

\subsection{\textcolor{Chapter }{IsLieAbelian}}
\logpage{[ 64, 5, 1 ]}\nobreak
\hyperdef{L}{X7F97D08F7B738ADE}{}
{\noindent\textcolor{FuncColor}{$\triangleright$\enspace\texttt{IsLieAbelian({\mdseries\slshape L})\index{IsLieAbelian@\texttt{IsLieAbelian}}
\label{IsLieAbelian}
}\hfill{\scriptsize (property)}}\\


 returns \texttt{true} if \mbox{\texttt{\mdseries\slshape L}} is a Lie algebra such that each product of elements in \mbox{\texttt{\mdseries\slshape L}} is zero, and \texttt{false} otherwise. 
\begin{Verbatim}[commandchars=!@|,fontsize=\small,frame=single,label=Example]
  !gapprompt@gap>| !gapinput@ T:= EmptySCTable( 5, 0, "antisymmetric" );;|
  !gapprompt@gap>| !gapinput@ L:= LieAlgebraByStructureConstants( Rationals, T );|
  <Lie algebra of dimension 5 over Rationals>
  !gapprompt@gap>| !gapinput@IsLieAbelian( L );|
  true
\end{Verbatim}
 }

 

\subsection{\textcolor{Chapter }{IsLieNilpotent}}
\logpage{[ 64, 5, 2 ]}\nobreak
\hyperdef{L}{X78452F4E875A62A8}{}
{\noindent\textcolor{FuncColor}{$\triangleright$\enspace\texttt{IsLieNilpotent({\mdseries\slshape L})\index{IsLieNilpotent@\texttt{IsLieNilpotent}}
\label{IsLieNilpotent}
}\hfill{\scriptsize (property)}}\\


 A Lie algebra \mbox{\texttt{\mdseries\slshape L}} is defined to be (Lie) \emph{nilpotent} when its (Lie) lower central series reaches the trivial subalgebra. 
\begin{Verbatim}[commandchars=!@|,fontsize=\small,frame=single,label=Example]
  !gapprompt@gap>| !gapinput@T:= EmptySCTable( 5, 0, "antisymmetric" );;|
  !gapprompt@gap>| !gapinput@L:= LieAlgebraByStructureConstants( Rationals, T );|
  <Lie algebra of dimension 5 over Rationals>
  !gapprompt@gap>| !gapinput@IsLieNilpotent( L );|
  true
\end{Verbatim}
 }

 

\subsection{\textcolor{Chapter }{IsLieSolvable}}
\logpage{[ 64, 5, 3 ]}\nobreak
\hyperdef{L}{X859FF1B3812B8FCC}{}
{\noindent\textcolor{FuncColor}{$\triangleright$\enspace\texttt{IsLieSolvable({\mdseries\slshape L})\index{IsLieSolvable@\texttt{IsLieSolvable}}
\label{IsLieSolvable}
}\hfill{\scriptsize (property)}}\\


 A Lie algebra \mbox{\texttt{\mdseries\slshape L}} is defined to be (Lie) \emph{solvable} when its (Lie) derived series reaches the trivial subalgebra. 
\begin{Verbatim}[commandchars=!@|,fontsize=\small,frame=single,label=Example]
  !gapprompt@gap>| !gapinput@T:= EmptySCTable( 5, 0, "antisymmetric" );;|
  !gapprompt@gap>| !gapinput@L:= LieAlgebraByStructureConstants( Rationals, T );|
  <Lie algebra of dimension 5 over Rationals>
  !gapprompt@gap>| !gapinput@IsLieSolvable( L );|
  true
\end{Verbatim}
 }

 }

  
\section{\textcolor{Chapter }{Semisimple Lie Algebras and Root Systems}}\label{Semisimple Lie Algebras and Root Systems}
\logpage{[ 64, 6, 0 ]}
\hyperdef{L}{X83F829017D46C544}{}
{
  This section contains some functions for dealing with semisimple Lie algebras
and their root systems. 

\subsection{\textcolor{Chapter }{SemiSimpleType}}
\logpage{[ 64, 6, 1 ]}\nobreak
\hyperdef{L}{X8401CDC2859F8A85}{}
{\noindent\textcolor{FuncColor}{$\triangleright$\enspace\texttt{SemiSimpleType({\mdseries\slshape L})\index{SemiSimpleType@\texttt{SemiSimpleType}}
\label{SemiSimpleType}
}\hfill{\scriptsize (attribute)}}\\


 Let \mbox{\texttt{\mdseries\slshape L}} be a semisimple Lie algebra, i.e., a direct sum of simple Lie algebras. Then \texttt{SemiSimpleType} returns the type of \mbox{\texttt{\mdseries\slshape L}}, i.e., a string containing the types of the simple summands of \mbox{\texttt{\mdseries\slshape L}}. The field of \mbox{\texttt{\mdseries\slshape L}} must not have characteristic 2 or 3, and the Killing form of \mbox{\texttt{\mdseries\slshape L}} must be nondegenerate. If \mbox{\texttt{\mdseries\slshape L}} has characteristic 0, the structure constants of \mbox{\texttt{\mdseries\slshape L}} must be rational numbers; Lie algebras over arbitrary characteristic 0 fields
are not supported. 
\begin{Verbatim}[commandchars=!@|,fontsize=\small,frame=single,label=Example]
  !gapprompt@gap>| !gapinput@L:= SimpleLieAlgebra( "E", 8, Rationals );;|
  !gapprompt@gap>| !gapinput@b:= BasisVectors( Basis( L ) );;|
  !gapprompt@gap>| !gapinput@K:= LieCentralizer(L, Subalgebra(L, [ b[61]+b[79]+b[101]+b[102] ]));|
  <Lie algebra of dimension 102 over Rationals>
  !gapprompt@gap>| !gapinput@lev:= LeviMalcevDecomposition(K);;|
  !gapprompt@gap>| !gapinput@SemiSimpleType( lev[1] );|
  "B3 A1"
\end{Verbatim}
 }

 

\subsection{\textcolor{Chapter }{ChevalleyBasis}}
\logpage{[ 64, 6, 2 ]}\nobreak
\hyperdef{L}{X82EBF10A7B3B6F6E}{}
{\noindent\textcolor{FuncColor}{$\triangleright$\enspace\texttt{ChevalleyBasis({\mdseries\slshape L})\index{ChevalleyBasis@\texttt{ChevalleyBasis}}
\label{ChevalleyBasis}
}\hfill{\scriptsize (attribute)}}\\


 Here \mbox{\texttt{\mdseries\slshape L}} must be a semisimple Lie algebra with a split Cartan subalgebra. Then \texttt{ChevalleyBasis(\mbox{\texttt{\mdseries\slshape L}})} returns a list consisting of three sublists. Together these sublists form a
Chevalley basis of \mbox{\texttt{\mdseries\slshape L}}. The first list contains the positive root vectors, the second list contains
the negative root vectors, and the third list the Cartan elements of the
Chevalley basis. 
\begin{Verbatim}[commandchars=!@|,fontsize=\small,frame=single,label=Example]
  !gapprompt@gap>| !gapinput@L:= SimpleLieAlgebra( "G", 2, Rationals );|
  <Lie algebra of dimension 14 over Rationals>
  !gapprompt@gap>| !gapinput@ChevalleyBasis( L );|
  [ [ v.1, v.2, v.3, v.4, v.5, v.6 ],
    [ v.7, v.8, v.9, v.10, v.11, v.12 ], [ v.13, v.14 ] ]
\end{Verbatim}
 }

 

\subsection{\textcolor{Chapter }{IsRootSystem}}
\logpage{[ 64, 6, 3 ]}\nobreak
\hyperdef{L}{X79B5D27681193625}{}
{\noindent\textcolor{FuncColor}{$\triangleright$\enspace\texttt{IsRootSystem({\mdseries\slshape obj})\index{IsRootSystem@\texttt{IsRootSystem}}
\label{IsRootSystem}
}\hfill{\scriptsize (Category)}}\\


 Category of root systems. }

 

\subsection{\textcolor{Chapter }{IsRootSystemFromLieAlgebra}}
\logpage{[ 64, 6, 4 ]}\nobreak
\hyperdef{L}{X7D64D49479CBB203}{}
{\noindent\textcolor{FuncColor}{$\triangleright$\enspace\texttt{IsRootSystemFromLieAlgebra({\mdseries\slshape obj})\index{IsRootSystemFromLieAlgebra@\texttt{IsRootSystemFromLieAlgebra}}
\label{IsRootSystemFromLieAlgebra}
}\hfill{\scriptsize (Category)}}\\


 Category of root systems that come from (semisimple) Lie algebras. They often
have special attributes such as \texttt{UnderlyingLieAlgebra} (\ref{UnderlyingLieAlgebra}), \texttt{PositiveRootVectors} (\ref{PositiveRootVectors}), \texttt{NegativeRootVectors} (\ref{NegativeRootVectors}), \texttt{CanonicalGenerators} (\ref{CanonicalGenerators}). }

 

\subsection{\textcolor{Chapter }{RootSystem}}
\logpage{[ 64, 6, 5 ]}\nobreak
\hyperdef{L}{X80D15C027BB8029B}{}
{\noindent\textcolor{FuncColor}{$\triangleright$\enspace\texttt{RootSystem({\mdseries\slshape L})\index{RootSystem@\texttt{RootSystem}}
\label{RootSystem}
}\hfill{\scriptsize (attribute)}}\\


 \texttt{RootSystem} calculates the root system of the semisimple Lie algebra \mbox{\texttt{\mdseries\slshape L}} with a split Cartan subalgebra. 
\begin{Verbatim}[commandchars=!@|,fontsize=\small,frame=single,label=Example]
  !gapprompt@gap>| !gapinput@L:= SimpleLieAlgebra( "G", 2, Rationals );|
  <Lie algebra of dimension 14 over Rationals>
  !gapprompt@gap>| !gapinput@R:= RootSystem( L );|
  <root system of rank 2>
  !gapprompt@gap>| !gapinput@IsRootSystem( R );|
  true
  !gapprompt@gap>| !gapinput@IsRootSystemFromLieAlgebra( R );|
  true
\end{Verbatim}
 }

 

\subsection{\textcolor{Chapter }{UnderlyingLieAlgebra}}
\logpage{[ 64, 6, 6 ]}\nobreak
\hyperdef{L}{X7CA021E28527763E}{}
{\noindent\textcolor{FuncColor}{$\triangleright$\enspace\texttt{UnderlyingLieAlgebra({\mdseries\slshape R})\index{UnderlyingLieAlgebra@\texttt{UnderlyingLieAlgebra}}
\label{UnderlyingLieAlgebra}
}\hfill{\scriptsize (attribute)}}\\


 For a root system \mbox{\texttt{\mdseries\slshape R}} coming from a semisimple Lie algebra \texttt{L}, returns the Lie algebra \texttt{L}. }

 

\subsection{\textcolor{Chapter }{PositiveRoots}}
\logpage{[ 64, 6, 7 ]}\nobreak
\hyperdef{L}{X7B6B0BBD8035D7E5}{}
{\noindent\textcolor{FuncColor}{$\triangleright$\enspace\texttt{PositiveRoots({\mdseries\slshape R})\index{PositiveRoots@\texttt{PositiveRoots}}
\label{PositiveRoots}
}\hfill{\scriptsize (attribute)}}\\


 The list of positive roots of the root system \mbox{\texttt{\mdseries\slshape R}}. }

 

\subsection{\textcolor{Chapter }{NegativeRoots}}
\logpage{[ 64, 6, 8 ]}\nobreak
\hyperdef{L}{X81F9E0E67DD2688F}{}
{\noindent\textcolor{FuncColor}{$\triangleright$\enspace\texttt{NegativeRoots({\mdseries\slshape R})\index{NegativeRoots@\texttt{NegativeRoots}}
\label{NegativeRoots}
}\hfill{\scriptsize (attribute)}}\\


 The list of negative roots of the root system \mbox{\texttt{\mdseries\slshape R}}. }

 

\subsection{\textcolor{Chapter }{PositiveRootVectors}}
\logpage{[ 64, 6, 9 ]}\nobreak
\hyperdef{L}{X829C78427A442C23}{}
{\noindent\textcolor{FuncColor}{$\triangleright$\enspace\texttt{PositiveRootVectors({\mdseries\slshape R})\index{PositiveRootVectors@\texttt{PositiveRootVectors}}
\label{PositiveRootVectors}
}\hfill{\scriptsize (attribute)}}\\


 A list of positive root vectors of the root system \mbox{\texttt{\mdseries\slshape R}} that comes from a Lie algebra \texttt{L}. This is a list in bijection with the list \texttt{PositiveRoots( L )} (see{\nobreakspace}\texttt{PositiveRoots} (\ref{PositiveRoots})). The root vector is a non\texttt{\symbol{45}}zero element of the root space
(in \texttt{L}) of the corresponding root. }

 

\subsection{\textcolor{Chapter }{NegativeRootVectors}}
\logpage{[ 64, 6, 10 ]}\nobreak
\hyperdef{L}{X7AB374DC87A39349}{}
{\noindent\textcolor{FuncColor}{$\triangleright$\enspace\texttt{NegativeRootVectors({\mdseries\slshape R})\index{NegativeRootVectors@\texttt{NegativeRootVectors}}
\label{NegativeRootVectors}
}\hfill{\scriptsize (attribute)}}\\


 A list of negative root vectors of the root system \mbox{\texttt{\mdseries\slshape R}} that comes from a Lie algebra \texttt{L}. This is a list in bijection with the list \texttt{NegativeRoots( L )} (see{\nobreakspace}\texttt{NegativeRoots} (\ref{NegativeRoots})). The root vector is a non\texttt{\symbol{45}}zero element of the root space
(in \texttt{L}) of the corresponding root. }

 

\subsection{\textcolor{Chapter }{SimpleSystem}}
\logpage{[ 64, 6, 11 ]}\nobreak
\hyperdef{L}{X7DBD179E7CCF6699}{}
{\noindent\textcolor{FuncColor}{$\triangleright$\enspace\texttt{SimpleSystem({\mdseries\slshape R})\index{SimpleSystem@\texttt{SimpleSystem}}
\label{SimpleSystem}
}\hfill{\scriptsize (attribute)}}\\


 A list of simple roots of the root system \mbox{\texttt{\mdseries\slshape R}}. }

 

\subsection{\textcolor{Chapter }{CartanMatrix}}
\logpage{[ 64, 6, 12 ]}\nobreak
\hyperdef{L}{X84E3FEF587CB66C3}{}
{\noindent\textcolor{FuncColor}{$\triangleright$\enspace\texttt{CartanMatrix({\mdseries\slshape R})\index{CartanMatrix@\texttt{CartanMatrix}}
\label{CartanMatrix}
}\hfill{\scriptsize (attribute)}}\\


 The Cartan matrix of the root system \mbox{\texttt{\mdseries\slshape R}}, relative to the simple roots in \texttt{SimpleSystem( \mbox{\texttt{\mdseries\slshape R}} )} (see{\nobreakspace}\texttt{SimpleSystem} (\ref{SimpleSystem})). }

 

\subsection{\textcolor{Chapter }{BilinearFormMat}}
\logpage{[ 64, 6, 13 ]}\nobreak
\hyperdef{L}{X878644D68571BF44}{}
{\noindent\textcolor{FuncColor}{$\triangleright$\enspace\texttt{BilinearFormMat({\mdseries\slshape R})\index{BilinearFormMat@\texttt{BilinearFormMat}}
\label{BilinearFormMat}
}\hfill{\scriptsize (attribute)}}\\


 The matrix of the bilinear form of the root system \mbox{\texttt{\mdseries\slshape R}}. If we denote this matrix by $B$, then we have $B(i,j) = (\alpha_i, \alpha_j)$, where the $\alpha_i$ are the simple roots of \mbox{\texttt{\mdseries\slshape R}}. }

 

\subsection{\textcolor{Chapter }{CanonicalGenerators}}
\logpage{[ 64, 6, 14 ]}\nobreak
\hyperdef{L}{X7FAE45B37C5779A0}{}
{\noindent\textcolor{FuncColor}{$\triangleright$\enspace\texttt{CanonicalGenerators({\mdseries\slshape R})\index{CanonicalGenerators@\texttt{CanonicalGenerators}}
\label{CanonicalGenerators}
}\hfill{\scriptsize (attribute)}}\\


 Here \mbox{\texttt{\mdseries\slshape R}} must be a root system coming from a semisimple Lie algebra \texttt{L}. This function returns $3l$ generators of \mbox{\texttt{\mdseries\slshape L}}, $x_1, \ldots, x_l, y_1, \ldots, y_l, h_1, \ldots, h_l$, where $x_i$ lies in the root space corresponding to the $i$\texttt{\symbol{45}}th simple root of the root system of \mbox{\texttt{\mdseries\slshape L}}, $y_i$ lies in the root space corresponding to $-$ the $i$\texttt{\symbol{45}}th simple root, and the $h_i$ are elements of the Cartan subalgebra. These elements satisfy the relations $h_i * h_j = 0$, $x_i * y_j = \delta_{ij} h_i$, $h_j * x_i = c_{ij} x_i$, $h_j * y_i = -c_{ij} y_i$, where $c_{ij}$ is the entry of the Cartan matrix on position $ij$. 

 Also if $a$ is a root of the root system \mbox{\texttt{\mdseries\slshape R}} (so $a$ is a list of numbers), then we have the relation $h_i * x = a[i] x$, where $x$ is a root vector corresponding to $a$. 
\begin{Verbatim}[commandchars=!@|,fontsize=\small,frame=single,label=Example]
  !gapprompt@gap>| !gapinput@L:= SimpleLieAlgebra( "G", 2, Rationals );;|
  !gapprompt@gap>| !gapinput@R:= RootSystem( L );;|
  !gapprompt@gap>| !gapinput@UnderlyingLieAlgebra( R );|
  <Lie algebra of dimension 14 over Rationals>
  !gapprompt@gap>| !gapinput@PositiveRoots( R );|
  [ [ 2, -1 ], [ -3, 2 ], [ -1, 1 ], [ 1, 0 ], [ 3, -1 ], [ 0, 1 ] ]
  !gapprompt@gap>| !gapinput@x:= PositiveRootVectors( R );|
  [ v.1, v.2, v.3, v.4, v.5, v.6 ]
  !gapprompt@gap>| !gapinput@g:=CanonicalGenerators( R );|
  [ [ v.1, v.2 ], [ v.7, v.8 ], [ v.13, v.14 ] ]
  !gapprompt@gap>| !gapinput@g[3][1]*x[1];|
  (2)*v.1
  !gapprompt@gap>| !gapinput@g[3][2]*x[1];|
  (-1)*v.1
  !gapprompt@gap>| !gapinput@# i.e., x[1] is the root vector belonging to the root [ 2, -1 ]|
  !gapprompt@gap>| !gapinput@BilinearFormMat( R );|
  [ [ 1/12, -1/8 ], [ -1/8, 1/4 ] ]
\end{Verbatim}
 }

 }

  
\section{\textcolor{Chapter }{Semisimple Lie Algebras and Weyl Groups of Root Systems}}\label{Semisimple Lie Algebras and Weyl Groups of Root Systems}
\logpage{[ 64, 7, 0 ]}
\hyperdef{L}{X7945D07786D1C4BB}{}
{
  This section deals with the Weyl group of a root system. A Weyl group is
represented by its action on the weight lattice. A \emph{weight} is by definition a linear function $\lambda: H \rightarrow F$ (where $F$ is the ground field), such that the values $\lambda(h_i)$ are all integers (where the $h_i$ are the Cartan elements of the \texttt{CanonicalGenerators} (\ref{CanonicalGenerators})). On the other hand each weight is determined by these values. Therefore we
represent a weight by a vector of integers; the $i$\texttt{\symbol{45}}th entry of this vector is the value $\lambda(h_i)$. Now the elements of the Weyl group are represented by matrices, and if \texttt{g} is an element of a Weyl group and \texttt{w} a weight, then \texttt{w*g} gives the result of applying \texttt{g} to \texttt{w}. Another way of applying the $i$\texttt{\symbol{45}}th simple reflection to a weight is by using the function \texttt{ApplySimpleReflection} (\ref{ApplySimpleReflection}). 

 A Weyl group is generated by the simple reflections. So \texttt{GeneratorsOfGroup} (\ref{GeneratorsOfGroup}) for a Weyl group \texttt{W} gives a list of matrices and the $i$\texttt{\symbol{45}}th entry of this list is the simple reflection
corresponding to the $i$\texttt{\symbol{45}}th simple root of the corresponding root system. 

\subsection{\textcolor{Chapter }{IsWeylGroup}}
\logpage{[ 64, 7, 1 ]}\nobreak
\hyperdef{L}{X82AA29DD7969A935}{}
{\noindent\textcolor{FuncColor}{$\triangleright$\enspace\texttt{IsWeylGroup({\mdseries\slshape G})\index{IsWeylGroup@\texttt{IsWeylGroup}}
\label{IsWeylGroup}
}\hfill{\scriptsize (property)}}\\


 A Weyl group is a group generated by reflections, with the attribute \texttt{SparseCartanMatrix} (\ref{SparseCartanMatrix}) set. }

 

\subsection{\textcolor{Chapter }{SparseCartanMatrix}}
\logpage{[ 64, 7, 2 ]}\nobreak
\hyperdef{L}{X81EF01E57E5DC18A}{}
{\noindent\textcolor{FuncColor}{$\triangleright$\enspace\texttt{SparseCartanMatrix({\mdseries\slshape W})\index{SparseCartanMatrix@\texttt{SparseCartanMatrix}}
\label{SparseCartanMatrix}
}\hfill{\scriptsize (attribute)}}\\


 This is a sparse form of the Cartan matrix of the corresponding root system.
If we denote the Cartan matrix by \texttt{C}, then the sparse Cartan matrix of \mbox{\texttt{\mdseries\slshape W}} is a list (of length equal to the length of the Cartan matrix), where the \texttt{i}\texttt{\symbol{45}}th entry is a list consisting of elements \texttt{[ j, C[i][j] ]}, where \texttt{j} is such that \texttt{C[i][j]} is non\texttt{\symbol{45}}zero. }

 

\subsection{\textcolor{Chapter }{WeylGroup}}
\logpage{[ 64, 7, 3 ]}\nobreak
\hyperdef{L}{X86BED5098322EBEF}{}
{\noindent\textcolor{FuncColor}{$\triangleright$\enspace\texttt{WeylGroup({\mdseries\slshape R})\index{WeylGroup@\texttt{WeylGroup}}
\label{WeylGroup}
}\hfill{\scriptsize (attribute)}}\\


 The Weyl group of the root system \mbox{\texttt{\mdseries\slshape R}}. It is generated by the simple reflections. A simple reflection is
represented by a matrix, and the result of letting a simple reflection \texttt{m} act on a weight \texttt{w} is obtained by \texttt{w*m}. 
\begin{Verbatim}[commandchars=!@|,fontsize=\small,frame=single,label=Example]
  !gapprompt@gap>| !gapinput@L:= SimpleLieAlgebra( "F", 4, Rationals );;|
  !gapprompt@gap>| !gapinput@R:= RootSystem( L );;|
  !gapprompt@gap>| !gapinput@W:= WeylGroup( R );|
  <matrix group with 4 generators>
  !gapprompt@gap>| !gapinput@IsWeylGroup( W );|
  true
  !gapprompt@gap>| !gapinput@SparseCartanMatrix( W );|
  [ [ [ 1, 2 ], [ 3, -1 ] ], [ [ 2, 2 ], [ 4, -1 ] ],
    [ [ 1, -1 ], [ 3, 2 ], [ 4, -1 ] ],
    [ [ 2, -1 ], [ 3, -2 ], [ 4, 2 ] ] ]
  !gapprompt@gap>| !gapinput@g:= GeneratorsOfGroup( W );;|
  !gapprompt@gap>| !gapinput@[ 1, 1, 1, 1 ]*g[2];|
  [ 1, -1, 1, 2 ]
\end{Verbatim}
 }

 

\subsection{\textcolor{Chapter }{ApplySimpleReflection}}
\logpage{[ 64, 7, 4 ]}\nobreak
\hyperdef{L}{X7829BC4D7F253649}{}
{\noindent\textcolor{FuncColor}{$\triangleright$\enspace\texttt{ApplySimpleReflection({\mdseries\slshape SC, i, wt})\index{ApplySimpleReflection@\texttt{ApplySimpleReflection}}
\label{ApplySimpleReflection}
}\hfill{\scriptsize (operation)}}\\


 Here \mbox{\texttt{\mdseries\slshape SC}} is the sparse Cartan matrix of a Weyl group. This function applies the \mbox{\texttt{\mdseries\slshape i}}\texttt{\symbol{45}}th simple reflection to the weight \mbox{\texttt{\mdseries\slshape wt}}, thus changing \mbox{\texttt{\mdseries\slshape wt}}. 
\begin{Verbatim}[commandchars=!@|,fontsize=\small,frame=single,label=Example]
  !gapprompt@gap>| !gapinput@L:= SimpleLieAlgebra( "F", 4, Rationals );;|
  !gapprompt@gap>| !gapinput@W:= WeylGroup( RootSystem( L ) );;|
  !gapprompt@gap>| !gapinput@C:= SparseCartanMatrix( W );;|
  !gapprompt@gap>| !gapinput@w:= [ 1, 1, 1, 1 ];;|
  !gapprompt@gap>| !gapinput@ApplySimpleReflection( C, 2, w );|
  !gapprompt@gap>| !gapinput@w;|
  [ 1, -1, 1, 2 ]
\end{Verbatim}
 }

 

\subsection{\textcolor{Chapter }{LongestWeylWordPerm}}
\logpage{[ 64, 7, 5 ]}\nobreak
\hyperdef{L}{X80A7204F7D40D80F}{}
{\noindent\textcolor{FuncColor}{$\triangleright$\enspace\texttt{LongestWeylWordPerm({\mdseries\slshape W})\index{LongestWeylWordPerm@\texttt{LongestWeylWordPerm}}
\label{LongestWeylWordPerm}
}\hfill{\scriptsize (attribute)}}\\


 Let $g_0$ be the longest element in the Weyl group \mbox{\texttt{\mdseries\slshape W}}, and let $\{ \alpha_1, \ldots, \alpha_l \}$ be a simple system of the corresponding root system. Then $g_0$ maps $\alpha_i$ to $-\alpha_{{\sigma(i)}}$, where $\sigma$ is a permutation of $(1, \ldots, l)$. This function returns that permutation. 
\begin{Verbatim}[commandchars=!@|,fontsize=\small,frame=single,label=Example]
  !gapprompt@gap>| !gapinput@L:= SimpleLieAlgebra( "E", 6, Rationals );;|
  !gapprompt@gap>| !gapinput@W:= WeylGroup( RootSystem( L ) );;|
  !gapprompt@gap>| !gapinput@LongestWeylWordPerm( W );|
  (1,6)(3,5)
\end{Verbatim}
 }

 

\subsection{\textcolor{Chapter }{ConjugateDominantWeight}}
\logpage{[ 64, 7, 6 ]}\nobreak
\hyperdef{L}{X7D4E213F82F73857}{}
{\noindent\textcolor{FuncColor}{$\triangleright$\enspace\texttt{ConjugateDominantWeight({\mdseries\slshape W, wt})\index{ConjugateDominantWeight@\texttt{ConjugateDominantWeight}}
\label{ConjugateDominantWeight}
}\hfill{\scriptsize (operation)}}\\
\noindent\textcolor{FuncColor}{$\triangleright$\enspace\texttt{ConjugateDominantWeightWithWord({\mdseries\slshape W, wt})\index{ConjugateDominantWeightWithWord@\texttt{ConjugateDominantWeightWithWord}}
\label{ConjugateDominantWeightWithWord}
}\hfill{\scriptsize (operation)}}\\


 Here \mbox{\texttt{\mdseries\slshape W}} is a Weyl group and \mbox{\texttt{\mdseries\slshape wt}} a weight (i.e., a list of integers). \texttt{ConjugateDominantWeight} returns the unique dominant weight conjugate to \mbox{\texttt{\mdseries\slshape wt}} under \mbox{\texttt{\mdseries\slshape W}}. 

 \texttt{ConjugateDominantWeightWithWord} returns a list of two elements. The first of these is the dominant weight
conjugate to \mbox{\texttt{\mdseries\slshape wt}}. The second element is a list of indices of simple reflections that have to
be applied to \mbox{\texttt{\mdseries\slshape wt}} in order to get the dominant weight conjugate to it. 
\begin{Verbatim}[commandchars=!@|,fontsize=\small,frame=single,label=Example]
  !gapprompt@gap>| !gapinput@L:= SimpleLieAlgebra( "E", 6, Rationals );;|
  !gapprompt@gap>| !gapinput@W:= WeylGroup( RootSystem( L ) );;|
  !gapprompt@gap>| !gapinput@C:= SparseCartanMatrix( W );;|
  !gapprompt@gap>| !gapinput@w:= [ 1, -1, 2, -2, 3, -3 ];;|
  !gapprompt@gap>| !gapinput@ConjugateDominantWeight( W, w );|
  [ 2, 1, 0, 0, 0, 0 ]
  !gapprompt@gap>| !gapinput@c:= ConjugateDominantWeightWithWord( W, w );|
  [ [ 2, 1, 0, 0, 0, 0 ], [ 2, 4, 2, 3, 6, 5, 4, 2, 3, 1 ] ]
  !gapprompt@gap>| !gapinput@for i in [1..Length(c[2])] do|
  !gapprompt@>| !gapinput@ApplySimpleReflection( C, c[2][i], w );|
  !gapprompt@>| !gapinput@od;|
  !gapprompt@gap>| !gapinput@w;|
  [ 2, 1, 0, 0, 0, 0 ]
\end{Verbatim}
 }

 

\subsection{\textcolor{Chapter }{WeylOrbitIterator}}
\logpage{[ 64, 7, 7 ]}\nobreak
\hyperdef{L}{X7E000FA97949BFD5}{}
{\noindent\textcolor{FuncColor}{$\triangleright$\enspace\texttt{WeylOrbitIterator({\mdseries\slshape W, wt})\index{WeylOrbitIterator@\texttt{WeylOrbitIterator}}
\label{WeylOrbitIterator}
}\hfill{\scriptsize (operation)}}\\


 Returns an iterator for the orbit of the weight \mbox{\texttt{\mdseries\slshape wt}} under the action of the Weyl group \mbox{\texttt{\mdseries\slshape W}}. 
\begin{Verbatim}[commandchars=!@|,fontsize=\small,frame=single,label=Example]
  !gapprompt@gap>| !gapinput@L:= SimpleLieAlgebra( "E", 6, Rationals );;|
  !gapprompt@gap>| !gapinput@W:= WeylGroup( RootSystem( L ) );;|
  !gapprompt@gap>| !gapinput@orb:= WeylOrbitIterator( W, [ 1, 1, 1, 1, 1, 1 ] );|
  <iterator>
  !gapprompt@gap>| !gapinput@NextIterator( orb );|
  [ 1, 1, 1, 1, 1, 1 ]
  !gapprompt@gap>| !gapinput@NextIterator( orb );|
  [ -1, -1, -1, -1, -1, -1 ]
  !gapprompt@gap>| !gapinput@orb:= WeylOrbitIterator( W, [ 1, 1, 1, 1, 1, 1 ] );|
  <iterator>
  !gapprompt@gap>| !gapinput@k:= 0;|
  0
  !gapprompt@gap>| !gapinput@while not IsDoneIterator( orb ) do|
  !gapprompt@>| !gapinput@w:= NextIterator( orb ); k:= k+1;|
  !gapprompt@>| !gapinput@od;|
  !gapprompt@gap>| !gapinput@k;  # this is the size of the Weyl group of E6|
  51840
\end{Verbatim}
 }

 }

  
\section{\textcolor{Chapter }{Restricted Lie algebras}}\label{Restricted Lie algebras}
\logpage{[ 64, 8, 0 ]}
\hyperdef{L}{X878080BB79BE3F2E}{}
{
  A Lie algebra $L$ over a field of characteristic $p>0$ is called restricted if there is a map $x \mapsto x^p$ from $L$ into $L$ (called a $p$\texttt{\symbol{45}}map) such that ad $x^p = ($ad$ x)^p$, $(\alpha x)^p = \alpha^p x^p$ and $(x+y)^p = x^p + y^p + \sum_{{i=1}}^{{p-1}} s_i(x,y)$, where $s_i: L \times L \rightarrow L$ are certain Lie polynomials in two variables. Using these relations we can
calculate $y^p$ for all $y \in L$, once we know $x^p$ for $x$ in a basis of $L$. Therefore a $p$\texttt{\symbol{45}}map is represented in \textsf{GAP}{\nobreakspace} by a list containing the images of the basis vectors of a
basis $B$ of $L$. For this reason this list is an attribute of the basis $B$. 

 

\subsection{\textcolor{Chapter }{IsRestrictedLieAlgebra}}
\logpage{[ 64, 8, 1 ]}\nobreak
\hyperdef{L}{X81F28B1D830F28EB}{}
{\noindent\textcolor{FuncColor}{$\triangleright$\enspace\texttt{IsRestrictedLieAlgebra({\mdseries\slshape L})\index{IsRestrictedLieAlgebra@\texttt{IsRestrictedLieAlgebra}}
\label{IsRestrictedLieAlgebra}
}\hfill{\scriptsize (property)}}\\


 Test whether \mbox{\texttt{\mdseries\slshape L}} is restricted. 
\begin{Verbatim}[commandchars=!@|,fontsize=\small,frame=single,label=Example]
  !gapprompt@gap>| !gapinput@L:= SimpleLieAlgebra( "W", [2], GF(5));|
  <Lie algebra of dimension 25 over GF(5)>
  !gapprompt@gap>| !gapinput@IsRestrictedLieAlgebra( L );|
  false
  !gapprompt@gap>| !gapinput@L:= SimpleLieAlgebra( "W", [1], GF(5));|
  <Lie algebra of dimension 5 over GF(5)>
  !gapprompt@gap>| !gapinput@IsRestrictedLieAlgebra( L );|
  true
\end{Verbatim}
 }

 

\subsection{\textcolor{Chapter }{PthPowerImages}}
\logpage{[ 64, 8, 2 ]}\nobreak
\hyperdef{L}{X7D7BD5908016461B}{}
{\noindent\textcolor{FuncColor}{$\triangleright$\enspace\texttt{PthPowerImages({\mdseries\slshape B})\index{PthPowerImages@\texttt{PthPowerImages}}
\label{PthPowerImages}
}\hfill{\scriptsize (attribute)}}\\


 Here \mbox{\texttt{\mdseries\slshape B}} is a basis of a restricted Lie algebra. This function returns the list of the
images of the basis vectors of \mbox{\texttt{\mdseries\slshape B}} under the $p$\texttt{\symbol{45}}map. 
\begin{Verbatim}[commandchars=!@|,fontsize=\small,frame=single,label=Example]
  !gapprompt@gap>| !gapinput@L:= SimpleLieAlgebra( "W", [1], GF(11) );|
  <Lie algebra of dimension 11 over GF(11)>
  !gapprompt@gap>| !gapinput@B:= Basis( L );|
  CanonicalBasis( <Lie algebra of dimension 11 over GF(11)> )
  !gapprompt@gap>| !gapinput@PthPowerImages( B );|
  [ 0*v.1, v.2, 0*v.1, 0*v.1, 0*v.1, 0*v.1, 0*v.1, 0*v.1, 0*v.1, 0*v.1,
    0*v.1 ]
\end{Verbatim}
 }

 

\subsection{\textcolor{Chapter }{PthPowerImage (for basis and element)}}
\logpage{[ 64, 8, 3 ]}\nobreak
\hyperdef{L}{X879BB01782E7D7A9}{}
{\noindent\textcolor{FuncColor}{$\triangleright$\enspace\texttt{PthPowerImage({\mdseries\slshape B, x})\index{PthPowerImage@\texttt{PthPowerImage}!for basis and element}
\label{PthPowerImage:for basis and element}
}\hfill{\scriptsize (operation)}}\\
\noindent\textcolor{FuncColor}{$\triangleright$\enspace\texttt{PthPowerImage({\mdseries\slshape x})\index{PthPowerImage@\texttt{PthPowerImage}!for element}
\label{PthPowerImage:for element}
}\hfill{\scriptsize (operation)}}\\
\noindent\textcolor{FuncColor}{$\triangleright$\enspace\texttt{PthPowerImage({\mdseries\slshape x, n})\index{PthPowerImage@\texttt{PthPowerImage}!for element and integer}
\label{PthPowerImage:for element and integer}
}\hfill{\scriptsize (operation)}}\\


 This function computes the image of an element \mbox{\texttt{\mdseries\slshape x}} of a restricted Lie algebra under its $p$\texttt{\symbol{45}}map. 

 In the first form, a basis of the Lie algebra is provided; this basis stores
the $p$th powers of its elements. It is the traditional form, provided for backwards
compatibility. 

 In its second form, only the element \mbox{\texttt{\mdseries\slshape x}} is provided. It is the only form for elements of Lie algebras with no
predetermined basis, such as those constructed by \texttt{LieObject} (\ref{LieObject}). 

 In its third form, an extra non\texttt{\symbol{45}}negative integer \mbox{\texttt{\mdseries\slshape n}} is specified; the $p$\texttt{\symbol{45}}mapping is iterated \mbox{\texttt{\mdseries\slshape n}} times on the element \mbox{\texttt{\mdseries\slshape x}}. 
\begin{Verbatim}[commandchars=!@|,fontsize=\small,frame=single,label=Example]
  !gapprompt@gap>| !gapinput@L:= SimpleLieAlgebra( "W", [1], GF(11) );;|
  !gapprompt@gap>| !gapinput@B:= Basis( L );;|
  !gapprompt@gap>| !gapinput@x:= B[1]+B[11];|
  v.1+v.11
  !gapprompt@gap>| !gapinput@PthPowerImage( B, x );|
  v.1+v.11
  !gapprompt@gap>| !gapinput@PthPowerImage( x, 2 );|
  v.1+v.11
  !gapprompt@gap>| !gapinput@f := FreeAssociativeAlgebra(GF(2),"x","y");|
  <algebra over GF(2), with 2 generators>
  !gapprompt@gap>| !gapinput@x := LieObject(f.1);; y := LieObject(f.2);;|
  !gapprompt@gap>| !gapinput@x*y; x^2; PthPowerImage(x);|
  LieObject( (Z(2)^0)*x*y+(Z(2)^0)*y*x )
  LieObject( <zero> of ... )
  LieObject( (Z(2)^0)*x^2 )
\end{Verbatim}
 }

 

\subsection{\textcolor{Chapter }{JenningsLieAlgebra}}
\logpage{[ 64, 8, 4 ]}\nobreak
\hyperdef{L}{X8692ADD581359CA1}{}
{\noindent\textcolor{FuncColor}{$\triangleright$\enspace\texttt{JenningsLieAlgebra({\mdseries\slshape G})\index{JenningsLieAlgebra@\texttt{JenningsLieAlgebra}}
\label{JenningsLieAlgebra}
}\hfill{\scriptsize (attribute)}}\\


 Let \mbox{\texttt{\mdseries\slshape G}} be a nontrivial $p$\texttt{\symbol{45}}group, and let $\mbox{\texttt{\mdseries\slshape G}} = G_1 \supset G_2 \supset \cdots \supset G_m = 1$ be its Jennings series (see{\nobreakspace}\texttt{JenningsSeries} (\ref{JenningsSeries})). Then the quotients $G_i / G_{{i+1}}$ are elementary abelian $p$\texttt{\symbol{45}}groups, i.e., they can be viewed as vector spaces over \texttt{GF}$(p)$. Now the Jennings\texttt{\symbol{45}}Lie algebra $L$ of \mbox{\texttt{\mdseries\slshape G}} is the direct sum of those vector spaces. The Lie bracket on $L$ is induced by the commutator in \mbox{\texttt{\mdseries\slshape G}}. Furthermore, the map $g \mapsto g^p$ in \mbox{\texttt{\mdseries\slshape G}} induces a $p$\texttt{\symbol{45}}map in $L$ making $L$ into a restricted Lie algebra. In the canonical basis of $L$ this $p$\texttt{\symbol{45}}map is added as an attribute. A Lie algebra created by \texttt{JenningsLieAlgebra} is naturally graded. The attribute \texttt{Grading} (\ref{Grading}) is set. }

 

\subsection{\textcolor{Chapter }{PCentralLieAlgebra}}
\logpage{[ 64, 8, 5 ]}\nobreak
\hyperdef{L}{X785251E879E1BFC6}{}
{\noindent\textcolor{FuncColor}{$\triangleright$\enspace\texttt{PCentralLieAlgebra({\mdseries\slshape G})\index{PCentralLieAlgebra@\texttt{PCentralLieAlgebra}}
\label{PCentralLieAlgebra}
}\hfill{\scriptsize (attribute)}}\\


 Here \mbox{\texttt{\mdseries\slshape G}} is a nontrivial $p$\texttt{\symbol{45}}group. \texttt{PCentralLieAlgebra( \mbox{\texttt{\mdseries\slshape G}} )} does the same as \texttt{JenningsLieAlgebra} (\ref{JenningsLieAlgebra}) except that the $p$\texttt{\symbol{45}}central series is used instead of the Jennings series
(see{\nobreakspace}\texttt{PCentralSeries} (\ref{PCentralSeries})). This function also returns a graded Lie algebra. However, it is not
necessarily restricted. 
\begin{Verbatim}[commandchars=!@|,fontsize=\small,frame=single,label=Example]
  !gapprompt@gap>| !gapinput@G:= SmallGroup( 3^6, 123 );|
  <pc group of size 729 with 6 generators>
  !gapprompt@gap>| !gapinput@L:= JenningsLieAlgebra( G );|
  <Lie algebra of dimension 6 over GF(3)>
  !gapprompt@gap>| !gapinput@HasPthPowerImages( Basis( L ) );|
  true
  !gapprompt@gap>| !gapinput@PthPowerImages( Basis( L ) );|
  [ v.6, 0*v.1, 0*v.1, 0*v.1, 0*v.1, 0*v.1 ]
  !gapprompt@gap>| !gapinput@g:= Grading( L );|
  rec( hom_components := function( d ) ... end, max_degree := 3,
    min_degree := 1, source := Integers )
  !gapprompt@gap>| !gapinput@List( [1,2,3], g.hom_components );|
  [ <vector space over GF(3), with 3 generators>,
    <vector space over GF(3), with 2 generators>,
    <vector space over GF(3), with 1 generator> ]
\end{Verbatim}
 }

 

\subsection{\textcolor{Chapter }{NaturalHomomorphismOfLieAlgebraFromNilpotentGroup}}
\logpage{[ 64, 8, 6 ]}\nobreak
\hyperdef{L}{X781ADBEC850C7DE7}{}
{\noindent\textcolor{FuncColor}{$\triangleright$\enspace\texttt{NaturalHomomorphismOfLieAlgebraFromNilpotentGroup({\mdseries\slshape L})\index{NaturalHomomorphismOfLieAlgebraFromNilpotentGroup@\texttt{Natural}\-\texttt{Homomorphism}\-\texttt{Of}\-\texttt{Lie}\-\texttt{Algebra}\-\texttt{From}\-\texttt{Nilpotent}\-\texttt{Group}}
\label{NaturalHomomorphismOfLieAlgebraFromNilpotentGroup}
}\hfill{\scriptsize (attribute)}}\\


 This is an attribute of Lie algebras created by \texttt{JenningsLieAlgebra} (\ref{JenningsLieAlgebra}) or \texttt{PCentralLieAlgebra} (\ref{PCentralLieAlgebra}). Then \mbox{\texttt{\mdseries\slshape L}} is the direct sum of quotients of successive terms of the Jennings, or $p$\texttt{\symbol{45}}central series of a $p$\texttt{\symbol{45}}group G. Let \texttt{Gi} be the $i$\texttt{\symbol{45}}th term in this series, and let \texttt{f = NaturalHomomorphismOfLieAlgebraFromNilpotentGroup( \mbox{\texttt{\mdseries\slshape L}} )}, then for \texttt{g} in \texttt{Gi}, \texttt{f( \mbox{\texttt{\mdseries\slshape g}}, \mbox{\texttt{\mdseries\slshape i}} )} returns the element of \mbox{\texttt{\mdseries\slshape L}} (lying in the $i$\texttt{\symbol{45}}th homogeneous component) corresponding to \texttt{g}. }

 }

  
\section{\textcolor{Chapter }{The Adjoint Representation}}\label{The Adjoint Representation}
\logpage{[ 64, 9, 0 ]}
\hyperdef{L}{X7C419FFA835EBE12}{}
{
  In this section we show functions for calculating with the adjoint
representation of a Lie algebra (and the corresponding trace form, called the
Killing form) (see also \texttt{AdjointBasis} (\ref{AdjointBasis}) and \texttt{IndicesOfAdjointBasis} (\ref{IndicesOfAdjointBasis})). 

 

\subsection{\textcolor{Chapter }{AdjointMatrix}}
\logpage{[ 64, 9, 1 ]}\nobreak
\hyperdef{L}{X786886D882795F78}{}
{\noindent\textcolor{FuncColor}{$\triangleright$\enspace\texttt{AdjointMatrix({\mdseries\slshape B, x})\index{AdjointMatrix@\texttt{AdjointMatrix}}
\label{AdjointMatrix}
}\hfill{\scriptsize (operation)}}\\


 is the matrix of the adjoint representation of the element \mbox{\texttt{\mdseries\slshape x}} w.r.t. the basis \mbox{\texttt{\mdseries\slshape B}}. The adjoint map is the left multiplication by \mbox{\texttt{\mdseries\slshape x}}. The $i$\texttt{\symbol{45}}th column of the resulting matrix represents the image of
the $i$\texttt{\symbol{45}}th basis vector of \mbox{\texttt{\mdseries\slshape B}} under left multiplication by \mbox{\texttt{\mdseries\slshape x}}. 
\begin{Verbatim}[commandchars=!@|,fontsize=\small,frame=single,label=Example]
  !gapprompt@gap>| !gapinput@L:= SimpleLieAlgebra( "A", 1, Rationals );;|
  !gapprompt@gap>| !gapinput@AdjointMatrix( Basis( L ), Basis( L )[1] );|
  [ [ 0, 0, -2 ], [ 0, 0, 0 ], [ 0, 1, 0 ] ]
\end{Verbatim}
 }

 

\subsection{\textcolor{Chapter }{AdjointAssociativeAlgebra}}
\logpage{[ 64, 9, 2 ]}\nobreak
\hyperdef{L}{X873A64307AC6C63E}{}
{\noindent\textcolor{FuncColor}{$\triangleright$\enspace\texttt{AdjointAssociativeAlgebra({\mdseries\slshape L, K})\index{AdjointAssociativeAlgebra@\texttt{AdjointAssociativeAlgebra}}
\label{AdjointAssociativeAlgebra}
}\hfill{\scriptsize (operation)}}\\


 is the associative matrix algebra (with 1) generated by the matrices of the
adjoint representation of the subalgebra \mbox{\texttt{\mdseries\slshape K}} on the Lie algebra \mbox{\texttt{\mdseries\slshape L}}. 
\begin{Verbatim}[commandchars=!@|,fontsize=\small,frame=single,label=Example]
  !gapprompt@gap>| !gapinput@L:= SimpleLieAlgebra( "A", 1, Rationals );;|
  !gapprompt@gap>| !gapinput@AdjointAssociativeAlgebra( L, L );|
  <algebra of dimension 9 over Rationals>
  !gapprompt@gap>| !gapinput@AdjointAssociativeAlgebra( L, CartanSubalgebra( L ) );|
  <algebra of dimension 3 over Rationals>
\end{Verbatim}
 }

 

\subsection{\textcolor{Chapter }{KillingMatrix}}
\logpage{[ 64, 9, 3 ]}\nobreak
\hyperdef{L}{X877CCFD5832E035D}{}
{\noindent\textcolor{FuncColor}{$\triangleright$\enspace\texttt{KillingMatrix({\mdseries\slshape B})\index{KillingMatrix@\texttt{KillingMatrix}}
\label{KillingMatrix}
}\hfill{\scriptsize (attribute)}}\\


 is the matrix of the Killing form $\kappa$ with respect to the basis \mbox{\texttt{\mdseries\slshape B}}, i.e., the matrix $( \kappa( b_i, b_j ) )$ where $b_1, b_2, \ldots$ are the basis vectors of \mbox{\texttt{\mdseries\slshape B}}. 
\begin{Verbatim}[commandchars=!@|,fontsize=\small,frame=single,label=Example]
  !gapprompt@gap>| !gapinput@L:= SimpleLieAlgebra( "A", 1, Rationals );;|
  !gapprompt@gap>| !gapinput@KillingMatrix( Basis( L ) );|
  [ [ 0, 4, 0 ], [ 4, 0, 0 ], [ 0, 0, 8 ] ]
\end{Verbatim}
 }

 

\subsection{\textcolor{Chapter }{KappaPerp}}
\logpage{[ 64, 9, 4 ]}\nobreak
\hyperdef{L}{X8234046083B60F6E}{}
{\noindent\textcolor{FuncColor}{$\triangleright$\enspace\texttt{KappaPerp({\mdseries\slshape L, U})\index{KappaPerp@\texttt{KappaPerp}}
\label{KappaPerp}
}\hfill{\scriptsize (operation)}}\\


 is the orthogonal complement of the subspace \mbox{\texttt{\mdseries\slshape U}} of the Lie algebra \mbox{\texttt{\mdseries\slshape L}} with respect to the Killing form $\kappa$, that is, the set $U^{{\perp}} = \{ x \in L; \kappa( x, y ) = 0 \hbox{ for all } y \in L \}$. 

 $U^{{\perp}}$ is a subspace of \mbox{\texttt{\mdseries\slshape L}}, and if \mbox{\texttt{\mdseries\slshape U}} is an ideal of \mbox{\texttt{\mdseries\slshape L}} then $U^{{\perp}}$ is a subalgebra of \mbox{\texttt{\mdseries\slshape L}}. 
\begin{Verbatim}[commandchars=!@|,fontsize=\small,frame=single,label=Example]
  !gapprompt@gap>| !gapinput@L:= SimpleLieAlgebra( "A", 1, Rationals );;|
  !gapprompt@gap>| !gapinput@b:= BasisVectors( Basis( L ) );;|
  !gapprompt@gap>| !gapinput@V:= VectorSpace( Rationals, [b[1],b[2]] );;|
  !gapprompt@gap>| !gapinput@KappaPerp( L, V );|
  <vector space of dimension 1 over Rationals>
\end{Verbatim}
 }

 

\subsection{\textcolor{Chapter }{IsNilpotentElement}}
\logpage{[ 64, 9, 5 ]}\nobreak
\hyperdef{L}{X7A00601387A060CF}{}
{\noindent\textcolor{FuncColor}{$\triangleright$\enspace\texttt{IsNilpotentElement({\mdseries\slshape L, x})\index{IsNilpotentElement@\texttt{IsNilpotentElement}}
\label{IsNilpotentElement}
}\hfill{\scriptsize (operation)}}\\


 \mbox{\texttt{\mdseries\slshape x}} is nilpotent in \mbox{\texttt{\mdseries\slshape L}} if its adjoint matrix is a nilpotent matrix. 
\begin{Verbatim}[commandchars=!@|,fontsize=\small,frame=single,label=Example]
  !gapprompt@gap>| !gapinput@L:= SimpleLieAlgebra( "A", 1, Rationals );;|
  !gapprompt@gap>| !gapinput@IsNilpotentElement( L, Basis( L )[1] );|
  true
\end{Verbatim}
 }

 

\subsection{\textcolor{Chapter }{NonNilpotentElement}}
\logpage{[ 64, 9, 6 ]}\nobreak
\hyperdef{L}{X86EF3E6F7BC0A8AD}{}
{\noindent\textcolor{FuncColor}{$\triangleright$\enspace\texttt{NonNilpotentElement({\mdseries\slshape L})\index{NonNilpotentElement@\texttt{NonNilpotentElement}}
\label{NonNilpotentElement}
}\hfill{\scriptsize (attribute)}}\\


 A non\texttt{\symbol{45}}nilpotent element of a Lie algebra \mbox{\texttt{\mdseries\slshape L}} is an element $x$ such that ad$x$ is not nilpotent. If \mbox{\texttt{\mdseries\slshape L}} is not nilpotent, then by Engel's theorem non\texttt{\symbol{45}}nilpotent
elements exist in \mbox{\texttt{\mdseries\slshape L}}. In this case this function returns a non\texttt{\symbol{45}}nilpotent
element of \mbox{\texttt{\mdseries\slshape L}}, otherwise (if \mbox{\texttt{\mdseries\slshape L}} is nilpotent) \texttt{fail} is returned. 
\begin{Verbatim}[commandchars=!@|,fontsize=\small,frame=single,label=Example]
  !gapprompt@gap>| !gapinput@L:= SimpleLieAlgebra( "G", 2, Rationals );;|
  !gapprompt@gap>| !gapinput@NonNilpotentElement( L );|
  v.13
  !gapprompt@gap>| !gapinput@IsNilpotentElement( L, last );|
  false
\end{Verbatim}
 }

 

\subsection{\textcolor{Chapter }{FindSl2}}
\logpage{[ 64, 9, 7 ]}\nobreak
\hyperdef{L}{X7A912D9E7B3BA874}{}
{\noindent\textcolor{FuncColor}{$\triangleright$\enspace\texttt{FindSl2({\mdseries\slshape L, x})\index{FindSl2@\texttt{FindSl2}}
\label{FindSl2}
}\hfill{\scriptsize (function)}}\\


 This function tries to find a subalgebra $S$ of the Lie algebra \mbox{\texttt{\mdseries\slshape L}} with $S$ isomorphic to $sl_2$ and such that the nilpotent element \mbox{\texttt{\mdseries\slshape x}} of \mbox{\texttt{\mdseries\slshape L}} is contained in $S$. If such an algebra exists then it is returned, otherwise \texttt{fail} is returned. 
\begin{Verbatim}[commandchars=!@|,fontsize=\small,frame=single,label=Example]
  !gapprompt@gap>| !gapinput@L:= SimpleLieAlgebra( "G", 2, Rationals );;|
  !gapprompt@gap>| !gapinput@b:= BasisVectors( Basis( L ) );;|
  !gapprompt@gap>| !gapinput@IsNilpotentElement( L, b[1] );|
  true
  !gapprompt@gap>| !gapinput@FindSl2( L, b[1] );|
  <Lie algebra of dimension 3 over Rationals>
\end{Verbatim}
 }

 }

  
\section{\textcolor{Chapter }{Universal Enveloping Algebras}}\label{Universal Enveloping Algebras}
\logpage{[ 64, 10, 0 ]}
\hyperdef{L}{X7875070C85DD4E8E}{}
{
  

\subsection{\textcolor{Chapter }{UniversalEnvelopingAlgebra}}
\logpage{[ 64, 10, 1 ]}\nobreak
\hyperdef{L}{X8226CD1680207A5F}{}
{\noindent\textcolor{FuncColor}{$\triangleright$\enspace\texttt{UniversalEnvelopingAlgebra({\mdseries\slshape L[, B]})\index{UniversalEnvelopingAlgebra@\texttt{UniversalEnvelopingAlgebra}}
\label{UniversalEnvelopingAlgebra}
}\hfill{\scriptsize (attribute)}}\\


 Returns the universal enveloping algebra of the Lie algebra \mbox{\texttt{\mdseries\slshape L}}. The elements of this algebra are written on a
Poincare\texttt{\symbol{45}}Birkhoff\texttt{\symbol{45}}Witt basis. 

 If a second argument \mbox{\texttt{\mdseries\slshape B}} is given, it must be a basis of \mbox{\texttt{\mdseries\slshape L}}, and an isomorphic copy of the universal enveloping algebra is returned,
generated by the images (in the universal enveloping algebra) of the elements
of \mbox{\texttt{\mdseries\slshape B}}. 
\begin{Verbatim}[commandchars=!@|,fontsize=\small,frame=single,label=Example]
  !gapprompt@gap>| !gapinput@L:= SimpleLieAlgebra( "A", 1, Rationals );;|
  !gapprompt@gap>| !gapinput@UL:= UniversalEnvelopingAlgebra( L );|
  <algebra-with-one of dimension infinity over Rationals>
  !gapprompt@gap>| !gapinput@g:= GeneratorsOfAlgebraWithOne( UL );|
  [ [(1)*x.1], [(1)*x.2], [(1)*x.3] ]
  !gapprompt@gap>| !gapinput@g[3]^2*g[2]^2*g[1]^2;|
  [(-4)*x.1*x.2*x.3^3+(1)*x.1^2*x.2^2*x.3^2+(2)*x.3^3+(2)*x.3^4]
\end{Verbatim}
 }

 }

  
\section{\textcolor{Chapter }{Finitely Presented Lie Algebras}}\label{Finitely Presented Lie Algebras}
\logpage{[ 64, 11, 0 ]}
\hyperdef{L}{X7B8C71E07F50B286}{}
{
  Finitely presented Lie algebras can be constructed from free Lie algebras by
using the \texttt{/} constructor, i.e., \texttt{FL/[r1, ..., rk]} is the quotient of the free Lie algebra \texttt{FL} by the ideal generated by the elements \texttt{r1, ..., rk} of \texttt{FL}. If the finitely presented Lie algebra \texttt{K} happens to be finite dimensional then an isomorphic structure constants Lie
algebra can be constructed by \texttt{NiceAlgebraMonomorphism(K)} (see{\nobreakspace}\texttt{NiceAlgebraMonomorphism} (\ref{NiceAlgebraMonomorphism})), which returns a surjective homomorphism. The structure constants Lie
algebra can then be accessed by calling \texttt{Range} (\ref{Range:of a general mapping}) for this map. Also limited computations with elements of the finitely
presented Lie algebra are possible. 

 
\begin{Verbatim}[commandchars=!@|,fontsize=\small,frame=single,label=Example]
  !gapprompt@gap>| !gapinput@L:= FreeLieAlgebra( Rationals, "s", "t" );|
  <Lie algebra over Rationals, with 2 generators>
  !gapprompt@gap>| !gapinput@gL:= GeneratorsOfAlgebra( L );; s:= gL[1];; t:= gL[2];;|
  !gapprompt@gap>| !gapinput@K:= L/[ s*(s*t), t*(t*(s*t)), s*(t*(s*t))-t*(s*t) ];|
  <Lie algebra over Rationals, with 2 generators>
  !gapprompt@gap>| !gapinput@h:= NiceAlgebraMonomorphism( K );|
  [ [(1)*s], [(1)*t] ] -> [ v.1, v.2 ]
  !gapprompt@gap>| !gapinput@U:= Range( h );|
  <Lie algebra of dimension 3 over Rationals>
  !gapprompt@gap>| !gapinput@IsLieNilpotent( U );|
  true
  !gapprompt@gap>| !gapinput@gK:= GeneratorsOfAlgebra( K );|
  [ [(1)*s], [(1)*t] ]
  !gapprompt@gap>| !gapinput@gK[1]*(gK[2]*gK[1]) = Zero( K );|
  true
\end{Verbatim}
 

 

\subsection{\textcolor{Chapter }{FpLieAlgebraByCartanMatrix}}
\logpage{[ 64, 11, 1 ]}\nobreak
\hyperdef{L}{X780A5B457A051110}{}
{\noindent\textcolor{FuncColor}{$\triangleright$\enspace\texttt{FpLieAlgebraByCartanMatrix({\mdseries\slshape C})\index{FpLieAlgebraByCartanMatrix@\texttt{FpLieAlgebraByCartanMatrix}}
\label{FpLieAlgebraByCartanMatrix}
}\hfill{\scriptsize (function)}}\\


 Here \mbox{\texttt{\mdseries\slshape C}} must be a Cartan matrix. The function returns the
finitely\texttt{\symbol{45}}presented Lie algebra over the field of rational
numbers defined by this Cartan matrix. By Serre's theorem, this Lie algebra is
a semisimple Lie algebra, and its root system has Cartan matrix \mbox{\texttt{\mdseries\slshape C}}. 
\begin{Verbatim}[commandchars=!@|,fontsize=\small,frame=single,label=Example]
  !gapprompt@gap>| !gapinput@C:= [ [ 2, -1 ], [ -3, 2 ] ];;|
  !gapprompt@gap>| !gapinput@K:= FpLieAlgebraByCartanMatrix( C );|
  <Lie algebra over Rationals, with 6 generators>
  !gapprompt@gap>| !gapinput@h:= NiceAlgebraMonomorphism( K );|
  [ [(1)*x1], [(1)*x2], [(1)*x3], [(1)*x4], [(1)*x5], [(1)*x6] ] ->
  [ v.1, v.2, v.3, v.4, v.5, v.6 ]
  !gapprompt@gap>| !gapinput@SemiSimpleType( Range( h ) );|
  "G2"
\end{Verbatim}
 }

 

\subsection{\textcolor{Chapter }{NilpotentQuotientOfFpLieAlgebra}}
\logpage{[ 64, 11, 2 ]}\nobreak
\hyperdef{L}{X79FD70C487EA9438}{}
{\noindent\textcolor{FuncColor}{$\triangleright$\enspace\texttt{NilpotentQuotientOfFpLieAlgebra({\mdseries\slshape FpL, max[, weights]})\index{NilpotentQuotientOfFpLieAlgebra@\texttt{NilpotentQuotientOfFpLieAlgebra}}
\label{NilpotentQuotientOfFpLieAlgebra}
}\hfill{\scriptsize (function)}}\\


 Here \mbox{\texttt{\mdseries\slshape FpL}} is a finitely presented Lie algebra. Let $K$ be the quotient of \mbox{\texttt{\mdseries\slshape FpL}} by the \mbox{\texttt{\mdseries\slshape max}}+1\texttt{\symbol{45}}th term of its lower central series. This function
calculates a surjective homomorphism from \mbox{\texttt{\mdseries\slshape FpL}} onto $K$. When called with the third argument \mbox{\texttt{\mdseries\slshape weights}}, the $k$\texttt{\symbol{45}}th generator of \mbox{\texttt{\mdseries\slshape FpL}} gets assigned the $k$\texttt{\symbol{45}}th element of the list \mbox{\texttt{\mdseries\slshape weights}}. In that case a quotient is calculated of \mbox{\texttt{\mdseries\slshape FpL}} by the ideal generated by all elements of weight \mbox{\texttt{\mdseries\slshape max}}+1. If the list \mbox{\texttt{\mdseries\slshape weights}} only consists of $1$'s then the two calls are equivalent. The default value of \mbox{\texttt{\mdseries\slshape weights}} is a list (of length equal to the number of generators of \mbox{\texttt{\mdseries\slshape FpL}}) consisting of $1$'s. 

 If the relators of \mbox{\texttt{\mdseries\slshape FpL}} are homogeneous, then the resulting algebra is naturally graded. 
\begin{Verbatim}[commandchars=!@|,fontsize=\small,frame=single,label=Example]
  !gapprompt@gap>| !gapinput@L:= FreeLieAlgebra( Rationals, "x", "y" );;|
  !gapprompt@gap>| !gapinput@g:= GeneratorsOfAlgebra(L);; x:= g[1]; y:= g[2];|
  (1)*x
  (1)*y
  !gapprompt@gap>| !gapinput@rr:=[ ((y*x)*x)*x-6*(y*x)*y,|
  !gapprompt@>| !gapinput@         3*((((y*x)*x)*x)*x)*x-20*(((y*x)*x)*x)*y ];|
  [ (-1)*(x*(x*(x*y)))+(6)*((x*y)*y),
    (-3)*(x*(x*(x*(x*(x*y)))))+(20)*(x*(x*((x*y)*y)))+(
      -20)*((x*(x*y))*(x*y)) ]
  !gapprompt@gap>| !gapinput@K:= L/rr;|
  <Lie algebra over Rationals, with 2 generators>
  !gapprompt@gap>| !gapinput@h:=NilpotentQuotientOfFpLieAlgebra(K, 50, [1,2] );|
  [ [(1)*x], [(1)*y] ] -> [ v.1, v.2 ]
  !gapprompt@gap>| !gapinput@L:= Range( h );|
  <Lie algebra of dimension 50 over Rationals>
  !gapprompt@gap>| !gapinput@Grading( L );|
  rec( hom_components := function( d ) ... end, max_degree := 50,
    min_degree := 1, source := Integers )
\end{Verbatim}
 }

 }

  
\section{\textcolor{Chapter }{Modules over Lie Algebras and Their Cohomology}}\label{Modules over Lie Algebras and Their Cohomology}
\logpage{[ 64, 12, 0 ]}
\hyperdef{L}{X7FBCB43C86BDD9C2}{}
{
  Representations of Lie algebras are dealt with in the same way as
representations of ordinary algebras (see \ref{Representations of Algebras}). In this section we mainly deal with modules over general Lie algebras and
their cohomology. The next section is devoted to modules over semisimple Lie
algebras. An $s$\texttt{\symbol{45}}cochain of a module $V$ over a Lie algebra $L$ is an $s$\texttt{\symbol{45}}linear map 
\[ c: L \times \cdots \times L \rightarrow V , \]
 with $s$ factors $L$, that is skew\texttt{\symbol{45}}symmetric (meaning that if any of the
arguments are interchanged, $c$ changes to $-c$). 

 Let $(x_1, \ldots, x_n)$ be a basis of $L$. Then any $s$\texttt{\symbol{45}}cochain is determined by the values $c( x_{{i_1}}, \ldots, x_{{i_s}} )$, where $1 \leq i_1 < i_2 < \cdots < i_s \leq \dim L$. Now this value again is a linear combination of basis elements of $V$: $c( x_{{i_1}}, \ldots, x_{{i_s}} ) = \sum \lambda^k_{{i_1,\ldots, i_s}} v_k$. Denote the dimension of $V$ by $r$. Then we represent an $s$\texttt{\symbol{45}}cocycle by a list of $r$ lists. The $j$\texttt{\symbol{45}}th of those lists consists of entries of the form 
\[ [ [ i_1, i_2, \ldots, i_s ], \lambda^j_{{i_1, \ldots, i_s}} ] \]
 where the coefficient on the second position is non\texttt{\symbol{45}}zero.
(We only store those entries for which this coefficient is
non\texttt{\symbol{45}}zero.) It follows that every $s$\texttt{\symbol{45}}tuple $(i_1, \ldots, i_s)$ gives rise to $r$ basis elements. 

 So the zero cochain is represented by a list of the form \texttt{[ [ ], [ ], \texttt{\symbol{92}}ldots, [ ] ]}. Furthermore, if $V$ is, e.g., $4$\texttt{\symbol{45}}dimensional, then the $2$\texttt{\symbol{45}}cochain represented by 

 
\begin{Verbatim}[commandchars=!@|,fontsize=\small,frame=single,label=Example]
  [ [ [ [1,2], 2] ], [ ], [ [ [1,2], 1/2 ] ], [ ] ]
\end{Verbatim}
 

 maps the pair $(x_1, x_2)$ to $2v_1 + 1/2 v_3$ (where $v_1$ is the first basis element of $V$, and $v_3$ the third), and all other pairs to zero. 

 By definition, $0$\texttt{\symbol{45}}cochains are constant maps $c( x ) = v_c \in V$ for all $x \in L$. So $0$\texttt{\symbol{45}}cochains have a different representation: they are just
represented by the list \texttt{[ v{\textunderscore}c ]}. 

 Cochains are constructed using the function \texttt{Cochain} (\ref{Cochain}), if \mbox{\texttt{\mdseries\slshape c}} is a cochain, then its corresponding list is returned by \texttt{ExtRepOfObj( \mbox{\texttt{\mdseries\slshape c}} )}. 

\subsection{\textcolor{Chapter }{IsCochain}}
\logpage{[ 64, 12, 1 ]}\nobreak
\hyperdef{L}{X82CC31CF79F59FEE}{}
{\noindent\textcolor{FuncColor}{$\triangleright$\enspace\texttt{IsCochain({\mdseries\slshape obj})\index{IsCochain@\texttt{IsCochain}}
\label{IsCochain}
}\hfill{\scriptsize (Category)}}\\
\noindent\textcolor{FuncColor}{$\triangleright$\enspace\texttt{IsCochainCollection({\mdseries\slshape obj})\index{IsCochainCollection@\texttt{IsCochainCollection}}
\label{IsCochainCollection}
}\hfill{\scriptsize (Category)}}\\


 Categories of cochains and of collections of cochains. }

 

\subsection{\textcolor{Chapter }{Cochain}}
\logpage{[ 64, 12, 2 ]}\nobreak
\hyperdef{L}{X79F3DF0D8791C2E3}{}
{\noindent\textcolor{FuncColor}{$\triangleright$\enspace\texttt{Cochain({\mdseries\slshape V, s, obj})\index{Cochain@\texttt{Cochain}}
\label{Cochain}
}\hfill{\scriptsize (operation)}}\\


 Constructs a \mbox{\texttt{\mdseries\slshape s}}\texttt{\symbol{45}}cochain given by the data in \mbox{\texttt{\mdseries\slshape obj}}, with respect to the Lie algebra module \mbox{\texttt{\mdseries\slshape V}}. If \mbox{\texttt{\mdseries\slshape s}} is non\texttt{\symbol{45}}zero, then \mbox{\texttt{\mdseries\slshape obj}} must be a list. 
\begin{Verbatim}[commandchars=!@|,fontsize=\small,frame=single,label=Example]
  !gapprompt@gap>| !gapinput@L:= SimpleLieAlgebra( "A", 1, Rationals );;|
  !gapprompt@gap>| !gapinput@V:= AdjointModule( L );|
  <3-dimensional left-module over <Lie algebra of dimension
  3 over Rationals>>
  !gapprompt@gap>| !gapinput@c1:= Cochain( V, 2,|
  !gapprompt@>| !gapinput@              [ [ [ [ 1, 3 ], -1 ] ], [ ], [ [ [ 2, 3 ], 1/2 ] ] ]);|
  <2-cochain>
  !gapprompt@gap>| !gapinput@ExtRepOfObj( c1 );|
  [ [ [ [ 1, 3 ], -1 ] ], [  ], [ [ [ 2, 3 ], 1/2 ] ] ]
  !gapprompt@gap>| !gapinput@c2:= Cochain( V, 0, Basis( V )[1] );|
  <0-cochain>
  !gapprompt@gap>| !gapinput@ExtRepOfObj( c2 );|
  v.1
  !gapprompt@gap>| !gapinput@IsCochain( c2 );|
  true
\end{Verbatim}
 }

 

\subsection{\textcolor{Chapter }{CochainSpace}}
\logpage{[ 64, 12, 3 ]}\nobreak
\hyperdef{L}{X7CF2919081600A3D}{}
{\noindent\textcolor{FuncColor}{$\triangleright$\enspace\texttt{CochainSpace({\mdseries\slshape V, s})\index{CochainSpace@\texttt{CochainSpace}}
\label{CochainSpace}
}\hfill{\scriptsize (operation)}}\\


 Returns the space of all \mbox{\texttt{\mdseries\slshape s}}\texttt{\symbol{45}}cochains with respect to \mbox{\texttt{\mdseries\slshape V}}. 
\begin{Verbatim}[commandchars=!@|,fontsize=\small,frame=single,label=Example]
  !gapprompt@gap>| !gapinput@L:= SimpleLieAlgebra( "A", 1, Rationals );;|
  !gapprompt@gap>| !gapinput@V:= AdjointModule( L );;|
  !gapprompt@gap>| !gapinput@C:=CochainSpace( V, 2 );|
  <vector space of dimension 9 over Rationals>
  !gapprompt@gap>| !gapinput@BasisVectors( Basis( C ) );|
  [ <2-cochain>, <2-cochain>, <2-cochain>, <2-cochain>, <2-cochain>,
    <2-cochain>, <2-cochain>, <2-cochain>, <2-cochain> ]
  !gapprompt@gap>| !gapinput@ExtRepOfObj( last[1] );|
  [ [ [ [ 1, 2 ], 1 ] ], [  ], [  ] ]
\end{Verbatim}
 }

 

\subsection{\textcolor{Chapter }{ValueCochain}}
\logpage{[ 64, 12, 4 ]}\nobreak
\hyperdef{L}{X7D6760DA84683011}{}
{\noindent\textcolor{FuncColor}{$\triangleright$\enspace\texttt{ValueCochain({\mdseries\slshape c, y1, y2, ..., ys})\index{ValueCochain@\texttt{ValueCochain}}
\label{ValueCochain}
}\hfill{\scriptsize (function)}}\\


 Here \mbox{\texttt{\mdseries\slshape c}} is an \texttt{s}\texttt{\symbol{45}}cochain. This function returns the value of \mbox{\texttt{\mdseries\slshape c}} when applied to the \texttt{s} elements \mbox{\texttt{\mdseries\slshape y1}} to \mbox{\texttt{\mdseries\slshape ys}} (that lie in the Lie algebra acting on the module corresponding to \mbox{\texttt{\mdseries\slshape c}}). It is also possible to call this function with two arguments: first \mbox{\texttt{\mdseries\slshape c}} and then the list containing \texttt{\mbox{\texttt{\mdseries\slshape y1}},...,\mbox{\texttt{\mdseries\slshape ys}}}. 
\begin{Verbatim}[commandchars=!@|,fontsize=\small,frame=single,label=Example]
  !gapprompt@gap>| !gapinput@L:= SimpleLieAlgebra( "A", 1, Rationals );;|
  !gapprompt@gap>| !gapinput@V:= AdjointModule( L );;|
  !gapprompt@gap>| !gapinput@C:= CochainSpace( V, 2 );;|
  !gapprompt@gap>| !gapinput@c:= Basis( C )[1];|
  <2-cochain>
  !gapprompt@gap>| !gapinput@ ValueCochain( c, Basis(L)[2], Basis(L)[1] );|
  (-1)*v.1
\end{Verbatim}
 }

 

\subsection{\textcolor{Chapter }{LieCoboundaryOperator}}
\logpage{[ 64, 12, 5 ]}\nobreak
\hyperdef{L}{X851F5EF47FA90CBC}{}
{\noindent\textcolor{FuncColor}{$\triangleright$\enspace\texttt{LieCoboundaryOperator({\mdseries\slshape c})\index{LieCoboundaryOperator@\texttt{LieCoboundaryOperator}}
\label{LieCoboundaryOperator}
}\hfill{\scriptsize (function)}}\\


 This is a function that takes an \texttt{s}\texttt{\symbol{45}}cochain \mbox{\texttt{\mdseries\slshape c}}, and returns an \texttt{s+1}\texttt{\symbol{45}}cochain. The coboundary operator is applied. 
\begin{Verbatim}[commandchars=!@|,fontsize=\small,frame=single,label=Example]
  !gapprompt@gap>| !gapinput@L:= SimpleLieAlgebra( "A", 1, Rationals );;|
  !gapprompt@gap>| !gapinput@V:= AdjointModule( L );;|
  !gapprompt@gap>| !gapinput@C:= CochainSpace( V, 2 );;|
  !gapprompt@gap>| !gapinput@c:= Basis( C )[1];;|
  !gapprompt@gap>| !gapinput@c1:= LieCoboundaryOperator( c );|
  <3-cochain>
  !gapprompt@gap>| !gapinput@c2:= LieCoboundaryOperator( c1 );|
  <4-cochain>
\end{Verbatim}
 }

 

\subsection{\textcolor{Chapter }{Cocycles (for Lie algebra module)}}
\logpage{[ 64, 12, 6 ]}\nobreak
\hyperdef{L}{X7FB815F38143939E}{}
{\noindent\textcolor{FuncColor}{$\triangleright$\enspace\texttt{Cocycles({\mdseries\slshape V, s})\index{Cocycles@\texttt{Cocycles}!for Lie algebra module}
\label{Cocycles:for Lie algebra module}
}\hfill{\scriptsize (operation)}}\\


 is the space of all \mbox{\texttt{\mdseries\slshape s}}\texttt{\symbol{45}}cocycles with respect to the Lie algebra module \mbox{\texttt{\mdseries\slshape V}}. That is the kernel of the coboundary operator when restricted to the space
of \mbox{\texttt{\mdseries\slshape s}}\texttt{\symbol{45}}cochains. }

 

\subsection{\textcolor{Chapter }{Coboundaries}}
\logpage{[ 64, 12, 7 ]}\nobreak
\hyperdef{L}{X7C4F372C7AE2F739}{}
{\noindent\textcolor{FuncColor}{$\triangleright$\enspace\texttt{Coboundaries({\mdseries\slshape V, s})\index{Coboundaries@\texttt{Coboundaries}}
\label{Coboundaries}
}\hfill{\scriptsize (operation)}}\\


 is the space of all \mbox{\texttt{\mdseries\slshape s}}\texttt{\symbol{45}}coboundaries with respect to the Lie algebra module \mbox{\texttt{\mdseries\slshape V}}. That is the image of the coboundary operator, when applied to the space of \mbox{\texttt{\mdseries\slshape s}}\texttt{\symbol{45}}1\texttt{\symbol{45}}cochains. By definition the space of
all 0\texttt{\symbol{45}}coboundaries is zero. 
\begin{Verbatim}[commandchars=!@|,fontsize=\small,frame=single,label=Example]
  !gapprompt@gap>| !gapinput@T:= EmptySCTable( 3, 0, "antisymmetric" );;|
  !gapprompt@gap>| !gapinput@SetEntrySCTable( T, 1, 2, [ 1, 3 ] );|
  !gapprompt@gap>| !gapinput@L:= LieAlgebraByStructureConstants( Rationals, T );;|
  !gapprompt@gap>| !gapinput@V:= FaithfulModule( L );|
  <left-module over <Lie algebra of dimension 3 over Rationals>>
  !gapprompt@gap>| !gapinput@Cocycles( V, 2 );|
  <vector space of dimension 7 over Rationals>
  !gapprompt@gap>| !gapinput@Coboundaries( V, 2 );|
  <vector space over Rationals, with 9 generators>
  !gapprompt@gap>| !gapinput@Dimension( last );|
  5
\end{Verbatim}
 }

 }

  
\section{\textcolor{Chapter }{Modules over Semisimple Lie Algebras}}\label{Modules over Semisimple Lie Algebras}
\logpage{[ 64, 13, 0 ]}
\hyperdef{L}{X78A201238137E822}{}
{
  This section contains functions for calculating information on representations
of semisimple Lie algebras. First we have some functions for calculating some
combinatorial data (set of dominant weights, the dominant character, the
decomposition of a tensor product, the dimension of a
highest\texttt{\symbol{45}}weight module). Then there is a function for
creating an admissible lattice in the universal enveloping algebra of a
semisimple Lie algebra. Finally we have a function for constructing a
highest\texttt{\symbol{45}}weight module over a semisimple Lie algebra. 

 

\subsection{\textcolor{Chapter }{DominantWeights}}
\logpage{[ 64, 13, 1 ]}\nobreak
\hyperdef{L}{X7D8522E37ED1024A}{}
{\noindent\textcolor{FuncColor}{$\triangleright$\enspace\texttt{DominantWeights({\mdseries\slshape R, maxw})\index{DominantWeights@\texttt{DominantWeights}}
\label{DominantWeights}
}\hfill{\scriptsize (operation)}}\\


 Returns a list consisting of two lists. The first of these contains the
dominant weights (written on the basis of fundamental weights) of the
irreducible highest\texttt{\symbol{45}}weight module, with highest weight \mbox{\texttt{\mdseries\slshape maxw}}, over the Lie algebra with the root system \mbox{\texttt{\mdseries\slshape R}}. The $i$\texttt{\symbol{45}}th element of the second list is the level of the $i$\texttt{\symbol{45}}th dominant weight. (Where the level is defined as
follows. For a weight $\mu$ we write $\mu = \lambda - \sum_i k_i \alpha_i$, where the $\alpha_i$ are the simple roots, and $\lambda$ the highest weight. Then the level of $\mu$ is $\sum_i k_i$.) }

 

\subsection{\textcolor{Chapter }{DominantCharacter (for a semisimple Lie algebra and a highest weight)}}
\logpage{[ 64, 13, 2 ]}\nobreak
\hyperdef{L}{X79AAC71E8267E9F8}{}
{\noindent\textcolor{FuncColor}{$\triangleright$\enspace\texttt{DominantCharacter({\mdseries\slshape L, maxw})\index{DominantCharacter@\texttt{DominantCharacter}!for a semisimple Lie algebra and a highest weight}
\label{DominantCharacter:for a semisimple Lie algebra and a highest weight}
}\hfill{\scriptsize (operation)}}\\
\noindent\textcolor{FuncColor}{$\triangleright$\enspace\texttt{DominantCharacter({\mdseries\slshape R, maxw})\index{DominantCharacter@\texttt{DominantCharacter}!for a root system and a highest weight}
\label{DominantCharacter:for a root system and a highest weight}
}\hfill{\scriptsize (operation)}}\\


 For a highest weight \mbox{\texttt{\mdseries\slshape maxw}} and a semisimple Lie algebra \mbox{\texttt{\mdseries\slshape L}}, this returns the dominant weights of the highest\texttt{\symbol{45}}weight
module over \mbox{\texttt{\mdseries\slshape L}}, with highest weight \mbox{\texttt{\mdseries\slshape maxw}}. The output is a list of two lists, the first list contains the dominant
weights; the second list contains their multiplicities. 

 The first argument can also be a root system, in which case the dominant
character of the highest\texttt{\symbol{45}}weight module over the
corresponding semisimple Lie algebra is returned. }

 

\subsection{\textcolor{Chapter }{DecomposeTensorProduct}}
\logpage{[ 64, 13, 3 ]}\nobreak
\hyperdef{L}{X7BE7129384B012DF}{}
{\noindent\textcolor{FuncColor}{$\triangleright$\enspace\texttt{DecomposeTensorProduct({\mdseries\slshape L, w1, w2})\index{DecomposeTensorProduct@\texttt{DecomposeTensorProduct}}
\label{DecomposeTensorProduct}
}\hfill{\scriptsize (operation)}}\\


 Here \mbox{\texttt{\mdseries\slshape L}} is a semisimple Lie algebra and \mbox{\texttt{\mdseries\slshape w1}}, \mbox{\texttt{\mdseries\slshape w2}} are dominant weights. Let $V_i$ be the irreducible highest\texttt{\symbol{45}}weight module over \mbox{\texttt{\mdseries\slshape L}} with highest weight $w_i$ for $i = 1, 2$. Let $W = V_1 \otimes V_2$. Then in general $W$ is a reducible \mbox{\texttt{\mdseries\slshape L}}\texttt{\symbol{45}}module. Now this function returns a list of two lists. The
first of these is the sorted list of highest weights of the irreducible
modules occurring in the decomposition of $W$ as a direct sum of irreducible modules. The second list contains the
multiplicities of these weights (i.e., the number of copies of the irreducible
module with the corresponding highest weight that occur in $W$). The algorithm uses Klimyk's formula (see{\nobreakspace}\cite{Klimyk68} or \cite{Klimyk66} for the original Russian version). }

 

\subsection{\textcolor{Chapter }{DimensionOfHighestWeightModule}}
\logpage{[ 64, 13, 4 ]}\nobreak
\hyperdef{L}{X7D67A9BC7E4714D9}{}
{\noindent\textcolor{FuncColor}{$\triangleright$\enspace\texttt{DimensionOfHighestWeightModule({\mdseries\slshape L, w})\index{DimensionOfHighestWeightModule@\texttt{DimensionOfHighestWeightModule}}
\label{DimensionOfHighestWeightModule}
}\hfill{\scriptsize (operation)}}\\


 Here \mbox{\texttt{\mdseries\slshape L}} is a semisimple Lie algebra, and \mbox{\texttt{\mdseries\slshape w}} a dominant weight. This function returns the dimension of the
highest\texttt{\symbol{45}}weight module over \mbox{\texttt{\mdseries\slshape L}} with highest weight \mbox{\texttt{\mdseries\slshape w}}. The algorithm uses Weyl's dimension formula. 
\begin{Verbatim}[commandchars=!@|,fontsize=\small,frame=single,label=Example]
  !gapprompt@gap>| !gapinput@L:= SimpleLieAlgebra( "F", 4, Rationals );;|
  !gapprompt@gap>| !gapinput@R:= RootSystem( L );;|
  !gapprompt@gap>| !gapinput@DominantWeights( R, [ 1, 1, 0, 0 ] );|
  [ [ [ 1, 1, 0, 0 ], [ 2, 0, 0, 0 ], [ 0, 0, 1, 0 ], [ 0, 1, 0, 0 ],
        [ 1, 0, 0, 0 ], [ 0, 0, 0, 0 ] ], [ 0, 3, 4, 8, 11, 19 ] ]
  !gapprompt@gap>| !gapinput@DominantCharacter( L, [ 1, 1, 0, 0 ] );|
  [ [ [ 1, 1, 0, 0 ], [ 2, 0, 0, 0 ], [ 0, 0, 1, 0 ], [ 0, 1, 0, 0 ],
        [ 1, 0, 0, 0 ], [ 0, 0, 0, 0 ] ], [ 1, 1, 4, 6, 14, 21 ] ]
  !gapprompt@gap>| !gapinput@DecomposeTensorProduct( L, [ 1, 0, 0, 0 ], [ 0, 0, 1, 0 ] );|
  [ [ [ 0, 0, 0, 1 ], [ 0, 0, 1, 0 ], [ 0, 1, 0, 0 ], [ 1, 0, 0, 0 ],
        [ 1, 0, 1, 0 ], [ 1, 1, 0, 0 ], [ 2, 0, 0, 0 ] ],
    [ 1, 1, 1, 1, 1, 1, 1 ] ]
  !gapprompt@gap>| !gapinput@DimensionOfHighestWeightModule( L, [ 1, 2, 3, 4 ] );|
  79316832731136
\end{Verbatim}
 }

 }

  
\section{\textcolor{Chapter }{Admissible Lattices in UEA}}\label{Admissible Lattices in UEA}
\logpage{[ 64, 14, 0 ]}
\hyperdef{L}{X840E5FAE7D2C2702}{}
{
  Let $L$ be a semisimple Lie algebra over a field of characteristic $0$, and let $R$ be its root system. For a positive root $\alpha$ we let $x_{\alpha}$ and $y_{\alpha}$ be positive and negative root vectors, respectively, both from a fixed
Chevalley basis of $L$. Furthermore, $h_1, \ldots, h_l$ are the Cartan elements from the same Chevalley basis. Also we set 
\[ x_{\alpha}^{(n)} = {{x_{\alpha}^n \over n!}}, y_{\alpha}^{(n)} =
{{y_{\alpha}^n \over n!}} . \]
 Furthermore, let $\alpha_1, \ldots, \alpha_s$ denote the positive roots of $R$. For multi\texttt{\symbol{45}}indices $N = (n_1, \ldots, n_s)$, $M = (m_1, \ldots, m_s)$ and $K = (k_1, \ldots, k_s)$ (where $n_i, m_i, k_i \geq 0$) set \begin{center}
\begin{tabular}{lcl}$x^N$&
=&
$x_{{\alpha_1}}^{(n_1)} \cdots x_{{\alpha_s}}^{(n_s)}$,\\
$y^M$&
=&
$y_{{\alpha_1}}^{(m_1)} \cdots y_{{\alpha_s}}^{(m_s)}$,\\
$h^K$&
=&
${{h_1 \choose k_1}} \cdots {{h_l \choose k_l}}$\\
\end{tabular}\\[2mm]
\end{center}

 Then by a theorem of Kostant, the $x_{\alpha}^{(n)}$ and $y_{\alpha}^{(n)}$ generate a subring of the universal enveloping algebra $U(L)$ spanned (as a free $Z$\texttt{\symbol{45}}module) by the elements 
\[ y^M h^K x^N \]
 (see, e.g., \cite{Hum72} or \cite[Section 26]{Hum78}) So by the Poincare\texttt{\symbol{45}}Birkhoff\texttt{\symbol{45}}Witt
theorem this subring is a lattice in $U(L)$. Furthermore, this lattice is invariant under the $x_{\alpha}^{(n)}$ and $y_{\alpha}^{(n)}$. Therefore, it is called an admissible lattice in $U(L)$. 

 The next functions enable us to construct the generators of such an admissible
lattice. 

\subsection{\textcolor{Chapter }{IsUEALatticeElement}}
\logpage{[ 64, 14, 1 ]}\nobreak
\hyperdef{L}{X86E6722379576746}{}
{\noindent\textcolor{FuncColor}{$\triangleright$\enspace\texttt{IsUEALatticeElement({\mdseries\slshape obj})\index{IsUEALatticeElement@\texttt{IsUEALatticeElement}}
\label{IsUEALatticeElement}
}\hfill{\scriptsize (Category)}}\\
\noindent\textcolor{FuncColor}{$\triangleright$\enspace\texttt{IsUEALatticeElementCollection({\mdseries\slshape obj})\index{IsUEALatticeElementCollection@\texttt{IsUEALatticeElementCollection}}
\label{IsUEALatticeElementCollection}
}\hfill{\scriptsize (Category)}}\\
\noindent\textcolor{FuncColor}{$\triangleright$\enspace\texttt{IsUEALatticeElementFamily({\mdseries\slshape fam})\index{IsUEALatticeElementFamily@\texttt{IsUEALatticeElementFamily}}
\label{IsUEALatticeElementFamily}
}\hfill{\scriptsize (Category)}}\\


 is the category of elements of an admissible lattice in the universal
enveloping algebra of a semisimple Lie algebra \texttt{L}. }

 

\subsection{\textcolor{Chapter }{LatticeGeneratorsInUEA}}
\logpage{[ 64, 14, 2 ]}\nobreak
\hyperdef{L}{X79F4F58B7888B0A5}{}
{\noindent\textcolor{FuncColor}{$\triangleright$\enspace\texttt{LatticeGeneratorsInUEA({\mdseries\slshape L})\index{LatticeGeneratorsInUEA@\texttt{LatticeGeneratorsInUEA}}
\label{LatticeGeneratorsInUEA}
}\hfill{\scriptsize (attribute)}}\\


 Here \mbox{\texttt{\mdseries\slshape L}} must be a semisimple Lie algebra of characteristic $0$. This function returns a list of generators of an admissible lattice in the
universal enveloping algebra of \mbox{\texttt{\mdseries\slshape L}}, relative to the Chevalley basis contained in \texttt{ChevalleyBasis( \mbox{\texttt{\mdseries\slshape L}} )} (see{\nobreakspace}\texttt{ChevalleyBasis} (\ref{ChevalleyBasis})). First are listed the negative root vectors (denoted by $y_1, \ldots, y_s$), then the positive root vectors (denoted by $x_1, \ldots, x_s$). At the end of the list there are the Cartan elements. They are printed as \texttt{( hi/1 )}, which means 
\[ {{h_i \choose 1}}. \]
 In general the printed form \texttt{( hi/ k )} means 
\[ {{h_i \choose k}}. \]
 

 Also $y_i^{(m)}$ is printed as \texttt{yi\texttt{\symbol{94}}(m)}, which means that entering \texttt{yi\texttt{\symbol{94}}m} at the \textsf{GAP} prompt results in the output \texttt{m!*yi\texttt{\symbol{94}}(m)}. 

 Products of lattice generators are collected using the following order: first
come the $y_i^{(m_i)}$ (in the same order as the positive roots), then the ${h_i \choose k_i}$, and then the $x_i^{(n_i)}$ (in the same order as the positive roots). }

 

\subsection{\textcolor{Chapter }{ObjByExtRep (for creating a UEALattice element)}}
\logpage{[ 64, 14, 3 ]}\nobreak
\hyperdef{L}{X875FD1627F3B72DB}{}
{\noindent\textcolor{FuncColor}{$\triangleright$\enspace\texttt{ObjByExtRep({\mdseries\slshape F, descr})\index{ObjByExtRep@\texttt{ObjByExtRep}!for creating a UEALattice element}
\label{ObjByExtRep:for creating a UEALattice element}
}\hfill{\scriptsize (method)}}\\


 An UEALattice element is represented by a list of the form \texttt{[ m1, c1, m2, c2, ... ]}, where the \texttt{c1}, \texttt{c2} etc. are coefficients, and the \texttt{m1}, \texttt{m2} etc. are monomials. A monomial is a list of the form \texttt{[ ind1, e1, ind2, e2, ... ]} where \texttt{ind1}, \texttt{ind2} are indices, and \texttt{e1}, \texttt{e2} etc. are exponents. Let $N$ be the number of positive roots of the underlying Lie algebra \texttt{L}. The indices lie between 1 and $dim(L)$. If an index lies between 1 and \texttt{N}, then it represents a negative root vector (corresponding to the root \texttt{NegativeRoots( R )[ind]}, where \texttt{R} is the root system of \texttt{L}; see{\nobreakspace}\texttt{NegativeRoots} (\ref{NegativeRoots})). This leads to a factor \texttt{yind1\texttt{\symbol{94}}(e1)} in the printed form of the monomial (which equals \texttt{z\texttt{\symbol{94}}e1/e1!}, where \texttt{z} is a basis element of \texttt{L}). If an index lies between $N+1$ and $2N$, then it represents a positive root vector. Finally, if ind lies between $2N+1$ and $2N+rank$, then it represents an element of the Cartan subalgebra. This is printed as $( h_1/ e_1 )$, meaning ${h_1 \choose e_1}$, where $h_1, \ldots, h_{rank}$ are the canonical Cartan generators. 

 The zero element is represented by the empty list, the identity element by the
list \texttt{[ [], 1 ]}. 

 
\begin{Verbatim}[commandchars=!@|,fontsize=\small,frame=single,label=Example]
  !gapprompt@gap>| !gapinput@L:= SimpleLieAlgebra( "G", 2, Rationals );;|
  !gapprompt@gap>| !gapinput@g:=LatticeGeneratorsInUEA( L );|
  [ y1, y2, y3, y4, y5, y6, x1, x2, x3, x4, x5, x6, ( h13/1 ),
    ( h14/1 ) ]
  !gapprompt@gap>| !gapinput@IsUEALatticeElement( g[1] );|
  true
  !gapprompt@gap>| !gapinput@g[1]^3;|
  6*y1^(3)
  !gapprompt@gap>| !gapinput@q:= g[7]*g[1]^2;|
  -2*y1+2*y1*( h13/1 )+2*y1^(2)*x1
  !gapprompt@gap>| !gapinput@ExtRepOfObj( q );|
  [ [ 1, 1 ], -2, [ 1, 1, 13, 1 ], 2, [ 1, 2, 7, 1 ], 2 ]
\end{Verbatim}
 }

 

\subsection{\textcolor{Chapter }{IsWeightRepElement}}
\logpage{[ 64, 14, 4 ]}\nobreak
\hyperdef{L}{X8248DB547B02B0FA}{}
{\noindent\textcolor{FuncColor}{$\triangleright$\enspace\texttt{IsWeightRepElement({\mdseries\slshape obj})\index{IsWeightRepElement@\texttt{IsWeightRepElement}}
\label{IsWeightRepElement}
}\hfill{\scriptsize (Category)}}\\
\noindent\textcolor{FuncColor}{$\triangleright$\enspace\texttt{IsWeightRepElementCollection({\mdseries\slshape obj})\index{IsWeightRepElementCollection@\texttt{IsWeightRepElementCollection}}
\label{IsWeightRepElementCollection}
}\hfill{\scriptsize (Category)}}\\
\noindent\textcolor{FuncColor}{$\triangleright$\enspace\texttt{IsWeightRepElementFamily({\mdseries\slshape fam})\index{IsWeightRepElementFamily@\texttt{IsWeightRepElementFamily}}
\label{IsWeightRepElementFamily}
}\hfill{\scriptsize (Category)}}\\


 Is a category of vectors, that is used to construct elements of
highest\texttt{\symbol{45}}weight modules (by \texttt{HighestWeightModule} (\ref{HighestWeightModule})). 

 \texttt{WeightRepElement}s are represented by a list of the form \texttt{[ v1, c1, v2, c2, ....]}, where the \texttt{vi} are basis vectors, and the \texttt{ci} are coefficients. Furthermore a basis vector \texttt{v} is a weight vector. It is represented by a list of the form \texttt{[ k, mon, wt ]}, where \texttt{k} is an integer (the basis vectors are numbered from $1$ to $\dim V$, where $V$ is the highest weight module), \texttt{mon} is an \texttt{UEALatticeElement} (which means that the result of applying \texttt{mon} to a highest weight vector is \texttt{v}; see{\nobreakspace}\texttt{IsUEALatticeElement} (\ref{IsUEALatticeElement})) and \texttt{wt} is the weight of \texttt{v}. A \texttt{WeightRepElement} is printed as \texttt{mon*v0}, where \texttt{v0} denotes a fixed highest weight vector. 

 If \texttt{v} is a \texttt{WeightRepElement}, then \texttt{ExtRepOfObj( v )} returns the corresponding list, and if \texttt{list} is such a list and \mbox{\texttt{\mdseries\slshape fam}} a \texttt{WeightRepElementFamily}, then \texttt{ObjByExtRep( \mbox{\texttt{\mdseries\slshape list}}, \mbox{\texttt{\mdseries\slshape fam}} )} returns the corresponding \texttt{WeightRepElement}. }

 

\subsection{\textcolor{Chapter }{HighestWeightModule}}
\logpage{[ 64, 14, 5 ]}\nobreak
\hyperdef{L}{X7FB14F7F80EFF33F}{}
{\noindent\textcolor{FuncColor}{$\triangleright$\enspace\texttt{HighestWeightModule({\mdseries\slshape L, wt})\index{HighestWeightModule@\texttt{HighestWeightModule}}
\label{HighestWeightModule}
}\hfill{\scriptsize (operation)}}\\


 returns the highest weight module with highest weight \mbox{\texttt{\mdseries\slshape wt}} of the semisimple Lie algebra \mbox{\texttt{\mdseries\slshape L}} of characteristic $0$. 

 Note that the elements of such a module lie in the category \texttt{IsLeftAlgebraModuleElement} (\ref{IsLeftAlgebraModuleElement}) (and in particular they do not lie in the category \texttt{IsWeightRepElement} (\ref{IsWeightRepElement})). However, if \texttt{v} is an element of such a module, then \texttt{ExtRepOfObj( v )} is a \texttt{WeightRepElement}. 

 Note that for the following examples of this chapter we increase the line
length limit from its default value 80 to 81 in order to make some long output
expressions fit into the lines. 

 
\begin{Verbatim}[commandchars=!@|,fontsize=\small,frame=single,label=Example]
  !gapprompt@gap>| !gapinput@K1:= SimpleLieAlgebra( "G", 2, Rationals );;|
  !gapprompt@gap>| !gapinput@K2:= SimpleLieAlgebra( "B", 2, Rationals );;|
  !gapprompt@gap>| !gapinput@L:= DirectSumOfAlgebras( K1, K2 );|
  <Lie algebra of dimension 24 over Rationals>
  !gapprompt@gap>| !gapinput@V:= HighestWeightModule( L, [ 0, 1, 1, 1 ] );|
  <224-dimensional left-module over <Lie algebra of dimension
  24 over Rationals>>
  !gapprompt@gap>| !gapinput@vv:= GeneratorsOfLeftModule( V );;|
  !gapprompt@gap>| !gapinput@vv[100];|
  y5*y7*y10*v0
  !gapprompt@gap>| !gapinput@e:= ExtRepOfObj( vv[100] );|
  y5*y7*y10*v0
  !gapprompt@gap>| !gapinput@ExtRepOfObj( e );|
  [ [ 100, y5*y7*y10, [ -3, 2, -1, 1 ] ], 1 ]
  !gapprompt@gap>| !gapinput@Basis(L)[17]^vv[100];|
  -1*y5*y7*y8*v0-1*y5*y9*v0
\end{Verbatim}
 }

 }

  
\section{\textcolor{Chapter }{Tensor Products and Exterior and Symmetric Powers of Algebra Modules}}\label{Tensor Products and Exterior and Symmetric Powers of Algebra Modules}
\logpage{[ 64, 15, 0 ]}
\hyperdef{L}{X8750BDBF7EA5E868}{}
{
  

\subsection{\textcolor{Chapter }{TensorProductOfAlgebraModules (for a list of algebra modules)}}
\logpage{[ 64, 15, 1 ]}\nobreak
\hyperdef{L}{X7A1E0AC4800E7FDA}{}
{\noindent\textcolor{FuncColor}{$\triangleright$\enspace\texttt{TensorProductOfAlgebraModules({\mdseries\slshape list})\index{TensorProductOfAlgebraModules@\texttt{TensorProductOfAlgebraModules}!for a list of algebra modules}
\label{TensorProductOfAlgebraModules:for a list of algebra modules}
}\hfill{\scriptsize (operation)}}\\
\noindent\textcolor{FuncColor}{$\triangleright$\enspace\texttt{TensorProductOfAlgebraModules({\mdseries\slshape V, W})\index{TensorProductOfAlgebraModules@\texttt{TensorProductOfAlgebraModules}!for two algebra modules}
\label{TensorProductOfAlgebraModules:for two algebra modules}
}\hfill{\scriptsize (operation)}}\\


 Here the elements of \mbox{\texttt{\mdseries\slshape list}} must be algebra modules. The tensor product is returned as an algebra module.
The two\texttt{\symbol{45}}argument version works in the same way and returns
the tensor product of its arguments. 
\begin{Verbatim}[commandchars=!@|,fontsize=\small,frame=single,label=Example]
  !gapprompt@gap>| !gapinput@L:= SimpleLieAlgebra("G",2,Rationals);;|
  !gapprompt@gap>| !gapinput@V:= HighestWeightModule( L, [ 1, 0 ] );;|
  !gapprompt@gap>| !gapinput@W:= TensorProductOfAlgebraModules( [ V, V, V ] );|
  <343-dimensional left-module over <Lie algebra of dimension
  14 over Rationals>>
  !gapprompt@gap>| !gapinput@w:= Basis(W)[1];|
  1*(1*v0<x>1*v0<x>1*v0)
  !gapprompt@gap>| !gapinput@Basis(L)[1]^w;|
  <0-tensor>
  !gapprompt@gap>| !gapinput@Basis(L)[7]^w;|
  1*(1*v0<x>1*v0<x>y1*v0)+1*(1*v0<x>y1*v0<x>1*v0)+1*(y
  1*v0<x>1*v0<x>1*v0)
\end{Verbatim}
 }

 

\subsection{\textcolor{Chapter }{ExteriorPowerOfAlgebraModule}}
\logpage{[ 64, 15, 2 ]}\nobreak
\hyperdef{L}{X7F4AB6A1863E8FB2}{}
{\noindent\textcolor{FuncColor}{$\triangleright$\enspace\texttt{ExteriorPowerOfAlgebraModule({\mdseries\slshape V, k})\index{ExteriorPowerOfAlgebraModule@\texttt{ExteriorPowerOfAlgebraModule}}
\label{ExteriorPowerOfAlgebraModule}
}\hfill{\scriptsize (operation)}}\\


 Here \mbox{\texttt{\mdseries\slshape V}} must be an algebra module, defined over a Lie algebra. This function returns
the \mbox{\texttt{\mdseries\slshape k}}\texttt{\symbol{45}}th exterior power of \mbox{\texttt{\mdseries\slshape V}} as an algebra module. 
\begin{Verbatim}[commandchars=!@|,fontsize=\small,frame=single,label=Example]
  !gapprompt@gap>| !gapinput@L:= SimpleLieAlgebra("G",2,Rationals);;|
  !gapprompt@gap>| !gapinput@V:= HighestWeightModule( L, [ 1, 0 ] );;|
  !gapprompt@gap>| !gapinput@W:= ExteriorPowerOfAlgebraModule( V, 3 );|
  <35-dimensional left-module over <Lie algebra of dimension
  14 over Rationals>>
  !gapprompt@gap>| !gapinput@w:= Basis(W)[1];|
  1*(1*v0/\y1*v0/\y3*v0)
  !gapprompt@gap>| !gapinput@Basis(L)[10]^w;|
  1*(1*v0/\y1*v0/\y6*v0)+1*(1*v0/\y3*v0/\y5*v0)+1*(y1*v0/\y3*v0/\y4*v0)
\end{Verbatim}
 }

 

\subsection{\textcolor{Chapter }{SymmetricPowerOfAlgebraModule}}
\logpage{[ 64, 15, 3 ]}\nobreak
\hyperdef{L}{X842DF85687D61A56}{}
{\noindent\textcolor{FuncColor}{$\triangleright$\enspace\texttt{SymmetricPowerOfAlgebraModule({\mdseries\slshape V, k})\index{SymmetricPowerOfAlgebraModule@\texttt{SymmetricPowerOfAlgebraModule}}
\label{SymmetricPowerOfAlgebraModule}
}\hfill{\scriptsize (operation)}}\\


 Here \mbox{\texttt{\mdseries\slshape V}} must be an algebra module. This function returns the \mbox{\texttt{\mdseries\slshape k}}\texttt{\symbol{45}}th symmetric power of \mbox{\texttt{\mdseries\slshape V}} (as an algebra module). 
\begin{Verbatim}[commandchars=!@|,fontsize=\small,frame=single,label=Example]
  !gapprompt@gap>| !gapinput@L:= SimpleLieAlgebra("G",2,Rationals);;|
  !gapprompt@gap>| !gapinput@V:= HighestWeightModule( L, [ 1, 0 ] );;|
  !gapprompt@gap>| !gapinput@W:= SymmetricPowerOfAlgebraModule( V, 3 );|
  <84-dimensional left-module over <Lie algebra of dimension
  14 over Rationals>>
  !gapprompt@gap>| !gapinput@w:= Basis(W)[1];|
  1*(1*v0.1*v0.1*v0)
  !gapprompt@gap>| !gapinput@Basis(L)[2]^w;|
  <0-symmetric element>
  !gapprompt@gap>| !gapinput@Basis(L)[7]^w;|
  3*(1*v0.1*v0.y1*v0)
\end{Verbatim}
 }

 }

 }

            
\chapter{\textcolor{Chapter }{Magma Rings}}\label{Magma Rings}
\logpage{[ 65, 0, 0 ]}
\hyperdef{L}{X825897DC7A16E07D}{}
{
  \index{group algebra} \index{group ring} Given a magma $M$ then the \emph{free magma ring} (or \emph{magma ring} for short) $RM$ of $M$ over a ring\texttt{\symbol{45}}with\texttt{\symbol{45}}one $R$ is the set of finite sums $\sum_{{i \in I}} r_i m_i$ with $r_i \in R$, and $m_i \in M$. With the obvious addition and $R$\texttt{\symbol{45}}action from the left, $RM$ is a free $R$\texttt{\symbol{45}}module with $R$\texttt{\symbol{45}}basis $M$, and with the usual convolution product, $RM$ is a ring. 

 Typical examples of free magma rings are 

 
\begin{itemize}
\item  (multivariate) polynomial rings (see{\nobreakspace}\ref{Polynomial Rings and Function Fields}), where the magma is a free abelian monoid generated by the indeterminates, 
\item  group rings (see{\nobreakspace}\texttt{IsGroupRing} (\ref{IsGroupRing})), where the magma is a group, 
\item  Laurent polynomial rings, which are group rings of the free abelian  groups generated by the indeterminates, 
\item  free algebras and free associative algebras, with or without one, where the
magma is a free magma or a free semigroup, or a free
magma\texttt{\symbol{45}}with\texttt{\symbol{45}}one or a free monoid,
respectively. 
\end{itemize}
 Note that formally, polynomial rings in \textsf{GAP} are not constructed as free magma rings. 

 Furthermore, a free Lie algebra is \emph{not} a magma ring, because of the additional relations given by the Jacobi
identity; see{\nobreakspace}\ref{Magma Rings modulo Relations} for a generalization of magma rings that covers such structures. 

 The coefficient ring $R$ and the magma $M$ cannot be regarded as subsets of $RM$, hence the natural \emph{embeddings} of $R$ and $M$ into $RM$ must be handled via explicit embedding maps (see{\nobreakspace}\ref{Natural Embeddings related to Magma Rings}). Note that in a magma ring, the addition of elements is in general different
from an addition that may be defined already for the elements of the magma;
for example, the addition in the group ring of a matrix group does in general \emph{not} coincide with the addition of matrices. 

 
\begin{Verbatim}[commandchars=!@|,fontsize=\small,frame=single,label=Example]
  !gapprompt@gap>| !gapinput@a:= Algebra( GF(2), [ [ [ Z(2) ] ] ] );;  Size( a );|
  2
  !gapprompt@gap>| !gapinput@rm:= FreeMagmaRing( GF(2), a );;|
  !gapprompt@gap>| !gapinput@emb:= Embedding( a, rm );;|
  !gapprompt@gap>| !gapinput@z:= Zero( a );;  o:= One( a );;|
  !gapprompt@gap>| !gapinput@imz:= z ^ emb;  IsZero( imz );|
  (Z(2)^0)*[ [ 0*Z(2) ] ]
  false
  !gapprompt@gap>| !gapinput@im1:= ( z + o ) ^ emb;|
  (Z(2)^0)*[ [ Z(2)^0 ] ]
  !gapprompt@gap>| !gapinput@im2:= z ^ emb + o ^ emb;|
  (Z(2)^0)*[ [ 0*Z(2) ] ]+(Z(2)^0)*[ [ Z(2)^0 ] ]
  !gapprompt@gap>| !gapinput@im1 = im2;|
  false
\end{Verbatim}
  
\section{\textcolor{Chapter }{Free Magma Rings}}\label{Free Magma Rings}
\logpage{[ 65, 1, 0 ]}
\hyperdef{L}{X8398F87F8231A163}{}
{
  

\subsection{\textcolor{Chapter }{FreeMagmaRing}}
\logpage{[ 65, 1, 1 ]}\nobreak
\hyperdef{L}{X7B9AF0A47F44E4B4}{}
{\noindent\textcolor{FuncColor}{$\triangleright$\enspace\texttt{FreeMagmaRing({\mdseries\slshape R, M})\index{FreeMagmaRing@\texttt{FreeMagmaRing}}
\label{FreeMagmaRing}
}\hfill{\scriptsize (function)}}\\


 is a free magma ring over the ring \mbox{\texttt{\mdseries\slshape R}}, free on the magma \mbox{\texttt{\mdseries\slshape M}}. }

 

\subsection{\textcolor{Chapter }{GroupRing}}
\logpage{[ 65, 1, 2 ]}\nobreak
\hyperdef{L}{X86D2CA90847C091B}{}
{\noindent\textcolor{FuncColor}{$\triangleright$\enspace\texttt{GroupRing({\mdseries\slshape R, G})\index{GroupRing@\texttt{GroupRing}}
\label{GroupRing}
}\hfill{\scriptsize (function)}}\\


 is the group ring of the group \mbox{\texttt{\mdseries\slshape G}}, over the ring \mbox{\texttt{\mdseries\slshape R}}. }

 

\subsection{\textcolor{Chapter }{IsFreeMagmaRing}}
\logpage{[ 65, 1, 3 ]}\nobreak
\hyperdef{L}{X7A24B95C8210BD09}{}
{\noindent\textcolor{FuncColor}{$\triangleright$\enspace\texttt{IsFreeMagmaRing({\mdseries\slshape D})\index{IsFreeMagmaRing@\texttt{IsFreeMagmaRing}}
\label{IsFreeMagmaRing}
}\hfill{\scriptsize (Category)}}\\


 A domain lies in the category \texttt{IsFreeMagmaRing} if it has been constructed as a free magma ring. In particular, if \mbox{\texttt{\mdseries\slshape D}} lies in this category then the operations \texttt{LeftActingDomain} (\ref{LeftActingDomain}) and \texttt{UnderlyingMagma} (\ref{UnderlyingMagma}) are applicable to \mbox{\texttt{\mdseries\slshape D}}, and yield the ring $R$ and the magma $M$ such that \mbox{\texttt{\mdseries\slshape D}} is the magma ring $RM$. 

 So being a magma ring in \textsf{GAP} includes the knowledge of the ring and the magma. Note that a magma ring $RM$ may abstractly be generated as a magma ring by a magma different from the
underlying magma $M$. For example, the group ring of the dihedral group of order $8$ over the field with $3$ elements is also spanned by a quaternion group of order $8$ over the same field. 

 
\begin{Verbatim}[commandchars=!@|,fontsize=\small,frame=single,label=Example]
  !gapprompt@gap>| !gapinput@d8:= DihedralGroup( 8 );|
  <pc group of size 8 with 3 generators>
  !gapprompt@gap>| !gapinput@rm:= FreeMagmaRing( GF(3), d8 );|
  <algebra-with-one over GF(3), with 3 generators>
  !gapprompt@gap>| !gapinput@emb:= Embedding( d8, rm );;|
  !gapprompt@gap>| !gapinput@gens:= List( GeneratorsOfGroup( d8 ), x -> x^emb );;|
  !gapprompt@gap>| !gapinput@x1:= gens[1] + gens[2];;|
  !gapprompt@gap>| !gapinput@x2:= ( gens[1] - gens[2] ) * gens[3];;|
  !gapprompt@gap>| !gapinput@x3:= gens[1] * gens[2] * ( One( rm ) - gens[3] );;|
  !gapprompt@gap>| !gapinput@g1:= x1 - x2 + x3;;|
  !gapprompt@gap>| !gapinput@g2:= x1 + x2;;|
  !gapprompt@gap>| !gapinput@q8:= Group( g1, g2 );;|
  !gapprompt@gap>| !gapinput@Size( q8 );|
  8
  !gapprompt@gap>| !gapinput@ForAny( [ d8, q8 ], IsAbelian );|
  false
  !gapprompt@gap>| !gapinput@List( [ d8, q8 ], g -> Number( AsList( g ), x -> Order( x ) = 2 ) );|
  [ 5, 1 ]
  !gapprompt@gap>| !gapinput@Dimension( Subspace( rm, q8 ) );|
  8
\end{Verbatim}
 }

 

\subsection{\textcolor{Chapter }{IsFreeMagmaRingWithOne}}
\logpage{[ 65, 1, 4 ]}\nobreak
\hyperdef{L}{X8382ED697A28CE67}{}
{\noindent\textcolor{FuncColor}{$\triangleright$\enspace\texttt{IsFreeMagmaRingWithOne({\mdseries\slshape obj})\index{IsFreeMagmaRingWithOne@\texttt{IsFreeMagmaRingWithOne}}
\label{IsFreeMagmaRingWithOne}
}\hfill{\scriptsize (Category)}}\\


 \texttt{IsFreeMagmaRingWithOne} is just a synonym for the meet of \texttt{IsFreeMagmaRing} (\ref{IsFreeMagmaRing}) and \texttt{IsMagmaWithOne} (\ref{IsMagmaWithOne}). }

 

\subsection{\textcolor{Chapter }{IsGroupRing}}
\logpage{[ 65, 1, 5 ]}\nobreak
\hyperdef{L}{X82C63644805EB1EE}{}
{\noindent\textcolor{FuncColor}{$\triangleright$\enspace\texttt{IsGroupRing({\mdseries\slshape obj})\index{IsGroupRing@\texttt{IsGroupRing}}
\label{IsGroupRing}
}\hfill{\scriptsize (property)}}\\


 A \emph{group ring} is a magma ring where the underlying magma is a group. }

 

\subsection{\textcolor{Chapter }{UnderlyingMagma}}
\logpage{[ 65, 1, 6 ]}\nobreak
\hyperdef{L}{X848D60417DFF7947}{}
{\noindent\textcolor{FuncColor}{$\triangleright$\enspace\texttt{UnderlyingMagma({\mdseries\slshape RM})\index{UnderlyingMagma@\texttt{UnderlyingMagma}}
\label{UnderlyingMagma}
}\hfill{\scriptsize (attribute)}}\\


 stores the underlying magma of a free magma ring. }

 

\subsection{\textcolor{Chapter }{AugmentationIdeal}}
\logpage{[ 65, 1, 7 ]}\nobreak
\hyperdef{L}{X7B21DB3E7CD80983}{}
{\noindent\textcolor{FuncColor}{$\triangleright$\enspace\texttt{AugmentationIdeal({\mdseries\slshape RG})\index{AugmentationIdeal@\texttt{AugmentationIdeal}}
\label{AugmentationIdeal}
}\hfill{\scriptsize (attribute)}}\\


 is the augmentation ideal of the group ring \mbox{\texttt{\mdseries\slshape RG}}, i.e., the kernel of the trivial representation of \mbox{\texttt{\mdseries\slshape RG}}. }

 }

  
\section{\textcolor{Chapter }{Elements of Free Magma Rings}}\label{Elements of Free Magma Rings}
\logpage{[ 65, 2, 0 ]}
\hyperdef{L}{X8402D3897F2C5955}{}
{
  In order to treat elements of free magma rings uniformly, also without an
external representation, the attributes \texttt{CoefficientsAndMagmaElements} (\ref{CoefficientsAndMagmaElements}) and \texttt{ZeroCoefficient} (\ref{ZeroCoefficient}) were introduced that allow one to ``take an element of an arbitrary magma ring into pieces''. 

 Conversely, for constructing magma ring elements from coefficients and magma
elements, \texttt{ElementOfMagmaRing} (\ref{ElementOfMagmaRing}) can be used. (Of course one can also embed each magma element into the magma
ring, see{\nobreakspace}\ref{Natural Embeddings related to Magma Rings}, and then form the linear combination, but many unnecessary intermediate
elements are created this way.) 

\subsection{\textcolor{Chapter }{IsMagmaRingObjDefaultRep}}
\logpage{[ 65, 2, 1 ]}\nobreak
\hyperdef{L}{X827B2D7D7E41780C}{}
{\noindent\textcolor{FuncColor}{$\triangleright$\enspace\texttt{IsMagmaRingObjDefaultRep({\mdseries\slshape obj})\index{IsMagmaRingObjDefaultRep@\texttt{IsMagmaRingObjDefaultRep}}
\label{IsMagmaRingObjDefaultRep}
}\hfill{\scriptsize (Representation)}}\\


 The default representation of a magma ring element is a list of length 2, at
first position the zero coefficient, at second position a list with the
coefficients at the even positions, and the magma elements at the odd
positions, with the ordering as defined for the magma elements. 

 It is assumed that arithmetic operations on magma rings produce only
normalized elements. }

 

\subsection{\textcolor{Chapter }{IsElementOfFreeMagmaRing}}
\logpage{[ 65, 2, 2 ]}\nobreak
\hyperdef{L}{X7D9C684A81E66310}{}
{\noindent\textcolor{FuncColor}{$\triangleright$\enspace\texttt{IsElementOfFreeMagmaRing({\mdseries\slshape obj})\index{IsElementOfFreeMagmaRing@\texttt{IsElementOfFreeMagmaRing}}
\label{IsElementOfFreeMagmaRing}
}\hfill{\scriptsize (Category)}}\\
\noindent\textcolor{FuncColor}{$\triangleright$\enspace\texttt{IsElementOfFreeMagmaRingCollection({\mdseries\slshape obj})\index{IsElementOfFreeMagmaRingCollection@\texttt{IsElementOfFreeMagmaRingCollection}}
\label{IsElementOfFreeMagmaRingCollection}
}\hfill{\scriptsize (Category)}}\\


 The category of elements of a free magma ring (See \texttt{IsFreeMagmaRing} (\ref{IsFreeMagmaRing})). }

 

\subsection{\textcolor{Chapter }{IsElementOfFreeMagmaRingFamily}}
\logpage{[ 65, 2, 3 ]}\nobreak
\hyperdef{L}{X869768AF7B444BF8}{}
{\noindent\textcolor{FuncColor}{$\triangleright$\enspace\texttt{IsElementOfFreeMagmaRingFamily({\mdseries\slshape Fam})\index{IsElementOfFreeMagmaRingFamily@\texttt{IsElementOfFreeMagmaRingFamily}}
\label{IsElementOfFreeMagmaRingFamily}
}\hfill{\scriptsize (Category)}}\\


 Elements of families in this category have trivial normalisation, i.e.,
efficient methods for \texttt{\texttt{\symbol{92}}=} and \texttt{\texttt{\symbol{92}}{\textless}}. }

 

\subsection{\textcolor{Chapter }{CoefficientsAndMagmaElements}}
\logpage{[ 65, 2, 4 ]}\nobreak
\hyperdef{L}{X843D1D8578C33513}{}
{\noindent\textcolor{FuncColor}{$\triangleright$\enspace\texttt{CoefficientsAndMagmaElements({\mdseries\slshape elm})\index{CoefficientsAndMagmaElements@\texttt{CoefficientsAndMagmaElements}}
\label{CoefficientsAndMagmaElements}
}\hfill{\scriptsize (attribute)}}\\


 is a list that contains at the odd positions the magma elements, and at the
even positions their coefficients in the element \mbox{\texttt{\mdseries\slshape elm}}. }

 

\subsection{\textcolor{Chapter }{ZeroCoefficient}}
\logpage{[ 65, 2, 5 ]}\nobreak
\hyperdef{L}{X78C3DB417E353390}{}
{\noindent\textcolor{FuncColor}{$\triangleright$\enspace\texttt{ZeroCoefficient({\mdseries\slshape elm})\index{ZeroCoefficient@\texttt{ZeroCoefficient}}
\label{ZeroCoefficient}
}\hfill{\scriptsize (attribute)}}\\


 For an element \mbox{\texttt{\mdseries\slshape elm}} of a magma ring (modulo relations) $RM$, \texttt{ZeroCoefficient} returns the zero element of the coefficient ring $R$. }

 

\subsection{\textcolor{Chapter }{ElementOfMagmaRing}}
\logpage{[ 65, 2, 6 ]}\nobreak
\hyperdef{L}{X8671DE0A81BEEFB0}{}
{\noindent\textcolor{FuncColor}{$\triangleright$\enspace\texttt{ElementOfMagmaRing({\mdseries\slshape Fam, zerocoeff, coeffs, mgmelms})\index{ElementOfMagmaRing@\texttt{ElementOfMagmaRing}}
\label{ElementOfMagmaRing}
}\hfill{\scriptsize (operation)}}\\


 \texttt{ElementOfMagmaRing} returns the element $\sum_{{i = 1}}^n c_i m_i'$, where $\mbox{\texttt{\mdseries\slshape coeffs}} = [ c_1, c_2, \ldots, c_n ]$ is a list of coefficients, $\mbox{\texttt{\mdseries\slshape mgmelms}} = [ m_1, m_2, \ldots, m_n ]$ is a list of magma elements, and $m_i'$ is the image of $m_i$ under an embedding of a magma containing $m_i$ into a magma ring whose elements lie in the family \mbox{\texttt{\mdseries\slshape Fam}}. \mbox{\texttt{\mdseries\slshape zerocoeff}} must be the zero of the coefficient ring containing the $c_i$. }

 }

  
\section{\textcolor{Chapter }{Natural Embeddings related to Magma Rings}}\label{Natural Embeddings related to Magma Rings}
\logpage{[ 65, 3, 0 ]}
\hyperdef{L}{X80366F1480ACD8DF}{}
{
  \index{Embedding@\texttt{Embedding}!for magma rings} Neither the coefficient ring $R$ nor the magma $M$ are regarded as subsets of the magma ring $RM$, so one has to use \emph{embeddings} (see{\nobreakspace}\texttt{Embedding} (\ref{Embedding:for two domains})) explicitly whenever one needs for example the magma ring element
corresponding to a given magma element. 

 
\begin{Verbatim}[commandchars=!@|,fontsize=\small,frame=single,label=Example]
  !gapprompt@gap>| !gapinput@f:= Rationals;;  g:= SymmetricGroup( 3 );;|
  !gapprompt@gap>| !gapinput@fg:= FreeMagmaRing( f, g );|
  <algebra-with-one over Rationals, with 2 generators>
  !gapprompt@gap>| !gapinput@Dimension( fg );|
  6
  !gapprompt@gap>| !gapinput@gens:= GeneratorsOfAlgebraWithOne( fg );|
  [ (1)*(1,2,3), (1)*(1,2) ]
  !gapprompt@gap>| !gapinput@( 3*gens[1] - 2*gens[2] ) * ( gens[1] + gens[2] );|
  (-2)*()+(3)*(2,3)+(3)*(1,3,2)+(-2)*(1,3)
  !gapprompt@gap>| !gapinput@One( fg );|
  (1)*()
  !gapprompt@gap>| !gapinput@emb:= Embedding( g, fg );;|
  !gapprompt@gap>| !gapinput@elm:= (1,2,3)^emb;  elm in fg;|
  (1)*(1,2,3)
  true
  !gapprompt@gap>| !gapinput@new:= elm + One( fg );|
  (1)*()+(1)*(1,2,3)
  !gapprompt@gap>| !gapinput@new^2;|
  (1)*()+(2)*(1,2,3)+(1)*(1,3,2)
  !gapprompt@gap>| !gapinput@emb2:= Embedding( f, fg );;|
  !gapprompt@gap>| !gapinput@elm:= One( f )^emb2;  elm in fg;|
  (1)*()
  true
\end{Verbatim}
 }

  
\section{\textcolor{Chapter }{Magma Rings modulo Relations}}\label{Magma Rings modulo Relations}
\logpage{[ 65, 4, 0 ]}
\hyperdef{L}{X81B002EE799B5E77}{}
{
  A more general construction than that of free magma rings allows one to create
rings that are not free $R$\texttt{\symbol{45}}modules on a given magma $M$ but arise from the magma ring $RM$ by factoring out certain identities. Examples for such structures are finitely
presented (associative) algebras and free Lie algebras (see{\nobreakspace}\texttt{FreeLieAlgebra} (\ref{FreeLieAlgebra:for ring, rank (and name)})).  

 In \textsf{GAP}, the use of magma rings modulo relations is limited to situations where a
normal form of the elements is known and where one wants to guarantee that all
elements actually constructed are in normal form. (In particular, the
computation of the normal form must be cheap.) This is because the methods for
comparing elements in magma rings modulo relations via \texttt{\texttt{\symbol{92}}=} and \texttt{\texttt{\symbol{92}}{\textless}} just compare the involved coefficients and magma elements, and also the vector
space functions regard those monomials as linearly independent over the
coefficients ring that actually occur in the representation of an element of a
magma ring modulo relations. 

 Thus only very special finitely presented algebras will be represented as
magma rings modulo relations, in general finitely presented algebras are dealt
with via the mechanism described in Chapter{\nobreakspace}\ref{Finitely Presented Algebras}. 

\subsection{\textcolor{Chapter }{IsElementOfMagmaRingModuloRelations}}
\logpage{[ 65, 4, 1 ]}\nobreak
\hyperdef{L}{X869D54847E881848}{}
{\noindent\textcolor{FuncColor}{$\triangleright$\enspace\texttt{IsElementOfMagmaRingModuloRelations({\mdseries\slshape obj})\index{IsElementOfMagmaRingModuloRelations@\texttt{IsElementOfMagmaRingModuloRelations}}
\label{IsElementOfMagmaRingModuloRelations}
}\hfill{\scriptsize (Category)}}\\
\noindent\textcolor{FuncColor}{$\triangleright$\enspace\texttt{IsElementOfMagmaRingModuloRelationsCollection({\mdseries\slshape obj})\index{IsElementOfMagmaRingModuloRelationsCollection@\texttt{IsElement}\-\texttt{Of}\-\texttt{Magma}\-\texttt{Ring}\-\texttt{Modulo}\-\texttt{Relations}\-\texttt{Collection}}
\label{IsElementOfMagmaRingModuloRelationsCollection}
}\hfill{\scriptsize (Category)}}\\


 This category is used, e.{\nobreakspace}g., for elements of free Lie algebras. }

 

\subsection{\textcolor{Chapter }{IsElementOfMagmaRingModuloRelationsFamily}}
\logpage{[ 65, 4, 2 ]}\nobreak
\hyperdef{L}{X875BEB1A840FFAA4}{}
{\noindent\textcolor{FuncColor}{$\triangleright$\enspace\texttt{IsElementOfMagmaRingModuloRelationsFamily({\mdseries\slshape Fam})\index{IsElementOfMagmaRingModuloRelationsFamily@\texttt{IsElement}\-\texttt{Of}\-\texttt{Magma}\-\texttt{Ring}\-\texttt{Modulo}\-\texttt{Relations}\-\texttt{Family}}
\label{IsElementOfMagmaRingModuloRelationsFamily}
}\hfill{\scriptsize (Category)}}\\


 The family category for the category \texttt{IsElementOfMagmaRingModuloRelations} (\ref{IsElementOfMagmaRingModuloRelations}). }

 

\subsection{\textcolor{Chapter }{NormalizedElementOfMagmaRingModuloRelations}}
\logpage{[ 65, 4, 3 ]}\nobreak
\hyperdef{L}{X85956ED27FA6AC68}{}
{\noindent\textcolor{FuncColor}{$\triangleright$\enspace\texttt{NormalizedElementOfMagmaRingModuloRelations({\mdseries\slshape F, descr})\index{NormalizedElementOfMagmaRingModuloRelations@\texttt{Normalized}\-\texttt{Element}\-\texttt{Of}\-\texttt{Magma}\-\texttt{Ring}\-\texttt{Modulo}\-\texttt{Relations}}
\label{NormalizedElementOfMagmaRingModuloRelations}
}\hfill{\scriptsize (operation)}}\\


 Let \mbox{\texttt{\mdseries\slshape F}} be a family of magma ring elements modulo relations, and \mbox{\texttt{\mdseries\slshape descr}} the description of an element in a magma ring modulo relations. \texttt{NormalizedElementOfMagmaRingModuloRelations} returns a description of the same element, but normalized
w.r.t.{\nobreakspace}the relations. So two elements are equal if and only if
the result of \texttt{NormalizedElementOfMagmaRingModuloRelations} is equal for their internal data, that is, \texttt{CoefficientsAndMagmaElements} (\ref{CoefficientsAndMagmaElements}) will return the same for the corresponding two elements. 

 \texttt{NormalizedElementOfMagmaRingModuloRelations} is allowed to return \mbox{\texttt{\mdseries\slshape descr}} itself, it need not make a copy. This is the case for example in the case of
free magma rings. }

 

\subsection{\textcolor{Chapter }{IsMagmaRingModuloRelations}}
\logpage{[ 65, 4, 4 ]}\nobreak
\hyperdef{L}{X804B5AAB813E184D}{}
{\noindent\textcolor{FuncColor}{$\triangleright$\enspace\texttt{IsMagmaRingModuloRelations({\mdseries\slshape obj})\index{IsMagmaRingModuloRelations@\texttt{IsMagmaRingModuloRelations}}
\label{IsMagmaRingModuloRelations}
}\hfill{\scriptsize (Category)}}\\


 A \textsf{GAP} object lies in the category \texttt{IsMagmaRingModuloRelations} if it has been constructed as a magma ring modulo relations. Each element of
such a ring has a unique normal form, so \texttt{CoefficientsAndMagmaElements} (\ref{CoefficientsAndMagmaElements}) is well\texttt{\symbol{45}}defined for it. 

 This category is not inherited to factor structures, which are in general best
described as finitely presented algebras, see Chapter{\nobreakspace}\ref{Finitely Presented Algebras}. }

 }

  
\section{\textcolor{Chapter }{Magma Rings modulo the Span of a Zero Element}}\label{Magma Rings modulo the Span of a Zero Element}
\logpage{[ 65, 5, 0 ]}
\hyperdef{L}{X7D859DBF81DFA751}{}
{
  

\subsection{\textcolor{Chapter }{IsElementOfMagmaRingModuloSpanOfZeroFamily}}
\logpage{[ 65, 5, 1 ]}\nobreak
\hyperdef{L}{X7B3D45A6802B695C}{}
{\noindent\textcolor{FuncColor}{$\triangleright$\enspace\texttt{IsElementOfMagmaRingModuloSpanOfZeroFamily({\mdseries\slshape Fam})\index{IsElementOfMagmaRingModuloSpanOfZeroFamily@\texttt{IsElement}\-\texttt{Of}\-\texttt{Magma}\-\texttt{Ring}\-\texttt{Modulo}\-\texttt{Span}\-\texttt{Of}\-\texttt{Zero}\-\texttt{Family}}
\label{IsElementOfMagmaRingModuloSpanOfZeroFamily}
}\hfill{\scriptsize (Category)}}\\


 We need this for the normalization method, which takes a family as first
argument. }

 

\subsection{\textcolor{Chapter }{IsMagmaRingModuloSpanOfZero}}
\logpage{[ 65, 5, 2 ]}\nobreak
\hyperdef{L}{X872713EE84DA9B72}{}
{\noindent\textcolor{FuncColor}{$\triangleright$\enspace\texttt{IsMagmaRingModuloSpanOfZero({\mdseries\slshape RM})\index{IsMagmaRingModuloSpanOfZero@\texttt{IsMagmaRingModuloSpanOfZero}}
\label{IsMagmaRingModuloSpanOfZero}
}\hfill{\scriptsize (Category)}}\\


 The category of magma rings modulo the span of a zero element. }

 

\subsection{\textcolor{Chapter }{MagmaRingModuloSpanOfZero}}
\logpage{[ 65, 5, 3 ]}\nobreak
\hyperdef{L}{X7A7F880D7D7D3722}{}
{\noindent\textcolor{FuncColor}{$\triangleright$\enspace\texttt{MagmaRingModuloSpanOfZero({\mdseries\slshape R, M, z})\index{MagmaRingModuloSpanOfZero@\texttt{MagmaRingModuloSpanOfZero}}
\label{MagmaRingModuloSpanOfZero}
}\hfill{\scriptsize (function)}}\\


 Let \mbox{\texttt{\mdseries\slshape R}} be a ring, \mbox{\texttt{\mdseries\slshape M}} a magma, and \mbox{\texttt{\mdseries\slshape z}} an element of \mbox{\texttt{\mdseries\slshape M}} with the property that $\mbox{\texttt{\mdseries\slshape z}} * m = \mbox{\texttt{\mdseries\slshape z}}$ holds for all $m \in M$. The element \mbox{\texttt{\mdseries\slshape z}} could be called a ``zero element'' of \mbox{\texttt{\mdseries\slshape M}}, but note that in general \mbox{\texttt{\mdseries\slshape z}} cannot be obtained as \texttt{Zero( }$m$\texttt{ )} for each $m \in M$, so this situation does not match the definition of \texttt{Zero} (\ref{Zero}). 

 \texttt{MagmaRingModuloSpanOfZero} returns the magma ring $\mbox{\texttt{\mdseries\slshape R}}\mbox{\texttt{\mdseries\slshape M}}$ modulo the relation given by the identification of \mbox{\texttt{\mdseries\slshape z}} with zero. This is an example of a magma ring modulo relations,
see{\nobreakspace}\ref{Magma Rings modulo Relations}. }

 }

  
\section{\textcolor{Chapter }{Technical Details about the Implementation of Magma Rings}}\label{Technical Details about the Implementation of Magma Rings}
\logpage{[ 65, 6, 0 ]}
\hyperdef{L}{X79889F017F2EB7ED}{}
{
  The \emph{family} containing elements in the magma ring $RM$ in fact contains all elements with coefficients in the family of elements of $R$ and magma elements in the family of elements of $M$. So arithmetic operations with coefficients outside $R$ or with magma elements outside $M$ might create elements outside $RM$. 

 It should be mentioned that each call of \texttt{FreeMagmaRing} (\ref{FreeMagmaRing}) creates a new family of elements, so for example the elements of two group
rings of permutation groups over the same ring lie in different families and
therefore are regarded as different. 

 
\begin{Verbatim}[commandchars=!@|,fontsize=\small,frame=single,label=Example]
  !gapprompt@gap>| !gapinput@g:= SymmetricGroup( 3 );;|
  !gapprompt@gap>| !gapinput@h:= AlternatingGroup( 3 );;|
  !gapprompt@gap>| !gapinput@IsSubset( g, h );|
  true
  !gapprompt@gap>| !gapinput@f:= GF(2);;|
  !gapprompt@gap>| !gapinput@fg:= GroupRing( f, g );|
  <algebra-with-one over GF(2), with 2 generators>
  !gapprompt@gap>| !gapinput@fh:= GroupRing( f, h );|
  <algebra-with-one over GF(2), with 1 generator>
  !gapprompt@gap>| !gapinput@IsSubset( fg, fh );|
  false
  !gapprompt@gap>| !gapinput@o1:= One( fh );  o2:= One( fg );  o1 = o2;|
  (Z(2)^0)*()
  (Z(2)^0)*()
  false
  !gapprompt@gap>| !gapinput@emb:= Embedding( g, fg );;|
  !gapprompt@gap>| !gapinput@im:= Image( emb, h );|
  <group of size 3 with 1 generator>
  !gapprompt@gap>| !gapinput@IsSubset( fg, im );|
  true
\end{Verbatim}
 

 There is \emph{no} generic \emph{external representation} for elements in an arbitrary free magma ring. For example, polynomials are
elements of a free magma ring, and they have an external representation
relying on the special form of the underlying monomials. On the other hand,
elements in a group ring of a permutation group do not admit such an external
representation. 

 For convenience, magma rings constructed with \texttt{FreeAlgebra} (\ref{FreeAlgebra:for ring, rank (and name)}), \texttt{FreeAssociativeAlgebra} (\ref{FreeAssociativeAlgebra:for ring, rank (and name)}), \texttt{FreeAlgebraWithOne} (\ref{FreeAlgebraWithOne:for ring, rank (and name)}), and \texttt{FreeAssociativeAlgebraWithOne} (\ref{FreeAssociativeAlgebraWithOne:for ring, rank (and name)}) support an external representation of their elements, which is defined as a
list of length 2, the first entry being the zero coefficient, the second being
a list with the external representations of the magma elements at the odd
positions and the corresponding coefficients at the even positions. 

 As the above examples show, there are several possible representations of
magma ring elements, the representations used for polynomials (see Chapter
{\nobreakspace}\ref{Polynomials and Rational Functions}) as well as the default representation \texttt{IsMagmaRingObjDefaultRep} (\ref{IsMagmaRingObjDefaultRep}) of magma ring elements. The latter simply stores the zero coefficient and a
list containing the coefficients of the element at the even positions and the
corresponding magma elements at the odd positions, where the succession is
compatible with the ordering of magma elements via \texttt{\texttt{\symbol{92}}{\textless}}. }

 }

             
\chapter{\textcolor{Chapter }{Polynomials and Rational Functions}}\label{Polynomials and Rational Functions}
\logpage{[ 66, 0, 0 ]}
\hyperdef{L}{X7A14A6588268810A}{}
{
  Let $R$ be a commutative ring\texttt{\symbol{45}}with\texttt{\symbol{45}}one. We call
a free associative algebra $A$ over $R$ a \emph{polynomial ring} over $R$. The free generators of $A$ are called \emph{indeterminates} (to avoid naming conflicts with the word \emph{variables} which will be used to denote \textsf{GAP} variables only) , they are usually denoted by $x_1, x_2, \ldots$. The number of indeterminates is called the \emph{rank} of $A$. The elements of $A$ are called \emph{polynomials}. Products of indeterminates are called \emph{monomials}, every polynomial can be expressed as a finite sum of products of monomials
with ring elements in a form like $r_{{1,0}} x_1 + r_{{1,1}} x_1 x_2 + r_{{0,1}} x_2 + \cdots$ with $r_{{i,j}} \in R$. 

 A polynomial ring of rank 1 is called an \emph{univariate} polynomial ring, its elements are \emph{univariate polynomials}. 

 Polynomial rings of smaller rank naturally embed in rings of higher rank; if $S$ is a subring of $R$ then a polynomial ring over $S$ naturally embeds in a polynomial ring over $R$ of the same rank. Note however that \textsf{GAP} does not consider $R$ as a subset of a polynomial ring over $R$; for example the zero of $R$ ($0$) and the zero of the polynomial ring ($0x^0$) are different objects. 

 Internally, indeterminates are represented by positive integers, but it is
possible to give names to them to have them printed in a nicer way. Beware,
however that there is not necessarily any relation between the way an
indeterminate is called and the way it is printed. See section \ref{Indeterminates} for details. 

 If $R$ is an integral domain, the polynomial ring $A$ over $R$ is an integral domain as well and one can therefore form its quotient field $Q$. This field is called a \emph{field of rational functions}. Again $A$ embeds naturally into $Q$ and \textsf{GAP} will perform this embedding implicitly. (In fact it implements the ring of
rational functions over $R$.) To avoid problems with leading coefficients, however, $R$ must be a unique factorization domain.  
\section{\textcolor{Chapter }{Indeterminates}}\label{Indeterminates}
\logpage{[ 66, 1, 0 ]}
\hyperdef{L}{X7A8FADCD875826DA}{}
{
  Internally, indeterminates are created for a \emph{family} of objects (for example all elements of finite fields in characteristic $3$ are in one family). Thus a variable ``x'' over the rationals is also an ``x'' over the integers, while an ``x'' over \texttt{GF(3)} is different. 

 Within one family, every indeterminate has a number \mbox{\texttt{\mdseries\slshape nr}} and as long as no other names have been assigned, this indeterminate will be
displayed as ``\texttt{x{\textunderscore}\mbox{\texttt{\mdseries\slshape nr}}}''. Indeterminate numbers can be arbitrary nonnegative integers. 

 It is possible to assign names to indeterminates; these names are strings and
only provide a means for printing the indeterminates in a nice way.
Indeterminates that have not been assigned a name will be printed as ``\texttt{x{\textunderscore}\mbox{\texttt{\mdseries\slshape nr}}}''. 

 (Because of this printing convention, the name \texttt{x{\textunderscore}\mbox{\texttt{\mdseries\slshape nr}}} is interpreted specially to always denote the variable with internal number \mbox{\texttt{\mdseries\slshape nr}}.) 

 The indeterminate names have not necessarily any relations to variable names:
this means that an indeterminate whose name is ``\texttt{x}'' cannot be accessed using the variable \texttt{x}, unless \texttt{x} was defined to be that indeterminate. 

 When asking for indeterminates with certain names, \textsf{GAP} usually will take the first (with respect to the internal numbering)
indeterminates that are not yet named, name these accordingly and return them.
Thus when asking for named indeterminates, no relation between names and
indeterminate numbers can be guaranteed. The attribute \texttt{IndeterminateNumberOfLaurentPolynomial(\mbox{\texttt{\mdseries\slshape indet}})} will return the number of the indeterminate \mbox{\texttt{\mdseries\slshape indet}}. 

 When asked to create an indeterminate with a name that exists already for the
family, \textsf{GAP} will by default return this existing indeterminate. If you explicitly want a \emph{new} indeterminate, distinct from the already existing one with the \emph{same} name, you can add the \texttt{new} option to the function call. (This is in most cases not a good idea.) 

 
\begin{Verbatim}[commandchars=!@|,fontsize=\small,frame=single,label=Example]
  !gapprompt@gap>| !gapinput@R:=PolynomialRing(GF(3),["x","y","z"]);|
  GF(3)[x,y,z]
  !gapprompt@gap>| !gapinput@List(IndeterminatesOfPolynomialRing(R),|
  !gapprompt@>| !gapinput@  IndeterminateNumberOfLaurentPolynomial);|
  [ 1, 2, 3 ]
  !gapprompt@gap>| !gapinput@R:=PolynomialRing(GF(3),["z"]);|
  GF(3)[z]
  !gapprompt@gap>| !gapinput@List(IndeterminatesOfPolynomialRing(R),|
  !gapprompt@>| !gapinput@  IndeterminateNumberOfLaurentPolynomial);|
  [ 3 ]
  !gapprompt@gap>| !gapinput@R:=PolynomialRing(GF(3),["x","y","z"]:new);|
  GF(3)[x,y,z]
  !gapprompt@gap>| !gapinput@List(IndeterminatesOfPolynomialRing(R),|
  !gapprompt@>| !gapinput@  IndeterminateNumberOfLaurentPolynomial);|
  [ 4, 5, 6 ]
  !gapprompt@gap>| !gapinput@R:=PolynomialRing(GF(3),["z"]);|
  GF(3)[z]
  !gapprompt@gap>| !gapinput@List(IndeterminatesOfPolynomialRing(R),|
  !gapprompt@>| !gapinput@  IndeterminateNumberOfLaurentPolynomial);|
  [ 3 ]
\end{Verbatim}
 
\subsection{\textcolor{Chapter }{Indeterminate}}\logpage{[ 66, 1, 1 ]}
\hyperdef{L}{X79D0380D7FA39F7D}{}
{
\noindent\textcolor{FuncColor}{$\triangleright$\enspace\texttt{Indeterminate({\mdseries\slshape R[, nr]})\index{Indeterminate@\texttt{Indeterminate}!for a ring (and a number)}
\label{Indeterminate:for a ring (and a number)}
}\hfill{\scriptsize (operation)}}\\
\noindent\textcolor{FuncColor}{$\triangleright$\enspace\texttt{Indeterminate({\mdseries\slshape R[, name][, avoid]})\index{Indeterminate@\texttt{Indeterminate}!for a ring (and a name, and an exclusion list)}
\label{Indeterminate:for a ring (and a name, and an exclusion list)}
}\hfill{\scriptsize (operation)}}\\
\noindent\textcolor{FuncColor}{$\triangleright$\enspace\texttt{Indeterminate({\mdseries\slshape fam, nr})\index{Indeterminate@\texttt{Indeterminate}!for a family and a number}
\label{Indeterminate:for a family and a number}
}\hfill{\scriptsize (operation)}}\\
\noindent\textcolor{FuncColor}{$\triangleright$\enspace\texttt{X({\mdseries\slshape R[, nr]})\index{X@\texttt{X}!for a ring (and a number)}
\label{X:for a ring (and a number)}
}\hfill{\scriptsize (operation)}}\\
\noindent\textcolor{FuncColor}{$\triangleright$\enspace\texttt{X({\mdseries\slshape R[, name][, avoid]})\index{X@\texttt{X}!for a ring (and a name, and an exclusion list)}
\label{X:for a ring (and a name, and an exclusion list)}
}\hfill{\scriptsize (operation)}}\\
\noindent\textcolor{FuncColor}{$\triangleright$\enspace\texttt{X({\mdseries\slshape fam, nr})\index{X@\texttt{X}!for a family and a number}
\label{X:for a family and a number}
}\hfill{\scriptsize (operation)}}\\


 returns the indeterminate number \mbox{\texttt{\mdseries\slshape nr}} over the ring \mbox{\texttt{\mdseries\slshape R}}. If \mbox{\texttt{\mdseries\slshape nr}} is not given it defaults to 1. If the number is not specified a list \mbox{\texttt{\mdseries\slshape avoid}} of indeterminates may be given. The function will return an indeterminate that
is guaranteed to be different from all the indeterminates in the list \mbox{\texttt{\mdseries\slshape avoid}}. The third usage returns an indeterminate called \mbox{\texttt{\mdseries\slshape name}} (also avoiding the indeterminates in \mbox{\texttt{\mdseries\slshape avoid}} if given). 

 \texttt{X} is simply a synonym for \texttt{Indeterminate}. 

 
\begin{Verbatim}[commandchars=!@|,fontsize=\small,frame=single,label=Example]
  !gapprompt@gap>| !gapinput@x:=Indeterminate(GF(3),"x");|
  x
  !gapprompt@gap>| !gapinput@y:=X(GF(3),"y");z:=X(GF(3),"X");|
  y
  X
  !gapprompt@gap>| !gapinput@X(GF(3),2);|
  y
  !gapprompt@gap>| !gapinput@X(GF(3),"x_3");|
  X
  !gapprompt@gap>| !gapinput@X(GF(3),[y,z]);|
  x
\end{Verbatim}
 }

 

\subsection{\textcolor{Chapter }{IndeterminateNumberOfUnivariateRationalFunction}}
\logpage{[ 66, 1, 2 ]}\nobreak
\hyperdef{L}{X816C8D797C804380}{}
{\noindent\textcolor{FuncColor}{$\triangleright$\enspace\texttt{IndeterminateNumberOfUnivariateRationalFunction({\mdseries\slshape rfun})\index{IndeterminateNumberOfUnivariateRationalFunction@\texttt{Indeterminate}\-\texttt{Number}\-\texttt{Of}\-\texttt{Univariate}\-\texttt{Rational}\-\texttt{Function}}
\label{IndeterminateNumberOfUnivariateRationalFunction}
}\hfill{\scriptsize (attribute)}}\\


 returns the number of the indeterminate in which the univariate rational
function \mbox{\texttt{\mdseries\slshape rfun}} is expressed. (This also provides a way to obtain the number of a given
indeterminate.) 

 A constant rational function might not possess an indeterminate number. In
this case \texttt{IndeterminateNumberOfUnivariateRationalFunction} will default to a value of 1. Therefore two univariate polynomials may be
considered to be in the same univariate polynomial ring if their
indeterminates have the same number or one if of them is constant. (see
also{\nobreakspace}\texttt{CIUnivPols} (\ref{CIUnivPols}) and{\nobreakspace}\texttt{IsLaurentPolynomialDefaultRep} (\ref{IsLaurentPolynomialDefaultRep})). }

 

\subsection{\textcolor{Chapter }{IndeterminateOfUnivariateRationalFunction}}
\logpage{[ 66, 1, 3 ]}\nobreak
\hyperdef{L}{X7A2FA46885EF403D}{}
{\noindent\textcolor{FuncColor}{$\triangleright$\enspace\texttt{IndeterminateOfUnivariateRationalFunction({\mdseries\slshape rfun})\index{IndeterminateOfUnivariateRationalFunction@\texttt{Indeterminate}\-\texttt{Of}\-\texttt{Univariate}\-\texttt{Rational}\-\texttt{Function}}
\label{IndeterminateOfUnivariateRationalFunction}
}\hfill{\scriptsize (attribute)}}\\


 returns the indeterminate in which the univariate rational function \mbox{\texttt{\mdseries\slshape rfun}} is expressed. (cf. \texttt{IndeterminateNumberOfUnivariateRationalFunction} (\ref{IndeterminateNumberOfUnivariateRationalFunction}).) 
\begin{Verbatim}[commandchars=!@|,fontsize=\small,frame=single,label=Example]
  !gapprompt@gap>| !gapinput@IndeterminateNumberOfUnivariateRationalFunction(z);|
  3
  !gapprompt@gap>| !gapinput@IndeterminateOfUnivariateRationalFunction(z^5+z);|
  X
\end{Verbatim}
 }

 

\subsection{\textcolor{Chapter }{IndeterminateName}}
\logpage{[ 66, 1, 4 ]}\nobreak
\hyperdef{L}{X7FD4AC807A1C8E27}{}
{\noindent\textcolor{FuncColor}{$\triangleright$\enspace\texttt{IndeterminateName({\mdseries\slshape fam, nr})\index{IndeterminateName@\texttt{IndeterminateName}}
\label{IndeterminateName}
}\hfill{\scriptsize (operation)}}\\
\noindent\textcolor{FuncColor}{$\triangleright$\enspace\texttt{HasIndeterminateName({\mdseries\slshape fam, nr})\index{HasIndeterminateName@\texttt{HasIndeterminateName}}
\label{HasIndeterminateName}
}\hfill{\scriptsize (operation)}}\\
\noindent\textcolor{FuncColor}{$\triangleright$\enspace\texttt{SetIndeterminateName({\mdseries\slshape fam, nr, name})\index{SetIndeterminateName@\texttt{SetIndeterminateName}}
\label{SetIndeterminateName}
}\hfill{\scriptsize (operation)}}\\


 \texttt{SetIndeterminateName} assigns the name \mbox{\texttt{\mdseries\slshape name}} to indeterminate \mbox{\texttt{\mdseries\slshape nr}} in the rational functions family \mbox{\texttt{\mdseries\slshape fam}}. It issues an error if the indeterminate was already named. 

 \texttt{IndeterminateName} returns the name of the \mbox{\texttt{\mdseries\slshape nr}}\texttt{\symbol{45}}th indeterminate (and returns \texttt{fail} if no name has been assigned). 

 \texttt{HasIndeterminateName} tests whether indeterminate \mbox{\texttt{\mdseries\slshape nr}} has already been assigned a name. 

 
\begin{Verbatim}[commandchars=!@|,fontsize=\small,frame=single,label=Example]
  !gapprompt@gap>| !gapinput@IndeterminateName(FamilyObj(x),2);|
  "y"
  !gapprompt@gap>| !gapinput@HasIndeterminateName(FamilyObj(x),4);|
  false
  !gapprompt@gap>| !gapinput@SetIndeterminateName(FamilyObj(x),10,"bla");|
  !gapprompt@gap>| !gapinput@Indeterminate(GF(3),10);|
  bla
\end{Verbatim}
 

 As a convenience there is a special method installed for \texttt{SetName} that will assign a name to an indeterminate. 

 
\begin{Verbatim}[commandchars=!@|,fontsize=\small,frame=single,label=Example]
  !gapprompt@gap>| !gapinput@a:=Indeterminate(GF(3),5);|
  x_5
  !gapprompt@gap>| !gapinput@SetName(a,"ah");|
  !gapprompt@gap>| !gapinput@a^5+a;|
  ah^5+ah
\end{Verbatim}
 }

 

\subsection{\textcolor{Chapter }{CIUnivPols}}
\logpage{[ 66, 1, 5 ]}\nobreak
\hyperdef{L}{X791A06E67F784328}{}
{\noindent\textcolor{FuncColor}{$\triangleright$\enspace\texttt{CIUnivPols({\mdseries\slshape upol1, upol2})\index{CIUnivPols@\texttt{CIUnivPols}}
\label{CIUnivPols}
}\hfill{\scriptsize (function)}}\\


 This function (whose name stands for ``common indeterminate of univariate polynomials'') takes two univariate polynomials as arguments. If both polynomials are given
in the same indeterminate number \mbox{\texttt{\mdseries\slshape indnum}} (in this case they are ``compatible'' as univariate polynomials) it returns \mbox{\texttt{\mdseries\slshape indnum}}. In all other cases it returns \texttt{fail}. \texttt{CIUnivPols} also accepts if either polynomial is constant but formally expressed in
another indeterminate, in this situation the indeterminate of the other
polynomial is selected. }

 }

  
\section{\textcolor{Chapter }{Operations for Rational Functions}}\label{Operations for Rational Functions}
\logpage{[ 66, 2, 0 ]}
\hyperdef{L}{X86A68FD582F4F757}{}
{
  The rational functions form a field, therefore all arithmetic operations are
applicable to rational functions. 

 \index{addition!rational functions} \index{subtraction!rational functions} \index{product!rational functions} \index{quotient!rational functions} \texttt{\mbox{\texttt{\mdseries\slshape f}} + \mbox{\texttt{\mdseries\slshape g}}} 

 \texttt{\mbox{\texttt{\mdseries\slshape f}} \texttt{\symbol{45}} \mbox{\texttt{\mdseries\slshape g}}} 

 \texttt{\mbox{\texttt{\mdseries\slshape f}} * \mbox{\texttt{\mdseries\slshape g}}} 

 \texttt{\mbox{\texttt{\mdseries\slshape f}} / \mbox{\texttt{\mdseries\slshape g}}} 

 
\begin{Verbatim}[commandchars=!@|,fontsize=\small,frame=single,label=Example]
  !gapprompt@gap>| !gapinput@x:=Indeterminate(Rationals,1);;y:=Indeterminate(Rationals,2);;|
  !gapprompt@gap>| !gapinput@f:=3+x*y+x^5;;g:=5+x^2*y+x*y^2;;|
  !gapprompt@gap>| !gapinput@a:=g/f;|
  (x_1^2*x_2+x_1*x_2^2+5)/(x_1^5+x_1*x_2+3)
\end{Verbatim}
 

 Note that the quotient \texttt{\mbox{\texttt{\mdseries\slshape f}}/\mbox{\texttt{\mdseries\slshape g}}} of two polynomials might be represented as a rational function again. If \mbox{\texttt{\mdseries\slshape g}} is known to divide \mbox{\texttt{\mdseries\slshape f}} the call \texttt{Quotient(\mbox{\texttt{\mdseries\slshape f}},\mbox{\texttt{\mdseries\slshape g}})} (see{\nobreakspace}\texttt{Quotient} (\ref{Quotient})) should be used instead. 

 \index{mod!Laurent polynomials} \texttt{\mbox{\texttt{\mdseries\slshape f}} mod \mbox{\texttt{\mdseries\slshape g}}}  

 For two Laurent polynomials \mbox{\texttt{\mdseries\slshape f}} and \mbox{\texttt{\mdseries\slshape g}}, \texttt{\mbox{\texttt{\mdseries\slshape f}} mod \mbox{\texttt{\mdseries\slshape g}}} is the Euclidean remainder (see{\nobreakspace}\texttt{EuclideanRemainder} (\ref{EuclideanRemainder})) of \mbox{\texttt{\mdseries\slshape f}} modulo \mbox{\texttt{\mdseries\slshape g}}.        

 As calculating a multivariate Gcd can be expensive, it is not guaranteed that
rational functions will always be represented as a quotient of coprime
polynomials. In certain unfortunate situations this might lead to a degree
explosion. To ensure cancellation you can use \texttt{Gcd} (\ref{Gcd:for (a ring and) several elements}) on the \texttt{NumeratorOfRationalFunction} (\ref{NumeratorOfRationalFunction}) and \texttt{DenominatorOfRationalFunction} (\ref{DenominatorOfRationalFunction}) values of a given rational function. 

 All polynomials as well as all the univariate polynomials in the same
indeterminate form subrings of this field. If two rational functions are known
to be in the same subring, the result will be expressed as element in this
subring. }

  
\section{\textcolor{Chapter }{Comparison of Rational Functions}}\label{Comparison of Rational Functions}
\logpage{[ 66, 3, 0 ]}
\hyperdef{L}{X824B6D328643CE04}{}
{
  \index{comparison!rational functions} \texttt{\mbox{\texttt{\mdseries\slshape f}} = \mbox{\texttt{\mdseries\slshape g}}} 

 Two rational functions \mbox{\texttt{\mdseries\slshape f}} and \mbox{\texttt{\mdseries\slshape g}} are equal if the product \texttt{Numerator(\mbox{\texttt{\mdseries\slshape f}}) * Denominator(\mbox{\texttt{\mdseries\slshape g}})} equals \texttt{Numerator(\mbox{\texttt{\mdseries\slshape g}}) * Denominator(\mbox{\texttt{\mdseries\slshape f}})}. 

 
\begin{Verbatim}[commandchars=!@|,fontsize=\small,frame=single,label=Example]
  !gapprompt@gap>| !gapinput@x:=Indeterminate(Rationals,"x");;y:=Indeterminate(Rationals,"y");;|
  !gapprompt@gap>| !gapinput@f:=3+x*y+x^5;;g:=5+x^2*y+x*y^2;;|
  !gapprompt@gap>| !gapinput@a:=g/f;|
  (x^2*y+x*y^2+5)/(x^5+x*y+3)
  !gapprompt@gap>| !gapinput@b:=(g*f)/(f^2);|
  (x^7*y+x^6*y^2+5*x^5+x^3*y^2+x^2*y^3+3*x^2*y+3*x*y^2+5*x*y+15)/(x^10+2\
  *x^6*y+6*x^5+x^2*y^2+6*x*y+9)
  !gapprompt@gap>| !gapinput@a=b;|
  true
\end{Verbatim}
 

 \index{smaller!rational functions} \texttt{\mbox{\texttt{\mdseries\slshape f}} {\textless} \mbox{\texttt{\mdseries\slshape g}}} 

 The ordering of rational functions is defined in several steps. Monomials
(products of indeterminates) are sorted first by degree, then
lexicographically (with $x_1>x_2$) (see{\nobreakspace}\texttt{MonomialGrlexOrdering} (\ref{MonomialGrlexOrdering})). Products of monomials with ring elements (``terms'') are compared first by their monomials and then by their coefficients. 

 
\begin{Verbatim}[commandchars=!@|,fontsize=\small,frame=single,label=Example]
  !gapprompt@gap>| !gapinput@x>y;|
  true
  !gapprompt@gap>| !gapinput@x^2*y<x*y^2;|
  false
  !gapprompt@gap>| !gapinput@x*y<x^2*y;|
  true
  !gapprompt@gap>| !gapinput@x^2*y < 5* y*x^2;|
  true
\end{Verbatim}
 Polynomials are compared by comparing the largest terms in turn until they
differ. 
\begin{Verbatim}[commandchars=!@|,fontsize=\small,frame=single,label=Example]
  !gapprompt@gap>| !gapinput@x+y<y;|
  false
  !gapprompt@gap>| !gapinput@x<x+1;|
  true
\end{Verbatim}
 Rational functions are compared by comparing the polynomial \texttt{Numerator(\mbox{\texttt{\mdseries\slshape f}}) * Denominator(\mbox{\texttt{\mdseries\slshape g}})} with the polynomial \texttt{Numerator(\mbox{\texttt{\mdseries\slshape g}}) * Denominator(\mbox{\texttt{\mdseries\slshape f}})}. (As the ordering of monomials used by \textsf{GAP} is invariant under multiplication this is independent of common factors in
numerator and denominator.) 
\begin{Verbatim}[commandchars=!@|,fontsize=\small,frame=single,label=Example]
  !gapprompt@gap>| !gapinput@f/g<g/f;|
  false
  !gapprompt@gap>| !gapinput@f/g<(g*g)/(f*g);|
  false
\end{Verbatim}
 

 For univariate polynomials this reduces to an ordering first by total degree
and then lexicographically on the coefficients. }

  
\section{\textcolor{Chapter }{Properties and Attributes of Rational Functions}}\label{Properties and Attributes of Rational Functions}
\logpage{[ 66, 4, 0 ]}
\hyperdef{L}{X7D871EA180E9486C}{}
{
  All these tests are applicable to \emph{every} rational function. Depending on the internal representation of the rational
function, however some of these tests (in particular, univariateness) might be
expensive in some cases. 

 For reasons of performance within algorithms it can be useful to use other
attributes, which give a slightly more technical representation. See
section{\nobreakspace}\ref{The Representations of Rational Functions} for details. 

\subsection{\textcolor{Chapter }{IsPolynomialFunction}}
\logpage{[ 66, 4, 1 ]}\nobreak
\hyperdef{L}{X86C92F677DA9347F}{}
{\noindent\textcolor{FuncColor}{$\triangleright$\enspace\texttt{IsPolynomialFunction({\mdseries\slshape obj})\index{IsPolynomialFunction@\texttt{IsPolynomialFunction}}
\label{IsPolynomialFunction}
}\hfill{\scriptsize (Category)}}\\
\noindent\textcolor{FuncColor}{$\triangleright$\enspace\texttt{IsRationalFunction({\mdseries\slshape obj})\index{IsRationalFunction@\texttt{IsRationalFunction}}
\label{IsRationalFunction}
}\hfill{\scriptsize (Category)}}\\


 A rational function is an element of the quotient field of a polynomial ring
over an UFD. It is represented as a quotient of two polynomials, its numerator
(see{\nobreakspace}\texttt{NumeratorOfRationalFunction} (\ref{NumeratorOfRationalFunction})) and its denominator (see{\nobreakspace}\texttt{DenominatorOfRationalFunction} (\ref{DenominatorOfRationalFunction})) 

 A polynomial function is an element of a polynomial ring (not necessarily an
UFD), or a rational function. 

 \textsf{GAP} considers \texttt{IsRationalFunction} as a subcategory of \texttt{IsPolynomialFunction}. }

 

\subsection{\textcolor{Chapter }{NumeratorOfRationalFunction}}
\logpage{[ 66, 4, 2 ]}\nobreak
\hyperdef{L}{X7D7D2667803D8D8A}{}
{\noindent\textcolor{FuncColor}{$\triangleright$\enspace\texttt{NumeratorOfRationalFunction({\mdseries\slshape ratfun})\index{NumeratorOfRationalFunction@\texttt{NumeratorOfRationalFunction}}
\label{NumeratorOfRationalFunction}
}\hfill{\scriptsize (attribute)}}\\


 returns the numerator of the rational function \mbox{\texttt{\mdseries\slshape ratfun}}. 

 As no proper multivariate gcd has been implemented yet, numerators and
denominators are not guaranteed to be reduced! }

 

\subsection{\textcolor{Chapter }{DenominatorOfRationalFunction}}
\logpage{[ 66, 4, 3 ]}\nobreak
\hyperdef{L}{X78DC1B5B866ADB6C}{}
{\noindent\textcolor{FuncColor}{$\triangleright$\enspace\texttt{DenominatorOfRationalFunction({\mdseries\slshape ratfun})\index{DenominatorOfRationalFunction@\texttt{DenominatorOfRationalFunction}}
\label{DenominatorOfRationalFunction}
}\hfill{\scriptsize (attribute)}}\\


 returns the denominator of the rational function \mbox{\texttt{\mdseries\slshape ratfun}}. 

 As no proper multivariate gcd has been implemented yet, numerators and
denominators are not guaranteed to be reduced! 
\begin{Verbatim}[commandchars=!@|,fontsize=\small,frame=single,label=Example]
  !gapprompt@gap>| !gapinput@x:=Indeterminate(Rationals,1);;y:=Indeterminate(Rationals,2);;|
  !gapprompt@gap>| !gapinput@DenominatorOfRationalFunction((x*y+x^2)/y);|
  y
  !gapprompt@gap>| !gapinput@NumeratorOfRationalFunction((x*y+x^2)/y);|
  x^2+x*y
\end{Verbatim}
 }

 

\subsection{\textcolor{Chapter }{IsPolynomial}}
\logpage{[ 66, 4, 4 ]}\nobreak
\hyperdef{L}{X7974B0707C8DAB6C}{}
{\noindent\textcolor{FuncColor}{$\triangleright$\enspace\texttt{IsPolynomial({\mdseries\slshape ratfun})\index{IsPolynomial@\texttt{IsPolynomial}}
\label{IsPolynomial}
}\hfill{\scriptsize (property)}}\\


 A polynomial is a rational function whose denominator is one. (If the
coefficients family forms a field this is equivalent to the denominator being
constant.) 

 If the base family is not a field, it may be impossible to represent the
quotient of a polynomial by a ring element as a polynomial again, but it will
have to be represented as a rational function. 
\begin{Verbatim}[commandchars=!@|,fontsize=\small,frame=single,label=Example]
  !gapprompt@gap>| !gapinput@IsPolynomial((x*y+x^2*y^3)/y);|
  true
  !gapprompt@gap>| !gapinput@IsPolynomial((x*y+x^2)/y);|
  false
\end{Verbatim}
 }

 

\subsection{\textcolor{Chapter }{AsPolynomial}}
\logpage{[ 66, 4, 5 ]}\nobreak
\hyperdef{L}{X7914771F7C6013EF}{}
{\noindent\textcolor{FuncColor}{$\triangleright$\enspace\texttt{AsPolynomial({\mdseries\slshape poly})\index{AsPolynomial@\texttt{AsPolynomial}}
\label{AsPolynomial}
}\hfill{\scriptsize (attribute)}}\\


 If \mbox{\texttt{\mdseries\slshape poly}} is a rational function that is a polynomial this attribute returns an equal
rational function $p$ such that $p$ is equal to its numerator and the denominator of $p$ is one. 
\begin{Verbatim}[commandchars=!@|,fontsize=\small,frame=single,label=Example]
  !gapprompt@gap>| !gapinput@AsPolynomial((x*y+x^2*y^3)/y);|
  x^2*y^2+x
\end{Verbatim}
 }

 

\subsection{\textcolor{Chapter }{IsUnivariateRationalFunction}}
\logpage{[ 66, 4, 6 ]}\nobreak
\hyperdef{L}{X8738F73583273FCA}{}
{\noindent\textcolor{FuncColor}{$\triangleright$\enspace\texttt{IsUnivariateRationalFunction({\mdseries\slshape ratfun})\index{IsUnivariateRationalFunction@\texttt{IsUnivariateRationalFunction}}
\label{IsUnivariateRationalFunction}
}\hfill{\scriptsize (property)}}\\


 A rational function is univariate if its numerator and its denominator are
both polynomials in the same one indeterminate. The attribute \texttt{IndeterminateNumberOfUnivariateRationalFunction} (\ref{IndeterminateNumberOfUnivariateRationalFunction}) can be used to obtain the number of this common indeterminate. }

 

\subsection{\textcolor{Chapter }{CoefficientsOfUnivariateRationalFunction}}
\logpage{[ 66, 4, 7 ]}\nobreak
\hyperdef{L}{X7F1F67527A35A9CE}{}
{\noindent\textcolor{FuncColor}{$\triangleright$\enspace\texttt{CoefficientsOfUnivariateRationalFunction({\mdseries\slshape rfun})\index{CoefficientsOfUnivariateRationalFunction@\texttt{Coefficients}\-\texttt{Of}\-\texttt{Univariate}\-\texttt{Rational}\-\texttt{Function}}
\label{CoefficientsOfUnivariateRationalFunction}
}\hfill{\scriptsize (attribute)}}\\


 if \mbox{\texttt{\mdseries\slshape rfun}} is a univariate rational function, this attribute returns a list \texttt{[ \mbox{\texttt{\mdseries\slshape ncof}}, \mbox{\texttt{\mdseries\slshape dcof}}, \mbox{\texttt{\mdseries\slshape val}} ]} where \mbox{\texttt{\mdseries\slshape ncof}} and \mbox{\texttt{\mdseries\slshape dcof}} are coefficient lists of univariate polynomials \mbox{\texttt{\mdseries\slshape n}} and \mbox{\texttt{\mdseries\slshape d}} and a valuation \mbox{\texttt{\mdseries\slshape val}} such that $\mbox{\texttt{\mdseries\slshape rfun}} = x^{\mbox{\texttt{\mdseries\slshape val}}} \cdot \mbox{\texttt{\mdseries\slshape n}} / \mbox{\texttt{\mdseries\slshape d}}$ where $x$ is the variable with the number given by \texttt{IndeterminateNumberOfUnivariateRationalFunction} (\ref{IndeterminateNumberOfUnivariateRationalFunction}). Numerator and denominator are guaranteed to be cancelled. }

 

\subsection{\textcolor{Chapter }{IsUnivariatePolynomial}}
\logpage{[ 66, 4, 8 ]}\nobreak
\hyperdef{L}{X86A2546685D0016B}{}
{\noindent\textcolor{FuncColor}{$\triangleright$\enspace\texttt{IsUnivariatePolynomial({\mdseries\slshape ratfun})\index{IsUnivariatePolynomial@\texttt{IsUnivariatePolynomial}}
\label{IsUnivariatePolynomial}
}\hfill{\scriptsize (property)}}\\


 A univariate polynomial is a polynomial in only one indeterminate. }

 

\subsection{\textcolor{Chapter }{CoefficientsOfUnivariatePolynomial}}
\logpage{[ 66, 4, 9 ]}\nobreak
\hyperdef{L}{X78C9524D7F2708C2}{}
{\noindent\textcolor{FuncColor}{$\triangleright$\enspace\texttt{CoefficientsOfUnivariatePolynomial({\mdseries\slshape pol})\index{CoefficientsOfUnivariatePolynomial@\texttt{CoefficientsOfUnivariatePolynomial}}
\label{CoefficientsOfUnivariatePolynomial}
}\hfill{\scriptsize (attribute)}}\\


 \texttt{CoefficientsOfUnivariatePolynomial} returns the coefficient list of the polynomial \mbox{\texttt{\mdseries\slshape pol}}, sorted in ascending order. (It returns the empty list if \mbox{\texttt{\mdseries\slshape pol}} is 0.) }

 

\subsection{\textcolor{Chapter }{IsLaurentPolynomial}}
\logpage{[ 66, 4, 10 ]}\nobreak
\hyperdef{L}{X79138FF28213B6EC}{}
{\noindent\textcolor{FuncColor}{$\triangleright$\enspace\texttt{IsLaurentPolynomial({\mdseries\slshape ratfun})\index{IsLaurentPolynomial@\texttt{IsLaurentPolynomial}}
\label{IsLaurentPolynomial}
}\hfill{\scriptsize (property)}}\\


 A Laurent polynomial is a univariate rational function whose denominator is a
monomial. Therefore every univariate polynomial is a Laurent polynomial. 

 The attribute \texttt{CoefficientsOfLaurentPolynomial} (\ref{CoefficientsOfLaurentPolynomial}) gives a compact representation as Laurent polynomial. }

 

\subsection{\textcolor{Chapter }{IsConstantRationalFunction}}
\logpage{[ 66, 4, 11 ]}\nobreak
\hyperdef{L}{X7F2A49208341C2A8}{}
{\noindent\textcolor{FuncColor}{$\triangleright$\enspace\texttt{IsConstantRationalFunction({\mdseries\slshape ratfun})\index{IsConstantRationalFunction@\texttt{IsConstantRationalFunction}}
\label{IsConstantRationalFunction}
}\hfill{\scriptsize (property)}}\\


 A constant rational function is a function whose numerator and denominator are
polynomials of degree 0. }

 

\subsection{\textcolor{Chapter }{IsPrimitivePolynomial}}
\logpage{[ 66, 4, 12 ]}\nobreak
\hyperdef{L}{X834B54947FAADEA4}{}
{\noindent\textcolor{FuncColor}{$\triangleright$\enspace\texttt{IsPrimitivePolynomial({\mdseries\slshape F, pol})\index{IsPrimitivePolynomial@\texttt{IsPrimitivePolynomial}}
\label{IsPrimitivePolynomial}
}\hfill{\scriptsize (operation)}}\\


 For a univariate polynomial \mbox{\texttt{\mdseries\slshape pol}} of degree $d$ in the indeterminate $X$, with coefficients in a finite field \mbox{\texttt{\mdseries\slshape F}} with $q$ elements, \texttt{IsPrimitivePolynomial} returns \texttt{true} if 
\begin{enumerate}
\item  \mbox{\texttt{\mdseries\slshape pol}} divides $X^{{q^d-1}} - 1$, and 
\item  for each prime divisor $p$ of $q^d - 1$, \mbox{\texttt{\mdseries\slshape pol}} does not divide $X^{{(q^d-1)/p}} - 1$, 
\end{enumerate}
 and \texttt{false} otherwise. }

 

\subsection{\textcolor{Chapter }{SplittingField}}
\logpage{[ 66, 4, 13 ]}\nobreak
\hyperdef{L}{X87531E03849391C1}{}
{\noindent\textcolor{FuncColor}{$\triangleright$\enspace\texttt{SplittingField({\mdseries\slshape f})\index{SplittingField@\texttt{SplittingField}}
\label{SplittingField}
}\hfill{\scriptsize (attribute)}}\\


 returns the smallest field which contains the coefficients of \mbox{\texttt{\mdseries\slshape f}} and the roots of \mbox{\texttt{\mdseries\slshape f}}. }

 }

  
\section{\textcolor{Chapter }{Univariate Polynomials}}\label{Univariate Polynomials}
\logpage{[ 66, 5, 0 ]}
\hyperdef{L}{X82E2F1707FC2E553}{}
{
  Some of the operations are actually defined on the larger domain of Laurent
polynomials (see \ref{Laurent Polynomials}). For this section you can simply ignore the word ``Laurent'' if it occurs in a description. 

 

\subsection{\textcolor{Chapter }{UnivariatePolynomial}}
\logpage{[ 66, 5, 1 ]}\nobreak
\hyperdef{L}{X8379F8CB7D0076BA}{}
{\noindent\textcolor{FuncColor}{$\triangleright$\enspace\texttt{UnivariatePolynomial({\mdseries\slshape ring, cofs[, ind]})\index{UnivariatePolynomial@\texttt{UnivariatePolynomial}}
\label{UnivariatePolynomial}
}\hfill{\scriptsize (operation)}}\\


 constructs an univariate polynomial over the ring \mbox{\texttt{\mdseries\slshape ring}} in the indeterminate \mbox{\texttt{\mdseries\slshape ind}} with the coefficients given by \mbox{\texttt{\mdseries\slshape cofs}}. }

 

\subsection{\textcolor{Chapter }{UnivariatePolynomialByCoefficients}}
\logpage{[ 66, 5, 2 ]}\nobreak
\hyperdef{L}{X85178A3E7B4F11E0}{}
{\noindent\textcolor{FuncColor}{$\triangleright$\enspace\texttt{UnivariatePolynomialByCoefficients({\mdseries\slshape fam, cofs, ind})\index{UnivariatePolynomialByCoefficients@\texttt{UnivariatePolynomialByCoefficients}}
\label{UnivariatePolynomialByCoefficients}
}\hfill{\scriptsize (operation)}}\\


 constructs an univariate polynomial over the coefficients family \mbox{\texttt{\mdseries\slshape fam}} and in the indeterminate \mbox{\texttt{\mdseries\slshape ind}} with the coefficients given by \mbox{\texttt{\mdseries\slshape cofs}}. This function should be used in algorithms to create polynomials as it
avoids overhead associated with \texttt{UnivariatePolynomial} (\ref{UnivariatePolynomial}). }

 

\subsection{\textcolor{Chapter }{DegreeOfLaurentPolynomial}}
\logpage{[ 66, 5, 3 ]}\nobreak
\hyperdef{L}{X78AF77C383245254}{}
{\noindent\textcolor{FuncColor}{$\triangleright$\enspace\texttt{DegreeOfLaurentPolynomial({\mdseries\slshape pol})\index{DegreeOfLaurentPolynomial@\texttt{DegreeOfLaurentPolynomial}}
\label{DegreeOfLaurentPolynomial}
}\hfill{\scriptsize (attribute)}}\\


 The degree of a univariate (Laurent) polynomial \mbox{\texttt{\mdseries\slshape pol}} is the largest exponent $n$ of a monomial $x^n$ of \mbox{\texttt{\mdseries\slshape pol}}. The degree of a zero polynomial is defined to be \texttt{\texttt{\symbol{45}}infinity}. 
\begin{Verbatim}[commandchars=!@|,fontsize=\small,frame=single,label=Example]
  !gapprompt@gap>| !gapinput@p:=UnivariatePolynomial(Rationals,[1,2,3,4],1);|
  4*x^3+3*x^2+2*x+1
  !gapprompt@gap>| !gapinput@UnivariatePolynomialByCoefficients(FamilyObj(1),[9,2,3,4],73);|
  4*x_73^3+3*x_73^2+2*x_73+9
  !gapprompt@gap>| !gapinput@CoefficientsOfUnivariatePolynomial(p);|
  [ 1, 2, 3, 4 ]
  !gapprompt@gap>| !gapinput@DegreeOfLaurentPolynomial(p);|
  3
  !gapprompt@gap>| !gapinput@DegreeOfLaurentPolynomial(Zero(p));|
  -infinity
  !gapprompt@gap>| !gapinput@IndeterminateNumberOfLaurentPolynomial(p);|
  1
  !gapprompt@gap>| !gapinput@IndeterminateOfLaurentPolynomial(p);|
  x
\end{Verbatim}
 }

 

\subsection{\textcolor{Chapter }{RootsOfPolynomial}}
\logpage{[ 66, 5, 4 ]}\nobreak
\hyperdef{L}{X7CBB760C87B04F75}{}
{\noindent\textcolor{FuncColor}{$\triangleright$\enspace\texttt{RootsOfPolynomial({\mdseries\slshape [R, ]p})\index{RootsOfPolynomial@\texttt{RootsOfPolynomial}}
\label{RootsOfPolynomial}
}\hfill{\scriptsize (function)}}\\


 For a univariate polynomial \mbox{\texttt{\mdseries\slshape p}}, this function returns all roots of \mbox{\texttt{\mdseries\slshape p}} over the ring \mbox{\texttt{\mdseries\slshape R}}. If the ring is not specified, it defaults to the ring specified by the
coefficients of \mbox{\texttt{\mdseries\slshape p}} via \texttt{DefaultRing} (\ref{DefaultRing:for ring elements})). 
\begin{Verbatim}[commandchars=!@|,fontsize=\small,frame=single,label=Example]
  !gapprompt@gap>| !gapinput@x:=X(Rationals,"x");;p:=x^4-1;|
  x^4-1
  !gapprompt@gap>| !gapinput@RootsOfPolynomial(p);|
  [ 1, -1 ]
  !gapprompt@gap>| !gapinput@RootsOfPolynomial(CF(4),p);|
  [ 1, -1, E(4), -E(4) ]
\end{Verbatim}
 }

 

\subsection{\textcolor{Chapter }{RootsOfUPol}}
\logpage{[ 66, 5, 5 ]}\nobreak
\hyperdef{L}{X80CEB10D7879767F}{}
{\noindent\textcolor{FuncColor}{$\triangleright$\enspace\texttt{RootsOfUPol({\mdseries\slshape [field, ]upol})\index{RootsOfUPol@\texttt{RootsOfUPol}}
\label{RootsOfUPol}
}\hfill{\scriptsize (function)}}\\


 This function returns a list of all roots of the univariate polynomial \mbox{\texttt{\mdseries\slshape upol}} in its default domain. If the optional argument \mbox{\texttt{\mdseries\slshape field}} is a field then the roots in this field are computed. If \mbox{\texttt{\mdseries\slshape field}} is the string \texttt{"split"} then the splitting field of the polynomial is taken. 
\begin{Verbatim}[commandchars=!@|,fontsize=\small,frame=single,label=Example]
  !gapprompt@gap>| !gapinput@RootsOfUPol(50-45*x-6*x^2+x^3);|
  [ 10, 1, -5 ]
\end{Verbatim}
 }

 

\subsection{\textcolor{Chapter }{QuotRemLaurpols}}
\logpage{[ 66, 5, 6 ]}\nobreak
\hyperdef{L}{X7887FBC78149BB0C}{}
{\noindent\textcolor{FuncColor}{$\triangleright$\enspace\texttt{QuotRemLaurpols({\mdseries\slshape left, right, mode})\index{QuotRemLaurpols@\texttt{QuotRemLaurpols}}
\label{QuotRemLaurpols}
}\hfill{\scriptsize (function)}}\\


 This internal function for euclidean division of polynomials takes two
polynomials \mbox{\texttt{\mdseries\slshape left}} and \mbox{\texttt{\mdseries\slshape right}} and computes their quotient. No test is performed whether the arguments indeed
are polynomials. Depending on the integer variable \mbox{\texttt{\mdseries\slshape mode}}, which may take values in a range from 1 to 4, it returns respectively: 
\begin{enumerate}
\item  the quotient (there might be some remainder), 
\item  the remainder, 
\item  a list \texttt{[\mbox{\texttt{\mdseries\slshape q}},\mbox{\texttt{\mdseries\slshape r}}]} of quotient and remainder, 
\item  the quotient if there is no remainder and \texttt{fail} otherwise. 
\end{enumerate}
 }

 

\subsection{\textcolor{Chapter }{UnivariatenessTestRationalFunction}}
\logpage{[ 66, 5, 7 ]}\nobreak
\hyperdef{L}{X7DDADF157879EFBF}{}
{\noindent\textcolor{FuncColor}{$\triangleright$\enspace\texttt{UnivariatenessTestRationalFunction({\mdseries\slshape f})\index{UnivariatenessTestRationalFunction@\texttt{UnivariatenessTestRationalFunction}}
\label{UnivariatenessTestRationalFunction}
}\hfill{\scriptsize (function)}}\\


 takes a rational function \mbox{\texttt{\mdseries\slshape f}} and tests whether it is univariate rational function (or even a Laurent
polynomial). It returns a list \texttt{[isunivariate, indet, islaurent, cofs]}. 

 If \mbox{\texttt{\mdseries\slshape f}} is a univariate rational function then \texttt{isunivariate} is \texttt{true} and \texttt{indet} is the number of the appropriate indeterminate. 

 Furthermore, if \mbox{\texttt{\mdseries\slshape f}} is a Laurent polynomial, then \texttt{islaurent} is also \texttt{true}. In this case the fourth entry, \texttt{cofs}, is the value of the attribute \texttt{CoefficientsOfLaurentPolynomial} (\ref{CoefficientsOfLaurentPolynomial}) for \mbox{\texttt{\mdseries\slshape f}}. 

 If \texttt{isunivariate} is \texttt{true} but \texttt{islaurent} is \texttt{false}, then \texttt{cofs} is the value of the attribute \texttt{CoefficientsOfUnivariateRationalFunction} (\ref{CoefficientsOfUnivariateRationalFunction}) for \mbox{\texttt{\mdseries\slshape f}}. 

 Otherwise, each entry of the returned list is equal to \texttt{fail}. As there is no proper multivariate gcd, this may also happen for the
rational function which may be reduced to univariate (see example). 
\begin{Verbatim}[commandchars=!@|,fontsize=\small,frame=single,label=Example]
  !gapprompt@gap>| !gapinput@UnivariatenessTestRationalFunction( 50-45*x-6*x^2+x^3 );|
  [ true, 1, true, [ [ 50, -45, -6, 1 ], 0 ] ]
  !gapprompt@gap>| !gapinput@UnivariatenessTestRationalFunction( (-6*y^2+y^3) / (y+1) );|
  [ true, 2, false, [ [ -6, 1 ], [ 1, 1 ], 2 ] ]
  !gapprompt@gap>| !gapinput@UnivariatenessTestRationalFunction( (-6*y^2+y^3) / (x+1));|
  [ false, fail, false, fail ]
  !gapprompt@gap>| !gapinput@UnivariatenessTestRationalFunction( ((y+2)*(x+1)) / ((y-1)*(x+1)) );|
  [ fail, fail, fail, fail ]
\end{Verbatim}
 }

 

\subsection{\textcolor{Chapter }{InfoPoly}}
\logpage{[ 66, 5, 8 ]}\nobreak
\hyperdef{L}{X7A3BC96B7A50DE98}{}
{\noindent\textcolor{FuncColor}{$\triangleright$\enspace\texttt{InfoPoly\index{InfoPoly@\texttt{InfoPoly}}
\label{InfoPoly}
}\hfill{\scriptsize (info class)}}\\


 is the info class for univariate polynomials. }

 We remark that some functions for multivariate polynomials (which will be
defined in the following sections) permit a different syntax for univariate
polynomials which drops the requirement to specify the indeterminate. Examples
are \texttt{Value} (\ref{Value:for rat. function, a list of indeterminates, a value (and a one)}), \texttt{Discriminant} (\ref{Discriminant}), \texttt{Derivative} (\ref{Derivative}), \texttt{LeadingCoefficient} (\ref{LeadingCoefficient}) and \texttt{LeadingMonomial} (\ref{LeadingMonomial}): 
\begin{Verbatim}[commandchars=!@|,fontsize=\small,frame=single,label=Example]
  !gapprompt@gap>| !gapinput@p:=UnivariatePolynomial(Rationals,[1,2,3,4],1);|
  4*x^3+3*x^2+2*x+1
  !gapprompt@gap>| !gapinput@Value(p,Z(5));|
  Z(5)^2
  !gapprompt@gap>| !gapinput@LeadingCoefficient(p);|
  4
  !gapprompt@gap>| !gapinput@Derivative(p);|
  12*x^2+6*x+2
\end{Verbatim}
 }

  
\section{\textcolor{Chapter }{Polynomials as Univariate Polynomials in one Indeterminate}}\label{Polynomials as Univariate Polynomials in one Indeterminate}
\logpage{[ 66, 6, 0 ]}
\hyperdef{L}{X81499B5A823E6EA3}{}
{
  

\subsection{\textcolor{Chapter }{DegreeIndeterminate}}
\logpage{[ 66, 6, 1 ]}\nobreak
\hyperdef{L}{X826B99B17ABD11BE}{}
{\noindent\textcolor{FuncColor}{$\triangleright$\enspace\texttt{DegreeIndeterminate({\mdseries\slshape pol, ind})\index{DegreeIndeterminate@\texttt{DegreeIndeterminate}}
\label{DegreeIndeterminate}
}\hfill{\scriptsize (operation)}}\\


 returns the degree of the polynomial \mbox{\texttt{\mdseries\slshape pol}} in the indeterminate (or indeterminate number) \mbox{\texttt{\mdseries\slshape ind}}. 
\begin{Verbatim}[commandchars=!@|,fontsize=\small,frame=single,label=Example]
  !gapprompt@gap>| !gapinput@f:=x^5+3*x*y+9*y^7+4*y^5*x+3*y+2;|
  9*y^7+4*x*y^5+x^5+3*x*y+3*y+2
  !gapprompt@gap>| !gapinput@DegreeIndeterminate(f,1);|
  5
  !gapprompt@gap>| !gapinput@DegreeIndeterminate(f,y);|
  7
\end{Verbatim}
 }

 

\subsection{\textcolor{Chapter }{PolynomialCoefficientsOfPolynomial}}
\logpage{[ 66, 6, 2 ]}\nobreak
\hyperdef{L}{X85646FD07F9C60F5}{}
{\noindent\textcolor{FuncColor}{$\triangleright$\enspace\texttt{PolynomialCoefficientsOfPolynomial({\mdseries\slshape pol, ind})\index{PolynomialCoefficientsOfPolynomial@\texttt{PolynomialCoefficientsOfPolynomial}}
\label{PolynomialCoefficientsOfPolynomial}
}\hfill{\scriptsize (operation)}}\\


 \texttt{PolynomialCoefficientsOfPolynomial} returns the coefficient list (whose entries are polynomials not involving the
indeterminate \mbox{\texttt{\mdseries\slshape ind}}) describing the polynomial \mbox{\texttt{\mdseries\slshape pol}} viewed as a polynomial in \mbox{\texttt{\mdseries\slshape ind}}. Instead of the indeterminate, \mbox{\texttt{\mdseries\slshape ind}} can also be an indeterminate number. 
\begin{Verbatim}[commandchars=!@|,fontsize=\small,frame=single,label=Example]
  !gapprompt@gap>| !gapinput@PolynomialCoefficientsOfPolynomial(f,2);|
  [ x^5+2, 3*x+3, 0, 0, 0, 4*x, 0, 9 ]
\end{Verbatim}
 }

 

\subsection{\textcolor{Chapter }{LeadingCoefficient}}
\logpage{[ 66, 6, 3 ]}\nobreak
\hyperdef{L}{X80710E9B7D8340BD}{}
{\noindent\textcolor{FuncColor}{$\triangleright$\enspace\texttt{LeadingCoefficient({\mdseries\slshape pol})\index{LeadingCoefficient@\texttt{LeadingCoefficient}}
\label{LeadingCoefficient}
}\hfill{\scriptsize (operation)}}\\


 returns the leading coefficient (that is the coefficient of the leading
monomial, see{\nobreakspace}\texttt{LeadingMonomial} (\ref{LeadingMonomial})) of the polynomial \mbox{\texttt{\mdseries\slshape pol}}. }

 

\subsection{\textcolor{Chapter }{LeadingMonomial}}
\logpage{[ 66, 6, 4 ]}\nobreak
\hyperdef{L}{X7B3EAE41795598A5}{}
{\noindent\textcolor{FuncColor}{$\triangleright$\enspace\texttt{LeadingMonomial({\mdseries\slshape pol})\index{LeadingMonomial@\texttt{LeadingMonomial}}
\label{LeadingMonomial}
}\hfill{\scriptsize (operation)}}\\


 returns the leading monomial (with respect to the ordering given by \texttt{MonomialExtGrlexLess} (\ref{MonomialExtGrlexLess})) of the polynomial \mbox{\texttt{\mdseries\slshape pol}} as a list containing indeterminate numbers and exponents. 
\begin{Verbatim}[commandchars=!@|,fontsize=\small,frame=single,label=Example]
  !gapprompt@gap>| !gapinput@LeadingCoefficient(f,1);|
  1
  !gapprompt@gap>| !gapinput@LeadingCoefficient(f,2);|
  9
  !gapprompt@gap>| !gapinput@LeadingMonomial(f);|
  [ 2, 7 ]
  !gapprompt@gap>| !gapinput@LeadingCoefficient(f);|
  9
\end{Verbatim}
 }

 

\subsection{\textcolor{Chapter }{Derivative}}
\logpage{[ 66, 6, 5 ]}\nobreak
\hyperdef{L}{X7B57CEE2780D0E0B}{}
{\noindent\textcolor{FuncColor}{$\triangleright$\enspace\texttt{Derivative({\mdseries\slshape ratfun[, ind]})\index{Derivative@\texttt{Derivative}}
\label{Derivative}
}\hfill{\scriptsize (attribute)}}\\


 If \mbox{\texttt{\mdseries\slshape ratfun}} is a univariate rational function then \texttt{Derivative} returns the \emph{derivative} of \mbox{\texttt{\mdseries\slshape ufun}} by its indeterminate. For a rational function \mbox{\texttt{\mdseries\slshape ratfun}}, the derivative by the indeterminate \mbox{\texttt{\mdseries\slshape ind}} is returned, regarding \mbox{\texttt{\mdseries\slshape ratfun}} as univariate in \mbox{\texttt{\mdseries\slshape ind}}. Instead of the desired indeterminate, also the number of this indeterminate
can be given as \mbox{\texttt{\mdseries\slshape ind}}. 
\begin{Verbatim}[commandchars=!@|,fontsize=\small,frame=single,label=Example]
  !gapprompt@gap>| !gapinput@Derivative(f,2);|
  63*y^6+20*x*y^4+3*x+3
\end{Verbatim}
 }

 

\subsection{\textcolor{Chapter }{Discriminant}}
\logpage{[ 66, 6, 6 ]}\nobreak
\hyperdef{L}{X7C7D790A7D6E11AD}{}
{\noindent\textcolor{FuncColor}{$\triangleright$\enspace\texttt{Discriminant({\mdseries\slshape pol[, ind]})\index{Discriminant@\texttt{Discriminant}}
\label{Discriminant}
}\hfill{\scriptsize (operation)}}\\


 If \mbox{\texttt{\mdseries\slshape pol}} is a univariate polynomial then \texttt{Discriminant} returns the \emph{discriminant} of \mbox{\texttt{\mdseries\slshape pol}} by its indeterminate. The two\texttt{\symbol{45}}argument form returns the
discriminant of a polynomial \mbox{\texttt{\mdseries\slshape pol}} by the indeterminate number \mbox{\texttt{\mdseries\slshape ind}}, regarding \mbox{\texttt{\mdseries\slshape pol}} as univariate in this indeterminate. Instead of the indeterminate number, the
indeterminate itself can also be given as \mbox{\texttt{\mdseries\slshape ind}}. 
\begin{Verbatim}[commandchars=!@|,fontsize=\small,frame=single,label=Example]
  !gapprompt@gap>| !gapinput@Discriminant(f,1);|
  20503125*y^28+262144*y^25+27337500*y^22+19208040*y^21+1474560*y^17+136\
  68750*y^16+18225000*y^15+6075000*y^14+1105920*y^13+3037500*y^10+648972\
  0*y^9+4050000*y^8+900000*y^7+62208*y^5+253125*y^4+675000*y^3+675000*y^\
  2+300000*y+50000
  !gapprompt@gap>| !gapinput@Discriminant(f,1) = Discriminant(f,x);|
  true
\end{Verbatim}
 }

 

\subsection{\textcolor{Chapter }{Resultant}}
\logpage{[ 66, 6, 7 ]}\nobreak
\hyperdef{L}{X857AD5587EF49029}{}
{\noindent\textcolor{FuncColor}{$\triangleright$\enspace\texttt{Resultant({\mdseries\slshape pol1, pol2, ind})\index{Resultant@\texttt{Resultant}}
\label{Resultant}
}\hfill{\scriptsize (operation)}}\\


 computes the resultant of the polynomials \mbox{\texttt{\mdseries\slshape pol1}} and \mbox{\texttt{\mdseries\slshape pol2}} with respect to the indeterminate \mbox{\texttt{\mdseries\slshape ind}}, or indeterminate number \mbox{\texttt{\mdseries\slshape ind}}. The resultant considers \mbox{\texttt{\mdseries\slshape pol1}} and \mbox{\texttt{\mdseries\slshape pol2}} as univariate in \mbox{\texttt{\mdseries\slshape ind}} and returns an element of the corresponding base ring (which might be a
polynomial ring). 
\begin{Verbatim}[commandchars=!@|,fontsize=\small,frame=single,label=Example]
  !gapprompt@gap>| !gapinput@Resultant(x^4+y,y^4+x,1);|
  y^16+y
  !gapprompt@gap>| !gapinput@Resultant(x^4+y,y^4+x,2);|
  x^16+x
\end{Verbatim}
 }

 }

  
\section{\textcolor{Chapter }{Multivariate Polynomials}}\label{Multivariate Polynomials}
\logpage{[ 66, 7, 0 ]}
\hyperdef{L}{X85ABC4687DF05777}{}
{
  
\subsection{\textcolor{Chapter }{Value}}\logpage{[ 66, 7, 1 ]}
\hyperdef{L}{X7A70769C7F52CD2D}{}
{
\noindent\textcolor{FuncColor}{$\triangleright$\enspace\texttt{Value({\mdseries\slshape ratfun, indets, vals[, one]})\index{Value@\texttt{Value}!for rat. function, a list of indeterminates, a value (and a one)}
\label{Value:for rat. function, a list of indeterminates, a value (and a one)}
}\hfill{\scriptsize (operation)}}\\
\noindent\textcolor{FuncColor}{$\triangleright$\enspace\texttt{Value({\mdseries\slshape upol, value[, one]})\index{Value@\texttt{Value}!for a univariate rat. function, a value (and a one)}
\label{Value:for a univariate rat. function, a value (and a one)}
}\hfill{\scriptsize (operation)}}\\


 The first variant takes a rational function \mbox{\texttt{\mdseries\slshape ratfun}} and specializes the indeterminates given in \mbox{\texttt{\mdseries\slshape indets}} to the values given in \mbox{\texttt{\mdseries\slshape vals}}, replacing the $i$\texttt{\symbol{45}}th entry in \mbox{\texttt{\mdseries\slshape indets}} by the $i$\texttt{\symbol{45}}th entry in \mbox{\texttt{\mdseries\slshape vals}}. If this specialization results in a constant polynomial, an element of the
coefficient ring is returned. If the specialization would specialize the
denominator of \mbox{\texttt{\mdseries\slshape ratfun}} to zero, an error is raised. 

 A variation is the evaluation at elements of another ring $R$, for which a multiplication with elements of the coefficient ring of \mbox{\texttt{\mdseries\slshape ratfun}} are defined. In this situation the identity element of $R$ may be given by a further argument \mbox{\texttt{\mdseries\slshape one}} which will be used for $x^0$ for any specialized indeterminate $x$. 

 The second version takes an univariate rational function and specializes the
value of its indeterminate to \mbox{\texttt{\mdseries\slshape val}}. Again, an optional argument \mbox{\texttt{\mdseries\slshape one}} may be given. \texttt{Value( \mbox{\texttt{\mdseries\slshape upol}}, \mbox{\texttt{\mdseries\slshape val}} )} can also be expressed as \texttt{upol( \mbox{\texttt{\mdseries\slshape val}} )}. 

 
\begin{Verbatim}[commandchars=!@|,fontsize=\small,frame=single,label=Example]
  !gapprompt@gap>| !gapinput@Value(x*y+y+x^7,[x,y],[5,7]);|
  78167
\end{Verbatim}
 

 Note that the default values for \mbox{\texttt{\mdseries\slshape one}} can lead to different results than one would expect: For example for a matrix $M$, the values $M+M^0$ and $M+1$ are \emph{different}. As \texttt{Value} defaults to the one of the coefficient ring, when evaluating matrices in
polynomials always the correct \mbox{\texttt{\mdseries\slshape one}} should be given! }

 }

  
\section{\textcolor{Chapter }{Minimal Polynomials}}\label{Minimal Polynomials}
\logpage{[ 66, 8, 0 ]}
\hyperdef{L}{X7ED3E7D17C7AC732}{}
{
  \index{MinimalPolynomial@\texttt{MinimalPolynomial}!over a ring} 

\subsection{\textcolor{Chapter }{MinimalPolynomial}}
\logpage{[ 66, 8, 1 ]}\nobreak
\hyperdef{L}{X8643915A8424DAF8}{}
{\noindent\textcolor{FuncColor}{$\triangleright$\enspace\texttt{MinimalPolynomial({\mdseries\slshape R, elm[, ind]})\index{MinimalPolynomial@\texttt{MinimalPolynomial}}
\label{MinimalPolynomial}
}\hfill{\scriptsize (operation)}}\\


 returns the \emph{minimal polynomial} of \mbox{\texttt{\mdseries\slshape elm}} over the ring \mbox{\texttt{\mdseries\slshape R}}, expressed in the indeterminate number \mbox{\texttt{\mdseries\slshape ind}}. If \mbox{\texttt{\mdseries\slshape ind}} is not given, it defaults to 1. 

 The minimal polynomial is the monic polynomial of smallest degree with
coefficients in \mbox{\texttt{\mdseries\slshape R}} that has value zero at \mbox{\texttt{\mdseries\slshape elm}}. 
\begin{Verbatim}[commandchars=!@|,fontsize=\small,frame=single,label=Example]
  !gapprompt@gap>| !gapinput@MinimalPolynomial(Rationals,[[2,0],[0,2]]);|
  x-2
\end{Verbatim}
 }

 }

  
\section{\textcolor{Chapter }{Cyclotomic Polynomials}}\label{Cyclotomic Polynomials}
\logpage{[ 66, 9, 0 ]}
\hyperdef{L}{X837B8E55832CDFEB}{}
{
  

\subsection{\textcolor{Chapter }{CyclotomicPolynomial}}
\logpage{[ 66, 9, 1 ]}\nobreak
\hyperdef{L}{X827FC7FE81EE4C02}{}
{\noindent\textcolor{FuncColor}{$\triangleright$\enspace\texttt{CyclotomicPolynomial({\mdseries\slshape F, n})\index{CyclotomicPolynomial@\texttt{CyclotomicPolynomial}}
\label{CyclotomicPolynomial}
}\hfill{\scriptsize (function)}}\\


 is the \mbox{\texttt{\mdseries\slshape n}}\texttt{\symbol{45}}th cyclotomic polynomial over the ring \mbox{\texttt{\mdseries\slshape F}}. 
\begin{Verbatim}[commandchars=!@|,fontsize=\small,frame=single,label=Example]
  !gapprompt@gap>| !gapinput@CyclotomicPolynomial(Rationals,5);|
  x^4+x^3+x^2+x+1
\end{Verbatim}
 }

 }

  
\section{\textcolor{Chapter }{Polynomial Factorization}}\label{Polynomial Factorization}
\logpage{[ 66, 10, 0 ]}
\hyperdef{L}{X8551EF5187509D69}{}
{
  At the moment \textsf{GAP} provides only methods to factorize polynomials over finite fields (see
Chapter{\nobreakspace}\ref{Finite Fields}), over subfields of cyclotomic fields (see Chapter{\nobreakspace}\ref{Abelian Number Fields}), and over algebraic extensions of these (see Chapter{\nobreakspace}\ref{Algebraic extensions of fields}). 

 

\subsection{\textcolor{Chapter }{Factors (of polynomial)}}
\logpage{[ 66, 10, 1 ]}\nobreak
\hyperdef{L}{X83511D517B544F36}{}
{\noindent\textcolor{FuncColor}{$\triangleright$\enspace\texttt{Factors({\mdseries\slshape [R, ]poly[, opt]})\index{Factors@\texttt{Factors}!of polynomial}
\label{Factors:of polynomial}
}\hfill{\scriptsize (method)}}\\


 returns a list of the irreducible factors of the polynomial \mbox{\texttt{\mdseries\slshape poly}} in the polynomial ring \mbox{\texttt{\mdseries\slshape R}}. (That is factors over the \texttt{CoefficientsRing} (\ref{CoefficientsRing}) value of \mbox{\texttt{\mdseries\slshape R}}.) 

 For univariate factorizations, it is possible to pass a record \mbox{\texttt{\mdseries\slshape opt}} as a third argument. This record can contain the following components: 
\begin{description}
\item[{\texttt{onlydegs}}]  is a set of positive integers. The factorization assumes that all irreducible
factors have a degree in this set. 
\item[{\texttt{stopdegs}}]  is a set of positive integers. The factorization will stop once a factor of
degree in \texttt{stopdegs} has been found and will return the factorization found so far. 
\end{description}
 

 
\begin{Verbatim}[commandchars=!@|,fontsize=\small,frame=single,label=Example]
  !gapprompt@gap>| !gapinput@f:= CyclotomicPolynomial( GF(2), 7 );|
  x_1^6+x_1^5+x_1^4+x_1^3+x_1^2+x_1+Z(2)^0
  !gapprompt@gap>| !gapinput@Factors( f );|
  [ x_1^3+x_1+Z(2)^0, x_1^3+x_1^2+Z(2)^0 ]
  !gapprompt@gap>| !gapinput@Factors( PolynomialRing( GF(8) ), f );|
  [ x_1+Z(2^3), x_1+Z(2^3)^2, x_1+Z(2^3)^3, x_1+Z(2^3)^4, x_1+Z(2^3)^5,
    x_1+Z(2^3)^6 ]
  !gapprompt@gap>| !gapinput@f:= MinimalPolynomial( Rationals, E(4) );|
  x^2+1
  !gapprompt@gap>| !gapinput@Factors( f );|
  [ x^2+1 ]
  !gapprompt@gap>| !gapinput@Factors( PolynomialRing( Rationals ), f );|
  [ x^2+1 ]
  !gapprompt@gap>| !gapinput@Factors( PolynomialRing( CF(4) ), f );|
  [ x+(-E(4)), x+E(4) ]
\end{Verbatim}
 }

 

\subsection{\textcolor{Chapter }{FactorsSquarefree}}
\logpage{[ 66, 10, 2 ]}\nobreak
\hyperdef{L}{X7F5A4ACB7AF9E329}{}
{\noindent\textcolor{FuncColor}{$\triangleright$\enspace\texttt{FactorsSquarefree({\mdseries\slshape pring, upol, opt})\index{FactorsSquarefree@\texttt{FactorsSquarefree}}
\label{FactorsSquarefree}
}\hfill{\scriptsize (operation)}}\\


 returns a factorization of the squarefree, monic, univariate polynomial \mbox{\texttt{\mdseries\slshape upol}} in the polynomial ring \mbox{\texttt{\mdseries\slshape pring}}; \mbox{\texttt{\mdseries\slshape opt}} must be a (possibly empty) record of options. \mbox{\texttt{\mdseries\slshape upol}} must not have zero as a root. This function is used by the factoring
algorithms. 

 The current method for multivariate factorization reduces to univariate
factorization by use of a reduction homomorphism of the form $f(x_1,x_2,x_3) \mapsto f(x,x^p,x^{{p^2}})$. It can be very time intensive for larger degrees. 

 
\begin{Verbatim}[commandchars=!@|,fontsize=\small,frame=single,label=Example]
  !gapprompt@gap>| !gapinput@Factors(x^10-y^10);|
  [ x-y, x+y, x^4-x^3*y+x^2*y^2-x*y^3+y^4, x^4+x^3*y+x^2*y^2+x*y^3+y^4 ]
\end{Verbatim}
 }

 }

  
\section{\textcolor{Chapter }{Polynomials over the Rationals}}\label{Polynomials over the Rationals}
\logpage{[ 66, 11, 0 ]}
\hyperdef{L}{X7F45E9E47EA2C18B}{}
{
  The following functions are only available to polynomials with rational
coefficients: 

\subsection{\textcolor{Chapter }{PrimitivePolynomial}}
\logpage{[ 66, 11, 1 ]}\nobreak
\hyperdef{L}{X7E66494B7B05A055}{}
{\noindent\textcolor{FuncColor}{$\triangleright$\enspace\texttt{PrimitivePolynomial({\mdseries\slshape f})\index{PrimitivePolynomial@\texttt{PrimitivePolynomial}}
\label{PrimitivePolynomial}
}\hfill{\scriptsize (operation)}}\\


 takes a polynomial \mbox{\texttt{\mdseries\slshape f}} with rational coefficients and computes a new polynomial with integral
coefficients, obtained by multiplying with the Lcm of the denominators of the
coefficients and casting out the content (the Gcd of the coefficients). The
operation returns a list [\mbox{\texttt{\mdseries\slshape newpol}},\mbox{\texttt{\mdseries\slshape coeff}}] with rational \mbox{\texttt{\mdseries\slshape coeff}} such that \texttt{\mbox{\texttt{\mdseries\slshape coeff}}*\mbox{\texttt{\mdseries\slshape newpol}}=\mbox{\texttt{\mdseries\slshape f}}}. }

 

\subsection{\textcolor{Chapter }{PolynomialModP}}
\logpage{[ 66, 11, 2 ]}\nobreak
\hyperdef{L}{X7A73A3877EB73566}{}
{\noindent\textcolor{FuncColor}{$\triangleright$\enspace\texttt{PolynomialModP({\mdseries\slshape pol, p})\index{PolynomialModP@\texttt{PolynomialModP}}
\label{PolynomialModP}
}\hfill{\scriptsize (function)}}\\


 for a rational polynomial \mbox{\texttt{\mdseries\slshape pol}} this function returns a polynomial over the field with \mbox{\texttt{\mdseries\slshape p}} elements, obtained by reducing the coefficients modulo \mbox{\texttt{\mdseries\slshape p}}. }

 

\subsection{\textcolor{Chapter }{GaloisType}}
\logpage{[ 66, 11, 3 ]}\nobreak
\hyperdef{L}{X7AB9A6257ED694EC}{}
{\noindent\textcolor{FuncColor}{$\triangleright$\enspace\texttt{GaloisType({\mdseries\slshape f})\index{GaloisType@\texttt{GaloisType}}
\label{GaloisType}
}\hfill{\scriptsize (attribute)}}\\


 Let \mbox{\texttt{\mdseries\slshape f}} be an irreducible polynomial with rational coefficients. This function returns
the type of Gal(\mbox{\texttt{\mdseries\slshape f}}) (considered as a transitive permutation group of the roots of \mbox{\texttt{\mdseries\slshape f}}). It returns a number \mbox{\texttt{\mdseries\slshape i}} if Gal(\mbox{\texttt{\mdseries\slshape f}}) is permutation isomorphic to \texttt{TransitiveGroup(\mbox{\texttt{\mdseries\slshape n}},\mbox{\texttt{\mdseries\slshape i}})} where \mbox{\texttt{\mdseries\slshape n}} is the degree of \mbox{\texttt{\mdseries\slshape f}}. 

 Identification is performed by factoring appropriate Galois resolvents as
proposed in \cite{MS85}. This function is provided for rational polynomials of degree up to 15.
However, in some cases the required calculations become unfeasibly large. 

 For a few polynomials of degree 14, a complete discrimination is not yet
possible, as it would require computations, that are not feasible with current
factoring methods. 

 This function requires the transitive groups library to be installed (see  (\textbf{transgrp: Transitive Permutation Groups})). }

 

\subsection{\textcolor{Chapter }{ProbabilityShapes}}
\logpage{[ 66, 11, 4 ]}\nobreak
\hyperdef{L}{X7EB610D37D156DC6}{}
{\noindent\textcolor{FuncColor}{$\triangleright$\enspace\texttt{ProbabilityShapes({\mdseries\slshape f})\index{ProbabilityShapes@\texttt{ProbabilityShapes}}
\label{ProbabilityShapes}
}\hfill{\scriptsize (function)}}\\


 Let \mbox{\texttt{\mdseries\slshape f}} be an irreducible polynomial with rational coefficients. This function returns
a list of the most likely type(s) of Gal(\mbox{\texttt{\mdseries\slshape f}}) (see \texttt{GaloisType} (\ref{GaloisType})), based on factorization modulo a set of primes. It is very fast, but the
result is only probabilistic. 

 This function requires the transitive groups library to be installed (see  (\textbf{transgrp: Transitive Permutation Groups})). 
\begin{Verbatim}[commandchars=!@|,fontsize=\small,frame=single,label=Example]
  !gapprompt@gap>| !gapinput@f:=x^9-9*x^7+27*x^5-39*x^3+36*x-8;;|
  !gapprompt@gap>| !gapinput@GaloisType(f);|
  25
  !gapprompt@gap>| !gapinput@TransitiveGroup(9,25);|
  [1/2.S(3)^3]3
  !gapprompt@gap>| !gapinput@ProbabilityShapes(f);|
  [ 25 ]
\end{Verbatim}
 }

 }

  
\section{\textcolor{Chapter }{Factorization of Polynomials over the Rationals}}\label{Factorization of Polynomials over the Rationals}
\logpage{[ 66, 12, 0 ]}
\hyperdef{L}{X7C178AB9866FDDE5}{}
{
  The following operations are used by \textsf{GAP} inside the factorization algorithm but might be of interest also in other
contexts. 

\subsection{\textcolor{Chapter }{BombieriNorm}}
\logpage{[ 66, 12, 1 ]}\nobreak
\hyperdef{L}{X8723075C81D2CCA6}{}
{\noindent\textcolor{FuncColor}{$\triangleright$\enspace\texttt{BombieriNorm({\mdseries\slshape pol})\index{BombieriNorm@\texttt{BombieriNorm}}
\label{BombieriNorm}
}\hfill{\scriptsize (function)}}\\


 computes weighted Norm [\mbox{\texttt{\mdseries\slshape pol}}]$_2$ of \mbox{\texttt{\mdseries\slshape pol}} which is a good measure for factor coefficients (see \cite{BTW93}). }

 

\subsection{\textcolor{Chapter }{MinimizedBombieriNorm}}
\logpage{[ 66, 12, 2 ]}\nobreak
\hyperdef{L}{X856D769D878AF7AE}{}
{\noindent\textcolor{FuncColor}{$\triangleright$\enspace\texttt{MinimizedBombieriNorm({\mdseries\slshape f})\index{MinimizedBombieriNorm@\texttt{MinimizedBombieriNorm}}
\label{MinimizedBombieriNorm}
}\hfill{\scriptsize (attribute)}}\\


 This function applies linear Tschirnhaus transformations ($x \mapsto x + i$) to the polynomial \mbox{\texttt{\mdseries\slshape f}}, trying to get the Bombieri norm of \mbox{\texttt{\mdseries\slshape f}} small. It returns a list \texttt{[\mbox{\texttt{\mdseries\slshape new{\textunderscore}polynomial}}, \mbox{\texttt{\mdseries\slshape i{\textunderscore}of{\textunderscore}transformation}}]}. }

 

\subsection{\textcolor{Chapter }{HenselBound}}
\logpage{[ 66, 12, 3 ]}\nobreak
\hyperdef{L}{X8139BB0F87399F2C}{}
{\noindent\textcolor{FuncColor}{$\triangleright$\enspace\texttt{HenselBound({\mdseries\slshape pol[, minpol, den]})\index{HenselBound@\texttt{HenselBound}}
\label{HenselBound}
}\hfill{\scriptsize (function)}}\\


 returns the Hensel bound of the polynomial \mbox{\texttt{\mdseries\slshape pol}}. If the computation takes place over an algebraic extension, then the minimal
polynomial \mbox{\texttt{\mdseries\slshape minpol}} and denominator \mbox{\texttt{\mdseries\slshape den}} must be given. }

 

\subsection{\textcolor{Chapter }{OneFactorBound}}
\logpage{[ 66, 12, 4 ]}\nobreak
\hyperdef{L}{X79CC9C8D7C9F6B6A}{}
{\noindent\textcolor{FuncColor}{$\triangleright$\enspace\texttt{OneFactorBound({\mdseries\slshape pol})\index{OneFactorBound@\texttt{OneFactorBound}}
\label{OneFactorBound}
}\hfill{\scriptsize (function)}}\\


 returns the coefficient bound for a single factor of the rational polynomial \mbox{\texttt{\mdseries\slshape pol}}. }

 }

  
\section{\textcolor{Chapter }{Laurent Polynomials}}\label{Laurent Polynomials}
\logpage{[ 66, 13, 0 ]}
\hyperdef{L}{X844B3C6C87A0E7E0}{}
{
  A univariate polynomial can be written in the form $r_0 + r_1 x + \cdots + r_n x^n$, with $r_i \in R$. Formally, there is no reason to start with 0, if $m$ is an integer, we can consider objects of the form $r_m x^m + r_{{m+1}} x^{{m+1}} + \cdots + r_n x^n$. We call these \emph{Laurent polynomials}. Laurent polynomials also can be considered as quotients of a univariate
polynomial by a power of the indeterminate. The addition and multiplication of
univariate polynomials extends to Laurent polynomials (though it might be
impossible to interpret a Laurent polynomial as a function) and many functions
for univariate polynomials extend to Laurent polynomials (or extended versions
for Laurent polynomials exist). 

\subsection{\textcolor{Chapter }{LaurentPolynomialByCoefficients}}
\logpage{[ 66, 13, 1 ]}\nobreak
\hyperdef{L}{X8467263B7EFA013E}{}
{\noindent\textcolor{FuncColor}{$\triangleright$\enspace\texttt{LaurentPolynomialByCoefficients({\mdseries\slshape fam, cofs, val[, ind]})\index{LaurentPolynomialByCoefficients@\texttt{LaurentPolynomialByCoefficients}}
\label{LaurentPolynomialByCoefficients}
}\hfill{\scriptsize (operation)}}\\


 constructs a Laurent polynomial over the coefficients family \mbox{\texttt{\mdseries\slshape fam}} and in the indeterminate \mbox{\texttt{\mdseries\slshape ind}} (defaulting to 1) with the coefficients given by \mbox{\texttt{\mdseries\slshape cofs}} and valuation \mbox{\texttt{\mdseries\slshape val}}. }

 

\subsection{\textcolor{Chapter }{CoefficientsOfLaurentPolynomial}}
\logpage{[ 66, 13, 2 ]}\nobreak
\hyperdef{L}{X86D58AB67F86469F}{}
{\noindent\textcolor{FuncColor}{$\triangleright$\enspace\texttt{CoefficientsOfLaurentPolynomial({\mdseries\slshape laurent})\index{CoefficientsOfLaurentPolynomial@\texttt{CoefficientsOfLaurentPolynomial}}
\label{CoefficientsOfLaurentPolynomial}
}\hfill{\scriptsize (attribute)}}\\


 For a Laurent polynomial \mbox{\texttt{\mdseries\slshape laurent}}, this function returns a pair \texttt{[\mbox{\texttt{\mdseries\slshape cof}}, \mbox{\texttt{\mdseries\slshape val}}]}, consisting of the coefficient list (in ascending order) \mbox{\texttt{\mdseries\slshape cof}} and the valuation \mbox{\texttt{\mdseries\slshape val}} of \mbox{\texttt{\mdseries\slshape laurent}}. 
\begin{Verbatim}[commandchars=!@|,fontsize=\small,frame=single,label=Example]
  !gapprompt@gap>| !gapinput@p:=LaurentPolynomialByCoefficients(FamilyObj(1),|
  !gapprompt@>| !gapinput@[1,2,3,4,5],-2);|
  5*x^2+4*x+3+2*x^-1+x^-2
  !gapprompt@gap>| !gapinput@NumeratorOfRationalFunction(p);DenominatorOfRationalFunction(p);|
  5*x^4+4*x^3+3*x^2+2*x+1
  x^2
  !gapprompt@gap>| !gapinput@CoefficientsOfLaurentPolynomial(p*p);|
  [ [ 1, 4, 10, 20, 35, 44, 46, 40, 25 ], -4 ]
\end{Verbatim}
 }

 

\subsection{\textcolor{Chapter }{IndeterminateNumberOfLaurentPolynomial}}
\logpage{[ 66, 13, 3 ]}\nobreak
\hyperdef{L}{X8381E1B582F38C85}{}
{\noindent\textcolor{FuncColor}{$\triangleright$\enspace\texttt{IndeterminateNumberOfLaurentPolynomial({\mdseries\slshape pol})\index{IndeterminateNumberOfLaurentPolynomial@\texttt{Indeterminate}\-\texttt{Number}\-\texttt{Of}\-\texttt{Laurent}\-\texttt{Polynomial}}
\label{IndeterminateNumberOfLaurentPolynomial}
}\hfill{\scriptsize (attribute)}}\\


 Is a synonym for \texttt{IndeterminateNumberOfUnivariateRationalFunction} (\ref{IndeterminateNumberOfUnivariateRationalFunction}). }

 }

  
\section{\textcolor{Chapter }{Univariate Rational Functions}}\label{Univariate Rational Functions}
\logpage{[ 66, 14, 0 ]}
\hyperdef{L}{X7C1708D27F97B05F}{}
{
  

\subsection{\textcolor{Chapter }{UnivariateRationalFunctionByCoefficients}}
\logpage{[ 66, 14, 1 ]}\nobreak
\hyperdef{L}{X83DD411179888783}{}
{\noindent\textcolor{FuncColor}{$\triangleright$\enspace\texttt{UnivariateRationalFunctionByCoefficients({\mdseries\slshape fam, ncof, dcof, val[, ind]})\index{UnivariateRationalFunctionByCoefficients@\texttt{Univariate}\-\texttt{Rational}\-\texttt{Function}\-\texttt{By}\-\texttt{Coefficients}}
\label{UnivariateRationalFunctionByCoefficients}
}\hfill{\scriptsize (operation)}}\\


 constructs a univariate rational function over the coefficients family \mbox{\texttt{\mdseries\slshape fam}} and in the indeterminate \mbox{\texttt{\mdseries\slshape ind}} (defaulting to 1) with numerator and denominator coefficients given by \mbox{\texttt{\mdseries\slshape ncof}} and \mbox{\texttt{\mdseries\slshape dcof}} and valuation \mbox{\texttt{\mdseries\slshape val}}. }

 

\subsection{\textcolor{Chapter }{TaylorSeriesRationalFunction}}
\logpage{[ 66, 14, 2 ]}\nobreak
\hyperdef{L}{X7B1EE4E07A9631C6}{}
{\noindent\textcolor{FuncColor}{$\triangleright$\enspace\texttt{TaylorSeriesRationalFunction({\mdseries\slshape ratfun, at, deg]})\index{TaylorSeriesRationalFunction@\texttt{TaylorSeriesRationalFunction}}
\label{TaylorSeriesRationalFunction}
}\hfill{\scriptsize (attribute)}}\\


 Computes the taylor series up to degree \mbox{\texttt{\mdseries\slshape deg}} of \mbox{\texttt{\mdseries\slshape ratfun}} at \mbox{\texttt{\mdseries\slshape at}}. 
\begin{Verbatim}[commandchars=!@|,fontsize=\small,frame=single,label=Example]
  !gapprompt@gap>| !gapinput@TaylorSeriesRationalFunction((x^5+3*x+7)/(x^5+x+1),0,11);|
  -50*x^11+36*x^10-26*x^9+22*x^8-18*x^7+14*x^6-10*x^5+4*x^4-4*x^3+4*x^2-4*x+7
\end{Verbatim}
 }

 }

  
\section{\textcolor{Chapter }{Polynomial Rings and Function Fields}}\label{Polynomial Rings and Function Fields}
\logpage{[ 66, 15, 0 ]}
\hyperdef{L}{X7C59471783C3FEDC}{}
{
  While polynomials depend only on the family of the coefficients, polynomial
rings $A$ are defined over a base ring $R$. A polynomial is an element of $A$ if and only if all its coefficients are contained in $R$. Besides providing domains and an easy way to create polynomials, polynomial
rings can affect the behavior of operations like factorization into
irreducibles. 

 If you need to work with a polynomial ring and its indeterminates the
following two approaches will produce a ring that contains given variables
(see section{\nobreakspace}\ref{Indeterminates} for details about the internal numbering): Either, first create the ring and
then get the indeterminates with \texttt{IndeterminatesOfPolynomialRing} (\ref{IndeterminatesOfPolynomialRing}). 

 
\begin{Verbatim}[commandchars=!@|,fontsize=\small,frame=single,label=Example]
  !gapprompt@gap>| !gapinput@r := PolynomialRing(Rationals,["a","b"]);;|
  !gapprompt@gap>| !gapinput@indets := IndeterminatesOfPolynomialRing(r);;|
  !gapprompt@gap>| !gapinput@a := indets[1]; a := indets[2];|
  a
  b
\end{Verbatim}
 

 Alternatively, first create the indeterminates and then create the ring
including these indeterminates. 

 
\begin{Verbatim}[commandchars=!@|,fontsize=\small,frame=single,label=Example]
  !gapprompt@gap>| !gapinput@a:=Indeterminate(Rationals,"a":old);;|
  !gapprompt@gap>| !gapinput@b:=Indeterminate(Rationals,"b":old);;|
  !gapprompt@gap>| !gapinput@PolynomialRing(Rationals,[a,b]);;|
\end{Verbatim}
 

 As a convenient shortcut, intended mainly for interactive working, the $i$\texttt{\symbol{45}}th indeterminate of a polynomial ring $R$ can be accessed as $R.i$, which corresponds exactly to \texttt{IndeterminatesOfPolynomialRing}$( R )[i]$ or, if it has the name \texttt{nam}, as $R$\texttt{.nam}. \emph{Note} that the number $i$ is in general \emph{not} the indeterminate number, but simply an index into the indeterminates list of $R$. 

 
\begin{Verbatim}[commandchars=!@|,fontsize=\small,frame=single,label=Example]
  !gapprompt@gap>| !gapinput@r := PolynomialRing(Rationals, ["a", "b"]:old );;|
  !gapprompt@gap>| !gapinput@r.1; r.2; r.a; r.b;|
  a
  b
  a
  b
  !gapprompt@gap>| !gapinput@IndeterminateNumberOfLaurentPolynomial(r.1);|
  3
\end{Verbatim}
 

 Polynomials as \textsf{GAP} objects can exist without a polynomial ring being defined and polynomials
cannot be associated to a particular polynomial ring. (For example dividing a
polynomial which is in a polynomial ring over the integers by another integer
will result in a polynomial over the rationals, not in a rational function
over the integers.) 

 
\subsection{\textcolor{Chapter }{PolynomialRing}}\logpage{[ 66, 15, 1 ]}
\hyperdef{L}{X7D2F16E480060330}{}
{
\noindent\textcolor{FuncColor}{$\triangleright$\enspace\texttt{PolynomialRing({\mdseries\slshape R, rank[, avoid]})\index{PolynomialRing@\texttt{PolynomialRing}!for a ring and a rank (and an exclusion list)}
\label{PolynomialRing:for a ring and a rank (and an exclusion list)}
}\hfill{\scriptsize (operation)}}\\
\noindent\textcolor{FuncColor}{$\triangleright$\enspace\texttt{PolynomialRing({\mdseries\slshape R, names[, avoid]})\index{PolynomialRing@\texttt{PolynomialRing}!for a ring and a list of names (and an exclusion list)}
\label{PolynomialRing:for a ring and a list of names (and an exclusion list)}
}\hfill{\scriptsize (operation)}}\\
\noindent\textcolor{FuncColor}{$\triangleright$\enspace\texttt{PolynomialRing({\mdseries\slshape R, indets})\index{PolynomialRing@\texttt{PolynomialRing}!for a ring and a list of indeterminates}
\label{PolynomialRing:for a ring and a list of indeterminates}
}\hfill{\scriptsize (operation)}}\\
\noindent\textcolor{FuncColor}{$\triangleright$\enspace\texttt{PolynomialRing({\mdseries\slshape R, indetnums})\index{PolynomialRing@\texttt{PolynomialRing}!for a ring and a list of indeterminate numbers}
\label{PolynomialRing:for a ring and a list of indeterminate numbers}
}\hfill{\scriptsize (operation)}}\\


 creates a polynomial ring over the ring \mbox{\texttt{\mdseries\slshape R}}. If a positive integer \mbox{\texttt{\mdseries\slshape rank}} is given, this creates the polynomial ring in \mbox{\texttt{\mdseries\slshape rank}} indeterminates. These indeterminates will have the internal index numbers 1 to \mbox{\texttt{\mdseries\slshape rank}}. The second usage takes a list \mbox{\texttt{\mdseries\slshape names}} of strings and returns a polynomial ring in indeterminates labelled by \mbox{\texttt{\mdseries\slshape names}}. These indeterminates have ``new'' internal index numbers as if they had been created by calls to \texttt{Indeterminate} (\ref{Indeterminate:for a ring (and a number)}). (If the argument \mbox{\texttt{\mdseries\slshape avoid}} is given it contains indeterminates that should be avoided, in this case
internal index numbers are incremented to skip these variables.) In the third
version, a list of indeterminates \mbox{\texttt{\mdseries\slshape indets}} is given. This creates the polynomial ring in the indeterminates \mbox{\texttt{\mdseries\slshape indets}}. Finally, the fourth version specifies indeterminates by their index numbers. 

 To get the indeterminates of a polynomial ring use \texttt{IndeterminatesOfPolynomialRing} (\ref{IndeterminatesOfPolynomialRing}). (Indeterminates created independently with \texttt{Indeterminate} (\ref{Indeterminate:for a ring (and a number)}) will usually differ, though they might be given the same name and display
identically, see Section{\nobreakspace}\ref{Indeterminates}.) }

 

\subsection{\textcolor{Chapter }{IndeterminatesOfPolynomialRing}}
\logpage{[ 66, 15, 2 ]}\nobreak
\hyperdef{L}{X80D585E1793D4552}{}
{\noindent\textcolor{FuncColor}{$\triangleright$\enspace\texttt{IndeterminatesOfPolynomialRing({\mdseries\slshape pring})\index{IndeterminatesOfPolynomialRing@\texttt{IndeterminatesOfPolynomialRing}}
\label{IndeterminatesOfPolynomialRing}
}\hfill{\scriptsize (attribute)}}\\
\noindent\textcolor{FuncColor}{$\triangleright$\enspace\texttt{IndeterminatesOfFunctionField({\mdseries\slshape ffield})\index{IndeterminatesOfFunctionField@\texttt{IndeterminatesOfFunctionField}}
\label{IndeterminatesOfFunctionField}
}\hfill{\scriptsize (attribute)}}\\


 returns a list of the indeterminates of the polynomial ring \mbox{\texttt{\mdseries\slshape pring}}, respectively the function field \mbox{\texttt{\mdseries\slshape ffield}}. }

 

\subsection{\textcolor{Chapter }{CoefficientsRing}}
\logpage{[ 66, 15, 3 ]}\nobreak
\hyperdef{L}{X8235D10781BE8003}{}
{\noindent\textcolor{FuncColor}{$\triangleright$\enspace\texttt{CoefficientsRing({\mdseries\slshape pring})\index{CoefficientsRing@\texttt{CoefficientsRing}}
\label{CoefficientsRing}
}\hfill{\scriptsize (attribute)}}\\


 returns the ring of coefficients of the polynomial ring \mbox{\texttt{\mdseries\slshape pring}}, that is the ring over which \mbox{\texttt{\mdseries\slshape pring}} was defined. 
\begin{Verbatim}[commandchars=!@|,fontsize=\small,frame=single,label=Example]
  !gapprompt@gap>| !gapinput@r:=PolynomialRing(GF(7));|
  GF(7)[x_1]
  !gapprompt@gap>| !gapinput@r:=PolynomialRing(GF(7),3);|
  GF(7)[x_1,x_2,x_3]
  !gapprompt@gap>| !gapinput@IndeterminatesOfPolynomialRing(r);|
  [ x_1, x_2, x_3 ]
  !gapprompt@gap>| !gapinput@r2:=PolynomialRing(GF(7),[5,7,12]);|
  GF(7)[x_5,x_7,x_12]
  !gapprompt@gap>| !gapinput@CoefficientsRing(r);|
  GF(7)
  !gapprompt@gap>| !gapinput@r:=PolynomialRing(GF(7),3);|
  GF(7)[x_1,x_2,x_3]
  !gapprompt@gap>| !gapinput@r2:=PolynomialRing(GF(7),3,IndeterminatesOfPolynomialRing(r));|
  GF(7)[x_4,x_5,x_6]
  !gapprompt@gap>| !gapinput@r:=PolynomialRing(GF(7),["x","y","z","z2"]);|
  GF(7)[x,y,z,z2]
\end{Verbatim}
 }

 

\subsection{\textcolor{Chapter }{IsPolynomialRing}}
\logpage{[ 66, 15, 4 ]}\nobreak
\hyperdef{L}{X7D631ACC86C584B7}{}
{\noindent\textcolor{FuncColor}{$\triangleright$\enspace\texttt{IsPolynomialRing({\mdseries\slshape pring})\index{IsPolynomialRing@\texttt{IsPolynomialRing}}
\label{IsPolynomialRing}
}\hfill{\scriptsize (Category)}}\\


 is the category of polynomial rings }

 

\subsection{\textcolor{Chapter }{IsFiniteFieldPolynomialRing}}
\logpage{[ 66, 15, 5 ]}\nobreak
\hyperdef{L}{X86F391237A76D804}{}
{\noindent\textcolor{FuncColor}{$\triangleright$\enspace\texttt{IsFiniteFieldPolynomialRing({\mdseries\slshape pring})\index{IsFiniteFieldPolynomialRing@\texttt{IsFiniteFieldPolynomialRing}}
\label{IsFiniteFieldPolynomialRing}
}\hfill{\scriptsize (Category)}}\\


 is the category of polynomial rings over a finite field (see
Chapter{\nobreakspace}\ref{Finite Fields}). }

 

\subsection{\textcolor{Chapter }{IsAbelianNumberFieldPolynomialRing}}
\logpage{[ 66, 15, 6 ]}\nobreak
\hyperdef{L}{X782D07F77BCF67C1}{}
{\noindent\textcolor{FuncColor}{$\triangleright$\enspace\texttt{IsAbelianNumberFieldPolynomialRing({\mdseries\slshape pring})\index{IsAbelianNumberFieldPolynomialRing@\texttt{IsAbelianNumberFieldPolynomialRing}}
\label{IsAbelianNumberFieldPolynomialRing}
}\hfill{\scriptsize (Category)}}\\


 is the category of polynomial rings over a field of cyclotomics (see the
chapters{\nobreakspace}\ref{Cyclotomic Numbers} and \ref{Abelian Number Fields}). }

 

\subsection{\textcolor{Chapter }{IsRationalsPolynomialRing}}
\logpage{[ 66, 15, 7 ]}\nobreak
\hyperdef{L}{X7D45213A8642033B}{}
{\noindent\textcolor{FuncColor}{$\triangleright$\enspace\texttt{IsRationalsPolynomialRing({\mdseries\slshape pring})\index{IsRationalsPolynomialRing@\texttt{IsRationalsPolynomialRing}}
\label{IsRationalsPolynomialRing}
}\hfill{\scriptsize (Category)}}\\


 is the category of polynomial rings over the rationals (see
Chapter{\nobreakspace}\ref{Rational Numbers}). 
\begin{Verbatim}[commandchars=!@|,fontsize=\small,frame=single,label=Example]
  !gapprompt@gap>| !gapinput@r := PolynomialRing(Rationals, ["a", "b"] );;|
  !gapprompt@gap>| !gapinput@IsPolynomialRing(r);|
  true
  !gapprompt@gap>| !gapinput@IsFiniteFieldPolynomialRing(r);|
  false
  !gapprompt@gap>| !gapinput@IsRationalsPolynomialRing(r);|
  true
\end{Verbatim}
 }

 
\subsection{\textcolor{Chapter }{FunctionField}}\logpage{[ 66, 15, 8 ]}
\hyperdef{L}{X812E801484E3624E}{}
{
\noindent\textcolor{FuncColor}{$\triangleright$\enspace\texttt{FunctionField({\mdseries\slshape R, rank[, avoid]})\index{FunctionField@\texttt{FunctionField}!for an integral ring and a rank (and an exclusion list)}
\label{FunctionField:for an integral ring and a rank (and an exclusion list)}
}\hfill{\scriptsize (operation)}}\\
\noindent\textcolor{FuncColor}{$\triangleright$\enspace\texttt{FunctionField({\mdseries\slshape R, names[, avoid]})\index{FunctionField@\texttt{FunctionField}!for an integral ring and a list of names (and an exclusion list)}
\label{FunctionField:for an integral ring and a list of names (and an exclusion list)}
}\hfill{\scriptsize (operation)}}\\
\noindent\textcolor{FuncColor}{$\triangleright$\enspace\texttt{FunctionField({\mdseries\slshape R, indets})\index{FunctionField@\texttt{FunctionField}!for an integral ring and a list of indeterminates}
\label{FunctionField:for an integral ring and a list of indeterminates}
}\hfill{\scriptsize (operation)}}\\
\noindent\textcolor{FuncColor}{$\triangleright$\enspace\texttt{FunctionField({\mdseries\slshape R, indetnums})\index{FunctionField@\texttt{FunctionField}!for an integral ring and a list of indeterminate numbers}
\label{FunctionField:for an integral ring and a list of indeterminate numbers}
}\hfill{\scriptsize (operation)}}\\


 creates a function field over the integral ring \mbox{\texttt{\mdseries\slshape R}}. If a positive integer \mbox{\texttt{\mdseries\slshape rank}} is given, this creates the function field in \mbox{\texttt{\mdseries\slshape rank}} indeterminates. These indeterminates will have the internal index numbers 1 to \mbox{\texttt{\mdseries\slshape rank}}. The second usage takes a list \mbox{\texttt{\mdseries\slshape names}} of strings and returns a function field in indeterminates labelled by \mbox{\texttt{\mdseries\slshape names}}. These indeterminates have ``new'' internal index numbers as if they had been created by calls to \texttt{Indeterminate} (\ref{Indeterminate:for a ring (and a number)}). (If the argument \mbox{\texttt{\mdseries\slshape avoid}} is given it contains indeterminates that should be avoided, in this case
internal index numbers are incremented to skip these variables.) In the third
version, a list of indeterminates \mbox{\texttt{\mdseries\slshape indets}} is given. This creates the function field in the indeterminates \mbox{\texttt{\mdseries\slshape indets}}. Finally, the fourth version specifies indeterminates by their index number. 

 To get the indeterminates of a function field use \texttt{IndeterminatesOfFunctionField} (\ref{IndeterminatesOfFunctionField}). (Indeterminates created independently with \texttt{Indeterminate} (\ref{Indeterminate:for a ring (and a number)}) will usually differ, though they might be given the same name and display
identically, see Section{\nobreakspace}\ref{Indeterminates}.) }

 

\subsection{\textcolor{Chapter }{IsFunctionField}}
\logpage{[ 66, 15, 9 ]}\nobreak
\hyperdef{L}{X8090C9EC85201AAC}{}
{\noindent\textcolor{FuncColor}{$\triangleright$\enspace\texttt{IsFunctionField({\mdseries\slshape ffield})\index{IsFunctionField@\texttt{IsFunctionField}}
\label{IsFunctionField}
}\hfill{\scriptsize (Category)}}\\


 is the category of function fields }

 }

  
\section{\textcolor{Chapter }{Univariate Polynomial Rings}}\label{Univariate Polynomial Rings}
\logpage{[ 66, 16, 0 ]}
\hyperdef{L}{X85CA757B844F12AE}{}
{
  
\subsection{\textcolor{Chapter }{UnivariatePolynomialRing}}\logpage{[ 66, 16, 1 ]}
\hyperdef{L}{X84DC2A59797A26DE}{}
{
\noindent\textcolor{FuncColor}{$\triangleright$\enspace\texttt{UnivariatePolynomialRing({\mdseries\slshape R[, nr]})\index{UnivariatePolynomialRing@\texttt{UnivariatePolynomialRing}!for a ring (and an indeterminate number)}
\label{UnivariatePolynomialRing:for a ring (and an indeterminate number)}
}\hfill{\scriptsize (operation)}}\\
\noindent\textcolor{FuncColor}{$\triangleright$\enspace\texttt{UnivariatePolynomialRing({\mdseries\slshape R[, name][, avoid]})\index{UnivariatePolynomialRing@\texttt{UnivariatePolynomialRing}!for a ring (and a name and an exclusion list)}
\label{UnivariatePolynomialRing:for a ring (and a name and an exclusion list)}
}\hfill{\scriptsize (operation)}}\\


 returns a univariate polynomial ring in the indeterminate \mbox{\texttt{\mdseries\slshape nr}} over the base ring \mbox{\texttt{\mdseries\slshape R}}. If \mbox{\texttt{\mdseries\slshape nr}} is not given it defaults to 1. 

 If the number is not specified a list \mbox{\texttt{\mdseries\slshape avoid}} of indeterminates may be given. Then the function will return a ring in an
indeterminate that is guaranteed to be different from all the indeterminates
in \mbox{\texttt{\mdseries\slshape avoid}}. 

 Also a string \mbox{\texttt{\mdseries\slshape name}} can be prescribed as the name of the indeterminate chosen (also avoiding the
indeterminates in the list \mbox{\texttt{\mdseries\slshape avoid}} if given). }

 

\subsection{\textcolor{Chapter }{IsUnivariatePolynomialRing}}
\logpage{[ 66, 16, 2 ]}\nobreak
\hyperdef{L}{X7A43D74B812401CA}{}
{\noindent\textcolor{FuncColor}{$\triangleright$\enspace\texttt{IsUnivariatePolynomialRing({\mdseries\slshape pring})\index{IsUnivariatePolynomialRing@\texttt{IsUnivariatePolynomialRing}}
\label{IsUnivariatePolynomialRing}
}\hfill{\scriptsize (Category)}}\\


 is the category of polynomial rings with one indeterminate. 
\begin{Verbatim}[commandchars=!@|,fontsize=\small,frame=single,label=Example]
  !gapprompt@gap>| !gapinput@r:=UnivariatePolynomialRing(Rationals,"p");|
  Rationals[p]
  !gapprompt@gap>| !gapinput@r2:=PolynomialRing(Rationals,["q"]);|
  Rationals[q]
  !gapprompt@gap>| !gapinput@IsUnivariatePolynomialRing(r);|
  true
  !gapprompt@gap>| !gapinput@IsUnivariatePolynomialRing(r2);|
  true
\end{Verbatim}
 }

 }

  
\section{\textcolor{Chapter }{Monomial Orderings}}\label{Monomial Orderings}
\logpage{[ 66, 17, 0 ]}
\hyperdef{L}{X86E2ADEA784AD163}{}
{
  It is often desirable to consider the monomials within a polynomial to be
arranged with respect to a certain ordering. Such an ordering is called a \emph{monomial ordering} if it is total, invariant under multiplication with other monomials and admits
no infinite descending chains. For details on monomial orderings see \cite{coxlittleoshea}. 

 In \textsf{GAP}, monomial orderings are represented by objects that provide a way to compare
monomials (as polynomials as well as {\textendash}for efficiency purposes
within algorithms{\textendash} in the internal representation as lists). 

 Normally the ordering chosen should be \emph{admissible}, i.e. it must be compatible with products: If $a < b$ then $ca < cb$ for all monomials $a, b$ and $c$. 

 Each monomial ordering provides the two functions \texttt{MonomialComparisonFunction} (\ref{MonomialComparisonFunction}) and \texttt{MonomialExtrepComparisonFun} (\ref{MonomialExtrepComparisonFun}) to compare monomials. These functions work as ``is less than'', i.e. they return \texttt{true} if and only if the left argument is smaller. 

 

\subsection{\textcolor{Chapter }{IsMonomialOrdering}}
\logpage{[ 66, 17, 1 ]}\nobreak
\hyperdef{L}{X79D4CBBF820EA204}{}
{\noindent\textcolor{FuncColor}{$\triangleright$\enspace\texttt{IsMonomialOrdering({\mdseries\slshape obj})\index{IsMonomialOrdering@\texttt{IsMonomialOrdering}}
\label{IsMonomialOrdering}
}\hfill{\scriptsize (Category)}}\\


 A monomial ordering is an object representing a monomial ordering. Its
attributes \texttt{MonomialComparisonFunction} (\ref{MonomialComparisonFunction}) and \texttt{MonomialExtrepComparisonFun} (\ref{MonomialExtrepComparisonFun}) are actual comparison functions. }

 

\subsection{\textcolor{Chapter }{LeadingMonomialOfPolynomial}}
\logpage{[ 66, 17, 2 ]}\nobreak
\hyperdef{L}{X7D052A017A73E91E}{}
{\noindent\textcolor{FuncColor}{$\triangleright$\enspace\texttt{LeadingMonomialOfPolynomial({\mdseries\slshape pol, ord})\index{LeadingMonomialOfPolynomial@\texttt{LeadingMonomialOfPolynomial}}
\label{LeadingMonomialOfPolynomial}
}\hfill{\scriptsize (operation)}}\\


 returns the leading monomial (with respect to the ordering \mbox{\texttt{\mdseries\slshape ord}}) of the polynomial \mbox{\texttt{\mdseries\slshape pol}}. 
\begin{Verbatim}[commandchars=!@|,fontsize=\small,frame=single,label=Example]
  !gapprompt@gap>| !gapinput@x:=Indeterminate(Rationals,"x");;|
  !gapprompt@gap>| !gapinput@y:=Indeterminate(Rationals,"y");;|
  !gapprompt@gap>| !gapinput@z:=Indeterminate(Rationals,"z");;|
  !gapprompt@gap>| !gapinput@lexord:=MonomialLexOrdering();grlexord:=MonomialGrlexOrdering();|
  MonomialLexOrdering()
  MonomialGrlexOrdering()
  !gapprompt@gap>| !gapinput@f:=2*x+3*y+4*z+5*x^2-6*z^2+7*y^3;|
  7*y^3+5*x^2-6*z^2+2*x+3*y+4*z
  !gapprompt@gap>| !gapinput@LeadingMonomialOfPolynomial(f,lexord);|
  x^2
  !gapprompt@gap>| !gapinput@LeadingMonomialOfPolynomial(f,grlexord);|
  y^3
\end{Verbatim}
 }

 

\subsection{\textcolor{Chapter }{LeadingTermOfPolynomial}}
\logpage{[ 66, 17, 3 ]}\nobreak
\hyperdef{L}{X7B6231137BA8B95F}{}
{\noindent\textcolor{FuncColor}{$\triangleright$\enspace\texttt{LeadingTermOfPolynomial({\mdseries\slshape pol, ord})\index{LeadingTermOfPolynomial@\texttt{LeadingTermOfPolynomial}}
\label{LeadingTermOfPolynomial}
}\hfill{\scriptsize (operation)}}\\


 returns the leading term (with respect to the ordering \mbox{\texttt{\mdseries\slshape ord}}) of the polynomial \mbox{\texttt{\mdseries\slshape pol}}, i.e. the product of leading coefficient and leading monomial. }

 

\subsection{\textcolor{Chapter }{LeadingCoefficientOfPolynomial}}
\logpage{[ 66, 17, 4 ]}\nobreak
\hyperdef{L}{X798E707D86141087}{}
{\noindent\textcolor{FuncColor}{$\triangleright$\enspace\texttt{LeadingCoefficientOfPolynomial({\mdseries\slshape pol, ord})\index{LeadingCoefficientOfPolynomial@\texttt{LeadingCoefficientOfPolynomial}}
\label{LeadingCoefficientOfPolynomial}
}\hfill{\scriptsize (operation)}}\\


 returns the leading coefficient (that is the coefficient of the leading
monomial, see{\nobreakspace}\texttt{LeadingMonomialOfPolynomial} (\ref{LeadingMonomialOfPolynomial})) of the polynomial \mbox{\texttt{\mdseries\slshape pol}}. 
\begin{Verbatim}[commandchars=!@|,fontsize=\small,frame=single,label=Example]
  !gapprompt@gap>| !gapinput@LeadingTermOfPolynomial(f,lexord);|
  5*x^2
  !gapprompt@gap>| !gapinput@LeadingTermOfPolynomial(f,grlexord);|
  7*y^3
  !gapprompt@gap>| !gapinput@LeadingCoefficientOfPolynomial(f,lexord);|
  5
\end{Verbatim}
 }

 

\subsection{\textcolor{Chapter }{MonomialComparisonFunction}}
\logpage{[ 66, 17, 5 ]}\nobreak
\hyperdef{L}{X7EDE941781BA7F8B}{}
{\noindent\textcolor{FuncColor}{$\triangleright$\enspace\texttt{MonomialComparisonFunction({\mdseries\slshape O})\index{MonomialComparisonFunction@\texttt{MonomialComparisonFunction}}
\label{MonomialComparisonFunction}
}\hfill{\scriptsize (attribute)}}\\


 If \mbox{\texttt{\mdseries\slshape O}} is an object representing a monomial ordering, this attribute returns a \emph{function} that can be used to compare or sort monomials (and polynomials which will be
compared by their monomials in decreasing order) in this order. 
\begin{Verbatim}[commandchars=!@|,fontsize=\small,frame=single,label=Example]
  !gapprompt@gap>| !gapinput@MonomialComparisonFunction(lexord);|
  function( a, b ) ... end
  !gapprompt@gap>| !gapinput@l:=[f,Derivative(f,x),Derivative(f,y),Derivative(f,z)];;|
  !gapprompt@gap>| !gapinput@Sort(l,MonomialComparisonFunction(lexord));l;|
  [ -12*z+4, 21*y^2+3, 10*x+2, 7*y^3+5*x^2-6*z^2+2*x+3*y+4*z ]
\end{Verbatim}
 }

 

\subsection{\textcolor{Chapter }{MonomialExtrepComparisonFun}}
\logpage{[ 66, 17, 6 ]}\nobreak
\hyperdef{L}{X7EDC3A457E7B591E}{}
{\noindent\textcolor{FuncColor}{$\triangleright$\enspace\texttt{MonomialExtrepComparisonFun({\mdseries\slshape O})\index{MonomialExtrepComparisonFun@\texttt{MonomialExtrepComparisonFun}}
\label{MonomialExtrepComparisonFun}
}\hfill{\scriptsize (attribute)}}\\


 If \mbox{\texttt{\mdseries\slshape O}} is an object representing a monomial ordering, this attribute returns a \emph{function} that can be used to compare or sort monomials \emph{in their external representation} (as lists). This comparison variant is used inside algorithms that manipulate
the external representation. }

 

\subsection{\textcolor{Chapter }{MonomialLexOrdering}}
\logpage{[ 66, 17, 7 ]}\nobreak
\hyperdef{L}{X852D7BB37ECE98E1}{}
{\noindent\textcolor{FuncColor}{$\triangleright$\enspace\texttt{MonomialLexOrdering({\mdseries\slshape [vari]})\index{MonomialLexOrdering@\texttt{MonomialLexOrdering}}
\label{MonomialLexOrdering}
}\hfill{\scriptsize (function)}}\\


 This function creates a lexicographic ordering for monomials. Monomials are
compared first by the exponents of the largest variable, then the exponents of
the second largest variable and so on. 

 The variables are ordered according to their (internal) index, i.e., $x_1$ is larger than $x_2$ and so on. If \mbox{\texttt{\mdseries\slshape vari}} is given, and is a list of variables or variable indices, instead this
arrangement of variables (in descending order; i.e. the first variable is
larger than the second) is used as the underlying order of variables. 
\begin{Verbatim}[commandchars=!@|,fontsize=\small,frame=single,label=Example]
  !gapprompt@gap>| !gapinput@l:=List(Tuples([1..3],3),i->x^(i[1]-1)*y^(i[2]-1)*z^(i[3]-1));|
  [ 1, z, z^2, y, y*z, y*z^2, y^2, y^2*z, y^2*z^2, x, x*z, x*z^2, x*y,
    x*y*z, x*y*z^2, x*y^2, x*y^2*z, x*y^2*z^2, x^2, x^2*z, x^2*z^2,
    x^2*y, x^2*y*z, x^2*y*z^2, x^2*y^2, x^2*y^2*z, x^2*y^2*z^2 ]
  !gapprompt@gap>| !gapinput@Sort(l,MonomialComparisonFunction(MonomialLexOrdering()));l;|
  [ 1, z, z^2, y, y*z, y*z^2, y^2, y^2*z, y^2*z^2, x, x*z, x*z^2, x*y,
    x*y*z, x*y*z^2, x*y^2, x*y^2*z, x*y^2*z^2, x^2, x^2*z, x^2*z^2,
    x^2*y, x^2*y*z, x^2*y*z^2, x^2*y^2, x^2*y^2*z, x^2*y^2*z^2 ]
  !gapprompt@gap>| !gapinput@Sort(l,MonomialComparisonFunction(MonomialLexOrdering([y,z,x])));l;|
  [ 1, x, x^2, z, x*z, x^2*z, z^2, x*z^2, x^2*z^2, y, x*y, x^2*y, y*z,
    x*y*z, x^2*y*z, y*z^2, x*y*z^2, x^2*y*z^2, y^2, x*y^2, x^2*y^2,
    y^2*z, x*y^2*z, x^2*y^2*z, y^2*z^2, x*y^2*z^2, x^2*y^2*z^2 ]
  !gapprompt@gap>| !gapinput@Sort(l,MonomialComparisonFunction(MonomialLexOrdering([z,x,y])));l;|
  [ 1, y, y^2, x, x*y, x*y^2, x^2, x^2*y, x^2*y^2, z, y*z, y^2*z, x*z,
    x*y*z, x*y^2*z, x^2*z, x^2*y*z, x^2*y^2*z, z^2, y*z^2, y^2*z^2,
    x*z^2, x*y*z^2, x*y^2*z^2, x^2*z^2, x^2*y*z^2, x^2*y^2*z^2 ]
\end{Verbatim}
 }

 

\subsection{\textcolor{Chapter }{MonomialGrlexOrdering}}
\logpage{[ 66, 17, 8 ]}\nobreak
\hyperdef{L}{X786C866C824D2688}{}
{\noindent\textcolor{FuncColor}{$\triangleright$\enspace\texttt{MonomialGrlexOrdering({\mdseries\slshape [vari]})\index{MonomialGrlexOrdering@\texttt{MonomialGrlexOrdering}}
\label{MonomialGrlexOrdering}
}\hfill{\scriptsize (function)}}\\


 This function creates a degree/lexicographic ordering. In this ordering
monomials are compared first by their total degree, then lexicographically
(see \texttt{MonomialLexOrdering} (\ref{MonomialLexOrdering})). 

 The variables are ordered according to their (internal) index, i.e., $x_1$ is larger than $x_2$ and so on. If \mbox{\texttt{\mdseries\slshape vari}} is given, and is a list of variables or variable indices, instead this
arrangement of variables (in descending order; i.e. the first variable is
larger than the second) is used as the underlying order of variables. }

 

\subsection{\textcolor{Chapter }{MonomialGrevlexOrdering}}
\logpage{[ 66, 17, 9 ]}\nobreak
\hyperdef{L}{X8094C733808D1799}{}
{\noindent\textcolor{FuncColor}{$\triangleright$\enspace\texttt{MonomialGrevlexOrdering({\mdseries\slshape [vari]})\index{MonomialGrevlexOrdering@\texttt{MonomialGrevlexOrdering}}
\label{MonomialGrevlexOrdering}
}\hfill{\scriptsize (function)}}\\


 This function creates a ``grevlex'' ordering. In this ordering monomials are compared first by total degree and
then backwards lexicographically. (This is different than ``grlex'' ordering with variables reversed.) 

 The variables are ordered according to their (internal) index, i.e., $x_1$ is larger than $x_2$ and so on. If \mbox{\texttt{\mdseries\slshape vari}} is given, and is a list of variables or variable indices, instead this
arrangement of variables (in descending order; i.e. the first variable is
larger than the second) is used as the underlying order of variables. 
\begin{Verbatim}[commandchars=!@|,fontsize=\small,frame=single,label=Example]
  !gapprompt@gap>| !gapinput@Sort(l,MonomialComparisonFunction(MonomialGrlexOrdering()));l;|
  [ 1, z, y, x, z^2, y*z, y^2, x*z, x*y, x^2, y*z^2, y^2*z, x*z^2,
    x*y*z, x*y^2, x^2*z, x^2*y, y^2*z^2, x*y*z^2, x*y^2*z, x^2*z^2,
    x^2*y*z, x^2*y^2, x*y^2*z^2, x^2*y*z^2, x^2*y^2*z, x^2*y^2*z^2 ]
  !gapprompt@gap>| !gapinput@Sort(l,MonomialComparisonFunction(MonomialGrevlexOrdering()));l;|
  [ 1, z, y, x, z^2, y*z, x*z, y^2, x*y, x^2, y*z^2, x*z^2, y^2*z,
    x*y*z, x^2*z, x*y^2, x^2*y, y^2*z^2, x*y*z^2, x^2*z^2, x*y^2*z,
    x^2*y*z, x^2*y^2, x*y^2*z^2, x^2*y*z^2, x^2*y^2*z, x^2*y^2*z^2 ]
  !gapprompt@gap>| !gapinput@Sort(l,MonomialComparisonFunction(MonomialGrlexOrdering([z,y,x])));l;|
  [ 1, x, y, z, x^2, x*y, y^2, x*z, y*z, z^2, x^2*y, x*y^2, x^2*z,
    x*y*z, y^2*z, x*z^2, y*z^2, x^2*y^2, x^2*y*z, x*y^2*z, x^2*z^2,
    x*y*z^2, y^2*z^2, x^2*y^2*z, x^2*y*z^2, x*y^2*z^2, x^2*y^2*z^2 ]
\end{Verbatim}
 }

 

\subsection{\textcolor{Chapter }{EliminationOrdering}}
\logpage{[ 66, 17, 10 ]}\nobreak
\hyperdef{L}{X84AC871283A74EC0}{}
{\noindent\textcolor{FuncColor}{$\triangleright$\enspace\texttt{EliminationOrdering({\mdseries\slshape elim[, rest]})\index{EliminationOrdering@\texttt{EliminationOrdering}}
\label{EliminationOrdering}
}\hfill{\scriptsize (function)}}\\


 This function creates an elimination ordering for eliminating the variables in \mbox{\texttt{\mdseries\slshape elim}}. Two monomials are compared first by the exponent vectors for the variables
listed in \mbox{\texttt{\mdseries\slshape elim}} (a lexicographic comparison with respect to the ordering indicated in \mbox{\texttt{\mdseries\slshape elim}}). If these submonomial are equal, the submonomials given by the other
variables are compared by a graded lexicographic ordering (with respect to the
variable order given in \mbox{\texttt{\mdseries\slshape rest}}, if called with two parameters). 

 Both \mbox{\texttt{\mdseries\slshape elim}} and \mbox{\texttt{\mdseries\slshape rest}} may be a list of variables or a list of variable indices. }

 

\subsection{\textcolor{Chapter }{PolynomialReduction}}
\logpage{[ 66, 17, 11 ]}\nobreak
\hyperdef{L}{X7C99593584D478D7}{}
{\noindent\textcolor{FuncColor}{$\triangleright$\enspace\texttt{PolynomialReduction({\mdseries\slshape poly, gens, order})\index{PolynomialReduction@\texttt{PolynomialReduction}}
\label{PolynomialReduction}
}\hfill{\scriptsize (function)}}\\


 reduces the polynomial \mbox{\texttt{\mdseries\slshape poly}} by the ideal generated by the polynomials in \mbox{\texttt{\mdseries\slshape gens}}, using the order \mbox{\texttt{\mdseries\slshape order}} of monomials. Unless \mbox{\texttt{\mdseries\slshape gens}} is a Gr{\"o}bner basis the result is not guaranteed to be unique. 

 The operation returns a list of length two, the first entry is the remainder
after the reduction. The second entry is a list of quotients corresponding to \mbox{\texttt{\mdseries\slshape gens}}. 

 Note that the strategy used by \texttt{PolynomialReduction} differs from the standard textbook reduction algorithm, which is provided by \texttt{PolynomialDivisionAlgorithm} (\ref{PolynomialDivisionAlgorithm}). }

 

\subsection{\textcolor{Chapter }{PolynomialReducedRemainder}}
\logpage{[ 66, 17, 12 ]}\nobreak
\hyperdef{L}{X7DE7D4467EBAD916}{}
{\noindent\textcolor{FuncColor}{$\triangleright$\enspace\texttt{PolynomialReducedRemainder({\mdseries\slshape poly, gens, order})\index{PolynomialReducedRemainder@\texttt{PolynomialReducedRemainder}}
\label{PolynomialReducedRemainder}
}\hfill{\scriptsize (function)}}\\


 this operation does the same way as \texttt{PolynomialReduction} (\ref{PolynomialReduction}) but does not keep track of the actual quotients and returns only the remainder
(it is therefore slightly faster). }

 

\subsection{\textcolor{Chapter }{PolynomialDivisionAlgorithm}}
\logpage{[ 66, 17, 13 ]}\nobreak
\hyperdef{L}{X7C8239057FD4EC03}{}
{\noindent\textcolor{FuncColor}{$\triangleright$\enspace\texttt{PolynomialDivisionAlgorithm({\mdseries\slshape poly, gens, order})\index{PolynomialDivisionAlgorithm@\texttt{PolynomialDivisionAlgorithm}}
\label{PolynomialDivisionAlgorithm}
}\hfill{\scriptsize (function)}}\\


 This function implements the division algorithm for multivariate polynomials
as given in \cite[Theorem 3 in Chapter 2]{coxlittleoshea}. (It might be slower than \texttt{PolynomialReduction} (\ref{PolynomialReduction}) but the remainders are guaranteed to agree with the textbook.) 

 The operation returns a list of length two, the first entry is the remainder
after the reduction. The second entry is a list of quotients corresponding to \mbox{\texttt{\mdseries\slshape gens}}. 
\begin{Verbatim}[commandchars=!@|,fontsize=\small,frame=single,label=Example]
  !gapprompt@gap>| !gapinput@bas:=[x^3*y*z,x*y^2*z,z*y*z^3+x];;|
  !gapprompt@gap>| !gapinput@pol:=x^7*z*bas[1]+y^5*bas[3]+x*z;;|
  !gapprompt@gap>| !gapinput@PolynomialReduction(pol,bas,MonomialLexOrdering());|
  [ -y*z^5, [ x^7*z, 0, y^5+z ] ]
  !gapprompt@gap>| !gapinput@PolynomialReducedRemainder(pol,bas,MonomialLexOrdering());|
  -y*z^5
  !gapprompt@gap>| !gapinput@PolynomialDivisionAlgorithm(pol,bas,MonomialLexOrdering());|
  [ -y*z^5, [ x^7*z, 0, y^5+z ] ]
\end{Verbatim}
 }

 

\subsection{\textcolor{Chapter }{MonomialExtGrlexLess}}
\logpage{[ 66, 17, 14 ]}\nobreak
\hyperdef{L}{X7A30E10B820311D1}{}
{\noindent\textcolor{FuncColor}{$\triangleright$\enspace\texttt{MonomialExtGrlexLess({\mdseries\slshape a, b})\index{MonomialExtGrlexLess@\texttt{MonomialExtGrlexLess}}
\label{MonomialExtGrlexLess}
}\hfill{\scriptsize (function)}}\\


 implements comparison of monomial in their external representation by a ``grlex'' order with $x_1>x_2$ (This is exactly the same as the ordering by \texttt{MonomialGrlexOrdering} (\ref{MonomialGrlexOrdering}), see{\nobreakspace} \ref{Monomial Orderings}). The function takes two monomials \mbox{\texttt{\mdseries\slshape a}} and \mbox{\texttt{\mdseries\slshape b}} in expanded form and returns whether the first is smaller than the second.
(This ordering is also used by \textsf{GAP} internally for representing polynomials as a linear combination of monomials.) 

 See section{\nobreakspace}\ref{The Defining Attributes of Rational Functions} for details on the expanded form of monomials. }

 }

  
\section{\textcolor{Chapter }{Groebner Bases}}\label{Groebner Bases}
\logpage{[ 66, 18, 0 ]}
\hyperdef{L}{X79BAB2937E6085D6}{}
{
  A \emph{Groebner Basis} of an ideal $I$i, in a polynomial ring $R$, with respect to a monomial ordering, is a set of ideal generators \mbox{\texttt{\mdseries\slshape G}} such that the ideal generated by the leading monomials of all polynomials in \mbox{\texttt{\mdseries\slshape G}} is equal to the ideal generated by the leading monomials of all polynomials in \mbox{\texttt{\mdseries\slshape I}}. 

 For more details on Groebner bases see \cite{coxlittleoshea}. 
\subsection{\textcolor{Chapter }{GroebnerBasis}}\logpage{[ 66, 18, 1 ]}
\hyperdef{L}{X7A43611E876B7560}{}
{
\noindent\textcolor{FuncColor}{$\triangleright$\enspace\texttt{GroebnerBasis({\mdseries\slshape L, O})\index{GroebnerBasis@\texttt{GroebnerBasis}!for a list and a monomial ordering}
\label{GroebnerBasis:for a list and a monomial ordering}
}\hfill{\scriptsize (operation)}}\\
\noindent\textcolor{FuncColor}{$\triangleright$\enspace\texttt{GroebnerBasis({\mdseries\slshape I, O})\index{GroebnerBasis@\texttt{GroebnerBasis}!for an ideal and a monomial ordering}
\label{GroebnerBasis:for an ideal and a monomial ordering}
}\hfill{\scriptsize (operation)}}\\
\noindent\textcolor{FuncColor}{$\triangleright$\enspace\texttt{GroebnerBasisNC({\mdseries\slshape L, O})\index{GroebnerBasisNC@\texttt{GroebnerBasisNC}}
\label{GroebnerBasisNC}
}\hfill{\scriptsize (function)}}\\


 Let \mbox{\texttt{\mdseries\slshape O}} be a monomial ordering and \mbox{\texttt{\mdseries\slshape L}} be a list of polynomials that generate an ideal \mbox{\texttt{\mdseries\slshape I}}. This operation returns a Groebner basis of \mbox{\texttt{\mdseries\slshape I}} with respect to the ordering \mbox{\texttt{\mdseries\slshape O}}. 

 \texttt{GroebnerBasisNC} works like \texttt{GroebnerBasis} with the only distinction that the first argument has to be a list of
polynomials and that no test is performed to check whether the ordering is
defined for all occurring variables. 

 Note that \textsf{GAP} at the moment only includes a na{\"\i}ve implementation of Buchberger's
algorithm (which is mainly intended as a teaching tool). It might not be
sufficient for serious problems. 
\begin{Verbatim}[commandchars=!@|,fontsize=\small,frame=single,label=Example]
  !gapprompt@gap>| !gapinput@l:=[x^2+y^2+z^2-1,x^2+z^2-y,x-y];;|
  !gapprompt@gap>| !gapinput@GroebnerBasis(l,MonomialLexOrdering());|
  [ x^2+y^2+z^2-1, x^2+z^2-y, x-y, -y^2-y+1, -z^2+2*y-1,
    1/2*z^4+2*z^2-1/2 ]
  !gapprompt@gap>| !gapinput@GroebnerBasis(l,MonomialLexOrdering([z,x,y]));|
  [ x^2+y^2+z^2-1, x^2+z^2-y, x-y, -y^2-y+1 ]
  !gapprompt@gap>| !gapinput@GroebnerBasis(l,MonomialGrlexOrdering());|
  [ x^2+y^2+z^2-1, x^2+z^2-y, x-y, -y^2-y+1, -z^2+2*y-1 ]
\end{Verbatim}
 }

 
\subsection{\textcolor{Chapter }{ReducedGroebnerBasis}}\logpage{[ 66, 18, 2 ]}
\hyperdef{L}{X7DEF286384967C9E}{}
{
\noindent\textcolor{FuncColor}{$\triangleright$\enspace\texttt{ReducedGroebnerBasis({\mdseries\slshape L, O})\index{ReducedGroebnerBasis@\texttt{ReducedGroebnerBasis}!for a list and a monomial ordering}
\label{ReducedGroebnerBasis:for a list and a monomial ordering}
}\hfill{\scriptsize (operation)}}\\
\noindent\textcolor{FuncColor}{$\triangleright$\enspace\texttt{ReducedGroebnerBasis({\mdseries\slshape I, O})\index{ReducedGroebnerBasis@\texttt{ReducedGroebnerBasis}!for an ideal and a monomial ordering}
\label{ReducedGroebnerBasis:for an ideal and a monomial ordering}
}\hfill{\scriptsize (operation)}}\\


 a Groebner basis $B$ (see{\nobreakspace}\texttt{GroebnerBasis} (\ref{GroebnerBasis:for a list and a monomial ordering})) is \emph{reduced} if no monomial in a polynomial in \mbox{\texttt{\mdseries\slshape B}} is divisible by the leading monomial of another polynomial in $B$. This operation computes a Groebner basis with respect to the monomial
ordering \mbox{\texttt{\mdseries\slshape O}} and then reduces it. 

 
\begin{Verbatim}[commandchars=!@|,fontsize=\small,frame=single,label=Example]
  !gapprompt@gap>| !gapinput@ReducedGroebnerBasis(l,MonomialGrlexOrdering());|
  [ x-y, z^2-2*y+1, y^2+y-1 ]
  !gapprompt@gap>| !gapinput@ReducedGroebnerBasis(l,MonomialLexOrdering());|
  [ z^4+4*z^2-1, -1/2*z^2+y-1/2, -1/2*z^2+x-1/2 ]
  !gapprompt@gap>| !gapinput@ReducedGroebnerBasis(l,MonomialLexOrdering([y,z,x]));|
  [ x^2+x-1, z^2-2*x+1, -x+y ]
\end{Verbatim}
 

 For performance reasons it can be advantageous to define monomial orderings
once and then to reuse them: 

 
\begin{Verbatim}[commandchars=!@|,fontsize=\small,frame=single,label=Example]
  !gapprompt@gap>| !gapinput@ord:=MonomialGrlexOrdering();;|
  !gapprompt@gap>| !gapinput@GroebnerBasis(l,ord);|
  [ x^2+y^2+z^2-1, x^2+z^2-y, x-y, -y^2-y+1, -z^2+2*y-1 ]
  !gapprompt@gap>| !gapinput@ReducedGroebnerBasis(l,ord);|
  [ x-y, z^2-2*y+1, y^2+y-1 ]
\end{Verbatim}
 }

 

\subsection{\textcolor{Chapter }{StoredGroebnerBasis}}
\logpage{[ 66, 18, 3 ]}\nobreak
\hyperdef{L}{X7FC1EFE78498C17C}{}
{\noindent\textcolor{FuncColor}{$\triangleright$\enspace\texttt{StoredGroebnerBasis({\mdseries\slshape I})\index{StoredGroebnerBasis@\texttt{StoredGroebnerBasis}}
\label{StoredGroebnerBasis}
}\hfill{\scriptsize (attribute)}}\\


 For an ideal \mbox{\texttt{\mdseries\slshape I}} in a polynomial ring, this attribute holds a list $[ B, O ]$ where $B$ is a Groebner basis for the monomial ordering $O$. this can be used to test membership or canonical coset representatives. }

 

\subsection{\textcolor{Chapter }{InfoGroebner}}
\logpage{[ 66, 18, 4 ]}\nobreak
\hyperdef{L}{X7C55702786D284A7}{}
{\noindent\textcolor{FuncColor}{$\triangleright$\enspace\texttt{InfoGroebner\index{InfoGroebner@\texttt{InfoGroebner}}
\label{InfoGroebner}
}\hfill{\scriptsize (info class)}}\\


 This info class gives information about Groebner basis calculations. }

 }

  
\section{\textcolor{Chapter }{Rational Function Families}}\label{Rational Function Families}
\logpage{[ 66, 19, 0 ]}
\hyperdef{L}{X8113DD9B781CA6C1}{}
{
  All rational functions defined over a ring lie in the same family, the
rational functions family over this ring. 

 In \textsf{GAP} therefore the family of a polynomial depends only on the family of the
coefficients, all polynomials whose coefficients lie in the same family are ``compatible''. 

\subsection{\textcolor{Chapter }{RationalFunctionsFamily}}
\logpage{[ 66, 19, 1 ]}\nobreak
\hyperdef{L}{X855DD73C78A90BC3}{}
{\noindent\textcolor{FuncColor}{$\triangleright$\enspace\texttt{RationalFunctionsFamily({\mdseries\slshape fam})\index{RationalFunctionsFamily@\texttt{RationalFunctionsFamily}}
\label{RationalFunctionsFamily}
}\hfill{\scriptsize (attribute)}}\\


 creates a family containing rational functions with coefficients in \mbox{\texttt{\mdseries\slshape fam}}. All elements of the \texttt{RationalFunctionsFamily} are rational functions (see{\nobreakspace}\texttt{IsRationalFunction} (\ref{IsRationalFunction})). }

 

\subsection{\textcolor{Chapter }{IsPolynomialFunctionsFamily}}
\logpage{[ 66, 19, 2 ]}\nobreak
\hyperdef{L}{X86E097307D188D3B}{}
{\noindent\textcolor{FuncColor}{$\triangleright$\enspace\texttt{IsPolynomialFunctionsFamily({\mdseries\slshape obj})\index{IsPolynomialFunctionsFamily@\texttt{IsPolynomialFunctionsFamily}}
\label{IsPolynomialFunctionsFamily}
}\hfill{\scriptsize (Category)}}\\
\noindent\textcolor{FuncColor}{$\triangleright$\enspace\texttt{IsRationalFunctionsFamily({\mdseries\slshape obj})\index{IsRationalFunctionsFamily@\texttt{IsRationalFunctionsFamily}}
\label{IsRationalFunctionsFamily}
}\hfill{\scriptsize (Category)}}\\


 \texttt{IsPolynomialFunctionsFamily} is the category of a family of polynomials. For families over an UFD, the
category becomes \texttt{IsRationalFunctionsFamily} (as rational functions and quotients are only provided for families over an
UFD.)  

 
\begin{Verbatim}[commandchars=!@|,fontsize=\small,frame=single,label=Example]
  !gapprompt@gap>| !gapinput@fam:=RationalFunctionsFamily(FamilyObj(1));|
  NewFamily( "RationalFunctionsFamily(...)", [ 618, 620 ],
  [ 82, 85, 89, 93, 97, 100, 103, 107, 111, 618, 620 ] )
\end{Verbatim}
 }

 

\subsection{\textcolor{Chapter }{CoefficientsFamily}}
\logpage{[ 66, 19, 3 ]}\nobreak
\hyperdef{L}{X7AADCA45826866FB}{}
{\noindent\textcolor{FuncColor}{$\triangleright$\enspace\texttt{CoefficientsFamily({\mdseries\slshape rffam})\index{CoefficientsFamily@\texttt{CoefficientsFamily}}
\label{CoefficientsFamily}
}\hfill{\scriptsize (attribute)}}\\


 If \mbox{\texttt{\mdseries\slshape rffam}} has been created as \texttt{RationalFunctionsFamily(\mbox{\texttt{\mdseries\slshape cfam}})} this attribute holds the coefficients family \mbox{\texttt{\mdseries\slshape cfam}}. 

 \textsf{GAP} does \emph{not} embed the base ring in the polynomial ring. While multiplication and addition
of base ring elements to rational functions return the expected results,
polynomials and rational functions are not equal. 
\begin{Verbatim}[commandchars=!@|,fontsize=\small,frame=single,label=Example]
  !gapprompt@gap>| !gapinput@1=Indeterminate(Rationals)^0;|
  false
\end{Verbatim}
 }

 }

  
\section{\textcolor{Chapter }{The Representations of Rational Functions}}\label{The Representations of Rational Functions}
\logpage{[ 66, 20, 0 ]}
\hyperdef{L}{X7E360788785DE530}{}
{
  \textsf{GAP} uses four representations of rational functions: Rational functions given by
numerator and denominator, polynomials, univariate rational functions (given
by coefficient lists for numerator and denominator and valuation) and Laurent
polynomials (given by coefficient list and valuation). 

 These representations do not necessarily reflect mathematical properties:
While an object in the Laurent polynomials representation must be a Laurent
polynomial it might turn out that a rational function given by numerator and
denominator is actually a Laurent polynomial and the property tests in
section{\nobreakspace}\ref{Properties and Attributes of Rational Functions} will find this out. 

 Each representation is associated one or several ``defining attributes'' that give an ``external'' representation (see{\nobreakspace}\ref{The Defining Attributes of Rational Functions}) of the representation in the form of lists and are the defining information
that tells a rational function what it is. 

 \textsf{GAP} also implements methods to compute these attributes for rational functions in \emph{other} representations, provided it would be possible to express an \emph{mathematically equal} rational function in the representation associated with the attribute. (That
is one can always get a numerator/denominator representation of a polynomial
while an arbitrary function of course can compute a polynomial representation
only if it is a polynomial.) 

 Therefore these attributes can be thought of as ``conceptual'' representations that allow us {\textendash}as far as possible{\textendash} to
consider an object as a rational function, a polynomial or a Laurent
polynomial, regardless of the way it is represented in the computer. 

 Functions thus usually do not need to care about the representation of a
rational function. Depending on its (known in the context or determined)
properties, they can access the attribute representing the rational function
in the desired way. 

 Consequently, methods for rational functions are installed for properties and
not for representations. 

 When \emph{creating} new rational functions however they must be created in one of the three
representations. In most cases this will be the representation for which the ``conceptual'' representation in which the calculation was done is the defining attribute. 

 Iterated operations (like forming the product over a list) therefore will tend
to stay in the most suitable representation and the calculation of another
conceptual representation (which may be comparatively expensive in certain
circumstances) is not necessary. }

  
\section{\textcolor{Chapter }{The Defining Attributes of Rational Functions}}\label{The Defining Attributes of Rational Functions}
\logpage{[ 66, 21, 0 ]}
\hyperdef{L}{X7F44CF87801DB965}{}
{
  In general, rational functions are given in terms of monomials. They are
represented by lists, using numbers (see{\nobreakspace}\ref{Indeterminates}) for the indeterminates. 

 \index{Expanded form of monomials} A monomial is a product of powers of indeterminates. A monomial is stored as a
list (we call this the \emph{expanded form} of the monomial) of the form \texttt{[\mbox{\texttt{\mdseries\slshape inum}},\mbox{\texttt{\mdseries\slshape exp}},\mbox{\texttt{\mdseries\slshape inum}},\mbox{\texttt{\mdseries\slshape exp}},...]} where each \mbox{\texttt{\mdseries\slshape inum}} is the number of an indeterminate and \mbox{\texttt{\mdseries\slshape exp}} the corresponding exponent. The list must be sorted according to the numbers
of the indeterminates. Thus for example, if $x$, $y$ and $z$ are the first three indeterminates, the expanded form of the monomial $x^5 z^8 = z^8 x^5$ is \texttt{[ 1, 5, 3, 8 ]}. \index{External representation of polynomials} The representation of a polynomials is a list of the form \texttt{[\mbox{\texttt{\mdseries\slshape mon}},\mbox{\texttt{\mdseries\slshape coeff}},\mbox{\texttt{\mdseries\slshape mon}},\mbox{\texttt{\mdseries\slshape coeff}},...]} where \mbox{\texttt{\mdseries\slshape mon}} is a monomial in expanded form (that is given as list) and \mbox{\texttt{\mdseries\slshape coeff}} its coefficient. The monomials must be sorted according to the total
degree/lexicographic order (This is the same as given by the ``grlex'' monomial ordering, see{\nobreakspace}\texttt{MonomialGrlexOrdering} (\ref{MonomialGrlexOrdering})). We call this the \emph{external representation} of a polynomial. (The reason for ordering is that addition of polynomials
becomes linear in the number of monomials instead of quadratic; the reason for
the particular ordering chose is that it is compatible with multiplication and
thus gives acceptable performance for quotient calculations.) 

 The attributes that give a representation of a rational function as a Laurent
polynomial are \texttt{CoefficientsOfLaurentPolynomial} (\ref{CoefficientsOfLaurentPolynomial}) and \texttt{IndeterminateNumberOfUnivariateRationalFunction} (\ref{IndeterminateNumberOfUnivariateRationalFunction}). 

 Algorithms should use only the attributes \texttt{ExtRepNumeratorRatFun} (\ref{ExtRepNumeratorRatFun}), \texttt{ExtRepDenominatorRatFun} (\ref{ExtRepDenominatorRatFun}), \texttt{ExtRepPolynomialRatFun} (\ref{ExtRepPolynomialRatFun}), \texttt{CoefficientsOfLaurentPolynomial} (\ref{CoefficientsOfLaurentPolynomial}) and {\textendash}if the univariate function is not constant{\textendash} \texttt{IndeterminateNumberOfUnivariateRationalFunction} (\ref{IndeterminateNumberOfUnivariateRationalFunction}) as the low\texttt{\symbol{45}}level interface to work with a polynomial. They
should not refer to the actual representation used. 

\subsection{\textcolor{Chapter }{IsRationalFunctionDefaultRep}}
\logpage{[ 66, 21, 1 ]}\nobreak
\hyperdef{L}{X791E16C67A352263}{}
{\noindent\textcolor{FuncColor}{$\triangleright$\enspace\texttt{IsRationalFunctionDefaultRep({\mdseries\slshape obj})\index{IsRationalFunctionDefaultRep@\texttt{IsRationalFunctionDefaultRep}}
\label{IsRationalFunctionDefaultRep}
}\hfill{\scriptsize (Representation)}}\\


 is the default representation of rational functions. A rational function in
this representation is defined by the attributes \texttt{ExtRepNumeratorRatFun} (\ref{ExtRepNumeratorRatFun}) and \texttt{ExtRepDenominatorRatFun} (\ref{ExtRepDenominatorRatFun}), the values of which are external representations of polynomials. }

 

\subsection{\textcolor{Chapter }{ExtRepNumeratorRatFun}}
\logpage{[ 66, 21, 2 ]}\nobreak
\hyperdef{L}{X7DF955C87CBFC12B}{}
{\noindent\textcolor{FuncColor}{$\triangleright$\enspace\texttt{ExtRepNumeratorRatFun({\mdseries\slshape ratfun})\index{ExtRepNumeratorRatFun@\texttt{ExtRepNumeratorRatFun}}
\label{ExtRepNumeratorRatFun}
}\hfill{\scriptsize (attribute)}}\\


 returns the external representation of the numerator polynomial of the
rational function \mbox{\texttt{\mdseries\slshape ratfun}}. Numerator and denominator are not guaranteed to be cancelled against each
other. }

 

\subsection{\textcolor{Chapter }{ExtRepDenominatorRatFun}}
\logpage{[ 66, 21, 3 ]}\nobreak
\hyperdef{L}{X8059E74D7DCABDBC}{}
{\noindent\textcolor{FuncColor}{$\triangleright$\enspace\texttt{ExtRepDenominatorRatFun({\mdseries\slshape ratfun})\index{ExtRepDenominatorRatFun@\texttt{ExtRepDenominatorRatFun}}
\label{ExtRepDenominatorRatFun}
}\hfill{\scriptsize (attribute)}}\\


 returns the external representation of the denominator polynomial of the
rational function \mbox{\texttt{\mdseries\slshape ratfun}}. Numerator and denominator are not guaranteed to be cancelled against each
other. }

 

\subsection{\textcolor{Chapter }{ZeroCoefficientRatFun}}
\logpage{[ 66, 21, 4 ]}\nobreak
\hyperdef{L}{X84F546F87B5ECFE0}{}
{\noindent\textcolor{FuncColor}{$\triangleright$\enspace\texttt{ZeroCoefficientRatFun({\mdseries\slshape ratfun})\index{ZeroCoefficientRatFun@\texttt{ZeroCoefficientRatFun}}
\label{ZeroCoefficientRatFun}
}\hfill{\scriptsize (operation)}}\\


 returns the zero of the coefficient ring. This might be needed to represent
the zero polynomial for which the external representation of the numerator is
the empty list. }

 

\subsection{\textcolor{Chapter }{IsPolynomialDefaultRep}}
\logpage{[ 66, 21, 5 ]}\nobreak
\hyperdef{L}{X833CE16579AB26E0}{}
{\noindent\textcolor{FuncColor}{$\triangleright$\enspace\texttt{IsPolynomialDefaultRep({\mdseries\slshape obj})\index{IsPolynomialDefaultRep@\texttt{IsPolynomialDefaultRep}}
\label{IsPolynomialDefaultRep}
}\hfill{\scriptsize (Representation)}}\\


 is the default representation of polynomials. A polynomial in this
representation is defined by the components and \texttt{ExtRepNumeratorRatFun} (\ref{ExtRepNumeratorRatFun}) where \texttt{ExtRepNumeratorRatFun} (\ref{ExtRepNumeratorRatFun}) is the external representation of the polynomial. }

 

\subsection{\textcolor{Chapter }{ExtRepPolynomialRatFun}}
\logpage{[ 66, 21, 6 ]}\nobreak
\hyperdef{L}{X8406EE2E8775FBAB}{}
{\noindent\textcolor{FuncColor}{$\triangleright$\enspace\texttt{ExtRepPolynomialRatFun({\mdseries\slshape polynomial})\index{ExtRepPolynomialRatFun@\texttt{ExtRepPolynomialRatFun}}
\label{ExtRepPolynomialRatFun}
}\hfill{\scriptsize (attribute)}}\\


 returns the external representation of a polynomial. The difference to \texttt{ExtRepNumeratorRatFun} (\ref{ExtRepNumeratorRatFun}) is that rational functions might know to be a polynomial but can still have a
non\texttt{\symbol{45}}vanishing denominator. In this case \texttt{ExtRepPolynomialRatFun} has to call a quotient routine. }

 

\subsection{\textcolor{Chapter }{IsLaurentPolynomialDefaultRep}}
\logpage{[ 66, 21, 7 ]}\nobreak
\hyperdef{L}{X7E1B98CC7BADAF56}{}
{\noindent\textcolor{FuncColor}{$\triangleright$\enspace\texttt{IsLaurentPolynomialDefaultRep({\mdseries\slshape obj})\index{IsLaurentPolynomialDefaultRep@\texttt{IsLaurentPolynomialDefaultRep}}
\label{IsLaurentPolynomialDefaultRep}
}\hfill{\scriptsize (Representation)}}\\


 This representation is used for Laurent polynomials and univariate
polynomials. It represents a Laurent polynomial via the attributes \texttt{CoefficientsOfLaurentPolynomial} (\ref{CoefficientsOfLaurentPolynomial}) and \texttt{IndeterminateNumberOfLaurentPolynomial} (\ref{IndeterminateNumberOfLaurentPolynomial}). }

 }

  
\section{\textcolor{Chapter }{Creation of Rational Functions}}\label{Creation of Rational Functions}
\logpage{[ 66, 22, 0 ]}
\hyperdef{L}{X791FADD278A2F32F}{}
{
  The operations \texttt{LaurentPolynomialByCoefficients} (\ref{LaurentPolynomialByCoefficients}), \texttt{PolynomialByExtRep} (\ref{PolynomialByExtRep}) and \texttt{RationalFunctionByExtRep} (\ref{RationalFunctionByExtRep}) are used to construct objects in the three basic representations for rational
functions. 

\subsection{\textcolor{Chapter }{RationalFunctionByExtRep}}
\logpage{[ 66, 22, 1 ]}\nobreak
\hyperdef{L}{X81297E4587A9F2A6}{}
{\noindent\textcolor{FuncColor}{$\triangleright$\enspace\texttt{RationalFunctionByExtRep({\mdseries\slshape rfam, num, den})\index{RationalFunctionByExtRep@\texttt{RationalFunctionByExtRep}}
\label{RationalFunctionByExtRep}
}\hfill{\scriptsize (function)}}\\
\noindent\textcolor{FuncColor}{$\triangleright$\enspace\texttt{RationalFunctionByExtRepNC({\mdseries\slshape rfam, num, den})\index{RationalFunctionByExtRepNC@\texttt{RationalFunctionByExtRepNC}}
\label{RationalFunctionByExtRepNC}
}\hfill{\scriptsize (function)}}\\


 constructs a rational function (in the representation \texttt{IsRationalFunctionDefaultRep} (\ref{IsRationalFunctionDefaultRep})) in the rational function family \mbox{\texttt{\mdseries\slshape rfam}}, the rational function itself is given by the external representations \mbox{\texttt{\mdseries\slshape num}} and \mbox{\texttt{\mdseries\slshape den}} for numerator and denominator. No cancellation takes place. 

 The variant \texttt{RationalFunctionByExtRepNC} does not perform any test of the arguments and thus potentially can create
illegal objects. It only should be used if speed is required and the arguments
are known to be in correct form. }

 

\subsection{\textcolor{Chapter }{PolynomialByExtRep}}
\logpage{[ 66, 22, 2 ]}\nobreak
\hyperdef{L}{X79E445AF7849F48A}{}
{\noindent\textcolor{FuncColor}{$\triangleright$\enspace\texttt{PolynomialByExtRep({\mdseries\slshape rfam, extrep})\index{PolynomialByExtRep@\texttt{PolynomialByExtRep}}
\label{PolynomialByExtRep}
}\hfill{\scriptsize (function)}}\\
\noindent\textcolor{FuncColor}{$\triangleright$\enspace\texttt{PolynomialByExtRepNC({\mdseries\slshape rfam, extrep})\index{PolynomialByExtRepNC@\texttt{PolynomialByExtRepNC}}
\label{PolynomialByExtRepNC}
}\hfill{\scriptsize (function)}}\\


 constructs a polynomial (in the representation \texttt{IsPolynomialDefaultRep} (\ref{IsPolynomialDefaultRep})) in the rational function family \mbox{\texttt{\mdseries\slshape rfam}}, the polynomial itself is given by the external representation \mbox{\texttt{\mdseries\slshape extrep}}. 

 The variant \texttt{PolynomialByExtRepNC} does not perform any test of the arguments and thus potentially can create
invalid objects. It only should be used if speed is required and the arguments
are known to be in correct form. 
\begin{Verbatim}[commandchars=!@|,fontsize=\small,frame=single,label=Example]
  !gapprompt@gap>| !gapinput@fam:=RationalFunctionsFamily(FamilyObj(1));;|
  !gapprompt@gap>| !gapinput@p:=PolynomialByExtRep(fam,[[1,2],1,[2,1,15,7],3]);|
  3*y*x_15^7+x^2
  !gapprompt@gap>| !gapinput@q:=p/(p+1);|
  (3*y*x_15^7+x^2)/(3*y*x_15^7+x^2+1)
  !gapprompt@gap>| !gapinput@ExtRepNumeratorRatFun(q);|
  [ [ 1, 2 ], 1, [ 2, 1, 15, 7 ], 3 ]
  !gapprompt@gap>| !gapinput@ExtRepDenominatorRatFun(q);|
  [ [  ], 1, [ 1, 2 ], 1, [ 2, 1, 15, 7 ], 3 ]
\end{Verbatim}
 }

 

\subsection{\textcolor{Chapter }{LaurentPolynomialByExtRep}}
\logpage{[ 66, 22, 3 ]}\nobreak
\hyperdef{L}{X7E2A46D68472F492}{}
{\noindent\textcolor{FuncColor}{$\triangleright$\enspace\texttt{LaurentPolynomialByExtRep({\mdseries\slshape fam, cofs, val, ind})\index{LaurentPolynomialByExtRep@\texttt{LaurentPolynomialByExtRep}}
\label{LaurentPolynomialByExtRep}
}\hfill{\scriptsize (function)}}\\
\noindent\textcolor{FuncColor}{$\triangleright$\enspace\texttt{LaurentPolynomialByExtRepNC({\mdseries\slshape fam, cofs, val, ind})\index{LaurentPolynomialByExtRepNC@\texttt{LaurentPolynomialByExtRepNC}}
\label{LaurentPolynomialByExtRepNC}
}\hfill{\scriptsize (function)}}\\


 creates a Laurent polynomial in the family \mbox{\texttt{\mdseries\slshape fam}} with [\mbox{\texttt{\mdseries\slshape cofs}},\mbox{\texttt{\mdseries\slshape val}}] as value of \texttt{CoefficientsOfLaurentPolynomial} (\ref{CoefficientsOfLaurentPolynomial}). No coefficient shifting is performed. This is the lowest level function to
create a Laurent polynomial but will rely on the coefficients being shifted
properly and will not perform any tests. Unless this is guaranteed for the
parameters, \texttt{LaurentPolynomialByCoefficients} (\ref{LaurentPolynomialByCoefficients}) should be used. }

 }

  
\section{\textcolor{Chapter }{Arithmetic for External Representations of Polynomials}}\label{Arithmetic for External Representations of Polynomials}
\logpage{[ 66, 23, 0 ]}
\hyperdef{L}{X809028CD7C0EA7CE}{}
{
  The following operations are used internally to perform the arithmetic for
polynomials in their ``external'' representation (see{\nobreakspace}\ref{The Defining Attributes of Rational Functions}) as lists. 

 Functions to perform arithmetic with the coefficient lists of Laurent
polynomials are described in Section{\nobreakspace}\ref{Coefficient List Arithmetic}. 

\subsection{\textcolor{Chapter }{ZippedSum}}
\logpage{[ 66, 23, 1 ]}\nobreak
\hyperdef{L}{X855094857A78ABF9}{}
{\noindent\textcolor{FuncColor}{$\triangleright$\enspace\texttt{ZippedSum({\mdseries\slshape z1, z2, czero, funcs})\index{ZippedSum@\texttt{ZippedSum}}
\label{ZippedSum}
}\hfill{\scriptsize (operation)}}\\


 computes the sum of two external representations of polynomials \mbox{\texttt{\mdseries\slshape z1}} and \mbox{\texttt{\mdseries\slshape z2}}. \mbox{\texttt{\mdseries\slshape czero}} is the appropriate coefficient zero and \mbox{\texttt{\mdseries\slshape funcs}} a list [ \mbox{\texttt{\mdseries\slshape monomial{\textunderscore}less}}, \mbox{\texttt{\mdseries\slshape coefficient{\textunderscore}sum}} ] containing a monomial comparison and a coefficient addition function. This
list can be found in the component \mbox{\texttt{\mdseries\slshape fam}}\texttt{!.zippedSum} of the rational functions family. 

 Note that \mbox{\texttt{\mdseries\slshape coefficient{\textunderscore}sum}} must be a proper ``summation'' function, not a function computing differences. }

 

\subsection{\textcolor{Chapter }{ZippedProduct}}
\logpage{[ 66, 23, 2 ]}\nobreak
\hyperdef{L}{X7B911136782F0F6D}{}
{\noindent\textcolor{FuncColor}{$\triangleright$\enspace\texttt{ZippedProduct({\mdseries\slshape z1, z2, czero, funcs})\index{ZippedProduct@\texttt{ZippedProduct}}
\label{ZippedProduct}
}\hfill{\scriptsize (operation)}}\\


 computes the product of two external representations of polynomials \mbox{\texttt{\mdseries\slshape z1}} and \mbox{\texttt{\mdseries\slshape z2}}. \mbox{\texttt{\mdseries\slshape czero}} is the appropriate coefficient zero and \mbox{\texttt{\mdseries\slshape funcs}} a list [ \mbox{\texttt{\mdseries\slshape monomial{\textunderscore}prod}}, \mbox{\texttt{\mdseries\slshape monomial{\textunderscore}less}}, \mbox{\texttt{\mdseries\slshape coefficient{\textunderscore}sum}}, \mbox{\texttt{\mdseries\slshape coefficient{\textunderscore}prod}}] containing functions to multiply and compare monomials, to add and to
multiply coefficients. This list can be found in the component \texttt{\mbox{\texttt{\mdseries\slshape fam}}!.zippedProduct} of the rational functions family. }

 

\subsection{\textcolor{Chapter }{QuotientPolynomialsExtRep}}
\logpage{[ 66, 23, 3 ]}\nobreak
\hyperdef{L}{X87E5EB8985AF04CD}{}
{\noindent\textcolor{FuncColor}{$\triangleright$\enspace\texttt{QuotientPolynomialsExtRep({\mdseries\slshape fam, a, b})\index{QuotientPolynomialsExtRep@\texttt{QuotientPolynomialsExtRep}}
\label{QuotientPolynomialsExtRep}
}\hfill{\scriptsize (function)}}\\


 Let \mbox{\texttt{\mdseries\slshape a}} and \mbox{\texttt{\mdseries\slshape b}} the external representations of two polynomials in the rational functions
family \mbox{\texttt{\mdseries\slshape fam}}. This function computes the external representation of the quotient of both
polynomials, it returns \texttt{fail} if the polynomial described by \mbox{\texttt{\mdseries\slshape b}} does not divide the polynomial described by \mbox{\texttt{\mdseries\slshape a}}. }

 }

  
\section{\textcolor{Chapter }{Cancellation Tests for Rational Functions}}\label{Cancellation Tests for Rational Functions}
\logpage{[ 66, 24, 0 ]}
\hyperdef{L}{X811B7D8E79E4BD46}{}
{
  The operation \texttt{Gcd} (\ref{Gcd:for (a ring and) several elements}) can be used to test for common factors of two polynomials. This however would
be too expensive to be done in the arithmetic, thus uses the following
operations internally to try to keep the denominators as small as possible 

\subsection{\textcolor{Chapter }{RationalFunctionByExtRepWithCancellation}}
\logpage{[ 66, 24, 1 ]}\nobreak
\hyperdef{L}{X878A1AC87B492E3D}{}
{\noindent\textcolor{FuncColor}{$\triangleright$\enspace\texttt{RationalFunctionByExtRepWithCancellation({\mdseries\slshape rfam, num, den})\index{RationalFunctionByExtRepWithCancellation@\texttt{Rational}\-\texttt{Function}\-\texttt{By}\-\texttt{Ext}\-\texttt{Rep}\-\texttt{With}\-\texttt{Cancellation}}
\label{RationalFunctionByExtRepWithCancellation}
}\hfill{\scriptsize (function)}}\\


 constructs a rational function as \texttt{RationalFunctionByExtRep} (\ref{RationalFunctionByExtRep}) does but tries to cancel out common factors of numerator and denominator,
calling \texttt{TryGcdCancelExtRepPolynomials} (\ref{TryGcdCancelExtRepPolynomials}). }

 

\subsection{\textcolor{Chapter }{TryGcdCancelExtRepPolynomials}}
\logpage{[ 66, 24, 2 ]}\nobreak
\hyperdef{L}{X7BFB55887A153003}{}
{\noindent\textcolor{FuncColor}{$\triangleright$\enspace\texttt{TryGcdCancelExtRepPolynomials({\mdseries\slshape fam, a, b})\index{TryGcdCancelExtRepPolynomials@\texttt{TryGcdCancelExtRepPolynomials}}
\label{TryGcdCancelExtRepPolynomials}
}\hfill{\scriptsize (function)}}\\


 Let \mbox{\texttt{\mdseries\slshape a}} and \mbox{\texttt{\mdseries\slshape b}} be the external representations of two polynomials. This function tries to
cancel common factors between the corresponding polynomials and returns a list $[ a', b' ]$ of external representations of cancelled polynomials. As there is no proper
multivariate GCD cancellation is not guaranteed to be optimal. }

 

\subsection{\textcolor{Chapter }{HeuristicCancelPolynomialsExtRep}}
\logpage{[ 66, 24, 3 ]}\nobreak
\hyperdef{L}{X8477D7337C4A98AB}{}
{\noindent\textcolor{FuncColor}{$\triangleright$\enspace\texttt{HeuristicCancelPolynomialsExtRep({\mdseries\slshape fam, ext1, ext2})\index{HeuristicCancelPolynomialsExtRep@\texttt{HeuristicCancelPolynomialsExtRep}}
\label{HeuristicCancelPolynomialsExtRep}
}\hfill{\scriptsize (operation)}}\\


 is called by \texttt{TryGcdCancelExtRepPolynomials} (\ref{TryGcdCancelExtRepPolynomials}) to perform the actual work. It will return either \texttt{fail} or a new list of external representations of cancelled polynomials. The
cancellation performed is not necessarily optimal. }

 }

 }

         
\chapter{\textcolor{Chapter }{Algebraic extensions of fields}}\label{Algebraic extensions of fields}
\logpage{[ 67, 0, 0 ]}
\hyperdef{L}{X85732CEF7ECFCA68}{}
{
  If we adjoin a root $\alpha$ of an irreducible polynomial $f \in K[x]$ to the field $K$ we get an \emph{algebraic extension} $K(\alpha)$, which is again a field. We call $K$ the \emph{base field} of $K(\alpha)$. 

 By Kronecker's construction, we may identify $K(\alpha)$ with the factor ring $K[x]/(f)$, an identification that also provides a method for computing in these
extension fields. 

 It is important to note that different extensions of the same field are
entirely different (and its elements lie in different families), even if
mathematically one could be embedded in the other one. 

 Currently \textsf{GAP} only allows extension fields of fields $K$, when $K$ itself is not an extension field.  
\section{\textcolor{Chapter }{Creation of Algebraic Extensions}}\label{Creation of Algebraic Extensions}
\logpage{[ 67, 1, 0 ]}
\hyperdef{L}{X7AD9B24E78ADC27F}{}
{
  

\subsection{\textcolor{Chapter }{AlgebraicExtension}}
\logpage{[ 67, 1, 1 ]}\nobreak
\hyperdef{L}{X7CDA90537D2BAC8A}{}
{\noindent\textcolor{FuncColor}{$\triangleright$\enspace\texttt{AlgebraicExtension({\mdseries\slshape K, f[, nam]})\index{AlgebraicExtension@\texttt{AlgebraicExtension}}
\label{AlgebraicExtension}
}\hfill{\scriptsize (operation)}}\\
\noindent\textcolor{FuncColor}{$\triangleright$\enspace\texttt{AlgebraicExtensionNC({\mdseries\slshape K, f[, nam]})\index{AlgebraicExtensionNC@\texttt{AlgebraicExtensionNC}}
\label{AlgebraicExtensionNC}
}\hfill{\scriptsize (operation)}}\\


 constructs an extension \mbox{\texttt{\mdseries\slshape L}} of the field \mbox{\texttt{\mdseries\slshape K}} by one root of the irreducible polynomial \mbox{\texttt{\mdseries\slshape f}}, using Kronecker's construction. \mbox{\texttt{\mdseries\slshape L}} is a field whose \texttt{LeftActingDomain} (\ref{LeftActingDomain}) value is \mbox{\texttt{\mdseries\slshape K}}. The polynomial \mbox{\texttt{\mdseries\slshape f}} is the \texttt{DefiningPolynomial} (\ref{DefiningPolynomial}) value of \mbox{\texttt{\mdseries\slshape L}} and the attribute \texttt{RootOfDefiningPolynomial} (\ref{RootOfDefiningPolynomial}) of \mbox{\texttt{\mdseries\slshape L}} holds a root of \mbox{\texttt{\mdseries\slshape f}} in \mbox{\texttt{\mdseries\slshape L}}. By default this root is printed as \texttt{a}, this string can be overwritten with the optional argument \mbox{\texttt{\mdseries\slshape nam}}. 

 The first version of the command checks that the polynomial \mbox{\texttt{\mdseries\slshape f}} is an irreducible polynomial over \mbox{\texttt{\mdseries\slshape K}}. This check is skipped with the \texttt{NC} variant. 
\begin{Verbatim}[commandchars=!@|,fontsize=\small,frame=single,label=Example]
  !gapprompt@gap>| !gapinput@x:=Indeterminate(Rationals,"x");;|
  !gapprompt@gap>| !gapinput@p:=x^4+3*x^2+1;;|
  !gapprompt@gap>| !gapinput@e:=AlgebraicExtension(Rationals,p);|
  <algebraic extension over the Rationals of degree 4>
  !gapprompt@gap>| !gapinput@IsField(e);|
  true
  !gapprompt@gap>| !gapinput@a:=RootOfDefiningPolynomial(e);|
  a
  !gapprompt@gap>| !gapinput@l := AlgebraicExtensionNC(Rationals, x^24+3*x^2+1, "alpha");;|
  !gapprompt@gap>| !gapinput@RootOfDefiningPolynomial(l)^50;|
  9*alpha^6+6*alpha^4+alpha^2
\end{Verbatim}
 }

 

\subsection{\textcolor{Chapter }{IsAlgebraicExtension}}
\logpage{[ 67, 1, 2 ]}\nobreak
\hyperdef{L}{X811F10217F12B3F9}{}
{\noindent\textcolor{FuncColor}{$\triangleright$\enspace\texttt{IsAlgebraicExtension({\mdseries\slshape obj})\index{IsAlgebraicExtension@\texttt{IsAlgebraicExtension}}
\label{IsAlgebraicExtension}
}\hfill{\scriptsize (Category)}}\\


 is the category of algebraic extensions of fields. 
\begin{Verbatim}[commandchars=!@|,fontsize=\small,frame=single,label=Example]
  !gapprompt@gap>| !gapinput@IsAlgebraicExtension(e);|
  true
  !gapprompt@gap>| !gapinput@IsAlgebraicExtension(Rationals);|
  false
\end{Verbatim}
 }

 }

  
\section{\textcolor{Chapter }{Elements in Algebraic Extensions}}\label{Elements in Algebraic Extensions}
\logpage{[ 67, 2, 0 ]}
\hyperdef{L}{X819C7E6F78817F1E}{}
{
  \index{Operations for algebraic elements} According to Kronecker's construction, the elements of an algebraic extension
are considered to be polynomials in the primitive element. The elements of the
base field are represented as polynomials of degree zero. \textsf{GAP} therefore displays elements of an algebraic extension as polynomials in an
indeterminate ``a'', which is a root of the defining polynomial of the extension. Polynomials of
degree zero are displayed with a leading exclamation mark to indicate that
they are different from elements of the base field. 

 The usual field operations are applicable to algebraic elements. 

 
\begin{Verbatim}[commandchars=@|A,fontsize=\small,frame=single,label=Example]
  @gapprompt|gap>A @gapinput|a^3/(a^2+a+1);A
  -1/2*a^3+1/2*a^2-1/2*a
  @gapprompt|gap>A @gapinput|a*(1/a);A
  !1
\end{Verbatim}
 

 The external representation of algebraic extension elements are the polynomial
coefficients in the primitive element \texttt{a}, the operations \texttt{ExtRepOfObj} (\ref{ExtRepOfObj}) and \texttt{ObjByExtRep} (\ref{ObjByExtRep}) can be used for conversion. 

 
\begin{Verbatim}[commandchars=!@|,fontsize=\small,frame=single,label=Example]
  !gapprompt@gap>| !gapinput@ExtRepOfObj(One(a));|
  [ 1, 0, 0, 0 ]
  !gapprompt@gap>| !gapinput@ExtRepOfObj(a^3+2*a-9);|
  [ -9, 2, 0, 1 ]
  !gapprompt@gap>| !gapinput@ObjByExtRep(FamilyObj(a),[3,19,-27,433]);|
  433*a^3-27*a^2+19*a+3
\end{Verbatim}
 

 \textsf{GAP} does \emph{not} embed the base field in its algebraic extensions and therefore lists which
contain elements of the base field and of the extension are not homogeneous
and thus cannot be used as polynomial coefficients or to form matrices. The
remedy is to multiply the list(s) with the value of the attribute \texttt{One} (\ref{One}) of the extension which will embed all entries in the extension. 

 
\begin{Verbatim}[commandchars=@|A,fontsize=\small,frame=single,label=Example]
  @gapprompt|gap>A @gapinput|m:=[[1,a],[0,1]];A
  [ [ 1, a ], [ 0, 1 ] ]
  @gapprompt|gap>A @gapinput|IsMatrix(m);A
  false
  @gapprompt|gap>A @gapinput|m:=m*One(e);A
  [ [ !1, a ], [ !0, !1 ] ]
  @gapprompt|gap>A @gapinput|IsMatrix(m);A
  true
  @gapprompt|gap>A @gapinput|m^2;A
  [ [ !1, 2*a ], [ !0, !1 ] ]
\end{Verbatim}
 

\subsection{\textcolor{Chapter }{IsAlgebraicElement}}
\logpage{[ 67, 2, 1 ]}\nobreak
\hyperdef{L}{X79695C887FD0AEAB}{}
{\noindent\textcolor{FuncColor}{$\triangleright$\enspace\texttt{IsAlgebraicElement({\mdseries\slshape obj})\index{IsAlgebraicElement@\texttt{IsAlgebraicElement}}
\label{IsAlgebraicElement}
}\hfill{\scriptsize (Category)}}\\


 is the category for elements of an algebraic extension. }

 }

  
\section{\textcolor{Chapter }{Finding Subfields}}\label{Finding Subfields}
\logpage{[ 67, 3, 0 ]}
\hyperdef{L}{X8529BB22865273B1}{}
{
  

\subsection{\textcolor{Chapter }{IdealDecompositionsOfPolynomial}}
\logpage{[ 67, 3, 1 ]}\nobreak
\hyperdef{L}{X7FCAEFBC87651BDD}{}
{\noindent\textcolor{FuncColor}{$\triangleright$\enspace\texttt{IdealDecompositionsOfPolynomial({\mdseries\slshape pol})\index{IdealDecompositionsOfPolynomial@\texttt{IdealDecompositionsOfPolynomial}}
\label{IdealDecompositionsOfPolynomial}
}\hfill{\scriptsize (function)}}\\


 Let $f$ be a univariate, rational, irreducible, polynomial. A pair $g$,$h$ of polynomials of degree strictly smaller than that of $f$, such that $f(x)|g(h(x))$ is called an ideal decomposition. In the context of field extensions, if $\alpha$ is a root of $f$ in a suitable extension and $Q$ the field of rational numbers. Such decompositions correspond to (proper)
subfields $Q < Q(\beta) < Q(\alpha)$, where $g$ is the minimal polynomial of $\beta$. This function determines such decompositions up to equality of the subfields $Q(\beta)$, thus determining subfields of a given algebraic extension. It returns a list
of pairs $[g,h]$ (and an empty list if no such decomposition exists). If the option \mbox{\texttt{\mdseries\slshape onlyone}} is given it returns at most one such decomposition (and performs faster). 
\begin{Verbatim}[commandchars=!@|,fontsize=\small,frame=single,label=Example]
  !gapprompt@gap>| !gapinput@x:=X(Rationals,"x");;pol:=x^8-24*x^6+144*x^4-288*x^2+144;;|
  !gapprompt@gap>| !gapinput@l:=IdealDecompositionsOfPolynomial(pol);|
  [ [ x^2+72*x+144, x^6-20*x^4+60*x^2-36 ],
    [ x^2-48*x+144, x^6-21*x^4+84*x^2-48 ],
    [ x^2+288*x+17280, x^6-24*x^4+132*x^2-288 ],
    [ x^4-24*x^3+144*x^2-288*x+144, x^2 ] ]
  !gapprompt@gap>| !gapinput@List(l,x->Value(x[1],x[2])/pol);|
  [ x^4-16*x^2-8, x^4-18*x^2+33, x^4-24*x^2+120, 1 ]
  !gapprompt@gap>| !gapinput@IdealDecompositionsOfPolynomial(pol:onlyone);|
  [ [ x^2+72*x+144, x^6-20*x^4+60*x^2-36 ] ]
\end{Verbatim}
 In this example the given polynomial is regular with Galois group $Q_8$, as expected we get four proper subfields. }

 }

 }

         
\chapter{\textcolor{Chapter }{p\texttt{\symbol{45}}adic Numbers (preliminary)}}\label{p-adic Numbers}
\logpage{[ 68, 0, 0 ]}
\hyperdef{L}{X7C6B3CBB873253E3}{}
{
  In this chapter $p$ is always a (fixed) prime integer. 

 The $p$\texttt{\symbol{45}}adic numbers $Q_p$ are the completion of the rational numbers with respect to the valuation $\nu_p( p^v \cdot a / b) = v$ if $p$ divides neither $a$ nor $b$. They form a field of characteristic 0 which nevertheless shows some
behaviour of the finite field with $p$ elements. 

 A $p$\texttt{\symbol{45}}adic numbers can be represented by a ``$p$\texttt{\symbol{45}}adic expansion'' which is similar to the decimal expansion used for the reals (but written from
left to right). So for example if $p = 2$, the numbers $1$, $2$, $3$, $4$, $1/2$, and $4/5$ are represented as $1(2)$, $0.1(2)$, $1.1(2)$, $0.01(2)$, $10(2)$, and the infinite periodic expansion $0.010110011001100...(2)$. $p$\texttt{\symbol{45}}adic numbers can be approximated by ignoring higher powers
of $p$, so for example with only 2 digits accuracy $4/5$ would be approximated as $0.01(2)$. This is different from the decimal approximation of real numbers in that $p$\texttt{\symbol{45}}adic approximation is a ring homomorphism on the subrings
of $p$\texttt{\symbol{45}}adic numbers whose valuation is bounded from below so that
rounding errors do not increase with repeated calculations. 

 In \textsf{GAP}, $p$\texttt{\symbol{45}}adic numbers are always represented by such
approximations. A family of approximated $p$\texttt{\symbol{45}}adic numbers consists of $p$\texttt{\symbol{45}}adic numbers with a fixed prime $p$ and a certain precision, and arithmetic with these numbers is done with this
precision.  
\section{\textcolor{Chapter }{Pure p\texttt{\symbol{45}}adic Numbers}}\label{Pure p-adic Numbers}
\logpage{[ 68, 1, 0 ]}
\hyperdef{L}{X7F81667C81655050}{}
{
  Pure $p$\texttt{\symbol{45}}adic numbers are the $p$\texttt{\symbol{45}}adic numbers described so far. 

\subsection{\textcolor{Chapter }{PurePadicNumberFamily}}
\logpage{[ 68, 1, 1 ]}\nobreak
\hyperdef{L}{X82D1AD1D872B480D}{}
{\noindent\textcolor{FuncColor}{$\triangleright$\enspace\texttt{PurePadicNumberFamily({\mdseries\slshape p, precision})\index{PurePadicNumberFamily@\texttt{PurePadicNumberFamily}}
\label{PurePadicNumberFamily}
}\hfill{\scriptsize (function)}}\\


 returns the family of pure $p$\texttt{\symbol{45}}adic numbers over the prime \mbox{\texttt{\mdseries\slshape p}} with \mbox{\texttt{\mdseries\slshape precision}} ``digits''. That is to say, the approximate value will differ from the correct value by
a multiple of $p^{digits}$. }

 

\subsection{\textcolor{Chapter }{PadicNumber (for pure padics)}}
\logpage{[ 68, 1, 2 ]}\nobreak
\hyperdef{L}{X84A79ED87B47CC07}{}
{\noindent\textcolor{FuncColor}{$\triangleright$\enspace\texttt{PadicNumber({\mdseries\slshape fam, rat})\index{PadicNumber@\texttt{PadicNumber}!for pure padics}
\label{PadicNumber:for pure padics}
}\hfill{\scriptsize (operation)}}\\


 returns the element of the $p$\texttt{\symbol{45}}adic number family \mbox{\texttt{\mdseries\slshape fam}} that approximates the rational number \mbox{\texttt{\mdseries\slshape rat}}. 

 $p$\texttt{\symbol{45}}adic numbers allow the usual operations for fields. 

 
\begin{Verbatim}[commandchars=!@|,fontsize=\small,frame=single,label=Example]
  !gapprompt@gap>| !gapinput@fam:=PurePadicNumberFamily(2,20);;|
  !gapprompt@gap>| !gapinput@a:=PadicNumber(fam,4/5);|
  0.010110011001100110011(2)
  !gapprompt@gap>| !gapinput@fam:=PurePadicNumberFamily(2,3);;|
  !gapprompt@gap>| !gapinput@a:=PadicNumber(fam,4/5);|
  0.0101(2)
  !gapprompt@gap>| !gapinput@3*a;|
  0.0111(2)
  !gapprompt@gap>| !gapinput@a/2;|
  0.101(2)
  !gapprompt@gap>| !gapinput@a*10;|
  0.001(2)
\end{Verbatim}
 See \texttt{PadicNumber} (\ref{PadicNumber:for a p-adic extension family and a rational}) for other methods for \texttt{PadicNumber}. }

 

\subsection{\textcolor{Chapter }{Valuation}}
\logpage{[ 68, 1, 3 ]}\nobreak
\hyperdef{L}{X80D67BB67A509A56}{}
{\noindent\textcolor{FuncColor}{$\triangleright$\enspace\texttt{Valuation({\mdseries\slshape obj})\index{Valuation@\texttt{Valuation}}
\label{Valuation}
}\hfill{\scriptsize (operation)}}\\


 The valuation is the $p$\texttt{\symbol{45}}part of the $p$\texttt{\symbol{45}}adic number. See also \texttt{PValuation} (\ref{PValuation}). }

 

\subsection{\textcolor{Chapter }{ShiftedPadicNumber}}
\logpage{[ 68, 1, 4 ]}\nobreak
\hyperdef{L}{X79059A9E876C8198}{}
{\noindent\textcolor{FuncColor}{$\triangleright$\enspace\texttt{ShiftedPadicNumber({\mdseries\slshape padic, int})\index{ShiftedPadicNumber@\texttt{ShiftedPadicNumber}}
\label{ShiftedPadicNumber}
}\hfill{\scriptsize (operation)}}\\


 \texttt{ShiftedPadicNumber} takes a $p$\texttt{\symbol{45}}adic number \mbox{\texttt{\mdseries\slshape padic}} and an integer \mbox{\texttt{\mdseries\slshape shift}} and returns the $p$\texttt{\symbol{45}}adic number $c$, that is \mbox{\texttt{\mdseries\slshape padic}} \texttt{*} $p$\texttt{\texttt{\symbol{94}}}\mbox{\texttt{\mdseries\slshape shift}}. }

 

\subsection{\textcolor{Chapter }{IsPurePadicNumber}}
\logpage{[ 68, 1, 5 ]}\nobreak
\hyperdef{L}{X7AD7FA3786AF9F0E}{}
{\noindent\textcolor{FuncColor}{$\triangleright$\enspace\texttt{IsPurePadicNumber({\mdseries\slshape obj})\index{IsPurePadicNumber@\texttt{IsPurePadicNumber}}
\label{IsPurePadicNumber}
}\hfill{\scriptsize (Category)}}\\


 The category of pure $p$\texttt{\symbol{45}}adic numbers. }

 

\subsection{\textcolor{Chapter }{IsPurePadicNumberFamily}}
\logpage{[ 68, 1, 6 ]}\nobreak
\hyperdef{L}{X83B2BA4586ECAA5C}{}
{\noindent\textcolor{FuncColor}{$\triangleright$\enspace\texttt{IsPurePadicNumberFamily({\mdseries\slshape fam})\index{IsPurePadicNumberFamily@\texttt{IsPurePadicNumberFamily}}
\label{IsPurePadicNumberFamily}
}\hfill{\scriptsize (Category)}}\\


 The family of pure $p$\texttt{\symbol{45}}adic numbers. }

 }

  
\section{\textcolor{Chapter }{Extensions of the p\texttt{\symbol{45}}adic Numbers}}\label{Extensions of the p-adic Numbers}
\logpage{[ 68, 2, 0 ]}
\hyperdef{L}{X83EEF8197D212075}{}
{
  The usual Kronecker construction with an irreducible polynomial can be used to
construct extensions of the $p$\texttt{\symbol{45}}adic numbers. Let $L$ be such an extension. Then there is a subfield $K < L$ such that $K$ is an unramified extension of the $p$\texttt{\symbol{45}}adic numbers and $L/K$ is purely ramified. 

 (For an explanation of ``ramification'' see for example \cite[Section II.7]{neukirch}, or another book on algebraic number theory. Essentially, an extension $L$ of the $p$\texttt{\symbol{45}}adic numbers generated by a rational polynomial $f$ is unramified if $f$ remains squarefree modulo $p$ and is completely ramified if modulo $p$ the polynomial $f$ is a power of a linear factor while remaining irreducible over the $p$\texttt{\symbol{45}}adic numbers.) 

 The representation of extensions of $p$\texttt{\symbol{45}}adic numbers in \textsf{GAP} uses the subfield $K$. 

\subsection{\textcolor{Chapter }{PadicExtensionNumberFamily}}
\logpage{[ 68, 2, 1 ]}\nobreak
\hyperdef{L}{X83EE630D7885DB3D}{}
{\noindent\textcolor{FuncColor}{$\triangleright$\enspace\texttt{PadicExtensionNumberFamily({\mdseries\slshape p, precision, unram, ram})\index{PadicExtensionNumberFamily@\texttt{PadicExtensionNumberFamily}}
\label{PadicExtensionNumberFamily}
}\hfill{\scriptsize (function)}}\\


 An extended $p$\texttt{\symbol{45}}adic field $L$ is given by two polynomials $h$ and $g$ with coefficient lists \mbox{\texttt{\mdseries\slshape unram}} (for the unramified part) and \mbox{\texttt{\mdseries\slshape ram}} (for the ramified part). Then $L$ is isomorphic to $Q_p[x,y]/(h(x),g(y))$. 

 This function takes the prime number \mbox{\texttt{\mdseries\slshape p}} and the two coefficient lists \mbox{\texttt{\mdseries\slshape unram}} and \mbox{\texttt{\mdseries\slshape ram}} for the two polynomials. The polynomial given by the coefficients in \mbox{\texttt{\mdseries\slshape unram}} must be a cyclotomic polynomial and the polynomial given by \mbox{\texttt{\mdseries\slshape ram}} must be either an Eisenstein polynomial or $1+x$. \emph{This is not checked by \textsf{GAP}.} 

 Every number in $L$ is represented as a coefficient list w. r. t. the basis $\{ 1, x, x^2, \ldots, y, xy, x^2 y, \ldots \}$ of $L$. The integer \mbox{\texttt{\mdseries\slshape precision}} is the number of ``digits'' that all the coefficients have. 

 \emph{A general comment:} 

 The polynomials with which \texttt{PadicExtensionNumberFamily} is called define an extension of $Q_p$. It must be ensured that both polynomials are really irreducible over $Q_p$! For example $x^2+x+1$ is \emph{not} irreducible over $Q_p$. Therefore the ``extension'' \texttt{PadicExtensionNumberFamily(3, 4, [1,1,1], [1,1])} contains non\texttt{\symbol{45}}invertible ``pseudo\texttt{\symbol{45}}p\texttt{\symbol{45}}adic numbers''. Conversely, if an ``extension'' contains noninvertible elements then one of the defining polynomials was not
irreducible. }

 

\subsection{\textcolor{Chapter }{PadicNumber (for a p-adic extension family and a rational)}}
\logpage{[ 68, 2, 2 ]}\nobreak
\hyperdef{L}{X7C6F2F018084AFC4}{}
{\noindent\textcolor{FuncColor}{$\triangleright$\enspace\texttt{PadicNumber({\mdseries\slshape fam, rat})\index{PadicNumber@\texttt{PadicNumber}!for a p-adic extension family and a rational}
\label{PadicNumber:for a p-adic extension family and a rational}
}\hfill{\scriptsize (operation)}}\\
\noindent\textcolor{FuncColor}{$\triangleright$\enspace\texttt{PadicNumber({\mdseries\slshape purefam, list})\index{PadicNumber@\texttt{PadicNumber}!for a pure p-adic numbers family and a list}
\label{PadicNumber:for a pure p-adic numbers family and a list}
}\hfill{\scriptsize (operation)}}\\
\noindent\textcolor{FuncColor}{$\triangleright$\enspace\texttt{PadicNumber({\mdseries\slshape extfam, list})\index{PadicNumber@\texttt{PadicNumber}!for a p-adic extension family and a list}
\label{PadicNumber:for a p-adic extension family and a list}
}\hfill{\scriptsize (operation)}}\\


 (see also \texttt{PadicNumber} (\ref{PadicNumber:for pure padics})). 

 \texttt{PadicNumber} creates a $p$\texttt{\symbol{45}}adic number in the $p$\texttt{\symbol{45}}adic numbers family \mbox{\texttt{\mdseries\slshape fam}}. The first form returns the $p$\texttt{\symbol{45}}adic number corresponding to the rational \mbox{\texttt{\mdseries\slshape rat}}. 

 The second form takes a pure $p$\texttt{\symbol{45}}adic numbers family \mbox{\texttt{\mdseries\slshape purefam}} and a list \mbox{\texttt{\mdseries\slshape list}} of length two, and returns the number $p$\texttt{\texttt{\symbol{94}}}\mbox{\texttt{\mdseries\slshape list}}\texttt{[1] * }\mbox{\texttt{\mdseries\slshape list}}\texttt{[2]}. It must be guaranteed that no entry of \mbox{\texttt{\mdseries\slshape list}}\texttt{[2]} is divisible by the prime $p$. (Otherwise precision will get lost.) 

 The third form creates a number in the family \mbox{\texttt{\mdseries\slshape extfam}} of a $p$\texttt{\symbol{45}}adic extension. The second argument must be a list \mbox{\texttt{\mdseries\slshape list}} of length two such that \mbox{\texttt{\mdseries\slshape list}}\texttt{[2]} is the list of coefficients w.r.t. the basis $\{ 1, \ldots, x^{{f-1}} \cdot y^{{e-1}} \}$ of the extended $p$\texttt{\symbol{45}}adic field and \mbox{\texttt{\mdseries\slshape list}}\texttt{[1]} is a common $p$\texttt{\symbol{45}}part of all these coefficients. 

 $p$\texttt{\symbol{45}}adic numbers admit the usual field operations. 
\begin{Verbatim}[commandchars=!@|,fontsize=\small,frame=single,label=Example]
  !gapprompt@gap>| !gapinput@efam:=PadicExtensionNumberFamily(3, 5, [1,1,1], [1,1]);;|
  !gapprompt@gap>| !gapinput@PadicNumber(efam,7/9);|
  padic(120(3),0(3))
\end{Verbatim}
 

 \emph{A word of warning:} 

 Depending on the actual representation of quotients, precision may seem to ``vanish''. For example in \texttt{PadicExtensionNumberFamily(3, 5, [1,1,1], [1,1])} the number \texttt{(1.2000, 0.1210)(3)} can be represented as \texttt{[ 0, [ 1.2000, 0.1210 ] ]} or as \texttt{[ \texttt{\symbol{45}}1, [ 12.000, 1.2100 ] ]} (here the coefficients have to be multiplied by $p^{{-1}}$). 

 So there may be a number \texttt{(1.2, 2.2)(3)} which seems to have only two digits of precision instead of the declared 5.
But internally the number is stored as \texttt{[ \texttt{\symbol{45}}3, [ 0.0012, 0.0022 ] ]} and so has in fact maximum precision. }

 

\subsection{\textcolor{Chapter }{IsPadicExtensionNumber}}
\logpage{[ 68, 2, 3 ]}\nobreak
\hyperdef{L}{X7923FC147BDCC810}{}
{\noindent\textcolor{FuncColor}{$\triangleright$\enspace\texttt{IsPadicExtensionNumber({\mdseries\slshape obj})\index{IsPadicExtensionNumber@\texttt{IsPadicExtensionNumber}}
\label{IsPadicExtensionNumber}
}\hfill{\scriptsize (Category)}}\\


 The category of elements of the extended $p$\texttt{\symbol{45}}adic field. 
\begin{Verbatim}[commandchars=!@|,fontsize=\small,frame=single,label=Example]
  !gapprompt@gap>| !gapinput@ efam:=PadicExtensionNumberFamily(3, 5, [1,1,1], [1,1]);;|
  !gapprompt@gap>| !gapinput@IsPadicExtensionNumber(PadicNumber(efam,7/9));|
  true
\end{Verbatim}
 }

 

\subsection{\textcolor{Chapter }{IsPadicExtensionNumberFamily}}
\logpage{[ 68, 2, 4 ]}\nobreak
\hyperdef{L}{X868807D487DAF713}{}
{\noindent\textcolor{FuncColor}{$\triangleright$\enspace\texttt{IsPadicExtensionNumberFamily({\mdseries\slshape fam})\index{IsPadicExtensionNumberFamily@\texttt{IsPadicExtensionNumberFamily}}
\label{IsPadicExtensionNumberFamily}
}\hfill{\scriptsize (Category)}}\\


 Family of elements of the extended $p$\texttt{\symbol{45}}adic field. 
\begin{Verbatim}[commandchars=!@|,fontsize=\small,frame=single,label=Example]
  !gapprompt@gap>| !gapinput@efam:=PadicExtensionNumberFamily(3, 5, [1,1,1], [1,1]);;|
  !gapprompt@gap>| !gapinput@IsPadicExtensionNumberFamily(efam);|
  true
\end{Verbatim}
 }

 }

 }

          
\chapter{\textcolor{Chapter }{The MeatAxe}}\label{The MeatAxe}
\logpage{[ 69, 0, 0 ]}
\hyperdef{L}{X7BF9D3CB81A8F8F9}{}
{
  The MeatAxe \cite{Par84} is a tool for the examination of submodules of a group algebra. It is a basic
tool for the examination of group actions on
finite\texttt{\symbol{45}}dimensional modules. 

 \textsf{GAP} uses the improved MeatAxe of Derek Holt and Sarah Rees, and also incorporates
further improvements of Ivanyos and Lux. 

 Please note that, consistently with the convention for group actions, the
action of the \textsf{GAP} MeatAxe is always that of matrices on row vectors by multiplication on the
right. If you want to investigate left modules you will have to transpose the
matrices.  
\section{\textcolor{Chapter }{MeatAxe Modules}}\label{MeatAxe Modules}
\logpage{[ 69, 1, 0 ]}
\hyperdef{L}{X85B05BBA78ED7BE2}{}
{
  
\subsection{\textcolor{Chapter }{GModuleByMats}}\logpage{[ 69, 1, 1 ]}
\hyperdef{L}{X801022027B066497}{}
{
\noindent\textcolor{FuncColor}{$\triangleright$\enspace\texttt{GModuleByMats({\mdseries\slshape gens, field})\index{GModuleByMats@\texttt{GModuleByMats}!for generators and a field}
\label{GModuleByMats:for generators and a field}
}\hfill{\scriptsize (function)}}\\
\noindent\textcolor{FuncColor}{$\triangleright$\enspace\texttt{GModuleByMats({\mdseries\slshape emptygens, dim, field})\index{GModuleByMats@\texttt{GModuleByMats}!for empty list, the dimension, and a field}
\label{GModuleByMats:for empty list, the dimension, and a field}
}\hfill{\scriptsize (function)}}\\


 creates a MeatAxe module over \mbox{\texttt{\mdseries\slshape field}} from a list of invertible matrices \mbox{\texttt{\mdseries\slshape gens}} which reflect a group's action. If the list of generators is empty, the
dimension must be given as second argument. 

 MeatAxe routines are on a level with Gaussian elimination. Therefore they do
not deal with \textsf{GAP} modules but essentially with lists of matrices. For the MeatAxe, a module is a
record with components 

 
\begin{description}
\item[{\texttt{generators}}]  A list of matrices which represent a group operation on a finite dimensional
row vector space. 
\item[{\texttt{dimension}}]  The dimension of the vector space (this is the common length of the row
vectors (see{\nobreakspace}\texttt{DimensionOfVectors} (\ref{DimensionOfVectors}))). 
\item[{\texttt{field}}]  The field over which the vector space is defined. 
\end{description}
 

 Once a module has been created its entries may not be changed. A MeatAxe may
create a new component \mbox{\texttt{\mdseries\slshape NameOfMeatAxe}} in which it can store private information. By a MeatAxe ``submodule'' or ``factor module'' we denote actually the \emph{induced action} on the submodule, respectively factor module. Therefore the submodules or
factor modules are again MeatAxe modules. The arrangement of \texttt{generators} is guaranteed to be the same for the induced modules, but to obtain the
complete relation to the original module, the bases used are needed as well. }

 }

  
\section{\textcolor{Chapter }{Module Constructions}}\label{Module Constructions}
\logpage{[ 69, 2, 0 ]}
\hyperdef{L}{X87B82250801A1BD0}{}
{
  
\subsection{\textcolor{Chapter }{NaturalGModule}}\logpage{[ 69, 2, 1 ]}
\hyperdef{L}{X860E128B7D388FBE}{}
{
\noindent\textcolor{FuncColor}{$\triangleright$\enspace\texttt{NaturalGModule({\mdseries\slshape group[, field]})\index{NaturalGModule@\texttt{NaturalGModule}!for matrix group and a field}
\label{NaturalGModule:for matrix group and a field}
}\hfill{\scriptsize (function)}}\\


 creates a MeatAxe module over \mbox{\texttt{\mdseries\slshape field}} from the generators of the matrix group \mbox{\texttt{\mdseries\slshape group}}. If \mbox{\texttt{\mdseries\slshape field}} is not provided then the value returned by \texttt{DefaultFieldOfMatrixGroup} (\ref{DefaultFieldOfMatrixGroup}) is used instead. }

 

\subsection{\textcolor{Chapter }{PermutationGModule}}
\logpage{[ 69, 2, 2 ]}\nobreak
\hyperdef{L}{X8233134A81D58DA3}{}
{\noindent\textcolor{FuncColor}{$\triangleright$\enspace\texttt{PermutationGModule({\mdseries\slshape G, F})\index{PermutationGModule@\texttt{PermutationGModule}}
\label{PermutationGModule}
}\hfill{\scriptsize (function)}}\\


 Called with a permutation group \mbox{\texttt{\mdseries\slshape G}} and a field \mbox{\texttt{\mdseries\slshape F}} (\mbox{\texttt{\mdseries\slshape F}} may be infinite), \texttt{PermutationGModule} returns the natural permutation module $M$ over \mbox{\texttt{\mdseries\slshape F}} for the group of permutation matrices that acts on the canonical basis of $M$ in the same way as \mbox{\texttt{\mdseries\slshape G}} acts on the points up to its largest moved point (see{\nobreakspace}\texttt{LargestMovedPoint} (\ref{LargestMovedPoint:for a list or collection of permutations})). }

 

\subsection{\textcolor{Chapter }{TrivialGModule}}
\logpage{[ 69, 2, 3 ]}\nobreak
\hyperdef{L}{X809B2F4585E7E7A5}{}
{\noindent\textcolor{FuncColor}{$\triangleright$\enspace\texttt{TrivialGModule({\mdseries\slshape G, F})\index{TrivialGModule@\texttt{TrivialGModule}}
\label{TrivialGModule}
}\hfill{\scriptsize (function)}}\\


 Called with a group \mbox{\texttt{\mdseries\slshape G}} and a field \mbox{\texttt{\mdseries\slshape F}} (\mbox{\texttt{\mdseries\slshape F}} may be infinite), \texttt{TrivialGModule} returns the trivial module over \mbox{\texttt{\mdseries\slshape F}}. }

 

\subsection{\textcolor{Chapter }{TensorProductGModule}}
\logpage{[ 69, 2, 4 ]}\nobreak
\hyperdef{L}{X80A50F717B206C98}{}
{\noindent\textcolor{FuncColor}{$\triangleright$\enspace\texttt{TensorProductGModule({\mdseries\slshape m1, m2})\index{TensorProductGModule@\texttt{TensorProductGModule}}
\label{TensorProductGModule}
}\hfill{\scriptsize (function)}}\\


 \texttt{TensorProductGModule} calculates the tensor product of the modules \mbox{\texttt{\mdseries\slshape m1}} and \mbox{\texttt{\mdseries\slshape m2}}. They are assumed to be modules over the same algebra so, in particular, they
should have the same number of generators. }

 

\subsection{\textcolor{Chapter }{DirectSumGModule}}
\logpage{[ 69, 2, 5 ]}\nobreak
\hyperdef{L}{X7A9B22E880F8F9DC}{}
{\noindent\textcolor{FuncColor}{$\triangleright$\enspace\texttt{DirectSumGModule({\mdseries\slshape m1, m2})\index{DirectSumGModule@\texttt{DirectSumGModule}}
\label{DirectSumGModule}
}\hfill{\scriptsize (function)}}\\


 \texttt{DirectSumGModule} calculates the direct sum of the modules \mbox{\texttt{\mdseries\slshape m1}} and \mbox{\texttt{\mdseries\slshape m2}}. They are assumed to be modules over the same algebra so, in particular, they
should have the same number of generators. }

 

\subsection{\textcolor{Chapter }{WedgeGModule}}
\logpage{[ 69, 2, 6 ]}\nobreak
\hyperdef{L}{X7ABC0E98832FEA69}{}
{\noindent\textcolor{FuncColor}{$\triangleright$\enspace\texttt{WedgeGModule({\mdseries\slshape module})\index{WedgeGModule@\texttt{WedgeGModule}}
\label{WedgeGModule}
}\hfill{\scriptsize (function)}}\\


 \texttt{WedgeGModule} calculates the wedge product of a \mbox{\texttt{\mdseries\slshape G}}\texttt{\symbol{45}}module. That is the action on antisymmetric tensors. }

 }

  
\section{\textcolor{Chapter }{Selecting a Different MeatAxe}}\label{Selecting a Different MeatAxe}
\logpage{[ 69, 3, 0 ]}
\hyperdef{L}{X7C77D22782C98D4E}{}
{
  

\subsection{\textcolor{Chapter }{MTX}}
\logpage{[ 69, 3, 1 ]}\nobreak
\hyperdef{L}{X7C2352A17B505AF6}{}
{\noindent\textcolor{FuncColor}{$\triangleright$\enspace\texttt{MTX\index{MTX@\texttt{MTX}}
\label{MTX}
}\hfill{\scriptsize (global variable)}}\\


 All MeatAxe routines are accessed via the global variable \texttt{MTX}, which is a record whose components hold the various functions. It is
possible to have several implementations of a MeatAxe available. Each MeatAxe
represents its routines in an own global variable and assigning \texttt{MTX} to this variable selects the corresponding MeatAxe. }

 }

  
\section{\textcolor{Chapter }{Accessing a Module}}\label{Accessing a Module}
\logpage{[ 69, 4, 0 ]}
\hyperdef{L}{X84AB808B7C543377}{}
{
  Even though a MeatAxe module is a record, its components should never be
accessed outside of MeatAxe functions. Instead the following operations should
be used: 

\subsection{\textcolor{Chapter }{MTX.Generators}}
\logpage{[ 69, 4, 1 ]}\nobreak
\hyperdef{L}{X78E61F7287BF1D0C}{}
{\noindent\textcolor{FuncColor}{$\triangleright$\enspace\texttt{MTX.Generators({\mdseries\slshape module})\index{MTX.Generators@\texttt{MTX.Generators}}
\label{MTX.Generators}
}\hfill{\scriptsize (function)}}\\


 returns a list of matrix generators of \mbox{\texttt{\mdseries\slshape module}}. }

 

\subsection{\textcolor{Chapter }{MTX.Dimension}}
\logpage{[ 69, 4, 2 ]}\nobreak
\hyperdef{L}{X7DF2D6C07D7B09CD}{}
{\noindent\textcolor{FuncColor}{$\triangleright$\enspace\texttt{MTX.Dimension({\mdseries\slshape module})\index{MTX.Dimension@\texttt{MTX.Dimension}}
\label{MTX.Dimension}
}\hfill{\scriptsize (function)}}\\


 returns the dimension in which the matrices act. }

 

\subsection{\textcolor{Chapter }{MTX.Field}}
\logpage{[ 69, 4, 3 ]}\nobreak
\hyperdef{L}{X830C00887CE9323C}{}
{\noindent\textcolor{FuncColor}{$\triangleright$\enspace\texttt{MTX.Field({\mdseries\slshape module})\index{MTX.Field@\texttt{MTX.Field}}
\label{MTX.Field}
}\hfill{\scriptsize (function)}}\\


 returns the field over which \mbox{\texttt{\mdseries\slshape module}} is defined. }

 }

  
\section{\textcolor{Chapter }{Irreducibility Tests}}\label{Irreducibility Tests}
\logpage{[ 69, 5, 0 ]}
\hyperdef{L}{X84D04C7E8423EB5D}{}
{
  

\subsection{\textcolor{Chapter }{MTX.IsIrreducible}}
\logpage{[ 69, 5, 1 ]}\nobreak
\hyperdef{L}{X83BEDF86784A6491}{}
{\noindent\textcolor{FuncColor}{$\triangleright$\enspace\texttt{MTX.IsIrreducible({\mdseries\slshape module})\index{MTX.IsIrreducible@\texttt{MTX.IsIrreducible}}
\label{MTX.IsIrreducible}
}\hfill{\scriptsize (function)}}\\


 tests whether the module \mbox{\texttt{\mdseries\slshape module}} is irreducible (i.e. contains no proper submodules.) }

 

\subsection{\textcolor{Chapter }{MTX.IsAbsolutelyIrreducible}}
\logpage{[ 69, 5, 2 ]}\nobreak
\hyperdef{L}{X876810D679926679}{}
{\noindent\textcolor{FuncColor}{$\triangleright$\enspace\texttt{MTX.IsAbsolutelyIrreducible({\mdseries\slshape module})\index{MTX.IsAbsolutelyIrreducible@\texttt{MTX.IsAbsolutelyIrreducible}}
\label{MTX.IsAbsolutelyIrreducible}
}\hfill{\scriptsize (function)}}\\


 A module is absolutely irreducible if it remains irreducible over the
algebraic closure of the field. (Formally: If the tensor product $L \otimes_K M$ is irreducible where $M$ is the module defined over $K$ and $L$ is the algebraic closure of $K$.) }

 

\subsection{\textcolor{Chapter }{MTX.DegreeSplittingField}}
\logpage{[ 69, 5, 3 ]}\nobreak
\hyperdef{L}{X7E84E1927EBFD483}{}
{\noindent\textcolor{FuncColor}{$\triangleright$\enspace\texttt{MTX.DegreeSplittingField({\mdseries\slshape module})\index{MTX.DegreeSplittingField@\texttt{MTX.DegreeSplittingField}}
\label{MTX.DegreeSplittingField}
}\hfill{\scriptsize (function)}}\\


 returns the degree of the splitting field as extension of the prime field. }

 }

  
\section{\textcolor{Chapter }{Decomposition of modules}}\label{Decomposition of modules}
\logpage{[ 69, 6, 0 ]}
\hyperdef{L}{X791BA495829669C4}{}
{
  A module is \emph{decomposable} if it can be written as the direct sum of two proper submodules (and \emph{indecomposable} if not). Obviously every finite dimensional module is a direct sum of its
indecomposable parts. The \emph{homogeneous components} of a module are the direct sums of isomorphic indecomposable components. They
are uniquely determined. 

 

\subsection{\textcolor{Chapter }{MTX.IsIndecomposable}}
\logpage{[ 69, 6, 1 ]}\nobreak
\hyperdef{L}{X7D9B5B4E7F5A5FBD}{}
{\noindent\textcolor{FuncColor}{$\triangleright$\enspace\texttt{MTX.IsIndecomposable({\mdseries\slshape module})\index{MTX.IsIndecomposable@\texttt{MTX.IsIndecomposable}}
\label{MTX.IsIndecomposable}
}\hfill{\scriptsize (function)}}\\


 returns whether \mbox{\texttt{\mdseries\slshape module}} is indecomposable. }

 

 

\subsection{\textcolor{Chapter }{MTX.Indecomposition}}
\logpage{[ 69, 6, 2 ]}\nobreak
\hyperdef{L}{X781772FD865B9F9C}{}
{\noindent\textcolor{FuncColor}{$\triangleright$\enspace\texttt{MTX.Indecomposition({\mdseries\slshape module})\index{MTX.Indecomposition@\texttt{MTX.Indecomposition}}
\label{MTX.Indecomposition}
}\hfill{\scriptsize (function)}}\\


 returns a decomposition of \mbox{\texttt{\mdseries\slshape module}} as a direct sum of indecomposable modules. It returns a list, each entry is a
list of form [\mbox{\texttt{\mdseries\slshape B}},\mbox{\texttt{\mdseries\slshape ind}}] where \mbox{\texttt{\mdseries\slshape B}} is a list of basis vectors for the indecomposable component and \mbox{\texttt{\mdseries\slshape ind}} the induced module action on this component. (Such a decomposition is not
unique.) }

 

 

\subsection{\textcolor{Chapter }{MTX.HomogeneousComponents}}
\logpage{[ 69, 6, 3 ]}\nobreak
\hyperdef{L}{X7F00E49484FBA7B8}{}
{\noindent\textcolor{FuncColor}{$\triangleright$\enspace\texttt{MTX.HomogeneousComponents({\mdseries\slshape module})\index{MTX.HomogeneousComponents@\texttt{MTX.HomogeneousComponents}}
\label{MTX.HomogeneousComponents}
}\hfill{\scriptsize (function)}}\\


 computes the homogeneous components of \mbox{\texttt{\mdseries\slshape module}} given as sums of indecomposable components. The function returns a list, each
entry of which is a record corresponding to one isomorphism type of
indecomposable components. The record has the following components. 

 
\begin{description}
\item[{\texttt{indices}}]  the index numbers of the indecomposable components, as given by \texttt{MTX.Indecomposition} (\ref{MTX.Indecomposition}), that are in the homogeneous component, 
\item[{\texttt{component}}]  one of the indecomposable components, 
\item[{\texttt{images}}]  a list of the remaining indecomposable components, each given as a record with
the components \texttt{component} (the component itself) and \texttt{isomorphism} (an isomorphism from the defining component to this one). 
\end{description}
 }

 }

  
\section{\textcolor{Chapter }{Finding Submodules}}\label{Finding Submodules}
\logpage{[ 69, 7, 0 ]}
\hyperdef{L}{X85A258567D96B9BE}{}
{
  

\subsection{\textcolor{Chapter }{MTX.SubmoduleGModule}}
\logpage{[ 69, 7, 1 ]}\nobreak
\hyperdef{L}{X80FFB229852B24E9}{}
{\noindent\textcolor{FuncColor}{$\triangleright$\enspace\texttt{MTX.SubmoduleGModule({\mdseries\slshape module, subspace})\index{MTX.SubmoduleGModule@\texttt{MTX.SubmoduleGModule}}
\label{MTX.SubmoduleGModule}
}\hfill{\scriptsize (function)}}\\
\noindent\textcolor{FuncColor}{$\triangleright$\enspace\texttt{MTX.SubGModule({\mdseries\slshape module, subspace})\index{MTX.SubGModule@\texttt{MTX.SubGModule}}
\label{MTX.SubGModule}
}\hfill{\scriptsize (function)}}\\


 \mbox{\texttt{\mdseries\slshape subspace}} should be a subspace of (or a vector in) the underlying vector space of \mbox{\texttt{\mdseries\slshape module}} i.e. the full row space of the same dimension and over the same field as \mbox{\texttt{\mdseries\slshape module}}. A normalized basis of the submodule of \mbox{\texttt{\mdseries\slshape module}} generated by \mbox{\texttt{\mdseries\slshape subspace}} is returned. }

 

\subsection{\textcolor{Chapter }{MTX.ProperSubmoduleBasis}}
\logpage{[ 69, 7, 2 ]}\nobreak
\hyperdef{L}{X81326D84845C206F}{}
{\noindent\textcolor{FuncColor}{$\triangleright$\enspace\texttt{MTX.ProperSubmoduleBasis({\mdseries\slshape module})\index{MTX.ProperSubmoduleBasis@\texttt{MTX.ProperSubmoduleBasis}}
\label{MTX.ProperSubmoduleBasis}
}\hfill{\scriptsize (function)}}\\


 returns the basis of a proper submodule of \mbox{\texttt{\mdseries\slshape module}} and \texttt{fail} if no proper submodule exists. }

 

\subsection{\textcolor{Chapter }{MTX.BasesSubmodules}}
\logpage{[ 69, 7, 3 ]}\nobreak
\hyperdef{L}{X84604D867983DD41}{}
{\noindent\textcolor{FuncColor}{$\triangleright$\enspace\texttt{MTX.BasesSubmodules({\mdseries\slshape module})\index{MTX.BasesSubmodules@\texttt{MTX.BasesSubmodules}}
\label{MTX.BasesSubmodules}
}\hfill{\scriptsize (function)}}\\


 returns a list containing a basis for every submodule. }

 

\subsection{\textcolor{Chapter }{MTX.BasesMinimalSubmodules}}
\logpage{[ 69, 7, 4 ]}\nobreak
\hyperdef{L}{X871D9AF87FABFB00}{}
{\noindent\textcolor{FuncColor}{$\triangleright$\enspace\texttt{MTX.BasesMinimalSubmodules({\mdseries\slshape module})\index{MTX.BasesMinimalSubmodules@\texttt{MTX.BasesMinimalSubmodules}}
\label{MTX.BasesMinimalSubmodules}
}\hfill{\scriptsize (function)}}\\


 returns a list of bases of all minimal submodules. }

 

\subsection{\textcolor{Chapter }{MTX.BasesMaximalSubmodules}}
\logpage{[ 69, 7, 5 ]}\nobreak
\hyperdef{L}{X864527B77A359195}{}
{\noindent\textcolor{FuncColor}{$\triangleright$\enspace\texttt{MTX.BasesMaximalSubmodules({\mdseries\slshape module})\index{MTX.BasesMaximalSubmodules@\texttt{MTX.BasesMaximalSubmodules}}
\label{MTX.BasesMaximalSubmodules}
}\hfill{\scriptsize (function)}}\\


 returns a list of bases of all maximal submodules. }

 

\subsection{\textcolor{Chapter }{MTX.BasisRadical}}
\logpage{[ 69, 7, 6 ]}\nobreak
\hyperdef{L}{X830500CE7ABF6039}{}
{\noindent\textcolor{FuncColor}{$\triangleright$\enspace\texttt{MTX.BasisRadical({\mdseries\slshape module})\index{MTX.BasisRadical@\texttt{MTX.BasisRadical}}
\label{MTX.BasisRadical}
}\hfill{\scriptsize (function)}}\\


 returns a basis of the radical of \mbox{\texttt{\mdseries\slshape module}}. }

 

\subsection{\textcolor{Chapter }{MTX.BasisSocle}}
\logpage{[ 69, 7, 7 ]}\nobreak
\hyperdef{L}{X86A5197D8154A63C}{}
{\noindent\textcolor{FuncColor}{$\triangleright$\enspace\texttt{MTX.BasisSocle({\mdseries\slshape module})\index{MTX.BasisSocle@\texttt{MTX.BasisSocle}}
\label{MTX.BasisSocle}
}\hfill{\scriptsize (function)}}\\


 returns a basis of the socle of \mbox{\texttt{\mdseries\slshape module}}. }

 

\subsection{\textcolor{Chapter }{MTX.BasesMinimalSupermodules}}
\logpage{[ 69, 7, 8 ]}\nobreak
\hyperdef{L}{X7F7FB6687ADE3FD8}{}
{\noindent\textcolor{FuncColor}{$\triangleright$\enspace\texttt{MTX.BasesMinimalSupermodules({\mdseries\slshape module, sub})\index{MTX.BasesMinimalSupermodules@\texttt{MTX.BasesMinimalSupermodules}}
\label{MTX.BasesMinimalSupermodules}
}\hfill{\scriptsize (function)}}\\


 returns a list of bases of all minimal supermodules of the submodule given by
the basis \mbox{\texttt{\mdseries\slshape sub}}. }

 

\subsection{\textcolor{Chapter }{MTX.BasesCompositionSeries}}
\logpage{[ 69, 7, 9 ]}\nobreak
\hyperdef{L}{X79B704998400B9FC}{}
{\noindent\textcolor{FuncColor}{$\triangleright$\enspace\texttt{MTX.BasesCompositionSeries({\mdseries\slshape module})\index{MTX.BasesCompositionSeries@\texttt{MTX.BasesCompositionSeries}}
\label{MTX.BasesCompositionSeries}
}\hfill{\scriptsize (function)}}\\


 returns a list of bases of submodules in a composition series in ascending
order. }

 

\subsection{\textcolor{Chapter }{MTX.CompositionFactors}}
\logpage{[ 69, 7, 10 ]}\nobreak
\hyperdef{L}{X7E77F9A97EA855E2}{}
{\noindent\textcolor{FuncColor}{$\triangleright$\enspace\texttt{MTX.CompositionFactors({\mdseries\slshape module})\index{MTX.CompositionFactors@\texttt{MTX.CompositionFactors}}
\label{MTX.CompositionFactors}
}\hfill{\scriptsize (function)}}\\


 returns a list of composition factors of \mbox{\texttt{\mdseries\slshape module}} in ascending order. }

 

\subsection{\textcolor{Chapter }{MTX.CollectedFactors}}
\logpage{[ 69, 7, 11 ]}\nobreak
\hyperdef{L}{X7E5038F384DBCAEC}{}
{\noindent\textcolor{FuncColor}{$\triangleright$\enspace\texttt{MTX.CollectedFactors({\mdseries\slshape module})\index{MTX.CollectedFactors@\texttt{MTX.CollectedFactors}}
\label{MTX.CollectedFactors}
}\hfill{\scriptsize (function)}}\\


 returns a list giving all irreducible composition factors with their
frequencies. }

 }

  
\section{\textcolor{Chapter }{Induced Actions}}\label{Induced Actions}
\logpage{[ 69, 8, 0 ]}
\hyperdef{L}{X7AE730FB81ED86FE}{}
{
  

\subsection{\textcolor{Chapter }{MTX.NormedBasisAndBaseChange}}
\logpage{[ 69, 8, 1 ]}\nobreak
\hyperdef{L}{X79EA05D4822C2668}{}
{\noindent\textcolor{FuncColor}{$\triangleright$\enspace\texttt{MTX.NormedBasisAndBaseChange({\mdseries\slshape sub})\index{MTX.NormedBasisAndBaseChange@\texttt{MTX.NormedBasisAndBaseChange}}
\label{MTX.NormedBasisAndBaseChange}
}\hfill{\scriptsize (function)}}\\


 returns a list \texttt{[\mbox{\texttt{\mdseries\slshape bas}}, \mbox{\texttt{\mdseries\slshape change}} ]} where \mbox{\texttt{\mdseries\slshape bas}} is a normed basis (i.e. in echelon form with pivots normed to 1) for \mbox{\texttt{\mdseries\slshape sub}} and \mbox{\texttt{\mdseries\slshape change}} is the base change from \mbox{\texttt{\mdseries\slshape bas}} to \mbox{\texttt{\mdseries\slshape sub}} (the basis vectors of \mbox{\texttt{\mdseries\slshape bas}} expressed in coefficients for \mbox{\texttt{\mdseries\slshape sub}}). }

 

\subsection{\textcolor{Chapter }{MTX.InducedActionSubmodule}}
\logpage{[ 69, 8, 2 ]}\nobreak
\hyperdef{L}{X7812D644850D7AED}{}
{\noindent\textcolor{FuncColor}{$\triangleright$\enspace\texttt{MTX.InducedActionSubmodule({\mdseries\slshape module, sub})\index{MTX.InducedActionSubmodule@\texttt{MTX.InducedActionSubmodule}}
\label{MTX.InducedActionSubmodule}
}\hfill{\scriptsize (function)}}\\
\noindent\textcolor{FuncColor}{$\triangleright$\enspace\texttt{MTX.InducedActionSubmoduleNB({\mdseries\slshape module, sub})\index{MTX.InducedActionSubmoduleNB@\texttt{MTX.InducedActionSubmoduleNB}}
\label{MTX.InducedActionSubmoduleNB}
}\hfill{\scriptsize (function)}}\\


 creates a new module corresponding to the action of \mbox{\texttt{\mdseries\slshape module}} on the non\texttt{\symbol{45}}trivial submodule \mbox{\texttt{\mdseries\slshape sub}}. In the \texttt{NB} version the basis \mbox{\texttt{\mdseries\slshape sub}} must be normed. (That is it must be in echelon form with pivots normed to 1,
see{\nobreakspace}\texttt{MTX.NormedBasisAndBaseChange} (\ref{MTX.NormedBasisAndBaseChange}).) }

 

\subsection{\textcolor{Chapter }{MTX.InducedActionFactorModule}}
\logpage{[ 69, 8, 3 ]}\nobreak
\hyperdef{L}{X7EAC61B381385A99}{}
{\noindent\textcolor{FuncColor}{$\triangleright$\enspace\texttt{MTX.InducedActionFactorModule({\mdseries\slshape module, sub[, compl]})\index{MTX.InducedActionFactorModule@\texttt{MTX.InducedActionFactorModule}}
\label{MTX.InducedActionFactorModule}
}\hfill{\scriptsize (function)}}\\


 creates a new module corresponding to the action of \mbox{\texttt{\mdseries\slshape module}} on the factor of the proper submodule \mbox{\texttt{\mdseries\slshape sub}}. If \mbox{\texttt{\mdseries\slshape compl}} is given, it has to be a basis of a (vector
space\texttt{\symbol{45}})complement of \mbox{\texttt{\mdseries\slshape sub}}. The action then will correspond to \mbox{\texttt{\mdseries\slshape compl}}. 

 The basis \mbox{\texttt{\mdseries\slshape sub}} has to be given in normed form. (That is it must be in echelon form with
pivots normed to 1, see{\nobreakspace}\texttt{MTX.NormedBasisAndBaseChange} (\ref{MTX.NormedBasisAndBaseChange})) }

 

\subsection{\textcolor{Chapter }{MTX.InducedActionSubMatrix}}
\logpage{[ 69, 8, 4 ]}\nobreak
\hyperdef{L}{X8753A03A7C7CBFF1}{}
{\noindent\textcolor{FuncColor}{$\triangleright$\enspace\texttt{MTX.InducedActionSubMatrix({\mdseries\slshape mat, sub})\index{MTX.InducedActionSubMatrix@\texttt{MTX.InducedActionSubMatrix}}
\label{MTX.InducedActionSubMatrix}
}\hfill{\scriptsize (function)}}\\
\noindent\textcolor{FuncColor}{$\triangleright$\enspace\texttt{MTX.InducedActionSubMatrixNB({\mdseries\slshape mat, sub})\index{MTX.InducedActionSubMatrixNB@\texttt{MTX.InducedActionSubMatrixNB}}
\label{MTX.InducedActionSubMatrixNB}
}\hfill{\scriptsize (function)}}\\
\noindent\textcolor{FuncColor}{$\triangleright$\enspace\texttt{MTX.InducedActionFactorMatrix({\mdseries\slshape mat, sub[, compl]})\index{MTX.InducedActionFactorMatrix@\texttt{MTX.InducedActionFactorMatrix}}
\label{MTX.InducedActionFactorMatrix}
}\hfill{\scriptsize (function)}}\\


 work the same way as the above functions for modules, but take as input only a
single matrix. }

 

\subsection{\textcolor{Chapter }{MTX.InducedAction}}
\logpage{[ 69, 8, 5 ]}\nobreak
\hyperdef{L}{X7B137BE5877A7FA1}{}
{\noindent\textcolor{FuncColor}{$\triangleright$\enspace\texttt{MTX.InducedAction({\mdseries\slshape module, sub[, type]})\index{MTX.InducedAction@\texttt{MTX.InducedAction}}
\label{MTX.InducedAction}
}\hfill{\scriptsize (function)}}\\


 Computes induced actions on submodules or factor modules and also returns the
corresponding bases. The action taken is binary encoded in \mbox{\texttt{\mdseries\slshape type}}: \texttt{1} stands for subspace action, \texttt{2} for factor action, and \texttt{4} for action of the full module on a subspace adapted basis. The routine returns
the computed results in a list in sequence (\mbox{\texttt{\mdseries\slshape sub}},\mbox{\texttt{\mdseries\slshape quot}},\mbox{\texttt{\mdseries\slshape both}},\mbox{\texttt{\mdseries\slshape basis}}) where \mbox{\texttt{\mdseries\slshape basis}} is a basis for the whole space, extending \mbox{\texttt{\mdseries\slshape sub}}. (Actions which are not computed are omitted, so the returned list may be
shorter.) If no \mbox{\texttt{\mdseries\slshape type}} is given, it is assumed to be \texttt{7}. The basis given in \mbox{\texttt{\mdseries\slshape sub}} must be normed! 

 All these routines return \texttt{fail} if \mbox{\texttt{\mdseries\slshape sub}} is not a proper subspace. }

 }

  
\section{\textcolor{Chapter }{Module Homomorphisms}}\label{Module Homomorphisms}
\logpage{[ 69, 9, 0 ]}
\hyperdef{L}{X8040270F791514C8}{}
{
  

\subsection{\textcolor{Chapter }{MTX.BasisModuleHomomorphisms}}
\logpage{[ 69, 9, 1 ]}\nobreak
\hyperdef{L}{X8292535D8533671C}{}
{\noindent\textcolor{FuncColor}{$\triangleright$\enspace\texttt{MTX.BasisModuleHomomorphisms({\mdseries\slshape module1, module2})\index{MTX.BasisModuleHomomorphisms@\texttt{MTX.BasisModuleHomomorphisms}}
\label{MTX.BasisModuleHomomorphisms}
}\hfill{\scriptsize (function)}}\\


 returns a basis of all module homomorphisms from \mbox{\texttt{\mdseries\slshape module1}} to \mbox{\texttt{\mdseries\slshape module2}}. Homomorphisms are by matrices, whose rows give the images of the standard
basis vectors of \mbox{\texttt{\mdseries\slshape module1}} in the standard basis of \mbox{\texttt{\mdseries\slshape module2}}. }

 

\subsection{\textcolor{Chapter }{MTX.BasisModuleEndomorphisms}}
\logpage{[ 69, 9, 2 ]}\nobreak
\hyperdef{L}{X78EE1274825D9E03}{}
{\noindent\textcolor{FuncColor}{$\triangleright$\enspace\texttt{MTX.BasisModuleEndomorphisms({\mdseries\slshape module})\index{MTX.BasisModuleEndomorphisms@\texttt{MTX.BasisModuleEndomorphisms}}
\label{MTX.BasisModuleEndomorphisms}
}\hfill{\scriptsize (function)}}\\


 returns a basis of all module homomorphisms from \mbox{\texttt{\mdseries\slshape module}} to \mbox{\texttt{\mdseries\slshape module}}. }

 

\subsection{\textcolor{Chapter }{MTX.IsomorphismModules}}
\logpage{[ 69, 9, 3 ]}\nobreak
\hyperdef{L}{X8519B3C486AC8C7E}{}
{\noindent\textcolor{FuncColor}{$\triangleright$\enspace\texttt{MTX.IsomorphismModules({\mdseries\slshape module1, module2})\index{MTX.IsomorphismModules@\texttt{MTX.IsomorphismModules}}
\label{MTX.IsomorphismModules}
}\hfill{\scriptsize (function)}}\\


 If \mbox{\texttt{\mdseries\slshape module1}} and \mbox{\texttt{\mdseries\slshape module2}} are isomorphic modules, this function returns an isomorphism from \mbox{\texttt{\mdseries\slshape module1}} to \mbox{\texttt{\mdseries\slshape module2}} in form of a matrix. It returns \texttt{fail} if the modules are not isomorphic. }

 

\subsection{\textcolor{Chapter }{MTX.ModuleAutomorphisms}}
\logpage{[ 69, 9, 4 ]}\nobreak
\hyperdef{L}{X8442D91F7C4D724F}{}
{\noindent\textcolor{FuncColor}{$\triangleright$\enspace\texttt{MTX.ModuleAutomorphisms({\mdseries\slshape module})\index{MTX.ModuleAutomorphisms@\texttt{MTX.ModuleAutomorphisms}}
\label{MTX.ModuleAutomorphisms}
}\hfill{\scriptsize (function)}}\\


 returns the module automorphisms of \mbox{\texttt{\mdseries\slshape module}} (the set of all isomorphisms from \mbox{\texttt{\mdseries\slshape module}} to itself) as a matrix group. }

 }

  
\section{\textcolor{Chapter }{Module Homomorphisms for irreducible modules}}\label{Module Homomorphisms for irreducible modules}
\logpage{[ 69, 10, 0 ]}
\hyperdef{L}{X850324FF7912A541}{}
{
  The following are lower\texttt{\symbol{45}}level functions that provide
homomorphism functionality for irreducible modules. Generic code should use
the functions in Section \ref{Module Homomorphisms} instead. 

\subsection{\textcolor{Chapter }{MTX.IsEquivalent}}
\logpage{[ 69, 10, 1 ]}\nobreak
\hyperdef{L}{X858D2B0D7AE032D5}{}
{\noindent\textcolor{FuncColor}{$\triangleright$\enspace\texttt{MTX.IsEquivalent({\mdseries\slshape module1, module2})\index{MTX.IsEquivalent@\texttt{MTX.IsEquivalent}}
\label{MTX.IsEquivalent}
}\hfill{\scriptsize (function)}}\\


 tests two irreducible modules for equivalence. }

 

\subsection{\textcolor{Chapter }{MTX.IsomorphismIrred}}
\logpage{[ 69, 10, 2 ]}\nobreak
\hyperdef{L}{X7E86F5B67CBD7C41}{}
{\noindent\textcolor{FuncColor}{$\triangleright$\enspace\texttt{MTX.IsomorphismIrred({\mdseries\slshape module1, module2})\index{MTX.IsomorphismIrred@\texttt{MTX.IsomorphismIrred}}
\label{MTX.IsomorphismIrred}
}\hfill{\scriptsize (function)}}\\


 returns an isomorphism from \mbox{\texttt{\mdseries\slshape module1}} to \mbox{\texttt{\mdseries\slshape module2}} (if one exists), and \texttt{fail} otherwise. It requires that one of the modules is known to be irreducible. It
implicitly assumes that the same group is acting, otherwise the results are
unpredictable. The isomorphism is given by a matrix $M$, whose rows give the images of the standard basis vectors of \mbox{\texttt{\mdseries\slshape module1}} in the standard basis of \mbox{\texttt{\mdseries\slshape module2}}. That is, conjugation of the generators of \mbox{\texttt{\mdseries\slshape module2}} with $M$ yields the generators of \mbox{\texttt{\mdseries\slshape module1}}. }

 

\subsection{\textcolor{Chapter }{MTX.Homomorphism}}
\logpage{[ 69, 10, 3 ]}\nobreak
\hyperdef{L}{X807AE3AC7E9B7CFF}{}
{\noindent\textcolor{FuncColor}{$\triangleright$\enspace\texttt{MTX.Homomorphism({\mdseries\slshape module1, module2, mat})\index{MTX.Homomorphism@\texttt{MTX.Homomorphism}}
\label{MTX.Homomorphism}
}\hfill{\scriptsize (function)}}\\


 \mbox{\texttt{\mdseries\slshape mat}} should be a \mbox{\texttt{\mdseries\slshape dim1}} $\times$ \mbox{\texttt{\mdseries\slshape dim2}} matrix defining a homomorphism from \mbox{\texttt{\mdseries\slshape module1}} to \mbox{\texttt{\mdseries\slshape module2}}. This function verifies that \mbox{\texttt{\mdseries\slshape mat}} really does define a module homomorphism, and then returns the corresponding
homomorphism between the underlying row spaces of the modules. This can be
used for computing kernels, images and pre\texttt{\symbol{45}}images. }

 

\subsection{\textcolor{Chapter }{MTX.Homomorphisms}}
\logpage{[ 69, 10, 4 ]}\nobreak
\hyperdef{L}{X7BC612D2860C582B}{}
{\noindent\textcolor{FuncColor}{$\triangleright$\enspace\texttt{MTX.Homomorphisms({\mdseries\slshape module1, module2})\index{MTX.Homomorphisms@\texttt{MTX.Homomorphisms}}
\label{MTX.Homomorphisms}
}\hfill{\scriptsize (function)}}\\


 returns a basis of the space of all homomorphisms from the irreducible module \mbox{\texttt{\mdseries\slshape module1}} to \mbox{\texttt{\mdseries\slshape module2}}. }

 

\subsection{\textcolor{Chapter }{MTX.Distinguish}}
\logpage{[ 69, 10, 5 ]}\nobreak
\hyperdef{L}{X81A6ECB078D4441C}{}
{\noindent\textcolor{FuncColor}{$\triangleright$\enspace\texttt{MTX.Distinguish({\mdseries\slshape cf, nr})\index{MTX.Distinguish@\texttt{MTX.Distinguish}}
\label{MTX.Distinguish}
}\hfill{\scriptsize (function)}}\\


 Let \mbox{\texttt{\mdseries\slshape cf}} be the output of \texttt{MTX.CollectedFactors} (\ref{MTX.CollectedFactors}). This routine tries to find a group algebra element that has nullity zero on
all composition factors except number \mbox{\texttt{\mdseries\slshape nr}}. }

 }

  
\section{\textcolor{Chapter }{MeatAxe Functionality for Invariant Forms}}\label{meataxe:Invariant Forms}
\logpage{[ 69, 11, 0 ]}
\hyperdef{L}{X7B426E4679C1AF25}{}
{
  The functions in this section can only be applied to an absolutely irreducible
MeatAxe module. 

\subsection{\textcolor{Chapter }{MTX.InvariantBilinearForm}}
\logpage{[ 69, 11, 1 ]}\nobreak
\hyperdef{L}{X78B114E78227EA37}{}
{\noindent\textcolor{FuncColor}{$\triangleright$\enspace\texttt{MTX.InvariantBilinearForm({\mdseries\slshape module})\index{MTX.InvariantBilinearForm@\texttt{MTX.InvariantBilinearForm}}
\label{MTX.InvariantBilinearForm}
}\hfill{\scriptsize (function)}}\\


 returns an invariant bilinear form, which may be symmetric or
anti\texttt{\symbol{45}}symmetric, of \mbox{\texttt{\mdseries\slshape module}}, or \texttt{fail} if no such form exists. }

 

\subsection{\textcolor{Chapter }{MTX.InvariantSesquilinearForm}}
\logpage{[ 69, 11, 2 ]}\nobreak
\hyperdef{L}{X7E1F430278A334E1}{}
{\noindent\textcolor{FuncColor}{$\triangleright$\enspace\texttt{MTX.InvariantSesquilinearForm({\mdseries\slshape module})\index{MTX.InvariantSesquilinearForm@\texttt{MTX.InvariantSesquilinearForm}}
\label{MTX.InvariantSesquilinearForm}
}\hfill{\scriptsize (function)}}\\


 returns an invariant hermitian (= self\texttt{\symbol{45}}adjoint)
sesquilinear form of \mbox{\texttt{\mdseries\slshape module}}, which must be defined over a finite field whose order is a square, or \texttt{fail} if no such form exists. }

 

\subsection{\textcolor{Chapter }{MTX.InvariantQuadraticForm}}
\logpage{[ 69, 11, 3 ]}\nobreak
\hyperdef{L}{X7ADE65997F16EE63}{}
{\noindent\textcolor{FuncColor}{$\triangleright$\enspace\texttt{MTX.InvariantQuadraticForm({\mdseries\slshape module})\index{MTX.InvariantQuadraticForm@\texttt{MTX.InvariantQuadraticForm}}
\label{MTX.InvariantQuadraticForm}
}\hfill{\scriptsize (function)}}\\


 returns either the matrix of an invariant quadratic form of the absolutely
irreducible module \mbox{\texttt{\mdseries\slshape module}}, or \texttt{fail}. 

 If the characteristic of \mbox{\texttt{\mdseries\slshape module}} is odd then \texttt{fail} is returned if there is no nonzero invariant bilinear form, otherwise a matrix
of the bilinear form divided by $2$ is returned; note that this matrix may be antisymmetric and thus describe the
zero quadratic form. If the characteristic of \mbox{\texttt{\mdseries\slshape module}} is $2$ then \texttt{fail} is returned if \mbox{\texttt{\mdseries\slshape module}} does not admit a nonzero quadratic form, otherwise a lower triangular matrix
describing the form is returned. 

 An error is signalled if \mbox{\texttt{\mdseries\slshape module}} is not absolutely irreducible. 

 
\begin{Verbatim}[commandchars=!@|,fontsize=\small,frame=single,label=Example]
  !gapprompt@gap>| !gapinput@g:= SO(-1, 4, 2);;|
  !gapprompt@gap>| !gapinput@m:= NaturalGModule( g );;|
  !gapprompt@gap>| !gapinput@Display( MTX.InvariantQuadraticForm( m ) );|
   . . . .
   1 . . .
   . . 1 .
   . . 1 1
  !gapprompt@gap>| !gapinput@g:= Sp(4, 2);;|
  !gapprompt@gap>| !gapinput@m:= NaturalGModule( g );;|
  !gapprompt@gap>| !gapinput@MTX.InvariantQuadraticForm( m );|
  fail
  !gapprompt@gap>| !gapinput@g:= Sp(4, 3);;|
  !gapprompt@gap>| !gapinput@m:= NaturalGModule( g );;|
  !gapprompt@gap>| !gapinput@q:= MTX.InvariantQuadraticForm( m );;|
  !gapprompt@gap>| !gapinput@q = - TransposedMat( q );  # antisymmetric inv. bilinear form|
  true
\end{Verbatim}
 }

 

\subsection{\textcolor{Chapter }{MTX.BasisInOrbit}}
\logpage{[ 69, 11, 4 ]}\nobreak
\hyperdef{L}{X78E60EFE802AEBC1}{}
{\noindent\textcolor{FuncColor}{$\triangleright$\enspace\texttt{MTX.BasisInOrbit({\mdseries\slshape module})\index{MTX.BasisInOrbit@\texttt{MTX.BasisInOrbit}}
\label{MTX.BasisInOrbit}
}\hfill{\scriptsize (function)}}\\


 returns a basis of the underlying vector space of \mbox{\texttt{\mdseries\slshape module}} which is contained in an orbit of the action of the generators of module on
that space. This is used by \texttt{MTX.InvariantQuadraticForm} (\ref{MTX.InvariantQuadraticForm}) in characteristic 2. }

 

\subsection{\textcolor{Chapter }{MTX.OrthogonalSign}}
\logpage{[ 69, 11, 5 ]}\nobreak
\hyperdef{L}{X8168EB348474046B}{}
{\noindent\textcolor{FuncColor}{$\triangleright$\enspace\texttt{MTX.OrthogonalSign({\mdseries\slshape module})\index{MTX.OrthogonalSign@\texttt{MTX.OrthogonalSign}}
\label{MTX.OrthogonalSign}
}\hfill{\scriptsize (function)}}\\


 Let \mbox{\texttt{\mdseries\slshape module}} be an absolutely irreducible $G$\texttt{\symbol{45}}module. If \mbox{\texttt{\mdseries\slshape module}} does not fix a nondegenerate quadratic form see \texttt{MTX.InvariantQuadraticForm} (\ref{MTX.InvariantQuadraticForm}) then \texttt{fail} is returned. Otherwise the sign $\epsilon \in \{ -1, 0, 1 \}$ is returned such that $G$ embeds into the general orthogonal group $GO^{\epsilon}(d, q)$ w.r.t. the invariant quadratic form, see \texttt{GeneralOrthogonalGroup} (\ref{GeneralOrthogonalGroup}). That is, \texttt{0} is returned if \mbox{\texttt{\mdseries\slshape module}} has odd dimension, and \texttt{1} or \texttt{\texttt{\symbol{45}}1} is returned if the orthogonal group has plus or minus type, respectively. 

 An error is signalled if \mbox{\texttt{\mdseries\slshape module}} is not absolutely irreducible. 

 The \texttt{SMTX} implementation uses an algorithm due to Jon Thackray. 

 
\begin{Verbatim}[commandchars=!@|,fontsize=\small,frame=single,label=Example]
  !gapprompt@gap>| !gapinput@mats:= GeneratorsOfGroup( GO(1,4,2) );;|
  !gapprompt@gap>| !gapinput@MTX.OrthogonalSign( GModuleByMats( mats, GF(2) ) );|
  1
  !gapprompt@gap>| !gapinput@mats:= GeneratorsOfGroup( GO(-1,4,2) );;|
  !gapprompt@gap>| !gapinput@MTX.OrthogonalSign( GModuleByMats( mats, GF(2) ) );|
  -1
  !gapprompt@gap>| !gapinput@mats:= GeneratorsOfGroup( GO(5,3) );;|
  !gapprompt@gap>| !gapinput@MTX.OrthogonalSign( GModuleByMats( mats, GF(3) ) );|
  0
  !gapprompt@gap>| !gapinput@mats:= GeneratorsOfGroup( SP(4,2) );;|
  !gapprompt@gap>| !gapinput@MTX.OrthogonalSign( GModuleByMats( mats, GF(2) ) );|
  fail
\end{Verbatim}
 }

 }

  
\section{\textcolor{Chapter }{The Smash MeatAxe}}\label{The Smash MeatAxe}
\logpage{[ 69, 12, 0 ]}
\hyperdef{L}{X87B0E3237BA056FC}{}
{
  The standard MeatAxe provided in the \textsf{GAP} library is based on the MeatAxe in the \textsf{GAP}{\nobreakspace}3 package \textsf{Smash}, originally written by Derek Holt and Sarah Rees \cite{HR94}. It is accessible via the variable \texttt{SMTX} to which \texttt{MTX} (\ref{MTX}) is assigned by default. For the sake of completeness the remaining sections
document more technical functions of this MeatAxe. 

\subsection{\textcolor{Chapter }{SMTX.RandomIrreducibleSubGModule}}
\logpage{[ 69, 12, 1 ]}\nobreak
\hyperdef{L}{X7E78525883E715E1}{}
{\noindent\textcolor{FuncColor}{$\triangleright$\enspace\texttt{SMTX.RandomIrreducibleSubGModule({\mdseries\slshape module})\index{SMTX.RandomIrreducibleSubGModule@\texttt{SMTX.RandomIrreducibleSubGModule}}
\label{SMTX.RandomIrreducibleSubGModule}
}\hfill{\scriptsize (function)}}\\


 returns the module action on a random irreducible submodule. }

 

\subsection{\textcolor{Chapter }{SMTX.GoodElementGModule}}
\logpage{[ 69, 12, 2 ]}\nobreak
\hyperdef{L}{X7EA698517A19D35B}{}
{\noindent\textcolor{FuncColor}{$\triangleright$\enspace\texttt{SMTX.GoodElementGModule({\mdseries\slshape module})\index{SMTX.GoodElementGModule@\texttt{SMTX.GoodElementGModule}}
\label{SMTX.GoodElementGModule}
}\hfill{\scriptsize (function)}}\\


 finds an element with minimal possible nullspace dimension if \mbox{\texttt{\mdseries\slshape module}} is known to be irreducible. }

 

\subsection{\textcolor{Chapter }{SMTX.SortHomGModule}}
\logpage{[ 69, 12, 3 ]}\nobreak
\hyperdef{L}{X811339547D341BBE}{}
{\noindent\textcolor{FuncColor}{$\triangleright$\enspace\texttt{SMTX.SortHomGModule({\mdseries\slshape module1, module2, homs})\index{SMTX.SortHomGModule@\texttt{SMTX.SortHomGModule}}
\label{SMTX.SortHomGModule}
}\hfill{\scriptsize (function)}}\\


 Function to sort the output of \texttt{Homomorphisms}. }

 

\subsection{\textcolor{Chapter }{SMTX.MinimalSubGModules}}
\logpage{[ 69, 12, 4 ]}\nobreak
\hyperdef{L}{X86B6092681221D7A}{}
{\noindent\textcolor{FuncColor}{$\triangleright$\enspace\texttt{SMTX.MinimalSubGModules({\mdseries\slshape module1, module2[, max]})\index{SMTX.MinimalSubGModules@\texttt{SMTX.MinimalSubGModules}}
\label{SMTX.MinimalSubGModules}
}\hfill{\scriptsize (function)}}\\


 returns (at most \mbox{\texttt{\mdseries\slshape max}}) bases of submodules of \mbox{\texttt{\mdseries\slshape module2}} which are isomorphic to the irreducible module \mbox{\texttt{\mdseries\slshape module1}}. }

 

\subsection{\textcolor{Chapter }{SMTX.Setter}}
\logpage{[ 69, 12, 5 ]}\nobreak
\hyperdef{L}{X87E49FCD867983B5}{}
{\noindent\textcolor{FuncColor}{$\triangleright$\enspace\texttt{SMTX.Setter({\mdseries\slshape string})\index{SMTX.Setter@\texttt{SMTX.Setter}}
\label{SMTX.Setter}
}\hfill{\scriptsize (function)}}\\


 returns a setter function for the component \texttt{smashMeataxe.(string)}. }

 

\subsection{\textcolor{Chapter }{SMTX.Getter}}
\logpage{[ 69, 12, 6 ]}\nobreak
\hyperdef{L}{X7E60EBC57FFDF7BD}{}
{\noindent\textcolor{FuncColor}{$\triangleright$\enspace\texttt{SMTX.Getter({\mdseries\slshape string})\index{SMTX.Getter@\texttt{SMTX.Getter}}
\label{SMTX.Getter}
}\hfill{\scriptsize (function)}}\\


 returns a getter function for the component \texttt{smashMeataxe.(string)}. }

 

\subsection{\textcolor{Chapter }{SMTX.IrreducibilityTest}}
\logpage{[ 69, 12, 7 ]}\nobreak
\hyperdef{L}{X808345D784E0AC85}{}
{\noindent\textcolor{FuncColor}{$\triangleright$\enspace\texttt{SMTX.IrreducibilityTest({\mdseries\slshape module})\index{SMTX.IrreducibilityTest@\texttt{SMTX.IrreducibilityTest}}
\label{SMTX.IrreducibilityTest}
}\hfill{\scriptsize (function)}}\\


 Tests for irreducibility and sets a subbasis if reducible. It neither sets an
irreducibility flag, nor tests it. Thus the routine also can simply be called
to obtain a random submodule. }

 

\subsection{\textcolor{Chapter }{SMTX.AbsoluteIrreducibilityTest}}
\logpage{[ 69, 12, 8 ]}\nobreak
\hyperdef{L}{X7E692DC97AFB661E}{}
{\noindent\textcolor{FuncColor}{$\triangleright$\enspace\texttt{SMTX.AbsoluteIrreducibilityTest({\mdseries\slshape module})\index{SMTX.AbsoluteIrreducibilityTest@\texttt{SMTX.AbsoluteIrreducibilityTest}}
\label{SMTX.AbsoluteIrreducibilityTest}
}\hfill{\scriptsize (function)}}\\


 Tests for absolute irreducibility and sets splitting field degree. It neither
sets an absolute irreducibility flag, nor tests it. }

 

\subsection{\textcolor{Chapter }{SMTX.MinimalSubGModule}}
\logpage{[ 69, 12, 9 ]}\nobreak
\hyperdef{L}{X80BC392285994DA8}{}
{\noindent\textcolor{FuncColor}{$\triangleright$\enspace\texttt{SMTX.MinimalSubGModule({\mdseries\slshape module, cf, nr})\index{SMTX.MinimalSubGModule@\texttt{SMTX.MinimalSubGModule}}
\label{SMTX.MinimalSubGModule}
}\hfill{\scriptsize (function)}}\\


 returns the basis of a minimal submodule of \mbox{\texttt{\mdseries\slshape module}} containing the indicated composition factor. It assumes \texttt{Distinguish} has been called already. }

 

\subsection{\textcolor{Chapter }{SMTX.MatrixSum}}
\logpage{[ 69, 12, 10 ]}\nobreak
\hyperdef{L}{X79EF16677C2EE095}{}
{\noindent\textcolor{FuncColor}{$\triangleright$\enspace\texttt{SMTX.MatrixSum({\mdseries\slshape matrices1, matrices2})\index{SMTX.MatrixSum@\texttt{SMTX.MatrixSum}}
\label{SMTX.MatrixSum}
}\hfill{\scriptsize (function)}}\\


 creates the direct sum of two matrix lists. }

 

\subsection{\textcolor{Chapter }{SMTX.CompleteBasis}}
\logpage{[ 69, 12, 11 ]}\nobreak
\hyperdef{L}{X7D1471077A774C81}{}
{\noindent\textcolor{FuncColor}{$\triangleright$\enspace\texttt{SMTX.CompleteBasis({\mdseries\slshape module, pbasis})\index{SMTX.CompleteBasis@\texttt{SMTX.CompleteBasis}}
\label{SMTX.CompleteBasis}
}\hfill{\scriptsize (function)}}\\


 extends the partial basis \mbox{\texttt{\mdseries\slshape pbasis}} to a basis of the full space by action of \mbox{\texttt{\mdseries\slshape module}}. It returns whether it succeeded. }

 }

  
\section{\textcolor{Chapter }{Smash MeatAxe Flags}}\label{Smash MeatAxe Flags}
\logpage{[ 69, 13, 0 ]}
\hyperdef{L}{X7FDF8F3F83B83336}{}
{
  The following getter routines access internal flags. For each routine, the
appropriate setter's name is prefixed with \texttt{Set}. 

\subsection{\textcolor{Chapter }{SMTX.Subbasis}}
\logpage{[ 69, 13, 1 ]}\nobreak
\hyperdef{L}{X84A93AC482A1946D}{}
{\noindent\textcolor{FuncColor}{$\triangleright$\enspace\texttt{SMTX.Subbasis({\mdseries\slshape module})\index{SMTX.Subbasis@\texttt{SMTX.Subbasis}}
\label{SMTX.Subbasis}
}\hfill{\scriptsize (function)}}\\


 Basis of a submodule. }

 

\subsection{\textcolor{Chapter }{SMTX.AlgEl}}
\logpage{[ 69, 13, 2 ]}\nobreak
\hyperdef{L}{X7ABCD69880772B2D}{}
{\noindent\textcolor{FuncColor}{$\triangleright$\enspace\texttt{SMTX.AlgEl({\mdseries\slshape module})\index{SMTX.AlgEl@\texttt{SMTX.AlgEl}}
\label{SMTX.AlgEl}
}\hfill{\scriptsize (function)}}\\


 list \texttt{[newgens,coefflist]} giving an algebra element used for chopping. }

 

\subsection{\textcolor{Chapter }{SMTX.AlgElMat}}
\logpage{[ 69, 13, 3 ]}\nobreak
\hyperdef{L}{X7D6C947A7C8C14B2}{}
{\noindent\textcolor{FuncColor}{$\triangleright$\enspace\texttt{SMTX.AlgElMat({\mdseries\slshape module})\index{SMTX.AlgElMat@\texttt{SMTX.AlgElMat}}
\label{SMTX.AlgElMat}
}\hfill{\scriptsize (function)}}\\


 matrix of \texttt{SMTX.AlgEl} (\ref{SMTX.AlgEl}). }

 

\subsection{\textcolor{Chapter }{SMTX.AlgElCharPol}}
\logpage{[ 69, 13, 4 ]}\nobreak
\hyperdef{L}{X8417F86A7A20F128}{}
{\noindent\textcolor{FuncColor}{$\triangleright$\enspace\texttt{SMTX.AlgElCharPol({\mdseries\slshape module})\index{SMTX.AlgElCharPol@\texttt{SMTX.AlgElCharPol}}
\label{SMTX.AlgElCharPol}
}\hfill{\scriptsize (function)}}\\


 minimal polynomial of \texttt{SMTX.AlgEl} (\ref{SMTX.AlgEl}). }

 

\subsection{\textcolor{Chapter }{SMTX.AlgElCharPolFac}}
\logpage{[ 69, 13, 5 ]}\nobreak
\hyperdef{L}{X79A82FED785BFB6D}{}
{\noindent\textcolor{FuncColor}{$\triangleright$\enspace\texttt{SMTX.AlgElCharPolFac({\mdseries\slshape module})\index{SMTX.AlgElCharPolFac@\texttt{SMTX.AlgElCharPolFac}}
\label{SMTX.AlgElCharPolFac}
}\hfill{\scriptsize (function)}}\\


 uses factor of \texttt{SMTX.AlgEl} (\ref{SMTX.AlgEl}). }

 

\subsection{\textcolor{Chapter }{SMTX.AlgElNullspaceVec}}
\logpage{[ 69, 13, 6 ]}\nobreak
\hyperdef{L}{X8367B4A17EC39ABD}{}
{\noindent\textcolor{FuncColor}{$\triangleright$\enspace\texttt{SMTX.AlgElNullspaceVec({\mdseries\slshape module})\index{SMTX.AlgElNullspaceVec@\texttt{SMTX.AlgElNullspaceVec}}
\label{SMTX.AlgElNullspaceVec}
}\hfill{\scriptsize (function)}}\\


 nullspace of the matrix evaluated under this factor. }

 

\subsection{\textcolor{Chapter }{SMTX.AlgElNullspaceDimension}}
\logpage{[ 69, 13, 7 ]}\nobreak
\hyperdef{L}{X877F1AB77DC1E12C}{}
{\noindent\textcolor{FuncColor}{$\triangleright$\enspace\texttt{SMTX.AlgElNullspaceDimension({\mdseries\slshape module})\index{SMTX.AlgElNullspaceDimension@\texttt{SMTX.AlgElNullspaceDimension}}
\label{SMTX.AlgElNullspaceDimension}
}\hfill{\scriptsize (function)}}\\


 dimension of the nullspace. }

 

\subsection{\textcolor{Chapter }{SMTX.CentMat}}
\logpage{[ 69, 13, 8 ]}\nobreak
\hyperdef{L}{X78A6B95686671067}{}
{\noindent\textcolor{FuncColor}{$\triangleright$\enspace\texttt{SMTX.CentMat({\mdseries\slshape module})\index{SMTX.CentMat@\texttt{SMTX.CentMat}}
\label{SMTX.CentMat}
}\hfill{\scriptsize (function)}}\\


 matrix centralising all generators which is computed as a byproduct of \texttt{SMTX.AbsoluteIrreducibilityTest} (\ref{SMTX.AbsoluteIrreducibilityTest}). }

 

\subsection{\textcolor{Chapter }{SMTX.CentMatMinPoly}}
\logpage{[ 69, 13, 9 ]}\nobreak
\hyperdef{L}{X7D199DB6804F5D8F}{}
{\noindent\textcolor{FuncColor}{$\triangleright$\enspace\texttt{SMTX.CentMatMinPoly({\mdseries\slshape module})\index{SMTX.CentMatMinPoly@\texttt{SMTX.CentMatMinPoly}}
\label{SMTX.CentMatMinPoly}
}\hfill{\scriptsize (function)}}\\


 minimal polynomial of \texttt{SMTX.CentMat} (\ref{SMTX.CentMat}). }

 }

 }

             
\chapter{\textcolor{Chapter }{Tables of Marks}}\label{Tables of Marks}
\logpage{[ 70, 0, 0 ]}
\hyperdef{L}{X84DBFB8287C3F1B4}{}
{
  The concept of a \emph{table of marks} was introduced by W.{\nobreakspace}Burnside in his book ``Theory of Groups of Finite Order'', see{\nobreakspace}\cite{Bur55}. Therefore a table of marks is sometimes called a \emph{Burnside matrix}. 

 The table of marks of a finite group $G$ is a matrix whose rows and columns are labelled by the conjugacy classes of
subgroups of $G$ and where for two subgroups $A$ and $B$ the $(A, B)$\texttt{\symbol{45}}entry is the number of fixed points of $B$ in the transitive action of $G$ on the cosets of $A$ in $G$. So the table of marks characterizes the set of all permutation
representations of $G$. 

 Moreover, the table of marks gives a compact description of the subgroup
lattice of $G$, since from the numbers of fixed points the numbers of conjugates of a
subgroup $B$ contained in a subgroup $A$ can be derived. 

 A table of marks of a given group $G$ can be constructed from the subgroup lattice of $G$ (see{\nobreakspace}\ref{Constructing Tables of Marks}). For several groups, the table of marks can be restored from the \textsf{GAP} library of tables of marks (see{\nobreakspace}\ref{The Library of Tables of Marks}). 

 Given the table of marks of $G$, one can display it (see{\nobreakspace}\ref{Printing Tables of Marks}) and derive information about $G$ and its Burnside ring from it (see{\nobreakspace}\ref{Attributes of Tables of Marks}, \ref{Properties of Tables of Marks}, \ref{Other Operations for Tables of Marks}). Moreover, tables of marks in \textsf{GAP} provide an easy access to the classes of subgroups of their underlying groups
(see{\nobreakspace}\ref{Accessing Subgroups via Tables of Marks}).    
\section{\textcolor{Chapter }{More about Tables of Marks}}\label{More about Tables of Marks}
\logpage{[ 70, 1, 0 ]}
\hyperdef{L}{X80883EC17968F442}{}
{
  Let $G$ be a finite group with $n$ conjugacy classes of subgroups $C_1, C_2, \ldots, C_n$ and representatives $H_i \in C_i$, $1 \leq i \leq n$. The \emph{table of marks} of $G$ is defined to be the $n \times n$ matrix $M = (m_{ij})$ where the \emph{mark} $m_{ij}$ is the number of fixed points of the subgroup $H_j$ in the action of $G$ on the right cosets of $H_i$ in $G$. 

 Since $H_j$ can only have fixed points if it is contained in a point stabilizer the matrix $M$ is lower triangular if the classes $C_i$ are sorted according to the condition that if $H_i$ is contained in a conjugate of $H_j$ then $i \leq j$. 

 Moreover, the diagonal entries $m_{ii}$ are nonzero since $m_{ii}$ equals the index of $H_i$ in its normalizer in $G$. Hence $M$ is invertible. Since any transitive action of $G$ is equivalent to an action on the cosets of a subgroup of $G$, one sees that the table of marks completely characterizes the set of all
permutation representations of $G$. 

 The marks $m_{ij}$ have further meanings. If $H_1$ is the trivial subgroup of $G$ then each mark $m_{i1}$ in the first column of $M$ is equal to the index of $H_i$ in $G$ since the trivial subgroup fixes all cosets of $H_i$. If $H_n = G$ then each $m_{nj}$ in the last row of $M$ is equal to $1$ since there is only one coset of $G$ in $G$. In general, $m_{ij}$ equals the number of conjugates of $H_i$ containing $H_j$, multiplied by the index of $H_i$ in its normalizer in $G$. Moreover, the number $c_{ij}$ of conjugates of $H_j$ which are contained in $H_i$ can be derived from the marks $m_{ij}$ via the formula 
\[ c_{ij} = ( m_{ij} m_{j1} ) / ( m_{i1} m_{jj} ) \]
. 

 Both the marks $m_{ij}$ and the numbers of subgroups $c_{ij}$ are needed for the functions described in this chapter. 

 A brief survey of properties of tables of marks and a description of
algorithms for the interactive construction of tables of marks using \textsf{GAP} can be found in{\nobreakspace}\cite{Pfe97}. }

  
\section{\textcolor{Chapter }{Table of Marks Objects in GAP}}\label{Table of Marks Objects in GAP}
\logpage{[ 70, 2, 0 ]}
\hyperdef{L}{X7D29539F7C14956D}{}
{
  A table of marks of a group $G$ in \textsf{GAP} is represented by an immutable (see{\nobreakspace}\ref{Mutability and Copyability}) object \mbox{\texttt{\mdseries\slshape tom}} in the category \texttt{IsTableOfMarks} (\ref{IsTableOfMarks}), with defining attributes \texttt{SubsTom} (\ref{SubsTom}) and \texttt{MarksTom} (\ref{MarksTom}). These two attributes encode the matrix of marks in a compressed form. The \texttt{SubsTom} (\ref{SubsTom}) value of \mbox{\texttt{\mdseries\slshape tom}} is a list where for each conjugacy class of subgroups the class numbers of its
subgroups are stored. These are exactly the positions in the corresponding row
of the matrix of marks which have nonzero entries. The marks themselves are
stored via the \texttt{MarksTom} (\ref{MarksTom}) value of \mbox{\texttt{\mdseries\slshape tom}}, which is a list that contains for each entry in \texttt{SubsTom( \mbox{\texttt{\mdseries\slshape tom}} )} the corresponding nonzero value of the table of marks. 

 It is possible to create table of marks objects that do not store a group,
moreover one can create a table of marks object from a matrix of marks
(see{\nobreakspace}\texttt{TableOfMarks} (\ref{TableOfMarks:for a matrix})). So it may happen that a table of marks object in \textsf{GAP} is in fact \emph{not} the table of marks of a group. To some extent, the consistency of a table of
marks object can be checked (see{\nobreakspace}\ref{Other Operations for Tables of Marks}), but \textsf{GAP} knows no general way to prove or disprove that a given matrix of nonnegative
integers is the matrix of marks for a group. Many functions for tables of
marks work well without access to the group {\textendash}this is one of the
arguments why tables of marks are so useful{\textendash}, but for example
normalizers (see{\nobreakspace}\texttt{NormalizerTom} (\ref{NormalizerTom})) and derived subgroups (see{\nobreakspace}\texttt{DerivedSubgroupTom} (\ref{DerivedSubgroupTom})) of subgroups are in general not uniquely determined by the matrix of marks. 

 \textsf{GAP} tables of marks are assumed to be in lower triangular form, that is, if a
subgroup from the conjugacy class corresponding to the $i$\texttt{\symbol{45}}th row is contained in a subgroup from the class
corresponding to the $j$\texttt{\symbol{45}}th row j then $i \leq j$. 

 The \texttt{MarksTom} (\ref{MarksTom}) information can be computed from the values of the attributes \texttt{NrSubsTom} (\ref{NrSubsTom}), \texttt{LengthsTom} (\ref{LengthsTom}), \texttt{OrdersTom} (\ref{OrdersTom}), and \texttt{SubsTom} (\ref{SubsTom}). \texttt{NrSubsTom} (\ref{NrSubsTom}) stores a list containing for each entry in the \texttt{SubsTom} (\ref{SubsTom}) value the corresponding number of conjugates that are contained in a subgroup, \texttt{LengthsTom} (\ref{LengthsTom}) a list containing for each conjugacy class of subgroups its length, and \texttt{OrdersTom} (\ref{OrdersTom}) a list containing for each class of subgroups their order. So the \texttt{MarksTom} (\ref{MarksTom}) value of \mbox{\texttt{\mdseries\slshape tom}} may be missing provided that the values of \texttt{NrSubsTom} (\ref{NrSubsTom}), \texttt{LengthsTom} (\ref{LengthsTom}), and \texttt{OrdersTom} (\ref{OrdersTom}) are stored in \mbox{\texttt{\mdseries\slshape tom}}. 

 Additional information about a table of marks is needed by some functions. The
class numbers of normalizers in $G$ and the number of the derived subgroup of $G$ can be stored via appropriate attributes (see{\nobreakspace}\texttt{NormalizersTom} (\ref{NormalizersTom}), \texttt{DerivedSubgroupTom} (\ref{DerivedSubgroupTom})). 

 If \mbox{\texttt{\mdseries\slshape tom}} stores its group $G$ and a bijection from the rows and columns of the matrix of marks of \mbox{\texttt{\mdseries\slshape tom}} to the classes of subgroups of $G$ then clearly normalizers, derived subgroup etc.{\nobreakspace}can be computed
from this information. But in general a table of marks need not have access to $G$, for example \mbox{\texttt{\mdseries\slshape tom}} might have been constructed from a generic table of marks (see{\nobreakspace}\ref{Generic Construction of Tables of Marks}), or as table of marks of a factor group from a given table of marks
(see{\nobreakspace}\texttt{FactorGroupTom} (\ref{FactorGroupTom})). Access to the group $G$ is provided by the attribute \texttt{UnderlyingGroup} (\ref{UnderlyingGroup:for tables of marks}) if this value is set. Access to the relevant information about conjugacy
classes of subgroups of $G$ {\textendash}compatible with the ordering of rows and columns of the marks in \mbox{\texttt{\mdseries\slshape tom}}{\textendash} is signalled by the filter \texttt{IsTableOfMarksWithGens} (\ref{IsTableOfMarksWithGens}). 

 Several examples in this chapter require the \textsf{GAP} package \textsf{TomLib} (the \textsf{GAP} Library of Tables of Marks) to be available. If it is not yet loaded then we
load it now. 

 
\begin{Verbatim}[commandchars=!@|,fontsize=\small,frame=single,label=Example]
  !gapprompt@gap>| !gapinput@LoadPackage( "tomlib" );|
  true
\end{Verbatim}
 }

  
\section{\textcolor{Chapter }{Constructing Tables of Marks}}\label{Constructing Tables of Marks}
\logpage{[ 70, 3, 0 ]}
\hyperdef{L}{X7B5E4B5F81AF6B00}{}
{
  

\subsection{\textcolor{Chapter }{TableOfMarks (for a group)}}
\logpage{[ 70, 3, 1 ]}\nobreak
\hyperdef{L}{X85B262AB7E219C34}{}
{\noindent\textcolor{FuncColor}{$\triangleright$\enspace\texttt{TableOfMarks({\mdseries\slshape G})\index{TableOfMarks@\texttt{TableOfMarks}!for a group}
\label{TableOfMarks:for a group}
}\hfill{\scriptsize (attribute)}}\\
\noindent\textcolor{FuncColor}{$\triangleright$\enspace\texttt{TableOfMarks({\mdseries\slshape string})\index{TableOfMarks@\texttt{TableOfMarks}!for a string}
\label{TableOfMarks:for a string}
}\hfill{\scriptsize (attribute)}}\\
\noindent\textcolor{FuncColor}{$\triangleright$\enspace\texttt{TableOfMarks({\mdseries\slshape matrix})\index{TableOfMarks@\texttt{TableOfMarks}!for a matrix}
\label{TableOfMarks:for a matrix}
}\hfill{\scriptsize (attribute)}}\\


 In the first form, \mbox{\texttt{\mdseries\slshape G}} must be a finite group, and \texttt{TableOfMarks} constructs the table of marks of \mbox{\texttt{\mdseries\slshape G}}. This computation requires the knowledge of the complete subgroup lattice of \mbox{\texttt{\mdseries\slshape G}} (see{\nobreakspace}\texttt{LatticeSubgroups} (\ref{LatticeSubgroups})). If the lattice is not yet stored then it will be constructed. This may take
a while if \mbox{\texttt{\mdseries\slshape G}} is large. The result has the \texttt{IsTableOfMarksWithGens} (\ref{IsTableOfMarksWithGens}) value \texttt{true}. 

 In the second form, \mbox{\texttt{\mdseries\slshape string}} must be a string, and \texttt{TableOfMarks} gets the table of marks with name \mbox{\texttt{\mdseries\slshape string}} from the \textsf{GAP} library (see \ref{The Library of Tables of Marks}). If no table of marks with this name is contained in the library then \texttt{fail} is returned. 

 In the third form, \mbox{\texttt{\mdseries\slshape matrix}} must be a matrix or a list of rows describing a lower triangular matrix where
the part above the diagonal is omitted. For such an argument \mbox{\texttt{\mdseries\slshape matrix}}, \texttt{TableOfMarks} returns a table of marks object (see{\nobreakspace}\ref{Table of Marks Objects in GAP}) for which \mbox{\texttt{\mdseries\slshape matrix}} is the matrix of marks. Note that not every matrix (containing only
nonnegative integers and having lower triangular shape) describes a table of
marks of a group. Necessary conditions are checked with \texttt{IsInternallyConsistent} (\ref{IsInternallyConsistent:for tables of marks}) (see{\nobreakspace}\ref{Other Operations for Tables of Marks}), and \texttt{fail} is returned if \mbox{\texttt{\mdseries\slshape matrix}} is proved not to describe a matrix of marks; but if \texttt{TableOfMarks} returns a table of marks object created from a matrix then it may still happen
that this object does not describe the table of marks of a group. 

 For an overview of operations for table of marks objects, see the introduction
to Chapter{\nobreakspace}\ref{Tables of Marks}. 

 
\begin{Verbatim}[commandchars=!@|,fontsize=\small,frame=single,label=Example]
  !gapprompt@gap>| !gapinput@tom:= TableOfMarks( AlternatingGroup( 5 ) );|
  TableOfMarks( Alt( [ 1 .. 5 ] ) )
  !gapprompt@gap>| !gapinput@TableOfMarks( "J5" );|
  fail
  !gapprompt@gap>| !gapinput@a5:= TableOfMarks( "A5" );|
  TableOfMarks( "A5" )
  !gapprompt@gap>| !gapinput@mat:=|
  !gapprompt@>| !gapinput@[ [ 60, 0, 0, 0, 0, 0, 0, 0, 0 ], [ 30, 2, 0, 0, 0, 0, 0, 0, 0 ],|
  !gapprompt@>| !gapinput@  [ 20, 0, 2, 0, 0, 0, 0, 0, 0 ], [ 15, 3, 0, 3, 0, 0, 0, 0, 0 ],|
  !gapprompt@>| !gapinput@  [ 12, 0, 0, 0, 2, 0, 0, 0, 0 ], [ 10, 2, 1, 0, 0, 1, 0, 0, 0 ],|
  !gapprompt@>| !gapinput@  [ 6, 2, 0, 0, 1, 0, 1, 0, 0 ], [ 5, 1, 2, 1, 0, 0, 0, 1, 0 ],|
  !gapprompt@>| !gapinput@  [ 1, 1, 1, 1, 1, 1, 1, 1, 1 ] ];;|
  !gapprompt@gap>| !gapinput@TableOfMarks( mat );|
  TableOfMarks( <9 classes> )
\end{Verbatim}
 

 The following \texttt{TableOfMarks} methods for a group are installed. 
\begin{itemize}
\item  If the group is known to be cyclic then \texttt{TableOfMarks} constructs the table of marks essentially without the group, instead the
knowledge about the structure of cyclic groups is used. 
\item  If the lattice of subgroups is already stored in the group then \texttt{TableOfMarks} computes the table of marks from the lattice (see{\nobreakspace}\texttt{TableOfMarksByLattice} (\ref{TableOfMarksByLattice})). 
\item  If the group is known to be solvable then \texttt{TableOfMarks} takes the lattice of subgroups (see{\nobreakspace}\texttt{LatticeSubgroups} (\ref{LatticeSubgroups})) of the group {\textendash}which means that the lattice is computed if it is
not yet stored{\textendash} and then computes the table of marks from it. This
method is also accessible via the global function \texttt{TableOfMarksByLattice} (\ref{TableOfMarksByLattice}). 
\item  If the group doesn't know its lattice of subgroups or its conjugacy classes of
subgroups then the table of marks and the conjugacy classes of subgroups are
computed at the same time by the cyclic extension method. Only the table of
marks is stored because the conjugacy classes of subgroups or the lattice of
subgroups can be easily read off (see{\nobreakspace}\texttt{LatticeSubgroupsByTom} (\ref{LatticeSubgroupsByTom})). 
\end{itemize}
 

 Conversely, the lattice of subgroups of a group with known table of marks can
be computed using the table of marks, via the function \texttt{LatticeSubgroupsByTom} (\ref{LatticeSubgroupsByTom}). This is also installed as a method for \texttt{LatticeSubgroups} (\ref{LatticeSubgroups}). }

 

\subsection{\textcolor{Chapter }{TableOfMarksByLattice}}
\logpage{[ 70, 3, 2 ]}\nobreak
\hyperdef{L}{X7B30FF3A79CCB0DF}{}
{\noindent\textcolor{FuncColor}{$\triangleright$\enspace\texttt{TableOfMarksByLattice({\mdseries\slshape G})\index{TableOfMarksByLattice@\texttt{TableOfMarksByLattice}}
\label{TableOfMarksByLattice}
}\hfill{\scriptsize (function)}}\\


 \texttt{TableOfMarksByLattice} computes the table of marks of the group \mbox{\texttt{\mdseries\slshape G}} from the lattice of subgroups of \mbox{\texttt{\mdseries\slshape G}}. This lattice is computed via \texttt{LatticeSubgroups} (\ref{LatticeSubgroups}) if it is not yet stored in \mbox{\texttt{\mdseries\slshape G}}. The function \texttt{TableOfMarksByLattice} is installed as a method for \texttt{TableOfMarks} (\ref{TableOfMarks:for a group}) for solvable groups and groups with stored subgroup lattice, and is available
as a global variable only in order to provide explicit access to this method. }

 

\subsection{\textcolor{Chapter }{LatticeSubgroupsByTom}}
\logpage{[ 70, 3, 3 ]}\nobreak
\hyperdef{L}{X79ABFA0A833DDCFE}{}
{\noindent\textcolor{FuncColor}{$\triangleright$\enspace\texttt{LatticeSubgroupsByTom({\mdseries\slshape G})\index{LatticeSubgroupsByTom@\texttt{LatticeSubgroupsByTom}}
\label{LatticeSubgroupsByTom}
}\hfill{\scriptsize (function)}}\\


 \texttt{LatticeSubgroupsByTom} computes the lattice of subgroups of \mbox{\texttt{\mdseries\slshape G}} from the table of marks of \mbox{\texttt{\mdseries\slshape G}}, using \texttt{RepresentativeTom} (\ref{RepresentativeTom}). }

 }

  
\section{\textcolor{Chapter }{Printing Tables of Marks}}\label{Printing Tables of Marks}
\logpage{[ 70, 4, 0 ]}
\hyperdef{L}{X7AC0FB9685DCBCFD}{}
{
  

\subsection{\textcolor{Chapter }{ViewObj (for a table of marks)}}
\logpage{[ 70, 4, 1 ]}\nobreak
\hyperdef{L}{X7DC656517D8335DC}{}
{\noindent\textcolor{FuncColor}{$\triangleright$\enspace\texttt{ViewObj({\mdseries\slshape tom})\index{ViewObj@\texttt{ViewObj}!for a table of marks}
\label{ViewObj:for a table of marks}
}\hfill{\scriptsize (method)}}\\


 The default \texttt{ViewObj} (\ref{ViewObj}) method for tables of marks prints the string \texttt{"TableOfMarks"}, followed by {\textendash}if known{\textendash} the identifier
(see{\nobreakspace}\texttt{Identifier} (\ref{Identifier:for tables of marks})) or the group of the table of marks enclosed in brackets; if neither group
nor identifier are known then just the number of conjugacy classes of
subgroups is printed instead. }

 

\subsection{\textcolor{Chapter }{PrintObj (for a table of marks)}}
\logpage{[ 70, 4, 2 ]}\nobreak
\hyperdef{L}{X86379C0D7D17AD92}{}
{\noindent\textcolor{FuncColor}{$\triangleright$\enspace\texttt{PrintObj({\mdseries\slshape tom})\index{PrintObj@\texttt{PrintObj}!for a table of marks}
\label{PrintObj:for a table of marks}
}\hfill{\scriptsize (method)}}\\


 The default \texttt{PrintObj} (\ref{PrintObj}) method for tables of marks does the same as \texttt{ViewObj} (\ref{ViewObj}), except that \texttt{PrintObj} (\ref{PrintObj}) is used for the group instead of \texttt{ViewObj} (\ref{ViewObj}). }

 

\subsection{\textcolor{Chapter }{Display (for a table of marks)}}
\logpage{[ 70, 4, 3 ]}\nobreak
\hyperdef{L}{X821F9438839F445D}{}
{\noindent\textcolor{FuncColor}{$\triangleright$\enspace\texttt{Display({\mdseries\slshape tom[, arec]})\index{Display@\texttt{Display}!for a table of marks}
\label{Display:for a table of marks}
}\hfill{\scriptsize (method)}}\\


 The default \texttt{Display} (\ref{Display}) method for a table of marks \mbox{\texttt{\mdseries\slshape tom}} produces a formatted output of the marks in \mbox{\texttt{\mdseries\slshape tom}}. Each line of output begins with the number of the corresponding class of
subgroups. This number is repeated if the output spreads over several pages.
The number of columns printed at one time depends on the actual line length,
which can be accessed and changed by the function \texttt{SizeScreen} (\ref{SizeScreen}). 

 An interactive variant of \texttt{Display} (\ref{Display}) is the \texttt{Browse} (\textbf{Browse: Browse}) method for tables of marks that is provided by the \textsf{GAP} package \textsf{Browse}, see \texttt{Browse} (\textbf{Browse: Browse for tables of marks}). 

 The optional second argument \mbox{\texttt{\mdseries\slshape arec}} of \texttt{Display} (\ref{Display}) can be used to change the default style for displaying a table of marks. \mbox{\texttt{\mdseries\slshape arec}} must be a record, its relevant components are the following. 

 
\begin{description}
\item[{\texttt{classes}}]  a list of class numbers to select only the rows and columns of the matrix that
correspond to this list for printing, 
\item[{\texttt{form}}]  one of the strings \texttt{"subgroups"}, \texttt{"supergroups"}; in the former case, at position $(i,j)$ of the matrix the number of conjugates of $H_j$ contained in $H_i$ is printed, and in the latter case, at position $(i,j)$ the number of conjugates of $H_i$ which contain $H_j$ is printed. 
\end{description}
 

 
\begin{Verbatim}[commandchars=!@|,fontsize=\small,frame=single,label=Example]
  !gapprompt@gap>| !gapinput@tom:= TableOfMarks( "A5" );;|
  !gapprompt@gap>| !gapinput@Display( tom );|
  1:  60
  2:  30 2
  3:  20 . 2
  4:  15 3 . 3
  5:  12 . . . 2
  6:  10 2 1 . . 1
  7:   6 2 . . 1 . 1
  8:   5 1 2 1 . . . 1
  9:   1 1 1 1 1 1 1 1 1
  
  !gapprompt@gap>| !gapinput@Display( tom, rec( classes:= [ 1, 2, 3, 4, 8 ] ) );|
  1:  60
  2:  30 2
  3:  20 . 2
  4:  15 3 . 3
  8:   5 1 2 1 1
  
  !gapprompt@gap>| !gapinput@Display( tom, rec( form:= "subgroups" ) );|
  1:  1
  2:  1  1
  3:  1  .  1
  4:  1  3  . 1
  5:  1  .  . . 1
  6:  1  3  1 . .  1
  7:  1  5  . . 1  . 1
  8:  1  3  4 1 .  . . 1
  9:  1 15 10 5 6 10 6 5 1
  
  !gapprompt@gap>| !gapinput@Display( tom, rec( form:= "supergroups" ) );|
  1:   1
  2:  15 1
  3:  10 . 1
  4:   5 1 . 1
  5:   6 . . . 1
  6:  10 2 1 . . 1
  7:   6 2 . . 1 . 1
  8:   5 1 2 1 . . . 1
  9:   1 1 1 1 1 1 1 1 1
  
\end{Verbatim}
 }

 }

  
\section{\textcolor{Chapter }{Sorting Tables of Marks}}\label{Sorting Tables of Marks}
\logpage{[ 70, 5, 0 ]}
\hyperdef{L}{X82385925797B5108}{}
{
  

\subsection{\textcolor{Chapter }{SortedTom}}
\logpage{[ 70, 5, 1 ]}\nobreak
\hyperdef{L}{X786A948E82C36F0E}{}
{\noindent\textcolor{FuncColor}{$\triangleright$\enspace\texttt{SortedTom({\mdseries\slshape tom, perm})\index{SortedTom@\texttt{SortedTom}}
\label{SortedTom}
}\hfill{\scriptsize (operation)}}\\


 \texttt{SortedTom} returns a table of marks where the rows and columns of the table of marks \mbox{\texttt{\mdseries\slshape tom}} are reordered according to the permutation \mbox{\texttt{\mdseries\slshape perm}}. 

 \emph{Note} that in each table of marks in \textsf{GAP}, the matrix of marks is assumed to have lower triangular shape
(see{\nobreakspace}\ref{Table of Marks Objects in GAP}). If the permutation \mbox{\texttt{\mdseries\slshape perm}} does \emph{not} have this property then the functions for tables of marks might return wrong
results when applied to the output of \texttt{SortedTom}. 

 The returned table of marks has only those attribute values stored that are
known for \mbox{\texttt{\mdseries\slshape tom}} and listed in \texttt{TableOfMarksComponents} (\ref{TableOfMarksComponents}). 

 
\begin{Verbatim}[commandchars=!@|,fontsize=\small,frame=single,label=Example]
  !gapprompt@gap>| !gapinput@tom:= TableOfMarksCyclic( 6 );;  Display( tom );|
  1:  6
  2:  3 3
  3:  2 . 2
  4:  1 1 1 1
  
  !gapprompt@gap>| !gapinput@sorted:= SortedTom( tom, (2,3) );;  Display( sorted );|
  1:  6
  2:  2 2
  3:  3 . 3
  4:  1 1 1 1
  
  !gapprompt@gap>| !gapinput@wrong:= SortedTom( tom, (1,2) );;  Display( wrong );|
  1:  3
  2:  . 6
  3:  . 2 2
  4:  1 1 1 1
  
\end{Verbatim}
 }

 

\subsection{\textcolor{Chapter }{PermutationTom}}
\logpage{[ 70, 5, 2 ]}\nobreak
\hyperdef{L}{X7EFD937D804662F6}{}
{\noindent\textcolor{FuncColor}{$\triangleright$\enspace\texttt{PermutationTom({\mdseries\slshape tom})\index{PermutationTom@\texttt{PermutationTom}}
\label{PermutationTom}
}\hfill{\scriptsize (attribute)}}\\


 For the table of marks \mbox{\texttt{\mdseries\slshape tom}} of the group $G$ stored as \texttt{UnderlyingGroup} (\ref{UnderlyingGroup:for tables of marks}) value of \mbox{\texttt{\mdseries\slshape tom}}, \texttt{PermutationTom} is a permutation $\pi$ such that the $i$\texttt{\symbol{45}}th conjugacy class of subgroups of $G$ belongs to the $i^\pi$\texttt{\symbol{45}}th column and row of marks in \mbox{\texttt{\mdseries\slshape tom}}. 

 This attribute value is bound only if \mbox{\texttt{\mdseries\slshape tom}} was obtained from another table of marks by permuting with \texttt{SortedTom} (\ref{SortedTom}), and there is no default method to compute its value. 

 The attribute is necessary because the original and the sorted table of marks
have the same identifier and the same group, and information computed from the
group may depend on the ordering of marks, for example the fusion from the
ordinary character table of $G$ into \mbox{\texttt{\mdseries\slshape tom}}. 

 
\begin{Verbatim}[commandchars=!@|,fontsize=\small,frame=single,label=Example]
  !gapprompt@gap>| !gapinput@MarksTom( tom )[2];|
  [ 3, 3 ]
  !gapprompt@gap>| !gapinput@MarksTom( sorted )[2];|
  [ 2, 2 ]
  !gapprompt@gap>| !gapinput@HasPermutationTom( sorted );|
  true
  !gapprompt@gap>| !gapinput@PermutationTom( sorted );|
  (2,3)
\end{Verbatim}
 }

 }

  
\section{\textcolor{Chapter }{Technical Details about Tables of Marks}}\label{Technical Details about Tables of Marks}
\logpage{[ 70, 6, 0 ]}
\hyperdef{L}{X82271C4F7FD21FAA}{}
{
  

\subsection{\textcolor{Chapter }{InfoTom}}
\logpage{[ 70, 6, 1 ]}\nobreak
\hyperdef{L}{X870985C58547FED4}{}
{\noindent\textcolor{FuncColor}{$\triangleright$\enspace\texttt{InfoTom\index{InfoTom@\texttt{InfoTom}}
\label{InfoTom}
}\hfill{\scriptsize (info class)}}\\


 is the info class for computations concerning tables of marks. }

 

\subsection{\textcolor{Chapter }{IsTableOfMarks}}
\logpage{[ 70, 6, 2 ]}\nobreak
\hyperdef{L}{X7AC1A73D8100C7EC}{}
{\noindent\textcolor{FuncColor}{$\triangleright$\enspace\texttt{IsTableOfMarks({\mdseries\slshape obj})\index{IsTableOfMarks@\texttt{IsTableOfMarks}}
\label{IsTableOfMarks}
}\hfill{\scriptsize (Category)}}\\


 Each table of marks belongs to this category. }

 

\subsection{\textcolor{Chapter }{TableOfMarksFamily}}
\logpage{[ 70, 6, 3 ]}\nobreak
\hyperdef{L}{X7ACF943D84BDF89E}{}
{\noindent\textcolor{FuncColor}{$\triangleright$\enspace\texttt{TableOfMarksFamily\index{TableOfMarksFamily@\texttt{TableOfMarksFamily}}
\label{TableOfMarksFamily}
}\hfill{\scriptsize (family)}}\\


 Each table of marks belongs to this family. }

 

\subsection{\textcolor{Chapter }{TableOfMarksComponents}}
\logpage{[ 70, 6, 4 ]}\nobreak
\hyperdef{L}{X87789FD27831B2A2}{}
{\noindent\textcolor{FuncColor}{$\triangleright$\enspace\texttt{TableOfMarksComponents\index{TableOfMarksComponents@\texttt{TableOfMarksComponents}}
\label{TableOfMarksComponents}
}\hfill{\scriptsize (global variable)}}\\


 The list \texttt{TableOfMarksComponents} is used when a table of marks object is created from a record via \texttt{ConvertToTableOfMarks} (\ref{ConvertToTableOfMarks}). \texttt{TableOfMarksComponents} contains at position $2i-1$ a name of an attribute and at position $2i$ the corresponding attribute getter function. }

 

\subsection{\textcolor{Chapter }{ConvertToTableOfMarks}}
\logpage{[ 70, 6, 5 ]}\nobreak
\hyperdef{L}{X8491CDBF8543A7D5}{}
{\noindent\textcolor{FuncColor}{$\triangleright$\enspace\texttt{ConvertToTableOfMarks({\mdseries\slshape record})\index{ConvertToTableOfMarks@\texttt{ConvertToTableOfMarks}}
\label{ConvertToTableOfMarks}
}\hfill{\scriptsize (function)}}\\


 \texttt{ConvertToTableOfMarks} converts a record with components from \texttt{TableOfMarksComponents} (\ref{TableOfMarksComponents}) into a table of marks object with the corresponding attributes. 

 
\begin{Verbatim}[commandchars=!@|,fontsize=\small,frame=single,label=Example]
  !gapprompt@gap>| !gapinput@record:= rec( MarksTom:= [ [ 4 ], [ 2, 2 ], [ 1, 1, 1 ] ],|
  !gapprompt@>| !gapinput@ SubsTom:= [ [ 1 ], [ 1, 2 ], [ 1, 2, 3 ] ] );;|
  !gapprompt@gap>| !gapinput@ConvertToTableOfMarks( record );;|
  !gapprompt@gap>| !gapinput@record;|
  TableOfMarks( <3 classes> )
\end{Verbatim}
 }

 }

  
\section{\textcolor{Chapter }{Attributes of Tables of Marks}}\label{Attributes of Tables of Marks}
\logpage{[ 70, 7, 0 ]}
\hyperdef{L}{X838D3B87827D6923}{}
{
  

\subsection{\textcolor{Chapter }{MarksTom}}
\logpage{[ 70, 7, 1 ]}\nobreak
\hyperdef{L}{X78F486A28561D006}{}
{\noindent\textcolor{FuncColor}{$\triangleright$\enspace\texttt{MarksTom({\mdseries\slshape tom})\index{MarksTom@\texttt{MarksTom}}
\label{MarksTom}
}\hfill{\scriptsize (attribute)}}\\
\noindent\textcolor{FuncColor}{$\triangleright$\enspace\texttt{SubsTom({\mdseries\slshape tom})\index{SubsTom@\texttt{SubsTom}}
\label{SubsTom}
}\hfill{\scriptsize (attribute)}}\\


 The matrix of marks (see{\nobreakspace}\ref{More about Tables of Marks}) of the table of marks \mbox{\texttt{\mdseries\slshape tom}} is stored in a compressed form where zeros are omitted, using the attributes \texttt{MarksTom} and \texttt{SubsTom}. If $M$ is the square matrix of marks of \mbox{\texttt{\mdseries\slshape tom}} (see{\nobreakspace}\texttt{MatTom} (\ref{MatTom})) then the \texttt{SubsTom} value of \mbox{\texttt{\mdseries\slshape tom}} is a list that contains at position $i$ the list of all positions of nonzero entries of the $i$\texttt{\symbol{45}}th row of $M$, and the \texttt{MarksTom} value of \mbox{\texttt{\mdseries\slshape tom}} is a list that contains at position $i$ the list of the corresponding marks. 

 \texttt{MarksTom} and \texttt{SubsTom} are defining attributes of tables of marks (see{\nobreakspace}\ref{Table of Marks Objects in GAP}). There is no default method for computing the \texttt{SubsTom} value, and the default \texttt{MarksTom} method needs the values of \texttt{NrSubsTom} (\ref{NrSubsTom}) and \texttt{OrdersTom} (\ref{OrdersTom}). 

 
\begin{Verbatim}[commandchars=!@|,fontsize=\small,frame=single,label=Example]
  !gapprompt@gap>| !gapinput@a5:= TableOfMarks( "A5" );|
  TableOfMarks( "A5" )
  !gapprompt@gap>| !gapinput@MarksTom( a5 );|
  [ [ 60 ], [ 30, 2 ], [ 20, 2 ], [ 15, 3, 3 ], [ 12, 2 ],
    [ 10, 2, 1, 1 ], [ 6, 2, 1, 1 ], [ 5, 1, 2, 1, 1 ],
    [ 1, 1, 1, 1, 1, 1, 1, 1, 1 ] ]
  !gapprompt@gap>| !gapinput@SubsTom( a5 );|
  [ [ 1 ], [ 1, 2 ], [ 1, 3 ], [ 1, 2, 4 ], [ 1, 5 ], [ 1, 2, 3, 6 ],
    [ 1, 2, 5, 7 ], [ 1, 2, 3, 4, 8 ], [ 1, 2, 3, 4, 5, 6, 7, 8, 9 ] ]
\end{Verbatim}
 }

 

\subsection{\textcolor{Chapter }{NrSubsTom}}
\logpage{[ 70, 7, 2 ]}\nobreak
\hyperdef{L}{X82E5DA217A5D1134}{}
{\noindent\textcolor{FuncColor}{$\triangleright$\enspace\texttt{NrSubsTom({\mdseries\slshape tom})\index{NrSubsTom@\texttt{NrSubsTom}}
\label{NrSubsTom}
}\hfill{\scriptsize (attribute)}}\\
\noindent\textcolor{FuncColor}{$\triangleright$\enspace\texttt{OrdersTom({\mdseries\slshape tom})\index{OrdersTom@\texttt{OrdersTom}}
\label{OrdersTom}
}\hfill{\scriptsize (attribute)}}\\


 Instead of storing the marks (see{\nobreakspace}\texttt{MarksTom} (\ref{MarksTom})) of the table of marks \mbox{\texttt{\mdseries\slshape tom}} one can use a matrix which contains at position $(i,j)$ the number of subgroups of conjugacy class $j$ that are contained in one member of the conjugacy class $i$. These values are stored in the \texttt{NrSubsTom} value in the same way as the marks in the \texttt{MarksTom} (\ref{MarksTom}) value. 

 \texttt{OrdersTom} returns a list that contains at position $i$ the order of a representative of the $i$\texttt{\symbol{45}}th conjugacy class of subgroups of \mbox{\texttt{\mdseries\slshape tom}}. 

 One can compute the \texttt{NrSubsTom} and \texttt{OrdersTom} values from the \texttt{MarksTom} (\ref{MarksTom}) value of \mbox{\texttt{\mdseries\slshape tom}} and vice versa. 

 
\begin{Verbatim}[commandchars=!@|,fontsize=\small,frame=single,label=Example]
  !gapprompt@gap>| !gapinput@NrSubsTom( a5 );|
  [ [ 1 ], [ 1, 1 ], [ 1, 1 ], [ 1, 3, 1 ], [ 1, 1 ], [ 1, 3, 1, 1 ],
    [ 1, 5, 1, 1 ], [ 1, 3, 4, 1, 1 ], [ 1, 15, 10, 5, 6, 10, 6, 5, 1 ]
   ]
  !gapprompt@gap>| !gapinput@OrdersTom( a5 );|
  [ 1, 2, 3, 4, 5, 6, 10, 12, 60 ]
\end{Verbatim}
 }

 

\subsection{\textcolor{Chapter }{LengthsTom}}
\logpage{[ 70, 7, 3 ]}\nobreak
\hyperdef{L}{X781AA1B28178AE9A}{}
{\noindent\textcolor{FuncColor}{$\triangleright$\enspace\texttt{LengthsTom({\mdseries\slshape tom})\index{LengthsTom@\texttt{LengthsTom}}
\label{LengthsTom}
}\hfill{\scriptsize (attribute)}}\\


 For a table of marks \mbox{\texttt{\mdseries\slshape tom}}, \texttt{LengthsTom} returns a list of the lengths of the conjugacy classes of subgroups. 

 
\begin{Verbatim}[commandchars=!@|,fontsize=\small,frame=single,label=Example]
  !gapprompt@gap>| !gapinput@LengthsTom( a5 );|
  [ 1, 15, 10, 5, 6, 10, 6, 5, 1 ]
\end{Verbatim}
 }

 

\subsection{\textcolor{Chapter }{ClassTypesTom}}
\logpage{[ 70, 7, 4 ]}\nobreak
\hyperdef{L}{X7A33C7C38083CC09}{}
{\noindent\textcolor{FuncColor}{$\triangleright$\enspace\texttt{ClassTypesTom({\mdseries\slshape tom})\index{ClassTypesTom@\texttt{ClassTypesTom}}
\label{ClassTypesTom}
}\hfill{\scriptsize (attribute)}}\\


 \texttt{ClassTypesTom} distinguishes isomorphism types of the classes of subgroups of the table of
marks \mbox{\texttt{\mdseries\slshape tom}} as far as this is possible from the \texttt{SubsTom} (\ref{SubsTom}) and \texttt{MarksTom} (\ref{MarksTom}) values of \mbox{\texttt{\mdseries\slshape tom}}. 

 Two subgroups are clearly not isomorphic if they have different orders.
Moreover, isomorphic subgroups must contain the same number of subgroups of
each type. 

 Each type is represented by a positive integer. \texttt{ClassTypesTom} returns the list which contains for each class of subgroups its corresponding
type. 

 
\begin{Verbatim}[commandchars=!@|,fontsize=\small,frame=single,label=Example]
  !gapprompt@gap>| !gapinput@a6:= TableOfMarks( "A6" );;|
  !gapprompt@gap>| !gapinput@ClassTypesTom( a6 );|
  [ 1, 2, 3, 3, 4, 5, 6, 6, 7, 7, 8, 9, 10, 11, 11, 12, 13, 13, 14, 15,
    15, 16 ]
\end{Verbatim}
 }

 

\subsection{\textcolor{Chapter }{ClassNamesTom}}
\logpage{[ 70, 7, 5 ]}\nobreak
\hyperdef{L}{X7A53E923819FE173}{}
{\noindent\textcolor{FuncColor}{$\triangleright$\enspace\texttt{ClassNamesTom({\mdseries\slshape tom})\index{ClassNamesTom@\texttt{ClassNamesTom}}
\label{ClassNamesTom}
}\hfill{\scriptsize (attribute)}}\\


 \texttt{ClassNamesTom} constructs generic names for the conjugacy classes of subgroups of the table
of marks \mbox{\texttt{\mdseries\slshape tom}}. In general, the generic name of a class of non\texttt{\symbol{45}}cyclic
subgroups consists of three parts and has the form \texttt{"(}\mbox{\texttt{\mdseries\slshape o}}\texttt{){\textunderscore}\texttt{\symbol{123}}}\mbox{\texttt{\mdseries\slshape t}}\texttt{\texttt{\symbol{125}}}\mbox{\texttt{\mdseries\slshape l}}\texttt{"}, where \mbox{\texttt{\mdseries\slshape o}} indicates the order of the subgroup, \mbox{\texttt{\mdseries\slshape t}} is a number that distinguishes different types of subgroups of the same order,
and \mbox{\texttt{\mdseries\slshape l}} is a letter that distinguishes classes of subgroups of the same type and
order. The type of a subgroup is determined by the numbers of its subgroups of
other types (see{\nobreakspace}\texttt{ClassTypesTom} (\ref{ClassTypesTom})). This is slightly weaker than isomorphism. 

 The letter is omitted if there is only one class of subgroups of that order
and type, and the type is omitted if there is only one class of that order.
Moreover, the braces \texttt{\texttt{\symbol{123}}\texttt{\symbol{125}}} around the type are omitted if the type number has only one digit. 

 For classes of cyclic subgroups, the parentheses round the order and the type
are omitted. Hence the most general form of their generic names is \texttt{"\mbox{\texttt{\mdseries\slshape o}},\mbox{\texttt{\mdseries\slshape l}}"}. Again, the letter is omitted if there is only one class of cyclic subgroups
of that order. 

 
\begin{Verbatim}[commandchars=!@|,fontsize=\small,frame=single,label=Example]
  !gapprompt@gap>| !gapinput@ClassNamesTom( a6 );|
  [ "1", "2", "3a", "3b", "5", "4", "(4)_2a", "(4)_2b", "(6)a", "(6)b",
    "(8)", "(9)", "(10)", "(12)a", "(12)b", "(18)", "(24)a", "(24)b",
    "(36)", "(60)a", "(60)b", "(360)" ]
\end{Verbatim}
 }

 

\subsection{\textcolor{Chapter }{FusionsTom}}
\logpage{[ 70, 7, 6 ]}\nobreak
\hyperdef{L}{X86B9891C788D5107}{}
{\noindent\textcolor{FuncColor}{$\triangleright$\enspace\texttt{FusionsTom({\mdseries\slshape tom})\index{FusionsTom@\texttt{FusionsTom}}
\label{FusionsTom}
}\hfill{\scriptsize (attribute)}}\\


 For a table of marks \mbox{\texttt{\mdseries\slshape tom}}, \texttt{FusionsTom} is a list of fusions into other tables of marks. Each fusion is a list of
length two, the first entry being the \texttt{Identifier} (\ref{Identifier:for tables of marks})) value of the image table, the second entry being the list of images of the
class positions of \mbox{\texttt{\mdseries\slshape tom}} in the image table. 

 This attribute is mainly used for tables of marks in the \textsf{GAP} library (see{\nobreakspace}\ref{The Library of Tables of Marks}). 

 
\begin{Verbatim}[commandchars=!@|,fontsize=\small,frame=single,label=Example]
  !gapprompt@gap>| !gapinput@fus:= FusionsTom( a6 );;|
  !gapprompt@gap>| !gapinput@fus[1];|
  [ "L3(4)",
    [ 1, 2, 3, 3, 14, 5, 9, 7, 15, 15, 24, 26, 27, 32, 33, 50, 57, 55,
        63, 73, 77, 90 ] ]
\end{Verbatim}
 }

 

\subsection{\textcolor{Chapter }{UnderlyingGroup (for tables of marks)}}
\logpage{[ 70, 7, 7 ]}\nobreak
\hyperdef{L}{X81E41D3880FA6C4C}{}
{\noindent\textcolor{FuncColor}{$\triangleright$\enspace\texttt{UnderlyingGroup({\mdseries\slshape tom})\index{UnderlyingGroup@\texttt{UnderlyingGroup}!for tables of marks}
\label{UnderlyingGroup:for tables of marks}
}\hfill{\scriptsize (attribute)}}\\


 \texttt{UnderlyingGroup} is used to access an underlying group that is stored on the table of marks \mbox{\texttt{\mdseries\slshape tom}}. There is no default method to compute an underlying group if it is not
stored. 

 
\begin{Verbatim}[commandchars=!@|,fontsize=\small,frame=single,label=Example]
  !gapprompt@gap>| !gapinput@UnderlyingGroup( a6 );|
  Group([ (1,2)(3,4), (1,2,4,5)(3,6) ])
\end{Verbatim}
 }

 

\subsection{\textcolor{Chapter }{IdempotentsTom}}
\logpage{[ 70, 7, 8 ]}\nobreak
\hyperdef{L}{X817238FB79A3462F}{}
{\noindent\textcolor{FuncColor}{$\triangleright$\enspace\texttt{IdempotentsTom({\mdseries\slshape tom})\index{IdempotentsTom@\texttt{IdempotentsTom}}
\label{IdempotentsTom}
}\hfill{\scriptsize (attribute)}}\\
\noindent\textcolor{FuncColor}{$\triangleright$\enspace\texttt{IdempotentsTomInfo({\mdseries\slshape tom})\index{IdempotentsTomInfo@\texttt{IdempotentsTomInfo}}
\label{IdempotentsTomInfo}
}\hfill{\scriptsize (attribute)}}\\


 \texttt{IdempotentsTom} encodes the idempotents of the integral Burnside ring described by the table
of marks \mbox{\texttt{\mdseries\slshape tom}}. The return value is a list $l$ of positive integers such that each row vector describing a primitive
idempotent has value $1$ at all positions with the same entry in $l$, and $0$ at all other positions. 

 According to A.{\nobreakspace}Dress{\nobreakspace}\cite{Dre69} (see also{\nobreakspace}\cite{Pfe97}), these idempotents correspond to the classes of perfect subgroups, and each
such idempotent is the characteristic function of all those subgroups that
arise by cyclic extension from the corresponding perfect subgroup
(see{\nobreakspace}\texttt{CyclicExtensionsTom} (\ref{CyclicExtensionsTom:for a prime})). 

 \texttt{IdempotentsTomInfo} returns a record with components \texttt{fixpointvectors} and \texttt{primidems}, both bound to lists. The $i$\texttt{\symbol{45}}th entry of the \texttt{fixpointvectors} list is the $0-1$\texttt{\symbol{45}}vector describing the $i$\texttt{\symbol{45}}th primitive idempotent, and the $i$\texttt{\symbol{45}}th entry of \texttt{primidems} is the decomposition of this idempotent in the rows of \mbox{\texttt{\mdseries\slshape tom}}. 

 
\begin{Verbatim}[commandchars=!@|,fontsize=\small,frame=single,label=Example]
  !gapprompt@gap>| !gapinput@IdempotentsTom( a5 );|
  [ 1, 1, 1, 1, 1, 1, 1, 1, 9 ]
  !gapprompt@gap>| !gapinput@IdempotentsTomInfo( a5 );|
  rec(
    fixpointvectors := [ [ 1, 1, 1, 1, 1, 1, 1, 1, 0 ],
        [ 0, 0, 0, 0, 0, 0, 0, 0, 1 ] ],
    primidems := [ [ 1, -2, -1, 0, 0, 1, 1, 1 ],
        [ -1, 2, 1, 0, 0, -1, -1, -1, 1 ] ] )
\end{Verbatim}
 }

 

\subsection{\textcolor{Chapter }{Identifier (for tables of marks)}}
\logpage{[ 70, 7, 9 ]}\nobreak
\hyperdef{L}{X810E53597B5BB4F8}{}
{\noindent\textcolor{FuncColor}{$\triangleright$\enspace\texttt{Identifier({\mdseries\slshape tom})\index{Identifier@\texttt{Identifier}!for tables of marks}
\label{Identifier:for tables of marks}
}\hfill{\scriptsize (attribute)}}\\


 The identifier of a table of marks \mbox{\texttt{\mdseries\slshape tom}} is a string. It is used for printing the table of marks (see{\nobreakspace}\ref{Printing Tables of Marks}) and in fusions between tables of marks (see{\nobreakspace}\texttt{FusionsTom} (\ref{FusionsTom})). 

 If \mbox{\texttt{\mdseries\slshape tom}} is a table of marks from the \textsf{GAP} library of tables of marks (see{\nobreakspace}\ref{The Library of Tables of Marks}) then it has an identifier, and if \mbox{\texttt{\mdseries\slshape tom}} was constructed from a group with \texttt{Name} (\ref{Name}) then this name is chosen as \texttt{Identifier} value. There is no default method to compute an identifier in all other cases. 

 
\begin{Verbatim}[commandchars=!@|,fontsize=\small,frame=single,label=Example]
  !gapprompt@gap>| !gapinput@Identifier( a5 );|
  "A5"
\end{Verbatim}
 }

 

\subsection{\textcolor{Chapter }{MatTom}}
\logpage{[ 70, 7, 10 ]}\nobreak
\hyperdef{L}{X8463272986781E17}{}
{\noindent\textcolor{FuncColor}{$\triangleright$\enspace\texttt{MatTom({\mdseries\slshape tom})\index{MatTom@\texttt{MatTom}}
\label{MatTom}
}\hfill{\scriptsize (attribute)}}\\


 \texttt{MatTom} returns the square matrix of marks (see{\nobreakspace}\ref{More about Tables of Marks}) of the table of marks \mbox{\texttt{\mdseries\slshape tom}} which is stored in a compressed form using the attributes \texttt{MarksTom} (\ref{MarksTom}) and \texttt{SubsTom} (\ref{SubsTom}) This may need substantially more space than the values of \texttt{MarksTom} (\ref{MarksTom}) and \texttt{SubsTom} (\ref{SubsTom}). 

 
\begin{Verbatim}[commandchars=!@|,fontsize=\small,frame=single,label=Example]
  !gapprompt@gap>| !gapinput@MatTom( a5 );|
  [ [ 60, 0, 0, 0, 0, 0, 0, 0, 0 ], [ 30, 2, 0, 0, 0, 0, 0, 0, 0 ],
    [ 20, 0, 2, 0, 0, 0, 0, 0, 0 ], [ 15, 3, 0, 3, 0, 0, 0, 0, 0 ],
    [ 12, 0, 0, 0, 2, 0, 0, 0, 0 ], [ 10, 2, 1, 0, 0, 1, 0, 0, 0 ],
    [ 6, 2, 0, 0, 1, 0, 1, 0, 0 ], [ 5, 1, 2, 1, 0, 0, 0, 1, 0 ],
    [ 1, 1, 1, 1, 1, 1, 1, 1, 1 ] ]
\end{Verbatim}
 }

 

\subsection{\textcolor{Chapter }{MoebiusTom}}
\logpage{[ 70, 7, 11 ]}\nobreak
\hyperdef{L}{X7D32C8B0786D16C1}{}
{\noindent\textcolor{FuncColor}{$\triangleright$\enspace\texttt{MoebiusTom({\mdseries\slshape tom})\index{MoebiusTom@\texttt{MoebiusTom}}
\label{MoebiusTom}
}\hfill{\scriptsize (attribute)}}\\


 \texttt{MoebiusTom} computes the M{\"o}bius values both of the subgroup lattice of the group $G$ with table of marks \mbox{\texttt{\mdseries\slshape tom}} and of the poset of conjugacy classes of subgroups of $G$. It returns a record where the component \texttt{mu} contains the M{\"o}bius values of the subgroup lattice, and the component \texttt{nu} contains the M{\"o}bius values of the poset. 

 Moreover, according to an observation of Isaacs et al. (see{\nobreakspace}\cite{HIO89}, \cite{Pah93}), the values on the subgroup lattice often can be derived from those of the
poset of conjugacy classes. These ``expected values'' are returned in the component \texttt{ex}, and the list of numbers of those subgroups where the expected value does not
coincide with the actual value are returned in the component \texttt{hyp}. For the computation of these values, the position of the derived subgroup of $G$ is needed (see{\nobreakspace}\texttt{DerivedSubgroupTom} (\ref{DerivedSubgroupTom})). If it is not uniquely determined then the result does not have the
components \texttt{ex} and \texttt{hyp}. 

 
\begin{Verbatim}[commandchars=!@|,fontsize=\small,frame=single,label=Example]
  !gapprompt@gap>| !gapinput@MoebiusTom( a5 );|
  rec( ex := [ -60, 4, 2,,, -1, -1, -1, 1 ], hyp := [  ],
    mu := [ -60, 4, 2,,, -1, -1, -1, 1 ],
    nu := [ -1, 2, 1,,, -1, -1, -1, 1 ] )
  !gapprompt@gap>| !gapinput@tom:= TableOfMarks( "M12" );;|
  !gapprompt@gap>| !gapinput@moebius:= MoebiusTom( tom );;|
  !gapprompt@gap>| !gapinput@moebius.hyp;|
  [ 1, 2, 4, 16, 39, 45, 105 ]
  !gapprompt@gap>| !gapinput@moebius.mu[1];  moebius.ex[1];|
  95040
  190080
\end{Verbatim}
 }

 

\subsection{\textcolor{Chapter }{WeightsTom}}
\logpage{[ 70, 7, 12 ]}\nobreak
\hyperdef{L}{X78525D04849A48EA}{}
{\noindent\textcolor{FuncColor}{$\triangleright$\enspace\texttt{WeightsTom({\mdseries\slshape tom})\index{WeightsTom@\texttt{WeightsTom}}
\label{WeightsTom}
}\hfill{\scriptsize (attribute)}}\\


 \texttt{WeightsTom} extracts the \emph{weights} from the table of marks \mbox{\texttt{\mdseries\slshape tom}}, i.e., the diagonal entries of the matrix of marks (see{\nobreakspace}\texttt{MarksTom} (\ref{MarksTom})), indicating the index of a subgroup in its normalizer. 

 
\begin{Verbatim}[commandchars=!@|,fontsize=\small,frame=single,label=Example]
  !gapprompt@gap>| !gapinput@wt:= WeightsTom( a5 );|
  [ 60, 2, 2, 3, 2, 1, 1, 1, 1 ]
\end{Verbatim}
 

 This information may be used to obtain the numbers of conjugate supergroups
from the marks. 
\begin{Verbatim}[commandchars=!@|,fontsize=\small,frame=single,label=Example]
  !gapprompt@gap>| !gapinput@marks:= MarksTom( a5 );;|
  !gapprompt@gap>| !gapinput@List( [ 1 .. 9 ], x -> marks[x] / wt[x] );|
  [ [ 1 ], [ 15, 1 ], [ 10, 1 ], [ 5, 1, 1 ], [ 6, 1 ], [ 10, 2, 1, 1 ],
    [ 6, 2, 1, 1 ], [ 5, 1, 2, 1, 1 ], [ 1, 1, 1, 1, 1, 1, 1, 1, 1 ] ]
\end{Verbatim}
 }

 }

  
\section{\textcolor{Chapter }{Properties of Tables of Marks}}\label{Properties of Tables of Marks}
\logpage{[ 70, 8, 0 ]}
\hyperdef{L}{X78A1B2E4826A9518}{}
{
  For a table of marks \mbox{\texttt{\mdseries\slshape tom}} of a group $G$, the following properties have the same meaning as the corresponding
properties for $G$. Additionally, if a positive integer \mbox{\texttt{\mdseries\slshape sub}} is given as the second argument then the value of the corresponding property
for the \mbox{\texttt{\mdseries\slshape sub}}\texttt{\symbol{45}}th class of subgroups of \mbox{\texttt{\mdseries\slshape tom}} is returned. 

 

\subsection{\textcolor{Chapter }{IsAbelianTom}}
\logpage{[ 70, 8, 1 ]}\nobreak
\hyperdef{L}{X7C93BAEC78B7C2B4}{}
{\noindent\textcolor{FuncColor}{$\triangleright$\enspace\texttt{IsAbelianTom({\mdseries\slshape tom[, sub]})\index{IsAbelianTom@\texttt{IsAbelianTom}}
\label{IsAbelianTom}
}\hfill{\scriptsize (property)}}\\
\noindent\textcolor{FuncColor}{$\triangleright$\enspace\texttt{IsCyclicTom({\mdseries\slshape tom[, sub]})\index{IsCyclicTom@\texttt{IsCyclicTom}}
\label{IsCyclicTom}
}\hfill{\scriptsize (property)}}\\
\noindent\textcolor{FuncColor}{$\triangleright$\enspace\texttt{IsNilpotentTom({\mdseries\slshape tom[, sub]})\index{IsNilpotentTom@\texttt{IsNilpotentTom}}
\label{IsNilpotentTom}
}\hfill{\scriptsize (property)}}\\
\noindent\textcolor{FuncColor}{$\triangleright$\enspace\texttt{IsPerfectTom({\mdseries\slshape tom[, sub]})\index{IsPerfectTom@\texttt{IsPerfectTom}}
\label{IsPerfectTom}
}\hfill{\scriptsize (property)}}\\
\noindent\textcolor{FuncColor}{$\triangleright$\enspace\texttt{IsSolvableTom({\mdseries\slshape tom[, sub]})\index{IsSolvableTom@\texttt{IsSolvableTom}}
\label{IsSolvableTom}
}\hfill{\scriptsize (property)}}\\


 
\begin{Verbatim}[commandchars=!@|,fontsize=\small,frame=single,label=Example]
  !gapprompt@gap>| !gapinput@tom:= TableOfMarks( "A5" );;|
  !gapprompt@gap>| !gapinput@IsAbelianTom( tom );  IsPerfectTom( tom );|
  false
  true
  !gapprompt@gap>| !gapinput@IsAbelianTom( tom, 3 );  IsNilpotentTom( tom, 7 );|
  true
  false
  !gapprompt@gap>| !gapinput@IsPerfectTom( tom, 7 );  IsSolvableTom( tom, 7 );|
  false
  true
  !gapprompt@gap>| !gapinput@for i in [ 1 .. 6 ] do|
  !gapprompt@>| !gapinput@Print( i, ": ", IsCyclicTom(a5, i), "  " );|
  !gapprompt@>| !gapinput@od;  Print( "\n" );|
  1: true  2: true  3: true  4: false  5: true  6: false
\end{Verbatim}
 }

 }

  
\section{\textcolor{Chapter }{Other Operations for Tables of Marks}}\label{Other Operations for Tables of Marks}
\logpage{[ 70, 9, 0 ]}
\hyperdef{L}{X7A40D99D7816F126}{}
{
  

\subsection{\textcolor{Chapter }{IsInternallyConsistent (for tables of marks)}}
\logpage{[ 70, 9, 1 ]}\nobreak
\hyperdef{L}{X7D8B4BE08094B137}{}
{\noindent\textcolor{FuncColor}{$\triangleright$\enspace\texttt{IsInternallyConsistent({\mdseries\slshape tom})\index{IsInternallyConsistent@\texttt{IsInternallyConsistent}!for tables of marks}
\label{IsInternallyConsistent:for tables of marks}
}\hfill{\scriptsize (method)}}\\


 For a table of marks \mbox{\texttt{\mdseries\slshape tom}}, \texttt{IsInternallyConsistent} decomposes all tensor products of rows of \mbox{\texttt{\mdseries\slshape tom}}. It returns \texttt{true} if all decomposition numbers are nonnegative integers, and \texttt{false} otherwise. This provides a strong consistency check for a table of marks. }

 

\subsection{\textcolor{Chapter }{DerivedSubgroupTom}}
\logpage{[ 70, 9, 2 ]}\nobreak
\hyperdef{L}{X8528D9397FFAF477}{}
{\noindent\textcolor{FuncColor}{$\triangleright$\enspace\texttt{DerivedSubgroupTom({\mdseries\slshape tom, sub})\index{DerivedSubgroupTom@\texttt{DerivedSubgroupTom}}
\label{DerivedSubgroupTom}
}\hfill{\scriptsize (operation)}}\\
\noindent\textcolor{FuncColor}{$\triangleright$\enspace\texttt{DerivedSubgroupsTom({\mdseries\slshape tom})\index{DerivedSubgroupsTom@\texttt{DerivedSubgroupsTom}}
\label{DerivedSubgroupsTom}
}\hfill{\scriptsize (function)}}\\


 For a table of marks \mbox{\texttt{\mdseries\slshape tom}} and a positive integer \mbox{\texttt{\mdseries\slshape sub}}, \texttt{DerivedSubgroupTom} returns either a positive integer $i$ or a list $l$ of positive integers. In the former case, the result means that the derived
subgroups of the subgroups in the \mbox{\texttt{\mdseries\slshape sub}}\texttt{\symbol{45}}th class of \mbox{\texttt{\mdseries\slshape tom}} lie in the $i$\texttt{\symbol{45}}th class. In the latter case, the class of the derived
subgroups could not be uniquely determined, and the position of the class of
derived subgroups is an entry of $l$. 

 Values computed with \texttt{DerivedSubgroupTom} are stored using the attribute \texttt{DerivedSubgroupsTomPossible} (\ref{DerivedSubgroupsTomPossible}). 

 \texttt{DerivedSubgroupsTom} is just the list of \texttt{DerivedSubgroupTom} values for all values of \mbox{\texttt{\mdseries\slshape sub}}. }

 

\subsection{\textcolor{Chapter }{DerivedSubgroupsTomPossible}}
\logpage{[ 70, 9, 3 ]}\nobreak
\hyperdef{L}{X7C29BD438127DFBE}{}
{\noindent\textcolor{FuncColor}{$\triangleright$\enspace\texttt{DerivedSubgroupsTomPossible({\mdseries\slshape tom})\index{DerivedSubgroupsTomPossible@\texttt{DerivedSubgroupsTomPossible}}
\label{DerivedSubgroupsTomPossible}
}\hfill{\scriptsize (attribute)}}\\
\noindent\textcolor{FuncColor}{$\triangleright$\enspace\texttt{DerivedSubgroupsTomUnique({\mdseries\slshape tom})\index{DerivedSubgroupsTomUnique@\texttt{DerivedSubgroupsTomUnique}}
\label{DerivedSubgroupsTomUnique}
}\hfill{\scriptsize (attribute)}}\\


 Let \mbox{\texttt{\mdseries\slshape tom}} be a table of marks. The value of the attribute \texttt{DerivedSubgroupsTomPossible} is a list in which the value at position $i$ {\textendash}if bound{\textendash} is a positive integer or a list; the
meaning of the entry is the same as in \texttt{DerivedSubgroupTom} (\ref{DerivedSubgroupTom}). 

 If the value of the attribute \texttt{DerivedSubgroupsTomUnique} is known for \mbox{\texttt{\mdseries\slshape tom}} then it is a list of positive integers, the value at position $i$ being the position of the class of derived subgroups of the $i$\texttt{\symbol{45}}th class of subgroups in \mbox{\texttt{\mdseries\slshape tom}}. 

 The derived subgroups are in general not uniquely determined by the table of
marks if no \texttt{UnderlyingGroup} (\ref{UnderlyingGroup:for tables of marks}) value is stored, so there is no default method for \texttt{DerivedSubgroupsTomUnique}. But in some cases the derived subgroups are explicitly set when the table of
marks is constructed. In this case, \texttt{DerivedSubgroupTom} (\ref{DerivedSubgroupTom}) does not set values in the \texttt{DerivedSubgroupsTomPossible} list. 

 The \texttt{DerivedSubgroupsTomUnique} value is automatically set when the last missing unique value is entered in
the \texttt{DerivedSubgroupsTomPossible} list by \texttt{DerivedSubgroupTom} (\ref{DerivedSubgroupTom}). }

 

\subsection{\textcolor{Chapter }{NormalizerTom}}
\logpage{[ 70, 9, 4 ]}\nobreak
\hyperdef{L}{X7CE6C45881F7F7D4}{}
{\noindent\textcolor{FuncColor}{$\triangleright$\enspace\texttt{NormalizerTom({\mdseries\slshape tom, sub})\index{NormalizerTom@\texttt{NormalizerTom}}
\label{NormalizerTom}
}\hfill{\scriptsize (operation)}}\\
\noindent\textcolor{FuncColor}{$\triangleright$\enspace\texttt{NormalizersTom({\mdseries\slshape tom})\index{NormalizersTom@\texttt{NormalizersTom}}
\label{NormalizersTom}
}\hfill{\scriptsize (attribute)}}\\


 Let \mbox{\texttt{\mdseries\slshape tom}} be the table of marks of a group $G$. \texttt{NormalizerTom} tries to find the conjugacy class of the normalizer $N$ in $G$ of a subgroup $U$ in the \mbox{\texttt{\mdseries\slshape sub}}\texttt{\symbol{45}}th class of \mbox{\texttt{\mdseries\slshape tom}}. The return value is either the list of class numbers of those subgroups that
have the right size and contain the subgroup and all subgroups that clearly
contain it as a normal subgroup, or the class number of the normalizer if it
is uniquely determined by these conditions. If \mbox{\texttt{\mdseries\slshape tom}} knows the subgroup lattice of $G$ (see{\nobreakspace}\texttt{IsTableOfMarksWithGens} (\ref{IsTableOfMarksWithGens})) then all normalizers are uniquely determined. \texttt{NormalizerTom} should never return an empty list. 

 \texttt{NormalizersTom} returns the list of positions of the classes of normalizers of subgroups in \mbox{\texttt{\mdseries\slshape tom}}. In addition to the criteria for a single class of subgroup used by \texttt{NormalizerTom}, the approximations of normalizers for several classes are used and thus \texttt{NormalizersTom} may return better approximations than \texttt{NormalizerTom}. 

 
\begin{Verbatim}[commandchars=!@|,fontsize=\small,frame=single,label=Example]
  !gapprompt@gap>| !gapinput@NormalizerTom( a5, 4 );|
  8
  !gapprompt@gap>| !gapinput@NormalizersTom( a5 );|
  [ 9, 4, 6, 8, 7, 6, 7, 8, 9 ]
\end{Verbatim}
 

 The example shows that a subgroup with class number 4 in $A_5$ (which is a Kleinian four group) is normalized by a subgroup in class 8. This
class contains the subgroups of $A_5$ which are isomorphic to $A_4$. }

 

\subsection{\textcolor{Chapter }{ContainedTom}}
\logpage{[ 70, 9, 5 ]}\nobreak
\hyperdef{L}{X7F87B2797827E5DE}{}
{\noindent\textcolor{FuncColor}{$\triangleright$\enspace\texttt{ContainedTom({\mdseries\slshape tom, sub1, sub2})\index{ContainedTom@\texttt{ContainedTom}}
\label{ContainedTom}
}\hfill{\scriptsize (operation)}}\\


 \texttt{ContainedTom} returns the number of subgroups in class \mbox{\texttt{\mdseries\slshape sub1}} of the table of marks \mbox{\texttt{\mdseries\slshape tom}} that are contained in one fixed member of the class \mbox{\texttt{\mdseries\slshape sub2}}. 

 
\begin{Verbatim}[commandchars=!@|,fontsize=\small,frame=single,label=Example]
  !gapprompt@gap>| !gapinput@ContainedTom( a5, 3, 5 );  ContainedTom( a5, 3, 8 );|
  0
  4
\end{Verbatim}
 }

 

\subsection{\textcolor{Chapter }{ContainingTom}}
\logpage{[ 70, 9, 6 ]}\nobreak
\hyperdef{L}{X7EE050FB87D6F0E7}{}
{\noindent\textcolor{FuncColor}{$\triangleright$\enspace\texttt{ContainingTom({\mdseries\slshape tom, sub1, sub2})\index{ContainingTom@\texttt{ContainingTom}}
\label{ContainingTom}
}\hfill{\scriptsize (operation)}}\\


 \texttt{ContainingTom} returns the number of subgroups in class \mbox{\texttt{\mdseries\slshape sub2}} of the table of marks \mbox{\texttt{\mdseries\slshape tom}} that contain one fixed member of the class \mbox{\texttt{\mdseries\slshape sub1}}. 

 
\begin{Verbatim}[commandchars=!@|,fontsize=\small,frame=single,label=Example]
  !gapprompt@gap>| !gapinput@ContainingTom( a5, 3, 5 );  ContainingTom( a5, 3, 8 );|
  0
  2
\end{Verbatim}
 }

 

\subsection{\textcolor{Chapter }{CyclicExtensionsTom (for a prime)}}
\logpage{[ 70, 9, 7 ]}\nobreak
\hyperdef{L}{X838DE06B823C19CA}{}
{\noindent\textcolor{FuncColor}{$\triangleright$\enspace\texttt{CyclicExtensionsTom({\mdseries\slshape tom, p})\index{CyclicExtensionsTom@\texttt{CyclicExtensionsTom}!for a prime}
\label{CyclicExtensionsTom:for a prime}
}\hfill{\scriptsize (operation)}}\\
\noindent\textcolor{FuncColor}{$\triangleright$\enspace\texttt{CyclicExtensionsTom({\mdseries\slshape tom[, list]})\index{CyclicExtensionsTom@\texttt{CyclicExtensionsTom}!for a list of primes}
\label{CyclicExtensionsTom:for a list of primes}
}\hfill{\scriptsize (attribute)}}\\


 According to A.{\nobreakspace}Dress{\nobreakspace}\cite{Dre69}, two columns of the table of marks \mbox{\texttt{\mdseries\slshape tom}} are equal modulo the prime \mbox{\texttt{\mdseries\slshape p}} if and only if the corresponding subgroups are connected by a chain of normal
extensions of order \mbox{\texttt{\mdseries\slshape p}}. 

 Called with \mbox{\texttt{\mdseries\slshape tom}} and \mbox{\texttt{\mdseries\slshape p}}, \texttt{CyclicExtensionsTom} returns the classes of this equivalence relation. 

 In the second form, \mbox{\texttt{\mdseries\slshape list}} must be a list of primes, and the return value is the list of classes of the
relation obtained by considering chains of normal extensions of prime order
where all primes are in \mbox{\texttt{\mdseries\slshape list}}. The default value for \mbox{\texttt{\mdseries\slshape list}} is the set of prime divisors of the order of the group of \mbox{\texttt{\mdseries\slshape tom}}. 

 (This information is \emph{not} used by \texttt{NormalizerTom} (\ref{NormalizerTom}) although it might give additional restrictions in the search of normalizers.) 

 
\begin{Verbatim}[commandchars=!@|,fontsize=\small,frame=single,label=Example]
  !gapprompt@gap>| !gapinput@CyclicExtensionsTom( a5, 2 );|
  [ [ 1, 2, 4 ], [ 3, 6 ], [ 5, 7 ], [ 8 ], [ 9 ] ]
\end{Verbatim}
 }

 

\subsection{\textcolor{Chapter }{DecomposedFixedPointVector}}
\logpage{[ 70, 9, 8 ]}\nobreak
\hyperdef{L}{X80890C247EB1E35C}{}
{\noindent\textcolor{FuncColor}{$\triangleright$\enspace\texttt{DecomposedFixedPointVector({\mdseries\slshape tom, fix})\index{DecomposedFixedPointVector@\texttt{DecomposedFixedPointVector}}
\label{DecomposedFixedPointVector}
}\hfill{\scriptsize (operation)}}\\


 Let \mbox{\texttt{\mdseries\slshape tom}} be the table of marks of a group $G$ and let \mbox{\texttt{\mdseries\slshape fix}} be a vector of fixed point numbers w.r.t.{\nobreakspace}an action of $G$, i.e., a vector which contains for each class of subgroups the number of
fixed points under the given action. \texttt{DecomposedFixedPointVector} returns the decomposition of \mbox{\texttt{\mdseries\slshape fix}} into rows of the table of marks. This decomposition corresponds to a
decomposition of the action into transitive constituents. Trailing zeros in \mbox{\texttt{\mdseries\slshape fix}} may be omitted. 

 
\begin{Verbatim}[commandchars=!@|,fontsize=\small,frame=single,label=Example]
  !gapprompt@gap>| !gapinput@DecomposedFixedPointVector( a5, [ 16, 4, 1, 0, 1, 1, 1 ] );|
  [ 0, 0, 0, 0, 0, 1, 1 ]
\end{Verbatim}
 

 The vector \mbox{\texttt{\mdseries\slshape fix}} may be any vector of integers. The resulting decomposition, however, will not
be integral, in general. 
\begin{Verbatim}[commandchars=!@|,fontsize=\small,frame=single,label=Example]
  !gapprompt@gap>| !gapinput@DecomposedFixedPointVector( a5, [ 0, 0, 0, 0, 1, 1 ] );|
  [ 2/5, -1, -1/2, 0, 1/2, 1 ]
\end{Verbatim}
 }

 

\subsection{\textcolor{Chapter }{EulerianFunctionByTom}}
\logpage{[ 70, 9, 9 ]}\nobreak
\hyperdef{L}{X7B1C1A7C867A4082}{}
{\noindent\textcolor{FuncColor}{$\triangleright$\enspace\texttt{EulerianFunctionByTom({\mdseries\slshape tom, n[, sub]})\index{EulerianFunctionByTom@\texttt{EulerianFunctionByTom}}
\label{EulerianFunctionByTom}
}\hfill{\scriptsize (operation)}}\\


 Called with two arguments, \texttt{EulerianFunctionByTom} computes the Eulerian function (see{\nobreakspace}\texttt{EulerianFunction} (\ref{EulerianFunction})) of the underlying group $G$ of the table of marks \mbox{\texttt{\mdseries\slshape tom}}, that is, the number of \mbox{\texttt{\mdseries\slshape n}}\texttt{\symbol{45}}tuples of elements in $G$ that generate $G$. If the optional argument \mbox{\texttt{\mdseries\slshape sub}} is given then \texttt{EulerianFunctionByTom} computes the Eulerian function of each subgroup in the \mbox{\texttt{\mdseries\slshape sub}}\texttt{\symbol{45}}th class of subgroups of \mbox{\texttt{\mdseries\slshape tom}}. 

 For a group $G$ whose table of marks is known, \texttt{EulerianFunctionByTom} is installed as a method for \texttt{EulerianFunction} (\ref{EulerianFunction}). 

 
\begin{Verbatim}[commandchars=!@|,fontsize=\small,frame=single,label=Example]
  !gapprompt@gap>| !gapinput@EulerianFunctionByTom( a5, 2 );|
  2280
  !gapprompt@gap>| !gapinput@EulerianFunctionByTom( a5, 3 );|
  200160
  !gapprompt@gap>| !gapinput@EulerianFunctionByTom( a5, 2, 3 );|
  8
\end{Verbatim}
 }

 

\subsection{\textcolor{Chapter }{IntersectionsTom}}
\logpage{[ 70, 9, 10 ]}\nobreak
\hyperdef{L}{X8224E51382FDB912}{}
{\noindent\textcolor{FuncColor}{$\triangleright$\enspace\texttt{IntersectionsTom({\mdseries\slshape tom, sub1, sub2})\index{IntersectionsTom@\texttt{IntersectionsTom}}
\label{IntersectionsTom}
}\hfill{\scriptsize (operation)}}\\


 The intersections of the groups in the \mbox{\texttt{\mdseries\slshape sub1}}\texttt{\symbol{45}}th conjugacy class of subgroups of the table of marks \mbox{\texttt{\mdseries\slshape tom}} with the groups in the \mbox{\texttt{\mdseries\slshape sub2}}\texttt{\symbol{45}}th conjugacy classes of subgroups of \mbox{\texttt{\mdseries\slshape tom}} are determined up to conjugacy by the decomposition of the tensor product of
their rows of marks. \texttt{IntersectionsTom} returns a list $l$ that describes this decomposition. The $i$\texttt{\symbol{45}}th entry in $l$ is the multiplicity of groups in the $i$\texttt{\symbol{45}}th conjugacy class as an intersection. 

 
\begin{Verbatim}[commandchars=!@|,fontsize=\small,frame=single,label=Example]
  !gapprompt@gap>| !gapinput@IntersectionsTom( a5, 8, 8 );|
  [ 0, 0, 1, 0, 0, 0, 0, 1 ]
\end{Verbatim}
 Any two subgroups of class number 8 ($A_4$) of $A_5$ are either equal and their intersection has again class number 8, or their
intersection has class number $3$, and is a cyclic subgroup of order 3. }

 

\subsection{\textcolor{Chapter }{FactorGroupTom}}
\logpage{[ 70, 9, 11 ]}\nobreak
\hyperdef{L}{X859F069C8428B598}{}
{\noindent\textcolor{FuncColor}{$\triangleright$\enspace\texttt{FactorGroupTom({\mdseries\slshape tom, n})\index{FactorGroupTom@\texttt{FactorGroupTom}}
\label{FactorGroupTom}
}\hfill{\scriptsize (operation)}}\\


 For a table of marks \mbox{\texttt{\mdseries\slshape tom}} of a group $G$ and the normal subgroup $N$ of $G$ corresponding to the \mbox{\texttt{\mdseries\slshape n}}\texttt{\symbol{45}}th class of subgroups of \mbox{\texttt{\mdseries\slshape tom}}, \texttt{FactorGroupTom} returns the table of marks of the factor group $G / N$. 

 
\begin{Verbatim}[commandchars=!@|,fontsize=\small,frame=single,label=Example]
  !gapprompt@gap>| !gapinput@s4:= TableOfMarks( SymmetricGroup( 4 ) );|
  TableOfMarks( Sym( [ 1 .. 4 ] ) )
  !gapprompt@gap>| !gapinput@LengthsTom( s4 );|
  [ 1, 3, 6, 4, 1, 3, 3, 4, 3, 1, 1 ]
  !gapprompt@gap>| !gapinput@OrdersTom( s4 );|
  [ 1, 2, 2, 3, 4, 4, 4, 6, 8, 12, 24 ]
  !gapprompt@gap>| !gapinput@s3:= FactorGroupTom( s4, 5 );|
  TableOfMarks( Group([ f1, f2 ]) )
  !gapprompt@gap>| !gapinput@Display( s3 );|
  1:  6
  2:  3 1
  3:  2 . 2
  4:  1 1 1 1
  
\end{Verbatim}
 }

 

\subsection{\textcolor{Chapter }{MaximalSubgroupsTom}}
\logpage{[ 70, 9, 12 ]}\nobreak
\hyperdef{L}{X8325811586C00ECF}{}
{\noindent\textcolor{FuncColor}{$\triangleright$\enspace\texttt{MaximalSubgroupsTom({\mdseries\slshape tom[, sub]})\index{MaximalSubgroupsTom@\texttt{MaximalSubgroupsTom}}
\label{MaximalSubgroupsTom}
}\hfill{\scriptsize (attribute)}}\\


 Called with a table of marks \mbox{\texttt{\mdseries\slshape tom}}, \texttt{MaximalSubgroupsTom} returns a list of length two, the first entry being the list of positions of
the classes of maximal subgroups of the whole group of \mbox{\texttt{\mdseries\slshape tom}}, the second entry being the list of class lengths of these groups. 

 Called with a table of marks \mbox{\texttt{\mdseries\slshape tom}} and a position \mbox{\texttt{\mdseries\slshape sub}}, the same information for the \mbox{\texttt{\mdseries\slshape sub}}\texttt{\symbol{45}}th class of subgroups is returned. 

 
\begin{Verbatim}[commandchars=!@|,fontsize=\small,frame=single,label=Example]
  !gapprompt@gap>| !gapinput@MaximalSubgroupsTom( s4 );|
  [ [ 10, 9, 8 ], [ 1, 3, 4 ] ]
  !gapprompt@gap>| !gapinput@MaximalSubgroupsTom( s4, 10 );|
  [ [ 5, 4 ], [ 1, 4 ] ]
\end{Verbatim}
 }

 

\subsection{\textcolor{Chapter }{MinimalSupergroupsTom}}
\logpage{[ 70, 9, 13 ]}\nobreak
\hyperdef{L}{X7923B19D7C47BF63}{}
{\noindent\textcolor{FuncColor}{$\triangleright$\enspace\texttt{MinimalSupergroupsTom({\mdseries\slshape tom, sub})\index{MinimalSupergroupsTom@\texttt{MinimalSupergroupsTom}}
\label{MinimalSupergroupsTom}
}\hfill{\scriptsize (operation)}}\\


 For a table of marks \mbox{\texttt{\mdseries\slshape tom}}, \texttt{MinimalSupergroupsTom} returns a list of length two, the first entry being the list of positions of
the classes containing the minimal supergroups of the groups in the \mbox{\texttt{\mdseries\slshape sub}}\texttt{\symbol{45}}th class of subgroups of \mbox{\texttt{\mdseries\slshape tom}}, the second entry being the list of class lengths of these groups. 

 
\begin{Verbatim}[commandchars=!@|,fontsize=\small,frame=single,label=Example]
  !gapprompt@gap>| !gapinput@MinimalSupergroupsTom( s4, 5 );|
  [ [ 9, 10 ], [ 3, 1 ] ]
\end{Verbatim}
 }

 }

  
\section{\textcolor{Chapter }{Accessing Subgroups via Tables of Marks}}\label{Accessing Subgroups via Tables of Marks}
\logpage{[ 70, 10, 0 ]}
\hyperdef{L}{X7FE9BE477A90199F}{}
{
  Let \mbox{\texttt{\mdseries\slshape tom}} be the table of marks of the group $G$, and assume that \mbox{\texttt{\mdseries\slshape tom}} has access to $G$ via the \texttt{UnderlyingGroup} (\ref{UnderlyingGroup:for tables of marks}) value. Then it makes sense to use \mbox{\texttt{\mdseries\slshape tom}} and its ordering of conjugacy classes of subgroups of $G$ for storing information for constructing representatives of these classes. The
group $G$ is in general not sufficient for this, \mbox{\texttt{\mdseries\slshape tom}} needs more information; this is available if and only if the \texttt{IsTableOfMarksWithGens} (\ref{IsTableOfMarksWithGens}) value of \mbox{\texttt{\mdseries\slshape tom}} is \texttt{true}. In this case, \texttt{RepresentativeTom} (\ref{RepresentativeTom}) can be used to get a subgroup of the $i$\texttt{\symbol{45}}th class, for all $i$. 

 \textsf{GAP} provides two different possibilities to store generators of the
representatives of classes of subgroups. The first is implemented by the
attribute \texttt{GeneratorsSubgroupsTom} (\ref{GeneratorsSubgroupsTom}), which uses explicit generators of the subgroups. The second, more general,
possibility is implemented by the attribute \texttt{StraightLineProgramsTom} (\ref{StraightLineProgramsTom}), which encodes the generators as straight line programs (see{\nobreakspace}\ref{Straight Line Programs}) that evaluate to the generators in question when applied to \emph{standard generators} of $G$.  This means that on the one hand, standard generators of $G$ must be known in order to use \texttt{StraightLineProgramsTom} (\ref{StraightLineProgramsTom}). On the other hand, the straight line programs allow one to compute easily
generators not only of a subgroup $U$ of $G$ but also generators of the image of $U$ in any representation of $G$, provided that one knows standard generators of the image of $G$ under this representation. See the manual of the package \textsf{TomLib} for details and an example. 

\subsection{\textcolor{Chapter }{GeneratorsSubgroupsTom}}
\logpage{[ 70, 10, 1 ]}\nobreak
\hyperdef{L}{X7B0B6FDD806E9734}{}
{\noindent\textcolor{FuncColor}{$\triangleright$\enspace\texttt{GeneratorsSubgroupsTom({\mdseries\slshape tom})\index{GeneratorsSubgroupsTom@\texttt{GeneratorsSubgroupsTom}}
\label{GeneratorsSubgroupsTom}
}\hfill{\scriptsize (attribute)}}\\


 Let \mbox{\texttt{\mdseries\slshape tom}} be a table of marks with \texttt{IsTableOfMarksWithGens} (\ref{IsTableOfMarksWithGens}) value \texttt{true}. Then \texttt{GeneratorsSubgroupsTom} returns a list of length two, the first entry being a list $l$ of elements of the group stored as \texttt{UnderlyingGroup} (\ref{UnderlyingGroup:for tables of marks}) value of \mbox{\texttt{\mdseries\slshape tom}}, the second entry being a list that contains at position $i$ a list of positions in $l$ of generators of a representative of a subgroup in class $i$. 

 The \texttt{GeneratorsSubgroupsTom} value is known for all tables of marks that have been computed with \texttt{TableOfMarks} (\ref{TableOfMarks:for a group}) from a group, and there is a method to compute the value for a table of marks
that admits \texttt{RepresentativeTom} (\ref{RepresentativeTom}). }

 

\subsection{\textcolor{Chapter }{StraightLineProgramsTom}}
\logpage{[ 70, 10, 2 ]}\nobreak
\hyperdef{L}{X7898BE7284E47FF3}{}
{\noindent\textcolor{FuncColor}{$\triangleright$\enspace\texttt{StraightLineProgramsTom({\mdseries\slshape tom})\index{StraightLineProgramsTom@\texttt{StraightLineProgramsTom}}
\label{StraightLineProgramsTom}
}\hfill{\scriptsize (attribute)}}\\


 For a table of marks \mbox{\texttt{\mdseries\slshape tom}} with \texttt{IsTableOfMarksWithGens} (\ref{IsTableOfMarksWithGens}) value \texttt{true}, \texttt{StraightLineProgramsTom} returns a list that contains at position $i$ either a list of straight line programs or a straight line program
(see{\nobreakspace}\ref{Straight Line Programs}), encoding the generators of a representative of the $i$\texttt{\symbol{45}}th conjugacy class of subgroups of \texttt{UnderlyingGroup( \mbox{\texttt{\mdseries\slshape tom}} )}; in the former case, each straight line program returns a generator, in the
latter case, the program returns the list of generators. 

 There is no default method to compute the \texttt{StraightLineProgramsTom} value of a table of marks if they are not yet stored. The value is known for
all tables of marks that belong to the \textsf{GAP} library of tables of marks (see{\nobreakspace}\ref{The Library of Tables of Marks}). }

 

\subsection{\textcolor{Chapter }{IsTableOfMarksWithGens}}
\logpage{[ 70, 10, 3 ]}\nobreak
\hyperdef{L}{X7889DB6D790593B9}{}
{\noindent\textcolor{FuncColor}{$\triangleright$\enspace\texttt{IsTableOfMarksWithGens({\mdseries\slshape tom})\index{IsTableOfMarksWithGens@\texttt{IsTableOfMarksWithGens}}
\label{IsTableOfMarksWithGens}
}\hfill{\scriptsize (filter)}}\\


 This filter shall express the union of the filters \texttt{IsTableOfMarks and HasStraightLineProgramsTom} and \texttt{IsTableOfMarks and HasGeneratorsSubgroupsTom}. If a table of marks \mbox{\texttt{\mdseries\slshape tom}} has this filter set then \mbox{\texttt{\mdseries\slshape tom}} can be asked to compute information that is in general not uniquely determined
by a table of marks, for example the positions of derived subgroups or
normalizers of subgroups (see{\nobreakspace}\texttt{DerivedSubgroupTom} (\ref{DerivedSubgroupTom}), \texttt{NormalizerTom} (\ref{NormalizerTom})). 

 
\begin{Verbatim}[commandchars=!@|,fontsize=\small,frame=single,label=Example]
  !gapprompt@gap>| !gapinput@a5:= TableOfMarks( "A5" );;  IsTableOfMarksWithGens( a5 );|
  true
  !gapprompt@gap>| !gapinput@HasGeneratorsSubgroupsTom( a5 );  HasStraightLineProgramsTom( a5 );|
  false
  true
  !gapprompt@gap>| !gapinput@alt5:= TableOfMarks( AlternatingGroup( 5 ) );;|
  !gapprompt@gap>| !gapinput@IsTableOfMarksWithGens( alt5 );|
  true
  !gapprompt@gap>| !gapinput@HasGeneratorsSubgroupsTom(alt5); HasStraightLineProgramsTom(alt5);|
  true
  false
  !gapprompt@gap>| !gapinput@progs:= StraightLineProgramsTom( a5 );;|
  !gapprompt@gap>| !gapinput@OrdersTom( a5 );|
  [ 1, 2, 3, 4, 5, 6, 10, 12, 60 ]
  !gapprompt@gap>| !gapinput@IsCyclicTom( a5, 4 );|
  false
  !gapprompt@gap>| !gapinput@Length( progs[4] );|
  2
  !gapprompt@gap>| !gapinput@progs[4][1];|
  <straight line program>
  !gapprompt@gap>| !gapinput@# first generator of an el. ab group of order 4:|
  !gapprompt@gap>| !gapinput@Display( progs[4][1] );|
  # input:
  r:= [ g1, g2 ];
  # program:
  r[3]:= r[2]*r[1];
  r[4]:= r[3]*r[2]^-1*r[1]*r[3]*r[2]^-1*r[1]*r[2];
  # return value:
  r[4]
  !gapprompt@gap>| !gapinput@x:= [ [ Z(2)^0, 0*Z(2) ], [ Z(2^2), Z(2)^0 ] ];;|
  !gapprompt@gap>| !gapinput@y:= [ [ Z(2^2), Z(2)^0 ], [ 0*Z(2), Z(2^2)^2 ] ];;|
  !gapprompt@gap>| !gapinput@res1:= ResultOfStraightLineProgram( progs[4][1], [ x, y ] );|
  [ [ Z(2)^0, 0*Z(2) ], [ Z(2^2)^2, Z(2)^0 ] ]
  !gapprompt@gap>| !gapinput@res2:= ResultOfStraightLineProgram( progs[4][2], [ x, y ] );|
  [ [ Z(2)^0, 0*Z(2) ], [ Z(2^2), Z(2)^0 ] ]
  !gapprompt@gap>| !gapinput@w:= y*x;;|
  !gapprompt@gap>| !gapinput@res1 = w*y^-1*x*w*y^-1*x*y;|
  true
  !gapprompt@gap>| !gapinput@subgrp:= Group( res1, res2 );;  Size( subgrp );  IsCyclic( subgrp );|
  4
  false
\end{Verbatim}
 }

 

\subsection{\textcolor{Chapter }{RepresentativeTom}}
\logpage{[ 70, 10, 4 ]}\nobreak
\hyperdef{L}{X7F625AB880B73AC3}{}
{\noindent\textcolor{FuncColor}{$\triangleright$\enspace\texttt{RepresentativeTom({\mdseries\slshape tom, sub})\index{RepresentativeTom@\texttt{RepresentativeTom}}
\label{RepresentativeTom}
}\hfill{\scriptsize (operation)}}\\
\noindent\textcolor{FuncColor}{$\triangleright$\enspace\texttt{RepresentativeTomByGenerators({\mdseries\slshape tom, sub, gens})\index{RepresentativeTomByGenerators@\texttt{RepresentativeTomByGenerators}}
\label{RepresentativeTomByGenerators}
}\hfill{\scriptsize (operation)}}\\
\noindent\textcolor{FuncColor}{$\triangleright$\enspace\texttt{RepresentativeTomByGeneratorsNC({\mdseries\slshape tom, sub, gens})\index{RepresentativeTomByGeneratorsNC@\texttt{RepresentativeTomByGeneratorsNC}}
\label{RepresentativeTomByGeneratorsNC}
}\hfill{\scriptsize (operation)}}\\


 Let \mbox{\texttt{\mdseries\slshape tom}} be a table of marks with \texttt{IsTableOfMarksWithGens} (\ref{IsTableOfMarksWithGens}) value \texttt{true}, and \mbox{\texttt{\mdseries\slshape sub}} a positive integer. \texttt{RepresentativeTom} returns a representative of the \mbox{\texttt{\mdseries\slshape sub}}\texttt{\symbol{45}}th conjugacy class of subgroups of \mbox{\texttt{\mdseries\slshape tom}}. 

 If the attribute \texttt{StraightLineProgramsTom} (\ref{StraightLineProgramsTom}) is set in \mbox{\texttt{\mdseries\slshape tom}} then methods for the operations \texttt{RepresentativeTomByGenerators} and \texttt{RepresentativeTomByGeneratorsNC} are available, which return a representative of the \mbox{\texttt{\mdseries\slshape sub}}\texttt{\symbol{45}}th conjugacy class of subgroups of \mbox{\texttt{\mdseries\slshape tom}}, as a subgroup of the group generated by \mbox{\texttt{\mdseries\slshape gens}}. This means that the standard generators of \mbox{\texttt{\mdseries\slshape tom}} are replaced by \mbox{\texttt{\mdseries\slshape gens}}. 

 \texttt{RepresentativeTomByGenerators} checks whether mapping the standard generators of \mbox{\texttt{\mdseries\slshape tom}} to \mbox{\texttt{\mdseries\slshape gens}} extends to a group isomorphism, and returns \texttt{fail} if not. \texttt{RepresentativeTomByGeneratorsNC} omits all checks. So \texttt{RepresentativeTomByGenerators} is thought mainly for debugging purposes; note that when several
representatives are constructed, it is cheaper to construct (and check) the
isomorphism once, and to map the groups returned by \texttt{RepresentativeTom} under this isomorphism. The idea behind \texttt{RepresentativeTomByGeneratorsNC}, however, is to avoid the overhead of using isomorphisms when \mbox{\texttt{\mdseries\slshape gens}} are known to be standard generators. In order to proceed like this, the
attribute \texttt{StraightLineProgramsTom} (\ref{StraightLineProgramsTom}) is needed. 

 
\begin{Verbatim}[commandchars=!@|,fontsize=\small,frame=single,label=Example]
  !gapprompt@gap>| !gapinput@RepresentativeTom( a5, 4 );|
  Group([ (2,3)(4,5), (2,4)(3,5) ])
\end{Verbatim}
 }

 }

  
\section{\textcolor{Chapter }{The Interface between Tables of Marks and Character Tables}}\label{The Interface between Tables of Marks and Character Tables}
\logpage{[ 70, 11, 0 ]}
\hyperdef{L}{X79ADA60880BE9C49}{}
{
  The following examples require the \textsf{GAP} Character Table Library to be available. If it is not yet loaded then we load
it now. 

 
\begin{Verbatim}[commandchars=!@|,fontsize=\small,frame=single,label=Example]
  !gapprompt@gap>| !gapinput@LoadPackage( "ctbllib" );|
  true
\end{Verbatim}
 

 

\subsection{\textcolor{Chapter }{FusionCharTableTom}}
\logpage{[ 70, 11, 1 ]}\nobreak
\hyperdef{L}{X7A82CB487DBDDC53}{}
{\noindent\textcolor{FuncColor}{$\triangleright$\enspace\texttt{FusionCharTableTom({\mdseries\slshape tbl, tom})\index{FusionCharTableTom@\texttt{FusionCharTableTom}}
\label{FusionCharTableTom}
}\hfill{\scriptsize (operation)}}\\
\noindent\textcolor{FuncColor}{$\triangleright$\enspace\texttt{PossibleFusionsCharTableTom({\mdseries\slshape tbl, tom[, options]})\index{PossibleFusionsCharTableTom@\texttt{PossibleFusionsCharTableTom}}
\label{PossibleFusionsCharTableTom}
}\hfill{\scriptsize (operation)}}\\


 Let \mbox{\texttt{\mdseries\slshape tbl}} be the ordinary character table of the group $G$, say, and \mbox{\texttt{\mdseries\slshape tom}} the table of marks of $G$. \texttt{FusionCharTableTom} determines the fusion of the classes of elements from \mbox{\texttt{\mdseries\slshape tbl}} to the classes of cyclic subgroups on \mbox{\texttt{\mdseries\slshape tom}}, that is, a list that contains at position $i$ the position of the class of cyclic subgroups in \mbox{\texttt{\mdseries\slshape tom}} that are generated by elements in the $i$\texttt{\symbol{45}}th conjugacy class of elements in \mbox{\texttt{\mdseries\slshape tbl}}. 

 Three cases are handled differently. 
\begin{enumerate}
\item  The fusion is explicitly stored on \mbox{\texttt{\mdseries\slshape tbl}}. Then nothing has to be done. This happens only if both \mbox{\texttt{\mdseries\slshape tbl}} and \mbox{\texttt{\mdseries\slshape tom}} are tables from the \textsf{GAP} library (see{\nobreakspace}\ref{The Library of Tables of Marks} and the manual of the \textsf{GAP} Character Table Library). 
\item  The \texttt{UnderlyingGroup} (\ref{UnderlyingGroup:for tables of marks}) values of \mbox{\texttt{\mdseries\slshape tbl}} and \mbox{\texttt{\mdseries\slshape tom}} are known and equal. Then the group is used to compute the fusion. 
\item  There is neither fusion nor group information available. In this case only
necessary conditions can be checked, and if they are not sufficient to
determine the fusion uniquely then \texttt{fail} is returned by \texttt{FusionCharTableTom}. 
\end{enumerate}
 

 \texttt{PossibleFusionsCharTableTom} computes the list of possible fusions from \mbox{\texttt{\mdseries\slshape tbl}} to \mbox{\texttt{\mdseries\slshape tom}}, according to the criteria that have been checked. So if \texttt{FusionCharTableTom} returns a unique fusion then the list returned by \texttt{PossibleFusionsCharTableTom} for the same arguments contains exactly this fusion, and if \texttt{FusionCharTableTom} returns \texttt{fail} then the length of this list is different from $1$.  

 The optional argument \mbox{\texttt{\mdseries\slshape options}} must be a record that may have the following components. 
\begin{description}
\item[{\texttt{fusionmap}}]  a parametrized map which is an approximation of the desired map, 
\item[{\texttt{quick}}]  a Boolean; if \texttt{true} then as soon as only one possibility remains this possibility is returned
immediately; the default value is \texttt{false}. 
\end{description}
 

 
\begin{Verbatim}[commandchars=!@|,fontsize=\small,frame=single,label=Example]
  !gapprompt@gap>| !gapinput@a5c:= CharacterTable( "A5" );;|
  !gapprompt@gap>| !gapinput@fus:= FusionCharTableTom( a5c, a5 );|
  [ 1, 2, 3, 5, 5 ]
\end{Verbatim}
 }

 

\subsection{\textcolor{Chapter }{PermCharsTom (via fusion map)}}
\logpage{[ 70, 11, 2 ]}\nobreak
\hyperdef{L}{X8016499282F0BA37}{}
{\noindent\textcolor{FuncColor}{$\triangleright$\enspace\texttt{PermCharsTom({\mdseries\slshape fus, tom})\index{PermCharsTom@\texttt{PermCharsTom}!via fusion map}
\label{PermCharsTom:via fusion map}
}\hfill{\scriptsize (operation)}}\\
\noindent\textcolor{FuncColor}{$\triangleright$\enspace\texttt{PermCharsTom({\mdseries\slshape tbl, tom})\index{PermCharsTom@\texttt{PermCharsTom}!from a character table}
\label{PermCharsTom:from a character table}
}\hfill{\scriptsize (operation)}}\\


 \texttt{PermCharsTom} returns the list of transitive permutation characters from the table of marks \mbox{\texttt{\mdseries\slshape tom}}. In the first form, \mbox{\texttt{\mdseries\slshape fus}} must be the fusion map from the ordinary character table of the group of \mbox{\texttt{\mdseries\slshape tom}} to \mbox{\texttt{\mdseries\slshape tom}} (see{\nobreakspace}\texttt{FusionCharTableTom} (\ref{FusionCharTableTom})). In the second form, \mbox{\texttt{\mdseries\slshape tbl}} must be the character table of the group of which \mbox{\texttt{\mdseries\slshape tom}} is the table of marks. If the fusion map is not uniquely determined
(see{\nobreakspace}\texttt{FusionCharTableTom} (\ref{FusionCharTableTom})) then \texttt{fail} is returned. 

 If the fusion map \mbox{\texttt{\mdseries\slshape fus}} is given as first argument then each transitive permutation character is
represented by its values list. If the character table \mbox{\texttt{\mdseries\slshape tbl}} is given then the permutation characters are class function objects (see
Chapter{\nobreakspace}\ref{Class Functions}). 

 
\begin{Verbatim}[commandchars=!@|,fontsize=\small,frame=single,label=Example]
  !gapprompt@gap>| !gapinput@PermCharsTom( a5c, a5 );|
  [ Character( CharacterTable( "A5" ), [ 60, 0, 0, 0, 0 ] ),
    Character( CharacterTable( "A5" ), [ 30, 2, 0, 0, 0 ] ),
    Character( CharacterTable( "A5" ), [ 20, 0, 2, 0, 0 ] ),
    Character( CharacterTable( "A5" ), [ 15, 3, 0, 0, 0 ] ),
    Character( CharacterTable( "A5" ), [ 12, 0, 0, 2, 2 ] ),
    Character( CharacterTable( "A5" ), [ 10, 2, 1, 0, 0 ] ),
    Character( CharacterTable( "A5" ), [ 6, 2, 0, 1, 1 ] ),
    Character( CharacterTable( "A5" ), [ 5, 1, 2, 0, 0 ] ),
    Character( CharacterTable( "A5" ), [ 1, 1, 1, 1, 1 ] ) ]
  !gapprompt@gap>| !gapinput@PermCharsTom( fus, a5 )[1];|
  [ 60, 0, 0, 0, 0 ]
\end{Verbatim}
 }

 }

  
\section{\textcolor{Chapter }{Generic Construction of Tables of Marks}}\label{Generic Construction of Tables of Marks}
\logpage{[ 70, 12, 0 ]}
\hyperdef{L}{X7CF66FAE7A8858E4}{}
{
  The following three operations construct a table of marks only from the data
given, i.e., without underlying group. 

\subsection{\textcolor{Chapter }{TableOfMarksCyclic}}
\logpage{[ 70, 12, 1 ]}\nobreak
\hyperdef{L}{X7CAA5B6C85CB9A8D}{}
{\noindent\textcolor{FuncColor}{$\triangleright$\enspace\texttt{TableOfMarksCyclic({\mdseries\slshape n})\index{TableOfMarksCyclic@\texttt{TableOfMarksCyclic}}
\label{TableOfMarksCyclic}
}\hfill{\scriptsize (operation)}}\\


 \texttt{TableOfMarksCyclic} returns the table of marks of the cyclic group of order \mbox{\texttt{\mdseries\slshape n}}. 

 A cyclic group of order \mbox{\texttt{\mdseries\slshape n}} has as its subgroups for each divisor $d$ of \mbox{\texttt{\mdseries\slshape n}} a cyclic subgroup of order $d$. 

 
\begin{Verbatim}[commandchars=!@|,fontsize=\small,frame=single,label=Example]
  !gapprompt@gap>| !gapinput@Display( TableOfMarksCyclic( 6 ) );|
  1:  6
  2:  3 3
  3:  2 . 2
  4:  1 1 1 1
  
\end{Verbatim}
 }

 

\subsection{\textcolor{Chapter }{TableOfMarksDihedral}}
\logpage{[ 70, 12, 2 ]}\nobreak
\hyperdef{L}{X7AADB47B8079C99E}{}
{\noindent\textcolor{FuncColor}{$\triangleright$\enspace\texttt{TableOfMarksDihedral({\mdseries\slshape n})\index{TableOfMarksDihedral@\texttt{TableOfMarksDihedral}}
\label{TableOfMarksDihedral}
}\hfill{\scriptsize (operation)}}\\


 \texttt{TableOfMarksDihedral} returns the table of marks of the dihedral group of order \mbox{\texttt{\mdseries\slshape m}}. 

 For each divisor $d$ of \mbox{\texttt{\mdseries\slshape m}}, a dihedral group of order $m = 2n$ contains subgroups of order $d$ according to the following rule. If $d$ is odd and divides $n$ then there is only one cyclic subgroup of order $d$. If $d$ is even and divides $n$ then there are a cyclic subgroup of order $d$ and two classes of dihedral subgroups of order $d$ (which are cyclic, too, in the case $d = 2$, see the example below). Otherwise (i.e., if $d$ does not divide $n$) there is just one class of dihedral subgroups of order $d$. 

 
\begin{Verbatim}[commandchars=!@|,fontsize=\small,frame=single,label=Example]
  !gapprompt@gap>| !gapinput@Display( TableOfMarksDihedral( 12 ) );|
   1:  12
   2:   6 6
   3:   6 . 2
   4:   6 . . 2
   5:   4 . . . 4
   6:   3 3 1 1 . 1
   7:   2 2 . . 2 . 2
   8:   2 . 2 . 2 . . 2
   9:   2 . . 2 2 . . . 2
  10:   1 1 1 1 1 1 1 1 1 1
  
\end{Verbatim}
 }

 

\subsection{\textcolor{Chapter }{TableOfMarksFrobenius}}
\logpage{[ 70, 12, 3 ]}\nobreak
\hyperdef{L}{X78E9DDF885E12687}{}
{\noindent\textcolor{FuncColor}{$\triangleright$\enspace\texttt{TableOfMarksFrobenius({\mdseries\slshape p, q})\index{TableOfMarksFrobenius@\texttt{TableOfMarksFrobenius}}
\label{TableOfMarksFrobenius}
}\hfill{\scriptsize (operation)}}\\


 \texttt{TableOfMarksFrobenius} computes the table of marks of a Frobenius group of order $p q$, where $p$ is a prime and $q$ divides $p-1$. 

 
\begin{Verbatim}[commandchars=!@|,fontsize=\small,frame=single,label=Example]
  !gapprompt@gap>| !gapinput@Display( TableOfMarksFrobenius( 5, 4 ) );|
  1:  20
  2:  10 2
  3:   5 1 1
  4:   4 . . 4
  5:   2 2 . 2 2
  6:   1 1 1 1 1 1
  
\end{Verbatim}
 }

 }

  
\section{\textcolor{Chapter }{The Library of Tables of Marks}}\label{The Library of Tables of Marks}
\logpage{[ 70, 13, 0 ]}
\hyperdef{L}{X794ABC7187A9285B}{}
{
  The \textsf{GAP} package \textsf{TomLib} provides access to several hundred tables of marks of almost simple groups and
their maximal subgroups. If this package is installed then the tables from
this database can be accessed via \texttt{TableOfMarks} (\ref{TableOfMarks:for a string}) with argument a string. If also the \textsf{GAP} Character Table Library is installed and contains the ordinary character table
of the group for which one wants to fetch the table of marks then one can also
call \texttt{TableOfMarks} (\ref{TableOfMarks:for a string})  with argument the character table. 

 A list of all names of tables of marks that are provided by the \textsf{TomLib} package can be obtained via \texttt{AllLibTomNames} (\textbf{TomLib: AllLibTomNames}). 

 
\begin{Verbatim}[commandchars=!@|,fontsize=\small,frame=single,label=Example]
  !gapprompt@gap>| !gapinput@names:= AllLibTomNames();;|
  !gapprompt@gap>| !gapinput@"A5" in names;|
  true
\end{Verbatim}
 }

 }

     
\chapter{\textcolor{Chapter }{Character Tables}}\label{Character Tables}
\logpage{[ 71, 0, 0 ]}
\hyperdef{L}{X7B7A9EE881E01C10}{}
{
  \index{tables} This chapter describes operations for \emph{character tables of finite groups}. 

 Operations for \emph{characters} (or, more general, \emph{class functions}) are described in Chapter{\nobreakspace}\ref{Class Functions}. 

 For a description of the \textsf{GAP} Library of Character Tables, see the separate manual for the \textsf{GAP} package \textsf{CTblLib}. 

 Several examples in this chapter require the \textsf{GAP} Character Table Library to be available. If it is not yet loaded then we load
it now. 

 
\begin{Verbatim}[commandchars=!@|,fontsize=\small,frame=single,label=Example]
  !gapprompt@gap>| !gapinput@LoadPackage( "ctbllib" );|
  true
\end{Verbatim}
  
\section{\textcolor{Chapter }{Some Remarks about Character Theory in \textsf{GAP}}}\label{Some Remarks about Character Theory in GAP}
\logpage{[ 71, 1, 0 ]}
\hyperdef{L}{X7B9FCBBC7B95F91B}{}
{
  It seems to be necessary to state some basic facts {\textendash}and maybe
warnings{\textendash} at the beginning of the character theory package. This
holds for people who are familiar with character theory because there is no
global reference on computational character theory, although there are many
papers on this topic, such as{\nobreakspace}\cite{NPP84} or{\nobreakspace}\cite{LP91}. It holds, however, also for people who are familiar with \textsf{GAP} because the general concept of domains (see Chapter{\nobreakspace}\ref{Domains}) plays no important role here {\textendash}we will justify this later in this
section. 

 Intuitively, \emph{characters} (or more generally, \emph{class functions}) of a finite group $G$ can be thought of as certain mappings defined on $G$, with values in the complex number field; the set of all characters of $G$ forms a semiring, with both addition and multiplication defined pointwise,
which is naturally embedded into the ring of \emph{generalized} (or \emph{virtual}) \emph{characters} in the natural way. A ${\ensuremath{\mathbb Z}}$\texttt{\symbol{45}}basis of this ring, and also a vector space basis of the
complex vector space of class functions of $G$, is given by the irreducible characters of $G$. 

 At this stage one could ask where there is a problem, since all these
algebraic structures are supported by \textsf{GAP}. But in practice, these structures are of minor importance, compared to
individual characters and the \emph{character tables} themselves (which are not domains in the sense of \textsf{GAP}). 

 For computations with characters of a finite group $G$ with $n$ conjugacy classes, we fix an ordering of the classes, and then identify each
class with its position according to this ordering. Each character of $G$ can be represented by a list of length $n$ in which the character value for elements of the $i$\texttt{\symbol{45}}th class is stored at the $i$\texttt{\symbol{45}}th position. Note that we need not know the conjugacy
classes of $G$ physically, even our knowledge of $G$ may be implicit in the sense that, e.g., we know how many classes of
involutions $G$ has, and which length these classes have, but we never have seen an element of $G$, or a presentation or representation of $G$. This allows us to work with the character tables of very large groups, e.g.,
of the so\texttt{\symbol{45}}called monster, where \textsf{GAP} has (currently) no chance to deal with the group. 

 As a consequence, also other information involving characters is given
implicitly. For example, we can talk about the kernel of a character not as a
group but as a list of classes (more exactly: a list of their positions
according to the chosen ordering of classes) forming this kernel; we can
deduce the group order, the contained cyclic subgroups and so on, but we do
not get the group itself. 

 So typical calculations with characters involve loops over lists of character
values. For example, the scalar product of two characters $\chi$, $\psi$ of $G$ given by 
\[ [ \chi, \psi ] = \left( \sum_{{g \in G}} \chi(g) \psi(g^{{-1}}) \right) / |G| \]
 can be written as 
\begin{Verbatim}[commandchars=!@|,fontsize=\small,frame=single,label=]
  Sum( [ 1 .. n ], i -> SizesConjugacyClasses( t )[i] * chi[i]
                            * ComplexConjugate( psi[i] ) ) / Size( t );
\end{Verbatim}
 where \texttt{t} is the character table of $G$, and \texttt{chi}, \texttt{psi} are the lists of values of $\chi$, $\psi$, respectively. 

 It is one of the advantages of character theory that after one has translated
a problem concerning groups into a problem concerning only characters, the
necessary calculations are mostly simple. For example, one can often prove
that a group is a Galois group over the rationals using calculations with
structure constants that can be computed from the character table, and
information about (the character tables of) maximal subgroups. When one deals
with such questions, the translation back to groups is just an interpretation
by the user, it does not take place in \textsf{GAP}. 

 \textsf{GAP} uses character \emph{tables} to store information such as class lengths, element orders, the irreducible
characters of $G$ etc.{\nobreakspace}in a consistent way; in the example above, we have seen
that \texttt{SizesConjugacyClasses} (\ref{SizesConjugacyClasses}) returns the list of class lengths of its argument. Note that the values of
these attributes rely on the chosen ordering of conjugacy classes, a character
table is not determined by something similar to generators of groups or rings
in \textsf{GAP} where knowledge could in principle be recovered from the generators but is
stored mainly for the sake of efficiency. 

 Note that the character table of a group $G$ in \textsf{GAP} must \emph{not} be mixed up with the list of complex irreducible characters of $G$. The irreducible characters are stored in a character table via the attribute \texttt{Irr} (\ref{Irr:for a group}). 

 Two further important instances of information that depends on the ordering of
conjugacy classes are \emph{power maps} and \emph{fusion maps}. Both are represented as lists of integers in \textsf{GAP}. The $k$\texttt{\symbol{45}}th power map maps each class to the class of $k$\texttt{\symbol{45}}th powers of its elements, the corresponding list contains
at each position the position of the image. A class fusion map between the
classes of a subgroup $H$ of $G$ and the classes of $G$ maps each class $c$ of $H$ to that class of $G$ that contains $c$, the corresponding list contains again the positions of image classes; if we
know only the character tables of $H$ and $G$ but not the groups themselves, this means with respect to a fixed embedding of $H$ into $G$. More about power maps and fusion maps can be found in Chapter{\nobreakspace}\ref{Maps Concerning Character Tables}. 

 So class functions, power maps, and fusion maps are represented by lists in \textsf{GAP}. If they are plain lists then they are regarded as class functions
etc.{\nobreakspace}of an appropriate character table when they are passed to \textsf{GAP} functions that expect class functions etc. For example, a list with all
entries equal to 1 is regarded as the trivial character if it is passed to a
function that expects a character. Note that this approach requires the
character table as an argument for such a function. 

 One can construct class function objects that store their underlying character
table and other attribute values (see Chapter{\nobreakspace}\ref{Class Functions}). This allows one to omit the character table argument in many functions, and
it allows one to use infix operations for tensoring or inducing class
functions. }

  
\section{\textcolor{Chapter }{History of Character Theory Stuff in GAP}}\label{History of Character Theory Stuff in GAP}
\logpage{[ 71, 2, 0 ]}
\hyperdef{L}{X7F8AB7CB7A46002F}{}
{
  \textsf{GAP} provides functions for dealing with group characters since the version \textsf{GAP}{\nobreakspace}3.1, which was released in March 1992. The reason for adding
this branch of mathematics to the topics of \textsf{GAP} was (apart from the usefulness of character theoretic computations in general)
the insight that \textsf{GAP} provides an ideal environment for developing the algorithms needed. In
particular, it had been decided at Lehrstuhl D f{\"u}r Mathematik that the \textsf{CAS} system (a standalone Fortran program together with a database of character
tables, see{\nobreakspace}\cite{NPP84}) should not be developed further and the functionality of \textsf{CAS} should be made available in \textsf{GAP}. The background was that extending \textsf{CAS} (by new Fortran code) had turned out to be much less flexible than writing
analogous \textsf{GAP} library code. 

 For integrating the existing character theory algorithms, \textsf{GAP}'s memory management and long integer arithmetic were useful as well as the
list handling {\textendash}it is an important feature of character theoretic
methods that questions about groups are translated into manipulations of
lists; on the other hand, the datatype of cyclotomics (see
Chapter{\nobreakspace}\texttt{Cyclotomics} (\ref{Cyclotomics})) was added to the \textsf{GAP} kernel because of the character theory algorithms. For developing further
code, also other areas of \textsf{GAP}'s library became interesting, such as permutation groups, finite fields, and
polynomials. 

 The development of character theory code for \textsf{GAP} has been supported by several DFG grants, in particular the project ``Representation Theory of Finite Groups and Finite Dimensional Algebras'' (until 1991), and the Schwerpunkt ``Algorithmische Zahlentheorie und Algebra'' (from 1991 until 1997). Besides that, several Diploma theses at Lehrstuhl D
were concerned with the development and/or implementation of algorithms
dealing with characters in \textsf{GAP}. 

 The major contributions can be listed as follows. 
\begin{itemize}
\item  The arithmetic for the cyclotomics data type, following{\nobreakspace}\cite{Zum89}, was first implemented by Marco van Meegen; an alternative approach was
studied in the diploma thesis of Michael Scherner (see{\nobreakspace}\cite{Scherner92}) but was not efficient enough; later Martin Sch{\"o}nert replaced the
implementation by a better one.  
\item  The basic routines for characters and character tables were written by Thomas
Breuer and G{\"o}tz Pfeiffer. 
\item  The lattice related functions, such as \texttt{LLL} (\ref{LLL}), \texttt{OrthogonalEmbeddings} (\ref{OrthogonalEmbeddings}), and \texttt{DnLattice} (\ref{DnLattice}), were implemented by Ansgar Kaup (see{\nobreakspace}\cite{Kaup92}). 
\item  Functions for computing possible class fusions, possible power maps, and table
automorphisms were written by Thomas Breuer (see{\nobreakspace}\cite{Bre91}). 
\item  Functions for computing possible permutation characters were written by Thomas
Breuer (see{\nobreakspace}\cite{Bre91}) and G{\"o}tz Pfeiffer (see{\nobreakspace}\cite{Pfe91}). 
\item  Functions for computing character tables from groups were written by Alexander
Hulpke (Dixon\texttt{\symbol{45}}Schneider algorithm, see{\nobreakspace}\cite{Hulpke93}) and Hans Ulrich Besche (Baum algorithm and Conlon algorithm,
see{\nobreakspace}\cite{Besche92}). 
\item  Functions for dealing with Clifford matrices were written by Ute Schiffer
(see{\nobreakspace}\cite{Schiffer94}). 
\item  Functions for monomiality questions were written by Thomas Breuer and
Erzs{\a'e}bet Horv{\a'a}th. 
\end{itemize}
 

 Since then, the code has been maintained and extended further by Alexander
Hulpke (code related to his implementation of the
Dixon\texttt{\symbol{45}}Schneider algorithm) and Thomas Breuer. 

  Currently \textsf{GAP} does not provide special functionality for computing Brauer character tables,
but there is an interface to the \textsf{MOC} system (see{\nobreakspace}\cite{HJLP92}), and the \textsf{GAP} Character Table Library contains many known Brauer character tables. }

  
\section{\textcolor{Chapter }{Creating Character Tables}}\label{Creating Character Tables}
\logpage{[ 71, 3, 0 ]}
\hyperdef{L}{X8701174D86B586AF}{}
{
  \index{tables} \index{character tables} \index{library tables} \index{character tables!access to} \index{character tables!calculate} \index{character tables!of groups} There are in general five different ways to get a character table in \textsf{GAP}. You can 

 
\begin{enumerate}
\item  compute the table from a group, 
\item  read a file that contains the table data, 
\item  construct the table using generic formulae, 
\item  derive it from known character tables, or 
\item  combine partial information about conjugacy classes, power maps of the group
in question, and about (character tables of) some subgroups and supergroups. 
\end{enumerate}
 

 In 1., the computation of the irreducible characters is the hardest part; the
different algorithms available for this are described in{\nobreakspace}\ref{Computing the Irreducible Characters of a Group}. Possibility 2.{\nobreakspace}is used for the character tables in the \textsf{GAP} Character Table Library, see the manual of this library. Generic character
tables {\textendash}as addressed by 3.{\textendash} are described
in{\nobreakspace} (\textbf{CTblLib: Generic Character Tables}). Several occurrences of 4.{\nobreakspace}are described in{\nobreakspace}\ref{Constructing Character Tables from Others}. The last of the above possibilities \emph{is currently not supported and will be described in a chapter of its own when
it becomes available}. 

 The operation \texttt{CharacterTable} (\ref{CharacterTable:for a group}) can be used for the cases 1. to 3. 
\subsection{\textcolor{Chapter }{CharacterTable}}\logpage{[ 71, 3, 1 ]}
\hyperdef{L}{X7FCA7A7A822BDA33}{}
{
\noindent\textcolor{FuncColor}{$\triangleright$\enspace\texttt{CharacterTable({\mdseries\slshape G[, p]})\index{CharacterTable@\texttt{CharacterTable}!for a group}
\label{CharacterTable:for a group}
}\hfill{\scriptsize (operation)}}\\
\noindent\textcolor{FuncColor}{$\triangleright$\enspace\texttt{CharacterTable({\mdseries\slshape ordtbl, p})\index{CharacterTable@\texttt{CharacterTable}!for an ordinary character table}
\label{CharacterTable:for an ordinary character table}
}\hfill{\scriptsize (operation)}}\\
\noindent\textcolor{FuncColor}{$\triangleright$\enspace\texttt{CharacterTable({\mdseries\slshape name[, param]})\index{CharacterTable@\texttt{CharacterTable}!for a string}
\label{CharacterTable:for a string}
}\hfill{\scriptsize (operation)}}\\


 Called with a group \mbox{\texttt{\mdseries\slshape G}}, \texttt{CharacterTable} calls the attribute \texttt{OrdinaryCharacterTable} (\ref{OrdinaryCharacterTable:for a group}). Called with first argument a group \mbox{\texttt{\mdseries\slshape G}} or an ordinary character table \mbox{\texttt{\mdseries\slshape ordtbl}}, and second argument a prime \mbox{\texttt{\mdseries\slshape p}}, \texttt{CharacterTable} calls the operation \texttt{BrauerTable} (\ref{BrauerTable:for a group, and a prime integer}). 

 Called with a string \mbox{\texttt{\mdseries\slshape name}} and perhaps optional parameters \mbox{\texttt{\mdseries\slshape param}}, \texttt{CharacterTable} tries to access a character table from the \textsf{GAP} Character Table Library. See the manual of the \textsf{GAP} package \textsf{CTblLib} for an overview of admissible arguments. An error is signalled if this \textsf{GAP} package is not loaded in this case. 

 Probably the most interesting information about the character table is its
list of irreducibles, which can be accessed as the value of the attribute \texttt{Irr} (\ref{Irr:for a character table}). If the argument of \texttt{CharacterTable} is a string \mbox{\texttt{\mdseries\slshape name}} then the irreducibles are just read from the library file, therefore the
returned table stores them already. However, if \texttt{CharacterTable} is called with a group \mbox{\texttt{\mdseries\slshape G}} or with an ordinary character table \mbox{\texttt{\mdseries\slshape ordtbl}}, the irreducible characters are \emph{not} computed by \texttt{CharacterTable}. They are only computed when the \texttt{Irr} (\ref{Irr:for a character table}) value is accessed for the first time, for example when \texttt{Display} (\ref{Display}) is called for the table (see{\nobreakspace}\ref{Printing Character Tables}). This means for example that \texttt{CharacterTable} returns its result very quickly, and the first call of \texttt{Display} (\ref{Display}) for this table may take some time because the irreducible characters must be
computed at that time before they can be displayed together with other
information stored on the character table. The value of the filter \texttt{HasIrr} indicates whether the irreducible characters have been computed already. 

 The reason why \texttt{CharacterTable} does not compute the irreducible characters is that there are situations where
one only needs the ``table head'', that is, the information about class lengths, power maps etc., but not the
irreducibles. For example, if one wants to inspect permutation characters of a
group then all one has to do is to induce the trivial characters of subgroups
one is interested in; for that, only class lengths and the class fusion are
needed. Or if one wants to compute the Molien series (see{\nobreakspace}\texttt{MolienSeries} (\ref{MolienSeries})) for a given complex matrix group, the irreducible characters of this group
are in general of no interest. 

 For details about different algorithms to compute the irreducible characters,
see{\nobreakspace}\ref{Computing the Irreducible Characters of a Group}. 

 If the group \mbox{\texttt{\mdseries\slshape G}} is given as an argument, \texttt{CharacterTable} accesses the conjugacy classes of \mbox{\texttt{\mdseries\slshape G}} and therefore causes that these classes are computed if they were not yet
stored (see \ref{The Interface between Character Tables and Groups}). }

 
\subsection{\textcolor{Chapter }{BrauerTable}}\logpage{[ 71, 3, 2 ]}
\hyperdef{L}{X8476B25A79D7A7FC}{}
{
\noindent\textcolor{FuncColor}{$\triangleright$\enspace\texttt{BrauerTable({\mdseries\slshape ordtbl, p})\index{BrauerTable@\texttt{BrauerTable}!for a character table, and a prime integer}
\label{BrauerTable:for a character table, and a prime integer}
}\hfill{\scriptsize (operation)}}\\
\noindent\textcolor{FuncColor}{$\triangleright$\enspace\texttt{BrauerTable({\mdseries\slshape G, p})\index{BrauerTable@\texttt{BrauerTable}!for a group, and a prime integer}
\label{BrauerTable:for a group, and a prime integer}
}\hfill{\scriptsize (operation)}}\\
\noindent\textcolor{FuncColor}{$\triangleright$\enspace\texttt{BrauerTableOp({\mdseries\slshape ordtbl, p})\index{BrauerTableOp@\texttt{BrauerTableOp}}
\label{BrauerTableOp}
}\hfill{\scriptsize (operation)}}\\
\noindent\textcolor{FuncColor}{$\triangleright$\enspace\texttt{ComputedBrauerTables({\mdseries\slshape ordtbl})\index{ComputedBrauerTables@\texttt{ComputedBrauerTables}}
\label{ComputedBrauerTables}
}\hfill{\scriptsize (attribute)}}\\


 Called with an ordinary character table \mbox{\texttt{\mdseries\slshape ordtbl}} or a group \mbox{\texttt{\mdseries\slshape G}}, \texttt{BrauerTable} returns its \mbox{\texttt{\mdseries\slshape p}}\texttt{\symbol{45}}modular character table if \textsf{GAP} can compute this table, and \texttt{fail} otherwise. 

 The \mbox{\texttt{\mdseries\slshape p}}\texttt{\symbol{45}}modular table can be computed in the following cases. 

 
\begin{itemize}
\item  The group is \mbox{\texttt{\mdseries\slshape p}}\texttt{\symbol{45}}solvable (see \texttt{IsPSolvable} (\ref{IsPSolvable}), apply the Fong\texttt{\symbol{45}}Swan Theorem); 
\item  the Sylow \mbox{\texttt{\mdseries\slshape p}}\texttt{\symbol{45}}subgroup of \mbox{\texttt{\mdseries\slshape G}} is cyclic, and all \mbox{\texttt{\mdseries\slshape p}}\texttt{\symbol{45}}modular Brauer characters of \mbox{\texttt{\mdseries\slshape G}} lift to ordinary characters (note that this situation can be detected from the
ordinary character table of \mbox{\texttt{\mdseries\slshape G}}); 
\item  the table \mbox{\texttt{\mdseries\slshape ordtbl}} stores information how it was constructed from other tables (as a direct
product or as an isoclinic variant, for example), and the Brauer tables of the
source tables can be computed; 
\item  \mbox{\texttt{\mdseries\slshape ordtbl}} is a table from the \textsf{GAP} character table library for which also the \mbox{\texttt{\mdseries\slshape p}}\texttt{\symbol{45}}modular table is contained in the table library. 
\end{itemize}
 

 The default method for a group and a prime delegates to \texttt{BrauerTable} for the ordinary character table of this group. The default method for \mbox{\texttt{\mdseries\slshape ordtbl}} uses the attribute \texttt{ComputedBrauerTables} for storing the computed Brauer table at position \mbox{\texttt{\mdseries\slshape p}}, and calls the operation \texttt{BrauerTableOp} for computing values that are not yet known. 

 So if one wants to install a new method for computing Brauer tables then it is
sufficient to install it for \texttt{BrauerTableOp}. 

 The \texttt{mod} operator for a character table and a prime (see{\nobreakspace}\ref{Operators for Character Tables}) delegates to \texttt{BrauerTable}. }

 

\subsection{\textcolor{Chapter }{CharacterTableRegular}}
\logpage{[ 71, 3, 3 ]}\nobreak
\hyperdef{L}{X85DB8AE7786A2DB5}{}
{\noindent\textcolor{FuncColor}{$\triangleright$\enspace\texttt{CharacterTableRegular({\mdseries\slshape tbl, p})\index{CharacterTableRegular@\texttt{CharacterTableRegular}}
\label{CharacterTableRegular}
}\hfill{\scriptsize (function)}}\\


 For an ordinary character table \mbox{\texttt{\mdseries\slshape tbl}} and a prime integer \mbox{\texttt{\mdseries\slshape p}}, \texttt{CharacterTableRegular} returns the ``table head'' of the \mbox{\texttt{\mdseries\slshape p}}\texttt{\symbol{45}}modular Brauer character table of \mbox{\texttt{\mdseries\slshape tbl}}. This is the restriction of \mbox{\texttt{\mdseries\slshape tbl}} to its \mbox{\texttt{\mdseries\slshape p}}\texttt{\symbol{45}}regular classes, like the return value of \texttt{BrauerTable} (\ref{BrauerTable:for a character table, and a prime integer}), but without the irreducible Brauer characters. (In general, these characters
are hard to compute, and \texttt{BrauerTable} (\ref{BrauerTable:for a character table, and a prime integer}) may return \texttt{fail} for the given arguments, for example if \mbox{\texttt{\mdseries\slshape tbl}} is a table from the \textsf{GAP} character table library.) 

 The returned table head can be used to create \mbox{\texttt{\mdseries\slshape p}}\texttt{\symbol{45}}modular Brauer characters, by restricting ordinary
characters, for example when one is interested in approximations of the
(unknown) irreducible Brauer characters. 

 
\begin{Verbatim}[commandchars=!@|,fontsize=\small,frame=single,label=Example]
  !gapprompt@gap>| !gapinput@g:= SymmetricGroup( 4 );|
  Sym( [ 1 .. 4 ] )
  !gapprompt@gap>| !gapinput@tbl:= CharacterTable( g );;  HasIrr( tbl );|
  false
  !gapprompt@gap>| !gapinput@tblmod2:= CharacterTable( tbl, 2 );|
  BrauerTable( Sym( [ 1 .. 4 ] ), 2 )
  !gapprompt@gap>| !gapinput@tblmod2 = CharacterTable( tbl, 2 );|
  true
  !gapprompt@gap>| !gapinput@tblmod2 = BrauerTable( tbl, 2 );|
  true
  !gapprompt@gap>| !gapinput@tblmod2 = BrauerTable( g, 2 );|
  true
  !gapprompt@gap>| !gapinput@libtbl:= CharacterTable( "M" );|
  CharacterTable( "M" )
  !gapprompt@gap>| !gapinput@CharacterTableRegular( libtbl, 2 );|
  BrauerTable( "M", 2 )
  !gapprompt@gap>| !gapinput@BrauerTable( libtbl, 2 );|
  fail
  !gapprompt@gap>| !gapinput@CharacterTable( "Symmetric", 4 );|
  CharacterTable( "Sym(4)" )
  !gapprompt@gap>| !gapinput@ComputedBrauerTables( tbl );|
  [ , BrauerTable( Sym( [ 1 .. 4 ] ), 2 ) ]
\end{Verbatim}
 }

 

\subsection{\textcolor{Chapter }{SupportedCharacterTableInfo}}
\logpage{[ 71, 3, 4 ]}\nobreak
\hyperdef{L}{X7DBEF4BF87F10CD6}{}
{\noindent\textcolor{FuncColor}{$\triangleright$\enspace\texttt{SupportedCharacterTableInfo\index{SupportedCharacterTableInfo@\texttt{SupportedCharacterTableInfo}}
\label{SupportedCharacterTableInfo}
}\hfill{\scriptsize (global variable)}}\\


 \texttt{SupportedCharacterTableInfo} is a list that contains at position $3i-2$ an attribute getter function, at position $3i-1$ the name of this attribute, and at position $3i$ a list containing a subset of \texttt{[ "character", "class", "mutable" ]}, depending on whether the attribute value relies on the ordering of
characters or classes, or whether the attribute value is a mutable list or
record. 

 When (ordinary or Brauer) character table objects are created from records,
using \texttt{ConvertToCharacterTable} (\ref{ConvertToCharacterTable}), \texttt{SupportedCharacterTableInfo} specifies those record components that shall be used as attribute values;
other record components are \emph{not} be regarded as attribute values in the conversion process. 

 New attributes and properties can be notified to \texttt{SupportedCharacterTableInfo} by creating them with \texttt{DeclareAttributeSuppCT} and \texttt{DeclarePropertySuppCT} instead of \texttt{DeclareAttribute} (\ref{DeclareAttribute}) and \texttt{DeclareProperty} (\ref{DeclareProperty}). }

 

\subsection{\textcolor{Chapter }{ConvertToCharacterTable}}
\logpage{[ 71, 3, 5 ]}\nobreak
\hyperdef{L}{X8195BC057B1DFAD5}{}
{\noindent\textcolor{FuncColor}{$\triangleright$\enspace\texttt{ConvertToCharacterTable({\mdseries\slshape record})\index{ConvertToCharacterTable@\texttt{ConvertToCharacterTable}}
\label{ConvertToCharacterTable}
}\hfill{\scriptsize (function)}}\\
\noindent\textcolor{FuncColor}{$\triangleright$\enspace\texttt{ConvertToCharacterTableNC({\mdseries\slshape record})\index{ConvertToCharacterTableNC@\texttt{ConvertToCharacterTableNC}}
\label{ConvertToCharacterTableNC}
}\hfill{\scriptsize (function)}}\\


 Let \mbox{\texttt{\mdseries\slshape record}} be a record. \texttt{ConvertToCharacterTable} converts \mbox{\texttt{\mdseries\slshape record}} into a component object (see{\nobreakspace}\ref{Component Objects}) representing a character table. The values of those components of \mbox{\texttt{\mdseries\slshape record}} whose names occur in \texttt{SupportedCharacterTableInfo} (\ref{SupportedCharacterTableInfo}) correspond to attribute values of the returned character table. All other
components of the record simply become components of the character table
object. 

 If inconsistencies in \mbox{\texttt{\mdseries\slshape record}} are detected, \texttt{fail} is returned. \mbox{\texttt{\mdseries\slshape record}} must have the component \texttt{UnderlyingCharacteristic} bound (cf.{\nobreakspace}\texttt{UnderlyingCharacteristic} (\ref{UnderlyingCharacteristic:for a character table})), since this decides about whether the returned character table lies in \texttt{IsOrdinaryTable} (\ref{IsOrdinaryTable}) or in \texttt{IsBrauerTable} (\ref{IsBrauerTable}). 

 \texttt{ConvertToCharacterTableNC} does the same except that all checks of \mbox{\texttt{\mdseries\slshape record}} are omitted. 

 An example of a conversion from a record to a character table object can be
found in Section{\nobreakspace}\texttt{PrintCharacterTable} (\ref{PrintCharacterTable}). }

 }

  
\section{\textcolor{Chapter }{Character Table Categories}}\label{Character Table Categories}
\logpage{[ 71, 4, 0 ]}
\hyperdef{L}{X789FAC077AEF088A}{}
{
  

\subsection{\textcolor{Chapter }{IsNearlyCharacterTable}}
\logpage{[ 71, 4, 1 ]}\nobreak
\hyperdef{L}{X82FF82C87CF82ADF}{}
{\noindent\textcolor{FuncColor}{$\triangleright$\enspace\texttt{IsNearlyCharacterTable({\mdseries\slshape obj})\index{IsNearlyCharacterTable@\texttt{IsNearlyCharacterTable}}
\label{IsNearlyCharacterTable}
}\hfill{\scriptsize (Category)}}\\
\noindent\textcolor{FuncColor}{$\triangleright$\enspace\texttt{IsCharacterTable({\mdseries\slshape obj})\index{IsCharacterTable@\texttt{IsCharacterTable}}
\label{IsCharacterTable}
}\hfill{\scriptsize (Category)}}\\
\noindent\textcolor{FuncColor}{$\triangleright$\enspace\texttt{IsOrdinaryTable({\mdseries\slshape obj})\index{IsOrdinaryTable@\texttt{IsOrdinaryTable}}
\label{IsOrdinaryTable}
}\hfill{\scriptsize (Category)}}\\
\noindent\textcolor{FuncColor}{$\triangleright$\enspace\texttt{IsBrauerTable({\mdseries\slshape obj})\index{IsBrauerTable@\texttt{IsBrauerTable}}
\label{IsBrauerTable}
}\hfill{\scriptsize (Category)}}\\
\noindent\textcolor{FuncColor}{$\triangleright$\enspace\texttt{IsCharacterTableInProgress({\mdseries\slshape obj})\index{IsCharacterTableInProgress@\texttt{IsCharacterTableInProgress}}
\label{IsCharacterTableInProgress}
}\hfill{\scriptsize (Category)}}\\


 Every ``character table like object'' in \textsf{GAP} lies in the category \texttt{IsNearlyCharacterTable}. There are four important subcategories, namely the \emph{ordinary} tables in \texttt{IsOrdinaryTable}, the \emph{Brauer} tables in \texttt{IsBrauerTable}, the union of these two in \texttt{IsCharacterTable}, and the \emph{incomplete ordinary} tables in \texttt{IsCharacterTableInProgress}. 

 We want to distinguish ordinary and Brauer tables because a Brauer table may
delegate tasks to the ordinary table of the same group, for example the
computation of power maps. A Brauer table is constructed from an ordinary
table and stores this table upon construction (see{\nobreakspace}\texttt{OrdinaryCharacterTable} (\ref{OrdinaryCharacterTable:for a group})). 

 Furthermore, \texttt{IsOrdinaryTable} and \texttt{IsBrauerTable} denote character tables that provide enough information to compute all power
maps and irreducible characters (and in the case of Brauer tables to get the
ordinary table), for example because the underlying group (see{\nobreakspace}\texttt{UnderlyingGroup} (\ref{UnderlyingGroup:for character tables})) is known or because the table is a library table (see the manual of the \textsf{GAP} Character Table Library). We want to distinguish these tables from partially
known ordinary tables that cannot be asked for all power maps or all
irreducible characters. 

 The character table objects in \texttt{IsCharacterTable} are always immutable (see{\nobreakspace}\ref{Mutability and Copyability}). This means mainly that the ordering of conjugacy classes used for the
various attributes of the character table cannot be changed;
see{\nobreakspace}\ref{Sorted Character Tables} for how to compute a character table with a different ordering of classes. 

 The \textsf{GAP} objects in \texttt{IsCharacterTableInProgress} represent incomplete ordinary character tables. This means that not all
irreducible characters, not all power maps are known, and perhaps even the
number of classes and the centralizer orders are known. Such tables occur when
the character table of a group $G$ is constructed using character tables of related groups and information about $G$ but for example without explicitly computing the conjugacy classes of $G$. An object in \texttt{IsCharacterTableInProgress} is first of all \emph{mutable}, so \emph{nothing is stored automatically} on such a table, since otherwise one has no control of
side\texttt{\symbol{45}}effects when a hypothesis is changed. Operations for
such tables may return more general values than for other tables, for example
class functions may contain unknowns (see Chapter{\nobreakspace}\ref{Unknowns}) or lists of possible values in certain positions, the same may happen also
for power maps and class fusions (see{\nobreakspace}\ref{Parametrized Maps}). \emph{Incomplete tables in this sense are currently not supported and will be
described in a chapter of their own when they become available.} Note that the term ``incomplete table'' shall express that \textsf{GAP} cannot compute certain values such as irreducible characters or power maps. A
table with access to its group is therefore always complete, also if its
irreducible characters are not yet stored. 

 
\begin{Verbatim}[commandchars=!@|,fontsize=\small,frame=single,label=Example]
  !gapprompt@gap>| !gapinput@g:= SymmetricGroup( 4 );;|
  !gapprompt@gap>| !gapinput@tbl:= CharacterTable( g );  modtbl:= tbl mod 2;|
  CharacterTable( Sym( [ 1 .. 4 ] ) )
  BrauerTable( Sym( [ 1 .. 4 ] ), 2 )
  !gapprompt@gap>| !gapinput@IsCharacterTable( tbl );  IsCharacterTable( modtbl );|
  true
  true
  !gapprompt@gap>| !gapinput@IsBrauerTable( modtbl );  IsBrauerTable( tbl );|
  true
  false
  !gapprompt@gap>| !gapinput@IsOrdinaryTable( tbl );  IsOrdinaryTable( modtbl );|
  true
  false
  !gapprompt@gap>| !gapinput@IsCharacterTable( g );  IsCharacterTable( Irr( g ) );|
  false
  false
\end{Verbatim}
 }

 

\subsection{\textcolor{Chapter }{InfoCharacterTable}}
\logpage{[ 71, 4, 2 ]}\nobreak
\hyperdef{L}{X7C6F3D947E5188D1}{}
{\noindent\textcolor{FuncColor}{$\triangleright$\enspace\texttt{InfoCharacterTable\index{InfoCharacterTable@\texttt{InfoCharacterTable}}
\label{InfoCharacterTable}
}\hfill{\scriptsize (info class)}}\\


 is the info class (see{\nobreakspace}\ref{Info Functions}) for computations with character tables. }

 

\subsection{\textcolor{Chapter }{NearlyCharacterTablesFamily}}
\logpage{[ 71, 4, 3 ]}\nobreak
\hyperdef{L}{X7FA867637EBB30F9}{}
{\noindent\textcolor{FuncColor}{$\triangleright$\enspace\texttt{NearlyCharacterTablesFamily\index{NearlyCharacterTablesFamily@\texttt{NearlyCharacterTablesFamily}}
\label{NearlyCharacterTablesFamily}
}\hfill{\scriptsize (family)}}\\


 Every character table like object lies in this family (see{\nobreakspace}\ref{Families}). }

 }

  
\section{\textcolor{Chapter }{Conventions for Character Tables}}\label{Conventions for Character Tables}
\logpage{[ 71, 5, 0 ]}
\hyperdef{L}{X829C4B6E83998F40}{}
{
  The following few conventions should be noted. 
\begin{itemize}
\item  The class of the \emph{identity element} is expected to be the first one; thus the degree of a character is the
character value at position $1$. 
\item  The \emph{trivial character} of a character table need not be the first in the list of irreducibles. 
\item  Most functions that take a character table as an argument and work with
characters expect these characters as an argument, too. For some functions,
the list of irreducible characters serves as the default, i.e, the value of
the attribute \texttt{Irr} (\ref{Irr:for a character table}); in these cases, the \texttt{Irr} (\ref{Irr:for a character table}) value is automatically computed if it was not yet known. 
\item  For a stored class fusion, the image table is denoted by its \texttt{Identifier} (\ref{Identifier:for character tables}) value; each library table has a unique identifier by which it can be accessed
(see{\nobreakspace} (\textbf{CTblLib: Accessing a Character Table from the Library}) in the manual for the \textsf{GAP} Character Table Library), tables constructed from groups get an identifier
that is unique in the current \textsf{GAP} session. 
\end{itemize}
 }

  
\section{\textcolor{Chapter }{The Interface between Character Tables and Groups}}\label{The Interface between Character Tables and Groups}
\logpage{[ 71, 6, 0 ]}
\hyperdef{L}{X793E0EBF84B07313}{}
{
  For a character table with underlying group (see{\nobreakspace}\texttt{UnderlyingGroup} (\ref{UnderlyingGroup:for character tables})), the interface between table and group consists of three attribute values,
namely the \emph{group}, the \emph{conjugacy classes} stored in the table (see \texttt{ConjugacyClasses} (\ref{ConjugacyClasses:for character tables}) below) and the \emph{identification} of the conjugacy classes of table and group (see{\nobreakspace}\texttt{IdentificationOfConjugacyClasses} (\ref{IdentificationOfConjugacyClasses}) below). 

 Character tables constructed from groups know these values upon construction,
and for character tables constructed without groups, these values are usually
not known and cannot be computed from the table. 

 However, given a group $G$ and a character table of a group isomorphic to $G$ (for example a character table from the \textsf{GAP} table library), one can tell \textsf{GAP} to compute a new instance of the given table and to use it as the character
table of $G$ (see{\nobreakspace}\texttt{CharacterTableWithStoredGroup} (\ref{CharacterTableWithStoredGroup})). 

 Tasks may be delegated from a group to its character table or vice versa only
if these three attribute values are stored in the character table. 

\subsection{\textcolor{Chapter }{UnderlyingGroup (for character tables)}}
\logpage{[ 71, 6, 1 ]}\nobreak
\hyperdef{L}{X7FF4826A82B667AF}{}
{\noindent\textcolor{FuncColor}{$\triangleright$\enspace\texttt{UnderlyingGroup({\mdseries\slshape ordtbl})\index{UnderlyingGroup@\texttt{UnderlyingGroup}!for character tables}
\label{UnderlyingGroup:for character tables}
}\hfill{\scriptsize (attribute)}}\\


 For an ordinary character table \mbox{\texttt{\mdseries\slshape ordtbl}} of a finite group, the group can be stored as value of \texttt{UnderlyingGroup}. 

 Brauer tables do not store the underlying group, they access it via the
ordinary table (see{\nobreakspace}\texttt{OrdinaryCharacterTable} (\ref{OrdinaryCharacterTable:for a character table})). }

 

\subsection{\textcolor{Chapter }{ConjugacyClasses (for character tables)}}
\logpage{[ 71, 6, 2 ]}\nobreak
\hyperdef{L}{X849A38F887F6EC86}{}
{\noindent\textcolor{FuncColor}{$\triangleright$\enspace\texttt{ConjugacyClasses({\mdseries\slshape tbl})\index{ConjugacyClasses@\texttt{ConjugacyClasses}!for character tables}
\label{ConjugacyClasses:for character tables}
}\hfill{\scriptsize (attribute)}}\\


 For a character table \mbox{\texttt{\mdseries\slshape tbl}} with known underlying group $G$, the \texttt{ConjugacyClasses} value of \mbox{\texttt{\mdseries\slshape tbl}} is a list of conjugacy classes of $G$. All those lists stored in the table that are related to the ordering of
conjugacy classes (such as sizes of centralizers and conjugacy classes, orders
of representatives, power maps, and all class functions) refer to the ordering
of this list. 

 This ordering need \emph{not} coincide with the ordering of conjugacy classes as stored in the underlying
group of the table (see{\nobreakspace}\ref{Sorted Character Tables}). One reason for this is that otherwise we would not be allowed to use a
library table as the character table of a group for which the conjugacy
classes are stored already. (Another, less important reason is that we can use
the same group as underlying group of character tables that differ only
w.r.t.{\nobreakspace}the ordering of classes.) 

 The class of the identity element must be the first class (see{\nobreakspace}\ref{Conventions for Character Tables}). 

 If \mbox{\texttt{\mdseries\slshape tbl}} was constructed from $G$ then the conjugacy classes have been stored at the same time when $G$ was stored. If $G$ and \mbox{\texttt{\mdseries\slshape tbl}} have been connected later than in the construction of \mbox{\texttt{\mdseries\slshape tbl}}, the recommended way to do this is via \texttt{CharacterTableWithStoredGroup} (\ref{CharacterTableWithStoredGroup}). So there is no method for \texttt{ConjugacyClasses} that computes the value for \mbox{\texttt{\mdseries\slshape tbl}} if it is not yet stored. 

 Brauer tables do not store the ($p$\texttt{\symbol{45}}regular) conjugacy classes, they access them via the
ordinary table (see{\nobreakspace}\texttt{OrdinaryCharacterTable} (\ref{OrdinaryCharacterTable:for a character table})) if necessary. }

 

\subsection{\textcolor{Chapter }{IdentificationOfConjugacyClasses}}
\logpage{[ 71, 6, 3 ]}\nobreak
\hyperdef{L}{X84DC12AA804C8085}{}
{\noindent\textcolor{FuncColor}{$\triangleright$\enspace\texttt{IdentificationOfConjugacyClasses({\mdseries\slshape tbl})\index{IdentificationOfConjugacyClasses@\texttt{IdentificationOfConjugacyClasses}}
\label{IdentificationOfConjugacyClasses}
}\hfill{\scriptsize (attribute)}}\\


 For an ordinary character table \mbox{\texttt{\mdseries\slshape tbl}} with known underlying group $G$, \texttt{IdentificationOfConjugacyClasses} returns a list of positive integers that contains at position $i$ the position of the $i$\texttt{\symbol{45}}th conjugacy class of \mbox{\texttt{\mdseries\slshape tbl}} in the \texttt{ConjugacyClasses} (\ref{ConjugacyClasses:for character tables}) value of $G$. 

 
\begin{Verbatim}[commandchars=!@|,fontsize=\small,frame=single,label=Example]
  !gapprompt@gap>| !gapinput@g:= SymmetricGroup( 4 );;|
  !gapprompt@gap>| !gapinput@repres:= [ (1,2), (1,2,3), (1,2,3,4), (1,2)(3,4), () ];;|
  !gapprompt@gap>| !gapinput@ccl:= List( repres, x -> ConjugacyClass( g, x ) );;|
  !gapprompt@gap>| !gapinput@SetConjugacyClasses( g, ccl );|
  !gapprompt@gap>| !gapinput@tbl:= CharacterTable( g );;   # the table stores already the values|
  !gapprompt@gap>| !gapinput@HasConjugacyClasses( tbl );  HasUnderlyingGroup( tbl );|
  true
  true
  !gapprompt@gap>| !gapinput@UnderlyingGroup( tbl ) = g;|
  true
  !gapprompt@gap>| !gapinput@HasIdentificationOfConjugacyClasses( tbl );|
  true
  !gapprompt@gap>| !gapinput@IdentificationOfConjugacyClasses( tbl );|
  [ 5, 1, 2, 3, 4 ]
\end{Verbatim}
 }

 

\subsection{\textcolor{Chapter }{CharacterTableWithStoredGroup}}
\logpage{[ 71, 6, 4 ]}\nobreak
\hyperdef{L}{X8788C6C7829C1ADE}{}
{\noindent\textcolor{FuncColor}{$\triangleright$\enspace\texttt{CharacterTableWithStoredGroup({\mdseries\slshape G, tbl[, info]})\index{CharacterTableWithStoredGroup@\texttt{CharacterTableWithStoredGroup}}
\label{CharacterTableWithStoredGroup}
}\hfill{\scriptsize (function)}}\\


 Let \mbox{\texttt{\mdseries\slshape G}} be a group and \mbox{\texttt{\mdseries\slshape tbl}} a character table of (a group isomorphic to) \mbox{\texttt{\mdseries\slshape G}}, such that \mbox{\texttt{\mdseries\slshape G}} does not store its \texttt{OrdinaryCharacterTable} (\ref{OrdinaryCharacterTable:for a group}) value. \texttt{CharacterTableWithStoredGroup} calls \texttt{CompatibleConjugacyClasses} (\ref{CompatibleConjugacyClasses}), trying to identify the classes of \mbox{\texttt{\mdseries\slshape G}} with the columns of \mbox{\texttt{\mdseries\slshape tbl}}. 

 If this identification is unique up to automorphisms of \mbox{\texttt{\mdseries\slshape tbl}} (see{\nobreakspace}\texttt{AutomorphismsOfTable} (\ref{AutomorphismsOfTable})) then \mbox{\texttt{\mdseries\slshape tbl}} is stored as \texttt{CharacterTable} (\ref{CharacterTable:for a group}) value of \mbox{\texttt{\mdseries\slshape G}}, and a new character table is returned that is equivalent to \mbox{\texttt{\mdseries\slshape tbl}}, is sorted in the same way as \mbox{\texttt{\mdseries\slshape tbl}}, and has the values of \texttt{UnderlyingGroup} (\ref{UnderlyingGroup:for character tables}), \texttt{ConjugacyClasses} (\ref{ConjugacyClasses:for character tables}), and \texttt{IdentificationOfConjugacyClasses} (\ref{IdentificationOfConjugacyClasses}) set. 

 Otherwise, i.e., if \textsf{GAP} cannot identify the classes of \mbox{\texttt{\mdseries\slshape G}} up to automorphisms of \mbox{\texttt{\mdseries\slshape tbl}}, \texttt{fail} is returned. 

 If a record is present as the third argument \mbox{\texttt{\mdseries\slshape info}}, its meaning is the same as the optional argument \mbox{\texttt{\mdseries\slshape arec}} for \texttt{CompatibleConjugacyClasses} (\ref{CompatibleConjugacyClasses}). 

 If a list is entered as third argument \mbox{\texttt{\mdseries\slshape info}} it is used as value of \texttt{IdentificationOfConjugacyClasses} (\ref{IdentificationOfConjugacyClasses}), relative to the \texttt{ConjugacyClasses} (\ref{ConjugacyClasses:for character tables}) value of \mbox{\texttt{\mdseries\slshape G}}, without further checking, and the corresponding character table is returned. }

 

\subsection{\textcolor{Chapter }{CompatibleConjugacyClasses}}
\logpage{[ 71, 6, 5 ]}\nobreak
\hyperdef{L}{X790019E87CFDDB98}{}
{\noindent\textcolor{FuncColor}{$\triangleright$\enspace\texttt{CompatibleConjugacyClasses({\mdseries\slshape [G, ccl, ]tbl[, arec]})\index{CompatibleConjugacyClasses@\texttt{CompatibleConjugacyClasses}}
\label{CompatibleConjugacyClasses}
}\hfill{\scriptsize (operation)}}\\


 If the arguments \mbox{\texttt{\mdseries\slshape G}} and \mbox{\texttt{\mdseries\slshape ccl}} are present then \mbox{\texttt{\mdseries\slshape ccl}} must be a list of the conjugacy classes of the group \mbox{\texttt{\mdseries\slshape G}}, and \mbox{\texttt{\mdseries\slshape tbl}} the ordinary character table of \mbox{\texttt{\mdseries\slshape G}}. Then \texttt{CompatibleConjugacyClasses} returns a list $l$ of positive integers that describes an identification of the columns of \mbox{\texttt{\mdseries\slshape tbl}} with the conjugacy classes \mbox{\texttt{\mdseries\slshape ccl}} in the sense that $l[i]$ is the position in \mbox{\texttt{\mdseries\slshape ccl}} of the class corresponding to the $i$\texttt{\symbol{45}}th column of \mbox{\texttt{\mdseries\slshape tbl}}, if this identification is unique up to automorphisms of \mbox{\texttt{\mdseries\slshape tbl}} (see{\nobreakspace}\texttt{AutomorphismsOfTable} (\ref{AutomorphismsOfTable})); if \textsf{GAP} cannot identify the classes, \texttt{fail} is returned. 

 If \mbox{\texttt{\mdseries\slshape tbl}} is the first argument then it must be an ordinary character table, and \texttt{CompatibleConjugacyClasses} checks whether the columns of \mbox{\texttt{\mdseries\slshape tbl}} can be identified with the conjugacy classes of a group isomorphic to that for
which \mbox{\texttt{\mdseries\slshape tbl}} is the character table; the return value is a list of all those sets of class
positions for which the columns of \mbox{\texttt{\mdseries\slshape tbl}} cannot be distinguished with the invariants used, up to automorphisms of \mbox{\texttt{\mdseries\slshape tbl}}. So the identification is unique if and only if the returned list is empty. 

 The usual approach is that one first calls \texttt{CompatibleConjugacyClasses} in the second form for checking quickly whether the first form will be
successful, and only if this is the case the more time consuming calculations
with both group and character table are done. 

 The following invariants are used. 
\begin{enumerate}
\item  element orders (see{\nobreakspace}\texttt{OrdersClassRepresentatives} (\ref{OrdersClassRepresentatives})), 
\item  class lengths (see{\nobreakspace}\texttt{SizesConjugacyClasses} (\ref{SizesConjugacyClasses})), 
\item  power maps (see{\nobreakspace}\texttt{PowerMap} (\ref{PowerMap}), \texttt{ComputedPowerMaps} (\ref{ComputedPowerMaps})), 
\item  symmetries of the table (see{\nobreakspace}\texttt{AutomorphismsOfTable} (\ref{AutomorphismsOfTable})). 
\end{enumerate}
 

 If the optional argument \mbox{\texttt{\mdseries\slshape arec}} is present then it must be a record whose components describe additional
information for the class identification. The following components are
supported. 
\begin{description}
\item[{\texttt{natchar} }]  if $G$ is a permutation group or matrix group then the value of this component is
regarded as the list of values of the natural character (see{\nobreakspace}\texttt{NaturalCharacter} (\ref{NaturalCharacter:for a group})) of \mbox{\texttt{\mdseries\slshape G}}, w.r.t.{\nobreakspace}the ordering of classes in \mbox{\texttt{\mdseries\slshape tbl}}, 
\item[{\texttt{bijection} }]  a list describing a partial bijection; the $i$\texttt{\symbol{45}}th entry, if bound, is the position of the $i$\texttt{\symbol{45}}th conjugacy class of \mbox{\texttt{\mdseries\slshape tbl}} in the list \mbox{\texttt{\mdseries\slshape ccl}}. 
\end{description}
 

 
\begin{Verbatim}[commandchars=!@|,fontsize=\small,frame=single,label=Example]
  !gapprompt@gap>| !gapinput@g:= AlternatingGroup( 5 );|
  Alt( [ 1 .. 5 ] )
  !gapprompt@gap>| !gapinput@tbl:= CharacterTable( "A5" );|
  CharacterTable( "A5" )
  !gapprompt@gap>| !gapinput@HasUnderlyingGroup( tbl );  HasOrdinaryCharacterTable( g );|
  false
  false
  !gapprompt@gap>| !gapinput@CompatibleConjugacyClasses( tbl );   # unique identification|
  [  ]
  !gapprompt@gap>| !gapinput@new:= CharacterTableWithStoredGroup( g, tbl );|
  CharacterTable( Alt( [ 1 .. 5 ] ) )
  !gapprompt@gap>| !gapinput@Irr( new ) = Irr( tbl );|
  true
  !gapprompt@gap>| !gapinput@HasConjugacyClasses( new );  HasUnderlyingGroup( new );|
  true
  true
  !gapprompt@gap>| !gapinput@IdentificationOfConjugacyClasses( new );|
  [ 1, 2, 3, 4, 5 ]
  !gapprompt@gap>| !gapinput@# Here is an example where the identification is not unique.|
  !gapprompt@gap>| !gapinput@CompatibleConjugacyClasses( CharacterTable( "J2" ) );|
  [ [ 17, 18 ], [ 9, 10 ] ]
\end{Verbatim}
 }

 }

  
\section{\textcolor{Chapter }{Operators for Character Tables}}\label{Operators for Character Tables}
\logpage{[ 71, 7, 0 ]}
\hyperdef{L}{X7CADCBC9824CB624}{}
{
  \index{*@\texttt{*}!for character tables} \index{/@\texttt{/}!for character tables} \index{mod@\texttt{mod}!for character tables} \index{character tables!infix operators} The following infix operators are defined for character tables. 
\begin{description}
\item[{\texttt{\mbox{\texttt{\mdseries\slshape tbl1}} * \mbox{\texttt{\mdseries\slshape tbl2}}}}]  the direct product of two character tables (see{\nobreakspace}\texttt{CharacterTableDirectProduct} (\ref{CharacterTableDirectProduct})), 
\item[{\texttt{\mbox{\texttt{\mdseries\slshape tbl}} / \mbox{\texttt{\mdseries\slshape list}}}}]  the table of the factor group modulo the normal subgroup spanned by the
classes in the list \mbox{\texttt{\mdseries\slshape list}} (see{\nobreakspace}\texttt{CharacterTableFactorGroup} (\ref{CharacterTableFactorGroup})), 
\item[{\texttt{\mbox{\texttt{\mdseries\slshape tbl}} mod \mbox{\texttt{\mdseries\slshape p}}}}]  the \mbox{\texttt{\mdseries\slshape p}}\texttt{\symbol{45}}modular Brauer character table corresponding to the
ordinary character table \mbox{\texttt{\mdseries\slshape tbl}} (see{\nobreakspace}\texttt{BrauerTable} (\ref{BrauerTable:for a character table, and a prime integer})), 
\item[{\texttt{\mbox{\texttt{\mdseries\slshape tbl}}.\mbox{\texttt{\mdseries\slshape name}}}}]  the position of the class with name \mbox{\texttt{\mdseries\slshape name}} in \mbox{\texttt{\mdseries\slshape tbl}} (see{\nobreakspace}\texttt{ClassNames} (\ref{ClassNames})). 
\end{description}
 }

  
\section{\textcolor{Chapter }{Attributes and Properties for Groups and Character Tables}}\label{Attributes and Properties of Character Tables}
\logpage{[ 71, 8, 0 ]}
\hyperdef{L}{X7F9D58208241D35E}{}
{
  Several \emph{attributes for groups} are valid also for character tables. 

 These are first those that have the same meaning for both the group and its
character table, and whose values can be read off or computed, respectively,
from the character table, such as \texttt{Size} (\ref{Size:for a character table}), \texttt{IsAbelian} (\ref{IsAbelian:for a character table}), or \texttt{IsSolvable} (\ref{IsSolvable:for a character table}). 

 Second, there are attributes whose meaning for character tables is different
from the meaning for groups, such as \texttt{ConjugacyClasses} (\ref{ConjugacyClasses:for character tables}). 
\subsection{\textcolor{Chapter }{CharacterDegrees}}\logpage{[ 71, 8, 1 ]}
\hyperdef{L}{X81FEFF768134481A}{}
{
\noindent\textcolor{FuncColor}{$\triangleright$\enspace\texttt{CharacterDegrees({\mdseries\slshape G[, p]})\index{CharacterDegrees@\texttt{CharacterDegrees}!for a group}
\label{CharacterDegrees:for a group}
}\hfill{\scriptsize (attribute)}}\\
\noindent\textcolor{FuncColor}{$\triangleright$\enspace\texttt{CharacterDegrees({\mdseries\slshape tbl})\index{CharacterDegrees@\texttt{CharacterDegrees}!for a character table}
\label{CharacterDegrees:for a character table}
}\hfill{\scriptsize (attribute)}}\\


 In the first form, \texttt{CharacterDegrees} returns a collected list of the degrees of the absolutely irreducible
characters of the group \mbox{\texttt{\mdseries\slshape G}}; the optional second argument \mbox{\texttt{\mdseries\slshape p}} must be either zero or a prime integer denoting the characteristic, the
default value is zero. In the second form, \mbox{\texttt{\mdseries\slshape tbl}} must be an (ordinary or Brauer) character table, and \texttt{CharacterDegrees} returns a collected list of the degrees of the absolutely irreducible
characters of \mbox{\texttt{\mdseries\slshape tbl}}. 

 (The default method for the call with only argument a group is to call the
operation with second argument \texttt{0}.) 

 For solvable groups, the default method is based on{\nobreakspace}\cite{Con90b}. 

 
\begin{Verbatim}[commandchars=!@|,fontsize=\small,frame=single,label=Example]
  !gapprompt@gap>| !gapinput@CharacterDegrees( SymmetricGroup( 4 ) );|
  [ [ 1, 2 ], [ 2, 1 ], [ 3, 2 ] ]
  !gapprompt@gap>| !gapinput@CharacterDegrees( SymmetricGroup( 4 ), 2 );|
  [ [ 1, 1 ], [ 2, 1 ] ]
  !gapprompt@gap>| !gapinput@CharacterDegrees( CharacterTable( "A5" ) );|
  [ [ 1, 1 ], [ 3, 2 ], [ 4, 1 ], [ 5, 1 ] ]
  !gapprompt@gap>| !gapinput@CharacterDegrees( CharacterTable( "A5" ) mod 2 );|
  [ [ 1, 1 ], [ 2, 2 ], [ 4, 1 ] ]
\end{Verbatim}
 }

 
\subsection{\textcolor{Chapter }{Irr}}\logpage{[ 71, 8, 2 ]}
\hyperdef{L}{X873B3CC57E9A5492}{}
{
\noindent\textcolor{FuncColor}{$\triangleright$\enspace\texttt{Irr({\mdseries\slshape G[, p]})\index{Irr@\texttt{Irr}!for a group}
\label{Irr:for a group}
}\hfill{\scriptsize (attribute)}}\\
\noindent\textcolor{FuncColor}{$\triangleright$\enspace\texttt{Irr({\mdseries\slshape tbl})\index{Irr@\texttt{Irr}!for a character table}
\label{Irr:for a character table}
}\hfill{\scriptsize (attribute)}}\\


 Called with a group \mbox{\texttt{\mdseries\slshape G}}, \texttt{Irr} returns the irreducible characters of the ordinary character table of \mbox{\texttt{\mdseries\slshape G}}. Called with a group \mbox{\texttt{\mdseries\slshape G}} and a prime integer \mbox{\texttt{\mdseries\slshape p}}, \texttt{Irr} returns the irreducible characters of the \mbox{\texttt{\mdseries\slshape p}}\texttt{\symbol{45}}modular Brauer table of \mbox{\texttt{\mdseries\slshape G}}. Called with an (ordinary or Brauer) character table \mbox{\texttt{\mdseries\slshape tbl}}, \texttt{Irr} returns the list of all complex absolutely irreducible characters of \mbox{\texttt{\mdseries\slshape tbl}}. 

 For a character table \mbox{\texttt{\mdseries\slshape tbl}} with underlying group, \texttt{Irr} may delegate to the group. For a group \mbox{\texttt{\mdseries\slshape G}}, \texttt{Irr} may delegate to its character table only if the irreducibles are already
stored there. 

 (If \mbox{\texttt{\mdseries\slshape G}} is \mbox{\texttt{\mdseries\slshape p}}\texttt{\symbol{45}}solvable (see{\nobreakspace}\texttt{IsPSolvable} (\ref{IsPSolvable})) then the \mbox{\texttt{\mdseries\slshape p}}\texttt{\symbol{45}}modular irreducible characters can be computed by the
Fong\texttt{\symbol{45}}Swan Theorem; in all other cases, there may be no
method.) 

 Note that the ordering of columns in the \texttt{Irr} matrix of the group \mbox{\texttt{\mdseries\slshape G}} refers to the ordering of conjugacy classes in the \texttt{CharacterTable} (\ref{CharacterTable:for a group}) value of \mbox{\texttt{\mdseries\slshape G}}, which may differ from the ordering of conjugacy classes in \mbox{\texttt{\mdseries\slshape G}} (see \ref{The Interface between Character Tables and Groups}). As an extreme example, for a character table obtained from sorting the
classes of the \texttt{CharacterTable} (\ref{CharacterTable:for a group}) value of \mbox{\texttt{\mdseries\slshape G}}, the ordering of columns in the \texttt{Irr} matrix respects the sorting of classes (see{\nobreakspace}\ref{Sorted Character Tables}), so the irreducibles of such a table will in general not coincide with the
irreducibles stored as the \texttt{Irr} value of \mbox{\texttt{\mdseries\slshape G}} although also the sorted table stores the group \mbox{\texttt{\mdseries\slshape G}}. 

 The ordering of the entries in the attribute \texttt{Irr} of a group need \emph{not} coincide with the ordering of its \texttt{IrreducibleRepresentations} (\ref{IrreducibleRepresentations}) value. 

 
\begin{Verbatim}[commandchars=!@|,fontsize=\small,frame=single,label=Example]
  !gapprompt@gap>| !gapinput@Irr( SymmetricGroup( 4 ) );|
  [ Character( CharacterTable( Sym( [ 1 .. 4 ] ) ), [ 1, -1, 1, 1, -1
       ] ), Character( CharacterTable( Sym( [ 1 .. 4 ] ) ),
      [ 3, -1, -1, 0, 1 ] ),
    Character( CharacterTable( Sym( [ 1 .. 4 ] ) ), [ 2, 0, 2, -1, 0 ] )
      , Character( CharacterTable( Sym( [ 1 .. 4 ] ) ),
      [ 3, 1, -1, 0, -1 ] ),
    Character( CharacterTable( Sym( [ 1 .. 4 ] ) ), [ 1, 1, 1, 1, 1 ] )
   ]
  !gapprompt@gap>| !gapinput@Irr( SymmetricGroup( 4 ), 2 );|
  [ Character( BrauerTable( Sym( [ 1 .. 4 ] ), 2 ), [ 1, 1 ] ),
    Character( BrauerTable( Sym( [ 1 .. 4 ] ), 2 ), [ 2, -1 ] ) ]
  !gapprompt@gap>| !gapinput@Irr( CharacterTable( "A5" ) );|
  [ Character( CharacterTable( "A5" ), [ 1, 1, 1, 1, 1 ] ),
    Character( CharacterTable( "A5" ),
      [ 3, -1, 0, -E(5)-E(5)^4, -E(5)^2-E(5)^3 ] ),
    Character( CharacterTable( "A5" ),
      [ 3, -1, 0, -E(5)^2-E(5)^3, -E(5)-E(5)^4 ] ),
    Character( CharacterTable( "A5" ), [ 4, 0, 1, -1, -1 ] ),
    Character( CharacterTable( "A5" ), [ 5, 1, -1, 0, 0 ] ) ]
  !gapprompt@gap>| !gapinput@Irr( CharacterTable( "A5" ) mod 2 );|
  [ Character( BrauerTable( "A5", 2 ), [ 1, 1, 1, 1 ] ),
    Character( BrauerTable( "A5", 2 ),
      [ 2, -1, E(5)+E(5)^4, E(5)^2+E(5)^3 ] ),
    Character( BrauerTable( "A5", 2 ),
      [ 2, -1, E(5)^2+E(5)^3, E(5)+E(5)^4 ] ),
    Character( BrauerTable( "A5", 2 ), [ 4, 1, -1, -1 ] ) ]
\end{Verbatim}
 }

 
\subsection{\textcolor{Chapter }{LinearCharacters}}\logpage{[ 71, 8, 3 ]}
\hyperdef{L}{X8549899A7DE206BA}{}
{
\noindent\textcolor{FuncColor}{$\triangleright$\enspace\texttt{LinearCharacters({\mdseries\slshape G[, p]})\index{LinearCharacters@\texttt{LinearCharacters}!for a group}
\label{LinearCharacters:for a group}
}\hfill{\scriptsize (attribute)}}\\
\noindent\textcolor{FuncColor}{$\triangleright$\enspace\texttt{LinearCharacters({\mdseries\slshape tbl})\index{LinearCharacters@\texttt{LinearCharacters}!for a character table}
\label{LinearCharacters:for a character table}
}\hfill{\scriptsize (attribute)}}\\


 \texttt{LinearCharacters} returns the linear (i.e., degree $1$) characters in the \texttt{Irr} (\ref{Irr:for a group}) list of the group \mbox{\texttt{\mdseries\slshape G}} or the character table \mbox{\texttt{\mdseries\slshape tbl}}, respectively. In the second form, \texttt{LinearCharacters} returns the \mbox{\texttt{\mdseries\slshape p}}\texttt{\symbol{45}}modular linear characters of the group \mbox{\texttt{\mdseries\slshape G}}. 

 For a character table \mbox{\texttt{\mdseries\slshape tbl}} with underlying group, \texttt{LinearCharacters} may delegate to the group. For a group \mbox{\texttt{\mdseries\slshape G}}, \texttt{LinearCharacters} may delegate to its character table only if the irreducibles are already
stored there. 

 The ordering of linear characters in \mbox{\texttt{\mdseries\slshape tbl}} need not coincide with the ordering of linear characters in the irreducibles
of \mbox{\texttt{\mdseries\slshape tbl}} (see{\nobreakspace}\texttt{Irr} (\ref{Irr:for a character table})). 

 
\begin{Verbatim}[commandchars=!@|,fontsize=\small,frame=single,label=Example]
  !gapprompt@gap>| !gapinput@LinearCharacters( SymmetricGroup( 4 ) );|
  [ Character( CharacterTable( Sym( [ 1 .. 4 ] ) ), [ 1, 1, 1, 1, 1 ] ),
    Character( CharacterTable( Sym( [ 1 .. 4 ] ) ), [ 1, -1, 1, 1, -1
       ] ) ]
\end{Verbatim}
 }

 
\subsection{\textcolor{Chapter }{OrdinaryCharacterTable}}\logpage{[ 71, 8, 4 ]}
\hyperdef{L}{X8011EEB684848039}{}
{
\noindent\textcolor{FuncColor}{$\triangleright$\enspace\texttt{OrdinaryCharacterTable({\mdseries\slshape G})\index{OrdinaryCharacterTable@\texttt{OrdinaryCharacterTable}!for a group}
\label{OrdinaryCharacterTable:for a group}
}\hfill{\scriptsize (attribute)}}\\
\noindent\textcolor{FuncColor}{$\triangleright$\enspace\texttt{OrdinaryCharacterTable({\mdseries\slshape modtbl})\index{OrdinaryCharacterTable@\texttt{OrdinaryCharacterTable}!for a character table}
\label{OrdinaryCharacterTable:for a character table}
}\hfill{\scriptsize (attribute)}}\\


 \texttt{OrdinaryCharacterTable} returns the ordinary character table of the group \mbox{\texttt{\mdseries\slshape G}} or the Brauer character table \mbox{\texttt{\mdseries\slshape modtbl}}, respectively. 

 Since Brauer character tables are constructed from ordinary tables, the
attribute value for \mbox{\texttt{\mdseries\slshape modtbl}} is already stored (cf.{\nobreakspace}\ref{Character Table Categories}). 

 
\begin{Verbatim}[commandchars=!@|,fontsize=\small,frame=single,label=Example]
  !gapprompt@gap>| !gapinput@OrdinaryCharacterTable( SymmetricGroup( 4 ) );|
  CharacterTable( Sym( [ 1 .. 4 ] ) )
  !gapprompt@gap>| !gapinput@tbl:= CharacterTable( "A5" );;  modtbl:= tbl mod 2;|
  BrauerTable( "A5", 2 )
  !gapprompt@gap>| !gapinput@OrdinaryCharacterTable( modtbl ) = tbl;|
  true
\end{Verbatim}
 }

 
\subsection{\textcolor{Chapter }{Group Operations Applicable to Character Tables}}\logpage{[ 71, 8, 5 ]}
\hyperdef{L}{X81EFD9FE804AC6EE}{}
{
\noindent\textcolor{FuncColor}{$\triangleright$\enspace\texttt{AbelianInvariants({\mdseries\slshape tbl})\index{AbelianInvariants@\texttt{AbelianInvariants}!for a character table}
\label{AbelianInvariants:for a character table}
}\hfill{\scriptsize (attribute)}}\\
\noindent\textcolor{FuncColor}{$\triangleright$\enspace\texttt{CommutatorLength({\mdseries\slshape tbl})\index{CommutatorLength@\texttt{CommutatorLength}!for a character table}
\label{CommutatorLength:for a character table}
}\hfill{\scriptsize (attribute)}}\\
\noindent\textcolor{FuncColor}{$\triangleright$\enspace\texttt{Exponent({\mdseries\slshape tbl})\index{Exponent@\texttt{Exponent}!for a character table}
\label{Exponent:for a character table}
}\hfill{\scriptsize (attribute)}}\\
\noindent\textcolor{FuncColor}{$\triangleright$\enspace\texttt{IsAbelian({\mdseries\slshape tbl})\index{IsAbelian@\texttt{IsAbelian}!for a character table}
\label{IsAbelian:for a character table}
}\hfill{\scriptsize (property)}}\\
\noindent\textcolor{FuncColor}{$\triangleright$\enspace\texttt{IsAlmostSimple({\mdseries\slshape tbl})\index{IsAlmostSimple@\texttt{IsAlmostSimple}!for a character table}
\label{IsAlmostSimple:for a character table}
}\hfill{\scriptsize (property)}}\\
\noindent\textcolor{FuncColor}{$\triangleright$\enspace\texttt{IsCyclic({\mdseries\slshape tbl})\index{IsCyclic@\texttt{IsCyclic}!for a character table}
\label{IsCyclic:for a character table}
}\hfill{\scriptsize (property)}}\\
\noindent\textcolor{FuncColor}{$\triangleright$\enspace\texttt{IsElementaryAbelian({\mdseries\slshape tbl})\index{IsElementaryAbelian@\texttt{IsElementaryAbelian}!for a character table}
\label{IsElementaryAbelian:for a character table}
}\hfill{\scriptsize (property)}}\\
\noindent\textcolor{FuncColor}{$\triangleright$\enspace\texttt{IsFinite({\mdseries\slshape tbl})\index{IsFinite@\texttt{IsFinite}!for a character table}
\label{IsFinite:for a character table}
}\hfill{\scriptsize (property)}}\\
\noindent\textcolor{FuncColor}{$\triangleright$\enspace\texttt{IsMonomial({\mdseries\slshape tbl})\index{IsMonomial@\texttt{IsMonomial}!for a character table}
\label{IsMonomial:for a character table}
}\hfill{\scriptsize (property)}}\\
\noindent\textcolor{FuncColor}{$\triangleright$\enspace\texttt{IsNilpotent({\mdseries\slshape tbl})\index{IsNilpotent@\texttt{IsNilpotent}!for a character table}
\label{IsNilpotent:for a character table}
}\hfill{\scriptsize (property)}}\\
\noindent\textcolor{FuncColor}{$\triangleright$\enspace\texttt{IsPerfect({\mdseries\slshape tbl})\index{IsPerfect@\texttt{IsPerfect}!for a character table}
\label{IsPerfect:for a character table}
}\hfill{\scriptsize (property)}}\\
\noindent\textcolor{FuncColor}{$\triangleright$\enspace\texttt{IsQuasisimple({\mdseries\slshape tbl})\index{IsQuasisimple@\texttt{IsQuasisimple}!for a character table}
\label{IsQuasisimple:for a character table}
}\hfill{\scriptsize (property)}}\\
\noindent\textcolor{FuncColor}{$\triangleright$\enspace\texttt{IsSimple({\mdseries\slshape tbl})\index{IsSimple@\texttt{IsSimple}!for a character table}
\label{IsSimple:for a character table}
}\hfill{\scriptsize (property)}}\\
\noindent\textcolor{FuncColor}{$\triangleright$\enspace\texttt{IsSolvable({\mdseries\slshape tbl})\index{IsSolvable@\texttt{IsSolvable}!for a character table}
\label{IsSolvable:for a character table}
}\hfill{\scriptsize (property)}}\\
\noindent\textcolor{FuncColor}{$\triangleright$\enspace\texttt{IsSporadicSimple({\mdseries\slshape tbl})\index{IsSporadicSimple@\texttt{IsSporadicSimple}!for a character table}
\label{IsSporadicSimple:for a character table}
}\hfill{\scriptsize (property)}}\\
\noindent\textcolor{FuncColor}{$\triangleright$\enspace\texttt{IsSupersolvable({\mdseries\slshape tbl})\index{IsSupersolvable@\texttt{IsSupersolvable}!for a character table}
\label{IsSupersolvable:for a character table}
}\hfill{\scriptsize (property)}}\\
\noindent\textcolor{FuncColor}{$\triangleright$\enspace\texttt{IsomorphismTypeInfoFiniteSimpleGroup({\mdseries\slshape tbl})\index{IsomorphismTypeInfoFiniteSimpleGroup@\texttt{Isomorphism}\-\texttt{Type}\-\texttt{Info}\-\texttt{Finite}\-\texttt{Simple}\-\texttt{Group}!for a character table}
\label{IsomorphismTypeInfoFiniteSimpleGroup:for a character table}
}\hfill{\scriptsize (attribute)}}\\
\noindent\textcolor{FuncColor}{$\triangleright$\enspace\texttt{NrConjugacyClasses({\mdseries\slshape tbl})\index{NrConjugacyClasses@\texttt{NrConjugacyClasses}!for a character table}
\label{NrConjugacyClasses:for a character table}
}\hfill{\scriptsize (attribute)}}\\
\noindent\textcolor{FuncColor}{$\triangleright$\enspace\texttt{Size({\mdseries\slshape tbl})\index{Size@\texttt{Size}!for a character table}
\label{Size:for a character table}
}\hfill{\scriptsize (attribute)}}\\


 These operations for groups are applicable to character tables and mean the
same for a character table as for its underlying group; see Chapter \ref{Groups} for the definitions. The operations are mainly useful for selecting character
tables with certain properties, also for character tables without access to a
group. 

 
\begin{Verbatim}[commandchars=!@|,fontsize=\small,frame=single,label=Example]
  !gapprompt@gap>| !gapinput@tables:= [ CharacterTable( CyclicGroup( 3 ) ),|
  !gapprompt@>| !gapinput@              CharacterTable( SymmetricGroup( 4 ) ),|
  !gapprompt@>| !gapinput@              CharacterTable( AlternatingGroup( 5 ) ),|
  !gapprompt@>| !gapinput@              CharacterTable( SL( 2, 5 ) ) ];;|
  !gapprompt@gap>| !gapinput@List( tables, AbelianInvariants );|
  [ [ 3 ], [ 2 ], [  ], [  ] ]
  !gapprompt@gap>| !gapinput@List( tables, CommutatorLength );|
  [ 1, 1, 1, 1 ]
  !gapprompt@gap>| !gapinput@List( tables, Exponent );|
  [ 3, 12, 30, 60 ]
  !gapprompt@gap>| !gapinput@List( tables, IsAbelian );|
  [ true, false, false, false ]
  !gapprompt@gap>| !gapinput@List( tables, IsAlmostSimple );|
  [ false, false, true, false ]
  !gapprompt@gap>| !gapinput@List( tables, IsCyclic );|
  [ true, false, false, false ]
  !gapprompt@gap>| !gapinput@List( tables, IsFinite );|
  [ true, true, true, true ]
  !gapprompt@gap>| !gapinput@List( tables, IsMonomial );|
  [ true, true, false, false ]
  !gapprompt@gap>| !gapinput@List( tables, IsNilpotent );|
  [ true, false, false, false ]
  !gapprompt@gap>| !gapinput@List( tables, IsPerfect );|
  [ false, false, true, true ]
  !gapprompt@gap>| !gapinput@List( tables, IsQuasisimple );|
  [ false, false, true, true ]
  !gapprompt@gap>| !gapinput@List( tables, IsSimple );|
  [ true, false, true, false ]
  !gapprompt@gap>| !gapinput@List( tables, IsSolvable );|
  [ true, true, false, false ]
  !gapprompt@gap>| !gapinput@List( tables, IsSupersolvable );|
  [ true, false, false, false ]
  !gapprompt@gap>| !gapinput@List( tables, NrConjugacyClasses );|
  [ 3, 5, 5, 9 ]
  !gapprompt@gap>| !gapinput@List( tables, Size );|
  [ 3, 24, 60, 120 ]
  !gapprompt@gap>| !gapinput@IsomorphismTypeInfoFiniteSimpleGroup( CharacterTable( "C5" ) );|
  rec( name := "Z(5)", parameter := 5, series := "Z", shortname := "C5"
   )
  !gapprompt@gap>| !gapinput@IsomorphismTypeInfoFiniteSimpleGroup( CharacterTable( "S3" ) );|
  fail
  !gapprompt@gap>| !gapinput@IsomorphismTypeInfoFiniteSimpleGroup( CharacterTable( "S6(3)" ) );|
  rec( name := "C(3,3) = S(6,3)", parameter := [ 3, 3 ], series := "C",
    shortname := "S6(3)" )
  !gapprompt@gap>| !gapinput@IsomorphismTypeInfoFiniteSimpleGroup( CharacterTable( "O7(3)" ) );|
  rec( name := "B(3,3) = O(7,3)", parameter := [ 3, 3 ], series := "B",
    shortname := "O7(3)" )
  !gapprompt@gap>| !gapinput@IsomorphismTypeInfoFiniteSimpleGroup( CharacterTable( "A8" ) );|
  rec( name := "A(8) ~ A(3,2) = L(4,2) ~ D(3,2) = O+(6,2)",
    parameter := 8, series := "A", shortname := "A8" )
  !gapprompt@gap>| !gapinput@IsomorphismTypeInfoFiniteSimpleGroup( CharacterTable( "L3(4)" ) );|
  rec( name := "A(2,4) = L(3,4)", parameter := [ 3, 4 ], series := "L",
    shortname := "L3(4)" )
\end{Verbatim}
 }

 }

  
\section{\textcolor{Chapter }{Attributes and Properties only for Character Tables}}\label{Attributes and Properties only for Character Tables}
\logpage{[ 71, 9, 0 ]}
\hyperdef{L}{X7995A2AD83BC58A0}{}
{
  The following three \emph{attributes for character tables} {\textendash}\texttt{OrdersClassRepresentatives} (\ref{OrdersClassRepresentatives}), \texttt{SizesCentralizers} (\ref{SizesCentralizers}), and \texttt{SizesConjugacyClasses} (\ref{SizesConjugacyClasses}){\textendash} would make sense also for groups but are in fact \emph{not} used for groups. This is because the values depend on the ordering of
conjugacy classes stored as the value of \texttt{ConjugacyClasses} (\ref{ConjugacyClasses:for character tables}), and this value may differ for a group and its character table (see \ref{The Interface between Character Tables and Groups}). Note that for character tables, the consistency of attribute values must be
guaranteed, whereas for groups, there is no need to impose such a consistency
rule. 

 The other attributes introduced in this section apply only to character
tables, not to groups. 

\subsection{\textcolor{Chapter }{OrdersClassRepresentatives}}
\logpage{[ 71, 9, 1 ]}\nobreak
\hyperdef{L}{X86F455DA7A9C30EE}{}
{\noindent\textcolor{FuncColor}{$\triangleright$\enspace\texttt{OrdersClassRepresentatives({\mdseries\slshape tbl})\index{OrdersClassRepresentatives@\texttt{OrdersClassRepresentatives}}
\label{OrdersClassRepresentatives}
}\hfill{\scriptsize (attribute)}}\\


 is a list of orders of representatives of conjugacy classes of the character
table \mbox{\texttt{\mdseries\slshape tbl}}, in the same ordering as the conjugacy classes of \mbox{\texttt{\mdseries\slshape tbl}}. 

 
\begin{Verbatim}[commandchars=!@|,fontsize=\small,frame=single,label=Example]
  !gapprompt@gap>| !gapinput@tbl:= CharacterTable( "A5" );;|
  !gapprompt@gap>| !gapinput@OrdersClassRepresentatives( tbl );|
  [ 1, 2, 3, 5, 5 ]
\end{Verbatim}
 }

 

\subsection{\textcolor{Chapter }{SizesCentralizers}}
\logpage{[ 71, 9, 2 ]}\nobreak
\hyperdef{L}{X7CF7907F790A5DE6}{}
{\noindent\textcolor{FuncColor}{$\triangleright$\enspace\texttt{SizesCentralizers({\mdseries\slshape tbl})\index{SizesCentralizers@\texttt{SizesCentralizers}}
\label{SizesCentralizers}
}\hfill{\scriptsize (attribute)}}\\
\noindent\textcolor{FuncColor}{$\triangleright$\enspace\texttt{SizesCentralisers({\mdseries\slshape tbl})\index{SizesCentralisers@\texttt{SizesCentralisers}}
\label{SizesCentralisers}
}\hfill{\scriptsize (attribute)}}\\


 is a list that stores at position $i$ the size of the centralizer of any element in the $i$\texttt{\symbol{45}}th conjugacy class of the character table \mbox{\texttt{\mdseries\slshape tbl}}. 

 
\begin{Verbatim}[commandchars=!@|,fontsize=\small,frame=single,label=Example]
  !gapprompt@gap>| !gapinput@tbl:= CharacterTable( "A5" );;|
  !gapprompt@gap>| !gapinput@SizesCentralizers( tbl );|
  [ 60, 4, 3, 5, 5 ]
\end{Verbatim}
 }

 

\subsection{\textcolor{Chapter }{SizesConjugacyClasses}}
\logpage{[ 71, 9, 3 ]}\nobreak
\hyperdef{L}{X7D9D2A45879A6A62}{}
{\noindent\textcolor{FuncColor}{$\triangleright$\enspace\texttt{SizesConjugacyClasses({\mdseries\slshape tbl})\index{SizesConjugacyClasses@\texttt{SizesConjugacyClasses}}
\label{SizesConjugacyClasses}
}\hfill{\scriptsize (attribute)}}\\


 is a list that stores at position $i$ the size of the $i$\texttt{\symbol{45}}th conjugacy class of the character table \mbox{\texttt{\mdseries\slshape tbl}}. 

 
\begin{Verbatim}[commandchars=!@|,fontsize=\small,frame=single,label=Example]
  !gapprompt@gap>| !gapinput@tbl:= CharacterTable( "A5" );;|
  !gapprompt@gap>| !gapinput@SizesConjugacyClasses( tbl );|
  [ 1, 15, 20, 12, 12 ]
\end{Verbatim}
 }

 

\subsection{\textcolor{Chapter }{AutomorphismsOfTable}}
\logpage{[ 71, 9, 4 ]}\nobreak
\hyperdef{L}{X7C2753DE8094F4BA}{}
{\noindent\textcolor{FuncColor}{$\triangleright$\enspace\texttt{AutomorphismsOfTable({\mdseries\slshape tbl})\index{AutomorphismsOfTable@\texttt{AutomorphismsOfTable}}
\label{AutomorphismsOfTable}
}\hfill{\scriptsize (attribute)}}\\


 is the permutation group of all column permutations of the character table \mbox{\texttt{\mdseries\slshape tbl}} that leave the set of irreducibles and each power map of \mbox{\texttt{\mdseries\slshape tbl}} invariant (see also{\nobreakspace}\texttt{TableAutomorphisms} (\ref{TableAutomorphisms})). 
\begin{Verbatim}[commandchars=!@|,fontsize=\small,frame=single,label=Example]
  !gapprompt@gap>| !gapinput@tbl:= CharacterTable( "Dihedral", 8 );;|
  !gapprompt@gap>| !gapinput@AutomorphismsOfTable( tbl );|
  Group([ (4,5) ])
  !gapprompt@gap>| !gapinput@OrdersClassRepresentatives( tbl );|
  [ 1, 4, 2, 2, 2 ]
  !gapprompt@gap>| !gapinput@SizesConjugacyClasses( tbl );|
  [ 1, 2, 1, 2, 2 ]
\end{Verbatim}
 }

 
\subsection{\textcolor{Chapter }{UnderlyingCharacteristic}}\logpage{[ 71, 9, 5 ]}
\hyperdef{L}{X7F58A82F7D88000A}{}
{
\noindent\textcolor{FuncColor}{$\triangleright$\enspace\texttt{UnderlyingCharacteristic({\mdseries\slshape tbl})\index{UnderlyingCharacteristic@\texttt{UnderlyingCharacteristic}!for a character table}
\label{UnderlyingCharacteristic:for a character table}
}\hfill{\scriptsize (attribute)}}\\
\noindent\textcolor{FuncColor}{$\triangleright$\enspace\texttt{UnderlyingCharacteristic({\mdseries\slshape psi})\index{UnderlyingCharacteristic@\texttt{UnderlyingCharacteristic}!for a character}
\label{UnderlyingCharacteristic:for a character}
}\hfill{\scriptsize (attribute)}}\\


 For an ordinary character table \mbox{\texttt{\mdseries\slshape tbl}}, the result is \texttt{0}, for a $p$\texttt{\symbol{45}}modular Brauer table \mbox{\texttt{\mdseries\slshape tbl}}, it is $p$. The underlying characteristic of a class function \mbox{\texttt{\mdseries\slshape psi}} is equal to that of its underlying character table. 

 The underlying characteristic must be stored when the table is constructed,
there is no method to compute it. 

 We cannot use the attribute \texttt{Characteristic} (\ref{Characteristic}) to denote this, since of course each Brauer character is an element of
characteristic zero in the sense of \textsf{GAP} (see Chapter{\nobreakspace}\ref{Class Functions}). 

 
\begin{Verbatim}[commandchars=!@|,fontsize=\small,frame=single,label=Example]
  !gapprompt@gap>| !gapinput@tbl:= CharacterTable( "A5" );;|
  !gapprompt@gap>| !gapinput@UnderlyingCharacteristic( tbl );|
  0
  !gapprompt@gap>| !gapinput@UnderlyingCharacteristic( tbl mod 17 );|
  17
\end{Verbatim}
 }

 
\subsection{\textcolor{Chapter }{Class Names and Character Names}}\logpage{[ 71, 9, 6 ]}
\hyperdef{L}{X804CFD597C795801}{}
{
\noindent\textcolor{FuncColor}{$\triangleright$\enspace\texttt{ClassNames({\mdseries\slshape tbl[, "ATLAS"]})\index{ClassNames@\texttt{ClassNames}}
\label{ClassNames}
}\hfill{\scriptsize (attribute)}}\\
\noindent\textcolor{FuncColor}{$\triangleright$\enspace\texttt{CharacterNames({\mdseries\slshape tbl})\index{CharacterNames@\texttt{CharacterNames}}
\label{CharacterNames}
}\hfill{\scriptsize (attribute)}}\\


 \texttt{ClassNames} and \texttt{CharacterNames} return lists of strings, one for each conjugacy class or irreducible
character, respectively, of the character table \mbox{\texttt{\mdseries\slshape tbl}}. These names are used when \mbox{\texttt{\mdseries\slshape tbl}} is displayed. 

 The default method for \texttt{ClassNames} computes class names consisting of the order of an element in the class and at
least one distinguishing letter. 

 The default method for \texttt{CharacterNames} returns the list \texttt{[ "X.1", "X.2", ... ]}, whose length is the number of irreducible characters of \mbox{\texttt{\mdseries\slshape tbl}}. 

 The position of the class with name \mbox{\texttt{\mdseries\slshape name}} in \mbox{\texttt{\mdseries\slshape tbl}} can be accessed as \texttt{\mbox{\texttt{\mdseries\slshape tbl}}.\mbox{\texttt{\mdseries\slshape name}}}. 

 When \texttt{ClassNames} is called with two arguments, the second being the string \texttt{"ATLAS"}, the class names returned obey the convention used in the \textsf{Atlas} of Finite Groups \cite[Chapter 7, Section 5]{CCN85}. If one is interested in ``relative'' class names of almost simple \textsf{Atlas} groups, one can use the function \texttt{AtlasClassNames} (\textbf{AtlasRep: AtlasClassNames}). 

 
\begin{Verbatim}[commandchars=!@|,fontsize=\small,frame=single,label=Example]
  !gapprompt@gap>| !gapinput@tbl:= CharacterTable( "A5" );;|
  !gapprompt@gap>| !gapinput@ClassNames( tbl );|
  [ "1a", "2a", "3a", "5a", "5b" ]
  !gapprompt@gap>| !gapinput@tbl.2a;|
  2
\end{Verbatim}
 }

 
\subsection{\textcolor{Chapter }{Class Parameters and Character Parameters}}\logpage{[ 71, 9, 7 ]}
\hyperdef{L}{X8333E8038308947E}{}
{
\noindent\textcolor{FuncColor}{$\triangleright$\enspace\texttt{ClassParameters({\mdseries\slshape tbl})\index{ClassParameters@\texttt{ClassParameters}}
\label{ClassParameters}
}\hfill{\scriptsize (attribute)}}\\
\noindent\textcolor{FuncColor}{$\triangleright$\enspace\texttt{CharacterParameters({\mdseries\slshape tbl})\index{CharacterParameters@\texttt{CharacterParameters}}
\label{CharacterParameters}
}\hfill{\scriptsize (attribute)}}\\


 The values of these attributes are lists containing a parameter for each
conjugacy class or irreducible character, respectively, of the character table \mbox{\texttt{\mdseries\slshape tbl}}. 

 It depends on \mbox{\texttt{\mdseries\slshape tbl}} what these parameters are, so there is no default to compute class and
character parameters. 

 For example, the classes of symmetric groups can be parametrized by
partitions, corresponding to the cycle structures of permutations. Character
tables constructed from generic character tables (see the manual of the \textsf{GAP} Character Table Library) usually have class and character parameters stored. 

 If \mbox{\texttt{\mdseries\slshape tbl}} is a $p$\texttt{\symbol{45}}modular Brauer table such that class parameters are stored
in the underlying ordinary table (see{\nobreakspace}\texttt{OrdinaryCharacterTable} (\ref{OrdinaryCharacterTable:for a character table})) of \mbox{\texttt{\mdseries\slshape tbl}} then \texttt{ClassParameters} returns the sublist of class parameters of the ordinary table, for $p$\texttt{\symbol{45}}regular classes.  }

 

\subsection{\textcolor{Chapter }{Identifier (for character tables)}}
\logpage{[ 71, 9, 8 ]}\nobreak
\hyperdef{L}{X79C40EE97890202F}{}
{\noindent\textcolor{FuncColor}{$\triangleright$\enspace\texttt{Identifier({\mdseries\slshape tbl})\index{Identifier@\texttt{Identifier}!for character tables}
\label{Identifier:for character tables}
}\hfill{\scriptsize (attribute)}}\\


 is a string that identifies the character table \mbox{\texttt{\mdseries\slshape tbl}} in the current \textsf{GAP} session. It is used mainly for class fusions into \mbox{\texttt{\mdseries\slshape tbl}} that are stored on other character tables. For character tables without group,
the identifier is also used to print the table; this is the case for library
tables, but also for tables that are constructed as direct products, factors
etc.{\nobreakspace}involving tables that may or may not store their groups. 

 The default method for ordinary tables constructs strings of the form \texttt{"CT\mbox{\texttt{\mdseries\slshape n}}"}, where \mbox{\texttt{\mdseries\slshape n}} is a positive integer. \texttt{LARGEST{\textunderscore}IDENTIFIER{\textunderscore}NUMBER} is a list containing the largest integer \mbox{\texttt{\mdseries\slshape n}} used in the current \textsf{GAP} session. 

 The default method for Brauer tables returns the concatenation of the
identifier of the ordinary table, the string \texttt{"mod"}, and the (string of the) underlying characteristic. 

 
\begin{Verbatim}[commandchars=!@|,fontsize=\small,frame=single,label=Example]
  !gapprompt@gap>| !gapinput@Identifier( CharacterTable( "A5" ) );|
  "A5"
  !gapprompt@gap>| !gapinput@tbl:= CharacterTable( Group( () ) );;|
  !gapprompt@gap>| !gapinput@Identifier( tbl );  Identifier( tbl mod 2 );|
  "CT8"
  "CT8mod2"
\end{Verbatim}
 }

 

\subsection{\textcolor{Chapter }{InfoText (for character tables)}}
\logpage{[ 71, 9, 9 ]}\nobreak
\hyperdef{L}{X7932C35180C80953}{}
{\noindent\textcolor{FuncColor}{$\triangleright$\enspace\texttt{InfoText({\mdseries\slshape tbl})\index{InfoText@\texttt{InfoText}!for character tables}
\label{InfoText:for character tables}
}\hfill{\scriptsize (method)}}\\


 is a mutable string with information about the character table \mbox{\texttt{\mdseries\slshape tbl}}. There is no default method to create an info text. 

 This attribute is used mainly for library tables (see the manual of the \textsf{GAP} Character Table Library). Usual parts of the information are the origin of the
table, tests it has passed (\texttt{1.o.r.} for the test of orthogonality, \texttt{pow[\mbox{\texttt{\mdseries\slshape p}}]} for the construction of the \mbox{\texttt{\mdseries\slshape p}}\texttt{\symbol{45}}th power map, \texttt{DEC} for the decomposition of ordinary into Brauer characters, \texttt{TENS} for the decomposition of tensor products of irreducibles), and choices made
without loss of generality. 

 
\begin{Verbatim}[commandchars=!@|,fontsize=\small,frame=single,label=Example]
  !gapprompt@gap>| !gapinput@Print( InfoText( CharacterTable( "A5" ) ), "\n" );|
  origin: ATLAS of finite groups, tests: 1.o.r., pow[2,3,5]
\end{Verbatim}
 }

 

\subsection{\textcolor{Chapter }{InverseClasses}}
\logpage{[ 71, 9, 10 ]}\nobreak
\hyperdef{L}{X7919E2897BE8234A}{}
{\noindent\textcolor{FuncColor}{$\triangleright$\enspace\texttt{InverseClasses({\mdseries\slshape tbl})\index{InverseClasses@\texttt{InverseClasses}}
\label{InverseClasses}
}\hfill{\scriptsize (attribute)}}\\


 For a character table \mbox{\texttt{\mdseries\slshape tbl}}, \texttt{InverseClasses} returns the list mapping each conjugacy class to its inverse class. This list
can be regarded as $(-1)$\texttt{\symbol{45}}st power map of \mbox{\texttt{\mdseries\slshape tbl}} (see{\nobreakspace}\texttt{PowerMap} (\ref{PowerMap})). 

 
\begin{Verbatim}[commandchars=!@|,fontsize=\small,frame=single,label=Example]
  !gapprompt@gap>| !gapinput@InverseClasses( CharacterTable( "A5" ) );|
  [ 1, 2, 3, 4, 5 ]
  !gapprompt@gap>| !gapinput@InverseClasses( CharacterTable( "Cyclic", 3 ) );|
  [ 1, 3, 2 ]
\end{Verbatim}
 }

 

\subsection{\textcolor{Chapter }{RealClasses}}
\logpage{[ 71, 9, 11 ]}\nobreak
\hyperdef{L}{X87FF547981456932}{}
{\noindent\textcolor{FuncColor}{$\triangleright$\enspace\texttt{RealClasses({\mdseries\slshape tbl})\index{RealClasses@\texttt{RealClasses}}
\label{RealClasses}
}\hfill{\scriptsize (attribute)}}\\


 \index{classes!real} For a character table \mbox{\texttt{\mdseries\slshape tbl}}, \texttt{RealClasses} returns the strictly sorted list of positions of classes in \mbox{\texttt{\mdseries\slshape tbl}} that consist of real elements. 

 An element $x$ is \emph{real} iff it is conjugate to its inverse $x^{{-1}} = x^{{o(x)-1}}$. 

 
\begin{Verbatim}[commandchars=!@|,fontsize=\small,frame=single,label=Example]
  !gapprompt@gap>| !gapinput@RealClasses( CharacterTable( "A5" ) );|
  [ 1, 2, 3, 4, 5 ]
  !gapprompt@gap>| !gapinput@RealClasses( CharacterTable( "Cyclic", 3 ) );|
  [ 1 ]
\end{Verbatim}
 }

 

\subsection{\textcolor{Chapter }{ClassOrbit}}
\logpage{[ 71, 9, 12 ]}\nobreak
\hyperdef{L}{X7ABB007C799F7C49}{}
{\noindent\textcolor{FuncColor}{$\triangleright$\enspace\texttt{ClassOrbit({\mdseries\slshape tbl, cc})\index{ClassOrbit@\texttt{ClassOrbit}}
\label{ClassOrbit}
}\hfill{\scriptsize (operation)}}\\


 is the list of positions of those conjugacy classes of the character table \mbox{\texttt{\mdseries\slshape tbl}} that are Galois conjugate to the \mbox{\texttt{\mdseries\slshape cc}}\texttt{\symbol{45}}th class. That is, exactly the classes at positions given
by the list returned by \texttt{ClassOrbit} contain generators of the cyclic group generated by an element in the \mbox{\texttt{\mdseries\slshape cc}}\texttt{\symbol{45}}th class. 

 This information is computed from the power maps of \mbox{\texttt{\mdseries\slshape tbl}}. 

 
\begin{Verbatim}[commandchars=!@|,fontsize=\small,frame=single,label=Example]
  !gapprompt@gap>| !gapinput@ClassOrbit( CharacterTable( "A5" ), 4 );|
  [ 4, 5 ]
\end{Verbatim}
 }

 

\subsection{\textcolor{Chapter }{ClassRoots}}
\logpage{[ 71, 9, 13 ]}\nobreak
\hyperdef{L}{X7F863B15804E0835}{}
{\noindent\textcolor{FuncColor}{$\triangleright$\enspace\texttt{ClassRoots({\mdseries\slshape tbl})\index{ClassRoots@\texttt{ClassRoots}}
\label{ClassRoots}
}\hfill{\scriptsize (attribute)}}\\


 For a character table \mbox{\texttt{\mdseries\slshape tbl}}, \texttt{ClassRoots} returns a list containing at position $i$ the list of positions of the classes of all nontrivial $p$\texttt{\symbol{45}}th roots, where $p$ runs over the prime divisors of the \texttt{Size} (\ref{Size:for a character table}) value of \mbox{\texttt{\mdseries\slshape tbl}}. 

 This information is computed from the power maps of \mbox{\texttt{\mdseries\slshape tbl}}. 

 
\begin{Verbatim}[commandchars=!@|,fontsize=\small,frame=single,label=Example]
  !gapprompt@gap>| !gapinput@ClassRoots( CharacterTable( "A5" ) );|
  [ [ 2, 3, 4, 5 ], [  ], [  ], [  ], [  ] ]
  !gapprompt@gap>| !gapinput@ClassRoots( CharacterTable( "Cyclic", 6 ) );|
  [ [ 3, 4, 5 ], [  ], [ 2 ], [ 2, 6 ], [ 6 ], [  ] ]
\end{Verbatim}
 }

 }

  
\section{\textcolor{Chapter }{Normal Subgroups Represented by Lists of Class Positions}}\label{Normal Subgroups Represented by Lists of Class Positions}
\logpage{[ 71, 10, 0 ]}
\hyperdef{L}{X79CEBC3C7E0E63DF}{}
{
  The following attributes for a character table \mbox{\texttt{\mdseries\slshape tbl}} correspond to attributes for the group $G$ of \mbox{\texttt{\mdseries\slshape tbl}}. But instead of a normal subgroup (or a list of normal subgroups) of $G$, they return a strictly sorted list of positive integers (or a list of such
lists) which are the positions {\textendash}relative to the \texttt{ConjugacyClasses} (\ref{ConjugacyClasses:for character tables}) value of \mbox{\texttt{\mdseries\slshape tbl}}{\textendash} of those classes forming the normal subgroup in question. 

\subsection{\textcolor{Chapter }{ClassPositionsOfNormalSubgroups}}
\logpage{[ 71, 10, 1 ]}\nobreak
\hyperdef{L}{X850C7D947B3DBFA2}{}
{\noindent\textcolor{FuncColor}{$\triangleright$\enspace\texttt{ClassPositionsOfNormalSubgroups({\mdseries\slshape ordtbl})\index{ClassPositionsOfNormalSubgroups@\texttt{ClassPositionsOfNormalSubgroups}}
\label{ClassPositionsOfNormalSubgroups}
}\hfill{\scriptsize (attribute)}}\\
\noindent\textcolor{FuncColor}{$\triangleright$\enspace\texttt{ClassPositionsOfMaximalNormalSubgroups({\mdseries\slshape ordtbl})\index{ClassPositionsOfMaximalNormalSubgroups@\texttt{Class}\-\texttt{Positions}\-\texttt{Of}\-\texttt{Maximal}\-\texttt{Normal}\-\texttt{Subgroups}}
\label{ClassPositionsOfMaximalNormalSubgroups}
}\hfill{\scriptsize (attribute)}}\\
\noindent\textcolor{FuncColor}{$\triangleright$\enspace\texttt{ClassPositionsOfMinimalNormalSubgroups({\mdseries\slshape ordtbl})\index{ClassPositionsOfMinimalNormalSubgroups@\texttt{Class}\-\texttt{Positions}\-\texttt{Of}\-\texttt{Minimal}\-\texttt{Normal}\-\texttt{Subgroups}}
\label{ClassPositionsOfMinimalNormalSubgroups}
}\hfill{\scriptsize (attribute)}}\\


 correspond to \texttt{NormalSubgroups} (\ref{NormalSubgroups}), \texttt{MaximalNormalSubgroups} (\ref{MaximalNormalSubgroups}), \texttt{MinimalNormalSubgroups} (\ref{MinimalNormalSubgroups}) for the group of the ordinary character table \mbox{\texttt{\mdseries\slshape ordtbl}}. 

 The entries of the result lists are sorted according to increasing length. (So
this total order respects the partial order of normal subgroups given by
inclusion.) 

 
\begin{Verbatim}[commandchars=!@|,fontsize=\small,frame=single,label=Example]
  !gapprompt@gap>| !gapinput@tbls4:= CharacterTable( "Symmetric", 4 );;|
  !gapprompt@gap>| !gapinput@ClassPositionsOfNormalSubgroups( tbls4 );|
  [ [ 1 ], [ 1, 3 ], [ 1, 3, 4 ], [ 1 .. 5 ] ]
\end{Verbatim}
 }

 

\subsection{\textcolor{Chapter }{ClassPositionsOfAgemo}}
\logpage{[ 71, 10, 2 ]}\nobreak
\hyperdef{L}{X8491DA0981D6F264}{}
{\noindent\textcolor{FuncColor}{$\triangleright$\enspace\texttt{ClassPositionsOfAgemo({\mdseries\slshape ordtbl, p})\index{ClassPositionsOfAgemo@\texttt{ClassPositionsOfAgemo}}
\label{ClassPositionsOfAgemo}
}\hfill{\scriptsize (operation)}}\\


 corresponds to \texttt{Agemo} (\ref{Agemo}) for the group of the ordinary character table \mbox{\texttt{\mdseries\slshape ordtbl}}. 

 
\begin{Verbatim}[commandchars=!@|,fontsize=\small,frame=single,label=Example]
  !gapprompt@gap>| !gapinput@tbls4:= CharacterTable( "Symmetric", 4 );;|
  !gapprompt@gap>| !gapinput@ClassPositionsOfAgemo( tbls4, 2 );|
  [ 1, 3, 4 ]
\end{Verbatim}
 }

 

\subsection{\textcolor{Chapter }{ClassPositionsOfCentre (for a character table)}}
\logpage{[ 71, 10, 3 ]}\nobreak
\hyperdef{L}{X7A6B1F8A84A495DC}{}
{\noindent\textcolor{FuncColor}{$\triangleright$\enspace\texttt{ClassPositionsOfCentre({\mdseries\slshape ordtbl})\index{ClassPositionsOfCentre@\texttt{ClassPositionsOfCentre}!for a character table}
\label{ClassPositionsOfCentre:for a character table}
}\hfill{\scriptsize (attribute)}}\\
\noindent\textcolor{FuncColor}{$\triangleright$\enspace\texttt{ClassPositionsOfCenter({\mdseries\slshape ordtbl})\index{ClassPositionsOfCenter@\texttt{ClassPositionsOfCenter}!for a character table}
\label{ClassPositionsOfCenter:for a character table}
}\hfill{\scriptsize (attribute)}}\\


 corresponds to \texttt{Centre} (\ref{Centre}) for the group of the ordinary character table \mbox{\texttt{\mdseries\slshape ordtbl}}. 

 
\begin{Verbatim}[commandchars=!@|,fontsize=\small,frame=single,label=Example]
  !gapprompt@gap>| !gapinput@tbld8:= CharacterTable( "Dihedral", 8 );;|
  !gapprompt@gap>| !gapinput@ClassPositionsOfCentre( tbld8 );|
  [ 1, 3 ]
\end{Verbatim}
 }

 

\subsection{\textcolor{Chapter }{ClassPositionsOfDirectProductDecompositions}}
\logpage{[ 71, 10, 4 ]}\nobreak
\hyperdef{L}{X7D53F60785AB22B1}{}
{\noindent\textcolor{FuncColor}{$\triangleright$\enspace\texttt{ClassPositionsOfDirectProductDecompositions({\mdseries\slshape tbl[, nclasses]})\index{ClassPositionsOfDirectProductDecompositions@\texttt{Class}\-\texttt{Positions}\-\texttt{Of}\-\texttt{Direct}\-\texttt{Product}\-\texttt{Decompositions}}
\label{ClassPositionsOfDirectProductDecompositions}
}\hfill{\scriptsize (attribute)}}\\


 Let \mbox{\texttt{\mdseries\slshape tbl}} be the ordinary character table of the group $G$, say. Called with the only argument \mbox{\texttt{\mdseries\slshape tbl}}, \texttt{ClassPositionsOfDirectProductDecompositions} returns the list of all those pairs $[ l_1, l_2 ]$ where $l_1$ and $l_2$ are lists of class positions of normal subgroups $N_1$, $N_2$ of $G$ such that $G$ is their direct product and $|N_1| \leq |N_2|$ holds. Called with second argument a list \mbox{\texttt{\mdseries\slshape nclasses}} of class positions of a normal subgroup $N$ of $G$, \texttt{ClassPositionsOfDirectProductDecompositions} returns the list of pairs describing the decomposition of $N$ as a direct product of two normal subgroups of $G$. }

 

\subsection{\textcolor{Chapter }{ClassPositionsOfDerivedSubgroup}}
\logpage{[ 71, 10, 5 ]}\nobreak
\hyperdef{L}{X79EE7BE17BD343D5}{}
{\noindent\textcolor{FuncColor}{$\triangleright$\enspace\texttt{ClassPositionsOfDerivedSubgroup({\mdseries\slshape ordtbl})\index{ClassPositionsOfDerivedSubgroup@\texttt{ClassPositionsOfDerivedSubgroup}}
\label{ClassPositionsOfDerivedSubgroup}
}\hfill{\scriptsize (attribute)}}\\


 corresponds to \texttt{DerivedSubgroup} (\ref{DerivedSubgroup}) for the group of the ordinary character table \mbox{\texttt{\mdseries\slshape ordtbl}}. 

 
\begin{Verbatim}[commandchars=!@|,fontsize=\small,frame=single,label=Example]
  !gapprompt@gap>| !gapinput@tbld8:= CharacterTable( "Dihedral", 8 );;|
  !gapprompt@gap>| !gapinput@ClassPositionsOfDerivedSubgroup( tbld8 );|
  [ 1, 3 ]
\end{Verbatim}
 }

 

\subsection{\textcolor{Chapter }{ClassPositionsOfElementaryAbelianSeries}}
\logpage{[ 71, 10, 6 ]}\nobreak
\hyperdef{L}{X86ABB2E179D7F6E1}{}
{\noindent\textcolor{FuncColor}{$\triangleright$\enspace\texttt{ClassPositionsOfElementaryAbelianSeries({\mdseries\slshape ordtbl})\index{ClassPositionsOfElementaryAbelianSeries@\texttt{Class}\-\texttt{Positions}\-\texttt{Of}\-\texttt{Elementary}\-\texttt{Abelian}\-\texttt{Series}}
\label{ClassPositionsOfElementaryAbelianSeries}
}\hfill{\scriptsize (attribute)}}\\


 corresponds to \texttt{ElementaryAbelianSeries} (\ref{ElementaryAbelianSeries:for a group}) for the group of the ordinary character table \mbox{\texttt{\mdseries\slshape ordtbl}}. 

 
\begin{Verbatim}[commandchars=!@|,fontsize=\small,frame=single,label=Example]
  !gapprompt@gap>| !gapinput@tbls4:= CharacterTable( "Symmetric", 4 );;|
  !gapprompt@gap>| !gapinput@tbla5:= CharacterTable( "A5" );;|
  !gapprompt@gap>| !gapinput@ClassPositionsOfElementaryAbelianSeries( tbls4 );|
  [ [ 1 .. 5 ], [ 1, 3, 4 ], [ 1, 3 ], [ 1 ] ]
  !gapprompt@gap>| !gapinput@ClassPositionsOfElementaryAbelianSeries( tbla5 );|
  fail
\end{Verbatim}
 }

 

\subsection{\textcolor{Chapter }{ClassPositionsOfFittingSubgroup}}
\logpage{[ 71, 10, 7 ]}\nobreak
\hyperdef{L}{X7D2A55A584F955DB}{}
{\noindent\textcolor{FuncColor}{$\triangleright$\enspace\texttt{ClassPositionsOfFittingSubgroup({\mdseries\slshape ordtbl})\index{ClassPositionsOfFittingSubgroup@\texttt{ClassPositionsOfFittingSubgroup}}
\label{ClassPositionsOfFittingSubgroup}
}\hfill{\scriptsize (attribute)}}\\


 corresponds to \texttt{FittingSubgroup} (\ref{FittingSubgroup}) for the group of the ordinary character table \mbox{\texttt{\mdseries\slshape ordtbl}}. 

 
\begin{Verbatim}[commandchars=!@|,fontsize=\small,frame=single,label=Example]
  !gapprompt@gap>| !gapinput@tbls4:= CharacterTable( "Symmetric", 4 );;|
  !gapprompt@gap>| !gapinput@ClassPositionsOfFittingSubgroup( tbls4 );|
  [ 1, 3 ]
\end{Verbatim}
 }

 

\subsection{\textcolor{Chapter }{ClassPositionsOfLowerCentralSeries}}
\logpage{[ 71, 10, 8 ]}\nobreak
\hyperdef{L}{X79AEFC4384769B72}{}
{\noindent\textcolor{FuncColor}{$\triangleright$\enspace\texttt{ClassPositionsOfLowerCentralSeries({\mdseries\slshape tbl})\index{ClassPositionsOfLowerCentralSeries@\texttt{ClassPositionsOfLowerCentralSeries}}
\label{ClassPositionsOfLowerCentralSeries}
}\hfill{\scriptsize (attribute)}}\\


 corresponds to \texttt{LowerCentralSeriesOfGroup} (\ref{LowerCentralSeriesOfGroup}) for the group of the ordinary character table \mbox{\texttt{\mdseries\slshape ordtbl}}. 

 
\begin{Verbatim}[commandchars=!@|,fontsize=\small,frame=single,label=Example]
  !gapprompt@gap>| !gapinput@tbls4:= CharacterTable( "Symmetric", 4 );;|
  !gapprompt@gap>| !gapinput@tbld8:= CharacterTable( "Dihedral", 8 );;|
  !gapprompt@gap>| !gapinput@ClassPositionsOfLowerCentralSeries( tbls4 );|
  [ [ 1 .. 5 ], [ 1, 3, 4 ] ]
  !gapprompt@gap>| !gapinput@ClassPositionsOfLowerCentralSeries( tbld8 );|
  [ [ 1 .. 5 ], [ 1, 3 ], [ 1 ] ]
\end{Verbatim}
 }

 

\subsection{\textcolor{Chapter }{ClassPositionsOfUpperCentralSeries}}
\logpage{[ 71, 10, 9 ]}\nobreak
\hyperdef{L}{X86065D217A36CD9B}{}
{\noindent\textcolor{FuncColor}{$\triangleright$\enspace\texttt{ClassPositionsOfUpperCentralSeries({\mdseries\slshape ordtbl})\index{ClassPositionsOfUpperCentralSeries@\texttt{ClassPositionsOfUpperCentralSeries}}
\label{ClassPositionsOfUpperCentralSeries}
}\hfill{\scriptsize (attribute)}}\\


 corresponds to \texttt{UpperCentralSeriesOfGroup} (\ref{UpperCentralSeriesOfGroup}) for the group of the ordinary character table \mbox{\texttt{\mdseries\slshape ordtbl}}. 

 
\begin{Verbatim}[commandchars=!@|,fontsize=\small,frame=single,label=Example]
  !gapprompt@gap>| !gapinput@tbls4:= CharacterTable( "Symmetric", 4 );;|
  !gapprompt@gap>| !gapinput@tbld8:= CharacterTable( "Dihedral", 8 );;|
  !gapprompt@gap>| !gapinput@ClassPositionsOfUpperCentralSeries( tbls4 );|
  [ [ 1 ] ]
  !gapprompt@gap>| !gapinput@ClassPositionsOfUpperCentralSeries( tbld8 );|
  [ [ 1, 3 ], [ 1, 2, 3, 4, 5 ] ]
\end{Verbatim}
 }

 

\subsection{\textcolor{Chapter }{ClassPositionsOfSolvableRadical}}
\logpage{[ 71, 10, 10 ]}\nobreak
\hyperdef{L}{X877FDE8A84A9F52C}{}
{\noindent\textcolor{FuncColor}{$\triangleright$\enspace\texttt{ClassPositionsOfSolvableRadical({\mdseries\slshape ordtbl})\index{ClassPositionsOfSolvableRadical@\texttt{ClassPositionsOfSolvableRadical}}
\label{ClassPositionsOfSolvableRadical}
}\hfill{\scriptsize (attribute)}}\\


 corresponds to \texttt{SolvableRadical} (\ref{SolvableRadical}) for the group of the ordinary character table \mbox{\texttt{\mdseries\slshape ordtbl}}. 

 
\begin{Verbatim}[commandchars=!@|,fontsize=\small,frame=single,label=Example]
  !gapprompt@gap>| !gapinput@ClassPositionsOfSolvableRadical( CharacterTable( "2.A5" ) );|
  [ 1, 2 ]
\end{Verbatim}
 }

 

\subsection{\textcolor{Chapter }{ClassPositionsOfSupersolvableResiduum}}
\logpage{[ 71, 10, 11 ]}\nobreak
\hyperdef{L}{X8392DD5B813250A4}{}
{\noindent\textcolor{FuncColor}{$\triangleright$\enspace\texttt{ClassPositionsOfSupersolvableResiduum({\mdseries\slshape ordtbl})\index{ClassPositionsOfSupersolvableResiduum@\texttt{Class}\-\texttt{Positions}\-\texttt{Of}\-\texttt{Supersolvable}\-\texttt{Residuum}}
\label{ClassPositionsOfSupersolvableResiduum}
}\hfill{\scriptsize (attribute)}}\\


 corresponds to \texttt{SupersolvableResiduum} (\ref{SupersolvableResiduum}) for the group of the ordinary character table \mbox{\texttt{\mdseries\slshape ordtbl}}. 

 
\begin{Verbatim}[commandchars=!@|,fontsize=\small,frame=single,label=Example]
  !gapprompt@gap>| !gapinput@tbls4:= CharacterTable( "Symmetric", 4 );;|
  !gapprompt@gap>| !gapinput@ClassPositionsOfSupersolvableResiduum( tbls4 );|
  [ 1, 3 ]
\end{Verbatim}
 }

 

\subsection{\textcolor{Chapter }{ClassPositionsOfPCore}}
\logpage{[ 71, 10, 12 ]}\nobreak
\hyperdef{L}{X7BBE7EBA7A64A6B0}{}
{\noindent\textcolor{FuncColor}{$\triangleright$\enspace\texttt{ClassPositionsOfPCore({\mdseries\slshape ordtbl, p})\index{ClassPositionsOfPCore@\texttt{ClassPositionsOfPCore}}
\label{ClassPositionsOfPCore}
}\hfill{\scriptsize (operation)}}\\


 corresponds to \texttt{PCore} (\ref{PCore}) for the group of the ordinary character table \mbox{\texttt{\mdseries\slshape ordtbl}}. 

 
\begin{Verbatim}[commandchars=!@|,fontsize=\small,frame=single,label=Example]
  !gapprompt@gap>| !gapinput@tbls4:= CharacterTable( "Symmetric", 4 );;|
  !gapprompt@gap>| !gapinput@ClassPositionsOfPCore( tbls4, 2 );|
  [ 1, 3 ]
  !gapprompt@gap>| !gapinput@ClassPositionsOfPCore( tbls4, 3 );|
  [ 1 ]
\end{Verbatim}
 }

 

\subsection{\textcolor{Chapter }{ClassPositionsOfNormalClosure}}
\logpage{[ 71, 10, 13 ]}\nobreak
\hyperdef{L}{X7FCF905D7FD7CC40}{}
{\noindent\textcolor{FuncColor}{$\triangleright$\enspace\texttt{ClassPositionsOfNormalClosure({\mdseries\slshape ordtbl, classes})\index{ClassPositionsOfNormalClosure@\texttt{ClassPositionsOfNormalClosure}}
\label{ClassPositionsOfNormalClosure}
}\hfill{\scriptsize (operation)}}\\


 is the sorted list of the positions of all conjugacy classes of the ordinary
character table \mbox{\texttt{\mdseries\slshape ordtbl}} that form the normal closure (see{\nobreakspace}\texttt{NormalClosure} (\ref{NormalClosure})) of the conjugacy classes at positions in the list \mbox{\texttt{\mdseries\slshape classes}}. 

 
\begin{Verbatim}[commandchars=!@|,fontsize=\small,frame=single,label=Example]
  !gapprompt@gap>| !gapinput@tbls4:= CharacterTable( "Symmetric", 4 );;|
  !gapprompt@gap>| !gapinput@ClassPositionsOfNormalClosure( tbls4, [ 1, 4 ] );|
  [ 1, 3, 4 ]
\end{Verbatim}
 }

 }

  
\section{\textcolor{Chapter }{Operations Concerning Blocks}}\label{Operations Concerning Blocks}
\logpage{[ 71, 11, 0 ]}
\hyperdef{L}{X8733F0EA801785D4}{}
{
  

\subsection{\textcolor{Chapter }{PrimeBlocks}}
\logpage{[ 71, 11, 1 ]}\nobreak
\hyperdef{L}{X7ACB9306804F4E3F}{}
{\noindent\textcolor{FuncColor}{$\triangleright$\enspace\texttt{PrimeBlocks({\mdseries\slshape ordtbl, p})\index{PrimeBlocks@\texttt{PrimeBlocks}}
\label{PrimeBlocks}
}\hfill{\scriptsize (operation)}}\\
\noindent\textcolor{FuncColor}{$\triangleright$\enspace\texttt{PrimeBlocksOp({\mdseries\slshape ordtbl, p})\index{PrimeBlocksOp@\texttt{PrimeBlocksOp}}
\label{PrimeBlocksOp}
}\hfill{\scriptsize (operation)}}\\
\noindent\textcolor{FuncColor}{$\triangleright$\enspace\texttt{ComputedPrimeBlockss({\mdseries\slshape tbl})\index{ComputedPrimeBlockss@\texttt{ComputedPrimeBlockss}}
\label{ComputedPrimeBlockss}
}\hfill{\scriptsize (attribute)}}\\


 For an ordinary character table \mbox{\texttt{\mdseries\slshape ordtbl}} and a prime integer \mbox{\texttt{\mdseries\slshape p}}, \texttt{PrimeBlocks} returns a record with the following components. 
\begin{description}
\item[{\texttt{block}}]  a list, the value $j$ at position $i$ means that the $i$\texttt{\symbol{45}}th irreducible character of \mbox{\texttt{\mdseries\slshape ordtbl}} lies in the $j$\texttt{\symbol{45}}th \mbox{\texttt{\mdseries\slshape p}}\texttt{\symbol{45}}block of \mbox{\texttt{\mdseries\slshape ordtbl}}, 
\item[{\texttt{defect}}]  a list containing at position $i$ the defect of the $i$\texttt{\symbol{45}}th block, 
\item[{\texttt{height}}]  a list containing at position $i$ the height of the $i$\texttt{\symbol{45}}th irreducible character of \mbox{\texttt{\mdseries\slshape ordtbl}} in its block, 
\item[{\texttt{relevant}}]  a list of class positions such that only the restriction to these classes need
be checked for deciding whether two characters lie in the same block, and 
\item[{\texttt{centralcharacter}}]  a list containing at position $i$ a list whose values at the positions stored in the component \texttt{relevant} are the values of a central character in the $i$\texttt{\symbol{45}}th block. 
\end{description}
 

 The components \texttt{relevant} and \texttt{centralcharacters} are used by \texttt{SameBlock} (\ref{SameBlock}). 

 If \texttt{InfoCharacterTable} (\ref{InfoCharacterTable}) has level at least 2, the defects of the blocks and the heights of the
characters are printed. 

 The default method uses the attribute \texttt{ComputedPrimeBlockss} for storing the computed value at position \mbox{\texttt{\mdseries\slshape p}}, and calls the operation \texttt{PrimeBlocksOp} for computing values that are not yet known. 

 Two ordinary irreducible characters $\chi, \psi$ of a group $G$ are said to lie in the same $p$\texttt{\symbol{45}}\emph{block} if the images of their central characters $\omega_{\chi}, \omega_{\psi}$ (see{\nobreakspace}\texttt{CentralCharacter} (\ref{CentralCharacter})) under the natural ring epimorphism $R \rightarrow R / M$ are equal, where $R$ denotes the ring of algebraic integers in the complex number field, and $M$ is a maximal ideal in $R$ with $pR \subseteq M$. (The distribution to $p$\texttt{\symbol{45}}blocks is in fact independent of the choice of $M$, see{\nobreakspace}\cite{Isa76}.) 

 For $|G| = p^a m$ where $p$ does not divide $m$, the \emph{defect} of a block is the integer $d$ such that $p^{{a-d}}$ is the largest power of $p$ that divides the degrees of all characters in the block. 

 The \emph{height} of a character $\chi$ in the block is defined as the largest exponent $h$ for which $p^h$ divides $\chi(1) / p^{{a-d}}$. 

 
\begin{Verbatim}[commandchars=!@|,fontsize=\small,frame=single,label=Example]
  !gapprompt@gap>| !gapinput@tbl:= CharacterTable( "L3(2)" );;|
  !gapprompt@gap>| !gapinput@pbl:= PrimeBlocks( tbl, 2 );|
  rec( block := [ 1, 1, 1, 1, 1, 2 ],
    centralcharacter := [ [ ,, 56,, 24 ], [ ,, -7,, 3 ] ],
    defect := [ 3, 0 ], height := [ 0, 0, 0, 1, 0, 0 ],
    relevant := [ 3, 5 ] )
\end{Verbatim}
 }

 

\subsection{\textcolor{Chapter }{SameBlock}}
\logpage{[ 71, 11, 2 ]}\nobreak
\hyperdef{L}{X7E80E35985275F35}{}
{\noindent\textcolor{FuncColor}{$\triangleright$\enspace\texttt{SameBlock({\mdseries\slshape p, omega1, omega2, relevant})\index{SameBlock@\texttt{SameBlock}}
\label{SameBlock}
}\hfill{\scriptsize (function)}}\\


 Let \mbox{\texttt{\mdseries\slshape p}} be a prime integer, \mbox{\texttt{\mdseries\slshape omega1}} and \mbox{\texttt{\mdseries\slshape omega2}} be two central characters (or their values lists) of a character table, and \mbox{\texttt{\mdseries\slshape relevant}} be a list of positions as is stored in the component \texttt{relevant} of a record returned by \texttt{PrimeBlocks} (\ref{PrimeBlocks}). 

 \texttt{SameBlock} returns \texttt{true} if \mbox{\texttt{\mdseries\slshape omega1}} and \mbox{\texttt{\mdseries\slshape omega2}} are equal modulo any maximal ideal in the ring of complex algebraic integers
containing the ideal spanned by \mbox{\texttt{\mdseries\slshape p}}, and \texttt{false} otherwise. 

 
\begin{Verbatim}[commandchars=!@|,fontsize=\small,frame=single,label=Example]
  !gapprompt@gap>| !gapinput@omega:= List( Irr( tbl ), CentralCharacter );;|
  !gapprompt@gap>| !gapinput@SameBlock( 2, omega[1], omega[2], pbl.relevant );|
  true
  !gapprompt@gap>| !gapinput@SameBlock( 2, omega[1], omega[6], pbl.relevant );|
  false
\end{Verbatim}
 }

 

\subsection{\textcolor{Chapter }{BlocksInfo}}
\logpage{[ 71, 11, 3 ]}\nobreak
\hyperdef{L}{X7FF4CE4A7A272F88}{}
{\noindent\textcolor{FuncColor}{$\triangleright$\enspace\texttt{BlocksInfo({\mdseries\slshape modtbl})\index{BlocksInfo@\texttt{BlocksInfo}}
\label{BlocksInfo}
}\hfill{\scriptsize (attribute)}}\\


 For a Brauer character table \mbox{\texttt{\mdseries\slshape modtbl}}, the value of \texttt{BlocksInfo} is a list of (mutable) records, the $i$\texttt{\symbol{45}}th entry containing information about the $i$\texttt{\symbol{45}}th block. Each record has the following components. 
\begin{description}
\item[{\texttt{defect}}]  the defect of the block, 
\item[{\texttt{ordchars}}]  the list of positions of the ordinary characters that belong to the block,
relative to \texttt{Irr( OrdinaryCharacterTable( \mbox{\texttt{\mdseries\slshape modtbl}} ) )}, 
\item[{\texttt{modchars}}]  the list of positions of the Brauer characters that belong to the block,
relative to \texttt{IBr( \mbox{\texttt{\mdseries\slshape modtbl}} )}. 
\end{description}
 Optional components are 
\begin{description}
\item[{\texttt{basicset}}]  a list of positions of ordinary characters in the block whose restriction to \mbox{\texttt{\mdseries\slshape modtbl}} is maximally linearly independent, relative to \texttt{Irr( OrdinaryCharacterTable( \mbox{\texttt{\mdseries\slshape modtbl}} ) )}, 
\item[{\texttt{decmat}}]  the decomposition matrix of the block, it is stored automatically when \texttt{DecompositionMatrix} (\ref{DecompositionMatrix}) is called for the block, 
\item[{\texttt{decinv}}]  inverse of the decomposition matrix of the block, restricted to the ordinary
characters described by \texttt{basicset}, 
\item[{\texttt{brauertree}}]  a list that describes the Brauer tree of the block, in the case that the block
is of defect $1$. 
\end{description}
 

 
\begin{Verbatim}[commandchars=!@|,fontsize=\small,frame=single,label=Example]
  !gapprompt@gap>| !gapinput@BlocksInfo( CharacterTable( "L3(2)" ) mod 2 );|
  [ rec( basicset := [ 1, 2, 3 ],
        decinv := [ [ 1, 0, 0 ], [ 0, 1, 0 ], [ 0, 0, 1 ] ],
        defect := 3, modchars := [ 1, 2, 3 ],
        ordchars := [ 1, 2, 3, 4, 5 ] ),
    rec( basicset := [ 6 ], decinv := [ [ 1 ] ], defect := 0,
        modchars := [ 4 ], ordchars := [ 6 ] ) ]
\end{Verbatim}
 }

 

\subsection{\textcolor{Chapter }{DecompositionMatrix}}
\logpage{[ 71, 11, 4 ]}\nobreak
\hyperdef{L}{X84701640811D2345}{}
{\noindent\textcolor{FuncColor}{$\triangleright$\enspace\texttt{DecompositionMatrix({\mdseries\slshape modtbl[, blocknr]})\index{DecompositionMatrix@\texttt{DecompositionMatrix}}
\label{DecompositionMatrix}
}\hfill{\scriptsize (operation)}}\\


 Let \mbox{\texttt{\mdseries\slshape modtbl}} be a Brauer character table. 

 Called with one argument, \texttt{DecompositionMatrix} returns the decomposition matrix of \mbox{\texttt{\mdseries\slshape modtbl}}, where the rows and columns are indexed by the irreducible characters of the
ordinary character table of \mbox{\texttt{\mdseries\slshape modtbl}} and the irreducible characters of \mbox{\texttt{\mdseries\slshape modtbl}}, respectively, 

 Called with two arguments, \texttt{DecompositionMatrix} returns the decomposition matrix of the block of \mbox{\texttt{\mdseries\slshape modtbl}} with number \mbox{\texttt{\mdseries\slshape blocknr}}; the matrix is stored as value of the \texttt{decmat} component of the \mbox{\texttt{\mdseries\slshape blocknr}}\texttt{\symbol{45}}th entry of the \texttt{BlocksInfo} (\ref{BlocksInfo}) list of \mbox{\texttt{\mdseries\slshape modtbl}}. 

 An ordinary irreducible character is in block $i$ if and only if all characters before the first character of the same block lie
in $i-1$ different blocks. An irreducible Brauer character is in block $i$ if it has nonzero scalar product with an ordinary irreducible character in
block $i$. 

 \texttt{DecompositionMatrix} is based on the more general function \texttt{Decomposition} (\ref{Decomposition}). 

 
\begin{Verbatim}[commandchars=!@|,fontsize=\small,frame=single,label=Example]
  !gapprompt@gap>| !gapinput@modtbl:= CharacterTable( "L3(2)" ) mod 2;|
  BrauerTable( "L3(2)", 2 )
  !gapprompt@gap>| !gapinput@DecompositionMatrix( modtbl );|
  [ [ 1, 0, 0, 0 ], [ 0, 1, 0, 0 ], [ 0, 0, 1, 0 ], [ 0, 1, 1, 0 ],
    [ 1, 1, 1, 0 ], [ 0, 0, 0, 1 ] ]
  !gapprompt@gap>| !gapinput@DecompositionMatrix( modtbl, 1 );|
  [ [ 1, 0, 0 ], [ 0, 1, 0 ], [ 0, 0, 1 ], [ 0, 1, 1 ], [ 1, 1, 1 ] ]
  !gapprompt@gap>| !gapinput@DecompositionMatrix( modtbl, 2 );|
  [ [ 1 ] ]
\end{Verbatim}
 }

 \index{LaTeX!for a decomposition matrix} 

\subsection{\textcolor{Chapter }{LaTeXStringDecompositionMatrix}}
\logpage{[ 71, 11, 5 ]}\nobreak
\hyperdef{L}{X83EC921380AF9B3B}{}
{\noindent\textcolor{FuncColor}{$\triangleright$\enspace\texttt{LaTeXStringDecompositionMatrix({\mdseries\slshape modtbl[, blocknr][, options]})\index{LaTeXStringDecompositionMatrix@\texttt{LaTeXStringDecompositionMatrix}}
\label{LaTeXStringDecompositionMatrix}
}\hfill{\scriptsize (function)}}\\


 is a string that contains La{\TeX} code to print a decomposition matrix (see{\nobreakspace}\texttt{DecompositionMatrix} (\ref{DecompositionMatrix})) nicely. 

 The optional argument \mbox{\texttt{\mdseries\slshape options}}, if present, must be a record with components \texttt{phi}, \texttt{chi} (strings used in each label for columns and rows), \texttt{collabels}, \texttt{rowlabels} (subscripts for the labels). The defaults for \texttt{phi} and \texttt{chi} are \texttt{"\texttt{\symbol{123}}\texttt{\symbol{92}}\texttt{\symbol{92}}tt
Y\texttt{\symbol{125}}"} and \texttt{"\texttt{\symbol{123}}\texttt{\symbol{92}}\texttt{\symbol{92}}tt
X\texttt{\symbol{125}}"}, the defaults for \texttt{collabels} and \texttt{rowlabels} are the lists of positions of the Brauer characters and ordinary characters in
the respective lists of irreducibles in the character tables. 

 The optional components \texttt{nrows} and \texttt{ncols} denote the maximal number of rows and columns per array; if they are present
then each portion of \texttt{nrows} rows and \texttt{ncols} columns forms an array of its own which is enclosed in \texttt{\texttt{\symbol{92}}[}, \texttt{\texttt{\symbol{92}}]}. 

 If the component \texttt{decmat} is bound in \mbox{\texttt{\mdseries\slshape options}} then it must be the decomposition matrix in question, in this case the matrix
is not computed from the information in \mbox{\texttt{\mdseries\slshape modtbl}}. 

 For those character tables from the \textsf{GAP} table library that belong to the \textsf{Atlas} of Finite Groups{\nobreakspace}\cite{CCN85}, \texttt{AtlasLabelsOfIrreducibles} (\textbf{CTblLib: AtlasLabelsOfIrreducibles}) constructs character labels that are compatible with those used in the \textsf{Atlas} (see{\nobreakspace} (\textbf{CTblLib: Atlas Tables}) in the manual of the \textsf{GAP} Character Table Library). 

 
\begin{Verbatim}[commandchars=!@A,fontsize=\small,frame=single,label=Example]
  !gapprompt@gap>A !gapinput@modtbl:= CharacterTable( "L3(2)" ) mod 2;;A
  !gapprompt@gap>A !gapinput@Print( LaTeXStringDecompositionMatrix( modtbl, 1 ) );A
  \[
  \begin{array}{r|rrr} \hline
   & {\tt Y}_{1}
   & {\tt Y}_{2}
   & {\tt Y}_{3}
   \rule[-7pt]{0pt}{20pt} \\ \hline
  {\tt X}_{1} & 1 & . & . \rule[0pt]{0pt}{13pt} \\
  {\tt X}_{2} & . & 1 & . \\
  {\tt X}_{3} & . & . & 1 \\
  {\tt X}_{4} & . & 1 & 1 \\
  {\tt X}_{5} & 1 & 1 & 1 \rule[-7pt]{0pt}{5pt} \\
  \hline
  \end{array}
  \]
  !gapprompt@gap>A !gapinput@options:= rec( phi:= "\\varphi", chi:= "\\chi" );;A
  !gapprompt@gap>A !gapinput@Print( LaTeXStringDecompositionMatrix( modtbl, 1, options ) );A
  \[
  \begin{array}{r|rrr} \hline
   & \varphi_{1}
   & \varphi_{2}
   & \varphi_{3}
   \rule[-7pt]{0pt}{20pt} \\ \hline
  \chi_{1} & 1 & . & . \rule[0pt]{0pt}{13pt} \\
  \chi_{2} & . & 1 & . \\
  \chi_{3} & . & . & 1 \\
  \chi_{4} & . & 1 & 1 \\
  \chi_{5} & 1 & 1 & 1 \rule[-7pt]{0pt}{5pt} \\
  \hline
  \end{array}
  \]
\end{Verbatim}
 }

 }

  
\section{\textcolor{Chapter }{Other Operations for Character Tables}}\label{Other Operations for Character Tables}
\logpage{[ 71, 12, 0 ]}
\hyperdef{L}{X873211618402ACF7}{}
{
  In the following, we list operations for character tables that are not
attributes. 

 

\subsection{\textcolor{Chapter }{Index (for two character tables)}}
\logpage{[ 71, 12, 1 ]}\nobreak
\hyperdef{L}{X8441983C845F2176}{}
{\noindent\textcolor{FuncColor}{$\triangleright$\enspace\texttt{Index({\mdseries\slshape tbl, subtbl})\index{Index@\texttt{Index}!for two character tables}
\label{Index:for two character tables}
}\hfill{\scriptsize (operation)}}\\


 For two character tables \mbox{\texttt{\mdseries\slshape tbl}} and \mbox{\texttt{\mdseries\slshape subtbl}}, \texttt{Index} returns the quotient of the \texttt{Size} (\ref{Size:for a character table}) values of \mbox{\texttt{\mdseries\slshape tbl}} and \mbox{\texttt{\mdseries\slshape subtbl}}. The containment of the underlying groups of \mbox{\texttt{\mdseries\slshape subtbl}} and \mbox{\texttt{\mdseries\slshape tbl}} is \emph{not} checked; so the distinction between \texttt{Index} (\ref{Index:for a group and its subgroup}) and \texttt{IndexNC} (\ref{IndexNC:for a group and its subgroup}) is not made for character tables. }

 

\subsection{\textcolor{Chapter }{IsInternallyConsistent (for character tables)}}
\logpage{[ 71, 12, 2 ]}\nobreak
\hyperdef{L}{X8123650E817926FC}{}
{\noindent\textcolor{FuncColor}{$\triangleright$\enspace\texttt{IsInternallyConsistent({\mdseries\slshape tbl})\index{IsInternallyConsistent@\texttt{IsInternallyConsistent}!for character tables}
\label{IsInternallyConsistent:for character tables}
}\hfill{\scriptsize (method)}}\\


 For an \emph{ordinary} character table \mbox{\texttt{\mdseries\slshape tbl}}, \texttt{IsInternallyConsistent} (\ref{IsInternallyConsistent}) checks the consistency of the following attribute values (if stored). 
\begin{itemize}
\item  \texttt{Size} (\ref{Size}), \texttt{SizesCentralizers} (\ref{SizesCentralizers}), and \texttt{SizesConjugacyClasses} (\ref{SizesConjugacyClasses}). 
\item  \texttt{SizesCentralizers} (\ref{SizesCentralizers}) and \texttt{OrdersClassRepresentatives} (\ref{OrdersClassRepresentatives}). 
\item  \texttt{ComputedPowerMaps} (\ref{ComputedPowerMaps}) and \texttt{OrdersClassRepresentatives} (\ref{OrdersClassRepresentatives}). 
\item  \texttt{SizesCentralizers} (\ref{SizesCentralizers}) and \texttt{Irr} (\ref{Irr:for a character table}). 
\item  \texttt{Irr} (\ref{Irr:for a character table}) (first orthogonality relation). 
\end{itemize}
 

 For a \emph{Brauer} table \mbox{\texttt{\mdseries\slshape tbl}}, \texttt{IsInternallyConsistent} checks the consistency of the following attribute values (if stored). 
\begin{itemize}
\item  \texttt{Size} (\ref{Size}), \texttt{SizesCentralizers} (\ref{SizesCentralizers}), and \texttt{SizesConjugacyClasses} (\ref{SizesConjugacyClasses}). 
\item  \texttt{SizesCentralizers} (\ref{SizesCentralizers}) and \texttt{OrdersClassRepresentatives} (\ref{OrdersClassRepresentatives}). 
\item  \texttt{ComputedPowerMaps} (\ref{ComputedPowerMaps}) and \texttt{OrdersClassRepresentatives} (\ref{OrdersClassRepresentatives}). 
\item  \texttt{Irr} (\ref{Irr:for a character table}) (closure under complex conjugation and Frobenius map). 
\end{itemize}
 

 If no inconsistency occurs, \texttt{true} is returned, otherwise each inconsistency is printed to the screen if the
level of \texttt{InfoWarning} (\ref{InfoWarning}) is at least $1$ (see{\nobreakspace}\ref{Info Functions}), and \texttt{false} is returned at the end. }

 

\subsection{\textcolor{Chapter }{IsPSolvableCharacterTable}}
\logpage{[ 71, 12, 3 ]}\nobreak
\hyperdef{L}{X7A0CBD1884276882}{}
{\noindent\textcolor{FuncColor}{$\triangleright$\enspace\texttt{IsPSolvableCharacterTable({\mdseries\slshape tbl, p})\index{IsPSolvableCharacterTable@\texttt{IsPSolvableCharacterTable}}
\label{IsPSolvableCharacterTable}
}\hfill{\scriptsize (operation)}}\\
\noindent\textcolor{FuncColor}{$\triangleright$\enspace\texttt{IsPSolubleCharacterTable({\mdseries\slshape tbl, p})\index{IsPSolubleCharacterTable@\texttt{IsPSolubleCharacterTable}}
\label{IsPSolubleCharacterTable}
}\hfill{\scriptsize (operation)}}\\
\noindent\textcolor{FuncColor}{$\triangleright$\enspace\texttt{IsPSolvableCharacterTableOp({\mdseries\slshape tbl, p})\index{IsPSolvableCharacterTableOp@\texttt{IsPSolvableCharacterTableOp}}
\label{IsPSolvableCharacterTableOp}
}\hfill{\scriptsize (operation)}}\\
\noindent\textcolor{FuncColor}{$\triangleright$\enspace\texttt{IsPSolubleCharacterTableOp({\mdseries\slshape tbl, p})\index{IsPSolubleCharacterTableOp@\texttt{IsPSolubleCharacterTableOp}}
\label{IsPSolubleCharacterTableOp}
}\hfill{\scriptsize (operation)}}\\
\noindent\textcolor{FuncColor}{$\triangleright$\enspace\texttt{ComputedIsPSolvableCharacterTables({\mdseries\slshape tbl})\index{ComputedIsPSolvableCharacterTables@\texttt{ComputedIsPSolvableCharacterTables}}
\label{ComputedIsPSolvableCharacterTables}
}\hfill{\scriptsize (attribute)}}\\
\noindent\textcolor{FuncColor}{$\triangleright$\enspace\texttt{ComputedIsPSolubleCharacterTables({\mdseries\slshape tbl})\index{ComputedIsPSolubleCharacterTables@\texttt{ComputedIsPSolubleCharacterTables}}
\label{ComputedIsPSolubleCharacterTables}
}\hfill{\scriptsize (attribute)}}\\


 \texttt{IsPSolvableCharacterTable} for the ordinary character table \mbox{\texttt{\mdseries\slshape tbl}} corresponds to \texttt{IsPSolvable} (\ref{IsPSolvable}) for the group of \mbox{\texttt{\mdseries\slshape tbl}}, \mbox{\texttt{\mdseries\slshape p}} must be either a prime integer or \texttt{0}. 

 The default method uses the attribute \texttt{ComputedIsPSolvableCharacterTables} for storing the computed value at position \mbox{\texttt{\mdseries\slshape p}}, and calls the operation \texttt{IsPSolvableCharacterTableOp} for computing values that are not yet known. 

 
\begin{Verbatim}[commandchars=!@|,fontsize=\small,frame=single,label=Example]
  !gapprompt@gap>| !gapinput@tbl:= CharacterTable( "Sz(8)" );;|
  !gapprompt@gap>| !gapinput@IsPSolvableCharacterTable( tbl, 2 );|
  false
  !gapprompt@gap>| !gapinput@IsPSolvableCharacterTable( tbl, 3 );|
  true
\end{Verbatim}
 }

 

\subsection{\textcolor{Chapter }{IsClassFusionOfNormalSubgroup}}
\logpage{[ 71, 12, 4 ]}\nobreak
\hyperdef{L}{X82F523E8784B3752}{}
{\noindent\textcolor{FuncColor}{$\triangleright$\enspace\texttt{IsClassFusionOfNormalSubgroup({\mdseries\slshape subtbl, fus, tbl})\index{IsClassFusionOfNormalSubgroup@\texttt{IsClassFusionOfNormalSubgroup}}
\label{IsClassFusionOfNormalSubgroup}
}\hfill{\scriptsize (function)}}\\


 For two ordinary character tables \mbox{\texttt{\mdseries\slshape tbl}} and \mbox{\texttt{\mdseries\slshape subtbl}} of a group $G$ and its subgroup $U$ and a list \mbox{\texttt{\mdseries\slshape fus}} of positive integers that describes the class fusion of $U$ into $G$, \texttt{IsClassFusionOfNormalSubgroup} returns \texttt{true} if $U$ is a normal subgroup of $G$, and \texttt{false} otherwise. 

 
\begin{Verbatim}[commandchars=!@|,fontsize=\small,frame=single,label=Example]
  !gapprompt@gap>| !gapinput@tblc2:= CharacterTable( "Cyclic", 2 );;|
  !gapprompt@gap>| !gapinput@tbld8:= CharacterTable( "Dihedral", 8 );;|
  !gapprompt@gap>| !gapinput@fus:= PossibleClassFusions( tblc2, tbld8 );|
  [ [ 1, 3 ], [ 1, 4 ], [ 1, 5 ] ]
  !gapprompt@gap>| !gapinput@List(fus, map -> IsClassFusionOfNormalSubgroup(tblc2, map, tbld8));|
  [ true, false, false ]
\end{Verbatim}
 }

 

\subsection{\textcolor{Chapter }{Indicator}}
\logpage{[ 71, 12, 5 ]}\nobreak
\hyperdef{L}{X7FD3D3047DE6381E}{}
{\noindent\textcolor{FuncColor}{$\triangleright$\enspace\texttt{Indicator({\mdseries\slshape tbl[, characters], n})\index{Indicator@\texttt{Indicator}}
\label{Indicator}
}\hfill{\scriptsize (operation)}}\\
\noindent\textcolor{FuncColor}{$\triangleright$\enspace\texttt{IndicatorOp({\mdseries\slshape tbl, characters, n})\index{IndicatorOp@\texttt{IndicatorOp}}
\label{IndicatorOp}
}\hfill{\scriptsize (operation)}}\\
\noindent\textcolor{FuncColor}{$\triangleright$\enspace\texttt{ComputedIndicators({\mdseries\slshape tbl})\index{ComputedIndicators@\texttt{ComputedIndicators}}
\label{ComputedIndicators}
}\hfill{\scriptsize (attribute)}}\\


 If \mbox{\texttt{\mdseries\slshape tbl}} is an ordinary character table then \texttt{Indicator} returns the list of \mbox{\texttt{\mdseries\slshape n}}\texttt{\symbol{45}}th Frobenius\texttt{\symbol{45}}Schur indicators of the
characters in the list \mbox{\texttt{\mdseries\slshape characters}}; the default of \mbox{\texttt{\mdseries\slshape characters}} is \texttt{Irr( \mbox{\texttt{\mdseries\slshape tbl}} )}. 

 The $n$\texttt{\symbol{45}}th Frobenius\texttt{\symbol{45}}Schur indicator $\nu_n(\chi)$ of an ordinary character $\chi$ of the group $G$ is given by $\nu_n(\chi) = ( \sum_{{g \in G}} \chi(g^n) ) / |G|$. 

 If \mbox{\texttt{\mdseries\slshape tbl}} is a Brauer table in characteristic $ \neq 2$ and $\mbox{\texttt{\mdseries\slshape n}} = 2$ then \texttt{Indicator} returns the second indicator. 

 The default method uses the attribute \texttt{ComputedIndicators} for storing the computed value at position \mbox{\texttt{\mdseries\slshape n}}, and calls the operation \texttt{IndicatorOp} for computing values that are not yet known. 

 
\begin{Verbatim}[commandchars=!@|,fontsize=\small,frame=single,label=Example]
  !gapprompt@gap>| !gapinput@tbl:= CharacterTable( "L3(2)" );;|
  !gapprompt@gap>| !gapinput@Indicator( tbl, 2 );|
  [ 1, 0, 0, 1, 1, 1 ]
\end{Verbatim}
 

 In nonzero characteristic $p$, the Frobenius\texttt{\symbol{45}}Schur indicator is defined only for
irreducible characters. For odd $p$, the indicator is computed using the Thompson\texttt{\symbol{45}}Willems
Theorem \cite[theorem on p. 227]{Tho86}. For $p = 2$, in general the indicator cannot be computed from the given character tables,
here the following necessary conditions are used. 

 
\begin{itemize}
\item  The trivial character has indicator $1$. 
\item  The indicator is $0$ if and only if the character is not real\texttt{\symbol{45}}valued. 
\item  Real characters outside the principal block (the $2$\texttt{\symbol{45}}block that contains the trivial character, see \texttt{PrimeBlocks} (\ref{PrimeBlocks})) have indicator $1$. 
\item  By \cite[Lemma 1.2]{GW95}, any real constituent with odd multiplicity in the $2$\texttt{\symbol{45}}modular restriction of an ordinary irreducible character
with indicator $1$ has indicator $1$, provided that the trivial character is not a constituent of the restriction. 
\end{itemize}
 

 For each $2$\texttt{\symbol{45}}modular Brauer characters where these conditions are not
sufficient to determine the indicator, an unknown value (see \texttt{Unknown} (\ref{Unknown})) is returned. }

 

\subsection{\textcolor{Chapter }{NrPolyhedralSubgroups}}
\logpage{[ 71, 12, 6 ]}\nobreak
\hyperdef{L}{X83AE05BF8085B3C8}{}
{\noindent\textcolor{FuncColor}{$\triangleright$\enspace\texttt{NrPolyhedralSubgroups({\mdseries\slshape tbl, c1, c2, c3})\index{NrPolyhedralSubgroups@\texttt{NrPolyhedralSubgroups}}
\label{NrPolyhedralSubgroups}
}\hfill{\scriptsize (function)}}\\


 \index{subgroups!polyhedral} returns the number and isomorphism type of polyhedral subgroups of the group
with ordinary character table \mbox{\texttt{\mdseries\slshape tbl}} which are generated by an element $g$ of class \mbox{\texttt{\mdseries\slshape c1}} and an element $h$ of class \mbox{\texttt{\mdseries\slshape c2}} with the property that the product $gh$ lies in class \mbox{\texttt{\mdseries\slshape c3}}. 

 According to \cite[p. 233]{NPP84}, the number of polyhedral subgroups of isomorphism type $V_4$, $D_{2n}$, $A_4$, $S_4$, and $A_5$ can be derived from the class multiplication coefficient (see{\nobreakspace}\texttt{ClassMultiplicationCoefficient} (\ref{ClassMultiplicationCoefficient:for character tables})) and the number of Galois conjugates of a class (see{\nobreakspace}\texttt{ClassOrbit} (\ref{ClassOrbit})). 

 The classes \mbox{\texttt{\mdseries\slshape c1}}, \mbox{\texttt{\mdseries\slshape c2}} and \mbox{\texttt{\mdseries\slshape c3}} in the parameter list must be ordered according to the order of the elements
in these classes. If elements in class \mbox{\texttt{\mdseries\slshape c1}} and \mbox{\texttt{\mdseries\slshape c2}} do not generate a polyhedral group then \texttt{fail} is returned. 

 
\begin{Verbatim}[commandchars=!@|,fontsize=\small,frame=single,label=Example]
  !gapprompt@gap>| !gapinput@NrPolyhedralSubgroups( tbl, 2, 2, 4 );|
  rec( number := 21, type := "D8" )
\end{Verbatim}
 }

 

\subsection{\textcolor{Chapter }{ClassMultiplicationCoefficient (for character tables)}}
\logpage{[ 71, 12, 7 ]}\nobreak
\hyperdef{L}{X7E2EA9FE7D3062D3}{}
{\noindent\textcolor{FuncColor}{$\triangleright$\enspace\texttt{ClassMultiplicationCoefficient({\mdseries\slshape tbl, i, j, k})\index{ClassMultiplicationCoefficient@\texttt{ClassMultiplicationCoefficient}!for character tables}
\label{ClassMultiplicationCoefficient:for character tables}
}\hfill{\scriptsize (operation)}}\\


 \index{ClassMultiplicationCoefficient@\texttt{ClassMultiplicationCoefficient}!for character tables} \index{class multiplication coefficient} \index{structure constant} returns the class multiplication coefficient of the classes \mbox{\texttt{\mdseries\slshape i}}, \mbox{\texttt{\mdseries\slshape j}}, and \mbox{\texttt{\mdseries\slshape k}} of the group $G$ with ordinary character table \mbox{\texttt{\mdseries\slshape tbl}}. 

 The class multiplication coefficient $c_{{i,j,k}}$ of the classes \mbox{\texttt{\mdseries\slshape i}}, \mbox{\texttt{\mdseries\slshape j}}, \mbox{\texttt{\mdseries\slshape k}} equals the number of pairs $(x,y)$ of elements $x, y \in G$ such that $x$ lies in class \mbox{\texttt{\mdseries\slshape i}}, $y$ lies in class \mbox{\texttt{\mdseries\slshape j}}, and their product $xy$ is a fixed element of class \mbox{\texttt{\mdseries\slshape k}}. 

 In the center of the group algebra of $G$, these numbers are found as coefficients of the decomposition of the product
of two class sums $K_i$ and $K_j$ into class sums: 
\[ K_i K_j = \sum_k c_{ijk} K_k . \]
 Given the character table of a finite group $G$, whose classes are $C_1, \ldots, C_r$ with representatives $g_i \in C_i$, the class multiplication coefficient $c_{ijk}$ can be computed with the following formula: 
\[ c_{ijk} = |C_i| \cdot |C_j| / |G| \cdot \sum_{{\chi \in Irr(G)}} \chi(g_i)
\chi(g_j) \chi(g_k^{{-1}}) / \chi(1). \]
 

 On the other hand the knowledge of the class multiplication coefficients
admits the computation of the irreducible characters of $G$, see{\nobreakspace}\texttt{IrrDixonSchneider} (\ref{IrrDixonSchneider}). }

 

\subsection{\textcolor{Chapter }{ClassStructureCharTable}}
\logpage{[ 71, 12, 8 ]}\nobreak
\hyperdef{L}{X7A19F56C7FD5EFC7}{}
{\noindent\textcolor{FuncColor}{$\triangleright$\enspace\texttt{ClassStructureCharTable({\mdseries\slshape tbl, classes})\index{ClassStructureCharTable@\texttt{ClassStructureCharTable}}
\label{ClassStructureCharTable}
}\hfill{\scriptsize (function)}}\\


 \index{class multiplication coefficient} \index{structure constant} 

 returns the so\texttt{\symbol{45}}called class structure of the classes in the
list \mbox{\texttt{\mdseries\slshape classes}}, for the character table \mbox{\texttt{\mdseries\slshape tbl}} of the group $G$. The length of \mbox{\texttt{\mdseries\slshape classes}} must be at least 2. 

 Let $C = (C_1, C_2, \ldots, C_n)$ denote the $n$\texttt{\symbol{45}}tuple of conjugacy classes of $G$ that are indexed by \mbox{\texttt{\mdseries\slshape classes}}. The class structure $n(C)$ equals the number of $n$\texttt{\symbol{45}}tuples $(g_1, g_2, \ldots, g_n)$ of elements $g_i \in C_i$ with $g_1 g_2 \cdots g_n = 1$. Note the difference to the definition of the class multiplication
coefficients in \texttt{ClassMultiplicationCoefficient} (\ref{ClassMultiplicationCoefficient:for character tables}). 

 $n(C_1, C_2, \ldots, C_n)$ is computed using the formula 
\[ n(C_1, C_2, \ldots, C_n) = |C_1| |C_2| \cdots |C_n| / |G| \cdot \sum_{{\chi
\in Irr(G)}} \chi(g_1) \chi(g_2) \cdots \chi(g_n) / \chi(1)^{{n-2}} . \]
 }

 

\subsection{\textcolor{Chapter }{MatClassMultCoeffsCharTable}}
\logpage{[ 71, 12, 9 ]}\nobreak
\hyperdef{L}{X809E67E57D4933B3}{}
{\noindent\textcolor{FuncColor}{$\triangleright$\enspace\texttt{MatClassMultCoeffsCharTable({\mdseries\slshape tbl, i})\index{MatClassMultCoeffsCharTable@\texttt{MatClassMultCoeffsCharTable}}
\label{MatClassMultCoeffsCharTable}
}\hfill{\scriptsize (function)}}\\


 \index{structure constant} \index{class multiplication coefficient} 

 For an ordinary character table \mbox{\texttt{\mdseries\slshape tbl}} and a class position \mbox{\texttt{\mdseries\slshape i}}, \texttt{MatClassMultCoeffsCharTable} returns the matrix $[ a_{ijk} ]_{{j,k}}$ of structure constants (see{\nobreakspace}\texttt{ClassMultiplicationCoefficient} (\ref{ClassMultiplicationCoefficient:for character tables})). 

 
\begin{Verbatim}[commandchars=!@|,fontsize=\small,frame=single,label=Example]
  !gapprompt@gap>| !gapinput@tbl:= CharacterTable( "L3(2)" );;|
  !gapprompt@gap>| !gapinput@ClassMultiplicationCoefficient( tbl, 2, 2, 4 );|
  4
  !gapprompt@gap>| !gapinput@ClassStructureCharTable( tbl, [ 2, 2, 4 ] );|
  168
  !gapprompt@gap>| !gapinput@ClassStructureCharTable( tbl, [ 2, 2, 2, 4 ] );|
  1848
  !gapprompt@gap>| !gapinput@MatClassMultCoeffsCharTable( tbl, 2 );|
  [ [ 0, 1, 0, 0, 0, 0 ], [ 21, 4, 3, 4, 0, 0 ], [ 0, 8, 6, 8, 7, 7 ],
    [ 0, 8, 6, 1, 7, 7 ], [ 0, 0, 3, 4, 0, 7 ], [ 0, 0, 3, 4, 7, 0 ] ]
\end{Verbatim}
 }

 }

  
\section{\textcolor{Chapter }{Printing Character Tables}}\label{Printing Character Tables}
\logpage{[ 71, 13, 0 ]}
\hyperdef{L}{X7C1941F17BE9FC21}{}
{
  

\subsection{\textcolor{Chapter }{ViewObj (for a character table)}}
\logpage{[ 71, 13, 1 ]}\nobreak
\hyperdef{L}{X7D45224B86D802E5}{}
{\noindent\textcolor{FuncColor}{$\triangleright$\enspace\texttt{ViewObj({\mdseries\slshape tbl})\index{ViewObj@\texttt{ViewObj}!for a character table}
\label{ViewObj:for a character table}
}\hfill{\scriptsize (method)}}\\


 The default \texttt{ViewObj} (\ref{ViewObj}) method for ordinary character tables prints the string \texttt{"CharacterTable"}, followed by the identifier (see{\nobreakspace}\texttt{Identifier} (\ref{Identifier:for character tables})) or, if known, the group of the character table enclosed in brackets. \texttt{ViewObj} (\ref{ViewObj}) for Brauer tables does the same, except that the first string is replaced by \texttt{"BrauerTable"}, and that the characteristic is also shown. }

 

\subsection{\textcolor{Chapter }{PrintObj (for a character table)}}
\logpage{[ 71, 13, 2 ]}\nobreak
\hyperdef{L}{X836554207C678D41}{}
{\noindent\textcolor{FuncColor}{$\triangleright$\enspace\texttt{PrintObj({\mdseries\slshape tbl})\index{PrintObj@\texttt{PrintObj}!for a character table}
\label{PrintObj:for a character table}
}\hfill{\scriptsize (method)}}\\


 The default \texttt{PrintObj} (\ref{PrintObj}) method for character tables does the same as \texttt{ViewObj} (\ref{ViewObj}), except that \texttt{PrintObj} (\ref{PrintObj}) is used for the group instead of \texttt{ViewObj} (\ref{ViewObj}). }

 

\subsection{\textcolor{Chapter }{Display (for a character table)}}
\logpage{[ 71, 13, 3 ]}\nobreak
\hyperdef{L}{X7B41F36478C47364}{}
{\noindent\textcolor{FuncColor}{$\triangleright$\enspace\texttt{Display({\mdseries\slshape tbl})\index{Display@\texttt{Display}!for a character table}
\label{Display:for a character table}
}\hfill{\scriptsize (method)}}\\


 There are various ways to customize the \texttt{Display} (\ref{Display}) output for character tables. First we describe the default behaviour,
alternatives are then described below. 

 The default \texttt{Display} (\ref{Display}) method prepares the data in \mbox{\texttt{\mdseries\slshape tbl}} for a columnwise output. The number of columns printed at one time depends on
the actual line length, which can be accessed and changed by the function \texttt{SizeScreen} (\ref{SizeScreen}). 

 An interesting variant of \texttt{Display} (\ref{Display}) is the function \texttt{PageDisplay} (\textbf{GAPDoc: PageDisplay}). Convenient ways to print the \texttt{Display} (\ref{Display}) format to a file are given by the function \texttt{PrintTo1} (\textbf{GAPDoc: PrintTo1}) or by using \texttt{PageDisplay} (\textbf{GAPDoc: PageDisplay}) and the facilities of the pager used, cf.{\nobreakspace}\texttt{Pager} (\ref{Pager}). 

 An interactive variant of \texttt{Display} (\ref{Display}) is the \texttt{Browse} (\textbf{Browse: Browse}) method for character tables that is provided by the \textsf{GAP} package \textsf{Browse}, see \texttt{Browse} (\textbf{Browse: Browse for character tables}). 

 \texttt{Display} (\ref{Display}) shows certain characters (by default all irreducible characters) of \mbox{\texttt{\mdseries\slshape tbl}}, together with the orders of the centralizers in factorized form and the
available power maps (see{\nobreakspace}\texttt{ComputedPowerMaps} (\ref{ComputedPowerMaps})). The \mbox{\texttt{\mdseries\slshape n}}\texttt{\symbol{45}}th displayed character is given the name \texttt{X.\mbox{\texttt{\mdseries\slshape n}}}. 

 The first lines of the output describe the order of the centralizer of an
element of the class factorized into its prime divisors. 

 The next line gives the name of each class. If no class names are stored on \mbox{\texttt{\mdseries\slshape tbl}}, \texttt{ClassNames} (\ref{ClassNames}) is called. 

 Preceded by a name \texttt{P\mbox{\texttt{\mdseries\slshape n}}}, the next lines show the \mbox{\texttt{\mdseries\slshape n}}th power maps of \mbox{\texttt{\mdseries\slshape tbl}} in terms of the former shown class names. 

 Every ambiguous or unknown (see Chapter{\nobreakspace}\ref{Unknowns}) value of the table is displayed as a question mark \texttt{?}. 

 Irrational character values are not printed explicitly because the lengths of
their printed representation might disturb the layout. Instead of that every
irrational value is indicated by a name, which is a string of at least one
capital letter. 

 Once a name for an irrational value is found, it is used all over the printed
table. Moreover the complex conjugate (see{\nobreakspace}\texttt{ComplexConjugate} (\ref{ComplexConjugate}), \texttt{GaloisCyc} (\ref{GaloisCyc:for a cyclotomic})) and the star of an irrationality (see{\nobreakspace}\texttt{StarCyc} (\ref{StarCyc})) are represented by that very name preceded by a \texttt{/} and a \texttt{*}, respectively. 

 The printed character table is then followed by a legend, a list identifying
the occurring symbols with their actual values. Occasionally this
identification is supplemented by a quadratic representation of the
irrationality (see{\nobreakspace}\texttt{Quadratic} (\ref{Quadratic})) together with the corresponding \textsf{Atlas} notation (see{\nobreakspace}\cite{CCN85}). 

 This default style can be changed by prescribing a record \mbox{\texttt{\mdseries\slshape arec}} of options, which can be given 

 
\begin{enumerate}
\item  as an optional argument in the call to \texttt{Display} (\ref{Display}), 
\item  as the value of the attribute \texttt{DisplayOptions} (\ref{DisplayOptions}) if this value is stored in the table, 
\item  as the value of the global variable \texttt{CharacterTableDisplayDefaults.User}, or 
\item  as the value of the global variable \texttt{CharacterTableDisplayDefaults.Global} 
\end{enumerate}
 

 (in this order of precedence). 

 The following components of \mbox{\texttt{\mdseries\slshape arec}} are supported. 

 
\begin{description}
\item[{\texttt{centralizers}}]  \texttt{false} to suppress the printing of the orders of the centralizers, or the string \texttt{"ATLAS"} to force the printing of non\texttt{\symbol{45}}factorized centralizer orders
in a style similar to that used in the \textsf{Atlas} of Finite Groups{\nobreakspace}\cite{CCN85}, 
\item[{\texttt{characterField}}]  \texttt{true} to show the degrees of the character fields over the prime field, in a column
with header \texttt{d}, 
\item[{\texttt{chars}}]  an integer or a list of integers to select a sublist of the irreducible
characters of \mbox{\texttt{\mdseries\slshape tbl}}, or a list of characters of \mbox{\texttt{\mdseries\slshape tbl}} (in the latter case, the default letter \texttt{"X"} in the character names is replaced by \texttt{"Y"}), 
\item[{\texttt{charnames}}]  a list of strings of length equal to the number of characters that shall be
shown; they are used as labels for the characters, 
\item[{\texttt{classes}}]  an integer or a list of integers to select a sublist of the classes of \mbox{\texttt{\mdseries\slshape tbl}}, 
\item[{\texttt{classnames}}]  a list of strings of length equal to the number of classes that shall be
shown; they are used as labels for the classes, 
\item[{\texttt{indicator}}]  \texttt{true} enables the printing of the second Frobenius Schur indicator, a list of
integers enables the printing of the corresponding indicators
(see{\nobreakspace}\texttt{Indicator} (\ref{Indicator})), 
\item[{\texttt{letter}}]  a single capital letter (e.{\nobreakspace}g.{\nobreakspace}\texttt{"P"} for permutation characters) to replace the default \texttt{"X"} in character names, 
\item[{\texttt{powermap}}]  an integer or a list of integers to select a subset of the available power
maps, \texttt{false} to suppress the printing of power maps, or the string \texttt{"ATLAS"} to force a printing of class names and power maps in a style similar to that
used in the \textsf{Atlas} of Finite Groups{\nobreakspace}\cite{CCN85} (the \texttt{"ATLAS"} variant works only if the function \texttt{CambridgeMaps} (\textbf{CTblLib: CambridgeMaps}) is available, which belongs to the \textsf{CTblLib} package), 
\item[{\texttt{Display}}]  the function that is actually called in order to display the table; the
arguments are the table and the optional record, whose components can be used
inside the \texttt{Display} function, 
\item[{\texttt{StringEntry}}]  a function that takes either a character value or a character value and the
return value of \texttt{StringEntryData} (see below), and returns the string that is actually displayed; it is called
for all character values to be displayed, and also for the displayed indicator
values (see above), 
\item[{\texttt{StringEntryData}}]  a unary function that is called once with argument \mbox{\texttt{\mdseries\slshape tbl}} before the character values are displayed; it returns an object that is used
as second argument of the function \texttt{StringEntry}, 
\item[{\texttt{Legend}}]  a function that takes the result of the \texttt{StringEntryData} call as its only argument, after the character table has been displayed; the
return value is a string that describes the symbols used in the displayed
table in a formatted way, it is printed below the displayed table. 
\end{description}
 }

 

\subsection{\textcolor{Chapter }{DisplayOptions}}
\logpage{[ 71, 13, 4 ]}\nobreak
\hyperdef{L}{X85E883A87A190AA4}{}
{\noindent\textcolor{FuncColor}{$\triangleright$\enspace\texttt{DisplayOptions({\mdseries\slshape tbl})\index{DisplayOptions@\texttt{DisplayOptions}}
\label{DisplayOptions}
}\hfill{\scriptsize (attribute)}}\\


  There is no default method to compute a value, one can set a value with \texttt{SetDisplayOptions}. 

 
\begin{Verbatim}[commandchars=!@|,fontsize=\small,frame=single,label=Example]
  !gapprompt@gap>| !gapinput@tbl:= CharacterTable( "A5" );;|
  !gapprompt@gap>| !gapinput@Display( tbl );|
  A5
  
       2  2  2  .  .  .
       3  1  .  1  .  .
       5  1  .  .  1  1
  
         1a 2a 3a 5a 5b
      2P 1a 1a 3a 5b 5a
      3P 1a 2a 1a 5b 5a
      5P 1a 2a 3a 1a 1a
  
  X.1     1  1  1  1  1
  X.2     3 -1  .  A *A
  X.3     3 -1  . *A  A
  X.4     4  .  1 -1 -1
  X.5     5  1 -1  .  .
  
  A = -E(5)-E(5)^4
    = (1-Sqrt(5))/2 = -b5
  !gapprompt@gap>| !gapinput@CharacterTableDisplayDefaults.User:= rec(|
  !gapprompt@>| !gapinput@       powermap:= "ATLAS", centralizers:= "ATLAS", chars:= false );;|
  !gapprompt@gap>| !gapinput@Display( CharacterTable( "A5" ) );|
  A5
  
      60  4  3  5  5
  
   p      A  A  A  A
   p'     A  A  A  A
      1A 2A 3A 5A B*
  
  !gapprompt@gap>| !gapinput@options:= rec( chars:= 4, classes:= [ tbl.3a .. tbl.5a ],|
  !gapprompt@>| !gapinput@                  centralizers:= false, indicator:= true,|
  !gapprompt@>| !gapinput@                  powermap:= [ 2 ] );;|
  !gapprompt@gap>| !gapinput@Display( tbl, options );|
  A5
  
            3a 5a
         2P 3a 5b
         2
  X.4    +   1 -1
  !gapprompt@gap>| !gapinput@SetDisplayOptions( tbl, options );  Display( tbl );|
  A5
  
            3a 5a
         2P 3a 5b
         2
  X.4    +   1 -1
  !gapprompt@gap>| !gapinput@Unbind( CharacterTableDisplayDefaults.User );|
\end{Verbatim}
 }

 

\subsection{\textcolor{Chapter }{PrintCharacterTable}}
\logpage{[ 71, 13, 5 ]}\nobreak
\hyperdef{L}{X79EC9603833AA2AB}{}
{\noindent\textcolor{FuncColor}{$\triangleright$\enspace\texttt{PrintCharacterTable({\mdseries\slshape tbl, varname})\index{PrintCharacterTable@\texttt{PrintCharacterTable}}
\label{PrintCharacterTable}
}\hfill{\scriptsize (function)}}\\


 Let \mbox{\texttt{\mdseries\slshape tbl}} be a nearly character table, and \mbox{\texttt{\mdseries\slshape varname}} a string. \texttt{PrintCharacterTable} prints those values of the supported attributes (see{\nobreakspace}\texttt{SupportedCharacterTableInfo} (\ref{SupportedCharacterTableInfo})) that are known for \mbox{\texttt{\mdseries\slshape tbl}}.  

 The output of \texttt{PrintCharacterTable} is \textsf{GAP} readable; actually reading it into \textsf{GAP} will bind the variable with name \mbox{\texttt{\mdseries\slshape varname}} to a character table that coincides with \mbox{\texttt{\mdseries\slshape tbl}} for all printed components. 

 This is used mainly for saving character tables to files. A more human
readable form is produced by \texttt{Display} (\ref{Display}).   

 
\begin{Verbatim}[commandchars=!@|,fontsize=\small,frame=single,label=Example]
  !gapprompt@gap>| !gapinput@PrintCharacterTable( CharacterTable( "Cyclic", 2 ), "tbl" );|
  tbl:= function()
  local tbl, i;
  tbl:=rec();
  tbl.Irr:=
  [ [ 1, 1 ], [ 1, -1 ] ];
  tbl.IsFinite:=
  true;
  tbl.NrConjugacyClasses:=
  2;
  tbl.Size:=
  2;
  tbl.OrdersClassRepresentatives:=
  [ 1, 2 ];
  tbl.SizesCentralizers:=
  [ 2, 2 ];
  tbl.UnderlyingCharacteristic:=
  0;
  tbl.ClassParameters:=
  [ [ 1, 0 ], [ 1, 1 ] ];
  tbl.CharacterParameters:=
  [ [ 1, 0 ], [ 1, 1 ] ];
  tbl.Identifier:=
  "C2";
  tbl.InfoText:=
  "computed using generic character table for cyclic groups";
  tbl.ComputedPowerMaps:=
  [ , [ 1, 1 ] ];
  ConvertToLibraryCharacterTableNC(tbl);
  return tbl;
  end;
  tbl:= tbl();
\end{Verbatim}
 }

 }

  
\section{\textcolor{Chapter }{Computing the Irreducible Characters of a Group}}\label{Computing the Irreducible Characters of a Group}
\logpage{[ 71, 14, 0 ]}
\hyperdef{L}{X79BC08C6846718D9}{}
{
  Several algorithms are available for computing the irreducible characters of a
finite group $G$. The default method for arbitrary finite groups is to use the
Dixon\texttt{\symbol{45}}Schneider algorithm (see{\nobreakspace}\texttt{IrrDixonSchneider} (\ref{IrrDixonSchneider})). For supersolvable groups, Conlon's algorithm can be used
(see{\nobreakspace}\texttt{IrrConlon} (\ref{IrrConlon})). For abelian\texttt{\symbol{45}}by\texttt{\symbol{45}}supersolvable groups,
the Baum\texttt{\symbol{45}}Clausen algorithm for computing the irreducible
representations (see{\nobreakspace}\texttt{IrreducibleRepresentations} (\ref{IrreducibleRepresentations})) can be used to compute the irreducible characters (see{\nobreakspace}\texttt{IrrBaumClausen} (\ref{IrrBaumClausen})). 

 These functions are installed in methods for \texttt{Irr} (\ref{Irr:for a group}), but explicitly calling one of them will \emph{not} set the \texttt{Irr} (\ref{Irr:for a group}) value of $G$. 

\subsection{\textcolor{Chapter }{IrrDixonSchneider}}
\logpage{[ 71, 14, 1 ]}\nobreak
\hyperdef{L}{X7ED39DB680BFEA96}{}
{\noindent\textcolor{FuncColor}{$\triangleright$\enspace\texttt{IrrDixonSchneider({\mdseries\slshape G})\index{IrrDixonSchneider@\texttt{IrrDixonSchneider}}
\label{IrrDixonSchneider}
}\hfill{\scriptsize (attribute)}}\\


 computes the irreducible characters of the finite group \mbox{\texttt{\mdseries\slshape G}}, using the Dixon\texttt{\symbol{45}}Schneider method (see{\nobreakspace}\ref{The Dixon-Schneider Algorithm}). It calls \texttt{DixonInit} (\ref{DixonInit}) and \texttt{DixonSplit} (\ref{DixonSplit}),  and finally returns the list returned by \texttt{DixontinI} (\ref{DixontinI}). See also the sections \ref{Components of a Dixon Record} and \ref{An Example of Advanced Dixon-Schneider Calculations}. }

 

\subsection{\textcolor{Chapter }{IrrConlon}}
\logpage{[ 71, 14, 2 ]}\nobreak
\hyperdef{L}{X7E81BCD686561DF0}{}
{\noindent\textcolor{FuncColor}{$\triangleright$\enspace\texttt{IrrConlon({\mdseries\slshape G})\index{IrrConlon@\texttt{IrrConlon}}
\label{IrrConlon}
}\hfill{\scriptsize (attribute)}}\\


 For a finite solvable group \mbox{\texttt{\mdseries\slshape G}}, \texttt{IrrConlon} returns a list of monomial irreducible characters of \mbox{\texttt{\mdseries\slshape G}}, among those all irreducibles that have the supersolvable residuum of \mbox{\texttt{\mdseries\slshape G}} in their kernels; so if \mbox{\texttt{\mdseries\slshape G}} is supersolvable, all irreducible characters of \mbox{\texttt{\mdseries\slshape G}} are returned. An error is signalled if \mbox{\texttt{\mdseries\slshape G}} is not solvable. 

 The characters are computed using Conlon's algorithm (see{\nobreakspace}\cite{Con90a} and{\nobreakspace}\cite{Con90b}). For each irreducible character in the returned list, the monomiality
information (see{\nobreakspace}\texttt{TestMonomial} (\ref{TestMonomial:for a group})) is stored. }

 

\subsection{\textcolor{Chapter }{IrrBaumClausen}}
\logpage{[ 71, 14, 3 ]}\nobreak
\hyperdef{L}{X7BF15729839203FC}{}
{\noindent\textcolor{FuncColor}{$\triangleright$\enspace\texttt{IrrBaumClausen({\mdseries\slshape G})\index{IrrBaumClausen@\texttt{IrrBaumClausen}}
\label{IrrBaumClausen}
}\hfill{\scriptsize (attribute)}}\\


 \texttt{IrrBaumClausen} returns the absolutely irreducible ordinary characters of the factor group of
the finite solvable group \mbox{\texttt{\mdseries\slshape G}} by the derived subgroup of its supersolvable residuum. 

 The characters are computed using the algorithm by Baum and Clausen
(see{\nobreakspace}\cite{BC94}). An error is signalled if \mbox{\texttt{\mdseries\slshape G}} is not solvable. 

 
\begin{Verbatim}[commandchars=!@|,fontsize=\small,frame=single,label=Example]
  !gapprompt@gap>| !gapinput@g:= SL(2,3);;|
  !gapprompt@gap>| !gapinput@irr1:= IrrDixonSchneider( g );|
  [ Character( CharacterTable( SL(2,3) ), [ 1, 1, 1, 1, 1, 1, 1 ] ),
    Character( CharacterTable( SL(2,3) ),
      [ 1, E(3)^2, E(3), 1, E(3), E(3)^2, 1 ] ),
    Character( CharacterTable( SL(2,3) ),
      [ 1, E(3), E(3)^2, 1, E(3)^2, E(3), 1 ] ),
    Character( CharacterTable( SL(2,3) ), [ 2, 1, 1, -2, -1, -1, 0 ] ),
    Character( CharacterTable( SL(2,3) ),
      [ 2, E(3)^2, E(3), -2, -E(3), -E(3)^2, 0 ] ),
    Character( CharacterTable( SL(2,3) ),
      [ 2, E(3), E(3)^2, -2, -E(3)^2, -E(3), 0 ] ),
    Character( CharacterTable( SL(2,3) ), [ 3, 0, 0, 3, 0, 0, -1 ] ) ]
  !gapprompt@gap>| !gapinput@irr2:= IrrConlon( g );|
  [ Character( CharacterTable( SL(2,3) ), [ 1, 1, 1, 1, 1, 1, 1 ] ),
    Character( CharacterTable( SL(2,3) ),
      [ 1, E(3), E(3)^2, 1, E(3)^2, E(3), 1 ] ),
    Character( CharacterTable( SL(2,3) ),
      [ 1, E(3)^2, E(3), 1, E(3), E(3)^2, 1 ] ),
    Character( CharacterTable( SL(2,3) ), [ 3, 0, 0, 3, 0, 0, -1 ] ) ]
  !gapprompt@gap>| !gapinput@irr3:= IrrBaumClausen( g );|
  [ Character( CharacterTable( SL(2,3) ), [ 1, 1, 1, 1, 1, 1, 1 ] ),
    Character( CharacterTable( SL(2,3) ),
      [ 1, E(3), E(3)^2, 1, E(3)^2, E(3), 1 ] ),
    Character( CharacterTable( SL(2,3) ),
      [ 1, E(3)^2, E(3), 1, E(3), E(3)^2, 1 ] ),
    Character( CharacterTable( SL(2,3) ), [ 3, 0, 0, 3, 0, 0, -1 ] ) ]
  !gapprompt@gap>| !gapinput@chi:= irr2[4];;  HasTestMonomial( chi );|
  true
\end{Verbatim}
 }

 

\subsection{\textcolor{Chapter }{IrreducibleRepresentations}}
\logpage{[ 71, 14, 4 ]}\nobreak
\hyperdef{L}{X7F29C5447B5DC102}{}
{\noindent\textcolor{FuncColor}{$\triangleright$\enspace\texttt{IrreducibleRepresentations({\mdseries\slshape G[, F]})\index{IrreducibleRepresentations@\texttt{IrreducibleRepresentations}}
\label{IrreducibleRepresentations}
}\hfill{\scriptsize (attribute)}}\\


 Called with a finite group \mbox{\texttt{\mdseries\slshape G}} and a field \mbox{\texttt{\mdseries\slshape F}}, \texttt{IrreducibleRepresentations} returns a list of representatives of the irreducible matrix representations of \mbox{\texttt{\mdseries\slshape G}} over \mbox{\texttt{\mdseries\slshape F}}, up to equivalence. 

 If \mbox{\texttt{\mdseries\slshape G}} is the only argument then \texttt{IrreducibleRepresentations} returns a list of representatives of the absolutely irreducible complex
representations of \mbox{\texttt{\mdseries\slshape G}}, up to equivalence. 

 At the moment, methods are available for the following cases: If \mbox{\texttt{\mdseries\slshape G}} is abelian by supersolvable the method of{\nobreakspace}\cite{BC94} is used. 

 Otherwise, if \mbox{\texttt{\mdseries\slshape F}} and \mbox{\texttt{\mdseries\slshape G}} are both finite, the regular module of \mbox{\texttt{\mdseries\slshape G}} is split by MeatAxe methods which can make this an expensive operation. 

 Finally, if \mbox{\texttt{\mdseries\slshape F}} is not given (i.e. it defaults to the cyclotomic numbers) and \mbox{\texttt{\mdseries\slshape G}} is a finite group, the method of \cite{Dix93} (see \texttt{IrreducibleRepresentationsDixon} (\ref{IrreducibleRepresentationsDixon})) is used. 

 For other cases no methods are implemented yet. 

 The representations obtained are \emph{not} guaranteed to be ``nice'' (for example preserving a unitary form) in any way. 

 See also \texttt{IrreducibleModules} (\ref{IrreducibleModules}), which provides efficient methods for solvable groups. 

 
\begin{Verbatim}[commandchars=!@|,fontsize=\small,frame=single,label=Example]
  !gapprompt@gap>| !gapinput@g:= AlternatingGroup( 4 );;|
  !gapprompt@gap>| !gapinput@repr:= IrreducibleRepresentations( g );|
  [ Pcgs([ (2,4,3), (1,3)(2,4), (1,2)(3,4) ]) ->
      [ [ [ 1 ] ], [ [ 1 ] ], [ [ 1 ] ] ],
    Pcgs([ (2,4,3), (1,3)(2,4), (1,2)(3,4) ]) ->
      [ [ [ E(3) ] ], [ [ 1 ] ], [ [ 1 ] ] ],
    Pcgs([ (2,4,3), (1,3)(2,4), (1,2)(3,4) ]) ->
      [ [ [ E(3)^2 ] ], [ [ 1 ] ], [ [ 1 ] ] ],
    Pcgs([ (2,4,3), (1,3)(2,4), (1,2)(3,4) ]) ->
      [ [ [ 0, 0, 1 ], [ 1, 0, 0 ], [ 0, 1, 0 ] ],
        [ [ -1, 0, 0 ], [ 0, 1, 0 ], [ 0, 0, -1 ] ],
        [ [ 1, 0, 0 ], [ 0, -1, 0 ], [ 0, 0, -1 ] ] ] ]
  !gapprompt@gap>| !gapinput@ForAll( repr, IsGroupHomomorphism );|
  true
  !gapprompt@gap>| !gapinput@Length( repr );|
  4
  !gapprompt@gap>| !gapinput@gens:= GeneratorsOfGroup( g );|
  [ (1,2,3), (2,3,4) ]
  !gapprompt@gap>| !gapinput@List( gens, x -> x^repr[1] );|
  [ [ [ 1 ] ], [ [ 1 ] ] ]
  !gapprompt@gap>| !gapinput@ List( gens, x -> x^repr[4] );|
  [ [ [ 0, 0, -1 ], [ 1, 0, 0 ], [ 0, -1, 0 ] ],
    [ [ 0, 1, 0 ], [ 0, 0, 1 ], [ 1, 0, 0 ] ] ]
\end{Verbatim}
 }

 

\subsection{\textcolor{Chapter }{IrreducibleRepresentationsDixon}}
\logpage{[ 71, 14, 5 ]}\nobreak
\hyperdef{L}{X8493ED7A86FFCB8A}{}
{\noindent\textcolor{FuncColor}{$\triangleright$\enspace\texttt{IrreducibleRepresentationsDixon({\mdseries\slshape G[, chi]})\index{IrreducibleRepresentationsDixon@\texttt{IrreducibleRepresentationsDixon}}
\label{IrreducibleRepresentationsDixon}
}\hfill{\scriptsize (function)}}\\


 Called with one argument, a group \mbox{\texttt{\mdseries\slshape G}}, \texttt{IrreducibleRepresentationsDixon} computes (representatives of) all irreducible complex representations for the
finite group \mbox{\texttt{\mdseries\slshape G}}, using the method of{\nobreakspace}\cite{Dix93}, which computes the character table and computes the representation as
constituent of an induced monomial representation of a subgroup. 

 This method can be quite expensive for larger groups, for example it might
involve calculation of the subgroup lattice of \mbox{\texttt{\mdseries\slshape G}}. 

 A character \mbox{\texttt{\mdseries\slshape chi}} of \mbox{\texttt{\mdseries\slshape G}} can be given as the second argument, in this case only a representation
affording \mbox{\texttt{\mdseries\slshape chi}} is returned. 

 The second argument can also be a list of characters of \mbox{\texttt{\mdseries\slshape G}}, in this case only representations for characters in this list are computed. 

 Note that this method might fail if for an irreducible representation there is
no subgroup in which its reduction has a linear constituent with multiplicity
one. 

 If the option \mbox{\texttt{\mdseries\slshape unitary}} is given, \textsf{GAP} tries, at extra cost, to find a unitary representation (and will issue an
error if it cannot do so). 
\begin{Verbatim}[commandchars=!@|,fontsize=\small,frame=single,label=Example]
  !gapprompt@gap>| !gapinput@a5:= AlternatingGroup( 5 );|
  Alt( [ 1 .. 5 ] )
  !gapprompt@gap>| !gapinput@char:= First( Irr( a5 ), x -> x[1] = 4 );|
  Character( CharacterTable( Alt( [ 1 .. 5 ] ) ), [ 4, 0, 1, -1, -1 ] )
  !gapprompt@gap>| !gapinput@hom:=IrreducibleRepresentationsDixon( a5, char: unitary );;|
  !gapprompt@gap>| !gapinput@Order( a5.1*a5.2 ) = Order( Image( hom, a5.1 )*Image( hom, a5.2 ) );|
  true
  !gapprompt@gap>| !gapinput@reps:= List( ConjugacyClasses( a5 ), Representative );;|
  !gapprompt@gap>| !gapinput@List( reps, g -> TraceMat( Image( hom, g ) ) );|
  [ 4, 0, 1, -1, -1 ]
\end{Verbatim}
 }

 }

  
\section{\textcolor{Chapter }{Representations Given by Modules}}\label{Representations given by modules}
\logpage{[ 71, 15, 0 ]}
\hyperdef{L}{X7E51AACD79CE0BC8}{}
{
  This section describes functions that return certain modules of a given group.
(Extensions by modules can be formed by the command \texttt{Extensions} (\ref{Extensions}).) 

\subsection{\textcolor{Chapter }{IrreducibleModules}}
\logpage{[ 71, 15, 1 ]}\nobreak
\hyperdef{L}{X87E82F8085745523}{}
{\noindent\textcolor{FuncColor}{$\triangleright$\enspace\texttt{IrreducibleModules({\mdseries\slshape G, F, dim})\index{IrreducibleModules@\texttt{IrreducibleModules}}
\label{IrreducibleModules}
}\hfill{\scriptsize (operation)}}\\


 returns a list of length 2. The first entry is a generating system of \mbox{\texttt{\mdseries\slshape G}}. The second entry is a list of all irreducible modules of \mbox{\texttt{\mdseries\slshape G}} over the field \mbox{\texttt{\mdseries\slshape F}} of dimension at most \mbox{\texttt{\mdseries\slshape dim}}, given as MeatAxe modules (see{\nobreakspace}\texttt{GModuleByMats} (\ref{GModuleByMats:for generators and a field})). Pass \texttt{0} for \mbox{\texttt{\mdseries\slshape dim}} to impose no dimension bound. }

 

\subsection{\textcolor{Chapter }{AbsolutelyIrreducibleModules}}
\logpage{[ 71, 15, 2 ]}\nobreak
\hyperdef{L}{X7D0BD5337D1C069B}{}
{\noindent\textcolor{FuncColor}{$\triangleright$\enspace\texttt{AbsolutelyIrreducibleModules({\mdseries\slshape G, F, dim})\index{AbsolutelyIrreducibleModules@\texttt{AbsolutelyIrreducibleModules}}
\label{AbsolutelyIrreducibleModules}
}\hfill{\scriptsize (operation)}}\\
\noindent\textcolor{FuncColor}{$\triangleright$\enspace\texttt{AbsoluteIrreducibleModules({\mdseries\slshape G, F, dim})\index{AbsoluteIrreducibleModules@\texttt{AbsoluteIrreducibleModules}}
\label{AbsoluteIrreducibleModules}
}\hfill{\scriptsize (operation)}}\\
\noindent\textcolor{FuncColor}{$\triangleright$\enspace\texttt{AbsolutIrreducibleModules({\mdseries\slshape G, F, dim})\index{AbsolutIrreducibleModules@\texttt{AbsolutIrreducibleModules}}
\label{AbsolutIrreducibleModules}
}\hfill{\scriptsize (operation)}}\\


 \texttt{AbsolutelyIrreducibleModules} returns a list of length 2. The first entry is a generating system of the
group \mbox{\texttt{\mdseries\slshape G}}. The second entry is a list of all those absolutely irreducible modules of \mbox{\texttt{\mdseries\slshape G}} that can be realized over the finite field \mbox{\texttt{\mdseries\slshape F}} and have dimension at most \mbox{\texttt{\mdseries\slshape dim}}, given as MeatAxe modules (see{\nobreakspace}\texttt{GModuleByMats} (\ref{GModuleByMats:for generators and a field})). 

 The other two names are just synonyms. }

 

\subsection{\textcolor{Chapter }{RegularModule}}
\logpage{[ 71, 15, 3 ]}\nobreak
\hyperdef{L}{X7EB88B2E87AF5556}{}
{\noindent\textcolor{FuncColor}{$\triangleright$\enspace\texttt{RegularModule({\mdseries\slshape G, F})\index{RegularModule@\texttt{RegularModule}}
\label{RegularModule}
}\hfill{\scriptsize (operation)}}\\


 returns a list of length 2. The first entry is a generating system of \mbox{\texttt{\mdseries\slshape G}}. The second entry is the regular module of \mbox{\texttt{\mdseries\slshape G}} over \mbox{\texttt{\mdseries\slshape F}}, given as a MeatAxe module (see{\nobreakspace}\texttt{GModuleByMats} (\ref{GModuleByMats:for generators and a field})). }

 }

  
\section{\textcolor{Chapter }{The Dixon\texttt{\symbol{45}}Schneider Algorithm}}\label{The Dixon-Schneider Algorithm}
\logpage{[ 71, 16, 0 ]}
\hyperdef{L}{X86CDA4007A5EF704}{}
{
  \index{Dixon\texttt{\symbol{45}}Schneider algorithm} The \textsf{GAP} library implementation of the Dixon\texttt{\symbol{45}}Schneider algorithm
first computes the linear characters, using the commutator factor group. If
irreducible characters are missing afterwards, they are computed using the
techniques described in \cite{Dix67}, \cite{Sch90} and \cite{Hulpke93}. 

 Called with a group $G$, the function \texttt{CharacterTable} (\ref{CharacterTable:for a group}) returns a character table object that stores already information such as class
lengths, but not the irreducible characters. The routines that compute the
irreducibles may use the information that is already contained in this table
object. In particular the ordering of classes in the computed characters
coincides with the ordering of classes in the character table of \mbox{\texttt{\mdseries\slshape G}} (see{\nobreakspace}\ref{The Interface between Character Tables and Groups}). Thus it is possible to combine computations using the group with character
theoretic computations (see{\nobreakspace}\ref{Advanced Methods for Dixon-Schneider Calculations} for details), for example one can enter known characters. Note that the user
is responsible for the correctness of the characters. (There is little use in
providing the trivial character to the routine.) 

 The computation of irreducible characters from the group needs to identify the
classes of group elements very often, so it can be helpful to store a class
list of all group elements. Since this is obviously limited by the group
order, it is controlled by the global function \texttt{IsDxLargeGroup} (\ref{IsDxLargeGroup}). 

 The routines compute in a prime field of size $p$, such that the exponent of the group divides $(p-1)$ and such that $2 \sqrt{{|G|}} < p$. Currently prime fields of size smaller than $65\,536$ are handled more efficiently than larger prime fields, so the runtime of the
character calculation depends on how large the chosen prime is. 

 The routine stores a Dixon record (see{\nobreakspace}\texttt{DixonRecord} (\ref{DixonRecord})) in the group that helps routines that identify classes, for example \texttt{FusionConjugacyClasses} (\ref{FusionConjugacyClasses:for two groups}), to work much faster. Note that interrupting
Dixon\texttt{\symbol{45}}Schneider calculations will prevent \textsf{GAP} from cleaning up the Dixon record; when the computation by \texttt{IrrDixonSchneider} (\ref{IrrDixonSchneider}) is complete, the possibly large record is shrunk to an acceptable size. }

  
\section{\textcolor{Chapter }{Advanced Methods for Dixon\texttt{\symbol{45}}Schneider Calculations}}\label{Advanced Methods for Dixon-Schneider Calculations}
\logpage{[ 71, 17, 0 ]}
\hyperdef{L}{X7C083207868066C1}{}
{
  \index{irreducible characters!computation} The computation of irreducible characters of very large groups may take quite
some time. On the other hand, for the expert only a few irreducible characters
may be needed, since the other ones can be computed using character theoretic
methods such as tensoring, induction, and restriction. Thus \textsf{GAP} provides also step\texttt{\symbol{45}}by\texttt{\symbol{45}}step routines for
doing the calculations. These routines allow one to compute some characters
and to stop before all are calculated. Note that there is no ``safety net'': The routines (being somehow internal) do no error checking, and assume the
information given is correct. 

 When the info level of \texttt{InfoCharacterTable} (\ref{InfoCharacterTable}) if positive, information about the progress of splitting is printed. (The
default value is zero.) 

\subsection{\textcolor{Chapter }{DixonRecord}}
\logpage{[ 71, 17, 1 ]}\nobreak
\hyperdef{L}{X7C398F2680C8616B}{}
{\noindent\textcolor{FuncColor}{$\triangleright$\enspace\texttt{DixonRecord({\mdseries\slshape G})\index{DixonRecord@\texttt{DixonRecord}}
\label{DixonRecord}
}\hfill{\scriptsize (attribute)}}\\


 The \texttt{DixonRecord} of a group contains information used by the routines to compute the
irreducible characters and related information via the
Dixon\texttt{\symbol{45}}Schneider algorithm such as class arrangement and
character spaces split obtained so far. Usually this record is passed as
argument to all subfunctions to avoid a long argument list. It has a component \texttt{conjugacyClasses} which contains the classes of \mbox{\texttt{\mdseries\slshape G}} \emph{ordered as the algorithm needs them}. }

 

\subsection{\textcolor{Chapter }{DixonInit}}
\logpage{[ 71, 17, 2 ]}\nobreak
\hyperdef{L}{X7E33C03E7BDDC9B0}{}
{\noindent\textcolor{FuncColor}{$\triangleright$\enspace\texttt{DixonInit({\mdseries\slshape G})\index{DixonInit@\texttt{DixonInit}}
\label{DixonInit}
}\hfill{\scriptsize (function)}}\\


 This function does all the initializations for the
Dixon\texttt{\symbol{45}}Schneider algorithm. This includes calculation of
conjugacy classes, power maps, linear characters and character morphisms. It
returns a record (see{\nobreakspace}\texttt{DixonRecord} (\ref{DixonRecord}) and Section \ref{Components of a Dixon Record}) that can be used when calculating the irreducible characters of \mbox{\texttt{\mdseries\slshape G}} interactively. }

 

\subsection{\textcolor{Chapter }{DixontinI}}
\logpage{[ 71, 17, 3 ]}\nobreak
\hyperdef{L}{X868476037907918F}{}
{\noindent\textcolor{FuncColor}{$\triangleright$\enspace\texttt{DixontinI({\mdseries\slshape D})\index{DixontinI@\texttt{DixontinI}}
\label{DixontinI}
}\hfill{\scriptsize (function)}}\\


 This function ends a Dixon\texttt{\symbol{45}}Schneider calculation. It sorts
the characters according to the degree and unbinds components in the Dixon
record that are not of use any longer. It returns a list of irreducible
characters. }

 

\subsection{\textcolor{Chapter }{DixonSplit}}
\logpage{[ 71, 17, 4 ]}\nobreak
\hyperdef{L}{X87ABE0B081DAD476}{}
{\noindent\textcolor{FuncColor}{$\triangleright$\enspace\texttt{DixonSplit({\mdseries\slshape D})\index{DixonSplit@\texttt{DixonSplit}}
\label{DixonSplit}
}\hfill{\scriptsize (function)}}\\


 This function performs one splitting step in the
Dixon\texttt{\symbol{45}}Schneider algorithm. It selects a class, computes the
(partial) class sum matrix, uses it to split character spaces and stores all
the irreducible characters obtained that way. 

 The class to use for splitting is chosen via \texttt{BestSplittingMatrix} (\ref{BestSplittingMatrix}) and the options described for this function apply here. 

 \texttt{DixonSplit} returns the number of the class that was used for splitting if a split was
performed, and \texttt{fail} otherwise. }

 

\subsection{\textcolor{Chapter }{BestSplittingMatrix}}
\logpage{[ 71, 17, 5 ]}\nobreak
\hyperdef{L}{X7BFD4C1A821731FB}{}
{\noindent\textcolor{FuncColor}{$\triangleright$\enspace\texttt{BestSplittingMatrix({\mdseries\slshape D})\index{BestSplittingMatrix@\texttt{BestSplittingMatrix}}
\label{BestSplittingMatrix}
}\hfill{\scriptsize (function)}}\\


 returns the number of the class sum matrix that is assumed to yield the best
(cost/earning ration) split. This matrix then will be the next one computed
and used. 

 The global option \texttt{maxclasslen} (defaulting to \texttt{infinity} (\ref{infinity})) is recognized by \texttt{BestSplittingMatrix}: Only classes whose length is limited by the value of this option will be
considered for splitting. If no usable class remains, \texttt{fail} is returned. }

 

\subsection{\textcolor{Chapter }{DxIncludeIrreducibles}}
\logpage{[ 71, 17, 6 ]}\nobreak
\hyperdef{L}{X7C85B56C80BFA2E3}{}
{\noindent\textcolor{FuncColor}{$\triangleright$\enspace\texttt{DxIncludeIrreducibles({\mdseries\slshape D, new[, newmod]})\index{DxIncludeIrreducibles@\texttt{DxIncludeIrreducibles}}
\label{DxIncludeIrreducibles}
}\hfill{\scriptsize (function)}}\\


 This function takes a list of irreducible characters \mbox{\texttt{\mdseries\slshape new}}, each given as a list of values (corresponding to the class arrangement in \mbox{\texttt{\mdseries\slshape D}}), and adds these to a partial computed list of irreducibles as maintained by
the Dixon record \mbox{\texttt{\mdseries\slshape D}}. This permits one to add characters in interactive use obtained from other
sources and to continue the Dixon\texttt{\symbol{45}}Schneider calculation
afterwards. If the optional argument \mbox{\texttt{\mdseries\slshape newmod}} is given, it must be a list of reduced characters, corresponding to \mbox{\texttt{\mdseries\slshape new}}. (Otherwise the function has to reduce the characters itself.) 

 The function closes the new characters under the action of Galois
automorphisms and tensor products with linear characters. }

 

\subsection{\textcolor{Chapter }{SplitCharacters}}
\logpage{[ 71, 17, 7 ]}\nobreak
\hyperdef{L}{X87A5B5C77F7F348E}{}
{\noindent\textcolor{FuncColor}{$\triangleright$\enspace\texttt{SplitCharacters({\mdseries\slshape D, list})\index{SplitCharacters@\texttt{SplitCharacters}}
\label{SplitCharacters}
}\hfill{\scriptsize (function)}}\\


 This routine decomposes the characters given in \mbox{\texttt{\mdseries\slshape list}} according to the character spaces found up to this point. By applying this
routine to tensor products etc., it may result in characters with smaller
norm, even irreducible ones. Since the recalculation of characters is only
possible if the degree is small enough, the splitting process is applied only
to characters of sufficiently small degree. }

 

\subsection{\textcolor{Chapter }{IsDxLargeGroup}}
\logpage{[ 71, 17, 8 ]}\nobreak
\hyperdef{L}{X8089009E7EA85BC8}{}
{\noindent\textcolor{FuncColor}{$\triangleright$\enspace\texttt{IsDxLargeGroup({\mdseries\slshape G})\index{IsDxLargeGroup@\texttt{IsDxLargeGroup}}
\label{IsDxLargeGroup}
}\hfill{\scriptsize (function)}}\\


 returns \texttt{true} if the order of the group \mbox{\texttt{\mdseries\slshape G}} is smaller than the current value of the global variable \texttt{DXLARGEGROUPORDER}, and \texttt{false} otherwise. In Dixon\texttt{\symbol{45}}Schneider calculations, for small
groups in the above sense a class map is stored, whereas for large groups,
each occurring element is identified individually. }

 }

  
\section{\textcolor{Chapter }{Components of a Dixon Record}}\label{Components of a Dixon Record}
\logpage{[ 71, 18, 0 ]}
\hyperdef{L}{X7C1153637E7D2133}{}
{
  The ``Dixon record'' \mbox{\texttt{\mdseries\slshape D}} returned by \texttt{DixonInit} (\ref{DixonInit}) stores all the information that is used by the
Dixon\texttt{\symbol{45}}Schneider routines while computing the irreducible
characters of a group. Some entries, however, may be useful to know about when
using the algorithm interactively, see \ref{An Example of Advanced Dixon-Schneider Calculations}. 
\begin{description}
\item[{\texttt{group}}]  the group $G$ of which the character table is to be computed, 
\item[{\texttt{conjugacyClasses}}]  classes of $G$ (all characters stored in the Dixon record correspond to this arrangement of
classes), 
\item[{\texttt{irreducibles}}]  the already known irreducible characters (given as lists of their values on
the conjugacy classes), 
\item[{\texttt{characterTable}}]  the \texttt{CharacterTable} (\ref{CharacterTable:for a group}) value of $G$ (whose irreducible characters are not yet known), 
\item[{\texttt{ClassElement( \mbox{\texttt{\mdseries\slshape D}}, \mbox{\texttt{\mdseries\slshape el}} )}}]  a function that returns the number of the class of $G$ that contains the element \mbox{\texttt{\mdseries\slshape el}}. 
\end{description}
 }

  
\section{\textcolor{Chapter }{An Example of Advanced Dixon\texttt{\symbol{45}}Schneider Calculations}}\label{An Example of Advanced Dixon-Schneider Calculations}
\logpage{[ 71, 19, 0 ]}
\hyperdef{L}{X782B5E37848786BC}{}
{
  First, we set the appropriate info level higher. 

 
\begin{Verbatim}[commandchars=!@|,fontsize=\small,frame=single,label=Example]
  !gapprompt@gap>| !gapinput@SetInfoLevel( InfoCharacterTable, 1 );|
\end{Verbatim}
 

 for printout of some internal results. We now define our group, which is
isomorphic to PSL$_4(3)$. 

 
\begin{Verbatim}[commandchars=!@|,fontsize=\small,frame=single,label=Example]
  !gapprompt@gap>| !gapinput@g:= PrimitiveGroup(40,5);|
  PSL(4, 3)
  !gapprompt@gap>| !gapinput@Size(g);|
  6065280
  !gapprompt@gap>| !gapinput@d:= DixonInit( g );;|
  #I  29 classes
  #I  choosing prime 65521
  !gapprompt@gap>| !gapinput@c:= d.characterTable;;|
\end{Verbatim}
 

 After the initialisation, one structure matrix is evaluated, yielding smaller
spaces and several irreducible characters. 

 
\begin{Verbatim}[commandchars=!@|,fontsize=\small,frame=single,label=Example]
  !gapprompt@gap>| !gapinput@DixonSplit( d );|
  #I  Matrix 2,Representative of Order 3,Centralizer: 5832
  #I  Dimensions: [ [ 1, 6 ], [ 2, 3 ], [ 4, 1 ], [ 12, 1 ] ]
  2
\end{Verbatim}
 

 In this case spaces of the listed dimensions are a result of the splitting
process. The three two dimensional spaces are split successfully by
combinatoric means. 

 We obtain several irreducible characters by tensor products and notify them to
the Dixon record. 

 
\begin{Verbatim}[commandchars=!@|,fontsize=\small,frame=single,label=Example]
  !gapprompt@gap>| !gapinput@asp:= AntiSymmetricParts( c, d.irreducibles, 2 );;|
  !gapprompt@gap>| !gapinput@ro:= ReducedCharacters( c, d.irreducibles, asp );;|
  !gapprompt@gap>| !gapinput@Length( ro.irreducibles );|
  3
  !gapprompt@gap>| !gapinput@DxIncludeIrreducibles( d, ro.irreducibles );|
\end{Verbatim}
 

 Finally we calculate the characters induced from all cyclic subgroups and
obtain the missing irreducibles by applying the
LLL\texttt{\symbol{45}}algorithm to them. 

 
\begin{Verbatim}[commandchars=!@|,fontsize=\small,frame=single,label=Example]
  !gapprompt@gap>| !gapinput@ic:= InducedCyclic( c, "all" );;|
  !gapprompt@gap>| !gapinput@ro:= ReducedCharacters( c, d.irreducibles, ic );;|
  !gapprompt@gap>| !gapinput@Length( ro.irreducibles );|
  0
  !gapprompt@gap>| !gapinput@l:= LLL( c, ro.remainders );;|
  !gapprompt@gap>| !gapinput@Length( l.irreducibles );|
  13
\end{Verbatim}
 

 The LLL returns class function objects (see Chapter{\nobreakspace}\ref{Class Functions}), and the Dixon record works with character values lists. So we convert them
to a list of values before feeding them in the machinery of the
Dixon\texttt{\symbol{45}}algorithm. 

 
\begin{Verbatim}[commandchars=!@|,fontsize=\small,frame=single,label=Example]
  !gapprompt@gap>| !gapinput@l.irreducibles[1];|
  Character( CharacterTable( PSL(4, 3) ),
   [ 640, -8, -8, -8, 1, 1, 1, 0, 0, 0, 0, 0, 0, 0, 0, 0, 0, 0, 0, 0,
    0, 0, 0, 0, E(13)^7+E(13)^8+E(13)^11, E(13)^4+E(13)^10+E(13)^12,
    E(13)^2+E(13)^5+E(13)^6, E(13)+E(13)^3+E(13)^9, 0 ] )
  !gapprompt@gap>| !gapinput@l:=List(l.irreducibles,ValuesOfClassFunction);;|
  !gapprompt@gap>| !gapinput@DxIncludeIrreducibles( d, l );|
  !gapprompt@gap>| !gapinput@Length( d.irreducibles );|
  29
  !gapprompt@gap>| !gapinput@Length( d.classes );|
  29
\end{Verbatim}
 

 It turns out we have found all irreducible characters. As the last step, we
obtain the irreducible characters and tell them to the group. This makes them
available also to the character table. 

 
\begin{Verbatim}[commandchars=!@|,fontsize=\small,frame=single,label=Example]
  !gapprompt@gap>| !gapinput@irrs:= DixontinI( d );;|
  #I  Total:1 matrices,[ 2 ]
  !gapprompt@gap>| !gapinput@SetIrr(g,irrs);|
  !gapprompt@gap>| !gapinput@Length(Irr(c));|
  29
  !gapprompt@gap>| !gapinput@SetInfoLevel( InfoCharacterTable, 0 );|
\end{Verbatim}
 }

  
\section{\textcolor{Chapter }{Constructing Character Tables from Others}}\label{Constructing Character Tables from Others}
\logpage{[ 71, 20, 0 ]}
\hyperdef{L}{X7C38C5067941D496}{}
{
  The following operations take one or more character table arguments, and
return a character table. This holds also for \texttt{BrauerTable} (\ref{BrauerTable:for a character table, and a prime integer}). Note that the return value of \texttt{BrauerTable} (\ref{BrauerTable:for a character table, and a prime integer}) will in general not know the irreducible Brauer characters, and \textsf{GAP} might be unable to compute these characters. 

 \emph{Note} that whenever fusions between input and output tables occur in these
operations, they are stored on the concerned tables, and the \texttt{NamesOfFusionSources} (\ref{NamesOfFusionSources}) values are updated. 

 (The interactive construction of character tables using character theoretic
methods and incomplete tables is not described here.) \emph{Currently it is not supported and will be described in a chapter of its own
when it becomes available}. 

\subsection{\textcolor{Chapter }{CharacterTableDirectProduct}}
\logpage{[ 71, 20, 1 ]}\nobreak
\hyperdef{L}{X7BE1572D7BBC6AC8}{}
{\noindent\textcolor{FuncColor}{$\triangleright$\enspace\texttt{CharacterTableDirectProduct({\mdseries\slshape tbl1, tbl2})\index{CharacterTableDirectProduct@\texttt{CharacterTableDirectProduct}}
\label{CharacterTableDirectProduct}
}\hfill{\scriptsize (operation)}}\\


 is the table of the direct product of the character tables \mbox{\texttt{\mdseries\slshape tbl1}} and \mbox{\texttt{\mdseries\slshape tbl2}}. 

 The matrix of irreducibles of this table is the Kronecker product
(see{\nobreakspace}\texttt{KroneckerProduct} (\ref{KroneckerProduct})) of the irreducibles of \mbox{\texttt{\mdseries\slshape tbl1}} and \mbox{\texttt{\mdseries\slshape tbl2}}. 

 Products of ordinary and Brauer character tables are supported. 

 In general, the result will not know an underlying group, so missing power
maps (for prime divisors of the result) and irreducibles of the input tables
may be computed in order to construct the table of the direct product. 

 The embeddings of the input tables into the direct product are stored, they
can be fetched with \texttt{GetFusionMap} (\ref{GetFusionMap}); if \mbox{\texttt{\mdseries\slshape tbl1}} is equal to \mbox{\texttt{\mdseries\slshape tbl2}} then the two embeddings are distinguished by their \texttt{specification} components \texttt{"1"} and \texttt{"2"}, respectively. 

 Analogously, the projections from the direct product onto the input tables are
stored, and can be distinguished by the \texttt{specification} components.  

 The attribute \texttt{FactorsOfDirectProduct} (\ref{FactorsOfDirectProduct}) is set to the lists of arguments. 

 The \texttt{*} operator for two character tables (see{\nobreakspace}\ref{Operators for Character Tables}) delegates to \texttt{CharacterTableDirectProduct}. 
\begin{Verbatim}[commandchars=!@|,fontsize=\small,frame=single,label=Example]
  !gapprompt@gap>| !gapinput@c2:= CharacterTable( "Cyclic", 2 );;|
  !gapprompt@gap>| !gapinput@s3:= CharacterTable( "Symmetric", 3 );;|
  !gapprompt@gap>| !gapinput@Display( CharacterTableDirectProduct( c2, s3 ) );|
  C2xSym(3)
  
       2  2  2  1  2  2  1
       3  1  .  1  1  .  1
  
         1a 2a 3a 2b 2c 6a
      2P 1a 1a 3a 1a 1a 3a
      3P 1a 2a 1a 2b 2c 2b
  
  X.1     1 -1  1  1 -1  1
  X.2     2  . -1  2  . -1
  X.3     1  1  1  1  1  1
  X.4     1 -1  1 -1  1 -1
  X.5     2  . -1 -2  .  1
  X.6     1  1  1 -1 -1 -1
\end{Verbatim}
 }

 

\subsection{\textcolor{Chapter }{FactorsOfDirectProduct}}
\logpage{[ 71, 20, 2 ]}\nobreak
\hyperdef{L}{X7C97CF727FBDFCAB}{}
{\noindent\textcolor{FuncColor}{$\triangleright$\enspace\texttt{FactorsOfDirectProduct({\mdseries\slshape tbl})\index{FactorsOfDirectProduct@\texttt{FactorsOfDirectProduct}}
\label{FactorsOfDirectProduct}
}\hfill{\scriptsize (attribute)}}\\


 For an ordinary character table that has been constructed via \texttt{CharacterTableDirectProduct} (\ref{CharacterTableDirectProduct}), the value of \texttt{FactorsOfDirectProduct} is the list of arguments in the \texttt{CharacterTableDirectProduct} (\ref{CharacterTableDirectProduct}) call. 

 Note that there is no default method for \emph{computing} the value of \texttt{FactorsOfDirectProduct}. }

 

\subsection{\textcolor{Chapter }{CharacterTableFactorGroup}}
\logpage{[ 71, 20, 3 ]}\nobreak
\hyperdef{L}{X7C3A4E5283B240BE}{}
{\noindent\textcolor{FuncColor}{$\triangleright$\enspace\texttt{CharacterTableFactorGroup({\mdseries\slshape tbl, classes})\index{CharacterTableFactorGroup@\texttt{CharacterTableFactorGroup}}
\label{CharacterTableFactorGroup}
}\hfill{\scriptsize (operation)}}\\


 is the character table of the factor group of the ordinary character table \mbox{\texttt{\mdseries\slshape tbl}} by the normal closure of the classes whose positions are contained in the list \mbox{\texttt{\mdseries\slshape classes}}. 

 The \texttt{/} operator for a character table and a list of class positions
(see{\nobreakspace}\ref{Operators for Character Tables}) delegates to \texttt{CharacterTableFactorGroup}. 

 
\begin{Verbatim}[commandchars=!@|,fontsize=\small,frame=single,label=Example]
  !gapprompt@gap>| !gapinput@s4:= CharacterTable( "Symmetric", 4 );;|
  !gapprompt@gap>| !gapinput@ClassPositionsOfNormalSubgroups( s4 );|
  [ [ 1 ], [ 1, 3 ], [ 1, 3, 4 ], [ 1 .. 5 ] ]
  !gapprompt@gap>| !gapinput@f:= CharacterTableFactorGroup( s4, [ 3 ] );|
  CharacterTable( "Sym(4)/[ 1, 3 ]" )
  !gapprompt@gap>| !gapinput@Display( f );|
  Sym(4)/[ 1, 3 ]
  
       2  1  1  .
       3  1  .  1
  
         1a 2a 3a
      2P 1a 1a 3a
      3P 1a 2a 1a
  
  X.1     1 -1  1
  X.2     2  . -1
  X.3     1  1  1
\end{Verbatim}
 }

 

\subsection{\textcolor{Chapter }{CharacterTableIsoclinic}}
\logpage{[ 71, 20, 4 ]}\nobreak
\hyperdef{L}{X85BE46F784B83938}{}
{\noindent\textcolor{FuncColor}{$\triangleright$\enspace\texttt{CharacterTableIsoclinic({\mdseries\slshape tbl[, arec]})\index{CharacterTableIsoclinic@\texttt{CharacterTableIsoclinic}}
\label{CharacterTableIsoclinic}
}\hfill{\scriptsize (operation)}}\\
\noindent\textcolor{FuncColor}{$\triangleright$\enspace\texttt{CharacterTableIsoclinic({\mdseries\slshape tbl[, classes][, centre]})\index{CharacterTableIsoclinic@\texttt{CharacterTableIsoclinic}!for a character table and one or two lists}
\label{CharacterTableIsoclinic:for a character table and one or two lists}
}\hfill{\scriptsize (operation)}}\\
\noindent\textcolor{FuncColor}{$\triangleright$\enspace\texttt{CharacterTableIsoclinic({\mdseries\slshape modtbl, ordiso})\index{CharacterTableIsoclinic@\texttt{CharacterTableIsoclinic}!for a Brauer table and an ordinary table}
\label{CharacterTableIsoclinic:for a Brauer table and an ordinary table}
}\hfill{\scriptsize (operation)}}\\
\noindent\textcolor{FuncColor}{$\triangleright$\enspace\texttt{SourceOfIsoclinicTable({\mdseries\slshape tbl})\index{SourceOfIsoclinicTable@\texttt{SourceOfIsoclinicTable}}
\label{SourceOfIsoclinicTable}
}\hfill{\scriptsize (attribute)}}\\


 Let \mbox{\texttt{\mdseries\slshape tbl}} be the (ordinary or modular) character table of a group $H$ with the structure $p.G.p$ for some prime $p$, that is, $H/Z$ has a normal subgroup $N$ of index $p$ and a central subgroup $Z$ of order $p$ contained in $N$. 

 Then \texttt{CharacterTableIsoclinic} returns the table of an isoclinic group in the sense of the \textsf{Atlas} of Finite Groups \cite[Chapter 6, Section 7]{CCN85}. 

 If $p = 2$ then also the case $H = 4.G.2$ is supported, that is, $Z$ has order four and $N$ has index two in $H$. 

 The optional arguments are needed if \mbox{\texttt{\mdseries\slshape tbl}} does not determine the class positions of $N$ or $Z$ uniquely, and in the case $p > 2$ if one wants to specify a ``variant number'' for the result. 

 
\begin{itemize}
\item  In general, the values can be specified via a record \mbox{\texttt{\mdseries\slshape arec}}. If $N$ is not uniquely determined then the positions of the classes forming $N$ must be entered as the value of the component \texttt{normalSubgroup}. If $Z$ is not unique inside $N$ then the class position of a generator of $Z$ must be entered as the value of the component \texttt{centralElement}. 
\item  If $p = 2$ then one may specify the positions of the classes forming $N$ via a list \mbox{\texttt{\mdseries\slshape classes}}, and the positions of the classes in $Z$ as a list \mbox{\texttt{\mdseries\slshape centre}}; if $Z$ has order $2$ then \mbox{\texttt{\mdseries\slshape centre}} can be also the position of the involution in $Z$. 
\end{itemize}
 

 Note that also if \mbox{\texttt{\mdseries\slshape tbl}} is a Brauer table then \texttt{normalSubgroup} and \texttt{centralElement}, resp.{\nobreakspace} \mbox{\texttt{\mdseries\slshape classes}} and \mbox{\texttt{\mdseries\slshape centre}}, denote class numbers w.r.t.{\nobreakspace}the \emph{ordinary} character table. 

 If $p$ is odd then the \textsf{Atlas} construction describes $p$ isoclinic variants that arise from $p.G.p$. (These groups need not be pairwise nonisomorphic.) Entering an integer $k \in \{ 1, 2, \ldots, p-1 \}$ as the value of the component \texttt{k} of \mbox{\texttt{\mdseries\slshape arec}} yields the $k$\texttt{\symbol{45}}th of the corresponding character tables; the default for \texttt{k} is $1$. 

 
\begin{Verbatim}[commandchars=!@|,fontsize=\small,frame=single,label=Example]
  !gapprompt@gap>| !gapinput@d8:= CharacterTable( "Dihedral", 8 );|
  CharacterTable( "Dihedral(8)" )
  !gapprompt@gap>| !gapinput@nsg:= ClassPositionsOfNormalSubgroups( d8 );|
  [ [ 1 ], [ 1, 3 ], [ 1 .. 3 ], [ 1, 3, 4 ], [ 1, 3 .. 5 ], [ 1 .. 5 ]
   ]
  !gapprompt@gap>| !gapinput@isod8:= CharacterTableIsoclinic( d8, nsg[3] );;|
  !gapprompt@gap>| !gapinput@Display( isod8 );|
  Isoclinic(Dihedral(8))
  
       2  3  2  3  2  2
  
         1a 4a 2a 4b 4c
      2P 1a 2a 1a 2a 2a
  
  X.1     1  1  1  1  1
  X.2     1  1  1 -1 -1
  X.3     1 -1  1  1 -1
  X.4     1 -1  1 -1  1
  X.5     2  . -2  .  .
  !gapprompt@gap>| !gapinput@t1:= CharacterTable( SmallGroup( 27, 3 ) );;|
  !gapprompt@gap>| !gapinput@t2:= CharacterTable( SmallGroup( 27, 4 ) );;|
  !gapprompt@gap>| !gapinput@nsg:= ClassPositionsOfNormalSubgroups( t1 );|
  [ [ 1 ], [ 1, 4, 8 ], [ 1, 2, 4, 5, 8 ], [ 1, 3, 4, 7, 8 ],
    [ 1, 4, 6, 8, 11 ], [ 1, 4, 8, 9, 10 ], [ 1 .. 11 ] ]
  !gapprompt@gap>| !gapinput@iso1:= CharacterTableIsoclinic( t1, rec( k:= 1,|
  !gapprompt@>| !gapinput@              normalSubgroup:= nsg[3] ) );;|
  !gapprompt@gap>| !gapinput@iso2:= CharacterTableIsoclinic( t1, rec( k:= 2,|
  !gapprompt@>| !gapinput@              normalSubgroup:= nsg[3] ) );;|
  !gapprompt@gap>| !gapinput@TransformingPermutationsCharacterTables( iso1, t1 ) <> fail;|
  false
  !gapprompt@gap>| !gapinput@TransformingPermutationsCharacterTables( iso1, t2 ) <> fail;|
  true
  !gapprompt@gap>| !gapinput@TransformingPermutationsCharacterTables( iso2, t2 ) <> fail;|
  true
\end{Verbatim}
 

 For an ordinary character table that has been constructed via \texttt{CharacterTableIsoclinic}, the value of \texttt{SourceOfIsoclinicTable} encodes this construction, and is defined as follows. If $p = 2$ then the value is the list with entries \mbox{\texttt{\mdseries\slshape tbl}}, \mbox{\texttt{\mdseries\slshape classes}}, the list of class positions of the nonidentity elements in $Z$, and the class position of a generator of $Z$. If $p$ is an odd prime then the value is a record with the following components. 

 
\begin{description}
\item[{\texttt{table}}]  the character table \mbox{\texttt{\mdseries\slshape tbl}}, 
\item[{\texttt{p}}]  the prime $p$, 
\item[{\texttt{k}}]  the variant number $k$, 
\item[{\texttt{outerClasses}}]  the list of length $p-1$ that contains at position $i$ the sorted list of class positions of the $i$\texttt{\symbol{45}}th coset of the normal subgroup $N$ 
\item[{\texttt{centralElement}}]  the class position of a generator of the central subgroup $Z$. 
\end{description}
 

 There is no default method for \emph{computing} the value of \texttt{SourceOfIsoclinicTable}. 

 
\begin{Verbatim}[commandchars=!@|,fontsize=\small,frame=single,label=Example]
  !gapprompt@gap>| !gapinput@SourceOfIsoclinicTable( isod8 );|
  [ CharacterTable( "Dihedral(8)" ), [ 1 .. 3 ], [ 3 ], 3 ]
  !gapprompt@gap>| !gapinput@SourceOfIsoclinicTable( iso1 );|
  rec( centralElement := 4, k := 1,
    outerClasses := [ [ 3, 6, 9 ], [ 7, 10, 11 ] ], p := 3,
    table := CharacterTable( <pc group of size 27 with 3 generators> ) )
\end{Verbatim}
 

 If the arguments of \texttt{CharacterTableIsoclinic} are a Brauer table \mbox{\texttt{\mdseries\slshape modtbl}} and an ordinary table \mbox{\texttt{\mdseries\slshape ordiso}} then the \texttt{SourceOfIsoclinicTable} value of \mbox{\texttt{\mdseries\slshape ordiso}} is assumed to be identical with the \texttt{OrdinaryCharacterTable} (\ref{OrdinaryCharacterTable:for a character table}) value of \mbox{\texttt{\mdseries\slshape modtbl}}, and the specified isoclinic table of \mbox{\texttt{\mdseries\slshape modtbl}} is returned. This variant is useful if one has already constructed \mbox{\texttt{\mdseries\slshape ordiso}} in advance. 

 
\begin{Verbatim}[commandchars=!@|,fontsize=\small,frame=single,label=Example]
  !gapprompt@gap>| !gapinput@g:= GL(2,3);;|
  !gapprompt@gap>| !gapinput@t:= CharacterTable( g );;|
  !gapprompt@gap>| !gapinput@iso:= CharacterTableIsoclinic( t );;|
  !gapprompt@gap>| !gapinput@t3:= t mod 3;;|
  !gapprompt@gap>| !gapinput@iso3:= CharacterTableIsoclinic( t3, iso );;|
  !gapprompt@gap>| !gapinput@TransformingPermutationsCharacterTables( iso3,|
  !gapprompt@>| !gapinput@       CharacterTableIsoclinic( t3 ) ) <> fail;|
  true
\end{Verbatim}
 

 \emph{Theoretical background:} Consider the central product $K$ of $H$ with a cyclic group $C$ of order $p^2$. That is, $K = H C$, $C \leq Z(K)$, and the central subgroup $Z$ of order $p$ in $H$ lies in $C$. There are $p+1$ subgroups of $K$ that contain the normal subgroup $N$ of index $p$ in $H$. One of them is the central product of $C$ with $N$, the others are $H_0 = H$ and its isoclinic variants $H_1, H_2, \ldots, H_{{p-1}}$. We fix $g \in H \setminus N$ and a generator $z$ of $C$, and get $H = N \cup N g \cup N g^2 \cup \cdots \cup N g^{{p-1}}$. Then $H_k$, $0 \leq k \leq p-1$, is given by $N \cup N gz^k \cup N (gz^k)^2 \cup \cdots \cup N (gz^k)^{{p-1}}$. The conjugacy classes of all $H_k$ are in bijection via multiplying the elements with suitable powers of $z$, and the irreducible characters of all $H_k$ extend to $K$ and are in bijection via multiplying the character values with suitable $p^2$\texttt{\symbol{45}}th roots of unity. }

 

\subsection{\textcolor{Chapter }{CharacterTableOfNormalSubgroup}}
\logpage{[ 71, 20, 5 ]}\nobreak
\hyperdef{L}{X806E55A58397B11B}{}
{\noindent\textcolor{FuncColor}{$\triangleright$\enspace\texttt{CharacterTableOfNormalSubgroup({\mdseries\slshape ordtbl, classes})\index{CharacterTableOfNormalSubgroup@\texttt{CharacterTableOfNormalSubgroup}}
\label{CharacterTableOfNormalSubgroup}
}\hfill{\scriptsize (function)}}\\


 Let \mbox{\texttt{\mdseries\slshape ordtbl}} be the ordinary character table of a group $G$, say, and \mbox{\texttt{\mdseries\slshape classes}} be a list of class positions for this table. If the classes given by \mbox{\texttt{\mdseries\slshape classes}} form a normal subgroup $N$, say, of $G$ and if these classes are conjugacy classes of $N$ then this function returns the character table of $N$. In all other cases, the function returns \texttt{fail}. 

 
\begin{Verbatim}[commandchars=!@|,fontsize=\small,frame=single,label=Example]
  !gapprompt@gap>| !gapinput@t:= CharacterTable( "Symmetric", 4 );|
  CharacterTable( "Sym(4)" )
  !gapprompt@gap>| !gapinput@nsg:= ClassPositionsOfNormalSubgroups( t );|
  [ [ 1 ], [ 1, 3 ], [ 1, 3, 4 ], [ 1 .. 5 ] ]
  !gapprompt@gap>| !gapinput@rest:= List( nsg, c -> CharacterTableOfNormalSubgroup( t, c ) );|
  [ CharacterTable( "Rest(Sym(4),[ 1 ])" ), fail, fail,
    CharacterTable( "Rest(Sym(4),[ 1 .. 5 ])" ) ]
\end{Verbatim}
 

 Here is a nontrivial example. We use \texttt{CharacterTableOfNormalSubgroup} for computing the two isoclinic variants of $2.A_5.2$. 

 
\begin{Verbatim}[commandchars=!@|,fontsize=\small,frame=single,label=Example]
  !gapprompt@gap>| !gapinput@g:= SchurCoverOfSymmetricGroup( 5, 3, 1 );;|
  !gapprompt@gap>| !gapinput@c:= CyclicGroup( 4 );;|
  !gapprompt@gap>| !gapinput@dp:= DirectProduct( g, c );;|
  !gapprompt@gap>| !gapinput@diag:= First( Elements( Centre( dp ) ),|
  !gapprompt@>| !gapinput@                 x -> Order( x ) = 2 and|
  !gapprompt@>| !gapinput@                      not x in Image( Embedding( dp, 1 ) ) and|
  !gapprompt@>| !gapinput@                      not x in Image( Embedding( dp, 2 ) ) );;|
  !gapprompt@gap>| !gapinput@fact:= Image( NaturalHomomorphismByNormalSubgroup( dp,|
  !gapprompt@>| !gapinput@                     Subgroup( dp, [ diag ] ) ));;|
  !gapprompt@gap>| !gapinput@t:= CharacterTable( fact );;|
  !gapprompt@gap>| !gapinput@Size( t );|
  480
  !gapprompt@gap>| !gapinput@nsg:= ClassPositionsOfNormalSubgroups( t );;|
  !gapprompt@gap>| !gapinput@rest:= List( nsg, c -> CharacterTableOfNormalSubgroup( t, c ) );;|
  !gapprompt@gap>| !gapinput@index2:= Filtered( rest, x -> x <> fail and Size( x ) = 240 );;|
  !gapprompt@gap>| !gapinput@Length( index2 );|
  2
  !gapprompt@gap>| !gapinput@tg:= CharacterTable( g );;|
  !gapprompt@gap>| !gapinput@SortedList(List(index2,x->IsRecord(|
  !gapprompt@>| !gapinput@      TransformingPermutationsCharacterTables(x,tg))));|
  [ true, false ]
\end{Verbatim}
 

 Alternatively, we could construct the character table of the central product
with character theoretic methods. Or we could use \texttt{CharacterTableIsoclinic} (\ref{CharacterTableIsoclinic}). }

 

\subsection{\textcolor{Chapter }{CharacterTableWreathSymmetric}}
\logpage{[ 71, 20, 6 ]}\nobreak
\hyperdef{L}{X79B75C8582426BC5}{}
{\noindent\textcolor{FuncColor}{$\triangleright$\enspace\texttt{CharacterTableWreathSymmetric({\mdseries\slshape tbl, n})\index{CharacterTableWreathSymmetric@\texttt{CharacterTableWreathSymmetric}}
\label{CharacterTableWreathSymmetric}
}\hfill{\scriptsize (function)}}\\


 returns the character table of the wreath product of a group $G$ with the full symmetric group on \mbox{\texttt{\mdseries\slshape n}} points, where \mbox{\texttt{\mdseries\slshape tbl}} is the character table of $G$. 

 The result has values for \texttt{ClassParameters} (\ref{ClassParameters}) and \texttt{CharacterParameters} (\ref{CharacterParameters}) stored, the entries in these lists are sequences of partitions. Note that this
parametrization prevents the principal character from being the first one in
the list of irreducibles. 

 
\begin{Verbatim}[commandchars=!@|,fontsize=\small,frame=single,label=Example]
  !gapprompt@gap>| !gapinput@c3:= CharacterTable( "Cyclic", 3 );;|
  !gapprompt@gap>| !gapinput@wr:= CharacterTableWreathSymmetric( c3, 2 );;|
  !gapprompt@gap>| !gapinput@Display( wr );|
  C3wrS2
  
       2  1   .   .   1  .   1  1   1   1
       3  2   2   2   2  2   2  1   1   1
  
         1a  3a  3b  3c 3d  3e 2a  6a  6b
      2P 1a  3b  3a  3e 3d  3c 1a  3c  3e
      3P 1a  1a  1a  1a 1a  1a 2a  2a  2a
  
  X.1     1   1   1   1  1   1 -1  -1  -1
  X.2     2   A  /A   B -1  /B  .   .   .
  X.3     2  /A   A  /B -1   B  .   .   .
  X.4     1 -/A  -A  -A  1 -/A -1  /A   A
  X.5     2  -1  -1   2 -1   2  .   .   .
  X.6     1  -A -/A -/A  1  -A -1   A  /A
  X.7     1   1   1   1  1   1  1   1   1
  X.8     1 -/A  -A  -A  1 -/A  1 -/A  -A
  X.9     1  -A -/A -/A  1  -A  1  -A -/A
  
  A = -E(3)^2
    = (1+Sqrt(-3))/2 = 1+b3
  B = 2*E(3)
    = -1+Sqrt(-3) = 2b3
  !gapprompt@gap>| !gapinput@CharacterParameters( wr )[1];|
  [ [ 1, 1 ], [  ], [  ] ]
\end{Verbatim}
 }

 

\subsection{\textcolor{Chapter }{CharacterValueWreathSymmetric}}
\logpage{[ 71, 20, 7 ]}\nobreak
\hyperdef{L}{X83E71B1F7FA70134}{}
{\noindent\textcolor{FuncColor}{$\triangleright$\enspace\texttt{CharacterValueWreathSymmetric({\mdseries\slshape tbl, n, beta, pi})\index{CharacterValueWreathSymmetric@\texttt{CharacterValueWreathSymmetric}}
\label{CharacterValueWreathSymmetric}
}\hfill{\scriptsize (function)}}\\


 Let \mbox{\texttt{\mdseries\slshape tbl}} be the ordinary character table of a group $G$. The aim of this function is to compute a single character value from the
character table of the wreath product of $G$ with the full symmetric group on \mbox{\texttt{\mdseries\slshape n}} points. 

 The conjugacy classes and the irreducible characters of this wreath product
are parametrized by $r$\texttt{\symbol{45}}tuples of partitions which together form a partition of \mbox{\texttt{\mdseries\slshape n}} (see \texttt{PartitionTuples} (\ref{PartitionTuples})), where $r$ is the number of conjugacy classes of $G$. 

 We describe the conjugacy class for which we want to compute the value by the $r$\texttt{\symbol{45}}tuple \mbox{\texttt{\mdseries\slshape pi}} of partitions in question, and describe the character for which we want to
compute the value by the $r$\texttt{\symbol{45}}tuple \mbox{\texttt{\mdseries\slshape beta}} of \texttt{BetaSet} (\ref{BetaSet}) values of the $r$\texttt{\symbol{45}}tuple of partitions in question. 

 
\begin{Verbatim}[commandchars=!@|,fontsize=\small,frame=single,label=Example]
  !gapprompt@gap>| !gapinput@n:= 4;;|
  !gapprompt@gap>| !gapinput@classpara:= [ [], [ 2, 1, 1 ] ];;|
  !gapprompt@gap>| !gapinput@charpara:= [ [ 2, 1 ], [ 1 ] ];;|
  !gapprompt@gap>| !gapinput@betas:= List( charpara, BetaSet );;|
  !gapprompt@gap>| !gapinput@c2:= CharacterTable( "Cyclic", 2 );;|
  !gapprompt@gap>| !gapinput@CharacterValueWreathSymmetric( c2, n, betas, classpara );|
  0
  !gapprompt@gap>| !gapinput@wr:= CharacterTableWreathSymmetric( c2, n );;|
  !gapprompt@gap>| !gapinput@classpos:= Position( ClassParameters( wr ), classpara );;|
  !gapprompt@gap>| !gapinput@charpos:= Position( CharacterParameters( wr ), charpara );;|
  !gapprompt@gap>| !gapinput@Irr( wr )[ charpos, classpos ];|
  0
\end{Verbatim}
 

 This function can be useful if one is interested in only a few character
values. If many character values are needed then it is probably faster to
compute the whole character table of the wreath product using \texttt{CharacterTableWreathSymmetric} (\ref{CharacterTableWreathSymmetric}), which uses intermediate results of recursive computations and therefore can
avoid repetitions. }

 }

  
\section{\textcolor{Chapter }{Sorted Character Tables}}\label{Sorted Character Tables}
\logpage{[ 71, 21, 0 ]}
\hyperdef{L}{X816FCD5A805F9FE8}{}
{
  

\subsection{\textcolor{Chapter }{CharacterTableWithSortedCharacters}}
\logpage{[ 71, 21, 1 ]}\nobreak
\hyperdef{L}{X7D9C4A7F8086F671}{}
{\noindent\textcolor{FuncColor}{$\triangleright$\enspace\texttt{CharacterTableWithSortedCharacters({\mdseries\slshape tbl[, perm]})\index{CharacterTableWithSortedCharacters@\texttt{CharacterTableWithSortedCharacters}}
\label{CharacterTableWithSortedCharacters}
}\hfill{\scriptsize (operation)}}\\


 is a character table that differs from \mbox{\texttt{\mdseries\slshape tbl}} only by the succession of its irreducible characters. This affects the values
of the attributes \texttt{Irr} (\ref{Irr:for a character table}) and \texttt{CharacterParameters} (\ref{CharacterParameters}). Namely, these lists are permuted by the permutation \mbox{\texttt{\mdseries\slshape perm}}. 

 If no second argument is given then a permutation is used that yields
irreducible characters of increasing degree for the result. For the succession
of characters in the result, see{\nobreakspace}\texttt{SortedCharacters} (\ref{SortedCharacters}). 

 The result has all those attributes and properties of \mbox{\texttt{\mdseries\slshape tbl}} that are stored in \texttt{SupportedCharacterTableInfo} (\ref{SupportedCharacterTableInfo}) and do not depend on the ordering of characters. }

 

\subsection{\textcolor{Chapter }{SortedCharacters}}
\logpage{[ 71, 21, 2 ]}\nobreak
\hyperdef{L}{X87E3CF317D8E4EC7}{}
{\noindent\textcolor{FuncColor}{$\triangleright$\enspace\texttt{SortedCharacters({\mdseries\slshape tbl, chars[, flag]})\index{SortedCharacters@\texttt{SortedCharacters}}
\label{SortedCharacters}
}\hfill{\scriptsize (operation)}}\\


 is a list containing the characters \mbox{\texttt{\mdseries\slshape chars}}, ordered as specified by the other arguments. 

 There are three possibilities to sort characters: They can be sorted according
to ascending norms (\mbox{\texttt{\mdseries\slshape flag}} is the string \texttt{"norm"}), to ascending degree (\mbox{\texttt{\mdseries\slshape flag}} is the string \texttt{"degree"}), or both (no third argument is given), i.e., characters with same norm are
sorted according to ascending degree, and characters with smaller norm precede
those with bigger norm. 

 Rational characters in the result precede other ones with same norm and/or
same degree. 

 The trivial character, if contained in \mbox{\texttt{\mdseries\slshape chars}}, will always be sorted to the first position. }

 

\subsection{\textcolor{Chapter }{CharacterTableWithSortedClasses}}
\logpage{[ 71, 21, 3 ]}\nobreak
\hyperdef{L}{X7E3DE0A47E62BE6B}{}
{\noindent\textcolor{FuncColor}{$\triangleright$\enspace\texttt{CharacterTableWithSortedClasses({\mdseries\slshape tbl[, flag]})\index{CharacterTableWithSortedClasses@\texttt{CharacterTableWithSortedClasses}}
\label{CharacterTableWithSortedClasses}
}\hfill{\scriptsize (operation)}}\\


 is a character table obtained by permutation of the classes of \mbox{\texttt{\mdseries\slshape tbl}}. If the second argument \mbox{\texttt{\mdseries\slshape flag}} is the string \texttt{"centralizers"} then the classes of the result are sorted according to descending centralizer
orders. If the second argument is the string \texttt{"representatives"} then the classes of the result are sorted according to ascending
representative orders. If no second argument is given then the classes of the
result are sorted according to ascending representative orders, and classes
with equal representative orders are sorted according to descending
centralizer orders. 

 If the second argument is a permutation then the classes of the result are
sorted by application of this permutation. 

 The result has all those attributes and properties of \mbox{\texttt{\mdseries\slshape tbl}} that are stored in \texttt{SupportedCharacterTableInfo} (\ref{SupportedCharacterTableInfo}) and do not depend on the ordering of classes. }

 

\subsection{\textcolor{Chapter }{SortedCharacterTable (w.r.t. a normal subgroup)}}
\logpage{[ 71, 21, 4 ]}\nobreak
\hyperdef{L}{X82DCAAA882416E24}{}
{\noindent\textcolor{FuncColor}{$\triangleright$\enspace\texttt{SortedCharacterTable({\mdseries\slshape tbl, kernel})\index{SortedCharacterTable@\texttt{SortedCharacterTable}!w.r.t. a normal subgroup}
\label{SortedCharacterTable:w.r.t. a normal subgroup}
}\hfill{\scriptsize (function)}}\\
\noindent\textcolor{FuncColor}{$\triangleright$\enspace\texttt{SortedCharacterTable({\mdseries\slshape tbl, normalseries})\index{SortedCharacterTable@\texttt{SortedCharacterTable}!w.r.t. a series of normal subgroups}
\label{SortedCharacterTable:w.r.t. a series of normal subgroups}
}\hfill{\scriptsize (function)}}\\
\noindent\textcolor{FuncColor}{$\triangleright$\enspace\texttt{SortedCharacterTable({\mdseries\slshape tbl, facttbl, kernel})\index{SortedCharacterTable@\texttt{SortedCharacterTable}!relative to the table of a factor group}
\label{SortedCharacterTable:relative to the table of a factor group}
}\hfill{\scriptsize (function)}}\\


 is a character table obtained on permutation of the classes and the
irreducibles characters of \mbox{\texttt{\mdseries\slshape tbl}}. 

 The first form sorts the classes at positions contained in the list \mbox{\texttt{\mdseries\slshape kernel}} to the beginning, and sorts all characters in \texttt{Irr( \mbox{\texttt{\mdseries\slshape tbl}} )} such that the first characters are those that contain \mbox{\texttt{\mdseries\slshape kernel}} in their kernel. 

 The second form does the same successively for all kernels $k_i$ in the list $\mbox{\texttt{\mdseries\slshape normalseries}} = [ k_1, k_2, \ldots, k_n ]$ where $k_i$ must be a sublist of $k_{{i+1}}$ for $1 \leq i \leq n-1$. 

 The third form computes the table $F$ of the factor group of \mbox{\texttt{\mdseries\slshape tbl}} modulo the normal subgroup formed by the classes whose positions are contained
in the list \mbox{\texttt{\mdseries\slshape kernel}}; $F$ must be permutation equivalent to the table \mbox{\texttt{\mdseries\slshape facttbl}}, in the sense of \texttt{TransformingPermutationsCharacterTables} (\ref{TransformingPermutationsCharacterTables}), otherwise \texttt{fail} is returned. The classes of \mbox{\texttt{\mdseries\slshape tbl}} are sorted such that the preimages of a class of $F$ are consecutive, and that the succession of preimages is that of \mbox{\texttt{\mdseries\slshape facttbl}}. The \texttt{Irr} (\ref{Irr:for a character table}) value of \mbox{\texttt{\mdseries\slshape tbl}} is sorted as with \texttt{SortedCharacterTable( \mbox{\texttt{\mdseries\slshape tbl}}, \mbox{\texttt{\mdseries\slshape kernel}} )}. 

 (\emph{Note} that the transformation is only unique up to table automorphisms of $F$, and this need not be unique up to table automorphisms of \mbox{\texttt{\mdseries\slshape tbl}}.) 

 All rearrangements of classes and characters are stable, i.e., the relative
positions of classes and characters that are not distinguished by any relevant
property is not changed. 

 The result has all those attributes and properties of \mbox{\texttt{\mdseries\slshape tbl}} that are stored in \texttt{SupportedCharacterTableInfo} (\ref{SupportedCharacterTableInfo}) and do not depend on the ordering of classes and characters. 

 The \texttt{ClassPermutation} (\ref{ClassPermutation}) value of \mbox{\texttt{\mdseries\slshape tbl}} is changed if necessary, see{\nobreakspace}\ref{Conventions for Character Tables}. 

 \texttt{SortedCharacterTable} uses \texttt{CharacterTableWithSortedClasses} (\ref{CharacterTableWithSortedClasses}) and \texttt{CharacterTableWithSortedCharacters} (\ref{CharacterTableWithSortedCharacters}). }

 

\subsection{\textcolor{Chapter }{ClassPermutation}}
\logpage{[ 71, 21, 5 ]}\nobreak
\hyperdef{L}{X8099FEDC7DE03AEE}{}
{\noindent\textcolor{FuncColor}{$\triangleright$\enspace\texttt{ClassPermutation({\mdseries\slshape tbl})\index{ClassPermutation@\texttt{ClassPermutation}}
\label{ClassPermutation}
}\hfill{\scriptsize (attribute)}}\\


 is a permutation $\pi$ of classes of the character table \mbox{\texttt{\mdseries\slshape tbl}}. If it is stored then class fusions into \mbox{\texttt{\mdseries\slshape tbl}} that are stored on other tables must be followed by $\pi$ in order to describe the correct fusion. 

 This attribute value is bound only if \mbox{\texttt{\mdseries\slshape tbl}} was obtained from another table by permuting the classes, using \texttt{CharacterTableWithSortedClasses} (\ref{CharacterTableWithSortedClasses}) or \texttt{SortedCharacterTable} (\ref{SortedCharacterTable:w.r.t. a normal subgroup}). 

 It is necessary because the original table and the sorted table have the same
identifier (and the same group if known), and hence the same fusions are valid
for the two tables. 

 
\begin{Verbatim}[commandchars=!@|,fontsize=\small,frame=single,label=Example]
  !gapprompt@gap>| !gapinput@tbl:= CharacterTable( "Symmetric", 4 );|
  CharacterTable( "Sym(4)" )
  !gapprompt@gap>| !gapinput@Display( tbl );|
  Sym(4)
  
       2  3  2  3  .  2
       3  1  .  .  1  .
  
         1a 2a 2b 3a 4a
      2P 1a 1a 1a 3a 2b
      3P 1a 2a 2b 1a 4a
  
  X.1     1 -1  1  1 -1
  X.2     3 -1 -1  .  1
  X.3     2  .  2 -1  .
  X.4     3  1 -1  . -1
  X.5     1  1  1  1  1
  !gapprompt@gap>| !gapinput@srt1:= CharacterTableWithSortedCharacters( tbl );|
  CharacterTable( "Sym(4)" )
  !gapprompt@gap>| !gapinput@List( Irr( srt1 ), Degree );|
  [ 1, 1, 2, 3, 3 ]
  !gapprompt@gap>| !gapinput@srt2:= CharacterTableWithSortedClasses( tbl );|
  CharacterTable( "Sym(4)" )
  !gapprompt@gap>| !gapinput@SizesCentralizers( tbl );|
  [ 24, 4, 8, 3, 4 ]
  !gapprompt@gap>| !gapinput@SizesCentralizers( srt2 );|
  [ 24, 8, 4, 3, 4 ]
  !gapprompt@gap>| !gapinput@nsg:= ClassPositionsOfNormalSubgroups( tbl );|
  [ [ 1 ], [ 1, 3 ], [ 1, 3, 4 ], [ 1 .. 5 ] ]
  !gapprompt@gap>| !gapinput@srt3:= SortedCharacterTable( tbl, nsg );|
  CharacterTable( "Sym(4)" )
  !gapprompt@gap>| !gapinput@nsg:= ClassPositionsOfNormalSubgroups( srt3 );|
  [ [ 1 ], [ 1, 2 ], [ 1 .. 3 ], [ 1 .. 5 ] ]
  !gapprompt@gap>| !gapinput@Display( srt3 );|
  Sym(4)
  
       2  3  3  .  2  2
       3  1  .  1  .  .
  
         1a 2a 3a 2b 4a
      2P 1a 1a 3a 1a 2a
      3P 1a 2a 1a 2b 4a
  
  X.1     1  1  1  1  1
  X.2     1  1  1 -1 -1
  X.3     2  2 -1  .  .
  X.4     3 -1  . -1  1
  X.5     3 -1  .  1 -1
  !gapprompt@gap>| !gapinput@ClassPermutation( srt3 );|
  (2,4,3)
\end{Verbatim}
 }

 }

  
\section{\textcolor{Chapter }{Automorphisms and Equivalence of Character Tables}}\label{Automorphisms and Equivalence of Character Tables}
\logpage{[ 71, 22, 0 ]}
\hyperdef{L}{X7B0A669484470D09}{}
{
  

\subsection{\textcolor{Chapter }{MatrixAutomorphisms}}
\logpage{[ 71, 22, 1 ]}\nobreak
\hyperdef{L}{X84353BB884AF0365}{}
{\noindent\textcolor{FuncColor}{$\triangleright$\enspace\texttt{MatrixAutomorphisms({\mdseries\slshape mat[, maps, subgroup]})\index{MatrixAutomorphisms@\texttt{MatrixAutomorphisms}}
\label{MatrixAutomorphisms}
}\hfill{\scriptsize (operation)}}\\


 For a matrix \mbox{\texttt{\mdseries\slshape mat}}, \texttt{MatrixAutomorphisms} returns the group of those permutations of the columns of \mbox{\texttt{\mdseries\slshape mat}} that leave the set of rows of \mbox{\texttt{\mdseries\slshape mat}} invariant. 

 If the arguments \mbox{\texttt{\mdseries\slshape maps}} and \mbox{\texttt{\mdseries\slshape subgroup}} are given, only the group of those permutations is constructed that
additionally fix each list in the list \mbox{\texttt{\mdseries\slshape maps}} under pointwise action \texttt{OnTuples} (\ref{OnTuples}), and \mbox{\texttt{\mdseries\slshape subgroup}} is a permutation group that is known to be a subgroup of this group of
automorphisms. 

 Each entry in \mbox{\texttt{\mdseries\slshape maps}} must be a list of same length as the rows of \mbox{\texttt{\mdseries\slshape mat}}. For example, if \mbox{\texttt{\mdseries\slshape mat}} is a list of irreducible characters of a group then the list of element orders
of the conjugacy classes (see{\nobreakspace}\texttt{OrdersClassRepresentatives} (\ref{OrdersClassRepresentatives})) may be an entry in \mbox{\texttt{\mdseries\slshape maps}}. }

 

\subsection{\textcolor{Chapter }{TableAutomorphisms}}
\logpage{[ 71, 22, 2 ]}\nobreak
\hyperdef{L}{X8082DD827C673138}{}
{\noindent\textcolor{FuncColor}{$\triangleright$\enspace\texttt{TableAutomorphisms({\mdseries\slshape tbl, characters[, info]})\index{TableAutomorphisms@\texttt{TableAutomorphisms}}
\label{TableAutomorphisms}
}\hfill{\scriptsize (operation)}}\\


 \texttt{TableAutomorphisms} returns the permutation group of those matrix automorphisms
(see{\nobreakspace}\texttt{MatrixAutomorphisms} (\ref{MatrixAutomorphisms})) of the list \mbox{\texttt{\mdseries\slshape characters}} that leave the element orders (see{\nobreakspace}\texttt{OrdersClassRepresentatives} (\ref{OrdersClassRepresentatives})) and all stored power maps (see{\nobreakspace}\texttt{ComputedPowerMaps} (\ref{ComputedPowerMaps})) of the character table \mbox{\texttt{\mdseries\slshape tbl}} invariant. 

 If \mbox{\texttt{\mdseries\slshape characters}} is closed under Galois conjugacy {\textendash}this is always fulfilled for the
list of all irreducible characters of ordinary character tables{\textendash}
the string \texttt{"closed"} may be entered as the third argument \mbox{\texttt{\mdseries\slshape info}}. Alternatively, a known subgroup of the table automorphisms can be entered as
the third argument \mbox{\texttt{\mdseries\slshape info}}. 

 The attribute \texttt{AutomorphismsOfTable} (\ref{AutomorphismsOfTable}) can be used to compute and store the table automorphisms for the case that \mbox{\texttt{\mdseries\slshape characters}} equals the \texttt{Irr} (\ref{Irr:for a character table}) value of \mbox{\texttt{\mdseries\slshape tbl}}. 

 
\begin{Verbatim}[commandchars=!@|,fontsize=\small,frame=single,label=Example]
  !gapprompt@gap>| !gapinput@tbld8:= CharacterTable( "Dihedral", 8 );;|
  !gapprompt@gap>| !gapinput@irrd8:= Irr( tbld8 );|
  [ Character( CharacterTable( "Dihedral(8)" ), [ 1, 1, 1, 1, 1 ] ),
    Character( CharacterTable( "Dihedral(8)" ), [ 1, 1, 1, -1, -1 ] ),
    Character( CharacterTable( "Dihedral(8)" ), [ 1, -1, 1, 1, -1 ] ),
    Character( CharacterTable( "Dihedral(8)" ), [ 1, -1, 1, -1, 1 ] ),
    Character( CharacterTable( "Dihedral(8)" ), [ 2, 0, -2, 0, 0 ] ) ]
  !gapprompt@gap>| !gapinput@orders:= OrdersClassRepresentatives( tbld8 );|
  [ 1, 4, 2, 2, 2 ]
  !gapprompt@gap>| !gapinput@MatrixAutomorphisms( irrd8 );|
  Group([ (4,5), (2,4) ])
  !gapprompt@gap>| !gapinput@MatrixAutomorphisms( irrd8, [ orders ], Group( () ) );|
  Group([ (4,5) ])
  !gapprompt@gap>| !gapinput@TableAutomorphisms( tbld8, irrd8 );|
  Group([ (4,5) ])
\end{Verbatim}
 }

 

\subsection{\textcolor{Chapter }{TransformingPermutations}}
\logpage{[ 71, 22, 3 ]}\nobreak
\hyperdef{L}{X7D721E3D7AA319F5}{}
{\noindent\textcolor{FuncColor}{$\triangleright$\enspace\texttt{TransformingPermutations({\mdseries\slshape mat1, mat2})\index{TransformingPermutations@\texttt{TransformingPermutations}}
\label{TransformingPermutations}
}\hfill{\scriptsize (operation)}}\\


 Let \mbox{\texttt{\mdseries\slshape mat1}} and \mbox{\texttt{\mdseries\slshape mat2}} be matrices. \texttt{TransformingPermutations} tries to construct a permutation $\pi$ that transforms the set of rows of the matrix \mbox{\texttt{\mdseries\slshape mat1}} to the set of rows of the matrix \mbox{\texttt{\mdseries\slshape mat2}} by permuting the columns. 

 If such a permutation exists, a record with the components \texttt{columns}, \texttt{rows}, and \texttt{group} is returned, otherwise \texttt{fail}. For \texttt{TransformingPermutations( \mbox{\texttt{\mdseries\slshape mat1}}, \mbox{\texttt{\mdseries\slshape mat2}} ) = \mbox{\texttt{\mdseries\slshape r}}} $\neq$ \texttt{fail}, we have \texttt{\mbox{\texttt{\mdseries\slshape mat2}} = Permuted( List( \mbox{\texttt{\mdseries\slshape mat1}}, x \texttt{\symbol{45}}{\textgreater} Permuted( x, \mbox{\texttt{\mdseries\slshape r}}.columns ) ), \mbox{\texttt{\mdseries\slshape r}}.rows )}. 

 \mbox{\texttt{\mdseries\slshape r}}\texttt{.group} is the group of matrix automorphisms of \mbox{\texttt{\mdseries\slshape mat2}} (see{\nobreakspace}\texttt{MatrixAutomorphisms} (\ref{MatrixAutomorphisms})). This group stabilizes the transformation in the sense that applying any of
its elements to the columns of \mbox{\texttt{\mdseries\slshape mat2}} preserves the set of rows of \mbox{\texttt{\mdseries\slshape mat2}}. }

 

\subsection{\textcolor{Chapter }{TransformingPermutationsCharacterTables}}
\logpage{[ 71, 22, 4 ]}\nobreak
\hyperdef{L}{X849731AA7EC9FA73}{}
{\noindent\textcolor{FuncColor}{$\triangleright$\enspace\texttt{TransformingPermutationsCharacterTables({\mdseries\slshape tbl1, tbl2})\index{TransformingPermutationsCharacterTables@\texttt{Transforming}\-\texttt{Permutations}\-\texttt{Character}\-\texttt{Tables}}
\label{TransformingPermutationsCharacterTables}
}\hfill{\scriptsize (operation)}}\\


 Let \mbox{\texttt{\mdseries\slshape tbl1}} and \mbox{\texttt{\mdseries\slshape tbl2}} be character tables. \texttt{TransformingPermutationsCharacterTables} tries to construct a permutation $\pi$ that transforms the set of rows of the matrix \texttt{Irr( \mbox{\texttt{\mdseries\slshape tbl1}} )} to the set of rows of the matrix \texttt{Irr( \mbox{\texttt{\mdseries\slshape tbl2}} )} by permuting the columns (see{\nobreakspace}\texttt{TransformingPermutations} (\ref{TransformingPermutations})), such that $\pi$ transforms also the power maps and the element orders. 

 If such a permutation $\pi$ exists then a record with the components \texttt{columns} ($\pi$), \texttt{rows} (the permutation of \texttt{Irr( \mbox{\texttt{\mdseries\slshape tbl1}} )} corresponding to $\pi$), and \texttt{group} (the permutation group of table automorphisms of \mbox{\texttt{\mdseries\slshape tbl2}}, see{\nobreakspace}\texttt{AutomorphismsOfTable} (\ref{AutomorphismsOfTable})) is returned. If no such permutation exists, \texttt{fail} is returned. 

 
\begin{Verbatim}[commandchars=!@|,fontsize=\small,frame=single,label=Example]
  !gapprompt@gap>| !gapinput@tblq8:= CharacterTable( "Quaternionic", 8 );;|
  !gapprompt@gap>| !gapinput@irrq8:= Irr( tblq8 );|
  [ Character( CharacterTable( "Q8" ), [ 1, 1, 1, 1, 1 ] ),
    Character( CharacterTable( "Q8" ), [ 1, 1, 1, -1, -1 ] ),
    Character( CharacterTable( "Q8" ), [ 1, -1, 1, 1, -1 ] ),
    Character( CharacterTable( "Q8" ), [ 1, -1, 1, -1, 1 ] ),
    Character( CharacterTable( "Q8" ), [ 2, 0, -2, 0, 0 ] ) ]
  !gapprompt@gap>| !gapinput@OrdersClassRepresentatives( tblq8 );|
  [ 1, 4, 2, 4, 4 ]
  !gapprompt@gap>| !gapinput@TransformingPermutations( irrd8, irrq8 );|
  rec( columns := (), group := Group([ (4,5), (2,4) ]), rows := () )
  !gapprompt@gap>| !gapinput@TransformingPermutationsCharacterTables( tbld8, tblq8 );|
  fail
  !gapprompt@gap>| !gapinput@tbld6:= CharacterTable( "Dihedral", 6 );;|
  !gapprompt@gap>| !gapinput@tbls3:= CharacterTable( "Symmetric", 3 );;|
  !gapprompt@gap>| !gapinput@TransformingPermutationsCharacterTables( tbld6, tbls3 );|
  rec( columns := (2,3), group := Group(()), rows := (1,3,2) )
\end{Verbatim}
 }

 

\subsection{\textcolor{Chapter }{FamiliesOfRows}}
\logpage{[ 71, 22, 5 ]}\nobreak
\hyperdef{L}{X8117D940835B0B47}{}
{\noindent\textcolor{FuncColor}{$\triangleright$\enspace\texttt{FamiliesOfRows({\mdseries\slshape mat, maps})\index{FamiliesOfRows@\texttt{FamiliesOfRows}}
\label{FamiliesOfRows}
}\hfill{\scriptsize (function)}}\\


 distributes the rows of the matrix \mbox{\texttt{\mdseries\slshape mat}} into families, as follows. Two rows of \mbox{\texttt{\mdseries\slshape mat}} belong to the same family if there is a permutation of columns that maps one
row to the other row. Each entry in the list \mbox{\texttt{\mdseries\slshape maps}} is regarded to form a family of length 1. 

 \texttt{FamiliesOfRows} returns a record with the components 
\begin{description}
\item[{\texttt{famreps}}]  the list of representatives for each family, 
\item[{\texttt{permutations}}]  the list that contains at position $i$ a list of permutations that map the members of the family with representative \texttt{famreps}$[i]$ to that representative, 
\item[{\texttt{families}}]  the list that contains at position $i$ the list of positions of members of the family of representative \texttt{famreps}$[i]$; (for the element \mbox{\texttt{\mdseries\slshape maps}}$[i]$ the only member of the family will get the number \texttt{Length( \mbox{\texttt{\mdseries\slshape mat}} ) + }$i$). 
\end{description}
 }

 }

  
\section{\textcolor{Chapter }{Storing Normal Subgroup Information}}\label{Storing Normal Subgroup Information}
\logpage{[ 71, 23, 0 ]}
\hyperdef{L}{X81272CEE79F13E7B}{}
{
  

\subsection{\textcolor{Chapter }{NormalSubgroupClassesInfo}}
\logpage{[ 71, 23, 1 ]}\nobreak
\hyperdef{L}{X7E66174C7C7A8C0C}{}
{\noindent\textcolor{FuncColor}{$\triangleright$\enspace\texttt{NormalSubgroupClassesInfo({\mdseries\slshape tbl})\index{NormalSubgroupClassesInfo@\texttt{NormalSubgroupClassesInfo}}
\label{NormalSubgroupClassesInfo}
}\hfill{\scriptsize (attribute)}}\\


 Let \mbox{\texttt{\mdseries\slshape tbl}} be the ordinary character table of the group $G$. Many computations for group characters of $G$ involve computations in normal subgroups or factor groups of $G$. 

 In some cases the character table \mbox{\texttt{\mdseries\slshape tbl}} is sufficient; for example questions about a normal subgroup $N$ of $G$ can be answered if one knows the conjugacy classes that form $N$, e.g., the question whether a character of $G$ restricts irreducibly to $N$. But other questions require the computation of $N$ or even more information, like the character table of $N$. 

 In order to do these computations only once, one stores in the group a record
with components to store normal subgroups, the corresponding lists of
conjugacy classes, and (if necessary) the factor groups, namely 

 
\begin{description}
\item[{\texttt{nsg}}]  list of normal subgroups of $G$, may be incomplete, 
\item[{\texttt{nsgclasses}}]  at position $i$, the list of positions of conjugacy classes of \mbox{\texttt{\mdseries\slshape tbl}} forming the $i$\texttt{\symbol{45}}th entry of the \texttt{nsg} component, 
\item[{\texttt{nsgfactors}}]  at position $i$, if bound, the factor group modulo the $i$\texttt{\symbol{45}}th entry of the \texttt{nsg} component. 
\end{description}
 

 \texttt{NormalSubgroupClasses} (\ref{NormalSubgroupClasses}), \texttt{FactorGroupNormalSubgroupClasses} (\ref{FactorGroupNormalSubgroupClasses}), and \texttt{ClassPositionsOfNormalSubgroup} (\ref{ClassPositionsOfNormalSubgroup}) each use these components, and they are the only functions to do so. 

 So if you need information about a normal subgroup for that you know the
conjugacy classes, you should get it using \texttt{NormalSubgroupClasses} (\ref{NormalSubgroupClasses}). If the normal subgroup was already used it is just returned, with all the
knowledge it contains. Otherwise the normal subgroup is added to the lists,
and will be available for the next call. 

 For example, if you are dealing with kernels of characters using the \texttt{KernelOfCharacter} (\ref{KernelOfCharacter}) function you make use of this feature because \texttt{KernelOfCharacter} (\ref{KernelOfCharacter}) calls \texttt{NormalSubgroupClasses} (\ref{NormalSubgroupClasses}). }

 

\subsection{\textcolor{Chapter }{ClassPositionsOfNormalSubgroup}}
\logpage{[ 71, 23, 2 ]}\nobreak
\hyperdef{L}{X7C2A87E085111090}{}
{\noindent\textcolor{FuncColor}{$\triangleright$\enspace\texttt{ClassPositionsOfNormalSubgroup({\mdseries\slshape tbl, N})\index{ClassPositionsOfNormalSubgroup@\texttt{ClassPositionsOfNormalSubgroup}}
\label{ClassPositionsOfNormalSubgroup}
}\hfill{\scriptsize (function)}}\\


 is the list of positions of conjugacy classes of the character table \mbox{\texttt{\mdseries\slshape tbl}} that are contained in the normal subgroup \mbox{\texttt{\mdseries\slshape N}} of the underlying group of \mbox{\texttt{\mdseries\slshape tbl}}. }

 

\subsection{\textcolor{Chapter }{NormalSubgroupClasses}}
\logpage{[ 71, 23, 3 ]}\nobreak
\hyperdef{L}{X87E7391F7F92377C}{}
{\noindent\textcolor{FuncColor}{$\triangleright$\enspace\texttt{NormalSubgroupClasses({\mdseries\slshape tbl, classes})\index{NormalSubgroupClasses@\texttt{NormalSubgroupClasses}}
\label{NormalSubgroupClasses}
}\hfill{\scriptsize (function)}}\\


 returns the normal subgroup of the underlying group $G$ of the ordinary character table \mbox{\texttt{\mdseries\slshape tbl}} that consists of those conjugacy classes of \mbox{\texttt{\mdseries\slshape tbl}} whose positions are in the list \mbox{\texttt{\mdseries\slshape classes}}. 

 If \texttt{NormalSubgroupClassesInfo( \mbox{\texttt{\mdseries\slshape tbl}} ).nsg} does not yet contain the required normal subgroup, and if \texttt{NormalSubgroupClassesInfo( \mbox{\texttt{\mdseries\slshape tbl}} ).normalSubgroups} is bound then the result will be identical to the group in \texttt{NormalSubgroupClassesInfo( \mbox{\texttt{\mdseries\slshape tbl}} ).normalSubgroups}. }

 

\subsection{\textcolor{Chapter }{FactorGroupNormalSubgroupClasses}}
\logpage{[ 71, 23, 4 ]}\nobreak
\hyperdef{L}{X79D451F0808EB252}{}
{\noindent\textcolor{FuncColor}{$\triangleright$\enspace\texttt{FactorGroupNormalSubgroupClasses({\mdseries\slshape tbl, classes})\index{FactorGroupNormalSubgroupClasses@\texttt{FactorGroupNormalSubgroupClasses}}
\label{FactorGroupNormalSubgroupClasses}
}\hfill{\scriptsize (function)}}\\


 is the factor group of the underlying group $G$ of the ordinary character table \mbox{\texttt{\mdseries\slshape tbl}} modulo the normal subgroup of $G$ that consists of those conjugacy classes of \mbox{\texttt{\mdseries\slshape tbl}} whose positions are in the list \mbox{\texttt{\mdseries\slshape classes}}. 

 
\begin{Verbatim}[commandchars=!@|,fontsize=\small,frame=single,label=Example]
  !gapprompt@gap>| !gapinput@g:= SymmetricGroup( 4 );|
  Sym( [ 1 .. 4 ] )
  !gapprompt@gap>| !gapinput@SetName( g, "S4" );|
  !gapprompt@gap>| !gapinput@tbl:= CharacterTable( g );|
  CharacterTable( S4 )
  !gapprompt@gap>| !gapinput@irr:= Irr( g );|
  [ Character( CharacterTable( S4 ), [ 1, -1, 1, 1, -1 ] ),
    Character( CharacterTable( S4 ), [ 3, -1, -1, 0, 1 ] ),
    Character( CharacterTable( S4 ), [ 2, 0, 2, -1, 0 ] ),
    Character( CharacterTable( S4 ), [ 3, 1, -1, 0, -1 ] ),
    Character( CharacterTable( S4 ), [ 1, 1, 1, 1, 1 ] ) ]
  !gapprompt@gap>| !gapinput@kernel:= KernelOfCharacter( irr[3] );;|
  !gapprompt@gap>| !gapinput@AsSet(kernel);|
  [ (), (1,2)(3,4), (1,3)(2,4), (1,4)(2,3) ]
  !gapprompt@gap>| !gapinput@SetName(kernel,"V4");|
  !gapprompt@gap>| !gapinput@HasNormalSubgroupClassesInfo( tbl );|
  true
  !gapprompt@gap>| !gapinput@NormalSubgroupClassesInfo( tbl );|
  rec( nsg := [ V4 ], nsgclasses := [ [ 1, 3 ] ], nsgfactors := [  ] )
  !gapprompt@gap>| !gapinput@ClassPositionsOfNormalSubgroup( tbl, kernel );|
  [ 1, 3 ]
  !gapprompt@gap>| !gapinput@G := FactorGroupNormalSubgroupClasses( tbl, [ 1, 3 ] );;|
  !gapprompt@gap>| !gapinput@NormalSubgroupClassesInfo( tbl ).nsgfactors[1] = G;|
  true
\end{Verbatim}
 }

 }

 }

             
\chapter{\textcolor{Chapter }{Class Functions}}\label{Class Functions}
\logpage{[ 72, 0, 0 ]}
\hyperdef{L}{X7C91D0D17850E564}{}
{
  \index{characters} \index{group characters} \index{virtual characters} \index{generalized characters} This chapter describes operations for \emph{class functions of finite groups}. For operations concerning \emph{character tables}, see Chapter{\nobreakspace}\ref{Character Tables}. 

 Several examples in this chapter require the \textsf{GAP} Character Table Library to be available. If it is not yet loaded then we load
it now. 

 
\begin{Verbatim}[commandchars=!@|,fontsize=\small,frame=single,label=Example]
  !gapprompt@gap>| !gapinput@LoadPackage( "ctbllib" );|
  true
\end{Verbatim}
  
\section{\textcolor{Chapter }{Why Class Functions?}}\label{Why Class Functions?}
\logpage{[ 72, 1, 0 ]}
\hyperdef{L}{X823319217E4B6852}{}
{
  In principle it is possible to represent group characters or more general
class functions by the plain lists of their values, and in fact many
operations for class functions work with plain lists of class function values.
But this has two disadvantages. 

 First, it is then necessary to regard a values list explicitly as a class
function of a particular character table, by supplying this character table as
an argument. In practice this means that with this setup, the user has the
task to put the objects into the right context. For example, forming the
scalar product or the tensor product of two class functions or forming an
induced class function or a conjugate class function then needs three
arguments in this case; this is particularly inconvenient in cases where infix
operations cannot be used because of the additional argument, as for tensor
products and induced class functions. 

 Second, when one says that ``$\chi$ is a character of a group $G$'' then this object $\chi$ carries a lot of information. $\chi$ has certain properties such as being irreducible or not. Several subgroups of $G$ are related to $\chi$, such as the kernel and the centre of $\chi$. Other attributes of characters are the determinant and the central
character. This knowledge cannot be stored in a plain list. 

 For dealing with a group together with its characters, and maybe also
subgroups and their characters, it is desirable that \textsf{GAP} keeps track of the interpretation of characters. On the other hand, for using
characters without accessing their groups, such as characters of tables from
the \textsf{GAP} table library, dealing just with values lists is often sufficient. In
particular, if one deals with incomplete character tables then it is often
necessary to specify the arguments explicitly, for example one has to choose a
fusion map or power map from a set of possibilities. 

 The main idea behind class function objects is that a class function object is
equal to its values list in the sense of \texttt{\texttt{\symbol{92}}=} (\ref{bSlash=}), so class function objects can be used wherever their values lists can be
used, but there are operations for class function objects that do not work
just with values lists.     \textsf{GAP} library functions prefer to return class function objects rather than
returning just values lists, for example \texttt{Irr} (\ref{Irr:for a group}) lists consist of class function objects, and \texttt{TrivialCharacter} (\ref{TrivialCharacter:for a group}) returns a class function object. 

 Here is an \emph{example} that shows both approaches. First we define some groups. 

 
\begin{Verbatim}[commandchars=!@|,fontsize=\small,frame=single,label=Example]
  !gapprompt@gap>| !gapinput@S4:= SymmetricGroup( 4 );;  SetName( S4, "S4" );|
  !gapprompt@gap>| !gapinput@D8:= SylowSubgroup( S4, 2 );; SetName( D8, "D8" );|
\end{Verbatim}
 

 We do some computations using the functions described later in this Chapter,
first with class function objects. 

 
\begin{Verbatim}[commandchars=!@|,fontsize=\small,frame=single,label=Example]
  !gapprompt@gap>| !gapinput@irrS4:= Irr( S4 );;|
  !gapprompt@gap>| !gapinput@irrD8:= Irr( D8 );;|
  !gapprompt@gap>| !gapinput@chi:= irrD8[4];|
  Character( CharacterTable( D8 ), [ 1, -1, 1, -1, 1 ] )
  !gapprompt@gap>| !gapinput@chi * chi;|
  Character( CharacterTable( D8 ), [ 1, 1, 1, 1, 1 ] )
  !gapprompt@gap>| !gapinput@ind:= chi ^ S4;|
  Character( CharacterTable( S4 ), [ 3, -1, -1, 0, 1 ] )
  !gapprompt@gap>| !gapinput@List( irrS4, x -> ScalarProduct( x, ind ) );|
  [ 0, 1, 0, 0, 0 ]
  !gapprompt@gap>| !gapinput@det:= Determinant( ind );|
  Character( CharacterTable( S4 ), [ 1, 1, 1, 1, 1 ] )
  !gapprompt@gap>| !gapinput@cent:= CentralCharacter( ind );|
  ClassFunction( CharacterTable( S4 ), [ 1, -2, -1, 0, 2 ] )
  !gapprompt@gap>| !gapinput@rest:= Restricted( cent, D8 );|
  ClassFunction( CharacterTable( D8 ), [ 1, -2, -1, -1, 2 ] )
\end{Verbatim}
 

 Now we repeat these calculations with plain lists of character values. Here we
need the character tables in some places. 

 
\begin{Verbatim}[commandchars=!@|,fontsize=\small,frame=single,label=Example]
  !gapprompt@gap>| !gapinput@tS4:= CharacterTable( S4 );;|
  !gapprompt@gap>| !gapinput@tD8:= CharacterTable( D8 );;|
  !gapprompt@gap>| !gapinput@chi:= ValuesOfClassFunction( irrD8[4] );|
  [ 1, -1, 1, -1, 1 ]
  !gapprompt@gap>| !gapinput@Tensored( [ chi ], [ chi ] )[1];|
  [ 1, 1, 1, 1, 1 ]
  !gapprompt@gap>| !gapinput@ind:= InducedClassFunction( tD8, chi, tS4 );|
  ClassFunction( CharacterTable( S4 ), [ 3, -1, -1, 0, 1 ] )
  !gapprompt@gap>| !gapinput@List( Irr( tS4 ), x -> ScalarProduct( tS4, x, ind ) );|
  [ 0, 1, 0, 0, 0 ]
  !gapprompt@gap>| !gapinput@det:= DeterminantOfCharacter( tS4, ind );|
  ClassFunction( CharacterTable( S4 ), [ 1, 1, 1, 1, 1 ] )
  !gapprompt@gap>| !gapinput@cent:= CentralCharacter( tS4, ind );|
  ClassFunction( CharacterTable( S4 ), [ 1, -2, -1, 0, 2 ] )
  !gapprompt@gap>| !gapinput@rest:= Restricted( tS4, cent, tD8 );|
  ClassFunction( CharacterTable( D8 ), [ 1, -2, -1, -1, 2 ] )
\end{Verbatim}
 

 If one deals with character tables from the \textsf{GAP} table library then one has no access to their groups, but often the tables
provide enough information for computing induced or restricted class
functions, symmetrizations etc., because the relevant class fusions and power
maps are often stored on library tables. In these cases it is possible to use
the tables instead of the groups as arguments. (If necessary information is
not uniquely determined by the tables then an error is signalled.) 

 
\begin{Verbatim}[commandchars=!@|,fontsize=\small,frame=single,label=Example]
  !gapprompt@gap>| !gapinput@s5 := CharacterTable( "A5.2" );; irrs5 := Irr( s5  );;|
  !gapprompt@gap>| !gapinput@m11:= CharacterTable( "M11"  );; irrm11:= Irr( m11 );;|
  !gapprompt@gap>| !gapinput@chi:= TrivialCharacter( s5 );|
  Character( CharacterTable( "A5.2" ), [ 1, 1, 1, 1, 1, 1, 1 ] )
  !gapprompt@gap>| !gapinput@chi ^ m11;|
  Character( CharacterTable( "M11" ), [ 66, 10, 3, 2, 1, 1, 0, 0, 0, 0
   ] )
  !gapprompt@gap>| !gapinput@Determinant( irrs5[4] );|
  Character( CharacterTable( "A5.2" ), [ 1, 1, 1, 1, -1, -1, -1 ] )
\end{Verbatim}
 

 Functions that compute \emph{normal} subgroups related to characters have counterparts that return the list of
class positions corresponding to these groups. 

 
\begin{Verbatim}[commandchars=!@|,fontsize=\small,frame=single,label=Example]
  !gapprompt@gap>| !gapinput@ClassPositionsOfKernel( irrs5[2] );|
  [ 1, 2, 3, 4 ]
  !gapprompt@gap>| !gapinput@ClassPositionsOfCentre( irrs5[2] );|
  [ 1, 2, 3, 4, 5, 6, 7 ]
\end{Verbatim}
 

 Non\texttt{\symbol{45}}normal subgroups cannot be described this way, so for
example inertia subgroups (see{\nobreakspace}\texttt{InertiaSubgroup} (\ref{InertiaSubgroup})) can in general not be computed from character tables without access to their
groups. 

\subsection{\textcolor{Chapter }{IsClassFunction}}
\logpage{[ 72, 1, 1 ]}\nobreak
\hyperdef{L}{X7E75A70F7BF00A0D}{}
{\noindent\textcolor{FuncColor}{$\triangleright$\enspace\texttt{IsClassFunction({\mdseries\slshape obj})\index{IsClassFunction@\texttt{IsClassFunction}}
\label{IsClassFunction}
}\hfill{\scriptsize (Category)}}\\


 \index{class function}\index{class function objects} A \emph{class function} (in characteristic $p$) of a finite group $G$ is a map from the set of ($p$\texttt{\symbol{45}}regular) elements in $G$ to the field of cyclotomics that is constant on conjugacy classes of $G$. 

 Each class function in \textsf{GAP} is represented by an \emph{immutable list}, where at the $i$\texttt{\symbol{45}}th position the value on the $i$\texttt{\symbol{45}}th conjugacy class of the character table of $G$ is stored. The ordering of the conjugacy classes is the one used in the
underlying character table. Note that if the character table has access to its
underlying group then the ordering of conjugacy classes in the group and in
the character table may differ (see \ref{The Interface between Character Tables and Groups}); class functions always refer to the ordering of classes in the character
table. 

 \emph{Class function objects} in \textsf{GAP} are not just plain lists, they store the character table of the group $G$ as value of the attribute \texttt{UnderlyingCharacterTable} (\ref{UnderlyingCharacterTable}). The group $G$ itself is accessible only via the character table and thus only if the
character table stores its group, as value of the attribute \texttt{UnderlyingGroup} (\ref{UnderlyingGroup:for character tables}). The reason for this is that many computations with class functions are
possible without using their groups, for example class functions of character
tables in the \textsf{GAP} character table library do in general not have access to their underlying
groups. 

 There are (at least) two reasons why class functions in \textsf{GAP} are \emph{not} implemented as mappings. First, we want to distinguish class functions in
different characteristics, for example to be able to define the Frobenius
character of a given Brauer character; viewed as mappings, the trivial
characters in all characteristics coprime to the order of $G$ are equal. Second, the product of two class functions shall be again a class
function, whereas the product of general mappings is defined as composition. 

 A further argument is that the typical operations for mappings such as \texttt{Image} (\ref{Image:set of images of the source of a general mapping}) and \texttt{PreImage} (\ref{PreImage:set of preimages of the range of a general mapping}) play no important role for class functions. }

 }

  
\section{\textcolor{Chapter }{Basic Operations for Class Functions}}\label{Basic Operations for Class Functions}
\logpage{[ 72, 2, 0 ]}
\hyperdef{L}{X8192EDDB84ADD46E}{}
{
  Basic operations for class functions are \texttt{UnderlyingCharacterTable} (\ref{UnderlyingCharacterTable}), \texttt{ValuesOfClassFunction} (\ref{ValuesOfClassFunction}), and the basic operations for lists (see{\nobreakspace}\ref{Basic Operations for Lists}). 

\subsection{\textcolor{Chapter }{UnderlyingCharacterTable}}
\logpage{[ 72, 2, 1 ]}\nobreak
\hyperdef{L}{X81B55C067D123B76}{}
{\noindent\textcolor{FuncColor}{$\triangleright$\enspace\texttt{UnderlyingCharacterTable({\mdseries\slshape psi})\index{UnderlyingCharacterTable@\texttt{UnderlyingCharacterTable}}
\label{UnderlyingCharacterTable}
}\hfill{\scriptsize (attribute)}}\\


 For a class function \mbox{\texttt{\mdseries\slshape psi}} of a group $G$ the character table of $G$ is stored as value of \texttt{UnderlyingCharacterTable}. The ordering of entries in the list \mbox{\texttt{\mdseries\slshape psi}} (see{\nobreakspace}\texttt{ValuesOfClassFunction} (\ref{ValuesOfClassFunction})) refers to the ordering of conjugacy classes in this character table. 

 If \mbox{\texttt{\mdseries\slshape psi}} is an ordinary class function then the underlying character table is the
ordinary character table of $G$ (see{\nobreakspace}\texttt{OrdinaryCharacterTable} (\ref{OrdinaryCharacterTable:for a group})), if \mbox{\texttt{\mdseries\slshape psi}} is a class function in characteristic $p \neq 0$ then the underlying character table is the $p$\texttt{\symbol{45}}modular Brauer table of $G$ (see{\nobreakspace}\texttt{BrauerTable} (\ref{BrauerTable:for a group, and a prime integer})). So the underlying characteristic of \mbox{\texttt{\mdseries\slshape psi}} can be read off from the underlying character table. }

 

\subsection{\textcolor{Chapter }{ValuesOfClassFunction}}
\logpage{[ 72, 2, 2 ]}\nobreak
\hyperdef{L}{X7FE14712843C6486}{}
{\noindent\textcolor{FuncColor}{$\triangleright$\enspace\texttt{ValuesOfClassFunction({\mdseries\slshape psi})\index{ValuesOfClassFunction@\texttt{ValuesOfClassFunction}}
\label{ValuesOfClassFunction}
}\hfill{\scriptsize (attribute)}}\\


 is the list of values of the class function \mbox{\texttt{\mdseries\slshape psi}}, the $i$\texttt{\symbol{45}}th entry being the value on the $i$\texttt{\symbol{45}}th conjugacy class of the underlying character table
(see{\nobreakspace}\texttt{UnderlyingCharacterTable} (\ref{UnderlyingCharacterTable})). 

 
\begin{Verbatim}[commandchars=!@|,fontsize=\small,frame=single,label=Example]
  !gapprompt@gap>| !gapinput@g:= SymmetricGroup( 4 );|
  Sym( [ 1 .. 4 ] )
  !gapprompt@gap>| !gapinput@psi:= TrivialCharacter( g );|
  Character( CharacterTable( Sym( [ 1 .. 4 ] ) ), [ 1, 1, 1, 1, 1 ] )
  !gapprompt@gap>| !gapinput@UnderlyingCharacterTable( psi );|
  CharacterTable( Sym( [ 1 .. 4 ] ) )
  !gapprompt@gap>| !gapinput@ValuesOfClassFunction( psi );|
  [ 1, 1, 1, 1, 1 ]
  !gapprompt@gap>| !gapinput@IsList( psi );|
  true
  !gapprompt@gap>| !gapinput@psi[1];|
  1
  !gapprompt@gap>| !gapinput@Length( psi );|
  5
  !gapprompt@gap>| !gapinput@IsBound( psi[6] );|
  false
  !gapprompt@gap>| !gapinput@Concatenation( psi, [ 2, 3 ] );|
  [ 1, 1, 1, 1, 1, 2, 3 ]
\end{Verbatim}
 }

 }

  
\section{\textcolor{Chapter }{Comparison of Class Functions}}\label{Comparison of Class Functions}
\logpage{[ 72, 3, 0 ]}
\hyperdef{L}{X829EFBF57FCB1A94}{}
{
  With respect to \texttt{\texttt{\symbol{92}}=} (\ref{bSlash=}) and \texttt{\texttt{\symbol{92}}{\textless}} (\ref{bSlash<}), class functions behave equally to their lists of values (see{\nobreakspace}\texttt{ValuesOfClassFunction} (\ref{ValuesOfClassFunction})). So two class functions are equal if and only if their lists of values are
equal, no matter whether they are class functions of the same character table,
of the same group but w.r.t.{\nobreakspace}different class ordering, or of
different groups. 

 
\begin{Verbatim}[commandchars=!@|,fontsize=\small,frame=single,label=Example]
  !gapprompt@gap>| !gapinput@grps:= Filtered( AllSmallGroups( 8 ), g -> not IsAbelian( g ) );|
  [ <pc group of size 8 with 3 generators>,
    <pc group of size 8 with 3 generators> ]
  !gapprompt@gap>| !gapinput@t1:= CharacterTable( grps[1] );  SetName( t1, "t1" );|
  CharacterTable( <pc group of size 8 with 3 generators> )
  !gapprompt@gap>| !gapinput@t2:= CharacterTable( grps[2] );  SetName( t2, "t2" );|
  CharacterTable( <pc group of size 8 with 3 generators> )
  !gapprompt@gap>| !gapinput@CharacterDegrees( t1 );|
  [ [ 1, 4 ], [ 2, 1 ] ]
  !gapprompt@gap>| !gapinput@irr1:= Irr( grps[1] );;|
  !gapprompt@gap>| !gapinput@irr2:= Irr( grps[2] );;|
  !gapprompt@gap>| !gapinput@irr1 = irr2;|
  true
  !gapprompt@gap>| !gapinput@irr1[1];|
  Character( t1, [ 1, 1, 1, 1, 1 ] )
  !gapprompt@gap>| !gapinput@irr1[5];|
  Character( t1, [ 2, 0, 0, -2, 0 ] )
  !gapprompt@gap>| !gapinput@irr2[1];|
  Character( t2, [ 1, 1, 1, 1, 1 ] )
  !gapprompt@gap>| !gapinput@IsSSortedList( irr1 );|
  false
  !gapprompt@gap>| !gapinput@irr1[1] < irr1[2];  # the triv. character has no '-1'|
  false
  !gapprompt@gap>| !gapinput@irr1[2] < irr1[5];|
  true
\end{Verbatim}
 }

  
\section{\textcolor{Chapter }{Arithmetic Operations for Class Functions}}\label{Arithmetic Operations for Class Functions}
\logpage{[ 72, 4, 0 ]}
\hyperdef{L}{X83B9F0C5871A5F7F}{}
{
  Class functions are \emph{row vectors} of cyclotomics. The \emph{additive} behaviour of class functions is defined such that they are equal to the plain
lists of class function values except that the results are represented again
as class functions whenever this makes sense. The \emph{multiplicative} behaviour, however, is different. This is motivated by the fact that the
tensor product of class functions is a more interesting operation than the
vector product of plain lists. (Another candidate for a multiplication of
compatible class functions would have been the inner product, which is
implemented via the function \texttt{ScalarProduct} (\ref{ScalarProduct:for characters}). In terms of filters, the arithmetic of class functions is based on the
decision that they lie in \texttt{IsGeneralizedRowVector} (\ref{IsGeneralizedRowVector}), with additive nesting depth $1$, but they do \emph{not} lie in \texttt{IsMultiplicativeGeneralizedRowVector} (\ref{IsMultiplicativeGeneralizedRowVector}). 

 More specifically, the scalar multiple of a class function with a cyclotomic
is a class function, and the sum and the difference of two class functions of
the same underlying character table (see{\nobreakspace}\texttt{UnderlyingCharacterTable} (\ref{UnderlyingCharacterTable})) are again class functions of this table. The sum and the difference of a
class function and a list that is \emph{not} a class function are plain lists, as well as the sum and the difference of two
class functions of different character tables. 

 
\begin{Verbatim}[commandchars=!@|,fontsize=\small,frame=single,label=Example]
  !gapprompt@gap>| !gapinput@g:= SymmetricGroup( 4 );;  tbl:= CharacterTable( g );;|
  !gapprompt@gap>| !gapinput@SetName( tbl, "S4" );  irr:= Irr( g );|
  [ Character( S4, [ 1, -1, 1, 1, -1 ] ),
    Character( S4, [ 3, -1, -1, 0, 1 ] ),
    Character( S4, [ 2, 0, 2, -1, 0 ] ),
    Character( S4, [ 3, 1, -1, 0, -1 ] ),
    Character( S4, [ 1, 1, 1, 1, 1 ] ) ]
  !gapprompt@gap>| !gapinput@2 * irr[5];|
  Character( S4, [ 2, 2, 2, 2, 2 ] )
  !gapprompt@gap>| !gapinput@irr[1] / 7;|
  ClassFunction( S4, [ 1/7, -1/7, 1/7, 1/7, -1/7 ] )
  !gapprompt@gap>| !gapinput@lincomb:= irr[3] + irr[1] - irr[5];|
  VirtualCharacter( S4, [ 2, -2, 2, -1, -2 ] )
  !gapprompt@gap>| !gapinput@lincomb:= lincomb + 2 * irr[5];|
  VirtualCharacter( S4, [ 4, 0, 4, 1, 0 ] )
  !gapprompt@gap>| !gapinput@IsCharacter( lincomb );|
  true
  !gapprompt@gap>| !gapinput@lincomb;|
  Character( S4, [ 4, 0, 4, 1, 0 ] )
  !gapprompt@gap>| !gapinput@irr[5] + 2;|
  [ 3, 3, 3, 3, 3 ]
  !gapprompt@gap>| !gapinput@irr[5] + [ 1, 2, 3, 4, 5 ];|
  [ 2, 3, 4, 5, 6 ]
  !gapprompt@gap>| !gapinput@zero:= 0 * irr[1];|
  VirtualCharacter( S4, [ 0, 0, 0, 0, 0 ] )
  !gapprompt@gap>| !gapinput@zero + Z(3);|
  [ Z(3), Z(3), Z(3), Z(3), Z(3) ]
  !gapprompt@gap>| !gapinput@irr[5] + TrivialCharacter( DihedralGroup( 8 ) );|
  [ 2, 2, 2, 2, 2 ]
\end{Verbatim}
 

 \index{class functions!as ring elements} The product of two class functions of the same character table is the tensor
product (pointwise product) of these class functions. Thus the set of all
class functions of a fixed group forms a ring, and for any field $F$ of cyclotomics, the $F$\texttt{\symbol{45}}span of a given set of class functions forms an algebra. 

 The product of two class functions of \emph{different} tables and the product of a class function and a list that is \emph{not} a class function are not defined, an error is signalled in these cases. Note
that in this respect, class functions behave differently from their values
lists, for which the product is defined as the standard scalar product. 

 
\begin{Verbatim}[commandchars=!@|,fontsize=\small,frame=single,label=Example]
  !gapprompt@gap>| !gapinput@tens:= irr[3] * irr[4];|
  Character( S4, [ 6, 0, -2, 0, 0 ] )
  !gapprompt@gap>| !gapinput@ValuesOfClassFunction( irr[3] ) * ValuesOfClassFunction( irr[4] );|
  4
\end{Verbatim}
 

 \index{inverse!of class function} Class functions without zero values are invertible, the \emph{inverse} is defined pointwise. As a consequence, for example groups of linear
characters can be formed. 

 
\begin{Verbatim}[commandchars=!@|,fontsize=\small,frame=single,label=Example]
  !gapprompt@gap>| !gapinput@tens / irr[1];|
  Character( S4, [ 6, 0, -2, 0, 0 ] )
\end{Verbatim}
 

 Other (somewhat strange) implications of the definition of arithmetic
operations for class functions, together with the general rules of list
arithmetic (see{\nobreakspace}\ref{Arithmetic for Lists}), apply to the case of products involving \emph{lists} of class functions. No inverse of the list of irreducible characters as a
matrix is defined; if one is interested in the inverse matrix then one can
compute it from the matrix of class function values. 

 
\begin{Verbatim}[commandchars=!@|,fontsize=\small,frame=single,label=Example]
  !gapprompt@gap>| !gapinput@Inverse( List( irr, ValuesOfClassFunction ) );|
  [ [ 1/24, 1/8, 1/12, 1/8, 1/24 ], [ -1/4, -1/4, 0, 1/4, 1/4 ],
    [ 1/8, -1/8, 1/4, -1/8, 1/8 ], [ 1/3, 0, -1/3, 0, 1/3 ],
    [ -1/4, 1/4, 0, -1/4, 1/4 ] ]
\end{Verbatim}
 

 Also the product of a class function with a list of class functions is \emph{not} a vector\texttt{\symbol{45}}matrix product but the list of pointwise products. 

 
\begin{Verbatim}[commandchars=!@|,fontsize=\small,frame=single,label=Example]
  !gapprompt@gap>| !gapinput@irr[1] * irr{ [ 1 .. 3 ] };|
  [ Character( S4, [ 1, 1, 1, 1, 1 ] ),
    Character( S4, [ 3, 1, -1, 0, -1 ] ),
    Character( S4, [ 2, 0, 2, -1, 0 ] ) ]
\end{Verbatim}
 

 And the product of two lists of class functions is \emph{not} the matrix product but the sum of the pointwise products. 

 
\begin{Verbatim}[commandchars=!@|,fontsize=\small,frame=single,label=Example]
  !gapprompt@gap>| !gapinput@irr * irr;|
  Character( S4, [ 24, 4, 8, 3, 4 ] )
\end{Verbatim}
 

 \index{character value!of group element using powering operator} \index{power!meaning for class functions} \index{^@\texttt{\texttt{\symbol{94}}}!for class functions} The \emph{powering} operator \texttt{\texttt{\symbol{92}}\texttt{\symbol{94}}} (\ref{bSlash^}) has several meanings for class functions. The power of a class function by a
nonnegative integer is clearly the tensor power. The power of a class function
by an element that normalizes the underlying group or by a Galois automorphism
is the conjugate class function. (As a consequence, the application of the
permutation induced by such an action cannot be denoted by \texttt{\texttt{\symbol{92}}\texttt{\symbol{94}}} (\ref{bSlash^}); instead one can use \texttt{Permuted} (\ref{Permuted}).) The power of a class function by a group or a character table is the
induced class function (see{\nobreakspace}\texttt{InducedClassFunction} (\ref{InducedClassFunction:for the character table of a supergroup})). The power of a group element by a class function is the class function
value at (the conjugacy class containing) this element. 

 
\begin{Verbatim}[commandchars=!@|,fontsize=\small,frame=single,label=Example]
  !gapprompt@gap>| !gapinput@irr[3] ^ 3;|
  Character( S4, [ 8, 0, 8, -1, 0 ] )
  !gapprompt@gap>| !gapinput@lin:= LinearCharacters( DerivedSubgroup( g ) );|
  [ Character( CharacterTable( Alt( [ 1 .. 4 ] ) ), [ 1, 1, 1, 1 ] ),
    Character( CharacterTable( Alt( [ 1 .. 4 ] ) ),
    [ 1, 1, E(3), E(3)^2 ] ),
    Character( CharacterTable( Alt( [ 1 .. 4 ] ) ),
    [ 1, 1, E(3)^2, E(3) ] ) ]
  !gapprompt@gap>| !gapinput@List( lin, chi -> chi ^ (1,2) );|
  [ Character( CharacterTable( Alt( [ 1 .. 4 ] ) ), [ 1, 1, 1, 1 ] ),
    Character( CharacterTable( Alt( [ 1 .. 4 ] ) ),
    [ 1, 1, E(3)^2, E(3) ] ),
    Character( CharacterTable( Alt( [ 1 .. 4 ] ) ),
    [ 1, 1, E(3), E(3)^2 ] ) ]
  !gapprompt@gap>| !gapinput@Orbit( GaloisGroup( CF(3) ), lin[2] );|
  [ Character( CharacterTable( Alt( [ 1 .. 4 ] ) ),
    [ 1, 1, E(3), E(3)^2 ] ),
    Character( CharacterTable( Alt( [ 1 .. 4 ] ) ),
    [ 1, 1, E(3)^2, E(3) ] ) ]
  !gapprompt@gap>| !gapinput@lin[1]^g;|
  Character( S4, [ 2, 0, 2, 2, 0 ] )
  !gapprompt@gap>| !gapinput@(1,2,3)^lin[2];|
  E(3)
\end{Verbatim}
 

\subsection{\textcolor{Chapter }{Characteristic (for a class function)}}
\logpage{[ 72, 4, 1 ]}\nobreak
\hyperdef{L}{X83AAD5527BBAFA03}{}
{\noindent\textcolor{FuncColor}{$\triangleright$\enspace\texttt{Characteristic({\mdseries\slshape chi})\index{Characteristic@\texttt{Characteristic}!for a class function}
\label{Characteristic:for a class function}
}\hfill{\scriptsize (attribute)}}\\


 The \emph{characteristic} of class functions is zero, as for all list of cyclotomics. For class
functions of a $p$\texttt{\symbol{45}}modular character table, such as Brauer characters, the
prime $p$ is given by the \texttt{UnderlyingCharacteristic} (\ref{UnderlyingCharacteristic:for a character table}) value of the character table. 

 
\begin{Verbatim}[commandchars=!@|,fontsize=\small,frame=single,label=Example]
  !gapprompt@gap>| !gapinput@Characteristic( irr[1] );|
  0
  !gapprompt@gap>| !gapinput@irrmod2:= Irr( g, 2 );|
  [ Character( BrauerTable( Sym( [ 1 .. 4 ] ), 2 ), [ 1, 1 ] ),
    Character( BrauerTable( Sym( [ 1 .. 4 ] ), 2 ), [ 2, -1 ] ) ]
  !gapprompt@gap>| !gapinput@Characteristic( irrmod2[1] );|
  0
  !gapprompt@gap>| !gapinput@UnderlyingCharacteristic( UnderlyingCharacterTable( irrmod2[1] ) );|
  2
\end{Verbatim}
 }

 

\subsection{\textcolor{Chapter }{ComplexConjugate (for a class function)}}
\logpage{[ 72, 4, 2 ]}\nobreak
\hyperdef{L}{X856AB97E785E0B04}{}
{\noindent\textcolor{FuncColor}{$\triangleright$\enspace\texttt{ComplexConjugate({\mdseries\slshape chi})\index{ComplexConjugate@\texttt{ComplexConjugate}!for a class function}
\label{ComplexConjugate:for a class function}
}\hfill{\scriptsize (attribute)}}\\
\noindent\textcolor{FuncColor}{$\triangleright$\enspace\texttt{GaloisCyc({\mdseries\slshape chi, k})\index{GaloisCyc@\texttt{GaloisCyc}!for a class function}
\label{GaloisCyc:for a class function}
}\hfill{\scriptsize (operation)}}\\
\noindent\textcolor{FuncColor}{$\triangleright$\enspace\texttt{Permuted({\mdseries\slshape chi, pi})\index{Permuted@\texttt{Permuted}!for a class function}
\label{Permuted:for a class function}
}\hfill{\scriptsize (method)}}\\


 The operations \texttt{ComplexConjugate}, \texttt{GaloisCyc}, and \texttt{Permuted} return a class function when they are called with a class function; The
complex conjugate of a class function that is known to be a (virtual)
character is again known to be a (virtual) character, and applying an
arbitrary Galois automorphism to an ordinary (virtual) character yields a
(virtual) character. 

 
\begin{Verbatim}[commandchars=!@|,fontsize=\small,frame=single,label=Example]
  !gapprompt@gap>| !gapinput@ComplexConjugate( lin[2] );|
  Character( CharacterTable( Alt( [ 1 .. 4 ] ) ),
   [ 1, 1, E(3)^2, E(3) ] )
  !gapprompt@gap>| !gapinput@GaloisCyc( lin[2], 5 );|
  Character( CharacterTable( Alt( [ 1 .. 4 ] ) ),
   [ 1, 1, E(3)^2, E(3) ] )
  !gapprompt@gap>| !gapinput@Permuted( lin[2], (2,3,4) );|
  ClassFunction( CharacterTable( Alt( [ 1 .. 4 ] ) ),
   [ 1, E(3)^2, 1, E(3) ] )
\end{Verbatim}
 }

 

 

\subsection{\textcolor{Chapter }{Order (for a class function)}}
\logpage{[ 72, 4, 3 ]}\nobreak
\hyperdef{L}{X7BCE99B88285EB39}{}
{\noindent\textcolor{FuncColor}{$\triangleright$\enspace\texttt{Order({\mdseries\slshape chi})\index{Order@\texttt{Order}!for a class function}
\label{Order:for a class function}
}\hfill{\scriptsize (attribute)}}\\


 By definition of \texttt{Order} (\ref{Order}) for arbitrary monoid elements, the return value of \texttt{Order} (\ref{Order}) for a character must be its multiplicative order. The \emph{determinantal order} (see{\nobreakspace}\texttt{DeterminantOfCharacter} (\ref{DeterminantOfCharacter})) of a character \mbox{\texttt{\mdseries\slshape chi}} can be computed as \texttt{Order( Determinant( \mbox{\texttt{\mdseries\slshape chi}} ) )}. 

 
\begin{Verbatim}[commandchars=!@|,fontsize=\small,frame=single,label=Example]
  !gapprompt@gap>| !gapinput@det:= Determinant( irr[3] );|
  Character( S4, [ 1, -1, 1, 1, -1 ] )
  !gapprompt@gap>| !gapinput@Order( det );|
  2
\end{Verbatim}
 }

 }

  
\section{\textcolor{Chapter }{Printing Class Functions}}\label{Printing Class Functions}
\logpage{[ 72, 5, 0 ]}
\hyperdef{L}{X828AD0C57EA57C21}{}
{
  

\subsection{\textcolor{Chapter }{ViewObj (for class functions)}}
\logpage{[ 72, 5, 1 ]}\nobreak
\hyperdef{L}{X7BDD2D4A7F7FB3B1}{}
{\noindent\textcolor{FuncColor}{$\triangleright$\enspace\texttt{ViewObj({\mdseries\slshape chi})\index{ViewObj@\texttt{ViewObj}!for class functions}
\label{ViewObj:for class functions}
}\hfill{\scriptsize (method)}}\\


 The default \texttt{ViewObj} (\ref{ViewObj}) methods for class functions print one of the strings \texttt{"ClassFunction"}, \texttt{"VirtualCharacter"}, \texttt{"Character"} (depending on whether the class function is known to be a character or virtual
character, see{\nobreakspace}\texttt{IsCharacter} (\ref{IsCharacter}), \texttt{IsVirtualCharacter} (\ref{IsVirtualCharacter})), followed by the \texttt{ViewObj} (\ref{ViewObj}) output for the underlying character table (see{\nobreakspace}\ref{Printing Character Tables}), and the list of values. The table is chosen (and not the group) in order to
distinguish class functions of different underlying characteristic
(see{\nobreakspace}\texttt{UnderlyingCharacteristic} (\ref{UnderlyingCharacteristic:for a character})). }

 

\subsection{\textcolor{Chapter }{PrintObj (for class functions)}}
\logpage{[ 72, 5, 2 ]}\nobreak
\hyperdef{L}{X871160B98595D4BA}{}
{\noindent\textcolor{FuncColor}{$\triangleright$\enspace\texttt{PrintObj({\mdseries\slshape chi})\index{PrintObj@\texttt{PrintObj}!for class functions}
\label{PrintObj:for class functions}
}\hfill{\scriptsize (method)}}\\


 The default \texttt{PrintObj} (\ref{PrintObj}) method for class functions does the same as \texttt{ViewObj} (\ref{ViewObj}), except that the character table is \texttt{Print} (\ref{Print})\texttt{\symbol{45}}ed instead of \texttt{View} (\ref{View})\texttt{\symbol{45}}ed. 

 \emph{Note} that if a class function is shown only with one of the strings \texttt{"ClassFunction"}, \texttt{"VirtualCharacter"}, it may still be that it is in fact a character; just this was not known at
the time when the class function was printed. 

 In order to reduce the space that is needed to print a class function, it may
be useful to give a name (see{\nobreakspace}\texttt{Name} (\ref{Name})) to the underlying character table. }

 

\subsection{\textcolor{Chapter }{Display (for class functions)}}
\logpage{[ 72, 5, 3 ]}\nobreak
\hyperdef{L}{X8430D31B8163C230}{}
{\noindent\textcolor{FuncColor}{$\triangleright$\enspace\texttt{Display({\mdseries\slshape chi})\index{Display@\texttt{Display}!for class functions}
\label{Display:for class functions}
}\hfill{\scriptsize (method)}}\\


 The default \texttt{Display} (\ref{Display}) method for a class function \mbox{\texttt{\mdseries\slshape chi}} calls \texttt{Display} (\ref{Display}) for its underlying character table (see{\nobreakspace}\ref{Printing Character Tables}), with \mbox{\texttt{\mdseries\slshape chi}} as the only entry in the \texttt{chars} list of the options record. 

 
\begin{Verbatim}[commandchars=!@|,fontsize=\small,frame=single,label=Example]
  !gapprompt@gap>| !gapinput@chi:= TrivialCharacter( CharacterTable( "A5" ) );|
  Character( CharacterTable( "A5" ), [ 1, 1, 1, 1, 1 ] )
  !gapprompt@gap>| !gapinput@Display( chi );|
  A5
  
       2  2  2  .  .  .
       3  1  .  1  .  .
       5  1  .  .  1  1
  
         1a 2a 3a 5a 5b
      2P 1a 1a 3a 5b 5a
      3P 1a 2a 1a 5b 5a
      5P 1a 2a 3a 1a 1a
  
  Y.1     1  1  1  1  1
\end{Verbatim}
 }

 }

  
\section{\textcolor{Chapter }{Creating Class Functions from Values Lists}}\label{Creating Class Functions from Values Lists}
\logpage{[ 72, 6, 0 ]}
\hyperdef{L}{X7BB90A8F86FFA456}{}
{
  

\subsection{\textcolor{Chapter }{ClassFunction (for a character table and a list)}}
\logpage{[ 72, 6, 1 ]}\nobreak
\hyperdef{L}{X78F4E23985FCA259}{}
{\noindent\textcolor{FuncColor}{$\triangleright$\enspace\texttt{ClassFunction({\mdseries\slshape tbl, values})\index{ClassFunction@\texttt{ClassFunction}!for a character table and a list}
\label{ClassFunction:for a character table and a list}
}\hfill{\scriptsize (operation)}}\\
\noindent\textcolor{FuncColor}{$\triangleright$\enspace\texttt{ClassFunction({\mdseries\slshape G, values})\index{ClassFunction@\texttt{ClassFunction}!for a group and a list}
\label{ClassFunction:for a group and a list}
}\hfill{\scriptsize (operation)}}\\


 In the first form, \texttt{ClassFunction} returns the class function of the character table \mbox{\texttt{\mdseries\slshape tbl}} with values given by the list \mbox{\texttt{\mdseries\slshape values}} of cyclotomics. In the second form, \mbox{\texttt{\mdseries\slshape G}} must be a group, and the class function of its ordinary character table is
returned. 

 Note that \mbox{\texttt{\mdseries\slshape tbl}} determines the underlying characteristic of the returned class function
(see{\nobreakspace}\texttt{UnderlyingCharacteristic} (\ref{UnderlyingCharacteristic:for a character})). }

 

\subsection{\textcolor{Chapter }{VirtualCharacter (for a character table and a list)}}
\logpage{[ 72, 6, 2 ]}\nobreak
\hyperdef{L}{X7805AFF77EFC3306}{}
{\noindent\textcolor{FuncColor}{$\triangleright$\enspace\texttt{VirtualCharacter({\mdseries\slshape tbl, values})\index{VirtualCharacter@\texttt{VirtualCharacter}!for a character table and a list}
\label{VirtualCharacter:for a character table and a list}
}\hfill{\scriptsize (operation)}}\\
\noindent\textcolor{FuncColor}{$\triangleright$\enspace\texttt{VirtualCharacter({\mdseries\slshape G, values})\index{VirtualCharacter@\texttt{VirtualCharacter}!for a group and a list}
\label{VirtualCharacter:for a group and a list}
}\hfill{\scriptsize (operation)}}\\


 \texttt{VirtualCharacter} returns the virtual character (see{\nobreakspace}\texttt{IsVirtualCharacter} (\ref{IsVirtualCharacter})) of the character table \mbox{\texttt{\mdseries\slshape tbl}} or the group \mbox{\texttt{\mdseries\slshape G}}, respectively, with values given by the list \mbox{\texttt{\mdseries\slshape values}}. 

 It is \emph{not} checked whether the given values really describe a virtual character. }

 

\subsection{\textcolor{Chapter }{Character (for a character table and a list)}}
\logpage{[ 72, 6, 3 ]}\nobreak
\hyperdef{L}{X849DD34C7968206C}{}
{\noindent\textcolor{FuncColor}{$\triangleright$\enspace\texttt{Character({\mdseries\slshape tbl, values})\index{Character@\texttt{Character}!for a character table and a list}
\label{Character:for a character table and a list}
}\hfill{\scriptsize (operation)}}\\
\noindent\textcolor{FuncColor}{$\triangleright$\enspace\texttt{Character({\mdseries\slshape G, values})\index{Character@\texttt{Character}!for a group and a list}
\label{Character:for a group and a list}
}\hfill{\scriptsize (operation)}}\\


 \texttt{Character} returns the character (see{\nobreakspace}\texttt{IsCharacter} (\ref{IsCharacter})) of the character table \mbox{\texttt{\mdseries\slshape tbl}} or the group \mbox{\texttt{\mdseries\slshape G}}, respectively, with values given by the list \mbox{\texttt{\mdseries\slshape values}}. 

 It is \emph{not} checked whether the given values really describe a character. 
\begin{Verbatim}[commandchars=!@|,fontsize=\small,frame=single,label=Example]
  !gapprompt@gap>| !gapinput@g:= DihedralGroup( 8 );  tbl:= CharacterTable( g );|
  <pc group of size 8 with 3 generators>
  CharacterTable( <pc group of size 8 with 3 generators> )
  !gapprompt@gap>| !gapinput@SetName( tbl, "D8" );|
  !gapprompt@gap>| !gapinput@phi:= ClassFunction( g, [ 1, -1, 0, 2, -2 ] );|
  ClassFunction( D8, [ 1, -1, 0, 2, -2 ] )
  !gapprompt@gap>| !gapinput@psi:= ClassFunction( tbl,|
  !gapprompt@>| !gapinput@             List( Irr( g ), chi -> ScalarProduct( chi, phi ) ) );|
  ClassFunction( D8, [ -3/8, 9/8, 5/8, 1/8, -1/4 ] )
  !gapprompt@gap>| !gapinput@chi:= VirtualCharacter( g, [ 0, 0, 8, 0, 0 ] );|
  VirtualCharacter( D8, [ 0, 0, 8, 0, 0 ] )
  !gapprompt@gap>| !gapinput@reg:= Character( tbl, [ 8, 0, 0, 0, 0 ] );|
  Character( D8, [ 8, 0, 0, 0, 0 ] )
\end{Verbatim}
 }

 

\subsection{\textcolor{Chapter }{ClassFunctionSameType}}
\logpage{[ 72, 6, 4 ]}\nobreak
\hyperdef{L}{X7B38035981D71B1B}{}
{\noindent\textcolor{FuncColor}{$\triangleright$\enspace\texttt{ClassFunctionSameType({\mdseries\slshape tbl, chi, values})\index{ClassFunctionSameType@\texttt{ClassFunctionSameType}}
\label{ClassFunctionSameType}
}\hfill{\scriptsize (function)}}\\


 Let \mbox{\texttt{\mdseries\slshape tbl}} be a character table, \mbox{\texttt{\mdseries\slshape chi}} a class function object (\emph{not} necessarily a class function of \mbox{\texttt{\mdseries\slshape tbl}}), and \mbox{\texttt{\mdseries\slshape values}} a list of cyclotomics. \texttt{ClassFunctionSameType} returns the class function $\psi$ of \mbox{\texttt{\mdseries\slshape tbl}} with values list \mbox{\texttt{\mdseries\slshape values}}, constructed with \texttt{ClassFunction} (\ref{ClassFunction:for a character table and a list}). 

 If \mbox{\texttt{\mdseries\slshape chi}} is known to be a (virtual) character then $\psi$ is also known to be a (virtual) character. 

 
\begin{Verbatim}[commandchars=!@|,fontsize=\small,frame=single,label=Example]
  !gapprompt@gap>| !gapinput@h:= Centre( g );;|
  !gapprompt@gap>| !gapinput@centbl:= CharacterTable( h );;  SetName( centbl, "C2" );|
  !gapprompt@gap>| !gapinput@ClassFunctionSameType( centbl, phi, [ 1, 1 ] );|
  ClassFunction( C2, [ 1, 1 ] )
  !gapprompt@gap>| !gapinput@ClassFunctionSameType( centbl, chi, [ 1, 1 ] );|
  VirtualCharacter( C2, [ 1, 1 ] )
  !gapprompt@gap>| !gapinput@ClassFunctionSameType( centbl, reg, [ 1, 1 ] );|
  Character( C2, [ 1, 1 ] )
\end{Verbatim}
 }

 }

  
\section{\textcolor{Chapter }{Creating Class Functions using Groups}}\label{Creating Class Functions using Groups}
\logpage{[ 72, 7, 0 ]}
\hyperdef{L}{X8727C2CB7ABEBC84}{}
{
  
\subsection{\textcolor{Chapter }{TrivialCharacter}}\logpage{[ 72, 7, 1 ]}
\hyperdef{L}{X86129DC37C55E4D6}{}
{
\noindent\textcolor{FuncColor}{$\triangleright$\enspace\texttt{TrivialCharacter({\mdseries\slshape tbl})\index{TrivialCharacter@\texttt{TrivialCharacter}!for a character table}
\label{TrivialCharacter:for a character table}
}\hfill{\scriptsize (attribute)}}\\
\noindent\textcolor{FuncColor}{$\triangleright$\enspace\texttt{TrivialCharacter({\mdseries\slshape G})\index{TrivialCharacter@\texttt{TrivialCharacter}!for a group}
\label{TrivialCharacter:for a group}
}\hfill{\scriptsize (attribute)}}\\


 is the \emph{trivial character} of the group \mbox{\texttt{\mdseries\slshape G}} or its character table \mbox{\texttt{\mdseries\slshape tbl}}, respectively. This is the class function with value equal to $1$ for each class. 

 
\begin{Verbatim}[commandchars=!@|,fontsize=\small,frame=single,label=Example]
  !gapprompt@gap>| !gapinput@TrivialCharacter( CharacterTable( "A5" ) );|
  Character( CharacterTable( "A5" ), [ 1, 1, 1, 1, 1 ] )
  !gapprompt@gap>| !gapinput@TrivialCharacter( SymmetricGroup( 3 ) );|
  Character( CharacterTable( Sym( [ 1 .. 3 ] ) ), [ 1, 1, 1 ] )
\end{Verbatim}
 }

 

\subsection{\textcolor{Chapter }{NaturalCharacter (for a group)}}
\logpage{[ 72, 7, 2 ]}\nobreak
\hyperdef{L}{X82C01DDB82D751A9}{}
{\noindent\textcolor{FuncColor}{$\triangleright$\enspace\texttt{NaturalCharacter({\mdseries\slshape G})\index{NaturalCharacter@\texttt{NaturalCharacter}!for a group}
\label{NaturalCharacter:for a group}
}\hfill{\scriptsize (attribute)}}\\
\noindent\textcolor{FuncColor}{$\triangleright$\enspace\texttt{NaturalCharacter({\mdseries\slshape hom})\index{NaturalCharacter@\texttt{NaturalCharacter}!for a homomorphism}
\label{NaturalCharacter:for a homomorphism}
}\hfill{\scriptsize (attribute)}}\\


 If the argument is a permutation group \mbox{\texttt{\mdseries\slshape G}} then \texttt{NaturalCharacter} returns the (ordinary) character of the natural permutation representation of \mbox{\texttt{\mdseries\slshape G}} on the set of moved points (see \texttt{MovedPoints} (\ref{MovedPoints:for a list or collection of permutations})), that is, the value on each class is the number of points among the moved
points of \mbox{\texttt{\mdseries\slshape G}} that are fixed by any permutation in that class. 

 If the argument is a matrix group \mbox{\texttt{\mdseries\slshape G}} in characteristic zero then \texttt{NaturalCharacter} returns the (ordinary) character of the natural matrix representation of \mbox{\texttt{\mdseries\slshape G}}, that is, the value on each class is the trace of any matrix in that class. 

 If the argument is a group homomorphism \mbox{\texttt{\mdseries\slshape hom}} whose image is a permutation group or a matrix group then \texttt{NaturalCharacter} returns the restriction of the natural character of the image of \mbox{\texttt{\mdseries\slshape hom}} to the preimage of \mbox{\texttt{\mdseries\slshape hom}}. 

 
\begin{Verbatim}[commandchars=!@|,fontsize=\small,frame=single,label=Example]
  !gapprompt@gap>| !gapinput@NaturalCharacter( SymmetricGroup( 3 ) );|
  Character( CharacterTable( Sym( [ 1 .. 3 ] ) ), [ 3, 1, 0 ] )
  !gapprompt@gap>| !gapinput@NaturalCharacter( Group( [ [ 0, -1 ], [ 1, -1 ] ] ) );|
  Character( CharacterTable( Group([ [ [ 0, -1 ], [ 1, -1 ] ] ]) ),
  [ 2, -1, -1 ] )
  !gapprompt@gap>| !gapinput@d8:= DihedralGroup( 8 );;  hom:= RegularActionHomomorphism( d8 );;|
  !gapprompt@gap>| !gapinput@NaturalCharacter( hom );|
  Character( CharacterTable( <pc group of size 8 with 3 generators> ),
  [ 8, 0, 0, 0, 0 ] )
\end{Verbatim}
 }

 
\subsection{\textcolor{Chapter }{PermutationCharacter}}\logpage{[ 72, 7, 3 ]}
\hyperdef{L}{X7938621F81B65E03}{}
{
\noindent\textcolor{FuncColor}{$\triangleright$\enspace\texttt{PermutationCharacter({\mdseries\slshape G, D, opr})\index{PermutationCharacter@\texttt{PermutationCharacter}!for a group, an action domain, and a function}
\label{PermutationCharacter:for a group, an action domain, and a function}
}\hfill{\scriptsize (operation)}}\\
\noindent\textcolor{FuncColor}{$\triangleright$\enspace\texttt{PermutationCharacter({\mdseries\slshape G, U})\index{PermutationCharacter@\texttt{PermutationCharacter}!for two groups}
\label{PermutationCharacter:for two groups}
}\hfill{\scriptsize (operation)}}\\


 Called with a group \mbox{\texttt{\mdseries\slshape G}}, an action domain or proper set \mbox{\texttt{\mdseries\slshape D}}, and an action function \mbox{\texttt{\mdseries\slshape opr}} (see Chapter{\nobreakspace}\ref{Group Actions}), \texttt{PermutationCharacter} returns the \emph{permutation character} of the action of \mbox{\texttt{\mdseries\slshape G}} on \mbox{\texttt{\mdseries\slshape D}} via \mbox{\texttt{\mdseries\slshape opr}}, that is, the value on each class is the number of points in \mbox{\texttt{\mdseries\slshape D}} that are fixed by an element in this class under the action \mbox{\texttt{\mdseries\slshape opr}}. 

 If the arguments are a group \mbox{\texttt{\mdseries\slshape G}} and a subgroup \mbox{\texttt{\mdseries\slshape U}} of \mbox{\texttt{\mdseries\slshape G}} then \texttt{PermutationCharacter} returns the permutation character of the action of \mbox{\texttt{\mdseries\slshape G}} on the right cosets of \mbox{\texttt{\mdseries\slshape U}} via right multiplication. 

 To compute the permutation character of a \emph{transitive permutation group} \mbox{\texttt{\mdseries\slshape G}} on the cosets of a point stabilizer \mbox{\texttt{\mdseries\slshape U}}, the attribute \texttt{NaturalCharacter} (\ref{NaturalCharacter:for a group}) of \mbox{\texttt{\mdseries\slshape G}} can be used instead of \texttt{PermutationCharacter( \mbox{\texttt{\mdseries\slshape G}}, \mbox{\texttt{\mdseries\slshape U}} )}. 

 More facilities concerning permutation characters are the transitivity test
(see Section{\nobreakspace}\ref{Operations for Class Functions}) and several tools for computing possible permutation characters
(see{\nobreakspace}\ref{Possible Permutation Characters}, \ref{Computing Possible Permutation Characters}). 

 
\begin{Verbatim}[commandchars=!@|,fontsize=\small,frame=single,label=Example]
  !gapprompt@gap>| !gapinput@PermutationCharacter( GL(2,2), AsSSortedList( GF(2)^2 ), OnRight );|
  Character( CharacterTable( SL(2,2) ), [ 4, 2, 1 ] )
  !gapprompt@gap>| !gapinput@s3:= SymmetricGroup( 3 );;  a3:= DerivedSubgroup( s3 );;|
  !gapprompt@gap>| !gapinput@PermutationCharacter( s3, a3 );|
  Character( CharacterTable( Sym( [ 1 .. 3 ] ) ), [ 2, 0, 2 ] )
\end{Verbatim}
 }

 }

  
\section{\textcolor{Chapter }{Operations for Class Functions}}\label{Operations for Class Functions}
\logpage{[ 72, 8, 0 ]}
\hyperdef{L}{X86DDB47A7C6B45D0}{}
{
  In the description of the following operations, the optional first argument \mbox{\texttt{\mdseries\slshape tbl}} is needed only if the argument \mbox{\texttt{\mdseries\slshape chi}} is a plain list and not a class function object. In this case, \mbox{\texttt{\mdseries\slshape tbl}} must always be the character table of which \mbox{\texttt{\mdseries\slshape chi}} shall be regarded as a class function. 

\subsection{\textcolor{Chapter }{IsCharacter}}
\logpage{[ 72, 8, 1 ]}\nobreak
\hyperdef{L}{X7FE3CD08794051F8}{}
{\noindent\textcolor{FuncColor}{$\triangleright$\enspace\texttt{IsCharacter({\mdseries\slshape [tbl, ]chi})\index{IsCharacter@\texttt{IsCharacter}}
\label{IsCharacter}
}\hfill{\scriptsize (property)}}\\


 \index{ordinary character} An \emph{ordinary character} of a group $G$ is a class function of $G$ whose values are the traces of a complex matrix representation of $G$. 

 \index{Brauer character} A \emph{Brauer character} of $G$ in characteristic $p$ is a class function of $G$ whose values are the complex lifts of a matrix representation of $G$ with image a finite field of characteristic $p$. }

 

\subsection{\textcolor{Chapter }{IsVirtualCharacter}}
\logpage{[ 72, 8, 2 ]}\nobreak
\hyperdef{L}{X788DD05C86CB7030}{}
{\noindent\textcolor{FuncColor}{$\triangleright$\enspace\texttt{IsVirtualCharacter({\mdseries\slshape [tbl, ]chi})\index{IsVirtualCharacter@\texttt{IsVirtualCharacter}}
\label{IsVirtualCharacter}
}\hfill{\scriptsize (property)}}\\


 \index{virtual character} A \emph{virtual character} is a class function that can be written as the difference of two proper
characters (see{\nobreakspace}\texttt{IsCharacter} (\ref{IsCharacter})). }

 

\subsection{\textcolor{Chapter }{IsIrreducibleCharacter}}
\logpage{[ 72, 8, 3 ]}\nobreak
\hyperdef{L}{X79A4B1D3870C8807}{}
{\noindent\textcolor{FuncColor}{$\triangleright$\enspace\texttt{IsIrreducibleCharacter({\mdseries\slshape [tbl, ]chi})\index{IsIrreducibleCharacter@\texttt{IsIrreducibleCharacter}}
\label{IsIrreducibleCharacter}
}\hfill{\scriptsize (property)}}\\


 \index{irreducible character} A character is \emph{irreducible} if it cannot be written as the sum of two characters. For ordinary characters
this can be checked using the scalar product of class functions
(see{\nobreakspace}\texttt{ScalarProduct} (\ref{ScalarProduct:for characters})). For Brauer characters there is no generic method for checking
irreducibility. 

 
\begin{Verbatim}[commandchars=!@|,fontsize=\small,frame=single,label=Example]
  !gapprompt@gap>| !gapinput@S4:= SymmetricGroup( 4 );;  SetName( S4, "S4" );|
  !gapprompt@gap>| !gapinput@psi:= ClassFunction( S4, [ 1, 1, 1, -2, 1 ] );|
  ClassFunction( CharacterTable( S4 ), [ 1, 1, 1, -2, 1 ] )
  !gapprompt@gap>| !gapinput@IsVirtualCharacter( psi );|
  true
  !gapprompt@gap>| !gapinput@IsCharacter( psi );|
  false
  !gapprompt@gap>| !gapinput@chi:= ClassFunction( S4, SizesCentralizers( CharacterTable( S4 ) ) );|
  ClassFunction( CharacterTable( S4 ), [ 24, 4, 8, 3, 4 ] )
  !gapprompt@gap>| !gapinput@IsCharacter( chi );|
  true
  !gapprompt@gap>| !gapinput@IsIrreducibleCharacter( chi );|
  false
  !gapprompt@gap>| !gapinput@IsIrreducibleCharacter( TrivialCharacter( S4 ) );|
  true
\end{Verbatim}
 }

 

\subsection{\textcolor{Chapter }{DegreeOfCharacter}}
\logpage{[ 72, 8, 4 ]}\nobreak
\hyperdef{L}{X7802BC157C69DD75}{}
{\noindent\textcolor{FuncColor}{$\triangleright$\enspace\texttt{DegreeOfCharacter({\mdseries\slshape chi})\index{DegreeOfCharacter@\texttt{DegreeOfCharacter}}
\label{DegreeOfCharacter}
}\hfill{\scriptsize (attribute)}}\\


 is the value of the character \mbox{\texttt{\mdseries\slshape chi}} on the identity element. This can also be obtained as \mbox{\texttt{\mdseries\slshape chi}}\texttt{[1]}. 

 
\begin{Verbatim}[commandchars=!@|,fontsize=\small,frame=single,label=Example]
  !gapprompt@gap>| !gapinput@List( Irr( S4 ), DegreeOfCharacter );|
  [ 1, 3, 2, 3, 1 ]
  !gapprompt@gap>| !gapinput@nat:= NaturalCharacter( S4 );|
  Character( CharacterTable( S4 ), [ 4, 2, 0, 1, 0 ] )
  !gapprompt@gap>| !gapinput@nat[1];|
  4
\end{Verbatim}
 }

 

\subsection{\textcolor{Chapter }{ScalarProduct (for characters)}}
\logpage{[ 72, 8, 5 ]}\nobreak
\hyperdef{L}{X855FD9F983D275CD}{}
{\noindent\textcolor{FuncColor}{$\triangleright$\enspace\texttt{ScalarProduct({\mdseries\slshape [tbl, ]chi, psi})\index{ScalarProduct@\texttt{ScalarProduct}!for characters}
\label{ScalarProduct:for characters}
}\hfill{\scriptsize (operation)}}\\
\textbf{\indent Returns:\ }
 the scalar product of the class functions \mbox{\texttt{\mdseries\slshape chi}} and \mbox{\texttt{\mdseries\slshape psi}}, which belong to the same character table \mbox{\texttt{\mdseries\slshape tbl}}. 



 \index{constituent!of a group character} \index{decompose!a group character} \index{multiplicity!of constituents of a group character} \index{inner product!of group characters} If \mbox{\texttt{\mdseries\slshape chi}} and \mbox{\texttt{\mdseries\slshape psi}} are class function objects, the argument \mbox{\texttt{\mdseries\slshape tbl}} is not needed, but \mbox{\texttt{\mdseries\slshape tbl}} is necessary if at least one of \mbox{\texttt{\mdseries\slshape chi}}, \mbox{\texttt{\mdseries\slshape psi}} is just a plain list. 

 The scalar product of two \emph{ordinary} class functions $\chi$, $\psi$ of a group $G$ is defined as 

 $( \sum_{{g \in G}} \chi(g) \psi(g^{{-1}}) ) / |G|$. 

 For two \emph{$p$\texttt{\symbol{45}}modular} class functions, the scalar product is defined as $( \sum_{{g \in S}} \chi(g) \psi(g^{{-1}}) ) / |G|$, where $S$ is the set of $p$\texttt{\symbol{45}}regular elements in $G$. }

 

\subsection{\textcolor{Chapter }{MatScalarProducts}}
\logpage{[ 72, 8, 6 ]}\nobreak
\hyperdef{L}{X858DF4E67EBB99DA}{}
{\noindent\textcolor{FuncColor}{$\triangleright$\enspace\texttt{MatScalarProducts({\mdseries\slshape [tbl, ]list[, list2]})\index{MatScalarProducts@\texttt{MatScalarProducts}}
\label{MatScalarProducts}
}\hfill{\scriptsize (operation)}}\\


 Called with two lists \mbox{\texttt{\mdseries\slshape list}}, \mbox{\texttt{\mdseries\slshape list2}} of class functions of the same character table (which may be given as the
argument \mbox{\texttt{\mdseries\slshape tbl}}), \texttt{MatScalarProducts} returns the matrix of scalar products (see \texttt{ScalarProduct} (\ref{ScalarProduct:for characters})) More precisely, this matrix contains in the $i$\texttt{\symbol{45}}th row the list of scalar products of $\mbox{\texttt{\mdseries\slshape list2}}[i]$ with the entries of \mbox{\texttt{\mdseries\slshape list}}. 

 If only one list \mbox{\texttt{\mdseries\slshape list}} of class functions is given then a lower triangular matrix of scalar products
is returned, containing (for $j \leq i$) in the $i$\texttt{\symbol{45}}th row in column $j$ the value \texttt{ScalarProduct}$( \mbox{\texttt{\mdseries\slshape tbl}}, \mbox{\texttt{\mdseries\slshape list}}[j], \mbox{\texttt{\mdseries\slshape list}}[i] )$. }

 

\subsection{\textcolor{Chapter }{Norm (for a class function)}}
\logpage{[ 72, 8, 7 ]}\nobreak
\hyperdef{L}{X8572B18A7BAED73E}{}
{\noindent\textcolor{FuncColor}{$\triangleright$\enspace\texttt{Norm({\mdseries\slshape [tbl, ]chi})\index{Norm@\texttt{Norm}!for a class function}
\label{Norm:for a class function}
}\hfill{\scriptsize (attribute)}}\\


 \index{Norm@\texttt{Norm}!of character} For an ordinary class function \mbox{\texttt{\mdseries\slshape chi}} of a group $G$ we have $\mbox{\texttt{\mdseries\slshape chi}} = \sum_{{\chi \in Irr(G)}} a_{\chi} \chi$, with complex coefficients $a_{\chi}$. The \emph{norm} of \mbox{\texttt{\mdseries\slshape chi}} is defined as $\sum_{{\chi \in Irr(G)}} a_{\chi} \overline{{a_{\chi}}}$. 

 
\begin{Verbatim}[commandchars=!@|,fontsize=\small,frame=single,label=Example]
  !gapprompt@gap>| !gapinput@tbl:= CharacterTable( "A5" );;|
  !gapprompt@gap>| !gapinput@ScalarProduct( TrivialCharacter( tbl ), Sum( Irr( tbl ) ) );|
  1
  !gapprompt@gap>| !gapinput@ScalarProduct( tbl, [ 1, 1, 1, 1, 1 ], Sum( Irr( tbl ) ) );|
  1
  !gapprompt@gap>| !gapinput@tbl2:= tbl mod 2;|
  BrauerTable( "A5", 2 )
  !gapprompt@gap>| !gapinput@chi:= Irr( tbl2 )[1];|
  Character( BrauerTable( "A5", 2 ), [ 1, 1, 1, 1 ] )
  !gapprompt@gap>| !gapinput@ScalarProduct( chi, chi );|
  3/4
  !gapprompt@gap>| !gapinput@ScalarProduct( tbl2, [ 1, 1, 1, 1 ], [ 1, 1, 1, 1 ] );|
  3/4
  !gapprompt@gap>| !gapinput@chars:= Irr( tbl ){ [ 2 .. 4 ] };;|
  !gapprompt@gap>| !gapinput@chars:= Set( Tensored( chars, chars ) );;|
  !gapprompt@gap>| !gapinput@MatScalarProducts( Irr( tbl ), chars );|
  [ [ 0, 0, 0, 1, 1 ], [ 1, 1, 0, 0, 1 ], [ 1, 0, 1, 0, 1 ],
    [ 0, 1, 0, 1, 1 ], [ 0, 0, 1, 1, 1 ], [ 1, 1, 1, 1, 1 ] ]
  !gapprompt@gap>| !gapinput@MatScalarProducts( tbl, chars );|
  [ [ 2 ], [ 1, 3 ], [ 1, 2, 3 ], [ 2, 2, 1, 3 ], [ 2, 1, 2, 2, 3 ],
    [ 2, 3, 3, 3, 3, 5 ] ]
  !gapprompt@gap>| !gapinput@List( chars, Norm );|
  [ 2, 3, 3, 3, 3, 5 ]
\end{Verbatim}
 }

 

\subsection{\textcolor{Chapter }{ConstituentsOfCharacter}}
\logpage{[ 72, 8, 8 ]}\nobreak
\hyperdef{L}{X78550D7087DB1181}{}
{\noindent\textcolor{FuncColor}{$\triangleright$\enspace\texttt{ConstituentsOfCharacter({\mdseries\slshape [tbl, ]chi})\index{ConstituentsOfCharacter@\texttt{ConstituentsOfCharacter}}
\label{ConstituentsOfCharacter}
}\hfill{\scriptsize (attribute)}}\\


 Let \mbox{\texttt{\mdseries\slshape chi}} be an ordinary or modular (virtual) character. If an ordinary or modular
character table \mbox{\texttt{\mdseries\slshape tbl}} is given then \mbox{\texttt{\mdseries\slshape chi}} may also be a list of character values. 

 \texttt{ConstituentsOfCharacter} returns the set of those irreducible characters that occur in the
decomposition of \mbox{\texttt{\mdseries\slshape chi}} with nonzero coefficient. 

 
\begin{Verbatim}[commandchars=!@|,fontsize=\small,frame=single,label=Example]
  !gapprompt@gap>| !gapinput@nat:= NaturalCharacter( S4 );|
  Character( CharacterTable( S4 ), [ 4, 2, 0, 1, 0 ] )
  !gapprompt@gap>| !gapinput@ConstituentsOfCharacter( nat );|
  [ Character( CharacterTable( S4 ), [ 1, 1, 1, 1, 1 ] ),
    Character( CharacterTable( S4 ), [ 3, 1, -1, 0, -1 ] ) ]
\end{Verbatim}
 }

 

\subsection{\textcolor{Chapter }{KernelOfCharacter}}
\logpage{[ 72, 8, 9 ]}\nobreak
\hyperdef{L}{X7E0A24498710F12B}{}
{\noindent\textcolor{FuncColor}{$\triangleright$\enspace\texttt{KernelOfCharacter({\mdseries\slshape [tbl, ]chi})\index{KernelOfCharacter@\texttt{KernelOfCharacter}}
\label{KernelOfCharacter}
}\hfill{\scriptsize (attribute)}}\\


 For a class function \mbox{\texttt{\mdseries\slshape chi}} of a group $G$, \texttt{KernelOfCharacter} returns the normal subgroup of $G$ that is formed by those conjugacy classes for which the value of \mbox{\texttt{\mdseries\slshape chi}} equals the degree of \mbox{\texttt{\mdseries\slshape chi}}. If the underlying character table of \mbox{\texttt{\mdseries\slshape chi}} does not store the group $G$ then an error is signalled. (See{\nobreakspace}\texttt{ClassPositionsOfKernel} (\ref{ClassPositionsOfKernel}) for a way to handle the kernel implicitly, by listing the positions of
conjugacy classes in the kernel.) 

 The returned group is the kernel of any representation of $G$ that affords \mbox{\texttt{\mdseries\slshape chi}}. 

 
\begin{Verbatim}[commandchars=!@|,fontsize=\small,frame=single,label=Example]
  !gapprompt@gap>| !gapinput@List( Irr( S4 ), chi -> StructureDescription(KernelOfCharacter(chi)) );|
  [ "A4", "1", "C2 x C2", "1", "S4" ]
\end{Verbatim}
 }

 

\subsection{\textcolor{Chapter }{ClassPositionsOfKernel}}
\logpage{[ 72, 8, 10 ]}\nobreak
\hyperdef{L}{X7B4708B47D9C05B3}{}
{\noindent\textcolor{FuncColor}{$\triangleright$\enspace\texttt{ClassPositionsOfKernel({\mdseries\slshape chi})\index{ClassPositionsOfKernel@\texttt{ClassPositionsOfKernel}}
\label{ClassPositionsOfKernel}
}\hfill{\scriptsize (attribute)}}\\


 is the list of positions of those conjugacy classes that form the kernel of
the character \mbox{\texttt{\mdseries\slshape chi}}, that is, those positions with character value equal to the character degree. 

 
\begin{Verbatim}[commandchars=!@|,fontsize=\small,frame=single,label=Example]
  !gapprompt@gap>| !gapinput@List( Irr( S4 ), ClassPositionsOfKernel );|
  [ [ 1, 3, 4 ], [ 1 ], [ 1, 3 ], [ 1 ], [ 1, 2, 3, 4, 5 ] ]
\end{Verbatim}
 }

 

\subsection{\textcolor{Chapter }{CentreOfCharacter}}
\logpage{[ 72, 8, 11 ]}\nobreak
\hyperdef{L}{X7E77D4147A0836D3}{}
{\noindent\textcolor{FuncColor}{$\triangleright$\enspace\texttt{CentreOfCharacter({\mdseries\slshape [tbl, ]chi})\index{CentreOfCharacter@\texttt{CentreOfCharacter}}
\label{CentreOfCharacter}
}\hfill{\scriptsize (attribute)}}\\


 \index{centre!of a character} For a character \mbox{\texttt{\mdseries\slshape chi}} of a group $G$, \texttt{CentreOfCharacter} returns the \emph{centre} of \mbox{\texttt{\mdseries\slshape chi}}, that is, the normal subgroup of all those elements of $G$ for which the quotient of the value of \mbox{\texttt{\mdseries\slshape chi}} by the degree of \mbox{\texttt{\mdseries\slshape chi}} is a root of unity. 

 If the underlying character table of \mbox{\texttt{\mdseries\slshape psi}} does not store the group $G$ then an error is signalled. (See{\nobreakspace}\texttt{ClassPositionsOfCentre} (\ref{ClassPositionsOfCentre:for a character}) for a way to handle the centre implicitly, by listing the positions of
conjugacy classes in the centre.) 

 
\begin{Verbatim}[commandchars=!@|,fontsize=\small,frame=single,label=Example]
  !gapprompt@gap>| !gapinput@List( Irr( S4 ), chi -> StructureDescription(CentreOfCharacter(chi)) );|
  [ "S4", "1", "C2 x C2", "1", "S4" ]
\end{Verbatim}
 }

 

\subsection{\textcolor{Chapter }{ClassPositionsOfCentre (for a character)}}
\logpage{[ 72, 8, 12 ]}\nobreak
\hyperdef{L}{X7CE5B4137B399274}{}
{\noindent\textcolor{FuncColor}{$\triangleright$\enspace\texttt{ClassPositionsOfCentre({\mdseries\slshape chi})\index{ClassPositionsOfCentre@\texttt{ClassPositionsOfCentre}!for a character}
\label{ClassPositionsOfCentre:for a character}
}\hfill{\scriptsize (attribute)}}\\


 is the list of positions of classes forming the centre of the character \mbox{\texttt{\mdseries\slshape chi}} (see{\nobreakspace}\texttt{CentreOfCharacter} (\ref{CentreOfCharacter})). 

 
\begin{Verbatim}[commandchars=!@|,fontsize=\small,frame=single,label=Example]
  !gapprompt@gap>| !gapinput@List( Irr( S4 ), ClassPositionsOfCentre );|
  [ [ 1, 2, 3, 4, 5 ], [ 1 ], [ 1, 3 ], [ 1 ], [ 1, 2, 3, 4, 5 ] ]
\end{Verbatim}
 }

 

\subsection{\textcolor{Chapter }{InertiaSubgroup}}
\logpage{[ 72, 8, 13 ]}\nobreak
\hyperdef{L}{X7C3187387C2D9938}{}
{\noindent\textcolor{FuncColor}{$\triangleright$\enspace\texttt{InertiaSubgroup({\mdseries\slshape [tbl, ]G, chi})\index{InertiaSubgroup@\texttt{InertiaSubgroup}}
\label{InertiaSubgroup}
}\hfill{\scriptsize (operation)}}\\


 Let \mbox{\texttt{\mdseries\slshape chi}} be a character of a group $H$ and \mbox{\texttt{\mdseries\slshape tbl}} the character table of $H$; if the argument \mbox{\texttt{\mdseries\slshape tbl}} is not given then the underlying character table of \mbox{\texttt{\mdseries\slshape chi}} (see{\nobreakspace}\texttt{UnderlyingCharacterTable} (\ref{UnderlyingCharacterTable})) is used instead. Furthermore, let \mbox{\texttt{\mdseries\slshape G}} be a group that contains $H$ as a normal subgroup. 

 \texttt{InertiaSubgroup} returns the stabilizer in \mbox{\texttt{\mdseries\slshape G}} of \mbox{\texttt{\mdseries\slshape chi}}, w.r.t.{\nobreakspace}the action of \mbox{\texttt{\mdseries\slshape G}} on the classes of $H$ via conjugation. In other words, \texttt{InertiaSubgroup} returns the group of all those elements $g \in \mbox{\texttt{\mdseries\slshape G}}$ that satisfy $\mbox{\texttt{\mdseries\slshape chi}}^g = \mbox{\texttt{\mdseries\slshape chi}}$. 

 
\begin{Verbatim}[commandchars=!@|,fontsize=\small,frame=single,label=Example]
  !gapprompt@gap>| !gapinput@der:= DerivedSubgroup( S4 );|
  Alt( [ 1 .. 4 ] )
  !gapprompt@gap>| !gapinput@List( Irr( der ), chi -> InertiaSubgroup( S4, chi ) );|
  [ S4, S4, Alt( [ 1 .. 4 ] ), Alt( [ 1 .. 4 ] ) ]
\end{Verbatim}
 }

 

\subsection{\textcolor{Chapter }{CycleStructureClass}}
\logpage{[ 72, 8, 14 ]}\nobreak
\hyperdef{L}{X8269BE0079A64D43}{}
{\noindent\textcolor{FuncColor}{$\triangleright$\enspace\texttt{CycleStructureClass({\mdseries\slshape [tbl, ]chi, class})\index{CycleStructureClass@\texttt{CycleStructureClass}}
\label{CycleStructureClass}
}\hfill{\scriptsize (operation)}}\\


 Let \mbox{\texttt{\mdseries\slshape permchar}} be a permutation character, and \mbox{\texttt{\mdseries\slshape class}} be the position of a conjugacy class of the character table of \mbox{\texttt{\mdseries\slshape permchar}}. \texttt{CycleStructureClass} returns a list describing the cycle structure of each element in class \mbox{\texttt{\mdseries\slshape class}} in the underlying permutation representation, in the same format as the result
of \texttt{CycleStructurePerm} (\ref{CycleStructurePerm}). 

 
\begin{Verbatim}[commandchars=!@|,fontsize=\small,frame=single,label=Example]
  !gapprompt@gap>| !gapinput@nat:= NaturalCharacter( S4 );|
  Character( CharacterTable( S4 ), [ 4, 2, 0, 1, 0 ] )
  !gapprompt@gap>| !gapinput@List( [ 1 .. 5 ], i -> CycleStructureClass( nat, i ) );|
  [ [  ], [ 1 ], [ 2 ], [ , 1 ], [ ,, 1 ] ]
\end{Verbatim}
 }

 

\subsection{\textcolor{Chapter }{IsTransitive (for a character)}}
\logpage{[ 72, 8, 15 ]}\nobreak
\hyperdef{L}{X86EDB4047C5AD6E7}{}
{\noindent\textcolor{FuncColor}{$\triangleright$\enspace\texttt{IsTransitive({\mdseries\slshape [tbl, ]chi})\index{IsTransitive@\texttt{IsTransitive}!for a character}
\label{IsTransitive:for a character}
}\hfill{\scriptsize (property)}}\\


 For a permutation character \mbox{\texttt{\mdseries\slshape chi}} of the group $G$ that corresponds to an action on the $G$\texttt{\symbol{45}}set $\Omega$ (see{\nobreakspace}\texttt{PermutationCharacter} (\ref{PermutationCharacter:for a group, an action domain, and a function})), \texttt{IsTransitive} (\ref{IsTransitive:for a group, an action domain, etc.}) returns \texttt{true} if the action of $G$ on $\Omega$ is transitive, and \texttt{false} otherwise. }

 

\subsection{\textcolor{Chapter }{Transitivity (for a character)}}
\logpage{[ 72, 8, 16 ]}\nobreak
\hyperdef{L}{X801DC07B8029841B}{}
{\noindent\textcolor{FuncColor}{$\triangleright$\enspace\texttt{Transitivity({\mdseries\slshape [tbl, ]chi})\index{Transitivity@\texttt{Transitivity}!for a character}
\label{Transitivity:for a character}
}\hfill{\scriptsize (attribute)}}\\


 For a permutation character \mbox{\texttt{\mdseries\slshape chi}} of the group $G$ that corresponds to an action on the $G$\texttt{\symbol{45}}set $\Omega$ (see{\nobreakspace}\texttt{PermutationCharacter} (\ref{PermutationCharacter:for a group, an action domain, and a function})), \texttt{Transitivity} returns the maximal nonnegative integer $k$ such that the action of $G$ on $\Omega$ is $k$\texttt{\symbol{45}}transitive. 

 
\begin{Verbatim}[commandchars=!@|,fontsize=\small,frame=single,label=Example]
  !gapprompt@gap>| !gapinput@IsTransitive( nat );  Transitivity( nat );|
  true
  4
  !gapprompt@gap>| !gapinput@Transitivity( 2 * TrivialCharacter( S4 ) );|
  0
\end{Verbatim}
 }

 

\subsection{\textcolor{Chapter }{CentralCharacter}}
\logpage{[ 72, 8, 17 ]}\nobreak
\hyperdef{L}{X7DD8FDCF7FB7834A}{}
{\noindent\textcolor{FuncColor}{$\triangleright$\enspace\texttt{CentralCharacter({\mdseries\slshape [tbl, ]chi})\index{CentralCharacter@\texttt{CentralCharacter}}
\label{CentralCharacter}
}\hfill{\scriptsize (attribute)}}\\


 \index{central character} For a character \mbox{\texttt{\mdseries\slshape chi}} of a group $G$, \texttt{CentralCharacter} returns the \emph{central character} of \mbox{\texttt{\mdseries\slshape chi}}. 

 The central character of $\chi$ is the class function $\omega_{\chi}$ defined by $\omega_{\chi}(g) = |g^G| \cdot \chi(g)/\chi(1)$ for each $g \in G$. }

 

\subsection{\textcolor{Chapter }{DeterminantOfCharacter}}
\logpage{[ 72, 8, 18 ]}\nobreak
\hyperdef{L}{X7A292A58827B95B8}{}
{\noindent\textcolor{FuncColor}{$\triangleright$\enspace\texttt{DeterminantOfCharacter({\mdseries\slshape [tbl, ]chi})\index{DeterminantOfCharacter@\texttt{DeterminantOfCharacter}}
\label{DeterminantOfCharacter}
}\hfill{\scriptsize (attribute)}}\\


 \index{determinant character} \texttt{DeterminantOfCharacter} returns the \emph{determinant character} of the character \mbox{\texttt{\mdseries\slshape chi}}. This is defined to be the character obtained by taking the determinant of
representing matrices of any representation affording \mbox{\texttt{\mdseries\slshape chi}}; the determinant can be computed using \texttt{EigenvaluesChar} (\ref{EigenvaluesChar}). 

 It is also possible to call \texttt{Determinant} (\ref{Determinant}) instead of \texttt{DeterminantOfCharacter}. 

 Note that the determinant character is well\texttt{\symbol{45}}defined for
virtual characters. 

 
\begin{Verbatim}[commandchars=!@|,fontsize=\small,frame=single,label=Example]
  !gapprompt@gap>| !gapinput@CentralCharacter( TrivialCharacter( S4 ) );|
  ClassFunction( CharacterTable( S4 ), [ 1, 6, 3, 8, 6 ] )
  !gapprompt@gap>| !gapinput@DeterminantOfCharacter( Irr( S4 )[3] );|
  Character( CharacterTable( S4 ), [ 1, -1, 1, 1, -1 ] )
\end{Verbatim}
 }

 

\subsection{\textcolor{Chapter }{EigenvaluesChar}}
\logpage{[ 72, 8, 19 ]}\nobreak
\hyperdef{L}{X861B435C7F68AE7D}{}
{\noindent\textcolor{FuncColor}{$\triangleright$\enspace\texttt{EigenvaluesChar({\mdseries\slshape [tbl, ]chi, class})\index{EigenvaluesChar@\texttt{EigenvaluesChar}}
\label{EigenvaluesChar}
}\hfill{\scriptsize (operation)}}\\


 Let \mbox{\texttt{\mdseries\slshape chi}} be a character of a group $G$. For an element $g \in G$ in the \mbox{\texttt{\mdseries\slshape class}}\texttt{\symbol{45}}th conjugacy class, of order $n$, let $M$ be a matrix of a representation affording \mbox{\texttt{\mdseries\slshape chi}}. 

 \texttt{EigenvaluesChar} returns the list of length $n$ where at position $k$ the multiplicity of \texttt{E}$(n)^k = \exp(2 \pi i k / n)$ as an eigenvalue of $M$ is stored. 

 We have \texttt{\mbox{\texttt{\mdseries\slshape chi}}[ \mbox{\texttt{\mdseries\slshape class}} ] = List( [ 1 .. n ], k \texttt{\symbol{45}}{\textgreater}
E(n)\texttt{\symbol{94}}k ) * EigenvaluesChar( \mbox{\texttt{\mdseries\slshape tbl}}, \mbox{\texttt{\mdseries\slshape chi}}, \mbox{\texttt{\mdseries\slshape class}} )}. 

 It is also possible to call \texttt{Eigenvalues} (\ref{Eigenvalues}) instead of \texttt{EigenvaluesChar}. 

 
\begin{Verbatim}[commandchars=!@|,fontsize=\small,frame=single,label=Example]
  !gapprompt@gap>| !gapinput@chi:= Irr( CharacterTable( "A5" ) )[2];|
  Character( CharacterTable( "A5" ),
  [ 3, -1, 0, -E(5)-E(5)^4, -E(5)^2-E(5)^3 ] )
  !gapprompt@gap>| !gapinput@List( [ 1 .. 5 ], i -> Eigenvalues( chi, i ) );|
  [ [ 3 ], [ 2, 1 ], [ 1, 1, 1 ], [ 0, 1, 1, 0, 1 ], [ 1, 0, 0, 1, 1 ] ]
\end{Verbatim}
 }

 

\subsection{\textcolor{Chapter }{Tensored}}
\logpage{[ 72, 8, 20 ]}\nobreak
\hyperdef{L}{X7A106BE281EFD953}{}
{\noindent\textcolor{FuncColor}{$\triangleright$\enspace\texttt{Tensored({\mdseries\slshape chars1, chars2})\index{Tensored@\texttt{Tensored}}
\label{Tensored}
}\hfill{\scriptsize (operation)}}\\


 Let \mbox{\texttt{\mdseries\slshape chars1}} and \mbox{\texttt{\mdseries\slshape chars2}} be lists of (values lists of) class functions of the same character table. \texttt{Tensored} returns the list of tensor products of all entries in \mbox{\texttt{\mdseries\slshape chars1}} with all entries in \mbox{\texttt{\mdseries\slshape chars2}}. 

 
\begin{Verbatim}[commandchars=!@|,fontsize=\small,frame=single,label=Example]
  !gapprompt@gap>| !gapinput@irra5:= Irr( CharacterTable( "A5" ) );;|
  !gapprompt@gap>| !gapinput@chars1:= irra5{ [ 1 .. 3 ] };;  chars2:= irra5{ [ 2, 3 ] };;|
  !gapprompt@gap>| !gapinput@Tensored( chars1, chars2 );|
  [ Character( CharacterTable( "A5" ),
      [ 3, -1, 0, -E(5)-E(5)^4, -E(5)^2-E(5)^3 ] ),
    Character( CharacterTable( "A5" ),
      [ 3, -1, 0, -E(5)^2-E(5)^3, -E(5)-E(5)^4 ] ),
    Character( CharacterTable( "A5" ),
      [ 9, 1, 0, -2*E(5)-E(5)^2-E(5)^3-2*E(5)^4,
        -E(5)-2*E(5)^2-2*E(5)^3-E(5)^4 ] ),
    Character( CharacterTable( "A5" ), [ 9, 1, 0, -1, -1 ] ),
    Character( CharacterTable( "A5" ), [ 9, 1, 0, -1, -1 ] ),
    Character( CharacterTable( "A5" ),
      [ 9, 1, 0, -E(5)-2*E(5)^2-2*E(5)^3-E(5)^4,
        -2*E(5)-E(5)^2-E(5)^3-2*E(5)^4 ] ) ]
\end{Verbatim}
 }

 

\subsection{\textcolor{Chapter }{TensorProduct (for characters)}}
\logpage{[ 72, 8, 21 ]}\nobreak
\hyperdef{L}{X7B83B4108490FD06}{}
{\noindent\textcolor{FuncColor}{$\triangleright$\enspace\texttt{TensorProduct({\mdseries\slshape chi, psi})\index{TensorProduct@\texttt{TensorProduct}!for characters}
\label{TensorProduct:for characters}
}\hfill{\scriptsize (operation)}}\\


 For two characters \mbox{\texttt{\mdseries\slshape chi}} and \mbox{\texttt{\mdseries\slshape psi}} afforded by the modules $V$ and $W$, \texttt{TensorProduct} returns the character that is afforded by the tensor product of $V$ and $W$. 

 The result can also be computed as \mbox{\texttt{\mdseries\slshape chi}}\texttt{*}\mbox{\texttt{\mdseries\slshape psi}}, see also \texttt{Tensored} (\ref{Tensored}). 

 
\begin{Verbatim}[commandchars=!@|,fontsize=\small,frame=single,label=Example]
  !gapprompt@gap>| !gapinput@t:= CharacterTable( "A5" );;|
  !gapprompt@gap>| !gapinput@chi:= Irr( t )[2];|
  Character( CharacterTable( "A5" ),
   [ 3, -1, 0, -E(5)-E(5)^4, -E(5)^2-E(5)^3 ] )
  !gapprompt@gap>| !gapinput@psi:= Irr( t )[3];|
  Character( CharacterTable( "A5" ),
   [ 3, -1, 0, -E(5)^2-E(5)^3, -E(5)-E(5)^4 ] )
  !gapprompt@gap>| !gapinput@TensorProduct( chi, psi );|
  Character( CharacterTable( "A5" ), [ 9, 1, 0, -1, -1 ] )
\end{Verbatim}
 }

 }

  
\section{\textcolor{Chapter }{Restricted and Induced Class Functions}}\label{Restricted and Induced Class Functions}
\logpage{[ 72, 9, 0 ]}
\hyperdef{L}{X854A4E3A85C5F89B}{}
{
  For restricting a class function of a group $G$ to a subgroup $H$ and for inducing a class function of $H$ to $G$, the \emph{class fusion} from $H$ to $G$ must be known (see{\nobreakspace}\ref{Class Fusions between Character Tables}). 

 \index{inflated class functions} If $F$ is the factor group of $G$ by the normal subgroup $N$ then each class function of $F$ can be naturally regarded as a class function of $G$, with $N$ in its kernel. For a class function of $F$, the corresponding class function of $G$ is called the \emph{inflated} class function. Restriction and inflation are in principle the same, namely
indirection of a class function by the appropriate fusion map, and thus no
extra operation is needed for this process. But note that contrary to the case
of a subgroup fusion, the factor fusion can in general not be computed from
the groups $G$ and $F$; either one needs the natural homomorphism, or the factor fusion to the
character table of $F$ must be stored on the table of $G$. This explains the different syntax for computing restricted and inflated
class functions. 

 In the following, the meaning of the optional first argument \mbox{\texttt{\mdseries\slshape tbl}} is the same as in Section{\nobreakspace}\ref{Operations for Class Functions}. 

\subsection{\textcolor{Chapter }{RestrictedClassFunction}}
\logpage{[ 72, 9, 1 ]}\nobreak
\hyperdef{L}{X86BABEA6841A40CF}{}
{\noindent\textcolor{FuncColor}{$\triangleright$\enspace\texttt{RestrictedClassFunction({\mdseries\slshape [tbl, ]chi, target})\index{RestrictedClassFunction@\texttt{RestrictedClassFunction}}
\label{RestrictedClassFunction}
}\hfill{\scriptsize (operation)}}\\


 Let \mbox{\texttt{\mdseries\slshape chi}} be a class function of a group $G$ and let \mbox{\texttt{\mdseries\slshape target}} be either a subgroup $H$ of $G$ or an injective homomorphism from $H$ to $G$ or the character table of \mbox{\texttt{\mdseries\slshape H}}. Then \texttt{RestrictedClassFunction} returns the class function of $H$ obtained by restricting \mbox{\texttt{\mdseries\slshape chi}} to $H$. 

 If \mbox{\texttt{\mdseries\slshape chi}} is a class function of a \emph{factor group} $G$ of $H$, where \mbox{\texttt{\mdseries\slshape target}} is either the group $H$ or a homomorphism from $H$ to $G$ or the character table of $H$ then the restriction can be computed in the case of the homomorphism; in the
other cases, this is possible only if the factor fusion from $H$ to $G$ is stored on the character table of $H$. }

 

\subsection{\textcolor{Chapter }{RestrictedClassFunctions}}
\logpage{[ 72, 9, 2 ]}\nobreak
\hyperdef{L}{X86DB64F08035D219}{}
{\noindent\textcolor{FuncColor}{$\triangleright$\enspace\texttt{RestrictedClassFunctions({\mdseries\slshape [tbl, ]chars, target})\index{RestrictedClassFunctions@\texttt{RestrictedClassFunctions}}
\label{RestrictedClassFunctions}
}\hfill{\scriptsize (operation)}}\\


 \texttt{RestrictedClassFunctions} is similar to \texttt{RestrictedClassFunction} (\ref{RestrictedClassFunction}), the only difference is that it takes a list \mbox{\texttt{\mdseries\slshape chars}} of class functions instead of one class function, and returns the list of
restricted class functions. 

 
\begin{Verbatim}[commandchars=!@|,fontsize=\small,frame=single,label=Example]
  !gapprompt@gap>| !gapinput@a5:= CharacterTable( "A5" );;  s5:= CharacterTable( "S5" );;|
  !gapprompt@gap>| !gapinput@RestrictedClassFunction( Irr( s5 )[2], a5 );|
  Character( CharacterTable( "A5" ), [ 1, 1, 1, 1, 1 ] )
  !gapprompt@gap>| !gapinput@RestrictedClassFunctions( Irr( s5 ), a5 );|
  [ Character( CharacterTable( "A5" ), [ 1, 1, 1, 1, 1 ] ),
    Character( CharacterTable( "A5" ), [ 1, 1, 1, 1, 1 ] ),
    Character( CharacterTable( "A5" ), [ 6, -2, 0, 1, 1 ] ),
    Character( CharacterTable( "A5" ), [ 4, 0, 1, -1, -1 ] ),
    Character( CharacterTable( "A5" ), [ 4, 0, 1, -1, -1 ] ),
    Character( CharacterTable( "A5" ), [ 5, 1, -1, 0, 0 ] ),
    Character( CharacterTable( "A5" ), [ 5, 1, -1, 0, 0 ] ) ]
  !gapprompt@gap>| !gapinput@hom:= NaturalHomomorphismByNormalSubgroup( S4, der );;|
  !gapprompt@gap>| !gapinput@RestrictedClassFunctions( Irr( Image( hom ) ), hom );|
  [ Character( CharacterTable( S4 ), [ 1, 1, 1, 1, 1 ] ),
    Character( CharacterTable( S4 ), [ 1, -1, 1, 1, -1 ] ) ]
\end{Verbatim}
 }

 
\subsection{\textcolor{Chapter }{InducedClassFunction}}\logpage{[ 72, 9, 3 ]}
\hyperdef{L}{X7FE39D3D78855D3B}{}
{
\noindent\textcolor{FuncColor}{$\triangleright$\enspace\texttt{InducedClassFunction({\mdseries\slshape [tbl, ]chi, H})\index{InducedClassFunction@\texttt{InducedClassFunction}!for a supergroup}
\label{InducedClassFunction:for a supergroup}
}\hfill{\scriptsize (operation)}}\\
\noindent\textcolor{FuncColor}{$\triangleright$\enspace\texttt{InducedClassFunction({\mdseries\slshape [tbl, ]chi, hom})\index{InducedClassFunction@\texttt{InducedClassFunction}!for a given monomorphism}
\label{InducedClassFunction:for a given monomorphism}
}\hfill{\scriptsize (operation)}}\\
\noindent\textcolor{FuncColor}{$\triangleright$\enspace\texttt{InducedClassFunction({\mdseries\slshape [tbl, ]chi, suptbl})\index{InducedClassFunction@\texttt{InducedClassFunction}!for the character table of a supergroup}
\label{InducedClassFunction:for the character table of a supergroup}
}\hfill{\scriptsize (operation)}}\\


 Let \mbox{\texttt{\mdseries\slshape chi}} be a class function of a group $G$ and let \mbox{\texttt{\mdseries\slshape target}} be either a supergroup $H$ of $G$ or an injective homomorphism from $H$ to $G$ or the character table of \mbox{\texttt{\mdseries\slshape H}}. Then \texttt{InducedClassFunction} returns the class function of $H$ obtained by inducing \mbox{\texttt{\mdseries\slshape chi}} to $H$. }

 

\subsection{\textcolor{Chapter }{InducedClassFunctions}}
\logpage{[ 72, 9, 4 ]}\nobreak
\hyperdef{L}{X8484C0F985AD2D28}{}
{\noindent\textcolor{FuncColor}{$\triangleright$\enspace\texttt{InducedClassFunctions({\mdseries\slshape [tbl, ]chars, target})\index{InducedClassFunctions@\texttt{InducedClassFunctions}}
\label{InducedClassFunctions}
}\hfill{\scriptsize (operation)}}\\


 \texttt{InducedClassFunctions} is similar to \texttt{InducedClassFunction} (\ref{InducedClassFunction:for a supergroup}), the only difference is that it takes a list \mbox{\texttt{\mdseries\slshape chars}} of class functions instead of one class function, and returns the list of
induced class functions. 

 
\begin{Verbatim}[commandchars=!@|,fontsize=\small,frame=single,label=Example]
  !gapprompt@gap>| !gapinput@InducedClassFunctions( Irr( a5 ), s5 );|
  [ Character( CharacterTable( "A5.2" ), [ 2, 2, 2, 2, 0, 0, 0 ] ),
    Character( CharacterTable( "A5.2" ), [ 6, -2, 0, 1, 0, 0, 0 ] ),
    Character( CharacterTable( "A5.2" ), [ 6, -2, 0, 1, 0, 0, 0 ] ),
    Character( CharacterTable( "A5.2" ), [ 8, 0, 2, -2, 0, 0, 0 ] ),
    Character( CharacterTable( "A5.2" ), [ 10, 2, -2, 0, 0, 0, 0 ] ) ]
\end{Verbatim}
 }

 

\subsection{\textcolor{Chapter }{InducedClassFunctionsByFusionMap}}
\logpage{[ 72, 9, 5 ]}\nobreak
\hyperdef{L}{X7C72003880743D28}{}
{\noindent\textcolor{FuncColor}{$\triangleright$\enspace\texttt{InducedClassFunctionsByFusionMap({\mdseries\slshape subtbl, tbl, chars, fusionmap})\index{InducedClassFunctionsByFusionMap@\texttt{InducedClassFunctionsByFusionMap}}
\label{InducedClassFunctionsByFusionMap}
}\hfill{\scriptsize (function)}}\\


 Let \mbox{\texttt{\mdseries\slshape subtbl}} and \mbox{\texttt{\mdseries\slshape tbl}} be two character tables of groups $H$ and $G$, such that $H$ is a subgroup of $G$, let \mbox{\texttt{\mdseries\slshape chars}} be a list of class functions of \mbox{\texttt{\mdseries\slshape subtbl}}, and let \mbox{\texttt{\mdseries\slshape fusionmap}} be a fusion map from \mbox{\texttt{\mdseries\slshape subtbl}} to \mbox{\texttt{\mdseries\slshape tbl}}. The function returns the list of induced class functions of \mbox{\texttt{\mdseries\slshape tbl}} that correspond to \mbox{\texttt{\mdseries\slshape chars}}, w.r.t. the given fusion map. 

 \texttt{InducedClassFunctionsByFusionMap} is the function that does the work for \texttt{InducedClassFunction} (\ref{InducedClassFunction:for the character table of a supergroup}) and \texttt{InducedClassFunctions} (\ref{InducedClassFunctions}). 

 
\begin{Verbatim}[commandchars=!@|,fontsize=\small,frame=single,label=Example]
  !gapprompt@gap>| !gapinput@fus:= PossibleClassFusions( a5, s5 );|
  [ [ 1, 2, 3, 4, 4 ] ]
  !gapprompt@gap>| !gapinput@InducedClassFunctionsByFusionMap( a5, s5, Irr( a5 ), fus[1] );|
  [ Character( CharacterTable( "A5.2" ), [ 2, 2, 2, 2, 0, 0, 0 ] ),
    Character( CharacterTable( "A5.2" ), [ 6, -2, 0, 1, 0, 0, 0 ] ),
    Character( CharacterTable( "A5.2" ), [ 6, -2, 0, 1, 0, 0, 0 ] ),
    Character( CharacterTable( "A5.2" ), [ 8, 0, 2, -2, 0, 0, 0 ] ),
    Character( CharacterTable( "A5.2" ), [ 10, 2, -2, 0, 0, 0, 0 ] ) ]
\end{Verbatim}
 }

 

\subsection{\textcolor{Chapter }{InducedCyclic}}
\logpage{[ 72, 9, 6 ]}\nobreak
\hyperdef{L}{X7C055F327C99CE71}{}
{\noindent\textcolor{FuncColor}{$\triangleright$\enspace\texttt{InducedCyclic({\mdseries\slshape tbl[, classes][, "all"]})\index{InducedCyclic@\texttt{InducedCyclic}}
\label{InducedCyclic}
}\hfill{\scriptsize (operation)}}\\


 \texttt{InducedCyclic} calculates characters induced up from cyclic subgroups of the ordinary
character table \mbox{\texttt{\mdseries\slshape tbl}} to \mbox{\texttt{\mdseries\slshape tbl}}, and returns the strictly sorted list of the induced characters. 

 If the string \texttt{"all"} is specified then all irreducible characters of these subgroups are induced,
otherwise only the permutation characters are calculated. 

 If a list \mbox{\texttt{\mdseries\slshape classes}} is specified then only those cyclic subgroups generated by these classes are
considered, otherwise all classes of \mbox{\texttt{\mdseries\slshape tbl}} are considered. 

 
\begin{Verbatim}[commandchars=!@|,fontsize=\small,frame=single,label=Example]
  !gapprompt@gap>| !gapinput@InducedCyclic( a5, "all" );|
  [ Character( CharacterTable( "A5" ), [ 12, 0, 0, 2, 2 ] ),
    Character( CharacterTable( "A5" ),
      [ 12, 0, 0, E(5)^2+E(5)^3, E(5)+E(5)^4 ] ),
    Character( CharacterTable( "A5" ),
      [ 12, 0, 0, E(5)+E(5)^4, E(5)^2+E(5)^3 ] ),
    Character( CharacterTable( "A5" ), [ 20, 0, -1, 0, 0 ] ),
    Character( CharacterTable( "A5" ), [ 20, 0, 2, 0, 0 ] ),
    Character( CharacterTable( "A5" ), [ 30, -2, 0, 0, 0 ] ),
    Character( CharacterTable( "A5" ), [ 30, 2, 0, 0, 0 ] ),
    Character( CharacterTable( "A5" ), [ 60, 0, 0, 0, 0 ] ) ]
\end{Verbatim}
 }

 }

  
\section{\textcolor{Chapter }{Reducing Virtual Characters}}\label{Reducing Virtual Characters}
\logpage{[ 72, 10, 0 ]}
\hyperdef{L}{X7C95F7937B752F48}{}
{
  The following operations are intended for the situation that one is given a
list of virtual characters of a character table and is interested in the
irreducible characters of this table. The idea is to compute virtual
characters of small norm from the given ones, hoping to get eventually virtual
characters of norm $1$. 

\subsection{\textcolor{Chapter }{ReducedClassFunctions}}
\logpage{[ 72, 10, 1 ]}\nobreak
\hyperdef{L}{X86F360D983343C2A}{}
{\noindent\textcolor{FuncColor}{$\triangleright$\enspace\texttt{ReducedClassFunctions({\mdseries\slshape [tbl, ][constituents, ]reducibles})\index{ReducedClassFunctions@\texttt{ReducedClassFunctions}}
\label{ReducedClassFunctions}
}\hfill{\scriptsize (operation)}}\\


 Let \mbox{\texttt{\mdseries\slshape reducibles}} be a list of ordinary virtual characters of a group $G$. If \mbox{\texttt{\mdseries\slshape constituents}} is given then it must also be a list of ordinary virtual characters of $G$, otherwise we have \mbox{\texttt{\mdseries\slshape constituents}} equal to \mbox{\texttt{\mdseries\slshape reducibles}} in the following. 

 \texttt{ReducedClassFunctions} returns a record with the components \texttt{remainders} and \texttt{irreducibles}, both lists of virtual characters of $G$. These virtual characters are computed as follows. 

 Let \texttt{rems} be the set of nonzero class functions obtained by subtraction of 
\[ \sum_{\chi} ( [\mbox{\texttt{\mdseries\slshape reducibles}}[i], \chi] / [\chi, \chi] ) \cdot \chi \]
 from $\mbox{\texttt{\mdseries\slshape reducibles}}[i]$, where the summation runs over \mbox{\texttt{\mdseries\slshape constituents}} and $[\chi, \psi]$ denotes the scalar product of $G$\texttt{\symbol{45}}class functions. Let \texttt{irrs} be the list of irreducible characters in \texttt{rems}. 

 We project \texttt{rems} into the orthogonal space of \texttt{irrs} and all those irreducibles found this way until no new irreducibles arise.
Then the \texttt{irreducibles} list is the set of all found irreducible characters, and the \texttt{remainders} list is the set of all nonzero remainders. }

 

\subsection{\textcolor{Chapter }{ReducedCharacters}}
\logpage{[ 72, 10, 2 ]}\nobreak
\hyperdef{L}{X7B7138ED8586F09E}{}
{\noindent\textcolor{FuncColor}{$\triangleright$\enspace\texttt{ReducedCharacters({\mdseries\slshape [tbl, ]constituents, reducibles})\index{ReducedCharacters@\texttt{ReducedCharacters}}
\label{ReducedCharacters}
}\hfill{\scriptsize (operation)}}\\


 \texttt{ReducedCharacters} is similar to \texttt{ReducedClassFunctions} (\ref{ReducedClassFunctions}), the only difference is that \mbox{\texttt{\mdseries\slshape constituents}} and \mbox{\texttt{\mdseries\slshape reducibles}} are assumed to be lists of characters. This means that only those scalar
products must be formed where the degree of the character in \mbox{\texttt{\mdseries\slshape constituents}} does not exceed the degree of the character in \mbox{\texttt{\mdseries\slshape reducibles}}. 

 
\begin{Verbatim}[commandchars=!@|,fontsize=\small,frame=single,label=Example]
  !gapprompt@gap>| !gapinput@tbl:= CharacterTable( "A5" );;|
  !gapprompt@gap>| !gapinput@chars:= Irr( tbl ){ [ 2 .. 4 ] };;|
  !gapprompt@gap>| !gapinput@chars:= Set( Tensored( chars, chars ) );;|
  !gapprompt@gap>| !gapinput@red:= ReducedClassFunctions( chars );|
  rec(
    irreducibles :=
      [ Character( CharacterTable( "A5" ), [ 1, 1, 1, 1, 1 ] ),
        Character( CharacterTable( "A5" ),
          [ 3, -1, 0, -E(5)-E(5)^4, -E(5)^2-E(5)^3 ] ),
        Character( CharacterTable( "A5" ),
          [ 3, -1, 0, -E(5)^2-E(5)^3, -E(5)-E(5)^4 ] ),
        Character( CharacterTable( "A5" ), [ 4, 0, 1, -1, -1 ] ),
        Character( CharacterTable( "A5" ), [ 5, 1, -1, 0, 0 ] ) ],
    remainders := [  ] )
\end{Verbatim}
 }

 

\subsection{\textcolor{Chapter }{IrreducibleDifferences}}
\logpage{[ 72, 10, 3 ]}\nobreak
\hyperdef{L}{X7D3289BB865BCF98}{}
{\noindent\textcolor{FuncColor}{$\triangleright$\enspace\texttt{IrreducibleDifferences({\mdseries\slshape tbl, reducibles, reducibles2[, scprmat]})\index{IrreducibleDifferences@\texttt{IrreducibleDifferences}}
\label{IrreducibleDifferences}
}\hfill{\scriptsize (function)}}\\


 \texttt{IrreducibleDifferences} returns the list of irreducible characters which occur as difference of an
element of \mbox{\texttt{\mdseries\slshape reducibles}} and an element of \mbox{\texttt{\mdseries\slshape reducibles2}}, where these two arguments are lists of class functions of the character
table \mbox{\texttt{\mdseries\slshape tbl}}. 

 If \mbox{\texttt{\mdseries\slshape reducibles2}} is the string \texttt{"triangle"} then the differences of elements in \mbox{\texttt{\mdseries\slshape reducibles}} are considered. 

 If \mbox{\texttt{\mdseries\slshape scprmat}} is not specified then it will be calculated, otherwise we must have \texttt{\mbox{\texttt{\mdseries\slshape scprmat}} = MatScalarProducts( \mbox{\texttt{\mdseries\slshape tbl}}, \mbox{\texttt{\mdseries\slshape reducibles}}, \mbox{\texttt{\mdseries\slshape reducibles2}} )} or \texttt{\mbox{\texttt{\mdseries\slshape scprmat}} = MatScalarProducts( \mbox{\texttt{\mdseries\slshape tbl}}, \mbox{\texttt{\mdseries\slshape reducibles}} )}, respectively. 

 
\begin{Verbatim}[commandchars=!@|,fontsize=\small,frame=single,label=Example]
  !gapprompt@gap>| !gapinput@IrreducibleDifferences( a5, chars, "triangle" );|
  [ Character( CharacterTable( "A5" ),
      [ 3, -1, 0, -E(5)-E(5)^4, -E(5)^2-E(5)^3 ] ),
    Character( CharacterTable( "A5" ),
      [ 3, -1, 0, -E(5)^2-E(5)^3, -E(5)-E(5)^4 ] ) ]
\end{Verbatim}
 }

 

\subsection{\textcolor{Chapter }{LLL}}
\logpage{[ 72, 10, 4 ]}\nobreak
\hyperdef{L}{X85B360C085B360C0}{}
{\noindent\textcolor{FuncColor}{$\triangleright$\enspace\texttt{LLL({\mdseries\slshape tbl, characters[, y][, "sort"][, "linearcomb"]})\index{LLL@\texttt{LLL}}
\label{LLL}
}\hfill{\scriptsize (function)}}\\


 \index{LLL algorithm!for virtual characters} \index{short vectors spanning a lattice} \index{lattice basis reduction!for virtual characters} \texttt{LLL} calls the LLL algorithm (see{\nobreakspace}\texttt{LLLReducedBasis} (\ref{LLLReducedBasis})) in the case of lattices spanned by the virtual characters \mbox{\texttt{\mdseries\slshape characters}} of the ordinary character table \mbox{\texttt{\mdseries\slshape tbl}} (see{\nobreakspace}\texttt{ScalarProduct} (\ref{ScalarProduct:for characters})). By finding shorter vectors in the lattice spanned by \mbox{\texttt{\mdseries\slshape characters}}, i.e., virtual characters of smaller norm, in some cases \texttt{LLL} is able to find irreducible characters. 

 \texttt{LLL} returns a record with at least components \texttt{irreducibles} (the list of found irreducible characters), \texttt{remainders} (a list of reducible virtual characters), and \texttt{norms} (the list of norms of the vectors in \texttt{remainders}). \texttt{irreducibles} together with \texttt{remainders} form a basis of the ${\ensuremath{\mathbb Z}}$\texttt{\symbol{45}}lattice spanned by \mbox{\texttt{\mdseries\slshape characters}}. 

 Note that the vectors in the \texttt{remainders} list are in general \emph{not} orthogonal (see{\nobreakspace}\texttt{ReducedClassFunctions} (\ref{ReducedClassFunctions})) to the irreducible characters in \texttt{irreducibles}. 

 Optional arguments of \texttt{LLL} are 

 
\begin{description}
\item[{\mbox{\texttt{\mdseries\slshape y}}}]  controls the sensitivity of the algorithm, see{\nobreakspace}\texttt{LLLReducedBasis} (\ref{LLLReducedBasis}), 
\item[{\mbox{\texttt{\mdseries\slshape "sort"}}}]  \texttt{LLL} sorts \mbox{\texttt{\mdseries\slshape characters}} and the \texttt{remainders} component of the result according to the degrees, 
\item[{\mbox{\texttt{\mdseries\slshape "linearcomb"}}}]  the returned record contains components \texttt{irreddecomp} and \texttt{reddecomp}, which are decomposition matrices of \texttt{irreducibles} and \texttt{remainders}, with respect to \mbox{\texttt{\mdseries\slshape characters}}. 
\end{description}
 

 
\begin{Verbatim}[commandchars=!@|,fontsize=\small,frame=single,label=Example]
  !gapprompt@gap>| !gapinput@s4:= CharacterTable( "Symmetric", 4 );;|
  !gapprompt@gap>| !gapinput@chars:= [ [ 8, 0, 0, -1, 0 ], [ 6, 0, 2, 0, 2 ],|
  !gapprompt@>| !gapinput@    [ 12, 0, -4, 0, 0 ], [ 6, 0, -2, 0, 0 ], [ 24, 0, 0, 0, 0 ],|
  !gapprompt@>| !gapinput@    [ 12, 0, 4, 0, 0 ], [ 6, 0, 2, 0, -2 ], [ 12, -2, 0, 0, 0 ],|
  !gapprompt@>| !gapinput@    [ 8, 0, 0, 2, 0 ], [ 12, 2, 0, 0, 0 ], [ 1, 1, 1, 1, 1 ] ];;|
  !gapprompt@gap>| !gapinput@LLL( s4, chars );|
  rec(
    irreducibles :=
      [ Character( CharacterTable( "Sym(4)" ), [ 2, 0, 2, -1, 0 ] ),
        Character( CharacterTable( "Sym(4)" ), [ 1, 1, 1, 1, 1 ] ),
        Character( CharacterTable( "Sym(4)" ), [ 3, 1, -1, 0, -1 ] ),
        Character( CharacterTable( "Sym(4)" ), [ 3, -1, -1, 0, 1 ] ),
        Character( CharacterTable( "Sym(4)" ), [ 1, -1, 1, 1, -1 ] ) ],
    norms := [  ], remainders := [  ] )
\end{Verbatim}
 }

 

\subsection{\textcolor{Chapter }{Extract}}
\logpage{[ 72, 10, 5 ]}\nobreak
\hyperdef{L}{X808D71A57D104ED7}{}
{\noindent\textcolor{FuncColor}{$\triangleright$\enspace\texttt{Extract({\mdseries\slshape tbl, reducibles, grammat[, missing]})\index{Extract@\texttt{Extract}}
\label{Extract}
}\hfill{\scriptsize (function)}}\\


 Let \mbox{\texttt{\mdseries\slshape tbl}} be an ordinary character table, \mbox{\texttt{\mdseries\slshape reducibles}} a list of characters of \mbox{\texttt{\mdseries\slshape tbl}}, and \mbox{\texttt{\mdseries\slshape grammat}} the matrix of scalar products of \mbox{\texttt{\mdseries\slshape reducibles}} (see{\nobreakspace}\texttt{MatScalarProducts} (\ref{MatScalarProducts})). \texttt{Extract} tries to find irreducible characters by drawing conclusions out of the scalar
products, using combinatorial and backtrack means. 

 The optional argument \mbox{\texttt{\mdseries\slshape missing}} is the maximal number of irreducible characters that occur as constituents of \mbox{\texttt{\mdseries\slshape reducibles}}. Specification of \mbox{\texttt{\mdseries\slshape missing}} may accelerate \texttt{Extract}. 

 \texttt{Extract} returns a record \mbox{\texttt{\mdseries\slshape ext}} with the components \texttt{solution} and \texttt{choice}, where the value of \texttt{solution} is a list $[ M_1, \ldots, M_n ]$ of decomposition matrices $M_i$ (up to permutations of rows) with the property that $M_i^{tr} \cdot X$ is equal to the sublist at the positions \mbox{\texttt{\mdseries\slshape ext}}\texttt{.choice[i]} of \mbox{\texttt{\mdseries\slshape reducibles}}, for a matrix $X$ of irreducible characters; the value of \texttt{choice} is a list of length $n$ whose entries are lists of indices. 

 So the $j$\texttt{\symbol{45}}th column in each matrix $M_i$ corresponds to $\mbox{\texttt{\mdseries\slshape reducibles}}[j]$, and each row in $M_i$ corresponds to an irreducible character. \texttt{Decreased} (\ref{Decreased}) can be used to examine the solution for computable irreducibles. 

 
\begin{Verbatim}[commandchars=!@|,fontsize=\small,frame=single,label=Example]
  !gapprompt@gap>| !gapinput@s4:= CharacterTable( "Symmetric", 4 );;|
  !gapprompt@gap>| !gapinput@red:= [ [ 5, 1, 5, 2, 1 ], [ 2, 0, 2, 2, 0 ], [ 3, -1, 3, 0, -1 ],|
  !gapprompt@>| !gapinput@           [ 6, 0, -2, 0, 0 ], [ 4, 0, 0, 1, 2 ] ];;|
  !gapprompt@gap>| !gapinput@gram:= MatScalarProducts( s4, red, red );|
  [ [ 6, 3, 2, 0, 2 ], [ 3, 2, 1, 0, 1 ], [ 2, 1, 2, 0, 0 ],
    [ 0, 0, 0, 2, 1 ], [ 2, 1, 0, 1, 2 ] ]
  !gapprompt@gap>| !gapinput@ext:= Extract( s4, red, gram, 5 );|
  rec( choice := [ [ 2, 5, 3, 4, 1 ] ],
    solution :=
      [
        [ [ 1, 1, 0, 0, 2 ], [ 1, 0, 1, 0, 1 ], [ 0, 1, 0, 1, 0 ],
            [ 0, 0, 1, 0, 1 ], [ 0, 0, 0, 1, 0 ] ] ] )
  !gapprompt@gap>| !gapinput@dec:= Decreased( s4, red, ext.solution[1], ext.choice[1] );|
  rec(
    irreducibles :=
      [ Character( CharacterTable( "Sym(4)" ), [ 1, 1, 1, 1, 1 ] ),
        Character( CharacterTable( "Sym(4)" ), [ 3, -1, -1, 0, 1 ] ),
        Character( CharacterTable( "Sym(4)" ), [ 1, -1, 1, 1, -1 ] ),
        Character( CharacterTable( "Sym(4)" ), [ 3, 1, -1, 0, -1 ] ),
        Character( CharacterTable( "Sym(4)" ), [ 2, 0, 2, -1, 0 ] ) ],
    matrix := [  ], remainders := [  ] )
\end{Verbatim}
 }

 

\subsection{\textcolor{Chapter }{OrthogonalEmbeddingsSpecialDimension}}
\logpage{[ 72, 10, 6 ]}\nobreak
\hyperdef{L}{X7F97B34A879D11BA}{}
{\noindent\textcolor{FuncColor}{$\triangleright$\enspace\texttt{OrthogonalEmbeddingsSpecialDimension({\mdseries\slshape tbl, reducibles, grammat[, "positive"], dim})\index{OrthogonalEmbeddingsSpecialDimension@\texttt{Orthogonal}\-\texttt{Embeddings}\-\texttt{Special}\-\texttt{Dimension}}
\label{OrthogonalEmbeddingsSpecialDimension}
}\hfill{\scriptsize (function)}}\\


 \texttt{OrthogonalEmbeddingsSpecialDimension} is a variant of \texttt{OrthogonalEmbeddings} (\ref{OrthogonalEmbeddings}) for the situation that \mbox{\texttt{\mdseries\slshape tbl}} is an ordinary character table, \mbox{\texttt{\mdseries\slshape reducibles}} is a list of virtual characters of \mbox{\texttt{\mdseries\slshape tbl}}, \mbox{\texttt{\mdseries\slshape grammat}} is the matrix of scalar products (see{\nobreakspace}\texttt{MatScalarProducts} (\ref{MatScalarProducts})), and \mbox{\texttt{\mdseries\slshape dim}} is an upper bound for the number of irreducible characters of \mbox{\texttt{\mdseries\slshape tbl}} that occur as constituents of \mbox{\texttt{\mdseries\slshape reducibles}}; if the vectors in \mbox{\texttt{\mdseries\slshape reducibles}} are known to be proper characters then the string \texttt{"positive"} may be entered as fourth argument. (See{\nobreakspace}\texttt{OrthogonalEmbeddings} (\ref{OrthogonalEmbeddings}) for information why this may help.) 

 \texttt{OrthogonalEmbeddingsSpecialDimension} first uses \texttt{OrthogonalEmbeddings} (\ref{OrthogonalEmbeddings}) to compute all orthogonal embeddings of \mbox{\texttt{\mdseries\slshape grammat}} into a standard lattice of dimension up to \mbox{\texttt{\mdseries\slshape dim}}, and then calls \texttt{Decreased} (\ref{Decreased}) in order to find irreducible characters of \mbox{\texttt{\mdseries\slshape tbl}}. 

 \texttt{OrthogonalEmbeddingsSpecialDimension} returns a record with the following components. 

 
\begin{description}
\item[{\texttt{irreducibles}}]  a list of found irreducibles, the intersection of all lists of irreducibles
found by \texttt{Decreased} (\ref{Decreased}), for all possible embeddings, and 
\item[{\texttt{remainders}}]  a list of remaining reducible virtual characters. 
\end{description}
 

 
\begin{Verbatim}[commandchars=!@|,fontsize=\small,frame=single,label=Example]
  !gapprompt@gap>| !gapinput@s6:= CharacterTable( "S6" );;|
  !gapprompt@gap>| !gapinput@red:= InducedCyclic( s6, "all" );;|
  !gapprompt@gap>| !gapinput@Add( red, TrivialCharacter( s6 ) );|
  !gapprompt@gap>| !gapinput@lll:= LLL( s6, red );;|
  !gapprompt@gap>| !gapinput@irred:= lll.irreducibles;|
  [ Character( CharacterTable( "A6.2_1" ),
      [ 1, 1, 1, 1, 1, 1, 1, 1, 1, 1, 1 ] ),
    Character( CharacterTable( "A6.2_1" ),
      [ 9, 1, 0, 0, 1, -1, -3, -3, 1, 0, 0 ] ),
    Character( CharacterTable( "A6.2_1" ),
      [ 16, 0, -2, -2, 0, 1, 0, 0, 0, 0, 0 ] ) ]
  !gapprompt@gap>| !gapinput@Set( Flat( MatScalarProducts( s6, irred, lll.remainders ) ) );|
  [ 0 ]
  !gapprompt@gap>| !gapinput@dim:= NrConjugacyClasses( s6 ) - Length( lll.irreducibles );|
  8
  !gapprompt@gap>| !gapinput@rem:= lll.remainders;;  Length( rem );|
  8
  !gapprompt@gap>| !gapinput@gram:= MatScalarProducts( s6, rem, rem );;  RankMat( gram );|
  8
  !gapprompt@gap>| !gapinput@emb1:= OrthogonalEmbeddings( gram, 8 );|
  rec( norms := [ 1, 1, 1, 1, 1, 1, 1, 1, 1, 1, 1, 1, 1, 1, 1, 1 ],
    solutions := [ [ 1, 2, 3, 7, 11, 12, 13, 15 ],
        [ 1, 2, 4, 8, 10, 12, 13, 14 ], [ 1, 2, 5, 6, 9, 12, 13, 16 ] ],
    vectors :=
      [ [ -1, 0, 1, 0, 1, 0, 1, 0 ], [ 1, 0, 0, 1, 0, 1, 0, 0 ],
        [ 0, 1, 1, 0, 0, 0, 1, 1 ], [ 0, 1, 1, 0, 0, 0, 1, 0 ],
        [ 0, 1, 1, 0, 0, 0, 0, 0 ], [ 0, 1, 0, 0, 0, 0, 1, 0 ],
        [ 0, -1, 0, 0, 0, 0, 0, 1 ], [ 0, 1, 0, 0, 0, 0, 0, 0 ],
        [ 0, 0, 1, 0, 0, 0, 1, 1 ], [ 0, 0, 1, 0, 0, 0, 0, 1 ],
        [ 0, 0, 1, 0, 0, 0, 0, 0 ], [ 0, 0, 0, -1, 1, 0, 0, 0 ],
        [ 0, 0, 0, 0, 0, 1, 0, 0 ], [ 0, 0, 0, 0, 0, 0, 1, 1 ],
        [ 0, 0, 0, 0, 0, 0, 1, 0 ], [ 0, 0, 0, 0, 0, 0, 0, 1 ] ] )
  !gapprompt@gap>| !gapinput@emb2:= OrthogonalEmbeddingsSpecialDimension( s6, rem, gram, 8 );|
  rec(
    irreducibles :=
      [ Character( CharacterTable( "A6.2_1" ),
          [ 5, 1, -1, 2, -1, 0, 1, -3, -1, 1, 0 ] ),
        Character( CharacterTable( "A6.2_1" ),
          [ 5, 1, 2, -1, -1, 0, -3, 1, -1, 0, 1 ] ),
        Character( CharacterTable( "A6.2_1" ),
          [ 10, -2, 1, 1, 0, 0, -2, 2, 0, 1, -1 ] ),
        Character( CharacterTable( "A6.2_1" ),
          [ 10, -2, 1, 1, 0, 0, 2, -2, 0, -1, 1 ] ) ],
    remainders :=
      [ VirtualCharacter( CharacterTable( "A6.2_1" ),
          [ 0, 0, 3, -3, 0, 0, 4, -4, 0, 1, -1 ] ),
        VirtualCharacter( CharacterTable( "A6.2_1" ),
          [ 6, 2, 3, 0, 0, 1, 2, -2, 0, -1, -2 ] ),
        VirtualCharacter( CharacterTable( "A6.2_1" ),
          [ 10, 2, 1, 1, 2, 0, 2, 2, -2, -1, -1 ] ),
        VirtualCharacter( CharacterTable( "A6.2_1" ),
          [ 14, 2, 2, -1, 0, -1, 6, 2, 0, 0, -1 ] ) ] )
\end{Verbatim}
 }

 

\subsection{\textcolor{Chapter }{Decreased}}
\logpage{[ 72, 10, 7 ]}\nobreak
\hyperdef{L}{X8799AB967C58C0E9}{}
{\noindent\textcolor{FuncColor}{$\triangleright$\enspace\texttt{Decreased({\mdseries\slshape tbl, chars, decompmat[, choice]})\index{Decreased@\texttt{Decreased}}
\label{Decreased}
}\hfill{\scriptsize (function)}}\\


 Let \mbox{\texttt{\mdseries\slshape tbl}} be an ordinary character table, \mbox{\texttt{\mdseries\slshape chars}} a list of virtual characters of \mbox{\texttt{\mdseries\slshape tbl}}, and \mbox{\texttt{\mdseries\slshape decompmat}} a decomposition matrix, that is, a matrix $M$ with the property that $M^{tr} \cdot X = \mbox{\texttt{\mdseries\slshape chars}}$ holds, where $X$ is a list of irreducible characters of \mbox{\texttt{\mdseries\slshape tbl}}. \texttt{Decreased} tries to compute the irreducibles in $X$ or at least some of them. 

 Usually \texttt{Decreased} is applied to the output of \texttt{Extract} (\ref{Extract}) or \texttt{OrthogonalEmbeddings} (\ref{OrthogonalEmbeddings}) or \texttt{OrthogonalEmbeddingsSpecialDimension} (\ref{OrthogonalEmbeddingsSpecialDimension}). In the case of \texttt{Extract} (\ref{Extract}), the choice component corresponding to the decomposition matrix must be
entered as argument \mbox{\texttt{\mdseries\slshape choice}} of \texttt{Decreased}. 

 \texttt{Decreased} returns \texttt{fail} if it can prove that no list $X$ of irreducible characters corresponding to the arguments exists; otherwise \texttt{Decreased} returns a record with the following components. 

 
\begin{description}
\item[{\texttt{irreducibles}}]  the list of found irreducible characters, 
\item[{\texttt{remainders}}]  the remaining reducible characters, and 
\item[{\texttt{matrix}}]  the decomposition matrix of the characters in the \texttt{remainders} component. 
\end{description}
 

 
\begin{Verbatim}[commandchars=!@|,fontsize=\small,frame=single,label=Example]
  !gapprompt@gap>| !gapinput@s4:= CharacterTable( "Symmetric", 4 );;|
  !gapprompt@gap>| !gapinput@x:= Irr( s4 );;|
  !gapprompt@gap>| !gapinput@red:= [ x[1]+x[2], -x[1]-x[3], -x[1]+x[3], -x[2]-x[4] ];;|
  !gapprompt@gap>| !gapinput@mat:= MatScalarProducts( s4, red, red );|
  [ [ 2, -1, -1, -1 ], [ -1, 2, 0, 0 ], [ -1, 0, 2, 0 ],
    [ -1, 0, 0, 2 ] ]
  !gapprompt@gap>| !gapinput@emb:= OrthogonalEmbeddings( mat );|
  rec( norms := [ 1, 1, 1, 1, 1, 1, 1, 1, 1, 1, 1, 1 ],
    solutions := [ [ 1, 6, 7, 12 ], [ 2, 5, 8, 11 ], [ 3, 4, 9, 10 ] ],
    vectors := [ [ -1, 1, 1, 0 ], [ -1, 1, 0, 1 ], [ 1, -1, 0, 0 ],
        [ -1, 0, 1, 1 ], [ -1, 0, 1, 0 ], [ -1, 0, 0, 1 ],
        [ 0, -1, 1, 0 ], [ 0, -1, 0, 1 ], [ 0, 1, 0, 0 ],
        [ 0, 0, -1, 1 ], [ 0, 0, 1, 0 ], [ 0, 0, 0, 1 ] ] )
  !gapprompt@gap>| !gapinput@dec:= Decreased( s4, red, emb.vectors{ emb.solutions[1] } );|
  rec(
    irreducibles :=
      [ Character( CharacterTable( "Sym(4)" ), [ 3, -1, -1, 0, 1 ] ),
        Character( CharacterTable( "Sym(4)" ), [ 1, -1, 1, 1, -1 ] ),
        Character( CharacterTable( "Sym(4)" ), [ 2, 0, 2, -1, 0 ] ),
        Character( CharacterTable( "Sym(4)" ), [ 3, 1, -1, 0, -1 ] ) ],
    matrix := [  ], remainders := [  ] )
  !gapprompt@gap>| !gapinput@Decreased( s4, red, emb.vectors{ emb.solutions[2] } );|
  fail
  !gapprompt@gap>| !gapinput@Decreased( s4, red, emb.vectors{ emb.solutions[3] } );|
  fail
\end{Verbatim}
 }

 

\subsection{\textcolor{Chapter }{DnLattice}}
\logpage{[ 72, 10, 8 ]}\nobreak
\hyperdef{L}{X85D510DC873A99B4}{}
{\noindent\textcolor{FuncColor}{$\triangleright$\enspace\texttt{DnLattice({\mdseries\slshape tbl, grammat, reducibles})\index{DnLattice@\texttt{DnLattice}}
\label{DnLattice}
}\hfill{\scriptsize (function)}}\\


 Let \mbox{\texttt{\mdseries\slshape tbl}} be an ordinary character table, and \mbox{\texttt{\mdseries\slshape reducibles}} a list of virtual characters of \mbox{\texttt{\mdseries\slshape tbl}}. 

 \texttt{DnLattice} searches for sublattices isomorphic to root lattices of type $D_n$, for $n \geq 4$, in the lattice that is generated by \mbox{\texttt{\mdseries\slshape reducibles}}; each vector in \mbox{\texttt{\mdseries\slshape reducibles}} must have norm $2$, and the matrix of scalar products (see{\nobreakspace}\texttt{MatScalarProducts} (\ref{MatScalarProducts})) of \mbox{\texttt{\mdseries\slshape reducibles}} must be entered as argument \mbox{\texttt{\mdseries\slshape grammat}}. 

 \texttt{DnLattice} is able to find irreducible characters if there is a lattice of type $D_n$ with $n > 4$. In the case $n = 4$, \texttt{DnLattice} may fail to determine irreducibles. 

 \texttt{DnLattice} returns a record with components 
\begin{description}
\item[{\texttt{irreducibles}}]  the list of found irreducible characters, 
\item[{\texttt{remainders}}]  the list of remaining reducible virtual characters, and 
\item[{\texttt{gram}}]  the Gram matrix of the vectors in \texttt{remainders}. 
\end{description}
 

 The \texttt{remainders} list is transformed in such a way that the \texttt{gram} matrix is a block diagonal matrix that exhibits the structure of the lattice
generated by the vectors in \texttt{remainders}. So \texttt{DnLattice} might be useful even if it fails to find irreducible characters. 

 
\begin{Verbatim}[commandchars=!@|,fontsize=\small,frame=single,label=Example]
  !gapprompt@gap>| !gapinput@s4:= CharacterTable( "Symmetric", 4 );;|
  !gapprompt@gap>| !gapinput@red:= [ [ 2, 0, 2, 2, 0 ], [ 4, 0, 0, 1, 2 ],|
  !gapprompt@>| !gapinput@           [ 5, -1, 1, -1, 1 ], [ -1, 1, 3, -1, -1 ] ];;|
  !gapprompt@gap>| !gapinput@gram:= MatScalarProducts( s4, red, red );|
  [ [ 2, 1, 0, 0 ], [ 1, 2, 1, -1 ], [ 0, 1, 2, 0 ], [ 0, -1, 0, 2 ] ]
  !gapprompt@gap>| !gapinput@dn:= DnLattice( s4, gram, red );|
  rec( gram := [  ],
    irreducibles :=
      [ Character( CharacterTable( "Sym(4)" ), [ 2, 0, 2, -1, 0 ] ),
        Character( CharacterTable( "Sym(4)" ), [ 1, -1, 1, 1, -1 ] ),
        Character( CharacterTable( "Sym(4)" ), [ 1, 1, 1, 1, 1 ] ),
        Character( CharacterTable( "Sym(4)" ), [ 3, -1, -1, 0, 1 ] ) ],
    remainders := [  ] )
\end{Verbatim}
 }

 

\subsection{\textcolor{Chapter }{DnLatticeIterative}}
\logpage{[ 72, 10, 9 ]}\nobreak
\hyperdef{L}{X78754D007F3572A7}{}
{\noindent\textcolor{FuncColor}{$\triangleright$\enspace\texttt{DnLatticeIterative({\mdseries\slshape tbl, reducibles})\index{DnLatticeIterative@\texttt{DnLatticeIterative}}
\label{DnLatticeIterative}
}\hfill{\scriptsize (function)}}\\


 Let \mbox{\texttt{\mdseries\slshape tbl}} be an ordinary character table, and \mbox{\texttt{\mdseries\slshape reducibles}} either a list of virtual characters of \mbox{\texttt{\mdseries\slshape tbl}} or a record with components \texttt{remainders} and \texttt{norms}, for example a record returned by \texttt{LLL} (\ref{LLL}). 

 \texttt{DnLatticeIterative} was designed for iterative use of \texttt{DnLattice} (\ref{DnLattice}). \texttt{DnLatticeIterative} selects the vectors of norm $2$ among the given virtual character, calls \texttt{DnLattice} (\ref{DnLattice}) for them, reduces the virtual characters with found irreducibles, calls \texttt{DnLattice} (\ref{DnLattice}) again for the remaining virtual characters, and so on, until no new
irreducibles are found. 

 \texttt{DnLatticeIterative} returns a record with the same components and meaning of components as \texttt{LLL} (\ref{LLL}). 

 
\begin{Verbatim}[commandchars=!@|,fontsize=\small,frame=single,label=Example]
  !gapprompt@gap>| !gapinput@s4:= CharacterTable( "Symmetric", 4 );;|
  !gapprompt@gap>| !gapinput@red:= [ [ 2, 0, 2, 2, 0 ], [ 4, 0, 0, 1, 2 ],|
  !gapprompt@>| !gapinput@           [ 5, -1, 1, -1, 1 ], [ -1, 1, 3, -1, -1 ] ];;|
  !gapprompt@gap>| !gapinput@dn:= DnLatticeIterative( s4, red );|
  rec(
    irreducibles :=
      [ Character( CharacterTable( "Sym(4)" ), [ 2, 0, 2, -1, 0 ] ),
        Character( CharacterTable( "Sym(4)" ), [ 1, -1, 1, 1, -1 ] ),
        Character( CharacterTable( "Sym(4)" ), [ 1, 1, 1, 1, 1 ] ),
        Character( CharacterTable( "Sym(4)" ), [ 3, -1, -1, 0, 1 ] ) ],
    norms := [  ], remainders := [  ] )
\end{Verbatim}
 }

 }

  
\section{\textcolor{Chapter }{Symmetrizations of Class Functions}}\label{Symmetrizations of Class Functions}
\logpage{[ 72, 11, 0 ]}
\hyperdef{L}{X87ED98F385B00D34}{}
{
  

\subsection{\textcolor{Chapter }{Symmetrizations}}
\logpage{[ 72, 11, 1 ]}\nobreak
\hyperdef{L}{X7E220413823330EC}{}
{\noindent\textcolor{FuncColor}{$\triangleright$\enspace\texttt{Symmetrizations({\mdseries\slshape [tbl, ]characters, n})\index{Symmetrizations@\texttt{Symmetrizations}}
\label{Symmetrizations}
}\hfill{\scriptsize (operation)}}\\


 \index{characters!symmetrizations of} \texttt{Symmetrizations} returns the list of symmetrizations of the characters \mbox{\texttt{\mdseries\slshape characters}} of the ordinary character table \mbox{\texttt{\mdseries\slshape tbl}} with the ordinary irreducible characters of the symmetric group of degree \mbox{\texttt{\mdseries\slshape n}}; instead of the integer \mbox{\texttt{\mdseries\slshape n}}, the character table of the symmetric group can be entered. 

 The symmetrization $\chi^{[\lambda]}$ of the character $\chi$ of \mbox{\texttt{\mdseries\slshape tbl}} with the character $\lambda$ of the symmetric group $S_n$ of degree $n$ is defined by 
\[ \chi^{[\lambda]}(g) = \left( \sum_{{\rho \in S_n}} \lambda(\rho)
\prod_{{k=1}}^n \chi(g^k)^{{a_k(\rho)}} \right) / n! , \]
 where $a_k(\rho)$ is the number of cycles of length $k$ in $\rho$. 

 \emph{Note} that the returned list may contain zero class functions, and duplicates are
not deleted.  

 For special kinds of symmetrizations, see{\nobreakspace}\texttt{SymmetricParts} (\ref{SymmetricParts}), \texttt{AntiSymmetricParts} (\ref{AntiSymmetricParts}), \texttt{MinusCharacter} (\ref{MinusCharacter}) and \texttt{OrthogonalComponents} (\ref{OrthogonalComponents}), \texttt{SymplecticComponents} (\ref{SymplecticComponents}), \texttt{ExteriorPower} (\ref{ExteriorPower:for a character}), \texttt{SymmetricPower} (\ref{SymmetricPower:for a character}). 

 
\begin{Verbatim}[commandchars=!@|,fontsize=\small,frame=single,label=Example]
  !gapprompt@gap>| !gapinput@tbl:= CharacterTable( "A5" );;|
  !gapprompt@gap>| !gapinput@Symmetrizations( Irr( tbl ){ [ 1 .. 3 ] }, 3 );|
  [ VirtualCharacter( CharacterTable( "A5" ), [ 0, 0, 0, 0, 0 ] ),
    VirtualCharacter( CharacterTable( "A5" ), [ 0, 0, 0, 0, 0 ] ),
    Character( CharacterTable( "A5" ), [ 1, 1, 1, 1, 1 ] ),
    Character( CharacterTable( "A5" ), [ 1, 1, 1, 1, 1 ] ),
    Character( CharacterTable( "A5" ),
      [ 8, 0, -1, -E(5)-E(5)^4, -E(5)^2-E(5)^3 ] ),
    Character( CharacterTable( "A5" ), [ 10, -2, 1, 0, 0 ] ),
    Character( CharacterTable( "A5" ), [ 1, 1, 1, 1, 1 ] ),
    Character( CharacterTable( "A5" ),
      [ 8, 0, -1, -E(5)^2-E(5)^3, -E(5)-E(5)^4 ] ),
    Character( CharacterTable( "A5" ), [ 10, -2, 1, 0, 0 ] ) ]
\end{Verbatim}
 }

 

\subsection{\textcolor{Chapter }{SymmetricParts}}
\logpage{[ 72, 11, 2 ]}\nobreak
\hyperdef{L}{X85CE68CA87CA383A}{}
{\noindent\textcolor{FuncColor}{$\triangleright$\enspace\texttt{SymmetricParts({\mdseries\slshape tbl, characters, n})\index{SymmetricParts@\texttt{SymmetricParts}}
\label{SymmetricParts}
}\hfill{\scriptsize (function)}}\\


 \index{symmetric power} is the list of symmetrizations of the characters \mbox{\texttt{\mdseries\slshape characters}} of the character table \mbox{\texttt{\mdseries\slshape tbl}} with the trivial character of the symmetric group of degree \mbox{\texttt{\mdseries\slshape n}} (see{\nobreakspace}\texttt{Symmetrizations} (\ref{Symmetrizations})). 

 
\begin{Verbatim}[commandchars=!@|,fontsize=\small,frame=single,label=Example]
  !gapprompt@gap>| !gapinput@tbl:= CharacterTable( "A5" );;|
  !gapprompt@gap>| !gapinput@SymmetricParts( tbl, Irr( tbl ), 3 );|
  [ Character( CharacterTable( "A5" ), [ 1, 1, 1, 1, 1 ] ),
    Character( CharacterTable( "A5" ), [ 10, -2, 1, 0, 0 ] ),
    Character( CharacterTable( "A5" ), [ 10, -2, 1, 0, 0 ] ),
    Character( CharacterTable( "A5" ), [ 20, 0, 2, 0, 0 ] ),
    Character( CharacterTable( "A5" ), [ 35, 3, 2, 0, 0 ] ) ]
\end{Verbatim}
 }

 

\subsection{\textcolor{Chapter }{AntiSymmetricParts}}
\logpage{[ 72, 11, 3 ]}\nobreak
\hyperdef{L}{X8329E934829FE965}{}
{\noindent\textcolor{FuncColor}{$\triangleright$\enspace\texttt{AntiSymmetricParts({\mdseries\slshape tbl, characters, n})\index{AntiSymmetricParts@\texttt{AntiSymmetricParts}}
\label{AntiSymmetricParts}
}\hfill{\scriptsize (function)}}\\


 \index{exterior power} is the list of symmetrizations of the characters \mbox{\texttt{\mdseries\slshape characters}} of the character table \mbox{\texttt{\mdseries\slshape tbl}} with the sign character of the symmetric group of degree \mbox{\texttt{\mdseries\slshape n}} (see{\nobreakspace}\texttt{Symmetrizations} (\ref{Symmetrizations})). 

 
\begin{Verbatim}[commandchars=!@|,fontsize=\small,frame=single,label=Example]
  !gapprompt@gap>| !gapinput@tbl:= CharacterTable( "A5" );;|
  !gapprompt@gap>| !gapinput@AntiSymmetricParts( tbl, Irr( tbl ), 3 );|
  [ VirtualCharacter( CharacterTable( "A5" ), [ 0, 0, 0, 0, 0 ] ),
    Character( CharacterTable( "A5" ), [ 1, 1, 1, 1, 1 ] ),
    Character( CharacterTable( "A5" ), [ 1, 1, 1, 1, 1 ] ),
    Character( CharacterTable( "A5" ), [ 4, 0, 1, -1, -1 ] ),
    Character( CharacterTable( "A5" ), [ 10, -2, 1, 0, 0 ] ) ]
\end{Verbatim}
 }

 

\subsection{\textcolor{Chapter }{ExteriorPower (for a character)}}
\logpage{[ 72, 11, 4 ]}\nobreak
\hyperdef{L}{X87BA59597CA21B3E}{}
{\noindent\textcolor{FuncColor}{$\triangleright$\enspace\texttt{ExteriorPower({\mdseries\slshape chi, n})\index{ExteriorPower@\texttt{ExteriorPower}!for a character}
\label{ExteriorPower:for a character}
}\hfill{\scriptsize (operation)}}\\


 For a character \mbox{\texttt{\mdseries\slshape chi}} afforded by the module $V$ and a positive integer \mbox{\texttt{\mdseries\slshape n}}, \texttt{ExteriorPower} returns the class function that is afforded by the \mbox{\texttt{\mdseries\slshape n}}\texttt{\symbol{45}}th exterior power of $V$. 

 This exterior power is the symmetrization of \mbox{\texttt{\mdseries\slshape chi}} with the sign character of the symmetric group of degree \mbox{\texttt{\mdseries\slshape n}}, see also \texttt{Symmetrizations} (\ref{Symmetrizations}) and \texttt{AntiSymmetricParts} (\ref{AntiSymmetricParts}). 

 
\begin{Verbatim}[commandchars=!@|,fontsize=\small,frame=single,label=Example]
  !gapprompt@gap>| !gapinput@t:= CharacterTable( "A5" );;|
  !gapprompt@gap>| !gapinput@List( Irr( t ), chi -> ExteriorPower( chi, 3 ) );|
  [ VirtualCharacter( CharacterTable( "A5" ), [ 0, 0, 0, 0, 0 ] ),
    Character( CharacterTable( "A5" ), [ 1, 1, 1, 1, 1 ] ),
    Character( CharacterTable( "A5" ), [ 1, 1, 1, 1, 1 ] ),
    Character( CharacterTable( "A5" ), [ 4, 0, 1, -1, -1 ] ),
    Character( CharacterTable( "A5" ), [ 10, -2, 1, 0, 0 ] ) ]
\end{Verbatim}
 }

 

\subsection{\textcolor{Chapter }{SymmetricPower (for a character)}}
\logpage{[ 72, 11, 5 ]}\nobreak
\hyperdef{L}{X855A7ECF80FCF0BA}{}
{\noindent\textcolor{FuncColor}{$\triangleright$\enspace\texttt{SymmetricPower({\mdseries\slshape chi, n})\index{SymmetricPower@\texttt{SymmetricPower}!for a character}
\label{SymmetricPower:for a character}
}\hfill{\scriptsize (operation)}}\\


 For a character \mbox{\texttt{\mdseries\slshape chi}} afforded by the module $V$ and a positive integer \mbox{\texttt{\mdseries\slshape n}}, \texttt{SymmetricPower} returns the class function that is afforded by the \mbox{\texttt{\mdseries\slshape n}}\texttt{\symbol{45}}th symmetric power of $V$. 

 This symmetric power is the symmetrization of \mbox{\texttt{\mdseries\slshape chi}} with the trivial character of the symmetric group of degree \mbox{\texttt{\mdseries\slshape n}}, see also \texttt{Symmetrizations} (\ref{Symmetrizations}) and \texttt{SymmetricParts} (\ref{SymmetricParts}). 

 
\begin{Verbatim}[commandchars=!@|,fontsize=\small,frame=single,label=Example]
  !gapprompt@gap>| !gapinput@t:= CharacterTable( "A5" );;|
  !gapprompt@gap>| !gapinput@List( Irr( t ), chi -> SymmetricPower( chi, 3 ) );|
  [ Character( CharacterTable( "A5" ), [ 1, 1, 1, 1, 1 ] ),
    Character( CharacterTable( "A5" ), [ 10, -2, 1, 0, 0 ] ),
    Character( CharacterTable( "A5" ), [ 10, -2, 1, 0, 0 ] ),
    Character( CharacterTable( "A5" ), [ 20, 0, 2, 0, 0 ] ),
    Character( CharacterTable( "A5" ), [ 35, 3, 2, 0, 0 ] ) ]
\end{Verbatim}
 }

 

\subsection{\textcolor{Chapter }{OrthogonalComponents}}
\logpage{[ 72, 11, 6 ]}\nobreak
\hyperdef{L}{X78648E367C65B1F1}{}
{\noindent\textcolor{FuncColor}{$\triangleright$\enspace\texttt{OrthogonalComponents({\mdseries\slshape tbl, chars, m})\index{OrthogonalComponents@\texttt{OrthogonalComponents}}
\label{OrthogonalComponents}
}\hfill{\scriptsize (function)}}\\


 \index{symmetrizations!orthogonal} \index{Frame} \index{Murnaghan components} If $\chi$ is a nonlinear character with indicator $+1$, a splitting of the tensor power $\chi^m$ is given by the so\texttt{\symbol{45}}called Murnaghan functions
(see{\nobreakspace}\cite{Mur58}). These components in general have fewer irreducible constituents than the
symmetrizations with the symmetric group of degree \mbox{\texttt{\mdseries\slshape m}} (see{\nobreakspace}\texttt{Symmetrizations} (\ref{Symmetrizations})). 

 \texttt{OrthogonalComponents} returns the Murnaghan components of the nonlinear characters of the character
table \mbox{\texttt{\mdseries\slshape tbl}} in the list \mbox{\texttt{\mdseries\slshape chars}} up to the power \mbox{\texttt{\mdseries\slshape m}}, where \mbox{\texttt{\mdseries\slshape m}} is an integer between 2 and 6. 

 The Murnaghan functions are implemented as in{\nobreakspace}\cite{Fra82}. 

 \emph{Note}: If \mbox{\texttt{\mdseries\slshape chars}} is a list of character objects (see{\nobreakspace}\texttt{IsCharacter} (\ref{IsCharacter})) then also the result consists of class function objects. It is not checked
whether all characters in \mbox{\texttt{\mdseries\slshape chars}} do really have indicator $+1$; if there are characters with indicator $0$ or $-1$, the result might contain virtual characters (see also{\nobreakspace}\texttt{SymplecticComponents} (\ref{SymplecticComponents})), therefore the entries of the result do in general not know that they are
characters. 

 
\begin{Verbatim}[commandchars=!@|,fontsize=\small,frame=single,label=Example]
  !gapprompt@gap>| !gapinput@tbl:= CharacterTable( "A8" );;  chi:= Irr( tbl )[2];|
  Character( CharacterTable( "A8" ), [ 7, -1, 3, 4, 1, -1, 1, 2, 0, -1,
    0, 0, -1, -1 ] )
  !gapprompt@gap>| !gapinput@OrthogonalComponents( tbl, [ chi ], 3 );|
  [ ClassFunction( CharacterTable( "A8" ),
      [ 21, -3, 1, 6, 0, 1, -1, 1, -2, 0, 0, 0, 1, 1 ] ),
    ClassFunction( CharacterTable( "A8" ),
      [ 27, 3, 7, 9, 0, -1, 1, 2, 1, 0, -1, -1, -1, -1 ] ),
    ClassFunction( CharacterTable( "A8" ),
      [ 105, 1, 5, 15, -3, 1, -1, 0, -1, 1, 0, 0, 0, 0 ] ),
    ClassFunction( CharacterTable( "A8" ),
      [ 35, 3, -5, 5, 2, -1, -1, 0, 1, 0, 0, 0, 0, 0 ] ),
    ClassFunction( CharacterTable( "A8" ),
      [ 77, -3, 13, 17, 2, 1, 1, 2, 1, 0, 0, 0, 2, 2 ] ) ]
\end{Verbatim}
 }

 

\subsection{\textcolor{Chapter }{SymplecticComponents}}
\logpage{[ 72, 11, 7 ]}\nobreak
\hyperdef{L}{X788B9AA17DD9418C}{}
{\noindent\textcolor{FuncColor}{$\triangleright$\enspace\texttt{SymplecticComponents({\mdseries\slshape tbl, chars, m})\index{SymplecticComponents@\texttt{SymplecticComponents}}
\label{SymplecticComponents}
}\hfill{\scriptsize (function)}}\\


 \index{symmetrizations!symplectic} \index{Murnaghan components} If $\chi$ is a (nonlinear) character with indicator $-1$, a splitting of the tensor power $\chi^m$ is given in terms of the so\texttt{\symbol{45}}called Murnaghan functions
(see{\nobreakspace}\cite{Mur58}). These components in general have fewer irreducible constituents than the
symmetrizations with the symmetric group of degree \mbox{\texttt{\mdseries\slshape m}} (see{\nobreakspace}\texttt{Symmetrizations} (\ref{Symmetrizations})). 

 \texttt{SymplecticComponents} returns the symplectic symmetrizations of the nonlinear characters of the
character table \mbox{\texttt{\mdseries\slshape tbl}} in the list \mbox{\texttt{\mdseries\slshape chars}} up to the power \mbox{\texttt{\mdseries\slshape m}}, where \mbox{\texttt{\mdseries\slshape m}} is an integer between $2$ and $5$. 

 \emph{Note}: If \mbox{\texttt{\mdseries\slshape chars}} is a list of character objects (see{\nobreakspace}\texttt{IsCharacter} (\ref{IsCharacter})) then also the result consists of class function objects. It is not checked
whether all characters in \mbox{\texttt{\mdseries\slshape chars}} do really have indicator $-1$; if there are characters with indicator $0$ or $+1$, the result might contain virtual characters (see also{\nobreakspace}\texttt{OrthogonalComponents} (\ref{OrthogonalComponents})), therefore the entries of the result do in general not know that they are
characters. 

 
\begin{Verbatim}[commandchars=!@|,fontsize=\small,frame=single,label=Example]
  !gapprompt@gap>| !gapinput@tbl:= CharacterTable( "U3(3)" );;  chi:= Irr( tbl )[2];|
  Character( CharacterTable( "U3(3)" ),
  [ 6, -2, -3, 0, -2, -2, 2, 1, -1, -1, 0, 0, 1, 1 ] )
  !gapprompt@gap>| !gapinput@SymplecticComponents( tbl, [ chi ], 3 );|
  [ ClassFunction( CharacterTable( "U3(3)" ),
      [ 14, -2, 5, -1, 2, 2, 2, 1, 0, 0, 0, 0, -1, -1 ] ),
    ClassFunction( CharacterTable( "U3(3)" ),
      [ 21, 5, 3, 0, 1, 1, 1, -1, 0, 0, -1, -1, 1, 1 ] ),
    ClassFunction( CharacterTable( "U3(3)" ),
      [ 64, 0, -8, -2, 0, 0, 0, 0, 1, 1, 0, 0, 0, 0 ] ),
    ClassFunction( CharacterTable( "U3(3)" ),
      [ 14, 6, -4, 2, -2, -2, 2, 0, 0, 0, 0, 0, -2, -2 ] ),
    ClassFunction( CharacterTable( "U3(3)" ),
      [ 56, -8, 2, 2, 0, 0, 0, -2, 0, 0, 0, 0, 0, 0 ] ) ]
\end{Verbatim}
 }

 }

  
\section{\textcolor{Chapter }{Molien Series}}\label{Molien Series}
\logpage{[ 72, 12, 0 ]}
\hyperdef{L}{X87B86B427A88CD25}{}
{
  

\subsection{\textcolor{Chapter }{MolienSeries}}
\logpage{[ 72, 12, 1 ]}\nobreak
\hyperdef{L}{X7D7F94D2820B1177}{}
{\noindent\textcolor{FuncColor}{$\triangleright$\enspace\texttt{MolienSeries({\mdseries\slshape [tbl, ]psi[, chi]})\index{MolienSeries@\texttt{MolienSeries}}
\label{MolienSeries}
}\hfill{\scriptsize (function)}}\\


 The \emph{Molien series} of the character $\psi$, relative to the character $\chi$, is the rational function given by the series $M_{{\psi,\chi}}(z) = \sum_{{d = 0}}^{\infty} [\chi,\psi^{[d]}] z^d$, where $\psi^{[d]}$ denotes the symmetrization of $\psi$ with the trivial character of the symmetric group $S_d$ (see{\nobreakspace}\texttt{SymmetricParts} (\ref{SymmetricParts})). 

 \texttt{MolienSeries} returns the Molien series of \mbox{\texttt{\mdseries\slshape psi}}, relative to \mbox{\texttt{\mdseries\slshape chi}}, where \mbox{\texttt{\mdseries\slshape psi}} and \mbox{\texttt{\mdseries\slshape chi}} must be characters of the same character table; this table must be entered as \mbox{\texttt{\mdseries\slshape tbl}} if \mbox{\texttt{\mdseries\slshape chi}} and \mbox{\texttt{\mdseries\slshape psi}} are only lists of character values. The default for \mbox{\texttt{\mdseries\slshape chi}} is the trivial character of \mbox{\texttt{\mdseries\slshape tbl}}. 

 The return value of \texttt{MolienSeries} stores a value for the attribute \texttt{MolienSeriesInfo} (\ref{MolienSeriesInfo}). This admits the computation of coefficients of the series with \texttt{ValueMolienSeries} (\ref{ValueMolienSeries}). Furthermore, this attribute gives access to numerator and denominator of the
Molien series viewed as rational function, where the denominator is a product
of polynomials of the form $(1-z^r)^k$; the Molien series is also displayed in this form. Note that such a
representation is not unique, one can use \texttt{MolienSeriesWithGivenDenominator} (\ref{MolienSeriesWithGivenDenominator}) to obtain the series with a prescribed denominator. 

 For more information about Molien series, see{\nobreakspace}\cite{NPP84}. 

 
\begin{Verbatim}[commandchars=!@|,fontsize=\small,frame=single,label=Example]
  !gapprompt@gap>| !gapinput@t:= CharacterTable( AlternatingGroup( 5 ) );;|
  !gapprompt@gap>| !gapinput@psi:= First( Irr( t ), x -> Degree( x ) = 3 );;|
  !gapprompt@gap>| !gapinput@mol:= MolienSeries( psi );|
  ( 1-z^2-z^3+z^6+z^7-z^9 ) / ( (1-z^5)*(1-z^3)*(1-z^2)^2 )
\end{Verbatim}
 }

 

\subsection{\textcolor{Chapter }{MolienSeriesInfo}}
\logpage{[ 72, 12, 2 ]}\nobreak
\hyperdef{L}{X82AC06A880EAA0AB}{}
{\noindent\textcolor{FuncColor}{$\triangleright$\enspace\texttt{MolienSeriesInfo({\mdseries\slshape ratfun})\index{MolienSeriesInfo@\texttt{MolienSeriesInfo}}
\label{MolienSeriesInfo}
}\hfill{\scriptsize (attribute)}}\\


 If the rational function \mbox{\texttt{\mdseries\slshape ratfun}} was constructed by \texttt{MolienSeries} (\ref{MolienSeries}), a representation as quotient of polynomials is known such that the
denominator is a product of terms of the form $(1-z^r)^k$. This information is encoded as value of \texttt{MolienSeriesInfo}. Additionally, there is a special \texttt{PrintObj} (\ref{PrintObj}) method for Molien series based on this. 

 \texttt{MolienSeriesInfo} returns a record that describes the rational function \mbox{\texttt{\mdseries\slshape ratfun}} as a Molien series. The components of this record are 
\begin{description}
\item[{\texttt{numer}}]  numerator of \mbox{\texttt{\mdseries\slshape ratfun}} (in general a multiple of the numerator one gets by \texttt{NumeratorOfRationalFunction} (\ref{NumeratorOfRationalFunction})), 
\item[{\texttt{denom}}]  denominator of \mbox{\texttt{\mdseries\slshape ratfun}} (in general a multiple of the denominator one gets by \texttt{NumeratorOfRationalFunction} (\ref{NumeratorOfRationalFunction})), 
\item[{\texttt{ratfun}}]  the rational function \mbox{\texttt{\mdseries\slshape ratfun}} itself, 
\item[{\texttt{numerstring}}]  string corresponding to the polynomial \texttt{numer}, expressed in terms of \texttt{z}, 
\item[{\texttt{denomstring}}]  string corresponding to the polynomial \texttt{denom}, expressed in terms of \texttt{z}, 
\item[{\texttt{denominfo}}]  a list of the form $[ [ r_1, k_1 ], \ldots, [ r_n, k_n ] ]$ such that \texttt{denom} is $\prod_{{i = 1}}^n (1-z^{{r_i}})^{{k_i}}$. 
\item[{\texttt{summands}}]  a list of records, each with the components \texttt{numer}, \texttt{r}, and \texttt{k}, describing the summand \texttt{numer}$ / (1-z^r)^k$, 
\item[{\texttt{pol}}]  a list of coefficients, describing a final polynomial which is added to those
described by \texttt{summands}, 
\item[{\texttt{size}}]  the order of the underlying matrix group, 
\item[{\texttt{degree}}]  the degree of the underlying matrix representation. 
\end{description}
 

 
\begin{Verbatim}[commandchars=!@|,fontsize=\small,frame=single,label=Example]
  !gapprompt@gap>| !gapinput@HasMolienSeriesInfo( mol );|
  true
  !gapprompt@gap>| !gapinput@MolienSeriesInfo( mol );|
  rec( degree := 3,
    denom := x_1^12-2*x_1^10-x_1^9+x_1^8+x_1^7+x_1^5+x_1^4-x_1^3-2*x_1^2\
  +1, denominfo := [ 5, 1, 3, 1, 2, 2 ],
    denomstring := "(1-z^5)*(1-z^3)*(1-z^2)^2",
    numer := -x_1^9+x_1^7+x_1^6-x_1^3-x_1^2+1,
    numerstring := "1-z^2-z^3+z^6+z^7-z^9", pol := [  ],
    ratfun := ( 1-z^2-z^3+z^6+z^7-z^9 ) / ( (1-z^5)*(1-z^3)*(1-z^2)^2 ),
    size := 60,
    summands := [ rec( k := 1, numer := [ -24, -12, -24 ], r := 5 ),
        rec( k := 1, numer := [ -20 ], r := 3 ),
        rec( k := 2, numer := [ -45/4, 75/4, -15/4, -15/4 ], r := 2 ),
        rec( k := 3, numer := [ -1 ], r := 1 ),
        rec( k := 1, numer := [ -15/4 ], r := 1 ) ] )
\end{Verbatim}
 }

 

\subsection{\textcolor{Chapter }{ValueMolienSeries}}
\logpage{[ 72, 12, 3 ]}\nobreak
\hyperdef{L}{X87083C4E7D11A02E}{}
{\noindent\textcolor{FuncColor}{$\triangleright$\enspace\texttt{ValueMolienSeries({\mdseries\slshape molser, i})\index{ValueMolienSeries@\texttt{ValueMolienSeries}}
\label{ValueMolienSeries}
}\hfill{\scriptsize (function)}}\\


 is the \mbox{\texttt{\mdseries\slshape i}}\texttt{\symbol{45}}th coefficient of the Molien series \mbox{\texttt{\mdseries\slshape series}} computed by \texttt{MolienSeries} (\ref{MolienSeries}). 

 
\begin{Verbatim}[commandchars=!@|,fontsize=\small,frame=single,label=Example]
  !gapprompt@gap>| !gapinput@List( [ 0 .. 20 ], i -> ValueMolienSeries( mol, i ) );|
  [ 1, 0, 1, 0, 1, 0, 2, 0, 2, 0, 3, 0, 4, 0, 4, 1, 5, 1, 6, 1, 7 ]
\end{Verbatim}
 }

 

\subsection{\textcolor{Chapter }{MolienSeriesWithGivenDenominator}}
\logpage{[ 72, 12, 4 ]}\nobreak
\hyperdef{L}{X86BAA3C487CE86D2}{}
{\noindent\textcolor{FuncColor}{$\triangleright$\enspace\texttt{MolienSeriesWithGivenDenominator({\mdseries\slshape molser, list})\index{MolienSeriesWithGivenDenominator@\texttt{MolienSeriesWithGivenDenominator}}
\label{MolienSeriesWithGivenDenominator}
}\hfill{\scriptsize (function)}}\\


 is a Molien series equal to \mbox{\texttt{\mdseries\slshape molser}} as rational function, but viewed as quotient with denominator $\prod_{{i = 1}}^n (1-z^{{r_i}})$, where $\mbox{\texttt{\mdseries\slshape list}} = [ r_1, r_2, \ldots, r_n ]$. If \mbox{\texttt{\mdseries\slshape molser}} cannot be represented this way, \texttt{fail} is returned. 

 
\begin{Verbatim}[commandchars=!@|,fontsize=\small,frame=single,label=Example]
  !gapprompt@gap>| !gapinput@MolienSeriesWithGivenDenominator( mol, [ 2, 6, 10 ] );|
  ( 1+z^15 ) / ( (1-z^10)*(1-z^6)*(1-z^2) )
\end{Verbatim}
 }

 }

  
\section{\textcolor{Chapter }{Possible Permutation Characters}}\label{Possible Permutation Characters}
\logpage{[ 72, 13, 0 ]}
\hyperdef{L}{X7D6336857E6BDF46}{}
{
  \index{characters!permutation} \index{candidates!for permutation characters} \index{possible permutation characters} \index{permutation characters!possible} For groups $H$ and $G$ with $H \leq G$, the induced character $(1_G)^H$ is called the \emph{permutation character} of the operation of $G$ on the right cosets of $H$. If only the character table of $G$ is available and not the group $G$ itself, one can try to get information about possible subgroups of $G$ by inspection of those $G$\texttt{\symbol{45}}class functions that might be permutation characters,
using that such a class function $\pi$ must have at least the following properties. (For details, see{\nobreakspace}\cite[Theorem 5.18.]{Isa76}), 
\begin{description}
\item[{(a)}]  $\pi$ is a character of $G$, 
\item[{(b)}]  $\pi(g)$ is a nonnegative integer for all $g \in G$, 
\item[{(c)}]  $\pi(1)$ divides $|G|$, 
\item[{(d)}]  $\pi(g^n) \geq \pi(g)$ for $g \in G$ and integers $n$, 
\item[{(e)}]  $[\pi, 1_G] = 1$, 
\item[{(f)}]  the multiplicity of any rational irreducible $G$\texttt{\symbol{45}}character $\psi$ as a constituent of $\pi$ is at most $\psi(1)/[\psi, \psi]$, 
\item[{(g)}]  $\pi(g) = 0$ if the order of $g$ does not divide $|G|/\pi(1)$, 
\item[{(h)}]  $\pi(1) |N_G(g)|$ divides $\pi(g) |G|$ for all $g \in G$, 
\item[{(i)}]  $\pi(g) \leq (|G| - \pi(1)) / (|g^G| |Gal_G(g)|)$ for all nonidentity $g \in G$, where $|Gal_G(g)|$ denotes the number of conjugacy classes of $G$ that contain generators of the group $\langle g \rangle$, 
\item[{(j)}]  if $p$ is a prime that divides $|G|/\pi(1)$ only once then $s/(p-1)$ divides $|G|/\pi(1)$ and is congruent to $1$ modulo $p$, where $s$ is the number of elements of order $p$ in the (hypothetical) subgroup $H$ for which $\pi = (1_H)^G$ holds. (Note that $s/(p-1)$ equals the number of Sylow $p$ subgroups in $H$.) 
\end{description}
 Any $G$\texttt{\symbol{45}}class function with these properties is called a \emph{possible permutation character} in \textsf{GAP}. 

 (Condition (d) is checked only for those power maps that are stored in the
character table of $G$; clearly (d) holds for all integers if it holds for all prime divisors of the
group order $|G|$.) 

 \textsf{GAP} provides some algorithms to compute possible permutation characters
(see{\nobreakspace}\texttt{PermChars} (\ref{PermChars})), and also provides functions to check a few more criteria whether a given
character can be a transitive permutation character (see{\nobreakspace}\texttt{TestPerm1} (\ref{TestPerm1})). 

 Some information about the subgroup $U$ can be computed from the permutation character $(1_U)^G$ using \texttt{PermCharInfo} (\ref{PermCharInfo}). \index{LaTeX!for permutation characters} 

\subsection{\textcolor{Chapter }{PermCharInfo}}
\logpage{[ 72, 13, 1 ]}\nobreak
\hyperdef{L}{X8477004C7A31D28C}{}
{\noindent\textcolor{FuncColor}{$\triangleright$\enspace\texttt{PermCharInfo({\mdseries\slshape tbl, permchars[, format]})\index{PermCharInfo@\texttt{PermCharInfo}}
\label{PermCharInfo}
}\hfill{\scriptsize (function)}}\\


 Let \mbox{\texttt{\mdseries\slshape tbl}} be the ordinary character table of the group $G$, and \mbox{\texttt{\mdseries\slshape permchars}} either the permutation character $(1_U)^G$, for a subgroup $U$ of $G$, or a list of such permutation characters. \texttt{PermCharInfo} returns a record with the following components. 
\begin{description}
\item[{\texttt{contained}:}]  a list containing, for each character $\psi = (1_U)^G$ in \mbox{\texttt{\mdseries\slshape permchars}}, a list containing at position $i$ the number $\psi[i] |U| /$ \texttt{SizesCentralizers( }\mbox{\texttt{\mdseries\slshape tbl}}\texttt{ )}$[i]$, which equals the number of those elements of $U$ that are contained in class $i$ of \mbox{\texttt{\mdseries\slshape tbl}}, 
\item[{\texttt{bound}:}]  a list containing, for each character $\psi = (1_U)^G$ in \mbox{\texttt{\mdseries\slshape permchars}}, a list containing at position $i$ the number $|U| / \gcd( |U|,$ \texttt{SizesCentralizers( \mbox{\texttt{\mdseries\slshape tbl}} )}$[i] )$, which divides the class length in $U$ of an element in class $i$ of \mbox{\texttt{\mdseries\slshape tbl}}, 
\item[{\texttt{display}:}]  a record that can be used as second argument of \texttt{Display} (\ref{Display}) to display each permutation character in \mbox{\texttt{\mdseries\slshape permchars}} and the corresponding components \texttt{contained} and \texttt{bound}, for those classes where at least one character of \mbox{\texttt{\mdseries\slshape permchars}} is nonzero, 
\item[{\texttt{ATLAS}:}]  a list of strings describing the decomposition of the permutation characters
in \mbox{\texttt{\mdseries\slshape permchars}} into the irreducible characters of \mbox{\texttt{\mdseries\slshape tbl}}, given in an \textsf{Atlas}\texttt{\symbol{45}}like notation. This means that the irreducible
constituents are indicated by their degrees followed by lower case letters \texttt{a}, \texttt{b}, \texttt{c}, $\ldots$, which indicate the successive irreducible characters of \mbox{\texttt{\mdseries\slshape tbl}} of that degree, in the order in which they appear in \texttt{Irr( }\mbox{\texttt{\mdseries\slshape tbl}}\texttt{ )}. A sequence of small letters (not necessarily distinct) after a single number
indicates a sum of irreducible constituents all of the same degree, an
exponent \mbox{\texttt{\mdseries\slshape n}} for the letter \mbox{\texttt{\mdseries\slshape lett}} means that \mbox{\texttt{\mdseries\slshape lett}} is repeated \mbox{\texttt{\mdseries\slshape n}} times. The default notation for exponentiation is \texttt{\mbox{\texttt{\mdseries\slshape lett}}\texttt{\symbol{94}}\texttt{\symbol{123}}\mbox{\texttt{\mdseries\slshape n}}\texttt{\symbol{125}}}, this is also chosen if the optional third argument \mbox{\texttt{\mdseries\slshape format}} is the string \texttt{"LaTeX"}; if the third argument is the string \texttt{"HTML"} then exponentiation is denoted by \texttt{\mbox{\texttt{\mdseries\slshape lett}}{\textless}sup{\textgreater}\mbox{\texttt{\mdseries\slshape n}}{\textless}/sup{\textgreater}}. 
\end{description}
 

 
\begin{Verbatim}[commandchars=!@|,fontsize=\small,frame=single,label=Example]
  !gapprompt@gap>| !gapinput@t:= CharacterTable( "A6" );;|
  !gapprompt@gap>| !gapinput@psi:= Sum( Irr( t ){ [ 1, 3, 6 ] } );|
  Character( CharacterTable( "A6" ), [ 15, 3, 0, 3, 1, 0, 0 ] )
  !gapprompt@gap>| !gapinput@info:= PermCharInfo( t, psi );|
  rec( ATLAS := [ "1a+5b+9a" ], bound := [ [ 1, 3, 8, 8, 6, 24, 24 ] ],
    contained := [ [ 1, 9, 0, 8, 6, 0, 0 ] ],
    display :=
      rec(
        chars := [ [ 15, 3, 0, 3, 1, 0, 0 ], [ 1, 9, 0, 8, 6, 0, 0 ],
            [ 1, 3, 8, 8, 6, 24, 24 ] ], classes := [ 1, 2, 4, 5 ],
        letter := "I" ) )
  !gapprompt@gap>| !gapinput@Display( t, info.display );|
  A6
  
       2  3  3  .  2
       3  2  .  2  .
       5  1  .  .  .
  
         1a 2a 3b 4a
      2P 1a 1a 3b 2a
      3P 1a 2a 1a 4a
      5P 1a 2a 3b 4a
  
  I.1    15  3  3  1
  I.2     1  9  8  6
  I.3     1  3  8  6
  !gapprompt@gap>| !gapinput@j1:= CharacterTable( "J1" );;|
  !gapprompt@gap>| !gapinput@psi:= TrivialCharacter( CharacterTable( "7:6" ) )^j1;|
  Character( CharacterTable( "J1" ), [ 4180, 20, 10, 0, 0, 2, 1, 0, 0,
    0, 0, 0, 0, 0, 0 ] )
  !gapprompt@gap>| !gapinput@PermCharInfo( j1, psi ).ATLAS;|
  [ "1a+56aabb+76aaab+77aabbcc+120aaabbbccc+133a^{4}bbcc+209a^{5}" ]
\end{Verbatim}
 }

 

\subsection{\textcolor{Chapter }{PermCharInfoRelative}}
\logpage{[ 72, 13, 2 ]}\nobreak
\hyperdef{L}{X7A8CB0298730D808}{}
{\noindent\textcolor{FuncColor}{$\triangleright$\enspace\texttt{PermCharInfoRelative({\mdseries\slshape tbl, tbl2, permchars})\index{PermCharInfoRelative@\texttt{PermCharInfoRelative}}
\label{PermCharInfoRelative}
}\hfill{\scriptsize (function)}}\\


 Let \mbox{\texttt{\mdseries\slshape tbl}} and \mbox{\texttt{\mdseries\slshape tbl2}} be the ordinary character tables of two groups $H$ and $G$, respectively, where $H$ is of index two in $G$, and \mbox{\texttt{\mdseries\slshape permchars}} either the permutation character $(1_U)^G$, for a subgroup $U$ of $G$, or a list of such permutation characters. \texttt{PermCharInfoRelative} returns a record with the same components as \texttt{PermCharInfo} (\ref{PermCharInfo}), the only exception is that the entries of the \texttt{ATLAS} component are names relative to \mbox{\texttt{\mdseries\slshape tbl}}. 

 More precisely, the $i$\texttt{\symbol{45}}th entry of the \texttt{ATLAS} component is a string describing the decomposition of the $i$\texttt{\symbol{45}}th entry in \mbox{\texttt{\mdseries\slshape permchars}}. The degrees and distinguishing letters of the constituents refer to the
irreducibles of \mbox{\texttt{\mdseries\slshape tbl}}, as follows. The two irreducible characters of \mbox{\texttt{\mdseries\slshape tbl2}} of degree $N$ that extend the irreducible character $N$ \texttt{a} of \mbox{\texttt{\mdseries\slshape tbl}} are denoted by $N$ \texttt{a}$^+$ and $N $\texttt{a}$^-$. The irreducible character of \mbox{\texttt{\mdseries\slshape tbl2}} of degree $2N$ whose restriction to \mbox{\texttt{\mdseries\slshape tbl}} is the sum of the irreducible characters $N$ \texttt{a} and $N$ \texttt{b} is denoted as $N$ \texttt{ab}. Multiplicities larger than $1$ of constituents are denoted by exponents. 

 (This format is useful mainly for multiplicity free permutation characters.) 

 
\begin{Verbatim}[commandchars=!@|,fontsize=\small,frame=single,label=Example]
  !gapprompt@gap>| !gapinput@t:= CharacterTable( "A5" );;|
  !gapprompt@gap>| !gapinput@t2:= CharacterTable( "A5.2" );;|
  !gapprompt@gap>| !gapinput@List( Irr( t2 ), x -> x[1] );|
  [ 1, 1, 6, 4, 4, 5, 5 ]
  !gapprompt@gap>| !gapinput@List( Irr( t ), x -> x[1] );|
  [ 1, 3, 3, 4, 5 ]
  !gapprompt@gap>| !gapinput@permchars:= List( [ [1], [1,2], [1,7], [1,3,4,4,6,6,7] ],|
  !gapprompt@>| !gapinput@                     l -> Sum( Irr( t2 ){ l } ) );|
  [ Character( CharacterTable( "A5.2" ), [ 1, 1, 1, 1, 1, 1, 1 ] ),
    Character( CharacterTable( "A5.2" ), [ 2, 2, 2, 2, 0, 0, 0 ] ),
    Character( CharacterTable( "A5.2" ), [ 6, 2, 0, 1, 0, 2, 0 ] ),
    Character( CharacterTable( "A5.2" ), [ 30, 2, 0, 0, 6, 0, 0 ] ) ]
  !gapprompt@gap>| !gapinput@info:= PermCharInfoRelative( t, t2, permchars );;|
  !gapprompt@gap>| !gapinput@info.ATLAS;|
  [ "1a^+", "1a^{\\pm}", "1a^++5a^-",
    "1a^++3ab+4(a^+)^{2}+5a^+a^{\\pm}" ]
\end{Verbatim}
 }

 }

  
\section{\textcolor{Chapter }{Computing Possible Permutation Characters}}\label{Computing Possible Permutation Characters}
\logpage{[ 72, 14, 0 ]}
\hyperdef{L}{X8330FDCE83D3DED3}{}
{
  

\subsection{\textcolor{Chapter }{PermChars}}
\logpage{[ 72, 14, 1 ]}\nobreak
\hyperdef{L}{X7D02541482C196A6}{}
{\noindent\textcolor{FuncColor}{$\triangleright$\enspace\texttt{PermChars({\mdseries\slshape tbl[, cond]})\index{PermChars@\texttt{PermChars}}
\label{PermChars}
}\hfill{\scriptsize (function)}}\\


 \textsf{GAP} provides several algorithms to determine possible permutation characters from
a given character table. They are described in detail in{\nobreakspace}\cite{BP98}. The algorithm is selected from the choice of the optional argument \mbox{\texttt{\mdseries\slshape cond}}. The user is encouraged to try different approaches, especially if one choice
fails to come to an end. 

 Regardless of the algorithm used in a specific case, \texttt{PermChars} returns a list of \emph{all} possible permutation characters with the properties described by \mbox{\texttt{\mdseries\slshape cond}}. There is no guarantee that a character of this list is in fact a permutation
character. But an empty list always means there is no permutation character
with these properties (e.g., of a certain degree). 

 Called with only one argument, a character table \mbox{\texttt{\mdseries\slshape tbl}}, \texttt{PermChars} returns the list of all possible permutation characters of the group with this
character table. This list might be rather long for big groups, and its
computation might take much time. The algorithm is described in \cite[Section 3.2]{BP98}; it depends on a preprocessing step, where the inequalities arising from the
condition $\pi(g) \geq 0$ are transformed into a system of inequalities that guides the search
(see{\nobreakspace}\texttt{Inequalities} (\ref{Inequalities})). So the following commands compute the list of 39 possible permutation
characters of the Mathieu group $M_{11}$. 

 
\begin{Verbatim}[commandchars=!@|,fontsize=\small,frame=single,label=Example]
  !gapprompt@gap>| !gapinput@m11:= CharacterTable( "M11" );;|
  !gapprompt@gap>| !gapinput@SetName( m11, "m11" );|
  !gapprompt@gap>| !gapinput@perms:= PermChars( m11 );;|
  !gapprompt@gap>| !gapinput@Length( perms );|
  39
\end{Verbatim}
 

 There are two different search strategies for this algorithm. The default
strategy simply constructs all characters with nonnegative values and then
tests for each such character whether its degree is a divisor of the order of
the group. The other strategy uses the inequalities to predict whether a
character of a certain degree can lie in the currently searched part of the
search tree. To choose this strategy, enter a record as the second argument of \texttt{PermChars}, and set its component \texttt{degree} to the range of degrees (which might also be a range containing all divisors
of the group order) you want to look for; additionally, the record component \texttt{ineq} can take the inequalities computed by \texttt{Inequalities} (\ref{Inequalities}) if they are needed more than once. 

 If a positive integer is given as the second argument \mbox{\texttt{\mdseries\slshape cond}}, \texttt{PermChars} returns the list of all possible permutation characters of \mbox{\texttt{\mdseries\slshape tbl}} that have degree \mbox{\texttt{\mdseries\slshape cond}}. For that purpose, a preprocessing step is performed where essentially the
rational character table is inverted in order to determine boundary points for
the simplex in which the possible permutation characters of the given degree
must lie (see{\nobreakspace}\texttt{PermBounds} (\ref{PermBounds})). The algorithm is described at the end of \cite[Section 3.2]{BP98}. Note that inverting big integer matrices needs a lot of time and space. So
this preprocessing is restricted to groups with less than 100 classes, say. 

 
\begin{Verbatim}[commandchars=!@|,fontsize=\small,frame=single,label=Example]
  !gapprompt@gap>| !gapinput@deg220:= PermChars( m11, 220 );|
  [ Character( m11, [ 220, 4, 4, 0, 0, 4, 0, 0, 0, 0 ] ),
    Character( m11, [ 220, 12, 4, 4, 0, 0, 0, 0, 0, 0 ] ),
    Character( m11, [ 220, 20, 4, 0, 0, 2, 0, 0, 0, 0 ] ) ]
\end{Verbatim}
 

 If a record is given as the second argument \mbox{\texttt{\mdseries\slshape cond}}, \texttt{PermChars} returns the list of all possible permutation characters that have the
properties described by the components of this record. One such situation has
been mentioned above. If \mbox{\texttt{\mdseries\slshape cond}} contains a degree as value of the record component \texttt{degree} then \texttt{PermChars} will behave exactly as if this degree was entered as \mbox{\texttt{\mdseries\slshape cond}}. 

 
\begin{Verbatim}[commandchars=!@|,fontsize=\small,frame=single,label=Example]
  !gapprompt@gap>| !gapinput@deg220 = PermChars( m11, rec( degree:= 220 ) );|
  true
\end{Verbatim}
 

 For the meaning of additional components of \mbox{\texttt{\mdseries\slshape cond}} besides \texttt{degree}, see{\nobreakspace}\texttt{PermComb} (\ref{PermComb}). 

 Instead of \texttt{degree}, \mbox{\texttt{\mdseries\slshape cond}} may have the component \texttt{torso} bound to a list that contains some known values of the required characters at
the right positions; at least the degree \mbox{\texttt{\mdseries\slshape cond}}\texttt{.torso[1]} must be an integer. In this case, the algorithm described in \cite[Section 3.3]{BP98} is chosen. The component \texttt{chars}, if present, holds a list of all those \emph{rational} irreducible characters of \mbox{\texttt{\mdseries\slshape tbl}} that might be constituents of the required characters. 

 (\emph{Note}: If \mbox{\texttt{\mdseries\slshape cond}}\texttt{.chars} is bound and does not contain \emph{all} rational irreducible characters of \mbox{\texttt{\mdseries\slshape tbl}}, \textsf{GAP} checks whether the scalar products of all class functions in the result list
with the omitted rational irreducible characters of \mbox{\texttt{\mdseries\slshape tbl}} are nonnegative; so there should be nontrivial reasons for excluding a
character that is known to be not a constituent of the desired possible
permutation characters.) 

 
\begin{Verbatim}[commandchars=!@|,fontsize=\small,frame=single,label=Example]
  !gapprompt@gap>| !gapinput@PermChars( m11, rec( torso:= [ 220 ] ) );|
  [ Character( m11, [ 220, 4, 4, 0, 0, 4, 0, 0, 0, 0 ] ),
    Character( m11, [ 220, 20, 4, 0, 0, 2, 0, 0, 0, 0 ] ),
    Character( m11, [ 220, 12, 4, 4, 0, 0, 0, 0, 0, 0 ] ) ]
  !gapprompt@gap>| !gapinput@PermChars( m11, rec( torso:= [ 220,,,,, 2 ] ) );|
  [ Character( m11, [ 220, 20, 4, 0, 0, 2, 0, 0, 0, 0 ] ) ]
\end{Verbatim}
 

 An additional restriction on the possible permutation characters computed can
be forced if \mbox{\texttt{\mdseries\slshape con}} contains, in addition to \texttt{torso}, the components \texttt{normalsubgroup} and \texttt{nonfaithful}, with values a list of class positions of a normal subgroup $N$ of the group $G$ of \mbox{\texttt{\mdseries\slshape tbl}} and a possible permutation character $\pi$ of $G$, respectively, such that $N$ is contained in the kernel of $\pi$. In this case, \texttt{PermChars} returns the list of those possible permutation characters $\psi$ of \mbox{\texttt{\mdseries\slshape tbl}} coinciding with \texttt{torso} wherever its values are bound and having the property that no irreducible
constituent of $\psi - \pi$ has $N$ in its kernel. If the component \texttt{chars} is bound in \mbox{\texttt{\mdseries\slshape cond}} then the above statements apply. An interpretation of the computed characters
is the following. Suppose there exists a subgroup $V$ of $G$ such that $\pi = (1_V)^G$; Then $N \leq V$, and if a computed character is of the form $(1_U)^G$, for a subgroup $U$ of $G$, then $V = UN$. 

 
\begin{Verbatim}[commandchars=!@|,fontsize=\small,frame=single,label=Example]
  !gapprompt@gap>| !gapinput@s4:= CharacterTable( "Symmetric", 4 );;|
  !gapprompt@gap>| !gapinput@nsg:= ClassPositionsOfDerivedSubgroup( s4 );;|
  !gapprompt@gap>| !gapinput@pi:= TrivialCharacter( s4 );;|
  !gapprompt@gap>| !gapinput@PermChars( s4, rec( torso:= [ 12 ], normalsubgroup:= nsg,|
  !gapprompt@>| !gapinput@                       nonfaithful:= pi ) );|
  [ Character( CharacterTable( "Sym(4)" ), [ 12, 2, 0, 0, 0 ] ) ]
  !gapprompt@gap>| !gapinput@pi:= Sum( Filtered( Irr( s4 ),|
  !gapprompt@>| !gapinput@             chi -> IsSubset( ClassPositionsOfKernel( chi ), nsg ) ) );|
  Character( CharacterTable( "Sym(4)" ), [ 2, 0, 2, 2, 0 ] )
  !gapprompt@gap>| !gapinput@PermChars( s4, rec( torso:= [ 12 ], normalsubgroup:= nsg,|
  !gapprompt@>| !gapinput@                       nonfaithful:= pi ) );|
  [ Character( CharacterTable( "Sym(4)" ), [ 12, 0, 4, 0, 0 ] ) ]
\end{Verbatim}
 

 The class functions returned by \texttt{PermChars} have the properties tested by \texttt{TestPerm1} (\ref{TestPerm1}), \texttt{TestPerm2} (\ref{TestPerm2}), and \texttt{TestPerm3} (\ref{TestPerm3}). So they are possible permutation characters. See{\nobreakspace}\texttt{TestPerm1} (\ref{TestPerm1}) for criteria whether a possible permutation character can in fact be a
permutation character. }

 
\subsection{\textcolor{Chapter }{TestPerm1, ..., TestPerm5}}\logpage{[ 72, 14, 2 ]}
\hyperdef{L}{X8127771D7EAB6EA7}{}
{
\noindent\textcolor{FuncColor}{$\triangleright$\enspace\texttt{TestPerm1({\mdseries\slshape tbl, char})\index{TestPerm1@\texttt{TestPerm1}}
\label{TestPerm1}
}\hfill{\scriptsize (function)}}\\
\noindent\textcolor{FuncColor}{$\triangleright$\enspace\texttt{TestPerm2({\mdseries\slshape tbl, char})\index{TestPerm2@\texttt{TestPerm2}}
\label{TestPerm2}
}\hfill{\scriptsize (function)}}\\
\noindent\textcolor{FuncColor}{$\triangleright$\enspace\texttt{TestPerm3({\mdseries\slshape tbl, chars})\index{TestPerm3@\texttt{TestPerm3}}
\label{TestPerm3}
}\hfill{\scriptsize (function)}}\\
\noindent\textcolor{FuncColor}{$\triangleright$\enspace\texttt{TestPerm4({\mdseries\slshape tbl, chars})\index{TestPerm4@\texttt{TestPerm4}}
\label{TestPerm4}
}\hfill{\scriptsize (function)}}\\
\noindent\textcolor{FuncColor}{$\triangleright$\enspace\texttt{TestPerm5({\mdseries\slshape tbl, chars, modtbl})\index{TestPerm5@\texttt{TestPerm5}}
\label{TestPerm5}
}\hfill{\scriptsize (function)}}\\


 The first three of these functions implement tests of the properties of
possible permutation characters listed in Section{\nobreakspace}\ref{Possible Permutation Characters}, The other two implement test of additional properties. Let \mbox{\texttt{\mdseries\slshape tbl}} be the ordinary character table of a group $G$, \mbox{\texttt{\mdseries\slshape char}} a rational character of \mbox{\texttt{\mdseries\slshape tbl}}, and \mbox{\texttt{\mdseries\slshape chars}} a list of rational characters of \mbox{\texttt{\mdseries\slshape tbl}}. For applying \texttt{TestPerm5}, the knowledge of a $p$\texttt{\symbol{45}}modular Brauer table \mbox{\texttt{\mdseries\slshape modtbl}} of $G$ is required. \texttt{TestPerm4} and \texttt{TestPerm5} expect the characters in \mbox{\texttt{\mdseries\slshape chars}} to satisfy the conditions checked by \texttt{TestPerm1} and \texttt{TestPerm2} (see below). 

 The return values of the functions were chosen parallel to the tests listed
in{\nobreakspace}\cite{NPP84}. 

 \texttt{TestPerm1} return \texttt{1} or \texttt{2} if \mbox{\texttt{\mdseries\slshape char}} fails because of (T1) or (T2), respectively; this corresponds to the criteria
(b) and (d). Note that only those power maps are considered that are stored on \mbox{\texttt{\mdseries\slshape tbl}}. If \mbox{\texttt{\mdseries\slshape char}} satisfies the conditions, \texttt{0} is returned. 

 \texttt{TestPerm2} returns \texttt{1} if \mbox{\texttt{\mdseries\slshape char}} fails because of the criterion (c), it returns \texttt{3}, \texttt{4}, or \texttt{5} if \mbox{\texttt{\mdseries\slshape char}} fails because of (T3), (T4), or (T5), respectively; these tests correspond to
(g), a weaker form of (h), and (j). If \mbox{\texttt{\mdseries\slshape char}} satisfies the conditions, \texttt{0} is returned. 

 \texttt{TestPerm3} returns the list of all those class functions in the list \mbox{\texttt{\mdseries\slshape chars}} that satisfy criterion (h); this is a stronger version of (T6). 

 \texttt{TestPerm4} returns the list of all those class functions in the list \mbox{\texttt{\mdseries\slshape chars}} that satisfy (T8) and (T9) for each prime divisor $p$ of the order of $G$; these tests use modular representation theory but do not require the
knowledge of decomposition matrices (cf.{\nobreakspace}\texttt{TestPerm5} below). 

 (T8) implements the test of the fact that in the case that $p$ divides $|G|$ and the degree of a transitive permutation character $\pi$ exactly once, the projective cover of the trivial character is a summand of $\pi$. (This test is omitted if the projective cover cannot be identified.) 

 Given a permutation character $\pi$ of a group $G$ and a prime integer $p$, the restriction $\pi_B$ to a $p$\texttt{\symbol{45}}block $B$ of $G$ has the following property, which is checked by (T9). For each $g \in G$ such that $g^n$ is a $p$\texttt{\symbol{45}}element of $G$, $\pi_B(g^n)$ is a nonnegative integer that satisfies $|\pi_B(g)| \leq \pi_B(g^n) \leq \pi(g^n)$. (This is \cite[Corollary A on p. 113]{Sco73}.) 

 \texttt{TestPerm5} requires the $p$\texttt{\symbol{45}}modular Brauer table \mbox{\texttt{\mdseries\slshape modtbl}} of $G$, for some prime $p$ dividing the order of $G$, and checks whether those characters in the list \mbox{\texttt{\mdseries\slshape chars}} whose degree is divisible by the $p$\texttt{\symbol{45}}part of the order of $G$ can be decomposed into projective indecomposable characters; \texttt{TestPerm5} returns the sublist of all those characters in \mbox{\texttt{\mdseries\slshape chars}} that either satisfy this condition or to which the test does not apply. 

       
\begin{Verbatim}[commandchars=!@|,fontsize=\small,frame=single,label=Example]
  !gapprompt@gap>| !gapinput@tbl:= CharacterTable( "A5" );;|
  !gapprompt@gap>| !gapinput@rat:= RationalizedMat( Irr( tbl ) );|
  [ Character( CharacterTable( "A5" ), [ 1, 1, 1, 1, 1 ] ),
    Character( CharacterTable( "A5" ), [ 6, -2, 0, 1, 1 ] ),
    Character( CharacterTable( "A5" ), [ 4, 0, 1, -1, -1 ] ),
    Character( CharacterTable( "A5" ), [ 5, 1, -1, 0, 0 ] ) ]
  !gapprompt@gap>| !gapinput@tup:= Filtered( Tuples( [ 0, 1 ], 4 ), x -> not IsZero( x ) );|
  [ [ 0, 0, 0, 1 ], [ 0, 0, 1, 0 ], [ 0, 0, 1, 1 ], [ 0, 1, 0, 0 ],
    [ 0, 1, 0, 1 ], [ 0, 1, 1, 0 ], [ 0, 1, 1, 1 ], [ 1, 0, 0, 0 ],
    [ 1, 0, 0, 1 ], [ 1, 0, 1, 0 ], [ 1, 0, 1, 1 ], [ 1, 1, 0, 0 ],
    [ 1, 1, 0, 1 ], [ 1, 1, 1, 0 ], [ 1, 1, 1, 1 ] ]
  !gapprompt@gap>| !gapinput@lincomb:= List( tup, coeff -> coeff * rat );;|
  !gapprompt@gap>| !gapinput@List( lincomb, psi -> TestPerm1( tbl, psi ) );|
  [ 1, 1, 1, 1, 1, 1, 1, 0, 0, 0, 0, 1, 0, 1, 0 ]
  !gapprompt@gap>| !gapinput@List( lincomb, psi -> TestPerm2( tbl, psi ) );|
  [ 0, 5, 1, 0, 1, 0, 0, 0, 0, 0, 0, 1, 0, 1, 1 ]
  !gapprompt@gap>| !gapinput@Set( TestPerm3(tbl, lincomb), x -> Position(lincomb, x) );|
  [ 1, 4, 6, 7, 8, 9, 10, 11, 13 ]
  !gapprompt@gap>| !gapinput@tbl:= CharacterTable( "A7" );|
  CharacterTable( "A7" )
  !gapprompt@gap>| !gapinput@perms:= PermChars( tbl, rec( degree:= 315 ) );|
  [ Character( CharacterTable( "A7" ), [ 315, 3, 0, 0, 3, 0, 0, 0, 0 ] )
      , Character( CharacterTable( "A7" ),
      [ 315, 15, 0, 0, 1, 0, 0, 0, 0 ] ) ]
  !gapprompt@gap>| !gapinput@TestPerm4( tbl, perms );|
  [ Character( CharacterTable( "A7" ), [ 315, 15, 0, 0, 1, 0, 0, 0, 0
       ] ) ]
  !gapprompt@gap>| !gapinput@perms:= PermChars( tbl, rec( degree:= 15 ) );|
  [ Character( CharacterTable( "A7" ), [ 15, 3, 0, 3, 1, 0, 0, 1, 1 ] ),
    Character( CharacterTable( "A7" ), [ 15, 3, 3, 0, 1, 0, 3, 1, 1 ] )
   ]
  !gapprompt@gap>| !gapinput@TestPerm5( tbl, perms, tbl mod 5 );|
  [ Character( CharacterTable( "A7" ), [ 15, 3, 0, 3, 1, 0, 0, 1, 1 ] )
   ]
\end{Verbatim}
 }

 

\subsection{\textcolor{Chapter }{PermBounds}}
\logpage{[ 72, 14, 3 ]}\nobreak
\hyperdef{L}{X879D2A127BE366A5}{}
{\noindent\textcolor{FuncColor}{$\triangleright$\enspace\texttt{PermBounds({\mdseries\slshape tbl, d})\index{PermBounds@\texttt{PermBounds}}
\label{PermBounds}
}\hfill{\scriptsize (function)}}\\


 Let \mbox{\texttt{\mdseries\slshape tbl}} be the ordinary character table of the group $G$. All $G$\texttt{\symbol{45}}characters $\pi$ satisfying $\pi(g) > 0$ and $\pi(1) = \mbox{\texttt{\mdseries\slshape d}}$, for a given degree \mbox{\texttt{\mdseries\slshape d}}, lie in a simplex described by these conditions. \texttt{PermBounds} computes the boundary points of this simplex for $d = 0$, from which the boundary points for any other \mbox{\texttt{\mdseries\slshape d}} are easily derived. (Some conditions from the power maps of \mbox{\texttt{\mdseries\slshape tbl}} are also involved.) For this purpose, a matrix similar to the rational
character table of $G$ has to be inverted. These boundary points are used by \texttt{PermChars} (\ref{PermChars}) to construct all possible permutation characters (see{\nobreakspace}\ref{Possible Permutation Characters}) of a given degree. \texttt{PermChars} (\ref{PermChars}) either calls \texttt{PermBounds} or takes this information from the \texttt{bounds} component of its argument record. }

 

\subsection{\textcolor{Chapter }{PermComb}}
\logpage{[ 72, 14, 4 ]}\nobreak
\hyperdef{L}{X7F11AFB783352903}{}
{\noindent\textcolor{FuncColor}{$\triangleright$\enspace\texttt{PermComb({\mdseries\slshape tbl, arec})\index{PermComb@\texttt{PermComb}}
\label{PermComb}
}\hfill{\scriptsize (function)}}\\


 \texttt{PermComb} computes possible permutation characters of the character table \mbox{\texttt{\mdseries\slshape tbl}} by the improved combinatorial approach described at the end of \cite[Section 3.2]{BP98}. 

 For computing the possible linear combinations \emph{without} prescribing better bounds (i.e., when the computation of bounds shall be
suppressed), enter 

 \texttt{\mbox{\texttt{\mdseries\slshape arec}}:= rec( degree := \mbox{\texttt{\mdseries\slshape degree}}, bounds := false )}, 

 where \mbox{\texttt{\mdseries\slshape degree}} is the character degree; this is useful if the multiplicities are expected to
be small, and if this is forced by high irreducible degrees. 

 A list of upper bounds on the multiplicities of the rational irreducibles
characters can be explicitly prescribed as a \texttt{maxmult} component in \mbox{\texttt{\mdseries\slshape arec}}. }

 

\subsection{\textcolor{Chapter }{Inequalities}}
\logpage{[ 72, 14, 5 ]}\nobreak
\hyperdef{L}{X866942167802E036}{}
{\noindent\textcolor{FuncColor}{$\triangleright$\enspace\texttt{Inequalities({\mdseries\slshape tbl, chars[, option]})\index{Inequalities@\texttt{Inequalities}}
\label{Inequalities}
}\hfill{\scriptsize (operation)}}\\


 Let \mbox{\texttt{\mdseries\slshape tbl}} be the ordinary character table of a group $G$. The condition $\pi(g) \geq 0$ for every possible permutation character $\pi$ of $G$ places restrictions on the multiplicities $a_i$ of the irreducible constituents $\chi_i$ of $\pi = \sum_{{i = 1}}^r a_i \chi_i$. For every element $g \in G$, we have $\sum_{{i = 1}}^r a_i \chi_i(g) \geq 0$. The power maps provide even stronger conditions. 

 This system of inequalities is kind of diagonalized, resulting in a system of
inequalities restricting $a_i$ in terms of $a_j$, $j < i$. These inequalities are used to construct characters with nonnegative values
(see{\nobreakspace}\texttt{PermChars} (\ref{PermChars})). \texttt{PermChars} (\ref{PermChars}) either calls \texttt{Inequalities} or takes this information from the \texttt{ineq} component of its argument record. 

 The number of inequalities arising in the process of diagonalization may grow
very strongly. 

 There are two ways to organize the projection. The first, which is chosen if
no \mbox{\texttt{\mdseries\slshape option}} argument is present, is the straight approach which takes the rational
irreducible characters in their original order and by this guarantees the
character with the smallest degree to be considered first. The other way,
which is chosen if the string \texttt{"small"} is entered as third argument \mbox{\texttt{\mdseries\slshape option}}, tries to keep the number of intermediate inequalities small by eventually
changing the order of characters. 

 
\begin{Verbatim}[commandchars=!@|,fontsize=\small,frame=single,label=Example]
  !gapprompt@gap>| !gapinput@tbl:= CharacterTable( "M11" );;|
  !gapprompt@gap>| !gapinput@PermComb( tbl, rec( degree:= 110 ) );|
  [ Character( CharacterTable( "M11" ),
      [ 110, 6, 2, 2, 0, 0, 2, 2, 0, 0 ] ),
    Character( CharacterTable( "M11" ),
      [ 110, 6, 2, 6, 0, 0, 0, 0, 0, 0 ] ),
    Character( CharacterTable( "M11" ), [ 110, 14, 2, 2, 0, 2, 0, 0, 0,
        0 ] ) ]
  !gapprompt@gap>| !gapinput@# Now compute only multiplicity free permutation characters.|
  !gapprompt@gap>| !gapinput@bounds:= List( RationalizedMat( Irr( tbl ) ), x -> 1 );;|
  !gapprompt@gap>| !gapinput@PermComb( tbl, rec( degree:= 110, maxmult:= bounds ) );|
  [ Character( CharacterTable( "M11" ),
      [ 110, 6, 2, 2, 0, 0, 2, 2, 0, 0 ] ) ]
\end{Verbatim}
 }

 }

  
\section{\textcolor{Chapter }{Operations for Brauer Characters}}\label{Operations for Brauer Characters}
\logpage{[ 72, 15, 0 ]}
\hyperdef{L}{X8204FB9F847340C8}{}
{
  

\subsection{\textcolor{Chapter }{FrobeniusCharacterValue}}
\logpage{[ 72, 15, 1 ]}\nobreak
\hyperdef{L}{X79BACBC47B4C413E}{}
{\noindent\textcolor{FuncColor}{$\triangleright$\enspace\texttt{FrobeniusCharacterValue({\mdseries\slshape value, p})\index{FrobeniusCharacterValue@\texttt{FrobeniusCharacterValue}}
\label{FrobeniusCharacterValue}
}\hfill{\scriptsize (function)}}\\


 Let \mbox{\texttt{\mdseries\slshape value}} be a cyclotomic whose coefficients over the rationals are in the ring ${\ensuremath{\mathbb Z}}_{\mbox{\texttt{\mdseries\slshape p}}}$ of \mbox{\texttt{\mdseries\slshape p}}\texttt{\symbol{45}}local numbers, where \mbox{\texttt{\mdseries\slshape p}} is a prime integer. Assume that \mbox{\texttt{\mdseries\slshape value}} lies in ${\ensuremath{\mathbb Z}}_{\mbox{\texttt{\mdseries\slshape p}}}[\zeta]$ for $\zeta = \exp(\mbox{\texttt{\mdseries\slshape p}}^n-1)$, for some positive integer $n$. 

 \texttt{FrobeniusCharacterValue} returns the image of \mbox{\texttt{\mdseries\slshape value}} under the ring homomorphism from ${\ensuremath{\mathbb Z}}_{\mbox{\texttt{\mdseries\slshape p}}}[\zeta]$ to the field with $\mbox{\texttt{\mdseries\slshape p}}^n$ elements that is defined with the help of Conway polynomials
(see{\nobreakspace}\texttt{ConwayPolynomial} (\ref{ConwayPolynomial})), more information can be found in \cite[Sections 2\texttt{\symbol{45}}5]{JLPW95}. 

 If \mbox{\texttt{\mdseries\slshape value}} is a Brauer character value in characteristic \mbox{\texttt{\mdseries\slshape p}} then the result can be described as the corresponding value of the Frobenius
character, that is, as the trace of a representing matrix with the given
Brauer character value. 

 If the result of \texttt{FrobeniusCharacterValue} cannot be expressed as an element of a finite field in \textsf{GAP} (see Chapter{\nobreakspace}\ref{Finite Fields}) then \texttt{FrobeniusCharacterValue} returns \texttt{fail}. 

 If the Conway polynomial of degree $n$ is required for the computation then it is computed only if \texttt{IsCheapConwayPolynomial} (\ref{IsCheapConwayPolynomial}) returns \texttt{true} when it is called with \mbox{\texttt{\mdseries\slshape p}} and $n$, otherwise \texttt{fail} is returned. }

 

\subsection{\textcolor{Chapter }{BrauerCharacterValue}}
\logpage{[ 72, 15, 2 ]}\nobreak
\hyperdef{L}{X8304B68E84511685}{}
{\noindent\textcolor{FuncColor}{$\triangleright$\enspace\texttt{BrauerCharacterValue({\mdseries\slshape mat})\index{BrauerCharacterValue@\texttt{BrauerCharacterValue}}
\label{BrauerCharacterValue}
}\hfill{\scriptsize (attribute)}}\\


 For an invertible matrix \mbox{\texttt{\mdseries\slshape mat}} over a finite field $F$, \texttt{BrauerCharacterValue} returns the Brauer character value of \mbox{\texttt{\mdseries\slshape mat}} if the order of \mbox{\texttt{\mdseries\slshape mat}} is coprime to the characteristic of $F$, and \texttt{fail} otherwise. 

 The \emph{Brauer character value} of a matrix is the sum of complex lifts of its eigenvalues. 

 
\begin{Verbatim}[commandchars=!@|,fontsize=\small,frame=single,label=Example]
  !gapprompt@gap>| !gapinput@g:= SL(2,4);;           # 2-dim. irreducible representation of A5|
  !gapprompt@gap>| !gapinput@ccl:= ConjugacyClasses( g );;|
  !gapprompt@gap>| !gapinput@rep:= List( ccl, Representative );;|
  !gapprompt@gap>| !gapinput@List( rep, Order );|
  [ 1, 2, 5, 5, 3 ]
  !gapprompt@gap>| !gapinput@phi:= List( rep, BrauerCharacterValue );|
  [ 2, fail, E(5)^2+E(5)^3, E(5)+E(5)^4, -1 ]
  !gapprompt@gap>| !gapinput@List( phi{ [ 1, 3, 4, 5 ] }, x -> FrobeniusCharacterValue( x, 2 ) );|
  [ 0*Z(2), Z(2^2), Z(2^2)^2, Z(2)^0 ]
  !gapprompt@gap>| !gapinput@List( rep{ [ 1, 3, 4, 5 ] }, TraceMat );|
  [ 0*Z(2), Z(2^2), Z(2^2)^2, Z(2)^0 ]
\end{Verbatim}
 }

 

\subsection{\textcolor{Chapter }{SizeOfFieldOfDefinition}}
\logpage{[ 72, 15, 3 ]}\nobreak
\hyperdef{L}{X8038FA0480B78243}{}
{\noindent\textcolor{FuncColor}{$\triangleright$\enspace\texttt{SizeOfFieldOfDefinition({\mdseries\slshape val, p})\index{SizeOfFieldOfDefinition@\texttt{SizeOfFieldOfDefinition}}
\label{SizeOfFieldOfDefinition}
}\hfill{\scriptsize (function)}}\\


 For a cyclotomic or a list of cyclotomics \mbox{\texttt{\mdseries\slshape val}}, and a prime integer \mbox{\texttt{\mdseries\slshape p}}, \texttt{SizeOfFieldOfDefinition} returns the size of the smallest finite field in characteristic \mbox{\texttt{\mdseries\slshape p}} that contains the \mbox{\texttt{\mdseries\slshape p}}\texttt{\symbol{45}}modular reduction of \mbox{\texttt{\mdseries\slshape val}} if this can be determined, and \texttt{fail} otherwise. 

 The latter happens if \mbox{\texttt{\mdseries\slshape val}} is not closed under Galois conjugacy and if the \mbox{\texttt{\mdseries\slshape p}}\texttt{\symbol{45}}modular reduction of some value cannot be determined via
the function \texttt{FrobeniusCharacterValue} (\ref{FrobeniusCharacterValue}). Note that the reduction map is defined as in{\nobreakspace}\cite{JLPW95}, that is, the complex $(\mbox{\texttt{\mdseries\slshape p}}^d-1)$\texttt{\symbol{45}}th root of unity $\exp(\mbox{\texttt{\mdseries\slshape p}}^d-1)$ is mapped to the residue class of the indeterminate, modulo the ideal spanned
by the Conway polynomial (see{\nobreakspace}\texttt{ConwayPolynomial} (\ref{ConwayPolynomial})) of degree $d$ over the field with $p$ elements. 

 If \mbox{\texttt{\mdseries\slshape val}} is closed under Galois conjugacy then the result can be determined without
explicitly computing the \mbox{\texttt{\mdseries\slshape p}}\texttt{\symbol{45}}modular reduction of \mbox{\texttt{\mdseries\slshape val}}. This happens for example if \mbox{\texttt{\mdseries\slshape val}} is a Brauer character. 

 If \mbox{\texttt{\mdseries\slshape val}} is an irreducible Brauer character then the value returned is the size of the
smallest finite field in characteristic \mbox{\texttt{\mdseries\slshape p}} over which the corresponding representation lives. }

 

\subsection{\textcolor{Chapter }{RealizableBrauerCharacters}}
\logpage{[ 72, 15, 4 ]}\nobreak
\hyperdef{L}{X782400277F6316A4}{}
{\noindent\textcolor{FuncColor}{$\triangleright$\enspace\texttt{RealizableBrauerCharacters({\mdseries\slshape matrix, q})\index{RealizableBrauerCharacters@\texttt{RealizableBrauerCharacters}}
\label{RealizableBrauerCharacters}
}\hfill{\scriptsize (function)}}\\


 For a list \mbox{\texttt{\mdseries\slshape matrix}} of absolutely irreducible Brauer characters in characteristic $p$, and a power \mbox{\texttt{\mdseries\slshape q}} of $p$, \texttt{RealizableBrauerCharacters} returns a duplicate\texttt{\symbol{45}}free list of sums of Frobenius
conjugates of the rows of \mbox{\texttt{\mdseries\slshape matrix}}, each irreducible over the field with \mbox{\texttt{\mdseries\slshape q}} elements. 

 
\begin{Verbatim}[commandchars=!@|,fontsize=\small,frame=single,label=Example]
  !gapprompt@gap>| !gapinput@irr:= Irr( CharacterTable( "A5" ) mod 2 );|
  [ Character( BrauerTable( "A5", 2 ), [ 1, 1, 1, 1 ] ),
    Character( BrauerTable( "A5", 2 ),
      [ 2, -1, E(5)+E(5)^4, E(5)^2+E(5)^3 ] ),
    Character( BrauerTable( "A5", 2 ),
      [ 2, -1, E(5)^2+E(5)^3, E(5)+E(5)^4 ] ),
    Character( BrauerTable( "A5", 2 ), [ 4, 1, -1, -1 ] ) ]
  !gapprompt@gap>| !gapinput@List( irr, phi -> SizeOfFieldOfDefinition( phi, 2 ) );|
  [ 2, 4, 4, 2 ]
  !gapprompt@gap>| !gapinput@RealizableBrauerCharacters( irr, 2 );|
  [ Character( BrauerTable( "A5", 2 ), [ 1, 1, 1, 1 ] ),
    ClassFunction( BrauerTable( "A5", 2 ), [ 4, -2, -1, -1 ] ),
    Character( BrauerTable( "A5", 2 ), [ 4, 1, -1, -1 ] ) ]
\end{Verbatim}
 }

 }

  
\section{\textcolor{Chapter }{Domains Generated by Class Functions}}\label{Domains Generated by Class Functions}
\logpage{[ 72, 16, 0 ]}
\hyperdef{L}{X7FEEDC0981A22850}{}
{
  \textsf{GAP} supports groups, vector spaces, and algebras generated by class functions.  }

 }

            
\chapter{\textcolor{Chapter }{Maps Concerning Character Tables}}\label{Maps Concerning Character Tables}
\logpage{[ 73, 0, 0 ]}
\hyperdef{L}{X7DF1ACDE7E9C6294}{}
{
  \index{maps} \index{parametrized maps} Besides the characters, \emph{power maps} are an important part of a character table, see Section{\nobreakspace}\ref{Power Maps}. Often their computation is not easy, and if the table has no access to the
underlying group then in general they cannot be obtained from the matrix of
irreducible characters; so it is useful to store them on the table. 

 If not only a single table is considered but different tables of a group and a
subgroup or of a group and a factor group are used, also \emph{class fusion maps} (see Section{\nobreakspace}\ref{Class Fusions between Character Tables}) must be known to get information about the embedding or simply to induce or
restrict characters, see Section{\nobreakspace}\ref{Restricted and Induced Class Functions}). 

 These are examples of functions from conjugacy classes which will be called \emph{maps} in the following. (This should not be confused with the term mapping, cf.
Chapter{\nobreakspace}\ref{Mappings}.) In \textsf{GAP}, maps are represented by lists. Also each character, each list of element
orders, of centralizer orders, or of class lengths are maps, and the list
returned by \texttt{ListPerm} (\ref{ListPerm}), when this function is called with a permutation of classes, is a map. 

 When maps are constructed without access to a group, often one only knows that
the image of a given class is contained in a set of possible images, e. g.,
that the image of a class under a subgroup fusion is in the set of all classes
with the same element order. Using further information, such as centralizer
orders, power maps and the restriction of characters, the sets of possible
images can be restricted further. In many cases, at the end the images are
uniquely determined. 

 Because of this approach, many functions in this chapter work not only with
maps but with \emph{parametrized maps} (or \emph{paramaps} for short). More about parametrized maps can be found in
Section{\nobreakspace}\ref{Parametrized Maps}. 

 The implementation follows{\nobreakspace}\cite{Bre91}, a description of the main ideas together with several examples can be found
in{\nobreakspace}\cite{Bre99}. 

 Several examples in this chapter require the \textsf{GAP} Character Table Library to be available. If it is not yet loaded then we load
it now. 

 
\begin{Verbatim}[commandchars=!@|,fontsize=\small,frame=single,label=Example]
  !gapprompt@gap>| !gapinput@LoadPackage( "ctbllib" );|
  true
\end{Verbatim}
  
\section{\textcolor{Chapter }{Power Maps}}\label{Power Maps}
\logpage{[ 73, 1, 0 ]}
\hyperdef{L}{X7FED949A86575949}{}
{
  The $n$\texttt{\symbol{45}}th power map of a character table is represented by a list
that stores at position $i$ the position of the class containing the $n$\texttt{\symbol{45}}th powers of the elements in the $i$\texttt{\symbol{45}}th class. The $n$\texttt{\symbol{45}}th power map can be composed from the power maps of the
prime divisors of $n$, so usually only power maps for primes are actually stored in the character
table. 

 For an ordinary character table \mbox{\texttt{\mdseries\slshape tbl}} with access to its underlying group $G$, the $p$\texttt{\symbol{45}}th power map of \mbox{\texttt{\mdseries\slshape tbl}} can be computed using the identification of the conjugacy classes of $G$ with the classes of \mbox{\texttt{\mdseries\slshape tbl}}. For an ordinary character table without access to a group, in general the $p$\texttt{\symbol{45}}th power maps (and hence also the element orders) for
prime divisors $p$ of the group order are not uniquely determined by the matrix of irreducible
characters. So only necessary conditions can be checked in this case, which in
general yields only a list of several possibilities for the desired power map.
Character tables of the \textsf{GAP} character table library store all $p$\texttt{\symbol{45}}th power maps for prime divisors $p$ of the group order. 

 Power maps of Brauer tables can be derived from the power maps of the
underlying ordinary tables. 

 For (computing and) accessing the $n$\texttt{\symbol{45}}th power map of a character table, \texttt{PowerMap} (\ref{PowerMap}) can be used; if the $n$\texttt{\symbol{45}}th power map cannot be uniquely determined then \texttt{PowerMap} (\ref{PowerMap}) returns \texttt{fail}. 

 The list of all possible $p$\texttt{\symbol{45}}th power maps of a table in the sense that certain
necessary conditions are satisfied can be computed with \texttt{PossiblePowerMaps} (\ref{PossiblePowerMaps}). This provides a default strategy, the subroutines are listed in
Section{\nobreakspace}\ref{Subroutines for the Construction of Power Maps}. 

\subsection{\textcolor{Chapter }{PowerMap}}
\logpage{[ 73, 1, 1 ]}\nobreak
\hyperdef{L}{X781FAA497E3B4D1A}{}
{\noindent\textcolor{FuncColor}{$\triangleright$\enspace\texttt{PowerMap({\mdseries\slshape tbl, n[, class]})\index{PowerMap@\texttt{PowerMap}}
\label{PowerMap}
}\hfill{\scriptsize (operation)}}\\
\noindent\textcolor{FuncColor}{$\triangleright$\enspace\texttt{PowerMapOp({\mdseries\slshape tbl, n[, class]})\index{PowerMapOp@\texttt{PowerMapOp}}
\label{PowerMapOp}
}\hfill{\scriptsize (operation)}}\\
\noindent\textcolor{FuncColor}{$\triangleright$\enspace\texttt{ComputedPowerMaps({\mdseries\slshape tbl})\index{ComputedPowerMaps@\texttt{ComputedPowerMaps}}
\label{ComputedPowerMaps}
}\hfill{\scriptsize (attribute)}}\\


 Called with first argument a character table \mbox{\texttt{\mdseries\slshape tbl}} and second argument an integer \mbox{\texttt{\mdseries\slshape n}}, \texttt{PowerMap} returns the \mbox{\texttt{\mdseries\slshape n}}\texttt{\symbol{45}}th power map of \mbox{\texttt{\mdseries\slshape tbl}}. This is a list containing at position $i$ the position of the class of \mbox{\texttt{\mdseries\slshape n}}\texttt{\symbol{45}}th powers of the elements in the $i$\texttt{\symbol{45}}th class of \mbox{\texttt{\mdseries\slshape tbl}}. 

 If the additional third argument \mbox{\texttt{\mdseries\slshape class}} is present then the position of \mbox{\texttt{\mdseries\slshape n}}\texttt{\symbol{45}}th powers of the \mbox{\texttt{\mdseries\slshape class}}\texttt{\symbol{45}}th class is returned. 

 If the \mbox{\texttt{\mdseries\slshape n}}\texttt{\symbol{45}}th power map is not uniquely determined by \mbox{\texttt{\mdseries\slshape tbl}} then \texttt{fail} is returned. This can happen only if \mbox{\texttt{\mdseries\slshape tbl}} has no access to its underlying group. 

 The power maps of \mbox{\texttt{\mdseries\slshape tbl}} that were computed already by \texttt{PowerMap} are stored in \mbox{\texttt{\mdseries\slshape tbl}} as value of the attribute \texttt{ComputedPowerMaps}, the $n$\texttt{\symbol{45}}th power map at position $n$. \texttt{PowerMap} checks whether the desired power map is already stored, computes it using the
operation \texttt{PowerMapOp} if it is not yet known, and stores it. So methods for the computation of power
maps can be installed for the operation \texttt{PowerMapOp}.  

 
\begin{Verbatim}[commandchars=!@|,fontsize=\small,frame=single,label=Example]
  !gapprompt@gap>| !gapinput@tbl:= CharacterTable( "L3(2)" );;|
  !gapprompt@gap>| !gapinput@ComputedPowerMaps( tbl );|
  [ , [ 1, 1, 3, 2, 5, 6 ], [ 1, 2, 1, 4, 6, 5 ],,,,
    [ 1, 2, 3, 4, 1, 1 ] ]
  !gapprompt@gap>| !gapinput@PowerMap( tbl, 5 );|
  [ 1, 2, 3, 4, 6, 5 ]
  !gapprompt@gap>| !gapinput@ComputedPowerMaps( tbl );|
  [ , [ 1, 1, 3, 2, 5, 6 ], [ 1, 2, 1, 4, 6, 5 ],, [ 1, 2, 3, 4, 6, 5 ],
    , [ 1, 2, 3, 4, 1, 1 ] ]
  !gapprompt@gap>| !gapinput@PowerMap( tbl, 137, 2 );|
  2
\end{Verbatim}
 }

 

\subsection{\textcolor{Chapter }{PossiblePowerMaps}}
\logpage{[ 73, 1, 2 ]}\nobreak
\hyperdef{L}{X7C7B292E80590BE0}{}
{\noindent\textcolor{FuncColor}{$\triangleright$\enspace\texttt{PossiblePowerMaps({\mdseries\slshape tbl, p[, options]})\index{PossiblePowerMaps@\texttt{PossiblePowerMaps}}
\label{PossiblePowerMaps}
}\hfill{\scriptsize (operation)}}\\


 For the ordinary character table \mbox{\texttt{\mdseries\slshape tbl}} of a group $G$ and a prime integer \mbox{\texttt{\mdseries\slshape p}}, \texttt{PossiblePowerMaps} returns the list of all maps that have the following properties of the $p$\texttt{\symbol{45}}th power map of \mbox{\texttt{\mdseries\slshape tbl}}. (Representative orders are used only if the \texttt{OrdersClassRepresentatives} (\ref{OrdersClassRepresentatives}) value of \mbox{\texttt{\mdseries\slshape tbl}} is known. 
\begin{enumerate}
\item  For class $i$, the centralizer order of the image is a multiple of the $i$\texttt{\symbol{45}}th centralizer order; if the elements in the $i$\texttt{\symbol{45}}th class have order coprime to $p$ then the centralizer orders of class $i$ and its image are equal. 
\item  Let $n$ be the order of elements in class $i$. If \mbox{\texttt{\mdseries\slshape prime}} divides $n$ then the images have order $n/p$; otherwise the images have order $n$. These criteria are checked in \texttt{InitPowerMap} (\ref{InitPowerMap}). 
\item  For each character $\chi$ of $G$ and each element $g$ in $G$, the values $\chi(g^p)$ and \texttt{GaloisCyc}$( \chi(g), p )$ are algebraic integers that are congruent modulo $p$; if $p$ does not divide the element order of $g$ then the two values are equal. This congruence is checked for the characters
specified below in the discussion of the \mbox{\texttt{\mdseries\slshape options}} argument; For linear characters $\lambda$ among these characters, the condition $\chi(g)^p = \chi(g^p)$ is checked. The corresponding function is \texttt{Congruences} (\ref{Congruences:for character tables}). 
\item  For each character $\chi$ of $G$, the kernel is a normal subgroup $N$, and $g^p \in N$ for all $g \in N$; moreover, if $N$ has index $p$ in $G$ then $g^p \in N$ for all $g \in G$, and if the index of $N$ in $G$ is coprime to $p$ then $g^p \not \in N$ for each $g \not \in N$. These conditions are checked for the kernels of all characters $\chi$ specified below, the corresponding function is \texttt{ConsiderKernels} (\ref{ConsiderKernels}). 
\item  If $p$ is larger than the order $m$ of an element $g \in G$ then the class of $g^p$ is determined by the power maps for primes dividing the residue of $p$ modulo $m$. If these power maps are stored in the \texttt{ComputedPowerMaps} (\ref{ComputedPowerMaps}) value of \mbox{\texttt{\mdseries\slshape tbl}} then this information is used. This criterion is checked in \texttt{ConsiderSmallerPowerMaps} (\ref{ConsiderSmallerPowerMaps}). 
\item  For each character $\chi$ of $G$, the symmetrization $\psi$ defined by $\psi(g) = (\chi(g)^p - \chi(g^p))/p$ is a character. This condition is checked for the kernels of all characters $\chi$ specified below, the corresponding function is \texttt{PowerMapsAllowedBySymmetrizations} (\ref{PowerMapsAllowedBySymmetrizations}). 
\end{enumerate}
 

 If \mbox{\texttt{\mdseries\slshape tbl}} is a Brauer table, the possibilities are computed from those for the
underlying ordinary table. 

 The optional argument \mbox{\texttt{\mdseries\slshape options}}, if given, must be a record that may have the following components: 
\begin{description}
\item[{\texttt{chars}:}]  a list of characters which are used for the check of the criteria 3., 4., and
6.; the default is \texttt{Irr( \mbox{\texttt{\mdseries\slshape tbl}} )}, 
\item[{\texttt{powermap}:}]  a parametrized map which is an approximation of the desired map 
\item[{\texttt{decompose}:}]  a Boolean; a \texttt{true} value indicates that all constituents of the symmetrizations of \texttt{chars} computed for criterion 6. lie in \texttt{chars}, so the symmetrizations can be decomposed into elements of \texttt{chars}; the default value of \texttt{decompose} is \texttt{true} if \texttt{chars} is not bound and \texttt{Irr( \mbox{\texttt{\mdseries\slshape tbl}} )} is known, otherwise \texttt{false}, 
\item[{\texttt{quick}:}]  a Boolean; if \texttt{true} then the subroutines are called with value \texttt{true} for the argument \mbox{\texttt{\mdseries\slshape quick}}; especially, as soon as only one candidate remains this candidate is returned
immediately; the default value is \texttt{false}, 
\item[{\texttt{parameters}:}]  a record with components \texttt{maxamb}, \texttt{minamb} and \texttt{maxlen} which control the subroutine \texttt{PowerMapsAllowedBySymmetrizations} (\ref{PowerMapsAllowedBySymmetrizations}); it only uses characters with current indeterminateness up to \texttt{maxamb}, tests decomposability only for characters with current indeterminateness at
least \texttt{minamb}, and admits a branch according to a character only if there is one with at
most \texttt{maxlen} possible symmetrizations. 
\end{description}
 

 
\begin{Verbatim}[commandchars=!@|,fontsize=\small,frame=single,label=Example]
  !gapprompt@gap>| !gapinput@tbl:= CharacterTable( "U4(3).4" );;|
  !gapprompt@gap>| !gapinput@PossiblePowerMaps( tbl, 2 );|
  [ [ 1, 1, 3, 4, 5, 2, 2, 8, 3, 4, 11, 12, 6, 14, 9, 1, 1, 2, 2, 3, 4,
        5, 6, 8, 9, 9, 10, 11, 12, 16, 16, 16, 16, 17, 17, 18, 18, 18,
        18, 20, 20, 20, 20, 22, 22, 24, 24, 25, 26, 28, 28, 29, 29 ] ]
\end{Verbatim}
 }

 

\subsection{\textcolor{Chapter }{ElementOrdersPowerMap}}
\logpage{[ 73, 1, 3 ]}\nobreak
\hyperdef{L}{X7E0289957E9D62EE}{}
{\noindent\textcolor{FuncColor}{$\triangleright$\enspace\texttt{ElementOrdersPowerMap({\mdseries\slshape powermap})\index{ElementOrdersPowerMap@\texttt{ElementOrdersPowerMap}}
\label{ElementOrdersPowerMap}
}\hfill{\scriptsize (function)}}\\


 Let \mbox{\texttt{\mdseries\slshape powermap}} be a nonempty list containing at position $p$, if bound, the $p$\texttt{\symbol{45}}th power map of a character table or group. \texttt{ElementOrdersPowerMap} returns a list of the same length as each entry in \mbox{\texttt{\mdseries\slshape powermap}}, with entry at position $i$ equal to the order of elements in class $i$ if this order is uniquely determined by \mbox{\texttt{\mdseries\slshape powermap}}, and equal to an unknown (see Chapter{\nobreakspace}\ref{Unknowns}) otherwise. 

 
\begin{Verbatim}[commandchars=!@|,fontsize=\small,frame=single,label=Example]
  !gapprompt@gap>| !gapinput@tbl:= CharacterTable( "U4(3).4" );;|
  !gapprompt@gap>| !gapinput@known:= ComputedPowerMaps( tbl );;|
  !gapprompt@gap>| !gapinput@Length( known );|
  7
  !gapprompt@gap>| !gapinput@sub:= ShallowCopy( known );;  Unbind( sub[7] );|
  !gapprompt@gap>| !gapinput@ElementOrdersPowerMap( sub );|
  [ 1, 2, 3, 3, 3, 4, 4, 5, 6, 6, Unknown(1), Unknown(2), 8, 9, 12, 2,
    2, 4, 4, 6, 6, 6, 8, 10, 12, 12, 12, Unknown(3), Unknown(4), 4, 4,
    4, 4, 4, 4, 8, 8, 8, 8, 12, 12, 12, 12, 12, 12, 20, 20, 24, 24,
    Unknown(5), Unknown(6), Unknown(7), Unknown(8) ]
  !gapprompt@gap>| !gapinput@ord:= ElementOrdersPowerMap( known );|
  [ 1, 2, 3, 3, 3, 4, 4, 5, 6, 6, 7, 7, 8, 9, 12, 2, 2, 4, 4, 6, 6, 6,
    8, 10, 12, 12, 12, 14, 14, 4, 4, 4, 4, 4, 4, 8, 8, 8, 8, 12, 12,
    12, 12, 12, 12, 20, 20, 24, 24, 28, 28, 28, 28 ]
  !gapprompt@gap>| !gapinput@ord = OrdersClassRepresentatives( tbl );|
  true
\end{Verbatim}
 }

 

\subsection{\textcolor{Chapter }{PowerMapByComposition}}
\logpage{[ 73, 1, 4 ]}\nobreak
\hyperdef{L}{X7C0F171F7DC846B7}{}
{\noindent\textcolor{FuncColor}{$\triangleright$\enspace\texttt{PowerMapByComposition({\mdseries\slshape tbl, n})\index{PowerMapByComposition@\texttt{PowerMapByComposition}}
\label{PowerMapByComposition}
}\hfill{\scriptsize (function)}}\\


 \mbox{\texttt{\mdseries\slshape tbl}} must be a nearly character table, and \mbox{\texttt{\mdseries\slshape n}} a positive integer. If the power maps for all prime divisors of \mbox{\texttt{\mdseries\slshape n}} are stored in the \texttt{ComputedPowerMaps} (\ref{ComputedPowerMaps}) list of \mbox{\texttt{\mdseries\slshape tbl}} then \texttt{PowerMapByComposition} returns the \mbox{\texttt{\mdseries\slshape n}}\texttt{\symbol{45}}th power map of \mbox{\texttt{\mdseries\slshape tbl}}. Otherwise \texttt{fail} is returned. 

 
\begin{Verbatim}[commandchars=!@|,fontsize=\small,frame=single,label=Example]
  !gapprompt@gap>| !gapinput@tbl:= CharacterTable( "U4(3).4" );;  exp:= Exponent( tbl );|
  2520
  !gapprompt@gap>| !gapinput@PowerMapByComposition( tbl, exp );|
  [ 1, 1, 1, 1, 1, 1, 1, 1, 1, 1, 1, 1, 1, 1, 1, 1, 1, 1, 1, 1, 1, 1,
    1, 1, 1, 1, 1, 1, 1, 1, 1, 1, 1, 1, 1, 1, 1, 1, 1, 1, 1, 1, 1, 1,
    1, 1, 1, 1, 1, 1, 1, 1, 1 ]
  !gapprompt@gap>| !gapinput@Length( ComputedPowerMaps( tbl ) );|
  7
  !gapprompt@gap>| !gapinput@PowerMapByComposition( tbl, 11 );|
  fail
  !gapprompt@gap>| !gapinput@PowerMap( tbl, 11 );;|
  !gapprompt@gap>| !gapinput@PowerMapByComposition( tbl, 11 );|
  [ 1, 2, 3, 4, 5, 6, 7, 8, 9, 10, 11, 12, 13, 14, 15, 16, 17, 18, 19,
    20, 21, 22, 23, 24, 26, 25, 27, 28, 29, 31, 30, 33, 32, 35, 34, 37,
    36, 39, 38, 41, 40, 43, 42, 45, 44, 47, 46, 49, 48, 51, 50, 53, 52 ]
\end{Verbatim}
 }

 }

  
\section{\textcolor{Chapter }{Orbits on Sets of Possible Power Maps}}\label{Orbits on Sets of Possible Power Maps}
\logpage{[ 73, 2, 0 ]}
\hyperdef{L}{X80980FF37F0D521B}{}
{
  The permutation group of matrix automorphisms (see{\nobreakspace}\texttt{MatrixAutomorphisms} (\ref{MatrixAutomorphisms})) acts on the possible power maps returned by \texttt{PossiblePowerMaps} (\ref{PossiblePowerMaps}) by permuting a list via \texttt{Permuted} (\ref{Permuted}) and then mapping the images via \texttt{OnPoints} (\ref{OnPoints}). Note that by definition, the group of \emph{table} automorphisms acts trivially. 

\subsection{\textcolor{Chapter }{OrbitPowerMaps}}
\logpage{[ 73, 2, 1 ]}\nobreak
\hyperdef{L}{X7ECB9DDE8608B9A9}{}
{\noindent\textcolor{FuncColor}{$\triangleright$\enspace\texttt{OrbitPowerMaps({\mdseries\slshape map, permgrp})\index{OrbitPowerMaps@\texttt{OrbitPowerMaps}}
\label{OrbitPowerMaps}
}\hfill{\scriptsize (function)}}\\


 returns the orbit of the power map \mbox{\texttt{\mdseries\slshape map}} under the action of the permutation group \mbox{\texttt{\mdseries\slshape permgrp}} via a combination of \texttt{Permuted} (\ref{Permuted}) and \texttt{OnPoints} (\ref{OnPoints}). }

 

\subsection{\textcolor{Chapter }{RepresentativesPowerMaps}}
\logpage{[ 73, 2, 2 ]}\nobreak
\hyperdef{L}{X8753F5217A570529}{}
{\noindent\textcolor{FuncColor}{$\triangleright$\enspace\texttt{RepresentativesPowerMaps({\mdseries\slshape listofmaps, permgrp})\index{RepresentativesPowerMaps@\texttt{RepresentativesPowerMaps}}
\label{RepresentativesPowerMaps}
}\hfill{\scriptsize (function)}}\\


 \index{matrix automorphisms} returns a list of orbit representatives of the power maps in the list \mbox{\texttt{\mdseries\slshape listofmaps}} under the action of the permutation group \mbox{\texttt{\mdseries\slshape permgrp}} via a combination of \texttt{Permuted} (\ref{Permuted}) and \texttt{OnPoints} (\ref{OnPoints}). 

 
\begin{Verbatim}[commandchars=!@|,fontsize=\small,frame=single,label=Example]
  !gapprompt@gap>| !gapinput@tbl:= CharacterTable( "3.McL" );;|
  !gapprompt@gap>| !gapinput@grp:= MatrixAutomorphisms( Irr( tbl ) );  Size( grp );|
  <permutation group with 5 generators>
  32
  !gapprompt@gap>| !gapinput@poss:= PossiblePowerMaps( CharacterTable( "3.McL" ), 3 );|
  [ [ 1, 1, 1, 4, 4, 4, 1, 1, 1, 1, 11, 11, 11, 14, 14, 14, 17, 17, 17,
        4, 4, 4, 4, 4, 4, 29, 29, 29, 26, 26, 26, 32, 32, 32, 9, 8, 37,
        37, 37, 40, 40, 40, 43, 43, 43, 11, 11, 11, 52, 52, 52, 49, 49,
        49, 14, 14, 14, 14, 14, 14, 37, 37, 37, 37, 37, 37 ],
    [ 1, 1, 1, 4, 4, 4, 1, 1, 1, 1, 11, 11, 11, 14, 14, 14, 17, 17, 17,
        4, 4, 4, 4, 4, 4, 29, 29, 29, 26, 26, 26, 32, 32, 32, 8, 9, 37,
        37, 37, 40, 40, 40, 43, 43, 43, 11, 11, 11, 52, 52, 52, 49, 49,
        49, 14, 14, 14, 14, 14, 14, 37, 37, 37, 37, 37, 37 ] ]
  !gapprompt@gap>| !gapinput@reps:= RepresentativesPowerMaps( poss, grp );|
  [ [ 1, 1, 1, 4, 4, 4, 1, 1, 1, 1, 11, 11, 11, 14, 14, 14, 17, 17, 17,
        4, 4, 4, 4, 4, 4, 29, 29, 29, 26, 26, 26, 32, 32, 32, 8, 9, 37,
        37, 37, 40, 40, 40, 43, 43, 43, 11, 11, 11, 52, 52, 52, 49, 49,
        49, 14, 14, 14, 14, 14, 14, 37, 37, 37, 37, 37, 37 ] ]
  !gapprompt@gap>| !gapinput@orb:= OrbitPowerMaps( reps[1], grp );|
  [ [ 1, 1, 1, 4, 4, 4, 1, 1, 1, 1, 11, 11, 11, 14, 14, 14, 17, 17, 17,
        4, 4, 4, 4, 4, 4, 29, 29, 29, 26, 26, 26, 32, 32, 32, 8, 9, 37,
        37, 37, 40, 40, 40, 43, 43, 43, 11, 11, 11, 52, 52, 52, 49, 49,
        49, 14, 14, 14, 14, 14, 14, 37, 37, 37, 37, 37, 37 ],
    [ 1, 1, 1, 4, 4, 4, 1, 1, 1, 1, 11, 11, 11, 14, 14, 14, 17, 17, 17,
        4, 4, 4, 4, 4, 4, 29, 29, 29, 26, 26, 26, 32, 32, 32, 9, 8, 37,
        37, 37, 40, 40, 40, 43, 43, 43, 11, 11, 11, 52, 52, 52, 49, 49,
        49, 14, 14, 14, 14, 14, 14, 37, 37, 37, 37, 37, 37 ] ]
  !gapprompt@gap>| !gapinput@Parametrized( orb );|
  [ 1, 1, 1, 4, 4, 4, 1, 1, 1, 1, 11, 11, 11, 14, 14, 14, 17, 17, 17,
    4, 4, 4, 4, 4, 4, 29, 29, 29, 26, 26, 26, 32, 32, 32, [ 8, 9 ],
    [ 8, 9 ], 37, 37, 37, 40, 40, 40, 43, 43, 43, 11, 11, 11, 52, 52,
    52, 49, 49, 49, 14, 14, 14, 14, 14, 14, 37, 37, 37, 37, 37, 37 ]
\end{Verbatim}
 }

 }

  
\section{\textcolor{Chapter }{Class Fusions between Character Tables}}\label{Class Fusions between Character Tables}
\logpage{[ 73, 3, 0 ]}
\hyperdef{L}{X806975FE81534444}{}
{
  \index{fusions}\index{subgroup fusions} For a group $G$ and a subgroup $H$ of $G$, the fusion map between the character table of $H$ and the character table of $G$ is represented by a list that stores at position $i$ the position of the $i$\texttt{\symbol{45}}th class of the table of $H$ in the classes list of the table of $G$. 

 For ordinary character tables \mbox{\texttt{\mdseries\slshape tbl1}} and \mbox{\texttt{\mdseries\slshape tbl2}} of $H$ and $G$, with access to the groups $H$ and $G$, the class fusion between \mbox{\texttt{\mdseries\slshape tbl1}} and \mbox{\texttt{\mdseries\slshape tbl2}} can be computed using the identifications of the conjugacy classes of $H$ with the classes of \mbox{\texttt{\mdseries\slshape tbl1}} and the conjugacy classes of $G$ with the classes of \mbox{\texttt{\mdseries\slshape tbl2}}. For two ordinary character tables without access to an underlying group, or
in the situation that the group stored in \mbox{\texttt{\mdseries\slshape tbl1}} is not physically a subgroup of the group stored in \mbox{\texttt{\mdseries\slshape tbl2}} but an isomorphic copy, in general the class fusion is not uniquely determined
by the information stored on the tables such as irreducible characters and
power maps. So only necessary conditions can be checked in this case, which in
general yields only a list of several possibilities for the desired class
fusion. Character tables of the \textsf{GAP} character table library store various class fusions that are regarded as
important, for example fusions from maximal subgroups (see{\nobreakspace}\texttt{ComputedClassFusions} (\ref{ComputedClassFusions}) and \texttt{Maxes} (\textbf{CTblLib: Maxes}) in the manual for the \textsf{GAP} Character Table Library). 

 Class fusions between Brauer tables can be derived from the class fusions
between the underlying ordinary tables. The class fusion from a Brauer table
to the underlying ordinary table is stored when the Brauer table is
constructed from the ordinary table, so no method is needed to compute such a
fusion. 

 For (computing and) accessing the class fusion between two character tables, \texttt{FusionConjugacyClasses} (\ref{FusionConjugacyClasses:for two character tables}) can be used; if the class fusion cannot be uniquely determined then \texttt{FusionConjugacyClasses} (\ref{FusionConjugacyClasses:for two character tables}) returns \texttt{fail}. 

 The list of all possible class fusion between two tables in the sense that
certain necessary conditions are satisfied can be computed with \texttt{PossibleClassFusions} (\ref{PossibleClassFusions}). This provides a default strategy, the subroutines are listed in Section \ref{Subroutines for the Construction of Class Fusions}. 

 It should be noted that all the following functions except \texttt{FusionConjugacyClasses} (\ref{FusionConjugacyClasses:for two character tables}) deal only with the situation of class fusions from subgroups. The computation
of \emph{factor fusions} from a character table to the table of a factor group is not dealt with here.
Since the ordinary character table of a group $G$ determines the character tables of all factor groups of $G$, the factor fusion to a given character table of a factor group of $G$ is determined up to table automorphisms (see{\nobreakspace}\texttt{AutomorphismsOfTable} (\ref{AutomorphismsOfTable})) once the class positions of the kernel of the natural epimorphism have been
fixed. 
\subsection{\textcolor{Chapter }{FusionConjugacyClasses}}\logpage{[ 73, 3, 1 ]}
\hyperdef{L}{X86CE53B681F13C63}{}
{
\noindent\textcolor{FuncColor}{$\triangleright$\enspace\texttt{FusionConjugacyClasses({\mdseries\slshape tbl1, tbl2})\index{FusionConjugacyClasses@\texttt{FusionConjugacyClasses}!for two character tables}
\label{FusionConjugacyClasses:for two character tables}
}\hfill{\scriptsize (operation)}}\\
\noindent\textcolor{FuncColor}{$\triangleright$\enspace\texttt{FusionConjugacyClasses({\mdseries\slshape H, G})\index{FusionConjugacyClasses@\texttt{FusionConjugacyClasses}!for two groups}
\label{FusionConjugacyClasses:for two groups}
}\hfill{\scriptsize (operation)}}\\
\noindent\textcolor{FuncColor}{$\triangleright$\enspace\texttt{FusionConjugacyClasses({\mdseries\slshape hom[, tbl1, tbl2]})\index{FusionConjugacyClasses@\texttt{FusionConjugacyClasses}!for a homomorphism}
\label{FusionConjugacyClasses:for a homomorphism}
}\hfill{\scriptsize (operation)}}\\
\noindent\textcolor{FuncColor}{$\triangleright$\enspace\texttt{FusionConjugacyClassesOp({\mdseries\slshape tbl1, tbl2})\index{FusionConjugacyClassesOp@\texttt{FusionConjugacyClassesOp}!for two character tables}
\label{FusionConjugacyClassesOp:for two character tables}
}\hfill{\scriptsize (operation)}}\\
\noindent\textcolor{FuncColor}{$\triangleright$\enspace\texttt{FusionConjugacyClassesOp({\mdseries\slshape hom})\index{FusionConjugacyClassesOp@\texttt{FusionConjugacyClassesOp}!for a homomorphism}
\label{FusionConjugacyClassesOp:for a homomorphism}
}\hfill{\scriptsize (attribute)}}\\


 Called with two character tables \mbox{\texttt{\mdseries\slshape tbl1}} and \mbox{\texttt{\mdseries\slshape tbl2}}, \texttt{FusionConjugacyClasses} returns the fusion of conjugacy classes between \mbox{\texttt{\mdseries\slshape tbl1}} and \mbox{\texttt{\mdseries\slshape tbl2}}. (If one of the tables is a Brauer table, it will delegate this task to the
underlying ordinary table.) 

 Called with two groups \mbox{\texttt{\mdseries\slshape H}} and \mbox{\texttt{\mdseries\slshape G}} where \mbox{\texttt{\mdseries\slshape H}} is a subgroup of \mbox{\texttt{\mdseries\slshape G}}, \texttt{FusionConjugacyClasses} returns the fusion of conjugacy classes between \mbox{\texttt{\mdseries\slshape H}} and \mbox{\texttt{\mdseries\slshape G}}. This is done by delegating to the ordinary character tables of \mbox{\texttt{\mdseries\slshape H}} and \mbox{\texttt{\mdseries\slshape G}}, since class fusions are stored only for character tables and not for groups. 

 Note that the returned class fusion refers to the ordering of conjugacy
classes in the character tables if the arguments are character tables and to
the ordering of conjugacy classes in the groups if the arguments are groups
(see{\nobreakspace}\texttt{ConjugacyClasses} (\ref{ConjugacyClasses:for character tables})). 

 Called with a group homomorphism \mbox{\texttt{\mdseries\slshape hom}}, \texttt{FusionConjugacyClasses} returns the fusion of conjugacy classes between the preimage and the image of \mbox{\texttt{\mdseries\slshape hom}}; contrary to the two cases above, also factor fusions can be handled by this
variant. If \mbox{\texttt{\mdseries\slshape hom}} is the only argument then the class fusion refers to the ordering of conjugacy
classes in the groups. If the character tables of preimage and image are given
as \mbox{\texttt{\mdseries\slshape tbl1}} and \mbox{\texttt{\mdseries\slshape tbl2}}, respectively (each table with its group stored), then the fusion refers to
the ordering of classes in these tables. 

 If no class fusion exists or if the class fusion is not uniquely determined, \texttt{fail} is returned; this may happen when \texttt{FusionConjugacyClasses} is called with two character tables that do not know compatible underlying
groups. 

 Methods for the computation of class fusions can be installed for the
operation \texttt{FusionConjugacyClassesOp}. 

 
\begin{Verbatim}[commandchars=!@|,fontsize=\small,frame=single,label=Example]
  !gapprompt@gap>| !gapinput@s4:= SymmetricGroup( 4 );|
  Sym( [ 1 .. 4 ] )
  !gapprompt@gap>| !gapinput@tbls4:= CharacterTable( s4 );;|
  !gapprompt@gap>| !gapinput@d8:= SylowSubgroup( s4, 2 );|
  Group([ (1,2), (3,4), (1,3)(2,4) ])
  !gapprompt@gap>| !gapinput@FusionConjugacyClasses( d8, s4 );|
  [ 1, 2, 3, 3, 5 ]
  !gapprompt@gap>| !gapinput@tbls5:= CharacterTable( "S5" );;|
  !gapprompt@gap>| !gapinput@FusionConjugacyClasses( CharacterTable( "A5" ), tbls5 );|
  [ 1, 2, 3, 4, 4 ]
  !gapprompt@gap>| !gapinput@FusionConjugacyClasses(CharacterTable("A5"), CharacterTable("J1"));|
  fail
  !gapprompt@gap>| !gapinput@PossibleClassFusions(CharacterTable("A5"), CharacterTable("J1"));|
  [ [ 1, 2, 3, 4, 5 ], [ 1, 2, 3, 5, 4 ] ]
\end{Verbatim}
 }

 

\subsection{\textcolor{Chapter }{ComputedClassFusions}}
\logpage{[ 73, 3, 2 ]}\nobreak
\hyperdef{L}{X7F71402285B7DE8E}{}
{\noindent\textcolor{FuncColor}{$\triangleright$\enspace\texttt{ComputedClassFusions({\mdseries\slshape tbl})\index{ComputedClassFusions@\texttt{ComputedClassFusions}}
\label{ComputedClassFusions}
}\hfill{\scriptsize (attribute)}}\\


 The class fusions from the character table \mbox{\texttt{\mdseries\slshape tbl}} that have been computed already by \texttt{FusionConjugacyClasses} (\ref{FusionConjugacyClasses:for two character tables}) or explicitly stored by \texttt{StoreFusion} (\ref{StoreFusion}) are stored in the \texttt{ComputedClassFusions} list of \mbox{\texttt{\mdseries\slshape tbl1}}. Each entry of this list is a record with the following components. 
\begin{description}
\item[{\texttt{name}}]  the \texttt{Identifier} (\ref{Identifier:for character tables}) value of the character table to which the fusion maps, 
\item[{\texttt{map}}]  the list of positions of image classes, 
\item[{\texttt{text} (optional)}]  a string giving additional information about the fusion map, for example
whether the map is uniquely determined by the character tables, 
\item[{\texttt{specification} (optional, rarely used)}]  a value that distinguishes different fusions between the same tables. 
\end{description}
 

 Note that stored fusion maps may differ from the maps returned by \texttt{GetFusionMap} (\ref{GetFusionMap}) and the maps entered by \texttt{StoreFusion} (\ref{StoreFusion}) if the table \mbox{\texttt{\mdseries\slshape destination}} has a nonidentity \texttt{ClassPermutation} (\ref{ClassPermutation}) value. So if one fetches a fusion map from a table \mbox{\texttt{\mdseries\slshape tbl1}} to a table \mbox{\texttt{\mdseries\slshape tbl2}} via access to the data in the \texttt{ComputedClassFusions} list of \mbox{\texttt{\mdseries\slshape tbl1}} then the stored value must be composed with the \texttt{ClassPermutation} (\ref{ClassPermutation}) value of \mbox{\texttt{\mdseries\slshape tbl2}} in order to obtain the correct class fusion. (If one handles fusions only via \texttt{GetFusionMap} (\ref{GetFusionMap}) and \texttt{StoreFusion} (\ref{StoreFusion}) then this adjustment is made automatically.) 

 Fusions are identified via the \texttt{Identifier} (\ref{Identifier:for character tables}) value of the destination table and not by this table itself because many
fusions between character tables in the \textsf{GAP} character table library are stored on library tables, and it is not desirable
to load together with a library table also all those character tables that
occur as destinations of fusions from this table. 

 For storing fusions and accessing stored fusions, see also{\nobreakspace}\texttt{GetFusionMap} (\ref{GetFusionMap}), \texttt{StoreFusion} (\ref{StoreFusion}). For accessing the identifiers of tables that store a fusion into a given
character table, see{\nobreakspace}\texttt{NamesOfFusionSources} (\ref{NamesOfFusionSources}). }

 

\subsection{\textcolor{Chapter }{GetFusionMap}}
\logpage{[ 73, 3, 3 ]}\nobreak
\hyperdef{L}{X8464DD23879431D9}{}
{\noindent\textcolor{FuncColor}{$\triangleright$\enspace\texttt{GetFusionMap({\mdseries\slshape source, destination[, specification]})\index{GetFusionMap@\texttt{GetFusionMap}}
\label{GetFusionMap}
}\hfill{\scriptsize (function)}}\\


 For two ordinary character tables \mbox{\texttt{\mdseries\slshape source}} and \mbox{\texttt{\mdseries\slshape destination}}, \texttt{GetFusionMap} checks whether the \texttt{ComputedClassFusions} (\ref{ComputedClassFusions}) list of \mbox{\texttt{\mdseries\slshape source}} contains a record with \texttt{name} component \texttt{Identifier( \mbox{\texttt{\mdseries\slshape destination}} )}, and returns the \texttt{map} component of the first such record. \texttt{GetFusionMap( \mbox{\texttt{\mdseries\slshape source}}, \mbox{\texttt{\mdseries\slshape destination}}, \mbox{\texttt{\mdseries\slshape specification}} )} fetches that fusion map for which the record additionally has the \texttt{specification} component \mbox{\texttt{\mdseries\slshape specification}}. 

 If both \mbox{\texttt{\mdseries\slshape source}} and \mbox{\texttt{\mdseries\slshape destination}} are Brauer tables, first the same is done, and if no fusion map was found then \texttt{GetFusionMap} looks whether a fusion map between the ordinary tables is stored; if so then
the fusion map between \mbox{\texttt{\mdseries\slshape source}} and \mbox{\texttt{\mdseries\slshape destination}} is stored on \mbox{\texttt{\mdseries\slshape source}}, and then returned. 

 If no appropriate fusion is found, \texttt{GetFusionMap} returns \texttt{fail}. For the computation of class fusions, see \texttt{FusionConjugacyClasses} (\ref{FusionConjugacyClasses:for two character tables}). }

 

\subsection{\textcolor{Chapter }{StoreFusion}}
\logpage{[ 73, 3, 4 ]}\nobreak
\hyperdef{L}{X808970FE87C3432F}{}
{\noindent\textcolor{FuncColor}{$\triangleright$\enspace\texttt{StoreFusion({\mdseries\slshape source, fusion, destination})\index{StoreFusion@\texttt{StoreFusion}}
\label{StoreFusion}
}\hfill{\scriptsize (function)}}\\


 For two character tables \mbox{\texttt{\mdseries\slshape source}} and \mbox{\texttt{\mdseries\slshape destination}}, \texttt{StoreFusion} stores the fusion \mbox{\texttt{\mdseries\slshape fusion}} from \mbox{\texttt{\mdseries\slshape source}} to \mbox{\texttt{\mdseries\slshape destination}} in the \texttt{ComputedClassFusions} (\ref{ComputedClassFusions}) list of \mbox{\texttt{\mdseries\slshape source}}, and adds the \texttt{Identifier} (\ref{Identifier:for character tables}) string of \mbox{\texttt{\mdseries\slshape destination}} to the \texttt{NamesOfFusionSources} (\ref{NamesOfFusionSources}) list of \mbox{\texttt{\mdseries\slshape destination}}. 

 \mbox{\texttt{\mdseries\slshape fusion}} can either be a fusion map (that is, the list of positions of the image
classes) or a record as described in{\nobreakspace}\texttt{ComputedClassFusions} (\ref{ComputedClassFusions}). 

 If fusions to \mbox{\texttt{\mdseries\slshape destination}} are already stored on \mbox{\texttt{\mdseries\slshape source}} then another fusion can be stored only if it has a record component \texttt{specification} that distinguishes it from the stored fusions. In the case of such an
ambiguity, \texttt{StoreFusion} raises an error. 

 
\begin{Verbatim}[commandchars=!@|,fontsize=\small,frame=single,label=Example]
  !gapprompt@gap>| !gapinput@tbld8:= CharacterTable( d8 );;|
  !gapprompt@gap>| !gapinput@ComputedClassFusions( tbld8 );|
  [ rec( map := [ 1, 2, 3, 3, 5 ], name := "CT1" ) ]
  !gapprompt@gap>| !gapinput@Identifier( tbls4 );|
  "CT1"
  !gapprompt@gap>| !gapinput@GetFusionMap( tbld8, tbls4 );|
  [ 1, 2, 3, 3, 5 ]
  !gapprompt@gap>| !gapinput@GetFusionMap( tbls4, tbls5 );|
  fail
  !gapprompt@gap>| !gapinput@poss:= PossibleClassFusions( tbls4, tbls5 );|
  [ [ 1, 5, 2, 3, 6 ] ]
  !gapprompt@gap>| !gapinput@StoreFusion( tbls4, poss[1], tbls5 );|
  !gapprompt@gap>| !gapinput@GetFusionMap( tbls4, tbls5 );|
  [ 1, 5, 2, 3, 6 ]
\end{Verbatim}
 }

 

\subsection{\textcolor{Chapter }{NamesOfFusionSources}}
\logpage{[ 73, 3, 5 ]}\nobreak
\hyperdef{L}{X7F6569D5786A9D49}{}
{\noindent\textcolor{FuncColor}{$\triangleright$\enspace\texttt{NamesOfFusionSources({\mdseries\slshape tbl})\index{NamesOfFusionSources@\texttt{NamesOfFusionSources}}
\label{NamesOfFusionSources}
}\hfill{\scriptsize (attribute)}}\\


 For a character table \mbox{\texttt{\mdseries\slshape tbl}}, \texttt{NamesOfFusionSources} returns the list of identifiers of all those character tables that are known
to have fusions to \mbox{\texttt{\mdseries\slshape tbl}} stored. The \texttt{NamesOfFusionSources} value is updated whenever a fusion to \mbox{\texttt{\mdseries\slshape tbl}} is stored using \texttt{StoreFusion} (\ref{StoreFusion}). 

 
\begin{Verbatim}[commandchars=!@|,fontsize=\small,frame=single,label=Example]
  !gapprompt@gap>| !gapinput@NamesOfFusionSources( tbls4 );|
  [ "CT2" ]
  !gapprompt@gap>| !gapinput@Identifier( CharacterTable( d8 ) );|
  "CT2"
\end{Verbatim}
 }

 

\subsection{\textcolor{Chapter }{PossibleClassFusions}}
\logpage{[ 73, 3, 6 ]}\nobreak
\hyperdef{L}{X7883271F7F26356E}{}
{\noindent\textcolor{FuncColor}{$\triangleright$\enspace\texttt{PossibleClassFusions({\mdseries\slshape subtbl, tbl[, options]})\index{PossibleClassFusions@\texttt{PossibleClassFusions}}
\label{PossibleClassFusions}
}\hfill{\scriptsize (operation)}}\\


 For two ordinary character tables \mbox{\texttt{\mdseries\slshape subtbl}} and \mbox{\texttt{\mdseries\slshape tbl}} of the groups $H$ and $G$, \texttt{PossibleClassFusions} returns the list of all maps that have the following properties of class
fusions from \mbox{\texttt{\mdseries\slshape subtbl}} to \mbox{\texttt{\mdseries\slshape tbl}}. 
\begin{enumerate}
\item  For class $i$, the centralizer order of the image in $G$ is a multiple of the $i$\texttt{\symbol{45}}th centralizer order in $H$, and the element orders in the $i$\texttt{\symbol{45}}th class and its image are equal. These criteria are
checked in \texttt{InitFusion} (\ref{InitFusion}). 
\item  The class fusion commutes with power maps. This is checked using \texttt{TestConsistencyMaps} (\ref{TestConsistencyMaps}). 
\item  If the permutation character of $G$ corresponding to the action of $G$ on the cosets of $H$ is specified (see the discussion of the \mbox{\texttt{\mdseries\slshape options}} argument below) then it prescribes for each class $C$ of $G$ the number of elements of $H$ fusing into $C$. The corresponding function is \texttt{CheckPermChar} (\ref{CheckPermChar}). 
\item  The table automorphisms of \mbox{\texttt{\mdseries\slshape tbl}} (see{\nobreakspace}\texttt{AutomorphismsOfTable} (\ref{AutomorphismsOfTable})) are used in order to compute only orbit representatives. (But note that the
list returned by \texttt{PossibleClassFusions} contains the full orbits.) 
\item  For each character $\chi$ of $G$, the restriction to $H$ via the class fusion is a character of $H$. This condition is checked for all characters specified below, the
corresponding function is \texttt{FusionsAllowedByRestrictions} (\ref{FusionsAllowedByRestrictions}). 
\item  The class multiplication coefficients in \mbox{\texttt{\mdseries\slshape subtbl}} do not exceed the corresponding coefficients in \mbox{\texttt{\mdseries\slshape tbl}}. This is checked in \texttt{ConsiderStructureConstants} (\ref{ConsiderStructureConstants}), see also the comment on the parameter \texttt{verify} below. 
\end{enumerate}
 

 If \mbox{\texttt{\mdseries\slshape subtbl}} and \mbox{\texttt{\mdseries\slshape tbl}} are Brauer tables then the possibilities are computed from those for the
underlying ordinary tables. 

 The optional argument \mbox{\texttt{\mdseries\slshape options}} must be a record that may have the following components: 
\begin{description}
\item[{\texttt{chars}}]  a list of characters of \mbox{\texttt{\mdseries\slshape tbl}} which are used for the check of{\nobreakspace}5.; the default is \texttt{Irr( \mbox{\texttt{\mdseries\slshape tbl}} )}, 
\item[{\texttt{subchars}}]  a list of characters of \mbox{\texttt{\mdseries\slshape subtbl}} which are constituents of the restrictions of \texttt{chars}, the default is \texttt{Irr( \mbox{\texttt{\mdseries\slshape subtbl}} )}, 
\item[{\texttt{fusionmap}}]  a parametrized map which is an approximation of the desired map, 
\item[{\texttt{decompose}}]  a Boolean; a \texttt{true} value indicates that all constituents of the restrictions of \texttt{chars} computed for criterion 5. lie in \texttt{subchars}, so the restrictions can be decomposed into elements of \texttt{subchars}; the default value of \texttt{decompose} is \texttt{true} if \texttt{subchars} is not bound and \texttt{Irr( \mbox{\texttt{\mdseries\slshape subtbl}} )} is known, otherwise \texttt{false}, 
\item[{\texttt{permchar}}]  (a values list of) a permutation character; only those fusions affording that
permutation character are computed, 
\item[{\texttt{quick}}]  a Boolean; if \texttt{true} then the subroutines are called with value \texttt{true} for the argument \mbox{\texttt{\mdseries\slshape quick}}; especially, as soon as only one possibility remains then this possibility is
returned immediately; the default value is \texttt{false} (note that in situations where the group of \mbox{\texttt{\mdseries\slshape tbl}} has no subgroups with character table \mbox{\texttt{\mdseries\slshape subtbl}}, it may happen that setting \texttt{quick} to \texttt{true} causes \texttt{PossibleClassFusions} to return solutions, whereas the value \texttt{false} yields an empty list), 
\item[{\texttt{verify}}]  a Boolean; if \texttt{false} then \texttt{ConsiderStructureConstants} (\ref{ConsiderStructureConstants}) is called only if more than one orbit of possible class fusions exists, under
the action of the groups of table automorphisms; the default value is \texttt{false} (because the computation of the structure constants is usually very time
consuming, compared with checking the other criteria), 
\item[{\texttt{parameters}}]  a record with components \texttt{maxamb}, \texttt{minamb} and \texttt{maxlen} (and perhaps some optional components) which control the subroutine \texttt{FusionsAllowedByRestrictions} (\ref{FusionsAllowedByRestrictions}); it only uses characters with current indeterminateness up to \texttt{maxamb}, tests decomposability only for characters with current indeterminateness at
least \texttt{minamb}, and admits a branch according to a character only if there is one with at
most \texttt{maxlen} possible restrictions. 
\end{description}
 

 
\begin{Verbatim}[commandchars=!@|,fontsize=\small,frame=single,label=Example]
  !gapprompt@gap>| !gapinput@subtbl:= CharacterTable( "U3(3)" );;  tbl:= CharacterTable( "J4" );;|
  !gapprompt@gap>| !gapinput@PossibleClassFusions( subtbl, tbl );|
  [ [ 1, 2, 4, 4, 5, 5, 6, 10, 12, 13, 14, 14, 21, 21 ],
    [ 1, 2, 4, 4, 5, 5, 6, 10, 13, 12, 14, 14, 21, 21 ],
    [ 1, 2, 4, 4, 6, 6, 6, 10, 12, 13, 15, 15, 22, 22 ],
    [ 1, 2, 4, 4, 6, 6, 6, 10, 12, 13, 16, 16, 22, 22 ],
    [ 1, 2, 4, 4, 6, 6, 6, 10, 13, 12, 15, 15, 22, 22 ],
    [ 1, 2, 4, 4, 6, 6, 6, 10, 13, 12, 16, 16, 22, 22 ] ]
\end{Verbatim}
 }

 

\subsection{\textcolor{Chapter }{ConsiderStructureConstants}}
\logpage{[ 73, 3, 7 ]}\nobreak
\hyperdef{L}{X7BCC5B4B7E9DF42C}{}
{\noindent\textcolor{FuncColor}{$\triangleright$\enspace\texttt{ConsiderStructureConstants({\mdseries\slshape subtbl, tbl, fusions, quick})\index{ConsiderStructureConstants@\texttt{ConsiderStructureConstants}}
\label{ConsiderStructureConstants}
}\hfill{\scriptsize (function)}}\\


 Let \mbox{\texttt{\mdseries\slshape subtbl}} and \mbox{\texttt{\mdseries\slshape tbl}} be ordinary character tables and \mbox{\texttt{\mdseries\slshape fusions}} be a list of possible class fusions from \mbox{\texttt{\mdseries\slshape subtbl}} to \mbox{\texttt{\mdseries\slshape tbl}}. \texttt{ConsiderStructureConstants} returns the list of those maps $\sigma$ in \mbox{\texttt{\mdseries\slshape fusions}} with the property that for all triples $(i,j,k)$ of class positions, \texttt{ClassMultiplicationCoefficient}$( \mbox{\texttt{\mdseries\slshape subtbl}}, i, j, k )$ is not bigger than \texttt{ClassMultiplicationCoefficient}$( \mbox{\texttt{\mdseries\slshape tbl}}, \sigma[i], \sigma[j], \sigma[k] )$; see{\nobreakspace}\texttt{ClassMultiplicationCoefficient} (\ref{ClassMultiplicationCoefficient:for character tables}) for the definition of class multiplication coefficients/structure constants. 

 The argument \mbox{\texttt{\mdseries\slshape quick}} must be a Boolean; if it is \texttt{true} then only those triples are checked for which at least two entries in \mbox{\texttt{\mdseries\slshape fusions}} have different images. }

 }

  
\section{\textcolor{Chapter }{Orbits on Sets of Possible Class Fusions}}\label{Orbits on Sets of Possible Class Fusions}
\logpage{[ 73, 4, 0 ]}
\hyperdef{L}{X7C34060278E4BFC4}{}
{
  The permutation groups of table automorphisms (see{\nobreakspace}\texttt{AutomorphismsOfTable} (\ref{AutomorphismsOfTable})) of the subgroup table \mbox{\texttt{\mdseries\slshape subtbl}} and the supergroup table \mbox{\texttt{\mdseries\slshape tbl}} act on the possible class fusions from \mbox{\texttt{\mdseries\slshape subtbl}} to \mbox{\texttt{\mdseries\slshape tbl}} that are returned by \texttt{PossibleClassFusions} (\ref{PossibleClassFusions}), the former by permuting a list via \texttt{Permuted} (\ref{Permuted}), the latter by mapping the images via \texttt{OnPoints} (\ref{OnPoints}). 

 If a set of possible fusions with certain properties was computed that are not
invariant under the full groups of table automorphisms then only a smaller
group acts on this set. This may happen for example if a permutation character
or if an explicit approximation of the fusion map was prescribed in the call
of \texttt{PossibleClassFusions} (\ref{PossibleClassFusions}). 

\subsection{\textcolor{Chapter }{OrbitFusions}}
\logpage{[ 73, 4, 1 ]}\nobreak
\hyperdef{L}{X79A0FE1C853302D2}{}
{\noindent\textcolor{FuncColor}{$\triangleright$\enspace\texttt{OrbitFusions({\mdseries\slshape subtblautomorphisms, fusionmap, tblautomorphisms})\index{OrbitFusions@\texttt{OrbitFusions}}
\label{OrbitFusions}
}\hfill{\scriptsize (function)}}\\


 returns the orbit of the class fusion map \mbox{\texttt{\mdseries\slshape fusionmap}} under the actions of the permutation groups \mbox{\texttt{\mdseries\slshape subtblautomorphisms}} and \mbox{\texttt{\mdseries\slshape tblautomorphisms}} of automorphisms of the character table of the subgroup and the supergroup,
respectively. }

 

\subsection{\textcolor{Chapter }{RepresentativesFusions}}
\logpage{[ 73, 4, 2 ]}\nobreak
\hyperdef{L}{X821D11D180B5D317}{}
{\noindent\textcolor{FuncColor}{$\triangleright$\enspace\texttt{RepresentativesFusions({\mdseries\slshape subtbl, listofmaps, tbl})\index{RepresentativesFusions@\texttt{RepresentativesFusions}}
\label{RepresentativesFusions}
}\hfill{\scriptsize (function)}}\\


 \index{table automorphisms} Let \mbox{\texttt{\mdseries\slshape listofmaps}} be a list of class fusions from the character table \mbox{\texttt{\mdseries\slshape subtbl}} to the character table \mbox{\texttt{\mdseries\slshape tbl}}. \texttt{RepresentativesFusions} returns a list of orbit representatives of the class fusions under the action
of maximal admissible subgroups of the table automorphism groups of these
character tables. 

 Instead of the character tables \mbox{\texttt{\mdseries\slshape subtbl}} and \mbox{\texttt{\mdseries\slshape tbl}}, also the permutation groups of their table automorphisms (see \texttt{AutomorphismsOfTable} (\ref{AutomorphismsOfTable})) may be entered. 

 
\begin{Verbatim}[commandchars=!@|,fontsize=\small,frame=single,label=Example]
  !gapprompt@gap>| !gapinput@fus:= GetFusionMap( subtbl, tbl );|
  [ 1, 2, 4, 4, 5, 5, 6, 10, 12, 13, 14, 14, 21, 21 ]
  !gapprompt@gap>| !gapinput@orb:= OrbitFusions( AutomorphismsOfTable( subtbl ), fus,|
  !gapprompt@>| !gapinput@             AutomorphismsOfTable( tbl ) );|
  [ [ 1, 2, 4, 4, 5, 5, 6, 10, 12, 13, 14, 14, 21, 21 ],
    [ 1, 2, 4, 4, 5, 5, 6, 10, 13, 12, 14, 14, 21, 21 ] ]
  !gapprompt@gap>| !gapinput@rep:= RepresentativesFusions( subtbl, orb, tbl );|
  [ [ 1, 2, 4, 4, 5, 5, 6, 10, 12, 13, 14, 14, 21, 21 ] ]
\end{Verbatim}
 }

 }

  
\section{\textcolor{Chapter }{Parametrized Maps}}\label{Parametrized Maps}
\logpage{[ 73, 5, 0 ]}
\hyperdef{L}{X7F18772E86F06179}{}
{
  \index{map!parametrized} \index{class functions} A \emph{parametrized map} is a list whose $i$\texttt{\symbol{45}}th entry is either unbound (which means that nothing is
known about the image(s) of the $i$\texttt{\symbol{45}}th class) or the image of the $i$\texttt{\symbol{45}}th class (i.e., an integer for fusion maps, power maps,
element orders etc., and a cyclotomic for characters), or a list of possible
images of the $i$\texttt{\symbol{45}}th class. In this sense, maps are special parametrized
maps. We often identify a parametrized map \mbox{\texttt{\mdseries\slshape paramap}} with the set of all maps \mbox{\texttt{\mdseries\slshape map}} with the property that either \texttt{\mbox{\texttt{\mdseries\slshape map}}[i] = \mbox{\texttt{\mdseries\slshape paramap}}[i]} or \texttt{\mbox{\texttt{\mdseries\slshape map}}[i]} is contained in the list \texttt{\mbox{\texttt{\mdseries\slshape paramap}}[i]}; we say then that \mbox{\texttt{\mdseries\slshape map}} is contained in \mbox{\texttt{\mdseries\slshape paramap}}. 

 This definition implies that parametrized maps cannot be used to describe sets
of maps where lists are possible images. An exception are strings which
naturally arise as images when class names are considered. So strings and
lists of strings are allowed in parametrized maps, and character constants
(see Chapter{\nobreakspace}\ref{Strings and Characters}) are not allowed in maps. 

\subsection{\textcolor{Chapter }{CompositionMaps}}
\logpage{[ 73, 5, 1 ]}\nobreak
\hyperdef{L}{X8740C1397C6A96C8}{}
{\noindent\textcolor{FuncColor}{$\triangleright$\enspace\texttt{CompositionMaps({\mdseries\slshape paramap2, paramap1[, class]})\index{CompositionMaps@\texttt{CompositionMaps}}
\label{CompositionMaps}
}\hfill{\scriptsize (function)}}\\


 The composition of two parametrized maps \mbox{\texttt{\mdseries\slshape paramap1}}, \mbox{\texttt{\mdseries\slshape paramap2}} is defined as the parametrized map \mbox{\texttt{\mdseries\slshape comp}} that contains all compositions $f_2 \circ f_1$ of elements $f_1$ of \mbox{\texttt{\mdseries\slshape paramap1}} and $f_2$ of \mbox{\texttt{\mdseries\slshape paramap2}}. For example, the composition of a character $\chi$ of a group $G$ by a parametrized class fusion map from a subgroup $H$ to $G$ is the parametrized map that contains all restrictions of $\chi$ by elements of the parametrized fusion map. 

 \texttt{CompositionMaps(\mbox{\texttt{\mdseries\slshape paramap2}}, \mbox{\texttt{\mdseries\slshape paramap1}})} is a parametrized map with entry \texttt{CompositionMaps(\mbox{\texttt{\mdseries\slshape paramap2}}, \mbox{\texttt{\mdseries\slshape paramap1}}, \mbox{\texttt{\mdseries\slshape class}})} at position \mbox{\texttt{\mdseries\slshape class}}. If \texttt{\mbox{\texttt{\mdseries\slshape paramap1}}[\mbox{\texttt{\mdseries\slshape class}}]} is an integer then \texttt{CompositionMaps(\mbox{\texttt{\mdseries\slshape paramap2}}, \mbox{\texttt{\mdseries\slshape paramap1}}, \mbox{\texttt{\mdseries\slshape class}})} is equal to \texttt{\mbox{\texttt{\mdseries\slshape paramap2}}[ \mbox{\texttt{\mdseries\slshape paramap1}}[ \mbox{\texttt{\mdseries\slshape class}} ] ]}. Otherwise it is the union of \texttt{\mbox{\texttt{\mdseries\slshape paramap2}}[\mbox{\texttt{\mdseries\slshape i}}]} for \mbox{\texttt{\mdseries\slshape i}} in \texttt{\mbox{\texttt{\mdseries\slshape paramap1}}[ \mbox{\texttt{\mdseries\slshape class}} ]}. 

 
\begin{Verbatim}[commandchars=!@|,fontsize=\small,frame=single,label=Example]
  !gapprompt@gap>| !gapinput@map1:= [ 1, [ 2 .. 4 ], [ 4, 5 ], 1 ];;|
  !gapprompt@gap>| !gapinput@map2:= [ [ 1, 2 ], 2, 2, 3, 3 ];;|
  !gapprompt@gap>| !gapinput@CompositionMaps( map2, map1 );|
  [ [ 1, 2 ], [ 2, 3 ], 3, [ 1, 2 ] ]
  !gapprompt@gap>| !gapinput@CompositionMaps( map1, map2 );|
  [ [ 1, 2, 3, 4 ], [ 2 .. 4 ], [ 2 .. 4 ], [ 4, 5 ], [ 4, 5 ] ]
\end{Verbatim}
 }

 

\subsection{\textcolor{Chapter }{InverseMap}}
\logpage{[ 73, 5, 2 ]}\nobreak
\hyperdef{L}{X7877EE167A711AB6}{}
{\noindent\textcolor{FuncColor}{$\triangleright$\enspace\texttt{InverseMap({\mdseries\slshape paramap})\index{InverseMap@\texttt{InverseMap}}
\label{InverseMap}
}\hfill{\scriptsize (function)}}\\


 For a parametrized map \mbox{\texttt{\mdseries\slshape paramap}}, \texttt{InverseMap} returns a mutable parametrized map whose $i$\texttt{\symbol{45}}th entry is unbound if $i$ is not in the image of \mbox{\texttt{\mdseries\slshape paramap}}, equal to $j$ if $i$ is (in) the image of \texttt{\mbox{\texttt{\mdseries\slshape paramap}}[\mbox{\texttt{\mdseries\slshape j}}]} exactly for $j$, and equal to the set of all preimages of $i$ under \mbox{\texttt{\mdseries\slshape paramap}} otherwise. 

 We have \texttt{CompositionMaps( \mbox{\texttt{\mdseries\slshape paramap}}, InverseMap( \mbox{\texttt{\mdseries\slshape paramap}} ) )} the identity map. 

 
\begin{Verbatim}[commandchars=!@|,fontsize=\small,frame=single,label=Example]
  !gapprompt@gap>| !gapinput@tbl:= CharacterTable( "2.A5" );;  f:= CharacterTable( "A5" );;|
  !gapprompt@gap>| !gapinput@fus:= GetFusionMap( tbl, f );|
  [ 1, 1, 2, 3, 3, 4, 4, 5, 5 ]
  !gapprompt@gap>| !gapinput@inv:= InverseMap( fus );|
  [ [ 1, 2 ], 3, [ 4, 5 ], [ 6, 7 ], [ 8, 9 ] ]
  !gapprompt@gap>| !gapinput@CompositionMaps( fus, inv );|
  [ 1, 2, 3, 4, 5 ]
  !gapprompt@gap>| !gapinput@# transfer a power map ``up'' to the factor group|
  !gapprompt@gap>| !gapinput@pow:= PowerMap( tbl, 2 );|
  [ 1, 1, 2, 4, 4, 8, 8, 6, 6 ]
  !gapprompt@gap>| !gapinput@CompositionMaps( fus, CompositionMaps( pow, inv ) );|
  [ 1, 1, 3, 5, 4 ]
  !gapprompt@gap>| !gapinput@last = PowerMap( f, 2 );|
  true
  !gapprompt@gap>| !gapinput@# transfer a power map of the factor group ``down'' to the group|
  !gapprompt@gap>| !gapinput@CompositionMaps( inv, CompositionMaps( PowerMap( f, 2 ), fus ) );|
  [ [ 1, 2 ], [ 1, 2 ], [ 1, 2 ], [ 4, 5 ], [ 4, 5 ], [ 8, 9 ],
    [ 8, 9 ], [ 6, 7 ], [ 6, 7 ] ]
\end{Verbatim}
 }

 

\subsection{\textcolor{Chapter }{ProjectionMap}}
\logpage{[ 73, 5, 3 ]}\nobreak
\hyperdef{L}{X82C0E76F804C3FF7}{}
{\noindent\textcolor{FuncColor}{$\triangleright$\enspace\texttt{ProjectionMap({\mdseries\slshape fusionmap})\index{ProjectionMap@\texttt{ProjectionMap}}
\label{ProjectionMap}
}\hfill{\scriptsize (function)}}\\


 For a map \mbox{\texttt{\mdseries\slshape fusionmap}}, \texttt{ProjectionMap} returns a parametrized map whose $i$\texttt{\symbol{45}}th entry is unbound if $i$ is not in the image of \mbox{\texttt{\mdseries\slshape fusionmap}}, and equal to $j$ if $j$ is the smallest position such that $i$ is the image of \mbox{\texttt{\mdseries\slshape fusionmap}}\texttt{[}$j$\texttt{]}. 

 We have \texttt{CompositionMaps( \mbox{\texttt{\mdseries\slshape fusionmap}}, ProjectionMap( \mbox{\texttt{\mdseries\slshape fusionmap}} ) )} the identity map, i.e., first projecting and then fusing yields the identity.
Note that \mbox{\texttt{\mdseries\slshape fusionmap}} must \emph{not} be a parametrized map. 

 
\begin{Verbatim}[commandchars=!@|,fontsize=\small,frame=single,label=Example]
  !gapprompt@gap>| !gapinput@ProjectionMap( [ 1, 1, 1, 2, 2, 2, 3, 4, 5, 5, 5, 6, 6, 6 ] );|
  [ 1, 4, 7, 8, 9, 12 ]
\end{Verbatim}
 }

 

\subsection{\textcolor{Chapter }{Indirected}}
\logpage{[ 73, 5, 4 ]}\nobreak
\hyperdef{L}{X7D9CA09385467EDE}{}
{\noindent\textcolor{FuncColor}{$\triangleright$\enspace\texttt{Indirected({\mdseries\slshape character, paramap})\index{Indirected@\texttt{Indirected}}
\label{Indirected}
}\hfill{\scriptsize (function)}}\\


 For a map \mbox{\texttt{\mdseries\slshape character}} and a parametrized map \mbox{\texttt{\mdseries\slshape paramap}}, \texttt{Indirected} returns a parametrized map whose entry at position $i$ is \mbox{\texttt{\mdseries\slshape character}}\texttt{[ }\mbox{\texttt{\mdseries\slshape paramap}}\texttt{[}$i$\texttt{] ]} if \mbox{\texttt{\mdseries\slshape paramap}}\texttt{[}$i$\texttt{]} is an integer, and an unknown (see Chapter{\nobreakspace}\ref{Unknowns}) otherwise. 

 
\begin{Verbatim}[commandchars=!@|,fontsize=\small,frame=single,label=Example]
  !gapprompt@gap>| !gapinput@tbl:= CharacterTable( "M12" );;|
  !gapprompt@gap>| !gapinput@fus:= [ 1, 3, 4, [ 6, 7 ], 8, 10, [ 11, 12 ], [ 11, 12 ],|
  !gapprompt@>| !gapinput@           [ 14, 15 ], [ 14, 15 ] ];;|
  !gapprompt@gap>| !gapinput@List( Irr( tbl ){ [ 1 .. 6 ] }, x -> Indirected( x, fus ) );|
  [ [ 1, 1, 1, 1, 1, 1, 1, 1, 1, 1 ],
    [ 11, 3, 2, Unknown(9), 1, 0, Unknown(10), Unknown(11), 0, 0 ],
    [ 11, 3, 2, Unknown(12), 1, 0, Unknown(13), Unknown(14), 0, 0 ],
    [ 16, 0, -2, 0, 1, 0, 0, 0, Unknown(15), Unknown(16) ],
    [ 16, 0, -2, 0, 1, 0, 0, 0, Unknown(17), Unknown(18) ],
    [ 45, -3, 0, 1, 0, 0, -1, -1, 1, 1 ] ]
\end{Verbatim}
 }

 

\subsection{\textcolor{Chapter }{Parametrized}}
\logpage{[ 73, 5, 5 ]}\nobreak
\hyperdef{L}{X7910BE5687DDAAF3}{}
{\noindent\textcolor{FuncColor}{$\triangleright$\enspace\texttt{Parametrized({\mdseries\slshape list})\index{Parametrized@\texttt{Parametrized}}
\label{Parametrized}
}\hfill{\scriptsize (function)}}\\


 For a list \mbox{\texttt{\mdseries\slshape list}} of (parametrized) maps of the same length, \texttt{Parametrized} returns the smallest parametrized map containing all elements of \mbox{\texttt{\mdseries\slshape list}}. 

 \texttt{Parametrized} is the inverse function to \texttt{ContainedMaps} (\ref{ContainedMaps}). 

 
\begin{Verbatim}[commandchars=!@|,fontsize=\small,frame=single,label=Example]
  !gapprompt@gap>| !gapinput@Parametrized( [ [ 1, 2, 3, 4, 5 ], [ 1, 3, 2, 4, 5 ],|
  !gapprompt@>| !gapinput@                   [ 1, 2, 3, 4, 6 ] ] );|
  [ 1, [ 2, 3 ], [ 2, 3 ], 4, [ 5, 6 ] ]
\end{Verbatim}
 }

 

\subsection{\textcolor{Chapter }{ContainedMaps}}
\logpage{[ 73, 5, 6 ]}\nobreak
\hyperdef{L}{X7917265684700B10}{}
{\noindent\textcolor{FuncColor}{$\triangleright$\enspace\texttt{ContainedMaps({\mdseries\slshape paramap})\index{ContainedMaps@\texttt{ContainedMaps}}
\label{ContainedMaps}
}\hfill{\scriptsize (function)}}\\


 For a parametrized map \mbox{\texttt{\mdseries\slshape paramap}}, \texttt{ContainedMaps} returns the set of all maps contained in \mbox{\texttt{\mdseries\slshape paramap}}. 

 \texttt{ContainedMaps} is the inverse function to \texttt{Parametrized} (\ref{Parametrized}) in the sense that \texttt{Parametrized( ContainedMaps( \mbox{\texttt{\mdseries\slshape paramap}} ) )} is equal to \mbox{\texttt{\mdseries\slshape paramap}}. 

 
\begin{Verbatim}[commandchars=!@|,fontsize=\small,frame=single,label=Example]
  !gapprompt@gap>| !gapinput@ContainedMaps( [ 1, [ 2, 3 ], [ 2, 3 ], 4, [ 5, 6 ] ] );|
  [ [ 1, 2, 2, 4, 5 ], [ 1, 2, 2, 4, 6 ], [ 1, 2, 3, 4, 5 ],
    [ 1, 2, 3, 4, 6 ], [ 1, 3, 2, 4, 5 ], [ 1, 3, 2, 4, 6 ],
    [ 1, 3, 3, 4, 5 ], [ 1, 3, 3, 4, 6 ] ]
\end{Verbatim}
 }

 

\subsection{\textcolor{Chapter }{UpdateMap}}
\logpage{[ 73, 5, 7 ]}\nobreak
\hyperdef{L}{X80C7328C85BFC20B}{}
{\noindent\textcolor{FuncColor}{$\triangleright$\enspace\texttt{UpdateMap({\mdseries\slshape character, paramap, indirected})\index{UpdateMap@\texttt{UpdateMap}}
\label{UpdateMap}
}\hfill{\scriptsize (function)}}\\


 Let \mbox{\texttt{\mdseries\slshape character}} be a map, \mbox{\texttt{\mdseries\slshape paramap}} a parametrized map, and \mbox{\texttt{\mdseries\slshape indirected}} a parametrized map that is contained in \texttt{CompositionMaps( \mbox{\texttt{\mdseries\slshape character}}, \mbox{\texttt{\mdseries\slshape paramap}} )}. 

 Then \texttt{UpdateMap} changes \mbox{\texttt{\mdseries\slshape paramap}} to the parametrized map containing exactly the maps whose composition with \mbox{\texttt{\mdseries\slshape character}} is equal to \mbox{\texttt{\mdseries\slshape indirected}}. 

 If a contradiction is detected then \texttt{false} is returned immediately, otherwise \texttt{true}. 

 
\begin{Verbatim}[commandchars=!@|,fontsize=\small,frame=single,label=Example]
  !gapprompt@gap>| !gapinput@subtbl:= CharacterTable("S4(4).2");; tbl:= CharacterTable("He");;|
  !gapprompt@gap>| !gapinput@fus:= InitFusion( subtbl, tbl );;|
  !gapprompt@gap>| !gapinput@fus;|
  [ 1, 2, 2, [ 2, 3 ], 4, 4, [ 7, 8 ], [ 7, 8 ], 9, 9, 9, [ 10, 11 ],
    [ 10, 11 ], 18, 18, 25, 25, [ 26, 27 ], [ 26, 27 ], 2, [ 6, 7 ],
    [ 6, 7 ], [ 6, 7, 8 ], 10, 10, 17, 17, 18, [ 19, 20 ], [ 19, 20 ] ]
  !gapprompt@gap>| !gapinput@chi:= Irr( tbl )[2];|
  Character( CharacterTable( "He" ), [ 51, 11, 3, 6, 0, 3, 3, -1, 1, 2,
    0, 3*E(7)+3*E(7)^2+3*E(7)^4, 3*E(7)^3+3*E(7)^5+3*E(7)^6, 2,
    E(7)+E(7)^2+2*E(7)^3+E(7)^4+2*E(7)^5+2*E(7)^6,
    2*E(7)+2*E(7)^2+E(7)^3+2*E(7)^4+E(7)^5+E(7)^6, 1, 1, 0, 0,
    -E(7)-E(7)^2-E(7)^4, -E(7)^3-E(7)^5-E(7)^6, E(7)+E(7)^2+E(7)^4,
    E(7)^3+E(7)^5+E(7)^6, 1, 0, 0, -1, -1, 0, 0, E(7)+E(7)^2+E(7)^4,
    E(7)^3+E(7)^5+E(7)^6 ] )
  !gapprompt@gap>| !gapinput@filt:= Filtered( Irr( subtbl ), x -> x[1] = 50 );|
  [ Character( CharacterTable( "S4(4).2" ),
      [ 50, 10, 10, 2, 5, 5, -2, 2, 0, 0, 0, 1, 1, 0, 0, 0, 0, -1, -1,
        10, 2, 2, 2, 1, 1, 0, 0, 0, -1, -1 ] ),
    Character( CharacterTable( "S4(4).2" ),
      [ 50, 10, 10, 2, 5, 5, -2, 2, 0, 0, 0, 1, 1, 0, 0, 0, 0, -1, -1,
        -10, -2, -2, -2, -1, -1, 0, 0, 0, 1, 1 ] ) ]
  !gapprompt@gap>| !gapinput@UpdateMap( chi, fus, filt[1] + TrivialCharacter( subtbl ) );|
  true
  !gapprompt@gap>| !gapinput@fus;|
  [ 1, 2, 2, 3, 4, 4, 8, 7, 9, 9, 9, 10, 10, 18, 18, 25, 25,
    [ 26, 27 ], [ 26, 27 ], 2, [ 6, 7 ], [ 6, 7 ], [ 6, 7 ], 10, 10,
    17, 17, 18, [ 19, 20 ], [ 19, 20 ] ]
\end{Verbatim}
 }

 

\subsection{\textcolor{Chapter }{MeetMaps}}
\logpage{[ 73, 5, 8 ]}\nobreak
\hyperdef{L}{X81A1A0E88570E42A}{}
{\noindent\textcolor{FuncColor}{$\triangleright$\enspace\texttt{MeetMaps({\mdseries\slshape paramap1, paramap2})\index{MeetMaps@\texttt{MeetMaps}}
\label{MeetMaps}
}\hfill{\scriptsize (function)}}\\


 For two parametrized maps \mbox{\texttt{\mdseries\slshape paramap1}} and \mbox{\texttt{\mdseries\slshape paramap2}}, \texttt{MeetMaps} changes \mbox{\texttt{\mdseries\slshape paramap1}} such that the image of class $i$ is the intersection of \mbox{\texttt{\mdseries\slshape paramap1}}\texttt{[}$i$\texttt{]} and \mbox{\texttt{\mdseries\slshape paramap2}}\texttt{[}$i$\texttt{]}. 

 If this implies that no images remain for a class, the position of such a
class is returned. If no such inconsistency occurs, \texttt{MeetMaps} returns \texttt{true}. 

 
\begin{Verbatim}[commandchars=!@|,fontsize=\small,frame=single,label=Example]
  !gapprompt@gap>| !gapinput@map1:= [ [ 1, 2 ], [ 3, 4 ], 5, 6, [ 7, 8, 9 ] ];;|
  !gapprompt@gap>| !gapinput@map2:= [ [ 1, 3 ], [ 3, 4 ], [ 5, 6 ], 6, [ 8, 9, 10 ] ];;|
  !gapprompt@gap>| !gapinput@MeetMaps( map1, map2 );  map1;|
  true
  [ 1, [ 3, 4 ], 5, 6, [ 8, 9 ] ]
\end{Verbatim}
 }

 

\subsection{\textcolor{Chapter }{CommutativeDiagram}}
\logpage{[ 73, 5, 9 ]}\nobreak
\hyperdef{L}{X8593A72A8193EC8B}{}
{\noindent\textcolor{FuncColor}{$\triangleright$\enspace\texttt{CommutativeDiagram({\mdseries\slshape paramap1, paramap2, paramap3, paramap4[, improvements]})\index{CommutativeDiagram@\texttt{CommutativeDiagram}}
\label{CommutativeDiagram}
}\hfill{\scriptsize (function)}}\\


 Let \mbox{\texttt{\mdseries\slshape paramap1}}, \mbox{\texttt{\mdseries\slshape paramap2}}, \mbox{\texttt{\mdseries\slshape paramap3}}, \mbox{\texttt{\mdseries\slshape paramap4}} be parametrized maps covering parametrized maps $f_1$, $f_2$, $f_3$, $f_4$ with the property that \texttt{CompositionMaps}$( f_2, f_1 )$ is equal to \texttt{CompositionMaps}$( f_4, f_3 )$. 

 \texttt{CommutativeDiagram} checks this consistency, and changes the arguments such that all possible
images are removed that cannot occur in the parametrized maps $f_i$. 

 The return value is \texttt{fail} if an inconsistency was found. Otherwise a record with the components \texttt{imp1}, \texttt{imp2}, \texttt{imp3}, \texttt{imp4} is returned, each bound to the list of positions where the corresponding
parametrized map was changed, 

 The optional argument \mbox{\texttt{\mdseries\slshape improvements}} must be a record with components \texttt{imp1}, \texttt{imp2}, \texttt{imp3}, \texttt{imp4}. If such a record is specified then only diagrams are considered where
entries of the $i$\texttt{\symbol{45}}th component occur as preimages of the $i$\texttt{\symbol{45}}th parametrized map. 

 When an inconsistency is detected, \texttt{CommutativeDiagram} immediately returns \texttt{fail}. Otherwise a record is returned that contains four lists \texttt{imp1}, $\ldots$, \texttt{imp4}: The $i$\texttt{\symbol{45}}th component is the list of classes where the $i$\texttt{\symbol{45}}th argument was changed. 

 
\begin{Verbatim}[commandchars=!@|,fontsize=\small,frame=single,label=Example]
  !gapprompt@gap>| !gapinput@map1:= [[ 1, 2, 3 ], [ 1, 3 ]];; map2:= [[ 1, 2 ], 1, [ 1, 3 ]];;|
  !gapprompt@gap>| !gapinput@map3:= [ [ 2, 3 ], 3 ];;  map4:= [ , 1, 2, [ 1, 2 ] ];;|
  !gapprompt@gap>| !gapinput@imp:= CommutativeDiagram( map1, map2, map3, map4 );|
  rec( imp1 := [ 2 ], imp2 := [ 1 ], imp3 := [  ], imp4 := [  ] )
  !gapprompt@gap>| !gapinput@map1;  map2;  map3;  map4;|
  [ [ 1, 2, 3 ], 1 ]
  [ 2, 1, [ 1, 3 ] ]
  [ [ 2, 3 ], 3 ]
  [ , 1, 2, [ 1, 2 ] ]
  !gapprompt@gap>| !gapinput@imp2:= CommutativeDiagram( map1, map2, map3, map4, imp );|
  rec( imp1 := [  ], imp2 := [  ], imp3 := [  ], imp4 := [  ] )
\end{Verbatim}
 }

 

\subsection{\textcolor{Chapter }{CheckFixedPoints}}
\logpage{[ 73, 5, 10 ]}\nobreak
\hyperdef{L}{X7B6EC10C7F7411E9}{}
{\noindent\textcolor{FuncColor}{$\triangleright$\enspace\texttt{CheckFixedPoints({\mdseries\slshape inside1, between, inside2})\index{CheckFixedPoints@\texttt{CheckFixedPoints}}
\label{CheckFixedPoints}
}\hfill{\scriptsize (function)}}\\


 Let \mbox{\texttt{\mdseries\slshape inside1}}, \mbox{\texttt{\mdseries\slshape between}}, \mbox{\texttt{\mdseries\slshape inside2}} be parametrized maps, where \mbox{\texttt{\mdseries\slshape between}} is assumed to map each fixed point of \mbox{\texttt{\mdseries\slshape inside1}} (that is, \mbox{\texttt{\mdseries\slshape inside1}}\texttt{[}$i$\texttt{] = }\mbox{\texttt{\mdseries\slshape i}}) to a fixed point of \mbox{\texttt{\mdseries\slshape inside2}} (that is, \mbox{\texttt{\mdseries\slshape between}}\texttt{[}$i$\texttt{]} is either an integer that is fixed by \mbox{\texttt{\mdseries\slshape inside2}} or a list that has nonempty intersection with the union of its images under \mbox{\texttt{\mdseries\slshape inside2}}). \texttt{CheckFixedPoints} changes \mbox{\texttt{\mdseries\slshape between}} and \mbox{\texttt{\mdseries\slshape inside2}} by removing all those entries violate this condition. 

 When an inconsistency is detected, \texttt{CheckFixedPoints} immediately returns \texttt{fail}. Otherwise the list of positions is returned where changes occurred. 

 
\begin{Verbatim}[commandchars=!@|,fontsize=\small,frame=single,label=Example]
  !gapprompt@gap>| !gapinput@subtbl:= CharacterTable( "L4(3).2_2" );;|
  !gapprompt@gap>| !gapinput@tbl:= CharacterTable( "O7(3)" );;|
  !gapprompt@gap>| !gapinput@fus:= InitFusion( subtbl, tbl );;  fus{ [ 48, 49 ] };|
  [ [ 54, 55, 56, 57 ], [ 54, 55, 56, 57 ] ]
  !gapprompt@gap>| !gapinput@CheckFixedPoints( ComputedPowerMaps( subtbl )[5], fus,|
  !gapprompt@>| !gapinput@       ComputedPowerMaps( tbl )[5] );|
  [ 48, 49 ]
  !gapprompt@gap>| !gapinput@fus{ [ 48, 49 ] };|
  [ [ 56, 57 ], [ 56, 57 ] ]
\end{Verbatim}
 }

 

\subsection{\textcolor{Chapter }{TransferDiagram}}
\logpage{[ 73, 5, 11 ]}\nobreak
\hyperdef{L}{X7AD5158E82AF1CD4}{}
{\noindent\textcolor{FuncColor}{$\triangleright$\enspace\texttt{TransferDiagram({\mdseries\slshape inside1, between, inside2[, improvements]})\index{TransferDiagram@\texttt{TransferDiagram}}
\label{TransferDiagram}
}\hfill{\scriptsize (function)}}\\


 Let \mbox{\texttt{\mdseries\slshape inside1}}, \mbox{\texttt{\mdseries\slshape between}}, \mbox{\texttt{\mdseries\slshape inside2}} be parametrized maps covering parametrized maps $m_1$, $f$, $m_2$ with the property that \texttt{CompositionMaps}$( m_2, f )$ is equal to \texttt{CompositionMaps}$( f, m_1 )$. 

 \texttt{TransferDiagram} checks this consistency, and changes the arguments such that all possible
images are removed that cannot occur in the parametrized maps $m_i$ and $f$. 

 So \texttt{TransferDiagram} is similar to \texttt{CommutativeDiagram} (\ref{CommutativeDiagram}), but \mbox{\texttt{\mdseries\slshape between}} occurs twice in each diagram checked. 

 If a record \mbox{\texttt{\mdseries\slshape improvements}} with fields \texttt{impinside1}, \texttt{impbetween}, and \texttt{impinside2} is specified, only those diagrams with elements of \texttt{impinside1} as preimages of \mbox{\texttt{\mdseries\slshape inside1}}, elements of \texttt{impbetween} as preimages of \mbox{\texttt{\mdseries\slshape between}} or elements of \texttt{impinside2} as preimages of \mbox{\texttt{\mdseries\slshape inside2}} are considered. 

 When an inconsistency is detected, \texttt{TransferDiagram} immediately returns \texttt{fail}. Otherwise a record is returned that contains three lists \texttt{impinside1}, \texttt{impbetween}, and \texttt{impinside2} of positions where the arguments were changed. 

 
\begin{Verbatim}[commandchars=!@|,fontsize=\small,frame=single,label=Example]
  !gapprompt@gap>| !gapinput@subtbl:= CharacterTable( "2F4(2)" );;  tbl:= CharacterTable( "Ru" );;|
  !gapprompt@gap>| !gapinput@fus:= InitFusion( subtbl, tbl );;|
  !gapprompt@gap>| !gapinput@permchar:= Sum( Irr( tbl ){ [ 1, 5, 6 ] } );;|
  !gapprompt@gap>| !gapinput@CheckPermChar( subtbl, tbl, fus, permchar );; fus;|
  [ 1, 2, 2, 4, 5, 7, 8, 9, 11, 14, 14, [ 13, 15 ], 16, [ 18, 19 ], 20,
    [ 25, 26 ], [ 25, 26 ], 5, 5, 6, 8, 14, [ 13, 15 ], [ 18, 19 ],
    [ 18, 19 ], [ 25, 26 ], [ 25, 26 ], 27, 27 ]
  !gapprompt@gap>| !gapinput@tr:= TransferDiagram(PowerMap( subtbl, 2), fus, PowerMap(tbl, 2));|
  rec( impbetween := [ 12, 23 ], impinside1 := [  ], impinside2 := [  ]
   )
  !gapprompt@gap>| !gapinput@tr:= TransferDiagram(PowerMap(subtbl, 3), fus, PowerMap( tbl, 3 ));|
  rec( impbetween := [ 14, 24, 25 ], impinside1 := [  ],
    impinside2 := [  ] )
  !gapprompt@gap>| !gapinput@tr:= TransferDiagram( PowerMap(subtbl, 3), fus, PowerMap(tbl, 3),|
  !gapprompt@>| !gapinput@            tr );|
  rec( impbetween := [  ], impinside1 := [  ], impinside2 := [  ] )
  !gapprompt@gap>| !gapinput@fus;|
  [ 1, 2, 2, 4, 5, 7, 8, 9, 11, 14, 14, 15, 16, 18, 20, [ 25, 26 ],
    [ 25, 26 ], 5, 5, 6, 8, 14, 13, 19, 19, [ 25, 26 ], [ 25, 26 ], 27,
    27 ]
\end{Verbatim}
 }

 

\subsection{\textcolor{Chapter }{TestConsistencyMaps}}
\logpage{[ 73, 5, 12 ]}\nobreak
\hyperdef{L}{X78487F03852A503B}{}
{\noindent\textcolor{FuncColor}{$\triangleright$\enspace\texttt{TestConsistencyMaps({\mdseries\slshape powermap1, fusionmap, powermap2[, fusimp]})\index{TestConsistencyMaps@\texttt{TestConsistencyMaps}}
\label{TestConsistencyMaps}
}\hfill{\scriptsize (function)}}\\


 Let \mbox{\texttt{\mdseries\slshape powermap1}} and \mbox{\texttt{\mdseries\slshape powermap2}} be lists of parametrized maps, and \mbox{\texttt{\mdseries\slshape fusionmap}} a parametrized map, such that for each $i$, the $i$\texttt{\symbol{45}}th entry in \mbox{\texttt{\mdseries\slshape powermap1}}, \mbox{\texttt{\mdseries\slshape fusionmap}}, and the $i$\texttt{\symbol{45}}th entry in \mbox{\texttt{\mdseries\slshape powermap2}} (if bound) are valid arguments for \texttt{TransferDiagram} (\ref{TransferDiagram}). So a typical situation for applying \texttt{TestConsistencyMaps} is that \mbox{\texttt{\mdseries\slshape fusionmap}} is an approximation of a class fusion, and \mbox{\texttt{\mdseries\slshape powermap1}}, \mbox{\texttt{\mdseries\slshape powermap2}} are the lists of power maps of the subgroup and the group. 

 \texttt{TestConsistencyMaps} repeatedly applies \texttt{TransferDiagram} (\ref{TransferDiagram}) to these arguments for all $i$ until no more changes occur. 

 If a list \mbox{\texttt{\mdseries\slshape fusimp}} is specified then only those diagrams with elements of \mbox{\texttt{\mdseries\slshape fusimp}} as preimages of \mbox{\texttt{\mdseries\slshape fusionmap}} are considered. 

 When an inconsistency is detected, \texttt{TestConsistencyMaps} immediately returns \texttt{false}. Otherwise \texttt{true} is returned. 

 
\begin{Verbatim}[commandchars=!@|,fontsize=\small,frame=single,label=Example]
  !gapprompt@gap>| !gapinput@subtbl:= CharacterTable( "2F4(2)" );;  tbl:= CharacterTable( "Ru" );;|
  !gapprompt@gap>| !gapinput@fus:= InitFusion( subtbl, tbl );;|
  !gapprompt@gap>| !gapinput@permchar:= Sum( Irr( tbl ){ [ 1, 5, 6 ] } );;|
  !gapprompt@gap>| !gapinput@CheckPermChar( subtbl, tbl, fus, permchar );; fus;|
  [ 1, 2, 2, 4, 5, 7, 8, 9, 11, 14, 14, [ 13, 15 ], 16, [ 18, 19 ], 20,
    [ 25, 26 ], [ 25, 26 ], 5, 5, 6, 8, 14, [ 13, 15 ], [ 18, 19 ],
    [ 18, 19 ], [ 25, 26 ], [ 25, 26 ], 27, 27 ]
  !gapprompt@gap>| !gapinput@TestConsistencyMaps( ComputedPowerMaps( subtbl ), fus,|
  !gapprompt@>| !gapinput@       ComputedPowerMaps( tbl ) );|
  true
  !gapprompt@gap>| !gapinput@fus;|
  [ 1, 2, 2, 4, 5, 7, 8, 9, 11, 14, 14, 15, 16, 18, 20, [ 25, 26 ],
    [ 25, 26 ], 5, 5, 6, 8, 14, 13, 19, 19, [ 25, 26 ], [ 25, 26 ], 27,
    27 ]
  !gapprompt@gap>| !gapinput@Indeterminateness( fus );|
  16
\end{Verbatim}
 }

 

\subsection{\textcolor{Chapter }{Indeterminateness}}
\logpage{[ 73, 5, 13 ]}\nobreak
\hyperdef{L}{X7DAD6EA585D74615}{}
{\noindent\textcolor{FuncColor}{$\triangleright$\enspace\texttt{Indeterminateness({\mdseries\slshape paramap})\index{Indeterminateness@\texttt{Indeterminateness}}
\label{Indeterminateness}
}\hfill{\scriptsize (function)}}\\


 For a parametrized map \mbox{\texttt{\mdseries\slshape paramap}}, \texttt{Indeterminateness} returns the number of maps contained in \mbox{\texttt{\mdseries\slshape paramap}}, that is, the product of lengths of lists in \mbox{\texttt{\mdseries\slshape paramap}} denoting lists of several images. 

 
\begin{Verbatim}[commandchars=!@|,fontsize=\small,frame=single,label=Example]
  !gapprompt@gap>| !gapinput@Indeterminateness([ 1, [ 2, 3 ], [ 4, 5 ], [ 6, 7, 8, 9, 10 ], 11 ]);|
  20
\end{Verbatim}
 }

 

\subsection{\textcolor{Chapter }{PrintAmbiguity}}
\logpage{[ 73, 5, 14 ]}\nobreak
\hyperdef{L}{X7888BDC88304BE5A}{}
{\noindent\textcolor{FuncColor}{$\triangleright$\enspace\texttt{PrintAmbiguity({\mdseries\slshape list, paramap})\index{PrintAmbiguity@\texttt{PrintAmbiguity}}
\label{PrintAmbiguity}
}\hfill{\scriptsize (function)}}\\


 For each map in the list \mbox{\texttt{\mdseries\slshape list}}, \texttt{PrintAmbiguity} prints its position in \mbox{\texttt{\mdseries\slshape list}}, the indeterminateness (see{\nobreakspace}\texttt{Indeterminateness} (\ref{Indeterminateness})) of the composition with the parametrized map \mbox{\texttt{\mdseries\slshape paramap}}, and the list of positions where a list of images occurs in this composition. 

 
\begin{Verbatim}[commandchars=!@|,fontsize=\small,frame=single,label=Example]
  !gapprompt@gap>| !gapinput@paramap:= [ 1, [ 2, 3 ], [ 3, 4 ], [ 2, 3, 4 ], 5 ];;|
  !gapprompt@gap>| !gapinput@list:= [ [ 1, 1, 1, 1, 1 ], [ 1, 1, 2, 2, 3 ], [ 1, 2, 3, 4, 5 ] ];;|
  !gapprompt@gap>| !gapinput@PrintAmbiguity( list, paramap );|
  1 1 [  ]
  2 4 [ 2, 4 ]
  3 12 [ 2, 3, 4 ]
\end{Verbatim}
 }

 

\subsection{\textcolor{Chapter }{ContainedSpecialVectors}}
\logpage{[ 73, 5, 15 ]}\nobreak
\hyperdef{L}{X7F957B1481E10A0C}{}
{\noindent\textcolor{FuncColor}{$\triangleright$\enspace\texttt{ContainedSpecialVectors({\mdseries\slshape tbl, chars, paracharacter, func})\index{ContainedSpecialVectors@\texttt{ContainedSpecialVectors}}
\label{ContainedSpecialVectors}
}\hfill{\scriptsize (function)}}\\
\noindent\textcolor{FuncColor}{$\triangleright$\enspace\texttt{IntScalarProducts({\mdseries\slshape tbl, chars, candidate})\index{IntScalarProducts@\texttt{IntScalarProducts}}
\label{IntScalarProducts}
}\hfill{\scriptsize (function)}}\\
\noindent\textcolor{FuncColor}{$\triangleright$\enspace\texttt{NonnegIntScalarProducts({\mdseries\slshape tbl, chars, candidate})\index{NonnegIntScalarProducts@\texttt{NonnegIntScalarProducts}}
\label{NonnegIntScalarProducts}
}\hfill{\scriptsize (function)}}\\
\noindent\textcolor{FuncColor}{$\triangleright$\enspace\texttt{ContainedPossibleVirtualCharacters({\mdseries\slshape tbl, chars, paracharacter})\index{ContainedPossibleVirtualCharacters@\texttt{ContainedPossibleVirtualCharacters}}
\label{ContainedPossibleVirtualCharacters}
}\hfill{\scriptsize (function)}}\\
\noindent\textcolor{FuncColor}{$\triangleright$\enspace\texttt{ContainedPossibleCharacters({\mdseries\slshape tbl, chars, paracharacter})\index{ContainedPossibleCharacters@\texttt{ContainedPossibleCharacters}}
\label{ContainedPossibleCharacters}
}\hfill{\scriptsize (function)}}\\


 Let \mbox{\texttt{\mdseries\slshape tbl}} be an ordinary character table, \mbox{\texttt{\mdseries\slshape chars}} a list of class functions (or values lists), \mbox{\texttt{\mdseries\slshape paracharacter}} a parametrized class function of \mbox{\texttt{\mdseries\slshape tbl}}, and \mbox{\texttt{\mdseries\slshape func}} a function that expects the three arguments \mbox{\texttt{\mdseries\slshape tbl}}, \mbox{\texttt{\mdseries\slshape chars}}, and a values list of a class function, and that returns either \texttt{true} or \texttt{false}. 

 \texttt{ContainedSpecialVectors} returns the list of all those elements \mbox{\texttt{\mdseries\slshape vec}} of \mbox{\texttt{\mdseries\slshape paracharacter}} that have integral norm, have integral scalar product with the principal
character of \mbox{\texttt{\mdseries\slshape tbl}}, and that satisfy \mbox{\texttt{\mdseries\slshape func}}\texttt{( }\mbox{\texttt{\mdseries\slshape tbl}}, \mbox{\texttt{\mdseries\slshape chars}}, \mbox{\texttt{\mdseries\slshape vec}} \texttt{) = }\texttt{true}. 

 Two special cases of \mbox{\texttt{\mdseries\slshape func}} are the check whether the scalar products in \mbox{\texttt{\mdseries\slshape tbl}} between the vector \mbox{\texttt{\mdseries\slshape vec}} and all lists in \mbox{\texttt{\mdseries\slshape chars}} are integers or nonnegative integers, respectively. These functions are
accessible as global variables \texttt{IntScalarProducts} and \texttt{NonnegIntScalarProducts}, and \texttt{ContainedPossibleVirtualCharacters} and \texttt{ContainedPossibleCharacters} provide access to these special cases of \texttt{ContainedSpecialVectors}. 

 
\begin{Verbatim}[commandchars=!@|,fontsize=\small,frame=single,label=Example]
  !gapprompt@gap>| !gapinput@subtbl:= CharacterTable( "HSM12" );;  tbl:= CharacterTable( "HS" );;|
  !gapprompt@gap>| !gapinput@fus:= InitFusion( subtbl, tbl );;|
  !gapprompt@gap>| !gapinput@rest:= CompositionMaps( Irr( tbl )[8], fus );|
  [ 231, [ -9, 7 ], [ -9, 7 ], [ -9, 7 ], 6, 15, 15, [ -1, 15 ],
    [ -1, 15 ], 1, [ 1, 6 ], [ 1, 6 ], [ 1, 6 ], [ 1, 6 ], [ -2, 0 ],
    [ 1, 2 ], [ 1, 2 ], [ 1, 2 ], 0, 0, 1, 0, 0, 0, 0 ]
  !gapprompt@gap>| !gapinput@irr:= Irr( subtbl );;|
  !gapprompt@gap>| !gapinput@# no further condition|
  !gapprompt@gap>| !gapinput@cont1:= ContainedSpecialVectors( subtbl, irr, rest,|
  !gapprompt@>| !gapinput@               function( tbl, chars, vec ) return true; end );;|
  !gapprompt@gap>| !gapinput@Length( cont1 );|
  24
  !gapprompt@gap>| !gapinput@# require scalar products to be integral|
  !gapprompt@gap>| !gapinput@cont2:= ContainedSpecialVectors( subtbl, irr, rest,|
  !gapprompt@>| !gapinput@               IntScalarProducts );|
  [ [ 231, 7, -9, -9, 6, 15, 15, -1, -1, 1, 6, 6, 1, 1, -2, 1, 2, 2, 0,
        0, 1, 0, 0, 0, 0 ],
    [ 231, 7, -9, 7, 6, 15, 15, -1, -1, 1, 6, 6, 1, 1, -2, 1, 2, 2, 0,
        0, 1, 0, 0, 0, 0 ],
    [ 231, 7, -9, -9, 6, 15, 15, 15, 15, 1, 6, 6, 1, 1, -2, 1, 2, 2, 0,
        0, 1, 0, 0, 0, 0 ],
    [ 231, 7, -9, 7, 6, 15, 15, 15, 15, 1, 6, 6, 1, 1, -2, 1, 2, 2, 0,
        0, 1, 0, 0, 0, 0 ] ]
  !gapprompt@gap>| !gapinput@# additionally require scalar products to be nonnegative|
  !gapprompt@gap>| !gapinput@cont3:= ContainedSpecialVectors( subtbl, irr, rest,|
  !gapprompt@>| !gapinput@               NonnegIntScalarProducts );|
  [ [ 231, 7, -9, -9, 6, 15, 15, -1, -1, 1, 6, 6, 1, 1, -2, 1, 2, 2, 0,
        0, 1, 0, 0, 0, 0 ],
    [ 231, 7, -9, 7, 6, 15, 15, -1, -1, 1, 6, 6, 1, 1, -2, 1, 2, 2, 0,
        0, 1, 0, 0, 0, 0 ] ]
  !gapprompt@gap>| !gapinput@cont2 = ContainedPossibleVirtualCharacters( subtbl, irr, rest );|
  true
  !gapprompt@gap>| !gapinput@cont3 = ContainedPossibleCharacters( subtbl, irr, rest );|
  true
\end{Verbatim}
 }

 

\subsection{\textcolor{Chapter }{CollapsedMat}}
\logpage{[ 73, 5, 16 ]}\nobreak
\hyperdef{L}{X84F87C2282EFB0EE}{}
{\noindent\textcolor{FuncColor}{$\triangleright$\enspace\texttt{CollapsedMat({\mdseries\slshape mat, maps})\index{CollapsedMat@\texttt{CollapsedMat}}
\label{CollapsedMat}
}\hfill{\scriptsize (function)}}\\


 is a record with the components 

 
\begin{description}
\item[{\texttt{fusion}}]  fusion that collapses those columns of \mbox{\texttt{\mdseries\slshape mat}} that are equal in \mbox{\texttt{\mdseries\slshape mat}} and also for all maps in the list \mbox{\texttt{\mdseries\slshape maps}}, 
\item[{\texttt{mat}}]  the image of \mbox{\texttt{\mdseries\slshape mat}} under that fusion. 
\end{description}
 

 
\begin{Verbatim}[commandchars=!@|,fontsize=\small,frame=single,label=Example]
  !gapprompt@gap>| !gapinput@mat:= [ [ 1, 1, 1, 1 ], [ 2, -1, 0, 0 ], [ 4, 4, 1, 1 ] ];;|
  !gapprompt@gap>| !gapinput@coll:= CollapsedMat( mat, [] );|
  rec( fusion := [ 1, 2, 3, 3 ],
    mat := [ [ 1, 1, 1 ], [ 2, -1, 0 ], [ 4, 4, 1 ] ] )
  !gapprompt@gap>| !gapinput@List( last.mat, x -> x{ last.fusion } ) = mat;|
  true
  !gapprompt@gap>| !gapinput@coll:= CollapsedMat( mat, [ [ 1, 1, 1, 2 ] ] );|
  rec( fusion := [ 1, 2, 3, 4 ],
    mat := [ [ 1, 1, 1, 1 ], [ 2, -1, 0, 0 ], [ 4, 4, 1, 1 ] ] )
\end{Verbatim}
 }

 

\subsection{\textcolor{Chapter }{ContainedDecomposables}}
\logpage{[ 73, 5, 17 ]}\nobreak
\hyperdef{L}{X81F1137A874EB962}{}
{\noindent\textcolor{FuncColor}{$\triangleright$\enspace\texttt{ContainedDecomposables({\mdseries\slshape constituents, moduls, parachar, func})\index{ContainedDecomposables@\texttt{ContainedDecomposables}}
\label{ContainedDecomposables}
}\hfill{\scriptsize (function)}}\\
\noindent\textcolor{FuncColor}{$\triangleright$\enspace\texttt{ContainedCharacters({\mdseries\slshape tbl, constituents, parachar})\index{ContainedCharacters@\texttt{ContainedCharacters}}
\label{ContainedCharacters}
}\hfill{\scriptsize (function)}}\\


 For these functions, let \mbox{\texttt{\mdseries\slshape constituents}} be a list of \emph{rational} class functions, \mbox{\texttt{\mdseries\slshape moduls}} a list of positive integers, \mbox{\texttt{\mdseries\slshape parachar}} a parametrized rational class function, \mbox{\texttt{\mdseries\slshape func}} a function that returns either \texttt{true} or \texttt{false} when called with (a values list of) a class function, and \mbox{\texttt{\mdseries\slshape tbl}} a character table. 

 \texttt{ContainedDecomposables} returns the set of all elements $\chi$ of \mbox{\texttt{\mdseries\slshape parachar}} that satisfy \mbox{\texttt{\mdseries\slshape func}}$( \chi ) =$ \texttt{true} and that lie in the ${\ensuremath{\mathbb Z}}$\texttt{\symbol{45}}lattice spanned by \mbox{\texttt{\mdseries\slshape constituents}}, modulo \mbox{\texttt{\mdseries\slshape moduls}}. The latter means they lie in the ${\ensuremath{\mathbb Z}}$\texttt{\symbol{45}}lattice spanned by \mbox{\texttt{\mdseries\slshape constituents}} and the set $\{ \mbox{\texttt{\mdseries\slshape moduls}}[i] \cdot e_i; 1 \leq i \leq n \}$ where $n$ is the length of \mbox{\texttt{\mdseries\slshape parachar}} and $e_i$ is the $i$\texttt{\symbol{45}}th standard basis vector. 

 One application of \texttt{ContainedDecomposables} is the following. \mbox{\texttt{\mdseries\slshape constituents}} is a list of (values lists of) rational characters of an ordinary character
table \mbox{\texttt{\mdseries\slshape tbl}}, \mbox{\texttt{\mdseries\slshape moduls}} is the list of centralizer orders of \mbox{\texttt{\mdseries\slshape tbl}} (see{\nobreakspace}\texttt{SizesCentralizers} (\ref{SizesCentralizers})), and \mbox{\texttt{\mdseries\slshape func}} checks whether a vector in the lattice mentioned above has nonnegative
integral scalar product in \mbox{\texttt{\mdseries\slshape tbl}} with all entries of \mbox{\texttt{\mdseries\slshape constituents}}. This situation is handled by \texttt{ContainedCharacters}. Note that the entries of the result list are \emph{not} necessary linear combinations of \mbox{\texttt{\mdseries\slshape constituents}}, and they are \emph{not} necessarily characters of \mbox{\texttt{\mdseries\slshape tbl}}. 

 
\begin{Verbatim}[commandchars=!@|,fontsize=\small,frame=single,label=Example]
  !gapprompt@gap>| !gapinput@subtbl:= CharacterTable( "HSM12" );;  tbl:= CharacterTable( "HS" );;|
  !gapprompt@gap>| !gapinput@rat:= RationalizedMat( Irr( subtbl ) );;|
  !gapprompt@gap>| !gapinput@fus:= InitFusion( subtbl, tbl );;|
  !gapprompt@gap>| !gapinput@rest:= CompositionMaps( Irr( tbl )[8], fus );|
  [ 231, [ -9, 7 ], [ -9, 7 ], [ -9, 7 ], 6, 15, 15, [ -1, 15 ],
    [ -1, 15 ], 1, [ 1, 6 ], [ 1, 6 ], [ 1, 6 ], [ 1, 6 ], [ -2, 0 ],
    [ 1, 2 ], [ 1, 2 ], [ 1, 2 ], 0, 0, 1, 0, 0, 0, 0 ]
  !gapprompt@gap>| !gapinput@# compute all vectors in the lattice|
  !gapprompt@gap>| !gapinput@ContainedDecomposables( rat, SizesCentralizers( subtbl ), rest,|
  !gapprompt@>| !gapinput@       ReturnTrue );|
  [ [ 231, 7, -9, -9, 6, 15, 15, -1, -1, 1, 6, 6, 1, 1, -2, 1, 2, 2, 0,
        0, 1, 0, 0, 0, 0 ],
    [ 231, 7, -9, -9, 6, 15, 15, 15, 15, 1, 6, 6, 1, 1, -2, 1, 2, 2, 0,
        0, 1, 0, 0, 0, 0 ],
    [ 231, 7, -9, 7, 6, 15, 15, -1, -1, 1, 6, 6, 1, 1, -2, 1, 2, 2, 0,
        0, 1, 0, 0, 0, 0 ],
    [ 231, 7, -9, 7, 6, 15, 15, 15, 15, 1, 6, 6, 1, 1, -2, 1, 2, 2, 0,
        0, 1, 0, 0, 0, 0 ] ]
  !gapprompt@gap>| !gapinput@# compute only those vectors that are characters|
  !gapprompt@gap>| !gapinput@ContainedDecomposables( rat, SizesCentralizers( subtbl ), rest,|
  !gapprompt@>| !gapinput@       x -> NonnegIntScalarProducts( subtbl, Irr( subtbl ), x ) );|
  [ [ 231, 7, -9, -9, 6, 15, 15, -1, -1, 1, 6, 6, 1, 1, -2, 1, 2, 2, 0,
        0, 1, 0, 0, 0, 0 ],
    [ 231, 7, -9, 7, 6, 15, 15, -1, -1, 1, 6, 6, 1, 1, -2, 1, 2, 2, 0,
        0, 1, 0, 0, 0, 0 ] ]
\end{Verbatim}
 }

 }

  
\section{\textcolor{Chapter }{Subroutines for the Construction of Power Maps}}\label{Subroutines for the Construction of Power Maps}
\logpage{[ 73, 6, 0 ]}
\hyperdef{L}{X86472A217D6C3CE7}{}
{
  In the argument lists of the functions \texttt{Congruences} (\ref{Congruences:for character tables}), \texttt{ConsiderKernels} (\ref{ConsiderKernels}), and \texttt{ConsiderSmallerPowerMaps} (\ref{ConsiderSmallerPowerMaps}), \mbox{\texttt{\mdseries\slshape tbl}} is an ordinary character table, \mbox{\texttt{\mdseries\slshape chars}} a list of (values lists of) characters of \mbox{\texttt{\mdseries\slshape tbl}}, \mbox{\texttt{\mdseries\slshape prime}} a prime integer, \mbox{\texttt{\mdseries\slshape approxmap}} a parametrized map that is an approximation for the \mbox{\texttt{\mdseries\slshape prime}}\texttt{\symbol{45}}th power map of \mbox{\texttt{\mdseries\slshape tbl}} (e.g., a list returned by \texttt{InitPowerMap} (\ref{InitPowerMap}), and \mbox{\texttt{\mdseries\slshape quick}} a Boolean. 

 The \mbox{\texttt{\mdseries\slshape quick}} value \texttt{true} means that only those classes are considered for which \mbox{\texttt{\mdseries\slshape approxmap}} lists more than one possible image. 

\subsection{\textcolor{Chapter }{InitPowerMap}}
\logpage{[ 73, 6, 1 ]}\nobreak
\hyperdef{L}{X85D068D77C3C041C}{}
{\noindent\textcolor{FuncColor}{$\triangleright$\enspace\texttt{InitPowerMap({\mdseries\slshape tbl, prime})\index{InitPowerMap@\texttt{InitPowerMap}}
\label{InitPowerMap}
}\hfill{\scriptsize (function)}}\\


 For an ordinary character table \mbox{\texttt{\mdseries\slshape tbl}} and a prime \mbox{\texttt{\mdseries\slshape prime}}, \texttt{InitPowerMap} returns a parametrized map that is a first approximation of the \mbox{\texttt{\mdseries\slshape prime}}\texttt{\symbol{45}}th powermap of \mbox{\texttt{\mdseries\slshape tbl}}, using the conditions 1.{\nobreakspace}and 2.{\nobreakspace}listed in the
description of \texttt{PossiblePowerMaps} (\ref{PossiblePowerMaps}). 

 If there are classes for which no images are possible, according to these
criteria, then \texttt{fail} is returned. 

 
\begin{Verbatim}[commandchars=!@|,fontsize=\small,frame=single,label=Example]
  !gapprompt@gap>| !gapinput@t:= CharacterTable( "U4(3).4" );;|
  !gapprompt@gap>| !gapinput@pow:= InitPowerMap( t, 2 );|
  [ 1, 1, 3, 4, 5, [ 2, 16 ], [ 2, 16, 17 ], 8, 3, [ 3, 4 ],
    [ 11, 12 ], [ 11, 12 ], [ 6, 7, 18, 19, 30, 31, 32, 33 ], 14,
    [ 9, 20 ], 1, 1, 2, 2, 3, [ 3, 4, 5 ], [ 3, 4, 5 ],
    [ 6, 7, 18, 19, 30, 31, 32, 33 ], 8, 9, 9, [ 9, 10, 20, 21, 22 ],
    [ 11, 12 ], [ 11, 12 ], 16, 16, [ 2, 16 ], [ 2, 16 ], 17, 17,
    [ 6, 18, 30, 31, 32, 33 ], [ 6, 18, 30, 31, 32, 33 ],
    [ 6, 7, 18, 19, 30, 31, 32, 33 ], [ 6, 7, 18, 19, 30, 31, 32, 33 ],
    20, 20, [ 9, 20 ], [ 9, 20 ], [ 9, 10, 20, 21, 22 ],
    [ 9, 10, 20, 21, 22 ], 24, 24, [ 15, 25, 26, 40, 41, 42, 43 ],
    [ 15, 25, 26, 40, 41, 42, 43 ], [ 28, 29 ], [ 28, 29 ], [ 28, 29 ],
    [ 28, 29 ] ]
\end{Verbatim}
 }

 

\subsection{\textcolor{Chapter }{Congruences (for character tables)}}
\logpage{[ 73, 6, 2 ]}\nobreak
\hyperdef{L}{X7B27749E7BF54EBB}{}
{\noindent\textcolor{FuncColor}{$\triangleright$\enspace\texttt{Congruences({\mdseries\slshape tbl, chars, approxmap, prime, quick})\index{Congruences@\texttt{Congruences}!for character tables}
\label{Congruences:for character tables}
}\hfill{\scriptsize (function)}}\\


 \texttt{Congruences} replaces the entries of \mbox{\texttt{\mdseries\slshape approxmap}} by improved values, according to condition 3.{\nobreakspace}listed in the
description of \texttt{PossiblePowerMaps} (\ref{PossiblePowerMaps}). 

 For each class for which no images are possible according to the tests, the
new value of \mbox{\texttt{\mdseries\slshape approxmap}} is an empty list. \texttt{Congruences} returns \texttt{true} if no such inconsistencies occur, and \texttt{false} otherwise. 

 
\begin{Verbatim}[commandchars=!@|,fontsize=\small,frame=single,label=Example]
  !gapprompt@gap>| !gapinput@Congruences( t, Irr( t ), pow, 2, false );  pow;|
  true
  [ 1, 1, 3, 4, 5, 2, 2, 8, 3, 4, 11, 12, [ 6, 7 ], 14, 9, 1, 1, 2, 2,
    3, 4, 5, [ 6, 7 ], 8, 9, 9, 10, 11, 12, 16, 16, 16, 16, 17, 17, 18,
    18, [ 18, 19 ], [ 18, 19 ], 20, 20, 20, 20, 22, 22, 24, 24,
    [ 25, 26 ], [ 25, 26 ], 28, 28, 29, 29 ]
\end{Verbatim}
 }

 

\subsection{\textcolor{Chapter }{ConsiderKernels}}
\logpage{[ 73, 6, 3 ]}\nobreak
\hyperdef{L}{X7D31B1548205E222}{}
{\noindent\textcolor{FuncColor}{$\triangleright$\enspace\texttt{ConsiderKernels({\mdseries\slshape tbl, chars, approxmap, prime, quick})\index{ConsiderKernels@\texttt{ConsiderKernels}}
\label{ConsiderKernels}
}\hfill{\scriptsize (function)}}\\


 \texttt{ConsiderKernels} replaces the entries of \mbox{\texttt{\mdseries\slshape approxmap}} by improved values, according to condition 4.{\nobreakspace}listed in the
description of \texttt{PossiblePowerMaps} (\ref{PossiblePowerMaps}). 

 \texttt{Congruences} (\ref{Congruences:for character tables}) returns \texttt{true} if the orders of the kernels of all characters in \mbox{\texttt{\mdseries\slshape chars}} divide the order of the group of \mbox{\texttt{\mdseries\slshape tbl}}, and \texttt{false} otherwise. 

 
\begin{Verbatim}[commandchars=!@|,fontsize=\small,frame=single,label=Example]
  !gapprompt@gap>| !gapinput@t:= CharacterTable( "A7.2" );;  init:= InitPowerMap( t, 2 );|
  [ 1, 1, 3, 4, [ 2, 9, 10 ], 6, 3, 8, 1, 1, [ 2, 9, 10 ], 3, [ 3, 4 ],
    6, [ 7, 12 ] ]
  !gapprompt@gap>| !gapinput@ConsiderKernels( t, Irr( t ), init, 2, false );|
  true
  !gapprompt@gap>| !gapinput@init;|
  [ 1, 1, 3, 4, 2, 6, 3, 8, 1, 1, 2, 3, [ 3, 4 ], 6, 7 ]
\end{Verbatim}
 }

 

\subsection{\textcolor{Chapter }{ConsiderSmallerPowerMaps}}
\logpage{[ 73, 6, 4 ]}\nobreak
\hyperdef{L}{X7DD1DCF3865E0017}{}
{\noindent\textcolor{FuncColor}{$\triangleright$\enspace\texttt{ConsiderSmallerPowerMaps({\mdseries\slshape tbl, approxmap, prime, quick})\index{ConsiderSmallerPowerMaps@\texttt{ConsiderSmallerPowerMaps}}
\label{ConsiderSmallerPowerMaps}
}\hfill{\scriptsize (function)}}\\


 \texttt{ConsiderSmallerPowerMaps} replaces the entries of \mbox{\texttt{\mdseries\slshape approxmap}} by improved values, according to condition 5.{\nobreakspace}listed in the
description of \texttt{PossiblePowerMaps} (\ref{PossiblePowerMaps}). 

 \texttt{ConsiderSmallerPowerMaps} returns \texttt{true} if each class admits at least one image after the checks, otherwise \texttt{false} is returned. If no element orders of \mbox{\texttt{\mdseries\slshape tbl}} are stored (see{\nobreakspace}\texttt{OrdersClassRepresentatives} (\ref{OrdersClassRepresentatives})) then \texttt{true} is returned without any tests. 

 
\begin{Verbatim}[commandchars=!@|,fontsize=\small,frame=single,label=Example]
  !gapprompt@gap>| !gapinput@t:= CharacterTable( "3.A6" );;  init:= InitPowerMap( t, 5 );|
  [ 1, [ 2, 3 ], [ 2, 3 ], 4, [ 5, 6 ], [ 5, 6 ], [ 7, 8 ], [ 7, 8 ],
    9, [ 10, 11 ], [ 10, 11 ], 1, [ 2, 3 ], [ 2, 3 ], 1, [ 2, 3 ],
    [ 2, 3 ] ]
  !gapprompt@gap>| !gapinput@Indeterminateness( init );|
  4096
  !gapprompt@gap>| !gapinput@ConsiderSmallerPowerMaps( t, init, 5, false );|
  true
  !gapprompt@gap>| !gapinput@Indeterminateness( init );|
  256
\end{Verbatim}
 }

 

\subsection{\textcolor{Chapter }{MinusCharacter}}
\logpage{[ 73, 6, 5 ]}\nobreak
\hyperdef{L}{X805B6C1C78AA5DB6}{}
{\noindent\textcolor{FuncColor}{$\triangleright$\enspace\texttt{MinusCharacter({\mdseries\slshape character, primepowermap, prime})\index{MinusCharacter@\texttt{MinusCharacter}}
\label{MinusCharacter}
}\hfill{\scriptsize (function)}}\\


 Let \mbox{\texttt{\mdseries\slshape character}} be (the list of values of) a class function $\chi$, \mbox{\texttt{\mdseries\slshape prime}} a prime integer $p$, and \mbox{\texttt{\mdseries\slshape primepowermap}} a parametrized map that is an approximation of the $p$\texttt{\symbol{45}}th power map for the character table of $\chi$. \texttt{MinusCharacter} returns the parametrized map of values of $\chi^{{p-}}$, which is defined by $\chi^{{p-}}(g) = ( \chi(g)^p - \chi(g^p) ) / p$. 

 
\begin{Verbatim}[commandchars=!@|,fontsize=\small,frame=single,label=Example]
  !gapprompt@gap>| !gapinput@tbl:= CharacterTable( "S7" );;  pow:= InitPowerMap( tbl, 2 );;|
  !gapprompt@gap>| !gapinput@pow;|
  [ 1, 1, 3, 4, [ 2, 9, 10 ], 6, 3, 8, 1, 1, [ 2, 9, 10 ], 3, [ 3, 4 ],
    6, [ 7, 12 ] ]
  !gapprompt@gap>| !gapinput@chars:= Irr( tbl ){ [ 2 .. 5 ] };;|
  !gapprompt@gap>| !gapinput@List( chars, x -> MinusCharacter( x, pow, 2 ) );|
  [ [ 0, 0, 0, 0, [ 0, 1 ], 0, 0, 0, 0, 0, [ 0, 1 ], 0, 0, 0, [ 0, 1 ] ]
      ,
    [ 15, -1, 3, 0, [ -2, -1, 0 ], 0, -1, 1, 5, -3, [ 0, 1, 2 ], -1, 0,
        0, [ 0, 1 ] ],
    [ 15, -1, 3, 0, [ -1, 0, 2 ], 0, -1, 1, 5, -3, [ 1, 2, 4 ], -1, 0,
        0, 1 ],
    [ 190, -2, 1, 1, [ 0, 2 ], 0, 1, 1, -10, -10, [ 0, 2 ], -1, -1, 0,
        [ -1, 0 ] ] ]
\end{Verbatim}
 }

 

\subsection{\textcolor{Chapter }{PowerMapsAllowedBySymmetrizations}}
\logpage{[ 73, 6, 6 ]}\nobreak
\hyperdef{L}{X808CCF6087D5B661}{}
{\noindent\textcolor{FuncColor}{$\triangleright$\enspace\texttt{PowerMapsAllowedBySymmetrizations({\mdseries\slshape tbl, subchars, chars, approxmap, prime, parameters})\index{PowerMapsAllowedBySymmetrizations@\texttt{PowerMapsAllowedBySymmetrizations}}
\label{PowerMapsAllowedBySymmetrizations}
}\hfill{\scriptsize (function)}}\\


 Let \mbox{\texttt{\mdseries\slshape tbl}} be an ordinary character table, \mbox{\texttt{\mdseries\slshape prime}} a prime integer, \mbox{\texttt{\mdseries\slshape approxmap}} a parametrized map that is an approximation of the \mbox{\texttt{\mdseries\slshape prime}}\texttt{\symbol{45}}th power map of \mbox{\texttt{\mdseries\slshape tbl}} (e.g., a list returned by \texttt{InitPowerMap} (\ref{InitPowerMap}), \mbox{\texttt{\mdseries\slshape chars}} and \mbox{\texttt{\mdseries\slshape subchars}} two lists of (values lists of) characters of \mbox{\texttt{\mdseries\slshape tbl}}, and \mbox{\texttt{\mdseries\slshape parameters}} a record with components \texttt{maxlen}, \texttt{minamb}, \texttt{maxamb} (three integers), \texttt{quick} (a Boolean), and \texttt{contained} (a function). Usual values of \texttt{contained} are \texttt{ContainedCharacters} (\ref{ContainedCharacters}) or \texttt{ContainedPossibleCharacters} (\ref{ContainedPossibleCharacters}). 

 \texttt{PowerMapsAllowedBySymmetrizations} replaces the entries of \mbox{\texttt{\mdseries\slshape approxmap}} by improved values, according to condition 6.{\nobreakspace}listed in the
description of \texttt{PossiblePowerMaps} (\ref{PossiblePowerMaps}). 

 More precisely, the strategy used is as follows. 

 First, for each $\chi \in \mbox{\texttt{\mdseries\slshape chars}}$, let \texttt{minus:= MinusCharacter(}$\chi$\texttt{, \mbox{\texttt{\mdseries\slshape approxmap}}, \mbox{\texttt{\mdseries\slshape prime}})}. 
\begin{itemize}
\item  If \texttt{Indeterminateness( minus )}$ = 1$ and \texttt{\mbox{\texttt{\mdseries\slshape parameters}}.quick = false} then the scalar products of \texttt{minus} with \mbox{\texttt{\mdseries\slshape subchars}} are checked; if not all scalar products are nonnegative integers then an empty
list is returned, otherwise $\chi$ is deleted from the list of characters to inspect. 
\item  Otherwise if \texttt{Indeterminateness( minus )} is smaller than \texttt{\mbox{\texttt{\mdseries\slshape parameters}}.minamb} then $\chi$ is deleted from the list of characters. 
\item  If \texttt{\mbox{\texttt{\mdseries\slshape parameters}}.minamb} $\leq$ \texttt{Indeterminateness( minus )} $\leq$ \texttt{\mbox{\texttt{\mdseries\slshape parameters}}.maxamb} then construct the list of contained class functions \texttt{poss:= \mbox{\texttt{\mdseries\slshape parameters}}.contained(\mbox{\texttt{\mdseries\slshape tbl}}, \mbox{\texttt{\mdseries\slshape subchars}}, minus)} and \texttt{Parametrized( poss )}, and improve the approximation of the power map using \texttt{UpdateMap} (\ref{UpdateMap}). 
\end{itemize}
 

 If this yields no further immediate improvements then we branch. If there is a
character from \mbox{\texttt{\mdseries\slshape chars}} left with less or equal \texttt{\mbox{\texttt{\mdseries\slshape parameters}}.maxlen} possible symmetrizations, compute the union of power maps allowed by these
possibilities. Otherwise we choose a class $C$ such that the possible symmetrizations of a character in \mbox{\texttt{\mdseries\slshape chars}} differ at $C$, and compute recursively the union of all allowed power maps with image at $C$ fixed in the set given by the current approximation of the power map. 

 
\begin{Verbatim}[commandchars=!@|,fontsize=\small,frame=single,label=Example]
  !gapprompt@gap>| !gapinput@tbl:= CharacterTable( "U4(3).4" );;|
  !gapprompt@gap>| !gapinput@pow:= InitPowerMap( tbl, 2 );;|
  !gapprompt@gap>| !gapinput@Congruences( tbl, Irr( tbl ), pow, 2 );;  pow;|
  [ 1, 1, 3, 4, 5, 2, 2, 8, 3, 4, 11, 12, [ 6, 7 ], 14, 9, 1, 1, 2, 2,
    3, 4, 5, [ 6, 7 ], 8, 9, 9, 10, 11, 12, 16, 16, 16, 16, 17, 17, 18,
    18, [ 18, 19 ], [ 18, 19 ], 20, 20, 20, 20, 22, 22, 24, 24,
    [ 25, 26 ], [ 25, 26 ], 28, 28, 29, 29 ]
  !gapprompt@gap>| !gapinput@PowerMapsAllowedBySymmetrizations( tbl, Irr( tbl ), Irr( tbl ),|
  !gapprompt@>| !gapinput@      pow, 2, rec( maxlen:= 10, contained:= ContainedPossibleCharacters,|
  !gapprompt@>| !gapinput@      minamb:= 2, maxamb:= infinity, quick:= false ) );|
  [ [ 1, 1, 3, 4, 5, 2, 2, 8, 3, 4, 11, 12, 6, 14, 9, 1, 1, 2, 2, 3, 4,
        5, 6, 8, 9, 9, 10, 11, 12, 16, 16, 16, 16, 17, 17, 18, 18, 18,
        18, 20, 20, 20, 20, 22, 22, 24, 24, 25, 26, 28, 28, 29, 29 ] ]
\end{Verbatim}
 }

 }

  
\section{\textcolor{Chapter }{Subroutines for the Construction of Class Fusions}}\label{Subroutines for the Construction of Class Fusions}
\logpage{[ 73, 7, 0 ]}
\hyperdef{L}{X7AF7305D80E1E5EF}{}
{
  

\subsection{\textcolor{Chapter }{InitFusion}}
\logpage{[ 73, 7, 1 ]}\nobreak
\hyperdef{L}{X7E2BC50C86A16604}{}
{\noindent\textcolor{FuncColor}{$\triangleright$\enspace\texttt{InitFusion({\mdseries\slshape subtbl, tbl})\index{InitFusion@\texttt{InitFusion}}
\label{InitFusion}
}\hfill{\scriptsize (function)}}\\


 For two ordinary character tables \mbox{\texttt{\mdseries\slshape subtbl}} and \mbox{\texttt{\mdseries\slshape tbl}}, \texttt{InitFusion} returns a parametrized map that is a first approximation of the class fusion
from \mbox{\texttt{\mdseries\slshape subtbl}} to \mbox{\texttt{\mdseries\slshape tbl}}, using condition{\nobreakspace}1.{\nobreakspace}listed in the description of \texttt{PossibleClassFusions} (\ref{PossibleClassFusions}). 

 If there are classes for which no images are possible, according to this
criterion, then \texttt{fail} is returned. 

 
\begin{Verbatim}[commandchars=!@|,fontsize=\small,frame=single,label=Example]
  !gapprompt@gap>| !gapinput@subtbl:= CharacterTable( "2F4(2)" );;  tbl:= CharacterTable( "Ru" );;|
  !gapprompt@gap>| !gapinput@fus:= InitFusion( subtbl, tbl );|
  [ 1, 2, 2, 4, [ 5, 6 ], [ 5, 6, 7, 8 ], [ 5, 6, 7, 8 ], [ 9, 10 ],
    11, 14, 14, [ 13, 14, 15 ], [ 16, 17 ], [ 18, 19 ], 20, [ 25, 26 ],
    [ 25, 26 ], [ 5, 6 ], [ 5, 6 ], [ 5, 6 ], [ 5, 6, 7, 8 ],
    [ 13, 14, 15 ], [ 13, 14, 15 ], [ 18, 19 ], [ 18, 19 ], [ 25, 26 ],
    [ 25, 26 ], [ 27, 28, 29 ], [ 27, 28, 29 ] ]
\end{Verbatim}
 }

 

\subsection{\textcolor{Chapter }{CheckPermChar}}
\logpage{[ 73, 7, 2 ]}\nobreak
\hyperdef{L}{X82F776A3850C6404}{}
{\noindent\textcolor{FuncColor}{$\triangleright$\enspace\texttt{CheckPermChar({\mdseries\slshape subtbl, tbl, approxmap, permchar})\index{CheckPermChar@\texttt{CheckPermChar}}
\label{CheckPermChar}
}\hfill{\scriptsize (function)}}\\


 \index{permutation character} \texttt{CheckPermChar} replaces the entries of the parametrized map \mbox{\texttt{\mdseries\slshape approxmap}} by improved values, according to
condition{\nobreakspace}3.{\nobreakspace}listed in the description of \texttt{PossibleClassFusions} (\ref{PossibleClassFusions}). 

 \texttt{CheckPermChar} returns \texttt{true} if no inconsistency occurred, and \texttt{false} otherwise. 

 
\begin{Verbatim}[commandchars=!@|,fontsize=\small,frame=single,label=Example]
  !gapprompt@gap>| !gapinput@permchar:= Sum( Irr( tbl ){ [ 1, 5, 6 ] } );;|
  !gapprompt@gap>| !gapinput@CheckPermChar( subtbl, tbl, fus, permchar ); fus;|
  true
  [ 1, 2, 2, 4, 5, 7, 8, 9, 11, 14, 14, [ 13, 15 ], 16, [ 18, 19 ], 20,
    [ 25, 26 ], [ 25, 26 ], 5, 5, 6, 8, 14, [ 13, 15 ], [ 18, 19 ],
    [ 18, 19 ], [ 25, 26 ], [ 25, 26 ], 27, 27 ]
\end{Verbatim}
 }

 

\subsection{\textcolor{Chapter }{ConsiderTableAutomorphisms}}
\logpage{[ 73, 7, 3 ]}\nobreak
\hyperdef{L}{X7C52CEDB7D98A6B8}{}
{\noindent\textcolor{FuncColor}{$\triangleright$\enspace\texttt{ConsiderTableAutomorphisms({\mdseries\slshape approxmap, grp})\index{ConsiderTableAutomorphisms@\texttt{ConsiderTableAutomorphisms}}
\label{ConsiderTableAutomorphisms}
}\hfill{\scriptsize (function)}}\\


 \index{table automorphisms} \texttt{ConsiderTableAutomorphisms} replaces the entries of the parametrized map \mbox{\texttt{\mdseries\slshape approxmap}} by improved values, according to
condition{\nobreakspace}4.{\nobreakspace}listed in the description of \texttt{PossibleClassFusions} (\ref{PossibleClassFusions}). 

 Afterwards exactly one representative of fusion maps (contained in \mbox{\texttt{\mdseries\slshape approxmap}}) in each orbit under the action of the permutation group \mbox{\texttt{\mdseries\slshape grp}} is contained in the modified parametrized map. 

 \texttt{ConsiderTableAutomorphisms} returns the list of positions where \mbox{\texttt{\mdseries\slshape approxmap}} was changed. 

 
\begin{Verbatim}[commandchars=!@|,fontsize=\small,frame=single,label=Example]
  !gapprompt@gap>| !gapinput@ConsiderTableAutomorphisms( fus, AutomorphismsOfTable( tbl ) );|
  [ 16 ]
  !gapprompt@gap>| !gapinput@fus;|
  [ 1, 2, 2, 4, 5, 7, 8, 9, 11, 14, 14, [ 13, 15 ], 16, [ 18, 19 ], 20,
    25, [ 25, 26 ], 5, 5, 6, 8, 14, [ 13, 15 ], [ 18, 19 ], [ 18, 19 ],
    [ 25, 26 ], [ 25, 26 ], 27, 27 ]
\end{Verbatim}
 }

 

\subsection{\textcolor{Chapter }{FusionsAllowedByRestrictions}}
\logpage{[ 73, 7, 4 ]}\nobreak
\hyperdef{L}{X85024BAE8585DB1C}{}
{\noindent\textcolor{FuncColor}{$\triangleright$\enspace\texttt{FusionsAllowedByRestrictions({\mdseries\slshape subtbl, tbl, subchars, chars, approxmap, parameters})\index{FusionsAllowedByRestrictions@\texttt{FusionsAllowedByRestrictions}}
\label{FusionsAllowedByRestrictions}
}\hfill{\scriptsize (function)}}\\


 Let \mbox{\texttt{\mdseries\slshape subtbl}} and \mbox{\texttt{\mdseries\slshape tbl}} be ordinary character tables, \mbox{\texttt{\mdseries\slshape subchars}} and \mbox{\texttt{\mdseries\slshape chars}} two lists of (values lists of) characters of \mbox{\texttt{\mdseries\slshape subtbl}} and \mbox{\texttt{\mdseries\slshape tbl}}, respectively, \mbox{\texttt{\mdseries\slshape approxmap}} a parametrized map that is an approximation of the class fusion of \mbox{\texttt{\mdseries\slshape subtbl}} in \mbox{\texttt{\mdseries\slshape tbl}}, and \mbox{\texttt{\mdseries\slshape parameters}} a record with the mandatory components \texttt{maxlen}, \texttt{minamb}, \texttt{maxamb} (three integers), \texttt{quick} (a Boolean), and \texttt{contained} (a function, usual values are \texttt{ContainedCharacters} (\ref{ContainedCharacters}) or \texttt{ContainedPossibleCharacters} (\ref{ContainedPossibleCharacters})); optional components of the \mbox{\texttt{\mdseries\slshape parameters}} record are \texttt{testdec} (the function that tests the decomposability, the default is \texttt{NonnegIntScalarProducts} (\ref{NonnegIntScalarProducts})), \texttt{powermaps} (the power paps of \mbox{\texttt{\mdseries\slshape subtbl}} that shall be used for compatibility checks, the default is the \texttt{ComputedPowerMaps} (\ref{ComputedPowerMaps}) value), \texttt{subpowermaps} (the power paps of \mbox{\texttt{\mdseries\slshape tbl}} that shall be used for compatibility checks, the default is the \texttt{ComputedPowerMaps} (\ref{ComputedPowerMaps}) value). 

 \texttt{FusionsAllowedByRestrictions} replaces the entries of \mbox{\texttt{\mdseries\slshape approxmap}} by improved values, according to condition 5.{\nobreakspace}listed in the
description of \texttt{PossibleClassFusions} (\ref{PossibleClassFusions}). 

 More precisely, the strategy used is as follows. 

 First, for each $\chi \in \mbox{\texttt{\mdseries\slshape chars}}$, let \texttt{restricted:= CompositionMaps( }$\chi$\texttt{, \mbox{\texttt{\mdseries\slshape approxmap}} )}. 
\begin{itemize}
\item  If \texttt{Indeterminateness( restricted )}$ = 1$ and \texttt{\mbox{\texttt{\mdseries\slshape parameters}}.quick = false} then the scalar products of \texttt{restricted} with \mbox{\texttt{\mdseries\slshape subchars}} are checked; if not all scalar products are nonnegative integers then an empty
list is returned, otherwise $\chi$ is deleted from the list of characters to inspect. 
\item  Otherwise if \texttt{Indeterminateness( minus )} is smaller than \texttt{\mbox{\texttt{\mdseries\slshape parameters}}.minamb} then $\chi$ is deleted from the list of characters. 
\item  If \texttt{\mbox{\texttt{\mdseries\slshape parameters}}.minamb} $\leq$ \texttt{Indeterminateness( restricted )} $\leq$ \texttt{\mbox{\texttt{\mdseries\slshape parameters}}.maxamb} then construct \texttt{poss:= \mbox{\texttt{\mdseries\slshape parameters}}.contained( \mbox{\texttt{\mdseries\slshape subtbl}}, \mbox{\texttt{\mdseries\slshape subchars}}, restricted )} and \texttt{Parametrized( poss )}, and improve the approximation of the fusion map using \texttt{UpdateMap} (\ref{UpdateMap}). 
\end{itemize}
      

 If this yields no further immediate improvements then we branch. If there is a
character from \mbox{\texttt{\mdseries\slshape chars}} left with less or equal \mbox{\texttt{\mdseries\slshape parameters}}\texttt{.maxlen} possible restrictions, compute the union of fusion maps allowed by these
possibilities. Otherwise we choose a class $C$ such that the possible restrictions of a character in \mbox{\texttt{\mdseries\slshape chars}} differ at $C$, and compute recursively the union of all allowed fusion maps with image at $C$ fixed in the set given by the current approximation of the fusion map. 

 
\begin{Verbatim}[commandchars=!@|,fontsize=\small,frame=single,label=Example]
  !gapprompt@gap>| !gapinput@subtbl:= CharacterTable( "U3(3)" );;  tbl:= CharacterTable( "J4" );;|
  !gapprompt@gap>| !gapinput@fus:= InitFusion( subtbl, tbl );;|
  !gapprompt@gap>| !gapinput@TestConsistencyMaps( ComputedPowerMaps( subtbl ), fus,|
  !gapprompt@>| !gapinput@       ComputedPowerMaps( tbl ) );|
  true
  !gapprompt@gap>| !gapinput@fus;|
  [ 1, 2, 4, 4, [ 5, 6 ], [ 5, 6 ], [ 5, 6 ], 10, [ 12, 13 ],
    [ 12, 13 ], [ 14, 15, 16 ], [ 14, 15, 16 ], [ 21, 22 ], [ 21, 22 ] ]
  !gapprompt@gap>| !gapinput@ConsiderTableAutomorphisms( fus, AutomorphismsOfTable( tbl ) );|
  [ 9 ]
  !gapprompt@gap>| !gapinput@fus;|
  [ 1, 2, 4, 4, [ 5, 6 ], [ 5, 6 ], [ 5, 6 ], 10, 12, [ 12, 13 ],
    [ 14, 15, 16 ], [ 14, 15, 16 ], [ 21, 22 ], [ 21, 22 ] ]
  !gapprompt@gap>| !gapinput@FusionsAllowedByRestrictions( subtbl, tbl, Irr( subtbl ),|
  !gapprompt@>| !gapinput@       Irr( tbl ), fus, rec( maxlen:= 10,|
  !gapprompt@>| !gapinput@       contained:= ContainedPossibleCharacters, minamb:= 2,|
  !gapprompt@>| !gapinput@       maxamb:= infinity, quick:= false ) );|
  [ [ 1, 2, 4, 4, 5, 5, 6, 10, 12, 13, 14, 14, 21, 21 ],
    [ 1, 2, 4, 4, 6, 6, 6, 10, 12, 13, 15, 15, 22, 22 ],
    [ 1, 2, 4, 4, 6, 6, 6, 10, 12, 13, 16, 16, 22, 22 ] ]
\end{Verbatim}
 }

 }

 }

            
\chapter{\textcolor{Chapter }{Unknowns}}\label{Unknowns}
\logpage{[ 74, 0, 0 ]}
\hyperdef{L}{X7C1FAB6280A02CCB}{}
{
  \index{data type!unknown} Sometimes the result of an operation does not allow further computations with
it. In many cases, then an error is signalled, and the computation is stopped. 

 This is not appropriate for some applications in character theory. For
example, if one wants to induce a character of a group to a supergroup
(see{\nobreakspace}\texttt{InducedClassFunction} (\ref{InducedClassFunction:for the character table of a supergroup})) but the class fusion is only a parametrized map (see Chapter{\nobreakspace}\ref{Maps Concerning Character Tables}), there may be values of the induced character which are determined by the
fusion map, whereas other values are not known. 

 For this and other situations, \textsf{GAP} provides the data type \emph{unknown}. An object of this type, further on called an \emph{unknown}, may stand for any cyclotomic (see Chapter{\nobreakspace}\ref{Cyclotomic Numbers}), in particular its family (see{\nobreakspace}\ref{Families}) is \texttt{CyclotomicsFamily}. 

 Unknowns are parametrized by positive integers. When a \textsf{GAP} session is started, no unknowns exist. 

 The only ways to create unknowns are to call the function \texttt{Unknown} (\ref{Unknown}) or a function that calls it, or to do arithmetical operations with unknowns. 

 \textsf{GAP} objects containing unknowns will contain \emph{fixed} unknowns when they are printed to files, i.e., function calls \texttt{Unknown(}$n$\texttt{)} instead of \texttt{Unknown()}. So be careful to read files printed in different \textsf{GAP} sessions, since there may be the same unknown at different places. 

 The rest of this chapter contains information about the unknown constructor,
the category, and comparison of and arithmetical operations for unknowns. More
is not known about unknowns in \textsf{GAP}.  
\section{\textcolor{Chapter }{More about Unknowns}}\label{sect:Unknown}
\logpage{[ 74, 1, 0 ]}
\hyperdef{L}{X85A1A27686C8D366}{}
{
  

\subsection{\textcolor{Chapter }{Unknown}}
\logpage{[ 74, 1, 1 ]}\nobreak
\hyperdef{L}{X79BAB8C48394779C}{}
{\noindent\textcolor{FuncColor}{$\triangleright$\enspace\texttt{Unknown({\mdseries\slshape [n]})\index{Unknown@\texttt{Unknown}}
\label{Unknown}
}\hfill{\scriptsize (operation)}}\\


 Called without argument, \texttt{Unknown} returns a new unknown value, i.e., the first one that is larger than all
unknowns which exist in the current \textsf{GAP} session. 

 Called with a positive integer \mbox{\texttt{\mdseries\slshape n}}, \texttt{Unknown} returns the \mbox{\texttt{\mdseries\slshape n}}\texttt{\symbol{45}}th unknown; if this did not exist yet, it is created. }

 

\subsection{\textcolor{Chapter }{LargestUnknown}}
\logpage{[ 74, 1, 2 ]}\nobreak
\hyperdef{L}{X7B38F63581D7A96A}{}
{\noindent\textcolor{FuncColor}{$\triangleright$\enspace\texttt{LargestUnknown\index{LargestUnknown@\texttt{LargestUnknown}}
\label{LargestUnknown}
}\hfill{\scriptsize (global variable)}}\\


 \texttt{LargestUnknown} is the largest \mbox{\texttt{\mdseries\slshape n}} that is used in any \texttt{Unknown( \mbox{\texttt{\mdseries\slshape n}} )} in the current \textsf{GAP} session. This is used in \texttt{Unknown} (\ref{Unknown}) which increments this value when asked to make a new unknown. }

 

\subsection{\textcolor{Chapter }{IsUnknown}}
\logpage{[ 74, 1, 3 ]}\nobreak
\hyperdef{L}{X828556067E069B6D}{}
{\noindent\textcolor{FuncColor}{$\triangleright$\enspace\texttt{IsUnknown({\mdseries\slshape obj})\index{IsUnknown@\texttt{IsUnknown}}
\label{IsUnknown}
}\hfill{\scriptsize (Category)}}\\


 is the category of unknowns in \textsf{GAP}. 
\begin{Verbatim}[commandchars=!@|,fontsize=\small,frame=single,label=Example]
  !gapprompt@gap>| !gapinput@Unknown();  List( [ 1 .. 20 ], i -> Unknown() );;|
  Unknown(1)
  !gapprompt@gap>| !gapinput@Unknown();   # note that we have already created 21 unknowns.|
  Unknown(22)
  !gapprompt@gap>| !gapinput@Unknown(2000);  Unknown();|
  Unknown(2000)
  Unknown(2001)
  !gapprompt@gap>| !gapinput@LargestUnknown;|
  2001
  !gapprompt@gap>| !gapinput@IsUnknown( Unknown );  IsUnknown( Unknown() );|
  false
  true
\end{Verbatim}
 }

 
\subsection{\textcolor{Chapter }{Comparison of Unknowns}}\logpage{[ 74, 1, 4 ]}
\hyperdef{L}{X7E6B0D62788BB464}{}
{
  Unknowns can be \emph{compared} via \texttt{=} and \texttt{{\textless}} with all cyclotomics and with certain other \textsf{GAP} objects (see{\nobreakspace}\ref{Comparisons}). We have \texttt{Unknown( \mbox{\texttt{\mdseries\slshape n}} ) {\textgreater}= Unknown( \mbox{\texttt{\mdseries\slshape m}} )} if and only if \texttt{\mbox{\texttt{\mdseries\slshape n}} {\textgreater}= \mbox{\texttt{\mdseries\slshape m}}} holds, unknowns are larger than all cyclotomics that are not unknowns. 

 
\begin{Verbatim}[commandchars=!@|,fontsize=\small,frame=single,label=Example]
  !gapprompt@gap>| !gapinput@Unknown() >= Unknown();  Unknown(2) < Unknown(3);|
  false
  true
  !gapprompt@gap>| !gapinput@Unknown() > 3;  Unknown() > E(3);|
  true
  true
  !gapprompt@gap>| !gapinput@Unknown() > Z(8);  Unknown() > [];|
  false
  false
\end{Verbatim}
 }

 
\subsection{\textcolor{Chapter }{Arithmetical Operations for Unknowns}}\logpage{[ 74, 1, 5 ]}
\hyperdef{L}{X81EFCA7C82E18EFF}{}
{
  The usual arithmetic operations \texttt{+}, \texttt{\texttt{\symbol{45}}}, \texttt{*} and \texttt{/} are defined for addition, subtraction, multiplication and division of unknowns
and cyclotomics. The result will be a new unknown except in one of the
following cases. 

 Multiplication with zero yields zero, and multiplication with one or addition
of zero yields the old unknown. \emph{Note} that division by an unknown causes an error, since an unknown might stand for
zero. 

 As unknowns are cyclotomics, dense lists of unknowns and other cyclotomics are
row vectors and they can be added and multiplied in the usual way.
Consequently, lists of such row vectors of equal length are (ordinary)
matrices (see{\nobreakspace}\texttt{IsOrdinaryMatrix} (\ref{IsOrdinaryMatrix})). }

 }

 }

             
\chapter{\textcolor{Chapter }{Monomiality Questions}}\label{Monomiality Questions}
\logpage{[ 75, 0, 0 ]}
\hyperdef{L}{X80D9CA647E680B19}{}
{
  This chapter describes functions dealing with the monomiality of finite
(solvable) groups and their characters. 

 All these functions assume \emph{characters} to be class function objects as described in Chapter{\nobreakspace}\ref{Class Functions}, lists of character \emph{values} are not allowed. 

 The usual \emph{property tests} of \textsf{GAP} that return either \texttt{true} or \texttt{false} are not sufficient for us. When we ask whether a group character $\chi$ has a certain property, such as quasiprimitivity, we usually want more
information than just yes or no. Often we are interested in the reason \emph{why} a group character $\chi$ was proved to have a certain property, e.g., whether monomiality of $\chi$ was proved by the observation that the underlying group is nilpotent, or
whether it was necessary to construct a linear character of a subgroup from
which $\chi$ can be induced. In the latter case we also may be interested in this linear
character. Therefore we need test functions that return a record containing
such useful information. For example, the record returned by the function \texttt{TestQuasiPrimitive} (\ref{TestQuasiPrimitive}) contains the component \texttt{isQuasiPrimitive} (which is the known boolean property flag), and additionally the component \texttt{comment}, a string telling the reason for the value of the \texttt{isQuasiPrimitive} component, and in the case that the argument $\chi$ was \emph{not} quasiprimitive also the component \texttt{character}, which is an irreducible constituent of a nonhomogeneous restriction of $\chi$ to a normal subgroup. Besides these test functions there are also the known
properties, e.g., the property \texttt{IsQuasiPrimitive} (\ref{IsQuasiPrimitive}) which will call the attribute \texttt{TestQuasiPrimitive} (\ref{TestQuasiPrimitive}), and return the value of the \texttt{isQuasiPrimitive} component of the result. 

 A few words about how to use the monomiality functions seem to be necessary.
Monomiality questions usually involve computations in many subgroups and
factor groups of a given group, and for these groups often expensive
calculations such as that of the character table are necessary. So one should
be careful not to construct the same group over and over again, instead the
same group object should be reused, such that its character table need to be
computed only once. For example, suppose you want to restrict a character to a
normal subgroup $N$ that was constructed as a normal closure of some group elements, and suppose
that you have already computed with normal subgroups (by calls to \texttt{NormalSubgroups} (\ref{NormalSubgroups}) or \texttt{MaximalNormalSubgroups} (\ref{MaximalNormalSubgroups})) and their character tables. Then you should look in the lists of known
normal subgroups whether $N$ is contained, and if so you can use the known character table. A mechanism
that supports this for normal subgroups is described in \ref{Storing Normal Subgroup Information}. 

 Also the following hint may be useful in this context. If you know that sooner
or later you will compute the character table of a group $G$ then it may be advisable to compute it as soon as possible. For example, if
you need the normal subgroups of $G$ then they can be computed more efficiently if the character table of $G$ is known, and they can be stored compatibly to the contained $G$\texttt{\symbol{45}}conjugacy classes. This correspondence of classes list and
normal subgroup can be used very often. 

 Several \emph{examples} in this chapter use the symmetric group $S_4$ and the special linear group $SL(2,3)$. For running the examples, you must first define the groups, for example as
follows. 

 
\begin{Verbatim}[commandchars=!@|,fontsize=\small,frame=single,label=Example]
  !gapprompt@gap>| !gapinput@S4:= SymmetricGroup( 4 );;  SetName( S4, "S4" );|
  !gapprompt@gap>| !gapinput@Sl23:= SL( 2, 3 );;|
\end{Verbatim}
  
\section{\textcolor{Chapter }{InfoMonomial (Info Class)}}\label{sect:InfoMonomial}
\logpage{[ 75, 1, 0 ]}
\hyperdef{L}{X7F2F753F7B354F09}{}
{
  

\subsection{\textcolor{Chapter }{InfoMonomial}}
\logpage{[ 75, 1, 1 ]}\nobreak
\hyperdef{L}{X8103DD607C7F2CD2}{}
{\noindent\textcolor{FuncColor}{$\triangleright$\enspace\texttt{InfoMonomial\index{InfoMonomial@\texttt{InfoMonomial}}
\label{InfoMonomial}
}\hfill{\scriptsize (info class)}}\\


 Most of the functions described in this chapter print some (hopefully useful) \emph{information} if the info level of the info class \texttt{InfoMonomial} is at least $1$, see{\nobreakspace}\ref{Info Functions} for details. }

 }

  
\section{\textcolor{Chapter }{Character Degrees and Derived Length}}\label{Character Degrees and Derived Length}
\logpage{[ 75, 2, 0 ]}
\hyperdef{L}{X842D10BC7CA9C2DE}{}
{
  

\subsection{\textcolor{Chapter }{Alpha}}
\logpage{[ 75, 2, 1 ]}\nobreak
\hyperdef{L}{X86A900897819E5AC}{}
{\noindent\textcolor{FuncColor}{$\triangleright$\enspace\texttt{Alpha({\mdseries\slshape G})\index{Alpha@\texttt{Alpha}}
\label{Alpha}
}\hfill{\scriptsize (attribute)}}\\


 For a group \mbox{\texttt{\mdseries\slshape G}}, \texttt{Alpha} returns a list whose $i$\texttt{\symbol{45}}th entry is the maximal derived length of groups $\mbox{\texttt{\mdseries\slshape G}} / \ker(\chi)$ for $\chi \in Irr(\mbox{\texttt{\mdseries\slshape G}})$ with $\chi(1)$ at most the $i$\texttt{\symbol{45}}th irreducible degree of \mbox{\texttt{\mdseries\slshape G}}. }

 

\subsection{\textcolor{Chapter }{Delta}}
\logpage{[ 75, 2, 2 ]}\nobreak
\hyperdef{L}{X82C33CF282FC5A73}{}
{\noindent\textcolor{FuncColor}{$\triangleright$\enspace\texttt{Delta({\mdseries\slshape G})\index{Delta@\texttt{Delta}}
\label{Delta}
}\hfill{\scriptsize (attribute)}}\\


 For a group \mbox{\texttt{\mdseries\slshape G}}, \texttt{Delta} returns the list $[ 1, alp[2] - alp[1], \ldots, alp[\mbox{\texttt{\mdseries\slshape n}}] - alp[\mbox{\texttt{\mdseries\slshape n}}-1] ]$, where $alp = $\texttt{Alpha( \mbox{\texttt{\mdseries\slshape G}} )} (see{\nobreakspace}\texttt{Alpha} (\ref{Alpha})). }

 
\subsection{\textcolor{Chapter }{IsBergerCondition}}\logpage{[ 75, 2, 3 ]}
\hyperdef{L}{X7D0D26267A9D37DD}{}
{
\noindent\textcolor{FuncColor}{$\triangleright$\enspace\texttt{IsBergerCondition({\mdseries\slshape G})\index{IsBergerCondition@\texttt{IsBergerCondition}!for a group}
\label{IsBergerCondition:for a group}
}\hfill{\scriptsize (property)}}\\
\noindent\textcolor{FuncColor}{$\triangleright$\enspace\texttt{IsBergerCondition({\mdseries\slshape chi})\index{IsBergerCondition@\texttt{IsBergerCondition}!for a character}
\label{IsBergerCondition:for a character}
}\hfill{\scriptsize (property)}}\\


 Called with an irreducible character \mbox{\texttt{\mdseries\slshape chi}} of a group $G$, \texttt{IsBergerCondition} returns \texttt{true} if \mbox{\texttt{\mdseries\slshape chi}} satisfies $M' \leq \ker(\chi)$ for every normal subgroup $M$ of $G$ with the property that $M \leq \ker(\psi)$ holds for all $\psi \in Irr(G)$ with $\psi(1) < \chi(1)$, and \texttt{false} otherwise. 

 Called with a group \mbox{\texttt{\mdseries\slshape G}}, \texttt{IsBergerCondition} returns \texttt{true} if all irreducible characters of \mbox{\texttt{\mdseries\slshape G}} satisfy the inequality above, and \texttt{false} otherwise. 

 For groups of odd order the result is always \texttt{true} by a theorem of T.{\nobreakspace}R.{\nobreakspace}Berger (see{\nobreakspace}\cite[Thm. 2.2]{Ber76}). 

 In the case that \texttt{false} is returned, \texttt{InfoMonomial} (\ref{InfoMonomial}) tells about a degree for which the inequality is violated. 

 
\begin{Verbatim}[commandchars=!@|,fontsize=\small,frame=single,label=Example]
  !gapprompt@gap>| !gapinput@Alpha( Sl23 );|
  [ 1, 3, 3 ]
  !gapprompt@gap>| !gapinput@Alpha( S4 );|
  [ 1, 2, 3 ]
  !gapprompt@gap>| !gapinput@Delta( Sl23 );|
  [ 1, 2, 0 ]
  !gapprompt@gap>| !gapinput@Delta( S4 );|
  [ 1, 1, 1 ]
  !gapprompt@gap>| !gapinput@IsBergerCondition( S4 );|
  true
  !gapprompt@gap>| !gapinput@IsBergerCondition( Sl23 );|
  false
  !gapprompt@gap>| !gapinput@List( Irr( Sl23 ), IsBergerCondition );|
  [ true, true, true, false, false, false, true ]
  !gapprompt@gap>| !gapinput@List( Irr( Sl23 ), Degree );|
  [ 1, 1, 1, 2, 2, 2, 3 ]
\end{Verbatim}
 }

 }

  
\section{\textcolor{Chapter }{Primitivity of Characters}}\label{Primitivity of Characters}
\logpage{[ 75, 3, 0 ]}
\hyperdef{L}{X82C21037806B52CE}{}
{
  

\subsection{\textcolor{Chapter }{TestHomogeneous}}
\logpage{[ 75, 3, 1 ]}\nobreak
\hyperdef{L}{X81FD26947924C500}{}
{\noindent\textcolor{FuncColor}{$\triangleright$\enspace\texttt{TestHomogeneous({\mdseries\slshape chi, N})\index{TestHomogeneous@\texttt{TestHomogeneous}}
\label{TestHomogeneous}
}\hfill{\scriptsize (function)}}\\


 For a group character \mbox{\texttt{\mdseries\slshape chi}} of a group $G$ and a normal subgroup \mbox{\texttt{\mdseries\slshape N}} of $G$, \texttt{TestHomogeneous} returns a record with information whether the restriction of \mbox{\texttt{\mdseries\slshape chi}} to \mbox{\texttt{\mdseries\slshape N}} is homogeneous, i.e., is a multiple of an irreducible character. 

 \mbox{\texttt{\mdseries\slshape N}} may be given also as list of conjugacy class positions
w.r.t.{\nobreakspace}the character table of $G$. 

 The components of the result are 

 
\begin{description}
\item[{\texttt{isHomogeneous}}]  \texttt{true} or \texttt{false}, 
\item[{\texttt{comment}}]  a string telling a reason for the value of the \texttt{isHomogeneous} component, 
\item[{\texttt{character}}]  irreducible constituent of the restriction, only bound if the restriction had
to be checked, 
\item[{\texttt{multiplicity}}]  multiplicity of the \texttt{character} component in the restriction of \mbox{\texttt{\mdseries\slshape chi}}. 
\end{description}
 

 
\begin{Verbatim}[commandchars=!@|,fontsize=\small,frame=single,label=Example]
  !gapprompt@gap>| !gapinput@n:= DerivedSubgroup( Sl23 );;|
  !gapprompt@gap>| !gapinput@chi:= Irr( Sl23 )[7];|
  Character( CharacterTable( SL(2,3) ), [ 3, 0, 0, 3, 0, 0, -1 ] )
  !gapprompt@gap>| !gapinput@test:= TestHomogeneous( chi, n );;|
  !gapprompt@gap>| !gapinput@test.isHomogeneous;  test.comment;  test.multiplicity;|
  false
  "restriction checked"
  1
  !gapprompt@gap>| !gapinput@Degree( test.character );|
  1
  !gapprompt@gap>| !gapinput@chi:= Irr( Sl23 )[4];|
  Character( CharacterTable( SL(2,3) ), [ 2, 1, 1, -2, -1, -1, 0 ] )
  !gapprompt@gap>| !gapinput@cln:= ClassPositionsOfNormalSubgroup( CharacterTable( Sl23 ), n );|
  [ 1, 4, 7 ]
  !gapprompt@gap>| !gapinput@TestHomogeneous( chi, cln );|
  rec( comment := "restricts irreducibly", isHomogeneous := true )
\end{Verbatim}
 }

 

\subsection{\textcolor{Chapter }{IsPrimitiveCharacter}}
\logpage{[ 75, 3, 2 ]}\nobreak
\hyperdef{L}{X7BC72ECE822D4245}{}
{\noindent\textcolor{FuncColor}{$\triangleright$\enspace\texttt{IsPrimitiveCharacter({\mdseries\slshape chi})\index{IsPrimitiveCharacter@\texttt{IsPrimitiveCharacter}}
\label{IsPrimitiveCharacter}
}\hfill{\scriptsize (property)}}\\


 For a character \mbox{\texttt{\mdseries\slshape chi}} of a group $G$, \texttt{IsPrimitiveCharacter} returns \texttt{true} if \mbox{\texttt{\mdseries\slshape chi}} is not induced from any proper subgroup, and \texttt{false} otherwise. This currently only works for characters of soluble groups. 

 
\begin{Verbatim}[commandchars=!@|,fontsize=\small,frame=single,label=Example]
  !gapprompt@gap>| !gapinput@IsPrimitiveCharacter( Irr( Sl23 )[4] );|
  true
  !gapprompt@gap>| !gapinput@IsPrimitiveCharacter( Irr( Sl23 )[7] );|
  false
\end{Verbatim}
 }

 

\subsection{\textcolor{Chapter }{TestQuasiPrimitive}}
\logpage{[ 75, 3, 3 ]}\nobreak
\hyperdef{L}{X82BFA6968415F308}{}
{\noindent\textcolor{FuncColor}{$\triangleright$\enspace\texttt{TestQuasiPrimitive({\mdseries\slshape chi})\index{TestQuasiPrimitive@\texttt{TestQuasiPrimitive}}
\label{TestQuasiPrimitive}
}\hfill{\scriptsize (attribute)}}\\
\noindent\textcolor{FuncColor}{$\triangleright$\enspace\texttt{IsQuasiPrimitive({\mdseries\slshape chi})\index{IsQuasiPrimitive@\texttt{IsQuasiPrimitive}}
\label{IsQuasiPrimitive}
}\hfill{\scriptsize (property)}}\\


 \texttt{TestQuasiPrimitive} returns a record with information about quasiprimitivity of the group
character \mbox{\texttt{\mdseries\slshape chi}}, i.e., whether \mbox{\texttt{\mdseries\slshape chi}} restricts homogeneously to every normal subgroup of its group. The result
record contains at least the components \texttt{isQuasiPrimitive} (with value either \texttt{true} or \texttt{false}) and \texttt{comment} (a string telling a reason for the value of the component \texttt{isQuasiPrimitive}). If \mbox{\texttt{\mdseries\slshape chi}} is not quasiprimitive then there is additionally a component \texttt{character}, with value an irreducible constituent of a nonhomogeneous restriction of \mbox{\texttt{\mdseries\slshape chi}}. 

 \texttt{IsQuasiPrimitive} returns \texttt{true} or \texttt{false}, depending on whether the character \mbox{\texttt{\mdseries\slshape chi}} is quasiprimitive. 

 Note that for solvable groups, quasiprimitivity is the same as primitivity
(see{\nobreakspace}\texttt{IsPrimitiveCharacter} (\ref{IsPrimitiveCharacter})). 

 
\begin{Verbatim}[commandchars=!@|,fontsize=\small,frame=single,label=Example]
  !gapprompt@gap>| !gapinput@chi:= Irr( Sl23 )[4];|
  Character( CharacterTable( SL(2,3) ), [ 2, 1, 1, -2, -1, -1, 0 ] )
  !gapprompt@gap>| !gapinput@TestQuasiPrimitive( chi );|
  rec( comment := "all restrictions checked", isQuasiPrimitive := true )
  !gapprompt@gap>| !gapinput@chi:= Irr( Sl23 )[7];|
  Character( CharacterTable( SL(2,3) ), [ 3, 0, 0, 3, 0, 0, -1 ] )
  !gapprompt@gap>| !gapinput@test:= TestQuasiPrimitive( chi );;|
  !gapprompt@gap>| !gapinput@test.isQuasiPrimitive;  test.comment;|
  false
  "restriction checked"
  !gapprompt@gap>| !gapinput@Degree( test.character );|
  1
\end{Verbatim}
 }

 

\subsection{\textcolor{Chapter }{TestInducedFromNormalSubgroup}}
\logpage{[ 75, 3, 4 ]}\nobreak
\hyperdef{L}{X84860E3A7FECDBA3}{}
{\noindent\textcolor{FuncColor}{$\triangleright$\enspace\texttt{TestInducedFromNormalSubgroup({\mdseries\slshape chi[, N]})\index{TestInducedFromNormalSubgroup@\texttt{TestInducedFromNormalSubgroup}}
\label{TestInducedFromNormalSubgroup}
}\hfill{\scriptsize (function)}}\\
\noindent\textcolor{FuncColor}{$\triangleright$\enspace\texttt{IsInducedFromNormalSubgroup({\mdseries\slshape chi})\index{IsInducedFromNormalSubgroup@\texttt{IsInducedFromNormalSubgroup}}
\label{IsInducedFromNormalSubgroup}
}\hfill{\scriptsize (property)}}\\


 \texttt{TestInducedFromNormalSubgroup} returns a record with information whether the irreducible character \mbox{\texttt{\mdseries\slshape chi}} of the group $G$ is induced from a proper normal subgroup of $G$. If the second argument \mbox{\texttt{\mdseries\slshape N}} is present, which must be a normal subgroup of $G$ or the list of class positions of a normal subgroup of $G$, it is checked whether \mbox{\texttt{\mdseries\slshape chi}} is induced from \mbox{\texttt{\mdseries\slshape N}}. 

 The result contains always the components \texttt{isInduced} (either \texttt{true} or \texttt{false}) and \texttt{comment} (a string telling a reason for the value of the component \texttt{isInduced}). In the \texttt{true} case there is a component \texttt{character} which contains a character of a maximal normal subgroup from which \mbox{\texttt{\mdseries\slshape chi}} is induced. 

 \texttt{IsInducedFromNormalSubgroup} returns \texttt{true} if \mbox{\texttt{\mdseries\slshape chi}} is induced from a proper normal subgroup of $G$, and \texttt{false} otherwise. 

 
\begin{Verbatim}[commandchars=!@|,fontsize=\small,frame=single,label=Example]
  !gapprompt@gap>| !gapinput@List( Irr( Sl23 ), IsInducedFromNormalSubgroup );|
  [ false, false, false, false, false, false, true ]
  !gapprompt@gap>| !gapinput@List( Irr( S4 ){ [ 1, 3, 4 ] },|
  !gapprompt@>| !gapinput@         TestInducedFromNormalSubgroup );|
  [ rec( comment := "linear character", isInduced := false ),
    rec( character := Character( CharacterTable( Alt( [ 1 .. 4 ] ) ),
          [ 1, 1, E(3), E(3)^2 ] ),
        comment := "induced from component '.character'",
        isInduced := true ),
    rec( comment := "all maximal normal subgroups checked",
        isInduced := false ) ]
\end{Verbatim}
 }

 }

  
\section{\textcolor{Chapter }{Testing Monomiality}}\label{Testing Monomiality}
\logpage{[ 75, 4, 0 ]}
\hyperdef{L}{X86567D7F781BBCAE}{}
{
  A character $\chi$ of a finite group $G$ is called \emph{monomial} if $\chi$ is induced from a linear character of a subgroup of $G$. A finite group $G$ is called \emph{monomial} (or \emph{$M$\texttt{\symbol{45}}group}) if each ordinary irreducible character of $G$ is monomial. 

  
\subsection{\textcolor{Chapter }{TestMonomial}}\logpage{[ 75, 4, 1 ]}
\hyperdef{L}{X84EB92B57DAF5C93}{}
{
\noindent\textcolor{FuncColor}{$\triangleright$\enspace\texttt{TestMonomial({\mdseries\slshape chi})\index{TestMonomial@\texttt{TestMonomial}!for a character}
\label{TestMonomial:for a character}
}\hfill{\scriptsize (attribute)}}\\
\noindent\textcolor{FuncColor}{$\triangleright$\enspace\texttt{TestMonomial({\mdseries\slshape G})\index{TestMonomial@\texttt{TestMonomial}!for a group}
\label{TestMonomial:for a group}
}\hfill{\scriptsize (attribute)}}\\
\noindent\textcolor{FuncColor}{$\triangleright$\enspace\texttt{TestMonomial({\mdseries\slshape chi, uselattice})\index{TestMonomial@\texttt{TestMonomial}!for a character and a Boolean}
\label{TestMonomial:for a character and a Boolean}
}\hfill{\scriptsize (operation)}}\\
\noindent\textcolor{FuncColor}{$\triangleright$\enspace\texttt{TestMonomial({\mdseries\slshape G, uselattice})\index{TestMonomial@\texttt{TestMonomial}!for a group and a Boolean}
\label{TestMonomial:for a group and a Boolean}
}\hfill{\scriptsize (operation)}}\\


 Called with a group character \mbox{\texttt{\mdseries\slshape chi}} of a group \mbox{\texttt{\mdseries\slshape G}}, \texttt{TestMonomial} returns a record containing information about monomiality of the group \mbox{\texttt{\mdseries\slshape G}} or the group character \mbox{\texttt{\mdseries\slshape chi}}, respectively. 

 If \texttt{TestMonomial} proves the character \mbox{\texttt{\mdseries\slshape chi}} to be monomial then the result contains components \texttt{isMonomial} (with value \texttt{true}), \texttt{comment} (a string telling a reason for monomiality), and if it was necessary to
compute a linear character from which \mbox{\texttt{\mdseries\slshape chi}} is induced, also a component \texttt{character}. 

 If \texttt{TestMonomial} proves \mbox{\texttt{\mdseries\slshape chi}} or \mbox{\texttt{\mdseries\slshape G}} to be nonmonomial then the value of the component \texttt{isMonomial} is \texttt{false}, and in the case of \mbox{\texttt{\mdseries\slshape G}} a nonmonomial character is the value of the component \texttt{character} if it had been necessary to compute it. 

 A Boolean can be entered as the second argument \mbox{\texttt{\mdseries\slshape uselattice}}; if the value is \texttt{true} then the subgroup lattice of the underlying group is used if necessary, if the
value is \texttt{false} then the subgroup lattice is used only for groups of order at most \texttt{TestMonomialUseLattice} (\ref{TestMonomialUseLattice}). The default value of \mbox{\texttt{\mdseries\slshape uselattice}} is \texttt{false}. 

 For a group whose lattice must not be used, it may happen that \texttt{TestMonomial} cannot prove or disprove monomiality; then the result record contains the
component \texttt{isMonomial} with value \texttt{"?"}. This case occurs in the call for a character \mbox{\texttt{\mdseries\slshape chi}} if and only if \mbox{\texttt{\mdseries\slshape chi}} is not induced from the inertia subgroup of a component of any reducible
restriction to a normal subgroup. It can happen that \mbox{\texttt{\mdseries\slshape chi}} is monomial in this situation. For a group, this case occurs if no irreducible
character can be proved to be nonmonomial, and if no decision is possible for
at least one irreducible character. 

 
\begin{Verbatim}[commandchars=!@|,fontsize=\small,frame=single,label=Example]
  !gapprompt@gap>| !gapinput@TestMonomial( S4 );|
  rec( comment := "abelian by supersolvable group", isMonomial := true )
  !gapprompt@gap>| !gapinput@TestMonomial( Sl23 );|
  rec( comment := "list Delta( G ) contains entry > 1",
    isMonomial := false )
\end{Verbatim}
 }

 

\subsection{\textcolor{Chapter }{TestMonomialUseLattice}}
\logpage{[ 75, 4, 2 ]}\nobreak
\hyperdef{L}{X787CCCBB7FC17F5E}{}
{\noindent\textcolor{FuncColor}{$\triangleright$\enspace\texttt{TestMonomialUseLattice\index{TestMonomialUseLattice@\texttt{TestMonomialUseLattice}}
\label{TestMonomialUseLattice}
}\hfill{\scriptsize (global variable)}}\\


 This global variable controls for which groups the operation \texttt{TestMonomial} (\ref{TestMonomial:for a group and a Boolean}) may compute the subgroup lattice. The value can be set to a positive integer
or \texttt{infinity} (\ref{infinity}), the default is $1000$. }

 

\subsection{\textcolor{Chapter }{IsMonomialNumber}}
\logpage{[ 75, 4, 3 ]}\nobreak
\hyperdef{L}{X8261B5AA7BCFFCC2}{}
{\noindent\textcolor{FuncColor}{$\triangleright$\enspace\texttt{IsMonomialNumber({\mdseries\slshape n})\index{IsMonomialNumber@\texttt{IsMonomialNumber}}
\label{IsMonomialNumber}
}\hfill{\scriptsize (property)}}\\


 For a positive integer \mbox{\texttt{\mdseries\slshape n}}, \texttt{IsMonomialNumber} returns \texttt{true} if every solvable group of order \mbox{\texttt{\mdseries\slshape n}} is monomial, and \texttt{false} otherwise. One can also use \texttt{IsMonomial} instead. \index{IsMonomial@\texttt{IsMonomial}!for positive integers} 

 Let $\nu_p(n)$ denote the multiplicity of the prime $p$ as factor of $n$, and $ord(p,q)$ the multiplicative order of $p \pmod{q}$. 

 Then there exists a solvable nonmonomial group of order $n$ if and only if one of the following conditions is satisfied. 

 
\begin{description}
\item[{1.}]  $\nu_2(n) \geq 2$ and there is a $p$ such that $\nu_p(n) \geq 3$ and $p \equiv -1 \pmod{4}$, 
\item[{2.}]  $\nu_2(n) \geq 3$ and there is a $p$ such that $\nu_p(n) \geq 3$ and $p \equiv 1 \pmod{4}$, 
\item[{3.}]  there are odd prime divisors $p$ and $q$ of $n$ such that $ord(p,q)$ is even and $ord(p,q) < \nu_p(n)$ (especially $\nu_p(n) \geq 3$), 
\item[{4.}]  there is a prime divisor $q$ of $n$ such that $\nu_2(n) \geq 2 ord(2,q) + 2$ (especially $\nu_2(n) \geq 4$), 
\item[{5.}]  $\nu_2(n) \geq 2$ and there is a $p$ such that $p \equiv 1 \pmod{4}$, $ord(p,q)$ is odd, and $2 ord(p,q) < \nu_p(n)$ (especially $\nu_p(n) \geq 3$). 
\end{description}
 

 These five possibilities correspond to the five types of solvable minimal
nonmonomial groups (see{\nobreakspace}\texttt{MinimalNonmonomialGroup} (\ref{MinimalNonmonomialGroup})) that can occur as subgroups and factor groups of groups of order \mbox{\texttt{\mdseries\slshape n}}. 

 
\begin{Verbatim}[commandchars=!@|,fontsize=\small,frame=single,label=Example]
  !gapprompt@gap>| !gapinput@Filtered( [ 1 .. 111 ], x -> not IsMonomial( x ) );|
  [ 24, 48, 72, 96, 108 ]
\end{Verbatim}
 }

 
\subsection{\textcolor{Chapter }{TestMonomialQuick}}\logpage{[ 75, 4, 4 ]}
\hyperdef{L}{X822E03EF7B8F92D3}{}
{
\noindent\textcolor{FuncColor}{$\triangleright$\enspace\texttt{TestMonomialQuick({\mdseries\slshape chi})\index{TestMonomialQuick@\texttt{TestMonomialQuick}!for a character}
\label{TestMonomialQuick:for a character}
}\hfill{\scriptsize (attribute)}}\\
\noindent\textcolor{FuncColor}{$\triangleright$\enspace\texttt{TestMonomialQuick({\mdseries\slshape G})\index{TestMonomialQuick@\texttt{TestMonomialQuick}!for a group}
\label{TestMonomialQuick:for a group}
}\hfill{\scriptsize (attribute)}}\\


 \texttt{TestMonomialQuick} does some cheap tests whether the irreducible character \mbox{\texttt{\mdseries\slshape chi}} or the group \mbox{\texttt{\mdseries\slshape G}}, respectively, is monomial. Here ``cheap'' means in particular that no computations of character tables are involved, and
it is \emph{not} checked whether \mbox{\texttt{\mdseries\slshape chi}} is a character and irreducible. The return value is a record with components 
\begin{description}
\item[{\texttt{isMonomial}}]  either \texttt{true} or \texttt{false} or the string \texttt{"?"}, depending on whether (non)monomiality could be proved, and 
\item[{\texttt{comment}}]  a string telling the reason for the value of the \texttt{isMonomial} component. 
\end{description}
 

 A group \mbox{\texttt{\mdseries\slshape G}} is proved to be monomial by \texttt{TestMonomialQuick} if \mbox{\texttt{\mdseries\slshape G}} is nilpotent or Sylow abelian by supersolvable, or if \mbox{\texttt{\mdseries\slshape G}} is solvable and its order is not divisible by the third power of a prime,
Nonsolvable groups are proved to be nonmonomial by \texttt{TestMonomialQuick}. 

 An irreducible character \mbox{\texttt{\mdseries\slshape chi}} is proved to be monomial if it is linear, or if its codegree is a prime power,
or if its group knows to be monomial, or if the factor group modulo the kernel
can be proved to be monomial by \texttt{TestMonomialQuick}. 

 
\begin{Verbatim}[commandchars=!@|,fontsize=\small,frame=single,label=Example]
  !gapprompt@gap>| !gapinput@TestMonomialQuick( Irr( S4 )[3] );|
  rec( comment := "whole group is monomial", isMonomial := true )
  !gapprompt@gap>| !gapinput@TestMonomialQuick( S4 );|
  rec( comment := "abelian by supersolvable group", isMonomial := true )
  !gapprompt@gap>| !gapinput@TestMonomialQuick( Sl23 );|
  rec( comment := "no decision by cheap tests", isMonomial := "?" )
\end{Verbatim}
 }

 
\subsection{\textcolor{Chapter }{TestSubnormallyMonomial}}\logpage{[ 75, 4, 5 ]}
\hyperdef{L}{X7E56A0EA868CC34A}{}
{
\noindent\textcolor{FuncColor}{$\triangleright$\enspace\texttt{TestSubnormallyMonomial({\mdseries\slshape G})\index{TestSubnormallyMonomial@\texttt{TestSubnormallyMonomial}!for a group}
\label{TestSubnormallyMonomial:for a group}
}\hfill{\scriptsize (attribute)}}\\
\noindent\textcolor{FuncColor}{$\triangleright$\enspace\texttt{TestSubnormallyMonomial({\mdseries\slshape chi})\index{TestSubnormallyMonomial@\texttt{TestSubnormallyMonomial}!for a character}
\label{TestSubnormallyMonomial:for a character}
}\hfill{\scriptsize (attribute)}}\\
\noindent\textcolor{FuncColor}{$\triangleright$\enspace\texttt{IsSubnormallyMonomial({\mdseries\slshape G})\index{IsSubnormallyMonomial@\texttt{IsSubnormallyMonomial}!for a group}
\label{IsSubnormallyMonomial:for a group}
}\hfill{\scriptsize (property)}}\\
\noindent\textcolor{FuncColor}{$\triangleright$\enspace\texttt{IsSubnormallyMonomial({\mdseries\slshape chi})\index{IsSubnormallyMonomial@\texttt{IsSubnormallyMonomial}!for a character}
\label{IsSubnormallyMonomial:for a character}
}\hfill{\scriptsize (property)}}\\


 An irreducible character of the group $G$ is called \emph{subnormally monomial} (\emph{SM} for short) if it is induced from a linear character of a subnormal subgroup of $G$. A group $G$ is called SM if all its irreducible characters are SM. 

 \texttt{TestSubnormallyMonomial} returns a record with information whether the group \mbox{\texttt{\mdseries\slshape G}} or the irreducible character \mbox{\texttt{\mdseries\slshape chi}} of \mbox{\texttt{\mdseries\slshape G}} is SM. 

 The result has the components \texttt{isSubnormallyMonomial} (either \texttt{true} or \texttt{false}) and \texttt{comment} (a string telling a reason for the value of the component \texttt{isSubnormallyMonomial}); in the case that the \texttt{isSubnormallyMonomial} component has value \texttt{false} there is also a component \texttt{character}, with value an irreducible character of $G$ that is not SM. 

 \texttt{IsSubnormallyMonomial} returns \texttt{true} if the group \mbox{\texttt{\mdseries\slshape G}} or the group character \mbox{\texttt{\mdseries\slshape chi}} is subnormally monomial, and \texttt{false} otherwise. 

 
\begin{Verbatim}[commandchars=!@|,fontsize=\small,frame=single,label=Example]
  !gapprompt@gap>| !gapinput@TestSubnormallyMonomial( S4 );|
  rec( character := Character( CharacterTable( S4 ), [ 3, -1, -1, 0, 1
       ] ), comment := "found non-SM character",
    isSubnormallyMonomial := false )
  !gapprompt@gap>| !gapinput@TestSubnormallyMonomial( Irr( S4 )[4] );|
  rec( comment := "all subnormal subgroups checked",
    isSubnormallyMonomial := false )
  !gapprompt@gap>| !gapinput@TestSubnormallyMonomial( DerivedSubgroup( S4 ) );|
  rec( comment := "all irreducibles checked",
    isSubnormallyMonomial := true )
\end{Verbatim}
 }

 
\subsection{\textcolor{Chapter }{TestRelativelySM}}\logpage{[ 75, 4, 6 ]}
\hyperdef{L}{X83EF7B8D7C1C2CA3}{}
{
\noindent\textcolor{FuncColor}{$\triangleright$\enspace\texttt{TestRelativelySM({\mdseries\slshape G})\index{TestRelativelySM@\texttt{TestRelativelySM}!for a group}
\label{TestRelativelySM:for a group}
}\hfill{\scriptsize (attribute)}}\\
\noindent\textcolor{FuncColor}{$\triangleright$\enspace\texttt{TestRelativelySM({\mdseries\slshape chi})\index{TestRelativelySM@\texttt{TestRelativelySM}!for a character}
\label{TestRelativelySM:for a character}
}\hfill{\scriptsize (attribute)}}\\
\noindent\textcolor{FuncColor}{$\triangleright$\enspace\texttt{TestRelativelySM({\mdseries\slshape G, N})\index{TestRelativelySM@\texttt{TestRelativelySM}!for a group and a normal subgroup}
\label{TestRelativelySM:for a group and a normal subgroup}
}\hfill{\scriptsize (operation)}}\\
\noindent\textcolor{FuncColor}{$\triangleright$\enspace\texttt{TestRelativelySM({\mdseries\slshape chi, N})\index{TestRelativelySM@\texttt{TestRelativelySM}!for a character and a normal subgroup}
\label{TestRelativelySM:for a character and a normal subgroup}
}\hfill{\scriptsize (operation)}}\\
\noindent\textcolor{FuncColor}{$\triangleright$\enspace\texttt{IsRelativelySM({\mdseries\slshape G})\index{IsRelativelySM@\texttt{IsRelativelySM}!for a group}
\label{IsRelativelySM:for a group}
}\hfill{\scriptsize (property)}}\\
\noindent\textcolor{FuncColor}{$\triangleright$\enspace\texttt{IsRelativelySM({\mdseries\slshape chi})\index{IsRelativelySM@\texttt{IsRelativelySM}!for a character}
\label{IsRelativelySM:for a character}
}\hfill{\scriptsize (property)}}\\


 In the first two cases, \texttt{TestRelativelySM} returns a record with information whether the argument, which must be a SM
group \mbox{\texttt{\mdseries\slshape G}} or an irreducible character \mbox{\texttt{\mdseries\slshape chi}} of a SM group $G$, is relatively SM with respect to every normal subgroup of \mbox{\texttt{\mdseries\slshape G}}. 

 In the second two cases, a normal subgroup \mbox{\texttt{\mdseries\slshape N}} of \mbox{\texttt{\mdseries\slshape G}} is the second argument. Here \texttt{TestRelativelySM} returns a record with information whether the first argument is relatively SM
with respect to \mbox{\texttt{\mdseries\slshape N}}, i.e, whether there is a subnormal subgroup $H$ of $G$ that contains \mbox{\texttt{\mdseries\slshape N}} such that the character \mbox{\texttt{\mdseries\slshape chi}} resp.{\nobreakspace}every irreducible character of $G$ is induced from a character $\psi$ of $H$ such that the restriction of $\psi$ to \mbox{\texttt{\mdseries\slshape N}} is irreducible. 

 The result record has the components \texttt{isRelativelySM} (with value either \texttt{true} or \texttt{false}) and \texttt{comment} (a string that describes a reason). If the argument is a group \mbox{\texttt{\mdseries\slshape G}} that is not relatively SM with respect to a normal subgroup then additionally
the component \texttt{character} is bound, with value a not relatively SM character of such a normal subgroup. 

 \texttt{IsRelativelySM} returns \texttt{true} if the SM group \mbox{\texttt{\mdseries\slshape G}} or the irreducible character \mbox{\texttt{\mdseries\slshape chi}} of the SM group \mbox{\texttt{\mdseries\slshape G}} is relatively SM with respect to every normal subgroup of \mbox{\texttt{\mdseries\slshape G}}, and \texttt{false} otherwise. 

 \emph{Note} that it is \emph{not} checked whether \mbox{\texttt{\mdseries\slshape G}} is SM. 

 
\begin{Verbatim}[commandchars=!@|,fontsize=\small,frame=single,label=Example]
  !gapprompt@gap>| !gapinput@IsSubnormallyMonomial( DerivedSubgroup( S4 ) );|
  true
  !gapprompt@gap>| !gapinput@TestRelativelySM( DerivedSubgroup( S4 ) );|
  rec(
    comment := "normal subgroups are abelian or have nilpotent factor gr\
  oup", isRelativelySM := true )
\end{Verbatim}
 }

 }

  
\section{\textcolor{Chapter }{Minimal Nonmonomial Groups}}\label{Minimal Nonmonomial Groups}
\logpage{[ 75, 5, 0 ]}
\hyperdef{L}{X7B5735897DE29BCB}{}
{
  

\subsection{\textcolor{Chapter }{IsMinimalNonmonomial}}
\logpage{[ 75, 5, 1 ]}\nobreak
\hyperdef{L}{X7D7E2667821A23CD}{}
{\noindent\textcolor{FuncColor}{$\triangleright$\enspace\texttt{IsMinimalNonmonomial({\mdseries\slshape G})\index{IsMinimalNonmonomial@\texttt{IsMinimalNonmonomial}}
\label{IsMinimalNonmonomial}
}\hfill{\scriptsize (property)}}\\


 A group \mbox{\texttt{\mdseries\slshape G}} is called \emph{minimal nonmonomial} if it is nonmonomial, and all proper subgroups and factor groups are monomial. 

 
\begin{Verbatim}[commandchars=!@|,fontsize=\small,frame=single,label=Example]
  !gapprompt@gap>| !gapinput@IsMinimalNonmonomial( Sl23 );  IsMinimalNonmonomial( S4 );|
  true
  false
\end{Verbatim}
 }

 

\subsection{\textcolor{Chapter }{MinimalNonmonomialGroup}}
\logpage{[ 75, 5, 2 ]}\nobreak
\hyperdef{L}{X7B416BBD80072079}{}
{\noindent\textcolor{FuncColor}{$\triangleright$\enspace\texttt{MinimalNonmonomialGroup({\mdseries\slshape p, factsize})\index{MinimalNonmonomialGroup@\texttt{MinimalNonmonomialGroup}}
\label{MinimalNonmonomialGroup}
}\hfill{\scriptsize (function)}}\\


 is a solvable minimal nonmonomial group described by the parameters \mbox{\texttt{\mdseries\slshape factsize}} and \mbox{\texttt{\mdseries\slshape p}} if such a group exists, and \texttt{false} otherwise. 

 Suppose that the required group $K$ exists. Then \mbox{\texttt{\mdseries\slshape factsize}} is the size of the Fitting factor $K / F(K)$, and this value is 4, 8, an odd prime, twice an odd prime, or four times an
odd prime. In the case that \mbox{\texttt{\mdseries\slshape factsize}} is twice an odd prime, the centre $Z(K)$ is cyclic of order $2^{{\mbox{\texttt{\mdseries\slshape p}}+1}}$. In all other cases \mbox{\texttt{\mdseries\slshape p}} is the (unique) prime that divides the order of $F(K)$. 

 The solvable minimal nonmonomial groups were classified by van der Waall,
see{\nobreakspace}\cite{vdW76}. 

 
\begin{Verbatim}[commandchars=!@|,fontsize=\small,frame=single,label=Example]
  !gapprompt@gap>| !gapinput@MinimalNonmonomialGroup(  2,  3 ); # the group SL(2,3)|
  2^(1+2):3
  !gapprompt@gap>| !gapinput@MinimalNonmonomialGroup(  3,  4 );|
  3^(1+2):4
  !gapprompt@gap>| !gapinput@MinimalNonmonomialGroup(  5,  8 );|
  5^(1+2):Q8
  !gapprompt@gap>| !gapinput@MinimalNonmonomialGroup( 13, 12 );|
  13^(1+2):2.D6
  !gapprompt@gap>| !gapinput@MinimalNonmonomialGroup(  1, 14 );|
  2^(1+6):D14
  !gapprompt@gap>| !gapinput@MinimalNonmonomialGroup(  2, 14 );|
  (2^(1+6)Y4):D14
\end{Verbatim}
 }

 }

 }

    
\chapter{\textcolor{Chapter }{Using and Developing \textsf{GAP} Packages}}\label{Using and Developing GAP Packages}
\logpage{[ 76, 0, 0 ]}
\hyperdef{L}{X79F76C1E834BFDCC}{}
{
  \index{package} The functionality of \textsf{GAP} can be extended by loading \textsf{GAP} packages. The \textsf{GAP} distribution already contains all currently redistributed \textsf{GAP} packages in the \texttt{gap\texttt{\symbol{45}}4.17dev/pkg} directory. 

 \textsf{GAP} packages are written by (groups of) \textsf{GAP} users who may not necessarily be members of the \textsf{GAP} developer team. The responsibility and copyright of a \textsf{GAP} package remains with the original author(s). 

 \textsf{GAP} packages have their own documentation which is smoothly integrated into the \textsf{GAP} help system. (When \textsf{GAP} is started, \texttt{LoadPackageDocumentation} is called for all packages.) 

 All \textsf{GAP} users who develop new code are invited to share the results of their efforts
with other \textsf{GAP} users by making the code and its documentation available in form of a package.
Guidance on how to do this is available from the \textsf{GAP} website (\href{https://www.gap-system.org} {\texttt{https://www.gap\texttt{\symbol{45}}system.org}}) and in the \textsf{GAP} package \textsf{Example} (see \href{https://www.gap-system.org/Packages/example.html} {\texttt{https://www.gap\texttt{\symbol{45}}system.org/Packages/example.html}}). 

 The \textsf{GAP} development team will assist in making any new package suitable for
distribution with \textsf{GAP}. It is also possible to submit a package to a formal refereeing process. 

 In this chapter we first describe how to use existing packages, and then
provide guidelines for writing a \textsf{GAP} package.  
\section{\textcolor{Chapter }{Installing a \textsf{GAP} Package}}\label{Installing a GAP Package}
\logpage{[ 76, 1, 0 ]}
\hyperdef{L}{X82473E4B8756C6CD}{}
{
  Before a package can be used it must be installed. A standard distribution of \textsf{GAP} already contains all the packages currently redistributed with \textsf{GAP}. This set of packages has been checked for compatibility with the system and
with each other during release preparation. Most of the packages can be used
immediately, but some of them may require further installation steps (see
below). 

 Also, since \textsf{GAP} packages are released independently of the main \textsf{GAP} system, it may sometimes be useful to upgrade or install new packages between
upgrades of your \textsf{GAP} installation, e.g. if a new version of a package adds new capabilities or bug
fixes that you need. 

 A package consists of a collection of files within a single directory that
must be a subdirectory of the \texttt{pkg} directory in one of the \textsf{GAP} root directories (see \ref{GAP Root Directories}). If you don't have access to the \texttt{pkg} directory in your main \textsf{GAP} installation you can add private root directories as explained in section \ref{GAP Root Directories}. 

 Whenever you download or clone an archive of a \textsf{GAP} package, it will contain a \texttt{README} file (or \texttt{README.md} etc.) that explains how it should be installed. Some packages just consist of \textsf{GAP} code and the installation is done by unpacking the archive in one of the
places described above. There are also packages that need further installation
steps, such as compilation or installing additional software to satisfy their
dependencies. If there are some external programs which have to be compiled,
this is often done by executing \texttt{./configure; make} inside the unpacked package directory (but check the individual \texttt{README} files). 

 Most of the packages that require compilation can be compiled in a single step
by changing to the \texttt{pkg} directory of your \textsf{GAP} installation and calling the \texttt{../bin/BuildPackages.sh} script. 

 Note that if you use Windows you may not be able to use some or all external
binaries. }

  
\section{\textcolor{Chapter }{Loading a \textsf{GAP} Package}}\label{Loading a GAP Package}
\logpage{[ 76, 2, 0 ]}
\hyperdef{L}{X825CBC5B86F8F811}{}
{
  If a package is not already loaded, it may be loaded using the function \texttt{LoadPackage} (\ref{LoadPackage}). 

 Some \textsf{GAP} packages are prepared for automatic loading, that is they will be loaded
automatically when \textsf{GAP} starts (see \ref{LoadPackageAutomatic}). 

\subsection{\textcolor{Chapter }{LoadPackage}}
\logpage{[ 76, 2, 1 ]}\nobreak
\hyperdef{L}{X79B373A77B29D1F5}{}
{\noindent\textcolor{FuncColor}{$\triangleright$\enspace\texttt{LoadPackage({\mdseries\slshape name[, version][, banner]})\index{LoadPackage@\texttt{LoadPackage}}
\label{LoadPackage}
}\hfill{\scriptsize (function)}}\\


 loads the \textsf{GAP} package with name \mbox{\texttt{\mdseries\slshape name}}. 

 As an example, the following loads the \textsf{GAP} package \textsf{SONATA} (case insensitive) which provides methods for the construction and analysis of
finite nearrings: 

 
\begin{Verbatim}[commandchars=!@|,fontsize=\small,frame=single,label=Example]
  !gapprompt@gap>| !gapinput@LoadPackage("sonata");|
  ... some more lines with package banner(s) ...
  true
\end{Verbatim}
 

 The package name is case insensitive and may be appropriately abbreviated. At
the time of writing, for example, \texttt{LoadPackage("semi");} will load the \textsf{Semigroups} package, and \texttt{LoadPackage("js");} will load the \textsf{json} package. If the abbreviation cannot be uniquely completed, a list of available
completions will be offered, and \texttt{LoadPackage} returns \texttt{fail}. Thus the names of \emph{all} installed packages can be shown by calling \texttt{LoadPackage("");}. 

 When the optional argument string \mbox{\texttt{\mdseries\slshape version}} is present, the package will only be loaded in a version number equal to or
greater than \mbox{\texttt{\mdseries\slshape version}} (see{\nobreakspace}\texttt{CompareVersionNumbers} (\ref{CompareVersionNumbers})). If the first character of \mbox{\texttt{\mdseries\slshape version}} is \texttt{=} then only that version will be loaded. 

 \texttt{LoadPackage} will return \texttt{true} if the package has been successfully loaded, and will return \texttt{fail} if the package could not be loaded. The latter may be the case if the package
is not installed, if necessary binaries have not been compiled, or if the
version number of the available version is too small. If the package cannot be
loaded, \texttt{TestPackageAvailability} (\ref{TestPackageAvailability}) can be used to find the reasons. Also, \texttt{DisplayPackageLoadingLog} (\ref{DisplayPackageLoadingLog}) can be used to find out more about the failure. To see the problems directly,
one can change the verbosity using the user preference \texttt{InfoPackageLoadingLevel}, see \texttt{InfoPackageLoading} (\ref{InfoPackageLoading}) for details. 

 If the package \mbox{\texttt{\mdseries\slshape name}} has already been loaded in a version number equal to or greater than \mbox{\texttt{\mdseries\slshape version}}, \texttt{LoadPackage} returns \texttt{true} without doing anything else. 

 If the optional argument \mbox{\texttt{\mdseries\slshape banner}} is present then it must be either \texttt{true} or \texttt{false}; in the latter case, the effect is that no package banner is printed. 

 After a package has been loaded, all its code becomes available to use with
the rest of the \textsf{GAP} library. }

 \index{automatic loading of GAP packages@automatic loading of \textsf{GAP} packages} \index{disable automatic loading} 
\subsection{\textcolor{Chapter }{Automatic loading of \textsf{GAP} packages}}\label{LoadPackageAutomatic}
\logpage{[ 76, 2, 2 ]}
\hyperdef{L}{X7E6767B485F23BFC}{}
{
  When \textsf{GAP} is started some packages are loaded automatically, and these belong to two
categories. The first are those packages which are needed to start \textsf{GAP} (at the present time, the only such package is \textsf{GAPDoc}). Their list is contained in \texttt{GAPInfo.Dependencies.NeededOtherPackages}. The second are packages which are loaded during \textsf{GAP} startup by default. The latter list may be obtained by calling \texttt{UserPreference("PackagesToLoad")} and is customisable as described in Section  (\textbf{Reference: Configuring User preferences}). 

 While \textsf{GAP} will not start if any of the packages from the former group is missing,
loading of the packages from the latter group may be suppressed by using the \texttt{\texttt{\symbol{45}}A} command line option (see \ref{Command Line Options}). 

 If for some reason you don't want certain packages to be automatically loaded, \textsf{GAP} provides three levels for disabling autoloading. 

 \index{NOAUTO@\texttt{NOAUTO}} The autoloading of specific packages can be overwritten \emph{for the whole \textsf{GAP} installation} by putting a file \texttt{NOAUTO} into a \texttt{pkg} directory that contains lines with the names of packages which should not be
automatically loaded. 

 Furthermore, \emph{individual users} can disable the autoloading of specific packages by putting the names of these
packages into the list that is assigned to the user preference ``ExcludeFromAutoload'', for example in the user's \texttt{gap.ini} file (see{\nobreakspace}\ref{subsect:gap.ini file}). 

 Using the \texttt{\texttt{\symbol{45}}A} command line option when starting \textsf{GAP} (see{\nobreakspace}\ref{Command Line Options}), automatic loading of packages is switched off \emph{for this \textsf{GAP} session}. 

 In any of the above three cases, the packages listed in \texttt{GAPInfo.Dependencies.NeededOtherPackages} are still loaded automatically, and an error is signalled if any of these
packages is unavailable. 

 See \texttt{SetPackagePath} (\ref{SetPackagePath}) for a way to force the loading of a prescribed package version. See also \texttt{ExtendRootDirectories} (\ref{ExtendRootDirectories}) and \texttt{ExtendPackageDirectories} (\ref{ExtendPackageDirectories}) for methods of adding directories containing packages \emph{after} \textsf{GAP} has been started. }

 

\subsection{\textcolor{Chapter }{SetPackagePath}}
\logpage{[ 76, 2, 3 ]}\nobreak
\hyperdef{L}{X858E8985840BFA72}{}
{\noindent\textcolor{FuncColor}{$\triangleright$\enspace\texttt{SetPackagePath({\mdseries\slshape pkgname, pkgpath})\index{SetPackagePath@\texttt{SetPackagePath}}
\label{SetPackagePath}
}\hfill{\scriptsize (function)}}\\


 This function can be used to force \textsf{GAP} to load a particular version of a package, even though newer versions of the
package are available. 

 Let \mbox{\texttt{\mdseries\slshape pkgname}} and \mbox{\texttt{\mdseries\slshape pkgpath}} be strings denoting the name of a \textsf{GAP} package and the path to a directory where a version of this package can be
found (i.{\nobreakspace}e., calling \texttt{Directory} (\ref{Directory}) with the argument \mbox{\texttt{\mdseries\slshape pkgpath}} will yield a directory that contains the file \texttt{PackageInfo.g} of the package). 

 If the package \mbox{\texttt{\mdseries\slshape pkgname}} is already loaded with an installation path different from \mbox{\texttt{\mdseries\slshape pkgpath}} then \texttt{SetPackagePath} signals an error. If the package \mbox{\texttt{\mdseries\slshape pkgname}} is not yet loaded then \texttt{SetPackagePath} erases the information about available versions of the package \mbox{\texttt{\mdseries\slshape pkgname}}, and stores the record that is contained in the \texttt{PackageInfo.g} file at \mbox{\texttt{\mdseries\slshape pkgpath}} instead, such that only the version installed at \mbox{\texttt{\mdseries\slshape pkgpath}} can be loaded with \texttt{LoadPackage} (\ref{LoadPackage}). 

 One should call \texttt{SetPackagePath} immediately before loading the package in question. Note that calling \texttt{ExtendPackageDirectories} (\ref{ExtendPackageDirectories}) or \texttt{ExtendRootDirectories} (\ref{ExtendRootDirectories}) may change the available versions of the package \mbox{\texttt{\mdseries\slshape pkgname}}. }

 

\subsection{\textcolor{Chapter }{ExtendRootDirectories}}
\logpage{[ 76, 2, 4 ]}\nobreak
\hyperdef{L}{X7CD0A2F27D19BA03}{}
{\noindent\textcolor{FuncColor}{$\triangleright$\enspace\texttt{ExtendRootDirectories({\mdseries\slshape paths})\index{ExtendRootDirectories@\texttt{ExtendRootDirectories}}
\label{ExtendRootDirectories}
}\hfill{\scriptsize (function)}}\\


 Let \mbox{\texttt{\mdseries\slshape paths}} be a list of strings that denote paths to intended \textsf{GAP} root directories (see \ref{GAP Root Directories}). The function \texttt{ExtendRootDirectories} adds these paths to the global list \texttt{GAPInfo.RootPaths} and calls the initialization of available \textsf{GAP} packages, such that later calls to \texttt{LoadPackage} (\ref{LoadPackage}) will find the \textsf{GAP} packages that are contained in \texttt{pkg} subdirectories of the directories given by \mbox{\texttt{\mdseries\slshape paths}}. 

 Note that the purpose of this function is to make \textsf{GAP} packages in the given directories available. It cannot be used to influence
the start of \textsf{GAP}, because the \textsf{GAP} library is loaded before \texttt{ExtendRootDirectories} can be called (and because \texttt{GAPInfo.RootPaths} is not used for reading the \textsf{GAP} library). }

 

\subsection{\textcolor{Chapter }{ExtendPackageDirectories}}
\logpage{[ 76, 2, 5 ]}\nobreak
\hyperdef{L}{X81BFC5C87EE2E02A}{}
{\noindent\textcolor{FuncColor}{$\triangleright$\enspace\texttt{ExtendPackageDirectories({\mdseries\slshape paths})\index{ExtendPackageDirectories@\texttt{ExtendPackageDirectories}}
\label{ExtendPackageDirectories}
}\hfill{\scriptsize (function)}}\\


 Let \mbox{\texttt{\mdseries\slshape paths}} be a list of strings that denote paths to intended \textsf{GAP} package directories (see \ref{GAP Package Directories}). The function \texttt{ExtendPackageDirectories} adds these paths to the global list \texttt{GAPInfo.PackageDirectories} and calls the initialization of available \textsf{GAP} packages, such that later calls to \texttt{LoadPackage} (\ref{LoadPackage}) will find the \textsf{GAP} packages that are contained in the directories given by \mbox{\texttt{\mdseries\slshape paths}}. }

 

\subsection{\textcolor{Chapter }{DisplayPackageLoadingLog}}
\logpage{[ 76, 2, 6 ]}\nobreak
\hyperdef{L}{X7D162DDF813D2BBA}{}
{\noindent\textcolor{FuncColor}{$\triangleright$\enspace\texttt{DisplayPackageLoadingLog({\mdseries\slshape [severity]})\index{DisplayPackageLoadingLog@\texttt{DisplayPackageLoadingLog}}
\label{DisplayPackageLoadingLog}
}\hfill{\scriptsize (function)}}\\
\noindent\textcolor{FuncColor}{$\triangleright$\enspace\texttt{InfoPackageLoading\index{InfoPackageLoading@\texttt{InfoPackageLoading}}
\label{InfoPackageLoading}
}\hfill{\scriptsize (info class)}}\\
\noindent\textcolor{FuncColor}{$\triangleright$\enspace\texttt{PACKAGE{\textunderscore}ERROR\index{PACKAGE{\textunderscore}ERROR@\texttt{PACKAGE{\textunderscore}ERROR}}
\label{PACKAGEuScoreERROR}
}\hfill{\scriptsize (global variable)}}\\
\noindent\textcolor{FuncColor}{$\triangleright$\enspace\texttt{PACKAGE{\textunderscore}WARNING\index{PACKAGE{\textunderscore}WARNING@\texttt{PACKAGE{\textunderscore}WARNING}}
\label{PACKAGEuScoreWARNING}
}\hfill{\scriptsize (global variable)}}\\
\noindent\textcolor{FuncColor}{$\triangleright$\enspace\texttt{PACKAGE{\textunderscore}INFO\index{PACKAGE{\textunderscore}INFO@\texttt{PACKAGE{\textunderscore}INFO}}
\label{PACKAGEuScoreINFO}
}\hfill{\scriptsize (global variable)}}\\
\noindent\textcolor{FuncColor}{$\triangleright$\enspace\texttt{PACKAGE{\textunderscore}DEBUG\index{PACKAGE{\textunderscore}DEBUG@\texttt{PACKAGE{\textunderscore}DEBUG}}
\label{PACKAGEuScoreDEBUG}
}\hfill{\scriptsize (global variable)}}\\
\noindent\textcolor{FuncColor}{$\triangleright$\enspace\texttt{LogPackageLoadingMessage({\mdseries\slshape severity, message[, name]})\index{LogPackageLoadingMessage@\texttt{LogPackageLoadingMessage}}
\label{LogPackageLoadingMessage}
}\hfill{\scriptsize (function)}}\\


 Whenever \textsf{GAP} considers loading a package, log messages are collected in a global list. The
messages for the current \textsf{GAP} session can be displayed with \texttt{DisplayPackageLoadingLog}. To each message, a ``severity'' is assigned, which is one of \texttt{PACKAGE{\textunderscore}ERROR}, \texttt{PACKAGE{\textunderscore}WARNING}, \texttt{PACKAGE{\textunderscore}INFO}, \texttt{PACKAGE{\textunderscore}DEBUG}, in increasing order. The function \texttt{DisplayPackageLoadingLog} shows only the messages whose severity is at most \mbox{\texttt{\mdseries\slshape severity}}, the default for \mbox{\texttt{\mdseries\slshape severity}} is \texttt{PACKAGE{\textunderscore}WARNING}. 

 The intended meaning of the severity levels is as follows. 

 
\begin{description}
\item[{PACKAGE{\textunderscore}ERROR}]  should be used whenever \textsf{GAP} will run into an error during package loading, where the reason of the error
shall be documented in the global list. 
\item[{PACKAGE{\textunderscore}WARNING}]  should be used whenever \textsf{GAP} has detected a reason why a package cannot be loaded, and where the message
describes how to solve this problem, for example if a package binary is
missing. 
\item[{PACKAGE{\textunderscore}INFO}]  should be used whenever \textsf{GAP} has detected a reason why a package cannot be loaded, and where it is not
clear how to solve this problem, for example if the package is not compatible
with other installed packages. 
\item[{PACKAGE{\textunderscore}DEBUG}]  should be used for other messages reporting what \textsf{GAP} does when it loads packages (checking dependencies, reading files, etc.). One
purpose is to record in which order packages have been considered for loading
or have actually been loaded. 
\end{description}
 

 The log messages are created either by the functions of \textsf{GAP}'s package loading mechanism or in the code of your package, for example in
the \texttt{AvailabilityTest} function of the package's \texttt{PackageInfo.g} file (see \ref{The PackageInfo.g File}), using \texttt{LogPackageLoadingMessage}. The arguments of this function are \mbox{\texttt{\mdseries\slshape severity}} (which must be one of the above severity levels), \mbox{\texttt{\mdseries\slshape message}} (which must be either a string or a list of strings), and optionally \mbox{\texttt{\mdseries\slshape name}} (which must be the name of the package to which the message belongs). The
argument \mbox{\texttt{\mdseries\slshape name}} is not needed if the function is called from a call of a package's \texttt{AvailabilityTest} function (see \ref{The PackageInfo.g File}) or is called from a package file that is read from \texttt{init.g} or \texttt{read.g}; in these cases, the name of the current package (stored in the record \texttt{GAPInfo.PackageCurrent}) is taken. According to the above list, the \mbox{\texttt{\mdseries\slshape severity}} argument of \texttt{LogPackageLoadingMessage} calls in a package's \texttt{AvailabilityTest} function is either \texttt{PACKAGE{\textunderscore}WARNING} or \texttt{PACKAGE{\textunderscore}INFO}. 

 If you want to see the log messages already during the package loading
process, you can set the level of the info class \texttt{InfoPackageLoading} to one of the severity values listed above; afterwards the messages with at
most this severity are shown immediately when they arise. In order to make
this work already for autoloaded packages, you can call \texttt{SetUserPreference("InfoPackageLoadingLevel", \mbox{\texttt{\mdseries\slshape lev}});} to set the desired severity level \mbox{\texttt{\mdseries\slshape lev}}. This can for example be done in your \texttt{gap.ini} file, see Section \ref{subsect:gap.ini file}. }

 }

  
\section{\textcolor{Chapter }{Functions for \textsf{GAP} Packages}}\label{Functions for GAP Packages}
\logpage{[ 76, 3, 0 ]}
\hyperdef{L}{X7C6CE28B7E142804}{}
{
  The following functions are mainly used in files contained in a package and
not by users of a package. They are needed to organise reading package files
into \textsf{GAP} in the right order, performing maintenance tasks like building documentation
and running package tests, checking package dependencies, etc. You will find
further information about their use in Section \ref{Guidelines for Writing a GAP Package} and subsequent sections. 

\subsection{\textcolor{Chapter }{ReadPackage}}
\logpage{[ 76, 3, 1 ]}\nobreak
\hyperdef{L}{X870954577B27DCAB}{}
{\noindent\textcolor{FuncColor}{$\triangleright$\enspace\texttt{ReadPackage({\mdseries\slshape [name, ]file})\index{ReadPackage@\texttt{ReadPackage}}
\label{ReadPackage}
}\hfill{\scriptsize (function)}}\\
\noindent\textcolor{FuncColor}{$\triangleright$\enspace\texttt{RereadPackage({\mdseries\slshape [name, ]file})\index{RereadPackage@\texttt{RereadPackage}}
\label{RereadPackage}
}\hfill{\scriptsize (function)}}\\


 Called with two strings \mbox{\texttt{\mdseries\slshape name}} and \mbox{\texttt{\mdseries\slshape file}}, \texttt{ReadPackage} reads the file \mbox{\texttt{\mdseries\slshape file}} of the \textsf{GAP} package \mbox{\texttt{\mdseries\slshape name}}, where \mbox{\texttt{\mdseries\slshape file}} is given as a path relative to the home directory of \mbox{\texttt{\mdseries\slshape name}}. Note that \mbox{\texttt{\mdseries\slshape file}} is read in the namespace of the package, see Section \ref{Namespaces} for details. 

 If only one argument \mbox{\texttt{\mdseries\slshape file}} is given, this should be the path of a file relative to the \texttt{pkg} subdirectory of \textsf{GAP} root paths (see{\nobreakspace}\ref{GAP Root Directories}). Note that in this case, the package name is assumed to be equal to the
first part of \mbox{\texttt{\mdseries\slshape file}}, \emph{so the one argument form is not recommended}. 

 The absolute path is determined as follows. If the package in question has
already been loaded then the file in the directory of the loaded version is
read. If the package is available but not yet loaded then the directory given
by \texttt{TestPackageAvailability} (\ref{TestPackageAvailability}) is used, without prescribed version number. (Note that the \texttt{ReadPackage} call does \emph{not} force the package to be loaded.) 

 If the file is readable then \texttt{true} is returned, otherwise a warning is displayed (for \texttt{ReadPackage}) or \texttt{false} is returned (for \texttt{RereadPackage}). 

 Each of \mbox{\texttt{\mdseries\slshape name}} and \mbox{\texttt{\mdseries\slshape file}} should be a string. The \mbox{\texttt{\mdseries\slshape name}} argument is case insensitive. 

 \texttt{RereadPackage} does the same as \texttt{ReadPackage}, except that also read\texttt{\symbol{45}}only global variables are
overwritten (cf.{\nobreakspace}\texttt{Reread} (\ref{Reread})). }

 

\subsection{\textcolor{Chapter }{TestPackageAvailability}}
\logpage{[ 76, 3, 2 ]}\nobreak
\hyperdef{L}{X8580DF257E4D7046}{}
{\noindent\textcolor{FuncColor}{$\triangleright$\enspace\texttt{TestPackageAvailability({\mdseries\slshape name[, version][, checkall]})\index{TestPackageAvailability@\texttt{TestPackageAvailability}}
\label{TestPackageAvailability}
}\hfill{\scriptsize (function)}}\\


 For strings \mbox{\texttt{\mdseries\slshape name}} and \mbox{\texttt{\mdseries\slshape version}}, this function tests whether the \textsf{GAP} package \mbox{\texttt{\mdseries\slshape name}} is available for loading in a version that is at least \mbox{\texttt{\mdseries\slshape version}}, or equal to \mbox{\texttt{\mdseries\slshape version}} if the first character of \mbox{\texttt{\mdseries\slshape version}} is \texttt{=} (see \texttt{CompareVersionNumbers} (\ref{CompareVersionNumbers}) for further details about version numbers). 

 The result is \texttt{true} if the package is already loaded, \texttt{fail} if it is not available, and the string denoting the \textsf{GAP} root path where the package resides if it is available, but not yet loaded. So
the package \mbox{\texttt{\mdseries\slshape name}} is available if the result of \texttt{TestPackageAvailability} is not equal to \texttt{fail}. 

 If the optional argument \mbox{\texttt{\mdseries\slshape checkall}} is \texttt{true} then all dependencies are checked, even if some have turned out to be not
satisfied. This is useful when one is interested in the reasons why the
package \mbox{\texttt{\mdseries\slshape name}} cannot be loaded. In this situation, calling first \texttt{TestPackageAvailability} and then \texttt{DisplayPackageLoadingLog} (\ref{DisplayPackageLoadingLog}) with argument \texttt{PACKAGE{\textunderscore}INFO} (\ref{PACKAGEuScoreINFO}) will give an overview of these reasons. 

 You should \emph{not} call \texttt{TestPackageAvailability} in the test function of a package (the value of the component \texttt{AvailabilityTest} in the \texttt{PackageInfo.g} file of the package, see \ref{The PackageInfo.g File}), because \texttt{TestPackageAvailability} calls this test function.  

 The argument \mbox{\texttt{\mdseries\slshape name}} is case insensitive. }

 

\subsection{\textcolor{Chapter }{IsPackageLoaded}}
\logpage{[ 76, 3, 3 ]}\nobreak
\hyperdef{L}{X7C8724C183E24665}{}
{\noindent\textcolor{FuncColor}{$\triangleright$\enspace\texttt{IsPackageLoaded({\mdseries\slshape name[, version]})\index{IsPackageLoaded@\texttt{IsPackageLoaded}}
\label{IsPackageLoaded}
}\hfill{\scriptsize (function)}}\\


 For strings \mbox{\texttt{\mdseries\slshape name}} and \mbox{\texttt{\mdseries\slshape version}}, this function tests whether the \textsf{GAP} package \mbox{\texttt{\mdseries\slshape name}} is already loaded in a version that is at least \mbox{\texttt{\mdseries\slshape version}}, or equal to \mbox{\texttt{\mdseries\slshape version}} if the first character of \mbox{\texttt{\mdseries\slshape version}} is \texttt{=} (see \texttt{CompareVersionNumbers} (\ref{CompareVersionNumbers}) for further details about version numbers). 

 The result is \texttt{true} if the package is already loaded, \texttt{false} otherwise. }

 

\subsection{\textcolor{Chapter }{IsPackageMarkedForLoading}}
\logpage{[ 76, 3, 4 ]}\nobreak
\hyperdef{L}{X8067348B836BAF37}{}
{\noindent\textcolor{FuncColor}{$\triangleright$\enspace\texttt{IsPackageMarkedForLoading({\mdseries\slshape name, version})\index{IsPackageMarkedForLoading@\texttt{IsPackageMarkedForLoading}}
\label{IsPackageMarkedForLoading}
}\hfill{\scriptsize (function)}}\\


 This function can be used in the code of a package $A$ for testing whether the package \mbox{\texttt{\mdseries\slshape name}} in version \mbox{\texttt{\mdseries\slshape version}} will be loaded after the \texttt{LoadPackage} (\ref{LoadPackage}) call for the package $A$ has been executed. This means that the package \mbox{\texttt{\mdseries\slshape name}} had been loaded before, or has been (directly or indirectly) requested as a
needed or suggested package of the package $A$ or of a package whose loading requested that $A$ was loaded. }

 

\subsection{\textcolor{Chapter }{TestPackage}}
\logpage{[ 76, 3, 5 ]}\nobreak
\hyperdef{L}{X866ADD4E814A54F0}{}
{\noindent\textcolor{FuncColor}{$\triangleright$\enspace\texttt{TestPackage({\mdseries\slshape pkgname})\index{TestPackage@\texttt{TestPackage}}
\label{TestPackage}
}\hfill{\scriptsize (function)}}\\


 It is recommended that a \textsf{GAP} package specifies a standard test in its \texttt{PackageInfo.g} file. If \mbox{\texttt{\mdseries\slshape pkgname}} is a string with the name of a \textsf{GAP} package, then \texttt{TestPackage(pkgname)} will check if this package is loadable and has the standard test, and will run
this test in the current \textsf{GAP} session.

 The output of the test depends on the particular package, and it also may
depend on the current \textsf{GAP} session (loaded packages, state of the random sources, defined global
variables etc.). }

 

\subsection{\textcolor{Chapter }{InstalledPackageVersion}}
\logpage{[ 76, 3, 6 ]}\nobreak
\hyperdef{L}{X7B79FEE57DBDBD71}{}
{\noindent\textcolor{FuncColor}{$\triangleright$\enspace\texttt{InstalledPackageVersion({\mdseries\slshape name})\index{InstalledPackageVersion@\texttt{InstalledPackageVersion}}
\label{InstalledPackageVersion}
}\hfill{\scriptsize (function)}}\\


 If the \textsf{GAP} package with name \mbox{\texttt{\mdseries\slshape name}} has already been loaded then \texttt{InstalledPackageVersion} returns the string denoting the version number of this version of the package.
If the package is available but has not yet been loaded then the version
number string for that version of the package that currently would be loaded.
(Note that loading \emph{another} package might force loading another version of the package \mbox{\texttt{\mdseries\slshape name}}, so the result of \texttt{InstalledPackageVersion} will be different afterwards.) If the package is not available then \texttt{fail} is returned. 

 The argument \mbox{\texttt{\mdseries\slshape name}} is case insensitive. }

 

\subsection{\textcolor{Chapter }{DirectoriesPackageLibrary}}
\logpage{[ 76, 3, 7 ]}\nobreak
\hyperdef{L}{X807D835C7B032D4E}{}
{\noindent\textcolor{FuncColor}{$\triangleright$\enspace\texttt{DirectoriesPackageLibrary({\mdseries\slshape name[, path]})\index{DirectoriesPackageLibrary@\texttt{DirectoriesPackageLibrary}}
\label{DirectoriesPackageLibrary}
}\hfill{\scriptsize (function)}}\\


 takes the string \mbox{\texttt{\mdseries\slshape name}}, a name of a \textsf{GAP} package, and returns a list that is either empty or contains one directory
object \texttt{dir} that describes the place where the library functions of this \textsf{GAP} package should be located. 

 In the latter case, \texttt{dir} is the \mbox{\texttt{\mdseries\slshape path}} subdirectory of a directory where the package \mbox{\texttt{\mdseries\slshape name}} is installed, where the default for \mbox{\texttt{\mdseries\slshape path}} is \texttt{"lib"}, and where the package directory belongs to the version of \mbox{\texttt{\mdseries\slshape name}} that is already loaded or is currently going to be loaded or would be the
first version \textsf{GAP} would try to load if no other version is explicitly prescribed. (If the
package \mbox{\texttt{\mdseries\slshape name}} is not yet loaded then we cannot guarantee that the directory belongs to a
version that really can be loaded.) 

 Note that \texttt{DirectoriesPackageLibrary} is likely to be called in the \texttt{AvailabilityTest} function in the package's \texttt{PackageInfo.g} file (see \ref{The PackageInfo.g File}). 

 As an example, the following returns a directory object for the library
functions of the \textsf{GAP} package \textsf{Example}: 

 
\begin{Verbatim}[commandchars=!@|,fontsize=\small,frame=single,label=Example]
  !gapprompt@gap>| !gapinput@DirectoriesPackageLibrary( "Example", "gap" );|
  [ dir("/home/werner/gap/4.0/pkg/example/gap/") ]
\end{Verbatim}
 

 Observe that we needed the second argument \texttt{"gap"} here, since \textsf{Example}'s library functions are in the subdirectory \texttt{gap} rather than \texttt{lib}. 

 In order to find a subdirectory deeper than one level in a package directory,
the second argument is again necessary whether or not the desired subdirectory
relative to the package's directory begins with \texttt{lib}. The directories in \mbox{\texttt{\mdseries\slshape path}} should be separated by \texttt{/} (even on systems, like Windows, which use \texttt{\texttt{\symbol{92}}} as the directory separator). For example, suppose there is a package \texttt{somepackage} with a subdirectory \texttt{m11} in the directory \texttt{data}, then we might expect the following: 

 
\begin{Verbatim}[commandchars=!@|,fontsize=\small,frame=single,label=Example]
  !gapprompt@gap>| !gapinput@DirectoriesPackageLibrary( "somepackage", "data/m11" );|
  [ dir("/home/werner/gap/4.0/pkg/somepackage/data/m11") ]
\end{Verbatim}
 }

 

\subsection{\textcolor{Chapter }{DirectoriesPackagePrograms}}
\logpage{[ 76, 3, 8 ]}\nobreak
\hyperdef{L}{X794508E5811D3BC9}{}
{\noindent\textcolor{FuncColor}{$\triangleright$\enspace\texttt{DirectoriesPackagePrograms({\mdseries\slshape name})\index{DirectoriesPackagePrograms@\texttt{DirectoriesPackagePrograms}}
\label{DirectoriesPackagePrograms}
}\hfill{\scriptsize (function)}}\\


 \index{GAPInfo.Architecture@\texttt{GAPInfo.Architecture}} returns a list that is either empty or contains one directory object \texttt{dir} that describes the place where external binaries of the \textsf{GAP} package \mbox{\texttt{\mdseries\slshape name}} should be located. 

 In the latter case, \texttt{dir} is the \texttt{bin/}\mbox{\texttt{\mdseries\slshape architecture}} subdirectory of a directory where the package \mbox{\texttt{\mdseries\slshape name}} is installed, where \mbox{\texttt{\mdseries\slshape architecture}} is the architecture on which \textsf{GAP} has been compiled (this can be accessed as \texttt{GAPInfo.Architecture}, see \texttt{GAPInfo} (\ref{GAPInfo})), and where the package directory belongs to the version of \mbox{\texttt{\mdseries\slshape name}} that is already loaded or is currently going to be loaded or would be the
first version \textsf{GAP} would try to load if no other version is explicitly prescribed. (If the
package \mbox{\texttt{\mdseries\slshape name}} is not yet loaded then we cannot guarantee that the directory belongs to a
version that really can be loaded.) 

 Note that \texttt{DirectoriesPackagePrograms} is likely to be called in the \texttt{AvailabilityTest} function in the package's \texttt{PackageInfo.g} file (see \ref{The PackageInfo.g File}). 

 
\begin{Verbatim}[commandchars=!@|,fontsize=\small,frame=single,label=Example]
  !gapprompt@gap>| !gapinput@DirectoriesPackagePrograms( "nq" );|
  [ dir("/home/gap/4.0/pkg/nq/bin/x86_64-pc-linux-gnu-default64-kv3/") ]
\end{Verbatim}
 }

 

\subsection{\textcolor{Chapter }{CompareVersionNumbers}}
\logpage{[ 76, 3, 9 ]}\nobreak
\hyperdef{L}{X787DFEB383545A49}{}
{\noindent\textcolor{FuncColor}{$\triangleright$\enspace\texttt{CompareVersionNumbers({\mdseries\slshape supplied, required[, "equal"]})\index{CompareVersionNumbers@\texttt{CompareVersionNumbers}}
\label{CompareVersionNumbers}
}\hfill{\scriptsize (function)}}\\


 A version number is a string which contains nonnegative integers separated by
non\texttt{\symbol{45}}numeric characters. Examples of valid version numbers
are for example: 

 
\begin{Verbatim}[commandchars=!@|,fontsize=\small,frame=single,label=Example]
  "1.0"   "3.141.59"  "2-7-8.3" "5 release 2 patchlevel 666"
\end{Verbatim}
 

 \texttt{CompareVersionNumbers} compares two version numbers, given as strings. They are split at
non\texttt{\symbol{45}}digit characters, the resulting integer lists are
compared lexicographically. The routine tests whether \mbox{\texttt{\mdseries\slshape supplied}} is at least as large as \mbox{\texttt{\mdseries\slshape required}}, and returns \texttt{true} or \texttt{false} accordingly. A version number ending in \texttt{dev} is considered to be infinite. }

 

\subsection{\textcolor{Chapter }{DeclareAutoreadableVariables}}
\logpage{[ 76, 3, 10 ]}\nobreak
\hyperdef{L}{X8495E5327D563AC3}{}
{\noindent\textcolor{FuncColor}{$\triangleright$\enspace\texttt{DeclareAutoreadableVariables({\mdseries\slshape pkgname, filename, varlist})\index{DeclareAutoreadableVariables@\texttt{DeclareAutoreadableVariables}}
\label{DeclareAutoreadableVariables}
}\hfill{\scriptsize (function)}}\\


 Let \mbox{\texttt{\mdseries\slshape pkgname}} be the name of a package, let \mbox{\texttt{\mdseries\slshape filename}} be the name of a file relative to the home directory of this package, and let \mbox{\texttt{\mdseries\slshape varlist}} be a list of strings that are the names of global variables which get bound
when the file is read. \texttt{DeclareAutoreadableVariables} notifies the names in \mbox{\texttt{\mdseries\slshape varlist}} such that the first attempt to access one of the variables causes the file to
be read. }

 
\subsection{\textcolor{Chapter }{Kernel modules in \textsf{GAP} packages}}\label{Kernel modules}
\logpage{[ 76, 3, 11 ]}
\hyperdef{L}{X85672DDD7D34D5F0}{}
{
  \index{gac@\texttt{gac}} If the package has a kernel module, then it can be compiled using the \textsf{gac} script. A kernel module is implemented in C and follows certain conventions to
comply with the \textsf{GAP} kernel interface, which we plan to document later. In the meantime, we advice
to get in touch with \textsf{GAP} developers if you plan to develop such a package. 

 To use the \textsf{gac} script to produce dynamically loadable modules, call it with the \texttt{\texttt{\symbol{45}}d} option, for example: 

 
\begin{Verbatim}[commandchars=!@|,fontsize=\small,frame=single,label=Example]
  $ gap4/gac -d test.c
\end{Verbatim}
 

 This will produce a file \texttt{test.so}, which then can be loaded into \textsf{GAP} with \texttt{LoadKernelExtension} (\ref{LoadKernelExtension}). If the kernel module is required for the package to work, then its \texttt{PackageInfo.g} should define a \texttt{AvailabilityTest} which calls \texttt{IsKernelExtensionAvailable} (\ref{IsKernelExtensionAvailable}), see \ref{Test for the Existence of GAP Package Binaries} for details. 

 Note that before \textsf{GAP} 4.12, \texttt{LoadDynamicModule} (\ref{LoadDynamicModule}) was used for this. It is still available and in fact \texttt{LoadKernelExtension} (\ref{LoadKernelExtension}) call it; but the latter provides a higher level abstraction and is more
convenient to use. }

 

\subsection{\textcolor{Chapter }{IsKernelExtensionAvailable}}
\logpage{[ 76, 3, 12 ]}\nobreak
\hyperdef{L}{X7D23646086F33A9E}{}
{\noindent\textcolor{FuncColor}{$\triangleright$\enspace\texttt{IsKernelExtensionAvailable({\mdseries\slshape pkgname[, modname]})\index{IsKernelExtensionAvailable@\texttt{IsKernelExtensionAvailable}}
\label{IsKernelExtensionAvailable}
}\hfill{\scriptsize (function)}}\\


 For use by packages: Search for a loadable kernel module inside package \mbox{\texttt{\mdseries\slshape pkgname}} with name \mbox{\texttt{\mdseries\slshape modname}} and return \texttt{true} if found, otherwise \texttt{false}. If \mbox{\texttt{\mdseries\slshape modname}} is omitted, then \mbox{\texttt{\mdseries\slshape pkgname}} is used instead. Note that package names are case insensitive, but \mbox{\texttt{\mdseries\slshape modname}} is not. 

 This function first appeared in GAP 4.12. It is typically called in the \texttt{AvailabilityTest} function of a package (see \ref{Test for the Existence of GAP Package Binaries}). 
\begin{Verbatim}[commandchars=!@|,fontsize=\small,frame=single,label=Example]
  !gapprompt@gap>| !gapinput@IsKernelExtensionAvailable("myPackageWithKernelExtension");|
  true
\end{Verbatim}
 }

 

\subsection{\textcolor{Chapter }{LoadKernelExtension}}
\logpage{[ 76, 3, 13 ]}\nobreak
\hyperdef{L}{X85F1B83079A11F89}{}
{\noindent\textcolor{FuncColor}{$\triangleright$\enspace\texttt{LoadKernelExtension({\mdseries\slshape pkgname[, modname]})\index{LoadKernelExtension@\texttt{LoadKernelExtension}}
\label{LoadKernelExtension}
}\hfill{\scriptsize (function)}}\\


 For use by packages: Search for a loadable kernel module inside package \mbox{\texttt{\mdseries\slshape pkgname}} with name \mbox{\texttt{\mdseries\slshape modname}}, and load it if found. If \mbox{\texttt{\mdseries\slshape modname}} is omitted, then \mbox{\texttt{\mdseries\slshape pkgname}} is used instead. Note that package names are case insensitive, but \mbox{\texttt{\mdseries\slshape modname}} is not. 

 This function first appeared in GAP 4.12. It is typically called in the \texttt{init.g} file of a package. 

 Previously, packages with a kernel module typically used code like this: 
\begin{Verbatim}[commandchars=!@|,fontsize=\small,frame=single,label=]
  path := Filename(DirectoriesPackagePrograms("SomePackage"), "SomePackage.so");
  if path <> fail then
    LoadDynamicModule(path);
  fi;
\end{Verbatim}
 That can now be replaced by the following, which also produces more helpful
error messages for the user: 
\begin{Verbatim}[commandchars=!@|,fontsize=\small,frame=single,label=]
  LoadKernelExtension("SomePackage");
\end{Verbatim}
 For packages where the name of the kernel extension is not identical to that
of the package, you can either rename the kernel extension to have a matching
name (recommended if you only have a single kernel extension in your package,
which is how we recommend to set up things anyway), or else use the two
argument version: 
\begin{Verbatim}[commandchars=!@|,fontsize=\small,frame=single,label=Example]
  LoadKernelExtension("SomePackage", "kext"); # this will look for kext.so
\end{Verbatim}
 }

 

\subsection{\textcolor{Chapter }{LoadDynamicModule}}
\logpage{[ 76, 3, 14 ]}\nobreak
\hyperdef{L}{X7C99782886B18C77}{}
{\noindent\textcolor{FuncColor}{$\triangleright$\enspace\texttt{LoadDynamicModule({\mdseries\slshape filename})\index{LoadDynamicModule@\texttt{LoadDynamicModule}}
\label{LoadDynamicModule}
}\hfill{\scriptsize (function)}}\\


 To load a compiled file, the command \texttt{LoadDynamicModule} is used. This command loads \mbox{\texttt{\mdseries\slshape filename}} as module. 

 
\begin{Verbatim}[commandchars=!@|,fontsize=\small,frame=single,label=Example]
  !gapprompt@gap>| !gapinput@LoadDynamicModule("./test.so");|
\end{Verbatim}
 

 On some operating systems, once you have loaded a dynamic module with a
certain filename, loading another with the same filename will have no effect,
even if the file on disk has changed. }

 
\subsection{\textcolor{Chapter }{The PackageInfo.g File}}\label{The PackageInfo.g File}
\logpage{[ 76, 3, 15 ]}
\hyperdef{L}{X85C8DE357EE424D8}{}
{
  Each package has the file \texttt{PackageInfo.g} which contains meta\texttt{\symbol{45}}information about the package (package
name, version, author(s), relations to other packages, homepage, download
archives, etc.). This file is used by the package loading mechanism, by the \textsf{GAP} webpages about packages, and also for the redistribution of a package with \textsf{GAP}. 

 A \texttt{PackageInfo.g} file contains a call to the function \texttt{SetPackageInfo}, with argument a record. The following components of this record are \emph{mandatory}. 
\begin{description}
\item[{\texttt{PackageName}}]  a nonempty string denoting the name of the package, 
\item[{\texttt{Subtitle}}]  a string that describes the package's contents, may be used by a default
banner or on a web page, should fit on one line, 
\item[{\texttt{Version}}]  a nonempty string that does not start with \texttt{=}, denoting the version number of the package (see Section \ref{Version Numbers}), 
\item[{\texttt{Date}}]  a string of the form \texttt{yyyy\texttt{\symbol{45}}mm\texttt{\symbol{45}}dd}  denoting the release date of the current version of the package (a date since
1999, when \textsf{GAP}{\nobreakspace}4 appeared), 
\item[{\texttt{License}}]  a nonempty string containing an SPDX ID (see Section \ref{Selecting a license for a GAP Package}), 
\item[{\texttt{ArchiveURL}}]  a string started with \texttt{http://}, \texttt{https://}, or \texttt{ftp://}, denoting an URL from where the current package archive can be downloaded,
but without the suffix describing the format (see the \texttt{ArchiveFormats} component), 
\item[{\texttt{ArchiveFormats}}]  a string that lists the supported formats (among \texttt{.tar.gz}, \texttt{.tar.bz2}, \texttt{\texttt{\symbol{45}}win.zip}), separated by whitespace or commas, 
\item[{\texttt{README{\textunderscore}URL}}]  a string started with \texttt{http://}, \texttt{https://}, or \texttt{ftp://}, denoting an URL from where the current \texttt{README.md} or \texttt{README} file of the package can be downloaded, 
\item[{\texttt{PackageInfoURL}}]  a string started with \texttt{http://}, \texttt{https://}, or \texttt{ftp://}, denoting an URL from where the current \texttt{PackageInfo.g} file of the package can be downloaded, 
\item[{\texttt{AbstractHTML}}]  a string that describes the package's contents in a few lines, in HTML format;
this text will be displayed on the package overview web page of \textsf{GAP}, 
\item[{\texttt{PackageWWWHome}}]  a string started with \texttt{http://}, \texttt{https://}, or \texttt{ftp://}, denoting the address of the package's home page, 
\item[{\texttt{PackageDoc}}]  a record or a list of records; each record describes a book of the package
documentation, with the following components 
\begin{description}
\item[{\texttt{BookName}}]  a string, the name of the book, 
\item[{\texttt{LongTitle}}]  a string shown by \texttt{?books}, 
\item[{\texttt{SixFile}}]  a string denoting a relative path to the \texttt{manual.six} file of the book, 
\item[{\texttt{HTMLStart}}]  a string denoting a relative path to the start file of the HTML version of the
book, 
\item[{\texttt{PDFFile}}]  a string denoting a relative path to the \texttt{.pdf} file of the book, 
\item[{\texttt{ArchiveURLSubset}}]  a list of strings denoting relative paths to those files and directories from
the archive that are needed for the online manual; typically, \texttt{[ "doc" ]} suffices, 
\end{description}
 
\end{description}
 The following components of the record are \emph{optional}. 
\begin{description}
\item[{\texttt{TextFiles} or \texttt{BinaryFiles} or \texttt{TextBinaryFilesPatterns}}]  a list of strings that specify which files in the archive are text files or
binary files (at most one of the three components can be available, each
string in \texttt{TextBinaryFilesPatterns} must start with \texttt{T} for text files and by \texttt{B} for binary files), 
\item[{\texttt{Persons}}]  a list of records, each with the mandatory components 
\begin{description}
\item[{\texttt{LastName}}]  a string, 
\item[{at least one of \texttt{IsAuthor} or \texttt{IsMaintainer}}]  \texttt{true} or \texttt{false}, 
\end{description}
 and optional components 
\begin{description}
\item[{\texttt{FirstNames}}]  a string (was mandatory before \textsf{GAP} 4.14), 
\item[{\texttt{Place}}]  a string, 
\item[{\texttt{Institution}}]  a string, 
\end{description}
 If the \texttt{IsMaintainer} value is \texttt{true} then also one of the following components is mandatory, otherwise these
components are optional. 
\begin{description}
\item[{\texttt{Email}}]  a string, 
\item[{\texttt{WWWHome}}]  a string denoting an URL, or 
\item[{\texttt{PostalAddress}}]  a string. 
\end{description}
 
\item[{\texttt{SourceRepository}}]  a record with the components \texttt{Type} (the version control system, e.g. \texttt{"git"} or \texttt{"hg"}) and \texttt{URL} (the URL of the repository), both strings, 
\item[{\texttt{IssueTrackerURL}}]  a string started with \texttt{http://}, \texttt{https://}, or \texttt{ftp://}, 
\item[{\texttt{SupportEmail}}]  a string denoting an e\texttt{\symbol{45}}mail address, 
\item[{\texttt{Dependencies}}]  a record describing the dependencies of the package (see Section \ref{Package dependencies}), with the following optional components 
\begin{description}
\item[{\texttt{GAP}}]  a string denoting the needed version of \textsf{GAP}, 
\item[{\texttt{NeededOtherPackages}}]  a list of pairs \texttt{[ pkgname, pkgversion ]} of strings, denoting the other packages which must be available if the current
package shall be loadable, 
\item[{\texttt{SuggestedOtherPackages}}]  a list of pairs \texttt{[ pkgname, pkgversion ]} of strings, denoting the other packages which shall be loaded together with
the current package if they are available, 
\item[{\texttt{TestPackages}}]  a list of pairs \texttt{[ pkgname, pkgversion ]} of strings, denoting packages that are needed only to test the package, but
not to load the package, 
\item[{\texttt{NeededSystemPackages}}]  a record whose contents describe external dependencies of the package in a
machine\texttt{\symbol{45}}readable form (see \ref{External dependencies} for more information), 
\item[{\texttt{ExternalConditions}}]  a list of strings or of pairs \texttt{[ text, URL ]} of strings, denoting conditions on external programs in a
human\texttt{\symbol{45}}readable form (see \ref{External dependencies} for more information), 
\end{description}
 
\item[{\texttt{AvailabilityTest}}]  a function with no arguments that returns \texttt{true} if the package is available, and \texttt{false} otherwise (can be \texttt{ReturnTrue} (\ref{ReturnTrue}) if the package consists only of \textsf{GAP} code; this is also the default value), 
\item[{\texttt{BannerString} or \texttt{BannerFunction}}]  a string or a function, respectively, that is used to create a package banner
different from the default banner (see Section \ref{The Banner}), 
\item[{\texttt{TestFile}}]  a string denoting a relative path to a readable file which contains tests of
the package's functionality (see Section \ref{Testing a GAP package}), 
\item[{\texttt{Keywords}}]  a list of strings that are keywords related to the topic of the package, 
\item[{\texttt{Extensions}}]  a list of records that describe conditional extensions of the package (see
Section \ref{Extensions Provided by a Package}). 
\end{description}
 Other components of the record can be supported; for example, \texttt{AutoDoc} is used by the \textsf{AutoDoc} package if applicable. }

 

\subsection{\textcolor{Chapter }{ValidatePackageInfo}}
\logpage{[ 76, 3, 16 ]}\nobreak
\hyperdef{L}{X79767C2482FF6F55}{}
{\noindent\textcolor{FuncColor}{$\triangleright$\enspace\texttt{ValidatePackageInfo({\mdseries\slshape info})\index{ValidatePackageInfo@\texttt{ValidatePackageInfo}}
\label{ValidatePackageInfo}
}\hfill{\scriptsize (function)}}\\


 This function is intended to support package authors who create or modify \texttt{PackageInfo.g} files. (It is \emph{not} called when these files are read during the startup of \textsf{GAP} or when packages are actually loaded.) 

 The argument \mbox{\texttt{\mdseries\slshape info}} must be either a record as is contained in a \texttt{PackageInfo.g} file or a string which describes the path to such a file. The result is \texttt{true} if the record or the contents of the file, respectively, has correct format,
and \texttt{false} otherwise; in the latter case information about the incorrect components is
printed. These diagnostic messages can be suppressed by setting the global
option \texttt{quiet} to \texttt{true}. 

 Note that the components used for package loading are checked as well as the
components that are needed for composing the package overview web page or for
updating the package archives. 

 If \mbox{\texttt{\mdseries\slshape info}} is a string then \texttt{ValidatePackageInfo} checks additionally whether those package files exist that are mentioned in
the file \texttt{info}, for example the \texttt{manual.six} file of the package documentation. }

 

\subsection{\textcolor{Chapter }{ShowPackageVariables}}
\logpage{[ 76, 3, 17 ]}\nobreak
\hyperdef{L}{X7D34AC3287611B15}{}
{\noindent\textcolor{FuncColor}{$\triangleright$\enspace\texttt{ShowPackageVariables({\mdseries\slshape pkgname[, version][, arec]})\index{ShowPackageVariables@\texttt{ShowPackageVariables}}
\label{ShowPackageVariables}
}\hfill{\scriptsize (function)}}\\
\noindent\textcolor{FuncColor}{$\triangleright$\enspace\texttt{PackageVariablesInfo({\mdseries\slshape pkgname, version})\index{PackageVariablesInfo@\texttt{PackageVariablesInfo}}
\label{PackageVariablesInfo}
}\hfill{\scriptsize (function)}}\\


 Let \mbox{\texttt{\mdseries\slshape pkgname}} be the name of a \textsf{GAP} package. If the package \mbox{\texttt{\mdseries\slshape pkgname}} is available but not yet loaded then \texttt{ShowPackageVariables} prints a list of global variables that become bound and of methods that become
installed when the package is loaded. (For that, \textsf{GAP} actually loads the package.) 

 If a version number \mbox{\texttt{\mdseries\slshape version}} is given (see Section{\nobreakspace} (\textbf{Reference: Version Numbers})) then this version of the package is considered. 

 An error message is printed if (the given version of) the package is not
available or already loaded. 

 Information is printed about new and redeclared global variables, and about
names of global variables introduced in the package that differ from existing
globals only by case; note that the \textsf{GAP} help system is case insensitive, so it is difficult to document identifiers
that differ only by case. 

 Info lines for undocumented variables are marked with an asterisk \texttt{*}. 

 The following entries are omitted from the list: default setter methods for
attributes and properties that are declared in the package, and \texttt{Set\mbox{\texttt{\mdseries\slshape attr}}} and \texttt{Has\mbox{\texttt{\mdseries\slshape attr}}} type variables where \mbox{\texttt{\mdseries\slshape attr}} is an attribute or property. 

 The output can be customized using the optional record \mbox{\texttt{\mdseries\slshape arec}}, the following components of this record are supported. 
\begin{description}
\item[{\texttt{show}}]  a list of strings describing those kinds of variables which shall be shown,
such as \texttt{"new global functions"}; the default are all kinds that appear in the package, 
\item[{\texttt{showDocumented}}]  \texttt{true} (the default) if documented variables shall be shown, and \texttt{false} otherwise, 
\item[{\texttt{showUndocumented}}]  \texttt{true} (the default) if undocumented variables shall be shown, and \texttt{false} otherwise, 
\item[{\texttt{showPrivate}}]  \texttt{true} (the default) if variables from the package's name space (see Section \ref{Namespaces}) shall be shown, and \texttt{false} otherwise, 
\item[{\texttt{Display}}]  a function that takes a string and shows it on the screen; the default is \texttt{Print} (\ref{Print}), another useful value is \texttt{Pager} (\ref{Pager}). 
\end{description}
 

 An interactive variant of \texttt{ShowPackageVariables} is the function \texttt{BrowsePackageVariables} (\textbf{Browse: BrowsePackageVariables}) that is provided by the \textsf{GAP} package \textsf{Browse}. For this function, it is not sensible to assume that the package \mbox{\texttt{\mdseries\slshape pkgname}} is not yet loaded before the function call, because one might be interested in
packages that must be loaded before \textsf{Browse} itself can be loaded. The solution is that \texttt{BrowsePackageVariables} (\textbf{Browse: BrowsePackageVariables}) takes the output of \texttt{PackageVariablesInfo} as its second argument. The function \texttt{PackageVariablesInfo} is used by both \texttt{ShowPackageVariables} and \texttt{BrowsePackageVariables} (\textbf{Browse: BrowsePackageVariables}) for collecting the information about the package in question, and can be
called before the package \textsf{Browse} is loaded. }

 

\subsection{\textcolor{Chapter }{BibEntry}}
\logpage{[ 76, 3, 18 ]}\nobreak
\hyperdef{L}{X79EA4BD37940AD25}{}
{\noindent\textcolor{FuncColor}{$\triangleright$\enspace\texttt{BibEntry({\mdseries\slshape pkgname[, key]})\index{BibEntry@\texttt{BibEntry}}
\label{BibEntry}
}\hfill{\scriptsize (function)}}\\
\textbf{\indent Returns:\ }
 a string in BibXMLext format (see  (\textbf{GAPDoc: The BibXMLext Format})) that can be used for referencing the \textsf{GAP} system or a \textsf{GAP} package. 



 If the argument \mbox{\texttt{\mdseries\slshape pkgname}} is the string \texttt{"GAP"}, the function returns an entry for the current version of \textsf{GAP}. 

 Otherwise, if a string \mbox{\texttt{\mdseries\slshape pkgname}} is given, which is the name of a \textsf{GAP} package, an entry for this package is returned; this entry is computed from
the \texttt{PackageInfo.g} file of \emph{the current version} of the package, see \texttt{InstalledPackageVersion} (\ref{InstalledPackageVersion}). If no package with name \mbox{\texttt{\mdseries\slshape pkgname}} is installed then the empty string is returned. 

 A string for \emph{a different version} of \textsf{GAP} or a package can be computed by entering, as the argument \mbox{\texttt{\mdseries\slshape pkgname}}, the desired record from the \texttt{PackageInfo.g} file. (One can access these records using the function \texttt{PackageInfo}.) 

 In each of the above cases, an optional argument \mbox{\texttt{\mdseries\slshape key}} can be given, a string which is then used as the key of the Bib{\TeX} entry instead of the default key that is generated from the system/package
name and the version number. 

 \texttt{BibEntry} requires the functions \texttt{FormatParagraph} (\textbf{GAPDoc: FormatParagraph}) and \texttt{NormalizedNameAndKey} (\textbf{GAPDoc: NormalizedNameAndKey}) from the \textsf{GAP} package \textsf{GAPDoc}. 

 The functions \texttt{ParseBibXMLextString} (\textbf{GAPDoc: ParseBibXMLextString}) and \texttt{StringBibXMLEntry} (\textbf{GAPDoc: StringBibXMLEntry}) can be used to create for example a Bib{\TeX} entry from the return value, as follows. 

 
\begin{Verbatim}[commandchars=!|D,fontsize=\small,frame=single,label=Example]
  !gapprompt|gap>D !gapinput|bib:= BibEntry( "GAP", "GAP4.5" );;D
  !gapprompt|gap>D !gapinput|Print( bib, "\n" );D
  <entry id="GAP4.5"><misc>
    <title><C>GAP</C> &ndash; <C>G</C>roups, <C>A</C>lgorithms,
           and <C>P</C>rogramming, <C>V</C>ersion 4.5.1</title>
    <howpublished><URL>https://www.gap-system.org</URL></howpublished>
    <key>GAP</key>
    <keywords>groups; *; gap; manual</keywords>
    <other type="organization">The GAP <C>G</C>roup</other>
  </misc></entry>
  !gapprompt|gap>D !gapinput|parse:= ParseBibXMLextString( bib );;D
  !gapprompt|gap>D !gapinput|Print( StringBibXMLEntry( parse.entries[1], "BibTeX" ) );D
  @misc{ GAP4.5,
    title =            {{GAP}   {\textendash}   {G}roups,   {A}lgorithms,  and
                        {P}rogramming, {V}ersion 4.5.1},
    organization =     {The GAP {G}roup},
    howpublished =     {\href                      {https://www.gap-system.org}
                        {\texttt{https://www.gap-system.org}}},
    key =              {GAP},
    keywords =         {groups; *; gap; manual}
  }
\end{Verbatim}
 }

 

\subsection{\textcolor{Chapter }{Cite}}
\logpage{[ 76, 3, 19 ]}\nobreak
\hyperdef{L}{X79637D9A7B1AD7F7}{}
{\noindent\textcolor{FuncColor}{$\triangleright$\enspace\texttt{Cite({\mdseries\slshape [pkgname[, key]]})\index{Cite@\texttt{Cite}}
\label{Cite}
}\hfill{\scriptsize (function)}}\\


 Used with no arguments or with argument \texttt{"GAP"} (case\texttt{\symbol{45}}insensitive), \texttt{Cite} displays instructions on citing the version of \textsf{GAP} that is being used. Suggestions are given in plain text, HTML, BibXML and
BibTeX formats. The same instructions are also contained in the \texttt{CITATION} file in the \textsf{GAP} root directory. 

 If \mbox{\texttt{\mdseries\slshape pkgname}} is the name of a \textsf{GAP} package, instructions on citing this package will be displayed. They will be
produced from the \texttt{PackageInfo.g} file of the working version of this package that must be available in the \textsf{GAP} installation being used. Otherwise, one will get a warning that no working
version of the package is available. 

 The optional 2nd argument \mbox{\texttt{\mdseries\slshape key}} has the same meaning as in \texttt{BibEntry} (\ref{BibEntry}). }

 }

 
\section{\textcolor{Chapter }{Guidelines for Writing a \textsf{GAP} Package}}\label{Guidelines for Writing a GAP Package}
\logpage{[ 76, 4, 0 ]}
\hyperdef{L}{X7EE8E5D97B0F8AFF}{}
{
  The remaining part of this chapter explains the basics of how to write a \textsf{GAP} package so that it integrates properly into \textsf{GAP}. 

 There are two basic aspects of creating a \textsf{GAP} package. 

 First, it is a convenient possibility to load additional functionality into \textsf{GAP} including a smooth integration of the package documentation. Second, a package
is a way to make your code available to other \textsf{GAP} users. 

 Moreover, the \textsf{GAP} Group may provide some help with redistributing your package via the \textsf{GAP} website after checking if the package provides some new or improved
functionality which looks interesting for other users, if it contains
reasonable documentation, and if it seems to work smoothly with the \textsf{GAP} library and other distributed packages. In this case the package can take part
in the \textsf{GAP} distribution update mechanism and becomes a \emph{deposited} package. 

 Furthermore, package authors are encouraged to check if the package would be
appropriate for the refereeing process and \emph{submit} it. If the refereeing has been successful, the package becomes an \emph{accepted} package. Check out \href{https://www.gap-system.org/Packages/Authors/authors.html} {\texttt{https://www.gap\texttt{\symbol{45}}system.org/Packages/Authors/authors.html}} on the \textsf{GAP} website for more details. 

 Below we start with a description how the directory structure of a \textsf{GAP} package should be constructed and then add remarks on certain aspects of
creating a package, some of these only apply to some packages. Finally, we
provide guidelines for the release preparation and its distribution. 

 }

  
\section{\textcolor{Chapter }{Structure of a \textsf{GAP} Package}}\label{Structure of a GAP Package}
\logpage{[ 76, 5, 0 ]}
\hyperdef{L}{X8383876782480702}{}
{
  \index{home directory!for a GAP package} A \textsf{GAP} package should have an alphanumeric name; mixed case is fine, but there should
be no whitespace characters. All files of a \textsf{GAP} package \mbox{\texttt{\mdseries\slshape packagename}} must be collected in a single directory \mbox{\texttt{\mdseries\slshape packagedir}}, where \mbox{\texttt{\mdseries\slshape packagedir}} should be just \mbox{\texttt{\mdseries\slshape packagename}} optionally converted to lowercase and optionally followed by the package
version (with or without hyphen to separate the version from \mbox{\texttt{\mdseries\slshape packagename}}). Let us call this directory the \emph{home directory} of the package. 

 To use the package with \textsf{GAP}, the directory \mbox{\texttt{\mdseries\slshape packagedir}} must be a subdirectory of a \texttt{pkg} directory in (one of) the \textsf{GAP} root directories (see \ref{GAP Root Directories}). For example, if \textsf{GAP} is installed in \texttt{/usr/local/gap4} then the files of the package \texttt{MyPack} may be placed in the directory \texttt{/usr/local/gap4/pkg/mypack}. The directory \mbox{\texttt{\mdseries\slshape packagedir}} preferably should have the following structure (below, a trailing \texttt{/} distinguishes directories from ordinary files): 

 \newpage 
\begin{Verbatim}[commandchars=!@|,fontsize=\small,frame=single,label=Example]
  packagedir/
    doc/
    lib/
    tst/
    CHANGES
    LICENSE
    README
    PackageInfo.g
    init.g
    read.g
\end{Verbatim}
 

 This layout of directories and files may be created manually, or automatically
using the tool called \textsf{PackageMaker}, available at \href{https://github.com/gap-system/PackageMaker} {\texttt{https://github.com/gap\texttt{\symbol{45}}system/PackageMaker}}. The \textsf{PackageMaker} asks several questions about the intended package and then creates a new
directory for it and populates it with all the files needed for a basic
package. 

 Packages that contain some code that requires compilation will usually have it
in the \texttt{src} subdirectory. They may also have extra files such as \texttt{configure}, \texttt{Makefile.in} etc. that automate the build procedure. There are three file names with a
special meaning in the home directory of a package: \texttt{PackageInfo.g} and \texttt{init.g} which must be present, and \texttt{read.g} which is optional. 

 On the other hand, the names of \texttt{CHANGES}, \texttt{LICENSE} and \texttt{README} files are not strictly fixed. They may have extensions \texttt{.txt} or \texttt{.md}, and instead of \texttt{LICENSE} one could use e.g. \texttt{COPYING} or \texttt{GPL} for packages distributed under the GNU General Public License, or use \texttt{HISTORY} instead of \texttt{CHANGES}. 

 We now describe the above files and directories in more details: 

 
\begin{description}
\item[{ \texttt{README}}]  \index{README@\texttt{README}!for a GAP package} The filename may optionally have an extension, e.g. \texttt{.txt} or \texttt{.md}. 

 This should contain ``how to get it'' instructions (covering the way of getting it with the \textsf{GAP} distribution and from the \textsf{GAP} website, if applicable), as well as installation instructions and names of the
package authors and their email addresses. These installation instructions
should be repeated or referenced from the package's documentation, which
should be in the \texttt{doc} directory (see \ref{Writing Documentation and Tools Needed}). Authors' names and addresses should be repeated both in the package's
documentation and in the \texttt{PackageInfo.g} (see below). 
\item[{ \texttt{CHANGES}}]  For further versions of the package, it will be also useful to have a \texttt{CHANGES} file that records the main changes between versions of the package. 

 The filename may optionally have an extension, e.g. \texttt{.txt} or \texttt{.md}. 
\item[{ \texttt{LICENSE}}]  The file which explains conditions on which the package is distributed. 

 We advise all package authors to make clear in the documentation of their
package the basis on which it is being distributed to users. Technically, this
is the terms of the license which you give the users to copy, modify and
redistribute your software (of which you presumably own the copyright) for
their purposes. 

 \textsf{GAP} itself is distributed under the GNU General Public License version 2, a
popular ``free software'' license which allows users to redistribute it freely under the same terms, and
requires that any software which incorporates \textsf{GAP} (technically, any ``derived work'') also be distributed under those terms. We would encourage you to consider
the GPL for your packages, but you might wish to be more restrictive (for
instance forbidding redistribution for profit) or less restrictive (allowing
your software to be incorporated into commercial software). 

 The filename may optionally have an extension, e.g. \texttt{.txt} or \texttt{.md}. Some packages also use different filenames, like \texttt{COPYING}. 
\item[{\texttt{configure}, \texttt{Makefile.in}}]  These files are typically only used by packages which have a
non\texttt{\symbol{45}}\textsf{GAP} component, e.g.{\nobreakspace}some C code (the files of which should be in the \texttt{src} directory). The \texttt{configure} and \texttt{Makefile.in} files of the \textsf{Example} package provide prototypes (or they may be created using the \textsf{PackageMaker} mentioned above). The \texttt{configure} file typically takes a path \mbox{\texttt{\mdseries\slshape path}} to the \textsf{GAP} root directory as argument and uses the value assigned to \texttt{GAParch} in the file \texttt{sysinfo.gap}, created when \textsf{GAP} was compiled to determine the compilation architecture, inserts this in place
of the string \texttt{@GAPARCH@} in \texttt{Makefile.in} and creates a file \texttt{Makefile}. When \texttt{make} is run (which, of course, reads the constructed \texttt{Makefile}), a directory \texttt{bin} (if necessary) and subdirectories of \texttt{bin} with the path equal to the string assigned to \texttt{GAParch} in the file \texttt{sysinfo.gap} should be created; any binaries constructed by compiling the code in \texttt{src} should end up in this subdirectory of \texttt{bin}.  
\item[{\texttt{PackageInfo.g}}]  \index{PackageInfo.g@\texttt{PackageInfo.g}!for a GAP package} Every \textsf{GAP} package \emph{must} have a \texttt{PackageInfo.g} file which contains meta\texttt{\symbol{45}}information about the package
(package name, version, author(s), relations to other packages, homepage,
download archives, etc.). This information is used by the package loading
mechanism and also for the redistribution of a package with \textsf{GAP}. The \textsf{Example} package's \texttt{PackageInfo.g} file is well\texttt{\symbol{45}}commented and can be used as a prototype (see
also \ref{The PackageInfo.g File} for further details). It may also be created using the \textsf{PackageMaker} mentioned above. 
\item[{\texttt{init.g}, \texttt{read.g}}]  \index{init.g@\texttt{init.g}!for a GAP package} \index{read.g@\texttt{read.g}!for a GAP package} A \textsf{GAP} package \emph{must} have a file \texttt{init.g}. Typical \texttt{init.g} and \texttt{read.g} files should normally consist entirely of \texttt{ReadPackage} (\ref{ReadPackage}) commands (and possibly also \texttt{Read} (\ref{Read}) commands) for reading further files of the package. If the ``declaration'' and ``implementation'' parts of the package are separated (and this is recommended), there should be
a \texttt{read.g} file. The ``declaration'' part of a package consists of function and variable \emph{name} declarations and these go in files with \texttt{.gd} extensions; these files are read in via \texttt{ReadPackage} commands in the \texttt{init.g} file. The ``implementation'' part of a package consists of the actual definitions of the functions and
variables whose names were declared in the ``declaration'' part, and these go in files with \texttt{.gi} extensions; these files are read in via \texttt{ReadPackage} commands in the \texttt{read.g} file. The reason for following the above dichotomy is that the \texttt{read.g} file is read \emph{after} the \texttt{init.g} file, thus enabling the possibility of a function's implementation to refer to
another function whose name is known but is not actually defined yet (see \ref{Declaration and Implementation Part of a Package} below for more details). 

 The \textsf{GAP} code (whether or not it is split into ``declaration'' and ``implementation'' parts) should go in the package's \texttt{lib} directory (see below). 
\item[{\texttt{doc}}]  \index{GAPDoc@GAPDoc format!for writing package documentation} This directory should contain the package's documentation, written in an
XML\texttt{\symbol{45}}based documentation format supported by the \textsf{GAP} package \textsf{GAPDoc} (see  (\textbf{GAPDoc: Introduction and Example})) which is used for the \textsf{GAP} documentation itself. 

 The \textsf{Example} package's documentation (see its \texttt{doc} directory) may be used as a prototype. It consists of the master file \texttt{main.xml}, further \texttt{.xml} files for manual chapters (included in the manual via \texttt{Include} directives in the master file) and the \textsf{GAP} input file \texttt{../makedocrel.g} which generates the manuals. Generally, one should also provide a \texttt{manual.bib} Bib{\TeX} database file or an \texttt{xml} file in the BibXMLext format (see  (\textbf{GAPDoc: The BibXMLext Format})). 

  One could also use the \textsf{AutoDoc} which simplifies writing documentation by generating most of the \textsf{GAPDoc} code automatically. 
\item[{\texttt{lib}}]  This is the preferred place for the \textsf{GAP} code of the package, i.e.{\nobreakspace}the \texttt{.g}, \texttt{.gd} and \texttt{.gi} files (other than \texttt{PackageInfo.g}, \texttt{init.g} and \texttt{read.g}). For some packages, the directory \texttt{gap} has been used instead of \texttt{lib}; \texttt{lib} has the advantage that it is the default subdirectory of a package directory
searched for by the \texttt{DirectoriesPackageLibrary} (\ref{DirectoriesPackageLibrary}) command. 
\item[{\texttt{src}}]  If the package contains non\texttt{\symbol{45}}\textsf{GAP} code, e.g.{\nobreakspace}C code, then this source code should go in the \texttt{src} directory. If there are \texttt{.h} ``include'' files you may prefer to put these all together in a separate \texttt{include} directory. There is one further rule for the location of kernel library
modules or external programs which is explained in \ref{Installation of GAP Package Binaries} below. 
\item[{\texttt{tst}}]  It is highly recommended that a package should have test files, which then
should go in the \texttt{tst} directory. For a deposited package, a test file with a basic test of the
package (for example, to check that it works as expected and/or that the
manual examples are correct) may be specified in the \texttt{PackageInfo.g} to be included in the \textsf{GAP} standard test suite and run as a part of the \textsf{GAP} release preparation. More specific and time consuming tests are not supposed
to be a part of the \textsf{GAP} standard test suite but may be placed in the \texttt{tst} directory with further instructions on how to run them. See Section \ref{Testing a GAP package} about the requirements to the test files formats and further recommendations. 
\end{description}
 All other files can be organised as you like. But we suggest that you have a
look at existing packages and use a similar scheme, for example, put examples
in the \texttt{examples} subdirectory, data libraries in extra subdirectories, and so on. 

 Sometimes there may be a need to include an empty directory in the package
distribution (for example, as a place to store some data that may appear at
runtime). In this case package authors are advised to put in this directory a
short \texttt{README} file describing its purpose to ensure that such directory will be included in
the redistribution. 

 Concerning the \textsf{GAP} code in packages, it is recommended to use only documented \textsf{GAP} functions, see \ref{Undocumented Variables}. In particular if you want to make your package available to other \textsf{GAP} users it is advisable to avoid using ``obsolete'' variables (see \ref{Replaced and Removed Command Names}). To test that the package does not use obsolete variables you can set the \texttt{ReadObsolete} component in your \texttt{gap.ini} file to \texttt{false} (see \ref{sect:gap.ini}) or start \textsf{GAP} with \texttt{\texttt{\symbol{45}}A \texttt{\symbol{45}}O} command line options (note that this may also cause problems with loading
other packages that use ``obsolete'' variables). }

  
\section{\textcolor{Chapter }{Writing Documentation and Tools Needed}}\label{Writing Documentation and Tools Needed}
\logpage{[ 76, 6, 0 ]}
\hyperdef{L}{X84164AA2859A195F}{}
{
  If you intend to make your package available to other users it is essential to
include documentation explaining how to install and use your programs. 

 Concerning the installation you should produce a \texttt{README} file which gives a short description of the purpose of the package and
contains proper instructions how to install your package. Again, check out
some existing packages to get an idea how this could look like. 

 Documentation for \textsf{GAP} package should be prepared in an XML\texttt{\symbol{45}}based documentation
format that is defined in and can be used with the \textsf{GAPDoc} package (see{\nobreakspace} (\textbf{GAPDoc: Introduction and Example})). 

 There should be at least a text version of your documentation provided for use
in the terminal running \textsf{GAP} and some nicely printable version in \texttt{.pdf} format. Many \textsf{GAP} users like to browse the documentation in HTML format via their web browser.
As a package author, you are not obliged to provide an HTML version of your
package manual,  but if you use the \textsf{GAPDoc} package you should have no trouble in producing one. 

 Moreover, using the \textsf{GAPDoc} package, it is also possible to produce HTML version of the documentation
supporting MathJax (\href{https://www.mathjax.org/} {\texttt{https://www.mathjax.org/}}) for the high quality rendering of mathematical symbols while viewing it
online. For example, if you are viewing the HTML version of the manual,
compare how this formula will look with MathJax turned on/off: 
\[ [ \chi, \psi ] = \left( \sum_{{g \in G}} \chi(g) \psi(g^{{-1}}) \right) / |G|. \]
 

 The manual of the \textsf{Example} package is written in the \textsf{GAPDoc} format, and commands needed to build it are contained in the file \texttt{makedocrel.g} (you don't need to re\texttt{\symbol{45}}build the manual since it is already
included in the package). You will also need to have certain {\TeX} tools installed: to produce manuals in the \texttt{.pdf} format, you need \texttt{pdflatex}. 

 In principle it is also possible to use alternative documentation systems.
Historically, there is one such {\TeX}\texttt{\symbol{45}}based system, which predates \textsf{GAPDoc}, and which is still in use by several packages. However, we do not recommend
using it for new packages. }

  
\section{\textcolor{Chapter }{An Example of a \textsf{GAP} Package}}\label{An Example of a GAP Package}
\logpage{[ 76, 7, 0 ]}
\hyperdef{L}{X79AB306684AC8E7A}{}
{
  We illustrate the creation of a \textsf{GAP} package by an example of a very basic package. 

 Create the following directories in your home directory: \texttt{.gap}, \texttt{.gap/pkg} and \texttt{.gap/pkg/test}. Then inside the directory \texttt{.gap/pkg/test} create an empty file \texttt{init.g}, and a file \texttt{PackageInfo.g} with the following contents: 

 
\begin{Verbatim}[commandchars=!@|,fontsize=\small,frame=single,label=Example]
  SetPackageInfo( rec(
    PackageName := "test",
    Version := "1.0",
    PackageDoc := rec(
        BookName  := "test",
        SixFile   := "doc/manual.six",
    ),
    Dependencies := rec(
        GAP       := "4.9",
        NeededOtherPackages := [ ["GAPDoc", "1.6"] ],
        SuggestedOtherPackages := [ ] ),
    AvailabilityTest := ReturnTrue ) );
\end{Verbatim}
 

 This file declares the \textsf{GAP} package with name ``test'' in version 1.0. The package documentation consists of one autoloaded book; the \texttt{SixFile} component is needed by the \textsf{GAP} help system. Package dependencies (picked for the purposes of this example)
require at least \textsf{GAP}{\nobreakspace}4.9 and \textsf{GAPDoc} package at version at least 1.6, and these conditions will be checked when the
package will be loaded (see \ref{Version Numbers}). Since there are no requirements that have to be tested, \texttt{AvailabilityTest} just uses \texttt{ReturnTrue} (\ref{ReturnTrue}). 

 Now start \textsf{GAP} (without using the \texttt{\texttt{\symbol{45}}r} option) and the \texttt{.gap} directory will be added to the \textsf{GAP} root directory to allow \textsf{GAP} to find the packages installed there (see \ref{GAP Root Directories}). 

 
\begin{Verbatim}[commandchars=!@|,fontsize=\small,frame=single,label=Example]
  !gapprompt@gap>| !gapinput@LoadPackage("test");|
  true
\end{Verbatim}
 

 This \textsf{GAP} package is too simple to be useful, but we have succeeded in loading it via \texttt{LoadPackage} (\ref{LoadPackage}), satisfying all specified dependencies. }

  
\section{\textcolor{Chapter }{File Structure}}\label{File Structure}
\logpage{[ 76, 8, 0 ]}
\hyperdef{L}{X7A61B1AE7D632E01}{}
{
  Package files may follow the style used for the \textsf{GAP} library. Every file in the \textsf{GAP} library starts with a header that lists the filename, copyright, a short
description of the file contents and the original authors of this file, and
ends with a comment line \texttt{\#E}. Indentation in functions and the use of decorative spaces in the code are
left to the decision of the authors of each file. Global (i.e.
re\texttt{\symbol{45}}used elsewhere) comments usually are indented by two
hash marks and two blanks, in particular, every declaration or method or
function installation which is not only of local scope is separated by a
header. 

 Facilities to distribute a document over several files to allow the
documentation for parts of some code to be stored in the same file as the code
itself are provided by the \textsf{GAPDoc} package (see  (\textbf{GAPDoc: Distributing a Document into Several Files})). The same approach is demonstrated by the \textsf{Example} package. E.g. \texttt{example/doc/example.xml} has the statement \texttt{{\textless}\#Include Label="ListDirectory"{\textgreater}} and \texttt{example/lib/files.gd} contains 
\begin{Verbatim}[commandchars=!@|,fontsize=\small,frame=single,label=Example]
  ##  <#GAPDoc Label="ListDirectory">
  ##  <ManSection>
  ##  <Func Name="ListDirectory" Arg="[dir]"/>
  ##
  ##  <Description>
  ##  lists the files in directory <A>dir</A> (a string)
  ##  or the current directory if called with no arguments.
  ##  </Description>
  ##  </ManSection>
  ##  <#/GAPDoc>
  DeclareGlobalFunction( "ListDirectory" );
\end{Verbatim}
 This is all put together in the file \texttt{example/makedocrel.g} which builds the package documentation, calling \texttt{MakeGAPDocDoc} (\textbf{GAPDoc: MakeGAPDocDoc}) with locations of library files containing parts of the documentation. 

 Alternatively, one could use the \textsf{AutoDoc}, which simplifies writing documentation by generating most of the \textsf{GAPDoc} code automatically. The equivalent of the fragment of the code above for \textsf{AutoDoc} would look like 
\begin{Verbatim}[commandchars=|BC,fontsize=\small,frame=single,label=Example]
  #! @Arguments [dir]
  #! @Description
  #!  lists the files in directory <A>dir</A> (a string)
  #!  or the current directory if called with no arguments.
  DeclareGlobalFunction( "ListDirectory" );
\end{Verbatim}
 }

  
\section{\textcolor{Chapter }{Creating the PackageInfo.g File}}\label{Creating the PackageInfo.g File}
\logpage{[ 76, 9, 0 ]}
\hyperdef{L}{X7A09C63685065B01}{}
{
  \index{ValidatePackageInfo@\texttt{ValidatePackageInfo}} While the minimalistic \texttt{PackageInfo.g} file described in \ref{An Example of a GAP Package} is enough to let \textsf{GAP} load the package, and check all specified dependencies, it is actually missing
many extra fields which become relevant if you want to distribute your
package: they contain lists of authors and/or maintainers including contact
information, URLs of the package archives and README files, status
information, text for a package overview webpage, and so on. All these details
are required for a package to be redistributed with \textsf{GAP}. 

 The command \texttt{ValidatePackageInfo} (\ref{ValidatePackageInfo}) can be used to get a quick idea about which fields are missing: 
\begin{Verbatim}[commandchars=!@|,fontsize=\small,frame=single,label=Example]
  !gapprompt@gap>| !gapinput@ValidatePackageInfo("PackageInfo.g");|
  #E  component `Subtitle' must be bound to a string
  #E  component `Date' must be bound to a string of the form `dd/mm/yyyy'
  #E  component `ArchiveURL' must be bound to a string started with http://, https:// or ftp://
  #E  component `ArchiveFormats' must be bound to a string
  #E  component `README_URL' must be bound to a string started with http://, https:// or ftp://
  #E  component `PackageInfoURL' must be bound to a string started with http://, https:// or ftp://
  #E  component `AbstractHTML' must be bound to a string
  #E  component `PackageWWWHome' must be bound to a string started with http://, https:// or ftp://
  #E  component `ArchiveURLSubset' must be bound to a list of strings denoting relative paths to readable files or directories
  #E  component `HTMLStart' must be bound to a string denoting a relative path to a readable file
  #E  component `PDFFile' must be bound to a string denoting a relative path to a readable file
  #E  component `SixFile' must be bound to a string denoting a relative path to a readable file
  #E  component `LongTitle' must be bound to a string
  false
\end{Verbatim}
 

 We suggest to create a \texttt{PackageInfo.g} file for your package by copying the one in the \textsf{Example} package, distributed with \textsf{GAP}, or using the \textsf{PackageMaker} (\href{https://github.com/gap-system/PackageMaker} {\texttt{https://github.com/gap\texttt{\symbol{45}}system/PackageMaker}}), and then adjusting it for your package. Within \textsf{GAP} you can look at this template file for a list and explanation of all
recognised entries by 
\begin{Verbatim}[commandchars=!@|,fontsize=\small,frame=single,label=Example]
  Pager(StringFile(Filename(DirectoriesLibrary(),
                            "../pkg/example/PackageInfo.g")));
\end{Verbatim}
 

 Instead of populating the rest of the \texttt{PackageInfo.g} by hands, you can also create a basic \textsf{GAP} package with the help of the tool called \textsf{PackageMaker}, available at \href{https://github.com/gap-system/PackageMaker} {\texttt{https://github.com/gap\texttt{\symbol{45}}system/PackageMaker}}. The \textsf{PackageMaker} asks several questions about the intended package and then creates a new
directory for it and populates it with all the files needed for a basic
package. }

  
\section{\textcolor{Chapter }{Functions and Variables and Choices of Their Names}}\label{Functions and Variables and Choices of Their Names}
\logpage{[ 76, 10, 0 ]}
\hyperdef{L}{X7DEACD9786DE29F1}{}
{
  In writing the \textsf{GAP} code for your package you need to be a little careful on just how you define
your functions and variables. 

 \emph{Firstly}, in general one should avoid defining functions and variables via assignment
statements in the way you would interactively, e.g. 

 
\begin{Verbatim}[commandchars=!@|,fontsize=\small,frame=single,label=Example]
  !gapprompt@gap>| !gapinput@Squared := x -> x^2;;|
  !gapprompt@gap>| !gapinput@Cubed := function(x) return x^3; end;;|
\end{Verbatim}
 

 The reason for this is that such functions and variables are \emph{easily overwritten} and what's more you are not warned about it when it happens. 

 To protect a function or variable against overwriting there is the function \texttt{BindGlobal} (\ref{BindGlobal}), or alternatively (and equivalently) you may define a global function via a \texttt{DeclareGlobalFunction} (\ref{DeclareGlobalFunction}) and \texttt{InstallGlobalFunction} (\ref{InstallGlobalFunction}) pair or a global variable via a \texttt{DeclareGlobalVariable} (\ref{DeclareGlobalVariable}) and \texttt{InstallValue} (\ref{InstallValue}) pair. There are also operations and their methods, and related objects like
attributes and filters which also have \texttt{Declare...} and \texttt{Install...} pairs. 

 \emph{Secondly}, it is a good idea to reduce the chance of accidental overwriting by choosing
names for your functions and variables that begin with a string that
identifies it with the package, e.g.{\nobreakspace}some of the undocumented
functions in the \textsf{Example} package begin with \texttt{Eg}. This is especially important in cases where you actually want the user to be
able to change the value of a function or variable defined by your package,
for which you have used direct assignments (for which the user will receive no
warning if she accidentally overwrites them). It is also important for
functions and variables defined via \texttt{BindGlobal}, \texttt{DeclareGlobalFunction}/\texttt{InstallGlobalFunction} and \texttt{DeclareGlobalVariable}/\texttt{InstallValue}, in order to avoid name clashes that may occur with (extensions of) the \textsf{GAP} library and other packages. 

 \index{local namespace!for a GAP package} Additionally, since \textsf{GAP}{\nobreakspace}4.5 a package may place global variables into a local namespace
as explained in \ref{Namespaces} in order to avoid name clashes and preserve compatibility. This new feature
allows you to define in your package global variables with the identifier
ending with the \texttt{@} symbol, e.g. \texttt{xYz@}. Such variables may be used in your package code safely, as they may be
accessed from outside the package only by their full name, i.e. \texttt{xYz@YourPackageName}. This helps to prevent clashes between different packages or between a
package and the \textsf{GAP} library because of the same variable names.  

 On the other hand, operations and their methods (defined via \texttt{DeclareOperation} (\ref{DeclareOperation}), \texttt{InstallMethod} (\ref{InstallMethod}) etc.{\nobreakspace}pairs) and their relatives do not need this consideration,
as they avoid name clashes by allowing for more than one ``method'' for the same\texttt{\symbol{45}}named object. 

 To demonstrate the definition of a function via a \texttt{DeclareOperation}/\texttt{InstallMethod} pair, the method \texttt{Recipe} (\textbf{Example: Recipe}) was included in the \textsf{Example} package; \texttt{Recipe( FruitCake );} gives a ``method'' for making a fruit cake (forgive the pun). 

 \emph{Thirdly}, functions or variables with \texttt{Set\mbox{\texttt{\mdseries\slshape XXX}}} or \texttt{Has\mbox{\texttt{\mdseries\slshape XXX}}} names (even if they are defined as operations) should be avoided as these may
clash with objects associated with attributes or properties (attributes and
properties \mbox{\texttt{\mdseries\slshape XXX}} declared via the \texttt{DeclareAttribute} and \texttt{DeclareProperty} commands have associated with them testers of form \texttt{Has\mbox{\texttt{\mdseries\slshape XXX}}} and setters of form \texttt{Set\mbox{\texttt{\mdseries\slshape XXX}}}). 

 \emph{Fourthly}, it is a good idea to have some convention for internal functions and
variables (i.e.{\nobreakspace}the functions and variables you don't intend for
the user to use). For example, they might be entirely CAPITALISED. 

 Additionally, there is a recommended naming convention that the \textsf{GAP} core system and \textsf{GAP} packages should not use global variables starting in the lowercase. This
allows to reserve variables with names starting in lowercase to the \textsf{GAP} user so they will never clash with the system. It is extremely important to
avoid using for package global variables very short names started in
lowercase. For example, such names like \texttt{cs}, \texttt{exp}, \texttt{ngens}, \texttt{pc}, \texttt{pow} which are perfectly fine for local variables, should never be used for
globals. Additionally, the package must not have writable global variables
with very short names even if they are starting in uppercase, for example, \texttt{C1} or \texttt{ORB}, since they also could be easily overwritten by the user. 

 It is a good practice to follow naming conventions used in \textsf{GAP} as explained in \ref{Naming Conventions} and  (\textbf{Tutorial: Changing the Structure}), which might help users to memorize or even guess names of functions provided
by the package. 

 \emph{Finally}, note the advantage of using \texttt{DeclareGlobalFunction}/\texttt{InstallGlobalFunction}, \texttt{DeclareGlobalVariable}/\texttt{InstallValue}, etc.{\nobreakspace}pairs (rather than \texttt{BindGlobal}) to define functions and variables, which allow the package author to
organise her function\texttt{\symbol{45}} and variable\texttt{\symbol{45}}
definitions in any order without worrying about any interdependence. The \texttt{Declare...} statements should go in files with \texttt{.gd} extensions and be loaded by \texttt{ReadPackage} statements in the package \texttt{init.g} file, and the \texttt{Install...} definitions should go in files with \texttt{.gi} extensions and be loaded by \texttt{ReadPackage} statements in the package \texttt{read.g} file; this ensures that the \texttt{.gi} files are read \emph{after} the \texttt{.gd} files. All other package code should go in \texttt{.g} files (other than the \texttt{init.g} and \texttt{read.g} files themselves) and be loaded via \texttt{ReadPackage} statements in the \texttt{init.g} file. 

 \index{ShowPackageVariables@\texttt{ShowPackageVariables}} In conclusion, here is some practical advice on how to check which variables
are used by the package. 

 Firstly, there is a function \texttt{ShowPackageVariables} (\ref{ShowPackageVariables}). If the package \mbox{\texttt{\mdseries\slshape pkgname}} is available but not yet loaded then \texttt{ShowPackageVariables( pkgname )} prints a list of global variables that become bound and of methods that become
installed when the package is loaded (for that, the package will be actually
loaded, so \texttt{ShowPackageVariables} can be called only once for the same package in the same \textsf{GAP} session.) The second optional argument \mbox{\texttt{\mdseries\slshape version}} may specify a particular package version to be loaded. An error message will
be printed if (the given version of) the package is not available or already
loaded. 

 Info lines for undocumented variables will be marked with an asterisk \texttt{*}. Note that the \textsf{GAP} help system is case insensitive, so it is difficult to document identifiers
that differ only by case. 

 The following entries are omitted from the list: default setter methods for
attributes and properties that are declared in the package, and \texttt{Set\mbox{\texttt{\mdseries\slshape attr}}} and \texttt{Has\mbox{\texttt{\mdseries\slshape attr}}} type variables where \mbox{\texttt{\mdseries\slshape attr}} is an attribute or property. 

 For example, for the \textsf{Example} package it may produce the output looking like this: 
\begin{Verbatim}[commandchars=!@|,fontsize=\small,frame=single,label=Example]
  !gapprompt@gap>| !gapinput@ShowPackageVariables("example");|
  ----------------------------------------------------------------
  Loading Example 3.3 (Example/Template of a GAP Package)
  by Werner Nickel (http://www.mathematik.tu-darmstadt.de/~nickel),
     Greg Gamble (http://www.math.rwth-aachen.de/~Greg.Gamble), and
     Alexander Konovalov (http://www.cs.st-andrews.ac.uk/~alexk/).
  ----------------------------------------------------------------
  new global functions:
    EgSeparatedString( str, c )*
    FindFile( dir, file )
    HelloWorld(  )
    ListDirectory( arg )
    LoadedPackages(  )
    WhereIsPkgProgram( prg )
    Which( prg )
  
  new global variables:
    FruitCake
  
  new operations:
    Recipe( arg )
  
  new methods:
    Recipe( cake )
\end{Verbatim}
 Another trick is to start \textsf{GAP} with \texttt{\texttt{\symbol{45}}r \texttt{\symbol{45}}A} options, immediately load your package and then call \texttt{NamesUserGVars} (\ref{NamesUserGVars}) which returns a list of the global variable names created since the library
was read, to which a value is currently bound. For example, for the \textsf{Example} it produces 
\begin{Verbatim}[commandchars=!@|,fontsize=\small,frame=single,label=Example]
  !gapprompt@gap>| !gapinput@NamesUserGVars();|
  [ "EgSeparatedString", "FindFile", "FruitCake", "HelloWorld", "ListDirectory",
    "LoadedPackages", "Recipe", "WhereIsPkgProgram", "Which" ]
\end{Verbatim}
 but for packages with dependencies it will also contain variables created by
other packages. Nevertheless, it may be a useful check to search for unwanted
variables appearing after package loading. A potentially dangerous situation
which should be avoided is when the package uses some simply named temporary
variables at the loading stage. Such ``phantom'' variables may then remain unnoticed and, as a result, there will be no
warnings if the user occasionally uses the same name as a local variable name
in a function. Even more dangerous is the case when the user variable with the
same name already exists before the package is loaded so it will be silently
overwritten. }

  
\section{\textcolor{Chapter }{Package Dependencies (Requesting one \textsf{GAP} Package from within Another)}}\label{Package dependencies}
\logpage{[ 76, 11, 0 ]}
\hyperdef{L}{X7928799186F9B2FE}{}
{
  \index{needed package} \index{suggested package} \index{dependencies!for a GAP package} It is possible for one \textsf{GAP} package \texttt{A} to require another package \texttt{B}. For that, one simply adds the name and the (least) version number of the
package \texttt{B} to the \texttt{NeededOtherPackages} component of the \texttt{Dependencies} component of the \texttt{PackageInfo.g} file of the package \texttt{A}. In this situation, loading the package \texttt{A} forces that also the package \texttt{B} is loaded, and that \texttt{A} cannot be loaded if \texttt{B} is not available. 

 If \texttt{B} is not essential for \texttt{A} but should be loaded if it is available (for example because \texttt{B} provides some improvements of the main system that are useful for \texttt{A}) then the name and the (least) version number of \texttt{B} should be added to the \texttt{SuggestedOtherPackages} component of the \texttt{Dependencies} component of the \texttt{PackageInfo.g} file of \texttt{A}. In this situation, loading \texttt{A} forces an attempt to load also \texttt{B}, but \texttt{A} is loaded even if \texttt{B} is not available. 

 If \texttt{B} is not essential for \texttt{A} but is used in the package tests, (for example because \texttt{B} is a library of groups used to provide testcases) then the name and the
(least) version number of \texttt{B} should be added to the \texttt{TestPackages} component of the \texttt{Dependencies} component of the \texttt{PackageInfo.g} file of \texttt{A}. Note that \textsf{GAP} itself does not use this information at all. In particular, no attempts are
made to actually load \texttt{B}, this should be done explicitly as part of the package tests. The benefit of
specifying this component is for automated test runners (e.g. as part of the \textsf{GAP} package distribution) which can use this information to ensure all packages
needed to run your package's tests are installed and ready. 

 All package dependencies must be documented explicitly in the \texttt{PackageInfo.g} file. It is important to properly identify package dependencies and make the
right decision whether the other package should be ``needed'' or ``suggested''. For example, declaring package as ``needed'' when ``suggested'' might be sufficient may prevent loading of packages under Windows for no good
reason. 

 It is not appropriate to explicitly call \texttt{LoadPackage} (\ref{LoadPackage}) \emph{when the package is loaded}, since this may distort the order of package loading and result in warning
messages. It is recommended to turn such dependencies into needed or suggested
packages. For example, a package can be designed in such a way that it can be
loaded with restricted functionality if another package (or standalone
program) is missing, and in this case the missing package (or binary) is \emph{suggested}. Alternatively, if the package author decides that loading the package in
this situation makes no sense, then the missing component is \emph{needed}. 

 On the other hand, if \texttt{LoadPackage} (\ref{LoadPackage}) is called inside functions of the package then there is no such problem,
provided that these functions are called only after the package has been
loaded, so it is not necessary to specify the other package as suggested. The
same applies to test files and manual examples, which may be simply extended
by calls to \texttt{LoadPackage} (\ref{LoadPackage}). 

 \index{OnlyNeeded@\texttt{OnlyNeeded}!option} It may happen that a package B that is listed as a suggested package of
package A is actually needed by A. If no explicit \texttt{LoadPackage} (\ref{LoadPackage}) calls for B occur in A at loading time, this can now be detected using the new
possibility to load a package without loading its suggested packages using the
global option \texttt{OnlyNeeded} which can be used to (recursively) suppress loading the suggested packages of
the package in question. Using this option, one can check whether errors or
warnings appear when B is not available (note that this option should be used
only for such checks to simulate the situation when package B is not
available; it is not supposed to be used in an actual \textsf{GAP} session when package B will be loaded later, since this may cause problems).
In case of any errors or warnings, their consequence can then be either
turning B into a needed package or (since apparently B was not intended to
become a needed package) changing the code accordingly. Only if package A
calls \texttt{LoadPackage} (\ref{LoadPackage}) for B at loading time (see above) then package B needs to be \emph{deinstalled} (i.e. removed) to test loading of A without B.  

 }

  
\section{\textcolor{Chapter }{External Dependencies (System packages needed by a \textsf{GAP} Package)}}\label{External dependencies}
\logpage{[ 76, 12, 0 ]}
\hyperdef{L}{X7DF775BF859B29D2}{}
{
  \index{external dependency} \index{system package} \index{dependencies!external} It is possible for a \textsf{GAP} package to require external software. If this software is not bundled together
with the package, then this dependency should be documented. To do so, one may
add either a string describing the dependency, or a list \texttt{[ description, URL ]}, to the \texttt{ExternalConditions} component of the \texttt{Dependencies} component of the \texttt{PackageInfo.g} file of the \textsf{GAP} package. 
\begin{Verbatim}[commandchars=!@|,fontsize=\small,frame=single,label=Example]
  ExternalConditions := [
      [ "needs the PARI/GP computer algebra system Version 2.5 or higher",
        "https://pari.math.u-bordeaux.fr/" ]
  ]
\end{Verbatim}
 Note that the data in \texttt{ExternalConditions} is purely informational and meant to be human\texttt{\symbol{45}}readable, not
machine\texttt{\symbol{45}}readable. \textsf{GAP} itself does not use it, but for packages distributed with \textsf{GAP}, the \textsf{GAP} website displays it in its package list. 

 If this dependency is available through a package manager such as \texttt{apt} or \texttt{brew}, this should also be documented. This can be done by adding the system
package to the relevant component (such as \texttt{Ubuntu} or \texttt{Homebrew}) of the \texttt{NeededSystemPackages} component of the \texttt{PackageInfo.g} file of the \textsf{GAP} package. Each system package is added by means of a list whose first component
is the name needed by the package manager. The other elements of this list are
reserved for future usage. 
\begin{Verbatim}[commandchars=!@|,fontsize=\small,frame=single,label=Example]
  NeededSystemPackages := rec(
      Ubuntu   := [["pari-gp"]],
      Homebrew := [["pari"]]
  )
\end{Verbatim}
 Note that this mechanism is primarily intended for consumption by the \textsf{GAP} package distribution and by other testing tooling. It is deliberately kept
very simple. For example, at this point there is no mechanism to prescribe a
distribution version (basically, "Ubuntu" means whatever
`ubuntu\texttt{\symbol{45}}latest` is referring to in GitHub actions), nor to
specify a minimal version of a system package. Despite these strong
limitations, this data may still be useful for people packaging \textsf{GAP} and its packages for a distribution, as it at least gives a hint which system
packages may be needed. }

  
\section{\textcolor{Chapter }{Extensions Provided by a Package}}\label{Extensions Provided by a Package}
\logpage{[ 76, 13, 0 ]}
\hyperdef{L}{X783A5F3D87F9AF78}{}
{
  Sometimes a package \texttt{A} can provide additional functionality, such as better methods or additional
data, if some other packages \texttt{B}, \texttt{C}, etc. are loaded. However, one would like package \texttt{A} to still be usable without these additional packages, and therefore \texttt{B}, \texttt{C}, etc. shall not be regarded as needed packages (see Section \ref{Package dependencies}) of \texttt{A}. 

 One way to deal with this situation is to put those parts of code of \texttt{A} that depend on \texttt{B}, \texttt{C}, etc., into files that get read only in the situation that the packages in
question have actually been loaded into the current \textsf{GAP} session. 

 However, this leaves the question when to load these files of a conditional \emph{extension} of \texttt{A}. In the past, the only option for \texttt{A} was to check for the presence of \texttt{B}, \texttt{C}, etc., while it itself was being loaded. With this setup, it depends on the
order in which packages get loaded whether some feature is available or not:
If \texttt{B} is loaded before \texttt{A}, the extension might be loaded as well; if \texttt{B} is loaded only after \texttt{A}, then the extension is not loaded. 

 To deal with this issue of conditional extensions of packages, \textsf{GAP} offers a dedicated mechanism: The \texttt{Extensions} component of the \texttt{PackageInfo.g} file of \texttt{A} is a list of declarations of conditional extension of \texttt{A}, each being a record with the following components. 
\begin{description}
\item[{\texttt{needed}}]  a list of the form \texttt{[ [ pkgname1, version1 ], [ pkgname2, version2 ], }$\ldots$\texttt{ ]}, meaning that the extension shall be loaded as soon as all packages \texttt{pkgname1}, \texttt{pkgname2}, $\ldots$, with versions (at least) \texttt{version1}, \texttt{version2}, $\ldots$, have been loaded, and 
\item[{\texttt{filename}}]  the path, relative to the package directory of \texttt{A}, of a file such that reading this file will load the code of the extension. 
\end{description}
 

 As an example suppose the following is part of the \texttt{PackageInfo.g}. Then \textsf{GAP} will load the file \texttt{fileForB.gd} as soon as package \texttt{B} is loaded in version 0.6 or newer, and \texttt{fileForCD.gi} once package \texttt{C} and \texttt{D} are loaded in version 1.2 and 0.1 or newer respectively. 
\begin{Verbatim}[commandchars=!@|,fontsize=\small,frame=single,label=Example]
  Extensions := [
    rec(
      needed := [ ["B", "0.6"] ],
      filename := "gap/fileForB.gd",
    ),
    rec(
      needed := [ ["C", "1.2"] , ["D", "0.1"] ],
      filename := "gap/fileForCD.gi",
    )
  ],
\end{Verbatim}
 Whenever \texttt{LoadPackage} (\ref{LoadPackage}) is called, \textsf{GAP} checks for package extensions whose conditions now are satisfied, and loads
them. 

 For example, package \texttt{A} can be loaded early in a \textsf{GAP} session, and declare in its \texttt{PackageInfo.g} the availability of an extension that requires package \texttt{B}. If \texttt{B} has not yet been loaded then this extension will not be loaded together with \texttt{A}. However, as soon as \texttt{B} gets (installed and) loaded later in the session, also the extension of \texttt{A} will automatically get loaded. 

 The contents of \texttt{Extensions} in a \texttt{PackageInfo.g} file does not affect the lists of needed or suggested packages. If an
extension of \texttt{A} is beneficial for the functions of \texttt{A} then it makes sense to list the packages needed for the extension among the
suggested packages of \texttt{A}, but this may not be the case if the extension is beneficial only for the
functions of its needed packages. }

  
\section{\textcolor{Chapter }{Declaration and Implementation Part of a Package}}\label{Declaration and Implementation Part of a Package}
\logpage{[ 76, 14, 0 ]}
\hyperdef{L}{X7A7835A5797AF766}{}
{
  When \textsf{GAP} packages require each other in a circular way, a ``bootstrapping'' problem arises of defining functions before they are called. The same problem
occurs in the \textsf{GAP} library, and it is resolved there by separating declarations (which define
global variables such as filters and operations) and implementations (which
install global functions and methods) in different files. Any implementation
file may use global variables defined in any declaration file. \textsf{GAP} initially reads all declaration files (in the library they have a \texttt{.gd} suffix) and afterwards reads all implementation files (which have a \texttt{.gi} suffix). 

 Something similar is possible for \textsf{GAP} packages: if a file \texttt{read.g} exists in the home directory of the package, this file is read only \emph{after} all the \texttt{init.g} files of all (implicitly) required \textsf{GAP} packages are read. Thus one can separate declaration and implementation for a \textsf{GAP} package in the same way as is done for the \textsf{GAP} library, by creating a file \texttt{read.g}, restricting the \texttt{ReadPackage} (\ref{ReadPackage}) statements in \texttt{init.g} to only read those files of the package that provide declarations, and to read
the implementation files from \texttt{read.g}. 

 \emph{Examples:} 

 Suppose that there are two packages \texttt{A} and \texttt{B}, each with files \texttt{init.g} and \texttt{read.g}. 

 
\begin{itemize}
\item  If package \texttt{A} suggests or needs package \texttt{B} and package \texttt{B} does not need or suggest any other package then first \texttt{init.g} of \texttt{B} is read, then \texttt{read.g} of \texttt{B}, then \texttt{init.g} of \texttt{A}, then \texttt{read.g} of \texttt{A}. 
\item  If package \texttt{A} suggests or needs package \texttt{B} and package \texttt{B} (or a package that is suggested or needed by \texttt{B}) suggests or needs package \texttt{A} then first the files \texttt{init.g} of \texttt{A} and \texttt{B} are read (in an unspecified order) and then the files \texttt{read.g} of \texttt{A} and \texttt{B} (in the same order). 
\end{itemize}
 

 In general, when \textsf{GAP} is asked to load a package then first the dependencies between this packages
and its needed and suggested packages are inspected (recursively), and a list
of package sets is computed such that no cyclic dependencies occur between
different package sets and such that no package in any of the package sets
needs any package in later package sets. Then \textsf{GAP} runs through the package sets and reads for each set first all \texttt{init.g} files and then all \texttt{read.g} files of the packages in the set. (There is one exception from this rule:
Whenever packages are autoloaded before the implementation part of the \textsf{GAP} library is read, only the \texttt{init.g} files of the packages are read; as soon as the \textsf{GAP} library has been read, the \texttt{read.g} files of these packages are also read, and afterwards the above rule holds.) 

 \index{IsPackageMarkedForLoading@\texttt{IsPackageMarkedForLoading}} It can happen that some code of a package depends on the availability of
suggested packages, i.e., different initialisations are performed depending on
whether a suggested package will eventually be loaded or not. One can test
this condition with the function \texttt{IsPackageMarkedForLoading} (\ref{IsPackageMarkedForLoading}). In particular, one should \emph{not} call (and use the value returned by this call) the function \texttt{LoadPackage} (\ref{LoadPackage}) inside package code that is read during package loading. Note that for
debugging purposes loading suggested packages may have been deliberately
disabled via the global option \texttt{OnlyNeeded}. 

 Note that the separation of the \textsf{GAP} code of packages into declaration part and implementation part does in general \emph{not} allow one to actually \emph{call} functions from a package when the implementation part is read. For example, in
the case of a ``cyclic dependency'' as in the second example above, suppose that \texttt{B} provides a new function \texttt{f} or a new global record \texttt{r} which are declared in the declaration part of \texttt{B}. Then the code in the implementation part of \texttt{A} may contain calls to the functions defined in the declaration part of \texttt{B}. However, the implementation part of \texttt{A} may be read \emph{before} the implementation part of \texttt{B}. So one can in general not assume that during the loading of \texttt{A}, the function \texttt{f} can be called, or that one can access components of the record \texttt{r}. 

 If one wants to call the function \texttt{f} or to access components of the record \texttt{r} in the code of the package \texttt{A} then the problem is that it may be not possible to determine a cyclic
dependency between \texttt{A} and \texttt{B} from the packages \texttt{A} and \texttt{B} alone. A safe solution is then to design \texttt{A} in such a way that the code that calls \texttt{f} or accesses \texttt{r} belongs to \emph{package extensions} of \texttt{A} that get loaded only after \texttt{B} has been loaded; see Section \ref{Extensions Provided by a Package} for details. 

 \index{autoreadable variables} In the case of cyclic dependencies, one solution for the above problem might
be to delay those computations (typically initialisations) in package \texttt{A} that require package \texttt{B} to be loaded until all required packages are completely loaded. This can be
done by moving the declaration and implementation of the variables that are
created in the initialisation into a separate file and to declare these
variables in the \texttt{init.g} file of the package, via a call to \texttt{DeclareAutoreadableVariables} (\ref{DeclareAutoreadableVariables}) (see also \ref{Autoreadable Variables}). 

 }

  
\section{\textcolor{Chapter }{Autoreadable Variables}}\label{Autoreadable Variables}
\logpage{[ 76, 15, 0 ]}
\hyperdef{L}{X7D7F236A78106358}{}
{
  Package files containing method installations must be read when the package is
loaded. For package files \emph{not} containing method installations (this applies, for example, to many data
files) another mechanism allows one to delay reading such files until the data
are actually accessed. See \ref{DeclareAutoreadableVariables} for further details. }

  
\section{\textcolor{Chapter }{Standalone Programs in a \textsf{GAP} Package}}\label{Standalone Programs in a GAP Package}
\logpage{[ 76, 16, 0 ]}
\hyperdef{L}{X7C8CCF057806EFD0}{}
{
   \textsf{GAP} packages that involve stand\texttt{\symbol{45}}alone programs are
fundamentally different from \textsf{GAP} packages that consist entirely of \textsf{GAP} code. 

 This difference is threefold: A user who installs the \textsf{GAP} package must also compile (or install) the package's binaries, the package
must check whether the binaries are indeed available, and finally the \textsf{GAP} code of the package has to start the external binary and to communicate with
it. We will cover these three points in the following sections. 

 If the package does not solely consist of an interface to an external binary
and if the external program called is not just
special\texttt{\symbol{45}}purpose code, but a generally available program,
chances are high that sooner or later other \textsf{GAP} packages might also require this program. We therefore strongly recommend the
provision of a documented \textsf{GAP} function that will call the external binary. We also suggest to create
actually two \textsf{GAP} packages; the first providing only the binary and the interface and the second
(requiring the first, see{\nobreakspace}\ref{Package dependencies}) being the actual \textsf{GAP} package. 

  
\subsection{\textcolor{Chapter }{Installation of \textsf{GAP} Package Binaries}}\label{Installation of GAP Package Binaries}
\logpage{[ 76, 16, 1 ]}
\hyperdef{L}{X7CD9ED5C86725ACF}{}
{
  \index{sysinfo.gap@\texttt{sysinfo.gap}} \index{external binaries!for a GAP package} The scheme for the installation of package binaries which is described further
on is intended to permit the installation on different architectures which
share a common file system (and share the architecture independent file). 

 A \textsf{GAP} package which includes external binaries contains a \texttt{bin} subdirectory. This subdirectory in turn contains subdirectories for the
different architectures on which the \textsf{GAP} package binaries are installed. The names of these directories must be the
same as the names of the architecture dependent subdirectories of the main \texttt{bin} directory. Unless you use a tool like \texttt{autoconf} yourself, you must obtain the correct name of the binary directory from the
main \textsf{GAP} branch. To help with this, the main \textsf{GAP} directory contains a file \texttt{sysinfo.gap} which assigns the shell variable \texttt{GAParch} to the proper name as determined by \textsf{GAP}'s \texttt{configure} process. For example on a Linux system, the file \texttt{sysinfo.gap} may look like this: 

 
\begin{Verbatim}[commandchars=!@|,fontsize=\small,frame=single,label=Example]
  GAParch=i586-unknown-linux2.0.31-gcc
\end{Verbatim}
 

 We suggest that your \textsf{GAP} package contains a file \texttt{configure} which is called with the path of the \textsf{GAP} root directory as parameter. This file then will read \texttt{sysinfo.gap} and set up everything for compiling under the given architecture (for example
creating a \texttt{Makefile} from \texttt{Makefile.in}). As initial templates, you may use installation scripts of the \textsf{Example} package or files generated with the help of \textsf{PackageMaker}. }

  
\subsection{\textcolor{Chapter }{Test for the Existence of GAP Package Binaries}}\label{Test for the Existence of GAP Package Binaries}
\logpage{[ 76, 16, 2 ]}
\hyperdef{L}{X7E4F39867CCC6026}{}
{
  If an external binary is essential for the workings of a \textsf{GAP} package, the function stored in the component \texttt{AvailabilityTest} of the \texttt{PackageInfo.g} file of the package should test whether the program has been compiled on the
architecture (and inhibit package loading if this is not the case). This is
especially important if the package is loaded automatically. 

 The easiest way to accomplish this is to use \texttt{Filename} (\ref{Filename:for a directory and a string}) for checking for the actual binaries in the path given by \texttt{DirectoriesPackagePrograms} (\ref{DirectoriesPackagePrograms}) for the respective package. For example the \textsf{example} \textsf{GAP} package could use the following function to test whether the binary \texttt{hello} has been compiled; it will issue a warning if not, and will only load the
package if the binary is indeed available: 

 
\begin{Verbatim}[commandchars=!@|,fontsize=\small,frame=single,label=Example]
  ...
  AvailabilityTest := function()
    local path,file;
      # test for existence of the compiled binary
      path:= DirectoriesPackagePrograms( "example" );
      file:= Filename( path, "hello" );
      if file = fail then
        LogPackageLoadingMessage( PACKAGE_WARNING,
            [ "The program `hello' is not compiled,",
              "`HelloWorld()' is thus unavailable.",
              "See the installation instructions;",
              "type: ?Installing the Example package" ] );
      fi;
      return file <> fail;
    end,
  ...
\end{Verbatim}
 

 However, if you look at the actual \texttt{PackageInfo.g} file of the \textsf{example} package, you will see that its \texttt{AvailabilityTest} function always returns \texttt{true}, and just logs the warning if the binary is not available (which may be later
viewed with \texttt{DisplayPackageLoadingLog} (\ref{DisplayPackageLoadingLog})). This means that the binary is not regarded as essential for this package. 

 You might also have to cope with the situation that external binaries will
only run under UNIX (and not e.g. under Windows), or may not compile with some
compilers or default compiler options. See{\nobreakspace}\ref{Testing for the System Architecture} for information on how to test for the architecture. 

 Package using a kernel module (see \ref{Kernel modules}), one may use a test like this: 
\begin{Verbatim}[commandchars=!@|,fontsize=\small,frame=single,label=]
  ...
  AvailabilityTest := function()
      # see if example.so exists and is a loadable kernel extension
      if not IsKernelExtensionAvailable("example") then
        LogPackageLoadingMessage( PACKAGE_WARNING,
            [ "The kernel extension `example' is unavailable,",
              "perhaps it needs to be recompiled?",
              "See the installation instructions;",
              "type: ?Installing the Example package" ] );
        return false;
      fi;
      return true;
    end,
  ...
\end{Verbatim}
 

 \index{LogPackageLoadingMessage@\texttt{LogPackageLoadingMessage}} Last but not least: do not print anything in the \texttt{AvailabilityTest} function of the package via \texttt{Print} or \texttt{Info}. Instead one should call \texttt{LogPackageLoadingMessage} (\ref{LogPackageLoadingMessage}) to store a message which may be viewed later with \texttt{DisplayPackageLoadingLog} (\ref{DisplayPackageLoadingLog}) (the latter two functions have been introduced in \textsf{GAP}{\nobreakspace}4.5) }

  
\subsection{\textcolor{Chapter }{Calling of and Communication with External Binaries}}\label{Calling of and Communication with External Binaries}
\logpage{[ 76, 16, 3 ]}
\hyperdef{L}{X8438685184FCEFEC}{}
{
  There are two reasons for this: the input data has to be passed on to the
stand\texttt{\symbol{45}}alone program and the stand\texttt{\symbol{45}}alone
program has to be started from within \textsf{GAP}. 

 There are two principal ways of doing this. 

 The first possibility is to write all the data for the
stand\texttt{\symbol{45}}alone to one or several files, then start the
stand\texttt{\symbol{45}}alone with \texttt{Process} (\ref{Process}) or \texttt{Exec} (\ref{Exec}) which then writes the output data to file, and finally read in the
standalone's output file. 

 The second way is interfacing via input\texttt{\symbol{45}}output streams, see
Section{\nobreakspace}\ref{Input-Output Streams}. 

 Some \textsf{GAP} packages use kernel modules (see \ref{Kernel modules}) instead of external binaries. A kernel module is implemented in C and
follows certain conventions to comply with the \textsf{GAP} kernel interface, which we plan to document later. In the meantime, we advise
you to look at existing examples of such packages and get in touch with \textsf{GAP} developers if you plan to develop such a package. }

 }

  
\section{\textcolor{Chapter }{Having an InfoClass}}\label{Having an InfoClass}
\logpage{[ 76, 17, 0 ]}
\hyperdef{L}{X78969BA778DDE385}{}
{
   \index{InfoClass@\texttt{InfoClass}!for a GAP package} It is a good idea to declare an \texttt{InfoClass} for your package. This gives the package user the opportunity to control the
verbosity of output and/or the possibility of receiving debugging information
(see{\nobreakspace}\ref{Info Functions}). Below, we give a quick overview of its utility. 

 An \texttt{InfoClass} is defined with a \texttt{DeclareInfoClass( \mbox{\texttt{\mdseries\slshape InfoPkgname}} );} statement and may be set to have an initial \texttt{InfoLevel} other than the zero default (which means no \texttt{Info} statement is to output information) via a \texttt{SetInfoLevel( \mbox{\texttt{\mdseries\slshape InfoPkgname}}, \mbox{\texttt{\mdseries\slshape level}} );} statement. An initial \texttt{InfoLevel} of 1 is typical. 

 \texttt{Info} statements have the form: \texttt{Info( \mbox{\texttt{\mdseries\slshape InfoPkgname}}, \mbox{\texttt{\mdseries\slshape level}}, \mbox{\texttt{\mdseries\slshape expr1}}, \mbox{\texttt{\mdseries\slshape expr2}}, ...);} where the expression list \texttt{\mbox{\texttt{\mdseries\slshape expr1}}, \mbox{\texttt{\mdseries\slshape expr2}}, ...} appears just like it would in a \texttt{Print} statement. The only difference is that the expression list is only printed (or
even executed) if the \texttt{InfoLevel} of \mbox{\texttt{\mdseries\slshape InfoPkgname}} is at least \mbox{\texttt{\mdseries\slshape level}}. }

  
\section{\textcolor{Chapter }{The Banner}}\label{The Banner}
\logpage{[ 76, 18, 0 ]}
\hyperdef{L}{X784E0A5A7DB88332}{}
{
  \index{banner!for a GAP package} When the package is loaded, \textsf{GAP} will display a default package banner, constructed from the package metadata
provided in the \texttt{PackageInfo.g} file. 

 Alternatively, the package may establish its own banner by assigning either a
string to the \texttt{BannerString} field of the record argument of \texttt{SetPackageInfo} in the \texttt{PackageInfo.g} file or a function to the \texttt{BannerFunction} field, which takes this record as its unique argument. The latter possibility
can be useful if the banner shall show information that is available only at
runtime. 

 If you will be designing a banner for your package, it is a good idea to
suggest there how to access package documentation. For example, the banner of
the \textsf{Example} package says: 

 
\begin{Verbatim}[commandchars=!@|,fontsize=\small,frame=single,label=Example]
  For help, type: ?Example package
\end{Verbatim}
 

 In order for this to display the introduction of the \textsf{Example} package the index\texttt{\symbol{45}}entry \texttt{{\textless}Index{\textgreater}Example package{\textless}/Index{\textgreater}} was added just before the first paragraph of the introductory section in the
file \texttt{doc/example.xml} of the \textsf{Example} package. }

  
\section{\textcolor{Chapter }{Version Numbers}}\label{Version Numbers}
\logpage{[ 76, 19, 0 ]}
\hyperdef{L}{X8180BCDA82587F41}{}
{
  \index{version number!for a GAP package} Version numbers are strings containing nonnegative integers separated by
non\texttt{\symbol{45}}numeric characters. They are compared by \texttt{CompareVersionNumbers} (\ref{CompareVersionNumbers}) which first splits them at non\texttt{\symbol{45}}digit characters and then
lexicographically compares the resulting integer lists. Thus version \texttt{"2\texttt{\symbol{45}}3"} is larger than version \texttt{"2\texttt{\symbol{45}}2\texttt{\symbol{45}}5"} but smaller than \texttt{"4r2p3"} or \texttt{"11.0"}. 

 It is possible for code to require \textsf{GAP} packages in certain versions. In this case, all versions, whose number is
equal or larger than the requested number are acceptable. It is the task of
the package author to provide upwards compatibility. 

 Loading a specific version of a package (that is, \emph{not} one with a larger version number) can be achieved by prepending \texttt{=} to the desired version number. For example, \texttt{LoadPackage( "example", "=1.0" )} will load version \texttt{"1.0"} of the package \texttt{"example"}, even if version \texttt{"1.1"} is available. As a consequence, version numbers must not start with \texttt{=}, so \texttt{"=1.0"} is not a valid version number. 

 Package authors should choose a version numbering scheme that admits a new
version number even after tiny changes to the package, and ensure that version
numbers of successive package versions increase. The automatic update of
package archives in the \textsf{GAP} distribution will only work if a package has a new version number. 

 It is a well\texttt{\symbol{45}}established custom to name package archives
like \texttt{name\texttt{\symbol{45}}version.tar.gz}, \texttt{name\texttt{\symbol{45}}version.tar.bz2} etc., where \texttt{name} is the lower case name, and \texttt{version} is the version (another custom is that the archive then should extract to a
directory that has exactly the name \texttt{name\texttt{\symbol{45}}version}). 

 It is very important that there should not ever be, for a given \textsf{GAP} package, two different archives with the same package version number. If you
make changes to your package and place a new archive of the package onto the
public server, please ensure that a new archive has a new version number. This
should be done even for very minor changes. 

 For most of the packages it will be inappropriate to re\texttt{\symbol{45}}use
the date of the release as a version number. It is much more obvious how big
are the changes between versions "4.4.12", "4.5.1" and "4.5.2" than between
versions "2008.12.17", "2011.04.15" and "2011.09.14". The concept of using
version numbers to convey the meaning of the status of the code and the way it
has been modified is known as ``Semantic Versioning'', see \href{https://semver.org/} {\texttt{https://semver.org/}} for further recommendations on its use. 

 Since version information is duplicated in several places throughout the
package documentation, for \textsf{GAPDoc}\texttt{\symbol{45}}based manuals you may define the version and the release
manual in the comments in \texttt{PackageInfo.g} file close to the place where you specified its \texttt{Version} and \texttt{Date} components, for example 
\begin{Verbatim}[commandchars=@|B,fontsize=\small,frame=single,label=Example]
  ##  <#GAPDoc Label="PKGVERSIONDATA">
  ##  <!ENTITY VERSION "3.3">
  ##  <!ENTITY RELEASEDATE "12/09/2017">
  ##  <!ENTITY RELEASEYEAR "2017">
  ##  <#/GAPDoc>
\end{Verbatim}
 notify \texttt{MakeGAPDocDoc} (\textbf{GAPDoc: MakeGAPDocDoc}) that a part of the document is stored in \texttt{PackageInfo.g} (see \texttt{example/makedocrel.g}), read this data into the header of the main document via \texttt{{\textless}\#Include Label="PKGVERSIONDATA"{\textgreater}} directive and then use them via \&VERSION; and \&RELEASEDATE; entities almost
everywhere where you need to refer to them (most commonly, in the title page
and installation instructions). }

  
\section{\textcolor{Chapter }{Testing a \textsf{GAP} package}}\label{Testing a GAP package}
\logpage{[ 76, 20, 0 ]}
\hyperdef{L}{X8559D1FF7C9B7D14}{}
{
  There are several aspects of testing a \textsf{GAP} package. 

 First, one should ensure that the package functionality works as expected.
Below we give an advice on creating test files for automated tests that may be
run by package authors, by \textsf{GAP} developers as part of the release preparation, and by package users interested
in checking that the package works. Such tests should be included in the
package distribution, and the responsibility for ensuring that they pass stays
with package authors. 

 Second, the package should cleanly integrate into the \textsf{GAP} system and other packages, and should not break their functionality. In
particular, all tests from the standard \textsf{GAP} testing suite should pass if the package is loaded. This is more comprehensive
and time consuming test, which \textsf{GAP} developers regularly run using special tools. They will report to you any
detected issues. Below we explain how to do several simple and less time
consuming checks which package authors are recommended to perform themselves. 
\subsection{\textcolor{Chapter }{Tests files for a GAP package}}\label{Tests files for a GAP package}
\logpage{[ 76, 20, 1 ]}
\hyperdef{L}{X85CA2F547CF87666}{}
{
  The (optional) \texttt{tst} directory of your package may contain as many tests of the package
functionality as appears appropriate. These tests should be organised into
test files similarly to those in the \texttt{tst} directory of the \textsf{GAP} distribution as documented in \ref{Test Files}. 

 For a deposited package, a test file with a basic test of the package (for
example, to check that it works as expected and/or that the manual examples
are correct) may be specified in the component \texttt{TestFile} in the \texttt{PackageInfo.g} to be included in the \textsf{GAP} standard test suite. This file can either consist of calls of \texttt{TestDirectory} (\ref{TestDirectory}) or \texttt{Test} (\ref{Test}) (in this case, it is common to call it \texttt{testall.g}) or be itself a test file having an extension \texttt{.tst} and supposed to be read via \texttt{Test} (\ref{Test}). It is assumed that the latter case occurs if and only if the file contains
the substring

 {\nobreakspace}\texttt{"gap{\textgreater} START{\textunderscore}TEST("} 

 \noindent (with exactly one space after the \textsf{GAP} prompt). 

 For deposited packages, these tests are run by the \textsf{GAP} Group regularly, as a part of the standard \textsf{GAP} test suite. For the efficient testing it is important that the test specified
in the \texttt{PackageInfo.g} file does not display any output (e.g. no test progress indicators) except
reporting discrepancies if such occur and the completion report as in the
example below: 
\begin{Verbatim}[commandchars=!@|,fontsize=\small,frame=single,label=Example]
  !gapprompt@gap>| !gapinput@Test("tst/testall.tst");|
  Example package: testall.tst
  true
\end{Verbatim}
 Tests which produce extended output and/or require substantial runtime are not
supposed to be a part of the \textsf{GAP} standard test suite but may be placed in the \texttt{tst} directory of the packages with further instructions on how to run them
elsewhere. 

 Because of different approaches to testing, used by different packages, it is
not always easy to identify whether an automated test passed or failed.
Presently, automated detection works fine if a package uses a single \texttt{.tst} file or uses \texttt{TestDirectory} (\ref{TestDirectory}) with the \texttt{exitGAP} option set to \texttt{true} to run a collection of tests and then exits \textsf{GAP} in a way that allows an automated test setup to determine whether the test
passed or failed: 
\begin{Verbatim}[commandchars=!@|,fontsize=\small,frame=single,label=Example]
  TestDirectory(DirectoriesPackageLibrary("packagename", "tst"), rec(exitGAP := true));
\end{Verbatim}
 If one needs a more sophisticated test file, then it should end with an
invocation of \texttt{ForceQuitGap} (\ref{ForceQuitGap}) with an argument that indicates whether the tests overall passed (\texttt{true}) or failed (\texttt{false} or \texttt{fail}). For example, if the test result is stored in a variable \texttt{testresult} then you can do this: 
\begin{Verbatim}[commandchars=!@|,fontsize=\small,frame=single,label=Example]
  ForceQuitGap(testresult);
\end{Verbatim}
 }

 
\subsection{\textcolor{Chapter }{Testing \textsf{GAP} package loading}}\label{Testing GAP package loading}
\logpage{[ 76, 20, 2 ]}
\hyperdef{L}{X84CD542B7C4C73A0}{}
{
  To test that your package may be loaded into \textsf{GAP} without any problems and conflicts with other packages, test that it may be
loaded in various configurations: 
\begin{itemize}
\item  starting \textsf{GAP} with no packages (except needed for \textsf{GAP}) using \texttt{\texttt{\symbol{45}}r \texttt{\symbol{45}}A} options and calling \texttt{LoadPackage("packagename");} 
\item  starting \textsf{GAP} with no packages (except needed for \textsf{GAP}) using \texttt{\texttt{\symbol{45}}r \texttt{\symbol{45}}A} options and calling \texttt{LoadPackage("packagename" : OnlyNeeded );} 
\item  starting \textsf{GAP} in the default configuration (with no options) and calling \texttt{LoadPackage("packagename");} 
\item  starting \textsf{GAP} in the default configuration (with no options) and calling \texttt{LoadPackage("packagename" : OnlyNeeded );} 
\item  finally, together with all other packages using \texttt{LoadAllPackages} (\ref{LoadAllPackages}) (see below) in four possible combinations of starting \textsf{GAP} with/without \texttt{\texttt{\symbol{45}}r \texttt{\symbol{45}}A} options and calling \texttt{LoadAllPackages} (\ref{LoadAllPackages}) with/without \texttt{reversed} option. 
\end{itemize}
 The test of loading all packages is the most subtle one. Quite often it
reveals problems which do not occur in the default configuration but may cause
difficulties to the users of specialised packages. 

 Additionally, we recommend using \texttt{ShowPackageVariables} (\ref{ShowPackageVariables}) to see information about variables created by your package to check if any
have either short names (no more than three characters) or names breaking a
recommended naming convention that the \textsf{GAP} core system. \textsf{GAP} packages also should not use global variables starting in the lowercase (see
Section \ref{Functions and Variables and Choices of Their Names}). }

 

\subsection{\textcolor{Chapter }{LoadAllPackages}}
\logpage{[ 76, 20, 3 ]}\nobreak
\hyperdef{L}{X80A0D21D78CF8494}{}
{\noindent\textcolor{FuncColor}{$\triangleright$\enspace\texttt{LoadAllPackages({\mdseries\slshape : reversed})\index{LoadAllPackages@\texttt{LoadAllPackages}}
\label{LoadAllPackages}
}\hfill{\scriptsize (function)}}\\


 loads all \textsf{GAP} packages from their list sorted in alphabetical order (needed and suggested
packages will be loaded when required). This is a technical function to check
packages compatibility, so it should NOT be used to run anything except tests;
it is known that \textsf{GAP} performance is slower if all packages are loaded. To introduce some variations
of the order in which packages will be loaded for testing purposes, \texttt{LoadAllPackages} accepts option \texttt{reversed} to load packages from their list sorted in the reverse alphabetical order. }

 
\subsection{\textcolor{Chapter }{Testing a \textsf{GAP} package with the \textsf{GAP} standard test suite}}\label{Testing a GAP package with the GAP standard test suite}
\logpage{[ 76, 20, 4 ]}
\hyperdef{L}{X7C90C8B87BF6EF0B}{}
{
  The \texttt{tst} directory of the \textsf{GAP} installation contains a selection of test files and scripts such as \texttt{testinstall.g} and \texttt{teststandard.g} which are a part of the \textsf{GAP} standard test suite. 

 It is important to check that your package does not break \textsf{GAP} standard tests. To perform a clean test and avoid interfering with other
packages, first you must start a new \textsf{GAP} session and then read either \texttt{testinstall.g} or \texttt{teststandard.g} as demonstrated below. 

 The quicker test, \texttt{testinstall.g}, should run in about a minute depending on your hardware speed. It may be
started with the command 
\begin{Verbatim}[commandchars=!@|,fontsize=\small,frame=single,label=Example]
  !gapprompt@gap>| !gapinput@Read( Filename( DirectoriesLibrary( "tst" ), "testinstall.g" ) );|
\end{Verbatim}
 You will get a large number of lines with output about the progress of the
tests, for example: 
\begin{Verbatim}[commandchars=!@|,fontsize=\small,frame=single,label=Example]
  You should start GAP4 using `gap -A -x 80 -r'.
  
  Architecture: SOMETHING-SOMETHING-gcc-default64
  
  testing: ..../gap-4.X.Y/tst/testinstall/alghom.tst
        84 ms (55 ms GC) and 2.90MB allocated for alghom.tst
  testing: ..../gap-4.X.Y/tst/testinstall/algmat.tst
       839 ms (114 ms GC) and 219MB allocated for algmat.tst
  [ further lines deleted ]
  testing: ..../gap-4.X.Y/tst/testinstall/zmodnze.tst
       127 ms (119 ms GC) and 1.29MB allocated for zmodnze.tst
  -----------------------------------
  total     62829 ms (24136 ms GC) and 8.61GB allocated
                0 failures in 252 files
  
  #I  No errors detected while testing
\end{Verbatim}
 

 (optionally, you may start \textsf{GAP} with the command line options which you will see in the test output, to run it
in a more conservative settings). 

 The more thorough test is \texttt{teststandard.g} which exercises more of \textsf{GAP}'s capabilities, also including all test files from \texttt{teststandard.g}. It runs quite a bit longer, maybe 10\texttt{\symbol{45}}20 minutes, and
produces an output similar to the testinstall.g test. To run it, also start a
new \textsf{GAP} session and then call 
\begin{Verbatim}[commandchars=!@|,fontsize=\small,frame=single,label=Example]
  !gapprompt@gap>| !gapinput@Read( Filename( DirectoriesLibrary( "tst" ), "teststandard.g" ) );|
\end{Verbatim}
 You may repeat the same check loading your package with \texttt{OnlyNeeded} option. Remember to perform each subsequent test in a new \textsf{GAP} session. 

 Also you may perform individual tests from the \texttt{tst} directory of the \textsf{GAP} installation loading them with \texttt{Test} (\ref{Test}). 

 }

 }

  
\section{\textcolor{Chapter }{Access to the \textsf{GAP} Development Version}}\label{Access to the GAP Development Version}
\logpage{[ 76, 21, 0 ]}
\hyperdef{L}{X81B52B657CA75BDA}{}
{
  We are aiming at providing a stable platform for package development and
testing with official \textsf{GAP} releases. We also invite everyone to contribute by submitting patches, pull
requests, and bug reports. We would like to make the contributing process as
easy as possible. 

 The main GAP development repository is hosted on GitHub at \href{https://github.com/gap-system/gap} {\texttt{https://github.com/gap\texttt{\symbol{45}}system/gap}}. Many \textsf{GAP} packages also have public repositories and issue trackers, and we are keeping
a list of such packages at \href{https://gap-packages.github.io/} {\texttt{https://gap\texttt{\symbol{45}}packages.github.io/}}. 

 For further information about contributing to the GAP development, please see \href{https://github.com/gap-system/gap/blob/master/CONTRIBUTING.md} {\texttt{https://github.com/gap\texttt{\symbol{45}}system/gap/blob/master/CONTRIBUTING.md}}. }

  
\section{\textcolor{Chapter }{Version control and continuous integration for \textsf{GAP} packages}}\label{Version control and continuous integration for GAP packages}
\logpage{[ 76, 22, 0 ]}
\hyperdef{L}{X836CDF8F7A846A1C}{}
{
  As we have mentioned above, many \textsf{GAP} packages have public repositories and issue trackers on GitHub, and we are
keeping a list of such packages at \href{https://gap-packages.github.io/} {\texttt{https://gap\texttt{\symbol{45}}packages.github.io/}}. We welcome establishing public repositories for new packages and migrating
existing package repositories there as well. Such repositories may be hosted
under their authors' accounts or under the gap\texttt{\symbol{45}}packages
organisation (\href{https://github.com/gap-packages/} {\texttt{https://github.com/gap\texttt{\symbol{45}}packages/}}). The latter has the benefit that while the authors will preserve their
deciding role on all aspects of the package development, the package will
become more visible for potential collaborators and \textsf{GAP} developers may help to set up \emph{continuous integration} for your package so that every commit to the repository will trigger automated
running of package tests and reporting any failures to package maintainers. }

  
\section{\textcolor{Chapter }{Selecting a license for a \textsf{GAP} Package}}\label{Selecting a license for a GAP Package}
\logpage{[ 76, 23, 0 ]}
\hyperdef{L}{X82EBCBC5829B6001}{}
{
  As it was mentioned in the description of the \texttt{LICENSE} file in Section \ref{Structure of a GAP Package}, it is advised to make clear in the documentation of the package the basis on
which it is being distributed to users. \textsf{GAP} itself is distributed under the GNU Public License version 2 (version 2 or
later). We would encourage you to consider the GPL license for your packages,
but you might wish to be more restrictive (for instance forbidding
redistribution for profit) or less restrictive (allowing your software to be
incorporated into commercial software). See ``Choosing a License for the Distribution of Your Package'' from \href{https://www.gap-system.org/Packages/Authors/authors.html} {\texttt{https://www.gap\texttt{\symbol{45}}system.org/Packages/Authors/authors.html}} and also \href{https://choosealicense.com/} {\texttt{https://choosealicense.com/}} for further details. 

 In the past many \textsf{GAP} packages used the text ``We adopt the copyright regulations of \textsf{GAP} as detailed in the copyright notice in the \textsf{GAP} manual'' or a similar statement. We now advise to be more explicit by making the exact
reference to the GPL license, for example: 

 \emph{ \textsf{packagename} is free software; you can redistribute it and/or modify it under the terms of
the \href{https://www.fsf.org/licenses/gpl.html} {GNU General Public License} as published by the Free Software Foundation; either version 2 of the License,
or (at your option) any later version. } and also including a copy of the full text of the license. }

  
\section{\textcolor{Chapter }{Releasing a GAP Package}}\label{Releasing a GAP Package}
\logpage{[ 76, 24, 0 ]}
\hyperdef{L}{X8074AAAE79911BE5}{}
{
  The \textsf{GAP} distribution provides archives in several different formats. 

 
\begin{description}
\item[{\texttt{.tar.gz}}]  a standard UNIX \texttt{tar} archive, compressed with \texttt{gzip} 
\item[{\texttt{.tar.bz2}}]  a standard UNIX \texttt{tar} archive, compressed with \texttt{bzip2} 
\item[{\texttt{.zip}}]  an archive in \texttt{zip} format, where text files should have UNIX style line breaks 
\end{description}
 

 For convenience of possible users it is sensible that you provide an archive
of your package in at least one of these formats. 

 For example, if you wish to supply a \texttt{.tar.gz} archive, you may create it with the command 

 {\nobreakspace}\texttt{tar \texttt{\symbol{45}}cvzf packagename\texttt{\symbol{45}}version.tar.gz
packagename} 

 \noindent Because the release of the \textsf{GAP} package is independent of the version of \textsf{GAP}, a \textsf{GAP} package should be wrapped up in separate file that can be installed onto any
version of \textsf{GAP}. In this way, a package can be upgraded any time without the need to wait for
new \textsf{GAP} releases. To ensure this, the package should be archived from the \textsf{GAP} \texttt{pkg} directory, that is all files are archived with the path starting at the
package's name. 

 \index{GAPDocManualLab@\texttt{GAPDocManualLab}} The archive of a \textsf{GAP} package should contain all files necessary for the package to work. In
particular there should be a compiled documentation, which includes the \texttt{manual.six}, \texttt{manual.toc} and \texttt{manual.lab} file in the documentation subdirectory which are created by \textsf{GAPDoc} while {\TeX}ing the documentation. (The first two files are needed by the \textsf{GAP} help system, and the \texttt{manual.lab} file is needed if the main manuals or another package is referring to your
package. Use the command \texttt{GAPDocManualLab( packagename );} to create this file for your help books if you use \textsf{GAPDoc}.) 

 Note that wrapping the \textsf{GAP} distribution as a single archive containing the core system and all currently
redistributed packages, will change file timestamps, so one should not rely on
them anywhere in the package. 

 For packages hosted on GitHub publishing package release and establishing its
website can be very efficiently automated using two tools: ReleaseTools (\href{https://github.com/gap-system/ReleaseTools} {\texttt{https://github.com/gap\texttt{\symbol{45}}system/ReleaseTools}}) and GitHubPagesForGAP (\href{https://github.com/gap-system/GitHubPagesForGAP} {\texttt{https://github.com/gap\texttt{\symbol{45}}system/GitHubPagesForGAP}}). }

  
\section{\textcolor{Chapter }{The homepage of a Package}}\label{The homepage of a Package}
\logpage{[ 76, 25, 0 ]}
\hyperdef{L}{X8232CC1385C4B1DD}{}
{
  If you want to distribute your package you should create its homepage
containing some basic information, archives for download, the \texttt{README} file with installation instructions, and a copy of the package's \texttt{PackageInfo.g} file. 

 The responsibility to maintain this homepage is with the package
authors/maintainers. 

 If you tell the \textsf{GAP} Group about your package (say, by mail to \href{mailto://support@gap-system.org} {\texttt{support@gap\texttt{\symbol{45}}system.org}}) we may consider either 
\begin{itemize}
\item  adding a link to your package homepage from the \textsf{GAP} website (thus, the package will be an \emph{undeposited contribution}); 
\item  or redistributing the current version of your package as a part of the \textsf{GAP} distribution (this, the package will be \emph{deposited}), also ; 
\end{itemize}
 Please also consider submitting your package to the \textsf{GAP} package refereeing process (see \href{https://www.gap-system.org/Contacts/submit.html} {\texttt{https://www.gap\texttt{\symbol{45}}system.org/Contacts/submit.html}} for further information). 

 For packages hosted on GitHub publishing package release and establishing its
website can be very efficiently automated using two tools: GitHubPagesForGAP (\href{https://github.com/gap-system/GitHubPagesForGAP} {\texttt{https://github.com/gap\texttt{\symbol{45}}system/GitHubPagesForGAP}}) and ReleaseTools (\href{https://github.com/gap-system/ReleaseTools} {\texttt{https://github.com/gap\texttt{\symbol{45}}system/ReleaseTools}}). }

  
\section{\textcolor{Chapter }{Some things to keep in mind}}\label{Some things to keep in mind}
\logpage{[ 76, 26, 0 ]}
\hyperdef{L}{X796D7F7583E845BE}{}
{
  
\begin{itemize}
\item  Some packages still use for their manuals the old ``gapmacro'' format, support for which may be discontinued in the future. We encourage
authors of those packages to eventually convert their documentation \textsf{GAPDoc}. New packages are recommended to use \textsf{GAPDoc}, which, for example, is capable of creating HTML documentation with MathJax
support, allows easy extraction of examples from documentation for testing,
etc. One could also use the \textsf{AutoDoc} which simplifies writing documentation by generating most of the \textsf{GAPDoc} code automatically. 
\item  The concept of an autoloaded package, which existed before \textsf{GAP}{\nobreakspace}4.5, has been integrated with the \emph{needed} and \emph{suggested} mechanism that exists between packages. \textsf{GAP} itself now ``needs'' certain packages (for instance \textsf{GAPDoc}) and ``suggests'' others (typically the packages that were autoloaded). The decisions which
packages \textsf{GAP} should need or suggest are made by developers based on technical criteria.
They can be easily overridden by a user using the new \texttt{gap.ini} (see \ref{sect:gap.ini}). The default file ensures that all formerly autoloaded packages are still
loaded if present. 
\item  Optional \texttt{\texttt{\symbol{126}}/.gap} directory for user's customisations which may contain e.g. locally installed
packages (see \ref{GAP Root Directories}). If package installation instructions explain how to install the package in
a non\texttt{\symbol{45}}standard location, they should mention this. 
\item  Packages loading mechanism allows to make loading packages more informative,
while avoiding confusing the user with warning and error messages for packages
they didn't know they were loading. For example, many messages are stored but
not displayed using the function \texttt{LogPackageLoadingMessage} (\ref{LogPackageLoadingMessage}) and there is a function \texttt{DisplayPackageLoadingLog} (\ref{DisplayPackageLoadingLog}) to show log messages that occur during package loading. Packages are
encouraged to use these mechanisms to report problems in loading (e.g.
binaries not compiled), rather than printing messages directly. 
\end{itemize}
 }

  
\section{\textcolor{Chapter }{Package release checklists}}\label{Package release checklists}
\logpage{[ 76, 27, 0 ]}
\hyperdef{L}{X82CE0A518440CCBB}{}
{
  The following checklists should be useful to package authors and maintainers,
as well as to everyone involved in the depositing and refereeing of \textsf{GAP} packages. 
\subsection{\textcolor{Chapter }{Checklist for releasing a new package}}\label{Checklist for releasing a new package}
\logpage{[ 76, 27, 1 ]}
\hyperdef{L}{X80E3926A7CF8B6DC}{}
{
  
\begin{itemize}
\item  Test that the package: 
\begin{itemize}
\item  does not break \texttt{testinstall.g} and \texttt{teststandard.g}, and does not slow them down noticeably (see \ref{Testing a GAP package with the GAP standard test suite}); 
\item  may be loaded in various configurations (see \ref{Testing GAP package loading}); 
\item  follows the guidelines of Section \ref{Functions and Variables and Choices of Their Names} about names of functions and variables; 
\end{itemize}
 
\item  \texttt{PackageInfo.g} file: 
\begin{itemize}
\item  correctly specifies package version, release date, and package authors; 
\item  passes validation using \texttt{ValidatePackageInfo} (\ref{ValidatePackageInfo}); 
\item  besides mandatory components, which are required to pass validation, also has
relevant optional components (such as, for example, URLs of public source code
repository and issue tracker; hints to distinguish binary and text files in
case of non\texttt{\symbol{45}}standard file names and extensions, etc.); 
\end{itemize}
 
\item  Package documentation: 
\begin{itemize}
\item  is built and included in the package archive together with its source files; 
\item  states the same version, release date and package authors as specified in the \texttt{PackageInfo.g} file; 
\item  has the same version, release date and package authors details as stated in
the \texttt{PackageInfo.g} file; 
\item  is searchable using the \textsf{GAP} help system in all formats (text, HTML and PDF); 
\item  is clear about the license under which the package is distributed, and refers
to the \texttt{LICENSE} file which should be included in the package; 
\end{itemize}
 
\item  Package archive(s): 
\begin{itemize}
\item  have correct permissions for all files and directories after their unpacking
(755 for directories and executables, if any; 644 for other files); 
\item  contain files with correct line breaks for the given format (see Section \ref{Releasing a GAP Package}); 
\item  contain no hidden system files and directories that are not supposed to be
included in the package, e.g. \texttt{.gitignore}, \texttt{.git} etc.; 
\end{itemize}
 
\item  Package availability: 
\begin{itemize}
\item  not only the package archive(s), but also the \texttt{PackageInfo.g} and \texttt{README} files are available online; 
\end{itemize}
 
\end{itemize}
 }

 
\subsection{\textcolor{Chapter }{Checklist for upgrading the package for the next major release of \textsf{GAP}}}\label{Checklist for upgrading the package for the next major release of GAP}
\logpage{[ 76, 27, 2 ]}
\hyperdef{L}{X820D4B207A41AEA6}{}
{
  \textsf{GAP} ecosystem is not static: both the core \textsf{GAP} system and packages redistributed with \textsf{GAP} are in constant development. \textsf{GAP} has a policy that changes that may have a disruptive effect on packages
redistributed with \textsf{GAP} should only be introduced in major \textsf{GAP} releases. When the next \textsf{GAP} major release is prepared, a beta version for package authors will be made
available in order to give them an opportunity to check and update, if
necessary, their packages for the public release of the next major version of \textsf{GAP}. 

 The following checklist will help you to check how well your package is ready
to work with the next major release of \textsf{GAP} 

 
\begin{itemize}
\item  Check that the package functionality works as expected, package tests run with
no discrepancies, and manual examples correspond to new version of \textsf{GAP}. This is a convenient opportunity to polish existing and add new tests, and
improve manual examples. 
\item  Revise package dependencies: check that the \texttt{PackageInfo.g} file has correct list of needed and suggested packages (see Section \ref{Package dependencies}). 
\item  Revise licensing information: check that the package states clearly under
which conditions it is distributed and includes a \texttt{LICENSING} file with the text of a license (see Section \ref{Selecting a license for a GAP Package}). 
\item  Rebuild the package documentation to update
cross\texttt{\symbol{45}}references to main \textsf{GAP} manuals and, if relevant, to the documentation of other \textsf{GAP} packages. This will ensure that cross\texttt{\symbol{45}}references from the
package manual to the main \textsf{GAP} manuals are correct and that the \textsf{GAP} help system will be able to navigate to the more precise location in the
package manual. This will also improve the layout of the package documentation
by picking up the changes in documenting tools. 
\item  Check if the package still relies on some obsolete variables (see Chapter \ref{Replaced and Removed Command Names}) and replace their usage by the new commands. To perform such check, start \textsf{GAP} with `\texttt{\symbol{45}}O` command line option to disable loading obsoletes,
and then load your package. 
\item  Check for any specific advice in release notes for the beta release for
package authors. 
\end{itemize}
 }

 }

 }

       
\chapter{\textcolor{Chapter }{Replaced and Removed Command Names}}\label{Replaced and Removed Command Names}
\logpage{[ 77, 0, 0 ]}
\hyperdef{L}{X78C85ED17F00DCC1}{}
{
  \index{obsolete} \index{deprecated} \index{legacy} In general we try to keep \textsf{GAP}{\nobreakspace}4 compatible with former releases as much as possible.
Nevertheless, from time to time it seems appropriate to remove some commands
or to change the names of some commands or variables. There are various
reasons for that: Some functionality was improved and got another (hopefully
better) interface, names turned out to be too special or too general for the
underlying functionality, or names are found to be unintuitive or inconsistent
with other names. 

 In this chapter we collect such old names while pointing to the sections which
explain how to substitute them. Usually, old names will be available for
several releases; they may be removed when they don't seem to be used any
more. 

 Information about obsolete names is printed by \texttt{Info} (\ref{Info}) using the \texttt{InfoObsolete} (\ref{InfoObsolete}) Info class. By default \texttt{InfoObsolete} (\ref{InfoObsolete}) is set to 1. Newly obsoleted identifiers should at first be outputted at info
level 2. Once they have been removed from all packages, they should then be
moved to info level 1, so they are visible to normal users, for at least one
major release before being removed. 

 The functions \texttt{DeclareObsoleteSynonym} and \texttt{DeclareObsoleteSynonymAttr} take an optional final parameter, specifying the info level at which the given
obsolete symbol should be reported. It defaults to 2 and 1, respectively. 

 The obsolete \textsf{GAP} code is collected in two library files, \texttt{lib/obsolete.gd} and \texttt{lib/obsolete.gi}. By default, these files are read when \textsf{GAP} is started. It may be useful to omit reading these files, for example in order
to make sure that one's own \textsf{GAP} code does not rely on the obsolete variables. For that, one can use the \texttt{\texttt{\symbol{45}}O} command line option (see \ref{Command Line Options}) or set the user preference \texttt{ReadObsolete} in the file \texttt{gap.ini} to \texttt{false} (see \ref{sect:gap.ini}). Note that the \texttt{\texttt{\symbol{45}}O} command line option overrides \texttt{ReadObsolete}. 

 (Note that the condition whether the library files with the obsolete \textsf{GAP} code shall be read has changed. In \textsf{GAP}{\nobreakspace}4.3 and 4.4, the global variables \texttt{GAP{\textunderscore}OBSOLESCENT} and \texttt{GAPInfo.ReadObsolete} {\textendash}to be set in the user's \texttt{.gaprc} file{\textendash} were used to control this behaviour.)  
\section{\textcolor{Chapter }{Group Actions {\textendash} Name Changes}}\label{Group Actions - Name Changes}
\logpage{[ 77, 1, 0 ]}
\hyperdef{L}{X7AA51AC9870D2360}{}
{
  \index{group operations} The concept of a group action is sometimes referred to as a ``group operation''. In \textsf{GAP}{\nobreakspace}3 as well as in older versions of \textsf{GAP}{\nobreakspace}4 the term \texttt{Operation} was used instead of \texttt{Action}. We decided to change the names to avoid confusion with the term ``operation'' as in \texttt{DeclareOperation} (\ref{DeclareOperation}) and ``operations for Xyz''. 

 Here are some examples of such name changes. 

 \begin{center}
\begin{tabular}{l|l}\emph{OLD}&
\emph{NOW USE}\\
\hline
\texttt{Operation}\index{Operation@\texttt{Operation}}&
\texttt{Action} (\ref{Action:for a group, an action domain, etc.})\\
\texttt{RepresentativeOperation}\index{RepresentativeOperation@\texttt{RepresentativeOperation}}&
\texttt{RepresentativeAction} (\ref{RepresentativeAction})\\
\texttt{OperationHomomorphism}\index{OperationHomomorphism@\texttt{OperationHomomorphism}}&
\texttt{ActionHomomorphism} (\ref{ActionHomomorphism:for a group, an action domain, etc.})\\
\texttt{FunctionOperation}\index{FunctionOperation@\texttt{FunctionOperation}}&
\texttt{FunctionAction} (\ref{FunctionAction})\\
\texttt{IsLexicographicallyLess}\index{IsLexicographicallyLess@\texttt{IsLexicographicallyLess}}&
\texttt{\texttt{\symbol{92}}{\textless}} (\ref{bSlash<})\\
\end{tabular}\\[2mm]
\end{center}

 }

  
\section{\textcolor{Chapter }{Package Interface {\textendash} Obsolete Functions and Name Changes}}\label{Package Interface - Obsolete Functions and Name Changes}
\logpage{[ 77, 2, 0 ]}
\hyperdef{L}{X831734077B00CB3B}{}
{
  With \textsf{GAP}{\nobreakspace}4.4 the package interface was changed. Thereby some functions
became obsolete and the names of some others were made more consistent. 

 \index{DeclarePackage@\texttt{DeclarePackage}} \index{DeclareAutoPackage@\texttt{DeclareAutoPackage}} \index{DeclarePackageDocumentation@\texttt{DeclarePackageDocumentation}} \index{DeclarePackageAutoDocumentation@\texttt{DeclarePackageAutoDocumentation}} The following functions are no longer needed: \texttt{DeclarePackage}, \texttt{DeclareAutoPackage}, \texttt{DeclarePackageDocumentation} and \texttt{DeclarePackageAutoDocumentation}. They are substituted by entries in the packages' \texttt{PackageInfo.g} files, see \ref{The PackageInfo.g File}. 

 Furthermore, the global variable \texttt{PACKAGES{\textunderscore}VERSIONS} is no longer needed, since this information is now contained in the \texttt{GAPInfo.PackagesInfo} record (see \ref{GAPInfo}). The global variable \texttt{Revisions} is also no longer needed, since the function \texttt{DisplayRevision} was made obsolete in \textsf{GAP}{\nobreakspace}4.5 and removed in \textsf{GAP}{\nobreakspace}4.7. 

 The following function names were changed. 

 \index{RequirePackage@\texttt{RequirePackage}} \index{ReadPkg@\texttt{ReadPkg}} \index{RereadPkg@\texttt{RereadPkg}} \begin{center}
\begin{tabular}{l|l}\emph{OLD}&
\emph{NOW USE}\\
\hline
\texttt{RequirePackage}&
\texttt{LoadPackage} (\ref{LoadPackage})\\
\texttt{ReadPkg}&
\texttt{ReadPackage} (\ref{ReadPackage})\\
\texttt{RereadPkg}&
\texttt{RereadPackage} (\ref{RereadPackage})\\
\end{tabular}\\[2mm]
\end{center}

 }

  
\section{\textcolor{Chapter }{Normal Forms of Integer Matrices {\textendash} Name Changes}}\label{Normal Forms of Integer Matrices - Name Changes}
\logpage{[ 77, 3, 0 ]}
\hyperdef{L}{X79676CD27EF0F096}{}
{
  \index{Smith normal form} \index{Hermite normal form} Former versions of \textsf{GAP}{\nobreakspace}4 documented several functions for computing the Smith or
Hermite normal form of integer matrices. Some of them were never implemented
and it was unclear which commands to use. The functionality of all of these
commands is now available with \texttt{NormalFormIntMat} (\ref{NormalFormIntMat}) and a few interface functions. }

  
\section{\textcolor{Chapter }{Miscellaneous Name Changes or Removed Names}}\label{Miscellaneous Name Changes or Removed Names}
\logpage{[ 77, 4, 0 ]}
\hyperdef{L}{X7F6A6CBC7C9E91E5}{}
{
  \index{QUIET@\texttt{QUIET}} \index{BANNER@\texttt{BANNER}} In former releases of \textsf{GAP}{\nobreakspace}4 there were some global variable names bound to general
information about the running \textsf{GAP}, such as path names or command line options. Although they were not
officially documented they were used by several users and in some packages. We
mention here \texttt{BANNER} and \texttt{QUIET}. This type of information is now collected in the global record \texttt{GAPInfo} (\ref{GAPInfo}). 

 Here are some further name changes. 

 \begin{center}
\begin{tabular}{l|l}\emph{OLD}&
\emph{NOW USE}\\
\hline
\texttt{MonomialTotalDegreeLess}\index{MonomialTotalDegreeLess@\texttt{MonomialTotalDegreeLess}}&
\texttt{MonomialExtGrlexLess} (\ref{MonomialExtGrlexLess})\\
\texttt{MultRowVector}\index{MultRowVector@\texttt{MultRowVector}}&
\texttt{MultVectorLeft} (\ref{MultVectorLeft})\\
\texttt{MutableCopyMat}\index{MutableCopyMat@\texttt{MutableCopyMat}}&
\texttt{MutableCopyMatrix} (\ref{MutableCopyMatrix:for a matrix object})\\
\texttt{MutableIdentityMat}\index{MutableIdentityMat@\texttt{MutableIdentityMat}}&
\texttt{IdentityMat} (\ref{IdentityMat})\\
\texttt{MutableNullMat}\index{MutableNullMat@\texttt{MutableNullMat}}&
\texttt{NullMat} (\ref{NullMat})\\
\texttt{NormedVectors}\index{NormedVectors@\texttt{NormedVectors}}&
\texttt{NormedRowVectors} (\ref{NormedRowVectors})\\
\texttt{RadicalGroup}\index{RadicalGroup@\texttt{RadicalGroup}}&
\texttt{SolvableRadical} (\ref{SolvableRadical})\\
\end{tabular}\\[2mm]
\end{center}

 
\begin{itemize}
\item  \index{PositionFirstComponent@\texttt{PositionFirstComponent}} The operation \texttt{PositionFirstComponent} has been deprecated in GAP 4.8 due to issues with its documentation and
implementation, and was removed in GAP 4.10. Instead of \texttt{PositionFirstComponent(list,obj)}, you may use \texttt{PositionSorted(list,[obj])} or \texttt{PositionProperty(list,x\texttt{\symbol{45}}{\textgreater}x[1]=obj)} as a replacement, depending on your specific use case. 
\item  \index{MultRowVector@\texttt{MultRowVector}} The five argument version of the operation \texttt{MultRowVector} has been deprecated in GAP 4.10 since it was unused and only available for
coefficient lists. Note that \texttt{MultRowVector} was also renamed to \texttt{MultVectorLeft}. 
\item  \index{TemporaryGlobalVarName@\texttt{TemporaryGlobalVarName}} The function \texttt{TemporaryGlobalVarName} has been deprecated in GAP 4.11. Code using it typically can be restructured
to not need it, e.g. by using \texttt{EvalString} (\ref{EvalString}). 
\item  A bunch of synonyms have been deprecated in GAP 4.13. The alternative names
have always existed and so are safe to use even in code which is meant to run
in older GAP versions. \begin{center}
\begin{tabular}{l|l}\emph{Deprecated}&
\emph{Recommended alternative}\\
\hline
\texttt{OneSM}\index{OneSM@\texttt{OneSM}}&
\texttt{OneSameMutability} (\ref{OneSameMutability})\\
\texttt{InverseSM}\index{InverseSM@\texttt{InverseSM}}&
\texttt{InverseSameMutability} (\ref{InverseSameMutability})\\
\texttt{ZeroSM}\index{ZeroSM@\texttt{ZeroSM}}&
\texttt{ZeroSameMutability} (\ref{ZeroSameMutability})\\
\texttt{AdditiveInverseSM}\index{AdditiveInverseSM@\texttt{AdditiveInverseSM}}&
\texttt{AdditiveInverseSameMutability} (\ref{AdditiveInverseSameMutability})\\
\texttt{OneAttr}\index{OneAttr@\texttt{OneAttr}}&
\texttt{One} (\ref{One}), \texttt{OneImmutable} (\ref{OneImmutable})\\
\texttt{InverseAttr}\index{InverseAttr@\texttt{InverseAttr}}&
\texttt{Inverse} (\ref{Inverse}), \texttt{InverseImmutable} (\ref{InverseImmutable})\\
\texttt{ZeroAttr}\index{ZeroAttr@\texttt{ZeroAttr}}&
\texttt{Zero} (\ref{Zero}), \texttt{ZeroImmutable} (\ref{ZeroImmutable})\\
\texttt{AdditiveInverseAttr}\index{AdditiveInverseAttr@\texttt{AdditiveInverseAttr}}&
\texttt{AdditiveInverse} (\ref{AdditiveInverse}), \texttt{AdditiveInverseImmutable} (\ref{AdditiveInverseImmutable})\\
\end{tabular}\\[2mm]
\end{center}

 
\end{itemize}
 

\subsection{\textcolor{Chapter }{InfoObsolete}}
\logpage{[ 77, 4, 1 ]}\nobreak
\hyperdef{L}{X87348614848EAD64}{}
{\noindent\textcolor{FuncColor}{$\triangleright$\enspace\texttt{InfoObsolete\index{InfoObsolete@\texttt{InfoObsolete}}
\label{InfoObsolete}
}\hfill{\scriptsize (info class)}}\\


 is an info class to display warnings when an obsolete variable is used. By
default, the info level for this class is set to 1, which will only show
variables which will be removed in the next major version of GAP. Setting it
to 2 will trigger further warnings, for variables which have alternative
names, or may be removed in future. This class can be set to 0 to disable all
obsolete warnings. 

 To check that the \textsf{GAP} code does not use any obsolete variables at parsing time, and not at a
runtime, use the \texttt{\texttt{\symbol{45}}O} command line option, see \ref{Command Line Options}. }

 }

  
\section{\textcolor{Chapter }{The former .gaprc file}}\label{obsolete-gaprc}
\logpage{[ 77, 5, 0 ]}
\hyperdef{L}{X7F2FF72A7AD60E0C}{}
{
  Up to \textsf{GAP}{\nobreakspace}4.4, a file \texttt{.gaprc} in the user's home directory (if available, and \textsf{GAP} was started without \texttt{\texttt{\symbol{45}}r} option) was read automatically during startup, early enough for influencing
the autoloading of packages and late enough for being allowed to execute any \textsf{GAP} code. On Windows machines this file was called \texttt{gap.rc}. 

 In \textsf{GAP}{\nobreakspace}4.5 the startup mechanism has changed, see \ref{sect:gap.ini} for details. These new configuration files are now contained in a directory \texttt{GAPInfo.UserGapRoot}. 

 For the sake of partial backwards compatibility, also the former file \texttt{\texttt{\symbol{126}}/.gaprc} is still supported for such initializations, but this file is read only if the
directory \texttt{GAPInfo.UserGapRoot} does not exist. In that case the \texttt{\texttt{\symbol{126}}/.gaprc} is read at the same time as \texttt{gaprc} would be read, i.{\nobreakspace}e., too late for influencing the startup of \textsf{GAP}. 

 As before, the command line option \texttt{\texttt{\symbol{45}}r} disables reading \texttt{\texttt{\symbol{126}}/.gaprc}, see \ref{Command Line Options}. 

 To migrate from the old setup to the new one introduced with \textsf{GAP}{\nobreakspace}4.5, first have a look at the function \texttt{WriteGapIniFile} (\ref{WriteGapIniFile}). Many users will find that all or most of what was set in the old \texttt{\texttt{\symbol{126}}/.gaprc} file can now be done via the user preferences in a \texttt{gap.ini} file. If you had code for new functions or abbreviations in your old \texttt{\texttt{\symbol{126}}/.gaprc} file or you were reading additional files, then move this into the file \texttt{gaprc} (without the leading dot, same name for all operating systems) in the
directory \texttt{GAPInfo.UserGapRoot}. }

  
\section{\textcolor{Chapter }{Semigroup properties}}\label{obsolete-semigroups-properties}
\logpage{[ 77, 6, 0 ]}
\hyperdef{L}{X7C89E03285799F40}{}
{
  Until Version 4.8 of \textsf{GAP} there was inconsistent use of the following properties of semigroups: \texttt{IsGroupAsSemigroup}, \texttt{IsMonoidAsSemigroup}, and \texttt{IsSemilatticeAsSemigroup}. \texttt{IsGroupAsSemigroup} was true for semigroups that mathematically defined a group, and for
semigroups in the category \texttt{IsGroup} (\ref{IsGroup}); \texttt{IsMonoidAsSemigroup} was true for semigroups that mathematically defined monoids, but did not
belong to the category \texttt{IsMonoid} (\ref{IsMonoid}); and \texttt{IsSemilatticeAsSemigroup} was simply a property of semigroups, there is no category \texttt{IsSemilattice}. 

 From Version 4.8 onwards, \texttt{IsMonoidAsSemigroup} returns \texttt{true} for semigroups in the category \texttt{IsMonoid} (\ref{IsMonoid}), and \texttt{IsSemilatticeAsSemigroup} has been moved to the \textsf{Semigroups} under the new name \texttt{IsSemilattice} (\textbf{Semigroups: IsSemilattice}). }

 }

        
\chapter{\textcolor{Chapter }{Method Selection}}\label{Method Selection}
\logpage{[ 78, 0, 0 ]}
\hyperdef{L}{X8058CC8187162644}{}
{
  \index{operation}\index{method} This chapter explains how \textsf{GAP} decides which function to call for which types of objects. It assumes that you
have read the chapters about objects (Chapter{\nobreakspace}\ref{Objects and Elements}) and types (Chapter{\nobreakspace}\ref{Types of Objects}). 

 An \emph{operation} is a special \textsf{GAP} function that bundles a set of functions, its \emph{methods}. 

 All methods of an operation compute the same result. But each method is
installed for specific types of arguments. 

 If an operation is called with a tuple of arguments, one of the applicable
methods is selected and called. 

 Special cases of methods are partial methods, immediate methods, and logical
implications.  
\section{\textcolor{Chapter }{Operations and Methods}}\label{Operations and Methods}
\logpage{[ 78, 1, 0 ]}
\hyperdef{L}{X7AEED9AB824CD4DA}{}
{
  

 Operations are functions in the category \texttt{IsOperation} (\ref{IsOperation}). 

 So on the one hand, \emph{operations} are \textsf{GAP} functions, that is, they can be applied to arguments and return a result or
cause a side\texttt{\symbol{45}}effect. 

 On the other hand, operations are more. Namely, an operation corresponds to a
set of \textsf{GAP} functions, called the \emph{methods} of the operation. 

 Each call of an operation causes a suitable method to be selected and then
called. The choice of which method to select is made according to the types of
the arguments, the underlying mechanism is described in the following
sections. 

 Examples of operations are the binary infix operators \texttt{=}, \texttt{+} etc., and \texttt{PrintObj} (\ref{PrintObj}) is the operation that is called for each argument of \texttt{Print} (\ref{Print}). 

 Also all attributes and properties are operations. Each attribute has a
special method which is called if the attribute value is already stored; this
method of course simply returns this value. 

 The setter of an attribute is called automatically if an attribute value has
been computed. Attribute setters are operations, too. They have a default
method that ignores the request to store the value. Depending on the type of
the object, there may be another method to store the value in a suitable way,
and then set the attribute tester for the object to \texttt{true}. 

\subsection{\textcolor{Chapter }{IsOperation}}
\logpage{[ 78, 1, 1 ]}\nobreak
\hyperdef{L}{X874C7C6D8650D648}{}
{\noindent\textcolor{FuncColor}{$\triangleright$\enspace\texttt{IsOperation({\mdseries\slshape obj})\index{IsOperation@\texttt{IsOperation}}
\label{IsOperation}
}\hfill{\scriptsize (Category)}}\\


 is the category of operations. Every operation is a function, but not vice
versa. 

 
\begin{Verbatim}[commandchars=!@|,fontsize=\small,frame=single,label=Example]
  !gapprompt@gap>| !gapinput@MinimalPolynomial;|
  <Operation "MinimalPolynomial">
  !gapprompt@gap>| !gapinput@IsOperation(MinimalPolynomial);|
  true
  !gapprompt@gap>| !gapinput@IsFunction(MinimalPolynomial);|
  true
  !gapprompt@gap>| !gapinput@Factorial;|
  function( n ) ... end
  !gapprompt@gap>| !gapinput@IsOperation(Factorial);|
  false
\end{Verbatim}
 }

 

\subsection{\textcolor{Chapter }{TypeOfOperation}}
\logpage{[ 78, 1, 2 ]}\nobreak
\hyperdef{L}{X813BE52887A3E0FA}{}
{\noindent\textcolor{FuncColor}{$\triangleright$\enspace\texttt{TypeOfOperation({\mdseries\slshape object})\index{TypeOfOperation@\texttt{TypeOfOperation}}
\label{TypeOfOperation}
}\hfill{\scriptsize (function)}}\\


 returns a string from the list \texttt{[ "Attribute", "Operation", "Property", "Category", "Representation",
"Filter", "Setter"]} reflecting which type of operation \mbox{\texttt{\mdseries\slshape op}} is. 

 (see{\nobreakspace}\ref{Categories}, \ref{Representation}, \ref{Attributes}, \ref{Setter and Tester for Attributes}, \ref{Properties}, \ref{Other Filters}) }

 

\subsection{\textcolor{Chapter }{ShowDeclarationsOfOperation}}
\logpage{[ 78, 1, 3 ]}\nobreak
\hyperdef{L}{X838E29D485D2BAB9}{}
{\noindent\textcolor{FuncColor}{$\triangleright$\enspace\texttt{ShowDeclarationsOfOperation({\mdseries\slshape oper})\index{ShowDeclarationsOfOperation@\texttt{ShowDeclarationsOfOperation}}
\label{ShowDeclarationsOfOperation}
}\hfill{\scriptsize (function)}}\\


 Displays information about all declarations of the operation \mbox{\texttt{\mdseries\slshape oper}}, including the location of each declaration and the argument filters. 
\begin{Verbatim}[commandchars=!@|,fontsize=\small,frame=single,label=Example]
  !gapprompt@gap>| !gapinput@ShowDeclarationsOfOperation(IsFinite);|
  Available declarations for operation <Property "IsFinite">:
    1: GAPROOT/lib/coll.gd:1451 with 1 argument, and filters [ IsListOrCollection ]
    2: GAPROOT/lib/float.gd:212 with 1 argument, and filters [ IsFloat ]
    3: GAPROOT/lib/ctbl.gd:1195 with 1 argument, and filters [ IsNearlyCharacterTable ]
\end{Verbatim}
 }

 

\subsection{\textcolor{Chapter }{NewOperation}}
\logpage{[ 78, 1, 4 ]}\nobreak
\hyperdef{L}{X85A9E019795B79D6}{}
{\noindent\textcolor{FuncColor}{$\triangleright$\enspace\texttt{NewOperation({\mdseries\slshape name, args\texttt{\symbol{45}}filts})\index{NewOperation@\texttt{NewOperation}}
\label{NewOperation}
}\hfill{\scriptsize (function)}}\\


 \texttt{NewOperation} returns an operation \mbox{\texttt{\mdseries\slshape opr}} with name \mbox{\texttt{\mdseries\slshape name}}. The list \mbox{\texttt{\mdseries\slshape args\texttt{\symbol{45}}filts}} describes requirements about the arguments of \mbox{\texttt{\mdseries\slshape opr}}, namely the number of arguments must be equal to the length of \mbox{\texttt{\mdseries\slshape args\texttt{\symbol{45}}filts}}, and the $i$\texttt{\symbol{45}}th argument must lie in the filter \mbox{\texttt{\mdseries\slshape args\texttt{\symbol{45}}filts}}$[i]$. 

 Each method that is installed for \mbox{\texttt{\mdseries\slshape opr}} via \texttt{InstallMethod} (\ref{InstallMethod}) must require that the $i$\texttt{\symbol{45}}th argument lies in the filter \mbox{\texttt{\mdseries\slshape args\texttt{\symbol{45}}filts}}$[i]$. 

 One can install methods for other argument tuples via \texttt{InstallOtherMethod} (\ref{InstallOtherMethod}), this way it is also possible to install methods for a different number of
arguments than the length of \mbox{\texttt{\mdseries\slshape args\texttt{\symbol{45}}filts}}. }

 

\subsection{\textcolor{Chapter }{DeclareOperation}}
\logpage{[ 78, 1, 5 ]}\nobreak
\hyperdef{L}{X843F48137B899BC3}{}
{\noindent\textcolor{FuncColor}{$\triangleright$\enspace\texttt{DeclareOperation({\mdseries\slshape name, filters})\index{DeclareOperation@\texttt{DeclareOperation}}
\label{DeclareOperation}
}\hfill{\scriptsize (function)}}\\


 does the same as \texttt{NewOperation} (\ref{NewOperation}) and then binds the new operation to the global variable \mbox{\texttt{\mdseries\slshape name}}. The variable must previously be writable, and is made
read\texttt{\symbol{45}}only by this function. }

 
\subsection{\textcolor{Chapter }{Tag Based Operations}}\logpage{[ 78, 1, 6 ]}
\hyperdef{L}{X799081B4854DC003}{}
{
\noindent\textcolor{FuncColor}{$\triangleright$\enspace\texttt{NewTagBasedOperation({\mdseries\slshape name, requirements})\index{NewTagBasedOperation@\texttt{NewTagBasedOperation}}
\label{NewTagBasedOperation}
}\hfill{\scriptsize (function)}}\\
\noindent\textcolor{FuncColor}{$\triangleright$\enspace\texttt{DeclareTagBasedOperation({\mdseries\slshape name, requirements})\index{DeclareTagBasedOperation@\texttt{DeclareTagBasedOperation}}
\label{DeclareTagBasedOperation}
}\hfill{\scriptsize (function)}}\\
\noindent\textcolor{FuncColor}{$\triangleright$\enspace\texttt{InstallTagBasedMethod({\mdseries\slshape oper[, tag], meth})\index{InstallTagBasedMethod@\texttt{InstallTagBasedMethod}}
\label{InstallTagBasedMethod}
}\hfill{\scriptsize (function)}}\\


 \texttt{NewTagBasedOperation} returns an operation with name \mbox{\texttt{\mdseries\slshape name}} that is declared as \emph{tag based} w.r.t. the list \mbox{\texttt{\mdseries\slshape requirements}} of filters for its arguments. If an operation with name \mbox{\texttt{\mdseries\slshape name}} exists already before the call then this operation is returned, otherwise a
new operation gets created. 

 \texttt{DeclareTagBasedOperation} does the same and additionally binds the returned operation to the global
variable \mbox{\texttt{\mdseries\slshape name}} if the operation is new. 

 Declaring the operation \mbox{\texttt{\mdseries\slshape oper}} as tag based w.r.t. \mbox{\texttt{\mdseries\slshape requirements}} means that \texttt{InstallTagBasedMethod} can be used to install the method \mbox{\texttt{\mdseries\slshape meth}} for \mbox{\texttt{\mdseries\slshape oper}}, a function whose arguments satisfy \mbox{\texttt{\mdseries\slshape requirements}}, with the following meaning. 

 
\begin{itemize}
\item  The method \mbox{\texttt{\mdseries\slshape meth}} is applicable if the first argument of the call to \mbox{\texttt{\mdseries\slshape oper}} is identical (in the sense of \texttt{IsIdenticalObj} (\ref{IsIdenticalObj})) with the tag \mbox{\texttt{\mdseries\slshape tag}} that has been specified in the \texttt{InstallTagBasedMethod} call. 
\item  If none of the tag based methods for \mbox{\texttt{\mdseries\slshape oper}} has a \mbox{\texttt{\mdseries\slshape tag}} that is identical with the first argument of the call to \mbox{\texttt{\mdseries\slshape oper}} and if there is a tag based method for \mbox{\texttt{\mdseries\slshape oper}} for which no \mbox{\texttt{\mdseries\slshape tag}} was specified then this method is applicable. 
\end{itemize}
 

 Thus at most \emph{one} tag based method for \mbox{\texttt{\mdseries\slshape oper}} is applicable, and if a method without \mbox{\texttt{\mdseries\slshape tag}} has been installed then it serves as the default method. This is in contrast
to the situation with constructors (see \ref{Constructors}) where the first argument is a filter that is used as a tag, but several
methods can be applicable in a call to a constructor and one cannot define a
default method for it. 

 Typical use cases for tag based operations are operations that shall create
objects in particular internal representations; the filters that define these
representations are then used as the first argument, and one wants that either
the unique method that belongs to this filter or a default method is called. 

 Currently it is possible to declare an operation as tag based only for \emph{one} list of requirements. 

 Installing methods with \texttt{InstallMethod} (\ref{InstallMethod}) for a tag based operation is possible. However, installing such methods with
the same requirements as the ones for the tag based methods will have no
effect because the handling of tag based methods gets installed with \texttt{InstallEarlyMethod} (\ref{InstallEarlyMethod}) and thus has higher priority. }

 }

  
\section{\textcolor{Chapter }{Constructors}}\label{Constructors}
\logpage{[ 78, 2, 0 ]}
\hyperdef{L}{X86EC0F0A78ECBC10}{}
{
  Constructors are a special type of operation used to make new objects. The key
difference compared to regular operations is that method selection works
slightly differently for them: The first argument in a call to a constructor
must always be a filter \ref{Filters}. The result of a method is expected to lie in that filter. Moreover, while
normally method selection matches on the type of each argument, for a
constructor the first argument is treated differently: instead of matching its
type, the argument itself (which is a filter) must be a subset of the filter
specified by the method which is being tested for match. In other words, the
argument filter must imply the method filter. 

 Also, method ranking works differently: instead of the sum of the ranks of the
types of all arguments, only the rank of the filter given as first argument is
considered; and for it, the normal ranking order is reversed. This ensures
that if multiple constructor methods match, the ``most general'' method is selected. 
\begin{Verbatim}[commandchars=!@|,fontsize=\small,frame=single,label=Example]
  !gapprompt@gap>| !gapinput@DeclareConstructor("XCons",[IsMagma,IsInt]);|
  !gapprompt@gap>| !gapinput@InstallMethod(XCons, [IsGroup, IsInt], function(t,x) return CyclicGroup(x); end);|
  !gapprompt@gap>| !gapinput@InstallMethod(XCons, [IsPermGroup, IsInt], function(t,x) return SymmetricGroup(x); end);|
  !gapprompt@gap>| !gapinput@InstallMethod(XCons, [IsSemigroup, IsInt], function(t,x) return FullTransformationMonoid(x); end);|
  !gapprompt@gap>| !gapinput@XCons(IsGroup,3);|
  <pc group of size 3 with 1 generator>
  !gapprompt@gap>| !gapinput@XCons(IsPermGroup,3);|
  Sym( [ 1 .. 3 ] )
  !gapprompt@gap>| !gapinput@XCons(IsSemigroup,4);|
  <full transformation monoid of degree 4>
  !gapprompt@gap>| !gapinput@# if multiple methods match, the most general is selected:|
  !gapprompt@gap>| !gapinput@XCons(IsMagma,3);|
  <full transformation monoid of degree 3>
\end{Verbatim}
 The example above shows some basic examples (usually a constructor will
produce isomorphic objects in different representations, not different objects
as in this case). 

 If no method has been installed which guarantees to produce a suitable
objecty, a "No Method Found" error will be returned. 
\begin{Verbatim}[commandchars=@|A,fontsize=\small,frame=single,label=Example]
  @gapprompt|gap>A @gapinput|XCons(IsFullTransformationMonoid,4);A
  Error, no method found! For debugging hints type ?Recovery from NoMethodFound
  Error, no 1st choice method found for `XCons' on 2 arguments
  called from read-eval loop at *stdin*:8
  you can enter 'quit;' to quit to outer loop
  @gapbrkprompt|brk>A @gapinput|quit;A
  @gapprompt|gap>A @gapinput|XCons(IsNilpotentGroup,4);A
  Error, no method found! For debugging hints type ?Recovery from NoMethodFound
  Error, no 1st choice method found for `XCons' on 2 arguments
  called from read-eval loop at *stdin*:8
  you can enter 'quit;' to quit to outer loop
  @gapbrkprompt|brk>A @gapinput|A
\end{Verbatim}
 Note that in both these cases there are methods that actually produce results
of the required types, but they have not been installed with this information,
so are not selected. 

\subsection{\textcolor{Chapter }{NewConstructor}}
\logpage{[ 78, 2, 1 ]}\nobreak
\hyperdef{L}{X783FA45E7858A8CF}{}
{\noindent\textcolor{FuncColor}{$\triangleright$\enspace\texttt{NewConstructor({\mdseries\slshape name, args\texttt{\symbol{45}}filts})\index{NewConstructor@\texttt{NewConstructor}}
\label{NewConstructor}
}\hfill{\scriptsize (function)}}\\


 \texttt{NewConstructor} returns a constructor \mbox{\texttt{\mdseries\slshape cons}} with name \mbox{\texttt{\mdseries\slshape name}}. The list \mbox{\texttt{\mdseries\slshape args\texttt{\symbol{45}}filts}} describes requirements about the arguments of \mbox{\texttt{\mdseries\slshape cons}}. Namely the number of arguments must be equal to the length of \mbox{\texttt{\mdseries\slshape args\texttt{\symbol{45}}filts}}, and the $i$\texttt{\symbol{45}}th argument must lie in the filter \mbox{\texttt{\mdseries\slshape args\texttt{\symbol{45}}filts}}$[i]$ for $i \neq 1$. A constructor expects the first argument to be a \emph{filter} instead of an object and it must be a subset of the filter \mbox{\texttt{\mdseries\slshape args\texttt{\symbol{45}}filts}}$[1]$. 

 Each method that is installed for \mbox{\texttt{\mdseries\slshape cons}} via \texttt{InstallMethod} (\ref{InstallMethod}) must require that the $i$\texttt{\symbol{45}}th argument lies in the filter \mbox{\texttt{\mdseries\slshape args\texttt{\symbol{45}}filts}}$[i]$ for $i \neq 1$. Its first argument is a filter and must be a subset of the filter \mbox{\texttt{\mdseries\slshape args\texttt{\symbol{45}}filts}}$[1]$. 

 One can install methods for other argument tuples via \texttt{InstallOtherMethod} (\ref{InstallOtherMethod}), this way it is also possible to install methods for a different number of
arguments than the length of \mbox{\texttt{\mdseries\slshape args\texttt{\symbol{45}}filts}}. 

 Note that the method selection for constructors works slightly differently
than for usual operations. As stated above, applicabilty to the first argument
in an argument tuple is tested by determining whether the
argument\texttt{\symbol{45}}filter is a \emph{subset} of \mbox{\texttt{\mdseries\slshape args\texttt{\symbol{45}}filts}}$[1]$. 

 The rank of a method installed for a constructor is determined solely by \mbox{\texttt{\mdseries\slshape args\texttt{\symbol{45}}filts}}$[1]$ of the method. Instead of taking the sum of the ranks of filters involved in
its \mbox{\texttt{\mdseries\slshape args\texttt{\symbol{45}}filts}}$[1]$, the sum of $-1$ times these values is taken. The result is added to the number \mbox{\texttt{\mdseries\slshape val}} used in the call of \texttt{InstallMethod} (\ref{InstallMethod}). 

 This has the following effects on the method selection for constructors. If \mbox{\texttt{\mdseries\slshape cons}} is called with an argument tuple whose first argument is the filter \mbox{\texttt{\mdseries\slshape filt}}, any method whose first argument is \emph{more} specific than \mbox{\texttt{\mdseries\slshape filt}} is applicable (if its other \mbox{\texttt{\mdseries\slshape args\texttt{\symbol{45}}filts}} also match). Then the method with the ``most general'' filter \mbox{\texttt{\mdseries\slshape args\texttt{\symbol{45}}filts}}$[1]$ is chosen, since the rank is computed by taking $-1$ times the ranks of the involved filters. Thus, a constructor is chosen which
returns an object in \mbox{\texttt{\mdseries\slshape filt}} using as few extra filters as possible, which presumably is both more flexible
to use and easier to construct. 

 The following example showcases this behaviour. Note that the argument \mbox{\texttt{\mdseries\slshape filter}} is only used for method dispatch. 
\begin{Verbatim}[commandchars=!@|,fontsize=\small,frame=single,label=Example]
  DeclareFilter( "IsMyObj" );
  DeclareFilter( "IsMyFilter" );
  DeclareFilter( "IsMyOtherFilter" );
  BindGlobal( "MyFamily", NewFamily( "MyFamily" ) );
  
  DeclareConstructor( "NewMyObj", [ IsMyObj ] );
  
  InstallMethod( NewMyObj,
  [ IsMyObj ],
  function( filter )
      local type;
      Print("General constructor\n");
      type := NewType( MyFamily, IsMyObj );
      return Objectify( type, [] );
  end );
  InstallMethod( NewMyObj,
  [ IsMyObj and IsMyFilter and IsMyOtherFilter ],
  function( filter )
      local type;
      Print("Special constructor\n");
      type := NewType( MyFamily, IsMyObj and IsMyFilter and IsMyOtherFilter );
      return Objectify( type, [] );
  end );
\end{Verbatim}
 If only IsMyObj is given, both methods are applicable and the general
constructor is called. If also IsMyFilter is given, only the special
constructor is applicable. 
\begin{Verbatim}[commandchars=!@|,fontsize=\small,frame=single,label=Example]
  !gapprompt@gap>| !gapinput@a := NewMyObj( IsMyObj );;|
  General constructor
  !gapprompt@gap>| !gapinput@IsMyOtherFilter(a);|
  false
  !gapprompt@gap>| !gapinput@b := NewMyObj( IsMyObj and IsMyFilter );;|
  Special constructor
  !gapprompt@gap>| !gapinput@IsMyOtherFilter(b);|
  true
  !gapprompt@gap>| !gapinput@c := NewMyObj( IsMyObj and IsMyFilter and IsMyOtherFilter );;|
  Special constructor
  !gapprompt@gap>| !gapinput@IsMyOtherFilter(c);|
  true
\end{Verbatim}
 }

 

\subsection{\textcolor{Chapter }{DeclareConstructor}}
\logpage{[ 78, 2, 2 ]}\nobreak
\hyperdef{L}{X7EB6830886F62CC0}{}
{\noindent\textcolor{FuncColor}{$\triangleright$\enspace\texttt{DeclareConstructor({\mdseries\slshape name, filters})\index{DeclareConstructor@\texttt{DeclareConstructor}}
\label{DeclareConstructor}
}\hfill{\scriptsize (function)}}\\


 does the same as \texttt{NewConstructor} (\ref{NewConstructor}) and then binds the result to the global variable \mbox{\texttt{\mdseries\slshape name}}. The variable must previously be writable, and is made
read\texttt{\symbol{45}}only by this function. 

 Note that for operations which are constructors special rules with respect to
applicability and rank of the corresponding methods apply (see section \texttt{NewConstructor} (\ref{NewConstructor})). }

 }

  
\section{\textcolor{Chapter }{Method Installation}}\label{Method Installation}
\logpage{[ 78, 3, 0 ]}
\hyperdef{L}{X795EE8257848B438}{}
{
  In order to describe what it means to select a method of an operation, we must
describe how the methods are connected to their operations. 

 For attributes and properties there is \texttt{InstallImmediateMethod} (\ref{InstallImmediateMethod}). 

 For declaring that a filter is implied by other filters there is \texttt{InstallTrueMethod} (\ref{InstallTrueMethod}). 

\subsection{\textcolor{Chapter }{InstallMethod}}
\logpage{[ 78, 3, 1 ]}\nobreak
\hyperdef{L}{X837EFDAB7BEF290B}{}
{\noindent\textcolor{FuncColor}{$\triangleright$\enspace\texttt{InstallMethod({\mdseries\slshape opr[, info][, famp], args\texttt{\symbol{45}}filts[, val], method})\index{InstallMethod@\texttt{InstallMethod}}
\label{InstallMethod}
}\hfill{\scriptsize (function)}}\\


 installs a function method \mbox{\texttt{\mdseries\slshape method}} for the operation \mbox{\texttt{\mdseries\slshape opr}}; \mbox{\texttt{\mdseries\slshape args\texttt{\symbol{45}}filts}} should be a list of requirements for the arguments, each entry being a filter;
if supplied \mbox{\texttt{\mdseries\slshape info}} should be a short but informative string that describes for what situation the
method is installed, \mbox{\texttt{\mdseries\slshape famp}} should be a function to be applied to the families of the arguments. \mbox{\texttt{\mdseries\slshape val}} should be an integer that measures the priority of the method, or a function
of no arguments which should return such an integer and will be called each
time method order is being recalculated (see \texttt{InstallTrueMethod} (\ref{InstallTrueMethod})). 

 The default values for \mbox{\texttt{\mdseries\slshape info}}, \mbox{\texttt{\mdseries\slshape famp}}, and \mbox{\texttt{\mdseries\slshape val}} are the empty string, the function \texttt{ReturnTrue} (\ref{ReturnTrue}), and the integer zero, respectively. 

 The exact meaning of the arguments \mbox{\texttt{\mdseries\slshape famp}}, \mbox{\texttt{\mdseries\slshape args\texttt{\symbol{45}}filts}}, and \mbox{\texttt{\mdseries\slshape val}} is explained in Section{\nobreakspace}\ref{Applicable Methods and Method Selection}. 

 \mbox{\texttt{\mdseries\slshape opr}} expects its methods to require certain filters for their arguments. For
example, the argument of a method for the operation \texttt{Zero} (\ref{Zero}) must be in the category \texttt{IsAdditiveElementWithZero} (\ref{IsAdditiveElementWithZero}). It is not possible to use \texttt{InstallMethod} to install a method for which the entries of \mbox{\texttt{\mdseries\slshape args\texttt{\symbol{45}}filts}} do not imply the respective requirements of the operation \mbox{\texttt{\mdseries\slshape opr}}. If one wants to override this restriction, one has to use \texttt{InstallOtherMethod} (\ref{InstallOtherMethod}) instead. }

 

\subsection{\textcolor{Chapter }{InstallOtherMethod}}
\logpage{[ 78, 3, 2 ]}\nobreak
\hyperdef{L}{X7D2C12DB841CE539}{}
{\noindent\textcolor{FuncColor}{$\triangleright$\enspace\texttt{InstallOtherMethod({\mdseries\slshape opr[, info][, famp], args\texttt{\symbol{45}}filts[, val], method})\index{InstallOtherMethod@\texttt{InstallOtherMethod}}
\label{InstallOtherMethod}
}\hfill{\scriptsize (function)}}\\


 installs a function method \mbox{\texttt{\mdseries\slshape method}} for the operation \mbox{\texttt{\mdseries\slshape opr}}, in the same way as for \texttt{InstallMethod} (\ref{InstallMethod}), but without the restriction that the number of arguments must match a
declaration of \mbox{\texttt{\mdseries\slshape opr}} and without the restriction that \mbox{\texttt{\mdseries\slshape args\texttt{\symbol{45}}filts}} imply the respective requirements of the operation \mbox{\texttt{\mdseries\slshape opr}}. }

 

\subsection{\textcolor{Chapter }{InstallEarlyMethod}}
\logpage{[ 78, 3, 3 ]}\nobreak
\hyperdef{L}{X7D54FB8A80C0AC3C}{}
{\noindent\textcolor{FuncColor}{$\triangleright$\enspace\texttt{InstallEarlyMethod({\mdseries\slshape opr, method})\index{InstallEarlyMethod@\texttt{InstallEarlyMethod}}
\label{InstallEarlyMethod}
}\hfill{\scriptsize (function)}}\\


 installs a special "early" function method \mbox{\texttt{\mdseries\slshape method}} for the operation \mbox{\texttt{\mdseries\slshape opr}}. An early method is special in that it bypasses method dispatch, and is
always the first method to be called when invoking the operation. 

 This can be used to avoid method selection overhead for certain special cases,
i.e., as an optimization. Overall, we recommend to use this feature very
sparingly, as it is tool with sharp edges: for example, any inputs that are
handled by an early method can not be intercepted by a regular method, no
matter how high its rank is; this can preclude other kinds of optimizations. 

 Also, unlike regular methods, no checks are performed on the arguments. Not
even the required filters for the operation are tested, so early methods must
be careful in validating their inputs. This also means that any operation can
have at most one such early method for each arity (i.e., one early method
taking 1 argument, one early method taking 2 arguments, etc.). 

 If an early method determines that it is not applicable, it can resume regular
method dispatch by invoking \texttt{TryNextMethod} (\ref{TryNextMethod}). 

 For an example application of early methods, they are used by \texttt{First} (\ref{First}) to deal with internal lists, for which computing the exact type (needed for
method selection) can be very expensive. }

 

\subsection{\textcolor{Chapter }{InstallMethodWithRandomSource}}
\logpage{[ 78, 3, 4 ]}\nobreak
\hyperdef{L}{X78CA646678B0539F}{}
{\noindent\textcolor{FuncColor}{$\triangleright$\enspace\texttt{InstallMethodWithRandomSource({\mdseries\slshape opr, info[, famp], args\texttt{\symbol{45}}filts[, val], method})\index{InstallMethodWithRandomSource@\texttt{InstallMethodWithRandomSource}}
\label{InstallMethodWithRandomSource}
}\hfill{\scriptsize (function)}}\\
\noindent\textcolor{FuncColor}{$\triangleright$\enspace\texttt{InstallOtherMethodWithRandomSource({\mdseries\slshape opr, info[, famp], args\texttt{\symbol{45}}filts[, val], method})\index{InstallOtherMethodWithRandomSource@\texttt{InstallOtherMethodWithRandomSource}}
\label{InstallOtherMethodWithRandomSource}
}\hfill{\scriptsize (function)}}\\


 These functions are designed to simplify adding new methods for \texttt{Random} (\ref{Random:for a list or collection}), \texttt{PseudoRandom} (\ref{PseudoRandom}), and \texttt{Randomize} (\ref{Randomize:for a vector object}) to \textsf{GAP} which can be called both with, and without, a random source. 

 They accept the same arguments as \texttt{InstallMethod} (\ref{InstallMethod}) and \texttt{InstallOtherMethod} (\ref{InstallOtherMethod}), with the extra requirement that the first member of \mbox{\texttt{\mdseries\slshape args\texttt{\symbol{45}}filts}} must be \texttt{IsRandomSource} (\ref{IsRandomSource}), and the \mbox{\texttt{\mdseries\slshape info}} argument is compulsory and must begin 'for a random source and'. 

 This function then installs two methods: first it calls \texttt{InstallMethod} (\ref{InstallMethod}) (or \texttt{InstallOtherMethod} (\ref{InstallOtherMethod})) with unchanged arguments. Then it calls \texttt{InstallMethod} (\ref{InstallMethod}) (or \texttt{InstallOtherMethod} (\ref{InstallOtherMethod})) a second time to install another method which lacks the initial random
source argument; this additional method simply invokes the original method,
with \texttt{GlobalMersenneTwister} (\ref{GlobalMersenneTwister}) added as first argument. }

 }

  
\section{\textcolor{Chapter }{Applicable Methods and Method Selection}}\label{Applicable Methods and Method Selection}
\logpage{[ 78, 4, 0 ]}
\hyperdef{L}{X851FC6387CA2B241}{}
{
  A method installed as above is \emph{applicable} for an arguments tuple if the following conditions are satisfied. 

 The number of arguments equals the length of the list \mbox{\texttt{\mdseries\slshape args\texttt{\symbol{45}}filts}}, the $i$\texttt{\symbol{45}}th argument lies in the filter \mbox{\texttt{\mdseries\slshape args\texttt{\symbol{45}}filts}}$[i]$, and \mbox{\texttt{\mdseries\slshape famp}} returns \texttt{true} when applied to the families of the arguments. The maximal number of arguments
supported for methods is six, one gets an error message if one tries to
install a method with at least seven arguments. 

 So \mbox{\texttt{\mdseries\slshape args\texttt{\symbol{45}}filt}} describes conditions for each argument, and \mbox{\texttt{\mdseries\slshape famp}} describes a relation between the arguments. 

 For unary operations such as attributes and properties, there is no such
relation to postulate, \mbox{\texttt{\mdseries\slshape famp}} is \texttt{ReturnTrue} (\ref{ReturnTrue}) for these operations, a function that always returns \texttt{true}. For binary operations, the usual value of \mbox{\texttt{\mdseries\slshape famp}} is \texttt{IsIdenticalObj} (\ref{IsIdenticalObj}), which means that both arguments must lie in the same family. 

 Note that any properties which occur among the filters in the filter list will \emph{not} be tested by the method selection if they are not yet known. (More exact: if \mbox{\texttt{\mdseries\slshape prop}} is a property then the filter implicitly uses not \mbox{\texttt{\mdseries\slshape prop}} but \texttt{Has\mbox{\texttt{\mdseries\slshape prop}} and \mbox{\texttt{\mdseries\slshape prop}}}.) If this is desired you must explicitly enforce a test (see
section{\nobreakspace}\ref{Redispatching}) below. 

 If no method is applicable, the error message ``no method found'' is signaled. 

 Otherwise, the applicable method with highest \emph{rank} is selected and then called. This rank is usually given by the sum of the
ranks of the filters in the list \mbox{\texttt{\mdseries\slshape args\texttt{\symbol{45}}filt}}, \emph{including involved filters}. This value can be modified by the optional \mbox{\texttt{\mdseries\slshape val}} argument used in the call of \texttt{InstallMethod} (\ref{InstallMethod}), as follows. 

 
\begin{itemize}
\item  If \mbox{\texttt{\mdseries\slshape val}} is an integer then add this value to the rank computed from the filters. 
\item  If \mbox{\texttt{\mdseries\slshape val}} is a function with zero arguments then add the value returned by this function
to the rank computed from the filters. 
\item  If \mbox{\texttt{\mdseries\slshape val}} is a list of the form \texttt{[ filters ]} or \texttt{[ filters, abs ]}, where \texttt{filters} is a list of filters and \texttt{abs}, if given, is an integer, then ignore the filters in \mbox{\texttt{\mdseries\slshape args\texttt{\symbol{45}}filt}} for the rank computation, and instead take the sum of the filters in \texttt{filters} (including involved filters) plus \texttt{abs}. 
\end{itemize}
 

 So the argument \mbox{\texttt{\mdseries\slshape val}} can be used to raise or lower the priority of a method relative to other
methods for \mbox{\texttt{\mdseries\slshape opr}}. The variant where \mbox{\texttt{\mdseries\slshape val}} is a function can be useful if one wants to shift the rank of the method by
the rank of certain filters; note that such ranks need not be constant during
a \textsf{GAP} session, they can grow due to implications, see \texttt{InstallTrueMethod} (\ref{InstallTrueMethod}). The variant where \mbox{\texttt{\mdseries\slshape val}} is a list can be useful for example if one wants that certain filters that are
involved in \mbox{\texttt{\mdseries\slshape args\texttt{\symbol{45}}filt}} are required but shall not count for the rank of the method. 

 Note that for operations which are constructors special rules with respect to
applicability and rank of the corresponding methods apply (see \texttt{NewConstructor} (\ref{NewConstructor})). 

 Note that from the applicable methods an efficient one shall be selected. This
is a method that needs only little time and storage for the computations. 

 It seems to be impossible for \textsf{GAP} to select an optimal method in all cases. The present ranking of methods is
based on the assumption that a method installed for a special situation shall
be preferred to a method installed for a more general situation. 

 For example, a method for computing a Sylow subgroup of a nilpotent group is
expected to be more efficient than a method for arbitrary groups. So the more
specific method will be selected if \textsf{GAP} knows that the group given as argument is nilpotent. 

 Of course there is no obvious way to decide between the efficiency of
incommensurable methods. For example, take an operation with one method for
permutation groups, another method for nilpotent groups, but no method for
nilpotent permutation groups, and call this operation with a permutation group
known to be nilpotent. }

  
\section{\textcolor{Chapter }{Partial Methods}}\label{Partial Methods}
\logpage{[ 78, 5, 0 ]}
\hyperdef{L}{X846865A681D4A623}{}
{
  

\subsection{\textcolor{Chapter }{TryNextMethod}}
\logpage{[ 78, 5, 1 ]}\nobreak
\hyperdef{L}{X7EED949B83046A7F}{}
{\noindent\textcolor{FuncColor}{$\triangleright$\enspace\texttt{TryNextMethod({\mdseries\slshape })\index{TryNextMethod@\texttt{TryNextMethod}}
\label{TryNextMethod}
}\hfill{\scriptsize (function)}}\\


 After a method has been selected and called, the method may recognize that it
cannot compute the desired result, and give up by calling \texttt{TryNextMethod()}. 

 In effect, the execution of the method is terminated, and the method selection
calls the next method that is applicable w.r.t. the original arguments. In
other words, the applicable method is called that is subsequent to the one
that called \texttt{TryNextMethod}, according to decreasing rank of the methods. 

 For example, since every finite group of odd order is solvable, one may
install a method for the property \texttt{IsSolvableGroup} (\ref{IsSolvableGroup}) that checks whether the size of the argument is an odd integer, returns \texttt{true} if so, and gives up otherwise. 

 Care is needed if a partial method might modify the type of one of its
arguments, for example by computing an attribute or property. If this happens,
and the type has really changed, then the method should not exit using \texttt{TryNextMethod()} but should call the operation again, as the new information in the type may
cause some methods previously judged inapplicable to be applicable. For
example, if the above method for \texttt{IsSolvableGroup} (\ref{IsSolvableGroup}) actually computes the size, (rather than just examining a stored size), then
it must take care to check whether the type of the group has changed. }

 }

  
\section{\textcolor{Chapter }{Redispatching}}\label{Redispatching}
\logpage{[ 78, 6, 0 ]}
\hyperdef{L}{X7B85DD797A907106}{}
{
  As mentioned above the method selection will not test unknown properties. In
situations, in which algorithms are only known (or implemented) under certain
conditions, however such a test might be actually desired. 

 One way to achieve this would be to install the method under weaker conditions
and explicitly test the properties first, exiting via \texttt{TryNextMethod} (\ref{TryNextMethod}) if some of them are not fulfilled. A problem of this approach however is that
such methods then automatically are ranked lower and that the code does not
look nice. 

 A much better way is to use redispatching: Before deciding that no method has
been found one tests these properties and if they turn out to be true the
method selection is started anew (and will then find a method). 

 This can be achieved via the following function: 

\subsection{\textcolor{Chapter }{RedispatchOnCondition}}
\logpage{[ 78, 6, 1 ]}\nobreak
\hyperdef{L}{X7D4A46CE7BCFCCF5}{}
{\noindent\textcolor{FuncColor}{$\triangleright$\enspace\texttt{RedispatchOnCondition({\mdseries\slshape oper[, info], fampred, reqs, cond, val})\index{RedispatchOnCondition@\texttt{RedispatchOnCondition}}
\label{RedispatchOnCondition}
}\hfill{\scriptsize (function)}}\\


 This function installs a method for the operation \mbox{\texttt{\mdseries\slshape oper}} under the conditions \mbox{\texttt{\mdseries\slshape fampred}} and \mbox{\texttt{\mdseries\slshape reqs}} which has absolute value \mbox{\texttt{\mdseries\slshape val}}; that is, the value of the filters \mbox{\texttt{\mdseries\slshape reqs}} is disregarded. \mbox{\texttt{\mdseries\slshape cond}} is a list of filters. If not all the values of properties involved in these
filters are already known for actual arguments of the method, they are
explicitly tested and if they are fulfilled \emph{and} stored after this test, the operation is dispatched again. Otherwise the
method exits with \texttt{TryNextMethod} (\ref{TryNextMethod}). If supplied, \mbox{\texttt{\mdseries\slshape info}} should be a short but informative string that describes these conditions. This
can be used to enforce tests like \texttt{IsFinite} (\ref{IsFinite}) in situations when all existing methods require this property. The list \mbox{\texttt{\mdseries\slshape cond}} may have unbound entries in which case the corresponding argument is ignored
for further tests. }

 }

  
\section{\textcolor{Chapter }{Immediate Methods}}\label{Immediate Methods}
\logpage{[ 78, 7, 0 ]}
\hyperdef{L}{X87D38D2584D0A8AF}{}
{
  Usually a method is called only if its operation has been called and if this
method has been selected, see \texttt{InstallMethod} (\ref{InstallMethod}). 

 For attributes and properties, one can install also \emph{immediate methods}. 

\subsection{\textcolor{Chapter }{InstallImmediateMethod}}
\logpage{[ 78, 7, 1 ]}\nobreak
\hyperdef{L}{X87B47AC0849611F8}{}
{\noindent\textcolor{FuncColor}{$\triangleright$\enspace\texttt{InstallImmediateMethod({\mdseries\slshape opr[, info], filter, rank, method})\index{InstallImmediateMethod@\texttt{InstallImmediateMethod}}
\label{InstallImmediateMethod}
}\hfill{\scriptsize (function)}}\\


 \texttt{InstallImmediateMethod} installs \mbox{\texttt{\mdseries\slshape method}} as an immediate method for \mbox{\texttt{\mdseries\slshape opr}}, which must be an attribute or a property, with requirement \mbox{\texttt{\mdseries\slshape filter}} and rank \mbox{\texttt{\mdseries\slshape rank}} (the rank can be omitted, in which case 0 is used as rank). The rank must be
an integer value that measures the priority of \mbox{\texttt{\mdseries\slshape method}} among the immediate methods for \mbox{\texttt{\mdseries\slshape opr}}. If supplied, \mbox{\texttt{\mdseries\slshape info}} should be a short but informative string that describes the situation in which
the method is called. 

 An immediate method is called automatically as soon as the object lies in \mbox{\texttt{\mdseries\slshape filter}}, provided that the value is not yet known. Afterwards the attribute setter is
called in order to store the value, unless the method exits via \texttt{TryNextMethod} (\ref{TryNextMethod}). 

 Note the difference to \texttt{InstallMethod} (\ref{InstallMethod}) that no family predicate occurs because \mbox{\texttt{\mdseries\slshape opr}} expects only one argument, and that \mbox{\texttt{\mdseries\slshape filter}} is not a list of requirements but the argument requirement itself. 

 Immediate methods are thought of as a possibility for objects to gain useful
knowledge. They must not be used to force the storing of ``defining information'' in an object. In other words, \textsf{GAP} should work even if all immediate methods are completely disabled. Therefore,
the call to \texttt{InstallImmediateMethod} installs \mbox{\texttt{\mdseries\slshape method}} also as an ordinary method for \mbox{\texttt{\mdseries\slshape opr}} with requirement \mbox{\texttt{\mdseries\slshape filter}}. 

 Note that in such a case \textsf{GAP} executes a computation for which it was not explicitly asked by the user. So
one should install only those methods as immediate methods that are \emph{extremely cheap}. To emphasize this, immediate methods are also called \emph{zero cost methods}. The time for their execution should really be approximately zero. 

 For example, the size of a permutation group can be computed very cheaply if a
stabilizer chain of the group is known. So it is reasonable to install an
immediate method for \texttt{Size} (\ref{Size}) with requirement \texttt{IsGroup and Tester( \mbox{\texttt{\mdseries\slshape stab}} )}, where \mbox{\texttt{\mdseries\slshape stab}} is the attribute corresponding to the stabilizer chain. 

 Another example would be the implementation of the conclusion that every
finite group of prime power order is nilpotent. This could be done by
installing an immediate method for the attribute \texttt{IsNilpotentGroup} (\ref{IsNilpotentGroup}) with requirement \texttt{IsGroup and Tester( Size )}. This method would then check whether the size is a finite prime power,
return \texttt{true} in this case and otherwise call \texttt{TryNextMethod} (\ref{TryNextMethod}). But this requires factoring of an integer, which cannot be guaranteed to be
very cheap, so one should not install this method as an immediate method. }

 

\subsection{\textcolor{Chapter }{IsNoImmediateMethodsObject}}
\logpage{[ 78, 7, 2 ]}\nobreak
\hyperdef{L}{X80C3A19E7B99BB41}{}
{\noindent\textcolor{FuncColor}{$\triangleright$\enspace\texttt{IsNoImmediateMethodsObject({\mdseries\slshape obj})\index{IsNoImmediateMethodsObject@\texttt{IsNoImmediateMethodsObject}}
\label{IsNoImmediateMethodsObject}
}\hfill{\scriptsize (filter)}}\\


 If this filter is set immediate methods will be ignored for \mbox{\texttt{\mdseries\slshape obj}}. This can be crucial for performance for objects like pcgs (see Section \ref{Polycyclic Generating Systems}), of which many are created, which are collections, but for which all those
immediate methods for \texttt{IsTrivial} (\ref{IsTrivial}) et cetera do not really make sense. Other examples of objects in \texttt{IsNoImmediateMethodsObject} are compressed vectors and matrices over small finite fields, see the sections \ref{Row Vectors over Finite Fields} and \ref{Matrices over Finite Fields}. }

 }

  
\section{\textcolor{Chapter }{Logical Implications}}\label{Logical Implications}
\logpage{[ 78, 8, 0 ]}
\hyperdef{L}{X7FB5016E83DB4349}{}
{
  

\subsection{\textcolor{Chapter }{InstallTrueMethod}}
\logpage{[ 78, 8, 1 ]}\nobreak
\hyperdef{L}{X860B8B707995CFE3}{}
{\noindent\textcolor{FuncColor}{$\triangleright$\enspace\texttt{InstallTrueMethod({\mdseries\slshape newfil, filt})\index{InstallTrueMethod@\texttt{InstallTrueMethod}}
\label{InstallTrueMethod}
}\hfill{\scriptsize (function)}}\\


 It may happen that a filter \mbox{\texttt{\mdseries\slshape newfil}} shall be implied by another filter \mbox{\texttt{\mdseries\slshape filt}}, which is usually a meet of other properties, or the meet of some properties
and some categories. Such a logical implication can be installed as an ``immediate method'' for \mbox{\texttt{\mdseries\slshape newfil}} that requires \mbox{\texttt{\mdseries\slshape filt}} and that always returns \texttt{true}. (This should not be mixed up with the methods installed via \texttt{InstallImmediateMethod} (\ref{InstallImmediateMethod}), which have to be called at runtime for the actual objects.) 

 \texttt{InstallTrueMethod} has the effect that \mbox{\texttt{\mdseries\slshape newfil}} becomes an implied filter of \mbox{\texttt{\mdseries\slshape filt}}, see{\nobreakspace}\ref{Filters}. 

 For example, each cyclic group is abelian, each finite vector space is finite
dimensional, and each division ring is integral. The first of these
implications is installed as follows. 

 
\begin{Verbatim}[commandchars=!@|,fontsize=\small,frame=single,label=Example]
  InstallTrueMethod( IsCommutative, IsGroup and IsCyclic );
\end{Verbatim}
 

 Contrary to the immediate methods installed with \texttt{InstallImmediateMethod} (\ref{InstallImmediateMethod}), logical implications cannot be switched off. This means that after the above
implication has been installed, one can rely on the fact that every object in
the filter \texttt{IsGroup and IsCyclic} whose type gets created after the installation of the implication will also be
in the filter \texttt{IsCommutative} (\ref{IsCommutative}). In particular, it may happen that an object which existed already before the
installation of the implication lies in \texttt{IsGroup and IsCyclic} but not in \texttt{IsCommutative} (\ref{IsCommutative}). Thus it is advisable to install all implications between filters before one
starts creating (types of) objects lying in these filters. 

 Adding logical implications can change the rank of filters (see \texttt{RankFilter} (\ref{RankFilter})) and consequently the rank, and so choice of methods for operations (see \ref{Applicable Methods and Method Selection}). By default \texttt{InstallTrueMethod} adjusts the method selection data structures to take care of this, but this
process can be time\texttt{\symbol{45}}consuming, so functions \texttt{SuspendMethodReordering} (\ref{SuspendMethodReordering}) and \texttt{ResumeMethodReordering} (\ref{ResumeMethodReordering}) are provided to allow control of this process. }

 

\subsection{\textcolor{Chapter }{SuspendMethodReordering}}
\logpage{[ 78, 8, 2 ]}\nobreak
\hyperdef{L}{X7B26BDF68754DF7A}{}
{\noindent\textcolor{FuncColor}{$\triangleright$\enspace\texttt{SuspendMethodReordering({\mdseries\slshape })\index{SuspendMethodReordering@\texttt{SuspendMethodReordering}}
\label{SuspendMethodReordering}
}\hfill{\scriptsize (function)}}\\
\noindent\textcolor{FuncColor}{$\triangleright$\enspace\texttt{ResumeMethodReordering({\mdseries\slshape })\index{ResumeMethodReordering@\texttt{ResumeMethodReordering}}
\label{ResumeMethodReordering}
}\hfill{\scriptsize (function)}}\\
\noindent\textcolor{FuncColor}{$\triangleright$\enspace\texttt{ResetMethodReordering({\mdseries\slshape })\index{ResetMethodReordering@\texttt{ResetMethodReordering}}
\label{ResetMethodReordering}
}\hfill{\scriptsize (function)}}\\


 These functions control whether the method reordering process described in \texttt{InstallTrueMethod} (\ref{InstallTrueMethod}) is invoked or not. Since this process can be comparatively
time\texttt{\symbol{45}}consuming, it is usually suspended when a lot of
implications are due to be installed, for instance when loading the library,
or a package. This is done by calling \texttt{SuspendMethodReordering} once the installations are done, \texttt{ResumeMethodReordering} should be called. These pairs of calls can be nested. When the outermost pair
is complete, method reordering takes place and is enabled in \texttt{InstallTrueMethod} (\ref{InstallTrueMethod}) thereafter. \texttt{ResetMethodReordering} effectively exits all nested suspensions, resuming reordering immediately.
This function is mainly provided for error recovery and similar purposes and
is called on quitting from a break loop. }

 }

  
\section{\textcolor{Chapter }{Operations and Mathematical Terms}}\label{Operations and Mathematical Terms}
\logpage{[ 78, 9, 0 ]}
\hyperdef{L}{X855FE25783FB0D4E}{}
{
  \index{overload} Usually an operation stands for a mathematical concept, and the name of the
operation describes this uniquely. Examples are the property \texttt{IsFinite} (\ref{IsFinite}) and the attribute \texttt{Size} (\ref{Size}). But there are cases where the same mathematical term is used to denote
different concepts, for example \texttt{Degree} is defined for polynomials, group characters, and permutation actions, and \texttt{Rank} is defined for matrices, free modules, $p$\texttt{\symbol{45}}groups, and transitive permutation actions. 

 It is in principle possible to install methods for the operation \texttt{Rank} that are applicable to the different types of arguments, corresponding to the
different contexts. But this is not the approach taken in the \textsf{GAP} library. Instead there are operations such as \texttt{RankMat} (\ref{RankMat}) for matrices and \texttt{DegreeOfCharacter} (\ref{DegreeOfCharacter}) (in fact these are attributes) which are installed as methods of the ``ambiguous'' operations \texttt{Rank} and \texttt{Degree}. 

 The idea is to distinguish between on the one hand different ways to compute
the same thing (e.g.{\nobreakspace}different methods for \texttt{\texttt{\symbol{92}}=} (\ref{bSlash=}), \texttt{Size} (\ref{Size}), etc.), and on the other hand genuinely different things (such as the degree
of a polynomial and a permutation action). 

 The former is the basic purpose of operations and attributes. The latter is
provided as a user convenience where mathematical usage forces it on us \emph{and} where no conflicts arise. In programming the library, we use the underlying
mathematically precise operations or attributes, such as \texttt{RankMat} (\ref{RankMat}) and \texttt{RankAction} (\ref{RankAction:for a group, an action domain, etc.}). These should be attributes if appropriate, and the only role of the
operation \texttt{Rank} is to decide which attribute the user meant. That way, stored information is
stored with ``full mathematical precision'' and is less likely to be retrieved for a wrong purpose later. 

 One word about possible conflicts. A typical example is the mathematical term ``centre'', which is defined as $\{ x \in M | a * x = x * a \forall a \in M \}$ for a magma $M$, and as $\{ x \in L | l * x = 0 \forall l \in L \}$ for a Lie algebra $L$. Here it is \emph{not} possible to introduce an operation \texttt{Centre} (\ref{Centre}) that delegates to attributes \texttt{CentreOfMagma} and \texttt{CentreOfLieAlgebra}, depending on the type of the argument. This is because any Lie algebra in \textsf{GAP} is also a magma, so both \texttt{CentreOfMagma} and \texttt{CentreOfLieAlgebra} would be defined for a Lie algebra, with different meaning if the
characteristic is two. So we cannot achieve that one operation in \textsf{GAP} corresponds to the mathematical term ``centre''. 

 ``Ambiguous'' operations such as \texttt{Rank} are declared in the library file \texttt{lib/overload.g}. }

 }

           
\chapter{\textcolor{Chapter }{Creating New Objects}}\label{Creating New Objects}
\logpage{[ 79, 0, 0 ]}
\hyperdef{L}{X83548994805AD1C9}{}
{
  This chapter is divided into three parts. 

 In the first part, it is explained how to create objects with given type
(see{\nobreakspace}\ref{Creating Objects}). 

 In the second part, first a few small examples are given, for dealing with the
usual cases of component objects (see{\nobreakspace}\ref{Component Objects}) and positional objects (see{\nobreakspace}\ref{Positional Objects}), and for the implementation of new kinds of lists (see{\nobreakspace}\ref{Implementing New List Objects} and{\nobreakspace}\ref{Arithmetic Issues in the Implementation of New Kinds of Lists}). Finally, the external representation of objects is introduced
(see{\nobreakspace}\ref{External Representation}), as a tool for representation independent access to an object. 

 The third part deals with some rules concerning the organization of the \textsf{GAP} library; namely, some commands for creating global variables are explained
(see{\nobreakspace}\ref{Global Variables in the Library}) that correspond to the ones discussed in the first part of the chapter, and
the idea of distinguishing declaration and implementation part of \textsf{GAP} packages is outlined (see{\nobreakspace}\ref{Declaration and Implementation Part}). 

 See also Chapter{\nobreakspace}\ref{An Example -- Residue Class Rings} for examples how the functions from the first part are used, and why it is
useful to have a declaration part and an implementation part.  
\section{\textcolor{Chapter }{Creating Objects}}\label{Creating Objects}
\logpage{[ 79, 1, 0 ]}
\hyperdef{L}{X82E86CF37B123FD4}{}
{
  

\subsection{\textcolor{Chapter }{Objectify}}
\logpage{[ 79, 1, 1 ]}\nobreak
\hyperdef{L}{X7CB5C12E813F512B}{}
{\noindent\textcolor{FuncColor}{$\triangleright$\enspace\texttt{Objectify({\mdseries\slshape type, data})\index{Objectify@\texttt{Objectify}}
\label{Objectify}
}\hfill{\scriptsize (function)}}\\


 New objects are created by \texttt{Objectify}. \mbox{\texttt{\mdseries\slshape data}} must be a plain list or a plain record, and \mbox{\texttt{\mdseries\slshape type}} is the type that the desired object shall have. \texttt{Objectify} turns \mbox{\texttt{\mdseries\slshape data}} into an object with type \mbox{\texttt{\mdseries\slshape type}}. That is, \mbox{\texttt{\mdseries\slshape data}} is changed, and afterwards it will not be a list or a record unless \mbox{\texttt{\mdseries\slshape type}} is of type list resp. record. 

 If \mbox{\texttt{\mdseries\slshape data}} is a list then \texttt{Objectify} turns it into a positional object, if \mbox{\texttt{\mdseries\slshape data}} is a record then \texttt{Objectify} turns it into a component object (for examples, see{\nobreakspace}\ref{Component Objects} and{\nobreakspace}\ref{Positional Objects}). 

 \texttt{Objectify} does also return the object that it made out of \mbox{\texttt{\mdseries\slshape data}}. 

 For examples where \texttt{Objectify} is used, see{\nobreakspace}\ref{Component Objects}, \ref{Positional Objects}, and the example in Chapter{\nobreakspace}\ref{An Example -- Residue Class Rings}. }

 

\subsection{\textcolor{Chapter }{ObjectifyWithAttributes}}
\logpage{[ 79, 1, 2 ]}\nobreak
\hyperdef{L}{X85377AC07E775066}{}
{\noindent\textcolor{FuncColor}{$\triangleright$\enspace\texttt{ObjectifyWithAttributes({\mdseries\slshape obj, type, attr1, val1, attr2, val2, ...})\index{ObjectifyWithAttributes@\texttt{ObjectifyWithAttributes}}
\label{ObjectifyWithAttributes}
}\hfill{\scriptsize (function)}}\\


 Attribute assignments will change the type of an object. If you create many
objects, code of the form 

 
\begin{Verbatim}[commandchars=!@|,fontsize=\small,frame=single,label=Example]
  o:=Objectify(type,rec());
  SetMyAttribute(o,value);
\end{Verbatim}
 

 will take a lot of time for type changes. You can avoid this by setting the
attributes immediately while the object is created, as follows. \texttt{ObjectifyWithAttributes} takes a plain list or record \mbox{\texttt{\mdseries\slshape obj}} and turns it an object just like \texttt{Objectify} (\ref{Objectify}) and sets attribute \mbox{\texttt{\mdseries\slshape attr1}} to \mbox{\texttt{\mdseries\slshape val1}}, sets attribute \mbox{\texttt{\mdseries\slshape attr2}} to \mbox{\texttt{\mdseries\slshape val2}} and so forth. 

 If the filter list of \mbox{\texttt{\mdseries\slshape type}} includes that these attributes are set (and the properties also include values
of the properties) and if no special setter methods are installed for any of
the involved attributes then they are set simultaneously without type changes.
This can produce a substantial speedup. 

 If the conditions of the last sentence are not fulfilled, an ordinary \texttt{Objectify} (\ref{Objectify}) with subsequent setter calls for the attributes is performed instead. }

 }

  
\section{\textcolor{Chapter }{Component Objects}}\label{Component Objects}
\logpage{[ 79, 2, 0 ]}
\hyperdef{L}{X866E223484649E5A}{}
{
  \index{"!.@\texttt{"!.}} A \emph{component object} is an object in the representation \texttt{IsComponentObjectRep} (\ref{IsComponentObjectRep}) or a subrepresentation of it. Such an object \mbox{\texttt{\mdseries\slshape cobj}} is built from subobjects that can be accessed via \texttt{\mbox{\texttt{\mdseries\slshape cobj}}!.\mbox{\texttt{\mdseries\slshape name}}}, similar to components of a record. Also analogously to records, values can
be assigned to components of \mbox{\texttt{\mdseries\slshape cobj}} via \texttt{\mbox{\texttt{\mdseries\slshape cobj}}!.\mbox{\texttt{\mdseries\slshape name}}:= \mbox{\texttt{\mdseries\slshape val}}}. For the creation of component objects, see{\nobreakspace}\ref{Creating Objects}. One must be \emph{very careful} when using the \texttt{!.} operator, in order to interpret the component in the right way, and even more
careful when using the assignment to components using \texttt{!.}, in order to keep the information stored in \mbox{\texttt{\mdseries\slshape cobj}} consistent. 

 First of all, in the access or assignment to a component as shown above, \mbox{\texttt{\mdseries\slshape name}} must be among the admissible component names for the representation of \mbox{\texttt{\mdseries\slshape cobj}}, see{\nobreakspace}\texttt{NewRepresentation} (\ref{NewRepresentation}). Second, preferably only few low level functions should use \texttt{!.}, whereas this operator should not occur in ``user interactions''. 

 Note that even if \mbox{\texttt{\mdseries\slshape cobj}} claims that it is immutable, i.e., if \mbox{\texttt{\mdseries\slshape cobj}} is not in the category \texttt{IsMutable} (\ref{IsMutable}), access and assignment via \texttt{!.} and \texttt{!.:=} work. This is necessary for being able to store newly discovered information
in immutable objects. 

 The following example shows the implementation of an iterator
(see{\nobreakspace}\ref{Iterators}) for the domain of integers, which is represented as component object.
See{\nobreakspace}\ref{Positional Objects} for an implementation using positional objects. (In practice, such an iterator
can be implemented more elegantly using \texttt{IteratorByFunctions} (\ref{IteratorByFunctions}), see{\nobreakspace}\ref{Example -- Constructing Iterators}.) 

 The used succession of integers is $0, 1, -1, 2, -2, 3, -3, \ldots$, that is, $a_n = n/2$ if $n$ is even, and $a_n = (1-n)/2$ otherwise. 

 
\begin{Verbatim}[commandchars=!@|,fontsize=\small,frame=single,label=Example]
  DeclareRepresentation( "IsIntegersIteratorCompRep",
      IsComponentObjectRep, [ "counter" ] );
\end{Verbatim}
 

 The above command creates a new representation (see{\nobreakspace}\texttt{NewRepresentation} (\ref{NewRepresentation})) \texttt{IsIntegersIteratorCompRep}, as a subrepresentation of \texttt{IsComponentObjectRep} (\ref{IsComponentObjectRep}), and with one admissible component \texttt{counter}. So no other components than \texttt{counter} will be needed. 

 
\begin{Verbatim}[commandchars=!@|,fontsize=\small,frame=single,label=Example]
  InstallMethod( Iterator,
      "method for `Integers'",
      [ IsIntegers ],
      function( Integers )
      return Objectify( NewType( IteratorsFamily,
                                     IsIterator
                                 and IsIntegersIteratorCompRep ),
                        rec( counter := 0 ) );
      end );
\end{Verbatim}
 

 After the above method installation, one can already ask for \texttt{Iterator( Integers )}. Note that exactly the domain of integers is described by the filter \texttt{IsIntegers} (\ref{IsIntegers}). 

 By the call to \texttt{NewType} (\ref{NewType}), the returned object lies in the family containing all iterators, which is \texttt{IteratorsFamily}, it lies in the category \texttt{IsIterator} (\ref{IsIterator}) and in the representation \texttt{IsIntegersIteratorCompRep}; furthermore, it has the component \texttt{counter} with value \texttt{0}. 

 What is missing now are methods for the two basic operations of iterators,
namely \texttt{IsDoneIterator} (\ref{IsDoneIterator}) and \texttt{NextIterator} (\ref{NextIterator}). The former must always return \texttt{false}, since there are infinitely many integers. The latter must return the next
integer in the iteration, and update the information stored in the iterator,
that is, increase the value of the component \texttt{counter}. 

 
\begin{Verbatim}[commandchars=@|A,fontsize=\small,frame=single,label=Example]
  InstallMethod( IsDoneIterator,
      "method for iterator of `Integers'",
      [ IsIterator and IsIntegersIteratorCompRep ],
      ReturnFalse );
  
  InstallMethod( NextIterator,
      "method for iterator of `Integers'",
      [ IsIntegersIteratorCompRep ],
      function( iter )
      iter!.counter:= iter!.counter + 1;
      if iter!.counter mod 2 = 0 then
        return iter!.counter / 2;
      else
        return ( 1 - iter!.counter ) / 2;
      fi;
      end );
\end{Verbatim}
 

\subsection{\textcolor{Chapter }{NamesOfComponents}}
\logpage{[ 79, 2, 1 ]}\nobreak
\hyperdef{L}{X823965BF7DFDACC9}{}
{\noindent\textcolor{FuncColor}{$\triangleright$\enspace\texttt{NamesOfComponents({\mdseries\slshape comobj})\index{NamesOfComponents@\texttt{NamesOfComponents}}
\label{NamesOfComponents}
}\hfill{\scriptsize (function)}}\\


 For a component object \mbox{\texttt{\mdseries\slshape comobj}}, \texttt{NamesOfComponents} returns a list of strings, which are the names of components currently bound
in \mbox{\texttt{\mdseries\slshape comobj}}. 

 For a record \mbox{\texttt{\mdseries\slshape comobj}} in internal representation, \texttt{NamesOfComponents} returns the result of \texttt{RecNames} (\ref{RecNames}). }

 }

  
\section{\textcolor{Chapter }{Positional Objects}}\label{Positional Objects}
\logpage{[ 79, 3, 0 ]}
\hyperdef{L}{X834893D07FAA6FD2}{}
{
  \index{"![]@\texttt{"![]}} A \emph{positional object} is an object in the representation \texttt{IsPositionalObjectRep} (\ref{IsPositionalObjectRep}) or a subrepresentation of it. Such an object \mbox{\texttt{\mdseries\slshape pobj}} is built from subobjects that can be accessed via \texttt{\mbox{\texttt{\mdseries\slshape pobj}}![\mbox{\texttt{\mdseries\slshape pos}}]}, similar to positions in a list. Also analogously to lists, values can be
assigned to positions of \mbox{\texttt{\mdseries\slshape pobj}} via \texttt{\mbox{\texttt{\mdseries\slshape pobj}}![\mbox{\texttt{\mdseries\slshape pos}}]:= \mbox{\texttt{\mdseries\slshape val}}}. For the creation of positional objects, see{\nobreakspace}\ref{Creating Objects}. 

 One must be \emph{very careful} when using the \texttt{![]} operator, in order to interpret the position in the right way, and even more
careful when using the assignment to positions using \texttt{![]}, in order to keep the information stored in \mbox{\texttt{\mdseries\slshape pobj}} consistent. 

 First of all, in the access or assignment to a position as shown above, \mbox{\texttt{\mdseries\slshape pos}} must be among the admissible positions for the representation of \mbox{\texttt{\mdseries\slshape pobj}}, see{\nobreakspace}\texttt{NewRepresentation} (\ref{NewRepresentation}). Second, preferably only few low level functions should use \texttt{![]}, whereas this operator should not occur in ``user interactions''. 

 Note that even if \mbox{\texttt{\mdseries\slshape pobj}} claims that it is immutable, i.e., if \mbox{\texttt{\mdseries\slshape pobj}} is not in the category \texttt{IsMutable} (\ref{IsMutable}), access and assignment via \texttt{![]} work. This is necessary for being able to store newly discovered information
in immutable objects. 

 The following example shows the implementation of an iterator
(see{\nobreakspace}\ref{Iterators}) for the domain of integers, which is represented as positional object.
See{\nobreakspace}\ref{Component Objects} for an implementation using component objects, and more details. 

 
\begin{Verbatim}[commandchars=!@|,fontsize=\small,frame=single,label=Example]
  DeclareRepresentation( "IsIntegersIteratorPosRep",
      IsPositionalObjectRep, [ 1 ] );
\end{Verbatim}
 

 The above command creates a new representation (see{\nobreakspace}\texttt{NewRepresentation} (\ref{NewRepresentation})) \texttt{IsIntegersIteratorPosRep}, as a subrepresentation of \texttt{IsPositionalObjectRep} (\ref{IsPositionalObjectRep}), and with only the first position being admissible for storing data. 

 
\begin{Verbatim}[commandchars=!@|,fontsize=\small,frame=single,label=Example]
  InstallMethod( Iterator,
      "method for `Integers'",
      [ IsIntegers ],
      function( Integers )
      return Objectify( NewType( IteratorsFamily,
                                     IsIterator
                                 and IsIntegersIteratorPosRep ),
                        [ 0 ] );
      end );
\end{Verbatim}
 

 After the above method installation, one can already ask for \texttt{Iterator( Integers )}. Note that exactly the domain of integers is described by the filter \texttt{IsIntegers} (\ref{IsIntegers}). 

 By the call to \texttt{NewType} (\ref{NewType}), the returned object lies in the family containing all iterators, which is \texttt{IteratorsFamily}, it lies in the category \texttt{IsIterator} (\ref{IsIterator}) and in the representation \texttt{IsIntegersIteratorPosRep}; furthermore, the first position has value \texttt{0}. 

 What is missing now are methods for the two basic operations of iterators,
namely \texttt{IsDoneIterator} (\ref{IsDoneIterator}) and \texttt{NextIterator} (\ref{NextIterator}). The former must always return \texttt{false}, since there are infinitely many integers. The latter must return the next
integer in the iteration, and update the information stored in the iterator,
that is, increase the value stored in the first position. 

 
\begin{Verbatim}[commandchars=@|A,fontsize=\small,frame=single,label=Example]
  InstallMethod( IsDoneIterator,
      "method for iterator of `Integers'",
      [ IsIterator and IsIntegersIteratorPosRep ],
      ReturnFalse );
  
  InstallMethod( NextIterator,
      "method for iterator of `Integers'",
      [ IsIntegersIteratorPosRep ],
      function( iter )
      iter![1]:= iter![1] + 1;
      if iter![1] mod 2 = 0 then
        return iter![1] / 2;
      else
        return ( 1 - iter![1] ) / 2;
      fi;
      end );
\end{Verbatim}
 

 It should be noted that one can of course install both the methods shown in
Section{\nobreakspace}\ref{Component Objects} and \ref{Positional Objects}. The call \texttt{Iterator( Integers )} will cause one of the methods to be selected, and for the returned iterator,
which will have one of the representations we constructed, the right \texttt{NextIterator} (\ref{NextIterator}) method will be chosen. }

  
\section{\textcolor{Chapter }{Implementing New List Objects}}\label{Implementing New List Objects}
\logpage{[ 79, 4, 0 ]}
\hyperdef{L}{X82309B3F81DD2237}{}
{
  This section gives some hints for the quite usual situation that one wants to
implement new objects that are lists. More precisely, one either wants to deal
with lists that have additional features, or one wants that some objects also
behave as lists. An example can be found in{\nobreakspace}\ref{Example -- Constructing Enumerators}. 

 A \emph{list} in \textsf{GAP} is an object in the category \texttt{IsList} (\ref{IsList}). Basic operations for lists are \texttt{Length} (\ref{Length}), \texttt{\texttt{\symbol{92}}[\texttt{\symbol{92}}]} (\ref{bSlash[bSlash]}), and \texttt{IsBound\texttt{\symbol{92}}[\texttt{\symbol{92}}]} (\ref{IsBoundbSlash[bSlash]}) (see{\nobreakspace}\ref{Basic Operations for Lists}). 

 Note that the access to the position \mbox{\texttt{\mdseries\slshape pos}} in the list \mbox{\texttt{\mdseries\slshape list}} via \texttt{\mbox{\texttt{\mdseries\slshape list}}[\mbox{\texttt{\mdseries\slshape pos}}]} is handled by the call \texttt{\texttt{\symbol{92}}[\texttt{\symbol{92}}]( \mbox{\texttt{\mdseries\slshape list}}, \mbox{\texttt{\mdseries\slshape pos}} )} to the operation \texttt{\texttt{\symbol{92}}[\texttt{\symbol{92}}]} (\ref{bSlash[bSlash]}). To explain the somewhat strange name \texttt{\texttt{\symbol{92}}[\texttt{\symbol{92}}]} of this operation, note that non\texttt{\symbol{45}}alphanumeric characters
like \texttt{[} and \texttt{]} may occur in \textsf{GAP} variable names only if they are escaped by a \texttt{\texttt{\symbol{92}}} character. 

 Analogously, the check \texttt{IsBound( \mbox{\texttt{\mdseries\slshape list}}[\mbox{\texttt{\mdseries\slshape pos}}] )} whether the position \mbox{\texttt{\mdseries\slshape pos}} of the list \mbox{\texttt{\mdseries\slshape list}} is bound is handled by the call \texttt{IsBound\texttt{\symbol{92}}[\texttt{\symbol{92}}]( \mbox{\texttt{\mdseries\slshape list}}, \mbox{\texttt{\mdseries\slshape pos}} )} to the operation \texttt{IsBound\texttt{\symbol{92}}[\texttt{\symbol{92}}]} (\ref{IsBoundbSlash[bSlash]}). 

 For mutable lists, also assignment to positions and unbinding of positions via
the operations \texttt{\texttt{\symbol{92}}[\texttt{\symbol{92}}]\texttt{\symbol{92}}:\texttt{\symbol{92}}=} (\ref{bSlash[bSlash]bSlash:bSlash=}) and \texttt{Unbind\texttt{\symbol{92}}[\texttt{\symbol{92}}]} (\ref{UnbindbSlash[bSlash]}) are basic operations. The assignment \texttt{\mbox{\texttt{\mdseries\slshape list}}[\mbox{\texttt{\mdseries\slshape pos}}]:= \mbox{\texttt{\mdseries\slshape val}}} is handled by the call \texttt{\texttt{\symbol{92}}[\texttt{\symbol{92}}]\texttt{\symbol{92}}:\texttt{\symbol{92}}=( \mbox{\texttt{\mdseries\slshape list}}, \mbox{\texttt{\mdseries\slshape pos}}, \mbox{\texttt{\mdseries\slshape val}} )}, and \texttt{Unbind( \mbox{\texttt{\mdseries\slshape list}}[\mbox{\texttt{\mdseries\slshape pos}}] )} is handled by the call \texttt{Unbind\texttt{\symbol{92}}[\texttt{\symbol{92}}]( \mbox{\texttt{\mdseries\slshape list}}, \mbox{\texttt{\mdseries\slshape pos}} )}. 

 All other operations for lists, e.g., \texttt{Add} (\ref{Add}), \texttt{Append} (\ref{Append}), \texttt{Sum} (\ref{Sum}), are based on these operations. This means that it is sufficient to install
methods for the new list objects only for the basic operations. 

 So if one wants to implement new list objects then one creates them as objects
in the category \texttt{IsList} (\ref{IsList}), and installs methods for \texttt{Length} (\ref{Length}), \texttt{\texttt{\symbol{92}}[\texttt{\symbol{92}}]} (\ref{bSlash[bSlash]}), and \texttt{IsBound\texttt{\symbol{92}}[\texttt{\symbol{92}}]} (\ref{IsBoundbSlash[bSlash]}). If the new lists shall be mutable, one needs to install also methods for \texttt{\texttt{\symbol{92}}[\texttt{\symbol{92}}]\texttt{\symbol{92}}:\texttt{\symbol{92}}=} (\ref{bSlash[bSlash]bSlash:bSlash=}) and \texttt{Unbind\texttt{\symbol{92}}[\texttt{\symbol{92}}]} (\ref{UnbindbSlash[bSlash]}). 

 One application for this is the implementation of \emph{enumerators} for domains. An enumerator for the domain $D$ is a dense list whose entries are in bijection with the elements of $D$. If $D$ is large then it is not useful to write down all elements. Instead one can
implement such a bijection implicitly. This works also for infinite domains. 

 In this situation, one implements a new representation of the lists that are
already available in \textsf{GAP}, in particular the family of such a list is the same as the family of the
domain $D$. 

 But it is also possible to implement new kinds of lists that lie in new
families, and thus are not equal to lists that were available in \textsf{GAP} before. An example for this is the implementation of matrices whose
multiplication via ``\texttt{*}'' is the Lie product of matrices. 

 In this situation, it makes no sense to put the new matrices into the same
family as the original matrices. Note that the product of two Lie matrices
shall be defined but not the product of an ordinary matrix and a Lie matrix.
So it is possible to have two lists that have the same entries but that are
not equal w.r.t. ``\texttt{=}'' because they lie in different families. }

  
\section{\textcolor{Chapter }{Example {\textendash} Constructing Enumerators}}\label{Example -- Constructing Enumerators}
\logpage{[ 79, 5, 0 ]}
\hyperdef{L}{X849D8BC278649EA5}{}
{
  When dealing with countable sets, a usual task is to define enumerations,
i.e., bijections to the positive integers. In \textsf{GAP}, this can be implemented via \emph{enumerators} (see{\nobreakspace}\ref{Enumerators}). These are lists containing the elements in a specified ordering, and the
operations \texttt{Position} (\ref{Position}) and list access via \texttt{\texttt{\symbol{92}}[\texttt{\symbol{92}}]} (\ref{bSlash[bSlash]}) define the desired bijection. For implementing such an enumerator, one mainly
needs to install the appropriate functions for these operations. 

 A general setup for creating such lists is given by \texttt{EnumeratorByFunctions} (\ref{EnumeratorByFunctions:for a domain and a record}). 

 If the set in question is a domain \mbox{\texttt{\mdseries\slshape D}} for which a \texttt{Size} (\ref{Size}) method is available then all one has to do is to write down the functions for
computing the $n$\texttt{\symbol{45}}th element of the list and for computing the position of a
given \textsf{GAP} object in the list, to put them into the components \texttt{ElementNumber} and \texttt{NumberElement} of a record, and to call \texttt{EnumeratorByFunctions} (\ref{EnumeratorByFunctions:for a domain and a record}) with the domain \mbox{\texttt{\mdseries\slshape D}} and this record as arguments. For example, the following lines of code install
an \texttt{Enumerator} (\ref{Enumerator}) method for the case that \mbox{\texttt{\mdseries\slshape D}} is the domain of rational integers. (Note that \texttt{IsIntegers} (\ref{IsIntegers}) is a filter that describes exactly the domain of rational integers.) 

 
\begin{Verbatim}[commandchars=!@|,fontsize=\small,frame=single,label=Example]
  InstallMethod( Enumerator,
      "for integers",
      [ IsIntegers ],
      Integers -> EnumeratorByFunctions( Integers, rec(
                      ElementNumber := function( e, n ) ... end,
                      NumberElement := function( e, x ) ... end ) ) );
\end{Verbatim}
 

 The bodies of the functions have been omitted above; here is the code that is
actually used in \textsf{GAP}. (The ordering coincides with that for the iterators for the domain of
rational integers that have been discussed in{\nobreakspace}\ref{Component Objects} and{\nobreakspace}\ref{Positional Objects}.) 

 
\begin{Verbatim}[commandchars=@|A,fontsize=\small,frame=single,label=Example]
  @gapprompt|gap>A @gapinput|enum:= Enumerator( Integers );A
  <enumerator of Integers>
  @gapprompt|gap>A @gapinput|Print( enum!.NumberElement, "\n" );A
  function ( e, x )
      local  pos;
      if not IsInt( x )  then
          return fail;
      elif 0 < x  then
          pos := 2 * x;
      else
          pos := -2 * x + 1;
      fi;
      return pos;
  end
  @gapprompt|gap>A @gapinput|Print( enum!.ElementNumber, "\n" );A
  function ( e, n )
      if n mod 2 = 0  then
          return n / 2;
      else
          return (1 - n) / 2;
      fi;
      return;
  end
\end{Verbatim}
 

 The situation becomes slightly more complicated if the set $S$ in question is not a domain. This is because one must provide also at least a
method for computing the length of the list, and because one has to determine
the family in which it lies (see{\nobreakspace}\ref{Creating Objects}). The latter should usually not be a problem since either $S$ is nonempty and all its elements lie in the same family {\textendash}in this
case one takes the collections family of any element in $S${\textendash} or the family of the enumerator must be \texttt{ListsFamily}. 

 An example in the \textsf{GAP} library is an enumerator for the set of $k$\texttt{\symbol{45}}tuples over a finite set; the function is called \texttt{EnumeratorOfTuples} (\ref{EnumeratorOfTuples}).   

 
\begin{Verbatim}[commandchars=@|A,fontsize=\small,frame=single,label=Example]
  @gapprompt|gap>A @gapinput|Print( EnumeratorOfTuples, "\n" );A
  function ( set, k )
      local  enum;
      if k = 0  then
          return Immutable( [ [  ] ] );
      elif IsEmpty( set )  then
          return Immutable( [  ] );
      fi;
      enum
       := EnumeratorByFunctions( CollectionsFamily( FamilyObj( set ) ),
         rec(
            ElementNumber := function ( enum, n )
                  local  nn, t, i;
                  nn := n - 1;
                  t := [  ];
                  for i  in [ 1 .. enum!.k ]  do
                      t[i] := RemInt( nn, Length( enum!.set ) ) + 1;
                      nn := QuoInt( nn, Length( enum!.set ) );
                  od;
                  if nn <> 0  then
                      Error( "<enum>[", n,
                       "] must have an assigned value" );
                  fi;
                  nn := enum!.set{Reversed( t )};
                  MakeImmutable( nn );
                  return nn;
              end,
            NumberElement := function ( enum, elm )
                  local  n, i;
                  if not IsList( elm )  then
                      return fail;
                  fi;
                  elm := List( elm, function ( x )
                          return Position( enum!.set, x );
                      end );
                  if fail in elm or Length( elm ) <> enum!.k  then
                      return fail;
                  fi;
                  n := 0;
                  for i  in [ 1 .. enum!.k ]  do
                      n := Length( enum!.set ) * n + elm[i] - 1;
                  od;
                  return n + 1;
              end,
            Length := function ( enum )
                  return Length( enum!.set ) ^ enum!.k;
              end,
            PrintObj := function ( enum )
                  Print( "EnumeratorOfTuples( ", enum!.set, ", ",
                   enum!.k, " )" );
                  return;
              end,
            set := Set( set ),
            k := k ) );
      SetIsSSortedList( enum, true );
      return enum;
  end
\end{Verbatim}
 

 We see that the enumerator is a homogeneous list that stores individual
functions \texttt{ElementNumber}, \texttt{NumberElement}, \texttt{Length}, and \texttt{PrintObj}; besides that, the data components $S$ and $k$ are contained. }

  
\section{\textcolor{Chapter }{Example {\textendash} Constructing Iterators}}\label{Example -- Constructing Iterators}
\logpage{[ 79, 6, 0 ]}
\hyperdef{L}{X7F6BF6CE7AD04EFC}{}
{
  Iterators are a kind of objects that is implemented for several collections in
the \textsf{GAP} library and which might be interesting also in other cases, see{\nobreakspace}\ref{Iterators}. A general setup for implementing new iterators is provided by \texttt{IteratorByFunctions} (\ref{IteratorByFunctions}). 

 All one has to do is to write down the functions for \texttt{NextIterator} (\ref{NextIterator}), \texttt{IsDoneIterator} (\ref{IsDoneIterator}), and \texttt{ShallowCopy} (\ref{ShallowCopy}), and to call \texttt{IteratorByFunctions} (\ref{IteratorByFunctions}) with this record as argument. For example, the following lines of code install
an \texttt{Iterator} (\ref{Iterator}) method for the case that the argument is the domain of rational integers. 

 (Note that \texttt{IsIntegers} (\ref{IsIntegers}) is a filter that describes exactly the domain of rational integers.) 

 
\begin{Verbatim}[commandchars=!@|,fontsize=\small,frame=single,label=Example]
  InstallMethod( Iterator,
      "for integers",
      [ IsIntegers ],
      Integers -> IteratorByFunctions( rec(
                      NextIterator:= function( iter ) ... end,
                      IsDoneIterator := ReturnFalse,
                      ShallowCopy := function( iter ) ... end ) ) );
\end{Verbatim}
 

 The bodies of two of the functions have been omitted above; here is the code
that is actually used in \textsf{GAP}. (The ordering coincides with that for the iterators for the domain of
rational integers that have been discussed in{\nobreakspace}\ref{Component Objects} and{\nobreakspace}\ref{Positional Objects}.) 

 
\begin{Verbatim}[commandchars=@|A,fontsize=\small,frame=single,label=Example]
  @gapprompt|gap>A @gapinput|iter:= Iterator( Integers );A
  <iterator of Integers at 0>
  @gapprompt|gap>A @gapinput|Print( iter!.NextIterator, "\n" );A
  function ( iter )
      iter!.counter := iter!.counter + 1;
      if iter!.counter mod 2 = 0  then
          return iter!.counter / 2;
      else
          return (1 - iter!.counter) / 2;
      fi;
      return;
  end
  @gapprompt|gap>A @gapinput|Print( iter!.ShallowCopy, "\n" );A
  function ( iter )
      return rec(
          counter := iter!.counter );
  end
\end{Verbatim}
 

 Note that the \texttt{ShallowCopy} component of the record must be a function that does not return an iterator
but a record that can be used as the argument of \texttt{IteratorByFunctions} (\ref{IteratorByFunctions}) in order to create the desired shallow copy. }

  
\section{\textcolor{Chapter }{Arithmetic Issues in the Implementation of New Kinds of Lists}}\label{Arithmetic Issues in the Implementation of New Kinds of Lists}
\logpage{[ 79, 7, 0 ]}
\hyperdef{L}{X829629E87E30090C}{}
{
  When designing a new kind of list objects in \textsf{GAP}, defining the arithmetic behaviour of these objects is an issue. 

 There are situations where arithmetic operations of list objects are
unimportant in the sense that adding two such lists need not be represented in
a special way. In such cases it might be useful either to support no
arithmetics at all for the new lists, or to enable the default arithmetic
methods. The former can be achieved by not setting the filters \texttt{IsGeneralizedRowVector} (\ref{IsGeneralizedRowVector}) and \texttt{IsMultiplicativeGeneralizedRowVector} (\ref{IsMultiplicativeGeneralizedRowVector}) in the types of the lists, the latter can be achieved by setting the filter \texttt{IsListDefault} (\ref{IsListDefault}). (for details, see{\nobreakspace}\ref{Filters Controlling the Arithmetic Behaviour of Lists}). An example for ``wrapped lists'' with default behaviour are vector space bases; they are lists with additional
properties concerning the computation of coefficients, but arithmetic
properties are not important. So it is no loss to enable the default methods
for these lists. 

 However, often the arithmetic behaviour of new list objects is important, and
one wants to keep these lists away from default methods for addition,
multiplication etc. For example, the sum and the product of (compatible) block
matrices shall be represented as a block matrix, so the default methods for
sum and product of matrices shall not be applicable, although the results will
be equal to those of the default methods in the sense that their entries at
corresponding positions are equal. 

 So one does not set the filter \texttt{IsListDefault} (\ref{IsListDefault}) in such cases, and thus one can implement one's own methods for arithmetic
operations.  (Of course ``can'' means on the other hand that one \emph{must} implement such methods if one is interested in arithmetics of the new lists.) 

 The specific binary arithmetic methods for the new lists will usually cover
the case that both arguments are of the new kind, and perhaps also the
interaction between a list of the new kind and certain other kinds of lists
may be handled if this appears to be useful. 

 For the last situation, interaction between a new kind of lists and other
kinds of lists, \textsf{GAP} provides already a setup. Namely, there are the categories \texttt{IsGeneralizedRowVector} (\ref{IsGeneralizedRowVector}) and \texttt{IsMultiplicativeGeneralizedRowVector} (\ref{IsMultiplicativeGeneralizedRowVector}), which are concerned with the additive and the multiplicative behaviour,
respectively, of lists. For lists in these filters, the structure of the
results of arithmetic operations is prescribed (see{\nobreakspace}\ref{Additive Arithmetic for Lists} and \ref{Multiplicative Arithmetic for Lists}). 

 For example, if one implements block matrices in \texttt{IsMultiplicativeGeneralizedRowVector} (\ref{IsMultiplicativeGeneralizedRowVector}) then automatically the product of such a block matrix and a (plain) list of
such block matrices will be defined as the obvious list of matrix products,
and a default method for plain lists will handle this multiplication. (Note
that this method will rely on a method for computing the product of the block
matrices, and of course no default method is available for that.) Conversely,
if the block matrices are not in \texttt{IsMultiplicativeGeneralizedRowVector} (\ref{IsMultiplicativeGeneralizedRowVector}) then the product of a block matrix and a (plain) list of block matrices is not
defined. (There is no default method for it, and one can define the result and
provide a method for computing it.) 

 Thus if one decides to set the filters \texttt{IsGeneralizedRowVector} (\ref{IsGeneralizedRowVector}) and \texttt{IsMultiplicativeGeneralizedRowVector} (\ref{IsMultiplicativeGeneralizedRowVector}) for the new lists, on the one hand one loses freedom in defining arithmetic
behaviour, but on the other hand one gains several default methods for a more
or less natural behaviour. 

 If a list in the filter \texttt{IsGeneralizedRowVector} (\ref{IsGeneralizedRowVector}) (\texttt{IsMultiplicativeGeneralizedRowVector} (\ref{IsMultiplicativeGeneralizedRowVector})) lies in \texttt{IsAttributeStoringRep} (\ref{IsAttributeStoringRep}), the values of additive (multiplicative) nesting depth is stored in the list
and need not be calculated for each arithmetic operation. One can then store
the value(s) already upon creation of the lists, with the effect that the
default arithmetic operations will access elements of these lists only if this
is unavoidable. For example, the sum of two plain lists of ``wrapped matrices'' with stored nesting depths are computed via the method for adding two such
wrapped lists, and without accessing any of their rows (which might be
expensive). In this sense, the wrapped lists are treated as black boxes. }

  
\section{\textcolor{Chapter }{External Representation}}\label{External Representation}
\logpage{[ 79, 8, 0 ]}
\hyperdef{L}{X7EBB961E7FE1B0EB}{}
{
  An operation is defined for elements rather than for objects in the sense that
if the arguments are replaced by objects that are equal to the old arguments
w.r.t. the equivalence relation ``\texttt{=}'' then the result must be equal to the old result w.r.t.{\nobreakspace}``\texttt{=}''. 

 But the implementation of many methods is representation dependent in the
sense that certain representation dependent subobjects are accessed. 

 For example, a method that implements the addition of univariate polynomials
may access coefficients lists of its arguments only if they are really stored,
while in the case of sparsely represented polynomials a different approach is
needed. 

 In spite of this, for many operations one does not want to write an own method
for each possible representations of each argument, for example because none
of the methods could in fact take advantage of the actually given
representations of the objects. Another reason could be that one wants to
install first a representation independent method, and then add specific
methods as they are needed to gain more efficiency, by really exploiting the
fact that the arguments have certain representations. 

 For the purpose of admitting representation independent code, one can define
an \emph{external representation} of objects in a given family, install methods to compute this external
representation for each representation of the objects, and then use this
external representation of the objects whenever they occur. 

 We cannot provide conversion functions that allow us to first convert any
object in question to one particular ``standard representation'', and then access the data in the way defined for this representation, simply
because it may be impossible to choose such a ``standard representation'' uniformly for all objects in the given family. 

 So the aim of an external representation of an object \mbox{\texttt{\mdseries\slshape obj}} is a different one, namely to describe the data from which \mbox{\texttt{\mdseries\slshape obj}} is composed. In particular, the external representation of \mbox{\texttt{\mdseries\slshape obj}} is \emph{not} one possible (``standard'') representation of \mbox{\texttt{\mdseries\slshape obj}}, in fact the external representation of \mbox{\texttt{\mdseries\slshape obj}} is in general different from \mbox{\texttt{\mdseries\slshape obj}} w.r.t.{\nobreakspace}``\texttt{=}'', first of all because the external representation of \mbox{\texttt{\mdseries\slshape obj}} does in general not lie in the same family as \mbox{\texttt{\mdseries\slshape obj}}. 

 For example the external representation of a rational function is a list of
length two or three, the first entry being the zero coefficient, the second
being a list describing the coefficients and monomials of the numerator, and
the third, if bound, being a list describing the coefficients and monomials of
the denominator. In particular, the external representation of a polynomial is
a list and not a polynomial. 

 The other way round, the external representation of \mbox{\texttt{\mdseries\slshape obj}} encodes \mbox{\texttt{\mdseries\slshape obj}} in such a way that from this data and the family of \mbox{\texttt{\mdseries\slshape obj}}, one can create an object that is equal to \mbox{\texttt{\mdseries\slshape obj}}. Usually the external representation of an object is a list or a record. 

 Although the external representation of \mbox{\texttt{\mdseries\slshape obj}} is by definition independent of the actually available representations for \mbox{\texttt{\mdseries\slshape obj}}, it is usual that a representation of \mbox{\texttt{\mdseries\slshape obj}} exists for which the computation of the external representation is obtained by
just ``unpacking'' \mbox{\texttt{\mdseries\slshape obj}}, in the sense that the desired data is stored in a component or a position of \mbox{\texttt{\mdseries\slshape obj}}, if \mbox{\texttt{\mdseries\slshape obj}} is a component object (see{\nobreakspace}\ref{Component Objects}) or a positional object (see{\nobreakspace}\ref{Positional Objects}). 

 To implement an external representation means to install methods for the
following two operations. 

\subsection{\textcolor{Chapter }{ExtRepOfObj}}
\logpage{[ 79, 8, 1 ]}\nobreak
\hyperdef{L}{X8542B32A8206118C}{}
{\noindent\textcolor{FuncColor}{$\triangleright$\enspace\texttt{ExtRepOfObj({\mdseries\slshape obj})\index{ExtRepOfObj@\texttt{ExtRepOfObj}}
\label{ExtRepOfObj}
}\hfill{\scriptsize (operation)}}\\
\noindent\textcolor{FuncColor}{$\triangleright$\enspace\texttt{ObjByExtRep({\mdseries\slshape fam, data})\index{ObjByExtRep@\texttt{ObjByExtRep}}
\label{ObjByExtRep}
}\hfill{\scriptsize (operation)}}\\


 \texttt{ExtRepOfObj} returns the external representation of its argument, and \texttt{ObjByExtRep} returns an object in the family \mbox{\texttt{\mdseries\slshape fam}} that has external representation \mbox{\texttt{\mdseries\slshape data}}. 

 Of course, \texttt{ObjByExtRep( FamilyObj( \mbox{\texttt{\mdseries\slshape obj}} ), ExtRepOfObj( \mbox{\texttt{\mdseries\slshape obj}} ) )} must be equal to \mbox{\texttt{\mdseries\slshape obj}} w.r.t. the operation \texttt{\texttt{\symbol{92}}=} (\ref{bSlash=}). But it is \emph{not} required that equal objects have equal external representations. 

 Note that if one defines a new representation of objects for which an external
representation does already exist then one \emph{must} install a method to compute this external representation for the objects in
the new representation. }

 }

  
\section{\textcolor{Chapter }{Mutability and Copying}}\label{Mutability and Copying}
\logpage{[ 79, 9, 0 ]}
\hyperdef{L}{X8090219A7C76AF55}{}
{
  Any \textsf{GAP} object is either mutable or immutable. This can be tested with the function \texttt{IsMutable} (\ref{IsMutable}). The intended meaning of (im)mutability is a mathematical one: an immutable
object should never change in such a way that it represents a different
Element. Objects \emph{may} change in other ways, for instance to store more information, or represent an
element in a different way. 

 Immutability is enforced in different ways for built\texttt{\symbol{45}}in
objects (like records, or lists) and for external objects (made using \texttt{Objectify} (\ref{Objectify})). 

 For built\texttt{\symbol{45}}in objects which are immutable, the kernel will
prevent you from changing them. Thus 

 
\begin{Verbatim}[commandchars=!@|,fontsize=\small,frame=single,label=Example]
  !gapprompt@gap>| !gapinput@l := [1,2,4];|
  [ 1, 2, 4 ]
  !gapprompt@gap>| !gapinput@MakeImmutable(l);|
  [ 1, 2, 4 ]
  !gapprompt@gap>| !gapinput@l[3] := 5;|
  Error, List Assignment: <list> must be a mutable list
\end{Verbatim}
 

 For external objects, the situation is different. An external object which
claims to be immutable (i.e. its type does not contain \texttt{IsMutable} (\ref{IsMutable})) should not admit any methods which change the element it represents. The
kernel does \emph{not} prevent the use of \texttt{!.} and \texttt{![} to change the underlying data structure. This is used for instance by the code
that stores attribute values for reuse. In general, these \texttt{!} operations should only be used in methods which depend on the representation
of the object. Furthermore, we would \emph{not} recommend users to install methods which depend on the representations of
objects created by the library or by \textsf{GAP} packages, as there is certainly no guarantee of the representations being the
same in future versions of \textsf{GAP}. 

 Here we see an immutable object (the group $S_4$), in which we improperly install a new component. 

 
\begin{Verbatim}[commandchars=@|A,fontsize=\small,frame=single,label=Example]
  @gapprompt|gap>A @gapinput|g := SymmetricGroup(IsPermGroup,4);A
  Sym( [ 1 .. 4 ] )
  @gapprompt|gap>A @gapinput|IsMutable(g);A
  false
  @gapprompt|gap>A @gapinput|NamesOfComponents(g);A
  [ "Size", "NrMovedPoints", "MovedPoints",
    "GeneratorsOfMagmaWithInverses" ]
  @gapprompt|gap>A @gapinput|g!.silly := "rubbish";A
  "rubbish"
  @gapprompt|gap>A @gapinput|NamesOfComponents(g);A
  [ "Size", "NrMovedPoints", "MovedPoints",
    "GeneratorsOfMagmaWithInverses", "silly" ]
  @gapprompt|gap>A @gapinput|g!.silly;A
  "rubbish"
\end{Verbatim}
 

 On the other hand, if we form an immutable externally represented list, we
find that \textsf{GAP} will not let us change the object. 

 
\begin{Verbatim}[commandchars=!@|,fontsize=\small,frame=single,label=Example]
  !gapprompt@gap>| !gapinput@e := Enumerator(g);|
  <enumerator of perm group>
  !gapprompt@gap>| !gapinput@IsMutable(e);|
  false
  !gapprompt@gap>| !gapinput@IsList(e);|
  true
  !gapprompt@gap>| !gapinput@e[3];|
  (1,2,4)
  !gapprompt@gap>| !gapinput@e[3] := false;|
  Error, The list you are trying to assign to is immutable
\end{Verbatim}
 

 When we consider copying objects, another filter \texttt{IsCopyable} (\ref{IsCopyable}), enters the game and we find that \texttt{ShallowCopy} (\ref{ShallowCopy}) and \texttt{StructuralCopy} (\ref{StructuralCopy}) behave quite differently. Objects can be divided for this purpose into three:
mutable objects, immutable but copyable objects, and
non\texttt{\symbol{45}}copyable objects (called constants). 

 A mutable or copyable object should have a method for the operation \texttt{ShallowCopy} (\ref{ShallowCopy}), which should make a new mutable object, sharing its
top\texttt{\symbol{45}}level subobjects with the original. The exact
definition of top\texttt{\symbol{45}}level subobject may be defined by the
implementor for new kinds of object. 

 \texttt{ShallowCopy} (\ref{ShallowCopy}) applied to a constant simply returns the constant. 

 \texttt{StructuralCopy} (\ref{StructuralCopy}) is expected to be much less used than \texttt{ShallowCopy} (\ref{ShallowCopy}). Applied to a mutable object, it returns a new mutable object which shares no
mutable sub\texttt{\symbol{45}}objects with the input. Applied to an immutable
object (even a copyable one), it just returns the object. It is not an
operation (indeed, it's a rather special kernel function). 

 
\begin{Verbatim}[commandchars=!@|,fontsize=\small,frame=single,label=Example]
  !gapprompt@gap>| !gapinput@e1 := StructuralCopy(e);|
  <enumerator of perm group>
  !gapprompt@gap>| !gapinput@IsMutable(e1);|
  false
  !gapprompt@gap>| !gapinput@e2 := ShallowCopy(e);|
  [ (), (1,4), (1,2,4), (1,3,4), (2,4), (1,4,2), (1,2), (1,3,4,2),
    (2,3,4), (1,4,2,3), (1,2,3), (1,3)(2,4), (3,4), (1,4,3), (1,2,4,3),
    (1,3), (2,4,3), (1,4,3,2), (1,2)(3,4), (1,3,2), (2,3), (1,4)(2,3),
    (1,2,3,4), (1,3,2,4) ]
\end{Verbatim}
 

 There are two other related functions: \texttt{Immutable} (\ref{Immutable}), which makes a new immutable object which shares no mutable subobjects with
its input and \texttt{MakeImmutable} (\ref{MakeImmutable}) which changes an object and its mutable subobjects \emph{in place} to be immutable. It should only be used on ``new'' objects that you have just created, and which cannot share mutable subobjects
with anything else. 

 Both \texttt{Immutable} (\ref{Immutable}) and \texttt{MakeImmutable} (\ref{MakeImmutable}) work on external objects by just resetting the \texttt{IsMutable} (\ref{IsMutable}) filter in the object's type. This should make ineligible any methods that
might change the object. As a consequence, you must allow for the possibility
of immutable versions of any objects you create. 

 So, if you are implementing your own external objects. The rules amount to the
following: 

 
\begin{enumerate}
\item  You decide if your objects should be mutable or copyable or constants, by
fixing whether their type includes \texttt{IsMutable} (\ref{IsMutable}) or \texttt{IsCopyable} (\ref{IsCopyable}). 
\item  You install methods for your objects respecting that decision: 
\begin{description}
\item[{for constants:}]  no methods change the underlying elements; 
\item[{for copyables:}]  you provide a method for \texttt{ShallowCopy} (\ref{ShallowCopy}); 
\item[{for mutables:}]  you may have methods that change the underlying elements and these should
explicitly require \texttt{IsMutable} (\ref{IsMutable}). 
\end{description}
 
\end{enumerate}
 }

  
\section{\textcolor{Chapter }{Global Variables in the Library}}\label{Global Variables in the Library}
\logpage{[ 79, 10, 0 ]}
\hyperdef{L}{X87E29BA57C8208A4}{}
{
  Global variables in the \textsf{GAP} library are usually read\texttt{\symbol{45}}only in order to prevent them from
being overwritten accidentally. See also Section \ref{More About Global Variables}. 

\subsection{\textcolor{Chapter }{DeclareGlobalName}}
\logpage{[ 79, 10, 1 ]}\nobreak
\hyperdef{L}{X828E14ED7EE39522}{}
{\noindent\textcolor{FuncColor}{$\triangleright$\enspace\texttt{DeclareGlobalName({\mdseries\slshape name})\index{DeclareGlobalName@\texttt{DeclareGlobalName}}
\label{DeclareGlobalName}
}\hfill{\scriptsize (function)}}\\


 For global variables, sometimes code needs to reference them before a value
can sensibly be assigned to them. For example, consider the following
definition of a recursive function: 
\begin{Verbatim}[commandchars=!@|,fontsize=\small,frame=single,label=Example]
  BindGlobal( "fun", function(n)
    if n > 0 then
      return 2*fun(n-1);
    fi;
    return 1;
  end );
\end{Verbatim}
 The problem with that code is that it triggers a syntax warning about access
to an unbound global variable, as \texttt{fun} only gets assigned after the complete statement has been parsed, yet that
statement references \texttt{fun}. 

 To resolve this, one can declare the variable with \texttt{DeclareGlobalName} before assigning it via \texttt{BindGlobal} (\ref{BindGlobal}). This informs \textsf{GAP} that a global variable with the specified name will eventually be defined, and
that thus no syntax warnings pertaining to its use should be printed. 

 We recommend using \texttt{DeclareGlobalName} instead of \texttt{DeclareGlobalVariable} (\ref{DeclareGlobalVariable}) or \texttt{DeclareGlobalFunction} (\ref{DeclareGlobalFunction}) whenever possible. 

 \texttt{DeclareGlobalName} shall be used in the declaration part of the respective package
(see{\nobreakspace}\ref{Declaration and Implementation Part}), values can then be assigned to the new variable as usual, preferably via \texttt{BindGlobal} (\ref{BindGlobal}), in the implementation part (again, see{\nobreakspace}\ref{Declaration and Implementation Part}). }

 

\subsection{\textcolor{Chapter }{DeclareGlobalVariable}}
\logpage{[ 79, 10, 2 ]}\nobreak
\hyperdef{L}{X8324B5DE8300E0F2}{}
{\noindent\textcolor{FuncColor}{$\triangleright$\enspace\texttt{DeclareGlobalVariable({\mdseries\slshape name[, description]})\index{DeclareGlobalVariable@\texttt{DeclareGlobalVariable}}
\label{DeclareGlobalVariable}
}\hfill{\scriptsize (function)}}\\


 For global variables that are \emph{not} functions, instead of using \texttt{BindGlobal} (\ref{BindGlobal}) one can also declare the variable with \texttt{DeclareGlobalVariable} which creates a new global variable named by the string \mbox{\texttt{\mdseries\slshape name}}. 

 In the past the main application of this was to allow access to variables
before they were assigned. Starting with \textsf{GAP} 4.12 we recommend to instead use \texttt{DeclareGlobalName} (\ref{DeclareGlobalName}) for this kind of problem. The main remaining application for \texttt{DeclareGlobalVariable} is when one needs flushable values. 

 If used at all, then \texttt{DeclareGlobalVariable} shall be used in the declaration part of the respective package
(see{\nobreakspace}\ref{Declaration and Implementation Part}), values can then be assigned to the new variable with \texttt{InstallValue} (\ref{InstallValue}), \texttt{InstallFlushableValue} (\ref{InstallFlushableValue}) or \texttt{InstallFlushableValueFromFunction} (\ref{InstallFlushableValueFromFunction}), in the implementation part (again, see{\nobreakspace}\ref{Declaration and Implementation Part}). }

 

\subsection{\textcolor{Chapter }{InstallValue}}
\logpage{[ 79, 10, 3 ]}\nobreak
\hyperdef{L}{X7A23F09886E936D2}{}
{\noindent\textcolor{FuncColor}{$\triangleright$\enspace\texttt{InstallValue({\mdseries\slshape gvar, value})\index{InstallValue@\texttt{InstallValue}}
\label{InstallValue}
}\hfill{\scriptsize (function)}}\\
\noindent\textcolor{FuncColor}{$\triangleright$\enspace\texttt{InstallFlushableValue({\mdseries\slshape gvar, value})\index{InstallFlushableValue@\texttt{InstallFlushableValue}}
\label{InstallFlushableValue}
}\hfill{\scriptsize (function)}}\\
\noindent\textcolor{FuncColor}{$\triangleright$\enspace\texttt{InstallFlushableValueFromFunction({\mdseries\slshape gvar, func})\index{InstallFlushableValueFromFunction@\texttt{InstallFlushableValueFromFunction}}
\label{InstallFlushableValueFromFunction}
}\hfill{\scriptsize (function)}}\\


 \texttt{InstallValue} assigns the value \mbox{\texttt{\mdseries\slshape value}} to the global variable \mbox{\texttt{\mdseries\slshape gvar}} if it was previously declared via \texttt{DeclareGlobalVariable} (\ref{DeclareGlobalVariable}). \texttt{InstallFlushableValue} does the same but additionally provides that each call of \texttt{FlushCaches} (\ref{FlushCaches}) will assign a structural copy of \mbox{\texttt{\mdseries\slshape value}} to \mbox{\texttt{\mdseries\slshape gvar}}. \texttt{InstallFlushableValueFromFunction} instead assigns the result of \mbox{\texttt{\mdseries\slshape func}} to \mbox{\texttt{\mdseries\slshape gvar}} (\mbox{\texttt{\mdseries\slshape func}} is re\texttt{\symbol{45}}evaluated for each invocation of \texttt{FlushCaches} (\ref{FlushCaches}) 

 \texttt{InstallValue} does \emph{not} work if \mbox{\texttt{\mdseries\slshape value}} is an ``immediate object'', i.e., an internally represented small integer or finite field element. It
also fails for booleans. In addition, the use of any other type of \mbox{\texttt{\mdseries\slshape value}} except for a list or record is deprecated. Furthermore, \texttt{InstallFlushableValue} works only if \mbox{\texttt{\mdseries\slshape value}} is a list or a record. (Note that \texttt{InstallFlushableValue} makes sense only for \emph{mutable} global variables.) }

 

\subsection{\textcolor{Chapter }{FlushCaches}}
\logpage{[ 79, 10, 4 ]}\nobreak
\hyperdef{L}{X87A4316C818B3DE3}{}
{\noindent\textcolor{FuncColor}{$\triangleright$\enspace\texttt{FlushCaches({\mdseries\slshape })\index{FlushCaches@\texttt{FlushCaches}}
\label{FlushCaches}
}\hfill{\scriptsize (operation)}}\\


 \texttt{FlushCaches} resets the value of each global variable that has been declared with \texttt{DeclareGlobalVariable} (\ref{DeclareGlobalVariable}) and for which the initial value has been set with \texttt{InstallFlushableValue} (\ref{InstallFlushableValue}) or \texttt{InstallFlushableValueFromFunction} (\ref{InstallFlushableValueFromFunction}) to this initial value. 

 \texttt{FlushCaches} should be used only for debugging purposes, since the involved global
variables include for example lists that store finite fields and cyclotomic
fields used in the current \textsf{GAP} session, in order to avoid that these fields are constructed anew in each call
to \texttt{GF} (\ref{GF:for field size}) and \texttt{CF} (\ref{CF:for (subfield and) conductor}). }

 

\subsection{\textcolor{Chapter }{DeclareGlobalFunction}}
\logpage{[ 79, 10, 5 ]}\nobreak
\hyperdef{L}{X834A8CC587A609BE}{}
{\noindent\textcolor{FuncColor}{$\triangleright$\enspace\texttt{DeclareGlobalFunction({\mdseries\slshape name})\index{DeclareGlobalFunction@\texttt{DeclareGlobalFunction}}
\label{DeclareGlobalFunction}
}\hfill{\scriptsize (function)}}\\
\noindent\textcolor{FuncColor}{$\triangleright$\enspace\texttt{InstallGlobalFunction({\mdseries\slshape oper, func})\index{InstallGlobalFunction@\texttt{InstallGlobalFunction}}
\label{InstallGlobalFunction}
}\hfill{\scriptsize (function)}}\\


 \textsf{GAP} functions that are not operations and that are intended to be called by users
should be notified to \textsf{GAP} via \texttt{DeclareGlobalFunction}. \texttt{DeclareGlobalFunction} returns a function that serves as a placeholder for the function that will be
installed later. The placeholder will print an error message if it is called.
See also{\nobreakspace}\texttt{DeclareSynonym} (\ref{DeclareSynonym}). 

 In the past the main application of this was to allow access to variables
before they were assigned. Starting with \textsf{GAP} 4.12 we recommend to use \texttt{DeclareGlobalName} (\ref{DeclareGlobalName})/\texttt{BindGlobal} (\ref{BindGlobal}) instead of \texttt{DeclareGlobalVariable} (\ref{DeclareGlobalVariable})/\texttt{InstallGlobalFunction} whenever possible. 

 If used at all, then \texttt{DeclareGlobalVariable} (\ref{DeclareGlobalVariable}) shall be used in the declaration part of the respective package
(see{\nobreakspace}\ref{Declaration and Implementation Part}). 

 A global function declared with \texttt{DeclareGlobalFunction} can be given its value \mbox{\texttt{\mdseries\slshape func}} via \texttt{InstallGlobalFunction}; \mbox{\texttt{\mdseries\slshape gvar}} is the global variable (or a string denoting its name) named with the \mbox{\texttt{\mdseries\slshape name}} argument of the call to \texttt{DeclareGlobalFunction}. For example, a declaration like 

 
\begin{Verbatim}[commandchars=!@|,fontsize=\small,frame=single,label=Example]
  DeclareGlobalFunction( "SumOfTwoCubes" );
\end{Verbatim}
 

 in the ``declaration part'' (see Section{\nobreakspace}\ref{Declaration and Implementation Part}) might have a corresponding ``implementation part'' of: 

 
\begin{Verbatim}[commandchars=!@|,fontsize=\small,frame=single,label=Example]
  InstallGlobalFunction( SumOfTwoCubes, function(x, y) return x^3 + y^3; end);
\end{Verbatim}
 

 }

 

\subsection{\textcolor{Chapter }{DeclareSynonym}}
\logpage{[ 79, 10, 6 ]}\nobreak
\hyperdef{L}{X851654DA87616207}{}
{\noindent\textcolor{FuncColor}{$\triangleright$\enspace\texttt{DeclareSynonym({\mdseries\slshape name, value})\index{DeclareSynonym@\texttt{DeclareSynonym}}
\label{DeclareSynonym}
}\hfill{\scriptsize (function)}}\\
\noindent\textcolor{FuncColor}{$\triangleright$\enspace\texttt{DeclareSynonymAttr({\mdseries\slshape name, value})\index{DeclareSynonymAttr@\texttt{DeclareSynonymAttr}}
\label{DeclareSynonymAttr}
}\hfill{\scriptsize (function)}}\\


 \texttt{DeclareSynonym} assigns the string \mbox{\texttt{\mdseries\slshape name}} to a global variable as a synonym for \mbox{\texttt{\mdseries\slshape value}}. Two typical intended usages are to declare an ``and\texttt{\symbol{45}}filter'', e.g. 

 
\begin{Verbatim}[commandchars=!@|,fontsize=\small,frame=single,label=Example]
  DeclareSynonym( "IsGroup", IsMagmaWithInverses and IsAssociative );
\end{Verbatim}
 

 and to provide a previously declared global function with an alternative name,
e.g. 

 
\begin{Verbatim}[commandchars=!@|,fontsize=\small,frame=single,label=Example]
  DeclareGlobalFunction( "SizeOfSomething" );
  DeclareSynonym( "OrderOfSomething", SizeOfSomething );
\end{Verbatim}
 

 \emph{Note:} Before using \texttt{DeclareSynonym} in the way of this second example, one should determine whether the synonym is
really needed. Perhaps an extra index entry in the documentation would be
sufficient. 

 When \mbox{\texttt{\mdseries\slshape value}} is actually an attribute then \texttt{DeclareSynonymAttr} should be used; this binds also globals variables \texttt{Set}\mbox{\texttt{\mdseries\slshape name}} and \texttt{Has}\mbox{\texttt{\mdseries\slshape name}} for its setter and tester, respectively. 

 
\begin{Verbatim}[commandchars=!@|,fontsize=\small,frame=single,label=Example]
  DeclareSynonymAttr( "IsField", IsDivisionRing and IsCommutative );
  DeclareAttribute( "GeneratorsOfDivisionRing", IsDivisionRing );
  DeclareSynonymAttr( "GeneratorsOfField", GeneratorsOfDivisionRing );
\end{Verbatim}
 }

 }

  
\section{\textcolor{Chapter }{Declaration and Implementation Part}}\label{Declaration and Implementation Part}
\logpage{[ 79, 11, 0 ]}
\hyperdef{L}{X7837CA9A83D93B38}{}
{
  Each package of \textsf{GAP} code consists of two parts, the \emph{declaration part} that defines the new categories and operations for the objects the package
deals with, and the \emph{implementation part} where the corresponding methods are installed. The declaration part should be
representation independent, representation dependent information should be
dealt with in the implementation part. 

 \textsf{GAP} functions that are not operations and that are intended to be called by users
should be notified to \textsf{GAP} in the declaration part via \texttt{DeclareGlobalFunction} (\ref{DeclareGlobalFunction}). Values for these functions can be installed in the implementation part via \texttt{InstallGlobalFunction} (\ref{InstallGlobalFunction}). 

 Calls to the following functions belong to the declaration part. 
\begin{itemize}
\item \texttt{DeclareAttribute} (\ref{DeclareAttribute}),
\item \texttt{DeclareCategory} (\ref{DeclareCategory}),
\item \texttt{DeclareFilter} (\ref{DeclareFilter}),
\item \texttt{DeclareOperation} (\ref{DeclareOperation}),
\item \texttt{DeclareGlobalFunction} (\ref{DeclareGlobalFunction}),
\item \texttt{DeclareGlobalName} (\ref{DeclareGlobalName}),
\item \texttt{DeclareGlobalVariable} (\ref{DeclareGlobalVariable}),
\item \texttt{DeclareSynonym} (\ref{DeclareSynonym}),
\item \texttt{DeclareSynonymAttr} (\ref{DeclareSynonymAttr}),
\item \texttt{DeclareProperty} (\ref{DeclareProperty}),
\item \texttt{InstallTrueMethod} (\ref{InstallTrueMethod}).
\end{itemize}
 Calls to the following functions belong to the implementation part. 
\begin{itemize}
\item \texttt{DeclareRepresentation} (\ref{DeclareRepresentation}),
\item \texttt{InstallGlobalFunction} (\ref{InstallGlobalFunction}),
\item \texttt{InstallValue} (\ref{InstallValue}),
\item \texttt{InstallMethod} (\ref{InstallMethod}),
\item \texttt{InstallImmediateMethod} (\ref{InstallImmediateMethod}),
\item \texttt{InstallOtherMethod} (\ref{InstallOtherMethod}),
\item \texttt{NewFamily} (\ref{NewFamily}),
\item \texttt{NewType} (\ref{NewType}),
\item \texttt{Objectify} (\ref{Objectify}).
\end{itemize}
 \index{DeclareRepresentation@\texttt{DeclareRepresentation}!belongs to implementation part} Whenever both a \texttt{New}\mbox{\texttt{\mdseries\slshape Something}} and a \texttt{Declare}\mbox{\texttt{\mdseries\slshape Something}} variant of a function exist (see{\nobreakspace}\ref{Global Variables in the Library}), the use of \texttt{Declare}\mbox{\texttt{\mdseries\slshape Something}} is recommended because this protects the variables in question from being
overwritten. Note that there are \emph{no} functions \texttt{DeclareFamily} and \texttt{DeclareType} since families and types are created dynamically, hence usually no global
variables are associated to them. Further note that \texttt{DeclareRepresentation} (\ref{DeclareRepresentation}) is regarded as belonging to the implementation part, because usually
representations of objects are accessed only in very few places, and all code
that involves a particular representation is contained in one file;
additionally, representations of objects are often not interesting for the
user, so there is no need to provide a user interface or documentation about
representations. 

 It should be emphasized that ``declaration'' means only an explicit notification of mathematical or technical terms or of
concepts to \textsf{GAP}. For example, declaring a category or property with name \texttt{IsInteresting} does of course not tell \textsf{GAP} what this shall mean, and it is necessary to implement possibilities to create
objects that know already that they lie in \texttt{IsInteresting} in the case that it is a category, or to install implications or methods in
order to compute for a given object whether \texttt{IsInteresting} is \texttt{true} or \texttt{false} for it in the case that \texttt{IsInteresting} is a property. }

 }

           
\chapter{\textcolor{Chapter }{Examples of Extending the System}}\label{Examples of Extending the System}
\logpage{[ 80, 0, 0 ]}
\hyperdef{L}{X8186831682A00097}{}
{
  This chapter gives a few examples of how one can extend the functionality of \textsf{GAP}. 

 They are arranged in ascending difficulty. We show how to install new methods,
add new operations and attributes and how to implement new features using
categories and representations. (As we do not introduce completely new kinds
of objects in these example it will not be necessary to declare any families.)
Finally we show a simple way how to create new objects with an own arithmetic. 

 The examples given are all very rudimentary {\textendash} no particular error
checks are performed and the user interface sometimes is quite clumsy. These
examples may be constructed for presentation purposes only and they do not
necessarily constitute parts of the \textsf{GAP} library. 

 Even more complex examples that create whole classes of objects anew will be
given in the following two chapters \ref{An Example -- Residue Class Rings} and \ref{An Example -- Designing Arithmetic Operations}.  
\section{\textcolor{Chapter }{Addition of a Method}}\label{Addition of a Method}
\logpage{[ 80, 1, 0 ]}
\hyperdef{L}{X7B42DF6E7CCF507D}{}
{
  The easiest case is the addition of a new algorithm as a method for an
existing operation for the existing structures. 

 For example, assume we wanted to implement a better method for computing the
exponent of a nilpotent group (it is the product of the exponents of the Sylow
subgroups). 

 The first task is to find which operation is used by \textsf{GAP} (it is \texttt{Exponent} (\ref{Exponent})) and how it is declared. We can find this in the Reference Manual (in our
particular case in section{\nobreakspace}\ref{Numerical Group Attributes}) and the declaration in the library file \texttt{lib/grp.gd}. The easiest way to find the place of the declaration is usually to \texttt{grep} over all \texttt{.gd} and \texttt{.g} files, see section{\nobreakspace}\ref{Library Files}. 

 In our example the declaration in the library is: 
\begin{Verbatim}[commandchars=!@|,fontsize=\small,frame=single,label=Example]
  DeclareAttribute("Exponent",IsGroup);
\end{Verbatim}
 

 Similarly we find that the filter \texttt{IsNilpotentGroup} (\ref{IsNilpotentGroup}) represents the concept of being nilpotent. 

 We then write a function that implements the new algorithm which takes the
right set of arguments and install it as a method. In our example this
installation would be: 
\begin{Verbatim}[commandchars=!@|,fontsize=\small,frame=single,label=Example]
  InstallMethod(Exponent,"for nilpotent groups",
    [IsGroup and IsNilpotent],
  function(G)
    [function body omitted]
  end);
\end{Verbatim}
 

 We have left out the optional rank argument of \texttt{InstallMethod} (\ref{InstallMethod}), which normally is a wise choice {\textendash}\textsf{GAP} automatically uses an internal ranking based on the filters that is only
offset by the given rank. So our method will certainly be regarded as ``better'' than a method that has been installed for mere groups or for solvable groups
but will be ranked lower than the library method for abelian groups. 

 That's all. Using \ref{ApplicableMethod} we can check for a nilpotent group that indeed our new method will be used. 

 When testing, remember that the method selection will not check for properties
that are not known. (This is done internally by checking the property tester
first.) Therefore the method would not be applicable for the group \texttt{g} in the following definition but only for the {\textendash}mathematically
identical but endowed with more knowledge by \textsf{GAP}{\textendash} group \texttt{h}. (Section{\nobreakspace}\ref{Enforcing Property Tests} shows a way around this.) 

 
\begin{Verbatim}[commandchars=!@|,fontsize=\small,frame=single,label=Example]
  !gapprompt@gap>| !gapinput@g:=Group((1,2),(1,3)(2,4));;|
  !gapprompt@gap>| !gapinput@h:=Group((1,2),(1,3)(2,4));;|
  !gapprompt@gap>| !gapinput@IsNilpotentGroup(h); # enforce test|
  true
  !gapprompt@gap>| !gapinput@HasIsNilpotentGroup(g);|
  false
  !gapprompt@gap>| !gapinput@HasIsNilpotentGroup(h);|
  true
\end{Verbatim}
 

 Let's now look at a slightly more complicated example: We want to implement a
better method for computing normalizers in a nilpotent permutation group.
(Such an algorithm can be found for example in \cite{luksrakocziwright97}.) 

 We already know \texttt{IsNilpotentGroup} (\ref{IsNilpotentGroup}), the filter \texttt{IsPermGroup} (\ref{IsPermGroup}) represents the concept of being a group of permutations. 

 \textsf{GAP} uses \texttt{Normalizer} (\ref{Normalizer:for two groups}) to compute normalizers, however the declaration is a bit more complicated. In
the library we find 

 
\begin{Verbatim}[commandchars=!@|,fontsize=\small,frame=single,label=Example]
  InParentFOA( "Normalizer", IsGroup, IsObject, NewAttribute );
\end{Verbatim}
 

 The full mechanism of \texttt{InParentFOA} (\ref{InParentFOA}) is described in chapter{\nobreakspace}\ref{Function-Operation-Attribute Triples}, however for our purposes it is sufficient to know that for such a function
the actual work is done by an operation \texttt{NormalizerOp}, an underlying operation for \texttt{Normalizer} (\ref{Normalizer:for two groups}) (and all the complications are just there to be able to remember certain
results) and that the declaration of this operation is given by the first
arguments, it would be: 
\begin{Verbatim}[commandchars=!@|,fontsize=\small,frame=single,label=Example]
  DeclareOperation( "NormalizerOp", [IsGroup, IsObject] );
\end{Verbatim}
 

 This time we decide to enter a non\texttt{\symbol{45}}default family predicate
in the call to \texttt{InstallMethod} (\ref{InstallMethod}). We could just leave it out as in the previous call; this would yield the
default value, the function \texttt{ReturnTrue} (\ref{ReturnTrue}) of arbitrary many arguments which always returns \texttt{true}. However, then the method might be called in some cases of inconsistent input
(for example matrix groups in different characteristics) that ought to fall
through the method selection to raise an error. 

 In our situation, we want the second group to be a subgroup of the first, so
necessarily both must have the same family and we can use \texttt{IsIdenticalObj} (\ref{IsIdenticalObj}) as family predicate. 

 Now we can install the method. Again this manual is lazy and does not show you
the actual code: 
\begin{Verbatim}[commandchars=!@|,fontsize=\small,frame=single,label=Example]
  InstallMethod(NormalizerOp,"for nilpotent permutation groups",IsIdenticalObj,
    [IsPermGroup and IsNilpotentGroup,
     IsPermGroup and IsNilpotentGroup],
  function(G,U)
    [ function body omitted ]
  end);
\end{Verbatim}
 }

  
\section{\textcolor{Chapter }{Extending the Range of Definition of an Existing Operation}}\label{Extending the Range of Definition of an Existing Operation}
\logpage{[ 80, 2, 0 ]}
\hyperdef{L}{X837CF3267EF0CFB3}{}
{
  It might be that the operation has been defined so far only for a set of
objects that is too restrictive for our purposes (or we want to install a
method that takes another number of arguments). If this is the case, the call
to \texttt{InstallMethod} (\ref{InstallMethod}) causes an error message. We can avoid this by using \texttt{InstallOtherMethod} (\ref{InstallOtherMethod}) instead. It is also possible to re\texttt{\symbol{45}}declare an operation
with another number of arguments and/or different filters for its arguments. }

  
\section{\textcolor{Chapter }{Enforcing Property Tests}}\label{Enforcing Property Tests}
\logpage{[ 80, 3, 0 ]}
\hyperdef{L}{X7D880DB779EBA8D5}{}
{
  As mentioned in Section \ref{Applicable Methods and Method Selection}, \textsf{GAP} does not check unknown properties to test whether a method might be
applicable. In some cases one wants to enforce this, however, because the gain
from knowing the property outweighs the cost of its determination. 

 In this situation one has to install a method \emph{without} the additional property (so it can be tried even if the property is not yet
known) and at high rank (so it will be used before other methods). The first
thing to do in the actual function then is to test the property and to bail
out with \texttt{TryNextMethod} (\ref{TryNextMethod}) if it turns out to be \texttt{false}. 

 The above \texttt{Exponent} (\ref{Exponent}) example thus would become:  
\begin{Verbatim}[commandchars=!@|,fontsize=\small,frame=single,label=Example]
  InstallMethod(Exponent,"test abelianity", [IsGroup],
    50,# enforced high rank
  function(G)
    if not IsAbelian(G) then
      TryNextMethod();
    fi;
    [remaining function body omitted]
  end);
\end{Verbatim}
 

 The value ``50'' used in this example is quite arbitrary. A better way is to use values that
are given by the system inherently: We want this method still to be ranked as
high, \emph{as if it had} the \texttt{IsAbelian} (\ref{IsAbelian}) requirement. So we have \textsf{GAP} compute the extra rank of this: 
\begin{Verbatim}[commandchars=!@|,fontsize=\small,frame=single,label=Example]
  InstallMethod(Exponent,"test abelianity", [IsGroup],
    # enforced absolute rank of `IsGroup and IsAbelian' installation: Subtract
    # the rank of `IsGroup' and add the rank of `IsGroup and IsAbelian':
    RankFilter(IsGroup and IsAbelian)
    -RankFilter(IsGroup),
  function(G)
\end{Verbatim}
 the slightly complicated construction of addition and subtraction is necessary
because \texttt{IsGroup} (\ref{IsGroup}) and \texttt{IsAbelian} (\ref{IsAbelian}) might imply the \emph{same} elementary filters which we otherwise would count twice. 

 A somehow similar situation occurs with matrix groups. Most methods for matrix
groups are only applicable if the group is known to be finite. 

 However we should not enforce a finiteness test early (someone else later
might install good methods for infinite groups while the finiteness test would
be too expensive) but just before \textsf{GAP} would give a ``no method found'' error. This is done by redispatching, see{\nobreakspace}\ref{Redispatching}. For example to enforce such a final finiteness test for normalizer
calculations could be done by: 
\begin{Verbatim}[commandchars=!@|,fontsize=\small,frame=single,label=Example]
  RedispatchOnCondition(NormalizerOp,IsIdenticalObj,
    [IsMatrixGroup,IsMatrixGroup],[IsFinite,IsFinite],0);
\end{Verbatim}
 }

  
\section{\textcolor{Chapter }{Adding a new Operation}}\label{Adding a new Operation}
\logpage{[ 80, 4, 0 ]}
\hyperdef{L}{X797545848520A44B}{}
{
  Next, we will consider how to add own operations. As an example we take the
Sylow normalizer in a group of a given prime. This operation gets two
arguments, the first has to be a group, the second a prime number. 

 There is a function \texttt{IsPrimeInt} (\ref{IsPrimeInt}), but no property for being prime (which would be pointless as integers cannot
store property values anyhow). So the second argument gets specified only as
positive integer: 
\begin{Verbatim}[commandchars=!@|,fontsize=\small,frame=single,label=Example]
  SylowNormalizer:=NewOperation("SylowNormalizer",[IsGroup,IsPosInt]);
\end{Verbatim}
 (Note that we are using \texttt{NewOperation} (\ref{NewOperation}) instead of \texttt{DeclareOperation} (\ref{DeclareOperation}) as used in the library. The only difference other than that \texttt{DeclareOperation} (\ref{DeclareOperation}) saves some typing, is that it also protects the variables against overwriting.
When testing code (when one probably wants to change things) this might be
restricting. If this does not bother you, you can use 

 
\begin{Verbatim}[commandchars=!@|,fontsize=\small,frame=single,label=Example]
  DeclareOperation("SylowNormalizer",[IsGroup,IsPosInt]);
\end{Verbatim}
 as well.) 

 The filters \texttt{IsGroup} (\ref{IsGroup}) and \texttt{IsPosInt} (\ref{IsPosInt}) given are \emph{only} used to test that \texttt{InstallMethod} (\ref{InstallMethod}) installs methods with suitable arguments and will be completely ignored when
using \texttt{InstallOtherMethod} (\ref{InstallOtherMethod}). Technically one could therefore simply use \texttt{IsObject} (\ref{IsObject}) for all arguments in the declaration. The main point of using more specific
filters here is to help documenting with which arguments the function is to be
used (so for example a call \texttt{SylowNormalizer(5,G)} would be invalid). 

 Of course initially there are no useful methods for newly declared operations;
you will have to write and install them yourself. 

 If the operation only takes one argument and has reproducible results without
side effects, it might be worth declaring it as an attribute instead; see
Section \ref{Adding a new Attribute}. }

  
\section{\textcolor{Chapter }{Adding a new Attribute}}\label{Adding a new Attribute}
\logpage{[ 80, 5, 0 ]}
\hyperdef{L}{X874AF11D864AEC1B}{}
{
  \index{NewAttribute@\texttt{NewAttribute}!example} \index{DeclareAttribute@\texttt{DeclareAttribute}!example} Now we look at an example of how to add a new attribute. As example we
consider the set of all primes that divide the size of a group. 

 First we have to declare the attribute: 

 
\begin{Verbatim}[commandchars=!@|,fontsize=\small,frame=single,label=Example]
  PrimesDividingSize:=NewAttribute("PrimesDividingSize",IsGroup);
\end{Verbatim}
 

 (See{\nobreakspace}\texttt{NewAttribute} (\ref{NewAttribute})). This implicitly declares attribute tester and setter, it is convenient
however to assign these to variables as well: 

 
\begin{Verbatim}[commandchars=!@|,fontsize=\small,frame=single,label=Example]
  HasPrimesDividingSize:=Tester(PrimesDividingSize);
  SetPrimesDividingSize:=Setter(PrimesDividingSize);
\end{Verbatim}
 

 Alternatively, there is a declaration command \texttt{DeclareAttribute} (\ref{DeclareAttribute}) that executes all three assignments simultaneously and protects the variables
against overwriting: 

 
\begin{Verbatim}[commandchars=!@|,fontsize=\small,frame=single,label=Example]
  DeclareAttribute("PrimesDividingSize",IsGroup);
\end{Verbatim}
 

 Next we have to install method(s) for the attribute that compute its value.
(This is not strictly necessary. We could use the attribute also without
methods only for storing and retrieving information, but calling it for
objects for which the value is not known would produce a ``no method found'' error.) For this purpose we can imagine the attribute simply as an
one\texttt{\symbol{45}}argument operation: 

 
\begin{Verbatim}[commandchars=!@|,fontsize=\small,frame=single,label=Example]
  InstallMethod(PrimesDividingSize,"for finite groups",
    [IsGroup and IsFinite],
  function(G)
    return PrimeDivisors(Size(G));
  end);
\end{Verbatim}
 

 The function installed \emph{must} always return a value (or call \texttt{TryNextMethod} (\ref{TryNextMethod})). If the object is in the representation \texttt{IsAttributeStoringRep} (\ref{IsAttributeStoringRep}) this return value once computed will be automatically stored and retrieved if
the attribute is called a second time. We don't have to call setter or tester
ourselves. (This storage happens by \textsf{GAP} internally calling the attribute setter with the return value of the function.
Retrieval is by a high\texttt{\symbol{45}}ranking method which is installed
under the condition \texttt{HasPrimesDividingSize}. This method was installed automatically when the attribute was declared.) }

  
\section{\textcolor{Chapter }{Adding a new Representation}}\label{Adding a new Representation}
\logpage{[ 80, 6, 0 ]}
\hyperdef{L}{X8111D831783C9ED6}{}
{
  \index{NewRepresentation@\texttt{NewRepresentation}!example} \index{DeclareRepresentation@\texttt{DeclareRepresentation}!example} Next, we look at the implementation of a new representation of existing
objects. In most cases we want to implement this representation only for
efficiency reasons while keeping all the existing functionality. 

 For example, assume we wanted (following \cite{wielandt69}) to implement permutation groups defined by relations. 

 Next, we have to decide a few basics about the representation. All existing
permutation groups in the library are attribute storing and we probably want
to keep this for our new objects. Thus the representation must be a
subrepresentation of \texttt{IsComponentObjectRep} (\ref{IsComponentObjectRep}) and \texttt{IsAttributeStoringRep} (\ref{IsAttributeStoringRep}). Furthermore we want each object to be a permutation group and we can imply
this directly in the representation. 

 We also decide that we store the degree (the largest point that might be
moved) in a component \texttt{degree} and the defining relations in a component \texttt{relations} (we do not specify the format of relations here. In an actual implementation
one would have to design this as well, but it does not affect the declarations
this chapter is about). 

 
\begin{Verbatim}[commandchars=!@|,fontsize=\small,frame=single,label=Example]
  DeclareRepresentation(
    "IsPermutationGroupByRelations",
    IsComponentObjectRep and IsAttributeStoringRep and IsPermGroup,
    ["degree","relations"]);
\end{Verbatim}
 

 (If we wanted to implement sparse matrices we might for example rather settle
for a positional object in which we store a list of the nonzero entries.) 

 We can make the new representation a subrepresentation of an existing one. In
such a case of course we have to provide all structure of this ``parent'' representation as well. 

 Next we need to check in which family our new objects will be. This will be
the same family as of every other permutation group, namely the \texttt{CollectionsFamily(PermutationsFamily)} (where the family \texttt{PermutationsFamily = FamilyObj((1,2,3))} has been defined already in the library). 

 Now we can write a function to create our new objects. Usually it is helpful
to look at functions from the library that are used in similar situations (for
example \texttt{GroupByGenerators} (\ref{GroupByGenerators}) in our case) to make sure we have not forgotten any further requirements in
the declaration we might have to add here. However in most cases the function
is straightforward: 

 
\begin{Verbatim}[commandchars=!@|,fontsize=\small,frame=single,label=Example]
  PermutationGroupByRelations:=function(degree,relations)
  local g
    g:=Objectify(NewType(CollectionsFamily(PermutationsFamily),
                 IsPermutationGroupByRelations),
                 rec(degree:=degree,relations:=relations));
    return g;
  end;
\end{Verbatim}
 

 It also is a good idea to install a \texttt{PrintObj} (\ref{PrintObj}) and possibly also a \texttt{ViewObj} (\ref{ViewObj}) method {\textendash}otherwise testing becomes quite hard: 

 
\begin{Verbatim}[commandchars=@|A,fontsize=\small,frame=single,label=Example]
  InstallMethod(PrintObj,"for perm grps. given by relations",
    [IsPermutationGroupByRelations],
  function(G)
    Print("PermutationGroupByRelations(", G!.degree,",",G!.relations,")");
  end);
\end{Verbatim}
 

 Next we have to write enough methods for the new representation so that the
existing algorithms can be used. In particular we will have to implement
methods for all operations for which library or kernel provides methods for
the existing (alternative) representations. In our particular case there are
no such methods. (If we would have implemented sparse matrices we would have
had to implement methods for the list access and assignment functions, see \ref{Basic Operations for Lists}.) However the existing way permutation groups are represented is by
generators. To be able to use the existing machinery we want to be able to
obtain a generating set also for groups in our new representation. This can be
done (albeit not very effectively) by a stabilizer calculation in the
symmetric group given by the \texttt{degree} component. The operation function to use is probably a bit complicated and
will depend on the format of the \texttt{relations} (we have not specified in this example). In the following method we use \texttt{operationfunction} as a placeholder; 

 
\begin{Verbatim}[commandchars=@|A,fontsize=\small,frame=single,label=Example]
  InstallMethod(GeneratorsOfGroup,"for perm grps. given by relations",
    [IsPermutationGroupByRelations],
  function(G)
  local S,U;
    S:=SymmetricGroup(G!.degree);
    U:=Stabilizer(S,G!.relations,  operationfunction );
    return GeneratorsOfGroup(U);
  end);
\end{Verbatim}
 

 This is all we \emph{must} do. Of course for performance reasons one might want to install methods for
further operations as well. }

  
\section{\textcolor{Chapter }{Components versus Attributes}}\label{Components versus Attributes}
\logpage{[ 80, 7, 0 ]}
\hyperdef{L}{X86AA65D4815CAE95}{}
{
  In the last section we introduced two new components, \texttt{G!.degree} and \texttt{G!.relations}. Technically, we could have used attributes instead. There is no clear
distinction which variant is to be preferred: An attribute expresses part of
the functionality available to certain objects (and thus could be computed
later and probably even for a wider class of objects), a component is just
part of the internal definition of an object. 

 So if the data is ``of general interest'', if we want the user to have access to it, attributes are preferable.
Moreover, attributes can be used by the method selection (by specifying the
filter \texttt{HasAttr} for an attribute \texttt{Attr}). They provide a clean interface and their immutability makes it safe to hand
the data to a user who potentially could corrupt a components entries. 

 On the other hand more ``technical'' data (say the encoding of a sparse matrix) is better hidden from the user in a
component, as declaring it as an attribute would not give any advantage. 

 Resource\texttt{\symbol{45}}wise, attributes need more memory (the attribute
setter and tester are implicitly declared, and one filter bit is required),
the attribute access is one further function call in the kernel, thus
components might be an immeasurable bit faster. }

  
\section{\textcolor{Chapter }{Adding new Concepts}}\label{Adding new Concepts}
\logpage{[ 80, 8, 0 ]}
\hyperdef{L}{X7E29DEC0813F8897}{}
{
  Now we look how to implement a new concept for existing objects and fit this
in the method selection. Three examples that will be made more explicit below
would be groups for which a ``length'' of elements (as a word in certain generators) is defined, groups that can be
decomposed as a semidirect product and M\texttt{\symbol{45}}groups. 

 In each case we have two possibilities for the declaration. We can either
declare it as a property or as a category. Both are eventually filter(s) and
in this way indistinguishable for the method selection. However, the value of
a property for a particular object can be unknown at first and later in the
session be computed (to be \texttt{true} or \texttt{false}). This is implemented by reserving two filters for each property, one
indicating whether the property value is known, and one, provided the value is
known, to indicate the actual boolean value. Contrary to this, the decision
whether or not an object lies in a category is taken at creation time and this
is implemented using a single filter. 

 
\begin{description}
\item[{Property:}]  Properties also are attributes: If a property value is not known for an
object, \textsf{GAP} tries to find a method to compute the property value. If no suitable method is
found, an error is raised. 
\item[{Category:}]  An object is in a category if it has been created in it. Testing the category
for an object simply returns this value. Existing objects cannot enter a new
category later in life. This means that in most cases one has to write own
code to create objects in a new category. 

 If we want to implement a completely new concept so that new operations are
defined only for the new objects {\textendash}for example bialgebras for which
a second scalar multiplication is defined{\textendash} usually a category is
chosen. 

 Technically, the behaviour of the category \texttt{IsXYZ}, declared as subcategory of \texttt{IsABC} is therefore exactly the same as if we would declare \texttt{IsXYZ} to be a property for \texttt{IsABC} and install the following method: 
\begin{Verbatim}[commandchars=!@|,fontsize=\small,frame=single,label=Example]
  InstallMethod(IsXYZ,"return false if not known",[IsABC],ReturnFalse);
\end{Verbatim}
 

 (The word \texttt{category} also has a well\texttt{\symbol{45}}defined mathematical meaning, but this does
not need to concern us at this point. The set of objects which is defined to
be a (\textsf{GAP}) category does not need to be a category in the mathematical sense, vice
versa not every mathematical category is declared as a (\textsf{GAP}) category.) 
\end{description}
 Eventually the choice between category and property often becomes a matter of
taste or style. 

 Sometimes there is even a third possibility (if you have \textsf{GAP}{\nobreakspace}3 experience this might reflect most closely ``an object whose operations record is \texttt{XYOps}''): We might want to indicate this new concept simply by the fact that certain
attributes are set. In this case we could simply use the respective attribute
tester(s). 

 The examples given below each give a short argument why the respective
solution was chosen, but one could argue as well for other choices.  
\subsection{\textcolor{Chapter }{Example: M\texttt{\symbol{45}}groups}}\label{Example: M-groups}
\logpage{[ 80, 8, 1 ]}
\hyperdef{L}{X7DC936877A3330D0}{}
{
    M\texttt{\symbol{45}}groups are finite groups for which all irreducible
complex representations are induced from linear representations of subgroups,
it turns out that they are all solvable and that every supersolvable group is
an M\texttt{\symbol{45}}group. See \cite{Isa76} for further details. 

 Solvability and supersolvability both are testable properties. We therefore
declare \texttt{IsMGroup} as a property for solvable groups: 
\begin{Verbatim}[commandchars=!@|,fontsize=\small,frame=single,label=Example]
  IsMGroup:=NewProperty("IsMGroup",IsSolvableGroup);
\end{Verbatim}
 The filter \texttt{IsSolvableGroup} (\ref{IsSolvableGroup}) in this declaration \emph{only} means that methods for \texttt{IsMGroup} by default can only be installed for groups that are (and know to be) solvable
(though they could be installed for more general situations using \texttt{InstallOtherMethod} (\ref{InstallOtherMethod})). It does not yet imply that M\texttt{\symbol{45}}groups are solvable. We
must do this deliberately via an implication and we use the same technique to
imply that every supersolvable group is an M\texttt{\symbol{45}}group. 
\begin{Verbatim}[commandchars=!@|,fontsize=\small,frame=single,label=Example]
  InstallTrueMethod(IsSolvableGroup,IsMGroup);
  InstallTrueMethod(IsMGroup,IsSupersolvableGroup);
\end{Verbatim}
    

 Now we might install a method that tests for solvable groups whether they are
M\texttt{\symbol{45}}groups: 
\begin{Verbatim}[commandchars=!@|,fontsize=\small,frame=single,label=Example]
  InstallMethod(IsMGroup,"for solvable groups",[IsSolvableGroup],
  function(G)
    [... code omitted. The function must return `true' or `false' ...]
  end);
\end{Verbatim}
 Note that this example of declaring the \texttt{IsMGroup} property for solvable groups is not a part of the \textsf{GAP} library, which uses a similar but different filter \texttt{IsMonomialGroup} (\ref{IsMonomialGroup}). }

  
\subsection{\textcolor{Chapter }{Example: Groups with a word length}}\label{Example: Groups with a word length}
\logpage{[ 80, 8, 2 ]}
\hyperdef{L}{X7CD762FD82DED051}{}
{
  Our second example is that of groups for whose elements a \emph{word length} is defined. (We assume that the word length is only defined in the context of
the group with respect to a preselected generating set but not for single
elements alone. However we will not delve into any details of how this length
is defined and how it could be computed.) 

 Having a word length is a feature which enables other operations (for example
a ``word length'' function). This is exactly what categories are intended for and therefore we
use one. 

 First, we declare the category. All objects in this category are groups and so
we inherit the supercategory \texttt{IsGroup} (\ref{IsGroup}): 

 
\begin{Verbatim}[commandchars=!@|,fontsize=\small,frame=single,label=Example]
  DeclareCategory("IsGroupWithWordLength",IsGroup);
\end{Verbatim}
 

 We also define the operation which is ``enabled'' by this category, the word length of a group element, which is defined for a
group and an element (remember that group elements are described by the
category \texttt{IsMultiplicativeElementWithInverse} (\ref{IsMultiplicativeElementWithInverse})): 

 
\begin{Verbatim}[commandchars=!@|,fontsize=\small,frame=single,label=Example]
  DeclareOperation("WordLengthOfElement",[IsGroupWithWordLength,
    IsMultiplicativeElementWithInverse]);
\end{Verbatim}
 

 We then would proceed by installing methods to compute the word length in
concrete cases and might for example add further operations to get shortest
words in cosets. }

  
\subsection{\textcolor{Chapter }{Example: Groups with a decomposition as semidirect product}}\label{Example: Groups with a decomposition as semidirect product}
\logpage{[ 80, 8, 3 ]}
\hyperdef{L}{X782AC35979925C71}{}
{
  

 \index{DeclareAttribute"!example@\texttt{DeclareAttribute"!example}} The third example is groups which have a (nontrivial) decomposition as a
semidirect product. If this information has been found out, we want to be able
to use it in algorithms. (Thus we do not only need the fact \emph{that} there is a decomposition, but also the decomposition itself.) 

 We also want this to be applicable to every group and not only for groups
which have been explicitly constructed via \texttt{SemidirectProduct} (\ref{SemidirectProduct:for acting group, action, and a group}). 

 Instead we simply declare an attribute \texttt{SemidirectProductDecomposition} for groups. (Again, in this manual we don't go in the details of how such an
decomposition would look like). 

 
\begin{Verbatim}[commandchars=!@|,fontsize=\small,frame=single,label=Example]
  DeclareAttribute("SemidirectProductDecomposition",IsGroup);
\end{Verbatim}
 

 If a decomposition has been found, it can be stored in a group using \texttt{SetSemidirectProductDecomposition}. (At the moment all groups in \textsf{GAP} are attribute storing.) 

 Methods that rely on the existence of such a decomposition then get installed
for the tester filter \texttt{HasSemidirectProductDecomposition}.   }

 }

  
\section{\textcolor{Chapter }{Creating Own Arithmetic Objects}}\label{Creating Own Arithmetic Objects}
\logpage{[ 80, 9, 0 ]}
\hyperdef{L}{X7BD325C5791C6A06}{}
{
  Finally let's look at a way to create new objects with a
user\texttt{\symbol{45}}defined arithmetic such that one can form for example
groups, rings or vector spaces of these elements. This topic is discussed in
much more detail in chapter{\nobreakspace}\ref{An Example -- Designing Arithmetic Operations}, in this section we present a simple approach that may be useful to get
started but does not permit you to exploit all potential features. 

 The basic design is that the user designs some way to represent her objects in
terms of \textsf{GAP}s built\texttt{\symbol{45}}in types, for example as a list or a record. We
call this the ``defining data'' of the new objects. Also provided are functions that perform arithmetic on
this ``defining data'', that is they take objects of this form and return objects that represent the
result of the operation. The function \texttt{ArithmeticElementCreator} (\ref{ArithmeticElementCreator}) then is called to provide a wrapping such that proper new \textsf{GAP}\texttt{\symbol{45}}objects are created which can be multiplied etc. with the
default infix operations such as \texttt{*}. 

 

\subsection{\textcolor{Chapter }{ArithmeticElementCreator}}
\logpage{[ 80, 9, 1 ]}\nobreak
\hyperdef{L}{X87A88E3D7F6E2A7C}{}
{\noindent\textcolor{FuncColor}{$\triangleright$\enspace\texttt{ArithmeticElementCreator({\mdseries\slshape spec})\index{ArithmeticElementCreator@\texttt{ArithmeticElementCreator}}
\label{ArithmeticElementCreator}
}\hfill{\scriptsize (function)}}\\


 offers a simple interface to create new arithmetic elements by providing
functions that perform addition, multiplication and so forth, conforming to
the specification \mbox{\texttt{\mdseries\slshape spec}}. \texttt{ArithmeticElementCreator} creates a new category, representation and family for the new arithmetic
elements being defined, and returns a function which takes the ``defining data'' of an element and returns the corresponding new arithmetic element. 

 \mbox{\texttt{\mdseries\slshape spec}} is a record with one or more of the following components: 
\begin{description}
\item[{\texttt{ElementName}}]  string used to identify the new type of object. A global identifier \texttt{Is\mbox{\texttt{\mdseries\slshape ElementName}}} will be defined to indicate a category for these now objects. (Therefore it is
not clever to have blanks in the name). Also a collections category is
defined. (You will get an error message if the identifier \texttt{Is\mbox{\texttt{\mdseries\slshape ElementName}}} is already defined.) 
\item[{\texttt{Equality}, \texttt{LessThan}, \texttt{One}, \texttt{Zero}, \texttt{Multiplication}, \texttt{Inverse}, \texttt{Addition}, \texttt{AdditiveInverse}}]  functions defining the arithmetic operations. The functions interface on the
level of ``defining data'', the actual methods installed will perform the unwrapping and wrapping as
objects. Components are optional, but of course if no multiplication is
defined elements cannot be multiplied and so forth. 

 There are default methods for \texttt{Equality} and \texttt{LessThan} which simply calculate on the defining data. If one is defined, it must be
ensured that the other is compatible (so that $a < b$ implies not $(a = b)$) 
\item[{\texttt{Print}}]  a function which prints the object. By default, just the defining data is
printed. 
\item[{\texttt{MathInfo}}]  filters determining the mathematical properties of the elements created. A
typical value is for example \texttt{IsMultiplicativeElementWithInverse} for group elements. 
\item[{\texttt{RepInfo}}]  filters determining the representational properties of the elements created.
The objects created are always component objects, so in most cases the only
reasonable option is \texttt{IsAttributeStoringRep} to permit the storing of attributes. 
\end{description}
 

 All components are optional and will be filled in with default values (though
of course an empty record will not result in useful objects). 

 Note that the resulting objects are \emph{not equal} to their defining data (even though by default they print as only the defining
data). The operation \texttt{UnderlyingElement} can be used to obtain the defining data of such an element. }

 
\subsection{\textcolor{Chapter }{Example: ArithmeticElementCreator}}\label{expl:ArithmeticElementCreator}
\logpage{[ 80, 9, 2 ]}
\hyperdef{L}{X79E535CC7B82BA47}{}
{
  As the first example we look at subsets of $\{ 1, \ldots, 4 \}$ and define an ``addition'' as union and ``multiplication'' as intersection. These operations are both commutative and we want the
resulting elements to know this. 

 We therefore use the following specification: 

 
\begin{Verbatim}[commandchars=!@|,fontsize=\small,frame=single,label=Example]
  !gapprompt@gap>| !gapinput@# the whole set|
  !gapprompt@gap>| !gapinput@w := [1,2,3,4];|
  [ 1, 2, 3, 4 ]
  !gapprompt@gap>| !gapinput@PosetElementSpec :=rec(|
  !gapprompt@>| !gapinput@  # name of the new elements|
  !gapprompt@>| !gapinput@  ElementName := "PosetOn4",|
  !gapprompt@>| !gapinput@  # arithmetic operations|
  !gapprompt@>| !gapinput@  One := a -> w,|
  !gapprompt@>| !gapinput@  Zero := a -> [],|
  !gapprompt@>| !gapinput@  Multiplication := function(a, b) return Intersection(a, b); end,|
  !gapprompt@>| !gapinput@  Addition := function(a, b) return Union(a, b); end,|
  !gapprompt@>| !gapinput@  # Mathematical properties of the elements|
  !gapprompt@>| !gapinput@  MathInfo := IsCommutativeElement and|
  !gapprompt@>| !gapinput@              IsAssociativeElement and|
  !gapprompt@>| !gapinput@              IsAdditivelyCommutativeElement|
  !gapprompt@>| !gapinput@);;|
  !gapprompt@gap>| !gapinput@mkposet := ArithmeticElementCreator(PosetElementSpec);|
  function( x ) ... end
\end{Verbatim}
 

 Now we can create new elements, perform arithmetic on them and form domains: 

 
\begin{Verbatim}[commandchars=!@|,fontsize=\small,frame=single,label=Example]
  !gapprompt@gap>| !gapinput@a := mkposet([1,2,3]);|
  [ 1, 2, 3 ]
  !gapprompt@gap>| !gapinput@CategoriesOfObject(a);|
  [ "IsExtAElement", "IsNearAdditiveElement",
    "IsNearAdditiveElementWithZero", "IsAdditiveElement",
    "IsExtLElement", "IsExtRElement", "IsMultiplicativeElement",
    "IsMultiplicativeElementWithOne", "IsAssociativeElement",
    "IsAdditivelyCommutativeElement", "IsCommutativeElement",
    "IsPosetOn4" ]
  !gapprompt@gap>| !gapinput@a=[1,2,3];|
  false
  !gapprompt@gap>| !gapinput@UnderlyingElement(a)=[1,2,3];|
  true
  !gapprompt@gap>| !gapinput@b:=mkposet([2,3,4]);|
  [ 2, 3, 4 ]
  !gapprompt@gap>| !gapinput@a+b;|
  [ 1 .. 4 ]
  !gapprompt@gap>| !gapinput@a*b;|
  [ 2, 3 ]
  !gapprompt@gap>| !gapinput@s:=Semigroup(a,b);|
  <commutative semigroup with 2 generators>
  !gapprompt@gap>| !gapinput@Size(s);|
  3
\end{Verbatim}
 

 The categories \texttt{IsPosetOn4} and \texttt{IsPosetOn4Collection} can be used to install methods specific to the new objects. 
\begin{Verbatim}[commandchars=!@|,fontsize=\small,frame=single,label=Example]
  !gapprompt@gap>| !gapinput@IsPosetOn4Collection(s);|
  true
\end{Verbatim}
 }

 }

 }

          
\chapter{\textcolor{Chapter }{An Example {\textendash} Residue Class Rings}}\label{An Example -- Residue Class Rings}
\logpage{[ 81, 0, 0 ]}
\hyperdef{L}{X8125CC6A87409887}{}
{
  In this chapter, we give an example how \textsf{GAP} can be extended by new data structures and new functionality. In order to
focus on the issues of the implementation, the mathematics in the example
chosen is trivial. Namely, we will discuss computations with elements of
residue class rings ${\ensuremath{\mathbb Z}} / n{\ensuremath{\mathbb Z}}$. 

 The first attempt is straightforward (see Section{\nobreakspace}\ref{A First Attempt to Implement Elements of Residue Class Rings}), it deals with the implementation of the necessary arithmetic operations.
Section{\nobreakspace}\ref{Why Proceed in a Different Way?} deals with the question why it might be useful to use an approach that
involves creating a new data structure and integrating the algorithms dealing
with these new \textsf{GAP} objects into the system. Section{\nobreakspace}\ref{A Second Attempt to Implement Elements of Residue Class Rings} shows how this can be done in our example, and Section{\nobreakspace}\ref{Compatibility of Residue Class Rings with Prime Fields}, the question of further compatibility of the new objects with known \textsf{GAP} objects is discussed. Finally, Section{\nobreakspace}\ref{Further Improvements in Implementing Residue Class Rings} gives some hints how to improve the implementation presented before.  
\section{\textcolor{Chapter }{A First Attempt to Implement Elements of Residue Class Rings}}\label{A First Attempt to Implement Elements of Residue Class Rings}
\logpage{[ 81, 1, 0 ]}
\hyperdef{L}{X81008A74838A792E}{}
{
  Suppose we want to do computations with elements of a ring ${\ensuremath{\mathbb Z}} / n{\ensuremath{\mathbb Z}}$, where $n$ is a positive integer. 

 First we have to decide how to represent the element $k + n{\ensuremath{\mathbb Z}}$ in \textsf{GAP}. If the modulus $n$ is fixed then we can use the integer $k$. More precisely, we can use any integer $k'$ such that $k - k'$ is a multiple of $n$. If different moduli are likely to occur then using a list of the form $[ k, n ]$, or a record of the form \texttt{rec( residue := \mbox{\texttt{\mdseries\slshape k}}, modulus := \mbox{\texttt{\mdseries\slshape n}} )} is more appropriate. In the following, let us assume the list representation $[ k, n ]$ is chosen. Moreover, we decide that the residue $k$ in all such lists satisfies $0 \leq k < n$, i.e., the result of adding two residue classes represented by $[ k_1, n ]$ and $[ k_2, n ]$ (of course with same modulus $n$) will be $[ k, n ]$ with $k_1 + k_2$ congruent to $k$ modulo $n$ and $0 \leq k < n$. 

 Now we can implement the arithmetic operations for residue classes. Note that
the result of the \texttt{mod} operator is normalized as required. The division by a noninvertible residue
class results in \texttt{fail}. 

 
\begin{Verbatim}[commandchars=!@|,fontsize=\small,frame=single,label=Example]
  !gapprompt@gap>| !gapinput@resclass_sum := function( c1, c2 )|
  !gapprompt@>| !gapinput@   if c1[2] <> c2[2] then Error( "different moduli" ); fi;|
  !gapprompt@>| !gapinput@   return [ ( c1[1] + c2[1] ) mod c1[2], c1[2] ];|
  !gapprompt@>| !gapinput@end;;|
  !gapprompt@gap>| !gapinput@|
  !gapprompt@gap>| !gapinput@resclass_diff := function( c1, c2 )|
  !gapprompt@>| !gapinput@   if c1[2] <> c2[2] then Error( "different moduli" ); fi;|
  !gapprompt@>| !gapinput@   return [ ( c1[1] - c2[1] ) mod c1[2], c1[2] ];|
  !gapprompt@>| !gapinput@end;;|
  !gapprompt@gap>| !gapinput@|
  !gapprompt@gap>| !gapinput@resclass_prod := function( c1, c2 )|
  !gapprompt@>| !gapinput@   if c1[2] <> c2[2] then Error( "different moduli" ); fi;|
  !gapprompt@>| !gapinput@   return [ ( c1[1] * c2[1] ) mod c1[2], c1[2] ];|
  !gapprompt@>| !gapinput@end;;|
  !gapprompt@gap>| !gapinput@|
  !gapprompt@gap>| !gapinput@resclass_quo := function( c1, c2 )|
  !gapprompt@>| !gapinput@   local quo;|
  !gapprompt@>| !gapinput@   if c1[2] <> c2[2] then Error( "different moduli" ); fi;|
  !gapprompt@>| !gapinput@   quo:= QuotientMod( c1[1], c2[1], c1[2] );|
  !gapprompt@>| !gapinput@   if quo <> fail then|
  !gapprompt@>| !gapinput@     quo:= [ quo, c1[2] ];|
  !gapprompt@>| !gapinput@   fi;|
  !gapprompt@>| !gapinput@   return quo;|
  !gapprompt@>| !gapinput@end;;|
\end{Verbatim}
 

 With these functions, we can in principle compute with residue classes. 

 
\begin{Verbatim}[commandchars=!@|,fontsize=\small,frame=single,label=Example]
  !gapprompt@gap>| !gapinput@list:= List( [ 0 .. 3 ], k -> [ k, 4 ] );|
  [ [ 0, 4 ], [ 1, 4 ], [ 2, 4 ], [ 3, 4 ] ]
  !gapprompt@gap>| !gapinput@resclass_sum( list[2], list[4] );|
  [ 0, 4 ]
  !gapprompt@gap>| !gapinput@resclass_diff( list[1], list[2] );|
  [ 3, 4 ]
  !gapprompt@gap>| !gapinput@resclass_prod( list[2], list[4] );|
  [ 3, 4 ]
  !gapprompt@gap>| !gapinput@resclass_prod( list[3], list[4] );|
  [ 2, 4 ]
  !gapprompt@gap>| !gapinput@List( list, x -> resclass_quo( list[2], x ) );|
  [ fail, [ 1, 4 ], fail, [ 3, 4 ] ]
\end{Verbatim}
 }

  
\section{\textcolor{Chapter }{Why Proceed in a Different Way?}}\label{Why Proceed in a Different Way?}
\logpage{[ 81, 2, 0 ]}
\hyperdef{L}{X78B6425787FDB0E5}{}
{
  It depends on the computations we intended to do with residue classes whether
or not the implementation described in the previous section is satisfactory
for us. 

 Probably we are mainly interested in more complex data structures than the
residue classes themselves, for example in matrix algebras or matrix groups
over a ring such as ${\ensuremath{\mathbb Z}} / 4{\ensuremath{\mathbb Z}}$. For this, we need functions to add, multiply, invert
etc.{\nobreakspace}matrices of residue classes. Of course this is not a
difficult task, but it requires to write additional \textsf{GAP} code. 

 And when we have implemented the arithmetic operations for matrices of residue
classes, we might be interested in domain operations such as computing the
order of a matrix group over ${\ensuremath{\mathbb Z}} / 4{\ensuremath{\mathbb Z}}$, a Sylow $2$ subgroup, and so on. The problem is that a residue class represented as a pair $[ k, n ]$ is not regarded as a group element by \textsf{GAP}. We have not yet discussed how a matrix of residue classes shall be
represented, but if we choose the obvious representation of a list of lists of
our residue classes then also this is not a valid group element in \textsf{GAP}. Hence we cannot apply the function \texttt{Group} (\ref{Group:for a list of generators (and an identity element)}) to create a group of residue classes or a group of matrices of residue
classes. This is because \textsf{GAP} assumes that group elements can be multiplied via the infix operator \texttt{*} (equivalently, via the operation \texttt{\texttt{\symbol{92}}*} (\ref{bSlash*})). Note that in fact the multiplication of two lists $[ k_1, n ]$, $[ k_2, n ]$ is defined, but we have $[ k_1, n ] * [ k_2, n ] = k_1 * k_2 + n * n$, the standard scalar product of two row vectors of same length. That is, the
multiplication with \texttt{*} is not compatible with the function \texttt{resclass{\textunderscore}prod} introduced in the previous section. Similarly, ring elements are assumed to be
added via the infix operator \texttt{+}; the addition of residue classes is not compatible with the available
addition of row vectors. 

 What we have done in the previous section can be described as implementation
of a ``standalone'' arithmetic for residue classes. In order to use the machinery of the \textsf{GAP} library for creating higher level objects such as matrices, polynomials, or
domains over residue class rings, we have to ``integrate'' this implementation into the \textsf{GAP} library. The key step will be to create a new kind of \textsf{GAP} objects. This will be done in the following sections; there we assume that
residue classes and residue class rings are not yet available in \textsf{GAP}; in fact they are available, and their implementation is very close to what
is described here. }

  
\section{\textcolor{Chapter }{A Second Attempt to Implement Elements of Residue Class Rings}}\label{A Second Attempt to Implement Elements of Residue Class Rings}
\logpage{[ 81, 3, 0 ]}
\hyperdef{L}{X85B914DD81732492}{}
{
  Faced with the problem to implement elements of the rings ${\ensuremath{\mathbb Z}} / n{\ensuremath{\mathbb Z}}$, we must define the \emph{types} of these elements as far as is necessary to distinguish them from other \textsf{GAP} objects. 

 As is described in Chapter{\nobreakspace}\ref{Types of Objects}, the type of an object comprises several aspects of information about this
object; the \emph{family} determines the relation of the object to other objects, the \emph{categories} determine what operations the object admits, the \emph{representation} determines how an object is actually represented, and the \emph{attributes} describe knowledge about the object. 

 First of all, we must decide about the \emph{family} of each residue class. A natural way to do this is to put the elements of each
ring ${\ensuremath{\mathbb Z}} / n{\ensuremath{\mathbb Z}}$ into a family of their own. This means that for example elements of ${\ensuremath{\mathbb Z}} / 3{\ensuremath{\mathbb Z}}$ and ${\ensuremath{\mathbb Z}} / 9{\ensuremath{\mathbb Z}}$ lie in different families. So the only interesting relation between the
families of two residue classes is equality; binary arithmetic operations with
two residue classes will be admissible only if their families are equal. Note
that in the naive approach in Section{\nobreakspace}\ref{A First Attempt to Implement Elements of Residue Class Rings}, we had to take care of different moduli by a check in each function; these
checks may disappear in the new approach because of our choice of families. 

 Note that we do not need to tell \textsf{GAP} anything about the above decision concerning the families of the objects that
we are going to implement, that is, the \emph{declaration part} (see{\nobreakspace}\ref{Declaration and Implementation Part}) of the little \textsf{GAP} package we are writing contains nothing about the distribution of the new
objects into families. (The actual construction of a family happens in the
function \texttt{MyZmodnZ} shown below.) 

 Second, we want to describe methods to add or multiply two elements in ${\ensuremath{\mathbb Z}} / n{\ensuremath{\mathbb Z}}$, and these methods shall be not applicable to other \textsf{GAP} objects. The natural way to do this is to create a new \emph{category} in which all elements of all rings ${\ensuremath{\mathbb Z}} / n{\ensuremath{\mathbb Z}}$ lie. This is done as follows. 

 
\begin{Verbatim}[commandchars=!@|,fontsize=\small,frame=single,label=Example]
  !gapprompt@gap>| !gapinput@DeclareCategory( "IsMyZmodnZObj", IsScalar );|
  !gapprompt@gap>| !gapinput@cat:= CategoryCollections( IsMyZmodnZObj );;|
  !gapprompt@gap>| !gapinput@cat:= CategoryCollections( cat );;|
  !gapprompt@gap>| !gapinput@cat:= CategoryCollections( cat );;|
\end{Verbatim}
 

 So all elements in the rings ${\ensuremath{\mathbb Z}} / n{\ensuremath{\mathbb Z}}$ will lie in the category \texttt{IsMyZmodnZObj}, which is a subcategory of \texttt{IsScalar} (\ref{IsScalar}). The latter means that one can add, subtract, multiply and divide two such
elements that lie in the same family, with the obvious restriction that the
second operand of a division must be invertible. (The name \texttt{IsMyZmodnZObj} is chosen because \texttt{IsZmodnZObj} (\ref{IsZmodnZObj}) is already defined in \textsf{GAP}, for an implementation of residue classes that is very similar to the one
developed in this manual chapter. Using this different name, one can simply
enter the \textsf{GAP} code of this chapter into a \textsf{GAP} session, either interactively or by reading a file with this code, and
experiment after each step whether the expected behaviour has been achieved,
and what is still missing.) 

 The next lines of \textsf{GAP} code above create the categories \texttt{CategoryCollections( IsMyZmodnZObj )} and two higher levels of collections categories of this, which will be needed
later; it is important to create these categories before collections of the
objects in \texttt{IsMyZmodnZObj} actually arise. 

 Note that the only difference between \texttt{DeclareCategory} (\ref{DeclareCategory}) and \texttt{NewCategory} (\ref{NewCategory}) is that in a call to \texttt{DeclareCategory} (\ref{DeclareCategory}), a variable corresponding to the first argument is set to the new category,
and this variable is read\texttt{\symbol{45}}only. The same holds for \texttt{DeclareRepresentation} (\ref{DeclareRepresentation}) and \texttt{NewRepresentation} (\ref{NewRepresentation}) etc. 

 There is no analogue of categories in the implementation in
Section{\nobreakspace}\ref{A First Attempt to Implement Elements of Residue Class Rings}, since there it was not necessary to distinguish residue classes from other \textsf{GAP} objects. Note that the functions there assumed that their arguments were
residue classes, and the user was responsible not to call them with other
arguments. Thus an important aspect of types is to describe arguments of
functions explicitly. 

 Third, we must decide about the \emph{representation} of our objects. This is something we know already from Section{\nobreakspace}\ref{A First Attempt to Implement Elements of Residue Class Rings}, where we chose a list of length two. Here we may choose between two
essentially different representations for the new \textsf{GAP} objects, namely as ``component object'' (record\texttt{\symbol{45}}like) or ``positional object'' (list\texttt{\symbol{45}}like). We decide to store the modulus of each residue
class in its family, and to encode the element $k + n{\ensuremath{\mathbb Z}}$ by the unique residue in the range $[ 0 .. n-1 ]$ that is congruent to $k$ modulo $n$, and the object itself is chosen to be a positional object with this residue
at the first and only position (see{\nobreakspace}\ref{Positional Objects}). 

 
\begin{Verbatim}[commandchars=!@|,fontsize=\small,frame=single,label=Example]
  !gapprompt@gap>| !gapinput@DeclareRepresentation("IsMyModulusRep", IsPositionalObjectRep, [1]);|
\end{Verbatim}
 

 The fourth ingredients of a type, \emph{attributes}, are usually of minor importance for element objects. In particular, we do
not need to introduce special attributes for residue classes. 

 Having defined what the new objects shall look like, we now declare a global
function (see{\nobreakspace}\ref{Declaration and Implementation Part}), to create an element when family and residue are given. 

 
\begin{Verbatim}[commandchars=!@|,fontsize=\small,frame=single,label=Example]
  !gapprompt@gap>| !gapinput@DeclareGlobalFunction( "MyZmodnZObj" );|
\end{Verbatim}
 

 Now we have declared what we need, and we can start to implement the missing
methods resp. functions; so the following command belongs to the \emph{implementation part} of our package (see{\nobreakspace}\ref{Declaration and Implementation Part}). 

 The probably most interesting function is the one to construct a residue
class. 

 
\begin{Verbatim}[commandchars=@|A,fontsize=\small,frame=single,label=Example]
  @gapprompt|gap>A @gapinput|InstallGlobalFunction( MyZmodnZObj, function( Fam, residue )A
  @gapprompt|>A @gapinput|   return Objectify( NewType( Fam, IsMyZmodnZObj and IsMyModulusRep ),A
  @gapprompt|>A @gapinput|                     [ residue mod Fam!.modulus ] );A
  @gapprompt|>A @gapinput|end );A
\end{Verbatim}
 

 Note that we normalize \texttt{residue} explicitly using \texttt{mod}; we assumed that the modulus is stored in \texttt{Fam}, so we must take care of this below. If \texttt{Fam} is a family of residue classes, and \texttt{residue} is an integer, \texttt{MyZmodnZObj} returns the corresponding object in the family \texttt{Fam}, which lies in the category \texttt{IsMyZmodnZObj} and in the representation \texttt{IsMyModulusRep}. 

 \texttt{MyZmodnZObj} needs an appropriate family as first argument, so let us see how to get our
hands on this. Of course we could write a handy function to create such a
family for given modulus, but we choose another way. In fact we do not really
want to call \texttt{MyZmodnZObj} explicitly when we want to create residue classes. For example, if we want to
enter a matrix of residues then usually we start with a matrix of
corresponding integers, and it is more elegant to do the conversion via
multiplying the matrix with the identity of the required ring ${\ensuremath{\mathbb Z}} / n{\ensuremath{\mathbb Z}}$; this is also done for the conversion of integral matrices to finite field
matrices. (Note that we will have to install a method for this.) So it is
often sufficient to access this identity, for example via \texttt{One( MyZmodnZ( \mbox{\texttt{\mdseries\slshape n}} ) )}, where \texttt{MyZmodnZ} returns a domain representing the ring ${\ensuremath{\mathbb Z}} / n{\ensuremath{\mathbb Z}}$ when called with the argument $n$. We decide that constructing this ring is a natural place where the creation
of the family can be hidden, and implement the function. (Note that the
declaration belongs to the declaration part, and the installation belongs to
the implementation part, see{\nobreakspace}\ref{Declaration and Implementation Part}). 

 
\begin{Verbatim}[commandchars=@|A,fontsize=\small,frame=single,label=Example]
  @gapprompt|gap>A @gapinput|DeclareGlobalFunction( "MyZmodnZ" );A
  @gapprompt|gap>A @gapinput|A
  @gapprompt|gap>A @gapinput|InstallGlobalFunction( MyZmodnZ, function( n )A
  @gapprompt|>A @gapinput|   local F, R;A
  @gapprompt|>A @gapinput|A
  @gapprompt|>A @gapinput|   if not IsPosInt( n ) thenA
  @gapprompt|>A @gapinput|     Error( "<n> must be a positive integer" );A
  @gapprompt|>A @gapinput|   fi;A
  @gapprompt|>A @gapinput|A
  @gapprompt|>A @gapinput|   # Construct the family of element objects of our ring.A
  @gapprompt|>A @gapinput|   F:= NewFamily( Concatenation( "MyZmod", String( n ), "Z" ),A
  @gapprompt|>A @gapinput|                  IsMyZmodnZObj );A
  @gapprompt|>A @gapinput|A
  @gapprompt|>A @gapinput|   # Install the data.A
  @gapprompt|>A @gapinput|   F!.modulus:= n;A
  @gapprompt|>A @gapinput|A
  @gapprompt|>A @gapinput|   # Make the domain.A
  @gapprompt|>A @gapinput|   R:= RingWithOneByGenerators( [ MyZmodnZObj( F, 1 ) ] );A
  @gapprompt|>A @gapinput|   SetIsWholeFamily( R, true );A
  @gapprompt|>A @gapinput|   SetName( R, Concatenation( "(Integers mod ", String(n), ")" ) );A
  @gapprompt|>A @gapinput|A
  @gapprompt|>A @gapinput|   # Return the ring.A
  @gapprompt|>A @gapinput|   return R;A
  @gapprompt|>A @gapinput|end );A
\end{Verbatim}
 

 Note that the modulus \texttt{n} is stored in the component \texttt{modulus} of the family, as is assumed by \texttt{MyZmodnZ}. Thus it is not necessary to store the modulus in each element. When storing \texttt{n} with the \texttt{!.} operator as value of the component \texttt{modulus}, we used that all families are in fact represented as component objects
(see{\nobreakspace}\ref{Component Objects}). 

 We see that we can use \texttt{RingWithOneByGenerators} (\ref{RingWithOneByGenerators}) to construct a ring with one if we have the appropriate generators. The
construction via \texttt{RingWithOneByGenerators} (\ref{RingWithOneByGenerators}) makes sure that \texttt{IsRingWithOne} (\ref{IsRingWithOne}) (and \texttt{IsRing} (\ref{IsRing})) is \texttt{true} for each output of \texttt{MyZmodnZ}. So the main problem is to create the identity element of the ring, which in
our case suffices to generate the ring. In order to create this element via \texttt{MyZmodnZObj}, we have to construct its family first, at each call of \texttt{MyZmodnZ}. 

 Also note that we may enter known information about the ring. Here we store
that it contains the whole family of elements; this is useful for example when
we want to check the membership of an element in the ring, which can be
decided from the type of the element if the ring contains its whole elements
family. Giving a name to the ring causes that it will be printed via printing
the name. (By the way: This name \texttt{(Integers mod \mbox{\texttt{\mdseries\slshape n}})} looks like a call to \texttt{\texttt{\symbol{92}}mod} (\ref{bSlashmod}) with the arguments \texttt{Integers} (\ref{Integers}) and \mbox{\texttt{\mdseries\slshape n}}; a construction of the ring via this call seems to be more natural than by
calling \texttt{MyZmodnZ}; later we shall install a \texttt{\texttt{\symbol{92}}mod} (\ref{bSlashmod}) method in order to admit this construction.) 

 Now we can read the above code into \textsf{GAP}, and the following works already. 

 
\begin{Verbatim}[commandchars=!@|,fontsize=\small,frame=single,label=Example]
  !gapprompt@gap>| !gapinput@R:= MyZmodnZ( 4 );|
  (Integers mod 4)
  !gapprompt@gap>| !gapinput@IsRing( R );|
  true
  !gapprompt@gap>| !gapinput@gens:= GeneratorsOfRingWithOne( R );|
  [ <object> ]
\end{Verbatim}
 

 But of course this means just to ask for the information we have explicitly
stored in the ring. Already the questions whether the ring is finite and how
many elements it has, cannot be answered by \textsf{GAP}. Clearly we know the answers, and we could store them in the ring, by setting
the value of the property \texttt{IsFinite} (\ref{IsFinite}) to \texttt{true} and the value of the attribute \texttt{Size} (\ref{Size}) to \mbox{\texttt{\mdseries\slshape n}} (the argument of the call to \texttt{MyZmodnZ}). If we do not want to do so then \textsf{GAP} could only try to find out the number of elements of the ring via forming the
closure of the generators under addition and multiplication, but up to now, \textsf{GAP} does not know how to add or multiply two elements of our ring. 

 So we must install some methods for arithmetic and other operations if the
elements are to behave as we want. 

 We start with a method for showing elements nicely on the screen. There are
different operations for this purpose. One of them is \texttt{PrintObj} (\ref{PrintObj}), which is called for each argument in an explicit call to \texttt{Print} (\ref{Print}). Another one is \texttt{ViewObj} (\ref{ViewObj}), which is called in the read\texttt{\symbol{45}}eval\texttt{\symbol{45}}print
loop for each object. \texttt{ViewObj} (\ref{ViewObj}) shall produce short and human readable information about the object in
question, whereas \texttt{PrintObj} (\ref{PrintObj}) shall produce information that may be longer and is (if reasonable) readable
by \textsf{GAP}. We cannot satisfy the latter requirement for a \texttt{PrintObj} (\ref{PrintObj}) method because there is no way to make a family \textsf{GAP} readable. So we decide to display the expression \texttt{( k mod n )} for an object that is given by the residue \texttt{k} and the modulus \texttt{n}, which would be fine as a \texttt{ViewObj} (\ref{ViewObj}) method. Since the default for \texttt{ViewObj} (\ref{ViewObj}) is to call \texttt{PrintObj} (\ref{PrintObj}), and since no other \texttt{ViewObj} (\ref{ViewObj}) method is applicable to our elements, we need only a \texttt{PrintObj} (\ref{PrintObj}) method. 

 
\begin{Verbatim}[commandchars=@|A,fontsize=\small,frame=single,label=Example]
  @gapprompt|gap>A @gapinput|InstallMethod( PrintObj,A
  @gapprompt|>A @gapinput|   "for element in Z/nZ (ModulusRep)",A
  @gapprompt|>A @gapinput|   [ IsMyZmodnZObj and IsMyModulusRep ],A
  @gapprompt|>A @gapinput|   function( x )A
  @gapprompt|>A @gapinput|   Print( "( ", x![1], " mod ", FamilyObj(x)!.modulus, " )" );A
  @gapprompt|>A @gapinput|   end );A
\end{Verbatim}
 

 So we installed a method for the operation \texttt{PrintObj} (\ref{PrintObj}) (first argument), and we gave it a suitable information message (second
argument), see{\nobreakspace}\ref{ApplicableMethod} and{\nobreakspace}\ref{Tracing Methods} for applications of this information string. The third argument tells \textsf{GAP} that the method is applicable for objects that lie in the category \texttt{IsMyZmodnZObj} and in the representation \texttt{IsMyModulusRep}. and the fourth argument is the method itself. More details about \texttt{InstallMethod} (\ref{InstallMethod}) can be found in{\nobreakspace}\ref{Method Installation}. 

 Note that the requirement \texttt{IsMyModulusRep} for the argument \texttt{x} allows us to access the residue as \texttt{x![1]}. Since the family of \texttt{x} has the component \texttt{modulus} bound if it is constructed by \texttt{MyZmodnZ}, we may access this component. We check whether the method installation has
some effect. 

 
\begin{Verbatim}[commandchars=!@|,fontsize=\small,frame=single,label=Example]
  !gapprompt@gap>| !gapinput@gens;|
  [ ( 1 mod 4 ) ]
\end{Verbatim}
 

 Next we install methods for the comparison operations. Note that we can assume
that the residues in the representation chosen are normalized. 

 
\begin{Verbatim}[commandchars=@|A,fontsize=\small,frame=single,label=Example]
  @gapprompt|gap>A @gapinput|InstallMethod( \=,A
  @gapprompt|>A @gapinput|   "for two elements in Z/nZ (ModulusRep)",A
  @gapprompt|>A @gapinput|   IsIdenticalObj,A
  @gapprompt|>A @gapinput|   [IsMyZmodnZObj and IsMyModulusRep, IsMyZmodnZObj and IsMyModulusRep],A
  @gapprompt|>A @gapinput|   function( x, y ) return x![1] = y![1]; end );A
  @gapprompt|gap>A @gapinput|A
  @gapprompt|gap>A @gapinput|InstallMethod( \<,A
  @gapprompt|>A @gapinput|   "for two elements in Z/nZ (ModulusRep)",A
  @gapprompt|>A @gapinput|   IsIdenticalObj,A
  @gapprompt|>A @gapinput|   [IsMyZmodnZObj and IsMyModulusRep, IsMyZmodnZObj and IsMyModulusRep],A
  @gapprompt|>A @gapinput|   function( x, y ) return x![1] < y![1]; end );A
\end{Verbatim}
 

 The third argument used in these installations specifies the required relation
between the families of the arguments (see{\nobreakspace}\ref{Families}). This argument of a method installation, if present, is a function that
shall be applied to the families of the arguments. \texttt{IsIdenticalObj} (\ref{IsIdenticalObj}) means that the methods are applicable only if both arguments lie in the same
family. (In installations for unary methods, obviously no relation is
required, so this argument is left out there.) 

 Up to now, we see no advantage of the new approach over the one in
Section{\nobreakspace}\ref{A First Attempt to Implement Elements of Residue Class Rings}. For a residue class represented as \texttt{[ \mbox{\texttt{\mdseries\slshape k}}, \mbox{\texttt{\mdseries\slshape n}} ]}, the way it is printed on the screen is sufficient, and equality and
comparison of lists are good enough to define equality and comparison of
residue classes if needed. But this is not the case in other situations. For
example, if we would have decided that the residue \mbox{\texttt{\mdseries\slshape k}} need not be normalized then we would have needed functions in
Section{\nobreakspace}\ref{A First Attempt to Implement Elements of Residue Class Rings} that compute whether two residue classes are equal, and which of two residue
classes is regarded as larger than another. Note that we are free to define
what ``larger'' means for objects that are newly introduced. 

 Next we install methods for the arithmetic operations, first for the additive
structure. 

 
\begin{Verbatim}[commandchars=@|B,fontsize=\small,frame=single,label=Example]
  @gapprompt|gap>B @gapinput|InstallMethod( \+,B
  @gapprompt|>B @gapinput|   "for two elements in Z/nZ (ModulusRep)",B
  @gapprompt|>B @gapinput|   IsIdenticalObj,B
  @gapprompt|>B @gapinput|   [IsMyZmodnZObj and IsMyModulusRep, IsMyZmodnZObj and IsMyModulusRep],B
  @gapprompt|>B @gapinput|   function( x, y )B
  @gapprompt|>B @gapinput|   return MyZmodnZObj( FamilyObj( x ), x![1] + y![1] );B
  @gapprompt|>B @gapinput|   end );B
  @gapprompt|gap>B @gapinput|B
  @gapprompt|gap>B @gapinput|InstallMethod( ZeroOp,B
  @gapprompt|>B @gapinput|   "for element in Z/nZ (ModulusRep)",B
  @gapprompt|>B @gapinput|   [ IsMyZmodnZObj ],B
  @gapprompt|>B @gapinput|   x -> MyZmodnZObj( FamilyObj( x ), 0 ) );B
  @gapprompt|gap>B @gapinput|B
  @gapprompt|gap>B @gapinput|InstallMethod( AdditiveInverseOp,B
  @gapprompt|>B @gapinput|   "for element in Z/nZ (ModulusRep)",B
  @gapprompt|>B @gapinput|   [ IsMyZmodnZObj and IsMyModulusRep ],B
  @gapprompt|>B @gapinput|   x -> MyZmodnZObj( FamilyObj( x ), AdditiveInverse( x![1] ) ) );B
\end{Verbatim}
 

 Here the new approach starts to pay off. The method for the operation \texttt{\texttt{\symbol{92}}+} (\ref{bSlash+}) allows us to use the infix operator \texttt{+} for residue classes. The method for \texttt{ZeroOp} (\ref{ZeroOp}) is used when we call this operation or the attribute \texttt{Zero} (\ref{Zero}) explicitly, and \texttt{ZeroOp} (\ref{ZeroOp}) it is also used when we ask for \texttt{0 * \mbox{\texttt{\mdseries\slshape rescl}}}, where \mbox{\texttt{\mdseries\slshape rescl}} is a residue class. 

 (Note that \texttt{Zero} (\ref{Zero}) and \texttt{ZeroOp} (\ref{ZeroOp}) are distinguished because \texttt{0 * \mbox{\texttt{\mdseries\slshape obj}}} is guaranteed to return a \emph{mutable} result whenever a mutable version of this result exists in \textsf{GAP} {\textendash}for example if \mbox{\texttt{\mdseries\slshape obj}} is a matrix{\textendash} whereas \texttt{Zero} (\ref{Zero}) is an attribute and therefore returns \emph{immutable} results; for our example there is no difference since the residue classes are
always immutable, nevertheless we have to install the method for \texttt{ZeroOp} (\ref{ZeroOp}). The same holds for \texttt{AdditiveInverse} (\ref{AdditiveInverse}), \texttt{One} (\ref{One}), and \texttt{Inverse} (\ref{Inverse}).) 

 Similarly, \texttt{AdditiveInverseOp} (\ref{AdditiveInverseOp}) can be either called directly or via the unary \texttt{\texttt{\symbol{45}}} operator; so we can compute the additive inverse of the residue class \mbox{\texttt{\mdseries\slshape rescl}} as \texttt{\texttt{\symbol{45}}\mbox{\texttt{\mdseries\slshape rescl}}}. 

 It is not necessary to install methods for subtraction, since this is handled
via addition of the additive inverse of the second argument if no other method
is installed. 

 Let us try what we can do with the methods that are available now. 

 
\begin{Verbatim}[commandchars=!@|,fontsize=\small,frame=single,label=Example]
  !gapprompt@gap>| !gapinput@x:= gens[1];  y:= x + x;|
  ( 1 mod 4 )
  ( 2 mod 4 )
  !gapprompt@gap>| !gapinput@0 * x;  -x;|
  ( 0 mod 4 )
  ( 3 mod 4 )
  !gapprompt@gap>| !gapinput@y = -y;  x = y;  x < y;  -x < y;|
  true
  false
  true
  false
\end{Verbatim}
 

 We might want to admit the addition of integers and elements in rings ${\ensuremath{\mathbb Z}} / n{\ensuremath{\mathbb Z}}$, where an integer is implicitly identified with its residue modulo $n$. To achieve this, we install methods to add an integer to an object in \texttt{IsMyZmodnZObj} from the left and from the right. 

 
\begin{Verbatim}[commandchars=@|A,fontsize=\small,frame=single,label=Example]
  @gapprompt|gap>A @gapinput|InstallMethod( \+,A
  @gapprompt|>A @gapinput|   "for element in Z/nZ (ModulusRep) and integer",A
  @gapprompt|>A @gapinput|   [ IsMyZmodnZObj and IsMyModulusRep, IsInt ],A
  @gapprompt|>A @gapinput|   function( x, y )A
  @gapprompt|>A @gapinput|   return MyZmodnZObj( FamilyObj( x ), x![1] + y );A
  @gapprompt|>A @gapinput|   end );A
  @gapprompt|gap>A @gapinput|A
  @gapprompt|gap>A @gapinput|InstallMethod( \+,A
  @gapprompt|>A @gapinput|   "for integer and element in Z/nZ (ModulusRep)",A
  @gapprompt|>A @gapinput|   [ IsInt, IsMyZmodnZObj and IsMyModulusRep ],A
  @gapprompt|>A @gapinput|   function( x, y )A
  @gapprompt|>A @gapinput|   return MyZmodnZObj( FamilyObj( y ), x + y![1] );A
  @gapprompt|>A @gapinput|   end );A
\end{Verbatim}
 

 Now we can do also the following. 

 
\begin{Verbatim}[commandchars=!@|,fontsize=\small,frame=single,label=Example]
  !gapprompt@gap>| !gapinput@2 + x;  7 - x;  y - 2;|
  ( 3 mod 4 )
  ( 2 mod 4 )
  ( 0 mod 4 )
\end{Verbatim}
 

 Similarly we install the methods dealing with the multiplicative structure. We
need methods to multiply two of our objects, and to compute identity and
inverse. The operation \texttt{OneOp} (\ref{OneOp}) is called when we ask for \texttt{\mbox{\texttt{\mdseries\slshape rescl}}\texttt{\symbol{94}}0}, and \texttt{InverseOp} (\ref{InverseOp}) is called when we ask for \texttt{\mbox{\texttt{\mdseries\slshape rescl}}\texttt{\symbol{94}}\texttt{\symbol{45}}1}. Note that the method for \texttt{InverseOp} (\ref{InverseOp}) returns \texttt{fail} if the argument is not invertible. 

 
\begin{Verbatim}[commandchars=@|A,fontsize=\small,frame=single,label=Example]
  @gapprompt|gap>A @gapinput|InstallMethod( \*,A
  @gapprompt|>A @gapinput|   "for two elements in Z/nZ (ModulusRep)",A
  @gapprompt|>A @gapinput|   IsIdenticalObj,A
  @gapprompt|>A @gapinput|   [IsMyZmodnZObj and IsMyModulusRep, IsMyZmodnZObj and IsMyModulusRep],A
  @gapprompt|>A @gapinput|   function( x, y )A
  @gapprompt|>A @gapinput|   return MyZmodnZObj( FamilyObj( x ), x![1] * y![1] );A
  @gapprompt|>A @gapinput|   end );A
  @gapprompt|gap>A @gapinput|A
  @gapprompt|gap>A @gapinput|InstallMethod( OneOp,A
  @gapprompt|>A @gapinput|   "for element in Z/nZ (ModulusRep)",A
  @gapprompt|>A @gapinput|   [ IsMyZmodnZObj ],A
  @gapprompt|>A @gapinput|   elm -> MyZmodnZObj( FamilyObj( elm ), 1 ) );A
  @gapprompt|gap>A @gapinput|A
  @gapprompt|gap>A @gapinput|InstallMethod( InverseOp,A
  @gapprompt|>A @gapinput|   "for element in Z/nZ (ModulusRep)",A
  @gapprompt|>A @gapinput|   [ IsMyZmodnZObj and IsMyModulusRep ],A
  @gapprompt|>A @gapinput|   function( elm )A
  @gapprompt|>A @gapinput|   local residue;A
  @gapprompt|>A @gapinput|   residue:= QuotientMod( 1, elm![1], FamilyObj( elm )!.modulus );A
  @gapprompt|>A @gapinput|   if residue <> fail thenA
  @gapprompt|>A @gapinput|     residue:= MyZmodnZObj( FamilyObj( elm ), residue );A
  @gapprompt|>A @gapinput|   fi;A
  @gapprompt|>A @gapinput|   return residue;A
  @gapprompt|>A @gapinput|   end );A
\end{Verbatim}
 

 To be able to multiply our objects with integers, we need not (but we may, and
we should if we are going for efficiency) install special methods. This is
because in general, \textsf{GAP} interprets the multiplication of an integer and an additive object as
abbreviation of successive additions, and there is one generic method for such
a multiplication that uses only additions and {\textendash}in the case of a
negative integer{\textendash} taking the additive inverse. Analogously, there
is a generic method for powering by integers that uses only multiplications
and taking the multiplicative inverse. 

 Note that we could also interpret the multiplication with an integer as a
shorthand for the multiplication with the corresponding residue class. We are
lucky that this interpretation is compatible with the one that is already
available. If this would not be the case then of course we would get into
trouble by installing a concurrent multiplication that computes something
different from the multiplication that is already defined, since \textsf{GAP} does not guarantee which of the applicable methods is actually chosen
(see{\nobreakspace}\ref{Applicable Methods and Method Selection}). 

 Now we have implemented methods for the arithmetic operations for our
elements, and the following calculations work. 

 
\begin{Verbatim}[commandchars=!@|,fontsize=\small,frame=single,label=Example]
  !gapprompt@gap>| !gapinput@y:= 2 * x;  z:= (-5) * x;|
  ( 2 mod 4 )
  ( 3 mod 4 )
  !gapprompt@gap>| !gapinput@y * z;  y * y;|
  ( 2 mod 4 )
  ( 0 mod 4 )
  !gapprompt@gap>| !gapinput@y^-1;  y^0;|
  fail
  ( 1 mod 4 )
  !gapprompt@gap>| !gapinput@z^-1;|
  ( 3 mod 4 )
\end{Verbatim}
 

 There are some other operations in \textsf{GAP} that we may want to accept our elements as arguments. An example is the
operation \texttt{Int} (\ref{Int}) that returns, e.g., the integral part of a rational number or the integer
corresponding to an element in a finite prime field. For our objects, we may
define that \texttt{Int} (\ref{Int}) returns the normalized residue. 

 Note that we \emph{define} this behaviour for elements but we \emph{implement} it for objects in the representation \texttt{IsMyModulusRep}. This means that if someone implements another representation of residue
classes then this person must be careful to implement \texttt{Int} (\ref{Int}) methods for objects in this new representation compatibly with our definition,
i.e., such that the result is independent of the representation. 

 
\begin{Verbatim}[commandchars=@|A,fontsize=\small,frame=single,label=Example]
  @gapprompt|gap>A @gapinput|InstallMethod( Int,A
  @gapprompt|>A @gapinput|   "for element in Z/nZ (ModulusRep)",A
  @gapprompt|>A @gapinput|   [ IsMyZmodnZObj and IsMyModulusRep ],A
  @gapprompt|>A @gapinput|   z -> z![1] );A
\end{Verbatim}
 

 Another example of an operation for which we might want to install a method is \texttt{\texttt{\symbol{92}}mod} (\ref{bSlashmod}). We make the ring print itself as \texttt{Integers} (\ref{Integers}) mod the modulus, and then it is reasonable to allow a construction this way,
which makes the \texttt{PrintObj} (\ref{PrintObj}) output of the ring \textsf{GAP} readable. 

 
\begin{Verbatim}[commandchars=@|A,fontsize=\small,frame=single,label=Example]
  @gapprompt|gap>A @gapinput|InstallMethod( PrintObj,A
  @gapprompt|>A @gapinput|   "for full collection Z/nZ",A
  @gapprompt|>A @gapinput|   [ CategoryCollections( IsMyZmodnZObj ) and IsWholeFamily ],A
  @gapprompt|>A @gapinput|   function( R )A
  @gapprompt|>A @gapinput|   Print( "(Integers mod ",A
  @gapprompt|>A @gapinput|          ElementsFamily( FamilyObj(R) )!.modulus, ")" );A
  @gapprompt|>A @gapinput|   end );A
  @gapprompt|gap>A @gapinput|A
  @gapprompt|gap>A @gapinput|InstallMethod( \mod,A
  @gapprompt|>A @gapinput|   "for `Integers', and a positive integer",A
  @gapprompt|>A @gapinput|   [ IsIntegers, IsPosRat and IsInt ],A
  @gapprompt|>A @gapinput|   function( Integers, n ) return MyZmodnZ( n ); end );A
\end{Verbatim}
 

 Let us try this. 

 
\begin{Verbatim}[commandchars=!@|,fontsize=\small,frame=single,label=Example]
  !gapprompt@gap>| !gapinput@Int( y );|
  2
  !gapprompt@gap>| !gapinput@Integers mod 1789;|
  (Integers mod 1789)
\end{Verbatim}
 

 Probably it is not necessary to emphasize that with the approach of
Section{\nobreakspace}\ref{A First Attempt to Implement Elements of Residue Class Rings}, installing methods for existing operations is usually not possible or at
least not recommended. For example, installing the function \texttt{resclass{\textunderscore}sum} defined in Section{\nobreakspace}\ref{A First Attempt to Implement Elements of Residue Class Rings} as a \texttt{\texttt{\symbol{92}}+} (\ref{bSlash+}) method for adding two lists of length two (with integer entries) would not be
compatible with the general definition of the addition of two lists of same
length. Installing a method for the operation \texttt{Int} (\ref{Int}) that takes a list \texttt{[ \mbox{\texttt{\mdseries\slshape k}}, \mbox{\texttt{\mdseries\slshape n}} ]} and returns \mbox{\texttt{\mdseries\slshape k}} would in principle be possible, since there is no \texttt{Int} (\ref{Int}) method for lists yet, but it is not sensible to do so because one can think of
other interpretations of such a list where different \texttt{Int} (\ref{Int}) methods could be installed with the same right. 

 As mentioned in Section{\nobreakspace}\ref{Why Proceed in a Different Way?}, one advantage of the new approach is that with the implementation we have up
to now, automatically also matrices of residue classes can be treated. 

 
\begin{Verbatim}[commandchars=!@|,fontsize=\small,frame=single,label=Example]
  !gapprompt@gap>| !gapinput@r:= Integers mod 16;|
  (Integers mod 16)
  !gapprompt@gap>| !gapinput@x:= One( r );|
  ( 1 mod 16 )
  !gapprompt@gap>| !gapinput@mat:= IdentityMat( 2 ) * x;|
  [ [ ( 1 mod 16 ), ( 0 mod 16 ) ], [ ( 0 mod 16 ), ( 1 mod 16 ) ] ]
  !gapprompt@gap>| !gapinput@mat[1][2]:= x;;|
  !gapprompt@gap>| !gapinput@mat;|
  [ [ ( 1 mod 16 ), ( 1 mod 16 ) ], [ ( 0 mod 16 ), ( 1 mod 16 ) ] ]
  !gapprompt@gap>| !gapinput@Order( mat );|
  16
  !gapprompt@gap>| !gapinput@mat + mat;|
  [ [ ( 2 mod 16 ), ( 2 mod 16 ) ], [ ( 0 mod 16 ), ( 2 mod 16 ) ] ]
  !gapprompt@gap>| !gapinput@last^4;|
  [ [ ( 0 mod 16 ), ( 0 mod 16 ) ], [ ( 0 mod 16 ), ( 0 mod 16 ) ] ]
\end{Verbatim}
 

 Such matrices, if they are invertible, are valid as group elements. One
technical problem is that the default algorithm for inverting matrices may
give up since Gaussian elimination need not be successful over rings
containing zero divisors. Therefore we install a simpleminded inversion method
that inverts an integer matrix. 

 
\begin{Verbatim}[commandchars=@|A,fontsize=\small,frame=single,label=Example]
  @gapprompt|gap>A @gapinput|InstallMethod( InverseOp,A
  @gapprompt|>A @gapinput|   "for an ordinary matrix over a ring Z/nZ",A
  @gapprompt|>A @gapinput|   [ IsMatrix and IsOrdinaryMatrixA
  @gapprompt|>A @gapinput|     and CategoryCollections( CategoryCollections( IsMyZmodnZObj ) ) ],A
  @gapprompt|>A @gapinput|   function( mat )A
  @gapprompt|>A @gapinput|   local one, modulus;A
  @gapprompt|>A @gapinput|A
  @gapprompt|>A @gapinput|   one:= One( mat[1][1] );A
  @gapprompt|>A @gapinput|   modulus:= FamilyObj( one )!.modulus;A
  @gapprompt|>A @gapinput|   mat:= InverseOp( List( mat, row -> List( row, Int ) ) );A
  @gapprompt|>A @gapinput|   if mat <> fail thenA
  @gapprompt|>A @gapinput|     mat:= ( mat mod modulus ) * one;A
  @gapprompt|>A @gapinput|   fi;A
  @gapprompt|>A @gapinput|   if not IsMatrix( mat ) thenA
  @gapprompt|>A @gapinput|     mat:= fail;A
  @gapprompt|>A @gapinput|   fi;A
  @gapprompt|>A @gapinput|   return mat;A
  @gapprompt|>A @gapinput|   end );A
\end{Verbatim}
 

 Additionally we install a method for finding a domain that contains the matrix
entries; this is used by some \textsf{GAP} library functions. 

 
\begin{Verbatim}[commandchars=!@|,fontsize=\small,frame=single,label=Example]
  !gapprompt@gap>| !gapinput@InstallMethod( DefaultFieldOfMatrixGroup,|
  !gapprompt@>| !gapinput@    "for a matrix group over a ring Z/nZ",|
  !gapprompt@>| !gapinput@    [ IsMatrixGroup and CategoryCollections( CategoryCollections(|
  !gapprompt@>| !gapinput@          CategoryCollections( IsMyZmodnZObj ) ) ) ],|
  !gapprompt@>| !gapinput@    G -> RingWithOneByGenerators([ One( Representative( G )[1][1] ) ]));|
\end{Verbatim}
 

 Now we can deal with matrix groups over residue class rings. 

 
\begin{Verbatim}[commandchars=!@|,fontsize=\small,frame=single,label=Example]
  !gapprompt@gap>| !gapinput@mat2:= IdentityMat( 2 ) * x;;|
  !gapprompt@gap>| !gapinput@mat2[2][1]:= x;;|
  !gapprompt@gap>| !gapinput@g:= Group( mat, mat2 );;|
  !gapprompt@gap>| !gapinput@Size( g );|
  3072
  !gapprompt@gap>| !gapinput@Factors( last );|
  [ 2, 2, 2, 2, 2, 2, 2, 2, 2, 2, 3 ]
  !gapprompt@gap>| !gapinput@syl3:= SylowSubgroup( g, 3 );;|
  !gapprompt@gap>| !gapinput@gens:= GeneratorsOfGroup( syl3 );|
  [ [ [ ( 1 mod 16 ), ( 7 mod 16 ) ], [ ( 11 mod 16 ), ( 14 mod 16 ) ]
       ] ]
  !gapprompt@gap>| !gapinput@Order( gens[1] );|
  3
\end{Verbatim}
 

 It should be noted that this way more involved methods for matrix groups may
not be available. For example, many questions about a finite matrix group can
be delegated to an isomorphic permutation group via a
so\texttt{\symbol{45}}called ``nice monomorphism''; this can be controlled by the filter \texttt{IsHandledByNiceMonomorphism} (\ref{IsHandledByNiceMonomorphism}). 

 By the way, also groups of (invertible) residue classes can be formed, but
this may be of minor interest. 
\begin{Verbatim}[commandchars=!@|,fontsize=\small,frame=single,label=Example]
  !gapprompt@gap>| !gapinput@g:= Group( x );;  Size( g );|
  #I  default `IsGeneratorsOfMagmaWithInverses' method returns `true' for
  [ ( 1 mod 16 ) ]
  1
  !gapprompt@gap>| !gapinput@g:= Group( 3*x );;  Size( g );|
  #I  default `IsGeneratorsOfMagmaWithInverses' method returns `true' for
  [ ( 3 mod 16 ) ]
  4
\end{Verbatim}
 

 (The messages above tell that \textsf{GAP} does not know a method for deciding whether the given elements are valid group
elements. We could add an appropriate \texttt{IsGeneratorsOfMagmaWithInverses} method if we would want.) 

 Having done enough for the elements, we may install some more methods for the
rings if we want to use them as arguments. These rings are finite, and there
are many generic methods that will work if they are able to compute the list
of elements of the ring, so we install a method for this. 

 
\begin{Verbatim}[commandchars=!@|,fontsize=\small,frame=single,label=Example]
  !gapprompt@gap>| !gapinput@InstallMethod( Enumerator,|
  !gapprompt@>| !gapinput@   "for full collection Z/nZ",|
  !gapprompt@>| !gapinput@   [ CategoryCollections( IsMyZmodnZObj ) and IsWholeFamily ],|
  !gapprompt@>| !gapinput@   function( R )|
  !gapprompt@>| !gapinput@   local F;|
  !gapprompt@>| !gapinput@   F:= ElementsFamily( FamilyObj(R) );|
  !gapprompt@>| !gapinput@   return List( [ 0 .. Size( R ) - 1 ], x -> MyZmodnZObj( F, x ) );|
  !gapprompt@>| !gapinput@   end );|
\end{Verbatim}
 

 Note that this method is applicable only to full rings ${\ensuremath{\mathbb Z}} / n{\ensuremath{\mathbb Z}}$, for proper subrings it would return a wrong result. Furthermore, it is not
required that the argument is a ring; in fact this method is applicable also
to the additive group formed by all elements in the family, provided that it
knows to contain the whole family. 

 Analogously, we install methods to compute the size, a random element, and the
units of full rings ${\ensuremath{\mathbb Z}} / n{\ensuremath{\mathbb Z}}$. 

 
\begin{Verbatim}[commandchars=@|A,fontsize=\small,frame=single,label=Example]
  @gapprompt|gap>A @gapinput|InstallMethod( Random,A
  @gapprompt|>A @gapinput|   "for full collection Z/nZ",A
  @gapprompt|>A @gapinput|   [ CategoryCollections( IsMyZmodnZObj ) and IsWholeFamily ],A
  @gapprompt|>A @gapinput|   R -> MyZmodnZObj( ElementsFamily( FamilyObj(R) ),A
  @gapprompt|>A @gapinput|                   Random( 0, Size( R ) - 1 ) ) );A
  @gapprompt|gap>A @gapinput|A
  @gapprompt|gap>A @gapinput|InstallMethod( Size,A
  @gapprompt|>A @gapinput|   "for full ring Z/nZ",A
  @gapprompt|>A @gapinput|   [ CategoryCollections( IsMyZmodnZObj ) and IsWholeFamily ],A
  @gapprompt|>A @gapinput|   R -> ElementsFamily( FamilyObj(R) )!.modulus );A
  @gapprompt|gap>A @gapinput|A
  @gapprompt|gap>A @gapinput|InstallMethod( Units,A
  @gapprompt|>A @gapinput|   "for full ring Z/nZ",A
  @gapprompt|>A @gapinput|   [     CategoryCollections( IsMyZmodnZObj )A
  @gapprompt|>A @gapinput|     and IsWholeFamily and IsRing ],A
  @gapprompt|>A @gapinput|   function( R )A
  @gapprompt|>A @gapinput|   local F;A
  @gapprompt|>A @gapinput|   F:= ElementsFamily( FamilyObj( R ) );A
  @gapprompt|>A @gapinput|   return List( PrimeResidues( Size(R) ), x -> MyZmodnZObj( F, x ) );A
  @gapprompt|>A @gapinput|   end );A
\end{Verbatim}
 

 The \texttt{Units} (\ref{Units}) method has the disadvantage that the result is returned as a list (in fact
this list is also strictly sorted). We could improve the implementation by
returning the units as a group; if we do not want to take the full list of
elements as generators, we can use the function \texttt{GeneratorsPrimeResidues} (\ref{GeneratorsPrimeResidues}). 

 
\begin{Verbatim}[commandchars=!@|,fontsize=\small,frame=single,label=Example]
  !gapprompt@gap>| !gapinput@InstallMethod( Units,|
  !gapprompt@>| !gapinput@   "for full ring Z/nZ",|
  !gapprompt@>| !gapinput@   [     CategoryCollections( IsMyZmodnZObj )|
  !gapprompt@>| !gapinput@     and IsWholeFamily and IsRing ],|
  !gapprompt@>| !gapinput@   function( R )|
  !gapprompt@>| !gapinput@   local G, gens;|
  !gapprompt@>| !gapinput@|
  !gapprompt@>| !gapinput@   gens:= GeneratorsPrimeResidues( Size( R ) ).generators;|
  !gapprompt@>| !gapinput@   if not IsEmpty( gens ) and gens[ 1 ] = 1 then|
  !gapprompt@>| !gapinput@     gens:= gens{ [ 2 .. Length( gens ) ] };|
  !gapprompt@>| !gapinput@   fi;|
  !gapprompt@>| !gapinput@   gens:= Flat( gens ) * One( R );|
  !gapprompt@>| !gapinput@   return GroupByGenerators( gens, One( R ) );|
  !gapprompt@>| !gapinput@   end );|
\end{Verbatim}
 

 Each ring ${\ensuremath{\mathbb Z}} / n{\ensuremath{\mathbb Z}}$ is finite, and we could install a method that returns \texttt{true} when \texttt{IsFinite} (\ref{IsFinite}) is called with ${\ensuremath{\mathbb Z}} / n{\ensuremath{\mathbb Z}}$ as argument. But we can do this more elegantly via installing a \emph{logical implication}. 

 
\begin{Verbatim}[commandchars=!@|,fontsize=\small,frame=single,label=Example]
  !gapprompt@gap>| !gapinput@InstallTrueMethod( IsFinite,|
  !gapprompt@>| !gapinput@   CategoryCollections( IsMyZmodnZObj ) and IsDomain );|
\end{Verbatim}
 

 In effect, every domain that consists of elements in \texttt{IsMyZmodnZObj} will automatically store that it is finite, even if \texttt{IsFinite} (\ref{IsFinite}) is not called for it. }

  
\section{\textcolor{Chapter }{Compatibility of Residue Class Rings with Prime Fields}}\label{Compatibility of Residue Class Rings with Prime Fields}
\logpage{[ 81, 4, 0 ]}
\hyperdef{L}{X83127B258512C436}{}
{
  The above implementation of residue classes and residue class rings has at
least two disadvantages. First, if $p$ is a prime then the ring ${\ensuremath{\mathbb Z}} / p{\ensuremath{\mathbb Z}}$ is in fact a field, but the return values of \texttt{MyZmodnZ} are never regarded as fields because they are not in the category \texttt{IsMagmaWithInversesIfNonzero} (\ref{IsMagmaWithInversesIfNonzero}). Second, and this makes the example really interesting, there are already
elements of finite prime fields implemented in \textsf{GAP}, and we may want to identify them with elements in ${\ensuremath{\mathbb Z}} / p{\ensuremath{\mathbb Z}}$. 

 To be more precise, elements of finite fields in \textsf{GAP} lie in the category \texttt{IsFFE} (\ref{IsFFE}), and there is already a representation, \texttt{IsInternalRep} (\ref{IsInternalRep}), of these elements via discrete logarithms. The aim of this section is to
make \texttt{IsMyModulusRep} an alternative representation of elements in finite prime fields. 

 Note that this is only one step towards the desired compatibility. Namely,
after having a second representation of elements in finite prime fields, we
may wish that the function \texttt{GF} (\ref{GF:for field size}) (which is the usual function to create finite fields in \textsf{GAP}) is able to return \texttt{MyZmodnZ( \mbox{\texttt{\mdseries\slshape p}} )} when \texttt{GF( \mbox{\texttt{\mdseries\slshape p}} )} is called for a prime \mbox{\texttt{\mdseries\slshape p}}. Moreover, then we have to decide about a default representation of elements
in \texttt{GF( \mbox{\texttt{\mdseries\slshape p}} )} for primes \mbox{\texttt{\mdseries\slshape p}} for which both representations are available. Of course we can force the new
representation by explicitly calling \texttt{MyZmodnZ} and \texttt{MyZmodnZObj} whenever we want, but it is not a priori clear in which situation which
representation is preferable. 

 The same questions will occur when we want to implement a new representation
for non\texttt{\symbol{45}}prime fields. The steps of this implementation will
be the same as described in this chapter, and we will have to achieve
compatibility with both the internal representation of elements in small
finite fields and the representation \texttt{IsMyModulusRep} of elements in arbitrary prime fields. 

 But let us now turn back to the task of this section. We first adjust the
setup of the declaration part of the previous section, and then repeat the
installations with suitable modifications. 

 (We should start a new \textsf{GAP} session for that, otherwise \textsf{GAP} will complain that the objects to be declared are already bound; additionally,
the methods installed above may be not compatible with the ones we want.) 

 
\begin{Verbatim}[commandchars=!@|,fontsize=\small,frame=single,label=Example]
  !gapprompt@gap>| !gapinput@DeclareCategory( "IsMyZmodnZObj", IsScalar );|
  !gapprompt@gap>| !gapinput@|
  !gapprompt@gap>| !gapinput@DeclareCategory( "IsMyZmodnZObjNonprime", IsMyZmodnZObj );|
  !gapprompt@gap>| !gapinput@|
  !gapprompt@gap>| !gapinput@DeclareSynonym( "IsMyZmodpZObj", IsMyZmodnZObj and IsFFE );|
  !gapprompt@gap>| !gapinput@|
  !gapprompt@gap>| !gapinput@DeclareRepresentation( "IsMyModulusRep", IsPositionalObjectRep, [ 1 ] );|
  !gapprompt@gap>| !gapinput@|
  !gapprompt@gap>| !gapinput@DeclareGlobalFunction( "MyZmodnZObj" );|
  !gapprompt@gap>| !gapinput@|
  !gapprompt@gap>| !gapinput@DeclareGlobalFunction( "MyZmodnZ" );|
\end{Verbatim}
 

 As in the previous section, all (newly introduced) elements of rings ${\ensuremath{\mathbb Z}} / n{\ensuremath{\mathbb Z}}$ lie in the category \texttt{IsMyZmodnZObj}. But now we introduce two subcategories, namely \texttt{IsMyZmodnZObjNonprime} for all elements in rings ${\ensuremath{\mathbb Z}} / n{\ensuremath{\mathbb Z}}$ where $n$ is not a prime, and \texttt{IsMyZmodpZObj} for elements in finite prime fields. All objects in the latter are
automatically known to lie in the category \texttt{IsFFE} (\ref{IsFFE}) of finite field elements. 

 It would be reasonable if also those internally represented elements in the
category \texttt{IsFFE} (\ref{IsFFE}) that do in fact lie in a prime field would also lie in the category \texttt{IsMyZmodnZObj} (and thus in fact in \texttt{IsMyZmodpZObj}). But this cannot be achieved because internally represented finite field
elements do in general not store whether they lie in a prime field. 

 As for the implementation part, again let us start with the definitions of \texttt{MyZmodnZObj} and \texttt{MyZmodnZ}. 

 
\begin{Verbatim}[commandchars=@|A,fontsize=\small,frame=single,label=Example]
  @gapprompt|gap>A @gapinput|InstallGlobalFunction( MyZmodnZObj, function( Fam, residue )A
  @gapprompt|>A @gapinput|   if IsFFEFamily( Fam ) thenA
  @gapprompt|>A @gapinput|     return Objectify( NewType( Fam, IsMyZmodpZObjA
  @gapprompt|>A @gapinput|                                 and IsMyModulusRep ),A
  @gapprompt|>A @gapinput|                   [ residue mod Characteristic( Fam ) ] );A
  @gapprompt|>A @gapinput|   elseA
  @gapprompt|>A @gapinput|     return Objectify( NewType( Fam, IsMyZmodnZObjNonprimeA
  @gapprompt|>A @gapinput|                                 and IsMyModulusRep ),A
  @gapprompt|>A @gapinput|                   [ residue mod Fam!.modulus ] );A
  @gapprompt|>A @gapinput|   fi;A
  @gapprompt|>A @gapinput|end );A
  
  @gapprompt|gap>A @gapinput|InstallGlobalFunction( MyZmodnZ, function( n )A
  @gapprompt|>A @gapinput|   local F, R;A
  @gapprompt|>A @gapinput|A
  @gapprompt|>A @gapinput|   if not ( IsInt( n ) and IsPosRat( n ) ) thenA
  @gapprompt|>A @gapinput|     Error( "<n> must be a positive integer" );A
  @gapprompt|>A @gapinput|   elif IsPrimeInt( n ) thenA
  @gapprompt|>A @gapinput|     # Construct the family of element objects of our field.A
  @gapprompt|>A @gapinput|     F:= FFEFamily( n );A
  @gapprompt|>A @gapinput|     # Make the domain.A
  @gapprompt|>A @gapinput|     R:= FieldOverItselfByGenerators( [ MyZmodnZObj( F, 1 ) ] );A
  @gapprompt|>A @gapinput|     SetIsPrimeField( R, true );A
  @gapprompt|>A @gapinput|   elseA
  @gapprompt|>A @gapinput|     # Construct the family of element objects of our ring.A
  @gapprompt|>A @gapinput|     F:= NewFamily( Concatenation( "MyZmod", String( n ), "Z" ),A
  @gapprompt|>A @gapinput|                    IsMyZmodnZObjNonprime );A
  @gapprompt|>A @gapinput|     # Install the data.A
  @gapprompt|>A @gapinput|     F!.modulus:= n;A
  @gapprompt|>A @gapinput|     # Make the domain.A
  @gapprompt|>A @gapinput|     R:= RingWithOneByGenerators( [ MyZmodnZObj( F, 1 ) ] );A
  @gapprompt|>A @gapinput|     SetIsWholeFamily( R, true );A
  @gapprompt|>A @gapinput|     SetName( R, Concatenation( "(Integers mod ",String(n),")" ) );A
  @gapprompt|>A @gapinput|   fi;A
  @gapprompt|>A @gapinput|A
  @gapprompt|>A @gapinput|   # Return the ring resp. field.A
  @gapprompt|>A @gapinput|   return R;A
  @gapprompt|>A @gapinput|end );A
\end{Verbatim}
 

 Note that the result of \texttt{MyZmodnZ} with a prime as argument is a field that does not contain the whole family of
its elements, since all finite field elements of a fixed characteristic lie in
the same family. Further note that we cannot expect a family of finite field
elements to have a component \texttt{modulus}, so we use \texttt{Characteristic} (\ref{Characteristic}) to get the modulus. Requiring that \texttt{Fam!.modulus} works also if \texttt{Fam} is a family of finite field elements would violate the rule that an extension
of \textsf{GAP} should not force changes in existing code, in this case code dealing with
families of finite field elements. 

 
\begin{Verbatim}[commandchars=@|A,fontsize=\small,frame=single,label=Example]
  @gapprompt|gap>A @gapinput|InstallMethod( PrintObj,A
  @gapprompt|>A @gapinput|   "for element in Z/nZ (ModulusRep)",A
  @gapprompt|>A @gapinput|   [ IsMyZmodnZObjNonprime and IsMyModulusRep ],A
  @gapprompt|>A @gapinput|   function( x )A
  @gapprompt|>A @gapinput|   Print( "( ", x![1], " mod ", FamilyObj(x)!.modulus, " )" );A
  @gapprompt|>A @gapinput|   end );A
  @gapprompt|gap>A @gapinput|A
  @gapprompt|gap>A @gapinput|InstallMethod( PrintObj,A
  @gapprompt|>A @gapinput|   "for element in Z/pZ (ModulusRep)",A
  @gapprompt|>A @gapinput|   [ IsMyZmodpZObj and IsMyModulusRep ],A
  @gapprompt|>A @gapinput|   function( x )A
  @gapprompt|>A @gapinput|   Print( "( ", x![1], " mod ", Characteristic(x), " )" );A
  @gapprompt|>A @gapinput|   end );A
  @gapprompt|gap>A @gapinput|A
  @gapprompt|gap>A @gapinput|InstallMethod( \=,A
  @gapprompt|>A @gapinput|   "for two elements in Z/nZ (ModulusRep)",A
  @gapprompt|>A @gapinput|   IsIdenticalObj,A
  @gapprompt|>A @gapinput|   [ IsMyZmodnZObj and IsMyModulusRep,A
  @gapprompt|>A @gapinput|     IsMyZmodnZObj and IsMyModulusRep ],A
  @gapprompt|>A @gapinput|   function( x, y ) return x![1] = y![1]; end );A
\end{Verbatim}
 

 The above method to check equality is independent of whether the arguments
have a prime or nonprime modulus, so we installed it for arguments in \texttt{IsMyZmodnZObj}. Now we install also methods to compare objects in \texttt{IsMyZmodpZObj} with the ``old'' finite field elements. 

 
\begin{Verbatim}[commandchars=@|A,fontsize=\small,frame=single,label=Example]
  @gapprompt|gap>A @gapinput|InstallMethod( \=,A
  @gapprompt|>A @gapinput|   "for element in Z/pZ (ModulusRep) and internal FFE",A
  @gapprompt|>A @gapinput|   IsIdenticalObj,A
  @gapprompt|>A @gapinput|   [ IsMyZmodpZObj and IsMyModulusRep, IsFFE and IsInternalRep ],A
  @gapprompt|>A @gapinput|   function( x, y )A
  @gapprompt|>A @gapinput|   return DegreeFFE( y ) = 1 and x![1] = IntFFE( y );A
  @gapprompt|>A @gapinput|   end );A
  @gapprompt|gap>A @gapinput|A
  @gapprompt|gap>A @gapinput|InstallMethod( \=,A
  @gapprompt|>A @gapinput|   "for internal FFE and element in Z/pZ (ModulusRep)",A
  @gapprompt|>A @gapinput|   IsIdenticalObj,A
  @gapprompt|>A @gapinput|   [ IsFFE and IsInternalRep, IsMyZmodpZObj and IsMyModulusRep ],A
  @gapprompt|>A @gapinput|   function( x, y )A
  @gapprompt|>A @gapinput|   return DegreeFFE( x ) = 1 and IntFFE( x ) = y![1];A
  @gapprompt|>A @gapinput|   end );A
\end{Verbatim}
 

 The situation with the operation \texttt{{\textless}} is more difficult. Of course we are free to define the comparison of objects
in \texttt{IsMyZmodnZObjNonprime}, but for the finite field elements, the comparison must be compatible with
the predefined comparison of the ``old'' finite field elements. The definition of the \texttt{{\textless}} comparison of internally represented finite field elements can be found in
Chapter{\nobreakspace}\ref{Finite Fields}. In situations where the documentation does not provide the required
information, one has to look it up in the \textsf{GAP} code; for example, the comparison in our case can be found in the appropriate
source code file of the \textsf{GAP} kernel. 

 
\begin{Verbatim}[commandchars=@|A,fontsize=\small,frame=single,label=Example]
  @gapprompt|gap>A @gapinput|InstallMethod( \<,A
  @gapprompt|>A @gapinput|   "for two elements in Z/nZ (ModulusRep, nonprime)",A
  @gapprompt|>A @gapinput|   IsIdenticalObj,A
  @gapprompt|>A @gapinput|   [ IsMyZmodnZObjNonprime and IsMyModulusRep,A
  @gapprompt|>A @gapinput|     IsMyZmodnZObjNonprime and IsMyModulusRep ],A
  @gapprompt|>A @gapinput|   function( x, y ) return x![1] < y![1]; end );A
  @gapprompt|gap>A @gapinput|A
  @gapprompt|gap>A @gapinput|InstallMethod( \<,A
  @gapprompt|>A @gapinput|   "for two elements in Z/pZ (ModulusRep)",A
  @gapprompt|>A @gapinput|   IsIdenticalObj,A
  @gapprompt|>A @gapinput|   [ IsMyZmodpZObj and IsMyModulusRep,A
  @gapprompt|>A @gapinput|     IsMyZmodpZObj and IsMyModulusRep ],A
  @gapprompt|>A @gapinput|   function( x, y )A
  @gapprompt|>A @gapinput|   local p, r;      # characteristic and primitive rootA
  @gapprompt|>A @gapinput|   if x![1] = 0 thenA
  @gapprompt|>A @gapinput|     return y![1] <> 0;A
  @gapprompt|>A @gapinput|   elif y![1] = 0 thenA
  @gapprompt|>A @gapinput|     return false;A
  @gapprompt|>A @gapinput|   elseA
  @gapprompt|>A @gapinput|     p:= Characteristic( x );A
  @gapprompt|>A @gapinput|     r:= PrimitiveRootMod( p );A
  @gapprompt|>A @gapinput|     return LogMod( x![1], r, p ) < LogMod( y![1], r, p );A
  @gapprompt|>A @gapinput|   fi;A
  @gapprompt|>A @gapinput|   end );A
  @gapprompt|gap>A @gapinput|A
  @gapprompt|gap>A @gapinput|InstallMethod( \<,A
  @gapprompt|>A @gapinput|   "for element in Z/pZ (ModulusRep) and internal FFE",A
  @gapprompt|>A @gapinput|   IsIdenticalObj,A
  @gapprompt|>A @gapinput|   [ IsMyZmodpZObj and IsMyModulusRep, IsFFE and IsInternalRep ],A
  @gapprompt|>A @gapinput|   function( x, y )A
  @gapprompt|>A @gapinput|   return x![1] * One( y ) < y;A
  @gapprompt|>A @gapinput|   end );A
  @gapprompt|gap>A @gapinput|A
  @gapprompt|gap>A @gapinput|InstallMethod( \<,A
  @gapprompt|>A @gapinput|   "for internal FFE and element in Z/pZ (ModulusRep)",A
  @gapprompt|>A @gapinput|   IsIdenticalObj,A
  @gapprompt|>A @gapinput|   [ IsFFE and IsInternalRep, IsMyZmodpZObj and IsMyModulusRep ],A
  @gapprompt|>A @gapinput|   function( x, y )A
  @gapprompt|>A @gapinput|   return x < y![1] * One( x );A
  @gapprompt|>A @gapinput|   end );A
\end{Verbatim}
 

 Now we install the same methods for the arithmetic operations \texttt{\texttt{\symbol{92}}+} (\ref{bSlash+}), \texttt{ZeroOp} (\ref{ZeroOp}), \texttt{AdditiveInverseOp} (\ref{AdditiveInverseOp}), \texttt{\texttt{\symbol{92}}\texttt{\symbol{45}}}, \texttt{\texttt{\symbol{92}}*} (\ref{bSlash*}), and \texttt{OneOp} (\ref{OneOp}) as in the previous section, without listing them below. Also the same \texttt{Int} (\ref{Int}) method is installed for objects in \texttt{IsMyZmodnZObj}. Note that it is compatible with the definition of \texttt{Int} (\ref{Int}) for finite field elements. And of course the same method for \texttt{\texttt{\symbol{92}}mod} (\ref{bSlashmod}) is installed. 

 We have to be careful, however, with the methods for \texttt{InverseOp} (\ref{InverseOp}), \texttt{\texttt{\symbol{92}}/} (\ref{bSlash/}), and \texttt{\texttt{\symbol{92}}\texttt{\symbol{94}}} (\ref{bSlash^}). These methods and the missing methods for arithmetic operations with one
argument in \texttt{IsMyModulusRep} and the other in \texttt{IsInternalRep} (\ref{IsInternalRep}) are given below. 

 
\begin{Verbatim}[commandchars=@|A,fontsize=\small,frame=single,label=Example]
  @gapprompt|gap>A @gapinput|InstallMethod( \+,A
  @gapprompt|>A @gapinput|   "for element in Z/pZ (ModulusRep) and internal FFE",A
  @gapprompt|>A @gapinput|   IsIdenticalObj,A
  @gapprompt|>A @gapinput|   [ IsMyZmodpZObj and IsMyModulusRep, IsFFE and IsInternalRep ],A
  @gapprompt|>A @gapinput|   function( x, y ) return x![1] + y; end );A
  @gapprompt|gap>A @gapinput|A
  @gapprompt|gap>A @gapinput|InstallMethod( \+,A
  @gapprompt|>A @gapinput|   "for internal FFE and element in Z/pZ (ModulusRep)",A
  @gapprompt|>A @gapinput|   IsIdenticalObj,A
  @gapprompt|>A @gapinput|   [ IsFFE and IsInternalRep, IsMyZmodpZObj and IsMyModulusRep ],A
  @gapprompt|>A @gapinput|   function( x, y ) return x + y![1]; end );A
  @gapprompt|gap>A @gapinput|A
  @gapprompt|gap>A @gapinput|InstallMethod( \*,A
  @gapprompt|>A @gapinput|   "for element in Z/pZ (ModulusRep) and internal FFE",A
  @gapprompt|>A @gapinput|   IsIdenticalObj,A
  @gapprompt|>A @gapinput|   [ IsMyZmodpZObj and IsMyModulusRep, IsFFE and IsInternalRep ],A
  @gapprompt|>A @gapinput|   function( x, y ) return x![1] * y; end );A
  @gapprompt|gap>A @gapinput|A
  @gapprompt|gap>A @gapinput|InstallMethod( \*,A
  @gapprompt|>A @gapinput|   "for internal FFE and element in Z/pZ (ModulusRep)",A
  @gapprompt|>A @gapinput|   IsIdenticalObj,A
  @gapprompt|>A @gapinput|   [ IsFFE and IsInternalRep, IsMyZmodpZObj and IsMyModulusRep ],A
  @gapprompt|>A @gapinput|   function( x, y ) return x * y![1]; end );A
  @gapprompt|gap>A @gapinput|A
  @gapprompt|gap>A @gapinput|InstallMethod( InverseOp,A
  @gapprompt|>A @gapinput|   "for element in Z/nZ (ModulusRep, nonprime)",A
  @gapprompt|>A @gapinput|   [ IsMyZmodnZObjNonprime and IsMyModulusRep ],A
  @gapprompt|>A @gapinput|   function( x )A
  @gapprompt|>A @gapinput|   local residue;A
  @gapprompt|>A @gapinput|   residue:= QuotientMod( 1, x![1], FamilyObj(x)!.modulus );A
  @gapprompt|>A @gapinput|   if residue <> fail thenA
  @gapprompt|>A @gapinput|     residue:= MyZmodnZObj( FamilyObj(x), residue );A
  @gapprompt|>A @gapinput|   fi;A
  @gapprompt|>A @gapinput|   return residue;A
  @gapprompt|>A @gapinput|   end );A
  @gapprompt|gap>A @gapinput|A
  @gapprompt|gap>A @gapinput|InstallMethod( InverseOp,A
  @gapprompt|>A @gapinput|   "for element in Z/pZ (ModulusRep)",A
  @gapprompt|>A @gapinput|   [ IsMyZmodpZObj and IsMyModulusRep ],A
  @gapprompt|>A @gapinput|   function( x )A
  @gapprompt|>A @gapinput|   local residue;A
  @gapprompt|>A @gapinput|   residue:= QuotientMod( 1, x![1], Characteristic( FamilyObj(x) ) );A
  @gapprompt|>A @gapinput|   if residue <> fail thenA
  @gapprompt|>A @gapinput|     residue:= MyZmodnZObj( FamilyObj(x), residue );A
  @gapprompt|>A @gapinput|   fi;A
  @gapprompt|>A @gapinput|   return residue;A
  @gapprompt|>A @gapinput|   end );A
\end{Verbatim}
 

 The operation \texttt{DegreeFFE} (\ref{DegreeFFE:for a FFE}) is defined for finite field elements, we need a method for objects in \texttt{IsMyZmodpZObj}. Note that we need not require \texttt{IsMyModulusRep} since no access to representation dependent data occurs. 

 
\begin{Verbatim}[commandchars=!@|,fontsize=\small,frame=single,label=Example]
  !gapprompt@gap>| !gapinput@InstallMethod( DegreeFFE,|
  !gapprompt@>| !gapinput@   "for element in Z/pZ",|
  !gapprompt@>| !gapinput@   [ IsMyZmodpZObj ],|
  !gapprompt@>| !gapinput@   z -> 1 );|
\end{Verbatim}
 

 The methods for \texttt{Enumerator} (\ref{Enumerator}), \texttt{Random} (\ref{Random:for a list or collection}), \texttt{Size} (\ref{Size}), and \texttt{Units} (\ref{Units}), that we had installed in the previous section had all assumed that their
argument contains the whole family of its elements. So these methods make
sense only for the nonprime case. For the prime case, there are already
methods for these operations with argument a field. 

 
\begin{Verbatim}[commandchars=@|A,fontsize=\small,frame=single,label=Example]
  @gapprompt|gap>A @gapinput|InstallMethod( Enumerator,A
  @gapprompt|>A @gapinput|   "for full ring Z/nZ",A
  @gapprompt|>A @gapinput|   [ CategoryCollections( IsMyZmodnZObjNonprime ) and IsWholeFamily ],A
  @gapprompt|>A @gapinput|   function( R )A
  @gapprompt|>A @gapinput|   local F;A
  @gapprompt|>A @gapinput|   F:= ElementsFamily( FamilyObj( R ) );A
  @gapprompt|>A @gapinput|   return List( [ 0 .. Size( R ) - 1 ], x -> MyZmodnZObj( F, x ) );A
  @gapprompt|>A @gapinput|   end );A
  @gapprompt|gap>A @gapinput|A
  @gapprompt|gap>A @gapinput|InstallMethod( Random,A
  @gapprompt|>A @gapinput|   "for full ring Z/nZ",A
  @gapprompt|>A @gapinput|   [ CategoryCollections( IsMyZmodnZObjNonprime ) and IsWholeFamily ],A
  @gapprompt|>A @gapinput|   R -> MyZmodnZObj( ElementsFamily( FamilyObj( R ) ),A
  @gapprompt|>A @gapinput|                   Random( 0, Size( R ) - 1 ) ) );A
  @gapprompt|gap>A @gapinput|A
  @gapprompt|gap>A @gapinput|InstallMethod( Size,A
  @gapprompt|>A @gapinput|   "for full ring Z/nZ",A
  @gapprompt|>A @gapinput|   [ CategoryCollections( IsMyZmodnZObjNonprime ) and IsWholeFamily ],A
  @gapprompt|>A @gapinput|   R -> ElementsFamily( FamilyObj( R ) )!.modulus );A
  @gapprompt|gap>A @gapinput|A
  @gapprompt|gap>A @gapinput|InstallMethod( Units,A
  @gapprompt|>A @gapinput|   "for full ring Z/nZ",A
  @gapprompt|>A @gapinput|   [     CategoryCollections( IsMyZmodnZObjNonprime )A
  @gapprompt|>A @gapinput|     and IsWholeFamily and IsRing ],A
  @gapprompt|>A @gapinput|   function( R )A
  @gapprompt|>A @gapinput|   local G, gens;A
  @gapprompt|>A @gapinput|A
  @gapprompt|>A @gapinput|   gens:= GeneratorsPrimeResidues( Size( R ) ).generators;A
  @gapprompt|>A @gapinput|   if not IsEmpty( gens ) and gens[ 1 ] = 1 thenA
  @gapprompt|>A @gapinput|     gens:= gens{ [ 2 .. Length( gens ) ] };A
  @gapprompt|>A @gapinput|   fi;A
  @gapprompt|>A @gapinput|   gens:= Flat( gens ) * One( R );A
  @gapprompt|>A @gapinput|   return GroupByGenerators( gens, One( R ) );A
  @gapprompt|>A @gapinput|   end );A
  @gapprompt|gap>A @gapinput|A
  @gapprompt|gap>A @gapinput|InstallTrueMethod( IsFinite,A
  @gapprompt|>A @gapinput|   CategoryCollections( IsMyZmodnZObjNonprime ) and IsDomain );A
\end{Verbatim}
 }

  
\section{\textcolor{Chapter }{Further Improvements in Implementing Residue Class Rings}}\label{Further Improvements in Implementing Residue Class Rings}
\logpage{[ 81, 5, 0 ]}
\hyperdef{L}{X81CA1C7087A815DE}{}
{
  There are of course many possibilities to improve the implementation. 

 With the setup as described above, subsequent calls \texttt{MyZmodnZ( \mbox{\texttt{\mdseries\slshape n}} )} with the same \mbox{\texttt{\mdseries\slshape n}} yield incompatible rings in the sense that elements of one ring cannot be
added to elements of an other one. The solution for this problem is to keep a
global list of all results of \texttt{MyZmodnZ} in the current \textsf{GAP} session, and to return the stored values whenever possible. Note that this
approach would admit \texttt{PrintObj} (\ref{PrintObj}) methods that produce \textsf{GAP} readable output. 

 One can improve the \texttt{Units} (\ref{Units}) method for the full ring in such a way that a group is returned and not only a
list of its elements; then the result of \texttt{Units} (\ref{Units}) can be used, e.{\nobreakspace}g., as input for the operation \texttt{SylowSubgroup} (\ref{SylowSubgroup}). 

 To make computations more efficient, one can install methods for \texttt{\texttt{\symbol{92}}\texttt{\symbol{45}}}, \texttt{\texttt{\symbol{92}}/} (\ref{bSlash/}), and \texttt{\texttt{\symbol{92}}\texttt{\symbol{94}}} (\ref{bSlash^}); one reason for doing so may be that this avoids the unnecessary construction
of the additive or multiplicative inverse, or of intermediate powers. 

 
\begin{Verbatim}[commandchars=!@|,fontsize=\small,frame=single,label=Example]
  InstallMethod( \-, "two elements in Z/nZ (ModulusRep)", ... );
  InstallMethod( \-, "Z/nZ-obj. (ModulusRep) and integer", ... );
  InstallMethod( \-, "integer and Z/nZ-obj. (ModulusRep)", ... );
  InstallMethod( \-, "Z/pZ-obj. (ModulusRep) and internal FFE", ... );
  InstallMethod( \-, "internal FFE and Z/pZ-obj. (ModulusRep)", ... );
  InstallMethod( \*, "Z/nZ-obj. (ModulusRep) and integer", ... );
  InstallMethod( \*, "integer and Z/nZ-obj. (ModulusRep)", ... );
  InstallMethod( \/, "two Z/nZ-objs. (ModulusRep, nonprime)", ... );
  InstallMethod( \/, "two Z/pZ-objs. (ModulusRep)", ... );
  InstallMethod( \/, "Z/nZ-obj. (ModulusRep) and integer", ... );
  InstallMethod( \/, "integer and Z/nZ-obj. (ModulusRep)", ... );
  InstallMethod( \/, "Z/pZ-obj. (ModulusRep) and internal FFE", ... );
  InstallMethod( \/, "internal FFE and Z/pZ-obj. (ModulusRep)", ... );
  InstallMethod( \^, "Z/nZ-obj. (ModulusRep, nonprime) & int.", ... );
  InstallMethod( \^, "Z/pZ-obj. (ModulusRep), and integer", ... );
\end{Verbatim}
 

 The call to \texttt{NewType} (\ref{NewType}) in \texttt{MyZmodnZObj} can be avoided by storing the required type, e.g., in the family. But note
that it is \emph{not} admissible to take the type of an existing object as first argument of \texttt{Objectify} (\ref{Objectify}). For example, suppose two objects in \texttt{IsMyZmodnZObj} shall be added. Then we must not use the type of one of the arguments in a
call of \texttt{Objectify} (\ref{Objectify}), because the argument may have knowledge that is not correct for the result
of the addition. One may think of the property \texttt{IsOne} (\ref{IsOne}) that may hold for both arguments but certainly not for their sum. 

 For comparing two objects in \texttt{IsMyZmodpZObj} via ``\texttt{{\textless}}'', we had to install a quite expensive method because of the compatibility with
the comparison of finite field elements that did already exist. In fact \textsf{GAP} supports finite fields with elements represented via discrete logarithms only
up to a given size. So in principle we have the freedom to define a cheaper
comparison via ``\texttt{{\textless}}'' for objects in \texttt{IsMyZmodpZObj} if the modulus is large enough. This is possible by introducing two categories \texttt{IsMyZmodpZObjSmall} and \texttt{IsMyZmodpZObjLarge}, which are subcategories of \texttt{IsMyZmodpZObj}, and to install different \texttt{\texttt{\symbol{92}}{\textless}} (\ref{bSlash<}) methods for pairs of objects in these categories. }

 }

          
\chapter{\textcolor{Chapter }{An Example {\textendash} Designing Arithmetic Operations}}\label{An Example -- Designing Arithmetic Operations}
\logpage{[ 82, 0, 0 ]}
\hyperdef{L}{X7E485C967A5778C9}{}
{
  In this chapter, we give a {\textendash}hopefully typical{\textendash} example
of extending \textsf{GAP} by new objects with prescribed arithmetic operations (for a simple approach
that may be useful to get started though does not permit to exploit all
potential features, see also \texttt{ArithmeticElementCreator} (\ref{ArithmeticElementCreator})).  
\section{\textcolor{Chapter }{New Arithmetic Operations vs. New Objects}}\label{New Arithmetic Operations vs. New Objects}
\logpage{[ 82, 1, 0 ]}
\hyperdef{L}{X7EA9422E7ACA7276}{}
{
  A usual procedure in mathematics is the definition of new operations for given
objects; here are a few typical examples. The Lie bracket defines an
interesting new multiplicative structure on a given (associative) algebra.
Forming a group ring can be viewed as defining a new addition for the elements
of the given group, and extending the multiplication to sums of group elements
in a natural way. Forming the exterior algebra of a given vector space can be
viewed as defining a new multiplication for the vectors in a natural way. 

 \textsf{GAP} does \emph{not} support such a procedure. The main reason for this is that in \textsf{GAP}, the multiplication in a group, a ring etc.{\nobreakspace}is always written
as \texttt{*}, and the addition in a vector space, a ring etc.{\nobreakspace} is always
written as \texttt{+}. Therefore it is not possible to define the Lie bracket as a ``second multiplication'' for the elements of a given algebra; in fact, the multiplication in Lie
algebras in \textsf{GAP} is denoted by \texttt{*}. Analogously, constructing the group ring as sketched above is impossible if
an addition is already defined for the elements; note the difference between
the usual addition of matrices and the addition in the group ring of a matrix
group! (See Chapter{\nobreakspace}\ref{Magma Rings} for an example.) Similarly, there is already a multiplication defined for row
vectors (yielding the standard scalar product), hence these vectors cannot be
regarded as elements of the exterior algebra of the space. 

 In situations such as the ones mentioned above, \textsf{GAP}'s way to deal with the structures in question is the following. Instead of
defining \emph{new} operations for the \emph{given} objects, \emph{new} objects are created to which the \emph{given} arithmetic operations \texttt{*} and \texttt{+} are then made applicable. 

 With this construction, matrix Lie algebras consist of matrices that are
different from the matrices with associative multiplication; technically, the
type of a matrix determines how it is multiplied with other matrices
(see{\nobreakspace}\texttt{IsMatrix} (\ref{IsMatrix})). A matrix with the Lie bracket as its multiplication can be created with the
function \texttt{LieObject} (\ref{LieObject}) from a matrix with the usual associative multiplication. 

 Group rings (more general: magma rings, see Chapter{\nobreakspace}\ref{Magma Rings}) can be constructed with \texttt{FreeMagmaRing} (\ref{FreeMagmaRing}) from a coefficient ring and a group. The elements of the group are not
contained in such a group ring, one has to use an embedding map for creating a
group ring element that corresponds to a given group element. 

 It should be noted that the \textsf{GAP} approach to the construction of Lie algebras from associative algebras is
generic in the sense that all objects in the filter \texttt{IsLieObject} (\ref{IsLieObject}) use the same methods for their addition, multiplication etc., by delegating to
the ``underlying'' objects of the associative algebra, no matter what these objects actually are.
Analogously, also the construction of group rings is generic. }

  
\section{\textcolor{Chapter }{Designing new Multiplicative Objects}}\label{Designing new Multiplicative Objects}
\logpage{[ 82, 2, 0 ]}
\hyperdef{L}{X7BE9D84482B421F9}{}
{
  The goal of this section is to implement objects with a prescribed
multiplication. Let us assume that we are given a field $F$, and that we want to define a new multiplication $*$ on $F$ that is given by $a * b = a b - a - b + 2$; here $a b$ denotes the ordinary product in $F$. 

 By the discussion in Section{\nobreakspace}\ref{New Arithmetic Operations vs. New Objects}, we know that we cannot define a new multiplication on $F$ itself but have to create new objects. 

 We want to distinguish these new objects from all other \textsf{GAP} objects, in order to describe for example the situation that two of our
objects shall be multiplied. This distinction is made via the \emph{type} of the objects. More precisely, we declare a new \emph{filter}, a function that will return \texttt{true} for our new objects, and \texttt{false} for all other \textsf{GAP} objects. This can be done by calling \texttt{DeclareFilter} (\ref{DeclareFilter}), but since our objects will know about the value already when they are
constructed, the filter can be created with \texttt{DeclareCategory} (\ref{DeclareCategory}) or \texttt{NewCategory} (\ref{NewCategory}). 

 
\begin{Verbatim}[commandchars=!@|,fontsize=\small,frame=single,label=Example]
  DeclareCategory( "IsMyObject", IsObject );
\end{Verbatim}
 

 The idea is that the new multiplication will be installed only for objects
that ``lie in the category \texttt{IsMyObject}''. 

 The next question is what internal data our new objects store, and how they
are accessed. The easiest solution is to store the ``underlying'' object from the field $F$. \textsf{GAP} provides two general possibilities how to store this, namely
record\texttt{\symbol{45}}like and list\texttt{\symbol{45}}like structures
(for examples, see{\nobreakspace}\ref{Component Objects} and{\nobreakspace}\ref{Positional Objects}). We decide to store the data in a list\texttt{\symbol{45}}like structure, at
position 1. This \emph{representation} is declared as follows. 

 
\begin{Verbatim}[commandchars=!@|,fontsize=\small,frame=single,label=Example]
  DeclareRepresentation( "IsMyObjectListRep", IsPositionalObjectRep, [ 1 ] );
\end{Verbatim}
 

 Of course we can argue that this declaration is superfluous because \emph{all} objects in the category \texttt{IsMyObject} will be represented this way; it is possible to proceed like that, but often
(in more complicated situations) it turns out to be useful that several
representations are available for ``the same element''. 

 For creating the type of our objects, we need to specify to which \emph{family} (see{\nobreakspace}\ref{Families}) the objects shall belong. For the moment, we need not say anything about
relations to other \textsf{GAP} objects, thus the only requirement is that all new objects lie in the \emph{same} family; therefore we create a \emph{new} family. Also we are not interested in properties that some of our objects have
and others do not have, thus we need only one type, and store it in a global
variable. 

 
\begin{Verbatim}[commandchars=!@|,fontsize=\small,frame=single,label=Example]
  MyType:= NewType( NewFamily( "MyFamily" ),
                    IsMyObject and IsMyObjectListRep );
\end{Verbatim}
 

 The next step is to write a function that creates a new object. It may look as
follows. 
\begin{Verbatim}[commandchars=!@|,fontsize=\small,frame=single,label=Example]
  MyObject:= val -> Objectify( MyType, [ Immutable( val ) ] );
\end{Verbatim}
 Note that we store an \emph{immutable copy} of the argument in the returned object; without doing so, for example if the
argument would be a mutable matrix then the corresponding new object would be
changed whenever the matrix is changed (see{\nobreakspace}\ref{Mutability and Copyability} for more details about mutability). 

 Having entered the above \textsf{GAP} code, we can create some of our objects. 
\begin{Verbatim}[commandchars=@|A,fontsize=\small,frame=single,label=Example]
  @gapprompt|gap>A @gapinput|a:= MyObject( 3 );  b:= MyObject( 5 );A
  <object>
  <object>
  @gapprompt|gap>A @gapinput|a![1];  b![1];A
  3
  5
\end{Verbatim}
 But clearly a lot is missing. Besides the fact that the desired multiplication
is not yet installed, we see that also the way how the objects are printed is
not satisfactory. 

 Let us improve the latter first. There are two \textsf{GAP} functions \texttt{View} (\ref{View}) and \texttt{Print} (\ref{Print}) for showing objects on the screen. \texttt{View} (\ref{View}) is thought to show a short and human readable form of the object, and \texttt{Print} (\ref{Print}) is thought to show a not necessarily short form that is \textsf{GAP} readable whenever this makes sense. We decide to show \texttt{a} as \texttt{\mbox{\texttt{\mdseries\slshape 3}}} by \texttt{View} (\ref{View}), and to show the construction \texttt{MyObject( 3 )} by \texttt{Print} (\ref{Print}); the methods are installed for the underlying operations \texttt{ViewObj} (\ref{ViewObj}) and \texttt{PrintObj} (\ref{PrintObj}). 
\begin{Verbatim}[commandchars=@|A,fontsize=\small,frame=single,label=Example]
  InstallMethod( ViewObj,
      "for object in `IsMyObject'",
      [ IsMyObject and IsMyObjectListRep ],
      function( obj )
      Print( "<", obj![1], ">" );
      end );
  
  InstallMethod( PrintObj,
      "for object in `IsMyObject'",
      [ IsMyObject and IsMyObjectListRep ],
      function( obj )
      Print( "MyObject( ", obj![1], " )" );
      end );
\end{Verbatim}
 

 This is the result of the above installations. 
\begin{Verbatim}[commandchars=!@|,fontsize=\small,frame=single,label=Example]
  !gapprompt@gap>| !gapinput@a; Print( a, "\n" );|
  <3>
  MyObject( 3 )
\end{Verbatim}
 

 And now we try to install the multiplication. 

 
\begin{Verbatim}[commandchars=@|A,fontsize=\small,frame=single,label=Example]
  InstallMethod( \*,
      "for two objects in `IsMyObject'",
      [ IsMyObject and IsMyObjectListRep,
        IsMyObject and IsMyObjectListRep ],
      function( a, b )
      return MyObject( a![1] * b![1] - a![1] - b![1] + 2 );
      end );
\end{Verbatim}
 

 When we enter the above code, \textsf{GAP} runs into an error. This is due to the fact that the operation \texttt{\texttt{\symbol{92}}*} (\ref{bSlash*}) is declared for two arguments that lie in the category \texttt{IsMultiplicativeElement} (\ref{IsMultiplicativeElement}). One could circumvent the check whether the method matches the declaration of
the operation, by calling \texttt{InstallOtherMethod} (\ref{InstallOtherMethod}) instead of \texttt{InstallMethod} (\ref{InstallMethod}). But it would make sense if our objects would lie in \texttt{IsMultiplicativeElement} (\ref{IsMultiplicativeElement}), for example because some generic methods for objects with multiplication
would be available then, such as powering by positive integers via repeated
squaring. So we want that \texttt{IsMyObject} implies \texttt{IsMultiplicativeElement} (\ref{IsMultiplicativeElement}). The easiest way to achieve such implications is to use the implied filter as
second argument of the \texttt{DeclareCategory} (\ref{DeclareCategory}) call; but since we do not want to start anew, we can also install the
implication afterwards. 

 
\begin{Verbatim}[commandchars=!@|,fontsize=\small,frame=single,label=Example]
  InstallTrueMethod( IsMultiplicativeElement, IsMyObject );
\end{Verbatim}
 

 Afterwards, installing the multiplication works without problems. Note that \texttt{MyType} and therefore also \texttt{a} and \texttt{b} are \emph{not} affected by this implication, so we construct them anew. 

 
\begin{Verbatim}[commandchars=!@|,fontsize=\small,frame=single,label=Example]
  !gapprompt@gap>| !gapinput@MyType:= NewType( NewFamily( "MyFamily" ),|
  !gapprompt@>| !gapinput@                     IsMyObject and IsMyObjectListRep );;|
  !gapprompt@gap>| !gapinput@a:= MyObject( 3 );;  b:= MyObject( 5 );;|
  !gapprompt@gap>| !gapinput@a*b;  a^27;|
  <9>
  <134217729>
\end{Verbatim}
 

 Powering the new objects by negative integers is not possible yet, because \textsf{GAP} does not know how to compute the inverse of an element $a$, say, which is defined as the unique element $a'$ such that both $a a'$ and $a' a$ are ``the unique multiplicative neutral element that belongs to $a$''. 

 And also this neutral element, if it exists, cannot be computed by \textsf{GAP} in our current situation. It does, however, make sense to ask for the
multiplicative neutral element of a given magma, and for inverses of elements
in the magma. 

 But before we can form domains of our objects, we must define when two objects
are regarded as equal; note that this is necessary in order to decide about
the uniqueness of neutral and inverse elements. In our situation, equality is
defined in the obvious way. For being able to form sets of our objects, also
an ordering via \texttt{\texttt{\symbol{92}}{\textless}} (\ref{bSlash<}) is defined for them. 

 
\begin{Verbatim}[commandchars=@|A,fontsize=\small,frame=single,label=Example]
  InstallMethod( \=,
      "for two objects in `IsMyObject'",
      [ IsMyObject and IsMyObjectListRep,
        IsMyObject and IsMyObjectListRep ],
      function( a, b )
      return a![1] = b![1];
      end );
  
  InstallMethod( \<,
      "for two objects in `IsMyObject'",
      [ IsMyObject and IsMyObjectListRep,
        IsMyObject and IsMyObjectListRep ],
      function( a, b )
      return a![1] < b![1];
      end );
\end{Verbatim}
 

 Let us look at an example. We start with finite field elements because then
the domains are finite, hence the generic methods for such domains will have a
chance to succeed. 

 
\begin{Verbatim}[commandchars=!@|,fontsize=\small,frame=single,label=Example]
  !gapprompt@gap>| !gapinput@a:= MyObject( Z(7) );|
  <Z(7)>
  !gapprompt@gap>| !gapinput@m:= Magma( a );|
  <magma with 1 generator>
  !gapprompt@gap>| !gapinput@e:= MultiplicativeNeutralElement( m );|
  <Z(7)^2>
  !gapprompt@gap>| !gapinput@elms:= AsList( m );|
  [ <Z(7)>, <Z(7)^2>, <Z(7)^5> ]
  !gapprompt@gap>| !gapinput@ForAll( elms, x -> ForAny( elms, y -> x*y = e and y*x = e ) );|
  true
  !gapprompt@gap>| !gapinput@List( elms, x -> First( elms, y -> x*y = e and y*x = e ) );|
  [ <Z(7)^5>, <Z(7)^2>, <Z(7)> ]
\end{Verbatim}
 

 So a multiplicative neutral element exists, in fact all elements in the magma \texttt{m} are invertible. But what about the following. 

 
\begin{Verbatim}[commandchars=!@|,fontsize=\small,frame=single,label=Example]
  !gapprompt@gap>| !gapinput@b:= MyObject( Z(7)^0 );  m:= Magma( a, b );|
  <Z(7)^0>
  <magma with 2 generators>
  !gapprompt@gap>| !gapinput@elms:= AsList( m );|
  [ <Z(7)^0>, <Z(7)>, <Z(7)^2>, <Z(7)^5> ]
  !gapprompt@gap>| !gapinput@e:= MultiplicativeNeutralElement( m );|
  <Z(7)^2>
  !gapprompt@gap>| !gapinput@ForAll( elms, x -> ForAny( elms, y -> x*y = e and y*x = e ) );|
  false
  !gapprompt@gap>| !gapinput@List( elms, x -> b * x );|
  [ <Z(7)^0>, <Z(7)^0>, <Z(7)^0>, <Z(7)^0> ]
\end{Verbatim}
 

 Here we found a multiplicative neutral element, but the element \texttt{b} does not have an inverse. If an addition would be defined for our elements
then we would say that \texttt{b} behaves like a zero element. 

 When we started to implement the new objects, we said that we wanted to define
the new multiplication for elements of a given field $F$. In principle, the current implementation would admit also something like \texttt{MyObject( 2 ) * MyObject( Z(7) )}. But if we decide that our initial assumption holds, we may define the
identity and the inverse of the object \texttt{{\textless}a{\textgreater}} as \texttt{{\textless}2*e{\textgreater}} and \texttt{{\textless}a/(a\texttt{\symbol{45}}e){\textgreater}}, respectively, where \texttt{e} is the identity element in $F$ and \texttt{/} denotes the division in $F$; note that the element \texttt{{\textless}e{\textgreater}} is not invertible, and that the above definitions are determined by the
multiplication defined for our objects. Further note that after the
installations shown below, also \texttt{One( MyObject( 1 ) )} is defined. 

 (For technical reasons, we do not install the intended methods for the
attributes \texttt{One} (\ref{One}) and \texttt{Inverse} (\ref{Inverse}) but for the operations \texttt{OneOp} (\ref{OneOp}) and \texttt{InverseOp} (\ref{InverseOp}). This is because for certain kinds of objects {\textendash}mainly
matrices{\textendash} one wants to support a method to compute a \emph{mutable} identity or inverse, and the attribute needs only a method that takes this
object, makes it immutable, and then returns this object. As stated above, we
only want to deal with immutable objects, so this distinction is not really
interesting for us.) 

 A more interesting point to note is that we should mark our objects as likely
to be invertible, since we add the possibility to invert them. Again, this
could have been part of the declaration of \texttt{IsMyObject}, but we may also formulate an implication for the existing category. 

 
\begin{Verbatim}[commandchars=@|A,fontsize=\small,frame=single,label=Example]
  InstallTrueMethod( IsMultiplicativeElementWithInverse, IsMyObject );
  
  InstallMethod( OneOp,
      "for an object in `IsMyObject'",
      [ IsMyObject and IsMyObjectListRep ],
      a -> MyObject( 2 * One( a![1] ) ) );
  
  InstallMethod( InverseOp,
      "for an object in `IsMyObject'",
      [ IsMyObject and IsMyObjectListRep ],
      a -> MyObject( a![1] / ( a![1] - One( a![1] ) ) ) );
\end{Verbatim}
 Now we can form groups of our (nonzero) elements. 
\begin{Verbatim}[commandchars=!@|,fontsize=\small,frame=single,label=Example]
  !gapprompt@gap>| !gapinput@MyType:= NewType( NewFamily( "MyFamily" ),|
  !gapprompt@>| !gapinput@                  IsMyObject and IsMyObjectListRep );;|
  gap>
  !gapprompt@gap>| !gapinput@a:= MyObject( Z(7) );|
  <Z(7)>
  !gapprompt@gap>| !gapinput@b:= MyObject( 0*Z(7) );  g:= Group( a, b );|
  <0*Z(7)>
  <group with 2 generators>
  !gapprompt@gap>| !gapinput@Size( g );|
  6
\end{Verbatim}
 

 We are completely free to define an \emph{addition} for our elements, a natural one is given by \texttt{{\textless}a{\textgreater} + {\textless}b{\textgreater} =
{\textless}a+b\texttt{\symbol{45}}1{\textgreater}}. As we did for the multiplication, we first change \texttt{IsMyObject} such that the additive structure is also known. 
\begin{Verbatim}[commandchars=!@|,fontsize=\small,frame=single,label=Example]
  InstallTrueMethod( IsAdditiveElementWithInverse, IsMyObject );
\end{Verbatim}
 Next we install the methods for the addition, and those to compute the
additive neutral element and the additive inverse. 

 
\begin{Verbatim}[commandchars=@|B,fontsize=\small,frame=single,label=Example]
  InstallMethod( \+,
      "for two objects in `IsMyObject'",
      [ IsMyObject and IsMyObjectListRep,
        IsMyObject and IsMyObjectListRep ],
      function( a, b )
      return MyObject( a![1] + b![1] - 1 );
      end );
  
  InstallMethod( ZeroOp,
      "for an object in `IsMyObject'",
      [ IsMyObject and IsMyObjectListRep ],
      a -> MyObject( One( a![1] ) ) );
  
  InstallMethod( AdditiveInverseOp,
      "for an object in `IsMyObject'",
      [ IsMyObject and IsMyObjectListRep ],
      a -> MyObject( a![1] / ( a![1] - One( a![1] ) ) ) );
\end{Verbatim}
 

 Let us try whether the addition works. 

 
\begin{Verbatim}[commandchars=!@|,fontsize=\small,frame=single,label=Example]
  !gapprompt@gap>| !gapinput@MyType:= NewType( NewFamily( "MyFamily" ),|
  !gapprompt@>| !gapinput@                  IsMyObject and IsMyObjectListRep );;|
  !gapprompt@gap>| !gapinput@a:= MyObject( Z(7) );;  b:= MyObject( 0*Z(7) );;|
  !gapprompt@gap>| !gapinput@m:= AdditiveMagma( a, b );|
  <additive magma with 2 generators>
  !gapprompt@gap>| !gapinput@Size( m );|
  7
\end{Verbatim}
 

 Similar as installing a multiplication automatically makes powering by
integers available, multiplication with integers becomes available with the
addition. 

 
\begin{Verbatim}[commandchars=!@|,fontsize=\small,frame=single,label=Example]
  !gapprompt@gap>| !gapinput@2 * a;|
  <Z(7)^5>
  !gapprompt@gap>| !gapinput@a+a;|
  <Z(7)^5>
  !gapprompt@gap>| !gapinput@MyObject( 2*Z(7)^0 ) * a;|
  <Z(7)>
\end{Verbatim}
 

 In particular we see that this multiplication does \emph{not} coincide with the multiplication of two of our objects, that is, an integer \emph{cannot} be used as a shorthand for one of the new objects in a multiplication. 

 (It should be possible to create a \emph{field} with the new multiplication and addition. Currently this fails, due to missing
methods for computing several kinds of generators from field generators, for
computing the characteristic in the case that the family does not know this in
advance, for checking with \texttt{AsField} (\ref{AsField}) whether a domain is in fact a field, for computing the closure as a field.) 

 It should be emphasized that the mechanism described above may be not suitable
for the situation that one wants to consider many different multiplications ``on the same set of objects'', since the installation of a new multiplication requires the declaration of
at least one new filter and the installation of several methods. But the
design of \textsf{GAP} is not suitable for such dynamic method installations. 

 Turning this argument the other way round, the implementation of the new
arithmetics defined by the above multiplication and addition is available for
any field $F$, one need not repeat it for each field one is interested in. 

 Similar to the above situation, the construction of a magma ring $RM$ from a coefficient ring $R$ and a magma $M$ is implemented only once, since the definition of the arithmetic operations
depends only on the given multiplication of $M$ and not on $M$ itself. So the addition is not implemented for the elements in $M$ or {\textendash}more precisely{\textendash} for an isomorphic copy. In some
sense, the addition is installed ``for the multiplication'', and as mentioned in Section{\nobreakspace}\ref{New Arithmetic Operations vs. New Objects}, there is only one multiplication \texttt{\texttt{\symbol{92}}*} (\ref{bSlash*}) in \textsf{GAP}. }

 }

          
\chapter{\textcolor{Chapter }{Library Files}}\label{Library Files}
\logpage{[ 83, 0, 0 ]}
\hyperdef{L}{X848C952A87FB36E2}{}
{
  This chapter describes some of the conventions used in the \textsf{GAP} library files. These conventions are intended as a help on how to read library
files and how to find information in them. So everybody is recommended to
follow these conventions, although they do not prescribe a compulsory
programming style {\textendash}\textsf{GAP} itself will not bother with the formatting of files. 

 Filenames have traditionally \textsf{GAP} adhered to the 8+3 convention (to make it possible to use the same filenames
even on a MS\texttt{\symbol{45}}DOS file system) and been in lower case
(systems that do not recognize lower case in file names will convert them
automatically to upper case). It is no longer so important to adhere to these
conventions, but at the very least filenames should adhere to a 16+5
convention, and be distinct even after identifying upper and lower case.
Directory names of packages, however, \emph{must} be in lower case (the \texttt{LoadPackage} (\ref{LoadPackage}) command assumes this).  
\section{\textcolor{Chapter }{File Types}}\label{File Types}
\logpage{[ 83, 1, 0 ]}
\hyperdef{L}{X7FF5DC397C79392C}{}
{
  The \textsf{GAP} library consists of the following types of files, distinguished by their
suffixes: 

 
\begin{description}
\item[{\texttt{.g} }]  Files which contain parts of the ``inner workings'' of \textsf{GAP}. These files usually do not contain mathematical functionality, except for
providing links to kernel functions. 
\item[{\texttt{.gd} }]  Declaration files. These files contain declarations of all categories,
attributes, operations, and global functions. These files also contain the
operation definitions in comments. 
\item[{\texttt{.gi} }]  Implementation files. These files contain all installations of methods and
global functions. Usually declarations of representations are also considered
to be part of the implementation and are therefore found in the \texttt{.gi} files. 
\end{description}
 

 As a rule of thumb, all \texttt{.gd} files are read in before the \texttt{.gi} files are read. Therefore a \texttt{.gi} file usually may use any operation or global function (it has been declared
before), and no care has to be taken towards the order in which the \texttt{.gi} files are read. }

  
\section{\textcolor{Chapter }{Finding Implementations in the Library}}\label{Finding Implementations in the Library}
\logpage{[ 83, 2, 0 ]}
\hyperdef{L}{X845CCBE082CDF4BB}{}
{
  For a concretely given function, you can use \texttt{FilenameFunc} (\ref{FilenameFunc}) and \texttt{StartlineFunc} (\ref{StartlineFunc}) for finding the file where this function is defined, and the line in this file
where the definition of this function starts. This does not work for arbitrary
functions, see Section \texttt{FilenameFunc} (\ref{FilenameFunc}) for the restrictions. 

 If you are interested in getting the function which implements a method for
specific arguments, you can use \ref{ApplicableMethod}. If \texttt{FilenameFunc} (\ref{FilenameFunc}) does not work for this method then setting the print level of \ref{ApplicableMethod} higher will give you the installation string for this method, which can be
used for searching in library files. 

 To find the occurrence of functions, methods, function names, and installation
strings in the library, one can use the \texttt{grep} tool under UNIX. To find a function, search for the function name in the \texttt{.gd} files; as global variables are usually declared only once, only few files will
show up. The function installation is likely to occur in the corresponding \texttt{.gi} file. 

 To find a method from the known operation name and the installation string,
search for the string ``\texttt{Method(}'' (this catches both \texttt{InstallMethod} (\ref{InstallMethod}) and \texttt{InstallOtherMethod} (\ref{InstallOtherMethod})) and the installation string or the operation name. 

 The following tools from the \textsf{GAP} package \textsf{Browse} can be used for accessing the code of functions. 

 
\begin{itemize}
\item  \texttt{BrowseGapMethods} (\textbf{Browse: BrowseGapMethods see BrowseGapData}) shows an overview of \textsf{GAP}'s operations and methods, and allows one to navigate through the files that
contain the implementations of the methods, using a pager. 
\item  \texttt{BrowseProfile} (\textbf{Browse: BrowseProfile}) shows profiling results (similar to \texttt{DisplayProfile} (\ref{DisplayProfile})) and allows one to navigate through the files that contain the
implementations of the functions that were actually used, using a pager. 
\end{itemize}
 }

  
\section{\textcolor{Chapter }{Undocumented Variables}}\label{Undocumented Variables}
\logpage{[ 83, 3, 0 ]}
\hyperdef{L}{X801428EB86E7113C}{}
{
  For several global variables in \textsf{GAP}, no information is available via the help system (see Section{\nobreakspace} (\textbf{Tutorial: Help}), for a quick overview of the help system, or Chapter{\nobreakspace}\ref{The Help System}, for details). There are various reasons for ``hiding'' a variable from the user; namely, the variable may be regarded as of minor
importance (for example, it may be a function called by documented \textsf{GAP} functions that first compute many input parameters for the undocumented
function), or it belongs to a part of \textsf{GAP} that is still experimental in the sense that the meaning of the variable has
not yet been fixed or even that it is not clear whether the variable will
vanish in a more developed version. 

 As a consequence, it is dangerous to use undocumented variables because they
are not guaranteed to exist or to behave the same in future versions of \textsf{GAP}. 

 Conversely, for \emph{documented} variables, the definitions in the \textsf{GAP} manual can be relied on for future \textsf{GAP} versions (unless they turn out to be erroneous); if the \textsf{GAP} developers find that some piece of minor, but documented functionality is an
insurmountable obstacle to important developments, they may make the smallest
possible incompatible change to the functionality at the time of a major
release. However, in any such case it will be announced clearly in the \textsf{GAP} Forum what has been changed and why. 

 So on the one hand, the developers of \textsf{GAP} want to keep the freedom of changing undocumented \textsf{GAP} code. On the other hand, users may be interested in using undocumented
variables. 

 In this case, whenever you write \textsf{GAP} code involving undocumented variables, and want to make sure that this code
will work in future versions of \textsf{GAP}, you may ask at \href{mailto://support@gap-system.org} {\texttt{support@gap\texttt{\symbol{45}}system.org}} for documentation about the variables in question. The \textsf{GAP} developers then decide whether these variables shall be documented or not, and
if yes, what the definitions shall be. 

 In the former case, the new documentation is added to the \textsf{GAP} manual, this means that from then on, this definition is protected against
changes. 

 In the latter case (which may occur for example if the variables in question
are still experimental), you may add the current values of these variables to
your private code if you want to be sure that nothing will be broken later due
to changes in \textsf{GAP}. }

 }

            
\chapter{\textcolor{Chapter }{Interface to the GAP Help System}}\label{Interface to the GAP Help System}
\logpage{[ 84, 0, 0 ]}
\hyperdef{L}{X79A6CE6C86A976AE}{}
{
  In this chapter we describe which information the help system needs about a
manual book and how to tell it this information. The code which implements
this interface can be found in \texttt{lib/helpbase.gi}. 

 If you are intending to use a documentation format that is already used by
some other help book you probably don't need to know anything from this
chapter. However, if you want to create a new format and make it available to \textsf{GAP} then hopefully you will find the necessary information here. 

 The basic idea of the help system is as follows: One tells \textsf{GAP} a directory which contains a file \texttt{manual.six}, see{\nobreakspace}\ref{Installing and Removing a Help Book}. When the \textsf{GAP} help is asked something about this book it reads in some basic information
from the file \texttt{manual.six}: strings like section headers, function names, and index entries to be
searched by the online help; information about the available formats of this
book like text, html, dvi, and pdf; the actual files containing the
documentation, corresponding section numbers, and page numbers: and so on.
See{\nobreakspace}\ref{The manual.six File} for a description of the format of the \texttt{manual.six} file. 

 It turns out that there is almost no restriction on the format of the \texttt{manual.six} file, except that it must provide a string, say \texttt{"myownformat"} which identifies the format of the help book. Then the basic actions on a help
book are delegated by the help system to handler functions stored in a record \texttt{HELP{\textunderscore}BOOK{\textunderscore}HANDLER.myownformat}. See{\nobreakspace}\ref{The Help Book Handler} for information which functions must be provided by the handler and what they
are supposed to do. The main work to teach \textsf{GAP} to use a new document format is to write these handler functions and to
produce an appropriate \texttt{manual.six} file.  
\section{\textcolor{Chapter }{Installing and Removing a Help Book}}\label{Installing and Removing a Help Book}
\logpage{[ 84, 1, 0 ]}
\hyperdef{L}{X7AFEAB6B84387635}{}
{
  

\subsection{\textcolor{Chapter }{HELP{\textunderscore}ADD{\textunderscore}BOOK}}
\logpage{[ 84, 1, 1 ]}\nobreak
\hyperdef{L}{X7CD0B8507A3D231D}{}
{\noindent\textcolor{FuncColor}{$\triangleright$\enspace\texttt{HELP{\textunderscore}ADD{\textunderscore}BOOK({\mdseries\slshape short, long, dir})\index{HELP{\textunderscore}ADD{\textunderscore}BOOK@\texttt{HEL}\-\texttt{P{\textunderscore}}\-\texttt{A}\-\texttt{D}\-\texttt{D{\textunderscore}}\-\texttt{BOOK}}
\label{HELPuScoreADDuScoreBOOK}
}\hfill{\scriptsize (function)}}\\


 This command tells \textsf{GAP} that in directory \mbox{\texttt{\mdseries\slshape dir}} (given as either a string describing the path relative to the \textsf{GAP} root directory \texttt{GAPInfo.RootPaths[1]} or as directory object) contains the basic information about a help book. The
string \mbox{\texttt{\mdseries\slshape short}} is used as an identifying name for that book by the online help. The string \mbox{\texttt{\mdseries\slshape long}} should be a short explanation of the content of the book. Both strings
together should easily fit on a line, since they are displayed with \texttt{?books}. 

 It is possible to reinstall a book with different strings \mbox{\texttt{\mdseries\slshape short}}, \mbox{\texttt{\mdseries\slshape long}}; (for example, documentation of a not\texttt{\symbol{45}}loaded \textsf{GAP} package indicates this in the string \mbox{\texttt{\mdseries\slshape short}} and if you later load the package, \textsf{GAP} quietly changes the string \mbox{\texttt{\mdseries\slshape short}} as it reinstalls its documentation). 

 The only condition necessary to make the installation of a book \emph{valid} is that the directory \mbox{\texttt{\mdseries\slshape dir}} must contain a file \texttt{manual.six}. The next section explains how this file must look. }

 

\subsection{\textcolor{Chapter }{HELP{\textunderscore}REMOVE{\textunderscore}BOOK}}
\logpage{[ 84, 1, 2 ]}\nobreak
\hyperdef{L}{X7BDEB25D7AFC4322}{}
{\noindent\textcolor{FuncColor}{$\triangleright$\enspace\texttt{HELP{\textunderscore}REMOVE{\textunderscore}BOOK({\mdseries\slshape short})\index{HELP{\textunderscore}REMOVE{\textunderscore}BOOK@\texttt{HEL}\-\texttt{P{\textunderscore}}\-\texttt{R}\-\texttt{E}\-\texttt{M}\-\texttt{O}\-\texttt{V}\-\texttt{E{\textunderscore}}\-\texttt{BOOK}}
\label{HELPuScoreREMOVEuScoreBOOK}
}\hfill{\scriptsize (function)}}\\


 This command tells \textsf{GAP} not to use the help book with identifying name \mbox{\texttt{\mdseries\slshape short}} any more. The book can be re\texttt{\symbol{45}}installed using \texttt{HELP{\textunderscore}ADD{\textunderscore}BOOK} (\ref{HELPuScoreADDuScoreBOOK}). 

 }

 }

  
\section{\textcolor{Chapter }{The manual.six File}}\label{The manual.six File}
\logpage{[ 84, 2, 0 ]}
\hyperdef{L}{X8713EEAE840CEDA3}{}
{
  The first non\texttt{\symbol{45}}empty line of \texttt{manual.six} should be of the form 

 \texttt{\#SIXFORMAT} \mbox{\texttt{\mdseries\slshape myownformat}} 

 where \mbox{\texttt{\mdseries\slshape myownformat}} is an identifying string for this format. The reading of the (remainder of
the) file is then delegated to the function \texttt{HELP{\textunderscore}BOOK{\textunderscore}HANDLER.\mbox{\texttt{\mdseries\slshape myownformat}}.ReadSix} which must exist. Thus there are no further regulations for the format of the \texttt{manual.six} file, other that what you yourself impose. If such a line is missing then it
is assumed that the \texttt{manual.six} file complies with the \texttt{gapmacro.tex} documentation format, which internally is referred to as the \texttt{default} format for historical reasons. In that case reading the file is delegated to \texttt{HELP{\textunderscore}BOOK{\textunderscore}HANDLER.default.ReadSix}. 

 Section \ref{The Help Book Handler} explains how the return value of \texttt{HELP{\textunderscore}BOOK{\textunderscore}HANDLER.\mbox{\texttt{\mdseries\slshape myownformat}}.ReadSix} should look like and which further function should be contained in \texttt{HELP{\textunderscore}BOOK{\textunderscore}HANDLER.\mbox{\texttt{\mdseries\slshape myownformat}}}. }

  
\section{\textcolor{Chapter }{The Help Book Handler}}\label{The Help Book Handler}
\logpage{[ 84, 3, 0 ]}
\hyperdef{L}{X7AD7541E7C30D5B3}{}
{
  \index{document formats!for help books} For each document format \mbox{\texttt{\mdseries\slshape myownformat}} there must be a record \texttt{HELP{\textunderscore}BOOK{\textunderscore}HANDLER.\mbox{\texttt{\mdseries\slshape myownformat}}} of functions with the following names and functionality. 

 An implementation example of such a set of handler functions is the \texttt{default} format, which is the format name used for the \texttt{gapmacro.tex} documentation format, and this is contained in the file \texttt{lib/helpdef.gi}. 

 The package \textsf{GAPDoc} (see Chapter{\nobreakspace} (\textbf{GAPDoc: Introduction and Example})) also defines a format (as it should) which is called: \texttt{GapDocGAP} (the case \emph{is} significant). 

 As you can see by the above two examples, the name for a document format can
be anything, but it should be in some way meaningful. 

 
\begin{description}
\item[{ \texttt{ReadSix( \mbox{\texttt{\mdseries\slshape stream}} )}}]  For an input text stream \mbox{\texttt{\mdseries\slshape stream}} to a \texttt{manual.six} file, this must return a record \mbox{\texttt{\mdseries\slshape info}} which has at least the following two components: \texttt{bookname} which is the short identifying name of the help book, and \texttt{entries}. Here \mbox{\texttt{\mdseries\slshape info}}\texttt{.entries} must be a list with one entry per search string (which can be a section
header, function name, index entry, or whatever seems sensible to be searched
for matching a help query). A \emph{match} for the \textsf{GAP} help is a pair (\mbox{\texttt{\mdseries\slshape info}}, \mbox{\texttt{\mdseries\slshape i}}) where \mbox{\texttt{\mdseries\slshape i}} refers to an index for the list \mbox{\texttt{\mdseries\slshape info}}\texttt{.entries} and this corresponds to a certain position in the document. There is one
further regulation for the format of the entries of \mbox{\texttt{\mdseries\slshape info}}\texttt{.entries}. They must be lists and the first element of such a list must be a string
which is printed by \textsf{GAP} for example when several matches are found for a query (so it should
essentially be the string which is searched for the match, except that it may
contain upper and lower case letters or some markup). There may be other
components in \mbox{\texttt{\mdseries\slshape info}} which are needed by the functions below, but their names and formats are not
prescribed. The \mbox{\texttt{\mdseries\slshape stream}} argument is typically generated using \texttt{InputTextFile} (\ref{InputTextFile}), e.g. 

 
\begin{Verbatim}[commandchars=!@|,fontsize=\small,frame=single,label=Example]
  !gapprompt@gap>| !gapinput@dirs := DirectoriesLibrary( "doc/ref" );;|
  !gapprompt@gap>| !gapinput@file := Filename( dirs, "manual.six" );;|
  !gapprompt@gap>| !gapinput@stream := InputTextFile( file );;|
\end{Verbatim}
 
\item[{\texttt{ShowChapters( \mbox{\texttt{\mdseries\slshape info}} )} }]  This must return a text string or list of text lines which contains the
chapter headers of the book \mbox{\texttt{\mdseries\slshape info}}\texttt{.bookname}. 
\item[{\texttt{ShowSection( \mbox{\texttt{\mdseries\slshape info}} )} }]  This must return a text string or list of text lines which contains the
section (and chapter) headers of the book \mbox{\texttt{\mdseries\slshape info}}\texttt{.bookname}. 
\item[{\texttt{SearchMatches( \mbox{\texttt{\mdseries\slshape info}}, \mbox{\texttt{\mdseries\slshape topic}}, \mbox{\texttt{\mdseries\slshape frombegin}} )} }]  This function must return a list of indices of \mbox{\texttt{\mdseries\slshape info}}\texttt{.entries} for entries which match the search string \mbox{\texttt{\mdseries\slshape topic}}. If \mbox{\texttt{\mdseries\slshape frombegin}} is \texttt{true} then those parts of \mbox{\texttt{\mdseries\slshape topic}} which are separated by spaces should be considered as the beginnings of words
to decide the matching. It \mbox{\texttt{\mdseries\slshape frombegin}} is \texttt{false}, a substring search should be performed. The string \mbox{\texttt{\mdseries\slshape topic}} can be assumed to be already normalized (transformed to lower case, and
whitespace normalized). The function must return a list with two entries \texttt{[exact, match]} where \texttt{exact} is the list of indices for exact matches and \texttt{match} a list of indices of the remaining matches. 
\item[{\texttt{MatchPrevChap( \mbox{\texttt{\mdseries\slshape info}}, \mbox{\texttt{\mdseries\slshape i}} )} }]  This should return the match [\mbox{\texttt{\mdseries\slshape info}}, \texttt{j}] which points to the beginning of the chapter containing match [\mbox{\texttt{\mdseries\slshape info}}, \mbox{\texttt{\mdseries\slshape i}}], respectively to the beginning of the previous chapter if [\mbox{\texttt{\mdseries\slshape info}}, \mbox{\texttt{\mdseries\slshape i}}] is already the beginning of a chapter. (Corresponds to \texttt{?{\textless}{\textless}}.) 
\item[{\texttt{MatchNextChap( \mbox{\texttt{\mdseries\slshape info}}, \mbox{\texttt{\mdseries\slshape i}} )} }]  Like the previous function except that it should return the match for the
beginning of the next chapter. (Corresponds to \texttt{?{\textgreater}{\textgreater}}.) 
\item[{\texttt{MatchPrev( \mbox{\texttt{\mdseries\slshape info}}, \mbox{\texttt{\mdseries\slshape i}} )} }]  This should return the previous section (or appropriate portion of the
document). (Corresponds to \texttt{?{\textless}}.) 
\item[{\texttt{MatchNext( \mbox{\texttt{\mdseries\slshape info}}, \mbox{\texttt{\mdseries\slshape i}} )} }]  Like the previous function except that it should return the next section (or
appropriate portion of the document). (Corresponds to \texttt{?{\textgreater}}.) 
\item[{\texttt{HelpData( \mbox{\texttt{\mdseries\slshape info}}, \mbox{\texttt{\mdseries\slshape i}}, \mbox{\texttt{\mdseries\slshape type}} )} }]  This returns for match [\mbox{\texttt{\mdseries\slshape info}}, \mbox{\texttt{\mdseries\slshape i}}] some data whose format depends on the string \mbox{\texttt{\mdseries\slshape type}}, or \texttt{fail} if these data are not available. The values of \mbox{\texttt{\mdseries\slshape type}} which currently must be handled and the corresponding result format are
described in the list below. 
\item[{\texttt{SubsectionNumber( \mbox{\texttt{\mdseries\slshape info}}, \mbox{\texttt{\mdseries\slshape i}} )} }]  This returns some \textsf{GAP} object that identifies the position in the book where the display of this
entry is started. This can be useful to detect if several help book entries
actually point to the same place. 
\end{description}
 

 The \texttt{HELP{\textunderscore}BOOK{\textunderscore}HANDLER.\mbox{\texttt{\mdseries\slshape myownformat}}.HelpData} function must recognize the following values of the \mbox{\texttt{\mdseries\slshape type}} argument. 

 
\begin{description}
\item[{ \texttt{"text"} }]  This must return a corresponding text string in a format which can be fed into
the \texttt{Pager}, see{\nobreakspace}\texttt{Pager} (\ref{Pager}). 
\item[{\texttt{"url"} }]  If the help book is available in HTML format this must return an URL as a
string (Probably a \texttt{file://} URL containing a label for the exact start position in that file). Otherwise
it returns \texttt{fail}. 
\item[{\texttt{"dvi"} }]  If the help book is available in dvi\texttt{\symbol{45}}format this must
return a record of form \texttt{rec( file := \mbox{\texttt{\mdseries\slshape filename}}, page := \mbox{\texttt{\mdseries\slshape pagenumber}} )}. Otherwise it returns \texttt{fail}. 
\item[{\texttt{"pdf"} }]  Same as case \texttt{"dvi"}, but for the corresponding pdf\texttt{\symbol{45}}file. 
\item[{\texttt{"secnr"} }]  This must return a pair like \texttt{[[3,3,1], "3.3.1"]} which gives the section number as chapter number, section number, subsection
number triple and a corresponding string (a chapter itself is encoded like \texttt{[[4,0,0], "4."]}). Useful for cross\texttt{\symbol{45}}referencing between help books. 
\end{description}
 }

  
\section{\textcolor{Chapter }{Introducing new Viewer for the Online Help}}\label{Introducing new Viewer for the Online Help}
\logpage{[ 84, 4, 0 ]}
\hyperdef{L}{X861927BF822FB162}{}
{
  To introduce a new viewer for the online help, one should extend the global
record \texttt{HELP{\textunderscore}VIEWER{\textunderscore}INFO} (\ref{HELPuScoreVIEWERuScoreINFO}), the structure of which is explained below. 

\subsection{\textcolor{Chapter }{HELP{\textunderscore}VIEWER{\textunderscore}INFO}}
\logpage{[ 84, 4, 1 ]}\nobreak
\hyperdef{L}{X84B011847A4D90F0}{}
{\noindent\textcolor{FuncColor}{$\triangleright$\enspace\texttt{HELP{\textunderscore}VIEWER{\textunderscore}INFO\index{HELP{\textunderscore}VIEWER{\textunderscore}INFO@\texttt{HEL}\-\texttt{P{\textunderscore}}\-\texttt{V}\-\texttt{I}\-\texttt{E}\-\texttt{W}\-\texttt{E}\-\texttt{R{\textunderscore}}\-\texttt{INFO}}
\label{HELPuScoreVIEWERuScoreINFO}
}\hfill{\scriptsize (global variable)}}\\


 The record \texttt{HELP{\textunderscore}VIEWER{\textunderscore}INFO} contains one component for each help viewer. Each such component is a record
with two components: \texttt{.type} and \texttt{.show}. 

 The component \texttt{.type} refers to one of the \texttt{type}s recognized by the \texttt{HelpData} handler function explained in the previous section (currently one of \texttt{"text"}, \texttt{"url"}, \texttt{"dvi"}, or \texttt{"pdf"}). 

 The component \texttt{.show} is a function which gets as input the result of a corresponding \texttt{HelpData} handler call, if it was not \texttt{fail}. This function has to perform the actual display of the data. (E.g., by
calling a function like \texttt{Pager} (\ref{Pager}) or by starting up an external viewer program.) }

 }

 }

           
\chapter{\textcolor{Chapter }{Function\texttt{\symbol{45}}Operation\texttt{\symbol{45}}Attribute Triples}}\label{Function-Operation-Attribute Triples}
\logpage{[ 85, 0, 0 ]}
\hyperdef{L}{X8350247A8501969F}{}
{
  \textsf{GAP} is eager to maintain information that it has gathered about an object,
possibly by lengthy calculations. The most important mechanism for information
maintenance is the automatic storage and look\texttt{\symbol{45}}up that takes
place for \emph{attributes}; and this was already mentioned in section{\nobreakspace} (\textbf{Tutorial: Attributes}). In this chapter we will describe further mechanisms that allow storage of
results that are not values of attributes. 

 \index{FOA triples} The idea which is common to all sections is that certain operations, which are
not themselves attributes, have an attribute associated with them. To
automatically delegate tasks to the attribute, \textsf{GAP} knows, in addition to the \emph{operation} and the \emph{attributes} also a \emph{function}, which is ``wrapped around'' the other two. This ``wrapper function'' is called by the user and decides whether to call the operation or the
attribute or possibly both. The whole \emph{f}unction\texttt{\symbol{45}}\emph{o}peration\texttt{\symbol{45}}\emph{a}ttribute triple (or \emph{FOA triple}) is set up by a single \textsf{GAP} command which writes the wrapper function and already installs some methods,
e.g., for the attribute to fall back on the operation. The idea is then that
subsequent methods, which perform the actual computation, are installed only
for the operation, whereas the wrapper function remains unaltered, and in
general no additional methods for the attribute are required either.  
\section{\textcolor{Chapter }{Key Dependent Operations}}\label{Key Dependent Operations}
\logpage{[ 85, 1, 0 ]}
\hyperdef{L}{X86F03E0D7C18C6B0}{}
{
  

\subsection{\textcolor{Chapter }{KeyDependentOperation}}
\logpage{[ 85, 1, 1 ]}\nobreak
\hyperdef{L}{X7CABFDAA8596757E}{}
{\noindent\textcolor{FuncColor}{$\triangleright$\enspace\texttt{KeyDependentOperation({\mdseries\slshape name, dom\texttt{\symbol{45}}req, key\texttt{\symbol{45}}req, key\texttt{\symbol{45}}test})\index{KeyDependentOperation@\texttt{KeyDependentOperation}}
\label{KeyDependentOperation}
}\hfill{\scriptsize (function)}}\\


 There are several functions that require as first argument a domain, e.g., a
group, and as second argument something much simpler, e.g., a prime. \texttt{SylowSubgroup} (\ref{SylowSubgroup}) is an example. Since its value depends on two arguments, it cannot be an
attribute, yet one would like to store the Sylow subgroups once they have been
computed. 

 The idea is to provide an attribute of the group, called \texttt{ComputedSylowSubgroups}, and to store the groups in this list. The name implies that the value of
this attribute may change in the course of a \textsf{GAP} session, whenever a newly\texttt{\symbol{45}}computed Sylow subgroup is put
into the list. Therefore, this is a \emph{mutable attribute} (see \ref{Attributes}). The list contains primes in each bound odd position and a corresponding
Sylow subgroup in the following even position. More precisely, if \texttt{\mbox{\texttt{\mdseries\slshape p}} = ComputedSylowSubgroups( \mbox{\texttt{\mdseries\slshape G}} )[ \mbox{\texttt{\mdseries\slshape even}} \texttt{\symbol{45}} 1 ]} then \texttt{ComputedSylowSubgroups( \mbox{\texttt{\mdseries\slshape G}} )[ \mbox{\texttt{\mdseries\slshape even}} ]} holds the value of \texttt{SylowSubgroup( \mbox{\texttt{\mdseries\slshape G}}, \mbox{\texttt{\mdseries\slshape p}} )}. The pairs are sorted in increasing order of \mbox{\texttt{\mdseries\slshape p}}, in particular at most one Sylow \mbox{\texttt{\mdseries\slshape p}} subgroup of \mbox{\texttt{\mdseries\slshape G}} is stored for each prime \mbox{\texttt{\mdseries\slshape p}}. This attribute value is maintained by the function \texttt{SylowSubgroup} (\ref{SylowSubgroup}), which calls the operation \texttt{SylowSubgroupOp( \mbox{\texttt{\mdseries\slshape G}}, \mbox{\texttt{\mdseries\slshape p}} )} to do the real work, if the prime \mbox{\texttt{\mdseries\slshape p}} cannot be found in the list. So methods that do the real work should be
installed for \texttt{SylowSubgroupOp} and not for \texttt{SylowSubgroup} (\ref{SylowSubgroup}). 

 The same mechanism works for other functions as well, e.g., for \texttt{PCore} (\ref{PCore}), but also for \texttt{HallSubgroup} (\ref{HallSubgroup}), where the second argument is not a prime but a set of primes. 

 \texttt{KeyDependentOperation} declares the two operations and the attribute as described above, with names \mbox{\texttt{\mdseries\slshape name}}, \mbox{\texttt{\mdseries\slshape name}}\texttt{Op}, and \texttt{Computed}\mbox{\texttt{\mdseries\slshape name}}\texttt{s}, as well as tester and setter operations \texttt{Has}\mbox{\texttt{\mdseries\slshape name}} and \texttt{Set}\mbox{\texttt{\mdseries\slshape name}}, respectively. Note, however, that the tester is not a filter. \mbox{\texttt{\mdseries\slshape dom\texttt{\symbol{45}}req}} and \mbox{\texttt{\mdseries\slshape key\texttt{\symbol{45}}req}} specify the required filters for the first and second argument of the
operation \mbox{\texttt{\mdseries\slshape name}}\texttt{Op}, which are needed to create this operation with \texttt{DeclareOperation} (\ref{DeclareOperation}). \mbox{\texttt{\mdseries\slshape dom\texttt{\symbol{45}}req}} is also the required filter for the corresponding attribute \texttt{Computed}\mbox{\texttt{\mdseries\slshape name}}\texttt{s}. The fourth argument \mbox{\texttt{\mdseries\slshape key\texttt{\symbol{45}}test}} is in general a function to which the second argument \mbox{\texttt{\mdseries\slshape info}} of \texttt{\mbox{\texttt{\mdseries\slshape name}}( \mbox{\texttt{\mdseries\slshape D}}, \mbox{\texttt{\mdseries\slshape info}} )} will be passed. This function can perform tests on \mbox{\texttt{\mdseries\slshape info}}, and raise an error if appropriate. 

 For example, to set up the three objects \texttt{SylowSubgroup} (\ref{SylowSubgroup}), \texttt{SylowSubgroupOp}, \texttt{ComputedSylowSubgroups} together, the declaration file \texttt{lib/grp.gd} contains the following line of code. 
\begin{Verbatim}[commandchars=!@|,fontsize=\small,frame=single,label=Example]
  KeyDependentOperation( "SylowSubgroup", IsGroup, IsPosInt, "prime" );
\end{Verbatim}
 In this example, \mbox{\texttt{\mdseries\slshape key\texttt{\symbol{45}}test}} has the value \texttt{"prime"}, which is silently replaced by a function that tests whether its argument is
a prime. 

 
\begin{Verbatim}[commandchars=!@|,fontsize=\small,frame=single,label=Example]
  !gapprompt@gap>| !gapinput@s4 := Group((1,2,3,4),(1,2));;|
  !gapprompt@gap>| !gapinput@SylowSubgroup( s4, 7 );;  ComputedSylowSubgroups( s4 );|
  [ 7, Group(()) ]
  !gapprompt@gap>| !gapinput@SylowSubgroup( s4, 2 );;  ComputedSylowSubgroups( s4 );|
  [ 2, Group([ (3,4), (1,4)(2,3), (1,3)(2,4) ]), 7, Group(()) ]
  !gapprompt@gap>| !gapinput@HasSylowSubgroup( s4, 5 );|
  false
  !gapprompt@gap>| !gapinput@SetSylowSubgroup( s4, 5, Group(()));; ComputedSylowSubgroups( s4 );|
  [ 2, Group([ (3,4), (1,4)(2,3), (1,3)(2,4) ]), 5, Group(()), 7, Group(()) ]
\end{Verbatim}
 

 
\begin{Verbatim}[commandchars=!|B,fontsize=\small,frame=single,label=Example]
  !gapprompt|gap>B !gapinput|SylowSubgroup( s4, 6 );B
  Error, SylowSubgroup: <p> must be a prime
  Stack trace:
  *[1] <<compiled GAP function>>
     @ GAPROOT/lib/oper1.g:964
   [2] <<compiled GAP function "SylowSubgroup default method">>
     @ GAPROOT/lib/oper1.g:1007
  <function "SylowSubgroup default method">( <arguments> )
   called from read-eval loop at *stdin*:3
  you can enter 'quit;' to quit to outer loop, or
  you can enter 'return;' to continue
  !gapbrkprompt|brk>B !gapinput|quit;B
\end{Verbatim}
 

 Thus the prime test need not be repeated in the methods for the operation \texttt{SylowSubgroupOp} (which are installed to do the real work). Note that no methods need be
installed for \texttt{SylowSubgroup} (\ref{SylowSubgroup}) and \texttt{ComputedSylowSubgroups}. If a method is installed with \texttt{InstallMethod} (\ref{InstallMethod}) for a wrapper operation such as \texttt{SylowSubgroup} (\ref{SylowSubgroup}) then a warning is signalled provided the \texttt{InfoWarning} (\ref{InfoWarning}) level is at least \texttt{1}. (Use \texttt{InstallMethod} (\ref{InstallMethod}) in order to suppress the warning.) }

 }

  
\section{\textcolor{Chapter }{In Parent Attributes}}\label{In Parent Attributes}
\logpage{[ 85, 2, 0 ]}
\hyperdef{L}{X78D4D0FF780C8A85}{}
{
  

\subsection{\textcolor{Chapter }{InParentFOA}}
\logpage{[ 85, 2, 1 ]}\nobreak
\hyperdef{L}{X7C0E62D8813A4EE6}{}
{\noindent\textcolor{FuncColor}{$\triangleright$\enspace\texttt{InParentFOA({\mdseries\slshape name, super, sub, AorP})\index{InParentFOA@\texttt{InParentFOA}}
\label{InParentFOA}
}\hfill{\scriptsize (function)}}\\


 This section describes how you can add new ``in parent attributes'' (see{\nobreakspace}\ref{Constructing Subdomains} and \ref{Parents}). As an example, we describe how \texttt{Index} (\ref{Index:for a group and its subgroup}) and its related functions are implemented. 

 There are two operations \texttt{Index} (\ref{Index:for a group and its subgroup}) and \texttt{IndexOp}, and an attribute \texttt{IndexInParent}. They are created together as shown below, and after they have been created,
methods need be installed only for \texttt{IndexOp}. In the creation process, \texttt{IndexInParent} already gets one default method installed (in addition to the usual system
getter of each attribute, see{\nobreakspace}\ref{Attributes}), namely \texttt{D \texttt{\symbol{45}}{\textgreater} IndexOp( Parent( D ), D )}. 

 The operation \texttt{Index} (\ref{Index:for a group and its subgroup}) proceeds as follows. 
\begin{itemize}
\item  If it is called with the two arguments \mbox{\texttt{\mdseries\slshape super}} and \mbox{\texttt{\mdseries\slshape sub}}, and if \texttt{HasParent( \mbox{\texttt{\mdseries\slshape sub}} )} and \texttt{IsIdenticalObj( \mbox{\texttt{\mdseries\slshape super}}, Parent( \mbox{\texttt{\mdseries\slshape sub}} ) )} are \texttt{true}, \texttt{IndexInParent} is called with argument \mbox{\texttt{\mdseries\slshape sub}}, and the result is returned. 
\item  Otherwise, \texttt{IndexOp} is called with the same arguments that \texttt{Index} (\ref{Index:for a group and its subgroup}) was called with, and the result is returned. 
\end{itemize}
 (Note that it is in principle possible to install even \texttt{Index} (\ref{Index:for a group and its subgroup}) and \texttt{IndexOp} methods for a number of arguments different from two, with \texttt{InstallOtherMethod} (\ref{InstallOtherMethod}), see \ref{Attributes}). 

 The call of \texttt{InParentFOA} declares the operations and the attribute as described above, with names \mbox{\texttt{\mdseries\slshape name}}, \mbox{\texttt{\mdseries\slshape name}}\texttt{Op}, and \mbox{\texttt{\mdseries\slshape name}}\texttt{InParent}. \mbox{\texttt{\mdseries\slshape super\texttt{\symbol{45}}req}} and \mbox{\texttt{\mdseries\slshape sub\texttt{\symbol{45}}req}} specify the required filters for the first and second argument of the
operation \mbox{\texttt{\mdseries\slshape name}}\texttt{Op}, which are needed to create this operation with \texttt{DeclareOperation} (\ref{DeclareOperation}). \mbox{\texttt{\mdseries\slshape sub\texttt{\symbol{45}}req}} is also the required filter for the corresponding attribute \mbox{\texttt{\mdseries\slshape name}}\texttt{InParent}; note that \texttt{HasParent} (\ref{HasParent}) is \emph{not} required for the argument \mbox{\texttt{\mdseries\slshape U}} of \mbox{\texttt{\mdseries\slshape name}}\texttt{InParent}, because even without a parent stored, \texttt{Parent( \mbox{\texttt{\mdseries\slshape U}} )} is legal, meaning \mbox{\texttt{\mdseries\slshape U}} itself (see{\nobreakspace}\ref{Parents}). The fourth argument must be \texttt{DeclareProperty} (\ref{DeclareProperty}) if \mbox{\texttt{\mdseries\slshape name}}\texttt{InParent} takes only boolean values (for example in the case \texttt{IsNormalInParent}), and \texttt{DeclareAttribute} (\ref{DeclareAttribute}) otherwise. 

 For example, to set up the three objects \texttt{Index} (\ref{Index:for a group and its subgroup}), \texttt{IndexOp}, and \texttt{IndexInParent} together, the declaration file \texttt{lib/domain.gd} contains the following line of code. 
\begin{Verbatim}[commandchars=!@|,fontsize=\small,frame=single,label=Example]
  InParentFOA( "Index", IsGroup, IsGroup, DeclareAttribute );
\end{Verbatim}
 

 Note that no methods need be installed for \texttt{Index} (\ref{Index:for a group and its subgroup}) and \texttt{IndexInParent}. }

 }

  
\section{\textcolor{Chapter }{Operation Functions}}\label{Operation Functions}
\logpage{[ 85, 3, 0 ]}
\hyperdef{L}{X7CD4A0867BD825F7}{}
{
  \index{ExternalSet@\texttt{ExternalSet}!computing orbits} \index{G-sets@$G$\texttt{\symbol{45}}sets!computing orbits} \index{Orbits@\texttt{Orbits}!as attributes for external sets} Chapter{\nobreakspace}\ref{Group Actions} and, in particular, the Section{\nobreakspace}\ref{About Group Actions}   explain that certain operations such as \ref{Orbits}), besides their usual usage with arguments \mbox{\texttt{\mdseries\slshape G}}, \mbox{\texttt{\mdseries\slshape D}}, and \mbox{\texttt{\mdseries\slshape opr}}, can also be applied to an external set ($G$\texttt{\symbol{45}}set), in which case they can be interpreted as attributes.
Moreover, they can also be interpreted as attributes for permutation groups,
meaning the natural action on the set of its moved points. 

 The definition of \ref{Orbits} says that a method should be a function with arguments \mbox{\texttt{\mdseries\slshape G}}, \mbox{\texttt{\mdseries\slshape D}}, \mbox{\texttt{\mdseries\slshape gens}}, \mbox{\texttt{\mdseries\slshape oprs}}, and \mbox{\texttt{\mdseries\slshape opr}}, as in the case of the operation \texttt{ExternalSet} (\ref{ExternalSet}) when specified via \mbox{\texttt{\mdseries\slshape gens}} and \mbox{\texttt{\mdseries\slshape oprs}} (see{\nobreakspace}\ref{External Sets}). All other syntax variants allowed for \ref{Orbits} (e.g., leaving out \mbox{\texttt{\mdseries\slshape gens}} and \mbox{\texttt{\mdseries\slshape oprs}}) are handled by default methods. 

 The default methods for \ref{Orbits} support the following behaviour. 
\begin{enumerate}
\item  If the only argument is an external set \mbox{\texttt{\mdseries\slshape xset}} and the attribute tester \texttt{HasOrbits( \mbox{\texttt{\mdseries\slshape xset}} )} returns \texttt{true}, the stored value of that attribute is returned. 
\item  If the only argument is an external set \mbox{\texttt{\mdseries\slshape xset}} and the attribute value is not known, the default arguments are obtained from
the data of \mbox{\texttt{\mdseries\slshape xset}}. 
\item  If \mbox{\texttt{\mdseries\slshape gens}} and \mbox{\texttt{\mdseries\slshape oprs}} are not specified, \mbox{\texttt{\mdseries\slshape gens}} is set to \texttt{Pcgs( \mbox{\texttt{\mdseries\slshape G}} )} if \texttt{CanEasilyComputePcgs( \mbox{\texttt{\mdseries\slshape G}} )} is \texttt{true}, and to \texttt{GeneratorsOfGroup( \mbox{\texttt{\mdseries\slshape G}} )} otherwise; \mbox{\texttt{\mdseries\slshape oprs}} is set to \mbox{\texttt{\mdseries\slshape gens}}. 
\item  The default value of \mbox{\texttt{\mdseries\slshape opr}} is \texttt{OnPoints} (\ref{OnPoints}). 
\item  In the case of an operation of a permutation group \mbox{\texttt{\mdseries\slshape G}} on \texttt{MovedPoints( \mbox{\texttt{\mdseries\slshape G}} )} via \texttt{OnPoints} (\ref{OnPoints}), if the attribute tester \texttt{HasOrbits( \mbox{\texttt{\mdseries\slshape G}} )} returns \texttt{true}, the stored attribute value is returned. 
\item  The operation is called as \texttt{\mbox{\texttt{\mdseries\slshape result}}:= Orbits( \mbox{\texttt{\mdseries\slshape G}}, \mbox{\texttt{\mdseries\slshape D}}, \mbox{\texttt{\mdseries\slshape gens}}, \mbox{\texttt{\mdseries\slshape oprs}}, \mbox{\texttt{\mdseries\slshape opr}} )}. 
\item  In the case of an external set \mbox{\texttt{\mdseries\slshape xset}} or a permutation group \mbox{\texttt{\mdseries\slshape G}} in its natural action, the attribute setter is called to store \mbox{\texttt{\mdseries\slshape result}}. 
\item  \mbox{\texttt{\mdseries\slshape result}} is returned. 
\end{enumerate}
 

 The declaration of operations that match the above pattern is done as follows. 

\subsection{\textcolor{Chapter }{OrbitsishOperation}}
\logpage{[ 85, 3, 1 ]}\nobreak
\hyperdef{L}{X7CA3826A7EBDE208}{}
{\noindent\textcolor{FuncColor}{$\triangleright$\enspace\texttt{OrbitsishOperation({\mdseries\slshape name, reqs, usetype, AorP})\index{OrbitsishOperation@\texttt{OrbitsishOperation}}
\label{OrbitsishOperation}
}\hfill{\scriptsize (function)}}\\


 declares an attribute \texttt{op}, with name \mbox{\texttt{\mdseries\slshape name}}. The second argument \mbox{\texttt{\mdseries\slshape reqs}} specifies the list of required filters for the usual
(five\texttt{\symbol{45}}argument) methods that do the real work. 

 If the third argument \mbox{\texttt{\mdseries\slshape usetype}} is \texttt{true}, the function call \texttt{op( xset )} will {\textendash}if the value of \texttt{op} for \texttt{xset} is not yet known{\textendash} delegate to the five\texttt{\symbol{45}}argument
call of \texttt{op} with second argument \texttt{xset} rather than with \texttt{D}. This allows certain methods for \texttt{op} to make use of the type of \texttt{xset}, in which the types of the external subsets of \texttt{xset} and of the external orbits in \texttt{xset} are stored. (This is used to avoid repeated calls of \texttt{NewType} (\ref{NewType}) in functions like \texttt{ExternalOrbits( xset )}, which call \texttt{ExternalOrbit( xset, pnt )} for several values of \texttt{pnt}.) 

 For property testing functions such as \texttt{IsTransitive} (\ref{IsTransitive:for a group, an action domain, etc.}), the fourth argument \mbox{\texttt{\mdseries\slshape AorP}} must be \texttt{NewProperty} (\ref{NewProperty}), otherwise it must be \texttt{NewAttribute} (\ref{NewAttribute}); in the former case, a property is returned, in the latter case an attribute
that is not a property. 

 For example, to set up the operation \texttt{Orbits} (\ref{Orbits:operation}), the declaration file \texttt{lib/oprt.gd} contains the following line of code: 
\begin{Verbatim}[commandchars=!@|,fontsize=\small,frame=single,label=Example]
  OrbitsishOperation( "Orbits", OrbitsishReq, false, NewAttribute );
\end{Verbatim}
 The global variable \texttt{OrbitsishReq} contains the standard requirements 
\begin{Verbatim}[commandchars=!@|,fontsize=\small,frame=single,label=Example]
  OrbitsishReq := [ IsGroup, IsList,
                    IsList,
                    IsList,
                    IsFunction ];
\end{Verbatim}
 which are usually entered in calls to \texttt{OrbitsishOperation}. 

 The new operation, e.g., \texttt{Orbits} (\ref{Orbits:operation}), can be called either as \texttt{Orbits( \mbox{\texttt{\mdseries\slshape xset}} )} for an external set \mbox{\texttt{\mdseries\slshape xset}}, or as \texttt{Orbits( \mbox{\texttt{\mdseries\slshape G}} )} for a permutation group \mbox{\texttt{\mdseries\slshape G}}, meaning the orbits on the moved points of \mbox{\texttt{\mdseries\slshape G}} via \texttt{OnPoints} (\ref{OnPoints}), or as 

 \texttt{Orbits( \mbox{\texttt{\mdseries\slshape G}}, \mbox{\texttt{\mdseries\slshape Omega}}[, \mbox{\texttt{\mdseries\slshape gens}}, \mbox{\texttt{\mdseries\slshape acts}}][, \mbox{\texttt{\mdseries\slshape act}}] )}, 

 with a group \mbox{\texttt{\mdseries\slshape G}}, a domain or list \mbox{\texttt{\mdseries\slshape Omega}}, generators \mbox{\texttt{\mdseries\slshape gens}} of \mbox{\texttt{\mdseries\slshape G}}, and corresponding elements \mbox{\texttt{\mdseries\slshape acts}} that act on \mbox{\texttt{\mdseries\slshape Omega}} via the function \mbox{\texttt{\mdseries\slshape act}}; the default of \mbox{\texttt{\mdseries\slshape gens}} and \mbox{\texttt{\mdseries\slshape acts}} is a list of group generators of \mbox{\texttt{\mdseries\slshape G}}, the default of \mbox{\texttt{\mdseries\slshape act}} is \texttt{OnPoints} (\ref{OnPoints}). 

 Only methods for the five\texttt{\symbol{45}}argument version need to be
installed for doing the real work. (And of course methods for one argument in
case one wants to define a new meaning of the attribute.) }

 

\subsection{\textcolor{Chapter }{OrbitishFO}}
\logpage{[ 85, 3, 2 ]}\nobreak
\hyperdef{L}{X7B23C48482ADB237}{}
{\noindent\textcolor{FuncColor}{$\triangleright$\enspace\texttt{OrbitishFO({\mdseries\slshape name, reqs, famrel, usetype, realenum})\index{OrbitishFO@\texttt{OrbitishFO}}
\label{OrbitishFO}
}\hfill{\scriptsize (function)}}\\


 is used to create operations like \texttt{Orbit} (\ref{Orbit}). This function is analogous to \texttt{OrbitsishOperation} (\ref{OrbitsishOperation}), but for operations \mbox{\texttt{\mdseries\slshape orbish}} like \texttt{Orbit( \mbox{\texttt{\mdseries\slshape G}}, \mbox{\texttt{\mdseries\slshape Omega}}, \mbox{\texttt{\mdseries\slshape pnt}} )}. Since the return values of these operations depend on the additional
argument \mbox{\texttt{\mdseries\slshape pnt}}, there is no associated attribute. 

 The call of \texttt{OrbitishFO} declares a wrapper function and its operation, with names \mbox{\texttt{\mdseries\slshape name}} and \mbox{\texttt{\mdseries\slshape name}}\texttt{Op}. 

 The second argument \mbox{\texttt{\mdseries\slshape reqs}} specifies the list of required filters for the operation \mbox{\texttt{\mdseries\slshape name}}\texttt{Op}. 

 The third argument \mbox{\texttt{\mdseries\slshape famrel}} is used to test the family relation between the second and third argument of \texttt{\mbox{\texttt{\mdseries\slshape name}}( \mbox{\texttt{\mdseries\slshape G}}, \mbox{\texttt{\mdseries\slshape D}}, \mbox{\texttt{\mdseries\slshape pnt}} )}. For example, \mbox{\texttt{\mdseries\slshape famrel}} is \texttt{IsCollsElms} in the case of \texttt{Orbit} (\ref{Orbit}) because \mbox{\texttt{\mdseries\slshape pnt}} must be an element of \mbox{\texttt{\mdseries\slshape D}}. Similarly, in the call \texttt{Blocks( \mbox{\texttt{\mdseries\slshape G}}, \mbox{\texttt{\mdseries\slshape D}}, \mbox{\texttt{\mdseries\slshape seed}} )}, \mbox{\texttt{\mdseries\slshape seed}} must be a subset of \mbox{\texttt{\mdseries\slshape D}}, and the family relation must be \texttt{IsIdenticalObj} (\ref{IsIdenticalObj}). 

 The fourth argument \mbox{\texttt{\mdseries\slshape usetype}} serves the same purpose as in the case of \texttt{OrbitsishOperation} (\ref{OrbitsishOperation}). \mbox{\texttt{\mdseries\slshape usetype}} can also be an attribute, such as \texttt{BlocksAttr} or \texttt{MaximalBlocksAttr}. In this case, if only one of the two arguments \mbox{\texttt{\mdseries\slshape Omega}} and \mbox{\texttt{\mdseries\slshape pnt}} is given, blocks with no seed are computed, they are stored as attribute
values according to the rules of \texttt{OrbitsishOperation} (\ref{OrbitsishOperation}). 

 If the 5th argument is set to \texttt{true}, the action for an external set should use the enumerator, otherwise it uses
the \texttt{HomeEnumerator} (\ref{HomeEnumerator}) value. This will make a difference for external orbits as part of a larger
domain. }

 
\subsection{\textcolor{Chapter }{Example: Orbit and OrbitOp}}\label{Example: Orbit and OrbitOp}
\logpage{[ 85, 3, 3 ]}
\hyperdef{L}{X834E92F07DD0BF04}{}
{
  For example, to setup the function \texttt{Orbit} (\ref{Orbit}) and its operation \texttt{OrbitOp}, the declaration file \texttt{lib/oprt.gd} contains the following line of code: 
\begin{Verbatim}[commandchars=!@|,fontsize=\small,frame=single,label=Example]
  OrbitishFO( "Orbit", OrbitishReq, IsCollsElms, false, false );
\end{Verbatim}
 The variable \texttt{OrbitishReq} contains the standard requirements 
\begin{Verbatim}[commandchars=!@|,fontsize=\small,frame=single,label=Example]
  OrbitishReq  := [ IsGroup, IsList, IsObject,
                    IsList,
                    IsList,
                    IsFunction ];
\end{Verbatim}
 which are usually entered in calls to \texttt{OrbitishFO} (\ref{OrbitishFO}). 

 The relation test via \mbox{\texttt{\mdseries\slshape famrel}} is used to provide a uniform construction of the wrapper functions created by \texttt{OrbitishFO} (\ref{OrbitishFO}), in spite of the different syntax of the specific functions. For example, \texttt{Orbit} (\ref{Orbit}) admits the calls \texttt{Orbit( \mbox{\texttt{\mdseries\slshape G}}, \mbox{\texttt{\mdseries\slshape D}}, \mbox{\texttt{\mdseries\slshape pnt}}, \mbox{\texttt{\mdseries\slshape opr}} )} and \texttt{Orbit( \mbox{\texttt{\mdseries\slshape G}}, \mbox{\texttt{\mdseries\slshape pnt}}, \mbox{\texttt{\mdseries\slshape opr}} )}, i.e., the second argument \mbox{\texttt{\mdseries\slshape D}} may be omitted; \texttt{Blocks} (\ref{Blocks:for a group, an action domain, etc.}) admits the calls \texttt{Blocks( \mbox{\texttt{\mdseries\slshape G}}, \mbox{\texttt{\mdseries\slshape D}}, \mbox{\texttt{\mdseries\slshape seed}}, \mbox{\texttt{\mdseries\slshape opr}} )} and \texttt{Blocks( \mbox{\texttt{\mdseries\slshape G}}, \mbox{\texttt{\mdseries\slshape D}}, \mbox{\texttt{\mdseries\slshape opr}} )}, i.e., the third argument may be omitted. The translation to the appropriate
call of \texttt{OrbitOp} or \texttt{BlocksOp}, for either operation with five or six arguments, is handled via \mbox{\texttt{\mdseries\slshape famrel}}. 

 As a consequence, there must not only be methods for \texttt{OrbitOp} with the six arguments corresponding to \texttt{OrbitishReq}, but also methods for only five arguments (i.e., without \mbox{\texttt{\mdseries\slshape D}}). Plenty of examples are contained in the implementation file \texttt{lib/oprt.gi}. 

 In order to handle a few special cases (currently \texttt{Blocks} (\ref{Blocks:for a group, an action domain, etc.}) and \texttt{MaximalBlocks} (\ref{MaximalBlocks:for a group, an action domain, etc.})), also the following form of \texttt{OrbitishFO} (\ref{OrbitishFO}) is supported. 

 \texttt{OrbitishFO( \mbox{\texttt{\mdseries\slshape name}}, \mbox{\texttt{\mdseries\slshape reqs}}, \mbox{\texttt{\mdseries\slshape famrel}}, \mbox{\texttt{\mdseries\slshape attr}} )} 

 The functions in question depend upon an argument \mbox{\texttt{\mdseries\slshape seed}}, so they cannot be regarded as attributes. However, they are most often
called without giving \mbox{\texttt{\mdseries\slshape seed}}, meaning ``choose any minimal resp. maximal block system''. In this case, the result can be stored as the value of the attribute \mbox{\texttt{\mdseries\slshape attr}} that was entered as fourth argument of \texttt{OrbitishFO} (\ref{OrbitishFO}). This attribute is considered by a call \texttt{Blocks( \mbox{\texttt{\mdseries\slshape G}}, \mbox{\texttt{\mdseries\slshape D}}, \mbox{\texttt{\mdseries\slshape opr}} )} (i.e., without \mbox{\texttt{\mdseries\slshape seed}}) in the same way as \texttt{Orbits} (\ref{Orbits:attribute}) considers \texttt{OrbitsAttr}. 

 To set this up, the declaration file \texttt{lib/oprt.gd} contains the following lines: 
\begin{Verbatim}[commandchars=!@|,fontsize=\small,frame=single,label=Example]
  DeclareAttribute( "BlocksAttr", IsExternalSet );
  OrbitishFO( "Blocks",
      [ IsGroup, IsList, IsList,
        IsList,
        IsList,
        IsFunction ], IsIdenticalObj, BlocksAttr, true );
\end{Verbatim}
 And this extraordinary FOA triple works as follows: 
\begin{Verbatim}[commandchars=!@|,fontsize=\small,frame=single,label=Example]
  !gapprompt@gap>| !gapinput@s4 := Group((1,2,3,4),(1,2));;|
  !gapprompt@gap>| !gapinput@Blocks( s4, MovedPoints(s4), [1,2] );|
  [ [ 1, 2, 3, 4 ] ]
  !gapprompt@gap>| !gapinput@Tester( BlocksAttr )( s4 );|
  false
  !gapprompt@gap>| !gapinput@Blocks( s4, MovedPoints(s4) );|
  [ [ 1, 2, 3, 4 ] ]
  !gapprompt@gap>| !gapinput@Tester( BlocksAttr )( s4 );  BlocksAttr( s4 );|
  true
  [ [ 1, 2, 3, 4 ] ]
\end{Verbatim}
 }

 }

 }

              
\chapter{\textcolor{Chapter }{Weak Pointers}}\label{Weak Pointers}
\logpage{[ 86, 0, 0 ]}
\hyperdef{L}{X86390538806F67CF}{}
{
  This chapter describes the use of the kernel feature of \emph{weak pointers}. This feature is primarily intended for use only in \textsf{GAP} internals, and should be used extremely carefully otherwise. 

 The garbage collector (see Section \ref{Garbage Collection}) is the part of the kernel that manages memory in the user's workspace. It
will normally only reclaim the storage used by an object when the object
cannot be reached as a subobject of any \textsf{GAP} variable, or from any reference in the kernel. We say that any link to object $a$ from object $b$ ``keeps object $a$ alive'', as long as $b$ is alive. It is occasionally convenient, however, to have a link to an object
which \emph{does not keep it alive}, and this is a weak pointer. The most common use is in caches, and similar
structures, where it is only necessary to remember how to solve problem $x$ as long as some other link to $x$ exists. 

 The following section \ref{Weak Pointer Objects} describes the semantics of the objects that contain weak pointers. Following
sections describe the functions available to manipulate them.  
\section{\textcolor{Chapter }{Weak Pointer Objects}}\label{Weak Pointer Objects}
\logpage{[ 86, 1, 0 ]}
\hyperdef{L}{X86D963DC7968899B}{}
{
  A \emph{weak pointer object} is similar to a mutable plain list, except that it does not keep its
subobjects alive during a garbage collection. From the \textsf{GAP} viewpoint this means that its entries may become unbound, apparently
spontaneously, at any time. Considerable care is therefore needed in
programming with such an object. 

\subsection{\textcolor{Chapter }{WeakPointerObj}}
\logpage{[ 86, 1, 1 ]}\nobreak
\hyperdef{L}{X8155EE1386F46063}{}
{\noindent\textcolor{FuncColor}{$\triangleright$\enspace\texttt{WeakPointerObj({\mdseries\slshape list})\index{WeakPointerObj@\texttt{WeakPointerObj}}
\label{WeakPointerObj}
}\hfill{\scriptsize (function)}}\\


 \texttt{WeakPointerObj} returns a weak pointer object which contains the same subobjects as the list \mbox{\texttt{\mdseries\slshape list}}, that is it returns a \emph{shallow} weak copy of \mbox{\texttt{\mdseries\slshape list}}. 

 
\begin{Verbatim}[commandchars=!@|,fontsize=\small,frame=single,label=Example]
  !gapprompt@gap>| !gapinput@w := WeakPointerObj( [ 1, , [2,3], fail, rec( a := 1) ] );|
  WeakPointerObj( [ 1, , [ 2, 3 ], fail, rec( a := 1 ) ] )
\end{Verbatim}
 

 After some computations involving garbage collections (but not necessarily in
the \emph{first} garbage collection after the above assignment), \textsf{GAP} will notice that the list and the record stored in \texttt{w} are not referenced by other objects than \texttt{w}, and that therefore these entries may disappear. 

 
\begin{Verbatim}[commandchars=!@|,fontsize=\small,frame=single,label=Example]
  !gapprompt@gap>| !gapinput@CollectGarbage( true );|
  
  ... (perhaps more computations and garbage collections) ...
  
  !gapprompt@gap>| !gapinput@CollectGarbage( true );|
  !gapprompt@gap>| !gapinput@w;|
  WeakPointerObj( [ 1, , , fail ] )
\end{Verbatim}
 

 Note that \texttt{w} has failed to keep its list and record subobjects alive during the garbage
collections. Certain subobjects, such as small integers and elements of small
finite fields, are not stored in the workspace, and so are not subject to
garbage collection, while certain other objects, such as the boolean values,
are always reachable from global variables or the kernel and so are never
garbage collected. 

 Subobjects reachable without going through a weak pointer object do not
evaporate, as in: 

 
\begin{Verbatim}[commandchars=!@|,fontsize=\small,frame=single,label=Example]
  !gapprompt@gap>| !gapinput@w := WeakPointerObj( [ 1, , , fail ] );|
  WeakPointerObj( [ 1, , , fail ] )
  !gapprompt@gap>| !gapinput@l := [1,2,3];;|
  !gapprompt@gap>| !gapinput@w[1] := l;;|
  !gapprompt@gap>| !gapinput@w;|
  WeakPointerObj( [ [ 1, 2, 3 ], , , fail ] )
  !gapprompt@gap>| !gapinput@CollectGarbage( true );|
  !gapprompt@gap>| !gapinput@w;|
  WeakPointerObj( [ [ 1, 2, 3 ], , , fail ] )
\end{Verbatim}
 

 Note also that the global variables \texttt{last}, \texttt{last2} and \texttt{last3} will keep things alive {\textendash}this can be confusing when debugging. }

 }

  
\section{\textcolor{Chapter }{Low Level Access Functions for Weak Pointer Objects}}\label{Low Level Access Functions for Weak Pointer Objects}
\logpage{[ 86, 2, 0 ]}
\hyperdef{L}{X7F4476958497F239}{}
{
  \index{ElmWPObj} 

\subsection{\textcolor{Chapter }{SetElmWPObj}}
\logpage{[ 86, 2, 1 ]}\nobreak
\hyperdef{L}{X7B9748ED7BAAA379}{}
{\noindent\textcolor{FuncColor}{$\triangleright$\enspace\texttt{SetElmWPObj({\mdseries\slshape wp, pos, val})\index{SetElmWPObj@\texttt{SetElmWPObj}}
\label{SetElmWPObj}
}\hfill{\scriptsize (function)}}\\
\noindent\textcolor{FuncColor}{$\triangleright$\enspace\texttt{UnbindElmWPObj({\mdseries\slshape wp, pos})\index{UnbindElmWPObj@\texttt{UnbindElmWPObj}}
\label{UnbindElmWPObj}
}\hfill{\scriptsize (function)}}\\
\noindent\textcolor{FuncColor}{$\triangleright$\enspace\texttt{ElmWPObj({\mdseries\slshape wp, pos})\index{ElmWPObj@\texttt{ElmWPObj}}
\label{ElmWPObj}
}\hfill{\scriptsize (function)}}\\
\noindent\textcolor{FuncColor}{$\triangleright$\enspace\texttt{IsBoundElmWPObj({\mdseries\slshape wp, pos})\index{IsBoundElmWPObj@\texttt{IsBoundElmWPObj}}
\label{IsBoundElmWPObj}
}\hfill{\scriptsize (function)}}\\
\noindent\textcolor{FuncColor}{$\triangleright$\enspace\texttt{LengthWPObj({\mdseries\slshape wp})\index{LengthWPObj@\texttt{LengthWPObj}}
\label{LengthWPObj}
}\hfill{\scriptsize (function)}}\\


 The functions \texttt{SetElmWPObj} and \texttt{UnbindElmWPObj} set and unbind entries in a weak pointer object. 

 The function \texttt{ElmWPObj} returns the element at position \mbox{\texttt{\mdseries\slshape pos}} of the weak pointer object \mbox{\texttt{\mdseries\slshape wp}}, if there is one, and \texttt{fail} otherwise. A return value of \texttt{fail} can thus arise either because (a) the value \texttt{fail} is stored at position \mbox{\texttt{\mdseries\slshape pos}}, or (b) no value is stored at position \mbox{\texttt{\mdseries\slshape pos}}. Since \texttt{fail} cannot vanish in a garbage collection, these two cases can safely be
distinguished by a \emph{subsequent} call to \texttt{IsBoundElmWPObj}, which returns \texttt{true} if there is currently a value bound at position \mbox{\texttt{\mdseries\slshape pos}} of \mbox{\texttt{\mdseries\slshape wp}} and \texttt{false} otherwise. 

 Note that it is \emph{not} safe to write: 

 \texttt{if IsBoundElmWPObj(w,i) then x:= ElmWPObj(w,i); fi;} 

 and treat \texttt{x} as reliably containing a value taken from \texttt{w}, as a badly timed garbage collection could leave \texttt{x} containing \texttt{fail}. Instead use 

 \texttt{x := ElmWPObj(w,i); if x {\textless}{\textgreater} fail or
IsBoundElmWPObj(w,i) then . . .}. 

 Here is an example. 

 
\begin{Verbatim}[commandchars=!@|,fontsize=\small,frame=single,label=Example]
  !gapprompt@gap>| !gapinput@w := WeakPointerObj( [ 1, , [2,3], fail, rec() ] );|
  WeakPointerObj( [ 1, , [ 2, 3 ], fail, rec(  ) ] )
  !gapprompt@gap>| !gapinput@SetElmWPObj(w,5,[]);|
  !gapprompt@gap>| !gapinput@w;|
  WeakPointerObj( [ 1, , [ 2, 3 ], fail, [  ] ] )
  !gapprompt@gap>| !gapinput@UnbindElmWPObj(w,1);|
  !gapprompt@gap>| !gapinput@w;|
  WeakPointerObj( [ , , [ 2, 3 ], fail, [  ] ] )
  !gapprompt@gap>| !gapinput@ElmWPObj(w,3);|
  [ 2, 3 ]
  !gapprompt@gap>| !gapinput@ElmWPObj(w,1);|
  fail
\end{Verbatim}
 

 Now after some computations and garbage collections $\ldots$ 

 
\begin{Verbatim}[commandchars=!@|,fontsize=\small,frame=single,label=Example]
  !gapprompt@gap>| !gapinput@2;; 3;; 4;; CollectGarbage( true );  # clear last, last2, last3|
\end{Verbatim}
 

 $\ldots$ we get the following. 

 
\begin{Verbatim}[commandchars=!@|,fontsize=\small,frame=single,label=Example]
  !gapprompt@gap>| !gapinput@ElmWPObj(w,3);|
  fail
  !gapprompt@gap>| !gapinput@w;|
  WeakPointerObj( [ , , , fail ] )
  !gapprompt@gap>| !gapinput@ElmWPObj(w,4);|
  fail
  !gapprompt@gap>| !gapinput@IsBoundElmWPObj(w,3);|
  false
  !gapprompt@gap>| !gapinput@IsBoundElmWPObj(w,4);|
  true
\end{Verbatim}
 }

 }

  
\section{\textcolor{Chapter }{Accessing Weak Pointer Objects as Lists}}\label{Accessing Weak Pointer Objects as Lists}
\logpage{[ 86, 3, 0 ]}
\hyperdef{L}{X8468DD647DDEFD82}{}
{
  Weak pointer objects are members of \texttt{ListsFamily} and the categories \texttt{IsList} (\ref{IsList}) and \texttt{IsMutable} (\ref{IsMutable}). Methods based on the low\texttt{\symbol{45}}level functions in the previous
section, are installed for the list access operations, enabling them to be
used as lists. However, it is \emph{not recommended} that these be used in programming. They are supplied mainly as a convenience
for interactive working, and may not be safe, since functions and methods for
lists may assume that after \texttt{IsBound(w[i])} returns \texttt{true}, access to \texttt{w[i]} is safe. }

  
\section{\textcolor{Chapter }{Copying Weak Pointer Objects}}\label{Copying Weak Pointer Objects}
\logpage{[ 86, 4, 0 ]}
\hyperdef{L}{X830918AC8702A189}{}
{
  A \texttt{ShallowCopy} (\ref{ShallowCopy}) method is installed, which makes a new weak pointer object containing the same
objects as the original. 

 It is possible to apply \texttt{StructuralCopy} (\ref{StructuralCopy}) to a weak pointer object, obtaining a new weak pointer object containing
copies of the objects in the original. This \emph{may not be safe} if a badly timed garbage collection occurs during copying. 

 Applying \texttt{Immutable} (\ref{Immutable}) to a weak pointer object produces an immutable plain list containing immutable
copies of the objects contained in the weak pointer object. An immutable weak
pointer object is a contradiction in terms. }

 }

           
\chapter{\textcolor{Chapter }{More about Stabilizer Chains}}\label{More about Stabilizer Chains}
\logpage{[ 87, 0, 0 ]}
\hyperdef{L}{X81F4282081027945}{}
{
  This chapter contains some rather technical complements to the material
handled in the chapters{\nobreakspace}\ref{Permutations} and \ref{Permutation Groups}.  
\section{\textcolor{Chapter }{Generalized Conjugation Technique}}\label{Generalized Conjugation Technique}
\logpage{[ 87, 1, 0 ]}
\hyperdef{L}{X870717BA831A0365}{}
{
  The command \texttt{ConjugateGroup( \mbox{\texttt{\mdseries\slshape G}}, \mbox{\texttt{\mdseries\slshape p}} )} (see{\nobreakspace}\texttt{ConjugateGroup} (\ref{ConjugateGroup})) for a permutation group \mbox{\texttt{\mdseries\slshape G}} with stabilizer chain equips its result also with a stabilizer chain, namely
with the chain of \mbox{\texttt{\mdseries\slshape G}} conjugate by \mbox{\texttt{\mdseries\slshape p}}. Conjugating a stabilizer chain by a permutation \mbox{\texttt{\mdseries\slshape p}} means replacing all the points which appear in the \texttt{orbit} components by their images under \mbox{\texttt{\mdseries\slshape p}} and replacing every permutation \mbox{\texttt{\mdseries\slshape g}} which appears in a \texttt{labels} or \texttt{transversal} component by its conjugate $g^p$. The conjugate $g^p$ acts on the mapped points exactly as \mbox{\texttt{\mdseries\slshape g}} did on the original points, i.e., $(pnt.p). g^p = (pnt.g).p$. Since the entries in the \texttt{translabels} components are integers pointing to positions of the \texttt{labels} list, the \texttt{translabels} lists just have to be permuted by \mbox{\texttt{\mdseries\slshape p}} for the conjugated stabilizer. Then \texttt{generators} is reconstructed as \texttt{labels\texttt{\symbol{123}} genlabels \texttt{\symbol{125}}} and \texttt{transversal\texttt{\symbol{123}} orbit \texttt{\symbol{125}}} as \texttt{labels\texttt{\symbol{123}} translabels\texttt{\symbol{123}} orbit
\texttt{\symbol{125}} \texttt{\symbol{125}}}. 

 \index{generalized conjugation technique} This conjugation technique can be generalized. Instead of mapping points and
permutations under the same permutation \mbox{\texttt{\mdseries\slshape p}}, it is sometimes desirable (e.g., in the context of permutation group
homomorphisms) to map the points with an arbitrary mapping $map$ and the permutations with a homomorphism $hom$ such that the compatibility of the actions is still valid: $map(pnt).hom(g) = map(pnt.g)$. (Of course the ordinary conjugation is a special case of this, with $map(pnt) = pnt.p$ and $hom(g) = g^p$.) 

 In the generalized case, the ``conjugated'' chain need not be a stabilizer chain for the image of $hom$, since the ``preimage'' of the stabilizer of $map(b)$ (where $b$ is a base point) need not fix $b$, but only fixes the preimage $map^{{-1}}( map(b) )$ setwise. Therefore the method can be applied only to one level and the next
stabilizer must be computed explicitly. But if $map$ is injective, we have $map(b).hom(g) = map(b)$ if and only if $b.g = b$, and if this holds, then $g = w(g_1, \ldots, g_n)$ is a word in the generators $g_1, \ldots, g_n$ of the stabilizer of{\nobreakspace}$b$ and  $hom(g) =^* w( hom(g_1), \ldots, hom(g_n) )$ is in the ``conjugated'' stabilizer. If, more generally, $hom$ is a right inverse to a homomorphism $\varphi$ (i.e., $\varphi(hom(g)) = g$ for all $g$), equality $*$ holds modulo the kernel of $\varphi$; in this case the ``conjugated'' chain can be made into a real stabilizer chain by extending each level with
the generators of the kernel and appending a proper stabilizer chain of the
kernel at the end. These special cases will occur in the algorithms for
permutation group homomorphisms (see{\nobreakspace}\ref{Group Homomorphisms}). 

 To ``conjugate'' the points (i.e., \texttt{orbit}) and permutations (i.e., \texttt{labels}) of the Schreier tree, a loop is set up over the \texttt{orbit} list constructed during the orbit algorithm, and for each vertex $b$ with unique edge $a(l)b$ ending at $b$, the label $l$ is mapped with $hom$ and $b$ with $map$. We assume that the \texttt{orbit} list was built w.r.t.{\nobreakspace}a certain ordering $<$ of the labels, where $l' < l$ means that every point in the orbit was mapped with $l'$ before it was mapped with $l$. This shape of the \texttt{orbit} list is guaranteed if the Schreier tree is extended only by \texttt{AddGeneratorsExtendSchreierTree} (\ref{AddGeneratorsExtendSchreierTree}), and it is then also guaranteed for the ``conjugated'' Schreier tree. (The ordering of the labels cannot be read from the Schreier
tree, however.) 

 In the generalized case, it can happen that the edge $a(l)b$ bears a label $l$ whose image is ``old'', i.e., equal to the image of an earlier label $l' < l$. Because of the compatibility of the actions we then have $map(b) = map(a).hom(l)^{{-1}} = map(a).hom(l')^{{-1}} = map(a{{l'}}^{{-1}})$, so $map(b)$ is already equal to the image of the vertex $a{{l'}}^{{-1}}$. This vertex must have been encountered before $b = al^{{-1}}$ because $l' < l$. We conclude that the image of a label can be ``old'' only if the vertex at the end of the corresponding edge has an ``old'' image, too, but then it need not be ``conjugated'' at all. A similar remark applies to labels which map under $hom$ to the identity. }

  
\section{\textcolor{Chapter }{The General Backtrack Algorithm with Ordered Partitions}}\label{The General Backtrack Algorithm with Ordered Partitions}
\logpage{[ 87, 2, 0 ]}
\hyperdef{L}{X8174E19F87C3A8AB}{}
{
  Section \ref{Backtrack} describes the basic functions for a backtrack search. The purpose of this
section is to document how the general backtrack algorithm is implemented in \textsf{GAP} and which parts you have to modify if you want to write your own backtrack
routines. 
\subsection{\textcolor{Chapter }{Internal representation of ordered partitions}}\label{Internal representation of ordered partitions}
\logpage{[ 87, 2, 1 ]}
\hyperdef{L}{X82E18F38824B5856}{}
{
  \index{ordered partitions!internal representation} \textsf{GAP} represents an ordered partition as a record with the following components. 
\begin{description}
\item[{\texttt{points}}]  a list of all points contained in the partition, such that the points of each
cell from lie consecutively, 
\item[{\texttt{cellno}}]  a list whose \mbox{\texttt{\mdseries\slshape i}}th entry is the number of the cell which contains the point \mbox{\texttt{\mdseries\slshape i}}, 
\item[{\texttt{firsts}}]  a list such that \texttt{points[firsts[\mbox{\texttt{\mdseries\slshape j}}]]} is the first point in \texttt{points} which is in cell \mbox{\texttt{\mdseries\slshape j}}, 
\item[{\texttt{lengths}}]  a list of the cell lengths. 
\end{description}
 Some of the information is redundant, e.g., the \texttt{lengths} could also be read off the \texttt{firsts} list, but since this need not be increasing, it would require some searching.
Similar for \texttt{cellno}, which could be replaced by a systematic search of \texttt{points}, keeping track of what cell is currently being traversed. With the above
components, the \mbox{\texttt{\mdseries\slshape m}}th cell of a partition \mbox{\texttt{\mdseries\slshape P}} is expressed as \texttt{\mbox{\texttt{\mdseries\slshape P}}.points\texttt{\symbol{123}} [ \mbox{\texttt{\mdseries\slshape P}}.firsts[\mbox{\texttt{\mdseries\slshape m}}] .. \mbox{\texttt{\mdseries\slshape P}}.firsts[\mbox{\texttt{\mdseries\slshape m}}] + \mbox{\texttt{\mdseries\slshape P}}.lengths[\mbox{\texttt{\mdseries\slshape m}}] \texttt{\symbol{45}} 1 ] \texttt{\symbol{125}}}. The most important operations, however, to be performed upon \mbox{\texttt{\mdseries\slshape P}} are the splitting of a cell and the reuniting of the two parts. Following the
strategy of J.{\nobreakspace}Leon,  this is done as follows: 

 
\begin{description}
\item[{(1)}]  The points which make up the cell that is to be split are sorted so that the
ones that remain inside occupy positions \texttt{[ \mbox{\texttt{\mdseries\slshape P}}.firsts[\mbox{\texttt{\mdseries\slshape m}}] .. \mbox{\texttt{\mdseries\slshape last}} ]} in the list \texttt{\mbox{\texttt{\mdseries\slshape P}}.points} (for a suitable value of \mbox{\texttt{\mdseries\slshape last}}). 
\item[{(2)}]  The points at positions \texttt{[ \mbox{\texttt{\mdseries\slshape last}} + 1 .. \mbox{\texttt{\mdseries\slshape P}}.firsts[\mbox{\texttt{\mdseries\slshape m}}] + \mbox{\texttt{\mdseries\slshape P}}.lengths[\mbox{\texttt{\mdseries\slshape m}}] \texttt{\symbol{45}} 1 ]} will form the additional cell. For this new cell requires additional entries
are added to the lists \texttt{\mbox{\texttt{\mdseries\slshape P}}.firsts} (namely, \texttt{\mbox{\texttt{\mdseries\slshape last}}+1}) and \texttt{\mbox{\texttt{\mdseries\slshape P}}.lengths} (namely, \texttt{\mbox{\texttt{\mdseries\slshape P}}.firsts[\mbox{\texttt{\mdseries\slshape m}}] + \mbox{\texttt{\mdseries\slshape P}}.lengths[\mbox{\texttt{\mdseries\slshape m}}] \texttt{\symbol{45}} \mbox{\texttt{\mdseries\slshape last}} \texttt{\symbol{45}} 1}). 
\item[{(3)}]  The entries of the sublist \texttt{\mbox{\texttt{\mdseries\slshape P}}.cellno\texttt{\symbol{123}} [ \mbox{\texttt{\mdseries\slshape last}}+1 .. \mbox{\texttt{\mdseries\slshape P}}.firsts[\mbox{\texttt{\mdseries\slshape m}}] + P.lengths[\mbox{\texttt{\mdseries\slshape m}}]\texttt{\symbol{45}}1 ] \texttt{\symbol{125}}} must be set to the number of the new cell. 
\item[{(4)}]  The entry \texttt{\mbox{\texttt{\mdseries\slshape P}}.lengths[\mbox{\texttt{\mdseries\slshape m}}]} must be reduced to \texttt{\mbox{\texttt{\mdseries\slshape last}} \texttt{\symbol{45}} \mbox{\texttt{\mdseries\slshape P}}.firsts[\mbox{\texttt{\mdseries\slshape m}}] + 1}. 
\end{description}
 

 Then reuniting the two cells requires only the reversal of
steps{\nobreakspace}2 to{\nobreakspace}4 above. The list \texttt{\mbox{\texttt{\mdseries\slshape P}}.points} need not be rearranged. }

 
\subsection{\textcolor{Chapter }{Functions for setting up an R\texttt{\symbol{45}}base}}\label{Functions for setting up an R-base}
\logpage{[ 87, 2, 2 ]}
\hyperdef{L}{X785508067969766B}{}
{
  This subsection explains some \textsf{GAP} functions which are local to the library file \texttt{lib/stbcbckt.gi} which contains the code for backtracking in permutation groups. They are
mentioned here because you might find them helpful when you want to implement
your own backtracking function based on the partition concept. An important
argument to most of the functions is the R\texttt{\symbol{45}}base $R$, which you should regard as a black box. We will tell you how to set it up,
how to maintain it and where to pass it as argument, but it is not necessary
for you to know its internal representation. However, if you insist to learn
the whole story: Here are the record components from which an
R\texttt{\symbol{45}}base is made up: 

 
\begin{description}
\item[{\texttt{domain}}]  the set $\Omega$ on which the group $G$ operates 
\item[{\texttt{base}}]  the sequence $(a_1, \ldots, a_r)$ of base points 
\item[{\texttt{partition}}]  an ordered partition, initially $\Pi_0$, this will be refined to $\Pi_1, \ldots, \Pi_r$ during the backtrack algorithm 
\item[{\texttt{where}}]  a list such that $a_i$ lies in cell number \texttt{where}$[i]$ of $\Pi_i$ 
\item[{\texttt{rfm}}]  a list whose $i$th entry is a list of refinements which take $\Sigma_i$ to $\Sigma_{{i+1}}$; the structure of a refinement is described below 
\item[{\texttt{chain}}]  a (copy of a) stabilizer chain for $G$ (not if $G$ is a symmetric group) 
\item[{\texttt{fix}}]  only if $G$ is a symmetric group: a list whose $i$ entry contains \texttt{Fixcells( }$\Pi_i$\texttt{ )} 
\item[{\texttt{level}}]  initially equal to \texttt{chain}, this will be changed to chains for the stabilizers $G_{{a_1 \ldots a_i}}$ for $i = 1, \ldots, r$ during the backtrack algorithm; if $G$ is a symmetric group, only the number of moved points is stored for each
stabilizer 
\item[{\texttt{lev}}]  a list whose $i$th entry remembers the \texttt{level} entry for $G_{{a_1 \ldots a_{{i-1}}}}$ 
\item[{\texttt{level2}, \texttt{lev2}}]  a similar construction for a second group (used in intersection calculations), \texttt{false} otherwise. This second group $H$ activated if the R\texttt{\symbol{45}}base is constructed as \texttt{EmptyRBase( }$[ G, H ], \Omega, \Pi_0$\texttt{ )} (if $G = H$, \textsf{GAP} sets \texttt{level2 = }\texttt{true} instead). 
\item[{\texttt{nextLevel}}]  this is described below 
\end{description}
 

 As our guiding example, we present code for the function \texttt{Centralizer} (\ref{Centralizer:for a magma and a submagma}) which calculates the centralizer of an element $g$ in the group $G$. (The real code is more general and has a few more subtleties.) 

 
\begin{Verbatim}[commandchars=!@|,fontsize=\small,frame=single,label=]
  Pi_0 := TrivialPartition( omega );
  R := EmptyRBase( G, omega, Pi_0 );
  R.nextLevel := function( Pi, rbase )
  local  fix, p, q, where;
  NextRBasePoint( Pi, rbase );
  fix := Fixcells( Pi );
  for p  in fix  do
    q := p ^ g;
    where := IsolatePoint( Pi, q );
    if where <> false  then
      Add( fix, q );
      ProcessFixpoint( R, q );
      AddRefinement( R, "Centralizer", [ Pi.cellno[ p ], q, where ] );
      if Pi.lengths[ where ] = 1  then
        p := FixpointCellNo( Pi, where );
        ProcessFixpoint( R, p );
        AddRefinement( R, "ProcessFixpoint", [ p, where ] );
      fi;
    fi;
  od;
  end;
  
  return PartitionBacktrack(
    G,
    c -> g ^ c = g,
    false,
    R,
    [ Pi_0, g ],
    L, R );
\end{Verbatim}
 

 The list numbers below refer to the line numbers of the code above. 

 
\begin{description}
\item[{1.}]  \texttt{omega} is the set on which \texttt{G} acts and \texttt{Pi{\textunderscore}0} is the first member of the decreasing sequence of partitions mentioned in \ref{Backtrack}. We set \texttt{Pi{\textunderscore}0 = omega}, which is constructed as \texttt{TrivialPartition( omega )}, but we could have started with a finer partition, e.g., into unions of \texttt{g}\texttt{\symbol{45}}cycles of the same length. 
\item[{2.}]  This statement sets up the R\texttt{\symbol{45}}base in the variable \texttt{R}. 
\item[{3.\texttt{\symbol{45}}21.}]  These lines define a function \texttt{R.nextLevel} which is called whenever an additional member in the sequence \texttt{Pi{\textunderscore}0 }$\geq \Pi_1 \geq \ldots$ of partitions is needed. If $\Pi_i$ does not yet contain enough base points in one\texttt{\symbol{45}}point cells, \textsf{GAP} will call \texttt{R.nextLevel( }$\Pi_i,$\texttt{ R )}, and this function will choose a new base point $a_{{i+1}}$, refine $\Pi_i$ to $\Pi_{{i+1}}$ (thereby \emph{changing} the first argument) and store all necessary information in{\nobreakspace}\texttt{R}. 
\item[{5.}]  This statement selects a new base point $a_{{i+1}}$, which is not yet in a one\texttt{\symbol{45}}point cell of $\Pi$ and still moved by the stabilizer $G_{{a_1 \ldots a_i}}$ of the earlier base points. If certain points of \texttt{omega} should be preferred as base point (e.g., because they belong to long cycles of \texttt{g}), a list of points starting with the most wanted ones, can be given as an
optional third argument to \texttt{NextRBasePoint} (actually, this is done in the real code for \texttt{Centralizer} (\ref{Centralizer:for a magma and a submagma})). 
\item[{6.}]  \texttt{Fixcells( }$\Pi$\texttt{ )} returns the list of points in one\texttt{\symbol{45}}point cells of $\Pi$ (ordered as the cells are ordered in $\Pi$). 
\item[{7.}]  For every point $p \in fix$, if we know the image $p$\texttt{\texttt{\symbol{94}}}$g$ under $c \in C_G(e)$, we also know $( p$\texttt{\texttt{\symbol{94}}}$g )$\texttt{\texttt{\symbol{94}}}$c = ( p$\texttt{\texttt{\symbol{94}}}$c )$\texttt{\texttt{\symbol{94}}}$g$. We therefore want to isolate these extra points in $\Pi$. 
\item[{9.}]  This statement puts point $q$ in a cell of its own, returning in \texttt{where} the number of the cell of $\Pi$ from which $q$ was taken. If $q$ was already the only point in its cell, \texttt{where = }\texttt{false} instead. 
\item[{12.}]  This command does the necessary bookkeeping for the extra base point $q$: It prescribes $q$ as next base in the stabilizer chain for $G$ (needed, e.g., in line{\nobreakspace}5) and returns \texttt{false} if $q$ was already fixed the stabilizer of the earlier base points (and \texttt{true} otherwise; this is not used here). Another call to \texttt{ProcessFixpoint} like this was implicitly made by the function \texttt{NextRBasePoint} to register the chosen base point. By contrast, the point $q$ was not chosen this way, so \texttt{ProcessFixpoint} must be called explicitly for{\nobreakspace}$q$. 
\item[{13.}]  This statement registers the function which will be used during the backtrack
search to perform the corresponding refinements on the ``image partition'' $\Sigma_i$ (to yield the refined $\Sigma_{{i+1}}$). After choosing an image $b_{{i+1}}$ for the base point $a_{{i+1}}$, \textsf{GAP} will compute $\Sigma_i \wedge (\{ b_{{i+1}} \}, \Omega \setminus \{ b_{{i+1}} \})$ and store this partition in $I$\texttt{.partition}, where $I$ is a black box similar to $R$, but corresponding to the current ``image partition'' (hence it is an ``R\texttt{\symbol{45}}image'' in analogy to the R\texttt{\symbol{45}}base). Then \textsf{GAP} will call the function \texttt{Refinements.Centralizer( R, I, Pi.cellno[ p ], p, where )}, with the then current values of $R$ and $I$, but where $\Pi$\texttt{.cellno}$[ p ]$, $p$, \texttt{where} still have the values they have at the time of this \texttt{AddRefinement} command. This function call will further refine $I$\texttt{.partition} to yield $\Sigma_{{i+1}}$ as it is programmed in the function \texttt{Refinements.Centralizer}, which is described below. (The global variable \texttt{Refinements} is a record which contains all refinement functions for all backtracking
procedures.) 
\item[{14.\texttt{\symbol{45}}19.}]  If the cell from which $q$ was taken out had only two points, we now have an additional
one\texttt{\symbol{45}}point cell. This condition is checked in
line{\nobreakspace}13 and if it is true, this extra fixpoint $p$ is taken (line{\nobreakspace}15), processed like $q$ before (line{\nobreakspace}16) and is then (line{\nobreakspace}17) passed to
another refinement function \texttt{Refinements.ProcessFixpoint( R, I, p, where )}, which is also described below. 
\item[{23.\texttt{\symbol{45}}29.}]  This command starts the backtrack search. Its result will be the centralizer
as a subgroup of $G$. Its arguments are 
\item[{24.}]  the group we want to run through, 
\item[{25.}]  the property we want to test, as a \textsf{GAP} function, 
\item[{26.}]  \texttt{false} if we are looking for a subgroup, \texttt{true} in the case of a representative search (when the result would be one
representative), 
\item[{27.}]  the R\texttt{\symbol{45}}base, 
\item[{28.}]  a list of data, to be stored in $I$\texttt{.data}, which has in position{\nobreakspace}1 the first member $\Sigma_0$ of the decreasing sequence of ``image partitions'' mentioned in \ref{Backtrack}. In the centralizer example, position{\nobreakspace}2 contains the element
that is to be centralized. In the case of a representative search, i.e., a
conjugacy test $g$\texttt{\texttt{\symbol{94}}}$c$\texttt{ ?= }$h$, we would have $h$ instead of $g$ here, and possibly a $\Sigma_0$ different from $\Pi_0$ (e.g., a partition into unions of $h$\texttt{\symbol{45}}cycles of same length). 
\item[{29.}]  two subgroups $L \leq C_G(g)$ and $R \leq C_G(h)$ known in advance (we have $L = R$ in the centralizer case). 
\end{description}
 }

 
\subsection{\textcolor{Chapter }{Refinement functions for the backtrack search}}\label{Refinement functions for the backtrack search}
\logpage{[ 87, 2, 3 ]}
\hyperdef{L}{X82427DA47D458224}{}
{
  The last subsection showed how the refinement process leading from $\Pi_i$ to $\Pi_{{i+1}}$ is coded in the function $R$\texttt{.nextLevel}, this has to be executed once the base point $a_{{i+1}}$. The analogous refinement step from $\Sigma_i$ to $\Sigma_{{i+1}}$ must be performed for each choice of an image $b_{{i+1}}$ for $a_{{i+1}}$, and it will depend on the corresponding value of $\Sigma_i \wedge (\{b_{{i+1}}\}, \Omega \setminus \{b_{{i+1}}\})$. But before we can continue our centralizer example, we must, for the
interested reader, document the record components of the other black box $I$, as we did above for the R\texttt{\symbol{45}}base black box $R$. Most of the components change as \textsf{GAP} walks up and down the levels of the search tree. 
\begin{description}
\item[{\texttt{data}}]  this will be mentioned below 
\item[{\texttt{depth}}]  the level $i$ in the search tree of the current node $\Sigma_i$ 
\item[{\texttt{bimg}}]  a list of images of the points in $R$\texttt{.base} 
\item[{\texttt{partition}}]  the partition $\Sigma_i$ of the current node 
\item[{\texttt{level}}]  the stabilizer chain $R$\texttt{.lev}$[i]$ at the current level 
\item[{\texttt{perm}}]  a permutation mapping \texttt{Fixcells}$( \Pi_i )$ to \texttt{Fixcells}$( \Sigma_i )$; this implies mapping $(a_1, \ldots, a_i)$ to $(b_1, \ldots, b_i)$ 
\item[{\texttt{level2}, \texttt{perm2}}]  a similar construction for the second stabilizer chain, \texttt{false} otherwise (and \texttt{true} if $R$\texttt{.level2 = }\texttt{true}) 
\end{description}
 

 As declared in the above code for \texttt{Centralizer} (\ref{Centralizer:for a magma and a submagma}), the refinement is performed by the function \texttt{Refinement.Centralizer}$( R, I, \Pi$\texttt{.cellno}$[p], p, where )$. The functions in the record \texttt{Refinement} always take two additional arguments before the ones specified in the \texttt{AddRefinement} call (in line{\nobreakspace}13 above), namely the R\texttt{\symbol{45}}base $R$ and the current value $I$ of the ``R\texttt{\symbol{45}}image''. In our example, $p$ is a fixpoint of $\Pi = \Pi_i \wedge (\{ a_{{i+1}} \}, \Omega \setminus \{ a_{{i+1}} \})$ such that $where = \Pi$\texttt{.cellno}$[ p^g ]$. The \texttt{Refinement} functions must return \texttt{false} if the refinement is unsuccessful (e.g., because it leads to $\Sigma_{{i+1}}$ having different cell sizes from $\Pi_{{i+1}}$) and \texttt{true} otherwise. Our particular function looks like this. 

 
\begin{Verbatim}[commandchars=!@|,fontsize=\small,frame=single,label=]
  Refinements.Centralizer := function( R, I, cellno, p, where )
  local  Sigma, q;
  Sigma := I.partition;
  q := FixpointCellNo( Sigma, cellno ) ^ I.data[ 2 ];
  return IsolatePoint( Sigma, q ) = where and ProcessFixpoint( I, p, q );
  end;
\end{Verbatim}
 

 The list numbers below refer to the line numbers of the code immediately
above. 

 
\begin{description}
\item[{3.}]  The current value of $\Sigma_i \wedge (\{ b_{{i+1}} \}, \Omega \setminus \{ b_{{i+1}} \})$ is always found in $I$\texttt{.partition}. 
\item[{4.}]  The image of the only point in cell number $cellno = \Pi_i$\texttt{.cellno}$[ p ]$ in $\Sigma$ under $g = I$\texttt{.data}$[ 2 ]$ is calculated. 
\item[{5.}]  The function returns \texttt{true} only if the image $q$ has the same cell number in $\Sigma$ as $p$ had in $\Pi$ (i.e., $where$) and if $q$ can be prescribed as an image for $p$ under the coset of the stabilizer $G_{{a_1 \ldots a_{{i+1}}}}.c$ where $c \in G$ is an (already constructed) element mapping the earlier base points $a_1, \ldots, a_{{i+1}}$ to the already chosen images $b_1, \ldots, b_{{i+1}}$. This latter condition is tested by \texttt{ProcessFixpoint}$( I, p, q )$ which, if successful, also does the necessary bookkeeping in $I$. In analogy to the remark about line{\nobreakspace}12 in the program above,
the chosen image $b_{{i+1}}$ for the base point $a_{{i+1}}$ has already been processed implicitly by the function \texttt{PartitionBacktrack}, and this processing includes the construction of an element $c \in G$ which maps \texttt{Fixcells}$( \Pi_i )$ to \texttt{Fixcells}$( \Sigma_i )$ and $a_{{i+1}}$ to $b_{{i+1}}$. By contrast, the extra fixpoints $p$ and $q$ in $\Pi_{{i+1}}$ and $\Sigma_{{i+1}}$ were not chosen automatically, so they require an explicit call of \texttt{ProcessFixpoint}, which replaces the element $c$ by some $c'.c$ (with $c' \in G_{{a_1 \ldots a_{{i+1}}}}$) which in addition maps $p$ to $q$, or returns \texttt{false} if this is impossible. 
\end{description}
 

 You should now be able to guess what \texttt{Refinements.ProcessFixpoint}$( R, I, p, where )$ does: it simply returns \texttt{ProcessFixpoint}$( I, p, $\texttt{FixpointCellNo}$( I$\texttt{.partition}$, where ) )$. 

 \emph{Summary.} 

 When you write your own backtrack functions using the partition technique, you
have to supply an R\texttt{\symbol{45}}base, including a component \texttt{nextLevel}, and the functions in the \texttt{Refinements} record which you need. Then you can start the backtrack by passing the
R\texttt{\symbol{45}}base and the additional data (for the \texttt{data} component of the ``R\texttt{\symbol{45}}image'') to \texttt{PartitionBacktrack}. }

 
\subsection{\textcolor{Chapter }{Functions for meeting ordered partitions}}\label{Functions for meeting ordered partitions}
\logpage{[ 87, 2, 4 ]}
\hyperdef{L}{X86CCA2B384A74856}{}
{
  A kind of refinement that occurs in particular in the normalizer calculation
involves computing the meet of $\Pi$ (cf.{\nobreakspace}lines{\nobreakspace}6ff.{\nobreakspace}above) with an
arbitrary other partition $\Lambda$, not just with one point. To do this efficiently, \textsf{GAP} uses the following two functions. 

 \texttt{StratMeetPartition( }$R$, $\Pi$, $\Lambda$ \texttt{[}, $g$ \texttt{] )} 

 \texttt{MeetPartitionStrat( }$R$, $I$\texttt{\texttt{\symbol{123}}, }$\Lambda'$\texttt{\texttt{\symbol{125}}[, \texttt{\symbol{123}}}$g'$\texttt{\texttt{\symbol{125}}]}, $strat$ \texttt{)} 

 \index{meet strategy} Such a \texttt{StratMeetPartition} command would typically appear in the function call $R$\texttt{.nextLevel}$( \Pi, R )$ (during the refinement of $\Pi_i$ to $\Pi_{{i+1}}$). This command replaces $\Pi$ by $\Pi \wedge \Lambda$ (thereby \emph{changing} the second argument) and returns a ``meet strategy'' $strat$. This is (for us) a black box which serves two purposes: First, it allows \textsf{GAP} to calculate faster the corresponding meet $\Sigma \wedge \Lambda'$, which must then appear in a \texttt{Refinements} function (during the refinement of $\Sigma_i$ to $\Sigma_{{i+1}}$). It is faster to compute $\Sigma \wedge \Lambda'$ with the ``meet strategy'' of $\Pi \wedge \Lambda$ because if the refinement of $\Sigma$ is successful at all, the intersection of a cell from the left hand side of
the $\wedge$ sign with a cell from the right hand side must have the same size in both
cases (and $strat$ records these sizes, so that only non\texttt{\symbol{45}}empty intersections
must be calculated for $\Sigma \wedge \Lambda'$). Second, if there is a discrepancy between the behaviour prescribed by $strat$ and the behaviour observed when refining $\Sigma$, the refinement can immediately be abandoned. 

 On the other hand, if you only want to meet a partition $\Pi$ with $\Lambda$ for a one\texttt{\symbol{45}}time use, without recording a strategy, you can
simply type \texttt{StratMeetPartition}$( \Pi, \Lambda )$ as in the following example, which also demonstrates some other
partition\texttt{\symbol{45}}related commands. 

                
\begin{Verbatim}[commandchars=!@|,fontsize=\small,frame=single,label=Example]
  !gapprompt@gap>| !gapinput@P := Partition( [[1,2],[3,4,5],[6]] );;  Cells( P );|
  [ [ 1, 2 ], [ 3, 4, 5 ], [ 6 ] ]
  !gapprompt@gap>| !gapinput@Q := Partition( OnTuplesTuples( last, (1,3,6) ) );;  Cells( Q );|
  [ [ 3, 2 ], [ 6, 4, 5 ], [ 1 ] ]
  !gapprompt@gap>| !gapinput@StratMeetPartition( P, Q );|
  [  ]
  !gapprompt@gap>| !gapinput@# The ``meet strategy'' was not recorded, ignore this result.|
  !gapprompt@gap>| !gapinput@Cells( P );|
  [ [ 1 ], [ 5, 4 ], [ 6 ], [ 2 ], [ 3 ] ]
\end{Verbatim}
 

 You can even say \texttt{StratMeetPartition}$( \Pi, \Delta )$ where $\Delta$ is simply a subset of $\Omega$, it will then be interpreted as the partition $(\Delta, \Omega \setminus \Delta)$. 

 \textsf{GAP} makes use of the advantages of a ``meet strategy'' if the refinement function in \texttt{Refinements} contains a \texttt{MeetPartitionStrat} command where $strat$ is the ``meet strategy'' calculated by \texttt{StratMeetPartition} before. Such a command replaces $I$\texttt{.partition} by its meet with $\Lambda'$, again changing the argument $I$. The necessary reversal of these changes when backtracking from a node (and
prescribing the next possible image for a base point) is automatically done by
the function \texttt{PartitionBacktrack}. 

 In all cases, an additional argument $g$ means that the meet is to be taken not with $\Lambda$, but instead with $\Lambda.{{g^{{-1}}}}$, where operation on ordered partitions is meant cellwise (and setwise on each
cell). (Analogously for the primed arguments.) 

 
\begin{Verbatim}[commandchars=!@|,fontsize=\small,frame=single,label=Example]
  !gapprompt@gap>| !gapinput@P := Partition( [[1,2],[3,4,5],[6]] );;|
  !gapprompt@gap>| !gapinput@StratMeetPartition( P, P, (1,6,3) );;  Cells( P );|
  [ [ 1 ], [ 5, 4 ], [ 6 ], [ 2 ], [ 3 ] ]
\end{Verbatim}
 Note that $P.(1,3,6) = Q$. }

 }

  
\section{\textcolor{Chapter }{Stabilizer Chains for Automorphisms Acting on Enumerators}}\label{Stabilizer Chains for Automorphisms Acting on Enumerators}
\logpage{[ 87, 3, 0 ]}
\hyperdef{L}{X7CA84E967B053C2C}{}
{
  This section describes a way of representing the automorphism group of a group
as permutation group, following \cite{Sims97}. The code however is not yet included in the \textsf{GAP} library. 

 In this section we present an example in which objects we already know
(namely, automorphisms of solvable groups) are equipped with the
permutation\texttt{\symbol{45}}like operations \texttt{\texttt{\symbol{94}}} and \texttt{/} for action on positive integers. To achieve this, we must define a new type of
objects which behave like permutations but are represented as automorphisms
acting on an enumerator. Our goal is to generalize the
Schreier\texttt{\symbol{45}}Sims algorithm for construction of a stabilizer
chain to groups of such new automorphisms. 
\subsection{\textcolor{Chapter }{An operation domain for automorphisms}}\label{An operation domain for automorphisms}
\logpage{[ 87, 3, 1 ]}
\hyperdef{L}{X864007907EA923FB}{}
{
  The idea we describe here is due to C.{\nobreakspace}Sims. We consider a group $A$ of automorphisms of a group $G$, given by generators, and we would like to know its order. Of course we could
follow the strategy of the Schreier\texttt{\symbol{45}}Sims algorithm
(described in \ref{Stabilizer Chains}) for $A$ acting on $G$. This would involve a call of \texttt{StabChainStrong( EmptyStabChain( [], One( \mbox{\texttt{\mdseries\slshape A}} ) ), GroupGenerators( \mbox{\texttt{\mdseries\slshape A}} ) )} where \texttt{StabChainStrong} is a function as the one described in the pseudo\texttt{\symbol{45}}code
below: 

 
\begin{Verbatim}[commandchars=!@|,fontsize=\small,frame=single,label=]
  StabChainStrong := function( S, newgens )
    Extend the Schreier tree of S with newgens.
    for sch  in  Schreier generators  do
      if not sch in S.stabilizer  then
        StabChainStrong( S.stabilizer, [ sch ] );
      fi;
    od;
  end;
\end{Verbatim}
 

 The membership test $sch \notin S$\texttt{.stabilizer} can be performed because the stabilizer chain of $S$\texttt{.stabilizer} is already correct at that moment. We even know a base in advance, namely any
generating set for $G$. Fix such a generating set $(g_1, \ldots, g_d)$ and observe that this base is generally very short compared to the degree $|G|$ of the operation. The problem with the Schreier\texttt{\symbol{45}}Sims
algorithm, however, is then that the length of the first basic orbit $g_1.A$ would already have the magnitude of $|G|$, and the basic orbits at deeper levels would not be much shorter. For the
advantage of a short base we pay the high price of long basic orbits, since
the product of the (few) basic orbit lengths must equal $|A|$. Such long orbits make the Schreier\texttt{\symbol{45}}Sims algorithm
infeasible, so we have to look for a longer base with shorter basic orbits. 

 Assume that $G$ is solvable and choose a characteristic series with elementary abelian
factors. For the sake of simplicity we assume that $N < G$ is an elementary abelian characteristic subgroup with elementary abelian
factor group $G/N$. Since $N$ is characteristic, $A$ also acts as a group of automorphisms on the factor group $G/N$, but of course not necessarily faithfully. To retain a faithful action, we
let $A$ act on the disjoint union $G/N$ with $G$, and choose as base $(g_1 N, \ldots, g_d N, g_1, \ldots, g_d)$. Now the first $d$ basic orbits lie inside $G/N$ and can have length at most $[G:N]$. Since the base points $g_1 N, \ldots, g_d N$ form a generating set for $G/N$, their iterated stabilizer $A^{(d+1)}$ acts trivially on the factor group $G/N$, i.e., it leaves the cosets $g_i N$ invariant. Accordingly, the next $d$ basic orbits lie inside $g_i N$ (for $i = 1, \ldots, d$) and can have length at most{\nobreakspace}$|N|$. 

 Generalizing this method to a characteristic series $G = N_0 > N_1 > \ldots > N_l = \{ 1 \}$ of length $l > 2$, we can always find a base of length $l.d$ such that each basic orbit is contained in a coset of a characteristic factor,
i.e. in a set of the form $g_i N_{{j-1}} / N_j$ (where $g_i$ is one of the generators of $G$ and $1 \leq j \leq l$). In particular, the length of the basic orbits is bounded by the size of the
corresponding characteristic factors. To implement a
Schreier\texttt{\symbol{45}}Sims algorithm for such a base, we must be able to
let automorphisms act on cosets of characteristic factors $g_i N_{{j-1}} / N_j$, for varying $i$ and $j$. We would like to translate each such action into an action on $\{ 1, \ldots, [ N_{{j-1}}:N_j] \}$, because then we need not enumerate the operation domain, which is the
disjoint union of $G / N_1$, $G / N_2 \ldots G / N_l$, as a whole. Enumerating it as a whole would result in basic orbits like \texttt{orbit}$ \subseteq \{ 1001, \ldots, 1100 \}$ with a \texttt{transversal} list whose first 1000 entries would be unbound, but still require
4{\nobreakspace}bytes of memory each (see{\nobreakspace}\ref{Stabilizer Chain Records}). 

 Identifying each coset $g_i N_{{j-1}} / N_j$ into $\{ 1, \ldots, [N_{{j-1}}:N_j] \}$ of course means that we have to change the action of the automorphisms on
every level of the stabilizer chain. Such flexibility is not possible with
permutations because their effect on positive integers is ``hardwired'' into them, but we can install new operations for automorphisms. }

 
\subsection{\textcolor{Chapter }{Enumerators for cosets of characteristic factors}}\label{Enumerators for cosets of characteristic factors}
\logpage{[ 87, 3, 2 ]}
\hyperdef{L}{X84A94914876C03F0}{}
{
  So far we have not used the fact that the characteristic factors are
elementary abelian, but we will do so from here on. Our first task is to
implement an enumerator (see \texttt{AsList} (\ref{AsList}) and \ref{Enumerators}) for a coset of a characteristic factor in a solvable group $G$. We assume that such a coset $g N/M$ is given by 

 
\begin{description}
\item[{(1)}]  a pcgs for the group $G$ (see \texttt{Pcgs} (\ref{Pcgs})), let $n = $\texttt{Length( }$pcgs$\texttt{ )}; 
\item[{(2)}]  a range $range = [ start .. stop ]$ indicating that $N = \langle pcgs\{ [ start .. n ] \} \rangle$ and $M = \langle pcgs\{ [ stop + 1 .. n ] \} \rangle$, i.e., the cosets of $pcgs\{ range \}$ form a base for the vector space $N/M$; 
\item[{(3)}]  the representative $g$. 
\end{description}
 

 We first define a new representation for such enumerators and then construct
them by simply putting these three pieces of data into a record object. The
enumerator should behave as a list of group elements (representing cosets
modulo $M$), consequently, its family will be the family of the $pcgs$ itself. 

 
\begin{Verbatim}[commandchars=!@|,fontsize=\small,frame=single,label=Example]
  DeclareRepresentation( "IsCosetSolvableFactorEnumeratorRep", IsEnumerator,
      [ "pcgs", "range", "representative" ] );
  
  EnumeratorCosetSolvableFactor := function( pcgs, range, g )
      return Objectify( NewType( FamilyObj( pcgs ),
                     IsCosetSolvableFactorEnumeratorRep ),
                     rec( pcgs := pcgs,
                         range := range,
                representative := g ) );
  end;
\end{Verbatim}
 

 The definition of the operations \texttt{Length} (\ref{Length}), \texttt{\texttt{\symbol{92}}[\texttt{\symbol{92}}]} (\ref{bSlash[bSlash]}) and \texttt{Position} (\ref{Position}) is now straightforward. The code has sometimes been abbreviated and is meant ``cum grano salis'', e.g., the declaration of the local variables has been left out. 

 
\begin{Verbatim}[commandchars=@|A,fontsize=\small,frame=single,label=Example]
  InstallMethod( Length, [ IsCosetSolvableFactorEnumeratorRep ],
      enum -> Product( RelativeOrdersPcgs( enum!.pcgs ){ enum!.range } ) );
  
  InstallMethod( \[\], [ IsCosetSolvableFactorEnumeratorRep,
          IsPosRat and IsInt ],
      function( enum, pos )
      elm := ();
      pos := pos - 1;
      for i  in Reversed( enum!.range )  do
          p := RelativeOrderOfPcElement( enum!.pcgs, i );
          elm := enum!.pcgs[ i ] ^ ( pos mod p ) * elm;
          pos := QuoInt( pos, p );
      od;
      return enum!.representative * elm;
  end );
  
  InstallMethod( Position, [ IsCosetSolvableFactorEnumeratorRep,
          IsObject, IsZeroCyc ],
      function( enum, elm, zero )
      exp := ExponentsOfPcElement( enum!.pcgs,
                     LeftQuotient( enum!.representative, elm ) );
      pos := 0;
      for i  in enum!.range  do
          pos := pos * RelativeOrderOfPcElement( pcgs, i ) + exp[ i ];
      od;
      return pos + 1;
  end );
\end{Verbatim}
 }

 
\subsection{\textcolor{Chapter }{Making automorphisms act on such enumerators}}\label{Making automorphisms act on such enumerators}
\logpage{[ 87, 3, 3 ]}
\hyperdef{L}{X79B146E9786FE153}{}
{
  Our next task is to make automorphisms of the solvable group $pcgs$\texttt{!.group} act on $[ 1 .. $\texttt{Length}$( enum ) ]$ for such an enumerator $enum$. We achieve this by introducing a new representation of automorphisms on
enumerators and by putting the enumerator together with the automorphism into
an object which behaves like a permutation. Turning an ordinary automorphism
into such a special automorphism requires then the construction of a new
object which has the new type. We provide an operation \texttt{PermOnEnumerator( \mbox{\texttt{\mdseries\slshape model}}, \mbox{\texttt{\mdseries\slshape aut}} )} which constructs such a new object having the same type as \mbox{\texttt{\mdseries\slshape model}}, but representing the automorphism \mbox{\texttt{\mdseries\slshape aut}}. So \mbox{\texttt{\mdseries\slshape aut}} can be either an ordinary automorphism or one which already has an enumerator
in its type, but perhaps different from the one we want (i.e. from the one in \mbox{\texttt{\mdseries\slshape model}}). 

 
\begin{Verbatim}[commandchars=@|B,fontsize=\small,frame=single,label=Example]
  DeclareCategory( "IsPermOnEnumerator",
      IsMultiplicativeElementWithInverse and IsPerm );
  
  DeclareRepresentation( "IsPermOnEnumeratorDefaultRep",
      IsPermOnEnumerator and IsAttributeStoringRep,
      [ "perm" ] );
  
  DeclareOperation( "PermOnEnumerator",
      [ IsEnumerator, IsObject ] );
  
  InstallMethod( PermOnEnumerator,
      [ IsEnumerator, IsObject ],
      function( enum, a )
      SetFilterObj( a, IsMultiplicativeElementWithInverse );
      a := Objectify( NewKind( PermutationsOnEnumeratorsFamily,
                   IsPermOnEnumeratorDefaultRep ),
                   rec( perm := a ) );
      SetEnumerator( a, enum );
      return a;
  end );
  
  InstallMethod( PermOnEnumerator,
      [ IsEnumerator, IsPermOnEnumeratorDefaultRep ],
      function( enum, a )
      a := Objectify( TypeObj( a ), rec( perm := a!.perm ) );
      SetEnumerator( a, enum );
      return a;
  end );
\end{Verbatim}
 

 Next we have to install new methods for the operations which calculate the
product of two automorphisms, because this product must again have the right
type. We also have to write a function which uses the enumerators to apply
such an automorphism to positive integers. 

 
\begin{Verbatim}[commandchars=@|A,fontsize=\small,frame=single,label=Example]
  InstallMethod( \*, IsIdenticalObj,
      [ IsPermOnEnumeratorDefaultRep, IsPermOnEnumeratorDefaultRep ],
      function( a, b )
      perm := a!.perm * b!.perm;
      SetIsBijective( perm, true );
      return PermOnEnumerator( Enumerator( a ), perm );
  end );
  
  InstallMethod( \^,
      [ IsPosRat and IsInt, IsPermOnEnumeratorDefaultRep ],
      function( p, a )
      return PositionCanonical( Enumerator( a ),
                     Enumerator( a )[ p ] ^ a!.perm );
  end );
\end{Verbatim}
 

 How the corresponding methods for \texttt{\mbox{\texttt{\mdseries\slshape p}} / \mbox{\texttt{\mdseries\slshape aut}}} and \texttt{\mbox{\texttt{\mdseries\slshape aut}} \texttt{\symbol{94}} \mbox{\texttt{\mdseries\slshape n}}} look like is obvious. 

 Now we can formulate the recursive procedure \texttt{StabChainStrong} which extends the stabilizer chain by adding in new generators $newgens$. We content ourselves again with pseudo\texttt{\symbol{45}}code, emphasizing
only the lines which set the \texttt{EnumeratorDomainPermutation}. We assume that initially $S$ is a stabilizer chain for the trivial subgroup with a level for each pair $(range,g)$ characterizing an enumerator (as described above). We also assume that the \texttt{identity} element at each level already has the type corresponding to that level. 

 
\begin{Verbatim}[commandchars=!@|,fontsize=\small,frame=single,label=]
  StabChainStrong := function( S, newgens )
    for i  in [ 1 .. Length( newgens ) ]  do
      newgens[ i ] := AutomorphismOnEnumerator( S.identity, newgens[ i ] );
    od;
    Extend the Schreier tree of S with newgens.
    for sch  in  Schreier generators  do
      if not sch in S.stabilizer  then
        StabChainStrong( S.stabilizer, [ sch ] );
      fi;
    od;
  end;
\end{Verbatim}
 }

 }

 }

    \def\bibname{References\logpage{[ "Bib", 0, 0 ]}
\hyperdef{L}{X7A6F98FD85F02BFE}{}
}

\bibliographystyle{alpha}
\bibliography{../manualbib.xml}

\addcontentsline{toc}{chapter}{References}

\def\indexname{Index\logpage{[ "Ind", 0, 0 ]}
\hyperdef{L}{X83A0356F839C696F}{}
}

\cleardoublepage
\phantomsection
\addcontentsline{toc}{chapter}{Index}


\printindex

\newpage
\immediate\write\pagenrlog{["End"], \arabic{page}];}
\immediate\closeout\pagenrlog
\end{document}
